\documentclass[11pt]{amsart}
\usepackage[margin=1.1in]{geometry}
\usepackage{amsmath,amssymb,mathrsfs}
\allowdisplaybreaks
\newcommand{\namedhead}{}
\theoremstyle{definition}
\newtheorem*{namedinner}{\namedhead}
\newenvironment{named}[1]{\renewcommand{\namedhead}{#1}\begin{namedinner}}{\end{namedinner}}
\renewcommand{\descriptionlabel}[1]{\hspace\labelsep\upshape\bfseries #1}
\pagestyle{plain}
\begin{document}
\title[Lattice QED in three dimensions: the chamber fermion]{Noncompact lattice QED in three dimensions: the chamber fermion, the renormalisation group analysis and the fourth cumulant}
\author{}
\date{}

\begin{abstract}
The block spin renormalisation group analysis that the first paper of this series makes for noncompact lattice QED on three-dimensional tori with Wilson fermions is extended to the chamber fermion of the second paper, which has one site in each chamber of the mirrors of $\mathbb{Z}^3$. The laws are those of the second paper: laws of the curvature of the noncompact photon field, all flux sectors summed, on antiperiodic cubic tori with the spacing $\eta=L^{-N}$, with the chamber determinant for the Wilson parameter $\rho\in(0,1)$ and the standard weights. There are bare mass shifts $m_{0,N}(e)$, $\left|m_{0,N}(e)\right|<\frac{1}{8}$, independent of the volume, with the following property. For every probe $\beta$ whose test form $f=dd^{*}\beta$ has $b(f)>0$, $b(f)/24$ being the quartic term of the free Dirac $\log\det$ in $d^{*}\beta$, there are $e_2>0$ and $C<\infty$ such that at every charge $e\in(0,e_2]$ the fourth cumulant of the curvature field smeared with $f$ differs from $e^4\cdot b(f)$ by at most $C\cdot\left(e^{4+\sigma}+\eta^{\frac{1}{4}}\right)$ at all cutoffs beyond thresholds that depend on $e$, and by at most $C\cdot e^{4+\sigma}$ in every continuum and infinite-volume limit point; the exponent $\sigma>0$ depends neither on $e$ nor on $\rho$ nor on the probe. Hence no limit point is Gaussian or a mixture of translation-invariant Gaussian laws with Gaussian domination. Moreover, by a bound on the expectations of products of translated Wick powers, in almost every ergodic component of every limit point the fourth cumulant of the curvature field smeared with $f$ differs from $e^4\cdot b(f)$ by at most $C\cdot e^{4+\sigma}$ as well, so that almost every ergodic component is non-Gaussian. With the second paper it follows that the limit points in which the volume tends to infinity before the spacing tends to zero are non-Gaussian laws, invariant under all Euclidean motions and reflection positive at every plane, whose Schwinger functions satisfy the Osterwalder--Schrader axioms (E0$'$)--(E3) of linear growth, Euclidean invariance, positivity and symmetry, in the form that the second paper gives them for a field of two-forms on $\mathbb{R}^3$; almost every ergodic component satisfies in addition the cluster property (E4).
\end{abstract}
\maketitle

\section*{Introduction}

This paper is the third of a series on noncompact lattice quantum electrodynamics in three dimensions, with the fermions integrated out. In the first paper \cite{QEDI} the fermion is the Wilson fermion, and two laws of the curvature field are studied, the harmonic-averaged law and the all-flux law, which sums the flux sectors. For every charge and every choice of the bare masses for which the lattice masses exceed $-\frac{1}{2}$, these laws are tight, and their limit points satisfy Gaussian domination, are invariant under the translations and the cubic group and are reflection positive at the coordinate planes. A block spin renormalisation group analysis at small charge shows that there are bare masses with the following property: if $f=dd^{*}\beta$, $\beta$ being a real smooth two-form of compact support, is a test form with $b(f)>0$, where $b(f)$ is a coefficient of the free Dirac field, then in every limit point the fourth cumulant of the curvature field smeared with $f$ is $e^4\cdot b(f)+O\left(e^{4+\sigma}\right)$, $e$ being the charge and $\sigma$ a positive exponent. As test forms with $b(f)>0$ exist \cite[Theorem~21.1(a)]{QEDI}, the limit points formed with these bare masses at small charge are not Gaussian \cite[Theorem~3.1]{QEDI}. In the second paper \cite{QEDII} the fermion is the chamber fermion, which has one site in each chamber of the mirrors of $\mathbb{Z}^3$. Its weight is reflection positive at the mirrors of all nine families, on the tori adapted to them \cite[Theorem~1.1]{QEDII}; the infinite-volume limits on the lattice are reflection positive at all these mirrors \cite[Theorem~1.2]{QEDII}; and the continuum limit points in which the volume tends to infinity before the spacing tends to zero are invariant under all Euclidean motions of $\mathbb{R}^3$ and reflection positive at every plane \cite[Theorems~1.3 and~1.4]{QEDII}. These results hold for every charge and every choice of the lattice masses, and they contain no lower bound: the point mass at the zero field has all the properties established there \cite[Section~1.3]{QEDII}.

The present paper carries the renormalisation group analysis of \cite{QEDI} over to the chamber fermion. Let $f=dd^{*}\beta$ be the test form of a probe $\beta$ with $b(f)>0$ (both are defined below). At small charge and for suitable bare masses, the fourth cumulant of the curvature field smeared with $f$ is shown to differ from $e^4\cdot b(f)$ by at most $C\cdot\left(e^{4+\sigma}+\eta^{\frac{1}{4}}\right)$ at every cutoff $(N,T)$ with $N\ge N_0(e)$ and $T\ge T^{*}(e)+1$, where $\eta$ is the spacing, $N_0(e)$ and $T^{*}(e)$ are thresholds that depend on the charge, and $C$ is a constant independent of the cutoffs; in every limit point it differs from $e^4\cdot b(f)$ by at most $C\cdot e^{4+\sigma}$. Consequently the limit points of \cite{QEDII} formed with these bare masses are not Gaussian, and the limit points in which the volume tends to infinity first are non-Gaussian laws that are invariant under the Euclidean group and reflection positive at every plane.

\subsection*{The lattice theory}

Let $L\ge5$ be an odd integer, the scale factor of the block averaging, which is fixed in Section~I.8, and let $\bar{m}\in(0,1]$ be the mass parameter. Put $h:=L^{\left\lfloor\log_L\bar{m}^{-1}\right\rfloor}$; then $\frac{1}{L}<\bar{m}\cdot h\le1$. A cutoff is a pair $(N,T)$ of integers $N,T\ge0$, and at the cutoff $(N,T)$
\begin{gather*}
\eta:=L^{-N},\qquad n:=8\cdot h\cdot L^{T+N},\qquad R:=n\cdot\eta=8\cdot h\cdot L^{T},\\
\Lambda_{\eta,R}:=\eta\cdot\mathbb{Z}^3/R\cdot\mathbb{Z}^3;
\end{gather*}
$\eta$ is the spacing, and $R$ is the side of the torus $\Lambda_{\eta,R}$, which has $n$ sites per side. Let $e>0$ be the charge, $q:=e\cdot\eta^{\frac{1}{2}}$ the lattice charge and $r_q:=\frac{2\pi}{q}$.

The photon field lives on the links of the torus $\mathbb{Z}^3/n\cdot\mathbb{Z}^3$, which is identified with $\Lambda_{\eta,R}$ through $x\mapsto\eta\cdot x$. A configuration is a class of pairs $(A,\mathfrak{z})$ of a real one-cochain $A$ and a closed integer two-cochain $\mathfrak{z}$, modulo $(A,\mathfrak{z})\sim\left(A+r_q\cdot k,\mathfrak{z}-dk\right)$ for integer one-cochains $k$. Its curvature is the closed two-cochain $F^{\mathrm{u}}:=dA+r_q\cdot\mathfrak{z}$, and the photon weight is the Maxwell weight $\exp\left(-\frac{1}{2}\left\|F^{\mathrm{u}}\right\|^2\right)$. The field $A$ is real valued, that is, the photon field is noncompact, and the sum over $\mathfrak{z}$ includes every flux sector \cite[Definition~2.6]{QEDII}.

The fermion is the chamber fermion of \cite{QEDII}. The mirrors of $\mathbb{Z}^3$ are the planes $\{x_\mu=\mathrm{po}\}$, $\mathrm{po}\in\frac{1}{2}\mathbb{Z}$, and $\{x_\mu\pm x_\nu=\mathrm{po}\}$, $\mathrm{po}\in\mathbb{Z}$, $\mu<\nu$; they fall into three axis families and six diagonal families. The group generated by the reflections in the mirrors is the group $W_{\mathrm{af}}=\mathbb{Z}^3\rtimes O_{\mathrm{cub}}$ of the isometries of $\mathbb{R}^3$ that map $\mathbb{Z}^3$ onto itself, $O_{\mathrm{cub}}$ being the cubic group of order $48$. The complement of the union of the mirrors consists of open simplices, the chambers, which $W_{\mathrm{af}}$ permutes simply transitively. The fermion sites are the points of the orbit under $W_{\mathrm{af}}$ of a point $v_0$ of the chamber $\mathrm{Chm}_0:=\left\{0<x_3<x_2<x_1<\frac{1}{2}\right\}$: there is one site in each chamber, there are $48$ sites in each cube of side one, and no site lies on a mirror. A site $x$ is joined by bonds to its mirror images $x^{(0)},\dots,x^{(3)}$ in the four walls of its chamber, and $\pi(x)\in\mathbb{Z}^3$ is the photon site nearest to $x$; one has $\pi\left(x^{(i)}\right)=\pi(x)$ for $i=0,1,2$, while $\pi\left(x^{(3)}\right)$ is a nearest neighbour of $\pi(x)$. Let $\gamma_\mu:=\operatorname{diag}\left(\sigma_\mu,-\sigma_\mu\right)$, $\mu=1,2,3$, with the Pauli matrices $\sigma_\mu$, and $\gamma\cdot v:=\sum_\mu v_\mu\cdot\gamma_\mu$ for $v\in\mathbb{R}^3$; let $\rho\in(0,1)$ be the Wilson parameter, let
\[
w=\left(w_0,w_1,w_2,w_3\right):=\left(6,6\sqrt{2},6\sqrt{2},6\right),\qquad\mathrm{sw}:=w_0+w_1+w_2+w_3=12+12\sqrt{2},
\]
be the standard weights, forty-eight times the areas of the walls of $\mathrm{Chm}_0$, and let $\boldsymbol{\mu}\in\mathbb{R}$ be the lattice mass. On functions $\psi$ on the sites with values in $\mathbb{C}^4$ the chamber operator acts by
\begin{multline*}
\left(\left(D_A+\boldsymbol{\mu}\right)\psi\right)(x)=M_{\mathrm{ch}}\cdot\psi(x)\\
-\sum_{i=0}^{3}\frac{w_i}{2}\cdot\left(\rho-\gamma\cdot\frac{x^{(i)}-x}{\left|x^{(i)}-x\right|_2}\right)\cdot\exp\left(iq\cdot A\left(\pi(x),\pi\left(x^{(i)}\right)\right)\right)\cdot\psi\left(x^{(i)}\right),
\end{multline*}
where $M_{\mathrm{ch}}:=\boldsymbol{\mu}+\frac{\rho}{2}\cdot\mathrm{sw}$, $A(c,c')=-A(c',c)$ is the value of $A$ on the link from $c$ to $c'$, and $A(c,c):=0$ \cite[Definition~2.10]{QEDII}. At every mirror the bonds that cross it are orthogonal to it, and the part of the action that couples the two sides is a sum of terms of positive type with the coefficients $\frac{w_i}{2}\cdot(1\pm\rho)\ge0$ \cite[Proposition~3.20]{QEDII}; this form of the crossing part gives the reflection positivity at the mirrors of all nine families. The sites of the Wilson fermion of \cite{QEDI}, by contrast, lie on the diagonal mirrors.

On the antiperiodic cubic torus $C(n)$ of \cite[Definition~2.11]{QEDII}, whose fermion fields change sign under the translation by $n$ along each coordinate axis, the determinant $\det\left(D_A+\boldsymbol{\mu}\right)$ is a non-negative function of the real one-cochain $A$ \cite[Proposition~3.18]{QEDII}, and it does not depend on the point $v_0$ (Proposition~46). For a bounded Borel function $\Phi$ of the curvature put
\[
Z_n[\Phi]:=\int\mathrm{d}(A,\mathfrak{z})\,\exp\left(-\frac{1}{2}\left\|F^{\mathrm{u}}\right\|^2\right)\cdot\det\left(D_A+\boldsymbol{\mu}\right)\cdot\Phi\left(F^{\mathrm{u}}\right),\qquad E^{\mathrm{af}}_n(\Phi):=\frac{Z_n[\Phi]}{Z_n[1]},
\]
where $\int\mathrm{d}(A,\mathfrak{z})$ sums over the closed $\mathfrak{z}$ and integrates the value of $A$ at each link over $[0,r_q)$; one has $0<Z_n[1]<\infty$ for every real lattice mass \cite[Theorem~4.26]{QEDII}. The probability measure $E^{\mathrm{af}}_n$ is the all-flux law of the curvature $F^{\mathrm{u}}$ on $C(n)$.

A mass assignment is a sequence $N\mapsto m_{0,N}$ of real numbers, the bare mass shifts; at the cutoff $(N,T)$ the lattice mass is $\boldsymbol{\mu}_N:=\bar{m}\cdot\eta+m_{0,N}$, which does not depend on $T$. The physical curvature is $F:=\eta^{-\frac{3}{2}}\cdot F^{\mathrm{u}}$, a two-cochain on $\Lambda_{\eta,R}$. Its step two-form $F^{\sharp}\in\mathcal{S}'\left(\mathbb{R}^3;\Lambda^2\right)$ has, on the cube of side $\eta$ centred at the centre of the plaquette $(x;\mu\nu)$ and on the translates of this cube by $R\cdot\mathbb{Z}^3$, the $\mu\nu$-coefficient $F(x;\mu\nu)$; for a test form $g\in\mathcal{S}\left(\mathbb{R}^3;\Lambda^2\right)$ the number $F^{\sharp}(g)=X_{\eta,R}(g)$ is the curvature field smeared with $g$. The law $P_{N,T}$ of $F^{\sharp}$ under $E^{\mathrm{af}}_n$, formed with $q=e\cdot\eta^{\frac{1}{2}}$ and $\boldsymbol{\mu}=\boldsymbol{\mu}_N$, is the law $P^{\mathrm{ch}}_{N,T}$ of \cite[Definition~2.36]{QEDII}; it is a Borel probability measure on the Hermite--Sobolev space $H_{-4}$ of \cite[Sections~2.4 and~7.1.5]{QEDI}. A limit point of the family $\left(P_{N,T}\right)$ is a weak limit on $H_{-4}$ along a sequence of cutoffs $\left(N_i,T_i\right)$ with $N_i\to\infty$ and $T_i\to\infty$ in any manner. For fixed $N$ the weak limits on $H_{-4}$ of $P_{N,T_j}$ along sequences $T_j\to\infty$ form the set $\mathcal{V}_N$ of the volume limits at the spacing $L^{-N}$, and a volume-first limit point is a weak limit on $H_{-4}$ of a sequence $\mu^{(N_i)}\in\mathcal{V}_{N_i}$ with $N_i\to\infty$ \cite[Definition~2.37]{QEDII}.

For every charge and every mass assignment one has $\int\mathrm{d}P_{N,T}(\omega)\,\|\omega\|_{-3}^2\le61$ at every cutoff, $\omega$ denoting the variable in $\mathcal{S}'\left(\mathbb{R}^3;\Lambda^2\right)$; the family $\left(P_{N,T}\right)$ is tight on $H_{-4}$, and every limit point $\mu$ satisfies Gaussian domination with the constant one,
\[
\int\mathrm{d}\mu(\omega)\,\exp\omega(g)\le\exp\left(\frac{1}{2}\|g\|_2^2\right)\qquad\text{for every real test form }g,
\]
is carried by the closed two-forms, is invariant under the translations of $\mathbb{R}^3$, under $O_{\mathrm{cub}}$ and under $\omega\mapsto-\omega$, and is reflection positive at every coordinate plane; volume-first limit points exist, and each of them is invariant under all Euclidean motions of $\mathbb{R}^3$ and reflection positive at every plane \cite[Theorems~1.3 and~1.4]{QEDII}. These properties are symmetries, positivity properties and bounds from above, and the point mass at the zero field has every one of them \cite[Section~1.3 and Example~12.43]{QEDII}.

\subsection*{The main result}

The constants of the analysis are fixed in Section~I.8, one after the other, in a prescribed order: the Wilson parameter $\rho$ and the weights; the H\"older order $\alpha$; the decay rate $\kappa$ of the polymer activities; the scale factor $L$; the side $M$ of the cubes of which the polymers are composed, a power of $L$; the exponents $r$, $p_0$, $p$ of the logarithmic thresholds, a small number $\varepsilon$ and a number $s\in(0,1)$ such that the exponent $\sigma$ below is positive; and the ceilings $e_0$ and $e_1$ of the charge; together with numbers that are functions of the constants fixed before them. Each constant is subject to conditions that contain, besides it, only constants fixed before it, and choices that satisfy all the conditions exist (Propositions~14.2 and~14.3). The exponent
\[
\sigma:=s\cdot\left(\frac{1}{4}-9\varepsilon\right)-12\varepsilon-\frac{2a_L}{N_{*}},\qquad a_L:=\frac{L^3-1}{(L-1)^3},\qquad N_{*}:=8L^3,
\]
satisfies $0<\sigma\le\frac{1}{4}$ (Proposition~14.2\,(f)). It is formed from $L$, $\varepsilon$ and $s$ alone; it depends neither on $\rho$ nor on $\bar{m}$ (Proposition~14.4\,(a), (b)), and, being fixed before the probe and the charge are chosen, it depends on neither of these (Proposition~14.2\,(a), (e)).

A probe is a real two-form $\beta\ne0$ in $C_c^{\infty}\left(\mathbb{R}^3;\Lambda^2\right)$; to it belong the one-form $a:=d^{*}\beta$, the test form $f:=dd^{*}\beta$ and the free coefficient
\[
b(f):=-6\operatorname{Tr}_{L^2(\mathbb{R}^3;\mathbb{C}^4)}\left[\left(S_{\bar{m}}\,i\gamma\cdot a\right)^4\right],\qquad S_{\bar{m}}:=\left(\bar{m}+\gamma\cdot\partial\right)^{-1},
\]
of \cite[Section~2.6]{QEDI}. It is a real number that depends on $f$ alone; it is the one-loop light-by-light coefficient of the free Dirac field of mass $\bar{m}$ in the external potential $a$; with $K:=S_{\bar{m}}\,i\gamma\cdot a$, the number $\frac{1}{24}b(f)=-\frac{1}{4}\operatorname{Tr}K^4$ is the coefficient of $t^4$ in the formal series $\sum_{n\ge1}\frac{(-1)^{n+1}}{n}\cdot t^n\cdot\operatorname{Tr}K^n$ of $\log\det\left(1+t\cdot K\right)$, of which only the term $n=4$ is used, $K^4$ being of trace class. For every $\bar{m}$ there are probes with $b(f)>0$ \cite[Theorem~21.1(a)]{QEDI}. Let $\beta$ be a probe with $b(f)>0$, let $r_\beta$ be the radius of a ball that contains the support of $\beta$, put $\|a\|_{\infty}:=\sup_y|a(y)|_2$, and let $x_0\in\mathbb{R}^3$ be a point with $d_0:=|x_0|_2-2r_\beta>0$ and
\[
\frac{6}{\bar{m}^2}\cdot\|a\|_{\infty}^4\cdot\left(\frac{4\pi\cdot r_{\beta}^3}{3}\right)^2\cdot\left(\bar{m}^2+\left(\bar{m}+\frac{1}{d_0}\right)^2\right)\cdot\frac{\exp\left(-2\bar{m}\cdot d_0\right)}{4\pi^2\cdot d_0^2}\le\frac{b(f)}{12},
\]
the condition of \cite[Section~2.8]{QEDI}; it holds as soon as $|x_0|_2$ is large enough. With the translate $\left(\tau_{x_0}f\right)(y):=f(y-x_0)$ put
\[
g_0:=f,\qquad g_{\pm}:=f\pm\tau_{x_0}f;
\]
these are the test forms $dd^{*}\beta'$ of the three probes $\beta'\in\left\{\beta,\beta+\tau_{x_0}\beta,\beta-\tau_{x_0}\beta\right\}$.

The main result is Theorem~44, stated in Part~II together with its data, and it asserts the following. There are real numbers $m_{0,N}(e)$, one for every integer $N\ge0$ and every $e>0$ with $e\cdot h^{\frac{1}{2}}\le e_1$, such that $\left|m_{0,N}(e)\right|<\frac{1}{8}$; they depend on $N$, $e$, $\bar{m}$ and $\rho$ and on the renormalisation group scheme, but not on $T$. The scheme consists of the constants of Section~I.8 up to the ceiling $e_1$ of the charge, among them $\rho$ and the weights, of the basis $\left\{e^{0}_{i}\right\}$ of \cite[Section~5.2.5]{QEDI}, and of the systems of primitives on the larger reference torus at the levels below its stop level (Section~I.6.2). There are numbers $e_2>0$ and $C_{\bar{\kappa}},C'_{\bar{\kappa}}<\infty$, which depend only on the constants of Section~I.8, among them $\rho$, on $\bar{m}$ and on $(\beta,x_0)$, such that the following holds for every charge $e\in(0,e_2]$, the laws being formed with the mass assignment $N\mapsto m_{0,N}(e)$, that is, with $q=e\cdot\eta^{\frac{1}{2}}$ and $\boldsymbol{\mu}_N=\bar{m}\cdot\eta+m_{0,N}(e)$.

\emph{At the cutoffs.} Let $N_0(e)$ and $T^{*}(e)\ge1$ be the thresholds of Section~I.6.2, which are given by the formulas of \cite[Section~2.7]{QEDI}, formed with the constants of Section~I.8; the smaller reference torus has the side $R^{*}(e)=h\cdot L^{T^{*}(e)}$. For every cutoff $(N,T)$ with $N\ge N_0(e)$ and $T\ge T^{*}(e)+1$, and for $g\in\left\{g_0,g_{+},g_{-}\right\}$,
\[
\left|\kappa_4\left(X_{\eta,R}(g);E^{\mathrm{af}}_n\right)-e^4\cdot b(g)\right|\le C_{\bar{\kappa}}\cdot\left(e^{4+\sigma}+\eta^{\frac{1}{4}}\right),
\]
where $\kappa_4\left(X_{\eta,R}(g);E^{\mathrm{af}}_n\right)$ is the fourth cumulant of $X_{\eta,R}(g)$ under $E^{\mathrm{af}}_n$, that is, of $\omega\mapsto\omega(g)$ under $P_{N,T}$. The right side does not depend on $T$ and depends on $N$ only through $\eta$: the bound is uniform in the cutoffs, and its second term tends to zero with the spacing.

\emph{In every limit point.} Every limit point $\mu$ of $\left(P_{N,T}\right)$ satisfies, for $g\in\left\{g_0,g_{+},g_{-}\right\}$,
\[
\left|\kappa_4^{\mu}(g)-e^4\cdot b(g)\right|\le C'_{\bar{\kappa}}\cdot e^{4+\sigma},\qquad\left|S_2^{\mu}(g,g)-\|g\|_2^2\right|\le C'_{\bar{\kappa}}\cdot e^2,
\]
where $\kappa_4^{\mu}(g)$ and $S_2^{\mu}(g,g)$ are the fourth cumulant and the second moment of $\omega\mapsto\omega(g)$ under $\mu$. Moreover, with the joint cumulants $\kappa_4^{\mu}\left(f_1,\dots,f_4\right)$,
\begin{gather*}
\kappa_4^{\mu}(f)\ge\frac{1}{2}e^4\cdot b(f)>0,\\
\mathrm{Ext}^{\mu}_{x_0}(f):=\kappa_4^{\mu}(f)-3\kappa_4^{\mu}\left(f,f,\tau_{x_0}f,\tau_{x_0}f\right)\ge\frac{1}{2}e^4\cdot b(f).
\end{gather*}
Consequently $\mu$ is not Gaussian, and it is not a mixture of translation-invariant Gaussian laws that satisfy Gaussian domination with the constant one. Let $\mu=\int\mathrm{d}\nu(\xi)\,\mu_{\xi}$ be the ergodic decomposition of $\mu$ under the translations of $\mathbb{R}^3$. For $\nu$-almost every $\xi$ the component $\mu_{\xi}$ is invariant and ergodic under the translations, satisfies Gaussian domination with the constant one, is carried by the closed two-forms, and satisfies
\begin{gather*}
\left|\kappa_4^{\xi}(g)-e^4\cdot b(g)\right|\le C'_{\bar{\kappa}}\cdot e^{4+\sigma}\qquad\left(g\in\left\{g_0,g_{+},g_{-}\right\}\right),\\
\kappa_4^{\xi}(f)\ge\frac{1}{2}e^4\cdot b(f)>0,\qquad\mathrm{Ext}^{\xi}_{x_0}(f)\ge\frac{1}{2}e^4\cdot b(f);
\end{gather*}
in particular $\mu_{\xi}$ is not Gaussian.

\subsection*{Limit points with the volume first}

Let $e\in(0,e_2]$, and let the laws be formed with the mass assignment $N\mapsto m_{0,N}(e)$. Then every set $\mathcal{V}_N$ is non-empty, volume-first limit points exist, and every volume-first limit point is a limit point, so that the assertions above apply to it. On the basis of \cite[Theorems~1.3 and~1.4]{QEDII}, Theorem~44 asserts in addition the following properties of a volume-first limit point $\mu$.

It lies in the class $\mathcal{C}_9$ of \cite{QEDII}: it is invariant under the translations of $\mathbb{R}^3$ and, at every plane $\{x_i=\mathrm{po}\}$, $\{x_i-x_l=\mathrm{po}\}$, $\{x_i+x_l=\mathrm{po}\}$, $i<l$, $\mathrm{po}\in\mathbb{R}$, of the nine families, invariant under the reflection in the plane and reflection positive at it. It is invariant under all Euclidean motions of $\mathbb{R}^3$, reflections included, and reflection positive at every plane. In the reconstruction in the direction of the $j$-th coordinate axis, $j=1,2,3$, the joint spectral measure of the energy $\mathsf{H}_j$ and of the momenta $\mathsf{P}_i$, $i\ne j$, is carried by the round forward cone $\left\{(E,p):E\ge|p|_2\right\}$, and the vectors of $L^2(\mu)$ that are invariant under the translations along one coordinate axis are invariant under all translations. Let $U(a)\mathrm{Ga}$ denote the translate by $a\in\mathbb{R}^3$ of a function $\mathrm{Ga}\in L^2(\mu)$, and $P^{123}_{\mu}$ the orthogonal projection of $L^2(\mu)$ on the vectors that are invariant under all translations. For every $\mathrm{vp}>0$, all $\mathrm{Ga},\mathrm{Gb}\in L^2(\mu)$ that are measurable with respect to the $\sigma$-algebra generated by the evaluations of $\omega$ at the real smooth test forms with support in the open cube $(-\mathrm{vp},\mathrm{vp})^3$, and all $a$ with $|a|_{\infty}\ge4\,\mathrm{vp}$,
\[
\left|\mu\left(\overline{\mathrm{Ga}}\cdot U(a)\mathrm{Gb}\right)-\mu\left(\overline{\mathrm{Ga}}\cdot P^{123}_{\mu}\mathrm{Gb}\right)\right|\le c_{\mathrm{cl}}\cdot\|\mathrm{Ga}\|_2\cdot\|\mathrm{Gb}\|_2\cdot\frac{\mathrm{vp}}{|a|_{\infty}},
\]
where $\mu(\mathrm{Ga}):=\int\mathrm{d}\mu(\omega)\,\mathrm{Ga}(\omega)$, $\|\mathrm{Ga}\|_2:=\mu\left(\left|\mathrm{Ga}\right|^2\right)^{\frac{1}{2}}$, and $c_{\mathrm{cl}}$ is an absolute constant. The two-point function has the K\"all\'en--Lehmann form, with a positive spectral measure of mass at most one. The Schwinger functions $S_n$ of $\mu$ are continuous linear functionals on the space $\mathrm{Sw}_n$ of the Schwartz functions on $\mathbb{R}^{3n}$ with values in the complexified $n$-th tensor power of $\Lambda^2$, determined by $S_n\left(f_1\otimes\dots\otimes f_n\right)=\int\mathrm{d}\mu(\omega)\,\omega(f_1)\cdots\omega(f_n)$, and they satisfy, coinciding arguments included, the following four conditions.

(E0$'$) One has $S_0=1$ and $\left|S_n(\Psi)\right|\le2\cdot66^n\cdot(n!)^{\frac{1}{2}}\cdot q_{6n}(\Psi)$ for $n\ge1$ and $\Psi\in\mathrm{Sw}_n$, where $q_m(\Psi):=\max_{|\iota|\le m}\sup_y\left(1+|y|_2^2\right)^{\frac{m}{2}}\cdot\left|\partial^{\iota}\Psi(y)\right|$, the maximum being taken over the multi-indices $\iota$ and the components of $\Psi$.

(E1) $S_n$ is invariant under every Euclidean motion of $\mathbb{R}^3$, reflections included, acting on each argument as on two-forms.

(E2) For every vector $\hat{n}\in\mathbb{R}^3$ of length one and every finite sequence $\Psi_0\in\mathbb{C}$, $\Psi_k\in\mathrm{Sw}_k$, such that $\Psi_k$ is supported in $\left\{\left(y_1,\dots,y_k\right):\hat{n}\cdot y_a\ge0\text{ for all }a\right\}$, one has $\sum_{k,l}S_{k+l}\left(\Theta^{\hat{n}}_k\Psi_k\otimes\Psi_l\right)\ge0$; here $\Theta^{\hat{n}}_k\Psi$ is the complex conjugate of the image of $\Psi$ under the reflection in the plane $\left\{\hat{n}\cdot y=0\right\}$, acting on each argument.

(E3) $S_n$ is invariant under the permutations of the $n$ arguments together with their tensor indices.

Let $\mu=\int\mathrm{d}\nu(\xi)\,\mu_{\xi}$ be the ergodic decomposition of the volume-first limit point $\mu$ under the translations of $\mathbb{R}^3$. For $\nu$-almost every $\xi$ the component $\mu_{\xi}$ is a law on the closed two-forms that satisfies Gaussian domination with the constant one; it lies in $\mathcal{C}_9$ and is ergodic under the translations; it is invariant under all Euclidean motions and reflection positive at every plane; the round forward cone and the K\"all\'en--Lehmann form hold for it; and in each of the three reconstructions the kernel of $\mathsf{H}_j$ is spanned by the vacuum vector. For $\mathrm{vp}$, $\mathrm{Ga}$, $\mathrm{Gb}$ and $a$ as above, with $\mu_{\xi}$ in the place of $\mu$,
\[
\left|\mu_{\xi}\left(\overline{\mathrm{Ga}}\cdot U(a)\mathrm{Gb}\right)-\overline{\mu_{\xi}(\mathrm{Ga})}\cdot\mu_{\xi}(\mathrm{Gb})\right|\le c_{\mathrm{cl}}\cdot\|\mathrm{Ga}\|_2\cdot\|\mathrm{Gb}\|_2\cdot\frac{\mathrm{vp}}{|a|_{\infty}},
\]
and $\mu_{\xi}\left(\overline{\mathrm{Ga}}\cdot U(a)\mathrm{Gb}\right)\to\overline{\mu_{\xi}(\mathrm{Ga})}\cdot\mu_{\xi}(\mathrm{Gb})$ as $|a|\to\infty$ in any manner, for all $\mathrm{Ga},\mathrm{Gb}\in L^2\left(\mu_{\xi}\right)$. The Schwinger functions of $\mu_{\xi}$ satisfy (E0$'$)--(E3) and the cluster property

(E4) $S_{k+l}\left(\Psi\otimes\tau_a^{\otimes l}\Phi\right)\to S_k(\Psi)\cdot S_l(\Phi)$ as $|a|\to\infty$ in any manner, for $\Psi\in\mathrm{Sw}_k$ and $\Phi\in\mathrm{Sw}_l$, where $\tau_a^{\otimes l}$ translates each of the $l$ arguments of $\Phi$ by $a$.

Finally $\kappa_4^{\xi}(\mathrm{Gi}\cdot f)=\kappa_4^{\xi}(f)\ge\frac{1}{2}e^4\cdot b(f)>0$ for every Euclidean motion $\mathrm{Gi}$, so that $\mu_{\xi}$ is neither Gaussian nor the point mass at the zero field. The five conditions (E0$'$)--(E4) are those of \cite[Theorem~1.4]{QEDII}: the axioms of Osterwalder and Schrader of linear growth, Euclidean invariance, positivity, symmetry and clustering, numbered as in \cite{OsterwalderSchrader1975}, in the form for a field of two-forms on $\mathbb{R}^3$.
The condition (E0$'$) is the literal analogue, for two-form valued tests on $\mathbb{R}^{3n}$, of the linear growth condition of \cite{OsterwalderSchrader1975}, the integer called $s$ there being $6$ and the constants being $2\cdot66^n\cdot(n!)^{\frac{1}{2}}$; \cite{OsterwalderSchrader1975} formulates its conditions for scalar tests on $\mathbb{R}^{4n}$, and Theorem~44 asserts (E0$'$) in the form displayed above.
Thus, for $0<e\le e_2$ and with the bare mass shifts $m_{0,N}(e)$, the volume-first limit points of \cite{QEDII} are non-Gaussian laws of a closed two-form field on $\mathbb{R}^3$, invariant under the Euclidean group and reflection positive at every plane, whose Schwinger functions satisfy (E0$'$)--(E3); and almost every ergodic component is a non-Gaussian law with a unique vacuum whose Schwinger functions satisfy (E0$'$)--(E4).

\subsection*{The range of Theorem~44}

Theorem~44 concerns the all-flux law on the antiperiodic cubic tori with $8\cdot h\cdot L^{T+N}$ sites per side, the charges $e\in(0,e_2]$ and the mass assignment $N\mapsto m_{0,N}(e)$; its bounds at the cutoffs hold for $N\ge N_0(e)$ and $T\ge T^{*}(e)+1$.
For the other charges, mass assignments and cutoffs the laws have the properties recalled above from \cite{QEDII}.
The properties stated in the preceding subsection concern the volume-first limit points, among them the limit points along the joint sequences of cutoffs that \cite[Theorem~13.18]{QEDII} shows to be volume-first.

\subsection*{The method}

The properties of the limit points that hold for every charge and every mass assignment are those of \cite{QEDII}; they rest on Gaussian domination \cite{FSS76}, on the symmetries of the lattice, and on reflection positivity and chessboard estimates \cite{FILS78}, without expansion. The bound on the fourth cumulant rests on an expansion at small charge: the block spin renormalisation group of Ba{\l}aban for lattice gauge fields \cite{Balaban} and of Ba{\l}aban, O'Carroll and Schor for lattice fermions in an external gauge field \cite{BalabanOCarrollSchor}, in Dimock's form for three-dimensional QED \cite{DimockUR,DimockMB,DimockUVS}, organised as in \cite[Section~5]{QEDI}: block averaging, axial gauge for the fluctuation fields, Landau gauge for the background fields, and fermionic polymer activities with H\"older derivatives. As in \cite{QEDI}, chessboard estimates reduce the expectations that are needed to two reference tori, of the sides $R^{*}(e)=h\cdot L^{T^{*}(e)}$ and $L\cdot R^{*}(e)$ at the spacing $\eta$, on which the renormalisation group, with an external connection in the determinant, has no large-field region: Gaussian domination estimates the complement of the small-field set.

The statements are linked as follows. Proposition~14.1 assembles the properties of the fermion operators of all levels that the renormalisation step uses. One renormalisation step without large fields, with an external connection, at every level is the content of Propositions~24.1--24.3. The flow on the reference tori and the tuned pair, whose first entry is the bare mass shift $m_{0,N}(e)$ for all $N$ at which the pair is defined ($m_{0,N}(e)=0$ for the finitely many others), are Propositions~28.1--28.3, and the partition functions are bounded above and below in Propositions~32.1--32.3. At the stop level, the level at which the flow on the reference torus ends, the fermionic operators are studied as functions of the connection (Propositions~33.1--33.4); the terminal fermion integral is written as a function of variables attached to the cells of the reference torus (Proposition~35.5), and the normalised generating function of the probes, with one probe in each cell, is represented in the ring of four-jets in these variables (Proposition~35.8). The Bloch symbol of the chamber operator and its inverse, and the two-photon and four-photon coefficients of the free ratio of the chamber determinant together with their limits, are the subject of Propositions~36 and~37.1--37.3. This yields the reference data of Proposition~43.1. The transfer inequality for Wick powers (Proposition~42.2) bounds expectations on the tori of the laws by numbers formed on the two reference tori; it rests on the determinant at complex connections (Propositions~38.1--38.3), on the chessboard estimate and the lengthening of one side for coordinate cells of half-integer width (Proposition~39), on the transfer to all tori (Proposition~40), on the position of the probes in the open cells (Proposition~41) and on the defect of the two reference tori (Proposition~42.1). Proposition~43.2, the ultraviolet proposition, combines the two. It contains the bound at the cutoffs and a bound on products of Wick powers. For $m\in\{2,4\}$ and each of the three probes, let $\mathrm{Pw}_m$ be the Wick power of order $m$ of the curvature field smeared with the lattice test form $g^{(\eta)}$ formed with the probe, divided by $e^m$ and diminished by the $m$-th derivative $c_m$ at $0$ of the free ratio of the chamber determinant on the reference torus. For every set of $J$ points $y$ of the lattice $\frac{R^{*}(e)}{L}\cdot\mathbb{Z}^3$ that are pairwise incongruent modulo $R\cdot\mathbb{Z}^3$, the modulus of the expectation of the product of the functions $\mathrm{Pw}_m$ formed with the translates $\tau_y g^{(\eta)}$ is at most $\mathsf{u}_m^{J}$, where $\mathsf{u}_4\le C_U\cdot e^{\sigma}$ and $\left(1+\left|c_2\right|\right)\cdot\mathsf{u}_2\le C_U\cdot e^{\sigma}$, with a constant $C_U$ that depends on the same quantities as $C_{\bar{\kappa}}$. The fourth cumulant at the cutoff differs from $e^4$ times the free coefficient $b_{\mathrm{ch}}$ of the chamber determinant on the reference torus by at most a constant times $e^{4+\sigma}+\eta^2$, and $b_{\mathrm{ch}}$ differs from $b(g)$ by at most a constant times $e+\eta^{\frac{1}{4}}$. Theorem~44 rests on Proposition~43.2, on the properties of the limit points established in \cite{QEDII}, and on statements on single laws on $\mathcal{S}'$ (Propositions~45.1--45.3 and~45.5--45.7): the ergodic decomposition in the form in which it is used, regularity, closedness and determinacy, the Schwinger functions on tensor tests, transfer semigroups, and the invariance and reflection positivity at the planes of one further direction, which the ergodic components inherit (for the ergodic components of reflection positive laws see \cite{Froehlich76} and \cite[Section~19.7]{GlimmJaffe}).

\subsection*{Comparison with the Wilson fermion}

The photon field and its operators are those of \cite{QEDI}, and statements of \cite{QEDI} that contain no fermion object are used as they are stated there. What changes is the fermion. The chamber operator has $48$ sites, and its fields $192$ components, at each photon site, and its Wilson parameter is $\rho\in(0,1)$ in the place of one; all fermion operators of the scheme are constructed and estimated for it in Propositions~1.1--14.1, and Proposition~14.1 is the counterpart of \cite[Theorem~5.72]{QEDI}. The free chamber operator is split into a light and a heavy part, with the Schur complement and its symbol, a coercivity bound and a Poincar\'e inequality (Propositions~1.1--1.4); its Green function, the blocked operator at zero field and the first step are treated in Propositions~1.5--1.7. The inverses of the operators of the levels are obtained at zero and at small fields through local inverses on mirror tori, a parametrix and expansions in walks (Propositions~2.1--2.8); the weakened operators and polymer propagators, the kernels and the localised families of the operators of the levels, their covariances and their forms on tori of different sides follow (Propositions~3--11), and then the integral representation and the polymer expansion of the logarithm of the one-step determinant (Propositions~12 and~13). The constants of these operators depend on $\rho$, and several of them degenerate as $\rho\to0$ (Proposition~14.4).

In \cite{QEDI} the laws live on tori with $h\cdot L^{T+N}$ sites per side, an odd number. Here they live on the tori of \cite{QEDII}, with $8\cdot h\cdot L^{T+N}$ sites per side, while the reference tori, with $h\cdot L^{T^{*}(e)+N}$ and $h\cdot L^{T^{*}(e)+1+N}$ sites per side, are odd. The chessboard estimates and the lengthening of one side are proved for the chamber determinant (Proposition~39), the invariances of the determinant and of the functionals are established on rectangular tori of either parity (Proposition~50), the odd reference torus is shown to be symmetric under the reflections in its walls (Proposition~49), the lattice probes are formed on every torus and the cells on the even tori (Propositions~47.1 and~47.2), and the probes lie in the open cells (Proposition~41). The bound at the cutoffs holds for $T\ge T^{*}(e)+1$, where \cite[Theorem~3.1]{QEDI} has $T\ge T^{*}(e)$.

The free coefficient $b_{\mathrm{ch}}$ of the chamber determinant on the reference torus differs from $b(g)$ by at most a constant times $e+\eta^{\frac{1}{4}}$; hence the term $\eta^{\frac{1}{4}}$ in the bound at the cutoffs, where \cite[Theorem~3.1]{QEDI} has $\eta^{\frac{1}{2}}$. The bound $\left|c_2\right|+\left|c_4\right|^{\frac{1}{2}}\le C\cdot e^{-\mathfrak{e}_2}$ on the two-photon and four-photon coefficients $c_2$, $c_4$ of the free ratio is taken from the free field (Propositions~36 and~37.3), with the exponent $\mathfrak{e}_2=0$ in the place of the exponent $\mathfrak{e}_2=6\varepsilon$ of \cite[Section~2.7]{QEDI}; the second moments of the three tests in a limit point lie within $C'_{\bar{\kappa}}\cdot e^2$ of their free values, whereas the corresponding bounds of \cite[Theorem~3.1(iii), (vii)]{QEDI} carry the factor $e^{2-6\varepsilon}$ and concern, in a limit point, the test $f$ alone, and the three tests only in almost every ergodic component. Only the all-flux law is treated, whereas \cite{QEDI} treats the harmonic-averaged law as well. In \cite{QEDI} the limit points are reflection positive at the coordinate planes and invariant under the cubic group; here, through \cite{QEDII}, the volume-first limit points are invariant under all Euclidean motions and reflection positive at every plane, and almost all of their ergodic components have a unique vacuum and satisfy (E0$'$)--(E4).

\subsection*{The plan of the paper and its numbering}

After this introduction, four choices that are made once and for all are stated; they are quoted as choice~(1) to choice~(4), and they fix, among other things, the exponent $\mathfrak{e}_2=0$ and the value $m_{0,N}(e)=0$ at the finitely many $N$ at which the tuned pair is not defined. The paper then consists of five parts.

Part~I, \emph{Conventions}, fixes the notation and the conventions of the whole paper: names, numbering and standing letters (Section~I.0); the lattices, the clusters and sites of the chamber fermion, measures, kernels and pairs of sites (I.1); the spin matrices, the hopping matrices, the chamber operator, mass and charge, the Bloch symbol and the dictionary (I.2); block averaging, the operators of a level and the fields (I.3); the symmetries (I.4); norms, Grassmann integration with weights, the spaces of activities and the renormalisation step (I.5); scales, tori, cells, boundary conditions and twists (I.6); the laws (I.7); the constants, in the order in which they are chosen (I.8); and the statements of \cite{QEDI} and \cite{QEDII} that are used, with the identifications under which they are read (I.9). Its subsections are numbered I.$n.m$, and its displays and named facts are quoted as (I.1.4) or (I.ke2). Of the statements that are numbered under the fifty headings $1$ to $50$, Part~I contains one, Proposition~46 (Section~I.2.9), the dictionary between the operators of this paper and those of \cite{QEDI} and \cite{QEDII}; it shows in particular that the determinant and the laws do not depend on the point $v_0$. Its clause~(h) is proved in Section~I.2.9, and its clauses (a)--(g+) are proved in Chapter~46. Besides Proposition~46, Part~I contains facts that are stated together with their arguments, among them the Consequence in Section~I.1.9 and the Lemma on the bound for the axial functional in Section~I.9.2. An appendix to Part~I contains computations, among them facts on finitely many explicit $4\times4$ matrices, each obtained by a direct computation.

Part~II, \emph{The statements, in the order of dependence}, contains the statements of the analysis without their proofs. They stand under fifty headings, numbered $1$ to $50$, each of which treats one step of the analysis; a statement under the heading $n$ is Proposition~$n$, and when the heading carries several statements they are Propositions~$n.1$, $n.2$, \dots. The headings $1$ to $44$ follow, roughly, the order of the corresponding steps of \cite{QEDI}, and the headings $45$ to $50$ treat further steps. The statement under the heading $44$ is Theorem~44; under the heading $45$ stand the Propositions~45.1--45.7, statements on single laws on $\mathcal{S}'$; and the statement under the heading $46$ is Proposition~46, which stands in Part~I. Under the heading $39$ stands, besides Proposition~39, Remark~39.R, which makes assertions of its own: the slice data of the chamber determinant for every length $n_1\ge1$ of the torus in the time direction, and the chessboard estimate and the lengthening of one side for functions of the curvature on the closed cells; it is proved in Chapter~39, Section~8, and in Chapter~39r, which is named after it. The statements are given in the order of their logical dependence, up to the exceptions described below with Part~IV; this order is not the numerical order, and the statements of one heading do not always stand together: a heading may be taken up again further on. The headings form six groups. The group G1 (headings~1--14) contains the fermion operators of all levels, the assembly of their properties and the ordered list of the constants with the conditions on them (Propositions~14.1--14.4); G2 (headings~15--20 and~48) Grassmann integration with weights, the extraction and the adjustment of activities, global fields, locality, scaling and contraction, the block covariance and the determinants of one step; G3 (headings~21--25) the identities and the two displays of one step, one renormalisation step without large fields with an external connection, and the formal locality lemma; G4 (headings~26--32 and~49) the projection to smaller tori, the agreement of the steps on two tori, the flow and the tuned pair, the terminal integral, free determinants and twists, the zero-flux sector and the complement of the small-field set, the partition functions, and the symmetry of the odd reference torus; G5 (headings~33--37) the stop level, the four-jet representation and the free coefficients; G6 (headings~38--47 and~50) the determinant at complex connections, the chessboard estimates and the transfer to all tori, the probes and their cells, the transfer inequality for Wick powers, the reference data and the ultraviolet proposition, the statements on single laws, and Theorem~44. Each group has a chapter of definitions, which stands in Part~II before the first statement of the group in Part~II (Proposition~46, of the group G6, stands in Part~I).
The sections of the group G$i$ are numbered with the prefix G$i$, and so are its labels, as (G2.r5) or (G6.T-14).
These sections stand in three kinds of places: in the chapter of definitions of the group, as Sections~G1.1, G4.0.2 and~G5.1.3; under the headings of Part~II, next to the Propositions that use them, namely Sections~G3.1.0, G3.2.0, G3.4.0--G3.4.2 and~G3.5.0, which are all the numbered sections of the group G3 in Part~II, and Sections~G5.4.0 and~G5.6.0; and in the annexes of Part~III, as Sections~G1.Z1 and~G2.9.
The data of a statement are the constants of Section~I.8, chosen in the order given there and subject to the conditions listed in Proposition~14.2, and the further data that it names.
A statement may be followed by a paragraph headed ``Conditions on the constants'', which names conditions on the constants that enter the statement; conditions shared by the Propositions with the same first number may be stated before the first of them.
The conditions on the constants that enter the proof of a statement are listed in the section ``Conditions on the constants'' of the chapter that proves it, described below under Part~V; besides the conditions named after the statement, this section may contain further ones, in particular ranges printed in the lines (C0)--(C10) of Section~I.8.2.
As a rule, a statement is followed by a remark that compares it with the corresponding statement of \cite{QEDI} or \cite{QEDII}.

Part~III, \emph{Annexes}, consists of Annex~A and Annex~B.
Annex~A restates clauses of Part~I, of Propositions under other headings and of \cite{QEDI} and \cite{QEDII} in the form in which they are applied, arranged by the Propositions in whose proofs they are applied; its lists also name further clauses that these proofs use in the form in which they are stated.
The list of Section~G1.Z1, for Propositions~1.1--14.4, contains in addition entries on which no proof of these Propositions rests: they show with which constants, numerals and conditions the later Propositions are stated, that is, where the rows of Proposition~14.2\,(a), (b) whose reference or whose indication ``Used in'' is a later Proposition come from.
Annex~B collects, group by group, conditions on the constants; the complete list of the conditions on the constants, in the order of Section~I.8, is that of Proposition~14.2.

Part~IV consists of Section~IV.1, on the order of the proofs.
The statements of Part~II are given in an order close to that of their logical dependence, but at a few places the proof of a statement uses a clause of a later one; Section~IV.1 lists these places and gives an order of the clauses in which every proof uses only what has been proved before.

Part~V contains the proofs, in chapters named after the statements that they prove: Chapter~24 proves Propositions~24.1--24.3, and Chapter~14a proves Proposition~14.1. Inside a chapter, the sections and the displays, and the lemmas, propositions, theorems, corollaries, definitions, remarks and facts of the chapter itself, carry local numbers or names; they are cited bare inside the chapter and together with the name of the chapter from outside it, as in ``Chapter~24, Section~5.1''.
Each of the headings $3$, $5$, $6$, $7$ and $9$ of Part~II carries a single statement, namely Propositions~3, 5, 6, 7 and~9; hence in Chapters~39 and~39r a bare Proposition~$n.m$ with $n\in\{3,5,6,7,9\}$ is a statement of the chapter itself.
Every chapter has a section headed ``Conditions on the constants''; it is the last section of the chapter, except in Chapters~18a and~42, where an appendix follows it.
This section lists the conditions on the constants of Section~I.8 that enter the proofs of the chapter.
For each condition it names the places in the proofs where the condition enters and the place where the condition is imposed: as a rule a line among (C0)--(C10) of Section~I.8.2, a row or a label of Proposition~14.2\,(a), (b), (d), or a line of Proposition~14.2\,(c), (f).
A chapter whose proofs use no condition on the constants says so in this section.
The role of the section is the one stated in Section~IV.1, paragraph~E: it shows that every inequality between the constants that is used in a proof is a row of Proposition~14.2\,(a), (b), (d) or follows from these rows.
In many chapters, sections or paragraphs on uniformity state in addition for which values of the data the Propositions hold and on which quantities their constants depend.

Statements of \cite{QEDI} and \cite{QEDII} are cited with the numbers that they carry there.
A label in brackets is a name: apart from a prefix I.\ or G1., \dots, G6., neither its letters nor its digits indicate a position in this paper or in \cite{QEDI} or \cite{QEDII}.
A label identifies its object only together with the place where it is introduced; this place is given by the prefix I.\ or G1., \dots, G6., or by the statement, section or chapter that is named with the label or for which the label is declared.
The same bare label may be introduced at several places: for instance, (8.2) labels a display of Definition~24.D\,(D3) and a display of Proposition~30.1\,(d); (A6) and (A7) are rows of Proposition~14.2\,(b) and clauses of Proposition~35.4; and (g1)--(g3) are conditions of Proposition~14.2\,(d) and the three parts of the clause~(g) of Proposition~30.3.
A bare label denotes the object so labelled inside the statement or section that introduces it, under the heading of that statement or section, and in a chapter that has it as one of its own labels; where a declaration fixes the meaning of a bare label for a heading, a range of Propositions or a chapter, the declaration decides; elsewhere the place where the label is introduced is named with it.
The labels (8.1)$^{\mathrm{ch}}$, (8.2), (8.3)$^{\mathrm{ch}}$, (8.4) and (8.6)$^{\mathrm{ch}}$ denote the displays of Definition~24.D (Section~G3.4.1), which define the step map, wherever the preceding rule assigns no other meaning to them, and the index $\mathrm{ch}$, where a label carries it, is part of the name.
In particular, in Proposition~30.1 and in Chapter~30 the label (8.2) denotes the display of Proposition~30.1\,(d); the form (G3.8.2) denotes the display of Definition~24.D everywhere.
The prefix of a label names only the place where the label is introduced, Part~I or the sections of the group G$i$; there the label is printed with or without the prefix, and the same label is cited in both forms, as (I.k2) and (k2), or as (G4.Sc3) and (Sc3).
A bare label in a chapter of Part~V that is neither one of the chapter's own labels nor declared in the conventions of the chapter is to be looked up at the place named with it.
Like the labels, the numbers of sections are names: within a chapter they increase, and they need not be consecutive.
Names of more than one letter in upright type, such as $\mathrm{sw}$, $\mathrm{Chm}_0$, $\mathrm{vp}$ or $\mathrm{Sw}_n$, denote single objects.

\section*{Four choices}

Four choices are made for the whole paper; they are cited as choice~(1) to choice~(4).
The objects that occur in them are defined in Part~I and at the places cited with them.

\begin{itemize}
\item[(1)] \emph{The exponent $\mathfrak{e}_2$.}
One puts $\mathfrak{e}_2:=0$ (Section~I.8.5).
The reference datum (R4) of Proposition~43.1 is then obtained from the two-photon and the four-photon coefficients of the free field (Propositions~36 and 37.3\,(b)).
The formulas are written with the letter $\mathfrak{e}_2$, whose value is $0$.
\item[(2)] \emph{The sequence of bare masses.}
For every integer $N\ge0$,
\[
m_{0,N}(e):=
\begin{cases}
\text{the first entry of the tuned pair } P_N(e), & N\ge N_{\min}\left(R_{\mathrm{top}}(e),e\right),\\
0, & 0\le N<N_{\min}\left(R_{\mathrm{top}}(e),e\right);
\end{cases}
\]
the first case is exactly the range of $N$ in which the pair $P_N(e)$ is defined (Section~I.6.2, Proposition~28.3).
In \cite[Theorem~21.3]{QEDI} the sequence is put equal to $0$ for all $N<N_0(e)$, with the number $N_0(e)$ of that theorem.
The two sequences differ for finitely many $N$ only.
\item[(3)] \emph{The two local pieces of $T_k(A)$.}
The local pieces $T^{(1)}_{\mathbf{Y}}$ and $T^{(2)}_{\mathbf{Y}}$ of $T_k(A)$ are operators on the whole unit lattice whose rows at the unit sites outside $\mathbf{Y}$ vanish ((I.5.4), Propositions~9 and 23.1).
\item[(4)] \emph{The line (C6$'$).}
The constants of \cite[Section~20.1.7]{QEDI} that contain neither $\bar{m}$ nor the probe are fixed at the end of the line (C6) of the ordered list of the constants, under the label (C6$'$) (Section~I.8.2, Proposition~14.2\,(a)).
The constant $\mu_-$ is fixed in the line (C4).
\end{itemize}

\section*{Part I. Conventions}

\subsection*{I.0. Scope, names, numbering, standing letters}

\subsubsection*{I.0.1. Scope}

The notation and the conventions fixed in Part~I are in force throughout the paper, in all statements and proofs, except where a chapter states a different convention at its beginning.
For what Part~I does not fix, the conventions of \cite[Section~5]{QEDI} that are listed in Section~I.9.2 apply in the chapters on the renormalisation group, which correspond to \cite[Sections~13--20]{QEDI},
and the conventions of \cite[Section~2]{QEDII} apply in the chapters that rest on \cite{QEDII} and on \cite[Sections~7--9]{QEDI};
both are read with the identifications of Section~I.9.

\subsubsection*{I.0.2. Companion papers}

The companion paper \cite{QEDI} treats Wilson fermions: the existence of limit points and their fourth cumulant.
The companion paper \cite{QEDII} treats chamber fermions: reflection positivity, all mirrors, and the Osterwalder--Schrader axioms.
Said of a law, the label (G) denotes Gaussian domination (\cite[Section~2.5]{QEDI}, \cite[Definition~2.5]{QEDII}, and, for laws on $\mathcal{S}'(\mathbb{R}^3;\Lambda^2)$, (G6.n3), the clause (n3) of the notation for laws on $\mathcal{S}'$ used in the Propositions~45.1--45.7).

\subsubsection*{I.0.3. Numbering of statements}

The statements are arranged under the numbers $1$ to $50$.
The title of the number $46$ stands in the heading of Section~I.2.9, in Part~I; the titles of all other numbers stand in headings of Part~II.
In Part~II a number may have several headings (a heading may recur, marked as continued, and the number $35$ has a first and a second part), and the numbers $36$ and $37$ share one heading.
A statement under the number $n$ is Proposition~$n$, with clauses (a), (b), \dots;
when several statements stand under the number $n$, they are Propositions~$n.1$, $n.2$, \dots.
The statement under the number $44$ is Theorem~44.
Under the number $45$ stand the Propositions~45.1--45.7.
They state, with proofs, \cite[Theorems~7.15, 7.17 and Corollary~7.18]{QEDI} for laws with the properties (L1)--(L4) of \cite[Section~7.5]{QEDI}, \cite[Lemma~21.36]{QEDI}, \cite[Theorem~8.8(v)]{QEDI}, a statement on tensor tests (Proposition~45.5), and \cite[Theorem~8.1]{QEDI} together with the display \cite[(8.1)]{QEDI} in its proof (\cite[Section~8.5]{QEDI}), which holds for every measurable $\Phi\ge0$ and every integrable $\Phi$.

Where a fact stated in Part~I is not accompanied by its argument, the statement from which it follows, or the place where it is proved, is named with it:
a Proposition, a statement of \cite{QEDI} or \cite{QEDII}, a section of a chapter of Part~V, or the literature.

\subsubsection*{I.0.4. Symbols}

For a complex number $z$, $\overline{z}$ is its complex conjugate; $O^{*}$ is the adjoint and $O^{T}$ the transpose of an operator $O$.
The transpose is always taken with respect to the bilinear pairings of Sections~I.1.6 and I.3.6.
Besides the transpose, the letter $T$ stands in an upper position in three ways.
In $D_A^{T}$ and $\det D_A^{T}$ of \cite[Definitions~2.11 and 4.1]{QEDII} it is the torus $T=(\Gamma,\chi)$, which is written $\mathbb{T}$ there, and $D_A^{T}$ is the operator of \cite[Definition~2.10]{QEDII} on the spinors of $T$ (Proposition~46\,(e)).
The Ursell coefficients $\mathrm{urs}^{T}(X_1,\dots,X_n)$ of Definition~24.D, written $\rho^{T}$ in \cite[Proposition~5.58]{QEDI}, are numbers attached to polymers; their upper index is part of the name and is not a transpose.
In the powers $L^{T}$ and $L^{T+N}$ the letter $T$ is the integer $T\ge0$ of Section~I.0.5, which forms with $N$ the cutoff $(N,T)$.
The symbols $\otimes$, $\oplus$ and $\rtimes$ denote the tensor product, the direct sum and the semidirect product.

The letter $e$ denotes the charge, $\exp(\cdot)$ denotes the exponential function, and $\mathrm{e}:=\exp(1)$.
A power $e^{u}$ of the italic letter $e$ is a power of the charge, and its exponent $u$ is formed only with integers (numerals, lengths of multi-indices and cardinalities included) and with the letters $m$, $N_*$, $a_L$, $\varepsilon$, $\mathfrak{e}$, $\mathfrak{e}_2$, $\sigma$.
A value of the exponential function is written $\exp(u)$ or as the power $\mathrm{e}^{u}$ of $\mathrm{e}$, whatever the exponent $u$ is, numerals included.
Hence a power of $e$ or of $\mathrm{e}$ whose exponent begins with $\pm i$, as $\mathrm{e}^{\pm i(\cdot)}$, or contains a cochain, a field, a test form, an operator, a pairing $\langle\cdot,\cdot\rangle$, a norm, a point or a momentum, is a power of $\mathrm{e}$ and a value of the exponential function.
Among the forms that occur are
\[
\mathrm{e}^{-\frac{1}{2}\|dA\|^2},\quad
\mathrm{e}^{-\frac{1}{2}\|F^{\mathrm{u}}\|^2},\quad
\mathrm{e}^{-\frac{1}{2}\|dA+r_q\mathfrak{z}\|^2},\quad
\mathrm{e}^{\langle da,F^{\mathrm{u}}\rangle},\quad
\mathrm{e}^{\frac{1}{2}\langle da,da\rangle},\quad
\mathrm{e}^{\langle g,F\rangle},\quad
\mathrm{e}^{\frac{1}{2}\|g\|^2},
\]
\[
\mathrm{e}^{X(f)},\quad
\mathrm{e}^{zX(f')},\quad
\mathrm{e}^{\frac{1}{2}\|f\|_2^2},\quad
\mathrm{e}^{-t\mathsf{H}_j},\quad
\mathrm{e}^{-S_T(A)},\quad
\mathrm{e}^{-\bar{\psi}\,\mathrm{Dm}\,\psi}.
\]
The following symbols are names with an upper index, not powers:
$e^F_0$ of \cite[Theorem~5.72]{QEDI};
$e^{+}_{n,j}$ and $e^{-}_{n,j}$ of \cite[Definition~2.7]{QEDII};
the vectors $e^{0}_{i}$ of \cite[Section~5.2.5]{QEDI}, which form the orthonormal basis, fixed there once for each $L$, of the kernel of $\tau$ in the space of the real functions on the interior bonds of the $L$-block $B(0)$;
and their translates $e^{y}_{i}$ by the points $y$ of the $L$-lattice, whose upper index is a point (so also in $e^{y+v}_{i}$).
The symbols $e_j^{\perp}$, the plane orthogonal to $e_j$, and $e_\mu^{T}$, the row vector of $e_\mu$ (as in $1-2e_\mu e_\mu^{T}$), are formed from the unit vectors and are not powers either.

$\mathbf{1}[\,\cdot\,]$ denotes an indicator.
The product of an indicator with a function that is defined only on the set of the indicator has the value $0$ off this set.
For a finite set $I$, $\#|I|$ denotes the number of its elements.
The set of the integers and half-integers is $\frac{1}{2}\mathbb{Z}$.

\subsubsection*{I.0.5. Standing letters}

The following letters denote numbers that are fixed once and for all.
Only their ranges are given here; the order in which they are chosen and the conditions on them are stated in Section~I.8.
\begin{itemize}
\item $\rho\in(0,1)$, the Wilson parameter.
\item $w=(w_0,w_1,w_2,w_3):=\left(6,6\sqrt{2},6\sqrt{2},6\right)$, the standard weights, and
\[
\mathrm{sw}:=w_0+w_1+w_2+w_3=12+12\sqrt{2}=28.97\ldots.
\]
\item $b:=1$, the letter being kept in the formulas, and $b_L:=b/L$.
\item $\alpha\in(0,1)$, the H\"older order.
\item $\kappa>0$, $\alpha_0>0$ and $\varepsilon>0$.
\item $L$, an odd integer with $L\ge5$.
\item $M$, a power of $L$ with $M\ge L^2$ and $M>4L$.
\item $\bar{m}\in(0,1]$, the fermion mass, and $h:=L^{\left\lfloor\log_L\frac{1}{\bar{m}}\right\rfloor}$.
\item $e>0$, the charge.
\item Integers $N\ge0$ and $T\ge0$, and $\eta:=L^{-N}$, the spacing.
\item $e_0,e_1,e_2>0$, ceilings for charges.
\end{itemize}
As numbers, $e_0$, $e_1$, $e_2$, with a numeral as index, always denote these ceilings and never running charges; as vectors, $e_1$, $e_2$, $e_3$ are the unit vectors of $\mathbb{R}^3$ (Section~I.1.1).
Which of the two is meant is decided by the kind of object, number or vector, as in \cite[Definition~5.3(g)]{QEDI}.

\subsection*{I.1. Lattices, clusters, sites, measure, kernels, positions, pairs}

\subsubsection*{I.1.1. Photon lattices}

The objects of this subsection are those of \cite[Sections~5.2.3 and 5.2.5]{QEDI}; only the letters for the sites differ.

\paragraph{Directions.}
The unit vectors of $\mathbb{R}^3$ are $e_1,e_2,e_3$ (as numbers, $e_1$ and $e_2$ are ceilings: Section~I.0.5), and in all chapters the directions are $\mu,\nu,\mathfrak{l}\in\{1,2,3\}$.
In \cite[Sections~5, 11, 13--20, 21.3.3 and 21.7]{QEDI} the directions are labelled by $\{0,1,2\}$, and the direction $\mu$ there is the direction $\mu+1$ here;
so the time direction $0$ of \cite{QEDI} is $e_1$, which is the time direction of \cite{QEDII}.
For a vector $x$ of $\mathbb{R}^3$, $|x|=|x|_\infty$ is the sup norm, $|x|_1$ the $\ell^1$ norm and $|x|_2$ the Euclidean norm.
This convention holds in all statements of this paper and in \cite[Sections~5 and 11--21]{QEDI}.

The following places are exceptions:
\cite[Sections~7--10]{QEDI} (in Section~10 the flux vector $w$ of \cite[Lemma~10.2(a)]{QEDI}),
\cite[Section~12.4]{QEDI} (the momentum $p$ of \cite[Lemma~12.4]{QEDI}),
\cite[Section~21.5]{QEDI},
and \cite[Theorem~1.4, Definitions~2.2, 2.4, 2.10, Remark~3.1, Lemma~3.3, Sections~12 and 14]{QEDII}.
In these places a modulus without index of a point $y$ or of a momentum $p$ of $\mathbb{R}^3$ is the Euclidean norm $|y|_2$, respectively $|p|_2$.
This applies to $p_R$ in \cite[Lemma~7.5]{QEDI}, to the multiplier $1-\frac{p\,p^{T}}{|p|^2}$ of $\Pi_{\mathrm{c}}$, and to the K\"all\'en--Lehmann form of \cite[Theorem~1.4(d)]{QEDII};
moreover, in the proof of \cite[Lemma~7.5]{QEDI} the modulus of the value $v$ of a two-form is $\left(\sum_{\mu<\nu}|v_{\mu\nu}|^2\right)^{\frac{1}{2}}$.
\cite[Definition~2.1]{QEDII} defines only $|x|_\infty$, with the index.
The moduli $|n_\ell|$, $|\mathrm{dlt}_i(x)|$ and $|x'-x|$ of \cite[Definitions~2.2, 2.4 and 2.10]{QEDII} are Euclidean:
$|n_\ell|^2\in\{1,2\}$, and $|\mathrm{dlt}_i|=2\operatorname{dist}(v_0,H_i)$, where $\mathrm{dlt}_i(x)=x^{(i)}-x$ is the displacement of the points of $\mathbb{R}^3$ along a bond of type $i$ and $H_i$ is the wall of type $i$ of the chamber (written $\Pi_i$ in \cite{QEDII}).
In Section~I.4.8, $|n_\ell|^2$ denotes the squared Euclidean length $|n_\ell|_2^2$, which is $1$ for an axis family and $2$ for a diagonal family \cite[Definition~2.2]{QEDII}.
Proposition~46\,(a) and the other formulas of this paper in which such a modulus occurs write it with the index~$2$, except where $n_\ell$ is a unit vector $e_\mu$, for which $|n_\ell|=|n_\ell|_2=1$.

\paragraph{Level, side.}
The level is an integer $k\ge0$; one puts $\eta_k:=L^{-k}$ and $\mathbb{L}_k:=\eta_k\mathbb{Z}^3$.
A side is an element $\ell$ of
\[
\mathrm{Cl}:=\left\{\ell:\ \ell/(LM)\ \text{is an odd integer}\ \ge13\right\}.
\]
The fine photon lattice is $T_k:=\mathbb{L}_k/\ell\mathbb{Z}^3$, the unit lattice is $T^0_k:=\mathbb{Z}^3/\ell\mathbb{Z}^3$, and the $L$-lattice is $L\mathbb{Z}^3/\ell\mathbb{Z}^3$.
The number of the unit sites is $n_k:=\ell^3$, and $d(\cdot,\cdot)$ is the sup distance on the torus.

\paragraph{Letters for sites.}
The letters $c,c',\tilde{c}$ denote sites of $\mathbb{L}_k$ or of $T_k$, the photon sites.
The letters $y,y',z,z'$ denote unit sites or sites of the $L$-lattice.
The letters $x,x',\tilde{x}$ denote fermion sites (Section~I.1.3).
In \cite{QEDI} a fine site is written $x$; in a statement of \cite{QEDI} on photon objects that is used here, this $x$ is read as $c$.

\paragraph{Cells.}
The cells are the sites;
the bonds, that is, the ordered pairs $(c,c')$ with $|c-c'|_2=\eta_k$;
the plaquettes $(c;\mu\nu)$, $\mu<\nu$, with the corners $c$, $c+\eta_ke_\mu$, $c+\eta_ke_\mu+\eta_ke_\nu$, $c+\eta_ke_\nu$ in this order;
and the cubes.
A function of the position is evaluated on a cell at the barycentre of the cell \cite[Definition~5.4]{QEDI}.

\paragraph{Bond functions.}
A bond function, or connection, is a complex function $A$ on the bonds with $A(c',c)=-A(c,c')$; one puts $A_\mu(c):=A(c,c+\eta_ke_\mu)$ and $A(c,c):=0$.
For functions $f$ and $\boldsymbol{\lambda}$ of the sites and a bond function $A$,
\[
(\partial_\mu f)(c):=\eta_k^{-1}\cdot\left(f(c+\eta_ke_\mu)-f(c)\right),\qquad
(d\boldsymbol{\lambda})_\mu:=\partial_\mu\boldsymbol{\lambda},\qquad
(dA)_{\mu\nu}:=\partial_\mu A_\nu-\partial_\nu A_\mu\quad(\mu<\nu);
\]
on $T^0_k$ the same formulas hold with $1$ in the place of $\eta_k$.
In lattice units (Sections~I.6 and I.7), $d$ is the coboundary of \cite[Definition~2.1]{QEDII}, which is also that of \cite[Sections~2.1 and 10.1.1]{QEDI}:
\begin{gather*}
(dA)_{\mu\nu}(c)=A_\mu(c)+A_\nu(c+e_\mu)-A_\mu(c+e_\nu)-A_\nu(c),\\
(dF)(c)=\left[F_{23}(c+e_1)-F_{23}(c)\right]-\left[F_{13}(c+e_2)-F_{13}(c)\right]+\left[F_{12}(c+e_3)-F_{12}(c)\right].
\end{gather*}
The two definitions agree.

\paragraph{Paths; the axial functional.}
For a lattice path $\boldsymbol{\Gamma}$, $A(\boldsymbol{\Gamma})$ is the sum of the values $A(b)$ over the bonds $b$ of $\boldsymbol{\Gamma}$, each bond oriented along the path, without weight.
For $c,c'\in\mathbb{L}_k$ and an order $\pi$ of the three coordinates, $\boldsymbol{\Gamma}^{\pi}(c,c')$ is the rectilinear path from $c$ to $c'$ that brings the coordinates to their final values in the order $\pi$.
The axial functional is
\[
(\tau A)(c,c'):=\frac{1}{6}\sum_{\pi}A\left(\boldsymbol{\Gamma}^{\pi}(c,c')\right),\qquad (\tau A)(c,c):=0,
\]
the sum extending over the six orders $\pi$.
On the torus the paths are formed on a lift with $|c-c'|<\frac{\ell}{2}$.

\paragraph{Blocks.}
For $0\le j\le k+1$ and $z\in L^j\eta_k\mathbb{Z}^3$ let
\[
\mathrm{Bl}_j(z):=\left\{c\in\mathbb{L}_k:\ |c-z|<\frac{L^j\eta_k}{2}\right\}.
\]
For a unit site $y$, the unit block is $\Delta_y=B_k(y):=\mathrm{Bl}_k(y)$, and $[c]$ denotes the unit site $y$ with $c\in\Delta_y$.
For $y'\in L\mathbb{Z}^3$, the $L$-block is
\[
B(y'):=\left\{y\ \text{a unit site}:\ |y-y'|<\frac{L}{2}\right\}.
\]

\paragraph{Cubes, polymers, parameter cubes.}
The $M$-cubes $\square$, centred at the points of $M\mathbb{Z}^3$, the $LM$-cubes, the sets $X^{\sim i}$ and $\square^{\dagger}:=\square^{\sim L}$, the set $D_k$ of the polymers, the numbers $|X|_M$, $\operatorname{Vol}(X):=M^3\cdot|X|_M$ and $d_M(X)$, the small and the large polymers, $\bar{X}^{L}$, the parameter cubes $\pi_k$, $\pi'_k$, $\pi''_k$, the relation ``linked'', the link-components, the polydisk $\mathrm{Pd}_k$, the numbers $N_c=(2L+1)^3$, $c_{l}=2L$, $d_U=2L+1$, and the balls $K_z:=\{c:\ |c-z|\le2L\}$ are those of \cite[Section~5.2.3 and Definitions~5.7, 5.8]{QEDI}.
They are sets of photon sites.
Two $M$-cubes have a common face if their centres differ by $Me_\mu$ or by $-Me_\mu$ for one direction $\mu$ (on the torus: modulo $\ell\mathbb{Z}^3$); such cubes are called adjacent in \cite[Section~13.3]{QEDI}.
For an integer $i\ge0$, the set $X^{\sim i}$, called in \cite{QEDI} the enlargement of $X$ by $i$ layers of $M$-cubes, is the union of the $M$-cubes whose centre has sup distance at most $iM$ from the centre of an $M$-cube of $X$.
Thus, for an $M$-cube $\square$ with the centre $z$,
\[
\square^{\sim i}=\left\{c:\ d(c,z)<\left(i+\frac{1}{2}\right)\cdot M\right\},
\]
a cube of the side $(2i+1)\cdot M$; \cite[Definition~5.7]{QEDI} uses $i=L$ and $i=5$.
The index of a parameter cube is written $\mathrm{pc}$, so that $s=(s_{\mathrm{pc}})$; in \cite{QEDI} it is written $s_c$.
A link-component of the union $Y$ of a set $\mathscr{Y}$ of parameter cubes of one kind is the union of the cubes of a link-component of $\mathscr{Y}$.
In the definition of $d_M$, a closed straight segment of $\mathbb{R}^3/\ell\mathbb{Z}^3$ is the image of a closed straight segment of $\mathbb{R}^3$, with the Euclidean length of that segment, and the length of a net refers to its representation as a union of segments.

\paragraph{Bonds and sets.}
A bond meets a set $X$ of photon sites if one of its end points lies in $X$;
it is contained in $X$ if both end points lie in $X$;
it is at distance $\le1$ from $X$ if both end points have sup distance $\le1$ from $X$
(\cite[Definition~5.24, (k2)]{QEDI}, \cite[Section~17.1.2]{QEDI}).

\paragraph{Faces.}
The numbers $L$ and $M$ are odd.
Hence the faces of all blocks, of all $M$-cubes and of all $LM$-cubes lie in the link planes $\left\{y\in\mathbb{R}^3:\ y_\mu\in\eta_k\cdot\left(\mathbb{Z}+\frac{1}{2}\right)\right\}$ of $\mathbb{L}_k$, and no photon site lies on a face (\cite[Section~13.1.2]{QEDI}, \cite[Section~17.4.1, (g0)]{QEDI}).

\paragraph{Sets that do not wrap; lifts.}
Let $Z$ be a set of sites of $T_k$.
A lift of $Z$ is a subset $\tilde{Z}$ of $\mathbb{L}_k$ that the covering map $\mathbb{L}_k\to T_k$ maps onto $Z$ and whose sup-diameter $\operatorname{diam}(\tilde{Z})$ satisfies
\[
\operatorname{diam}(\tilde{Z})+\eta_k<\ell .
\]
The set $Z$ does not wrap if it has a lift.
In this case the covering map is injective on $\tilde{Z}$, because two distinct points of $\mathbb{L}_k$ with the same image differ by a non-zero vector of $\ell\mathbb{Z}^3$, whose sup norm is at least $\ell$.
For $c\in Z$ let $\tilde{c}$ denote its preimage in $\tilde{Z}$.
Let $\mathrm{sd}>0$ satisfy $\operatorname{diam}(\tilde{Z})+\mathrm{sd}<\ell$, and let $c,c'\in Z$ with $d(c,c')\le\mathrm{sd}$.
Then $\tilde{c}-\tilde{c}'$ is the only representative of $c-c'$ modulo $\ell\mathbb{Z}^3$ of sup norm $\le\mathrm{sd}$:
a representative $\tilde{c}-\tilde{c}'+\ell n$, $n\in\mathbb{Z}^3$, of sup norm $\le\mathrm{sd}$ satisfies $|\ell n|\le\mathrm{sd}+\operatorname{diam}(\tilde{Z})<\ell$, hence $n=0$, and such a representative exists because $d(c,c')\le\mathrm{sd}$.
In particular $|\tilde{c}-\tilde{c}'|=d(c,c')$.
Consequently, for $\mathrm{sd}=\eta_k$ the bonds of $Z$ are the images of the bonds of $\tilde{Z}$, and for $\mathrm{sd}=M$ the pairs at distance $\le M$ inside $Z$, together with their six paths, are those of $\tilde{Z}$.
If $Z$ is connected through nearest neighbours (as balls, cubes and polymers are), two lifts of $Z$ differ by a vector of $\ell\mathbb{Z}^3$.
A lift of a union of $M$-cubes is a union of $M$-cubes of $\mathbb{L}_k$ with the same common faces.
The definition of a lift and of a set that does not wrap agrees with \cite[Definition~5.5(9)]{QEDI}: there too a lift is a subset of $\mathbb{L}_k$ that the covering map $\mathbb{L}_k\to T_k$ maps onto the set and whose sup-diameter is less than $\ell-\eta_k$, that is, $\operatorname{diam}(\tilde{Z})+\eta_k<\ell$.
The consequences stated above, namely the injectivity of the covering map on a lift, the rule for the representatives of $c-c'$ with the margin $\mathrm{sd}$ and its cases $\mathrm{sd}=\eta_k$ and $\mathrm{sd}=M$, the uniqueness of the lift of a set connected through nearest neighbours up to a vector of $\ell\mathbb{Z}^3$, and the lifts of unions of $M$-cubes, are stated there as well, with the letters $W$, $x$, $l$, $s$ in the place of $Z$, $c$, $\ell$, $\mathrm{sd}$.

\subsubsection*{I.1.2. The point group, the walls}

Let
\[
O_{\mathrm{cub}}:=O(3)\cap\mathrm{GL}(3,\mathbb{Z}),
\]
the group of the $48$ signed permutation matrices, which is the group $O_{\mathrm{cub}}$ of \cite[Section~2.5]{QEDI}, and let $\{\pm1\}^3$ be its subgroup of the $8$ diagonal matrices.
Put
\[
\hat{n}_0:=e_3,\qquad
\hat{n}_1:=\frac{e_2-e_3}{\sqrt{2}},\qquad
\hat{n}_2:=\frac{e_1-e_2}{\sqrt{2}},\qquad
\hat{n}_3:=-e_1,
\]
and $s_i:=1-2\hat{n}_i\hat{n}_i^{T}\in O_{\mathrm{cub}}$ for $i=0,1,2,3$.
Thus $s_0=\operatorname{diag}(1,1,-1)$, $s_1$ is the transposition $e_2\leftrightarrow e_3$, $s_2$ is the transposition $e_1\leftrightarrow e_2$, and $s_3=\operatorname{diag}(-1,1,1)$.
One has $s_i^2=1$, $\det s_i=-1$ and $s_i\hat{n}_i=-\hat{n}_i$, and $s_0,s_1,s_2$ generate $O_{\mathrm{cub}}$:
the transpositions $s_1$, $s_2$ generate the permutation matrices, and the conjugates of $s_0$ under them are the three sign changes of a single coordinate.
The standard weights satisfy the closure relation
\begin{equation}\tag{I.1.1}
\sum_{i=0}^{3}w_i\cdot\hat{n}_i=6e_3+6(e_2-e_3)+6(e_1-e_2)-6e_1=0 .
\end{equation}
It holds for the standard weights; it is a choice of the weights and not a consequence of a symmetry (Section~I.4.6).
In \cite{QEDII} the $s_i$ are the affine reflections in the walls of the chamber.
Their linear parts, written $\vec{s}_i$ in \cite[Remark~3.5]{QEDII}, are the matrices $s_i$ above, and the affine reflection $s_3$ of \cite{QEDII} is the map $y\mapsto e_1+s_3y$.

\subsubsection*{I.1.3. Fermion sites: clusters, group coordinates, neighbours}

Let
\[
\mathbb{L}^{\mathrm{f}}_k:=\mathbb{L}_k\times O_{\mathrm{cub}},\qquad
T^{\mathrm{f}}_k:=T_k\times O_{\mathrm{cub}}.
\]
A fermion site is a pair $x=(c,g)$; here $c=\pi(x)$ is the photon site under $x$, and $g\in O_{\mathrm{cub}}$ is the label of $x$.
The cluster of $c$ is $\{c\}\times O_{\mathrm{cub}}$, the set of the $48$ fermion sites over $c$; ``$x$ lies over $c$'' means $\pi(x)=c$.
The group $\ell\mathbb{Z}^3$ and all translations act on $c$ only: $(c,g)+v:=(c+v,g)$.
For $i=0,1,2,3$ put
\begin{equation}\tag{I.1.2}
\begin{aligned}
&x^{(i)}:=(c,gs_i)\quad(i=0,1,2),\qquad x^{(3)}:=\left(c+\eta_k\cdot ge_1,\ gs_3\right),\\
&c_3=c_3(x):=\pi\left(x^{(3)}\right)=c+\eta_k\cdot ge_1,\qquad \hat{d}_i(x):=-g\hat{n}_i\quad\left(\text{so that }\hat{d}_3(x)=ge_1\right).
\end{aligned}
\end{equation}
The pair $\left\{x,x^{(i)}\right\}$ is the bond of type $i$ at $x$.
The following facts are read off from (I.1.2) and from $s_i\hat{n}_i=-\hat{n}_i$.
\begin{description}
\item[(I.s1)] $\left(x^{(i)}\right)^{(i)}=x$; the map $x\mapsto x^{(i)}$ is a bijection of $\mathbb{L}^{\mathrm{f}}_k$ and of $T^{\mathrm{f}}_k$ that commutes with the translations.
\item[(I.s2)] $\hat{d}_i\left(x^{(i)}\right)=-\hat{d}_i(x)$.
\item[(I.s3)] $\sum_iw_i\cdot\hat{d}_i(x)=-g\sum_iw_i\cdot\hat{n}_i=0$, by (I.1.1).
\item[(I.s4)] The bonds of the types $0$, $1$, $2$ stay in the cluster; a bond of type $3$ joins the clusters of the nearest neighbours $c$ and $c_3$.
For an oriented photon bond $(c,c+\eta_ka)$, $a\in\{\pm e_\mu\}$, the sites $x$ over $c$ with $c_3(x)=c+\eta_ka$ are the $8$ sites $(c,g)$ with $ge_1=a$, and $x\mapsto x^{(3)}$ maps them onto the $8$ sites over $c+\eta_ka$ with $ge_1=-a$.
\item[(I.s5)] $\det g$ changes its sign along every bond.
\end{description}
Indeed, $s_i^2=1$ gives (I.s1) for $i\le2$, and $s_3e_1=-e_1$ gives $\left(x^{(3)}\right)^{(3)}=\left(c+\eta_k\cdot ge_1+\eta_k\cdot gs_3e_1,\ gs_3^2\right)=x$;
next, $\hat{d}_i\left(x^{(i)}\right)=-gs_i\hat{n}_i=g\hat{n}_i$;
for each $a\in\{\pm e_\mu\}$ exactly $8$ signed permutation matrices $g$ have the first column $ge_1=a$, and $(gs_3)e_1=-ge_1$;
finally, $\det(gs_i)=-\det g$.

\paragraph{The geometric picture.}
The fundamental chamber is
\[
\mathrm{Chm}_0:=\left\{y\in\mathbb{R}^3:\ 0<y_3<y_2<y_1<\frac{1}{2}\right\}
\]
(\cite[Definition~2.3]{QEDII}, where it is written $C_0$), and $v_0$ is a point of $\mathrm{Chm}_0$.
In lattice units, $\mathrm{geo}(c,g):=c+gv_0$ is the point of $\mathbb{R}^3$ for which the fermion site $(c,g)$ stands; the set of the fermion sites of \cite[Definition~2.4]{QEDII} is $\mathrm{geo}\left(\mathbb{L}^{\mathrm{f}}_0\right)$.
Apart from the map $\mathrm{geo}$ and the objects formed with it, namely the symbol $\mathrm{Dh}(p)$ and the operator $Y$ of Section~I.2.6, no object defined in Part~I depends on $v_0$:
written in the coordinates $(c,g)$, the sites, the bonds and their types, the map $\pi$ and the directions $\hat{d}_i$ (Proposition~46\,(a)), the operator (Proposition~46\,(b)) and the determinant (Proposition~46\,(e)) do not contain $v_0$ (\cite[Remark~3.5]{QEDII}).

\subsubsection*{I.1.4. Where positions are read}

\begin{description}
\item[(I.POS)] Every function of the position, every membership in a set of sites, every support and every distance in an exponent is read at the photon site $c=\pi(x)$.
\end{description}
In particular, for $x=(c,g)$ and $x'=(c',g')$:
\begin{itemize}
\item $\operatorname{supp}x:=\{c\}$;
\item $x$ lies in, or over, a set $Z$ of photon sites (a block, an $M$-cube, a polymer, a ball, a union of cubes) if and only if $c\in Z$;
\item $\square(x)$ is the $M$-cube of $c$, and $K_x:=K_c$, $[x]:=[c]$, $d(x,x'):=d(c,c')$, $|x-x'|:=|c-c'|$;
\item a function $\psi$ of the position acts on spinors by $(\psi f)(x):=\psi(c)\cdot f(x)$; this applies to the gauge functions $\boldsymbol{\lambda}$, to the cut-off functions $h_\square$ and to the exponential weights;
\item $\mathbf{1}_Z$ is the indicator of $Z\times O_{\mathrm{cub}}$; in particular $\mathbf{1}_{\Delta_y}$ is the indicator of $\Delta_y\times O_{\mathrm{cub}}$.
\end{itemize}
By the paragraph on faces in Section~I.1.1, the sets $Z\times O_{\mathrm{cub}}$, where $Z$ ranges over the blocks of one side (over the $M$-cubes, over the $LM$-cubes), form a partition of $\mathbb{L}^{\mathrm{f}}_k$:
no block, cube or polymer cuts a cluster, and only bonds of type $3$ join two different blocks or cubes.
One has $\#\left|\Delta_y\times O_{\mathrm{cub}}\right|=48\cdot\eta_k^{-3}$.
Only one distance distinguishes two sites that differ in the label; it is defined in Section~I.1.7 and occurs in H\"older quotients only.

\subsubsection*{I.1.5. Spinors, boundary condition}

For $a\in\{\eta_k,1,L\}$ let
\begin{gather*}
\mathfrak{S}(a,\ell):=\left\{f:a\mathbb{Z}^3\to\mathbb{C}^4:\ f(c+\ell e_\mu)=-f(c)\ \text{for}\ \mu=1,2,3\right\},\\
\mathfrak{S}^{\mathrm{ch}}(a,\ell):=\left\{f:a\mathbb{Z}^3\times O_{\mathrm{cub}}\to\mathbb{C}^4:\ f(c+\ell e_\mu,g)=-f(c,g)\ \text{for}\ \mu=1,2,3\right\}.
\end{gather*}
The elements of $\mathfrak{S}(a,\ell)$ have $4$ components, and $\dim\mathfrak{S}(a,\ell)=4\cdot\left(\ell/a\right)^3$; moreover $\dim\mathfrak{S}^{\mathrm{ch}}(a,\ell)=192\cdot\left(\ell/a\right)^3$.

\paragraph{One norm on $\mathbb{C}^4$.}
The numbered properties carry the names of statements of \cite{QEDI} that are listed in Section~I.9.1:
W1--W18 are the properties (1)--(18) of a system of primitives in \cite[Section~5.4.2]{QEDI}, FE1--FE8 are \cite[Propositions~5.42--5.49]{QEDI}, and AC8 is \cite[Proposition~5.65]{QEDI}.
The upper index $\mathrm{ch}$, as in AC8$^{\mathrm{ch}}$, marks the form of a property for the chamber operator (Section~I.9.1), and Section~I.5.8 states how W1--W18 and FE1--FE8 are read for the chamber fermion.
The norm $|v|_c:=\max_\beta|v_\beta|$ is the norm on $\mathbb{C}^4$ in all numbered properties and in all sup norms:
$|f|_\infty:=\sup_x|f(x)|_c$, and the norm $|m|_c$ of a $4\times4$ block $m$ is the operator norm for $|\cdot|_c$, that is, the largest row sum of the moduli of its entries;
every entry of $m$ has modulus $\le|m|_c$.
The Euclidean norm is written $|v|_2$, and the corresponding operator norm $|m|_2$; one has
\[
|v|_c\le|v|_2\le2|v|_c .
\]
If a linear operator has norm at most $C$ for the sup norms formed with $|\cdot|_2$, then its norm for the sup norms formed with $|\cdot|_c$ is at most $2C$ \cite[Section~16.2.3]{QEDI};
accordingly, a bound that is proved with a constant $C$ in the Euclidean norm is stated with the constant $2C$ in the norm $|\cdot|_c$.
$L^2$ norms are formed with $|\cdot|_2$.
The one exception to the use of $|\cdot|_c$ in the numbered properties is the bound of the blocks $D_k(A)^{-1}(y,y')$ in AC8$^{\mathrm{ch}}$, that is, Proposition~19\,(e), for a real bond function $A$, $k\ge1$ and $\mu\le12\mathrm{e}$; it is stated in the norm $|\cdot|_2$:
\[
\left|D_k(A)^{-1}(y,y')\right|_2\le b_k^{-1}\cdot\mathbf{1}[y=y']+\frac{2\mathrm{e}}{\mu}\cdot\exp\left(-\frac{\mu}{12\mathrm{e}}\cdot d(y,y')\right),
\]
with the numbers $2\mathrm{e}$ and $12\mathrm{e}$ of \cite[Proposition~5.65]{QEDI}.
For rows and columns one has: the $\ell^1$ norm of a row or of a column of a $4\times4$ matrix $m$ is at most $2|m|_2$, by the Cauchy--Schwarz inequality, since the Euclidean norm of a row or of a column of $m$ is at most $|m|_2$.
The constants $C_5$, $C_{00}$, $a'$, $C'_V$ and the numerals of \cite[Sections~16.5--16.9]{QEDI} are constants for the Euclidean norm.
Accordingly, the Propositions~1.1--1.7, 2.1--2.8, 3--6 and 9\,(c), and the walk sums in Proposition~3, are stated with the Euclidean sup norm $|\cdot|_{\mathrm{su}}$ (Section~G1.3), while every clause that is a numbered property (a property W$n$ or FE$n$) is stated with $|\cdot|_c$ and the doubled constant ($2C_1$, $C^F$, $C_T$).

Spinors are antiperiodic in $c$ in the three directions, for every label.
A spinor on $T^{\mathrm{f}}_k$ is an element of $\mathfrak{S}^{\mathrm{ch}}(\eta_k,\ell)$.
At level $0$, in lattice units, $\mathfrak{S}^{\mathrm{ch}}(1,\ell)$ is, through the map $\mathrm{geo}$ (Proposition~46\,(e)), the space of the spinors of the cubic torus $C(\ell)=C(\ell,\ell,\ell)$ of \cite[Definition~2.11]{QEDII}, whose character satisfies $\chi(\ell e_\mu)=-1$ for $\mu=1,2,3$.
On the tori of the flow (Section~I.6.2), $\ell=\ell_0(R)=R/\eta$ is the number $n$ of Section~I.6.1.
Functions of pairs are introduced in Section~I.1.7.
On the auxiliary block tori of \cite[Definition~16.1]{QEDI} the sign $-1$ is replaced by $\boldsymbol{\sigma}\in\{1,-1\}$, and the lattice may be shifted;
the label is a spectator: $f(c+Re_\mu,g)=\boldsymbol{\sigma}\cdot f(c,g)$.

\subsubsection*{I.1.6. Measure, pairings, kernels}

The measure on the fermion sites is
\begin{equation}
\int\mathrm{d}x\,F(x):=\eta_k^3\cdot\frac{1}{48}\sum_{c}\ \sum_{g\in O_{\mathrm{cub}}}F(c,g),
\tag{I.1.3}
\end{equation}
where $c$ ranges over a fundamental domain of $\ell\mathbb{Z}^3$ in $\mathbb{L}_k$.
It is the measure on the fermion sites of the fine lattice at every level $k\ge0$, and at level $0$ it is also the measure on ``the unit lattice'' $T^0_0\times O_{\mathrm{cub}}=T^{\mathrm{f}}_0$, since $\eta_0=1$.
For a set $Z$ of photon sites, $\int_Z\mathrm{d}x$ denotes the integral over $Z\times O_{\mathrm{cub}}$.
Thus $\int_X\mathrm{d}x=\operatorname{Vol}(X)$ for a union $X$ of $M$-cubes, and $\int_{\Delta_y}\mathrm{d}x=1$ for every unit site $y$.
On the spaces of four-component spinors the measures are those of \cite{QEDI}: $\eta_k^3\sum_c$ on $\mathfrak{S}(\eta_k,\ell)$, $\sum_y$ on $\mathfrak{S}(1,\ell)$, and $L^3\sum_{y'}$ on $\mathfrak{S}(L,\ell)$.

\emph{Pairings.}
For $f$, $f'$ in $\mathfrak{S}^{\mathrm{ch}}(\eta_k,\ell)$ put
\[
\langle f',f\rangle^{\mathrm{ch}}:=\int\mathrm{d}x\,f'(x)\cdot f(x),\qquad
(f',f)^{\mathrm{ch}}:=\int\mathrm{d}x\,\overline{f'(x)}\cdot f(x),
\]
where $f'\cdot f:=\sum_{\beta}f'_{\beta}f_{\beta}$; the first pairing is bilinear, the second sesquilinear.
For $f$, $f'$ in $\mathfrak{S}(a,\ell)$ put
\[
\langle f',f\rangle_a:=\sum a^3\,f'\cdot f,\qquad
(f',f)_a:=\left\langle\overline{f'},f\right\rangle_a,
\]
again a bilinear and a sesquilinear pairing, the sum extending over a fundamental domain of $\ell\mathbb{Z}^3$ in $a\mathbb{Z}^3$.
In each of these pairings the summand is $\ell\mathbb{Z}^3$-periodic, so that the choice of the fundamental domain is immaterial.
Transposes $O^T$ and adjoints $O^{*}$ of operators are always taken with respect to these pairings.

\emph{Kernels.}
Let $O$ be an operator whose domain is a space of spinors on fermion sites.
Its kernel $O(t,x')$, a $4\times4$ matrix, is always taken with respect to the measure (I.1.3):
\begin{equation}
(Of)(t)=\int\mathrm{d}x'\,O(t,x')\,f(x'),
\tag{I.1.4}
\end{equation}
where $x'$ ranges over a fundamental domain, and $O(t,x'+\ell e_{\mu})=-O(t,x')$.
If the domain of $O$ is $\mathfrak{S}(1,\ell)$, the kernel is defined as in \cite[Section~16.2.3]{QEDI}:
\[
(Ov)(t)=\sum_{y'}O(t,y')\,v(y'),
\]
with $y'$ ranging over a fundamental domain of $\ell\mathbb{Z}^3$ in $\mathbb{Z}^3$, and $O(t,y'+\ell e_{\mu})=-O(t,y')$.
The following facts are consequences of these definitions.
\begin{itemize}
\item[(I.ke1)] The kernel of the identity $1$ on $\mathfrak{S}^{\mathrm{ch}}(\eta_k,\ell)$ is $48\,\eta_k^{-3}\,\mathbf{1}[x=x']\,1_4$.
The kernel of the multiplication by a matrix-valued function $\mathrm{Mx}(x)$ is $48\,\eta_k^{-3}\,\mathbf{1}[x=x']\,\mathrm{Mx}(x)$.
\item[(I.ke2)] $\operatorname{Tr}O=\int\mathrm{d}x\,\operatorname{tr}O(x,x)$.
\item[(I.ke3)] The norm of $O$ with respect to the sup norms of Section~I.1.5 is
\[
|O|_{1,\infty}:=\sup_t\ \sup_{\beta}\ \sum_{\beta'}\int\mathrm{d}x'\,\left|O(t,x')_{\beta\beta'}\right|,
\]
and it contains no factor $48$.
If the domain of $O$ is $\mathfrak{S}(1,\ell)$, one puts
\[
|O|_{1,\infty}:=\sup_t\ \sup_{\beta}\ \sum_{y'}\sum_{\beta'}\left|O(t,y')_{\beta\beta'}\right|,
\]
with $y'$ over a fundamental domain, every unit site carrying the weight $1$.
\item[(I.ke4)] Let $k=0$.
If the norm of $\mathbf{1}_{\{c\}}\,O\,\mathbf{1}_{\{c'\}}$ with respect to the sup norms is at most $a$, then $|O(x,x')|_c\le48\,a$ for every $x$ over $c$ and every $x'$ over $c'$.
Thus, at level $0$, a pointwise bound on a kernel contains the factor $48$ of the inverse weight, whereas a bound on row sums does not.
\item[(I.ke5)] (Kernel norm; level $0$ only.)
Let $\gamma\ge0$, and let $O$ be an operator on $\mathfrak{S}^{\mathrm{ch}}(1,\ell)$ or on the spinors of a block torus of level $0$.
Define
\[
N^{\mathrm{ke}}_{\gamma}(O):=\sup_z\sum_{z'}\exp\left(\gamma\,d(z,z')\right)\cdot\sup\left\{|O(x,x')|_2:\ x\text{ over }z,\ x'\text{ over }z'\right\}
\]
and
\[
N_{\gamma}(O):=\sup_z\sum_{z'}\exp\left(\gamma\,d(z,z')\right)\cdot\left\|\mathbf{1}_{\{z\}}\,O\,\mathbf{1}_{\{z'\}}\right\|,
\]
where $\|\cdot\|$ is the operator norm with respect to the sup norms formed with $|\cdot|_2$.
Then
\begin{gather*}
N_{\gamma}(O)\le N^{\mathrm{ke}}_{\gamma}(O),\qquad
N^{\mathrm{ke}}_{\gamma}(O_1O_2)\le N_{\gamma}(O_1)\cdot N^{\mathrm{ke}}_{\gamma}(O_2),\qquad
N^{\mathrm{ke}}_{\gamma}(\psi\,O\,\psi')\le|\psi|_{\infty}\cdot|\psi'|_{\infty}\cdot N^{\mathrm{ke}}_{\gamma}(O),\\
|O(x,x')|_2\le\exp\left(-\gamma\,d(x,x')\right)\cdot N^{\mathrm{ke}}_{\gamma}(O),\qquad
N^{\mathrm{ke}}_{\gamma}(1)=48,
\end{gather*}
the third inequality holding for $\ell\mathbb{Z}^3$-periodic functions $\psi$, $\psi'$ of $c$ (on a block torus: for $R\mathbb{Z}^3$-periodic ones).
At level $0$, Neumann series and sums over walks are estimated in $N^{\mathrm{ke}}_{\gamma}$; in each product exactly one factor, namely the last inverse at zero field, is measured by its kernel.
For this reason the factor $48$ appears only once, in the constant $C_{00}$ (Section~I.8.2).
\end{itemize}

\begin{named}{Remark}
In \cite[Section~16.2.3]{QEDI} the kernel is defined for every linear operator $O$ from the spinors on one of the three lattices of a block torus, the fine lattice, the unit lattice and the $L$-lattice, to the spinors on one of them, by
\[
(Of)(x)=\sum_{x'}O(x,x')\,f(x'),
\]
where the sum extends over a period cell and carries no weight.
The formulas $(Of)(y)=\sum_{y'}O(y,y')\,f(y')$ and $(Of)(x)=\sum_{y'}O(x,y')\,f(y')$ given there are the cases of an operator on the spinors on the unit lattice and of an operator from these spinors to the spinors on the fine lattice.
Consequently, for an operator acting on the spinors on the fine lattice, the weight $\eta^3$ of that lattice is part of the kernel of \cite[Section~16.2.3]{QEDI}.
The same holds in \cite[Section~19.3.1]{QEDI}, where the factors $\eta_k^3$ that weight the sums over the fine lattice are part of the kernel, $(Of)(x)=\sum_{x'}O(x,x')\,f(x')$; on the torus $T_k$, for the operators between the spaces of spinors on its fine lattice and on its unit lattice, the kernel of \cite[Section~16.2.3]{QEDI} is identified there with the kernel of \cite[Section~19.3.1]{QEDI}.
In \cite[Section~12.1]{QEDI}, on the other hand, the weight stands outside the kernel of an operator from spinor functions on a fine lattice to spinor functions on a fine lattice, $(Of)(x)=\eta^3\sum_{x'}O(x,x')\,f(x')$.
The kernel (I.1.4) has the weight outside.
A kernel in the normalisation of \cite[Sections~16.2.3 and~19.3.1]{QEDI}, that is, the coefficient of $f(x')$ in $(Of)(t)$, equals $\frac{\eta_k^3}{48}$ times the kernel (I.1.4).
\end{named}

\subsubsection*{I.1.7. The distance $d_X$; pairs; kinds of pairs}

On the fermion sites of level $k$ let
\begin{equation}
d_X(x,x'):=\max\left(|c-c'|,\eta_k\right)\qquad\left(x\ne x'\text{ in }\mathbb{L}^{\mathrm{f}}_k\right).
\tag{I.1.5}
\end{equation}
The point $v_0$ does not enter $d_X$, and $d_X$ is a distance on the covering lattice.
The \emph{fermion graph} of level $k$ is the graph with the vertex set $\mathbb{L}^{\mathrm{f}}_k$ (on the torus, $T^{\mathrm{f}}_k$) whose edges are the bonds $\left\{x,x^{(i)}\right\}$, $i=0,1,2,3$, of Section~I.1.3.
A path of $n$ bonds is a sequence $x^0,\dots,x^n$ of fermion sites with $x^{j+1}=\left(x^j\right)^{(i_j)}$ and $i_j\in\{0,1,2,3\}$ for $0\le j<n$.
\begin{itemize}
\item[(I.s6)] Two sites $x\ne\tilde x$ of $\mathbb{L}^{\mathrm{f}}_k$ are joined by a path of at most $39\,d_X(x,\tilde x)/\eta_k$ bonds of the fermion graph whose photon sites lie in the lattice box spanned by $c$ and $\tilde c$.
A proof is given in Chapter~1b, Section~0.4.
\end{itemize}

Pairs are formed on the covering lattice.
Let
\[
\mathrm{Pzt}_k:=\left\{(x,x'):\ x,x'\in\mathbb{L}^{\mathrm{f}}_k,\ x\ne x',\ d_X(x,x')\le M\right\}
\]
be the set of the ordered pairs of distinct points of $\mathbb{L}^{\mathrm{f}}_k$ at distance at most $M$.
The group $\ell\mathbb{Z}^3$ acts on $\mathrm{Pzt}_k$ by $(x,x')\mapsto(x+v,x'+v)$, and $\mathrm{Pz}_k:=\mathrm{Pzt}_k/\ell\mathbb{Z}^3$.
Put $\ell_H:=M$.
Since $2M<\ell$, a pair of distinct sites of $T^{\mathrm{f}}_k$ whose photon sites are at distance at most $M$ has exactly one class of lifts in $\mathrm{Pzt}_k$, namely the class of its shortest lift.
A \emph{function of pairs} is a function $F$ on $\mathrm{Pzt}_k$ with
\[
F(x+\ell e_{\mu},x'+\ell e_{\mu})=-F(x,x');
\]
its modulus is a function on $\mathrm{Pz}_k$, and $F$ is localised in the $M$-cube of its first point $x$.
The support of a pair is $\operatorname{supp}(x,x'):=\{c,c'\}$.

\emph{Pairs inside a cluster:} $c=c'$ and $g\ne g'$.
For these $d_X=\eta_k$, and over every photon site there are $48\cdot47$ ordered pairs of this kind.
\emph{Short pairs:} $d_X(x,x')<1$.
At level $0$ there is no short pair, since $d_X\ge\eta_0=1$ for all pairs; for $k\ge1$ every pair inside a cluster is short.
\emph{Long pairs:} $1\le d_X(x,x')\le M$.

\emph{Kinds.}
Let $Z$ be a union of $M$-cubes.
A pair $(x,x')$ of $\mathrm{Pz}_k$ is \emph{admissible} for $Z$ if it is of one of the following two kinds.
\begin{itemize}
\item[(I)] The sites $c$ and $c'$ lie in $Z$, and so do the six paths $\boldsymbol{\Gamma}^{\pi}(c,c')$.
Here a path lies in $Z$ if all its sites do; for $c=c'$ the condition reads $c\in Z$.
\item[(II)] $x'=x^{(3)}$, and $c$ or $c_3$ lies in $Z$.
\end{itemize}
The set of the pairs admissible for $Z$ is denoted by $\mathrm{Pz}_k(Z)$.
Consequently every pair inside a cluster over $Z$, in particular every bond $(x,x^{(i)})$ with $i\le2$, is of kind (I); every bond of type $3$ that leaves $Z$ is of kind (II); and no other pair with a point outside $Z$ is admissible.
The two kinds are not exclusive: a pair $(x,x^{(3)})$ with $c$ and $c_3$ in $Z$ is of both kinds.

\begin{named}{Remark}
The pairs and their kinds take the place of those of \cite[Definition~5.22]{QEDI}.
There a pair consists of two photon sites $x$, $x'$ of $T_k$ with $0<|x-x'|\le\ell_H$, and kind (II) admits every pair of nearest neighbours one of which lies in the given set.
Here the points of a pair are fermion sites, pairs inside one cluster occur and fall under kind (I), and kind (II) is restricted to the bonds $(x,x^{(3)})$ of type $3$.
\end{named}

\begin{named}{Remark}
The pairs that occur in the construction are the pairs $(x,x^{(3)})$ in the two displays of the step (Section~I.5.9); pairs with equal labels in all links of the chain of the adjustment except the last; and pairs with arbitrary labels, possibly inside one cluster, in the last link of that chain (Section~I.5.6).
\end{named}

\subsubsection*{I.1.8. Running charge of a level; phases; H\"older quotients}

\emph{Scale-free setting.}
As in \cite[Section~5.4]{QEDI}, the data of a level are a quadruple $(k,\ell,\mu,e_k)$, consisting of an integer $k\ge0$, a side $\ell\in\mathrm{Cl}$, a mass $\mu\in(0,1]$ and a number $e_k\in(0,e_0]$, between which no relation is assumed.
(For the data along the flow see Section~I.6.1.)
For a bond function $B$ on $T_k$ define
\begin{equation}
\begin{aligned}
u_B(c,c')&:=\exp\left\{i\cdot e_k\cdot\eta_k\cdot B(c,c')\right\}\qquad\text{(bonds)},\\
U_B(c,c')&:=\exp\left\{i\cdot e_k\cdot\eta_k\cdot(\tau B)(c,c')\right\}\qquad\left(|c-c'|\le M\right),\qquad
U_B(c,c):=1,\\
U_B(x,x')&:=U_B\left(\pi(x),\pi(x')\right),\qquad
\mathfrak{i}(B):=e_k\cdot|\operatorname{Im}B|_{\infty}.
\end{aligned}
\tag{I.1.6}
\end{equation}
For nearest neighbours $c$, $c'$ each of the six paths is the bond itself, so that $U_B(c,c')=u_B(c,c')$.
Hence $U_B(x,x^{(i)})=1$ for $i=0,1,2$ and $U_B(x,x^{(3)})=u_B(c,c_3)$: one symbol serves both the pairs and the bonds of the fermion graph.

\emph{H\"older quotient.}
For $f\in\mathfrak{S}^{\mathrm{ch}}(\eta_k,\ell)$ and $(x,x')\in\mathrm{Pzt}_k$ let
\begin{equation}
(\delta_{\alpha,B}f)(x,x'):=d_X(x,x')^{-\alpha}\cdot\left(U_B(x,x')\cdot f(x')-f(x)\right);
\tag{I.1.7}
\end{equation}
on barred fields $\delta_{\alpha,-B}$ is used in its place.
In (I.1.7), $f(x)$ and $f(x')$ are the values of the antiperiodic function $f$ at the two points of the pair in $\mathbb{L}^{\mathrm{f}}_k$, and the paths are those between the photon sites of the same two points.
The quotient $\delta_{\alpha,B}f$ is a function of pairs in the sense of Section~I.1.7.
Inside a cluster
\[
(\delta_{\alpha,B}f)(x,x')=\eta_k^{-\alpha}\left(f(x')-f(x)\right),
\]
without a phase, and at a bond of type $3$
\[
\eta_k^{\alpha}\cdot(\delta_{\alpha,B}f)(x,x^{(3)})=u_B(c,c_3)\cdot f(x^{(3)})-f(x).
\]
For every pair $\left|U_B(x,x')^{\pm1}\right|\le\exp\left(3M\cdot\mathfrak{i}(B)\right)$, because a path has at most $3M/\eta_k$ bonds.
Hence, for a long pair,
\[
\left|(\delta_{\alpha,B}f)(x,x')\right|_c\le\left(1+\exp\left(3M\cdot\mathfrak{i}(B)\right)\right)\cdot|f|_{\infty}.
\]

\begin{named}{Remark}
The H\"older quotient (I.1.7) takes the place of the H\"older derivative of \cite[Definition~5.22]{QEDI}.
For a pair with $c\ne c'$ one has $d_X(x,x')=|c-c'|$, and (I.1.7) is the formula there, formed at the photon sites $c$, $c'$ with the values of $f$ at the fermion sites $x$, $x'$.
The pairs inside a cluster, with the factor $\eta_k^{-\alpha}$ and without a phase, have no counterpart in \cite[Definition~5.22]{QEDI}.
\end{named}

\emph{Names of the standing data.}
\begin{itemize}
\item[(SF)] $(k,\ell,\mu,e_k)$: $k\ge0$ an integer, $\ell\in\mathrm{Cl}$, $\mu\in(0,1]$, $e_k\in(0,e_0]$; no relation between them is assumed.
\item[(SF+)] (SF), together with $\ell/L\in\mathrm{Cl}$, $\ell\ge M\,r_k$ (for $r_k$ see Section~I.6.1) and $e_{k+1}:=L^{\frac{1}{2}}e_k\le e_0$.
\item[(OS-)] $k\ge0$; $\ell\in\mathrm{Cl}$ and $\ell/L\in\mathrm{Cl}$; $\mu\in(0,1/L]$; $e_k>0$ with $e_{k+1}:=L^{\frac{1}{2}}e_k\le e_0$.
No primitives are given.
\item[(OS)] (OS-), $\ell\ge M\,r_k$, and a system of primitives of level $k$ on $T_k$ with the properties W1--W17$^{\mathrm{ch}}$ (Section~I.5.8) for the parameters $\mathrm{par}$ (Section~I.8.2).
The level $k+1$ has the side $\ell/L$, the mass $L\mu$ and the charge $e_{k+1}$.
\item[(FL)] $(\bar{m},e,R,N)$: $\bar{m}\in(0,1]$; $e>0$ with $e\cdot h^{\frac{1}{2}}\le e_1$; $R\in\mathrm{Tor}(e)$; $N\ge N_{\min}(R,e)$ an integer; together with the primitives of Sections~I.6.2 and I.6.3, that is, the systems of Proposition~26.5 on the tori $T_k(R_{\mathrm{top}})$, $k<k_{\mathrm{top}}$, and their projections.
The primitives are formed from the quadruple $(\bar{m},e,R,N)$ alone and involve no further datum; by Proposition~26.6\,(e) they are defined under the four conditions.
Hence (FL) and the condition (S-4) below are the same condition on $(\bar{m},e,R,N)$.
\item[(ST*)$^{\mathrm{ch}}$] (FL) with $R=R^{*}(e)$ and $N\ge N_0(e)$, together with the pair $P_N(e)$.
\end{itemize}

\emph{Dictionary.}
Further hypotheses on the data carry the following names.
\begin{itemize}
\item[(S-1)] The hypothesis (SF) without the mass.
\item[(S-2)] The hypothesis (S-1), together with $\ell/L\in\mathrm{Cl}$ and $e_{k+1}\le e_0$.
\item[(S-3)] The hypothesis (SF) with $k\ge1$.
\item[(S-4)] The hypothesis (FL) without the primitives.
\item[(S22)] The hypothesis (OS-).
\item[(ST1)$^{\mathrm{ch}}$] The hypothesis (OS) with W1--W8, W10--W17 in place of W1--W17; W9 is required at no level.
At $k=0$ the row W17 is taken without the property W9 that it states for the families of level $0$, that is, with W8, W10, W11 in place of its W8--W11 and with W6--W8 in place of its W6--W9, in the form stated in Proposition~48.6\,(a).
Hence, at every level, (OS) implies (ST1)$^{\mathrm{ch}}$, but not conversely.
\item[(S21)] $k\ge0$, $\ell\in\mathrm{Cl}$, $\mu$ complex, $e_k>0$.
\item[(TT)] The hypothesis (FL) with $R<R_{\mathrm{top}}$, together with a level $k\le k_{*}(R)$.
\end{itemize}
Finally, the label (ST) denotes the table of orientation signs of \cite[Lemma~12.3]{QEDII} (see Section~I.4.7).

\subsubsection*{I.1.9. ``Given by one formula on $\mathbb{L}^{\mathrm{f}}_k$''}

\emph{Index sets.}
An index set of level $k$ is a set $\mathrm{Ind}$ on which the translations by the vectors of $LM\mathbb{Z}^3$ act, together with a map $\operatorname{supp}$ from $\mathrm{Ind}$ to the finite subsets of $\mathbb{L}_k$ such that $\operatorname{supp}(t+v)=\operatorname{supp}t+v$.
The index sets that occur are the following; each of them may carry in addition a spinor index $\beta\in\{1,\dots,4\}$ and a bar index $\omega\in\{0,1\}$.
\begin{itemize}
\item[(I.i1)] The cells (sites, bonds, plaquettes, cubes) of $\mathbb{L}_k$, of $\mathbb{Z}^3$ and of $L\mathbb{Z}^3$.
The support of a site is the site itself, that of a bond the set of its two end points \cite[Section~5.2.3]{QEDI}, and that of a plaquette or of a cube the set of its corners.
As the index set of the bond functions, which are odd under reversal, one takes the positively oriented bonds, with $f(c',c):=-f(c,c')$.
\item[(I.i2)] The basis elements $\Upsilon$ of \cite[Section~5.2.5]{QEDI}, with $\operatorname{supp}\Upsilon$ the set of the end points of the bonds at which $\Upsilon$ does not vanish (\cite[Section~5.2.3]{QEDI}, Section~G4.0.4).
\item[(I.i3)] The pairs $(y',y)$ with $y'\in L\mathbb{Z}^3$ and $y\in B(y')$, with $\operatorname{supp}(y',y)=\{y',y\}$, and the pairs $(y,c)$ with $y\in\mathbb{Z}^3$ and $c\in\Delta_y$, with $\operatorname{supp}(y,c)=\{y,c\}$.
\item[(I.i4)] $\mathbb{L}^{\mathrm{f}}_k$, with $\operatorname{supp}x=\{c\}$.
\item[(I.i5)] $\mathrm{Pzt}_k$, with $\operatorname{supp}(x,x')=\{c,c'\}$.
\end{itemize}
The translations act on functions by $(T_vf)(t):=f(t-v)$.
On the torus of side $\ell$ the function spaces are the $\ell\mathbb{Z}^3$-periodic functions on $\mathrm{Ind}$ (photon objects) and the antiperiodic ones, $f(t+\ell e_{\mu})=-f(t)$ (spinor objects; for pairs both points are translated).
Forms enter through their kernels, and constant matrices enter as operators that act at each site.

\begin{named}{Definition}
Let $\mathrm{Op}$ be a linear operator between such function spaces, defined for every side $\ell$ of a class of multiples of $LM$; it may depend on a bond function $B$ on $T_k$ and on a union $Y$ of $M$-cubes of $T_k$.
The operator $\mathrm{Op}$ is \emph{given by one formula on $\mathbb{L}^{\mathrm{f}}_k$, of range $\mathrm{rg}<\infty$}, if there are linear operators $\mathrm{Op}_{\mathbb{L}}(\tilde B,\tilde Y)$ on the functions on the index sets, one for every bond function $\tilde B$ on $\mathbb{L}_k$ and every union $\tilde Y$ of $M$-cubes of $\mathbb{L}_k$, with the following properties.
\begin{itemize}
\item[(I.f1)] Every row of $\mathrm{Op}_{\mathbb{L}}(\tilde B,\tilde Y)$ has finite support.
\item[(I.f2)] For all $v\in LM\mathbb{Z}^3$,
\[
T_v\,\mathrm{Op}_{\mathbb{L}}(\tilde B,\tilde Y)=\mathrm{Op}_{\mathbb{L}}(T_v\tilde B,\tilde Y+v)\,T_v .
\]
\item[(I.f3)] For every $\ell$ of the class, $\mathrm{Op}(B,Y)$ is the restriction of $\mathrm{Op}_{\mathbb{L}}(\tilde B,\tilde Y)$ to the periodic, respectively antiperiodic, functions, where $\tilde B$ and $\tilde Y$ are the $\ell\mathbb{Z}^3$-periodic extensions of $B$ and $Y$.
\item[(I.f4)] (Locality.)
For all $\tilde B$, $\tilde Y$ and every row index $t$, the row at $t$ vanishes at the columns $t'$ with $\operatorname{dist}(\operatorname{supp}t,\operatorname{supp}t')>\mathrm{rg}$, and it reads $\tilde B$ on the bonds, and $\tilde Y$ through the sites, within the distance $\mathrm{rg}$ of $\operatorname{supp}t$.
\end{itemize}
Here $\operatorname{dist}(\operatorname{supp}t,\operatorname{supp}t')$ is the smallest sup-distance of a point of $\operatorname{supp}t$ from a point of $\operatorname{supp}t'$, as in \cite[Section~5.2.3]{QEDI}; the word ``reads'' has the meaning fixed in \cite[Definition~5.5\,(1)]{QEDI}; and a bond lies within the distance $\mathrm{rg}$ of a set of sites if both of its end points have sup-distance at most $\mathrm{rg}$ from that set.
\end{named}

\begin{named}{Consequence}
Let $\ell_1$ and $\ell_2$ belong to the class, and for $i=1,2$ let $(B_i,Y_i)$ be data on the torus of side $\ell_i$ whose lifts are equal within the distance $\mathrm{rg}$ of $\operatorname{supp}t$.
Then the rows at $t$ of $\mathrm{Op}_{\mathbb{L}}$ for the two periodic extensions coincide.
If moreover $t$ and $t'$ are contained in a set $U$ whose sup-diameter is smaller than $\ell_i-\mathrm{rg}$, then the kernel of $\mathrm{Op}(B_i,Y_i)$ at $(t,t')$ equals $\mathrm{Op}_{\mathbb{L}}(t,t')$, for $i=1$ and for $i=2$.
\end{named}

Here, for an operator $\mathrm{Op}_{\mathbb{L}}$ with rows of finite support on the functions on the covering index sets, $\mathrm{Op}_{\mathbb{L}}(t,t')$ denotes its kernel, namely the coefficient of $f(t')$ in $(\mathrm{Op}_{\mathbb{L}}f)(t)$ divided by the weight of the index $t'$.
The weights are $\eta_k^3/48$ on $\mathbb{L}^{\mathrm{f}}_k$, $\left(\eta_k^3/48\right)^2$ on $\mathrm{Pzt}_k$, $\eta_k^3$ on the sites of $\mathbb{L}_k$, $1$ on $\mathbb{Z}^3$, and $L^3$ on $L\mathbb{Z}^3$ as the index set of $\mathfrak{S}(L,\ell)$ (Sections~I.1.6 and I.3.6); they are $1$ on the unit bonds and on the basis elements (Section~I.5.8, the definition of $c_{\gamma}$).
On the bonds of the fine lattices no weight is fixed; for an operator with columns on such bonds a weight on them is chosen together with the operator, and its kernel refers to that weight.
Nor is a weight fixed on the plaquettes, the cubes and the pairs (I.i3); these never occur as columns.
For the domain $\mathbb{L}^{\mathrm{f}}_k$ this kernel is the kernel (I.1.4), with the integral extended over all of $\mathbb{L}^{\mathrm{f}}_k$.

\begin{proof}[Proof of the Consequence]
For $i=1,2$ let $\tilde B_i$ and $\tilde Y_i$ be the $\ell_i\mathbb{Z}^3$-periodic extensions of $B_i$ and $Y_i$.
By (I.f4), for every bond function $\tilde B$ on $\mathbb{L}_k$ and every union $\tilde Y$ of $M$-cubes of $\mathbb{L}_k$, the row at $t$ of $\mathrm{Op}_{\mathbb{L}}(\tilde B,\tilde Y)$ reads $\tilde B$ on the bonds, and $\tilde Y$ through the sites, within the distance $\mathrm{rg}$ of $\operatorname{supp}t$.
In the sense of \cite[Definition~5.5\,(1)]{QEDI} this means that, the other datum being held fixed, the row takes equal values at any two bond functions that coincide on these bonds, and at any two unions of $M$-cubes that have the same intersection with the set of these sites.
By hypothesis $\tilde B_1=\tilde B_2$ on these bonds, and $\tilde Y_1$ and $\tilde Y_2$ have the same intersection with these sites.
Replacing first $\tilde B_1$ by $\tilde B_2$, with $\tilde Y_1$ held fixed, and then $\tilde Y_1$ by $\tilde Y_2$, one finds that the rows at $t$ of $\mathrm{Op}_{\mathbb{L}}(\tilde B_1,\tilde Y_1)$ and of $\mathrm{Op}_{\mathbb{L}}(\tilde B_2,\tilde Y_2)$ are equal.
This is the first assertion.

For the second assertion fix $i$, put $\ell:=\ell_i$, and write $\mathrm{Op}_{\mathbb{L}}$ for $\mathrm{Op}_{\mathbb{L}}(\tilde B_i,\tilde Y_i)$.
Let $f$ belong to the function space of the torus of side $\ell$ on the index set of the columns.
Since $\ell$ is a multiple of $LM$, the translations by the vectors $\ell n$, $n\in\mathbb{Z}^3$, act on this index set, and
\[
f(t''+\ell n)=\epsilon(n)\,f(t'')\qquad\left(n\in\mathbb{Z}^3\right),
\]
where $\epsilon(n):=1$ for the periodic and $\epsilon(n):=(-1)^{n_1+n_2+n_3}$ for the antiperiodic functions.
Denote by $\varpi(t'')$ the weight of a column index $t''$.
By (I.f3) and (I.f1),
\[
\left(\mathrm{Op}(B_i,Y_i)f\right)(t)=\left(\mathrm{Op}_{\mathbb{L}}f\right)(t)=\sum_{t''}\varpi(t'')\cdot\mathrm{Op}_{\mathbb{L}}(t,t'')\cdot f(t''),
\]
a finite sum over the column indices $t''$ of the covering index set.
Let $\mathscr{F}$ be a fundamental domain of $\ell\mathbb{Z}^3$ in this index set, and collect the columns that are congruent modulo $\ell\mathbb{Z}^3$, writing $t''=t'+\ell n$ with $t'\in\mathscr{F}$ and $n\in\mathbb{Z}^3$.
The weights are invariant under translations, $\varpi(t'+\ell n)=\varpi(t')$; it follows that
\[
\left(\mathrm{Op}(B_i,Y_i)f\right)(t)=\sum_{t'\in\mathscr{F}}\varpi(t')\cdot K(t,t')\cdot f(t'),\qquad
K(t,t'):=\sum_{n\in\mathbb{Z}^3}\epsilon(n)\,\mathrm{Op}_{\mathbb{L}}(t,t'+\ell n).
\]
For every column index $t'$ only finitely many terms of the last sum differ from zero, by (I.f1), and $K(t,t'+\ell n')=\epsilon(n')\,K(t,t')$ for $n'\in\mathbb{Z}^3$, because $\epsilon(n-n')=\epsilon(n)\,\epsilon(n')$.
Hence $K$ is the kernel of $\mathrm{Op}(B_i,Y_i)$ in the sense of (I.1.4), formed with the columns over a fundamental domain and extended periodically, respectively antiperiodically, and the formula for $K(t,t')$ holds at every column index $t'$.

Now let $\operatorname{supp}t$ and $\operatorname{supp}t'$ be contained in $U$, and let $n\ne0$.
The points of $\operatorname{supp}(t'+\ell n)=\operatorname{supp}t'+\ell n$ are the points $c'+\ell n$ with $c'\in\operatorname{supp}t'$.
For $c\in\operatorname{supp}t$ one has, with the sup norm $|\cdot|_{\infty}$ of $\mathbb{R}^3$,
\[
\left|c-c'-\ell n\right|_{\infty}\ge\ell\cdot|n|_{\infty}-\left|c-c'\right|_{\infty}\ge\ell-\left|c-c'\right|_{\infty}>\mathrm{rg},
\]
since $|c-c'|_{\infty}$ is at most the sup-diameter of $U$, which is smaller than $\ell-\mathrm{rg}$.
Hence $\operatorname{dist}\left(\operatorname{supp}t,\operatorname{supp}(t'+\ell n)\right)>\mathrm{rg}$, and $\mathrm{Op}_{\mathbb{L}}(t,t'+\ell n)=0$ by (I.f4).
Only the term $n=0$ remains, so that the kernel of $\mathrm{Op}(B_i,Y_i)$ at $(t,t')$ equals $\mathrm{Op}_{\mathbb{L}}(t,t')$.
By the first assertion the rows at $t$ of $\mathrm{Op}_{\mathbb{L}}(\tilde B_1,\tilde Y_1)$ and of $\mathrm{Op}_{\mathbb{L}}(\tilde B_2,\tilde Y_2)$ coincide, and so do their kernels at $(t,t')$; thus this kernel is the same for $i=1$ and for $i=2$.
\end{proof}

\emph{Objects attached to a polymer.}
Consider a family $O_{\ell}(X;\mathrm{data})$ of numbers, kernels, elements of Grassmann algebras, linear functionals on kernels, integrations, or homomorphisms of Grassmann algebras.
It is indexed by the sides $\ell$ of the class, by the polymers $X$ of $T_k$ that have a lift $\tilde X$ of sup-diameter smaller than $\ell-2\,\mathrm{rg}-M$, and by data, which are functions of one or of several indices within the distance $\mathrm{rg}$ of $X$ (antisymmetric kernels included) and of auxiliary parameters $(y,t)$.
The family is \emph{given by one formula of range $\mathrm{rg}$} if there is one map $\Phi$, the same for all $\ell$, defined on the pairs consisting of a set $\tilde X$ and of data on the $\mathrm{rg}$-neighbourhood of $\tilde X$, where $\tilde X$ ranges over the non-empty finite unions of $M$-cubes of $\mathbb{L}_k$ that are connected through faces, such that
\[
\Phi\left(\tilde X+v,\,T_v\,\mathrm{data}\right)=T_v\,\Phi\left(\tilde X,\mathrm{data}\right)\quad\left(v\in LM\mathbb{Z}^3\right),
\qquad
O_{\ell}(X;\mathrm{data})=\Phi\left(\tilde X,\,\mathrm{data}\circ\mathrm{pr}\right),
\]
$\mathrm{pr}$ being the covering map.
The data composed with the covering map are understood as follows.
A bond function $B$ on $T_k$ becomes the function $b\mapsto B(\mathrm{pr}\,b)$ on the bonds within the distance $\mathrm{rg}$ of $\tilde X$.
A kernel, which by (I.AP) (Section~I.3.6) is a function on the covering index sets, antiperiodic in each argument, becomes its restriction to the indices over the $\mathrm{rg}$-neighbourhood $\tilde U$ of $\tilde X$.
An element $F$ of the algebra of the torus that occurs among the data, and whose kernels vanish unless all $\xi$ lie over $X$ and all $\boldsymbol{\zeta}$ are admissible for $X$, is replaced, when $\mathrm{rg}\ge\eta_k$, by the element of $\mathrm{Grc}_k(\tilde U)$ with these restricted kernels; $\mathrm{cov}_{\ell}$ maps this element to $F$.
For elements the second identity is read as
\[
O_{\ell}(X;\mathrm{data})=\mathrm{cov}_{\ell}\left(\Phi\left(\tilde X,\mathrm{data}\circ\mathrm{pr}\right)\right),
\]
and for homomorphisms and integrations as
\[
\mathrm{cov}_{\ell}\circ\Phi\left(\tilde X,\mathrm{data}\circ\mathrm{pr}\right)=O_{\ell}(X;\mathrm{data})\circ\mathrm{cov}_{\ell}.
\]

In this definition, for a finite set $\tilde U$ of sites of $\mathbb{L}_k$, $\mathrm{Grc}_k(\tilde U)$ is the Grassmann algebra over $\mathbb{C}$ with the independent generators $\psi(x,\beta,\omega)$ with $c\in\tilde U$, $\chi(x,x',\beta,\omega)$ with $(x,x')\in\mathrm{Pzt}_k$ and $c,c'\in\tilde U$, and $\Psi(u)$, $W(u)$, where $u=(y,\beta,\omega)$ with $y\in\tilde U\cap\mathbb{Z}^3$ if $k\ge1$, and $u=(x,\beta,\omega)$ with $c\in\tilde U$ if $k=0$.
For $\ell$ in the class such that the sup-diameter of $\tilde U$ is smaller than $\ell-M$, the map
\[
\mathrm{cov}_{\ell}:\mathrm{Grc}_k(\tilde U)\to\mathrm{Gr}[\psi,\chi,\Psi,W]
\]
into the algebra of the torus of side $\ell$ is the homomorphism that sends every generator to the generator of the torus with the same index, in the sense of (I.AP).
It is injective.
If $\tilde U$ is replaced by its translate by $\ell w$, with $w\in\mathbb{Z}^3$, the image of every generator is multiplied by $(-1)^{w_1+w_2+w_3}$; here $w$ denotes a vector of $\mathbb{Z}^3$, as in Section~I.6.3, and not the weight quadruple of Section~I.0.5.
The same rule holds at the level $k+1$ for the side $\ell/L$, $\ell$ in the class (Section~G3.5.0).
In the Propositions~27.m one has $\mathrm{Gr}_i(\tilde X)=\mathrm{cov}_{\ell_i}\left(\mathrm{Grc}_k(N(\tilde X))\right)$ for $k\ge1$; for $k=0$, $\mathrm{Gr}_i(\tilde X)$ is the image under $\mathrm{cov}_{\ell_i}$ of the subalgebra of $\mathrm{Grc}_0(N(\tilde X))$ generated by $\psi$, $\chi$, $W$; and $\mathrm{idn}[\tilde X]=\mathrm{cov}_{\ell_2}\circ\mathrm{cov}_{\ell_1}^{-1}$ on $\mathrm{Gr}_1(\tilde X)$.
For the scaling homomorphism $S_L$ of Section~I.3.7 the margin condition reads: the sup-diameter of $\tilde U$ is smaller than $\ell-LM$.
The translation $T_v$ acts trivially on numbers, on kernels by translating all arguments, and on homomorphisms by relabelling the generators.
For the scaling $S_L$ a translation by $v$ at level $k$ corresponds to the translation by $v/L$ at level $k+1$.

\begin{named}{Remark}
The definitions of this section take the place of \cite[Definition~5.5\,(2)]{QEDI}.
They agree with that clause in content, except in five points.
The notion of an index set with supports is the same as there, and the conditions (I.f1)--(I.f3) are the conditions (i)--(iii) there.
The index sets (I.i1), (I.i2) and the pairs $(y',y)$ of (I.i3) are among the index sets listed in \cite[Definition~5.5\,(2)]{QEDI}, and pairs of sites occur there as well: the pairs $(x,x')$ of sites of $\mathbb{L}_k$ with $0<|x-x'|\le M$ form an index set, the support of such a pair being $\{x,x'\}$.
The locality (I.f4) agrees with the requirement of locality of \cite[Definition~5.5\,(2)]{QEDI} in both of its parts: the row at $t$ vanishes at the columns $t'$ with $\operatorname{dist}(\operatorname{supp}t,\operatorname{supp}t')>\mathrm{rg}$, and it reads $\tilde B$ on the bonds both of whose end points have sup-distance at most $\mathrm{rg}$ from $\operatorname{supp}t$, and $\tilde Y$ on the sites at sup-distance at most $\mathrm{rg}$ from $\operatorname{supp}t$.
Thus $\mathrm{rg}$ is a range of the formula in the sense of that definition, and the Consequence above is the comparison of two tori stated there, with the same margin: the sup-diameter of $U$ plus the range is smaller than the side.
A family of objects attached to polymers is given by one formula there in the same way as here: by one map $\Phi$, the same for all sides of the class, defined on the pairs formed by a non-empty finite union $\tilde X$ of $M$-cubes of $\mathbb{L}_k$ connected through faces and by data on the $\mathrm{rg}$-neighbourhood of $\tilde X$, for the polymers that have a lift of sup-diameter smaller than $\ell-2\,\mathrm{rg}-M$.
The map $\Phi$ is covariant under the translations by the vectors of $LM\mathbb{Z}^3$, and $O_{\ell}(X;\mathrm{data})$ is the value of $\Phi$ at the lift $\tilde X$ and at the transported data, which are the restrictions of the data to the indices over the $\mathrm{rg}$-neighbourhood of $\tilde X$.
For the scaling homomorphism $S_L$ a translation by $v$ at level $k$ is matched there, as here, with the translation by $v/L$ at level $k+1$.
The five points in which the definitions differ are the following.
First, the fermion sites $\mathbb{L}^{\mathrm{f}}_k$, with $\operatorname{supp}x=\{c\}$, take here the place of the sites of $\mathbb{L}_k$ as the indices of the spinors; the pairs $\mathrm{Pzt}_k$ of fermion sites, with $\operatorname{supp}(x,x')=\{c,c'\}$ and including the pairs inside a cluster, take the place of the pairs of sites of $\mathbb{L}_k$; and the pairs $(y,c)$ with $y\in\mathbb{Z}^3$ and $c\in\Delta_y$ of (I.i3) are not among the index sets listed there.
Second, the kernel $\mathrm{Op}_{\mathbb{L}}(t,t')$ is the coefficient of $f(t')$ in $(\mathrm{Op}_{\mathbb{L}}f)(t)$ divided by the weight of the column index $t'$, whereas in \cite[Definition~5.5\,(2)]{QEDI} the weights are counted into the entries.
Third, linear functionals on kernels and integrations are admitted here as objects attached to a polymer, and elements $F$ of the algebra of the torus are admitted among the data when $\mathrm{rg}\ge\eta_k$.
Fourth, elements, homomorphisms and integrations are compared here through the algebra $\mathrm{Grc}_k(\tilde U)$ with the independent generators $\psi$, $\chi$, $\Psi$, $W$ and through the injective homomorphisms $\mathrm{cov}_{\ell}$, defined when the sup-diameter of $\tilde U$ is smaller than $\ell-M$; in \cite[Definition~5.5\,(2)]{QEDI} elements and homomorphisms are compared through the labelling of the generators of the torus by the indices contained in the $\mathrm{rg}$-neighbourhood of the lift.
Fifth, for the scaling homomorphism $S_L$ the margin condition that the sup-diameter of $\tilde U$ be smaller than $\ell-LM$ is stated here; no margin condition for $S_L$ is stated there.
\end{named}

\subsection*{I.2. Spin table, hopping matrices, the operator, mass and charge, the dictionary}

\subsubsection*{I.2.1. The spin table}

Let $\sigma_1$, $\sigma_2$, $\sigma_3$ be the Pauli matrices and $1_2$ the unit $2\times2$ matrix.
Define the $4\times4$ matrices
\begin{gather*}
\gamma_{\mu}:=\sigma_3\otimes\sigma_{\mu}=\operatorname{diag}\left(\sigma_{\mu},-\sigma_{\mu}\right)\quad(\mu=1,2,3),\qquad
\gamma\cdot v:=\sum_{\mu}v_{\mu}\gamma_{\mu},\\
\gamma_5:=\sigma_1\otimes1_2=\left(\begin{array}{cc}0&1_2\\1_2&0\end{array}\right),\qquad
\Sigma_3:=\sigma_3\otimes1_2=\operatorname{diag}(1,1,-1,-1),\\
\mathsf{C}:=i\gamma_2=i\,\sigma_3\otimes\sigma_2,\qquad
\varkappa:=\gamma_5\cdot\mathsf{C}=\sigma_2\otimes\sigma_2,\qquad
\mathrm{Cc}:=1_2\otimes\sigma_2=\operatorname{diag}\left(\sigma_2,\sigma_2\right).
\end{gather*}
Explicitly,
\[
\mathsf{C}=\left(\begin{array}{rrrr}0&1&0&0\\-1&0&0&0\\0&0&0&-1\\0&0&1&0\end{array}\right),\qquad
\varkappa=\left(\begin{array}{rrrr}0&0&0&-1\\0&0&1&0\\0&1&0&0\\-1&0&0&0\end{array}\right);
\]
so $\varkappa$ is the antidiagonal matrix with the entries $-1$, $1$, $1$, $-1$ from the upper right to the lower left corner.

Multiplying out the matrices one obtains the following relations.
\begin{itemize}
\item[(I.c1)] $\gamma_{\mu}$, $\gamma_5$ and $\Sigma_3$ are Hermitian; $\gamma_{\mu}\gamma_{\nu}+\gamma_{\nu}\gamma_{\mu}=2\delta_{\mu\nu}$; $\gamma_5\cdot\gamma_{\mu}=-\gamma_{\mu}\cdot\gamma_5$; $\gamma_5^2=\Sigma_3^2=1$; $(\gamma\cdot v)^2=|v|_2^2$ for real $v$.
\item[(I.c2)] $\gamma_1\gamma_2\gamma_3=i\,\Sigma_3$; $\Sigma_3$ commutes with $\gamma_1$, $\gamma_2$, $\gamma_3$ and anticommutes with $\gamma_5$; $i\,\gamma_5\cdot\Sigma_3=\sigma_2\otimes1_2$.
\item[(I.c3)] $\mathsf{C}$ is real and $\mathsf{C}^T=\mathsf{C}^{-1}=-\mathsf{C}$; moreover $-\gamma_{\mu}^T=\mathsf{C}^{-1}\cdot\gamma_{\mu}\cdot\mathsf{C}$ for $\mu=1,2,3$, and $-\gamma_5^T=\mathsf{C}^{-1}\cdot\gamma_5\cdot\mathsf{C}$.
\item[(I.c4)] $\varkappa$ is real and $\varkappa^T=\varkappa=\varkappa^{-1}$; moreover $\overline{\gamma_{\mu}}=\varkappa^{-1}\cdot\gamma_{\mu}\cdot\varkappa$ for $\mu=1,2,3$, $\overline{\gamma_5}=-\varkappa^{-1}\cdot\gamma_5\cdot\varkappa$, and $\varkappa\cdot\mathsf{C}=-\mathsf{C}\cdot\varkappa$.
\item[(I.c5)] $\mathsf{C}=i\,\Sigma_3\cdot\mathrm{Cc}$; $\mathrm{Cc}\cdot\gamma_{\mu}^T\cdot\mathrm{Cc}^{-1}=-\gamma_{\mu}$ for $\mu=1,2,3$, but $\mathrm{Cc}\cdot\gamma_5^T\cdot\mathrm{Cc}^{-1}=+\gamma_5$.
\item[(I.c6)] The eight matrices $\gamma_1$, $\gamma_2$, $\gamma_3$, $\gamma_5$, $\Sigma_3$, $\mathsf{C}$, $\varkappa$, $\mathrm{Cc}$ are monomial: every row and every column contains exactly one non-zero entry, and this entry has modulus one.
The matrices $\gamma_3$ and $\Sigma_3$ are diagonal.
\end{itemize}

\emph{Names in the three papers.}
Wherever \cite{QEDI} and \cite{QEDII} define the spin matrices, the matrices in one row of the following table are equal; Proposition~46\,(c) lists the places at which these papers define no matrix of a row.
The papers differ only in the names and in the numbering of the directions; no sign differs.

\medskip
\noindent
\begin{tabular}{|p{1.5cm}|p{2.9cm}|p{3.3cm}|p{2.5cm}|}
\hline
this paper & \cite[Section~5.2.4]{QEDI} and \cite[Sections~11, 13--20]{QEDI} & \cite[Sections~1--4, 6--10, 12, 21]{QEDI} & \cite[Definition~2.7]{QEDII} \\
\hline
$\gamma_1,\gamma_2,\gamma_3$ & $\gamma_0,\gamma_1,\gamma_2$ & $\gamma_1,\gamma_2,\gamma_3$ & $\gamma_1,\gamma_2,\gamma_3$ \\
\hline
$\gamma_5$ & $\gamma_3$ & $\hat{\Gamma}$ \cite[Sections~12.1, 21.1.1]{QEDI}, $\gamma_4$ \cite[Section~7.1.2]{QEDI} & $\gamma_5$ \\
\hline
$\Sigma_3$ & no letter; it is $-i\gamma_0\gamma_1\gamma_2$ & no letter & no letter; it is $+1$ on the first and $-1$ on the second summand of $\mathbb{C}^2\oplus\mathbb{C}^2$ \\
\hline
$\mathsf{C}$ & $\mathsf{C}=i\gamma_1$ & --- & --- \\
\hline
$\varkappa$ & $\varkappa=\gamma_3\cdot\mathsf{C}$ & --- & --- \\
\hline
$\mathrm{Cc}$ & --- & $\mathbf{C}$ \cite[Section~21.3.2]{QEDI}; $\mathsf{C}$ in the proof of \cite[Lemma~7.7\,(c)]{QEDI}, where this letter denotes $\operatorname{diag}(\sigma_2,\sigma_2)=\mathrm{Cc}=-i\,\Sigma_3\cdot\mathsf{C}$, not the matrix of the row $\mathsf{C}$ & $\sigma_2\oplus\sigma_2$ \\
\hline
\end{tabular}
\medskip

\emph{Choice.}
The charge conjugation matrix is $\mathsf{C}$.
It is the matrix that is needed together with $\gamma_5$: with $\mathsf{C}$ the resolvent parameter changes its sign under transposition (Section~I.4.5).
The matrix $\mathrm{Cc}$ is the charge conjugation matrix of \cite[Definition~2.7 and Proposition~3.16\,(c)]{QEDII}; the translation between $\mathsf{C}$ and $\mathrm{Cc}$ is given in Section~I.4.5.

\subsubsection*{I.2.2. Hopping matrices}

For $i=0,\dots,3$ and $x=(c,g)$ let
\begin{equation}
\begin{aligned}
M_{i}(x)&:=\frac{w_i}{2}\left(\rho-\gamma\cdot\hat{d}_i(x)\right)=\frac{w_i}{2}\left(\rho+\gamma\cdot(g\hat n_i)\right),\\
M_{3}(x)&=3\left(\rho-\gamma\cdot(ge_1)\right),\qquad
M_{3}(x^{(3)})=3\left(\rho+\gamma\cdot(ge_1)\right).
\end{aligned}
\tag{I.2.1}
\end{equation}
The matrix $M_{i}(x)$ depends only on the label $g$ of $x$; it contains neither the mesh nor the level.
It is the hopping matrix $w(x,x^{(i)})$ of \cite[Definition~2.10]{QEDII}.
It takes the place of the projections $\frac{1}{2}\left(1\mp\gamma_{\mu}\right)$ of \cite{QEDI}, denoted there $P^{\boldsymbol{\sigma}}_{\nu}$, $\Pi^{\mp}_{\mu}$ and $\mathrm{Pr}^{\boldsymbol{\sigma}}_{\mu}$, which do not occur here.
From (I.s2), (I.s3) and (I.c1)--(I.c5) one obtains the following facts.
\begin{itemize}
\item[(I.m1)] $\sum_iM_{i}(x)=\frac{\rho}{2}\cdot\mathrm{sw}$, a scalar; $M_{i}(x^{(i)})=\frac{w_i}{2}\left(\rho+\gamma\cdot\hat{d}_i(x)\right)$; $M_{i}(x^{(i)})-M_{i}(x)=w_i\,\gamma\cdot\hat{d}_i(x)$; $\sum_i\left(M_{i}(x^{(i)})-M_{i}(x)\right)=0$.
\item[(I.m2)] $\mathsf{C}^{-1}\cdot M_{i}(x)^T\cdot\mathsf{C}=M_{i}(x^{(i)})=\mathrm{Cc}\cdot M_{i}(x)^T\cdot\mathrm{Cc}^{-1}$.
\item[(I.m3)] $\overline{M_{i}(x)}=\varkappa^{-1}\cdot M_{i}(x)\cdot\varkappa$.
\item[(I.m4)] The covariance of the hopping matrices under the spin group is stated in Section~I.4.3.
\item[(I.m5)] $M_{i}(x)^{*}=M_{i}(x)$; $\gamma_5\cdot M_{i}(x)\cdot\gamma_5=M_{i}(x^{(i)})$; $\Sigma_3\cdot M_{i}(x)=M_{i}(x)\cdot\Sigma_3$.
\item[(I.m6)] The eigenvalues of $M_{i}(x)$ are $\frac{w_i}{2}(\rho-1)$ and $\frac{w_i}{2}(\rho+1)$, each with multiplicity two; $|M_{i}(x)|_2=\frac{w_i}{2}(1+\rho)$.
Since $\rho<1$, the matrix $M_{i}(x)$ is invertible; it is not a projection.
\item[(I.m7)] $|M_{0}(x)|_c=|M_{3}(x)|_c=3(1+\rho)\le6$, and $|M_{1}(x)|_c,|M_{2}(x)|_c\le3\sqrt{2}\left(\rho+\sqrt{2}\right)$.
The moduli of the sixteen entries of $M_{3}(x)$ add up to at most $12(1+\rho)\le24$: to $12(1+\rho)$ if $ge_1=\pm e_1$ or $ge_1=\pm e_2$, and to $12$ if $ge_1=\pm e_3$.
\end{itemize}

\subsubsection*{I.2.3. The operator}

Let $A$ be a complex bond function on $T_k$, and let $U_A$ be the phase of (I.1.6).
For $f\in\mathfrak{S}^{\mathrm{ch}}(\eta_k,\ell)$ put
\begin{align}
\left(R_i(A)f\right)(x)&:=U_A\left(x,x^{(i)}\right)\cdot f\left(x^{(i)}\right),\tag{I.2.2}\\
\left(D_Af\right)(x)&:=-\eta_k^{-1}\sum_{i=0}^{3}M_i(x)\cdot\left[U_A\left(x,x^{(i)}\right)\cdot f\left(x^{(i)}\right)-f(x)\right].\tag{I.2.3}
\end{align}
Here $U_A(x,x^{(i)})=1$ for $i=0,1,2$, and $U_A(x,x^{(3)})=u_A(c,c_3)=\exp\left(i\,e_k\eta_k\cdot A(c,c_3)\right)$.

The operator $D_A$ contains no mass.
By (I.m1) it can be written as
\begin{equation}
\begin{aligned}
&D_A=\eta_k^{-1}\left[\frac{\rho}{2}\,\nabla_A^{*}\nabla_A+\frac{1}{2}\,J_A\right],\qquad
\nabla_A^{*}\nabla_A:=\sum_iw_i\left(1-R_i(A)\right),\\
&\left(J_Af\right)(x):=\sum_iw_i\,\gamma\cdot\hat{d}_i(x)\cdot U_A\left(x,x^{(i)}\right)\cdot f\left(x^{(i)}\right).
\end{aligned}
\tag{I.2.4}
\end{equation}

\begin{description}
\item[(I.BARE)] (Bare operator.)
At every level, the operator that stands in all propagators is this operator, formed with the bare $\rho$ and the bare weights;
only the mass and the vacuum energy change from level to level.
In every variant of the operator the weights are the bare weights $w_i$:
the relation (I.1.1) enters (I.m1), (I.2.4), (I.O8) and the bound on the operator $T_k(A)$ in Proposition~9,
and no symmetry protects it (Section~I.4.6).
\end{description}

\subsubsection*{I.2.4. Properties that are read off (I.2.3)}

\begin{description}
\item[(I.O0)] (Scope of the facts of Sections~I.1.3 and I.2.2--I.2.5.)
The displays (I.1.2) and (I.2.1)--(I.2.5) are formed on the covering lattice.
The arguments given with the facts use of $T_k$ and of $\mathfrak{S}^{\mathrm{ch}}(\eta_k,\ell)$ only the following:
the map $x\mapsto x^{(i)}$ commutes with the translations (I.s1);
$M_i(x)$ depends on the label only;
$A$ is periodic;
$|f(x)|_2^2$, respectively $f'(x)\cdot f(x)$, is periodic, so that it may be summed over a period cell and $x$ may be replaced by $x^{(i)}$ under the sum.
Hence the displays and the facts may be read for the following data, and for these data the facts hold with the same arguments:
\begin{itemize}
\item a mesh $\mathrm{me}>0$ in the place of $\eta_k$;
\item a translate of $\mathrm{me}\,\mathbb{Z}^3$ in the place of $\mathbb{L}_k$;
\item the periods $n_\mu\,\mathrm{me}\,e_\mu$, with integers $n_\mu\ge1$, in the place of $\ell e_\mu$; all sites $x^{(i)}$ and all bonds $(c,c_3)$ are formed on the covering lattice, and sites and bonds are counted modulo the periods (this matters for $n_\mu\le2$);
\item a number $e_k>0$;
\item a periodic bond function $A$, complex, and real where the fact says so;
\item $\theta\in\mathbb{R}^3$ and the spinors with $f(c+n_\mu\,\mathrm{me}\,e_\mu,g)=-\exp(i\theta_\mu)\cdot f(c,g)$;
\item the measure $\int\mathrm{d}x:=\mathrm{me}^3\cdot\frac{1}{48}\sum_{c,g}$, the sum running over a period cell;
\item a real number $\rho$.
\end{itemize}
For all these data the facts (I.s1)--(I.s5), (I.m1)--(I.m3), (I.m5), the displays (I.2.4), (I.2.5), the facts (I.O1)--(I.O4), (I.O6) and the second half of (I.O7) hold.
Here (I.2.5) is stated at $\mathrm{me}=1$, in lattice units; at a mesh $\mathrm{me}$ the same argument gives
\[
D_A=\mathrm{me}^{-1}\left[\frac{\rho}{2}\,\mathrm{sw}-\sum_iM_i(\,\cdot\,)\cdot R_i(A)\right];
\]
in (I.O4) the functions $d\boldsymbol{\lambda}$ and $\exp(i\,e_k\boldsymbol{\lambda})$ have the periods $n_\mu\,\mathrm{me}\,e_\mu$.
The fact (I.O5) holds if every $\exp(i\theta_\mu)$ is $1$ or $-1$; for the other $\theta$ the pairing $\langle\,\cdot\,,\cdot\,\rangle^{\mathrm{ch}}$ pairs the spinors of $-\theta$ with those of $\theta$, the operator $D_A^T$ acts on the spinors of $-\theta$, and so does $D_{\overline{A}}$ in $\overline{D_{\overline{A}}}$.
For $\rho\notin(0,1)$ the following changes occur:
in (I.m6) one has $|M_i(x)|_2=\frac{w_i}{2}\left(1+|\rho|\right)$, and $M_i(x)$ is invertible if and only if $\rho^2\ne1$;
in (I.O2) the norm of $\mathrm{me}\,D_A$ is at most $\left(1+|\rho|\right)\mathrm{sw}$, and $\operatorname{Re}D_A\ge0$ holds for $\rho\ge0$ only;
in (I.O6) one has $|M_3|_2=3\left(1+|\rho|\right)$;
(I.m7) and the numerals $6$ and $58$ are stated for $\rho\in(0,1)$.
Instances of these data are:
the block tori of Sections~I.1.5 and G1.2 ($\theta=0$ for $\boldsymbol{\sigma}=-1$, and $\theta=(\pi,\pi,\pi)$ for $\boldsymbol{\sigma}=1$);
the spaces of Section~G1.4 with $R<\infty$;
the fibre $p\in\mathbb{R}^3$ ($\mathrm{me}=1$, $n_\mu=1$, $\theta_\mu=p_\mu-\pi$, $A=0$);
the space $\mathfrak{S}^{\mathrm{ch}}(1;n)$ of Proposition~46 ($\mathrm{me}=1$, $\theta=0$), hence all tori $C(n)$ of Sections~I.6 and I.7, of either parity.
\item[(I.O1)] (Rows.)
$D_A$ maps $\mathfrak{S}^{\mathrm{ch}}(\eta_k,\ell)$ into itself.
The row at $x$ reads $f$ at $x$, $x^{(0)}$, $x^{(1)}$, $x^{(2)}$, which lie over $c$, and at $x^{(3)}$, which lies over the nearest neighbour $c_3$; it reads $A$ on the single bond $(c,c_3)$.
Its entries are entire functions of the numbers $A(b)$.
\item[(I.O2)] (Real $A$.)
Let $A$ be real.
Then $R_i(A)$ is a unitary involution for $(\,\cdot\,,\cdot\,)^{\mathrm{ch}}$; $\nabla_A^{*}\nabla_A=\left(\nabla_A^{*}\nabla_A\right)^{*}\ge0$, with
\[
\left(f,\nabla_A^{*}\nabla_Af\right)^{\mathrm{ch}}=\sum_{\{x,x^{(i)}\}}\frac{\eta_k^3}{48}\cdot w_i\cdot\left|U_A\left(x,x^{(i)}\right)\cdot f\left(x^{(i)}\right)-f(x)\right|_2^2,
\]
the sum running over the unoriented bonds $\{x,x^{(i)}\}$; and $J_A^{*}=-J_A$.
Consequently
\[
\operatorname{Re}D_A=\frac{\rho}{2\eta_k}\,\nabla_A^{*}\nabla_A\ge0,
\]
and $D_A^{*}=\gamma_5\cdot D_A\cdot\gamma_5$ by (I.m5).
The norm of $\eta_kD_A$ on $L^2$ is at most $(1+\rho)\,\mathrm{sw}\le58$.
\item[(I.O3)] (Blocks.)
$D_A$ commutes with $\Sigma_3$.
For real $A$ one has $D_A=\operatorname{diag}\left(D_A^{\boldsymbol{\sigma}},\left(D_A^{\boldsymbol{\sigma}}\right)^{*}\right)$ in the two eigenspaces of $\Sigma_3$; hence
\[
\det\left(D_A+t\right)=\left|\det\left(D_A^{\boldsymbol{\sigma}}+t\right)\right|^2\ge0
\]
for every real number $t$.
This is \cite[Proposition~3.18]{QEDII}.
\item[(I.O4)] (Gauge.)
Let $\boldsymbol{\lambda}$ be a complex function of $c$ such that $d\boldsymbol{\lambda}$ and $\exp(i\,e_k\boldsymbol{\lambda})$ are $\ell\mathbb{Z}^3$-periodic.
Then
\[
D_{A-d\boldsymbol{\lambda}}\cdot\exp(i\,e_k\boldsymbol{\lambda})=\exp(i\,e_k\boldsymbol{\lambda})\cdot D_A;
\]
this is a single relation, on the hop of type $3$, the three other hops staying over $c$.
Moreover $D_{A+\frac{2\pi}{e_k\eta_k}m}=D_A$ for every integer valued bond function $m$.
\item[(I.O5)] (Transposition, conjugation; complex $A$.)
By (I.m2) and (I.m1),
\[
D_A^T=\mathsf{C}^{-1}\cdot D_{-A}\cdot\mathsf{C}=\mathrm{Cc}^{-1}\cdot D_{-A}\cdot\mathrm{Cc},
\]
and by (I.m3)
\[
\overline{D_{\overline{A}}}=\varkappa^{-1}\cdot D_{-A}\cdot\varkappa,
\]
where $\overline{O}$ denotes the operator $f\mapsto\overline{O\,\overline{f}}$.
\item[(I.O6)] (One hop.)
For a function $\psi$ of $c$ and bond functions $A$, $B$,
\begin{align*}
\left([D_A,\psi]f\right)(x)&=-\eta_k^{-1}\cdot M_3(x)\cdot u_A(c,c_3)\cdot\left(\psi(c_3)-\psi(c)\right)\cdot f\left(x^{(3)}\right),\\
\left((D_A-D_B)f\right)(x)&=-\eta_k^{-1}\cdot M_3(x)\cdot\left(u_A-u_B\right)(c,c_3)\cdot f\left(x^{(3)}\right).
\end{align*}
The map $x\mapsto x^{(3)}$ is a bijection, and $|M_3|_2,|M_3|_c\le6$.
Thus one hop of norm at most $6$ takes the place of the six hops of norm one of the Wilson operator of \cite{QEDI}.
For this reason the following numerals remain the same for $D_A$:
$12\mathrm{e}$ in \cite[Proposition~5.65]{QEDI};
$12$ in the proof of \cite[Lemma~20.18\,(b)]{QEDI};
$48$ in the condition $48\,e_{*}^{2\varepsilon}\le\mu_{-}$ of \cite[Section~5.6.2]{QEDI}, where the running mass $\mu_{-}$ is written $\mathfrak{m}_{-}$;
and $12M^3$ in the proof of \cite[Proposition~17.16\,(b)]{QEDI}.
\item[(I.O7)] (One formula.)
$D_A$ is given by one formula on $\mathbb{L}^{\mathrm{f}}_k$, of range $\eta_k$, in the sense of Section~I.1.9; it commutes with all translations by vectors of $\eta_k\mathbb{Z}^3$, the bond function $A$ being translated with it.
\item[(I.O8), (I.O9)] (Light and heavy part; scaling.)
These are stated in Sections~I.3.1 and I.3.7.
The lattice covariance of $D_A$ is stated in Section~I.4.4.
\end{description}

\subsubsection*{I.2.5. Lattice units; charge, lattice mass, mass parameter}

At the cutoff $N$, that is, for $\eta=L^{-N}$, the lattice model is written at level $0$ in lattice units:
$\eta_0=1$, the photon lattice is $\mathbb{Z}^3/n\mathbb{Z}^3$, and the charge of level $0$ is
\[
q:=e\,\eta^{\frac{1}{2}},\qquad r_q:=\frac{2\pi}{q}.
\]
As in \cite[Definition~5.3]{QEDI}, the symbols $e_0$, $e_1$, $e_2$ with a numeral as index, read as numbers, denote the ceilings of the lines (C8), (C9), (C10) of Section~I.8.2, and never running charges;
read as vectors, $e_1$, $e_2$, $e_3$ are the unit vectors of Section~I.1.1, as in $ge_1$ and $n_\mu e_\mu$.
In a formula written for a general level $k$, the symbols $e_k$ and $e_{k+1}$ at the levels $0$, $1$, $2$ stand for the running charges $q$, $L^{\frac{1}{2}}q$, $Lq$, and these are written as such.

The lattice mass $\boldsymbol{\mu}$ is a real number, namely the number that is added to $D_A$ in lattice units.
For a real $m_{\mathrm{in}}$ one has $\boldsymbol{\mu}=\bar{m}\eta+m_{\mathrm{in}}$;
for a real sequence $N\mapsto m_{0,N}$ one puts $\boldsymbol{\mu}_N:=\bar{m}\eta+m_{0,N}$.
The mass parameter $M_{\mathrm{ch}}$ is defined by the first identity of
\begin{equation}
M_{\mathrm{ch}}:=\boldsymbol{\mu}+\frac{\rho}{2}\,\mathrm{sw}=\bar{m}\eta+m_{\mathrm{in}}+\frac{\rho}{2}\sum_iw_i,\qquad
D_A+\boldsymbol{\mu}=M_{\mathrm{ch}}-\sum_iM_i(\,\cdot\,)\cdot R_i(A),
\tag{I.2.5}
\end{equation}
and the second identity follows from (I.m1).
The number $M_{\mathrm{ch}}$ is the number denoted $M$ in \cite[Definition~2.10]{QEDII}; in this paper the letter $M$ denotes the side of the $M$-cubes.

In the flow the two parts of $\boldsymbol{\mu}$ play different roles.
The part $\bar{m}\eta$, the running mass of level $0$, stands in the propagators, whereas $m_{\mathrm{in}}$ is a coupling, the coefficient of the mass term in the exponent of Section~I.5.7.
At all levels the mass term is uniform on the $48$ sites of a cluster.

\begin{description}
\item[(I.K10)] (Every real lattice mass.)
Every real lattice mass $\boldsymbol{\mu}$ is admitted:
the statements hold without any of the conditions $\boldsymbol{\mu}>0$, $\boldsymbol{\mu}_N>0$, $\boldsymbol{\mu}>-\frac{1}{2}$, $\boldsymbol{\mu}>-1$.
The results \cite[Proposition~3.18, Theorem~4.26 and Proposition~13.11]{QEDII} hold for every real lattice mass.
Proposition~28.3 does not determine the sign of the shift $m_{0,N}$, so that $\boldsymbol{\mu}_N<0$ may occur.
\end{description}

\subsubsection*{I.2.6. Fourier transform; Bloch symbol; continuum operator}

\emph{Fourier transform.}
For functions on $\mathbb{R}^3$,
\[
\hat{f}(p):=\int\mathrm{d}y\,\exp(-ip\cdot y)\,f(y),\qquad
f(y)=(2\pi)^{-3}\int\mathrm{d}p\,\exp(ip\cdot y)\,\hat{f}(p)\qquad\left(y\in\mathbb{R}^3\right),
\]
and for a $2\pi\mathbb{Z}^3$-periodic function $f$ of $p$,
\[
\mathrm{fh}(z):=(2\pi)^{-3}\int_{[-\pi,\pi]^3}\mathrm{d}^3p\,f(p)\,\exp(-ip\cdot z)\qquad\left(z\in\mathbb{Z}^3\right).
\]
Thus $\partial_\mu$ corresponds to $ip_\mu$.
These are the conventions of \cite[Section~21.1.3 and the proof of Lemma~19.18]{QEDI} and of \cite[Sections~12.8.1 and 14.1]{QEDII}, with one exception, declared in \cite[Section~14.1]{QEDII}: the transform of the transverse profile in \cite[Definition~14.24]{QEDII} carries the opposite sign in the exponent.

\emph{Continuum.}
$S_{\bar{m}}:=\left(\bar{m}+\gamma\cdot\partial\right)^{-1}$ on $L^2(\mathbb{R}^3;\mathbb{C}^4)$, that is, the Fourier multiplier $\left(\bar{m}+i\gamma\cdot p\right)^{-1}$.

\emph{Bloch symbol} (lattice units, $A=0$).
Let ${V_{\mathrm{v}}}:=\left\{u\colon O_{\mathrm{cub}}\to\mathbb{C}^4\right\}$, a space of dimension $192$, and $M_i(g):=M_i(c,g)$.
On the functions $f(c,g)=\exp(ip\cdot c)\,u(g)$ the operator $D_0+\boldsymbol{\mu}$ acts by $u\mapsto D^{(p)}u$, where
\begin{equation}
\left(D^{(p)}u\right)(g):=\boldsymbol{\mu}\,u(g)-\sum_{i\le2}M_i(g)\cdot\left(u(gs_i)-u(g)\right)-M_3(g)\cdot\left(\exp\left(ip\cdot(ge_1)\right)\cdot u(gs_3)-u(g)\right);
\tag{I.2.6}
\end{equation}
it is an entire, $2\pi\mathbb{Z}^3$-periodic function of $p\in\mathbb{C}^3$.
Here, as in \cite{QEDI} and \cite{QEDII}, the forward hop in a direction $a$ carries matrices of the type $\frac{1}{2}\left(\rho-\gamma\cdot a\right)$ and the phase $\exp(+ip\cdot a)$.

\emph{The symbol in the coordinates of the points $\mathrm{geo}(x)$.}
If one writes $f(x)=\exp\left(ip\cdot\mathrm{geo}(x)\right)u(g)$ instead, the operator $D_0+\boldsymbol{\mu}$ acts by a symbol $\mathrm{Dh}(p)$ which depends on $v_0$ and is not periodic.
The two symbols are related by
\[
D^{(p)}=Y\cdot\mathrm{Dh}(p)\cdot Y^{-1},\qquad (Yu)(g):=\exp\left(ip\cdot(gv_0)\right)\cdot u(g),
\]
and $Y$ is unitary for real $p$.
Thus the symbol $\mathrm{Dh}(p)$ and the map $Y$ depend on $v_0$;
both are formed with the map $\mathrm{geo}$ of Section~I.1.3 and Proposition~46\,(a), which depends on $v_0$.
The image sets of Proposition~46 depend on $v_0$ as well:
the set $W_{\mathrm{af}}v_0$ of the fermion sites of \cite[Definition~2.4]{QEDII}, and the set $Y\subset W_{\mathrm{af}}v_0$ of Proposition~46\,(e), the image under $\mathrm{geo}$ of a fundamental domain, which is a set of points and not the map $Y$ above.
Apart from the map $\mathrm{geo}$ and the objects formed with it, no object defined in Part~I depends on $v_0$;
the Berezin integral $\int_Y$ of Proposition~46\,(e), written in the generators $\psi_\alpha(x)$, $\bar{\psi}_\alpha(x)$, contains no $v_0$.
The symbol $\mathrm{Dh}(p)$ and the map $Y$ serve only to compare the two forms of the symbol.

\emph{Momenta.}
For an integer $n\ge1$ and $\theta\in\mathbb{R}^3$ (the antiperiodic torus of $n^3$ photon sites; $\theta$ is the twist of Section~I.6.6) put
\[
P_{n,\theta}:=\left\{p:\ p_\mu=\frac{2\pi j_\mu+\pi+\theta_\mu}{n},\ j\in\{0,\dots,n-1\}^3\right\}.
\]

\subsubsection*{I.2.7. Names of the operator and of its parts in the three papers}

The following table gives, for each object, its name in this paper, in \cite{QEDI} and in \cite{QEDII}.

\medskip
\noindent
\begin{tabular}{|p{2.3cm}|p{3.3cm}|p{3.4cm}|p{3.6cm}|}
\hline
object & this paper & \cite{QEDI} & \cite{QEDII} \\
\hline
operator without mass, mesh $\eta_k$ & $D_A$ & $D_A$, the Wilson operator; it is replaced here by (I.2.3) & $\eta_k^{-1}$ times $D_A$ at $m=0$ \\
\hline
operator with mass, lattice units & $D_A+\boldsymbol{\mu}$ & $D_A+\boldsymbol{\mu}$; $D[qA]+\boldsymbol{\mu}$ & $D_A$; it contains $M$ \\
\hline
lattice mass & $\boldsymbol{\mu}$; $\boldsymbol{\mu}_N$ & $\boldsymbol{\mu}$, $\mathbf{m}$, $\boldsymbol{\mu}_N$ & $m$, $m_N$ \\
\hline
mass parameter & $M_{\mathrm{ch}}$ & --- & $M$ \\
\hline
hopping matrix & $M_i(x)$ & the projections $\frac{1}{2}\left(1\mp\gamma_{\mu}\right)$, written $P^{\boldsymbol{\sigma}}_{\nu}$, $\Pi^{\mp}_{\mu}$, $\mathrm{Pr}^{\boldsymbol{\sigma}}_{\mu}$; they do not occur here & $w\left(x,x^{(i)}\right)$ \\
\hline
hop & $R_i(A)$ & $T^{\boldsymbol{\sigma}}_{\mu}(A)$; it does not occur here & --- \\
\hline
Hermitian part & $\nabla_A^{*}\nabla_A$ & --- & $\nabla_A^{*}\nabla_A=\sum_iw_i-H_A$ \\
\hline
anti-Hermitian part & $J_A$ & --- & $J_A$ \\
\hline
energy & $\eta_k^{-2}\left(f,\nabla_A^{*}\nabla_Af\right)^{\mathrm{ch}}$, without a letter & the square of the norm of $\nabla^{+}u$ & $\langle\psi,\nabla_A^{*}\nabla_A\psi\rangle$ (lattice units, counting pairing); it is $48$ times $\left(f,\nabla_A^{*}\nabla_Af\right)^{\mathrm{ch}}$ \\
\hline
normals & $\hat{n}_i$; $\hat{d}_i(x)=-g\hat{n}_i$ & --- & $\hat{n}_i$; $\hat{d}_i(x)$ \\
\hline
label, photon site & $g$, $c$ & --- & $\vec{g}$, $k=\pi(x)$ \\
\hline
determinant in lattice units & $\mathrm{Dt}^{\mathrm{ch}}(B):=\det\left(D^{\mathrm{ch}}[B]+\boldsymbol{\mu}\right)$, with $D^{\mathrm{ch}}$ of Proposition~46\,(b) & $\mathrm{Dt}(B)$ & $\det D^{T}_A$; the upper index is the torus $T$ \\
\hline
\end{tabular}
\medskip

\subsubsection*{I.2.8. The light/heavy Schur complement}

At a connection $A$, that is, in the functional integral and in the steps of the flow, the heavy components are not integrated out first:
the determinant of the heavy block $D^{\mathrm{hh}}_{A}+\boldsymbol{\mu}\cdot\mathrm{Phv}$ of (I.J6) in Section~I.3.1 is not split off from $\det(D_A+\boldsymbol{\mu})$, and neither the Schur complement $\mathrm{Dsc}_{A}$ of (I.J6) nor the Schur complement of $D_A+\boldsymbol{\mu}$ enters a statement or a proof.
At a connection these complements are only written down, in (I.J6) and in the identity for $J^{*}\cdot S_k(A)\cdot J$ below;
that identity holds wherever $D^{\mathrm{hh}}_{A}+\mu\cdot\mathrm{Phv}$ is invertible on the range of $\mathrm{Phv}$, and it serves to compare the operators of a level $k\ge1$ with those of \cite[Section~5.2.6]{QEDI}.
The four blocks of (I.J6) occur in every field in Proposition~1.1 and in its proof.

In the statements and in their proofs the Schur complement with respect to the light and the heavy part enters only at zero field, fibre by fibre, as the $4\times4$ Schur symbol $\mathrm{Dsc}_{\mathbf{m}}(p)$ of the Bloch symbol $D^{(p)}$ of (I.2.6) at $\boldsymbol{\mu}=\mathbf{m}$ (Proposition~1.2\,(S3)).
It occurs in Proposition~1.2, in the proof of Proposition~1.3, in Proposition~1.5, in Proposition~14.2\,(a), row~$[19]$, in the proofs of Propositions~14.2 and 14.4 (Chapters~14b and 14c), in the proof of Proposition~30.1\,(c), and in the proof of Proposition~37.1\,(f), where it is written $\mathrm{ssc}$.
In Proposition~14.2 and in Chapters~14b and 14c it enters only through the constant
\[
K_{\mathrm{Str}}=\sup\left\{\left|\mathrm{Dsc}_{\mathbf{m}}(p)\right|_2:\ |\operatorname{Im}p_{\mu}|\le\theta_{0},\ 0\le\mathbf{m}\le1\right\}
\]
of Proposition~1.2\,(S3), which row~$[19]$ repeats.
At $p=0$ one has $\sin0=0$ and $\mathrm{tw}(0)=0$, so that (S3) gives $\mathrm{Dsc}_{\mathbf{m}}(0)=\mathbf{m}$ and hence $K_{\mathrm{Str}}\ge1$;
and $\mathrm{Dsc}_{\mathbf{m}}(p)$, hence $K_{\mathrm{Str}}$, is formed with the symbol $D^{(p)}$ of (I.2.6), which contains $\rho$ and $w$.
In the proof of Proposition~30.1\,(c), in Chapter~30, the determinant of the heavy block $D^{\mathrm{hh}}(p)$ of the symbol is split off fibre by fibre, by the factorisation
\[
\det D^{(p)}=\det D^{\mathrm{hh}}(p)\cdot\det\mathrm{Dsc}_{\mathbf{m}}(p)
\]
of Proposition~1.2\,(S8);
the constants $C_{\mathrm{lH}}$ and $C_{\mathrm{lS}}$ of Proposition~30.1, from which $C_6$ is formed, belong to the heavy and to the light factor.

The operators $D_k(A)$, $M_k(A)$, $\dot{M}_k(A)$ and $\Gamma_k(A)$ of a level $k\ge1$, which act on the unit-lattice spinors $\mathfrak{S}(1,\ell)$ (table of Section~I.3.5), are formed with the operator $S_k(A)$ of Section~I.3.5, $\mu$ being the mass of the data $(k,\ell,\mu,e_k)$.
Since $J^{*}\cdot J=1$ and $J^{*}\cdot\mathrm{Phv}=0=\mathrm{Phv}\cdot J$, the operator
\[
S_k(A)^{-1}=D_A+\mu+b_k\cdot J\cdot P_k(A)\cdot J^{*}
\]
has, with respect to the light and the heavy part, the blocks $D^{\mathrm{ll}}_{A}+\mu+b_k\cdot P_k(A)$, $D^{\mathrm{lh}}_{A}$, $D^{\mathrm{hl}}_{A}$ and $D^{\mathrm{hh}}_{A}+\mu\cdot\mathrm{Phv}$.
Hence, wherever $D^{\mathrm{hh}}_{A}+\mu\cdot\mathrm{Phv}$ is invertible on the range of $\mathrm{Phv}$,
\[
J^{*}\cdot S_k(A)\cdot J=\left(D^{\mathrm{ll}}_{A}+\mu-D^{\mathrm{lh}}_{A}\cdot\left(D^{\mathrm{hh}}_{A}+\mu\cdot\mathrm{Phv}\right)^{-1}\cdot D^{\mathrm{hl}}_{A}+b_k\cdot P_k(A)\right)^{-1}.
\]
So $D_k(A)$, $M_k(A)$ and $\Gamma_k(A)$ are the operators of \cite[Section~5.2.6]{QEDI} in which $D_A+\mathfrak{m}$, formed there with the Wilson operator, is replaced by the Schur complement
\[
D^{\mathrm{ll}}_{A}+\mu-D^{\mathrm{lh}}_{A}\cdot\left(D^{\mathrm{hh}}_{A}+\mu\cdot\mathrm{Phv}\right)^{-1}\cdot D^{\mathrm{hl}}_{A}
\]
of $D_A+\mu$, and not by $\mathrm{Dsc}_{A}+\mu$.
For $\dot{M}_k(A)$ this replacement does not give the operator.
Indeed, $J\cdot J^{*}+\mathrm{Phv}=1$ gives
\[
J^{*}\cdot S_k(A)^2\cdot J=\left(J^{*}\cdot S_k(A)\cdot J\right)^2+J^{*}\cdot S_k(A)\cdot\mathrm{Phv}\cdot S_k(A)\cdot J,
\]
so that $\dot{M}_k(A)=b_k^2\cdot Q_k(A)\cdot J^{*}\cdot S_k(A)^2\cdot J\cdot Q_k^{T}(-A)$ contains the second term, whereas the replacement in the operator $\dot{M}_k(A)$ of \cite[Section~5.2.6]{QEDI} gives $b_k^2\cdot Q_k(A)\cdot\left(J^{*}\cdot S_k(A)\cdot J\right)^2\cdot Q_k^{T}(-A)$.

\subsubsection*{I.2.9. The dictionary}

\begin{named}{Proposition 46}[The operator dictionary]
Let $\rho$ be real, $w$ the standard weights, $\boldsymbol{\mu}$ real, $q>0$, $M_{\mathrm{ch}}:=\boldsymbol{\mu}+\frac{\rho}{2}\,\mathrm{sw}$, $v_0\in\mathrm{Chm}_0$, $n=(n_1,n_2,n_3)$ with integers $n_\mu\ge1$, $|n|:=n_1n_2n_3$, and let $A$ be a complex one-cochain on $\mathbb{Z}^3$ that is periodic under $\sum_{\mu}n_{\mu}\mathbb{Z}\,e_{\mu}$.
Lattice units are used: Sections~I.1, I.2 and I.3.1 are read at $k=0$ with $\eta_0=1$, with $q$ in the place of $e_k$, and with the conditions $f(c+n_\mu e_\mu,g)=-f(c,g)$, respectively $\xi(c+n_\mu e_\mu)=-\xi(c)$, in the place of those of Section~I.1.5.
These are the data of (I.O0) with $\mathrm{me}=1$, $\theta=0$ and the periods $n_\mu e_\mu$, and for them the facts listed in (I.O0) hold:
(I.s1)--(I.s5), (I.m1)--(I.m3), (I.m5), (I.2.4), (I.2.5), (I.O1)--(I.O4), (I.O6), the second half of (I.O7), not its first half, and, since $\theta=0$, also (I.O5);
for $\rho\notin(0,1)$ with the changes listed in (I.O0).
The space of these spinors, with $192$, respectively $4$, components per photon site, is denoted $\mathfrak{S}^{\mathrm{ch}}(1;n)$, respectively $\mathfrak{S}(1;n)$, and $J\colon\mathfrak{S}(1;n)\to\mathfrak{S}^{\mathrm{ch}}(1;n)$, $J^{*}\colon\mathfrak{S}^{\mathrm{ch}}(1;n)\to\mathfrak{S}(1;n)$.
All sites $x^{(i)}$ and all bonds $(c,c+ge_1)$ are formed in $\mathbb{Z}^3$, and no display is formed on the quotient; for $n_\mu\le2$ the rule of Section~I.1.1 on lifts with $|c-c'|<\ell/2$ does not apply to nearest neighbours.
The displays of \cite{QEDI} and \cite{QEDII} named in the clauses are definitions.
They are read with the data of this statement, also where they are stated for a smaller range:
for $\rho\in(0,1)$ in \cite[Definitions~2.10 and 2.36]{QEDII};
for $\rho\in[-1,1]$ in \cite[Sections~13.3 and 13.4]{QEDII}, in the second paragraph of \cite[Definition~2.13]{QEDII}, which defines $Z_T[G]$ on two further classes, the space spanned by the products read on $T$ and the class $\mathrm{Ig}(T)$ of (h), in \cite[Definition~5.11]{QEDII}, which introduces these two classes under the hypothesis $-1\le\rho\le1$ that stands throughout \cite[Section~5]{QEDII}, and in \cite[Proposition~4.23]{QEDII}, of which only the definition $C(n):=\left(\sum_\mu n_\mu\mathbb{Z}\,e_\mu,\chi\right)$, $\chi(n_\mu e_\mu)=-1$, enters, a definition that contains no $\rho$ and is read for every real $\rho$;
for real $A$ in \cite[Definition~4.1]{QEDII};
for a lattice mass $>-1$ or $>-\frac{1}{2}$ in \cite[Sections~10.1.2 and 21.1.4]{QEDI};
for the sides and the mesh in \cite[Sections~10.1.1 and 21.1.4, and Section~7.1.1 with Section~21.1.3]{QEDI};
for a probe $\beta$ that fits $(N,T^{*})$, for $\eta=L^{-N}$ and for a placement $v$ in the paragraph on the placement of the probe in \cite[Section~21.1.6]{QEDI}, where $\hat{\mathfrak{a}}=\eta^{\frac{1}{2}}a_\eta$ is defined.
Finally, $W_{\mathrm{af}}:=\mathbb{Z}^3\rtimes O_{\mathrm{cub}}$ is the group $W$ of \cite[Definition~2.2]{QEDII}, with the elements $g_{\mathrm{w}}=(a,r)\colon y\mapsto ry+a$.
\begin{itemize}
\item[(a)] (Sites.)
The map
\[
\mathrm{geo}\colon\mathbb{Z}^3\times O_{\mathrm{cub}}\to W_{\mathrm{af}}v_0,\qquad\mathrm{geo}(c,g)=c+gv_0,
\]
is a bijection onto the set $W_{\mathrm{af}}v_0$ of the fermion sites of \cite[Definition~2.4]{QEDII}.
It intertwines the action $g_{\mathrm{w}}(c,g):=(rc+a,rg)$ with the action of $W_{\mathrm{af}}$ on $\mathbb{R}^3$; one has $\pi(\mathrm{geo}\,x)=c$ for the map $\pi$ of \cite[Definition~2.4]{QEDII}; the point $\mathrm{geo}(x^{(i)})$ is the neighbour of $\mathrm{geo}(x)$ across the wall of type $i$ in the sense of \cite[Definition~2.4]{QEDII}; and
\[
\frac{\mathrm{geo}\left(x^{(i)}\right)-\mathrm{geo}(x)}{\left|\mathrm{geo}\left(x^{(i)}\right)-\mathrm{geo}(x)\right|_2}=-g\hat{n}_i=\hat{d}_i(x).
\]
Hence the sites, the bonds, their types, the map $\pi$ and the directions $\hat{d}_i$, written in the coordinates $(c,g)$, do not contain $v_0$.
\item[(b)] (One operator.)
For a complex one-cochain $B$ on $\mathbb{Z}^3$, read as a function on the oriented bonds with $B(c',c)=-B(c,c')$, let $D^{\mathrm{ch}}[B]$ be the operator on the functions $f\colon\mathbb{Z}^3\times O_{\mathrm{cub}}\to\mathbb{C}^4$ given by
\begin{multline*}
\left(D^{\mathrm{ch}}[B]f\right)(c,g):=\sum_{i=0}^{2}\frac{w_i}{2}\left(\rho+\gamma\cdot(g\hat{n}_i)\right)\cdot\left(f(c,g)-f(c,gs_i)\right)\\
+\frac{w_3}{2}\left(\rho-\gamma\cdot(ge_1)\right)\cdot\left(f(c,g)-\exp\left(iB(c,c+ge_1)\right)\cdot f(c+ge_1,gs_3)\right).
\end{multline*}
Transported by $\mathrm{geo}$, the following operators on the functions $\mathbb{Z}^3\times O_{\mathrm{cub}}\to\mathbb{C}^4$ are equal, for every complex $A$:
the operator $D_A$ of \cite[Definition~2.10]{QEDII} with the parameters $(\rho,w,M=M_{\mathrm{ch}},q)$;
the operator $D_A+\boldsymbol{\mu}$, with $D_A$ the display (I.2.3);
and the operator $D^{\mathrm{ch}}[qA]+\boldsymbol{\mu}$.
At level $k$ the operator (I.2.3) is $\eta_k^{-1}D^{\mathrm{ch}}[B]$ in the variable $c/\eta_k$, with $B(c/\eta_k,c'/\eta_k):=e_k\eta_k\cdot A(c,c')$.
\item[(c)] (Spin matrices.)
The relations (I.c1)--(I.c6) hold.
At every place of \cite{QEDI} and \cite{QEDII} named in the table of Section~I.2.1 at which the matrices of the lattice model are defined, the matrix is the one in the first column of its row.
At the following places no matrix of its row is defined.
\cite[Section~6.1.2]{QEDI} (its first sentence) and \cite[Section~12.1]{QEDI} admit every choice with the relations listed there, namely a Hermitian Clifford triple and, in \cite[Section~12.1]{QEDI}, also a Hermitian matrix $\hat{\Gamma}$ with $\hat{\Gamma}^2=1$; the matrices $\gamma_1$, $\gamma_2$, $\gamma_3$, $\gamma_5$ have these relations by (I.c1).
\cite[Section~11.1.2]{QEDI} names the three matrices without their order; the triple is taken in the order given by the numbering $\mu\mapsto\mu+1$ of the directions (Section~I.1.1).
In \cite[Section~6]{QEDI}, from the paragraph on the spinor basis adapted to the time direction in \cite[Section~6.3.2]{QEDI} on, $\gamma_\nu$ and $\gamma_i$ denote $w\gamma_\nu w^{*}$ and $w\gamma_iw^{*}$, with a unitary matrix $w$ chosen there (for $s=4$: $w\gamma_\nu w^{*}=\Sigma_3$).
From that paragraph on these letters do not denote the matrices of the row $\gamma_1,\gamma_2,\gamma_3$ of the table; the triple $\left(w\gamma_\mu w^{*}\right)_\mu$ is unitarily conjugate to the triple of \cite[Section~6.1.2]{QEDI}, which is that of Section~I.2.1.
\item[(d)] (Parameters of a cutoff.)
Let $\eta=L^{-N}$, $q=e\,\eta^{\frac{1}{2}}$ and $\boldsymbol{\mu}=\bar{m}\eta+m_{\mathrm{in}}$.
The lattice mass $m$ of \cite[Definition~2.10]{QEDII}, called $m_N$ in \cite[Definition~2.36]{QEDII} and $\boldsymbol{\mu}_N$ in \cite[Section~21.1.4]{QEDI} for $m_{\mathrm{in}}=m_{0,N}$, is $\boldsymbol{\mu}$.
The mass parameter $M$ of \cite[Definitions~2.10 and 2.36]{QEDII} is $M_{\mathrm{ch}}=\bar{m}\eta+m_{\mathrm{in}}+\frac{\rho}{2}\sum_iw_i$.
\item[(e)] (Determinant; Berezin integral.)
$\det(D_A+\boldsymbol{\mu})$ is the determinant of a linear map of the space $\mathfrak{S}^{\mathrm{ch}}(1;n)$, of dimension $192|n|$, into itself; it depends on no basis, on no pairing and not on $v_0$.
For $n_\mu\ge4$ it is $\det D^{T}_A$ of \cite[Definition~2.11]{QEDII} for $T=C(n)$.
For all $n_\mu\ge1$ it is $\det D^{T}_A$ for $T=C(n)$ in the sense of \cite[Definition~4.1 and Proposition~4.23]{QEDII}, where every subgroup of $\mathbb{Z}^3$ of rank three is admitted; \cite[Definitions~2.6 and 2.11]{QEDII} define $C(n)$ for $n_\mu\ge4$ only, and for complex $A$ the display is that of \cite[Definition~2.10]{QEDII}.
It is $\mathrm{Dt}^{\mathrm{ch}}(qA)$ of Section~I.2.7, and, for $n=(n,n,n)$, $\boldsymbol{\mu}=\bar{m}\eta$ and $qA=B_t$, it is the numerator $\det\left(D^{\mathrm{ch}}[B_t]+\bar{m}\eta\right)$ of the left side of the identity in (g+).
Name the generators of \cite[Definition~2.8]{QEDII} at the site $\mathrm{geo}(x)$ by $\psi_\alpha(x)$, $\bar{\psi}_\alpha(x)$, $\alpha=1,\dots,4$.
Let $Y\subset W_{\mathrm{af}}v_0$ be the image under $\mathrm{geo}$ of a fundamental domain of $\sum_\mu n_\mu\mathbb{Z}\,e_\mu$ in $\mathbb{Z}^3\times O_{\mathrm{cub}}$, and let $\int_Y$ be the Berezin integral of \cite[Definition~2.8]{QEDII} over the generators $\psi_\alpha(x)$, $\bar{\psi}_\alpha(x)$ with $\mathrm{geo}(x)\in Y$.
For every linear map $O$ of $\mathfrak{S}^{\mathrm{ch}}(1;n)$,
\[
\int_Y\exp\left(-\sum_x\bar{\psi}(x)\cdot(O\psi)(x)\right)=\det O,\qquad
\int_Y\exp\left(-\left\langle\bar{\psi},O\psi\right\rangle^{\mathrm{ch}}\right)=48^{-192|n|}\det O,
\]
$x$ running over the fundamental domain.
Hence $\mathrm{D}\Psi:=48^{192|n|}\int_Y$ is the integral with
\[
\int\mathrm{D}\Psi\,\exp\left(-\left\langle\bar{\Psi},O\Psi\right\rangle^{\mathrm{ch}}\right)=\det O
\]
for all $O$.
The constant contains neither $A$ nor $\boldsymbol{\mu}$; it cancels in every quotient of two such integrals, hence in the density of Section~I.7.3 and in all laws.
\item[(f)] (Laws.)
Let $n_\mu\ge4$, and let $\Phi$ be a Borel function of the curvature $F^{\mathrm{u}}$ with $\left|\Phi\left(F^{\mathrm{u}}\right)\right|\le\mathrm{cf}\cdot\exp\left(\mathrm{cf}\cdot\|F^{\mathrm{u}}\|\right)$ for some $\mathrm{cf}$.
Then the number $Z_T[\Phi]$ of \cite[Definition~2.13]{QEDII} for $T=C(n)$ and $M=M_{\mathrm{ch}}$ equals the number $Z_n[\Phi]$ of \cite[Section~5.2.1]{QEDI}, formed with $\det(D_A+\boldsymbol{\mu})$, $D_A$ of (I.2.3), in the place of the Wilson determinant.
(\cite[Section~5.2.1]{QEDI} states this definition for odd cubic tori and refers to \cite[Definition~10.1]{QEDI} for the other $n$; the number is $\boldsymbol{\omega}_n(\Phi)$ of \cite[Section~21.1.4]{QEDI}.)
\item[(f+)] (Sectors.)
In the setting of (f), let $w=(w_{23},w_{13},w_{12})\in\mathbb{Z}^3$.
Then
\[
Z^{(w)}_T[\Phi]=Z^{(w)}_n[\Phi]=\boldsymbol{\omega}^{(w)}_n(\Phi),
\]
where the left side is the sector functional of \cite[Definition~13.7]{QEDII}, the middle term that of Section~I.7.3, and the right side that of \cite[Section~21.1.4]{QEDI}, formed with $\det(D_A+\boldsymbol{\mu})$.
Moreover $P_T(w)=P_n(w)$, where $P_T(w)$ is the weight of the sector $w$ in \cite[Definition~13.7]{QEDII};
the law $\mu^{\mathrm{av}}_T$ of \cite[Definition~13.7]{QEDII} is the law of the periodic extension of $F^{\mathrm{u}}$ under $E^{\mathrm{av}}_n$;
the cochain $h_w$ of \cite[Definition~13.7]{QEDII} is the cochain $h_w$ of Section~I.7.1; and
\[
\exp\left(-\frac{\|h_w\|^2}{2}\right)=\mathfrak{q}_w(n)=\exp\left[-\frac{r_q^2}{2|n|}\sum_{\mu<\nu}n_{\mathfrak{l}}^2\,w_{\mu\nu}^2\right],
\]
$\mathfrak{l}$ being the third direction.
In \cite[Definition~13.7]{QEDII} the flux is written $k$, the cochain $h_w$ is written $\mathfrak{h}_k$, and the weight of a sector is written as the expectation of its indicator.
\item[(g)] (Free ratio.)
Let the data be those of (g+) with $v=0$, in particular $e>0$ and $q=e\,\eta^{\frac{1}{2}}$, and let $a_\eta$ and $\hat{\mathfrak{a}}=\eta^{\frac{1}{2}}a_\eta$ be formed by the displays of \cite[Section~21.1.6]{QEDI} with $\beta$ itself, that is, with $\tau_0\beta=\beta$ in the place of $\tau_v\beta$;
the conditions of the paragraph on the placement of the probe there ($\beta$ fits $(N,T^{*})$, $\eta=L^{-N}$, and the inequality for $v$) are not imposed.
Then the phase field $B_t$ of (g+) is $q\cdot\frac{t}{e}\cdot\hat{\mathfrak{a}}$, and on $\mathfrak{S}^{\mathrm{ch}}(1;(n,n,n))$
\[
\frac{\det\left(D^{\mathrm{ch}}[B_t]+\bar{m}\eta\right)}{\det\left(D^{\mathrm{ch}}[0]+\bar{m}\eta\right)}
=\frac{\det\left(D_{(t/e)\hat{\mathfrak{a}}}+\bar{m}\eta\right)}{\det\left(D_0+\bar{m}\eta\right)}.
\]
\item[(g+)] (Free ratio, every $v\in\eta\mathbb{Z}^3$, in particular every placement.)
Let $0<\eta\le1$, let $R/\eta=n$ be an integer $\ge2$, let $e>0$ and $q=e\,\eta^{\frac{1}{2}}$, let $\beta$ be a probe (Section~I.7.6) and $v\in\eta\mathbb{Z}^3$, and let $a_\eta$, $\hat{\mathfrak{a}}=\eta^{\frac{1}{2}}a_\eta$ be the lattice probe of Section~I.7.6 formed with $\tau_v\beta$ on $\Lambda_{\eta,R}$; let $u$ be the periodised cochain of Section~I.7.6 formed with $\tau_v\beta$.
For complex $t$ let $B_t$ be the phase field on the oriented bonds of $\mathbb{Z}^3$ given by
\[
B_t(c,c+e_\nu):=-t\sum_{\mu}\left(u_{\mu\nu}(\eta c)-u_{\mu\nu}(\eta c-\eta e_\mu)\right),\qquad
B_t(c+e_\nu,c):=-B_t(c,c+e_\nu)\qquad\left(c\in\mathbb{Z}^3\right).
\]
Then $B_t=q\cdot\frac{t}{e}\cdot\hat{\mathfrak{a}}$, both sides being $t\eta\cdot a_\eta$, and
\[
\frac{\det\left(D^{\mathrm{ch}}[B_t]+\bar{m}\eta\right)}{\det\left(D^{\mathrm{ch}}[0]+\bar{m}\eta\right)}
=N_0(t)=\frac{\det\left(D_{(t/e)\hat{\mathfrak{a}}}+\bar{m}\eta\right)}{\det\left(D_0+\bar{m}\eta\right)}
\]
on $C(n,n,n)$; for $n=2,3$ this is the lattice $n\mathbb{Z}^3$ with $\chi(ne_\mu)=-1$, in the sense of \cite[Proposition~4.23]{QEDII}.
In (g), in (g+) and in the density of Section~I.7.3 named in (e), the denominator $\det(D_0+\bar{m}\eta)$ is positive for $\rho\ge0$: it is not $0$, because $\operatorname{Re}(D_0+\bar{m}\eta)\ge\bar{m}\eta>0$ by (I.O2), and it is non-negative by (I.O3), with the scope (I.O0).
For $\rho<0$ the numerators are equal and the denominators are equal, and the quotients are read where the denominator is not $0$.
The phase field $B_t$ and the two quotients are defined by these displays, with the periodised $u$ of Section~I.7.6, for every $R$ as above.
In particular they are defined for $R\le D_{\mathrm{be}}$, where $D_{\mathrm{be}}$ is the least side of a closed coordinate cube that contains $\operatorname{supp}\tau_v\beta$; for a coordinate cube, the condition that $\operatorname{supp}\tau_v\beta$ lie in the interior of a cube of side $R$ reads $D_{\mathrm{be}}<R$, and it is not imposed.
Proposition~36 is stated on this range too.
\item[(h)] (The Grassmann structure of \cite{QEDII}.)
Let the generators of \cite[Definition~2.8]{QEDII} at the site $\mathrm{geo}(x)$ be named $\psi_\alpha(x)$, $\bar{\psi}_\alpha(x)$, where $x\in\mathbb{Z}^3\times O_{\mathrm{cub}}$ and $\alpha=1,\dots,4$ is the spin index of \cite{QEDII}, written $\beta$ in Section~I.4.8 and in Chapter~46.
For the other objects of this clause \cite{QEDII} uses the following letters.
A torus $T$ is written $\mathbb{T}$ \cite[Definition~2.11]{QEDII}.
A monomial $\mathrm{mm}$ in the generators and its charge density $\rho_{\mathrm{mm}}$ are written $\Psi$ and $\mathfrak{q}(\Psi)$ \cite[Definition~2.10]{QEDII}.
A box $\mathrm{Bx}$, the set $X(\mathrm{Bx})$ of its fermion sites, the class $\mathrm{Ad}(\mathrm{Bx})$ and the coefficients $c_{\mathrm{mm}}$ of an integrand are written $B$, $\mathbb{X}(B)$, $\mathcal{G}(B)$ and $a_\Psi$ \cite[Definition~2.12]{QEDII}.
The class $\mathrm{Pl}$ is written $\mathcal{P}$ \cite[Definition~2.15]{QEDII}, and the class $\mathrm{Ig}(T)$ is written $\mathcal{I}(\mathbb{T})$ \cite[Definitions~2.13 and 5.11]{QEDII}.
In these names the following holds.
\begin{itemize}
\item The reflection $\Theta_m$ of \cite[Definition~2.9]{QEDII}, $m=(\ell,\mathrm{po})$, is the antilinear map that reverses the order of products and acts on the generators by
\[
\Theta_m\psi(x)=-\bar{\psi}(\theta x)\,\gamma\cdot\hat{n},\qquad
\Theta_m\bar{\psi}(x)=-\gamma\cdot\hat{n}\,\psi(\theta x),
\]
where $\theta x$ is the action (a) of $\theta=\theta_{\ell,\mathrm{po}}\in W_{\mathrm{af}}$ and $\hat{n}=n_\ell/|n_\ell|_2$; it is the map of Section~I.4.8.
\item The identities \cite[Lemma~3.10\,(c), (d)]{QEDII}, the reflection of integrands and the action $G\mapsto G^{\mathrm{gw}}$ of \cite[Section~2.3]{QEDII} are formed with this map, with the action (a) on the sites and with the pull-backs of cochains.
\item The gauge transformation of \cite[Definition~2.10]{QEDII} is $\psi(x)\mapsto\exp\left(iq\boldsymbol{\lambda}(c)\right)\psi(x)$, $\bar{\psi}(x)\mapsto\exp\left(-iq\boldsymbol{\lambda}(c)\right)\bar{\psi}(x)$, and the charge density $\rho_{\mathrm{mm}}(c')$ of a monomial $\mathrm{mm}$ is the number of its factors $\psi(x)$ with $c=c'$ minus the number of its factors $\bar{\psi}(x)$ with $c=c'$.
\item The set $X(\mathrm{Bx})$ of the fermion sites of a box $\mathrm{Bx}$ in \cite[Definition~2.12]{QEDII} is $\left\{x:\ c+\left[-\frac{1}{2},\frac{1}{2}\right]^3\subset\mathrm{Bx}\right\}$.
\item The classes $\mathrm{Ad}(\mathrm{Bx})$ and $\mathcal{A}(U)$ of \cite[Definition~2.12]{QEDII} and the class $\mathrm{Pl}$ of \cite[Definition~2.15]{QEDII} are formed with these objects and with photon objects only.
\end{itemize}
Let $T=C(n)$ with $n_\mu\ge4$, and $M=M_{\mathrm{ch}}$.
The algebra $\mathrm{Gr}_T$ of \cite[Definition~2.11]{QEDII} is generated by the $\psi_\alpha(x)$, $\bar{\psi}_\alpha(x)$ subject to $\psi(c+n_\mu e_\mu,g)=-\psi(c,g)$ and $\bar{\psi}(c+n_\mu e_\mu,g)=-\bar{\psi}(c,g)$; $\int_T=\int_Y$, with $Y$ as in (e); $S_T(A)=\sum_x\bar{\psi}(x)\cdot\left((D_A+\boldsymbol{\mu})\psi\right)(x)$; and for every $G$ in the three classes on which the first two paragraphs of \cite[Definition~2.13]{QEDII} define the number by one display, that is, for $G\in\mathcal{A}(\mathbb{R}^3)$ that embeds in $T$ or is read on $T$, and for $G\in\mathrm{Ig}(T)$,
\[
Z_T[G]=\sum_{\mathfrak{z}:\,d\mathfrak{z}=0}\ \int_{[0,r_q)^{\mathrm{links}}}\mathrm{d}A\,\exp\left(-\frac{\|F^{\mathrm{u}}\|^2}{2}\right)\int_Y\exp\left(-\sum_x\bar{\psi}(x)\cdot\left((D_A+\boldsymbol{\mu})\psi\right)(x)\right)G,
\]
where $x$ runs over a fundamental domain and $F^{\mathrm{u}}=dA+r_q\mathfrak{z}$; for $G=\Phi(F^{\mathrm{u}})$ this is the number of (f).
None of these expressions contains $v_0$; in this sense Sections~I.4.8, I.6.4, I.7.2 and I.9.3 read these objects in the coordinates $(c,g)$.
\end{itemize}
\end{named}

\begin{proof}
The proof of (a)--(g+) is given in Chapter~46, and the paragraphs on (a)--(g+) below are an outline of it.
The paragraph on (h) is the proof of (h).

(a)
One has $|gv_0|_\infty<\frac{1}{2}$, so that $c$ is the point of $\mathbb{Z}^3$ nearest to $c+gv_0$; the coordinates of $v_0$ have pairwise different non-zero moduli, so that $gv_0=g'v_0$ implies $g=g'$.
This gives the injectivity and $\pi(\mathrm{geo}\,x)=c$.
The surjectivity is the definition of $W_{\mathrm{af}}v_0$, since every element of $W_{\mathrm{af}}$ is a map $y\mapsto c+gy$.
Next, $g_{\mathrm{w}}(c+gv_0)=(rc+a)+(rg)v_0$.
In \cite[Definition~2.4]{QEDII} the neighbour of $x=\tilde{g}v_0$, $\tilde{g}=(c,g)$, across the wall of type $i$ is $\tilde{g}\tilde{s}_iv_0$, where $\tilde{s}_i$ is the reflection in the wall $H_i$ of $\mathrm{Chm}_0$: $\tilde{s}_i=s_i$ for $i\le2$, these walls passing through $0$, and $\tilde{s}_3y=e_1+s_3y$, the wall of type $3$ being $\{y_1=\frac{1}{2}\}$.
Hence $\tilde{g}\tilde{s}_iv_0=c+gs_iv_0=\mathrm{geo}(c,gs_i)$ for $i\le2$, and $\tilde{g}\tilde{s}_3v_0=c+ge_1+gs_3v_0=\mathrm{geo}(c+ge_1,gs_3)$.
The vector $\hat{n}_i$ is the inward unit normal of $H_i$; therefore $\tilde{s}_iv_0-v_0=-2\operatorname{dist}(v_0,H_i)\,\hat{n}_i$, and
\[
\mathrm{geo}\left(x^{(i)}\right)-\mathrm{geo}(x)=-2\operatorname{dist}(v_0,H_i)\,g\hat{n}_i,\qquad\operatorname{dist}(v_0,H_i)>0.
\]
This is \cite[Lemma~3.3 and Remark~3.5]{QEDII}.

(b)
By (a), the hopping matrix of the bond $\{x,x^{(i)}\}$ in \cite[Definition~2.10]{QEDII} is $M_i(x)$, and the phase $U_A(x,x^{(i)})=\exp\left(iqA(\pi x,\pi x^{(i)})\right)$ of \cite[Definition~2.10]{QEDII} is $U_A(x,x^{(i)})$ of (I.1.6).
So the operator of \cite[Definition~2.10]{QEDII} reads $M_{\mathrm{ch}}f(x)-\sum_iM_i(x)\cdot U_A(x,x^{(i)})\cdot f(x^{(i)})$.
By (I.m1), $\sum_iM_i(x)=\frac{\rho}{2}\,\mathrm{sw}=M_{\mathrm{ch}}-\boldsymbol{\mu}$; subtracting and adding $\sum_iM_i(x)f(x)$ one obtains
\[
\boldsymbol{\mu}f(x)-\sum_iM_i(x)\cdot\left[U_A\left(x,x^{(i)}\right)\cdot f\left(x^{(i)}\right)-f(x)\right],
\]
which is (I.2.3) plus $\boldsymbol{\mu}$.
In the display that defines $D^{\mathrm{ch}}$, the term of index $i\le2$ is
\[
\frac{w_i}{2}\left(\rho+\gamma\cdot(g\hat{n}_i)\right)\left(f(x)-f\left(x^{(i)}\right)\right)=-M_i(x)\left[f\left(x^{(i)}\right)-f(x)\right],
\]
and for $B=qA$ the last term is
\[
\frac{w_3}{2}\left(\rho-\gamma\cdot(ge_1)\right)\left(f(x)-\exp\left(iB(c,c+ge_1)\right)f\left(x^{(3)}\right)\right)=-M_3(x)\left[U_A\left(x,x^{(3)}\right)f\left(x^{(3)}\right)-f(x)\right].
\]
The sentence on the level $k$ is the definition (I.2.3).

(c)
Multiplication of $4\times4$ matrices.

(d)
The definitions, together with (b).

(e)
That $\det(D_A+\boldsymbol{\mu})$ is the determinant of a linear map of the space $\mathfrak{S}^{\mathrm{ch}}(1;n)$, of dimension $192|n|$, into itself, and that it depends on no basis, on no pairing and not on $v_0$, is linear algebra.
The map $\mathrm{geo}$ identifies $\mathfrak{S}^{\mathrm{ch}}(1;n)$ with the spinors of $C(n)$, for which $\chi(n_\mu e_\mu)=-1$ by \cite[Definition~2.11]{QEDII}, and (b) identifies the operators.
By \cite[Definition~2.8]{QEDII}, $\int_Y\exp\left(-\bar{\psi}\,\mathrm{Dm}\,\psi\right)=\det\mathrm{Dm}$ for a matrix $\mathrm{Dm}$.
Now $\sum_x\bar{\psi}(x)\cdot(O\psi)(x)=\bar{\psi}\,\mathrm{Dm}\,\psi$ with $\mathrm{Dm}$ the matrix of $O$ in the basis of the delta functions, and $\langle\bar{\psi},O\psi\rangle^{\mathrm{ch}}=\frac{1}{48}\sum_x\bar{\psi}(x)\cdot(O\psi)(x)$ belongs to the matrix $\mathrm{Dm}/48$, of size $192|n|$.

(f)
The two definitions have the same configurations $(A,\mathfrak{z})$, the same measure, the same curvature $F^{\mathrm{u}}=dA+r_q\mathfrak{z}$ and the same Maxwell weight; indeed, by the paragraph that follows \cite[Definition~2.6]{QEDII}, the integral over the configurations of that definition is, for $n_\mu\ge4$, the sum over all flux sectors of the integrals with which \cite[Definition~10.1]{QEDI} forms the partition functions.
The determinants agree by (e).
In \cite[Section~21.1.4]{QEDI} the integral over $A$ stands outside the sum over $\mathfrak{z}$, while in the other definitions the sum stands outside; the two orders give the same number, because sum and integral converge absolutely.
Indeed, for real $A$,
\[
\left|\det(D_A+\boldsymbol{\mu})\right|\le\left(|M_{\mathrm{ch}}|+\frac{1+|\rho|}{2}\,\mathrm{sw}\right)^{192|n|},
\]
and
\[
\left|\Phi\left(F^{\mathrm{u}}\right)\right|\cdot\exp\left(-\frac{\|F^{\mathrm{u}}\|^2}{2}\right)\le\mathrm{cf}\cdot\exp\left(\mathrm{cf}^2\right)\cdot\exp\left(-\frac{\|F^{\mathrm{u}}\|^2}{4}\right),
\]
which is summable and integrable by the majorant of \cite[Lemma~3.14]{QEDII}.
The action $S_T(A)$ of \cite[Definition~2.11]{QEDII}, formed by the display of \cite[Definition~2.10]{QEDII} with the generators and with the rule $\psi(x+v)=\chi(v)\psi(x)$, is the counting form $\sum_x\bar{\psi}(x)\cdot\left((D_A+\boldsymbol{\mu})\psi\right)(x)$ of the map $D_A+\boldsymbol{\mu}$; so $\int_T\exp\left(-S_T(A)\right)=\det(D_A+\boldsymbol{\mu})$ by the first display of (e).
The scalar function $\Phi(dA+r_q\mathfrak{z})$ lies in the class $\mathrm{Ig}(T)$ of \cite[Definitions~2.13 and 5.11]{QEDII}.

(f+)
In Section~I.7.1, in \cite[Definition~13.7]{QEDII} and in \cite[Section~21.1.4]{QEDI} the components of the flux are labelled by the pairs $\mu<\nu$ of directions;
so they are in \cite[Definition~11.1]{QEDI}, whose directions $0,1,2$ become $1,2,3$ under the renumbering $\mu\mapsto\mu+1$ (Section~I.1.1).
In each of these places the component with the label $\mu\nu$ is the sum of $\mathfrak{z}$ over a coordinate two-torus in the directions $\mu$, $\nu$, and this sum does not depend on the choice of the two-torus, because $d\mathfrak{z}=0$ (Chapter~46, Section~9.2).
Hence components with the same label are equal, and in all three papers the sector of $w$ is the set of the configurations whose flux has the components $w_{\mu\nu}$.
The components are listed in the order $(23,13,12)$ in Section~I.7.1 and in \cite[Definition~13.7]{QEDII}, and, after the renumbering, in \cite[Definition~11.1]{QEDI}, where the triple is $\left(\mathfrak{n}_{12},\mathfrak{n}_{02},\mathfrak{n}_{01}\right)$;
\cite[(10.2)]{QEDI} and \cite[Section~10.1.4]{QEDI} list them in the order $(12,13,23)$, and \cite[Section~21.1.4]{QEDI} fixes no order.
Only the order of listing differs.
Moreover
\[
\|h_w\|^2=r_q^2\sum_{\mu<\nu}w_{\mu\nu}^2\cdot\frac{n_{\mathfrak{l}}}{n_\mu n_\nu}=\frac{r_q^2}{|n|}\sum_{\mu<\nu}n_{\mathfrak{l}}^2\,w_{\mu\nu}^2 .
\]

(g)
Since $q=e\,\eta^{\frac{1}{2}}$, one has $q\cdot\frac{t}{e}\cdot\eta^{\frac{1}{2}}a_\eta=t\eta\cdot a_\eta$.
Moreover
\[
-\eta^{-1}\sum_\mu\left(u_{\mu\nu}(y)-u_{\mu\nu}(y-\eta e_\mu)\right)=\left(d_\eta^{*}u\right)_\nu(y)=a_{\eta,\nu}(y)
\]
for the lattice probe $a_\eta$ of \cite[Section~21.1.6]{QEDI}, because
\[
\left(d^{*}u\right)_\nu(c)=-\sum_\mu\left(u_{\mu\nu}(c)-u_{\mu\nu}(c-e_\mu)\right)
\]
for the coboundary of Section~I.1.1.
Then (b) gives the identity of the two quotients.

(g+)
The computation of (g) does not use $v=0$.

(h)
By (a), $\mathrm{geo}$ is a bijection of $\mathbb{Z}^3\times O_{\mathrm{cub}}$ onto $W_{\mathrm{af}}v_0$.
So the names index all generators of \cite[Definition~2.8]{QEDII}.
For a set $Z\subset W_{\mathrm{af}}v_0$, the algebra $\mathrm{Gr}(Z)$ of \cite[Definition~2.8]{QEDII} is the Grassmann algebra generated by the $\psi_\alpha(x)$, $\bar{\psi}_\alpha(x)$ with $\mathrm{geo}(x)\in Z$, and for finite $Z$ the Berezin integral $\int_Z$ is the coefficient of the product of the $\psi_\alpha(x)\bar{\psi}_\alpha(x)$ over these $x$ and $\alpha=1,\dots,4$.
The reflection $\theta=\theta_{\ell,\mathrm{po}}$ lies in $W=W_{\mathrm{af}}$ (\cite[Definition~2.2]{QEDII}), and $\theta(\mathrm{geo}\,x)=\mathrm{geo}(\theta x)$ by (a); the unit normal $\hat{n}_\ell$ of \cite[Definition~2.2]{QEDII} is $n_\ell/|n_\ell|_2=\hat{n}$ (Section~I.1.1).
So the display of \cite[Definition~2.9]{QEDII} is the display of (h), and an antilinear map that reverses products is determined by its values on the generators.
The identities \cite[Lemma~3.10\,(c), (d)]{QEDII} are derived there from the display of \cite[Definition~2.9]{QEDII}; the sites $\theta x$, $\theta x'$, $\theta_m\theta_{m'}x$ in them are the images of $\mathrm{geo}\,x$, $\mathrm{geo}\,x'$ under elements of $W_{\mathrm{af}}$, hence formed by the action of (a).
The reflection of integrands of \cite[Section~2.3]{QEDII} contains, besides $\Theta_m(\mathrm{mm})$, only the pull-back $\theta_m^{*}$ of one-cochains, and the action $G\mapsto G^{\mathrm{gw}}$ is a composition of reflections of integrands, which are formed with the maps $\Theta_m$ of \cite[Definition~2.9]{QEDII}.
By (a), $\pi(\mathrm{geo}\,x)=c$; this gives the displays of the gauge transformation, of $\rho_{\mathrm{mm}}$ and of $X(\mathrm{Bx})$.
The definitions of an integrand, of $\mathrm{Ad}(\mathrm{Bx})$, of a local lift and of $\mathcal{A}(U)$ in \cite[Definition~2.12]{QEDII} and of $\mathrm{Pl}$ in \cite[Definition~2.15]{QEDII} contain, besides $\mathrm{Gr}(X(\mathrm{Bx}))$, the gauge transformation and $\rho_{\mathrm{mm}}$, only boxes, links, plaquettes, cochains and bounds on the coefficients $c_{\mathrm{mm}}$.

Let $T=C(n)$ with $n_\mu\ge4$.
For $v\in\mathbb{Z}^3$ one has $\mathrm{geo}(x)+v=\mathrm{geo}(x+v)$.
So the rule of \cite[Definition~2.11]{QEDII} for the fields on $T$, according to which the field at a site translated by a period $v$ is $\chi(v)$ times the field at the site, with $\chi(n_\mu e_\mu)=-1$, reads $\psi(c+n_\mu e_\mu,g)=-\psi(c,g)$, and the same for $\bar{\psi}$.
A fundamental domain of $\sum_\mu n_\mu\mathbb{Z}\,e_\mu$ in $\mathbb{Z}^3\times O_{\mathrm{cub}}$ represents the classes of the sites of $T$, and $\int_T$ is $\int_Y$ over it.
The identity $S_T(A)=\sum_x\bar{\psi}(x)\cdot\left((D_A+\boldsymbol{\mu})\psi\right)(x)$ is the display (46.7b) of Chapter~46; it rests on (b) and on this rule.
The sum over $\mathfrak{z}$ and the integral over $A$ are those of (f), and the curvature $dA+r_q\mathfrak{z}$ of \cite[Definition~2.13]{QEDII} is $F^{\mathrm{u}}$.
Inserted in the display of \cite[Definition~2.13]{QEDII}, which is one display for the three classes of its first two paragraphs, this gives the display of (h).
For $G=\Phi(F^{\mathrm{u}})$ the scalar $G$ leaves the Berezin integral, which is $\det(D_A+\boldsymbol{\mu})$ by (e).
No right side contains $v_0$.

For every real $\rho$ the display of (h) has a meaning on all three classes, coefficient by coefficient in a monomial basis of $\mathrm{Gr}_T$: the sum over $\mathfrak{z}$ and the integral over $A$ converge absolutely.
For $G\in\mathcal{A}(\mathbb{R}^3)$ that embeds in $T$ this is stated in the first paragraph of \cite[Definition~2.13]{QEDII}, for every real $\rho$.
For the two further classes one argues as follows.
By \cite[Lemma~3.13\,(b)]{QEDII}, which is stated for all real $\rho$ and $M$, the coefficients of $\exp\left(-S_T(A)\right)$ are polynomials in the link phases $\exp\left(iqA(b)\right)$ and in their complex conjugates;
hence they are bounded Borel functions on the configurations.
Let $G\in\mathrm{Ig}(T)$.
The definition of this class in \cite[Definition~5.11]{QEDII} contains no $\rho$:
$G$ is a Borel map of the pairs $(A,\mathfrak{z})$, invariant under the shifts $(A,\mathfrak{z})\mapsto\left(A+r_q\,m,\ \mathfrak{z}-dm\right)$ with integer valued one-cochains $m$ on the torus, and there are constants $\mathrm{Cq}<\infty$ and $\varepsilon_G\in\left(0,\frac{1}{2}\right]$ such that every coefficient of $G(A,\mathfrak{z})$ has modulus at most $\mathrm{Cq}\cdot\exp\left(\left(\frac{1}{2}-\varepsilon_G\right)\|F^{\mathrm{u}}\|^2\right)$.
The product of two monomials of $\mathrm{Gr}_T$ is $0$ or, up to a sign, a monomial.
Hence every coefficient of $\exp\left(-S_T(A)\right)\cdot G(A,\mathfrak{z})$ is a finite sum of products, each taken with a sign, of a coefficient of $\exp\left(-S_T(A)\right)$ and a coefficient of $G(A,\mathfrak{z})$,
and there is a constant $\mathrm{Cpr}$ such that every coefficient of
\[
\exp\left(-\frac{\|F^{\mathrm{u}}\|^2}{2}\right)\cdot\exp\left(-S_T(A)\right)\cdot G(A,\mathfrak{z})
\]
has modulus at most $\mathrm{Cpr}\cdot\exp\left(-\varepsilon_G\|F^{\mathrm{u}}\|^2\right)$.
Since $\|F^{\mathrm{u}}\|^2$ is the sum of the squares of the values of $F^{\mathrm{u}}$ on the plaquettes of the torus, the function $\exp\left(-\varepsilon_G\|F^{\mathrm{u}}\|^2\right)$ is the product over the plaquettes of the exponentials of $-\varepsilon_G$ times these squares, that is, the integrand of the second display of \cite[Lemma~3.14]{QEDII}, taken with $\varepsilon=\varepsilon_G$;
that display bounds the integral of this product over the configurations.
By that bound, which holds for every $\varepsilon>0$ and is stated for all real $\rho$ and $M$, this function is integrable over the configurations, that is, the sum over the closed $\mathfrak{z}$ of its integrals over $A\in[0,r_q)^{\mathrm{links}}$ is finite.
Finally, an element $G$ of the space spanned by the products read on $T$ enters the display through the function $(A,\mathfrak{z})\mapsto G(A,\mathfrak{z})$ of \cite[Lemma~5.14\,(a), (c)]{QEDII}, which lies in $\mathrm{Ig}(T)$;
this function is built from the cochains and the integrands alone, without an operator and hence without $\rho$, and the preceding argument applies to it.
\end{proof}

\begin{named}{Remark}
Proposition~46 corresponds to the display of the Wilson operator in \cite[Section~21.1.4]{QEDI} and to the sentence there that identifies this operator with the other definitions of the Wilson operator in \cite{QEDI}.
It differs from them in the following points.
The Wilson operator is replaced by one operator, given in the equal forms of (b).
The mass parameter $M_{\mathrm{ch}}$ is distinguished from the side $M$ of the $M$-cubes.
Clause (f+) has no counterpart there; in it the flux is labelled by the pairs $\mu<\nu$ in the order $(23,13,12)$, and the expression $r_q^2\sum_{\mu<\nu}|n|\,w_{\mu\nu}^2/(n_\mu n_\nu)^2$ of $\|h_w\|^2$ in \cite[Definition~13.7\,(d)]{QEDII}, which is $r_q^2\sum_{\mu<\nu}w_{\mu\nu}^2\,n_{\mathfrak{l}}/(n_\mu n_\nu)$ because $|n|=n_\mu n_\nu n_{\mathfrak{l}}$, equals the expression $r_q^2\sum_{\mu<\nu}n_{\mathfrak{l}}^2\,w_{\mu\nu}^2/|n|$ in \cite[Section~21.1.4]{QEDI}.
Clause (g+) states (g) for every $v\in\eta\mathbb{Z}^3$, in particular for every placement, because Section~I.7.6 forms $\hat{\mathfrak{a}}$ with $\tau_v\beta$; the independence of $v$ is Proposition~50\,(h).
Clause (h), which has no counterpart there either, reads the reflections, the classes of integrands and $Z_T[G]$ of \cite{QEDII} in the coordinates $(c,g)$, as Sections~I.4.8, I.6.4, I.7.2 and I.9.3 do.
\end{named}

\subsection*{I.3. Cluster maps, block averaging, operators of a level, fields, scaling, level 0}

\subsubsection*{I.3.1. The cluster maps}

Let $a\in\{\eta_k,1\}$.
The cluster maps $J\colon\mathfrak{S}(a,\ell)\to\mathfrak{S}^{\mathrm{ch}}(a,\ell)$ and $J^{*}\colon\mathfrak{S}^{\mathrm{ch}}(a,\ell)\to\mathfrak{S}(a,\ell)$ are defined by
\begin{equation}
(J\xi)(c,g):=\xi(c),\qquad
\left(J^{*}f\right)(c):=\frac{1}{48}\sum_{g\in O_{\mathrm{cub}}}f(c,g).
\tag{I.3.1}
\end{equation}
On $\mathfrak{S}^{\mathrm{ch}}(a,\ell)$ the measure, the pairings $\langle\cdot,\cdot\rangle^{\mathrm{ch}}$ and $(\cdot,\cdot)^{\mathrm{ch}}$, and the kernels (I.1.4) are formed with the weight $a^3/48$, that is, with
\[
\int\mathrm{d}x\,f(x):=a^3\cdot\frac{1}{48}\sum_{c}\sum_{g\in O_{\mathrm{cub}}}f(c,g),
\]
where $c$ ranges over a fundamental domain of $\ell\mathbb{Z}^3$ in $a\mathbb{Z}^3$;
for $a=\eta_k$ this is the measure (I.1.3) of Section~I.1.6.
In the averaging operators (I.3.4) the value $a=1$ occurs only at $k=0$, where $1=\eta_0$.
The cluster maps have the following properties.
\begin{itemize}
\item[(I.J1)]
One has $J^{*}\cdot J=1$.
For the pairings of Section~I.1.6 the map $J^{*}$ is at the same time the transpose and the adjoint of $J$:
\[
\left\langle f,J\xi\right\rangle^{\mathrm{ch}}=\left\langle J^{*}f,\xi\right\rangle_a .
\]
The map $J$ is an isometry of $L^2$.
\item[(I.J2)]
With respect to the sup norms, $J$ and $J^{*}$ have norm one; $J^{*}$ is a mean.
\item[(I.J3)]
The maps $J$ and $J^{*}$ commute with constant spin matrices, with the multiplication by functions of $c$, and with translations.
The same holds for every map defined later that does not mix the labels, or that permutes them:
for the scaling, (I.O9);
for the action of the symmetry groups, (I.a3) in Section~I.4.3;
for the periodic extension, Section~I.6.3;
and for the maps $\mathrm{ext}$ and $\mathrm{rs}$ of \cite[Section~16.6]{QEDI}, in which the label is a spectator:
\[
(\mathrm{ext}\,f)(u,g):=
\begin{cases}
f\left(\mathrm{Fo}(u),g\right)&\text{for }u\in\mathrm{Ch},\\
0&\text{otherwise},
\end{cases}
\qquad
(\mathrm{rs}\,v)(c,g):=v\left(\mathrm{emb}(c),g\right).
\]
\item[(I.J4)]
The operator $\mathrm{Plt}:=J\cdot J^{*}$ is the orthogonal projection onto the spinors that are constant on clusters (the light part), and $\mathrm{Phv}:=1-\mathrm{Plt}$ (the heavy part).
\item[(I.J5)]
If $O(c,c')$ is the kernel of an operator $O$ on $\mathfrak{S}(a,\ell)$ with respect to $a^3\sum_{c'}$, then the kernel (I.1.4) of $J\cdot O\cdot J^{*}$ is $O(c,c')$, constant on the two clusters.
\item[(I.J6)]
In the Propositions~1.m the blocks of $D_A$ are written
\begin{align*}
D^{\mathrm{ll}}_{A}&:=J^{*}\cdot D_A\cdot J,&
D^{\mathrm{lh}}_{A}&:=J^{*}\cdot D_A\cdot\mathrm{Phv},\\
D^{\mathrm{hl}}_{A}&:=\mathrm{Phv}\cdot D_A\cdot J,&
D^{\mathrm{hh}}_{A}&:=\mathrm{Phv}\cdot D_A\cdot\mathrm{Phv},
\end{align*}
and
\[
\mathrm{Dsc}_{A}:=D^{\mathrm{ll}}_{A}-D^{\mathrm{lh}}_{A}\cdot\left(D^{\mathrm{hh}}_{A}\right)^{-1}\cdot D^{\mathrm{hl}}_{A}
\]
is the Schur complement, $\left(D^{\mathrm{hh}}_{A}\right)^{-1}$ denoting the inverse of $D^{\mathrm{hh}}_{A}$ on the range of $\mathrm{Phv}$.
For the operator $D_A+\boldsymbol{\mu}$, which contains the mass $\boldsymbol{\mu}$, the four blocks are $D^{\mathrm{ll}}_{A}+\boldsymbol{\mu}$, $D^{\mathrm{lh}}_{A}$, $D^{\mathrm{hl}}_{A}$ and $D^{\mathrm{hh}}_{A}+\boldsymbol{\mu}\cdot\mathrm{Phv}$, and its Schur complement is
\[
D^{\mathrm{ll}}_{A}+\boldsymbol{\mu}-D^{\mathrm{lh}}_{A}\cdot\left(D^{\mathrm{hh}}_{A}+\boldsymbol{\mu}\cdot\mathrm{Phv}\right)^{-1}\cdot D^{\mathrm{hl}}_{A},
\]
the inverse being taken, again, on the range of $\mathrm{Phv}$.
At $\boldsymbol{\mu}=0$ this expression is $\mathrm{Dsc}_{A}$, which contains no mass.
Where both inverses exist, the resolvent identity
\[
\left(D^{\mathrm{hh}}_{A}\right)^{-1}-\left(D^{\mathrm{hh}}_{A}+\boldsymbol{\mu}\cdot\mathrm{Phv}\right)^{-1}=\boldsymbol{\mu}\cdot\left(D^{\mathrm{hh}}_{A}\right)^{-1}\cdot\left(D^{\mathrm{hh}}_{A}+\boldsymbol{\mu}\cdot\mathrm{Phv}\right)^{-1}
\]
on the range of $\mathrm{Phv}$ shows that the Schur complement of $D_A+\boldsymbol{\mu}$ is the sum of $\mathrm{Dsc}_{A}$ and
\[
\boldsymbol{\mu}\cdot\left[1+D^{\mathrm{lh}}_{A}\cdot\left(D^{\mathrm{hh}}_{A}\right)^{-1}\cdot\left(D^{\mathrm{hh}}_{A}+\boldsymbol{\mu}\cdot\mathrm{Phv}\right)^{-1}\cdot D^{\mathrm{hl}}_{A}\right];
\]
for $\boldsymbol{\mu}\ne0$ it is therefore a different expression, in general neither $\mathrm{Dsc}_{A}$ nor $\mathrm{Dsc}_{A}+\boldsymbol{\mu}$.
When the argument is a momentum, the index of $\mathrm{Dsc}$ is a mass:
the matrix $\mathrm{Dsc}_{\mathbf{m}}(p)$ of Proposition~1.2\,(S3) is formed in the same way from the blocks of the symbol $D^{(p)}$ of (I.2.6) at $\boldsymbol{\mu}=\mathbf{m}$.
\item[(I.O8)]
(Light and heavy part of $D_A$.)
For every complex $A$ and every fermion site $x=(c,g)$,
\[
\left(D_A\cdot J\xi\right)(c,g)=-\eta_k^{-1}\cdot M_{3}(x)\cdot\left[u_A(c,c_{3})\cdot\xi(c_{3})-\xi(c)\right],
\]
and, since every $a\in\{\pm e_\mu\}$ equals $g\,e_1$ for exactly $8$ of the $48$ labels $g$,
\begin{equation}
\begin{aligned}
&J^{*}\cdot D_A\cdot J=D^{W,\rho}_A,\\
&\left(D^{W,\rho}_A\xi\right)(c):=-\eta_k^{-1}\sum_{a\in\{\pm e_\mu\}}\frac{1}{2}\left(\rho-\gamma\cdot a\right)\cdot\left[u_A(c,c+\eta_ka)\cdot\xi(c+\eta_ka)-\xi(c)\right],
\end{aligned}
\tag{I.3.2}
\end{equation}
which is the Wilson operator with the Wilson parameter $\rho$, in the same field.
Let $\lambda_{\mathrm{int}}>0$ be the smallest non-zero eigenvalue of the operator $\sum_{i\le2}w_i\cdot\left(1-R_i\right)$ on $\ell^2(O_{\mathrm{cub}})$, where $(R_iu)(g):=u(g\,s_i)$;
this operator is positive semidefinite, and, since $s_0$, $s_1$, $s_2$ generate $O_{\mathrm{cub}}$, its kernel consists of the constant functions;
numerically, $\lambda_{\mathrm{int}}=1.99092\ldots$.
Put $\hat{g}:=\rho\cdot\lambda_{\mathrm{int}}/2$.
For real $A$ and $f=\mathrm{Phv}\,f$,
\[
\operatorname{Re}\left(f,D_Af\right)^{\mathrm{ch}}\ge\frac{\hat{g}}{\eta_k}\cdot(f,f)^{\mathrm{ch}}.
\]
In the fibre $p$,
\[
J^{*}\cdot D^{(p)}\cdot J=\boldsymbol{\mu}+i\gamma\cdot\sin p+\rho\sum_{\mu}\left(1-\cos p_\mu\right),\qquad\sin p:=\left(\sin p_\mu\right)_\mu ;
\]
at $p=0$ and $\boldsymbol{\mu}=0$ one has $D^{(0)}\cdot J=0$ and $J^{*}\cdot D^{(0)}=0$.
\end{itemize}

\subsubsection*{I.3.2. Phases of the averages (photon objects)}

For $c\in\Delta_y$ let $c=y_0,y_1,\dots,y_k=y$ be the centres of the blocks of the sides $\eta_k,L\eta_k,\dots,1$ that contain $c$, the paths being taken on the lift, with $|y_{j+1}-y_j|<1/2$.
Put
\begin{equation}
\begin{aligned}
(\tau_kA)(y,c)&:=\sum_{j<k}(\tau A)(y_{j+1},y_j),\qquad
\varphi_A(y,c):=\exp\left(i\,e_k\cdot\eta_k\cdot(\tau_kA)(y,c)\right)\quad(c\in\Delta_y),\\
\varphi^{+}_A(y',c)&:=U_A(y',y)\cdot\varphi_A(y,c)\quad\left(c\in\Delta_y,\ y\in B(y')\right).
\end{aligned}
\tag{I.3.3}
\end{equation}
In $U_A(y',y)$, where $y'\in L\mathbb{Z}^3$ and $y\in B(y')$, the paths lie on the unit lattice, and $A$ is summed over their fine bonds.
The phase $\varphi_A(y,c)$ is the phase $\boldsymbol{\phi}_A(y,c)$ of \cite[Section~16.2.4]{QEDI};
it is written $\phi_{B,y}(c)$ in \cite[Section~17.1.2]{QEDI} and $\boldsymbol{\phi}^B_k(y,x)$ in \cite[Section~14.6.2]{QEDI}.
At the stop level, where the function $\sigma_A$ of \cite[Section~20.2.2]{QEDI} is formed, one has $\varphi_A(y,c)=\exp\left(i\,\sigma_A(y,c)\right)$.

\subsubsection*{I.3.3. Averaging operators}

The operators $Q_k(A)$, $Q_k^{T}(-A)$, $P_k(A)$, $Q(A)$, $Q^{T}(-A)$, $P(A)$, $P_{k+1}(A)$ have the meaning of \cite[Sections~5.2.6, 5.3.6 and 16.2.4]{QEDI}:
they act between spaces of $4$-component spinors,
\begin{align*}
&\left(Q_k(A)\xi\right)(y):=\eta_k^3\sum_{c\in\Delta_y}\varphi_A(y,c)\cdot\xi(c),\\
&\left(Q_k^{T}(-A)v\right)(c):=\varphi_A(y,c)^{-1}\cdot v(y)\quad(c\in\Delta_y),\\
&P_k(A):=Q_k^{T}(-A)\cdot Q_k(A)\qquad\left(\text{between }\mathfrak{S}(\eta_k,\ell)\text{ and }\mathfrak{S}(1,\ell)\right);\\
&\left(Q(A)v\right)(y'):=L^{-3}\sum_{y\in B(y')}U_A(y',y)\cdot v(y),\\
&\left(Q^{T}(-A)v'\right)(y):=U_A(y',y)^{-1}\cdot v'(y')\quad\left(y\in B(y')\right),\\
&P(A):=Q^{T}(-A)\cdot Q(A)\qquad\left(\text{between }\mathfrak{S}(1,\ell)\text{ and }\mathfrak{S}(L,\ell)\right);\\
&P_{k+1}(A):=Q_k^{T}(-A)\cdot P(A)\cdot Q_k(A).
\end{align*}
For an argument of either sign, $Q_k^{T}(A)$ and $Q^{T}(A)$ denote the operators given by the formulas for $Q_k^{T}(-A)$ and $Q^{T}(-A)$ with $-A$ in the place of $A$:
\[
\left(Q_k^{T}(A)v\right)(c)=\varphi_{-A}(y,c)^{-1}\cdot v(y),\qquad
\left(Q^{T}(A)v'\right)(y)=U_{-A}(y',y)^{-1}\cdot v'(y').
\]
The meaning of $P_{k+1}(\cdot)$ is determined by its argument.
For a connection $A$ on $T_k$ it is the operator $Q_k^{T}(-A)\cdot P(A)\cdot Q_k(A)$ on $\mathfrak{S}(\eta_k,\ell)$ defined above.
For a connection $A'$ of level $k+1$, that is, on $T_{k+1}=\mathbb{L}_{k+1}/(\ell/L)\mathbb{Z}^3$ (Section~I.3.7),
\[
P_{k+1}(A'):=Q_{k+1}^{T}(-A')\cdot Q_{k+1}(A')\qquad\text{on }\mathfrak{S}(\eta_{k+1},\ell/L)
\]
is the operator $P_k(\cdot)$ of the level $k+1$.
None of the seven symbols $Q_k(A)$, $Q_k^{T}(-A)$, $P_k(A)$, $Q(A)$, $Q^{T}(-A)$, $P(A)$, $P_{k+1}(A)$ contains a cluster map:
for $k\ge1$ they are the operators between spaces of $4$-component spinors displayed above, and at level $0$ one has $Q_0=P_0=1$ by (I.3.4) and $P_1(A)=P(A)$ on $\mathfrak{S}(1,\ell)$ by Section~I.3.9.
In a product of one of them with $J$ or $J^{*}$ the cluster maps are written out, as in (I.3.4).
An operator formed with $J$ or $J^{*}$ carries a symbol of its own only through an explicit definition, given where the symbol is introduced.
Among such symbols are $\mathrm{Plt}=J\cdot J^{*}$ and $\mathrm{Phv}$ in (I.J4);
the blocks of $D_A$ in (I.J6);
$\operatorname{Pen}_y(A)$, $H_k(A)$, $\mathfrak{H}_k(A)$, $T_k(A)$, $D_k(A)$ and $\dot{M}_k(A)$ in Section~I.3.5, and $\operatorname{Pen}_y(A)$ and $\Gamma_0(A)$ of level $0$ in Section~I.3.9;
the abbreviations $\mathrm{Qk}:=Q_k(A)\cdot J^{*}$ and $\mathrm{Qt}:=J\cdot Q_k^{T}(-A)$ in Proposition~4 and in Chapters~4 and~5 (Section~G1.10);
and the same abbreviations with $B$ in the place of $A$ in Chapter~21a.
The averaging operators of the scheme are
\begin{equation}
\begin{aligned}
k\ge1:&\quad Q_k(A)\cdot J^{*},\quad J\cdot Q_k^{T}(-A),\quad J\cdot P_k(A)\cdot J^{*},\quad J\cdot P_{k+1}(A)\cdot J^{*}\quad\text{on the fine side};\\
&\quad Q(A),\quad Q^{T}(-A),\quad P(A)\quad\text{between the unit lattice and the $L$-lattice,}\\
&\quad\text{without cluster maps};\\
k=0:&\quad Q_0:=1,\quad P_0:=1\quad\text{on }\mathfrak{S}^{\mathrm{ch}}(1,\ell);\qquad Q(A)\cdot J^{*},\quad J\cdot Q^{T}(-A),\quad J\cdot P(A)\cdot J^{*}.
\end{aligned}
\tag{I.3.4}
\end{equation}
Explicitly,
\[
\left(Q_k(A)\cdot J^{*}f\right)(y)=\int_{\Delta_y}\mathrm{d}x\,\varphi_A(y,c)\cdot f(x),\qquad
\left(J\cdot Q_k^{T}(-A)v\right)(x)=\varphi_A(y,c)^{-1}\cdot v(y)
\]
for $x=(c,g)$ with $c\in\Delta_y$:
the average extends over the photon block and over the $48$ sites of each cluster.
\begin{itemize}
\item[(I.Q1)]
By (I.J1),
\[
\left(Q_k(A)\cdot J^{*}\right)\cdot\left(J\cdot Q_k^{T}(-A)\right)=1,\qquad
\left(Q(A)\cdot J^{*}\right)\cdot\left(J\cdot Q^{T}(-A)\right)=1 .
\]
Hence $J\cdot P_k(A)\cdot J^{*}$, $J\cdot P_{k+1}(A)\cdot J^{*}$ and $J\cdot P(A)\cdot J^{*}$ are idempotent; their kernels are multiples of $1_4$.
The weight $1/48$ in (I.1.3) is forced by (I.Q1): with the counting weight these products would equal $1/48$ or $48$, and the numbers of Section~I.3.4 would change.
\item[(I.Q2)]
For complex $A$,
\[
\left(Q_k(A)\cdot J^{*}\right)^{T}=J\cdot Q_k^{T}(A),\qquad
\left(Q(A)\cdot J^{*}\right)^{T}=J\cdot Q^{T}(A);
\]
for real $A$ one has $\left(Q_k(A)\cdot J^{*}\right)^{*}=J\cdot Q_k^{T}(-A)$, and $J\cdot P_k(A)\cdot J^{*}$ and $J\cdot P(A)\cdot J^{*}$ are orthogonal projections that commute with all spin matrices.
\item[(I.Q3)]
The sup norms of $Q_k(A)\cdot J^{*}$ and of $J\cdot Q_k^{T}(-A)$ are at most $\exp\left(\frac{3}{2}\,\mathfrak{i}(A)\right)$, by \cite[Lemma~16.2(a)]{QEDI} and (I.J2).
\end{itemize}

\subsubsection*{I.3.4. Numbers of the blocking}

The numbers of the blocking are
\begin{gather*}
b_k:=\frac{b\cdot\left(1-L^{-1}\right)}{1-L^{-k}}\quad\left(k\ge1;\ b_1=b,\ b_{k+1}=\frac{b\cdot b_k}{b_k+b_L}\right),\\
b_-:=b\cdot\left(1-L^{-1}\right)\le b_k\le b_+:=b;\qquad
b_0\text{ is not defined};\\
n_{k+1}:=n_k/L^3;\qquad
c_k:=b^{-4n_{k+1}},\quad c'_k:=b^{4n_{k+1}}\quad(k\ge0);\\
c^{\mathrm{d}}_k:=4\left[\left(1-L^{-3}\right)\cdot\log b_k+L^{-3}\cdot\log\left(b_k+b_L\right)\right]\quad(k\ge1),\qquad
c^{\mathrm{d}}_0:=0 .
\end{gather*}
These numbers count the block spinors of the levels $k\ge1$, which keep $4$ components per unit site.
The number $b_L=b\cdot L^{-1}$ is written $b\cdot L^{-1}$ in \cite[Section~5.2.2]{QEDI} and $\tilde{b}$ in \cite[Section~14.6.2]{QEDI}.

\subsubsection*{I.3.5. The operators of a level $k\ge1$ (data $(k,\ell,\mu,e_k)$; connection $A$ on $T_k$)}

The resolvent parameter is $y\in Y_k:=\mathbb{R}\cup\{\infty\}$, with $-\infty:=\infty$.
As functions of the matrix $\gamma_5$ put
\begin{gather*}
\boldsymbol{\alpha}_y:=b_k\cdot i\gamma_5\cdot y\cdot\left(b_k+i\gamma_5\cdot y\right)^{-1},\qquad
\boldsymbol{\beta}_y:=b_k^2\cdot b_L\cdot\left(b_k+b_L+i\gamma_5\cdot y\right)^{-1}\cdot\left(b_k+i\gamma_5\cdot y\right)^{-1},\\
\boldsymbol{\alpha}_\infty:=b_k,\qquad\boldsymbol{\beta}_\infty:=0\qquad\left(\boldsymbol{\alpha}_0=0,\ \boldsymbol{\beta}_0=b_{k+1}/L\right),
\end{gather*}
and define
\begin{align*}
\operatorname{Pen}_y(A)&:=\boldsymbol{\alpha}_y\cdot J\cdot P_k(A)\cdot J^{*}+\boldsymbol{\beta}_y\cdot J\cdot P_{k+1}(A)\cdot J^{*},\qquad
F_y(A):=D_A+\mu+\operatorname{Pen}_y(A),\\
S_y(A)&:=F_y(A)^{-1},\qquad w_k(y):=1;\\
S_k(A)&:=S_\infty(A)=\left(D_A+\mu+b_k\cdot J\cdot P_k(A)\cdot J^{*}\right)^{-1},\\
S_{k,y}(A)&:=S_y(A),\qquad
S^{0}_{k+1}(A):=S_0(A);\\
H_k(A)&:=b_k\cdot S_k(A)\cdot J\cdot Q_k^{T}(-A),\qquad
\mathfrak{H}_k(A):=b_k\cdot Q_k(A)\cdot J^{*}\cdot S_k(A),\\
T_k(A)&:=b_k^{-1}\cdot Q_k(A)\cdot J^{*}\cdot\left(D_A+\mu+b_k\cdot J\cdot P_k(A)\cdot J^{*}\right)\\
&\phantom{:}=b_k^{-1}\cdot Q_k(A)\cdot J^{*}\cdot\left(D_A+\mu\right)+Q_k(A)\cdot J^{*};\\
D_k(A)&:=b_k-b_k^2\cdot Q_k(A)\cdot J^{*}\cdot S_k(A)\cdot J\cdot Q_k^{T}(-A),\qquad
M_k(A):=b_k-D_k(A),\\
\dot{M}_k(A)&:=b_k^2\cdot Q_k(A)\cdot J^{*}\cdot S_k(A)^2\cdot J\cdot Q_k^{T}(-A);\\
\Gamma_k(A)&:=\left(D_k(A)+b_L\cdot P(A)\right)^{-1},\qquad
\delta Z_k(A):=\det\left(D_k(A)+b_L\cdot P(A)\right),\\
\Gamma_{k,y}(A)&:=\left(D_k(A)+b_L\cdot P(A)+i\gamma_5\cdot y\right)^{-1}\quad(\text{for real }y),\\
B_{k,y}(A)&:=\left(b_k+b_L\cdot P(A)+i\gamma_5\cdot y\right)^{-1}\quad(\text{for real }y),\\
B_k(A)&:=B_{k,0}(A);\\
\bar{H}_k(A)&:=\mathsf{C}^{-1}\cdot H_k(-A)\cdot\mathsf{C},\qquad
\bar{T}_k(A):=\mathsf{C}^{-1}\cdot T_k(-A)\cdot\mathsf{C}.
\end{align*}
Further, $H_k^{\#}(A):=\left(H_k(A),\,\delta_{\alpha,A}H_k(A)\right)$, and on barred fields $H_k^{\#}$ stands for the pair $\left(\bar{H}_k(A),\,\delta_{\alpha,-A}\bar{H}_k(A)\right)$;
$\left|H_k^{\#}(A)v\right|_{\infty}$ denotes the larger of the two numbers
\[
\left|H_k(A)v\right|_{\infty}
\qquad\text{and}\qquad
\sup_{(x,x')\in\mathrm{Pzt}_k}\left|\left(\delta_{\alpha,A}H_k(A)v\right)(x,x')\right|_c ,
\]
and the same definition is made for the barred pair.
These are the operators of \cite[Sections~5.2.6, 14.6.2 and 16.2.5]{QEDI}, formed with the operator $D_A$ of (I.2.3), with the averaging operators (I.3.4), and with $\gamma_5$ in the place of the matrix $\gamma_3$ of \cite[Section~5.2.4]{QEDI}.
The existence of the inverses is established in the Propositions~2.m and in Proposition~4.
The operator $X_k(B)$ of \cite[Section~14.6.2]{QEDI} is here $F_\infty(B)$.
The operators act between the following spaces.

\noindent
\begin{tabular}{p{0.33\textwidth}p{0.17\textwidth}p{0.17\textwidth}p{0.19\textwidth}}
operator & from & to & components\\
\hline
$D_A$, $\operatorname{Pen}_y(A)$, $F_y(A)$, $S_k(A)$, $S_{k,y}(A)$, $S^{0}_{k+1}(A)$, $J\cdot P_k(A)\cdot J^{*}$, $J\cdot P_{k+1}(A)\cdot J^{*}$
& $\mathfrak{S}^{\mathrm{ch}}(\eta_k,\ell)$ & $\mathfrak{S}^{\mathrm{ch}}(\eta_k,\ell)$ & $192$ per photon site\\
$H_k(A)$, $\bar{H}_k(A)$, $J\cdot Q_k^{T}(-A)$
& $\mathfrak{S}(1,\ell)$ & $\mathfrak{S}^{\mathrm{ch}}(\eta_k,\ell)$ & $4\to192$\\
$\delta_{\alpha,A}$
& $\mathfrak{S}^{\mathrm{ch}}(\eta_k,\ell)$ & functions of pairs (I.1.7) & --\\
$T_k(A)$, $\bar{T}_k(A)$, $\mathfrak{H}_k(A)$, $Q_k(A)\cdot J^{*}$
& $\mathfrak{S}^{\mathrm{ch}}(\eta_k,\ell)$ & $\mathfrak{S}(1,\ell)$ & $192\to4$\\
$D_k(A)$, $M_k(A)$, $\dot{M}_k(A)$, $\Gamma_k(A)$, $\Gamma_{k,y}(A)$, $B_{k,y}(A)$, $P(A)$
& $\mathfrak{S}(1,\ell)$ & $\mathfrak{S}(1,\ell)$ & $4$ per unit site; $\delta Z_k(A)$ is a determinant of size $4n_k$\\
$Q(A)$; $Q^{T}(-A)$
& $\mathfrak{S}(1,\ell)$; $\mathfrak{S}(L,\ell)$ & $\mathfrak{S}(L,\ell)$; $\mathfrak{S}(1,\ell)$ & $4$
\end{tabular}


\subsubsection*{I.3.6. Fields: generators, their index sets and weights; the Berezin integral}

The independent Grassmann generators of level $k$ are listed in the following table.
Here $\beta\in\{1,\dots,4\}$, $\omega=0$ for unbarred and $\omega=1$ for barred generators, and for a family $\phi$ of generators $\phi_\beta(t):=\phi(t,\beta,0)$ and $\bar{\phi}_\beta(t):=\phi(t,\beta,1)$.

\noindent
\begin{tabular}{p{0.25\textwidth}p{0.2\textwidth}p{0.13\textwidth}p{0.11\textwidth}p{0.17\textwidth}}
generators & index & range & weight $\nu$ & components\\
\hline
$\psi(\xi)$ & $\xi=(x,\beta,\omega)=(c,g,\beta,\omega)$ & $x\in T^{\mathrm{f}}_k$ & $\eta_k^3/48$ & $192$ per photon site, every $k\ge0$\\
$\chi(\boldsymbol{\zeta})$ & $\boldsymbol{\zeta}=(x,x',\beta,\omega)$ & $(x,x')\in\mathrm{Pz}_k$ & $\left(\eta_k^3/48\right)^2$ & --\\
$\Psi(u)$, $W(u)$, $k\ge1$ & $u=(y,\beta,\omega)$ & $y\in T^0_k$ & $1$ & $4$ per unit site\\
$\Psi(u)$, $W(u)$, $k=0$ & $u=(x,\beta,\omega)$ & $x\in T^{\mathrm{f}}_0$ & $1/48$ & $192$ per unit site\\
$\Psi'(u')$, the block field of level $k+1$ in the step $k\to k+1$, every $k\ge0$ & $u'=(y',\beta,\omega)$ & $y'\in T^0_{k+1}$ & $1$ & $4$ per unit site
\end{tabular}

\noindent
The integrals over the index sets are the weighted sums
\begin{gather*}
\int\mathrm{d}\xi\,F(\xi):=\sum_{\xi}\nu(\xi)\cdot F(\xi),\qquad
\int\mathrm{d}\boldsymbol{\zeta}\,F(\boldsymbol{\zeta}):=\sum_{\boldsymbol{\zeta}}\nu(\boldsymbol{\zeta})\cdot F(\boldsymbol{\zeta}),\\
\int\mathrm{d}u\,F(u):=\sum_{u}\nu(u)\cdot F(u),
\end{gather*}
and $\psi^{\#}:=(\psi,\chi)$.
The symbols $\mathrm{Gr}[\psi,\chi]$, $\mathrm{Gr}[\Psi]$, \dots\ denote the Grassmann algebras over $\mathbb{C}$ generated by the fields named.
A homomorphism (automorphism) of Grassmann algebras is a homomorphism of rings with unit;
it is linear over $\mathbb{C}$, except for $\mathbf{K}$ of Section~I.4.5 and the compositions with it, which are antilinear.
\begin{itemize}
\item[(I.AP)]
The generators are indexed by the points of the covering lattices, with
\[
\psi(x+\ell e_\mu,\cdot)=-\psi(x,\cdot),\qquad
\chi(x+\ell e_\mu,x'+\ell e_\mu,\cdot)=-\chi(x,x',\cdot),
\]
and in the same way for $\Psi$ and $W$;
independent generators are obtained by restricting $x$, the first point, to a fundamental domain.
Kernels are antiperiodic in each argument;
their moduli and all norms do not depend on the fundamental domain;
even elements do not change under $\psi\to-\psi$ (\cite[Section~13.1.3]{QEDI}, with $x=(c,g)$).
Kernels that are integrated without generators (in the extraction and in the adjustment, Section~I.5.6) are evaluated with all arguments over one lift of the small polymer.
\end{itemize}

\paragraph{Forms.}
Put
\begin{gather*}
\left\langle\bar{\psi},\psi\right\rangle:=\int\mathrm{d}x\,\bar{\psi}(x)\cdot\psi(x),\qquad
\left\langle\bar{\psi},\psi\right\rangle_{\square}:=\int_{\square}\mathrm{d}x\,\bar{\psi}(x)\cdot\psi(x);\\
\left\langle\bar{\Psi},O\Psi\right\rangle:=\sum_{y,y'}\bar{\Psi}(y)\cdot O(y,y')\cdot\Psi(y')\quad\text{for }k\ge1,\\
\left\langle\bar{\Psi},O\Psi\right\rangle:=\left\langle\bar{\Psi},O\Psi\right\rangle^{\mathrm{ch}}=\int\mathrm{d}x\,\bar{\Psi}(x)\cdot(O\Psi)(x)\quad\text{for }k=0.
\end{gather*}

\paragraph{Berezin integral on an index set with weights.}
The following definition serves for all levels and for every finite index set.
Let $W_\beta(t)$, $\bar{W}_\beta(t)$ be generators indexed by a finite set $\{(t,\beta)\}$, with weights $\nu(t)>0$, and put
\[
\left\langle\bar{W},W\right\rangle_\nu:=\sum_{t,\beta}\nu(t)\cdot\bar{W}_\beta(t)\cdot W_\beta(t).
\]
For a linear map $O$ of the functions on $\{(t,\beta)\}$, with the kernel given by $(Of)(t)=\sum_{t'}\nu(t')\cdot O(t,t')\cdot f(t')$, put
\[
\left\langle\bar{W},OW\right\rangle_\nu:=\sum_{t}\nu(t)\cdot\bar{W}(t)\cdot(OW)(t),
\]
and let $\mathrm{Top}:=\prod_{(t,\beta)}W_\beta(t)\cdot\bar{W}_\beta(t)$; the factors are even, and their order is immaterial.
Define
\begin{equation}
\begin{aligned}
\int\mathrm{D}W\,F&:=\left(\prod_{(t,\beta)}\nu(t)^{-1}\right)\cdot\left(\text{the coefficient of }\mathrm{Top}\text{ in }F\right),\\
\mathrm{d}\mu_I(W)&:=\exp\left(-\left\langle\bar{W},W\right\rangle_\nu\right)\mathrm{D}W.
\end{aligned}
\tag{I.3.5}
\end{equation}
In an algebra with further generators, the coefficient of $\mathrm{Top}$ in $F$ is the element $[F]_{\mathrm{Top}}$ of the subalgebra generated by the generators outside the family for which
\[
F=[F]_{\mathrm{Top}}\cdot\mathrm{Top}+\left(\text{terms of lower degree in the family}\right);
\]
since $\mathrm{Top}$ is even, left and right coefficients coincide, and $\int\mathrm{D}W$ is a linear map into this subalgebra (it is the map of Proposition~15\,(D3)).
Then the following hold.
\begin{itemize}
\item[(I.B1)]
$\int\mathrm{D}W\,\exp\left(-\left\langle\bar{W},OW\right\rangle_\nu\right)=\det O$ for every $O$, where $\det O$ is the determinant of the linear map $O$.
\item[(I.B2)]
$\int\mathrm{d}\mu_I(W)=1$.
\item[(I.B3)]
$\int\mathrm{d}\mu_I(W)\,W_\beta(t)\cdot\bar{W}_{\beta'}(t')=\mathbf{1}[t=t',\,\beta=\beta']/\nu(t)$.
\item[(I.B4)]
For a set $J$ of generators of the family let $W(J)$ be the product of the elements of $J$ in a fixed order; another order changes $W(J)$ at most by a sign.
Then $\int\mathrm{d}\mu_I\,W(J)$ vanishes unless $J$ is a union of pairs $\{W_\beta(t),\bar{W}_\beta(t)\}$, and it is then equal to $\pm\prod\nu(t)^{-1}$, the product extending over these pairs.
When the generators are identified with their indices, $W(J)$ is the monomial $W(J)$ of Proposition~15\,(D2), (D3), (e), and a union of pairs is a paired set in the sense of Proposition~15\,(D2).
\end{itemize}
\begin{proof}
For one mode, $\exp\left(-\nu\cdot\bar{W}\cdot W\right)=1+\nu\cdot W\cdot\bar{W}$.
Hence
\[
\exp\left(-\left\langle\bar{W},W\right\rangle_\nu\right)=\prod_{(t,\beta)}\left(1+\nu(t)\cdot W_\beta(t)\cdot\bar{W}_\beta(t)\right),
\]
and the product of this element with a monomial $W(J)$ has a non-zero component along $\mathrm{Top}$ only if $J$ is a union of pairs $\{W_\beta(t),\bar{W}_\beta(t)\}$;
this component is then $\pm$ the product of the weights $\nu(t)$ of the pairs that do not belong to $J$, times $\mathrm{Top}$, with the sign $+$ if $W(J)$ is a product of the pairs $W_\beta(t)\cdot\bar{W}_\beta(t)$.
After multiplication by $\prod_{(t,\beta)}\nu(t)^{-1}$ this gives (I.B2), (I.B3) and (I.B4).
The matrix of the form $\left\langle\bar{W},OW\right\rangle_\nu$ is $\operatorname{diag}(\nu)$ times the matrix of the map $O$.
The coefficient of $\mathrm{Top}$ in $\exp\left(-\left\langle\bar{W},OW\right\rangle_\nu\right)$ is the determinant of the matrix of the form, and the factor $\prod_{(t,\beta)}\nu(t)^{-1}$ cancels $\det\operatorname{diag}(\nu)$; this gives (I.B1).
\end{proof}
For $\nu\equiv1$ the map $\int\mathrm{D}W$ is the map $\int\mathrm{D}W$ of \cite[Section~13.2]{QEDI} and the Berezin integral of \cite[Definition~2.8]{QEDII}.
The normalisations of \cite[Section~13.2]{QEDI} and of \cite[Definition~2.8]{QEDII} agree: the coefficient of $\psi\bar{\psi}$ is $1$, and $\int\exp\left(-\bar{\psi}\,\mathrm{Dm}\,\psi\right)=\det\mathrm{Dm}$ for a matrix $\mathrm{Dm}$.
For $\nu\equiv1/48$ the map $\int\mathrm{D}W$ is the integral $\int\mathrm{D}\Psi$ of Proposition~46\,(e).

\subsubsection*{I.3.7. Scaling (step $k\to k+1$; side $\ell$ at level $k$, side $\ell/L$ at level $k+1$)}

Put $x/L:=(c/L,g)$; this is a bijection $\mathbb{L}^{\mathrm{f}}_k\to\mathbb{L}^{\mathrm{f}}_{k+1}$.
One has $e_{k+1}=L^{\frac{1}{2}}e_k$, and the mass $\mu$ at level $k$ becomes $L\mu$ at level $k+1$.
Let
\begin{align*}
&A_L(b):=L^{-\frac{1}{2}}A(b/L),\qquad
f_L(x):=L^{-1}f(x/L),\\
&\psi_L(x):=L^{-1}\psi(x/L),\qquad
\chi_L(x,x'):=L^{-1-\alpha}\chi(x/L,x'/L),\\
&\Psi'_L(y):=L^{-1}\Psi'(y/L)\quad\left(y\in L\mathbb{Z}^3/\ell\mathbb{Z}^3\right),\\
&v_L(y'):=L^{-1}v(y'/L)\quad\left(v\in\mathfrak{S}(1,\ell/L),\ y'\in L\mathbb{Z}^3;\ v_L\in\mathfrak{S}(L,\ell)\right),\\
&\xi_L(c):=L^{-1}\xi(c/L)\quad\left(\xi\in\mathfrak{S}(\eta_{k+1},\ell/L),\ c\in\mathbb{L}_k;\ \xi_L\in\mathfrak{S}(\eta_k,\ell)\right),\\
&F_L(x,x'):=L^{-1-\alpha}F(x/L,x'/L)\quad(\text{functions of pairs}).
\end{align*}
The label is not affected.
One has $(x/L)^{(i)}=x^{(i)}/L$ and $M_{i}(x/L)=M_{i}(x)$;
$d_X(x/L,x'/L)$, formed at level $k+1$, equals $d_X(x,x')/L$, because $\eta_{k+1}=\eta_k/L$;
a pair of $\mathrm{Pz}_k$ divided by $L$ is a pair of $\mathrm{Pz}_{k+1}$ with $d_X\le M/L$;
a pair that is admissible for $X\in D_k$, where $X\subset LY$ and $Y\in D_{k+1}$, is mapped to a pair of the same kind that is admissible for $Y$ (Proposition~18.3\,(e));
and $\eta_k^3/48=L^3\eta_{k+1}^3/48$.
Let $S_L$ be the homomorphism of $\mathrm{Gr}[\psi,\chi]$ of level $k$ into $\mathrm{Gr}[\psi,\chi]$ of level $k+1$ with
\[
S_L\psi(x,\beta,\omega):=L^{-1}\psi(x/L,\beta,\omega),\qquad
S_L\chi(x,x',\beta,\omega):=L^{-1-\alpha}\chi(x/L,x'/L,\beta,\omega);
\]
thus $S_LF(\psi,\chi)=F(\psi_L,\chi_L)$, and
\[
(LE)(Y;\mathscr{A},\mathscr{C})=\sum_{X:\,\bar{X}^{L}=LY}S_L\left[E(X;\mathscr{A}_L,\mathscr{C}_L)\right],
\]
which is the map $L$ of Section~I.5.6 on the families of Section~I.5.4.
The definition of $S_L$ requires only that $\ell/L$ be a multiple of $M$ larger than $2M$ (Proposition~16.1\,(d)).
For a function $\Lambda$ of the sites of $T_{k+1}$ put $\Lambda_L(c):=L^{\frac{1}{2}}\Lambda(c/L)$; then $(d\Lambda)_L=d(\Lambda_L)$.
\begin{itemize}
\item[(I.O9)]
For $A=A'_L$ and $f=f'_L$,
\begin{gather*}
D_Af'_L=L^{-1}\left(D_{A'}f'\right)_L,\qquad
\delta_{\alpha,A}f'_L=\left(\delta_{\alpha,A'}f'\right)_L,\\
\left\langle f''_L,f'_L\right\rangle^{\mathrm{ch}}\text{ at level }k=L\cdot\left\langle f'',f'\right\rangle^{\mathrm{ch}}\text{ at level }k+1,\\
(J\xi)_L=J\xi_L,\qquad\left(J^{*}f\right)_L=J^{*}f_L .
\end{gather*}
\end{itemize}
\begin{proof}
One has $e_k\cdot\eta_k\cdot A'_L(b)=e_{k+1}\cdot\eta_{k+1}\cdot A'(b/L)$ and $\eta_k^{-1}=L^{-1}\eta_{k+1}^{-1}$, and the matrices $M_i$ contain no mesh.
Hence the phases of $A$ at $(x,x')$ are those of $A'$ at $(x/L,x'/L)$, and the first two identities follow from (I.2.3) and (I.1.7), since $d_X(x,x')=L\cdot d_X(x/L,x'/L)$.
The third follows from $\eta_k^3/48=L^3\eta_{k+1}^3/48$, and the last two hold because the scaling does not affect the label.
\end{proof}

\subsubsection*{I.3.8. Levels}

``Level $k$'' means that the block fields live on the unit lattice $T^0_k$ and the fine fields on the lattices of spacing $\eta_k=L^{-k}$;
one unit of length of level $k$ is the physical length $L^{k-N}$ (denoted $a_k$ in Section~I.6.1).
The step $k\to k+1$ integrates out the fluctuation of the block fields of level $k$ at fixed block fields of level $k+1$, and rescales by $L$.
The step $0\to1$ is the first step: in it the number of components of the block spinor drops from $192$ to $4$.

\subsubsection*{I.3.9. Level 0, first part: operators and fields}

At $k=0$ one has $\eta_0=1$;
the fine lattice and the unit lattice coincide, and $T^{\mathrm{f}}_0=T^0_0\times O_{\mathrm{cub}}$;
for spinors, the ``unit block'' of a site $y$ is its cluster $\{y\}\times O_{\mathrm{cub}}$;
in the flow $\mu=\bar{m}\cdot\eta$;
and $e_k$ is read as $q$.
The objects of level $0$ are
\begin{align*}
&Q_0=P_0:=1,\qquad
H_0(A)=\bar{H}_0(A)=T_0(A)=\bar{T}_0(A):=1,\qquad
\dot{M}_0(A):=1\quad\text{on }\mathfrak{S}^{\mathrm{ch}}(1,\ell);\\
&\psi=\Psi;\qquad
D_0(A):=D_A+\mu;\\
&Y_0:=\mathbb{R},\qquad
\operatorname{Pen}_y(A):=b_L\cdot J\cdot P(A)\cdot J^{*}+i\gamma_5\cdot y,\\
&F_y(A):=D_A+\mu+\operatorname{Pen}_y(A),\qquad
S_y(A):=F_y(A)^{-1},\\
&w_0(y):=\left(1+|y|\right)^{-1};\\
&\Gamma_0(A)=S^{0}_1(A):=F_0(A)^{-1}=\left(D_A+\mu+b_L\cdot J\cdot P(A)\cdot J^{*}\right)^{-1},\qquad
S_{0,y}(A):=S_y(A),\\
&\delta Z_0(A):=\det F_0(A)\quad\left(\text{a determinant of size }192\,n_0\right).
\end{align*}
All these operators act on $\mathfrak{S}^{\mathrm{ch}}(1,\ell)$.
On the four-component space $\mathfrak{S}(1,\ell)$ put $P_1(A):=P(A)=Q^{T}(-A)\cdot Q(A)$;
the operator on $\mathfrak{S}^{\mathrm{ch}}(1,\ell)$ that takes the place of $J\cdot P_{k+1}(A)\cdot J^{*}$ at level $0$ is $J\cdot P_1(A)\cdot J^{*}=J\cdot P(A)\cdot J^{*}$, written with the cluster maps.
The objects $b_0$, $S_k(A)=S_\infty(A)$ at $k=0$, $\mathfrak{H}_0(A)$, $M_0(A)$, $B_{0,y}(A)$ and $\Gamma_{0,y}(A)$ are not defined.
The operator $S_y(A)$ at $y=0$, written $S_0(A)$ in Section~I.3.5, is defined at level $0$: it is $\Gamma_0(A)$.
The determinant $\delta Z_0(A)$ is taken on the whole of $\mathfrak{S}^{\mathrm{ch}}(1,\ell)$, the range of $\mathrm{Phv}$ included.
All differences between level $0$ and the levels $k\ge1$ are listed in Section~I.5.10.

\subsection*{I.4. Symmetries}

\subsubsection*{I.4.1. The spin group $\mathrm{Sp}$}

Let $\mathrm{Sp}$ be the group generated by $\gamma_1$, $\gamma_2$, $\gamma_3$, $\gamma_5$;
it has $32$ elements, namely $\pm$ the products of distinct generators.
Define $r\colon\mathrm{Sp}\to\{\pm1\}^3$ by
\[
r(S):=\operatorname{diag}\left(r_1(S),r_2(S),r_3(S)\right),\qquad
\gamma_\nu\cdot S=r_\nu(S)\cdot S\cdot\gamma_\nu\quad(\nu=1,2,3).
\]
Then $r$ is a homomorphism onto $\{\pm1\}^3$, and
\begin{equation}
S\cdot(\gamma\cdot v)\cdot S^{-1}=\gamma\cdot\left(r(S)v\right)\qquad\text{for all }v\in\mathbb{C}^3\text{ and }S\in\mathrm{Sp}.
\tag{I.4.1}
\end{equation}
The values of $r$ are the following.

\noindent
\begin{tabular}{p{0.08\textwidth}p{0.21\textwidth}p{0.12\textwidth}p{0.23\textwidth}p{0.18\textwidth}}
$S$ & $\gamma_\mu$ & $\gamma_5$ & $\gamma_\mu\gamma_5$ & $\gamma_1\gamma_2\gamma_3=i\Sigma_3$\\
\hline
$r(S)$ & $+1$ at $\mu$, $-1$ at the two others (rotation by $\pi$ about $e_\mu$) & $-1$ (inversion) & $-1$ at $\mu$, $+1$ at the two others (reflection $e_\mu\to-e_\mu$) & $1$
\end{tabular}

\noindent
One has $\ker r=\{1,-1,i\Sigma_3,-i\Sigma_3\}$.
Let $\mathrm{Sp}_3:=\{S\in\mathrm{Sp}:\ S\cdot\gamma_5=\gamma_5\cdot S\}$;
it has $16$ elements, and it is the group $\mathrm{Sp}_3$ of \cite[Section~5.2.4]{QEDI}, the index $3$ referring to the matrix $\gamma_3$ of \cite[Section~5.2.4]{QEDI}, which is $\gamma_5$ here.
One has $r(\mathrm{Sp}_3)=\{\pm1\}^3$, and an element of $\mathrm{Sp}$ outside $\mathrm{Sp}_3$ anticommutes with $\gamma_5$.

\paragraph{Lifts of reflections.}
For a unit vector $\hat{n}\in\mathbb{R}^3$ put $S_{\mathrm{nv}}:=(\gamma\cdot\hat{n})\cdot\gamma_5$.
Then
\[
S_{\mathrm{nv}}\cdot(\gamma\cdot v)\cdot S_{\mathrm{nv}}^{-1}=\gamma\cdot\left(v-2(v\cdot\hat{n})\,\hat{n}\right),\qquad
S_{\mathrm{nv}}^2=-1 .
\]
For $\hat{n}=e_\mu$ it is the element $\gamma_\mu\gamma_5$ of $\mathrm{Sp}$.
The same lift appears in \cite[Proposition~3.16(b)]{QEDII}.
\begin{itemize}
\item[(I.Sch)]
A $4\times4$ matrix that commutes with $\gamma_1$, $\gamma_2$, $\gamma_3$, $\gamma_5$ is a multiple of $1$.
For $S\in\mathrm{Sp}$ one has $\left(S^{-1}\right)^{T}=\pm\mathsf{C}^{-1}\cdot S\cdot\mathsf{C}$ and $\overline{S}=\pm\varkappa^{-1}\cdot S\cdot\varkappa$.
Every $S\in\mathrm{Sp}$ is monomial, so that $|Sv|_c=|v|_c$.
\end{itemize}

\subsubsection*{I.4.2. The groups of the scheme}

Let
\begin{gather*}
\mathbb{G}_k:=\left\{\tilde{g}=(S,a):\ S\in\mathrm{Sp},\ a\in M\mathbb{Z}^3/\ell\mathbb{Z}^3\right\},\qquad
(S,a)\cdot(S',a'):=\left(S\cdot S',\,a+r(S)a'\right);\\
\mathbb{G}_k^{L}:=\left\{(S,a)\in\mathbb{G}_k:\ a\in LM\mathbb{Z}^3/\ell\mathbb{Z}^3\right\}\quad(\text{the step group});\\
\mathbb{G}_{k,3},\ \mathbb{G}^{L}_{k,3}:\ \text{the elements of }\mathbb{G}_k,\ \mathbb{G}_k^{L}\text{ with }S\in\mathrm{Sp}_3;\\
\tilde{g}_z:=\left(S,\,(1-r(S))z\right)\quad\text{for }z\in M\mathbb{Z}^3\quad(\text{it fixes }z).
\end{gather*}
Elements of these groups are written $\tilde{g}$.
This is the letter of \cite[Section~20.1.3]{QEDI}, where $\tilde{g}$ denotes a lattice isometry, that is, a map of the points;
an element of $\mathbb{G}_0^{L}$ is written $\tilde{g}$ in \cite[Section~20.1.5]{QEDI}, in the proof of \cite[Lemma~20.4(b)]{QEDI}.
In \cite[Sections~5 and 13--19]{QEDI} the elements of these groups are written $g$, a letter that here denotes the label.

\subsubsection*{I.4.3. Actions of $\tilde{g}=(S,a)$; $r:=r(S)$ ($r^{-1}=r$)}

The element $\tilde{g}$ acts as follows.
\begin{description}
\item[Photon sites, cells, sets, polymers:]
$\tilde{g}c:=rc+a$; $\tilde{g}X$, $\tilde{g}\square$, \dots\ are the images.
\item[Fermion sites:]
\begin{equation}
\tilde{g}(c,g):=\left(rc+a,\,rg\right),
\tag{I.4.2}
\end{equation}
that is, the label is multiplied from the left by $r\in\{\pm1\}^3$, a subgroup of $O_{\mathrm{cub}}$.
\item[Bond functions:]
on each of the lattices $L^{-i}\mathbb{Z}^3/\ell\mathbb{Z}^3$, $-1\le i\le k$,
$(\tilde{g}A)(c,c'):=A\left(\tilde{g}^{-1}c,\tilde{g}^{-1}c'\right)$.
\item[Scalars; plaquette functions:]
$(\tilde{g}\boldsymbol{\lambda})(c):=\boldsymbol{\lambda}(\tilde{g}^{-1}c)$; $(\tilde{g}F)(p):=F(\tilde{g}^{-1}p)$, the plaquette functions being odd under reversal of the orientation.
\item[Spinors:]
$(\tilde{g}f)(x):=S\cdot f(\tilde{g}^{-1}x)$, that is, $(\tilde{g}f)(c,g)=S\cdot f\left(r(c-a),rg\right)$ on $\mathfrak{S}^{\mathrm{ch}}$;
$(\tilde{g}v)(y)=S\cdot v(\tilde{g}^{-1}y)$ on $\mathfrak{S}$.
\item[Functions of pairs:]
$(\tilde{g}F)(x,x'):=S\cdot F\left(\tilde{g}^{-1}x,\tilde{g}^{-1}x'\right)$.
\item[Operators, kernels:]
$O\to\tilde{g}\cdot O\cdot\tilde{g}^{-1}$; its kernel is $S\cdot O\left(\tilde{g}^{-1}t,\tilde{g}^{-1}t'\right)\cdot S^{-1}$ (the measure (I.1.3) is invariant).
\item[Parameters:]
$(\tilde{g}s)_{\mathrm{pc}}:=s_{\tilde{g}^{-1}\mathrm{pc}}$ (for the family $\pi''_k$: $\tilde{g}\in\mathbb{G}_k^{L}$).
\item[Generators:]
by the automorphism $a_{\tilde{g}}$ of (I.T-g) in Section~I.4.5.
\end{description}
On antiperiodic functions $a$ is taken in $M\mathbb{Z}^3$;
the translations by $\ell\mathbb{Z}^3$ and the element $-1$ of $\mathrm{Sp}$ act by signs, hence trivially on operators (by conjugation) and on even elements.
On elements that are not even, $a_{\tilde{g}}$ depends on the vector $a$ of $M\mathbb{Z}^3$ that represents the class:
if $a$ is replaced by $a+\ell n$, $n\in\mathbb{Z}^3$, then $a_{\tilde{g}}$ is composed with the automorphism that multiplies every generator by $(-1)^{n_1+n_2+n_3}$ (Proposition~16.1\,(b)).
Every relation that contains $a_{\tilde{g}}$ on elements that need not be even is asserted for every vector $a\in M\mathbb{Z}^3$ that represents the class.
Among the relations of this kind are
Proposition~16.1\,(d4);
the first relation of the display (3.4) in Proposition~17.1\,(e);
the clause (I.k3) of Section~I.5.4, in Proposition~16.3\,(e);
and the clause (k3)$^{\mathrm{o}}$ of Section~I.5.11 and of the Definition in Proposition~48.1, in Proposition~48.2\,(e).
Where $\tilde{g}'=(S,a/L)$ stands beside $\tilde{g}=(S,a)$, the vector of $\tilde{g}'$ is $a/L$.
\begin{itemize}
\item[(I.a1)]
The map $x\mapsto\tilde{g}x$ is a bijection of $\mathbb{L}^{\mathrm{f}}_k$ and of $T^{\mathrm{f}}_k$ with $\pi(\tilde{g}x)=\tilde{g}\pi(x)$;
it preserves $\int\mathrm{d}x$, $d_X$ and $\mathrm{Pz}_k$, and it maps the pairs of kind (I), (II) for $Z$ to those for $\tilde{g}Z$.
\item[(I.a2)]
$(\tilde{g}x)^{(i)}=\tilde{g}\left(x^{(i)}\right)$ for $i=0,\dots,3$ (left and right multiplication commute; $(rg)e_1=r(ge_1)$), and $\hat{d}_i(\tilde{g}x)=r\,\hat{d}_i(x)$.
\item[(I.m4)]
$S\cdot M_{i}(x)\cdot S^{-1}=M_{i}(\tilde{g}x)$, by (I.4.1) and (I.a2).
\item[(I.a3)]
$J\cdot\tilde{g}=\tilde{g}\cdot J$ and $J^{*}\cdot\tilde{g}=\tilde{g}\cdot J^{*}$, since $\tilde{g}$ permutes the labels of a cluster and maps clusters to clusters.
\item[(I.a4)]
For $\tilde{g}=(S,a)\in\mathbb{G}_k^{L}$ put $\tilde{g}':=(S,a/L)\in\mathbb{G}_{k+1}$ (side $\ell/L$).
Then
\[
(\tilde{g}x)/L=\tilde{g}'(x/L),\qquad
\left(\tilde{g}'A'\right)_L=\tilde{g}\left(A'_L\right),\qquad
\left(\tilde{g}'f'\right)_L=\tilde{g}\left(f'_L\right).
\]
\end{itemize}

\subsubsection*{I.4.4. Covariance of the elementary operators}

Let $\tilde{g}=(S,a)$ and $r=r(S)$; the matrix $r$ is diagonal.
The facts (I.a1)--(I.a3) and (I.m4), together with the invariance of paths and blocks under $c\mapsto rc+a$, give
\begin{gather*}
D_{\tilde{g}A}\cdot\tilde{g}=\tilde{g}\cdot D_A,\\
Q_k(\tilde{g}A)\cdot J^{*}\cdot\tilde{g}=\tilde{g}\cdot Q_k(A)\cdot J^{*},\qquad
J\cdot Q_k^{T}(-\tilde{g}A)\cdot\tilde{g}=\tilde{g}\cdot J\cdot Q_k^{T}(-A),
\end{gather*}
the same relations for $Q(A)\cdot J^{*}$ and $J\cdot Q^{T}(-A)$ at $k=0$ and for $Q(A)$, $Q^{T}(-A)$, and
\[
\delta_{\alpha,\tilde{g}A}\left(\tilde{g}f\right)=\tilde{g}\left(\delta_{\alpha,A}f\right).
\]
Moreover
\[
\operatorname{Pen}_{y'}(\tilde{g}A)\cdot\tilde{g}=\tilde{g}\cdot\operatorname{Pen}_{y}(A),\qquad
F_{y'}(\tilde{g}A)\cdot\tilde{g}=\tilde{g}\cdot F_{y}(A),
\]
where $y'=y$ for $S\in\mathrm{Sp}_3$ and $y'=-y$ otherwise.
The covariance relations of the composed operators are the content of Proposition~10.
The group under which the activities are required to be invariant is fixed by (I.k3) in Section~I.5.4; for the primitives it is fixed in Section~I.5.8.
Coordinate permutations, and translations by vectors outside $M\mathbb{Z}^3$, are symmetries of $D_A$ as well; they belong to the group $W_{\mathrm{af}}$ of Section~I.4.7.

\subsubsection*{I.4.5. The four transformations of the generators; charge conjugation; conjugation}

The four maps below are defined on $\mathrm{Gr}[\psi,\chi]$.
On the families $\Psi$ and $W$ they are defined by the same rules: for $k\ge1$ with the unit sites $y$ in place of $x$, and for $k=0$ exactly as on $\psi$.
Spinors are written as columns.
\begin{itemize}
\item[(I.T-g)]
$a_{\tilde{g}}$ is the automorphism with
\begin{align*}
a_{\tilde{g}}\psi(x)&:=S\cdot\psi(\tilde{g}^{-1}x),&
a_{\tilde{g}}\bar{\psi}(x)&:=(S^{-1})^{T}\cdot\bar{\psi}(\tilde{g}^{-1}x),\\
a_{\tilde{g}}\chi(x,x')&:=S\cdot\chi(\tilde{g}^{-1}x,\tilde{g}^{-1}x'),&
a_{\tilde{g}}\bar{\chi}(x,x')&:=(S^{-1})^{T}\cdot\bar{\chi}(\tilde{g}^{-1}x,\tilde{g}^{-1}x').
\end{align*}
\item[(I.T-G)]
For a real function $\Lambda$ on the sites of $T_k$, read on the covering lattice $\mathbb{L}_k$ as an $\ell\mathbb{Z}^3$-periodic function of $c$, $G_{\Lambda}$ is the automorphism with
\begin{align*}
\psi(x)&\mapsto\exp\left(ie_k\Lambda(c)\right)\cdot\psi(x),&
\bar{\psi}(x)&\mapsto\exp\left(-ie_k\Lambda(c)\right)\cdot\bar{\psi}(x),\\
\chi(x,x')&\mapsto\exp\left(ie_k\Lambda(c)\right)\cdot\chi(x,x'),&
\bar{\chi}(x,x')&\mapsto\exp\left(-ie_k\Lambda(c)\right)\cdot\bar{\chi}(x,x'),
\end{align*}
where $x=(c,g)$; the phase is that of the first point.
All $48$ sites of a cluster receive the same phase, as in the gauge transformation of \cite[Definition~2.10]{QEDII}.
\item[(I.T-C)]
$\mathrm{Cg}$ is the automorphism with
\[
\psi\mapsto\mathsf{C}\cdot\bar{\psi},\qquad
\bar{\psi}\mapsto-\mathsf{C}\cdot\psi,\qquad
\chi\mapsto\mathsf{C}\cdot\bar{\chi},\qquad
\bar{\chi}\mapsto-\mathsf{C}\cdot\chi;
\]
it goes together with the substitution $A\mapsto-A$.
In \cite[Section~13.1.5]{QEDI} this automorphism is denoted $\mathbf{C}$.
\item[(I.T-K)]
$\mathbf{K}$ is the antilinear ring automorphism that replaces every coefficient by its complex conjugate and $\psi$, $\bar{\psi}$, $\chi$, $\bar{\chi}$ by $\varkappa\cdot\psi$, $\varkappa\cdot\bar{\psi}$, $\varkappa\cdot\chi$, $\varkappa\cdot\bar{\chi}$;
it goes together with the substitution $A\mapsto-\overline{A}$.
\end{itemize}
The maps (I.T-G), (I.T-C) and (I.T-K) move neither sites nor labels.
The matrices $S$, $(S^{-1})^{T}$, $\mathsf{C}$ and $\varkappa$ are monomial.
Consequently each of the four maps sends every generator to a generator multiplied by a number of modulus one, bijectively up to these factors, and it preserves the weights $\nu$, the bidegree and all norms of Section~I.5.1.
For the four families $\psi$, $\chi$, $\Psi$, $W$ this is proved in Proposition~16.1\,(b), together with the compatibility with (I.AP).

\paragraph{The matrices $\mathsf{C}$ and $\mathrm{Cc}$.}
By (I.O5),
\[
D_A^{T}=\mathsf{C}^{-1}\cdot D_{-A}\cdot\mathsf{C}=\mathrm{Cc}^{-1}\cdot D_{-A}\cdot\mathrm{Cc};
\]
both identities hold because $\mathsf{C}=i\,\Sigma_3\cdot\mathrm{Cc}$ and $\Sigma_3$ commutes with $D_A$.
The two matrices behave differently on $\gamma_5$.
For the operators that contain the resolvent parameter one has
\[
F_y(A)^{T}=\mathsf{C}^{-1}\cdot F_{-y}(-A)\cdot\mathsf{C}
\qquad\text{and}\qquad
F_y(A)^{T}=\mathrm{Cc}^{-1}\cdot F_{y}(-A)\cdot\mathrm{Cc}.
\]
Transposition relations are written in the first form, which is the form of \cite[Proposition~5.47(c)]{QEDI}.
The relation $\mathrm{Cc}\cdot(D_A)^{T}\cdot\mathrm{Cc}^{-1}=D_{-A}$ of \cite[Proposition~3.16(c)]{QEDII}, where $\mathrm{Cc}$ appears as $\sigma_2\oplus\sigma_2$, is brought into the first form by means of $\mathrm{Cc}=-i\,\Sigma_3\cdot\mathsf{C}$.
For the laws the two forms lead to the same conclusion, namely $\det\left(D_{-A}+\boldsymbol{\mu}\right)=\det\left(D_A+\boldsymbol{\mu}\right)$, that is, invariance under $F\mapsto-F$.

\subsubsection*{I.4.6. The internal symmetry; which bilinears are allowed; the seven couplings}

\paragraph{Internal symmetries of $\langle\bar{\psi},(D_A+\boldsymbol{\mu})\psi\rangle$.}
There are two of them:
(i) the global phase;
(ii) $\psi\mapsto\exp\left(it\,\Sigma_3\right)\cdot\psi$, $\bar{\psi}\mapsto\bar{\psi}\cdot\exp\left(-it\,\Sigma_3\right)$ with real $t$, which is a symmetry by (I.O3).
At $t=\pi/2$ the second one is the element $i\,\Sigma_3=\gamma_1\gamma_2\gamma_3$ of $\mathrm{Sp}$, for which $r=1$; it moves no site.

\paragraph{Why the group $\mathrm{Sp}$, and not only the lifts of the lattice group.}
The matrices that commute with the lifts $S_{\mathrm{nv}}$ (Section~I.4.1) of all nine reflections in $O_{\mathrm{cub}}$, and equally the matrices that commute only with the lifts $\gamma_\mu\gamma_5$ of the three axis reflections, form the two-dimensional space spanned by $1$ and by $\sigma_2\otimes1_2=i\,\gamma_5\cdot\Sigma_3$.
The second matrix anticommutes with $\Sigma_3$; it is excluded only by the internal element $i\,\Sigma_3$.
Once this element is included, (I.Sch) applies: there is one mass matrix, and it is a multiple of $1$.

These facts follow from (I.c1) and (I.c2).
The generators $\gamma_1$, $\gamma_2$, $\gamma_3$, $\gamma_5$ anticommute pairwise and have square $1$.
Hence every product of distinct generators other than $1$ anticommutes with one of the generators, and therefore has trace $0$.
The product of two different products of distinct generators is, up to sign, a product of distinct generators other than $1$;
so $\operatorname{tr}\left(\Gamma^{-1}\cdot\Gamma'\right)=4\,\delta_{\Gamma\Gamma'}$ for any two of the sixteen products $\Gamma$, $\Gamma'$ of distinct generators, and these sixteen products form a basis of the $4\times4$ matrices.
Each of them commutes or anticommutes with each $\gamma_\mu\gamma_5$, and it commutes with $\gamma_\mu\gamma_5$ if and only if it contains either both $\gamma_\mu$ and $\gamma_5$ or neither of them.
Consequently it commutes with all three axis lifts if and only if it contains none or all four of the generators, that is, if it is $1$ or
\[
\gamma_1\gamma_2\gamma_3\gamma_5=i\,\Sigma_3\cdot\gamma_5=-\sigma_2\otimes1_2 .
\]
Conjugation by $\gamma_\mu\gamma_5$ multiplies every element of this basis by $\pm1$; therefore the matrices that commute with the three axis lifts are the linear combinations of $1$ and $\sigma_2\otimes1_2$.
The unit normal of a diagonal reflection is $\hat{n}=\left(e_\mu\pm e_\nu\right)/\sqrt{2}$ with $\mu\neq\nu$, and its lift $S_{\mathrm{nv}}=(\gamma\cdot\hat{n})\cdot\gamma_5=\left(\gamma_\mu\gamma_5\pm\gamma_\nu\gamma_5\right)/\sqrt{2}$ is a linear combination of two axis lifts.
Hence the matrices that commute with all nine lifts form the same space.
Finally, $\sigma_2\otimes1_2=i\,\gamma_5\cdot\Sigma_3$ anticommutes with $\Sigma_3$, since $\Sigma_3$ anticommutes with $\gamma_5$.

\begin{center}
\begin{tabular}{p{4.3cm}cp{7.4cm}}
bilinear $\bar{\psi}\,\Gamma\,\psi$ & number & excluded by\\
\hline
$\Gamma=1$ & $1$ & allowed; the coefficient is real by the reality clause (I.k7) of Section~I.5.4\\
$\Sigma_3$ (parity-odd mass) & $1$ & $S=\gamma_\mu\gamma_5$ (a reflection), or $S=\gamma_5$ (the inversion)\\
$\gamma_5$ and $\sigma_2\otimes1_2=i\,\gamma_5\cdot\Sigma_3$ (masses that mix the two blocks) & $2$ & $S=i\,\Sigma_3$ (both). The matrix $\gamma_5$ is odd also under the lifts $\gamma_\mu\gamma_5$ of the reflections. The matrix $\sigma_2\otimes1_2$ commutes with the nine lifts $S_{\mathrm{nv}}$ of Section~I.4.1 and with the group generated by them; it is excluded only by elements of $\mathrm{Sp}$ outside this group ($i\,\Sigma_3$, $\gamma_\mu$, $\gamma_5$).\\
$\gamma_\mu$ & $3$ & $S=\gamma_\nu$, $\nu\neq\mu$\\
$1_2\otimes\sigma_\mu=-i\,\gamma_\nu\gamma_{\mathfrak{l}}$, $(\mu,\nu,\mathfrak{l})$ cyclic & $3$ & $S=\gamma_\nu$\\
$\gamma_\mu\gamma_5$, $\gamma_\mu\gamma_5\cdot\Sigma_3$ & $6$ & $S=i\,\Sigma_3$
\end{tabular}
\end{center}

The sixteen products of the generators $\gamma_1,\gamma_2,\gamma_3,\gamma_5$ transform under conjugation by $\mathrm{Sp}$ according to sixteen different characters.
The table applies to the sum of $\mathrm{al}_{2}(E,X)$ over the small polymers $X$ that contain a fixed $M$-cube (Section~I.5.6; Proposition~16.3\,(a)).
For a single polymer one has only $\mathrm{al}_{2}(E,\tilde{g}X)=S\cdot\mathrm{al}_{2}(E,X)\cdot S^{-1}$, and $\mathrm{al}_{2}(E,X)$ need not be a multiple of $1$; see the Remark after Proposition~16.2, and the end of the proof of Proposition~16.2 in Chapter~16, Section~2.

\paragraph{The seven light--heavy couplings.}
An operator $\mathrm{Ocl}$ that acts inside the clusters and has all symmetries of the lattice model (the lattice group, $\Sigma_3$, charge conjugation, $\gamma_5$-Hermiticity) acts on cluster constants by
\[
(\mathrm{Ocl}\,J\xi)(c,g)=\left[\mathrm{co}_0+i\det(g)\,\mathrm{co}_1\,\Sigma_3+\gamma\cdot(g\,\mathrm{cu})+i\det(g)\,(1_2\otimes\sigma)\cdot(g\,\mathrm{cv})\right]\xi(c)
\]
with real numbers $\mathrm{co}_0$, $\mathrm{co}_1$ and vectors $\mathrm{cu},\mathrm{cv}\in\mathbb{R}^3$.
In particular $J^{*}\cdot\mathrm{Ocl}\cdot J=\mathrm{co}_0$ is a single real number.
Here $\mathrm{Ocl}$ is an operator on $\mathfrak{S}^{\mathrm{ch}}(\eta_k,\ell)$ of the form
\[
(\mathrm{Ocl}\,f)(c,g)=\frac{1}{48}\sum_{g'}\mathrm{Ng}(c;g,g')\cdot f(c,g')
\]
with $4\times4$ matrices $\mathrm{Ng}(c;g,g')$, and the four symmetries are the following.
\begin{itemize}
\item[(I.oc1)]
$\mathrm{Ocl}$ commutes with the translations by $\eta_k\mathbb{Z}^3$ and with the maps $\mathrm{Us}_r$, $r\in O_{\mathrm{cub}}$,
\[
(\mathrm{Us}_rf)(c,g):=S_r\cdot f(r^{-1}c,r^{-1}g),
\]
where $S_r$ is a product of lifts $S_{\mathrm{nv}}$ (Section~I.4.1) along a word of reflections for $r$; these are the matrix $S_{\mathrm{gw}}$ and the map $U_{\mathrm{gw}}$ of Section~I.4.7 for $g_{\mathrm{w}}=(0,r)$.
\item[(I.oc2)]
$\mathrm{Ocl}\cdot\Sigma_3=\Sigma_3\cdot\mathrm{Ocl}$.
\item[(I.oc3)]
$\mathrm{Ocl}^{T}=\mathsf{C}^{-1}\cdot\mathrm{Ocl}\cdot\mathsf{C}$.
\item[(I.oc4)]
$\mathrm{Ocl}^{*}=\gamma_5\cdot\mathrm{Ocl}\cdot\gamma_5$.
\end{itemize}

\begin{proof}
By (I.oc1) the matrix $\mathrm{Ng}$ does not depend on $c$.
Comparing the kernels of $\mathrm{Us}_r\cdot\mathrm{Ocl}$ and of $\mathrm{Ocl}\cdot\mathrm{Us}_r$ one finds
\[
S_r\cdot\mathrm{Ng}(r^{-1}g,r^{-1}g')=\mathrm{Ng}(g,g')\cdot S_r,
\qquad\text{that is,}\qquad
S_r\cdot\mathrm{Ng}(g,g')\cdot S_r^{-1}=\mathrm{Ng}(rg,rg').
\]
Moreover $(\mathrm{Ocl}\,J\xi)(c,g)=\mathrm{Nl}(g)\,\xi(c)$ with the row sum
\[
\mathrm{Nl}(g):=\frac{1}{48}\sum_{g'}\mathrm{Ng}(g,g');
\]
hence $S_r\cdot\mathrm{Nl}(g)\cdot S_r^{-1}=\mathrm{Nl}(rg)$, that is, $\mathrm{Nl}(g)=S_g\cdot\mathrm{Nl}(1)\cdot S_g^{-1}$.
By (I.oc2) the matrix $\mathrm{Nl}(1)$ commutes with $\Sigma_3=\sigma_3\otimes1_2$, so it is block diagonal.
The eight matrices $1$, $\Sigma_3$, $\gamma_\mu=\sigma_3\otimes\sigma_\mu$ and $1_2\otimes\sigma_\mu=\Sigma_3\cdot\gamma_\mu$ form a basis of the block diagonal matrices.
It follows that
\[
\mathrm{Nl}(1)=\mathrm{co}_0+i\,\mathrm{co}_1\,\Sigma_3+\gamma\cdot\mathrm{cu}+i\,(1_2\otimes\sigma)\cdot\mathrm{cv}
\]
with unique complex numbers $\mathrm{co}_0$, $\mathrm{co}_1$ and unique vectors $\mathrm{cu},\mathrm{cv}\in\mathbb{C}^3$.

In terms of kernels, (I.oc3) and (I.oc4) read
\[
\mathrm{Ng}(g',g)^{T}=\mathsf{C}^{-1}\cdot\mathrm{Ng}(g,g')\cdot\mathsf{C},\qquad
\mathrm{Ng}(g',g)^{*}=\gamma_5\cdot\mathrm{Ng}(g,g')\cdot\gamma_5;
\]
each of the two relations exchanges rows and columns.
The complex conjugate of the first relation ($\mathsf{C}$ is real), combined with the second, gives $\mathsf{C}^{-1}\cdot\overline{\mathrm{Ng}(g,g')}\cdot\mathsf{C}=\gamma_5\cdot\mathrm{Ng}(g,g')\cdot\gamma_5$, that is,
\[
\overline{\mathrm{Ng}(g,g')}=\varkappa^{-1}\cdot\mathrm{Ng}(g,g')\cdot\varkappa
\qquad\left(\varkappa=\gamma_5\cdot\mathsf{C},\ \mathsf{C}^{-1}=-\mathsf{C}\right);
\]
hence $\overline{\mathrm{Nl}(1)}=\varkappa^{-1}\cdot\mathrm{Nl}(1)\cdot\varkappa$.
By (I.c4), $\overline{\gamma_\mu}=\varkappa^{-1}\cdot\gamma_\mu\cdot\varkappa$.
The matrix $\Sigma_3$ is real, commutes with $\mathsf{C}=i\gamma_2$ and anticommutes with $\gamma_5$ (I.c2); therefore $\overline{\Sigma_3}=-\varkappa^{-1}\cdot\Sigma_3\cdot\varkappa$ and $\overline{1_2\otimes\sigma_\mu}=-\varkappa^{-1}\cdot\left(1_2\otimes\sigma_\mu\right)\cdot\varkappa$.
Consequently
\[
\overline{\mathrm{Nl}(1)}=\varkappa^{-1}\cdot\left[\overline{\mathrm{co}_0}+i\,\overline{\mathrm{co}_1}\,\Sigma_3+\gamma\cdot\overline{\mathrm{cu}}+i\,(1_2\otimes\sigma)\cdot\overline{\mathrm{cv}}\right]\cdot\varkappa,
\]
and by the uniqueness of the coefficients the numbers $\mathrm{co}_0$, $\mathrm{co}_1$ and the vectors $\mathrm{cu}$, $\mathrm{cv}$ are real.

By Section~I.4.1, $S_{\mathrm{nv}}\cdot(\gamma\cdot v)\cdot S_{\mathrm{nv}}^{-1}=\gamma\cdot\left(v-2(v\cdot\hat{n})\,\hat{n}\right)$.
Moreover $S_{\mathrm{nv}}\cdot\Sigma_3\cdot S_{\mathrm{nv}}^{-1}=-\Sigma_3$, because $\Sigma_3$ commutes with $\gamma\cdot\hat{n}$ and anticommutes with $\gamma_5$.
Along a word for $g$ one obtains
\[
S_g\cdot(\gamma\cdot v)\cdot S_g^{-1}=\gamma\cdot(gv),\qquad
S_g\cdot\Sigma_3\cdot S_g^{-1}=\det(g)\,\Sigma_3,\qquad
S_g\cdot\left((1_2\otimes\sigma)\cdot v\right)\cdot S_g^{-1}=\det(g)\,(1_2\otimes\sigma)\cdot(gv).
\]
Inserted into $\mathrm{Nl}(g)=S_g\cdot\mathrm{Nl}(1)\cdot S_g^{-1}$, these rules give the display.
Finally,
\[
J^{*}\cdot\mathrm{Ocl}\cdot J=\frac{1}{48}\sum_g\mathrm{Nl}(g)=\mathrm{co}_0 .
\]
Indeed, $\sum_g\det(g)=0$ (substitute $g\mapsto gs_0$), $\sum_gg=0$ (substitute $g\mapsto-g$), and the matrix $\sum_g\det(g)\,g$ changes its sign under multiplication from the right by $s_0$ and by $s_3$ (substitute $g\mapsto gs_0$, resp.\ $g\mapsto gs_3$).
So this matrix is annihilated by $1+s_0=\operatorname{diag}(2,2,0)$ and by $1+s_3=\operatorname{diag}(0,2,2)$, and it vanishes.
\end{proof}

The remaining seven numbers, $\mathrm{co}_1$ and the components of $\mathrm{cu}$ and $\mathrm{cv}$, couple the light part to the heavy part at zero momentum.
The symmetry does not exclude them.
For all real $\mathrm{co}_0$, $\mathrm{co}_1$ and all $\mathrm{cu},\mathrm{cv}\in\mathbb{R}^3$ the function
\[
\mathrm{Ng}=\mathrm{Ng}_0+\mathrm{Ng}_1+\mathrm{Ng}_u+\mathrm{Ng}_v
\]
of Proposition~18.6\,(f) has the row sum of the display, and the operator $\mathrm{Ocl}$ formed with it has (I.oc1)--(I.oc4), not only the properties (n1)--(n3) of Proposition~18.6.

\begin{proof}
In terms of kernels, (I.oc2) is (n1) at $S=i\,\Sigma_3$, and (I.oc3) is (n2), since $\mathsf{C}^{-1}=-\mathsf{C}$.
The relation (I.oc4) follows from (n2) and (n3), since $\mathsf{C}$ is real and $\varkappa^{-1}=\varkappa$:
\[
\mathrm{Ng}(g',g)^{*}=\mathsf{C}^{-1}\cdot\overline{\mathrm{Ng}(g,g')}\cdot\mathsf{C}=\left(\varkappa\cdot\mathsf{C}\right)^{-1}\cdot\mathrm{Ng}(g,g')\cdot\left(\varkappa\cdot\mathsf{C}\right),
\qquad
\varkappa\cdot\mathsf{C}=\gamma_5\cdot\mathsf{C}^{2}=-\gamma_5 .
\]
As to (I.oc1): the four functions do not depend on $c$, and
\[
S_r\cdot\mathrm{Ng}(g,g')\cdot S_r^{-1}=\mathrm{Ng}(rg,rg')
\]
holds for every $r\in O_{\mathrm{cub}}$ by the three rules for $S_g$ above, because $\mathbf{1}[\,rg'=rgs_i\,]=\mathbf{1}[\,g'=gs_i\,]$ and $\det(rg)=\det(r)\det(g)$.
\end{proof}

Consider the bare operator in the form (I.2.5), at $A=0$ and at the momentum $p=0$ of (I.2.6), that is,
\[
(\mathrm{Ocl}\,f)(c,g):=M_{\mathrm{ch}}\,f(c,g)-\sum_{i=0}^{3}M_{i}(x)\cdot f(c,gs_i).
\]
For arbitrary weights $w_i$ the formula (I.2.1) gives
\[
\mathrm{Nl}(g)=M_{\mathrm{ch}}-\frac{\rho}{2}\sum_iw_i-\frac{1}{2}\,\gamma\cdot\left(g\sum_iw_i\hat{n}_i\right);
\]
thus $\mathrm{co}_1=0$, $\mathrm{cv}=0$ and $\mathrm{cu}=-\frac{1}{2}\sum_iw_i\hat{n}_i$, and $\mathrm{cu}=0$ is the closure relation (I.1.1).
Three consequences for the scheme follow.
\begin{itemize}
\item[(I.E1)]
The extracted two-point kernel is integrated over both cluster labels (Section~I.5.6): it is a $4\times4$ matrix, hence it yields one real mass, uniform on the $48$ sites.
\item[(I.E2)]
Differences inside a cluster are H\"older legs with $d_X=\eta_k$, and they are never extracted.
\item[(I.E3)]
At every level the operator in the propagators is formed with the bare $\rho$ and the bare weights $w_i$; this is the rule (I.BARE) of Section~I.2.3.
\end{itemize}

\subsubsection*{I.4.7. The full lattice group}

Lattice units are used in this section.
Let
\[
W_{\mathrm{af}}:=\mathbb{Z}^3\rtimes O_{\mathrm{cub}};
\]
this is the group $W$ of \cite[Definition~2.2]{QEDII}, which is generated by the reflections in the nine families of mirrors.
Its elements are written $g_{\mathrm{w}}=(a,r)\colon y\mapsto ry+a$.
Let $\mathrm{Gc}:=\mathbb{Z}^3\rtimes\{\pm1\}^3$.
An element $g_{\mathrm{w}}$ is called even or odd according to the sign of $\det r$.
On fermion sites $g_{\mathrm{w}}(c,g):=(rc+a,rg)$ (Proposition~46\,(a)), and
\[
(g_{\mathrm{w}}x)^{(i)}=g_{\mathrm{w}}\left(x^{(i)}\right),\qquad
\hat{d}_i(g_{\mathrm{w}}x)=r\,\hat{d}_i(x).
\]

\paragraph{Lifts.}
$S_{\mathrm{gw}}$ is a product of matrices $S_{\mathrm{nv}}$ (Section~I.4.1) along a word of reflections for $r$; it is determined up to a sign.
Put $(U_{\mathrm{gw}}f)(x):=S_{\mathrm{gw}}\cdot f(g_{\mathrm{w}}^{-1}x)$.
Then, by \cite[Proposition~3.16(b)]{QEDII},
\[
U_{\mathrm{gw}}\cdot D_A\cdot U_{\mathrm{gw}}^{-1}=D_{(g_{\mathrm{w}}^{-1})^{*}A}.
\]

\paragraph{Pull-backs; orientation.}
On cochains put
\[
(g_{\mathrm{w}}^{*}u)(\boldsymbol{\sigma}):=\epsilon(g_{\mathrm{w}},\boldsymbol{\sigma})\,u(g_{\mathrm{w}}\boldsymbol{\sigma}),
\]
where $\epsilon(g_{\mathrm{w}},\boldsymbol{\sigma})=\pm1$ compares orientations; for a plaquette $p=(c;\mu\nu)$ it is given by $re_\mu\wedge re_\nu=\epsilon(g_{\mathrm{w}},p)\,e_{\mu'}\wedge e_{\nu'}$, where $\mu'<\nu'$ are the directions of $g_{\mathrm{w}}p$.
One has $(g_{\mathrm{w}}g_{\mathrm{w}}')^{*}=(g_{\mathrm{w}}')^{*}g_{\mathrm{w}}^{*}$; put $(g_{\mathrm{w}})_{*}:=(g_{\mathrm{w}}^{-1})^{*}$.
The maps $g_{\mathrm{w}}^{*}$ and $(g_{\mathrm{w}})_{*}$ commute with $d$, and they are orthogonal for the counting pairings, on the finitely supported cochains of $\mathbb{Z}^3$ as well as between the cochains of a torus and those of its image (\cite[Definition~2.1]{QEDII}; \cite[Lemma~7.6]{QEDI}; the data of Proposition~50).
On one-cochains, $(g_{\mathrm{w}}^{*}A)(c,c')=A(g_{\mathrm{w}}c,g_{\mathrm{w}}c')$.
Hence the action $\tilde{g}A$ of Section~I.4.3 is $\tilde{g}_{*}A$, the element $\tilde{g}$ acting on the photon sites as an element of $W_{\mathrm{af}}$; both companion papers use this convention (\cite[Sections~5.2.9 and~21.1.5]{QEDI}; \cite[Definitions~2.1 and~12.2]{QEDII}).
The signs of the nine reflections on two-forms form the table (ST), which is the list in \cite[Lemma~12.3(b)]{QEDII}.
For the axis families the sign is $-1$ exactly when the plaquette contains the direction of the normal (\cite[Section~7.1.4]{QEDI}; \cite[Definition~2.26]{QEDII}).

\paragraph{Relation to $\mathbb{G}_0$.}
For $\tilde{g}=(S,a)$ the map of the sites is that of $g_{\mathrm{w}}=(a,r(S))\in\mathrm{Gc}$, and $S$ is one of the four matrices $\pm S_{\mathrm{gw}}$, $\pm i\,\Sigma_3\cdot S_{\mathrm{gw}}$.
Hence $a_{\tilde{g}}$ is the lift of a lattice symmetry, possibly composed with the internal element $i\,\Sigma_3$.

\subsubsection*{I.4.8. Mirrors and reflections}

The conventions are those of \cite[Definitions~2.2 and~2.9]{QEDII}.

\paragraph{Mirrors.}
$\mathcal{L}$ is the set of the nine linear forms on $\mathbb{R}^3$
\[
\ell=y_\mu\quad\text{(axis families)},\qquad
\ell=y_\mu-y_\nu,\ y_\mu+y_\nu,\ \mu<\nu\quad\text{(diagonal families)};
\]
in \cite[Definition~2.2]{QEDII} the coordinate forms are written $x_\mu$.
One writes $\ell(y)=n_\ell\cdot y$, and $|n_\ell|^{2}$ denotes the squared Euclidean length $|n_\ell|_2^2$.
A mirror is a pair $(\ell,\mathrm{po})$ with $\mathrm{po}\in\operatorname{Adm}(\ell)$, where $\operatorname{Adm}(\ell):=\frac{1}{2}\mathbb{Z}$ for an axis family ($\mathrm{po}\in\mathbb{Z}$: site plane; $\mathrm{po}\in\mathbb{Z}+\frac{1}{2}$: link plane) and $\operatorname{Adm}(\ell):=\mathbb{Z}$ for a diagonal family.
The reflection in the mirror is
\[
\theta_{\ell,\mathrm{po}}(y):=y-\frac{2}{|n_\ell|^{2}}\left(\ell(y)-\mathrm{po}\right)n_\ell .
\]
In \cite[Definition~2.2]{QEDII} the position $\mathrm{po}$ is denoted $c$; here $c$ denotes the photon site.

\paragraph{In group coordinates.}
The reflection $\theta_{\ell,\mathrm{po}}$ is an element of $W_{\mathrm{af}}$ and acts as in Section~I.4.7.
For an axis mirror
\[
\theta_{y_\mu,\mathrm{po}}(c,g)=\left(c-2(c_\mu-\mathrm{po})\,e_\mu,\ r_\mu g\right),\qquad r_\mu:=1-2e_\mu e_\mu^{T}.
\]
No fermion site is fixed: no site lies on a mirror.
A site $x=(c,g)$ lies above $(y_\mu,\mathrm{po})$ if $c_\mu>\mathrm{po}$, or if $c_\mu=\mathrm{po}$ and the non-zero entry of the row $\mu$ of $g$ is $+1$; this is the sign of $\mathrm{geo}(x)_\mu-\mathrm{po}$, whatever $v_0$ is.

\paragraph{Bonds that cross a mirror.}
By \cite[Proposition~3.6(b), (c)]{QEDII}, a bond that joins the two sides of a mirror is of the form $\{x,\theta x\}$.
A site plane is crossed by bonds of type $0$ only: eight over each photon site of the plane, without phase.
A link plane is crossed by bonds of type $3$ only: the eight bonds over each photon link that it cuts, with the phase of this link.
A diagonal mirror is crossed by bonds of the types $1$ and $2$.
Hence every reflection is of link type for the fermion.

\paragraph{The reflections $\Theta_m$.}
Let $m=(\ell,\mathrm{po})$ be a mirror, $\theta:=\theta_{\ell,\mathrm{po}}$, and $\hat{n}:=n_\ell/|n_\ell|_2$ its unit normal.
Then $\Theta_m$ is the antilinear map of the Grassmann algebra that reverses products and acts on the generators by
\[
\Theta_m\psi(x)=-\bar{\psi}(\theta x)\,(\gamma\cdot\hat{n}),\qquad
\Theta_m\bar{\psi}(x)=-(\gamma\cdot\hat{n})\,\psi(\theta x),
\]
$\psi$ being a column and $\bar{\psi}$ a row \cite[Definition~2.9]{QEDII}.
The sign rules \cite[Lemma~3.10(c), (d)]{QEDII} and the definitions of the reflection of integrands and of the action of $W$ on even elements in \cite[Section~2.3]{QEDII}, which follow \cite[Definition~2.12]{QEDII}, hold for these maps as stated there.
In \cite[Lemma~3.10(d)]{QEDII}, for two mirrors $m=(\ell,\mathrm{po})$ and $m'=(\ell,\mathrm{po}')$ of one family, $\Theta_m\Theta_{m'}$ is the translation $\psi(x)\mapsto\psi(x+v)$, $\bar{\psi}(x)\mapsto\bar{\psi}(x+v)$ by the vector
\[
v=\frac{2}{|n_\ell|^{2}}\left(\mathrm{po}-\mathrm{po}'\right)n_\ell .
\]
In the coordinates $(c,g)$ these rules are written out in Proposition~46\,(h).
They apply to the maps defined here for the following reason.
The generators $\psi_\beta(x)$, $\bar{\psi}_\beta(x)$ at $x=(c,g)$ are the generators of \cite[Definition~2.8]{QEDII} at the site $\mathrm{geo}(x)$.
By Proposition~46\,(a) the map $\mathrm{geo}$ intertwines the actions of $W_{\mathrm{af}}$: $\mathrm{geo}(\theta x)=\theta\,\mathrm{geo}(x)$ and $\mathrm{geo}(g_{\mathrm{w}}x)=g_{\mathrm{w}}\,\mathrm{geo}(x)$.
So the displays of \cite[Definition~2.9]{QEDII} and their forms in the coordinates $(c,g)$ define the same maps.
On functions of the curvature,
\[
(\Theta_mG)\left(F\right)=\overline{G\left(\theta^{*}F\right)}.
\]
For $g_{\mathrm{w}}\in W_{\mathrm{af}}$ and an even $G$ put
\[
G^{\mathrm{gw}}:=\Theta_{m_1}\cdots\Theta_{m_r}G\qquad\text{for a word }g_{\mathrm{w}}=\theta_{m_1}\cdots\theta_{m_r};
\]
by \cite[Proposition~3.16(f$''$)]{QEDII} this does not depend on the word.
On functions of $F$,
\[
G^{\mathrm{gw}}\left(F\right)=G\left(g_{\mathrm{w}}^{*}F\right)\quad\text{for even }g_{\mathrm{w}},\qquad
G^{\mathrm{gw}}\left(F\right)=\overline{G\left(g_{\mathrm{w}}^{*}F\right)}\quad\text{for odd }g_{\mathrm{w}}.
\]

\paragraph{Crossing coefficients.}
In the lattice units of Section~I.2.5, that is, at $k=0$ with $\eta_0=1$ and with the charge $q$ in the place of $e_k$, the fermion action is
\[
\bar{\psi}\left(D_A+\boldsymbol{\mu}\right)\psi=\sum_x\bar{\psi}(x)\cdot\left(\left(D_A+\boldsymbol{\mu}\right)\psi\right)(x);
\]
by Proposition~46\,(b), (h) it is the action $S_F$ of \cite[Definition~2.10]{QEDII} with the parameters $(\rho,w,M=M_{\mathrm{ch}},q)$.
Let $\mathrm{Xc}_m$ be the sum of those of its terms that belong to the oriented bonds $(x,x')$ and $(x',x)$, where $\{x,x'\}$ ranges over the bonds that cross $m$; the term of the oriented bond $(x,x')$ is the part of $\bar{\psi}(x)\cdot\left(\left(D_A+\boldsymbol{\mu}\right)\psi\right)(x)$ that contains $\psi(x')$.
This is the sum $S^{\times}_{\Pi}$ of \cite[Proposition~3.20]{QEDII} for the mirror $\Pi=m$.
The terms of the oriented bonds contain no mass: by (I.2.5) the part of $D_A+\boldsymbol{\mu}$ that is diagonal in the sites is the number $M_{\mathrm{ch}}$, and the term of an oriented bond $(x,x')$ of type $i$ is $-\bar{\psi}(x)\cdot M_i(x)\cdot U_A(x,x')\cdot\psi(x')$.
If $A$ is real, as in \cite[Definition~2.27 and Proposition~3.20]{QEDII}, then
\[
-\mathrm{Xc}_m=\sum_ic_i\,u_i\,\Theta_m(C_i)\,C_i,
\]
with the phases $u_i$, the elements $C_i$ and the coefficients $c_i$ of \cite[Definition~2.27 and Proposition~3.20]{QEDII}, where they are denoted $u_i$, $\zeta_i$ and $\omega_i$; each $c_i$ is one of the numbers $\frac{w_j}{2}(1+\rho)$, $\frac{w_j}{2}(1-\rho)$, $j$ being the type of the bond.
For complex $A$ the same identity holds if, for every bond $\{x,\theta x\}$ that crosses $m$, $x$ being its lower end as in \cite[Definition~2.27]{QEDII}, the phase $\overline{U_A(x,\theta x)}$ of the modes of the groups~$3$ and~$4$ is replaced by $U_A(\theta x,x)=U_A(x,\theta x)^{-1}$.
All coefficients are $>0$ for $0<\rho<1$.
Positivity at a mirror requires $|\rho|\le1$, and no condition on the mass.

\paragraph{Two further maps denoted by $\Theta$.}
The letter $\mathrm{lw}$ denotes a law on $\mathcal{S}'$, that is, a probability measure on $(\mathcal{S}',\Sigma)$, $\Sigma$ being the cylinder $\sigma$-algebra;
for a function $\Phi$ on $\mathcal{S}'$, $\mathrm{lw}(\Phi)$ is the integral of $\Phi$ with respect to $\mathrm{lw}$, and $\langle\Phi,\Phi'\rangle_{\mathrm{lw}}:=\mathrm{lw}\left(\overline{\Phi}\cdot\Phi'\right)$.
This agrees with the clause (n1) of the notation for laws on $\mathcal{S}'$ in Part~II and with \cite[Definition~2.38]{QEDII}; the law $\mathrm{lw}$ takes the place of the law $\rho$ of \cite[Section~8.1.2]{QEDI} and of the law $\mu$ of \cite[Definition~2.38]{QEDII}.
For $\nu\in\{1,2,3\}$ and real $\mathrm{po}$ let $r_{e_\nu,\mathrm{po}}$ be the reflection of $\mathbb{R}^3$ in the plane $\{y_\nu=\mathrm{po}\}$, and for $j\in\{1,2,3\}$ let $\theta_j:=r_{e_j,0}$ be the reflection $y_j\mapsto-y_j$.
As in the clause (n2) of that notation and in \cite[Section~8.1.2]{QEDI}, a reflection $r$ of $\mathbb{R}^3$ with linear part $r_0$ acts on test forms, on $\mathcal{S}'$ and on functions on $\mathcal{S}'$ by
\[
(r\cdot f)(y):=(\wedge^2r_0)\,f\left(r^{-1}y\right),\qquad
(r\omega)(f):=\omega\left(r^{-1}\cdot f\right),\qquad
(\hat{r}\Phi)(\omega):=\Phi\left(r^{-1}\omega\right).
\]
The reflection of test forms in the plane $\{y_\nu=\mathrm{po}\}$ is $\mathsf{R}_{\nu,\mathrm{po}}f:=r_{e_\nu,\mathrm{po}}\cdot f$, and $\mathsf{R}^{*}_{\nu,\mathrm{po}}$ denotes the map of $\mathcal{S}'$ given by $\left(\mathsf{R}^{*}_{\nu,\mathrm{po}}\omega\right)(f):=\omega\left(\mathsf{R}_{\nu,\mathrm{po}}f\right)$;
these are the maps $\mathsf{R}_{\nu,c}$ and $\mathsf{R}^{*}_{\nu,c}$ of \cite[Section~2.5]{QEDI} with $c=\mathrm{po}$.
Since $r_{e_\nu,\mathrm{po}}$ is an involution, $\mathsf{R}^{*}_{\nu,\mathrm{po}}\omega=r_{e_\nu,\mathrm{po}}\omega$.
In \cite[Section~2.5]{QEDI} and \cite[Theorem~7.15(7)]{QEDI}, $\Theta_{\nu,\mathrm{po}}$ is the antilinear map $A\mapsto\overline{A\circ\mathsf{R}^{*}_{\nu,\mathrm{po}}}$ on functions $A$ on $\mathcal{S}'$.
Hence, for every real $\mathrm{po}$,
\[
\Theta_{\nu,\mathrm{po}}A=\overline{\hat{r}_{e_\nu,\mathrm{po}}A},\qquad
\mathrm{lw}\left(\left(\Theta_{\nu,\mathrm{po}}A\right)\cdot B\right)=\left\langle\hat{r}_{e_\nu,\mathrm{po}}A,B\right\rangle_{\mathrm{lw}},
\]
the second identity whenever $\hat{r}_{e_\nu,\mathrm{po}}A$ and $B$ lie in $L^2(\mathrm{lw})$.
In this paper, as in \cite[Section~8.1.2]{QEDI}, $\Theta_j:=\hat{\theta}_j$ is the linear map $F\mapsto F\circ\theta_j^{-1}$ on functions on $\mathcal{S}'$, and the reflected form is
\[
(F,G)\mapsto\left\langle\Theta_jF,G\right\rangle_{\mathrm{lw}}=\mathrm{lw}\left(\overline{\Theta_jF}\cdot G\right).
\]
The letter $\Theta_j$ has another meaning in \cite[Definition~2.38]{QEDII}, and wherever it occurs in \cite[Theorem~12.21, Propositions~12.32 and~12.33, Theorem~12.39]{QEDII} and in their proofs: there it denotes the complex conjugate of this map, that is, the antilinear map $F\mapsto\overline{F\circ\theta_j^{-1}}$, while the linear map $F\mapsto F\circ\theta_j^{-1}$ is written $\hat{\theta}_j$, as here.
The reflected form of \cite[Definition~2.38]{QEDII}, which is written there with the antilinear map and equals $(F,G)\mapsto\left\langle\hat{\theta}_jF,G\right\rangle_{\mathrm{lw}}$, is therefore the form $(F,G)\mapsto\left\langle\Theta_jF,G\right\rangle_{\mathrm{lw}}$.
For $\nu=j$ and $\mathrm{po}=0$ one has $\Theta_{j,0}F=\overline{\Theta_jF}$ \cite[Section~8.1.2]{QEDI}, so that the map denoted $\Theta_j$ in \cite[Definition~2.38]{QEDII} is $\Theta_{j,0}$.
In this paper the maps $\Theta_{\nu,\mathrm{po}}$ and $\Theta_j$ are different: the first is antilinear, the second linear.
The numbers $\Theta^{(0)}_m$ of Section~I.7.6 are distinguished from the maps $\Theta_m$ by their indices.

\subsection*{I.5. Norms, weighted Grassmann integration, spaces of activities, the step, level~0}

\subsubsection*{I.5.1. Norms on Grassmann algebras}

Let $\mathrm{Ind}_1,\dots,\mathrm{Ind}_p$ be finite index sets, each with a total order, and let $\mathrm{Gr}[\phi_1,\dots,\phi_p]$ be the Grassmann algebra with the independent generators $\phi_i(t)$, $t\in\mathrm{Ind}_i$.
For $I\subset\mathrm{Ind}_i$ let $\phi_i(I)$ be the product of the $\phi_i(t)$, $t\in I$, in increasing order, and write
\[
F=\sum c_F(I_1,\dots,I_p)\,\phi_1(I_1)\cdots\phi_p(I_p).
\]
For a weight vector $\mathsf{w}=(\mathsf{w}_1,\dots,\mathsf{w}_p)$ with $\mathsf{w}_i>0$ put
\begin{equation}
\|F\|_{\mathsf{w}}:=\sum_{I_1,\dots,I_p}\left|c_F(I_1,\dots,I_p)\right|\,\mathsf{w}_1^{\#|I_1|}\cdots\mathsf{w}_p^{\#|I_p|}.
\tag{I.5.1}
\end{equation}
This is the norm of \cite[Section~13.1.4]{QEDI}, where the weight vector is denoted $g$.
The element $F$ \emph{contains} a generator if the generator is a factor of a monomial $\phi_1(I_1)\cdots\phi_p(I_p)$ with $c_F(I_1,\dots,I_p)\neq0$; by (I.AP), a change of the fundamental domain changes signs only.
A function with values in the algebra contains a generator if one of its values does.

Two different notions of weight occur, and they must be kept apart: $\nu(t)$, the weight of an \emph{index} (a measure on the index set, Section~I.3.6), and $\mathsf{w}_i$, the weight of a \emph{family} of generators in a norm.
The norm (I.5.1) does not contain $\nu$; the norm of a single generator $\phi_i(t)$ is $\mathsf{w}_i$.
If $F$ is written with kernels with respect to the measures $\nu_i$,
\[
F=\sum_n\left(\prod_in_i!\right)^{-1}\int\mathrm{d}t\,K_n\,\phi_1\cdots\phi_p,
\]
then
\[
\|F\|_{\mathsf{w}}\le\sum_n\prod_i\frac{\mathsf{w}_i^{n_i}}{n_i!}\int\mathrm{d}t\,\left|K_n\right|,
\]
with equality for antisymmetric kernels, by (I.N0).

\paragraph{The norms of the scheme.}
For kernels $E_{nm}(\xi_1,\dots,\xi_n;\boldsymbol{\zeta}_1,\dots,\boldsymbol{\zeta}_m)$ that are antisymmetric in the $\xi$ and in the $\boldsymbol{\zeta}$ separately put
\[
\|E_{nm}\|:=\int\mathrm{d}\xi\,\mathrm{d}\boldsymbol{\zeta}\,\left|E_{nm}\right|,\qquad
\|E\|_{\mathsf{w}}=\sum_{n,m}\frac{\mathsf{w}_1^{n}\,\mathsf{w}_2^{m}}{n!\,m!}\,\|E_{nm}\|\quad\text{for }\mathsf{w}=(\mathsf{w}_1,\mathsf{w}_2),
\]
\[
h_k=(h_{k,1},h_{k,2}):=\left(e_k^{-\frac{1}{4}},\,e_k^{-\frac{1}{4}+\varepsilon}\right),
\]
and let $|F|_a$ denote the norm (I.5.1) of $F\in\mathrm{Gr}[\Psi]$ with the single weight $a>0$.
The norm $\|E\|_{\mathsf{w}}$ is (I.5.1) for the two families $(\psi,\chi)$.
In \cite[Section~5.3.3]{QEDI} it is written $\|E\|_h$ with $h=(h_1,h_2)$; here the letter $h$ without index denotes the scale of Section~I.0.5.
The integrals contain the weights $\eta_k^3/48$ and $(\eta_k^3/48)^2$ of Section~I.3.6.
Hence all powers of $\eta_k$ and all volumes have the same values as in \cite[Section~5.3.3]{QEDI}; for instance
\[
\left\|\langle\bar{\psi},\psi\rangle_{\square}\right\|_{h_k}=4M^3h_{k,1}^{2}.
\]

\paragraph{Matrix notation.}
$E_{20}(X;\mathscr{A},\mathscr{C};x,x')$ is the $4\times4$ matrix whose entry in row $\beta$ and column $\beta'$ is $E_{20}\left(X;\mathscr{A},\mathscr{C};(x,\beta,1),(x',\beta',0)\right)$.
If every monomial of $E(X;\mathscr{A},\mathscr{C})$ contains as many barred generators as unbarred ones, that is, if $E$ is balanced in the sense of (G2.r8), then, by Proposition~15\,(n) (the display (G2.0.1)),
\[
E^{[2,0]}(X)=\int\mathrm{d}x\,\mathrm{d}x'\,\bar{\psi}(x)^{T}\,E_{20}(X;x,x')\,\psi(x');
\]
$F^{[n,m]}$ is the part of bidegree $(n,m)$ in $(\psi,\chi)$.
For a constant matrix $\mathrm{Mx}$, the kernel of $\int_X\mathrm{d}x\,\bar{\psi}^{T}\mathrm{Mx}\,\psi$ is $48\,\eta_k^{-3}\,\mathbf{1}[\,x=x'\,]\,\mathrm{Mx}$ on $X$, and its integral is $\operatorname{Vol}(X)\,\mathrm{Mx}$.

\subsubsection*{I.5.2. The rules (I.N0)--(I.N6), (I.S6c); integration on an index set with weights}

The rules (I.N0), (I.N1) (products), (I.N2) (linear substitutions), (I.N4) (factorisation), (I.N4$'$) (invariance), (I.N5) (contour integrals), (I.N6) (monotonicity) and (I.S6c)\,(i), (iii), (iv) are, in this order, the clauses (a), (b), (d), (g), (h), (j), (k), (l) of Proposition~15.
In (I.N2) the norm $|O|_{1,\infty}$ is that of (I.ke3); for an operator whose source is a space of functions of pairs, such as $T^{(2)}$ of Section~I.5.9, $|O|_{1,\infty}$ is the supremum over the rows of the sum of the moduli of the coefficients of the values $F(x,x')_\beta$, and this is the same number as for the kernel with respect to $\int\mathrm{d}\boldsymbol{\zeta}$.
The rules are those of \cite[Theorem~13.1]{QEDI}, whose paragraphs are the clauses (N0)--(N7) of \cite[Proposition~5.57]{QEDI}; (I.S6c) is the clause (N7) there.
They are read for arbitrary finite index sets with weights $\nu$ and with the integral (I.3.5), and with the following changes of wording:
$\mathsf{w}$ in place of $g$;
$\mathrm{rad}$ in place of the radius $\rho$;
in (I.N4) the unit sites $y$ are replaced by the points $t$ of the index set of $W$;
in (I.N4$'$), $U$ is an invertible linear map of the functions on the index set, and $(U^{-1})^{T}$ is its inverse transpose with respect to $\langle\cdot,\cdot\rangle_\nu$;
in (I.S6c)\,(iii), $\mathrm{Cov}$ is an invertible linear map, and
\[
\int\mathrm{d}\mu_{\mathrm{Cov}}(W)\,F:=(\det\mathrm{Cov})\int\mathrm{D}W\,\exp\left(-\langle\bar{W},\mathrm{Cov}^{-1}W\rangle_\nu\right)F.
\]

\paragraph{The rule (I.N3)$^{\mathrm{w}}$.}
This is the form in which integration is used; it is proved in Proposition~15.
Let $W$ be a family of generators on an index set with weights $\nu$, and let $\mathrm{d}\mu_I(W)$ be as in (I.3.5).
For $F\in\mathrm{Gr}[\phi_1,\dots,\phi_p,W]$,
\begin{equation}
\left\|\int\mathrm{d}\mu_I(W)\,F\right\|_{(\mathsf{w}_1,\dots,\mathsf{w}_p)}\le\|F\|_{(\mathsf{w}_1,\dots,\mathsf{w}_p,\mathsf{w}_W)},\qquad
\mathsf{w}_W:=\max_t\nu(t)^{-\frac{1}{2}}.
\tag{I.5.2}
\end{equation}
Each integrated generator, $W$ or $\bar{W}$, carries the weight $\nu^{-\frac{1}{2}}$: the form is symmetric in $W$ and $\bar{W}$.
Accordingly (I.S6c)\,(ii) reads
\[
\left|\int\mathrm{d}\mu_I(\Psi)\,F\right|\le|F|_{\mathsf{w}_W}.
\]
For $k\ge1$ one has $\nu=1$ and $\mathsf{w}_W=1$, as in \cite[Theorem~13.1]{QEDI}.
For $k=0$ one has $\nu=1/48$ and $\mathsf{w}_W=\sqrt{48}<7$; the number $7$ is used as a bound for it in the constant $C_{W}$ (I.L0.19), and (I.5.2) holds with $7$ in place of $\mathsf{w}_W$, by (I.N6).

Contour integrals are written
\[
\oint_{\mathrm{rad}}f:=\frac{1}{2\pi i}\oint_{|t|=\mathrm{rad}}\mathrm{d}t\,\frac{f(t)}{t(t-1)},\qquad
\oint^{(2)}_{\mathrm{rad}}f:=\frac{1}{2\pi i}\oint_{|t|=\mathrm{rad}}\mathrm{d}t\,\frac{f(t)}{t^{2}(t-1)},
\]
the circle being oriented counterclockwise.

\subsubsection*{I.5.3. Field domains (sets of bond functions)}

The sets $R_k$, $\mathrm{cf}\,\tilde{R}_k$ ($0<\mathrm{cf}\le1$), $\tilde{R}_k^{-}$, $G_k$, $\tilde{G}_k$, $\mathfrak{X}_k$, $\mathrm{Dom}_k$ and the ``three norms of $A$'' are defined word for word as in \cite[Section~5.3.1 and Definition~5.16]{QEDI}, with $c$ for the fine sites.
In particular
\[
G_k:=\left\{A\text{ complex}:\ |A|_\infty<e_k^{-1+\varepsilon}\right\};
\]
$\tilde{G}_k$ is the set of the bond functions $A_0+A_1$ with $A_1\in G_k$ and a real bond function $A_0$ with the following property:
for every $M$-cube $\square$ there are a real function $\boldsymbol{\lambda}$ on the sites of $\square^{\sim5}$ and a real function $A'$ on the bonds contained in $\square^{\sim5}$ (Section~I.1.1) such that $A_0=A'+d\boldsymbol{\lambda}$ on each bond contained in $\square^{\sim5}$ and $|A'|_\infty<e_k^{-1+\varepsilon}$, where $|A'|_\infty$ denotes the supremum of $|A'|$ over the bonds contained in $\square^{\sim5}$;
the functions $\boldsymbol{\lambda}$ and $A'$ are defined on $\square^{\sim5}$ only, and they may depend on $\square$.
Finally,
\[
\mathfrak{X}_k:=\left\{\mathscr{C}\text{ complex}:\ |\mathscr{C}|_\infty<e_k^{-1+2\varepsilon},\ |\partial\mathscr{C}|_\infty<e_k^{-1+3\varepsilon}\right\}.
\]
The properties DF1--DF4 of the domains are \cite[Lemmas~17.4--17.7]{QEDI}.
The property DF5, which is \cite[Lemma~20.7]{QEDI}, is the clause (DF5) of Proposition~35.4, in the setting (ST*)$^{\mathrm{ch}}$ of Section~G5.1.1.
The domain of all fermionic operators and of all primitives is $\tilde{G}_k$.
On it $\mathfrak{i}(A)<e_k^{\varepsilon}$.
\begin{itemize}
\item[(I.Lc2)$^{\mathrm{ch}}$]
For $B\in\tilde{G}_k$ and every pair $(x,x')$ of $\mathrm{Pz}_k$ one has $\left|U_B(x,x')^{\pm1}\right|\le\mathrm{e}$; for a long pair, $\left|(\delta_{\alpha,B}f)(x,x')\right|_c\le4\,|f|_\infty$.
This rests on the bounds of Section~I.1.8 and on the condition $3Me_0^{\varepsilon}\le1$.
\end{itemize}

\subsubsection*{I.5.4. Polymer activities with and without slot}

Let $k\ge0$ and $\ell\in\mathrm{Cl}$.
The space $K^{x,C}_{k}=K^{x,C}_{k}(\ell)$ is the complex Banach space of the families $\mathscr{E}=\left\{\mathscr{E}(X;\mathscr{A},\mathscr{C},\psi,\chi):X\in D_k\right\}$ of even elements of $\mathrm{Gr}[\psi,\chi]$ of the form
\begin{multline*}
\mathscr{E}(X;\mathscr{A},\mathscr{C},\psi,\chi)=\sum_{n,m}\frac{1}{n!\cdot m!}\int\mathrm{d}\xi\,\mathrm{d}\boldsymbol{\zeta}\;\mathscr{E}_{nm}\left(X;\mathscr{A},\mathscr{C};\xi_1,\dots,\xi_n;\boldsymbol{\zeta}_1,\dots,\boldsymbol{\zeta}_m\right)\\
\cdot\psi(\xi_1)\cdots\psi(\xi_n)\cdot\chi(\boldsymbol{\zeta}_1)\cdots\chi(\boldsymbol{\zeta}_m),
\end{multline*}
with kernels $\mathscr{E}_{nm}$ that are antisymmetric in $\xi_1,\dots,\xi_n$ and, separately, in $\boldsymbol{\zeta}_1,\dots,\boldsymbol{\zeta}_m$,
in which every monomial has as many barred as unbarred legs, and which have the following properties.
\begin{itemize}
\item[(k1)] (Analyticity.)
The kernels are analytic functions of $(\mathscr{A},\mathscr{C})$ on $\tilde{R}_k\times\mathfrak{X}_k$.
\item[(k2)] (Locality.)
The kernels vanish unless every $\xi_i$ lies over $X$ and every pair $\boldsymbol{\zeta}_j$ is admissible for $X$ in the sense of Section~I.1.7.
They depend on $\mathscr{A}(b)$ and $\mathscr{C}(b)$ only for the bonds $b$ that meet $X$.
\item[(k3)] (Lattice symmetry.)
For every $\tilde{g}\in\mathbb{G}_k$, acting on the sites by (I.4.2), one has
\[
a_{\tilde{g}}\left[\mathscr{E}(\tilde{g}X;\tilde{g}\mathscr{A},\tilde{g}\mathscr{C})\right]=\mathscr{E}(X;\mathscr{A},\mathscr{C}).
\]
\item[(k4)] (Gauge invariance, local form.)
Let $\boldsymbol{\lambda}$ and $\boldsymbol{\lambda}'$ be real functions of $c$, and let $(\mathscr{A},\mathscr{C})$ and $(\mathscr{A}^1,\mathscr{C}^1)$ be two pairs of connections in $\tilde{R}_k\times\mathfrak{X}_k$ with $\mathscr{A}^1=\mathscr{A}-d\boldsymbol{\lambda}$ and $\mathscr{C}^1=\mathscr{C}-d\boldsymbol{\lambda}'$ on the bonds at distance $\le1$ from $X$.
Then
\[
G_{\boldsymbol{\lambda}+\boldsymbol{\lambda}'}\left[\mathscr{E}(X;\mathscr{A}^1,\mathscr{C}^1)\right]=\mathscr{E}(X;\mathscr{A},\mathscr{C}).
\]
\item[(k5)] (Charge conjugation.)
$\mathrm{Cg}\left[\mathscr{E}(X;-\mathscr{A},-\mathscr{C})\right]=\mathscr{E}(X;\mathscr{A},\mathscr{C})$.
\item[(k6)] (Norm.)
The number
\[
\|\mathscr{E}\|^x_k:=\sup_X\exp\left(\kappa\cdot d_M(X)\right)\cdot\sup_{(\mathscr{A},\mathscr{C})\in\tilde{R}_k\times\mathfrak{X}_k}\left\|\mathscr{E}(X;\mathscr{A},\mathscr{C})\right\|_{h_k}
\]
is finite.
\end{itemize}
The space $K^{x}_{k}$ is the real subspace of $K^{x,C}_{k}$ formed by the families with the property
\begin{itemize}
\item[(k7)] (Reality.)
$\mathbf{K}\left[\mathscr{E}(X;-\overline{\mathscr{A}},-\overline{\mathscr{C}})\right]=\mathscr{E}(X;\mathscr{A},\mathscr{C})$.
\end{itemize}
The spaces $K_{k}^{C}$, $K_{k}$ and the norm $\|\cdot\|_k$ are defined in the same way for families $E(X;\mathscr{A},\psi,\chi)$ on $\tilde{R}_k$ alone, that is, without a second slot.
By Proposition~15\,(a), which is the rule (I.N0) of Section~I.5.2, every element of $\mathrm{Gr}[\psi,\chi]$ has exactly one system of kernels with this antisymmetry.
Hence the kernels of (k1) and (k2), and the kernels $E_{00}$ and $E_{20}$ of Section~I.5.6, are determined by the family, and the norm in (k6) is
\[
\left\|\mathscr{E}(X;\mathscr{A},\mathscr{C})\right\|_{h_k}=\sum_{n,m}\frac{h_{k,1}^{n}\cdot h_{k,2}^{m}}{n!\cdot m!}\cdot\left\|\mathscr{E}_{nm}(X;\mathscr{A},\mathscr{C})\right\|,
\]
which is the case of equality in Section~I.5.1.
The same holds for the spaces $K_{k}^{C}$ and $K_k$.
For $\mathscr{E}\in K^{x,C}_{k}$ put $\mathscr{E}|_0:=\mathscr{E}(\,\cdot\,;\,\cdot\,,0,\,\cdot\,)$.
By Proposition~48.1\,(a), $\mathscr{E}|_0$ lies in $K_{k}^{C}$, and in $K_k$ if $\mathscr{E}\in K^{x}_{k}$, and $\|\mathscr{E}|_0\|_k\le\|\mathscr{E}\|^x_k$.

\begin{named}{Remark}
The space $K^{x,C}_{k}$ corresponds to the space of activities of \cite[Definition~5.24]{QEDI}.
As there, the kernels are antisymmetric in the $\xi_i$ and, separately, in the $\boldsymbol{\zeta}_j$.
The difference is that here the kernels are functions on the index sets of the fermion lattice $T^{\mathrm{f}}_k$.
\end{named}

\paragraph{Fields to be inserted.}
The sum $\mathscr{T}:=\mathscr{A}+\mathscr{C}$ is called the total connection.
The fields that are inserted into an activity are
\[
\psi=H_k(\mathscr{T})\Psi,\qquad\chi=\delta_{\alpha,\mathscr{T}}\psi,\qquad\bar{\psi}=\bar{H}_k(\mathscr{T})\bar{\Psi},\qquad\bar{\chi}=\delta_{\alpha,-\mathscr{T}}\bar{\psi}.
\]
Put
\[
w^{\mathrm{m}}_k:=e_k^{\frac{3}{4}-8\varepsilon},\qquad w^{\mathrm{E}}_k:=e_k^{\frac{1}{4}-7\varepsilon},
\]
and, for $k\ge1$ and a number $\mathrm{cf}$,
\[
\mathrm{Ball}_k(\mathrm{cf}):=\left\{(m,E)\in\mathbb{C}\times K_{k}^{C}:\ |m|\le\mathrm{cf}\cdot w^{\mathrm{m}}_k,\ \|E\|_k\le\mathrm{cf}\cdot w^{\mathrm{E}}_k\right\};
\]
for $k\ge1$ the set $\mathrm{Ball}^{x}_{k}(\mathrm{cf})$ is defined in the same way with $K^{x,C}_{k}$ and $\|\cdot\|^x_k$ in the places of $K_{k}^{C}$ and $\|\cdot\|_k$.
For $k=0$ both sets are, by definition, equal to
\[
\left\{(m,0):\ m\in\mathbb{C},\ |m|\le\mathrm{cf}\cdot w^{\mathrm{m}}_0\right\},
\]
where $w^{\mathrm{m}}_0$ is formed with the charge of level 0.
A leg is a factor $\psi(\xi)$ or $\chi(\boldsymbol{\zeta})$ of a monomial.
It is barred if its generator is barred, and a leg $\chi(\boldsymbol{\zeta})$ is called a H\"older leg.

\subsubsection*{I.5.5. Legs inside a cluster}

The values of an activity are elements of $\mathrm{Gr}[\psi,\chi]$, whose only generators are the legs $\psi(\xi)$ and $\chi(\boldsymbol{\zeta})$ of Section~I.5.4.
The heavy part $\mathrm{Phv}\,\psi=\psi-J\cdot J^{*}\psi$ of (I.J4) is expressed through these legs, as follows.
At $\chi=\delta_{\alpha,\mathscr{T}}\psi$ a difference inside a cluster is the H\"older leg
\[
\psi(x')-\psi(x)=\eta_k^{\alpha}\cdot\chi(x,x')
\]
of length $d_X=\eta_k$;
indeed, a pair inside a cluster has $c=c'$, $g\ne g'$ and $d_X=\eta_k$ (Section~I.1.7), and the H\"older quotient (I.1.7) carries no phase at such a pair.
Since $J^{*}$ is the mean over the cluster (I.3.1), $(\mathrm{Phv}\,\psi)(x)$ is the mean of the differences $\psi(x)-\psi(x')$ over the 48 sites $x'$ of the cluster of $x$.
Hence, at $\chi=\delta_{\alpha,\mathscr{T}}\psi$,
\[
(\mathrm{Phv}\,\psi)(x)=-\frac{\eta_k^{\alpha}}{48}\sum_{x'}\chi(x,x'),
\]
where $x'$ ranges over the sites $x'\ne x$ of the cluster of $x$.
Thus the heavy part enters an activity through $\psi$ and through the H\"older legs of the pairs inside a cluster.
In the norm $\|\cdot\|_{h_k}$ of Section~I.5.1 every H\"older leg, whatever its pair, carries the weight $h_{k,2}=e_k^{-\frac{1}{4}+\varepsilon}$;
the factor $\eta_k^{\alpha}$ of the first display is not part of the norm.
What enters the estimates is the factor $L^{-1-\alpha}$ by which $S_L$ multiplies every H\"older leg, whatever its pair (Section~I.3.7 and Proposition~18.4\,(a), (b)).
In the contraction bound, Proposition~18.5\,(d), the norm $\|\cdot\|_{h_{k+1}}$ of the image under $S_L$ of the part of bidegree $(1,1)$, respectively $(0,2)$, of $(R\mathscr{E})^{\mathrm{nat}}(X;\mathscr{A}_L,\mathscr{C}_L)$ for a small polymer $X$ is the norm $\|\cdot\|_{h_k}$ of that part times the factor $L^{-2-\alpha-\frac{1}{4}+\frac{\varepsilon}{2}}$, respectively $L^{-2-2\alpha-\frac{1}{4}+\varepsilon}$;
these factors do not depend on the labels and contain no factor $\eta_k^{\alpha}$.

\subsubsection*{I.5.6. Extraction, adjustment, reblocking}

Let $E$ be a family of activities of level $k$.
For a small polymer $X$ put
\[
\mathrm{al}_0(E,X):=\operatorname{Vol}(X)^{-1}\cdot E_{00}(X;0,0),\qquad\mathrm{al}_2(E,X):=\operatorname{Vol}(X)^{-1}\int\mathrm{d}x\int\mathrm{d}x'\,E_{20}(X;0,0;x,x').
\]
Thus $\mathrm{al}_2(E,X)$ is a $4\times4$ matrix.
Each of the two integrals, over $x$ and over $x'$, contains the mean $\frac{1}{48}\sum$ over the label of its site;
thus the two-point kernel is integrated over both cluster labels, in accordance with (I.E1) of Section~I.4.6.
For an $M$-cube $\square$ let
\[
\varepsilon({E}):=-\sum_{X\text{ small},\,X\supset\square}\mathrm{al}_0(E,X),\qquad m({E})\,1:=-\sum_{X\text{ small},\,X\supset\square}\mathrm{al}_2(E,X),
\]
and define
\begin{align*}
(RE)(X)&:=E(X)-\mathrm{al}_0(E,X)\cdot\operatorname{Vol}(X)-\int_X\mathrm{d}x\,\bar{\psi}(x)^T\,\mathrm{al}_2(E,X)\,\psi(x)&&(X\text{ small}),\\
(RE)(X)&:=E(X)&&(X\text{ large}),\\
E^{\mathrm{nat}}(X)&:=E(X)-E^{[2,0]}(X)+\boldsymbol{\Theta}_{X,\mathscr{T}}\left(E^{[2,0]}(X)\right)&&(X\text{ small}),\\
E^{\mathrm{nat}}(X)&:=E(X)&&(X\text{ large}),
\end{align*}
\[
(LE)(Y;\mathscr{A},\mathscr{C},\psi,\chi):=\sum_{X:\,\bar{X}^{L}=LY}E\left(X;\mathscr{A}_L,\mathscr{C}_L,\psi_L,\chi_L\right),\qquad LRE:=L\left((RE)^{\mathrm{nat}}\right).
\]
That the sum which defines $m({E})\,1$ is a multiple of the unit matrix, with a number $m({E})$ that is real for $E\in K^{x}_{k}$ (for $E\in K_k$), is part of Proposition~16.3\,(a); for families without slot it is part of Proposition~48.2\,(a).
The extracted bilinear term is $-m({E})\cdot\langle\bar{\psi},\psi\rangle$ (Proposition~16.3\,(c)); together with the term $-m\cdot\langle\bar{\psi},\psi\rangle$ of $I^{x}_{k}$ (Section~I.5.7) it gives $-\left(m+m({E})\right)\cdot\langle\bar{\psi},\psi\rangle$.
It is uniform on the 48 sites of every cluster, and it has the form of the mass assignment of \cite[Definition~2.36]{QEDII}.

\begin{named}{Remark}
The operations $\mathrm{al}_0$, $\mathrm{al}_2$, $\varepsilon$, $m$, $R$, $(\cdot)^{\mathrm{nat}}$ and $L$ correspond to those of \cite[Section~5.3.5]{QEDI}.
Here the integrals over the sites are taken with the measure $\int\mathrm{d}x$ of Section~I.1.6.
\end{named}

\paragraph{(I.SY6)$^{\mathrm{ch}}$: the chain of the adjustment.}
For points $p$, $p'$ let $u(p,p'):=U_{\mathscr{T}}(p,p')$ be the transport (I.1.6) formed with the total connection.
Thus $u(p,p')=1$ if $p$ and $p'$ lie over the same photon site.
Let $X$ be a small polymer, let $x=(c,g)$ and $\tilde{x}=(\tilde{c},\tilde{g})$ lie over $X$, and let $(\square_0,\dots,\square_n)$ be a shortest chain of $M$-cubes of $X$ from the cube of $c$ to the cube of $\tilde{c}$ in which consecutive cubes have a common face.
Let $z_i$ be the centre of $\square_i$; it is a site of $\mathbb{L}_k$.
The points of the chain are
\[
p_0:=x=(c,g),\qquad p_{i+1}:=(z_i,g)\quad(0\le i\le n),\qquad p_{n+2}:=\tilde{x}=(\tilde{c},\tilde{g});
\]
so the label $g$ of the starting point is kept at all intermediate points.
Put
\[
U_i:=\prod_{j<i}u(p_j,p_{j+1}),
\]
and let $\operatorname{Mn}$ denote the mean over the shortest chains.
The homomorphism $\boldsymbol{\Theta}=\boldsymbol{\Theta}_{X,\mathscr{T}}$ from the subalgebra generated by the $\psi(\xi)$, $\xi$ over $X$, into $\mathrm{Gr}[\psi,\chi]$ is defined by
\begin{equation}\tag{I.5.3}
\boldsymbol{\Theta}\psi(x):=\operatorname{Vol}(X)^{-1}\int_X\mathrm{d}\tilde{x}\,\operatorname{Mn}\left\{U_{n+2}\cdot\psi(\tilde{x})-\sum_{i=0}^{n+1}U_i\cdot d_X(p_i,p_{i+1})^{\alpha}\cdot\chi(p_i,p_{i+1})\right\},
\end{equation}
and $\boldsymbol{\Theta}\bar{\psi}(x)$ is defined by the same formula with $u^{-1}$, $\bar{\psi}$, $\bar{\chi}$ in the places of $u$, $\psi$, $\chi$.
The terms with $p_i=p_{i+1}$ are left out.
The integral $\int_X\mathrm{d}\tilde{x}$ contains the mean $\frac{1}{48}\sum_{\tilde{g}}$ over the label of $\tilde{x}$, so that no label is distinguished.
All points are chosen over one lift of $X$.

The following facts are read off from these definitions.
Every pair $(p_i,p_{i+1})$ is of kind (I) for $X$.
The first $n+1$ pairs join points with equal labels.
The last pair carries the arbitrary labels $(g,\tilde{g})$; if $\tilde{c}=z_n$, it lies inside a cluster, and then $d_X=\eta_k$ and there is no phase.
For every pair of the chain one has $d_X^{\alpha}\le\max(1,d_X)\le M$.
An element $(S,a)$ of $\mathbb{G}_k$ maps the points of the chain that belongs to $(X,x,\tilde{x})$ and to the chain $(\square_0,\dots,\square_n)$ onto the points of the chain that belongs to the images of $X$, $x$, $\tilde{x}$ and of $(\square_0,\dots,\square_n)$, because it maps $(z_i,g)$ to the site formed by the centre of the image of $\square_i$ and the label $r(S)g$.

\begin{named}{Remark}
(I.SY6)$^{\mathrm{ch}}$ corresponds to the transport to the mean of \cite[Definition~13.6]{QEDI}, where the points $x$, $\tilde{x}$ are written $y$, $y_0$ and the points of the chain are not fermion sites.
Here every point of the chain is a fermion site, and the symbol $x^{(0)}$ denotes the neighbour of type 0 of $x$.
\end{named}

\subsubsection*{I.5.7. Densities, block transformation, step map, flows}

The photon objects $H_k$, $\Delta_k$, $C$, $C_{k}^{1/2}$, $\mathrm{Sq}_k$, $A^{\min}_k$, $\omega_k$, $V_k$, $\{\Upsilon\}$, $\mathrm{Box}_k$, $\mu_I(\zeta)$, $\hat{\mu}$, $\delta\varepsilon^{0}_{k}$, the slice $\mathrm{Sl}$, $\mathfrak{G}$, $\mathfrak{G}^{-}$, $A_j$, $\zeta_j$, $\mathrm{d}\mathsf{m}_{*}$, $\chi_{*}$ and $\mathrm{Af}_0$ are those of \cite[Section~5.2.5]{QEDI} and \cite[Definitions~5.10 and~5.11]{QEDI}.
Here $k_{*}=k_{*}(R)$ denotes the stop level of Section~I.6.2, which is defined under (FL), that is, for $R\in\mathrm{Tor}(e)$ and $N\ge N_{\min}(R,e)$, and $T^0_{k_{*}}$ is the unit lattice of that level on the torus of side $R$.
For a real bond function $A$ on $T^0_{k_{*}}$ and a family $\zeta=(\zeta_j)_{j<k_{*}}$ with $\zeta_j\in V_j$, the field $\mathrm{Af}_0(A,\zeta)$ is the field $A^{\downarrow}_0(A,\zeta)$ of \cite[Definition~5.11]{QEDI}, the last of the fields given by the downward recursion
\[
A^{\downarrow}_{k_{*}}:=A,\qquad A^{\downarrow}_j:=A^{\min}_j\left(A^{\downarrow}_{j+1}\right)+\mathrm{Sq}_j\zeta_j\quad(j=k_{*}-1,\dots,0).
\]
The photon minimiser is written $H_kA$, without a connection as argument; $H_k(A)$, with a connection as argument, is the fermion operator of Section~I.3.5.
In both cases $H_0:=1$.

\paragraph{Settings.}
The objects $I^{x}_{k}$, $\mathrm{Act}^x_k$, $\mathscr{F}^x_k$ and $y^{x}_{k}$ are formed under (SF) (Section~I.1.8).
There $\mathscr{T}=\mathscr{A}+\mathscr{C}$ and $-\mathscr{T}$ lie in $\tilde{G}_k$ by \cite[Lemmas~17.4(a) and~17.6]{QEDI}, so that $H_k(\mathscr{T})$, $\bar{H}_k(\mathscr{T})=\mathsf{C}^{-1}H_k(-\mathscr{T})\mathsf{C}$ and $D_k(\mathscr{T})$ exist (Proposition~14.1\,(i)).
By the definition of $\mathrm{Dom}_k$, $H_kA$ lies in $\frac{1}{4}\tilde{R}_k$ for $A\in\mathrm{Dom}_k$.
Moreover $\det D_k(0)\ne0$ (Proposition~19\,(d); for $k=0$, Proposition~19\,(a)).
For $\mathbf{T}^{x}_{k}$ and for every object of level $k+1$ one requires in addition $\ell/L\in\mathrm{Cl}$.
The functions $\acute{E}_k$, $E^{d}_k$, $E_k^{\mathrm{det}}$ and ${\delta E^{\mathrm{det}}}_k$ are formed under (SF), with a family (p5) that has the properties named in Proposition~13.
The step map is formed under (ST1)$^{\mathrm{ch}}$ (Sections~I.1.8 and~G3.4.0), with the system of primitives of this setting.
The setting (ST1)$^{\mathrm{ch}}$ consists of the data $(k,\ell,\mu,e_k)$ of (OS) together with a system of primitives of level $k$ on $T_k$ that has the properties W1--W8 and W10--W17 for $\mathrm{par}$; W9 is required at no level, and at $k=0$ the property W17 is taken without W9, that is, with W8, W10, W11 in the place of its W8--W11 and with W6--W8 in the place of its W6--W9.
The data of the step map are $(m,\mathscr{E})\in\mathrm{Ball}^{x}_{k}(4)$, with $\mathscr{E}=0$ at $k=0$, $A\in\tilde{R}_k^{-}$ and $\mathscr{C}\in\mathfrak{X}_k$; for them every display of Definition~24.D is defined (Proposition~24.1\,(a)).
Since (OS) implies (ST1)$^{\mathrm{ch}}$, the step map is defined under (OS), and hence at every level $k<k_{*}(R)$ of a flow (Proposition~26.6\,(b)).
It is not formed under (SF), (SF+) or (OS-), which contain no system of primitives.
The flows are formed under (FL).
At every level $k<k_{*}(R)$ their data are (OS), and at $k=k_{*}(R)$ they are (SF) (Proposition~26.6\,(b), (c)).

\paragraph{The density and the block transformation.}
Let $(\mathscr{A},\mathscr{C})\in\tilde{R}_k\times\mathfrak{X}_k$, $\mathscr{T}=\mathscr{A}+\mathscr{C}$, let $m$ be a complex number and $\mathscr{E}\in K^{x,C}_{k}$.
With the fields of Section~I.5.4 inserted, that is, $(\psi,\chi)=H_k^{\#}(\mathscr{T})\Psi$ and the corresponding barred fields, put
\begin{align*}
I^{x}_{k}[m,\mathscr{E}](\mathscr{A},\mathscr{C};\Psi)&:=-m\cdot\langle\bar{\psi},\psi\rangle+\sum_X\mathscr{E}(X;\mathscr{A},\mathscr{C},\psi,\chi),\\
\mathrm{Act}^x_k[m,\mathscr{E}](\mathscr{A},\mathscr{C};\Psi)&:=-\left\langle\bar{\Psi},D_k(\mathscr{T})\Psi\right\rangle+I^{x}_{k}[m,\mathscr{E}](\mathscr{A},\mathscr{C};\Psi).
\end{align*}
Let $\mathfrak{P}_k=\left\{\mathfrak{P}_k(X,\cdot):X\in D_k\right\}$ be a family of complex functions on $\mathfrak{X}_k$.
It is a datum of $\mathscr{F}^x_k$ and of $y^{x}_{k}$ that the notation does not show; in statements on a single step it is arbitrary, and in a flow it is given by $\mathfrak{P}_0=0$ and the recursion of the paragraph on flows below.
For $A\in\mathrm{Dom}_k$ put
\begin{align*}
\mathscr{F}^x_k[\varepsilon^{\mathrm{v}},m,\mathscr{E}](A,\mathscr{C};\Psi)&:=\det D_k(0)^{-1}\cdot\exp\left\{\mathrm{Act}^x_k[m,\mathscr{E}](H_kA,\mathscr{C};\Psi)-\varepsilon^{\mathrm{v}}n_k+\sum_X\mathfrak{P}_k(X,\mathscr{C})\right\},\\
y^{x}_{k}&:=\int\mathrm{D}\Psi\,\mathscr{F}^x_k.
\end{align*}
For $k\ge1$ the block transformation is
\begin{multline*}
\left(\mathbf{T}^{x}_{k}F\right)(A',\mathscr{C}';\Psi'):=c_k\cdot\int_{\mathrm{Box}_k}\mathrm{d}\mu_I(\zeta)\int\mathrm{D}\Psi\,\exp\left\{-b_L\left\langle\exp\left(-ie_k\omega_k\right)\bar{\Psi}'_L-Q(-\mathscr{B}_k-\mathscr{C})\bar{\Psi},\right.\right.\\
\left.\left.\exp\left(ie_k\omega_k\right)\Psi'_L-Q(\mathscr{B}_k+\mathscr{C})\Psi\right\rangle_L\right\}\cdot F\left(A^{\min}_k(A')+\mathrm{Sq}_k\zeta,\mathscr{C};\Psi\right),
\end{multline*}
and $\mathbf{T}^{x}_{0}F$ is given by the same display with $Q(-\mathscr{B}_0-\mathscr{C})J^{*}\bar{\Psi}$ and $Q(\mathscr{B}_0+\mathscr{C})J^{*}\Psi$ in the places of $Q(-\mathscr{B}_k-\mathscr{C})\bar{\Psi}$ and $Q(\mathscr{B}_k+\mathscr{C})\Psi$.
Here
\[
\mathscr{C}:=\mathscr{C}'_L,\qquad\mathscr{B}_k:=H_kA^{\min}_k(A'),\qquad\omega_k=\omega_k(A'),\qquad\langle v'',v'\rangle_L:=L^3\sum_{y'\in L\mathbb{Z}^3/\ell\mathbb{Z}^3}v''(y')\cdot v'(y').
\]
The forms $\langle\bar{\psi},\psi\rangle$, $\langle\bar{\Psi},O\Psi\rangle$ and the integral $\int\mathrm{D}\Psi$ are those of Section~I.3.6.
The mass term $-m\cdot\langle\bar{\psi},\psi\rangle$ extends over all 192 components.
At $\mathscr{C}=0$ (for $\mathbf{T}^{x}_{k}$: at $\mathscr{C}'=0$) the objects $\mathscr{F}^x_k$, $y^{x}_{k}$, $\mathbf{T}^{x}_{k}$ are written $\mathscr{F}_k$, $y_k$, $\mathbf{T}_k$, as in \cite[Section~5.3.6]{QEDI}.

\paragraph{Level 0 integrates to the determinant.}
Consider the flow, that is, the side $\ell=n$ of Section~I.6.2, the mass $\mu=\bar{m}\eta$ and the family $\mathfrak{P}_0=0$.
By (I.B1), for every real pair $(m_{\mathrm{in}},\varepsilon^{\mathrm{v}}_{\mathrm{in}})$, every $A_0\in\mathrm{Dom}_0$ and every $a\in\mathfrak{X}_0$,
\[
y^{x}_{0}[\varepsilon^{\mathrm{v}}_{\mathrm{in}},m_{\mathrm{in}},0](A_0;a)=\exp\left(-\varepsilon^{\mathrm{v}}_{\mathrm{in}}n_0\right)\cdot\frac{\det\left(D_{A_0+a}+\bar{m}\eta+m_{\mathrm{in}}\right)}{\det\left(D_0+\bar{m}\eta\right)};
\]
here $D_0$ denotes the operator $D_A$ at $A=0$, and not the operator $D_0(A)$ of (I.L0.7).
For an arbitrary family $\mathfrak{P}_0$ the right side carries in addition the factor $\exp\left\{\sum_X\mathfrak{P}_0(X,a)\right\}$.
This is Proposition~22.2\,(c).

\paragraph{Step map.}
The step map is defined in Definition~24.D (Section~G3.4.1), by the rule (SLOT) and the displays (8.1)$^{\mathrm{ch}}$, (8.2), (8.3)$^{\mathrm{ch}}$, (8.4), (8.6)$^{\mathrm{ch}}$ there, by the two displays of Section~I.5.9, and by the contours
\[
\mathrm{rad}_1:=\frac{1}{2}\,e_k^{-\frac{3}{4}+\frac{9\varepsilon}{2}},\qquad\mathrm{rad}_2:=e_k^{-\frac{3}{4}+7\varepsilon},\qquad\mathrm{rad}_4:=e_k^{-\frac{1}{4}+3\varepsilon}.
\]
The weakened fluctuation field is $\mathscr{Z}(s'):=H_k(s')\,C\,C_{k}^{1/2}(s')\,\mathsf{Z}$, where $\mathsf{Z}\in\mathrm{Box}_k$ is the fluctuation variable of (D1) of Definition~24.D.
The step map with slot is denoted $E^{*,x}_k$; it is given by (8.6)$^{\mathrm{ch}}$ of Definition~24.D.
The step map without slot is denoted $E^{*}_k$; it is given by the reading (D7) of Definition~24.D.

\begin{named}{Remark}
The displays of Definition~24.D take the places of \cite[Definition~5.25]{QEDI} and of the corresponding displays of \cite[Section~17.4]{QEDI}.
The two displays of Section~I.5.9 take the places of the display of $T^{(1)}$, $T^{(2)}$ and of the display of $G^{(1)}$ at level zero in \cite[Section~17.4.4]{QEDI}.
The radii $\mathrm{rad}_1$, $\mathrm{rad}_2$, $\mathrm{rad}_4$ are denoted $\rho_1$, $\rho_2$, $\rho_4$ in \cite[Definition~5.25]{QEDI}.
\end{named}

\paragraph{Determinant.}
Let
\begin{align*}
\acute{E}_k(X,A)&:=-i\,b_k^2\int_0^{\infty}\mathrm{d}y\,\operatorname{Tr}\left[\gamma_5\cdot B_{k,y}(A)\cdot Q_k(A)J^{*}\cdot S_{k,y}(X,A)\cdot JQ_k^{T}(-A)\cdot B_{k,y}(A)\right]&&(k\ge1),\\
\acute{E}_0(X,A)&:=\int_1^{\infty}\frac{\mathrm{d}y}{y}\,\operatorname{Tr}\left[F_0(A)\cdot S_{0,y}(X,A)\right]-i\int_0^1\mathrm{d}y\,\operatorname{Tr}\left[\gamma_5\cdot S_{0,y}(X,A)\right].
\end{align*}
For $k\ge1$ the trace is the trace $\operatorname{Tr}O=\sum_z\operatorname{tr}O(z,z)$ on $\mathfrak{S}(1,\ell)$; for $k=0$ both traces are traces on $\mathfrak{S}^{\mathrm{ch}}(1,\ell)$, formed by (I.ke2).
Let $E^{d}_k(X,A)$ be the mean of the four numbers $\acute{E}_k(X,A)$, $\acute{E}_k(X,-A)$, $\overline{\acute{E}_k(X,\overline{A})}$ and $\overline{\acute{E}_k(X,-\overline{A})}$, and put
\[
E_k^{\mathrm{det}}(X,A):=E^{d}_k(X,A)-E^{d}_k(X,0),\qquad{\delta E^{\mathrm{det}}}_k(X;\mathscr{A},\mathscr{C}):=E^{d}_k(X,\mathscr{A}+\mathscr{C})-E^{d}_k(X,\mathscr{C}).
\]
These functions correspond to those of \cite[Definitions~5.26 and~5.28]{QEDI}.

\paragraph{Flows.}
Let $(m_{\mathrm{in}},\varepsilon^{\mathrm{v}}_{\mathrm{in}})$ be a real pair.
Put $\varepsilon^{\mathrm{v},t}_0=\varepsilon^{\mathrm{v}}_{\mathrm{in}}$, let the mass coupling of level 0 be $m_{\mathrm{in}}$, and put $E_0=0$, $\mathscr{E}_0=0$, $\mathfrak{P}_0=0$.
The higher levels are given by the recursion of \cite[Section~5.3.8]{QEDI}:
\begin{align*}
\varepsilon^{\mathrm{v},t}_{k+1}&=L^3\left(\varepsilon^{\mathrm{v},t}_k+\varepsilon(E_k)+\delta\varepsilon^{0}_{k}\right),\qquad m_{k+1}=L\left(m_k+m(E_k)\right),\\
E_{k+1}&=LRE_k+LE_k^{\mathrm{det}}+E^{*}_k(m_k,E_k),\\
\mathscr{E}_{k+1}&=LR\mathscr{E}_k+L{\delta E^{\mathrm{det}}}_k+E^{*,x}_k(m_k,\mathscr{E}_k),\\
\mathfrak{P}_{k+1}(Y,\mathscr{C})&=\sum_{X:\,\bar{X}^{L}=LY}\left(\mathfrak{P}_k(X,\mathscr{C}_L)+E_k^{\mathrm{det}}(X,\mathscr{C}_L)\right)\qquad\left(Y\in D_{k+1},\ \mathscr{C}\in\mathfrak{X}_{k+1}\right).
\end{align*}
In statements on a single step $\mathfrak{P}_k$ is an arbitrary family.

\subsubsection*{I.5.8. How the numbered properties of the primitives are read}

\paragraph{Families.}
The families of primitives are those of \cite[Section~5.4.1]{QEDI}, with the following spaces.
\begin{itemize}
\item[(p1)] $H_k(s')$, the photon family of \cite[Section~5.4.1]{QEDI}.
\item[(p2)] $C_{k}^{1/2}(s')$, the photon family of \cite[Section~5.4.1]{QEDI}.
\item[(p3)] $H_k(s'',A):\mathfrak{S}(1,\ell)\to\mathfrak{S}^{\mathrm{ch}}(\eta_k,\ell)$, together with $\bar{H}_k(s'',A):=\mathsf{C}^{-1}H_k(s'',-A)\mathsf{C}$.
\item[(p4)] $\Gamma_k(s'',A)$ and $M_k(X,A)$, operators on $\mathfrak{S}(1,\ell)$.
\item[(p5)] $S_{k,y}(X,A)$, operators on $\mathfrak{S}^{\mathrm{ch}}(\eta_k,\ell)$ for $y\ge0$, Borel measurable in $y$.
\end{itemize}
A system of level 0 consists of a family (p2) and of operators $\Gamma_0(s'',A)$ and $S_{0,y}(X,A)$ on $\mathfrak{S}^{\mathrm{ch}}(1,\ell)$, the latter Borel measurable in $y$.
At level 0 the families (p1) and (p3) are the diagonal families
\[
H_0(s')(b,b'):=\mathbf{1}[b=b']\cdot\prod\left\{s'_{\mathrm{pc}}:\ \mathrm{pc}\in\pi'_0,\ b\text{ contained in }\mathrm{pc}\right\}
\]
and $H_0(s'',A)$, the multiplication by
\[
\prod\left\{s''_{\mathrm{pc}}:\ \mathrm{pc}\in\pi_0,\ \mathrm{pc}\text{ meets }K_c\right\},
\]
which is the same factor on the 48 sites over $c$.

\paragraph{Reading of W1--W18, (Lo), FE1--FE8.}
In this paper W1--W18 denote the properties 1--18 of a system of primitives in \cite[Section~5.4.2]{QEDI}, (Lo) denotes \cite[Proposition~5.30]{QEDI}, and FE1--FE8 denote \cite[Propositions~5.42--5.49]{QEDI}.
For the chamber fermion they are read as follows.
A fine index of a fermionic kernel is a fermion site $x=(c,g)$, with $\operatorname{supp}x=\{c\}$, $K_x=K_c$ and $\square(x)=\square(c)$ (I.POS).
The absolute value of a kernel entry is the norm $|\cdot|_c$ of the $4\times4$ block at the two indices, and a supremum extends over sites or pairs and over components; kernels are those of (I.1.4).
Where these properties speak of pairs at distance smaller than 1, the short pairs, $d_X<1$, are meant, the pairs inside a cluster included; where they speak of all pairs of $\mathrm{Pz}_k$, the pairs inside a cluster are among them.
The row sums in (I.N2) and in (I.S6c)\,(iv) of Section~I.5.2, with $\chi\to\delta_{\alpha,B}H\Psi$, are suprema over all pairs of $\mathrm{Pz}_k$.

\paragraph{W7.}
For $k\ge1$, W7 is the operator bound of \cite[Section~5.4.2]{QEDI}, for all $y\ge0$.
At $k=0$, wherever W7 is proved or assumed, it is the kernel bound
\[
\left|S_{0,y}(X,A)(x,x')\right|_c\le C\cdot(1+y)^{-1}\cdot\exp\left(-\gamma\cdot d(x,x')-\mathrm{kq}\cdot d_M(X)\right),
\]
together with the Borel measurability in $y$.
Here $C$, $\gamma$ and $\mathrm{kq}$ are the entries of the set of parameters for which W7 is stated.
The rate per unit of $d_M$ is denoted $\mathrm{kq}$, as in Sections~G1.10 and~G1.11, while $\kappa'$ denotes the number $\kappa_{\mathrm{w}}+\kappa_0$ of Section~I.8.2.
The factor $(1+y)^{-1}$ cannot be dispensed with: without it the determinant bound (d1) of FE7 fails at $k=0$.

\paragraph{Required symmetry.}
Of the activities the group $\mathbb{G}_k$ is required, by clause (k3) of Section~I.5.4.
Of the primitives (p1)--(p4) the step group $\mathbb{G}_k^{L}$ is required.
Of the polymer resolvents (p5), $S_{k,y}(X,A)$, only the group $\mathbb{G}^{L}_{k,3}$ is required, because an element $S$ outside $\mathrm{Sp}_3$ turns $y$ into $-y$ (Section~I.4.4).
Since $r(\mathrm{Sp}_3)=\{\pm1\}^3$, the maps of the torus are the same in both cases.
Coordinate permutations, unit translations and the group $W_{\mathrm{af}}$ are not required.
Property W14 is read accordingly.

\begin{named}{Remark}
This requirement agrees with \cite[Definition~5.9]{QEDI} in all its parts.
There, too, the group $\mathbb{G}_k$ is required of the activities, the step group $\mathbb{G}_k^{L}$ of the families (p1)--(p4), and only the subgroup $\mathbb{G}^{L}_{k,3}$ of the polymer resolvents (p5), for the same reason: conjugation by an element $S$ outside $\mathrm{Sp}_3$ turns $y$ into $-y$, whereas (p5) is indexed by $y\ge0$.
In both papers the maps of the torus are the same for the step group and for this subgroup, because $r$ maps $\mathrm{Sp}_3$ onto the whole group of signs, and neither coordinate permutations nor unit translations are required.
Particular to this paper are the objects, which are those of the chamber fermion: the groups act on the fermion sites, and the matrix $\gamma_5$ takes the place of the matrix $\gamma_3$ of \cite[Section~5.2.4]{QEDI}.
Moreover the group $W_{\mathrm{af}}$, which does not occur in \cite{QEDI}, is named here among the symmetries that are not required.
\end{named}

\paragraph{(I.P6)$^{\mathrm{ch}}$: bounds of blocks.}
For $k\ge1$, a bound $|O(y,y')|_c\le\mathrm{cf}$ for an operator $O$ on $\mathfrak{S}(1,\ell)$ bounds the row sums of the block, hence each of its 16 entries.
At level 0 a block of an operator is its $4\times4$ block at a pair $(x,x')$ of fermion sites, and not its $192\times192$ block at a pair of photon sites.
This is the convention of \cite[Section~13.1.7]{QEDI} on bounds of $4\times4$ blocks, read for the chamber fermion.

\paragraph{The number $c_\gamma$.}
Put
\[
c_{\gamma}:=\sup_t\sum_{t'}\nu(t')\cdot\exp\left(-\gamma_{\mathrm{w}}\cdot d(t,t')\right).
\]
The supremum extends over all levels $k\ge0$, all sides $\ell\in\mathrm{Cl}$, all indices $t$ in the sense of Section~G4.0.4 (fine bonds, unit bonds, basis elements $\Upsilon$, unit sites, fermion sites; their supports have sup-diameter at most $L$; pairs are not indices here), and over the following three index sets in which $t'$ may range:
\begin{itemize}
\item[(i)] the pairs $(y,\beta)$, $y\in T^0_k$, $\beta\in\{1,\dots,4\}$, with $\nu=1$, for $k\ge1$, and at $k=0$ the pairs $(x,\beta)$, $x\in T^{\mathrm{f}}_0$, with $\nu=\frac{1}{48}$;
\item[(ii)] the positively oriented unit bonds, with $\nu=1$;
\item[(iii)] the basis $\{\Upsilon\}$, with $\nu=1$.
\end{itemize}
The bar index $\omega$ of Section~I.3.6 is not summed over.
Thus $c_\gamma$ is given by the formula of \cite[Sections~5.6.1 and~17.1.2]{QEDI}.
Its value differs from the value in \cite{QEDI}, because $\gamma_{\mathrm{w}}=\min\left(\gamma^{F},\gamma^{P}\right)$ depends on $\rho$ (Proposition~14.2\,(a)).
With kernels taken by (I.1.4), the factor 48 of level 0 enters only the constants of the kernel bounds W6 and W7 (through the constant $C_{00}$ of (I.L0.13)), and the factor $\sqrt{48}$ enters only the weight $\mathsf{w}_W$.

\subsubsection*{I.5.9. The two displays of the step that contain rows of $D_A$}

\paragraph{The operators of the substitution (8.3a)$^{\mathrm{ch}}$ ($k\ge1$).}
Let $\mathbf{Y}$ be a union of $M$-cubes of $T_k$ (in the step map, a polymer), let $B$ be a complex bond function on $T_k$ (in the step map, $B\in\tilde{G}_k$), and let $y$ be a unit site of $\mathbf{Y}$.
With $c=\pi(x)$ put
\begin{equation}\tag{I.5.4}
\begin{split}
\left(T^{(1)}(B)\psi\right)(y)&:=\int_{\Delta_y}\mathrm{d}x\,\varphi_B(y,c)\left\{\left(1+\frac{\mu}{b_k}\right)\psi(x)\right.\\
&\qquad\qquad\left.-(b_k\eta_k)^{-1}\sum_{i:\ x^{(i)}\text{ over }\mathbf{Y}}M_i(x)\left[U_B\left(x,x^{(i)}\right)\psi\left(x^{(i)}\right)-\psi(x)\right]\right\},\\
\left(T^{(2)}(B)\chi\right)(y)&:=-b_k^{-1}\eta_k^{\alpha-1}\int_{\Delta_y}\mathrm{d}x\,\mathbf{1}[c_3\notin\mathbf{Y}]\cdot\varphi_B(y,c)\cdot M_3(x)\cdot\chi\left(x,x^{(3)}\right),\\
\bar{T}^{(i)}(B)&:=\mathsf{C}^{-1}T^{(i)}(-B)\mathsf{C}.
\end{split}
\end{equation}
For $i\le2$ the condition that $x^{(i)}$ lies over $\mathbf{Y}$ always holds; so differences inside a cluster are always written with $\psi$.
If $\chi=\delta_{\alpha,B}\psi$, then $T^{(1)}(B)\psi+T^{(2)}(B)\chi=T_k(B)\psi$ at the unit sites of $\mathbf{Y}$, because $d_X(x,x^{(3)})=\eta_k$.

The set $\mathbf{Y}$ may be any union of $M$-cubes of $T_k$, for instance of the form $\mathrm{pj}^{-1}\mathbf{Y}$ as in Proposition~11\,(b).
At the unit sites $y$ outside $\mathbf{Y}$ the rows of $T^{(i)}(B)$ and of $\bar{T}^{(i)}(B)$ are zero.
Thus $T^{(1)}(B),\bar{T}^{(1)}(B):\mathfrak{S}^{\mathrm{ch}}(\eta_k,\ell)\to\mathfrak{S}(1,\ell)$ and $T^{(2)}(B),\bar{T}^{(2)}(B)$, which map the functions of pairs to $\mathfrak{S}(1,\ell)$, are operators between the spaces of Section~I.1.9, and the norm $|\cdot|_{1,\infty}$ of each of them is the supremum over the rows at the unit sites of $\mathbf{Y}$.
They are also written $T^{(i)}(B,\mathbf{Y})$ and $T^{(i)}_{\mathbf{Y}}(B)$.
For a polymer $\mathbf{Y}$ these operators enter the substitution (8.3a)$^{\mathrm{ch}}$ that precedes Proposition~23.1: the homomorphism of algebras that fixes the generators $\psi(\xi)$, $\chi(\boldsymbol{\zeta})$ and maps, at the unit sites $y$ of $\mathbf{Y}$,
\begin{gather*}
\Psi(y)\mapsto\left(T^{(1)}(B,\mathbf{Y})\psi\right)(y)+\left(T^{(2)}(B,\mathbf{Y})\chi\right)(y),\\
\bar{\Psi}(y)\mapsto\left(\bar{T}^{(1)}(B,\mathbf{Y})\bar{\psi}\right)(y)+\left(\bar{T}^{(2)}(B,\mathbf{Y})\bar{\chi}\right)(y).
\end{gather*}

\begin{named}{Remark}
The display (I.5.4) takes the places of the display of $T^{(1)}$, $T^{(2)}$ in \cite[Section~17.4.4]{QEDI} and of the display of $T_k$ in \cite[Section~16.11]{QEDI}.
\end{named}

\paragraph{The localised function $G^{(1)}$ at $k=0$: the display (8.3$_0$)$^{\mathrm{ch}}$.}
For an $M$-cube $\square$ put
\begin{equation}\tag{I.5.5}
\begin{split}
G^{(1)}(s,\square):=\frac{1}{2}\int_{\square}\mathrm{d}x\,&\left\{\bar{\psi}(x)M_3(x)\left(\exp\left(iq\,\mathscr{Z}(s')(c,c_3)\right)-1\right)\left(\psi(x)+\chi\left(x,x^{(3)}\right)\right)\right.\\
&\left.+\left(\bar{\psi}(x)+\bar{\chi}\left(x,x^{(3)}\right)\right)M_3\left(x^{(3)}\right)\left(\exp\left(-iq\,\mathscr{Z}(s')(c,c_3)\right)-1\right)\psi(x)\right\}.
\end{split}
\end{equation}
Here $\int_{\square}\mathrm{d}x$ is $\frac{1}{48}$ times the sum over the sites $x$ over $\square$, and $\bar{\psi}$ is a row.
Only bonds of type 3 occur.
Insert arbitrary fields on $T^{\mathrm{f}}_0$ for $\psi$ and $\bar{\psi}$, and $\delta_{\alpha,\mathscr{T}}\psi$, $\delta_{\alpha,-\mathscr{T}}\bar{\psi}$ for $\chi$ and $\bar{\chi}$.
Then the sum of (I.5.5) over the cubes becomes
\[
-\left\langle\bar{\psi},\left(D_{\mathscr{T}+\mathscr{Z}(s')}-D_{\mathscr{T}}\right)\psi\right\rangle^{\mathrm{ch}};
\]
in the step the inserted fields are $\Psi^{f}$, $\bar{\Psi}^{f}$ of Proposition~22.3, at $s=1$.
The identity, the four symmetries, the locality and the bound
\[
\left\|G^{(1)}(s,\square)\right\|_{h_0/2}\le12M^3q^{\frac{1}{2}-\frac{\varepsilon}{2}},
\]
valid if $q\le1$ and $q\left|\mathscr{Z}(s')(b)\right|\le\frac{1}{2}q^{1-\frac{\varepsilon}{2}}$ on the bonds $b$ that meet $\square$, are the clauses (a)--(d) of Proposition~23.2.

\begin{named}{Remark}
The display (I.5.5) takes the place of the display of $G^{(1)}$ at level zero in \cite[Section~17.4.4]{QEDI}.
\end{named}

\subsubsection*{I.5.10. Level 0: the differences}

At level 0 every object and every statement is the same as at the levels $k\ge1$, or as in \cite{QEDI} for level 0, except for the following differences (I.L0.1)--(I.L0.27).

\paragraph{A. Lattices, fields, integrals.}
\begin{itemize}
\item[(I.L0.1)] $\eta_0=1$.
The fine lattice and the unit lattice coincide: $T^{\mathrm{f}}_0=T^0_0\times O_{\mathrm{cub}}$.
For spinors the unit block of $y$ is the cluster $\{y\}\times O_{\mathrm{cub}}$, which has 48 sites and measure 1, and the $L$-block of $y'$ is $B(y')\times O_{\mathrm{cub}}$, which has $48L^3$ sites.
\item[(I.L0.2)] At $k=0$ the symbol $e_k$ denotes the charge of level 0, written $q$; it is never the ceiling $e_0$.
In the flow $q=e\eta^{\frac{1}{2}}$ and $\mu=\bar{m}\eta$; in the statements under (SF), (OS-), (OS) and (ST1)$^{\mathrm{ch}}$ the numbers $q$ and $\mu$ are data of the setting.
\item[(I.L0.3)] The fields $\Psi$, $\bar{\Psi}$, $W$, $\bar{W}$, the fields $W'$ and $\Psi^{c}$ of Proposition~21.4\,(a) and (21.9), and the field $\Psi^{f}$ of Proposition~22.3, together with their barred partners, are indexed by $T^{\mathrm{f}}_0$.
They have 192 components per unit site and $2\cdot192\,n_0$ generators in each family, with the weight $\nu=\frac{1}{48}$.
One has $\psi=\Psi$.
The block field $\Psi'$ of level 1 has 4 components.
\item[(I.L0.4)] The pairing of the unit lattice is $\langle\,,\,\rangle^{\mathrm{ch}}$, that is,
\[
\langle\bar{\Psi},O\Psi\rangle:=\int\mathrm{d}x\,\bar{\Psi}(x)\cdot(O\Psi)(x);
\]
kernels of operators on the unit lattice are taken by (I.1.4).
\item[(I.L0.5)] The integrals $\int\mathrm{D}\Psi$ and $\int\mathrm{D}W$ are given by (I.3.5) with $\nu=\frac{1}{48}$; that is, they are $48^{192n_0}$ times the Berezin integral of \cite[Definition~2.8]{QEDII}, and
\[
\int\mathrm{D}\Psi\,\exp\left(-\langle\bar{\Psi},O\Psi\rangle^{\mathrm{ch}}\right)=\det O.
\]
The measure $\mathrm{d}\mu_I(W)=\exp\left(-\langle\bar{W},W\rangle^{\mathrm{ch}}\right)\mathrm{D}W$ satisfies
\[
\int\mathrm{d}\mu_I(W)\,W_{\beta}(x)\,\bar{W}_{\beta'}(x')=48\cdot\mathbf{1}[x=x',\ \beta=\beta'],
\]
with the same number on every torus.
\item[(I.L0.6)] In the change of variables (Ber3) of \cite[Section~14.6.1]{QEDI}, with phases $u(c)$ that are the same on the 48 sites over $c$, the exponent 4 becomes 192.
In the rule (Ber4) of \cite[Section~14.6.1]{QEDI} the family that is replaced has $192n_0$ pairs of generators.
\end{itemize}

\paragraph{B. Operators.}
\begin{itemize}
\item[(I.L0.7)] On $\mathfrak{S}^{\mathrm{ch}}(1,\ell)$ (Section~I.3.9) one has
\begin{gather*}
Q_0=P_0=1,\qquad H_0(A)=\bar{H}_0(A)=T_0(A)=\bar{T}_0(A)=1,\qquad\dot{M}_0=1,\qquad D_0(A)=D_A+\mu,\\
F_y(A)=D_A+\mu+b_L\,JP(A)J^{*}+i\gamma_5y\quad\left(y\in Y_0=\mathbb{R}\right),\qquad\Gamma_0(A)=S^{0}_{1}(A)=F_0(A)^{-1},\\
S_{0,y}(A)=F_y(A)^{-1},\qquad\delta Z_0(A)=\det F_0(A),\qquad w_0(y)=(1+|y|)^{-1}.
\end{gather*}
Of the objects $\bar{H}_0$, $\bar{T}_0$, $S^{0}_{1}$, $P_1$ and $\dot{M}_0$ of level 0, the first three are listed in the display of the operators of level 0 in \cite[Section~5.2.6]{QEDI}, with the same values $\bar{H}_0=\bar{T}_0=1$ and $S^{0}_{1}=\Gamma_0$; there they act on the spinors on $T_0$, which have four components, and here on $\mathfrak{S}^{\mathrm{ch}}(1,\ell)$.
The operators $P_1$ and $\dot{M}_0$ of level 0 are not listed in that section: $\dot{M}_0:=1$ is set in the paragraph on level zero of \cite[Section~14.6.2]{QEDI}, and $P_1=P$ at $k=0$ stands in \cite[Proposition~5.55(d)]{QEDI} and in \cite[Section~16.2.4]{QEDI}.
\item[(I.L0.8)] The averaging operators of one step carry the cluster maps: they are $Q(A)J^{*}$, $JQ^{T}(-A)$ and $JP(A)J^{*}=J\cdot P_1(A)\cdot J^{*}$, where $P_1(A)=P(A)$ is the operator on $\mathfrak{S}(1,\ell)$ of Section~I.3.9.
For $k\ge1$ they do not carry them.
\item[(I.L0.9)] The following objects are not defined: $b_0$; $S_k(A)$ at $k=0$, that is, $S_{\infty}(A)$; $\mathfrak{H}_0(A)$; $M_0(A)$; $M_0(X,A)$; $B_{0,y}(A)$; $\Gamma_{0,y}(A)$; $T^{(1)}$ and $T^{(2)}$.
\item[(I.L0.10)] Sizes.
$\delta Z_0(A)$ is a determinant of size $192n_0$.
Since $192n_0\cdot\frac{i\pi}{2}=96\pi i\,n_0$ lies in $2\pi i\mathbb{Z}$ and $\det\left(\gamma_5\otimes1_{48n_0}\right)=1$, the value $c^{\mathrm{d}}_{0}=0$ and the display of $\acute{E}_0$ keep their form.
Further $c_0=b^{-4n_1}$ and $c'_0=b^{4n_1}$.
Let $B'$ be a complex bond function on the lattice $T_1$ of level 1, of side $\ell/L$, let $B:=B'_L$ (Section~I.3.7), and let $F_0(B)$ be invertible; then
\[
\det D_1(B')=b^{4n_1}\cdot\frac{\det D_0(B)}{\delta Z_0(B)},
\]
by Proposition~21.4\,(c) at $k=0$.
For a constant matrix $\mathrm{Mx}$,
\[
\operatorname{Tr}\left[\mathrm{Mx}\cdot JP(A)J^{*}\right]=n_1\operatorname{tr}\mathrm{Mx},\qquad\operatorname{Tr}\left[\mathrm{Mx}\cdot\left(1-JP(A)J^{*}\right)\right]=(48n_0-n_1)\operatorname{tr}\mathrm{Mx}.
\]
\item[(I.L0.11)] $\delta Z_0(A)$ contains the determinant of the 188 heavy components per photon site; it is expanded as one object.
\item[(I.L0.12)] Row sums.
For $A\in\tilde{G}_0$,
\[
\sup_x\sup_{\beta}\sum_{\beta'}\int\mathrm{d}x'\,\left|F_0(A)(x,x')_{\beta\beta'}\right|\le58\exp\left(\mathfrak{i}(A)\right)+1+b\,\mathrm{e};
\]
this uses $3(L-1)\,\mathfrak{i}(A)\le1$, which follows from the condition (Z2.s1) of Section~I.8.4, and $\mu\le1$.
In the proof of the bound of \cite[Proposition~16.31(b)]{QEDI} for $k=0$ it takes the place of the bound $(3+\mu)+6\exp\left(\mathfrak{i}(A)\right)+b_L\,\mathrm{e}$.
The entry $F_0(A)(x,x')$ reads $A$ on the bond $(c,c')$ or inside the $L$-block of $c$.
\item[(I.L0.13)] Input at zero field: the bound on the real part
\[
\operatorname{Re}\left(f,F_y(0)f\right)^{\mathrm{ch}}\ge\frac{c_{\mathrm{re}}}{L}\cdot(f,f)^{\mathrm{ch}}
\]
for the operator $F_y(0)$ of Section~I.3.9, $y$ real, and every spinor $f$ in $\mathfrak{S}^{\mathrm{ch}}(1,\ell)$ or on a block torus of level 0.
It is the bound of Proposition~1.7, because $F_y(0)$ is the operator $\mathrm{Tst}$ there at $\mathrm{me}=1$, with $\mathrm{Wd}$ the whole lattice and $b_E=b$.
At the levels $k\ge1$ the input at zero field is the regularity theorem for the blocked operator at zero field, Proposition~1.6, with its bounds of supremum, H\"older and exponential type; at level 0 that theorem is not the input, and the bound above takes its place.
The constants are
\begin{gather*}
c_{*}:=\frac{1}{18}\min\left(w_3,\lambda_{\mathrm{int}}\right),\qquad C_{\mathrm{Po}}:=\frac{1}{8c_{*}}+\frac{1}{\lambda_{\mathrm{int}}},\qquad c_{\mathrm{re}}:=\min\left\{b\left(1-\frac{1}{L}\right),\frac{\rho}{2C_{\mathrm{Po}}L}\right\},\\
a_z:=\min\left\{\frac{1}{3L},\frac{c_{\mathrm{re}}}{2\mathrm{e}L(6+3b)}\right\},\qquad k_1:=60+b,\qquad C_{00}:=48\cdot\frac{2L}{c_{\mathrm{re}}}\cdot(1+2k_1).
\end{gather*}
\end{itemize}

\paragraph{C. Pairs.}
\begin{itemize}
\item[(I.L0.14)] There is no short pair: $d_X\ge1$ for all pairs.
A pair inside a cluster has $d_X=1$ and no phase, and $(\delta_{\alpha,B}f)(x,x')=f(x')-f(x)$.
A bond of type 3 has $d_X=1$, and $\chi\left(x,x^{(3)}\right)=u_B(c,c_3)\psi\left(x^{(3)}\right)-\psi(x)$.
\item[(I.L0.15)] At $k=0$, FE4 is the single bound $\left|H_0^{\#}(A)f\right|_{\infty}\le(1+\mathrm{e})|f|_{\infty}\le4|f|_{\infty}$, together with its barred form.
The localised clauses of FE4 are stated for $k\ge1$ only.
\end{itemize}

\paragraph{D. Primitives.}
\begin{itemize}
\item[(I.L0.16)] A system of level 0 consists of (p2), $\Gamma_0(s'',A)$ and $S_{0,y}(X,A)$, with the diagonal families $H_0(s')$, $H_0(s'',A)$ of Section~I.5.8; $M_0(X,A)$ does not occur.
\item[(I.L0.17)] W6 and W7 are kernel bounds at pairs of fermion sites, with kernels taken by (I.1.4); their constant contains the factor 48, by (I.ke4).
W7 has the factor $(1+y)^{-1}$, and $S_{0,y}(X,A)$ is Borel measurable in $y$ (Section~I.5.8).
\item[(I.L0.18)] At level 0 the lists of exact operators in FE6 and FE8 contain $D_A$, $Q(A)J^{*}$, $JQ^{T}(-A)$, $JP(A)J^{*}$, $H_0=\bar{H}_0=T_0=\bar{T}_0=1$ (in FE8 also $\dot{M}_0=1$), $D_0(A)$, $\operatorname{Pen}_y$, $F_y$, $S_{0,y}(A)$, $\Gamma_0(A)$ and $\delta_{\alpha,A}$; the list in FE6 contains also $\delta Z_0(A)$.
\end{itemize}

\paragraph{E. Norms.}
\begin{itemize}
\item[(I.L0.19)] $\mathsf{w}_W=\sqrt{48}\le7$ in (I.N3)$^{\mathrm{w}}$ of Section~I.5.2.
The constant $C_W:=28\,C_w^2\cdot c_{\gamma}$ is used at all levels; in \cite[Section~17.1.2]{QEDI} one has $C_W=4C_w^2\cdot c_{\gamma}$.
The condition (A6) reads $112\,C_w^2\cdot c_{\gamma}\cdot e_k^{2\varepsilon}\le3$; it is a condition on $e_0$ only.
\item[(I.L0.20)] $c_{\gamma}$ is formed with the weight $\frac{1}{48}$ and keeps its formula of Section~I.5.8.
\end{itemize}

\paragraph{F. The first step.}
\begin{itemize}
\item[(I.L0.21)] The data are pairs $(m,0)$: the set $\mathrm{Ball}_0=\mathrm{Ball}^{x}_{0}$ consists of such pairs, and $E_0=\mathscr{E}_0=0$, $\mathfrak{P}_0=0$.
Nothing is extracted or adjusted: $\varepsilon(\mathscr{E}_0)=0$ and $m(\mathscr{E}_0)=0$, so that the number $m':=m+m(\mathscr{E})$ and the family $\mathscr{E}^{r}:=R\mathscr{E}$ of (D1) of Definition~24.D and of Proposition~22.3 are $m'=m$ and $\mathscr{E}^{r}=0$ at $k=0$.
The definitions of Section~I.5.6 make sense at $k=0$, and they are applied to the family $0$.
\item[(I.L0.22)] Since $H_0A=A$,
\begin{multline*}
\mathscr{F}^x_0[\varepsilon^{\mathrm{v}},m,0](A,\mathscr{C};\Psi)=\det D_0(0)^{-1}\cdot\exp\left\{\sum_{X\in D_0}\mathfrak{P}_0(X,\mathscr{C})\right\}\\
\cdot\exp\left\{-\left\langle\bar{\Psi},D_0(A+\mathscr{C})\Psi\right\rangle^{\mathrm{ch}}-m\langle\bar{\Psi},\Psi\rangle^{\mathrm{ch}}-\varepsilon^{\mathrm{v}}n_0\right\};
\end{multline*}
in the flow $\mathfrak{P}_0=0$ by (I.L0.21).
\item[(I.L0.23)] $\mathbf{T}^{x}_{0}$ contains $Q(\cdot)J^{*}$ (Section~I.5.7).
The Gaussian in $\Psi'$ has 4 components, and $c_0$ is the number of Section~I.3.4.
\item[(I.L0.24)] Completion of the square at $B=\mathscr{T}$, where $\mathscr{T}=\mathscr{T}'_L$ and $\mathscr{T}'=\mathscr{A}'+\mathscr{C}'$ is the total connection of level 1, as in Propositions~21.4 and~22.3: the field
\[
\Psi^{c}=b_L\,\Gamma_0(\mathscr{T})\,JQ^{T}(-\mathscr{T})\,\Psi'_L=\left(H_1(\mathscr{T}')\Psi'\right)_L
\]
has 192 components and is not constant on clusters.
One has $\Psi=\Psi^{c}+\Gamma_0(\mathscr{T})W$ and $\bar{\Psi}=\bar{\Psi}^{c}+\bar{W}$, with $\Psi^{c}$, $\bar{\Psi}^{c}$ as in Proposition~21.4 and (21.9); the two sums are the fields $\Psi^{f}$, $\bar{\Psi}^{f}$ of Proposition~22.3.
The quadratic form becomes $\left\langle\bar{\Psi}',D_1(\mathscr{T}')\Psi'\right\rangle+\langle\bar{W},W\rangle^{\mathrm{ch}}$, and the factor is $\delta Z_0(\mathscr{T})$.
\item[(I.L0.25)] The localised function $G^{(1)}$ is (I.5.5); it contains no term with $M_k(X,\cdot)$.
The display of \cite[Definition~5.26]{QEDI} that expresses $G^{(1)}$ through $M_k\left(X,\mathscr{T}+t\cdot\mathscr{Z}(s')\right)$ is a display for the levels $k\ge1$.
\item[(I.L0.26)] The determinant enters through $\acute{E}_0$ of Section~I.5.7.
In the bound (d1) of FE7 the number of sites of $X$ is $\int_X\mathrm{d}x=M^3|X|_M$, the same number as the count $\#|X\cap T^0_0|=M^3\cdot|X|_M$ of unit sites that enters the proof of \cite[Proposition~16.31(b)]{QEDI} at $k=0$; the constant $C_d$ is affected by (I.L0.12) and (I.L0.17).
\item[(I.L0.27)] Output.
$E_1$ and $\mathscr{E}_1$ are functions of $(\mathscr{A}',\mathscr{C}',\psi',\chi')$ with $\psi'$ on $T^{\mathrm{f}}_1$.
The H\"older legs that come from (I.5.5) are the pairs
\[
\left(x/L,x^{(3)}/L\right)=\left(x/L,(x/L)^{(3)}\right)
\]
of length $\eta_1$.
Further $m_1=Lm$ and $\varepsilon^{\mathrm{v},t}_1=L^3\left(\varepsilon^{\mathrm{v}}+\delta\varepsilon^{0}_{0}\right)$.
\end{itemize}

\paragraph{G. Places of \cite{QEDI} at which level 0 enters.}
The following places of \cite{QEDI} are read at level 0 with the entries named in parentheses after them.
Some of them treat level 0 separately.
The others are statements for all $k\ge0$ that do not mention level 0 but contain an object that the entries change at $k=0$: the weight 1 of the field $W$, in whose place $\mathsf{w}_W=\sqrt{48}$ stands, its covariance 1, in whose place 48 stands, or the unit sites as index set, in whose place the fermion sites $T^{\mathrm{f}}_0$ stand.
Statements of this second kind are \cite[Corollary~17.15]{QEDI}, the weights $(h_k/4;1)$ at the beginning of \cite[Section~17.8]{QEDI}, the hypothesis and Step~2 of the proof of \cite[Lemma~17.17]{QEDI}, clause (1) of \cite[Section~17.10.1]{QEDI}, \cite[Lemma~17.20]{QEDI} and \cite[Proposition~5.45]{QEDI}.
At the places of \cite{QEDI} at level 0 that are not named below, the entries (I.L0.1)--(I.L0.27) are applied in the same way.
\begin{itemize}
\item \cite[Section~5.2.6]{QEDI} ((I.L0.1), (I.L0.2), (I.L0.7)--(I.L0.11)); \cite[Definition~5.14]{QEDI} ((I.L0.16), (I.L0.21), (I.L0.25), (I.L0.26)); property~17 of \cite[Section~5.4.2]{QEDI} ((I.L0.16), (I.L0.17)); \cite[Proposition~5.44]{QEDI} ((I.L0.4), (I.L0.7), (I.L0.9), (I.L0.13), (I.L0.17)); \cite[Proposition~5.45]{QEDI} ((I.L0.14), (I.L0.15)); \cite[Proposition~5.47]{QEDI} ((I.L0.18)); \cite[Proposition~5.48]{QEDI} ((I.L0.10)--(I.L0.12), (I.L0.17), (I.L0.26)).
\item \cite[Lemma~13.19]{QEDI} at $k=0$ and the step $j=0$ of \cite[Theorem~13.20]{QEDI} ((I.L0.10)).
\item The rules (Ber1)--(Ber5) of \cite[Section~14.6.1]{QEDI}; the paragraph on level zero in \cite[Section~14.6.2]{QEDI}; \cite[Lemma~14.19(a)]{QEDI}; \cite[Theorem~14.21]{QEDI} at $k=0$; \cite[Lemma~14.22(b), (d)]{QEDI}; \cite[Lemma~14.23(c)]{QEDI}; \cite[Theorem~14.24]{QEDI} at $k=0$; \cite[Remark~14.25(ii)]{QEDI} ((I.L0.3)--(I.L0.10), (I.L0.21)--(I.L0.25)).
\item The last lines of \cite[Section~16.2.4]{QEDI}; \cite[Section~16.2.5]{QEDI}; the proof of \cite[Lemma~16.4]{QEDI}; \cite[Lemma~16.7]{QEDI}; the paragraph that follows the proof of \cite[Theorem~16.16]{QEDI}; the part of \cite[Proposition~16.21]{QEDI} that follows its clause~(c) and treats the case $k=0$, to the end of the statement, including the closing sentence, which defines $H_k^{\#}(A)$ and $\left|H_k^{\#}(A)f\right|_{\infty}$; \cite[Definition~16.24]{QEDI}; \cite[Theorem~16.25]{QEDI}; the last lines of \cite[Propositions~16.27 and~16.28]{QEDI}; \cite[Proposition~16.29]{QEDI}; \cite[Lemma~16.30]{QEDI}; \cite[Proposition~16.31(b)--(d)]{QEDI} ((I.L0.7)--(I.L0.18)).
\item The sentence on level 0 in the proof of \cite[Lemma~17.6]{QEDI}, which concerns the photon operator $H_0=1$ and is read with paragraph~H below; \cite[Lemma~17.11(q1), (q3), (q4)]{QEDI}; \cite[Lemma~17.14(a)]{QEDI} with Steps (1), (2), (4), (5) of its proof; \cite[Corollary~17.15]{QEDI}; the weights $(h_k/4;1)$ at the beginning of \cite[Section~17.8]{QEDI}; \cite[Proposition~17.16(b)]{QEDI} for $i=1$, and \cite[Proposition~17.16(c)]{QEDI}; the hypothesis and Step~2 of the proof of \cite[Lemma~17.17]{QEDI}; clause (1) of \cite[Section~17.10.1]{QEDI}; \cite[Lemma~17.20]{QEDI}; the places in the proofs of the four parts (a)--(d) of \cite[Lemma~17.21]{QEDI} that treat level 0, and Step~3 of the proof of part (c) of that lemma; Steps~0--2 of the proof in \cite[Section~17.12]{QEDI}; \cite[Proposition~17.24(b)--(d)]{QEDI} ((I.L0.3)--(I.L0.5), (I.L0.14)--(I.L0.25)).
\item \cite[Theorem~19.5(c)]{QEDI} with the paragraph on property~17 in its proof, and \cite[Lemma~19.7]{QEDI} ((I.L0.16), (I.L0.17)); Steps~2 and~4 of the proof of \cite[Proposition~19.12]{QEDI} ((I.L0.5), (I.L0.12)).
\item \cite[Lemma~20.4(b)]{QEDI}, clause (SL4) of \cite[Proposition~20.6]{QEDI} at $j=0$ and the display in part (vi) of that clause, which involve level 0 through the statements of \cite{QEDI} that they cite and are read with the entries that apply to those statements.
In the proof of \cite[Lemma~20.4(a)]{QEDI} the computation with the Wilson operator of level 0 is replaced by (I.O2), that is, $\operatorname{Re}D_A\ge0$ for real $A$, whence $\operatorname{Re}\left(D_a+z\right)\ge z$ for real $a$ and real $z$, and by (I.O3), that is, $\det\left(D_A+t\right)\ge0$ for real $A$ and real $t$ \cite[Proposition~3.18]{QEDII}, as in the proofs of Propositions~31.1 and~31.2.
\end{itemize}

\paragraph{H. Not different at level 0.}
All photon objects ($H_0=1$, (p2), $\mathrm{Box}_0$, $A^{\min}_0$, $\mathrm{Sq}_0$, $\omega_0$); the $M$-cubes, polymers and domains; the exponents of Section~I.8.5; the numbers $b_L$, $c_0$, $c'_0$ and $c^{\mathrm{d}}_{0}=0$; the bound $12M^3q^{\frac{1}{2}-\frac{\varepsilon}{2}}$ on $\left\|G^{(1)}(s,\square)\right\|_{h_0/2}$ (Section~I.5.9); the second part of (A3).

\subsubsection*{I.5.11. Statements without the second slot (conventions for the Propositions 48.m)}

A family on $\tilde{R}_k$ alone is not a special case of a family with two slots.
Therefore each of the Propositions 16.m, 17.m and 18.m that contains a family of activities has two forms with the same constants: one with slot ($K^{x,C}_{k}$, $K^{x}_{k}$) and one without slot ($K_{k}^{C}$, $K_k$).
Propositions~15, 16.1, 17.1, 18.1\,(a), (c), 18.3\,(a), (b) and 18.4 cover both forms; the forms without slot are the Propositions 48.m.

\paragraph{Definitions without slot.}
A family without slot of level $k$ is a family $E=\left\{E(X;\mathscr{A}):X\in D_k,\ \mathscr{A}\in\tilde{R}_k\right\}$ of elements of $\mathrm{Gr}[\psi,\chi]$.
The clauses (k1)$^{\mathrm{o}}$--(k7)$^{\mathrm{o}}$ are the clauses (k1)--(k7) of Section~I.5.4 with $\mathscr{C}$, $\mathscr{C}^1$, $\boldsymbol{\lambda}'$ and $\mathfrak{X}_k$ left out.
For such a family the definitions of Section~I.5.6 are read with $E(X;0)$, $E(X;\mathscr{A})$ and $\mathscr{T}:=\mathscr{A}$ in the places of $E(X;0,0)$, $E(X;\mathscr{A},\mathscr{C})$ and $\mathscr{T}=\mathscr{A}+\mathscr{C}$.
Thus, for small $X$,
\[
\mathrm{al}_0(E,X):=\operatorname{Vol}(X)^{-1}\cdot E_{00}(X;0),\qquad\mathrm{al}_2(E,X):=\operatorname{Vol}(X)^{-1}\int\mathrm{d}x\int\mathrm{d}x'\,E_{20}(X;0;x,x');
\]
for an $M$-cube $\square$ the quantities $\varepsilon(E,\square)$ and $m(E,\square)$ are the negatives of the sums of these over the small $X$ that contain $\square$,
\[
\varepsilon(E,\square):=-\sum_{X\text{ small},\,X\supset\square}\mathrm{al}_0(E,X),\qquad m(E,\square):=-\sum_{X\text{ small},\,X\supset\square}\mathrm{al}_2(E,X),
\]
so that $m(E,\square)$ is a $4\times4$ matrix;
for small $X$,
\begin{align*}
(RE)(X;\mathscr{A})&:=E(X;\mathscr{A})-\mathrm{al}_0(E,X)\cdot\operatorname{Vol}(X)-\int_X\mathrm{d}x\,\bar{\psi}(x)^T\,\mathrm{al}_2(E,X)\,\psi(x),\\
E^{\mathrm{nat}}(X;\mathscr{A})&:=E(X;\mathscr{A})-E^{[2,0]}(X;\mathscr{A})+\boldsymbol{\Theta}_{X,\mathscr{A}}\left(E^{[2,0]}(X;\mathscr{A})\right);
\end{align*}
for large $X$, $(RE)(X):=E(X)$ and $E^{\mathrm{nat}}(X):=E(X)$; and
\[
(LE)(Y;\mathscr{A}):=\sum_{X:\,\bar{X}^{L}=LY}S_L\left[E(X;\mathscr{A}_L)\right],\qquad LRE:=L\left((RE)^{\mathrm{nat}}\right).
\]
In the transition to global fields (Proposition~18.2\,(o)), $\mathscr{C}=0$ gives $\mathscr{C}^{b}=0$.
The recursion of Section~I.5.7 has the invariant
\[
E_k=\mathscr{E}_k|_0
\]
for all $k$ for which the trajectory exists (Proposition~48.6).
The objects of level $k+1$ that occur in a statement under the sole hypothesis $\ell\in\mathrm{Cl}$ (clause (d) of the group lemma, Proposition~16.1) need only that $\ell/L$ is a multiple of $M$ and that $\ell/L>2M$.

\subsection*{I.6. Scales and tori; cells; boundary conditions; twists}

\subsubsection*{I.6.1. Scales}

The bounds on the fermion operators hold for every mass in $(0,1]$, with constants that do not depend on the mass (Proposition~14.1\,(v)).
Consequently the mass range ends at $1$.
The scales are those of \cite[Section~5.2.2]{QEDI}, with three letters removed:
wherever \cite{QEDI} uses the upper end $\mathfrak{m}_0$ of the mass range, the integer $m_0$ or the number $h'$, one reads
\[
\mathfrak{m}_0=1,\qquad m_0=0,\qquad h'=h .
\]
For every level $k$ put
\begin{align*}
a_k&:=L^{k-N},&e_k&:=e\cdot a_k^{\frac{1}{2}},&v_k&:=-\log e_k,\\
\mathrm{lg}&:=\frac{1}{2}\log L,&\mu_k&:=\bar{m}\cdot a_k .&&
\end{align*}
Here $a_k$ is the physical length of one lattice unit of level $k$, and $a_0=\eta$;
at $k=0$ one has $e_k=e\cdot\eta^{\frac{1}{2}}=q$, and the charge of level~$0$ is always written $q$ (Section~I.2.5);
the numbers of level~$0$ defined in this section are formed with $q$, and the symbol $e_0$ is not used for this charge: it denotes the ceiling of the line (C8) of Section~I.8.2;
$\mu_k$ is the running mass, and at $k=0$ it is written $\bar{m}\cdot\eta$.
Apart from its use as a local letter in Chapter~39r (declared in Section~0.2 there), the symbol $a_0$ has three meanings, and no formula contains two of them.
As the scale $a_k$ at $k=0$ it equals $\eta=L^{-N}$, and it is written $a_0$ only together with this value:
in this section, in Chapter~26b, Section~2, in Chapter~35a, Section~1, and in Chapter~34, in (P1)\,(a) of Section~0.3 and in the proof of Proposition~34.3;
everywhere else the scale of level~$0$ is written $\eta$.
In the Remark after Proposition~1.6, in the quotient $a_0/(4L)$ and where that Remark speaks of the number $a_0$ of \cite[Section~12.2]{QEDI}, the symbol denotes that number,
\[
a_0=\min\left\{\frac{1}{3L},\frac{1}{2C_I\cdot c_3}\right\}\le\frac{1}{3L},
\]
which is formed with the constants $C_I$ and $c_3$ defined in \cite[Section~12.2]{QEDI}.
At all other places $a_0$ denotes the constant
\[
a_0=\min\left\{\frac{1}{3},\frac{1}{2C_I\cdot c_{\mathrm{wt}}}\right\}
\]
of Proposition~1.6, which is the row [62] of Proposition~14.2\,(a):
in the line (C4r) of Section~I.8.2, in Proposition~1.6 and in the rest of the Remark after it, in the Remark ``Where the mass enters'' that follows Proposition~14.1,
in Proposition~14.2 (the rows [62] and [66] and the paragraph ``Ranges'' of (f)), and in Chapters~1b, 2a, 14a, 14b and~14c.
This holds also in Proposition~14.2\,(d) and in Chapters~14b and~14c, where the scales $a_k$ and $a_*$ occur beside the constant;
the scale of level~$0$ is not written $a_0$ there.
Further,
\[
K:=N+\log_L h,\qquad\text{so that}\qquad a_K=h,\qquad e_K=e\cdot h^{\frac{1}{2}},\qquad\mu_K=\bar{m}\cdot h\in\left(\frac{1}{L},1\right],
\]
and
\[
r_k:=\min\left\{L^i:\ i\ge0\text{ an integer},\ L^i\ge v_k^r\right\},\qquad p_k:=v_k^p,\qquad p_{0,k}:=v_k^{p_0},
\]
where $r$, $p$ and $p_0$ are the integers of the line (C7) of Section~I.8.2.
For a torus of the physical side $R$ put
\[
\ell_k=\ell_k(R):=\frac{R}{a_k},\qquad n_k:=\ell_k^3,\qquad\mathfrak{n}_k:=\log_L\left(\frac{\ell_k}{M\cdot r_k}\right),\qquad T_k(R):=\mathbb{L}_k/\ell_k(R)\mathbb{Z}^3 .
\]

\paragraph{Coordinates of level $k$.}
One unit of length of level $k$ equals $a_k/\eta=L^k$ lattice units.
A photon site $c^0\in\mathbb{Z}^3$, given in lattice units, has at level $k$ the coordinates $c=L^{-k}\cdot c^0\in\mathbb{L}_k$.
Let $R$ be the physical side of a torus and $n:=R/\eta$ its number of sites per side.
Then $n=L^k\cdot\ell_k(R)$;
hence $c^0\mapsto L^{-k}\cdot c^0$ maps $n\mathbb{Z}^3$ onto $\ell_k(R)\mathbb{Z}^3$ and identifies $\mathbb{Z}^3/n\mathbb{Z}^3$ with $T_k(R)$.
A map $g_{\mathrm{w}}$ of $\mathbb{Z}^3$ (of $\mathbb{Z}^3/n\mathbb{Z}^3$), read in the coordinates of level $k$, is the map $c\mapsto L^{-k}\cdot g_{\mathrm{w}}\left(L^k\cdot c\right)$ of $\mathbb{L}_k$ (of $T_k(R)$).
The range of the charge is always expressed by the condition $e\cdot h^{\frac{1}{2}}\le e_1$.

\subsubsection*{I.6.2. The odd tori of the flow; the two reference tori}

The tori of the flow and the reference tori are those of \cite[Definition~5.6]{QEDI}, with $h'=h$.
Put
\begin{gather*}
R_{\mathrm{top}}(e):=h\cdot M\cdot r_K(e),\qquad\mathrm{Tor}(e):=\left\{R=h\cdot L^T:\ T\ge0\text{ an integer},\ R\le R_{\mathrm{top}}(e)\right\},\\
R^{*}(e):=\frac{R_{\mathrm{top}}(e)}{L}=h\cdot L^{T^{*}(e)},\qquad T^{*}(e):=\log_L\left(M\cdot r_K(e)\right)-1 ;
\end{gather*}
the number $T^{*}(e)$ is an integer $\ge1$.
The two reference tori are the tori of the sides $R^{*}(e)$ and $L\cdot R^{*}(e)=R_{\mathrm{top}}(e)$.
For $R\in\mathrm{Tor}(e)$ the number $n=R/\eta=h\cdot L^{T+N}$ of photon sites per side is odd.

Let $R\in\mathrm{Tor}(e)$.
The number $N_{\min}(R,e)$ is the least integer $N_1\ge0$ such that $R\cdot L^N\ge L\cdot M\cdot r_0$ for all $N\ge N_1$;
for $e\cdot h^{\frac{1}{2}}\le e_1$ it exists, by \cite[Proposition~19.1(b), (c)]{QEDI}, and for $R=R^{*}(e)$ also by Proposition~35.1.
The stop level $k_{*}=k_{*}(R)$ is the first level $k$ with $\mathfrak{n}_k=0$;
for $e\cdot h^{\frac{1}{2}}\le e_1$ and $N\ge N_{\min}(R,e)$ it exists and satisfies $1\le k_{*}\le K$, by \cite[Proposition~19.1(c)]{QEDI}, which is (Sc1) of Section~G4.0.2.
A subscript $*$ denotes the value at $k=k_{*}(R)$: $\ell_*=M\cdot r_{*}$, $e_{*}$, $\mu_*$, and so on.
Finally,
\[
k_{\mathrm{top}}:=k_{*}\left(R_{\mathrm{top}}\right),\qquad N_0(e):=N_{\min}\left(R^{*}(e),e\right).
\]

Let $e\cdot h^{\frac{1}{2}}\le e_1$, $R\in\mathrm{Tor}(e)$, $N\ge N_{\min}(R,e)$ and $k\le k_{*}(R)$.
Then $\ell_k(R)$ lies in $\mathrm{Cl}$.
Indeed, $\ell_k/(L\cdot M)=r_k\cdot L^{\mathfrak{n}_k-1}$ with $\mathfrak{n}_k\ge0$ and $r_k\ge r_{*}$.
The power $r_{*}$ of $L$ is at least $L^2$, by the condition (G4.2.1) of Section~I.8.4, $r_{*}>31L+16$, which holds by the row (e1) of the definition of $e_1$.
So $\ell_k/(L\cdot M)$ is a power of the odd number $L$ with an exponent $\ge1$, and $L\ge13$
(\cite[Proposition~19.1(e), (f)]{QEDI}, which are (Sc3) and (Sc4) of Section~G4.0.2).
The data (FL) are assumed, directly or through the settings (TT) and (ST*)$^{\mathrm{ch}}$ of Section~I.1.8 and the data (DAT) of Section~G5.4.0,
in Propositions~27.1--27.6, 28.1, 28.2, 28.3\,(a)--(c), 29, 31.2, 32.1--32.3, 33.1--33.4, 34.1--34.4 and~35.2,
in the clauses of Proposition~35.3 other than (b1) and (c1) and in the second display of its clause (b1), and in Propositions~35.4--35.9.
The data (S-4), that is, the quadruple $(\bar{m},e,R,N)$ of (FL) without the primitives, are assumed in Propositions~20.3, 26.6 and~48.6\,(b) and in the part ``In the flow'' of Proposition~30.3\,(f);
with $R=R^{*}(e)$ and $N\ge N_0(e)$ they are assumed in the clauses (a)--(f) of Proposition~35.1 and in the clause (b) of Proposition~49.
In all these statements the torus is one of the tori $T_k(R)$ just described.
Proposition~49 has the data $\bar{m}\in(0,1]$, $e>0$ with $e\cdot h^{\frac{1}{2}}\le e_1$, an integer $N\ge N_0(e)$ and an integer $\mathrm{kt}\ge1$,
and in parts of its clauses (d)--(f) a probe and a placement.
Since $R^{*}(e)\in\mathrm{Tor}(e)$ (Proposition~35.1) and $N_0(e)=N_{\min}\left(R^{*}(e),e\right)$, the quadruple $\left(\bar{m},e,R^{*}(e),N\right)$ satisfies the conditions of (S-4), as in the clauses (a)--(f) of Proposition~35.1.
The clause (b) of Proposition~49 is the case $\mathrm{kt}=1$;
it is a statement on the tori $T_j\left(R^{*}(e)\right)$, $0\le j\le k_{*}\left(R^{*}(e)\right)$.
The clauses (a) and (c)--(f) of Proposition~49 are made in lattice units on the cubic torus with $\mathrm{kt}\cdot n_*$ sites per side, where $n_*=R^{*}(e)/\eta$ (Section~I.6.4),
that is, on the torus of the physical side $\mathrm{kt}\cdot R^{*}(e)$, for every integer $\mathrm{kt}\ge1$, even or odd.
Exactly two of these sides lie in $\mathrm{Tor}(e)$: the side $R^{*}(e)$, for $\mathrm{kt}=1$, and the side $L\cdot R^{*}(e)=R_{\mathrm{top}}(e)=h\cdot L^{T^{*}(e)+1}$, for $\mathrm{kt}=L$.
Indeed, a side $\mathrm{kt}\cdot R^{*}(e)=h\cdot L^{T}$ of $\mathrm{Tor}(e)$ has $\mathrm{kt}=L^{T-T^{*}(e)}$, and $\mathrm{kt}\cdot R^{*}(e)\le R_{\mathrm{top}}(e)=L\cdot R^{*}(e)$ gives $\mathrm{kt}\le L$.
For both sides the cutoff satisfies the bound of (S-4), since $N\ge N_0(e)=N_{\min}\left(R^{*}(e),e\right)$ and $N_0(e)\ge N_{\min}\left(R_{\mathrm{top}}(e),e\right)$ \cite[Proposition~19.1(d)]{QEDI}.
The following statements do not assume the quadruple of (FL) and (S-4):
no side $R\in\mathrm{Tor}(e)$ together with a cutoff bound $N\ge N_{\min}(R,e)$ is among their data.
The first assertion of Proposition~35.1, which states that $R^{*}(e)\in\mathrm{Tor}(e)$ and that $N_0(e)$ is finite, has the data $\bar{m}$ and $e$.
Proposition~28.3\,(d) has the data $\bar{m}$ and $e$, and its bound is asserted for every integer $N\ge0$.
The clause (c1) of Proposition~35.3 has the data $\bar{m}$ and a probe.
Proposition~35.10 has the data formed by the constants of the lines (C0)--(C9) and (C6$'$), by $\bar{m}$ and by a probe $\beta$;
its conclusions are asserted for every $e$ with $0<e\le e_{\sharp}(\beta)$ and every $N\ge N_0(e)$.
The statements under the data (SF), (SF+), (OS-) or (OS), which stand among the Propositions numbered $j$ or $j.i$ with $2\le j\le25$ or $j=48$ and among Propositions~26.1--26.5,
are made for every side that the respective data admit, with no relation to $\mathrm{Tor}(e)$ (Section~I.1.8):
under (SF) for every $\ell\in\mathrm{Cl}$;
under (OS-) for every $\ell$ with $\ell\in\mathrm{Cl}$ and $\ell/L\in\mathrm{Cl}$;
under (SF+) and (OS) for every $\ell$ with $\ell\in\mathrm{Cl}$, $\ell/L\in\mathrm{Cl}$ and $\ell\ge M\cdot r_k$.
In Propositions~30.1, 30.2 and~31.1, and in Proposition~30.3 apart from the part ``In the flow'' of its clause (f), the sides are those named there.
The same holds for the clause (b1) of Proposition~35.3 apart from its second display:
it is made for every probe, every mesh $\eta\in(0,1]$ and every side $R>0$ such that $R/\eta$ is an integer $\ge2$,
that is, on the tori $\Lambda_{\eta,R}=\eta\mathbb{Z}^3/R\mathbb{Z}^3$ of Section~I.7.5, with no relation to $\mathrm{Tor}(e)$.
The auxiliary block tori of \cite[Definition~16.1]{QEDI}, mirror tori of the even side $6M$ with periodic spinors, serve only in the construction of the local inverses (Propositions~1.1--1.7, 2.1--2.8 and~3).
The flow stops at $k_{*}(R)\ge1$;
so the terminal integral is an integral over $4\ell_*^3$ pairs of generators.

\paragraph{The tuned pair.}
For $e\cdot h^{\frac{1}{2}}\le e_1$ and $N\ge N_{\min}\left(R_{\mathrm{top}}(e),e\right)$ let
\[
P_N(e)=\left(m_{0,N}(e),\varepsilon^{\mathrm{v}}_{0,N}(e)\right)
\]
be the real pair of Proposition~28.3.
For these $N$ the number $m_{0,N}(e)$ is its first entry, and
\[
m_{0,N}(e):=0\qquad\text{for }0\le N<N_{\min}\left(R_{\mathrm{top}}(e),e\right).
\]
One has $N_{\min}\left(R_{\mathrm{top}}(e),e\right)\le N_0(e)$ by \cite[Proposition~19.1(d)]{QEDI}.

\begin{named}{Remark}
In \cite[Theorem~21.3]{QEDI} one puts $m_{0,N}(e):=0$ for all $N<N_0(e)$.
Here the number $m_{0,N}(e)$ is put equal to $0$ only for $N<N_{\min}\left(R_{\mathrm{top}}(e),e\right)$:
for the finitely many $N$ with $N_{\min}\left(R_{\mathrm{top}}(e),e\right)\le N<N_0(e)$ the definition of \cite[Theorem~21.3]{QEDI} gives $0$,
whereas here $m_{0,N}(e)$ is the first entry of $P_N(e)$.
The following clauses of Propositions~43.1, 43.2 and Theorem~44 range over all $N\ge0$ and therefore involve these $N$:
(R0) of Proposition~43.1;
the clause of Proposition~43.2 by which $N\mapsto m_{0,N}(e)$ is a mass assignment;
the definition of the laws $P_{N,T}$, $N,T\ge0$, on which Theorem~44 is stated,
the first sentence of Theorem~44, and its clauses (B0), (B1) (whose property (P8) involves the cutoffs of a defining sequence) and (C).
Of $m_{0,N}(e)$ these clauses involve only that it is real, of modulus $<\frac{1}{8}$ and independent of $T$
(the clause of Proposition~43.2 only that it is real and independent of $T$),
and, in (R0) and in the first sentence of Theorem~44, that it is a function of $(N,e,\bar{m},\rho)$ and of the scheme, with no torus and no twist among its arguments.
That (B0), (B1) and (C) involve no more than this is seen as follows:
(B0) and (B1) hold for every mass assignment of Section~I.7.5,
and so does (C), except for the properties (B2)--(B5) contained in its assertion that (B) holds for every volume-first limit point;
these properties rest on the cutoffs $N\ge N_0(e)$ of a defining sequence only.
The properties of $m_{0,N}(e)$ listed above are given by Proposition~28.3\,(c), (d).
The first three are Proposition~28.3\,(d).
For $N\ge N_{\min}\left(R_{\mathrm{top}}(e),e\right)$ the fourth is Proposition~28.3\,(c), applied with $R=R_{\mathrm{top}}(e)$;
for $N<N_{\min}\left(R_{\mathrm{top}}(e),e\right)$ it holds because $m_{0,N}(e)=0$ by definition and the threshold $N_{\min}\left(R_{\mathrm{top}}(e),e\right)$ is formed from $e$, $h$, $L$, $M$ and $r$ only.
All other clauses in which a cutoff $N$ occurs begin at $N_0(e)$, among them all clauses that contain a bound in $e$:
(R2)--(R5) of Proposition~43.1, the clauses of Proposition~43.2 other than the one on the mass assignment, and Theorem~44\,(A).
\end{named}

The pair $P_N(e)$ depends on $(N,e,\bar{m},\rho)$ and on the scheme, not on the torus.
The scheme consists of the constants of the lines (C0)--(C9) of Section~I.8 (among them $\rho$ and the weights),
the basis $\left\{e^{0}_{i}\right\}$ of \cite[Section~5.2.5]{QEDI}
and the systems of primitives on $T_k(R_{\mathrm{top}})$, $k<k_{\mathrm{top}}$.

\subsubsection*{I.6.3. Periodisation inside $\mathrm{Tor}(e)$}

This is the periodisation of \cite[Definition~5.13]{QEDI}, where $\mathrm{Jr}$ is written $J$.
Let $\ell_1,\ell_2\in\mathrm{Cl}$ with $\ell_2=\mathrm{Jr}\cdot\ell_1$, $\mathrm{Jr}$ odd (in the scheme a power of $L$),
let $T_i:=\mathbb{L}_k/\ell_i\mathbb{Z}^3$ for $i=1,2$, and let $\mathrm{pj}\colon T_2\to T_1$ be the covering map.
Functions on the cells of $T_i$, elements of $V_k(T_i)$ and parameter families are identified with $\ell_i\mathbb{Z}^3$-periodic objects on $\mathbb{L}_k$;
spinors and functions of pairs on $T_i\times O_{\mathrm{cub}}$ are identified with the antiperiodic ones (Sections~I.1.5 and~I.1.7).
Since $\mathrm{Jr}$ is odd, every object on $T_1$ is an object on $T_2$;
this inclusion is $\mathrm{pj}^{*}$.
The label is a spectator: $\left(\mathrm{pj}^{*}f\right)(c,g)=f(c,g)$.
One has
\[
J\,\mathrm{pj}^{*}=\mathrm{pj}^{*}J,\qquad J^{*}\,\mathrm{pj}^{*}=\mathrm{pj}^{*}J^{*},
\]
and the pairing $\left\langle\mathrm{pj}^{*}f',\mathrm{pj}^{*}f\right\rangle^{\mathrm{ch}}$ on $T_2$ equals $\mathrm{Jr}^3\cdot\left\langle f',f\right\rangle^{\mathrm{ch}}$ on $T_1$.
The projection of a system of primitives is the one defined with Proposition~26.1.
In its kernels, that is, in the display (8.0) of Proposition~26.3\,(a), one has
\[
\boldsymbol{\chi}(w):=(-1)^{w_1+w_2+w_3}
\]
for spinors;
the product $\boldsymbol{\chi}(w)\cdot\widetilde{\mathrm{Op}}(\,\cdot\,)\left(t,t'+\ell_1\cdot w\right)$ is a function of $w$ modulo $\mathrm{Jr}$, whereas neither of the two factors is.

\emph{An even multiple of a side is never reached by a map of spinors}:
an antiperiodic function of period $\ell$ is periodic, not antiperiodic, with respect to an even multiple of $\ell$.
The even tori of Section~I.6.5 are related to the reference tori only through numbers,
namely partition functions and expectations of cell observables (Propositions~39, 40, 41, 42.1 and~42.2).

\subsubsection*{I.6.4. Cells of half-integer width}

Throughout this subsection the cutoff satisfies $N\ge N_0(e)$, and lengths are measured in lattice units.
Put
\begin{gather*}
n_*:=\frac{R^{*}(e)}{\eta}=h\cdot L^{T^{*}(e)+N},\qquad\mathrm{per}:=\frac{n_*}{L}=2m_{\mathrm{p}}+1,\qquad\mathrm{hc}:=\frac{\mathrm{per}}{2}=m_{\mathrm{p}}+\frac{1}{2},\\
\ell_c:=\mathrm{hc}\cdot\eta=\frac{R^{*}(e)}{2L},\qquad N_{*}:=8L^3 ;
\end{gather*}
the identity $\mathrm{per}=2m_{\mathrm{p}}+1$ defines the integer $m_{\mathrm{p}}\ge0$.
The number $n_*$ is odd; it is $\ell_0\left(R^{*}(e)\right)$, the number of sites per side.
The asterisk is part of the name: $n_*$ is not the number $n_k$ of Section~I.6.1 at $k=k_{*}$, which is $\ell_*^3$.
The number $\mathrm{per}$ is odd, $\mathrm{hc}$ is the width of a cell in lattice units, and $\ell_c$ is its physical width.

\paragraph{Cell structure.}
Let
\[
\mathcal{Q}_{*}:=\left(\mathfrak{F}_{\mathrm{c}};0,0,0;\mathrm{hc}\right),
\]
a cell structure of the coordinate frame in the sense of \cite[Definition~2.18]{QEDII}, whose width lies in $\frac{1}{2}\mathbb{Z}$ (see (I.h1) below).
Its walls are the planes $\left\{y_\mu=j\cdot\mathrm{hc}\right\}$, $j\in\mathbb{Z}$.
For even $j$ they are site planes ($y_\mu\in\mathrm{per}\,\mathbb{Z}$), for odd $j$ link planes ($y_\mu\in\mathrm{per}\,\mathbb{Z}+m_{\mathrm{p}}+\frac{1}{2}$);
the two kinds alternate.
All walls are mirrors (Section~I.4.8).
The open cells are
\[
Q^{\circ}_{\mathfrak{c}}:=\left\{y\in\mathbb{R}^3:\ \mathfrak{c}_\mu\cdot\mathrm{hc}<y_\mu<\left(\mathfrak{c}_\mu+1\right)\cdot\mathrm{hc},\ \mu=1,2,3\right\},\qquad\mathfrak{c}\in\mathbb{Z}^3 .
\]
The index of a cell is always denoted $\mathfrak{c}$; in \cite{QEDI} it is written $c$, in \cite{QEDII} it is written $\kappa$.
The photon sites of the closed cell with the index $\mathfrak{c}$ form the set $\mathrm{cc}_{\mathfrak{c}}$ (the set $\mathfrak{c}_c$ of \cite[Section~5.2.9]{QEDI}),
\begin{gather*}
\mathrm{cc}_{\mathfrak{c}}:=\prod_{\mu}I_{\mathfrak{c}_\mu},\\
I_{2i}:=\left\{i\cdot\mathrm{per},\dots,i\cdot\mathrm{per}+m_{\mathrm{p}}\right\},\qquad
I_{2i+1}:=\left\{i\cdot\mathrm{per}+m_{\mathrm{p}}+1,\dots,(i+1)\cdot\mathrm{per}\right\}.
\end{gather*}
Let $W(\mathcal{Q}_{*})$ be the group generated by the reflections in the walls.
It acts simply transitively on the cells:
in each coordinate the reflections in the points $j\cdot\mathrm{hc}$, $j\in\mathbb{Z}$, generate an infinite dihedral group, which acts simply transitively on the open intervals between consecutive points,
and reflections in walls of different directions commute.
Let $\mathsf{g}_{\mathfrak{c}}$ be the element of $W(\mathcal{Q}_{*})$ with $\mathsf{g}_{\mathfrak{c}}Q^{\circ}_{0}=Q^{\circ}_{\mathfrak{c}}$.
In the coordinate $\mu$ it acts by
\[
y_\mu\mapsto i\cdot\mathrm{per}+y_\mu\quad\text{if }\mathfrak{c}_\mu=2i,\qquad
y_\mu\mapsto(i+1)\cdot\mathrm{per}-y_\mu\quad\text{if }\mathfrak{c}_\mu=2i+1 .
\]
The element $\mathsf{g}_{\mathfrak{c}}$ is even if and only if $\sum_\mu\mathfrak{c}_\mu$ is even;
for $\mathfrak{c}\in2\mathbb{Z}^3$ it is the translation by $\mathrm{hc}\cdot\mathfrak{c}$, a vector of $\mathrm{per}\,\mathbb{Z}^3$.

\paragraph{Tori that carry the cells.}
The torus $C(n_1,n_2,n_3)$ carries $\mathcal{Q}_{*}$ if and only if $\mathrm{per}$ divides every $n_\mu$,
that is, $n_\mu\in2\,\mathrm{hc}\,\mathbb{Z}$, so that every side consists of an even number $n_\mu/\mathrm{hc}$ of cells.
Sides of either parity occur.
For an integer $\mathrm{kt}\ge1$, the torus $\mathrm{kt}$ is the cubic torus with $\mathrm{kt}\cdot n_*$ sites per side; in \cite[Section~21.1.5]{QEDI} it is called the torus $k$.
It has $\left(2\cdot\mathrm{kt}\cdot L\right)^3$ cells.
The torus $1$ has $N_{*}$ cells, and their set is denoted $\mathrm{Cel}$;
the second reference torus is the torus $L$.

\paragraph{Cell algebra.}
For $\mathfrak{c}\in\mathbb{Z}^3$ the cell algebra $\mathcal{A}\left(Q^{\circ}_{\mathfrak{c}}\right)$ is the algebra $\mathcal{A}(U)$ of \cite[Definition~2.12]{QEDII} for the open set $U=Q^{\circ}_{\mathfrak{c}}$,
written in the coordinates $(c,g)$ by Proposition~46\,(h).
Its elements read in boxes contained in the open cell;
nothing that meets a wall is read.
In \cite[Definitions~2.9, 2.12 and~2.13]{QEDII} a fermion site is a point of the set $W_{\mathrm{af}}v_0$, denoted $\mathbb{X}$ there;
here it is written $(c,g)$ by means of $\mathrm{geo}$.
The reflections $\Theta_m$, the action $G\mapsto G^{\mathrm{gw}}$, the sets of the fermion sites in a box, the classes $\mathrm{Ad}(\mathrm{Bx})$ (written $\mathcal{G}(B)$ in \cite[Definition~2.12]{QEDII}) and $\mathcal{A}(U)$ and the functional $Z_{\mathbb{T}}[G]$, also for Grassmann valued $G$,
involve of the fermion lattice only the set of the sites, the map $\pi$, the action of $W_{\mathrm{af}}$, the operator and the Berezin integral,
and $\mathrm{geo}$ carries these over (Proposition~46\,(a), (b), (e)).
The cell algebras take the place of the objects $P^o(\mathfrak{c}_c)$ and $\mathfrak{A}(\mathfrak{c}_c)$ attached to closed cells in \cite[Sections~5.2.9 and~21.1.5]{QEDI}.
In the proof of Proposition~42.2, where Propositions~39 and~40 are applied,
the cell contents are the functions $O_m=e^{-m}\cdot H_m(X)-c_m$, $m\in\{2,4\}$, of Proposition~41\,(d), their mirror copies and, in the remaining cells, the constant~$1$;
all of them are functions of the curvature.

\paragraph{Mirror copies.}
For $O\in\mathcal{A}\left(Q^{\circ}_{0}\right)$ put $O^{[\mathfrak{c}]}:=O^{\mathsf{g}_{\mathfrak{c}}}$ (Section~I.4.8).
For a function $O$ of two-cochains $u$ (in the applications, of the curvature) this reads
\[
O^{[\mathfrak{c}]}(u)=O\left(\mathsf{g}_{\mathfrak{c}}^{*}u\right)\quad\text{if }\sum_\mu\mathfrak{c}_\mu\text{ is even},\qquad
O^{[\mathfrak{c}]}(u)=\overline{O\left(\mathsf{g}_{\mathfrak{c}}^{*}u\right)}\quad\text{if }\sum_\mu\mathfrak{c}_\mu\text{ is odd},
\]
as in \cite[Section~21.1.5]{QEDI} and in \cite[Definition~2.19]{QEDII}, the latter with the action of the group $W$ on even integrands defined in \cite[Section~2]{QEDII}.
Further,
\[
O^{\sharp}:=\prod_{\mathfrak{c}}O^{[\mathfrak{c}]},
\]
the product extending over all cells of the torus under consideration;
it is the fold of $O$ on this torus in the sense of \cite[Definition~2.19]{QEDII}.

\paragraph{Fermion sites and cells.}
The following assignment serves only for reading \cite[Sections~5 and~6]{QEDII}.
A fermion site $x$ belongs to the open cell that contains $\mathrm{geo}(x)$;
for a photon site $c$ on a wall $\left\{y_\mu=j\cdot\mathrm{hc}\right\}$ with $j$ even, the $48$ sites $(c,g)$ over $c$ are divided according to the sign of the row $\mu$ of $g$ (Section~I.4.8):
this sign decides on which side of the wall the site $(c,g)$ lies,
and if $c$ lies on walls of several directions $\mu$, this holds for each of them.

\paragraph{Widths in $\frac{1}{2}\mathbb{Z}$ in \cite{QEDII}.}
Widths in $\frac{1}{2}\mathbb{Z}$ occur for the coordinate frame only.
The five points (I.h1)--(I.h5) below extend statements of \cite{QEDII} to such widths.
For \cite[Definitions~2.18 and~6.13]{QEDII} on the tori that carry the cell structure, and for \cite[Lemma~6.5(d)]{QEDII}, \cite[Theorem~6.16(a)]{QEDII}, \cite[Proposition~6.23(a)]{QEDII} and \cite[Proposition~12.10(f)]{QEDII},
the extension is justified by Proposition~39\,(a)--(f).
The other extensions are justified in Chapter~39, Sections~5, 7.2 and~7.3;
among them are \cite[Definition~6.13, Lemma~6.14 and Proposition~6.15]{QEDII} on $C(n)$,
the number $|\mathbb{T}|/|Q|$ of the cells of a torus $\mathbb{T}$ which is cut by the cell structure but does not carry it,
the statements on coarser cell structures, the extension to the images of the frame $\mathfrak{F}_{\mathrm{c}}$ under $O_{\mathrm{cub}}$,
\cite[Lemma~6.17]{QEDII}, which is Lemma~7.1 there, and \cite[Proposition~6.23(b)]{QEDII}, which is Proposition~7.2\,(b$'$) there.
In (I.h1)--(I.h5) the notation is that of \cite{QEDII}; in particular $h$ denotes the width of a cell structure.
\begin{itemize}
\item[(I.h1)]
In \cite[Definition~2.18]{QEDII} the width $h$ of a cell structure may be, besides a positive integer,
for the frame $\mathfrak{F}_{\mathrm{c}}$ and its images under $O_{\mathrm{cub}}$, any number $h\in\frac{1}{2}\mathbb{Z}$ with $h>0$.
Then the volume $|Q|=h^3$ of an open cell of the coordinate frame is a number and no longer a number of photon sites,
while the number $|\mathbb{T}|/|Q|$ of the cells of a torus $\mathbb{T}$ remains an integer.
Widths in $\frac{1}{2}\mathbb{Z}$ are admitted in this sense in \cite[Definitions~2.18, 2.19 and~6.7]{QEDII},
in \cite[Lemma~6.5(d)]{QEDII},
in \cite[Definition~6.13]{QEDII},
on $C(n)$ in \cite[Lemma~6.14 and Proposition~6.15]{QEDII},
and in \cite[Theorem~6.16(a)]{QEDII}, \cite[Lemma~6.17]{QEDII}, \cite[Proposition~6.23]{QEDII} and \cite[Proposition~12.10(f)]{QEDII}.
In \cite[Lemma~2.33(b)]{QEDII}, \cite[Theorem~6.16(c)]{QEDII} and \cite[Sections~8, 9 and~13]{QEDII}
(in particular in \cite[Proposition~13.9(b)]{QEDII}, whose proof works with planes of the same kind) the width remains an integer.
\item[(I.h2)]
In \cite[Definition~6.13]{QEDII} the positions $o_a+j\cdot h$ lie in $\operatorname{Adm}(\ell^a)$
because $o_a$ and $h$ lie in $\operatorname{Adm}(\ell^a)=\frac{1}{2}\mathbb{Z}$ for the frame $\mathfrak{F}_{\mathrm{c}}$ and its images under $O_{\mathrm{cub}}$, and because $h$ is an integer for the other frames.
The same reason gives the relation $o_a+j\cdot h\in\operatorname{Adm}(\ell^a)$ in \cite[Definition~2.18]{QEDII},
and the inclusion of the planes $o_a+j\cdot h$, $j\in\mathbb{Z}$, in $\operatorname{Adm}(\ell^a)$ in the paragraph on the trace in the proof of \cite[Proposition~6.23(a)]{QEDII}.
\item[(I.h3)]
In the paragraph following \cite[Definition~6.22]{QEDII} the numbers $p_a:=h\cdot m_a$ are formed from the even entries $m_a$ of a shape, and it is shown that the vectors $\left(p_a/|n_a|^2\right)n_a$ are integral.
They remain integral for widths in $\frac{1}{2}\mathbb{Z}$:
$h\cdot m_a$ is an integer, because $m_a$ is even and $2h$ is an integer, and $|n_a|^2=1$ for the coordinate frame.
For such widths the number $p_a$ need not be even.
\item[(I.h4)]
For widths in $\frac{1}{2}\mathbb{Z}$ the sides of the tori formed with the shapes need not be even:
for the width $\mathrm{hc}=m_{\mathrm{p}}+\frac{1}{2}$ odd sides occur, and for the width $\frac{1}{2}$ and an even integer $n$ one has
\[
\left(\frac{n}{2},2n,n\right)=\frac{1}{2}\cdot\left(n,4n,2n\right)
\]
with even entries.
\item[(I.h5)]
In \cite[Proposition~12.10(f)]{QEDII} the width $h$ may be any number $h\in\frac{1}{2}\mathbb{Z}$ with $h>0$ in the place of a positive integer;
the hypothesis $n\in2h\mathbb{Z}$ remains.
\end{itemize}
Let $e\cdot h^{\frac{1}{2}}\le e_1$ and $T\ge T^{*}(e)$.
On the side $n=8\cdot h\cdot L^{T+N}$ of Section~I.6.5 the cells of \cite[Section~21.1.5]{QEDI} exist, because $\mathrm{per}$ divides $n$:
one has $n/\mathrm{per}=8\cdot L^{T-T^{*}(e)+1}$ (Section~I.6.5).
For $0\le T\le T^{*}(e)-2$ the number $\mathrm{per}$ does not divide $n$, since $L>8$, and $C(n,n,n)$ does not carry $\mathcal{Q}_{*}$.
The diagonal frames are not affected: there the width remains an integer.

\subsubsection*{I.6.5. The even tori of the theorem}

For integers $N,T\ge0$ put
\[
n:=8\cdot h\cdot L^{T+N},\qquad R:=n\cdot\eta=8\cdot h\cdot L^T
\]
(\cite[Definition~2.36]{QEDII}, where $h$ is written out as $L^{\lfloor\log_L\bar{m}^{-1}\rfloor}$).
Let $e\cdot h^{\frac{1}{2}}\le e_1$, $N\ge N_0(e)$ and $T\ge T^{*}(e)$.
Then
\[
n=\mathrm{kt}\cdot n_*\quad\text{with}\quad\mathrm{kt}=8\cdot L^{T-T^{*}(e)},\qquad\frac{n}{\mathrm{per}}=\mathrm{kt}\cdot L ;
\]
so $C(n,n,n)$ carries $\mathcal{Q}_{*}$, with $2\cdot\mathrm{kt}\cdot L$ cells per side.
Moreover $\mathrm{kt}\ge L$ if and only if $T\ge T^{*}(e)+1$, since $L>8$ by (C3).
The transfer (Propositions~40, 42.1 and~42.2) covers the side $n_*$ and the sides that are multiples of $\mathrm{per}=n_*/L$ and at least $L\cdot n_*$;
hence the statements at finite cutoff hold for $T\ge T^{*}(e)+1$.
Note that
\[
\frac{R^{*}(e)}{L}=\mathrm{per}\cdot\eta=2\ell_c .
\]
All sides are at least $4$, so that $n\mathbb{Z}^3$ is a period lattice in the sense of \cite[Definition~2.6]{QEDII}.
The sequence of the sides $8h\cdot L^{T+N}$, $T\to\infty$, satisfies the condition (Div) of \cite[Definition~2.21]{QEDII}.
For the cubic tori of these sides the theorem on the nine families of mirrors is \cite[Theorem~1.2]{QEDII}, which holds for every sequence of integers $n_i\ge4$ with $n_i\to\infty$.
By (Div) and \cite[Lemma~11.14(c)]{QEDII}, the sequence of the sides is a companion of itself in the sense of \cite[Definition~11.12]{QEDII};
hence in the proof of \cite[Theorem~1.2]{QEDII}, that is, in \cite[Lemma~11.18]{QEDII}, the companion can be taken equal to the sequence of the sides,
and the reduced sequence of \cite[Definition~11.13]{QEDII} does not occur.

\subsubsection*{I.6.6. Boundary condition, twist, external connection}

\paragraph{Boundary condition.}
On all tori the spinors are antiperiodic in the three directions:
\[
C(n)=\left(\sum_{\mu}n_{\mu}\mathbb{Z}\,e_{\mu};\ \chi=-1\text{ on the generators}\right)
\]
in the sense of \cite[Definition~2.11]{QEDII}.
The numbers $Z_n[\Phi]$ and the laws do not depend on the sign character \cite[Proposition~3.16(h)]{QEDII};
the flow, the periodisation of Section~I.6.3 and the closing sign of the trace formulas are formed with $\chi=-1$.

\paragraph{Twist.}
Let $\theta\in\mathbb{R}^3$.
In lattice units $H_{\theta}$ is the constant real one-cochain with
\[
H_{\theta}\left(c,c+e_\mu\right):=\frac{\theta_\mu}{q\cdot n_\mu}.
\]
At level $k$ on a torus of the flow, $c^{(k)}_{\theta}$ is the constant connection with the components $\theta_\mu/\left(e_k\cdot\ell_k\right)$; thus
\[
c^{(0)}_{\theta}=H_{\theta},\qquad L^{-\frac{1}{2}}\cdot c^{(k+1)}_{\theta}=c^{(k)}_{\theta},\qquad c_{\theta}:=c^{(k_{*})}_{\theta}.
\]
A constant connection enters $D_A$ only through the phases of the bonds of type $3$:
\[
U_{A+H_{\theta}}\left(x,x^{(3)}\right)=\exp\left(i\,\frac{\theta\cdot(g\,e_1)}{n}\right)\cdot U_A\left(x,x^{(3)}\right).
\]
It amounts to a change of the character \cite[Lemma~4.24(a)]{QEDII}:
\begin{equation}\tag{I.6.1}
\begin{aligned}
&D_A\cdot\exp\left(i\,\frac{\theta\cdot c}{n}\right)=\exp\left(i\,\frac{\theta\cdot c}{n}\right)\cdot D_{A+H_{\theta}},\\
&f\mapsto\exp\left(i\,\frac{\theta\cdot c}{n}\right)\cdot f\ \text{ maps the antiperiodic spinors onto the spinors with}\\
&\qquad f\left(c+n\,e_\mu,g\right)=-\exp\left(i\theta_\mu\right)\cdot f(c,g).
\end{aligned}
\end{equation}
The identity (I.6.1) and the display before it are written for the cubic torus.
On $C(n_1,n_2,n_3)$ the quotient $\theta_\mu/n_\mu$ stands in the place of $\theta_\mu/n$:
$\theta\cdot(g\,e_1)/n$ and $\theta\cdot c/n$ are read $\sum_\mu\theta_\mu\,(g\,e_1)_\mu/n_\mu$ and $\sum_\mu\theta_\mu c_\mu/n_\mu$, and $c+n\,e_\mu$ is read $c+n_\mu e_\mu$.
By (I.6.1) the twist enters the spinors and the determinants through the numbers $\exp\left(i\theta_\mu\right)$ only;
everything is $2\pi$-periodic in each $\theta_\mu$.
The twist connections $c^{(k)}_{\theta}$, $k\ge1$ (in the first slot, as the harmonic part of a block field), are defined on the tori of $\mathrm{Tor}(e)$ only.
The cochain $H_{\theta}$ and the mean over $\theta$ are defined on every torus of Sections~I.7.3 and~I.7.4, of either parity;
they are formed on rectangular tori in Propositions~38.1--38.3 and for every $n\ge1$ in Propositions~30.1--30.3,
and in the proof of Theorem~44 on $R^{*}(e)$ and $L\cdot R^{*}(e)$ only (Section~I.7.3).
For $\theta\in\mathbb{C}^3$ the connections $H_{\theta}=c^{(0)}_{\theta}$ and $c_{\theta}$ are defined by the same formulas;
complex $\theta$ is admitted in Proposition~38.1\,(e) for $H_{\theta}$ and in Proposition~34.4 for $c_{\theta}$.

\paragraph{Two slots.}
As in \cite[Definition~5.12]{QEDI}, the twist enters as the harmonic part of the block field, in the first slot;
the terminal block field is $A_{*}+c_{\theta}$.
An external connection $a$, a complex one-cochain in lattice units, enters in the second slot:
\[
\mathscr{C}=a^{[k]},\qquad a^{[k]}(b):=L^{\frac{k}{2}}\cdot a\left(L^k\cdot b\right);
\]
it is a parameter and is not integrated.
The twist is the vector $\theta$, and the twist connection is $c^{(k)}_{\theta}$.

\subsection*{I.7. Laws}

\subsubsection*{I.7.1. Configurations, curvature, flux}

All objects of this subsection are taken in lattice units;
the configurations are those of \cite[Definition~10.1]{QEDI} and \cite[Definition~2.6]{QEDII}.
Let $n=(n_1,n_2,n_3)$ with $n_\mu\ge4$, and $|n|:=n_1n_2n_3$;
a single letter $n$ also denotes the cubic torus $(n,n,n)$.
Let $C^{j}(n)$ be the space of the real $j$-cochains of $\mathbb{Z}^3/\sum_{\mu}n_{\mu}\mathbb{Z}\,e_{\mu}$, with the counting pairing $\langle\,\cdot\,,\,\cdot\,\rangle$ and the norm $\|\cdot\|$.
One has
\[
C^{1}(n)={V_{\mathrm{g}}}\oplus{V_{\mathrm{h}}}\oplus{V_{\mathrm{t}}},
\]
the spaces of the exact cochains, of the constant cochains and the orthogonal complement of both, and $\sigma$ denotes the Lebesgue measure of ${V_{\mathrm{t}}}$.
A configuration is a pair $(A,\mathfrak{z})$ with $A\in[0,r_q)^{\mathrm{links}}$, $r_q=2\pi/q$, and $\mathfrak{z}\in C^{2}(n;\mathbb{Z})$ with $d\mathfrak{z}=0$.
Put
\[
F^{\mathrm{u}}:=dA+r_q\cdot\mathfrak{z},\qquad F:=\eta^{-\frac{3}{2}}\cdot F^{\mathrm{u}} ;
\]
$F^{\mathrm{u}}$ is the curvature in lattice units, a closed two-cochain, and $F$ is the physical curvature.
In \cite[Definitions~2.6--2.34]{QEDII} the letter $F$ stands for $F^{\mathrm{u}}$.

\paragraph{Flux.}
For $\mu<\nu$ let $\Pi_{\mu\nu}$ be the set of the $n_\mu n_\nu$ plaquettes $(c;\mu\nu)$ in one plane, and put
\[
w_{\mu\nu}:=\sum_{p\in\Pi_{\mu\nu}}\mathfrak{z}_p\in\mathbb{Z},\qquad w:=\left(w_{23},w_{13},w_{12}\right)\in\mathbb{Z}^3 .
\]
The components are labelled by the pairs $\mu<\nu$, as in \cite[Section~21.1.4]{QEDI} and in \cite[Definition~13.7]{QEDII};
in \cite[Section~21.7]{QEDI} the directions are labelled $0,1,2$ and the flux is written $(w_{12},w_{02},w_{01})$.
Let $h_{w}$ be the two-cochain with the constant components $r_q\cdot w_{\mu\nu}/\left(n_\mu n_\nu\right)$.
For a closed real two-cochain $u\in C^{2}(n)$ put
\[
\mathrm{fl}_{\mu\nu}(u):=r_q^{-1}\cdot\sum_{p\in\Pi_{\mu\nu}}u_p,\qquad
\mathrm{fl}(u):=\left(\mathrm{fl}_{23}(u),\mathrm{fl}_{13}(u),\mathrm{fl}_{12}(u)\right);
\]
the real number $\mathrm{fl}_{\mu\nu}(u)$ does not depend on the plane.
The flux of $u$ is said to be $w$ if $\mathrm{fl}(u)=w$.
The vector $\mathrm{fl}\left(dA+r_q\cdot\mathfrak{z}\right)$ is the flux $w$ of $\mathfrak{z}$;
so the flux is a function of $F^{\mathrm{u}}$, and the indicator $\mathbf{1}[\,\text{flux}=w\,]$ of Section~I.7.3 is the function $\Phi_w$ of \cite[Definition~13.7]{QEDII}, where the flux is written $k$.

\subsubsection*{I.7.2. The all-flux law (af): the law of the theorem}

Let $\boldsymbol{\mu}$ be a real lattice mass, and let $G$ be a Borel function of $F^{\mathrm{u}}$ with
\[
|G|\le\mathrm{cf}\cdot\exp\left(\mathrm{cf}\cdot\left\|F^{\mathrm{u}}\right\|\right)
\]
for some number $\mathrm{cf}$;
as in \cite[Section~21.1.4]{QEDI}, the class of these functions is denoted $\mathfrak{F}(n)$.
Put
\begin{equation}\tag{I.7.1}
Z_n[G]:=\sum_{\mathfrak{z}:\ d\mathfrak{z}=0}\ \int_{[0,r_q)^{\mathrm{links}}}\mathrm{d}A\,\exp\left(-\frac{1}{2}\left\|F^{\mathrm{u}}\right\|^2\right)\cdot\det\left(D_A+\boldsymbol{\mu}\right)\cdot G\left(F^{\mathrm{u}}\right),\qquad
E^{\mathrm{af}}_n:=\frac{Z_n[\,\cdot\,]}{Z_n[1]}.
\end{equation}
One has $0<Z_n[1]<\infty$ for every real $\boldsymbol{\mu}$ \cite[Theorem~4.26]{QEDII}, and $\det\left(D_A+\boldsymbol{\mu}\right)\ge0$ by (I.O3);
so $E^{\mathrm{af}}_n$ is the expectation of a probability law of $F^{\mathrm{u}}$.
For the class $\mathcal{A}(\mathbb{R}^3)$ of \cite[Definition~2.12]{QEDII} and the class $\mathrm{Ig}(\mathbb{T})$, written $\mathcal{I}(\mathbb{T})$ in \cite[Definition~5.11]{QEDII}, the functional is the functional $Z_{\mathbb{T}}$ of \cite[Definition~2.13]{QEDII} with $\mathbb{T}=C(n)$ (Proposition~46\,(h)).

The corresponding notation of \cite{QEDI} and \cite{QEDII} is the following.

\medskip
\begin{tabular}{p{2.7cm}p{6.6cm}p{4.6cm}}
here & in \cite{QEDI} & in \cite{QEDII}\\
\hline
$Z_n[G]$ &
$Z_n[G]$ \cite[Sections~5.2.1 and~10]{QEDI};
$\boldsymbol{\omega}_n(G)$ \cite[Section~21.1.4]{QEDI};
$\omega^{\mathrm{af}}_n$ \cite[Definition~11.1]{QEDI} &
$Z_{\mathbb{T}}[G]$, $\mathbb{T}=C(n)$\\
$E^{\mathrm{af}}_n$ &
$E_n$ \cite[Sections~5.2.1 and~10.1.3]{QEDI};
$E^{\mathrm{af}}_n$ \cite[Section~21.1.4]{QEDI};
$E^{\mathrm{af}}_k$ on the torus $k$ \cite[Section~21.1.5]{QEDI} &
$\mathbb{E}_{\mathbb{T}}$; the law $\mu_{\mathbb{T}}$ is the law of the periodic extension of $F^{\mathrm{u}}$\\
$Z^{(w)}_n[G]$, $P_n(w)$ &
$\boldsymbol{\omega}^{(w)}_n$ \cite[Section~21.1.4]{QEDI};
$Z_{n,w}$ and $P_n(w)$ \cite[Section~10.1.4]{QEDI} &
$Z^{(k)}_{\mathbb{T}}[G]$ and the weight $\mathbb{E}_{\mathbb{T}}[\Phi_k]$ of the sector, the flux being written $k$ \cite[Definition~13.7]{QEDII}\\
$E^{\mathrm{av}}_n$ &
$E^{\mathrm{av}}_n$ \cite[Section~21.1.4]{QEDI};
$\overline{\mathsf{E}}$ \cite[Section~6.1.5]{QEDI};
$\bar{\mu}$ \cite[Section~5.2.1]{QEDI};
the law $\bar{\mathsf{P}}$ of \cite[Section~7.1.2]{QEDI} is the law of $F=\eta^{-\frac{3}{2}}F^{\mathrm{u}}$ when $F^{\mathrm{u}}$ has the law $E^{\mathrm{av}}_n$ &
$\mu^{\mathrm{av}}_{\mathbb{T}}$ \cite[Definition~13.7]{QEDII}\\
$\exp\left(-\frac{1}{2}\langle h_{w},h_{w}\rangle\right)$ &
$\mathfrak{q}_w(n)$ \cite[Section~21.1.4]{QEDI} &
the same number, written with the norm of $h_{w}$, which is denoted $\mathfrak{h}_k$ there\\
\end{tabular}
\medskip

\subsubsection*{I.7.3. The zero-flux sector, the harmonic average, the measures of the flow}

Put
\begin{gather*}
Z^{(w)}_n[G]:=Z_n\left[\mathbf{1}[\,\text{flux}=w\,]\cdot G\right],\qquad
P_n(w):=\frac{Z^{(w)}_n[1]}{Z_n[1]},\\
E^{\mathrm{av}}_n:=\frac{Z^{(0)}_n[\,\cdot\,]}{Z^{(0)}_n[1]}=E^{\mathrm{af}}_n\left(\,\cdot\mid\text{flux}=0\right).
\end{gather*}
The functional $E^{\mathrm{av}}_n$ is the zero-flux sector; its law is called (av).
The flux is a function of the curvature.
Indeed, the sum of $dA$ over $\Pi_{\mu\nu}$ vanishes, so that $\sum_{p\in\Pi_{\mu\nu}}F^{\mathrm{u}}_p=r_q\cdot w_{\mu\nu}$;
and for a closed $u\in C^{2}(n)$ the sums $\sum_{p\in\Pi_{\mu\nu}}u_p$ do not depend on the plane,
since the sum of $du=0$ over the cubes based at the sites of one plane is, up to its sign, the difference of the sums of $u$ over two neighbouring planes.
Hence
\[
\mathbf{1}[\,\text{flux}=w\,]=\prod_{\mu<\nu}\mathbf{1}\left[\,\sum_{p\in\Pi_{\mu\nu}}F^{\mathrm{u}}_p=r_q\cdot w_{\mu\nu}\right]
\]
is a bounded Borel function of $F^{\mathrm{u}}$, and $\mathbf{1}[\,\text{flux}=w\,]\cdot G$ lies in the class of Section~I.7.2 together with $G$
(\cite[Definition~13.7]{QEDII}; \cite[(10.2)]{QEDI}).

Let now $n\ge4$ be the number of sites per side of a cubic torus, of either parity (in the flow $n$ is odd);
let $\left(m_{\mathrm{in}},\varepsilon^{\mathrm{v}}_{\mathrm{in}}\right)$ be a real pair, $\theta\in\mathbb{R}^3$ and $a\in C^{1}(n)\otimes\mathbb{C}$.
Let $F_N:=dC^{1}(n)$, let $\gamma_N$ be its standard Gaussian measure,
and for $F^{\mathrm{u}}\in F_N$ let $A(F^{\mathrm{u}})$ be the element of ${V_{\mathrm{t}}}$ with $dA(F^{\mathrm{u}})=F^{\mathrm{u}}$.
Put
\begin{equation}\tag{I.7.2}
\mathrm{dens}^{\theta,a}_N\left(F^{\mathrm{u}}\right):=\exp\left(-\varepsilon^{\mathrm{v}}_{\mathrm{in}}\cdot n^3\right)\cdot\frac{\det\left(D_{A(F^{\mathrm{u}})+H_{\theta}+a}+\bar{m}\cdot\eta+m_{\mathrm{in}}\right)}{\det\left(D_0+\bar{m}\cdot\eta\right)},
\end{equation}
and
\begin{gather*}
\nu^{\theta,a}:=\mathrm{dens}^{\theta,a}_N\cdot\gamma_N,\qquad\nu^{\theta}:=\nu^{\theta,0},\qquad
\bar{\nu}^{a}:=\int_{[-\pi,\pi]^3}\frac{\mathrm{d}\theta}{(2\pi)^3}\,\nu^{\theta,a},\qquad\bar{\nu}:=\bar{\nu}^{0},\qquad\bar{\mu}:=\frac{\bar{\nu}}{\bar{\nu}(1)},\\
\mathrm{zr}_n(q,\mathbf{m}):=\int_{[-\pi,\pi]^3}\frac{\mathrm{d}\theta}{(2\pi)^3}\int\mathrm{d}\gamma_N\left(F^{\mathrm{u}}\right)\,\det\left(D_{A(F^{\mathrm{u}})+H_{\theta}}+\mathbf{m}\right),\\
\mathrm{zr}^{\mathrm{free}}_n:=\int_{[-\pi,\pi]^3}\frac{\mathrm{d}\theta}{(2\pi)^3}\,\det\left(D_{H_{\theta}}+\bar{m}\cdot\eta\right).
\end{gather*}
The measure $\bar{\mu}$ is the harmonic average.
In \cite[Section~5.2.1]{QEDI} the objects $\mathrm{dens}^{\theta,a}_N$, $\mathrm{zr}_n(q,\mathbf{m})$ and $\mathrm{zr}^{\mathrm{free}}_n$ are denoted $\rho^{\theta,a}_N$, $\rho_n(q,\mathbf{m})$ and $\rho^{\mathrm{free}}_n$.
The definition of $\mathrm{zr}^{\mathrm{free}}_n$ contains no gauge field;
it makes sense for every integer $n\ge1$, $H_{\theta}$ being read on $\mathbb{Z}^3$ as in the setting of Proposition~46,
and Proposition~30.2\,(b) concerns it in this generality.
By the definitions,
\[
\mathrm{zr}_n\left(q,\bar{m}\cdot\eta+m_{\mathrm{in}}\right)=\exp\left(\varepsilon^{\mathrm{v}}_{\mathrm{in}}\cdot n^3\right)\cdot\det\left(D_0+\bar{m}\cdot\eta\right)\cdot\bar{\nu}(1).
\]

\paragraph{Identification.}
Let $\Delta_0:=d^{*}d$ on $C^{0}(n)$, where $d^{*}$ is the adjoint of $d$ for the counting pairing;
it is the Laplacian on the functions of the sites.
Let $\operatorname{det}'\Delta_0$ be the product of its non-zero eigenvalues, and
\[
\mathrm{Nc}_n:=r_q^{|n|+2}\cdot\left(\operatorname{det}'\Delta_0\right)^{\frac{1}{2}}.
\]
This operator $\Delta_0$ acts on functions of the sites;
it is neither the photon operator $\Delta_k$ of Section~I.5.7 at $k=0$, which acts on bond functions, nor the unit block $\Delta_y$ at $y=0$.
By Proposition~38.2\,(b),
\[
Z^{(0)}_n[\Phi]=\mathrm{Nc}_n\cdot\int_{[-\pi,\pi]^3}\frac{\mathrm{d}\theta}{(2\pi)^3}\int_{{V_{\mathrm{t}}}}\mathrm{d}\sigma(A)\,\Phi(dA)\cdot\exp\left(-\frac{1}{2}\|dA\|^2\right)\cdot\det\left(D_{A+H_{\theta}}+\boldsymbol{\mu}\right);
\]
consequently, for $\boldsymbol{\mu}=\bar{m}\cdot\eta+m_{\mathrm{in}}$, the measure $\bar{\mu}$ is the law of $F^{\mathrm{u}}$ under $E^{\mathrm{av}}_n$.

\paragraph{The role of the averaged law.}
The law of the main theorem is (af); the averaged law is an auxiliary law, and in the proof of Theorem~44 it is taken on the two reference tori only.
The statements on the averaged law hold on the following tori:
Propositions~31.2, 32.2 and the displays (10.1), (10.2) of Proposition~32.3 on every torus of $\mathrm{Tor}(e)$;
Propositions~35.7--35.9 on $R^{*}(e)$;
Propositions~38.2, 38.3 and Proposition~50\,(e)--(g$''$) on every torus with $n_\mu\ge4$;
Propositions~31.1 and~42.1 on every cubic torus with $n\ge4$;
Proposition~49\,(e), (f) on the tori $n=\mathrm{kt}\cdot n_*$, $\mathrm{kt}\ge1$.
Apart from the identification above and the shift identity of Section~I.7.4, the averaged law has the following properties:
$\nu^{\theta}\ge0$ by (I.O3);
$\bar{\nu}(1)>0$, since every flux term is positive \cite[Theorem~4.26(b)]{QEDII}, so that $Z^{(0)}_n[1]>0$ (\cite[Definition~13.7(b)]{QEDII}; \cite[Lemma~13.12(b)]{QEDII});
Gaussian domination with the constant one holds on every $C(n)$, of either parity \cite[Proposition~13.11]{QEDII};
the weights of the flux sectors are bounded as in \cite[Lemma~13.12]{QEDII}.
For the averaged law no reflection positivity and no chessboard estimate is available on odd sides.

\subsubsection*{I.7.4. External connection; shift identity (Propositions~38.1--38.3)}

Let $a\in C^1(n)\otimes\mathbb{C}$ be a complex one-cochain, and put
\begin{align*}
Z^{\mathrm{x},(w)}_n(a)&:=\sum_{\substack{\mathfrak{z}\ \text{closed}\\ \text{of flux }w}}\int_{[0,r_q)^{\mathrm{links}}}\mathrm{d}A\,\exp\left(-\frac{\|F^{\mathrm{u}}\|^2}{2}\right)\cdot\det\left(D_{A+a}+\boldsymbol{\mu}\right),\\
\mathbf{Z}^{\mathrm{x}}_n(a)&:=\int\frac{\mathrm{d}\theta}{(2\pi)^3}\int_{{V_{\mathrm{t}}}}\mathrm{d}\sigma(A)\,\exp\left(-\frac{\|dA\|^2}{2}\right)\cdot\det\left(D_{A+H_\theta+a}+\boldsymbol{\mu}\right),
\end{align*}
where $F^{\mathrm{u}}=dA+r_q\cdot\mathfrak{z}$.
At a complex connection the determinant is understood as the entire function of (I.O1).
With the pairing $\langle\cdot,\cdot\rangle$ extended bilinearly to complex cochains, the shift identity reads
\[
Z^{(w)}_n\left[\exp\langle da,F^{\mathrm{u}}\rangle\right]=\exp\left(\frac{\langle da,da\rangle}{2}\right)\cdot Z^{\mathrm{x},(w)}_n(a),
\]
and in the sector of flux zero one has
\[
Z^{\mathrm{x},(0)}_n(a)=\mathrm{Nc}_n\cdot\mathbf{Z}^{\mathrm{x}}_n(a).
\]

\subsubsection*{I.7.5. Physical units, test map, the laws of the theorem}

\begin{named}{Remark}
The cells, the pairings, the test map and the Hermite--Sobolev scale are the same in \cite[Sections~7.1.1, 7.1.3, 7.1.5, 21.1.3 and the first paragraph of Section~7.1.4]{QEDI} and in \cite[Sections~12.1 and~12.2]{QEDII}.
The two papers differ in $n$ and $R$;
in the number $m=\frac{n-1}{2}$ of \cite[Section~7.1.1]{QEDI}, against the number $\frac{n}{2}$ that replaces it in \cite[Remark~12.11(a)]{QEDII};
in the set $\mathcal{P}^{+\circ}_\nu$ of \cite[Section~7.1.4]{QEDI}, against the set $\Sigma^{++}_\nu$ of \cite[Lemma~12.8]{QEDII};
in the condition (Ad3) of \cite[Definition~7.1]{QEDI}, against the condition (Ad3$'$) of \cite[Proposition~12.10]{QEDII};
in the operator;
and in the laws.
The statements of \cite{QEDII} that replace those of \cite{QEDI} on account of these differences are named in \cite[Remark~12.11(a)]{QEDII} and in the first paragraph of the proof of \cite[Proposition~12.12]{QEDII}.
For the laws of the theorem, \cite[Section~2.4]{QEDI} and \cite[Definition~2.36]{QEDII} differ in three points:
the odd side $n=h\cdot L^{T+N}$ against $n=8h\cdot L^{T+N}$;
mass assignments with $\boldsymbol{\mu}_N>-\frac{1}{2}$ against arbitrary real sequences;
the two laws (af) and (av) against the law (af) alone.
In all three points the definitions below follow \cite{QEDII}.
\end{named}

Let $\Lambda_{\eta,R}:=\eta\mathbb{Z}^3/R\mathbb{Z}^3$, where $R=n\eta$.
The map $c\mapsto\eta c$ is a bijection of $\mathbb{Z}^3/n\mathbb{Z}^3$ onto $\Lambda_{\eta,R}$ which carries cells to cells;
the cochains of the two lattices are identified through it, and $u(\eta\boldsymbol{\sigma})$ is written $u(\boldsymbol{\sigma})$.
On $\Lambda_{\eta,R}$ put
\[
\langle u,v\rangle_\eta:=\eta^3\sum u\cdot v,\qquad|u|_\eta:=\langle u,u\rangle_\eta^{\frac{1}{2}},\qquad d_\eta:=\eta^{-1}d,
\]
where $d$ is the coboundary of Section~I.1.1 in lattice units,
and let $d_\eta^{*}$ be the adjoint of $d_\eta$ with respect to $\langle\cdot,\cdot\rangle_\eta$;
thus $d_\eta^{*}=\eta^{-1}d^{*}$ (\cite[Section~7.1.1]{QEDI}; Section~G5.1.3).
For a two-cochain $u$ on $\Lambda_{\eta,R}$ the cochain $u^{\mathrm{u}}:=\eta^{\frac{3}{2}}u$ is $u$ in lattice units, and
\[
\langle u,F\rangle_\eta=\langle u^{\mathrm{u}},F^{\mathrm{u}}\rangle,\qquad|u|_\eta=\|u^{\mathrm{u}}\|.
\]
For a plaquette $\boldsymbol{\sigma}=(c;\mu\nu)$ of $\Lambda_{\eta,R}$, that is, with $c\in\Lambda_{\eta,R}$, let
\[
\mathrm{ce}(\boldsymbol{\sigma}):=c+\frac{\eta}{2}\left(e_\mu+e_\nu\right)
\]
be its centre, a point of $\mathbb{R}^3/R\mathbb{Z}^3$ in physical units, and put $Q(\boldsymbol{\sigma}):=\mathrm{ce}(\boldsymbol{\sigma})+\left[-\frac{\eta}{2},\frac{\eta}{2}\right)^3$.
If $\boldsymbol{\sigma}'$ is a plaquette of $\mathbb{Z}^3$ whose class modulo $n\mathbb{Z}^3$ is carried to $\boldsymbol{\sigma}$ by the bijection of $\mathbb{Z}^3/n\mathbb{Z}^3$ onto $\Lambda_{\eta,R}$ above, then $Q(\boldsymbol{\sigma})$ is the image in $\mathbb{R}^3/R\mathbb{Z}^3$ of the cube $Q_\eta(\boldsymbol{\sigma}')$ of \cite[Definition~12.1]{QEDII}.
The identification of the two lattices concerns the values of cochains, such as $f_\eta(\boldsymbol{\sigma})$ and $F(\boldsymbol{\sigma})$ below, and not $\mathrm{ce}$ and $Q$.
Let $\mathcal{S}:=\mathcal{S}(\mathbb{R}^3;\Lambda^2)$, let $\mathcal{S}'$ be its dual, and let $H_s$, $s\in\mathbb{R}$, be the Hermite--Sobolev scale.
For $f\in\mathcal{S}$ put
\[
f^{R}(y):=\sum_{k\in\mathbb{Z}^3}f(y+R\cdot k),\qquad f_\eta(\boldsymbol{\sigma}):=\eta^{-3}\int_{Q(\boldsymbol{\sigma})}\mathrm{d}y\,f^{R}_{\mu\nu}(y),\qquad X_{\eta,R}(f):=\langle f_\eta,F\rangle_\eta,
\]
and let $F^{\sharp}$ be the $R\mathbb{Z}^3$-periodic step two-form that takes the value $F(\boldsymbol{\sigma})$ on $Q(\boldsymbol{\sigma})$.
Then $F^{\sharp}(f)=X_{\eta,R}(f)$.
These definitions, and \cite[Lemmas~12.5, 12.6, 12.7 and~12.9]{QEDII}, hold for $n$ of either parity;
\cite[Lemma~12.5(c)]{QEDII} holds for a period cube with an arbitrary centre.
The tests are acted upon by $\tau_b$, $b\in\mathbb{R}^3$, by $U_r$, $r\in O_{\mathrm{cub}}$, and by $\mathsf{R}_\nu$, $\nu\in\{1,2,3\}$.
For an isometry $\mathrm{Gi}$ of $\mathbb{R}^3$ with linear part $\mathrm{Rl}$ one puts
\[
(\mathrm{Gi}\cdot f)(y):=\left(\wedge^2\mathrm{Rl}\right)f\left(\mathrm{Gi}^{-1}y\right),
\]
and $\tau_b$, $U_r$, $\mathsf{R}_\nu$ are the maps $f\mapsto\mathrm{Gi}\cdot f$ for the translation $\mathrm{Gi}(y)=y+b$, for $\mathrm{Gi}=r$, and for the reflection $\mathrm{Gi}$ of $\mathbb{R}^3$ in the plane $\{y_\nu=0\}$, respectively;
thus
\[
(\tau_bf)(y)=f(y-b),\qquad(U_rf)(y)=\left(\wedge^2r\right)f\left(r^{-1}y\right).
\]

\paragraph{The laws.}
Let $P_{N,T}$ be the law on $\mathcal{S}'$ (a Borel probability measure on $H_{-3}$) of $F^{\sharp}$ when $F^{\mathrm{u}}$ has the law $E^{\mathrm{af}}_n$ with
\[
n=8h\cdot L^{T+N},\qquad q=e\cdot\eta^{\frac{1}{2}},\qquad\boldsymbol{\mu}=\boldsymbol{\mu}_N=\bar{m}\cdot\eta+m_{0,N}.
\]
It is the law $P^{\mathrm{ch}}_{N,T}$ of \cite[Definition~2.36]{QEDII}.
A \emph{mass assignment} is an arbitrary real sequence $N\mapsto m_{0,N}$ that does not depend on $T$;
in the theorem it is the sequence $m_{0,N}(e)$ of Section~I.6.2.

\paragraph{Limit points.}
The following notions are those of \cite[Definition~2.37]{QEDII}.
A \emph{limit point} of $(P_{N,T})$ is a weak limit on $H_{-4}$ along a sequence $(N_i,T_i)$ with $N_i\to\infty$ and $T_i\to\infty$.
Let $\mathcal{V}_N$ be the set of the weak limits on $H_{-4}$ of the sequences $P_{N,T_j}$ with $T_j\to\infty$.
A \emph{volume-first limit point} is a weak limit on $H_{-4}$ of a sequence $\mu^{(N_i)}\in\mathcal{V}_{N_i}$ with $N_i\to\infty$.
Let $\Sigma$ be the cylinder $\sigma$-algebra of $\mathcal{S}'$, that is, the $\sigma$-algebra generated by the maps $\omega\mapsto\omega(f)$, $f\in\mathcal{S}$.
The set $H_{-4}$ belongs to $\Sigma$, and the trace of $\Sigma$ on $H_{-4}$ is the Borel $\sigma$-algebra of $H_{-4}$ (\cite[Section~7.1.5]{QEDI}).
Accordingly a Borel probability measure $\mu$ on $H_{-4}$ is read as the probability measure $\mathrm{Ev}\mapsto\mu\left(\mathrm{Ev}\cap H_{-4}\right)$ on $(\mathcal{S}',\Sigma)$, and a Borel probability measure on $H_{-3}$ is read in the same way.
For a limit point $\mu$ let $\nu$ be the measure of its ergodic decomposition under the translations of $\mathbb{R}^3$,
\[
\mu=\int\mathrm{d}\nu(\xi)\,\mu_\xi
\]
(\cite[Theorem~8.1]{QEDI}, with the index set $\{1,2,3\}$).
Each almost sure statement on the components holds outside a $\nu$-null set of its own, which may depend on the parameters of the statement.
In this sense the properties 1--6 of \cite[Corollary~8.9]{QEDI} hold: in property~5 the exceptional null set depends on the event, and part of property~6 is a statement on $\nu$ itself.

\subsubsection*{I.7.6. Probes, free ratio, reference numbers}

The objects of this section are those of \cite[Sections~21.1.6, 21.1.7 and~21.1.9]{QEDI}, formed with the determinant of the chamber fermion.

\paragraph{Probe.}
A \emph{probe} is a real two-form $\beta\in C_c^{\infty}(\mathbb{R}^3;\Lambda^2)$ with $\beta\ne0$.
One puts $\beta_{\nu\mu}:=-\beta_{\mu\nu}$ and $\beta_{\mu\mu}:=0$; the same convention is used for all two-forms and two-cochains.
To a probe belong
\[
a:=d^{*}\beta,\quad a_\nu=-\sum_\mu\partial_\mu\beta_{\mu\nu};\qquad f:=dd^{*}\beta;\qquad b(f):=-6\operatorname{Tr}\left[\left(S_{\bar{m}}\,i\gamma\cdot a\right)^4\right].
\]
Let $D_\beta$ be the least side of a closed coordinate cube
\[
Q_\beta=y_\beta+\left[-\frac{D_\beta}{2},\frac{D_\beta}{2}\right]^3
\]
that contains $\operatorname{supp}\beta+[-2,2]^3$, in physical units.
The number $D_\beta$ is determined by $\beta$; the centre $y_\beta$ is determined by $\beta$ only if the least coordinate box that contains $\operatorname{supp}\beta$ is a cube.
With each probe one centre $y_\beta$, hence one cube $Q_\beta$, is chosen once and for all, and with $\tau_b\beta$ the cube $b+Q_\beta$.
The following objects are formed with this choice: the placements $v$ defined below;
the cube $\mathrm{Qb}:=v+Q_\beta$ of (G5.1.1) in Section~G5.1.3;
and the number $r_g:=|y_\beta|_\infty+\frac{D_\beta}{2}$ of Propositions~47.2\,(e) and~43.2.
In \cite[Definition~5.88]{QEDI} one has $D_\beta=2\rho_\beta$.
Let $D_{\mathrm{be}}$ be the least side of a closed coordinate cube that contains $\operatorname{supp}\beta$; then $D_\beta=D_{\mathrm{be}}+4$.
The probe $\beta$ \emph{fits} $(N,T^{*})$ if
\[
D_\beta+3\eta\le\ell_c.
\]

\paragraph{Lattice probe.}
For a function $\mathrm{fs}$ on $\mathbb{R}^3$ let
\[
(J_{\eta}\mathrm{fs})(y):=\eta^{-3}\int_{\left[-\frac{\eta}{2},\frac{\eta}{2}\right]^3}\mathrm{d}y'\,\mathrm{fs}(y+y'),
\]
and put
\[
u_{\mu\nu}(c):=\left(J_{\eta}\beta_{\mu\nu}\right)\left(c+\frac{\eta}{2}\left(e_\mu+e_\nu\right)\right)\qquad\left(c\in\Lambda_{\eta,R}\right),
\]
periodised: a function $\mathrm{fs}$ with compact support on $\mathbb{R}^3$ is evaluated at a point of $\mathbb{R}^3/R\mathbb{Z}^3$ through
\[
\mathrm{fs}^{R}(y):=\sum_{z\in\mathbb{Z}^3}\mathrm{fs}(y+R\cdot z),
\]
and the same convention holds in Proposition~47.1\,(b1).
Let $a_{\eta}:=d_\eta^{*}u$, with $d_\eta$ and $d_\eta^{*}$ as in Section~I.7.5.
A \emph{placement} is a vector $v\in\eta\mathbb{Z}^3$ with
\[
\left|v+y_\beta-\frac{\ell_c}{2}(1,1,1)\right|\le\frac{\eta}{2},
\]
and $a_{\eta}$ is formed with $\tau_v\beta$.
There may be several placements.
The quantities $\boldsymbol{\sigma}^2$, $N_0$, $c_m$, $a_{\mathrm{ch}}$, $b_{\mathrm{ch}}$ defined below do not depend on the placement (Proposition~47.1\,(b5$'$), Proposition~50\,(h));
the function $\mathring{\Phi}$, the numbers $\Theta^{(0)}_m$ and the numbers $\mathsf{u}_m$ of Proposition~42.2 may depend on it, because
\[
(\mathsf{g}_{\mathfrak{c}})_{*}\tau_w=\tau_{r_{\mathfrak{c}}w}(\mathsf{g}_{\mathfrak{c}})_{*},
\]
where $r_{\mathfrak{c}}$ denotes the linear part of $\mathsf{g}_{\mathfrak{c}}$.
Put
\[
\hat{\mathfrak{a}}:=\eta^{\frac{1}{2}}a_{\eta},\qquad\boldsymbol{\sigma}^2:=\|d\hat{\mathfrak{a}}\|^2,\qquad X:=\langle d\hat{\mathfrak{a}},F^{\mathrm{u}}\rangle,\qquad\hat{\mathfrak{a}}_{\mathfrak{c}}:=(\mathsf{g}_{\mathfrak{c}})_{*}\hat{\mathfrak{a}},\qquad X_{\mathfrak{c}}:=\langle d\hat{\mathfrak{a}}_{\mathfrak{c}},F^{\mathrm{u}}\rangle;
\]
thus $\hat{\mathfrak{a}}$ is a one-cochain in lattice units.
For a real variable $z$ let
\[
H_2(z):=z^2-\boldsymbol{\sigma}^2,\qquad H_4(z):=z^4-6\boldsymbol{\sigma}^2\cdot z^2+3\boldsymbol{\sigma}^4.
\]
Modulo the periods $n\mathbb{Z}^3$, every plaquette of $\operatorname{supp}d\hat{\mathfrak{a}}$ lies in the box
\[
\mathrm{Bx}_0:=\frac{\mathrm{hc}}{2}(1,1,1)+[-r_0,r_0]^3,\qquad r_0:=\frac{\mathrm{hc}}{2}-\frac{1}{\eta}\left(2-\frac{3\eta}{2}\right),
\]
which lies inside the open cell $Q^{\circ}_{0}$ (Proposition~41);
hence $X$ and its polynomials belong to $\mathcal{A}(Q^{\circ}_{0})$.
Here and in Proposition~41, $r_0$ denotes the half side of this box;
in $N_{\min}$ (Section~I.6.2) and in Proposition~31.2, $r_0$ is the number $r_k$ of Section~I.6.1 at $k=0$.

\paragraph{Free ratio.}
On the torus under consideration, which is $C(n_*)$ unless another torus is named, let
\[
N_0(t):=\frac{\det\left(D_{e^{-1}t\hat{\mathfrak{a}}}+\bar{m}\cdot\eta\right)}{\det\left(D_0+\bar{m}\cdot\eta\right)},\qquad c_{m}:=N_0^{(m)}(0),
\]
and
\[
a_{\mathrm{ch}}(\eta,R):=c_{2},\qquad b_{\mathrm{ch}}(\eta,R):=\left.\frac{\mathrm{d}^4}{\mathrm{d}t^4}\log N_0(t)\right|_{t=0}=c_4-3c_2^2,
\]
the last identity because $N_0$ is even.
In \cite[Section~2.6]{QEDI} the corresponding objects are written $N_0(t)$ and $b_{\eta,R}(0)$.
The function $N_0$, and with it $c_m$, $a_{\mathrm{ch}}$ and $b_{\mathrm{ch}}$, does not depend on the vector $v\in\eta\mathbb{Z}^3$ with which $\hat{\mathfrak{a}}$ is formed, and it does not change when $\hat{\mathfrak{a}}$ is replaced by $\hat{\mathfrak{a}}_{\mathfrak{c}}$ (Proposition~50\,(h), Proposition~49\,(d), Proposition~36\,(h)).
Consequently the function $N_0$ formed here with a centred placement, the one of Section~G5.1.5, formed with the placement at hand, and the one of Section~G5.6.0, formed with $v=0$, coincide.

\paragraph{Reference numbers.}
They are formed at the cutoff $N$, with the charge $e$ and the lattice mass $\boldsymbol{\mu}_N$, and with $c_m=N_0^{(m)}(0)$.
Where a statement admits an arbitrary real lattice mass $\boldsymbol{\mu}$, or arbitrary real numbers $c_2$, $c_4$ (Proposition~38.3\,(e), Proposition~40\,(e), Proposition~42.2), the same displays are formed with that $\boldsymbol{\mu}$ and with those $c_m$.
Let
\begin{gather*}
D^{\mathrm{mat}}:=L^3\log\mathrm{zr}_{n_*}-\log\mathrm{zr}_{Ln_*}\quad\text{with }\mathrm{zr}_n:=\mathrm{zr}_n(q,\boldsymbol{\mu}_N),\qquad D_L:=L^3\log Z_{n_*}[1]-\log Z_{Ln_*}[1],\\
\delta:=\frac{D_L}{(L-1)^3\cdot N_{*}},\qquad a_{L}:=\frac{L^3-1}{(L-1)^3}.
\end{gather*}
On $C(n_*)$, for $t=(t_{\mathfrak{c}})\in\mathbb{C}^{N_{*}}$, let
\begin{gather*}
\bar{Z}^{\mathrm{ext}}(t):=\int\frac{\mathrm{d}\theta}{(2\pi)^3}\int\mathrm{d}\gamma_N(F^{\mathrm{u}})\,\det\left(D_{A(F^{\mathrm{u}})+H_\theta+e^{-1}\sum_{\mathfrak{c}}t_{\mathfrak{c}}\hat{\mathfrak{a}}_{\mathfrak{c}}}+\boldsymbol{\mu}_N\right),\qquad\mathring{\Phi}(t):=\frac{\bar{Z}^{\mathrm{ext}}(t)}{\bar{Z}^{\mathrm{ext}}(0)},\\
\Theta^{(0)}_m:=\left|\left.\prod_{\mathfrak{c}}\left(\partial_{t_{\mathfrak{c}}}^{m}-c_m\right)\mathring{\Phi}(t)\right|_{t=0}\right|^{\frac{1}{N_{*}}}\qquad(m=2,4).
\end{gather*}
For real variables $z$, $z'$, $\mathrm{cf}$ let
\[
\mathrm{Pw}_2(z;z',\mathrm{cf}):=e^{-2}\cdot\left(z^2-z'\right)-\mathrm{cf},\qquad\mathrm{Pw}_4(z;z',\mathrm{cf}):=e^{-4}\cdot\left(z^4-6z'\cdot z^2+3z'^2\right)-\mathrm{cf}.
\]

\paragraph{Reference data.}
The seven clauses of Proposition~43.1 are (R0)--(R6):
(R0) concerns the mass assignment, and (R1)--(R6) are six ultraviolet statements on the two reference tori.
Together with Propositions~42.1 and~42.2 they yield the ultraviolet proposition, Proposition~43.2, with the rate $\eta^{\frac{1}{4}}$ in the spacing.

\paragraph{Tests of the theorem.}
Let $x_0$ be a point with the inequality of \cite[Section~2.8]{QEDI} (see (Pb) in Section~I.8.4), and put
\[
g_0:=f,\qquad g_\pm:=f\pm\tau_{x_0}f.
\]
Where these tests occur, the letter $g$ does not denote the label of a fermion site.

\subsection*{I.8. Constants, in the order in which they are chosen}

\subsubsection*{I.8.1. Rules}

(I.U4)$^{\mathrm{ch}}$\quad A \emph{constant} is a number that depends only on quantities chosen earlier in the list of Section~I.8.2.
No constant of the lines (C0)--(C9) depends on $N$, on the level $k$, on the side $\ell$, on the twist $\theta$, on the running mass in $(0,1]$, on the pair $(m_{\mathrm{in}},\varepsilon^{\mathrm{v}}_{\mathrm{in}})$, on $\mathscr{C}$, or on $\bar{m}\in(0,1]$.
The constants of the line (C10) may depend on $\bar{m}$ and on the probes.

The number $\rho$ is chosen first.
Every constant below the line (C0) may depend on $\rho$ and on the weights; this dependence is not indicated again.
All statements are made at a fixed $\rho$, and no uniformity in $\rho$ is asserted or required:
some of the fermionic constants that contain $\rho$ degenerate as $\rho\to0$ (Proposition~14.4\,(d)), and \cite[Theorem~4.13(c)]{QEDII} and the bounds of \cite[Section~9]{QEDII} degenerate as $\rho\to1$.

The constants are chosen in the order of Section~I.8.2, each of them large or small enough for all the statements to hold.
The conditions of Section~I.8.4 are of two kinds.
\begin{itemize}
\item[(i)] Conditions on the choice of the constants: (A0)--(A10), (Z2.m), (Z2.s1)--(Z2.s5), (Z3.s1), (Z8.s1), (Z10.s1), (Z10.s2), (Z10.s4), and the condition (Pb) on the choice of $(\beta,x_0)$.
\item[(ii)] Conditions on the data $(e,T,N)$:
the conditions (G4.2.1), (fl) and (T1)--(T6), which hold whenever $e\cdot h^{\frac{1}{2}}\le e_1$, $R\in\mathrm{Tor}(e)$ and $N\ge N_{\min}(R,e)$ (Proposition~14.2\,(c));
the conditions (W3) and (Z8.w), which hold whenever $e\le e_G$;
and the conditions (g1)--(g9), of which (Z9.s1)--(Z9.s3) are parts, which hold whenever $e\le e_\sharp(\beta)$ (Proposition~14.2\,(d), Section~G5.1.9).
\end{itemize}
The conditions of the second kind hold by the definitions of $e_1$, $e_G$ and $e_\sharp$ as soon as the charge $e$ lies below the ceiling that the proposition in question names.
Accordingly the hypotheses of a proposition concern its data only, and neither kind of condition is among them:
the constants in a proposition are those of Section~I.8, and the conditions of the second kind enter it through the ceiling on $e$ that it names.
Propositions~14.1--14.4 contain the complete list of the conditions and the proof that admissible $L$, $M$, $e_0$ and $e_1$ exist.

\subsubsection*{I.8.2. The ordered list}

In the following table the constants are arranged by lines (C$j$);
a line consists of one or more rows, and each row names constants chosen at that stage, the quantities on which they may depend, and their relation to \cite[Definition~5.97]{QEDI}, in which the line (C$j$) is stage~$j+1$.

\begin{description}
\item[(C0)]
$\rho\in(0,1)$; the weights $w=\left(6,6\sqrt{2},6\sqrt{2},6\right)$ and $\mathrm{sw}$.
They depend on nothing.
They have no counterpart among the constants of \cite[Definition~5.97]{QEDI}; $\rho$ and the weights are those of \cite[Definition~2.10]{QEDII}.

\item[(C0)]
The absolute numbers
\[
\bar{\varepsilon}:=\frac{1}{200};\qquad c_{X}:=54,\quad\kappa_{0}:=54\log72,\quad K_{0}:=\frac{\exp(\kappa_0)}{36},\quad c_{\mathrm{rb}}:=c_X\cdot K_0,\quad c_{\mathrm{fl}}:=61\,c_X\cdot K_0;\qquad b:=1;
\]
\[
c_{\mathrm{cl}}:=\frac{96\cdot\left(1+\sqrt{2}\right)}{4\pi^2}\cdot\left(128+64\cdot\left(\frac{16}{\pi}+\frac{64}{\pi^2}\right)\right).
\]
The numbers $\bar{\varepsilon}$ and $b$ are those of stage~1 of \cite[Definition~5.97]{QEDI}, and $c_X$, $\kappa_0$, $K_0$, $c_{\mathrm{rb}}$, $c_{\mathrm{fl}}$ those of \cite[Section~13.1.6]{QEDI}.
The number $c_{\mathrm{cl}}$ is the constant $C_0$ of \cite[Theorem~12.39]{QEDII}; it contains no constant of the list, and no condition is imposed on it.

\item[(C0r)]
$\lambda_{\mathrm{int}}=1.99092\ldots$ (Section~I.3.1);
$\lambda_1:=6+9\sqrt{2}-\sqrt{198-36\sqrt{2}}=6.599925\ldots$, the gap of $\sum_{i\le3}w_i\cdot\left(1-R_i\right)$ at $p=0$.
They depend on the weights.
Neither of them is a constant of \cite[Definition~5.97]{QEDI}.
The number $\lambda_{\mathrm{int}}$ occurs in Proposition~1.1, which has no counterpart in \cite{QEDI} (Remark after Proposition~1.1).

\item[(C0r)]
$\hat{g}:=\frac{\rho\cdot\lambda_{\mathrm{int}}}{2}$;
$c_{\mathrm{hl}}:=\sqrt{3}\cdot(1+\rho)$;
$c_{*}:=\frac{1}{18}\min\left(w_3,\lambda_{\mathrm{int}}\right)$;
the number $\theta_{0}$ with $3(1+\rho)\cdot\left(\exp(\theta_0)-1\right)=\frac{\hat{g}}{2}$;
$K_{\rho}:=\rho+\frac{2c_{\mathrm{hl}}^2}{\hat{g}}$;
$t_{0}:=\frac{1}{4K_\rho^2}$;
$c_{\mathrm{el}}:=\min\left\{0.22,\frac{\rho\cdot c_*\cdot t_0}{\pi\sqrt{3}}\right\}$;
$K_{\mathrm{Str}}$, the bound of the Schur symbol on the strip of width $\theta_0$;
$\kappa_{(0)}:=\min\left\{\frac{\theta_0}{2},\frac{c_{\mathrm{el}}\cdot\theta_0}{12K_{\mathrm{Str}}}\right\}$;
the constants $A_{\mathrm{ak}}$, $A'_{\mathrm{ak}}$ of Proposition~1.2\,(S6), (S7);
$c_{b}$ (Proposition~1.3\,(b));
$C_{\mathrm{Po}}:=\frac{1}{8c_*}+\frac{1}{\lambda_{\mathrm{int}}}$;
$C_{G}$ (Proposition~1.5);
$C_{6}$, whose formula is given with Proposition~30.1 and in Proposition~14.2\,(a);
$c_{\rho}$, the largest number with which, at the given $\rho$, the inequality $D_0^{*}D_0\ge c_\rho\cdot\nabla^{*}\nabla$ holds on all the spaces of spinors named in the row $c_\rho$ of Proposition~14.2\,(a).
They depend on $\rho$.
None of them is a constant of \cite[Definition~5.97]{QEDI}.
Each of them other than $c_\rho$ occurs in one or more of Propositions~1.1--1.5, 30.1 and~30.3, and its relation to \cite{QEDI} is the one stated in the Remarks after these propositions:
Proposition~1.1 has no counterpart in \cite{QEDI};
Propositions~1.2, 1.3, 1.4 and~1.5 correspond, in this order, to \cite[Lemma~12.5]{QEDI} together with \cite[Remark~12.2]{QEDI}, to \cite[Lemma~12.1]{QEDI}, to \cite[Lemma~12.3]{QEDI}, and to \cite[Proposition~12.6]{QEDI} together with \cite[Lemma~12.7(b)]{QEDI};
in the function $\varepsilon_5$ of Proposition~30.1 the constants $C_6$ and $\kappa_{(0)}$ stand in the place of the numerals $312$ and $\frac{1}{4}$ of \cite{QEDI};
and Proposition~30.3 corresponds to \cite[Section~19.8.2]{QEDI} and to \cite[Lemmas~19.20 and~19.21]{QEDI}.

\item[(C0r)]
$C_{E}$ (Proposition~37.1\,(f));
$\mathrm{Csy}_{\mathrm{ak}}$ and $\mathrm{Cpv}_0$ (Proposition~37.1\,(g) and the bounds of Proposition~36\,(c));
\[
c_{\mathrm{N3}}:=\min\left(\frac{\rho}{\pi},\frac{\rho^3\cdot\lambda_{\mathrm{int}}}{98\sqrt{3}\,\pi^3}\right).
\]
They depend on $\rho$.
The constants $C_E$, $\mathrm{Csy}_{\mathrm{ak}}$, $\mathrm{Cpv}_0$ are obtained by compactness arguments and are not explicit.
None of the constants of this row is a constant of \cite[Definition~5.97]{QEDI}.
The constants $C_E$ and $\mathrm{Csy}_{\mathrm{ak}}$ occur in Proposition~37.1, which corresponds to the first and the third bound of \cite[Lemma~21.8]{QEDI}, and $\mathrm{Cpv}_0$ occurs in Proposition~36, which has no counterpart in \cite{QEDI} (Remarks after Propositions~36 and~37.1).

\item[(C0r)]
The functions
\[
\mathfrak{d}_{\mathrm{ma}}(p):=\det D^{(p)}\ \text{at}\ \boldsymbol{\mu}=\mathbf{m},\qquad p^{0}_{\mathrm{ch}}(\mathbf{m}):=\int_{[-\pi,\pi]^3}\frac{\mathrm{d}^3p}{(2\pi)^3}\,\log\mathfrak{d}_{\mathrm{ma}}(p),
\]
where in the index of $\mathfrak{d}_{\mathrm{ma}}$ the letters $\mathrm{ma}$ stand for the mass $\mathbf{m}$, and
\[
\varepsilon_5(x):=C_6\cdot\log\left(\frac{1}{1-\exp\left(-\kappa_{(0)}\cdot x\right)}\right).
\]
They depend on $\rho$.
In \cite[Section~19.8.1]{QEDI} the corresponding functions are $w_{\mathbf{m}}$ and $p^0$, and in \cite[Section~5.6.2]{QEDI} one has $\varepsilon_5(x)=312\log\left(\left(1-\exp\left(-\frac{x}{4}\right)\right)^{-1}\right)$.

\item[(C1)]
$\alpha\in\left(1-\bar{\varepsilon},1\right)$, as at stage~2 of \cite[Definition~5.97]{QEDI}.

\item[(C2)]
$\kappa\ge7\kappa_0+7$.
It may depend on the constants of the line (C0) other than $\rho$.
This is the choice of stage~3 of \cite[Definition~5.97]{QEDI}.

\item[(C3)]
An odd integer $L$ that satisfies the five conditions imposed on $L$ at stage~4 of \cite[Definition~5.97]{QEDI}; the thresholds are $L^{F}=L^{P}=5$.
It may depend on $\bar{\varepsilon}$, $c_{\mathrm{fl}}$, $c_{\mathrm{rb}}$, $\kappa$, $\kappa_0$, as in \cite[Definition~5.97]{QEDI}.
The line (C3) contains no fermionic constant and not $\rho$.
It gives $L\ge4^8$.

\item[(C4)]
$b_{L}$, $b_{-}$, $b_{+}$;
$N_{*}=8L^3$, $a_L$, $G_{L}$, $K_{L}$;
$N_c$, $c_{l}$, $d_U$;
$c_{C}$ (\cite[Proposition~5.35]{QEDI}), $C_{Q}$;
$\alpha^{P}$, $\alpha^{F}:=\frac{1}{54}$, $\alpha_{0}$;
the photon constants $c_{P}$, $\gamma_{P}$, $\gamma^{P}$, $C_{P}$;
$\kappa_{\mathrm{w}}:=\kappa+\kappa_0+3$;
$\kappa':=\kappa_{\mathrm{w}}+\kappa_0$;
$\mu_-:=L^{-3}$;
$\hat{\kappa}:=N_c\cdot\kappa+\log\left(6c_X\right)$;
$C'_{V}:=2\mathrm{e}\cdot(42+6b)$;
\[
C_{T}:=2\mathrm{e}^{\frac{3}{2}}\cdot\left[b_-^{-1}\cdot\left(54L\cdot\mathrm{e}^6+12\mathrm{e}+1\right)+1\right].
\]
They may depend on $b$, $\alpha$, $\kappa$, $L$, not on $\rho$.
Their relation to stage~5 of \cite[Definition~5.97]{QEDI} is as follows.
The constants $b_{-}$, $b_{+}$, $N_{*}$, $a_L$, $G_{L}$, $K_{L}$, $N_c$, $c_{l}$, $d_U$, $c_{C}$, $C_{Q}$, $C_{T}$, $\alpha^{F}$, $\alpha^{P}$, $\alpha_{0}$, $\gamma_{P}$, $\gamma^{P}$, $C_{P}$, $c_{P}$, $\kappa_{\mathrm{w}}$ and $\hat{\kappa}$ of the line (C4) are constants of stage~5, with the same definitions; there $N_{*}$ is written $N^{*}$.
The other constants of stage~5 are $\gamma^{F}$ and $\gamma_{\mathrm{w}}$, which stand in the line (C4r) because here they depend on $\rho$, and $\mathfrak{m}_0$, $m_0$, $h'$, which are read as $\mathfrak{m}_0=1$, $m_0=0$, $h'=h$ (Section~I.6.1).
Thus the line (C4), together with $\gamma^{F}$ and $\gamma_{\mathrm{w}}$ of the line (C4r), takes the place of stage~5.
The constants $b_{L}$, $\kappa'$, $\mu_-$ and $C'_V$ of the line (C4) are not named at stage~5:
$b_{L}=\frac{b}{L}$ is written $b\cdot L^{-1}$ in \cite[Section~5.2.2]{QEDI} and is listed in \cite[Section~16.14]{QEDI};
$\kappa'=\kappa_{\mathrm{w}}+\kappa_0$ is the value at which $M^{F}$ is taken at stage~6;
$\mu_-=L^{-3}$ is the number $\mathfrak{m}_{-}=L^{-m_0-3}$ of \cite[Section~5.5.8]{QEDI} at $m_0=0$;
and $C'_V$ takes the place of the constant $C_V=12\mathrm{e}\cdot(2+b)$ of \cite[Lemma~16.4]{QEDI}, which is not itself a constant of \cite[Definition~5.97]{QEDI} and enters it through $M^{F}$ at stage~6 and through the condition (Z2.s2) of Section~I.8.4 at stage~9.
The constant $C_T$ has the value of \cite[Proposition~16.27]{QEDI}: the bound for which it serves here carries the number $12+6\mathrm{e}$ in place of $12\mathrm{e}$, and $12+6\mathrm{e}\le12\mathrm{e}$.

\item[(C4r)]
$K_{1}$, $C_{I}$, $c_{\mathrm{wt}}$;
$a_{0}:=\min\left\{\frac{1}{3},\frac{1}{2C_I\cdot c_{\mathrm{wt}}}\right\}$;
$C_{R}$, $C_{3}$;
for $K_1$, $C_I$, $c_{\mathrm{wt}}$, $C_R$ and $C_3$ the formulas are given with Proposition~1.6 and in Proposition~14.2\,(a);
$C_{5}:=\mathrm{e}^2\cdot L\cdot C_3$, which is $\ge1$;
$\gamma_{(5)}:=\frac{a_0}{L}$;
\[
c_{\mathrm{re}}:=\min\left\{b\cdot\left(1-\frac{1}{L}\right),\frac{\rho}{2C_{\mathrm{Po}}\cdot L}\right\},\qquad a_{z}:=\min\left\{\frac{1}{3L},\frac{c_{\mathrm{re}}}{2\mathrm{e}\cdot L\cdot(6+3b)}\right\},
\]
\[
k_{1}:=60+b,\qquad C_{00}:=48\cdot\frac{2L}{c_{\mathrm{re}}}\cdot\left(1+2k_1\right);
\]
$C_z:=\max\left(C_5,C_{00}\right)$, $\gamma_{z}:=\min\left(\gamma_{(5)},a_z\right)$;
$\gamma^{F}:=\frac{\gamma_z}{2}$, which is $\le\frac{1}{6L}$;
for $\gamma>0$ let
\[
c(\gamma):=\sum_{z\in\mathbb{Z}^3}\exp\left(-\gamma\cdot|z|\right),\qquad c_{\mathrm{bd}}(\gamma):=3\exp(3\gamma)\cdot c(\gamma),
\]
with $|z|$ the sup norm, as in Proposition~14.2\,(a) and at stage~7 of \cite[Definition~5.97]{QEDI};
$a':=C_z\cdot c(\gamma^F)$, which is $\ge1$;
$\gamma_{\mathrm{w}}$.
They may depend on $b$, $\alpha$, $L$, $\rho$.
The formulas are those of \cite[Section~16.4]{QEDI}, with these differences: in \cite{QEDI} the constants $C_5$ and $\gamma_{(5)}$ are those of \cite[Section~12]{QEDI}, $c_{\mathrm{re}}=\min\left\{b\cdot\left(1-\frac{1}{L}\right),\frac{4}{L}\right\}$, $k_1=14+b$, $C_{00}$ is defined without the factor $48$, and $a'$, $\gamma^F$ are written $a$, $\gamma$.

\item[(C5)]
$M$, a power of $L$, with
\[
M\ge L^2,\quad M>4L,\quad\gamma_{\mathrm{w}}\cdot M\ge16,\quad M^{\alpha_0}\ge1+\exp(\hat{\kappa}),\quad M\ge M^{P},
\]
\[
M\ge M^{F}:=\max\left\{4L+1,\left(108\pi\cdot C'_V\cdot a'\right)^2,108\pi\cdot C'_V\cdot a'\cdot\exp\left(27\kappa'\right)\right\}\quad\text{at }\kappa'=\kappa_{\mathrm{w}}+\kappa_0.
\]
It may depend on the constants of (C0)--(C4r).
The formulas are those of stage~6 of \cite[Definition~5.97]{QEDI}, with $C'_V$ and $a'$.
All conditions on $M$ are lower bounds; none of them contains $C_w$, $\varepsilon$ or $e$.

\item[(C6)]
$C_{\mathrm{By}}$, $C_{\Gamma}$ (\cite[Section~16.9.2]{QEDI}), $C_{1}$ (\cite[Theorem~16.23]{QEDI});
$C^{F}$, which contains the factor $2$ of Section~I.1.5 and, through $C_{00}$, the factor $48$ of level~$0$;
$C^{P}$, $C_{w}^{(0)}$, $C_{w}^{(1)}$, $C_{w}:=64\,K_0\cdot C_w^{(1)}$;
\[
\mathrm{par}^{0}:=\left(\alpha_0,\gamma_{\mathrm{w}},\kappa_{\mathrm{w}}+\kappa_0;C_w^{(1)}\right),\qquad\mathrm{par}:=\left(\alpha_0,\gamma_{\mathrm{w}},\kappa_{\mathrm{w}};C_w\right).
\]
They may depend on the constants of (C0)--(C5).
The formulas are those of \cite[Section~16.14]{QEDI} for $C^F$, with $a'$ in place of $a$, of \cite[Section~15.10.3]{QEDI} for $C^P$, of \cite[Section~5.5.3]{QEDI} for $C_w^{(0)}$ and $C_w^{(1)}$, and of \cite[Section~5.4.1]{QEDI} for $C_w$, $\mathrm{par}^0$ and $\mathrm{par}$.
The constants written $C_F$, $\gamma_F$ in \cite[Section~5.4.4]{QEDI}, \cite[Section~13]{QEDI} and \cite[Section~20]{QEDI} are $C^F$, $\gamma^F$.

\item[(C6)]
$c_{\gamma}$ (Section~I.5.8, with the formula of stage~7 of \cite[Definition~5.97]{QEDI}; through $\gamma_{\mathrm{w}}$ it depends on $\rho$);
$C_{W}:=28\,C_w^2\cdot c_\gamma$;
$C_{B}$;
$C_{d}$ (Proposition~13; it contains the row sum (I.L0.12));
$c_{\kappa}$;
\[
C_{0}:=\max\left\{C_P\cdot c_{\mathrm{bd}}(\gamma_P),\ C^F,\ 4b_++16\,C^F\cdot c(\gamma^F)\right\};
\]
$\gamma_{\Sigma}$, $C_{\Sigma}$, whose formulas are given with Proposition~30.3 and in Proposition~14.2\,(a);
$c_{S}:=10^5\,c_X^5\cdot K_0^5\cdot\exp(6\kappa)$.
They may depend on the constants of (C0)--(C5); $\gamma_\Sigma$ and $C_\Sigma$ depend on $(b_-,b_+,\rho)$.
In \cite[Lemma~17.11]{QEDI} one has $C_W=4C_w^2\cdot c_\gamma$, and $\gamma_\Sigma$ is written $\gamma_1$ in \cite[Lemma~19.21]{QEDI}.

\item[(C6$'$)]
\begin{gather*}
\mathrm{ssum}(\gamma):=\left(\frac{12}{\gamma}\right)^3;\qquad\mathbf{S}:=C^F\cdot\mathrm{ssum}\left(\frac{\gamma^F}{2}\right);\qquad\mathrm{hL}:=6\,\mathbf{S};\\
C_K:=226\,\mathbf{S}^2;\qquad C_K':=1328\,\mathbf{S}^2;\qquad C_H:=203\,M\cdot\left(C^F\right)^2;\\
\mathrm{CVj}:=10^3\,C_H^2\cdot\mathrm{ssum}\left(\frac{\gamma^F}{2}\right)^4\cdot K_0;\qquad\theta_{\mathrm{CT}}:=\frac{\mu_-}{12\mathrm{e}};\qquad\mathrm{vt}:=\min\left\{\frac{\theta_{\mathrm{CT}}}{2},\frac{\gamma^F}{8},\frac{1}{M}\right\};\\
\mathrm{gGa}:=\left(b_-^{-1}+\mu_-^{-1}\right)^{\frac{1}{2}};\qquad h_{G}:=1+\mathrm{gGa};\qquad\mathrm{gGa}_{1}:=b_-^{-1}+\frac{4}{\mu_-};\\
\mathrm{alG}:=2b_-^{-1}+4\mathrm{e}\cdot\mu_-^{-1}\cdot\mathrm{ssum}\left(\frac{\theta_{\mathrm{CT}}}{2}\right);\qquad c_{\tau}:=\min\left\{\frac{1}{8},\frac{\mu_-}{48},\frac{1}{292\,\mathbf{S}}\right\};\\
\mathrm{CWj}:=3\,\mathrm{CVj}\cdot h_G^2+2\,\mathrm{alG}\cdot\mathrm{CVj}^2\cdot h_G^4+4\mathrm{e}^6\cdot C_d\cdot M^3\cdot K_0.
\end{gather*}
They may depend on the constants of (C0)--(C6), not on $\bar{m}$, $\beta$, $e$, $N$.
In \cite[Section~20.1.7]{QEDI} one has $C_K=100\,\mathbf{S}^2$, $C_K'=400\,\mathbf{S}^2$, $C_H=100\,M\cdot C_F^2$, and the constant $c_1$ there has $\frac{1}{80\,\mathbf{S}}$ in place of $\frac{1}{292\,\mathbf{S}}$;
the constants $\mathrm{gGa}$, $\mathrm{gGa}_1$, $\mathrm{alG}$, $\mathrm{CVj}$, $\mathrm{CWj}$ are written there $g$, $g_1$, $\boldsymbol{\alpha}$, $C_V$, $C_W$.

\item[(C7)]
Integers $r\ge2$, $p_0>3r+2$, $p>2p_0$;
$\varepsilon\in\left(1-\alpha,\bar{\varepsilon}\right]$;
$s\in(0,1)$ with
\[
\sigma:=s\cdot\left(\frac{1}{4}-9\varepsilon\right)-12\varepsilon-\frac{2a_L}{N_{*}}>0.
\]
They may depend on $\alpha$ and $L$, not on $\rho$.
These are the choices of stage~8 of \cite[Definition~5.97]{QEDI}.

\item[(C8)]
$e_0>0$, subject to all the conditions of Section~I.8.4 that contain $e_0$.
It may depend on the constants of (C0)--(C7).
This is the choice of stage~9 of \cite[Definition~5.97]{QEDI}, except that (A6) carries the number $112$ and that the inequality $e_0\le e^{F}_0$ of the condition (A10) of \cite[Section~5.6.2]{QEDI} is not imposed, because the number $e^{F}_0$ does not occur here and (Z2.s1)--(Z2.s5) are imposed on $e_0$ itself.
The four inequalities that this condition (A10) imposes for every $e_k\in(0,e_0]$ are, with $v=v_k=-\log e_k$, the four inequalities of (A10) in Section~I.8.4, and its last inequality, $8\,c_{\mathrm{rb}}\cdot L^3\cdot C_d\cdot M^3\cdot e_0^{2\varepsilon}\le1$, is (Z2.s5), formed with the constant $C_d$ of the line (C6).

\item[(C9)]
$e_1\in\left(0,\min\left(e_0,\frac{1}{3}\right)\right]$, the number defined by the conditions (e1)--(e8) of \cite[Section~19.1.2]{QEDI} with the constants above;
the definition is that of Proposition~14.2\,(a), line (C9), and (b), rows (e1)--(e8).
It may depend on the constants of (C0)--(C8).
The definition is that of \cite{QEDI}; the number $e_1$ does not depend on $\bar{m}$.

\item[(C10)]
$\bar{m}\in(0,1]$, hence $h$, and
\[
e_G:=\sup\left\{e:\ e\cdot h^{\frac{1}{2}}\le e_1,\ e\cdot h^{\frac{1}{2}}\cdot\log h\le1\right\};
\]
the probe $\beta$ and the point $x_0$, with (Pb) of Proposition~14.2\,(d) (the three probes $\beta_{\mathrm{ix}}$, namely $\beta$ and $\beta\pm\tau_{x_0}\beta$);
then $D_{\mathrm{be}}$, $D_\beta=D_{\mathrm{be}}+4$, $C_{\beta}$, $\mathbf{D}$, $\mathrm{Cb}$;
$\mathrm{Cah}_i$, $\mathrm{Cpv}_1$, $\mathrm{Cpv}_2$, $\mathrm{epv}$, $\mathrm{Cpv}_\beta$ (Proposition~36\,(e), (f), (g));
$C_{\mathrm{ct}}$, $C_{1}^{\mathrm{N2}}$, $C_{2}^{\mathrm{N2}}$ (Proposition~37.2\,(c), (d), (e)), where $C_1^{\mathrm{N2}}=C_1^{\mathrm{N2}}(\rho,\bar{m},\beta)$ is the constant at $\eta^{\frac{1}{4}}$ and $C_2^{\mathrm{N2}}=C_2^{\mathrm{N2}}(\bar{m},\beta)$ the constant at $\exp\left(-\bar{m}\cdot\left(R-D_{\mathrm{be}}\right)\right)$;
$C_{\mathrm{fr}}$;
then, in this order, $e_\sharp(\beta')$ for the three probes (Proposition~35.10) and $C^{G}_{\beta'}$ (Proposition~43.1);
$C_{\mathrm{mat}}:=\left(L^3+1\right)\cdot\left(\varepsilon_5\left(\frac{1}{L}\right)+1\right)$;
$s_{B}$, $C_{3\mathrm{p}}$, $C_{\Gamma}^{(16)}$, $C_{U}$,
\[
C_{\bar{\kappa}}:=2\cdot\left(8C_U+3C_U^2\right),\qquad C'_{\bar{\kappa}}:=8C_{\bar{\kappa}}+3C_{\bar{\kappa}}^2
\]
(Proposition~43.2);
$e_2$.
They may depend on everything above.
The line (C10) takes the place of stage~11 of \cite[Definition~5.97]{QEDI}.
Its constants are defined at the places named in it and in Proposition~14.2\,(a), line (C10).
They differ from those of stage~11 of \cite[Definition~5.97]{QEDI} in the following points.
\begin{itemize}
\item The constant $C_Y$ of \cite{QEDI} does not occur, and the condition on $e_2$ contains $C'_{\bar{\kappa}}$ (Section~I.8.4).
\item The constants of \cite[Section~20.1.7]{QEDI} which contain neither $\bar{m}$ nor $\beta$ belong to the line (C6$'$).
\item The constants $D_{\mathrm{be}}$, $\mathrm{Cah}_i$, $\mathrm{Cpv}_1$, $\mathrm{Cpv}_2$, $\mathrm{epv}$, $\mathrm{Cpv}_\beta$, $C_{\mathrm{ct}}$, $C_U$ and $C'_{\bar{\kappa}}$ have no counterpart in \cite{QEDI}.
\item The constants $\mathrm{Cb}$, $C_{\mathrm{fr}}$, $C_{\mathrm{mat}}$, $s_B$ and $C_{3\mathrm{p}}$ are defined in \cite{QEDI} outside stage~11:
$\mathrm{Cb}$ in \cite[Theorem~5.92]{QEDI}, where it is written $\bar{C}=1+h'^{\frac{1}{2}}\cdot C_\beta$, with $h'$ in place of $h$;
$C_{\mathrm{fr}}$ in \cite[Section~20.1.7]{QEDI}, $C_{\mathrm{mat}}$ in \cite[Theorem~5.95]{QEDI} and \cite[Theorem~21.3]{QEDI}, and $s_B$ and $C_{3\mathrm{p}}$, the latter written $C_B$ there, in \cite[Section~21.10.1]{QEDI}, the proof of Part~II of \cite[Theorem~3.1]{QEDI}.
\item The definitions of the following constants differ from those of \cite{QEDI}:
$C_1^{\mathrm{N2}}$ and $C_2^{\mathrm{N2}}$, which take the place of the constants $C_1$, $C_2$ of stage~11, the first being the constant at $\eta^{\frac{1}{4}}$ and depending on $\rho$, the second being formed with $D_{\mathrm{be}}$;
$e_\sharp$, which is formed with (g1)--(g9);
and $C^{G}_{\beta}$ and $C_{\bar{\kappa}}$, which are given by explicit formulas.
Moreover $h$ takes the place of $h'$.
\end{itemize}

\item[(C11)]
$e$; then $T$, that is, the torus; then $N$.
\end{description}

\paragraph{The chain that carries $\rho$.}
The constants that depend on $\rho$ are linked by the chain
\[
\rho\to\text{(C0r)}\to\left(c_{\mathrm{re}},C_{00},a',\gamma^F,\gamma_{\mathrm{w}}\right)\to M\to\left(C^F,C_w,C_W,C_d,C_0\right)\to e_0\to e_1\to\left(e_\sharp,e_2\right).
\]
Beside this chain one has $\rho\to\left(\gamma_\Sigma,C_\Sigma\right)\to e_1$, where $\gamma_\Sigma$ and $C_\Sigma$ are functions of $(b_-,b_+,\rho,w)$ and do not depend on $M$;
and $\text{(C0r)}\to\left(C_{\mathrm{mat}},C^G_{\beta'}\right)\to C_{\bar{\kappa}}\to e_2$, where $C_{\mathrm{mat}}$ is a function of $(L,\varepsilon_5)$ and none of $C_{\mathrm{mat}}$, $C^G_{\beta'}$, $C_{\bar{\kappa}}$ depends on $M$, $e_0$, $e_1$.
The numbers $\alpha$, $\kappa$, $L$, $N_*$, $a_L$, $\varepsilon$, $s$, $\sigma$, all the exponents in the first nine rows of the table of Section~I.8.5, and the rates $\kappa_{\mathrm{w}}+\kappa_0$, $\kappa_{\mathrm{w}}$, $\kappa$ do not depend on $\rho$;
$\gamma_{\mathrm{w}}$ and $\ell_H=M$, in the last row of that table, do.
No condition bounds $M$ from above.
The conditions (T2)--(T5) hold by the rows (e3)--(e6) of the definition of $e_1$, and (W3) holds by the row (e8) together with (Z8.w), which is contained in $e\le e_G$ (line (C10)); all of these come after $M$.
Every constant, and every condition on it, involves only constants that precede it in the list (Proposition~14.2\,(e)).

\subsubsection*{I.8.3. Numerals that come from the rows of the operator}

The following numerals of \cite{QEDI} hold also for the chamber fermion, because one hop of norm at most $6$ takes the place of six hops of norm one ((I.O6)), or because they count block spinors of levels $\ge1$:
$12\mathrm{e}$ and $2\mathrm{e}$ in \cite[Proposition~5.65]{QEDI}, and $144\mathrm{e}$ in \cite[Theorem~13.18]{QEDI};
$12$ in \cite[Lemma~20.18(b)]{QEDI};
$48$ in (Z9.s3);
$8$ in $|\partial_BD_A|(z,z)$;
$12M^3$ in \cite[Proposition~17.16(b)]{QEDI};
$6+3b$ in $a_z$;
$c_{k}$, $c'_{k}$, $c^{\mathrm{d}}_{k}$;
the numbers $2$, $18$, $19$, $16$, $152$, $32$, $304$, $27$, $8$, $54\pi$, $108\pi$, $1600$ of \cite[Sections~16.5--16.7]{QEDI};
$\frac{1}{54}$;
$L^F=5$.

For the chamber fermion other numerals hold (the Propositions~33.m).
For it the majorants of \cite[Lemma~20.14(b)]{QEDI} are $\mathfrak{d}(z,z)=15$, as in \cite{QEDI}, and $\mathfrak{d}(z,z')=8$, where \cite{QEDI} has $\frac{3}{2}$, and the sums over rows and over columns are at most $146$, where \cite{QEDI} has $40$.
Consequently, in the bounds of \cite[Section~20.4]{QEDI} that rest on these majorants and in the constants of \cite[Section~20.1.7]{QEDI}, the following expressions hold for the chamber fermion;
after each of them the corresponding expression of \cite{QEDI} stands in brackets:
\begin{itemize}
\item $149\left(C^F\right)^2$ \quad $\left[43\,C_F^2\right]$;
\item $160\left(C^F\right)^2$ and $203\,M\cdot\left(C^F\right)^2$ \quad $\left[54\,C_F^2\text{ and }62\,M\cdot C_F^2\right]$;
\item $21645\left(C^F\right)^3$ \quad $\left[1690\,C_F^3\right]$;
\item $226\left(C^F\right)^2$ \quad $\left[67\,C_F^2\right]$;
\item $13.5\,\mathbf{S}+1314\,\mathbf{S}^2$ \quad $\left[13.5\,\mathbf{S}+360\,\mathbf{S}^2\right]$;
\item $\frac{1}{292\,\mathbf{S}}$ \quad $\left[\frac{1}{80\,\mathbf{S}}\right]$;
\end{itemize}
and
\[
C_K:=226\,\mathbf{S}^2,\qquad C_K':=1328\,\mathbf{S}^2,\qquad C_H:=203\,M\cdot\left(C^F\right)^2.
\]
The constant $c_1=\min\left\{\frac{1}{8},\frac{\mu_-}{48},\frac{1}{80\,\mathbf{S}}\right\}$ of \cite[Section~20.1.7]{QEDI} is replaced by
\[
c_\tau:=\min\left\{\frac{1}{8},\frac{\mu_-}{48},\frac{1}{292\,\mathbf{S}}\right\},\qquad\tau_s:=c_\tau\cdot\ell_*.
\]
Of the constants of \cite[Section~20.1.7]{QEDI}, those written there $\boldsymbol{\theta}_{\mathrm{CT}}$, $\vartheta$, $g$, $h_G$, $g_1$, $\boldsymbol{\alpha}$, $c_1$, $C_V$, $C_W$ (here $\theta_{\mathrm{CT}}$, $\mathrm{vt}$, $\mathrm{gGa}$, $h_G$, $\mathrm{gGa}_1$, $\mathrm{alG}$, $c_\tau$, $\mathrm{CVj}$, $\mathrm{CWj}$, line (C6$'$)) do not depend on $e$;
the quantities $\tau_s$, $\bar{r}$, $B_1$, $B_0$ and $\varpi_{\mathrm{sep}}$ do.

\subsubsection*{I.8.4. Conditions}

The following conditions are those of \cite[Section~5.6.2]{QEDI}, formed with the constants of Section~I.8.2, with $h$ in place of $h'$ and with $0$ in place of $m_0$:
the condition (G4.2.1), that is, $r_*>31L+16$;
(A0)--(A9);
(fl);
(T1)--(T6);
(W3);
(Z8.w);
(g1)--(g9), in the form of Proposition~14.2\,(d);
(Z2.m);
(Z2.s1)--(Z2.s5);
(Z3.s1);
(Z8.s1);
(Z9.s1)--(Z9.s3);
(Z10.s1), (Z10.s2), and (Z10.s4), the last one with $C'_{\bar{\kappa}}$.
The conditions (A0)--(A10), (fl) and (T1)--(T6) bear these names in \cite[Section~5.6.2]{QEDI}, and (g1)--(g8) are the conditions (sm1)--(sm8) of \cite[Section~20.1.8]{QEDI}.
The remaining names denote displays of \cite[Section~5.6.2]{QEDI} that carry no name there:
(G4.2.1) is the first condition on $(R,e,N)$;
(Z2.m) is the condition $M\ge M^F$;
(Z2.s1)--(Z2.s5) are, in this order, the five conditions on $e_0$ that follow it;
(Z3.s1) is the condition $e_0\le\exp(-1)$;
(Z8.s1) is the condition that defines $e_1$;
(W3) is the condition $6N_*\cdot e^2\cdot R^{*}\cdot\log\frac{1}{e}\le\pi^2$, and (Z8.w) the condition $e_K\cdot\log h'\le1$;
(Z9.s1)--(Z9.s3) are, in this order, the three conditions which are described there as parts of (sm1), (sm3) and (sm4);
(Z10.s1)--(Z10.s4) are, in this order, the four conditions on $e_2$, so that (Z10.s3) is the condition
\[
\frac{C_Y+1}{6\log\left(e_2^{-1}\right)}+\frac{\mathfrak{e}_2}{3}\le1.
\]
The conditions imposed here are stated, in their present form, in Proposition~14.2\,(a)--(d).
They differ from the conditions of \cite{QEDI} in the following points.
\begin{itemize}
\item[(A6)]
$112\,C_w^2\cdot c_\gamma\cdot e_k^{2\varepsilon}\le3$, with $112$ in place of the number $16$ of \cite{QEDI} (see (I.L0.19)).
\item
The conditions (Z2.m) and (Z2.s2) are formed with $C'_V$ and $a'$, and (Z2.s5) with the constant $C_d$ of the line (C6).
\item
The condition (T5) and the error term $\varepsilon_2$ are formed with $\gamma_\Sigma$.
\item
The clause $w^{\mathrm{m}}_0<1$ of (T1) is kept; every real lattice mass is admitted (I.K10).
\item[(Z10.s4)]
\[
\frac{5}{2}\,C'_{\bar{\kappa}}\cdot e_2^{\sigma}\le\frac{b(f)}{4},
\]
with $C'_{\bar{\kappa}}$ in place of the constant $C_{\bar{\kappa}}$, which is written $C_{\mathrm{kap}}$ in \cite[Section~5.6.2]{QEDI}.
The condition (Z10.s3) is not imposed, and the constant $C_Y$ of \cite[Section~5.6.2]{QEDI} does not occur.
\item
The conditions (g2), (g5), (g8), (g9) have the forms of Proposition~14.2\,(d), that is, of Section~G5.1.9.
Unlike (sm2) and (sm5) of \cite[Section~20.1.8]{QEDI}, the conditions (g2) and (g5) involve no placement;
the label (g9) does not occur in \cite{QEDI}.
\item[(Pb)]
This is a condition on the choice of $(\beta,x_0)$ and of the radius $r_\beta$ of a ball that contains $\operatorname{supp}\beta$:
$b(f)>0$, $d_0>0$, and the inequality of \cite[Section~2.8]{QEDI}, in the form of Proposition~14.2\,(d).
With $a=d^{*}\beta$, $\|a\|_\infty:=\sup_y|a(y)|_2$ and $d_0:=|x_0|_2-2r_\beta$, this inequality reads
\[
\frac{6}{\bar{m}^2}\cdot\|a\|_\infty^4\cdot\left(\frac{4\pi\cdot r_\beta^3}{3}\right)^2\cdot\left(\bar{m}^2+\left(\bar{m}+\frac{1}{d_0}\right)^2\right)\cdot\frac{\exp\left(-2\bar{m}\cdot d_0\right)}{4\pi^2\cdot d_0^2}\le\frac{b(f)}{12}.
\]
For every $\bar{m}$ there are admissible pairs $(\beta,x_0)$ (Proposition~14.3\,(g)).
\item[(A10)]
Here the condition (A10) consists of the following four inequalities, imposed for all $v\ge-\log e_0$ (Proposition~14.2\,(b)):
\begin{gather*}
v\ge9,\qquad8M\cdot\left(\frac{v}{4}+4\right)\cdot\exp(-\varepsilon v)\le1,\\
8\,C_B\cdot v^p\le\exp\left(\left(\frac{3}{4}-3\varepsilon\right)v\right),\qquad8\exp\left(-\left(\frac{3}{4}-\frac{7\varepsilon}{2}\right)v\right)\le1.
\end{gather*}
With $v=v_k=-\log e_k$ and $p_k=v_k^p$, these are exactly the four inequalities that the condition (A10) of \cite[Section~5.6.2]{QEDI} imposes for every $e_k\in(0,e_0]$;
in this point the two conditions agree, and (A10) is there, as here, a condition on the constants alone.
In \cite{QEDI} the first two of the four inequalities are the second and the third inequality of \cite[Proposition~5.48(d4)]{QEDI}, and the last two are the two displayed inequalities of \cite[Proposition~5.38]{QEDI}.
The requirement $\ell\ge M\cdot r_k$ of \cite[Proposition~5.48(d4)]{QEDI}, the requirements $\ell/L\in\mathrm{Cl}$ and $L^{\frac{1}{2}}\cdot e_k\le e_0$ of \cite[Proposition~5.48(d5), (d6)]{QEDI}, and the requirement $\ell/L\in\mathrm{Cl}$ of \cite[Proposition~5.38(ii)]{QEDI} are requirements on the data $(k,\ell,e_k)$;
they belong to the hypothesis for one step, \cite[Definition~5.69]{QEDI}, and they are no part of (A10), there as here.
The condition (A10) of \cite{QEDI} contains two further inequalities, $e_0\le e^{F}_0$ and $8\,c_{\mathrm{rb}}\cdot L^3\cdot C_d\cdot M^3\cdot e_0^{2\varepsilon}\le1$.
The first is not imposed here: the number $e^{F}_0$ of \cite{QEDI} does not occur, and (Z2.s1)--(Z2.s5) are conditions on $e_0$ itself.
The second is (Z2.s5), formed with the constant $C_d$ of Proposition~13.
\item
The error terms are
\[
\varepsilon_1:=\frac{1}{2}e_*^{\frac{1}{4}-9\varepsilon},\qquad\varepsilon_2:=C_\Sigma\cdot\ell_*^{-1}\cdot\exp\left(-\gamma_\Sigma\cdot\ell_*\right),
\]
the numbers $\varepsilon_3$, $\varepsilon_4$ of \cite[Section~5.6.2]{QEDI}, and the function $\varepsilon_5$ of Section~I.8.2.
\end{itemize}
The enumeration of the conditions on the constants in the proof of \cite[Theorem~17.2(h)]{QEDI}, in \cite[Section~17.13.3]{QEDI}, names among others the conditions $0<\alpha$ and $1-\varepsilon<\alpha$, the first inequality of stage~4 of \cite[Definition~5.97]{QEDI}, which is (L1), the conditions $M\ge M^{F}$ and $M\ge M^{P}$, which are (M5) and (M6), the condition $r\ge2$, and the relations $C_w\ge C_P$ and $\gamma_{\mathrm{w}}\le\gamma_P$.
For the constants of Section~I.8.2 all of them hold (Proposition~14.2\,(a), (b), (f));
here, as in (Z2.m), the number $M^{F}$ is formed with $C'_V$ and $a'$ in place of $C_V$ and $a$.

\subsubsection*{I.8.5. Fixed exponents}

The following exponents are fixed.
Except in the rows $\mathfrak{e}_2$ and ``rate in the spacing'', they are the exponents of \cite{QEDI}.
The bound on the fluctuation field, namely that each of the three norms of $\mathscr{Z}(s')$ is at most $\frac{1}{2}e_k^{-\frac{\varepsilon}{2}}$ on $\mathrm{Pd}_k\times\mathrm{Box}_{k}$ (compare \cite[Proposition~5.38(iii)]{QEDI}), is that of Proposition~24.3\,(q1).
The bounds (F1)--(F4) of the row ``one step'' are those of \cite[Proposition~5.59]{QEDI}, \cite[Proposition~5.62]{QEDI}, \cite[Proposition~5.48(d6)]{QEDI} and \cite[Theorem~17.2(a), (b)]{QEDI}, respectively:
(F4) is the bound of \cite[Theorem~17.2(a)]{QEDI} in the case with the slot and of \cite[Theorem~17.2(b)]{QEDI} in the case without it, and the bound $\delta\varepsilon^{0}_{k}\le e_k^7$ of the same row is part of \cite[Theorem~17.2(a)]{QEDI}.

\noindent
\begin{tabular}{p{0.27\textwidth}p{0.66\textwidth}}
object & value\\
\hline
thresholds of $R_k$; $G_k$; $\mathfrak{X}_k$ &
$e_k^{-\frac{3}{4}+\varepsilon}$, $e_k^{-\frac{3}{4}+2\varepsilon}$, $e_k^{-\frac{3}{4}+3\varepsilon}$; $e_k^{-1+\varepsilon}$; $e_k^{-1+2\varepsilon}$, $e_k^{-1+3\varepsilon}$\\
weights $h_{k}$; windows &
$h_k=\left(e_k^{-\frac{1}{4}},e_k^{-\frac{1}{4}+\varepsilon}\right)$; $w^{\mathrm{m}}_k=e_k^{\frac{3}{4}-8\varepsilon}$, $w^{\mathrm{E}}_k=e_k^{\frac{1}{4}-7\varepsilon}$\\
parameter polydisk; contours &
$M^{\alpha_0}$; $\mathrm{rad}_1=\frac{1}{2}e_k^{-\frac{3}{4}+\frac{9\varepsilon}{2}}$, $\mathrm{rad}_2=e_k^{-\frac{3}{4}+7\varepsilon}$, $\mathrm{rad}_4=e_k^{-\frac{1}{4}+3\varepsilon}$\\
one step &
(F1): $c_{\mathrm{fl}}$, $c_{\mathrm{fl}}\cdot e_k^{\frac{1}{2}}$; (F2): $c_{\mathrm{fl}}\cdot L^{-\frac{1}{4}+2\varepsilon}$; (F3): $e_{k+1}^{\frac{1}{4}-2\varepsilon}$; (F4): $e_k^{\frac{1}{4}-6\varepsilon}$; $\delta\varepsilon^{0}_{k}\le e_k^7$\\
trajectory &
the norms of $E_k$ and $\mathscr{E}_k$ are at most $4e_k^{\varepsilon}\cdot w^{\mathrm{E}}_k$; $|m_k|\le\frac{1}{8}w^{\mathrm{m}}_k$; $\left|\varepsilon^{\mathrm{v},t}_k\right|\le5\,c_{\mathrm{fl}}\cdot e_k^{\varepsilon}\cdot w^{\mathrm{E}}_k$\\
terminal &
$\varepsilon_1=\frac{1}{2}e_*^{\frac{1}{4}-9\varepsilon}$; in (10.1) of Proposition~32.3: $3e_*^{\frac{1}{4}-9\varepsilon}$; $\bar{\kappa}_1=6e_*^{\frac{1}{4}-9\varepsilon}$; $t_x$ proportional to $e_*^{3\varepsilon}$; $\mathbf{r}=\frac{t_x}{4}$\\
final &
$\mathfrak{e}:=s\cdot\left(\frac{1}{4}-9\varepsilon\right)-12\varepsilon$; $\sigma:=\mathfrak{e}-\frac{2a_L}{N_*}$\\
$\mathfrak{e}_2$ &
$\mathfrak{e}_2:=0$ (Proposition~36\,(g), Proposition~37.3\,(b)): the bound on $c_2$ is uniform in $(\eta,R)$ and has no rate. In \cite[Section~2.7]{QEDI} one has $\mathfrak{e}_2=6\varepsilon$.\\
rate in the spacing &
$\eta^{\frac{1}{4}}$, in place of the rate $\eta^{\frac{1}{2}}$ of \cite[Theorem~21.1(c2)]{QEDI}, \cite[Theorem~21.3, part~2\,(c)]{QEDI} and \cite[Theorem~3.1, Part~II\,(ii)]{QEDI}. This rate occurs at finite cutoff only; it does not affect the limit points.\\
rates &
$\gamma_{\mathrm{w}}$ per unit length; $\kappa_{\mathrm{w}}+\kappa_0$ (for $\mathrm{par}^0$), $\kappa_{\mathrm{w}}$ (for $\mathrm{par}$) and $\kappa$ (for the activities) per unit of $d_M$; $\ell_H=M$\\
\end{tabular}

\subsubsection*{I.8.6. Lower bounds that are not uniform in the spacing; the sign of the mass}

The lower bounds on partition functions in \cite[Theorem~4.13(c)]{QEDII} and \cite[Proposition~6.19(c)]{QEDII} are not uniform in the spacing; in \cite{QEDII} they serve \cite[Theorem~6.20]{QEDII} at fixed parameters.
The bounds on propagators and the constants of Section~I.8 do not depend on the sign of the lattice mass $\boldsymbol{\mu}$, and every real $\boldsymbol{\mu}$ is admitted (I.K10).

\subsection*{I.9. The statements of the two companion papers}

\subsubsection*{I.9.1. Constants, classical facts, continuum objects, names of statements of \cite{QEDI}}

\paragraph{Constants.}
The constants that occur in the statements of this paper are those of Section~I.8.
The conditions imposed on them there are part of the choice of the constants,
and they hold for the constants as chosen.

\paragraph{Classical facts.}
The labels (cl1), (cl2), (cl3) and (cl5) are those under which \cite[Section~13]{QEDI} displays the following classical facts,
and they denote these facts.
\begin{itemize}
\item[(cl1)] Cauchy's integral formula and the residue theorem for holomorphic functions of one complex variable,
applied to each component of a function with values in a finite-dimensional complex vector space
\cite[Chapter~4, Sections~2.2, 5.1]{Ahlfors1979}.
\item[(cl2)] The identity theorem for holomorphic functions of several complex variables
\cite[Chapter~I, Section~A]{GunningRossi1965}.
\item[(cl3)] The two formulas for the determinant of a block matrix with square diagonal blocks $A_1$, $A_4$
\cite[Section~0.8.5]{HornJohnson2013}:
\begin{align*}
\det\left(\begin{array}{cc}A_1&A_2\\A_3&A_4\end{array}\right)&=\det A_1\cdot\det\left(A_4-A_3\cdot A_1^{-1}\cdot A_2\right)&&\text{if $A_1$ is invertible},\\
\det\left(\begin{array}{cc}A_1&A_2\\A_3&A_4\end{array}\right)&=\det A_4\cdot\det\left(A_1-A_2\cdot A_4^{-1}\cdot A_3\right)&&\text{if $A_4$ is invertible}.
\end{align*}
\item[(cl5)] The Leibniz formula in a Grassmann algebra;
it is proved in Chapter~15, Section~0.4.
\end{itemize}

\paragraph{Continuum objects.}
Consider a statement of \cite{QEDI} that concerns the continuum Dirac operator alone and no lattice,
that is, one whose objects are the matrices $\gamma_\mu$, the propagator $S_{\bar{m}}$ and the coefficient $b(\cdot)$,
together with the operators $S_{\mathrm{Mq}}$ and the numbers $q_{4,\mathrm{Mq}}$, $\mathrm{Mq}>0$, of \cite[Section~21.1.9]{QEDI}, where the mass $\mathrm{Mq}$ is written $\mathsf{M}$.
Such a statement holds as stated for the objects of this paper.
Indeed, Sections~I.2.1, I.2.6 and I.7.6 define the same $\gamma_\mu$, $S_{\bar{m}}$ and $b(\cdot)$,
and by Propositions~37.1\,(f) and 37.2\,(f) these are the continuum data of the chamber operator.

\paragraph{Names of statements of \cite{QEDI}.}
In this paper the following names denote the listed statements of \cite{QEDI}, all of which stand in \cite[Section~5]{QEDI}:
\begin{itemize}
\item W1--W18: the properties (1)--(18) of a system of primitives in \cite[Section~5.4.2]{QEDI}, in this order;
\item (Lo): \cite[Proposition~5.30]{QEDI};
\item PH1--PH11: \cite[Propositions~5.31--5.41]{QEDI};
\item FE1--FE8: \cite[Propositions~5.42--5.49]{QEDI};
\item ID1a, ID1b, ID2, ID3, ID4: \cite[Propositions~5.50--5.54]{QEDI};
\item AC1, AC2: \cite[Propositions~5.57, 5.58]{QEDI}; AC7, AC8, AC9: \cite[Propositions~5.64, 5.65, 5.66]{QEDI};
\item DF1--DF5: \cite[Lemmas~5.17--5.21]{QEDI};
\item Z1, Z2: \cite[Theorems~5.71, 5.72]{QEDI}.
\end{itemize}
For a property that involves fermion objects the superscript $\mathrm{ch}$,
as in FE4$^{\mathrm{ch}}$, W7$^{\mathrm{ch}}$, AC8$^{\mathrm{ch}}$ and Z1$^{\mathrm{ch}}$,
denotes the property formulated for the chamber operator in the conventions of Part~I;
it is a clause of the proposition that Section~I.9.4 assigns to the property.

The properties PH1--PH11, DF1--DF4, ID1a--ID4 and AC2 carry no superscript.
The objects in them that depend on the fermion, that is, on $\rho$, are the following:
\begin{itemize}
\item the group $\mathbb{G}_k$, which enters through $r(S)$ and $a$;
\item the numbers $C_{w}$, $\gamma_{\mathrm{w}}$ and $c_{\gamma}$ in the condition (A2) of \cite[Section~5.6.2]{QEDI};
\item the number $M$, which enters through the class $\mathrm{Cl}$, through the $M$-cubes and the enlarged cubes in the domains of \cite[Definition~5.16]{QEDI} and in the polymers of AC2,
through $s_d=(2L+1)\cdot M$ and $C_{B}$ in PH7 and PH8, through the inequality $3M\cdot e_k^{\varepsilon}\le1$ of the condition (A1),
and through the stop level $k_{*}$ of \cite[Definition~5.6]{QEDI}, written $k^{*}$ there, which is formed with $M$ and occurs in ID3 and ID4
(for these two properties see the paragraph on $k_{*}(R)$, $\mathrm{Tor}(e)$ and $N_{\min}$ below);
\item the number $e_0$, which enters through the range $0<e_k\le e_0$ of the charge,
in DF1--DF3 and in the scale-free setting at the head of \cite[Section~5.4]{QEDI}, in which PH1--PH11 are stated,
and through the conditions (A9) in PH5.
\end{itemize}
This group and these numbers are those of Sections~I.4.2 and~I.8.2, and (A1), (A2), (A8), (A9) and (A10) are conditions of Section~I.8.4.
In \cite{QEDI} these properties are proved for free values of the parameters;
for the values of Section~I.8.2 they hold in the following forms.
\begin{itemize}
\item PH1--PH11 hold in the form of \cite[Theorem~15.37]{QEDI}, that is,
for every $\alpha\in\left(1-\bar{\varepsilon},1\right)$, every odd $L\ge5$, every $\alpha_{0}\in\left(0,\frac{1}{4}\right]$, every $M=L^m\ge M^{P}(\alpha,L,\alpha_{0})$, every $k\ge0$ and every $\ell\in\mathrm{Cl}$;
of the number $e_0$ only the inequality $e_0\le\exp(-1)$, which is (Z3.s1) of Section~I.8.4, enters \cite[Section~15]{QEDI};
and PH8 holds, for every $\varepsilon\in\left(0,\bar{\varepsilon}\right]$, under the conditions stated in it, namely (A1), (A2) and the two inequalities displayed there,
which are the third and the fourth inequality of (A10) in Section~I.8.4 at $v=-\log e_k$.
\item ID1a--ID4 hold in the form of \cite[Sections~14.1--14.5]{QEDI}, that is, for every odd $L\ge3$, every $k\ge0$ and every $\ell\in L\mathbb{N}$,
and in \cite[Section~14.5]{QEDI} for every $k\ge1$ and every $n\in L^k\mathbb{N}$;
neither $M$ nor $e_0$ occurs there, except in the sentences on the case $k=k_{*}$ under \cite[Definition~5.70]{QEDI}.
\item AC2 holds in the form of \cite[Theorem~13.2]{QEDI}, that is, for every $k\ge0$ and every $\ell\in\mathrm{Cl}$,
with $\ell/L\in\mathrm{Cl}$ in addition in the clauses (t5) and (t8).
Of the numbers $L$, $M$ and $\ell$ the proof of \cite[Theorem~13.2]{QEDI} uses only the following:
$L$ is odd, in (t5) and, through (t5), in (t8);
$L\ge1$;
and $\ell\in\mathrm{Cl}$, the class being formed with the same $L$ and $M$,
that is, $\ell$ is a multiple of $L\cdot M$, in (t5) and (t10), and $\ell\ge13L\cdot M$,
whence $2M<\ell$, in (t2), (t4) and (t10), and $2M\cdot(L+2)<\ell$, in (t7) and (t10).
No condition on $M$ enters, and the numbers $c_X$, $\kappa_{0}$, $K_0$ and $c_{\mathrm{rb}}$ are absolute.
Hence AC2 holds with $\mathrm{Cl}$, the $M$-cubes, $D_k$ and $d_M$ formed with the numbers $L$ and $M$ of Section~I.8.2;
that $L$ is odd is part of the row (L5) of Proposition~14.2\,(b).
\item DF1--DF4 hold in the form of \cite[Lemmas~17.4--17.7]{QEDI}, with the data stated in the head of \cite[Section~17.3]{QEDI}:
$\ell\in\mathrm{Cl}$ and $0<e_k\le e_0$ in DF1 and DF3;
$\ell\in\mathrm{Cl}$, $\ell/L\in\mathrm{Cl}$ and $0<e_k\le e_0$ in DF2;
$\ell_1,\ell_2\in\mathrm{Cl}$ and $e_k>0$ in DF4.
Of the conditions (A1)--(A10), (Z2.s1)--(Z2.s5) and (Z3.s1) on $e_0$, the head of \cite[Section~17.3]{QEDI} refers only to (A1) and (A8), written for $e_k$,
and the proofs of \cite[Lemmas~17.4--17.7]{QEDI} use only the first and the fourth inequality of (A8) at $v=-\log e_k$, that is,
\[
e_k^{\varepsilon}\le\frac{1}{4},\qquad e_k^{\frac{1}{4}}+e_k^{\varepsilon}\le1.
\]
The other conditions on the constants that enter these proofs are the following;
they are taken from the rows (L1), (L5), (M1), (S1) of Proposition~14.2\,(b) and from the line (C1) of Section~I.8.2.
\begin{itemize}
\item $L$ is odd, and $M$ is a power of $L$ with $M>1$ ((L5), (M1)).
These properties enter \cite[Lemmas~17.4(d) and 17.5(a)]{QEDI} through the fact (g0) of \cite[Section~17.4.1]{QEDI}, whose proof uses no property of $L$ and $M$ other than these and $\ell\in\mathrm{Cl}$;
the inequality $M>1$ enters \cite[Lemma~17.4(d)]{QEDI} also directly.
\item $L\ge5$ ((L5)), in \cite[Lemma~17.4(a)]{QEDI}.
\item $L^{-\frac{3}{4}}\le\frac{1}{2}$, a consequence of (L1), in \cite[Lemma~17.5(b)]{QEDI}.
\item $0<\varepsilon\le\frac{1}{200}$ ((S1), together with $\alpha<1$ and $\bar{\varepsilon}=\frac{1}{200}$), in \cite[Lemmas~17.4(b) and 17.5(a), (c)]{QEDI}.
\item $\alpha>0$ ((C1)), in \cite[Lemma~17.5(a)]{QEDI}.
\end{itemize}
\end{itemize}
Every condition used in these forms holds for the constants of Sections~I.8.2 and~I.8.4, by Proposition~14.2\,(a), (b), (f).

The properties PH7, PH8, and W3, W4 for the families (p1), (p2) of primitives of \cite[Section~5.4.1]{QEDI}
contain $M$, $\alpha_{0}$, $C_{w}$, $\gamma_{\mathrm{w}}$, $c_{\gamma}$, $C_{B}$ and $C^{P}$.
They are taken in the parametric forms in which \cite[Section~15]{QEDI} states them:
\begin{itemize}
\item PH7 for every $M$ that is a power of $L$, with $\ell\in\mathrm{Cl}$ (\cite[Proposition~15.29]{QEDI});
\item PH8 for every $\varepsilon\in\left(0,\bar{\varepsilon}\right]$, every $L$ with $L^{\frac{3}{4}}\ge\frac{8}{3}$, every $e_k\in(0,1)$, every $C_{w}\ge C_{P}$ and every $\gamma_{\mathrm{w}}\le\gamma_{P}$, with $c_{\gamma}$ formed with $\gamma_{\mathrm{w}}$,
whenever (A1), (A2) and the two inequalities displayed in PH8 hold at $e_k$ (\cite[Proposition~15.30]{QEDI});
\item W3 and W4 for every $\alpha_{0}\le\frac{1}{4}$ and every $M=L^m\ge M^{P}$, with the parameters $\left(\alpha_{0},\gamma^{P},\cdot\,;C^{P}\right)$ (\cite[Theorems~15.35 and~15.37]{QEDI}).
\end{itemize}
As W3 and W4 are bounds of the form $C\cdot\exp(-\gamma\cdot d)$ with distances $d\ge0$, they hold then for every rate $\gamma\le\gamma^{P}$ and every constant $C\ge C^{P}$;
in particular they hold for $\mathrm{par}^{0}$ and for $\mathrm{par}$, since $\gamma_{\mathrm{w}}\le\gamma^{P}$ and $C^{P}\le C_{w}^{(1)}\le C_{w}$ by Proposition~14.2\,(a), (f).

In \cite{QEDI} the hypothesis for one step, \cite[Definition~5.69]{QEDI}, is attached to the assertion of DF3 on $\mathrm{Dom}_k$,
and, in its form without primitives, to ID1b and ID2 (\cite[Theorem~5.78(b)]{QEDI}) and to the clauses (t5), (t8) of AC2 (\cite[Theorem~5.76(a), (b)]{QEDI}).
These are taken without that hypothesis, as follows.
\begin{itemize}
\item The assertion of DF3 on $\mathrm{Dom}_k$ holds for every $k\ge0$, every $\ell\in\mathrm{Cl}$ and every $e_k$ with $0<e_k\le e_0$.
Indeed, let $A\in\mathrm{Dom}_k$ and $\tilde{g}\in\mathbb{G}_k$; the bond functions $\tilde{g}A$ and $-A$ are real.
By PH6 and PH1 the photon minimiser $H_k$ commutes with $\tilde{g}$ and is linear;
hence $H_k\tilde{g}A=\tilde{g}H_kA$, and $H_k$ maps $-A$ to $-H_kA$.
By the first assertion of DF3 the set $\frac{1}{4}\tilde{R}_k$ is invariant under $\mathbb{G}_k$ and under $B\mapsto-B$;
hence $\tilde{g}A$ and $-A$ lie in $\mathrm{Dom}_k$.
Finally $\overline{A}=A$, since $A$ is real.
\item ID1b and ID2 hold in the form of \cite[Proposition~14.14]{QEDI} and \cite[Theorem~14.15]{QEDI}, that is, for every odd $L\ge3$, every $k\ge0$ and every $\ell\in L\mathbb{N}$.
\item The clauses (t5) and (t8) of AC2 hold as in \cite[Theorem~13.2]{QEDI}, that is, for every $k\ge0$ and all $\ell$ with $\ell\in\mathrm{Cl}$ and $\ell/L\in\mathrm{Cl}$.
\end{itemize}

The quantities $k_{*}(R)$, $\mathrm{Tor}(e)$ and $N_{\min}$ that enter the hypothesis for the flow, \cite[Definition~5.70]{QEDI},
are formed in \cite{QEDI} with the numbers $M$, $r$ and $e_1$ of the Wilson operator.
For this reason ID3 and ID4 are taken without the restriction to this hypothesis that \cite[Theorem~5.78(c)]{QEDI} attaches to them,
ID4 and \cite[Corollary~14.18]{QEDI} are taken without their last sentences,
which begin with ``In particular'' and concern the twist $c_{\theta}$ and the constant bond function $H_{\theta}$,
and \cite[Section~14.5]{QEDI} is taken without its sentences on the case $k=k_{*}$ under this hypothesis.
The role of the hypothesis there is played by the relations
\[
k_{*}\ge1,\qquad n=L^{k_{*}}\cdot\ell_{*}\quad\text{((Sc1) of Section~G4.0.2, Proposition~26.6)},\qquad L^{-\frac{k_{*}}{2}}\cdot c_{\theta}=H_{\theta};
\]
the last one follows from Sections~I.6.1 and I.6.6, since $e_{*}\cdot\ell_{*}=L^{-\frac{k_{*}}{2}}\cdot q\cdot n$.

The property DF5, that is, \cite[Lemma~5.21]{QEDI} in the form of \cite[Lemma~20.7]{QEDI},
is the clause (DF5) of Proposition~35.4.

\subsubsection*{I.9.2. The statements of \cite{QEDI}}

\paragraph{Statements that hold as stated.}
Each statement of \cite{QEDI} in the following list contains no fermion object,
or is a statement on single laws,
or concerns the continuum Dirac operator only.
It holds for the objects of this paper as it is stated in \cite{QEDI},
with the modifications of Section~I.9.5.
\begin{itemize}
\item In \cite[Section~7]{QEDI}:
the definitions of \cite[Section~7.1.1]{QEDI} (the lattice, its cells and cochains, the pairings, $d_{\eta}$ and $d_{\eta}^{*}$), for every integer $n=R/\eta$;
\cite[Section~7.1.3]{QEDI};
the first paragraph of \cite[Section~7.1.4]{QEDI}, with $\hat{\mathsf{R}}_{\nu}$, $v_{*}$ and $r_{*}$;
the definitions of the properties (Ad1), (Ad2) in \cite[Definition~7.1]{QEDI};
\cite[Section~7.1.5]{QEDI} with its displayed formulas;
\cite[Lemmas~7.5, 7.6, 7.7(d), 7.8]{QEDI} and \cite[Proposition~7.10]{QEDI};
the first clause of \cite[Theorem~7.12(2)]{QEDI}, together with the last sentence of its proof;
\cite[Lemmas~7.14, 7.16]{QEDI};
the heads of \cite[Sections~7.5, 7.6]{QEDI};
and \cite[Theorems~7.15, 7.17]{QEDI} and \cite[Corollary~7.18]{QEDI} as statements on limit points,
in the form in which \cite[Section~12]{QEDII} applies them,
without the last sentence of \cite[Theorem~7.15(9)]{QEDI}.
For single laws with the properties (L1)--(L4) of \cite[Section~12.4]{QEDII},
these three statements and \cite[Proposition~7.19]{QEDI} are replaced by Propositions~45.1--45.3.
In \cite{QEDI} the notion of a cutoff in \cite[Section~7.1.1]{QEDI},
\cite[Definition~7.1]{QEDI} with its third property (Ad3),
\cite[Lemma~7.9]{QEDI}, \cite[Theorems~7.11, 7.12]{QEDI}
and, through the notion of an admissible family, \cite[Theorems~7.15, 7.17]{QEDI} and \cite[Corollary~7.18]{QEDI}
are formulated for the odd lattices
\[
n=h\cdot L^{T+N},\qquad R=h\cdot L^{T}.
\]
On the tori of Section~I.7.5 their counterparts are
\cite[Lemmas~12.5--12.9]{QEDII},
\cite[Proposition~12.10(a)--(e)]{QEDII}, with the property (Ad3$'$) of \cite[Proposition~12.10(e)]{QEDII} in the place of (Ad3),
\cite[Proposition~12.12]{QEDII}
and \cite[Proposition~12.13]{QEDII}.
\item In \cite[Section~8]{QEDI}:
all statements, with \cite[Theorem~8.1]{QEDI} in the form of Proposition~45.6
and \cite[Theorem~8.8(v)]{QEDI} in the form of Proposition~45.7.
Proposition~45.6 is stated for the index set $\{1,2,3\}$.
In the proof of \cite[Theorem~8.3(iv)]{QEDI},
\cite[Theorem~8.1(a), (c)]{QEDI} and the first numbered display of the proof of \cite[Theorem~8.1]{QEDI} in \cite[Section~8.5]{QEDI}
are applied to the index set $J=\{j\}$, as they are stated in \cite{QEDI}.
In the notation of \cite[Section~8]{QEDI}, that display states that for a law $m$ with $(T_J)$ the kernel $(m_{\omega})_{\omega\in Q}$ constructed there satisfies
\[
m_{\omega}(\Phi)=E_m\left[\Phi\mid\mathcal{I}^J\right](\omega)\quad\text{for $m$-almost every }\omega,\qquad
m(\Phi)=\int\mathrm{d}m(\omega)\,m_{\omega}(\Phi),
\]
for every $\Sigma$-measurable function $\Phi\colon Q\to[0,\infty]$ and for every $\Phi\in L^1(m)$.
The one consequence drawn from them in the proof of \cite[Theorem~8.3(iv)]{QEDI},
namely $P^{j}_{\xi}X(f)=0$ for a law $\xi$ with (G) that is invariant under the translations by the vectors of $\mathbb{R}e_j$,
follows also without a decomposition, from \cite[Lemma~8.13(iii)]{QEDI} and \cite[Lemma~8.15(i)]{QEDI}.
In the proofs of \cite[Theorem~8.7]{QEDI} and of \cite[Theorem~8.8(i)]{QEDI},
\cite[Theorem~8.1]{QEDI} is applied to the index set $\{j\}$ in the same way.
The last paragraph of \cite[Section~8.1.5]{QEDI}, on the limit points of the lattice laws,
is replaced by Theorem~44\,(B1) together with the display (2.1) of Chapter~44,
which rest on \cite[Proposition~12.10(a)--(e)]{QEDII} and \cite[Proposition~12.12]{QEDII}.
\item In \cite[Section~9]{QEDI}:
all statements, with \cite[Proposition~9.7(b)]{QEDI} in the form of Proposition~45.6\,(a$'$), (d0), (d),
and the first assertion of \cite[Corollary~9.10]{QEDI}, on the set $\Xi_0$ of full measure, in the form of Proposition~45.6\,(c);
the assertion of \cite[Theorem~9.17]{QEDI} on mixtures is stated and proved within the proof of Theorem~44.
The last sentence of \cite[Section~9.2]{QEDI}, on the limit points of the lattice laws, is replaced by Theorem~44\,(B2), and, for the ergodic components, by Theorem~44\,(B5).
Of the objects that belong to the fermion, \cite[Section~9.7]{QEDI} contains only the continuum objects $\gamma_{\mu}$, $S_{\bar{m}}$ and $b(\cdot)$ of Sections~I.2.1, I.2.6 and I.7.6.
\item In \cite[Section~10]{QEDI}:
the two facts displayed in \cite[Section~10.1.1]{QEDI},
namely the Hodge decomposition \cite[(10.1)]{QEDI}, that is, the orthogonal decomposition of the closed cochains into exact and harmonic ones,
and the statement \cite[(10.2)]{QEDI} on the flux of a closed integer two-cochain, the latter also on a two-dimensional torus;
\cite[Lemma~10.2(a)]{QEDI};
\cite[Lemmas~10.6, 10.10]{QEDI}.
\item In \cite[Section~12]{QEDI}:
\cite[Lemmas~12.3, 12.4, 12.7(a)]{QEDI}.
\item In \cite[Section~13]{QEDI}:
of \cite[Section~13.1.6]{QEDI}, only the absolute constants $c_{X}$, $\kappa_{0}$, $K_{0}$, $c_{\mathrm{rb}}$, $c_{\mathrm{fl}}$ and the number $c(\gamma)$ with its bounds;
the constant $C_0^{\mathrm{unit}}$ of \cite[Section~13.1.6]{QEDI} is the third entry of $C_{0}$ in the line (C6) of Section~I.8.2.
The paragraphs of \cite[Section~13.1.6]{QEDI} on the twist and on the first step concern statements whose place is taken by statements of this paper;
at $k=0$ the pairs are those of Section~I.1.7 and of (I.L0.14).
Further:
the conventions of \cite[Section~13.1.7]{QEDI} that define $d_M$ through nets and the set $\bar{X}^L$,
and the convention there by which a cut-off function is evaluated at the midpoints of bonds;
\cite[Section~13.3]{QEDI};
\cite[Lemmas~13.9, 13.10]{QEDI}, with the three displayed properties (f1)--(f3) of the cut-off functions in \cite[Section~13.6]{QEDI};
and, from \cite[Section~13.10]{QEDI}, the two facts displayed in the proof of \cite[Theorem~13.18]{QEDI}:
a matrix $T$ with $\operatorname{Re}(f,Tf)\ge c\cdot|f|^2$ for all $f$, where $c>0$, is invertible, with $\left|T^{-1}\right|\le c^{-1}$ and $\operatorname{Re}T^{-1}\ge0$;
and the Woodbury identity.
\item In \cite[Section~14]{QEDI}:
\cite[Sections~14.1--14.5]{QEDI},
without the last sentence of \cite[Corollary~14.18]{QEDI}, the one that begins with ``In particular'',
and without the sentences of \cite[Section~14.5]{QEDI} on the case $k=k_{*}$ under the hypothesis of \cite[Definition~5.70]{QEDI}.
\item \cite[Section~15]{QEDI}.
\item In \cite[Section~16]{QEDI}:
the photon part of \cite[Definition~16.1]{QEDI};
the bounds on $\boldsymbol{\alpha}_{y}$ and $\boldsymbol{\beta}_{y}$ in \cite[Section~16.2.5]{QEDI};
the phases of \cite[Section~16.2.4]{QEDI};
the conditions (G1)--(G3) of \cite[Section~16.2.7]{QEDI};
\cite[Lemma~16.2]{QEDI};
of \cite[Section~16.6.1]{QEDI}, its head, \cite[Definition~16.9]{QEDI}, \cite[Lemma~16.10]{QEDI},
\cite[Definition~16.11]{QEDI} as far as $A^W$ and $\boldsymbol{\lambda}^W$,
and the sentences after it on $A^W$, $\boldsymbol{\lambda}^W$ and $G^W(\boldsymbol{\delta})$;
the maps $\mathrm{ext}$, $\mathrm{rs}$ and the operators $F^W_y$, $\mathrm{Sl}_{\square,y}$ of that definition,
and the identities $\mathrm{rs}\cdot\mathrm{ext}=1$, $\mathrm{rs}=\mathrm{ext}^{T}$ stated in the paragraph that follows it,
are those of Section~G1.5, (I.J3) and Proposition~2.5;
the properties of the quadratic partition of unity in \cite[Section~16.7.1]{QEDI};
\cite[Lemma~16.26]{QEDI}.
\item In \cite[Section~17]{QEDI}:
the last paragraph of \cite[Section~17.1.1]{QEDI}, with (I.Lc2)$^{\mathrm{ch}}$ in the place of \cite[Lemma~17.8]{QEDI};
the head of \cite[Section~17.3]{QEDI}, with its conventions and the hypotheses of that section,
and \cite[Lemmas~17.3--17.7]{QEDI}, the assertion of \cite[Lemma~17.6]{QEDI} on $\mathrm{Dom}_k$ holding for every $k\ge0$, every $\ell\in\mathrm{Cl}$ and every $0<e_k\le e_0$ (Section~I.9.1);
\cite[Section~17.4.1]{QEDI} with (g0)--(g4);
\cite[Section~17.4.3]{QEDI};
\cite[Lemmas~17.12, 17.13]{QEDI}.
\item In \cite[Section~19]{QEDI}:
\cite[Section~19.2]{QEDI}, apart from \cite[Proposition~19.2]{QEDI}.
\item In \cite[Section~20]{QEDI}:
\cite[Section~20.3]{QEDI}.
\item In \cite[Section~21]{QEDI}:
the first paragraph of \cite[Section~21.1.9]{QEDI}, with $S_{\mathrm{Mq}}$, $q_{4,\mathrm{Mq}}(a)$, $b(f)$ and the fact that $f$ determines $a$;
\cite[Theorem~21.1(a)]{QEDI};
\cite[Section~21.4.1]{QEDI} with (K1), (K2);
\cite[Lemmas~21.5, 21.31, 21.33]{QEDI};
\cite[Remark~21.17]{QEDI}, without its last sentence, on \cite[Theorem~3.1, Part~I\,(d)]{QEDI},
and read for an arbitrary family of Borel probability measures on $H_{-4}$ indexed by the cutoffs $(N,T)$:
an iterated weak limit of such a family, taken in either order, is a weak limit of the family along a sequence of cutoffs $(N_i,T_i)$ with $N_i\to\infty$ and $T_i\to\infty$;
\cite[Section~21.5]{QEDI}.
Of these, the first paragraph of \cite[Section~21.1.9]{QEDI}, \cite[Theorem~21.1(a)]{QEDI}, \cite[Section~21.4.1]{QEDI} and \cite[Section~21.5]{QEDI}
concern the continuum Dirac operator only:
$b(f)$ is defined in Section~I.7.6 without a lattice,
and by Propositions~37.1\,(f) and 37.2\,(f) the lattice coefficient of the chamber operator converges to it.
\end{itemize}
Beside the two facts from the proof of \cite[Theorem~13.18]{QEDI} stands the following bound on the axial functional $\tau_k$ of \cite[Section~5.2.6]{QEDI};
it is also derived in the proof of \cite[Lemma~16.2]{QEDI}.

\begin{named}{Lemma}[Bound for the axial functional]
Let $k\ge1$, let $A$ be a bond function on $T_k$, let $y$ be a unit site and let $x\in B_k(y)$.
Then
\[
\eta_k\cdot\left|(\tau_kA)(y,x)\right|\le\frac{3}{2}\cdot|A|_{\infty}.
\]
\end{named}

\begin{proof}
For $0\le j\le k$ let $y_j$ be the centre of the block of side $L^j\cdot\eta_k$ that contains $x$;
then $y_0=x$, $y_k=y$ and
\[
(\tau_kA)(y,x)=\sum_{j=0}^{k-1}(\tau A)(y_{j+1},y_j).
\]
Let $0\le j\le k-1$.
The block of side $L^j\cdot\eta_k$ with the centre $y_j$ lies in the block of side $L^{j+1}\cdot\eta_k$ with the centre $y_{j+1}$, and $L$ is odd;
hence every coordinate of $y_j-y_{j+1}$ has modulus at most $\frac{L-1}{2}\cdot L^j\cdot\eta_k$, and
\[
|y_{j+1}-y_j|_1\le\frac{3}{2}\cdot(L-1)\cdot L^j\cdot\eta_k .
\]
Each of the six rectilinear paths from $y_{j+1}$ to $y_j$ that enter the definition of $\tau$ consists of $|y_{j+1}-y_j|_1/\eta_k$ bonds,
and $\eta_k\cdot(\tau A)(y_{j+1},y_j)$ is the mean, over these six paths, of the sums of $\eta_k\cdot A(b)$ over their bonds $b$.
Therefore
\[
\eta_k\cdot\left|(\tau A)(y_{j+1},y_j)\right|\le|y_{j+1}-y_j|_1\cdot|A|_{\infty}\le\frac{3}{2}\cdot(L-1)\cdot L^j\cdot\eta_k\cdot|A|_{\infty}.
\]
Summation over $j$, with $\sum_{j=0}^{k-1}(L-1)\cdot L^j=L^k-1$ and $L^k\cdot\eta_k=1$, gives
\[
\eta_k\cdot\left|(\tau_kA)(y,x)\right|\le\frac{3}{2}\cdot\left(1-\eta_k\right)\cdot|A|_{\infty}\le\frac{3}{2}\cdot|A|_{\infty}.
\]
\end{proof}

\paragraph{The integration lemma.}
The statement \cite[Lemma~17.17]{QEDI} involves the Grassmann algebra $\mathrm{Gr}[\psi,\chi,W]$ and the measure $\mathrm{d}\mu_I(W)$,
and its proof in \cite{QEDI} rests on AC1 and on the last sentence of \cite[Corollary~17.15]{QEDI},
whose places are taken here by Proposition~15 and Proposition~24.3.
For every $k\ge0$ the lemma is stated and proved as Lemma~5.C of Chapter~24, within the proof of Proposition~24.1\,(a).
For $k\ge1$ that lemma is \cite[Lemma~17.17]{QEDI} as stated:
on the index sets of this paper the rules (I.N1), (I.N3), (I.N4) are Proposition~15\,(b), (f), (g),
and the family $W$ has the weight $\nu=1$.
At $k=0$ the lemma takes the following form.
Its hypothesis is
\[
\left\|V(X)\right\|_{(\mathsf{w};\mathsf{w}_W)}\le\mathbf{H}_{0}\cdot\exp\left(-\mathrm{rt}_2\cdot d_M(X)\right),
\qquad\text{with }\mathsf{w}_W=\sqrt{48}\text{ as in (I.5.2)},
\]
where $\mathrm{rt}_2=\kappa-2\kappa_{0}-2$ is the rate of Chapter~24, Section~0.3, which \cite[Section~17.1.2]{QEDI} denotes by $\kappa_2$,
and $\mathbf{H}_{0}=36\,c_X^2\cdot K_0^2\cdot e_k^{\frac{1}{4}-5\varepsilon}$ is the number of the condition (A7) of \cite[Section~5.6.2]{QEDI};
the rule (I.N3)$^{\mathrm{w}}$ takes the place of (I.N3),
and (I.N4) is applied at the points of $T^{\mathrm{f}}_{0}$;
the hypothesis that $V(X)$ contain no generator $W(u)$ whose unit site lies outside $X$ is read as follows:
$V(X)$ contains a generator $W(u)$, $u=(x,\beta,\omega)$, only if the fermion site $x$ lies over $X$, that is, only if $\pi(x)\in X$.
The conclusion and the constants are those of \cite[Lemma~17.17]{QEDI}.

\paragraph{Statements of \cite{QEDI} that are replaced.}
All further statements of \cite[Section~5]{QEDI} are replaced by statements of this paper.
The definitions of \cite[Section~5]{QEDI} that this paper retains for photon objects appear in Part~I:
\cite[Section~5.2.3]{QEDI} and \cite[Definitions~5.4, 5.7, 5.8]{QEDI} in Section~I.1.1;
\cite[Section~5.3.1]{QEDI} and \cite[Definition~5.16]{QEDI} in Section~I.5.3;
\cite[Section~5.2.5]{QEDI}, \cite[Definitions~5.10, 5.11]{QEDI} and the recursion of \cite[Section~5.3.8]{QEDI} in Section~I.5.7;
\cite[Section~5.2.2]{QEDI} and \cite[Definition~5.6]{QEDI} in Sections~I.6.1 and I.6.2;
\cite[Definition~5.12]{QEDI} in Section~I.6.6.
Besides the exact root, \cite[Definition~5.10]{QEDI} contains the definition of a controlled pair;
this is not a photon object, and it is restated in Section~G4.0.5.
Moreover the following remain in force as in \cite{QEDI}:
\cite[Definition~5.5(1)]{QEDI}, which fixes the meaning of the word ``reads'', the word being used in this sense in (I.f4) of Section~I.1.9;
\cite[Section~5.2.4]{QEDI}, as far as it describes the action of $\mathbb{G}_k$ on photon objects (Sections~I.4.2, I.4.3);
the formulation of W1--W4, W8--W10, W14 and W16 for the families (p1), (p2) in \cite[Section~5.4.2]{QEDI} (Section~G4.0.4);
and the formulas of \cite[Sections~5.4.1, 5.5.3, 5.5.8]{QEDI} that Proposition~14.2\,(a) names.

In the other sections of \cite{QEDI} the following statements are replaced;
the statements of this paper that take their places are given in parentheses.
\begin{itemize}
\item \cite[Section~7.1.2]{QEDI} (Sections~I.7.2, I.7.3; Propositions~50\,(a), (e), 38.1\,(d) and 31.1\,(a)--(c));
\cite[Proposition~7.2]{QEDI} and the last sentence of \cite[Theorem~7.11]{QEDI} (Proposition~31.1\,(d) and \cite[Proposition~12.10(c)]{QEDII});
\cite[Lemma~7.7(a)]{QEDI} ((I.6.1) and Proposition~38.1\,(e));
\cite[Lemma~7.7(b)]{QEDI} (Proposition~50\,(d), (f), (g$''$));
\cite[Lemma~7.7(c)]{QEDI} (Proposition~50\,(c), (g$'$), (g$''$));
\cite[Definition~7.1]{QEDI} with (Ad3), \cite[Lemma~7.9]{QEDI}, \cite[Corollary~7.4]{QEDI} and \cite[Theorems~7.11, 7.12]{QEDI}
(\cite[Proposition~12.10]{QEDII}, \cite[Lemma~12.8]{QEDII}, \cite[Proposition~12.12]{QEDII} and \cite[Proposition~12.13(a)]{QEDII});
\cite[Definition~7.13]{QEDI} (Section~I.7.5 and \cite[Proposition~12.13(d)]{QEDII}).
\item \cite[Definition~12.8]{QEDI} with (B3) and (Pen1)--(Pen3), \cite[Definition~12.9]{QEDI} and \cite[Lemma~12.10]{QEDI}.
These concern $L^2(E;\mathbb{C}^4)$, with kernels taken with respect to $\eta^3\sum_{E}$;
their place is taken by Proposition~1.6, whose proof establishes the corresponding properties of its penalty on $L^2(E\times O_{\mathrm{cub}};\mathbb{C}^4)$ (Chapter~1b, Lemma~6.1).
\item \cite[Section~13]{QEDI}, apart from the statements listed above:
\cite[Section~13.2]{QEDI} (Proposition~15) and the convention of \cite[Section~13.1.7]{QEDI} on pairs admissible for a union of $M$-cubes (Section~I.1.7).
\item \cite[Sections~14.6--14.8]{QEDI}, as far as Section~I.9.4 assigns them to the Propositions~21.m and~22.m.
\item \cite[Section~16]{QEDI} and \cite[Section~17]{QEDI}, apart from the statements listed above;
in particular \cite[Lemma~17.8]{QEDI} ((I.Lc2)$^{\mathrm{ch}}$) and \cite[Lemma~17.17]{QEDI} (Chapter~24, Lemma~5.C).
\item \cite[Proposition~19.2]{QEDI}.
It involves $C_{0}$, $c_{\kappa}$, $\gamma_1$, $C_{\Sigma}$ and the number $e_1$ of \cite[Section~19.1.2]{QEDI}.
Its place is taken
by Proposition~14.2\,(c) for the condition $r_{*}>31L+16$ of \cite[Section~5.6.2]{QEDI} and for the conditions (fl), (T1)--(T6) there,
by Proposition~14.2\,(d) for the condition (W3) with (Z8.w) of Section~I.8.4,
by Proposition~14.2\,(f) for $e\le\frac{1}{3}$ and $v_K\ge8$,
and by the row (e1) of Proposition~14.2\,(b) at $v=v_K$, together with Proposition~14.3\,(f), for
\[
\left(1+\frac{3\cdot\mathrm{lg}}{v_K}\right)^{r}<L.
\]
Apart from (W3), this is the condition (Sc6) of Section~G4.0.2;
the condition (W3) is not contained in (Sc6), and it is restated in (G6.T-14)\,(v).
\item \cite[Sections~19.1, 19.3--19.10]{QEDI}.
\item \cite[Section~20]{QEDI}, apart from \cite[Section~20.3]{QEDI};
in particular \cite[Sections~20.7, 20.8, 20.9.1]{QEDI} (Propositions~35.7--35.9).
\item \cite[Section~21.1]{QEDI}, apart from the first paragraph of \cite[Section~21.1.9]{QEDI};
\cite[Section~21.2]{QEDI}, apart from \cite[Theorem~21.1(a)]{QEDI};
\cite[Section~21.3.1]{QEDI};
\cite[Sections~21.8--21.10]{QEDI}, apart from \cite[Lemmas~21.31, 21.33]{QEDI};
\cite[Lemma~21.36]{QEDI} (Section~I.9.4).
\end{itemize}
Of \cite[Section~6]{QEDI}, the statements whose place is taken by Proposition~31.1 are those named in Section~I.9.4.
The definitions of \cite[Sections~10.1.1, 10.1.4]{QEDI} and \cite[Definition~10.1]{QEDI}
are used with the determinant of the chamber operator, in Sections~I.1.1, I.2.9 and I.7.1.
Likewise \cite[Definition~11.1]{QEDI} and \cite[Section~11.2.1]{QEDI} are definitions that are used with the determinant of the chamber operator,
while \cite[Lemmas~11.2, 11.7]{QEDI}, together with the properties of the determinant on which they rest,
are replaced by the Propositions~38.m and by Proposition~50.
The statements of \cite[Section~12]{QEDI} other than those listed above are replaced by the Propositions~1.m.

\paragraph{The opening sections of \cite{QEDI}.}
The following parts of \cite[Sections~1--4]{QEDI} hold as notation and as statements on the continuum:
\cite[Section~2.1]{QEDI} (cochains, $h$);
\cite[Section~2.4]{QEDI} (cells, the test map, the Hermite scale, the notion of a limit point);
\cite[Section~2.5]{QEDI};
the paragraph of \cite[Section~2.6]{QEDI} on the coefficient on $\mathbb{R}^3$, which is the first paragraph of \cite[Section~21.1.9]{QEDI};
and the existence of probes with $b(f)>0$, which is \cite[Theorem~21.1(a)]{QEDI}.
\cite[Sections~2.2, 2.3]{QEDI}, the laws and the mass assignment of \cite[Section~2.4]{QEDI},
the paragraph of \cite[Section~2.6]{QEDI} on the coefficient on a lattice torus,
\cite[Sections~2.7, 2.8]{QEDI}, both parts of \cite[Theorem~3.1]{QEDI} and \cite[Corollaries~3.2, 3.3]{QEDI}
are replaced by Sections~I.2.3, I.6.1, I.6.2 and I.7.1--I.7.6, by the Propositions~37.m and by Theorem~44.

\paragraph{Directions, matrices, sites.}
In \cite[Sections~5, 11, 13--20]{QEDI}, in \cite[Section~21.3.3]{QEDI} and in \cite[Section~21.7]{QEDI}
the coordinate directions are numbered $0,1,2$;
the direction $\mu$ there is the direction $\mu+1$ of this paper (Section~I.1.1).
The matrices of \cite{QEDI} are related to those of this paper by Section~I.2.1.
A fine site $x$ of \cite{QEDI} is the photon site $c$ in statements on photon objects,
and the fermion site $(c,g)$ in statements on spinors.

\subsubsection*{I.9.3. The statements of \cite{QEDII}}

\paragraph{Statements that hold as stated.}
Every statement of \cite[Sections~1--14]{QEDII} holds as it is stated there,
read in the setting of this paper as described below.
In them the five modifications (I.h1)--(I.h5) of Section~I.6.4, which admit cell widths in $\frac{1}{2}\mathbb{Z}$ for the coordinate frame, are made at the places named there.

\paragraph{How a statement of \cite{QEDII} is read.}
In \cite[Section~1.1]{QEDII}, in \cite[Section~1.2]{QEDII} up to the paragraph on the continuum, with \cite[Theorems~1.1 and~1.2]{QEDII},
in \cite[Definitions~2.1--2.35]{QEDII} and in \cite[Sections~3--11, 13.1, 13.3 and~13.4]{QEDII}
the quantities of \cite{QEDII} are in lattice units and the pairing is the counting pairing;
the curvature $F$ there is $F^{\mathrm{u}}$ or, as an element of the space $\Omega$ of the closed two-cochains on $\mathbb{Z}^3$, its periodic extension.
In \cite[Definitions~2.36--2.39]{QEDII}, in \cite[Section~1.2]{QEDII} from the paragraph on the continuum on, with \cite[Theorems~1.3 and~1.4]{QEDII},
and in \cite[Sections~12, 13.2, 13.5, 13.6 and~14]{QEDII},
wherever the letters $\eta$, $R$, $e$, $\bar{m}$ denote the data of a cutoff
and $F^{\mathrm{u}}$, $F=\eta^{-\frac{3}{2}}\cdot F^{\mathrm{u}}$, $\langle\cdot,\cdot\rangle_\eta$ denote the curvature in lattice units, the physical curvature and the pairing at a cutoff,
they have the meaning of Sections~I.7.1 and~I.7.5, and they are read unchanged.
\cite[Section~14]{QEDII} treats laws on $\mathcal{S}'$ in the notation of \cite[Definition~2.38]{QEDII}, $F$ being the variable of $\mathcal{S}'$;
data of a cutoff occur there only in \cite[Section~14.14]{QEDII}, where the spacing is written $\eta_i=L^{-N_i}$ and where $\bar{m}$ and $e$ occur.
Elsewhere in \cite[Section~14]{QEDII} the letters $\eta$ and $R$ have the meanings declared at the places where they occur:
$\eta$ denotes a function, in \cite[Sections~14.2, 14.6, 14.7 and~14.13]{QEDII};
among the objects written with $R$ are the rotations $R^{ji}_t$ and $R_t$, the complex linear maps $R_\zeta$, $\zeta\in\mathbb{C}$, of $\mathbb{C}^2$,
and the condition $(R_j)$ on laws.
The operator $D_A$ of \cite[Definition~2.10]{QEDII} is $D_A+\boldsymbol{\mu}$,
its mass parameter $M$ is $M_{\mathrm{ch}}$,
and its lattice mass $m$ and the mass $m_N$ of \cite[Definition~2.36]{QEDII} are $\boldsymbol{\mu}$ and $\boldsymbol{\mu}_N$.
Its group $W$ is $W_{\mathrm{af}}$.
A site $x\in\mathbb{X}$ of \cite{QEDII} is the point $\mathrm{geo}(c,g)$;
by Proposition~46\,(a), (h) nothing depends on the point $v_0$.
The torus $\mathbb{T}=C(n)$ of \cite{QEDII} is the antiperiodic torus of Section~I.1.5,
and $Z_{\mathbb{T}}$, $\mathbb{E}_{\mathbb{T}}$ are $Z_n$, $E^{\mathrm{af}}_n$.
The pairing of \cite{QEDII} is the counting pairing,
and its Berezin integral, normalised in \cite[Definition~2.8]{QEDII}, is the integral (I.3.5) with $\nu=1$; hence
\[
\int_{\mathbb{T}}\exp\left(-S_{\mathbb{T}}(A)\right)=\det\left(D_A+\boldsymbol{\mu}\right)
\]
in both normalisations (Proposition~46\,(e)).
The following objects of \cite{QEDII} are transported by $\mathrm{geo}$:
the reflections $\Theta_\Pi$ of \cite[Definition~2.9]{QEDII}, written $\Theta_m$ in this paper;
the action $G\mapsto G^g$ of the group $W$ on even elements, defined after \cite[Definition~2.12]{QEDII}, in the paragraph on the action of the group $W$, and written $G\mapsto G^{\mathrm{gw}}$ in this paper;
the classes $\mathcal{G}(B)$ and $\mathcal{A}(U)$ of \cite[Definition~2.12]{QEDII} and the class $\mathcal{P}$ of \cite[Definition~2.15]{QEDII},
written $\mathrm{Ad}(\mathrm{Bx})$, $\mathcal{A}(U)$ and $\mathrm{Pl}$ in this paper;
and the functional $Z_{\mathbb{T}}[G]$ of \cite[Definition~2.13]{QEDII} for $G$ with values in the Grassmann algebra.
By Proposition~46\,(a) the map $\mathrm{geo}$ intertwines the actions of $W_{\mathrm{af}}$,
and it carries the projection $\pi$, the types of the bonds and the directions $\hat d_i$ of this paper onto the objects of the same names of \cite[Definition~2.4]{QEDII}:
the map $\pi$ of \cite[Definition~2.4]{QEDII} satisfies $\pi(\mathrm{geo}(c,g))=c$,
$\mathrm{geo}$ carries a bond of type $i$ onto a bond of type $i$,
and, for a fermion site $x=(c,g)$ of this paper, $\hat d_i(x)$ is the unit vector $\hat d_i(\mathrm{geo}(x))$ of \cite[Definition~2.4]{QEDII}.
In the coordinates $(c,g)$ the transported objects are those of Sections~I.4.8 and I.6.4.
For the cells and walls of Section~I.6.4,
the width $h$ and the cell index $\kappa$ of a cell structure in \cite[Definition~2.18]{QEDII} are $\mathrm{hc}$ and $\mathfrak{c}$,
and the position $c$ of a mirror $\Pi=(\ell,c)$ of \cite[Definition~2.2]{QEDII},
in particular the position $o_a+j\cdot h$ of a wall $(\ell^a,o_a+j\cdot h)$ of \cite[Definition~2.18]{QEDII}, is $\mathrm{po}$.

\paragraph{Statements of \cite{QEDI} applied in \cite{QEDII}.}
In \cite{QEDII} the following statements of \cite{QEDI} are applied in the form in which they stand in \cite{QEDI}:
\cite[Theorems~7.15, 7.17]{QEDI} and \cite[Corollary~7.18]{QEDI}, through \cite[Proposition~7.19]{QEDI};
\cite[Lemma~21.36]{QEDI};
\cite[Theorem~8.1]{QEDI};
and \cite[Theorem~8.8(v)]{QEDI}.
As they are applied there, they are statements on single laws without a fermion object, and they hold as stated;
the corresponding statements of this paper are Propositions~45.1--45.4, 45.6 and~45.7.

\paragraph{Slice data from \cite{QEDII}.}
Let $n_2,n_3\ge4$, and let $\mathrm{Fk}$, $\mathrm{Fk}'$, $\mathrm{Qs}(A)\colon\mathrm{Fk}'\to\mathrm{Fk}$ and $\mathrm{Wq}(a)$
be the objects $\mathcal{F}$, $\mathcal{F}'$, $\widehat{\mathfrak{N}}(A)$ and $u(a)$ of \cite[Definition~2.27]{QEDII}
for the cross-section with the sides $n_2$, $n_3$, the time direction being $e_1$.
For an arbitrary unitary map $\mathrm{Uf}\colon\mathrm{Fk}\to\mathrm{Fk}'$ put
\[
\mathrm{Qf}(A):=\mathrm{Qs}(A)\cdot\mathrm{Uf},\qquad s:=1,\qquad\mathrm{Gmf}:=\mathrm{Wq}.
\]
In \cite[Proposition~5.33]{QEDII} the time length $n_1$ is at least $4$.
For $\Upsilon_j=u(a_j)$, $j\in\mathbb{Z}/n_1\mathbb{Z}$, the display of that proposition together with its last sentence gives
\[
\det\nolimits_n\left(D_A+\boldsymbol{\mu}\right)=\operatorname{tr}_{\mathcal{F}'}\prod_{j=n_1-1,\dots,0}\left(\widehat{\mathfrak{N}}(A_{j+1})^{*}\cdot u(a_j)\cdot\widehat{\mathfrak{N}}(A_j)\right),\qquad A_{n_1}=A_0,
\]
with $A_j$ and $a_j$ as in \cite[Definition~2.26]{QEDII} and the later times to the left.
By Remark~39.R (Chapter~39, Lemma~8.1) this identity holds for every time length $n_1\ge1$.
With it and with \cite[Lemma~5.35(a)]{QEDII}, the data $\mathrm{Fk}$, $s$, $\mathrm{Qf}$ and $\mathrm{Gmf}$,
in the roles of $\mathcal{F}$, $s$, $\mathfrak{Q}$ and $\mathsf{G}$ of \cite[Lemma~21.4]{QEDI},
are slice data with the properties (a), (b) stated there,
the function $A\mapsto\det\nolimits_n\left(D_A+\boldsymbol{\mu}\right)$ standing in (b) in the place of $\mathrm{Dt}(qA)$.
Neither an intertwining property of $\mathrm{Uf}$ nor the gauge covariance of \cite[Lemma~5.35(b)]{QEDII} is needed for this.

\subsubsection*{I.9.4. Concordance}

For every number $n$, the following list names the statements of \cite{QEDI} to which the Propositions $n.m$ correspond.
Where the Propositions $n.m$ provide, for the chamber operator, properties stated in \cite[Section~5]{QEDI}, these properties are named after the statements.
Here property~$j$ of \cite[Section~5.4.2]{QEDI} is the $j$-th entry of the list of properties of systems of primitives printed in that section, and (p3)--(p5) are the families of primitives of \cite[Section~5.4.1]{QEDI}.
The six properties of the fermion weight that are displayed as (10.3)--(10.8) in \cite[Section~10.1.2]{QEDI} are denoted by (W1)--(W6), in this order.
These six names refer to the fermion weight only; the condition (W3) of Section~I.8.4 and the properties W1--W18 of systems of primitives are other objects.

\begin{description}
\item[Propositions 1.1--1.7]
They correspond to the third display of \cite[Theorem~12.20]{QEDI}, which is the sup and H\"older bound, to \cite[Proposition~12.21]{QEDI}, \cite[Lemmas~12.1, 12.3 and 12.5]{QEDI}, \cite[Remark~12.2]{QEDI}, \cite[Proposition~12.6]{QEDI}, \cite[Lemma~12.7(b)]{QEDI}, to the condition (B3) of \cite[Definition~12.8]{QEDI} together with (Pen1)--(Pen3), and to \cite[Definition~12.9]{QEDI} and \cite[Lemma~12.10]{QEDI}.
\item[Propositions 2.1--2.8]
They correspond to \cite[Lemmas~16.3, 16.4, 16.5 and 16.7]{QEDI}, \cite[Propositions~16.6 and 16.8]{QEDI}, to \cite[Definition~16.11]{QEDI} as far as the objects $\mathrm{ext}$, $\mathrm{rs}$, $F^{W}_{y}$ and $\mathrm{Sl}_{\square,y}$ are concerned, to \cite[Proposition~16.12]{QEDI}, \cite[Section~16.6.2]{QEDI}, \cite[Section~16.7.2]{QEDI} with \cite[Lemma~16.13]{QEDI}, and to \cite[Definition~16.14]{QEDI}, \cite[Lemma~16.15]{QEDI} and \cite[Theorem~16.16]{QEDI}.
They provide \cite[Propositions~5.42 and 5.43]{QEDI}, which concern the levels $k\ge1$, and the part of \cite[Proposition~5.44]{QEDI} that concerns the level~$0$; Proposition~2.8 states them in the Euclidean norms.
Their forms for the chamber operator, in the norms of Section~I.1.5 and with the constant $C^{F}$, are clauses of Proposition~14.1\,(i): the forms of \cite[Propositions~5.42 and 5.43]{QEDI} for $k\ge1$, and the form of \cite[Proposition~5.44]{QEDI} at $k=0$.
\item[Proposition 3]
It corresponds to \cite[Proposition~16.18]{QEDI}.
\item[Proposition 4]
It corresponds to \cite[Proposition~16.19]{QEDI} and provides the algebraic part of \cite[Proposition~5.44]{QEDI}.
\item[Proposition 5]
It corresponds to \cite[Proposition~16.20]{QEDI} and provides the bounds of \cite[Proposition~5.44]{QEDI} for $k\ge1$.
\item[Proposition 6]
It corresponds to \cite[Proposition~16.21]{QEDI} and provides \cite[Proposition~5.45]{QEDI}.
\item[Proposition 7]
It corresponds to \cite[Theorem~16.23]{QEDI} and provides the properties 1, 2 and 5--16 of \cite[Section~5.4.2]{QEDI} for the families (p3)--(p5) at the levels $k\ge1$.
\item[Proposition 8]
It corresponds to \cite[Theorem~16.25]{QEDI} and provides property~17 of \cite[Section~5.4.2]{QEDI}.
\item[Proposition 9]
It corresponds to \cite[Proposition~16.27]{QEDI} and provides \cite[Proposition~5.46]{QEDI}.
\item[Proposition 10]
It corresponds to \cite[Proposition~16.28]{QEDI} and provides \cite[Proposition~5.47]{QEDI}.
\item[Proposition 11]
It corresponds to \cite[Proposition~16.29]{QEDI} and provides \cite[Proposition~5.49]{QEDI}.
\item[Proposition 12]
It corresponds to \cite[Lemma~16.30]{QEDI}.
\item[Proposition 13]
It corresponds to \cite[Proposition~16.31]{QEDI} and provides \cite[Proposition~5.48]{QEDI}.
\item[Propositions 14.1--14.4]
They correspond to \cite[Theorem~16.32]{QEDI}, to \cite[Sections~5.6.1 and 5.6.2]{QEDI}, \cite[Section~17.13.3]{QEDI} and \cite[Section~19.1.2]{QEDI}, to \cite[Proposition~19.2]{QEDI} (through Proposition~14.2\,(c), (d), (f)), and to \cite[Sections~20.1.7 and 20.1.8]{QEDI}.
Of \cite[Theorem~5.72]{QEDI} they provide the clauses (i) and (ii) in the form of Proposition~14.1\,(i), (ii), at the fixed rate $\kappa'=\kappa_{\mathrm{w}}+\kappa_0$ and for $e_k\in(0,e_0]$;
the statement after its clause (iii) by which \cite[Proposition~5.48]{QEDI} applies, as stated, to the system of its clause (ii) for every $e_k\in(0,e_0]$, through Proposition~14.1\,(ii), (iii);
and, among the values stated at the end of that theorem, $\mathfrak{m}_0=1$ and $\alpha^{F}=\frac{1}{54}$ (Proposition~14.1).
Clause (iii) of \cite[Theorem~5.72]{QEDI} has no counterpart here (Section~I.9.5).
They provide further the stages 1--12 of \cite[Definition~5.97]{QEDI}.
\item[Proposition 15]
It corresponds to \cite[Section~13.1.4]{QEDI}, to the head of \cite[Section~13.2]{QEDI} and to all clauses of \cite[Theorem~13.1]{QEDI}, and provides \cite[Proposition~5.57]{QEDI}.
\item[Propositions 16.1--16.3]
They correspond to \cite[Section~13.1.5]{QEDI}, to the groups and their actions in \cite[Section~13.4.1]{QEDI}, and to \cite[Proposition~13.3]{QEDI}, \cite[Lemma~13.4]{QEDI} and \cite[Theorem~13.5]{QEDI}.
They provide \cite[Proposition~5.59]{QEDI}.
\item[Propositions 17.1 and 17.2]
They correspond to \cite[Definition~13.6]{QEDI} and \cite[Theorem~13.7]{QEDI}.
They provide the transport of \cite[Definition~13.6]{QEDI} in the form (I.SY6)$^{\mathrm{ch}}$, and \cite[Proposition~5.60]{QEDI}.
\item[Propositions 18.1--18.6]
They correspond to \cite[Lemma~13.8]{QEDI}, \cite[Theorem~13.11]{QEDI}, \cite[Theorem~13.13(a)--(e)]{QEDI}, \cite[Lemma~13.14]{QEDI}, \cite[Theorem~13.15]{QEDI} and \cite[Remark~13.16]{QEDI}.
They provide \cite[Propositions~5.61, 5.67 and 5.62]{QEDI}.
\item[Proposition 19]
It corresponds to the head of \cite[Section~13.10]{QEDI} and to \cite[Theorem~13.18]{QEDI}, and provides \cite[Proposition~5.65]{QEDI}.
\item[Propositions 20.1--20.3]
They correspond to \cite[Theorem~13.17]{QEDI}, \cite[Lemma~13.19]{QEDI} and \cite[Theorem~13.20]{QEDI}, and provide \cite[Propositions~5.64 and 5.66]{QEDI}.
\item[Propositions 21.1--21.5]
They correspond to \cite[Lemma~14.19]{QEDI}, \cite[Proposition~14.20]{QEDI}, \cite[Theorem~14.21]{QEDI} and \cite[Lemma~14.22]{QEDI}, and provide \cite[Proposition~5.55]{QEDI}.
\item[Propositions 22.1--22.4]
They correspond to the Berezin rules of \cite[Section~14.6.1]{QEDI}, to \cite[Lemma~14.23]{QEDI} and to \cite[Theorem~14.24]{QEDI}.
They provide the identity of \cite[Lemma~14.23(b)]{QEDI} and \cite[Proposition~5.56]{QEDI}.
\item[Propositions 23.1 and 23.2]
They correspond to the display of $D_B$, $Q_k(B)$ and $T_k(B)$ in \cite[Section~17.1.2]{QEDI};
to the two displays of \cite[Section~17.4.4]{QEDI}, in their forms with the slot, that define $T^{(1)}(B)$, $T^{(2)}(B)$, $\bar{T}^{(i)}(B)$ and, at $k=0$, $G^{(1)}(s,\square)$;
to \cite[Lemma~17.9]{QEDI}, \cite[Proposition~17.16(b)]{QEDI} and \cite[Lemma~17.21]{QEDI};
and to Steps~1 and~2 of the proof in \cite[Section~17.12]{QEDI}.
They provide the two displays that contain rows of $D_A$ and enter \cite[Definitions~5.14 and 5.25]{QEDI}: the display of $T^{(1)}$, $T^{(2)}$ and $\bar{T}^{(i)}$ with the substitution, and the display of $G^{(1)}$ at $k=0$.
As a whole, \cite[Definition~5.25]{QEDI} corresponds to Definition~24.D (Section~G3.4.1);
the other parts of \cite[Definition~5.14]{QEDI} stand in Section~I.5.8 (the diagonal families), in Section~I.5.7 (the function $\acute{E}_0$) and in Definition~24.D\,(D1) (the datum $(m,0)$).
\item[Propositions 24.1--24.3]
They correspond to \cite[Theorem~17.2]{QEDI}, \cite[Lemmas~17.20 and 17.21]{QEDI}, to the condition (adm) of \cite[Section~17.7]{QEDI}, to \cite[Lemma~17.11(q1), (q5)]{QEDI}, \cite[Lemma~17.14]{QEDI}, \cite[Corollary~17.15]{QEDI}, \cite[Lemma~17.17]{QEDI}, to the assertion~(2), on locality, of \cite[Section~17.10.1]{QEDI}, and to \cite[Proposition~17.24(a), (c), (d)]{QEDI}.
They provide \cite[Theorem~5.71(a)--(f), (h)]{QEDI}.
Of \cite[Theorem~5.71(g)]{QEDI} they provide the first paragraph (Proposition~24.1\,(g)) and the first, third, fourth and fifth assertions of its list, through the clauses (al), (ga), (ze) and (de) of Proposition~24.3, in this order.
For the second assertion of that list, Proposition~24.3\,(be) gives only the places at which the exact operators and forms stand.
Definition~24.D (Section~G3.4.1), which defines the step map treated in them, corresponds to \cite[Definition~5.25]{QEDI}.
\item[Propositions 25.1--25.3]
They correspond to \cite[Proposition~17.24(b)]{QEDI} and to clause~(2) of \cite[Definition~5.5]{QEDI}.
They provide the second assertion of the list in \cite[Theorem~5.71(g)]{QEDI}: the exact operators and forms in the displays of the step are given by one formula in the sense of Section~I.1.9, with their ranges.
\item[Propositions 26.1--26.6]
They correspond to \cite[Section~19.3.1]{QEDI}, \cite[Lemma~19.3]{QEDI}, \cite[Definition~19.4]{QEDI}, \cite[Theorem~19.5]{QEDI}, \cite[Lemmas~19.6 and 19.7]{QEDI} and \cite[Proposition~19.8]{QEDI}.
They provide \cite[Definition~5.13]{QEDI}, property~18 of \cite[Section~5.4.2]{QEDI} and \cite[Proposition~5.30]{QEDI}, and a part of \cite[Proposition~5.68]{QEDI}, namely the conditions (G4.Sc3), (G4.Sc5) of Section~G4.0.2 and Proposition~26.6\,(b)--(d).
\item[Propositions 27.1--27.6]
They correspond to \cite[Definition~19.9]{QEDI}, \cite[Lemmas~19.10 and 19.11]{QEDI} and \cite[Proposition~19.12]{QEDI}, and provide \cite[Theorem~5.80]{QEDI}.
\item[Propositions 28.1--28.3]
They correspond to \cite[Propositions~19.13 and 19.14]{QEDI} and \cite[Theorem~19.16]{QEDI}, and provide \cite[Theorems~5.81 and 5.83]{QEDI}.
\item[Proposition 29]
It corresponds to \cite[Theorem~19.17]{QEDI} and provides \cite[Theorem~5.84]{QEDI}.
\item[Propositions 30.1--30.3]
They correspond to \cite[Lemmas~19.18, 19.20 and 19.21]{QEDI} and \cite[Corollary~19.19]{QEDI}, and provide \cite[Theorem~5.85(d)]{QEDI}.
\item[Propositions 31.1 and 31.2]
They correspond to \cite[Lemma~19.23]{QEDI}; to \cite[Lemma~20.4(a)]{QEDI} and \cite[Proposition~20.5]{QEDI}; to the last sentence of \cite[Theorem~6.1]{QEDI}, to the second numbered display of \cite[Theorem~6.3]{QEDI} together with the positivity $\bar{Z}[1]>0$ stated there, and to \cite[Corollary~6.4(iii)]{QEDI}; and to a part of \cite[Section~7.1.2]{QEDI}, to \cite[Proposition~7.2]{QEDI} and to the last sentence of \cite[Theorem~7.11]{QEDI}.
In the notation of \cite[Theorem~6.3]{QEDI}, this display is the Gaussian bound
\[
\overline{\mathsf{E}}\left[\exp\langle g,\cdot\rangle\right]\le\exp\left(\frac{1}{2}\|Pg\|^2\right)\le\exp\left(\frac{1}{2}\|g\|^2\right)\qquad\left(g\in C^2\right)
\]
for the harmonic-averaged law, where $P$ is the orthogonal projection onto the range of $d_1$.
They provide \cite[Theorem~5.85(b)]{QEDI}.
\item[Propositions 32.1--32.3]
They correspond to \cite[Proposition~19.22]{QEDI} and \cite[Theorems~19.24 and 19.25]{QEDI}.
They provide \cite[Theorem~5.85(a), (c)]{QEDI}, \cite[Theorem~5.86]{QEDI} and the reference datum (R2) of Proposition~43.1.
\item[Propositions 33.1--33.4]
They correspond to \cite[Section~20.2.2]{QEDI}, \cite[Lemmas~20.14, 20.15, 20.17 and 20.18]{QEDI}, to the paragraph on the legs in \cite[Section~20.4]{QEDI}, and to \cite[Corollary~20.16]{QEDI}.
\item[Propositions 34.1--34.4]
They correspond to the first paragraph of \cite[Section~20.5.1]{QEDI}, to \cite[Lemma~20.21]{QEDI}, \cite[Section~20.6.1]{QEDI}, \cite[Lemmas~20.22 and 20.23]{QEDI} and \cite[Proposition~20.24]{QEDI}, and provide \cite[Theorem~5.91]{QEDI}.
\item[Propositions 35.1--35.10]
They correspond to \cite[Lemma~20.1(a)--(c)]{QEDI}, to \cite[Lemma~20.2(a), (b)]{QEDI} (together with Proposition~49), \cite[Lemma~20.3]{QEDI}, \cite[Sections~20.1.4 and 20.1.8]{QEDI};
to the paragraph of \cite[Section~20.1.5]{QEDI} that immediately precedes \cite[Lemma~20.4]{QEDI}, by which the entries of the matrix of the Wilson operator are polynomials in the link phases and $N_0$ and the functions $t\mapsto\bar{\nu}^{a(t)}(\mathrm{Ev})$ are entire (its counterpart is Proposition~35.7\,(a));
to \cite[Proposition~20.6]{QEDI}, \cite[Lemma~20.7]{QEDI}, to the bounds (Gr6), (Gr7) of \cite[Section~20.2.4]{QEDI}, which are the clauses (A6), (A7) of Proposition~35.4, to the head of \cite[Section~20.4]{QEDI}, to \cite[Section~20.5.1]{QEDI}, \cite[Lemmas~20.19 and 20.20]{QEDI}, \cite[Proposition~20.25]{QEDI}, \cite[Lemmas~20.26 and 20.27]{QEDI}, \cite[Proposition~20.28]{QEDI}, \cite[Theorem~20.29]{QEDI}, \cite[Proposition~20.31]{QEDI}, \cite[Corollary~20.32]{QEDI} and \cite[Lemma~20.33]{QEDI}.
They provide \cite[Theorems~5.89, 5.90, 5.92 and 5.93]{QEDI}, through the Propositions~35.4, 35.8, 35.9 and 35.10 in this order, the reference datum (R3) of Proposition~43.1, and a part of \cite[Proposition~5.68]{QEDI}, namely Proposition~35.1\,(a), (d), (f).
\cite[Theorem~5.91]{QEDI} is provided by Proposition~34.3 alone.
\cite[Theorem~5.94]{QEDI}, which is proved there as \cite[Corollary~20.30]{QEDI}, has no counterpart in this paper and is not used in it.
\item[Proposition 36]
Apart from its clause (h), it has no counterpart in \cite{QEDI}.
Clause (h) corresponds to the last sentence of \cite[Theorem~21.1(d5)]{QEDI}, by which $b_{\eta,R}(0)$ does not change when $a_\eta$ is replaced by $g_*a_\eta$.
Proposition~36 provides the reference datum (R4) of Proposition~43.1 with $\mathfrak{e}_2=0$.
\item[Propositions 37.1--37.3]
They correspond to two of the definitions of \cite[Section~21.1.9]{QEDI}:
that of the coefficient $b_R(f)$ on the antiperiodic torus of side $R$,
and that of the lattice quantities $N_0(t)$ and $b_{\eta,R}(0)$, formed with a one-cochain $a_\eta$ on $\Lambda_{\eta,R}$.
They correspond further to \cite[Theorem~21.1(c1)--(c3)]{QEDI}, \cite[Section~21.4.3]{QEDI}, \cite[Lemma~21.8]{QEDI}, and to clause~2(b), (c) of \cite[Theorem~21.3]{QEDI}.
They provide \cite[Theorem~5.96(c)]{QEDI} and the reference datum (R5) of Proposition~43.1.
\item[Propositions 38.1--38.3]
They correspond to \cite[Theorem~21.1(d1)]{QEDI}, that is, to (W2)--(W4), the displays (10.4)--(10.6) of \cite[Section~10.1.2]{QEDI}; to \cite[Lemma~10.2(b)]{QEDI}; to the first of the three displayed properties of the determinant in \cite[Section~11.1.2]{QEDI} and to \cite[Lemmas~11.2 and 11.7]{QEDI}; and to the bound of the determinant in \cite[Section~7.1.2]{QEDI} and to \cite[Lemma~7.7(a)]{QEDI}.
They provide a part of \cite[Theorem~5.96(d)]{QEDI}.
\item[Proposition 39]
It corresponds to \cite[Definition~21.18]{QEDI}, \cite[Lemma~21.20]{QEDI} and \cite[Propositions~21.22 and 21.24(a)]{QEDI}; its clause (f) corresponds to \cite[Proposition~12.10(f)]{QEDII} for cell widths in $\frac{1}{2}\mathbb{Z}$, read with (I.h5).
It provides a part of \cite[Theorem~21.2]{QEDI} and of the statements of \cite[Sections~21.7 and 21.8]{QEDI} on which that theorem rests.
\item[Proposition 40]
It corresponds to \cite[Theorem~21.25]{QEDI}, and provides a further part of \cite[Theorem~21.2]{QEDI} and of the statements of \cite[Sections~21.7 and 21.8]{QEDI} on which that theorem rests.
\item[Proposition 41]
It corresponds to a clause of \cite[Theorem~21.1(b2)]{QEDI}, and provides a part of \cite[Theorem~5.96(b)]{QEDI}.
\item[Propositions 42.1 and 42.2]
They correspond to \cite[Theorem~21.2]{QEDI} for the all-flux law, to \cite[Lemma~21.32]{QEDI}, and to Steps 1--4 of \cite[Section~21.8.2]{QEDI}.
They provide a further part of \cite[Theorem~21.2]{QEDI} and of the statements of \cite[Sections~21.7 and 21.8]{QEDI} on which it rests.
\item[Propositions 43.1 and 43.2]
They correspond to \cite[Theorem~21.3]{QEDI}, to Steps~0--3 of the proof in \cite[Section~21.10.1]{QEDI}, and to the beginning of its Step~6 with the display there that bounds the modulus of the expectation of a product over finitely many cells by the corresponding power of $\mathsf{u}_m(N)$;
this display is Proposition~43.2\,(U-b), which rests on the display (2.4) of Proposition~42.2.
They provide \cite[Theorem~5.95]{QEDI}.
\item[Theorem 44]
It corresponds to both parts of \cite[Theorem~3.1]{QEDI}, to \cite[Corollary~3.2]{QEDI}, to Steps~4 and~5 of the proof in \cite[Section~21.10.1]{QEDI} and to the part of its Step~6 that follows the bound by the power of $\mathsf{u}_m(N)$, and to \cite[Section~21.10.2]{QEDI}.
Its range is the all-flux law on the even tori, with the charges $e\in(0,e_2]$ and the mass assignment $N\mapsto m_{0,N}(e)$.
In this range it contains the following assertions of \cite[Theorem~3.1]{QEDI}:
Part~I\,(a), apart from the tightness on $\mathcal{S}'$;
Part~I\,(b);
of Part~I\,(c), the hypothesis (H) for the limit point and, for its ergodic components, the invariance and ergodicity under $\mathbb{R}^3$, the condition (G) and the condition (dF);
Part~I\,(d) without its last sentence, that is, its assertion for the iterated limit in which $T\to\infty$ is taken first;
and Part~II\,(i)--(iv), (vi) and (vii), the clause (i) with the same restriction as Part~I\,(c), and the clause (ii) for $T\ge T^{*}(e)+1$ and with $\eta^{\frac{1}{4}}$ in the place of $\eta^{\frac{1}{2}}$.
In the same range it provides the deduction of Part~II of \cite[Theorem~3.1]{QEDI} and of \cite[Corollary~3.2]{QEDI} in \cite[Section~21.10]{QEDI}, for the clauses that have a counterpart.
For an arbitrary limit point, the other assertions of Part~I\,(c), and hence of Part~II\,(i), are provided by Proposition~45.4\,(b)--(d), applied to the limit point, which is soft by the clause (B1) of Theorem~44, together with \cite[Theorem~8.3(i)--(iv)]{QEDI} and \cite[Proposition~8.19]{QEDI}, that is, by \cite[Corollary~8.9]{QEDI}, whose quantifiers over $\xi$ are recalled in the Remark that follows Proposition~45.4.
These assertions are the invariance of the components under $\{\pm1\}^3$, their reflection positivity at every plane $\{y_j=\mathrm{po}\}$, the moment bound of property~1, the reconstruction and the statement on its invariant vectors, the clustering, the properties~5 and~6, and the assertion that the statements of Part~I\,(c) on the components, other than the properties~5 and~6, hold together for every $\xi$ of one set of full $\nu$-measure.
The last of them is obtained in the same way: the set of the points $\xi$ at which the six statements of Proposition~45.4\,(b) hold has full $\nu$-measure, and at each of its points Proposition~45.4\,(c), together with \cite[Theorem~8.3(i)--(iv)]{QEDI} and \cite[Proposition~8.19]{QEDI}, gives all these statements.
In the same way one obtains the assertions of \cite[Corollary~3.2(a)]{QEDI} on the invariance under $\{\pm1\}^3$, on reflection positivity, on the reconstruction and on clustering.
For the volume-first limit points, the clause (D) of Theorem~44 gives the invariance, the reflection positivity, the statement on the invariant vectors and the clustering in a stronger form.
Theorem~44 has no counterpart of the tightness on $\mathcal{S}'$ in Part~I\,(a);
of the last sentence of Part~I\,(d), in which the two limits are taken in the other order, and which holds by \cite[Remark~21.17]{QEDI} in the reading of Section~I.9.2;
of the clause (v) of Part~II of \cite[Theorem~3.1]{QEDI};
of \cite[Corollary~3.2(b)]{QEDI};
and of \cite[Corollary~3.3]{QEDI}.
\item[Propositions 45.1--45.7]
They correspond to \cite[Theorems~7.15 and 7.17]{QEDI}, \cite[Corollary~7.18]{QEDI}, \cite[Proposition~7.19]{QEDI}, \cite[Lemma~21.36]{QEDI}, \cite[Theorem~8.8(v)]{QEDI}, to \cite[Theorem~8.1]{QEDI} together with the first numbered display of its proof in \cite[Section~8.5]{QEDI}, to \cite[Proposition~9.7(b)]{QEDI}, and to the first assertion of \cite[Corollary~9.10]{QEDI}: there is a set $\Xi_0$ of full $\nu$-measure such that, for every $\xi\in\Xi_0$, the law $\mu_\xi$ is translation invariant and ergodic and satisfies the condition (H1) of \cite[Section~9.2]{QEDI} for all $f\in\mathcal{S}$ simultaneously.
The counterpart of this assertion is Proposition~45.6\,(c).
Proposition~45.4 corresponds, besides \cite[Lemma~21.36]{QEDI}, to Part~I\,(c) of \cite[Theorem~3.1]{QEDI}, that is, to the properties 1--6 of \cite[Corollary~8.9]{QEDI}.
Proposition~45.4\,(b)--(d) keeps apart the almost sure statements, the statements on a single component and the statements whose null set depends on the event or which concern $\nu$;
in this it follows the arrangement of Part~I\,(c) of \cite[Theorem~3.1]{QEDI} and of \cite[Corollary~8.9]{QEDI}, where the properties 1--4 hold for every $\xi$ of one set of full $\nu$-measure, property~5 and the part of property~6 on the events of $\mu$-measure $1$ hold with a null set that depends on the event, and the invariance of $\nu$ under $\xi\mapsto w_{*}\mu_\xi$ is a statement on $\nu$.
Proposition~45.5 corresponds to \cite[Lemma~12.22(a)--(f)]{QEDII}, whose clause (f) is stated only for the normals $e_j$ and for $\mathrm{po}=0$;
of the statements of \cite{QEDI}, \cite[Lemma~8.16(3)]{QEDI}, transported to the plane $\{\hat{n}\cdot y=0\}$, corresponds to the first inequality of Proposition~45.5\,(d).
\item[Proposition 46]
It corresponds to \cite[Section~21.1.4]{QEDI}, where the several written forms of the operator are identified with one another; it stands in Section~I.2.9.
\item[Propositions 47.1 and 47.2]
They correspond to \cite[Theorem~21.1(b1)--(b5)]{QEDI};
to the cells of \cite[Section~21.1.5]{QEDI}, taken on the tori $n=\mathrm{kt}\cdot n_*$ with $\mathrm{kt}$ even, and to the part of Step~6 of the proof in \cite[Section~21.10.1]{QEDI} that precedes the bound on the expectation of the product (Proposition~47.2);
and to the assertion of \cite[Lemma~12.5(c)]{QEDII} on tests supported in a period cube $y_0+\left(-\frac{R}{2},\frac{R}{2}\right)^3$ with an arbitrary centre $y_0\in\mathbb{R}^3$, which is Proposition~47.1\,(b0).
They provide \cite[Theorem~5.96(b)]{QEDI}.
\item[Propositions 48.1--48.6]
They correspond to \cite[Theorems~13.5, 13.7, 13.11, 13.13 and 13.15(c)]{QEDI} and to \cite[Proposition~13.3]{QEDI}, all without the second slot; to \cite[Section~5.3.8]{QEDI}; and to the paragraph on the space in the proof of \cite[Proposition~19.14]{QEDI}.
They provide \cite[Propositions~5.59, 5.60 and 5.62]{QEDI} without the second slot.
\item[Proposition 49]
It corresponds to \cite[Section~20.1.3]{QEDI}, to \cite[Lemma~20.2(a), (b)]{QEDI} (together with Proposition~35.2), and to \cite[Lemma~20.4(b)]{QEDI}.
\item[Proposition 50]
It corresponds to \cite[Theorem~21.1(d2)--(d5)]{QEDI};
to the properties (W1)--(W6), of which (W1) is the clause (d2), (W5) the clause (d3) and (W6) the clause (d4) of \cite[Theorem~21.1]{QEDI}, while (W2)--(W4) form its clause (d1) (Propositions~38.1--38.3);
to \cite[Section~10.1.3]{QEDI}, in which the absolute convergence of the series and the positivity of $Z_n[1]$ are Proposition~50\,(e), the invariances under the sign change and under the isometries, also for the parts of flux zero, are the clause (d5), that is, Proposition~50\,(f), (g), (g$'$), (g$''$), and the product formula for the free Wilson determinant has no counterpart;
to \cite[Section~10.1.4]{QEDI}; to the three displayed properties of the determinant in \cite[Section~11.1.2]{QEDI}; and to a part of \cite[Section~7.1.2]{QEDI} and to \cite[Lemma~7.7(b), (c)]{QEDI}.
It provides \cite[Theorem~5.96(d)]{QEDI}.
\end{description}

Three statements of \cite{QEDI} correspond to statements of \cite{QEDII}.
The trace representation \cite[Proposition~11.5]{QEDI} corresponds to \cite[Corollary~5.30]{QEDII}, \cite[Propositions~5.33, 5.43 and 5.44]{QEDII}, \cite[Lemma~5.35]{QEDII}, \cite[Corollary~5.41]{QEDII} and \cite[Lemma~13.8]{QEDII}.
The Gaussian domination of the all-flux parent, \cite[Theorem~11.6]{QEDI}, corresponds to \cite[Theorem~6.27(a), (b)]{QEDII}, and for the averaged law to \cite[Proposition~13.11]{QEDII}.
The bound on the weight of a flux sector, \cite[Lemma~11.8]{QEDI}, is needed only in its case $g=0$, which is also \cite[Corollary~10.8(d)]{QEDI}; in this case it corresponds to \cite[Lemma~13.12(a)]{QEDII}.

\subsubsection*{I.9.5. Statements of the first paper in the form in which they are used}

The following statements of \cite{QEDI} are used in the form stated here.
\begin{itemize}
\item As \cite[Definition~5.9]{QEDI} states, of the polymer resolvents, the family (p5) of the operators $S_{k,y}(X,A)$, only the symmetry under the subgroup $\mathbb{G}^{L}_{k,3}$ of the elements $(S,a)$ with $S\in\mathrm{Sp}_3$ is required;
Section~I.5.8 requires the same.
\item At $k=0$, property~7 of \cite[Section~5.4.2]{QEDI} is, by its own wording, the kernel bound of the resolvent with the factor $(1+y)^{-1}$ that is displayed in property~17 of \cite[Section~5.4.2]{QEDI}, for all $y\ge0$;
\cite[Proposition~5.48]{QEDI}, its clause (d1) included, assumes property~7 at $k=0$ in this form;
and the Borel measurability of the resolvent in $y$ is part of the definition of a system of primitives of level~$0$ in \cite[Definition~5.29]{QEDI}.
The property W7 at $k=0$ of Section~I.5.8 is the same requirement, read with the norm $|\cdot|_c$ of the $4\times4$ blocks at pairs of fermion sites and with the constants of the chamber fermion.
\item The four inequalities of Section~I.8.4, imposed for all $v\ge-\log e_0$, are, with $v=-\log e_k$, the four inequalities that the condition (A10) of \cite[Section~5.6.2]{QEDI} imposes for every $e_k\in(0,e_0]$;
like (A10) itself, they are conditions on the constants only.
Of the two further inequalities of (A10), the inequality $e_0\le e^{F}_{0}$ is not imposed here, the number $e^{F}_{0}$ not occurring, and the last one is (Z2.s5), formed here with the constant $C_{d}$ of Proposition~13.
See also the Propositions~14.1--14.4.
\item As \cite[Proposition~5.55(d)]{QEDI} states, the inverses $S_{k+1}(B')$ and $S_{k+1}(-B')$ of level $k+1$ are obtained from the invertibility of the operators of level $k$ at $B$ and at $-B$, through the scaling relation of Step~1 of the proof of \cite[Theorem~14.21]{QEDI} and the bijectivity of $f\mapsto f_L$ (\cite[Theorem~14.21(a)]{QEDI}, and Step~4 of its proof for $-B'$);
of $B'$ that proposition requires only that $B'_L$ lie in $\tilde{G}_k$, and $B'$ need not lie in $\tilde{G}_{k+1}$.
For $k\ge1$ the existence of the inverses of level $k$ at $B$ and at $-B$ is contained also in \cite[Propositions~5.42--5.44]{QEDI}.
The Propositions~21.1--21.5 obtain these inverses in the same way (Proposition~21.4\,(b)).
\item By \cite[Remark~14.25(iii)]{QEDI}, a constant $c'$ contained in $A'$ produces the constant $c'$ in $\mathscr{A}'$ and the constant $L^{-\frac{1}{2}}\cdot c'$ in $A$, and on the right side of the identity of \cite[Theorem~14.24]{QEDI} it occurs only through these two constants, because there $A'$ enters only through $\mathscr{A}'=H_{k+1}A'$;
on the left side, and in Steps~(0) and~(b) of the proof of \cite[Theorem~14.24]{QEDI}, $A'$ enters in addition through $A^{\min}_k(A')$ and through $\omega_k(A')$, so that $c'$ occurs there also through $A^{\min}_k(c')$ and $\omega_k(c')$.
The first assertion of clause (d) of Proposition~22.3 corresponds to this Remark and agrees with it: the right side of the formula of Proposition~22.3\,(c) contains $A'$ only through $H_{k+1}A'$.
The second assertion of clause (d), that both sides are unchanged when $\omega_k$ is replaced by $\omega_k+t$ with a real number $t$, corresponds to \cite[Remark~14.25(i)]{QEDI}: the additive constant of $\omega_k$ is immaterial, and this replacement leaves Steps~(b)--(f) of the proof of \cite[Theorem~14.24]{QEDI} unchanged.
As \cite[Remark~14.25(iii)]{QEDI} states, for $c'\ne0$ the constant enters the left side also through the function $\omega_k(c')$, which is not constant on $T_k$;
for this reason the first assertion of clause (d) of Proposition~22.3, that $A'$ enters only through $H_{k+1}A'$, is stated for the right side alone.
\item \cite[Theorem~21.1(b2)]{QEDI} states the assertions on the supports of $\hat{\mathfrak{a}}$ and of $d\hat{\mathfrak{a}}$, on the disjointness of the supports of the cochains $d\hat{\mathfrak{a}}_{\mathfrak{c}}$ and the identity $\|d\hat{\mathfrak{a}}_{\mathfrak{c}}\|^2=\boldsymbol{\sigma}^2$ under the hypotheses that $\eta=L^{-N}$, that the probe $\beta$ fits $(N,T^{*})$ and that $v$ is a placement, on the tori $R=k\cdot R^{*}$, $R^{*}=h\cdot L^{T^{*}}$, which carry the cells;
for $T^{*}=T^{*}(e)$ these are the tori with $n=k\cdot n_*$ sites per side, and under these hypotheses the form used here agrees with \cite[Theorem~21.1(b2)]{QEDI}.
Here (Proposition~41) these assertions are used on the larger class of the cubic tori with $n$ sites per side, $n\in\mathrm{per}\,\mathbb{Z}$ and $n\ge n_*$, which contains the tori $n=k\cdot n_*$ and the tori $n=\mathrm{kt}\cdot n_*$ of Proposition~47.2.
The last assertion of (b2), $X=X_{\eta,R}\left(\tau_vf^{(\eta)}\right)$, and the clause (b5) hold for every $0<\eta\le1$, every $R>D_\beta$ with $R/\eta$ an integer $\ge2$ and every $v\in\mathbb{R}^3$, without the hypothesis that $\beta$ fits (Proposition~47.1\,(b2$'$), (b5)).
\end{itemize}

Further, the following holds.
\begin{itemize}
\item A family of activities on $\tilde{R}_k$ alone is not a special case of a family with two slots.
Hence the forms without the second slot of \cite[Theorems~13.7 and 13.15]{QEDI} do not follow from the forms with the slot by specialisation;
they are stated separately, with the same constants, in the Propositions~48.1--48.6 (Section~I.5.11).
The assertions of \cite[Theorem~13.15(c)]{QEDI} without the second slot hold with the constants printed there; they are Proposition~48.5\,(a), (b) and the last inequality of Proposition~48.5\,(d).
Their proof is the part of the proof of \cite[Theorem~13.15]{QEDI} that concerns the families without the second slot, at the end of that proof: it is the proof of \cite{QEDI} with the second connection put equal to $0$ and the total connection equal to $\mathscr{A}$.
The statements to which it refers include the forms for one slot of \cite[Theorem~13.7(b), (e)]{QEDI}, \cite[Theorem~13.11(ii), (iii)]{QEDI} and \cite[Theorem~13.13(c)--(e)]{QEDI}, which \cite{QEDI} states beside the forms with the second slot: at the end of \cite[Theorem~13.7]{QEDI}, before (i) and at the end of (iii) in \cite[Theorem~13.11]{QEDI}, and at the end of \cite[Theorem~13.13]{QEDI}.
For the chamber fermion, the forms of these statements for one slot, with the same constants, are the last sentence of Proposition~48.2\,(d), which is the form of \cite[Theorem~13.13(c)]{QEDI}, and the Propositions~48.3 and~48.4;
they are proved in Chapter~48a, and Proposition~48.5 is proved from them in Chapter~48b.
\item In \cite[Proposition~13.3(d)]{QEDI} the side $\ell/L$ need not belong to $\mathrm{Cl}$: the objects of level $k+1$ that occur there require only that $\ell/L$ be a multiple of $M$ and that $\ell/L>2M$ (Section~I.5.11).
\item The recursion of \cite[Section~5.3.8]{QEDI} has the invariant $E_k=\mathscr{E}_k$ at $\mathscr{C}=0$, for every $k$ for which the trajectory exists (Section~I.5.11).
\item \cite[Proposition~5.48]{QEDI}, its clauses (d4)--(d6) included, is stated for the parameter set $\mathrm{par}$ alone, with the constant $C_{d}$ formed with the constant $C_{w}$ of $\mathrm{par}$, and for the ceiling $e_0$;
in this point Proposition~13 agrees with it.
Clause (iii) of \cite[Theorem~5.72]{QEDI} concerns the parameter set $(\alpha_0,\gamma^{F},\kappa_{\mathrm{w}};C^{F})$, with $C_{d}$ formed with $C_{w}=C^{F}$ and with $e^{F}_{0}$ in the place of $e_0$; it has no counterpart here.
After its clause (iii), \cite[Theorem~5.72]{QEDI} states that, for $\kappa'=\kappa_{\mathrm{w}}+\kappa_0$, the system of \cite[Theorem~5.72(ii)]{QEDI} has the properties required in \cite[Proposition~5.48]{QEDI} also for $\mathrm{par}$, so that \cite[Proposition~5.48]{QEDI} applies to it as stated, for every $e_k\in(0,e_0]$;
the counterpart of this statement is Proposition~14.1\,(ii), (iii).
The ceilings of the running charge differ:
in \cite{QEDI} the running charge $e_k$ has the ceiling $e^{F}_{0}$ in \cite[Theorem~5.72]{QEDI} besides the ceiling $e_0$ of \cite[Proposition~5.48]{QEDI}, and these ceilings are tied by the inequality $e_0\le e^{F}_{0}$ of the condition (A10) of \cite[Section~5.6.2]{QEDI},
whereas here the number $e^{F}_{0}$ does not occur, this inequality is not imposed, and in the data (SF) and (SF+) of Propositions~13 and~14.1 the running charges $e_k$ and $e_{k+1}=L^{\frac{1}{2}}\cdot e_k$ have the ceiling $e_0$;
the charge $e$ is bounded by further ceilings named in the propositions, among them $e_1$ through $e\cdot h^{\frac{1}{2}}\le e_1$, $e_{G}$ and $e_{\sharp}(\beta)$ (Section~I.8.1\,(ii)), and $e_2$ (Theorem~44).
As in stage~9 of \cite[Definition~5.97]{QEDI}, the number $e_0$ itself is subject to the conditions (Z2.s1)--(Z2.s5), (Z2.s5) with the constant $C_{d}$ formed with the constant $C_{w}$ of $\mathrm{par}$ (Section~I.8.4, Propositions~14.1 and 14.2);
there these conditions are formed with the constants of \cite{QEDI}, and Section~I.8.4 states in which constants they differ.
\item The quantities of \cite[Section~20.1.7]{QEDI} that depend on $e$ are $\tau_s$, $\bar{\kappa}_1$, $\bar{r}$, $B_1$, $B_{\mathrm{wr}}$, $B_0$, $\varpi_{\mathrm{sep}}$, $\varpi_0$, $t_2$, $\varpi_L$ and $\lambda_L$; all of them stand in the display of Proposition~14.2\,(d).
Besides them, that display contains the number $\varpi_{\mathrm{sep}}^{+}$, which is not a quantity of \cite{QEDI};
the numbers $\mathrm{pp}$, $t_x$ and $\mathbf{r}$ of \cite[Section~20.1.4]{QEDI}, which also depend on $e$;
and the number $\mathrm{mxL}$, with which $\lambda_L$ is formed.
The constants of the line (C6$'$), the constant $C_{\mathrm{fr}}$ of the line (C10), $\mathrm{mxL}$ (the number $M_L$ of \cite{QEDI}) and $z_2$ do not depend on $e$.
Section~I.8.3 names, among these quantities, only those into which the numerals listed there enter.
\item In \cite[Section~17.1.3]{QEDI} the openness of $\mathfrak{X}_k$ is used, and of the pairs considered there only the pair $(\{0\},\{0\})$ is needed (Propositions~24.1--24.3).
\item The rules (T-G) and (Pr1)--(Pr3) of \cite[Section~17.11.1]{QEDI} express an invariance relative to a polymer (Propositions~24.1--24.3).
\item The Osterwalder--Schrader condition (E0$'$), which occurs in \cite[Theorem~7.17]{QEDI}, is taken in the wording of Theorem~44.
\item \cite[Corollary~3.3]{QEDI} makes its first assertion for each integer $N$ with $\boldsymbol{\mu}_N>0$, and its second for the integers $N$ of an infinite set of integers with $\boldsymbol{\mu}_N>0$;
it requires the positivity of $\boldsymbol{\mu}_N$ only at the integers $N$ concerned, and it leaves open whether such an infinite set exists.
Here no sign of $\boldsymbol{\mu}_N$ is assumed ((I.K10)).
\end{itemize}

The following points of \cite{QEDI} are read with the conventions of Part~I:
the numbering of the directions, the lifts, and the words ``does not wrap'' (Section~I.1.1);
the words ``given by one formula'' (Section~I.1.9);
the kernels of the fine operators (Section~I.1.6);
the values of a function that enter a H\"older quotient (Section~I.1.8);
the norm on $\mathbb{C}^4$ (Section~I.1.5);
the objects of level~$0$ (Section~I.5.10);
\cite[Proposition~5.45]{QEDI} at $k=0$ ((I.L0.15));
the Berezin integral on a general index set (Section~I.3.6);
the family $\mathfrak{P}_k$ of the complex functions $\mathfrak{P}_k(X,\cdot)$ of the external connection $\mathscr{C}\in\mathfrak{X}_k$, which is arbitrary in the statements on one step, while in a flow it is given by $\mathfrak{P}_0=0$ and the recursion of Section~I.5.7;
the rate $\kappa'=\kappa_{\mathrm{w}}+\kappa_{0}$ in $M^{F}$ and $C^{F}$ (Section~I.8.2);
the number $m_0$ (Section~I.6.1);
the twist and the external connection (Section~I.6.6);
the sign $\boldsymbol{\chi}(w)$ in the kernels of a projection (Section~I.6.3).
The quotients $Z_n[\,\cdot\,]/Z_n[1]$ and $Z^{(0)}_n[\,\cdot\,]/Z^{(0)}_n[1]$ that define the laws $E^{\mathrm{af}}_n$ and $E^{\mathrm{av}}_n$ are defined for every real lattice mass, since the partition functions are positive for every real mass (\cite[Theorem~4.26]{QEDII}; Proposition~50\,(e)).

\subsubsection*{I.9.6. Lattice mass, weights, fermion sites, cells; the factor $48$}

The following conventions hold throughout the paper.

The lattice mass $\boldsymbol{\mu}$ is an arbitrary real number ((I.K10)):
the statements hold without any of the conditions $\boldsymbol{\mu}>0$, $\boldsymbol{\mu}_N>0$, $\boldsymbol{\mu}>-\frac{1}{2}$, $\boldsymbol{\mu}>-1$.
The weights and $\rho$ are bare ((I.BARE)): neither the weights nor $\rho$ are renormalised.

Fermion sites are the points $x=(c,g)$ of Section~I.1.3.
The fields $\psi$, $\bar{\psi}$, $\chi$, $\bar{\chi}$, the spinors of $\mathfrak{S}^{\mathrm{ch}}(\eta_k,\ell)$ and the fine indices of fermionic kernels are evaluated at fermion sites;
at level~$0$ the same holds for the fields $\Psi$, $\bar{\Psi}$, $W$, $\bar{W}$, for the spinors of $\mathfrak{S}^{\mathrm{ch}}(1,\ell)$, and for the operators on this space and their kernels ((I.L0.1), (I.L0.3), (I.L0.4), Section~I.3.9).
A photon site, the centre of a cube or a point of $\eta_k\mathbb{Z}^3$ is not a fermion site.
At every level $k\ge0$ the spinors of $\mathfrak{S}(L,\ell)$, the block field of the next level in its scaled form $\Psi'_L$, $\bar{\Psi}'_L$, and the indices on the $L$-lattice, among them those of the pairing $\langle\cdot,\cdot\rangle_L$, are evaluated at sites of $L\mathbb{Z}^3$.
For $k\ge1$ the block fields $\Psi$, $\bar{\Psi}$, $W$, $\bar{W}$ of level~$k$, the spinors of $\mathfrak{S}(1,\ell)$ and the unit indices of kernels are evaluated at unit sites.
The extraction of Section~I.5.6 removes, besides the constant part $\mathrm{al}_{0}(E,X)\cdot\operatorname{Vol}(X)$ of a small polymer $X$, whose numbers $\mathrm{al}_{0}(E,X)$ give the vacuum-energy number $\varepsilon(E)=-\sum_{X\text{ small},\,X\supset\square}\mathrm{al}_{0}(E,X)$, with $\square$ an $M$-cube,
only the two-point kernel integrated over both cluster labels, which gives one real mass, uniform on the $48$ sites of a cluster;
no other bilinear is extracted, and differences inside a cluster are H\"older legs ((I.E2)).

The even tori are related to the reference tori only through numbers, namely partition functions and expectations of cell observables (Section~I.6.3), and the content of a cell is read in the open cell (Section~I.6.4).

For the norms on $\mathbb{C}^4$ the convention of Section~I.1.5 holds.
In all numbered properties, with the one exception of Proposition~19\,(e), which is stated in the Euclidean norm $|\cdot|_2$, and in all sup norms $|\cdot|_\infty$, the norm on $\mathbb{C}^4$ is $|\cdot|_c$;
a bound that is proved with a constant $C$ in the Euclidean norm is stated in the norm $|\cdot|_c$ with the constant $2C$.
The $L^2$ norms are formed with $|\cdot|_2$, and so is the Euclidean sup norm $|\cdot|_{\mathrm{su}}$, in which the Propositions~1.1--1.7, 2.1--2.8, 3--6 and 9\,(c) and the walk sums in Proposition~3 are stated.
Kernels are formed by (I.1.4).
In the numbered properties of level~$0$, the factor $48$ of the inverse weight enters only through the constant $C_{00}$, and the factor $\sqrt{48}$ only through $\mathsf{w}_W$.
The constants $C_{G}$ and $C_{3}$ of the lines (C0r) and (C4r) contain the numbers $48$ and $7\ge\sqrt{48}$ at all levels:
$C_{G}$ through the comparison of a kernel with the operator norm on one cluster (Proposition~1.5),
and $C_{3}$ through the comparison of the value at one site with the $L^2$ norm (Proposition~1.6).

\subsubsection*{I.9.7. Six conventions: choice and translation}

For each of the following six conventions, the list gives its form in \cite{QEDI} and in \cite{QEDII}, the form chosen here, and the translation between the forms.

\begin{description}
\item[Numbering of the directions]
In \cite{QEDI} the directions are numbered $0,1,2$ in \cite[Sections~5, 11 and 13--20]{QEDI}, in \cite[Section~21.3.3]{QEDI} and in \cite[Section~21.7]{QEDI}, and $1,2,3$ elsewhere.
In \cite{QEDII} they are numbered $1,2,3$, the time direction being $e_1$.
Here they are numbered $1,2,3$ everywhere.
A direction $\mu$ of \cite{QEDI} in the sections just named is the direction $\mu+1$ here (Section~I.1.1).
\item[Gamma matrices]
Where \cite{QEDI} defines the matrices of the lattice model, they are the following.
In \cite[Sections~5.2.4, 11 and 13--20]{QEDI}:
\[
\gamma_\mu=\sigma_3\otimes\sigma_{\mu+1}\quad(\mu=0,1,2),\qquad\gamma_3=\sigma_1\otimes1 .
\]
In \cite[Sections~1--4, 6--10, 12 and 21]{QEDI}: $\gamma_\mu=\operatorname{diag}(\sigma_\mu,-\sigma_\mu)$ for $\mu=1,2,3$, and the matrix $\sigma_1\otimes1$ is written $\gamma_4$ in \cite[Section~7.1.2]{QEDI} and $\hat{\Gamma}$ in \cite[Section~21.1.1]{QEDI}.
In \cite{QEDII}: $\gamma_\mu=\operatorname{diag}(\sigma_\mu,-\sigma_\mu)$, and $\gamma_5$.
Here: the matrices of \cite{QEDII}, with $\Sigma_3$ added.
At the places where both papers define them, the matrices of the two papers are equal, and no sign differs.
In the first group of sections of \cite{QEDI} the matrix $\gamma_\mu$ is $\gamma_{\mu+1}$ here, and $\gamma_3$ is $\gamma_5$; in the second group the matrices $\gamma_1$, $\gamma_2$, $\gamma_3$ are those of this paper, and $\gamma_4$ and $\hat{\Gamma}$ are $\gamma_5$ (Section~I.2.1).
By Proposition~46\,(c), the following places of \cite{QEDI} do not define these matrices.
\cite[Section~6.1.2]{QEDI}, in its first sentence, admits every Hermitian Clifford triple; the four-component representation defined after that sentence is the triple $\gamma_1$, $\gamma_2$, $\gamma_3$.
\cite[Section~12.1]{QEDI} admits every Hermitian Clifford triple of $4\times4$ matrices and, as a free datum, every Hermitian matrix $\hat{\Gamma}$ with $\hat{\Gamma}^2=1$;
here the triple is taken to be $\gamma_1$, $\gamma_2$, $\gamma_3$ and $\hat{\Gamma}$ to be $\gamma_5$, as \cite{QEDI} itself does in the proofs of \cite[Proposition~16.6]{QEDI} and \cite[Lemma~16.7]{QEDI}.
\cite[Section~11.1.2]{QEDI} names the three Pauli matrices without fixing their order and defines no $\gamma_3$; their order is the one given by the numbering $\mu\mapsto\mu+1$ of the directions (Section~I.1.1).
In \cite[Section~6]{QEDI}, from the paragraph on the spinor basis adapted to the time direction in \cite[Section~6.3.2]{QEDI} on, the letters $\gamma_\nu$ and $\gamma_i$ denote $w\gamma_\nu w^{*}$ and $w\gamma_iw^{*}$, with a unitary matrix $w$ chosen there, and not the matrices $\gamma_1$, $\gamma_2$, $\gamma_3$ of Section~I.2.1;
the triple $\left(w\gamma_\mu w^{*}\right)_\mu$ is unitarily conjugate to the triple used before that paragraph, and for $s=4$ one has $w\gamma_\nu w^{*}=\Sigma_3$.
\item[Charge conjugation]
In \cite{QEDI}: $\mathsf{C}=i\gamma_1$, also for $\gamma_3$, acting on the fields by $\psi\mapsto\mathsf{C}\bar{\psi}$ and $\bar{\psi}\mapsto-\mathsf{C}\psi$.
In \cite{QEDII}: $\mathrm{Cc}=\operatorname{diag}(\sigma_2,\sigma_2)$, acting on the operator only.
Here: $\mathsf{C}=i\gamma_2$ in the numbering of the directions fixed above, together with the automorphism $\mathrm{Cg}$.
The two matrices are related by $\mathsf{C}=i\,\Sigma_3\cdot\mathrm{Cc}$.
Under transposition with $\mathsf{C}$ the resolvent parameter changes its sign; with $\mathrm{Cc}$ it does not (Section~I.4.5).
\item[Orientation and pull-backs]
In \cite{QEDI}:
\[
(gA)(x,x')=A\left(g^{-1}x,g^{-1}x'\right),\qquad g_*=\left(g^{-1}\right)^*,\qquad(dA)_{\mu\nu}=\partial_\mu A_\nu-\partial_\nu A_\mu .
\]
In \cite{QEDII}: $(g^*u)(\boldsymbol{\sigma})=\epsilon(g,\boldsymbol{\sigma})\cdot u(g\boldsymbol{\sigma})$, with the same coboundary $d$.
Here both conventions are used; they agree.
One has $\tilde{g}A=\tilde{g}_*A$.
The mirror copies $O^{[\mathfrak{c}]}$ are those of \cite[Definition~2.19]{QEDII}, formed with the action $G\mapsto G^{g}$ defined after \cite[Definition~2.12]{QEDII}, and the signs are those of \cite[Lemma~12.3]{QEDII} (Sections~I.1.1, I.4.7 and I.6.4).
\item[Fourier transform and Bloch form]
In \cite{QEDI} the Fourier transform is $\hat{f}(p)=\int\mathrm{d}y\,\exp(-ip\cdot y)\,f(y)$, and the symbols are formed with $\exp(ip\cdot c)$.
\cite{QEDII} uses the same transform.
Here the transform is $\hat{f}(p)=\int\mathrm{d}y\,\exp(-ip\cdot y)\,f(y)$, and the Bloch form is $\exp(ip\cdot c)\cdot u(g)$.
In the other Bloch form, $\exp\left(ip\cdot\mathrm{geo}(x)\right)\cdot u(g)$, the symbol $\mathrm{Dh}(p)$ is related to the symbol $D^{(p)}$ of (I.2.6) by
\[
D^{(p)}=Y\cdot\mathrm{Dh}(p)\cdot Y^{-1},\qquad(Yu)(g)=\exp\left(ip\cdot(gv_0)\right)\cdot u(g)
\]
(Section~I.2.6).
\item[Berezin integral]
In \cite{QEDI} the Berezin integral is the coefficient of $\prod W\bar{W}$, so that $\int\exp\left(-\langle\bar{\Psi},O\Psi\rangle\right)=\det O$, with the counting pairing (\cite[Section~13.2]{QEDI}).
In \cite{QEDII} it is normalised by
\[
\int\mathrm{d}\bar{\psi}_\alpha(x)\,\mathrm{d}\psi_\alpha(x)\;\psi_\alpha(x)\,\bar{\psi}_\alpha(x)=1,\qquad\text{so that}\qquad\int_Y\exp\left(-\bar{\psi}D\psi\right)=\det D,
\]
again with the counting pairing (\cite[Definition~2.8]{QEDII}).
Here the normalisation is (I.3.5), which is the same for $\nu=1$; for the pairing with the weight $\frac{1}{48}$ the integral is normalised anew by
\[
\int\mathrm{D}\Psi\,\exp\left(-\langle\bar{\Psi},O\Psi\rangle^{\mathrm{ch}}\right)=\det O .
\]
At level~$0$, $\mathrm{D}\Psi$ is $48^{192n_0}$ times the Berezin integral of \cite[Definition~2.8]{QEDII}; the covariance of $\mathrm{d}\mu_I(W)$ is $48$; and every generator costs a factor $\sqrt{48}$ (Sections~I.3.6 and I.5.2, Proposition~46\,(e)).
\end{description}

\subsection*{Appendix to Part I}

\subsubsection*{A1. Computations}

\paragraph{Spin matrices.}
Each of the following facts concerns finitely many explicit $4\times4$ matrices and follows from a direct computation with them.
\begin{itemize}
\item The relations (I.c1)--(I.c6) of Section~I.2.1.
\item $\varkappa=\sigma_2\otimes\sigma_2$.
\item $\#|\mathrm{Sp}|=32$, $\ker r=\{\pm1,\pm i\Sigma_3\}$, $\#|\mathrm{Sp}_3|=16$ and $r(\mathrm{Sp}_3)=\{\pm1\}^3$.
\item The table of Section~I.4.1, and the identity (I.4.1) for every $S\in\mathrm{Sp}$.
\item $\left(S^{-1}\right)^{T}=\pm\mathsf{C}^{-1}\cdot S\cdot\mathsf{C}$ for every $S\in\mathrm{Sp}$.
\item For the unit normal $\hat{n}$ of each of the nine reflections contained in $O_{\mathrm{cub}}$ one has $S_{\mathrm{nv}}^2=-1$ and, for every $v\in\mathbb{C}^3$,
\[
S_{\mathrm{nv}}\cdot(\gamma\cdot v)\cdot S_{\mathrm{nv}}^{-1}=\gamma\cdot\left(v-2(v\cdot\hat{n})\,\hat{n}\right).
\]
\item The $4\times4$ matrices that commute with the nine lifts $S_{\mathrm{nv}}$ form a space of dimension $2$, and $\sigma_2\otimes1_2$ belongs to it.
Those that commute with the three lifts $\gamma_\mu\gamma_5$ of the reflections in the coordinate planes also form a space of dimension $2$;
if $i\Sigma_3$ is adjoined to these three lifts, the dimension becomes $1$.
The matrices that commute with every element of $\mathrm{Sp}$ form a space of dimension $1$.
\end{itemize}

\paragraph{Hopping matrices.}
If $ge_1=\pm e_\mu$, then $M_3(x)=3\left(\rho\mp\gamma_\mu\right)$ by (I.2.1).
For $\mu=1,2$ the matrix $\gamma_\mu$ has exactly four non-zero entries, all of modulus one and none on the diagonal, whereas $\gamma_3=\operatorname{diag}(1,-1,-1,1)$ (Section~I.2.1).
As $0<\rho<1$, the moduli of the sixteen entries of $M_3(x)$ therefore add up to $4\cdot3\rho+4\cdot3=12(1+\rho)$ if $ge_1\in\{\pm e_1,\pm e_2\}$, and to $2\cdot3(1-\rho)+2\cdot3(1+\rho)=12$ if $ge_1=\pm e_3$, in accordance with (I.m7).

\section*{Part II. The statements, in the order of dependence}

A law $\mathrm{lw}$ on $\mathcal{S}'$ has the property (G), Gaussian domination with the constant one, if
\[
\mathrm{lw}\left(\exp\left(X(f)\right)\right)\le\exp\left(\frac{1}{2}\|f\|_2^2\right)\qquad\text{for every real }f\in\mathcal{S};
\]
here $X(f)$ is the function $F\mapsto F(f)$ on $\mathcal{S}'$, and $\mathrm{lw}(\Phi)$ is the integral of a function $\Phi$ on $\mathcal{S}'$ with respect to $\mathrm{lw}$.
This is condition~(G) of \cite[Section~2.5]{QEDI} and of \cite[Definition~2.38]{QEDII};
for laws on the closed two-cochains of $\mathbb{Z}^3$ the corresponding condition is condition~(G) of \cite[Definition~2.5]{QEDII}.

The data of each statement of this Part are the constants and the further data that the statement names.
The constants are those of Section~I.8, as Proposition~14.2 lists them.
They are chosen in the order of that list, each of them sufficiently large or sufficiently small,
and the conditions on the constants that are stated in this Part hold for this choice.

\subsection*{G1. Definitions for Propositions 1.1--14.4 that are not in Part I}

\subsubsection*{G1.1. Data}

\begin{description}
\item[(SF)] (scale-free data; Section~I.1.8)
A quadruple $(k,\ell,\mu,e_k)$ formed by an integer $k\ge0$, a side $\ell\in\mathrm{Cl}$, a mass $\mu\in(0,1]$ and a number $e_k\in(0,e_0]$.
No relation among the four entries is assumed.
All objects of Sections~I.1--I.5 are formed with these data, and so are the numbers
\[
v_k:=-\log e_k,\qquad r_k:=\min\left\{L^i:\ i\ge0\ \text{an integer},\ L^i\ge v_k^r\right\},
\]
which are taken with the $e_k$ of the quadruple.
\item[(SF+)] (data of a step)
The data (SF) together with $\ell/L\in\mathrm{Cl}$, $\ell\ge M\cdot r_k$ and $e_{k+1}:=L^{\frac{1}{2}}\cdot e_k\le e_0$.
At the level $k+1$ the side is $\ell/L$ and the charge is $e_{k+1}$.
\end{description}
Throughout, the constants are those of Section~I.8, with the definitions and the conditions of Proposition~14.2.
The word ``constant'' has the meaning of (I.U4)$^{\mathrm{ch}}$ in Section~I.8.1;
in particular no constant depends on $(k,\ell,\mu,e_k)$, on $A$, on $y$, on $s$ or on $X$.
In all statements, $\kappa':=\kappa_{\mathrm{w}}+\kappa_0$.

\subsubsection*{G1.2. Block tori}

The photon part of the following definition coincides with \cite[Definition~16.1]{QEDI}; the spinors are those of Section~I.1.5.

A \emph{block torus of level $k\ge0$} is a triple $\mathbb{T}=(t,R,\boldsymbol{\sigma})$ in which $t\in\left\{0,\frac{\eta_k}{2}(1,1,1)\right\}$, $R$ is an integer multiple of $L$ with $R\ge3L$, and $\boldsymbol{\sigma}\in\{1,-1\}$.
Its lattice is $\mathrm{Lab}:=t+\eta_k\mathbb{Z}^3$.
For $0\le j\le k+1$ let $t_j:=0$ if $t=0$ and $t_j:=\frac{L^j\eta_k}{2}(1,1,1)$ otherwise, let $\mathrm{Lab}_j:=t_j+L^j\eta_k\mathbb{Z}^3$, and put
\[
\mathrm{Bl}_j(z):=\left\{c\in\mathrm{Lab}:\ |c-z|<\frac{L^j\eta_k}{2}\right\}\qquad\left(z\in\mathrm{Lab}_j\right).
\]
The unit lattice is $\mathrm{Lab}^0:=\mathrm{Lab}_k$, with the unit blocks $\Delta_z:=\mathrm{Bl}_k(z)$ for $z\in\mathrm{Lab}^0$;
the $L$-lattice is $\mathrm{Lab}_{k+1}$, and $B(z'):=\mathrm{Lab}^0\cap\mathrm{Bl}_{k+1}(z')$ for $z'\in\mathrm{Lab}_{k+1}$.
Since $L$ is odd, $\mathrm{Lab}_{j+1}\subset\mathrm{Lab}_j$;
the blocks of one side form a partition of $\mathrm{Lab}$, the blocks of the different sides are nested,
and all these objects are invariant under the translations by $R\mathbb{Z}^3$.
The distance $d(\cdot,\cdot)$ is the sup distance modulo $R\mathbb{Z}^3$.

\emph{Fermion sites.}
The fermion sites of $\mathbb{T}$ are the points of $\mathrm{Lab}\times O_{\mathrm{cub}}$.
The neighbours $x^{(i)}$, the site $c_3$, the directions $\hat{d}_i$, the matrices $M_i(x)$, the distance $d_X$ and the rule (I.POS) are those of Sections~I.1.3, I.1.4, I.1.7 and~I.2.2.
These displays contain $c$ only through differences, so that the shift $t$ plays no role.

\emph{Spinors, connections, scalars.}
A spinor on $\mathbb{T}$ is a function $f:\mathrm{Lab}\times O_{\mathrm{cub}}\to\mathbb{C}^4$ with $f(c+Re_\mu,g)=\boldsymbol{\sigma}\cdot f(c,g)$.
A spinor on the unit lattice, respectively on the $L$-lattice, is a function $v$ from $\mathrm{Lab}^0$, respectively from $\mathrm{Lab}_{k+1}$, to $\mathbb{C}^4$ with $v(z+Re_\mu)=\boldsymbol{\sigma}\cdot v(z)$;
at $k=0$ the spinors on the unit lattice are the spinors on $\mathbb{T}$ (Section~I.3.9).
The connections are the $R\mathbb{Z}^3$-periodic bond functions that change sign under the reversal of the bond, and the scalars are the $R\mathbb{Z}^3$-periodic functions of $c$.

\emph{Operators.}
Let $A$ be a connection on $\mathbb{T}$, let $e_k>0$ and $\mu\ge0$.
The operators
\begin{gather*}
D_A,\quad Q_k(A)J^{*},\quad JQ_k^{T}(-A),\quad Q(A),\quad Q^{T}(-A),\quad P(A),\\
JP_k(A)J^{*},\quad JP_{k+1}(A)J^{*},\quad\operatorname{Pen}_y(A),\quad F_y(A),\quad S_y(A),\quad\delta_{\alpha,A}
\end{gather*}
are formed on $\mathbb{T}$ by the display (I.2.3), by Sections~I.3.2, I.3.3, I.3.5 and~I.3.9, and by the display (I.1.7).
The function $\delta_{\alpha,A}f$ is formed at every ordered pair $(x,x')$ of distinct points of the covering lattice $\mathrm{Lab}\times O_{\mathrm{cub}}$, with the paths of these two points;
these pairs include the short pairs of Section~G1.3 and, on $T_k$, the pairs of $\mathrm{Pzt}_k$.
The measure, the pairings, the transposes and the adjoints are those of Section~I.1.6, formed with a period cell of $R\mathbb{Z}^3$.
By (I.O0), applied with $\theta=0$ if $\boldsymbol{\sigma}=-1$ and with $\theta=(\pi,\pi,\pi)$ if $\boldsymbol{\sigma}=1$,
the facts (I.s1)--(I.s5), (I.m1)--(I.m3), (I.m5)--(I.m7), (I.2.4) and (I.O1)--(I.O6) hold on $\mathbb{T}$.

\emph{Kernels.}
The kernel of an operator on $\mathbb{T}$ is taken as in (I.1.4), with the sign rule of $\mathbb{T}$:
\[
(Of)(t)=\int\mathrm{d}x'\,O(t,x')\,f(x'),
\]
where $x'$ ranges over one period cell, for every choice of the cell, and
\[
O(t+Re_\mu,x')=O(t,x'+Re_\mu)=\boldsymbol{\sigma}\cdot O(t,x').
\]
Consequently the modulus of $O(t,t')$, and the question whether it vanishes, depend only on the classes of $t$ and $t'$ modulo $R\mathbb{Z}^3$.
A phrase of the form ``the kernel vanishes at $(t,t')$ unless $t$ and $t'$ satisfy a relation'' always means that it vanishes unless the relation holds between $t$ and $t'+Rn$ for some $n\in\mathbb{Z}^3$.
For scalars $\psi$, $\psi'$, that is, for $R\mathbb{Z}^3$-periodic functions of position,
\[
\left(\psi O\psi'\right)(t,t')=\psi(t)\cdot O(t,t')\cdot\psi'(t').
\]

\emph{Morphisms.}
Let $\mathbb{T}'$ be a block torus of level $k$ with the same $R$ and the same $\boldsymbol{\sigma}$, and let $\mathrm{Lab}'_j$ be its lattices.
A morphism $\mathbb{T}\to\mathbb{T}'$ is a map $\mathrm{mph}(c)=rc+\mathrm{tv}$ of $\mathbb{R}^3$, with $r\in\{\pm1\}^3$ and $\mathrm{tv}\in\mathbb{R}^3$, such that $\mathrm{mph}(\mathrm{Lab}_j)=\mathrm{Lab}'_j$ for $0\le j\le k+1$.
On fermion sites it acts by $\mathrm{mph}(c,g):=(\mathrm{mph}\,c,rg)$.
For $S\in\mathrm{Sp}$ with $r(S)=r$, the pair $\mathrm{mS}:=(S,\mathrm{mph})$ acts by
\[
(\mathrm{mS}\,f)(x):=S\,f\left(\mathrm{mph}^{-1}x\right),\qquad(\mathrm{mS}\,A)(c,c'):=A\left(\mathrm{mph}^{-1}c,\mathrm{mph}^{-1}c'\right),
\]
and on $4$-component spinors on $\mathrm{Lab}$, on $\mathrm{Lab}^0$ and on $\mathrm{Lab}_{k+1}$ by $(\mathrm{mS}\,v)(z):=S\,v\left(\mathrm{mph}^{-1}z\right)$.

The torus $T_k$ is the block torus $(0,\ell,-1)$.
Every $\tilde{g}=(S,a)\in\mathbb{G}_k$ induces the morphism $c\mapsto r(S)c+a$ of $T_k$ onto itself,
and the action of the corresponding pair $\mathrm{mS}$ is the action of $\tilde{g}$ described in Section~I.4.3.

\subsubsection*{G1.3. Norms}

\emph{Euclidean sup norms.}
Let $|f|_{\mathrm{su}}:=\sup|f(\cdot)|_2$, the supremum extending over the sites (fermion sites or unit sites) or over the pairs,
and let $\|O\|_{\mathrm{su}}$ be the norm of a linear operator $O$ with respect to these norms.
Let $\|u\|_{L^2}:=\left((u,u)^{\mathrm{ch}}\right)^{\frac{1}{2}}$, and let $\|O\|_{L^2}$ be the operator norm on $L^2$.
By Section~I.1.5 one has $|f|_\infty\le|f|_{\mathrm{su}}\le2|f|_\infty$.
Hence a bound $\|O\|_{\mathrm{su}}\le C$ implies $|Of|_\infty\le2C\cdot|f|_\infty$, that is, $|O|_{1,\infty}\le2C$ in the notation of (I.ke3);
moreover $|m|_c\le2|m|_2$ for every $4\times4$ block $m$.

\emph{Short pairs.}
The short pairs on $\mathbb{T}$ are the ordered pairs $(x,x')$ of distinct points of the covering lattice $\mathrm{Lab}\times O_{\mathrm{cub}}$ with $d_X(x,x')<1$.
For $k\ge1$ every pair inside a cluster is a short pair; at $k=0$ there is no short pair.
A function of short pairs is a $\mathbb{C}^4$-valued function $F$ on the short pairs with $F(x+Re_\mu,x'+Re_\mu)=\boldsymbol{\sigma}\cdot F(x,x')$;
it is measured by the norm $|F|_{\mathrm{su}}$.
In $N_\gamma(\delta_{\alpha,A}O)$ the H\"older quotient is always restricted to the short pairs.

\emph{Localisation.}
On spinors on $\mathbb{T}$, $\mathbf{1}_{\Delta_z}$ is the multiplication by the indicator function of the orbit of $\Delta_z\times O_{\mathrm{cub}}$ under the translations by $R\mathbb{Z}^3$.
On functions of pairs it restricts to the pairs whose first point lies over the orbit of $\Delta_z$;
on spinors on the unit lattice ($k\ge1$) it is the multiplication by the indicator function of the class of $z$.
A function $\psi$ of $c$ acts on functions of pairs by $(\psi F)(x,x'):=\psi(c)\cdot F(x,x')$.
For $\mathrm{vv}\in\mathbb{C}^4$ and a unit site $z'$ of the torus, $\mathbf{1}_{z'}\mathrm{vv}$ denotes the spinor on the unit lattice that vanishes outside the class of $z'$ and takes the value $\mathrm{vv}$ at a chosen representative of $z'$;
another choice of the representative multiplies it by $1$ or by $-1$, which leaves the moduli of its images unchanged.

\emph{Localised norm.}
Let $\gamma\ge0$, and let $O$ be a linear operator between any two of the following spaces: the spinors on $\mathbb{T}$, the spinors on the unit lattice, the functions of short pairs.
Put
\[
N_\gamma(O):=\sup_z\sum_{z'}\exp\left(\gamma\cdot d(z,z')\right)\cdot\left\|\mathbf{1}_{\Delta_z}O\mathbf{1}_{\Delta_{z'}}\right\|_{\mathrm{su}},
\]
where $z$ and $z'$ range over $\mathrm{Lab}^0$ modulo $R\mathbb{Z}^3$.
For all such operators the following holds.
\begin{itemize}
\item[(G1.N1)] $\|O\|_{\mathrm{su}}\le N_\gamma(O)$, and $\left\|\mathbf{1}_{\Delta_z}O\mathbf{1}_{\Delta_{z'}}\right\|_{\mathrm{su}}\le\exp\left(-\gamma\cdot d(z,z')\right)\cdot N_\gamma(O)$.
\item[(G1.N2)] $N_\gamma(O_1+O_2)\le N_\gamma(O_1)+N_\gamma(O_2)$ and $N_\gamma(O_1O_2)\le N_\gamma(O_1)\cdot N_\gamma(O_2)$.
\item[(G1.N3)] Let $\psi$ be a scalar, that is, an $R\mathbb{Z}^3$-periodic function of $c$ (Section~G1.2).
For the multiplication by $\psi$ on the spinors on $\mathbb{T}$, which at $k=0$ are the spinors on the unit lattice, one has $N_\gamma(\psi)=|\psi|_\infty$.
For the multiplication by $\psi$ on the spinors on the unit lattice at $k\ge1$, where $\psi$ acts at a unit site $z$ by its value $\psi(z)$, one has $N_\gamma(\psi)=\sup_{z\in\mathrm{Lab}^0}|\psi(z)|\le|\psi|_\infty$.
Moreover $N_\gamma(\psi O)\le|\psi|_\infty\cdot N_\gamma(O)$ for every $O$, including the operators with values in the functions of short pairs.
\item[(G1.Nc)] A series $\sum_jO_j$ converges absolutely in $N_\gamma$ if $\sum_jN_\gamma(O_j)<\infty$.
\end{itemize}

\begin{proof}[Proof of (G1.N1)--(G1.N3)]
(G1.N1): On the source space $\sum_{z'}\mathbf{1}_{\Delta_{z'}}=1$, and $\exp\left(\gamma\cdot d(z,z')\right)\ge1$.
Hence $\left\|\mathbf{1}_{\Delta_z}O\right\|_{\mathrm{su}}\le\sum_{z'}\left\|\mathbf{1}_{\Delta_z}O\mathbf{1}_{\Delta_{z'}}\right\|_{\mathrm{su}}\le N_\gamma(O)$ for every $z$,
and $\|O\|_{\mathrm{su}}$ is the supremum over $z$ of the left side.
The second inequality holds because $\exp\left(\gamma\cdot d(z,z')\right)\cdot\left\|\mathbf{1}_{\Delta_z}O\mathbf{1}_{\Delta_{z'}}\right\|_{\mathrm{su}}$ is one of the non-negative terms of the sum that defines $N_\gamma(O)$.

(G1.N2): The first inequality follows from the triangle inequality for the operator norm.
For the second, insert $\sum_{z''}\mathbf{1}_{\Delta_{z''}}=1$ between the two factors:
\[
\mathbf{1}_{\Delta_z}O_1O_2\mathbf{1}_{\Delta_{z'}}=\sum_{z''}\left(\mathbf{1}_{\Delta_z}O_1\mathbf{1}_{\Delta_{z''}}\right)\left(\mathbf{1}_{\Delta_{z''}}O_2\mathbf{1}_{\Delta_{z'}}\right).
\]
Together with $d(z,z')\le d(z,z'')+d(z'',z')$ and $\gamma\ge0$ this bounds the sum over $z'$, for fixed $z$, by
$\sum_{z''}\exp\left(\gamma\cdot d(z,z'')\right)\cdot\left\|\mathbf{1}_{\Delta_z}O_1\mathbf{1}_{\Delta_{z''}}\right\|_{\mathrm{su}}\cdot N_\gamma(O_2)\le N_\gamma(O_1)\cdot N_\gamma(O_2)$.

(G1.N3): The scalar $\psi$ is $R\mathbb{Z}^3$-periodic on $\mathbb{T}$, so that it commutes with the indicators;
on functions of pairs it multiplies by its value at the first point, which is the point read by the indicators.
Hence $\mathbf{1}_{\Delta_z}\psi\mathbf{1}_{\Delta_{z'}}=0$ for $z\ne z'$.
On the spinors on $\mathbb{T}$ the operator $\mathbf{1}_{\Delta_z}\psi\mathbf{1}_{\Delta_z}$ has the norm $\sup_{\Delta_z}|\psi|$, and the blocks $\Delta_z$ cover $\mathrm{Lab}$;
this gives $N_\gamma(\psi)=|\psi|_\infty$.
On the spinors on the unit lattice at $k\ge1$ the operator $\mathbf{1}_{\Delta_z}\psi\mathbf{1}_{\Delta_z}$ is the multiplication by the number $\psi(z)$ on the class of $z$, of norm $|\psi(z)|$;
this gives $N_\gamma(\psi)=\sup_{z\in\mathrm{Lab}^0}|\psi(z)|$, which is at most $|\psi|_\infty$ because $\mathrm{Lab}^0\subset\mathrm{Lab}$.
Further $\mathbf{1}_{\Delta_z}\psi O\mathbf{1}_{\Delta_{z'}}=\psi\,\mathbf{1}_{\Delta_z}O\mathbf{1}_{\Delta_{z'}}$, whose norm is at most $|\psi|_\infty\cdot\left\|\mathbf{1}_{\Delta_z}O\mathbf{1}_{\Delta_{z'}}\right\|_{\mathrm{su}}$.
\end{proof}

The same arguments give these inequalities row by row.
Let
\[
\mathrm{Rw}_z(O):=\sum_{z'}\exp\left(\gamma\cdot d(z,z')\right)\cdot\left\|\mathbf{1}_{\Delta_z}O\mathbf{1}_{\Delta_{z'}}\right\|_{\mathrm{su}},
\]
so that $N_\gamma(O)=\sup_z\mathrm{Rw}_z(O)$.
Then
\begin{gather*}
\mathrm{Rw}_z(O_1+O_2)\le\mathrm{Rw}_z(O_1)+\mathrm{Rw}_z(O_2),\\
\mathrm{Rw}_z(O_1O_2)\le\mathrm{Rw}_z(O_1)\cdot N_\gamma(O_2)\quad\text{if }O_2\text{ takes values in spinors},\\
\mathrm{Rw}_z\left(\psi O\psi'\right)\le\sup_{\Delta_z}|\psi|\cdot|\psi'|_\infty\cdot\mathrm{Rw}_z(O).
\end{gather*}

\begin{named}{Remark}
The norm $N_\gamma$ is that of \cite[Section~16.2.6]{QEDI}, where (G1.N1), (G1.N2) and the equality $N_\gamma(\psi)=|\psi|_\infty$ of (G1.N3) on the spinors on $\mathbb{T}$ appear as (EW1)--(EW3).
The multipliers are also those of \cite[Section~16.2.6]{QEDI}:
the action of a scalar on the functions of short pairs through its value at the first point, and the inequality $N_\gamma(\psi O)\le|\psi|_\infty\cdot N_\gamma(O)$ of (G1.N3), also for the operators $O$ with values in the functions of short pairs, stand there in the paragraph that follows the proof of (EW3);
for operators with values in the functions of short pairs, this inequality is the bound applied in the proof of \cite[Lemma~16.13(d)]{QEDI}.
Here, in addition, $O$ may be an operator between any two of the three spaces, so that the functions of short pairs are admitted also as a source space;
in \cite[Section~16.2.6]{QEDI} the source of $O$ is a space of spinors.
\end{named}

\emph{Kernel norm at level $0$.}
Let $k=0$, and let $O$ be an operator on the spinors on $\mathbb{T}$, with its kernel as in Section~G1.2.
Put
\[
N^{\mathrm{ke}}_\gamma(O):=\sup_z\sum_{z'}\exp\left(\gamma\cdot d(z,z')\right)\cdot\sup\left\{|O(x,x')|_2:\ x\text{ over }z,\ x'\text{ over }z'\right\}.
\]
\begin{itemize}
\item[(G1.Nk1)] $N_\gamma(O)\le N^{\mathrm{ke}}_\gamma(O)$, since a cluster has the measure one.
\item[(G1.Nk2)] $N^{\mathrm{ke}}_\gamma(O_1O_2)\le N_\gamma(O_1)\cdot N^{\mathrm{ke}}_\gamma(O_2)$.
\item[(G1.Nk3)] $N^{\mathrm{ke}}_\gamma\left(\psi O\psi'\right)\le|\psi|_\infty\cdot|\psi'|_\infty\cdot N^{\mathrm{ke}}_\gamma(O)$.
\item[(G1.Nk4)] $|O(x,x')|_2\le\exp\left(-\gamma\cdot d(x,x')\right)\cdot N^{\mathrm{ke}}_\gamma(O)$.
\item[(G1.Nk5)] $|O(x,x')|_2\le48\cdot\|O\|_{L^2}$.
Indeed, a site of a period cell carries the weight $\frac{1}{48}$.
The spinor $u$ that carries the column at $x'$, that is, the spinor with $(Ou)(x)=O(x,x')\,\mathrm{vv}$ for a given $\mathrm{vv}\in\mathbb{C}^4$, has the norm $\|u\|_{L^2}=\sqrt{48}\cdot|\mathrm{vv}|_2$,
and every spinor $w$ satisfies $|w(x)|_2\le\sqrt{48}\cdot\|w\|_{L^2}$;
so $|O(x,x')\,\mathrm{vv}|_2\le\sqrt{48}\cdot\|Ou\|_{L^2}\le48\cdot\|O\|_{L^2}\cdot|\mathrm{vv}|_2$.
\end{itemize}
Note that $N^{\mathrm{ke}}_\gamma(1)=48$.

\begin{proof}[Proof of (G1.Nk1)--(G1.Nk4)]
(G1.Nk1): At $k=0$ the block $\Delta_{z'}$ is the single site $z'$.
For $x$ over $z$,
\begin{align*}
\left|\left(O\mathbf{1}_{\Delta_{z'}}f\right)(x)\right|_2
&=\left|\frac{1}{48}\sum_{g'}O\left(x,(z',g')\right)f(z',g')\right|_2\\
&\le\sup\left\{|O(x,x')|_2:\ x'\text{ over }z'\right\}\cdot|f|_{\mathrm{su}}.
\end{align*}
(G1.Nk2): For $\mathrm{vv}\in\mathbb{C}^4$ one has $(O_1O_2)(x,x')\,\mathrm{vv}=\left(O_1\left[O_2(\cdot,x')\,\mathrm{vv}\right]\right)(x)$;
the argument then proceeds as for (G1.N2).
(G1.Nk3): $\left(\psi O\psi'\right)(x,x')=\psi(c)\cdot O(x,x')\cdot\psi'(c')$ by Section~G1.2.
(G1.Nk4): The right side dominates a single term of the sum that defines $N^{\mathrm{ke}}_\gamma(O)$.

The same arguments apply row by row.
With
\[
\mathrm{Rk}_z(O):=\sum_{z'}\exp\left(\gamma\cdot d(z,z')\right)\cdot\sup\left\{|O(x,x')|_2:\ x\text{ over }z,\ x'\text{ over }z'\right\}
\]
one obtains
\begin{gather*}
\mathrm{Rw}_z(O)\le\mathrm{Rk}_z(O),\\
\mathrm{Rk}_z(O_1O_2)\le\mathrm{Rw}_z(O_1)\cdot N^{\mathrm{ke}}_\gamma(O_2),\\
\mathrm{Rk}_z\left(\psi O\psi'\right)\le|\psi(z)|\cdot|\psi'|_\infty\cdot\mathrm{Rk}_z(O).
\end{gather*}
\end{proof}

For $\gamma>0$ let
\[
c(\gamma):=\sum_{n\in\mathbb{Z}^3}\exp\left(-\gamma\cdot|n|\right).
\]
On every block torus, $\sum_{z'}\exp\left(-\gamma\cdot d(z,z')\right)\le c(\gamma)$.

\subsubsection*{G1.4. Spinors with a character at an arbitrary mesh}

Let $\mathrm{me}>0$ be a mesh, let $R\in\mathrm{me}\,\mathbb{N}:=\{\mathrm{me}\,n:\ n=1,2,\dots\}$ or $R=\infty$, and let $\theta\in\mathbb{R}^3$.
Put
\[
\mathfrak{S}^{\mathrm{ch}}(\mathrm{me},R;\theta):=\left\{f:\mathrm{me}\,\mathbb{Z}^3\times O_{\mathrm{cub}}\to\mathbb{C}^4\ :\ f(c+Re_\mu,g)=-\exp(i\theta_\mu)\cdot f(c,g)\right\},
\]
where no condition is imposed if $R=\infty$, and let $\mathrm{HH}(\mathrm{me},R;\theta)$ be the set of the elements $f$ of this space with $\|f\|_{L^2}<\infty$;
if $R<\infty$, it is the whole space.
For $\theta=0$ the spinors are antiperiodic, for $\theta=(\pi,\pi,\pi)$ periodic.
This is the convention of (I.6.1); in the notation of \cite[Section~12.1]{QEDI} the character of the same boundary condition is $(\theta_\mu+\pi)_\mu$.

The neighbours are $x^{(i)}=(c,gs_i)$ for $i\le2$ and $x^{(3)}=(c+\mathrm{me}\cdot ge_1,gs_3)$.
The measure is $\int\mathrm{d}x:=\mathrm{me}^3\cdot\frac{1}{48}\sum_{c,g}$, the sum extending over a period cell;
$(f',f)^{\mathrm{ch}}$ is the corresponding sesquilinear pairing on $\mathrm{HH}(\mathrm{me},R;\theta)$,
and the bilinear pairing $\langle f',f\rangle^{\mathrm{ch}}$ pairs $\mathfrak{S}^{\mathrm{ch}}(\mathrm{me},R;-\theta)$ with $\mathfrak{S}^{\mathrm{ch}}(\mathrm{me},R;\theta)$.
Kernels are taken with respect to $\int\mathrm{d}x$, and $O(t,x'+Re_\mu)=-\exp(-i\theta_\mu)\cdot O(t,x')$.
The operators
\begin{gather*}
(Df)(x):=-\mathrm{me}^{-1}\sum_{i=0}^{3}M_i(x)\left[f(x^{(i)})-f(x)\right],\\
\nabla^{*}\nabla:=\sum_iw_i\,(1-R_i),\qquad(R_if)(x):=f(x^{(i)})
\end{gather*}
are the operators (I.2.3) and (I.2.4) at $A=0$ and at the mesh $\mathrm{me}$.
For $x\ne x'$ let $d_X(x,x'):=\max\left(|c-c'|,\mathrm{me}\right)$; this is a distance of the covering lattice, also if $R<\infty$.
In exponents, $|c-c'|$ denotes the sup distance modulo $R\mathbb{Z}^3$.

\emph{Blocks.}
For $j\ge0$, a block of level $j$ is a set
\[
E=\left\{c\in\mathrm{me}\,\mathbb{Z}^3:\ |c-z|<\frac{L^j\mathrm{me}}{2}\right\},\qquad z\in L^j\mathrm{me}\,\mathbb{Z}^3;
\]
its side is $\ell_E:=L^j\mathrm{me}$.
For a spinor $u$, $P_Eu$ is the mean of $u$ over $E\times O_{\mathrm{cub}}$, regarded as a constant function on $E\times O_{\mathrm{cub}}$; its kernel is $\ell_E^{-3}$ on $(E\times O_{\mathrm{cub}})^2$.
The norms $\|u\|_{L^2,E}$ and $|u|_{\mathrm{su},E}$ are formed over $E\times O_{\mathrm{cub}}$,
and $(u,\nabla^{*}\nabla u)_E$ is $\mathrm{me}^3\cdot\frac{1}{48}$ times the sum of $w_i\cdot\left|u(x^{(i)})-u(x)\right|_2^2$ over the unoriented bonds $\{x,x^{(i)}\}$ whose two ends lie over $E$.
The same displays are read for every cube $E$ of $n^3$ sites of $\mathrm{me}\,\mathbb{Z}^3$, $n\ge1$, with $\ell_E:=n\,\mathrm{me}$; these are the data of Proposition~1.4.

\subsubsection*{G1.5. Mirror torus and local inverse of an $M$-cube}

The photon objects of this section are those of \cite[Section~16.6.1]{QEDI}; on fermion sites the label is a spectator, as in (I.J3).

Let the data (SF) be given, let $\square$ be an $M$-cube of $T_k$, and let $z_{\square}\in M\mathbb{Z}^3$ be a lift of its centre.
Put
\begin{gather*}
\tilde{W}:=\left\{c\in\mathbb{L}_k:\ |c-z_{\square}|<\frac{3M}{2}\right\},\\
W^{-}:=\left\{c\in\tilde{W}:\ |c-z_{\square}|<\frac{3M}{2}-\eta_k\right\},\\
o:=z_{\square}-\frac{3M}{2}(1,1,1);
\end{gather*}
the set $\tilde{W}$ is a lift of $\square^{\sim1}$.
The bonds of $\tilde{W}$ are the bonds with both ends in $\tilde{W}$.
The mirror torus is the block torus $\mathbb{T}^{W}:=\left(\frac{\eta_k}{2}(1,1,1),6M,+1\right)$, which carries periodic spinors; its lattice is denoted by $\mathrm{Lab}^W$.
Let $\mathrm{fo}(t):=\operatorname{dist}(t,6M\mathbb{Z})$ for $t\in\mathbb{R}$, and
\[
\mathrm{Fo}(u):=o+\left(\mathrm{fo}(u_1),\mathrm{fo}(u_2),\mathrm{fo}(u_3)\right),\qquad\mathrm{emb}(c):=c-o\ \text{ modulo }6M\mathbb{Z}^3,
\]
\[
\mathrm{Ch}:=\left\{u\in\mathrm{Lab}^W:\ 0<u_\mu<3M\ \text{modulo }6M\ \text{for }\mu=1,2,3\right\}.
\]
For a complex antisymmetric function $A$ on the bonds of $\tilde{W}$ and a function $\boldsymbol{\lambda}$ on $\tilde{W}$ let
\[
A^W(u,u'):=\begin{cases}A(\mathrm{Fo}\,u,\mathrm{Fo}\,u')&\text{if }\mathrm{Fo}\,u\ne\mathrm{Fo}\,u',\\0&\text{otherwise},\end{cases}\qquad\boldsymbol{\lambda}^W:=\boldsymbol{\lambda}\circ\mathrm{Fo}.
\]
For a function $f$ on $\tilde{W}\times O_{\mathrm{cub}}$ and a spinor $v$ on $\mathbb{T}^{W}$, that is, a $6M\mathbb{Z}^3$-periodic function on $\mathrm{Lab}^W\times O_{\mathrm{cub}}$, let
\[
(\mathrm{ext}\,f)(u,g):=\begin{cases}f(\mathrm{Fo}\,u,g)&\text{if }u\in\mathrm{Ch},\\0&\text{otherwise},\end{cases}\qquad(\mathrm{rs}\,v)(c,g):=v(\mathrm{emb}\,c,g).
\]
Define
\[
F^W_y(A):=F_y(A^W)\quad\text{on the spinors on }\mathbb{T}^{W},\qquad\mathrm{Sloc}_{\square,y}(A):=\mathrm{rs}\cdot F^W_y(A)^{-1}\cdot\mathrm{ext};
\]
the local inverse $\mathrm{Sloc}_{\square,y}(A)$ is an operator on the functions on $\tilde{W}\times O_{\mathrm{cub}}$.

For $\delta>0$ let $\mathrm{Gloc}(\delta)$ be the set of the functions $A=B+d\boldsymbol{\lambda}$ on the bonds of $\tilde{W}$ with $\boldsymbol{\lambda}$ real and $e_k\cdot|B|_\infty<\delta$;
it is an open subset of the space of the complex antisymmetric functions on the bonds of $\tilde{W}$.
The local inverse $\mathrm{Sloc}_{\square,y}(A)$ is defined wherever $F^W_y(A)$ is bijective.
By Proposition~2.5\,(a) this is the case for all $y\in Y_k$ and all $A\in\mathrm{Gloc}(\delta)$ provided that
\[
3L\cdot\delta\le1\qquad\text{and}\qquad2C'_{V}\cdot a'\cdot\delta\le1,
\]
which are the conditions on $\delta$ in Proposition~2.4.
The number $\delta=2e_k^{\varepsilon}$ of Section~G1.6 satisfies both conditions;
this follows from $e_k\le e_0$ and from the conditions (Z2.s1) and (Z2.s2) of Proposition~14.2\,(b).

The assertions of \cite[Lemma~16.10]{QEDI}, and those of the paragraph that follows \cite[Definition~16.11]{QEDI}
(the map $A\mapsto A^W$ is complex linear, $(d\boldsymbol{\lambda})^W=d\left(\boldsymbol{\lambda}^W\right)$, $\left(\overline{A}\right)^W=\overline{A^W}$, and $|A^W|_\infty=|A|_\infty$),
concern photon sites and bonds only; they hold here as stated there, the directions being numbered $\mu=1,2,3$.
Moreover $\mathrm{rs}\cdot\mathrm{ext}=1$, and $\mathrm{rs}=\mathrm{ext}^{T}$ for the pairings with the weight $\frac{1}{48}$.

The label $g$ is not reflected in the seven other chambers: $F_y(A^W)$ is the operator of Section~G1.2 in the folded connection $A^W$.

\subsubsection*{G1.6. Operators over $\tilde{W}$ as operators on $T_k$}

Let $\square$ be an $M$-cube of $T_k$, with the sets of Section~G1.5; compare \cite[Section~16.6.2]{QEDI}.
Let $O$ be a linear operator on the functions on $\tilde{W}\times O_{\mathrm{cub}}$, with values either in the functions on $\tilde{W}\times O_{\mathrm{cub}}$ or in the functions on the short pairs whose two points lie over $\tilde{W}$.
Its kernel is taken with respect to $\int_{\tilde{W}}\mathrm{d}x$, that is,
\[
(Of)(x)=\frac{\eta_k^3}{48}\sum_{c'\in\tilde{W}}\ \sum_{g'\in O_{\mathrm{cub}}}O(x,x')\cdot f(c',g'),\qquad x'=(c',g').
\]
The norm $N_{\gamma}(O)$ and, at $k=0$, the kernel norm $N^{\mathrm{ke}}_{\gamma}(O)$ are formed with the unit sites of $\tilde{W}$ and with the distance $|z_1-z_2|$.

For $f\in\mathfrak{S}^{\mathrm{ch}}(\eta_k,\ell)$ the spinor $\mathrm{Ecb}(O)f$ is defined by
\[
\left(\mathrm{Ecb}(O)f\right)(x):=\left(O\left(f|_{\tilde{W}\times O_{\mathrm{cub}}}\right)\right)(x)\qquad(x\ \text{over}\ \tilde{W}),
\]
by
\[
\left(\mathrm{Ecb}(O)f\right)(x+\ell n):=(-1)^{n_1+n_2+n_3}\cdot\left(\mathrm{Ecb}(O)f\right)(x)\qquad\left(n\in\mathbb{Z}^3\right),
\]
and by zero off the orbit of $\tilde{W}$.
For functions of pairs both points are shifted, and $\mathrm{Ecb}(O)f$ vanishes on the short pairs that do not lie over one translate of $\tilde{W}$.
Where the cube has to be indicated, $\mathrm{Ecb}$ is written $E_{\square}$; thus $E_{\tilde{g}\square}$ and $E_{\square_i}$ denote the same map for the cubes $\tilde{g}\square$ and $\square_i$, formed with their own sets $\tilde{W}$.

Let $A\in\tilde{G}_k$.
By the relation (G1) of \cite[Section~16.2.7]{QEDI}, the restriction $A|_{\tilde{W}}$ of $A$ to the bonds of $\tilde{W}$ lies in $\mathrm{Gloc}\left(2e_k^{\varepsilon}\right)$; the symbol $\mathrm{Sloc}_{\square,y}(A)$ also denotes the operator $\mathrm{Ecb}\left(\mathrm{Sloc}_{\square,y}\left(A|_{\tilde{W}}\right)\right)$ on $T_k$.
On $T_k$ the symbol $\delta_{\alpha,A}\mathrm{Sloc}_{\square,y}(A)$ always denotes $\mathrm{Ecb}$ applied to the operator $\delta_{\alpha,A}\mathrm{Sloc}_{\square,y}\left(A|_{\tilde{W}}\right)$ over $\tilde{W}$, and never the composition of $\delta_{\alpha,A}$ with $\mathrm{Sloc}_{\square,y}(A)$ on $T_k$.
Every other symbol $\delta_{\alpha,A}O$ with an operator $O$ on $\mathfrak{S}^{\mathrm{ch}}(\eta_k,\ell)$, for instance $\delta_{\alpha,A}h_{\square}\cdot\mathrm{Sloc}_{\square}$, $\delta_{\alpha,A}S^{\#}_{y,\mathrm{wk}}(A)$, $\delta_{\alpha,A}S^{\#}_y(s,A)$ and $\delta_{\alpha,A}S_y(A)$, denotes the composition of $O$ with the map (I.1.7) on $T_k$, restricted to the short pairs.

\subsubsection*{G1.7. Cut-off functions}

Compare \cite[Section~16.7.1]{QEDI}.
Let
\[
\mathrm{cu}(t):=
\begin{cases}
1, & |t|\le\frac{1}{4},\\
\cos\left(\frac{\pi}{2}\left(2|t|-\frac{1}{2}\right)\right), & \frac{1}{4}\le|t|\le\frac{3}{4},\\
0, & |t|\ge\frac{3}{4}.
\end{cases}
\]
For an $M$-cube $\square$ of $T_k$, with $\tilde{W}$ and $z_{\square}$ as in Section~G1.5, put
\begin{gather*}
h_{\square}(c):=\prod_{\mu}\mathrm{cu}\left(\frac{c_{\mu}-z_{\square,\mu}}{M}\right)\quad(c\in\tilde{W}),\\
h_{\square}(c+\ell n):=h_{\square}(c)\quad\left(n\in\mathbb{Z}^3\right),
\end{gather*}
so that $h_{\square}$ is $\ell\mathbb{Z}^3$-periodic, and let $h_{\square}:=0$ off the orbit of $\tilde{W}$.
On fermion sites, $h_{\square}(x):=h_{\square}(c)$.
Then
\begin{equation}\tag{6.1}
\begin{aligned}
&\sum_{\square}h_{\square}^2=1,\qquad\operatorname{supp}h_{\square}\subset\left\{c\in\mathbb{L}_k:\ |c-z_{\square}|<\frac{3M}{4}\right\}+\ell\mathbb{Z}^3,\\
&\left|h_{\square}(c)-h_{\square}(c')\right|\le\frac{\pi}{M}\,|c-c'|_1,\qquad h_{\tilde{g}\square}(\tilde{g}c)=h_{\square}(c)\quad\left(\tilde{g}\in\mathbb{G}_k\right).
\end{aligned}
\end{equation}
A unit block meets $\operatorname{supp}h_{\square}$ for at most $8$ cubes $\square$.
Finally,
\[
\mathrm{Cm}_{\square,y}(A):=-\left[F_y(A),h_{\square}\right].
\]

\subsubsection*{G1.8. Walks}

Compare \cite[Definition~16.14]{QEDI}.
Let $n\ge0$ be an integer.
A walk of length $|\mathrm{wk}|=n$ is a sequence $\mathrm{wk}=(\square_0,\dots,\square_n)$ of $M$-cubes of $T_k$ with $d\left(z_{\square_{i-1}},z_{\square_i}\right)\le M$ for $1\le i\le n$.
The walks of length $n$ form the set $\mathrm{Wk}_n$, and $\mathrm{Wk}$ is the union of these sets.
With $h_i:=h_{\square_i}$, $\mathrm{Sloc}_i:=\mathrm{Sloc}_{\square_i,y}(A)$ and $\mathrm{Cm}_i:=\mathrm{Cm}_{\square_i,y}(A)$ one defines
\begin{align*}
S_{y,\mathrm{wk}}(A)&:=\left(h_0\cdot\mathrm{Sloc}_0\cdot h_0\right)\cdot\left(\mathrm{Cm}_1\cdot\mathrm{Sloc}_1\cdot h_1\right)\cdot\left(\mathrm{Cm}_2\cdot\mathrm{Sloc}_2\cdot h_2\right)\cdots\\
&\qquad\cdots\left(\mathrm{Cm}_n\cdot\mathrm{Sloc}_n\cdot h_n\right),\\
S'_{y,\mathrm{wk}}(A)&:=\left(h_n\cdot\mathrm{Sloc}_n\right)\cdot\left(-\mathrm{Cm}_n\cdot h_{n-1}\cdot\mathrm{Sloc}_{n-1}\right)\cdot\left(-\mathrm{Cm}_{n-1}\cdot h_{n-2}\cdot\mathrm{Sloc}_{n-2}\right)\cdots\\
&\qquad\cdots\left(-\mathrm{Cm}_1\cdot h_0\cdot\mathrm{Sloc}_0\right)\cdot h_0,\\
X_{\mathrm{wk}}&:=\bigcup_{0\le i\le n}\square_i^{\sim 1}.
\end{align*}
The set $X_{\mathrm{wk}}$ is a polymer, and $|X_{\mathrm{wk}}|_M\le27(n+1)$, $d_M(X_{\mathrm{wk}})\le27(n+1)$ (Proposition~3\,(d)).

\subsubsection*{G1.9. Weakened operators and polymer propagators on the fine lattice}

Compare \cite[Definition~16.17]{QEDI}.
Let $s=(s_{\mathrm{pc}})_{\mathrm{pc}\in\pi_k}\in\mathrm{Pd}_k$, $X\in D_k$, $y\in Y_k$ and $A\in\tilde{G}_k$.
Put
\[
s_{\mathrm{wk}}:=\prod\left\{s_{\mathrm{pc}}:\ \mathrm{pc}\in\pi_k,\ \mathrm{pc}\subset X_{\mathrm{wk}}\right\},
\]
\[
S_y(s,A):=\sum_{\mathrm{wk}\in\mathrm{Wk}}s_{\mathrm{wk}}\cdot S_{y,\mathrm{wk}}(A),\qquad S_y(X,A):=\sum_{\mathrm{wk}:\ X_{\mathrm{wk}}=X}S_{y,\mathrm{wk}}(A),
\]
and let $S'_y(s,A)$, $S'_y(X,A)$ be formed in the same way with $S'_{y,\mathrm{wk}}(A)$ in place of $S_{y,\mathrm{wk}}(A)$.
For both families,
\[
S^{\mathrm{sy}}_y:=\frac{1}{2}\left(S_y+S'_y\right),\qquad S^{\mathrm{sy}}_{y,\mathrm{wk}}:=\frac{1}{2}\left(S_{y,\mathrm{wk}}+S'_{y,\mathrm{wk}}\right).
\]
The sums over walks are series.
They converge absolutely in $N_{\gamma^{F}}$, in the sense of (G1.Nc) of Section~G1.3, uniformly on $\mathrm{Pd}_k\times\tilde{G}_k$; this is Proposition~3\,(a).
The symbol $S^{\#}$ stands for $S$, $S'$ or $S^{\mathrm{sy}}$.
Every walk, those of length zero included, carries the parameters of all cubes of $X_{\mathrm{wk}}$.
A value $y$ is called neutral if $F_y$ contains no $\gamma_5$: this means $y\in\{0,\infty\}$ for $k\ge1$ and $y=0$ for $k=0$.
One sets $-\infty:=\infty$.

\subsubsection*{G1.10. Notation}

In Sections~G1.1--G1.11 and in the Propositions $n$ and $n.m$ with $1\le n\le14$ the following notation is used.
An arbitrary mesh is denoted $\mathrm{me}>0$.
The set of the positive integers is $\mathbb{N}=\{1,2,3,\dots\}$; accordingly $\mathrm{me}\,\mathbb{N}$ (Section~G1.4), $L\,\mathbb{N}$ and $\ell_1\mathbb{N}$ (Propositions~1.6 and~1.7) are the sets of the positive integer multiples of $\mathrm{me}$, $L$ and $\ell_1$.
For a momentum $p$ and a mass $\mathbf{m}$,
\[
\mathrm{tw}(p):=\sum_{\mu}\left(1-\cos p_{\mu}\right),\qquad\mathrm{rp}(p):=\left(|p|_2^2+\mathbf{m}^2\right)^{\frac{1}{2}}.
\]
A multi-index is denoted $\mathrm{ak}$, and $\mathrm{dc}_{A,a}$ is the covariant difference quotient of a $4$-component spinor (Proposition~1.1).
The indices $\mathrm{su}$ and $L^2$ of norms and the kernel norm $N^{\mathrm{ke}}_{\gamma}$ are those of Section~G1.3.
A morphism of block tori is denoted $\mathrm{mph}$, its translation part $\mathrm{tv}$, and its action $\mathrm{mS}$ (Section~G1.2).
The point $z_{\square}$ is the lift of the centre of $\square$ of Section~G1.5.
The letters $\mathrm{Sloc}$, $\mathrm{Gloc}$, $\mathrm{Ecb}$ denote the local inverse, its domain, and the map that turns an operator over $\tilde{W}$ into an operator on $T_k$ (Sections~G1.5 and~G1.6).
The cut-off function $\mathrm{cu}$ of Section~G1.7 is a function of a real variable; in Section~I.4.6 and in Proposition~18.6, $\mathrm{cu}$ is a vector of $\mathbb{R}^3$.
The letters $\mathrm{wk}$, $\mathrm{Wk}_n$, $\mathrm{Wk}$, $X_{\mathrm{wk}}$, $s_{\mathrm{wk}}$ refer to walks (Sections~G1.8 and~G1.9).
A weight is a number $\mathrm{wp}\ge1$, and
\[
\mathrm{qq}(\mathrm{wp}):=\frac{54\pi\cdot C'_{V}\cdot a'\cdot\mathrm{wp}}{M}.
\]
A real function of $c$ is denoted $\psi$, and $\mathrm{Lip}$ is its Lipschitz constant.
A periodic complex, or matrix valued, function on the short pairs is denoted $\mathrm{wq}$, and $E_{\mathrm{wq}}$ is the operator formed with it (Proposition~2.1\,(c)).
A closed real connection is denoted $\mathrm{Acl}$, its phase function $\mathrm{uf}$, and $\mathrm{uf}^0$, $\mathrm{uf}^1$ are restrictions of $\mathrm{uf}$.
For $\tilde{g}=(S,a)\in\mathbb{G}_k$ the action $\mathrm{glb}$ of $\tilde{g}$ on barred fields is
\[
(\mathrm{glb}\,f)(x):=\left(S^{-1}\right)^T\cdot f\left(\tilde{g}^{-1}x\right),
\]
and the same matrix acts on functions of pairs and on $\mathfrak{S}(1,\ell)$:
\[
(\mathrm{glb}\,F)(x,x'):=\left(S^{-1}\right)^T\cdot F\left(\tilde{g}^{-1}x,\tilde{g}^{-1}x'\right),\qquad(\mathrm{glb}\,v)(y):=\left(S^{-1}\right)^T\cdot v\left(\tilde{g}^{-1}y\right).
\]
In Proposition~4 and in the proofs of Propositions~4 and~5 (Chapters~4 and~5), $\mathrm{Qk}$ and $\mathrm{Qt}$ abbreviate $Q_k(A)\cdot J^{*}$ and $J\cdot Q_k^{T}(-A)$.
The lattices of a block torus are denoted $\mathrm{Lab}$, $\mathrm{Lab}_j$, $\mathrm{Lab}^0$, $\mathrm{Lab}^W$ (Sections~G1.2 and~G1.5); in Part~I, $\Lambda$ is a gauge function and $\Lambda_{\eta,R}=\eta\mathbb{Z}^3/R\mathbb{Z}^3$ is the photon torus of Section~I.7.5, and in \cite[Section~2.1]{QEDI} the letter $\Lambda_n$ denotes the lattice torus $\prod_{\mu=1}^{3}\mathbb{Z}/n_{\mu}\mathbb{Z}$.
The operator of Proposition~1.7 is denoted $\mathrm{Tst}$.
The rate per unit of $d_M$ in a set of parameters (Section~G1.11) is denoted $\mathrm{kq}$; \cite[Section~5.4.1]{QEDI} writes $\kappa'$ for it, whereas here $\kappa'$ is the number $\kappa_{\mathrm{w}}+\kappa_{0}$ of Section~G1.1.

The penalty constants $c_1$, $C_1$, $C_2$ of \cite[Section~12.2]{QEDI} are written out: they are $b_{-}$, $L\cdot b$ and $2L^4\cdot b$.
The translation part of an element of $\mathbb{G}_k$, which \cite{QEDI} denotes $a'$, is denoted $a$; the letter $a'$ is the constant of Proposition~2.3.

Several letters of Part~I also have the following meanings, each in the places named with it.
The letter $t$ denotes the shift of a block torus $\mathbb{T}=(t,R,\boldsymbol{\sigma})$ (Section~G1.2) and of the lattice $\mathrm{Lab}$ in Proposition~1.1\,(f), and a real variable in $\mathrm{fo}(t)$ (Section~G1.5), in $\mathrm{cu}(t)$ (Section~G1.7) and in Proposition~1.1\,(c); as the argument of a kernel or of a spinor, as in $O(t,t')$, $O(t,x')$, $(Of)(t)$, it is an index, as in (I.1.4).
The letter $R$ denotes the side of a block torus (Section~G1.2; the unit lattice $\mathrm{Lab}^0$ has the spacing $1$) and the period, finite or infinite, of the spinors of Section~G1.4 and of Propositions~1.1--1.7.
The letter $C$ denotes a number in a set of parameters $(\alpha_0,\gamma,\mathrm{kq};C)$ and in W5--W7 (Section~G1.11), and in the bound $\|O\|_{\mathrm{su}}\le C$ (Section~G1.3).
In $\mathrm{Dsc}_{\mathbf{m}}(p)$ (Proposition~1.2), and in $\mathrm{Dsc}_0(p)$ (Chapter~1a, Sections~2.9 and~3, the proof of Proposition~1.3\,(b)), the index is the mass and $p$ is the momentum; the symbol $\mathrm{Dsc}_{A}$ of (I.J6) carries a connection as its index and has no argument.
In Proposition~1.3\,(b) and in Proposition~14.2\,(a), row $c_b$, the letters $A_1$, $A_2$ denote the two numbers, functions of $\rho$, with which $c_b$ is formed.
In Chapter~1a, Sections~2.6 and~2.7, $A_0$ and $A'_0$ denote the members with $\mathrm{ak}=0$ of the families of constants $A_{\mathrm{ak}}$, $A'_{\mathrm{ak}}$ (Proposition~1.2\,(S6), (S7); Section~I.8.2, line (C0r)), whose index is a multi-index.

\subsubsection*{G1.11. The properties W1, W2, W5--W17 of the fermionic primitives}

Let the data (SF) of Section~G1.1 be given, and let $s=(s_{\mathrm{pc}})_{\mathrm{pc}\in\pi_k}\in\mathrm{Pd}_k$ (in Section~I.5.8 this parameter is written $s''$), $A\in\tilde{G}_k$, $X\in D_k$ and $y\in[0,\infty)$.
An index $t$ is either a fermion site $x=(c,g)$, with $\operatorname{supp}x=\{c\}$, or a unit site $z$.
Kernels are taken by (I.1.4) when the source is a space of spinors on fermion sites, and by $(Ov)(t)=\sum_{z'}O(t,z')\cdot v(z')$ when the source is $\mathfrak{S}(1,\ell)$.
Conditions on indices refer to the sites of the torus, that is, to classes modulo $\ell\mathbb{Z}^3$ (Section~G1.2).

\emph{The fermionic part of a system of primitives of level $k\ge1$} consists of families of linear operators:
\begin{itemize}
\item[(p3)] $H_k(s,A):\mathfrak{S}(1,\ell)\to\mathfrak{S}^{\mathrm{ch}}(\eta_k,\ell)$, together with $\bar{H}_k(s,A):=\mathsf{C}^{-1}\cdot H_k(s,-A)\cdot\mathsf{C}$;
\item[(p4)] $\Gamma_k(s,A)$ and $M_k(X,A)$ on $\mathfrak{S}(1,\ell)$;
\item[(p5)] $S_{k,y}(X,A)$ on $\mathfrak{S}^{\mathrm{ch}}(\eta_k,\ell)$, Borel in $y$.
\end{itemize}
\emph{At level $0$} it consists of $\Gamma_0(s,A)$ and $S_{0,y}(X,A)$ on $\mathfrak{S}^{\mathrm{ch}}(1,\ell)$, the second Borel in $y$ on $[0,\infty)$.
There (p3) is the diagonal family: $H_0(s,A)$ is the multiplication by
\[
\prod\left\{s_{\mathrm{pc}}:\ \mathrm{pc}\in\pi_0,\ \mathrm{pc}\ \text{meets}\ K_c\right\},
\]
the same factor on the $48$ sites over $c$, and $\bar{H}_0(s,A):=\mathsf{C}^{-1}\cdot H_0(s,-A)\cdot\mathsf{C}=H_0(s,A)$.
There is no $M_0(X,A)$.

A family has a property \emph{for $(\alpha_0,\gamma,\mathrm{kq};C)$}, where $M^{\alpha_0}$ is the radius of $\mathrm{Pd}_k$, $\gamma$ a rate per unit of length, $\mathrm{kq}$ a rate per unit of $d_M$ and $C$ a constant, if the following holds.
All moduli are $|\cdot|_c$, and all sup norms are the norms $|\cdot|_{\infty}$ of Section~I.1.5.
\begin{description}
\item[W1.] The symbol $1$ also denotes the family $\left(s_{\mathrm{pc}}\right)_{\mathrm{pc}\in\pi_k}$ with $s_{\mathrm{pc}}=1$ for every $\mathrm{pc}\in\pi_k$;
it lies in $\mathrm{Pd}_k$, because the radius $M^{\alpha_0}$ of $\mathrm{Pd}_k$ is at least $1$.
One has $H_k(1,A)=H_k(A)$, $\Gamma_k(1,A)=\Gamma_k(A)$, $\sum_{X}M_k(X,A)=M_k(A)$ and $\sum_{X}S_{k,y}(X,A)=S_{k,y}(A)$.
\item[W2.] $H_k(s,A)$ and $\Gamma_k(s,A)$ are continuous on $\mathrm{Pd}_k\times\tilde{G}_k$ and analytic in $(s,A)$ in its interior, and for every fixed $s\in\mathrm{Pd}_k$, the boundary included, they are analytic in $A\in\tilde{G}_k$; $M_k(X,\cdot\,)$ and $S_{k,y}(X,\cdot\,)$ are analytic on $\tilde{G}_k$.
\item[W5.] $|H_k(s,A)\,v|_{\infty}\le C\cdot|v|_{\infty}$, and $\left|\left(\delta_{\alpha,A}H_k(s,A)\,v\right)(x,x')\right|_c\le C\cdot|v|_{\infty}$ for all \emph{short} pairs ($d_X<1$; the pairs inside a cluster included).
\item[W6.] One has
\begin{gather*}
\left|\Gamma_k(s,A)(t,t')\right|_c\le C\cdot\exp\left(-\gamma\cdot d(t,t')\right),\\
\left|M_k(X,A)(z,z')\right|_c\le C\cdot\exp\left(-\gamma\cdot d(z,z')-\mathrm{kq}\cdot d_M(X)\right).
\end{gather*}
For $k\ge1$, $t$ and $t'$ are unit sites; for $k=0$, they are fermion sites and the kernel is that of (I.1.4).
\item[W7.] For $k\ge1$ and all $y\ge0$,
\[
\left|\mathbf{1}_{\Delta_z}\cdot S_{k,y}(X,A)\cdot\mathbf{1}_{\Delta_{z'}}\,f\right|_{\infty}\le C\cdot\exp\left(-\gamma\cdot d(z,z')-\mathrm{kq}\cdot d_M(X)\right)\cdot|f|_{\infty}.
\]
For $k=0$ and all $y\ge0$, with the kernel of (I.1.4),
\[
\left|S_{0,y}(X,A)(x,x')\right|_c\le C\cdot(1+y)^{-1}\cdot\exp\left(-\gamma\cdot d(x,x')-\mathrm{kq}\cdot d_M(X)\right).
\]
\item[W8.] \emph{(Rows.)} Let $\mathscr{Y}$ be a subset of $\pi_k$ with union $Y$, let $s\in\mathrm{Pd}_k$ satisfy $s_{\mathrm{pc}}=0$ for $\mathrm{pc}\notin\mathscr{Y}$, let $Y_1$ be a link-component of $Y$, and let $\mathrm{Op}(s,A)$ be one of $H_k(s,A)$, $\Gamma_k(s,A)$.
If the row index $t$ is contained in $Y_1$, then $\mathrm{Op}(s,A)(t,t')=0$ unless $t'$ is contained in $Y_1$, and $\mathrm{Op}(s,A)(t,\cdot\,)=\mathrm{Op}(\tilde{s},\tilde{A})(t,\cdot\,)$ whenever $\tilde{s}\in\mathrm{Pd}_k$ vanishes outside $\mathscr{Y}$ and equals $s$ on the cubes of $Y_1$, and $\tilde{A}\in\tilde{G}_k$ equals $A$ on the bonds that meet $Y_1$.
\item[W9.] \emph{(Columns.)} Let $\mathscr{Y}$, $Y$, $s$, $Y_1$, $\mathrm{Op}$, $\tilde{s}$, $\tilde{A}$ be as in W8.
If the column index $t'$ is contained in $Y_1$, then $\mathrm{Op}(s,A)(t,t')=0$ unless $t$ is contained in $Y_1$, and $\mathrm{Op}(s,A)(\cdot\,,t')=\mathrm{Op}(\tilde{s},\tilde{A})(\cdot\,,t')$.
\item[W10.] For $\mathscr{Y}$, $Y$, $s$ as in W8, the row $H_k(s,A)(x,\cdot\,)$ vanishes if $x$ is contained in no link-component of $Y$.
\item[W11.] For $\mathscr{Y}$, $Y$, $s$ as in W8, the row of $H_k(s,A)$ at a fermion site $x$ vanishes unless $\mathscr{Y}$ contains all cubes of $\pi_k$ that meet $K_x=K_c$.
\item[W12.] The kernels of $M_k(X,A)$ and of $S_{k,y}(X,A)$ vanish unless both indices lie in (over) $X$; and $\mathrm{Op}(X,A)=\mathrm{Op}(X,\tilde{A})$ for these two families whenever $A,\tilde{A}\in\tilde{G}_k$ coincide on the bonds that meet $X$.
\item[W13.] Let $\boldsymbol{\lambda}$ be a real function of $c$ on $T_k$ and $\boldsymbol{\lambda}^{\mathrm{u}}$ its restriction to $T^0_k$.
Then
\begin{gather*}
H_k(s,A-d\boldsymbol{\lambda})=\exp\left(ie_k\cdot\boldsymbol{\lambda}\right)\cdot H_k(s,A)\cdot\exp\left(-ie_k\cdot\boldsymbol{\lambda}^{\mathrm{u}}\right),\\
O(\cdot\,,A-d\boldsymbol{\lambda})=\exp\left(ie_k\cdot\boldsymbol{\lambda}^{\mathrm{u}}\right)\cdot O(\cdot\,,A)\cdot\exp\left(-ie_k\cdot\boldsymbol{\lambda}^{\mathrm{u}}\right)\qquad(O=\Gamma_k,M_k),\\
S_{k,y}(X,A-d\boldsymbol{\lambda})=\exp\left(ie_k\cdot\boldsymbol{\lambda}\right)\cdot S_{k,y}(X,A)\cdot\exp\left(-ie_k\cdot\boldsymbol{\lambda}\right);
\end{gather*}
at $k=0$ the right side of the second relation is $\exp\left(ie_k\cdot\boldsymbol{\lambda}\right)\cdot\Gamma_0(s,A)\cdot\exp\left(-ie_k\cdot\boldsymbol{\lambda}\right)$.
\item[W14.] $\mathrm{Op}(\tilde{g}s,\tilde{g}A)\cdot\tilde{g}=\tilde{g}\cdot\mathrm{Op}(s,A)$ and $\mathrm{Op}(\tilde{g}X,\tilde{g}A)\cdot\tilde{g}=\tilde{g}\cdot\mathrm{Op}(X,A)$ for $\tilde{g}\in\mathbb{G}_k^{L}$; for $S_{k,y}(X,A)$ this is required for $\tilde{g}\in\mathbb{G}^{L}_{k,3}$ only (Section~I.5.8).
\item[W15.] $\Gamma_k(s,A)^T=\mathsf{C}^{-1}\cdot\Gamma_k(s,-A)\cdot\mathsf{C}$ and $M_k(X,A)^T=\mathsf{C}^{-1}\cdot M_k(X,-A)\cdot\mathsf{C}$, the transposes being taken for the pairings of Sections~I.1.6 and~I.3.6, at $k=0$ for $\langle\cdot\,,\cdot\rangle^{\mathrm{ch}}$.
\item[W16.] $\overline{O\left(\overline{s},\overline{A}\right)}=\varkappa^{-1}\cdot O(s,-A)\cdot\varkappa$ for $O=H_k,\Gamma_k$, and $\overline{M_k\left(X,\overline{A}\right)}=\varkappa^{-1}\cdot M_k(X,-A)\cdot\varkappa$;
here $\overline{O}$ is the operator $f\mapsto\overline{O\,\overline{f}}$, as in (I.O5),
and $\overline{s}:=\left(\overline{s_{\mathrm{pc}}}\right)_{\mathrm{pc}\in\pi_k}$ is the family of the complex conjugates of the $s_{\mathrm{pc}}$, which lies again in $\mathrm{Pd}_k$, because $\left|\overline{s_{\mathrm{pc}}}\right|=\left|s_{\mathrm{pc}}\right|$.
\end{description}

\emph{Barred family.}
The following assertions are consequences of the definition $\bar{H}_k(s,A)=\mathsf{C}^{-1}\cdot H_k(s,-A)\cdot\mathsf{C}$ and require no further property.
Let $H_k(s,A)$ have W1, W2, W5, W8--W11, W13, W14 and W16 for $(\alpha_0,\gamma,\mathrm{kq};C)$.
Then $\bar{H}_k(1,A)=\bar{H}_k(A)$; W2 and W8--W11 hold for $\bar{H}_k(s,A)$ as stated; W5 holds for $\bar{H}_k(s,A)$ with $\delta_{\alpha,-A}$ in place of $\delta_{\alpha,A}$ and with the same $C$; and
\begin{gather*}
\bar{H}_k(s,A-d\boldsymbol{\lambda})=\exp\left(-ie_k\cdot\boldsymbol{\lambda}\right)\cdot\bar{H}_k(s,A)\cdot\exp\left(ie_k\cdot\boldsymbol{\lambda}^{\mathrm{u}}\right),\\
\bar{H}_k(\tilde{g}s,\tilde{g}A)\cdot\mathrm{glb}=\mathrm{glb}\cdot\bar{H}_k(s,A)\quad\text{for the }\tilde{g}\text{ of W14 (}\mathrm{glb}\text{: Section~G1.10)},\\
\overline{\bar{H}_k\left(\overline{s},\overline{A}\right)}=\varkappa^{-1}\cdot\bar{H}_k(s,-A)\cdot\varkappa.
\end{gather*}
\begin{proof}
One uses the properties of $H_k$ at the point $-A$, which lies in $\tilde{G}_k$ by the relation (G2) of \cite[Section~16.2.7]{QEDI}, together with the following facts: $\mathsf{C}$ is real and monomial, and $\mathsf{C}^{-1}=\mathsf{C}^T$ ((I.c3), (I.c6)); the phase in (I.1.7) is a scalar, so that $\delta_{\alpha,-A}$ commutes with constant matrices; $\left(S^{-1}\right)^T=\pm\mathsf{C}^{-1}\cdot S\cdot\mathsf{C}$ (Section~I.4.1); and $\varkappa\cdot\mathsf{C}=-\mathsf{C}\cdot\varkappa$ ((I.c4)).
\end{proof}

\begin{description}
\item[W17.] \emph{(Level $0$, fermionic part.)} $H_0(s,A)$ has W1, W2, W8--W11, W13, W14, W16, and W5 with the constant $4\exp(2\gamma)\cdot M^{27\alpha_0}$; $\Gamma_0(s,A)$ and $S_{0,y}(X,A)$ have W1, W2, W6--W9 and W12--W16, each as far as it concerns the family, with W6 and W7 in their forms for $k=0$.
\end{description}
The two sets of parameters are
\[
\mathrm{par}^{0}:=\left(\alpha_0,\gamma_{\mathrm{w}},\kappa_{\mathrm{w}}+\kappa_{0};C_{w}^{(1)}\right),\qquad\mathrm{par}:=\left(\alpha_0,\gamma_{\mathrm{w}},\kappa_{\mathrm{w}};C_{w}\right).
\]

\begin{named}{Remark}
For the families (p3)--(p5) the definitions and properties of Section~G1.11 take the place of those of \cite[Sections~5.4.1 and~5.4.2]{QEDI}; W1, W2 and W5--W17 correspond to the properties 1, 2 and 5--17 of \cite[Section~5.4.2]{QEDI}.
They are read together with Section~I.5.8.
\end{named}

\subsection*{1. The free operator}

Throughout Propositions~1.1--1.7 the following data are fixed: $\rho\in(0,1)$, the standard weights and the spin table of Section~I.2.1; $b$, $\alpha$ and $L$ are as in Section~I.8.
The constants of Propositions~1.1--1.5 depend only on $\rho$ and the weights, and those of Propositions~1.6 and~1.7 only on $(b,\alpha,L,\rho)$.
The mass $\mathbf{m}$, resp.\ $\mu$, of Propositions~1.1--1.7 is a running mass; it is not the lattice mass $\boldsymbol{\mu}$ of (I.K10).

\begin{named}{Proposition 1.1}[Light and heavy part of $D_A$, in every field]
Let either $\mathbb{T}$ be a block torus of level $k\ge0$ (Section~G1.2), with the lattice $\mathrm{Lab}$, or let $\mathrm{Lab}=\mathbb{L}_k$ without periodicity, in which case the spinors are square summable.
Let $e_k>0$, let $A$ be a complex connection, let $\xi$ be a $4$-component spinor on $\mathrm{Lab}$, with the same sign $\boldsymbol{\sigma}$ as the spinors of $\mathbb{T}$, and let $f$ be a spinor on $\mathrm{Lab}\times O_{\mathrm{cub}}$.
For $a\in\{\pm e_{\mu}\}$ put
\[
\left(\mathrm{dc}_{A,a}\xi\right)(c):=\eta_k^{-1}\left(u_A(c,c+\eta_k a)\cdot\xi(c+\eta_k a)-\xi(c)\right),\qquad\|\mathrm{dc}_A\xi\|^2:=\sum_{\mu=1}^{3}\left\|\mathrm{dc}_{A,e_{\mu}}\xi\right\|_{L^2}^2,
\]
the $L^2$ norm being formed with $\eta_k^3\sum_c$.
The blocks $D^{\mathrm{ll}}_{A}$, $D^{\mathrm{lh}}_{A}$, $D^{\mathrm{hl}}_{A}$, $D^{\mathrm{hh}}_{A}$ and the projections $\mathrm{Plt}$, $\mathrm{Phv}$ are those of (I.J4) and (I.J6).
\begin{itemize}
\item[(a)] For every complex $A$, $J^{*}\cdot D_A\cdot J=D^{W,\rho}_A$, the operator of (I.3.2).
\item[(b)] For every complex $A$ and $x=(c,g)$,
\[
\left(D_A\cdot J\,\xi\right)(c,g)=-M_{3}(x)\cdot\left(\mathrm{dc}_{A,ge_1}\xi\right)(c).
\]
For real $A$,
\begin{gather*}
\left\|D^{\mathrm{hl}}_{A}\xi\right\|_{L^2}\le c_{\mathrm{hl}}\cdot\|\mathrm{dc}_A\xi\|,\qquad\left|\left(\xi,D^{\mathrm{lh}}_{A}f\right)\right|\le c_{\mathrm{hl}}\cdot\|\mathrm{dc}_A\xi\|\cdot\|f\|_{L^2},\\
c_{\mathrm{hl}}:=\sqrt{3}\,(1+\rho).
\end{gather*}
\item[(c)] For real $A$ and $f=\mathrm{Phv}\,f$,
\[
\operatorname{Re}(f,D_Af)^{\mathrm{ch}}\ge\frac{\hat{g}}{\eta_k}\,(f,f)^{\mathrm{ch}},\qquad\hat{g}:=\frac{\rho\cdot\lambda_{\mathrm{int}}}{2}.
\]
Consequently, for every real $t>-\hat{g}/\eta_k$ the operator $D^{\mathrm{hh}}_{A}+t$ is invertible on the range of $\mathrm{Phv}$, and
\[
\left\|\left(D^{\mathrm{hh}}_{A}+t\right)^{-1}\right\|_{L^2}\le\frac{\eta_k}{\eta_k\cdot t+\hat{g}}.
\]
\item[(d)] For real $A$, every $f$ and $\xi:=J^{*}f$,
\[
\eta_k^{-2}\cdot\left(f,\nabla_A^{*}\nabla_A f\right)^{\mathrm{ch}}\ge c_{*}\cdot\|\mathrm{dc}_A\xi\|^2,\qquad c_{*}:=\frac{\min\left(w_3,\lambda_{\mathrm{int}}\right)}{18},
\]
and
\[
\left(f,\nabla_A^{*}\nabla_A f\right)^{\mathrm{ch}}\ge\lambda_{\mathrm{int}}\cdot\|\mathrm{Phv}\,f\|_{L^2}^2.
\]
Both inequalities hold also locally.
Let $E$ be a set of photon sites; let $\|\mathrm{dc}_A\xi\|_E^2$ be the sum of $\eta_k^3\left|\left(\mathrm{dc}_{A,a}\xi\right)(c)\right|_2^2$ over the unoriented photon bonds $(c,c+\eta_k a)$ with both end points in $E$; let $\left(f,\nabla_A^{*}\nabla_A f\right)_E$ be the part of $\left(f,\nabla_A^{*}\nabla_A f\right)^{\mathrm{ch}}$ (see (I.O2)) that comes from the bonds $\{x,x^{(i)}\}$ with both end points over $E$; and let $\|\cdot\|_{L^2,E}$ be the norm over $E\times O_{\mathrm{cub}}$.
Then
\[
\eta_k^{-2}\cdot\left(f,\nabla_A^{*}\nabla_A f\right)_E\ge c_{*}\cdot\|\mathrm{dc}_A\xi\|_E^2,\qquad\left(f,\nabla_A^{*}\nabla_A f\right)_E\ge\lambda_{\mathrm{int}}\cdot\|\mathrm{Phv}\,f\|_{L^2,E}^2.
\]
\item[(e)] $J$, $J^{*}$, $\mathrm{Plt}$ and $\mathrm{Phv}$ commute with the multiplication by functions of $c$, with the translations of $\mathrm{Lab}$, with the maps $\mathrm{mS}$ of Section~G1.2 and with constant spin matrices.
\item[(f)] \emph{(The same for every mesh, side and character, and on finite sets.)}
Consider the following data (GS): a mesh $\mathrm{me}>0$; a translate $\mathrm{Lab}=t+\mathrm{me}\,\mathbb{Z}^3$ of $\mathrm{me}\,\mathbb{Z}^3$; $R\in\mathrm{me}\,\mathbb{N}$ or $R=\infty$; $\theta\in\mathbb{R}^3$; a number $e_k>0$; a complex function $A$ on the oriented bonds of $\mathrm{Lab}$ with $A(c',c)=-A(c,c')$ and, if $R<\infty$, $A(c+Re_{\mu},c'+Re_{\mu})=A(c,c')$.
The spinors are functions $f:\mathrm{Lab}\times O_{\mathrm{cub}}\to\mathbb{C}^4$ with $f(c+Re_{\mu},g)=-\exp(i\theta_{\mu})\cdot f(c,g)$ and $\xi:\mathrm{Lab}\to\mathbb{C}^4$ with $\xi(c+Re_{\mu})=-\exp(i\theta_{\mu})\cdot\xi(c)$; for $R=\infty$ there is no condition.
Put $u_A(c,c'):=\exp\left(ie_k\cdot\mathrm{me}\cdot A(c,c')\right)$.
The objects $x^{(i)}$, $c_{3}$, $\hat{d}_i$, $M_{i}(x)$, $U_A(x,x^{(i)})$, $R_i(A)$, $D_A$, $\nabla_A^{*}\nabla_A$, $J_A$, $J$, $J^{*}$, $\mathrm{Plt}$, $\mathrm{Phv}$, the four blocks, $\mathrm{dc}_{A,a}$, $\|\mathrm{dc}_A\xi\|$, $\int\mathrm{d}x$, $(\cdot\,,\cdot)^{\mathrm{ch}}$ and the $L^2$ norms are formed by the displays (I.1.2), (I.2.1)--(I.2.4), (I.3.1), (I.J4), (I.J6), by those of Section~I.1.6 and by those of the data above, with $\mathrm{me}$ and $\mathrm{Lab}$ in place of $\eta_k$ and $\mathbb{L}_k$.
Sites, bonds $\{x,x^{(i)}\}$ and photon bonds are counted modulo $R\mathbb{Z}^3$, as classes of sites and bonds of the covering lattice; this matters for $R=\mathrm{me}$ and $R=2\,\mathrm{me}$.
Wherever a norm or a pairing occurs, $f$ and $\xi$ are square summable.

For all data (GS), (I.s1)--(I.s4) and (I.2.4) hold; (a)--(d) hold with $\mathrm{me}$ in place of $\eta_k$, and (e) holds without the maps $\mathrm{mS}$, with the same $\lambda_{\mathrm{int}}$, $\hat{g}$, $c_{\mathrm{hl}}$, $c_{*}$.
In the local form of (d) the set $E$ is either an $R\mathbb{Z}^3$-periodic subset of $\mathrm{Lab}$, the sums being taken modulo $R\mathbb{Z}^3$, or a finite subset of $\mathrm{Lab}$; in the second case $f$ is an arbitrary function on $E\times O_{\mathrm{cub}}$, and the sums extend over the sites and bonds of $\mathrm{Lab}$ themselves.

For real $A$, moreover, the following holds.
$R_i(A)$ is unitary and self-adjoint.
$\left(f,\nabla_A^{*}\nabla_A f\right)^{\mathrm{ch}}$ is the sum over bonds displayed in (I.O2), with the weight $\mathrm{me}^3/48$, so that $\nabla_A^{*}\nabla_A=\left(\nabla_A^{*}\nabla_A\right)^{*}\ge0$.
$J_A^{*}=-J_A$, and
\begin{gather*}
\operatorname{Re}(f,D_Af)^{\mathrm{ch}}=\frac{\rho}{2\,\mathrm{me}}\left(f,\nabla_A^{*}\nabla_A f\right)^{\mathrm{ch}},\\
D_A^{*}=\gamma_5\cdot D_A\cdot\gamma_5,\qquad\|D_A\|_{L^2}\le\frac{(1+\rho)\cdot\mathrm{sw}}{\mathrm{me}}.
\end{gather*}
$J^{*}\cdot J=1$, $J^{*}$ is the adjoint of $J$, $J$ is an isometry, $\mathrm{Plt}$ and $\mathrm{Phv}$ are orthogonal projections, and for every $c$
\[
\frac{1}{48}\sum_{g}|f(c,g)|_2^2=\left|(J^{*}f)(c)\right|_2^2+\frac{1}{48}\sum_{g}\left|(\mathrm{Phv}\,f)(c,g)\right|_2^2.
\]
Further,
\[
\left(f,\nabla_A^{*}\nabla_A f\right)^{\mathrm{ch}}\le3\,\mathrm{me}^2\cdot\|\mathrm{dc}_A\,J^{*}f\|^2+\left(30+24\sqrt{2}\right)\cdot\|\mathrm{Phv}\,f\|_{L^2}^2,
\]
and $0<\lambda_{\mathrm{int}}\le12+24\sqrt{2}<48$.

The data (GS) comprise the following cases: the data at the head of this proposition, with $\mathrm{me}=\eta_k$, $\theta=0$ for $\boldsymbol{\sigma}=-1$ and $\theta=(\pi,\pi,\pi)$ for $\boldsymbol{\sigma}=+1$; the spaces of Section~G1.4, with $\mathrm{Lab}=\mathrm{me}\,\mathbb{Z}^3$ and $A=0$; and the fibre $p\in\mathbb{R}^3$, with $\mathrm{me}=R=1$, $\mathrm{Lab}=\mathbb{Z}^3$, $A=0$ and $\theta_{\mu}=p_{\mu}-\pi$.
In the fibre, $f(c,g)=\exp(ip\cdot c)\,u(g)$ and $\xi(c)=\exp(ip\cdot c)\,v$; the operator $D_0$ acts as $D^{(p)}\big|_{\boldsymbol{\mu}=0}$ of (I.2.6), and $\|\mathrm{dc}_0\xi\|^2=2\,\mathrm{tw}(p)\cdot|v|_2^2$.
\end{itemize}
No constant depends on $A$, $e_k$, $k$ or on the block torus $\mathbb{T}$, and in (f) none depends on $\mathrm{me}$, $\mathrm{Lab}$, $R$, $\theta$ or $E$.
\end{named}

\begin{named}{Remark}
The operator of \cite{QEDI} has no heavy part, and Proposition~1.1 has no counterpart there.
It contains (I.O8).
\end{named}

\begin{named}{Proposition 1.2}[Bloch symbol and its $4\times4$ Schur symbol]
Lattice units are used.
Let $\mathbf{m}\in[0,1]$, let $p\in\mathbb{C}^3$, and let $D^{(p)}$ be the matrix of (I.2.6) at $\boldsymbol{\mu}=\mathbf{m}$, acting on ${V_{\mathrm{v}}}$ with the pairing $\frac{1}{48}\sum_{g}$.
One has ${V_{\mathrm{v}}}=\operatorname{ran}J\oplus\ker J^{*}$; put
\begin{gather*}
D^{\mathrm{ll}}(p):=J^{*}\cdot D^{(p)}\cdot J,\qquad D^{\mathrm{lh}}(p):=J^{*}\cdot D^{(p)}\cdot\mathrm{Phv},\qquad D^{\mathrm{hl}}(p):=\mathrm{Phv}\cdot D^{(p)}\cdot J,\\
D^{\mathrm{hh}}(p):=\mathrm{Phv}\cdot D^{(p)}\cdot\mathrm{Phv}\ \text{on}\ \ker J^{*},
\end{gather*}
all formed with the mass $\mathbf{m}$, and
\[
\mathrm{tw}(p):=\sum_{\mu}\left(1-\cos p_{\mu}\right),\qquad\sin p:=\left(\sin p_{\mu}\right)_{\mu},\qquad\mathrm{rp}(p):=\left(|p|_2^2+\mathbf{m}^2\right)^{\frac{1}{2}}\quad\text{for real }p.
\]
The constants are functions of $\rho$, chosen in the following order:
\begin{gather*}
\theta_{0}>0\ \text{with}\ 3(1+\rho)\cdot\left(\exp(\theta_{0})-1\right)=\frac{\hat{g}}{2};\qquad K_{\rho}:=\rho+\frac{2c_{\mathrm{hl}}^2}{\hat{g}};\qquad t_{0}:=\frac{1}{4K_{\rho}^2}\quad\left(K_{\rho}>1,\ t_{0}\le\tfrac{1}{4},\ c_{*}\le\tfrac{1}{3}\right);\\
c_{\mathrm{el}}:=\min\left\{0.22,\ \frac{\rho\cdot c_{*}\cdot t_{0}}{\pi\sqrt{3}}\right\};\qquad K_{\mathrm{Str}}\ \text{(defined in (S3))};\qquad\kappa_{(0)}:=\min\left\{\frac{\theta_{0}}{2},\ \frac{c_{\mathrm{el}}\cdot\theta_{0}}{12K_{\mathrm{Str}}}\right\}.
\end{gather*}
\begin{itemize}
\item[(S1)] $D^{\mathrm{hh}}(p)$ is invertible for $|\operatorname{Im}p_{\mu}|\le\theta_{0}$ ($\mu=1,2,3$); there $\left|D^{\mathrm{hh}}(p)^{-1}\right|\le2/\hat{g}$, and $\left|D^{\mathrm{hh}}(p)^{-1}\right|\le1/(\hat{g}+\mathbf{m})$ for real $p$ (operator norms on $\ker J^{*}$).
\item[(S2)] $D^{\mathrm{lh}}(0)=0$ and $D^{\mathrm{hl}}(0)=0$; for real $p$
\[
\left|D^{\mathrm{lh}}(p)\right|,\ \left|D^{\mathrm{hl}}(p)\right|\le c_{\mathrm{hl}}\cdot\left(2\,\mathrm{tw}(p)\right)^{\frac{1}{2}}\le c_{\mathrm{hl}}\cdot|p|_2.
\]
\item[(S3)] The $4\times4$ matrix
\[
\mathrm{Dsc}_{\mathbf{m}}(p):=D^{\mathrm{ll}}(p)-D^{\mathrm{lh}}(p)\cdot D^{\mathrm{hh}}(p)^{-1}\cdot D^{\mathrm{hl}}(p)
\]
is holomorphic and $2\pi\mathbb{Z}^3$-periodic on $\left\{|\operatorname{Im}p_{\mu}|\le\theta_{0}\right\}$, and
\[
K_{\mathrm{Str}}:=\sup\left\{\left|\mathrm{Dsc}_{\mathbf{m}}(p)\right|_2:\ |\operatorname{Im}p_{\mu}|\le\theta_{0},\ 0\le\mathbf{m}\le1\right\}<\infty.
\]
For real $p$,
\[
\mathrm{Dsc}_{\mathbf{m}}(p)=\mathbf{m}+i\gamma\cdot\sin p+\mathrm{Er}_{\mathbf{m}}(p),\qquad\left|\mathrm{Er}_{\mathbf{m}}(p)\right|_2\le K_{\rho}\cdot\mathrm{tw}(p).
\]
\item[(S4)] For real $p$ and $v\in\mathbb{C}^4$,
\[
\operatorname{Re}\left(v,\mathrm{Dsc}_{\mathbf{m}}(p)\,v\right)\ge\left(\mathbf{m}+\rho\cdot c_{*}\cdot\mathrm{tw}(p)\right)\cdot|v|_2^2.
\]
\item[(S5)] \emph{(Ellipticity.)} For $p\in[-\pi,\pi]^3$,
\[
\left|\mathrm{Dsc}_{\mathbf{m}}(p)\,v\right|_2\ge c_{\mathrm{el}}\cdot\mathrm{rp}(p)\cdot|v|_2.
\]
\item[(S6)] For every multi-index $\mathrm{ak}$ there is $A_{\mathrm{ak}}=A_{\mathrm{ak}}(\rho)$ with
\[
\left|\partial^{\mathrm{ak}}\mathrm{Dsc}_{\mathbf{m}}(p)^{-1}\right|_2\le A_{\mathrm{ak}}\cdot\mathrm{rp}(p)^{-1-|\mathrm{ak}|}\qquad\text{on }[-\pi,\pi]^3\setminus\{0\}.
\]
\item[(S7)] \emph{(Strip of a width proportional to the mass.)} Let $\mathbf{m}>0$ and $p=p'+ip''$ with $p'\in[-\pi,\pi]^3$ and $|p''_{\mu}|\le\kappa_{(0)}\cdot\mathbf{m}$.
Then
\[
\left|\mathrm{Dsc}_{\mathbf{m}}(p)\,v\right|_2\ge\frac{c_{\mathrm{el}}}{2}\cdot\mathrm{rp}(p')\cdot|v|_2,\qquad\left|\partial^{\mathrm{ak}}\mathrm{Dsc}_{\mathbf{m}}(p)^{-1}\right|_2\le A'_{\mathrm{ak}}\cdot\mathrm{rp}(p')^{-1-|\mathrm{ak}|},
\]
with $A'_{\mathrm{ak}}=A'_{\mathrm{ak}}(\rho)$.
\item[(S8)] Wherever $\mathrm{Dsc}_{\mathbf{m}}(p)$ is invertible ($|\operatorname{Im}p_{\mu}|\le\theta_{0}$), put $\mathrm{Sc}:=\mathrm{Dsc}_{\mathbf{m}}(p)^{-1}$ and $\mathrm{Hi}:=D^{\mathrm{hh}}(p)^{-1}$.
Then, in the decomposition $\operatorname{ran}J\oplus\ker J^{*}$,
\begin{gather*}
\left(D^{(p)}\right)^{-1}=\left(\begin{array}{cc}\mathrm{Sc} & -\mathrm{Sc}\cdot D^{\mathrm{lh}}\cdot\mathrm{Hi}\\ -\mathrm{Hi}\cdot D^{\mathrm{hl}}\cdot\mathrm{Sc} & \mathrm{Hi}+\mathrm{Hi}\cdot D^{\mathrm{hl}}\cdot\mathrm{Sc}\cdot D^{\mathrm{lh}}\cdot\mathrm{Hi}\end{array}\right),\\
\det D^{(p)}=\det D^{\mathrm{hh}}(p)\cdot\det\mathrm{Dsc}_{\mathbf{m}}(p).
\end{gather*}
\item[(S9)] For real $p$ the matrices $D^{(p)}$, $D^{\mathrm{hh}}(p)$ and $\mathrm{Dsc}_{\mathbf{m}}(p)$ commute with $\Sigma_3$, and in the two eigenspaces of $\Sigma_3$
\[
D^{\mathrm{hh}}(p)=\operatorname{diag}\left(\mathrm{Hh}(p),\mathrm{Hh}(p)^{*}\right),\qquad\mathrm{Dsc}_{\mathbf{m}}(p)=\operatorname{diag}\left(\mathrm{Ds}(p),\mathrm{Ds}(p)^{*}\right)
\]
with matrices $\mathrm{Hh}(p)$, $\mathrm{Ds}(p)$.
Hence $\det D^{\mathrm{hh}}(p)=\left|\det\mathrm{Hh}(p)\right|^2>0$ for all real $p$ and all $\mathbf{m}\in[0,1]$, and $\det\mathrm{Dsc}_{\mathbf{m}}(p)=\left|\det\mathrm{Ds}(p)\right|^2\ge0$, with strict inequality unless $\mathbf{m}=0$ and $p\in2\pi\mathbb{Z}^3$.
\end{itemize}
All bounds are uniform in $\mathbf{m}\in[0,1]$.
\end{named}

\begin{named}{Remark}
Proposition~1.2 corresponds to \cite[Lemma~12.5]{QEDI} together with \cite[Remark~12.2]{QEDI}.
In \cite[Lemma~12.5]{QEDI} the symbol is the normal $4\times4$ matrix $i\gamma\cdot\sin p+w+m$, with the scalar denominator $N(p)$ and the inverse $a_m(p)=n(p)/N(p)$; these have no counterpart here, since the $192\times192$ symbol $D^{(p)}$ is not normal, and the block form (S8) together with the properties (S3)--(S7) of the Schur symbol takes their place.
The identity of \cite[Remark~12.2]{QEDI} has no counterpart here.
The constants $A_{\mathrm{ak}}$ depend on $\rho$, whereas those of \cite[Lemma~12.5]{QEDI} are absolute.
The clauses (S1), (S2) and (S7) have no counterpart in \cite[Lemma~12.5]{QEDI}.
The Bloch convention is $f(c,g)=\exp(ip\cdot c)\,u(g)$, with periodic symbols (Section~I.2.6); in this form the symbols are $2\pi\mathbb{Z}^3$-periodic, as \cite[Lemma~12.4]{QEDI} requires.
\end{named}

\subsubsection*{Proposition~1.3 (three facts: real part, coercivity, commutator; adjoint and transpose)}

\begin{named}{Proposition 1.3}
Let $\mathrm{me}>0$, $R\in\mathrm{me}\,\mathbb{N}\cup\{\infty\}$ and $\theta\in\mathbb{R}^3$ be data as in Section~G1.4, and let $u\in\mathrm{HH}(\mathrm{me},R;\theta)$.
\begin{itemize}
\item[(a)] The real part of the quadratic form of $D$ is
\[
\operatorname{Re}(u,Du)^{\mathrm{ch}}=\frac{\rho}{2\,\mathrm{me}}\cdot\left(u,\nabla^{*}\nabla u\right)^{\mathrm{ch}}.
\]
\item[(b)] (Coercivity.)
One has
\[
\|Du\|_{L^2}^2\ge c_b\cdot\mathrm{me}^{-2}\cdot\left(u,\nabla^{*}\nabla u\right)^{\mathrm{ch}},
\]
where
\begin{gather*}
c_b:=\left[3A_1^2+\left(30+24\sqrt{2}\right)\cdot A_2^2\right]^{-1},\\
A_1:=\frac{2\sqrt{2}\,K_{\rho}}{\rho\cdot c_{*}}\cdot\left(1+\frac{6(1+\rho)}{\hat{g}}\right),\qquad
A_2:=\frac{1+c_{\mathrm{hl}}\cdot A_1}{\hat{g}}.
\end{gather*}
\item[(c)] Let $\psi$ be an arbitrary function of $c$, and let $u$ be an arbitrary function on $\mathrm{me}\,\mathbb{Z}^3\times O_{\mathrm{cub}}$; no periodicity and no summability is required of either.
Then for every $x=(c,g)$
\[
\left([D,\psi]\,u\right)(x)=-\mathrm{me}^{-1}\cdot M_{3}(x)\cdot\left(\psi(c_{3})-\psi(c)\right)\cdot u\left(x^{(3)}\right).
\]
Here $x\mapsto x^{(3)}$ is a bijection, $|c_{3}-c|=\mathrm{me}$, and $|M_{3}(x)|_2=|M_{3}(x)|_c=3(1+\rho)\le6$.
\item[(d)] On $\mathrm{HH}(\mathrm{me},R;\theta)$ the adjoint of $D$ is $D^{*}=\gamma_5\,D\,\gamma_5$.
The transpose is $D^{T}=\mathsf{C}^{-1}D\,\mathsf{C}$, where $D^{T}$ acts on $\mathfrak{S}^{\mathrm{ch}}(\mathrm{me},R;-\theta)$ (Section~G1.4).
For $R=\infty$ the pairing is not defined on all pairs of functions; there $D^{T}f'$ denotes the function such that $\left\langle D^{T}f',f\right\rangle^{\mathrm{ch}}=\left\langle f',Df\right\rangle^{\mathrm{ch}}$ for every $f$ of finite support, and on $\mathrm{HH}$ it is the transpose with respect to the pairing of $\mathrm{HH}$ with $\mathrm{HH}$.
\end{itemize}
\end{named}

\begin{named}{Remark}
The counterpart of Proposition~1.3 is \cite[Lemma~12.1(a)--(d)]{QEDI}.
In (a) the factor $\rho$ appears, and $\mathrm{me}^{-2}\cdot\left(u,\nabla^{*}\nabla u\right)^{\mathrm{ch}}$ stands in the place of the squared norm of $\nabla^{+}u$.
The identity of \cite[Lemma~12.1(b)]{QEDI} holds because the Hermitian and the anti-Hermitian part of the free Wilson operator with Wilson parameter one commute; for $D$ the corresponding identity is false, and (b) asserts only the inequality, with the constant $c_b$, which depends on $\rho$.
In (c) a single term takes the place of the six terms of \cite[Lemma~12.1(c)]{QEDI}, and the inequality of \cite[Lemma~12.1(c)]{QEDI} for the projections $\Pi^{\pm}_{\mu}$ of \cite[Section~12.1]{QEDI} has no counterpart.
In (d) the description of $D^{*}$ and of $D^{T}$ as the operators formed with $-\gamma_{\mu}$ and with $-\gamma_{\mu}^{T}$ is replaced by the conjugation with the matrices $\gamma_5$ and $\mathsf{C}$ of Section~I.2.1, and the character $-\theta$ of the space on which $D^{T}$ acts is specified.
\end{named}

\subsubsection*{Proposition~1.4 (Poincar\'e inequality on a block times the cluster)}

\begin{named}{Proposition 1.4}
In the setting of Section~G1.4, let $E$ be a cube of $n^3$ sites of $\mathrm{me}\,\mathbb{Z}^3$, where $n\ge1$, put $\ell_E:=n\cdot\mathrm{me}$, and let $u$ be a function on $E\times O_{\mathrm{cub}}$.
Then
\[
\|u-P_Eu\|_{L^2,E}^2\le C_{\mathrm{Po}}\cdot\ell_E^2\cdot\mathrm{me}^{-2}\cdot\left(u,\nabla^{*}\nabla u\right)_E,\qquad
C_{\mathrm{Po}}:=\frac{1}{8c_{*}}+\frac{1}{\lambda_{\mathrm{int}}}.
\]
\end{named}

\begin{named}{Remark}
In \cite[Lemma~12.3]{QEDI} the inequality is stated for functions of $c$ alone, with the constant $\frac{1}{8}$ in the place of $C_{\mathrm{Po}}$, and with the sum $\|\nabla u\|_E^2$ of the squared difference quotients of $u$ along the bonds of $\mathrm{me}\,\mathbb{Z}^3$ inside $E$ in the place of $\mathrm{me}^{-2}\cdot\left(u,\nabla^{*}\nabla u\right)_E$.
In Proposition~1.4 the right-hand side is the energy of the bonds of all four types inside $E\times O_{\mathrm{cub}}$, and the constant is $C_{\mathrm{Po}}$.
\end{named}

\subsubsection*{Proposition~1.5 (free Green function)}

\begin{named}{Proposition 1.5}
First let the mesh be one.
For $\mathbf{m}\in[0,1]$ and for $x=(c,g)$, $x'=(c',g')$ in $\mathbb{Z}^3\times O_{\mathrm{cub}}$ let
\[
G_{\mathbf{m}}(x,x'):=\int_{[-\pi,\pi]^3}\frac{\mathrm{d}^3p}{(2\pi)^3}\,\exp\left(i\,p\cdot(c-c')\right)\left[\left(D^{(p)}\right)^{-1}\right](g,g'),
\]
where $D^{(p)}$ is the matrix (I.2.6) at $\boldsymbol{\mu}=\mathbf{m}$, and where $[O](g,g')$ is the kernel of an operator $O$ on ${V_{\mathrm{v}}}$ with respect to $\frac{1}{48}\sum_{g'}$, that is, $(Ou)(g)=\frac{1}{48}\sum_{g'}[O](g,g')\,u(g')$.
For $n\in\mathbb{Z}^3$ let
\[
\mathrm{Sc}^{v}_{\mathbf{m}}(n):=\int_{[-\pi,\pi]^3}\frac{\mathrm{d}^3p}{(2\pi)^3}\,\exp(i\,p\cdot n)\,\mathrm{Dsc}_{\mathbf{m}}(p)^{-1},
\]
$\mathrm{Dsc}_{\mathbf{m}}(p)$ being the Schur symbol of Proposition~1.2\,(S3).
There is a constant $C_G=C_G(\rho)$ such that the following holds for all $\mathbf{m}\in[0,1]$ and, in (f), for every $\alpha\in[0,1]$.
\begin{itemize}
\item[(a)] Let $D_0$ be the operator (I.2.3) at $k=0$ and $A=0$, acting on the columns.
Then
\[
\left((D_0+\mathbf{m})\,G_{\mathbf{m}}(\cdot,x')\right)(x)=48\cdot\mathbf{1}[x=x']\cdot\mathbf{1}_4.
\]
\item[(b)] Write
\[
G_{\mathbf{m}}(x,x')=\mathrm{Sc}^{v}_{\mathbf{m}}(c-c')+G^{h}_{\mathbf{m}}(x,x'),
\]
which defines $G^{h}_{\mathbf{m}}$.
Then
\begin{gather*}
\left|\mathrm{Sc}^{v}_{\mathbf{m}}(n)\right|_2\le C_G\cdot(1+|n|)^{-2},\qquad
\left|\mathrm{Sc}^{v}_{\mathbf{m}}(n+a)-\mathrm{Sc}^{v}_{\mathbf{m}}(n)\right|_2\le C_G\cdot(1+|n|)^{-3}\quad\left(a\in\{\pm e_{\mu}\}\right),\\
\left|G^{h}_{\mathbf{m}}(x,x')\right|_2\le C_G\cdot(1+|c-c'|)^{-3}.
\end{gather*}
It follows that
\begin{gather*}
\left|G_{\mathbf{m}}(x,x')\right|_2\le2C_G\cdot(1+|c-c'|)^{-2},\\
\left|G_{\mathbf{m}}\left(x^{(i)},x'\right)-G_{\mathbf{m}}(x,x')\right|_2\le20\,C_G\cdot(1+|c-c'|)^{-3}\qquad(i=0,1,2,3).
\end{gather*}
\item[(c)] Every finitely supported $v$ satisfies
\[
v(x)=\frac{1}{48}\sum_{x'}G_{\mathbf{m}}(x,x')\,\left((D_0+\mathbf{m})\,v\right)(x').
\]
\item[(d)] If $\mathbf{m}>0$, every bound of (b) holds with the additional factor $\exp\left(-\kappa_{(0)}\cdot\mathbf{m}\cdot|c-c'|\right)$, where $\kappa_{(0)}$ is the constant of Proposition~1.2; in the bounds on $\mathrm{Sc}^{v}_{\mathbf{m}}$ the number $|n|$ stands in the place of $|c-c'|$.
\item[(e)] Let $\mathrm{me}>0$ be a mesh, let $D$ be the operator of Section~G1.4, let $0\le\mu\le1/\mathrm{me}$, and put
\[
G(x,x'):=\mathrm{me}^{-2}\cdot G_{\mu\cdot\mathrm{me}}(x/\mathrm{me},x'/\mathrm{me}),\qquad x/\mathrm{me}:=(c/\mathrm{me},g).
\]
Then $(D+\mu)\,G(\cdot,x')$ is the kernel of the identity with respect to $\int\mathrm{d}x$ (Section~G1.4), namely $48\,\mathrm{me}^{-3}\cdot\mathbf{1}[x=x']\cdot\mathbf{1}_4$; moreover
\begin{gather*}
|G(x,x')|_2\le2C_G\cdot(\mathrm{me}+|c-c'|)^{-2},\\
\left|G\left(x^{(i)},x'\right)-G(x,x')\right|_2\le20\,C_G\cdot\mathrm{me}\cdot(\mathrm{me}+|c-c'|)^{-3}\qquad(i=0,\dots,3),
\end{gather*}
and every finitely supported $v$ satisfies
\[
v=\int\mathrm{d}x'\,G(\cdot,x')\,\left((D+\mu)\,v\right)(x').
\]
\item[(f)] (H\"older interpolation.)
For $x\ne\tilde{x}$ let $\mathrm{hd}:=d_X(x,\tilde{x})$.
Then
\[
\left|G(x,x')-G(\tilde{x},x')\right|_2\le6240\,C_G\cdot\mathrm{hd}^{\alpha}\cdot\left[(\mathrm{me}+|c-c'|)^{-2-\alpha}+(\mathrm{me}+|\tilde{c}-c'|)^{-2-\alpha}\right].
\]
\end{itemize}
In (e) and (f) all points are points of the covering lattice $\mathrm{me}\,\mathbb{Z}^3\times O_{\mathrm{cub}}$, and all distances are taken in it.
\end{named}

\begin{named}{Remark}
Proposition~1.5 corresponds to \cite[Proposition~12.6]{QEDI} together with \cite[Lemma~12.7(b)]{QEDI}, with the following differences.
The Green function is a function of $(c-c',g,g')$, not of a difference of sites.
The bounds on differences are stated along the bonds of all four types; along the types $0$, $1$, $2$ the part $\mathrm{Sc}^{v}_{\mathbf{m}}$ cancels.
The constant $C_G$ depends on $\rho$, and the factor $24\,C_G$ of \cite[Lemma~12.7(b)]{QEDI} becomes $6240\,C_G$, since two distinct sites $x$, $\tilde{x}$ are joined by a path of at most $39\,d_X(x,\tilde{x})/\mathrm{me}$ bonds, by (I.s6).
Clause (d) has no counterpart in \cite[Proposition~12.6]{QEDI}.
The delta function carries the normalisation $48$ of (I.ke1).
\end{named}

\subsubsection*{Proposition~1.6 (the blocked operator at zero field, levels $k\ge1$: sup, H\"older and exponential bounds)}

\begin{named}{Proposition 1.6}
Let $k\ge1$ be an integer and $\mathrm{me}=\eta_k=L^{-k}$; let $R\in L\,\mathbb{N}\cup\{\infty\}$, $\theta\in\mathbb{R}^3$, $\mu\in[0,1]$ and $y\in\mathbb{R}\cup\{\infty\}$.
Let $\boldsymbol{\alpha}_{y}$, $\boldsymbol{\beta}_{y}$ be the matrices of Section~I.3.5, formed with $b_{k}$.
Let $J P_k(0) J^{*}$ be the orthogonal projection of $\mathrm{HH}(\eta_k,R;\theta)$ onto the functions that are constant on each set $E\times O_{\mathrm{cub}}$, with $E$ ranging over the blocks of side $1$, that is, the blocks of level $k$ of Section~G1.4; for $z\in\mathbb{Z}^3$ the block of side $1$ with centre $z$ is denoted by $\Delta_z$.
Let $J P_{k+1}(0) J^{*}$ be the orthogonal projection onto the functions that are constant on each set $E\times O_{\mathrm{cub}}$, with $E$ ranging over the blocks of side $L$, that is, of level $k+1$, whose centres lie in $L\mathbb{Z}^3$.
On $\mathrm{HH}(\eta_k,R;\theta)$ let
\[
F_{y}(0):=D+\mu+\boldsymbol{\alpha}_{y}\,J P_k(0) J^{*}+\boldsymbol{\beta}_{y}\,J P_{k+1}(0) J^{*};
\]
for $y=\infty$, where $\boldsymbol{\alpha}_{\infty}=b_{k}$ and $\boldsymbol{\beta}_{\infty}=0$, this is $F_{\infty}(0)=D+\mu+b_{k}\,J P_k(0) J^{*}$.
The following constants are defined in this order:
\begin{gather*}
K_{1}:=L^{\frac{1}{2}}+(L\cdot b)^{\frac{1}{2}},\qquad
C_{I}:=\frac{1}{b_{-}}+\frac{2C_{\mathrm{Po}}\cdot K_{1}^2}{c_b}+\left(\frac{2C_{\mathrm{Po}}}{c_b}\right)^{\frac{1}{2}},\qquad
c_{\mathrm{wt}}:=\mathrm{e}\cdot\left(3(1+\rho)+6L^4\cdot b\right),\\
a_{0}:=\min\left\{\frac{1}{3},\frac{1}{2C_{I}\cdot c_{\mathrm{wt}}}\right\},\qquad
\gamma_{(5)}:=\frac{a_{0}}{L},\qquad
C_{R}:=\frac{18\,000\,000\,C_{G}}{1-\alpha},\\
C_{3}:=\left(1+(1+\mathrm{e})\cdot L^2\right)\cdot\max\left\{C_{R}\cdot\left(\frac{\mathrm{e}^{\frac{1}{3}}}{4}+\left(L^4\cdot b+46\right)\cdot\mathrm{e}^{1.3}\cdot C_{I}\right),\ 14L^3\cdot C_{I}\right\},\qquad
C_{5}:=\mathrm{e}^2\cdot L\cdot C_{3},
\end{gather*}
where $C_G$ is the constant of Proposition~1.5.
One has $0<\gamma_{(5)}\le\frac{1}{3L}$, and $C_5$ is a number $\ge1$ that depends on $(b,\alpha,L,\rho)$ only.
\begin{itemize}
\item[(a)] $F_{y}(0)$ is a bijection of $\mathrm{HH}(\eta_k,R;\theta)$.
On this space $F_{y}(0)^{*}=\gamma_5\,F_{-y}(0)\,\gamma_5$; on $\mathrm{HH}(\eta_k,R;-\theta)$ one has $F_{y}(0)^{T}=\mathsf{C}^{-1}F_{-y}(0)\,\mathsf{C}$; and
\[
\left(F_{y}(0)^{*}\right)^{-1}=\left(F_{y}(0)^{-1}\right)^{*},\qquad
\left(F_{y}(0)^{T}\right)^{-1}=\left(F_{y}(0)^{-1}\right)^{T}.
\]
Here $-\infty:=\infty$, as in Section~I.3.5.
\item[(b)] $\left\|F_{y}(0)^{-1}\right\|_{L^2}\le L\cdot C_{I}$.
\item[(c)] Let $z,z'\in\mathbb{Z}^3$, let $f\in\mathfrak{S}^{\mathrm{ch}}(\eta_k,R;\theta)$ be bounded and supported over the orbit of $\Delta_{z'}$, and put $v:=F_{y}(0)^{-1}f$.
Let $x$ be a point over $\Delta_z$, and let $\tilde{x}$ be a point of $\eta_k\mathbb{Z}^3\times O_{\mathrm{cub}}$ with $0<d_X(x,\tilde{x})\le1$; the pairs inside a cluster, for which $d_X=\eta_k$, are admitted.
Then
\[
|v(x)|_2+d_X(x,\tilde{x})^{-\alpha}\cdot\left|v(x)-v(\tilde{x})\right|_2\le C_{5}\cdot\exp\left(-\gamma_{(5)}\cdot|z-z'|\right)\cdot|f|_{\mathrm{su}},
\]
where $|z-z'|$ is the sup distance modulo $R\mathbb{Z}^3$.
\item[(d)] By (a), the bound of (c) holds also for the adjoint and for the transpose of $F_{y}(0)^{-1}$, the transpose being taken on the character $-\theta$.
\end{itemize}
The numbers $C_5$, $\gamma_{(5)}$ and $C_I$ do not depend on $k$, $R$, $\theta$, $\mu$ or $y$.
\end{named}

\begin{named}{Remark}
The bijectivity and the bound of \cite[Theorem~12.20]{QEDI}, in the form in which they enter \cite[Proposition~16.6]{QEDI}, are the counterpart of Proposition~1.6; the operator $T^{*}_y$ and the integer $i$ of \cite[Theorem~12.20]{QEDI} are here $F_y(0)$ and $k$.
The functions are defined on $\eta_k\mathbb{Z}^3\times O_{\mathrm{cub}}$, and the block projections are composed with $J$ and $J^{*}$.
The H\"older quotients are formed with $d_X$, and pairs inside a cluster are admitted.
The rate is $\gamma_{(5)}=a_0/L$, with the number $a_0$ defined above, whereas in \cite[Theorem~12.20]{QEDI} it is $a_0/(4L)$, with the number $a_0$ of \cite[Section~12.2]{QEDI}, and \cite[Proposition~16.6]{QEDI} records the bound $a_0/(4L)\le1/(12L^2)$.
The constants depend on $\rho$ and contain $1/c_b$.
Proposition~1.6 contains the inequality $C_5\ge1$, which \cite[Theorem~12.20]{QEDI} does not state and \cite[Proposition~16.6]{QEDI} does.
Clause (a) specifies the spaces on which the adjoint and the transpose act, and the inverses of both.
\end{named}

\subsubsection*{Proposition~1.7 (the first step: lower bound of the real part at the mesh $\ell_1/L$)}

\begin{named}{Proposition 1.7}
Let $\mathrm{me}>0$ and $\ell_1:=L\cdot\mathrm{me}$; let $R\in\ell_1\mathbb{N}\cup\{\infty\}$ and $\theta\in\mathbb{R}^3$.
Let $\mathrm{Wd}$ be a non-empty $R\mathbb{Z}^3$-periodic union of blocks $E$ of level one (Section~G1.4), and for each such block $E$ let a number $b_E\in[b_{-},b_{+}]$ be given, with $b_{E+Re_{\mu}}=b_E$ ($\mu=1,2,3$) when $R<\infty$.
Let $\mu\ge0$ and $y\in\mathbb{R}$.
For the elements $u$ of $\mathrm{HH}(\mathrm{me},R;\theta)$ that vanish outside $\mathrm{Wd}\times O_{\mathrm{cub}}$ put
\[
\mathrm{Tst}\,u:=\mathbf{1}_{\mathrm{Wd}}\,(D+\mu)\,u+\sum_{E}b_E\cdot\ell_1^{-1}\cdot P_Eu+i\,\gamma_5\,y\,u,
\]
the sum extending over the blocks $E$ of level one contained in $\mathrm{Wd}$.
Then every such $u$ satisfies
\[
\operatorname{Re}(u,\mathrm{Tst}\,u)^{\mathrm{ch}}\ge c_{\mathrm{re}}\cdot\ell_1^{-1}\cdot(u,u)^{\mathrm{ch}},\qquad
c_{\mathrm{re}}:=\min\left\{b\left(1-\frac{1}{L}\right),\frac{\rho}{2C_{\mathrm{Po}}\cdot L}\right\};
\]
the operator $\mathrm{Tst}$ is a bijection of the space of these $u$, and $\left\|\mathrm{Tst}^{-1}\right\|_{L^2}\le\ell_1/c_{\mathrm{re}}$.

For $\mathrm{me}=1$, $\mathrm{Wd}$ the whole lattice and $b_E=b$ for every $E$, the operator $\mathrm{Tst}$ is the operator $F_y(0)$ of Section~I.3.9, because in this case $b\cdot\ell_1^{-1}\sum_{E}P_E=b_{L}\,J P(0) J^{*}$.
\end{named}

\begin{named}{Remark}
The corresponding statement is \cite[Proposition~12.21]{QEDI}, which is applied in the proof of \cite[Lemma~16.7]{QEDI}; the constant $c_0$ of \cite[Proposition~12.21]{QEDI} is called $c_{\mathrm{re}}$ in \cite[Lemma~16.7]{QEDI}.
In this constant the number $4/L$ is replaced by $\rho/(2C_{\mathrm{Po}}\cdot L)$.
A block of level one carries $192L^3$ components, and on each block the penalty $b_E\cdot\ell_1^{-1}\cdot P_E$ has rank $4$.
The matrix $\hat{\Gamma}$ of \cite[Proposition~12.21]{QEDI} is here $\gamma_5$.
\end{named}

\subsection*{2. Small fields: the inverse of $D_A+\mu+\operatorname{Pen}_{y}(A)$}

The constants that enter the Propositions~2.m are among those of (C4) and (C4r) in Section~I.8.2;
their defining formulas are those of (C4) and (C4r) in Proposition~14.2\,(a).
They are
\begin{gather*}
C'_{V}:=2\mathrm{e}\cdot(42+6b),\qquad c_{\mathrm{re}}\ \text{(the constant of Proposition~1.7)},\\
a_{z}:=\min\left\{\frac{1}{3L},\ \frac{c_{\mathrm{re}}}{2\mathrm{e}\cdot L\cdot(6+3b)}\right\},\qquad k_{1}:=60+b,\qquad C_{00}:=48\cdot\frac{2L}{c_{\mathrm{re}}}\cdot\left(1+2k_{1}\right),\\
C_{z}:=\max\left(C_{5},C_{00}\right),\qquad\gamma_{z}:=\min\left(\gamma_{(5)},a_{z}\right),\qquad\gamma^{F}:=\frac{\gamma_{z}}{2},\qquad a':=C_{z}\cdot c(\gamma^{F}).
\end{gather*}
They satisfy
\[
\gamma^{F}L\le\frac{1}{6}<\log2,\qquad a'\ge1,\qquad C'_{V}\ge24\mathrm{e}.
\]
The Propositions~2.m use the following conditions, each of which is a row of Proposition~14.2\,(b).
The condition (L5) concerns the constant $L$ of line (C3) of Section~I.8.2;
$L$ is chosen before the constants of (C4) and (C4r), which contain $L$.
The conditions (M1), (M2) and the inequality $M\ge108\pi\cdot C'_{V}\cdot a'$, which is contained in (M5), concern the constant $M$ of line (C5);
the conditions (Z2.s1) and (Z2.s2) concern the constant $e_0$ of line (C8), and they contain the exponent $\varepsilon$ of line (C7).
Of the constants so constrained, only $M$ and $e_0$ are chosen after (C4r).
\begin{itemize}
\item[(L5)] $L$ is odd and $L\ge5$.
\item[(M1)] $M$ is a power of $L$, and $M\ge L^{2}$.
\item[(M2)] $M>4L$.
\item[] $M\ge108\pi\cdot C'_{V}\cdot a'$.
\item[(Z2.s1)] $6L\cdot e_{0}^{\varepsilon}\le1$.
\item[(Z2.s2)] $4C'_{V}\cdot a'\cdot e_{0}^{\varepsilon}\le1$.
\end{itemize}

\begin{named}{Proposition 2.1}[Perturbations, commutators, H\"older quotient of a product]
Let $\mathbb{T}$ be a block torus of level $k\ge0$, let $e_k>0$, $\mu\ge0$ and $y\in Y_{k}$, let $\gamma$ be a number with $0\le\gamma L\le\log2$, and let $A$ and $B$ be connections on $\mathbb{T}$.
\begin{itemize}
\item[(a)] If $e_k|B|_{\infty}\le\delta$ and $3L\cdot\delta\le1$, then
\[
N_{\gamma}\left(F_{y}(B)-F_{y}(0)\right)\le C'_{V}\cdot\delta.
\]
\item[(b)] Let $3L\cdot\mathfrak{i}(A)\le1$, and let $\psi$ be a real scalar such that $\left|\psi(c)-\psi(c')\right|\le\mathrm{Lip}\cdot|c-c'|_{1}$ for all $c$, $c'$ with $|c-c'|\le L$.
Then
\[
N_{\gamma}\left(\left[F_{y}(A),\psi\right]\right)\le C'_{V}\cdot\mathrm{Lip}.
\]
The kernel of $\left[F_{y}(A),\psi\right]$ vanishes at $(x,x')$ unless $\psi(c)\neq\psi(c')$ and, in addition, either $x'=x^{(3)}$ or $c$ and $c'$ lie in one block of side $L$;
this is meant in the sense of Section~G1.2, that is, modulo $R\mathbb{Z}^{3}$.
\item[(c)] For a spinor $f$, a scalar $\psi$ and a short pair $(x,x')$ one has
\begin{align*}
\left(\delta_{\alpha,A}(\psi f)\right)(x,x')&=\psi(c)\left(\delta_{\alpha,A}f\right)(x,x')+d_X(x,x')^{-\alpha}\left(\psi(c')-\psi(c)\right)U_A(x,x')\,f(x'),\\
\left(\delta_{\alpha,B}f-\delta_{\alpha,0}f\right)(x,x')&=d_X(x,x')^{-\alpha}\left(U_B(x,x')-1\right)f(x').
\end{align*}
At pairs inside a cluster the second term on the right of the first identity vanishes, and so does the right side of the second identity.
Further, let $O$ be an operator with values in the spinors on $\mathbb{T}$ and with $N_{\gamma}(O)<\infty$;
let $\mathrm{wq}$ be a function on the short pairs, with complex values or with values in the $4\times4$ matrices, such that $\mathrm{wq}(x+Re_{\mu},x'+Re_{\mu})=\mathrm{wq}(x,x')$ and $\left|\mathrm{wq}(x,x')\right|\le\mathrm{wq}_{0}$, where in the matrix case the bound reads $\left|\mathrm{wq}(x,x')\right|_{2}\le\mathrm{wq}_{0}$;
and let $E_{\mathrm{wq}}O$ be defined by $\left(E_{\mathrm{wq}}Of\right)(x,x'):=\mathrm{wq}(x,x')\,(Of)(x')$.
Then
\[
N_{\gamma}\left(E_{\mathrm{wq}}O\right)\le27\exp(\gamma)\cdot\mathrm{wq}_{0}\cdot N_{\gamma}(O).
\]
\end{itemize}
\end{named}

\begin{named}{Remark}
Proposition~2.1 corresponds to \cite[Lemma~16.4]{QEDI}.
The constant of that lemma is $C_V=12\mathrm{e}\cdot(2+b)$;
here it is $C'_{V}$, because the operator $D_A$ has one hop of norm at most $6$, which reaches the unit block of the site and the six adjacent unit blocks, in place of six hops of norm one.
The function $h$, its constant $L_h$ and the function $g$ of \cite[Lemma~16.4]{QEDI} are denoted here by $\psi$, $\mathrm{Lip}$ and $\mathrm{wq}$, and $d_X(x,x')$ takes the place of $|x-x'|$.
In (b) the alternative that $x$ and $x'$ be nearest neighbours becomes $x'=x^{(3)}$.
The assertion on the kernel is read modulo $R\mathbb{Z}^{3}$, as in \cite[Lemma~16.4(b)]{QEDI}, by the convention of \cite[Section~16.2.3]{QEDI} on kernels supported in a set of pairs;
this is the reading fixed in Section~G1.2.
At $k=0$ the penalty is $b_{L}\cdot JP(A)J^{*}+i\gamma_5y$.
\end{named}

\begin{named}{Proposition 2.2}[Algebra, rows and covariances of the elementary operators on block tori]
Let $\mathbb{T}$ and $\mathbb{T}'$ be block tori of level $k\ge0$, let $e_k>0$, $\mu\ge0$ and $y\in Y_{k}$, and let $A$ be a complex connection on $\mathbb{T}$.
At $k=0$ the operator $JP_k(A)J^{*}$ is read as $P_0=1$, and $JP_{k+1}(A)J^{*}$ as $JP(A)J^{*}$ (see (I.3.4) and Section~I.3.9);
the relations that contain $Q_k(A)J^{*}$ or $JQ_k^{T}(-A)$, and those that contain $Q(A)$, $Q^{T}(-A)$ or $P(A)$ without $J$, $J^{*}$, are asserted for $k\ge1$.
In (a) and (b), $(x,x')$ is an arbitrary ordered pair of distinct points of $\mathrm{Lab}\times O_{\mathrm{cub}}$ (Section~G1.2).
\begin{itemize}
\item[(f0)] \emph{Algebra and rows.}
For $k\ge1$,
\[
\left(Q_k(A)J^{*}\right)\left(JQ_k^{T}(-A)\right)=1,\qquad Q(A)\,Q^{T}(-A)=1,
\]
and for $k=0$, $\left(Q(A)J^{*}\right)\left(JQ^{T}(-A)\right)=1$.
The operators $JP_k(A)J^{*}$, $P(A)$ and $JP_{k+1}(A)J^{*}$ (at $k=0$: $JP(A)J^{*}$) are idempotent, and
\[
\left(JP_k(A)J^{*}\right)\left(JP_{k+1}(A)J^{*}\right)=\left(JP_{k+1}(A)J^{*}\right)\left(JP_k(A)J^{*}\right)=JP_{k+1}(A)J^{*};
\]
their kernels are multiples of $\mathbf{1}_{4}$.
These idempotents and the six operators $Q_k(A)J^{*}$, $JQ_k^{T}(-A)$, $Q(A)$, $Q^{T}(-A)$, $Q(A)J^{*}$, $JQ^{T}(-A)$ commute with every constant $4\times4$ matrix.
Moreover $\operatorname{Tr}P(A)=4\left(\frac{R}{L}\right)^{3}$ for $k\ge1$, and $\operatorname{Tr}JP(A)J^{*}=4\left(\frac{R}{L}\right)^{3}$ for $k=0$.
The operator $JP_k(A)J^{*}$ couples only sites over one unit block;
the operators $P(A)$, $JP_{k+1}(A)J^{*}$ and $JP(A)J^{*}$ couple only sites in (over) one block of side $L$;
the rows at $x$ read $A$ on bonds inside the block of $c$.
The row of $D_A$ at $x$ reads $f$ at $x$, $x^{(0)}$, $x^{(1)}$, $x^{(2)}$, $x^{(3)}$, and $A$ on the bond $(c,c_{3})$ (see (I.O1)).
\item[(a)] \emph{Gauge.}
Let $\mathrm{Acl}$ be a real connection on $\mathbb{T}$ with $d\mathrm{Acl}=0$, and let $\mathrm{uf}$ be an $R\mathbb{Z}^{3}$-periodic function on $\mathrm{Lab}$ with $|\mathrm{uf}|=1$ and
\[
\mathrm{uf}(c')=\mathrm{uf}(c)\cdot\exp\left(ie_k\eta_k\cdot\mathrm{Acl}(c,c')\right)
\]
on every bond $(c,c')$;
let $\mathrm{uf}^{0}$ and $\mathrm{uf}^{1}$ be its restrictions to $\mathrm{Lab}^{0}$ and to $\mathrm{Lab}_{k+1}$.
Then
\begin{gather*}
D_{A-\mathrm{Acl}}=\mathrm{uf}\,D_A\,\mathrm{uf}^{-1},\qquad Q_k(A-\mathrm{Acl})J^{*}=\mathrm{uf}^{0}\,Q_k(A)J^{*}\,\mathrm{uf}^{-1},\\
JQ_k^{T}\left(-(A-\mathrm{Acl})\right)=\mathrm{uf}\,JQ_k^{T}(-A)\left(\mathrm{uf}^{0}\right)^{-1},\qquad Q(A-\mathrm{Acl})=\mathrm{uf}^{1}\,Q(A)\left(\mathrm{uf}^{0}\right)^{-1},\\
Q^{T}\left(-(A-\mathrm{Acl})\right)=\mathrm{uf}^{0}\,Q^{T}(-A)\left(\mathrm{uf}^{1}\right)^{-1},\\
O(A-\mathrm{Acl})=\mathrm{uf}\,O(A)\,\mathrm{uf}^{-1}\quad\text{for }O=JP_kJ^{*},\ JP_{k+1}J^{*},\ \operatorname{Pen}_{y},\ F_{y},\\
P(A-\mathrm{Acl})=\mathrm{uf}^{0}\,P(A)\left(\mathrm{uf}^{0}\right)^{-1},\\
\left(\delta_{\alpha,A-\mathrm{Acl}}(\mathrm{uf}\,f)\right)(x,x')=\mathrm{uf}(c)\left(\delta_{\alpha,A}f\right)(x,x');
\end{gather*}
at $k=0$,
\begin{gather*}
Q(A-\mathrm{Acl})J^{*}=\mathrm{uf}^{1}\,Q(A)J^{*}\,\mathrm{uf}^{-1},\qquad JQ^{T}\left(-(A-\mathrm{Acl})\right)=\mathrm{uf}\,JQ^{T}(-A)\left(\mathrm{uf}^{1}\right)^{-1},\\
JP(A-\mathrm{Acl})J^{*}=\mathrm{uf}\,JP(A)J^{*}\,\mathrm{uf}^{-1}.
\end{gather*}
For example, $\mathrm{Acl}=d\boldsymbol{\lambda}$ and $\mathrm{uf}=\exp\left(ie_k\boldsymbol{\lambda}\right)$, with a real scalar $\boldsymbol{\lambda}$, satisfy these requirements.
\item[(b)] \emph{Morphisms.}
Let $\mathrm{mph}\colon\mathbb{T}\to\mathbb{T}'$ be a morphism with the linear part $r$, let $S\in\mathrm{Sp}$ with $r(S)=r$, and put $\mathrm{mS}:=(S,\mathrm{mph})$.
Then
\begin{gather*}
D_{\mathrm{mS}A}\,\mathrm{mS}=\mathrm{mS}\,D_A,\\
Q_k(\mathrm{mS}A)J^{*}\,\mathrm{mS}=\mathrm{mS}\,Q_k(A)J^{*},\qquad JQ_k^{T}(-\mathrm{mS}A)\,\mathrm{mS}=\mathrm{mS}\,JQ_k^{T}(-A),
\end{gather*}
the corresponding relations hold for $Q$ and $Q^{T}$ (at $k=0$: for $QJ^{*}$ and $JQ^{T}$), for $JP_kJ^{*}$, $P$ and $JP_{k+1}J^{*}$, and
\[
\left(\delta_{\alpha,\mathrm{mS}A}\,\mathrm{mS}f\right)\left(\mathrm{mph}\,x,\mathrm{mph}\,x'\right)=S\left(\delta_{\alpha,A}f\right)(x,x'),\qquad F_{y'}(\mathrm{mS}A)\,\mathrm{mS}=\mathrm{mS}\,F_{y}(A),
\]
where $y'=y$ if $S\in\mathrm{Sp}_3$ or if $y$ is neutral (Section~G1.9), and $y'=-y$ otherwise.
\item[(c)] \emph{Transposition.}
One has
\[
D_A^{T}=\mathsf{C}^{-1}D_{-A}\,\mathsf{C},\qquad\left(Q_k(A)J^{*}\right)^{T}=JQ_k^{T}(A),\qquad Q(A)^{T}=Q^{T}(A),
\]
and at $k=0$ one has $\left(Q(A)J^{*}\right)^{T}=JQ^{T}(A)$;
further $\left(JP_k(A)J^{*}\right)^{T}=JP_k(-A)J^{*}$, and the same relation holds for $P$, $JP_{k+1}J^{*}$ and $JPJ^{*}$;
finally
\[
F_y(A)^{T}=\mathsf{C}^{-1}F_{-y}(-A)\,\mathsf{C}.
\]
\item[(d)] \emph{Conjugation.}
One has
\[
\overline{O\left(\overline{A}\right)}=\varkappa^{-1}O(-A)\,\varkappa
\]
for $O(A)=D_A$, $Q_k(A)J^{*}$, $JQ_k^{T}(-A)$, $Q(A)$, $Q^{T}(-A)$ (at $k=0$: $Q(A)J^{*}$, $JQ^{T}(-A)$), $JP_k(A)J^{*}$, $P(A)$, $JP_{k+1}(A)J^{*}$, $JP(A)J^{*}$, and for $O(A)=F_y(A)$ with the same $y$ on both sides.
\item[(e)] \emph{Adjoints.}
For real $A$,
\[
D_A^{*}=\gamma_5D_A\gamma_5,\qquad\left(Q_k(A)J^{*}\right)^{*}=JQ_k^{T}(-A),\qquad Q(A)^{*}=Q^{T}(-A),
\]
and at $k=0$, $\left(Q(A)J^{*}\right)^{*}=JQ^{T}(-A)$;
the operators $JP_k(A)J^{*}$, $P(A)$, $JP_{k+1}(A)J^{*}$ and $JP(A)J^{*}$ are orthogonal projections;
and
\[
F_y(A)^{*}=\gamma_5F_{-y}(A)\gamma_5.
\]
\end{itemize}
\end{named}

\begin{named}{Remark}
Proposition~2.2 corresponds to \cite[Lemma~16.3]{QEDI} and \cite[Lemma~16.5]{QEDI}.
Here the cluster maps $J$, $J^{*}$ are written out, a morphism acts also on the label, $\mathrm{mph}(c,g)=\left(\mathrm{mph}\,c,rg\right)$, and $\gamma_5$ takes the place of the matrix $\gamma_3$ of \cite[Section~16.2.1]{QEDI}.
The hopping matrices of $D_A$ depend on the site;
the relations for $D_A$ rest on the identities
\begin{gather*}
M_{i}(\mathrm{mph}\,x)\,S=S\,M_{i}(x),\qquad\mathsf{C}^{-1}M_{i}(x)^{T}\mathsf{C}=M_{i}\left(x^{(i)}\right),\\
\overline{M_{i}(x)}=\varkappa^{-1}M_{i}(x)\,\varkappa,\qquad\gamma_5M_{i}(x)\gamma_5=M_{i}\left(x^{(i)}\right),
\end{gather*}
and, for the diagonal term, on $\sum_{i}M_{i}(x)=\frac{\rho}{2}\cdot\mathrm{sw}$.
The last relation of (e) has no counterpart in \cite[Lemma~16.5]{QEDI}.
\end{named}

\begin{named}{Proposition 2.3}[Zero field on every block torus]
Let $\mathbb{T}=(t,R,\boldsymbol{\sigma})$ be a block torus of level $k\ge0$, let $\mu\in[0,1]$ and $y\in Y_{k}$.
\begin{itemize}
\item[(a)] ($k\ge1$.)
The operator $F_y(0)$ is bijective on the spinors on $\mathbb{T}$.
Let $z_1,z_2\in\mathrm{Lab}^{0}$, let $f$ be a spinor supported over the orbit of $\Delta_{z_2}$, and put $v:=S_y(0)f$.
Then for every $x$ over $\Delta_{z_1}$ and every $\tilde{x}$ with $0<d_X(x,\tilde{x})\le1$,
\[
\left|v(x)\right|_{2}+d_X(x,\tilde{x})^{-\alpha}\left|v(x)-v(\tilde{x})\right|_{2}\le C_{5}\cdot\exp\left(-\gamma_{(5)}\cdot d(z_1,z_2)\right)\cdot|f|_{\mathrm{su}}.
\]
\item[(b)] ($k=0$.)
The operator
\[
F_y(0)=D_0+\mu+b_{L}\cdot JP(0)J^{*}+i\gamma_5y
\]
is bijective, and the kernel of $S_y(0)$, in the sense of Section~G1.2, satisfies
\[
\left|S_y(0)(x,x')\right|_{2}\le C_{00}\cdot\left(1+|y|\right)^{-1}\cdot\exp\left(-a_{z}\cdot d(x,x')\right).
\]
\item[(c)] For every $k\ge0$,
\[
N_{\gamma^{F}}\left(S_y(0)\right)\le a'\cdot w_{k}(y),\qquad N_{\gamma^{F}}\left(\delta_{\alpha,0}S_y(0)\right)\le a',
\]
where the H\"older quotient is taken at the short pairs (at $k=0$ there is no short pair, and the second bound is void);
at $k=0$, moreover,
\[
N^{\mathrm{ke}}_{\gamma^{F}}\left(S_y(0)\right)\le a'\cdot w_{0}(y).
\]
\end{itemize}
\end{named}

\begin{named}{Remark}
Proposition~2.3 corresponds to \cite[Proposition~16.6]{QEDI}, to \cite[Lemma~16.7]{QEDI}, and to the two displays of \cite[Section~16.4]{QEDI} that define $C_z$, $\gamma$, $a$ and bound $N_{\gamma}\left(S_y(0)\right)$ and $N_{\gamma}\left(\delta_{\alpha,0}S_y(0)\right)$.
The constants $C_5$ and $\gamma_{(5)}$ are those of Proposition~1.6, and $\gamma_{(5)}$ lies in $\left(0,\frac{1}{3L}\right]$ instead of $\left(0,\frac{1}{12L^{2}}\right]$;
the constant $c_{\mathrm{re}}$ is that of Proposition~1.7.
The number $k_1=14+b$ of \cite[Lemma~16.7]{QEDI} becomes $60+b$, since $\left\|D_0\right\|_{L^{2}}\le58$ by (I.O2).
The constant $C_{00}$ contains the factor $48$ of the kernel (I.1.4):
by (G1.Nk5) of Section~G1.3, the step from the $L^{2}$ operator norm to the kernel costs the factor $\sqrt{48}$ twice.
The number $6+3b$ in $a_z$ is the same as in \cite[Lemma~16.7]{QEDI} (one hop of norm at most $6a_z\mathrm{e}$).
The third bound of (c), a bound in the kernel norm of (I.ke5), has no counterpart in \cite{QEDI}.
\end{named}

\begin{named}{Proposition 2.4}[Small fields on a block torus]
Let $\mathbb{T}$ be a block torus of level $k\ge0$, let $e_k>0$, $\mu\in[0,1]$ and $y\in Y_{k}$.
Let $\delta>0$ be a number with $3L\cdot\delta\le1$ and $2C'_{V}\cdot a'\cdot\delta\le1$, let $B$ be a connection on $\mathbb{T}$ with $e_k|B|_{\infty}\le\delta$, and put $V:=F_y(B)-F_y(0)$.
\begin{itemize}
\item[(a)] $F_y(B)$ is bijective, and
\[
S_y(B)=S_y(0)\sum_{n\ge0}\left(-V\,S_y(0)\right)^{n}=\sum_{n\ge0}\left(-S_y(0)\,V\right)^{n}S_y(0),
\]
both series converging in $N_{\gamma^{F}}$, uniformly in $B$.
\item[(b)] $N_{\gamma^{F}}\left(S_y(B)\right)\le2a'\cdot w_{k}(y)$.
\item[(c)] $N_{\gamma^{F}}\left(\delta_{\alpha,B}S_y(B)\right)\le18a'$, the H\"older quotient being taken at the short pairs.
\item[(d)] ($k=0$.) $N^{\mathrm{ke}}_{\gamma^{F}}\left(S_y(B)\right)\le2a'\cdot w_{0}(y)$.
\end{itemize}
\end{named}

\begin{named}{Remark}
Proposition~2.4 corresponds to \cite[Proposition~16.8]{QEDI}, with the constants $C'_{V}$, $a'$ and $\gamma^{F}$.
The second form of the series in (a) and the clause (d), which concerns the kernel norm of (I.ke5), have no counterpart there.
The H\"older quotients are taken at the short pairs, formed with $d_X$;
inside a cluster $U_B=1$.
\end{named}

\begin{named}{Proposition 2.5}[Local inverses on mirror tori]
Let the data (SF) be given, except that $e_k>0$ is arbitrary (the bound $e_k\le e_0$ is not required).
Let $\square$ be an $M$-cube, with the objects of Section~G1.5, let $\delta$ be a number as in Proposition~2.4, let $y\in Y_{k}$, and let $A\in\mathrm{Gloc}(\delta)$.
\begin{itemize}
\item[(a)] $F^{W}_{y}(A)$ is bijective.
The map $\mathrm{Sloc}_{\square,y}(\cdot)$ is analytic on $\mathrm{Gloc}(\delta)$, and it is continuous in $(y,A)$ for $y\in\mathbb{R}$.
\item[(b)] Let $N_{\gamma}$ be formed with the unit sites of $\tilde{W}$ and the distance $|z_1-z_2|$.
Then
\[
N_{\gamma^{F}}\left(\mathrm{Sloc}_{\square,y}(A)\right)\le2a'\cdot w_{k}(y),\qquad N_{\gamma^{F}}\left(\delta_{\alpha,A}\mathrm{Sloc}_{\square,y}(A)\right)\le18a',
\]
the H\"older quotient being taken at the short pairs with both points over $\tilde{W}$;
at $k=0$, moreover,
\[
N^{\mathrm{ke}}_{\gamma^{F}}\left(\mathrm{Sloc}_{\square,y}(A)\right)\le2a'\cdot w_{0}(y).
\]
\item[(c)] Let $F_y(A)$ be formed on $\mathbb{L}^{\mathrm{f}}_k$ with an arbitrary continuation of $A$, by the display (I.2.3) and the displays of Sections~I.3.3, I.3.5 and I.3.9, applied to arbitrary functions on $\mathbb{L}^{\mathrm{f}}_k$;
these are finite sums.
The rows of $F_y(A)$ at the sites $x$ over $W^{-}$ read the function only over $\tilde{W}$, and $A$ only on the bonds of $\tilde{W}$.
For every function $f$ on $\tilde{W}\times O_{\mathrm{cub}}$, with $v:=\mathrm{Sloc}_{\square,y}(A)f$ continued by zero, one has $\left(F_y(A)v\right)(x)=f(x)$ for all $x$ over $W^{-}$.
For every function $\mathrm{fs}$ supported over $W^{-}$, the function $F_y(A)\mathrm{fs}$ is supported over $\tilde{W}$, and
\[
\mathrm{Sloc}_{\square,y}(A)\left(F_y(A)\mathrm{fs}\right)=\mathrm{fs}.
\]
\item[(d)] For every real function $\boldsymbol{\lambda}$ on $\tilde{W}$,
\[
\mathrm{Sloc}_{\square,y}(A-d\boldsymbol{\lambda})=\exp\left(ie_k\boldsymbol{\lambda}\right)\mathrm{Sloc}_{\square,y}(A)\exp\left(-ie_k\boldsymbol{\lambda}\right).
\]
\item[(e)] Let $\tilde{g}=(S,a)\in\mathbb{G}_k$, and let $y'=y$ if $S\in\mathrm{Sp}_3$ or if $y$ is neutral, and $y'=-y$ otherwise.
Then
\[
\mathrm{Sloc}_{\tilde{g}\square,y'}(\tilde{g}A)\,\tilde{g}=\tilde{g}\,\mathrm{Sloc}_{\square,y}(A);
\]
here $\tilde{g}$ maps functions over $\tilde{W}$ to functions over the lift $r(S)\tilde{W}+a$ of $\left(\tilde{g}\square\right)^{\sim1}$, with the action (I.4.2) on the sites.
\item[(f)] With the transpose taken for the pairing $\int_{\tilde{W}}\mathrm{d}x\,f'(x)\cdot f(x)$,
\[
\mathrm{Sloc}_{\square,y}(A)^{T}=\mathsf{C}^{-1}\mathrm{Sloc}_{\square,-y}(-A)\,\mathsf{C},\qquad\overline{\mathrm{Sloc}_{\square,y}\left(\overline{A}\right)}=\varkappa^{-1}\mathrm{Sloc}_{\square,y}(-A)\,\varkappa.
\]
\item[(g)] Let $A$ be the restriction to the bonds of $\tilde{W}$ of a complex antisymmetric function $\tilde{A}$ on the bonds of $\mathbb{L}_k$ that is invariant under the translations by $\ell\mathbb{Z}^{3}$, and let $n\in\mathbb{Z}^{3}$.
The point $z_{\square}+\ell n$ lies in $M\mathbb{Z}^{3}$, since $\ell$ is a multiple of $LM$;
let the objects of Section~G1.5 be formed with the lift $z_{\square}+\ell n$ in place of $z_{\square}$, so that $\tilde{W}+\ell n$, $W^{-}+\ell n$ and $o+\ell n$ take the places of $\tilde{W}$, $W^{-}$ and $o$,
and let $A_n$ be the restriction of $\tilde{A}$ to the bonds of $\tilde{W}+\ell n$.
Then $A_n(c+\ell n,c'+\ell n)=A(c,c')$, so that $A_n$ lies in the set $\mathrm{Gloc}(\delta)$ formed with the lift $z_{\square}+\ell n$;
the folded connection of $A_n$ formed with this lift is $A^{W}$, so that $F^{W}_{y}$ does not change;
and the local inverse formed with the lift $z_{\square}+\ell n$ at $A_n$ is the translate
\[
T_{\ell n}\,\mathrm{Sloc}_{\square,y}(A)\,T_{-\ell n}
\]
of $\mathrm{Sloc}_{\square,y}(A)$, which sends a function $f$ on $\left(\tilde{W}+\ell n\right)\times O_{\mathrm{cub}}$ to the function $x\mapsto\left(\mathrm{Sloc}_{\square,y}(A)\,f(\cdot+\ell n)\right)(x-\ell n)$.
\end{itemize}
\end{named}

\begin{named}{Remark}
Proposition~2.5 corresponds to \cite[Definition~16.11]{QEDI} (the maps $\mathrm{ext}$ and $\mathrm{rs}$, the operator $F^{W}_{y}$ and the local inverse), to the relations $\mathrm{rs}\cdot\mathrm{ext}=1$ and $\mathrm{ext}^{T}=\mathrm{rs}$ in the paragraph that follows that definition, and to \cite[Proposition~16.12]{QEDI}.
Here the local inverse acts on functions on $\tilde{W}\times O_{\mathrm{cub}}$, and $\mathrm{ext}$, $\mathrm{rs}$ leave the label unchanged.
In (c) the row of $D_A$ at $x=(c,g)$ reads $x$, $x^{(0)}$, $x^{(1)}$, $x^{(2)}$ over $c$ and $x^{(3)}$ over the nearest neighbour $c_3$, and $A$ on the bond $(c,c_3)$.
In the proof of (e) (Chapter~2b, Section~1) the element $\tilde{g}=(S,a)$ induces on the mirror torus the map
\[
\mathrm{mph}^{W}(u):=r\left(u-\frac{3M}{2}(1,1,1)\right)+\frac{3M}{2}(1,1,1),\qquad r:=r(S),
\]
a morphism of $\mathbb{T}^{W}$ onto itself with the linear part $r$ that maps $\mathrm{Ch}$ onto $\mathrm{Ch}$;
for $c\in\mathbb{R}^{3}$ one has $\mathrm{mph}^{W}(c-o)=rc+a-o'$, where $o':=rz_{\square}+a-\frac{3M}{2}(1,1,1)$ is the point $o$ formed with the lift $rz_{\square}+a$ of the centre of $\tilde{g}\square$.
On the fermion sites of $\mathbb{T}^{W}$ this morphism acts also on the label, $(u,g)\mapsto\left(\mathrm{mph}^{W}u,rg\right)$, as every morphism does by Section~G1.2.
The constants are $a'$ and $\gamma^{F}$;
the kernel clause in (b), which concerns the kernel norm of (I.ke5), has no counterpart in \cite[Proposition~16.12]{QEDI}.
\end{named}

The constants in Proposition~2.5 are subject to the conditions (L5) and (M1).

\begin{named}{Proposition 2.6}[Operators over $\tilde{W}$ as operators on $T_k$: four compatibilities]
Let the data (SF) be given, let $\square$ be an $M$-cube, and let $O$ and $O'$ be operators as in Section~G1.6.
Of the condition $\ell\in\mathrm{Cl}$ only the inequality $3M<\ell$ enters, and, for the equalities of norms in (a), the inequality $3M-1\le\frac{\ell}{2}$.
\begin{itemize}
\item[(a)] $E_{\square}(O)$ is a linear operator on $\mathfrak{S}^{\mathrm{ch}}(\eta_k,\ell)$, with values in $\mathfrak{S}^{\mathrm{ch}}(\eta_k,\ell)$ or in the functions of short pairs, respectively.
If $\tilde{W}$ is replaced by $\tilde{W}+\ell n$ and $O$ by its translate, then $E_{\square}(O)$ does not change.
The map $O\mapsto E_{\square}(O)$ is linear;
$E_{\square}(OO')=E_{\square}(O)\,E_{\square}(O')$ for operators with values in the functions on $\tilde{W}\times O_{\mathrm{cub}}$;
and $N_{\gamma}\left(E_{\square}(O)\right)=N_{\gamma}(O)$, the right side being formed over $\tilde{W}$ as in Proposition~2.5\,(b);
at $k=0$ the same holds for $N^{\mathrm{ke}}_{\gamma}$.
\item[(b)] For an $\ell\mathbb{Z}^{3}$-periodic function $\mathrm{uf}$ of $c$ without zeros and a constant invertible $4\times4$ matrix $\mathrm{Mx}$,
\[
E_{\square}\left(\mathrm{uf}\,O\,\mathrm{uf}^{-1}\right)=\mathrm{uf}\,E_{\square}(O)\,\mathrm{uf}^{-1},\qquad E_{\square}\left(\mathrm{Mx}\,O\,\mathrm{Mx}^{-1}\right)=\mathrm{Mx}\,E_{\square}(O)\,\mathrm{Mx}^{-1};
\]
the restriction of $A-d\boldsymbol{\lambda}$ to $\tilde{W}$ is the restriction of $A$ minus $d$ of the restriction of $\boldsymbol{\lambda}$.
\item[(c)] For $\tilde{g}\in\mathbb{G}_k$,
\[
\tilde{g}\,E_{\square}(O)\,\tilde{g}^{-1}=E_{\tilde{g}\square}\left(\tilde{g}\,O\,\tilde{g}^{-1}\right),
\]
where $E_{\tilde{g}\square}$ is the map $E_{\square}$ of Section~G1.6 for the cube $\tilde{g}\square$, formed with $\tilde{g}\tilde{W}$;
the restriction of $\tilde{g}f$ to $\tilde{g}\tilde{W}$ is $\tilde{g}$ applied to the restriction of $f$ to $\tilde{W}$, and the restriction of $\tilde{g}A$ to $\tilde{g}\tilde{W}$ is $\tilde{g}$ applied to the restriction of $A$.
\item[(d)] For $O$ with values in the functions on $\tilde{W}\times O_{\mathrm{cub}}$,
\[
E_{\square}(O)^{T}=E_{\square}\left(O^{T}\right),
\]
where the transpose on the left is taken for $\langle\cdot,\cdot\rangle^{\mathrm{ch}}$ on $T_k$, and the one on the right for $\int_{\tilde{W}}\mathrm{d}x\,f'(x)\cdot f(x)$.
\item[(e)] $\overline{E_{\square}(O)}=E_{\square}\left(\overline{O}\right)$.
\end{itemize}
\end{named}

\begin{named}{Remark}
Proposition~2.6 corresponds to \cite[Section~16.6.2]{QEDI}.
The clauses (b)--(e) correspond to the four identities proved in the second to fourth paragraphs of that section:
(b) to the conjugation by an $\ell\mathbb{Z}^{3}$-periodic function without zeros and by a constant invertible $4\times4$ matrix,
(c) to the identity for an element of $\mathbb{G}_k$,
(d) to the transposition,
and (e) to the complex conjugation.
The assertions of (b) and (c) on the restrictions of $A-d\boldsymbol{\lambda}$, of $\tilde{g}f$ and of $\tilde{g}A$ correspond to the statements on restrictions made there, in the proof of the identity for an element of $\mathbb{G}_k$ (for $\tilde{g}f$) and in the last paragraph of that section (for $A-d\boldsymbol{\lambda}$ and $\tilde{g}A$).
The clauses (b)--(e) are applied in Chapter~2b, in the proof of Corollary~6.A, and in Chapter~3, in the proofs of Proposition~3\,(e)--(h) (gauge, lattice symmetries, transposition and conjugation), which correspond to \cite[Proposition~16.18(e)--(h)]{QEDI}.
Here the functions live on $\tilde{W}\times O_{\mathrm{cub}}$, and $\tilde{g}$ acts also on the label;
the identities of \cite[Section~16.6.2]{QEDI} are stated for operators with values in the functions on $\tilde{W}$, whereas (b), (c) and (e) hold here also for operators with values in the functions of short pairs.
The pairing $\langle\cdot,\cdot\rangle^{\mathrm{ch}}$ carries the weight $\frac{1}{48}$.
\end{named}

The constants in Proposition~2.6 are subject to the conditions (L5) and (M1).

\begin{named}{Proposition 2.7}[Parametrix; sums over walks]
Let the data (SF) be given, let $A\in\tilde{G}_k$ and $y\in Y_{k}$, and let $F:=F_y(A)$ on $T_k$;
in this statement $F$ denotes this operator, not the curvature.
For $M$-cubes $\square$ and $\square'$ put $\mathrm{Sloc}_{\square}:=\mathrm{Sloc}_{\square,y}(A)$ and $\mathrm{Cm}_{\square}:=\mathrm{Cm}_{\square,y}(A)$ (Sections~G1.6 and G1.7).
\begin{itemize}
\item[(a)] $F\,h_{\square}\mathrm{Sloc}_{\square}h_{\square}=h_{\square}^{2}-\mathrm{Cm}_{\square}\mathrm{Sloc}_{\square}h_{\square}$.
\item[(b)] $h_{\square}\mathrm{Sloc}_{\square}h_{\square}\,F=h_{\square}^{2}+h_{\square}\mathrm{Sloc}_{\square}\mathrm{Cm}_{\square}$.
\item[(c)] $\mathrm{Cm}_{\square}(x,x')=\left(h_{\square}(c)-h_{\square}(c')\right)F(x,x')$.
This kernel vanishes unless $|c-z_{\square}|<\frac{3M}{4}+L$ and $|c'-z_{\square}|<\frac{3M}{4}+L$ (for the lift), and these sites lie in $W^{-}$.
The operator $\mathrm{Cm}_{\square}$ reads $A$ only on the bonds of $\tilde{W}$.
Moreover $h_{\square'}\mathrm{Cm}_{\square}=0=\mathrm{Cm}_{\square}h_{\square'}$ unless $d(z_{\square},z_{\square'})\le M$.
\item[(d)] One has
\begin{gather*}
N_{\gamma^{F}}\left(\mathrm{Cm}_{\square}\right)\le\frac{\pi\cdot C'_{V}}{M},\qquad N_{\gamma^{F}}\left(h_{\square}\mathrm{Sloc}_{\square}h_{\square}\right)\le2a'\cdot w_{k}(y),\\
N_{\gamma^{F}}\left(\mathrm{Cm}_{\square}\mathrm{Sloc}_{\square}h_{\square}\right),\ N_{\gamma^{F}}\left(h_{\square}\mathrm{Sloc}_{\square}\mathrm{Cm}_{\square}\right),\ N_{\gamma^{F}}\left(\mathrm{Cm}_{\square}h_{\square'}\mathrm{Sloc}_{\square'}\right)\le\frac{2\pi\cdot C'_{V}\cdot a'}{M},\\
N_{\gamma^{F}}\left(\delta_{\alpha,A}h_{\square}\mathrm{Sloc}_{\square}\right)\le19a'.
\end{gather*}
\item[(e)] \emph{Sums over walks.}
Let $\mathrm{wp}\ge1$ be a number with
\[
\mathrm{qq}(\mathrm{wp}):=\frac{54\pi\cdot C'_{V}\cdot a'\cdot\mathrm{wp}}{M}\le\frac{1}{2}.
\]
For every $n\ge0$, with $S^{\#}_{\mathrm{wk}}$ standing for $S_{y,\mathrm{wk}}(A)$ or for $S'_{y,\mathrm{wk}}(A)$,
\begin{align*}
&\sup_{z}\sum_{z'}\exp\left(\gamma^{F}d(z,z')\right)\sum_{\mathrm{wk}\in\mathrm{Wk}_{n}}\mathrm{wp}^{n+1}\left\|\mathbf{1}_{\Delta_z}S^{\#}_{\mathrm{wk}}\mathbf{1}_{\Delta_{z'}}\right\|_{\mathrm{su}}\\
&\qquad\le16a'\cdot\mathrm{wp}\cdot w_{k}(y)\cdot\mathrm{qq}(\mathrm{wp})^{n},\\
&\sup_{z}\sum_{z'}\exp\left(\gamma^{F}d(z,z')\right)\sum_{\mathrm{wk}\in\mathrm{Wk}_{n}}\mathrm{wp}^{n+1}\left\|\mathbf{1}_{\Delta_z}\delta_{\alpha,A}S^{\#}_{\mathrm{wk}}\mathbf{1}_{\Delta_{z'}}\right\|_{\mathrm{su}}\\
&\qquad\le152a'\cdot\mathrm{wp}\cdot\mathrm{qq}(\mathrm{wp})^{n}.
\end{align*}
\item[(f)] ($k=0$.)
The first bound of (e) holds with
\[
\sup\left\{\left|S^{\#}_{\mathrm{wk}}(x,x')\right|_{2}:\ x\text{ over }z,\ x'\text{ over }z'\right\}
\]
in place of $\left\|\mathbf{1}_{\Delta_z}S^{\#}_{\mathrm{wk}}\mathbf{1}_{\Delta_{z'}}\right\|_{\mathrm{su}}$.
\end{itemize}
\end{named}

\begin{named}{Remark}
Proposition~2.7 corresponds to \cite[Section~16.7.2]{QEDI} with \cite[Lemma~16.13]{QEDI}, to \cite[Definition~16.14]{QEDI} and to \cite[Lemma~16.15]{QEDI}, with the constants $C'_{V}$, $a'$ and $\gamma^{F}$.
Since $h_{\square}$ is a function of $c$, only the hops of type $3$ and the penalty contribute to $\mathrm{Cm}_{\square}$.
The hypothesis $M\ge1600$ of \cite[Lemma~16.13(d)]{QEDI} is absent from (d);
it follows from $M\ge108\pi\cdot C'_{V}\cdot a'$, $C'_{V}\ge24\mathrm{e}$ and $a'\ge1$.
The clause (f), which concerns the kernel norm of (I.ke5), has no counterpart in \cite{QEDI}.
The third assertion of Proposition~2.5\,(c), like that of \cite[Proposition~16.12(c)]{QEDI}, concerns a function on one lift that is supported over $W^{-}$,
whereas for $f\in\mathfrak{S}^{\mathrm{ch}}(\eta_k,\ell)$ the spinor $h_{\square}f$ is antiperiodic and supported over the orbit of $W^{-}$.
In the proof of (b) in Chapter~2b, Section~3, the identity $\mathrm{Sloc}_{\square}F\left(h_{\square}f\right)=h_{\square}f$ is therefore obtained by reducing the spinor $h_{\square}f$ to the function $\mathbf{1}_{W^{-}}h_{\square}f$ on one lift, to which that assertion applies.
In the last bound of (d) the function $h_{\square}$ multiplies also functions of short pairs, through its value at the first point;
this multiplication is estimated by (G1.N3) of Section~G1.3, that is, $N_{\gamma}(\psi O)\le|\psi|_{\infty}\cdot N_{\gamma}(O)$ for a scalar $\psi$ and every operator $O$, including the operators with values in the functions of short pairs.
\end{named}

The constants in Proposition~2.7 are subject to the conditions (L5), (M1), (M2), $M\ge108\pi\cdot C'_{V}\cdot a'$, (Z2.s1) and (Z2.s2), and to $\alpha<1$.

\begin{named}{Proposition 2.8}[FE1$^{\mathrm{ch}}$, FE2$^{\mathrm{ch}}$; the level-$0$ part of FE3$^{\mathrm{ch}}$]
Let the data (SF) be given, let $y\in Y_{k}$ and $A\in\tilde{G}_k$.
\begin{itemize}
\item[(a)] $F_y(A)$ is bijective on $\mathfrak{S}^{\mathrm{ch}}(\eta_k,\ell)$, and
\[
S_y(A)=\sum_{\mathrm{wk}\in\mathrm{Wk}}S_{y,\mathrm{wk}}(A)=\sum_{\mathrm{wk}\in\mathrm{Wk}}S'_{y,\mathrm{wk}}(A),
\]
where both series converge absolutely in $N_{\gamma^{F}}$, uniformly on $\tilde{G}_k$.
\item[(b)] The map $A\mapsto S_y(A)$ is analytic on $\tilde{G}_k$.
\item[(c)] $N_{\gamma^{F}}\left(S_y(A)\right)\le32a'\cdot w_{k}(y)$ and $N_{\gamma^{F}}\left(\delta_{\alpha,A}S_y(A)\right)\le304a'$.
In particular, for unit sites $z$, $z'$ and every $f\in\mathfrak{S}^{\mathrm{ch}}(\eta_k,\ell)$,
\begin{align*}
\left|\mathbf{1}_{\Delta_z}S_y(A)\mathbf{1}_{\Delta_{z'}}f\right|_{\mathrm{su}}&\le32a'\cdot w_{k}(y)\cdot\exp\left(-\gamma^{F}d(z,z')\right)\cdot|f|_{\mathrm{su}},\\
\left|\mathbf{1}_{\Delta_z}\delta_{\alpha,A}S_y(A)\mathbf{1}_{\Delta_{z'}}f\right|_{\mathrm{su}}&\le304a'\cdot\exp\left(-\gamma^{F}d(z,z')\right)\cdot|f|_{\mathrm{su}},
\end{align*}
where in the second inequality the H\"older quotient is taken at the short pairs, the pairs inside a cluster included.
\item[(d)] ($k=0$.)
For the kernel (I.1.4) and every real $y$,
\[
\left|S_{0,y}(A)(x,x')\right|_{2}\le32a'\cdot\left(1+|y|\right)^{-1}\cdot\exp\left(-\gamma^{F}d(x,x')\right);
\]
for $y=0$ this is a bound of the kernel of $\Gamma_0(A)=S^{0}_{1}(A)=S_{0,0}(A)$.
In particular, $F_0(A)=D_A+\mu+b_{L}\cdot JP(A)J^{*}$ is invertible, and $\delta Z_0(A)\neq0$.
\end{itemize}
For $k\ge1$ one has $S_k(A)=S_{\infty}(A)$, $S_{k,y}(A)=S_y(A)$ for real $y$, and $S^{0}_{k+1}(A)=S_0(A)$.
With these identities, (a)--(c) at $y=\infty$ give FE1$^{\mathrm{ch}}$, and (a)--(c) at real $y$ give FE2$^{\mathrm{ch}}$, in the form stated in Proposition~14.1\,(i), once $|\cdot|_{\mathrm{su}}$ is replaced by $|\cdot|_{\infty}$ at the cost of the factor $2$ of Section~I.1.5.
\end{named}

\begin{named}{Remark}
Proposition~2.8 corresponds to \cite[Theorem~16.16]{QEDI}, to \cite[Proposition~5.42]{QEDI} and \cite[Proposition~5.43]{QEDI}, and to the level-$0$ part of \cite[Proposition~5.44]{QEDI}.
The constants are $a'$ and $\gamma^{F}$, the spinors have $192$ components per photon site, and the short pairs are formed with $d_X$.
In \cite{QEDI} the kernel bound at level $0$ is stated after the proof of \cite[Theorem~16.16]{QEDI} and follows there from (c), because at level $0$ the unit blocks of \cite{QEDI} are the sites (see the proof of \cite[Theorem~16.25(a)]{QEDI}).
On the lattice of fermion sites this deduction fails, since the unit block of a site is its cluster of $48$ sites;
its place is taken by (d), a bound of the kernel (I.1.4) at pairs of fermion sites (see (I.P6)$^{\mathrm{ch}}$ and (I.L0.17)), which is obtained through the kernel norm $N^{\mathrm{ke}}_{\gamma^{F}}$.
The invertibility of $F_0(A)$ is stated explicitly in (d).
\end{named}

The constants in Proposition~2.8 are subject to the conditions (L5), (M1), (M2), $M\ge108\pi\cdot C'_{V}\cdot a'$ (which gives $\mathrm{qq}(1)\le\frac{1}{2}$), (Z2.s1) and (Z2.s2), and to $\alpha<1$.

\subsection*{3. Weakened operators and polymer propagators on the fine lattice}

\subsubsection*{Proposition 3}

\begin{named}{Proposition 3}
Let data (SF) be given, together with $y\in Y_k$, $A\in\tilde{G}_k$, $s\in\mathrm{Pd}_k$ and $X\in D_k$.
The families are those of Section~G1.9, and $S^{\#}$ denotes any one of $S$, $S'$, $S^{\mathrm{sy}}$.
Put $N:=N_{\gamma^{F}}$, and recall that $\kappa'=\kappa_{\mathrm{w}}+\kappa_0$ and $\alpha_0=1/54$.
\begin{itemize}
\item[(a)] (Bounds.)
The series of Section~G1.9 converge absolutely in $N$, uniformly on $\mathrm{Pd}_k\times\tilde{G}_k$.
For all unit sites $z$, $z'$ one has
\begin{gather*}
N\left(S^{\#}_y(s,A)\right)\le 32a'\cdot M^{27\alpha_0}\cdot w_k(y),\qquad
N\left(\delta_{\alpha,A}S^{\#}_y(s,A)\right)\le 304a'\cdot M^{27\alpha_0}\quad\text{(on the short pairs)},\\
\left\|\mathbf{1}_{\Delta_z}\cdot S^{\#}_y(X,A)\cdot\mathbf{1}_{\Delta_{z'}}\right\|_{\mathrm{su}}\le 32a'\cdot\exp(27\kappa')\cdot w_k(y)\cdot\exp\left(-\gamma^{F}\cdot d(z,z')-\kappa'\cdot d_M(X)\right).
\end{gather*}
If $k=0$, then for the kernels in the sense of (I.1.4) and for all fermion sites $x$, $x'$
\begin{gather*}
\left|S^{\#}_y(s,A)(x,x')\right|_2\le 32a'\cdot M^{27\alpha_0}\cdot\left(1+|y|\right)^{-1}\cdot\exp\left(-\gamma^{F}\cdot d(x,x')\right),\\
\left|S^{\#}_y(X,A)(x,x')\right|_2\le 32a'\cdot\exp(27\kappa')\cdot\left(1+|y|\right)^{-1}\cdot\exp\left(-\gamma^{F}\cdot d(x,x')-\kappa'\cdot d_M(X)\right).
\end{gather*}
\item[(b)] (Normalisation.)
Let $1$ denote here the family $\left(s_{\mathrm{pc}}\right)_{\mathrm{pc}\in\pi_k}$ with $s_{\mathrm{pc}}=1$ for every $\mathrm{pc}\in\pi_k$;
it lies in $\mathrm{Pd}_k$, because $M^{\alpha_0}\ge1$.
One has $S^{\#}_y(1,A)=S_y(A)$ and $\sum_{X\in D_k}S^{\#}_y(X,A)=S_y(A)$.
\item[(c)] (Analyticity.)
The operator $S^{\#}_y(s,A)$ is a continuous function of $(s,A)$ on $\mathrm{Pd}_k\times\tilde{G}_k$;
it is analytic in $(s,A)$ on the interior of this set, and, for every fixed $s\in\mathrm{Pd}_k$, the points of the boundary included, analytic in $A$ on $\tilde{G}_k$.
The map $S^{\#}_y(X,\cdot\,)$ is analytic on $\tilde{G}_k$.
For $y\in\mathbb{R}$ all these operators depend continuously on $y$.
\item[(d)] (Locality of a walk.)
Let $\mathrm{wk}\in\mathrm{Wk}_n$.
Then $X_{\mathrm{wk}}$ is a polymer, and
\[
\left|X_{\mathrm{wk}}\right|_M\le 27(n+1),\qquad d_M\left(X_{\mathrm{wk}}\right)\le 27(n+1).
\]
The kernels of $S_{y,\mathrm{wk}}(A)$ and of $S'_{y,\mathrm{wk}}(A)$ vanish unless both arguments lie over $X_{\mathrm{wk}}$,
and $S^{\#}_{y,\mathrm{wk}}(A)=S^{\#}_{y,\mathrm{wk}}(\tilde{A})$ whenever $A,\tilde{A}\in\tilde{G}_k$ coincide on the bonds both of whose end points lie in $X_{\mathrm{wk}}$.
If $x$ is the row argument of an entry that does not vanish, then every $M$-cube that meets $K_x=K_c$ is contained in $X_{\mathrm{wk}}$;
the same holds for the column argument.
\item[(e)] (Gauge.)
For every real function $\boldsymbol{\lambda}$ of $c$ on $T_k$,
\[
S^{\#}_{y,\mathrm{wk}}(A-d\boldsymbol{\lambda})=\exp\left(i e_k\boldsymbol{\lambda}\right)\cdot S^{\#}_{y,\mathrm{wk}}(A)\cdot\exp\left(-i e_k\boldsymbol{\lambda}\right);
\]
consequently the same relation holds for $S^{\#}_y(s,A)$ and for $S^{\#}_y(X,A)$.
\item[(f)] (Lattice symmetries.)
Let $\tilde{g}=(S,a)\in\mathbb{G}_k$.
For $\mathrm{wk}=(\square_0,\ldots,\square_n)$ put $\tilde{g}\,\mathrm{wk}:=(\tilde{g}\square_0,\ldots,\tilde{g}\square_n)$;
let $(\tilde{g}s)_{\mathrm{pc}}:=s_{\tilde{g}^{-1}\mathrm{pc}}$;
and let $y'=y$ if $S\in\mathrm{Sp}_3$ or $y$ is neutral, $y'=-y$ otherwise.
Then
\[
X_{\tilde{g}\mathrm{wk}}=\tilde{g}X_{\mathrm{wk}},\qquad(\tilde{g}s)_{\tilde{g}\mathrm{wk}}=s_{\mathrm{wk}},\qquad S^{\#}_{y',\tilde{g}\mathrm{wk}}(\tilde{g}A)\cdot\tilde{g}=\tilde{g}\cdot S^{\#}_{y,\mathrm{wk}}(A),
\]
and therefore
\[
S^{\#}_{y'}(\tilde{g}s,\tilde{g}A)\cdot\tilde{g}=\tilde{g}\cdot S^{\#}_y(s,A),\qquad S^{\#}_{y'}(\tilde{g}X,\tilde{g}A)\cdot\tilde{g}=\tilde{g}\cdot S^{\#}_y(X,A).
\]
\item[(g)] (Transposition.)
\[
S_{y,\mathrm{wk}}(A)^{T}=\mathsf{C}^{-1}\cdot S'_{-y,\mathrm{wk}}(-A)\cdot\mathsf{C},\qquad S'_{y,\mathrm{wk}}(A)^{T}=\mathsf{C}^{-1}\cdot S_{-y,\mathrm{wk}}(-A)\cdot\mathsf{C}.
\]
Hence, for neutral $y$,
\[
S^{\mathrm{sy}}_y(s,A)^{T}=\mathsf{C}^{-1}\cdot S^{\mathrm{sy}}_y(s,-A)\cdot\mathsf{C},\qquad S^{\mathrm{sy}}_y(X,A)^{T}=\mathsf{C}^{-1}\cdot S^{\mathrm{sy}}_y(X,-A)\cdot\mathsf{C}.
\]
\item[(h)] (Conjugation.)
For $s\in\mathrm{Pd}_k$ let $\overline{s}:=\left(\overline{s_{\mathrm{pc}}}\right)_{\mathrm{pc}\in\pi_k}$;
since $\left|\overline{s_{\mathrm{pc}}}\right|=\left|s_{\mathrm{pc}}\right|$, the family $\overline{s}$ lies again in $\mathrm{Pd}_k$.
One has
\[
\overline{S^{\#}_{y,\mathrm{wk}}\left(\overline{A}\right)}=\varkappa^{-1}\cdot S^{\#}_{y,\mathrm{wk}}(-A)\cdot\varkappa;
\]
hence
\[
\overline{S^{\#}_y\left(\overline{s},\overline{A}\right)}=\varkappa^{-1}\cdot S^{\#}_y(s,-A)\cdot\varkappa,\qquad\overline{S^{\#}_y\left(X,\overline{A}\right)}=\varkappa^{-1}\cdot S^{\#}_y(X,-A)\cdot\varkappa.
\]
\item[(i)] (Rows and columns are exactly local where the parameters vanish.)
Let $\mathscr{Y}$ be a subset of $\pi_k$, let $Y$ be its union, and let $s\in\mathrm{Pd}_k$ satisfy $s_{\mathrm{pc}}=0$ for every $\mathrm{pc}$ outside $\mathscr{Y}$.
Then $s_{\mathrm{wk}}=0$ unless $X_{\mathrm{wk}}$ is contained in $Y$;
and if $X_{\mathrm{wk}}$ is contained in $Y$, then $X_{\mathrm{wk}}$, which is connected through faces, lies in a single link-component of $Y$.
Consequently, for every link-component $Y_1$ of $Y$ and every fermion site $x$ that lies over $Y_1$, one has $S^{\#}_y(s,A)(x,x')=0$ unless $x'$ lies over $Y_1$, and
\[
S^{\#}_y(s,A)(x,\cdot\,)=S^{\#}_y(\tilde{s},\tilde{A})(x,\cdot\,)
\]
for every $\tilde{s}\in\mathrm{Pd}_k$ that vanishes outside $\mathscr{Y}$ and coincides with $s$ on the cubes of $Y_1$, and for every $\tilde{A}\in\tilde{G}_k$ that coincides with $A$ on the bonds that meet $Y_1$.
The same holds with rows and columns interchanged.
Finally, for every fermion site $x=(c,g)$, the row of $S^{\#}_y(s,A)$ at $x$ vanishes unless $\mathscr{Y}$ contains all cubes of $\pi_k$ that meet $K_c$;
in particular it vanishes if $c$ does not lie in $Y$.
\end{itemize}
\end{named}

\emph{Conditions on the constants.}
The conditions that enter here, and that are collected in Proposition~14.2, are the following:
the conditions listed before Proposition~2.1;
the range $\alpha<1$ of the line (C1) of Section~I.8.2, which enters through the second bound of Proposition~2.7\,(e) and through Proposition~2.8;
(M1), of which only the part that $M$ is a power of $L$ enters;
(M5) in full;
and (M2).
Within (M5), with the number $\mathrm{qq}$ of Proposition~2.7\,(e), the inequality $M\ge\left(108\pi\cdot C'_V\cdot a'\right)^2$ gives $\mathrm{qq}\left(M^{27\alpha_0}\right)\le 1/2$, because $27\alpha_0\le 1/2$;
and $M\ge 108\pi\cdot C'_V\cdot a'\cdot\exp(27\kappa')$ gives $\mathrm{qq}\left(\exp(27\kappa')\right)\le 1/2$.

\begin{named}{Remark}
Proposition~3 corresponds to \cite[Definition~16.17]{QEDI} and \cite[Proposition~16.18]{QEDI}.
The constants are $a'$ and $\gamma^{F}$, and the rate $\kappa'=\kappa_{\mathrm{w}}+\kappa_0$ is a fixed number.
Membership in $X_{\mathrm{wk}}$ is read at the photon site $c$ of a fermion site $x=(c,g)$, by (I.POS), and $K_x=K_c$;
the elements $\tilde{g}$ act also on the label $g$;
the transposes are taken with respect to the pairing $\langle\,\cdot\,,\cdot\,\rangle^{\mathrm{ch}}$;
the kernels on the torus are those of Section~G1.2.
The kernel bounds of (a) at $k=0$, which are obtained through the kernel norm $N^{\mathrm{ke}}_{\gamma^{F}}$ of (I.ke5), have no counterpart in \cite[Proposition~16.18]{QEDI}.
Clause (i) is the first of the two facts on locality in the proof of \cite[Theorem~16.23]{QEDI}, together with its consequences.
\end{named}

\subsection*{4. Resolvent algebra between the fine and the unit lattice}

\subsubsection*{Proposition 4}

\begin{named}{Proposition 4}
Fix data (SF) with $k\ge1$, a connection $A\in\tilde{G}_k$ and a real number $y$.
Write
\[
\mathrm{Qk}:=Q_k(A)\cdot J^{*},\qquad\mathrm{Qt}:=J\cdot Q_k^{T}(-A),\qquad P:=P(A),
\]
and let $H_k$, $\mathfrak{H}_k$, $M_k$, $D_k$, $B_{k,y}$, $\Gamma_{k,y}$ and $\delta Z_k$ be as in Section~I.3.5.
\begin{itemize}
\item[(a)] The operator $b_k+b_L\cdot P+i\gamma_5\cdot y$ is invertible, and its inverse is
\[
B_{k,y}(A)=\left(b_k+i\gamma_5\cdot y\right)^{-1}\cdot(1-P)+\left(b_k+b_L+i\gamma_5\cdot y\right)^{-1}\cdot P.
\]
Moreover
\[
b_k-b_k^2\cdot B_{k,y}(A)=\boldsymbol{\alpha}_y+\boldsymbol{\beta}_y\cdot P,\qquad F_y(A)=F_{\infty}(A)-b_k^2\cdot\mathrm{Qt}\cdot B_{k,y}(A)\cdot\mathrm{Qk}.
\]
\item[(b)] The operator $D_k(A)+b_L\cdot P(A)+i\gamma_5\cdot y$ is a bijection of $\mathfrak{S}(1,\ell)$ onto itself, with the inverse
\[
\Gamma_{k,y}(A)=B_{k,y}(A)+b_k^2\cdot B_{k,y}(A)\cdot\mathrm{Qk}\cdot S_y(A)\cdot\mathrm{Qt}\cdot B_{k,y}(A).
\]
\item[(c)] $S_y(A)=S_k(A)+H_k(A)\cdot\Gamma_{k,y}(A)\cdot\mathfrak{H}_k(A)$.
\item[(d)] The inverse $\Gamma_k(A)=\left(D_k(A)+b_L\cdot P(A)\right)^{-1}$ exists, one has
\[
\Gamma_k(A)=B_k(A)+b_k^2\cdot B_k(A)\cdot\mathrm{Qk}\cdot S^{0}_{k+1}(A)\cdot\mathrm{Qt}\cdot B_k(A),
\]
and $\delta Z_k(A)\ne0$.
\end{itemize}
\end{named}

\begin{named}{Remark}
Proposition~4 corresponds to the opening paragraph of \cite[Section~16.9]{QEDI} and to \cite[Proposition~16.19]{QEDI}.
Here $\mathrm{Qk}$ and $\mathrm{Qt}$ contain the cluster maps, and the matrix $\gamma_5$ stands in the place of the matrix $\gamma_3$ of \cite[Section~16.9]{QEDI}.
The identities $\mathrm{Qt}\cdot\mathrm{Qk}=J\cdot P_k(A)\cdot J^{*}$ and $\mathrm{Qt}\cdot P\cdot\mathrm{Qk}=J\cdot P_{k+1}(A)\cdot J^{*}$ follow from the definitions of Section~I.3.3,
and $\mathrm{Qk}$, $\mathrm{Qt}$ and $P$ commute with every constant $4\times4$ matrix, since their coefficients are numbers.
The second identity of (a) is the step that leads from the first identity of (a) to the third.
\end{named}

\subsection*{5. Kernels of the unit-lattice operators at the levels $k\ge1$}

\subsubsection*{Proposition 5}

\begin{named}{Proposition 5}
Define the constants
\[
C_{\mathrm{By}}:=\left(1+2\mathrm{e}\right)\cdot\sqrt{2}\cdot\max\left(1,\frac{1}{b_-}\right),\qquad
C_{\Gamma}:=\frac{2\left(1+2\mathrm{e}\right)}{b_-}+128\,\mathrm{e}\cdot\left(1+2\mathrm{e}\right)^2\cdot a'.
\]
Let data (SF) with $k\ge1$ be given, together with $A\in\tilde{G}_k$, $y\in\mathbb{R}$ and sites $z$, $z'$ of $T^0_k$;
the indicator $\mathbf{1}[z=z']$ and the words ``one $L$-block'' are understood on the torus.
\begin{itemize}
\item[(a)] If $z$ and $z'$ lie in one set $B(y')$, then
\[
P(A)(z,z')=L^{-3}\cdot U_A(y',z)^{-1}\cdot U_A(y',z')\cdot 1_4,
\]
and $P(A)(z,z')=0$ otherwise.
In either case
\[
\left|P(A)(z,z')\right|_2=\left|P(A)(z,z')\right|_c\le\mathrm{e}\cdot L^{-3},
\]
both norms being equal to the modulus of the number that multiplies $1_4$.
Assertion (a) holds also for data (SF) with $k=0$ and $A\in\tilde{G}_0$, the operator $P(A)$ being that of Section~I.3.3 on $\mathfrak{S}(1,\ell)$;
in that case the kernel (I.1.4) of $J\cdot P(A)\cdot J^{*}$ at $(x,x')$ is $P(A)(c,c')$, and
\[
\sup_x\int\mathrm{d}x'\,\left|b_L\cdot\left(J\cdot P(A)\cdot J^{*}\right)(x,x')\right|_c\le b_L\cdot\mathrm{e}\le b\cdot\mathrm{e},
\]
which is the last term of (I.L0.12).
\item[(b)] One has $B_{k,y}(A)(z,z')=0$ unless $z$ and $z'$ lie in one $L$-block, and
\[
\left|B_{k,y}(A)(z,z')\right|_2\le\left(b_k^2+y^2\right)^{-\frac{1}{2}}\cdot\left(\mathbf{1}[z=z']+2\mathrm{e}\cdot L^{-3}\right)\le C_{\mathrm{By}}\cdot\left(1+|y|\right)^{-1}.
\]
Each sum over a row and each sum over a column of $\left|B_{k,y}(A)(z,z')\right|_2$ is at most $\left(1+2\mathrm{e}\right)\cdot\left(b_k^2+y^2\right)^{-\frac{1}{2}}$.
\item[(c)] For every operator $O$ on $\mathfrak{S}^{\mathrm{ch}}(\eta_k,\ell)$,
\[
\left|\left(Q_k(A)\cdot J^{*}\cdot O\cdot J\cdot Q_k^{T}(-A)\right)(z,z')\right|_2\le\mathrm{e}\cdot\left\|\mathbf{1}_{\Delta_z}\cdot O\cdot\mathbf{1}_{\Delta_{z'}}\right\|_{\mathrm{su}}.
\]
\item[(d)] The following bounds hold:
\begin{gather*}
\left|M_k(A)(z,z')\right|_2\le 32\,\mathrm{e}\cdot b^2\cdot a'\cdot\exp\left(-\gamma^{F}\cdot d(z,z')\right),\qquad
\left|\dot{M}_k(A)(z,z')\right|_2\le\mathrm{e}\cdot b^2\cdot\left(32a'\right)^2\cdot\exp\left(-\gamma^{F}\cdot d(z,z')\right),\\
\left|\left(\Gamma_{k,y}(A)-B_{k,y}(A)\right)(z,z')\right|_2\le 128\,\mathrm{e}\cdot\left(1+2\mathrm{e}\right)^2\cdot a'\cdot b_k^2\cdot\left(b_k^2+y^2\right)^{-1}\cdot\exp\left(-\gamma^{F}\cdot d(z,z')\right),\\
\left|\Gamma_k(A)(z,z')\right|_2\le C_{\Gamma}\cdot\exp\left(-\gamma^{F}\cdot d(z,z')\right).
\end{gather*}
\item[(e)] (The property FE3$^{\mathrm{ch}}$.)
With $|\cdot|_c$ in place of $|\cdot|_2$, all these bounds hold with the constants doubled.
In (c) this is to be read with $|\cdot|_c$ on the left, $\|\cdot\|_{\mathrm{su}}$ on the right, and the constant $2\mathrm{e}$;
if the operator norm for $|\cdot|_{\infty}$ is taken on the right, the bound holds with the constant $\mathrm{e}$.
The doubled constants do not exceed $C^{F}$, which is defined in Proposition~14.1.
For $k\ge1$ this is FE3$^{\mathrm{ch}}$:
\begin{gather*}
\left|B_{k,y}(A)(z,z')\right|_c\le C^{F}\cdot\left(1+|y|\right)^{-1}\cdot\mathbf{1}[\,z\text{ and }z'\text{ lie in one $L$-block}\,],\\
\left|M_k(A)(z,z')\right|_c,\ \left|\dot{M}_k(A)(z,z')\right|_c,\ \left|\Gamma_k(A)(z,z')\right|_c\le C^{F}\cdot\exp\left(-\gamma^{F}\cdot d(z,z')\right);
\end{gather*}
hence
\[
\left|D_k(A)(z,z')\right|_c\le b_+\cdot\mathbf{1}[z=z']+C^{F}\cdot\exp\left(-\gamma^{F}\cdot d(z,z')\right).
\]
\end{itemize}
\end{named}

\emph{Conditions on the constants.}
The conditions that enter directly are (Z2.s1) and, of (L5) and (M1), only the facts that $L$ is odd, that $L>1$, and that $M$ is an integer with $M\ge(L-1)/2$;
the inequality $L>1$ gives $b_-=b\cdot(1-1/L)>0$ and $b_L=b/L\le b$.
Clauses (a)--(c) require these conditions only.
Clauses (d) and (e) rest in addition, through Proposition~2.8\,(a), (c) and Proposition~4\,(b), on the conditions of Proposition~2.8, that is, on (L5), (M1), (M2), (Z2.s1), (Z2.s2), $\alpha<1$, and the inequality $M\ge108\pi\cdot C'_V\cdot a'$, which is part of (M5).

\begin{named}{Remark}
Proposition~5 corresponds to the constants defined in \cite[Section~16.9.2]{QEDI} and to \cite[Proposition~16.20]{QEDI};
for $k\ge1$, clause (e) corresponds to \cite[Proposition~5.44]{QEDI}.
The constants are $a'$ and $\gamma^{F}$.
In (c) the test function is $J\left[\varphi_A(z',\cdot\,)^{-1}\cdot\mathbf{1}_{\Delta_{z'}}\cdot v\right]$, and $J$ does not change the sup norm, by (I.J2).
The indicator $\mathbf{1}[\,\cdot\,]$ is that of Section~I.0.4, read modulo $\ell\mathbb{Z}^3$.
Clause (e) fixes the norm in which the bounds of \cite[Proposition~5.44]{QEDI} are read.
The operator $\dot{M}_k(A)$ in (d) is that of Section~I.3.5.
\end{named}

\subsection*{6. The field $\psi=H_k(A)\Psi$ and all its H\"older quotients}

\subsubsection*{Proposition 6}

\begin{named}{Proposition 6}
Consider data (SF), a connection $A\in\tilde{G}_k$, sites $z$, $z'$, $z''$ of $T^0_k$, a vector $\mathrm{vv}\in\mathbb{C}^4$, spinors $v\in\mathfrak{S}(1,\ell)$ and $f\in\mathfrak{S}^{\mathrm{ch}}(\eta_k,\ell)$, and a pair $(x,x')\in\mathrm{Pzt}_k$.
\begin{itemize}
\item[(a)] (Localised bounds on the short pairs; $k\ge1$.)
\begin{gather*}
\left\|\mathbf{1}_{\Delta_z}\cdot H_k(A)\cdot\mathbf{1}_{z'}\right\|_{\mathrm{su}}\le 32\,\mathrm{e}\cdot b\cdot a'\cdot\exp\left(-\gamma^{F}\cdot d(z,z')\right),\\
\left\|\mathbf{1}_{\Delta_z}\cdot\delta_{\alpha,A}H_k(A)\cdot\mathbf{1}_{z'}\right\|_{\mathrm{su}}\le 304\,\mathrm{e}\cdot b\cdot a'\cdot\exp\left(-\gamma^{F}\cdot d(z,z')\right),
\end{gather*}
the second bound on the short pairs, that is, the pairs with $d_X<1$, among them the pairs inside a cluster.
\item[(b)] (All $k\ge0$.)
For every pair of $\mathrm{Pzt}_k$, short or long, $\left|U_A(x,x')^{\pm1}\right|\le\mathrm{e}$.
For a long pair, $1\le d_X(x,x')\le M$, one has
\[
\left|\left(\delta_{\alpha,A}f\right)(x,x')\right|_2\le\mathrm{e}\cdot\left|f(x')\right|_2+\left|f(x)\right|_2\le 4\left|f\right|_{\mathrm{su}},
\]
and the same holds with $|\cdot|_c$ and $|f|_{\infty}$; this is (I.Lc2)$^{\mathrm{ch}}$ on $\tilde{G}_k$.
If $k\ge1$ and the long pair has $x$ over $\Delta_z$ and $x'$ over $\Delta_{z''}$, then
\[
\left|\left(\delta_{\alpha,A}H_k(A)\cdot\mathbf{1}_{z'}\mathrm{vv}\right)(x,x')\right|_2\le 128\,\mathrm{e}\cdot b\cdot a'\cdot\exp\left(-\gamma^{F}\cdot\min\left(d(z,z'),d(z'',z')\right)\right)\cdot\left|\mathrm{vv}\right|_2.
\]
\item[(c)] (Global bounds; $k\ge1$.)
\[
\left|H_k(A)v\right|_{\mathrm{su}}\le 32\,\mathrm{e}\cdot b\cdot a'\cdot|v|_{\mathrm{su}},\qquad
\sup_{(x,x')\in\mathrm{Pzt}_k}\left|\left(\delta_{\alpha,A}H_k(A)v\right)(x,x')\right|_2\le 304\,\mathrm{e}\cdot b\cdot a'\cdot|v|_{\mathrm{su}},
\]
the supremum ranging over all pairs: those inside a cluster, the short ones, and the long ones up to the distance $M$.
\item[(d)] (Barred fields.)
The assertions (a)--(c) hold for $\bar{H}_k(A)=\mathsf{C}^{-1}\cdot H_k(-A)\cdot\mathsf{C}$, with $\delta_{\alpha,-A}$ in place of $\delta_{\alpha,A}$.
\item[(e)] (Level $0$.)
One has $H_0(A)=\bar{H}_0(A)=1$ on $\mathfrak{S}^{\mathrm{ch}}(1,\ell)$, and there is no short pair;
for a pair inside a cluster $\left(\delta_{\alpha,A}f\right)(x,x')=f(x')-f(x)$;
and
\[
\left|H_0^{\#}(A)f\right|_{\infty}\le\left(1+\mathrm{e}\right)\cdot|f|_{\infty}\le 4|f|_{\infty},
\]
together with the same bound for the pair $\left(\bar{H}_0(A),\delta_{\alpha,-A}\bar{H}_0(A)\right)$.
The localised assertions, namely (a) and the second part of (b), concern $k\ge1$ only.
\item[(f)] (The property FE4$^{\mathrm{ch}}$.)
For $k\ge1$, in the norms of Section~I.1.5,
\[
\left|H_k(A)\right|_{1,\infty},\ \left|\delta_{\alpha,A}H_k(A)\right|_{1,\infty}\le C^{F},
\]
the rows being those at all fermion sites, respectively at all pairs of $\mathrm{Pzt}_k$, and hence at the pairs of every $\mathrm{Pz}_k(Z)$;
\[
\left|\mathbf{1}_{\Delta_z}\cdot H_k(A)\cdot\mathbf{1}_{z'}\mathrm{vv}\right|_{\infty},\ \left|\mathbf{1}_{\Delta_z}\cdot\delta_{\alpha,A}H_k(A)\cdot\mathbf{1}_{z'}\mathrm{vv}\right|_{\infty}\le C^{F}\cdot\exp\left(-\gamma^{F}\cdot d(z,z')\right)\cdot|\mathrm{vv}|_c,
\]
the second on the short pairs;
on the long pairs the bound of (b) holds with $C^{F}$ in place of $128\,\mathrm{e}\cdot b\cdot a'$;
and the same bounds hold for the barred operators.
For $k=0$ the property FE4$^{\mathrm{ch}}$ consists of (e), together with $4\le C^{F}$.
\end{itemize}
\end{named}

\emph{Conditions on the constants.}
The conditions are (Z2.s1), (Z2.s3), (L5) and (M1), of which only $L$ odd, $L>1$ and $M\ge1$ enter from the last two;
those of Proposition~2.8;
and $\alpha>0$ and $\rho>0$, which are the ranges of the lines (C1) and (C0).

\begin{named}{Remark}
Proposition~6 corresponds to \cite[Proposition~16.21]{QEDI} and to \cite[Proposition~5.45]{QEDI};
the first part of (b) corresponds to \cite[Lemma~5.23]{QEDI}, on $\tilde{G}_k$.
The constants are $a'$ and $\gamma^{F}$, the distance is $d_X$, the pairs inside a cluster are included, and $H_k(A)$ contains the cluster map $J$.
The statement of \cite[Proposition~5.45]{QEDI} carries no restriction on the level:
by \cite[Theorem~16.32(i)]{QEDI} it holds for every $k\ge0$, the localised bounds and the bounds of the barred field included.
At level $0$ the property FE4$^{\mathrm{ch}}$ consists of the one bound of (e), for both fields, and its localised assertions are restricted to $k\ge1$; this is (I.L0.15).
At $k=0$ one has $H_0(A)=\bar{H}_0(A)=1$, the operator $\bar{H}_0$ being that of Section~I.3.9;
with this, the first localised bound and the bound on the long pairs of \cite[Proposition~5.45]{QEDI} are immediate, and the bounds of the barred field are the same;
the second localised bound is void at $k=0$, since at that level there is no short pair, by (I.L0.14).
Clause (f) fixes the norms in which the bounds of \cite[Proposition~5.45]{QEDI} are read: the norm $|\cdot|_c$ on $\mathbb{C}^4$ and the sup norm $|\cdot|_{\infty}$ of Section~I.1.5 and, for operators, the norm $|\cdot|_{1,\infty}$ of (I.ke3), which is the operator norm for these sup norms.
\end{named}

\subsection*{7. The localised operator families at the levels $k\ge1$}

\subsubsection*{Proposition 7 (W1, W2, W5--W16)}

\begin{named}{Proposition 7}[W1, W2, W5--W16]
Put
\[
C_1:=\max\left\{304\,\mathrm{e}\cdot\max(1,b)\cdot a'\cdot M^{27\alpha_0},\ C_{\Gamma}\cdot M^{27\alpha_0},\ 32\,\mathrm{e}\cdot\max\left(1,b^2\right)\cdot a'\cdot\exp(27\kappa')\right\}.
\]
Let data (SF) with $k\ge1$ be given.
For $s\in\mathrm{Pd}_k$, $A\in\tilde{G}_k$, $X\in D_k$ and $y\ge0$ define, with the families of Section~G1.9,
\begin{equation}\tag{p3}
H_k(s,A):=b_k\cdot S_{\infty}(s,A)\cdot J\cdot Q_k^{T}(-A),\qquad\bar{H}_k(s,A):=\mathsf{C}^{-1}\cdot H_k(s,-A)\cdot\mathsf{C};
\end{equation}
\begin{equation}\tag{p4}
\begin{split}
\Gamma_k(s,A)&:=B_k(A)+b_k^2\cdot B_k(A)\cdot Q_k(A)\cdot J^{*}\cdot S^{\mathrm{sy}}_0(s,A)\cdot J\cdot Q_k^{T}(-A)\cdot B_k(A),\\
M_k(X,A)&:=b_k^2\cdot Q_k(A)\cdot J^{*}\cdot S^{\mathrm{sy}}_{\infty}(X,A)\cdot J\cdot Q_k^{T}(-A);
\end{split}
\end{equation}
\begin{equation}\tag{p5}
S_{k,y}(X,A):=S_y(X,A).
\end{equation}
These families have the properties W1, W2, W5--W16 of Section~G1.11 for $\left(\alpha_0,\gamma^{F},\kappa';2C_1\right)$.
The relations of W14 hold for every $\tilde{g}\in\mathbb{G}_k$, and, for the family (p5), for every $\tilde{g}\in\mathbb{G}_{k,3}$.
The map $y\mapsto S_{k,y}(X,A)$ is continuous on $[0,\infty)$.
In the Euclidean norms the properties W5--W7 hold with the constant $C_1$:
\begin{gather*}
\left|H_k(s,A)v\right|_{\mathrm{su}}\le 32\,\mathrm{e}\cdot b\cdot a'\cdot M^{27\alpha_0}\cdot|v|_{\mathrm{su}},\qquad
\left|\delta_{\alpha,A}H_k(s,A)v\right|_{\mathrm{su}}\le 304\,\mathrm{e}\cdot b\cdot a'\cdot M^{27\alpha_0}\cdot|v|_{\mathrm{su}}\quad\text{(on the short pairs)},\\
\left|\Gamma_k(s,A)(z,z')\right|_2\le C_{\Gamma}\cdot M^{27\alpha_0}\cdot\exp\left(-\gamma^{F}\cdot d(z,z')\right),\\
\left|M_k(X,A)(z,z')\right|_2\le 32\,\mathrm{e}\cdot b^2\cdot a'\cdot\exp(27\kappa')\cdot\exp\left(-\gamma^{F}\cdot d(z,z')-\kappa'\cdot d_M(X)\right),
\end{gather*}
and for W7 the bound is that of Proposition~3\,(a).
\end{named}

\emph{Conditions on the constants.}
The conditions are those of Proposition~3, and (M1), which enters through the fact (g0) of \cite[Section~17.4.1]{QEDI}.

\begin{named}{Remark}
Proposition~7 corresponds to \cite[Definition~16.22]{QEDI} and \cite[Theorem~16.23]{QEDI};
for the families (p3)--(p5), the properties W1, W2, W5--W16 correspond to the properties 1, 2, 5--16 of \cite[Section~5.4.2]{QEDI}.
The cluster maps $J$ and $J^{*}$ enter the definitions;
the constants are $a'$ and $\gamma^{F}$, and $\kappa'$ is fixed.
The constant is $2C_1$ for the norms of Section~I.1.5: a bound proved with a constant $C$ for the Euclidean sup norms holds with the constant $2C$ in these norms.
W5 covers the pairs inside a cluster;
``contained in'' for an index $x=(c,g)$ is read at $c$;
W11 is stated with $K_x=K_c$.
Of the family (p5) the step requires the covariance under $\mathbb{G}^{L}_{k,3}$ only (Section~I.5.8).
The nesting of the unit blocks, the $L$-blocks and the $M$-cubes is the fact (g0) of \cite[Section~17.4.1]{QEDI}.
\end{named}

\subsection*{8. The localised operator families at level 0 (W17)}

\subsubsection*{Proposition 8}

\begin{named}{Proposition 8}
Let data (SF) with $k=0$ be given.
For $s\in\mathrm{Pd}_0$, $A\in\tilde{G}_0$, $X\in D_0$ and $y\ge0$ define the operators
\[
\Gamma_0(s,A):=S^{\mathrm{sy}}_0(s,A),\qquad S_{0,y}(X,A):=S_y(X,A)
\]
on $\mathfrak{S}^{\mathrm{ch}}(1,\ell)$, which has $192$ components per photon site;
let $H_0(s,A)$ be the multiplication by
\[
\prod\left\{s_{\mathrm{pc}}:\ \mathrm{pc}\in\pi_0,\ \mathrm{pc}\text{ meets }K_c\right\},
\]
and let $\bar{H}_0(s,A):=\mathsf{C}^{-1}\cdot H_0(s,-A)\cdot\mathsf{C}=H_0(s,A)$.
\begin{itemize}
\item[(a)] The families $\Gamma_0(s,A)$ and $S_{0,y}(X,A)$ have the properties W1, W2, W6--W9 and W12--W16 of Section~G1.11, each as far as it concerns the family, for $\left(\alpha_0,\gamma^{F},\kappa';2C_1\right)$;
here W6 and W7 are the following bounds of the kernels (I.1.4) at pairs of fermion sites, valid for all $y\ge0$:
\begin{gather*}
\left|\Gamma_0(s,A)(x,x')\right|_c\le 2C_1\cdot\exp\left(-\gamma^{F}\cdot d(x,x')\right),\\
\left|S_{0,y}(X,A)(x,x')\right|_c\le 2C_1\cdot\left(1+y\right)^{-1}\cdot\exp\left(-\gamma^{F}\cdot d(x,x')-\kappa'\cdot d_M(X)\right).
\end{gather*}
The map $y\mapsto S_{0,y}(X,A)$ is continuous on $[0,\infty)$, hence Borel.
W14 holds for every $\tilde{g}\in\mathbb{G}_0$, and, for $S_{0,y}(X,A)$, for every $\tilde{g}\in\mathbb{G}_{0,3}$.
W1 takes the form
\[
\Gamma_0(1,A)=\Gamma_0(A)=F_0(A)^{-1},\qquad\sum_{X}S_{0,y}(X,A)=S_{0,y}(A)=\left(F_0(A)+i\gamma_5\cdot y\right)^{-1}.
\]
\item[(b)] The family $H_0(s,A)$ has W1, W2, W8--W11, W13, W14 (for every $\tilde{g}\in\mathbb{G}_0$) and W16, and it has W5 with the constant
\[
M^{8\alpha_0}\le 4\exp(2\gamma)\cdot M^{27\alpha_0},
\]
for every $\gamma\ge0$; the second part of W5 is void, because there is no short pair.
The operator $H_0(s,A)$ does not depend on $A$, is a scalar in the spin components, and is the same at the $48$ sites of a cluster.
\end{itemize}
\end{named}

\emph{Conditions on the constants.}
The conditions are those of Proposition~3.

\begin{named}{Remark}
Proposition~8 corresponds to \cite[Definition~16.24]{QEDI} and \cite[Theorem~16.25]{QEDI}, and to the fermionic part of property~17 of \cite[Section~5.4.2]{QEDI}.
The spinors have $192$ components per photon site.
In the proof of \cite[Theorem~16.25(a)]{QEDI} the unit blocks at $k=0$ are the sites, and $\left\|\mathbf{1}_{\{x\}}\cdot O\cdot\mathbf{1}_{\{x'\}}\right\|=\left|O(x,x')\right|$;
on the lattice of fermion sites of Section~I.1.3 the unit block of a site at $k=0$ is its cluster of $48$ sites, and this identity fails.
Here W6 and W7 are bounds of the kernel (I.1.4);
they follow from the second display of Proposition~3\,(a), whose bounds are obtained through the kernel norm $N^{\mathrm{ke}}_{\gamma^{F}}$ of (I.ke5).
The constant is $2C_1$, by Section~I.1.5.
The Borel measurability in $y$ is part of the assertion at level $0$.
At $k=0$ the bound of W7 belongs to the property W7 of Section~G1.11, so that W7 for $\mathrm{par}$ means at level $0$ the bound with the factor $(1+y)^{-1}$.
\end{named}

\subsection*{G2. Definitions for the Propositions 15--20.3 and 48.1--48.6 that are not in Part I}

\subsubsection*{G2.0.4. Settings}

The constants are those of Section~I.8.
They are chosen in the order of Section~I.8.2, each sufficiently large or sufficiently small, so that the conditions collected in Section~G2.9 (Annex~B) hold;
these are the conditions on the constants that enter the proofs of the Propositions~15--20.3 and 48.1--48.6.

The constants that occur in these statements are, in the order in which they are chosen, the following.
\begin{description}
\item[(C0)] $\rho$ and the weights $w_i$; the absolute numbers $c_X$, $\kappa_0$, $K_0$, $c_{\mathrm{rb}}$, $c_{\mathrm{fl}}$, $b$.
\item[(C1)] $\alpha$.
\item[(C2)] $\kappa$.
\item[(C3)] $L$.
\item[(C4)] $b_L$, $b_-$, $b_+$, which are functions of $b$ and $L$;
the number $\alpha_0=1/54$;
the rate $\kappa_{\mathrm{w}}=\kappa+\kappa_{0}+3$.
\item[(C4r)] $\gamma^{F}$, a function of $b$, $\alpha$, $L$ and $\rho$;
the rate $\gamma_{\mathrm{w}}=\min\left(\gamma^{F},\gamma^{P}\right)$, where $\gamma^{P}$ is the rate of the line (C4) of Section~I.8.2;
through $\gamma^{F}$ it depends on $\rho$.
\item[(C5)] $M$.
\item[(C6)] $C^{F}$ and $C_0$, functions of the constants of the lines (C0)--(C5), $\rho$ among them;
the constant $C_{w}=64K_{0}\cdot C_{w}^{(1)}$, formed with the constant $C_{w}^{(1)}$ of the same line (Section~I.8.2);
the set of parameters $\mathrm{par}=\left(\alpha_0,\gamma_{\mathrm{w}},\kappa_{\mathrm{w}};C_{w}\right)$.
\item[(C7)] $\varepsilon$;
the integer $r\ge2$.
\item[(C8)] $e_0$.
\item[(C9)] $e_1$.
\item[(C10)] $\bar{m}$ and $h$.
\end{description}
The set $\mathrm{par}$ occurs only in the hypotheses on the primitives in Propositions~20.3 and~48.6;
the integer $r$ enters only through the numbers $r_k$ of Section~I.6.1, in Proposition~48.6\,(a), and through $\mathrm{Tor}(e)$, $N_{\min}$ and $k_{*}$ in (S-4).
For $k\ge1$ the number $b_k$ is not a constant of the line (C4):
it is a number attached to the level, a function of $(b,L,k)$ with $b_-\le b_k\le b_+$ (Section~I.3.4).
The numerical factors, exponents and dimensions that occur in the statements are absolute numbers.
No constant of the lines (C0)--(C9) depends on $N$, $k$, $\ell$, the twist, the running mass, the pair $(m_{\mathrm{in}},\varepsilon^{\mathrm{v}}_{\mathrm{in}})$, $\mathscr{C}$ or $\bar{m}$ ((I.U4)$^{\mathrm{ch}}$);
the exceptions are $\bar{m}$ and $h$, a function of $(L,\bar{m})$, which stand in the line (C10).
The settings (S-1)--(S-3) involve the constants of the lines (C0)--(C8) only;
$e_1$ and the condition that defines it enter in (S-4) only.

The settings are the following.
\begin{description}
\item[(S-abs)] No datum of the scheme (Proposition~15\,(a)--(m)).
\item[(S-1)] \emph{One level.}
An integer $k\ge0$; a side $\ell\in\mathrm{Cl}$; a number $e_k\in(0,e_0]$, which at $k=0$ is the number called $q$ in the flow (Section~I.2.5).
The objects $T_k$, $T^{\mathrm{f}}_k$, $\mathrm{Pz}_k$, $D_k$, the generators $\psi$, $\chi$ (Section~I.3.6), the domains (Section~I.5.3) and the spaces of activities (Section~I.5.4) are formed with $(k,\ell,e_k)$.
Neither a mass nor a primitive occurs.
\item[(S-2)] \emph{Two levels.}
(S-1), and moreover $\ell/L\in\mathrm{Cl}$ and $e_{k+1}:=L^{\frac{1}{2}}e_k\le e_0$.
The level $k$ lives on $T_k=\mathbb{L}_k/\ell\mathbb{Z}^3$, the level $k+1$ on $T_{k+1}=\mathbb{L}_{k+1}/(\ell/L)\mathbb{Z}^3$, with the charge $e_{k+1}$.
\item[(S-3)] \emph{Scale free, $k\ge1$.}
An integer $k\ge1$; $\ell\in\mathrm{Cl}$; $\mu\in(0,1]$; $e_k\in(0,e_0]$.
No relation among them is assumed.
\item[(S-4)] \emph{Flow.}
$\bar{m}\in(0,1]$; $e>0$ with $e\cdot h^{\frac{1}{2}}\le e_1$; $R\in\mathrm{Tor}(e)$; an integer $N\ge N_{\min}(R,e)$; $\eta=L^{-N}$ and $n=R/\eta$;
the numbers $a_k$, $e_k$, $\mu_k$, $\ell_k$, $n_k$ and $k_{*}(R)$ are those of Sections~I.6.1 and~I.6.2.
\end{description}

\emph{Level $0$.}
The settings (S-1) and (S-2) include $k=0$.
There $\eta_0=1$;
the generator $\psi$ is indexed by $T^{\mathrm{f}}_0$, with the weight $1/48$;
every pair has $d_X\ge1$;
and a pair inside a cluster has $d_X=1$ and carries no phase ((I.L0.14)).
The clauses of the Propositions~16.1--18.6 and~48.1--48.5 include $k=0$, with no change.
In the flow, at $k=0$, these statements are applied only to the family $0$ ((I.L0.21)).
Proposition~19\,(b)--(f) and Proposition~20.1 concern $k\ge1$ only.
The level $0$ is treated separately in Proposition~15\,(n), Proposition~19\,(a), Proposition~20.2\,(iii0), (iv), Proposition~20.3 and Proposition~48.6.

\emph{Tori.}
All sides lie in $\mathrm{Cl}$.
The properties of the sides on which the Propositions~15--20.3 and 48.1--48.6 depend are these:
$\ell$ is a multiple of $LM$, $\ell\ge 13LM$, and $L$ and $M$ are odd.
The oddness of $\ell/(LM)$ is needed only for the assertion that $\ell/L=M\cdot\left(\ell/(LM)\right)$ is an odd multiple of $M$,
which is made in Proposition~16.1\,(d) and in the part \emph{Scaling} of Proposition~20.2.

\subsubsection*{G2.0.5. Recalled definitions}

\begin{itemize}
\item[(G2.r1)]
A union $Z$ of $M$-cubes is identified with the set of the photon sites that it contains.
A path lies in $Z$ if each of its sites lies in $Z$.
A bond meets $Z$ if at least one of its end points lies in $Z$;
it is at distance at most $1$ from $Z$ if each of its two end points has sup-distance at most $1$ from $Z$.
A fermion site $x=(c,g)$ lies over $Z$ if and only if $c\in Z$ (I.POS).
\item[(G2.r2)]
For a photon site $\tilde{c}$ and a number $r$ with $0<2r+\eta_k<\ell$ put
\[
K(\tilde{c},r):=\left\{c\in T_k:d(c,\tilde{c})\le r\right\}.
\]
This set does not wrap in the sense of Section~I.1.1.
For $c\in K(\tilde{c},r)$ the symbol $c-\tilde{c}$ denotes the vector of $\mathbb{L}_k$ of sup norm at most $r$ that represents the difference.
\item[(G2.r3)]
A polymer $X\in D_k$ is called small if $d_M(X)\le L$, and large otherwise.
This is the definition; two facts are attached to it.
\begin{itemize}
\item[(i)] A small polymer has a lift $\tilde{X}$ (Section~I.1.1), unique up to a translation by a vector of $\ell\mathbb{Z}^3$, whose sup-diameter is at most $M\cdot(L+2)$;
it is the lift of \cite[Theorem~13.2(t10)]{QEDI};
and $M\cdot(L+2)+M+2\eta_k<\ell$.
\item[(ii)] A polymer of $T_k$ that has a lift $\tilde{X}$ is small if and only if $d_M(\tilde{X})\le L$, where $d_M(\tilde{X})$ is formed with nets in $\mathbb{R}^3$, as in \cite[Theorem~13.2(t10)]{QEDI}.
\end{itemize}
Both facts follow from the notion of a lift in Section~I.1.1, from the clauses (t7) and (t10) of \cite[Theorem~13.2]{QEDI}, and from $\ell\ge13LM$;
they are proved in Chapter~16, Lemma~0.2\,(b)--(d), with the fifth inequality of Chapter~16, (Ar3).
\item[(G2.r4)]
The three norms of a bond function $A$ are $|A|_\infty$, $|\partial A|_\infty$ and $|\delta_\alpha\partial A|_\infty$ (Section~I.5.3; \cite[Section~5.3.1]{QEDI}).
On a subset $Z$ of sites they are formed with the bonds, the difference quotients and the H\"older differences inside $Z$.
Here a bond lies inside $Z$ if both of its end points belong to $Z$;
the quotient $(\partial_\nu A_\mu)(c)$ lies inside $Z$ if the four sites $c$, $c+\eta_k\cdot e_\mu$, $c+\eta_k\cdot e_\nu$ and $c+\eta_k\cdot e_\mu+\eta_k\cdot e_\nu$ belong to $Z$;
and a H\"older difference lies inside $Z$ if both of its quotients do.
The thresholds of $R_k$ are
\[
\mathrm{t1}_k:=e_k^{-\frac{3}{4}+\varepsilon},\qquad\mathrm{t2}_k:=e_k^{-\frac{3}{4}+2\varepsilon},\qquad\mathrm{t3}_k:=e_k^{-\frac{3}{4}+3\varepsilon},
\]
and those of $\mathfrak{X}_k$ are
\[
\mathrm{td1}_k:=e_k^{-1+2\varepsilon},\qquad\mathrm{td2}_k:=e_k^{-1+3\varepsilon}.
\]
\item[(G2.r5)]
The objects defined here concern two levels.
The homomorphism $S_L$, the functions $\Lambda_L$ and the indices $L\xi$, $L\boldsymbol{\zeta}$, together with the objects of the level $k+1$ that enter them, such as the torus $T_{k+1}$ and the generators $\psi$, $\chi$ of the level $k+1$, require only that $\ell/L$ is a multiple of $M$ and that $\ell/L>2M$ (Proposition~16.1\,(d); Section~I.3.7).
Let $S_L$ be the homomorphism of the algebra $\mathrm{Gr}[\psi,\chi]$ of the level $k$ into the algebra $\mathrm{Gr}[\psi,\chi]$ of the level $k+1$ with
\[
S_L\psi(x,\beta,\omega):=L^{-1}\cdot\psi(x/L,\beta,\omega),\qquad S_L\chi(x,x',\beta,\omega):=L^{-1-\alpha}\cdot\chi(x/L,x'/L,\beta,\omega),
\]
where $x/L:=(c/L,g)$ for $x=(c,g)$.
Then $S_LF(\psi,\chi)=F(\psi_L,\chi_L)$ (Section~I.3.7).
The reblocked family $LE$ of a family $E$ of level $k$ (Section~I.5.6) involves, in addition, the polymers $D_{k+1}$ and the domains $\tilde{R}_{k+1}$, $\mathfrak{X}_{k+1}$ of the level $k+1$;
these objects, and hence $LE$, require that $\ell/L\in\mathrm{Cl}$, which holds in the setting (G2.S-2) (Proposition~16.1\,(d); Section~I.5.7).
For $\ell/L\in\mathrm{Cl}$, $Y\in D_{k+1}$ and $(\mathscr{A},\mathscr{C})\in\tilde{R}_{k+1}\times\mathfrak{X}_{k+1}$ one has
\[
(LE)(Y;\mathscr{A},\mathscr{C})=\sum_{X:\bar{X}^{L}=LY}S_L\left[E(X;\mathscr{A}_L,\mathscr{C}_L)\right].
\]
For a function $\Lambda$ on the sites of $T_{k+1}$ put $\Lambda_L(c):=L^{\frac{1}{2}}\cdot\Lambda(c/L)$.
For indices $\xi=(c,g,\beta,\omega)$ and $\boldsymbol{\zeta}=((c,g),(c',g'),\beta,\omega)$ of the level $k+1$ put
\[
L\xi:=(Lc,g,\beta,\omega),\qquad L\boldsymbol{\zeta}:=((Lc,g),(Lc',g'),\beta,\omega).
\]
\item[(G2.r6)]
For $\gamma>0$ let
\[
c(\gamma):=\sum_{z\in\mathbb{Z}^3}\exp\left(-\gamma\cdot|z|\right).
\]
On every torus this number dominates $\sup_y\sum_{y'}\exp\left(-\gamma\cdot d(y,y')\right)$, and $c(\gamma)\le\left(\frac{12}{\gamma}\right)^3$ for $\gamma\le6$ \cite[Section~13.1.6]{QEDI}.
\item[(G2.r7)]
For $F\in\mathrm{Gr}[\psi,\chi]$ the part of $F$ of bidegree $(n,m)$ is denoted by $F^{[n,m]}$, and one has $\|F\|_{\mathsf{w}}=\sum_{n,m}\|F^{[n,m]}\|_{\mathsf{w}}$ (Proposition~15\,(n)).
The matrix notation $E_{20}(X;\cdot\,;x,x')$ is that of Section~I.5.1.
By (N0), that is by Proposition~15\,(a), and since every monomial has as many barred as unbarred legs, the part of bidegree $(2,0)$ of $E(X)$ and its norm are given by (Proposition~15\,(n))
\begin{equation}\tag{G2.0.1}
\begin{aligned}
E^{[2,0]}(X)&=\int\mathrm{d}x\int\mathrm{d}x'\,\bar{\psi}(x)^T\cdot E_{20}(X;x,x')\cdot\psi(x'),\\
\frac{\mathsf{w}_1^2}{2}\cdot\left\|E_{20}(X)\right\|&=\left\|E^{[2,0]}(X)\right\|_{\mathsf{w}}=\mathsf{w}_1^2\cdot\sum_{\beta,\beta'}\int\mathrm{d}x\int\mathrm{d}x'\,\left|E_{20}(X;x,x')_{\beta\beta'}\right|.
\end{aligned}
\end{equation}
\item[(G2.r8)]
A family of level $k$ (with slot) is a family
\[
E=\left\{E(X;\mathscr{A},\mathscr{C}):X\in D_k,\ (\mathscr{A},\mathscr{C})\in\tilde{R}_k\times\mathfrak{X}_k\right\}
\]
of elements of the algebra $\mathrm{Gr}[\psi,\chi]$ of the level $k$.
Such a family is said to have (I.k2) if it has the property stated in that clause of Section~I.5.4, and in the same way for the other clauses (I.k1)--(I.k7).
It is called balanced if every monomial of every $E(X;\mathscr{A},\mathscr{C})$ has as many barred as unbarred legs.
\end{itemize}

\subsubsection*{G2.0.6. Notation}

The following notation is used in Sections~G2.0.4--G2.0.6 and in the Propositions~15--20.3 and 48.1--48.6;
for each letter the letter of \cite[Section~13]{QEDI} in whose place it stands is indicated.
A vector of weights of a norm is written $\mathsf{w}=(\mathsf{w}_1,\dots,\mathsf{w}_p)$, the radius of a contour is written $\mathrm{rad}$, and a bound for row sums is written $\mathrm{Cr}$;
in \cite[Theorem~13.1]{QEDI} these are $g=(g_1,\dots,g_p)$, $\rho$, and the number $C_0$ of \cite[Proposition~5.57(N7)(iv)]{QEDI}, an arbitrary number $C_0>0$ in that clause.
The letter $\rho$ is reserved for the Wilson parameter (Section~I.0.5).
The generators $W$ are indexed by an arbitrary finite set $\mathrm{Mo}$, whose elements $\mathrm{mo}$ are called modes, instead of the set of the pairs $(y,\beta)$ of \cite[Section~13.2]{QEDI}.
The single weight of the norm on $\mathrm{Gr}[\Psi]$ is written $\mathrm{aw}$, as in $|F|_{\mathrm{aw}}$, and the numbers $a_1$, $a_2$ of \cite[Theorem~13.17]{QEDI} are written $\mathrm{aw}_1$, $\mathrm{aw}_2$;
the letter $a$ denotes the external connection, and $a_k=L^{k-N}$ (Section~I.6.1).
The elements $g$ and $g_c$ of $\mathbb{G}_k$, the automorphism $a_g$ and the charge conjugation $\mathbf{C}$ of \cite[Sections~13.1.5 and~13.4]{QEDI} are written $\tilde{g}$, $\tilde{g}_z$, $a_{\tilde{g}}$ and $\mathrm{Cg}$.
A constant $4\times4$ matrix, written $N$ in \cite[Sections~13.4 and~13.5]{QEDI}, is written $\mathrm{Mx}$, since $N$ is the index of the spacing $\eta=L^{-N}$.
The sites $y$, $y_0$ and $x$, $x_0$ of \cite[Section~13.5]{QEDI} are written $x$, $\tilde{x}$, as in Section~I.5.6;
the transports to the mean and the remainders $\mathbf{T}$, $\mathbf{T}^{-}$, $\operatorname{Rm}$, $\operatorname{Rm}^{-}$ defined there are written $\mathrm{tm}$, $\mathrm{tm}^{-}$, $\mathrm{rm}$, $\mathrm{rm}^{-}$.
The abbreviations $c'=3c_X$ of \cite[Section~13.1.6]{QEDI} and $x_k=c'\cdot L\cdot M\cdot e_k^{\varepsilon}$ of \cite[Theorem~13.7]{QEDI} are not used;
the corresponding expressions are written out, $c'$ denoting a photon site.
The base site $x_0$, the radii $\rho=2LM$ and $\rho'$ and the constant fields $c_1$, $c_2$ of \cite[Theorem~13.11]{QEDI}, and the radius $\rho_t$ of \cite[Theorem~13.15]{QEDI}, are written $\tilde{c}$, $\mathrm{rad}=2LM$, $\mathrm{rad}'$, $\mathrm{kc}_1$, $\mathrm{kc}_2$ and $\mathrm{rad}_t$.
The thresholds $x_{1,k}$, $x_{2,k}$ of $\mathfrak{X}_k$, introduced in \cite[Section~13.1.1]{QEDI}, are written $\mathrm{td1}_k$, $\mathrm{td2}_k$, as in (G2.r4), and the numbers $v_1$, $w_1$ of \cite[Lemma~13.8]{QEDI} are written out.
The subalgebra $\mathfrak{A}_k(Z)$ and the set $\mathrm{Dm}$ of \cite[Definition~13.12]{QEDI} are written $\mathrm{Gr}_k[Z]$ and $\mathrm{Dc}$;
sets of sites that \cite[Section~13]{QEDI} writes $W$, $W'$ are written $Z$, $Z'$, the letter $W$ being reserved for the integration generators.
The matrices $P^{\boldsymbol{\sigma}}_{\mu}$, the operators $T^{\boldsymbol{\sigma}}_{\mu}(A)$ and the formula for $D_A$ in terms of them in \cite[Section~13.10]{QEDI} are replaced by the hopping matrices $M_i(x)$, the hops $R_i(A)$ and (I.2.3), (I.2.4).
In \cite[Section~13.11.1]{QEDI} the complex mass $z$, the operator $R_z$, the numbers $\mathring{a}$, $\mathring{a}'$, $N_f$ and the maps $U$, $U_0$ are replaced by $\mathrm{zc}$, $\mathrm{Rz}$, $b_{L}$, $b_{k+1}/L$, $(\ell/\eta_k)^3$ and the scalings $f\mapsto f_L$, $v\mapsto v_L$;
the operator $D^{0}_{k+1}$ of \cite[Lemma~13.19]{QEDI} is $L\cdot D^{0}_{k+1}(B)$, where $D^{0}_{k+1}(B)$ is normalised as in \cite[Theorem~14.21]{QEDI}.
Finally, the constants $C_F$, $\gamma_F$ and $C_0^{\mathrm{unit}}$ of \cite[Section~13.1.6]{QEDI} are written $C^{F}$, $\gamma^{F}$ and $4b_{+}+16\,C^{F}\cdot c(\gamma^{F})$.

\subsection*{15. Grassmann algebras, norms, Gaussian integration on an index set with weights}

\subsubsection*{Proposition 15}

\begin{named}{Proposition 15}
The clauses (a)--(m) are stated in the setting (G2.S-abs), that is, without any datum of the scheme.
The data of (n) are an integer $k\ge0$ and a side $\ell$ (Section~I.1.1), with which the index sets and the weights of Section~I.3.6 are formed;
the sentence of (n) on row sums involves, besides them, a number $e_k>0$, a complex bond function $B$ on $T_k$, and a linear map $H$ of $\mathfrak{S}(1,\ell)$ (for $k=0$: of $\mathfrak{S}^{\mathrm{ch}}(1,\ell)$) to $\mathfrak{S}^{\mathrm{ch}}(\eta_k,\ell)$, all three arbitrary;
the H\"older quotient $\delta_{\alpha,B}$ is that of (I.1.7), formed with $e_k$.

\emph{Data.}
\begin{itemize}
\item[(D1)]
An integer $p\ge0$, and finite, totally ordered index sets $\mathrm{Ind}_1,\dots,\mathrm{Ind}_p$ with weights $\nu_i(t)>0$, $t\in\mathrm{Ind}_i$.
For a function $f$ on $\mathrm{Ind}_i$ put $\int\mathrm{d}t\,f(t):=\sum_{t\in\mathrm{Ind}_i}\nu_i(t)\cdot f(t)$, and for $I\subset\mathrm{Ind}_i$ put $\nu_i(I):=\prod_{t\in I}\nu_i(t)$.
The Grassmann algebra $\mathrm{Gr}[\phi_1,\dots,\phi_p]$ over $\mathbb{C}$ has the independent generators $\phi_i(t)$, $t\in\mathrm{Ind}_i$, $1\le i\le p$.
The products $\phi_i(I)$, the coefficients $c_F(I_1,\dots,I_p)$ and the weight vectors $\mathsf{w}=(\mathsf{w}_1,\dots,\mathsf{w}_p)$, $\mathsf{w}_i>0$, are those of Section~I.5.1, and the norm is (I.5.1):
\[
\|F\|_{\mathsf{w}}=\sum_{I_1,\dots,I_p}\left|c_F(I_1,\dots,I_p)\right|\cdot\mathsf{w}_1^{\#|I_1|}\cdots\mathsf{w}_p^{\#|I_p|}.
\]
\item[(D2)]
A finite non-empty set $\mathrm{Mo}$, whose elements are called modes and are written $\mathrm{mo}$, $\mathrm{mo}'$, with weights $\nu(\mathrm{mo})>0$;
and a further family of independent generators $W(u)$, $u=(\mathrm{mo},\omega)\in\mathrm{Mo}\times\{0,1\}$.
The family $W$ is a family in the sense of (D1):
the set $\mathrm{Mo}\times\{0,1\}$ carries a total order, chosen arbitrarily and then kept fixed, and in $\mathrm{Gr}[\phi,W]$ the generators $W(u)$ come after the generators $\phi_i(t)$.
The norms, $\int\mathrm{D}W$, $\mathrm{d}\mu_I(W)$ and $\mathrm{d}\mu_{\mathrm{Cov}}$ do not depend on these orders;
the numbers $\mathrm{sn}(J)$ and $m(J)$ of (e) depend only on the order of the elements of $J$ among themselves.
For empty $\mathrm{Mo}$ one has $\mathrm{Gr}[\phi,W]=\mathrm{Gr}[\phi]$ and $\mathrm{Top}=1$, the map $\int\mathrm{D}W$ is the identity, and all assertions below hold with an arbitrary number $\mathsf{w}_W>0$, except the sentence of (f) according to which $\mathsf{w}_W$ cannot be lowered.
Put
\[
W(\mathrm{mo}):=W(\mathrm{mo},0),\qquad\bar{W}(\mathrm{mo}):=W(\mathrm{mo},1),\qquad\nu(\mathrm{mo},\omega):=\nu(\mathrm{mo}),
\]
and for a subset $J$ of $\mathrm{Mo}\times\{0,1\}$ put $\nu^{-\frac{1}{2}}(J):=\prod_{u\in J}\nu(u)^{-\frac{1}{2}}$;
the set $J$ is called paired if it contains $(\mathrm{mo},0)$ exactly when it contains $(\mathrm{mo},1)$.
For complex functions $f'$, $f$ on $\mathrm{Mo}$ let $\langle f',f\rangle_\nu:=\sum_{\mathrm{mo}}\nu(\mathrm{mo})\cdot f'(\mathrm{mo})\cdot f(\mathrm{mo})$.
A linear map $O$ of the space of the complex functions on $\mathrm{Mo}$ has the kernel $O(\mathrm{mo},\mathrm{mo}')$, given by
\[
(Of)(\mathrm{mo})=\sum_{\mathrm{mo}'}\nu(\mathrm{mo}')\cdot O(\mathrm{mo},\mathrm{mo}')\cdot f(\mathrm{mo}'),
\]
the determinant $\det O$ of the linear map, and the transpose $O^T$ with respect to $\langle\,\cdot\,,\cdot\,\rangle_\nu$.
Put
\[
(OW)(\mathrm{mo}):=\sum_{\mathrm{mo}'}\nu(\mathrm{mo}')\cdot O(\mathrm{mo},\mathrm{mo}')\cdot W(\mathrm{mo}'),\qquad\left\langle\bar{W},OW\right\rangle_\nu:=\sum_{\mathrm{mo}}\nu(\mathrm{mo})\cdot\bar{W}(\mathrm{mo})\cdot(OW)(\mathrm{mo}).
\]
In the same way, for families $\mathrm{el}=(\mathrm{el}(\mathrm{mo}))$ and $\mathrm{el}'=(\mathrm{el}'(\mathrm{mo}))$ of elements of a Grassmann algebra,
\[
(O\,\mathrm{el})(\mathrm{mo}):=\sum_{\mathrm{mo}'}\nu(\mathrm{mo}')\cdot O(\mathrm{mo},\mathrm{mo}')\cdot\mathrm{el}(\mathrm{mo}'),\qquad\left\langle\mathrm{el}',\mathrm{el}\right\rangle_\nu:=\sum_{\mathrm{mo}}\nu(\mathrm{mo})\cdot\mathrm{el}'(\mathrm{mo})\cdot\mathrm{el}(\mathrm{mo}),
\]
the factors being taken in this order;
and for the operators $H^{(i)}$ of (l)\,(iv)
\[
(H^{(i)}\Psi)(s):=\sum_{u}H^{(i)}(s,u)\cdot\nu(u)\cdot\Psi(u).
\]
\item[(D3)]
Let
\[
\mathrm{Top}:=\prod_{\mathrm{mo}\in\mathrm{Mo}}W(\mathrm{mo})\cdot\bar{W}(\mathrm{mo});
\]
its factors are even, so that their order is immaterial.
$\int\mathrm{D}W$ is the linear map of $\mathrm{Gr}[\phi,W]$ to $\mathrm{Gr}[\phi]$ that sends $\phi_1(I_1)\cdots\phi_p(I_p)\cdot W(J)$ to $\phi_1(I_1)\cdots\phi_p(I_p)$, multiplied by $\prod_{\mathrm{mo}}\nu(\mathrm{mo})^{-1}$ and by the coefficient of $\mathrm{Top}$ in $W(J)$.
Further,
\[
\int\mathrm{d}\mu_I(W)\,F:=\int\mathrm{D}W\,\exp\left(-\left\langle\bar{W},W\right\rangle_\nu\right)\cdot F,
\]
and for an invertible linear map $\mathrm{Cov}$
\[
\int\mathrm{d}\mu_{\mathrm{Cov}}(W)\,F:=(\det\mathrm{Cov})\cdot\int\mathrm{D}W\,\exp\left(-\left\langle\bar{W},\mathrm{Cov}^{-1}W\right\rangle_\nu\right)\cdot F;
\]
finally $\mathsf{w}_W:=\max_{\mathrm{mo}}\nu(\mathrm{mo})^{-\frac{1}{2}}$.
This is the integral of (I.3.5).
\item[(D4)]
For an operator $O$ from the functions on an index set $(\mathrm{Ind},\nu)$ to the functions on an index set $\mathrm{Ind}'$, acting by $(Of)(t)=\sum_sO(t,s)\cdot f(s)\cdot\nu(s)$, put
\[
|O|_{1,\infty}:=\sup_t\sum_s\left|O(t,s)\right|\cdot\nu(s),\qquad|f|_\infty:=\sup_t|f(t)|.
\]
\item[(D5)]
$\oint_{\mathrm{rad}}f$ and $\oint^{(2)}_{\mathrm{rad}}f$ are the contour integrals of Section~I.5.2.
For a family $\Psi$ on $\mathrm{Mo}\times\{0,1\}$ with the weights $\nu$ and a number $\mathrm{aw}>0$, $|F|_{\mathrm{aw}}$ is the norm (I.5.1) of $F\in\mathrm{Gr}[\Psi]$ with the single weight $\mathrm{aw}$.
$\int\mathrm{D}\Psi$, $\mathrm{d}\mu_I(\Psi)$ and $\mathrm{d}\mu_{\mathrm{Cov}}(\Psi)$ are the objects of (D3), formed with the family $\Psi$ in the place of $W$.
\end{itemize}

\emph{Statement.}
\begin{itemize}
\item[(a)]
\emph{Kernels} (N0).
Let
\begin{multline*}
F=\sum_{n}\left(\prod_{i=1}^{p}\frac{1}{n_i!}\right)\int\mathrm{d}t\,K_n\left(t^1_1,\dots,t^1_{n_1};\dots;t^p_1,\dots,t^p_{n_p}\right)\\
\cdot\phi_1(t^1_1)\cdots\phi_1(t^1_{n_1})\cdots\phi_p(t^p_1)\cdots\phi_p(t^p_{n_p})
\end{multline*}
with arbitrary complex kernels $K_n$;
here $\int\mathrm{d}t$ integrates every argument $t^i_j$ over $\mathrm{Ind}_i$ with the weight $\nu_i$, and the sum extends over all $n=(n_1,\dots,n_p)\in\{0,1,2,\dots\}^p$, the terms with $n_i>\#|\mathrm{Ind}_i|$ for some $i$ being zero.
Then
\[
\|F\|_{\mathsf{w}}\le\sum_{n}\left(\prod_{i=1}^{p}\frac{\mathsf{w}_i^{n_i}}{n_i!}\right)\int\mathrm{d}t\,\left|K_n\right|,
\]
where the right side may be $+\infty$, with equality if every $K_n$ is antisymmetric in each group of arguments.
Every $F\in\mathrm{Gr}[\phi_1,\dots,\phi_p]$ can be written in this form with kernels that are antisymmetric in each group of arguments, and such kernels are determined by $F$:
\[
K_n(I_1;\dots;I_p)\cdot\prod_{i=1}^{p}\nu_i(I_i)=c_F(I_1,\dots,I_p)\qquad\left(\#|I_i|=n_i,\ 1\le i\le p\right),
\]
the elements of each $I_i$ being inserted in increasing order.
A single generator $\phi_i(t)$ has the norm $\mathsf{w}_i$;
the norm does not contain the weights $\nu_i$.
\item[(b)]
\emph{Products} (N1).
One has $\|F\cdot G\|_{\mathsf{w}}\le\|F\|_{\mathsf{w}}\cdot\|G\|_{\mathsf{w}}$.
The even elements form a commutative Banach algebra with unit, and every even $F$ satisfies
\[
\left\|\exp F-1\right\|_{\mathsf{w}}\le\exp\left(\|F\|_{\mathsf{w}}\right)-1.
\]
\item[(c)]
\emph{Row sums.}
For an operator $O$ as in (D4), $|O|_{1,\infty}=\sup\left\{|Of|_\infty:|f|_\infty\le1\right\}$.
\item[(d)]
\emph{Linear substitutions} (N2).
Let $\theta_1,\dots,\theta_q$ be families of generators on index sets $(\mathrm{Ind}'_j,\nu'_j)$;
let $O_{ij}$ be operators from the functions on $\mathrm{Ind}'_j$ to the functions on $\mathrm{Ind}_i$, and $C_{ij}\ge0$ numbers with $|O_{ij}|_{1,\infty}\le C_{ij}$;
and let $F\mapsto F'$ be the homomorphism of $\mathrm{Gr}[\phi]$ to $\mathrm{Gr}[\theta]$ with
\[
\phi_i(t)\mapsto\sum_{j}(O_{ij}\theta_j)(t)=\sum_{j}\sum_{s}O_{ij}(t,s)\cdot\nu'_j(s)\cdot\theta_j(s).
\]
If $\sum_jC_{ij}\cdot\mathsf{w}'_j\le\mathsf{w}_i$ for all $i$, then $\|F'\|_{\mathsf{w}'}\le\|F\|_{\mathsf{w}}$.
The same holds on subalgebras with the row sums of the rows that occur:
if $\mathrm{Ind}^0_i\subset\mathrm{Ind}_i$ for every $i$ and $F$ lies in the subalgebra generated by the $\phi_i(t)$, $t\in\mathrm{Ind}^0_i$, then $|O_{ij}|_{1,\infty}$ may be replaced by the supremum of the row sums over the rows $t\in\mathrm{Ind}^0_i$ alone.
\item[(e)]
\emph{Moments} (m)$^{\mathrm{w}}$, (I.B1)--(I.B4).
For $J\subset\mathrm{Mo}\times\{0,1\}$ put $m(J):=\int\mathrm{d}\mu_I(W)\,W(J)$.
If $J$ is not paired, then $m(J)=0$.
If $J$ is paired, then
\[
W(J)=\mathrm{sn}(J)\cdot\prod_{(\mathrm{mo},0)\in J}W(\mathrm{mo})\cdot\bar{W}(\mathrm{mo}),\qquad m(J)=\mathrm{sn}(J)\cdot\prod_{(\mathrm{mo},0)\in J}\nu(\mathrm{mo})^{-1},
\]
with a sign $\mathrm{sn}(J)\in\{1,-1\}$.
Consequently all moments are real, and $|m(J)|=\nu^{-\frac{1}{2}}(J)$ for paired $J$.
In particular,
\begin{gather*}
\int\mathrm{d}\mu_I(W)\,1=1,\qquad\int\mathrm{d}\mu_I(W)\,W(\mathrm{mo})\cdot\bar{W}(\mathrm{mo}')=\frac{\mathbf{1}[\mathrm{mo}=\mathrm{mo}']}{\nu(\mathrm{mo})},\\
\int\mathrm{d}\mu_I(W)\,W(\mathrm{mo})\cdot W(\mathrm{mo}')=\int\mathrm{d}\mu_I(W)\,\bar{W}(\mathrm{mo})\cdot\bar{W}(\mathrm{mo}')=0.
\end{gather*}
For every linear map $O$,
\[
\int\mathrm{D}W\,\exp\left(-\left\langle\bar{W},OW\right\rangle_\nu\right)=\det O.
\]
\item[(f)]
\emph{Integration} (N3)$^{\mathrm{w}}$.
Let
\[
F=\sum_{I,J}c_F(I_1,\dots,I_p,J)\cdot\phi_1(I_1)\cdots\phi_p(I_p)\cdot W(J)\in\mathrm{Gr}[\phi_1,\dots,\phi_p,W].
\]
Then
\[
\int\mathrm{d}\mu_I(W)\,F=\sum_{I}\left(\sum_{J}c_F(I_1,\dots,I_p,J)\cdot m(J)\right)\cdot\phi_1(I_1)\cdots\phi_p(I_p),
\]
and
\begin{multline*}
\left\|\int\mathrm{d}\mu_I(W)\,F\right\|_{(\mathsf{w}_1,\dots,\mathsf{w}_p)}\le\sum_{I,J}\left|c_F(I_1,\dots,I_p,J)\right|\cdot\mathsf{w}_1^{\#|I_1|}\cdots\mathsf{w}_p^{\#|I_p|}\cdot\nu^{-\frac{1}{2}}(J)\\
\le\|F\|_{(\mathsf{w}_1,\dots,\mathsf{w}_p,\mathsf{w}_W)}.
\end{multline*}
The weight $\mathsf{w}_W$ cannot be lowered:
for a mode $\mathrm{mo}$ with $\nu(\mathrm{mo})^{-\frac{1}{2}}=\mathsf{w}_W$ and for $F=W(\mathrm{mo})\cdot\bar{W}(\mathrm{mo})$ both inequalities are equalities.
The integral does not see the modes that do not occur.
Namely, let $\mathrm{Mo}^0$ be a subset of $\mathrm{Mo}$, with the weights $\nu(\mathrm{mo})$, and let $\mathrm{Mo}^0\times\{0,1\}$ carry the induced order.
For $J\subset\mathrm{Mo}^0\times\{0,1\}$ the number $m(J)$ formed with $\mathrm{Mo}^0$ equals the number $m(J)$ formed with $\mathrm{Mo}$;
and for $F$ in the subalgebra generated by the $\phi_i(t)$ and the $W(u)$, $u\in\mathrm{Mo}^0\times\{0,1\}$, the element $\int\mathrm{d}\mu_I(W)\,F$ of $\mathrm{Gr}[\phi]$ is the same, whether it is formed with $\mathrm{Mo}$ or with $\mathrm{Mo}^0$;
for empty $\mathrm{Mo}^0$ it equals $F$.
\item[(g)]
\emph{Factorisation} (N4).
Let $F_1,F_2\in\mathrm{Gr}[\phi,W]$, and suppose that no mode $\mathrm{mo}$ occurs, as $W(\mathrm{mo})$ or as $\bar{W}(\mathrm{mo})$, both in a monomial of $F_1$ and in a monomial of $F_2$.
Then
\[
\int\mathrm{d}\mu_I(W)\,\left(F_1\cdot F_2\right)=\left(\int\mathrm{d}\mu_I(W)\,F_1\right)\cdot\left(\int\mathrm{d}\mu_I(W)\,F_2\right).
\]
In the scheme the modes are the pairs $\mathrm{mo}=(t,\beta)$;
there the hypothesis is satisfied if the set of the points $t$ at which $F_1$ contains a generator $W(t,\cdot,\cdot)$ is disjoint from the corresponding set for $F_2$.
\item[(h)]
\emph{Invariance} (N4$'$).
For an invertible linear map $U$ of the functions on $\mathrm{Mo}$ let $s_U$ be the homomorphism with
\[
W(\mathrm{mo})\mapsto(UW)(\mathrm{mo}),\qquad\bar{W}(\mathrm{mo})\mapsto\left((U^{-1})^T\bar{W}\right)(\mathrm{mo})
\]
that fixes the other generators.
Then, for all $F\in\mathrm{Gr}[\phi,W]$,
\[
\left\langle(U^{-1})^T\bar{W},UW\right\rangle_\nu=\left\langle\bar{W},W\right\rangle_\nu,\qquad\int\mathrm{D}W\,s_U\left(F\right)=\int\mathrm{D}W\,F,\qquad\int\mathrm{d}\mu_I(W)\,s_U\left(F\right)=\int\mathrm{d}\mu_I(W)\,F.
\]
\item[(i)]
\emph{Wick's theorem.}
Let $v_1,\dots,v_r$ lie in the linear span of the generators $W(u)$, $u\in\mathrm{Mo}\times\{0,1\}$, and let $F_0\in\mathrm{Gr}[\phi]$.
If $r$ is even, then $\int\mathrm{d}\mu_I(W)\,F_0\cdot v_1\cdots v_r$ is the product of $F_0$ with the Pfaffian of the antisymmetric $r\times r$ matrix whose entries above the diagonal are the numbers $\int\mathrm{d}\mu_I(W)\,v_i\cdot v_j$, $i<j$;
for $r=0$ the Pfaffian is $1$.
If $r$ is odd, the integral vanishes.
\item[(j)]
\emph{Contour integrals} (N5).
Let $\mathrm{rad}>1$, and let $f$ be analytic in a neighbourhood of $\{t\in\mathbb{C}:|t|\le\mathrm{rad}\}$, with values in a normed space of finite dimension.
Then
\begin{align*}
\oint_{\mathrm{rad}}f&=f(1)-f(0),&\left\|\oint_{\mathrm{rad}}f\right\|&\le\frac{\sup_{|t|=\mathrm{rad}}\|f(t)\|}{\mathrm{rad}-1},\\
\oint^{(2)}_{\mathrm{rad}}f&=f(1)-f(0)-f'(0),&\left\|\oint^{(2)}_{\mathrm{rad}}f\right\|&\le\frac{\sup_{|t|=\mathrm{rad}}\|f(t)\|}{\mathrm{rad}\cdot(\mathrm{rad}-1)}.
\end{align*}
Let $F\in\mathrm{Gr}[\phi]$, let $\theta=(\theta_1,\dots,\theta_p)$ be a family of generators on the index sets of $\phi$, with $\theta_i$ indexed by $\mathrm{Ind}_i$, let $F(\phi+\theta)$ be the image of $F$ under $\phi_i(t)\mapsto\phi_i(t)+\theta_i(t)$, and let the weights $\mathsf{w}'$ of $\phi$ and $\mathsf{w}''$ of $\theta$ satisfy $\mathsf{w}'_i+\mathrm{rad}\cdot\mathsf{w}''_i\le\mathsf{w}_i$ for all $i$.
Then
\[
\left\|F(\phi+\theta)-F(\phi)\right\|_{(\mathsf{w}',\mathsf{w}'')}\le\frac{\|F\|_{\mathsf{w}}}{\mathrm{rad}-1}.
\]
\item[(k)]
\emph{Monotonicity, degree} (N6).
If $0<\mathsf{w}'_i\le\mathsf{w}_i$ for all $i$, then $\|F\|_{\mathsf{w}'}\le\|F\|_{\mathsf{w}}$.
If, in addition, every monomial of $F$ has degree at least $r_i$ in $\phi_i$ for every $i$, then
\[
\|F\|_{\mathsf{w}'}\le\prod_{i}\left(\frac{\mathsf{w}'_i}{\mathsf{w}_i}\right)^{r_i}\cdot\|F\|_{\mathsf{w}}.
\]
\item[(l)]
\emph{The algebra $\mathrm{Gr}[\Psi]$} (S6c).
Let $F,G,f\in\mathrm{Gr}[\Psi]$ and $\mathrm{aw}>0$.
\begin{itemize}
\item[(i)] $|F\cdot G|_{\mathrm{aw}}\le|F|_{\mathrm{aw}}\cdot|G|_{\mathrm{aw}}$.
\item[(ii)] $\left|\int\mathrm{d}\mu_I(\Psi)\,F\right|\le|F|_{\mathsf{w}_W}$.
\item[(iii)] For every invertible $\mathrm{Cov}$,
\[
\int\mathrm{d}\mu_{\mathrm{Cov}}(\Psi)\,f(\bar{\Psi},\Psi)=\int\mathrm{d}\mu_I(\Psi)\,f(\bar{\Psi},\mathrm{Cov}\,\Psi),
\]
where $f(\bar{\Psi},\mathrm{Cov}\,\Psi)$ is the image of $f$ under the homomorphism with $\Psi(\mathrm{mo})\mapsto(\mathrm{Cov}\,\Psi)(\mathrm{mo})$ and $\bar{\Psi}(\mathrm{mo})\mapsto\bar{\Psi}(\mathrm{mo})$.
Moreover $\int\mathrm{d}\mu_{\mathrm{Cov}}(\Psi)\,1=1$;
the two-point function
\[
\int\mathrm{d}\mu_{\mathrm{Cov}}(\Psi)\,\Psi(\mathrm{mo})\cdot\bar{\Psi}(\mathrm{mo}')=\mathrm{Cov}(\mathrm{mo},\mathrm{mo}')
\]
is the kernel of $\mathrm{Cov}$;
and
\[
\int\mathrm{D}\Psi\,\exp\left(-\left\langle\bar{\Psi},\mathrm{Cov}^{-1}\Psi\right\rangle_\nu\right)\cdot F=\det\left(\mathrm{Cov}^{-1}\right)\cdot\int\mathrm{d}\mu_{\mathrm{Cov}}(\Psi)\,F.
\]
The first and the last identity of (iii) hold also for $f,F\in\mathrm{Gr}[\phi,\Psi]$, the homomorphism fixing the generators $\phi_i(t)$.
\item[(iv)] Let $E\in\mathrm{Gr}[\psi,\chi]$, where $\psi$ and $\chi$ are two families on arbitrary index sets;
let $H^{(1)}$, $H^{(2)}$ be operators from the functions on $\mathrm{Mo}\times\{0,1\}$ to the functions on the index set of $\psi$, resp.\ of $\chi$;
and let $\mathrm{Cr}\ge0$ be a number with $|H^{(i)}|_{1,\infty}\le\mathrm{Cr}$ for $i=1,2$.
Let $E(H^{(1)}\Psi,H^{(2)}\Psi)$ be the image of $E$ under $\psi\mapsto H^{(1)}\Psi$, $\chi\mapsto H^{(2)}\Psi$.
If $\mathrm{Cr}\cdot\mathrm{aw}\le\min(\mathsf{w}_1,\mathsf{w}_2)$, then
\[
\left|E\left(H^{(1)}\Psi,H^{(2)}\Psi\right)\right|_{\mathrm{aw}}\le\|E\|_{(\mathsf{w}_1,\mathsf{w}_2)}.
\]
\end{itemize}
\item[(m)]
\emph{Reduction to the weight one.}
Put $\mathrm{Wr}(u):=\nu(u)^{\frac{1}{2}}\cdot W(u)$, and $\mathrm{Wr}(\mathrm{mo}):=\mathrm{Wr}(\mathrm{mo},0)$, $\overline{\mathrm{Wr}}(\mathrm{mo}):=\mathrm{Wr}(\mathrm{mo},1)$.
Then
\[
\left\langle\bar{W},W\right\rangle_\nu=\sum_{\mathrm{mo}}\overline{\mathrm{Wr}}(\mathrm{mo})\cdot\mathrm{Wr}(\mathrm{mo}),
\]
and $\int\mathrm{D}W\,F$ is the coefficient of $\prod_{\mathrm{mo}}\mathrm{Wr}(\mathrm{mo})\cdot\overline{\mathrm{Wr}}(\mathrm{mo})$ in $F$.
Thus $\int\mathrm{D}W$ and $\mathrm{d}\mu_I(W)$ are the integral and the functional of \cite[Section~13.2]{QEDI}, written in the generators $\mathrm{Wr}$.
\item[(n)]
\emph{Values in the scheme.}
The index sets and the weights are those of Section~I.3.6:
the set $\{\xi\}$ carries the weight $\eta_k^3/48$, the set $\{\boldsymbol{\zeta}\}$ the weight $(\eta_k^3/48)^2$, and for $k\ge1$ the set $\{u\}$ has the modes $\mathrm{mo}=(y,\beta)$, with $\nu=1$ and $\mathsf{w}_W=1$.

\emph{Level $0$.}
The modes are $\mathrm{mo}=(x,\beta)$, $x\in T^{\mathrm{f}}_0$, with $\nu=1/48$ and $\mathsf{w}_W=\sqrt{48}\le7$.
One has
\begin{gather*}
\int\mathrm{d}\mu_I(W)\,W_{\beta}(x)\cdot\bar{W}_{\beta'}(x')=48\cdot\mathbf{1}[x=x',\ \beta=\beta'],\\
\int\mathrm{D}\Psi\,\exp\left(-\left\langle\bar{\Psi},O\Psi\right\rangle^{\mathrm{ch}}\right)=\det O\qquad\text{for every linear map $O$ of $\mathfrak{S}^{\mathrm{ch}}(1,\ell)$},\\
\int\mathrm{d}\mu_{\mathrm{Cov}}(\Psi)\,\Psi(x)\cdot\bar{\Psi}(x')^T=\mathrm{Cov}(x,x'),
\end{gather*}
the right side of the last identity being the kernel (I.1.4).

\emph{Row sums.}
In (l)\,(iv), and in every application of (d) with the substitution $\chi\mapsto\delta_{\alpha,B}H\Psi$, the row sums are suprema over all pairs of $\mathrm{Pz}_k$, the pairs inside a cluster included.

\emph{Levels $k\ge1$.}
For unit sites $y$, $y'$ one has
\begin{gather*}
\int\mathrm{d}\mu_I(W)\,W_{\beta}(y)\cdot\bar{W}_{\beta'}(y')=\mathbf{1}[y=y',\ \beta=\beta'],\\
\int\mathrm{D}\Psi\,\exp\left(-\left\langle\bar{\Psi},O\Psi\right\rangle\right)=\det O\qquad\text{for every linear map $O$ of $\mathfrak{S}(1,\ell)$},\\
\int\mathrm{d}\mu_{\mathrm{Cov}}(\Psi)\,\Psi(y)\cdot\bar{\Psi}(y')^T=\mathrm{Cov}(y,y').
\end{gather*}

\emph{The number $7$.}
At every level the two inequalities of (f) remain valid with $7$ in the place of $\mathsf{w}_W$:
\begin{multline*}
\left\|\int\mathrm{d}\mu_I(W)\,F\right\|_{(\mathsf{w}_1,\dots,\mathsf{w}_p)}\le\sum_{I,J}\left|c_F(I_1,\dots,I_p,J)\right|\cdot\mathsf{w}_1^{\#|I_1|}\cdots\mathsf{w}_p^{\#|I_p|}\cdot\nu^{-\frac{1}{2}}(J)\\
\le\|F\|_{(\mathsf{w}_1,\dots,\mathsf{w}_p,7)}.
\end{multline*}

\emph{The fundamental domain; the covering lattices} (every $k\ge0$).
A change of the fundamental domain of $\ell\mathbb{Z}^3$ multiplies $W(\mathrm{mo})$ and $\bar{W}(\mathrm{mo})$, and likewise $\Psi(\mathrm{mo})$ and $\bar{\Psi}(\mathrm{mo})$, by one and the same sign, and it multiplies $\psi(\xi)$ and $\chi(\boldsymbol{\zeta})$ by signs.
The following objects do not depend on the choice:
the norms (I.5.1);
$\mathrm{Top}$, $\langle\bar{W},W\rangle_\nu$, $\int\mathrm{D}W$ and $\mathrm{d}\mu_I(W)$;
the kernel of a linear map $O$ of $\mathfrak{S}^{\mathrm{ch}}(1,\ell)$ (for $k=0$), resp.\ of $\mathfrak{S}(1,\ell)$ (for $k\ge1$), which is antiperiodic in each of its two arguments;
the form $\langle\bar{\Psi},O\Psi\rangle$, the determinant $\det O$, the element $\int\mathrm{D}\Psi\,\exp\left(-\langle\bar{\Psi},O\Psi\rangle\right)\cdot F$, and $\mathrm{d}\mu_{\mathrm{Cov}}(\Psi)$.
The identities for $\int\mathrm{d}\mu_I(W)\,W\cdot\bar{W}$ and for $\int\mathrm{d}\mu_{\mathrm{Cov}}(\Psi)\,\Psi\cdot\bar{\Psi}^T$ hold at all points of the covering lattices, both sides being continued antiperiodically in each argument;
the right side of the first one is then the kernel of the identity map.

\emph{The identities of (G2.r7).}
For $F\in\mathrm{Gr}[\psi,\chi]$ one has $\|F\|_{\mathsf{w}}=\sum_{n,m}\|F^{[n,m]}\|_{\mathsf{w}}$.
For a balanced element $E$ of $\mathrm{Gr}[\psi,\chi]$ the display (G2.0.1) holds.
\end{itemize}
\end{named}

\begin{named}{Remark}
Proposition~15 corresponds to \cite[Proposition~5.57]{QEDI} and to \cite[Theorem~13.1]{QEDI}, taken together with \cite[Section~13.1.4]{QEDI}, with the definitions at the head of \cite[Section~13.2]{QEDI}, including the display (m) there, and with the paragraph on index sets that follows the proof of \cite[Theorem~13.1]{QEDI}.
Its clauses (a), (b), (d), (f), (g), (h), (j), (k) and (l) correspond, in this order, to the clauses (N0), (N1), (N2), (N3), (N4), (N4$'$), (N5), (N6) and (N7) of \cite[Proposition~5.57]{QEDI}.
Clause (e) corresponds to the display (m) of \cite[Section~13.2]{QEDI}, and clause (n) to the paragraph on index sets.
The content of \cite[Proposition~5.57]{QEDI} is recalled in \cite[Section~17.1.4]{QEDI}, and in the displays (Gr1)--(Gr4) of \cite[Section~20.2.4]{QEDI}.
The mathematical differences are the following.
\begin{itemize}
\item[(1)] The index set of $W$ is an arbitrary finite set with weights, and $\int\mathrm{D}W$ contains the factor $\prod_{\mathrm{mo}}\nu(\mathrm{mo})^{-1}$.
The normalisation $\int\mathrm{D}\Psi\,\exp\left(-\langle\bar{\Psi},O\Psi\rangle_\nu\right)=\det O$ is asserted for every finite index set with weights;
in this form it applies to the Gaussian integral $\mathrm{d}\mu_C$ of \cite[Lemma~20.9(a)]{QEDI}.
\item[(2)] The moment of a paired set is $\pm\prod_{(\mathrm{mo},0)\in J}\nu(\mathrm{mo})^{-1}$ (clause (e)), not $\pm1$ as in the display (m) of \cite[Section~13.2]{QEDI}.
\item[(3)] In the bound of (f) the weight $\mathsf{w}_W$ of the integrated family takes the place of the weight $1$ of (N3).
\item[(4)] In (g) the modes, that is the points of the index set, take the place of the unit sites of (N4).
\item[(5)] In (h) and in (l)\,(iii) the matrices of (N4$'$) and of (N7)\,(iii) are replaced by linear maps, with kernels and transposes taken with respect to $\langle\,\cdot\,,\cdot\,\rangle_\nu$;
the two-point function of $\mathrm{d}\mu_{\mathrm{Cov}}$ is the kernel of $\mathrm{Cov}$.
\item[(6)] In (l)\,(ii) the bound is $|F|_{\mathsf{w}_W}$, in the place of the bound $|F|_1$ of (N7)\,(ii).
\item[(7)] In (N7)\,(iv) of \cite[Proposition~5.57]{QEDI}, and in clause (iv) of the paragraph \emph{Norms on $\mathrm{Gr}[\Psi]$} of \cite[Theorem~13.1]{QEDI}, the bound on the row sums $\left|H^{(1)}\right|_{1,\infty}$ and $\left|H^{(2)}\right|_{1,\infty}$ is an arbitrary number $C_0>0$;
in (l)\,(iv) it is the arbitrary number $\mathrm{Cr}\ge0$ of the data.
Apart from the letter, $C_0$ denoting in this paper the constant of the line (C6), and apart from the admission of $\mathrm{Cr}=0$, there is no difference.
In (d) the bounds $C_{ij}$ are numbers $C_{ij}\ge0$.
\item[(8)] The identity (c) has no counterpart among the clauses of \cite[Proposition~5.57]{QEDI} and of \cite[Theorem~13.1]{QEDI};
it is the identity stated at the end of \cite[Section~13.1.4]{QEDI} and again in \cite[Section~17.1.4]{QEDI}.
\item[(9)] Wick's theorem (i) and the reality of the moments in (e) have no counterpart among the clauses of \cite[Proposition~5.57]{QEDI} and of \cite[Theorem~13.1]{QEDI};
both facts are stated in \cite[Section~17.1.4]{QEDI}, and Wick's theorem is applied in the proof of \cite[Lemma~17.20]{QEDI}.
\item[(10)] Clause (m) has no counterpart in \cite{QEDI}.
\end{itemize}
\end{named}

\subsection*{16. The group lemma, the low kernels, the extraction}

\subsubsection*{Proposition 16.1 (group lemma)}

\begin{named}{Proposition 16.1}[Group lemma]
Let $k\ge0$ be an integer, let $\ell\in\mathrm{Cl}$ be a side, and let $e_k>0$ be a number; $e_k$ enters only through $U_B$, through $G_{\Lambda}$ and through the domains.
Let $\tilde{g}=(S,a)\in\mathbb{G}_k$ and put $r:=r(S)$ (Sections~I.4.1--I.4.3).
No slot occurs.
\begin{itemize}
\item[(a1)] \emph{Photon lattice.}
The map $c\mapsto\tilde{g}c=rc+a$ is an isometry of $\mathbb{R}^3$ for the sup norm, for the $\ell^1$ norm and for the Euclidean norm, and
\[
\tilde{g}(c+\ell n)=\tilde{g}c+\ell\, rn\qquad\left(n\in\mathbb{Z}^3\right).
\]
Consequently it induces on each torus $L^{-i}\mathbb{Z}^3/\ell\mathbb{Z}^3$, $-1\le i\le k$, a bijection which preserves $d(\cdot,\cdot)$.
It carries a lift (Section~I.1.1) of a set $Z$ of sites onto a lift of $\tilde{g}Z$, and the shortest lift of a pair of sites at distance $<\ell/2$ onto the shortest lift of the image pair.
The block $\left\{c:|c-z|<\frac{1}{2}L^{-i}\right\}$ with centre $z\in L^{-i}\mathbb{Z}^3$ is mapped onto the block with centre $\tilde{g}z$; in particular $\mathrm{Bl}_{j}(z)$ goes onto $\mathrm{Bl}_{j}(\tilde{g}z)$, $\Delta_y$ onto $\Delta_{\tilde{g}y}$, and $B(y')$ onto $B(\tilde{g}y')$.
For sites $y$, $c$ with $d(y,c)<\ell/2$ the path $\boldsymbol{\Gamma}^{\pi}(y,c)$, formed on the shortest lift, is mapped onto the path $\boldsymbol{\Gamma}^{\pi}(\tilde{g}y,\tilde{g}c)$ with the same order $\pi$.
Segments go onto segments, $M$-cubes onto $M$-cubes, centres onto centres, cubes with a common face onto cubes with a common face, and nets onto nets of the same length.
Hence polymers go onto polymers and small polymers onto small polymers, and
\[
|\tilde{g}X|_M=|X|_M,\qquad d_M(\tilde{g}X)=d_M(X),\qquad(\tilde{g}X)^{\sim i}=\tilde{g}\left(X^{\sim i}\right),\qquad(\tilde{g}\square)^{\dagger}=\tilde{g}\left(\square^{\dagger}\right);
\]
the cubes of $\pi_k$, resp.\ of $\pi'_k$, go onto cubes of the same kind, and linked cubes onto linked cubes;
one has $K_{\tilde{g}z}=\tilde{g}K_z$;
and the shortest chains of $M$-cubes of $X$ between two of its cubes go onto the shortest chains of $M$-cubes of $\tilde{g}X$ between the two image cubes.
If $\tilde{g}\in\mathbb{G}_k^{L}$, then $LM$-cubes go onto $LM$-cubes, the cubes of $\pi''_k$ onto cubes of $\pi''_k$, and
\[
\overline{\tilde{g}X}^{L}=\tilde{g}\left(\bar{X}^{L}\right).
\]
The rule $(\tilde{g}s)_{\mathrm{pc}}=s_{\tilde{g}^{-1}\mathrm{pc}}$ for parameters indexed by $\pi''_k$ is laid down for $\tilde{g}\in\mathbb{G}_k^{L}$ only.
\item[(a2)] \emph{Fermion sites} (the facts (I.a1), (I.a2)).
The map $x=(c,g)\mapsto\tilde{g}x:=(rc+a,\,rg)$ is a bijection of $\mathbb{L}^{\mathrm{f}}_k$ with $\tilde{g}(x+\ell n)=\tilde{g}x+\ell\, rn$, and it induces a bijection of $T^{\mathrm{f}}_k$.
One has $\pi(\tilde{g}x)=\tilde{g}\,\pi(x)$, and the cluster of $c$ goes onto the cluster of $\tilde{g}c$.
For $i=0,1,2,3$,
\[
(\tilde{g}x)^{(i)}=\tilde{g}\left(x^{(i)}\right),\qquad\hat{d}_i(\tilde{g}x)=r\,\hat{d}_i(x).
\]
Let $f$ be an $\ell\mathbb{Z}^3$-periodic function on $\mathbb{L}^{\mathrm{f}}_k$ with values in a complex vector space.
Then
\[
\int\mathrm{d}x\,f(\tilde{g}x)=\int\mathrm{d}x\,f(x),
\qquad
\int_{\tilde{g}Z}\mathrm{d}x\,f(x)=\int_{Z}\mathrm{d}x\,f(\tilde{g}x)
\]
for every set $Z$ of sites of $T_k$.
On the covering lattice, put
\[
\int_{\mathsf{Z}}\mathrm{d}x\,f(x):=\eta_k^3\cdot\frac{1}{48}\sum_{c\in\mathsf{Z}}\ \sum_{g\in O_{\mathrm{cub}}}f(c,g)
\]
for a finite subset $\mathsf{Z}$ of $\mathbb{L}_k$.
Then
\[
\int_{\tilde{g}\mathsf{Z}}\mathrm{d}x\,f(x)=\int_{\mathsf{Z}}\mathrm{d}x\,f(\tilde{g}x)
\]
for every function $f$ on $\mathbb{L}^{\mathrm{f}}_k$ with values in a complex vector space, periodic or not; this is the form in which antiperiodic kernels are integrated over a lift.
Moreover $d_X(\tilde{g}x,\tilde{g}x')=d_X(x,x')$.
The map $(x,x')\mapsto(\tilde{g}x,\tilde{g}x')$ maps $\mathrm{Pzt}_k$ onto $\mathrm{Pzt}_k$ and induces a bijection of $\mathrm{Pz}_k$.
It maps pairs inside a cluster, short pairs, long pairs and bonds of type $i$ to pairs of the same sort, and, for every union $Z$ of $M$-cubes, a pair admissible for $Z$ of kind (I), resp.\ (II), onto a pair admissible for $\tilde{g}Z$ of the same kind.
Let $X$ be a small polymer, let $x=(c,g)$ and $\tilde{x}$ be sites over $X$, and let $(\square_0,\dots,\square_n)$ be a shortest chain with the centres $z_i$, as in Section~I.5.6.
Then
\[
\tilde{g}(z_i,g)=\left(\text{the centre of }\tilde{g}\square_i,\ rg\right),
\]
and therefore the chain points of $\left(\tilde{g}X,\tilde{g}x,\tilde{g}\tilde{x},(\tilde{g}\square_i)\right)$ are the images under $\tilde{g}$ of the chain points of $\left(X,x,\tilde{x},(\square_i)\right)$.
\item[(b)] \emph{Fields.}
The map $A\mapsto\tilde{g}A$ preserves the sup, $\ell^1$ and $\ell^2$ norms, the pairings of bond functions and of plaquette functions, and the three norms of (G2.r4); for a subset $Z$, the three norms of $\tilde{g}A$ on $\tilde{g}Z$ coincide with those of $A$ on $Z$.
One has
\[
(\tau\,\tilde{g}A)(\tilde{g}y,\tilde{g}c)=(\tau A)(y,c),\qquad\tilde{g}(d\boldsymbol{\lambda})=d\left(\boldsymbol{\lambda}\circ\tilde{g}^{-1}\right),
\]
and
\[
U_{\tilde{g}B}(\tilde{g}x,\tilde{g}x')=U_B(x,x')\qquad\left((x,x')\in\mathrm{Pzt}_k\right).
\]
Multiplication by an element of $\mathrm{Sp}$, by $\mathsf{C}$ or by $\varkappa$ preserves $|v|_c$ on $\mathbb{C}^4$, as well as $|f|_{\infty}$ on spinors and on functions of pairs.
Each of the four maps $a_{\tilde{g}}$, $G_{\Lambda}$ ($\Lambda$ real), $\mathrm{Cg}$, $\mathbf{K}$ of Section~I.4.5 is compatible with (I.AP).
Each of them sends every generator $\psi(\xi)$, $\chi(\boldsymbol{\zeta})$ to a generator of the same family multiplied by a number of modulus one, and, these factors being disregarded, it does so bijectively on the set of the generators; for $a_{\tilde{g}}$ this bijection contains the permutation $g\mapsto rg$ of the labels.
Each of them preserves the weights $\nu$ and the bidegree, and
\begin{equation}
\left\|a_{\tilde{g}}F\right\|_{\mathsf{w}}=\left\|G_{\Lambda}F\right\|_{\mathsf{w}}=\left\|\mathrm{Cg}\,F\right\|_{\mathsf{w}}=\left\|\mathbf{K}F\right\|_{\mathsf{w}}=\|F\|_{\mathsf{w}}
\tag{2.1}
\end{equation}
for all $F\in\mathrm{Gr}[\psi,\chi]$ and all $\mathsf{w}=(\mathsf{w}_1,\mathsf{w}_2)$.

\emph{The families $\Psi$, $W$.}
On $\Psi$ and on $W$ the four maps act by the same rules (Section~I.4.5), on the index sets and with the weights $\nu$ of Section~I.3.6.
For $k\ge1$ the unit sites $y$ of $T^0_k=\mathbb{Z}^3/\ell\mathbb{Z}^3$ take the place of $x$, so that $\tilde{g}^{-1}y$ and $\Lambda(y)$ occur, and $\nu=1$; for $k=0$ the index set is that of $\psi$ ($x\in T^{\mathrm{f}}_0$), and $\nu=\frac{1}{48}$.
On $\mathrm{Gr}[\psi,\chi,\Psi,W]$ each of the four maps is compatible with (I.AP) and is an automorphism of rings, antilinear in the case of $\mathbf{K}$.
It sends every generator $\Psi(u)$, $W(u)$ to a generator of the same family multiplied by a number of modulus one, bijectively on the set of the generators of that family when these factors are disregarded; it preserves the weights $\nu$ and the degree in each of the four families; and (2.1) holds for all $F\in\mathrm{Gr}[\psi,\chi,\Psi,W]$ and all $\mathsf{w}=(\mathsf{w}_1,\mathsf{w}_2,\mathsf{w}_3,\mathsf{w}_4)$, one weight for each family (see (I.5.1)).
The same is true of a map that acts by its rule on some of the four families and fixes the generators of the others (Section~G3.4.2, Definition~24.I: $\mathrm{Cg}$ on $\psi$, $\chi$, $\Psi$, and the identity on $W$).

\emph{The vector $a$.}
On the covering lattices, on antiperiodic objects and in $a_{\tilde{g}}$, a vector $a$ of $M\mathbb{Z}^3$ in the class $a$ is used (Section~I.4.3).
If this vector is replaced by $a+\ell n$ with $n\in\mathbb{Z}^3$, then $a_{\tilde{g}}$ is replaced by its composition with the automorphism that multiplies every generator by $(-1)^{n_1+n_2+n_3}$; this automorphism fixes the even elements.
All assertions of (a1)--(d) hold for every vector of the class; in (d) the vector of $\tilde{g}'=(S,a/L)$ is $a/L$.
\item[(c)] \emph{Domains} (an assertion on bond functions alone).
Each of the sets $R_k$, $\mathrm{cf}\,\tilde{R}_k$ $(0<\mathrm{cf}\le1)$, $\tilde{R}_k^{-}$, $G_k$, $\tilde{G}_k$, $\mathfrak{X}_k$ is invariant under $A\mapsto\tilde{g}A$, under $A\mapsto-A$ and under $A\mapsto\overline{A}$.
The sets $\mathrm{cf}\,\tilde{R}_k$, $\tilde{R}_k^{-}$ and $\tilde{G}_k$ are, moreover, invariant under $A\mapsto A-d\boldsymbol{\lambda}$, where $\boldsymbol{\lambda}$ is a real function of $c$, and under the addition of real constant bond functions.
(The data admit $e_k\ge1$.
For such $e_k$ the factor $1-e_k^{\varepsilon}$ in $\tilde{R}_k^{-}$ is $\le0$, so that, by Section~I.5.3, that is by \cite[Definition~5.16(a), (c)]{QEDI}, the set $\left(1-e_k^{\varepsilon}\right)\cdot R_k$ is empty;
hence $\tilde{R}_k^{-}$ is empty, and the assertions of (c) on $\tilde{R}_k^{-}$ hold trivially.)
\item[(d)] \emph{Two levels, under the sole hypothesis $\ell\in\mathrm{Cl}$.}
The number $\ell/L$ is an odd multiple of $M$, and $\ell/L\ge13M>2M$.
The following objects of level $k+1$ on the side $\ell/L$ are given by the formulas of Sections~I.1.1, I.1.3, I.1.5, I.1.7, I.3.6, I.4.2, I.4.3 and I.4.5, which use of the side only that it is a multiple of $M$ larger than $2M$:
the torus $T_{k+1}=\mathbb{L}_{k+1}/(\ell/L)\mathbb{Z}^3$ and $T^{\mathrm{f}}_{k+1}$;
the lattices $L^{-i}\mathbb{Z}^3/(\ell/L)\mathbb{Z}^3$, $-1\le i\le k+1$, with their bonds and their bond functions;
the $M$-cubes;
the set $\mathrm{Pz}_{k+1}$ with its shortest lifts;
the generators $\psi$, $\chi$ of level $k+1$;
the spaces $\mathfrak{S}^{\mathrm{ch}}(\eta_{k+1},\ell/L)$ and $\mathfrak{S}(1,\ell/L)$;
the group
\[
\mathbb{G}_{k+1}=\left\{(S,a'):\ a'\in M\mathbb{Z}^3/(\ell/L)\mathbb{Z}^3\right\}
\]
with the actions of Section~I.4.3 of its elements on these sites, bond functions and spinors;
the four maps of Section~I.4.5 on these generators, namely $a_{(S,a')}$ for $(S,a')\in\mathbb{G}_{k+1}$, $G_{\Lambda}$ with the function $\Lambda$ and the number $e_{k+1}$ of (d4), $\mathrm{Cg}$ and $\mathbf{K}$;
and the homomorphism $S_L$ of (G2.r5).
(The condition $\ell/L\in\mathrm{Cl}$ is needed only for the polymers $D_{k+1}$ together with the assertions of \cite[Proposition~5.58]{QEDI} on them, for the domains of level $k+1$, and for reblocked families.)
Then the following hold.
\begin{itemize}
\item[(d1)] The map $\tilde{g}=(S,a)\mapsto\tilde{g}':=(S,a/L)$ is an isomorphism of $\mathbb{G}_k^{L}$ onto $\mathbb{G}_{k+1}$.
\item[(d2)] One has
\[
(\tilde{g}c)/L=\tilde{g}'(c/L),\qquad(\tilde{g}x)/L=\tilde{g}'(x/L),
\]
and $L(\tilde{g}'Y)=\tilde{g}(LY)$ for every set $Y$ of sites of $T_{k+1}$.
\item[(d3)] Let $A$ be a bond function on one of the lattices $L^{-i}\mathbb{Z}^3/(\ell/L)\mathbb{Z}^3$, $0\le i\le k+1$, and let
\[
A_L(b):=L^{-\frac{1}{2}}A(b/L)
\]
on the bonds $b$ of $L^{-(i-1)}\mathbb{Z}^3/\ell\mathbb{Z}^3$.
Then $(\tilde{g}'A)_L=\tilde{g}(A_L)$.
Moreover $(\tilde{g}'f')_L=\tilde{g}(f'_L)$ for $f'\in\mathfrak{S}^{\mathrm{ch}}(\eta_{k+1},\ell/L)$ and for $f'\in\mathfrak{S}(1,\ell/L)$.
\item[(d4)] One has
\[
S_L\circ a_{\tilde{g}}=a_{\tilde{g}'}\circ S_L,\qquad S_L\circ\mathrm{Cg}=\mathrm{Cg}\circ S_L,\qquad S_L\circ\mathbf{K}=\mathbf{K}\circ S_L,
\]
and, with $e_{k+1}:=L^{\frac{1}{2}}e_k$ and a real function $\Lambda$ on the sites of $T_{k+1}$,
\[
G_{\Lambda}\circ S_L=S_L\circ G_{\Lambda_L},\qquad(d\Lambda)_L=d\left(\Lambda_L\right).
\]
In each of these relations, a map written to the left of $S_L$ is the map of level $k+1$, and a map written to its right is the map of level $k$.
\end{itemize}
\end{itemize}
\end{named}

\begin{named}{Remark}
Proposition~16.1 corresponds to the facts (a)--(d) of \cite[Proposition~5.59]{QEDI}, to \cite[Proposition~13.3]{QEDI}, to the part of \cite[Section~13.4.1]{QEDI} that precedes \cite[Proposition~13.3]{QEDI}, from the definition of $r(S)$ onwards, and to \cite[Section~13.1.5]{QEDI} together with the equalities of norms displayed there.
The part of \cite[Section~13.4.1]{QEDI} just named contains the homomorphism $r$ and the group $\mathrm{Sp}_3$; the relations $\left(S^{-1}\right)^{T}=\pm\mathsf{C}^{-1}\cdot S\cdot\mathsf{C}$ and $\overline{S}=\pm\varkappa^{-1}\cdot S\cdot\varkappa$ for $S\in\mathrm{Sp}$, with signs that depend on $S$; the groups $\mathbb{G}_k$ and $\mathbb{G}_k^{L}$; the actions of their elements; and the use, for antiperiodic functions, of a vector $a$ of $M\mathbb{Z}^3$ in the place of its class.
It differs from these statements in the following points.
\begin{itemize}
\item[(1)] In (a2), the assertion that $(x,x')\mapsto(\tilde{g}x,\tilde{g}x')$ induces a bijection of $\mathrm{Pz}_k$ and carries a pair admissible for $Z$ of kind (I), resp.\ (II), onto a pair admissible for $\tilde{g}Z$ of the same kind is the form, for pairs of fermion sites, of the corresponding sentence of \cite[Proposition~5.59(a)]{QEDI} and of \cite[Proposition~13.3(a)]{QEDI}.
The rest of (a2) has no counterpart in \cite{QEDI}: the action on the labels and on the clusters, on the neighbours $x^{(i)}$ and on $\hat{d}_i$, on the measure and the rules of substitution, among them the rule over finite subsets of $\mathbb{L}_k$ for arbitrary $f$, on $d_X$, on $\mathrm{Pzt}_k$, on the sorts of pairs, and on the chain points.
The two paragraphs that follow (2.1), on the families $\Psi$, $W$ and on the vector $a$, have no counterpart among the clauses of \cite[Proposition~5.59]{QEDI} and of \cite[Proposition~13.3]{QEDI}.
\item[(2)] In (b), $U_B$ is taken at pairs of fermion sites, and in (2.1) the map $a_{\tilde{g}}$ contains the permutation of the labels; $|f|_{\infty}$ is formed with $|\cdot|_c$ (Section~I.1.5).
\item[(3)] In (a1), the pairs for which a path on the torus is meant are specified, and the assertion on shortest lifts is stated.
The rule by which an element that acts on families of parameters indexed by cubes of $\pi''_k$ is taken in $\mathbb{G}_k^{L}$, which \cite[Section~13.4.1]{QEDI} lays down among the definitions of the actions, is stated here as the last sentence of (a1).
\item[(4)] In (b), the three norms on a subset are those defined in (G2.r4).
\item[(5)] Clause (d) is stated under the sole hypothesis $\ell\in\mathrm{Cl}$, with the list of the objects of level $k+1$ that it covers, as in \cite{QEDI}.
In \cite[Proposition~13.3(d)]{QEDI} the side is subject only to the condition $\ell\in\mathrm{Cl}$;
the whole of \cite[Proposition~13.3]{QEDI}, its clause (d) included, is placed under this sole hypothesis in \cite[Section~13.1.1]{QEDI}, where $\ell/L\in\mathrm{Cl}$ is required only for the polymers of $D_{k+1}$, for the reblocked family and for the domains of level $k+1$;
and by \cite[Theorem~5.76(b)]{QEDI} the assertions (a)--(d) of \cite[Proposition~5.59]{QEDI} hold for every integer $k\ge0$ under the sole condition $\ell\in\mathrm{Cl}$.
Each of these three places notes, as (d) does, that $\ell/L$ is $M$ times an odd integer $\ge13$, hence a multiple of $M$ larger than $2M$, and names the objects of level $k+1$, which are formed with the side $\ell/L$;
\cite[Proposition~13.3(d)]{QEDI} and \cite[Theorem~5.76(b)]{QEDI} add that $\ell/L$ need not lie in $\mathrm{Cl}$.
The list of (d) contains in addition the objects of the present fermion, namely the torus $T^{\mathrm{f}}_{k+1}$, the spaces $\mathfrak{S}^{\mathrm{ch}}(\eta_{k+1},\ell/L)$ and $\mathfrak{S}(1,\ell/L)$ and the actions of the elements of $\mathbb{G}_{k+1}$ on fermion sites, and, besides the automorphisms $a_{(S,a')}$, the maps $G_{\Lambda}$, $\mathrm{Cg}$ and $\mathbf{K}$ on the generators of level $k+1$.
\item[(6)] The first assertion of (d3) agrees with \cite[Proposition~13.3(d)]{QEDI}, which is stated for the same lattices $L^{-i}\mathbb{Z}^3/(\ell/L)\mathbb{Z}^3$, $0\le i\le k+1$, and with the same function $A_L$;
the case $i=0$, that of the unit lattice $T^0_{k+1}$, is the one applied in \cite[Section~15]{QEDI}.
Among bond functions, \cite[Proposition~5.59(d)]{QEDI} concerns only those of the level $k+1$, with the scalings of \cite[Section~5.3.5]{QEDI}.
The second assertion of (d3), the relation $(\tilde{g}'f')_L=\tilde{g}(f'_L)$ for $f'\in\mathfrak{S}^{\mathrm{ch}}(\eta_{k+1},\ell/L)$ and for $f'\in\mathfrak{S}(1,\ell/L)$, is not among the three relations displayed in \cite[Proposition~13.3(d)]{QEDI}.
\item[(7)] The relations of (d4) for $\mathrm{Cg}$, $\mathbf{K}$ and $G_{\Lambda}$, which \cite{QEDI} establishes within the proof of \cite[Theorem~13.5(e)]{QEDI}, are clauses of the statement.
\item[(8)] The proof of (c) rests on the inequalities $(2L+1)M+\eta_k<\ell$ and $11M+\eta_k<\ell$.
\end{itemize}
\end{named}

\subsection*{9. The operator $T_k(A)$}

\subsubsection*{Proposition 9}

\begin{named}{Proposition 9}
Let
\[
C_T:=2\mathrm{e}^{\frac{3}{2}}\left[b_{-}^{-1}\left(54L\mathrm{e}^{6}+12\mathrm{e}+1\right)+1\right],
\]
a function of $(b_{-},L)$; this is the value of the constant in \cite[Proposition~16.27(d)]{QEDI}.
Let the data (SF) be given.
Throughout this Proposition $y$ denotes a unit site; no resolvent parameter occurs.
The operators $T_k$ and $\bar{T}_k$ are those of Sections~I.3.5 and I.3.9.
For $k\ge1$, for a union $\mathbf{Y}$ of $M$-cubes (in the step map, a polymer) and for a unit site $y$ of $\mathbf{Y}$, the operators
\[
T^{(1)}(A)=T^{(1)}_{\mathbf{Y}}(A),\qquad T^{(2)}(A)=T^{(2)}_{\mathbf{Y}}(A),\qquad\bar{T}^{(i)}(A)
\]
are those of (I.5.4).
Write $\mathfrak{i}:=\mathfrak{i}(A)$ and $\mathrm{qd}:=e_k|dA|_{\infty}$.
\begin{itemize}
\item[(a)] Let $k\ge1$ and $A\in\tilde{G}_k$.
Then
\[
T_k(A)\cdot H_k(A)=1\qquad\text{and}\qquad\bar{T}_k(A)\cdot\bar{H}_k(A)=1\qquad\text{on }\mathfrak{S}(1,\ell).
\]
At $k=0$, for every connection $A$, the four operators $T_0(A)$, $\bar{T}_0(A)$, $H_0(A)$, $\bar{H}_0(A)$ are the identity of $\mathfrak{S}^{\mathrm{ch}}(1,\ell)$, by (I.L0.7) and Section~I.3.9.
Hence both relations hold on $\mathfrak{S}^{\mathrm{ch}}(1,\ell)$, which at $k=0$ takes the place of $\mathfrak{S}(1,\ell)$.
\item[(b)] Let $k\ge1$, and let $A$ be an arbitrary connection on $T_k$.
Then
\begin{multline*}
\left(T_k(A)\psi\right)(y)=\int_{\Delta_y}\mathrm{d}x\,\varphi_A(y,c)\cdot\left\{\vphantom{\sum_{i=0}^{3}}\left(1+\frac{\mu}{b_k}\right)\psi(x)\right.\\
\left.-(b_k\eta_k)^{-1}\sum_{i=0}^{3}M_i(x)\cdot\left[U_A\left(x,x^{(i)}\right)\cdot\psi\left(x^{(i)}\right)-\psi(x)\right]\right\},
\end{multline*}
where $U_A\left(x,x^{(i)}\right)=1$ for $i\le2$.
Let $\mathbf{Y}$ be a union of $M$-cubes and $y$ a unit site of $\mathbf{Y}$, and suppose that
\[
\chi\left(x,x^{(3)}\right)=\left(\delta_{\alpha,A}\psi\right)\left(x,x^{(3)}\right)
\]
for those $x$ over $\Delta_y$ with $c_3\notin\mathbf{Y}$.
Then
\[
\left(T^{(1)}(A)\psi\right)(y)+\left(T^{(2)}(A)\chi\right)(y)=\left(T_k(A)\psi\right)(y).
\]
In $\left(T^{(1)}(A)\psi\right)(y)$ the field $\psi$ occurs only at sites over $\mathbf{Y}$, namely at the sites over $\Delta_y$ and at the sites $x^{(3)}$ over $\mathbf{Y}$ with $x$ over $\Delta_y$.
In $\left(T^{(2)}(A)\chi\right)(y)$ the function $\chi$ occurs only at the pairs $\left(x,x^{(3)}\right)$ with $x$ over $\Delta_y$ and $c_3\notin\mathbf{Y}$; these are pairs of the kind (II) of $\mathrm{Pz}_k(\mathbf{Y})$, of length $d_X=\eta_k$.
Differences inside a cluster ($i\le2$) are always written with $\psi$.
Both rows read $A$ on the bonds that have an end point in $\Delta_y$.
\item[(c)] Let $k\ge1$, and let $A$ be an arbitrary connection.
For the operator norms from $|\cdot|_{\mathrm{su}}$ over the sites, resp.\ over the pairs, to $|\cdot|_{\mathrm{su}}$ over the unit sites of $\mathbf{Y}$ (Section~G1.3),
\begin{align*}
\left\|T^{(1)}(A)\right\|_{\mathrm{su}},\ \left\|T_k(A)\right\|_{\mathrm{su}}
&\le\exp\left(\frac{3\mathfrak{i}}{2}\right)\cdot\left[b_{-}^{-1}\left(54L\cdot\mathrm{qd}\cdot\exp\left(6\,\mathrm{qd}\right)+12+6\exp(\mathfrak{i})+\mu\right)+1\right],\\
\left\|T^{(2)}(A)\right\|_{\mathrm{su}}
&\le6b_{-}^{-1}\cdot\exp\left(\frac{3\mathfrak{i}}{2}\right)\cdot\eta_k^{\alpha}.
\end{align*}
\item[(d)] Let $k\ge1$ and $A\in\tilde{R}_k+\mathfrak{X}_k$.
Then $\mathfrak{i}\le1$ and $\mathrm{qd}\le1$, and
\[
\left|T^{(i)}(A)\right|_{1,\infty}\le C_T\quad(i=1,2),\qquad\left|T_k(A)\right|_{1,\infty}\le C_T;
\]
the same bounds hold for $\bar{T}^{(i)}(A)$ and for $\bar{T}_k(A)$.
Here $|\cdot|_{1,\infty}$ is the norm of (I.ke3) and Section~I.5.2; for $T^{(2)}$ the rows are summed over the coefficients of the values $\chi\left(x,x^{(3)}\right)_{\beta}$.
\item[(e)] \emph{Covariances.}
Let $k\ge1$, and let $A$ be an arbitrary connection.
For a real function $\boldsymbol{\lambda}$ of $c$ let $\boldsymbol{\lambda}^{\mathrm{u}}$ be its restriction to $T^0_k$, and put
\[
\chi^{\boldsymbol{\lambda}}(x,x'):=\exp\left(ie_k\boldsymbol{\lambda}(c)\right)\cdot\chi(x,x').
\]
Then
\begin{align*}
T^{(1)}(A-d\boldsymbol{\lambda})\cdot\exp\left(ie_k\boldsymbol{\lambda}\right)\psi&=\exp\left(ie_k\boldsymbol{\lambda}^{\mathrm{u}}\right)\cdot T^{(1)}(A)\psi,\\
T^{(2)}(A-d\boldsymbol{\lambda})\chi^{\boldsymbol{\lambda}}&=\exp\left(ie_k\boldsymbol{\lambda}^{\mathrm{u}}\right)\cdot T^{(2)}(A)\chi.
\end{align*}
For $\tilde{g}=(S,a)\in\mathbb{G}_k$, with $(\tilde{g}\chi)(\tilde{g}x,\tilde{g}x'):=S\chi(x,x')$,
\[
T^{(i)}_{\tilde{g}\mathbf{Y}}(\tilde{g}A)\cdot\tilde{g}=\tilde{g}\cdot T^{(i)}_{\mathbf{Y}}(A),
\]
and
\[
\overline{T^{(i)}\left(\overline{A}\right)}=\varkappa^{-1}\cdot T^{(i)}(-A)\cdot\varkappa.
\]
For the barred operators, with the action $\mathrm{glb}$ of $\tilde{g}$ on barred fields (Section~G1.10),
\begin{align*}
\bar{T}^{(1)}(A-d\boldsymbol{\lambda})\cdot\exp\left(-ie_k\boldsymbol{\lambda}\right)&=\exp\left(-ie_k\boldsymbol{\lambda}^{\mathrm{u}}\right)\cdot\bar{T}^{(1)}(A),\\
\bar{T}^{(2)}(A-d\boldsymbol{\lambda})\chi^{-\boldsymbol{\lambda}}&=\exp\left(-ie_k\boldsymbol{\lambda}^{\mathrm{u}}\right)\cdot\bar{T}^{(2)}(A)\chi,\\
\bar{T}^{(i)}_{\tilde{g}\mathbf{Y}}(\tilde{g}A)\cdot\mathrm{glb}&=\mathrm{glb}\cdot\bar{T}^{(i)}_{\mathbf{Y}}(A),\\
\overline{\bar{T}^{(i)}\left(\overline{A}\right)}&=\varkappa^{-1}\cdot\bar{T}^{(i)}(-A)\cdot\varkappa.
\end{align*}
\end{itemize}
\end{named}

The conditions on the constants that enter Proposition~9 are the following:
in (a), those of Proposition~2.8;
the range $\rho\in(0,1)$ of the line (C0) of Section~I.8.2, which is used in the form $0\le\rho\le1$: it enters (c), and through (c) also (d), the barred operators included, by way of the equalities $|M_3(x)|_2=|M_3(x)|_c=3(1+\rho)$ of (I.m6), (I.m7) and the bound $3(1+\rho)\le6$;
the value $b=1$ of the line (C0), which together with $L>1$ from (L5) makes the numbers $b_k$ of Section~I.3.4 defined, with $0<b_{-}\le b_k$;
the standard weights $w$ of the line (C0), through the closure relation (I.1.1) and the value $w_3=6$;
the condition (Z2.s4), which enters through the display (G3) of \cite[Section~16.2.7]{QEDI};
the row (S1) of Proposition~14.2\,(b), that is, the range $1-\alpha<\varepsilon\le\bar{\varepsilon}$ of $\varepsilon$ in the line (C7) of Section~I.8.2, which enters (d) through the same display;
the range $\alpha\in(1-\bar{\varepsilon},1)$ of the line (C1) of Section~I.8.2, in two forms: as $\alpha\ge0$ it enters (d) through $\eta_k^{\alpha}\le1$ in the bound of $T^{(2)}$, and as $0<\alpha<1$ it is among the data of \cite[Section~16.2.1]{QEDI} under which the display (G3) stands;
(L5);
and (M1), which together with (L5) makes $M$ odd, so that $\Delta_y$ lies in the $M$-cube of $y$.
The display (G3) stands under the data of \cite[Section~16.2.1]{QEDI}, among them $\varepsilon\in(0,\bar{\varepsilon}]$, and its proof uses $\varepsilon\ge0$;
together with $\alpha<1$, the row (S1) gives $0<\varepsilon\le\bar{\varepsilon}$, and, as $\bar{\varepsilon}=\frac{1}{200}$, also $\alpha>1-\bar{\varepsilon}>0$.

\begin{named}{Remark}
Proposition~9 corresponds to \cite[Proposition~5.46]{QEDI} and to \cite[Proposition~16.27]{QEDI}, together with the display of $\left(T_k(A)\psi\right)(y)$ and the definitions of $T^{(1)}$, $T^{(2)}$ in \cite[Section~16.11]{QEDI}.
In \cite[Lemma~16.26]{QEDI}, a statement on bond functions alone, the coordinates $0,1,2$ are read as the coordinates $1,2,3$, in this order; thus the rectilinear paths of that lemma change the coordinates in the order $1,2,3$.
The differences are the following.
\begin{itemize}
\item[(1)] The rows of $D_A$ contain the four types of bonds, with the hopping matrices $M_i(x)$; only a bond of type $3$ can leave the block or the polymer.
\item[(2)] In the proof the bulk terms cancel by
\[
w_i\,\gamma\cdot\hat{d}_i(x)=M_i\left(x^{(i)}\right)-M_i(x)
\]
and by the closure relation $\sum_iw_i\hat{d}_i(x)=0$ of (I.s3); in \cite{QEDI} they cancel because $\mathrm{Pr}^{+}_{\nu}+\mathrm{Pr}^{-}_{\nu}=1$.
\item[(3)] In (c) the term $12\exp(\mathfrak{i})$ of \cite[Proposition~16.27(c)]{QEDI} becomes $12+6\exp(\mathfrak{i})$; the constant $C_T$ keeps its value, because $12+6\mathrm{e}\le12\mathrm{e}$.
\item[(4)] As in the last sentence of \cite[Proposition~16.27]{QEDI}, the two relations $T_0\cdot H_0=1$ and $\bar{T}_0\cdot\bar{H}_0=1$ are included in (a), the four operators $T_0(A)$, $H_0(A)$, $\bar{T}_0(A)$, $\bar{H}_0(A)$ being the identity for every connection $A$.
Here they are the identity of $\mathfrak{S}^{\mathrm{ch}}(1,\ell)$, which has $192$ components at each unit site and at $k=0$ takes the place of $\mathfrak{S}(1,\ell)$.
The operators $T^{(1)}$, $T^{(2)}$ are those of (I.5.4): they are operators on the whole unit lattice whose rows vanish at the unit sites outside $\mathbf{Y}$ (choice~(3)), and the pairs of $T^{(2)}$ are of the kind (II).
\item[(5)] Clause (e), with its barred forms, has no counterpart in \cite[Proposition~16.27]{QEDI}.
\end{itemize}
\end{named}

\subsection*{10. Covariances of all exact operators}

\subsubsection*{Proposition 10}

\begin{named}{Proposition 10}
Let the data (SF) be given, let $A\in\tilde{G}_k$, and let $y\in Y_k$, with $y\in\mathbb{R}$ wherever $B_{k,y}$ or $\Gamma_{k,y}$ occurs.
The clauses (a)--(f) concern $k\ge1$.
\begin{itemize}
\item[(a)] \emph{Gauge covariance, for closed and not only for exact connections.}
Let $\mathrm{Acl}$ be a real connection on $T_k$ with $d\,\mathrm{Acl}=0$, and let $\mathrm{uf}$ be a function of $c$ on $T_k$ with $|\mathrm{uf}|=1$ and
\[
\mathrm{uf}(c')=\mathrm{uf}(c)\cdot\exp\left\{ie_k\eta_k\,\mathrm{Acl}(c,c')\right\}
\]
on every bond $(c,c')$; let $\mathrm{uf}^0$ be the restriction of $\mathrm{uf}$ to $T^0_k$.
Then $A-\mathrm{Acl}$ lies in $\tilde{G}_k$, and in addition to the relations of Proposition~2.2\,(a) one has
\begin{gather*}
S_y(A-\mathrm{Acl})=\mathrm{uf}\cdot S_y(A)\cdot\mathrm{uf}^{-1},\qquad H_k(A-\mathrm{Acl})=\mathrm{uf}\cdot H_k(A)\cdot\left(\mathrm{uf}^0\right)^{-1},\\
\mathfrak{H}_k(A-\mathrm{Acl})=\mathrm{uf}^0\cdot\mathfrak{H}_k(A)\cdot\mathrm{uf}^{-1},\qquad T_k(A-\mathrm{Acl})=\mathrm{uf}^0\cdot T_k(A)\cdot\mathrm{uf}^{-1},\\
O(A-\mathrm{Acl})=\mathrm{uf}^0\cdot O(A)\cdot\left(\mathrm{uf}^0\right)^{-1}\qquad\text{for }O=D_k,\ M_k,\ \dot{M}_k,\ B_{k,y},\ \Gamma_{k,y},\\
\delta Z_k(A-\mathrm{Acl})=\delta Z_k(A),\\
\bar{H}_k(A-\mathrm{Acl})=\mathrm{uf}^{-1}\cdot\bar{H}_k(A)\cdot\mathrm{uf}^0,\qquad\bar{T}_k(A-\mathrm{Acl})=\left(\mathrm{uf}^0\right)^{-1}\cdot\bar{T}_k(A)\cdot\mathrm{uf}.
\end{gather*}
Examples: $\mathrm{Acl}=d\boldsymbol{\lambda}$ and $\mathrm{uf}=\exp\left(ie_k\boldsymbol{\lambda}\right)$, where $\boldsymbol{\lambda}$ is a real function of $c$; or a real constant connection $\mathrm{Acl}$ with $e_k\ell\,\mathrm{Acl}_{\mu}\in2\pi\mathbb{Z}$, together with $\mathrm{uf}(c)=\exp\left(ie_k\,\mathrm{Acl}\cdot c\right)$.
\item[(b)] \emph{Lattice symmetries.}
Let $\tilde{g}=(S,a)\in\mathbb{G}_k$.
Then
\[
O(\tilde{g}A)\cdot\tilde{g}=\tilde{g}\cdot O(A)
\]
for each of
\begin{gather*}
O(A)=D_A,\ Q_k(A)J^{*},\ JQ_k^{T}(-A),\ Q(A),\ Q^{T}(-A),\ JP_k(A)J^{*},\ P(A),\ JP_{k+1}(A)J^{*},\ S_k(A),\ S^{0}_{k+1}(A),\\
H_k(A),\ \mathfrak{H}_k(A),\ T_k(A),\ D_k(A),\ M_k(A),\ \dot{M}_k(A),\ B_k(A),\ \Gamma_k(A);
\end{gather*}
further
\[
O_{y'}(\tilde{g}A)\cdot\tilde{g}=\tilde{g}\cdot O_y(A)\qquad\text{for }O_y=\operatorname{Pen}_{y},\ F_{y},\ S_{y},\ B_{k,y},\ \Gamma_{k,y},
\]
where $y'=y$ if $S\in\mathrm{Sp}_3$ or if $y$ is neutral (Section~G1.9), and $y'=-y$ otherwise.
Moreover
\[
\delta Z_k(\tilde{g}A)=\delta Z_k(A),\qquad\bar{H}_k(\tilde{g}A)\cdot\mathrm{glb}=\mathrm{glb}\cdot\bar{H}_k(A),\qquad\bar{T}_k(\tilde{g}A)\cdot\mathrm{glb}=\mathrm{glb}\cdot\bar{T}_k(A),
\]
where $\mathrm{glb}$ is the action of $\tilde{g}$ on barred fields (Section~G1.10), and
\[
\left(\delta_{\alpha,\tilde{g}A}\,\tilde{g}f\right)(\tilde{g}x,\tilde{g}x')=S\left(\delta_{\alpha,A}f\right)(x,x').
\]
\item[(c)] \emph{Transposition.}
One has
\[
O(A)^{T}=\mathsf{C}^{-1}\cdot O(-A)\cdot\mathsf{C}
\]
for
\begin{gather*}
O(A)=D_A,\ JP_k(A)J^{*},\ JP_{k+1}(A)J^{*},\ S_k(A),\ S^{0}_{k+1}(A),\ P(A),\\
D_k(A),\ M_k(A),\ \dot{M}_k(A),\ B_k(A),\ \Gamma_k(A),
\end{gather*}
and
\[
O_y(A)^{T}=\mathsf{C}^{-1}\cdot O_{-y}(-A)\cdot\mathsf{C}\qquad\text{for }O_y=\operatorname{Pen}_{y},\ F_{y},\ S_{y},\ B_{k,y},\ \Gamma_{k,y};
\]
moreover
\begin{gather*}
\left(Q_k(A)J^{*}\right)^{T}=JQ_k^{T}(A),\qquad Q(A)^{T}=Q^{T}(A),\qquad H_k(A)^{T}=\mathsf{C}^{-1}\cdot\mathfrak{H}_k(-A)\cdot\mathsf{C},\\
\bar{H}_k(A)=b_k\cdot S_k(A)^{T}\cdot JQ_k^{T}(A),\qquad\bar{T}_k(A)=b_k^{-1}\cdot Q_k(-A)J^{*}\cdot\left(D_A+\mu+b_k\cdot JP_k(A)J^{*}\right)^{T}.
\end{gather*}
\item[(d)] \emph{Conjugation.}
For all operators named in (b), that is, for the two lists of (b) and for $\delta_{\alpha,\cdot}$, and also for $\bar{H}_k$ and $\bar{T}_k$, with the same $y$,
\[
\overline{O\left(\overline{A}\right)}=\varkappa^{-1}\cdot O(-A)\cdot\varkappa.
\]
(For the number $\delta Z_k$ see (e).)
\item[(e)] One has
\[
\delta Z_k(-A)=\delta Z_k(A)\qquad\text{and}\qquad\overline{\delta Z_k\left(\overline{A}\right)}=\delta Z_k(A).
\]
\item[(f)] \emph{Adjoints.}
Let $A\in\tilde{G}_k$ be real.
Then
\[
O(A)^{*}=\gamma_5\cdot O(A)\cdot\gamma_5
\]
for
\[
O(A)=D_A,\ S_k(A),\ S^{0}_{k+1}(A),\ D_k(A),\ M_k(A),\ \dot{M}_k(A),\ B_k(A),\ \Gamma_k(A),
\]
further
\[
O_y(A)^{*}=\gamma_5\cdot O_{-y}(A)\cdot\gamma_5\qquad\text{for }O_y=F_{y},\ S_{y},\ B_{k,y},\ \Gamma_{k,y},
\]
and $H_k(A)^{*}=\gamma_5\cdot\mathfrak{H}_k(A)\cdot\gamma_5$.
\item[(g)] \emph{Level 0.}
The assertions (a)--(f) hold at $k=0$ for the operators
\begin{gather*}
D_A,\quad Q(A)J^{*},\quad JQ^{T}(-A),\quad JP(A)J^{*},\quad H_0=\bar{H}_0=T_0=\bar{T}_0=1,\quad D_0(A)=D_A+\mu,\\
\operatorname{Pen}_{y},\quad F_{y},\quad S_{0,y}(A)\ (y\text{ real}),\quad\Gamma_0(A),\quad\delta_{\alpha,A}
\end{gather*}
and for the number $\delta Z_0(A)=\det F_0(A)$.
In them $\mathrm{uf}$ stands in the place of $\mathrm{uf}^0$, since at $k=0$ the unit lattice is the fine lattice and the operators of the unit lattice act on $\mathfrak{S}^{\mathrm{ch}}(1,\ell)$; the restriction $\mathrm{uf}^1$ of $\mathrm{uf}$ to $L\mathbb{Z}^3/\ell\mathbb{Z}^3$ remains on $\mathfrak{S}(L,\ell)$.
Explicitly,
\begin{gather*}
S_{0,y}(A-\mathrm{Acl})=\mathrm{uf}\cdot S_{0,y}(A)\cdot\mathrm{uf}^{-1},\qquad\Gamma_0(A-\mathrm{Acl})=\mathrm{uf}\cdot\Gamma_0(A)\cdot\mathrm{uf}^{-1},\\
\delta Z_0(A-\mathrm{Acl})=\delta Z_0(A)=\delta Z_0(\tilde{g}A)=\delta Z_0(-A)=\overline{\delta Z_0\left(\overline{A}\right)}.
\end{gather*}
The clauses are read at $k=0$ as follows.
The operators $Q(A)J^{*}$, $JQ^{T}(-A)$, $JP(A)J^{*}$, $D_0(A)$, $S_{0,y}(A)$, $\Gamma_0(A)$ and $1$ stand in the places of, respectively,
$Q(A)$;
$Q^{T}(-A)$;
$P(A)$ and $JP_{k+1}(A)J^{*}$;
$D_k(A)$;
$S_y(A)$;
$S^{0}_{k+1}(A)$ and $\Gamma_k(A)$;
and $H_k=\bar{H}_k=T_k=\bar{T}_k$.
In (c) the relation $\left(Q(A)J^{*}\right)^{T}=JQ^{T}(A)$ stands; $\operatorname{Pen}_{y}$ is included in (f); and the four relations of (c) and (f) that contain $\mathfrak{H}_k$, $b_k$ or $S_k$ have no form at $k=0$ (see (I.L0.9)).
\end{itemize}
\end{named}

\begin{named}{Remark}
Proposition~10 corresponds to \cite[Proposition~16.28]{QEDI} and to \cite[Proposition~5.47]{QEDI}.
As in \cite[Proposition~16.28(a), (c)]{QEDI} and in \cite[Proposition~5.47(a), (c)]{QEDI}, the connection in (a) is closed, and under transposition $y$ goes to $-y$, with the forms of $\bar{H}_k$ and $\bar{T}_k$ in (c).
The rule $y\mapsto-y$ outside $\mathrm{Sp}_3$ in (b) is that of \cite[Proposition~16.28(b)]{QEDI}.
The differences are the following.
\begin{itemize}
\item[(1)] The cluster maps $J$, $J^{*}$ are written out; the matrix written $\gamma_3$ in \cite{QEDI} is $\gamma_5$ here; and $\tilde{g}$ acts on the label.
\item[(2)] The lists include the operators $S_y$ for all $y\in Y_k$, so that $y=\infty$ is covered in (b) and (d), as in \cite[Proposition~16.28]{QEDI}; in \cite[Proposition~5.47]{QEDI} the parameter $y$ is real.
\item[(3)] The operator $\mathfrak{H}_k$, the barred operators and $\delta_{\alpha,\cdot}$ are added to the lists, and clause (f) is added.
\item[(4)] The list at level 0 is that of (I.L0.18).
Compared with the list of \cite{QEDI} at level 0, which consists of $D_A$, $Q$, $P$, $S_{0,y}(A)$, $\Gamma_0(A)$ and $\delta Z_0(A)$, it contains in addition $JQ^{T}(-A)$, $H_0$, $\bar{H}_0$, $T_0$, $\bar{T}_0$, $D_0(A)$, $\operatorname{Pen}_{y}$, $F_{y}$ and $\delta_{\alpha,A}$.
\end{itemize}
\end{named}

\subsection*{11. Exact operators on the tori of the sides $\ell$ and $\mathrm{Jr}\,\ell$}

\subsubsection*{Proposition 11}

\begin{named}{Proposition 11}
Let $k\ge0$ be a level, and let $\ell_1$ and $\ell_2=\mathrm{Jr}\cdot\ell_1$ be sides in $\mathrm{Cl}$, where $\mathrm{Jr}$ is odd.
Let $\mu\in(0,1]$ and $e_k\in(0,e_0]$ be given, the same on both tori, and put $T_i:=\mathbb{L}_k/\ell_i\mathbb{Z}^3$ for $i=1,2$.
The map $\mathrm{pj}^{*}$ is the inclusion of Section~I.6.3.
A spinor on $T_1$ (on fermion sites, on unit sites or on the sites of $L\mathbb{Z}^3$), a function of pairs (antiperiodic under the simultaneous shift of both points), a connection and a scalar on $T_1$ are functions on the covering objects that are antiperiodic, respectively periodic, under $\ell_1\mathbb{Z}^3$;
since $\mathrm{Jr}$ is odd, they are then antiperiodic, respectively periodic, under $\ell_2\mathbb{Z}^3$ as well.
For an operator $O$, the operator $O$ formed on $T_i$ is denoted by $O^{(i)}$.
The upper index of $O^{(i)}$ refers to the torus, whereas the upper index of $T^{(i)}_{\mathbf{Y}}$ and of $\bar{T}^{(i)}_{\mathbf{Y}}$ refers to the piece, as in (I.5.4);
the lower index of $T_1$ and $T_2$ refers to the side and not to a level.
\begin{itemize}
\item[(a)] (One formula.)
In the sense of Section~I.1.9, for the class $\mathrm{Cl}$ of sides, each of the operators listed below is given by one formula on $\mathbb{L}^{\mathrm{f}}_k$, of the range $\mathrm{rg}$ indicated.
In this list $y\in Y_k$ is a parameter, which is real in $B_{k,y}(A)$.
The formulas are the displays of (I.2.3), of Sections~I.3.3, I.3.5 and~I.3.9, of (I.1.7) and of (I.5.4); for $B_{k,y}(A)$ the formula is
\[
\left(b_k+i\gamma_5 y\right)^{-1}\cdot\left(1-P(A)\right)+\left(b_k+b_L+i\gamma_5 y\right)^{-1}\cdot P(A),
\]
which, for every complex bond function $A$, is the inverse of $b_k+b_L\cdot P(A)+i\gamma_5 y$ (Section~I.3.5).
\begin{center}
\begin{tabular}{p{0.62\textwidth}l}
operators at the levels $k\ge1$ & range $\mathrm{rg}$\\
\hline
$D_A$ & $\eta_k$\\
$Q_k(A)\cdot J^{*}$ & $1/2$\\
$J\cdot Q_k^{T}(-A)$, $J\cdot P_k(A)\cdot J^{*}$, $T_k(A)$, $\bar{T}_k(A)$, $T^{(i)}_{\mathbf{Y}}(A)$, $\bar{T}^{(i)}_{\mathbf{Y}}(A)$ (the last two with zero rows at the unit sites outside $\mathbf{Y}$) & $1$\\
$Q(A)$ & $L/2$\\
$Q^{T}(-A)$, $P(A)$, $J\cdot P_{k+1}(A)\cdot J^{*}$, $\operatorname{Pen}_{y}(A)$, $F_{y}(A)$, $B_{k,y}(A)$ & $L$\\
$\delta_{\alpha,A}$ (rows indexed by $\mathrm{Pzt}_k$) & $M$\\
\hline
operators at the level $k=0$ & range $\mathrm{rg}$\\
\hline
$D_A$, $D_0(A)$ & $1$\\
$Q(A)\cdot J^{*}$ & $L/2$\\
$J\cdot Q^{T}(-A)$, $J\cdot P(A)\cdot J^{*}$, $\operatorname{Pen}_{y}(A)$, $F_{y}(A)$ & $L$\\
$\delta_{\alpha,A}$ & $M$\\
$H_0$, $\bar{H}_0$, $T_0$, $\bar{T}_0$, $\dot{M}_0$ (all equal to $1$) & $0$
\end{tabular}
\end{center}
The inverses $S_{y}(A)$ and $\Gamma_{k,y}(A)$, and the operators of the lists in (b) that are built from them ($H_k$, $\bar{H}_k$, $\mathfrak{H}_k$, $D_k$, $M_k$, $\dot{M}_k$; at $k=0$: $S_{0,y}(A)$, $\Gamma_0(A)$), are not among the operators of (a).
\item[(b)] For every $A\in\tilde{G}_k(T_1)$ the connection $\mathrm{pj}^{*}A$ lies in $\tilde{G}_k(T_2)$, and
\[
O^{(2)}\left(\mathrm{pj}^{*}A\right)\cdot\mathrm{pj}^{*}=\mathrm{pj}^{*}\cdot O^{(1)}(A)
\]
for each of the following operators $O$.
\begin{itemize}
\item[] For $k\ge1$:
$D_A$, $Q_k(A)\cdot J^{*}$, $J\cdot Q_k^{T}(-A)$, $J\cdot P_k\cdot J^{*}$, $J\cdot P_{k+1}\cdot J^{*}$, $Q$, $Q^{T}$, $P$, $\operatorname{Pen}_{y}$, $F_{y}$, $S_{y}$ ($y\in Y_k$; this covers $S_k$, $S_{k,y}$ and $S^{0}_{k+1}$), $H_k$, $\bar{H}_k$, $\mathfrak{H}_k$, $D_k$, $M_k$, $\dot{M}_k$, $B_{k,y}$, $\Gamma_{k,y}$ ($y$ real; this covers $B_k$ and $\Gamma_k$), $T_k$, $\bar{T}_k$, $\delta_{\alpha,A}$, and $T^{(i)}_{\mathbf{Y}}$, $\bar{T}^{(i)}_{\mathbf{Y}}$, where the polymer $\mathbf{Y}$ of $T_1$ is replaced on $T_2$ by the union $\mathrm{pj}^{-1}\mathbf{Y}$ of $M$-cubes.
\item[] For $k=0$:
$D_A$, $Q(A)\cdot J^{*}$, $J\cdot Q^{T}(-A)$, $J\cdot P(A)\cdot J^{*}$, $H_0$, $\bar{H}_0$, $T_0$, $\bar{T}_0$, $\dot{M}_0$, $D_0(A)$, $\operatorname{Pen}_{y}$, $F_{y}$, $S_{0,y}(A)$ ($y$ real), $\Gamma_0(A)$, $\delta_{\alpha,A}$.
\end{itemize}
The numbers $\delta Z_k$ are not concerned by this assertion.
\end{itemize}
\end{named}

The conditions on the constants are those of Proposition~2.8.

\begin{named}{Remark}
Proposition~11 corresponds to \cite[Proposition~16.29]{QEDI} and to \cite[Proposition~5.49]{QEDI}.
The differences are the following.
\begin{itemize}
\item[(1)] The formulas are formulas on $\mathbb{L}_k\times O_{\mathrm{cub}}$, and the deck translations act on the photon site $c$ only.
\item[(2)] The cluster maps $J$ and $J^{*}$ commute with $\mathrm{pj}^{*}$, by (I.J3).
\item[(3)] The lists of (b) contain more operators than those of \cite[Proposition~16.29]{QEDI}.
At $k\ge1$ the list of (b) consists of the operators of the list of Proposition~10\,(b), together with the operators $T^{(i)}_{\mathbf{Y}}$ and $\bar{T}^{(i)}_{\mathbf{Y}}$ of Proposition~9, whose covariances are stated in Proposition~9\,(e);
at $k=0$ it consists of the operators of the list of Proposition~10\,(g), together with $\dot{M}_0=1$, and it is the list (I.L0.18) in its form that contains $\dot{M}_0=1$.
\item[(4)] The map $\mathrm{pj}^{*}$ is defined on functions of pairs as well.
\item[(5)] Clause (a) is the observation (i) of the proof of \cite[Proposition~16.29]{QEDI}, formulated as a clause in the terms of Section~I.1.9.
\item[(6)] The odd integer $J$ of \cite[Proposition~16.29]{QEDI} is written $\mathrm{Jr}$, the letter $J$ denoting the cluster map.
\end{itemize}
\end{named}

\begin{named}{Remark}
In (b) the set $\mathrm{pj}^{-1}\mathbf{Y}$ is a union of $M$-cubes of $T_2$ which need not be connected.
The definition (I.5.4) makes sense for every union of $M$-cubes, and clause (a) is stated for such unions.
\end{named}

\subsection*{12. Integral representation of the logarithm of the one-step determinant (real fields)}

\subsubsection*{Proposition 12}

\begin{named}{Proposition 12}
Let the data be (SF), and let $A\in\tilde{G}_k$ be real.
Traces are formed by (I.ke2) if $k=0$, and by $\operatorname{Tr}O=\sum_{z}\operatorname{tr}O(z,z)$ on $\mathfrak{S}(1,\ell)$ if $k\ge1$; further $n_k=\ell^3$.
Then all the integrals below converge absolutely.
If $k\ge1$, then
\[
\delta Z_k(A)=\exp\left\{c^{\mathrm{d}}_k\cdot n_k-i\,b_k^2\int_0^{\infty}\mathrm{d}y\,\operatorname{Tr}\left[\gamma_5\cdot B_{k,y}(A)\cdot Q_k(A)\cdot J^{*}\cdot S_{k,y}(A)\cdot J\cdot Q_k^{T}(-A)\cdot B_{k,y}(A)\right]\right\}.
\]
If $k=0$, then
\[
\delta Z_0(A)=\exp\left\{\int_1^{\infty}\frac{\mathrm{d}y}{y}\,\operatorname{Tr}\left[F_0(A)\cdot S_{0,y}(A)\right]-i\int_0^1\mathrm{d}y\,\operatorname{Tr}\left[\gamma_5\cdot S_{0,y}(A)\right]\right\}.
\]
Moreover, put $K:=D_k(A)+b_L\cdot P(A)$ if $k\ge1$, a matrix of size $4\,n_k$, and $K:=F_0(A)$ if $k=0$, a matrix of size $192\,n_0$.
Then $K^{*}=\gamma_5\cdot K\cdot\gamma_5$;
the operator $K\cdot\gamma_5$ is invertible and Hermitian with respect to the sesquilinear pairing of the space on which $K$ acts;
$\det\left(K\cdot\gamma_5\right)=\det K$;
and, as $y\to\infty$, one has $\left|\operatorname{Tr}\left[F_0(A)\cdot S_{0,y}(A)\right]\right|=O\left((1+y)^{-1}\right)$, while the trace in the first display is $O\left(\left(b_k^2+y^2\right)^{-1}\right)$.
\end{named}

\begin{named}{Remark}
Proposition~12 corresponds to \cite[Lemma~16.30]{QEDI}, read with the definitions at the beginning of \cite[Section~16.13]{QEDI}.
The differences are the following.
\begin{itemize}
\item[(1)] The matrix $\gamma_5$ stands in the place of the matrix $\gamma_3$ of \cite[Lemma~16.30]{QEDI}, and the cluster maps $J$, $J^{*}$ are written out.
\item[(2)] At $k=0$ the matrix has size $192\,n_0$.
Summed over the eigenvalues, the second of the two elementary identities in the proof of \cite[Lemma~16.30]{QEDI} produces the additive constant $192\,n_0\cdot\frac{i\pi}{2}=96\pi i\,n_0\in2\pi i\mathbb{Z}$, in the place of $2\pi i\,n_0$ there, and $\det\left(\gamma_5\otimes1_{48n_0}\right)=1$.
Hence $c^{\mathrm{d}}_0=0$, and the display for $k=0$ keeps its form (I.L0.10).
\item[(3)] Adjoints and traces refer to $(\cdot,\cdot)^{\mathrm{ch}}$ and to (I.ke2).
\end{itemize}
\end{named}

\subsection*{13. Polymer expansion of the one-step determinant}

\subsubsection*{Proposition 13}

\begin{named}{Proposition 13}
Put
\[
C_d:=16\,c_X\cdot C_w\cdot\max\left\{\frac{\pi}{2}\cdot b\cdot\mathrm{e}\cdot\left(1+2\,\mathrm{e}\right),\ 2+58\,\mathrm{e}+b\cdot\mathrm{e}\right\},
\]
where $C_w$ is the constant of $\mathrm{par}$.
Let the data be (SF), and let a family (p5) of level $k$ on $T_k$ be given, that is, operators $S_{k,y}(X,A)$ on $\mathfrak{S}^{\mathrm{ch}}(\eta_k,\ell)$, for $X\in D_k$, $A\in\tilde{G}_k$ and $y\in[0,\infty)$, Borel measurable in $y$, which have the properties W1, W2, W7, W12, W13 and W14 of Section~G1.11 for $\mathrm{par}=(\alpha_0,\gamma_{\mathrm{w}},\kappa_{\mathrm{w}};C_w)$;
here W7 is taken, when $k=0$, in its form for $k=0$, and W14 is required for $\tilde{g}\in\mathbb{G}^{L}_{k,3}$.
The functions $\acute{E}_k$, $E^{d}_k$, $E_k^{\mathrm{det}}$ and ${\delta E^{\mathrm{det}}}_k$ are formed with this family by the displays of Section~I.5.7.
Then the following hold for $X\in D_k$.
\begin{itemize}
\item[(d0)] For $A\in\tilde{G}_k$ the integrals in the displays of $\acute{E}_k(X,A)$ (Section~I.5.7) converge absolutely, and
\[
\delta Z_k(A)=\exp\left\{c^{\mathrm{d}}_k\cdot n_k+\sum_{X\in D_k}\acute{E}_k(X,A)\right\},\qquad\frac{\delta Z_k(A)}{\delta Z_k(0)}=\exp\left\{\sum_{X\in D_k}E_k^{\mathrm{det}}(X,A)\right\}.
\]
At $k=0$ the determinant $\delta Z_0(A)=\det F_0(A)$ has size $192\,n_0$; it contains the determinant of the $188$ heavy components at each photon site, and it is expanded as a single object.
\item[(d1)] The functions $\acute{E}_k(X,\cdot)$ and $E^{d}_k(X,\cdot)$ are analytic on $\tilde{G}_k$, and for $A\in\tilde{G}_k$ both $\left|\acute{E}_k(X,A)\right|$ and $\left|E^{d}_k(X,A)\right|$ are bounded by
\[
C_d\cdot M^3\cdot\exp\left(-(\kappa+1)\cdot d_M(X)\right).
\]
\item[(d2)] If $A,\tilde{A}\in\tilde{G}_k$ coincide on the bonds that meet $X$, then $\acute{E}_k(X,A)=\acute{E}_k(X,\tilde{A})$ and $E^{d}_k(X,A)=E^{d}_k(X,\tilde{A})$.
\item[(d3)] Let $A\in\tilde{G}_k$. For every real function $\boldsymbol{\lambda}$ of $c$ one has $\acute{E}_k(X,A-d\boldsymbol{\lambda})=\acute{E}_k(X,A)$, and
\[
E^{d}_k(X,A-d\boldsymbol{\lambda})=E^{d}_k(X,A),\qquad E^{d}_k(X,-A)=E^{d}_k(X,A),\qquad\overline{E^{d}_k\left(X,\overline{A}\right)}=E^{d}_k(X,A);
\]
moreover $E^{d}_k(\tilde{g}X,\tilde{g}A)=E^{d}_k(X,A)$ for \emph{every} $\tilde{g}\in\mathbb{G}_k^{L}$.
\item[(d4)] If the side satisfies $\ell\ge M\cdot r_k$, then for $(\mathscr{A},\mathscr{C})\in\tilde{R}_k\times\mathfrak{X}_k$
\[
\left|{\delta E^{\mathrm{det}}}_k(X;\mathscr{A},\mathscr{C})\right|\le8\,C_d\cdot M^3\cdot e_k^{\frac{1}{4}}\cdot\exp\left(-\kappa\cdot d_M(X)\right);
\]
in particular $\left|E_k^{\mathrm{det}}(X,\mathscr{A})\right|$ satisfies the same bound.
\item[(d5)] For data (SF+), the activity $L\,E_k^{\mathrm{det}}$ lies in $K_{k+1}$ and $L\,{\delta E^{\mathrm{det}}}_k$ lies in $K^{x}_{k+1}$ (Sections~I.5.4, I.5.6 and~I.5.11); these are activities without fermion legs on the torus of side $\ell/L$.
\item[(d6)] For data (SF+),
\[
\left\|L\,E_k^{\mathrm{det}}\right\|_{k+1}\le\left\|L\,{\delta E^{\mathrm{det}}}_k\right\|^{x}_{k+1}\le e_{k+1}^{\frac{1}{4}-2\varepsilon}.
\]
\end{itemize}
\end{named}

The conditions on the constants, which belong to the list of Proposition~14.2, are the following:
$\rho<1$ (C0);
(L5) and (M1), which enter through the fact (g0) of \cite[Section~17.4.1]{QEDI}, through $b_->0$ and through $\eta_k\le1$;
(Z2.s1);
the conditions of Proposition~2.8, that is, besides (L5), (M1) and (Z2.s1), the conditions (M2), the part $M\ge108\pi\cdot C'_{V}\cdot a'$ of (M5), (Z2.s2) and $\alpha<1$ (range of the line (C1));
these enter the property W1 and (d0) through Proposition~2.8\,(a), (b), by which $S_{k,y}(A)$ and $S_k(A)$ exist and are analytic on $\tilde{G}_k$, so that $\delta Z_k(A)$ is defined for complex $A$ when $k\ge1$;
they enter (d0) also through Propositions~12 and~10\,(e), (g), and they enter (d3) through Proposition~10\,(b), (g);
in (d4)--(d6), the condition (Z2.s4), which enters through the inequalities (G3) of \cite[Section~16.2.7]{QEDI}, and in (d4) also through its first inequality;
in (d4), and through (d4) also in (d5) and (d6), the condition $r\ge2$ (C7) and the first two inequalities of (A10), namely $v\ge9$ and $8M\cdot\left(\frac{v}{4}+4\right)\cdot\exp(-\varepsilon v)\le1$;
in (d5) and (d6), $\varepsilon\le\frac{1}{4}$ (from (S1)) and $\alpha\ge0$ (range of the line (C1)), which enter through \cite[Lemma~17.5(a), (c)]{QEDI};
in (d5) and (d6), $L\cdot(\kappa-\kappa_0-1)\ge\kappa$ (from (L3)) and $\kappa\ge\kappa_0$ (C2), which enter through the bound of the activity $L\,{\delta E^{\mathrm{det}}}_k$ at a polymer $Y\in D_{k+1}$, weighted with $\exp\left(\kappa\cdot d_M(Y)\right)$, that follows from (d4) by summation over the polymers $X$ with $\bar{X}^{L}=LY$;
the clause (k6) of Section~I.5.4 in (d5) follows from this bound, and (d6) is derived from it;
in (d6), (Z2.s5) with the above $C_d$, and $\varepsilon\le\frac{1}{8}$ (from (S1));
finally $\kappa_{\mathrm{w}}-1\ge\kappa+1$.

\begin{named}{Remark}
Proposition~13 corresponds to \cite[Proposition~16.31]{QEDI} together with the definitions that precede it in \cite[Section~16.13]{QEDI}, and to \cite[Proposition~5.48]{QEDI}; the functions $\acute{E}_k$, $E^{d}_k$, $E_k^{\mathrm{det}}$, ${\delta E^{\mathrm{det}}}_k$ are those of \cite[Definitions~5.26 and~5.28]{QEDI}.
The differences are the following.
\begin{itemize}
\item[(1)] The matrix $\gamma_5$ and the cluster maps $J$, $J^{*}$ occur, and traces are taken by (I.ke2).
\item[(2)] In $C_d$ the second entry of the maximum is $2+58\,\mathrm{e}+b\cdot\mathrm{e}$ instead of $5+6\,\mathrm{e}+b\cdot\mathrm{e}$; it is the row sum (I.L0.12) at $\exp\left(\mathfrak{i}(A)\right)\le\mathrm{e}$, increased by $1$.
With the kernels (I.1.4) the number of sites of $X$ is $\int_X\mathrm{d}x=M^3\cdot|X|_M$, and no further factor $48$ enters, because this factor is contained in the constant $C_w$ of W7 (I.L0.26).
Since $\left|\operatorname{tr}\left(m\cdot m'\right)\right|$ is at most the sum of the row sums of $m$ times the largest modulus of an entry of $m'$, the factor $16$ in $C_d$ could be replaced by $4$.
\item[(3)] Like \cite[Proposition~5.48]{QEDI}, Proposition~13 is stated for $\mathrm{par}$ only, with the $C_w$ of $\mathrm{par}$ in $C_d$ (Section~I.8.4).
In this it differs from \cite[Proposition~16.31]{QEDI}, in which the numbers $\gamma_{\mathrm{w}}>0$ and $C_w>0$ of the parameter set $(\alpha_0,\gamma_{\mathrm{w}},\kappa_{\mathrm{w}};C_w)$ are arbitrary and $C_d$ is formed with the $C_w$ at hand.
\item[(4)] The hypotheses (Z2.s4), $v_k\ge9$, $8M\cdot\left(\frac{v_k}{4}+4\right)\cdot e_k^{\varepsilon}\le1$, $L\cdot(\kappa-\kappa_0-1)\ge\kappa$ and (Z2.s5) of \cite[Proposition~16.31(e), (g)]{QEDI} are conditions on the constants and are imposed in Proposition~14.2, while $\ell\ge M\cdot r_k$, $\ell/L\in\mathrm{Cl}$ and $e_{k+1}\le e_0$ are data.
\item[(5)] The numbers $r_k$ and $v_k$ are formed with the $e_k$ of the data (Section~G1.1).
\item[(6)] In the Borel measurability in $y$ and in the factor $(1+y)^{-1}$ at $k=0$, the data of Proposition~13 agree with the hypotheses of \cite[Propositions~5.48 and~16.31]{QEDI}.
There too the family (p5) is Borel measurable in $y$ at every level, the level $0$ included (\cite[Definition~5.29]{QEDI}), and at $k=0$ property~7 of \cite[Section~5.4.2]{QEDI} is the kernel bound with the factor $(1+y)^{-1}$, the form in which it enters the hypotheses of these two propositions.
Here the bound W7 at $k=0$ is read for the kernel (I.1.4) at pairs of fermion sites and in the norm $|\cdot|_c$ (Section~G1.11).
\item[(7)] In (d2) at $k=0$: for $x$, $x'$ over $X$ the entry $F_0(A)(x,x')$ reads $A$ on the bond $(c,c_3)$ or inside the $L$-block of $c$, and this block lies in $X$ by the fact (g0) of \cite[Section~17.4.1]{QEDI}.
\item[(8)] In (d3) the property W14 is assumed for $\mathbb{G}^{L}_{k,3}$ only; but $\tilde{g}X$ and $\tilde{g}A$ depend on $(r(S),a)$ only, and $r(\mathrm{Sp}_3)=\{\pm1\}^3$.
\end{itemize}
\end{named}

\subsection*{19. The block covariance at real fields}

\subsubsection*{Proposition 19}

\begin{named}{Proposition 19}
Let $\rho\in(0,1)$, and let the weights be the standard ones (Section~I.0.5).
Let $\ell\in\mathrm{Cl}$ be a side, $e_k>0$ a number, $\mu>0$ a real number and $A$ a real bond function on $T_k$;
in (a) let $k\ge0$ be an integer, and in (b)--(f) let $k\ge1$.
No upper bound is imposed on $e_k$, on $\mu$ (except where one is stated) or on $A$, and no smallness is required.
The space $\mathfrak{S}^{\mathrm{ch}}(\eta_k,\ell)$ carries the scalar product $(\cdot,\cdot)^{\mathrm{ch}}$ and $\mathfrak{S}(1,\ell)$ the counting scalar product, both formed with $|\cdot|_2$ on $\mathbb{C}^4$;
$|\cdot|_{\mathrm{op}}$ denotes the operator norm with respect to these scalar products;
for a $4\times4$ block $m$, $|m|_2$ is its Euclidean operator norm and $|m|_c$ its largest row sum (Section~I.1.5).
Throughout, $y$ and $y'$ are unit sites.
\begin{itemize}
\item[(a)] (Every level $k\ge0$; at $k=0$ one has $\eta_0=1$, the charge is $q$, and the space is $\mathfrak{S}^{\mathrm{ch}}(1,\ell)$.)
For every $f$,
\[
\operatorname{Re}\left(f,(D_A+\mu)f\right)^{\mathrm{ch}}=\mu\cdot(f,f)^{\mathrm{ch}}+\frac{\rho}{2\eta_k}\cdot\left(f,\nabla_A^{*}\nabla_A f\right)^{\mathrm{ch}}\ge\mu\cdot(f,f)^{\mathrm{ch}}.
\]
Consequently $D_A+\mu$ is invertible, $\left|(D_A+\mu)^{-1}\right|_{\mathrm{op}}\le\frac{1}{\mu}$, $\operatorname{Re}(D_A+\mu)^{-1}\ge0$ and $\det(D_A+\mu)\ne0$.
\item[(b)] One has
\begin{gather*}
\left(Q_k(A)\cdot J^{*}\right)^{*}=J\cdot Q_k^{T}(-A),\\
\left(Q_k(A)\cdot J^{*}\right)\cdot\left(J\cdot Q_k^{T}(-A)\right)=1,\\
\left|Q_k(A)\cdot J^{*}\right|_{\mathrm{op}}=\left|J\cdot Q_k^{T}(-A)\right|_{\mathrm{op}}=1,
\end{gather*}
and $J\cdot P_k(A)\cdot J^{*}$ is an orthogonal projection.
\item[(c)] The operator $D_A+\mu+b_k\cdot J\cdot P_k(A)\cdot J^{*}$ has real part $\ge\mu$ and is invertible; hence $S_k(A)$ and $D_k(A)$ are defined.
\item[(d)] The operator $D_k(A)$ is invertible, and
\[
D_k(A)^{-1}=b_k^{-1}+Q_k(A)\cdot J^{*}\cdot(D_A+\mu)^{-1}\cdot J\cdot Q_k^{T}(-A),\qquad\operatorname{Re}D_k(A)^{-1}\ge b_k^{-1},
\]
\[
\left|D_k(A)^{-1}\right|_{\mathrm{op}}\le b_k^{-1}+\mu^{-1},\qquad\left|D_k(A)\right|_{\mathrm{op}}\le b_k.
\]
\item[(e)] Let $\mu\le12\,\mathrm{e}$. Then
\[
\left|D_k(A)^{-1}(y,y')\right|_2\le b_k^{-1}\cdot\mathbf{1}[y=y']+\frac{2\,\mathrm{e}}{\mu}\cdot\exp\left(-\frac{\mu}{12\,\mathrm{e}}\cdot d(y,y')\right),
\]
and hence
\[
\left|D_k(A)^{-1}(y,y')\right|_c\le2\,b_k^{-1}\cdot\mathbf{1}[y=y']+\frac{4\,\mathrm{e}}{\mu}\cdot\exp\left(-\frac{\mu}{12\,\mathrm{e}}\cdot d(y,y')\right);
\]
every entry of the block $D_k(A)^{-1}(y,y')$ has modulus at most the first of these two bounds.
\item[(f)] Let $\mu\le12\,\mathrm{e}$. Then every row and every column of the matrix of $D_k(A)^{-1}$, whose rows and columns are indexed by the pairs $(y,\beta)$, has $\ell^1$ norm at most
\[
2\,b_k^{-1}+4\,\mathrm{e}\cdot\mu^{-1}\cdot\left(\frac{144\,\mathrm{e}}{\mu}\right)^{3}.
\]
\end{itemize}
All these bounds are uniform in $A$, in particular in the twist, which is a real constant connection in the first slot (Section~I.6.6).
\end{named}

\begin{named}{Remark}
Proposition~19 corresponds to the opening paragraphs of \cite[Section~13.10]{QEDI} and to \cite[Theorem~13.18]{QEDI}, of which \cite[Proposition~5.65]{QEDI} is a special case.
The differences are the following.
\begin{itemize}
\item[(1)] In \cite[Section~13.10]{QEDI} the operator $D_A$ is written in the form
\[
D_A=\eta^{-1}\sum_{\mu=0}^{2}\sum_{\boldsymbol{\sigma}=\pm1}P^{\boldsymbol{\sigma}}_{\mu}\cdot\left(1-T^{\boldsymbol{\sigma}}_{\mu}(A)\right)
\]
with orthogonal projections $P^{\boldsymbol{\sigma}}_{\mu}$, the index $\mu$ being a direction there.
This representation does not hold for the operator $D_A$ considered here; its place is taken by (I.2.3) and (I.2.4), which hold for every $k\ge0$, every side, every complex $A$ and every charge.
\item[(2)] The operators $Q_k(A)\cdot J^{*}$, $J\cdot Q_k^{T}(-A)$ and $J\cdot P_k(A)\cdot J^{*}$ stand in the places of $Q_k(A)$, $Q_k^{T}(-A)$ and $P_k(A)$, and $(\cdot,\cdot)^{\mathrm{ch}}$ in the place of the scalar product of \cite[Section~13.10]{QEDI}.
The product in (b) equals $1$ because $\Delta_y\times O_{\mathrm{cub}}$ has $48\,\eta_k^{-3}$ sites, each of weight $\eta_k^3/48$.
\item[(3)] In the decay estimate, the six terms with $\left|P^{\boldsymbol{\sigma}}_{\mu}\right|=\left|T^{\boldsymbol{\sigma}}_{\mu}\right|=1$ of the proof of \cite[Theorem~13.18]{QEDI} are replaced by a single term: for a real function $\psi$ of $c$,
\[
\left(\exp(\psi)\cdot D_A\cdot\exp(-\psi)\,f-D_A f\right)(x)=-\eta_k^{-1}\cdot M_{3}(x)\cdot u_A(c,c_3)\cdot\left(\exp\left(\psi(c)-\psi(c_3)\right)-1\right)\cdot f\left(x^{(3)}\right),
\]
where $\left|M_{3}(x)\right|_2=3\,(1+\rho)\le6$, the map $x\mapsto x^{(3)}$ is a bijection that preserves the measure, and $|c_3-c|=\eta_k$.
Therefore $6\,\eta_k^{-1}\cdot\left(\exp\left(\frac{\mu\,\eta_k}{12\,\mathrm{e}}\right)-1\right)\le\frac{\mu}{2}$, as in \cite[Theorem~13.18]{QEDI}, and the numbers $12\,\mathrm{e}$, $2\,\mathrm{e}$ and $144\,\mathrm{e}$ are the same.
\item[(4)] The identity for $\operatorname{Re}(D_A+\mu)$, the bound $\operatorname{Re}D_k(A)^{-1}\ge b_k^{-1}$ and the level $0$ in (a) are stated as clauses; in the proof of \cite[Theorem~13.20]{QEDI} the level $0$ is reached through the observation that the first lines of the proof of \cite[Theorem~13.18]{QEDI} hold at level $0$ as well.
\item[(5)] The norms of a block are the norms $|\cdot|_2$ and $|\cdot|_c$ of Section~I.1.5, and $|\cdot|_2$ is the norm $|B|$ of a $4\times4$ matrix $B$ in \cite[Section~13.10]{QEDI}.
The first bound of (e) is the bound on the blocks in \cite[Theorem~13.18]{QEDI};
new in (e) are the bound in $|\cdot|_c$, whose two constants are doubled, and the bound on the entries.
Clause (f) follows from the bound in $|\cdot|_2$, since a row or a column of a $4\times4$ block $m$ has $\ell^1$ norm at most $2\,|m|_2$;
its bound is the bound on the rows and the columns in \cite[Theorem~13.18]{QEDI}.
\item[(6)] Here $e_k$ is any positive number and $\mu$ any positive real number, with $\mu\le12\,\mathrm{e}$ in (e) and (f).
In \cite[Section~13.10]{QEDI} the setting is the setting (Sc) of \cite[Section~13.1.1]{QEDI}, that is, $k\ge1$, $\ell\in\mathrm{Cl}$ and $e_k\in(0,e_0]$, with the one change that the mass, which (Sc) bounds by a constant at most $1$, may be any real number $>0$;
\cite[Theorem~13.18]{QEDI} assumes a real mass $>0$ and a real bond function $A$, and a mass at most $12\,\mathrm{e}$ for the bound on the blocks and for the sums over the rows and the columns.
Thus the range of $\mu$ is that of \cite[Section~13.10]{QEDI} and of \cite[Theorem~13.18]{QEDI}, and the ranges differ only in that $e_k$ is not restricted to $(0,e_0]$.
\item[(7)] The parameter $\rho$ enters only through $\rho\ge0$ (accretivity) and $\rho\le1$ (the number $6$).
\end{itemize}
\end{named}

The mass $\mu$ of Proposition~19 is the running mass $\bar{m}\cdot a_k>0$ of the propagators, and not a lattice mass $\boldsymbol{\mu}$.

\subsection*{G4. Definitions for the Propositions 26.m--32.m and 49.m that are not in Part~I}

\subsubsection*{G4.0.0. Notation; the word admissible}

In the Propositions 26.m--32.m and 49.m several letters of \cite[Section~19]{QEDI} are replaced by other letters.
The isomorphism $\mathfrak{i}$ of \cite[Definition~19.9(c)]{QEDI} is written $\mathrm{idn}$, since $\mathfrak{i}(A)$ denotes $e_k\cdot|\operatorname{Im}A|_{\infty}$.
The admissible set $\tilde{V}$ of \cite[Lemma~19.10(c)]{QEDI} and of Step~0 of the proof of \cite[Proposition~19.12]{QEDI} is written $\tilde{Y}$, since $V_{\mathrm{t}}$ denotes a Hodge part.
The union $W$ of $M$-cubes in \cite[Section~19.3.4]{QEDI} is written $Z$.
The remainder $\vartheta$ of \cite[Proposition~19.22]{QEDI} is written $\mathrm{vrem}$; in \cite[Section~20.1.7]{QEDI} the letter $\vartheta$ denotes a rate.
The objects $R_{\theta}$, $D(c)$ and $T_c$ of \cite[Section~19.8.2]{QEDI} are written $\mathrm{Rtw}_{\theta}$, $\mathrm{Dfl}_{\theta}$ and $T_{\theta}$, and their values at $\theta=0$ carry the index $\theta=0$, because $T_0(A)$ and $D_0(A)$ denote operators of level $0$.
The letters $z$ and $\boldsymbol{\sigma}$ of Step~3 of the proof of \cite[Proposition~19.12]{QEDI} are written $\mathrm{zval}$ and $\mathrm{sval}$;
an element of $\ell_1\mathbb{Z}^3$, denoted $g$ in \cite[Definition~19.9]{QEDI}, is written $v$;
and the numbers $a_1$, $a_2$ of \cite[Proposition~5.64]{QEDI} are written $\mathrm{aw}_1$, $\mathrm{aw}_2$, as in Section~G2.
Furthermore, $k_{\mathrm{up}}$ denotes a level (Section~G4.0.5), $D^{(p';\mathrm{ma})}$ is the symbol (I.2.6) with its two arguments displayed, and in Proposition~30.3 the symbol $[c]$ denotes the centre of the unit block of $c$.

The word admissible has two meanings in the Propositions 27.m:
a pair is admissible for a polymer in the sense of Section~I.1.7, and a union of $M$-cubes of the covering lattice is admissible in the sense of Proposition~27.1, as in \cite[Definition~19.9(a)]{QEDI}.

A superscript $\mathrm{ch}$ at the name of a property of \cite[Section~5.4.2]{QEDI}, as in W7$^{\mathrm{ch}}$, indicates that the property is read with the conventions of Part~I.

\subsubsection*{G4.0.1. Three kinds of data}

\begin{description}
\item[(SF)] (Scale-free data; Section~I.1.8.)
A quadruple $(k,\ell,\mu,e_k)$ consisting of an integer $k\ge0$, a side $\ell\in\mathrm{Cl}$, a mass $\mu\in(0,1]$ and a number $e_k\in(0,e_0]$, with no relation between them.
\item[(OS)] (Data of one step; compare the hypothesis of \cite[Definition~5.69]{QEDI}.)
A quadruple $(k,\ell,\mu,e_k)$ such that
$k\ge0$;
$\ell$ and $\ell/L$ lie in $\mathrm{Cl}$;
$\mu\in(0,1/L]$;
$e_k>0$ and $e_{k+1}:=L^{\frac{1}{2}}\cdot e_k\le e_0$;
$v_k:=-\log e_k$, $r_k$, $p_k$, $p_{0,k}$, $w^{\mathrm{m}}_k$ and $w^{\mathrm{E}}_k$ are formed with $e_k$ (Sections~I.6.1 and~I.5.4);
$\ell\ge M\cdot r_k$;
together with a system of primitives of level $k$ on $T_k$ (Section~G4.0.4) that has the properties W1--W17$^{\mathrm{ch}}$ for $\mathrm{par}$.
The level $k+1$ has the side $\ell/L$, the mass $L\mu$ and the charge $e_{k+1}$.
\item[(FL)] (Flow data; compare the hypothesis of \cite[Definition~5.70]{QEDI}.)
A quadruple $(\bar{m},e,R,N)$ with $\bar{m}\in(0,1]$, with $e>0$ such that $e_K=e\cdot h^{\frac{1}{2}}\le e_1$, with $R\in\mathrm{Tor}(e)$, and with an integer $N\ge N_{\min}(R,e)$.
These four determine $\eta=L^{-N}$, $q=e\cdot\eta^{\frac{1}{2}}$, $n:=R/\eta$, $a_k$, $e_k$, $v_k$, $\mu_k$, $K$, $r_k$, $p_k$, $p_{0,k}$, $\ell_k=\ell_k(R)$, $n_k$, $\mathfrak{n}_k$, $k_{*}=k_{*}(R)$, $k_{\mathrm{top}}$ and the index $*$ (Sections~I.6.1 and~I.6.2), and one puts $T_k(R):=\mathbb{L}_k/\ell_k(R)\mathbb{Z}^3$.
The primitives that belong to (FL) are, on $T_k(R_{\mathrm{top}})$ for $k<k_{\mathrm{top}}$, the systems of Proposition~26.5, and, on $T_k(R)$ for $R<R_{\mathrm{top}}$ and $k<k_{*}(R)$, their projections (Proposition~26.1) with $\mathrm{Jr}=R_{\mathrm{top}}/R$.
By Proposition~26.6\,(e), these primitives are defined for every quadruple that satisfies the four conditions.
\end{description}

\subsubsection*{G4.0.2. Scales}

The scales are those of \cite[Section~19.2]{QEDI}, read with $h'=h$, $m_0=0$ and with $1$ as the upper end of the mass range (Section~I.6.1).
There is one exception: in the place of \cite[Proposition~19.2]{QEDI} stands the assertion (Sc6) below, which is Proposition~14.2\,(c), (f).

Throughout this section (S-4) is assumed, that is, (FL) without the primitives (Section~I.1.8):
$\bar{m}\in(0,1]$;
$e>0$ with $e_K=e\cdot h^{\frac{1}{2}}\le e_1$;
$R\in\mathrm{Tor}(e)$;
and $N\ge N_{\min}(R,e)$ is an integer.
The assertions (Sc0)--(Sc6) below involve only the numbers $\bar{m}$, $e$, $R$, $N$, the constants and the scales formed from them; they hold under (S-4) alone, before the primitives of (FL) are defined.
The assertion of (Sc5) that $R^{*}$ and $R_{\mathrm{top}}$ belong to $\mathrm{Tor}(e)$ and the first three assertions of (Sc6) do not involve $R$ or $N$; they hold for every $e>0$ with $e_K\le e_1$.

\begin{itemize}
\item[(Sc0)] Of the data, only $\bar{m}\in(0,1]$ and $e>0$ with $e\cdot h^{\frac{1}{2}}\le e_1$ enter here; no torus and no $N$ are given.
One has
\[
e\le\frac{1}{3},\qquad v_K\ge 8,\qquad \left(1+\frac{\mathrm{lg}}{v_K}\right)^{r}\le\left(1+\frac{3\,\mathrm{lg}}{v_K}\right)^{r}<L .
\]
The sides $R^{*}(e)$ and $R_{\mathrm{top}}(e)=L\cdot R^{*}(e)$ belong to $\mathrm{Tor}(e)$, because $M$ is a power of $L$ with $M\ge L^2$ and $r_K$ is a power of $L$ with an exponent $\ge 0$.
For every $R=h\cdot L^{T}\in\mathrm{Tor}(e)$ the number $N_{\min}(R,e)$ of Section~I.6.2 exists, and
\[
N_{\min}(R,e)=\max\left\{0,\ \mathrm{js}(R,e)+1-\log_L h\right\},
\]
where $\mathrm{js}(R,e)\ge 0$ is the integer $j_{*}(R,e)$ of \cite[Proposition~19.1(b)]{QEDI}; it is formed with $\left(e\cdot h^{\frac{1}{2}},L,M,r,T\right)$ and does not contain $N$.
Consequently $N_0(e)=N_{\min}\left(R^{*}(e),e\right)$ and $N_{\min}\left(R_{\mathrm{top}}(e),e\right)$ are defined, and the conditions of (FL) can be formed at $R^{*}(e)$ and at $R_{\mathrm{top}}(e)$.
The inequalities of the first display are the first three assertions of (Sc6), that is, Proposition~14.2\,(f) and the row (e1) of Proposition~14.2\,(b) at $v=v_K$ (see the comment on (Sc6) below).
The remaining assertions of (Sc0) are \cite[Proposition~19.1(b), (c)]{QEDI} and the opening lines of the proof of \cite[Proposition~19.1(g)]{QEDI}.

\item[(Sc1)] One has $1\le k_{*}\le K$.
The numbers $K-k_{*}$, $a_{*}$, $e_{*}=e\cdot a_{*}^{\frac{1}{2}}$, $\ell_*=M\cdot r_{*}$ and $\mu_*=\bar{m}\cdot a_{*}$ do not depend on $N\ge N_{\min}(R,e)$.
This is \cite[Proposition~19.1(c)]{QEDI}.

\item[(Sc2)] Let $R\le R'$ be elements of $\mathrm{Tor}(e)$.
Then $N_{\min}(R',e)\le N_{\min}(R,e)$, and $k_{*}(R)\le k_{*}(R')$ for $N\ge N_{\min}(R,e)$.
In particular $N\ge N_{\min}(R_{\mathrm{top}},e)$ and $k_{*}(R)\le k_{\mathrm{top}}$.
This is \cite[Proposition~19.1(d)]{QEDI}.

\item[(Sc3)] For $k\le k_{*}$ one has
\[
\ell_k\ge M\cdot r_k,\qquad r_k\ge r_{*},\qquad v_k\ge v_{*}\ge v_K,\qquad e_k\le e_{*}\le e_K\le e_{*}\cdot\ell_*^{\frac{1}{2}},
\]
and for $k<k_{*}$ one has $\ell_k\ge L\cdot M\cdot r_k$.
Further, $\mu_{k+1}=L\cdot\mu_k$, and $0<\mu_k\le\mu_*\le\mu_K$ for $k\le k_{*}$.
Finally,
\begin{gather*}
\frac{1}{L}<\mu_K\le 1,\qquad \mu_*\cdot\ell_*=\bar{m}\cdot R\ge\mu_K,\qquad \mu_*^{-1}<L\cdot\ell_*,\\
M\cdot v_{*}^{r}\le\ell_*<L\cdot M\cdot v_{*}^{r},\qquad e^2\cdot R=e_{*}^2\cdot\ell_* .
\end{gather*}
This is \cite[Proposition~19.1(e)]{QEDI}.

\item[(Sc4)] For every $k\le k_{*}(R)$ the side $\ell_k(R)$ lies in $\mathrm{Cl}$.
This is \cite[Proposition~19.1(f)]{QEDI}, read with the hypothesis $N\ge N_{\min}(R,e)$.
The hypothesis of \cite[Proposition~19.1(f)]{QEDI}, namely the condition (2.1) of Section~G4.0.3, $r_{*}>31L+16$, is supplied by (Sc6), that is, by Proposition~14.2\,(c), under the datum $e\cdot h^{\frac{1}{2}}\le e_1$; thus (Sc4) rests on (Sc6).

\item[(Sc5)] One has $\mathrm{js}(R^{*},e)\in\{1,2\}$ and $\mathrm{js}(R_{\mathrm{top}},e)\in\{0,1\}$.
Consequently $k_{*}(R^{*})\in\{K-1,K-2\}$ for $N\ge N_0(e)=N_{\min}\left(R^{*}(e),e\right)$, and $k_{\mathrm{top}}\in\{K,K-1\}$ for $N\ge N_{\min}(R_{\mathrm{top}},e)$; for $N\ge N_0(e)$ one has on $R^{*}$
\[
a_{*}\in\left\{\frac{h}{L},\frac{h}{L^2}\right\},\qquad r_{*}\in\left\{r_K,\ L\cdot r_K\right\}.
\]
Under (S-4) the condition $N\ge N_{\min}(R_{\mathrm{top}},e)$ holds by (Sc2).
If $R\ne R_{\mathrm{top}}$, then $R\le R^{*}$, since the elements of $\mathrm{Tor}(e)$ are the numbers $h\cdot L^{T}\le R_{\mathrm{top}}=L\cdot R^{*}$ with an integer $T\ge 0$, and $R^{*}$ is one of them; in this case $N\ge N_{\min}(R,e)\ge N_0(e)$ by (Sc2).
If $R=R_{\mathrm{top}}$, the condition $N\ge N_0(e)$ is a hypothesis in addition to (S-4).
The sides $R^{*}$ and $L\cdot R^{*}=R_{\mathrm{top}}$ belong to $\mathrm{Tor}(e)$.
This is \cite[Proposition~19.1(g)]{QEDI}.

\item[(Sc6)] One has
\[
e\le\frac{1}{3},\qquad v_K\ge 8,\qquad \left(1+\frac{3\,\mathrm{lg}}{v_K}\right)^{r}<L,
\]
and the conditions of Section~G4.0.3 hold for every $R\in\mathrm{Tor}(e)$ and every $N\ge N_{\min}(R,e)$.

This is \cite[Proposition~19.2]{QEDI}, read with the constants and the number $e_1$ of Section~I.8.2; in this reading it is Proposition~14.2\,(c), (f).
Its third assertion is stated in Proposition~14.2\,(f); it is the row (e1) of Proposition~14.2\,(b) at $v=v_K\ge-\log e_1$.
That the inequalities (e1)--(e8) hold for all $v\ge-\log e_1$ is the first sentence of Proposition~14.2\,(c), and it is also contained in Proposition~14.3\,(f).
It follows from the row $e_1$ of Proposition~14.2\,(a) together with the datum $0<e_K\le e_1$.
Indeed, the supremum in that row equals $2e_1>e_1$; hence there is a number $e'>e_1$ such that (e1)--(e8) hold for all $v\ge-\log e'$, and $-\log e'<-\log e_1$.
\end{itemize}

\subsubsection*{G4.0.3. Conditions on $(R,e,N)$ and error terms}

The conditions and the error terms below have the form of those of \cite[Sections~5.6.2 and~19.1.2]{QEDI}; they are formed with the constants of Section~I.8.2 and with the changes listed in Section~I.8.4.
In particular, the rate $\gamma_{\Sigma}$ enters (T5) and $\varepsilon_2$ where \cite[Sections~5.6.2 and~19.1.2]{QEDI} have $\gamma_1$, and $\varepsilon_5$ is the function of the line (C0r) of Section~I.8.2, which takes the place of the function $x\mapsto 312\cdot\log\left(\left(1-\exp\left(-\frac{x}{4}\right)\right)^{-1}\right)$ of \cite[Section~5.6.2]{QEDI}.
With $p_{0*}:=p_{0,k_{*}}$, the conditions read
\begin{align}
&r_{*}>31L+16;\tag{2.1}\\
&\mathrm{sm}_k:=\left(c_{\mathrm{fl}}+2\right)\cdot e_k^{\varepsilon},\qquad \sum_{k=0}^{k_{*}}\mathrm{sm}_k\le\frac{1}{16};\qquad v_{*}\ge 8;\tag{fl}\\
&v_{*}\ge 8,\qquad p_0\cdot v_{*}^{p_0-1}\ge 28,\qquad w^{\mathrm{m}}_0<1\quad\left(w^{\mathrm{m}}_0=q^{\frac{3}{4}-8\varepsilon}\right);\tag{T1}\\
&\begin{aligned}[t]
&8\,C_{0}^2\cdot\ell_*^6\cdot\left(b_{-}^{-1}+\mu_*^{-1}\right)\le e_{*}^{-\frac{\varepsilon}{2}},\qquad 4\left(2C_{0}+c_{\kappa}+5c_{\mathrm{fl}}+6\right)\cdot\ell_*^3\le e_{*}^{-\varepsilon},\\
&e_{*}^{\frac{1}{4}-9\varepsilon}\le\frac{1}{8};
\end{aligned}\tag{T2}\\
&4\,C_{0}\cdot\ell_*^2\cdot p_{*}<e_{*}^{-\varepsilon};\tag{T3}\\
&12\,\ell_*^3\cdot\left(\exp\left(-\frac{p_{0*}}{2}\right)+\exp\left(-\frac{p_{*}^2}{2C_{0}^2}\right)\right)\le e_{*};\tag{T4}\\
&C_{\Sigma}\cdot\exp\left(-\gamma_{\Sigma}\cdot\ell_*\right)\le e_{*};\tag{T5}\\
&L^2\cdot e^2\cdot R\le 1,\tag{T6}
\end{align}
and the error terms are
\begin{align*}
&\varepsilon_1:=\frac{1}{2}\,e_{*}^{\frac{1}{4}-9\varepsilon},\qquad \varepsilon_2:=C_{\Sigma}\cdot\ell_*^{-1}\cdot\exp\left(-\gamma_{\Sigma}\cdot\ell_*\right),\\
&\varepsilon_3:=6\,\ell_*^3\cdot\exp\left(-\frac{p_{0*}}{2}\right)+6\,\ell_*^3\cdot\exp\left(-\frac{p_{*}^2}{2C_{0}^2}\right),\qquad \varepsilon_4:=6\,\ell_*^3\cdot\exp\left(-\frac{p_{*}^2}{2C_{0}^2}\right),\\
&\varepsilon_5(x):=C_{6}\cdot\log\left(\frac{1}{1-\exp\left(-\kappa_{(0)}\cdot x\right)}\right)\qquad(x>0).
\end{align*}
Here $\varepsilon_5$ is a decreasing function on $(0,\infty)$, not a number, and $C_{6}$, $\kappa_{(0)}$ in its definition are the constants of the line (C0r) of Section~I.8.2.
The number $\kappa_{(0)}=\min\left\{\frac{\theta_{0}}{2},\frac{c_{\mathrm{el}}\cdot\theta_{0}}{12K_{\mathrm{Str}}}\right\}$ is a constant other than the number $\kappa_{0}=54\log 72$ of the line (C0), which enters $\mathrm{par}^{0}$ in Section~G4.0.4.

These are conditions on the triple $(R,e,N)$: the condition (fl) and the number $w^{\mathrm{m}}_0$ contain $N$.
By (Sc6) they are consequences of $e_K\le e_1$.
The number $e_1$ is defined in Proposition~14.2\,(a), (b) by the inequalities (e1)--(e8) of \cite[Section~19.1.2]{QEDI}, formed with the constants $C_{0}$, $c_{\kappa}$, $\gamma_{\Sigma}$, $C_{\Sigma}$ of Section~I.8.2 and with $L\cdot\mathrm{lh}(v)$ in the place of $L^{m_0+1}\cdot\mathrm{lh}(v)$ in (e3); here $\mathrm{lh}(v)=L\cdot M\cdot v^{r}$ (Proposition~14.2\,(a)).

\subsubsection*{G4.0.4. Systems of primitives; the properties W1--W17$^{\mathrm{ch}}$}

The systems of primitives and their properties are those of \cite[Sections~5.4.1 and~5.4.2]{QEDI}, read as in Section~I.5.8; the properties W1--W17 are the properties 1--17 of the list in \cite[Section~5.4.2]{QEDI}.
For the fermionic families they are established in Proposition~7 ($k\ge 1$) and Proposition~8 ($k=0$); for the families (p1), (p2) they are \cite[Theorem~15.35 and Proposition~15.36]{QEDI}.

Let $(k,\ell,\mu,e_k)$ be data (SF).

\emph{Systems of level $k\ge 1$ on $T_k$.}
Such a system consists of families of linear operators, defined for $s\in\mathrm{Pd}_k$, $A\in\tilde{G}_k$, $X\in D_k$ and for the resolvent parameter $y\in[0,\infty)$:
\begin{itemize}
\item[(p1)] $H_{k}(s')$, $s'=\left(s'_{\mathrm{pc}}\right)_{\mathrm{pc}\in\pi'_k}$, from the bond functions on $T^0_k$ to the bond functions on $T_k$;
\item[(p2)] $C_{k}^{1/2}(s')$, $s'=\left(s'_{\mathrm{pc}}\right)_{\mathrm{pc}\in\pi''_k}$, on $V_{k}$ (complexified for complex $s'$);
\item[(p3)] $H_{k}(s'',A)$, $s''=\left(s''_{\mathrm{pc}}\right)_{\mathrm{pc}\in\pi_k}$, from $\mathfrak{S}(1,\ell)$ to $\mathfrak{S}^{\mathrm{ch}}(\eta_k,\ell)$, together with $\bar{H}_{k}(s'',A):=\mathsf{C}^{-1}\cdot H_{k}(s'',-A)\cdot\mathsf{C}$;
\item[(p4)] $\Gamma_k(s'',A)$, in the same variables $s''$, and $M_k(X,A)$, both on $\mathfrak{S}(1,\ell)$;
\item[(p5)] $S_{k,y}(X,A)$ on $\mathfrak{S}^{\mathrm{ch}}(\eta_k,\ell)$, Borel in $y$.
\end{itemize}

\emph{Systems of level $0$.}
A system of level $0$ consists of a family (p2) and of the operators $\Gamma_0(s'',A)$ and $S_{0,y}(X,A)$ on $\mathfrak{S}^{\mathrm{ch}}(1,\ell)$, indexed as in (p4), (p5), the latter Borel in $y$.
The families (p1), (p3) of level $0$ are the diagonal families $H_{0}(s')$, $H_{0}(s'',A)$ of Section~I.5.8.
There is no family $M_0$.

\emph{Indices.}
An index $t$ is a bond of $T_k$ (a fine bond $b$), a bond of $T^0_k$ (a unit bond $b'$), a basis element $\Upsilon$, a unit site $z$, or a fermion site $x=(c,g)$.
Its support $\operatorname{supp}t$ is the set of the end points if $t$ is a bond, the set of the end points of the bonds in the support of $\Upsilon$ if $t=\Upsilon$, the set $\{z\}$ if $t=z$, and the set $\{c\}$ if $t=x=(c,g)$ (I.POS).
The relations ``$t$ is contained in $Z$'' and ``$t$ meets $Z$'', the distance $d(t,t')$ and the set $\square(t)$, the union of the $M$-cubes that $t$ meets, are those of \cite[Section~5.2.3]{QEDI}, formed with $\operatorname{supp}t$.
All supports have sup-diameter at most $L$.

\emph{Kinds, rows and columns.}
The family (p1) is of kind $\pi'_k$; its rows are fine bonds and its columns unit bonds.
The family (p2) is of kind $\pi''_k$; its rows and columns are basis elements.
The family (p3) is of kind $\pi_k$; its rows are the fermion sites of $T^{\mathrm{f}}_k$, and its columns are unit sites for $k\ge 1$ and fermion sites of $T^{\mathrm{f}}_0$ for $k=0$.
The family $\Gamma_k$ is of kind $\pi_k$; its rows and columns are unit sites for $k\ge 1$ and fermion sites of $T^{\mathrm{f}}_0$ for $k=0$.
The rows and columns of $M_k(X,A)$ are unit sites, those of $S_{k,y}(X,A)$ fermion sites of $T^{\mathrm{f}}_k$.

\emph{Kernels.}
If the source of $O$ is a space of bond functions, then $(OA)(b)=\sum_{b'}O(b,b')\cdot A(b')$, where $b'$ ranges over the positively oriented unit bonds.
For (p2) the kernel is the matrix in the basis $\{\Upsilon\}$.
If the source is $\mathfrak{S}(1,\ell)$, then $(Ov)(t)=\sum_{z'}O(t,z')\cdot v(z')$.
If the source is $\mathfrak{S}^{\mathrm{ch}}(\eta_k,\ell)$, the kernel is that of (I.1.4).
On the covering lattice
\[
O(t+\ell e_{\mu},t')=O(t,t'+\ell e_{\mu})=\boldsymbol{\chi}(e_{\mu})\cdot O(t,t'),
\]
with $\boldsymbol{\chi}=1$ for photon objects and $\boldsymbol{\chi}(e_{\mu})=-1$ for spinors.

\emph{The properties.}
Let $\mathrm{par}'=(\alpha_{0},\gamma,\mathrm{kq};C)$ be a set of parameters.
Here $\mathrm{kq}$ denotes the third entry of a general set of parameters, the rate per unit of $d_M$, as in the paragraph W7 of Section~I.5.8 and in Sections~G1.10 and~G1.11; it equals $\kappa_{\mathrm{w}}+\kappa_{0}$ for $\mathrm{par}^{0}$ and $\kappa_{\mathrm{w}}$ for $\mathrm{par}$.
The letter $\kappa'$ is reserved for the number $\kappa_{\mathrm{w}}+\kappa_{0}$ of the line (C4) of Section~I.8.2 and does not occur in W1--W17.
A family has the property W$n$ for $\mathrm{par}'$ if the assertion W$n$ below holds with the numbers of $\mathrm{par}'$, for all $s\in\mathrm{Pd}_k$, $A\in\tilde{G}_k$, $X\in D_k$ and $y\ge 0$; at the end of each property the families concerned are named.
\begin{description}
\item[W1] One has
\[
H_{k}(1)=H_{k},\qquad C_{k}^{1/2}(1)=C_{k}^{1/2},\qquad H_{k}(1,A)=H_{k}(A),\qquad \Gamma_k(1,A)=\Gamma_k(A),
\]
\[
\sum_{X}M_k(X,A)=M_k(A)\quad(k\ge 1),\qquad \sum_{X}S_{k,y}(X,A)=S_{k,y}(A),
\]
where $H_{k}$ and $C_{k}^{1/2}$ on the right of the first two equations are the photon operators of \cite[Section~5.2.5]{QEDI}, as in Section~I.5.7, with $H_{0}=1$, and the other four right-hand sides are the exact fermion operators of Section~I.3.5 for $k\ge 1$ and of Section~I.3.9 for $k=0$.
Families: all.
\item[W2] The families (p1)--(p3) and $\Gamma_k$ are continuous on $\mathrm{Pd}_k\times\tilde{G}_k$, analytic in $(s,A)$ in the interior, and analytic in $A\in\tilde{G}_k$ for every fixed $s\in\mathrm{Pd}_k$, the boundary included; the maps $M_k(X,\cdot\,)$ and $S_{k,y}(X,\cdot\,)$ are analytic on $\tilde{G}_k$.
Families: all.
\item[W3] Let $b'$ be a unit bond, let $\mathfrak{l}\in\{1,2,3\}$, and let $b=(c,c+\eta_k e_{\mu})$, $\tilde{b}=(\tilde{c},\tilde{c}+\eta_k e_{\mu})$ be fine bonds of the same direction with $0<|c-\tilde{c}|\le 1$.
Then each of the three numbers
\[
\left|H_{k}(s')(b,b')\right|,\qquad \left|\left(\partial_{\mathfrak{l}}H_{k}(s')\right)(b,b')\right|,\qquad \left|\left(\delta_{\alpha}\partial_{\mathfrak{l}}H_{k}(s')\right)\left((b,\tilde{b});b'\right)\right|
\]
is at most $C\cdot\exp\left(-\gamma\,d(b,b')\right)$, where
\begin{align*}
\left(\partial_{\mathfrak{l}}O\right)(b,b')&:=\eta_k^{-1}\cdot\left(O\left(b+\eta_k e_{\mathfrak{l}},b'\right)-O(b,b')\right),\\
\left(\delta_{\alpha}\partial_{\mathfrak{l}}O\right)\left((b,\tilde{b});b'\right)&:=|c-\tilde{c}|^{-\alpha}\cdot\left(\left(\partial_{\mathfrak{l}}O\right)(\tilde{b},b')-\left(\partial_{\mathfrak{l}}O\right)(b,b')\right);
\end{align*}
these are the difference quotients of \cite[Section~5.4.2]{QEDI}, written in the notation of Part~I.
Families: (p1).
\item[W4] For all basis elements $\Upsilon$, $\Upsilon'$,
\[
\left|C_{k}^{1/2}(s')(\Upsilon,\Upsilon')\right|\le C\cdot\exp\left(-\gamma\,d(\Upsilon,\Upsilon')\right).
\]
Families: (p2).
\item[W5] One has
\[
\left|H_{k}(s'',A)\,v\right|_{\infty}\le C\,|v|_{\infty},\qquad \sup_{(x,x')}\left|\left(\delta_{\alpha,A}H_{k}(s'',A)\,v\right)(x,x')\right|_{c}\le C\,|v|_{\infty},
\]
where the supremum extends over the short pairs, that is, the pairs with $d_X(x,x')<1$, the pairs inside a cluster included.
Families: (p3).
\item[W6] One has
\[
\left|\Gamma_k(s'',A)(t,t')\right|_{c}\le C\cdot\exp\left(-\gamma\,d(t,t')\right),\qquad \left|M_k(X,A)(z,z')\right|_{c}\le C\cdot\exp\left(-\gamma\,d(z,z')-\mathrm{kq}\cdot d_M(X)\right).
\]
Families: (p4).
\item[W7] For $k\ge 1$ and all $y\ge 0$,
\[
\left|\mathbf{1}_{\Delta_z}\,S_{k,y}(X,A)\,\mathbf{1}_{\Delta_{z'}}f\right|_{\infty}\le C\cdot\exp\left(-\gamma\,d(z,z')-\mathrm{kq}\cdot d_M(X)\right)\cdot|f|_{\infty},
\]
where $\mathbf{1}_{\Delta_z}$ is the indicator of $\Delta_z\times O_{\mathrm{cub}}$.
For $k=0$,
\[
\left|S_{0,y}(X,A)(x,x')\right|_{c}\le C\cdot(1+y)^{-1}\cdot\exp\left(-\gamma\,d(x,x')-\mathrm{kq}\cdot d_M(X)\right).
\]
Families: (p5).
\item[W8] Let $\mathrm{Op}$ be one of the families concerned, let $\mathscr{Y}$ be a set of parameter cubes of the kind of $\mathrm{Op}$ and $Y$ its union, let $s\in\mathrm{Pd}_k$ satisfy $s_{\mathrm{pc}}=0$ for every cube $\mathrm{pc}$ outside $\mathscr{Y}$, and let $Y_1$ be a link-component of $Y$.
If the row index $t$ is contained in $Y_1$, then $\mathrm{Op}(s,A)(t,t')=0$ unless $t'$ is contained in $Y_1$, and
\[
\mathrm{Op}(s,A)(t,\cdot\,)=\mathrm{Op}(\tilde{s},\tilde{A})(t,\cdot\,)
\]
whenever $\tilde{s}\in\mathrm{Pd}_k$ vanishes outside $\mathscr{Y}$ and coincides with $s$ on the cubes of $Y_1$, and $\tilde{A}\in\tilde{G}_k$ coincides with $A$ on the bonds that meet $Y_1$.
Families: (p1)--(p3), $\Gamma_k$.
\item[W9] The assertion of W8 with rows and columns interchanged.
Families: (p1)--(p3), $\Gamma_k$.
\item[W10] For $\mathscr{Y}$, $Y$, $s$ as in W8, the row $\mathrm{Op}(s,A)(t,\cdot\,)$ vanishes if $t$ is contained in no link-component of $Y$.
Families: (p1)--(p3).
\item[W11] For $\mathscr{Y}$, $Y$, $s$ as in W8, the row at a fermion site $x$ vanishes unless $\mathscr{Y}$ contains all cubes of $\pi_k$ that meet $K_{x}=K_{c}$.
Families: (p3).
\item[W12] The kernels vanish unless both indices lie over $X$, and $\mathrm{Op}(X,A)=\mathrm{Op}(X,\tilde{A})$ whenever $A,\tilde{A}\in\tilde{G}_k$ coincide on the bonds that meet $X$.
Families: $M_k(X,A)$, $S_{k,y}(X,A)$.
\item[W13] Let $\boldsymbol{\lambda}$ be a real function of $c$ on $T_k$ and $\boldsymbol{\lambda}^{\mathrm{u}}$ its restriction to $T^0_k$; for $k=0$, $\boldsymbol{\lambda}^{\mathrm{u}}=\boldsymbol{\lambda}$, acting on $\mathfrak{S}^{\mathrm{ch}}(1,\ell)$ by (I.POS).
Then
\begin{align*}
H_{k}(s'',A-d\boldsymbol{\lambda})&=\exp\left(i e_k\boldsymbol{\lambda}\right)\cdot H_{k}(s'',A)\cdot\exp\left(-i e_k\boldsymbol{\lambda}^{\mathrm{u}}\right),\\
O(\,\cdot\,,A-d\boldsymbol{\lambda})&=\exp\left(i e_k\boldsymbol{\lambda}^{\mathrm{u}}\right)\cdot O(\,\cdot\,,A)\cdot\exp\left(-i e_k\boldsymbol{\lambda}^{\mathrm{u}}\right)\qquad\text{for the families (p4)},\\
S_{k,y}(X,A-d\boldsymbol{\lambda})&=\exp\left(i e_k\boldsymbol{\lambda}\right)\cdot S_{k,y}(X,A)\cdot\exp\left(-i e_k\boldsymbol{\lambda}\right).
\end{align*}
Families: (p3)--(p5).
\item[W14] For $\tilde{g}\in\mathbb{G}_k^{L}$,
\[
\mathrm{Op}(\tilde{g}s,\tilde{g}A)\cdot\tilde{g}=\tilde{g}\cdot\mathrm{Op}(s,A),\qquad \mathrm{Op}(\tilde{g}X,\tilde{g}A)\cdot\tilde{g}=\tilde{g}\cdot\mathrm{Op}(X,A);
\]
for the family (p5) this is required for $\tilde{g}\in\mathbb{G}^{L}_{k,3}$ only.
The actions are those of Section~I.4.3, on fermion sites that of (I.4.2).
Families: all.
\item[W15] One has
\[
\Gamma_k(s'',A)^{T}=\mathsf{C}^{-1}\cdot\Gamma_k(s'',-A)\cdot\mathsf{C},\qquad M_k(X,A)^{T}=\mathsf{C}^{-1}\cdot M_k(X,-A)\cdot\mathsf{C},
\]
the transposes being taken for the pairings of Section~I.1.6, for $k=0$ for $\langle\,\cdot\,,\cdot\,\rangle^{\mathrm{ch}}$.
Families: (p4).
\item[W16] For the families (p3), (p4),
\[
\overline{O\left(\overline{s},\overline{A}\right)}=\varkappa^{-1}\cdot O(s,-A)\cdot\varkappa ;
\]
the families (p1), (p2) are real for real $s$.
Families: (p1)--(p4).
\item[W17] (Level $0$.) The diagonal families have W1, W2, W8--W11, W13, W14, W16, and W3, W5 with the constant $4\exp(2\gamma)\cdot M^{27\alpha_{0}}$.
The families (p2), $\Gamma_0$ and $S_{0,y}$ have W1, W2, W4, W6--W9, W10 (the family (p2)), W12--W16, with W7 in the form for $k=0$.
Families: those of level $0$.
\end{description}
The two sets of parameters that occur are
\[
\mathrm{par}^{0}:=\left(\alpha_{0},\gamma_{\mathrm{w}},\kappa_{\mathrm{w}}+\kappa_{0};C_{w}^{(1)}\right),\qquad \mathrm{par}:=\left(\alpha_{0},\gamma_{\mathrm{w}},\kappa_{\mathrm{w}};C_{w}\right),\qquad C_{w}:=64\,K_{0}\cdot C_{w}^{(1)}
\]
(Section~I.8.2, line (C6)).

\subsubsection*{G4.0.5. Trajectories}

Let (FL) hold, and let $(m_{\mathrm{in}},\varepsilon^{\mathrm{v}}_{\mathrm{in}})$ be a real pair.
Its \emph{trajectory on $R$}, $\left(\varepsilon^{\mathrm{v},t}_k,m_k,E_k\right)$, and its \emph{trajectory with slot}, $\left(\mathscr{E}_k,\mathfrak{P}_k\right)$, are given by the recursion of Section~I.5.7 (paragraph Flows), formed with the step maps of the levels $k<k_{*}(R)$ on $T_k(R)$ and with the primitives of (FL); the mass coupling of level $0$ is $m_{\mathrm{in}}$.
At $k=0$ the symbol $m_k$ denotes $m_{\mathrm{in}}$, and the mass coupling of level $0$ is never written $m_0$: the letter $m_0$ denotes the integer of \cite[Definition~5.98]{QEDI}, which is $0$ here (Section~I.6.1).
The recursion defines the data of level $k+1$ as long as $(m_k,\mathscr{E}_k)$ lies in $\mathrm{Ball}^{x}_{k}(4)$.
The trajectory \emph{exists up to level $k_{\mathrm{up}}$} if this holds for all $k<k_{\mathrm{up}}$.

Proposition~48.6\,(a), (b) shows that the recursion defines the data of level $k+1$ under this condition.
Precisely, let $k<k_{*}(R)$, and let $(m_j,\mathscr{E}_j)$ lie in $\mathrm{Ball}^{x}_{j}(4)$ for $j=0,\dots,k$, these conditions being read in this order.
Then $E_k=\mathscr{E}_k|_{0}$, the pair $(m_k,E_k)$ lies in $\mathrm{Ball}_{k}(4)$, the domain of $E^{*}_k$, and $\varepsilon(E_k)=\varepsilon(\mathscr{E}_k)$, $m(E_k)=m(\mathscr{E}_k)$; so all five entries of level $k+1$ are defined.
The step maps of the levels $k<k_{*}(R)$ exist because the data are (OS) at these levels (Proposition~26.6\,(b)).

The pair is \emph{controlled on $R$} if its trajectory on $R$ exists up to $k_{*}(R)$ and $|m_k|\le w^{\mathrm{m}}_k$ for $0\le k\le k_{*}(R)$.

\subsubsection*{G4.0.6. Photon objects of the flow}

The following photon objects of the flow are those of \cite[Section~5.2.5 and Definition~5.11]{QEDI}, as in Section~I.5.7:
the slice $\mathrm{Sl}$;
the field $A_{0}(F^{\mathrm{u}})$;
the fields $A_{j}$ and $A_{*}:=A_{k_{*}}$;
the elements $\zeta_j\in V_{j}$, which are elements of $V_{j}$ and not fluctuation fields in the sense of \cite[Definition~5.10]{QEDI};
the law $\mathrm{d}\mathsf{m}_{*}$ of $A_{*}$ under $\gamma_N$, which is carried by ${V_{\mathrm{t}}}(T^0_{k_{*}})$, the coclosed bond functions on $T^0_{k_{*}}$ with mean zero;
the function $\chi_{*}(A):=\mathbf{1}\left[\,|dA|_{\infty}\le p_{*}\right]$;
the sets $\mathfrak{G}$ and $\mathfrak{G}^{-}$;
the field $\mathrm{Af}_0(A,\zeta)$, written $A^{\downarrow}_0(A,\zeta)$ in \cite[Definition~5.11]{QEDI};
the sets $\mathrm{Box}_{j}$;
the measure $\mu_I$;
the minimisers $A^{\min}_{j}$;
and the operators $\mathrm{Sq}_{j}$.
In lattice units the curvature is written $F^{\mathrm{u}}$ (Section~I.7.1).

\subsection*{30. Free determinants and twists}

The constants $C_{6}$ and $\kappa_{(0)}$ (line (C0r) of Section~I.8.2) depend on $\rho$ only, and the constants $\gamma_{\Sigma}$ and $C_{\Sigma}$ (line (C6)) depend on $(b_{-},b_{+},\rho)$ only.
The function $\varepsilon_5$ of line (C0r) is formed from $C_{6}$ and $\kappa_{(0)}$.
The four numbers $C_{6}$, $\kappa_{(0)}$, $\gamma_{\Sigma}$, $C_{\Sigma}$ are finite, and $\kappa_{(0)}>0$, $\gamma_{\Sigma}>0$.

\subsubsection*{Proposition 30.1. The free determinant on a cubic torus: bulk form, uniformly in the twist and in the spacing}

Lattice units are used, in the setting of Proposition~46 at $(n,n,n)$.
Let $n\ge1$ be an integer, $q>0$, $\theta\in\mathbb{R}^3$ and $\mathbf{m}\in(0,1]$, and let
\[
\mathfrak{S}^{\mathrm{ch}}(1;n)=\left\{f\colon\mathbb{Z}^3\times O_{\mathrm{cub}}\to\mathbb{C}^4:\ f(c+n\,e_{\mu},g)=-f(c,g)\right\},
\]
a space of dimension $192\,n^3$.
$H_{\theta}$ denotes the constant real one-cochain on $\mathbb{Z}^3$ with $H_{\theta}(c,c+e_{\mu})=\theta_\mu/(q\cdot n)$, and $D_{H_{\theta}}$ the operator (I.2.3) at $\eta_0=1$ with the charge $q$.
Its only phases are the numbers $\exp\left(i\,\theta\cdot(g\,e_1)/n\right)$, so that $q$ drops out.
Let $D^{(p)}$ be the Bloch symbol (I.2.6) on ${V_{\mathrm{v}}}$, a $192\times192$ matrix, and put
\[
\mathfrak{d}_{\mathrm{ma}}(p):=\det D^{(p)}\ \text{ at }\ \boldsymbol{\mu}=\mathbf{m},\qquad p^{0}_{\mathrm{ch}}(\mathbf{m}):=\int_{[-\pi,\pi]^3}\frac{\mathrm{d}^3p}{(2\pi)^3}\,\log\mathfrak{d}_{\mathrm{ma}}(p).
\]
In the index of $\mathfrak{d}_{\mathrm{ma}}$ the letters $\mathrm{ma}$ stand for the mass $\mathbf{m}$;
for another value of the mass, that value is written in the place of $\mathrm{ma}$ (as in Proposition~30.3\,(g2)).
The set of momenta $P_{n,\theta}$ is that of Section~I.2.6.

\begin{named}{Proposition 30.1}
For all integers $n\ge1$, all $\theta\in\mathbb{R}^3$, all $\mathbf{m}$ with $0<\mathbf{m}\le1$ and all $q>0$ the following assertions hold.
\begin{itemize}
\item[(a)] The function $\mathfrak{d}_{\mathrm{ma}}$ is real analytic and $2\pi\mathbb{Z}^3$-periodic on $\mathbb{R}^3$, and $\mathfrak{d}_{\mathrm{ma}}(p)>0$ for real $p$.
Consequently $p^{0}_{\mathrm{ch}}(\mathbf{m})$ is a real number.
\item[(b)] One has
\[
\det\left(D_{H_{\theta}}+\mathbf{m}\right)=\prod_{p\in P_{n,\theta}}\mathfrak{d}_{\mathrm{ma}}(p)>0.
\]
\item[(c)] One has the following bound (8.1) for the deviation of $\log\det\left(D_{H_{\theta}}+\mathbf{m}\right)$ from the bulk term $p^{0}_{\mathrm{ch}}(\mathbf{m})\cdot n^3$:
\begin{equation*}
\left|\log\det\left(D_{H_{\theta}}+\mathbf{m}\right)-p^{0}_{\mathrm{ch}}(\mathbf{m})\cdot n^3\right|\le C_{6}\cdot\log\frac{1}{1-\exp\left(-\kappa_{(0)}\cdot\mathbf{m}\cdot n\right)}=\varepsilon_5(\mathbf{m}\cdot n).
\tag{8.1}
\end{equation*}
\item[(d)] (Poisson summation.)
Let $\mathrm{fper}$ be a $2\pi\mathbb{Z}^3$-periodic function with absolutely convergent Fourier series; that is, with
\[
\mathrm{fperh}(z):=\int_{[-\pi,\pi]^3}\frac{\mathrm{d}^3p}{(2\pi)^3}\,\mathrm{fper}(p)\,\exp(-i\,p\cdot z)\qquad(z\in\mathbb{Z}^3)
\]
one has $\sum_{z\in\mathbb{Z}^3}|\mathrm{fperh}(z)|<\infty$, and $\mathrm{fper}(p)=\sum_{z\in\mathbb{Z}^3}\mathrm{fperh}(z)\,\exp(i\,p\cdot z)$ at every real $p$.
(For a continuous $\mathrm{fper}$ the second property follows from the first.)
Then the two relations of the display (8.2) hold, the Poisson summation formula over $P_{n,\theta}$ and the bound for the deviation of the sum from $n^3\cdot\mathrm{fperh}(0)$:
\begin{equation*}
\begin{aligned}
&\sum_{p\in P_{n,\theta}}\mathrm{fper}(p)=n^3\cdot\sum_{w\in\mathbb{Z}^3}\mathrm{fperh}(n\cdot w)\cdot(-1)^{w_1+w_2+w_3}\cdot\exp(i\,\theta\cdot w),\\
&\left|\sum_{p\in P_{n,\theta}}\mathrm{fper}(p)-n^3\cdot\mathrm{fperh}(0)\right|\le n^3\cdot\sum_{w\in\mathbb{Z}^3\setminus\{0\}}\left|\mathrm{fperh}(n\cdot w)\right|.
\end{aligned}
\tag{8.2}
\end{equation*}
\end{itemize}
The constants are these.
The number $\kappa_{(0)}>0$ is that of Proposition~1.2, and
\[
C_{6}:=26\left[\mathrm{cF13}\cdot C_{\mathrm{lS}}+\frac{16\,C_{\mathrm{lH}}}{(\mathrm{e}\cdot\theta_{0})^2}\right]<\infty,
\]
where $\mathrm{cF13}$ is the number $c(\mathrm{kk},N)$ of \cite[Lemma~12.4]{QEDI} at $(\mathrm{kk},N)=(1,3)$, an absolute number, and
\begin{gather*}
C_{\mathrm{lH}}:=188\left(1+\frac{6\,(1+\rho)}{\hat{g}}\right),\\
C_{\mathrm{lS}}:=\max_{|\mathrm{ak}|\le3}\,96\,K_{\mathrm{Str}}\cdot\sum_{\mathrm{ak}'\le\mathrm{ak}}\binom{\mathrm{ak}}{\mathrm{ak}'}\cdot A'_{\mathrm{ak}'}\cdot\left(\frac{2}{\theta_{0}}\right)^{|\mathrm{ak}|-|\mathrm{ak}'|+1}\cdot\left(3\pi^2+1\right)^{\frac{|\mathrm{ak}|-|\mathrm{ak}'|}{2}}.
\end{gather*}
Here $\hat{g}$, $\theta_{0}$, $K_{\mathrm{Str}}$ and the numbers $A'_{\mathrm{ak}}$ are those of Proposition~1.2;
for a multi-index $\mathrm{ak}=(\mathrm{ak}_1,\mathrm{ak}_2,\mathrm{ak}_3)$ one puts $|\mathrm{ak}|:=\mathrm{ak}_1+\mathrm{ak}_2+\mathrm{ak}_3$;
the relation $\mathrm{ak}'\le\mathrm{ak}$ is meant componentwise, and $\binom{\mathrm{ak}}{\mathrm{ak}'}:=\prod_{\mu}\binom{\mathrm{ak}_\mu}{\mathrm{ak}'_\mu}$.
The constants $C_{6}$ and $\kappa_{(0)}$ depend on $\rho$ only; they do not depend on $n$, $\theta$, $\mathbf{m}$, $q$.
The bound (8.1) depends on $(n,\mathbf{m})$ through the product $\mathbf{m}\cdot n$ only.
\end{named}

\begin{named}{Remark}
Proposition~30.1 corresponds to \cite[Lemma~19.18]{QEDI}, together with the definitions of \cite[Section~19.8.1]{QEDI}, and to the first assertion of \cite[Theorem~5.85(d)]{QEDI}.
It differs from them as follows.
The factor $w_{\mathbf{m}}(p)^2$ of the product formula there is replaced by $\mathfrak{d}_{\mathrm{ma}}(p)$, and the number $p^0(\mathbf{m})=2\int_{[-\pi,\pi]^3}\frac{\mathrm{d}^3p}{(2\pi)^3}\,\log w_{\mathbf{m}}(p)$ by $p^{0}_{\mathrm{ch}}(\mathbf{m})=\int_{[-\pi,\pi]^3}\frac{\mathrm{d}^3p}{(2\pi)^3}\,\log\mathfrak{d}_{\mathrm{ma}}(p)$.
The function $x\mapsto-312\log\left(1-\exp\left(-\frac{x}{4}\right)\right)$ is replaced by the function $\varepsilon_5(x)=C_{6}\cdot\log\frac{1}{1-\exp\left(-\kappa_{(0)}\cdot x\right)}$ of $x>0$, whose constants are not numerical.
The factorisation of $w_{\mathbf{m}}$ through the energy $\mathrm{En}$ and the lower bound for $\mathrm{En}$, used in the proof of \cite[Lemma~19.18]{QEDI}, have no analogue here: $\mathfrak{d}_{\mathrm{ma}}$ is a polynomial of degree up to $32$ in $\exp(i\,p_\mu)$ and in $\exp(-i\,p_\mu)$, and no integration in closed form is available.
The space of spinors has dimension $192\,n^3$.
For $n\le3$ the cochain $H_{\theta}$ is read on $\mathbb{Z}^3$, in the setting of Proposition~46, as \cite[Lemma~19.18]{QEDI} does for $n\le2$.
The directions are numbered $1,2,3$, and there is no momentum coordinate $p_0$.
\end{named}

\subsubsection*{Proposition 30.2. Dependence on the twist; the free partition number}

\begin{named}{Proposition 30.2}
\begin{itemize}
\item[(a)] In the setting of Proposition~30.1, for all $\theta\in\mathbb{R}^3$,
\[
\left|\log\frac{\det\left(D_{H_{\theta}}+\mathbf{m}\right)}{\det\left(D_0+\mathbf{m}\right)}\right|\le2\,\varepsilon_5(\mathbf{m}\cdot n).
\]
\item[(b)] Let $\bar{m}\in(0,1]$, let $N\ge0$ be an integer, $\eta=L^{-N}$, and let $R>0$ be such that $n:=R/\eta$ is an integer $\ge1$.
Let
\[
\mathrm{zr}^{\mathrm{free}}_n:=\int_{[-\pi,\pi]^3}\frac{\mathrm{d}\theta}{(2\pi)^3}\,\det\left(D_{H_{\theta}}+\bar{m}\cdot\eta\right)
\]
(Section~I.7.3).
Then $\mathrm{zr}^{\mathrm{free}}_n>0$, and
\[
\left|\log\mathrm{zr}^{\mathrm{free}}_n-p^{0}_{\mathrm{ch}}(\bar{m}\cdot\eta)\cdot n^3\right|\le\varepsilon_5(\bar{m}\cdot R).
\]
\end{itemize}
\end{named}

\begin{named}{Remark}
Proposition~30.2 corresponds to \cite[Corollary~19.19]{QEDI} and to the second assertion of \cite[Theorem~5.85(d)]{QEDI}.
The number $\mathrm{zr}^{\mathrm{free}}_n$ corresponds to the number $\rho^{\mathrm{free}}_n$ there;
both are means over $\theta$ of $\det\left(D_{H_{\theta}}+\bar{m}\cdot\eta\right)$, formed here with the operator (I.2.3) and there with the Wilson operator.
It differs from them through the number $p^{0}_{\mathrm{ch}}$ and the function $\varepsilon_5$ of Proposition~30.1.
\end{named}

\subsubsection*{Proposition 30.3. The blocked operator at a constant connection: inverse formula, positivity, twist dependence}

Let $\ell\ge1$ and $\mathfrak{N}\ge3$ be integers, $\mathfrak{N}$ odd, and put $\mathrm{et}:=1/\mathfrak{N}$; let $\mu\in(0,1]$, $\mathsf{b}>0$ and $\theta\in\mathbb{R}^3$.
Define
\[
\mathscr{H}_{\theta}:=\left\{f\colon\mathrm{et}\,\mathbb{Z}^3\times O_{\mathrm{cub}}\to\mathbb{C}^4:\ f(c+\ell\,e_{\mu},g)=-\exp(i\,\theta_\mu)\,f(c,g)\right\},
\]
with the scalar product $(f',f):=\mathrm{et}^3\,\frac{1}{48}\sum\overline{f'(x)}\cdot f(x)$, summed over a period cell; its dimension is $192\,(\ell\cdot\mathfrak{N})^3$.
Define further
\[
\mathscr{H}^{\mathrm{u}}_{\theta}:=\left\{v\colon\mathbb{Z}^3\to\mathbb{C}^4:\ v(y+\ell\,e_{\mu})=-\exp(i\,\theta_\mu)\,v(y)\right\},
\]
with the scalar product $(v',v):=\sum\overline{v'(y)}\cdot v(y)$ over a period cell; its dimension is $4\,\ell^3$.
For a vector, $|\cdot|$ denotes the norm defined by these scalar products; the norm $\|O\|$ of an operator $O$ is the operator norm formed with them, and $\operatorname{Re}O:=(O+O^{*})/2$.
For $y\in\mathbb{Z}^3$ let
\[
\Delta_y:=\left\{c\in\mathrm{et}\,\mathbb{Z}^3:\ |c-y|<1/2\right\};
\]
it consists of $\mathfrak{N}^3$ sites, and these blocks tile $\mathrm{et}\,\mathbb{Z}^3$ because $\mathfrak{N}$ is odd.
For $c\in\mathrm{et}\,\mathbb{Z}^3$, $[c]$ denotes the point $y$ with $c\in\Delta_y$.
For $x=(c,g)$ let
\[
(Df)(x):=-\mathrm{et}^{-1}\sum_{i=0}^{3}M_{i}(x)\left[f\left(x^{(i)}\right)-f(x)\right],\qquad x^{(i)}=(c,g\,s_i)\ \ (i\le2),\qquad x^{(3)}=\left(c+\mathrm{et}\cdot g\,e_1,\ g\,s_3\right);
\]
this is the operator (I.2.3) at the mesh $\mathrm{et}$ and $A=0$, and it maps $\mathscr{H}_{\theta}$ into itself.
Let
\[
(Qf)(y):=\mathrm{et}^3\,\frac{1}{48}\sum_{c\in\Delta_y}\sum_{g\in O_{\mathrm{cub}}}f(c,g)\quad\left(\mathscr{H}_{\theta}\to\mathscr{H}^{\mathrm{u}}_{\theta}\right),\qquad(Q^{*}v)(c,g):=v([c])\quad\left(\mathscr{H}^{\mathrm{u}}_{\theta}\to\mathscr{H}_{\theta}\right),
\]
and
\[
T_{\theta}:=D+\mu+\mathsf{b}\cdot Q^{*}Q\ \text{ on }\mathscr{H}_{\theta},\qquad\mathrm{Dfl}_{\theta}:=\mathsf{b}-\mathsf{b}^2\cdot Q\,T_{\theta}^{-1}Q^{*}\ \text{ on }\mathscr{H}^{\mathrm{u}}_{\theta},\qquad\mathrm{Rtw}_{\theta}:=\frac{\det\mathrm{Dfl}_{\theta}}{\det\mathrm{Dfl}_{\theta=0}}.
\]
An index $\theta=0$ indicates the value of the object at $\theta=0$.

\begin{named}{Proposition 30.3}
For all data as above the following assertions hold.
\begin{itemize}
\item[(a)] (Structure.)
$Q^{*}$ is the adjoint of $Q$; $Q\,Q^{*}=1$; and $Q^{*}Q$ is the orthogonal projection onto the functions that are constant on each set $\Delta_y\times O_{\mathrm{cub}}$.
One has $\operatorname{Re}(D+\mu)\ge\mu$ and $\operatorname{Re}T_{\theta}\ge\mu$; hence $D+\mu$ and $T_{\theta}$ are bijective, and their inverses have norm at most $1/\mu$.
In the two eigenspaces of the matrix $\Sigma_3$ of Section~I.2.1 (see (I.O3); for an operator $O$ that commutes with $\Sigma_3$, $O^{+}$ denotes its part in the first of them, the eigenspace of the eigenvalue $1$) one has
\[
D+\mu=(D+\mu)^{+}\oplus\left((D+\mu)^{+}\right)^{*},\qquad T_{\theta}=T^{+}\oplus\left(T^{+}\right)^{*}.
\]
Hence
\[
\det\left(D+\mu\ \text{ on }\mathscr{H}_{\theta}\right)=\left|\det(D+\mu)^{+}\right|^2>0,\qquad\det T_{\theta}=\left|\det T^{+}\right|^2>0.
\]
\item[(b)] (Inverse formula.)
$\mathrm{Dfl}_{\theta}$ is invertible, and
\[
\mathrm{Dfl}_{\theta}^{-1}=\frac{1}{\mathsf{b}}+Q\,(D+\mu)^{-1}Q^{*},\qquad\operatorname{Re}\mathrm{Dfl}_{\theta}^{-1}\ge\frac{1}{\mathsf{b}},\qquad\left\|\mathrm{Dfl}_{\theta}^{-1}\right\|\le\frac{1}{\mathsf{b}}+\frac{1}{\mu}.
\]
\item[(c)] (Positivity.)
One has
\[
\det\mathrm{Dfl}_{\theta}=\mathsf{b}^{4\ell^3}\cdot\frac{\det\left(D+\mu\ \text{ on }\mathscr{H}_{\theta}\right)}{\det T_{\theta}}>0.
\]
In particular $\mathrm{Rtw}_{\theta}>0$.
\item[(d)] One has
\[
\mathrm{Rtw}_{\theta}=\frac{\det\left(D_{H_{\theta}}+\mathrm{et}\cdot\mu\right)}{\det\left(D_0+\mathrm{et}\cdot\mu\right)}\cdot\frac{\det T_{\theta=0}}{\det T_{\theta}},
\]
where the first ratio is that of Proposition~30.1 on the torus with $n=\ell\cdot\mathfrak{N}$ sites per side, at the mass $\mathbf{m}=\mathrm{et}\cdot\mu$ (so that $\mathbf{m}\cdot n=\mu\cdot\ell$).
\item[(e)] (Twist dependence, uniformly in the mass $\mu$ and in the mesh.)
Let $0<b_{-}\le\mathsf{b}\le b_{+}$.
There are constants $\gamma_{\Sigma}\in(0,1]$ and $C_{\Sigma}<\infty$, depending on $(b_{-},b_{+},\rho)$ only, such that for all integers $\ell\ge1$, all odd integers $\mathfrak{N}\ge3$, all $\mu\in(0,1]$ and all $\theta\in\mathbb{R}^3$ the two twist bounds of the display (8.5) hold, the first for the ratio $\det T_{\theta}/\det T_{\theta=0}$ and the second for $\mathrm{Rtw}_{\theta}$:
\begin{equation*}
\begin{aligned}
&\left|\log\frac{\det T_{\theta}}{\det T_{\theta=0}}\right|\le C_{\Sigma}\cdot\ell^{-1}\exp\left(-\gamma_{\Sigma}\cdot\ell\right),\\
&\left|\log\mathrm{Rtw}_{\theta}\right|\le2\,\varepsilon_5(\mu\cdot\ell)+C_{\Sigma}\cdot\ell^{-1}\exp\left(-\gamma_{\Sigma}\cdot\ell\right).
\end{aligned}
\tag{8.5}
\end{equation*}
\item[(f)] (The picture with a connection.)
Let $\mathrm{ec}>0$, let $c_{\theta}$ be the constant real bond function on $\mathrm{et}\,\mathbb{Z}^3$ with the components $\theta_\mu/(\mathrm{ec}\cdot\ell)$, let $\mathfrak{S}^{\mathrm{ch}}(\mathrm{et},\ell)$ and $\mathfrak{S}(1,\ell)$ be the antiperiodic spaces of Section~I.1.5 with $\mathrm{et}$ in the place of $\eta_k$, and put
\[
(\mathrm{Uth}\,f)(c,g):=\exp\left(\frac{i\,\theta\cdot c}{\ell}\right)f(c,g),\qquad\left(\mathrm{Uth}^{\mathrm{u}}v\right)(y):=\exp\left(\frac{i\,\theta\cdot y}{\ell}\right)v(y).
\]
Then $\mathrm{Uth}$ maps $\mathfrak{S}^{\mathrm{ch}}(\mathrm{et},\ell)$ unitarily onto $\mathscr{H}_{\theta}$, $\mathrm{Uth}^{\mathrm{u}}$ maps $\mathfrak{S}(1,\ell)$ onto $\mathscr{H}^{\mathrm{u}}_{\theta}$, and $D_{c_{\theta}}=\mathrm{Uth}^{-1}D\,\mathrm{Uth}$ for the operator (I.2.3) with $(\eta_k,e_k)=(\mathrm{et},\mathrm{ec})$ (see (I.6.1)).
If $\mathfrak{N}=L^k$ with $k\ge1$, $\ell\in\mathrm{Cl}$, $\mathrm{ec}=e_k$ and $\mathsf{b}=b_{k}$, then for the operators of Sections~I.3.3 and I.3.5 at the mass $\mu$
\begin{gather*}
Q_k(c_{\theta})\,J^{*}=\left(\mathrm{Uth}^{\mathrm{u}}\right)^{-1}Q\,\mathrm{Uth},\qquad J\,Q_k^{T}(-c_{\theta})=\mathrm{Uth}^{-1}Q^{*}\,\mathrm{Uth}^{\mathrm{u}},\qquad S_k(c_{\theta})=\mathrm{Uth}^{-1}T_{\theta}^{-1}\,\mathrm{Uth},\\
D_k(c_{\theta})=\left(\mathrm{Uth}^{\mathrm{u}}\right)^{-1}\mathrm{Dfl}_{\theta}\,\mathrm{Uth}^{\mathrm{u}},\qquad\frac{\det D_k(c_{\theta})}{\det D_k(0)}=\mathrm{Rtw}_{\theta}.
\end{gather*}
(For points $y$, $c$ of $\mathrm{et}\,\mathbb{Z}^3$ with $c\in\Delta_y$, and for the rectilinear paths of the covering lattice that make up $\tau_k$, one has $\mathrm{et}\cdot(\tau_k c_{\theta})(y,c)=c_{\theta}\cdot(c-y)$, where $c_{\theta}\cdot v:=\sum_\mu\theta_\mu v_\mu/(\mathrm{ec}\cdot\ell)$; hence $\varphi_{c_{\theta}}(y,c)=\exp\left(i\,\theta\cdot(c-y)/\ell\right)$.)

In the flow:
under (S-4) of Section~I.1.8, that is, for $\bar{m}\in(0,1]$, $e>0$ with $e\cdot h^{\frac{1}{2}}\le e_1$, $R\in\mathrm{Tor}(e)$ and an integer $N\ge N_{\min}(R,e)$ (this is (FL) without the primitives), and with $k_{*}$, $\ell_*$, $\mu_*$, $e_{*}$ as in Sections~I.6.1 and I.6.2,
one has $k_{*}\ge1$ (so that $L^{k_{*}}$ is an odd integer $\ge3$, and $b_{k_{*}}$ is defined), $\ell_*\in\mathrm{Cl}$, $0<\mu_*\le1$, $b_{-}\le b_{k_{*}}\le b_{+}$ (Section~I.3.4) and $e_{*}>0$,
so that $(\ell,\mathfrak{N},\mu,\mathsf{b},\mathrm{ec})=(\ell_*,L^{k_{*}},\mu_*,b_{k_{*}},e_{*})$ are data of this proposition and of (e) and, with $k=k_{*}$, satisfy the hypotheses $\mathfrak{N}=L^k$, $k\ge1$, $\ell\in\mathrm{Cl}$, $\mathrm{ec}=e_k$, $\mathsf{b}=b_{k}$ of the identification above;
for these data $c_{\theta}=c^{(k_{*})}_{\theta}$ (Section~I.6.6), $D_k(\cdot)=D_{*}(\cdot)$ (Section~I.6.2), $\mathrm{et}\cdot\mu=\bar{m}\cdot\eta$ and $\ell\cdot\mathfrak{N}=R/\eta$.
\item[(g)] (Symbols.)
Let $\mathrm{Mf}:=\left\{-(\mathfrak{N}-1)/2,\dots,(\mathfrak{N}-1)/2\right\}^3$.
For $p\in\mathbb{R}^3$ and $\mathrm{jf}\in\mathrm{Mf}$ put $p_{\mathrm{jf}}:=p+2\pi\,\mathrm{jf}$, and let
\[
u_1(z):=\mathfrak{N}^{-1}\sum_{j\in\mathbb{Z},\ |j|\le(\mathfrak{N}-1)/2}\exp\left(-\frac{i\,z\,j}{\mathfrak{N}}\right),\qquad\mathrm{uh}(p):=\prod_{\mu}u_1(p_\mu);
\]
$\mathrm{uh}(p)$ is real for real $p$.
Let $D^{(p';\mathrm{ma})}$ denote the symbol (I.2.6) at the momentum $p'$ and at $\boldsymbol{\mu}=\mathbf{m}$;
as in the index of $\mathfrak{d}_{\mathrm{ma}}$, the letters $\mathrm{ma}$ stand for the mass $\mathbf{m}$, and a value of the mass is written in their place.
The space ${V_{\mathrm{v}}}$ of Section~I.2.6 carries the scalar product $(u',u):=\frac{1}{48}\sum_{g\in O_{\mathrm{cub}}}\overline{u'(g)}\cdot u(g)$, the space $\mathbb{C}^{\mathrm{Mf}}\otimes{V_{\mathrm{v}}}=\bigoplus_{\mathrm{jf}\in\mathrm{Mf}}{V_{\mathrm{v}}}$ the scalar product $\sum_{\mathrm{jf}\in\mathrm{Mf}}(u'_{\mathrm{jf}},u_{\mathrm{jf}})$, and $\mathbb{C}^4$ the standard one, $\overline{v'}\cdot v$.
On ${V_{\mathrm{v}}}$ let $(Jv)(g):=v$ for $v\in\mathbb{C}^4$, $J^{*}u:=\frac{1}{48}\sum_{g\in O_{\mathrm{cub}}}u(g)$, and $\mathrm{Plt}:=J\,J^{*}$.
Let
\begin{gather*}
T(p):=\bigoplus_{\mathrm{jf}\in\mathrm{Mf}}\mathrm{et}^{-1}D^{(\mathrm{et}\,p_{\mathrm{jf}};\,\mathrm{et}\,\mu)}+\mathsf{b}\cdot\left(\mathrm{uh}(p_{\mathrm{jf}})\cdot\mathrm{uh}(p_{\mathrm{jf}'})\right)_{\mathrm{jf},\mathrm{jf}'}\otimes\mathrm{Plt}\qquad\text{on }\mathbb{C}^{\mathrm{Mf}}\otimes{V_{\mathrm{v}}},\\
\mathrm{Ch}(p):=\frac{1}{\mathsf{b}}+\sum_{\mathrm{jf}\in\mathrm{Mf}}\mathrm{uh}(p_{\mathrm{jf}})^2\cdot J^{*}\left(\mathrm{et}^{-1}D^{(\mathrm{et}\,p_{\mathrm{jf}};\,\mathrm{et}\,\mu)}\right)^{-1}J\qquad(4\times4).
\end{gather*}
\begin{itemize}
\item[(g1)] $T_{\theta}$ is unitarily equivalent to the direct sum of the operators $T(p)$, and $\mathrm{Dfl}_{\theta}^{-1}$ to the direct sum of the matrices $\mathrm{Ch}(p)$, over $p\in P_{\ell,\theta}$.
\item[(g2)] For real $p$
\[
\det T(p)=\mathsf{b}^4\cdot\det\mathrm{Ch}(p)\cdot\prod_{\mathrm{jf}\in\mathrm{Mf}}\mathrm{et}^{-192}\cdot\mathfrak{d}_{\mathrm{et}\,\mu}(\mathrm{et}\,p_{\mathrm{jf}})>0;
\]
moreover $\operatorname{Re}\mathrm{Ch}(p)\ge1/\mathsf{b}$, and $\mathrm{Ch}(p)=\alpha-i\,\gamma\cdot\beta$ with a real number $\alpha\ge1/\mathsf{b}$ and a vector $\beta\in\mathbb{R}^3$; finally $\det\mathrm{Ch}(p)=\left(\alpha^2+|\beta|_2^2\right)^2$.
\item[(g3)] The function $\mathrm{Flg}(p):=\log\det T(p)$ is real analytic, $2\pi\mathbb{Z}^3$-periodic, invariant under $p\mapsto r\,p$ for $r\in O_{\mathrm{cub}}$, and independent of $(\ell,\theta)$.
With the Fourier coefficients $\mathrm{Flgh}(z)$, $z\in\mathbb{Z}^3$, of $\mathrm{Flg}$ (Section~I.2.6),
\[
\log\frac{\det T_{\theta}}{\det T_{\theta=0}}=\ell^3\cdot\sum_{w\in\mathbb{Z}^3\setminus\{0\}}\mathrm{Flgh}(\ell\cdot w)\cdot(-1)^{w_1+w_2+w_3}\cdot\left(\exp(i\,\theta\cdot w)-1\right).
\]
\end{itemize}
\end{itemize}
The constants, in the order in which they are formed, are the following; $\hat{g}$ and $c_{\mathrm{N3}}$ are those of line (C0r) of Section~I.8.2.
\begin{gather*}
K_{\mathrm{iv}}:=b_{-}^{-\frac{1}{2}}+\frac{8}{\pi^3}\,(3\,\hat{g})^{-\frac{1}{2}}+\left(1+\frac{8}{\pi^3}\right)\frac{2+b_{+}^{\frac{1}{2}}}{\pi\,c_{\mathrm{N3}}},\qquad C_{\mathrm{iv}}:=\max\left\{1,\left(\frac{\pi^3K_{\mathrm{iv}}}{8}\right)^2\right\},\\
\gamma_{\Sigma}:=\min\left\{1,\frac{1}{2\,\mathrm{e}\cdot(6+b_{+})\cdot C_{\mathrm{iv}}}\right\},\qquad C_{\mathrm{sv}}:=2\left(1+(1+b_{+})\,C_{\mathrm{iv}}\right),\qquad C_{\mathrm{uv}}:=\left(1+3\pi^2c_{\mathrm{N3}}^2\right)^{\frac{1}{2}},\\
S_{\mathrm{jf}}:=1+\frac{4\pi^2+\pi^4/45}{(\pi\,c_{\mathrm{N3}})^4},\\
C_{\mathrm{Flg}}:=26\left[9^4\cdot192\cdot C_{\mathrm{sv}}^4\cdot C_{\mathrm{uv}}^3\cdot S_{\mathrm{jf}}+300\,\mathrm{e}\cdot b_{+}\cdot(2\,C_{\mathrm{iv}})^4\cdot\max(9,\mathrm{e}\cdot b_{+})^3\right],\qquad C_{\Sigma}:=2\left(4\pi^2+\frac{\pi^4}{45}\right)C_{\mathrm{Flg}};
\end{gather*}
here $4\pi^2+\pi^4/45=\sum_{w\in\mathbb{Z}^3\setminus\{0\}}|w|_\infty^{-4}$.
The numbers $\gamma_{\Sigma}$ and $C_{\Sigma}$ are explicit functions of $(b_{-},b_{+},\rho,\lambda_{\mathrm{int}})$.
In the scheme $b_{-}=b\,(1-1/L)$ and $b_{+}=b$ (Section~I.3.4).
The constants do not depend on $\ell$, $\mathfrak{N}$, $\mu$, $\mathsf{b}$, $\theta$, $\mathrm{ec}$.
\end{named}

\begin{named}{Remark}
Proposition~30.3 corresponds to the definitions, the statement on invertibility and the formulas in momentum space of \cite[Section~19.8.2]{QEDI}, to \cite[Lemma~19.20]{QEDI}, to \cite[Lemma~19.21]{QEDI}, and to the last assertion of \cite[Theorem~5.85(d)]{QEDI}.
The operators $D(c)$, $T_c$ and the number $R_\theta$ there correspond to $\mathrm{Dfl}_{\theta}$, $T_{\theta}$ and $\mathrm{Rtw}_{\theta}$ (by (f)).
The differences are the following.
The fine spinors carry the label $g\in O_{\mathrm{cub}}$, and $Q$ contains the mean $\frac{1}{48}\sum_{g}$ over the labels, that is, $Q$ corresponds to $Q_k(\cdot)\,J^{*}$.
The operators are defined on spaces with the twisted boundary condition and without connection; the description with a constant connection is clause (f).
The fine space has dimension $192\,(\ell\cdot\mathfrak{N})^3$, and the factor $\mathrm{et}^{-192\,n^3}$ cancels in (d); the factors $\mathsf{b}^{4\ell^3}$ and $\mathsf{b}^4$, and the rank $4$ of the perturbation of the block-diagonal part of $T(p)$, are as there.
The explicit symbols $\hat{D}(q)$, $a(q)$, $s_\nu(q)$, $w(q)$ and the matrices $R_m$ of \cite[Section~19.8.2]{QEDI} are replaced by the symbol $\mathrm{et}^{-1}D^{(\mathrm{et}\,p_{\mathrm{jf}};\,\mathrm{et}\,\mu)}$ and by the compression $J^{*}\left(\mathrm{et}^{-1}D^{(\mathrm{et}\,p_{\mathrm{jf}};\,\mathrm{et}\,\mu)}\right)^{-1}J$ of its inverse to the range of $J$.
The constants $\gamma_1\in(0,1/16]$ and $C_{\Sigma}=86\,C_4$ of \cite[Lemma~19.21]{QEDI}, functions of $(b_{-},b_{+})$, are replaced by $\gamma_{\Sigma}\in(0,1]$ and $C_{\Sigma}$, functions of $(b_{-},b_{+},\rho)$; the factorisation of the scalar function $w(q_m)$ as a power of the mesh times $\acute{M}\cdot\exp(\mathrm{En})\left(1-z_{+}\right)\left(1-z_{-}\right)$, on which the proof of \cite[Lemma~19.21]{QEDI} rests, has no analogue here.
The identification with the operators of the renormalisation scheme in (f) is made only for $\ell\in\mathrm{Cl}$ and $\mathfrak{N}=L^k$, and the identity for the phases is stated on the covering lattice, for the rectilinear paths that make up $\tau_k$.
\end{named}

\subsection*{14. Assembly of the bounds on the fermion operators; the complete ordered list of constants and conditions}

\subsubsection*{Proposition 14.1 (Z2$^{\mathrm{ch}}$)}

\begin{named}{Proposition 14.1}[Z2$^{\mathrm{ch}}$]
\emph{Constants.}
Put
\[
C^{F}:=2\max\left\{C_{1},\ C_{\mathrm{By}},\ 304\,\mathrm{e}\cdot\max(1,b)\cdot a',\ C_{\Gamma},\ \mathrm{e}\cdot b^{2}\cdot(32a')^{2}\right\},
\]
where $C_{1}$ is the constant of Proposition~7, $C_{\mathrm{By}}$ and $C_{\Gamma}$ are the constants of Proposition~5, and $a'$ is the constant of Proposition~2.3.
Further, $\alpha^{F}:=\frac{1}{54}$ and $L^{F}:=5$;
$\gamma^{F}$ is the number of Proposition~2.3, and $0<\gamma^{F}\le\frac{1}{6L}$;
$M^{F}$ is the number of the entry $M^{F}$ of Proposition~14.2\,(a);
$C_{T}$ is the constant of Proposition~9.

Let data (SF) of Section~G1.1 be given, that is, an integer $k\ge0$, a side $\ell\in\mathrm{Cl}$, a mass $\mu\in(0,1]$ and a number $e_k\in(0,e_0]$.
All norms are those of Section~I.1.5, and all kernels are taken in the sense of (I.1.4).
Then the following assertions hold.
\begin{itemize}
\item[(i)] The properties FE1$^{\mathrm{ch}}$--FE6$^{\mathrm{ch}}$ and FE8$^{\mathrm{ch}}$ hold in the following form.
\begin{description}
\item[FE1$^{\mathrm{ch}}$ ($k\ge1$).]
For $A\in\tilde{G}_k$ the operator $D_A+\mu+b_k\,J P_k(A)J^{*}$ is bijective on $\mathfrak{S}^{\mathrm{ch}}(\eta_k,\ell)$, and its inverse $S_k(A)$ is an analytic function of $A$ on $\tilde{G}_k$.
For all unit sites $z,z'$ and all $f\in\mathfrak{S}^{\mathrm{ch}}(\eta_k,\ell)$ one has
\begin{gather*}
\left|\mathbf{1}_{\Delta_z}S_k(A)\mathbf{1}_{\Delta_{z'}}f\right|_{\infty}\le C^{F}\exp\left(-\gamma^{F}d(z,z')\right)|f|_{\infty},\\
\left|\mathbf{1}_{\Delta_z}\delta_{\alpha,A}S_k(A)\mathbf{1}_{\Delta_{z'}}f\right|_{\infty}\le C^{F}\exp\left(-\gamma^{F}d(z,z')\right)|f|_{\infty},
\end{gather*}
where in the second bound the H\"older quotient is taken at the short pairs, the pairs inside a cluster included.
This is Proposition~2.8 at $y=\infty$.
\item[FE2$^{\mathrm{ch}}$ ($k\ge1$).]
The same assertions hold for $F_y(A)$ and $S_{k,y}(A)$, for every real $y$, with constants that do not depend on $y$;
in particular they hold for $S^{0}_{k+1}(A)=S_{k,0}(A)$.
This is Proposition~2.8.
\item[FE3$^{\mathrm{ch}}$.]
For $k\ge1$: the assertions (a)--(d) of Proposition~4 and (e) of Proposition~5.
For $k=0$ and $A\in\tilde{G}_0$: the operator $F_y(A)$ is bijective for every real $y$,
\[
\left|S_{0,y}(A)(x,x')\right|_{c}\le C^{F}\,(1+|y|)^{-1}\exp\left(-\gamma^{F}d(x,x')\right),
\]
in particular $\left|\Gamma_0(A)(x,x')\right|_{c}\le C^{F}\exp\left(-\gamma^{F}d(x,x')\right)$, and $\delta Z_0(A)\neq0$.
This is Proposition~2.8\,(a), (d).
\item[FE4$^{\mathrm{ch}}$:] the assertions (e), (f) of Proposition~6.
\item[FE5$^{\mathrm{ch}}$:] the assertions (a), (b), (d), (e) of Proposition~9, with the constant $C_T$.
\item[FE6$^{\mathrm{ch}}$:] the assertions (a)--(g) of Proposition~10.
\item[FE8$^{\mathrm{ch}}$:] the assertion (b) of Proposition~11.
\end{description}
\item[(ii)] The families of Proposition~7 (for $k\ge1$) and of Proposition~8 (for $k=0$) are the fermionic part (p3), (p4), (p5) of a system of primitives of level $k$ on $T_k$ (Section~G1.11).
They have the properties W1, W2, W5--W17 of Section~G1.11 for the parameters $(\alpha_0,\gamma^{F},\kappa_{\mathrm{w}}+\kappa_0;C^{F})$, because $2C_1\le C^{F}$.
Since
\[
\gamma_{\mathrm{w}}\le\gamma^{F},\qquad C^{F}\le C_{w}^{(0)}\le C_{w}^{(1)}\le C_{w},\qquad M^{8\alpha_0}\le4\exp(2\gamma_{\mathrm{w}})\cdot M^{27\alpha_0}\le C_{w}^{(1)},
\]
they have these properties also for $\mathrm{par}^{0}$ and for $\mathrm{par}$.
The property W14 holds for the whole group $\mathbb{G}_k$ (for (p5): for $\mathbb{G}_{k,3}$), hence for the translations by $\ell_1\mathbb{Z}^3$, for every $\ell_1\in\mathrm{Cl}$ that divides $\ell$.
\item[(iii)] FE7$^{\mathrm{ch}}$: the assertions (d0)--(d6) of Proposition~13 hold for every family (p5) that has, for $\mathrm{par}$, the properties required there;
in particular they hold for the family of (ii).
\item[(iv)] None of the constants depends on $k$, $\ell$, $\mu$, $e_k$, $A$, $y$, $s$, $X$.
The sets $\tilde{G}_k$ and $\tilde{R}_k$ are invariant under the addition of real constant connections (twists), and $\tilde{R}_k+\mathfrak{X}_k\subset\tilde{G}_k$.
Consequently all assertions hold at every connection $\mathscr{T}$ that is the sum of a connection $\mathscr{A}\in\tilde{R}_k$, a connection $\mathscr{C}\in\mathfrak{X}_k$ and a real constant connection, the constant being counted with $\mathscr{A}$.
\item[(v)] (Range of the mass.)
All assertions of (i)--(iii) hold for every $\mu\in(0,1]$, with constants that do not depend on $\mu$, and no relation between $\mu$ and $(k,\ell,e_k)$ is required.
The operator $F_y(A)$ is bijective also at $\mu=0$, for every $y\in Y_k$ and every $A\in\tilde{G}_k$, and at $\mu=0$ the bounds of FE1$^{\mathrm{ch}}$, FE2$^{\mathrm{ch}}$ and those of FE3$^{\mathrm{ch}}$ at $k=0$ hold with the same $C^{F}$ and $\gamma^{F}$.
Hence the upper end of the range of the mass is $1$: for the upper end $\mathfrak{m}_0$ of the mass range, the integer $m_0$ and the number $h'$ of \cite[Definition~5.98]{QEDI} one may take
\[
\mathfrak{m}_0=1,\qquad m_0=0,\qquad h'=h,
\]
and these are the values used wherever \cite{QEDI} uses these three numbers (Sections~I.6.1 and I.8.4).
\end{itemize}
\end{named}

\begin{named}{Remark}[The projected families]
Let $(\bar{m},e,R,N)$ be (FL) of Sections~I.1.8 and G4.0.1, which by Proposition~26.6\,(e) is the same condition as (S-4), and let $R<R_{\mathrm{top}}$.
For a level $k<k_{*}(R)$, the projected family (p5) is the family (p5) of the projection to $T_k(R)$, with $\mathrm{Jr}=R_{\mathrm{top}}/R$, of the system of Proposition~26.5 on $T_k(R_{\mathrm{top}})$, the projection being that of Proposition~26.1;
it is the family (p5) of the primitives of (FL) on $T_k(R)$.
At every level $k<k_{*}(R)$ the assertions of (iii) hold for the projected family (p5), with the data $(k,\ell_k(R),\mu_k,e_k)$.
Indeed, by Proposition~26.6\,(b) these data, together with the primitives of (FL), are (OS).
Data (OS) are data (SF+), since $\mu_k\le\frac{1}{L}\le1$ and $e_k<L^{\frac{1}{2}}e_k\le e_0$, while the conditions $\ell_k(R)\in\mathrm{Cl}$, $\ell_k(R)/L\in\mathrm{Cl}$, $\ell_k(R)\ge M\cdot r_k$ and $L^{\frac{1}{2}}e_k\le e_0$ are common to (OS) and (SF+).
By Proposition~26.3\,(b), the projected family (p5) is Borel measurable in $y$ and has the properties W1, W2, W7, W12, W13 and W14 for $\mathrm{par}$, which are the properties required in Proposition~13
(see the Remark that follows the proof of Proposition~26.6 in Chapter~26b).
For $R=R_{\mathrm{top}}$ and $k<k_{\mathrm{top}}$, the family (p5) of the primitives of (FL) on $T_k(R_{\mathrm{top}})$ is that of (ii).
Proposition~13 is applied to the projected families in Proposition~24.1, through (G3.T-13), with the clauses (d0), (d5), (d6);
in Proposition~20.3, with the clause (d0);
in Proposition~28.1, with the clause (d1);
in Proposition~28.2, with the clauses (d5), (d6);
in Proposition~48.6, with the clause (d5);
and in Proposition~35.4\,(SF4), with the clauses (d1)--(d3).
\end{named}

\begin{named}{Remark}[Where the mass enters]
The bound $\mu\le1$ is used in the summand $L^{\frac{1}{2}}$ of $K_1$, hence in $C_I$, $a_0$, $\gamma_{(5)}$ and $C_5$, that is, in $\gamma^{F}$ and in $a'$ (Proposition~1.6);
in $k_1$ (Proposition~2.3\,(b));
in the summand $1$ of the inner bracket of $C_T$ (Proposition~9\,(c), (d));
and in the summand $1$ of (I.L0.12), hence in $C_d$ (Proposition~13\,(d1) at $k=0$).
Differences and commutators of $F_y$ contain no mass.
The relations with complex conjugation and with adjoints (Propositions~1.6\,(a), (d), 2.2\,(d), (e), 2.5\,(f), 3\,(h), the property W16, Propositions~9\,(e), 10\,(d)--(g), and Proposition~12) rest on the reality of the mass.
The inequality $\mathbf{m}>0$ is a hypothesis of the clause (S7) of Proposition~1.2 and of the clause (d) of Proposition~1.5.
\end{named}

\begin{named}{Remark}
Proposition~14.1 corresponds to \cite[Theorem~16.32]{QEDI} (see \cite[Section~16.14]{QEDI}) and to \cite[Theorem~5.72]{QEDI} together with \cite[Section~5.4.4]{QEDI};
the properties FE1$^{\mathrm{ch}}$--FE8$^{\mathrm{ch}}$ correspond to \cite[Propositions~5.42--5.49]{QEDI}.
It differs from these statements in the following points.
All constants are those of the one ordered list of Proposition~14.2, there is one bound $e_0$ on the charge, and $\kappa'=\kappa_{\mathrm{w}}+\kappa_0$ is fixed (Section~I.8.2, line (C4), and Section~I.8.4).
In \cite[Theorem~16.32]{QEDI} the number $\kappa'>0$ is left free, and the charge is bounded by a number $e^{F}_0$ of its own, which is defined through conditions in which $C_d$ is formed with $C_w=C^{F}$.
Here no such number occurs: the conditions (Z2.s1)--(Z2.s5) are conditions on $e_0$, the condition (Z2.s5) being formed with the constant $C_d$ of Proposition~13, which contains the constant $C_w$ of $\mathrm{par}$;
and the condition (A10) of \cite[Section~5.6.2]{QEDI} is a condition on the constants alone.
It consists of the inequality $e_0\le e^{F}_0$, of the four inequalities
\begin{gather*}
v_k\ge9,\qquad8M\cdot\left(\frac{v_k}{4}+4\right)\cdot e_k^{\varepsilon}\le1,\\
8\,C_B\cdot v_k^{p}\le e_k^{-\frac{3}{4}+3\varepsilon},\qquad8\,e_k^{\frac{3}{4}-\frac{7\varepsilon}{2}}\le1,
\end{gather*}
which are imposed for every $e_k\in(0,e_0]$, with $v_k=-\log e_k$, and of the inequality
\[
8\,c_{\mathrm{rb}}\cdot L^{3}\cdot C_d\cdot M^{3}\cdot e_0^{2\varepsilon}\le1,
\]
in which $C_d$ is formed with the constant $C_w$ of $\mathrm{par}$.
The four inequalities imposed for every $e_k\in(0,e_0]$ are the four inequalities (A10) of Proposition~14.2\,(b), written in the variable $v=-\log e_k$ and with the constants of Proposition~14.2\,(a);
the last inequality, formed with the constant $C_d$ of Proposition~13, is the condition (Z2.s5), since the left side of (Z2.s5) decreases in $v$, so that (Z2.s5) holds for all $v\ge-\log e_0$ as soon as it holds at $v=-\log e_0$;
the inequality $e_0\le e^{F}_0$ is not imposed, since no number $e^{F}_0$ occurs.
In the order of Proposition~14.2\,(a), the constants $C_w$ and $C_d$ stand in the line (C6), and the bound $e_0$, which is subject to (A10) and (Z2.s1)--(Z2.s5), stands in the line (C8).
The numbered properties are stated in the one norm $|\cdot|_c$ of Section~I.1.5.
At the level $k=0$ the property FE4$^{\mathrm{ch}}$ is one bound, which concerns both fields:
by Proposition~6\,(e), (f), for $A\in\tilde{G}_0$ one has $H_0(A)=\bar{H}_0(A)=1$ on $\mathfrak{S}^{\mathrm{ch}}(1,\ell)$, and for every $f\in\mathfrak{S}^{\mathrm{ch}}(1,\ell)$
\[
\left|H_0^{\#}(A)f\right|_{\infty}\le\left(1+\mathrm{e}\right)\cdot|f|_{\infty}\le4|f|_{\infty};
\]
the same bound holds for the pair $\left(\bar{H}_0(A),\delta_{\alpha,-A}\bar{H}_0(A)\right)$, and $4\le C^{F}$.
The localised bounds of FE4$^{\mathrm{ch}}$ are stated for $k\ge1$ only.
In \cite[Theorem~16.32(i)]{QEDI}, by contrast, all bounds of \cite[Proposition~5.45]{QEDI}, the localised ones and those of the barred field included, are asserted for every $k\ge0$.
At $k=0$ these bounds follow from $H_0(A)=\bar{H}_0(A)=1$,
and the second localised bound is void, since at $k=0$ there is no short pair.
As for the assertion (v), in \cite{QEDI} the numbers $\mathfrak{m}_0$, $m_0$, $h'$ of \cite[Definition~5.98]{QEDI} have the same values $1$, $0$, $h$:
one has $\mathfrak{m}_0=1$ by \cite[Theorems~5.72 and 16.32]{QEDI} and \cite[Definition~5.98]{QEDI}, and hence $m_0=0$ and $h'=h$, as noted in \cite[Section~5.4.7]{QEDI}.
The first sentence of (v) gives the same range $(0,1]$ of the mass for the operators of Proposition~14.1, so that these values stand in the place of the three symbols (Section~I.6.1).
Only the second sentence of (v) has no counterpart in \cite[Theorems~5.72 and 16.32]{QEDI}:
it asserts, at $\mu=0$, the bijectivity of $F_y(A)$ and the bounds of FE1$^{\mathrm{ch}}$, FE2$^{\mathrm{ch}}$ and those of FE3$^{\mathrm{ch}}$ at $k=0$.
\end{named}

\subsubsection*{Proposition 14.2 (the ordered list)}

\begin{named}{Proposition 14.2}[The ordered list]
\begin{itemize}
\item[(a)] \emph{The constants, in the order in which they are chosen.}
The entries of (a) are the rows of a table and are numbered $[1]$--$[136]$ in the order in which they stand.
The row of a constant gives its definition (the column ``definition''), the letters on which it depends (the column ``depends on'', after the words ``Depends on'') and the statement or section to which the constant belongs, its reference (after the word ``Reference:'').
An entry that prescribes a choice gives the set in which its constant is chosen.
Every other constant is fixed in one of the following two ways.
In the first case it is defined by what its entry gives, namely numerals, a formula, or a description as an eigenvalue, a supremum, or a least or a largest element, or by the formula at the place that its entry names.
These places are the display of \cite[Section~15.7]{QEDI} that precedes \cite[Proposition~15.22]{QEDI}, for $c_P$;
the display of \cite[Section~15.7]{QEDI} that precedes \cite[Proposition~15.24]{QEDI}, for $\gamma_P$ and $C_P$;
the display of \cite[Section~15.9]{QEDI}, formed with $s_d=(2L+1)M$, for $C_B$;
and, for the numbers $C_{\mathrm{lS}}$ in $C_6$, $C_{\mathrm{iv}}$ in $\gamma_{\Sigma}$ and $C_{\mathrm{Flg}}$ in $C_{\Sigma}$, the formulas of Propositions~30.1 and 30.3.
The number $c_P$ belongs to the first case, as the number defined by its display, although \cite[Proposition~5.33]{QEDI} is worded as an existence statement.
In the second case the constant is the one whose existence is asserted by the statement named as its reference;
this is the case for the entries $A_{\mathrm{ak}},A'_{\mathrm{ak}}$; $C_G$; $C_E$; $\mathrm{Csy}_{\mathrm{ak}},\mathrm{Cpv}_0$; $\mathrm{Cah}_i,\mathrm{Cpv}_1,\mathrm{Cpv}_2,\mathrm{epv},\mathrm{Cpv}_{\beta}$; $C_{\mathrm{ct}},C_1^{\mathrm{N2}},C_2^{\mathrm{N2}}$, and the constant is then the number that is given at the place of Part~V named in the entry, which is taken as its definition.
The constant $c_\rho$ is the largest number with its property, and the absolute numbers $c(1,3)$, $c(0,4)$ of \cite[Lemma~12.4]{QEDI}, which enter $C_G$ and $C_6$, are fixed once and for all.
Hence every entry that does not prescribe a choice determines one value as soon as the letters of the preceding entries have values.
The letters named after ``Depends on'' determine, together with $\rho$ and $w$ (Section~I.8.1), the constant, and for a constant that is chosen they determine the set in which it is chosen.
Where the property that defines a constant contains a letter of a later entry ($\alpha$ in Proposition~1.5\,(f), $\bar{m}$ in Proposition~36\,(c)), one number serves for the whole range of that letter.
The sets in which $L$, $M$, $(r,p_0,p,\varepsilon,s)$ and $e_0$ are chosen are not empty, and $e_1$, $e_G$, $e_2$ are positive, by Proposition~14.3.
The constants of the entries $C_E$, $\mathrm{Csy}_{\mathrm{ak}}$, $\mathrm{Cpv}_0$ and of the entries of the line (C10) from $D_{\mathrm{be}}$ on exist and are finite, respectively positive, by the statements named as their references, and, for $e_2$, by Proposition~14.3\,(g).
The other entries up to $\sigma$ whose reference is a later statement, among them the whole line (C6'), are defined by the formulas given, and their ranges are asserted in (f).
The assertions (c)--(f) hold for every choice that (a) and (b) admit.

Put, for $\gamma>0$,
\begin{gather*}
\mathrm{lg}:=\frac{1}{2}\log L,\qquad c(\gamma):=\sum_{z\in\mathbb{Z}^3}\exp\left(-\gamma|z|\right),\qquad c_{\mathrm{bd}}(\gamma):=3\exp(3\gamma)\cdot c(\gamma),\\
\mathrm{ssum}(\gamma):=\left(\frac{12}{\gamma}\right)^{3},\qquad \mathrm{lh}(v):=L\cdot M\cdot v^{r}.
\end{gather*}
The lines (C0r) and (C4r) consist of the entries listed under these marks: the constants of the lines (C0r), (C4r) of Section~I.8.2, and the numbers $M^{F}$, $M^{P}$, which enter the conditions (M5), (M6) on $M$ of the line (C5) of Section~I.8.2.
The constants $K_1$, $k_1$ and $M^{P}$ depend neither on $\rho$ nor on the weights $w$;
every other constant of these two lines depends on $\rho$ or on $w$, the number $M^{F}$ through $a'$.
The line (C6') is the final part of the line (C6);
it consists of the constants that correspond to those of \cite[Section~20.1.7]{QEDI}.

\smallskip\noindent\emph{Line (C0).}
\begin{description}
\item[{[1]}\ $\rho$] A number in $(0,1)$.
Reference: Section~I.0.5.
\item[{[2]}\ $w$] $w=(6,6\sqrt{2},6\sqrt{2},6)$, and $\mathrm{sw}=12+12\sqrt{2}$.
Reference: Section~I.0.5.
\item[{[3]}\ $\bar{\varepsilon}$] $\bar{\varepsilon}=\frac{1}{200}$.
Reference: \cite[Definition~5.97]{QEDI}, stage~1; Section~I.8.2.
\item[{[4]}\ $c_X$] $c_X=54$.
Reference: \cite[Section~13.1.6]{QEDI}.
\item[{[5]}\ $\kappa_0$] $\kappa_0=54\log72$.
Reference: \cite[Section~13.1.6]{QEDI}.
\item[{[6]}\ $K_0$] $K_0=\frac{\exp(\kappa_0)}{36}$.
Depends on $\kappa_0$.
Reference: \cite[Section~13.1.6]{QEDI}.
\item[{[7]}\ $c_{\mathrm{rb}}$] $c_{\mathrm{rb}}=c_X\cdot K_0$.
Depends on $c_X$, $K_0$.
Reference: \cite[Section~13.1.6]{QEDI}.
\item[{[8]}\ $c_{\mathrm{fl}}$] $c_{\mathrm{fl}}=61\,c_X\cdot K_0$.
Depends on $c_X$, $K_0$.
Reference: \cite[Section~13.1.6]{QEDI}; the Propositions~16.m and 18.m.
\item[{[9]}\ $b$] $b=1$, kept as a letter.
Reference: Section~I.0.5.
\end{description}

\smallskip\noindent\emph{Line (C0r).}
\begin{description}
\item[{[10]}\ $\lambda_{\mathrm{int}}$] The least non-zero eigenvalue of $\sum_{i\le2}w_i(1-R_i)$ on $\ell^{2}(O_{\mathrm{cub}})$.
Depends on $w$.
Reference: (I.O8); Proposition~1.1.
\item[{[11]}\ $\lambda_1$] $\lambda_1=6+9\sqrt{2}-\sqrt{198-36\sqrt{2}}$.
Depends on $w$.
Reference: Proposition~37.1\,(b).
\item[{[12]}\ $\hat{g}$] $\hat{g}=\frac{\rho\cdot\lambda_{\mathrm{int}}}{2}$.
Depends on $\rho$, $\lambda_{\mathrm{int}}$.
Reference: Proposition~1.1.
\item[{[13]}\ $c_{\mathrm{hl}}$] $c_{\mathrm{hl}}=\sqrt{3}\,(1+\rho)$.
Depends on $\rho$.
Reference: Proposition~1.1.
\item[{[14]}\ $c_{*}$] $c_{*}=\frac{\min(w_3,\lambda_{\mathrm{int}})}{18}$.
Depends on $w$, $\lambda_{\mathrm{int}}$.
Reference: Proposition~1.1.
\item[{[15]}\ $\theta_0$] The number $\theta_0>0$ with $3(1+\rho)\left(\exp(\theta_0)-1\right)=\frac{\hat{g}}{2}$.
Depends on $\rho$, $\hat{g}$.
Reference: Proposition~1.2.
\item[{[16]}\ $K_\rho$] $K_\rho=\rho+\frac{2c_{\mathrm{hl}}^{2}}{\hat{g}}$.
Depends on $\rho$, $c_{\mathrm{hl}}$, $\hat{g}$.
Reference: Proposition~1.2.
\item[{[17]}\ $t_0$] $t_0=\frac{1}{4K_\rho^{2}}$.
Depends on $K_\rho$.
Reference: Proposition~1.2.
\item[{[18]}\ $c_{\mathrm{el}}$] $c_{\mathrm{el}}=\min\left\{0.22,\ \frac{\rho\cdot c_{*}\cdot t_0}{\pi\sqrt{3}}\right\}$.
Depends on $\rho$, $c_{*}$, $t_0$.
Reference: Proposition~1.2.
\item[{[19]}\ $K_{\mathrm{Str}}$] The supremum of the norm of $\mathrm{Dsc}_{\mathbf{m}}(p)$ over the strip of width $\theta_0$ and over $0\le\mathbf{m}\le1$.
Depends on $\rho$, $w$, $\theta_0$, $\hat{g}$.
Reference: Proposition~1.2\,(S3).
\item[{[20]}\ $\kappa_{(0)}$] $\kappa_{(0)}=\min\left\{\frac{\theta_0}{2},\ \frac{c_{\mathrm{el}}\cdot\theta_0}{12K_{\mathrm{Str}}}\right\}$.
Depends on $\theta_0$, $c_{\mathrm{el}}$, $K_{\mathrm{Str}}$.
Reference: Proposition~1.2\,(S7).
\item[{[21]}\ $A_{\mathrm{ak}},A'_{\mathrm{ak}}$] The constants of Proposition~1.2\,(S6), (S7), one for every multi-index $\mathrm{ak}$: the numbers given by the recursions of Chapter~1a, Sections~2.6 and 2.7.
Depend on $c_{\mathrm{el}}$, $K_{\mathrm{Str}}$, $\theta_0$.
Reference: Proposition~1.2\,(S6), (S7).
\item[{[22]}\ $c_b$] $c_b=\left[3A_1^{2}+(30+24\sqrt{2})A_2^{2}\right]^{-1}$, where
\[
A_1=\frac{2\sqrt{2}\,K_\rho}{\rho\cdot c_{*}}\left(1+\frac{6(1+\rho)}{\hat{g}}\right),\qquad A_2=\frac{1+c_{\mathrm{hl}}\cdot A_1}{\hat{g}}.
\]
Depends on $\rho$, $c_{*}$, $K_\rho$, $\hat{g}$, $c_{\mathrm{hl}}$.
Reference: Proposition~1.3\,(b).
\item[{[23]}\ $C_{\mathrm{Po}}$] $C_{\mathrm{Po}}=\frac{1}{8c_{*}}+\frac{1}{\lambda_{\mathrm{int}}}$.
Depends on $c_{*}$, $\lambda_{\mathrm{int}}$.
Reference: Proposition~1.4.
\item[{[24]}\ $C_G$] The constant of Proposition~1.5: the number $C_G$ of Chapter~1b, Section~1.4.
Depends on $\rho$, $c_{\mathrm{el}}$, $K_{\mathrm{Str}}$, $\theta_0$, $\hat{g}$, $c_{\mathrm{hl}}$, $\kappa_{(0)}$, $A_{\mathrm{ak}}$, $A'_{\mathrm{ak}}$.
Reference: Proposition~1.5.
\item[{[25]}\ $C_6$] $C_6=26\left[\mathrm{cF13}\cdot C_{\mathrm{lS}}+\frac{16\,C_{\mathrm{lH}}}{(\mathrm{e}\cdot\theta_0)^{2}}\right]$, with the constants of Proposition~30.1:
$\mathrm{cF13}=c(1,3)$ is the absolute number of \cite[Lemma~12.4]{QEDI},
$C_{\mathrm{lH}}=188\left(1+\frac{6(1+\rho)}{\hat{g}}\right)$,
and $C_{\mathrm{lS}}$ is formed with $K_{\mathrm{Str}}$, $\theta_0$ and the numbers $A'_{\mathrm{ak}}$, $|\mathrm{ak}|\le3$, of Proposition~1.2\,(S7);
the constant $c_{\mathrm{el}}$ enters $C_6$ only through these numbers.
Depends on $\rho$, $c_{\mathrm{el}}$, $K_{\mathrm{Str}}$, $\theta_0$, $\hat{g}$, $\kappa_{(0)}$, $A'_{\mathrm{ak}}$.
Reference: the Propositions~30.m.
\item[{[26]}\ $c_\rho$] The largest number with which, at the given $\rho$, the inequality
\[
D_0^{*}D_0\ge c_\rho\,\nabla^{*}\nabla
\]
holds on $\ell^{2}(\mathbb{L}^{\mathrm{f}}_0;\mathbb{C}^4)$ and on the spinors of every torus $(\Gamma,\chi)$, with $\Gamma$ a subgroup of $\mathbb{Z}^3$ of rank three and $\chi$ a unitary character;
here $D_0$ is the operator at the connection $0$, without mass and with the standard weights, and $\nabla^{*}\nabla$ is the weighted covariant Laplacian at the connection $0$ \cite[Definition~3.17]{QEDII}.
The inequality does not contain the point $v_0$ of Section~I.1.3.
Depends on $\rho$, $w$.
Reference: Section~I.8.2, line (C0r);
Proposition~14.2\,(f), paragraph ``Ranges'' (the number $c_\rho$ exists, and $0<c_\rho<\infty$; proved in Chapter~14b, rows~24--26);
Proposition~1.3\,(b), which gives $c_\rho\ge c_b$ (carried to all pairs $(\Gamma,\chi)$ in Chapter~14b, rows~24--26).
The number $c_\rho$ enters no other entry of (a) and no condition of Proposition~14.2, that is, none of the conditions of (b), (c) and (d);
no statement other than Proposition~14.2 contains it.
\item[{[27]}\ $C_E$] The constant of Proposition~37.1\,(f): the number [1.10] of Chapter~37.
Depends on $\rho$, $w$.
Reference: Proposition~37.1.
\item[{[28]}\ $\mathrm{Csy}_{\mathrm{ak}},\mathrm{Cpv}_0$] The constants of Proposition~37.1\,(g) ($|\mathrm{ak}|\le4$) and of Proposition~36\,(c): the numbers [1.12] of Chapter~37, and
\[
\mathrm{Cpv}_0=10^{8}\max\left\{1,\ \max\left\{\mathrm{Csy}_{\mathrm{ak}}:\ |\mathrm{ak}|\le4\right\}\right\}^{2}
\]
(Chapter~36, Section~4.7).
Depend on $\rho$, $w$.
Reference: Propositions~37.1 and 36.
\item[{[29]}\ $c_{\mathrm{N3}}$] $c_{\mathrm{N3}}=\min\left(\frac{\rho}{\pi},\ \frac{\rho^{3}\cdot\lambda_{\mathrm{int}}}{98\sqrt{3}\,\pi^{3}}\right)$.
Depends on $\rho$, $\lambda_{\mathrm{int}}$.
Reference: the Propositions~30.m.
\item[{[30]}\ $\varepsilon_5$] The function
\[
\varepsilon_5(x)=C_6\cdot\log\frac{1}{1-\exp\left(-\kappa_{(0)}\,x\right)},\qquad x>0.
\]
Depends on $C_6$, $\kappa_{(0)}$.
Reference: the Propositions~30.m.
\end{description}

\smallskip\noindent\emph{Line (C1).}
\begin{description}
\item[{[31]}\ $\alpha$] A number in $(1-\bar{\varepsilon},1)$.
Depends on $\bar{\varepsilon}$.
Reference: Section~I.8.2.
\end{description}

\smallskip\noindent\emph{Line (C2).}
\begin{description}
\item[{[32]}\ $\kappa$] A number $\kappa\ge7\kappa_0+7$.
Depends on $\kappa_0$.
Reference: Section~I.8.2.
\end{description}

\smallskip\noindent\emph{Line (C3).}
\begin{description}
\item[{[33]}\ $L$] An odd integer with (L1)--(L5).
Depends on $\bar{\varepsilon}$, $c_{\mathrm{fl}}$, $c_{\mathrm{rb}}$, $\kappa$, $\kappa_0$.
Reference: Section~I.8.2.
\end{description}

\smallskip\noindent\emph{Line (C4).}
\begin{description}
\item[{[34]}\ $b_L$] $b_L=\frac{b}{L}$.
Depends on $b$, $L$.
Reference: Section~I.3.4.
\item[{[35]}\ $b_{-}$] $b_{-}=b\left(1-\frac{1}{L}\right)$.
Depends on $b$, $L$.
Reference: Section~I.3.4.
\item[{[36]}\ $b_{+}$] $b_{+}=b$.
Depends on $b$.
Reference: Section~I.3.4.
\item[{[37]}\ $N_{*}$] $N_{*}=8L^{3}$.
Depends on $L$.
Reference: Section~I.6.4.
\item[{[38]}\ $a_L$] $a_L=\frac{L^{3}-1}{(L-1)^{3}}$.
Depends on $L$.
Reference: Section~I.7.6.
\item[{[39]}\ $G_L$] $G_L=\left(1+\frac{L}{3.2}\right)\frac{2L^{2}}{2.56}$.
Depends on $L$.
Reference: \cite[Section~21.1.7]{QEDI}.
\item[{[40]}\ $K_L$] $K_L=\left(1+L+L^{2}\right)\left(\log L+L\cdot G_L\right)$.
Depends on $L$, $G_L$.
Reference: \cite[Section~21.1.7]{QEDI}.
\item[{[41]}\ $N_c$] $N_c=(2L+1)^{3}$.
Depends on $L$.
Reference: Section~I.1.1.
\item[{[42]}\ $c_l$] $c_l=2L$.
Depends on $L$.
Reference: Section~I.1.1.
\item[{[43]}\ $d_U$] $d_U=2L+1$.
Depends on $L$.
Reference: Section~I.1.1.
\item[{[44]}\ $c_C$] $c_C=3L^{\frac{9}{2}}$.
Depends on $L$.
Reference: \cite[Propositions~5.35 and 15.21]{QEDI}.
\item[{[45]}\ $C_Q$] $C_Q=1$.
Reference: \cite[Proposition~5.50]{QEDI}.
\item[{[46]}\ $\alpha^{P}$] $\alpha^{P}=\frac{1}{4}$.
Reference: \cite[Section~15.11.2]{QEDI}.
\item[{[47]}\ $\alpha^{F}$] $\alpha^{F}=\frac{1}{54}$.
Reference: Proposition~3.
\item[{[48]}\ $\alpha_0$] $\alpha_0=\min\left(\alpha^{F},\alpha^{P}\right)=\frac{1}{54}$.
Depends on $\alpha^{F}$, $\alpha^{P}$.
Reference: \cite[Section~5.5.3]{QEDI}.
\item[{[49]}\ $c_P$] The constant $c_P$ defined at the beginning of \cite[Section~15.7]{QEDI}.
Depends on $L$, $C_Q$.
Reference: \cite[Proposition~5.33]{QEDI}.
\item[{[50]}\ $\gamma_P$] The constant $\gamma_P$ defined in \cite[Section~15.7]{QEDI}.
Depends on $L$, $c_P$.
Reference: \cite[Propositions~5.32--5.34]{QEDI}.
\item[{[51]}\ $\gamma^{P}$] $\gamma^{P}=\frac{\gamma_P}{2}$.
Depends on $\gamma_P$.
Reference: \cite[Section~15.10.3]{QEDI}.
\item[{[52]}\ $C_P$] The constant $C_P$ defined in \cite[Section~15.7]{QEDI}.
Depends on $L$, $\alpha$, $c_P$.
Reference: \cite[Propositions~5.32--5.34]{QEDI}.
\item[{[53]}\ $\kappa_{\mathrm{w}}$] $\kappa_{\mathrm{w}}=\kappa+\kappa_0+3$.
Depends on $\kappa$, $\kappa_0$.
Reference: Section~I.8.2.
\item[{[54]}\ $\kappa'$] $\kappa'=\kappa_{\mathrm{w}}+\kappa_0$.
Depends on $\kappa_{\mathrm{w}}$, $\kappa_0$.
Reference: Section~I.8.2; Propositions~3, 7, 8.
\item[{[55]}\ $\hat{\kappa}$] $\hat{\kappa}=N_c\cdot\kappa+\log(6c_X)$.
Depends on $N_c$, $\kappa$, $c_X$.
Reference: Section~I.8.2.
\item[{[56]}\ $C'_V$] $C'_V=2\mathrm{e}\cdot(42+6b)$.
Depends on $b$.
Reference: Proposition~2.1.
\item[{[57]}\ $C_T$] $C_T=2\mathrm{e}^{\frac{3}{2}}\left[b_{-}^{-1}\left(54L\cdot\mathrm{e}^{6}+12\mathrm{e}+1\right)+1\right]$.
Depends on $b_{-}$, $L$.
Reference: Proposition~9\,(d).
\item[{[58]}\ $\mu_{-}$] $\mu_{-}=L^{-3}$.
Depends on $L$.
Reference: \cite[Section~5.5.8]{QEDI}, with $m_0=0$.
\end{description}

\smallskip\noindent\emph{Line (C4r).}
\begin{description}
\item[{[59]}\ $K_1$] $K_1=L^{\frac{1}{2}}+(Lb)^{\frac{1}{2}}$.
Depends on $L$, $b$.
Reference: Proposition~1.6.
\item[{[60]}\ $C_I$] $C_I=\frac{1}{b_{-}}+\frac{2C_{\mathrm{Po}}\cdot K_1^{2}}{c_b}+\left(\frac{2C_{\mathrm{Po}}}{c_b}\right)^{\frac{1}{2}}$.
Depends on $b_{-}$, $C_{\mathrm{Po}}$, $K_1$, $c_b$.
Reference: Proposition~1.6\,(b).
\item[{[61]}\ $c_{\mathrm{wt}}$] $c_{\mathrm{wt}}=\mathrm{e}\left(3(1+\rho)+6L^{4}b\right)$.
Depends on $\rho$, $L$, $b$.
Reference: Proposition~1.6.
\item[{[62]}\ $a_0$] $a_0=\min\left\{\frac{1}{3},\ \frac{1}{2C_I\cdot c_{\mathrm{wt}}}\right\}$.
Depends on $C_I$, $c_{\mathrm{wt}}$.
Reference: Proposition~1.6.
\item[{[63]}\ $C_R$] $C_R=\frac{18\,000\,000\,C_G}{1-\alpha}$.
Depends on $C_G$, $\alpha$.
Reference: Proposition~1.6.
\item[{[64]}\ $C_3$]
\[
C_3=\left(1+(1+\mathrm{e})L^{2}\right)\max\left\{C_R\left(\frac{\mathrm{e}^{\frac{1}{3}}}{4}+\left(L^{4}b+46\right)\mathrm{e}^{1.3}\cdot C_I\right),\ 14L^{3}\cdot C_I\right\}.
\]
Depends on $L$, $b$, $C_I$, $C_R$.
Reference: Proposition~1.6.
\item[{[65]}\ $C_5$] $C_5=\mathrm{e}^{2}L\cdot C_3$ (and $C_5\ge1$).
Depends on $L$, $C_3$.
Reference: Proposition~1.6\,(c).
\item[{[66]}\ $\gamma_{(5)}$] $\gamma_{(5)}=\frac{a_0}{L}$.
Depends on $a_0$, $L$.
Reference: Proposition~1.6\,(c).
\item[{[67]}\ $c_{\mathrm{re}}$] $c_{\mathrm{re}}=\min\left\{b\left(1-\frac{1}{L}\right),\ \frac{\rho}{2C_{\mathrm{Po}}\cdot L}\right\}$.
Depends on $b$, $L$, $\rho$, $C_{\mathrm{Po}}$.
Reference: Proposition~1.7.
\item[{[68]}\ $a_z$] $a_z=\min\left\{\frac{1}{3L},\ \frac{c_{\mathrm{re}}}{2\mathrm{e}\cdot L\,(6+3b)}\right\}$.
Depends on $L$, $c_{\mathrm{re}}$, $b$.
Reference: Proposition~2.3\,(b).
\item[{[69]}\ $k_1$] $k_1=60+b$.
Depends on $b$.
Reference: Proposition~2.3\,(b).
\item[{[70]}\ $C_{00}$] $C_{00}=48\cdot\frac{2L}{c_{\mathrm{re}}}\left(1+2k_1\right)$.
Depends on $L$, $c_{\mathrm{re}}$, $k_1$.
Reference: Proposition~2.3\,(b).
\item[{[71]}\ $C_z$] $C_z=\max\left(C_5,C_{00}\right)$.
Depends on $C_5$, $C_{00}$.
Reference: Proposition~2.3\,(c).
\item[{[72]}\ $\gamma_z$] $\gamma_z=\min\left(\gamma_{(5)},a_z\right)$.
Depends on $\gamma_{(5)}$, $a_z$.
Reference: Proposition~2.3\,(c).
\item[{[73]}\ $\gamma^{F}$] $\gamma^{F}=\frac{\gamma_z}{2}$ (and $\gamma^{F}\le\frac{1}{6L}$).
Depends on $\gamma_z$.
Reference: Proposition~2.3\,(c).
\item[{[74]}\ $a'$] $a'=C_z\cdot c(\gamma^{F})$ (and $a'\ge1$).
Depends on $C_z$, $\gamma^{F}$.
Reference: Proposition~2.3\,(c).
\item[{[75]}\ $\gamma_{\mathrm{w}}$] $\gamma_{\mathrm{w}}=\min\left(\gamma^{F},\gamma^{P}\right)$ (and $\gamma_{\mathrm{w}}>0$).
Depends on $\gamma^{F}$, $\gamma^{P}$.
Reference: \cite[Section~5.5.3]{QEDI}.
\item[{[76]}\ $M^{F}$]
\[
M^{F}=\max\left\{4L+1,\ \left(108\pi\cdot C'_V\cdot a'\right)^{2},\ 108\pi\cdot C'_V\cdot a'\cdot\exp(27\kappa')\right\}.
\]
Depends on $L$, $C'_V$, $a'$, $\kappa'$.
Reference: Propositions~2.8 and 3.
\item[{[77]}\ $M^{P}$] The least power $M_0\ge L^{2}$ of $L$ such that $2\alpha_0\cdot402\cdot\log x\le\gamma_P\cdot x$ for all $x\ge M_0$.
Depends on $L$, $\alpha_0$, $\gamma_P$.
Reference: \cite[Section~15.10.3]{QEDI}.
\end{description}

\smallskip\noindent\emph{Line (C5).}
\begin{description}
\item[{[78]}\ $M$] A power of $L$ with (M1)--(M6).
Depends on $L$, $\gamma_{\mathrm{w}}$, $\alpha_0$, $\hat{\kappa}$, $M^{F}$, $M^{P}$.
Reference: Section~I.8.2.
\end{description}

\smallskip\noindent\emph{Line (C6).}
\begin{description}
\item[{[79]}\ $C_{\mathrm{By}}$] $C_{\mathrm{By}}=(1+2\mathrm{e})\sqrt{2}\,\max\left(1,\frac{1}{b_{-}}\right)$.
Depends on $b_{-}$.
Reference: Proposition~5\,(b).
\item[{[80]}\ $C_{\Gamma}$] $C_{\Gamma}=\frac{2(1+2\mathrm{e})}{b_{-}}+128\,\mathrm{e}\,(1+2\mathrm{e})^{2}a'$.
Depends on $b_{-}$, $a'$.
Reference: Proposition~5\,(d).
\item[{[81]}\ $C_1$]
\[
C_1=\max\left\{304\,\mathrm{e}\cdot\max(1,b)\cdot a'\cdot M^{27\alpha_0},\ C_{\Gamma}\cdot M^{27\alpha_0},\ 32\,\mathrm{e}\cdot\max\left(1,b^{2}\right)\cdot a'\cdot\exp(27\kappa')\right\}.
\]
Depends on $b$, $a'$, $M$, $\alpha_0$, $C_{\Gamma}$, $\kappa'$.
Reference: Propositions~7 and 8.
\item[{[82]}\ $C^{F}$]
\[
C^{F}=2\max\left\{C_1,\ C_{\mathrm{By}},\ 304\,\mathrm{e}\cdot\max(1,b)\cdot a',\ C_{\Gamma},\ \mathrm{e}\cdot b^{2}\cdot(32a')^{2}\right\}.
\]
Depends on $C_1$, $C_{\mathrm{By}}$, $b$, $a'$, $C_{\Gamma}$.
Reference: Proposition~14.1.
\item[{[83]}\ $C^{P}$] $C^{P}=116\,C_P\cdot M^{402\alpha_0}$.
Depends on $C_P$, $M$, $\alpha_0$.
Reference: \cite[Section~15.10.3]{QEDI}.
\item[{[84]}\ $C_{w}^{(0)}$] $C_{w}^{(0)}=\max\left(C^{F},C^{P}\right)$.
Depends on $C^{F}$, $C^{P}$.
Reference: \cite[Section~5.5.3]{QEDI}.
\item[{[85]}\ $C_{w}^{(1)}$] $C_{w}^{(1)}=\max\left(C_{w}^{(0)},\ 4\exp(2\gamma_{\mathrm{w}})\cdot M^{27\alpha_0}\right)$.
Depends on $C_{w}^{(0)}$, $\gamma_{\mathrm{w}}$, $M$, $\alpha_0$.
Reference: \cite[Section~5.5.3]{QEDI}.
\item[{[86]}\ $C_{w}$] $C_{w}=64\,K_0\cdot C_{w}^{(1)}$.
Depends on $K_0$, $C_{w}^{(1)}$.
Reference: \cite[Section~5.4.1]{QEDI}; the Propositions~26.m.
\item[{[87]}\ $c_{\gamma}$] $c_{\gamma}=\sup_{t}\sum_{t'}\nu(t')\exp\left(-\gamma_{\mathrm{w}}\,d(t,t')\right)$ (Section~I.5.8);
the letter $M$, which enters through the class $\mathrm{Cl}$, drops out.
Depends on $\gamma_{\mathrm{w}}$, $L$.
Reference: Section~I.5.8.
\item[{[88]}\ $C_W$] $C_W=28\,C_{w}^{2}\cdot c_{\gamma}$.
Depends on $C_{w}$, $c_{\gamma}$.
Reference: (I.L0.19); the Propositions~24.m.
\item[{[89]}\ $C_B$] The constant $C_B$ defined in \cite[Section~15.9]{QEDI}, with $s_d=(2L+1)M$.
Depends on $C_P$, $\gamma_P$, $L$, $M$.
Reference: \cite[Proposition~5.37]{QEDI}.
\item[{[90]}\ $C_d$] $C_d=16\,c_X\cdot C_{w}\cdot\max\left\{\frac{\pi}{2}\,b\cdot\mathrm{e}\,(1+2\mathrm{e}),\ 2+58\,\mathrm{e}+b\cdot\mathrm{e}\right\}$.
Depends on $c_X$, $C_{w}$, $b$.
Reference: Proposition~13.
\item[{[91]}\ $c_{\kappa}$] $c_{\kappa}=\sup_{k,\ell,\square}\sum_{X\in D_k,\ X\supseteq\square}\exp\left(-\kappa\,d_M(X)\right)$;
one has $c_{\kappa}\le K_0$, the letters $M$ and $k$ drop out, and $L$ remains through the class $\mathrm{Cl}$.
Depends on $\kappa$, $L$.
Reference: Proposition~14.2; Proposition~32.1.
\item[{[92]}\ $C_0$] $C_0=\max\left\{C_P\cdot c_{\mathrm{bd}}(\gamma_P),\ C^{F},\ 4b_{+}+16\,C^{F}\cdot c(\gamma^{F})\right\}$.
Depends on $C_P$, $\gamma_P$, $C^{F}$, $b_{+}$, $\gamma^{F}$.
Reference: Proposition~14.2; the Propositions~20.m.
\item[{[93]}\ $\gamma_{\Sigma}$] $\gamma_{\Sigma}=\min\left\{1,\ \frac{1}{2\mathrm{e}\,(6+b_{+})\,C_{\mathrm{iv}}}\right\}$, with the constant $C_{\mathrm{iv}}$ of Proposition~30.3;
it is the rate of the dependence on the twist for the interval $[b_{-},b_{+}]$.
Depends on $b_{-}$, $b_{+}$, $\rho$, $w$.
Reference: the Propositions~30.m.
\item[{[94]}\ $C_{\Sigma}$] $C_{\Sigma}=2\left(4\pi^{2}+\frac{\pi^{4}}{45}\right)C_{\mathrm{Flg}}$, with the constant $C_{\mathrm{Flg}}$ of Proposition~30.3;
it is the constant of the dependence on the twist for the interval $[b_{-},b_{+}]$.
Depends on $b_{-}$, $b_{+}$, $\rho$, $w$.
Reference: the Propositions~30.m.
\item[{[95]}\ $c_S$] $c_S=10^{5}\,c_X^{5}\cdot K_0^{5}\cdot\exp(6\kappa)$.
Depends on $c_X$, $K_0$, $\kappa$.
Reference: \cite[Section~17.1.1]{QEDI}.
\end{description}

\smallskip\noindent\emph{Line (C6').}
\begin{description}
\item[{[96]}\ $\mathbf{S}$] $\mathbf{S}=C^{F}\cdot\mathrm{ssum}\left(\frac{\gamma^{F}}{2}\right)$, where $\mathrm{ssum}(\gamma)=\left(\frac{12}{\gamma}\right)^{3}$.
Depends on $C^{F}$, $\gamma^{F}$.
Reference: Propositions~35.4\,(SF2) and 33.3.
\item[{[97]}\ $\mathrm{hL}$] $\mathrm{hL}=6\,\mathbf{S}$.
Depends on $\mathbf{S}$.
Reference: Proposition~33.2\,(L1).
\item[{[98]}\ $C_K$] $C_K=226\,\mathbf{S}^{2}$.
Depends on $\mathbf{S}$.
Reference: the Propositions~33.m;
in \cite[Section~20.1.7]{QEDI} the corresponding constant is $100\,\mathbf{S}^{2}$.
\item[{[99]}\ $C_K'$] $C_K'=1328\,\mathbf{S}^{2}$.
Depends on $\mathbf{S}$.
Reference: the Propositions~33.m;
in \cite[Section~20.1.7]{QEDI} the corresponding constant is $400\,\mathbf{S}^{2}$.
\item[{[100]}\ $C_H$] $C_H=203\,M\cdot\left(C^{F}\right)^{2}$.
Depends on $M$, $C^{F}$.
Reference: the Propositions~33.m;
in \cite[Section~20.1.7]{QEDI} the corresponding constant is $100\,M\cdot C_F^{2}$.
\item[{[101]}\ $\mathrm{CVj}$] $\mathrm{CVj}=10^{3}\,C_H^{2}\cdot\mathrm{ssum}\left(\frac{\gamma^{F}}{2}\right)^{4}\cdot K_0$.
Depends on $C_H$, $\gamma^{F}$, $K_0$.
Reference: Proposition~35.5;
in \cite[Section~20.1.7]{QEDI} this constant is denoted $C_V$.
\item[{[102]}\ $\theta_{\mathrm{CT}}$] $\theta_{\mathrm{CT}}=\frac{\mu_{-}}{12\mathrm{e}}$.
Depends on $\mu_{-}$.
Reference: Proposition~33.4.
\item[{[103]}\ $\mathrm{vt}$] $\mathrm{vt}=\min\left\{\frac{\theta_{\mathrm{CT}}}{2},\ \frac{\gamma^{F}}{8},\ \frac{1}{M}\right\}$.
Depends on $\theta_{\mathrm{CT}}$, $\gamma^{F}$, $M$.
Reference: Proposition~33.4.
\item[{[104]}\ $\mathrm{gGa}$] $\mathrm{gGa}=\left(b_{-}^{-1}+\mu_{-}^{-1}\right)^{\frac{1}{2}}$.
Depends on $b_{-}$, $\mu_{-}$.
Reference: Proposition~33.4;
in \cite[Section~20.1.7]{QEDI} this constant is denoted $g$.
\item[{[105]}\ $h_G$] $h_G=1+\mathrm{gGa}$.
Depends on $\mathrm{gGa}$.
Reference: Propositions~35.5 and 35.6.
\item[{[106]}\ $\mathrm{gGa}_1$] $\mathrm{gGa}_1=b_{-}^{-1}+\frac{4}{\mu_{-}}$.
Depends on $b_{-}$, $\mu_{-}$.
Reference: Proposition~33.4.
\item[{[107]}\ $\mathrm{alG}$] $\mathrm{alG}=2b_{-}^{-1}+4\mathrm{e}\cdot\mu_{-}^{-1}\cdot\mathrm{ssum}\left(\frac{\theta_{\mathrm{CT}}}{2}\right)$.
Depends on $b_{-}$, $\mu_{-}$, $\theta_{\mathrm{CT}}$.
Reference: Proposition~33.4;
in \cite[Section~20.1.7]{QEDI} this constant is denoted $\boldsymbol{\alpha}$.
\item[{[108]}\ $c_{\tau}$] $c_{\tau}=\min\left\{\frac{1}{8},\ \frac{\mu_{-}}{48},\ \frac{1}{292\,\mathbf{S}}\right\}$.
Depends on $\mu_{-}$, $\mathbf{S}$.
Reference: the Propositions~33.m;
in \cite[Section~20.1.7]{QEDI} the corresponding constant is $c_1$, with $\frac{1}{80\,\mathbf{S}}$ in place of $\frac{1}{292\,\mathbf{S}}$.
\item[{[109]}\ $\mathrm{CWj}$]
\[
\mathrm{CWj}=3\,\mathrm{CVj}\cdot h_G^{2}+2\,\mathrm{alG}\cdot\mathrm{CVj}^{2}\cdot h_G^{4}+4\exp(6)\cdot C_d\cdot M^{3}\cdot K_0.
\]
Depends on $\mathrm{CVj}$, $h_G$, $\mathrm{alG}$, $C_d$, $M$, $K_0$.
Reference: Proposition~35.6;
in \cite[Section~20.1.7]{QEDI} this constant is denoted $C_W$.
\end{description}

\smallskip\noindent\emph{Line (C7).}
\begin{description}
\item[{[110]}\ $r$] An integer $r\ge2$.
Reference: Section~I.8.2.
\item[{[111]}\ $p_0$] An integer $p_0>3r+2$.
Depends on $r$.
Reference: Section~I.8.2.
\item[{[112]}\ $p$] An integer $p>2p_0$.
Depends on $p_0$.
Reference: Section~I.8.2.
\item[{[113]}\ $\varepsilon$] A number in $(1-\alpha,\bar{\varepsilon}]$.
Depends on $\alpha$, $\bar{\varepsilon}$.
Reference: Section~I.8.2.
\item[{[114]}\ $s$] A number in $(0,1)$ with (S2).
Depends on $\varepsilon$, $a_L$, $N_{*}$.
Reference: Section~I.8.2.
\item[{[115]}\ $\mathfrak{e}$] $\mathfrak{e}=s\left(\frac{1}{4}-9\varepsilon\right)-12\varepsilon$.
Depends on $s$, $\varepsilon$.
Reference: Section~I.8.5.
\item[{[116]}\ $\sigma$] $\sigma=\mathfrak{e}-\frac{2a_L}{N_{*}}$.
Depends on $\mathfrak{e}$, $a_L$, $N_{*}$.
Reference: Section~I.8.5.
\end{description}

\smallskip\noindent\emph{Line (C8).}
\begin{description}
\item[{[117]}\ $e_0$] A number $e_0>0$ with (A1)--(A10), (Z2.s1)--(Z2.s5) and (Z3.s1).
Depends on $\varepsilon$, $p_0$, $p$, $c_{\mathrm{fl}}$, $M$, $L$, $C_{w}$, $c_{\gamma}$, $c_C$, $c_X$, $C_T$, $c_S$, $K_0$, $C'_V$, $a'$, $c_{\mathrm{rb}}$, $C_d$, $C_B$.
Reference: Proposition~14.2.
\end{description}

\smallskip\noindent\emph{Line (C9).}
\begin{description}
\item[{[118]}\ $e_1$] One half of the supremum of the numbers $e'\in\left(0,\min\left(e_0,\frac{1}{3}\right)\right]$ such that (e1)--(e8) hold for all $v\ge-\log e'$;
$e_1:=0$ if there is no such $e'$.
Depends on $e_0$, $L$, $M$, $b_{-}$, $r$, $p_0$, $p$, $\varepsilon$, $c_{\mathrm{fl}}$, $c_{\kappa}$, $C_0$, $\gamma_{\Sigma}$, $C_{\Sigma}$, $N_{*}$.
Reference: Proposition~14.2; Proposition~14.3\,(f).
\end{description}

\smallskip\noindent\emph{Line (C10).}
\begin{description}
\item[{[119]}\ $\bar{m}$] A number in $(0,1]$.
Reference: Section~I.0.5.
\item[{[120]}\ $h$] $h=L^{\left\lfloor\log_L\frac{1}{\bar{m}}\right\rfloor}$.
Depends on $L$, $\bar{m}$.
Reference: Section~I.0.5.
\item[{[121]}\ $e_G$] $e_G=\sup\left\{e>0:\ e\,h^{\frac{1}{2}}\le e_1,\ e\,h^{\frac{1}{2}}\log h\le1\right\}$;
$e_G:=0$ if there is no such $e$.
Depends on $e_1$, $h$.
Reference: Section~I.8.2.
\item[{[122]}\ $\beta,x_0$] A probe $\beta$ in the sense of Section~I.7.6 and a point $x_0\in\mathbb{R}^3$, the vector of the translation $\tau_{x_0}$ with which, in Proposition~43.2, the probes $\beta_{\pm}=\beta\pm\tau_{x_0}\beta$ and the tests $g_{\pm}=f\pm\tau_{x_0}f$, where $f=dd^{*}\beta$, are formed;
the pair is chosen subject to (Pb) of (d).
Depends on $\bar{m}$.
Reference: Section~I.7.6; Proposition~43.2.
\item[{[123]}\ $D_{\mathrm{be}},D_{\beta},C_{\beta},\mathbf{D},\mathrm{Cb}$] For each of the three probes $\beta'$:
$D_{\mathrm{be}}$ is the least side of a closed coordinate cube that contains $\operatorname{supp}\beta'$;
$D_{\beta}=D_{\mathrm{be}}+4$;
$C_{\beta}$ and $\mathbf{D}$ are the suprema over $\mathrm{Ind}_{\beta}$ of Section~G5.1.3;
$\mathrm{Cb}=1+h^{\frac{1}{2}}C_{\beta}$.
Depend on $\beta$, $x_0$, $h$, $L$, $\alpha$.
Reference: Proposition~35.3\,(c1).
\item[{[124]}\ $\mathrm{Cah}_i,\mathrm{Cpv}_1,\mathrm{Cpv}_2,\mathrm{epv},\mathrm{Cpv}_{\beta}$] The constants and the function of Proposition~36\,(e), (f), (g): the entries of the table of constants at the end of Chapter~36.
Depend on $\rho$, $\bar{m}$, $\beta$, $x_0$, $\mathrm{Csy}_{\mathrm{ak}}$, $\mathrm{Cpv}_0$.
Reference: Proposition~36.
\item[{[125]}\ $C_{\mathrm{ct}},C_1^{\mathrm{N2}},C_2^{\mathrm{N2}}$] The constants of Proposition~37.2\,(c), (d), (e): the numbers [2.1] of Chapter~37.
Depend on $\rho$, $\bar{m}$, $\beta$, $x_0$, $C_E$.
Reference: Proposition~37.2.
\item[{[126]}\ $C_{\mathrm{fr}}$] $C_{\mathrm{fr}}=2592\,C_d\cdot K_0\cdot C_{\beta}^{2}\left(M^{2}\cdot D_{\beta}^{3}\cdot L^{4}\cdot h^{-2}+8M^{5}\cdot h\right)$.
Depends on $C_d$, $K_0$, $D_{\beta}$, $C_{\beta}$, $M$, $L$, $h$.
Reference: Proposition~34.3\,(b).
\item[{[127]}\ $e_{\sharp}$] One half of the supremum of the numbers $e'\in\left(0,e_1h^{-\frac{1}{2}}\right]$ such that (g1)--(g9) of (d) hold for every $e\in(0,e']$;
$e_{\sharp}:=0$ if there is no such $e'$, as in Section~G5.1.9.
There is one such number for each of the three probes, and it is positive by Proposition~35.10.
Depends on $e_1$, $h$, $D_{\beta}$, $C_{\beta}$, $\mathbf{D}$, $\mathrm{Cb}$, $C_{\mathrm{fr}}$, $\mathrm{hL}$, $h_G$, $\mathrm{alG}$, $C_K$, $C_K'$, $\mathrm{gGa}_1$, $C_0$, $K_0$, $M$, $\mu_{-}$, $c_{\tau}$, $\mathrm{CWj}$, $\mathrm{vt}$, $N_{*}$, $s$, $\varepsilon$, $C_d$, $r$, $p_0$, $p$, $L$.
Reference: Proposition~35.10.
\item[{[128]}\ $C^{G}_{\beta}$]
\[
C^{G}_{\beta}=\max\left\{C_{35},\ C_{\mathrm{R4}},\ C_1^{\mathrm{N2}},\ C_2^{\mathrm{N2}}\cdot\exp\left(\bar{m}\cdot D_{\mathrm{be}}\right)\cdot h^{\frac{1}{2}}\right\},
\]
where
\[
C_{35}:=3\cdot10^{5}\,\mathrm{Cb}^{4}\cdot h^{\frac{1}{8}},\qquad C_{\mathrm{R4}}:=\mathrm{Cpv}_{\beta}+\left(\left|b(dd^{*}\beta)\right|+C_1^{\mathrm{N2}}+C_2^{\mathrm{N2}}+3\,\mathrm{Cpv}_{\beta}^{2}\right)^{\frac{1}{2}};
\]
there is one such number for each of the three probes.
Depends on $\mathrm{Cb}$, $h$, $\mathrm{Cpv}_{\beta}$, $C_1^{\mathrm{N2}}$, $C_2^{\mathrm{N2}}$, $D_{\mathrm{be}}$, $\bar{m}$, $\beta$, $x_0$.
Reference: Proposition~43.1.
\item[{[129]}\ $C_{\mathrm{mat}}$] $C_{\mathrm{mat}}=(L^{3}+1)\left(\varepsilon_5\left(\frac{1}{L}\right)+1\right)$.
Depends on $L$, $\varepsilon_5$.
Reference: the Propositions~32.m.
\item[{[130]}\ $s_B$] $s_B=\max_{\mathrm{ix}}\left(\|g_{\mathrm{ix}}\|_2+\|\beta_{\mathrm{ix}}\|_{H^{4}}\right)$, $\mathrm{ix}\in\{0,+,-\}$;
here $\beta_{\mathrm{ix}}$ and $g_{\mathrm{ix}}=dd^{*}\beta_{\mathrm{ix}}$ are those of the data of Proposition~43.2, and $\|\cdot\|_{H^{4}}$ is the norm of Proposition~47.1 at $s=4$: for a two-form $u$,
\[
\|u\|_{H^{s}}^{2}=\sum_{\mu<\nu}\frac{1}{(2\pi)^{3}}\int\mathrm{d}p\,\left(1+|p|_2^{2}\right)^{s}\cdot\left|\hat{u}_{\mu\nu}(p)\right|^{2},
\]
and $\|\cdot\|_2=\|\cdot\|_{H^{0}}$.
In this definition $s$ is the order of the norm, a real number that is taken here only at the values $4$ and $0$, and $p\in\mathbb{R}^3$ is the variable of integration, a momentum;
neither is the constant $s$ or $p$ of the line (C7), and of the letters of (a) the definition of $s_B$ contains only $\beta$ and $x_0$.
This norm is not that of the Hermite--Sobolev scale $H_{s}$ of Section~I.7.5.
Depends on $\beta$, $x_0$.
Reference: Proposition~43.2.
\item[{[131]}\ $C_{3\mathrm{p}}$] $C_{3\mathrm{p}}=\max_{\mathrm{ix}}C^{G}_{\beta_{\mathrm{ix}}}$.
Depends on $C^{G}_{\beta}$.
Reference: Proposition~43.2.
\item[{[132]}\ $C_{\Gamma}^{(16)}$] $C_{\Gamma}^{(16)}=(2\pi)^{\frac{a_L}{N_{*}}}\exp\left(\frac{K_L+C_{\mathrm{mat}}+1}{(L-1)^{3}N_{*}}\right)$.
Depends on $a_L$, $N_{*}$, $K_L$, $C_{\mathrm{mat}}$, $L$.
Reference: Proposition~43.2.
\item[{[133]}\ $C_U$]
\begin{multline*}
C_U=\max\Bigl\{C_{\Gamma}^{(16)}\left(C_{3\mathrm{p}}+1200\,N_{*}^{2}\cdot s_B^{4}+2C_{3\mathrm{p}}^{2}\right),\ C_{\Gamma}^{(16)}\left(C_{3\mathrm{p}}+2\left(12\,N_{*}\cdot s_B^{2}+C_{3\mathrm{p}}\right)\left(1+C_{3\mathrm{p}}\right)\right),\\
C_{3\mathrm{p}},\ 60\,s_B^{4}\Bigr\}.
\end{multline*}
Depends on $C_{\Gamma}^{(16)}$, $C_{3\mathrm{p}}$, $s_B$, $N_{*}$.
Reference: Proposition~43.2.
\item[{[134]}\ $C_{\bar{\kappa}}$] $C_{\bar{\kappa}}=2\left(8C_U+3C_U^{2}\right)$.
Depends on $C_U$.
Reference: Proposition~43.2.
\item[{[135]}\ $C'_{\bar{\kappa}}$] $C'_{\bar{\kappa}}=8C_{\bar{\kappa}}+3C_{\bar{\kappa}}^{2}$.
Depends on $C_{\bar{\kappa}}$.
Reference: Proposition~43.2.
\item[{[136]}\ $e_2$] The largest number with (Z10.s1), (Z10.s2) and (Z10.s4).
Depends on $e_G$, $e_{\sharp}$, $C'_{\bar{\kappa}}$, $\sigma$, $\beta$, $x_0$, $\bar{m}$.
Reference: Proposition~43.2.
\end{description}

\smallskip\noindent\emph{Line (C11).}
The charge $e\in(0,e_2]$; then the torus, through the integer $T$; then $N$.

\smallskip
The parameter sets are
\[
\mathrm{par}^{0}:=\left(\alpha_0,\gamma_{\mathrm{w}},\kappa_{\mathrm{w}}+\kappa_0;C_{w}^{(1)}\right),\qquad\mathrm{par}:=\left(\alpha_0,\gamma_{\mathrm{w}},\kappa_{\mathrm{w}};C_{w}\right),
\]
and $\mathfrak{e}_2=0$ (Section~I.8.5).
The functions $\mathfrak{d}_{\mathrm{ma}}$ and $p^{0}_{\mathrm{ch}}$ of the line (C0r) of Section~I.8.2 belong to the Propositions~30.m.
The absolute number
\[
c_{\mathrm{cl}}=\frac{96\left(1+\sqrt{2}\right)\left(128+64\left(\frac{16}{\pi}+\frac{64}{\pi^{2}}\right)\right)}{4\pi^{2}}
\]
of the data of Theorem~44 (Section~I.8.2, line (C0)) depends on none of the constants of (a), and no condition of Proposition~14.2 contains it, that is, none of the conditions of (b), (c) and (d).

\item[(b)] \emph{The conditions on the constants of the lines (C3), (C5), (C7), (C8), (C9).}
The conditions on $e_0$ are imposed for all $v\ge-\log e_0$ (that is, for $v=v_k=-\log e_k$ with $e_k\in(0,e_0]$), and those on $e_1$ for all $v\ge-\log e_1$.
The entries of (b) are the rows of a table, one for each condition.
The row of a condition gives the constant on which it is imposed (the column ``condition on'', after the word ``On''), the inequality (the column ``inequality''), the letters that it contains (the column ``letters contained'') and places where it is used (the column ``used in'').
The conditions (2.1), (fl), (T1)--(T6) named below are those of (c), and (W3), (Z8.w) are those of (d).
\begin{description}
\item[(L1)] On $L$:
\[
L^{\frac{1}{8}-\frac{3\bar{\varepsilon}}{2}}\ge2400.
\]
Letters contained: $\bar{\varepsilon}$.
Used in: Propositions~18.1\,(c), 18.2\,(ii), (iii);
with the consequences $L^{\frac{1}{4}}\ge2$, $L^{\frac{3}{4}}\ge110>2\mathrm{e}^{4}$, $L^{3}\exp(-L)\le L^{-\frac{1}{4}}$: Propositions~18.5\,(b), (d), 48.5\,(b), (d);
Proposition~48.4\,(b);
with $L\ge4^{8}$: Proposition~24.1;
with $L^{-\frac{3}{4}}\le\frac{1}{2}$ in \cite[Lemma~5.18]{QEDI}: Propositions~22.3, 24.1, 27.6, 28.1;
with $L^{\frac{3}{4}}\ge110$ in \cite[Proposition~5.38(ii)]{QEDI}: Propositions~22.3, 24.1;
with $L^{-\frac{5}{8}}\le\frac{1}{4}$: Proposition~28.2;
with $L\ge13$ (the class $\mathrm{Cl}$) in \cite[Proposition~19.1(f)]{QEDI}, which is (G4.Sc4) of Section~G4.0.2: Propositions~20.3, 48.6, 26.6, 30.3\,(f), 32.1--32.3, 49\,(b), 35.1\,(d);
Proposition~47.2\,(a) ($L>8$).
\item[(L2)] On $L$:
\[
c_{\mathrm{fl}}\cdot L^{-\frac{3}{8}+6\bar{\varepsilon}}\le\frac{1}{16}.
\]
Letters contained: $c_{\mathrm{fl}}$, $\bar{\varepsilon}$.
Used in: Propositions~28.1, 28.2 (the contraction, measured relative to the windows).
\item[(L3)] On $L$:
\[
L\left(\kappa-6\kappa_0-4\right)-1\ge\kappa.
\]
Letters contained: $\kappa$, $\kappa_0$.
Used in: the Propositions~24.m (reblocking).
It gives $L\left(\kappa-\kappa_0-2\right)\ge\kappa$, used in Propositions~18.5\,(b), 48.5\,(b), and $L\left(\kappa-\kappa_0-1\right)\ge\kappa$, used in Propositions~13\,(d6), 34.2.
\item[(L4)] On $L$:
\[
c_{\mathrm{rb}}\cdot L^{3}\exp\left(-\left(L\left(\kappa-\kappa_0-1\right)-\kappa\right)\right)\le\frac{1}{2}.
\]
Letters contained: $c_{\mathrm{rb}}$, $\kappa$, $\kappa_0$.
Used in: Proposition~34.2 (\cite[Section~20.1.8]{QEDI}, last paragraph; \cite[Lemma~20.22]{QEDI}).
\item[(L5)] On $L$: $L$ is odd, and $L\ge5$.
Used in: Propositions~1.6, 1.7, the Propositions~2.m, Propositions~3, 4, 5, 6, 7, 8\,(b), 9, 10 ($L$ odd), 11, 12 ($L>1$), 13, 14.1, \cite[Section~15]{QEDI};
Proposition~15\,(n), the Propositions~16.m and 17.m, Propositions~18.1--18.4, 48.1--48.4;
$L$ odd: Propositions~19\,(b), 20.2 (no site lies on a face of a block; $L^{k}$, $L^{3}$ sites of a block), 18.5, 18.6\,(e), 20.3\,(c), 48.5 (an $LM$-cube is the union of $L^{3}$ $M$-cubes: \cite[Theorem~13.2]{QEDI}, assertion (t5));
the Propositions~21.m--25.m ($L$ odd, $L>1$);
the Propositions~26.m, 27.m, 28.m, Propositions~29, 30.3\,(f), 31.2, the Propositions~32.m, Proposition~49;
the Propositions~33.m ($L$ odd, $L>1$);
Propositions~34.2, 34.3\,(b) ($L>1$, $L\ge3$; with (S1): $L^{1-\varepsilon}\ge2$);
Propositions~35.1 ($L$ odd), 35.3, 35.4, 35.7, 35.9, 35.10;
Propositions~37.2\,(f), 37.3;
Propositions~41, 47.2\,(a) ($\mathrm{per}$ is odd, $\mathrm{hc}=m_{\mathrm{p}}+\frac{1}{2}$), 38.3\,(e), 40\,(e), 42.2, 43.1 ($L\ge5$), Theorem~44 ($L$ odd: a datum of \cite[Definition~2.36]{QEDII}).
\item[(M1)] On $M$: $M$ is a power of $L$, and $M\ge L^{2}$.
Letters contained: $L$.
Used in: Propositions~2.5--2.8, 3, 4, 5, 6, 7, 8\,(b), 9, 10 ($M$ odd; $M\mathbb{Z}^3\subset L\mathbb{Z}^3$), 11, 13, 14.1;
\cite[Section~15]{QEDI}, \cite[Proposition~19.1(g)]{QEDI};
the Propositions~26.m and 27.m, Propositions~28.3, 29, 31.2, the Propositions~32.m, Proposition~49;
the Propositions~16.m and 17.m, Propositions~18.1--18.4, 48.1--48.4;
Propositions~15\,(n), 48.5\,(e) (with (L5): $\eta_k\le1\le M$);
Propositions~18.5\,(f), 18.6\,(a), (b), 48.5\,(f) (with (L5): $ML^{k}$ is odd, and an $M$-cube has $\left(\frac{M}{\eta_k}\right)^{3}$ photon sites);
Proposition~20.2 ($\frac{L-1}{2}\le M$);
the Propositions~21.m--25.m ($M$ a power of $L$, $M$ odd, $L-1<M$);
the Propositions~33.m ($M$ an integer);
Propositions~34.1\,(b), 34.3\,(b);
Propositions~35.1, 35.3, 35.4, 35.5 ($M$ odd), 35.10;
Propositions~38.3\,(e), 41, 47.2\,(a), 40\,(e), 42.2, 50\,(h) ($T^{*}(e)$ is an integer $\ge1$, so that $\mathrm{per}$ is an integer), 43.1.
\item[(M2)] On $M$: $M>4L$.
Letters contained: $L$.
Used in: Propositions~2.7, 2.8\,(a), 3\,(d), 8\,(b);
the facts (g1), (g4) of \cite[Section~17.4.1]{QEDI};
as a standing condition of \cite[Section~17.4.1]{QEDI}: Propositions~7, 26.2--26.4.
\item[(M3)] On $M$: $\gamma_{\mathrm{w}}\cdot M\ge16$.
Letters contained: $\gamma_{\mathrm{w}}$.
Used in: the Propositions~26.m ($\gamma_{\mathrm{w}}M\ge1$);
\cite[Proposition~5.38]{QEDI} ($\gamma_{\mathrm{w}}>0$), and through it Propositions~29, 32.1;
the Propositions~35.m ($\gamma^{F}M\ge16$).
\item[(M4)] On $M$: $M^{\alpha_0}\ge1+\exp(\hat{\kappa})$.
Letters contained: $\alpha_0$, $\hat{\kappa}$.
Used in: the Propositions~24.m (\cite[Lemma~17.13]{QEDI}); Proposition~27.5\,(c6).
\item[(M5)] On $M$: $M\ge M^{F}$.
Letters contained: $M^{F}$.
Used in: Propositions~2.7, 2.8 ($M\ge108\pi\cdot C'_V\cdot a'$), 3\,(a); Proposition~26.5.
\item[(M6)] On $M$: $M\ge M^{P}$.
Letters contained: $M^{P}$.
Used in: \cite[Theorem~15.35]{QEDI} (in Proposition~26.5) and \cite[Proposition~15.36]{QEDI} (in Proposition~26.3\,(c)).
\item[(S1)] On $\varepsilon$: $1-\alpha<\varepsilon\le\bar{\varepsilon}$.
Letters contained: $\alpha$, $\bar{\varepsilon}$.
Used in: the Propositions~18.m ($1-\varepsilon<\alpha<1$); Proposition~13\,(d6) ($\varepsilon\le\frac{1}{8}$); Proposition~48.5\,(b), (d);
with $\varepsilon>0$: Propositions~2.7, 20.1, 20.3\,(a), 17.1\,(c2), 17.2\,(d), 48.3\,(d), 35.3\,(f);
with $\varepsilon\le\frac{1}{200}$ alone: Proposition~18.1;
with $0<\varepsilon\le\frac{1}{200}$: Propositions~9\,(d), 17.1\,(a), 18.2, 48.1\,(c), 48.4\,(b), \cite[Proposition~5.38(ii)]{QEDI}, Propositions~22.3, 23.2\,(d), 24.1, 24.3\,(q1), 27.5, the Propositions~28.m, Proposition~29, the Propositions~32.m, Propositions~34.1--34.4, 35.4\,(SF4), (DF5), (SF5), 35.5, 35.9, 35.10.
\item[(S2)] On $s$:
\[
\sigma=s\left(\frac{1}{4}-9\varepsilon\right)-12\varepsilon-\frac{2a_L}{N_{*}}>0.
\]
Letters contained: $\varepsilon$, $a_L$, $N_{*}$.
Used in: Proposition~35.9\,(b) ($\mathfrak{e}>0$); the Propositions~43.m, Theorem~44.
\item[(A1)] On $e_0$:
\begin{gather*}
v^{p_0}\le\exp\left(\frac{\varepsilon v}{4}\right),\qquad12\exp\left(-\frac{v^{2p_0}}{2}\right)\le\exp(-7v),\qquad\exp\left(-\frac{5\varepsilon v}{2}\right)\le\frac{1}{2},\\
c_{\mathrm{fl}}\exp(-\varepsilon v)\le1,\qquad3M\exp(-\varepsilon v)\le1.
\end{gather*}
Letters contained: $p_0$, $\varepsilon$, $c_{\mathrm{fl}}$, $M$.
Used in: the Propositions~24.m; \cite[Proposition~5.38]{QEDI}, (I.Lc2)$^{\mathrm{ch}}$;
the third inequality: Proposition~27.5\,(c2);
the first inequality, through \cite[Proposition~5.38]{QEDI}: Propositions~22.3, 24.1, 29, 32.1, 35.4\,(SF5).
\item[(A2)] On $e_0$:
\[
2\,C_{w}^{2}\cdot c_{\gamma}^{2}\cdot c_C\exp\left(-\frac{\varepsilon v}{4}\right)\le1.
\]
Letters contained: $C_{w}$, $c_{\gamma}$, $c_C$, $\varepsilon$.
Used in: (q1) of the Propositions~24.m; \cite[Proposition~5.38]{QEDI}, and through it Propositions~22.3, 24.1, 29, 32.1, 35.4\,(SF5).
\item[(A3)] On $e_0$:
\[
16\,c_X\cdot c_{\gamma}\cdot C_{w}\cdot C_T^{2}\cdot M^{3}\exp\left(-\frac{\varepsilon v}{2}\right)\le1,\qquad96\,M^{3}\exp\left(-\frac{v}{4}\right)\le1.
\]
Letters contained: $c_X$, $c_{\gamma}$, $C_{w}$, $C_T$, $M$, $\varepsilon$.
Used in: the Propositions~24.m (the first part for $k\ge1$; the second part for $k=0$, together with the Propositions~23.m).
\item[(A4)] On $e_0$:
\[
32\,C_{w}\cdot C_T\exp\left(-\frac{3\varepsilon v}{2}\right)\le1.
\]
Letters contained: $C_{w}$, $C_T$, $\varepsilon$.
Used in: the Propositions~24.m.
\item[(A5)] On $e_0$:
\[
88\,M^{3}\exp\left(-\left(\frac{1}{4}-6\varepsilon\right)v\right)\le1.
\]
Letters contained: $M$, $\varepsilon$.
Used in: the Propositions~24.m; Proposition~27.5\,(c2).
\item[(A6)] On $e_0$:
\[
112\,C_{w}^{2}\cdot c_{\gamma}\exp(-2\varepsilon v)\le3.
\]
Letters contained: $C_{w}$, $c_{\gamma}$, $\varepsilon$.
Used in: the Propositions~24.m, together with Proposition~15 ($\mathsf{w}_W\le7$).
\item[(A7)] On $e_0$:
\begin{gather*}
c_S\exp\left(-\left(\frac{1}{4}-5\varepsilon\right)v\right)\le1,\qquad c_S\cdot L^{3}\exp(-\varepsilon v)\le1,\\
36\,c_X^{2}\cdot K_0^{2}\exp\left(-\left(\frac{1}{4}-5\varepsilon\right)v\right)\le1.
\end{gather*}
Letters contained: $c_S$, $L$, $c_X$, $K_0$, $\varepsilon$.
Used in: the Propositions~24.m.
\item[(A8)] On $e_0$:
\begin{gather*}
\exp(-\varepsilon v)\le\frac{1}{4},\qquad3M\left(\exp\left(-\left(\frac{1}{4}+\varepsilon\right)v\right)+\exp(-2\varepsilon v)\right)\le1,\\
3M\exp(-\varepsilon v)\le1,\qquad\exp\left(-\frac{v}{4}\right)+\exp(-\varepsilon v)\le1.
\end{gather*}
Letters contained: $M$, $\varepsilon$.
Used in: \cite[Lemmas~5.17 and 5.18]{QEDI}, (I.Lc2)$^{\mathrm{ch}}$, \cite[Proposition~5.38]{QEDI} (the first inequality: $e_k<1$);
the Propositions~17.m, Proposition~18.1\,(c), the Propositions~20.m, Propositions~27.5, 48.1\,(c);
the first inequality: Propositions~22.3, 24.1 ($e_0<1$, $e_0\le4^{-200}$), 28.1, 28.2, 34.1--34.4, 35.5;
the first, third and fourth inequalities: Proposition~35.4;
through \cite[Lemma~5.17 and Proposition~5.38]{QEDI}: Proposition~22.3, the Propositions~24.m, Propositions~29, 32.1.
\item[(A9)] On $e_0$:
\begin{gather*}
160\,L^{1+\frac{3\varepsilon}{2}}\cdot M\exp(-\varepsilon v)\le1,\qquad90\,c_X\cdot L\cdot M\exp(-\varepsilon v)\le1,\\
9M\cdot L^{\varepsilon}\exp(-2\varepsilon v)\le1,\qquad v^{2}\ge10L.
\end{gather*}
Letters contained: $L$, $M$, $c_X$, $\varepsilon$.
Used in: the Propositions~17.m and 18.m, Propositions~48.3\,(d), 48.4\,(b) (the first three inequalities);
the fourth inequality follows from (A8) and (M4), by (f).
\item[(A10)] On $e_0$:
\begin{gather*}
v\ge9,\qquad8M\left(\frac{v}{4}+4\right)\exp(-\varepsilon v)\le1,\\
8\,C_B\cdot v^{p}\le\exp\left(\left(\frac{3}{4}-3\varepsilon\right)v\right),\qquad8\exp\left(-\left(\frac{3}{4}-\frac{7\varepsilon}{2}\right)v\right)\le1.
\end{gather*}
Letters contained: $M$, $\varepsilon$, $C_B$, $p$.
Used in: Proposition~13\,(d4) (the first two inequalities); \cite[Proposition~5.38]{QEDI} (the last two), and through it Propositions~22.3, 24.1, 29, 32.1, 35.4\,(SF5).
\item[(Z2.s1)] On $e_0$:
\[
6L\exp(-\varepsilon v)\le1.
\]
Letters contained: $L$, $\varepsilon$.
Used in: Propositions~2.7, 2.8, 3, 5, 6, 7, 8, 13.
\item[(Z2.s2)] On $e_0$:
\[
4\,C'_V\cdot a'\exp(-\varepsilon v)\le1.
\]
Letters contained: $C'_V$, $a'$, $\varepsilon$.
Used in: Propositions~2.7, 2.8, 3, 8.
\item[(Z2.s3)] On $e_0$:
\[
3M\exp(-\varepsilon v)\le1.
\]
Letters contained: $M$, $\varepsilon$.
Used in: Proposition~6\,(b)--(f).
\item[(Z2.s4)] On $e_0$:
\[
\exp(-\varepsilon v)\le\frac{1}{4},\qquad\exp\left(-\frac{v}{4}\right)+\exp(-\varepsilon v)\le1,\qquad4\exp\left(-\left(\frac{1}{4}+2\varepsilon\right)v\right)+2\exp(-3\varepsilon v)\le1.
\]
Letters contained: $\varepsilon$.
Used in: Propositions~9\,(d), 13\,(d4)--(d6); the assertions (G3) of \cite[Section~16.2.7]{QEDI}.
\item[(Z2.s5)] On $e_0$:
\[
8\,c_{\mathrm{rb}}\cdot L^{3}\cdot C_d\cdot M^{3}\exp(-2\varepsilon v)\le1.
\]
Letters contained: $c_{\mathrm{rb}}$, $L$, $C_d$, $M$, $\varepsilon$.
Used in: Proposition~13\,(d6).
\item[(Z3.s1)] On $e_0$: $v\ge1$.
Used in: \cite[Section~15]{QEDI}.
\item[(e1)] On $e_1$:
\[
v\ge8,\qquad v^{r}>31L+16,\qquad\left(1+\frac{3\,\mathrm{lg}}{v}\right)^{r}<L,\qquad p_0\cdot v^{p_0-1}\ge28.
\]
Letters contained: $r$, $L$, $p_0$.
Used in: (2.1), (fl), (T1); \cite[Propositions~19.1 and 19.2]{QEDI}; Propositions~35.1, 49\,(b);
through (G4.Sc1)--(G4.Sc6) of Section~G4.0.2: Propositions~20.3, 48.6, 26.6, 28.2, 28.3, 30.3\,(f), 31.2, the Propositions~32.m, Proposition~49\,(b).
\item[(e2)] On $e_1$:
\[
\left(c_{\mathrm{fl}}+2\right)\exp(-\varepsilon v)\le\frac{1}{16}\left(1-L^{-\frac{\varepsilon}{2}}\right).
\]
Letters contained: $c_{\mathrm{fl}}$, $\varepsilon$, $L$.
Used in: (fl).
\item[(e3)] On $e_1$:
\begin{gather*}
8\,C_0^{2}\cdot\mathrm{lh}(v)^{6}\left(b_{-}^{-1}+L\cdot\mathrm{lh}(v)\right)\le\exp\left(\frac{\varepsilon v}{2}\right),\qquad4\left(2C_0+c_{\kappa}+5c_{\mathrm{fl}}+6\right)\mathrm{lh}(v)^{3}\le\exp(\varepsilon v),\\
\exp\left(-\left(\frac{1}{4}-9\varepsilon\right)v\right)\le\frac{1}{8}.
\end{gather*}
Letters contained: $C_0$, $L$, $M$, $r$, $b_{-}$, $c_{\kappa}$, $c_{\mathrm{fl}}$, $\varepsilon$.
Used in: (T2); the third inequality also at $v=v_K$: Proposition~32.3\,(10.3), (G6.T-14)\,(vi).
\item[(e4)] On $e_1$:
\[
4\,C_0\cdot\mathrm{lh}(v)^{2}\cdot v^{p}<\exp(\varepsilon v).
\]
Letters contained: $C_0$, $L$, $M$, $r$, $p$, $\varepsilon$.
Used in: (T3); Proposition~35.1\,(e).
\item[(e5)] On $e_1$:
\[
12\,\mathrm{lh}(v)^{3}\left(\exp\left(-\frac{v^{p_0}}{2}\right)+\exp\left(-\frac{v^{2p}}{2C_0^{2}}\right)\right)\le\exp(-v).
\]
Letters contained: $L$, $M$, $r$, $p_0$, $p$, $C_0$.
Used in: (T4).
\item[(e6)] On $e_1$:
\[
C_{\Sigma}\exp\left(-\gamma_{\Sigma}\cdot M\cdot v^{r}\right)\le\exp(-v).
\]
Letters contained: $C_{\Sigma}$, $\gamma_{\Sigma}$, $M$, $r$.
Used in: (T5).
\item[(e7)] On $e_1$:
\[
L^{2}\exp(-2v)\cdot\mathrm{lh}(v)\le1.
\]
Letters contained: $L$, $M$, $r$.
Used in: (T6).
\item[(e8)] On $e_1$:
\[
6N_{*}\cdot M\cdot v^{r}\exp(-v)\left(v\exp(-v)+\frac{1}{2}\right)\le\pi^{2}.
\]
Letters contained: $N_{*}$, $M$, $r$.
Used in: (W3), together with (Z8.w).
\end{description}
\end{itemize}
\end{named}

\begin{named}{Proposition 14.2}[The ordered list, continued]
Apart from those ranges printed in (a) that (b) does not repeat, namely
\begin{gather*}
\rho\in(0,1),\qquad\alpha\in\left(1-\bar{\varepsilon},1\right),\qquad\kappa\ge7\kappa_{0}+7,\\
\text{integers}\ r\ge2,\ p_0>3r+2,\ p>2p_0,\qquad s\in(0,1),\qquad e_0>0,\\
e'\in\left(0,\min\left(e_0,\frac{1}{3}\right)\right]\qquad\text{(the variable $e'$ of the entry $e_1$)},
\end{gather*}
the conditions on the constants of the lines (C0)--(C9) are the conditions of (b).
Each condition on the constants that the proof of a proposition adds is of the type ``$e_0$ (respectively $e_1$, $e_{\sharp}$, $e_2$) small depending on earlier constants'' and stands in this list (Section~G1.Z1, (G1.T-12));
thus the list contains every condition on the constants that the proofs require.
The conditions (A0), (Z2.m), (Z8.s1), (Z9.s1)--(Z9.s3) and (G4.2.1) of Section~I.8.4 are expressed through those of this proposition as follows.
The condition (A0) consists of (L3), the definition of $\kappa_{\mathrm{w}}$, (M1), (M2), the inequality $\gamma_{\mathrm{w}}\cdot M\ge1$ (which is implied by (M3)) and (M4).
The condition (Z2.m) is (M5).
The condition (Z8.s1) is the definition of $e_1$, that is, the entry $e_1$ of (a);
the range $e'\in\left(0,\min\left(e_0,\frac{1}{3}\right)\right]$ of the variable $e'$ is part of this definition and restricts none of the constants that are chosen,
and it gives $e_1\le e_0$ and $e_1\le\frac{1}{3}$, as recorded in (f).
The conditions (Z9.s1), (Z9.s2), (Z9.s3) are the second inequality of (g1) and the first inequalities of (g3) and of (g4) in (d).
The condition (G4.2.1) is the condition (2.1) of (c).
\begin{itemize}
\item[(c)] \emph{Conditions on $(e,T,N)$.}
If the set in the definition of $e_1$ (entry $e_1$ of (a)) is not empty, then (e1)--(e8) hold for all $v\ge-\log e_1$.
The conditions of this clause are formed with the quantities $k_{*}$, $\ell_*$, $e_{*}$, $\mu_*$, $v_{*}=-\log e_{*}$, $r_{*}$, $p_{*}=v_{*}^{p}$, $p_{0*}=v_{*}^{p_0}$ of Section~I.6 and of \cite[Section~5.2.2 and Definition~5.6]{QEDI}, and with
\[
w^{\mathrm{m}}_0=q^{\frac{3}{4}-8\varepsilon},\qquad q=e\cdot\eta^{\frac{1}{2}}\qquad\text{(Section~I.5.4)}.
\]
They are not conditions on constants.
They hold for every $\bar{m}\in(0,1]$, every $e>0$ with $e\cdot h^{\frac{1}{2}}\le e_1$, every integer $T\ge0$ with $R=h\cdot L^{T}\in\mathrm{Tor}(e)$ and every integer $N\ge N_{\min}(R,e)$; this follows from (e1)--(e7) and \cite[Proposition~19.1(c), (e)]{QEDI}.
The conditions are
\begin{gather*}
r_{*}>31L+16;\tag{2.1}\\
\sum_{k\le k_{*}}\mathrm{sm}_k\le\frac{1}{16},\quad\text{where}\quad\mathrm{sm}_k:=\left(c_{\mathrm{fl}}+2\right)\cdot e_k^{\varepsilon};\qquad v_{*}\ge8;\tag{fl}\\
v_{*}\ge8,\qquad p_0\cdot v_{*}^{p_0-1}\ge28,\qquad w^{\mathrm{m}}_0<1;\tag{T1}\\
\begin{gathered}
8\,C_{0}^2\cdot\ell_*^6\cdot\left(b_{-}^{-1}+\mu_*^{-1}\right)\le e_{*}^{-\frac{\varepsilon}{2}},\qquad
4\left(2C_{0}+c_{\kappa}+5c_{\mathrm{fl}}+6\right)\cdot\ell_*^3\le e_{*}^{-\varepsilon},\\
e_{*}^{\frac{1}{4}-9\varepsilon}\le\frac{1}{8};
\end{gathered}\tag{T2}\\
4\,C_{0}\cdot\ell_*^2\cdot p_{*}<e_{*}^{-\varepsilon};\tag{T3}\\
12\,\ell_*^3\cdot\left(\exp\left(-\frac{p_{0*}}{2}\right)+\exp\left(-\frac{p_{*}^2}{2C_{0}^2}\right)\right)\le e_{*};\tag{T4}\\
C_{\Sigma}\cdot\exp\left(-\gamma_{\Sigma}\cdot\ell_*\right)\le e_{*};\tag{T5}\\
L^2\cdot e^2\cdot R\le1.\tag{T6}
\end{gather*}
The error terms are
\begin{gather*}
\varepsilon_1:=\frac{1}{2}\,e_{*}^{\frac{1}{4}-9\varepsilon},\qquad
\varepsilon_2:=C_{\Sigma}\cdot\ell_*^{-1}\cdot\exp\left(-\gamma_{\Sigma}\cdot\ell_*\right),\\
\varepsilon_3:=6\,\ell_*^3\cdot\exp\left(-\frac{p_{0*}}{2}\right)+6\,\ell_*^3\cdot\exp\left(-\frac{p_{*}^2}{2C_{0}^2}\right),\qquad
\varepsilon_4:=6\,\ell_*^3\cdot\exp\left(-\frac{p_{*}^2}{2C_{0}^2}\right),
\end{gather*}
together with the function $\varepsilon_5(x)$ of the line (C0r).
\item[(d)] \emph{Conditions of the line (C10).}
These conditions contain $\bar{m}$ or the probes.
The statements that need them are Propositions~34.1, 34.3, 34.4, Propositions~35.5--35.10, Propositions~42.2, 43.1, 43.2 and Theorem~44.
For Propositions~34.1, 34.3, 34.4 and~35.5--35.9 the conditions of this clause that enter are named in the paragraph on the conditions on the constants that follows the statement, and Proposition~35.10 names (g1)--(g9) in its statement.
Propositions~42.2 and~43.1 contain conditions of this clause through the bounds $e\le e_G$ and $e\le\min\left(e_G,e_{\sharp}(\beta)\right)$ on the charge, respectively; the condition (W3) is named in the Remark that follows Proposition~42.2 and in Proposition~43.1\,(R1).
For Proposition~43.2 and Theorem~44 the conditions (Pb), (Z10.s1), (Z10.s2), (Z10.s4) are written out in the data and the constants of Proposition~43.2, and Section~G6.Z4 names them by their labels.
\begin{gather*}
6\,N_{*}\cdot e^2\cdot R^{*}\cdot\log\frac{1}{e}\le\pi^2;\tag{W3}\\
e\cdot h^{\frac{1}{2}}\cdot\log h\le1.\tag{Z8.w}
\end{gather*}
Together with $e\cdot h^{\frac{1}{2}}\le e_1$ and (e8), the condition (Z8.w) gives (W3); moreover (Z8.w) is contained in $e\le e_G$.

In (d) the index $*$ means ``at $k=k_{*}\left(R^{*}(e)\right)$, for $N\ge N_0(e)$'', in the setting (ST*)$^{\mathrm{ch}}$ of Section~G5.1.1, with $e>0$ and $e\cdot h^{\frac{1}{2}}\le e_1$.
The numbers $a_{*}$, $e_{*}$, $\ell_*$ do not depend on $N\ge N_0(e)$, by \cite[Proposition~19.1(c)]{QEDI}, which is (G4.Sc1);
the numbers $r_{*}=\ell_*/M$, $v_{*}=-\log e_{*}$, $p_{*}=v_{*}^{p}$ and $p_{0*}=v_{*}^{p_0}$ are functions of these and of constants, so that they do not depend on $N\ge N_0(e)$ either (compare Proposition~35.1\,(c)).
Consequently (g1)--(g9) are conditions on $e$ alone: of the data $(e,T,N)$ they contain only $e$.
With the exception of $\mathrm{mxL}$, the quantities defined in the following display depend on $e$;
the number $\mathrm{mxL}$, which enters $\lambda_L$, is a function of the constant $\mathbf{D}$ of the line (C10) alone, and it does not depend on $e$.
The quantities are written in the letters of (a), with $h$ in place of $h'$ and with $C^{F}$, $\gamma^{F}$ in place of $C_F$, $\gamma_F$;
the quantities $\tau_s$, $\bar{\kappa}_1$, $\bar{r}$, $B_1$, $B_{\mathrm{wr}}$, $B_0$, $\varpi_{\mathrm{sep}}$, $\varpi_0$, $t_2$, $\varpi_L$, $\mathrm{mxL}$, $\lambda_L$ are those of \cite[Section~20.1.7]{QEDI}, and $\mathrm{pp}$, $t_x$, $\mathbf{r}$ are those of \cite[Section~20.1.4]{QEDI}:
\begin{align*}
\tau_s&:=c_{\tau}\cdot\ell_*,\qquad\bar{\kappa}_1:=6\,e_{*}^{\frac{1}{4}-9\varepsilon},\qquad
\bar{r}:=4\,h_{G}^2\cdot C_K\cdot e_{*}^{2\varepsilon}+650\,K_{0}\cdot e_{*}^{\frac{1}{4}-7\varepsilon},\\
B_1&:=16\,\mathrm{gGa}_{1}\cdot C_K'\cdot\mathrm{pp}\cdot e_{*}^{2\varepsilon},\\
B_{\mathrm{wr}}&:=16\exp(6)\cdot C_{d}\cdot M^3\cdot K_{0}\cdot\mathrm{pp}\cdot\exp\left(-r_{*}\right),\\
B_0&:=C_{\mathrm{fr}}\cdot e_{*}^{4\varepsilon}+B_{\mathrm{wr}}+B_1,\\
\varpi_{\mathrm{sep}}&:=2^{N_{*}}\cdot\mathrm{CWj}\cdot\mathrm{pp}\cdot\exp\left(-\mathrm{vt}\cdot\left(d_{\mathrm{sep}}-4M-6\right)\right),\\
\varpi_{\mathrm{sep}}^{+}&:=2^{N_{*}}\cdot\mathrm{CWj}\cdot\mathrm{pp}\cdot\exp\left(-\mathrm{vt}\cdot\left(\frac{M\cdot r_{*}}{8L}-4M-7\right)\right),\\
\varpi_0&:=18\,\ell_*^3\cdot\exp\left(-\frac{p_{0*}}{2}\right)+6\,\ell_*^3\cdot\exp\left(-\frac{\left(p_{*}-1\right)^2}{2C_{0}^2}\right),\\
t_2&:=e\cdot\min\left\{N_{*}^{-\frac{1}{2}},\ \frac{1}{C_{\beta}},\ \frac{1}{2\sqrt{N_{*}\cdot\mathbf{D}}}\right\},\\
t_x&:=\frac{e_{*}^{3\varepsilon}}{\mathrm{Cb}},\qquad\mathbf{r}:=\frac{t_x}{4},\\
\varpi_L&:=\varpi_0^{\frac{1}{2}}\cdot\exp\left(\frac{\mathbf{D}}{2}\right),\qquad\mathrm{mxL}:=2+\exp\left(\frac{\mathbf{D}}{2}\right),\\
\lambda_L&:=\left(\varpi_L^{\frac{1}{3}}\cdot\mathrm{mxL}^{\frac{2}{3}}\right)^{\frac{1}{N_{*}}}\cdot5\cdot\left(\frac{2N_{*}\cdot\mathbf{r}}{t_2}\right)^4,\\
\mathrm{pp}&:=\left(\frac{D_\beta}{a_{*}}+2M+4\right)^3.
\end{align*}
Here $d_{\mathrm{sep}}$ is the number of Section~G5.1.4; through it $\varpi_{\mathrm{sep}}$ depends on the placement and on $N$.
The number $\varpi_{\mathrm{sep}}^{+}$ is not a quantity of \cite{QEDI}: it arises from $\varpi_{\mathrm{sep}}$ when $d_{\mathrm{sep}}$ is replaced by its lower bound $\frac{M\cdot r_{*}}{8L}-1$ of Proposition~35.3\,(e), since
\[
\left(\frac{M\cdot r_{*}}{8L}-1\right)-4M-6=\frac{M\cdot r_{*}}{8L}-4M-7;
\]
thus $\varpi_{\mathrm{sep}}^{+}$ is the right side of the bound on $\varpi_{\mathrm{sep}}$ in \cite[Theorem~20.29]{QEDI}.
In the minimum that defines $t_2$, an entry whose denominator is $0$ is read as $+\infty$.
The number $\mathrm{pp}$ is the one of Section~G5.1.4, formed on $R=R^{*}(e)$, so that
\[
\mathrm{pp}\le\left(\frac{D_\beta\cdot L^2}{h}+2M+4\right)^3.
\]
The conditions on $e$ are the following:
\begin{gather*}
5\,\mathrm{hL}\cdot h_{G}\le e_{*}^{-\frac{1}{4}+\varepsilon},\qquad\left(4+\frac{M^3}{2}\right)\cdot e_{*}^{\varepsilon}\le1,\qquad\mathrm{alG}\cdot\bar{r}\le\frac{1}{2};\tag{g1}\\
\frac{M\cdot r_{*}}{8L}-1\ge6M+8;\tag{g2}\\
\begin{gathered}
4\,C_{0}\cdot\ell_*^3\cdot\mathrm{gGa}_{1}\le e_{*}^{-\frac{1}{4}+\varepsilon},\qquad4\,\ell_*^3\cdot\left(2C_K+K_{0}+4\right)\le e_{*}^{-2\varepsilon},\\
e_{*}^{\frac{1}{4}-9\varepsilon}\le\frac{1}{8};
\end{gathered}\tag{g3}\\
\begin{gathered}
48\,e_{*}^{2\varepsilon}\le\mu_-,\qquad2\,\mathrm{gGa}_{1}\cdot C_K'\cdot e_{*}^{2\varepsilon}\le\frac{1}{2},\qquad\tau_s\ge1,\\
52\,B_1\cdot\exp\left(-\tau_s\right)\le e_{*}^{\frac{1}{4}-9\varepsilon},\qquad B_0\le\frac{1}{4};
\end{gathered}\tag{g4}\\
2\,\varpi_{\mathrm{sep}}^{+}\le\bar{\kappa}_1^{N_{*}};\tag{g5}\\
N_{*}\cdot\bar{\kappa}_1^{1-s}\le\frac{1}{2};\tag{g6}\\
\varpi_0\le\frac{1}{4},\qquad8\,N_{*}\cdot t_2\le t_x;\tag{g7}\\
\lambda_L\le\bar{\kappa}_1^{s};\tag{g8}\\
\frac{D_\beta}{2}+\frac{1}{2}\le\frac{3}{8}\,\ell_c,\qquad D_\beta+3\le\ell_c,\qquad\text{where}\quad\ell_c=\frac{h\cdot M\cdot r_K(e)}{2L^2}.\tag{g9}
\end{gather*}
The second inequality of (g1) is (Z9.s1), the first inequality of (g3) is (Z9.s2), and the first inequality of (g4) is (Z9.s3).
Concerning (g2): by Proposition~35.3\,(e) one has $d_{\mathrm{sep}}\ge\frac{M\cdot r_{*}}{8L}-1$ for every admitted placement, and that the placement is admitted is a datum.

The condition (Pb) restricts the choice of $(\beta,x_0)$, which belong to the data of Proposition~43.2.
Let $f=dd^{*}\beta$ and $a=d^{*}\beta$, put $\|a\|_{\infty}:=\sup_{y}|a(y)|_2$, and let $r_\beta>0$ be a number such that $\operatorname{supp}\beta$ lies in a ball of radius $r_\beta$.
The number $r_\beta$ is a datum chosen together with $(\beta,x_0)$ (data of Proposition~43.2); unlike $\beta$ and $x_0$, it is not a letter of the list (a), and no constant depends on it.
The condition reads
\begin{equation*}
\begin{gathered}
b(f)>0,\qquad d_0:=|x_0|_2-2r_\beta>0,\\
\frac{6}{\bar{m}^2}\cdot\|a\|_{\infty}^4\cdot\left(\frac{4\pi\cdot r_\beta^3}{3}\right)^2\cdot\left(\bar{m}^2+\left(\bar{m}+\frac{1}{d_0}\right)^2\right)\cdot\frac{\exp\left(-2\bar{m}\cdot d_0\right)}{4\pi^2\cdot d_0^2}\le\frac{b(f)}{12};
\end{gathered}\tag{Pb}
\end{equation*}
its last inequality is the inequality of \cite[Section~2.8]{QEDI}.
Finally, the number $e_2$ is subject to
\begin{gather*}
e_2\le e_G;\tag{Z10.s1}\\
e_2\le e_{\sharp}(\beta')\quad\text{for the three probes}\quad\beta'=\beta,\ \beta+\tau_{x_0}\beta,\ \beta-\tau_{x_0}\beta;\tag{Z10.s2}\\
\frac{5}{2}\,C'_{\bar{\kappa}}\cdot e_2^{\sigma}\le\frac{b(f)}{4},\qquad C'_{\bar{\kappa}}=8\,C_{\bar{\kappa}}+3\,C_{\bar{\kappa}}^2.\tag{Z10.s4}
\end{gather*}
The condition (Z10.s4) implies
\[
\frac{5}{2}\,C_{\bar{\kappa}}\cdot e_2^{\sigma}\le\frac{b(f)}{4}\qquad\text{and}\qquad C'_{\bar{\kappa}}\cdot e_2^{\sigma}\le\frac{b(f)}{10}.
\]
\item[(e)] \emph{Order.}
In the list (a) every constant is defined through constants that stand before it, and every letter named after ``depends on'' in its entry stands before it.
Every condition of (b) on a constant contains, besides this constant and the variable $v$, only constants that stand before it.
The conditions of (c) and (d) contain only constants of (a) and the data $(e,T,N)$; besides these, (Pb) contains the radius $r_\beta$, a datum of Proposition~43.2 that is chosen together with $(\beta,x_0)$.
None of these conditions contains a placement.
\item[(f)] \emph{Consequences.}
From (L1) one obtains $L\ge2400^8$, and hence
\begin{gather*}
L\ge5,\qquad L>8,\qquad L\ge4^8,\qquad L^{\frac{1}{4}}\ge2,\qquad L^{\frac{3}{4}}\ge110,\\
L^{-\frac{5}{8}}\le\frac{1}{4},\qquad L^{-\frac{3}{4}}\le\frac{1}{2},\qquad L^{-3+\frac{1}{8}}\le\frac{1}{2},\qquad L^3\cdot\exp(-L)\le L^{-\frac{1}{4}},\\
L^{\frac{1}{4}-\frac{3\varepsilon}{2}}\ge13.6,\qquad L^{-\frac{7}{4}}+4\,L^{-\frac{7}{8}+\frac{\varepsilon}{2}}<1\qquad\left(\varepsilon\le\bar{\varepsilon}\right).
\end{gather*}
From (L3) and (C2) one obtains
\[
L\cdot\left(\kappa-\kappa_{0}-1\right)\ge\kappa,\qquad\kappa\ge\kappa_{0},\qquad\kappa_{\mathrm{w}}-1\ge\kappa+1.
\]
Moreover the following hold:
\begin{gather*}
\gamma^{F}\cdot L\le\frac{1}{6}<\log2,\qquad\gamma_{\mathrm{w}}\le\gamma^{F},\qquad\gamma_{\mathrm{w}}\le\gamma^{P}\le\gamma_{P}\le1,\\
\gamma_{\mathrm{w}}\cdot M\ge1,\qquad\gamma^{F}\cdot M\ge16,\qquad27\,\alpha_{0}=\frac{1}{2};\\
a'\ge1,\qquad C_{5}\ge1,\qquad C'_{V}\ge24\,\mathrm{e},\qquad M\ge108\pi\cdot C'_{V}\cdot a'\ge1600,\qquad M\ge9;\\
\mathrm{qq}\left(M^{27\alpha_{0}}\right)\le\frac{1}{2},\qquad\mathrm{qq}\left(\exp\left(27\kappa'\right)\right)\le\frac{1}{2};\\
2\,C_{1}\le C^{F},\qquad4\le C^{F}\le C_{w}^{(0)}\le C_{w}^{(1)}\le C_{w},\qquad C_{P}\le C^{P}\le C_{w},\\
4\exp\left(2\gamma_{\mathrm{w}}\right)\cdot M^{27\alpha_{0}}\le C_{w}^{(1)},\qquad C_{0}\ge\max\left(1,C^{F}\right),\qquad c_{\kappa}\le K_{0},\\
C_{T}\ge1,\qquad c_{\gamma}\ge1,\qquad\kappa_{0}\ge c_{X}\quad\left(54\log72\ge54\right);\\
0<\alpha<1,\qquad1-\varepsilon<\alpha,\qquad\varepsilon\le\frac{1}{200};\qquad p_0\ge2,\qquad2p>p_0.
\end{gather*}
Here $\mathrm{qq}$ is the function of Section~G1.10; the two inequalities on $\mathrm{qq}$ are those that enter Proposition~2.7\,(e).
Since $\exp(-\varepsilon\cdot v)\le\frac{1}{4}$ and $\varepsilon\le\frac{1}{200}$, one has $e_0\le4^{-200}$.
Hence $v\ge277$ for every $v=-\log e_k$ with $e_k\in(0,e_0]$, so that the inequality $v\ge9$ in (A10), the condition (Z3.s1) and the inequality $p_{0,k}\ge1$ hold.
Moreover $e_1\le e_0$ and $e_1\le\frac{1}{3}$.

\emph{Ranges.}
Let every constant be chosen as (a) and (b) prescribe.
Then every constant of the entries $\rho$, \dots, $\sigma$ of (a), apart from $C_{E}$, $\mathrm{Csy}_{\mathrm{ak}}$ and $\mathrm{Cpv}_0$, is defined and finite.
(These three are given by Propositions~37.1 and~36; they occur in no entry up to $e_1$ other than their own entries, and in no condition of (b).)
All these constants are positive, except that $A_{\mathrm{ak}}$, $A'_{\mathrm{ak}}$, $C_{G}$, $C_{6}$, $C_{R}$, $C_{\Sigma}$ and the values of $\varepsilon_5$ are non-negative.
Furthermore
\begin{gather*}
\lambda_{\mathrm{int}}\le2\left(w_0+w_1+w_2\right),\qquad c_{*}\le\frac{1}{3},\qquad c_{\mathrm{el}}\le0.22,\qquad K_{\mathrm{Str}}\ge1,\qquad\kappa_{(0)}\le\frac{\theta_{0}}{2},\\
a_{L}>1,\qquad a_{0}\le\frac{1}{3},\qquad\gamma_{(5)},\ a_{z},\ \gamma_{z}\le\frac{1}{3L},\qquad c_{\mathrm{re}}\le b_{-}<1,\qquad C_{00}\ge1,\\
\gamma_{P}\le1,\qquad C_{P}\ge4\,\mathrm{e}^2,\qquad\gamma_{\Sigma}\le1,\qquad c_{\gamma}\le2L^3\cdot\exp\left(\left(2L+1\right)\cdot\gamma_{\mathrm{w}}\right)\cdot c\left(\gamma_{\mathrm{w}}\right),\\
c_{\kappa}\ge1,\qquad p_0\ge9,\qquad p\ge19;\\
K_{0}\ge1,\qquad c_{\mathrm{fl}}\ge1,\qquad M\ \text{odd},\qquad L-1<M,\qquad12\left(1+\rho\right)\le24,\qquad b_{k}\ge b_{-}\ \text{for}\ k\ge1,\\
C^{F}\ge1,\qquad C_{w}\ge1,\qquad0<\gamma^{F}\le1,\qquad\kappa-1\ge\kappa_{0},\qquad\kappa-\kappa_{0}\ge2,\\
e_0\le1,\qquad6\varepsilon<1,\qquad0<\sigma<\mathfrak{e}\le\frac{1}{4},\qquad\mathrm{vt}\le\min\left(\frac{\gamma^{F}}{8},\frac{1}{M}\right),\qquad h\ge1.
\end{gather*}
If $e\cdot h^{\frac{1}{2}}\le e_1$, then
\[
e\le\frac{1}{3},\qquad v_K\ge8,\qquad\left(1+\frac{3\,\mathrm{lg}}{v_K}\right)^{r}<L,\qquad3\left(e\cdot h^{\frac{1}{2}}\right)^{\frac{1}{4}-9\varepsilon}\le1;
\]
if moreover $R=R^{*}(e)$ and $N\ge N_0(e)$, then $\mu_-<\mu_*$.
The inequalities between the constants of the line (C6$'$) that are printed in Annex~B, Section~G5.9.3, hold; the last of them, which contains $\mu_*$, holds under the hypotheses of the preceding sentence.

\emph{Conditions that are imposed and implied by others.}
The condition (L1) is implied by (L2), and the inequality $L\ge5$ in (L5) by (L1).
The inequality $\gamma_{\mathrm{w}}\cdot M\ge1$ is implied by (M3), and (M2) by (M5).
The condition (Z2.s3) coincides with the third inequality of (A8) and with the fifth inequality of (A1).
The first and the second inequality of (Z2.s4) coincide with the first and the fourth inequality of (A8).
The second inequality of (A8) is implied by its first and its third inequality.
The condition (Z3.s1) and the first inequality of (A10) are implied by the first inequality of (A8).
The fourth inequality of (A9) is implied by the third inequality of (A8) together with (M4).
Indeed, the third inequality of (A8) gives $\varepsilon\cdot v\ge\log(3M)$, hence $v\ge200\log M$; and (M4) gives $\log M\ge54\,\hat{\kappa}\ge54\left(2L+1\right)^3\cdot\kappa$.
Therefore
\[
v\ge10800\left(2L+1\right)^3\cdot\kappa\ge\max\left(1,10L\right),\qquad v^2\ge v\ge10L.
\]
\end{itemize}
\end{named}

\begin{named}{Remark}
Proposition~14.2 corresponds to the following places of \cite{QEDI}:
the order of choice and the conditions, \cite[Sections~5.6.1 and~5.6.2]{QEDI};
the paragraph on the conditions on the constants at the beginning of \cite[Section~16]{QEDI}, and the paragraph on the order of the constants in \cite[Section~16.14]{QEDI};
\cite[Section~17.1.1]{QEDI}, except its closing passage on \cite[Lemmas~17.3--17.8]{QEDI} at the level $k+1$;
the enumeration of the conditions in \cite[Section~17.13.3]{QEDI};
\cite[Section~19.1.2]{QEDI};
\cite[Sections~20.1.7 and~20.1.8]{QEDI};
and the conditions on $e_2$ in \cite[Section~21.1.2]{QEDI}.
Clause (c) takes the place of \cite[Proposition~19.2]{QEDI}, which is formulated there with $\gamma_1$, with $L^{m_0+1}\cdot\mathrm{lh}(v)$, and with the constants and the number $e_1$ of \cite[Section~19.1.2]{QEDI}.

The conditions of (b), together with the ranges printed in (a), are the conditions of \cite[Sections~5.6.1 and~5.6.2]{QEDI} on the constants of the lines (C0)--(C9), in full, with the changes of Section~I.8.4 and with the following ones.
The parameter $\rho$ and the weights are chosen first.
No stage of the order of choice in \cite[Definition~5.97]{QEDI} corresponds to the lines (C0r) and (C4r);
of their constants, $\gamma^{F}$ and $\gamma_{\mathrm{w}}$ occur at stage~5 of \cite[Definition~5.97]{QEDI} and $M^{F}$, $M^{P}$ at stage~6, and for the others the relation to \cite{QEDI} is the one stated with the lines (C0r) and (C4r) in Section~I.8.2.
The conditions (M5) and (Z2.s2) are formed with $C'_{V}$ and $a'$ in place of $C_V$ and $a$.
The condition $M\ge L^{m_0+2}$ of \cite{QEDI} is (M1).
In (A6) the numeral $112$ stands in place of $16$.
Of the condition (A10) of \cite[Section~5.6.2]{QEDI}, which is a condition on the constants alone, (b) lists under (A10) the four inequalities that it imposes for every $e_k\in(0,e_0]$, in the same form, as in Section~I.8.4;
its last inequality,
\[
8\,c_{\mathrm{rb}}\cdot L^3\cdot C_{d}\cdot M^3\cdot e_0^{2\varepsilon}\le1,
\]
stands in (b) as (Z2.s5), and its first inequality, $e_0\le e^{F}_0$, has no counterpart in (b).
The condition (Z2.s5) is formed with the constant $C_{d}$ of Proposition~13, computed with the constant $C_{w}$ of $\mathrm{par}$.
The number $e^{F}_0$ does not occur (Section~I.8.4).
In (e3) the factor $L\cdot\mathrm{lh}(v)$ stands in place of $L^{m_0+1}\cdot\mathrm{lh}(v)$.
The condition (e6) is formed with $\gamma_{\Sigma}$ in place of $\gamma_1$.
The upper bound $\min\left(e_0,\frac{1}{3}\,L^{-\frac{m_0}{2}}\right)$ of $e_1$ becomes $\min\left(e_0,\frac{1}{3}\right)$.
In the line (C3), that is, in the choice of $L$, the words ``for every $\varepsilon\in(0,\bar{\varepsilon}]$'' are omitted, since the five conditions on $L$ contain $\bar{\varepsilon}$ and not $\varepsilon$.
The constants $C_K$, $C_K'$, $C_H$, $c_{\tau}$ are those of the Propositions~33.m.

Furthermore, (a) records for every constant the letters on which it depends, and (b) records for every condition the letters that it contains.
The conditions on the constants of the lines (C0)--(C9) stand in (a), (b) and (f) in full.
Among them are, with the changes described above (the inequality $e_0\le e^{F}_0$ being omitted), all the conditions on the constants that are enumerated in \cite[Section~17.13.3]{QEDI}, in particular
\[
0<\alpha,\qquad1-\varepsilon<\alpha,\qquad\text{(L1)},\qquad\text{(M5)},\qquad\text{(M6)},\qquad r\ge2,\qquad C_{w}\ge C_{P},\qquad\gamma_{\mathrm{w}}\le\gamma_{P}.
\]
That enumeration contains the conditions on which \cite[Theorem~17.2(a)--(g)]{QEDI} rests, whereas Proposition~14.2 contains the conditions on all constants of the lines (C0)--(C9).
The running charge $e_k$ is a datum (Section~G1.1), not a constant.
The number $\kappa'$ is the number $\kappa_{\mathrm{w}}+\kappa_{0}$ of the line (C4).
Of the numbers defined in \cite[Section~20.1.7]{QEDI}, those that do not depend on $e$ are entries of (a), namely the constants of the line (C6$'$) and the constant $C_{\mathrm{fr}}$ of the line (C10), with two exceptions: $M_L$ is the number $\mathrm{mxL}$ of (d), and $z_2$ is no entry, the number $t_2$ of (d) being $e\cdot z_2$.
The quantities formed with $e$ stand in (d).
The constants of the Propositions~1.m stand in the lines (C0r) and (C4r) in the order of their definitions.
The third condition on $e_2$ in \cite[Section~21.1.2]{QEDI}, which contains the constant $C_Y$ and serves the averaged law in all volumes, has no counterpart here, and $C_Y$ does not occur.
\end{named}

\subsubsection*{Proposition~14.3 (admissible constants exist)}

\begin{named}{Proposition 14.3}[Admissible constants exist]
A choice of the lines (C0)--(C$j$) is called \emph{admissible} if every constant of these lines is chosen as Proposition~14.2\,(a) prescribes, subject to the conditions of Proposition~14.2\,(b) that restrict it.
For given values of the constants that precede a chosen constant, a value of this constant is \emph{admissible} if it lies in the range printed in its entry of Proposition~14.2\,(a) and satisfies the conditions of Proposition~14.2\,(b) that restrict this constant, formed with these values.
Several constants of one line are treated one after the other.
Thus, in (d), the assertion ``$s=\frac{1}{2}$ is admissible for $\varepsilon$'' means $0<\frac{1}{2}<1$ together with (S2) for this $\varepsilon$.
\begin{itemize}
\item[(a)] (The lines (C0)--(C2).)
The ranges are not empty.
\item[(b)] (The choice of $L$.)
For every admissible $\kappa$ there is a number $L_{\min}<\infty$, a function of $\left(\bar{\varepsilon},c_{\mathrm{fl}},c_{\mathrm{rb}},\kappa,\kappa_{0}\right)$, such that every odd integer $L\ge L_{\min}$ satisfies (L1)--(L5).
Indeed, the exponents
\[
\frac{1}{8}-\frac{3\bar{\varepsilon}}{2}=0.1175\qquad\text{and}\qquad\frac{3}{8}-6\bar{\varepsilon}=0.345
\]
are positive, $\kappa-6\kappa_{0}-4\ge\kappa_{0}+3>0$, and $\kappa-\kappa_{0}-1>0$.
\item[(c)] (The choice of $M$.)
For every admissible choice of the lines (C0)--(C4r) the numbers
\[
\frac{16}{\gamma_{\mathrm{w}}},\qquad\left(1+\exp\left(\hat{\kappa}\right)\right)^{\frac{1}{\alpha_{0}}},\qquad M^{F},\qquad M^{P}
\]
are finite, and every power $M$ of $L$ with
\[
M\ge M_{\min}:=\max\left\{L^2,\ 4L+1,\ \frac{16}{\gamma_{\mathrm{w}}},\ \left(1+\exp\left(\hat{\kappa}\right)\right)^{54},\ M^{F},\ M^{P}\right\}
\]
satisfies (M1)--(M6).
All conditions on $M$ are lower bounds.
\item[(d)] (The line (C7).)
For every admissible $\alpha$ and $L$ the set of the admissible $(r,p_0,p,\varepsilon,s)$ is not empty.
The value $s=\frac{1}{2}$ is admissible for every $\varepsilon\in(0,\bar{\varepsilon}]$, since
\[
\frac{1}{2}\left(\frac{1}{4}-9\varepsilon\right)-12\varepsilon\ge0.0425>\frac{2a_{L}}{N_{*}}.
\]
\item[(e)] (The choice of $e_0$.)
For every admissible choice of the lines (C0)--(C7) the set of the numbers $e_0>0$ that satisfy (A1)--(A10), (Z2.s1)--(Z2.s5) and (Z3.s1) is an interval with left end point $0$ and of positive length.
Indeed, each inequality of the conditions (A1)--(A10), (Z2.s1)--(Z2.s5) and (Z3.s1) is of one of the following four kinds, and it holds for all $v$ beyond some number:
a sum of terms $(\text{constant})\cdot\exp\left(-\mathrm{cf}\cdot v\right)$ compared with $1$, $\frac{1}{2}$, $\frac{1}{4}$ or $3$, where
\[
\mathrm{cf}\in\left\{\frac{\varepsilon}{4},\ \frac{\varepsilon}{2},\ \varepsilon,\ \frac{3\varepsilon}{2},\ 2\varepsilon,\ \frac{5\varepsilon}{2},\ 3\varepsilon,\ \frac{1}{4},\ \frac{1}{4}+\varepsilon,\ \frac{1}{4}+2\varepsilon,\ \frac{1}{4}-5\varepsilon,\ \frac{1}{4}-6\varepsilon,\ \frac{3}{4}-\frac{7\varepsilon}{2}\right\},
\]
all these numbers being positive because $\varepsilon\le\frac{1}{200}$;
a polynomial in $v$ times $\exp\left(-\mathrm{cf}\cdot v\right)$ with $\mathrm{cf}>0$ (the first inequality of (A1), and the second and the third inequality of (A10));
the inequality $12\exp\left(7v-\frac{v^{2p_0}}{2}\right)\le1$, where $2p_0>1$;
a lower bound on $v$.
The inequalities are imposed for all $v\ge-\log e_0$; hence they remain valid for every smaller value of $e_0$.
\item[(f)] (The number $e_1$.)
For every admissible choice of the lines (C0)--(C8) the supremum in the definition of $e_1$ is positive.
Hence $e_1>0$, $e_1\le\frac{1}{2}\min\left(e_0,\frac{1}{3}\right)$, and (e1)--(e8) hold for all $v\ge-\log e_1$.
The number $e_1$ is a function of $e_0$ and of
\[
\left(L,M,b_{-},r,p_0,p,\varepsilon,c_{\mathrm{fl}},c_{\kappa},C_{0},\gamma_{\Sigma},C_{\Sigma},N_{*}\right);
\]
it does not depend on $\bar{m}$, $N$, $T$, $\theta$.
\item[(g)] (The line (C10).)
For every $\bar{m}\in(0,1]$ one has $e_G>0$, and both conditions in the definition of $e_G$ remain valid for every smaller $e$.
That $e_{\sharp}(\beta')>0$ is Proposition~35.10.
Moreover $e_2>0$: given $e_G>0$, $e_{\sharp}(\beta_{\mathrm{ix}})>0$ for $\mathrm{ix}=0,+,-$, $C'_{\bar{\kappa}}<\infty$, $b(f)>0$ and $\sigma>0$, the conditions (Z10.s1), (Z10.s2), (Z10.s4) hold for all sufficiently small positive numbers and remain valid for every smaller number, and $e_2$ is the largest number that satisfies them.
Pairs $(\beta,x_0)$ with (Pb) of Proposition~14.2\,(d) exist.
Indeed, probes with $b(f)>0$ exist by \cite[Theorem~21.1(a)]{QEDI}; the left side of the last inequality in (Pb), the inequality of \cite[Section~2.8]{QEDI}, tends to $0$ as $|x_0|_2\to\infty$, while its right side $\frac{b(f)}{12}$ is positive.
\end{itemize}
\end{named}

\begin{named}{Remark}
Proposition~14.3 corresponds to the assertions of existence in \cite[Sections~5.6.1, 17.13.3 and~19.1.2]{QEDI}.
Its clauses (b), (c) and (e) correspond to the assertions of existence on the odd integers $L$ of stage~4, on the powers $M$ of stage~6 and on the numbers $e_0$ of stage~9 of \cite[Definition~5.97]{QEDI}, among those that follow \cite[Definition~5.98]{QEDI} in \cite[Section~5.6.1]{QEDI}, and these stand in the second paragraph after that definition;
its clause (d) corresponds to the assertion of existence of admissible numbers $s$ in the first paragraph after \cite[Definition~5.98]{QEDI}, with the same value $s=\frac{1}{2}$ and the same bound $0.0425>\frac{2a_{L}}{N_{*}}$;
the clause (e) corresponds also to the end of the paragraph that follows the condition (A10) in \cite[Section~5.6.2]{QEDI}.
The arguments are the same: for $L$, the positivity of the same four numbers; for $M$, lower bounds by finite numbers; for $e_0$, inequalities that hold for all sufficiently large $v$.
The differences are the following.
The conditions are those of Proposition~14.2\,(b), formed with the constants of this paper: the numeral $112$ stands in (A6), $C'_{V}$ and $a'$ enter $M^{F}$ and (Z2.s2), the constant $C_{d}$ is that of Proposition~13, the number $e^{F}_0$ does not occur, and, since $m_0=0$ (Proposition~14.1\,(v)), the condition $M\ge L^{m_0+2}$ of stage~6 is (M1).
The clauses (b) and (c) name thresholds $L_{\min}$ and $M_{\min}$, the first as a function of $\left(\bar{\varepsilon},c_{\mathrm{fl}},c_{\mathrm{rb}},\kappa,\kappa_{0}\right)$, the second by a formula;
and the clause (e) states that the admissible values of $e_0$ form an interval with left end point $0$.
The constants $a'$, $\gamma^{F}$, $C^{F}$, $C_{w}$, $C_{d}$, $C_{0}$, $\gamma_{\Sigma}$, $C_{\Sigma}$, which are defined differently from the corresponding constants of \cite{QEDI}, enter (c), (e) and (f) only through being finite and positive.
\end{named}

\subsubsection*{Proposition~14.4 (what does not depend on what)}

\begin{named}{Proposition 14.4}[What does not depend on what]
\begin{itemize}
\item[(a)] No constant of the lines (C0)--(C9) depends on $N$, on the level $k$, on the side $\ell$, on the twist $\theta$, on the running mass in $(0,1]$, on the pair $\left(m_{\mathrm{in}},\varepsilon^{\mathrm{v}}_{\mathrm{in}}\right)$, on $\mathscr{C}$, or on $\bar{m}\in(0,1]$.
\item[(b)] The following do not depend on $\rho$:
$\bar{\varepsilon}$, $c_{X}$, $\kappa_{0}$, $K_{0}$, $c_{\mathrm{rb}}$, $c_{\mathrm{fl}}$, $b$;
$\lambda_{\mathrm{int}}$, $\lambda_1$, $c_{*}$, $C_{\mathrm{Po}}$ (functions of the weights);
$\alpha$, $\kappa$, $L$;
$b_{L}$, $b_{-}$, $b_{+}$, $N_{*}$, $a_{L}$, $G_{L}$, $K_{L}$, $N_c$, $c_{l}$, $d_U$, $c_{C}$, $C_{Q}$, $\alpha^{P}$, $\alpha^{F}$, $\alpha_{0}$, $c_{P}$, $\gamma_{P}$, $\gamma^{P}$, $C_{P}$, $\kappa_{\mathrm{w}}$, $\kappa'$, $\hat{\kappa}$, $C'_{V}$, $C_{T}$, $\mu_-$, $K_{1}$, $k_{1}$, $M^{P}$;
$C_{\mathrm{By}}$, $c_{\kappa}$, $c_{S}$, $\theta_{\mathrm{CT}}$, $\mathrm{gGa}$, $h_{G}$, $\mathrm{gGa}_{1}$, $\mathrm{alG}$;
$r$, $p_0$, $p$, $\varepsilon$, $s$, $\mathfrak{e}$, $\sigma$;
all exponents in the first nine rows of the table of Section~I.8.5 (the numbers to which $e_k$, $e_{k+1}$, $e_{*}$, $L$, $M$, $\eta$ are raised there, and $\mathfrak{e}$, $\sigma$, $\mathfrak{e}_2$), and the rates $\kappa_{\mathrm{w}}+\kappa_{0}$, $\kappa_{\mathrm{w}}$, $\kappa$ of its last row (but not $\gamma_{\mathrm{w}}$ and $\ell_H=M$, which carry $\rho$; see (c), (d)).
In particular (L1)--(L5), (S1), (S2) contain neither $\rho$ nor a fermionic constant.
\item[(c)] The parameter $\rho$ enters $M$ only through the lower bounds (M3), (M5), that is, through $\gamma^{F}$ and $a'$:
in every chain $\rho=Y_0,\dots,Y_m=M$ in which $Y_{i-1}$ is a parent of $Y_i$ for every $i$, the letter $Y_{m-1}$ appears in (M3) or in (M5) and in no other condition on $M$, these two conditions bound $M$ from below, and the chain contains $\gamma^{F}$ or $a'$.
No condition on $M$ contains a constant of the lines (C6)--(C10), and no condition bounds $M$ from above.
The constants $C_{w}$, $C_{T}$, $c_{\gamma}$, $C_{B}$ enter conditions on $e_0$ only; $C_{d}$ enters conditions on $e_0$ and conditions of the line (C10); $C_{0}$, $\gamma_{\Sigma}$, $C_{\Sigma}$, $c_{\kappa}$ enter conditions on $e_1$ and conditions of the line (C10) only.
\item[(d)] As $\rho\to0$, the constants $c_{b}$, $c_{\mathrm{re}}$, $\gamma^{F}$ tend to $0$, and $a'$, $M_{\min}$ tend to infinity ($M_{\min}$ as in Proposition~14.3\,(c)).
\end{itemize}
In (a)--(d) the following terms are used.
The \emph{letters} are those that head the entries of Proposition~14.2\,(a); the abbreviations $\mathrm{lg}$, $c$, $c_{\mathrm{bd}}$, $\mathrm{ssum}$, $\mathrm{lh}$ are read with their definitions.
A letter is \emph{chosen} if its entry prescribes a set and not a value (these are $\rho$, $\alpha$, $\kappa$, $L$, $M$, $r$, $p_0$, $p$, $\varepsilon$, $s$, $e_0$; $\bar{m}$, $\beta$, $x_0$); otherwise it is \emph{defined}.
The \emph{parents} of a letter are the letters named after ``depends on'' in its entry; its \emph{ancestors} are its parents, their parents, and so on.
The phrase ``$X$ does not depend on $Q$'' means the following: if $Q$ is a letter of the list, then $Q$ is not an ancestor of $X$; if $Q$ is not a letter of the list ($N$, a level, a side, a twist, a running mass, a pair, a connection), then $Q$ is a free variable neither in the definition of $X$ nor in the definition of an ancestor of $X$.
A defined letter has one value as soon as the letters before it have values (see the paragraph that opens Proposition~14.2\,(a)).
For some constants of the lines (C0)--(C9) the definition, read together with the places that the entry names, contains letters of the list that are not named after ``depends on''; these letters are the following:
$w$, a vector of numerals;
$\rho$ and $\kappa_{(0)}$ in the entry $A_{\mathrm{ak}},A'_{\mathrm{ak}}$ (here $\rho$ is an ancestor, through $c_{\mathrm{el}}$, and $\kappa_{(0)}$ is a function of $\theta_{0}$, $c_{\mathrm{el}}$, $K_{\mathrm{Str}}$);
$\lambda_1$ in the entry $C_{E}$, and $C_{E}$ in the entry $\mathrm{Csy}_{\mathrm{ak}},\mathrm{Cpv}_0$ (functions of $\rho$, $w$);
$\hat{g}$, $c_{\mathrm{N3}}$ in the entries $\gamma_{\Sigma}$, $C_{\Sigma}$ (functions of $\rho$, $w$);
$M$ in the entries $c_{\gamma}$, $c_{\kappa}$ (it drops out);
$\alpha$ in the entry $C_{G}$, and $\bar{m}$ in the entry $\mathrm{Csy}_{\mathrm{ak}},\mathrm{Cpv}_0$ (letters of later entries, over whose ranges the property is quantified; see the paragraph that opens Proposition~14.2\,(a)).
With these conventions, the value of $X$ (for a chosen letter $X$: the set of its admissible values) is a function of the values of the chosen letters among the ancestors of $X$.
In (c) a \emph{condition} is a range printed in an entry of Proposition~14.2\,(a), a condition of Proposition~14.2\,(b), or a display of Proposition~14.2\,(d) (these displays are the conditions of the line (C10)); the displays (T2)--(T5) of Proposition~14.2\,(c), which follow from (e3)--(e6), are counted among the conditions on $e_1$.
A condition \emph{contains} a letter $Y$, and $Y$ \emph{enters} the condition, if $Y$ appears in it after the abbreviations, the quantities that depend on $e$, and the number $\mathrm{mxL}$ of Proposition~14.2\,(d) have been replaced by their definitions.
A \emph{fermionic constant} is a letter of the lines (C0r), (C4r), a letter whose entry in Proposition~14.2\,(a) names as its reference one of the Propositions~1.m and~2.m, one of Propositions~3--13, Proposition~14.1, one of the Propositions~30.m and~33.m, Proposition~36 or one of the Propositions~37.m, or a letter that has $\rho$, $w$ or one of these letters among its ancestors.
\end{named}

\begin{named}{Remark}
Proposition~14.4 corresponds to \cite[Definition~5.1]{QEDI}, to the closing sentences of the stage of \cite[Definition~5.97]{QEDI} in which $M$ is chosen, and to the clause on uniformity of \cite[Theorem~17.2]{QEDI}.
Its clauses (b) and (d), and the assertions of its clause (c) on $\rho$, have no counterpart in \cite{QEDI}.
For the constants of \cite{QEDI}, the second sentence of (c), and the assertion of its third sentence on $C_{w}$, $C_{T}$ and $c_{\gamma}$, have counterparts in the closing sentence of stage~6 of \cite[Definition~5.97]{QEDI}, in the passage of \cite[Section~5.6.1]{QEDI} on the existence of the powers $M$, and in the third paragraph of the proof of \cite[Theorem~17.2(h)]{QEDI} in \cite[Section~17.13.3]{QEDI}.
\end{named}

\subsection*{16 (continued). The group lemma, the low kernels, the extraction}

\subsubsection*{Proposition 16.2 (the low kernels)}

\begin{named}{Proposition 16.2}[the low kernels]
Assume (G2.S-1).
Let $\mathscr{E}$ be a family of elements $\mathscr{E}(X';\mathscr{A},\mathscr{C})$ of $\mathrm{Gr}[\psi,\chi]$, indexed by the polymers $X'\in D_k$, which is balanced in the sense of (G2.r8) and defined at $(\mathscr{A},\mathscr{C})=(0,0)$, and whose kernels $\mathscr{E}_{20}(X';0,0;x,x')$ vanish unless both $x$ and $x'$ lie over $X'$.
Let $X$ be a small polymer, and let $\tilde{g}=(S,a)\in\mathbb{G}_k$.
Define
\[
\mathrm{al}_0(\mathscr{E},X):=\operatorname{Vol}(X)^{-1}\cdot\mathscr{E}_{00}(X;0,0),\qquad
\mathrm{al}_2(\mathscr{E},X):=\operatorname{Vol}(X)^{-1}\int\mathrm{d}x\,\mathrm{d}x'\;\mathscr{E}_{20}(X;0,0;x,x').
\]
Here $\mathrm{al}_2(\mathscr{E},X)$ is a $4\times4$ matrix;
each of the two integrations carries the factor $\frac{1}{48}$ together with the sum over the label,
and all arguments range over one lift of $X$.
\begin{itemize}
\item[(o)] Neither $\mathrm{al}_0(\mathscr{E},X)$ nor $\mathrm{al}_2(\mathscr{E},X)$ depends on the lift of $X$ over which the arguments are taken.
\item[(i)] Suppose that $\mathscr{E}$ has the property (I.k3) at $(0,0)$.
Then
\[
\mathrm{al}_0(\mathscr{E},\tilde{g}X)=\mathrm{al}_0(\mathscr{E},X),\qquad
\mathrm{al}_2(\mathscr{E},\tilde{g}X)=S\cdot\mathrm{al}_2(\mathscr{E},X)\cdot S^{-1},
\]
and, for the kernels,
\[
S^{-1}\cdot\mathscr{E}_{20}(\tilde{g}X;0,0;\tilde{g}x,\tilde{g}x')\cdot S=\mathscr{E}_{20}(X;0,0;x,x').
\]
\item[(ii)] Suppose that $\mathscr{E}$ has the property (I.k5) at $(0,0)$.
Then
\[
\mathscr{E}_{20}(X;0,0;x,x')=\mathsf{C}^{-1}\cdot\mathscr{E}_{20}(X;0,0;x',x)^T\cdot\mathsf{C},\qquad
\mathrm{al}_2(\mathscr{E},X)=\mathsf{C}^{-1}\cdot\mathrm{al}_2(\mathscr{E},X)^T\cdot\mathsf{C}.
\]
\item[(iii)] Suppose that $\mathscr{E}$ has the property (I.k7) at $(0,0)$.
Then $\mathrm{al}_0(\mathscr{E},X)$ is real, and
\[
\mathscr{E}_{20}(X;0,0;x,x')=\varkappa^{-1}\cdot\overline{\mathscr{E}_{20}(X;0,0;x,x')}\cdot\varkappa,\qquad
\mathrm{al}_2(\mathscr{E},X)=\varkappa^{-1}\cdot\overline{\mathrm{al}_2(\mathscr{E},X)}\cdot\varkappa.
\]
\item[(iv)] Let $\mathrm{Mx}$ be a $4\times4$ matrix and $Z$ a union of $M$-cubes.
The kernel of the element $\int_Z\mathrm{d}x\;\bar{\psi}(x)^T\cdot\mathrm{Mx}\cdot\psi(x)$ is $48\cdot\eta_k^{-3}\cdot\mathbf{1}[\,x=x'\text{ over }Z\,]\cdot\mathrm{Mx}$.
This is its value at points $x$, $x'$ over one set of sites of $\mathbb{L}_k$ on which the covering map $\mathrm{pr}$ is injective;
as a function on $\mathbb{L}^{\mathrm{f}}_k\times\mathbb{L}^{\mathrm{f}}_k$ that is antiperiodic in each argument, it is
\[
48\cdot\eta_k^{-3}\cdot\mathbf{1}[\,\mathrm{pr}\,c\in Z\,]\cdot\sum_{n\in\mathbb{Z}^3}(-1)^{n_1+n_2+n_3}\cdot\mathbf{1}[\,x'=x+\ell n\,]\cdot\mathrm{Mx}\qquad\left(x=(c,g)\right).
\]
The double integral of this kernel equals $\operatorname{Vol}(Z)\cdot\mathrm{Mx}$, both arguments ranging over one fundamental domain or, if $Z$ has a lift, over one lift;
moreover
\[
\left\|\int_Z\mathrm{d}x\;\bar{\psi}(x)^T\cdot\mathrm{Mx}\cdot\psi(x)\right\|_{\mathsf{w}}=\mathsf{w}_1^2\cdot\operatorname{Vol}(Z)\cdot\sum_{\beta,\beta'}\left|\mathrm{Mx}_{\beta\beta'}\right|.
\]
For every real function $\Lambda$ of the photon site $c$ one has, with $\psi=\psi(x)$ and $\bar{\psi}=\bar{\psi}(x)$,
\begin{align*}
a_{\tilde{g}}\left[\int_{\tilde{g}Z}\mathrm{d}x\;\bar{\psi}^T\cdot S\cdot\mathrm{Mx}\cdot S^{-1}\cdot\psi\right]&=\int_Z\mathrm{d}x\;\bar{\psi}^T\cdot\mathrm{Mx}\cdot\psi,\\
G_{\Lambda}\left[\int_Z\mathrm{d}x\;\bar{\psi}^T\cdot\mathrm{Mx}\cdot\psi\right]&=\int_Z\mathrm{d}x\;\bar{\psi}^T\cdot\mathrm{Mx}\cdot\psi,\\
\mathrm{Cg}\left[\int_Z\mathrm{d}x\;\bar{\psi}^T\cdot\mathrm{Mx}\cdot\psi\right]&=\int_Z\mathrm{d}x\;\bar{\psi}^T\cdot\mathsf{C}^{-1}\cdot\mathrm{Mx}^T\cdot\mathsf{C}\cdot\psi,\\
\mathbf{K}\left[\int_Z\mathrm{d}x\;\bar{\psi}^T\cdot\mathrm{Mx}\cdot\psi\right]&=\int_Z\mathrm{d}x\;\bar{\psi}^T\cdot\varkappa^{-1}\cdot\overline{\mathrm{Mx}}\cdot\varkappa\cdot\psi.
\end{align*}
\item[(v)] \emph{General kernels.}
Let $K$ be a function on $\mathbb{L}^{\mathrm{f}}_k\times\mathbb{L}^{\mathrm{f}}_k$ with values in the complex $4\times4$ matrices which is antiperiodic in each argument,
\[
K(x+\ell e_\mu,x')=K(x,x'+\ell e_\mu)=-K(x,x')\qquad(\mu=1,2,3),
\]
and put
\[
\mathrm{Bil}[K]:=\int\mathrm{d}x\,\mathrm{d}x'\;\bar{\psi}(x)^T\cdot K(x,x')\cdot\psi(x').
\]
The element $\mathrm{Bil}[K]$ is balanced and of bidegree $(2,0)$, and in the matrix notation of Section~I.5.1 its kernel is $K$;
consequently $\mathrm{Bil}[K]=\mathrm{Bil}[K']$ implies $K=K'$.
Every balanced element of bidegree $(2,0)$ is of the form $\mathrm{Bil}[K]$, and
\[
\left\|\mathrm{Bil}[K]\right\|_{\mathsf{w}}=\mathsf{w}_1^2\cdot\sum_{\beta,\beta'}\int\mathrm{d}x\,\mathrm{d}x'\;\left|K(x,x')_{\beta\beta'}\right|;
\]
this is (G2.0.1).
Let $\tilde{g}=(S,a)\in\mathbb{G}_k$, let $\Lambda$ be a real function of $c$, and write $x=(c,g)$, $x'=(c',g')$.
Then
\begin{align*}
a_{\tilde{g}}\,\mathrm{Bil}[K]&=\mathrm{Bil}[K^{\tilde{g}}],&K^{\tilde{g}}(x,x')&:=S^{-1}\cdot K(\tilde{g}x,\tilde{g}x')\cdot S,\\
G_{\Lambda}\,\mathrm{Bil}[K]&=\mathrm{Bil}[K^{\Lambda}],&K^{\Lambda}(x,x')&:=\exp\left\{i\,e_k\cdot\left(\Lambda(c')-\Lambda(c)\right)\right\}\cdot K(x,x'),\\
\mathrm{Cg}\,\mathrm{Bil}[K]&=\mathrm{Bil}[K^{\mathrm{Cg}}],&K^{\mathrm{Cg}}(x,x')&:=\mathsf{C}^{-1}\cdot K(x',x)^T\cdot\mathsf{C},\\
\mathbf{K}\,\mathrm{Bil}[K]&=\mathrm{Bil}[K^{\mathbf{K}}],&K^{\mathbf{K}}(x,x')&:=\varkappa^{-1}\cdot\overline{K(x,x')}\cdot\varkappa,
\end{align*}
and each of the four functions $K^{\tilde{g}}$, $K^{\Lambda}$, $K^{\mathrm{Cg}}$, $K^{\mathbf{K}}$ is again antiperiodic in each argument.
\end{itemize}
\end{named}

\begin{named}{Remark}
For a single small polymer the matrix $\mathrm{al}_2(\mathscr{E},X)$ need not be a multiple of $\mathbf{1}$.
Indeed, there exist a family $\mathscr{E}\in K^{x}_{k}$ and a small polymer $X$ with
\[
\mathrm{al}_2(\mathscr{E},X)=i\,\Sigma_3
\]
(Chapter~16, Section~2, paragraph ``Verification of the Remark'').
\end{named}

\begin{named}{Remark}
Proposition~16.2 corresponds to \cite[Lemma~13.4]{QEDI} and to \cite[Proposition~5.59(e)]{QEDI}.
In contrast to these, $x$ and $x'$ are fermion sites and $\int\mathrm{d}x$ is the measure (I.1.3);
so $\mathrm{al}_2$ is integrated over both labels, in accordance with (I.E1).
The kernel of a local quadratic term is $48\cdot\eta_k^{-3}\cdot\mathbf{1}[\,x=x'\,]\cdot\mathrm{Mx}$, whereas in \cite[Section~13.4.2]{QEDI} it is $\eta^{-3}\cdot\mathbf{1}[\,x=y\,]\cdot N$.
The assertions (i)--(iii) of \cite[Lemma~13.4]{QEDI} concern $\mathrm{al}_0$ and $\mathrm{al}_2$ only;
the relations between the kernels in (i)--(iii) occur in the proof of that lemma.
Clause (o) has no counterpart among the assertions of \cite[Lemma~13.4]{QEDI}; for kernels with an even number of arguments the corresponding fact is stated in \cite[Section~13.1.3]{QEDI}.
Clause (v) has no such counterpart either:
three of its four rules, namely those for $a_{\tilde{g}}$, $\mathrm{Cg}$ and $\mathbf{K}$, occur in the proof of \cite[Lemma~13.4]{QEDI} as identities for an arbitrary kernel.
\end{named}

\subsection*{18. Global fields, locality, scaling, contraction}

\subsubsection*{Proposition 18.1 (scaling of the thresholds; a statement on photon fields; it contains the case $e_{k+1}=L^{\frac{1}{2}}\cdot e_k\le e_0$ of \cite[Lemma~5.18]{QEDI})}

\begin{named}{Proposition 18.1}[{scaling of the thresholds; a statement on photon fields; it contains the case $e_{k+1}=L^{\frac{1}{2}}\cdot e_k\le e_0$ of \cite[Lemma~5.18]{QEDI}}]
Assume (G2.S-2).
For a bond function $A$ of the level $k+1$ let $A_L$ be given by $A_L(b)=L^{-\frac{1}{2}}\cdot A(b/L)$.
The thresholds $\mathrm{t1}_k$, $\mathrm{t2}_k$, $\mathrm{t3}_k$ of $R_k$ and $\mathrm{td1}_k$, $\mathrm{td2}_k$ of $\mathfrak{X}_k$ are those of (G2.r4).
\begin{itemize}
\item[(a)] Let $Z'$ be a subset of $T_{k+1}$ and $A$ a bond function on $Z'$.
On $LZ'$ one has
\[
(A_L)_\mu(c)=L^{-\frac{1}{2}}\cdot A_\mu(c/L),\qquad
\left(\partial_\nu(A_L)_\mu\right)(c)=L^{-\frac{3}{2}}\cdot\left(\partial_\nu A_\mu\right)(c/L),
\]
and the H\"older quotient of $\partial A_L$ at $(c,c')$ equals $L^{-\frac{3}{2}-\alpha}$ times the H\"older quotient of $\partial A$ at $(c/L,c'/L)$.
If the three norms of $A$ on $Z'$ are smaller than the thresholds of $R_{k+1}$, then on $LZ'$
\[
|A_L|<L^{-\frac{7}{8}+\frac{\varepsilon}{2}}\cdot\mathrm{t1}_k,\qquad
|\partial A_L|<L^{-\frac{15}{8}+\varepsilon}\cdot\mathrm{t2}_k,\qquad
|\delta_\alpha\partial A_L|<L^{-\alpha}\cdot L^{-\frac{15}{8}+\frac{3\varepsilon}{2}}\cdot\mathrm{t3}_k.
\]
\item[(b)] If $\mathscr{C}\in\mathfrak{X}_{k+1}$, then
\[
|\mathscr{C}_L|<L^{-1+\varepsilon}\cdot\mathrm{td1}_k,\qquad
|\partial\mathscr{C}_L|<L^{-2+\frac{3\varepsilon}{2}}\cdot\mathrm{td2}_k;
\]
hence $\mathscr{C}_L$ lies in $L^{-\frac{3}{4}}\cdot\mathfrak{X}_k$, a subset of $\mathfrak{X}_k$.
\item[(c)] Let $\mathscr{A}=A_0+A_1\in\tilde{R}_{k+1}$ be split as in \cite[Definition~5.16]{QEDI}, let $\square'$ be an $M$-cube of the level $k+1$, and let $A_0=A'+d\boldsymbol{\lambda}$ on $\square'^{\dagger}$, with a real function $\boldsymbol{\lambda}$ and a real bond function $A'$ as in that definition.
Let $z'$ be the centre of the $LM$-cube $L\square'$, and let $\boldsymbol{\lambda}_L(c)=L^{\frac{1}{2}}\cdot\boldsymbol{\lambda}(c/L)$ as in (G2.r5).
Then
\[
(A_0)_L=(A')_L+d(\boldsymbol{\lambda}_L)\quad\text{on}\quad L\square'^{\dagger}=\left\{c:\ |c-z'|<\left(L+\frac{1}{2}\right)\cdot LM\right\};
\]
the bond function $(A')_L$ satisfies the bounds of (a) on $L\square'^{\dagger}$, and $(A_1)_L$ satisfies them on $T_k$.
For every $M$-cube $\square$ of the level $k$ that is contained in $L\square'$, the set $\square^{\dagger}$ is contained in $L\square'^{\dagger}$.
Consequently $\mathscr{A}_L$ lies in $L^{-\frac{3}{4}}\cdot\tilde{R}_k$, and this set is contained both in $\tilde{R}_k$ and in $\tilde{R}_k^{-}$.
\end{itemize}
\end{named}

\begin{named}{Remark}
Proposition~18.1 corresponds to \cite[Lemma~13.8]{QEDI}.
The content differs only in the last assertion of (b), the inclusion of $L^{-\frac{3}{4}}\cdot\mathfrak{X}_k$ in $\mathfrak{X}_k$, which is contained in \cite[Lemma~17.5(c)]{QEDI}.
\end{named}

\subsection*{16 (continued). The group lemma, the low kernels, the extraction}

\subsubsection*{Proposition 16.3 (extraction; the chamber form of \cite[Proposition~5.59]{QEDI})}

\begin{named}{Proposition 16.3}[{extraction; the chamber form of \cite[Proposition~5.59]{QEDI}}]
Assume (G2.S-1), and in (e) assume (G2.S-2).
Let $\mathscr{E}\in K^{x,C}_{k}$.
For an $M$-cube $\square$ put
\[
\varepsilon(\mathscr{E},\square):=-\sum_{X\text{ small},\ X\supset\square}\mathrm{al}_0(\mathscr{E},X),\qquad
m(\mathscr{E},\square):=-\sum_{X\text{ small},\ X\supset\square}\mathrm{al}_2(\mathscr{E},X),
\]
and let $R\mathscr{E}$ be the family of Section~I.5.6.
\begin{itemize}
\item[(a)] Neither the number $\varepsilon(\mathscr{E},\square)$ nor the matrix $m(\mathscr{E},\square)$ depends on $\square$;
let $\varepsilon(\mathscr{E})$ denote the common value of the former.
Moreover $m(\mathscr{E},\square)=m(\mathscr{E})\cdot\mathbf{1}$ with a complex number $m(\mathscr{E})$.
The maps $\varepsilon(\cdot)$ and $m(\cdot)$ are complex linear on $K^{x,C}_{k}$, and
\[
\varepsilon(\mathscr{E})=\varepsilon(\mathscr{E}|_0),\qquad m(\mathscr{E})=m(\mathscr{E}|_0).
\]
If $\mathscr{E}\in K^{x}_{k}$, both numbers are real.
\item[(b)] One has
\begin{align*}
|\varepsilon(\mathscr{E})|&\le K_0\cdot M^{-3}\cdot\|\mathscr{E}\|^x_k\le c_{\mathrm{fl}}\cdot\|\mathscr{E}\|^x_k,\\
|m(\mathscr{E})|&\le K_0\cdot M^{-3}\cdot e_k^{\frac{1}{2}}\cdot\|\mathscr{E}\|^x_k\le c_{\mathrm{fl}}\cdot e_k^{\frac{1}{2}}\cdot\|\mathscr{E}\|^x_k,
\end{align*}
and every small polymer $X$ satisfies the pair of inequalities (2.2):
\begin{equation}
\begin{aligned}
\operatorname{Vol}(X)\cdot\left|\mathrm{al}_0(\mathscr{E},X)\right|&\le\left\|\mathscr{E}(X;0,0)\right\|_{h_k},\\
\operatorname{Vol}(X)\cdot\sum_{\beta,\beta'}\left|\mathrm{al}_2(\mathscr{E},X)_{\beta\beta'}\right|&\le e_k^{\frac{1}{2}}\cdot\left\|\mathscr{E}^{[2,0]}(X;0,0)\right\|_{h_k}.
\end{aligned}
\tag{2.2}
\end{equation}
\item[(c)] For all $(\mathscr{A},\mathscr{C})\in\tilde{R}_k\times\mathfrak{X}_k$ the identity
\[
\sum_X\mathscr{E}(X;\mathscr{A},\mathscr{C})=-\varepsilon(\mathscr{E})\cdot n_k-m(\mathscr{E})\cdot\langle\bar{\psi},\psi\rangle+\sum_X(R\mathscr{E})(X;\mathscr{A},\mathscr{C})
\]
holds in $\mathrm{Gr}[\psi,\chi]$, where $n_k=\ell^3$ and $\langle\bar{\psi},\psi\rangle=\int\mathrm{d}x\;\bar{\psi}(x)\cdot\psi(x)$.
In particular the extracted mass term is the same at each of the $48$ sites of a cluster.
\item[(d)] The map $R$ is linear, and it maps $K^{x,C}_{k}$ into itself and $K^{x}_{k}$ into itself.
For every $X\in D_k$ and every $(\mathscr{A},\mathscr{C})\in\tilde{R}_k\times\mathfrak{X}_k$
\[
\left\|(R\mathscr{E})(X;\mathscr{A},\mathscr{C})\right\|_{h_k}\le\left\|\mathscr{E}(X;\mathscr{A},\mathscr{C})\right\|_{h_k}+\left\|\mathscr{E}(X;0,0)\right\|_{h_k},
\]
and $\|R\mathscr{E}\|^x_k\le2\cdot\|\mathscr{E}\|^x_k\le c_{\mathrm{fl}}\cdot\|\mathscr{E}\|^x_k$.
For every small $X$
\[
(R\mathscr{E})_{00}(X;0,0)=0,\qquad\int\mathrm{d}x\,\mathrm{d}x'\;(R\mathscr{E})_{20}(X;0,0;x,x')=0.
\]
Furthermore $(R\mathscr{E})|_0=R(\mathscr{E}|_0)$, and if a family has the property (I.k2), then its image under $R$ has the property (I.k2) as well.
\item[(e)] \emph{Reblocking preserves the symmetries.}
Let $E'$ be a family of the level $k$ in the sense of (G2.r8), and put
\[
(L\,E')(Y;\mathscr{A},\mathscr{C}):=\sum_{X:\,\bar{X}^L=LY}S_L\left[E'(X;\mathscr{A}_L,\mathscr{C}_L)\right]\qquad\left(Y\in D_{k+1},\ (\mathscr{A},\mathscr{C})\in\tilde{R}_{k+1}\times\mathfrak{X}_{k+1}\right).
\]
If $E'$ has the property (I.k3) for the elements of $\mathbb{G}_k^{L}$, then $L\,E'$ has the property (I.k3) at the level $k+1$, for the group $\mathbb{G}_{k+1}$.
If $E'$ has the property (I.k4), resp.\ (I.k5), resp.\ (I.k7), then $L\,E'$ has the same property at the level $k+1$.
\item[(f)] \emph{One formula.}
The sides range over the class $\mathrm{Cl}$, and $\mathrm{pr}\colon\mathbb{L}_k\to T_k$ denotes the covering map.
\begin{itemize}
\item[(f1)] For an $M$-cube $\tilde{\square}$ of $\mathbb{L}_k$ let $\mathrm{Sm}(\tilde{\square})$ be the set of the non-empty finite unions $\tilde{X}$ of $M$-cubes of $\mathbb{L}_k$ which are connected through faces, contain $\tilde{\square}$ and satisfy $d_M(\tilde{X})\le L$, the number $d_M(\tilde{X})$ being formed with nets in $\mathbb{R}^3$.
The set $\mathrm{Sm}(\tilde{\square})$ does not depend on $\ell$, and
\[
\mathrm{Sm}(\tilde{\square}+v)=\mathrm{Sm}(\tilde{\square})+v\qquad(v\in M\mathbb{Z}^3).
\]
The map $\tilde{X}\mapsto\mathrm{pr}(\tilde{X})$ is a bijection of $\mathrm{Sm}(\tilde{\square})$ onto the set of the small polymers of $T_k$ that contain $\mathrm{pr}(\tilde{\square})$, and $\tilde{X}$ is the lift of $\mathrm{pr}(\tilde{X})$ in the sense of Section~I.1.1.
\item[(f2)] In the sense of Section~I.1.9 (objects attached to a polymer), the numbers $\mathrm{al}_0(\mathscr{E},X)$ and the matrices $\mathrm{al}_2(\mathscr{E},X)$, $X$ small, as well as the elements $(R\mathscr{E})(X;\mathscr{A},\mathscr{C})$, are given by one formula of range $\mathrm{rg}=\eta_k$.
The data are the antisymmetric kernels of $\mathscr{E}(X;0,0)$ and of $\mathscr{E}(X;\mathscr{A},\mathscr{C})$, with all arguments over the $\eta_k$-neighbourhood of $\tilde{X}$.
Let $K$ stand for the kernels at $(0,0)$ and $F$ for the element at $(\mathscr{A},\mathscr{C})$, let $|\tilde{X}|_M$ be the number of the $M$-cubes of $\tilde{X}$, and put $\operatorname{Vol}(\tilde{X}):=M^3\cdot|\tilde{X}|_M$;
the integrals over subsets of $\mathbb{L}_k$ carry the weight $\eta_k^3/48$ at each fermion site, as in Proposition~16.1\,(a2).
The maps are
\[
\operatorname{Vol}(\tilde{X})^{-1}\cdot K_{00},\qquad
\operatorname{Vol}(\tilde{X})^{-1}\int_{\tilde{X}}\mathrm{d}x\int_{\tilde{X}}\mathrm{d}x'\;K_{20}(x,x'),
\]
and
\[
F-K_{00}-\int_{\tilde{X}}\mathrm{d}x\;\bar{\psi}(x)^T\cdot\left[\operatorname{Vol}(\tilde{X})^{-1}\int_{\tilde{X}}\mathrm{d}x'\int_{\tilde{X}}\mathrm{d}x''\;K_{20}(x',x'')\right]\cdot\psi(x)
\]
if $d_M(\tilde{X})\le L$, while the third map is $F$ otherwise.
\item[(f3)] \emph{Two tori.}
For a side $\ell$ let $T_k(\ell)$ denote the torus $T_k$ formed with the side $\ell$.
Let $\ell_1,\ell_2\in\mathrm{Cl}$, let $\mathrm{pr}_i\colon\mathbb{L}_k\to T_k(\ell_i)$ be the covering maps, and let $\mathscr{E}^{(i)}$ be a family on $T_k(\ell_i)$ as in Proposition~16.2, $i=1,2$.
For an $M$-cube $\square$ of $T_k(\ell_i)$ the number $\varepsilon(\mathscr{E}^{(i)},\square)$ and the $4\times4$ matrix $m(\mathscr{E}^{(i)},\square)$ are defined by the same two sums as above,
that is, as minus the sums of $\mathrm{al}_0(\mathscr{E}^{(i)},X)$ and of $\mathrm{al}_2(\mathscr{E}^{(i)},X)$ over the small polymers $X$ of $T_k(\ell_i)$ that contain $\square$;
for a family as in Proposition~16.2 the matrix $m(\mathscr{E}^{(i)},\square)$ need not be a multiple of $\mathbf{1}$, and both quantities may depend on $\square$.
Let $\tilde{\square}$ be an $M$-cube of $\mathbb{L}_k$.
Suppose that for every $\tilde{X}\in\mathrm{Sm}(\tilde{\square})$ the two numbers $\mathscr{E}^{(i)}_{00}(\mathrm{pr}_i\tilde{X};0,0)$, $i=1,2$, coincide, and that the two kernels $\mathscr{E}^{(i)}_{20}(\mathrm{pr}_i\tilde{X};0,0;x,x')$, $i=1,2$, coincide for all $x$, $x'$ over $\tilde{X}$.
Then
\[
\mathrm{al}_0(\mathscr{E}^{(1)},\mathrm{pr}_1\tilde{X})=\mathrm{al}_0(\mathscr{E}^{(2)},\mathrm{pr}_2\tilde{X}),\qquad
\mathrm{al}_2(\mathscr{E}^{(1)},\mathrm{pr}_1\tilde{X})=\mathrm{al}_2(\mathscr{E}^{(2)},\mathrm{pr}_2\tilde{X})
\]
for every $\tilde{X}\in\mathrm{Sm}(\tilde{\square})$, and
\[
\varepsilon(\mathscr{E}^{(1)},\mathrm{pr}_1\tilde{\square})=\varepsilon(\mathscr{E}^{(2)},\mathrm{pr}_2\tilde{\square}),\qquad
m(\mathscr{E}^{(1)},\mathrm{pr}_1\tilde{\square})=m(\mathscr{E}^{(2)},\mathrm{pr}_2\tilde{\square}).
\]
\end{itemize}
\end{itemize}
\end{named}

\begin{named}{Remark}
Proposition~16.3 corresponds to \cite[Theorem~13.5]{QEDI}, to the first half of \cite[Remark~13.16]{QEDI}, to the definitions of $\mathrm{al}_0$, $\mathrm{al}_2$, $\varepsilon$, $m$ and $R$ in \cite[Section~5.3.5]{QEDI}, and to the parts (f) and (g) of \cite[Proposition~5.59]{QEDI}.
The last sentence of \cite[Theorem~13.5]{QEDI}, which concerns the families without the second slot, corresponds to Proposition~48.2.
The mathematical differences are the following.
\begin{itemize}
\item[(1)] Sites, measure and kernels are those of Proposition~16.2.
The extracted term $m(\mathscr{E})\cdot\langle\bar{\psi},\psi\rangle=m(\mathscr{E})\cdot\int\mathrm{d}x\;\bar{\psi}(x)\cdot\psi(x)$ extends over all $192$ components at a photon site;
this is the mass assignment of \cite[Definition~2.36]{QEDII}, as in Section~I.5.6.
\item[(2)] In the proof of (a) the property (I.k3) enters through the translations $(\mathbf{1},a)$, $a\in M\mathbb{Z}^3$, which do not move the labels, and through the elements $\tilde{g}_z=\left(S,(1-r(S))z\right)$, where $z$ is the centre of $\square$ and $S\in\{\gamma_1,\gamma_2,\gamma_3,\gamma_5\}$;
these elements act on the labels by $r(S)$.
\item[(3)] Part (b) contains in addition the inequalities (2.2) for each single small polymer; in \cite[Theorem~13.5]{QEDI} they occur in the proof of part (b).
\item[(4)] The last sentence of (d) holds for every family with the property (I.k2); it is \cite[Theorem~13.13(c)]{QEDI}.
\item[(5)] Assertion (e) rests, in addition, on the invariance of the domains of the level $k+1$ (Proposition~16.1\,(c) at the level $k+1$).
\item[(6)] Part (f) states the content of the first half of \cite[Remark~13.16]{QEDI} as one formula in the sense of Section~I.1.9.
In (f2) the range is $\mathrm{rg}=\eta_k$, whereas in \cite[Remark~13.16]{QEDI} the range is $1$.
The data of (f2) are the antisymmetric kernels of $\mathscr{E}(X;0,0)$ and of $\mathscr{E}(X;\mathscr{A},\mathscr{C})$, with all arguments over the $\eta_k$-neighbourhood of $\tilde{X}$;
in \cite[Remark~13.16]{QEDI} the datum consists of the kernels of $\mathscr{E}(X;\cdot\,,\cdot)$ at the arguments within distance $1$ of $\tilde{X}$ and of the values of $\mathscr{A}+\mathscr{C}$ on the bonds of $\tilde{X}$,
and these values are needed only for the adjustment and the scaling, which belong to the second half of that remark.
The three maps of (f2) are written out, and (f3) states the comparison of two tori.
As in the converse part of \cite[Theorem~13.2(t10)]{QEDI}, the unions of (f1) are non-empty.
Moreover (f1) contains the identification of $\tilde{X}$ with the lift of $\mathrm{pr}(\tilde{X})$;
this is the identification of the lift of a small polymer with the lift of \cite[Theorem~13.2(t10)]{QEDI} that is asserted in \cite[Section~13.1.3]{QEDI}, and it lies on the way to the independence of the mass of the volume.
\item[(7)] The setting is that of (G2.S-1) and (G2.S-2) in Section~G2.0.4, in which all sides lie in $\mathrm{Cl}$;
the constants that enter are those chosen in (C0)--(C8).
\item[(8)] By the Remark after Proposition~16.2, the matrix $\mathrm{al}_2(\mathscr{E},X)$ of a single small polymer $X$ need not be a multiple of $\mathbf{1}$;
what is a multiple of $\mathbf{1}$ is the sum over the small polymers that contain $\square$.
The reality of $m(\mathscr{E})$ rests on (I.k7) alone, through
\[
m(\mathscr{E})\cdot\mathbf{1}=\varkappa^{-1}\cdot\overline{m(\mathscr{E})}\cdot\varkappa.
\]
\end{itemize}
\end{named}

\subsection*{17. The adjustment}

\subsubsection*{Proposition 17.1 (transports, chains, the homomorphism $\boldsymbol{\Theta}$; (I.SY6)$^{\mathrm{ch}}$)}

\begin{named}{Proposition 17.1}[Transports, chains, the homomorphism $\boldsymbol{\Theta}$; (I.SY6)$^{\mathrm{ch}}$]
Assume (S-1).
Let $X$ be a small polymer with a lift $\tilde{X}$, and let $B$ be a complex bond function on $T_k$; in the scheme, $B=\mathscr{T}$ is the total connection.
Let further $\tilde{g}\in\mathbb{G}_k$, and let $\Lambda$ be a real function of the photon sites.
All points $p$, $p'$, $x$, $\tilde{x}$ that occur below are points of $\mathbb{L}^{\mathrm{f}}_k$ over $\tilde{X}$.
In (u3) and in (e) the points $\tilde{g}p$ lie over $\tilde{g}\tilde{X}$, which by Proposition~16.1\,(a1) is a lift of $\tilde{g}X$.
\begin{itemize}
\item[(a)] For $(\mathscr{A},\mathscr{C})\in\tilde{R}_k\times\mathfrak{X}_k$ and $\mathscr{T}=\mathscr{A}+\mathscr{C}$,
\[
e_k\cdot\left|\operatorname{Im}\mathscr{T}\right|_\infty<e_k^{\frac{1}{4}+\varepsilon}+e_k^{2\varepsilon}\le\frac{1}{2}\,e_k^{\varepsilon},
\]
and for $\mathscr{A}\in\tilde{R}_k$
\[
e_k\cdot\left|\operatorname{Im}\mathscr{A}\right|_\infty<e_k^{\frac{1}{4}+\varepsilon}\le\frac{1}{4}\,e_k^{\varepsilon}.
\]
\item[(b)] (Transports.)
For $(p,p')\in\mathrm{Pzt}_k$ let
\[
u_B(p,p'):=U_B(p,p')=\exp\left\{i\,e_k\cdot\eta_k\cdot(\tau B)(\pi p,\pi p')\right\}
\]
be the transport of (I.1.6), and put $u_B(p,p):=1$.
\begin{itemize}
\item[(u1)] One has
\[
\left|u_B(p,p')^{\pm1}\right|\le\exp\left(e_k\cdot\left|\operatorname{Im}B\right|_\infty\cdot\left|\pi p-\pi p'\right|_1\right),
\]
where $\left|\pi p-\pi p'\right|_1$ is the $\ell^1$ norm of the difference of the two points $\pi p$, $\pi p'$ of $\mathbb{L}_k$.
\item[(u2)] The number $u_B(p,p')$ depends on $B$ only through the values of $B$ on the bonds of the six paths $\boldsymbol{\Gamma}^{\pi}(\pi p,\pi p')$, and it is an entire function of these values.
If $\pi p=\pi p'$, then $u_B(p,p')=1$.
\item[(u3)] One has $u_{\tilde{g}B}(\tilde{g}p,\tilde{g}p')=u_B(p,p')$.
If $B^1=B-d\Lambda$ on the bonds of the six paths, then
\[
u_{B^1}(p,p')=u_B(p,p')\cdot\exp\left\{-i\,e_k\cdot\left(\Lambda(\pi p')-\Lambda(\pi p)\right)\right\}.
\]
Moreover $u_{-B}=u_B^{-1}$ and $\overline{u_{-\overline{B}}}=u_B$.
\item[(u4)] Let $(p,p')\in\mathrm{Pzt}_k$, and let $\psi(p)$, $\bar{\psi}(p)$, $\chi(p,p')$, $\bar{\chi}(p,p')$ be elements of an algebra with
\[
\begin{aligned}
\chi(p,p')&=d_X(p,p')^{-\alpha}\cdot\left(u_B(p,p')\cdot\psi(p')-\psi(p)\right),\\
\bar{\chi}(p,p')&=d_X(p,p')^{-\alpha}\cdot\left(u_B(p,p')^{-1}\cdot\bar{\psi}(p')-\bar{\psi}(p)\right),
\end{aligned}
\]
that is, $\chi=\delta_{\alpha,B}\psi$ and $\bar{\chi}=\delta_{\alpha,-B}\bar{\psi}$.
Then
\begin{equation}\tag{3.1}
\begin{aligned}
\psi(p)&=u_B(p,p')\cdot\psi(p')-d_X(p,p')^{\alpha}\cdot\chi(p,p'),\\
\bar{\psi}(p)&=u_B(p,p')^{-1}\cdot\bar{\psi}(p')-d_X(p,p')^{\alpha}\cdot\bar{\chi}(p,p').
\end{aligned}
\end{equation}
\end{itemize}
\item[(c)] (Chains.)
One has $|X|_M\le c_X\cdot(L+1)$.
Let $x=(c,g)$, let $\tilde{x}$ be a second point, and let $\tilde{c}:=\pi\tilde{x}$ be its photon site.
Let $\mathrm{Ch}(x,\tilde{x})$ denote the set of the shortest chains $(\square_0,\dots,\square_n)$ of $M$-cubes of $X$ that lead from the cube of $c$ to the cube of $\tilde{c}$ and in which any two consecutive cubes have a common face.
This set is finite and not empty, it depends only on the two cubes, and $n+1\le|X|_M$.
For a chain in $\mathrm{Ch}(x,\tilde{x})$ let $z_i$ be the centre of $\square_i$, and put
\[
p_0:=x=(c,g),\qquad p_{i+1}:=(z_i,g)\quad(0\le i\le n),\qquad p_{n+2}:=\tilde{x},
\]
\[
U_i:=\prod_{j<i}u_B(p_j,p_{j+1}),\qquad U^{-}_i:=\prod_{j<i}u_B(p_j,p_{j+1})^{-1}.
\]
\begin{itemize}
\item[(c1)] For every $i$ with $p_i\ne p_{i+1}$ the pair $(p_i,p_{i+1})$ belongs to $\mathrm{Pzt}_k$ and is admissible for $X$ of kind (I).
For $0\le i\le n$ the two points $p_i$, $p_{i+1}$ carry the same label $g$.
One has $|c-z_0|<M/2$, and $d_X(p_0,p_1)=|c-z_0|$ if $p_0\ne p_1$.
For $1\le i\le n$ one has $d_X(p_i,p_{i+1})=M$, and each of the six paths of $(p_i,p_{i+1})$ is the straight segment joining the two centres.
The last pair $(p_{n+1},p_{n+2})=((z_n,g),\tilde{x})$ has arbitrary labels.
If $\tilde{c}=z_n$ and the label of $\tilde{x}$ differs from $g$, this pair lies inside a cluster, and $d_X=\eta_k$, $u_B=1$ for it; if $\tilde{c}=z_n$ and the label of $\tilde{x}$ is $g$, its two points coincide.
\item[(c2)] One has
\[
\sum_{i=0}^{n+1}\left|\pi p_i-\pi p_{i+1}\right|_1\le(n+3)\cdot M\le3\,c_X\cdot L\cdot M.
\]
If $e_k\cdot\left|\operatorname{Im}B\right|_\infty\le\frac{1}{2}\,e_k^{\varepsilon}$, then
\[
\left|U_i\right|,\ \left|U^{-}_i\right|\le\exp\left(\frac{3\,c_X\cdot L\cdot M\cdot e_k^{\varepsilon}}{2}\right)\le\mathrm{e}.
\]
\item[(c3)] One has
\[
\sum_{i=0}^{n+1}d_X(p_i,p_{i+1})^{\alpha}\le(n+2)\cdot M\le3\,c_X\cdot L\cdot M;
\]
here $d_X^{\alpha}\le\max(1,d_X)\le M$, and the terms with $p_i=p_{i+1}$, for which $d_X$ is not defined, are left out, as in (d).
\end{itemize}
\item[(d)] (Definition.)
Let $\operatorname{Mn}$ denote the mean over $\mathrm{Ch}(x,\tilde{x})$.
In the sums below the terms with $p_i=p_{i+1}$ are left out, and the spinor index goes through unchanged.
Define
\begin{equation}\tag{3.2}
\begin{aligned}
\mathrm{tm}(x,\tilde{x})&:=\operatorname{Mn}U_{n+2},&
\mathrm{rm}(x,\tilde{x})[\chi]&:=\operatorname{Mn}\sum_{i=0}^{n+1}U_i\cdot d_X(p_i,p_{i+1})^{\alpha}\cdot\chi(p_i,p_{i+1}),\\
\mathrm{tm}^{-}(x,\tilde{x})&:=\operatorname{Mn}U^{-}_{n+2},&
\mathrm{rm}^{-}(x,\tilde{x})[\bar{\chi}]&:=\operatorname{Mn}\sum_{i=0}^{n+1}U^{-}_i\cdot d_X(p_i,p_{i+1})^{\alpha}\cdot\bar{\chi}(p_i,p_{i+1}),
\end{aligned}
\end{equation}
and
\begin{equation}\tag{3.3}
\begin{aligned}
\boldsymbol{\Theta}\psi(x)&:=\operatorname{Vol}(X)^{-1}\int_X\mathrm{d}\tilde{x}\,\left(\mathrm{tm}(x,\tilde{x})\cdot\psi(\tilde{x})-\mathrm{rm}(x,\tilde{x})[\chi]\right),\\
\boldsymbol{\Theta}\bar{\psi}(x)&:=\operatorname{Vol}(X)^{-1}\int_X\mathrm{d}\tilde{x}\,\left(\mathrm{tm}^{-}(x,\tilde{x})\cdot\bar{\psi}(\tilde{x})-\mathrm{rm}^{-}(x,\tilde{x})[\bar{\chi}]\right).
\end{aligned}
\end{equation}
By (I.1.3) the integral $\int_X\mathrm{d}\tilde{x}$ contains the factor $\frac{1}{48}$ and the sum over the label of $\tilde{x}$; the definition (3.3) is (I.5.3).
There is exactly one homomorphism $\boldsymbol{\Theta}=\boldsymbol{\Theta}_{X,B}$ from the subalgebra of $\mathrm{Gr}[\psi,\chi]$ generated by the $\psi(\xi)$, $\xi$ over $X$, into $\mathrm{Gr}[\psi,\chi]$ that satisfies (3.3).
It does not depend on the lift.
Indeed, a translation of $\tilde{X}$ by $\ell w$, $w\in\mathbb{Z}^3$, multiplies $\psi(x)$, $\psi(\tilde{x})$ and $\chi(p_i,p_{i+1})$ by the common sign $(-1)^{w_1+w_2+w_3}$, and it leaves the transports, the distances, the chains and $\operatorname{Vol}(X)$ unchanged.
The homomorphism $\boldsymbol{\Theta}$ maps every unbarred generator to a linear combination of unbarred generators, and every barred generator to a linear combination of barred ones.
\item[(e)] On the subalgebras on which both sides are defined,
\begin{equation}\tag{3.4}
\begin{aligned}
a_{\tilde{g}}\circ\boldsymbol{\Theta}_{\tilde{g}X,\tilde{g}B}&=\boldsymbol{\Theta}_{X,B}\circ a_{\tilde{g}},\\
G_{\Lambda}\circ\boldsymbol{\Theta}_{X,B^1}&=\boldsymbol{\Theta}_{X,B}\circ G_{\Lambda}\qquad\text{if $B^1=B-d\Lambda$ on the bonds with both ends in $X$},\\
\mathrm{Cg}\circ\boldsymbol{\Theta}_{X,-B}&=\boldsymbol{\Theta}_{X,B}\circ\mathrm{Cg},\\
\mathbf{K}\circ\boldsymbol{\Theta}_{X,-\overline{B}}&=\boldsymbol{\Theta}_{X,B}\circ\mathbf{K}.
\end{aligned}
\end{equation}
\item[(f)] Let $K(x,x')$, for $x$, $x'$ over $X$, be a kernel with values in the $4\times4$ matrices; all integrals below extend over $X$.
Then
\[
\boldsymbol{\Theta}\left(\int\mathrm{d}x\int\mathrm{d}x'\;\bar{\psi}(x)^T\cdot K(x,x')\cdot\psi(x')\right)
\]
is the sum of the four terms
\begin{align*}
&\operatorname{Vol}(X)^{-2}\int\mathrm{d}\tilde{x}'\int\mathrm{d}\tilde{x}\;\bar{\psi}(\tilde{x}')^T\cdot\left[\int\mathrm{d}x\int\mathrm{d}x'\;\mathrm{tm}^{-}(x,\tilde{x}')\cdot K(x,x')\cdot\mathrm{tm}(x',\tilde{x})\right]\cdot\psi(\tilde{x}),\\
&-\operatorname{Vol}(X)^{-2}\int\mathrm{d}x\,\mathrm{d}x'\,\mathrm{d}\tilde{x}'\,\mathrm{d}\tilde{x}\;\bar{\psi}(\tilde{x}')^T\cdot\mathrm{tm}^{-}(x,\tilde{x}')\cdot K(x,x')\cdot\mathrm{rm}(x',\tilde{x})[\chi],\\
&-\operatorname{Vol}(X)^{-2}\int\mathrm{d}x\,\mathrm{d}x'\,\mathrm{d}\tilde{x}'\,\mathrm{d}\tilde{x}\;\mathrm{rm}^{-}(x,\tilde{x}')[\bar{\chi}]^T\cdot K(x,x')\cdot\mathrm{tm}(x',\tilde{x})\cdot\psi(\tilde{x}),\\
&+\operatorname{Vol}(X)^{-2}\int\mathrm{d}x\,\mathrm{d}x'\,\mathrm{d}\tilde{x}'\,\mathrm{d}\tilde{x}\;\mathrm{rm}^{-}(x,\tilde{x}')[\bar{\chi}]^T\cdot K(x,x')\cdot\mathrm{rm}(x',\tilde{x})[\chi];
\end{align*}
in the last three terms the integration extends over the four points $x$, $x'$, $\tilde{x}'$, $\tilde{x}$.
At $B=0$ the first term equals
\[
\bar{\psi}_{\mathrm{av}}^T\cdot\left[\int\mathrm{d}x\int\mathrm{d}x'\;K(x,x')\right]\cdot\psi_{\mathrm{av}},\qquad
\psi_{\mathrm{av}}:=\operatorname{Vol}(X)^{-1}\int_X\mathrm{d}x\;\psi(x),
\]
with $\bar{\psi}_{\mathrm{av}}$ formed in the same way from $\bar{\psi}$; these are means over the sites and over the labels.
\end{itemize}
\end{named}

\begin{named}{Remark}
Proposition~17.1 corresponds to the opening paragraphs of \cite[Section~13.5]{QEDI}, whose first display is the first bound of (a), to \cite[Definition~13.6]{QEDI}, to the four relations displayed in the proof of \cite[Theorem~13.7]{QEDI}, which are (3.4), and to the definition of the homomorphism in \cite[Section~5.3.5]{QEDI}.
It differs from them in the following points.
In \cite[Section~13.5]{QEDI} the chain point $p_{i+1}$ is the centre of $\square_i$.
Here the centre $z_i$ is a photon site and not a fermion site, and the chain point is $p_{i+1}=(z_i,g)$, the fermion site over $z_i$ that carries the label $g$ of the starting point, as in (I.SY6)$^{\mathrm{ch}}$.
The distance $d_X(p,p')$ stands in place of $|p-p'|$ in (u4), in (3.1), in (3.2) and in (c3).
Pairs inside a cluster are admitted, in the last link of a chain, as stated in (c1).
The mean over $\tilde{x}$ in (3.3) includes the mean over the 48 labels, so that no label is distinguished.
The homomorphism $\boldsymbol{\Theta}_{X,B}$ and the relations (3.4) are formulated for an arbitrary complex bond function $B$; the second slot enters only through $B=\mathscr{A}+\mathscr{C}$.
Finally, the independence of $\boldsymbol{\Theta}$ of the lift is part of the statement (d).
\end{named}

\subsubsection*{Proposition 17.2 (adjustment)}

\begin{named}{Proposition 17.2}[Adjustment]
Assume (S-1).
Let $X$ be small, let $(\mathscr{A},\mathscr{C})\in\tilde{R}_k\times\mathfrak{X}_k$ and $\mathscr{T}:=\mathscr{A}+\mathscr{C}$, and let $E$ be a family of level $k$ (G2.r8) with (k2).
For small $X$ define
\[
E^{\mathrm{nat}}(X;\mathscr{A},\mathscr{C}):=E(X;\mathscr{A},\mathscr{C})-E^{[2,0]}(X;\mathscr{A},\mathscr{C})+\boldsymbol{\Theta}_{X,\mathscr{T}}\left(E^{[2,0]}(X;\mathscr{A},\mathscr{C})\right),
\]
and for large $X$ put $E^{\mathrm{nat}}(X):=E(X)$.
\begin{itemize}
\item[(a)] (Identity.)
Let $\mathrm{ev}$ be a homomorphism of $\mathrm{Gr}[\psi,\chi]$ into an associative algebra, and assume that the images of $\psi$, $\bar{\psi}$, $\chi$, $\bar{\chi}$ satisfy the relations of (u4) with $B=\mathscr{T}$ at all pairs of kind (I) for $X$, the pairs inside the clusters over $X$ included.
Then
\[
\mathrm{ev}\left(E^{\mathrm{nat}}(X;\mathscr{A},\mathscr{C})\right)=\mathrm{ev}\left(E(X;\mathscr{A},\mathscr{C})\right).
\]
In particular $E^{\mathrm{nat}}(X;\mathscr{A},\mathscr{C},\psi,\delta_{\alpha,\mathscr{T}}\psi)=E(X;\mathscr{A},\mathscr{C},\psi,\delta_{\alpha,\mathscr{T}}\psi)$, where on the barred fields $\delta_{\alpha,-\mathscr{T}}$ is inserted.
\item[(b)] Let $E$ be balanced.
If
\[
\int\mathrm{d}x\int\mathrm{d}x'\;E_{20}(X;0,0;x,x')=0,
\]
then $E^{\mathrm{nat}}_{20}(X;0,0;\cdot,\cdot)=0$ as a kernel on $(X\times O_{\mathrm{cub}})^2$.
For $(n,m)\notin\{(2,0),(1,1),(0,2)\}$ one has $E^{\mathrm{nat}}_{nm}=E_{nm}$; in particular $E^{\mathrm{nat}}_{00}=E_{00}$.
\item[(c)] (Locality, leg by leg.)
The element $E^{\mathrm{nat}}(X;\mathscr{A},\mathscr{C})$ contains generators $\psi(\xi)$ only with $\xi$ over $X$, and generators $\chi(\boldsymbol{\zeta})$ only with $\boldsymbol{\zeta}$ admissible for $X$.
Let $\chi(\boldsymbol{\zeta})$, $\boldsymbol{\zeta}=(p,p',\beta,\omega)$, be a generator that occurs in a monomial of $E^{\mathrm{nat}}(X;\mathscr{A},\mathscr{C})$ with a non-zero coefficient and does not occur in $E(X;\mathscr{A},\mathscr{C})$.
Then $(p,p')=(p_i,p_{i+1})$, $p_i\ne p_{i+1}$, is a chain pair of Proposition~17.1\,(c); hence it is of kind (I) for $X$, $d_X(p,p')\le M$, and $p$, $p'$ have the same label unless $(p,p')$ is the last pair of its chain.
The coefficients of $E^{\mathrm{nat}}(X;\mathscr{A},\mathscr{C})$ depend on $(\mathscr{A},\mathscr{C})$ only through the kernels $E_{nm}(X;\mathscr{A},\mathscr{C})$ and through the values of $\mathscr{T}$ on the bonds both of whose ends lie in $X$.
Consequently $E^{\mathrm{nat}}$ has (k2).
\item[(d)] (Bounds.)
Let $E$ be balanced.
With all kernels taken at $(\mathscr{A},\mathscr{C})$,
\begin{align*}
\left\|E^{\mathrm{nat}}_{20}(X)\right\|&\le\mathrm{e}^2\cdot\left\|E_{20}(X)\right\|,\\
\left\|E^{\mathrm{nat}}_{11}(X)\right\|&\le\left\|E_{11}(X)\right\|+\mathrm{e}^2\cdot\left(3\,c_X\cdot L\cdot M\right)\cdot\left\|E_{20}(X)\right\|,\\
\left\|E^{\mathrm{nat}}_{02}(X)\right\|&\le\left\|E_{02}(X)\right\|+\mathrm{e}^2\cdot\left(3\,c_X\cdot L\cdot M\right)^2\cdot\left\|E_{20}(X)\right\|,
\end{align*}
and
\begin{align*}
\left\|E^{\mathrm{nat}}(X;\mathscr{A},\mathscr{C})\right\|_{h_k}&\le\mathrm{e}^2\cdot\left(1+3\,c_X\cdot L\cdot M\cdot e_k^{\varepsilon}\right)^2\cdot\left\|E(X;\mathscr{A},\mathscr{C})\right\|_{h_k}\\
&\le10\cdot\left\|E(X;\mathscr{A},\mathscr{C})\right\|_{h_k},
\end{align*}
where $3\,c_X\cdot L\cdot M\cdot e_k^{\varepsilon}\le\frac{1}{30}$.
\item[(e)] (Symmetries.)
If $E$ lies in $K^{x,C}_k$, then so does $E^{\mathrm{nat}}$; if $E$ lies in $K^{x}_k$, then so does $E^{\mathrm{nat}}$.
One has $\left\|E^{\mathrm{nat}}\right\|^{x}_k\le10\cdot\left\|E\right\|^{x}_k$.
The map $E\mapsto E^{\mathrm{nat}}$ is linear, and $E^{\mathrm{nat}}|_0=(E|_0)^{\mathrm{nat}}$, where on the right the adjustment is that of Proposition~48.3.
\item[(f)] (One formula; the class of sides $\mathrm{Cl}$.)
In the sense of Section~I.1.9, $\boldsymbol{\Theta}_{X,B}$ ($X$ small) is given by one formula of range $\mathrm{rg}=0$, the data being the values of $B$ on the bonds with both ends in $\tilde{X}$.
The corresponding map assigns to $(\tilde{X},\tilde{B})$ the homomorphism (3.3) formed with the cubes of $\tilde{X}$, their centres and common faces, the labels, the measure (I.1.3) and the transports; it commutes with the translations by the vectors of $M\mathbb{Z}^3$.
The element $E^{\mathrm{nat}}(X;\mathscr{A},\mathscr{C})$ is given by one formula of range $\mathrm{rg}=\eta_k$ for every polymer $X$ of $T_k$, small or large, that has a lift $\tilde{X}$ whose sup-diameter is less than $\ell-2\eta_k-M$; the data are the kernels of $E(X;\mathscr{A},\mathscr{C})$ and the values of $\mathscr{T}$ on the bonds with both ends in $\tilde{X}$.
The map is
\[
F\mapsto F-F^{[2,0]}+\boldsymbol{\Theta}_{\tilde{X},\tilde{B}}\left(F^{[2,0]}\right)\quad\text{if }d_M(\tilde{X})\le L,\qquad F\mapsto F\quad\text{otherwise},
\]
where $F$ is the element with these kernels, $\tilde{B}$ is the lift of $\mathscr{T}$, and $\boldsymbol{\Theta}_{\tilde{X},\tilde{B}}$ is the homomorphism (3.3) for $(\tilde{X},\tilde{B})$.
As a consequence, the following holds.
Let $\ell_1,\ell_2\in\mathrm{Cl}$, let $T_k(\ell)$ denote $T_k$ formed with the side $\ell$, and let $\mathrm{pr}_i\colon\mathbb{L}_k\to T_k(\ell_i)$ be the covering map.
Let $\tilde{X}\in\mathrm{Sm}(\tilde{\square})$ (Proposition~16.3\,(f1)), and let $B_i$ be bond functions on $T_k(\ell_i)$ whose lifts coincide on the bonds with both ends in $\tilde{X}$.
Then, for every $x$ over $\tilde{X}$, the elements $\boldsymbol{\Theta}_{\mathrm{pr}_1\tilde{X},B_1}\psi(x,\cdot)$ and $\boldsymbol{\Theta}_{\mathrm{pr}_2\tilde{X},B_2}\psi(x,\cdot)$ are the same linear combination of the generators $\psi(\tilde{x},\cdot)$, $\chi(p,p',\cdot)$ with indices over $\tilde{X}$; the same holds for $\bar{\psi}$.
\end{itemize}
\end{named}

\begin{named}{Remark}
Proposition~17.2 corresponds to \cite[Theorem~13.7]{QEDI} and to \cite[Proposition~5.60]{QEDI}.
The definition of $E^{\mathrm{nat}}$ is that of \cite[Definition~13.6]{QEDI} and of the display defining the adjustment in \cite[Section~5.3.5]{QEDI}; clause (c) corresponds also to \cite[Theorem~13.13(d)]{QEDI}, and clause (f) to the second half of \cite[Remark~13.16]{QEDI}.
The differences in the homomorphism $\boldsymbol{\Theta}$ are those listed after Proposition~17.1.
In (b) the vanishing holds as a kernel on $(X\times O_{\mathrm{cub}})^2$.
By Proposition~17.1\,(f) the part of bidegree $(2,0)$ of $E^{\mathrm{nat}}(X;0,0)$ is
\[
\bar{\psi}_{\mathrm{av}}^T\cdot\left[\int\mathrm{d}x\int\mathrm{d}x'\;E_{20}(X;0,0;x,x')\right]\cdot\psi_{\mathrm{av}},
\]
with means over the sites and over the labels; a bilinear inside a cluster whose mean over both labels vanishes is therefore turned into terms with at least one H\"older leg ((I.E2); Proposition~18.6).
Clauses (a) and (c) hold for every family with (k2), clauses (b) and (d) for every balanced family with (k2); \cite[Theorem~13.7]{QEDI} is stated for $E\in K^{x,C}_k$, while \cite[Theorem~13.13(d)]{QEDI} asserts the locality of $E^{\mathrm{nat}}$ for every family with (k2).
In (c) the new H\"older legs are located at the chain pairs, as in \cite[Theorem~13.7(c)]{QEDI}.
Of the assertions of (c) on these legs, only the bound $d_X(p,p')\le M$ and the clause on the labels of $p$, $p'$ have no counterpart in \cite[Theorem~13.7(c)]{QEDI}; the first of these takes the place of the bound $M$ on the length of these legs in \cite[Theorem~13.13(d)]{QEDI}.
The forms without the second slot are Proposition~48.3.
\end{named}

\subsection*{18. Global fields, locality, scaling, contraction (continued)}

\subsubsection*{Proposition 18.2 (global fields on a small polymer)}

\begin{named}{Proposition 18.2}[Global fields on a small polymer]
Assume (S-2).
Let $X\in D_k$ be small, let $\tilde{c}$ be a photon site of $X$, let $(\mathscr{A},\mathscr{C})\in\tilde{R}_{k+1}\times\mathfrak{X}_{k+1}$, and put $\mathrm{rad}:=2L\cdot M$.
There exist complex bond functions $\mathrm{bh}$, $\mathscr{C}^{b}$ on $T_k$, a real function $\mathrm{mg}$ of the sites of $T_k$ and constant complex bond functions $\mathrm{kc}_1$, $\mathrm{kc}_2$ such that the following holds.
\begin{itemize}
\item[(o)] (What they are made of.)
The functions $\mathrm{bh}$, $\mathrm{kc}_1$, $\mathrm{mg}$ are constructed from $\mathscr{A}$ and $\tilde{c}$ alone, through a splitting of $\mathscr{A}$ and a representative on $L\square'^{\dagger}$, where $L\square'$ is the $LM$-cube of $\tilde{c}$; the functions $\mathscr{C}^{b}$, $\mathrm{kc}_2$ are constructed from $\mathscr{C}$ and $\tilde{c}$ alone.
None of them depends on a family of activities.
The constant field $\mathrm{kc}_2$ has the components $(\mathscr{C}_L)_\mu(\tilde{c})$.
If $\mathscr{C}=0$, then $\mathscr{C}^{b}=0$ and $\mathrm{kc}_2=0$.
\item[(i)] On the bonds with both ends in $K(\tilde{c},\mathrm{rad})$,
\[
\mathscr{A}_L-d\,\mathrm{mg}=\mathrm{bh}+\mathrm{kc}_1,\qquad\mathscr{C}_L=\mathscr{C}^{b}+\mathrm{kc}_2;
\]
every bond that has an end at distance at most $1$ from $X$ is among these bonds.
\item[(ii)] The three norms (G2.r4) of $\mathrm{bh}$ obey
\[
\begin{aligned}
\left|\mathrm{bh}\right|_\infty&\le1382\,L^{-\frac{15}{8}+\frac{3\varepsilon}{2}}\cdot\mathrm{t1}_k,\\
\left|\partial\,\mathrm{bh}\right|_\infty&\le1382\,L^{-\frac{15}{8}+\frac{3\varepsilon}{2}}\cdot\mathrm{t2}_k,\\
\left|\delta_\alpha\partial\,\mathrm{bh}\right|_\infty&\le1382\,L^{-\frac{15}{8}+\frac{3\varepsilon}{2}}\cdot\mathrm{t3}_k,
\end{aligned}
\]
and
\[
\left|\mathscr{C}^{b}\right|_\infty<L^{-\frac{7}{4}}\cdot\mathrm{td1}_k,\qquad\left|\partial\mathscr{C}^{b}\right|_\infty<L^{-\frac{7}{4}}\cdot\mathrm{td2}_k;
\]
hence $(t\cdot\mathrm{bh},t\cdot\mathscr{C}^{b})\in R_k\times\mathfrak{X}_k$ for every complex number $t$ with $|t|\le L^{\frac{7}{4}}$.
Moreover
\[
\left|\mathrm{kc}_1\right|\le L^{-\frac{7}{8}+\frac{\varepsilon}{2}}\cdot\mathrm{t1}_k,\qquad\left|\mathrm{kc}_2\right|\le L^{-1+\varepsilon}\cdot\mathrm{td1}_k.
\]
\item[(iii)] Let $E$ be a family of level $k$ with (k1), (k2), (k4).
For all $(n,m)$ and all arguments $\xi=(\xi_1,\dots,\xi_n)$, $\boldsymbol{\zeta}=(\boldsymbol{\zeta}_1,\dots,\boldsymbol{\zeta}_m)$,
\[
\begin{aligned}
\left|E_{nm}(X;\mathscr{A}_L,\mathscr{C}_L;\xi,\boldsymbol{\zeta})\right|&\le\exp\left(2(n+m)\right)\cdot\left|E_{nm}(X;\mathrm{bh},\mathscr{C}^{b};\xi,\boldsymbol{\zeta})\right|,\\
E_{00}(X;\mathscr{A}_L,\mathscr{C}_L)&=E_{00}(X;\mathrm{bh},\mathscr{C}^{b}).
\end{aligned}
\]
The modulus on the left is the modulus on the right times a product of factors, one for each argument and each slot, each of modulus at most $\mathrm{e}$; the phase of an argument $\xi_j=(x,\cdot)$, resp.\ $\boldsymbol{\zeta}_j=(x,x',\cdot)$, is read at the photon site $c=\pi(x)$ of its site, resp.\ of its first site.
Consequently
\[
\left\|E_{nm}(X;\mathscr{A}_L,\mathscr{C}_L)\right\|\le\exp\left(2(n+m)\right)\cdot\left\|E_{nm}(X;\mathrm{bh},\mathscr{C}^{b})\right\|.
\]
\end{itemize}
\end{named}

\begin{named}{Remark}
Proposition~18.2 corresponds to \cite[Theorem~13.11]{QEDI} and to \cite[Proposition~5.61]{QEDI}.
Clauses (i) and (ii) concern the photon field only and coincide with (i), (ii) of \cite[Theorem~13.11]{QEDI}.
Clause (o) corresponds to the sentences that precede (i) in \cite[Theorem~13.11]{QEDI} and in \cite[Proposition~5.61]{QEDI}, where the base site is a site $x_0$ of $X$; here it is a photon site $\tilde{c}$ of $X$.
On this point the statements agree: $\mathrm{bh}$, $\mathrm{kc}_1$, $\mathrm{mg}$ are formed from $\mathscr{A}$ and the base site alone, $\mathscr{C}^{b}$ and $\mathrm{kc}_2$ from $\mathscr{C}$ and the base site alone, and $\mathscr{C}=0$ gives $\mathscr{C}^{b}=0$ and $\mathrm{kc}_2=0$.
Like (o), \cite[Theorem~13.11]{QEDI} states that the choice does not depend on the families $E$ of (iii) and that $\mathrm{kc}_2$ has the components $(\mathscr{C}_L)_\mu$ at the base site.
Clause (o) names in addition how $\mathrm{bh}$, $\mathrm{kc}_1$, $\mathrm{mg}$ are built, namely through a splitting of $\mathscr{A}$ and a representative on $L\square'^{\dagger}$.
In \cite[Theorem~13.11]{QEDI} a formula for $\mathscr{C}^{b}$ is given in addition: $\mathscr{C}^{b}$ is obtained from $\mathscr{C}_L-\mathrm{kc}_2$ by the truncation of \cite[Lemma~13.9]{QEDI}, formed with the base site and with $r=\mathrm{rad}$.
For the pair $(\mathscr{A},0)$ clause (o) thus yields $\mathscr{C}^{b}=0$ and $\mathrm{kc}_2=0$, while $\mathrm{bh}$, $\mathrm{kc}_1$, $\mathrm{mg}$ depend on $(\mathscr{A},\tilde{c})$ only;
these are the objects of Proposition~48.4\,(a), (b), the counterparts of (i)--(iii) for the families without the second slot.
In (iii) the arguments are indices over $T^{\mathrm{f}}_k$ and $\mathrm{Pz}_k$, and the phases are read at the photon site $c$ of the site, resp.\ of the first site, of an argument.
By (k2) and Section~I.1.7 this photon site has distance at most $\eta_k$ from $X$ (for a pair of kind (II) one has $x'=x^{(3)}$), so that $|c-\tilde{c}|_1\le9L\cdot M$, which is the bound used in the proof of \cite[Theorem~13.11]{QEDI}.
\end{named}

\subsubsection*{Proposition 18.3 (stability of the locality clause)}

\begin{named}{Proposition 18.3}[Stability of the locality clause]
The data of the definition below and of (a), (b) are an integer $k\ge0$, a side $\ell\in\mathrm{Cl}$, a number $e_k>0$, and unions $Z$, $Z'$, $Z_i$ of $M$-cubes of $T_k$.
These data concern one level; the definition and (a), (b) therefore hold also for $(k+1,\ell/L,e_{k+1})$ in the setting (S-2).
Besides $\psi$ and $\chi$, any finite families of further independent generators may be present; they are called spectators (in the scheme, $\Psi$ and $W$).
For (c) and (d) the setting is (S-1), as in Propositions~16.3 and~17.2; for (e) it is (S-2).

\emph{Definition.}
Let $\mathrm{Gr}_k[Z]$ be the subalgebra that contains $1$ and is generated by the $\psi(\xi)$ with $\xi$ over $Z$, the $\chi(\boldsymbol{\zeta})$ with $\boldsymbol{\zeta}$ admissible for $Z$ (Section~I.1.7), and the spectators.
Let $\mathrm{Dc}$ be a set of pairs of complex bond functions on $T_k$.
A map $F$ from $\mathrm{Dc}$ into the algebra is called \emph{local in $Z$} if
\begin{itemize}
\item[(i)] $F(\mathscr{A},\mathscr{C})\in\mathrm{Gr}_k[Z]$ for all $(\mathscr{A},\mathscr{C})\in\mathrm{Dc}$, and
\item[(ii)] $F(\mathscr{A},\mathscr{C})=F(\tilde{\mathscr{A}},\tilde{\mathscr{C}})$ whenever $(\mathscr{A},\mathscr{C})$ and $(\tilde{\mathscr{A}},\tilde{\mathscr{C}})$ lie in $\mathrm{Dc}$ and coincide, in both components, on the bonds that meet $Z$.
\end{itemize}
In these terms (k2) for a family $\mathscr{E}$ says that $\mathscr{E}(X;\cdot,\cdot)$ is local in $X$ for every $X\in D_k$, with $\mathrm{Dc}=\tilde{R}_k\times\mathfrak{X}_k$ and without spectators.
\begin{itemize}
\item[(a)] (Monotonicity.)
Let $Z\subset Z'$.
A pair admissible for $Z$ is admissible for $Z'$, of the same kind; a bond that meets $Z$ meets $Z'$; $\mathrm{Gr}_k[Z]\subset\mathrm{Gr}_k[Z']$; and a map local in $Z$ is local in $Z'$.
\item[(b)] (Algebra.)
On a fixed set $\mathrm{Dc}$ the maps local in $Z$ form an algebra, which is closed under pointwise limits.
This algebra is stable under $\exp$ and under every norm-convergent series; under derivatives and integrals with respect to auxiliary parameters, whenever these exist, in particular under $\partial/\partial s_P$, $\int_{[0,1]^P}\mathrm{d}s_P$, $\oint_{\mathrm{rad}}$ and $\int\mathrm{d}\hat{\mu}(\mathsf{Z})$; and under $\int\mathrm{d}\mu_I(W)$ if $W$ is one of the spectators.
Here $P$ is a finite set of parameter cubes (Section~I.1.1), and
\[
s_P:=\left(s_{\mathrm{pc}}\right)_{\mathrm{pc}\in P},\qquad
\mathrm{d}s_P:=\prod_{\mathrm{pc}\in P}\mathrm{d}s_{\mathrm{pc}},\qquad
\frac{\partial}{\partial s_P}:=\prod_{\mathrm{pc}\in P}\frac{\partial}{\partial s_{\mathrm{pc}}};
\]
thus $\partial/\partial s_P$ is the iterated partial derivative in the parameters $s_{\mathrm{pc}}$, $\mathrm{pc}\in P$, and $\int_{[0,1]^P}\mathrm{d}s_P$ is the integral over the cube $[0,1]^P$ with respect to Lebesgue measure, the other parameters being fixed.
In the scheme $P$ is a set of parameter cubes as in Definition~24.D\,(D0), and these are the two operations of the display (8.2) in Definition~24.D\,(D3); they are defined in the same way in \cite[Section~17.4.3]{QEDI}.
If $F_i$ is local in $Z_i$ for $i=1,\dots,n$, then $\prod_iF_i$ is local in $Z_1\cup\dots\cup Z_n$.
\item[(c)] If a family has (k2), then so has its image under $R$ (Proposition~16.3\,(d)).
\item[(d)] If a family has (k2), then so has its image under the adjustment $E\mapsto E^{\mathrm{nat}}$ (Proposition~17.2\,(c)).
\item[(e)] (Reblocking, scaling.)
Let $\mathrm{Dc}$ be a set of pairs of complex bond functions on $T_k$, and let $\mathrm{Dc}'$ be a set of pairs of complex bond functions on $T_{k+1}$ such that $(\mathscr{A}_L,\mathscr{C}_L)\in\mathrm{Dc}$ for every $(\mathscr{A},\mathscr{C})\in\mathrm{Dc}'$.
By Proposition~18.1\,(b), (c) this condition is satisfied by each of the pairs
\[
\begin{aligned}
(\mathrm{Dc},\mathrm{Dc}')&=\left(\tilde{R}_k\times\mathfrak{X}_k,\ \tilde{R}_{k+1}\times\mathfrak{X}_{k+1}\right),\\
(\mathrm{Dc},\mathrm{Dc}')&=\left(\tilde{R}_k^{-}\times\mathfrak{X}_k,\ \tilde{R}_{k+1}\times\mathfrak{X}_{k+1}\right),\\
(\mathrm{Dc},\mathrm{Dc}')&=\left(\tilde{R}_k\times\{0\},\ \tilde{R}_{k+1}\times\{0\}\right).
\end{aligned}
\]
Let $Y\in D_{k+1}$, let $X\in D_k$ lie in $LY$, and let $F$ be local in $X$ on $\mathrm{Dc}$, without spectators.
Then the map $(\mathscr{A},\mathscr{C})\mapsto S_L\left[F(\mathscr{A}_L,\mathscr{C}_L)\right]$ is local in $Y$ on $\mathrm{Dc}'$, at level $k+1$.
Under $S_L$ a leg at a pair of kind (I) for $X$ with $d_X=d$ becomes a leg at a pair of kind (I) for $Y$ with $d_X=d/L$; a pair inside a cluster, with $d=\eta_k$, goes to a pair inside a cluster, with $d/L=\eta_{k+1}$.
A leg at a pair $(x,x^{(3)})$ of kind (II) for $X$ becomes a leg at $(x/L,(x/L)^{(3)})$, which is of kind (II) for $Y$.
Consequently, if for every $X\in D_k$ the map $E'(X;\cdot,\cdot)$ is local in $X$ on $\mathrm{Dc}$, without spectators, then for every $Y\in D_{k+1}$ the map
\[
(\mathscr{A},\mathscr{C})\mapsto\sum_{X:\ \bar{X}^{L}=LY}S_L\left[E'(X;\mathscr{A}_L,\mathscr{C}_L)\right]
\]
is local in $Y$ on $\mathrm{Dc}'$.
In particular, if the family $E'$ has (k2) at level $k$, then $L\,E'$ has (k2) at level $k+1$.
\end{itemize}
\end{named}

\begin{named}{Remark}
Proposition~18.3 corresponds to \cite[Definition~13.12]{QEDI}, to \cite[Theorem~13.13(a)--(e)]{QEDI} and to \cite[Proposition~5.67]{QEDI}.
The kinds (I), (II) are those of Section~I.1.7: a pair of kind (II) is a pair $(x,x^{(3)})$, where \cite[Section~13.7]{QEDI} has a pair of nearest neighbours in $T_k$; one has $x^{(3)}/L=(x/L)^{(3)}$, and $d_X$ stands in place of the length of a pair.
The definition and (a), (b) concern a single torus, so that (b) applies at level $k+1$ as well.
The spectators are arbitrary finite families of independent generators.
The derivatives and integrals in (b) are those that exist.
For them the stability rests on two facts: $\mathrm{Gr}_k[Z]$ is a closed linear subspace of the algebra, and two points of $\mathrm{Dc}$ that coincide on the bonds meeting $Z$ give the same function of the auxiliary parameter.
Hence such a derivative or integral of a map local in $Z$ takes its values in $\mathrm{Gr}_k[Z]$ and satisfies condition (ii) of the definition.
In \cite[Theorem~13.13]{QEDI} clause (e) is stated for the pair $(\mathrm{Dc},\mathrm{Dc}')=\left(\tilde{R}_k\times\mathfrak{X}_k,\ \tilde{R}_{k+1}\times\mathfrak{X}_{k+1}\right)$, and the last assertion of that theorem gives it for the pair $\left(\tilde{R}_k\times\{0\},\ \tilde{R}_{k+1}\times\{0\}\right)$, a family without the second slot being regarded there as the family of the maps $(\mathscr{A},0)\mapsto E(X;\mathscr{A})$; these are the first and the third pair of the display in clause (e) of Proposition~18.3.
Proposition~18.3\,(e) differs from \cite[Theorem~13.13(e)]{QEDI} in that $\mathrm{Dc}$, $\mathrm{Dc}'$ are arbitrary subject to the inclusion required there; this covers in particular the second pair $\left(\tilde{R}_k^{-}\times\mathfrak{X}_k,\ \tilde{R}_{k+1}\times\mathfrak{X}_{k+1}\right)$.
The second and the third pair in (e) are formed with the following domains.
The functions $E^{\#,x}_k$ of the Propositions~24.m are defined on $\tilde{R}_k^{-}\times\mathfrak{X}_k$, with the background in $\tilde{R}_k^{-}$ and the second slot in $\mathfrak{X}_k$ (Definition~24.D\,(D1)); this set is $\mathrm{Dc}$ for the second pair.
For a family $E$ without the second slot, Proposition~48.4\,(c) reads $E(X;\cdot)$ as the map $(\mathscr{A},0)\mapsto E(X;\mathscr{A})$ on $\tilde{R}_k\times\{0\}$; this set is $\mathrm{Dc}$ for the third pair.
\end{named}

\subsubsection*{Proposition~18.4 (scaling of norms)}

\begin{named}{Proposition 18.4}[Scaling of norms]
Let $k\ge0$ and $\ell\in\mathrm{Cl}$, so that $S_L$ is defined (Proposition~16.1\,(d)).
Let $e_k>0$ and $e_{k+1}=L^{\frac{1}{2}}\cdot e_k$, and let $F$ be an element of $\mathrm{Gr}[\psi,\chi]$ of level $k$.
The weights of the two levels are related by
\[
h_{k+1,1}=L^{-\frac{1}{8}}\cdot h_{k,1},\qquad h_{k+1,2}=L^{-\frac{1}{8}+\frac{\varepsilon}{2}}\cdot h_{k,2}.
\]
\begin{itemize}
\item[(a)] The map $x\mapsto x/L$ is a bijection of $T^{\mathrm{f}}_k$ onto $T^{\mathrm{f}}_{k+1}$.
The map $(x,x')\mapsto(x/L,x'/L)$ is a bijection of $\mathrm{Pz}_k$ onto the set of the pairs of $\mathrm{Pz}_{k+1}$ with $d_X\le M/L$; every pair inside a cluster belongs to this set.
The homomorphism $S_L$ maps the basis monomials of level $k$ injectively to multiples of basis monomials of level $k+1$, and on a monomial of bidegree $(n,m)$ the factor is $\mathrm{sn}\cdot L^{-n-(1+\alpha)m}$ with a sign $\mathrm{sn}\in\{1,-1\}$.
If the fundamental domain of level $k+1$ is the image of that of level $k$ under $c\mapsto c/L$, and if the maps $\xi\mapsto\xi/L$ and $\boldsymbol{\zeta}\mapsto\boldsymbol{\zeta}/L$ are increasing for the total orders of the index sets of the two levels (Section~I.3.6, (I.AP), and Section~I.5.1), then $\mathrm{sn}=1$ for all monomials.
The assertions (b) and (c) do not depend on the signs $\mathrm{sn}$.
\item[(b)] For every pair of weights $\mathsf{w}=(\mathsf{w}_1,\mathsf{w}_2)$ and all $n,m\ge0$,
\begin{gather*}
\left\|S_LF\right\|_{(\mathsf{w}_1,\mathsf{w}_2)}=\left\|F\right\|_{\left(L^{-1}\cdot\mathsf{w}_1,\,L^{-1-\alpha}\cdot\mathsf{w}_2\right)},\\
\left\|S_LF^{[n,m]}\right\|_{h_{k+1}}=L^{-n-(1+\alpha)m}\cdot L^{-\frac{n+m}{8}+\frac{m\varepsilon}{2}}\cdot\left\|F^{[n,m]}\right\|_{h_k}.
\end{gather*}
\item[(c)] With $L\xi$ and $L\boldsymbol{\zeta}$ as in (G2.r5), the antisymmetric kernels of $S_LF$ are
\[
(S_LF)_{nm}(\xi;\boldsymbol{\zeta})=L^{3n+6m}\cdot L^{-n-(1+\alpha)m}\cdot F_{nm}(L\xi;L\boldsymbol{\zeta})
\]
if all pairs of $\boldsymbol{\zeta}$ have $d_X\le M/L$, and $(S_LF)_{nm}(\xi;\boldsymbol{\zeta})=0$ otherwise.
Consequently $\left\|(S_LF)_{nm}\right\|=L^{-n-(1+\alpha)m}\cdot\left\|F_{nm}\right\|$.
\end{itemize}
\end{named}

\begin{named}{Remark}
Proposition~18.4 corresponds to \cite[Lemma~13.14]{QEDI}, together with the relations between $h_{k+1}$ and $h_k$ stated at the beginning of \cite[Section~13.8]{QEDI}.
Here the sites are those of $T^{\mathrm{f}}_k$, the pairs those of $\mathrm{Pz}_k$, and the length of a pair is measured by $d_X$.
The scaling $x\mapsto x/L$ does not change the label of a fermion site.
The kernels carry the weights $\left(\eta_k^3/48\right)^{n+2m}$ at level $k$ and $\left(\eta_{k+1}^3/48\right)^{n+2m}$ at level $k+1$;
since $\eta_k^3/48=L^3\cdot\eta_{k+1}^3/48$, the factor $L^{3n+6m}$ in (c) is the same as in \cite[Lemma~13.14]{QEDI}.
In (a) the factor is given together with its sign $\mathrm{sn}$ and with a condition under which $\mathrm{sn}=1$.
\end{named}

\subsubsection*{Proposition~18.5 (contraction)}

\begin{named}{Proposition 18.5}[Contraction]
Assume the setting (S-2) of Section~G2.0.4, and let $\mathscr{E}\in K^{x,C}_k$.
Put $L\,R\,\mathscr{E}:=L\left((R\,\mathscr{E})^{\mathrm{nat}}\right)$, that is,
\[
(L\,R\,\mathscr{E})(Y;\mathscr{A},\mathscr{C})=\sum_{X:\ \bar{X}^L=LY}S_L\left[(R\,\mathscr{E})^{\mathrm{nat}}\left(X;\mathscr{A}_L,\mathscr{C}_L\right)\right]
\qquad\left(Y\in D_{k+1},\ (\mathscr{A},\mathscr{C})\in\tilde{R}_{k+1}\times\mathfrak{X}_{k+1}\right).
\]
\begin{itemize}
\item[(a)] The map $L\,R$ is linear; it maps $K^{x,C}_k(\ell)$ into $K^{x,C}_{k+1}(\ell/L)$ and $K^x_k(\ell)$ into $K^x_{k+1}(\ell/L)$.
\item[(b)] One has
\[
\left\|L\,R\,\mathscr{E}\right\|^x_{k+1}\le c_{\mathrm{fl}}\cdot L^{-\frac{1}{4}+2\varepsilon}\cdot\|\mathscr{E}\|^x_k.
\]
In the units of the windows, $w^{\mathrm{E}}_{k+1}=L^{\frac{1}{8}-\frac{7\varepsilon}{2}}\cdot w^{\mathrm{E}}_k$, and therefore
\[
\frac{\left\|L\,R\,\mathscr{E}\right\|^x_{k+1}}{w^{\mathrm{E}}_{k+1}}\le c_{\mathrm{fl}}\cdot L^{-\frac{3}{8}+\frac{11\varepsilon}{2}}\cdot\frac{\|\mathscr{E}\|^x_k}{w^{\mathrm{E}}_k}.
\]
\item[(c)] $(L\,R\,\mathscr{E})|_0=L\,R\,(\mathscr{E}|_0)$, where $L\,R$ on the right is the map of Proposition~48.5.
\item[(d)] (Small polymers.) Let $E:=(R\,\mathscr{E})^{\mathrm{nat}}$.
For every $X\in D_k$,
\[
\sup_{(\mathscr{A},\mathscr{C})\in\tilde{R}_k\times\mathfrak{X}_k}\left\|E(X;\mathscr{A},\mathscr{C})\right\|_{h_k}\le
\begin{cases}
20\,\|\mathscr{E}\|^x_k\cdot\exp\left(-\kappa\cdot d_M(X)\right)&\text{if $X$ is small},\\
\|\mathscr{E}\|^x_k\cdot\exp\left(-\kappa\cdot d_M(X)\right)&\text{if $X$ is large}.
\end{cases}
\]
Let $X$ be small.
Then $E_{00}(X;0,0)=0$ and $E_{20}(X;0,0;\cdot\,,\cdot)=0$, and for every $(\mathscr{A},\mathscr{C})\in\tilde{R}_{k+1}\times\mathfrak{X}_{k+1}$
\[
\left\|S_L\left[E\left(X;\mathscr{A}_L,\mathscr{C}_L\right)\right]\right\|_{h_{k+1}}\le3\,L^{-\frac{13}{4}+\frac{3\varepsilon}{2}}\cdot\sup_{\tilde{R}_k\times\mathfrak{X}_k}\left\|E(X;\cdot\,,\cdot)\right\|_{h_k}.
\]
Bidegree by bidegree:
the part of bidegree $(0,0)$ of $S_L\left[E(X;\mathscr{A}_L,\mathscr{C}_L)\right]$ has norm at most $2L^{-\frac{7}{2}}$ times the supremum on the right, and the part of bidegree $(2,0)$ has norm at most $2\mathrm{e}^4\cdot L^{-4}$ times it.
For the other bidegrees $(n,m)$ the norm of the part of bidegree $(n,m)$ is the norm $\|\cdot\|_{h_k}$ of the part of the same bidegree of $E(X;\mathscr{A}_L,\mathscr{C}_L)$ times a factor;
this factor is $L^{-2-\alpha-\frac{1}{4}+\frac{\varepsilon}{2}}\le L^{-\frac{13}{4}+\frac{3\varepsilon}{2}}$ for $(1,1)$, it is $L^{-2-2\alpha-\frac{1}{4}+\varepsilon}$ for $(0,2)$, and it is at most $L^{-4}$ for $n+m\ge4$.
These factors do not depend on the labels, and they contain no gain $\eta_k^{\alpha}$ from a leg inside a cluster.
\item[(e)] (Sum over the polymers.) For every family $E'$ of level $k$ and every $(\mathscr{A},\mathscr{C})\in\tilde{R}_{k+1}\times\mathfrak{X}_{k+1}$,
\[
\sum_{Y\in D_{k+1}}(L\,E')(Y;\mathscr{A},\mathscr{C})=S_L\left[\sum_{X\in D_k}E'\left(X;\mathscr{A}_L,\mathscr{C}_L\right)\right].
\]
Let $\mathscr{T}=\mathscr{A}+\mathscr{C}$, and let $\mathrm{ev}'$ be a homomorphism of $\mathrm{Gr}[\psi,\chi]$ of level $k+1$ into an associative algebra such that the images $\psi'$, $\bar{\psi}'$, $\chi'$, $\bar{\chi}'$ of the generators satisfy
\[
\chi'=\delta_{\alpha,\mathscr{T}}\psi',\qquad\bar{\chi}'=\delta_{\alpha,-\mathscr{T}}\bar{\psi}'
\]
at all pairs of $\mathrm{Pz}_{k+1}$.
Then $\chi'_L=\delta_{\alpha,\mathscr{T}_L}\psi'_L$ and $\bar{\chi}'_L=\delta_{\alpha,-\mathscr{T}_L}\bar{\psi}'_L$ at all pairs of $\mathrm{Pz}_k$, and
\[
\mathrm{ev}'\left[\sum_{Y\in D_{k+1}}(L\,R\,\mathscr{E})(Y;\mathscr{A},\mathscr{C})\right]=\mathrm{ev}'\,S_L\left[\sum_{X\in D_k}(R\,\mathscr{E})\left(X;\mathscr{A}_L,\mathscr{C}_L\right)\right].
\]
\item[(f)] (One formula; the class of sides consists of the $\ell$ with $\ell\in\mathrm{Cl}$ and $\ell/L\in\mathrm{Cl}$.)
In the sense of Section~I.1.9, $S_L$ is given on the generators indexed by $\mathbb{L}^{\mathrm{f}}_k$ and by $\mathrm{Pzt}_k$ by the one formula (G2.r5), of range $0$;
it intertwines the translation by $v\in LM\mathbb{Z}^3$ at level $k$ with the translation by $v/L$ at level $k+1$.
Let $X$ be a polymer with a lift $\tilde{X}$ such that $\operatorname{sup-diameter}(\tilde{X})+2\eta_k+M<\ell$.
Then the kernels of $(R\,\mathscr{E})^{\mathrm{nat}}(X;\mathscr{A}_L,\mathscr{C}_L)$ and of $S_L\left[(R\,\mathscr{E})^{\mathrm{nat}}(X;\mathscr{A}_L,\mathscr{C}_L)\right]$ are given by one formula of range $\mathrm{rg}=\eta_k$.
The data of this formula are:
the kernels of $\mathscr{E}(X;0,0)$ and of $\mathscr{E}(X;\mathscr{A}_L,\mathscr{C}_L)$, with their arguments over the $\eta_k$-neighbourhood of $\tilde{X}$;
the cubes of $\tilde{X}$, their centres and their common faces;
and $\mathscr{T}_L$ on the bonds with both ends in $\tilde{X}$.
The formula is defined on the data in which the kernels of bidegree $(2,0)$ vanish unless both of their arguments lie over $\tilde{X}$;
by (k2) of Section~I.5.4, the data that arise in the scheme are of this kind.
Let $\tilde{U}$ be the $\eta_k$-neighbourhood of $\tilde{X}$.
If $\operatorname{sup-diameter}(\tilde{U})+LM<\ell$, the margin under which $S_L$ is taken in Section~I.1.9 and in Proposition~25.2\,(o7), then $\mathrm{cov}_{\ell/L}$ is the map of Section~I.1.9 on all of $\mathrm{Grc}_{k+1}(\tilde{U}/L)$.
If $\operatorname{sup-diameter}(\tilde{U})+\eta_k<\ell$, which holds under the margin $\operatorname{sup-diameter}(\tilde{X})+2\eta_k+M<\ell$ of this clause, then on the subalgebras generated by $\psi$ and $\chi$ the rule of Section~I.1.9 defines injective homomorphisms $\mathrm{cov}_{\ell}$ and $\mathrm{cov}_{\ell/L}$, and
\[
\mathrm{cov}_{\ell/L}\circ S_L=S_L\circ\mathrm{cov}_{\ell}.
\]
The two margins $\operatorname{sup-diameter}(\tilde{X})+2\eta_k+M<\ell$ and $\operatorname{sup-diameter}(\tilde{U})+LM<\ell$ hold for small $X$;
for the two sides $\ell=\ell_1$ and $\ell=\ell_2$ of the standing data (TT) of the Propositions~27.m, they also hold for the sets that are admissible in the sense of the Definitions in Proposition~27.1.
\end{itemize}
\end{named}

\begin{named}{Remark}
Clauses (a)--(d) of Proposition~18.5 correspond to \cite[Theorem~13.15]{QEDI} and to \cite[Proposition~5.62]{QEDI};
the bound in the units of the windows in (b) stands in the remark that follows \cite[Proposition~5.62]{QEDI} and in the remarks after the proof of \cite[Theorem~13.15]{QEDI};
clause (f) takes the place of \cite[Remark~13.16]{QEDI}.
All objects are formed on $T^{\mathrm{f}}_k$ and on $\mathrm{Pz}_k$, and the constants $c_{\mathrm{fl}}$, $20$, $3$, $60$ and $61\,c_{\mathrm{rb}}$ are those of \cite[Theorem~13.15]{QEDI}.
In (d), the bound on the suprema over the connections and the factors by bidegree are assertions of the proposition, whereas they stand in the proof of \cite[Theorem~13.15]{QEDI};
the last sentence of (d) has no counterpart in \cite{QEDI} (see Section~I.5.5).
In (c), the last assertion of \cite[Theorem~13.15(c)]{QEDI}, on the families without the second slot, is replaced by Proposition~48.5.
In clause (e), the relations for $\chi'_L$ and $\bar{\chi}'_L$ are (I.O9), with the distance $d_X$ in the H\"older quotients.
\end{named}

\subsubsection*{Proposition~18.6 (bilinears inside a cluster; the seven light-heavy couplings through extraction, adjustment and scaling)}

\begin{named}{Proposition 18.6}[Bilinears inside a cluster; the seven light-heavy couplings through extraction, adjustment and scaling]
Assume the setting (S-2) of Section~G2.0.4.
Let $\mathrm{Ng}$ be a function on $O_{\mathrm{cub}}\times O_{\mathrm{cub}}$ with values in the complex $4\times4$ matrices such that, for all $g,g'\in O_{\mathrm{cub}}$,
\begin{align}
S\cdot\mathrm{Ng}(g,g')\cdot S^{-1}&=\mathrm{Ng}\left(r(S)g,\,r(S)g'\right)\qquad\text{for }S\in\mathrm{Sp},\tag{n1}\\
\mathsf{C}^{-1}\cdot\mathrm{Ng}(g',g)^T\cdot\mathsf{C}&=\mathrm{Ng}(g,g'),\tag{n2}
\end{align}
and, where this is stated,
\begin{equation}
\varkappa^{-1}\cdot\overline{\mathrm{Ng}(g,g')}\cdot\varkappa=\mathrm{Ng}(g,g').\tag{n3}
\end{equation}
Put
\[
\mathrm{Nl}(g):=\frac{1}{48}\sum_{g'}\mathrm{Ng}(g,g'),\qquad
\mathrm{Nav}:=\left(\frac{1}{48}\right)^2\sum_{g,g'}\mathrm{Ng}(g,g'),\qquad
|\mathrm{Ng}|_1:=\left(\frac{1}{48}\right)^2\sum_{g,g'}\ \sum_{\beta,\beta'}\left|\mathrm{Ng}(g,g')_{\beta\beta'}\right|.
\]
Let $E^{\mathrm{Ng}}$ be the family with
\[
E^{\mathrm{Ng}}(\square;\mathscr{A},\mathscr{C}):=\int_{\square}\mathrm{d}x\,\frac{1}{48}\sum_{g'}\bar{\psi}(c,g)^T\cdot\mathrm{Ng}(g,g')\cdot\psi(c,g')\qquad\left(x=(c,g)\right)
\]
for an $M$-cube $\square$, and $E^{\mathrm{Ng}}(X):=0$ if $|X|_M\ge2$.
This definition is used for every function $\mathrm{Ng}$ on $O_{\mathrm{cub}}\times O_{\mathrm{cub}}$, also for one without (n1) and (n2) (clause (g)).
On $\mathfrak{S}^{\mathrm{ch}}(\eta_k,\ell)$ let
\[
(\mathrm{Ocl}\,f)(c,g):=\frac{1}{48}\sum_{g'}\mathrm{Ng}(g,g')\cdot f(c,g').
\]
\begin{itemize}
\item[(a)] $E^{\mathrm{Ng}}$ lies in $K^{x,C}_k$, and in $K^x_k$ if (n3) holds; it does not depend on $(\mathscr{A},\mathscr{C})$.
Its only kernel is
\[
E^{\mathrm{Ng}}_{20}(\square;x,x')=\eta_k^{-3}\cdot\mathbf{1}[c=c'\in\square]\cdot\mathrm{Ng}(g,g'),
\]
for both arguments over one lift of $\square$; this kernel is continued antiperiodically in each argument (Section~I.3.6, (I.AP)).
Moreover
\[
\left\|E^{\mathrm{Ng}}\right\|^x_k=M^3\cdot e_k^{-\frac{1}{2}}\cdot|\mathrm{Ng}|_1,\qquad
\sum_{\square}E^{\mathrm{Ng}}(\square)=\left\langle\bar{\psi},\mathrm{Ocl}\,\psi\right\rangle^{\mathrm{ch}}.
\]
\item[(b)] $\mathrm{al}_{0}(E^{\mathrm{Ng}},\square)=0$, and $\mathrm{al}_{2}(E^{\mathrm{Ng}},\square)=\mathrm{Nav}=\mathrm{na}\cdot\mathbf{1}$ with a complex number $\mathrm{na}$, which is real if (n3) holds.
Further $\varepsilon(E^{\mathrm{Ng}})=0$ and $m(E^{\mathrm{Ng}})=-\mathrm{na}$.
\item[(c)] For $\xi\in\mathfrak{S}(\eta_k,\ell)$ one has $(\mathrm{Ocl}\,J\xi)(c,g)=\mathrm{Nl}(g)\cdot\xi(c)$, and $J^{*}\cdot\mathrm{Ocl}\cdot J=\mathrm{na}$:
of the action on the cluster constants only the mean over the label is extracted.
The map $\mathrm{Phv}\,\mathrm{Ocl}\,J\xi=\left(\mathrm{Nl}(\cdot)-\mathrm{na}\right)\xi$ is the coupling of the light part to the heavy part at zero momentum.
\item[(d)] $(R\,E^{\mathrm{Ng}})(\square)=E^{\mathrm{Ng}'}(\square)$, where
\[
\mathrm{Ng}'(g,g'):=\mathrm{Ng}(g,g')-48\cdot\mathbf{1}[g=g']\cdot\mathrm{na};
\]
the function $\mathrm{Ng}'$ has (n1) and (n2), and (n3) if $\mathrm{Ng}$ has it, and $\mathrm{Nav}'=0$.
\item[(e)] Let $E:=(R\,E^{\mathrm{Ng}})^{\mathrm{nat}}$.
Then $E_{20}(\square;0,0;\cdot\,,\cdot)=0$, every monomial of $E(\square;0,0)$ contains at least one H\"older leg, and
\[
\sup_{\tilde{R}_k\times\mathfrak{X}_k}\left\|E(\square;\cdot\,,\cdot)\right\|_{h_k}\le20\,M^3\cdot e_k^{-\frac{1}{2}}\cdot|\mathrm{Ng}|_1.
\]
For $(\mathscr{A},\mathscr{C})\in\tilde{R}_{k+1}\times\mathfrak{X}_{k+1}$,
\[
\left\|S_L\left[E\left(\square;\mathscr{A}_L,\mathscr{C}_L\right)\right]\right\|_{h_{k+1}}\le60\,L^{-\frac{13}{4}+\frac{3\varepsilon}{2}}\cdot M^3\cdot e_k^{-\frac{1}{2}}\cdot|\mathrm{Ng}|_1.
\]
Moreover $(L\,R\,E^{\mathrm{Ng}})(Y)=0$ unless $Y$ is an $M$-cube of level $k+1$, and
\[
\left\|L\,R\,E^{\mathrm{Ng}}\right\|^x_{k+1}\le60\,L^{-\frac{1}{4}+\frac{3\varepsilon}{2}}\cdot M^3\cdot e_k^{-\frac{1}{2}}\cdot|\mathrm{Ng}|_1=60\,L^{-\frac{1}{4}+\frac{3\varepsilon}{2}}\cdot\left\|E^{\mathrm{Ng}}\right\|^x_k.
\]
\item[(f)] (The seven couplings of Section~I.4.6.)
Let $\mathrm{co}_1$ be a real number and $\mathrm{cu},\mathrm{cv}\in\mathbb{R}^3$.
Write $\mathrm{cu}=\sum_{i\le2}\mathrm{ru}_i\,\hat{n}_i$ with real numbers $\mathrm{ru}_i$ (the vectors $\hat{n}_0,\hat{n}_1,\hat{n}_2$ form a basis of $\mathbb{R}^3$), and $\mathrm{cv}=\mathrm{cv}_0+\mathrm{cv}_1+\mathrm{cv}_2$ with vectors $\mathrm{cv}_i\in\mathbb{R}^3$ such that $\mathrm{cv}_i\cdot\hat{n}_i=0$ (the three planes orthogonal to $\hat{n}_0,\hat{n}_1,\hat{n}_2$ span $\mathbb{R}^3$).
The three functions
\begin{gather*}
\mathrm{Ng}_1(g,g'):=48\cdot\mathbf{1}[g=g']\cdot i\,\mathrm{co}_1\cdot\det(g)\cdot\Sigma_3,\qquad
\mathrm{Ng}_u(g,g'):=48\sum_{i\le2}\mathrm{ru}_i\cdot\mathbf{1}[g'=gs_i]\cdot\gamma\cdot(g\hat{n}_i),\\
\mathrm{Ng}_v(g,g'):=48\sum_{j\le2}\mathbf{1}[g'=gs_j]\cdot i\det(g)\cdot(\mathbf{1}_2\otimes\sigma)\cdot(g\,\mathrm{cv}_j)
\end{gather*}
(in $\mathrm{Ng}_1$ and $\mathrm{Ng}_v$ the factor $i$ is the imaginary unit $\sqrt{-1}$, the same in the three terms of the sum over $j$)
have (n1), (n2), (n3) and $\mathrm{Nav}=0$, and the functions $\mathrm{Nl}$ formed with them are
\[
\mathrm{Nl}_1(g)=i\,\mathrm{co}_1\cdot\det(g)\cdot\Sigma_3,\qquad
\mathrm{Nl}_u(g)=\gamma\cdot(g\,\mathrm{cu}),\qquad
\mathrm{Nl}_v(g)=i\det(g)\cdot(\mathbf{1}_2\otimes\sigma)\cdot(g\,\mathrm{cv}).
\]
So (a)--(e) hold for each of them with $\mathrm{na}=0$, and $R\,E^{\mathrm{Ng}}=E^{\mathrm{Ng}}$: nothing is extracted.
Furthermore
\[
E^{\mathrm{Ng}_u}(\square)=-\sum_{i\le2}\mathrm{ru}_i\int_{\square}\mathrm{d}x\,\bar{\psi}(x)^T\left(\gamma\cdot\hat{d}_i(x)\right)\psi\left(x^{(i)}\right),
\qquad|\mathrm{Ng}_1|_1=4\,|\mathrm{co}_1|,\qquad|\mathrm{Ng}_u|_1\le4\sqrt{2}\sum_i|\mathrm{ru}_i|.
\]
Let $\mathrm{co}_0$ be a real number and $\mathrm{Ng}_0(g,g'):=48\cdot\mathbf{1}[g=g']\cdot\mathrm{co}_0$;
this function has (n1), (n2) and (n3), with $\mathrm{na}=\mathrm{co}_0$ and $m(E^{\mathrm{Ng}_0})=-\mathrm{co}_0$.
The sum $\mathrm{Ng}_0+\mathrm{Ng}_1+\mathrm{Ng}_u+\mathrm{Ng}_v$ has (n1), (n2) and (n3), and for it
\[
\mathrm{Nl}(g)=\mathrm{co}_0+i\det(g)\cdot\mathrm{co}_1\cdot\Sigma_3+\gamma\cdot(g\,\mathrm{cu})+i\det(g)\cdot(\mathbf{1}_2\otimes\sigma)\cdot(g\,\mathrm{cv}),
\]
which is the display of Section~I.4.6.
\item[(g)] The on-site function $48\cdot\mathbf{1}[g=g']\cdot\gamma\cdot(g\,\mathrm{cu})$ has (n1) and (n3), and $\mathrm{Nl}(g)=\gamma\cdot(g\,\mathrm{cu})$ for it.
Unless $\mathrm{cu}=0$ it violates (n2), and the corresponding family violates (k5) of Section~I.5.4.
\end{itemize}
\end{named}

\begin{named}{Remark}
Section~I.4.6 describes the seven couplings through the action of $\mathrm{Ocl}$ on the cluster constants, that is, through $\mathrm{Nl}$ alone.
In Proposition~18.6 the datum is the function $\mathrm{Ng}$ on $O_{\mathrm{cub}}\times O_{\mathrm{cub}}$, and (a)--(e) hold for every $\mathrm{Ng}$ with (n1) and (n2).
The couplings $\mathrm{cu}$ and $\mathrm{cv}$ are realised by functions supported on $g'=gs_i$, $i\le2$ (clause (f)), and the on-site function with $\mathrm{Nl}(g)=\gamma\cdot(g\,\mathrm{cu})$ is excluded by (n2) (clause (g)).
\end{named}

\subsection*{20. Terminal bounds; the free determinant}

\subsubsection*{Proposition~20.1 (terminal bounds)}

\begin{named}{Proposition 20.1}[Terminal bounds]
Assume the setting (S-3) of Section~G2.0.4; thus $k\ge1$, and $\Psi$, $\bar{\Psi}$ are indexed by $T^0_k$, with $4$ components per unit site and the weight $\nu=1$.
Let $C_0$ be the constant of (C6), so that $C_0\ge C^F$ and $C_0\ge4b_{+}+16\,C^F\cdot c(\gamma^F)$ (Section~I.8.2), and put
\[
\mathrm{aw}_1:=e_k^{-\frac{1}{4}}\ \left(=h_{k,1}\right),\qquad
\mathrm{aw}_2:=C_0^{-1}\cdot e_k^{-\frac{1}{4}+\varepsilon}\ \left(=C_0^{-1}\cdot h_{k,2}\le\mathrm{aw}_1\right).
\]
For $F\in\mathrm{Gr}[\psi,\chi]$ and $B\in\tilde{G}_k$, let $F\left(H_k^{\#}(B)\Psi\right)$ be the image of $F$ under the homomorphism with
\[
\psi\mapsto H_k(B)\Psi,\qquad\chi\mapsto\delta_{\alpha,B}H_k(B)\Psi,\qquad\bar{\psi}\mapsto\bar{H}_k(B)\bar{\Psi},\qquad\bar{\chi}\mapsto\delta_{\alpha,-B}\bar{H}_k(B)\bar{\Psi}.
\]
\begin{itemize}
\item[(a)] For $A\in\tilde{G}_k$ the operator $D_k(A)$ is defined, its entries are analytic functions on $\tilde{G}_k$, and
\[
\left|\left\langle\bar{\Psi},D_k(A)\Psi\right\rangle\right|_{\mathrm{aw}_1}\le C_0\cdot e_k^{-\frac{1}{2}}\cdot\ell^3.
\]
\item[(b)] For $A\in\tilde{G}_k$, at $\psi=H_k(A)\Psi$ and $\bar{\psi}=\bar{H}_k(A)\bar{\Psi}$,
\[
\left|\left\langle\bar{\psi},\psi\right\rangle\right|_{\mathrm{aw}_2}\le4e_k^{-\frac{1}{2}+2\varepsilon}\cdot\ell^3\le4e_k^{-\frac{1}{2}}\cdot\ell^3.
\]
\item[(c)] For $(\mathscr{A},\mathscr{C})\in\tilde{R}_k\times\mathfrak{X}_k$ the total connection $\mathscr{T}=\mathscr{A}+\mathscr{C}$ lies in $\tilde{G}_k$, and
\[
\left|F\left(H_k^{\#}(\mathscr{T})\Psi\right)\right|_{\mathrm{aw}_2}\le\|F\|_{h_k}
\qquad\text{for every }F\in\mathrm{Gr}[\psi,\chi];
\]
in particular
\[
\left|\mathscr{E}\left(X;\mathscr{A},\mathscr{C},H_k^{\#}(\mathscr{T})\Psi\right)\right|_{\mathrm{aw}_2}\le\left\|\mathscr{E}(X;\mathscr{A},\mathscr{C})\right\|_{h_k}
\qquad\left(\mathscr{E}\in K^{x,C}_k,\ X\in D_k\right).
\]
\end{itemize}
\end{named}

\begin{named}{Remark}
Proposition~20.1 corresponds to \cite[Theorem~13.17]{QEDI} and to \cite[Proposition~5.64]{QEDI}; its clauses (b) and (c) are recalled, at the stop level, as (Gr6) in \cite[Section~20.2.4]{QEDI}.
The differences are the following.
(1) The three hypotheses of \cite[Theorem~13.17]{QEDI}, which are assertions of \cite[Propositions~5.42, 5.44 and~5.45]{QEDI}, are not hypotheses of Proposition~20.1.
(2) $H_k(A)$ maps $\mathfrak{S}(1,\ell)$ into $\mathfrak{S}^{\mathrm{ch}}(\eta_k,\ell)$;
the row sums of $H_k^{\#}(A)$ extend over the fermion sites and over all pairs of $\mathrm{Pz}_k$, the $48\cdot47$ ordered pairs inside each cluster (with $d_X=\eta_k$) included (Section~I.5.8).
(3) In $\mathrm{Gr}[\psi]$ the element $\langle\bar{\psi},\psi\rangle=\int\mathrm{d}x\,\bar{\psi}\cdot\psi$ has the norm $4\mathsf{w}_1^2\int\mathrm{d}x=4\ell^3\mathsf{w}_1^2$ for the weight $\mathsf{w}_1$;
this is the origin of the factor $4\ell^3$ in (b).
(4) $C_0$ is the constant of (C6), not a free number, and the constants $(C_F,\gamma_F)$ of \cite[Theorem~13.17]{QEDI} are $(C^F,\gamma^F)$.
(5) The norm on $\mathbb{C}^4$ in \cite[Proposition~5.45]{QEDI} is $|\cdot|_c$;
in (a), a bound $\left|M_k(A)(y,y')\right|_c\le\mathrm{cf}$ bounds the four row sums of the block, so that the estimate by $4b_k+16\,C^F\cdot\sum\exp(\cdot)$ remains valid.
(6) Clause (c) is stated for every element $F$ of $\mathrm{Gr}[\psi,\chi]$, not only for the values of a family.
(7) The weights $\mathrm{aw}_1$, $\mathrm{aw}_2$ are the weights $a_1$, $a_2$ of \cite[Theorem~13.17]{QEDI}; here $a_1$, $a_2$ are the lengths $a_k$ at $k=1,2$ (Section~I.6.1).
\end{named}

\subsubsection*{Proposition~20.2 (determinants of one step; the algebra)}

\begin{named}{Proposition 20.2}[Determinants of one step; the algebra]
Let $k\ge0$ be an integer, $\ell\in\mathrm{Cl}$ a side, $e_k>0$ a number, $B$ a complex bond function on $T_k$, and $\mathrm{zc}$ a complex number, which takes the place of the mass in all operators of Sections~I.3.5 and~I.3.9.
Put
\[
\mathrm{Rz}:=D_B+\mathrm{zc}\qquad\text{on }\mathfrak{S}^{\mathrm{ch}}(\eta_k,\ell),
\]
a space of dimension $192\cdot(\ell/\eta_k)^3$.
Here $n_k=\ell^3$ and $n_{k+1}=\ell^3/L^3$;
$b_{k+1}/L=b_L\cdot b_k/(b_k+b_L)$ for $k\ge1$, and $b_1/L=b_L$;
and $\det$ is the determinant of a linear map of a finite-dimensional space into itself.

\emph{The levels $k\ge1$.}
Let $(B,\mathrm{zc})$ be such that $\mathrm{Rz}+b_k\cdot JP_k(B)J^{*}$ is invertible, let $S_k(B)$ be its inverse, and put
\[
D_k(B):=b_k-b_k^2\cdot Q_k(B)J^{*}\cdot S_k(B)\cdot JQ_k^T(-B)\quad\text{on }\mathfrak{S}(1,\ell),\qquad
\delta Z_k(B):=\det\left(D_k(B)+b_L\cdot P(B)\right).
\]
Then the following hold.
\begin{itemize}
\item[(i)] $\det D_k(B)\cdot\det\left(\mathrm{Rz}+b_k\cdot JP_k(B)J^{*}\right)=b_k^{4n_k}\cdot\det\mathrm{Rz}$.
\item[(ii)] One has
\[
\delta Z_k(B)\cdot\det\left(\mathrm{Rz}+b_k\cdot JP_k(B)J^{*}\right)=b_k^{4(n_k-n_{k+1})}\cdot\left(b_k+b_L\right)^{4n_{k+1}}\cdot\det\left(\mathrm{Rz}+\frac{b_{k+1}}{L}\cdot JP_{k+1}(B)J^{*}\right).
\]
\item[(iii)] If $\delta Z_k(B)\ne0$, then $\mathrm{Rz}+\frac{b_{k+1}}{L}\cdot JP_{k+1}(B)J^{*}$ is invertible.
Let $S^{0}_{k+1}(B)$ be its inverse, and let
\[
D^{0}_{k+1}(B):=\frac{b_{k+1}}{L}-\left(\frac{b_{k+1}}{L}\right)^2\cdot Q(B)\cdot Q_k(B)J^{*}\cdot S^{0}_{k+1}(B)\cdot JQ_k^T(-B)\cdot Q^T(-B)\qquad\text{on }\mathfrak{S}(L,\ell);
\]
this operator equals $b_L-b_L^2\cdot Q(B)\cdot\left(D_k(B)+b_L\cdot P(B)\right)^{-1}\cdot Q^T(-B)$.
Then
\[
L^{4n_{k+1}}\cdot\det D^{0}_{k+1}(B)=b^{4n_{k+1}}\cdot\frac{\det D_k(B)}{\delta Z_k(B)}.
\]
\end{itemize}

\emph{The level $k=0$.}
Put $D_0(B):=\mathrm{Rz}$ on $\mathfrak{S}^{\mathrm{ch}}(1,\ell)$, of size $192n_0$, and $\delta Z_0(B):=\det\left(\mathrm{Rz}+b_L\cdot JP(B)J^{*}\right)$, of size $192n_0$.
\begin{itemize}
\item[(iii0)] Let $\delta Z_0(B)\ne0$, and let
\[
D^{0}_1(B):=b_L-b_L^2\cdot Q(B)J^{*}\cdot\left(\mathrm{Rz}+b_L\cdot JP(B)J^{*}\right)^{-1}\cdot JQ^T(-B)\qquad\text{on }\mathfrak{S}(L,\ell),
\]
of size $4n_1$.
Then
\[
L^{4n_1}\cdot\det D^{0}_1(B)=b^{4n_1}\cdot\frac{\det D_0(B)}{\delta Z_0(B)}.
\]
\end{itemize}

\emph{Scaling ($k\ge0$).}
Since $\ell\in\mathrm{Cl}$, the side $\ell/L=M\cdot\ell/(LM)$ is an odd multiple of $M$, and $\ell/L\ge13M$.
Let $e_{k+1}:=L^{\frac{1}{2}}\cdot e_k$, let $B'$ be a complex bond function on $T_{k+1}=\mathbb{L}_{k+1}/(\ell/L)\mathbb{Z}^3$, let $B=B'_L$, and let $\mathrm{zc}':=L\cdot\mathrm{zc}$.
The operators with the argument $B'$ that occur are
\[
D_{B'},\quad Q_{k+1}(B')J^{*},\quad JQ_{k+1}^T(-B'),\quad JP_{k+1}(B')J^{*},\quad S_{k+1}(B'),\quad D_{k+1}(B');
\]
none of them contains a one-step average (Section~G3.1.0, setting (S21)).
They are given by the displays of Sections~I.3.3 and~I.3.5 at the level $k+1$, on the side $\ell/L$, with the charge $e_{k+1}$ and with $\mathrm{zc}'$ in the place of the mass.
The operator $P_{k+1}(B'):=Q_{k+1}^T(-B')\cdot Q_{k+1}(B')$ is the operator $P_k(\cdot)$ of Section~I.3.3 at the level $k+1$;
it is distinguished by its argument from the operator $P_{k+1}(B)=Q_k^T(-B)\cdot P(B)\cdot Q_k(B)$ of (ii) and (iii).
In what follows $f'$ ranges over $\mathfrak{S}^{\mathrm{ch}}(\eta_{k+1},\ell/L)$ and $v$ over $\mathfrak{S}(1,\ell/L)$;
then $f'_L\in\mathfrak{S}^{\mathrm{ch}}(\eta_k,\ell)$ and $v_L\in\mathfrak{S}(L,\ell)$ (Section~I.3.7).
\begin{itemize}
\item[(iv)] One has
\begin{gather*}
\left(D_B+\mathrm{zc}\right)f'_L=L^{-1}\left(\left(D_{B'}+\mathrm{zc}'\right)f'\right)_L,\\
Q(B)Q_k(B)J^{*}f'_L=\left(Q_{k+1}(B')J^{*}f'\right)_L,\qquad
JQ_k^T(-B)Q^T(-B)v_L=\left(JQ_{k+1}^T(-B')v\right)_L
\end{gather*}
(at $k=0$ with $Q(B)J^{*}$ and $JQ^T(-B)$ on the left).
Hence $S_{k+1}(B')$ exists if and only if $\mathrm{Rz}+\frac{b_{k+1}}{L}\cdot JP_{k+1}(B)J^{*}$ (at $k=0$: $\mathrm{Rz}+b_L\cdot JP(B)J^{*}$) is invertible, and then
\begin{gather*}
S^{0}_{k+1}(B)f'_L=L\left(S_{k+1}(B')f'\right)_L,\qquad
D^{0}_{k+1}(B)v_L=L^{-1}\left(D_{k+1}(B')v\right)_L,\\
\det D^{0}_{k+1}(B)=L^{-4n_{k+1}}\cdot\det D_{k+1}(B').
\end{gather*}
Consequently, whenever $S_k(B)$ exists (for $k\ge1$) and $\delta Z_k(B)\ne0$, the determinants of the levels $k+1$ and $k$ are related by the identity (6.1) of one step:
\begin{equation}
\det D_{k+1}(B')=b^{4n_{k+1}}\cdot\frac{\det D_k(B)}{\delta Z_k(B)},\tag{6.1}
\end{equation}
with the mass $\mathrm{zc}'=L\cdot\mathrm{zc}$ at level $k+1$ and the mass $\mathrm{zc}$ at level $k$;
this holds for every $k\ge0$, and at $k=0$ one has $D_0(B)=D_B+\mathrm{zc}$.
\end{itemize}
The exponents $4n_k$ and $4n_{k+1}$ are the dimensions of $\mathfrak{S}(1,\ell)$ (for $k\ge1$) and of $\mathfrak{S}(L,\ell)$, and $4n_{k+1}$ is the rank of $P(B)$, resp.\ of $JP(B)J^{*}$; the dimension of the fine space does not enter.
\end{named}

\begin{named}{Remark}
Proposition~20.2 corresponds to the beginning of \cite[Section~13.11.1]{QEDI} and to \cite[Lemma~13.19]{QEDI}, including the identity displayed at the end of that lemma, which is (6.1);
with $c'_k=b^{4n_{k+1}}$, the identity (6.1) is also that of \cite[Proposition~5.55(c)]{QEDI}.
The differences are the following.
(1) The averaging operators on the fine side are $Q_k(B)J^{*}$ and $JQ_k^T(-B)$ (at $k=0$: $Q(B)J^{*}$ and $JQ^T(-B)$);
the relations $\left(Q_k(B)J^{*}\right)\left(JQ_k^T(-B)\right)=1$ and $\left(Q(B)J^{*}\right)\left(JQ^T(-B)\right)=1$ hold by $J^{*}J=1$, also for complex $B$.
(2) The fine space has the dimension $192\cdot(\ell/\eta_k)^3$; this number does not enter the identities.
(3) At level $0$, $\det D_0(B)$ and $\delta Z_0(B)$ are determinants of size $192n_0$, the operator $JP(B)J^{*}$ has the rank $4n_1$, and $c'_0=b^{4n_1}$, as in \cite{QEDI}; see (I.L0.10).
(4) $D^{0}_{k+1}(B)$ has the normalisation of \cite[Theorem~14.21(a)]{QEDI}; the operator denoted $D^0_{k+1}$ in \cite[Lemma~13.19]{QEDI} is $L\cdot D^{0}_{k+1}(B)$.
(5) In (iv) the scaling of $D_B$ is (I.O9); the maps $U$ and $U_0$ of \cite[Lemma~13.19(iv)]{QEDI} are $f'\mapsto L\,f'_L$ and $v\mapsto L\,v_L$.
(6) The charges are data, subject only to $e_{k+1}=L^{\frac{1}{2}}\cdot e_k$.
(7) In \cite[Lemma~13.19]{QEDI} the side $\ell$ is any positive multiple of $L$; here $\ell\in\mathrm{Cl}$ (Section~G2.0.4, Tori), as in Proposition~21.4, so that $\tau$, $\tau_k$, the operators of level $k+1$ and the change of scale are formed only on sides on which they are defined.
\end{named}

\subsubsection*{Proposition~20.3 (the free determinant telescopes)}

\begin{named}{Proposition 20.3}[The free determinant telescopes]
Assume the setting (S-4) of Section~G2.0.4, and let $k$ be an integer with $0\le k\le k_{*}(R)$.
For every $j<k$ let a family (p5) $S_{j,y}(X,A)$ of level $j$ on $T_j(R)$ be given (Section~I.5.8; $X\in D_j$, $A\in\tilde{G}_j$, $y\ge0$, Borel in $y$) which has the properties W1, W2, W7$^{\mathrm{ch}}$, W12, W13 and W14 for $\mathrm{par}$;
here W1, W2, W7, W12, W13, W14 are the properties 1, 2, 7, 12, 13, 14 of the systems of primitives in \cite[Section~5.4.2]{QEDI}, read as in Section~I.5.8.
(In the scheme these families are the primitives of (FL): Sections~I.1.8 and~G4.0.1.)
Let $\acute{E}_j$, $E^{d}_j$ and $E_j^{\mathrm{det}}$ be formed with these families (Section~I.5.7).
Put $\mathfrak{P}_0:=0$ and
\[
\mathfrak{P}_{j+1}(Y,\mathscr{C}):=\sum_{X:\ \bar{X}^L=LY}\left(\mathfrak{P}_j\left(X,\mathscr{C}_L\right)+E_j^{\mathrm{det}}\left(X,\mathscr{C}_L\right)\right)\qquad(j<k);
\]
in this proposition the recursion is evaluated only at $\mathscr{C}=a^{[j+1]}$, for which $\mathscr{C}_L=a^{[j]}$ lies in $\mathfrak{X}_j$ by the data.
Let $a\in C^1(n)\otimes\mathbb{C}$ be a complex one-cochain with $a^{[j]}\in\mathfrak{X}_j$ for all $j\le k$.
The operators of level $j$ are formed with $(j,\ell_j,\mu_j,e_j)$, where $\mu_j=\bar{m}\cdot a_j$ (at $j=0$: the mass $\bar{m}\cdot\eta$ and the charge $q$).
$D_a$ and $D_0$ denote the operator (I.2.3) at level $0$ at the connections $a$ and $0$, on $\mathfrak{S}^{\mathrm{ch}}(1,n)$.
\begin{itemize}
\item[(a)] $a^{[j]}$ lies in $\tilde{G}_j$ for $0\le j\le k$, and $\left(a^{[j+1]}\right)_L=a^{[j]}$ and $\mu_{j+1}=L\mu_j$ for $0\le j<k$.
The operator $S_j\left(a^{[j]}\right)$, hence $D_j\left(a^{[j]}\right)$, is defined for $1\le j\le k$;
$\delta Z_j\left(a^{[j]}\right)\ne0$ and $\delta Z_j(0)\ne0$ for $j<k$;
and $\det D_j(0)\ne0$ for $0\le j\le k$.
\item[(b)] For $0\le j<k$,
\[
\frac{\det D_{j+1}\left(a^{[j+1]}\right)}{\det D_{j+1}(0)}=\frac{\det D_j\left(a^{[j]}\right)}{\det D_j(0)}\cdot\exp\left\{-\sum_{X\in D_j}E_j^{\mathrm{det}}\left(X,a^{[j]}\right)\right\}.
\]
In the step $j=0$, $D_0\left(a^{[0]}\right)=D_a+\bar{m}\cdot\eta$ has the size $192n^3$ and $D_1(\cdot)$ has the size $4n_1$;
$E_0^{\mathrm{det}}$ is formed with $\acute{E}_0$ and contains the determinant of the $188$ heavy components per photon site.
\item[(c)] One has
\begin{gather*}
\frac{\det\left(D_a+\bar{m}\cdot\eta\right)}{\det\left(D_0+\bar{m}\cdot\eta\right)}=\exp\left\{\sum_{X\in D_k}\mathfrak{P}_k\left(X,a^{[k]}\right)\right\}\cdot\frac{\det D_k\left(a^{[k]}\right)}{\det D_k(0)},\\
\sum_{X\in D_k}\mathfrak{P}_k\left(X,a^{[k]}\right)=\sum_{j<k}\ \sum_{X\in D_j}E_j^{\mathrm{det}}\left(X,a^{[j]}\right).
\end{gather*}
\end{itemize}
At $a=(t/e)\,\hat{\mathfrak{a}}$ the left side of (c) is the free ratio $N_0(t)$ of Section~I.7.6.
The operators contain no mass other than the running masses $\mu_j=\bar{m}\cdot a_j$, and $\mu_j>0$ is the only positivity of a mass that enters; no lattice mass $\boldsymbol{\mu}$ occurs (I.K10).
\end{named}

\begin{named}{Remark}
Proposition~20.3 corresponds to \cite[Theorem~13.20]{QEDI} and to \cite[Proposition~5.66]{QEDI}.
The differences are the following.
(1) In \cite[Theorem~13.20]{QEDI} the existence of the operators $S_j(A)$, in the form of \cite[Proposition~5.42]{QEDI} stated there, and the assertion (d0) of \cite[Proposition~5.48]{QEDI} for the families of the levels $j<k$ are hypotheses.
They are not part of Proposition~20.3: its only datum on the primitives is, for each $j<k$, one family (p5) of Section~I.5.8, with the properties under which Proposition~13\,(d0) is stated, and $\acute{E}_j$, $E^{d}_j$ and $E_j^{\mathrm{det}}$ are formed with this same family;
the setting (S-4) contains no system of primitives.
In \cite[Theorem~13.20]{QEDI} two sets of families occur:
the systems of primitives of the scheme, \cite[Definition~5.13]{QEDI}, which the theorem assumes through the last clause of the hypothesis for the flow, \cite[Definition~5.70]{QEDI};
and the families (p5) of the levels $j<k$ that the theorem lets be given, for which \cite[Proposition~5.48(d0)]{QEDI} is assumed and with which $E_j^{\mathrm{det}}$ and the functions $\mathfrak{P}_j$ are formed.
(2) At level $0$ the determinants have the size $192n^3$, and W7 at $j=0$ is W7$^{\mathrm{ch}}$, the kernel bound of Section~I.5.8 for the kernels (I.1.4), whose indices are fermion sites.
The factor $(1+y)^{-1}$ of this bound and the Borel measurability in $y$ stand also in \cite{QEDI}:
by \cite[Section~5.4.2]{QEDI}, property~7 at $k=0$ is the kernel bound displayed in property~17, which carries the factor $(1+y)^{-1}$;
and by \cite[Section~5.4.1]{QEDI}, the operators $S_{0,y}(X,A)$ of a system of primitives of level~$0$ are Borel measurable functions of $y$.
(3) The clauses (a) and (b) are steps of the proof of \cite[Theorem~13.20]{QEDI}.
(4) At level $0$ the accretivity is that of Proposition~19\,(a).
(5) The upper end $\mathfrak{m}_0$ of the mass range, the integer $m_0$ and the number $h'$ of \cite[Definition~5.98]{QEDI} do not occur (Section~I.6.1); one has $0<\mu_j\le\mu_K\le1$.
\end{named}

\subsection*{48. The statements without the second slot; the invariant at $\mathscr{C}=0$}

The following statements hold for the families with the second slot and for the families without it alike:
the group lemma, Proposition~16.1;
the abstract Proposition~15;
Proposition~17.1, in which the bond function $B$ is arbitrary;
Proposition~18.1\,(a), (c);
the definition and the clauses (a), (b) of Proposition~18.3;
Proposition~18.4, Proposition~19, Proposition~20.1 and Proposition~20.2.
The second connection $\mathscr{C}$ enters them only in Proposition~17.1\,(a), in the definition of Proposition~18.3 and in Proposition~20.1\,(c).
Proposition~17.1\,(a) bounds $e_k\cdot|\operatorname{Im}\mathscr{T}|_{\infty}$ both for $\mathscr{T}=\mathscr{A}+\mathscr{C}$ with $(\mathscr{A},\mathscr{C})\in\tilde{R}_k\times\mathfrak{X}_k$ and for $\mathscr{T}=\mathscr{A}\in\tilde{R}_k$.
In the definition of Proposition~18.3 the set $\mathrm{Dc}$ of pairs is arbitrary; for the families without slot one takes $\mathrm{Dc}=\tilde{R}_k\times\{0\}$ (Proposition~48.4\,(c)).
Proposition~20.1\,(c) holds at every pair $(\mathscr{A},\mathscr{C})$ and for every $F\in\mathrm{Gr}[\psi,\chi]$, hence also at $\mathscr{C}=0$, which lies in $\mathfrak{X}_k$ by Proposition~48.1\,(a).
In particular, in the setting (G2.S-3), Proposition~20.1\,(c) at $\mathscr{C}=0$ and $F=E(X;\mathscr{A})$ gives
\[
\left|E\left(X;\mathscr{A},H_k^{\#}(\mathscr{A})\Psi\right)\right|_{\mathrm{aw}_2}\le\left\|E(X;\mathscr{A})\right\|_{h_k}
\]
for every family $E$ without slot of level $k$, every $X\in D_k$ and every $\mathscr{A}\in\tilde{R}_k$.

\subsubsection*{Proposition 48.1 (the spaces without slot)}

\begin{named}{Proposition 48.1}[The spaces without slot]
Assume the setting (G2.S-1).

\emph{Definition.}
A \emph{family without slot} of level $k$ is a family
\[
E=\left\{E(X;\mathscr{A}):\ X\in D_k,\ \mathscr{A}\in\tilde{R}_k\right\}
\]
of elements of $\mathrm{Gr}[\psi,\chi]$.
Let $K_k^{C}=K_k^{C}(\ell)$ be the space of those families without slot whose elements are even and balanced, written with their kernels $E_{nm}$ as
\begin{multline*}
E(X;\mathscr{A})=\sum_{n,m\ge0}\frac{1}{n!\cdot m!}\int\mathrm{d}\xi\,\mathrm{d}\boldsymbol{\zeta}\;E_{nm}\left(X;\mathscr{A};\xi_1,\dots,\xi_n;\boldsymbol{\zeta}_1,\dots,\boldsymbol{\zeta}_m\right)\\
\cdot\psi(\xi_1)\cdots\psi(\xi_n)\cdot\chi(\boldsymbol{\zeta}_1)\cdots\chi(\boldsymbol{\zeta}_m),
\end{multline*}
which have the following properties.
\begin{itemize}
\item[(k1)$^{\mathrm{o}}$] The kernels are analytic functions of $\mathscr{A}\in\tilde{R}_k$.
\item[(k2)$^{\mathrm{o}}$] The kernels vanish unless every $\xi_i$ lies over $X$ and the pair of every $\boldsymbol{\zeta}_j$ is admissible for $X$;
they depend on $\mathscr{A}$ only through the values $\mathscr{A}(b)$ at the bonds $b$ that meet $X$.
\item[(k3)$^{\mathrm{o}}$] $a_{\tilde{g}}\left[E(\tilde{g}X;\tilde{g}\mathscr{A})\right]=E(X;\mathscr{A})$ for every $\tilde{g}\in\mathbb{G}_k$.
\item[(k4)$^{\mathrm{o}}$] $G_{\boldsymbol{\lambda}}\left[E(X;\mathscr{A}^1)\right]=E(X;\mathscr{A})$ whenever $\mathscr{A}$ and $\mathscr{A}^1$ lie in $\tilde{R}_k$, $\boldsymbol{\lambda}$ is a real function of the photon sites $c$, and $\mathscr{A}^1=\mathscr{A}-d\boldsymbol{\lambda}$ on the bonds at distance at most $1$ from $X$.
\item[(k5)$^{\mathrm{o}}$] $\mathrm{Cg}\left[E(X;-\mathscr{A})\right]=E(X;\mathscr{A})$.
\item[(k6)$^{\mathrm{o}}$] The number
\[
\|E\|_k:=\sup_{X\in D_k}\exp\left(\kappa\cdot d_M(X)\right)\cdot\sup_{\mathscr{A}\in\tilde{R}_k}\left\|E(X;\mathscr{A})\right\|_{h_k}
\]
is finite.
\end{itemize}
The real subspace $K_k$ of $K_k^{C}$ consists of the families that have, in addition, the property
\begin{itemize}
\item[(k7)$^{\mathrm{o}}$] $\mathbf{K}\left[E\left(X;-\overline{\mathscr{A}}\right)\right]=E(X;\mathscr{A})$.
\end{itemize}
For a family without slot the total connection is $\mathscr{T}:=\mathscr{A}$.
The fields to be inserted are those of Section~I.5.4 at $\mathscr{C}=0$.
At $k=0$ one has $H_0=\bar{H}_0=1$ (Section~I.3.9), no mass enters, and these fields are
\[
\psi=\Psi,\qquad\chi=\delta_{\alpha,\mathscr{A}}\Psi,\qquad\bar{\psi}=\bar{\Psi},\qquad\bar{\chi}=\delta_{\alpha,-\mathscr{A}}\bar{\Psi}.
\]
For $k\ge1$ they are formed when, besides $(k,\ell,e_k)$, a mass $\mu\in(0,1]$ is given, so that $(k,\ell,\mu,e_k)$ are data (SF) of Section~I.1.8 with $k\ge1$, that is, data of the setting (G2.S-3).
Then $\mathscr{A}$ and $-\mathscr{A}$ lie in $\tilde{R}_k$ (\cite[Lemma~17.6]{QEDI}), which is contained in $\tilde{G}_k$ (\cite[Lemma~17.4(a)]{QEDI});
hence the operators $H_k(\mathscr{A})$ and $\bar{H}_k(\mathscr{A})=\mathsf{C}^{-1}\cdot H_k(-\mathscr{A})\cdot\mathsf{C}$ of Section~I.3.5 exist by Proposition~14.1\,(i) (the paragraph ``Settings'' of Section~I.5.7), and the fields to be inserted are
\[
\psi=H_k(\mathscr{A})\Psi,\qquad\chi=\delta_{\alpha,\mathscr{A}}\psi,\qquad\bar{\psi}=\bar{H}_k(\mathscr{A})\bar{\Psi},\qquad\bar{\chi}=\delta_{\alpha,-\mathscr{A}}\bar{\psi}.
\]
The spaces $K_k^{C}$, $K_k$ and the clauses (a)--(d) below involve neither a mass nor these fields.
\begin{itemize}
\item[(a)] One has $0\in\mathfrak{X}_k$.
For $\mathscr{E}\in K^{x,C}_k$ the family $\mathscr{E}|_0$ given by $\mathscr{E}|_0(X;\mathscr{A}):=\mathscr{E}(X;\mathscr{A},0)$ lies in $K_k^{C}$, and it lies in $K_k$ if $\mathscr{E}\in K^{x}_k$.
The map $\mathscr{E}\mapsto\mathscr{E}|_0$ is linear, and $\left\|\mathscr{E}|_0\right\|_k\le\|\mathscr{E}\|^{x}_k$.
\item[(b)] For a family $E$ without slot let $E^{\mathrm{c}}(X;\mathscr{A},\mathscr{C}):=E(X;\mathscr{A})$ be the extension of $E$ that is constant in $\mathscr{C}$.
The family $E^{\mathrm{c}}$ has the property (k2) of Section~I.5.4 if and only if $E$ has (k2)$^{\mathrm{o}}$, and $E^{\mathrm{c}}|_0=E$.
\item[(c)] (Extension to the second slot.)
The family $E^{\mathrm{c}}$ need not have the property (k4) of Section~I.5.4:
for $\boldsymbol{\lambda}=0$ and a non-constant $\boldsymbol{\lambda}'$ the transformation in (k4) multiplies a term $\bar{\psi}(x)^T\cdot K(x,x')\cdot\psi(x')$ with $\pi(x)\ne\pi(x')$ by the factor
\[
\exp\left(i\cdot e_k\cdot\left(\boldsymbol{\lambda}'\left(\pi(x')\right)-\boldsymbol{\lambda}'\left(\pi(x)\right)\right)\right).
\]
The expression $E(X;\mathscr{A}+\mathscr{C})$ is not defined on all of $\tilde{R}_k\times\mathfrak{X}_k$:
the set $\tilde{R}_k+\mathfrak{X}_k$ is not contained in $\tilde{R}_k$, since the thresholds of (G2.r4) satisfy $\frac{1}{2}\cdot\mathrm{td1}_k>\mathrm{t1}_k$ and the constant field $\frac{i}{2}\cdot\mathrm{td1}_k$ lies in $\mathfrak{X}_k$.
\item[(d)] (Completeness.)
The space $K_k^{C}$ with the norm $\|\cdot\|_k$ and the space $K^{x,C}_k$ with the norm $\|\cdot\|^{x}_k$ are complex Banach spaces;
$K_k$ and $K^{x}_k$ are closed real subspaces of them, and hence real Banach spaces.
\end{itemize}
\end{named}

\begin{named}{Remark}
The spaces $K_k^{C}$, $K_k$ and the norm $\|\cdot\|_k$ correspond to those of \cite[Definition~5.24]{QEDI}, where they are defined through the requirements on the families with slot, taken with $\mathscr{C}=0$ throughout;
Section~I.5.4 defines them in the same way for families on $\tilde{R}_k$ alone, without a second slot.
In both places $K^{x,C}_k$ and $K_k^{C}$ are complex Banach spaces, and $K^{x}_k$, $K_k$ are real subspaces of them.
As for $K^{x,C}_k$ (see the Remark in Section~I.5.4), the two definitions differ in that here the kernels are functions on the index sets of the fermion lattice $T^{\mathrm{f}}_k$.
The completeness of $K^{x,C}_k$ is derived in \cite[Definition~5.24]{QEDI}:
the set $D_k$ is finite, a uniform limit of analytic maps from the open set $\tilde{R}_k\times\mathfrak{X}_k$ into the finite-dimensional algebra $\mathrm{Gr}[\psi,\chi]$ is analytic, and the requirements other than analyticity and the finiteness of the norm are linear conditions that persist under pointwise limits.
The completeness of $K_k$ is proved in the first paragraph of the proof of \cite[Proposition~19.14]{QEDI}.
Clause (d) gives the completeness of all four spaces, together with the closedness of $K^{x}_k$ in $K^{x,C}_k$ and of $K_k$ in $K_k^{C}$.
\end{named}

\subsubsection*{Proposition 48.2 (low kernels and extraction without slot)}

\begin{named}{Proposition 48.2}[Low kernels and extraction without slot]
Assume the setting (G2.S-1), and in (e) the setting (G2.S-2).
Let $E\in K_k^{C}$;
in (o), $E$ may be any balanced family without slot that is defined at $\mathscr{A}=0$ and whose kernels $E_{20}(X;0;x,x')$ vanish unless $x$ and $x'$ lie over $X$.
For a small polymer $X$ let
\[
\mathrm{al}_0(E,X):=\operatorname{Vol}(X)^{-1}\cdot E_{00}(X;0),\qquad\mathrm{al}_2(E,X):=\operatorname{Vol}(X)^{-1}\cdot\int\mathrm{d}x\,\mathrm{d}x'\;E_{20}(X;0;x,x'),
\]
for an $M$-cube $\square$ let
\[
\varepsilon(E,\square):=-\sum_{X\ \mathrm{small},\ X\supset\square}\mathrm{al}_0(E,X),\qquad m(E,\square):=-\sum_{X\ \mathrm{small},\ X\supset\square}\mathrm{al}_2(E,X),
\]
and let
\[
(R\,E)(X;\mathscr{A}):=\begin{cases}E(X;\mathscr{A})-\mathrm{al}_0(E,X)\cdot\operatorname{Vol}(X)-\displaystyle\int_X\mathrm{d}x\;\bar{\psi}(x)^T\cdot\mathrm{al}_2(E,X)\cdot\psi(x)&\text{if $X$ is small},\\[1ex] E(X;\mathscr{A})&\text{if $X$ is large}.\end{cases}
\]
\begin{itemize}
\item[(o)] The numbers $\mathrm{al}_0(E,X)$ and the matrices $\mathrm{al}_2(E,X)$ do not depend on the lift of $X$.
Let $\tilde{g}=(S,a)\in\mathbb{G}_k$.
If (k3)$^{\mathrm{o}}$ holds at $\mathscr{A}=0$, then
\[
\mathrm{al}_0(E,\tilde{g}X)=\mathrm{al}_0(E,X),\qquad\mathrm{al}_2(E,\tilde{g}X)=S\cdot\mathrm{al}_2(E,X)\cdot S^{-1}.
\]
If (k5)$^{\mathrm{o}}$ holds at $\mathscr{A}=0$, then $\mathrm{al}_2(E,X)=\mathsf{C}^{-1}\cdot\mathrm{al}_2(E,X)^T\cdot\mathsf{C}$.
If (k7)$^{\mathrm{o}}$ holds at $\mathscr{A}=0$, then $\mathrm{al}_0(E,X)$ is real and $\mathrm{al}_2(E,X)=\varkappa^{-1}\cdot\overline{\mathrm{al}_2(E,X)}\cdot\varkappa$.
\item[(a)] Neither $\varepsilon(E,\square)$ nor $m(E,\square)$ depends on $\square$, and $m(E,\square)=m({E})\cdot\mathbf{1}$ with a complex number $m({E})$, where $\mathbf{1}$ is the $4\times4$ identity matrix;
the common value of the numbers $\varepsilon(E,\square)$ is denoted $\varepsilon({E})$.
The maps $E\mapsto\varepsilon({E})$ and $E\mapsto m({E})$ are complex linear on $K_k^{C}$ and take real values on $K_k$.
For $\mathscr{E}\in K^{x,C}_k$ and every small $X$,
\[
\mathrm{al}_0(\mathscr{E},X)=\mathrm{al}_0(\mathscr{E}|_0,X),\quad\mathrm{al}_2(\mathscr{E},X)=\mathrm{al}_2(\mathscr{E}|_0,X),\quad\varepsilon(\mathscr{E})=\varepsilon(\mathscr{E}|_0),\quad m(\mathscr{E})=m(\mathscr{E}|_0).
\]
\item[(b)] One has
\begin{align*}
|\varepsilon({E})|&\le K_0\cdot M^{-3}\cdot\|E\|_k\le c_{\mathrm{fl}}\cdot\|E\|_k,\\
|m({E})|&\le K_0\cdot M^{-3}\cdot e_k^{\frac{1}{2}}\cdot\|E\|_k\le c_{\mathrm{fl}}\cdot e_k^{\frac{1}{2}}\cdot\|E\|_k,
\end{align*}
and, for every small $X$,
\[
\operatorname{Vol}(X)\cdot\left|\mathrm{al}_0(E,X)\right|\le\left\|E(X;0)\right\|_{h_k},\qquad\operatorname{Vol}(X)\cdot\sum_{\beta,\beta'}\left|\mathrm{al}_2(E,X)_{\beta\beta'}\right|\le e_k^{\frac{1}{2}}\cdot\left\|E^{[2,0]}(X;0)\right\|_{h_k}.
\]
\item[(c)] For every $\mathscr{A}\in\tilde{R}_k$,
\[
\sum_{X\in D_k}E(X;\mathscr{A})=-\varepsilon({E})\cdot n_k-m({E})\cdot\langle\bar{\psi},\psi\rangle+\sum_{X\in D_k}(R\,E)(X;\mathscr{A}).
\]
\item[(d)] The map $R$ is linear and maps $K_k^{C}$ into $K_k^{C}$ and $K_k$ into $K_k$.
For all $X\in D_k$ and $\mathscr{A}\in\tilde{R}_k$,
\[
\left\|(R\,E)(X;\mathscr{A})\right\|_{h_k}\le\left\|E(X;\mathscr{A})\right\|_{h_k}+\left\|E(X;0)\right\|_{h_k},
\]
and $\|R\,E\|_k\le2\,\|E\|_k\le c_{\mathrm{fl}}\cdot\|E\|_k$.
For every small $X$,
\[
(R\,E)_{00}(X;0)=0,\qquad\int\mathrm{d}x\,\mathrm{d}x'\;(R\,E)_{20}(X;0;x,x')=0.
\]
If a family without slot has (k2)$^{\mathrm{o}}$, so has its image under $R$.
\item[(e)] Let $E'$ be a family without slot of level $k$, and for $Y\in D_{k+1}$ and $\mathscr{A}\in\tilde{R}_{k+1}$ let
\[
(L\,E')(Y;\mathscr{A}):=\sum_{X:\ \bar{X}^{L}=LY}S_L\left[E'(X;\mathscr{A}_L)\right].
\]
If $E'$ has (k3)$^{\mathrm{o}}$ for the elements of $\mathbb{G}_k^{L}$, then $L\,E'$ has (k3)$^{\mathrm{o}}$ at the level $k+1$, for the group $\mathbb{G}_{k+1}$.
If $E'$ has (k4)$^{\mathrm{o}}$, resp. (k5)$^{\mathrm{o}}$, (k7)$^{\mathrm{o}}$, then $L\,E'$ has the same property.
\item[(f)] The following hold.
\begin{itemize}
\item[(f1)] Proposition~16.3\,(f1), as stated there.
\item[(f2)] Proposition~16.3\,(f2), with $E(X;0)$ and $E(X;\mathscr{A})$ in place of $\mathscr{E}(X;0,0)$ and $\mathscr{E}(X;\mathscr{A},\mathscr{C})$.
\item[(f3)] Proposition~16.3\,(f3), with the same replacement and with ``a family without slot on $T_k(\ell_i)$ as in (o)'' in place of ``a family on $T_k(\ell_i)$ as in Proposition~16.2''; here $T_k(\ell_i)$ is the lattice $T_k$ formed with the side $\ell_i$.
\end{itemize}
\end{itemize}
\end{named}

\begin{named}{Remark}
Proposition~48.2 corresponds to the last sentence of \cite[Theorem~13.5]{QEDI} and to the last sentence of \cite[Proposition~5.59]{QEDI}, according to which the assertions of these statements remain true for the families without the second slot;
the numbers $\varepsilon({E})$, $m({E})$ and the map $R$ without slot are those that enter the recursion of \cite[Section~5.3.8]{QEDI}.
The proof of \cite[Theorem~13.5]{QEDI} applies to the families without slot as it stands, since \cite[Lemma~13.4]{QEDI} requires of the family only that it be defined at the zero field.
The second sentence of \cite[Theorem~13.15(c)]{QEDI} asserts in the same way that the assertions of \cite[Theorem~13.15]{QEDI} hold without the second slot.
These assertions are true with the constants printed in \cite[Theorem~13.15]{QEDI}: without slot they are Proposition~48.5\,(a), (b) and the last inequality of Proposition~48.5\,(d).
Their proof is the part of the proof of \cite[Theorem~13.15]{QEDI} on the families without the second slot, which carries out the argument with $\mathscr{C}=0$ and with the total connection $\mathscr{T}=\mathscr{A}$;
besides the last sentence of \cite[Theorem~13.5]{QEDI}, it refers to the forms for one slot of \cite[Theorem~13.7(b), (e)]{QEDI}, \cite[Theorem~13.11(ii), (iii)]{QEDI} and \cite[Theorem~13.13(c)--(e)]{QEDI}.
\cite{QEDI} states these three theorems for the pairs $(\mathscr{A},\mathscr{C})$ and the families with slot and, beside these forms, for $\mathscr{A}$ alone, that is at $\mathscr{C}=0$, and for the families without slot:
in the last assertion of \cite[Theorem~13.7]{QEDI}, in the last sentence before (i) and in the second half of (iii) of \cite[Theorem~13.11]{QEDI}, and in the last assertion of \cite[Theorem~13.13]{QEDI}.
For the present fermion, the form for one slot of \cite[Theorem~13.13(c)]{QEDI} is the last sentence of Proposition~48.2\,(d);
the forms for one slot of the other statements of \cite[Theorem~13.7]{QEDI}, \cite[Theorem~13.11]{QEDI} and \cite[Theorem~13.13]{QEDI} just named are contained in Propositions~48.3 and~48.4.
In this respect Propositions~48.2\,(d), 48.3 and~48.4 agree with \cite{QEDI};
they differ from it in the fermion and in the wider class of families admitted in clauses (a), (b), (d) and in part of clause (c) of Proposition~48.3 (see the Remark after Proposition~48.3 and, for the forms with slot, the Remark after Proposition~17.2).
Propositions~48.1--48.4 are proved in Chapter~48a, and Proposition~48.5 is proved with them in Chapter~48b.
\end{named}

\subsubsection*{Proposition 48.3 (adjustment without slot)}

\begin{named}{Proposition 48.3}[Adjustment without slot]
Assume the setting (G2.S-1).
Let $X$ be a small polymer, let $\mathscr{A}\in\tilde{R}_k$, and let $E$ be a family without slot with (k2)$^{\mathrm{o}}$.
Its adjustment $E^{\mathrm{nat}}$ is the family with
\[
E^{\mathrm{nat}}(X;\mathscr{A}):=E(X;\mathscr{A})-E^{[2,0]}(X;\mathscr{A})+\boldsymbol{\Theta}_{X,\mathscr{A}}\left(E^{[2,0]}(X;\mathscr{A})\right)
\]
for small $X$, where $\boldsymbol{\Theta}_{X,\mathscr{A}}$ is the homomorphism $\boldsymbol{\Theta}_{X,B}$ of Proposition~17.1 at $B=\mathscr{A}$, and with $E^{\mathrm{nat}}(X):=E(X)$ for large $X$.
\begin{itemize}
\item[(a)] Let $\mathrm{ev}$ be a homomorphism of $\mathrm{Gr}[\psi,\chi]$ into an associative algebra such that the images of the generators satisfy the relations (u4) of Proposition~17.1, with $B=\mathscr{A}$, at all pairs of kind (I) for $X$.
Then
\[
\mathrm{ev}\left(E^{\mathrm{nat}}(X;\mathscr{A})\right)=\mathrm{ev}\left(E(X;\mathscr{A})\right).
\]
In particular $E^{\mathrm{nat}}$ and $E$ agree after the insertion of $\chi=\delta_{\alpha,\mathscr{A}}\psi$ and $\bar{\chi}=\delta_{\alpha,-\mathscr{A}}\bar{\psi}$.
\item[(b)] Let $E$ be balanced.
If $\int\mathrm{d}x\,\mathrm{d}x'\;E_{20}(X;0;x,x')=0$, then $E^{\mathrm{nat}}_{20}(X;0;\cdot,\cdot)=0$.
For $(n,m)\notin\{(2,0),(1,1),(0,2)\}$ one has $E^{\mathrm{nat}}_{nm}=E_{nm}$; in particular $E^{\mathrm{nat}}_{00}=E_{00}$.
\item[(c)] The element $E^{\mathrm{nat}}(X;\mathscr{A})$ contains the generators $\psi(\xi)$ only for $\xi$ over $X$ and the generators $\chi(\boldsymbol{\zeta})$ only for $\boldsymbol{\zeta}$ admissible for $X$.
Every generator $\chi(\boldsymbol{\zeta})$ that occurs in $E^{\mathrm{nat}}(X;\mathscr{A})$ but not in $E(X;\mathscr{A})$ sits at a chain pair $(p_i,p_{i+1})$ of Proposition~17.1\,(c) with $p_i\ne p_{i+1}$;
hence it is of kind (I) for $X$, with $d_X\le M$.
The coefficients depend on $\mathscr{A}$ through the kernels $E_{nm}(X;\mathscr{A})$ and through the values of $\mathscr{A}$ on the bonds with both end points in $X$.
The family $E^{\mathrm{nat}}$ has (k2)$^{\mathrm{o}}$.
\item[(d)] Let $E$ be balanced.
With all kernels taken at $\mathscr{A}$,
\begin{align*}
\left\|E^{\mathrm{nat}}_{20}(X)\right\|&\le\mathrm{e}^2\cdot\left\|E_{20}(X)\right\|,\\
\left\|E^{\mathrm{nat}}_{11}(X)\right\|&\le\left\|E_{11}(X)\right\|+\mathrm{e}^2\cdot(3c_X\cdot L\cdot M)\cdot\left\|E_{20}(X)\right\|,\\
\left\|E^{\mathrm{nat}}_{02}(X)\right\|&\le\left\|E_{02}(X)\right\|+\mathrm{e}^2\cdot(3c_X\cdot L\cdot M)^2\cdot\left\|E_{20}(X)\right\|,
\end{align*}
and
\[
\left\|E^{\mathrm{nat}}(X;\mathscr{A})\right\|_{h_k}\le\mathrm{e}^2\cdot\left(1+3c_X\cdot L\cdot M\cdot e_k^{\varepsilon}\right)^2\cdot\left\|E(X;\mathscr{A})\right\|_{h_k}\le10\cdot\left\|E(X;\mathscr{A})\right\|_{h_k}.
\]
\item[(e)] If $E$ lies in $K_k^{C}$, resp. in $K_k$, then so does $E^{\mathrm{nat}}$, and $\|E^{\mathrm{nat}}\|_k\le10\cdot\|E\|_k$.
The map $E\mapsto E^{\mathrm{nat}}$ is linear.
For $\mathscr{E}\in K^{x,C}_k$ one has $(\mathscr{E}^{\mathrm{nat}})|_0=(\mathscr{E}|_0)^{\mathrm{nat}}$.
\item[(f)] Proposition~17.2\,(f) holds with $B=\mathscr{A}$ and with the kernels of $E(X;\mathscr{A})$.
\end{itemize}
\end{named}

\begin{named}{Remark}
\cite[Theorem~13.7]{QEDI} and \cite[Proposition~5.60]{QEDI} each end with their form without the second slot:
their assertions (a)--(d), and those of their clause (e) other than the identity $(\mathscr{E}^{\mathrm{nat}})|_0=(\mathscr{E}|_0)^{\mathrm{nat}}$, hold for $E\in K_k^{C}$, a small polymer $X$ and $\mathscr{A}\in\tilde{R}_k$;
in these last assertions the adjustment of a family without the second slot is defined with the total connection $\mathscr{T}=\mathscr{A}$.
This form is proved in the last part of the proof of \cite[Theorem~13.7]{QEDI}.
Proposition~48.3 is the form of this assertion for the present fermion, with the same constants.
Apart from the fermion, the two differ in the following points.
Clauses (a) and (c) hold for every family without slot with (k2)$^{\mathrm{o}}$, and clauses (b) and (d) for every balanced such family, whereas the last assertions of \cite[Theorem~13.7]{QEDI} and of \cite[Proposition~5.60]{QEDI} take $E\in K_k^{C}$.
The locality part of (c), namely its first and its last sentence and the assertion that the new generators are of kind (I) for $X$, is asserted in \cite{QEDI} for the same class of families:
the form of \cite[Theorem~13.13(d)]{QEDI} without the second slot, in the last assertion of that theorem, states for every family without slot with (k2)$^{\mathrm{o}}$ that its adjustment, formed with $\mathscr{T}=\mathscr{A}$, has (k2)$^{\mathrm{o}}$, and that every generator $\chi$ which occurs in $E^{\mathrm{nat}}(X;\mathscr{A})$ but not in $E(X;\mathscr{A})$ is of kind (I) for $X$ and of length at most $M$.
The wider class is therefore a difference from \cite{QEDI} for the clauses (a), (b), (d) and for the remaining assertions of (c), that is, for the position of the new generators at the chain pairs and for the dependence of the coefficients on $\mathscr{A}$.
In (c) the new generators sit at the chain pairs of Proposition~17.1\,(c), as in \cite[Theorem~13.7(c)]{QEDI}, and the bound $d_X\le M$ takes the place of the bound $M$ on their length, which \cite{QEDI} gives without the second slot in the last assertion of \cite[Theorem~13.13]{QEDI}.
Clause (e) contains, in addition, the identity $(\mathscr{E}^{\mathrm{nat}})|_0=(\mathscr{E}|_0)^{\mathrm{nat}}$ for $\mathscr{E}\in K^{x,C}_k$, which \cite[Theorem~13.7(e)]{QEDI} states among the assertions with slot.
Clause (f) has no counterpart in the last assertions of \cite[Theorem~13.7]{QEDI} and of \cite[Proposition~5.60]{QEDI}.
It is the form without slot of Proposition~17.2\,(f), whose counterpart in \cite{QEDI} is the second half of \cite[Remark~13.16]{QEDI};
that remark is stated in \cite{QEDI} with the second slot only, for $\mathscr{E}\in K^{x,C}_k$, for pairs $(\mathscr{A},\mathscr{C})\in\tilde{R}_k\times\mathfrak{X}_k$ and with the total connection $\mathscr{T}=\mathscr{A}+\mathscr{C}$.
For the place of Proposition~48.3 in the contraction without slot see the Remark after Proposition~48.2.
\end{named}

\subsubsection*{Proposition 48.4 (global fields and locality without slot)}

\begin{named}{Proposition 48.4}[Global fields and locality without slot]
Assume the setting (G2.S-2).
In (a) and (b) let $X\in D_k$ be small, let $\tilde{c}$ be a photon site of $X$, let $\mathscr{A}\in\tilde{R}_{k+1}$, and let $\mathrm{bh}$, $\mathrm{kc}_1$, $\mathrm{mg}$ be the objects of Proposition~18.2 for the pair $(\mathscr{A},0)$;
by Proposition~18.2\,(o) they depend on $(\mathscr{A},\tilde{c})$ only, and $\mathscr{C}^b=0$, $\mathrm{kc}_2=0$.
In (c) let $X$ be any polymer of $D_k$, small or large, and let $Y\in D_{k+1}$;
let $E$ and $E'$ be families without slot of level $k$;
the first assertion of (c), that (k2)$^{\mathrm{o}}$ states that $E(X;\cdot)$ is local in $X$, is meant also at the level $k+1$, with $\mathrm{Dc}=\tilde{R}_{k+1}\times\{0\}$, for a family of level $k+1$ and a polymer of $D_{k+1}$;
let $F$ be a map from $\tilde{R}_k\times\{0\}$ to the algebra $\mathrm{Gr}[\psi,\chi]$ of level $k$;
$F(\mathscr{A}_L)$ stands for $F(\mathscr{A}_L,0)$, and $E(X;\cdot)$ for the map $(\mathscr{A},0)\mapsto E(X;\mathscr{A})$.
\begin{itemize}
\item[(a)] On the bonds with both end points in $K(\tilde{c},2L\cdot M)$ one has
\[
\mathscr{A}_L-d\,\mathrm{mg}=\mathrm{bh}+\mathrm{kc}_1.
\]
For every complex $t$ with $|t|\le L^{\frac{7}{4}}$ the bond function $t\cdot\mathrm{bh}$ lies in $R_k$, and hence in $\tilde{R}_k$.
\item[(b)] Let $E$ be a family without slot of level $k$ with (k1)$^{\mathrm{o}}$, (k2)$^{\mathrm{o}}$ and (k4)$^{\mathrm{o}}$.
For all $(n,m)$ and all arguments $\xi$, $\boldsymbol{\zeta}$,
\[
\left|E_{nm}(X;\mathscr{A}_L;\xi,\boldsymbol{\zeta})\right|\le\exp(n+m)\cdot\left|E_{nm}(X;\mathrm{bh};\xi,\boldsymbol{\zeta})\right|\le\exp\left(2(n+m)\right)\cdot\left|E_{nm}(X;\mathrm{bh};\xi,\boldsymbol{\zeta})\right|,
\]
the first inequality arising from one factor of modulus at most $\mathrm{e}$ for each argument;
moreover
\[
E_{00}(X;\mathscr{A}_L)=E_{00}(X;\mathrm{bh}),\qquad\left\|E_{nm}(X;\mathscr{A}_L)\right\|\le\exp\left(2(n+m)\right)\cdot\left\|E_{nm}(X;\mathrm{bh})\right\|.
\]
\item[(c)] Take $\mathrm{Dc}=\tilde{R}_k\times\{0\}$, resp. $\mathrm{Dc}=\tilde{R}_{k+1}\times\{0\}$, in the definition of Proposition~18.3.
Then (k2)$^{\mathrm{o}}$ states that $E(X;\cdot)$ is local in $X$.
If $X\subset LY$ and $F$ is local in $X$ on $\tilde{R}_k\times\{0\}$, without spectators, then the map $\mathscr{A}\mapsto S_L\left[F(\mathscr{A}_L)\right]$ is local in $Y$ on $\tilde{R}_{k+1}\times\{0\}$.
Consequently $L\,E'$ has (k2)$^{\mathrm{o}}$ at the level $k+1$ if $E'$ has (k2)$^{\mathrm{o}}$ at the level $k$.
\end{itemize}
\end{named}

\begin{named}{Remark}
The corresponding statements of \cite{QEDI} are \cite[Theorem~13.11]{QEDI}, for (a) and (b), and \cite[Theorem~13.13(e)]{QEDI}, for (c).
Each of them is stated there for the pairs $(\mathscr{A},\mathscr{C})$ and, in addition, for one slot:
\cite[Theorem~13.11]{QEDI}, and likewise \cite[Proposition~5.61]{QEDI}, in the last sentence before (i), according to which $\mathscr{C}^b=0$ and $c_2=0$ if $\mathscr{C}=0$, and in the second half of (iii), for the families defined on $\tilde{R}_k$ alone;
\cite[Theorem~13.13(e)]{QEDI} in the last assertion of that theorem, for the maps on $\tilde{R}_k\times\{0\}$ and $\tilde{R}_{k+1}\times\{0\}$.
Proposition~48.4 agrees with these forms for one slot:
(a) is the case $\mathscr{C}=0$ of \cite[Theorem~13.11(i), (ii)]{QEDI};
(b) is the second half of \cite[Theorem~13.11(iii)]{QEDI}, with the same inequalities and the same factors $\exp(n+m)\le\exp\left(2(n+m)\right)$;
(c) is the form of \cite[Theorem~13.13(e)]{QEDI} without the second slot given in the last assertion of that theorem, with the same domains $\tilde{R}_k\times\{0\}$ and $\tilde{R}_{k+1}\times\{0\}$.
The difference lies in the fermion, as described in the Remark after Proposition~18.2:
the base site is a photon site $\tilde{c}$ of $X$, where \cite{QEDI} takes a site $x_0$ of $X$;
the arguments are indices over $T^{\mathrm{f}}_k$ and $\mathrm{Pz}_k$;
and the phases are read at the photon site of the site, resp.\ of the first site, of an argument.
The forms for one slot of \cite[Theorem~13.13(c)]{QEDI} and of \cite[Theorem~13.13(d)]{QEDI}, which stand in the same last assertion of that theorem, are, for the present fermion, the last sentence of Proposition~48.2\,(d) and Proposition~48.3\,(c), respectively.
For the place of Proposition~48.4 in the contraction without slot see the Remark after Proposition~48.2.
\end{named}

\subsubsection*{Proposition 48.5 (contraction without slot)}

\begin{named}{Proposition 48.5}[Contraction without slot]
Assume the setting (G2.S-2), and let $E\in K_k^{C}$.
Put $L\,R\,E:=L\left((R\,E)^{\mathrm{nat}}\right)$, that is,
\[
(L\,R\,E)(Y;\mathscr{A})=\sum_{X:\ \bar{X}^{L}=LY}S_L\left[(R\,E)^{\mathrm{nat}}(X;\mathscr{A}_L)\right]\qquad\left(Y\in D_{k+1},\ \mathscr{A}\in\tilde{R}_{k+1}\right).
\]
\begin{itemize}
\item[(a)] The map $L\,R$ is linear; it maps $K_k^{C}(\ell)$ into $K_{k+1}^{C}(\ell/L)$ and $K_k(\ell)$ into $K_{k+1}(\ell/L)$.
\item[(b)] One has
\[
\|L\,R\,E\|_{k+1}\le c_{\mathrm{fl}}\cdot L^{-\frac{1}{4}+2\varepsilon}\cdot\|E\|_k.
\]
\item[(c)] For $\mathscr{E}\in K^{x,C}_k$ one has $(L\,R\,\mathscr{E})|_0=L\,R\,(\mathscr{E}|_0)$.
\item[(d)] Let $E^{a}:=(R\,E)^{\mathrm{nat}}$.
Then
\[
\sup_{\mathscr{A}\in\tilde{R}_k}\left\|E^{a}(X;\mathscr{A})\right\|_{h_k}\le\begin{cases}20\cdot\|E\|_k\cdot\exp\left(-\kappa\cdot d_M(X)\right)&\text{for small $X$},\\[1ex]\|E\|_k\cdot\exp\left(-\kappa\cdot d_M(X)\right)&\text{for large $X$}.\end{cases}
\]
For every small $X$ one has $E^{a}_{00}(X;0)=0$ and $E^{a}_{20}(X;0;\cdot,\cdot)=0$, and, for every $\mathscr{A}\in\tilde{R}_{k+1}$,
\[
\left\|S_L\left[E^{a}(X;\mathscr{A}_L)\right]\right\|_{h_{k+1}}\le3L^{-\frac{13}{4}+\frac{3\varepsilon}{2}}\cdot\sup_{\tilde{R}_k}\left\|E^{a}(X;\cdot)\right\|_{h_k}.
\]
\item[(e)] Proposition~18.5\,(f) holds with the kernels of $E(X;0)$ and of $E(X;\mathscr{A}_L)$, and with $\mathscr{A}_L$ on the bonds with both end points in $\tilde{X}$.
In particular, both margins of Proposition~18.5\,(f) hold for small $X$ and, for the two sides $\ell=\ell_1$ and $\ell=\ell_2$ of the standing data (TT) of the Propositions~27.m, for the sets $\tilde{X}$ that are admissible in the sense of the Definitions in Proposition~27.1.
\item[(f)] Proposition~18.6 holds without slot:
the family $E^{\mathrm{Ng}}$, read as a family without slot, lies in $K_k^{C}$, and in $K_k$ under (n3);
one has $\|E^{\mathrm{Ng}}\|_k=M^3\cdot e_k^{-\frac{1}{2}}\cdot|\mathrm{Ng}|_1$;
and the clauses (b)--(g) of Proposition~18.6 hold with the norm $\|\cdot\|_{k+1}$.
Here Proposition~18.6\,(a)--(g) is read, for every function $\mathrm{Ng}$ that occurs there, with the following objects in the places of the objects with slot:
\begin{itemize}
\item the family $E^{\mathrm{Ng}}|_0$, given by $E^{\mathrm{Ng}}|_0(X;\mathscr{A}):=E^{\mathrm{Ng}}(X;\mathscr{A},0)$, in place of $E^{\mathrm{Ng}}$;
\item $\mathscr{A}$ in place of $(\mathscr{A},\mathscr{C})$, and $0$ in place of $(0,0)$;
\item $\tilde{R}_k$ and $\tilde{R}_{k+1}$ in place of $\tilde{R}_k\times\mathfrak{X}_k$ and $\tilde{R}_{k+1}\times\mathfrak{X}_{k+1}$;
\item $K_k^{C}$ and $K_k$ in place of $K^{x,C}_k$ and $K^{x}_k$;
\item $\|\cdot\|_k$ and $\|\cdot\|_{k+1}$ in place of $\|\cdot\|^{x}_k$ and $\|\cdot\|^{x}_{k+1}$;
\item (k5)$^{\mathrm{o}}$ in place of (k5);
\end{itemize}
and $R$, $\mathrm{nat}$ and $L\,R$ are the maps without slot of Proposition~48.2, of Proposition~48.3 and of the present proposition.
\item[(g)] (Sum over the polymers: the first sentence of Proposition~18.5\,(e), without slot.)
For every family $E'$ without slot of level $k$ and every $\mathscr{A}\in\tilde{R}_{k+1}$, with $L\,E'$ as in Proposition~48.2\,(e),
\[
\sum_{Y\in D_{k+1}}(L\,E')(Y;\mathscr{A})=S_L\left[\sum_{X\in D_k}E'(X;\mathscr{A}_L)\right].
\]
\end{itemize}
\end{named}

\begin{named}{Remark}
Proposition~48.5 corresponds to the second sentence of \cite[Theorem~13.15(c)]{QEDI} together with its proof, the part of the proof of \cite[Theorem~13.15]{QEDI} on the families without the second slot, which carries the argument through with the forms for one slot of \cite[Theorem~13.5]{QEDI}, \cite[Theorem~13.7]{QEDI}, \cite[Theorem~13.11]{QEDI} and \cite[Theorem~13.13]{QEDI} (see the Remark after Proposition~48.2);
it corresponds also to the sentence of \cite[Proposition~5.62]{QEDI} according to which the same assertions hold without the second slot;
the map $L\,R$ without slot is the one that enters the recursion of \cite[Section~5.3.8]{QEDI}.
Clause (d) agrees with the bounds and identities that this part of the proof of \cite[Theorem~13.15]{QEDI} establishes.
In the argument for (d), as in that part of the proof, with $\mathrm{bh}$ as in Proposition~48.4, the disc is $\left\{t\cdot\mathrm{bh}:\ t\in\mathbb{C},\ |t|\le L^{\frac{7}{4}}\right\}$, the functions are $t\mapsto E^{a}_{00}(X;t\cdot\mathrm{bh})$, which is even in $t$ by (k5)$^{\mathrm{o}}$, and $t\mapsto E^{a}_{20}(X;t\cdot\mathrm{bh};\cdot,\cdot)$, and the suprema extend over $\tilde{R}_k$.
The constants $3$, $20$, $60$ and $61\,c_{\mathrm{rb}}=c_{\mathrm{fl}}$ are those of \cite[Theorem~13.15]{QEDI} and of its proof.
\end{named}

\subsection*{G3. Definitions for the Propositions 21.m--25.m}

The settings and definitions on which the Propositions 21.m--25.m rest are stated next to these propositions:
the setting (S21) in Section~G3.1.0,
the setting (S22) in Section~G3.2.0,
the setting (ST1)$^{\mathrm{ch}}$ in Section~G3.4.0,
Definition~24.D (the step map) in Section~G3.4.1,
Definition~24.I (invariance) in Section~G3.4.2,
and the covering objects in Section~G3.5.0.

\paragraph{Notation.}
In Section~G3 and in the Propositions 21.m--25.m, $J_k^{\#}$ denotes the operator of \cite[Section~17.4.2]{QEDI}, and the bare letter $J$ denotes the cluster map.
The components without H\"older derivative of $J_k^{\#}$ on the unbarred and on the barred fields, which \cite[Section~17.4.4]{QEDI} writes $J$ and $\bar{J}$, are written $J_k$ and $\bar{J}_k$.
The sources $J$, $\bar{J}$ of the Gaussian integration formula of \cite[Section~14.6.1]{QEDI} are written $\mathrm{Sr}$ and $\overline{\mathrm{Sr}}$.
The objects $\mathrm{Grc}_k$ and $\mathrm{cov}_{\ell}$ are those of Section~I.1.9; they are recalled in Section~G3.5.0.
The letters $\mathrm{ff}$, $\overline{\mathrm{ff}}$, $\mathrm{chf}$, $\mathrm{chfb}$, $\mathrm{Bf}$, $\mathrm{hf}$ and $\mathrm{gf}$ denote generic fields and bond functions, used as dummy variables.

\subsection*{21. Identities of one step (completion of the square, scaling relations, determinants), first step included}

\subsubsection*{G3.1.0. Setting (S21), used in Propositions 21.1--21.5}

\paragraph{Data of level $k$.}
The data of level $k$ are an integer $k\ge0$, a side $\ell\in\mathrm{Cl}$, a complex number $\mu$, a number $e_k>0$ and a complex bond function $B$ on $T_k$.
All operators of Section~I.3.3, of Section~I.3.5 (for $k\ge1$) and of Section~I.3.9 (for $k=0$) are formed with $(k,\ell,\mu,e_k)$, the complex number $\mu$ taking the place of the mass.
No relation among the four data $k$, $\ell$, $\mu$, $e_k$ is assumed.
Among the constants of Section~I.8, the objects of this setting and the Propositions~21.1--21.5 depend only on $L$, $M$, $b$, $\rho$, the weights $w$ and $\alpha$, on $M$ through $\mathrm{Cl}$, through $\mathrm{Pzt}_k$, through the range $|c-c'|\le M$ of $U_B$ and through $\mathbb{G}_k$, with one exception:
Proposition~21.4\,(d2) depends, in addition, on $e_0$ and, through $\tilde{G}_k$, on the exponent $\varepsilon$.
For an operator $O(B)$ that Section~I.3.5 or Section~I.3.9 defines through an inverse, the words ``$O(B)$ exists'' mean that this inverse exists.

\paragraph{Data of level $k+1$.}
At the level $k+1$ the side is $\ell/L$, the mass is $L\cdot\mu$ and the charge is $e_{k+1}:=L^{\frac{1}{2}}\cdot e_k$;
$B'$ is a complex bond function on $T_{k+1}=\mathbb{L}_{k+1}/(\ell/L)\mathbb{Z}^3$, and $B=B'_L$ (Section~I.3.7).
Since $\ell\in\mathrm{Cl}$, the number $\ell/L$ is an odd multiple of $M$ and $\ell/L\ge13M$;
hence $\ell/L$ is a multiple of $L$, $2M<\ell/L$, and the pairs of $\mathrm{Pz}_{k+1}$ at the side $\ell/L$ have unique shortest lifts.
The side $\ell/L$ need not lie in $\mathrm{Cl}$.
The operators of level $k+1$ that occur are those that contain no one-step average, namely
\begin{gather*}
D_{B'},\qquad Q_{k+1}(B')\cdot J^{*},\qquad J\cdot Q_{k+1}^{T}(-B'),\qquad J\cdot P_{k+1}(B')\cdot J^{*},\\
F_{\infty}(B')=D_{B'}+L\cdot\mu+b_{k+1}\cdot J\cdot P_{k+1}(B')\cdot J^{*},\\
S_{k+1}(B'),\qquad H_{k+1}(B'),\qquad\bar{H}_{k+1}(B'),\qquad D_{k+1}(B'),\qquad\delta_{\alpha,\pm B'}.
\end{gather*}
They are given by the displays of Section~I.3.3, of Section~I.3.5 and by (I.1.7), with $k+1$ in place of $k$, and they are distinguished from the operators of level $k$ by their argument $B'$.
The operators $Q(B')$, $Q^{T}(-B')$, $P(B')$, $\Gamma_{k+1}(B')$ and $\delta Z_{k+1}(B')$ do not occur.

\paragraph{Further notation.}
As in \cite[Lemma~14.19(b)]{QEDI} and \cite[Theorem~14.21]{QEDI}, let, for $k\ge1$,
\begin{align*}
Q_{k+1}(B)&:=Q(B)\cdot Q_k(B)\colon\ \mathfrak{S}(\eta_k,\ell)\to\mathfrak{S}(L,\ell),\\
Q_{k+1}^{T}(-B)&:=Q_k^{T}(-B)\cdot Q^{T}(-B)\colon\ \mathfrak{S}(L,\ell)\to\mathfrak{S}(\eta_k,\ell),
\end{align*}
and, for $k=0$, let $Q_1(B):=Q(B)$ and $Q_1^{T}(-B):=Q^{T}(-B)$, which act between $\mathfrak{S}(1,\ell)$ and $\mathfrak{S}(L,\ell)$.
Then
\[
P_{k+1}(B)=Q_{k+1}^{T}(-B)\cdot Q_{k+1}(B)\qquad(k\ge0);
\]
for $k=0$ the operator $J\cdot P_{1}(B)\cdot J^{*}$ is $J\cdot P(B)\cdot J^{*}$, the operator $J\cdot P_1(A)\cdot J^{*}$ of Section~I.3.9 with $B$ in place of $A$, where $P_1(A):=P(A)$ on $\mathfrak{S}(1,\ell)$.
Further, let
\begin{align*}
D^{0}_{k+1}(B)&:=b_L-b_L^2\cdot Q(B)\cdot\Gamma_k(B)\cdot Q^{T}(-B)&&(k\ge1),\\
D^{0}_{1}(B)&:=b_L-b_L^2\cdot Q(B)\cdot J^{*}\cdot\Gamma_0(B)\cdot J\cdot Q^{T}(-B)&&(k=0);
\end{align*}
these are operators on $\mathfrak{S}(L,\ell)$.

\paragraph{Fields.}
Let $\mathrm{Gr}$ be a finite-dimensional Grassmann algebra over $\mathbb{C}$.
A field on one of the lattices $T^{\mathrm{f}}_k$, $T^0_k$, $L\mathbb{Z}^3/\ell\mathbb{Z}^3$, $T^0_{k+1}$, $T^{\mathrm{f}}_{k+1}$, or on $\mathrm{Pz}_k$ or $\mathrm{Pz}_{k+1}$, is a map from the covering index set, together with the spinor index $\beta$, to the odd part of $\mathrm{Gr}$ that has the antiperiodicity (I.AP) of Section~I.3.6.
Linear operators act on fields by linearity.
The pairings $\langle\cdot,\cdot\rangle^{\mathrm{ch}}$, $\langle\cdot,\cdot\rangle_1$ and $\langle\cdot,\cdot\rangle_L$ of Section~I.1.6 are formed with the factors in the order in which they are written, and for fields
\[
\langle f'',O f'\rangle=\langle O^{T}f'',f'\rangle,
\]
no factors being exchanged.
The scalings $f\mapsto f_L$, $v\mapsto v_L$ and $F\mapsto F_L$ of Section~I.3.7 are applied to fields in the same way;
thus
\[
v_L(y'):=L^{-1}\cdot v(y'/L)\quad\left(y'\in L\mathbb{Z}^3\right),\qquad F_L(x,x'):=L^{-1-\alpha}\cdot F(x/L,x'/L)
\]
for functions $F$ of pairs.
On $T^0_k$ and on $T^{\mathrm{f}}_0$ these definitions are used for every integer side $\ell\ge1$, and on $L\mathbb{Z}^3/\ell\mathbb{Z}^3$ and on $T^0_{k+1}$ for every $\ell\in L\mathbb{N}$.
A field is then a map $f$ from $\mathbb{Z}^3$, $\mathbb{Z}^3\times O_{\mathrm{cub}}$, $L\mathbb{Z}^3$, resp. $\mathbb{Z}^3$, together with the index $\beta$, to the odd part of $\mathrm{Gr}$ such that $f(t+\ell\cdot e_{\mu})=-f(t)$, and on $T^0_{k+1}$ such that $f\left(t+(\ell/L)\cdot e_{\mu}\right)=-f(t)$;
the pairings $\langle\cdot,\cdot\rangle_1$, $\langle\cdot,\cdot\rangle^{\mathrm{ch}}$, $\langle\cdot,\cdot\rangle_L$ and the scaling $v\mapsto v_L$ are given by the same displays.

\subsubsection*{Proposition~21.1. Elementary relations}

\begin{named}{Proposition 21.1}
Let the data be those of the setting (S21).
In clauses (c), (d) and (e), as in clause (a), a relation that contains $Q_k$ or $P_k$ is asserted for $k\ge1$, and a relation that contains $Q(\,\cdot\,)\cdot J^*$, $J\cdot Q^T(\,\cdot\,)$ or $J\cdot P(\,\cdot\,)\cdot J^*$ is asserted for $k=0$;
the relations for $Q$ and $Q^T$ between $\mathfrak{S}(1,\ell)$ and $\mathfrak{S}(L,\ell)$ hold for every $k\ge0$.
\begin{itemize}
\item[(a)] \emph{Algebra.}
Let $B$ be arbitrary.
For $k\ge1$ one has
\begin{gather*}
\left(Q_k(B)\cdot J^*\right)\cdot\left(J\cdot Q_k^T(-B)\right)=1\ \text{ on }\mathfrak{S}(1,\ell),
\qquad
Q(B)\cdot Q^T(-B)=1\ \text{ on }\mathfrak{S}(L,\ell),\\
P(B)\cdot Q^T(-B)=Q^T(-B),
\qquad
Q(B)\cdot P(B)=Q(B),
\end{gather*}
and for $k=0$ one has
\begin{gather*}
\left(Q(B)\cdot J^*\right)\cdot\left(J\cdot Q^T(-B)\right)=1\ \text{ on }\mathfrak{S}(L,\ell),\\
\left(J\cdot P(B)\cdot J^*\right)\cdot\left(J\cdot Q^T(-B)\right)=J\cdot Q^T(-B),
\qquad
\left(Q(B)\cdot J^*\right)\cdot\left(J\cdot P(B)\cdot J^*\right)=Q(B)\cdot J^*.
\end{gather*}
For $k\ge1$ the operators $J\cdot P_k(B)\cdot J^*$, $J\cdot P_{k+1}(B)\cdot J^*$ and $P(B)$, and for $k=0$ the operator $J\cdot P(B)\cdot J^*$, are idempotent;
their kernels are multiples of $1_4$, and they commute with every constant $4\times4$ matrix.
If $k\ge1$, then in addition
\[
\left(J\cdot P_k(B)\cdot J^*\right)\cdot\left(J\cdot P_{k+1}(B)\cdot J^*\right)
=\left(J\cdot P_{k+1}(B)\cdot J^*\right)\cdot\left(J\cdot P_k(B)\cdot J^*\right)
=J\cdot P_{k+1}(B)\cdot J^*.
\]
Let $\mathrm{Mx}$ be a constant $4\times4$ matrix.
For $k\ge1$, the traces being taken on $\mathfrak{S}(1,\ell)$,
\[
\operatorname{Tr}\left[\mathrm{Mx}\cdot P(B)\right]=n_{k+1}\cdot\operatorname{tr}\mathrm{Mx},
\qquad
\operatorname{Tr}\left[\mathrm{Mx}\cdot\left(1-P(B)\right)\right]=\left(n_k-n_{k+1}\right)\cdot\operatorname{tr}\mathrm{Mx};
\]
for $k=0$, the traces being taken on $\mathfrak{S}^{\mathrm{ch}}(1,\ell)$,
\[
\operatorname{Tr}\left[\mathrm{Mx}\cdot J\cdot P(B)\cdot J^*\right]=n_1\cdot\operatorname{tr}\mathrm{Mx},
\qquad
\operatorname{Tr}\left[\mathrm{Mx}\cdot\left(1-J\cdot P(B)\cdot J^*\right)\right]=\left(48\,n_0-n_1\right)\cdot\operatorname{tr}\mathrm{Mx}.
\]
If $k\ge1$ and $S_k(B)$ exists, then $T_k(B)\cdot H_k(B)=1$ on $\mathfrak{S}(1,\ell)$;
for $k=0$ the relation $T_0(B)\cdot H_0(B)=1$ holds by Section~I.3.9.
\item[(b)] \emph{Composition.}
Let $k\ge0$, $y'\in L\mathbb{Z}^3$, $f\in\mathfrak{S}^{\mathrm{ch}}(\eta_k,\ell)$ and $v'\in\mathfrak{S}(L,\ell)$.
Let $\varphi^+_B$ be the phase of (I.3.3), which at $k=0$ is $\varphi^+_B(y',c)=U_B(y',c)$, and let
\[
\mathrm{Bl}_{k+1}(y')=\left\{c\in\mathbb{L}_k:\ |c-y'|<\frac{L}{2}\right\}.
\]
Then, $c=\pi(x)$ denoting the photon site under the fermion site $x$,
\begin{gather*}
\left(Q_{k+1}(B)\cdot J^*f\right)(y')=L^{-3}\cdot\int_{\mathrm{Bl}_{k+1}(y')}\mathrm{d}x\,\varphi^+_B(y',c)\cdot f(x),\\
\left(J\cdot Q_{k+1}^T(-B)\,v'\right)(x)=\varphi^+_B(y',c)^{-1}\cdot v'(y')\qquad\text{for }c\in\mathrm{Bl}_{k+1}(y'),\\
J\cdot P_{k+1}(B)\cdot J^*=\left(J\cdot Q_{k+1}^T(-B)\right)\cdot\left(Q_{k+1}(B)\cdot J^*\right).
\end{gather*}
\item[(c)] \emph{Transposition.}
Transposes are taken with respect to the pairings of Section~I.1.6.
One has
\begin{gather*}
D_B^T=\mathsf{C}^{-1}\cdot D_{-B}\cdot\mathsf{C};
\qquad
\left(Q_k(B)\cdot J^*\right)^T=J\cdot Q_k^T(B),\\
Q(B)^T=Q^T(B),
\qquad
\left(Q(B)\cdot J^*\right)^T=J\cdot Q^T(B);
\qquad
\left(J\cdot P_k(B)\cdot J^*\right)^T=J\cdot P_k(-B)\cdot J^*,
\end{gather*}
and the last relation holds in the same form for $J\cdot P_{k+1}(B)\cdot J^*$, for $P(B)$ and for $J\cdot P(B)\cdot J^*$.
Let $k\ge1$.
Then $F_\infty(B)$ is invertible if and only if $F_\infty(-B)$ is invertible;
in this case $S_k(B)$ and $S_k(-B)$ exist, and $D_k(B)+b_L\cdot P(B)$ is invertible if and only if $D_k(-B)+b_L\cdot P(-B)$ is invertible.
For $k=0$ the operator $F_0(B)$ is invertible if and only if $F_0(-B)$ is invertible.
Wherever the operators exist, one has
\[
O(B)^T=\mathsf{C}^{-1}\cdot O(-B)\cdot\mathsf{C}
\]
for $O=F_\infty,S_k,D_k,M_k,\dot{M}_k,\Gamma_k$ when $k\ge1$, and for $O=D_0,F_0,\Gamma_0$ when $k=0$.
Moreover, for $k\ge1$ and again wherever the operators exist,
\begin{gather*}
\bar{H}_k(B)=\mathfrak{H}_k(B)^T,
\qquad
\bar{T}_k(B)=b_k^{-1}\cdot Q_k(-B)\cdot J^*\cdot F_\infty(B)^T,\\
\bar{T}_k(B)\cdot\bar{H}_k(B)=1,
\qquad
\dot{M}_k(B)=\mathfrak{H}_k(B)\cdot H_k(B),\\
\left\langle\bar{H}_k(B)\bar{\Psi},H_k(B)\Psi\right\rangle^{\mathrm{ch}}=\left\langle\bar{\Psi},\dot{M}_k(B)\Psi\right\rangle_1
\quad\text{for all fields }\Psi,\ \bar{\Psi}\text{ on }T^0_k,
\end{gather*}
while for $k=0$ one has $\bar{H}_0=\bar{T}_0=\dot{M}_0=1$ on $\mathfrak{S}^{\mathrm{ch}}(1,\ell)$.
\item[(d0)] \emph{Phases.}
Let $B$ be a complex bond function on $T_k$.
Of the following relations, those that contain $\varphi$ are asserted for $k\ge1$, for every unit site $y$ and every $c\in\Delta_y$;
those that contain $U$ are asserted for all $c,c'\in\mathbb{L}_k$ with $|c-c'|\le M$.
For every $\tilde{g}\in\mathbb{G}_k$,
\[
\varphi_{\tilde{g}B}(\tilde{g}y,\tilde{g}c)=\varphi_B(y,c),
\qquad
U_{\tilde{g}B}(\tilde{g}c,\tilde{g}c')=U_B(c,c').
\]
For every complex function $\boldsymbol{\lambda}$ on $\mathbb{L}_k$ such that $d\boldsymbol{\lambda}$ is $\ell\mathbb{Z}^3$-periodic,
\begin{gather*}
\varphi_{B-d\boldsymbol{\lambda}}(y,c)=\exp\left(i\cdot e_k\cdot\left(\boldsymbol{\lambda}(y)-\boldsymbol{\lambda}(c)\right)\right)\cdot\varphi_B(y,c),\\
U_{B-d\boldsymbol{\lambda}}(c,c')=\exp\left(i\cdot e_k\cdot\left(\boldsymbol{\lambda}(c)-\boldsymbol{\lambda}(c')\right)\right)\cdot U_B(c,c').
\end{gather*}
Finally,
\[
\overline{\varphi_{\overline{B}}(y,c)}=\varphi_{-B}(y,c),
\qquad
\overline{U_{\overline{B}}(c,c')}=U_{-B}(c,c'),
\qquad
\varphi_{-B}=\varphi_B^{-1},
\qquad
U_{-B}=U_B^{-1}.
\]
\item[(d)] \emph{Gauge.}
Let $\boldsymbol{\lambda}$ be a complex function on $\mathbb{L}_k$ such that $d\boldsymbol{\lambda}$ and $\exp\left(i\cdot e_k\cdot\boldsymbol{\lambda}\right)$ are $\ell\mathbb{Z}^3$-periodic.
The symbol $\exp\left(i\cdot e_k\cdot\boldsymbol{\lambda}\right)$ also denotes the multiplication by $\exp\left(i\cdot e_k\cdot\boldsymbol{\lambda}(c)\right)$ on spinors on $T^{\mathrm{f}}_k$, and the multiplication by the restriction of this function on spinors on $\mathbb{Z}^3$ and on $L\mathbb{Z}^3$.
Then
\begin{gather*}
D_{B-d\boldsymbol{\lambda}}\cdot\exp\left(i\cdot e_k\cdot\boldsymbol{\lambda}\right)=\exp\left(i\cdot e_k\cdot\boldsymbol{\lambda}\right)\cdot D_B,\\
Q_k(B-d\boldsymbol{\lambda})\cdot J^*\cdot\exp\left(i\cdot e_k\cdot\boldsymbol{\lambda}\right)=\exp\left(i\cdot e_k\cdot\boldsymbol{\lambda}\right)\cdot Q_k(B)\cdot J^*,\\
J\cdot Q_k^T\left(-(B-d\boldsymbol{\lambda})\right)\cdot\exp\left(i\cdot e_k\cdot\boldsymbol{\lambda}\right)=\exp\left(i\cdot e_k\cdot\boldsymbol{\lambda}\right)\cdot J\cdot Q_k^T(-B),\\
Q(B-d\boldsymbol{\lambda})\cdot\exp\left(i\cdot e_k\cdot\boldsymbol{\lambda}\right)=\exp\left(i\cdot e_k\cdot\boldsymbol{\lambda}\right)\cdot Q(B),\\
Q^T\left(-(B-d\boldsymbol{\lambda})\right)\cdot\exp\left(i\cdot e_k\cdot\boldsymbol{\lambda}\right)=\exp\left(i\cdot e_k\cdot\boldsymbol{\lambda}\right)\cdot Q^T(-B),
\end{gather*}
where at $k=0$ the last two relations are taken with $Q(\,\cdot\,)\cdot J^*$ and $J\cdot Q^T(-\,\cdot\,)$, and
\[
\left(\delta_{\alpha,B-d\boldsymbol{\lambda}}\left(\exp\left(i\cdot e_k\cdot\boldsymbol{\lambda}\right)f\right)\right)(x,x')
=\exp\left(i\cdot e_k\cdot\boldsymbol{\lambda}(c)\right)\cdot\left(\delta_{\alpha,B}f\right)(x,x')
\qquad\text{for }(x,x')\in\mathrm{Pzt}_k,
\]
with $c=\pi(x)$.
It follows that, wherever the operators exist (they exist at $B-d\boldsymbol{\lambda}$ if and only if they exist at $B$),
\[
O(B-d\boldsymbol{\lambda})\cdot\exp\left(i\cdot e_k\cdot\boldsymbol{\lambda}\right)=\exp\left(i\cdot e_k\cdot\boldsymbol{\lambda}\right)\cdot O(B)
\]
for every operator $O$ without bar of Sections~I.3.5 and I.3.9;
on the second entry of $H_k^{\#}$, which is a function of pairs, the factor is $\exp\left(i\cdot e_k\cdot\boldsymbol{\lambda}(c)\right)$ with $c$ the photon site under the first point of the pair.
Furthermore,
\[
\bar{O}(B-d\boldsymbol{\lambda})\cdot\exp\left(-i\cdot e_k\cdot\boldsymbol{\lambda}\right)=\exp\left(-i\cdot e_k\cdot\boldsymbol{\lambda}\right)\cdot\bar{O}(B)
\qquad\text{for }\bar{O}=\bar{H}_k,\ \bar{T}_k,
\]
and
\[
\det D_k(B-d\boldsymbol{\lambda})=\det D_k(B),
\qquad
\delta Z_k(B-d\boldsymbol{\lambda})=\delta Z_k(B).
\]
\item[(e)] \emph{Real bond functions.}
Let $B$ be real.
Then
\[
\left(Q_k(B)\cdot J^*\right)^*=J\cdot Q_k^T(-B),
\qquad
Q(B)^*=Q^T(-B),
\qquad
\left(Q(B)\cdot J^*\right)^*=J\cdot Q^T(-B);
\]
the operators $J\cdot P_k(B)\cdot J^*$, $J\cdot P_{k+1}(B)\cdot J^*$, $P(B)$ and $J\cdot P(B)\cdot J^*$ are orthogonal projections;
and
\[
D_B^*=\gamma_5\cdot D_B\cdot\gamma_5,
\qquad
\operatorname{Re}D_B=\frac{\rho}{2\eta_k}\cdot\nabla_B^{*}\nabla_B\ge0.
\]
\item[(f)] \emph{Scaling.}
Let $B=B'_L$.
Then $e_k\cdot\eta_k\cdot B(b)=e_{k+1}\cdot\eta_{k+1}\cdot B'(b/L)$ for every bond $b$ of $T_k$.
For $f,f''\in\mathfrak{S}^{\mathrm{ch}}(\eta_{k+1},\ell/L)$ and $v,v''\in\mathfrak{S}(1,\ell/L)$ one has
\begin{gather*}
\left(D_B+\mu\right)f_L=L^{-1}\cdot\left(\left(D_{B'}+L\cdot\mu\right)f\right)_L,
\qquad
Q_{k+1}(B)\cdot J^*f_L=\left(Q_{k+1}(B')\cdot J^*f\right)_L,\\
J\cdot Q_{k+1}^T(-B)\,v_L=\left(J\cdot Q_{k+1}^T(-B')\,v\right)_L,
\qquad
J\cdot P_{k+1}(B)\cdot J^*f_L=\left(J\cdot P_{k+1}(B')\cdot J^*f\right)_L,\\
\left\langle f''_L,f_L\right\rangle^{\mathrm{ch}}=L\cdot\left\langle f'',f\right\rangle^{\mathrm{ch}},
\qquad
\left\langle v''_L,v_L\right\rangle_L=L\cdot\left\langle v'',v\right\rangle_1,
\qquad
\delta_{\alpha,B}f_L=\left(\delta_{\alpha,B'}f\right)_L\ \text{ on }\mathrm{Pzt}_k,
\end{gather*}
the pairing $\langle\,,\rangle^{\mathrm{ch}}$ on the left being formed at level $k$ and the one on the right at level $k+1$.
The maps $f\mapsto f_L$ and $v\mapsto v_L$ are bijections of $\mathfrak{S}^{\mathrm{ch}}(\eta_{k+1},\ell/L)$ onto $\mathfrak{S}^{\mathrm{ch}}(\eta_k,\ell)$ and of $\mathfrak{S}(1,\ell/L)$ onto $\mathfrak{S}(L,\ell)$.
The same relations hold for fields.
\end{itemize}
\end{named}

\begin{named}{Remark}
Proposition~21.1 corresponds to \cite[Lemma~14.19\,(a)--(f)]{QEDI};
the definitions of \cite[Sections~14.6.1 and~14.6.2]{QEDI} on which that lemma rests are replaced by those of Sections~I.1.5, I.1.6 and I.3.1--I.3.9.
The fine space is $\mathfrak{S}^{\mathrm{ch}}(\eta_k,\ell)$, with the pairing $\langle\,,\rangle^{\mathrm{ch}}$;
at $k=0$ it is also the unit space.
Every averaging operator on the fine side contains explicitly the cluster map $J^*$ or $J$, and the relations of the form $Q\cdot Q^T=1$ rest on $J^*\cdot J=1$.
The transposition relations of (c) are obtained by re-indexing through the involutions $x\mapsto x^{(i)}$ together with (I.m2);
the diagonal term requires (I.m1), that is, the closure relation (I.1.1).
At $k=0$ the second trace formula of (a) carries the factor $48\,n_0-n_1$.
The matrix $\gamma_3$ of \cite[Lemma~14.19]{QEDI} is $\gamma_5$ here, and the operators $K(B)$ and $X_k(B)$ of that lemma are written $D_B+\mu$ and $F_\infty(B)$.
In (f) the operators of level $k+1$ on the torus of side $\ell/L$ are those listed in (S21), and the pairs of level $k+1$ are formed on that torus;
the inequality $2M<\ell/L$ needed for them follows from $\ell\in\mathrm{Cl}$.

Clause (d0) has no counterpart among the clauses of \cite[Lemma~14.19]{QEDI}.
In \cite{QEDI} its relations are used inside proofs:
the gauge relations in the proofs of \cite[Lemma~14.19\,(d)]{QEDI} and \cite[Lemma~16.5\,(a)]{QEDI};
the relations $\varphi_{-B}=\varphi_B^{-1}$ and $U_{-B}=U_B^{-1}$ in the proofs of \cite[Lemma~14.19\,(a) and~(d)]{QEDI} and \cite[Lemma~16.5\,(c)]{QEDI};
the relations under complex conjugation in the proof of \cite[Lemma~16.5\,(d)]{QEDI}, and, for real $B$, in that of \cite[Lemma~14.19\,(e)]{QEDI};
the covariance under $\mathbb{G}_k$ in the proof of \cite[Lemma~16.5\,(b)]{QEDI}.

The following assertions of Proposition~21.1 are not contained in the statement of \cite[Lemma~14.19]{QEDI}.
\begin{itemize}
\item[(a)] The idempotency of $J\cdot P_{k+1}(B)\cdot J^*$, and the relation
\[
\left(J\cdot P_k\cdot J^*\right)\cdot\left(J\cdot P_{k+1}\cdot J^*\right)=\left(J\cdot P_{k+1}\cdot J^*\right)\cdot\left(J\cdot P_k\cdot J^*\right)=J\cdot P_{k+1}\cdot J^*;
\]
both are part of \cite[Lemma~16.3]{QEDI}.
\item[(b)] The formula for $J\cdot Q_{k+1}^T(-B)\,v'$, which in \cite{QEDI} is derived in the proof of \cite[Lemma~14.19\,(f)]{QEDI}.
\item[(c)] The transposes of the averaging operators, which in \cite{QEDI} are given in \cite[Section~14.6.2]{QEDI} and in \cite[Lemma~16.5\,(c)]{QEDI};
the transposition relation for $J\cdot P_{k+1}(B)\cdot J^*$, which is part of \cite[Lemma~16.5\,(c)]{QEDI};
the three equivalences of invertibility, whereas \cite[Lemma~14.19]{QEDI} assumes that every inverse written there exists;
the operator $\dot{M}_k$ in the list of the operators $O$.
The relation $\dot{M}_k(B)=\mathfrak{H}_k(B)\cdot H_k(B)$ of (c) is the definition of $\dot{M}_k$ in \cite[Section~14.6.2]{QEDI}.
\item[(d)] The existence of the operators at $B-d\boldsymbol{\lambda}$ if and only if they exist at $B$;
the relation for the second entry of $H_k^{\#}$;
the relation for every operator without bar of Sections~I.3.5 and I.3.9, where \cite[Lemma~14.19\,(d)]{QEDI} lists twelve operators;
the relation $\det D_k(B-d\boldsymbol{\lambda})=\det D_k(B)$.
\item[(e)] The assertion that $J\cdot P_{k+1}(B)\cdot J^*$ is an orthogonal projection, which is part of \cite[Lemma~16.5\,(e)]{QEDI};
the formula $\operatorname{Re}D_B=\frac{\rho}{2\eta_k}\cdot\nabla_B^{*}\nabla_B$, where \cite[Lemma~14.19\,(e)]{QEDI} asserts $\operatorname{Re}D_B\ge0$.
\item[(f)] The relation for $J\cdot P_{k+1}(B)\cdot J^*f_L$, which in \cite{QEDI} is derived in Step~1 of the proof of \cite[Theorem~14.21]{QEDI};
the bijectivity of $f\mapsto f_L$ and of $v\mapsto v_L$, which is stated in the paragraph on scaling of \cite[Section~14.6.2]{QEDI}.
\end{itemize}
\end{named}

\subsubsection*{Proposition~21.2. Two quadratic identities}

\begin{named}{Proposition 21.2}
Assume the setting (S21).
\begin{itemize}
\item[(a)] Let $k\ge1$, and suppose that $S_k(B)$ exists.
Let $\Psi$, $\bar{\Psi}$ be fields on $T^0_k$, let $\bar{\psi}$ be a field on $T^{\mathrm{f}}_k$, and put $\psi:=H_k(B)\Psi$.
Then
\[
b_k\cdot\left\langle\bar{\Psi}-Q_k(-B)\cdot J^*\bar{\psi},\,\Psi-Q_k(B)\cdot J^*\psi\right\rangle_1+\left\langle\bar{\psi},\left(D_B+\mu\right)\psi\right\rangle^{\mathrm{ch}}=\left\langle\bar{\Psi},D_k(B)\Psi\right\rangle_1.
\]
\item[(b)] Let $k\ge0$, and suppose that the inverse
\[
S^0_{k+1}(B)=\left(D_B+\mu+\frac{b_{k+1}}{L}\cdot J\cdot P_{k+1}(B)\cdot J^*\right)^{-1}
\]
exists.
Let $\Phi$, $\bar{\Phi}$ be fields on $L\mathbb{Z}^3/\ell\mathbb{Z}^3$, let $\bar{\psi}$ be a field on $T^{\mathrm{f}}_k$, and put
\[
\psi:=\frac{b_{k+1}}{L}\cdot S^0_{k+1}(B)\cdot J\cdot Q_{k+1}^T(-B)\Phi.
\]
Then
\begin{multline*}
\frac{b_{k+1}}{L}\cdot\left\langle\bar{\Phi}-Q_{k+1}(-B)\cdot J^*\bar{\psi},\,\Phi-Q_{k+1}(B)\cdot J^*\psi\right\rangle_L+\left\langle\bar{\psi},\left(D_B+\mu\right)\psi\right\rangle^{\mathrm{ch}}\\
=\left\langle\bar{\Phi},\left(\frac{b_{k+1}}{L}-\left(\frac{b_{k+1}}{L}\right)^2\cdot Q_{k+1}(B)\cdot J^*\cdot S^0_{k+1}(B)\cdot J\cdot Q_{k+1}^T(-B)\right)\Phi\right\rangle_L.
\end{multline*}
\end{itemize}
\end{named}

\begin{named}{Remark}
No statement of \cite{QEDI} corresponds to Proposition~21.2;
Propositions~21.1, 21.3, 21.4 and~21.5 correspond to \cite[Lemma~14.19]{QEDI}, \cite[Proposition~14.20]{QEDI}, \cite[Theorem~14.21]{QEDI} and \cite[Lemma~14.22]{QEDI}, respectively.
Proposition~21.2 gives the values of the quadratic forms on the left sides of (a) and (b) at $\psi=H_k(B)\Psi$, respectively at $\psi=\frac{b_{k+1}}{L}\cdot S^0_{k+1}(B)\cdot J\cdot Q_{k+1}^T(-B)\Phi$, for an arbitrary field $\bar{\psi}$;
the right sides do not depend on $\bar{\psi}$.
The operator between $\bar{\Phi}$ and $\Phi$ on the right side of (b) is the right side of (21.7a);
under the supposition of Proposition~21.4 it is therefore $D^0_{k+1}(B)$.
\end{named}

\subsubsection*{Proposition~21.3. Resolvent identities between the fine lattice and the unit lattice, $k\ge1$}

\begin{named}{Proposition 21.3}
Assume the setting (S21) with $k\ge1$, let $y$ be a real number, and suppose that $F_\infty(B)$ is invertible.
All operators are taken at $B$.
Then $b_k+b_L\cdot P+i\gamma_5\cdot y$ is invertible on $\mathfrak{S}(1,\ell)$, and
\[
B_{k,y}=\left(b_k+i\gamma_5\cdot y\right)^{-1}\cdot\left(1-P\right)+\left(b_k+b_L+i\gamma_5\cdot y\right)^{-1}\cdot P,
\qquad
b_k-b_k^2\cdot B_{k,y}=\boldsymbol{\alpha}_y+\boldsymbol{\beta}_y\cdot P,
\]
\begin{equation}\tag{21.1}
F_y=F_\infty-b_k^2\cdot J\cdot Q_k^T(-B)\cdot B_{k,y}\cdot Q_k\cdot J^*.
\end{equation}
The following two conditions are equivalent:
\begin{itemize}
\item[(i)] $D_k+b_L\cdot P+i\gamma_5\cdot y$ is invertible on $\mathfrak{S}(1,\ell)$;
\item[(ii)] $F_y$ is invertible on $\mathfrak{S}^{\mathrm{ch}}(\eta_k,\ell)$.
\end{itemize}
If they hold, let $\Gamma_{k,y}=\left(D_k+b_L\cdot P+i\gamma_5\cdot y\right)^{-1}$ and $S_{k,y}=F_y^{-1}$;
then
\begin{align}
\Gamma_{k,y}&=B_{k,y}+b_k^2\cdot B_{k,y}\cdot Q_k\cdot J^*\cdot S_{k,y}\cdot J\cdot Q_k^T(-B)\cdot B_{k,y},\tag{21.2}\\
S_{k,y}&=S_k+H_k\cdot\Gamma_{k,y}\cdot\mathfrak{H}_k,\tag{21.3}\\
H_k\cdot\Gamma_{k,y}&=b_k\cdot S_{k,y}\cdot J\cdot Q_k^T(-B)\cdot B_{k,y},
\qquad
\Gamma_{k,y}\cdot\mathfrak{H}_k=b_k\cdot B_{k,y}\cdot Q_k\cdot J^*\cdot S_{k,y},\tag{21.4}\\
\det\left(D_k+b_L\cdot P+i\gamma_5\cdot y\right)\cdot\det F_\infty&=\det\left(b_k+b_L\cdot P+i\gamma_5\cdot y\right)\cdot\det F_y,\tag{21.5}
\end{align}
where the four determinants in (21.5) are taken, from left to right, on $\mathfrak{S}(1,\ell)$, $\mathfrak{S}^{\mathrm{ch}}(\eta_k,\ell)$, $\mathfrak{S}(1,\ell)$ and $\mathfrak{S}^{\mathrm{ch}}(\eta_k,\ell)$.
At $y=0$,
\[
\Gamma_{k,0}=\Gamma_k,
\qquad
F_0=D_B+\mu+\frac{b_{k+1}}{L}\cdot J\cdot P_{k+1}\cdot J^*,
\qquad
S_{k,0}=S^0_{k+1}.
\]
\end{named}

\begin{named}{Remark}
Proposition~21.3 corresponds to \cite[Proposition~14.20]{QEDI}, together with the definitions that precede it in \cite[Section~14.6.3]{QEDI}.
Here the averaging operators on the fine side contain $J^*$ and $J$, and the block operator matrix on which the identities rest acts on $\mathfrak{S}(1,\ell)\oplus\mathfrak{S}^{\mathrm{ch}}(\eta_k,\ell)$.
The parameter $y$ ranges over all real numbers;
in \cite[Proposition~14.20]{QEDI} it is restricted to $y\ge0$, and the proof there makes no use of the sign of $y$.
The objects $a_y$, $\boldsymbol{\alpha}_{k,y}$, $\boldsymbol{\beta}_{k,y}$, $X_{k,y}$ and $X^0_{k+1}$ of \cite[Section~14.6.3]{QEDI} are the objects $B_{k,y}^{-1}$, $\boldsymbol{\alpha}_y$, $\boldsymbol{\beta}_y$, $F_y$ and $F_0$ of Section~I.3.5.
\end{named}

\subsubsection*{Proposition~21.5. Real backgrounds}

\begin{named}{Proposition 21.5}
In the setting (S21) let $B$ be real and $\mu>0$.
All operators are taken at $B$, and $\operatorname{Re}O:=\frac{1}{2}\left(O+O^*\right)$, the adjoint referring to the sesquilinear pairings of Section~I.1.6.
\begin{itemize}
\item[(a)] One has $\operatorname{Re}\left(D_B+\mu\right)\ge\mu$, and for $k\ge1$ also $\operatorname{Re}F_\infty\ge\mu$.
Hence $D_B+\mu$ and $F_\infty$ are invertible.
\item[(b)] The operator $D_k$ is invertible and $\operatorname{Re}D_k>0$, so that $\det D_k\ne0$.
For $k\ge1$,
\[
D_k^{-1}=b_k^{-1}+Q_k\cdot J^*\cdot\left(D_B+\mu\right)^{-1}\cdot J\cdot Q_k^T(-B).
\]
\item[(c)] For $k\ge1$,
\[
\operatorname{Re}\left(D_k+b_L\cdot P\right)>0,
\qquad
\left(D_k+b_L\cdot P\right)^*=\gamma_5\cdot\left(D_k+b_L\cdot P\right)\cdot\gamma_5;
\]
for $k=0$,
\[
\operatorname{Re}F_0\ge\mu,
\qquad
F_0^*=\gamma_5\cdot F_0\cdot\gamma_5.
\]
Consequently $\Gamma_k$ exists and $\delta Z_k\ne0$.
\item[(d)] For every real $y$ the operator $D_k+b_L\cdot P+i\gamma_5\cdot y$ if $k\ge1$, respectively $F_y=F_0+i\gamma_5\cdot y$ if $k=0$, is invertible;
for $y\ne0$ the operator norm of its inverse with respect to the $L^2$ norms is at most $|y|^{-1}$.
For $k\ge1$ the operator $S_{k,y}$ exists.
\end{itemize}
\end{named}

\begin{named}{Remark}
Proposition~21.5 corresponds to \cite[Lemma~14.22\,(a)--(d)]{QEDI}.
Here the averaging operators contain $J^*$ and $J$, the matrix $\gamma_3$ of \cite[Lemma~14.22]{QEDI} is $\gamma_5$, and at $k=0$ the operators act on $\mathfrak{S}^{\mathrm{ch}}(1,\ell)$.
The condition $\rho\ge0$ enters (a) through (I.O2).
Clause (c) of \cite[Lemma~14.22]{QEDI} contains in addition the identity $\det\left(\left(D_k+b_L\cdot P\right)\cdot\gamma_5\right)=\delta Z_k(B)$, written there with $\gamma_3$;
this identity rests on $\det\left(\gamma_5\otimes1\right)=1$, and for the data (SF) of Section~I.1.8 and real $B\in\tilde{G}_k$ it follows from Proposition~12.
\end{named}

\subsubsection*{Proposition~21.4. Completion of the square, operators under scaling, determinants}

\begin{named}{Proposition 21.4}
Assume the setting (S21), and let $B=B'_L$.
Clauses (a), (b) and (c) are asserted for every $B'$ such that $\Gamma_k(B)$ exists and, if $k\ge1$, $S_k(B)$ exists.
\begin{itemize}
\item[(b)] \emph{Operators under scaling.}
The operator $S^0_{k+1}(B)$ exists; for $k=0$ it is $\Gamma_0(B)$.
The operators $S_{k+1}(B')$ and $S_{k+1}(-B')$ exist, and hence so do $H_{k+1}(B')$, $\bar{H}_{k+1}(B')$ and $D_{k+1}(B')$.
It is not required that $B'$ lie in $\tilde{G}_{k+1}$.
For $f\in\mathfrak{S}^{\mathrm{ch}}(\eta_{k+1},\ell/L)$ and $v\in\mathfrak{S}(1,\ell/L)$,
\begin{align}
S^0_{k+1}(B)f_L&=L\cdot\left(S_{k+1}(B')f\right)_L,\tag{21.6}\\
D^0_{k+1}(B)&=\frac{b_{k+1}}{L}-\left(\frac{b_{k+1}}{L}\right)^2\cdot Q_{k+1}(B)\cdot J^*\cdot S^0_{k+1}(B)\cdot J\cdot Q_{k+1}^T(-B),\tag{21.7a}\\
D^0_{k+1}(B)v_L&=L^{-1}\cdot\left(D_{k+1}(B')v\right)_L,\tag{21.7b}
\end{align}
\begin{equation}\tag{21.8a}
\begin{aligned}
b_L\cdot H_k(B)\cdot\Gamma_k(B)\cdot Q^T(-B)&=\frac{b_{k+1}}{L}\cdot S^0_{k+1}(B)\cdot J\cdot Q_{k+1}^T(-B)&&\text{if }k\ge1,\\
b_L\cdot\Gamma_0(B)\cdot J\cdot Q^T(-B)&=\frac{b_1}{L}\cdot S^0_1(B)\cdot J\cdot Q_1^T(-B)&&\text{if }k=0,
\end{aligned}
\end{equation}
where for $k=0$ the two sides are one and the same expression, and
\begin{equation}\tag{21.8b}
\begin{aligned}
b_L\cdot H_k(B)\cdot\Gamma_k(B)\cdot Q^T(-B)v_L&=\left(H_{k+1}(B')v\right)_L&&\text{if }k\ge1,\\
b_L\cdot\Gamma_0(B)\cdot J\cdot Q^T(-B)v_L&=\left(H_1(B')v\right)_L&&\text{if }k=0.
\end{aligned}
\end{equation}
\item[(a)] \emph{Completion of the square.}
Let $\Psi'$, $\bar{\Psi}'$ be fields on $T^0_{k+1}$, with $4$ components;
let $W'$, $\bar{W}'$ be fields on $T^0_k$, with $4$ components, if $k\ge1$, and fields on $T^{\mathrm{f}}_0$, with $192$ components per unit site, if $k=0$.
Put
\begin{equation}\tag{21.9}
\begin{aligned}
\Psi^c&:=b_L\cdot\Gamma_k(B)\cdot Q^T(-B)\Psi'_L,&\bar{\Psi}^c&:=b_L\cdot\Gamma_k(B)^T\cdot Q^T(B)\bar{\Psi}'_L&&\text{if }k\ge1,\\
\Psi^c&:=b_L\cdot\Gamma_0(B)\cdot J\cdot Q^T(-B)\Psi'_L,&\bar{\Psi}^c&:=b_L\cdot\Gamma_0(B)^T\cdot J\cdot Q^T(B)\bar{\Psi}'_L&&\text{if }k=0;
\end{aligned}
\end{equation}
these are fields on $T^0_k$ if $k\ge1$, and on $T^{\mathrm{f}}_0$ if $k=0$.
Then, for $\Psi=\Psi^c+W'$ and $\bar{\Psi}=\bar{\Psi}^c+\bar{W}'$,
\begin{equation}\tag{21.10}
\begin{aligned}
&\left\langle\bar{\Psi},D_k(B)\Psi\right\rangle_1+b_L\cdot\left\langle\bar{\Psi}'_L-Q(-B)\bar{\Psi},\,\Psi'_L-Q(B)\Psi\right\rangle_L\\
&\qquad=\left\langle\bar{\Psi}',D_{k+1}(B')\Psi'\right\rangle_1+\left\langle\bar{W}',\left(D_k(B)+b_L\cdot P(B)\right)W'\right\rangle_1
\qquad\text{if }k\ge1,\\
&\left\langle\bar{\Psi},\left(D_B+\mu\right)\Psi\right\rangle^{\mathrm{ch}}+b_L\cdot\left\langle\bar{\Psi}'_L-Q(-B)\cdot J^*\bar{\Psi},\,\Psi'_L-Q(B)\cdot J^*\Psi\right\rangle_L\\
&\qquad=\left\langle\bar{\Psi}',D_1(B')\Psi'\right\rangle_1+\left\langle\bar{W}',\left(D_B+\mu+b_L\cdot J\cdot P(B)\cdot J^*\right)W'\right\rangle^{\mathrm{ch}}
\qquad\text{if }k=0.
\end{aligned}
\end{equation}
Moreover, for every $k\ge0$,
\begin{equation}\tag{21.11}
\begin{aligned}
&\psi^0:=H_k(B)\Psi^c=\left(H_{k+1}(B')\Psi'\right)_L,
\qquad
\Psi^c=T_k(B)\psi^0,\\
&\bar{\psi}^0:=\bar{H}_k(B)\bar{\Psi}^c=\left(\bar{H}_{k+1}(B')\bar{\Psi}'\right)_L,
\qquad
\bar{\Psi}^c=\bar{T}_k(B)\bar{\psi}^0,\\
&\delta_{\alpha,B}\psi^0=\left(\delta_{\alpha,B'}H_{k+1}(B')\Psi'\right)_L,
\qquad
\delta_{\alpha,-B}\bar{\psi}^0=\left(\delta_{\alpha,-B'}\bar{H}_{k+1}(B')\bar{\Psi}'\right)_L
\quad\text{on }\mathrm{Pz}_k.
\end{aligned}
\end{equation}
At $k=0$ one has $H_0=\bar{H}_0=T_0=\bar{T}_0=1$, so that $\psi^0=\Psi^c$ and $\bar{\psi}^0=\bar{\Psi}^c$;
these fields are in general not constant on clusters:
at $B=0$, if $F_0(0)$ is invertible, $\Psi^c$ is constant on clusters only if $\Psi'=0$, and $\bar{\Psi}^c$ only if $\bar{\Psi}'=0$.
For $k\ge1$,
\[
\Psi^c=Q_k(B)\cdot J^*\psi^0+\frac{b_L}{b_k+b_L}\cdot Q^T(-B)\left(\Psi'_L-Q_{k+1}(B)\cdot J^*\psi^0\right).
\]
\item[(c)] \emph{Determinants.}
For every $k\ge0$,
\[
\det D_{k+1}(B')=c'_k\cdot\frac{\det D_k(B)}{\delta Z_k(B)},
\qquad
c'_k=b^{4n_{k+1}}.
\]
Here $D_{k+1}(B')$ acts on $\mathfrak{S}(1,\ell/L)$, of dimension $4\,n_{k+1}$;
for $k\ge1$ the operators $D_k(B)$ and $D_k(B)+b_L\cdot P(B)$ act on $\mathfrak{S}(1,\ell)$, of dimension $4\,n_k$;
for $k=0$ the operators $D_0(B)=D_B+\mu$ and $F_0(B)$ act on $\mathfrak{S}^{\mathrm{ch}}(1,\ell)$, of dimension $192\,n_0$.
\end{itemize}
Clause (d) is asserted without the supposition under which (a)--(c) are stated.
\begin{itemize}
\item[(d)] \emph{Where the supposition holds; analyticity.}
Let $\mathrm{Iv}_k(\ell,\mu,e_k)$ denote the set of the complex bond functions $B$ on $T_k$ at which $\Gamma_k(B)$ exists and, if $k\ge1$, $S_k(B)$ exists.
For $B$ in this set, $B'$ is the bond function on $T_{k+1}$ with $B'_L=B$, that is, $B'(b')=L^{\frac{1}{2}}\cdot B(L\,b')$.
\begin{itemize}
\item[(d1)] The set $\mathrm{Iv}_k(\ell,\mu,e_k)$ is open in the space of the complex bond functions and invariant under $B\mapsto-B$;
all operators and determinants that occur in (a)--(c) are analytic functions of $B$ on it, and (a)--(c) are identities between analytic functions on it.
\item[(d2)] If $\mu\in(0,1]$ and $e_k\in(0,e_0]$, the constants being those of Section~I.8, then $\tilde{G}_k\subset\mathrm{Iv}_k(\ell,\mu,e_k)$.
\item[(d3)] If $\mu>0$, then every real $B$ lies in $\mathrm{Iv}_k(\ell,\mu,e_k)$.
\end{itemize}
The inverses that occur are obtained in three ways.
Those of level $k$ at $\pm B$ come from (d2) or (d3), together with the invariance of $\tilde{G}_k$ under $B\mapsto-B$;
at $k=0$ one has $\Gamma_0(B)=S^0_1(B)=F_0(B)^{-1}$.
Those of level $k+1$ at $\pm B'$ come from (b), that is, from the relation
\[
F_0(\pm B)f_L=L^{-1}\cdot\left(F_\infty(\pm B')f\right)_L
\]
and from the bijectivity of $f\mapsto f_L$.
\end{itemize}
\end{named}

\begin{named}{Remark}
Proposition~21.4 corresponds to \cite[Theorem~14.21]{QEDI} and to \cite[Proposition~5.55]{QEDI}:
clause (b) of Proposition~21.4 corresponds to clause (a) of these two statements, its clause (a) to their clause (b), and its clauses (c) and (d) to their clauses (c) and (d).
At $k=0$ the fields $W'$, $\bar{W}'$, $\Psi^c$ and $\bar{\Psi}^c$ have $192$ components per unit site, the pairing is $\langle\,,\rangle^{\mathrm{ch}}$, and the one-step averages are $Q(\,\cdot\,)\cdot J^*$, $J\cdot Q^T(\,\cdot\,)$ and $J\cdot P(\,\cdot\,)\cdot J^*$, whereas $\Psi'$ keeps $4$ components;
the constant $c'_0=b^{4n_1}$ is the same as in \cite[Theorem~14.21]{QEDI}, since its exponent is the dimension of $\mathfrak{S}(L,\ell)$.
For $k\ge1$ the cluster maps $J^*$ and $J$ occur on the fine side only.
Clause (b) contains the existence of $S_{k+1}(\pm B')$, which follows, as stated in (d), from the relation $F_0(\pm B)f_L=L^{-1}\cdot\left(F_\infty(\pm B')f\right)_L$ and the bijectivity of $f\mapsto f_L$.
In \cite[Proposition~5.55\,(d)]{QEDI} the inverses of level $k$ at $\pm B$ are obtained from the last paragraph of \cite[Section~5.2.6]{QEDI}, applied at $B$ and at $-B$, which lies in $\tilde{G}_k$ together with $B$ by \cite[Lemma~5.19]{QEDI};
for $k\ge1$ they are contained also in \cite[Propositions~5.42--5.44]{QEDI}, applied at $B$ and at $-B$, and for $k=0$ they are $\Gamma_0(\pm B)$, denoted there also $S^0_1(\pm B)$ and $S_{0,0}(\pm B)$.
The inverses of level $k+1$ at $\pm B'$ are derived there, as in clause (d) of Proposition~21.4, from the relation between the operators of the two levels, which is established in Step~1 of the proof of \cite[Theorem~14.21]{QEDI}, and from the bijectivity of $f\mapsto f_L$.
In this the two statements agree, the operators of that relation being here $F_0(\pm B)$ and $F_\infty(\pm B')$, which contain the cluster maps $J^*$ and $J$.
Neither statement applies the bounds of level $k+1$ at $B'$, since $B'_L\in\tilde{G}_k$ does not imply $B'\in\tilde{G}_{k+1}$.
At $k=0$ the invertibility of $F_0(B)$ for $B\in\tilde{G}_0$ is the case $y=0$ of Proposition~2.8.
As in \cite[Theorem~14.21]{QEDI}, the supposition that $\Gamma_k(B)$ exists and, if $k\ge1$, that $S_k(B)$ exists is made for the clauses (a)--(c) only;
clause (d), which quantifies over $B$, with $B'$ defined by $B'_L=B$, is stated without it.
At $k=0$ the operators $\bar{H}_0$, $\bar{T}_0$, $S^0_1$, $P_1$ and $\dot{M}_0$ are those defined in Section~I.3.9 and in (I.L0.7);
there $P_1(B)=P(B)$ acts on $\mathfrak{S}(1,\ell)$, and the operator on $\mathfrak{S}^{\mathrm{ch}}(1,\ell)$ that takes the place of $J\cdot P_{k+1}(B)\cdot J^*$ is $J\cdot P_1(B)\cdot J^*=J\cdot P(B)\cdot J^*$.
The operator $D^0_{k+1}(B)$ is normalised as in (21.7a);
the operator denoted $D^0_{k+1}$ in \cite[Lemma~13.19]{QEDI} is $L$ times this operator.
Clause (d2) is not subject to the hypotheses of \cite[Theorem~16.16]{QEDI}, under which \cite[Theorem~14.21\,(d)]{QEDI} is stated;
it holds under the conditions on the constants under which (G3.T-2a) holds, and these conditions are part of those of the Propositions~14.m.
\end{named}

\subsection*{22. The one-step formula for the density with external connection, all levels}

\subsubsection*{G3.2.0. The setting (S22) of the Propositions 22.2--22.4}

\paragraph{Constants.}
The constants are those of Section~I.8, lines (C0)--(C8).

\paragraph{Data.}
The data are an integer $k\ge0$, a side $\ell$ with $\ell\in\mathrm{Cl}$ and $\ell/L\in\mathrm{Cl}$, a mass $\mu\in\left(0,\frac{1}{L}\right]$, and a number $e_k>0$ such that $e_{k+1}:=L^{\frac{1}{2}}\cdot e_k\le e_0$.
For $k=0$ the number $e_k$ is the charge of level $0$, which in the flow is written $q$ (Section~I.2.5).
At level $k+1$ the side is $\ell/L$, the mass is $L\cdot\mu$ and the charge is $e_{k+1}$; moreover $n_k=\ell^3$ and $n_{k+1}=(\ell/L)^3$.
The numbers $v_k=-\log e_k$ and $p_{0,k}=v_k^{p_0}$ are formed with $e_k$.
No further relation between $k$, $\ell$, $\mu$ and $e_k$ is assumed, and no system of primitives enters.

\paragraph{Objects.}
The photon objects $H_k$, $H_{k+1}$, $A^{\min}_{k}$, $\omega_{k}$, $C$, $C_{k}^{1/2}$, $\mathrm{Sq}_{k}$, $V_{k}$, $\mathrm{Box}_{k}$, the measures $\mu_I$ (in the variable $\zeta$) and $\hat{\mu}$, the number $\delta\varepsilon^{0}_{k}$ and the domains are those of Sections~I.5.3 and~I.5.7.
The spaces $K^{x,C}_{k}$ and $K^{x}_{k}$ are those of Section~I.5.4, and $\varepsilon(\mathscr{E})$, $m(\mathscr{E})$, $R$ are those of Section~I.5.6.
Finally, $I^{x}_{k}$, $\mathrm{Act}^{x}_{k}$, $\mathscr{F}^{x}_{k}$, $y^{x}_{k}$ and $\mathbf{T}^{x}_{k}$ are the objects of Section~I.5.7.

\paragraph{The family $\mathfrak{P}_k$.}
Here $\mathfrak{P}_k=\left\{\mathfrak{P}_k(X,\cdot\,):X\in D_k\right\}$ is an arbitrary family of complex functions on $\mathfrak{X}_k$, and $\mathscr{F}^{x}_{k}$ is formed with this family.
The families of the flow (Section~I.5.7) are special cases.

\paragraph{Why the definitions of Section~I.5.7 make sense.}
Let $(\mathscr{A},\mathscr{C})\in\tilde{R}_k\times\mathfrak{X}_k$ and $\mathscr{T}=\mathscr{A}+\mathscr{C}$.
Then $\mathscr{T}$ and $-\mathscr{T}$ belong to $\tilde{G}_k$, by \cite[Lemma~5.17]{QEDI} and \cite[Lemma~5.19]{QEDI}.
Consequently $S_k(\pm\mathscr{T})$ exist for $k\ge1$, and with them $H_{k}(\mathscr{T})$, $\bar{H}_{k}(\mathscr{T})$ and $D_k(\mathscr{T})$, by Proposition~21.4\,(d2).
Moreover $\det D_k(0)\ne0$ by Proposition~21.5\,(b).

\paragraph{Pairs of fields in the last slot.}
In Section~G3.2.0 and in the Propositions 22.2--22.4 a field is understood in the sense of the paragraph on fields of Section~G3.1.0.
Whenever a rule of Proposition~22.1 requires the values of a field to lie in a subalgebra (the narrower sense of the word in that proposition), this requirement is stated in the proposition concerned or verified in its proof.
Let $(\Phi,\bar{\Phi})$ be a pair of fields on the unit index set of level $k$, that is, on $T^0_k$ if $k\ge1$ and on $T^{\mathrm{f}}_0$ if $k=0$.
Then $I^{x}_{k}[m,\mathscr{E}](\mathscr{A},\mathscr{C};\Phi)$ and $\mathrm{Act}^{x}_{k}[m,\mathscr{E}](\mathscr{A},\mathscr{C};\Phi)$ denote the images of the elements $I^{x}_{k}[m,\mathscr{E}](\mathscr{A},\mathscr{C};\Psi)$ and $\mathrm{Act}^{x}_{k}[m,\mathscr{E}](\mathscr{A},\mathscr{C};\Psi)$ under the homomorphism which replaces the pair of generators $(\Psi,\bar{\Psi})$ by the pair $(\Phi,\bar{\Phi})$.
In other words, one inserts
\[
\psi=H_{k}(\mathscr{T})\Phi,\qquad\bar{\psi}=\bar{H}_{k}(\mathscr{T})\bar{\Phi},\qquad\chi=\delta_{\alpha,\mathscr{T}}\psi,\qquad\bar{\chi}=\delta_{\alpha,-\mathscr{T}}\bar{\psi}.
\]

\paragraph{The class of maps $F$ in $\mathbf{T}^{x}_{k}$.}
Let $A'\in\mathrm{Dom}_{k+1}$ and $\mathscr{C}'\in\mathfrak{X}_{k+1}$.
The element $\left(\mathbf{T}^{x}_{k}F\right)(A',\mathscr{C}';\Psi')$ is defined for every map $F(\,\cdot\,,\mathscr{C}'_L;\Psi)$ from $\mathrm{Dom}_k$ into the even part of $\mathrm{Gr}[\Psi]$ such that the map $\zeta\mapsto F\left(A^{\min}_{k}(A')+\mathrm{Sq}_{k}\zeta,\mathscr{C}'_L;\Psi\right)$ is integrable over $\mathrm{Box}_{k}$ with respect to $\mu_I$, coefficient by coefficient.
For $\zeta\in\mathrm{Box}_{k}$ the point $A^{\min}_{k}(A')+\mathrm{Sq}_{k}\zeta$ lies in $\mathrm{Dom}_k$ by \cite[Proposition~5.38(ii)]{QEDI}.

\paragraph{The function $\omega_k$.}
The letter $\omega_{k}=\omega_{k}(A')$ denotes the real function on the sites of $T_k$ of \cite[Proposition~5.52]{QEDI}, see \cite[Theorem~14.15(b)]{QEDI}.
In the display that defines $\mathbf{T}^{x}_{k}$ it is evaluated at the sites of $L\mathbb{Z}^3$.

\subsubsection*{Proposition~22.1 (Berezin rules on an index set with weights)}

\begin{named}{Proposition 22.1}
Let $\{(t,\beta)\}$ be a finite index set with weights $\nu(t)>0$, and let $\langle\cdot,\cdot\rangle_{\nu}$, $\int\mathrm{D}W$ and $\mathrm{d}\mu_I(W)$ be as in Section~I.3.6 and (I.3.5).
Let $\mathrm{Gr}$ be a Grassmann algebra of finite dimension over $\mathbb{C}$ whose generators include a family $(W,\bar{W})$ indexed by $\{(t,\beta)\}$.
A second family on this index set occurs only among the data of (G4)$^{\mathrm{w}}$, and (G5)$^{\mathrm{w}}$ has data of its own.
Within this proposition a \emph{field} is a map from $\{(t,\beta)\}$ into the odd part of the subalgebra generated by the generators that do not belong to the family over which one integrates.
At the instances listed below such a field is the restriction to a fundamental domain of a field in the sense of the paragraph on fields of Section~G3.1.0 whose values lie in this subalgebra.
For $G\in\mathrm{Gr}$, a family $(W,\bar{W})$ and maps $\bar{h}$, $\mathrm{hf}$ from $\{(t,\beta)\}$ into the odd part of $\mathrm{Gr}$ (fields, and also $\bar{W}+\bar{g}$, $W+\mathrm{gf}$, $u^{-1}\bar{W}$, $uW$, $O^{-1}W$), let
\[
G\left[\bar{W}\to\bar{h},\ W\to\mathrm{hf}\right]
\]
be the image of $G$ under the homomorphism of $\mathrm{Gr}$ which sends $\bar{W}_{\beta}(t)$ to $\bar{h}_{\beta}(t)$ and $W_{\beta}(t)$ to $\mathrm{hf}_{\beta}(t)$ and leaves every other generator fixed.
Here $O$ is a linear map of the functions on $\{(t,\beta)\}$, $O^{T}$ is its transpose with respect to $\langle\cdot,\cdot\rangle_{\nu}$, and $O$ acts on fields by linearity.
Then the following rules hold.
\begin{itemize}
\item[(G0)$^{\mathrm{w}}$] For every $O$,
\[
\int\mathrm{D}W\,\exp\left(-\left\langle\bar{W},OW\right\rangle_{\nu}\right)=\det O,
\]
the integral being determined by the single requirement $\int\mathrm{D}W\,\exp\left(-\left\langle\bar{W},W\right\rangle_{\nu}\right)=1$.
These are the identities (I.B1) and (I.B2).
\item[(G1)$^{\mathrm{w}}$] For all fields $\mathrm{gf}$ and $\bar{g}$,
\[
\int\mathrm{D}W\,G\left[\bar{W}\to\bar{W}+\bar{g},\ W\to W+\mathrm{gf}\right]=\int\mathrm{D}W\,G.
\]
\item[(G2)$^{\mathrm{w}}$] For invertible $O$ and fields $\mathrm{Sr}$, $\overline{\mathrm{Sr}}$,
\[
\int\mathrm{D}W\,\exp\left(-\left\langle\bar{W},OW\right\rangle_{\nu}+\left\langle\bar{W},\mathrm{Sr}\right\rangle_{\nu}+\left\langle\overline{\mathrm{Sr}},W\right\rangle_{\nu}\right)=\det O\cdot\exp\left\langle\overline{\mathrm{Sr}},O^{-1}\mathrm{Sr}\right\rangle_{\nu}.
\]
\item[(G3)$^{\mathrm{w}}$] For a function $u$ on $\{t\}$ without zeros,
\[
\int\mathrm{D}W\,G=\int\mathrm{D}W\,G\left[\bar{W}\to u^{-1}\bar{W},\ W\to uW\right].
\]
\item[(G4)$^{\mathrm{w}}$] Let $(W',\bar{W}')$ be a second family of generators of $\mathrm{Gr}$, different from $(W,\bar{W})$, indexed by the same set $\{(t,\beta)\}$ and with the same weights.
For invertible $O$ and for every $G$ in the subalgebra generated by all generators other than those of $W$, $\bar{W}$,
\[
\int\mathrm{D}W'\,\exp\left(-\left\langle\bar{W}',OW'\right\rangle_{\nu}\right)\cdot G=\det O\cdot\int\mathrm{d}\mu_I(W)\,G\left[\bar{W}'\to\bar{W},\ W'\to O^{-1}W\right].
\]
\item[(G5)$^{\mathrm{w}}$] Let $(W^1,\bar{W}^1)$ and $(W^2,\bar{W}^2)$ be two different families of generators of $\mathrm{Gr}$, each indexed by a finite set of its own and with weights of its own (the index sets and the weights may differ).
Then
\[
\int\mathrm{D}W^1\int\mathrm{D}W^2\,F=\int\mathrm{D}W^2\int\mathrm{D}W^1\,F\qquad\left(F\in\mathrm{Gr}\right).
\]
Moreover, a Berezin integral commutes with the integral, with respect to a finite positive measure, of a function with values in $\mathrm{Gr}$ whose coefficients are integrable.
\end{itemize}
\emph{Instances.}
Here $k\ge0$ is an arbitrary integer and the side $\ell$ an arbitrary integer $\ge1$, which for the family $(\Psi',\bar{\Psi}')$ is a multiple of $L$.
The sets $T^0_k=\mathbb{Z}^3/\ell\mathbb{Z}^3$, $T^0_{k+1}=\mathbb{Z}^3/(\ell/L)\mathbb{Z}^3$, $T^{\mathrm{f}}_0=T^0_0\times O_{\mathrm{cub}}$ and the numbers $n_k=\ell^3$, $n_{k+1}=(\ell/L)^3$ are given by the displays of Sections~I.1.1 and~I.1.3, read at these sides.
\begin{itemize}
\item[] For $(\Psi',\bar{\Psi}')$: $t=y'\in T^0_{k+1}$, $\nu=1$, and there are $4n_{k+1}$ pairs of generators, for every $k\ge0$.
\item[] For $(\Psi,\bar{\Psi})$, $(W,\bar{W})$ and $(W',\bar{W}')$ with $k\ge1$: $t=y\in T^0_k$, $\nu=1$, and there are $4n_k$ pairs.
\item[] For the same three families with $k=0$: $t=x\in T^{\mathrm{f}}_0$, $\nu=\frac{1}{48}$, there are $192\,n_0$ pairs, and $\langle\cdot,\cdot\rangle_{\nu}=\langle\cdot,\cdot\rangle^{\mathrm{ch}}$.
\end{itemize}
The generators are antiperiodic, (I.AP); the index $t$ ranges over a fundamental domain, and nothing depends on the choice of this domain.
At these instances the function $u$ in (G3)$^{\mathrm{w}}$ is an $\ell\mathbb{Z}^3$-periodic function of the photon site without zeros.
For $k=0$ it takes the same value at the $48$ sites over $c$.
For the family $(\Psi',\bar{\Psi}')$ it is an $(\ell/L)\mathbb{Z}^3$-periodic function of $y'$, that is, an $\ell\mathbb{Z}^3$-periodic function of the site $Ly'$ of $L\mathbb{Z}^3$.
\end{named}

\begin{named}{Remark}
Proposition~22.1 corresponds to the Berezin integral of \cite[Section~14.6.1]{QEDI} together with the rules (Ber1)--(Ber5) stated in that section; for $1\le i\le5$ the rule (G$i$)$^{\mathrm{w}}$ is the counterpart of (Ber$i$), and (G0)$^{\mathrm{w}}$ is the counterpart of the normalisation fixed there.
It differs from them in the following points.
\begin{itemize}
\item[(1)] The index sets carry weights.
At $k=0$ the family consists of $2\cdot192\,n_0$ generators, and under the substitution of (G3)$^{\mathrm{w}}$ the top monomial is multiplied by a product of factors $u\cdot u^{-1}$, with $192$ factors for each photon site where the computation behind (Ber3) has four; compare (I.L0.6).
\item[(2)] In \cite[Section~14.6.1]{QEDI} the Berezin integral is normalised by the identity $\int\mathrm{D}\Psi\,\exp\left(-\left\langle\bar{\Psi},O\Psi\right\rangle_1\right)=\det O$ for every endomorphism $O$ of the space of spinors.
Here the normalisation is the single requirement $\int\mathrm{D}W\,\exp\left(-\left\langle\bar{W},W\right\rangle_{\nu}\right)=1$, and the identity for every $O$ is the assertion (G0)$^{\mathrm{w}}$.
\item[(3)] Every rule names the family that is replaced; in (G4)$^{\mathrm{w}}$ it is the primed family.
\item[(4)] The fields denoted $J$, $\bar{J}$ in (Ber2) are denoted $\mathrm{Sr}$, $\overline{\mathrm{Sr}}$ here, since $J$ is the cluster map.
\end{itemize}
\end{named}

\subsubsection*{Proposition~22.2 (normalisation of the block transformation; the density of level 0)}

\begin{named}{Proposition 22.2}
\begin{itemize}
\item[(a)] Let $k\ge0$, $\ell\in L\mathbb{N}$, $e_k>0$ and $b>0$.
Let $\omega$ be a real $\ell\mathbb{Z}^3$-periodic function on $L\mathbb{Z}^3$, and let $\mathrm{ff}$, $\overline{\mathrm{ff}}$ be fields on $L\mathbb{Z}^3/\ell\mathbb{Z}^3$ with four components whose values lie in the subalgebra generated by the generators other than $\Psi'$, $\bar{\Psi}'$.
Then
\[
c_{k}\cdot\int\mathrm{D}\Psi'\,\exp\left\{-b_{L}\cdot\left\langle\exp\left(-ie_k\cdot\omega\right)\bar{\Psi}'_L-\overline{\mathrm{ff}},\ \exp\left(ie_k\cdot\omega\right)\Psi'_L-\mathrm{ff}\right\rangle_L\right\}=1,\qquad c_{k}=b^{-4n_{k+1}}.
\]
\item[(b)] In the setting (S22) let $A'\in\mathrm{Dom}_{k+1}$ and $\mathscr{C}'\in\mathfrak{X}_{k+1}$, and let $F(\,\cdot\,,\mathscr{C}'_L;\Psi)$ be a map of the class described in Section~G3.2.0.
Thus $F(\,\cdot\,,\mathscr{C}'_L;\Psi)$ maps $\mathrm{Dom}_k$ into the even part of $\mathrm{Gr}[\Psi]$, and the coefficients of $\zeta\mapsto F\left(A^{\min}_{k}(A')+\mathrm{Sq}_{k}\zeta,\mathscr{C}'_L;\Psi\right)$ are integrable over $\mathrm{Box}_{k}$ with respect to $\mu_I$; a bounded Borel map is an example.
Then the identity (CV.3x)$^{\mathrm{ch}}$ holds; it expresses the Berezin integral over $\Psi'$ of $\mathbf{T}^{x}_{k}F$ as the integral over $\mathrm{Box}_{k}$, with respect to $\mu_I$, of the Berezin integral of $F$ over $\Psi$:
\begin{equation}
\int\mathrm{D}\Psi'\,\left(\mathbf{T}^{x}_{k}F\right)(A',\mathscr{C}';\Psi')=\int_{\mathrm{Box}_{k}}\mathrm{d}\mu_I(\zeta)\int\mathrm{D}\Psi\,F\left(A^{\min}_{k}(A')+\mathrm{Sq}_{k}\zeta,\mathscr{C}'_L;\Psi\right).
\tag*{(CV.3x)$^{\mathrm{ch}}$}
\end{equation}
For $k=0$, $\int\mathrm{D}\Psi$ denotes the integral over the $192\,n_0$ pairs of generators on $T^{\mathrm{f}}_0$ with the weight $\frac{1}{48}$.
\item[(c)] (Level $0$.)
Let $\ell\ge1$ be an integer, $n_0=\ell^3$, $q>0$ and $\mu>0$, let $\varepsilon^{\mathrm{v}}$ and $m$ be complex numbers, and let $\mathrm{Bf}$ be a complex bond function on $T_0=\mathbb{Z}^3/\ell\mathbb{Z}^3$ (for $\ell=1,2$: an $\ell\mathbb{Z}^3$-periodic bond function on $\mathbb{Z}^3$, as in Proposition~46).
Then $\det\left(D_0+\mu\right)\ne0$, where $D_0$ denotes here the operator $D_A$ of (I.2.3) at $A=0$, and
\begin{equation*}
\begin{split}
&\int\mathrm{D}\Psi\,\left(\det\left(D_0+\mu\right)\right)^{-1}\cdot\exp\left\{-\left\langle\bar{\Psi},\left(D_{\mathrm{Bf}}+\mu\right)\Psi\right\rangle^{\mathrm{ch}}-m\cdot\left\langle\bar{\Psi},\Psi\right\rangle^{\mathrm{ch}}-\varepsilon^{\mathrm{v}}\cdot n_0\right\}\\
&\qquad=\exp\left(-\varepsilon^{\mathrm{v}}\cdot n_0\right)\cdot\frac{\det\left(D_{\mathrm{Bf}}+\mu+m\right)}{\det\left(D_0+\mu\right)}.
\end{split}
\end{equation*}
In particular, in the flow ($\ell=n$, $\mu=\bar{m}\cdot\eta$, $\mathfrak{P}_0=0$; Sections~I.5.7 and~I.7.3) one has, for every real pair $(m_{\mathrm{in}},\varepsilon^{\mathrm{v}}_{\mathrm{in}})$, every $A_0\in\mathrm{Dom}_0$ and every $a\in\mathfrak{X}_0$,
\[
y^{x}_{0}\left[\varepsilon^{\mathrm{v}}_{\mathrm{in}},m_{\mathrm{in}},0\right](A_0;a)=\exp\left(-\varepsilon^{\mathrm{v}}_{\mathrm{in}}\cdot n^3\right)\cdot\frac{\det\left(D_{A_0+a}+\bar{m}\cdot\eta+m_{\mathrm{in}}\right)}{\det\left(D_0+\bar{m}\cdot\eta\right)}.
\]
No condition is imposed on the sign or on the size of $\bar{m}\cdot\eta+m_{\mathrm{in}}$; compare (I.K10).
For an arbitrary family $\mathfrak{P}_0$ the right side acquires the factor $\exp\left\{\sum_X\mathfrak{P}_0(X,a)\right\}$.
\end{itemize}
\end{named}

\begin{named}{Remark}
Proposition~22.2 corresponds to \cite[Lemma~14.23(a)--(c)]{QEDI}; the identity (CV.3x)$^{\mathrm{ch}}$ of (b) and the formula for $y^{x}_{0}$ in (c) correspond to the identity for $\int\mathrm{D}\Psi'\,\mathbf{T}^{x}_{k}F$ and to the formula for $y^{x}_{0}$ in \cite[Section~5.3.6]{QEDI}.
The comparison, point by point, covers agreements as well as differences.
\begin{itemize}
\item[(1)] In (b) and (c) at $k=0$ the family $(\Psi,\bar{\Psi})$ lives on $T^{\mathrm{f}}_0$ with the weight $\frac{1}{48}$; the determinants in (c) are of size $192\,n_0$; and the mass term involves all $192$ components, so that it is $m$ times the identity.
\item[(2)] In the block transformation of (b) the fields $\mathrm{ff}$ of (a) are $Q(\mathscr{B}_{k}+\mathscr{C})\Psi$ for $k\ge1$ and $Q(\mathscr{B}_{0}+\mathscr{C})J^{*}\Psi$ for $k=0$.
\item[(3)] Each clause carries its own hypotheses: (a) requires only $\ell\in L\mathbb{N}$, (b) the setting (S22), and (c) data of level $0$ only.
\item[(4)] In (c), as in \cite[Lemma~14.23(c)]{QEDI}, the symbol $D_0$ denotes the operator $D_A$ at $A=0$, and not the operator $D_0(A)=D_A+\mu$ of level $0$.
The same convention is fixed in \cite[Section~5.3.6]{QEDI}, where the mass $\mu$ is written $\mathfrak{m}$: there $D_0$ without argument denotes the operator $D_A$ at $A=0$, which is neither the operator $D_0(A)$ of level $0$ nor the set of the polymers of level $0$, and $D_0(0)=D_0+\mathfrak{m}$, with $\mathfrak{m}=\bar{m}\cdot\eta$ for the data of level $0$ of the flow.
\item[(5)] Clause (b) is stated for the whole class of maps $F$ for which $\mathbf{T}^{x}_{k}F$ is defined.
It therefore applies to $F=\mathscr{F}^{x}_{k}[\varepsilon^{\mathrm{v}},m,\mathscr{E}]$ by the continuity in $\zeta$ asserted in Proposition~22.3\,(a), without any continuity of $F$ on a subset of $\mathrm{Dom}_k$ having to be shown.
\end{itemize}
\end{named}

\begin{named}{Remark}
By Proposition~46\,(b), (d) and (e), the numerator $\det\left(D_{A_0+a}+\bar{m}\cdot\eta+m_{\mathrm{in}}\right)$ in (c) is the determinant $\det D^{\mathbb{T}}_{A}$, at $A=A_0+a$, of \cite[Definition~2.11]{QEDII} and \cite[Definition~2.13]{QEDII} on the torus $\mathbb{T}=C(n,n,n)$, formed with the mass parameter
\[
M_{\mathrm{ch}}=\bar{m}\cdot\eta+m_{\mathrm{in}}+\frac{\rho}{2}\cdot\mathrm{sw}.
\]
\end{named}

\subsubsection*{Proposition~22.3 (the one-step formula with slot)}

\begin{named}{Proposition 22.3}
In the setting (S22) let $\varepsilon^{\mathrm{v}}$ and $m$ be complex numbers, let $\mathscr{E}\in K^{x,C}_{k}$, with $\mathscr{E}=0$ if $k=0$, let $\mathfrak{P}_k$ be an arbitrary family as in Section~G3.2.0, let $A'\in\mathrm{Dom}_{k+1}$, whose harmonic part is arbitrary, and let $\mathscr{C}'\in\mathfrak{X}_{k+1}$.
Define
\[
m':=m+m(\mathscr{E}),\qquad\mathscr{E}^{r}:=R\mathscr{E},\qquad m^{+}:=L\cdot m',\qquad\varepsilon^{\mathrm{v},+}:=L^3\cdot\left(\varepsilon^{\mathrm{v}}+\varepsilon(\mathscr{E})+\delta\varepsilon^{0}_{k}\right)
\]
(for $k=0$ this means $m'=m$, $\mathscr{E}^{r}=0$, $m^{+}=L\cdot m$ and $\varepsilon^{\mathrm{v},+}=L^3\cdot\left(\varepsilon^{\mathrm{v}}+\delta\varepsilon^{0}_{0}\right)$);
\begin{gather*}
\mathscr{A}':=H_{k+1}A',\qquad\mathscr{T}':=\mathscr{A}'+\mathscr{C}',\\
\psi':=H_{k+1}(\mathscr{T}')\Psi',\qquad\bar{\psi}':=\bar{H}_{k+1}(\mathscr{T}')\bar{\Psi}',\qquad\chi':=\delta_{\alpha,\mathscr{T}'}\psi',\qquad\bar{\chi}':=\delta_{\alpha,-\mathscr{T}'}\bar{\psi}',
\end{gather*}
which are fields on $T^{\mathrm{f}}_{k+1}$ and on $\mathrm{Pz}_{k+1}$;
\begin{gather*}
A:=\mathscr{A}'_L,\qquad\mathscr{C}:=\mathscr{C}'_L,\qquad\mathscr{T}:=A+\mathscr{C}=\mathscr{T}'_L,\\
\psi^{0}:=\psi'_L,\qquad\bar{\psi}^{0}:=\bar{\psi}'_L,\qquad\chi^{0}:=\chi'_L,\qquad\bar{\chi}^{0}:=\bar{\chi}'_L,\qquad\psi^{0,\#}:=\left(\psi^{0},\chi^{0}\right),
\end{gather*}
the last one together with the barred fields;
\[
\Psi^{c}:=T_k(\mathscr{T})\psi^{0},\qquad\bar{\Psi}^{c}:=\bar{T}_k(\mathscr{T})\bar{\psi}^{0}
\]
(for $k=0$: $\Psi^{c}=\psi^{0}$ and $\bar{\Psi}^{c}=\bar{\psi}^{0}$, which are fields on $T^{\mathrm{f}}_0$ with $192$ components per unit site);
\[
\mathscr{Z}:=H_k\,C\,C_{k}^{1/2}\mathsf{Z}\qquad\left(\mathsf{Z}\in\mathrm{Box}_{k}\right)
\]
(for $k=0$: $\mathscr{Z}=C\,C_{0}^{1/2}\mathsf{Z}$);
\[
\Psi^{f}:=\Psi^{c}+\Gamma_k(\mathscr{T})W,\qquad\bar{\Psi}^{f}:=\bar{\Psi}^{c}+\bar{W},
\]
where $W$, $\bar{W}$ are the generators of Section~I.3.6 (for $k=0$ they live on $T^{\mathrm{f}}_0$ with the weight $\frac{1}{48}$, and $\left(\Gamma_0(\mathscr{T})W\right)(x)=\int\mathrm{d}x'\,\Gamma_0(\mathscr{T})(x,x')W(x')$);
\begin{equation*}
\begin{split}
V^{\mathrm{ex},x}_{k}:={}&\mathrm{Act}^{x}_{k}[m',\mathscr{E}^{r}](A+\mathscr{Z},\mathscr{C};\Psi^{f})-\mathrm{Act}^{x}_{k}[m',\mathscr{E}^{r}](A,\mathscr{C};\Psi^{f})\\
&+I^{x}_{k}[m',\mathscr{E}^{r}](A,\mathscr{C};\Psi^{f})-I^{x}_{k}[m',\mathscr{E}^{r}](A,\mathscr{C};\Psi^{c}),
\end{split}
\end{equation*}
and
\[
\Xi^{x}_{k}:=\int\mathrm{d}\mu_I(W)\,\mathrm{d}\hat{\mu}(\mathsf{Z})\,\exp\left(V^{\mathrm{ex},x}_{k}\right).
\]
Then the following hold.
\begin{itemize}
\item[(a)] (All objects are defined.)
For every $\mathsf{Z}\in\mathrm{Box}_{k}$ the pairs $(A,\mathscr{C})$ and $(A+\mathscr{Z},\mathscr{C})$ lie in $\tilde{R}_k\times\mathfrak{X}_k$, and $A\in\tilde{R}_k^{-}$.
The connections $\mathscr{T}$ and $\mathscr{T}+\mathscr{Z}$ and their negatives lie in $\tilde{G}_k$.
The operators $S_k$ (for $k\ge1$), $H_k$, $\bar{H}_k$, $D_k$ and $\Gamma_k$ exist at $\mathscr{T}$, at $\mathscr{T}+\mathscr{Z}$ and at $0$, and the operators $S_{k+1}$, $H_{k+1}$, $\bar{H}_{k+1}$, $D_{k+1}$ exist at $\mathscr{T}'$ and at $0$.
The numbers $\det D_k(0)$, $\delta Z_k(0)$ and $\det D_{k+1}(0)$ are different from $0$.
The integrand of $\left(\mathbf{T}^{x}_{k}\mathscr{F}^{x}_{k}[\varepsilon^{\mathrm{v}},m,\mathscr{E}]\right)(A',\mathscr{C}';\Psi')$ is a continuous function of $\zeta$ on the compact set $\mathrm{Box}_{k}$, and $V^{\mathrm{ex},x}_{k}$ is an even element which depends continuously on $\mathsf{Z}\in\mathrm{Box}_{k}$.
\item[(b)] (The fields.)
One has
\[
\psi^{0}=H_k(\mathscr{T})\Psi^{c},\qquad\bar{\psi}^{0}=\bar{H}_k(\mathscr{T})\bar{\Psi}^{c},\qquad\chi^{0}=\delta_{\alpha,\mathscr{T}}\psi^{0},\qquad\bar{\chi}^{0}=\delta_{\alpha,-\mathscr{T}}\bar{\psi}^{0};
\]
in other words, $\psi^{0,\#}=H_k^{\#}(\mathscr{T})\Psi^{c}$ on the unbarred and on the barred fields.
\item[(c)] (The formula.)
One has the following identity, the one-step formula with slot, denoted (2.4a$^{x}$)$^{\mathrm{ch}}$:
\begin{equation}
\begin{split}
&\left(\mathbf{T}^{x}_{k}\mathscr{F}^{x}_{k}[\varepsilon^{\mathrm{v}},m,\mathscr{E}]\right)(A',\mathscr{C}';\Psi')\\
&\quad=\left(\det D_{k+1}(0)\right)^{-1}\cdot\exp\Bigl\{-\left\langle\bar{\Psi}',D_{k+1}(\mathscr{T}')\Psi'\right\rangle_1-m^{+}\cdot\left\langle\bar{\psi}',\psi'\right\rangle-\varepsilon^{\mathrm{v},+}\cdot n_{k+1}\\
&\qquad+\sum_{X\in D_k}\mathscr{E}^{r}\left(X;A,\mathscr{C},\psi^{0,\#}\right)+\sum_{X\in D_k}\mathfrak{P}_k(X,\mathscr{C})\Bigr\}\cdot\frac{\delta Z_k(\mathscr{T})}{\delta Z_k(0)}\cdot\Xi^{x}_{k}.
\end{split}
\tag*{(2.4a$^{x}$)$^{\mathrm{ch}}$}
\end{equation}
Here $\left\langle\bar{\psi}',\psi'\right\rangle=\int\mathrm{d}x\,\bar{\psi}'(x)\cdot\psi'(x)$ is the form of Section~I.3.6 at level $k+1$, which involves all $192$ components over each photon site of $T_{k+1}$, and $\mathscr{E}^{r}\left(X;A,\mathscr{C},\psi^{0,\#}\right)$ is the image of $\mathscr{E}^{r}(X;A,\mathscr{C},\psi,\chi)$ under the substitution $\psi\to\psi^{0}$, $\chi\to\chi^{0}$, $\bar{\psi}\to\bar{\psi}^{0}$, $\bar{\chi}\to\bar{\chi}^{0}$.
Written out at level $0$, the formula reads
\begin{equation*}
\begin{split}
&\left(\mathbf{T}^{x}_{0}\mathscr{F}^{x}_{0}[\varepsilon^{\mathrm{v}},m,0]\right)(A',\mathscr{C}';\Psi')\\
&\quad=\left(\det D_1(0)\right)^{-1}\cdot\exp\Bigl\{-\left\langle\bar{\Psi}',D_1(\mathscr{T}')\Psi'\right\rangle_1-L\cdot m\cdot\left\langle\bar{\psi}',\psi'\right\rangle-L^3\cdot\left(\varepsilon^{\mathrm{v}}+\delta\varepsilon^{0}_{0}\right)\cdot n_1\\
&\qquad+\sum_X\mathfrak{P}_0(X,\mathscr{C})\Bigr\}\cdot\frac{\delta Z_0(\mathscr{T})}{\delta Z_0(0)}\cdot\Xi^{x}_{0},
\end{split}
\end{equation*}
with
\[
\delta Z_0(B)=\det\left(D_B+\mu+b_{L}\cdot JP(B)J^{*}\right)\qquad\left(\text{a determinant of size }192\,n_0\right)
\]
and
\begin{equation*}
\begin{split}
\Xi^{x}_{0}={}&\int\mathrm{d}\mu_I(W)\,\mathrm{d}\hat{\mu}(\mathsf{Z})\,\exp\Bigl\{-\left\langle\bar{\Psi}^{f},\left(D_{\mathscr{T}+\mathscr{Z}}-D_{\mathscr{T}}\right)\Psi^{f}\right\rangle^{\mathrm{ch}}\\
&-m\cdot\left[\left\langle\bar{\Psi}^{f},\Psi^{f}\right\rangle^{\mathrm{ch}}-\left\langle\bar{\Psi}^{c},\Psi^{c}\right\rangle^{\mathrm{ch}}\right]\Bigr\}.
\end{split}
\end{equation*}
\item[(d)] The right side of (2.4a$^{x}$)$^{\mathrm{ch}}$ depends on $A'$ only through $\mathscr{A}'=H_{k+1}A'$.
Neither side changes when $\omega_k$ is replaced by $\omega_k+t$ with a real number $t$.
\end{itemize}
\end{named}

\begin{named}{Remark}
Proposition~22.3 corresponds to \cite[Theorem~14.24]{QEDI}, which is \cite[Proposition~5.56]{QEDI}, together with the setting of \cite[Section~14.8.1]{QEDI} and with \cite[Remark~14.25]{QEDI}.
Point by point, the agreements and the differences are the following.
\begin{itemize}
\item[(1)] Level $0$ is written out: $\Psi$, $W$, $\Psi^{c}$ and $\Psi^{f}$ live on $T^{\mathrm{f}}_0$; $\mathrm{d}\mu_I(W)=\exp\left(-\left\langle\bar{W},W\right\rangle^{\mathrm{ch}}\right)\mathrm{D}W$; the block transformation $\mathbf{T}^{x}_{0}$ contains $Q(\,\cdot\,)J^{*}$; and $\delta Z_0$ is a determinant of size $192\,n_0$ which contains the determinant of the heavy components.
\item[(2)] For $k\ge1$ the fields $\psi$ and $\chi$ live on $T^{\mathrm{f}}_k$ and on $\mathrm{Pz}_k$, and the forms are those of Section~I.3.6.
\item[(3)] Here, as in \cite[Section~14.8.1]{QEDI}, \cite[Section~5.3.6]{QEDI} and \cite[Proposition~5.56]{QEDI}, the family $\mathfrak{P}_k$ with which $\mathscr{F}^{x}_{k}$ is formed is an arbitrary family of complex functions on $\mathfrak{X}_k$, and no system of primitives enters Proposition~22.3, just as none enters \cite[Proposition~5.56]{QEDI}; along a flow, \cite[Section~5.3.6]{QEDI} takes for $\mathfrak{P}_k$ the family of \cite[Section~5.3.8]{QEDI}.
\item[(4)] By \cite[Remark~14.25(iii)]{QEDI}, a constant contained in $A'$ occurs on the right side of the identity of \cite[Theorem~14.24]{QEDI} only through the constants which it produces in $\mathscr{A}'$ and in $A$, because there $A'$ enters only through $\mathscr{A}'=H_{k+1}A'$; on the left side $A'$ enters in addition through $A^{\min}_{k}(A')$, through $\mathscr{B}_{k}$ and through $\omega_{k}(A')$, the function $\omega_{k}$ of a nonzero constant is not constant on $T_k$, and $\omega_{k}(A')$ is removed by gauge covariance, whatever the harmonic part of $A'$ is.
Proposition~22.3 agrees with this: that $A'$ enters the right side of (2.4a$^{x}$)$^{\mathrm{ch}}$ only through $H_{k+1}A'$ is the first sentence of clause (d), and the second sentence of clause (d) corresponds to \cite[Remark~14.25(i)]{QEDI}.
On the left side of (2.4a$^{x}$)$^{\mathrm{ch}}$ a constant bond function $\mathrm{cst}$ contained in $A'$ enters through $A^{\min}_{k}(A')$ and through $\mathscr{B}_{k}=H_kA^{\min}_{k}(A')=\mathscr{A}'_L-d\omega_{k}(A')$.
The phases $\exp\left(\pm ie_k\cdot\omega_k\right)$ at the sites of $L\mathbb{Z}^3$ do not contain it, since $\omega_{k}(\mathrm{cst})$ vanishes on $L\mathbb{Z}^3$.
To see this, let $y'\in L\mathbb{Z}^3$.
The point reflection $c\mapsto 2y'-c$ (at level $k+1$: $c\mapsto 2y'/L-c$) maps each of the lattices $\mathbb{L}_k$, $\mathbb{Z}^3$ and $L\mathbb{Z}^3$ onto itself and blocks onto blocks, and it maps the paths of the averaging operators and of $\tau$ onto paths of the same kind.
Hence the minimisation problems which define $H_k$, $H_{k+1}$ and $A^{\min}_{k}$, together with their side conditions, are invariant under it, and their unique solutions commute with it.
The reflection maps $\mathrm{cst}$ to $-\mathrm{cst}$.
Since it maps blocks onto blocks, the condition $Q_{k+1}\omega=0$ on a function $\omega$ on the sites of $T_k$ is preserved when $\omega$ is replaced by $c\mapsto\omega(2y'-c)$.
By \cite[Theorem~14.15(b)]{QEDI}, $\omega_{k}(A')$ is the unique function on the sites of $T_k$ with $Q_{k+1}\omega_{k}(A')=0$ for which $H_kA^{\min}_{k}(A')=\left(H_{k+1}A'\right)_L-d\omega_{k}(A')$, and it depends linearly on $A'$.
It follows that $\omega_{k}(\mathrm{cst})(2y'-c)=-\omega_{k}(\mathrm{cst})(c)$, and the choice $c=y'$ gives $\omega_{k}(\mathrm{cst})(y')=0$.
On $T_k$, on the other hand, the function $\omega_{k}(\mathrm{cst})$ is not constant if $\mathrm{cst}\ne0$.
Otherwise one would have $H_kA^{\min}_{k}(\mathrm{cst})=\left(H_{k+1}\mathrm{cst}\right)_L=\mathrm{cst}_L$ and $A^{\min}_{k}(\mathrm{cst})=Q_kH_kA^{\min}_{k}(\mathrm{cst})=Q_k\mathrm{cst}_L=\mathrm{cst}_L$, by \cite[Proposition~14.4(c)]{QEDI} and \cite[Proposition~14.8(ii), (iii)]{QEDI}, $Q_k$ being the photon average; this contradicts $\tau A^{\min}_{k}(\mathrm{cst})=0$.
In the proof of Proposition~22.3 the function $\omega_{k}(A')$ on $T_k$ is eliminated by means of (G3)$^{\mathrm{w}}$, Proposition~21.1\,(d) and (I.k4), for an arbitrary harmonic part of $A'$.
\item[(5)] Unlike \cite[Theorem~14.24]{QEDI}, Proposition~22.3 assumes the setting (S22) and no hypotheses of other statements.
In \cite[Lemma~5.19]{QEDI} the invariance of $\mathrm{Dom}_k$ is asserted under the hypothesis of \cite[Definition~5.69]{QEDI}, which includes a system of primitives; this assertion does not enter here.
The other assertions of \cite[Lemma~5.19]{QEDI}, on the sets $R_k$, $c\cdot\tilde{R}_k$ with $0<c\le1$, $\tilde{R}_k^{-}$, $G_k$, $\tilde{G}_k$ and $\mathfrak{X}_k$, require only $k\ge0$, $\ell\in\mathrm{Cl}$ and $0<e_k\le e_0$; hence they hold in the setting (S22), which contains no system of primitives.
\item[(6)] In \cite[Theorem~14.24]{QEDI} the last argument $\Psi^{f}$, respectively $\Psi^{c}$, of $\mathrm{Act}^{x}_{k}$ and $I^{x}_{k}$ has the meaning fixed in \cite[Proposition~5.56]{QEDI}: the pair $(\Psi,\bar{\Psi})$ is replaced by $(\Psi^{f},\bar{\Psi}^{f})$, respectively by $(\Psi^{c},\bar{\Psi}^{c})$.
The insertion of a pair of fields in the last slot of $I^{x}_{k}$ and $\mathrm{Act}^{x}_{k}$ by the homomorphism of Section~G3.2.0 (at $k=0$, of a pair of fields on $T^{\mathrm{f}}_0$) agrees with that of \cite[Theorem~14.24]{QEDI}, where the pair is inserted by the homomorphism of \cite[Section~14.6.1]{QEDI} which maps $\bar{\Psi}$ and $\Psi$ to the fields with the same superscript and fixes the remaining generators; in both cases the operators remain those of the total connection formed from the first two arguments.
In the proof of Proposition~22.3 the rules (G3)$^{\mathrm{w}}$ and (G1)$^{\mathrm{w}}$ act as homomorphisms on the family $(\Psi,\bar{\Psi})$ itself, and no new family of generators of integration occurs, whereas the proof of \cite[Theorem~14.24]{QEDI} introduces new generators $\Psi^{n}$, $\bar{\Psi}^{n}$ and $W'$, $\bar{W}'$.
The family $(\Psi,\bar{\Psi})$ then takes the place of $(W',\bar{W}')$ in Proposition~21.4\,(a) and in (G4)$^{\mathrm{w}}$.
\item[(7)] Where \cite[Theorem~14.24]{QEDI} asserts that all objects are well defined, clause (a) names the connections at which each operator exists and the numbers that are different from $0$.
The class of maps $F$ on which $\mathbf{T}^{x}_{k}$ acts and the domain of $\omega_k$, the sites of $T_k$, are fixed in Section~G3.2.0.
\item[(8)] As in \cite[Theorem~14.24]{QEDI}, the requirement $\mathscr{E}=0$ for $k=0$ is made also for complex data.
The complex data admitted here are those of the last sentence of \cite[Proposition~5.56]{QEDI} and of \cite[Theorem~14.24]{QEDI}: complex numbers $\varepsilon^{\mathrm{v}}$ and $m$ and $\mathscr{E}\in K^{x,C}_{k}$, with $\mathscr{E}=0$ if $k=0$.
\item[(9)] Clause (b) corresponds to identities of \cite[Section~17.1.4]{QEDI} and of \cite[Section~17.12]{QEDI}.
\end{itemize}
\end{named}

\subsection*{23. The new displays of the step}

In \cite[Section~17.1.2]{QEDI} the operators $D_B$, $Q_k(B)$ and $T_k(B)$ of the Wilson fermion are written out in a single display.
For the chamber fermion its place is taken by the operator $D_B$ of (I.2.3), by the operator $Q_k(B)\cdot J^{*}$ of Section~I.3.3, by the operator
\[
T_k(B)=b_k^{-1}\cdot Q_k(B)\cdot J^{*}\cdot\left(D_B+\mu\right)+Q_k(B)\cdot J^{*}\qquad(k\ge1)
\]
of Section~I.3.5, and, at level $0$, by $T_0=1$ of Section~I.3.9.
The three relations of \cite[Lemma~17.1(d)]{QEDI} for the matrices $P^{\boldsymbol{\sigma}}_{\nu}$ are replaced by (I.m2), (I.m3) and (I.m4), and the relation $P^{+}_{\nu}+P^{-}_{\nu}=1$ of \cite[Section~17.1.2]{QEDI} is replaced by (I.m1).
The two displays (I.5.4) and (I.5.5) have the following properties.

\subsubsection*{Proposition 23.1 (levels $k\ge1$: the rows of $T_k$, its two local pieces, the substitution)}

\begin{named}{Definition}
Let $k\ge1$ be an integer, $\ell\in\mathrm{Cl}$, $\mu$ a complex number, $e_k>0$, $B$ a complex bond function on $T_k$, and $Y$ a union of $M$-cubes of $T_k$; in the step map, $Y$ is a polymer $\mathbf{Y}$.
For every unit site $y$ of $Y$, every $f\in\mathfrak{S}^{\mathrm{ch}}(\eta_k,\ell)$ and every function $F$ of pairs (Section~I.1.7) put, with $x=(c,g)$ and $c_3=c_3(x)$,
\begin{equation}
\begin{aligned}
\left(T^{(1)}(B,Y)\,f\right)(y)&:=\int_{\Delta_y}\mathrm{d}x\,\varphi_B(y,c)\cdot\left\{\left(1+\frac{\mu}{b_k}\right)f(x)-\left(b_k\cdot\eta_k\right)^{-1}\left[\sum_{i\le2}M_i(x)\left(f\left(x^{(i)}\right)-f(x)\right)\right.\right.\\
&\qquad\qquad\left.\left.{}+\mathbf{1}[c_3\in Y]\cdot M_3(x)\left(u_B(c,c_3)\cdot f\left(x^{(3)}\right)-f(x)\right)\right]\right\},\\
\left(T^{(2)}(B,Y)\,F\right)(y)&:=-b_k^{-1}\cdot\eta_k^{\alpha-1}\int_{\Delta_y}\mathrm{d}x\,\mathbf{1}[c_3\notin Y]\cdot\varphi_B(y,c)\cdot M_3(x)\cdot F\left(x,x^{(3)}\right),\\
\bar{T}^{(i)}(B,Y)&:=\mathsf{C}^{-1}\cdot T^{(i)}(-B,Y)\cdot\mathsf{C}\qquad(i=1,2).
\end{aligned}
\tag{23.1}
\end{equation}
At the unit sites $y$ outside $Y$ all four operators have zero rows.
\end{named}

Thus $T^{(1)}(B,Y)$ and $\bar{T}^{(1)}(B,Y)$ map $\mathfrak{S}^{\mathrm{ch}}(\eta_k,\ell)$ to $\mathfrak{S}(1,\ell)$, and $T^{(2)}(B,Y)$ and $\bar{T}^{(2)}(B,Y)$ map the functions of pairs to $\mathfrak{S}(1,\ell)$.
For $y\in Y$ the rows (23.1) are the displays (I.5.4) with $\mathbf{Y}=Y$: there $\Delta_y$ lies in $Y$, so that the condition that $x^{(i)}$ lie over $\mathbf{Y}$ is satisfied for $i\le2$.
The operator $T^{(i)}(B)$ of (I.5.4) is $T^{(i)}(B,\mathbf{Y})$.
The rows at the unit sites outside $Y$ vanish (Section~I.5.9; Proposition~11\,(a)); consequently $T^{(i)}(B,Y)$ is an operator between the spaces of Section~I.1.9, and $\left|T^{(i)}(B,Y)\right|_{1,\infty}$ is the supremum over the rows at the unit sites of $Y$.
The step map involves these rows only.

\begin{named}{Definition}
For a polymer $\mathbf{Y}$ let $\mathrm{Gr}[\psi,\chi,\Psi\,|\,\mathbf{Y}]$ denote the subalgebra of $\mathrm{Gr}[\psi,\chi,\Psi]$ generated by all $\psi(\xi)$, all $\chi(\boldsymbol{\zeta})$, and the $\Psi(u)$ with $u=(y,\beta,\omega)$, $y$ a unit site of $\mathbf{Y}$.
The substitution
\[
\mathrm{sub}_{\mathbf{Y},B}\colon\mathrm{Gr}[\psi,\chi,\Psi\,|\,\mathbf{Y}]\longrightarrow\mathrm{Gr}[\psi,\chi]
\]
is the homomorphism of algebras which fixes every $\psi(\xi)$ and every $\chi(\boldsymbol{\zeta})$ and which acts, for the unit sites $y$ of $\mathbf{Y}$, by
\begin{equation}
\begin{aligned}
\Psi(y)&\longmapsto\left(T^{(1)}(B,\mathbf{Y})\,\psi\right)(y)+\left(T^{(2)}(B,\mathbf{Y})\,\chi\right)(y),\\
\bar{\Psi}(y)&\longmapsto\left(\bar{T}^{(1)}(B,\mathbf{Y})\,\bar{\psi}\right)(y)+\left(\bar{T}^{(2)}(B,\mathbf{Y})\,\bar{\chi}\right)(y).
\end{aligned}
\tag*{(8.3a)$^{\mathrm{ch}}$}
\end{equation}
The display (8.3a)$^{\mathrm{ch}}$ is called the substitution rule; the operators $T^{(i)}(B,\mathbf{Y})$ and $\bar{T}^{(i)}(B,\mathbf{Y})$ in it are those of (23.1).
This homomorphism exists and is unique, because the images are odd and have the antiperiodicity of $\Psi$ and $\bar{\Psi}$.
If the coefficients of an element $F$ of $\mathrm{Gr}[\psi,\chi,\Psi\,|\,\mathbf{Y}]$ depend on $(A,\mathscr{C},\mathsf{Z})$, then $B=\mathscr{T}$ is taken at the point at which $F$ is evaluated.
\end{named}

\begin{named}{Proposition 23.1}
With these data the following assertions hold.
\begin{itemize}
\item[(a)] \emph{Rows of $T_k$; identity.}
For every $f\in\mathfrak{S}^{\mathrm{ch}}(\eta_k,\ell)$ and every unit site $y$,
\begin{equation}
\begin{aligned}
\left(T_k(B)\,f\right)(y)=\int_{\Delta_y}\mathrm{d}x\,\varphi_B(y,c)\cdot&\left\{\left(1+\frac{\mu}{b_k}\right)f(x)\right.\\
&\left.{}-\left(b_k\cdot\eta_k\right)^{-1}\sum_{i=0}^{3}M_i(x)\left[U_B\left(x,x^{(i)}\right)\cdot f\left(x^{(i)}\right)-f(x)\right]\right\},
\end{aligned}
\tag{23.2}
\end{equation}
and at the unit sites of $Y$
\[
T^{(1)}(B,Y)\,f+T^{(2)}(B,Y)\,\delta_{\alpha,B}f=T_k(B)\,f,\qquad\bar{T}^{(1)}(B,Y)\,f+\bar{T}^{(2)}(B,Y)\,\delta_{\alpha,-B}f=\bar{T}_k(B)\,f.
\]
The same holds for fields.
\item[(b)] \emph{What a row contains and what it reads.}
Let $y$ be a unit site of $Y$.
The row of $T^{(1)}(B,Y)$ at $y$ contains $f$ at the sites $x$ over $\Delta_y$ and at the sites $x^{(3)}$ with $c\in\Delta_y$ and $c_3\in Y$; the row of $T^{(2)}(B,Y)$ at $y$ contains $F$ at the pairs $\left(x,x^{(3)}\right)$ with $c\in\Delta_y$ and $c_3\notin Y$.
These sites lie over $Y$, and these pairs are of kind (II) for $Y$, with the first point over $Y$.
In the sense of Section~I.1.9, both rows read $B$ on the bonds contained in $\Delta_y$ and on the bonds $(c,c_3)$ with $c\in\Delta_y$; each of these bonds has an end point in $\Delta_y$.
The rows read $Y$ through the site $y$ and through the sites $c_3$ with $c\in\Delta_y$.
All the sites, pairs and bonds just named lie within sup-distance $\frac{1+\eta_k}{2}$ of $y$.
The entries of the rows are entire functions of the numbers $B(b)$.
The same holds for $\bar{T}^{(i)}(B,Y)$.
\item[(c)] \emph{Covariances of the operators} ($i=1,2$).
\emph{Gauge.}
Let $\boldsymbol{\lambda}$ be as in Proposition~21.1\,(d), and let $\left(\exp\left(i\,e_k\cdot\boldsymbol{\lambda}\right)F\right)(x,x'):=\exp\left(i\,e_k\cdot\boldsymbol{\lambda}(c)\right)F(x,x')$ on functions of pairs.
Then
\begin{gather*}
T^{(i)}\left(B-d\boldsymbol{\lambda},Y\right)\exp\left(i\,e_k\cdot\boldsymbol{\lambda}\right)=\exp\left(i\,e_k\cdot\boldsymbol{\lambda}\right)T^{(i)}(B,Y),\\
\bar{T}^{(i)}\left(B-d\boldsymbol{\lambda},Y\right)\exp\left(-i\,e_k\cdot\boldsymbol{\lambda}\right)=\exp\left(-i\,e_k\cdot\boldsymbol{\lambda}\right)\bar{T}^{(i)}(B,Y).
\end{gather*}
\emph{Lattice.}
Let $\tilde{g}=(S,a)\in\mathbb{G}_k$, let $\tilde{g}$ act as in Section~I.4.3, and put
\[
\left(\mathrm{glb}\,f\right)(x):=\left(S^{-1}\right)^{T}f\left(\tilde{g}^{-1}x\right),\qquad\left(\mathrm{glb}\,F\right)(x,x'):=\left(S^{-1}\right)^{T}F\left(\tilde{g}^{-1}x,\tilde{g}^{-1}x'\right),
\]
and, on $\mathfrak{S}(1,\ell)$, $\left(\mathrm{glb}\,v\right)(y):=\left(S^{-1}\right)^{T}v\left(\tilde{g}^{-1}y\right)$.
Then
\[
T^{(i)}\left(\tilde{g}B,\tilde{g}Y\right)\tilde{g}=\tilde{g}\,T^{(i)}(B,Y),\qquad\bar{T}^{(i)}\left(\tilde{g}B,\tilde{g}Y\right)\mathrm{glb}=\mathrm{glb}\,\bar{T}^{(i)}(B,Y).
\]
\emph{Charge conjugation.}
\[
T^{(i)}(-B,Y)\cdot\mathsf{C}=\mathsf{C}\cdot\bar{T}^{(i)}(B,Y),\qquad\bar{T}^{(i)}(-B,Y)\cdot\mathsf{C}=\mathsf{C}\cdot T^{(i)}(B,Y).
\]
\emph{Conjugation.}
If $\mu$ is real, then
\[
\overline{T^{(i)}\left(\overline{B},Y\right)}=\varkappa^{-1}\cdot T^{(i)}(-B,Y)\cdot\varkappa,\qquad\overline{\bar{T}^{(i)}\left(\overline{B},Y\right)}=\varkappa^{-1}\cdot\bar{T}^{(i)}(-B,Y)\cdot\varkappa,
\]
where $\overline{O}$ denotes the operator $f\mapsto\overline{O\,\overline{f}}$.
\item[(d)] \emph{The substitution: locality and parity.}
For $F\in\mathrm{Gr}[\psi,\chi,\Psi\,|\,\mathbf{Y}]$ the element $\mathrm{sub}_{\mathbf{Y},B}\,F$ contains, apart from the generators $\psi(\xi)$ and $\chi(\boldsymbol{\zeta})$ contained in $F$, only generators $\psi(\xi)$ with $\xi$ over $\mathbf{Y}$ and generators $\chi(\boldsymbol{\zeta})$ at pairs of kind (II) for $\mathbf{Y}$ whose first point lies over $\mathbf{Y}$.
Its coefficients are polynomials in the coefficients of $F$ and in entries of the rows of (b) at the unit sites of $\mathbf{Y}$, and these entries read $B$ on bonds that meet $\mathbf{Y}$.
The homomorphism $\mathrm{sub}_{\mathbf{Y},B}$ maps every generator to a linear combination of generators with the same bar index.
Hence it maps every monomial of $\mathrm{Gr}[\psi,\chi,\Psi\,|\,\mathbf{Y}]$ to a linear combination of monomials with the same number of unbarred factors and the same number of barred factors, all factors $\psi(\xi)$, $\chi(\boldsymbol{\zeta})$ and $\Psi(u)$ being counted.
In particular it preserves the difference of these two numbers and the parity of their sum, and it maps even elements to even elements.
\item[(e)] \emph{The substitution: evaluation.}
Let $\mathrm{ev}$ be a homomorphism, linear over $\mathbb{C}$, of $\mathrm{Gr}[\psi,\chi]$ into a Grassmann algebra such that $\mathrm{ev}(\chi)=\delta_{\alpha,B}\,\mathrm{ev}(\psi)$ and $\mathrm{ev}(\bar{\chi})=\delta_{\alpha,-B}\,\mathrm{ev}(\bar{\psi})$ at the pairs of kind (II) for $\mathbf{Y}$.
Then $\mathrm{ev}\circ\mathrm{sub}_{\mathbf{Y},B}$ is the homomorphism, linear over $\mathbb{C}$, of $\mathrm{Gr}[\psi,\chi,\Psi\,|\,\mathbf{Y}]$ which coincides with $\mathrm{ev}$ on $\psi$ and $\chi$ and maps $\Psi(y)$ to $\left(T_k(B)\,\mathrm{ev}(\psi)\right)(y)$ and $\bar{\Psi}(y)$ to $\left(\bar{T}_k(B)\,\mathrm{ev}(\bar{\psi})\right)(y)$.
If, moreover, $S_k(B)$ exists and $\mathrm{ev}(\psi)=H_k(B)\,\Phi$, $\mathrm{ev}(\bar{\psi})=\bar{H}_k(B)\,\bar{\Phi}$ for fields $\Phi$, $\bar{\Phi}$ on $T^0_k$, then this homomorphism maps $\Psi(y)$ to $\Phi(y)$ and $\bar{\Psi}(y)$ to $\bar{\Phi}(y)$.
\item[(f)] \emph{The substitution: four covariances.}
For the automorphisms of Section~I.4.5, extended to $\Psi$ as described there, one has
\begin{align*}
a_{\tilde{g}}\circ\mathrm{sub}_{\tilde{g}\mathbf{Y},\tilde{g}B}&=\mathrm{sub}_{\mathbf{Y},B}\circ a_{\tilde{g}}\quad\text{on }\mathrm{Gr}[\psi,\chi,\Psi\,|\,\tilde{g}\mathbf{Y}],\ \text{for }\tilde{g}\in\mathbb{G}_k;\\
G_{\Lambda}\circ\mathrm{sub}_{\mathbf{Y},B^{1}}&=\mathrm{sub}_{\mathbf{Y},B}\circ G_{\Lambda}\quad\text{for every real function }\Lambda\text{ of }c\text{ and every }B^{1}\\
&\qquad\qquad\text{with }B^{1}=B-d\Lambda\text{ on the bonds that meet }\mathbf{Y};\\
\mathrm{Cg}\circ\mathrm{sub}_{\mathbf{Y},-B}&=\mathrm{sub}_{\mathbf{Y},B}\circ\mathrm{Cg};\\
\mathbf{K}\circ\mathrm{sub}_{\mathbf{Y},-\overline{B}}&=\mathrm{sub}_{\mathbf{Y},B}\circ\mathbf{K}\quad\text{if }\mu\text{ is real}.
\end{align*}
\item[(g)] \emph{The substitution: norm.}
Let the constants be those of Section~I.8, and let $\mu\in(0,1]$, $e_k\in(0,e_0]$ and $B\in\tilde{R}_k+\mathfrak{X}_k$.
For weights $\mathsf{w}_1,\mathsf{w}_2,\mathsf{w}_3>0$ with $C_T\cdot\left(\mathsf{w}_1+\mathsf{w}_2\right)\le\mathsf{w}_3$ and for every $F\in\mathrm{Gr}[\psi,\chi,\Psi\,|\,\mathbf{Y}]$,
\[
\left\|\mathrm{sub}_{\mathbf{Y},B}\,F\right\|_{(\mathsf{w}_1,\mathsf{w}_2)}\le\left\|F\right\|_{(\mathsf{w}_1,\mathsf{w}_2;\mathsf{w}_3)},
\]
where $\mathsf{w}_1$, $\mathsf{w}_2$, $\mathsf{w}_3$ are the weights of $\psi$, $\chi$, $\Psi$ in the norm (I.5.1).
\end{itemize}
\end{named}

\begin{named}{Remark}
For the Wilson fermion the corresponding objects and statements are the following:
the display of $T_k(B)$ in \cite[Section~17.1.2]{QEDI};
the definition of $T^{(1)}$, $T^{(2)}$, $\bar{T}^{(i)}$ and of the substitution in \cite[Section~17.4.4]{QEDI};
\cite[Lemma~17.9]{QEDI};
the paragraph on the substitution in the proof of \cite[Proposition~17.16(b)]{QEDI};
the four passages on the operators $T^{(i)}$, $\bar{T}^{(i)}$ and the substitution in the proofs of the clauses (a)--(d) of \cite[Lemma~17.21]{QEDI}, one passage in the proof of each clause;
and Step~2 of the proof in \cite[Section~17.12]{QEDI}.
Mathematically they differ as follows.
\begin{enumerate}
\item[(1)] The six hops with the constant projections $P^{\boldsymbol{\sigma}}_{\nu}$ become four bonds with the matrices $M_i(x)$, which depend on the label; only the bond of type~3 carries a phase and can leave $\mathbf{Y}$, and the internal differences ($i\le2$) are always written with $\psi$.
\item[(2)] The measure is $\int\mathrm{d}x$, and the prefactor of $T^{(2)}$ is $b_k^{-1}\cdot\eta_k^{\alpha-1}$ in front of $\int\mathrm{d}x$.
\item[(3)] The relations that enter are (I.m2)--(I.m4) and $\left(\tilde{g}x\right)^{(i)}=\tilde{g}\left(x^{(i)}\right)$; no sign $\boldsymbol{\sigma}$ is renamed.
\item[(4)] In the gauge covariance four kinds of terms occur: $\varphi_B(y,c)\cdot f(x)$; $\varphi_B(y,c)\cdot f\left(x^{(i)}\right)$ with $i\le2$, over the same photon site $c$; $\varphi_B(y,c)\cdot u_B(c,c_3)\cdot f\left(x^{(3)}\right)$; and $\varphi_B(y,c)\cdot F\left(x,x^{(3)}\right)$.
\item[(5)] In \cite[Section~17.4.4]{QEDI} the substitution is the homomorphism of the algebra $\mathrm{Gr}[\psi,\chi,\Psi]$ which maps, for every unit site $y$ of $\mathbf{Y}$, the generator $\Psi(y)$ to $\left(T^{(1)}\psi\right)(y)+\left(T^{(2)}\chi\right)(y)$ and the generator $\bar{\Psi}(y)$ to $\left(\bar{T}^{(1)}\bar{\psi}\right)(y)+\left(\bar{T}^{(2)}\bar{\chi}\right)(y)$, and which fixes all other generators;
a function of $(A,\mathscr{C},\mathsf{Z})$ with values in $\mathrm{Gr}[\psi,\chi,\Psi]$ is mapped at every point of its domain, the operators being taken at the connection $\mathscr{T}=A+\mathscr{C}$ of that point.
The homomorphism $\mathrm{sub}_{\mathbf{Y},B}$, with $B=\mathscr{T}$ taken at the point of evaluation, is the same homomorphism, formed with the operators (23.1), on the subalgebra $\mathrm{Gr}[\psi,\chi,\Psi\,|\,\mathbf{Y}]$;
there its values lie in $\mathrm{Gr}[\psi,\chi]$, as \cite[Section~17.4.4]{QEDI} notes for the elements that contain $\Psi$ and $\bar{\Psi}$ only at unit sites of $\mathbf{Y}$.
Apart from the operators, the two definitions thus differ in the domain only.
\item[(6)] $T^{(i)}$ is an operator on the whole unit lattice, given by one display at the unit sites of $Y$ and with zero rows outside $Y$.
\item[(7)] The rule (N2) of Section~I.5.2 is applied on $\mathrm{Gr}[\psi,\chi,\Psi\,|\,\mathbf{Y}]$.
\item[(8)] The preservation of parity is part of clause (d).
\item[(9)] At coinciding points the axial functional is $(\tau A)(c,c):=0$ (Section~I.1.1).
\end{enumerate}
\end{named}

\subsubsection*{Proposition 23.2 (level $0$: the vertex of the first step)}

\begin{named}{Definition}
Let $k=0$, so that $\eta_0=1$; let $\ell\in\mathrm{Cl}$, let $q>0$ be the charge of level $0$, let $\mathrm{Zf}$ be a complex bond function on $T_0$, and let $\square$ be an $M$-cube of $T_0$.
Put, in $\mathrm{Gr}[\psi,\chi]$,
\begin{equation}
\begin{aligned}
\mathrm{Vx}(\mathrm{Zf},\square):=\frac{1}{2}\int_{\square}\mathrm{d}x\,&\left\{\bar{\psi}(x)\cdot M_3(x)\left(\exp\left(i\,q\cdot\mathrm{Zf}(c,c_3)\right)-1\right)\left(\psi(x)+\chi\left(x,x^{(3)}\right)\right)\right.\\
&\left.{}+\left(\bar{\psi}(x)+\bar{\chi}\left(x,x^{(3)}\right)\right)\cdot M_3\left(x^{(3)}\right)\left(\exp\left(-i\,q\cdot\mathrm{Zf}(c,c_3)\right)-1\right)\psi(x)\right\},
\end{aligned}
\tag{23.3}
\end{equation}
where $\int_{\square}\mathrm{d}x$ is $\frac{1}{48}$ times the sum over the $48\,M^3$ fermion sites $x=(c,g)$ over $\square$, where $v\cdot m\,v':=\sum_{\beta,\beta'}v_{\beta}\,m_{\beta\beta'}\,v'_{\beta'}$, and where
\[
M_3(x)=3\left(\rho-\gamma\cdot(ge_1)\right),\qquad M_3\left(x^{(3)}\right)=3\left(\rho+\gamma\cdot(ge_1)\right).
\]
In this notation the display (I.5.5) takes the following form, which is denoted (8.3$_0$)$^{\mathrm{ch}}$:
\begin{equation}
G^{(1)}(s,\square)=\mathrm{Vx}\left(\mathscr{Z}(s'),\square\right).
\tag*{(8.3$_0$)$^{\mathrm{ch}}$}
\end{equation}
\end{named}

\begin{named}{Proposition 23.2}
With these data the following assertions hold.
\begin{itemize}
\item[(a)] \emph{Identity.}
Let $\mathscr{T}$ be a complex bond function on $T_0$, let $\mathrm{ff}$ and $\overline{\mathrm{ff}}$ be fields on $T^{\mathrm{f}}_0$, and let $\mathrm{ev}$ be the homomorphism of $\mathrm{Gr}[\psi,\chi]$ with
\[
\psi\mapsto\mathrm{ff},\qquad\bar{\psi}\mapsto\overline{\mathrm{ff}},\qquad\chi\mapsto\delta_{\alpha,\mathscr{T}}\,\mathrm{ff},\qquad\bar{\chi}\mapsto\delta_{\alpha,-\mathscr{T}}\,\overline{\mathrm{ff}}.
\]
Then the two halves of (23.3), each summed over the $M$-cubes, have the same image under $\mathrm{ev}$, and
\[
\mathrm{ev}\left(\sum_{\square}\mathrm{Vx}(\mathrm{Zf},\square)\right)=-\left\langle\overline{\mathrm{ff}},\left(D_{\mathscr{T}+\mathrm{Zf}}-D_{\mathscr{T}}\right)\mathrm{ff}\right\rangle^{\mathrm{ch}}.
\]
\item[(b)] \emph{Four symmetries.}
\begin{gather*}
a_{\tilde{g}}\left[\mathrm{Vx}\left(\tilde{g}\,\mathrm{Zf},\tilde{g}\square\right)\right]=\mathrm{Vx}(\mathrm{Zf},\square)\quad\text{for every }\tilde{g}\in\mathbb{G}_0,\\
G_{\Lambda}\left[\mathrm{Vx}(\mathrm{Zf},\square)\right]=\mathrm{Vx}(\mathrm{Zf},\square)\quad\text{for every real function }\Lambda\text{ of }c,\\
\mathrm{Cg}\left[\mathrm{Vx}(-\mathrm{Zf},\square)\right]=\mathrm{Vx}(\mathrm{Zf},\square),\qquad\mathbf{K}\left[\mathrm{Vx}\left(-\overline{\mathrm{Zf}},\square\right)\right]=\mathrm{Vx}(\mathrm{Zf},\square).
\end{gather*}
Under $\mathrm{Cg}$ the first half of $\mathrm{Vx}(-\mathrm{Zf},\square)$ goes to the second half of $\mathrm{Vx}(\mathrm{Zf},\square)$, and conversely; each of the three other maps preserves each half.
\item[(c)] \emph{Locality.}
$\mathrm{Vx}(\mathrm{Zf},\square)$ contains no generators other than $\psi$, $\bar{\psi}$ at the sites over $\square$ and $\chi$, $\bar{\chi}$ at the pairs $\left(x,x^{(3)}\right)$ with $c\in\square$.
Such a pair is admissible for $\square$: it satisfies the condition of kind (II), and also that of kind (I) if $c_3$ lies in $\square$.
The element reads $\mathrm{Zf}$ on the bonds $(c,c_3)$ with $c\in\square$, which are the oriented bonds whose initial point lies in $\square$;
together with their reverses $(c_3,c)$, on which $\mathrm{Zf}(c_3,c)=-\mathrm{Zf}(c,c_3)$, they are the bonds that meet $\square$.
It contains neither a connection $A$ or $\mathscr{C}$ nor the mass.
Every monomial has one barred and one unbarred generator.
The coefficients are entire functions of the numbers $\mathrm{Zf}(b)$.
\item[(d)] \emph{Norm.}
For all weights $\mathsf{w}=(\mathsf{w}_1,\mathsf{w}_2)$,
\begin{multline*}
\left\|\mathrm{Vx}(\mathrm{Zf},\square)\right\|_{\mathsf{w}}\le12\,(1+\rho)\cdot M^3\cdot\mathsf{w}_1\cdot\left(\mathsf{w}_1+\mathsf{w}_2\right)\\
\cdot\sup\left\{\left|\exp\left(i\,q\cdot\mathrm{Zf}(b)\right)-1\right|:b\text{ an oriented bond that meets }\square\right\}.
\end{multline*}
If $q\le1$ and $q\cdot\left|\mathrm{Zf}(b)\right|\le\frac{1}{2}\,q^{1-\frac{\varepsilon}{2}}$ on these bonds, then
\[
\left\|\mathrm{Vx}(\mathrm{Zf},\square)\right\|_{h_0/2}\le6\,(1+\rho)\cdot M^3\cdot q^{\frac{1}{2}-\frac{\varepsilon}{2}}\le12\,M^3\cdot q^{\frac{1}{2}-\frac{\varepsilon}{2}},\qquad h_0/2=\left(\tfrac{1}{2}\,q^{-\frac{1}{4}},\ \tfrac{1}{2}\,q^{-\frac{1}{4}+\varepsilon}\right).
\]
\item[(e)] \emph{In the step map.}
In the setting (ST1)$^{\mathrm{ch}}$ of Section~G3.4.0 at $k=0$, the following holds.
For $s\in\mathrm{Pd}_0$ and $\mathsf{Z}\in\mathrm{Box}_0$ the field $\mathrm{Zf}=\mathscr{Z}(s')$ satisfies the condition of (d) (Proposition~24.3\,(q1));
the function $G^{(1)}(\cdot,\square)$ contains $s$ only through $\mathrm{Zf}=\mathscr{Z}(s')$, so that it does not depend on $s''$;
it is continuous in $(s',\mathsf{Z})$ on $\mathrm{Pd}_0\times\mathrm{Box}_0$ and analytic in $s'$ in the interior; and
\[
\left\|G^{(1)}(s,\square)\right\|_{h_0/2}\le12\,M^3\cdot q^{\frac{1}{2}-\frac{\varepsilon}{2}}\le q^{\frac{1}{4}-5\varepsilon}.
\]
\end{itemize}
\end{named}

\begin{named}{Remark}
For the Wilson fermion the corresponding objects and statements are the following:
the display that defines $G^{(1)}(s,\square)$ at $k=0$ in \cite[Section~17.4.4]{QEDI}, which \cite[Definition~5.14]{QEDI} takes as the localised function of the first step;
the two sentences on level $0$ with which Step~(5) of the proof of \cite[Lemma~17.14]{QEDI} ends;
the estimate of $G^{(1)}$ for $k=0$ in the proof of \cite[Proposition~17.16(b)]{QEDI};
the passages on level $0$ in the proofs of the four clauses (a)--(d) of \cite[Lemma~17.21]{QEDI};
and Step~1 of the proof in \cite[Section~17.12]{QEDI} for $k=0$.
Mathematically they differ as follows.
\begin{enumerate}
\item[(1)] The sum over $(\nu,\boldsymbol{\sigma})$ at a site becomes the sum over the $48$ sites of the cluster ($8$ for each of the six directions $ge_1$), with the weight $\frac{1}{48}$; $P^{\boldsymbol{\sigma}}_{\nu}$ becomes $M_3(x)$, and $P^{-\boldsymbol{\sigma}}_{\nu}$ becomes $M_3\left(x^{(3)}\right)$; the bonds of the types $0$, $1$, $2$ carry no phase and do not occur.
\item[(2)] The count of $6\,M^3$ terms, the entries of $P^{\boldsymbol{\sigma}}_{\nu}$ adding up to at most $4$, becomes a count of $48\,M^3$ terms with the weight $\frac{1}{48}$, the moduli of the entries of $M_3$ adding up to at most $12\,(1+\rho)$ (I.m7); equality holds only for $ge_1=\pm e_1,\pm e_2$.
The number $12\,M^3$ and the second part of (A3) do not change.
\item[(3)] The statement is made for an arbitrary bond function $\mathrm{Zf}$, so that (a)--(d) contain no primitive.
\item[(4)] The charge $e_k$ is written $q$.
\item[(5)] \cite[Definition~5.26]{QEDI} fixes the sign of $G^{(1)}$ for $k\ge1$ by its display, and for $k=0$ it takes for $G^{(1)}$ the function of the single $M$-cubes of \cite[Section~17.4.4]{QEDI};
the sum of this function over the $M$-cubes, taken at $\chi=\delta_{\alpha,\mathscr{T}}\psi$ and $\bar{\chi}=\delta_{\alpha,-\mathscr{T}}\bar{\psi}$, is $-\left\langle\bar{\psi},\left(D_{\mathscr{T}+\mathscr{Z}(s')}-D_{\mathscr{T}}\right)\psi\right\rangle$, written as half the sum of two re-indexed forms of this pairing, so that its overall sign is $+$.
For the chamber fermion this function is (8.3$_0$)$^{\mathrm{ch}}$, with the overall sign $+$ of (a).
\end{enumerate}
\end{named}

\subsection*{22 (continued). The one-step formula for the density with external connection, all levels}

\subsubsection*{Proposition 22.4 (the exponent of the first step)}

\begin{named}{Proposition 22.4}
\begin{itemize}
\item[(a)] In the situation of Proposition~22.3 at $k=0$ one has $V^{\mathrm{ex},x}_0=\mathrm{Ev}^{(1)}+\mathrm{Ev}^{(5)}$, where
\[
\mathrm{Ev}^{(1)}:=-\left\langle\bar{\Psi}^f,\left(D_{\mathscr{T}+\mathscr{Z}}-D_{\mathscr{T}}\right)\Psi^f\right\rangle^{\mathrm{ch}},\qquad\mathrm{Ev}^{(5)}:=-m\cdot\left[\left\langle\bar{\Psi}^f,\Psi^f\right\rangle^{\mathrm{ch}}-\left\langle\bar{\Psi}^c,\Psi^c\right\rangle^{\mathrm{ch}}\right].
\]
\item[(b)] Let $\ell\in\mathrm{Cl}$ and $q>0$, let $\mathscr{T}$ and $\mathrm{Zf}$ be complex bond functions on $T_0$, and let $\mathrm{ff}$, $\overline{\mathrm{ff}}$ be fields on $T^{\mathrm{f}}_0$ (Section~G3.1.0).
Put $\mathrm{chf}:=\delta_{\alpha,\mathscr{T}}\,\mathrm{ff}$ and $\mathrm{chfb}:=\delta_{\alpha,-\mathscr{T}}\,\overline{\mathrm{ff}}$.
Then only the bonds of type~3 contribute, and
\begin{equation}
\begin{aligned}
&{-}\left\langle\overline{\mathrm{ff}},\left(D_{\mathscr{T}+\mathrm{Zf}}-D_{\mathscr{T}}\right)\mathrm{ff}\right\rangle^{\mathrm{ch}}\\
&\qquad=\int\mathrm{d}x\,\overline{\mathrm{ff}}(x)\cdot M_3(x)\left(\exp\left(i\,q\cdot\mathrm{Zf}(c,c_3)\right)-1\right)u_{\mathscr{T}}(c,c_3)\cdot\mathrm{ff}\left(x^{(3)}\right)\\
&\qquad=\frac{1}{2}\int\mathrm{d}x\,\left\{\overline{\mathrm{ff}}(x)\cdot M_3(x)\left(\exp\left(i\,q\cdot\mathrm{Zf}(c,c_3)\right)-1\right)\left(\mathrm{ff}(x)+\mathrm{chf}\left(x,x^{(3)}\right)\right)\right.\\
&\qquad\qquad\left.{}+\left(\overline{\mathrm{ff}}(x)+\mathrm{chfb}\left(x,x^{(3)}\right)\right)\cdot M_3\left(x^{(3)}\right)\left(\exp\left(-i\,q\cdot\mathrm{Zf}(c,c_3)\right)-1\right)\mathrm{ff}(x)\right\},
\end{aligned}
\tag{22.1}
\end{equation}
where the integrals extend over a fundamental domain and $v\cdot m\,v':=\sum_{\beta,\beta'}v_{\beta}\,m_{\beta\beta'}\,v'_{\beta'}$.
The last expression is the image of $\sum_{\square}\mathrm{Vx}(\mathrm{Zf},\square)$ (Proposition~23.2) under $\psi\mapsto\mathrm{ff}$, $\bar{\psi}\mapsto\overline{\mathrm{ff}}$, $\chi\mapsto\mathrm{chf}$, $\bar{\chi}\mapsto\mathrm{chfb}$.
\end{itemize}
\end{named}

\begin{named}{Remark}
For the Wilson fermion the identity (22.1) at level $0$ is obtained in Step~1 of the proof in \cite[Section~17.12]{QEDI}, for the exponent $V^{\mathrm{ex},x}_0$ of \cite[Theorem~14.24]{QEDI}.
There six hops per site with the matrices $P^{\boldsymbol{\sigma}}_{\nu}$ occur; here there is one hop of type~3 per fermion site, with $M_3(x)$, respectively $M_3\left(x^{(3)}\right)$, and with the weight $\frac{1}{48}$.
The re-indexing $(x,\nu,\boldsymbol{\sigma})\mapsto\left(x^{\boldsymbol{\sigma}},\nu,-\boldsymbol{\sigma}\right)$ used there becomes $x\mapsto x^{(3)}$.
\end{named}

\subsection*{24. One renormalisation step without large fields, with external connection, every level $k\ge0$: setting and definitions}

Under the heading 24 and its continuation, and under the heading 25, the labels (8.1)$^{\mathrm{ch}}$, (8.2), (8.3)$^{\mathrm{ch}}$, (8.4) and (8.6)$^{\mathrm{ch}}$ denote the displays of Definition~24.D (Section~G3.4.1) in (D2), (D3), (D4), (D5) and (D6), which define the fields, the localisation $\mathrm{Loc}$, the localised functions, the fluctuation integral and the step.
The labels (8.3a)$^{\mathrm{ch}}$ and (8.3$_0$)$^{\mathrm{ch}}$ denote the substitution rule of Proposition~23.1 and the form of (I.5.5) given before Proposition~23.2.
These seven labels are also written with the prefix G3, as (G3.8.1)$^{\mathrm{ch}}$, (G3.8.2), (G3.8.3)$^{\mathrm{ch}}$, (G3.8.4), (G3.8.6)$^{\mathrm{ch}}$, (G3.8.3a)$^{\mathrm{ch}}$ and (G3.8.3$_0$)$^{\mathrm{ch}}$; both forms denote the same displays.

\subsubsection*{G3.4.0. Setting (ST1)$^{\mathrm{ch}}$}

\paragraph{Constants.}
The constants are those of Section~I.8, lines (C0)--(C8), chosen in this order, each of them depending on the earlier lines only.
Those which occur in the statements of this section are $\rho$, $w$; $c_X$, $\kappa_0$, $K_0$, $c_{\mathrm{fl}}$, $b$ (line (C0), which consists of $\rho$ and absolute numbers); $\alpha$ (C1); $\kappa$ (C2); $L$ (C3); $b_L$, $\alpha_0$, $\kappa_{\mathrm{w}}$, $C_T$ (C4); $\gamma_{\mathrm{w}}$ (C4r; it depends on $\rho$); $M$ (C5); $C_w$, $c_{\gamma}$, $C_W$ (C6); $r$, $p_0$, $p$, $\varepsilon$ (C7); $e_0$ (C8).
None of them depends on $(k,\ell,\mu,e_k)$, on the system, on $(m,\mathscr{E})$, or on $A$, $\mathscr{C}$, $A'$, $\mathscr{C}'$.

\paragraph{Data.}
An integer $k\ge0$; a side $\ell$ with $\ell\in\mathrm{Cl}$, $\ell/L\in\mathrm{Cl}$ and $\ell\ge M\cdot r_k$; a mass $\mu\in(0,1/L]$; a number $e_k>0$ with $e_{k+1}:=L^{\frac{1}{2}}\cdot e_k\le e_0$.
Here $e_0$ is the ceiling (C8); at $k=0$ the symbol $e_k$ is the charge of level $0$, which is written $q$ in the flow.
The quantities $v_k=-\log e_k$, $r_k$, $p_k$, $p_{0,k}$, $w^{\mathrm{m}}_k$, $w^{\mathrm{E}}_k$, $h_k=(h_{k,1},h_{k,2})$, $\mathrm{rad}_1$, $\mathrm{rad}_2$, $\mathrm{rad}_4$ (Sections~I.5.1, I.5.4, I.5.7 and I.6.1) are formed with $e_k$.
The level $k+1$ has the side $\ell/L$, the mass $L\cdot\mu$ and the charge $e_{k+1}$.
No other relation between $k$, $\ell$, $\mu$, $e_k$ is assumed.

\paragraph{The system.}
On $T_k$ a system of primitives of level $k$ in the sense of Section~I.5.8 is given: the families (p1)--(p5) for $k\ge1$; for $k=0$ the family (p2), $\Gamma_0(s'',A)$, $S_{0,y}(X,A)$ and the two diagonal families $H_0(s')$, $H_0(s'',A)$.
It has the properties W1--W8 and W10--W17 of systems of primitives of \cite[Section~5.4.2]{QEDI}, which are numbered 1--17 there, for $\mathrm{par}=\left(\alpha_0,\gamma_{\mathrm{w}},\kappa_{\mathrm{w}};C_w\right)$, read as in Section~I.5.8.
At the places where a fermion index occurs, the exact forms of these properties are those of (G3.T-W) in Section~G3.0.4.
The property W9 is not required, not even at $k=0$: there the property W17 is taken without W9, that is, with W8, W10, W11 in place of its W8--W11 and with W6--W8 in place of its W6--W9.

\paragraph{Numbers.}
\begin{gather*}
\mathsf{w}_W=\max\nu^{-\frac{1}{2}}\quad\left(\text{(I.5.2), Proposition~15\,(D3)}\right),\qquad\text{that is,}\\
\mathsf{w}_W=1\ \text{for }k\ge1,\qquad\mathsf{w}_W=\sqrt{48}\ \text{for }k=0;
\end{gather*}
the bound $\sqrt{48}\le7$ is the factor $7$ of $C_W$ (I.L0.19).
Further,
\[
C_W=28\,C_w^2\cdot c_{\gamma},\qquad\mathbf{H}_0=36\,c_X^2\cdot K_0^2\cdot e_k^{\frac{1}{4}-5\varepsilon}.
\]

The setting (S22) is (ST1)$^{\mathrm{ch}}$ without the condition $\ell\ge M\cdot r_k$ and without the system.

\subsubsection*{G3.4.1. Definition 24.D (the step map of level $k$)}

Definition~24.D is the counterpart, in the notation of Part~I, of \cite[Definition~5.25]{QEDI} and of the definition of the step map in \cite[Section~17.4]{QEDI}.

\paragraph{(D0) Geometry.}
All sets in (D0) are sets of photon sites; the notions are those of \cite[Section~5.2.3]{QEDI}, \cite[Definition~5.8]{QEDI}, \cite[Section~17.4.1]{QEDI} (the facts (g0)--(g4)) and \cite[Section~17.4.3]{QEDI}.
A parameter cube is contained in a set $Z$ of sites if all its lattice sites lie in $Z$; it meets an $M$-cube if the two have a lattice site in common.
For a polymer $X$, $P_0(X)$ denotes the set of the parameter cubes, of the kinds that occur, which are contained in $X$.
For a set $P$ of parameter cubes of these kinds outside $P_0(X)$, $\mathrm{Yp}(X,P)$ denotes the union of the cubes of $P_0(X)\cup P$, and $Y'(X,P)$ the union of the $M$-cubes that meet $X\cup\mathrm{Yp}(X,P)$.
A component of a union of $M$-cubes is a maximal subset that is connected through faces.

\paragraph{(D1) Data, background, variables, parameters.}
\emph{Data:} $(m,\mathscr{E})\in\mathrm{Ball}^{x}_{k}(4)$, where $\mathscr{E}=0$ and $|m|\le4\,w^{\mathrm{m}}_0$ for $k=0$; and
\[
m':=m+m(\mathscr{E}),\qquad\mathscr{E}^{r}:=R\,\mathscr{E}\qquad\left(\text{for }k=0:\ m'=m,\ \mathscr{E}^{r}=0\right).
\]
\emph{Background:} $A\in\tilde{R}_k^{-}$ in the first slot, $\mathscr{C}\in\mathfrak{X}_k$ in the second slot, and $\mathscr{T}:=A+\mathscr{C}$.
\emph{Fluctuation variables:} $\mathsf{Z}\in\mathrm{Box}_k$, and the generators $W$ of Section~I.3.6 (for $k\ge1$ on $T^0_k$, with the weight $1$; for $k=0$ on $T^{\mathrm{f}}_0$, with the weight $\frac{1}{48}$).
\emph{Parameters:} $s=(s',s'')$.
Here $s'=\left(s'_{\mathrm{pc}}\right)$ with $\mathrm{pc}\in\pi'_k\cup\pi''_k$; $H_k(s')$ contains the parameters of $\pi'_k$, and $C_{k}^{1/2}(s')$ those of $\pi''_k$.
Further $s''=\left(s''_{\mathrm{pc}}\right)$ with $\mathrm{pc}\in\pi_k$; the same variables stand in $H_k(s'',\cdot)$, $\bar{H}_k(s'',\cdot)$ and $\Gamma_k(s'',\cdot)$.
The parameters of different cubes and of different kinds are independent complex variables, each in the disk $\left|s_{\mathrm{pc}}\right|\le M^{\alpha_0}$; thus $s\in\mathrm{Pd}_k$.
For a union $U$ of $M$-cubes, $\mathbf{1}_U$ denotes the family with the value $1$ at the cubes contained in $U$ and the value $0$ at the others.
For complex $s'$ the maps $H_k(s')$, $C$ and $C_{k}^{1/2}(s')$ act between the complexified spaces.

\emph{Rule (SLOT).}
Every fermionic operator, exact or weakened, is evaluated at its first-slot argument plus $\mathscr{C}$.
The activity enters as $\mathscr{E}^{r}(X;B_1,\mathscr{C},\cdot)$, where $B_1$ is the first-slot argument, $A$ or $A+t\cdot\mathscr{Z}(s')$.
The photon operators, the measures, the contours, $\mathrm{Loc}$ and the cluster expansion do not contain $\mathscr{C}$.

\emph{Adding and inserting fields.}
Let $v^{\#}$ be a family of odd elements indexed like $\psi^{\#}=(\psi,\chi)$, that is, by sites and by pairs, barred and unbarred.
It is added to the generators by the homomorphism $\psi(\xi)\mapsto\psi(\xi)+v(\xi)$, $\chi(\boldsymbol{\zeta})\mapsto\chi(\boldsymbol{\zeta})+v(\boldsymbol{\zeta})$, and the image of $F$ under this homomorphism is written $F\left(\psi^{\#}+v^{\#}\right)$.

\paragraph{(D2) Fields.}
Let $B\in\tilde{G}_k$, let $\mathrm{Zf}$ be a bond function with $B+\mathrm{Zf}\in\tilde{G}_k$, and let $(\Phi,\bar{\Phi})$ be a pair of fields on the unit index set of level $k$.
Put
\begin{equation}
\begin{aligned}
&\mathscr{Z}(s'):=H_k(s')\cdot C\cdot C_{k}^{1/2}(s')\cdot\mathsf{Z}\\
&\qquad\left(\text{for }k=0,\ H_0(s')\text{ is the diagonal family of Section~I.5.8}\right);\\
&H_k^{\#}(s'',B)\,\Phi:=\left(H_k(s'',B)\,\Phi,\ \delta_{\alpha,B}H_k(s'',B)\,\Phi;\ \bar{H}_k(s'',B)\,\bar{\Phi},\ \delta_{\alpha,-B}\bar{H}_k(s'',B)\,\bar{\Phi}\right),\\
&\qquad\bar{H}_k(s'',B):=\mathsf{C}^{-1}\cdot H_k(s'',-B)\cdot\mathsf{C};\\
&J_k^{\#}(s'',B,\mathrm{Zf}):=H_k^{\#}(s'',B+\mathrm{Zf})-H_k^{\#}(s'',B);\\
&\qquad J_k(s'',B,\mathrm{Zf}),\ \bar{J}_k(s'',B,\mathrm{Zf})\text{: its two components without H\"older quotient};\\
&\tilde{\Gamma}_k(s'',B)\,(W,\bar{W}):=\left(\Gamma_k(s'',B)\,W,\ \bar{W}\right),\\
&W_k^{\#}(s'',B):=H_k^{\#}(s'',B)\cdot\tilde{\Gamma}_k(s'',B)\cdot W;\\
&\qquad\mathrm{Wc},\ \overline{\mathrm{Wc}}\text{: its two components without H\"older quotient}.
\end{aligned}
\tag*{(8.1)$^{\mathrm{ch}}$}
\end{equation}
The H\"older quotient $\delta_{\alpha,B}$ is taken at all pairs of $\mathrm{Pz}_k$ (I.1.7).
At $k=0$, $H_0(s'',B)$ is the multiplication by $\prod\left\{s''_{\mathrm{pc}}:\mathrm{pc}\in\pi_0,\ \mathrm{pc}\text{ meets }K_c\right\}$, and
\[
\left(\Gamma_0(s'',B)\,W\right)(x)=\int\mathrm{d}x'\,\Gamma_0(s'',B)(x,x')\,W(x').
\]

\paragraph{(D3) Localisation.}
Let the kinds of the parameter cubes be fixed.
For every $X\in D_k$ let $G(\cdot,X)$ be a map from $\mathrm{Pd}_k$, the parameters being of these kinds, to one finite-dimensional complex vector space, continuous on $\mathrm{Pd}_k$ and analytic in its interior.
The cube $[0,1]^P$ lies in this interior, because $M^{\alpha_0}>1$.
For $\mathbf{Y}\in D_k$ put
\begin{equation}
\mathrm{Loc}[G](\mathbf{Y}):=\sum_{(X,P):\,Y'(X,P)=\mathbf{Y}}\ \int_{[0,1]^P}\mathrm{d}s_P\,\frac{\partial}{\partial s_P}\left[G(s,X)\right]_{s=1\text{ on }P_0(X),\ s=0\text{ outside }P_0(X)\cup P}.
\tag{8.2}
\end{equation}
Pairs $(X,P)$ for which $Y'(X,P)$ is not connected belong to no $\mathrm{Loc}[G](\mathbf{Y})$.
A function $G(\cdot,\square)$ that is defined for single $M$-cubes is set equal to $0$ on the other polymers.

\paragraph{(D4) The localised functions.}
$\oint_{\mathrm{rad}}$ is the functional of Section~I.5.2, taken on the circle $|t|=\mathrm{rad}$, counterclockwise.
For $k\ge1$ put
\begin{equation}
\begin{aligned}
G^{(1)}(s,X)&:=\oint_{\mathrm{rad}_1}\left\langle\bar{\Psi},M_k\left(X,\mathscr{T}+t\cdot\mathscr{Z}(s')\right)\Psi\right\rangle_1,\\
G^{(2)}(s,X)&:=\oint_{\mathrm{rad}_2}\mathscr{E}^{r}\left(X;A+t\cdot\mathscr{Z}(s'),\mathscr{C},\psi^{\#}+t\cdot J_k^{\#}\left(s'',\mathscr{T},\mathscr{Z}(s')\right)\Psi\right),\\
G^{(3)}(s,\square)&:=\oint_{\mathrm{rad}_2}\left(-m'\right)\left\langle\bar{\psi}+t\cdot\bar{J}_k\left(s'',\mathscr{T},\mathscr{Z}(s')\right)\bar{\Psi},\right.\\
&\qquad\qquad\left.\psi+t\cdot J_k\left(s'',\mathscr{T},\mathscr{Z}(s')\right)\Psi\right\rangle_{\square},\\
E^{(i),*}&:=\mathrm{Loc}\left[G^{(i)}\right]\quad\left(\text{kinds }\pi_k,\pi'_k,\pi''_k\right)\qquad(i=1,2,3),\\
E^{(i),x}(\mathbf{Y})&:=\mathrm{sub}_{\mathbf{Y},\mathscr{T}}\left(E^{(i),*}(\mathbf{Y})\right)\quad\left(\text{Proposition~23.1}\right)\qquad(i=1,2,3),\\
E^{(\le3),x}&:=E^{(1),x}+E^{(2),x}+E^{(3),x},\\
G^{(\le3)}(s'',X)&:=E^{(\le3),x}\left(X;A,\mathscr{C},\mathsf{Z},\psi^{\#}+W_k^{\#}(s'',\mathscr{T})\right),\\
G^{(4)}(s'',X)&:=\oint_{\mathrm{rad}_4}\mathscr{E}^{r}\left(X;A,\mathscr{C},\psi^{\#}+t\cdot W_k^{\#}(s'',\mathscr{T})\right),\\
G^{(5)}(s'',\square)&:=\oint_{\mathrm{rad}_4}\left(-m'\right)\left\langle\bar{\psi}+t\cdot\overline{\mathrm{Wc}},\ \psi+t\cdot\mathrm{Wc}\right\rangle_{\square},\\
E^{(\le5),x}&:=\mathrm{Loc}\left[G^{(\le3)}+G^{(4)}+G^{(5)}\right]\quad\left(\text{kind }\pi_k\right).
\end{aligned}
\tag*{(8.3)$^{\mathrm{ch}}$}
\end{equation}
For $k=0$ there is no field $\Psi$ and no substitution, $G^{(2)}=G^{(3)}=G^{(4)}=0$, and
\begin{gather*}
G^{(1)}(s,\square):=\mathrm{Vx}\left(\mathscr{Z}(s'),\square\right)\quad\left(\text{Proposition~23.2, (8.3$_0$)$^{\mathrm{ch}}$}\right),\\
E^{(1),*}=E^{(1),x}=E^{(\le3),x}:=\mathrm{Loc}\left[G^{(1)}\right]\quad\left(\text{kinds }\pi'_0,\pi''_0,\pi_0\right),\\
G^{(\le3)}(s'',X):=E^{(\le3),x}\left(X;\mathsf{Z},\psi^{\#}+W_0^{\#}(s'',\mathscr{T})\right),\\
G^{(5)}(s'',\square):=\oint_{\mathrm{rad}_4}\left(-m\right)\left\langle\bar{\psi}+t\cdot\overline{\mathrm{Wc}},\ \psi+t\cdot\mathrm{Wc}\right\rangle_{\square},\\
E^{(\le5),x}:=\mathrm{Loc}\left[G^{(\le3)}+G^{(5)}\right]\quad\left(\text{kind }\pi_0\right).
\end{gather*}

\paragraph{(D5) Fluctuation integral.}
$\mathrm{urs}^{T}$ is the Ursell coefficient of \cite[Proposition~5.58]{QEDI}, where it is denoted $\rho^T$.
For a union $X$ of $M$-cubes put, with $\mathrm{Km}(\emptyset):=1$,
\begin{equation}
\begin{aligned}
&\mathrm{Km}(X):=\sum_{\{X_i\}:\,\bigcup_iX_i=X}\ \prod_i\left(\exp\left(E^{(\le5),x}(X_i)\right)-1\right)\\
&\qquad\qquad\left(\{X_i\}\text{ a set of distinct polymers}\right),\\
&\mathrm{Hm}(X):=\int\mathrm{d}\mu_I(W)\,\mathrm{d}\hat{\mu}(\mathsf{Z})\,\mathrm{Km}(X),\\
&E^{\#,x}_k(m,\mathscr{E})(X;A,\mathscr{C},\psi,\chi)\\
&\qquad:=\sum_{n\ge1}\frac{1}{n!}\sum_{(X_1,\dots,X_n):\,\bigcup_iX_i=X}\mathrm{urs}^{T}(X_1,\dots,X_n)\prod_i\mathrm{Hm}(X_i)\qquad(X\in D_k).
\end{aligned}
\tag{8.4}
\end{equation}

\paragraph{(D6) The step.}
Let $Y\in D_{k+1}$ and $(\mathscr{A},\mathscr{C})\in\tilde{R}_{k+1}\times\mathfrak{X}_{k+1}$ (side $\ell/L$), let $\varepsilon^{\mathrm{v}}$ be a complex number, and let $\mathfrak{P}_k=\left\{\mathfrak{P}_k(X,\cdot):X\in D_k\right\}$ be an arbitrary family of complex functions on $\mathfrak{X}_k$ (Section~I.5.7, Section~G3.2.0).
With the operations $L$, $R$, $(\cdot)^{\mathrm{nat}}$ of Section~I.5.6 and the functions $E^{d}_k$, $E_k^{\mathrm{det}}$, ${\delta E^{\mathrm{det}}}_k$ of Section~I.5.7, formed with the family (p5), put
\begin{equation}
\begin{aligned}
E^{*,x}_k(m,\mathscr{E})&:=L\,E^{\#,x}_k(m,\mathscr{E}),\quad\text{that is,}\\
E^{*,x}_k(m,\mathscr{E})(Y;\mathscr{A},\mathscr{C},\psi,\chi)&=\sum_{X:\,\bar{X}^{L}=LY}E^{\#,x}_k(m,\mathscr{E})\left(X;\mathscr{A}_L,\mathscr{C}_L,\psi_L,\chi_L\right),\\
\mathscr{E}^{+}&:=L\,R\,\mathscr{E}+L\,{\delta E^{\mathrm{det}}}_k+E^{*,x}_k(m,\mathscr{E}),\\
\varepsilon^{\mathrm{v},+}&:=L^3\cdot\left(\varepsilon^{\mathrm{v}}+\varepsilon(\mathscr{E})+\delta\varepsilon^{0}_{k}\right),\qquad m^{+}:=L\cdot\left(m+m(\mathscr{E})\right),\\
\mathfrak{P}_{k+1}(Y,\mathscr{C})&:=\sum_{X:\,\bar{X}^{L}=LY}\left(\mathfrak{P}_k(X,\mathscr{C}_L)+E_k^{\mathrm{det}}(X,\mathscr{C}_L)\right).
\end{aligned}
\tag*{(8.6)$^{\mathrm{ch}}$}
\end{equation}
For $k=0$ this reads $\mathscr{E}^{+}=L\,{\delta E^{\mathrm{det}}}_0+E^{*,x}_0(m,0)$, $\varepsilon^{\mathrm{v},+}=L^3\cdot\left(\varepsilon^{\mathrm{v}}+\delta\varepsilon^{0}_{0}\right)$ and $m^{+}=L\cdot m$.
For $(\mathscr{A},\mathscr{C})\in\tilde{R}_{k+1}\times\mathfrak{X}_{k+1}$ the bond function $\mathscr{A}_L$ lies in $\tilde{R}_k^{-}$ and $\mathscr{C}_L$ lies in $\mathfrak{X}_k$ \cite[Lemma~5.18]{QEDI}.
Moreover $\mathscr{A}_L+\mathscr{C}_L$ and $\mathscr{C}_L$ lie in $\tilde{G}_k$, since $0$ lies in $\tilde{R}_k$ and $\tilde{R}_k+\mathfrak{X}_k$ is contained in $\tilde{G}_k$ \cite[Lemma~5.17]{QEDI}.
Hence $E^{d}_k$, $E_k^{\mathrm{det}}$ and ${\delta E^{\mathrm{det}}}_k$ are evaluated inside their domain.

\paragraph{(D7) Reading (o) (no second slot).}
The step map without slot is defined by (D1)--(D6), read as follows.
The second slot is suppressed: $\mathscr{C}$ is not a variable, $\mathscr{T}=A$, and $\boldsymbol{\lambda}'=0$.
Every product $\mathrm{Dm}\times\mathfrak{X}_j$ with $\mathrm{Dm}=\tilde{R}_k,\tilde{R}_k^{-},\tilde{R}_{k+1}$ is read $\mathrm{Dm}$.
The symbols $K^{x,C}_{j}$, $K^{x}_{j}$, $\|\cdot\|^{x}_{j}$, $\mathrm{Ball}^{x}_{k}$ are read $K^{C}_{j}$, $K_j$, $\|\cdot\|_j$, $\mathrm{Ball}_k$; ${\delta E^{\mathrm{det}}}_k$ is read $E_k^{\mathrm{det}}$; and $\varepsilon$, $m$, $R$, $(\cdot)^{\mathrm{nat}}$, $L$ are the operations of Section~I.5.6 on families without slot.
Every assertion of continuity or analyticity in $(A,\mathscr{C})$ is the same assertion in $A$ alone, on the open set $\mathrm{Dm}$.
Locality in $\mathscr{C}$ and comparisons of two values of $\mathscr{C}$ are void; so are the family $\mathfrak{P}_k$ and the entry $\mathfrak{P}_{k+1}$ of (8.6)$^{\mathrm{ch}}$, which are functions of the suppressed variable alone.
The objects of this reading are written $E^{(i)}$, $E^{\#}_k(m,E)$, $E^{*}_k(m,E)$.
The reading (o) is the case in which the pair $(\{0\},\{0\})$ stands in the place of $\left(\mathfrak{X}_k,\mathfrak{X}_{k+1}\right)$, a product $\mathrm{Dm}\times\{0\}$ being read $\mathrm{Dm}$; it is the counterpart of the case without slot of \cite[Section~17.1.3]{QEDI}.
That section admits in the place of $\left(\mathfrak{X}_k,\mathfrak{X}_{k+1}\right)$ every pair of sets of complex bond functions on $T_k$ and on $T_{k+1}$ which are open in the space of all complex bond functions and have the five properties (a)--(e) listed there;
the pair $(\{0\},\{0\})$, which is not open, is treated there separately, by a reading that belongs to the definition of the case without slot and takes the place of openness, and this reading agrees with (D7).
Definition~24.D, on the other hand, forms the step map with no pair other than $\left(\mathfrak{X}_k,\mathfrak{X}_{k+1}\right)$ and $(\{0\},\{0\})$.

\paragraph{(D8) The objects attached to $\mathbf{Y}$.}
Let $\mathbf{Y}\in D_k$.
The objects attached to $\mathbf{Y}$ are
\[
E^{(i),*}(\mathbf{Y}),\quad E^{(i),x}(\mathbf{Y})\ \ (i\le3),\quad E^{(\le3),x}(\mathbf{Y}),\quad E^{(\le5),x}(\mathbf{Y}),\quad \mathrm{Km}(\mathbf{Y}),\quad \mathrm{Hm}(\mathbf{Y}),\quad E^{\#,x}_k(m,\mathscr{E})(\mathbf{Y}).
\]
A term $(X,P)$ of (G3.8.2) is said to be \emph{attributed to} $\mathbf{Y}$ if $Y'(X,P)=\mathbf{Y}$.

\paragraph{Order in which the displays acquire a meaning.}
The displays of Definition~24.D become meaningful in the order listed below, and each step of the list is a clause of Proposition~24.1\,(a) or of Proposition~24.3:
(G3.8.1)$^{\mathrm{ch}}$;
Proposition~24.3\,(q1), (q5);
$G^{(1)}$, $G^{(2)}$, $G^{(3)}$;
the first localisation;
(G3.8.3a)$^{\mathrm{ch}}$;
$G^{(\le3)}$, $G^{(4)}$, $G^{(5)}$;
the second localisation;
(G3.8.4);
(G3.8.6)$^{\mathrm{ch}}$.

\paragraph{Order of proof in the Propositions~23.m and~24.m.}
The clauses of the Propositions~23.m and~24.m are proved in the following order.
Among these clauses, the proof of each step uses only clauses of the steps before it.
\begin{itemize}
\item[(1)] Proposition~23.1\,(a)--(g) and Proposition~23.2\,(a)--(d); their proofs rest on no clause of the Propositions~24.m.
\item[(2)] Proposition~24.3\,(adm), (q1), from the properties W2--W4 of the system of primitives (at level $0$: W17), from \cite[Proposition~5.35]{QEDI}, \cite[Lemmas~17.3, 17.4(a), 17.6]{QEDI}, and from the conditions (A1), (A2), (A5) of Section~I.8.4.
Here W$n$ is the property numbered $n$ in \cite[Section~5.4.2]{QEDI}, read as in Section~I.5.8.
No clause of the Propositions~23.m enters this step.
\item[(3)] Proposition~23.2\,(e), from Proposition~23.2\,(c), (d), from Proposition~24.3\,(q1), from the property W2 of the family (p2), read as in Section~I.5.8 (regularity in $s'$), and from the second part of (A3).
\item[(4)] Proposition~24.3\,(q5), (al), (be), (ga), (de), (ze), (et).
At this step the following are established only for the objects $E^{(i),*}(\mathbf{Y})$, $E^{(i),x}(\mathbf{Y})$, $E^{(\le3),x}(\mathbf{Y})$, $E^{(\le5),x}(\mathbf{Y})$, $\mathrm{Km}(\mathbf{Y})$, $\mathrm{Hm}(\mathbf{Y})$, which are finite sums and integrals, and for the terms of the series (G3.8.4):
the clause (ze);
the assertion of (al) on two data, namely that the objects attached to $\mathbf{Y}$ that are formed with two data $(m_1,\mathscr{E}_1),(m_2,\mathscr{E}_2)\in\mathrm{Ball}^{x}_{k}(4)$ coincide if $m_1+m(\mathscr{E}_1)=m_2+m(\mathscr{E}_2)$ and $(R\mathscr{E}_1)(X)=(R\mathscr{E}_2)(X)$ for all polymers $X\subseteq\mathbf{Y}$;
and the last sentence of (ga4).
\item[(5)] The first assertion of Proposition~24.1\,(a), namely that every display of Definition~24.D is defined, that the series (G3.8.4) converge absolutely in $\|\cdot\|_{h_k/4}$, and that the bound displayed there for $E^{\#,x}_k(m,\mathscr{E})(X;A,\mathscr{C})$, $X\in D_k$, holds;
and the first sentence of Proposition~24.1\,(b), which is the same assertion in the reading (o).
From this step on, the convergence of the series (G3.8.4) is available.
\item[(6)] Proposition~24.3\,(ze), the assertion of (al) on two data $(m_1,\mathscr{E}_1)$, $(m_2,\mathscr{E}_2)$ and the last sentence of (ga4), for $E^{\#,x}_k(m,\mathscr{E})(\mathbf{Y})$;
the closing sentences of Proposition~24.3, on the two other parts $(R\mathscr{E})^{\mathrm{nat}}(X)$ and ${\delta E^{\mathrm{det}}}_k(X)$ of $\mathscr{E}^{+}(Y)$;
their proof uses only the display (G3.8.6)$^{\mathrm{ch}}$ of Definition~24.D, which gives $\mathscr{E}^{+}(Y)$, Sections~I.5.6 and~I.5.7 and the property W12, and no clause of the Propositions~23.m and~24.m enters it;
Proposition~24.2.
\item[(7)] Proposition~24.1\,(f), (g);
then the remaining sentences of (a);
then (b) from its second sentence on, the first clause of (d) (that $\mathscr{E}^{+}$ lies in $K^{x}_{k+1}$, which follows from Proposition~18.5\,(a), Proposition~13\,(d5) and Proposition~24.1\,(a)), (c), the rest of (d), (e) and (i).
\end{itemize}

\subsubsection*{G3.4.2. Definition 24.I (the four transformations on arguments; invariance relative to a polymer and to a set of arguments)}

\begin{named}{Definition 24.I}
The objects of Definition~24.D are elements of a Grassmann algebra which depend on arguments
\[
\mathrm{ar}=(X,P,s,t;A,\mathscr{C},\mathsf{Z}).
\]
On the generators $\psi$, $\chi$, $\Psi$, $W$ the four transformations act by the maps of Section~I.4.5, with a single exception: $\mathrm{Cg}$ is the map of Section~I.4.5 on $\psi$, $\chi$, $\Psi$ only, and it acts as the identity on $W$ (see (T-C) below).
On the arguments they act as follows.
\begin{itemize}
\item[(T-g)] For $\tilde{g}\in\mathbb{G}_k^{L}$,
\[
\mathrm{ar}^{\theta}:=\left(\tilde{g}X,\tilde{g}P,\tilde{g}s,t;\tilde{g}A,\tilde{g}\mathscr{C},\tilde{g}\mathsf{Z}\right),\qquad \theta:=a_{\tilde{g}}.
\]
\item[(T-G)] Let $\boldsymbol{\lambda}$, $\boldsymbol{\lambda}'$ be real functions of the site $c$, put $\Lambda:=\boldsymbol{\lambda}+\boldsymbol{\lambda}'$, and let a polymer $\mathbf{Y}$ be given; this transformation is taken relative to $\mathbf{Y}$.
Then $\mathrm{ar}^{\theta}$ is any argument $(X,P,s,t;A^1,\mathscr{C}^1,\mathsf{Z})$ such that
\[
A^1=A-d\boldsymbol{\lambda},\qquad \mathscr{C}^1=\mathscr{C}-d\boldsymbol{\lambda}'
\]
on the bonds at distance at most $1$ from $\mathbf{Y}$, that is, on the bonds both of whose ends have sup-distance at most $1$ from $\mathbf{Y}$;
and $\theta:=G_{\Lambda}$.
\item[(T-C)] $\mathrm{ar}^{\theta}:=(X,P,s,t;-A,-\mathscr{C},-\mathsf{Z})$, and $\theta$ is $\mathrm{Cg}$ on $\psi$, $\chi$, $\Psi$, extended by the identity on $W$.
\item[(T-K)] $\mathrm{ar}^{\theta}:=\left(X,P,\overline{s},\overline{t};-\overline{A},-\overline{\mathscr{C}},-\mathsf{Z}\right)$, and $\theta:=\mathbf{K}$.
\end{itemize}
Let $\mathrm{Ar}$ be a set of arguments which contains, together with each $\mathrm{ar}$, every admitted transform $\mathrm{ar}^{\theta}$.
An object $F$ is \emph{invariant on} $\mathrm{Ar}$ if
\[
\theta\left(F\left(\mathrm{ar}^{\theta}\right)\right)=F(\mathrm{ar})\qquad\text{for all }\mathrm{ar}\in\mathrm{Ar}.
\]
A family $v$ of odd elements standing in the place of generators is \emph{covariant on} $\mathrm{Ar}$ (at an index) if $\theta\left(v\left(\mathrm{ar}^{\theta}\right)(\cdot)\right)$ is the combination of the elements $v(\mathrm{ar})(\cdot)$ which $\theta$ assigns to the generators in whose place $v$ stands.
If $v$ is covariant, then so is $t\cdot v$; for (T-K) this uses $t^{\theta}=\overline{t}$ and the antilinearity of $\theta$.

For (T-g), (T-C) and (T-K), $\mathrm{Ar}$ is the set of all arguments in the domain of $F$.
For (T-G), the polymer $\mathbf{Y}$ fixes the admitted transforms, and nothing else: $\mathrm{ar}^{\theta}$ is \emph{admitted relative to} $\mathbf{Y}$ if it has the form given under (T-G) and $(A^1,\mathscr{C}^1)$ lies in the domain of $F$.
The set $\mathrm{Ar}$ is then specified separately: $F$ is \emph{invariant on $\mathrm{Ar}$ relative to $\mathbf{Y}$}, and $v$ is \emph{covariant on $\mathrm{Ar}$ relative to $\mathbf{Y}$}, if $\mathrm{Ar}$ contains, together with each $\mathrm{ar}$, every $\mathrm{ar}^{\theta}$ admitted relative to $\mathbf{Y}$, and if the relation of the preceding paragraph holds for all $\mathrm{ar}\in\mathrm{Ar}$ and all these $\mathrm{ar}^{\theta}$.
The sets that occur are, for $\mathbf{Y}\in D_k$,
\[
\mathrm{Ar}(\mathbf{Y}):=\left\{\mathrm{ar}:\
\begin{array}{l}
Y'(X,P)=\mathbf{Y};\ s\in\mathrm{Pd}_k\text{ vanishes outside }P_0(X)\cup P;\\
|t|\le\mathrm{rad}_1;\ \mathsf{Z}\in\mathrm{Box}_k;\ (A,\mathscr{C})\text{ lies in the domain of }F
\end{array}
\right\}.
\]

The polymer that belongs to an object is the index $\mathbf{Y}$ for the objects attached to $\mathbf{Y}$;
it is $X$ for $\mathscr{E}(X)$, $\mathscr{E}^{r}(X)$ and $\left\langle\bar{\Psi},M_k(X,\cdot)\Psi\right\rangle_1$;
and it is $Y'(X,P)$ for $G(s,X)$ inside the term $(X,P)$ of (G3.8.2) and for this term.
The set $\mathrm{Ar}$ for (T-G) is $\mathrm{Ar}(\mathbf{Y})$ for the objects attached to $\mathbf{Y}$;
it is $\mathrm{Ar}\left(Y'(X,P)\right)$ for $G(s,X)$ inside the term $(X,P)$ of (G3.8.2), for this term, and for every object that stands inside it ($\mathscr{E}^{r}(X)$, $\left\langle\bar{\Psi},M_k(X,\cdot)\Psi\right\rangle_1$, $E^{(\le3),x}(X)$ and the inserted fields).
For an object that does not contain $(P,s,t)$ or $\mathsf{Z}$, these entries of $\mathrm{ar}$ are dummy variables.
Invariance (covariance) on $\mathrm{Ar}$ relative to $\mathbf{Y}$ implies invariance (covariance) on the same set $\mathrm{Ar}$ relative to every polymer that contains $\mathbf{Y}$.
Indeed, a bond at distance at most $1$ from $\mathbf{Y}$ is at distance at most $1$ from the larger polymer, so that every transform admitted relative to the larger polymer is admitted relative to $\mathbf{Y}$.

The domain in the two connection slots is part of the notion.
It is $\tilde{R}_k\times\mathfrak{X}_k$ for $\mathscr{E}(X)$ and $\mathscr{E}^{r}(X)$ (clause (k4) of Section~I.5.4);
it is $\left\{(A_1,\mathscr{C}):\ A_1+\mathscr{C}\in\tilde{G}_k\right\}$ for $\left\langle\bar{\Psi},M_k(X,A_1+\mathscr{C})\Psi\right\rangle_1$ and for the primitives;
and it is $\tilde{R}_k^{-}\times\mathfrak{X}_k$ for the functions of Definition~24.D\,(D4), (D5).
\end{named}

\subsection*{25. The formal locality lemma}

\subsubsection*{G3.5.0. Setting, recalled clauses of Section~I.1.9, two local notions}

\paragraph{Data.}
Let $k\ge0$ be an integer, $\mu$ a complex number and $e_k>0$ a number, and let $L$, $M$, $b$, $\alpha$, $\rho$, $w$ be the constants of Section~I.8.
For Proposition~25.2 only (and for Proposition~25.3\,(b5), which takes up (o8) and (o10) of Proposition~25.2), the data comprise in addition $p_0$ and $\varepsilon$ of the line (C7) and $e_0$ of the line (C8).
Here $p_0$ enters through $p_{0,k}$ (Proposition~25.2\,(o10)), $\varepsilon$ through $\mathrm{rad}_1$, $\mathrm{rad}_2$, $\mathrm{rad}_4$ (the paragraph ``Completeness'' of Proposition~25.2), and $e_0$ through (SF) (the last two sentences of Proposition~25.2\,(o8)).
The class of sides is $\mathrm{Cl}$.
The side $\ell$ is not a datum: the statements assert that certain formulas do not contain it.
Objects of level $k+1$ have the mass $L\mu$, the charge $L^{\frac{1}{2}}e_k$ and, on a torus, the side $\ell/L$.

\paragraph{Covering objects.}
In what follows, $\tilde{B}$ is a complex function on the oriented bonds of $\mathbb{L}_k$ which is odd under reversal of the orientation and subject to no periodicity;
$\tilde{Y}$ is a union of $M$-cubes of $\mathbb{L}_k$;
and $\tilde{X}$, $\tilde{\mathbf{Y}}$, $\tilde{\square}$ are finite non-empty unions of $M$-cubes of $\mathbb{L}_k$ that are connected through faces, $\tilde{\square}$ being a single $M$-cube.
The number $d_M(\tilde{X})$ is defined as $d_M$ of a polymer (Section~I.1.1), with nets in $\mathbb{R}^3$ (\cite[Proposition~5.58(t10), Theorem~13.2]{QEDI}).
For $\ell\in\mathrm{Cl}$, a polymer of $T_k$ with the lift $\tilde{X}$ is small if and only if $d_M(\tilde{X})\le L$.
Indeed, by (t10) and $2M\cdot(L+2)<13L\cdot M\le\ell$, the number $d_M(\tilde{X})$ equals $d_M$ of the polymer as soon as one of the two numbers is at most $L$.
The displays of Part~I for phases, paths, blocks, neighbours and matrices ((I.1.2), (I.1.6), (I.2.1), (I.3.3), Section~I.1.1) are displays on $\mathbb{L}_k$, $\mathbb{L}^{\mathrm{f}}_k$, $\mathrm{Pzt}_k$; there they are read with $\tilde{B}$ in the place of $B$.

\paragraph{Clauses recalled from Section~I.1.9.}
\begin{itemize}
\item[(G3.r1)] A bond lies within the distance $\mathrm{rg}$ of a set of sites if both of its ends have sup-distance at most $\mathrm{rg}$ from that set.
\item[(G3.r2)] The data of an object attached to a polymer may be functions of one or of several indices (kernels), and they may depend on auxiliary parameters ($y$, $t$).
The phrase ``data composed with the covering map'' means the following: for a bond function $B$ on $T_k$, the function $b\mapsto B(\text{image of }b)$; for a kernel of the torus, which by (I.AP) is a function on the covering index sets, antiperiodic in each argument, its values at the arguments in (over) the $\mathrm{rg}$-neighbourhood of the lift.
\item[(G3.r3)] An element of a Grassmann algebra and a linear functional on kernels are admitted as objects attached to a polymer, in the sense that follows.
\end{itemize}

\paragraph{The covering Grassmann algebra.}
For a finite set $\tilde{U}$ of sites of $\mathbb{L}_k$, let $\mathrm{Grc}_k(\tilde{U})$ denote the Grassmann algebra over $\mathbb{C}$ with the following independent generators, $c$ and $c'$ being the photon sites of $x$ and $x'$:
\begin{itemize}
\item $\psi(x,\beta,\omega)$ with $c\in\tilde{U}$;
\item $\chi(x,x',\beta,\omega)$ with $(x,x')\in\mathrm{Pzt}_k$ and $c,c'\in\tilde{U}$;
\item $\Psi(u)$ and $W(u)$, where $u=(y,\beta,\omega)$ with $y\in\tilde{U}\cap\mathbb{Z}^3$ if $k\ge1$, and $u=(x,\beta,\omega)$ with $c\in\tilde{U}$ if $k=0$.
\end{itemize}
For $v\in\mathbb{Z}^3$ (on the subalgebra generated by $\psi$, $\chi$: for $v\in\eta_k\mathbb{Z}^3$), the translation $T_v$ relabels the generators and maps $\mathrm{Grc}_k(\tilde{U})$ to $\mathrm{Grc}_k(\tilde{U}+v)$.

\paragraph{The map $\mathrm{cov}_{\ell}$.}
Let $\ell\in\mathrm{Cl}$ with $\text{sup-diameter}(\tilde{U})+M<\ell$.
The homomorphism
\[
\mathrm{cov}_{\ell}:\mathrm{Grc}_k(\tilde{U})\to\mathrm{Gr}[\psi,\chi,\Psi,W]\qquad(\text{torus of side }\ell)
\]
maps every generator to the generator of the torus which has the same index in the sense of (I.AP), Section~I.3.6.
It is injective, because the covering map is injective on $\tilde{U}$ and a pair of $\mathrm{Pzt}_k$ with both points over $\tilde{U}$ is the shortest lift of its image (Section~I.1.1 with $\mathrm{sd}=M$).
If $\tilde{U}$ is replaced by $\tilde{U}+\ell n$, $n\in\mathbb{Z}^3$, then $\mathrm{cov}_{\ell}\circ T_{\ell n}$ multiplies every generator by $(-1)^{n_1+n_2+n_3}$; consequently it coincides with $\mathrm{cov}_{\ell}$ on even elements.
The same rule is applied at the level $k+1$, and for every side $\ell'$ which is a multiple of $M$ and larger than $2M$ (Sections~I.3.7, I.5.11);
in particular it is applied for $\ell'=\ell/L$ with $\ell\in\mathrm{Cl}$, where $\ell/L$ need not lie in $\mathrm{Cl}$:
\[
\mathrm{cov}_{\ell/L}:\mathrm{Grc}_{k+1}(\tilde{U}')\to\mathrm{Gr}[\psi,\chi,\Psi,W]\qquad(\text{torus of side }\ell/L).
\]
The rule defines an injective homomorphism as soon as $\text{sup-diameter}(\tilde{U})<\ell'$; the margin $M$ is needed for the assertion on the pairs.

\paragraph{Neighbourhoods.}
For a set $\tilde{U}$ of sites of $\mathbb{L}_k$ and a number $a$,
\[
\tilde{U}^{+a}:=\left\{c\in\mathbb{L}_k:\ \mathrm{dist}(c,\tilde{U})\le a\right\}\qquad(\text{sup-distance}).
\]
At the level $k+1$, for a set $\tilde{U}'$ of sites of $\mathbb{L}_{k+1}$ (for instance $\tilde{Y}'$ in Proposition~25.2\,(o7) and in Proposition~25.3\,(b5)), the neighbourhood is formed in $\mathbb{L}_{k+1}$:
\[
\tilde{U}'^{+a}:=\left\{c'\in\mathbb{L}_{k+1}:\ \mathrm{dist}(c',\tilde{U}')\le a\right\}.
\]

\subsubsection*{Proposition 25.1 (the exact operators of the step: one formula, range, what a row reads)}

\begin{named}{Proposition 25.1}
Let the data be as in Section~G3.5.0.
Each operator $\mathrm{Op}$ of the following table is defined by the display named in the table (a display of Part~I, one of the displays (23.1), (23.2), or, for the row no.~12, the formula of $C$ in \cite[Proposition~15.21(iv)]{QEDI}), for every side $\ell\in\mathrm{Cl}$, every complex bond function $B$ on $T_k$ and, where $Y$ occurs, every union $Y$ of $M$-cubes of $T_k$.
Read on the covering lattice, the same display defines, for all $\tilde{B}$ and $\tilde{Y}$, a linear operator $\mathrm{Op}_{\mathbb{L}}(\tilde{B},\tilde{Y})$ between the spaces of all functions on the index sets named.
These operators have the properties (f1)--(f4) of Section~I.1.9 with the range $\mathrm{rg}$ given in the table; that is, $\mathrm{Op}$ is given by one formula on $\mathbb{L}^{\mathrm{f}}_k$, of range $\mathrm{rg}$.
The property (f2) holds for all $v\in L\mathbb{Z}^3$ if $\mathrm{Op}$ does not depend on $Y$, and for all $v\in M\mathbb{Z}^3$ otherwise.
The formulas contain $k$, $\eta_k$, $\mu$, $e_k$, $b_k$, $b_L$, $\alpha$, $\rho$, $w$, $L$, $M$, the constant matrices of Section~I.2.1, the basis $\{e_i^0\}$ of \cite[Section~5.2.5]{QEDI} (chosen once for each $L$, and the same for all $k$ and $\ell$) and the resolvent parameter; they do not contain $\ell$.
The entries of every row are entire functions of those numbers $\tilde{B}(b)$ which the row reads.

For each row the table gives the operator, the levels, the index sets between which it acts (those of Section~I.1.9, with the spinor index where spinors occur; the operator of the row no.~12 and the map $A'\mapsto A'_L$ of the row no.~14 are scalar), the defining display, the range $\mathrm{rg}$, and what the row at the index $t$ reads.
\begin{description}
\item[No.~1.] Multiplication by $M_i(\cdot)$, $i=0,\dots,3$; multiplication by a constant matrix ($\mathsf{C}$, $\gamma_5$, $\varkappa$, $\mathrm{Mx}$).
All levels.
$\mathbb{L}^{\mathrm{f}}_k\to\mathbb{L}^{\mathrm{f}}_k$; for a constant matrix, any index set to itself.
(I.2.1), Section~I.2.1.
$\mathrm{rg}=0$.
The row reads nothing; $M_i(x)$ depends on the label $g$ only.
\item[No.~2.] $J$; $J^{*}$.
All levels.
Sites of $\mathbb{L}_k\to\mathbb{L}^{\mathrm{f}}_k$, and back.
(I.3.1).
$\mathrm{rg}=0$.
The row reads nothing.
\item[No.~3.] $R_i(B)$, $i=0,\dots,3$; $D_B$.
All levels.
$\mathbb{L}^{\mathrm{f}}_k\to\mathbb{L}^{\mathrm{f}}_k$.
(I.2.2), (I.2.3).
$\mathrm{rg}=\eta_k$.
The row reads $B$ on the one bond $(c,c_3)$.
\item[No.~4.] $Q_k(B)\cdot J^{*}$; $J\cdot Q_k^{T}(-B)$.
Levels $k\ge1$.
$\mathbb{L}^{\mathrm{f}}_k\to\mathbb{Z}^3$, and back.
Section~I.3.3.
$\mathrm{rg}=(1-\eta_k)/2$.
The row reads $B$ on the bonds of the paths $\boldsymbol{\Gamma}^{\pi}(y_{j+1},y_j)$, $j<k$, of the towers of Section~I.3.2: in the row of $Q_k(B)\cdot J^{*}$ at $y$, those of the towers of all $c\in\Delta_y$; in the row of $J\cdot Q_k^{T}(-B)$ at $x$, those of the tower of $c$.
These bonds are contained in $\Delta_y$, and both of their ends lie within $(1-\eta_k)/2$ of the row index, since each coordinate of a point of such a path lies between the corresponding coordinates of $y_{j+1}$ and $y_j$.
\item[No.~5.] $Q(B)$, $Q^{T}(-B)$ for $k\ge1$; $Q(B)\cdot J^{*}$, $J\cdot Q^{T}(-B)$ for $k=0$.
All levels.
$\mathbb{Z}^3$ (for $k=0$: $\mathbb{L}^{\mathrm{f}}_0$) $\to L\mathbb{Z}^3$, and back.
Section~I.3.3, (I.3.4).
$\mathrm{rg}=(L-1)/2$.
The row at $y'$ reads $B$ on the fine bonds of the paths $\boldsymbol{\Gamma}^{\pi}(y',y)$, $y\in B(y')$; the row of $Q^{T}(-B)$ at $y$ reads $B$ on the fine bonds of the paths $\boldsymbol{\Gamma}^{\pi}(y',y)$ for this $y$.
A path lies in the box spanned by $y'$ and $y$.
\item[No.~6.] $P(B)$ and $B_{k,y}(B)$, $y$ real, for $k\ge1$; $J\cdot P(B)\cdot J^{*}$ for $k=0$.
All levels.
$\mathbb{Z}^3\to\mathbb{Z}^3$ (for $k=0$: $\mathbb{L}^{\mathrm{f}}_0\to\mathbb{L}^{\mathrm{f}}_0$).
Sections~I.3.3, I.3.5.
$\mathrm{rg}=L-1$.
The row reads $B$ as in the row no.~5, $y'$ being the centre of the $L$-block of the row.
\item[No.~7.] $J\cdot P_k(B)\cdot J^{*}$; $J\cdot P_{k+1}(B)\cdot J^{*}$, $\operatorname{Pen}_{y}(B)$, $F_{y}(B)$ ($y\in Y_{k}$).
Levels $k\ge1$.
$\mathbb{L}^{\mathrm{f}}_k\to\mathbb{L}^{\mathrm{f}}_k$.
Sections~I.3.3, I.3.5.
$\mathrm{rg}=1-\eta_k$ for the first operator, and $\mathrm{rg}=L-\eta_k$ for the three others.
The row of the first operator reads $B$ on the bonds contained in $\Delta_{[c]}$ (the towers of all $c'\in\Delta_{[c]}$);
the row of each of the three others reads $B$ on the bonds contained in $\mathrm{Bl}_{k+1}(y')$, where $c\in\mathrm{Bl}_{k+1}(y')$ (the towers, and the paths of the row no.~5), and, for $F_{y}(B)$, also on $(c,c_3)$.
\item[No.~8.] $F_{y}(B)=D_B+\mu+b_L\cdot J\cdot P(B)\cdot J^{*}+i\gamma_5 y$, $y$ real, in particular $F_0(B)$.
Level $k=0$.
$\mathbb{L}^{\mathrm{f}}_0\to\mathbb{L}^{\mathrm{f}}_0$.
Section~I.3.9.
$\mathrm{rg}=L-1$.
The row reads $B$ on the bond $(c,c_3)$, and as in the row no.~6.
\item[No.~9.] $T_k(B)$, $\bar{T}_k(B)$; $T^{(1)}(B,Y)$, $\bar{T}^{(1)}(B,Y)$.
Levels $k\ge1$.
$\mathbb{L}^{\mathrm{f}}_k\to\mathbb{Z}^3$.
Section~I.3.5; (23.1), (23.2).
$\mathrm{rg}=(1+\eta_k)/2$.
The row at $y$ reads $B$ on the bonds contained in $\Delta_y$ and on the bonds $(c,c_3)$, $c\in\Delta_y$; it reads $Y$ at $y$ and at the sites $c_3$, $c\in\Delta_y$.
\item[No.~10.] $T^{(2)}(B,Y)$, $\bar{T}^{(2)}(B,Y)$.
Levels $k\ge1$.
$\mathrm{Pzt}_k\to\mathbb{Z}^3$.
(23.1).
$\mathrm{rg}=(1+\eta_k)/2$.
The row at $y$ reads $B$ on the bonds contained in $\Delta_y$, and $Y$ as in the row no.~9.
\item[No.~11.] $\delta_{\alpha,B}$.
All levels.
$\mathbb{L}^{\mathrm{f}}_k\to\mathrm{Pzt}_k$.
(I.1.7).
$\mathrm{rg}=M$.
The row at $(x,x')$ reads $B$ on the bonds of the six paths $\boldsymbol{\Gamma}^{\pi}(c,c')$ (on no bond if $c=c'$); the columns of this row are $x$ and $x'$.
\item[No.~12.] $C$.
All levels.
$\{\Upsilon\}\to$ bonds of $\mathbb{Z}^3$.
\cite[Section~5.2.5, Proposition~15.21(iv)]{QEDI}.
$\mathrm{rg}=L-1$.
The row reads nothing; $C(b',\Upsilon)=0$ unless $\operatorname{supp}\Upsilon$ lies in the union of the $L$-blocks of the ends of $b'$.
\item[No.~13.] $\mathbf{1}_Y$, the multiplication by the indicator function of $Y\times O_{\mathrm{cub}}$.
All levels.
$\mathbb{L}^{\mathrm{f}}_k\to\mathbb{L}^{\mathrm{f}}_k$.
Section~I.1.4.
$\mathrm{rg}=0$.
The row reads $Y$ at the site $c$.
\item[No.~14.] $A'\mapsto A'_L$, $f'\mapsto f'_L$, $F'\mapsto F'_L$, $v'\mapsto v'_L$.
From level $k+1$ to level $k$.
From the bond functions on $\mathbb{L}_{k+1}$ and the functions on $\mathbb{L}^{\mathrm{f}}_{k+1}$, $\mathrm{Pzt}_{k+1}$, $\mathbb{Z}^3$ to the bond functions on $\mathbb{L}_k$ and the functions on $\mathbb{L}^{\mathrm{f}}_k$, $\mathrm{Pzt}_k$, $L\mathbb{Z}^3$.
Section~I.3.7.
$\mathrm{rg}=0$ (the row at $t$ has the single column $t/L$).
The row reads nothing.
The translation $T_v$ at level $k$ goes with $T_{v/L}$ at level $k+1$, $v\in L\mathbb{Z}^3$; $\operatorname{supp}(t/L)=\operatorname{supp}(t)/L$; a column index $t'$ of level $k+1$ is contained in $\tilde{U}$ if $L\operatorname{supp}t'$ lies in $\tilde{U}$; and the side at level $k+1$ is $\ell/L$.
\end{description}

The kernels are taken by (I.1.4) where the operator acts on spinors on fermion sites, and with the sum over $y'$ where it acts on $\mathfrak{S}(1,\ell)$.
On the covering lattice they are as follows:
\begin{align*}
\left(Q_k(\tilde{B})\cdot J^{*}\right)(y,x)&=\mathbf{1}[c\in\Delta_y]\cdot\varphi_{\tilde{B}}(y,c)\cdot1_4,\qquad
\left(J\cdot Q_k^{T}(-\tilde{B})\right)(x,y)=\mathbf{1}[c\in\Delta_y]\cdot\varphi_{\tilde{B}}(y,c)^{-1}\cdot1_4;\tag{K1}\\
P(\tilde{B})(y_1,y_2)&=
\begin{cases}
L^{-3}\cdot U_{\tilde{B}}(y',y_1)^{-1}\cdot U_{\tilde{B}}(y',y_2)\cdot1_4 & \text{if }y_1,y_2\text{ lie in one block }B(y'),\\
0 & \text{otherwise};
\end{cases}\tag{K2}
\end{align*}
for $k=0$ the kernel of $J\cdot P(\tilde{B})\cdot J^{*}$ at $(x_1,x_2)$ is $P(\tilde{B})(c_1,c_2)$;
\begin{align*}
B_{k,y}(\tilde{B})&=\left(b_k+i\gamma_5 y\right)^{-1}\cdot\left(1-P(\tilde{B})\right)+\left(b_k+b_L+i\gamma_5 y\right)^{-1}\cdot P(\tilde{B}),\\
B_{k,y}(\tilde{B})\cdot\gamma_5\cdot B_{k,y}(\tilde{B})&=\gamma_5\cdot\left[\left(b_k+i\gamma_5 y\right)^{-2}\cdot\left(1-P(\tilde{B})\right)+\left(b_k+b_L+i\gamma_5 y\right)^{-2}\cdot P(\tilde{B})\right];\tag{K3}
\end{align*}
and for $k=0$
\begin{equation}
F_0(\tilde{B})(x,x')=48\cdot\left\{\mathbf{1}[x'=x]\cdot\left(\mu+\frac{\rho}{2}\cdot\mathrm{sw}\right)\cdot1_4-\sum_{i=0}^{3}\mathbf{1}\left[x'=x^{(i)}\right]\cdot M_i(x)\cdot U_{\tilde{B}}\left(x,x^{(i)}\right)\right\}+b_L\cdot P(\tilde{B})(c,c').\tag{K4}
\end{equation}

None of the ten operators $S_k(A)$, $S_{k,y}(A)$, $H_{k}(A)$, $\bar{H}_{k}(A)$, $\mathfrak{H}_{k}(A)$, $D_k(A)$, $M_k(A)$, $\dot{M}_k(A)$, $\Gamma_k(A)$, $\Gamma_{k,y}(A)$, and none of the four photon operators $H_{k}$, $C_{k}^{1/2}$, $\mathrm{Sq}_{k}$, $A^{\min}_{k}$, is an entry of the table; nothing is asserted on their range.
None of these fourteen operators enters Definition~24.D, the definitions of Section~I.5.6 or the definition of $\acute{E}_k$.
Where $H_{k}$, $C_{k}^{1/2}$, $\bar{H}_{k}$, $\Gamma_k$, $M_k$, $S_{k,y}$ occur in these definitions, they carry a parameter family or a polymer as first argument, and the operators so formed are members of the families (p1)--(p5) of Section~I.5.8.
\end{named}

\begin{named}{Remark}
Proposition~25.1 corresponds to \cite[Proposition~17.24(b)]{QEDI} with its proof;
to the first paragraph of \cite[Definition~5.5(2)]{QEDI}, the notion defined there being replaced by that of Section~I.1.9;
to the observation~(i) in the proof of \cite[Proposition~16.29]{QEDI};
and to the assertions on the exact operators in Step~2 and in the paragraph on exact operators and forms of Step~3 of the proof of \cite[Proposition~19.12]{QEDI}.
The relation to these places is as follows.
\begin{itemize}
\item[(1)] The index sets are $\mathbb{L}^{\mathrm{f}}_k=\mathbb{L}_k\times O_{\mathrm{cub}}$ and $\mathrm{Pzt}_k$; translations act on $c$ only; the hopping matrices depend on the label.
\item[(2)] The notion of Section~I.1.9 agrees in its conditions with that of \cite[Definition~5.5(2)]{QEDI}:
(f1)--(f3) are the conditions (i)--(iii), and (f4) is the requirement of locality which is stated there after these conditions.
Both require, for a number $r\ge0$ (in (f4) the number $\mathrm{rg}$), that for all $\tilde{B}$, $\tilde{Y}$ and every row index $t$ the row at $t$ vanish at the columns $t'$ for which the smallest sup-distance of a point of $\operatorname{supp}t$ from a point of $\operatorname{supp}t'$ exceeds $r$;
and that this row read $\tilde{B}$ on the bonds both of whose end points have sup-distance at most $r$ from $\operatorname{supp}t$, and $\tilde{Y}$ on the sites at sup-distance at most $r$ from $\operatorname{supp}t$.
Hence a number $\mathrm{rg}$ as in (f4) is a range in the sense of \cite[Definition~5.5(2)]{QEDI}.
Likewise, \cite[Proposition~17.24(b)]{QEDI} states for its operators a range, namely the common range $M$, and, operator by operator, the columns of the rows and the bonds and sites which the rows read;
the table of Proposition~25.1 states, for each of its operators, a range of its own and the bonds and sites which the rows read.
The operator $\mathrm{Op}_{\mathbb{L}}(\tilde{B},\tilde{Y}):=\left(\text{upper density of }\tilde{Y}\right)\cdot1$, the operator $\mathrm{Op}(B,Y)$ on the torus being the multiplication by the proportion of the sites of $T_k$ that lie in $Y$, satisfies the conditions (i)--(iii) of \cite[Definition~5.5(2)]{QEDI}, that is, (f1)--(f3), and each of its rows has no column other than the row index.
It violates, however, the second part of the requirement of locality of \cite[Definition~5.5(2)]{QEDI}, and in the same way (f4).
Indeed, let $r\ge0$, let $t$ be an index, and let $\tilde{Y}_r$ be the union of the $M$-cubes of $\mathbb{L}_k$ that contain no site at sup-distance at most $r$ from $\operatorname{supp}t$.
The sets $\emptyset$ and $\tilde{Y}_r$ contain the same sites at sup-distance at most $r$ from $\operatorname{supp}t$, namely none;
but the row at $t$ vanishes at $\tilde{Y}=\emptyset$ and does not vanish at $\tilde{Y}=\tilde{Y}_r$, since $\tilde{Y}_r$ has a finite complement and hence the upper density $1$.
Thus the family of the operators $\mathrm{Op}_{\mathbb{L}}(\tilde{B},\tilde{Y})$ has the properties (f1)--(f3) but not (f4);
it exhibits $\mathrm{Op}$ as given by one formula neither in the sense of Section~I.1.9 nor in that of \cite[Definition~5.5(2)]{QEDI}, and the requirement of locality does not follow from the conditions (i)--(iii).
\item[(3)] The operators $T^{(i)}$, $\bar{T}^{(i)}$ map into the functions on the whole unit lattice; they are formed with $Y$ at the unit sites of $Y$ and have zero rows at the other unit sites.
In this respect Proposition~25.1 agrees with \cite[Proposition~17.24(b)]{QEDI}, in which the operators $T^{(i)}(\tilde{B})$ and $\bar{T}^{(i)}(\tilde{B})$ map into the spinor functions on $\mathbb{Z}^3$, are formed with $\tilde{Y}$ at the unit sites of $\tilde{Y}$ and have zero rows at the other unit sites.
\item[(4)] The table contains operators which \cite[Proposition~17.24(b)]{QEDI} does not name: those of the rows nos.~2--8 (the operators of $\acute{E}_k$ and of the averaging) and of the row no.~14.
\item[(5)] The matrices $P^{\boldsymbol{\sigma}}_{\nu}$ of \cite[Proposition~17.24(b)]{QEDI} are replaced by $M_i(x)$.
\item[(6)] The formula of $C$ is that of \cite[Section~5.2.5, Proposition~15.21(iv)]{QEDI}.
\end{itemize}
\end{named}

\subsubsection*{Proposition 25.2 (forms, functionals, homomorphisms and integrations attached to a polymer)}

\begin{named}{Proposition 25.2}
Let the data be as in Section~G3.5.0.
Each of the families $O_{\ell}(X;\mathrm{data})$ of (o1)--(o9) below is given by one formula of the range $\mathrm{rg}$ named there, in the sense of Section~I.1.9 together with (G3.r1)--(G3.r3).
The object (o10) is not attached to a polymer; for it, the assertion is the statement made under (o10).
This means: there is one map $\Phi$, which does not contain $\ell$ and is defined for all $\tilde{X}$ (as restricted below) and for all data on the $\mathrm{rg}$-neighbourhood of $\tilde{X}$, such that
\[
\Phi\left(\tilde{X}+v,T_v\,\mathrm{data}\right)=T_v\,\Phi\left(\tilde{X},\mathrm{data}\right)\qquad\text{for }v\in LM\mathbb{Z}^3
\]
(for (o1)--(o6), (o8), (o9) also for $v\in M\mathbb{Z}^3$), and such that for every $\ell\in\mathrm{Cl}$ and every polymer $X$ of $T_k$ with a lift $\tilde{X}$ satisfying
\[
\text{sup-diameter}(\tilde{X})+2\,\mathrm{rg}+M<\ell
\]
one has $O_{\ell}(X;\mathrm{data})=\Phi\left(\tilde{X},\text{data composed with the covering map}\right)$;
for elements this reads $O_{\ell}(X;\mathrm{data})=\mathrm{cov}_{\ell}\left(\Phi(\dots)\right)$, and for homomorphisms $\mathrm{cov}_{\ell}\circ\Phi(\dots)=O_{\ell}(X;\mathrm{data})\circ\mathrm{cov}_{\ell}$.
For the operation $L$ of (o7), in which $D_{k+1}$ occurs, the side $\ell$ is any side with $\ell\in\mathrm{Cl}$ and $\ell/L\in\mathrm{Cl}$, and a polymer $Y$ of level $k+1$, its lift $\tilde{Y}'$ and the side $\ell/L$ stand in the places of $X$, $\tilde{X}$ and $\ell$.
The value is independent of the choice of the lift.
Three of the objects are not attached to a polymer:
$S_L$ of (o7) and the integration of (o9) are indexed by a finite set $\tilde{U}$ of sites of $\mathbb{L}_k$, which stands in the place of $\tilde{X}$, under the condition on $(\tilde{U},\ell)$ stated there;
and (o10) contains neither a polymer nor a range.
The object $\bar{X}^{L}$ of (o7) is a set of sites, namely the image of $\Phi(\tilde{X})$ under the covering map.
\begin{itemize}
\item[(o1)] (Volume and forms; $\mathrm{rg}=0$; no data.)
One has $\operatorname{Vol}(X)=M^3\cdot|\tilde{X}|_M$.
For a constant $4\times4$ matrix $\mathrm{Mx}$,
\[
\int_X\mathrm{d}x\,\bar{\psi}(x)\cdot\mathrm{Mx}\cdot\psi(x)=\mathrm{cov}_{\ell}\left(\int_{\tilde{X}}\mathrm{d}x\,\bar{\psi}(x)\cdot\mathrm{Mx}\cdot\psi(x)\right),
\]
where the integral on the right is the finite sum $\frac{\eta_k^3}{48}\sum_x$ over the fermion sites $x$ over $\tilde{X}$, an element of $\mathrm{Grc}_k(\tilde{X})$.
In particular this applies to $\langle\bar{\psi},\psi\rangle_{\square}$, with $\mathrm{Mx}=1$ and $\tilde{X}=\tilde{\square}$.
The kernel of this element is $48\,\eta_k^{-3}\cdot\mathbf{1}[x=x']\cdot\mathbf{1}[c\in\tilde{X}]\cdot\mathrm{Mx}$.
The sum extends over the sites over $\tilde{X}$, not over the periodic extension of $X$.
\item[(o2)] (Pairing of the unit lattice; $\mathrm{rg}=0$; data: a kernel $O(u,u')$ at the unit indices in $\tilde{X}$.)
The element is
\[
\sum_{y,y'\in\tilde{X}}\bar{\Psi}(y)\cdot O(y,y')\cdot\Psi(y')\quad(k\ge1),\qquad
\int_{\tilde{X}}\mathrm{d}x\int_{\tilde{X}}\mathrm{d}x'\,\bar{\Psi}(x)\cdot O(x,x')\cdot\Psi(x')\quad(k=0).
\]
For an operator $O$ of the torus whose kernel vanishes unless both indices lie in (over) $X$, the image of this element under $\mathrm{cov}_{\ell}$ is $\langle\bar{\Psi},O\Psi\rangle$.
In Definition~24.D one has $O=M_k(X,\cdot)$, $k\ge1$.
\item[(o3)] (Extraction; $\mathrm{rg}=0$; $\tilde{X}$ with $d_M(\tilde{X})\le L$; data: a number $\mathrm{Ec}_{00}$ and a matrix kernel $\mathrm{Ec}_{20}(x,x')$ at the fermion sites over $\tilde{X}$.)
Put
\[
\mathrm{al}_{0}:=\operatorname{Vol}(\tilde{X})^{-1}\cdot\mathrm{Ec}_{00},\qquad
\mathrm{al}_{2}:=\operatorname{Vol}(\tilde{X})^{-1}\cdot\int_{\tilde{X}}\mathrm{d}x\int_{\tilde{X}}\mathrm{d}x'\,\mathrm{Ec}_{20}(x,x').
\]
For a family $E$ with (k2) and a small polymer $X$, the numbers $\mathrm{al}_{0}(E,X)$, $\mathrm{al}_{2}(E,X)$ of Section~I.5.6 are these numbers at $\mathrm{Ec}_{00}=E_{00}(X;0,0)$ and at $\mathrm{Ec}_{20}$ equal to the kernel $E_{20}(X;0,0;\cdot,\cdot)$ (Section~I.5.1) with both arguments over the lift $\tilde{X}$; that is, they are formed with the parts of bidegree $(0,0)$ and $(2,0)$ of $\mathrm{cov}_{\ell}^{-1}\left(E(X;0,0)\right)$.
The element
\[
(RE)(X)=E(X)-\mathrm{al}_{0}\cdot\operatorname{Vol}(X)-\int_X\mathrm{d}x\,\bar{\psi}(x)\cdot\mathrm{al}_{2}\cdot\psi(x)
\]
is formed with (o1).
\item[(o4)] (Adjustment; $\mathrm{rg}=0$; $\tilde{X}$ with $d_M(\tilde{X})\le L$; data: $\tilde{B}$ on the bonds contained in $\tilde{X}$.)
Let $\boldsymbol{\Theta}_{\tilde{X},\tilde{B}}$ be the homomorphism of the subalgebra of $\mathrm{Grc}_k(\tilde{X})$ generated by the $\psi(\xi)$ into $\mathrm{Grc}_k(\tilde{X})$ which is defined by the display (I.5.3) with the chains of $M$-cubes of $\tilde{X}$, the centres $z_i$, the points $p_i=(z_{i-1},g)$, the transports $u(p,p')=U_{\tilde{B}}(p,p')$ and $\int_{\tilde{X}}\mathrm{d}\tilde{x}$.
Then
\[
\mathrm{cov}_{\ell}\circ\boldsymbol{\Theta}_{\tilde{X},\tilde{B}}=\boldsymbol{\Theta}_{X,\mathscr{T}}\circ\mathrm{cov}_{\ell}\qquad\text{if }\tilde{B}\text{ is the lift of }\mathscr{T}.
\]
All pairs $(p_i,p_{i+1})$ are of kind (I) for $\tilde{X}$.
\item[(o5)] (Substitution; $k\ge1$; $\mathrm{rg}=\eta_k$; data: $\tilde{B}$ on the bonds that meet $\tilde{\mathbf{Y}}$.)
Let $\mathrm{sub}_{\tilde{\mathbf{Y}},\tilde{B}}$ be the homomorphism (G3.8.3a)$^{\mathrm{ch}}$ formed with the rows nos.~9, 10 of Proposition~25.1 at the unit sites of $\tilde{\mathbf{Y}}$ and with $\tilde{Y}=\tilde{\mathbf{Y}}$.
It maps the subalgebra of $\mathrm{Grc}_k\left(\tilde{\mathbf{Y}}^{+\eta_k}\right)$ generated by $\psi$, $\chi$ and the $\Psi(u)$, $u\in\tilde{\mathbf{Y}}$, into the subalgebra generated by $\psi$, $\chi$.
\item[(o6)] (Vertex; $k=0$; $\mathrm{rg}=1$; data: $\mathrm{Zf}$ on the bonds that meet $\tilde{\square}$.)
The element $\mathrm{Vx}(\mathrm{Zf},\tilde{\square})$ is the display (23.3) with $\int_{\tilde{\square}}\mathrm{d}x$; it lies in $\mathrm{Grc}_0\left(\tilde{\square}^{+1}\right)$.
\item[(o7)] (Reblocking.)
\emph{The homomorphism $S_L$} ($\mathrm{rg}=0$; no data).
$S_L$ is the homomorphism of the subalgebra of $\mathrm{Grc}_k(\tilde{U})$ generated by $\psi$, $\chi$ into the subalgebra of $\mathrm{Grc}_{k+1}(\tilde{U}/L)$ generated by $\psi$, $\chi$ given by
\[
\psi(x,\cdot)\mapsto L^{-1}\cdot\psi(x/L,\cdot),\qquad \chi(x,x',\cdot)\mapsto L^{-1-\alpha}\cdot\chi(x/L,x'/L,\cdot).
\]
One has $\mathrm{cov}_{\ell/L}\circ S_L=S_L\circ\mathrm{cov}_{\ell}$ if $\text{sup-diameter}(\tilde{U})+LM<\ell$, with $\mathrm{cov}_{\ell/L}$ as in Section~G3.5.0; with the injective homomorphisms of Section~G3.5.0 this identity holds as soon as $\text{sup-diameter}(\tilde{U})<\ell$.
Moreover $T_{v/L}\circ S_L=S_L\circ T_v$ for $v\in LM\mathbb{Z}^3$.

\emph{The set $\bar{X}^{L}$} ($\mathrm{rg}=0$; no data).
$\bar{X}^{L}$ is the image of the union of those $LM$-cubes of $\mathbb{L}_k$ which contain an $M$-cube of $\tilde{X}$.

\emph{The operation $L$} (an object attached to a polymer $Y$ of level $k+1$; $\mathrm{rg}=\eta_{k+1}$ at level $k+1$; $\ell\in\mathrm{Cl}$ and $\ell/L\in\mathrm{Cl}$).
For a finite non-empty union $\tilde{Y}'$ of $M$-cubes of $\mathbb{L}_{k+1}$ which is connected through faces, let $\mathrm{Ix}(\tilde{Y}')$ be the set of the non-empty unions $\tilde{X}$ of $M$-cubes of $\mathbb{L}_k$ which are contained in $L\tilde{Y}'$, are connected through faces, and for which the union of the $LM$-cubes containing an $M$-cube of $\tilde{X}$ is $L\tilde{Y}'$.
This set does not contain $\ell$, and $\mathrm{Ix}(\tilde{Y}'+v/L)=\mathrm{Ix}(\tilde{Y}')+v$ for $v\in LM\mathbb{Z}^3$.
Let $E$ be a family of even elements of $\mathrm{Gr}[\psi,\chi]$ of level $k$ with (k2), let $Y\in D_{k+1}$ have a lift $\tilde{Y}'$ with $\text{sup-diameter}(\tilde{Y}')+2\eta_{k+1}+M<\ell/L$, and let $(\mathscr{A},\mathscr{C})$ be a pair of bond functions on $T_{k+1}$ for which the elements $E(X;\mathscr{A}_L,\mathscr{C}_L)$ are defined.
Then
\[
(LE)(Y;\mathscr{A},\mathscr{C},\cdot)=\sum_{X:\ \bar{X}^{L}=LY}S_L\left[E(X;\mathscr{A}_L,\mathscr{C}_L,\cdot)\right]=\mathrm{cov}_{\ell/L}\left(\Phi_L\right),\qquad
\Phi_L:=\sum_{\tilde{X}\in\mathrm{Ix}(\tilde{Y}')}S_L\,\mathrm{Ecv}(\tilde{X}),
\]
and $\Phi_L$ is an element of the subalgebra of $\mathrm{Grc}_{k+1}\left(\tilde{Y}'^{+\eta_{k+1}}\right)$ generated by $\psi$, $\chi$.
The data are the elements $\mathrm{Ecv}(\tilde{X}):=\mathrm{cov}_{\ell}^{-1}\left(E(X;\mathscr{A}_L,\mathscr{C}_L)\right)$, where $X$ is the image of $\tilde{X}$ in $T_k$ and $\tilde{X}\in\mathrm{Ix}(\tilde{Y}')$;
a translation by $v\in LM\mathbb{Z}^3$ at level $k$ goes with the translation by $v/L$ at level $k+1$;
and the value is independent of the choice of the lift.
Suppose now that $\ell,\ell/L\in\mathrm{Cl}$, that $\tilde{Y}'$ is as in the definition of $\mathrm{Ix}(\tilde{Y}')$, that $E$ and $(\mathscr{A},\mathscr{C})$ are as before, and that in the place of the margin $\text{sup-diameter}(\tilde{Y}')+2\eta_{k+1}+M<\ell/L$ only
\[
\text{sup-diameter}(L\tilde{Y}')+2\eta_k+M<\ell
\]
is given.
Then the image $Y$ of $\tilde{Y}'$ in $T_{k+1}$ is a polymer of level $k+1$ with the lift $\tilde{Y}'$;
the polymers $X$ with $\bar{X}^{L}=LY$ are the images of the $\tilde{X}\in\mathrm{Ix}(\tilde{Y}')$, one for each $\tilde{X}$;
the elements $\mathrm{Ecv}(\tilde{X})$ are defined;
and the same identity holds, $\mathrm{cov}_{\ell/L}$ being the injective homomorphism of Section~G3.5.0 on $\mathrm{Grc}_{k+1}\left(\tilde{Y}'^{+\eta_{k+1}}\right)$.
\item[(o8)] (Traces of $\acute{E}_k$; $\mathrm{rg}=0$; data: $\tilde{B}$ on the bonds contained in $\tilde{X}$; a kernel $O(x,x')$, with values in the $4\times4$ matrices, at the fermion sites over $\tilde{X}$; $y$.)
For $k\ge1$ put
\begin{multline*}
\mathrm{Trk}_k(\tilde{X};\tilde{B},O,y):=\sum_{y_1,y_2\in\tilde{X}\cap\mathbb{Z}^3}\operatorname{tr}\left[\vphantom{\int_{\Delta_{y_1}}}\left(B_{k,y}(\tilde{B})\cdot\gamma_5\cdot B_{k,y}(\tilde{B})\right)(y_2,y_1)\right.\\
\left.\cdot\int_{\Delta_{y_1}}\mathrm{d}x\int_{\Delta_{y_2}}\mathrm{d}x'\,\varphi_{\tilde{B}}(y_1,c)\cdot O(x,x')\cdot\varphi_{\tilde{B}}(y_2,c')^{-1}\right],
\end{multline*}
and for $k=0$ put
\[
\mathrm{Trk}_0'(\tilde{X};\tilde{B},O):=\int_{\tilde{X}}\mathrm{d}x\int_{\tilde{X}}\mathrm{d}x'\,\operatorname{tr}\left[F_0(\tilde{B})(x,x')\cdot O(x',x)\right],\qquad
\mathrm{Trk}_0''(\tilde{X};O):=\int_{\tilde{X}}\mathrm{d}x\,\operatorname{tr}\left[\gamma_5\cdot O(x,x)\right].
\]
For every operator $O$ on $\mathfrak{S}^{\mathrm{ch}}(\eta_k,\ell)$ whose kernel (I.1.4) vanishes unless both arguments lie over $X$,
\begin{gather*}
\operatorname{Tr}\left[\gamma_5\cdot B_{k,y}(B)\cdot Q_k(B)\cdot J^{*}\cdot O\cdot J\cdot Q_k^{T}(-B)\cdot B_{k,y}(B)\right]=\mathrm{Trk}_k(\tilde{X};\tilde{B},O,y)\qquad(k\ge1),\\
\operatorname{Tr}\left[F_0(B)\cdot O\right]=\mathrm{Trk}_0'(\tilde{X};\tilde{B},O),\qquad \operatorname{Tr}\left[\gamma_5\cdot O\right]=\mathrm{Trk}_0''(\tilde{X};O)\qquad(k=0),
\end{gather*}
where $\tilde{B}$ and $O$ on the right are the lifts.
The kernels (K2)--(K4) of Proposition~25.1 at indices in (over) $\tilde{X}$ read $\tilde{B}$ on bonds contained in $\tilde{X}$, and, for $F_0(\tilde{B})\left(x,x^{(3)}\right)$, on the bond $(c,c_3)$, both of whose ends lie in $\tilde{X}$.
Consequently, in the setting in which Section~I.5.7 forms these numbers (data (SF), $B\in\tilde{G}_k$, and a family (p5) as in the data of Proposition~13, of whose properties only the first half of W12 is required), the number $\acute{E}_k(X,B)$ of Section~I.5.7 is the value of a functional that does not contain $\ell$, taken at the lifts of $B$ on the bonds contained in $X$ and of the kernels $S_{k,y}(X,B)(x,x')$, $x,x'$ over $X$, $y\ge0$.
The quantities $E^{d}_k$, $E_k^{\mathrm{det}}$, ${\delta E^{\mathrm{det}}}_k$ are formed from this functional by $B\mapsto-B$, $B\mapsto\overline{B}$, complex conjugation and differences.
\item[(o9)] (Integration over $W$; $\mathrm{rg}=0$; no data; for every finite set $\tilde{U}$ of sites of $\mathbb{L}_k$ and every $\ell\in\mathrm{Cl}$ with $\text{sup-diameter}(\tilde{U})+M<\ell$.)
Let $\int_{\mathrm{c}}$ be the integral (I.3.5) with $\mathrm{d}\mu_I$ on the index set of the generators $W(u)$ of $\mathrm{Grc}_k(\tilde{U})$, with $\nu=1$ for $k\ge1$ and $\nu=1/48$ for $k=0$.
Then
\[
\int\mathrm{d}\mu_I(W)\,\mathrm{cov}_{\ell}\left(F\right)=\mathrm{cov}_{\ell}\left(\int_{\mathrm{c}}F\right)\qquad\text{for all }F\in\mathrm{Grc}_k(\tilde{U}).
\]
The moments are those of (I.B4); they contain $\nu$ and nothing else.
\item[(o10)] (Integration over $\mathsf{Z}$; no polymer, no range, no data; for $p_{0,k}>0$, which follows from $e_k<1$.)
In the coordinates $\mathsf{Z}(\Upsilon)$, the measure $\hat{\mu}$ is the product, over the basis, of copies of one probability measure on $\mathbb{R}$, which depends on $p_{0,k}$ only; and
\[
\delta\varepsilon^{0}_{k}=-d_V\cdot\log\mathbb{P}\left(|\mathrm{gn}|\le p_{0,k}\right),
\]
where $\mathrm{gn}$ is a standard normal random variable and $d_V=2-2L^{-3}$.
\end{itemize}

\emph{Completeness.}
Apart from the primitives, the data $(m,\mathscr{E})$, $\varepsilon^{\mathrm{v}}$ and the family $\mathfrak{P}_k$, the connections with their pointwise sums and their multiples by numbers, the parameters, the variables of integration and the generators, and apart from the operations of a Grassmann algebra over $\mathbb{C}$ (finite sums, products, multiples by numbers, $\exp$, the series of (G3.8.4), the part $F^{[2,0]}$ of bidegree $(2,0)$, the application of a row of an operator to a family of generators, and the homomorphism of Definition~24.D\,(D1) which adds a family of odd elements to the generators; $\mathrm{cov}_{\ell}$ commutes with all of them), the displays (G3.8.1)$^{\mathrm{ch}}$--(G3.8.6)$^{\mathrm{ch}}$ of Definition~24.D, the operations $\varepsilon$, $m$, $R$, $\mathrm{nat}$, $L$ of Section~I.5.6 and the displays of $\acute{E}_k$, $E^{d}_k$, $E_k^{\mathrm{det}}$, ${\delta E^{\mathrm{det}}}_k$ of Section~I.5.7 contain the following objects and operations, and no others.
\begin{itemize}
\item In (G3.8.1)$^{\mathrm{ch}}$: $C$, $\delta_{\alpha,\pm B}$, $\mathsf{C}$ and $B\mapsto-B$.
\item In (G3.8.3)$^{\mathrm{ch}}$: (o2) (in $G^{(1)}$, $k\ge1$), (o1) (in $G^{(3)}$, $G^{(5)}$), (o5), (o6), the functionals $\oint_{\mathrm{rad}}$, and $m'=m+m(\mathscr{E})$, $\mathscr{E}^{r}=R\mathscr{E}$.
\item In (G3.8.2): the combinatorial data $P_0(X)$, $\mathrm{Yp}(X,P)$, $Y'(X,P)$, the sum over the pairs $(X,P)$ with $Y'(X,P)=\mathbf{Y}$, the operation $\int\mathrm{d}s_P\,\frac{\partial}{\partial s_P}$, and the distinction between single $M$-cubes and other polymers made in Definition~24.D\,(D3).
\item In (G3.8.4): (o9), (o10), $\mathrm{urs}^{T}$, the index sets (the sets $\{X_i\}$ with the union $X$, and the families $(X_1,\dots,X_n)$ with the union $X$) and the numbers $1/n!$.
\item In (G3.8.6)$^{\mathrm{ch}}$: (o7) together with the row no.~14 of Proposition~25.1, the operations $R$, $\mathrm{nat}$ and the displays of $\acute{E}_k$, the numbers $L^3$, $L$, $\varepsilon(\mathscr{E})$, $m(\mathscr{E})$ and $\delta\varepsilon^{0}_{k}$ ((o10)).
\item In $R$: the distinction between small and large polymers, (o1), (o3).
\item In $\varepsilon$ and $m$: (o3) and the sum over the small polymers that contain a given $M$-cube; for a lift $\tilde{\square}$ of the cube these polymers are the images of the finite non-empty unions $\tilde{X}$ of $M$-cubes of $\mathbb{L}_k$ which are connected through faces, contain $\tilde{\square}$ and satisfy $d_M(\tilde{X})\le L$, and this set of unions does not contain $\ell$.
\item In $\mathrm{nat}$: the distinction between small and large polymers, (o4).
\item In $L$: (o7), that is, the set $\bar{X}^{L}$, the index set of the sum over the polymers $X$ with $\bar{X}^{L}=LY$, and the homomorphism $S_L$;
and the row no.~14 of Proposition~25.1, that is, $\mathscr{A}_L$ and $\mathscr{C}_L$.
\item In $\acute{E}_k$: (o8) and the integrations over $y$ with the factors $-i b_k^2$, $1/y$, $-i$.
\item In $E^{d}_k$, $E_k^{\mathrm{det}}$, ${\delta E^{\mathrm{det}}}_k$: $B\mapsto-B$, $B\mapsto\overline{B}$, complex conjugation, the mean of four numbers, and differences.
\end{itemize}
\end{named}

\begin{named}{Remark}
Proposition~25.2 corresponds to the paragraph of \cite[Definition~5.5(2)]{QEDI} on objects attached to a polymer;
to the forms $\langle\bar{\psi},\psi\rangle_{\square}$ of \cite[Proposition~17.24(b)]{QEDI}, to \cite[Proposition~17.24(d)]{QEDI} as far as the measures are concerned, and to the last sentence of \cite[Proposition~17.24]{QEDI};
to the assertions on $\boldsymbol{\Theta}_{X,\mathscr{T}}$, on $S_L$, on the traces and on the moments of $W$ in Steps~0, 1, 2 and~4 of the proof of \cite[Proposition~19.12]{QEDI};
and to \cite[Remark~13.16]{QEDI}.
The comparison with these places is as follows.
\begin{itemize}
\item[(1)] The notion of Section~I.1.9 for objects attached to a polymer agrees with the notion of the last paragraph of \cite[Definition~5.5(2)]{QEDI}.
Both concern a family $O_{\ell}(X;\mathrm{data})$ with a range $\mathrm{rg}$, indexed by the sides of the class and by the polymers $X$ that have a lift $\tilde{X}$ with
\[
\text{sup-diameter}\left(\tilde{X}\right)+2\,\mathrm{rg}+M<\ell ;
\]
both require one map $\Phi$, the same for all sides of the class, defined on the non-empty finite unions $\tilde{X}$ of $M$-cubes of $\mathbb{L}_k$ that are connected through faces, together with the data on the $\mathrm{rg}$-neighbourhood of $\tilde{X}$, with the covariance
\[
\Phi\left(\tilde{X}+v,T_v\,\mathrm{data}\right)=T_v\,\Phi\left(\tilde{X},\mathrm{data}\right)\qquad\text{for }v\in LM\mathbb{Z}^3;
\]
and for $S_L$ both match a translation by $v$ at the level $k$ with the translation by $v/L$ at the level $k+1$.
The two notions differ in the following respects.
Elements and homomorphisms are compared there through the labelling of the generators of the torus by the indices contained in $\tilde{X}^{+\mathrm{rg}}$, and here through the covering algebras $\mathrm{Grc}_k(\tilde{U})$ and the homomorphisms $\mathrm{cov}_{\ell}$ of Section~G3.5.0 (at the level $k+1$, with the side $\ell/L$, by the homomorphisms $\mathrm{cov}_{\ell/L}$).
Linear functionals on kernels and integrations are admitted here as objects, and elements of the Grassmann algebra as data (Section~I.1.9 together with (G3.r1)--(G3.r3)).
The index sets contain the fermion sites of $\mathbb{L}^{\mathrm{f}}_k$ and the pairs of $\mathrm{Pzt}_k$.
\item[(2)] The forms are objects attached to a polymer that has a lift; they are not operators restricted from a periodic extension.
\item[(3)] The pairing $\langle\bar{\Psi},O\Psi\rangle$ of (o2) is one of the objects.
\item[(4)] The chain of the adjustment consists of the fermion sites $(z_i,g)$.
\item[(5)] At $k=0$ the moments of $W$ contain the number $48$, which is the same on every torus (see (I.L0.5)), and the traces are traces of operators on $\mathfrak{S}^{\mathrm{ch}}(1,\ell)$.
\item[(6)] The value is independent of the choice of the lift, and this is part of the statement.
The same holds in the last paragraph of \cite[Definition~5.5(2)]{QEDI}, where the identity that defines the notion is required for every lift of $X$, in \cite[Remark~13.16]{QEDI}, and in Step~1 of the proof of \cite[Proposition~19.12]{QEDI};
on this point there is no difference.
\item[(7)] In (o8) the bonds are specified on which the kernels (K2)--(K4) of Proposition~25.1 read the bond function $\tilde{B}$:
at indices in (over) $\tilde{X}$ these kernels read $\tilde{B}$ on bonds contained in $\tilde{X}$, and $F_0(\tilde{B})\left(x,x^{(3)}\right)$ reads it on the bond $(c,c_3)$, both of whose ends lie in $\tilde{X}$.
Step~2 of the proof of \cite[Proposition~19.12]{QEDI} specifies the bonds read by the corresponding kernels in the same way, the entries of $P(\tilde{B})$ and of $B_{k,y}(\tilde{B})$ reading $\tilde{B}$ on the fine bonds of the paths $\boldsymbol{\Gamma}^{\pi}$ from the centre of an $L$-block to the two arguments, and the entry $F_0(\tilde{B})(x,x')$ reading it, besides, on the bond $(x,x')$;
it draws the same conclusion: at indices in a polymer $X$ all these bonds have both ends in $X$.
What remains different is the additional bond read by $F_0(\tilde{B})$:
here the bond $(c,c_3)$, read by the entries $F_0(\tilde{B})\left(x,x^{(3)}\right)$;
there the bonds $(x,x\pm e_{\mu})$, read by the entries $F_0(\tilde{B})(x,x\pm e_{\mu})$ through the term $D_{\tilde{B}}$.
\end{itemize}
\end{named}

\subsubsection*{Proposition~25.3 (consequence for two tori)}

\begin{named}{Proposition 25.3}
Let the data be as in Section~G3.5.0, let $\ell_1,\ell_2\in\mathrm{Cl}$, and put $T_i:=\mathbb{L}_k/\ell_i\mathbb{Z}^3$ for $i=1,2$, with the same $k$, $\mu$ and $e_k$ on both tori.
\begin{itemize}
\item[(a)] (Rows.)
Let $\mathrm{Op}$ be one of the operators in the table of Proposition~25.1, let $\mathrm{rg}$ be its range, and let $t$ be a row index.
For $i=1,2$ let $B_i$ be a complex bond function on $T_i$ and $Y_i$ a union of $M$-cubes of $T_i$, and assume that the periodic extensions $\left(\tilde{B}_1,\tilde{Y}_1\right)$ and $\left(\tilde{B}_2,\tilde{Y}_2\right)$ coincide on the bonds, respectively at the sites, that lie within the distance $\mathrm{rg}$ of $\operatorname{supp}t$.
Then the rows of $\mathrm{Op}_{\mathbb{L}}\left(\tilde{B}_1,\tilde{Y}_1\right)$ and of $\mathrm{Op}_{\mathbb{L}}\left(\tilde{B}_2,\tilde{Y}_2\right)$ at $t$ coincide.
If, moreover, $t$ and $t'$ are contained in a set $\tilde{U}$ with
\[
\text{sup-diameter}\left(\tilde{U}\right)+\mathrm{rg}<\ell_i\qquad(i=1,2),
\]
then the kernel of $\mathrm{Op}(B_1,Y_1)$ on $T_1$ and the kernel of $\mathrm{Op}(B_2,Y_2)$ on $T_2$ at $(t,t')$ are both equal to $\mathrm{Op}_{\mathbb{L}}\left(\tilde{B}_1,\tilde{Y}_1\right)(t,t')$.
\item[(b)] (The operators of Definition~24.D near a polymer.)
Let $\tilde{\mathbf{Y}}$ be a finite non-empty union of $M$-cubes of $\mathbb{L}_k$ that is connected through faces, and assume
\[
\text{sup-diameter}\left(\tilde{\mathbf{Y}}\right)+3M<\ell_i\qquad(i=1,2).
\]
Let $\mathbf{Y}_i$ be the image of $\tilde{\mathbf{Y}}$ in $T_i$.
For $i=1,2$ let $B_i$ be a complex bond function on $T_i$, and let $B_1$ and $B_2$ have equal lifts on the bonds that meet $\tilde{\mathbf{Y}}$, that is, their compositions with the covering maps of $\mathbb{L}_k$ onto $T_1$ and onto $T_2$ coincide on these bonds.
Then the following assertions hold, all indices being taken in $\tilde{\mathbf{Y}}^{+\eta_k}$.
\begin{itemize}
\item[(b1)] The polymers contained in $\mathbf{Y}_i$ are the images of the non-empty connected unions $\tilde{X}$ of $M$-cubes contained in $\tilde{\mathbf{Y}}$.
The pairs of $\mathrm{Pz}_k(T_i)$ that are admissible for the image of such a set $\tilde{X}$ are the images of the pairs of $\mathrm{Pzt}_k$ that are admissible for $\tilde{X}$, and each of the latter has for $\tilde{X}$ the kind that its image has for the image of $\tilde{X}$.
Here the two kinds of Section~I.1.7 are read on the covering lattice: for a union $\tilde{X}$ of $M$-cubes of $\mathbb{L}_k$, a pair $(x,x')$ of $\mathrm{Pzt}_k$ is of kind~(I) for $\tilde{X}$ if $c$, $c'$ and all sites of the six paths $\boldsymbol{\Gamma}^{\pi}(c,c')$ lie in $\tilde{X}$ (for $c=c'$: if $c$ lies in $\tilde{X}$);
it is of kind~(II) for $\tilde{X}$ if $x'=x^{(3)}$ and $c$ or $c_{3}$ lies in $\tilde{X}$;
and it is admissible for $\tilde{X}$ if it is of one of the two kinds.
A site at distance at most $\eta_k$ from $\tilde{\mathbf{Y}}$ lies in the periodic extension of $\mathbf{Y}_i$ if and only if it lies in $\tilde{\mathbf{Y}}$.
\item[(b2)] At these pairs the rows of $\delta_{\alpha,B_1}$ and of $\delta_{\alpha,B_2}$ coincide, and so do the rows of $\delta_{\alpha,-B_1}$ and of $\delta_{\alpha,-B_2}$.
\item[(b3)] Let $k\ge1$.
For $j=1,2$, the rows of $T^{(j)}(B_1,\mathbf{Y}_1)$ and of $T^{(j)}(B_2,\mathbf{Y}_2)$ at the unit sites of $\tilde{\mathbf{Y}}$ coincide, and so do the rows of $\bar{T}^{(j)}(B_1,\mathbf{Y}_1)$ and of $\bar{T}^{(j)}(B_2,\mathbf{Y}_2)$ at these sites.
Consequently, with $\tilde{B}$ the common lift of $B_1$ and $B_2$ on the bonds that meet $\tilde{\mathbf{Y}}$,
\[
\mathrm{cov}_{\ell_1}^{-1}\circ\mathrm{sub}_{\mathbf{Y}_1,B_1}\circ\mathrm{cov}_{\ell_1}
=\mathrm{cov}_{\ell_2}^{-1}\circ\mathrm{sub}_{\mathbf{Y}_2,B_2}\circ\mathrm{cov}_{\ell_2}
=\mathrm{sub}_{\tilde{\mathbf{Y}},\tilde{B}}
\]
on the subalgebra on which $\mathrm{sub}_{\tilde{\mathbf{Y}},\tilde{B}}$ is defined in (o5).
\item[(b4)] The entries $C(b',\Upsilon)$ with $b'$ and $\Upsilon$ contained in $\tilde{\mathbf{Y}}$ are the same on $T_1$ and on $T_2$.
\item[(b5)] Let $\tilde{X}$ be contained in $\tilde{\mathbf{Y}}$.
The objects (o1)--(o4), (o6), the set $\bar{X}^{L}$ of (o7), and the objects (o8), (o9) of Proposition~25.2, formed on $T_1$ and on $T_2$ at data that have equal lifts (the data composed with the covering maps in the sense of (G3.r2) coincide), are, in the sense of Proposition~25.2, the images under $\mathrm{cov}_{\ell_1}$ and under $\mathrm{cov}_{\ell_2}$ of one and the same object on the covering lattice; where these objects are numbers, the two numbers are equal.
If $p_{0,k}>0$, then in the coordinates $\mathsf{Z}(\Upsilon)$ the measure $\hat{\mu}$ is, on each of the two tori, the product over the basis of that torus of copies of one and the same probability measure on $\mathbb{R}$, and $\delta\varepsilon^{0}_{k}$ is the same number on both tori (see (o10)).

For the homomorphism $S_L$ of (o7) and a finite set $\tilde{U}$ contained in $\tilde{\mathbf{Y}}^{+\eta_k}$ one has, for $i=1,2$,
\[
\mathrm{cov}_{\ell_i/L}\circ S_L=S_L\circ\mathrm{cov}_{\ell_i}
\]
on the subalgebra of $\mathrm{Grc}_k(\tilde{U})$ generated by $\psi$ and $\chi$.
Here $\mathrm{cov}_{\ell_i/L}$ is the map of Section~G3.5.0 when $\text{sup-diameter}(\tilde{U})+L\cdot M<\ell_i$;
under the condition $\text{sup-diameter}(\tilde{\mathbf{Y}})+3M<\ell_i$ alone, the identity holds with the injective homomorphisms $\mathrm{cov}_{\ell_i/L}$ of Section~G3.5.0.

For the operation $L$ of (o7), let $\ell_1/L$ and $\ell_2/L$ lie in $\mathrm{Cl}$, and let $\tilde{Y}'$ be a finite non-empty union of $M$-cubes of $\mathbb{L}_{k+1}$ that is connected through faces and satisfies $L\tilde{Y}'\subset\tilde{\mathbf{Y}}$.
For $i=1,2$ let $Y'_i$ be the image of $\tilde{Y}'$ in $\mathbb{L}_{k+1}/(\ell_i/L)\mathbb{Z}^3$, let $\mathrm{Ef}^{(i)}$ be a family of even elements of $\mathrm{Gr}[\psi,\chi]$ of level $k$ on $T_i$ that has the property (k2) of Section~I.5.4, and let $(\mathscr{A}_i,\mathscr{C}_i)$ be a pair of bond functions on $\mathbb{L}_{k+1}/(\ell_i/L)\mathbb{Z}^3$ for which the elements $\mathrm{Ef}^{(i)}\left(X;(\mathscr{A}_i)_L,(\mathscr{C}_i)_L\right)$ are defined.
For $\tilde{X}\in\mathrm{Ix}(\tilde{Y}')$ let $X_i$ denote the image of $\tilde{X}$ in $T_i$.
If the elements
\[
\mathrm{cov}_{\ell_i}^{-1}\left(\mathrm{Ef}^{(i)}\left(X_i;(\mathscr{A}_i)_L,(\mathscr{C}_i)_L\right)\right),\qquad\tilde{X}\in\mathrm{Ix}(\tilde{Y}'),
\]
do not depend on $i$, then
\[
\left(L\,\mathrm{Ef}^{(i)}\right)\left(Y'_i;\mathscr{A}_i,\mathscr{C}_i\right)=\mathrm{cov}_{\ell_i/L}\left(\Phi_L\right)\qquad(i=1,2),
\]
with one index set $\mathrm{Ix}(\tilde{Y}')$ and one and the same element $\Phi_L$ for $i=1$ and for $i=2$.
Here $\mathrm{cov}_{\ell_i/L}$ is the map of Section~G3.5.0 at the level $k+1$ if $\text{sup-diameter}(\tilde{Y}')+2\eta_{k+1}+M<\ell_i/L$, and the injective homomorphism of Section~G3.5.0 otherwise.

Let the setting described at the end of (o8) be given on each of the two tori, with a family (p5), written $S^{(i)}_{k,y}$, on $T_i$ for $i=1,2$;
let $\tilde{X}$ be contained in $\tilde{\mathbf{Y}}$, let $X_i$ be the image of $\tilde{X}$ in $T_i$, and let every connection on $T_i$ that occurs below lie in $\tilde{G}_k(T_i)$.
For a pair $(B'_1,B'_2)$ of such connections consider the two conditions:
(i) $B'_1$ and $B'_2$ have equal lifts on the bonds contained in $\tilde{X}$;
(ii) for every $y\ge0$ the kernels $S^{(i)}_{k,y}(X_i,B'_i)(x,x')$, $x$ and $x'$ over $\tilde{X}$, do not depend on $i$.
By (o8), $\acute{E}_k(X_i,B'_i)$ is the value of one functional at the lift of $B'_i$ on the bonds contained in $\tilde{X}$ and at the lifts of these kernels;
hence (i) and (ii) give $\acute{E}_k(X_1,B'_1)=\acute{E}_k(X_2,B'_2)$.
At the five pairs $(B_1,B_2)$, $(-B_1,-B_2)$, $\left(\overline{B_1},\overline{B_2}\right)$, $\left(-\overline{B_1},-\overline{B_2}\right)$ and $(0,0)$ condition (i) holds by the hypothesis of (b), since a bond contained in $\tilde{X}$ meets $\tilde{\mathbf{Y}}$.
Consequently $\acute{E}_k(X_i,B_i)$, $E^{d}_k(X_i,B_i)$ and $E_k^{\mathrm{det}}(X_i,B_i)$ do not depend on $i$ if (ii) holds at each pair of connections at which $\acute{E}_k$ is evaluated in them:
at $(B_1,B_2)$ for $\acute{E}_k$;
at the first four of these pairs for $E^{d}_k$;
at all five for $E_k^{\mathrm{det}}$.
Let further $(A_i,C_i)$ be a pair of connections on $T_i$ for $i=1,2$, so that
\[
{\delta E^{\mathrm{det}}}_k(X_i;A_i,C_i)=E^{d}_k(X_i,A_i+C_i)-E^{d}_k(X_i,C_i)
\]
is evaluated at $\pm(A_i+C_i)$, $\pm\overline{A_i+C_i}$, $\pm C_i$ and $\pm\overline{C_i}$.
Assume that each of the pairs $(A_1+C_1,A_2+C_2)$ and $(C_1,C_2)$ has equal lifts on the bonds contained in $\tilde{X}$, and that (ii) holds at the pairs $(B'_1,B'_2)$, $(-B'_1,-B'_2)$, $\left(\overline{B'_1},\overline{B'_2}\right)$ and $\left(-\overline{B'_1},-\overline{B'_2}\right)$ for $B'_i=A_i+C_i$ and for $B'_i=C_i$.
Then (i) holds at these eight pairs as well, so that $E^{d}_k(X_i,A_i+C_i)$ and $E^{d}_k(X_i,C_i)$ do not depend on $i$, and ${\delta E^{\mathrm{det}}}_k(X_1;A_1,C_1)={\delta E^{\mathrm{det}}}_k(X_2;A_2,C_2)$.

The sums that define $\varepsilon$ and $m$ in Section~I.5.6 behave as follows, for all $\ell_1,\ell_2\in\mathrm{Cl}$ and without reference to the set $\tilde{\mathbf{Y}}$.
Let $\tilde{\square}$ be an $M$-cube of $\mathbb{L}_k$, and let $\mathrm{Sm}(\tilde{\square})$ be the set of the finite non-empty unions $\tilde{X}$ of $M$-cubes of $\mathbb{L}_k$ that are connected through faces, contain $\tilde{\square}$ and satisfy $d_M(\tilde{X})\le L$.
For $i=1,2$ let $\mathrm{Ef}^{(i)}$ be a family on $T_i$ with (k2).
Assume that for every $\tilde{X}\in\mathrm{Sm}(\tilde{\square})$, with $X_i$ the image of $\tilde{X}$ in $T_i$, neither the number $\mathrm{Ef}^{(i)}_{00}(X_i;0,0)$ nor the kernel $\mathrm{Ef}^{(i)}_{20}(X_i;0,0;x,x')$, $x$ and $x'$ over $\tilde{X}$, depends on $i$.
Then $\mathrm{al}_{0}\left(\mathrm{Ef}^{(i)},X_i\right)$ and $\mathrm{al}_{2}\left(\mathrm{Ef}^{(i)},X_i\right)$, $\tilde{X}\in\mathrm{Sm}(\tilde{\square})$, do not depend on $i$, and neither do the sums
\[
-\sum_{X}\mathrm{al}_{0}\left(\mathrm{Ef}^{(i)},X\right)
\qquad\text{and}\qquad
-\sum_{X}\mathrm{al}_{2}\left(\mathrm{Ef}^{(i)},X\right),
\]
in which $X$ ranges over the small polymers of $T_i$ that contain the image of $\tilde{\square}$.
These sums are the right sides of the display of Section~I.5.6 that defines $\varepsilon$ and $m$, formed with the family $\mathrm{Ef}^{(i)}$ and with the cube $\square_i$, the image of $\tilde{\square}$ in $T_i$;
in the notation of Proposition~16.3 and of Section~I.5.11 they are the number $\varepsilon\left(\mathrm{Ef}^{(i)},\square_i\right)$ and the $4\times4$ matrix $m\left(\mathrm{Ef}^{(i)},\square_i\right)$.
If only (k2) is assumed, the matrix $m\left(\mathrm{Ef}^{(i)},\square_i\right)$ need not be a multiple of $\mathbf{1}_4$, and both quantities may depend on the cube.
Where Proposition~16.3\,(a) applies to $\mathrm{Ef}^{(1)}$ and to $\mathrm{Ef}^{(2)}$ (for families without slot, Proposition~48.2\,(a)), one has $\varepsilon\left(\mathrm{Ef}^{(i)},\square_i\right)=\varepsilon\left(\mathrm{Ef}^{(i)}\right)$ and $m\left(\mathrm{Ef}^{(i)},\square_i\right)=m\left(\mathrm{Ef}^{(i)}\right)\cdot\mathbf{1}_4$, and the assertion on the two sums reads
\[
\varepsilon\left(\mathrm{Ef}^{(1)}\right)=\varepsilon\left(\mathrm{Ef}^{(2)}\right),\qquad m\left(\mathrm{Ef}^{(1)}\right)=m\left(\mathrm{Ef}^{(2)}\right).
\]
\end{itemize}
\item[(c)] (Periodisation.)
Let $\ell_2=\mathrm{Jr}\cdot\ell_1$ with $\mathrm{Jr}$ odd, and let $\mathrm{pj}\colon T_2\to T_1$ and $\mathrm{pj}^{*}$ be as in Section~I.6.3.
Then
\[
\mathrm{Op}\left(\mathrm{pj}^{*}B,\mathrm{pj}^{-1}Y\right)\circ\mathrm{pj}^{*}=\mathrm{pj}^{*}\circ\mathrm{Op}(B,Y)
\]
for every operator $\mathrm{Op}$ of the table of Proposition~25.1, every complex bond function $B$ on $T_1$ and every union $Y$ of $M$-cubes of $T_1$.
For the operator of the row no.~14 of that table, whose source consists of objects of level $k+1$, the map $\mathrm{pj}^{*}$ to the right of $\mathrm{Op}$ is the inclusion of the periodic objects (bond functions), respectively of the antiperiodic objects (spinors and functions of pairs), of level $k+1$ and side $\ell_1/L$ into those of side $\ell_2/L=\mathrm{Jr}\cdot\ell_1/L$;
this is the definition of Section~I.6.3, read at the level $k+1$ with these two sides, which need not lie in $\mathrm{Cl}$.
Since $\mathrm{Jr}$ is odd, an object that is antiperiodic for the side $\ell_1/L$ is antiperiodic for the side $\ell_2/L$.
\end{itemize}
\end{named}

\begin{named}{Remark}
Clause~(c) is the identity
\[
\mathrm{Op}\left(\mathrm{pj}^{*}B,\mathrm{pj}^{-1}Y\right)\circ\mathrm{pj}^{*}=\mathrm{pj}^{*}\circ\mathrm{Op}(B,Y)
\]
of the second paragraph of \cite[Definition~5.5(2)]{QEDI}, which concerns the pair $\left(\mathrm{pj}^{*}B,\mathrm{pj}^{-1}Y\right)$ only.
Clause~(a) is the comparison of two tori in the third paragraph of that definition, formulated for the index sets of Section~I.1.9 and for the operators of the table of Proposition~25.1:
in both, the periodic extensions of the data $(B_i,Y_i)$ are required to agree only on the bonds, respectively at the sites, within the distance $\mathrm{rg}$ of the support of the row index, and the kernels on the two tori are compared under the same margin $\text{sup-diameter}\left(\tilde{U}\right)+\mathrm{rg}<\ell_i$, $i=1,2$.
Clause~(b) compares the two tori for the operators of Definition~24.D near a polymer under a hypothesis weaker than the one which (a) imposes at the same rows:
$B_1$ and $B_2$ are required to have equal lifts only on the bonds that meet $\tilde{\mathbf{Y}}$.
Within the distance $\mathrm{rg}$ of the support of a row index of (b2) or (b3), where $\mathrm{rg}=M$ for $\delta_{\alpha,B}$ (no.~11 of the table of Proposition~25.1) and $\mathrm{rg}=(1+\eta_k)/2$ for $T^{(j)}$ and $\bar{T}^{(j)}$ (nos.~9 and~10), there may lie bonds that do not meet $\tilde{\mathbf{Y}}$;
on such bonds the hypothesis of (b) says nothing, so that the hypothesis of (a) does not follow from it.
In (b2) and (b3) the rows are compared through the bonds and the sites which they read according to the table of Proposition~25.1, in the place of their $\mathrm{rg}$-neighbourhoods, and in (b5) the objects are compared through Proposition~25.2.
The clauses (b1)--(b4), together with the assertion of (b5) on the forms of (o1), correspond to the four assertions that follow the heading \emph{Exact operators and forms} in Step~3 of the proof of \cite[Proposition~19.12]{QEDI};
the remaining assertions of (b5) are Proposition~25.2 read on two tori (for their counterparts see the Remark after Proposition~25.2).
Clause~(c) alone does not suffice for (b):
the bond function $B_2=\mathscr{T}_2+t\cdot\mathscr{Z}^{(2)}$ on $T_2$, formed with the total connection and the weakened fluctuation field of $T_2$, is not of the form $\mathrm{pj}^{*}B_1$, and the set $\mathbf{Y}_2$, the image of a single copy of $\tilde{\mathbf{Y}}$, is not $\mathrm{pj}^{-1}\mathbf{Y}_1$.
The same data occur in Step~3 of the proof of \cite[Proposition~19.12]{QEDI}.
There $B_1$ and $B_2$ have equal lifts on the bonds contained in $\tilde{\mathbf{Y}}^{\sim(1+d_U)}$, the enlargement of $\tilde{\mathbf{Y}}$ by $1+d_U$ layers of $M$-cubes, where $d_U=2L+1$, and the rows of the first, second and fourth of the four assertions are compared by the comparison of two tori of \cite[Definition~5.5(2)]{QEDI}, with the range $M$.
This is the difference from (b), which requires equal lifts only on the bonds that meet $\tilde{\mathbf{Y}}$.
\end{named}

\begin{named}{Remark}
Let $\ell_2=\mathrm{Jr}\cdot\ell_1$ with $\mathrm{Jr}$ odd, as for the two tori of Propositions~27.1 and~27.2, and let $\tilde{\mathbf{Y}}$ be contained in a set that is admissible in the sense of Proposition~27.1.
The largest side of the box of an admissible set is at most $\ell_1-14\,d_U\cdot M$, where $d_U=2L+1$.
Hence $\text{sup-diameter}(\tilde{\mathbf{Y}})+3M<\ell_i$, and $\text{sup-diameter}(\tilde{U})+L\cdot M<\ell_i$ for every finite set $\tilde{U}$ contained in $\tilde{\mathbf{Y}}^{+\eta_k}$, $i=1,2$;
so clause~(b) of Proposition~25.3 applies to $\tilde{\mathbf{Y}}$, and the identity for $S_L$ in (b5) holds with the maps $\mathrm{cov}_{\ell_i/L}$ of Section~G3.5.0.
The homomorphism $\mathrm{idn}$ of Proposition~27.2\,(D2) is $\mathrm{cov}_{\ell_2}\circ\mathrm{cov}_{\ell_1}^{-1}$.
\end{named}

\subsection*{24 (continued). One renormalisation step without large fields, with external connection, every level $k\ge0$: the statements}

\subsubsection*{Proposition 24.1}

\begin{named}{Proposition 24.1}[One step without large fields, with an external slot]
Assume (ST1)$^{\mathrm{ch}}$ (Section~G3.4.0), and let the step map be the one of Definition~24.D.
\begin{itemize}
\item[(a)]
For every $(m,\mathscr{E})\in\mathrm{Ball}^{x}_{k}(4)$, with $\mathscr{E}=0$ when $k=0$, and for all $A\in\tilde{R}_k^{-}$ and $\mathscr{C}\in\mathfrak{X}_k$, each display of Definition~24.D is defined, in the order fixed at the end of that definition;
the series (8.4) of Definition~24.D converge absolutely in the norm $\|\cdot\|_{h_k/4}$;
and for every $X\in D_k$
\[
\left\|E^{\#,x}_k(m,\mathscr{E})(X;A,\mathscr{C})\right\|_{h_k/4}
\le14\,c_X^2\cdot K_0^2\cdot\exp(3\kappa)\cdot\mathbf{H}_0\cdot\exp\left(-(\kappa-5\kappa_0-4)\cdot d_M(X)\right).
\]
The family $E^{*,x}_k(m,\mathscr{E})$ is an element of $K^{x,C}_{k+1}=K^{x,C}_{k+1}(\ell/L)$, and
\[
\left\|E^{*,x}_k(m,\mathscr{E})\right\|^{x}_{k+1}\le e_k^{\frac{1}{4}-6\varepsilon}.
\]
Let $(m_a,\mathscr{E}_a)$ and $(m_b,\mathscr{E}_b)$ be elements of $\mathbb{C}\times K^{x,C}_k$, with $\mathscr{E}_a=\mathscr{E}_b=0$ when $k=0$.
The complex numbers $t$ for which $\left(m_a+t\cdot m_b,\,\mathscr{E}_a+t\cdot\mathscr{E}_b\right)$ belongs to $\mathrm{Ball}^{x}_{k}(4)$ form a closed set, and the map
\[
t\longmapsto E^{*,x}_k\left(m_a+t\cdot m_b,\,\mathscr{E}_a+t\cdot\mathscr{E}_b\right)\in K^{x,C}_{k+1}
\]
is continuous on this set and analytic in its interior:
near every interior point it is the sum of a power series in $t$ with coefficients in $K^{x,C}_{k+1}$ that converges in $\|\cdot\|^{x}_{k+1}$.
If $m$ is real and $\mathscr{E}\in K^{x}_{k}$, then $E^{*,x}_k(m,\mathscr{E})$ lies in $K^{x}_{k+1}$.
Finally,
\[
0\le\delta\varepsilon^{0}_{k}\le e_k^{7}.
\]
\item[(b)]
In the reading~(o) of Definition~24.D\,(D7) the first assertion of (a) holds:
for every $(m,E)\in\mathrm{Ball}_{k}(4)$, with $E=0$ when $k=0$, and every $A\in\tilde{R}_k^{-}$, all displays of Definition~24.D, read according to (D7), are defined, in the order fixed there,
and the series (8.4) converge absolutely in $\|\cdot\|_{h_k/4}$, with the bound displayed in (a).
The map $E^{*}_k$ of the reading~(o) sends $\mathrm{Ball}_{k}(4)$ into $K^{C}_{k+1}$, and the pairs $(m,E)\in\mathrm{Ball}_{k}(4)$ with real $m$ and $E\in K_k$ into $K_{k+1}$;
it satisfies
\[
\left\|E^{*}_k(m,E)\right\|_{k+1}\le e_k^{\frac{1}{4}-6\varepsilon},
\]
and it is analytic in the sense of (a), with $\mathbb{C}\times K^{C}_{k}$, $\mathrm{Ball}_{k}(4)$, $K^{C}_{k+1}$ and $\|\cdot\|_{k+1}$ in the places of $\mathbb{C}\times K^{x,C}_{k}$, $\mathrm{Ball}^{x}_{k}(4)$, $K^{x,C}_{k+1}$ and $\|\cdot\|^{x}_{k+1}$.
For $(m,\mathscr{E})\in\mathrm{Ball}^{x}_{k}(4)$ the pair $(m,\mathscr{E}|_0)$ lies in $\mathrm{Ball}_{k}(4)$, and
\[
E^{*,x}_k(m,\mathscr{E})|_0=E^{*}_k\left(m,\mathscr{E}|_0\right).
\]
\item[(c)] \emph{Identity of the step.}
Let $\varepsilon^{\mathrm{v}}$ and $m$ be real numbers with $|m|\le w^{\mathrm{m}}_k$, and let $\mathscr{E}\in K^{x}_{k}$ with $\|\mathscr{E}\|^{x}_{k}\le w^{\mathrm{E}}_k$, where $\mathscr{E}=0$ if $k=0$.
Let $\mathfrak{P}_k=\left\{\mathfrak{P}_k(X,\cdot\,):X\in D_k\right\}$ be an arbitrary family of complex functions on $\mathfrak{X}_k$;
let $A'\in\mathrm{Dom}_{k+1}$, with an arbitrary harmonic part, and $\mathscr{C}'\in\mathfrak{X}_{k+1}$.
Let $(\varepsilon^{\mathrm{v},+},m^{+},\mathscr{E}^{+})$ and $\mathfrak{P}_{k+1}$ be given by (8.6)$^{\mathrm{ch}}$ of Definition~24.D,
and form $\mathscr{F}^{x}_{k}$ with the family $\mathfrak{P}_k$ and $\mathscr{F}^{x}_{k+1}$ with the family $\mathfrak{P}_{k+1}$ (Section~I.5.7).
Then one has the identity of the step, denoted (CV.4x)$^{\mathrm{ch}}$, and the identity for $y^{x}$, denoted (2.4x)$^{\mathrm{ch}}$:
\begin{align}
\mathbf{T}^{x}_{k}\mathscr{F}^{x}_{k}[\varepsilon^{\mathrm{v}},m,\mathscr{E}](A',\mathscr{C}';\,\cdot\,)
&=\mathscr{F}^{x}_{k+1}[\varepsilon^{\mathrm{v},+},m^{+},\mathscr{E}^{+}](A',\mathscr{C}';\,\cdot\,),\tag*{(CV.4x)$^{\mathrm{ch}}$}\\
y^{x}_{k+1}[\varepsilon^{\mathrm{v},+},m^{+},\mathscr{E}^{+}](A';\mathscr{C}')
&=\int_{\mathrm{Box}_k}\mathrm{d}\mu_I(\zeta)\,y^{x}_{k}[\varepsilon^{\mathrm{v}},m,\mathscr{E}]\left(A^{\min}_{k}(A')+\mathrm{Sq}_{k}\,\zeta;\,\mathscr{C}'_L\right).\tag*{(2.4x)$^{\mathrm{ch}}$}
\end{align}
At level $0$ the first identity takes the form
\[
\mathbf{T}^{x}_{0}\mathscr{F}^{x}_{0}[\varepsilon^{\mathrm{v}},m,0](A',\mathscr{C}';\,\cdot\,)
=\mathscr{F}^{x}_{1}\left[L^3\cdot\left(\varepsilon^{\mathrm{v}}+\delta\varepsilon^{0}_{0}\right),\,L\cdot m,\,L\,{\delta E^{\mathrm{det}}}_0+E^{*,x}_0(m,0)\right](A',\mathscr{C}';\,\cdot\,);
\]
its left side is an integral over $192\,n_0$ pairs of generators $\Psi$ on $T^{\mathrm{f}}_0$, and its right side is an element of $\mathrm{Gr}[\Psi']$ with $4\,n_1$ pairs of generators.
For a family with $\mathfrak{P}_k(\cdot\,,0)=0$, the two identities at $\mathscr{C}'=0$ are the identities for the objects $\mathscr{F}_k$, $y_k$, $\mathbf{T}_k$ without slot, with the data $(\varepsilon^{\mathrm{v}},m,\mathscr{E}|_0)$ and their image $(\varepsilon^{\mathrm{v},+},m^{+},\mathscr{E}^{+}|_0)$;
moreover $\mathfrak{P}_{k+1}(\cdot\,,0)=0$ in this case.
\item[(d)] \emph{The full step.}
In the situation of (c) the family $\mathscr{E}^{+}$ lies in $K^{x}_{k+1}$,
\[
\mathscr{E}^{+}|_0=L R\,(\mathscr{E}|_0)+L\,E_k^{\mathrm{det}}+E^{*}_k\left(m,\mathscr{E}|_0\right),
\]
and
\begin{gather*}
\left\|\mathscr{E}^{+}\right\|^{x}_{k+1}\le c_{\mathrm{fl}}\cdot L^{-\frac{1}{4}+2\varepsilon}\cdot\|\mathscr{E}\|^{x}_{k}+e_{k+1}^{\frac{1}{4}-2\varepsilon}+e_k^{\frac{1}{4}-6\varepsilon},\\
\left|m^{+}-L\cdot m\right|\le L\cdot c_{\mathrm{fl}}\cdot e_k^{\frac{1}{2}}\cdot\|\mathscr{E}\|^{x}_{k},\qquad
\left|\varepsilon^{\mathrm{v},+}-L^3\cdot\varepsilon^{\mathrm{v}}\right|\le L^3\cdot\left(c_{\mathrm{fl}}\cdot\|\mathscr{E}\|^{x}_{k}+e_k^{7}\right).
\end{gather*}
For $k=0$ one has
\[
\left\|\mathscr{E}^{+}\right\|^{x}_{1}\le\left(L^{\frac{1}{2}}\cdot q\right)^{\frac{1}{4}-2\varepsilon}+q^{\frac{1}{4}-6\varepsilon},\qquad
m^{+}=L\cdot m,\qquad
0\le\varepsilon^{\mathrm{v},+}-L^3\cdot\varepsilon^{\mathrm{v}}\le L^3\cdot q^{7}.
\]
\item[(e)] \emph{Dependence on the couplings.}
For $(m_1,\mathscr{E}_1)$ and $(m_2,\mathscr{E}_2)$ in $\mathrm{Ball}^{x}_{k}(1)$,
\[
\left\|E^{*,x}_k(m_1,\mathscr{E}_1)-E^{*,x}_k(m_2,\mathscr{E}_2)\right\|^{x}_{k+1}
\le e_k^{\frac{1}{4}-6\varepsilon}\cdot\left(\frac{|m_1-m_2|}{w^{\mathrm{m}}_k}+\frac{\|\mathscr{E}_1-\mathscr{E}_2\|^{x}_{k}}{w^{\mathrm{E}}_k}\right).
\]
The same bound holds for $E^{*}_k$ and for $(m_1,E_1)$, $(m_2,E_2)$ in $\mathrm{Ball}_{k}(1)$, with the norms $\|\cdot\|_k$ and $\|\cdot\|_{k+1}$.
(For $k=0$ the second of the two quotients in these bounds is absent.)
\item[(f)] \emph{Symmetries.}
Let $(m,\mathscr{E})\in\mathrm{Ball}^{x}_{k}(4)$.
At level $k+1$, and in the terms of Section~I.5.4, the family $E^{*,x}_k(m,\mathscr{E})$ has clause (k3) for the group $\mathbb{G}_{k+1}$ of the side $\ell/L$, which acts on the fermion sites by (I.4.2);
it has clause (k4) in both connections, and clause (k5).
Every monomial that occurs in it has as many barred legs as unbarred legs.
If $m$ is real and $\mathscr{E}\in K^{x}_{k}$, it has clause (k7) as well.
\item[(g)] \emph{Locality, analyticity.}
Let $(m,\mathscr{E})\in\mathrm{Ball}^{x}_{k}(4)$.
At level $k+1$ the family $E^{*,x}_k(m,\mathscr{E})$ has the clauses (k1) and (k2):
its kernels are analytic on $\tilde{R}_{k+1}\times\mathfrak{X}_{k+1}$;
they vanish unless all their sites lie over the polymer and all their pairs are of kind (I) or of kind (II) for it (Section~I.1.7);
and they depend on the two connections through the bonds that meet the polymer.
Let $Y\in D_{k+1}$.
Each H\"older leg of $E^{*,x}_k(m,\mathscr{E})(Y)$ sits at the image, under $(x,x')\mapsto(x/L,x'/L)$, either of a pair at which some $\mathscr{E}^{r}(X;B_1,\mathscr{C})$ with $X\subseteq LY$ and $(B_1,\mathscr{C})\in\tilde{R}_k\times\mathfrak{X}_k$ has a H\"older leg,
or of a pair $(x,x^{(3)})$, $x=(c,g)$, with $c\in LY$;
the pairs of the second sort arise from (8.3a)$^{\mathrm{ch}}$ when $k\ge1$ and from (8.3$_0$)$^{\mathrm{ch}}$ when $k=0$ (Section~I.5.9).
In particular, every H\"older leg of $E^{*,x}_0(m,0)$ is a pair $(x',{x'}^{(3)})$ of $\mathrm{Pz}_1$, of length $d_X=\eta_1$.
The structure of the objects attached to a polymer of level $k$ is described in Proposition~24.3, and the exact operators that occur in them in Propositions~25.1--25.3.
\item[(i)] \emph{Uniformity.}
Apart from numerals, from the norms of the data and from $d_M(X)$, and apart from $\varepsilon$, which enters through exponents only, the bounds displayed in (a)--(e) contain the datum $e_k$ (also by way of $e_{k+1}=L^{\frac{1}{2}}\cdot e_k$, $w^{\mathrm{m}}_k$, $w^{\mathrm{E}}_k$ and $\mathbf{H}_0$) and the constants $L$, $c_{\mathrm{fl}}$, $c_X$, $K_0$, $\kappa$, $\kappa_0$, and nothing else.
With these numbers they hold for all data $(k,\ell,\mu,e_k)$ of (ST1)$^{\mathrm{ch}}$, for every system of primitives with the properties required there, for every harmonic part of $A'$, for every $\mathscr{C}'$ and for all $(m,\mathscr{E})$ in the balls named.
\end{itemize}
\end{named}

\begin{named}{Remark}
Proposition~24.1 corresponds to \cite[Theorem~17.2]{QEDI}, which is \cite[Theorem~5.71]{QEDI}, with the hypothesis \cite[Definition~5.69]{QEDI}, the step map of \cite[Definitions~5.25 and~5.26]{QEDI} and \cite[Section~17.4]{QEDI}, the two cases of \cite[Section~17.1.3]{QEDI}, and \cite[Section~17.11.1]{QEDI}.
Its clauses (a)--(g) correspond to the clauses with the same letters there, the identities (CV.4x)$^{\mathrm{ch}}$ and (2.4x)$^{\mathrm{ch}}$ of (c) to the two identities displayed in \cite[Theorem~17.2(c)]{QEDI}, and its clause (i) to clause (h) there; the letter (h) is not used here.
It differs from these places in the following points.
\begin{itemize}
\item[(1)] The fine fields $\psi$ and $\chi$ live on $T^{\mathrm{f}}_k$ and on $\mathrm{Pz}_k$ (Sections~I.1.3 and~I.1.7), with the measure $\int\mathrm{d}x$ of (I.1.3);
all norms, windows, contours and exponents have the values of \cite{QEDI}.
\item[(2)] The operators $T^{(i)}$, $\bar{T}^{(i)}$ of \cite[Section~17.4.4]{QEDI} are replaced by those of (23.1), that is, of (I.5.4), and the substitution defined there by (8.3a)$^{\mathrm{ch}}$;
for $k=0$ the function $G^{(1)}$ of \cite[Section~17.4.4]{QEDI} is replaced by (8.3$_0$)$^{\mathrm{ch}}$, which is (I.5.5) written with the element $\mathrm{Vx}$ of (23.3) (Propositions~23.1 and~23.2, Section~I.5.9).
\item[(3)] At level $0$ the fields $\Psi$ and $W$ have $192$ components per unit site and the weight $\frac{1}{48}$;
the weight of $W$ in the norms is $\mathsf{w}_W=\sqrt{48}\le7$ (rule (I.N3)$^{\mathrm{w}}$ of Section~I.5.2);
$C_W=28\,C_w^2\cdot c_{\gamma}$, whereas \cite[Lemma~17.11]{QEDI} has $4\,C_w^2\cdot c_{\gamma}$;
condition (A6) of \cite[Section~5.6.2]{QEDI} carries the number $112$ in place of $16$;
and the kernels of $\Gamma_0(s'',A)$ are taken by (I.1.4).
The bounds of the proposition are the same as in \cite{QEDI}:
the first display of clause (a) is the bound of assertion (1) of \cite[Section~17.10.1]{QEDI}, in which $c'_S=14\,c_X^2\cdot K_0^2\cdot\exp(3\kappa)$ is the constant of \cite[Lemma~17.17]{QEDI} and $\kappa_3=\kappa-5\kappa_0-4$ is the number of \cite[Section~17.1.2]{QEDI};
the other bounds are those of \cite[Theorem~17.2(a), (b), (d), (e)]{QEDI}.
\item[(4)] The groups act on the labels of the fermion sites, by (I.4.2), and the kinds of pairs are those of Section~I.1.7;
a pair $(x,x')$ of kind (II) has $x'=x^{(3)}$.
\item[(5)] In \cite[Definition~5.69]{QEDI} the mass lies in $\left(0,\mathfrak{m}_0/L\right]$, where $\mathfrak{m}_0$ denotes the upper end of the mass range.
Here this upper end has the value $1$ (Section~I.6.1; Proposition~14.1\,(v)), so that the mass $\mu$ of (ST1)$^{\mathrm{ch}}$ ranges over $(0,1/L]$.
\item[(6)] The case without slot is the reading~(o) of Definition~24.D\,(D7), in which the pair $(\{0\},\{0\})$ stands in the place of $(\mathfrak{X}_k,\mathfrak{X}_{k+1})$;
this reading agrees with the one by which \cite[Section~17.1.3]{QEDI} defines the case without slot.
In addition, (D7) declares the family $\mathfrak{P}_k$ and the entry $\mathfrak{P}_{k+1}$ of (8.6)$^{\mathrm{ch}}$ void, these being functions of the suppressed variable alone.
In Chapter~24, Sections~2, 3, 5, 6, 8, 9.1--9.4 and~9.8, which treat the case with slot and the reading~(o) together, the sets $\mathfrak{X}_k$ and $\mathfrak{X}_{k+1}$ enter the arguments by way of the following six properties:
both sets contain $0$;
$\tilde{R}_k+\mathfrak{X}_k$ is contained in $\tilde{G}_k$ and in the set on which (G3.T-9) holds;
invariance under $\mathbb{G}_j$, under $\mathscr{C}\mapsto-\mathscr{C}$ and under $\mathscr{C}\mapsto\overline{\mathscr{C}}$, at both levels $j=k$ and $j=k+1$;
the map $\mathscr{C}'\mapsto\mathscr{C}'_L$ sends $\mathfrak{X}_{k+1}$ into $\mathfrak{X}_k$;
the clauses (k1)--(k7) for the data and for the values, together with (G3.T-16);
and openness.
The first five of them are the properties (a)--(e) of \cite[Section~17.1.3]{QEDI}, with (G3.T-9) in the place of the bound of \cite[Proposition~5.46]{QEDI} and (G3.T-16) in the place of \cite[Proposition~5.59]{QEDI};
in particular, property (c) there requires the invariance, as here, at both levels $j=k$ and $j=k+1$.
Openness is the second way in which, according to that section, the two sets enter.
That section says this of \cite[Sections~17.4--17.11]{QEDI} and of \cite[Proposition~17.24]{QEDI} with its proof.
Since no derivative, integral, contour integral or identity theorem in the variable $\mathscr{C}$ occurs in the parts of Chapter~24 named above, openness serves there only to give a meaning to analyticity in $(A,\mathscr{C})$, as in \cite[Section~17.1.3]{QEDI}.
Besides the six properties, these parts apply statements on families with slot and, in the reading~(o), their forms without slot;
among them are Proposition~16.3\,(a), (b), (d), which enters the fifth property and whose form without slot is Proposition~48.2\,(a), (b), (d),
and, for the display (8.6)$^{\mathrm{ch}}$ in Chapter~24, Section~5.4, Proposition~18.5\,(a), whose form without slot is Proposition~48.5\,(a), and Proposition~13\,(d5), whose assertion on $L\,{\delta E^{\mathrm{det}}}_k$ has the assertion on $L\,E_k^{\mathrm{det}}$ as its form without slot.
This corresponds to the paragraph of \cite[Section~17.1.3]{QEDI} on \cite[Propositions~5.59, 5.60, 5.62 and~5.48]{QEDI}, which in the form stated in \cite[Section~17.1.4]{QEDI} include the case without slot.
The pair $(\{0\},\{0\})$ has the first five properties, the fifth with (G3.T-48) in the place of (G3.T-16), and the reading (D7) stands in the place of openness, as in \cite[Section~17.1.3]{QEDI};
with the forms without slot named above, the arguments of these parts are therefore valid in the reading~(o).
Moreover, \cite[Section~17.1.3]{QEDI} draws from its list a principle for every pair of open sets with the properties (a)--(e), whereas Definition~24.D uses only the pairs $(\mathfrak{X}_k,\mathfrak{X}_{k+1})$ and $(\{0\},\{0\})$.
The last sentence of clause~(b), with its display, and the clauses~(c) and~(d) concern the case with slot.
Their proofs, in Chapter~24, Section~9.5 for that sentence and Sections~9.6 and~9.7, apply in addition statements that contain the second slot,
among them Proposition~48.1\,(a), the identities $m(\mathscr{E})=m(\mathscr{E}|_0)$, $\varepsilon(\mathscr{E})=\varepsilon(\mathscr{E}|_0)$ and $(R\mathscr{E})|_0=R(\mathscr{E}|_0)$ of Proposition~16.3\,(a), (d), and Propositions~17.2\,(a), 22.2\,(b), 22.3\,(a)--(c), 22.4\,(a), 18.5\,(b), (c) and~13\,(d6);
likewise, \cite[Section~17.12]{QEDI} concerns data with slot.
\item[(7)] The invariances are those of Proposition~24.2, in the sense of Definition~24.I.
Invariance and covariance relative to a polymer, on a set $\mathrm{Ar}$ of arguments and with a domain in the two connection slots are also the notions of \cite[Section~17.11.1]{QEDI};
the sets $\mathrm{Ar}(\mathbf{Y})$ are defined there by the same conditions on $X$, $P$, $s$, $t$ and $\mathsf{Z}$, the polymer that belongs to each object is chosen in the same way, and the monotonicity in the polymer holds there as here.
On the other hand, $\mathrm{Cg}$ is extended here by the identity on the generators $W$, whereas in \cite[Section~17.11.1]{QEDI} the generators $W$ lie outside the domain of the charge conjugation (T-C).
The rule that, together with a covariant family $v$, also $t\cdot v$ is covariant is part of Definition~24.I.
In $\mathrm{Ar}(\mathbf{Y})$ the pair $(A,\mathscr{C})$ ranges here over the domain of the object considered, and there over $\tilde{R}_k^{-}\times\mathfrak{X}_k$.
The domain of the form $\left\langle\bar{\Psi},M_k(X,A_1+\mathscr{C})\Psi\right\rangle_1$ in the two connection slots is here $\left\{(A_1,\mathscr{C}):A_1+\mathscr{C}\in\tilde{G}_k\right\}$, and there $\tilde{R}_k\times\mathfrak{X}_k$;
Definition~24.I gives the primitives the same domain $\left\{(A_1,\mathscr{C}):A_1+\mathscr{C}\in\tilde{G}_k\right\}$.
\item[(8)] The assertion of \cite[Theorem~5.71(g)]{QEDI} that the objects named there read the connections on the bonds that meet $\mathbf{Y}$, the fourth assertion of the list in that clause, is contained in Proposition~24.3\,(ze).
\item[(9)] Where \cite[Theorem~17.2(h)]{QEDI} states that the constants and the conditions on them depend on none of the data, clause (i) states which quantities the displayed bounds contain; in it $e_k$ is a datum.
The conditions on the constants on which Proposition~24.1 rests are those of Section~G3.6.1;
they enter the ordered list of constants and conditions of the Propositions~14.m.
\item[(10)] Besides the absolute convergence of the series (8.4), clause (a) asserts that each display of Definition~24.D is defined, in the order fixed at the end of that definition;
the hypothesis on the functions $G(\cdot\,,X)$ to which $\mathrm{Loc}$ is applied is stated in Definition~24.D\,(D3).
Since there is no operator $M_0$, the entries of $M_k$ enter Proposition~24.3\,(q5) for $k\ge1$ only.
Clause (g) states both vanishing conditions of clause (k2), the one on the sites and the one on the pairs, as \cite[Theorem~17.2(g)]{QEDI} does;
\cite[Theorem~5.71(g)]{QEDI} states only the condition on the pairs.
For $(\mathscr{A},\mathscr{C})\in\tilde{R}_{k+1}\times\mathfrak{X}_{k+1}$ the bond functions $\mathscr{A}_L+\mathscr{C}_L$ and $\mathscr{C}_L$ lie in $\tilde{G}_k$, as stated in Definition~24.D\,(D6).
\end{itemize}
\end{named}

\subsubsection*{Proposition 24.2 (invariances of $E^{\#,x}_k$ at level $k$)}

\begin{named}{Proposition 24.2}
Assume (ST1)$^{\mathrm{ch}}$.
Let $(m,\mathscr{E})\in\mathrm{Ball}^{x}_{k}(4)$, $X\in D_k$ and $(A,\mathscr{C})\in\tilde{R}_k^{-}\times\mathfrak{X}_k$.
The following identities hold in $\mathrm{Gr}[\psi,\chi]$.
\begin{itemize}
\item[(a)] For every $\tilde{g}\in\mathbb{G}_k^{L}$,
\[
a_{\tilde{g}}\left[E^{\#,x}_k(m,\mathscr{E})\left(\tilde{g}X;\tilde{g}A,\tilde{g}\mathscr{C}\right)\right]=E^{\#,x}_k(m,\mathscr{E})(X;A,\mathscr{C}).
\]
\item[(b)] Let $(A^1,\mathscr{C}^1)\in\tilde{R}_k^{-}\times\mathfrak{X}_k$, let $\boldsymbol{\lambda}$ and $\boldsymbol{\lambda}'$ be real functions of $c$, and let $A^1=A-d\boldsymbol{\lambda}$ and $\mathscr{C}^1=\mathscr{C}-d\boldsymbol{\lambda}'$ on the bonds at distance at most $1$ from $X$.
Then
\[
G_{\boldsymbol{\lambda}+\boldsymbol{\lambda}'}\left[E^{\#,x}_k(m,\mathscr{E})(X;A^1,\mathscr{C}^1)\right]=E^{\#,x}_k(m,\mathscr{E})(X;A,\mathscr{C});
\]
this is the invariance under (T-G) relative to the polymer $X$.
\item[(c)]
\[
\mathrm{Cg}\left[E^{\#,x}_k(m,\mathscr{E})(X;-A,-\mathscr{C})\right]=E^{\#,x}_k(m,\mathscr{E})(X;A,\mathscr{C}).
\]
\item[(d)] If $m$ is real and $\mathscr{E}\in K^{x}_{k}$, then
\[
\mathbf{K}\left[E^{\#,x}_k(m,\mathscr{E})\left(X;-\overline{A},-\overline{\mathscr{C}}\right)\right]=E^{\#,x}_k(m,\mathscr{E})(X;A,\mathscr{C}).
\]
\end{itemize}
The same assertions hold in the reading~(o).
\end{named}

\begin{named}{Remark}
Proposition~24.2 corresponds to \cite[Lemma~17.21]{QEDI}, together with \cite[Lemma~17.20]{QEDI} and \cite[Section~17.11.1]{QEDI};
here the invariances are understood in the sense of Definition~24.I, relative to a polymer, to a set of arguments and to a domain, with $\mathrm{Cg}$ extended by the identity on $W$.
Rule (Pr5) of Chapter~24, Section~8.3, is a statement on the same algebra $\mathrm{Gr}[\psi,\chi,W]$ as \cite[Lemma~17.20(b)]{QEDI}.
It differs from these places in the following points.
\begin{itemize}
\item[(1)] The parts of the argument of \cite[Lemma~17.21]{QEDI} that concern the operators $T^{(i)}$ and level $0$ are replaced by Propositions~23.1\,(c), (f) and~23.2\,(b).
\item[(2)] The group acts on the labels of the fermion sites.
\item[(3)] At $k=0$ the invertible matrices $U$ of \cite[Lemma~17.20(b)]{QEDI} act on the index set $T^{\mathrm{f}}_0\times\{\beta\}$, and transposes refer to the pairing $\langle\,\cdot\,,\cdot\,\rangle^{\mathrm{ch}}$.
In Step~3 of the proof of \cite[Lemma~17.21(c)]{QEDI}, where $U$ denotes a union of $M$-cubes, the two-point function of $\left(H_k(1_U,\mathscr{T})\cdot\Gamma_k(1_U,\mathscr{T})\,W,\ \bar{H}_k(1_U,\mathscr{T})\,\bar{W}\right)$ is the kernel (I.1.4) of $H_k\cdot\Gamma_k\cdot\bar{H}_k^{T}$.
\item[(4)] That the elements of $\mathbb{G}^{L}_{k}$ map centres of cubes to centres, and links of cubes of $\pi''_k$ to links, follows from $\tilde{g}\,c=r(S)\,c+a$ with $a\in LM\mathbb{Z}^3$.
\end{itemize}
\end{named}

\subsubsection*{Proposition 24.3 (structure of the step map: what the objects attached to a polymer contain and read)}

\begin{named}{Proposition 24.3}
Assume (ST1)$^{\mathrm{ch}}$.
Let $(m,\mathscr{E})\in\mathrm{Ball}^{x}_{k}(4)$, $A\in\tilde{R}_k^{-}$, $\mathscr{C}\in\mathfrak{X}_k$, $\mathscr{T}=A+\mathscr{C}$, $\mathsf{Z}\in\mathrm{Box}_k$ and $\mathbf{Y}\in D_k$.
\begin{itemize}
\item[(adm)]
Let a pair be admissible for a polymer $X$ (Section~I.1.7).
Then it is admissible, and of the same kind, for every polymer that contains $X$.
The paths of a pair of kind (I) consist of bonds contained in $X$; a pair inside a cluster has no path.
A pair $(x,x^{(3)})$ of kind (II), $x=(c,g)$, has a single path, the photon bond $(c,c_3)$, and this bond meets $X$.
The photon sites of the two points of an admissible pair have sup-distance at most $\eta_k$ from $X$.
\item[(q1)] \emph{Domains.}
One has $2\le\mathrm{rad}_2\le\mathrm{rad}_1$ and $2\le\mathrm{rad}_4\le\mathrm{rad}_1$.
For $s\in\mathrm{Pd}_k$ each of the three norms (Section~I.5.3) of $\mathscr{Z}(s')$ is at most $\frac{1}{2}\,e_k^{-\frac{\varepsilon}{2}}$.
Let $\mathrm{Zf}$ be a complex bond function each of whose three norms is at most $\frac{1}{2}\,e_k^{-\frac{\varepsilon}{2}}$, and let $|t|\le\mathrm{rad}_1$.
Then $A+t\cdot\mathrm{Zf}$ lies in $\tilde{R}_k$, and $\mathscr{T}+t\cdot\mathrm{Zf}$ lies in $\tilde{R}_k+\mathfrak{X}_k$, a set contained in $\tilde{G}_k$.
\item[(q5)] \emph{Regularity.}
The field $\mathscr{Z}(s')$ is linear in $\mathsf{Z}$, continuous in $s'$ on $\mathrm{Pd}_k$, and analytic in the interior of $\mathrm{Pd}_k$.
At $B=\mathscr{T}+t\cdot\mathscr{Z}(s')$ the entries of $H_k^{\#}(s'',B)$, of $\Gamma_k(s'',B)$ and, for $k\ge1$, of $M_k(X,B)$ are continuous functions of $(s,t,A,\mathscr{C},\mathsf{Z})$ on
\[
\mathrm{Pd}_k\times\left\{|t|\le\mathrm{rad}_1\right\}\times\tilde{R}_k^{-}\times\mathfrak{X}_k\times\mathrm{Box}_k,
\]
analytic in $(s,t,A,\mathscr{C})$ at the interior points, and, for every fixed $s\in\mathrm{Pd}_k$, analytic in $(t,A,\mathscr{C})$ on a neighbourhood of $\left\{|t|\le\mathrm{rad}_1\right\}\times\tilde{R}_k^{-}\times\mathfrak{X}_k$.
Every function $G(\cdot\,,X)$ to which $\mathrm{Loc}$ is applied in Definition~24.D\,(D4) satisfies the hypothesis of Definition~24.D\,(D3).
\item[(al)] \emph{Sub-polymers; the datum.}
For every term $(X,P)$ of (8.2) of Definition~24.D that is attributed to $\mathbf{Y}$ (Definition~24.D\,(D8)), the polymer $X$ and every cube of $P_0(X)\cup P$ are contained in $\mathbf{Y}$.
The sets $\{X_i\}$ and the families $(X_1,\dots,X_n)$ of (8.4) whose union is $\mathbf{Y}$ consist of polymers contained in $\mathbf{Y}$.
For every $Y\in D_{k+1}$ the polymers $X\in D_k$ with $\bar{X}^{L}=LY$ (Definition~24.D\,(D6)) are contained in $LY$.
Let $(m_1,\mathscr{E}_1)$ and $(m_2,\mathscr{E}_2)$ in $\mathrm{Ball}^{x}_{k}(4)$ satisfy $m_1+m(\mathscr{E}_1)=m_2+m(\mathscr{E}_2)$ and $(R\mathscr{E}_1)(X)=(R\mathscr{E}_2)(X)$ for all polymers $X\subseteq\mathbf{Y}$;
then the objects attached to $\mathbf{Y}$ that are formed with the two data coincide.
(By Section~I.5.6, $m(\mathscr{E})$ and $(R\mathscr{E})(X)-\mathscr{E}(X)$ depend on $\mathscr{E}$ only through the kernels $\mathscr{E}_{00}(X;0,0)$ and $\mathscr{E}_{20}(X;0,0;\cdot\,,\cdot\,)$ of small polymers $X$.)
\item[(be)] \emph{Where the exact operators and forms stand.}
In the displays that define an object attached to $\mathbf{Y}$, the H\"older quotients $\delta_{\alpha,B}$ and $\delta_{\alpha,-B}$ occur at the pairs that are admissible for polymers contained in $\mathbf{Y}$, with $B=\mathscr{T}+\mathscr{Z}(s')$ or $B=\mathscr{T}$.
No substitution occurs in $E^{(i),*}(\mathbf{Y})$.
The substitution is $\mathrm{sub}_{\mathbf{Y},\mathscr{T}}$ in $E^{(i),x}(\mathbf{Y})$ and in $E^{(\le3),x}(\mathbf{Y})$;
in $E^{(\le5),x}(\mathbf{Y})$, $\mathrm{Km}(\mathbf{Y})$, $\mathrm{Hm}(\mathbf{Y})$ and $E^{\#,x}_k(m,\mathscr{E})(\mathbf{Y})$ it is $\mathrm{sub}_{X,\mathscr{T}}$ with polymers $X\subseteq\mathbf{Y}$, and it stands inside $G^{(\le3)}(s'',X)=E^{(\le3),x}(X;\,\cdot\,)$.
Hence, for $k\ge1$, the operators $T^{(i)}(\mathscr{T},\cdot\,)$ and $\bar{T}^{(i)}(\mathscr{T},\cdot\,)$ enter through their rows at unit sites of $\mathbf{Y}$.
The forms $\langle\bar{\psi},\psi\rangle_{\square}$ and, at $k=0$, the elements $\mathrm{Vx}(\mathscr{Z}(s'),\square)$ occur with $M$-cubes $\square\subseteq\mathbf{Y}$.
The map $C$ occurs in (24.1).
The rows just named read the connection, and $\mathrm{Vx}$ reads $\mathscr{Z}(s')$, on bonds that meet $\mathbf{Y}$;
$C$, $M_i(\cdot)$, $\mathsf{C}$ and $\langle\bar{\psi},\psi\rangle_{\square}$ contain no connection.
(That each of them is given by one formula, with its range, is the content of Propositions~25.1--25.3.)
\item[(ga)] \emph{Primitives.}
Let $s=(s',s'')\in\mathrm{Pd}_k$ vanish at every parameter cube that is not contained in $\mathbf{Y}$.
(The families $s$ that occur in a term of (8.2) attributed to $\mathbf{Y}$ have this property.)
\begin{itemize}
\item[(ga1)]
Let $b$ be a fine bond that meets $\mathbf{Y}$; then $H_k(s')(b,b')=0$ unless the unit bond $b'$ is contained in $\mathbf{Y}$.
Let $b'$ be a unit bond contained in $\mathbf{Y}$; then $C(b',\Upsilon)=0$ unless $\Upsilon$ is contained in $\mathbf{Y}$.
Let $\Upsilon$ be contained in $\mathbf{Y}$; then $C_k^{1/2}(s')(\Upsilon,\Upsilon')=0$ unless $\Upsilon'$ is contained in $\mathbf{Y}$.
The rows of $H_k(s')$ at the fine bonds that are not contained in $\mathbf{Y}$ vanish, and so do the rows of $C_k^{1/2}(s')$ at the basis elements that are not contained in $\mathbf{Y}$.
Consequently, for every fine bond $b$ that meets $\mathbf{Y}$,
\begin{equation}
\mathscr{Z}(s')(b)=\sum_{b',\Upsilon,\Upsilon'\subseteq\mathbf{Y}}H_k(s')(b,b')\cdot C(b',\Upsilon)\cdot C_k^{1/2}(s')(\Upsilon,\Upsilon')\cdot\mathsf{Z}(\Upsilon'),\tag{24.1}
\end{equation}
the sum extending over the unit bonds $b'$ and the basis elements $\Upsilon$, $\Upsilon'$ that are contained in $\mathbf{Y}$.
\item[(ga2)]
Let $B\in\tilde{G}_k$ and let $x=(c,g)$ be a fermion site of $T^{\mathrm{f}}_k$.
The rows of $H_k(s'',B)$ and of $\bar{H}_k(s'',B)$ at $x$ vanish unless every cube of $\pi_k$ that meets $K_c$ is contained in $\mathbf{Y}$, in particular unless $c\in\mathbf{Y}$;
and their columns lie in $\mathbf{Y}$ (for $k\ge1$ they are unit sites of $\mathbf{Y}$; for $k=0$ the only column is $x$ itself).
Let $u$ be a unit index in $\mathbf{Y}$, that is, an element of the unit index set of level $k$ (Section~G3.2.0) in $\mathbf{Y}$:
for $k\ge1$ a unit site $y\in\mathbf{Y}$, for $k=0$ a fermion site $x$ over $\mathbf{Y}$.
Then the row of $\Gamma_k(s'',B)$ at $u$ has its columns in (over) $\mathbf{Y}$.
(In the rows of (ga2) and in the entries of (ga4) the letter $u$ denotes this point, and an entry is a $4\times4$ block;
in $\Psi(u)$ and $W(u)$ of (de) and (ze) it denotes the index $(y,\beta,\omega)$, respectively $(x,\beta,\omega)$, of Section~I.3.6, which lies in, respectively over, $\mathbf{Y}$ if its point does.)
For $k\ge1$ and a polymer $X\subseteq\mathbf{Y}$ the kernel of $M_k(X,B)$ vanishes unless both of its unit sites lie in $X$.
\item[(ga3)]
The rows named in (ga2) are the same for any two $B,\tilde{B}\in\tilde{G}_k$ that coincide on the bonds that meet $\mathbf{Y}$;
and $M_k(X,B)=M_k(X,\tilde{B})$ whenever $B$ and $\tilde{B}$ coincide on the bonds that meet $X$.
\item[(ga4)] \emph{Which entries stand in the displays.}
After the omission of the terms that vanish by (ga1) and (ga2), the displays (8.1)$^{\mathrm{ch}}$--(8.4) for an object attached to $\mathbf{Y}$ contain the primitives only through the following entries, $s$ ranging over the families named above:
\begin{gather*}
H_k(s')(b,b'),\ b\ \text{a fine bond that meets}\ \mathbf{Y},\ b'\subseteq\mathbf{Y};\qquad
C_k^{1/2}(s')(\Upsilon,\Upsilon'),\ \Upsilon,\Upsilon'\subseteq\mathbf{Y};\\
H_k(s'',B)(x,u),\ \bar{H}_k(s'',B)(x,u),\ x=(c,g),\ c\in\mathbf{Y},\ u\in\mathbf{Y},\ \text{at}\ B=\mathscr{T}+\mathscr{Z}(s')\ \text{and}\ B=\mathscr{T};\\
\Gamma_k(s'',\mathscr{T})(u,u'),\ u,u'\in\mathbf{Y};\qquad
\text{for}\ k\ge1\colon\ M_k\left(X,\mathscr{T}+t\cdot\mathscr{Z}(s')\right)(y,y'),\ X\subseteq\mathbf{Y},\ y,y'\in X,\ |t|=\mathrm{rad}_1.
\end{gather*}
The entries of the second line stand at the sites over polymers $X\subseteq\mathbf{Y}$ and at the two points of the pairs admissible for such $X$;
the rows at points outside $\mathbf{Y}$ vanish.
The first slot of $\mathscr{E}^{r}(X;\cdot\,,\mathscr{C},\cdot\,)$, $X\subseteq\mathbf{Y}$, is $A+t\cdot\mathscr{Z}(s')$ with $|t|=\mathrm{rad}_2$, or $A$.
By (24.1) and (q1), on the bonds that meet $\mathbf{Y}$ all these connections are determined by the values of $A$ and $\mathscr{C}$ on these bonds, by the coordinates $\mathsf{Z}(\Upsilon)$ with $\Upsilon\subseteq\mathbf{Y}$, and by the entries of $H_k(s')$ and $C_k^{1/2}(s')$ in the first line.
Consequently, if two systems of primitives on $T_k$ with the properties of (ST1)$^{\mathrm{ch}}$ coincide in these entries (each entry being taken at the connections $\mathscr{T}+t\cdot\mathscr{Z}(s')$ and $\mathscr{T}$ formed with the respective system, for all $s$ as above, all $(A,\mathscr{C})\in\tilde{R}_k^{-}\times\mathfrak{X}_k$ and all $\mathsf{Z}\in\mathrm{Box}_k$;
this is the case if the entries of the first line coincide and the other entries coincide as functions of $(s'',B)$ on $\mathrm{Pd}_k\times\tilde{G}_k$),
then the objects attached to $\mathbf{Y}$ that are formed with the two systems coincide.
\end{itemize}
\item[(de)] \emph{Integrations.}
The integrations in the displays for an object attached to $\mathbf{Y}$ are the following:
in $t$, over the circles $|t|=\mathrm{rad}_i$;
in $s_P$, over $[0,1]^P$, where $P$ is a set of parameter cubes contained in $\mathbf{Y}$;
in the generators $W(u)$ with $u$ in $\mathbf{Y}$ (for $k=0$: over $\mathbf{Y}$), with respect to $\mathrm{d}\mu_I(W)$, whose moments are given by (I.B4) with $\nu=1$ for $k\ge1$ and $\nu=\frac{1}{48}$ for $k=0$, and which factorises over the points of the index set (rule (I.N4) of Section~I.5.2);
and in the coordinates $\mathsf{Z}(\Upsilon)$ with $\Upsilon\subseteq\mathbf{Y}$, with respect to the product of copies of a single probability measure on $\mathbb{R}$, namely the standard normal law restricted to $[-p_{0,k},p_{0,k}]$ and normalised.
\item[(ze)] \emph{Generators; what is read.}
For $i\le3$ every generator contained in $E^{(i),*}(\mathbf{Y})$ is a $\psi(\xi)$ with $\xi$ over $\mathbf{Y}$, a $\chi(\boldsymbol{\zeta})$ with $\boldsymbol{\zeta}$ admissible for $\mathbf{Y}$ or, for $k\ge1$, a $\Psi(u)$ with $u\in\mathbf{Y}$; in particular $E^{(i),*}(\mathbf{Y})$ contains no generator $W$, and for $k=0$ no generator $\Psi$.
The elements $E^{(i),x}(\mathbf{Y})$ and $E^{(\le3),x}(\mathbf{Y})$ contain only generators $\psi(\xi)$ with $\xi$ over $\mathbf{Y}$ and $\chi(\boldsymbol{\zeta})$ with $\boldsymbol{\zeta}$ admissible for $\mathbf{Y}$;
$E^{(\le5),x}(\mathbf{Y})$ and $\mathrm{Km}(\mathbf{Y})$ contain, besides these, only generators $W(u)$ with $u$ in (over) $\mathbf{Y}$;
$\mathrm{Hm}(\mathbf{Y})$ and $E^{\#,x}_k(m,\mathscr{E})(\mathbf{Y})$ contain only $\psi$ over $\mathbf{Y}$ and $\chi$ at pairs admissible for $\mathbf{Y}$.
All these elements are even.
The objects $E^{(i),*}(\mathbf{Y})$, $E^{(i),x}(\mathbf{Y})$, $E^{(\le3),x}(\mathbf{Y})$, $E^{(\le5),x}(\mathbf{Y})$ and $\mathrm{Km}(\mathbf{Y})$ depend on $\mathsf{Z}$ only through the coordinates $\mathsf{Z}(\Upsilon)$ with $\Upsilon\subseteq\mathbf{Y}$.
Every object attached to $\mathbf{Y}$ takes the same value at any two pairs $(A,\mathscr{C})$ and $(A^1,\mathscr{C}^1)$ of $\tilde{R}_k^{-}\times\mathfrak{X}_k$ that coincide on the bonds that meet $\mathbf{Y}$, at every $\mathsf{Z}\in\mathrm{Box}_k$.
The coefficients of $E^{(i),*}(\mathbf{Y})$, $E^{(i),x}(\mathbf{Y})$, $E^{(\le3),x}(\mathbf{Y})$, $E^{(\le5),x}(\mathbf{Y})$ and $\mathrm{Km}(\mathbf{Y})$ are continuous in $(A,\mathscr{C},\mathsf{Z})$ on $\tilde{R}_k^{-}\times\mathfrak{X}_k\times\mathrm{Box}_k$ and, at every $\mathsf{Z}\in\mathrm{Box}_k$, analytic in $(A,\mathscr{C})$ on $\tilde{R}_k^{-}\times\mathfrak{X}_k$;
the coefficients of $\mathrm{Hm}(\mathbf{Y})$ are analytic in $(A,\mathscr{C})$ on $\tilde{R}_k^{-}\times\mathfrak{X}_k$.
Two polymers without a common $M$-cube have no fermion site, no unit site and no contained basis element in common.
\item[(et)] \emph{Terms that vanish; expansion.}
Let $G(\cdot\,,X)$ be one of the functions to which $\mathrm{Loc}$ is applied in Definition~24.D\,(D4), and let $P$ be a set of parameter cubes, of the kinds that occur, outside $P_0(X)$.
If $Y'(X,P)$ is not connected, there is a cube $\mathrm{pc}\in P$ such that $G(s,X)=G(\tilde{s},X)$ for all $s,\tilde{s}\in\mathrm{Pd}_k$ that vanish outside $P_0(X)\cup P$ and differ at $\mathrm{pc}$ only;
hence the term $(X,P)$ of (8.2) vanishes.
Consequently
\[
\sum_{\mathbf{Y}\in D_k}\mathrm{Loc}[G](\mathbf{Y})=\sum_{X\in D_k}G(1,X),
\]
where $1$ is the family with the value $1$ at all parameter cubes (\cite[Lemmas~17.12 and~17.13]{QEDI}).
The elements $E^{(\le5),x}(\mathbf{Y})$ are even and therefore commute, and for every union $U$ of $M$-cubes
\[
\exp\left\{\sum_{\mathbf{Y}\in D_k,\ \mathbf{Y}\subseteq U}E^{(\le5),x}(\mathbf{Y})\right\}=\sum_{Z}\mathrm{Km}(Z),
\]
where $Z$ ranges over the unions of $M$-cubes contained in $U$, the empty union included, and where $\mathrm{Km}(Z)$ is the product of the $\mathrm{Km}(Z_a)$ over the components $Z_a$ of $Z$.
\end{itemize}
The two other parts of $\mathscr{E}^{+}(Y)$ fall outside (ga):
$(R\mathscr{E})^{\mathrm{nat}}(X)$ contains the numbers $\mathrm{al}_{0}$, $\mathrm{al}_{2}$ and the homomorphism $\boldsymbol{\Theta}_{X,\mathscr{T}}$ of Section~I.5.6;
${\delta E^{\mathrm{det}}}_k(X)$ contains the family (p5) of Section~I.5.8 through the kernel of $S_{k,y}(X,B)$ at indices over $X$ (by W12, the locality property~12 of the systems of primitives in \cite[Section~5.4.2]{QEDI}), and it contains the exact operators in the displays of $\acute{E}_k$ (Section~I.5.7).
The exact operators of all three parts are described in Propositions~25.1--25.3.

The same assertions hold in the reading~(o).
\end{named}

\begin{named}{Remark}
Among the places of \cite{QEDI} to which Proposition~24.3 corresponds are the display (adm) in \cite[Section~17.7]{QEDI}, the clauses (q1) and (q5) of \cite[Lemma~17.11]{QEDI}, \cite[Lemma~17.14]{QEDI} and \cite[Corollary~17.15]{QEDI}, assertion (2) of \cite[Section~17.10.1]{QEDI}, \cite[Proposition~17.24(a), (c), (d)]{QEDI} and the paragraph that follows the list of its clauses, together with the part of \cite[Proposition~17.24(b)]{QEDI} that says where the exact operators enter, and the second half of \cite[Theorem~5.71(g)]{QEDI}.
Further places correspond to single sentences of the proposition; among them are the following.
The evenness asserted in (ze), and the continuity and the analyticity asserted there, correspond to \cite[Proposition~17.16(e)]{QEDI} and to the paragraphs on continuity and analyticity in the proof of \cite[Proposition~17.16]{QEDI};
the last sentence of (q5) corresponds to the proofs of \cite[Proposition~17.16(b), (c)]{QEDI}, in which the functions to which $\mathrm{Loc}$ is applied are shown to be continuous on $\mathrm{Pd}_k$ and analytic in its interior;
the analyticity of the coefficients of $\mathrm{Hm}(\mathbf{Y})$ asserted in (ze) corresponds to the proof of assertion (3) of \cite[Section~17.10.1]{QEDI};
the first display of (et) corresponds to \cite[Proposition~17.16(d)]{QEDI};
and the second display of (et), with the product of the $\mathrm{Km}(Z_a)$ over the components $Z_a$ of $Z$, corresponds to Step~1 (Mayer expansion) of the proof of \cite[Lemma~17.17]{QEDI}.
In \cite[Proposition~17.16(e)]{QEDI} evenness, continuity and analyticity are stated for $E^{(\le3),x}(\mathbf{Y})$ and $E^{(\le5),x}(\mathbf{Y})$,
whereas (ze) asserts the evenness also for $E^{(i),*}(\mathbf{Y})$, $E^{(i),x}(\mathbf{Y})$, $\mathrm{Km}(\mathbf{Y})$, $\mathrm{Hm}(\mathbf{Y})$ and $E^{\#,x}_k(m,\mathscr{E})(\mathbf{Y})$, the continuity and the analyticity also for $E^{(i),*}(\mathbf{Y})$, $E^{(i),x}(\mathbf{Y})$ and $\mathrm{Km}(\mathbf{Y})$, and the analyticity of the coefficients also for $\mathrm{Hm}(\mathbf{Y})$.
In Step~1 of the proof of \cite[Lemma~17.17]{QEDI} the expansion is written for the sum over the polymers of $D_k$, whereas (et) states it for every union $U$ of $M$-cubes.
The last sentence of (q5) refers to the hypothesis of Definition~24.D\,(D3);
in \cite[Section~17.4.3]{QEDI} the map $\mathrm{Loc}$ is defined by the same formula as in (8.2) of Definition~24.D, for a function $G(s,X)$ of the parameters and of a polymer.
Its relation to these places is the following.
\begin{itemize}
\item[(1)] A fine index is a fermion site $x=(c,g)$, read at $c$ by (I.POS) of Section~I.1.4;
for fermion sites ``in $\mathbf{Y}$'' becomes ``over $\mathbf{Y}$'';
at $k=0$ the unit indices are fermion sites.
\item[(2)] Where the exact operators stand is clause (be);
that they are given by one formula, their ranges and their list are the content of Propositions~25.1--25.3.
\item[(3)] Clause (ga) states the vanishing of entries ((ga1), (ga2)), the reduced display (24.1) and the list (ga4) of the entries that stand in the displays.
\item[(4)] The assertion on the pairs $(A,\mathscr{C})$ and the assertion on $E^{(i),*}(\mathbf{Y})$ are part of (ze).
The first has the form of the paragraph that follows the list of clauses in \cite[Proposition~17.24]{QEDI}: each object attached to $\mathbf{Y}$ takes the same value at any two pairs $(A,\mathscr{C})$ and $(A^1,\mathscr{C}^1)$ of $\tilde{R}_k^{-}\times\mathfrak{X}_k$ that coincide on the bonds that meet $\mathbf{Y}$, at every $\mathsf{Z}\in\mathrm{Box}_k$;
(ze) asserts it in addition for $E^{(i),*}(\mathbf{Y})$ and $E^{(\le3),x}(\mathbf{Y})$, which are not among the objects attached to $\mathbf{Y}$ in \cite[Proposition~17.24]{QEDI}.
\item[(5)] The family (p3) is taken at the connections $\mathscr{T}+\mathscr{Z}(s')$ and $\mathscr{T}$.
\item[(6)] The product form of $\hat{\mu}$ in (de) follows, as in the proof of \cite[Proposition~17.24(d)]{QEDI}, from the definitions of $\mu_I$, $\mathrm{Box}_k$ and $\hat{\mu}$ and from the orthonormality of the basis $\{\Upsilon\}$.
\item[(7)] \cite[Lemma~17.14(b)]{QEDI} is clause (et), in the following form: on the families that vanish outside $P_0(X)\cup P$, the function $G(\cdot\,,X)$ does not depend on $s_{\mathrm{pc}}$ for some $\mathrm{pc}\in P$.
\item[(8)] The operators $M_k$ occur only for $k\ge1$.
\item[(9)] The inequalities between the radii are part of (q1).
\end{itemize}
\end{named}

\subsection*{48 (continued). The statements without the second slot; the invariant at $\mathscr{C}=0$}

\subsubsection*{Proposition~48.6 (the invariant $E_k=\mathscr{E}_k$ at $\mathscr{C}=0$)}

\begin{named}{Proposition 48.6}[The invariant $E_k=\mathscr{E}_k$ at $\mathscr{C}=0$]
\begin{itemize}
\item[(a)] (One step.)
The data are those of (S-2), together with the data for which the step map of level $k$ is defined (Section~G3.4.0 and the Propositions~24.m; the notation of the step map and of the recursion is that of Section~I.5.7):
$\ell\ge M\cdot r_k$, $\mu\in(0,1/L]$, and a system of primitives of level $k$ on $T_k$ with the properties W1--W8 and W10--W17 for $\mathrm{par}$.
Here W1--W17 are the properties~1--17 of \cite[Section~5.4.2]{QEDI}, read as in Section~I.5.8, and the system is that of (ST1)$^{\mathrm{ch}}$ in Section~G3.4.0.
The property W9 is not required, not even at $k=0$: at that level the row W17 is taken without W9, that is, with W8, W10, W11 in place of its W8--W11 and with W6--W8 in place of its W6--W9.
Further, let $(m,\mathscr{E})\in\mathrm{Ball}^{x}_{k}(4)$, where $\mathscr{E}=0$ if $k=0$, and put $E:=\mathscr{E}|_0$.
Then $(m,E)$ lies in $\mathrm{Ball}_{k}(4)$, the domain of $E^{*}_k$, and
\[
\varepsilon({E})=\varepsilon(\mathscr{E}),\qquad m({E})=m(\mathscr{E}).
\]
For every $X\in D_k$ and every $\mathscr{A}\in\tilde{R}_k$ one has ${\delta E^{\mathrm{det}}}_k(X;\mathscr{A},0)=E_k^{\mathrm{det}}(X,\mathscr{A})$.
Finally, the identity
\[
\left.\left(L R\mathscr{E}+L\,{\delta E^{\mathrm{det}}}_k+E^{*,x}_k(m,\mathscr{E})\right)\right|_0=L R E+L E_k^{\mathrm{det}}+E^{*}_k(m,E)
\]
holds in $K_{k+1}^{C}$.
At level $0$ one has $\mathscr{E}=0$ and $E=0$, and the identity reads
\[
\left.\left(L\,{\delta E^{\mathrm{det}}}_0+E^{*,x}_0(m,0)\right)\right|_0=L E_0^{\mathrm{det}}+E^{*}_0(m,0).
\]
\item[(b)] (Trajectory.)
Assume (S-4), and for every level $k<k_{*}(R)$ let a system of primitives of level $k$ on $T_k(R)$ be given which has the properties W1--W8 and W10--W17 for $\mathrm{par}$, understood as in~(a).
Let $(m_{\mathrm{in}},\varepsilon^{\mathrm{v}}_{\mathrm{in}})$ be a real pair.
Consider the recursion of Section~I.5.7, with the step map and the determinant functions of level $k$ formed with the given system of level $k$, with $\varepsilon^{\mathrm{v},t}_0=\varepsilon^{\mathrm{v}}_{\mathrm{in}}$, with the mass coupling $m_{\mathrm{in}}$ at level $0$, and with $E_0=0$, $\mathscr{E}_0=0$, $\mathfrak{P}_0=0$; at $k=0$ the symbol $m_k$ denotes $m_{\mathrm{in}}$ (Section~G4.0.5).
Let $k_{\mathrm{up}}$ be an integer with $0\le k_{\mathrm{up}}\le k_{*}(R)$ such that the trajectory exists up to level $k_{\mathrm{up}}$, that is,
\[
(m_k,\mathscr{E}_k)\in\mathrm{Ball}^{x}_{k}(4)\qquad\text{for all }k<k_{\mathrm{up}};
\]
these conditions are read in the order $k=0,1,\dots$, the entries of level $k$ being defined as soon as the conditions at all levels below $k$ hold.
Then, for every $k\le k_{\mathrm{up}}$, all five entries of level $k$ are defined, $\varepsilon^{\mathrm{v},t}_k$ and $m_k$ are real, $\mathscr{E}_k\in K^{x}_{k}$,
\[
E_k=\mathscr{E}_k|_0\in K_{k},
\]
where at $k=0$ the symbol $K_{0}$ denotes the space of level $0$ and not the constant $K_{0}$ of (C0) in Section~I.8.2, and $\mathfrak{P}_k(X,0)=0$ for every $X\in D_k$.
If moreover $k<k_{\mathrm{up}}$, then
\[
(m_k,E_k)\in\mathrm{Ball}_{k}(4),\qquad\varepsilon(E_k)=\varepsilon(\mathscr{E}_k),\qquad m(E_k)=m(\mathscr{E}_k).
\]
Hence the first three of the five equations of the recursion, those for $\varepsilon^{\mathrm{v},t}_{k+1}$, $m_{k+1}$ and $E_{k+1}$, are the restriction to $\mathscr{C}=0$ of the step $(\varepsilon^{\mathrm{v}},m,\mathscr{E})\mapsto(\varepsilon^{\mathrm{v},+},m^{+},\mathscr{E}^{+})$ of the Propositions~24.m:
for $k<k_{\mathrm{up}}$ and $(\varepsilon^{\mathrm{v}},m,\mathscr{E})=(\varepsilon^{\mathrm{v},t}_k,m_k,\mathscr{E}_k)$ one has
\[
\varepsilon^{\mathrm{v},+}=\varepsilon^{\mathrm{v},t}_{k+1},\qquad m^{+}=m_{k+1},\qquad \mathscr{E}^{+}|_0=E_{k+1}.
\]
\end{itemize}
\end{named}

\begin{named}{Remark}
Proposition~48.6 corresponds to the last paragraphs of \cite[Section~5.3.8]{QEDI}:
the sentence according to which the recursion gives the terms of level $k+1$ whenever $(m_k,\mathscr{E}_k)$ lies in $\mathrm{Ball}^{x}_{k}(4)$;
the assertion $E_k=\mathscr{E}_k|_0$ for every $k\le k_{\mathrm{up}}$, when the trajectory exists up to level $k_{\mathrm{up}}$, together with its consequences $(m_k,E_k)\in\mathrm{Ball}_{k}(4)$, $\varepsilon(E_k)=\varepsilon(\mathscr{E}_k)$ and $m(E_k)=m(\mathscr{E}_k)$ for $k<k_{\mathrm{up}}$;
the proof of this assertion, an induction on $k$ whose step is the identity of clause~(a);
and the identity $\mathfrak{P}_k(X,0)=0$.
It corresponds further to \cite[Proposition~19.13]{QEDI}, which contains the identity $\mathscr{E}_k|_0=E_k$ under the hypothesis $|m_i|\le w^{\mathrm{m}}_i$, together with bounds on the norms of $E_k$ and $\mathscr{E}_k$, and to \cite[Theorem~5.71(b), (d)]{QEDI}.
As in \cite[Section~5.3.8]{QEDI}, the invariant is stated where the recursion is defined, that is, on all of $\mathrm{Ball}^{x}_{k}(4)$.
This is needed because the equation of the recursion for $E_{k+1}$ requires $(m_k,E_k)\in\mathrm{Ball}_{k}(4)$, and the equations for $\varepsilon^{\mathrm{v},t}_{k+1}$ and $m_{k+1}$ contain $\varepsilon(E_k)$ and $m(E_k)$, whereas the step map contains $\varepsilon(\mathscr{E}_k)$ and $m(\mathscr{E}_k)$.
Proposition~48.6 differs from these places in the following points.
\begin{itemize}
\item[(1)] Clause~(a) states the identity of one step as an assertion of its own, for every $(m,\mathscr{E})\in\mathrm{Ball}^{x}_{k}(4)$;
in \cite[Section~5.3.8]{QEDI} it is the induction step of the proof, and \cite[Theorem~5.71(d)]{QEDI} states it under the hypotheses of \cite[Theorem~5.71(c)]{QEDI}.
\item[(2)] Clause~(b) asserts in addition, for every $k\le k_{\mathrm{up}}$, that $\varepsilon^{\mathrm{v},t}_k$ and $m_k$ are real, that $\mathscr{E}_k\in K^{x}_{k}$ and that $E_k\in K_{k}$.
\cite[Section~5.3.8]{QEDI} asserts $\mathscr{E}_k\in K^{x,C}_{k}$ and $E_k\in K^{C}_{k}$;
the reality of the two numbers and the real spaces stand in \cite[Proposition~19.13]{QEDI}, under the hypothesis $|m_i|\le w^{\mathrm{m}}_i$.
\item[(3)] The systems of primitives need the properties W1--W8 and W10--W17 only, and the property W9 is not required;
\cite[Section~5.3.8]{QEDI} assumes \cite[Definition~5.69]{QEDI}, which requires the properties~1--17 of \cite[Section~5.4.2]{QEDI}.
\item[(4)] The fermions are chamber fermions, whereas \cite{QEDI} treats Wilson fermions.
\end{itemize}
\end{named}

\subsection*{26. Projection of the localised operators to smaller tori}

\subsubsection*{Proposition~26.1 (definition: two tori; projection of a system of primitives)}

\begin{named}{Proposition 26.1}[Two tori; projection of a system of primitives]
\emph{Data.}
Let $k\ge0$ be a level, and let $\ell_1,\ell_2\in\mathrm{Cl}$ be sides with $\ell_2=\mathrm{Jr}\cdot\ell_1$ for an odd integer $\mathrm{Jr}\ge1$.
Let $\mu\in(0,1]$ and $e_k\in(0,e_0]$ be given, the same on both tori; thus $(k,\ell_i,\mu,e_k)$ is (SF) for $i=1,2$.

\emph{Notation.}
For $i=1,2$ put
\[
T_i:=\mathbb{L}_k/\ell_i\mathbb{Z}^3,\qquad T^{\mathrm{f}}_i:=T_i\times O_{\mathrm{cub}},\qquad T^0_i:=\mathbb{Z}^3/\ell_i\mathbb{Z}^3.
\]
The covering map $\mathbb{L}_k\to T_i$ is denoted $\mathrm{pr}_i$.
It is extended to the index sets (I.i1)--(I.i5) of Section~I.1.9: on fermion sites $\mathrm{pr}_i(c,g):=(\mathrm{pr}_i\,c,g)$, and on pairs $\mathrm{pr}_i(x,x')$ is the class of $(x,x')$ in $\mathrm{Pz}_k(T_i)$.
Let $\mathrm{pj}\colon T_2\to T_1$ be the covering map with $\mathrm{pj}\circ\mathrm{pr}_2=\mathrm{pr}_1$, and let $d_i$ denote the sup distance on $T_i$.
Objects on $T_i$ are identified with the $\ell_i\mathbb{Z}^3$-periodic, respectively antiperiodic, objects on the covering index sets, and $\mathrm{pj}^{*}$ is the inclusion of the objects on $T_1$ into the objects on $T_2$ (Section~I.6.3).
For $w\in\mathbb{Z}^3$ let $\boldsymbol{\chi}(w):=1$ for bond functions, scalars, parameter families and elements of $V_{k}$, and $\boldsymbol{\chi}(w):=(-1)^{w_1+w_2+w_3}$ for spinors and for functions of pairs.
For $v\in\ell_1\mathbb{Z}^3$ the translation $T_v$ is $(T_v f)(t):=f(t-v)$.
The image of $\mathrm{pj}^{*}$ is the space
\[
\mathrm{Per}:=\left\{f\text{ on }T_2:\ f(t+\ell_1\cdot w)=\boldsymbol{\chi}(w)\cdot f(t)\text{ for all }w\in\mathbb{Z}^3\right\}.
\]

\emph{Facts.}
The following properties are read off from Section~I.6.3.
The map $\mathrm{pj}^{*}$ is linear and injective, and $|\mathrm{pj}^{*}f|_{\infty}=|f|_{\infty}$.
For each of the bilinear pairings $\langle\,\cdot\,,\cdot\,\rangle^{\mathrm{ch}}$ on $\mathfrak{S}^{\mathrm{ch}}(\eta_k,\cdot\,)$ and $\langle\,\cdot\,,\cdot\,\rangle_1$ on $\mathfrak{S}(1,\cdot\,)$ of Section~I.1.6, and for the pairings of bond functions, of scalars and of elements of $V_{k}$ of \cite[Section~5.2.5]{QEDI}, the pairing $\langle\mathrm{pj}^{*}f',\mathrm{pj}^{*}f\rangle$ on $T_2$ equals $\mathrm{Jr}^3\cdot\langle f',f\rangle$ on $T_1$.
For parameter families $(\mathrm{pj}^{*}s)_{\mathrm{pc}}=s_{\mathrm{pj}\,\mathrm{pc}}$; consequently $\mathrm{pj}^{*}$ maps $\mathrm{Pd}_k(T_1)$ into $\mathrm{Pd}_k(T_2)$, and interior points to interior points.
By \cite[Lemma~17.7]{QEDI} the map $\mathrm{pj}^{*}$ sends $\tilde{G}_k(T_1)$ into $\tilde{G}_k(T_2)$.
One has $J\,\mathrm{pj}^{*}=\mathrm{pj}^{*}J$ and $J^{*}\mathrm{pj}^{*}=\mathrm{pj}^{*}J^{*}$, and $\mathrm{pj}^{*}$ commutes with $A\mapsto-A$, with complex conjugation and with constant spin matrices.

\emph{Definition.}
Let a system of primitives of level $k$ on $T_2$ with the property W14$^{\mathrm{ch}}$ be given, and write $\mathrm{Op}^{(2)}$ for its families.
For $s\in\mathrm{Pd}_k(T_1)$, $A\in\tilde{G}_k(T_1)$, $X\in D_k(T_1)$, $y\ge0$ and $f$ in the source space on $T_1$ set
\begin{align*}
\mathrm{pj}^{*}\left(\mathrm{Op}^{\mathrm{pj}}(s,A)f\right)&:=\mathrm{Op}^{(2)}\left(\mathrm{pj}^{*}s,\mathrm{pj}^{*}A\right)\mathrm{pj}^{*}f&&\text{for (p1)--(p3) and }\Gamma_k,\\
\mathrm{pj}^{*}\left(\mathrm{Op}^{\mathrm{pj}}(X,A)f\right)&:=\sum_{X'\in D_k(T_2):\ \mathrm{pj}\,X'=X}\mathrm{Op}^{(2)}\left(X',\mathrm{pj}^{*}A\right)\mathrm{pj}^{*}f&&\text{for }M_k\ (k\ge1)\text{ and }S_{k,y},\\
\bar{H}_{k}^{\mathrm{pj}}(s'',A)&:=\mathsf{C}^{-1}H_{k}^{\mathrm{pj}}(s'',-A)\,\mathsf{C}.
\end{align*}
For the last family one then has $\mathrm{pj}^{*}\left(\bar{H}_{k}^{\mathrm{pj}}(s'',A)f\right)=\bar{H}_{k}^{(2)}\left(\mathrm{pj}^{*}s'',\mathrm{pj}^{*}A\right)\mathrm{pj}^{*}f$, so that it is the projection of $\bar{H}_{k}^{(2)}$.
Write
\[
\widetilde{\mathrm{Op}}(s,A):=\mathrm{Op}^{(2)}\left(\mathrm{pj}^{*}s,\mathrm{pj}^{*}A\right),\qquad
\widetilde{\mathrm{Op}}(X,A):=\sum_{X':\ \mathrm{pj}\,X'=X}\mathrm{Op}^{(2)}\left(X',\mathrm{pj}^{*}A\right).
\]
The family $\mathrm{Op}^{\mathrm{pj}}$ is called the \emph{projection} of the system to $T_1$.
That the right-hand sides of the first two definitions lie in $\mathrm{Per}$ is the first assertion of Proposition~26.3\,(a); this assertion concerns $\widetilde{\mathrm{Op}}$ only, and it holds for every system on $T_2$ with W14$^{\mathrm{ch}}$, as a consequence of W14$^{\mathrm{ch}}$ at the translations by $\ell_1\mathbb{Z}^3$ and of the facts listed above.

\emph{Level $0$.}
The families $\Gamma_0$, $S_{0,y}$ and $H_0(s'',A)$ have $\mathfrak{S}^{\mathrm{ch}}(1,\ell_i)$ as source and target space, and there is no family $M_0$.
\end{named}

\begin{named}{Remark}
Proposition~26.1 corresponds to \cite[Section~19.3.1]{QEDI}, \cite[Definition~19.4]{QEDI} and \cite[Definition~5.13]{QEDI}.
It differs from them in the following points.
Fine spinors are functions on $\mathbb{L}^{\mathrm{f}}_k$, and the label is not affected by $\mathrm{pr}_i$, $\mathrm{pj}$, $\mathrm{pj}^{*}$ and $T_v$.
The kernels of operators whose source is a space $\mathfrak{S}^{\mathrm{ch}}$ are taken as in (I.1.4), whereas in \cite[Section~19.3.1]{QEDI} the weights $\eta_k^3$ are included in the kernel.
The fine pairing carries the weight $\eta_k^3/48$, and the factor $\mathrm{Jr}^3$ is the same.
At $k=0$ the unit lattice is $T^{\mathrm{f}}_0$.
\end{named}

\subsubsection*{Proposition~26.2 (cover cubes)}

\begin{named}{Proposition 26.2}[Cover cubes]
Let $k\ge0$ and $\ell\in\mathrm{Cl}$, let $\mathscr{Y}$ be a set of parameter cubes of one kind on $T_k=\mathbb{L}_k/\ell\mathbb{Z}^3$, and let $Y$ be its union.
Let $t$ be an index that belongs to this kind and satisfies $\operatorname{supp}t\subset Y$; the indices of the kind $\pi'_k$ are the fine bonds and the unit bonds, those of $\pi''_k$ the basis elements, and those of $\pi_k$ the unit sites and the fermion sites.
\begin{itemize}
\item[(a)] Any two cubes of $\mathscr{Y}$ that contain a point of $\operatorname{supp}t$ are equal or linked.
\item[(b)] The index $t$ is contained in exactly one link-component of $Y$.
\end{itemize}
Consequently an index that belongs to the kind is contained in no link-component of $Y$ if and only if its support is not contained in $Y$.
\end{named}
\emph{Conditions on the constants.}
$L$ is odd (C3), and $M$ is a power of $L$ with $M>4L$ (C5); this is the setting of (g0) and (g3) in \cite[Section~17.4.1]{QEDI}.

\begin{named}{Remark}
Proposition~26.2 corresponds to \cite[Lemma~19.3]{QEDI}.
The fermion sites are added to the indices of the kind $\pi_k$, with $\operatorname{supp}x=\{c\}$.
The lifts to $\mathbb{L}_k$ in the proof of \cite[Lemma~19.3]{QEDI} keep the difference vectors of the points of the support of the index: two lifts chosen within the sup-distance $2L\cdot M+L<\ell/2$ of each other differ by the unique representative, of sup norm less than $\ell/2$, of the difference of the two points on $T_k$.
In particular the lifted end points of a bond are $p$ and $p+a\cdot e_{\nu}$.
\end{named}

\subsubsection*{Proposition~26.3 (the projection keeps all properties: W18$^{\mathrm{ch}}$)}

\begin{named}{Proposition 26.3}[The projection keeps all properties: W18$^{\mathrm{ch}}$]
Let the data and the notation be those of Proposition~26.1, and let the system $\mathrm{Op}^{(2)}$ on $T_2$ have the properties W1--W17$^{\mathrm{ch}}$ of Section~G4.0.4 for
\[
\mathrm{par}^{0}=\left(\alpha_{0},\gamma_{\mathrm{w}},\kappa_{\mathrm{w}}+\kappa_{0};C_{w}^{(1)}\right).
\]
Clause~(a), together with (8.0), holds for every system on $T_2$ with W14$^{\mathrm{ch}}$, and it requires W14$^{\mathrm{ch}}$ only at the elements $(\mathbf{1}_4,v)$, $v\in\ell_1\mathbb{Z}^3$.
Clauses~(d) and~(e) hold without any assumption on a system.
\begin{itemize}
\item[(a)] The operators $\widetilde{\mathrm{Op}}(s,A)$ and $\widetilde{\mathrm{Op}}(X,A)$ commute with all $T_v$, $v\in\ell_1\mathbb{Z}^3$, and map $\mathrm{Per}$ into $\mathrm{Per}$.
Hence the projection is well defined.
Its kernels, taken as in Section~G4.0.4, are
\begin{equation}
\mathrm{Op}^{\mathrm{pj}}(\,\cdot\,)(t,t')=\sum_{w\in(\mathbb{Z}/\mathrm{Jr}\mathbb{Z})^3}\boldsymbol{\chi}(w)\cdot\widetilde{\mathrm{Op}}(\,\cdot\,)\left(t,t'+\ell_1\cdot w\right).\tag{8.0}
\end{equation}
The product $\boldsymbol{\chi}(w)\cdot\widetilde{\mathrm{Op}}(\,\cdot\,)\left(t,t'+\ell_1\cdot w\right)$ depends only on the class of $w$ modulo $\mathrm{Jr}$.
Indeed, for the spinor families each of the two factors changes its sign under $w\mapsto w+\mathrm{Jr}\cdot e_{\mu}$, and for (p1), (p2), where $\boldsymbol{\chi}=1$, each factor is a function of the class.
\item[(b)] $\mathrm{Op}^{\mathrm{pj}}$ is a system of primitives of level $k$ on $T_1$, with $S^{\mathrm{pj}}_{k,y}(X,A)$ Borel in $y$, and it has the properties W1--W17$^{\mathrm{ch}}$ for
\[
\mathrm{par}=\left(\alpha_{0},\gamma_{\mathrm{w}},\kappa_{\mathrm{w}};C_{w}\right),\qquad C_{w}=64K_{0}\cdot C_{w}^{(1)}.
\]
At $k=0$ the bounds W6 for $\Gamma_0^{\mathrm{pj}}$ and W7 for $S^{\mathrm{pj}}_{0,y}$ are bounds on the kernels (I.1.4) at pairs of fermion sites, the second one with the factor $(1+y)^{-1}$.
\item[(c)] At level $0$ the projections of the two diagonal families $H_{0}(s')$ and $H_{0}(s'',A)$ on $T_2$ are the diagonal families on $T_1$.
They have W3 and W5 with the constant
\[
4\exp(2\gamma_{\mathrm{w}})\cdot M^{27\alpha_{0}}\le C_{w}^{(1)},
\]
that is, without the factor $64$.
\item[(d)] For all indices $t$, $t'$ in the sense of Section~G4.0.4,
\begin{equation}
\sum_{w\in(\mathbb{Z}/\mathrm{Jr}\mathbb{Z})^3}\exp\left(-\gamma_{\mathrm{w}}\cdot d_2\left(t,t'+\ell_1\cdot w\right)\right)\le64\exp\left(-\gamma_{\mathrm{w}}\cdot d_1(t,t')\right).\tag{8.0$'$}
\end{equation}
\item[(e)] (Transposes.)
Let $O^{(2)}$ be an operator on $T_2$ that commutes with the $T_v$, $v\in\ell_1\mathbb{Z}^3$, and let $O^{(2)}\mathrm{pj}^{*}=\mathrm{pj}^{*}O^{(1)}$.
Then $\left(O^{(2)}\right)^{T}\mathrm{pj}^{*}=\mathrm{pj}^{*}\left(O^{(1)}\right)^{T}$, the transposes being formed with the pairings of Proposition~26.1.
\end{itemize}
The constants are $K_{0}$ (C0), $\gamma_{\mathrm{w}}$ (C4r), and $C_{w}^{(1)}$, $C_{w}$ (C6); none of them depends on $(k,\ell_1,\ell_2,\mathrm{Jr},\mu,e_k)$.
\end{named}
\emph{Conditions on the constants.}
$L$ is odd (C3);
$M$ is a power of $L$ with $M\ge L^2$, $M>4L$ and $\gamma_{\mathrm{w}}\cdot M\ge1$ (the three inequalities belong to (A0) of Section~I.8.4; all four conditions follow from (C5));
$M\ge M^{P}$ ((C5), (M6)), as in the setting of \cite[Proposition~15.36]{QEDI};
$0<\alpha_{0}\le1/4$ and $\kappa_{\mathrm{w}}\ge0$ (the definitions of the line (C4), with (C2) and (C0));
$K_{0}\ge1$ (C0);
$C_{w}^{(1)}\ge4\exp(2\gamma_{\mathrm{w}})\cdot M^{27\alpha_{0}}$ (the definition of $C_{w}^{(1)}$ in (C6)).
Clause~(b) requires in addition, at W1 and through Proposition~11\,(b), the conditions of Proposition~11, which are those of Proposition~2.8:
(L5), (M1), (M2), the part $M\ge108\pi\cdot C'_{V}\cdot a'$ of (M5), (Z2.s1), (Z2.s2), and $\alpha<1$.
Clause~(c), and through it W17 in clause~(b), requires through Proposition~8\,(b) the conditions of Proposition~8, which are those of Proposition~3: the same conditions, with (M5) in full.
These conditions belong to the lines (C1), (C3), (C5), (C8).

\begin{named}{Remark}
Proposition~26.3 corresponds to \cite[Theorem~19.5]{QEDI} and to property~18 of \cite[Section~5.4.2]{QEDI}.
The display (8.0) is the formula for the kernels in \cite[Theorem~19.5(a)]{QEDI}.
Clauses~(d) and~(e) correspond to two displays of the proof of \cite[Theorem~19.5]{QEDI}: the bound (8.0$'$) is the first display of the paragraph ``An estimate for sums over $(\mathbb{Z}/J\mathbb{Z})^3$'', and the identity of~(e) is the first display of the paragraph ``Transposition and periodic extension''; there the ratio $\mathrm{Jr}$ is written $J$, as in \cite[Section~19.3.1]{QEDI}.
The differences are the following.
The index sets, the kernels and the pairings are those of Proposition~26.1.
In W5 the short pairs are the pairs with $d_X<1$, the pairs inside a cluster included.
For $k\ge1$ the property W7 is formulated with the indicators of $\Delta_z\times O_{\mathrm{cub}}$.
At $k=0$ the properties W6 and W7 are bounds on the kernels (I.1.4); their constant contains, through $C^{F}$, the factor $48$ of (I.ke4), and this factor plays no role in the step from $\mathrm{par}^{0}$ to $\mathrm{par}$.
In W11 and in (c) one has $K_{x}=K_c$, and the diagonal family carries the same factor on the $48$ sites over $c$.
In W13 the gauge functions are functions of $c$.
In W14 the groups are, as in property~14 of \cite[Section~5.4.2]{QEDI}, $\mathbb{G}_k^{L}$ for (p1)--(p4) and $\mathbb{G}^{L}_{k,3}$ for (p5); they act on the fermion sites by (I.4.2).
At $k=0$ the transposes in W15 are formed with $\langle\,\cdot\,,\cdot\,\rangle^{\mathrm{ch}}$.
The formula (8.0) is asserted under W14$^{\mathrm{ch}}$ alone, whereas property~18 of \cite[Section~5.4.2]{QEDI} and \cite[Theorem~19.5(a)]{QEDI} assert it for a system with the properties~1--17 for $\mathrm{par}^{0}$.
In the following points Proposition~26.3 agrees with \cite{QEDI}.
The assertion of clause~(a) that the product in (8.0) depends only on the class of $w$ modulo $\mathrm{Jr}$ is also part of property~18 of \cite[Section~5.4.2]{QEDI}; \cite[Theorem~19.5(a)]{QEDI} does not state it.
As in \cite[Definition~5.29]{QEDI}, Borel measurability in $y$ is part of the definition of a system of level $0$.
The well-definedness of the projection is obtained from W14$^{\mathrm{ch}}$ alone, as it is obtained from property~14 alone in \cite[Definition~5.13]{QEDI}.
The numbers $27.3$ and $64$ of \cite[Theorem~19.5]{QEDI} and of its proof, and the relation $C_{w}=64K_{0}\cdot C_{w}^{(1)}$, are the same here, since the distances are those of photon sites.
\end{named}

\subsubsection*{Proposition~26.4 (rows well inside a component: (Lo)$^{\mathrm{ch}}$)}

\begin{named}{Proposition 26.4}[Rows well inside a component: (Lo)$^{\mathrm{ch}}$]
\emph{Definitions.}
Let $Z$ be a union of $M$-cubes.
A parameter family $s$ \emph{vanishes off} $Z$ if $s_{\mathrm{pc}}=0$ for every parameter cube $\mathrm{pc}$ that is not contained in $Z$.
The family $s\mathbf{1}_Z$ is given by $(s\mathbf{1}_Z)_{\mathrm{pc}}:=s_{\mathrm{pc}}$ if $\mathrm{pc}$ is contained in $Z$, and $(s\mathbf{1}_Z)_{\mathrm{pc}}:=0$ otherwise.
A \emph{component} of a union of $M$-cubes is a maximal subset that is connected through faces.

\emph{Statement.}
Let $(k,\ell,\mu,e_k)$ be (SF), and let a system of level $k$ on $T_k$ have the properties W8$^{\mathrm{ch}}$ and W12$^{\mathrm{ch}}$.
Let $\mathrm{Op}$ be one of the families (p1)--(p3) and $\Gamma_k$, let $Z$ be a union of $M$-cubes, let $s\in\mathrm{Pd}_k$ be a family that vanishes off $Z$, and let $Z_1$ be a component of $Z$.
Let $t$ be a row index of $\mathrm{Op}$ with $\square(t)^{\sim c_{l}}\subset Z_1$, where $c_{l}=2L$, and let $A,A_1\in\tilde{G}_k$ coincide on the bonds that meet $Z_1$.
Then
\[
\mathrm{Op}(s,A)(t,t')=0\ \text{ unless }t'\text{ is contained in }Z_1,\qquad
\mathrm{Op}(s,A)(t,\cdot\,)=\mathrm{Op}\left(s\mathbf{1}_{Z_1},A_1\right)(t,\cdot\,).
\]
The kernels of $M_k(X,A)$, $k\ge1$, and of $S_{k,y}(X,A)$ vanish unless both indices lie over $X$; and if $A$ and $A_1$ coincide on the bonds that meet $X$, then $M_k(X,A)=M_k(X,A_1)$ and $S_{k,y}(X,A)=S_{k,y}(X,A_1)$.
\end{named}
\emph{Conditions on the constants.}
$L$ is odd (C3), and $M$ is a power of $L$ with $M>4L$ (C5); this is the setting of (g0), (g1) and (g2) in \cite[Section~17.4.1]{QEDI}.

\begin{named}{Remark}
Proposition~26.4 corresponds to \cite[Section~19.3.4]{QEDI} with \cite[Lemma~19.6]{QEDI}, and to \cite[Proposition~5.30]{QEDI}.
The row indices of (p3), at every level, and those of $\Gamma_0$ are fermion sites, and $\square(x)=\square(c)$.
\end{named}

\subsubsection*{Proposition~26.5 (the systems on the top torus)}

\begin{named}{Proposition 26.5}[The systems on the top torus]
Let $(k,\ell,\mu,e_k)$ be (SF).
Consider the families (p1) and (p2) defined by the two displays for $H_{k}(s')$ and $C_{k}^{1/2}(s')$ in \cite[Section~15.10.3]{QEDI} (at $k=0$: the family (p2) and the diagonal family $H_{0}(s')$).
Together with them consider, for $k\ge1$, the families (p3), (p4), (p5) constructed in Proposition~7, and, for $k=0$, the families $\Gamma_0(s'',A)$, $S_{0,y}(X,A)$, $H_{0}(s'',A)$ of Proposition~8, in both cases taken at $\kappa'=\kappa_{\mathrm{w}}+\kappa_{0}$.
These families form a system of primitives of level $k$ on $T_k$ which has the properties W1--W17$^{\mathrm{ch}}$ for $\mathrm{par}^{0}$, and hence for $\mathrm{par}$.
The constants are
\begin{align*}
\alpha_{0}&:=\min\left(\alpha^{F},\alpha^{P}\right),\\
\gamma_{\mathrm{w}}&:=\min\left(\gamma^{F},\gamma^{P}\right),\\
C_{w}^{(0)}&:=\max\left(C^{F},C^{P}\right)\quad\text{at }\kappa'=\kappa_{\mathrm{w}}+\kappa_{0},\\
C_{w}^{(1)}&:=\max\left(C_{w}^{(0)},4\exp(2\gamma_{\mathrm{w}})\cdot M^{27\alpha_{0}}\right),
\end{align*}
as in the lines (C4), (C4r) and (C6) of Section~I.8.2.
\end{named}
\emph{Conditions on the constants.}
$M\ge M^{F}$ at $\kappa'=\kappa_{\mathrm{w}}+\kappa_{0}$ and $M\ge M^{P}$ (C5);
(Z2.s1)--(Z2.s4) (C8);
and the four conditions $0\le\alpha<1$ (C1), $L$ odd with $L\ge3$ (C3), $0<\alpha_{0}\le1/4$ (C4), and $M$ a power of $L$ with $M\ge L^2$ (C5).
The last four conditions are those of the setting stated at the beginning of \cite[Section~15.10]{QEDI} and in \cite[Theorem~15.35]{QEDI}; the theorem assumes in addition the condition $M\ge M^{P}$ of the first group.

\begin{named}{Remark}
Proposition~26.5 corresponds to \cite[Lemma~19.7]{QEDI}.
The fermionic families are those of Propositions~7 and~8.
In \cite[Lemma~19.7]{QEDI} the mass ranges over $(0,\mathfrak{m}_0]$, with the number $\mathfrak{m}_0$ of \cite[Definition~5.98]{QEDI}, whose value is $1$; here the symbol $\mathfrak{m}_0$ is not used and is replaced by this value (Section~I.6.1).
The number $e^{F}_0$ of \cite[Theorem~5.72]{QEDI} does not occur here (Proposition~14.1).
\end{named}

\subsubsection*{Proposition~26.6 (the data of the flow are data of one step)}

\begin{named}{Proposition 26.6}[The data of the flow are data of one step]
Let (S-4) hold, that is, (FL) without the primitives (Section~I.1.8).
Then the primitives named in (FL) are defined, by clause~(e), so that (FL) holds; clauses~(a) and~(b) are understood with these primitives.
\begin{itemize}
\item[(a)] The quadruple $(\bar{m},e,R_{\mathrm{top}},N)$ is (FL).
\item[(b)] At every level $k<k_{*}(R)$ the data $(k,\ell_k(R),\mu_k,e_k)$, together with the primitives of (FL), are (OS).
The next level has the side $\ell_{k+1}(R)=\ell_k(R)/L$, the mass $\mu_{k+1}=L\cdot\mu_k$ and the charge $e_{k+1}=L^{\frac{1}{2}}\cdot e_k$.
\item[(c)] At $k=k_{*}(R)$ the data $(k_{*},\ell_*,\mu_*,e_*)$ are (SF), and $k_{*}\ge1$.
\item[(d)] Let $R<R_{\mathrm{top}}$.
For every $k\le k_{*}(R)$ the sides $\ell_k(R)$ and $\ell_k(R_{\mathrm{top}})$ lie in $\mathrm{Cl}$, and $\ell_k(R_{\mathrm{top}})=\mathrm{Jr}\cdot\ell_k(R)$, where $\mathrm{Jr}=R_{\mathrm{top}}/R$ is a power of $L$, hence odd.
Both $(k,\ell_k(R),\mu_k,e_k)$ and $(k,\ell_k(R_{\mathrm{top}}),\mu_k,e_k)$ are (SF).
Thus the data of Proposition~26.1 are given at every level $k\le k_{*}(R)$.
\item[(e)] (The primitives of (FL) are defined.)
One has $R_{\mathrm{top}}\in\mathrm{Tor}(e)$, $N\ge N_{\min}(R_{\mathrm{top}},e)$ and $k_{*}(R)\le k_{\mathrm{top}}$.
For every $k<k_{\mathrm{top}}$ the quadruple $(k,\ell_k(R_{\mathrm{top}}),\mu_k,e_k)$ is (SF); therefore the systems of Proposition~26.5 on $T_k(R_{\mathrm{top}})$, $k<k_{\mathrm{top}}$, are defined, and they have the properties W1--W17$^{\mathrm{ch}}$ for $\mathrm{par}^{0}$.
If $R<R_{\mathrm{top}}$ and $k<k_{*}(R)$, then $k<k_{\mathrm{top}}$, the data of Proposition~26.1 are given by~(d), and the system of Proposition~26.5 on $T_k(R_{\mathrm{top}})$ has W14$^{\mathrm{ch}}$.
So its projection to $T_k(R)$, with $\mathrm{Jr}=R_{\mathrm{top}}/R$, is defined by Proposition~26.1 together with Proposition~26.3\,(a), and by Proposition~26.3\,(b) it is a system of primitives of level $k$ on $T_k(R)$.
Hence every quadruple with (S-4) carries the primitives of (FL).
They are objects formed from the quadruple, and no further datum enters; thus (S-4) and (FL) are the same condition on $(\bar{m},e,R,N)$.
\end{itemize}
No primitive occurs in~(c) and~(d).
\end{named}
In particular, statements whose data are (OS) apply at every level $k<k_{*}(R)$, on $R$ and on $R_{\mathrm{top}}$, and statements whose data are (SF) with $k\ge1$ apply at $k=k_{*}(R)$.

\begin{named}{Remark}
Proposition~26.6 corresponds to \cite[Proposition~19.8]{QEDI}, to the last paragraph of \cite[Section~19.3.5]{QEDI}, and to \cite[Theorem~5.79(c)]{QEDI}.
The numbers $\mathfrak{m}_0$, $m_0$ and $h'$ of \cite[Definition~5.98]{QEDI} occur in \cite[Section~19.3.5]{QEDI}, and $\mathfrak{m}_0$ occurs also in \cite[Theorem~5.79(c)]{QEDI}.
Their values are $\mathfrak{m}_0=1$, $m_0=0$ and $h'=h$; here the three symbols are not used, and these values stand in their place (Section~I.6.1).
Clause~(d) corresponds to the part of the opening paragraph of \cite[Section~19.4.1]{QEDI} in which the hypotheses of \cite[Section~19.3.1]{QEDI} are verified at every level $k\le k_{*}(R)$;
it gives, at every such level, the setting of \cite[Section~19.3.1]{QEDI} in which \cite[Section~19.4.1]{QEDI} works: $\ell_k(R)$ and $\ell_k(R_{\mathrm{top}})$ lie in $\mathrm{Cl}$, $\mathrm{Jr}=R_{\mathrm{top}}/R$ is a power of $L$ and hence odd, and the mass and the charge lie in their ranges.
Clause~(e) corresponds to the third paragraph of \cite[Definition~5.13]{QEDI}, which fixes the systems of primitives required in the last clause of \cite[Definition~5.70]{QEDI}, and to the passages on the primitives in the proofs of \cite[Proposition~19.8(a), (b)]{QEDI}.
In \cite{QEDI} the primitives of the scheme are fixed in a definition and form part of the hypothesis for the flow, and their availability is verified in that definition and in those proofs;
among the facts verified there are $R_{\mathrm{top}}\in\mathrm{Tor}(e)$, $N\ge N_{\min}(R_{\mathrm{top}},e)$, $k_{*}(R)\le k_{\mathrm{top}}$, the applicability of \cite[Lemma~19.7]{QEDI} on $T_k(R_{\mathrm{top}})$, the oddness of $R_{\mathrm{top}}/R$ and the well-definedness of the projection.
Here the hypothesis is (S-4), that is, (FL) without the primitives, and that the primitives are defined is asserted in clause~(e).
The last sentence of~(e) concludes that (S-4) and (FL) are the same condition on $(\bar{m},e,R,N)$;
this conclusion is not stated in \cite[Definitions~5.13 and~5.70]{QEDI} or in the proof of \cite[Proposition~19.8]{QEDI}.
\end{named}

\subsection*{27. Agreement of the steps on two tori}

\paragraph{Standing data (TT).}
In the Propositions 27.$m$ the following data are fixed; together they are referred to as (TT).
The quadruple of (FL) is given, with $R<R_{\mathrm{top}}$.
Put $k_m:=k_{*}(R)$; by (Sc2), $k_m\le k_{\mathrm{top}}$.
Let $k\le k_m$ be a level, and put
\[
\ell_1:=\ell_k(R),\qquad \ell_2:=\ell_k(R_{\mathrm{top}})=\mathrm{Jr}\cdot\ell_1,\qquad \mathrm{Jr}:=R_{\mathrm{top}}/R .
\]
The tori $T_i$, $T^{\mathrm{f}}_i$, $T^0_i$, the maps $\mathrm{pr}_i$, $\mathrm{pj}$, $\mathrm{pj}^{*}$ and the distances $d_i$ are those of Proposition~26.1; the data of that proposition are supplied by Proposition~26.6\,(d).
From (Sc3) and (G4.2.1) one obtains
\[
\ell_1\ge\ell_{k_m}(R)=M\cdot r_{*}>(31L+16)\cdot M .
\]
The numbers $e_k$, $\mu_k$, $b_k$, $\eta_k$, $p_{0,k}$, $\mathrm{rad}_1$, $\mathrm{rad}_2$, $\mathrm{rad}_4$, $w^{\mathrm{m}}_k$, $w^{\mathrm{E}}_k$ take the same values on both tori.
For $k<k_m$ the symbol $\mathrm{Op}^{(2)}$ denotes the system of Proposition~26.5 on $T_2$, and $\mathrm{Op}^{(1)}$ denotes its projection to $T_1$.
An object formed on $T_i$ carries the upper index $(i)$.
The upper indices of $T^{(j)}$, $\bar{T}^{(j)}$ in (I.5.4) and of $G^{(j)}$, $E^{(j),x}$, $E^{(j)}$ in Definition~24.D\,(D4), (D7) are not indices of tori: they number the two pieces of the operators $T_k(\,\cdot\,)$, $\bar{T}_k(\,\cdot\,)$ and the localised functions.
When such an object is marked with its torus, the index of the torus is written last, as in $G^{(j),(i)}$ and $E^{(j),x,(i)}$.
Data without slot on $T_i$ carry a lower index, as in $E_{(i)}$.
Finally, $d_U=2L+1$ and $c_l=2L$.

\subsubsection*{Proposition 27.1 (admissible sets; lifts of indices; transfer of relations)}

\begin{named}{Definition}[Admissible sets, flat polymers, indices, relations of range below $2d_U\cdot M$]
Let $\tilde{X}$ be a non-empty finite union of $M$-cubes of $\mathbb{L}_k$.
Denote by $\mathrm{box}(\tilde{X})$ the smallest product of three intervals that is a union of $M$-cubes and contains $\tilde{X}$, by $\mathrm{diam}(\tilde{X})$ the largest side of $\mathrm{box}(\tilde{X})$, and put
\[
N(\tilde{X}):=\mathrm{box}(\tilde{X})^{\sim 6d_U}.
\]
The set $\tilde{X}$ is called \emph{admissible} (at level $k$) if
\[
\mathrm{diam}(\tilde{X})\le\ell_1-14d_U\cdot M .
\]
A polymer $X$ of $T_i$ is called \emph{flat} if $X=\mathrm{pr}_i(\tilde{X})$ for some admissible $\tilde{X}$ that is connected through faces; such a set $\tilde{X}$ is a \emph{lift} of $X$.

In this context an \emph{index} is one of the following objects of the covering lattice, or its image on $T_i$: a site, a bond, a unit site, a unit bond, a basis element, a fermion site, a pair of $\mathrm{Pzt}_k$, a block $\mathrm{Bl}_{j}(z)$, an $M$-cube, an $LM$-cube or a parameter cube.
The support of an index is the set of its photon sites (for a bond and for a basis element $\Upsilon$ see Section~G4.0.4, for a pair Section~I.1.7).
A block or a cube is determined by its sort and by its centre $z$, and $\mathrm{pr}_i$ maps it to the block or cube of $T_i$ of the same sort whose centre is $\mathrm{pr}_i(z)$.
This extends the map $\mathrm{pr}_i$ of Proposition~26.1, which is defined on the index sets (I.i1)--(I.i5) of Section~I.1.9 only.

A \emph{relation of range below $2d_U\cdot M$} is a set $\mathrm{Rel}$ of $p$-tuples of indices of the covering lattice, for some $p\ge2$, with two properties: $\mathrm{Rel}$ is invariant under the common translations of all entries by the vectors of $\ell_1\mathbb{Z}^3$, and for every tuple of $\mathrm{Rel}$ the union of the supports of its entries has sup-diameter less than $2d_U\cdot M=(4L+2)\cdot M$.
For such a relation put
\[
\mathrm{Rel}_i:=\left\{(\mathrm{pr}_i t_1,\dots,\mathrm{pr}_i t_p):(t_1,\dots,t_p)\in\mathrm{Rel}\right\}.
\]
\end{named}

\begin{named}{Proposition 27.1}
Assume (TT), let $\tilde{X}$ be admissible, and let $i\in\{1,2\}$.
\begin{itemize}
\item[(a)] $N(\tilde{X})$ is a box of $M$-cubes all of whose sides are at most $\ell_1-2d_U\cdot M$.
Every non-empty union $\tilde{X}'$ of $M$-cubes with $\tilde{X}'\subset\tilde{X}$ is admissible, and $N(\tilde{X}')\subset N(\tilde{X})$.
\item[(b)] For every $v\in\ell_1\mathbb{Z}^3$ with $v\ne0$ the sets $N(\tilde{X})$ and $N(\tilde{X})+v$ are disjoint.
Consequently $\mathrm{pr}_i$ is injective on the sites of $N(\tilde{X})$ and on the indices contained in $N(\tilde{X})$.
\item[(c)] Let $u$, $u'$ be sites of $N(\tilde{X})$ with $d_i(\mathrm{pr}_i u,\mathrm{pr}_i u')<2d_U\cdot M$.
Then $|u-u'|=d_i(\mathrm{pr}_i u,\mathrm{pr}_i u')$, and $u-u'$ is the unique representative of its class modulo $\ell_i\mathbb{Z}^3$ whose sup norm is less than $2d_U\cdot M$.
\item[(d)] (Lifts of indices.) Let $\tilde{S}$ be a set of sites contained in $N(\tilde{X})$.
Every index of $T_i$ that is contained in $\mathrm{pr}_i(\tilde{S})$ has exactly one lift that is contained in $\tilde{S}$.
For a pair of $\mathrm{Pz}_k(T_i)$ this lift is the shortest lift of the pair, an element of $\mathrm{Pzt}_k$ with the same $d_X$ and the same six paths.
Moreover, $\mathrm{pr}_i(\tilde{X}^{\sim a})=(\mathrm{pr}_i\tilde{X})^{\sim a}$ for every integer $a\ge0$.
\item[(e)] (Transfer.) Let $\mathrm{Rel}$ be a relation of range below $2d_U\cdot M$, and let $t_1,\dots,t_p$ be indices contained in $N(\tilde{X})$.
Then $(t_1,\dots,t_p)\in\mathrm{Rel}$ if and only if $(\mathrm{pr}_i t_1,\dots,\mathrm{pr}_i t_p)\in\mathrm{Rel}_i$.
\item[(f)] (The relations in use.) Each of the relations listed below is a relation of range below $2d_U\cdot M$ and is invariant under the translations by $(L\cdot M)\mathbb{Z}^3$; for each of them, $\mathrm{Rel}_i$ is the relation that defines the corresponding object on $T_i$ (every extent that occurs is less than $(4L+2)\cdot M<\ell_i/2$).
\begin{itemize}
\item Nearest neighbours of $\mathbb{L}_k$ and of $\mathbb{Z}^3$.
\item $x'=x^{(j)}$, for $j=0,1,2,3$ (see (I.1.2)).
\item $\pi(x)=\pi(x')$.
\item $(x,x')$ is a pair of $\mathrm{Pzt}_k$.
\item The bond $b$ lies on the path $\boldsymbol{\Gamma}^{\pi}(c,c')$ of the pair $(x,x')$.
\item $c$ lies in $\Delta_y$, in $\mathrm{Bl}_{j}(z)$, in an $M$-cube, in an $LM$-cube, in a parameter cube.
\item $y$ lies in $B(y')$.
\item $z$ is the centre of the $M$-cube $\square$.
\item Two $M$-cubes have a common face.
\item Two parameter cubes of one kind are linked.
\item A parameter cube meets an $M$-cube.
\item A bond belongs to the support of $\Upsilon$.
\end{itemize}
\end{itemize}
\end{named}

\emph{Conditions on the constants.}
The inequality $\ell_1>(31L+16)\cdot M$ of the standing data (TT) rests on the condition (G4.2.1), that is, on the condition (2.1) of Proposition~14.2\,(c),
\[
r_{*}>31L+16,
\]
since $\ell_1\ge\ell_{k_m}(R)=M\cdot r_{*}$.
The condition (G4.2.1) is a condition on the data $(R,e,N)$, not on a constant; it holds for all data (FL), hence under (TT).
It follows through (Sc6) of Section~G4.0.2 from the row (e1) of Proposition~14.2\,(b), the condition on the number $e_1$ of the line (C9), together with the datum $e\cdot h^{\frac{1}{2}}\le e_1$.
Indeed, with $e_K=e\cdot h^{\frac{1}{2}}$ this datum gives $v_K=-\log e_K\ge-\log e_1$; by (Sc3), $v_{*}\ge v_K$ and $M\cdot v_{*}^{r}\le\ell_{*}$, and $\ell_{*}=M\cdot r_{*}$ by (Sc1); and the row (e1), which holds for all $v\ge-\log e_1$ by Proposition~14.3\,(f), gives at $v=v_{*}$ its second inequality $v_{*}^{r}>31L+16$, so that $r_{*}\ge v_{*}^{r}>31L+16$.
Of the rows (e1)--(e8) of Proposition~14.2\,(b), which concern the number $e_1$, this derivation uses the row (e1) alone.
Its second inequality enters at $v=v_{*}$, as above.
Its first and third inequalities enter at $v=v_K$, where they read $v_K\ge8$ and $\left(1+\frac{3\,\mathrm{lg}}{v_K}\right)^{r}<L$;
since $r>0$, they give the hypotheses $v_K\ge1$ and $\left(1+\frac{\mathrm{lg}}{v_K}\right)^{r}\le L$ of \cite[Proposition~19.1]{QEDI}, whose clauses (c) and (e) are (Sc1) and (Sc3) of Section~G4.0.2.
Through (Sc1) and (Sc3) there enter, moreover, the condition $r>0$, which follows from $r\ge2$ (line (C7)), and the condition that $M$ is a power of $L$, which is part of the row (M1) of Proposition~14.2\,(b); both are used in the proof of \cite[Proposition~19.1]{QEDI} (Chapter~26b, Section~3.2; Annex~B, Section~G4.A.6).
The proof of Proposition~27.1 rests, besides, directly on the fact (0.1c) of Chapter~27a: $L$ is odd, by the row (L5) of Proposition~14.2\,(b), and $M$ is a power of $L$, by the row (M1);
the same row gives, moreover, $M\ge L^2$, and in particular $M\ge1$.
These two conditions, $L$ odd and $M$ a power of $L$, make $a/\eta_k$ an odd integer for $a=L^j\cdot\eta_k$ ($0\le j\le k+1$), $M$ and $L\cdot M$;
the quotient $M/L$ is a power of $L$ as well, and the inequality $M\ge L^2$ of the row (M1) makes it at least $L$, so that it is an odd integer (for a power $M$ of $L$, already $M>1$ would suffice for this);
hence the blocks and the cubes of each of these sizes partition $\mathbb{L}_k$ and are nested ((0.2a), (0.2b) of Chapter~27a);
moreover $\eta_k=L^{-k}\le1\le M$, the second inequality being $M\ge1$.
Through Proposition~26.6\,(d), which supplies in (TT) the data of Proposition~26.1, among them $\ell_1,\ell_2\in\mathrm{Cl}$ and $\mathrm{Jr}$ odd, and through (Sc4), on which that clause rests, there enter further the condition $L\ge13$, which follows from the row (L1) of Proposition~14.2\,(b), and the condition $e_1\le e_0$ of the line (C9).
The conditions on the constants that enter Propositions~27.1--27.4, among them those that enter through the Propositions applied in their proofs, are set out in Chapter~27a, in the section headed ``Conditions on the constants''; see also Annex~B, Section~G4.A.6.

\begin{named}{Remark}
The Definition above and Proposition~27.1 correspond to the paragraph that opens \cite[Section~19.4.1]{QEDI} and to \cite[Definition~19.9(a), (b)]{QEDI}, together with the displays (a1), (a2) printed there.
In three features the notion of a relation of range below $2d_U\cdot M$ and the transfer principle, clause (e), agree with \cite[Definition~19.9(a)]{QEDI}.
First, in both papers the relations are relations among $p\ge2$ objects; here they are the sets $\mathrm{Rel}$ of $p$-tuples of indices.
Second, in both papers a relation is required to be invariant only under the common translations by $\ell_1\mathbb{Z}^3$, and for every tuple of the relation any two points of its entries (of their supports) have sup distance less than $2d_U\cdot M$.
Third, in both papers the relation on $T_i$ is the image under $\mathrm{pr}_i$ of the relation on the covering lattice; here it is $\mathrm{Rel}_i$, the image of $\mathrm{Rel}$.
The path relation of (f) is a relation among three indices; \cite[Definition~19.9(a)]{QEDI} lists it as well among the relations of range less than $2d_U\cdot M$, as a relation among a bond and two sites, so that it falls under the notion of both papers.
They differ in the following points.
Fermion sites and the pairs of $\mathrm{Pzt}_k$, pairs inside a cluster included, are indices, and the list (f) contains the relations $x'=x^{(j)}$ and $\pi(x)=\pi(x')$.
In \cite[Definition~19.9(a)]{QEDI} the transfer principle stands, together with its proof, inside the definition; here it is clause (e) of a proposition.
Existence and uniqueness of the lift of an index whose support has several points is clause (d).
In clause (a) the unions of $M$-cubes contained in $\tilde{X}$ are required to be non-empty; the corresponding sentence at the end of \cite[Definition~19.9(c)]{QEDI} contains no such condition.
The arithmetic facts on which the setting rests, namely $\ell_1,\ell_2\in\mathrm{Cl}$ and $\mathrm{Jr}$ odd, are Proposition~26.6\,(d).
\end{named}

\subsubsection*{Proposition 27.2 (generators near a flat set; local families; agreement; criterion; small polymers; sub-objects)}

\begin{named}{Definition}[Equal lifts, generators, local families, agreement]
Assume (TT).
\begin{itemize}
\item[(D1)] (Equal lifts.) Bond functions $B_1$ on $T_1$ and $B_2$ on $T_2$ have \emph{equal lifts} on a set $\tilde{S}$ of bonds of $\mathbb{L}_k$ if $B_1(\mathrm{pr}_1 b)=B_2(\mathrm{pr}_2 b)$ for every $b\in\tilde{S}$.
Equal lifts \emph{on a set $Z$ of sites} means equal lifts on the bonds contained in $Z$.
Elements $\mathsf{Z}^{(i)}$ of $V_k(T_i)$ have equal lifts on $Z$ if, for every basis element $\Upsilon$ contained in $Z$, the coordinate of $\mathsf{Z}^{(1)}$ at $\mathrm{pr}_1(\Upsilon)$ equals the coordinate of $\mathsf{Z}^{(2)}$ at $\mathrm{pr}_2(\Upsilon)$.
Parameter families on $T_1$ and on $T_2$ have equal lifts on $Z$ if, for every parameter cube $\widetilde{\mathrm{pc}}$ contained in $Z$, the value of the first at $\mathrm{pr}_1(\widetilde{\mathrm{pc}})$ equals the value of the second at $\mathrm{pr}_2(\widetilde{\mathrm{pc}})$.
Real functions $\Lambda_1$, $\Lambda_2$ of the sites of $T_1$, $T_2$ (gauge functions) have equal lifts on $Z$ if $\Lambda_1(\mathrm{pr}_1 c)=\Lambda_2(\mathrm{pr}_2 c)$ for every site $c$ of $Z$.
\item[(D2)] (Generators.) The generators on $T_i$ are indexed by the points of the covering index sets, as in (I.AP) of Section~I.3.6.
For an admissible $\tilde{X}$ let $\mathrm{Gr}_i(\tilde{X})$ be the subalgebra generated by the following generators of $T_i$:
$\psi(x,\beta,\omega)$, with $x$ over $N(\tilde{X})$;
$\chi(x,x',\beta,\omega)$, with $(x,x')\in\mathrm{Pzt}_k$ and $c,c'\in N(\tilde{X})$;
for $k\ge1$, $\Psi(y,\beta,\omega)$ and $W(y,\beta,\omega)$, with $y$ a unit site of $N(\tilde{X})$;
for $k=0$, $W(x,\beta,\omega)$, with $x$ over $N(\tilde{X})$.
Let
\[
\mathrm{idn}=\mathrm{idn}[\tilde{X}]\colon\mathrm{Gr}_1(\tilde{X})\to\mathrm{Gr}_2(\tilde{X})
\]
be the homomorphism of algebras that sends each of these generators to the generator of $T_2$ with the same index.
\item[(D3)] (Local families.) $\mathrm{Dm}$ stands for $\tilde{R}_k\times\mathfrak{X}_k$ or for $\tilde{R}_k^{-}\times\mathfrak{X}_k$.
A \emph{local family} on $T_i$ is a family
\[
F=\left\{F(X;\mathscr{A},\mathscr{C}):X\in D_k(T_i),\ (\mathscr{A},\mathscr{C})\in\mathrm{Dm}(T_i)\right\}
\]
of even elements of $\mathrm{Gr}[\psi,\chi]$ with the locality property (I.k2) of Section~I.5.4.
$\mathrm{Dm}^{\mathrm{o}}$ stands for $\tilde{R}_k$ or for $\tilde{R}_k^{-}$.
A \emph{local family without slot} on $T_i$ is a family
\[
F=\left\{F(X;\mathscr{A}):X\in D_k(T_i),\ \mathscr{A}\in\mathrm{Dm}^{\mathrm{o}}(T_i)\right\}
\]
of even elements of $\mathrm{Gr}[\psi,\chi]$ with the property (k2)$^{\mathrm{o}}$ of the Definition in Proposition~48.1, read on $\mathrm{Dm}^{\mathrm{o}}$.
\item[(D4)] (Agreement.) Local families $F^{(1)}$ on $T_1$ and $F^{(2)}$ on $T_2$ \emph{agree on flat polymers}, in symbols $F^{(1)}\sim F^{(2)}$, if
\[
\mathrm{idn}[\tilde{X}]\,F^{(1)}(\mathrm{pr}_1\tilde{X};\mathscr{A}_1,\mathscr{C}_1)=F^{(2)}(\mathrm{pr}_2\tilde{X};\mathscr{A}_2,\mathscr{C}_2)
\]
for every connected admissible $\tilde{X}$ and all $(\mathscr{A}_i,\mathscr{C}_i)\in\mathrm{Dm}(T_i)$ that have equal lifts on $N(\tilde{X})$.
Families $\mathfrak{P}^{(i)}=\left\{\mathfrak{P}^{(i)}(X,\,\cdot\,):X\in D_k(T_i)\right\}$ of complex functions on $\mathfrak{X}_k(T_i)$ \emph{agree}, in symbols $\mathfrak{P}^{(1)}\sim\mathfrak{P}^{(2)}$, if
\[
\mathfrak{P}^{(2)}(\mathrm{pr}_2\tilde{X},\mathrm{pj}^{*}\mathscr{C})=\mathfrak{P}^{(1)}(\mathrm{pr}_1\tilde{X},\mathscr{C})
\]
for every connected admissible $\tilde{X}$ and every $\mathscr{C}\in\mathfrak{X}_k(T_1)$.
Local families without slot $F^{(1)}$ on $T_1$ and $F^{(2)}$ on $T_2$ \emph{agree on flat polymers}, $F^{(1)}\sim F^{(2)}$, if
\[
\mathrm{idn}[\tilde{X}]\,F^{(1)}(\mathrm{pr}_1\tilde{X};\mathscr{A}_1)=F^{(2)}(\mathrm{pr}_2\tilde{X};\mathscr{A}_2)
\]
for every connected admissible $\tilde{X}$ and all $\mathscr{A}_i\in\mathrm{Dm}^{\mathrm{o}}(T_i)$ that have equal lifts on $N(\tilde{X})$.
\end{itemize}
\end{named}

\begin{named}{Proposition 27.2}
Assume (TT).
\begin{itemize}
\item[(a)] The generators listed in (D2) are independent on $T_i$; therefore $\mathrm{idn}[\tilde{X}]$ is an isomorphism.
If $\tilde{X}$ is contained in the admissible set $\tilde{X}'$, then $\mathrm{Gr}_i(\tilde{X})\subset\mathrm{Gr}_i(\tilde{X}')$, and $\mathrm{idn}[\tilde{X}]$ is the restriction of $\mathrm{idn}[\tilde{X}']$.
The map $\mathrm{idn}$ preserves the bidegree and all norms (I.5.1), and it commutes with the maps (I.T-G), (I.T-C), (I.T-K) of Section~I.4.5 formed with gauge functions that have equal lifts on $N(\tilde{X})$.
\item[(b)] Let $F$ be a local family on $T_i$, and let $X=\mathrm{pr}_i(\tilde{X})$ be flat.
Then $F(X;\mathscr{A},\mathscr{C})$ lies in $\mathrm{Gr}_i(\tilde{X})$.
In detail: a fermion site over $X$ has exactly one lift over $\tilde{X}$; a pair that is admissible for $X$ has exactly one lift whose two photon sites lie in $\tilde{X}^{\sim1}$; this lift is the shortest lift of the pair, and it has for $\tilde{X}$ the kinds that the pair has for $X$.
If the pair is of kind (I) for $X$, the lift is of kind (I) for $\tilde{X}$ (both photon sites and the six paths lie in $\tilde{X}$).
If the pair is of kind (II) for $X$, the lift is of kind (II) for $\tilde{X}$ ($x'=x^{(3)}$, and $c$ or $c_{3}$ lies in $\tilde{X}$).
Hence, if the pair satisfies both conditions for $X$, its lift satisfies both for $\tilde{X}$.
\item[(c)] (Criterion.) Local families $F^{(1)}$ and $F^{(2)}$ agree if and only if
\[
\mathrm{idn}\,F^{(1)}(\mathrm{pr}_1\tilde{X};\mathscr{A},\mathscr{C})=F^{(2)}(\mathrm{pr}_2\tilde{X};\mathrm{pj}^{*}\mathscr{A},\mathrm{pj}^{*}\mathscr{C})
\]
for every connected admissible $\tilde{X}$ and all $(\mathscr{A},\mathscr{C})\in\mathrm{Dm}(T_1)$.
In that case
\[
\mathrm{idn}\,F^{(1)}(\mathrm{pr}_1\tilde{X};\mathscr{A}_1,\mathscr{C}_1)=F^{(2)}(\mathrm{pr}_2\tilde{X};\mathscr{A}_2,\mathscr{C}_2)
\]
whenever the pairs $(\mathscr{A}_i,\mathscr{C}_i)\in\mathrm{Dm}(T_i)$ have equal lifts on the bonds that meet $\tilde{X}$.
\item[(d)] (Small polymers.) Every connected non-empty union $\tilde{X}$ of $M$-cubes of $\mathbb{L}_k$ with $d_M(\tilde{X})\le L$ is admissible.
Every small polymer of $T_i$ is flat.
For a fixed $M$-cube $\tilde{\square}$ of $\mathbb{L}_k$, the map $\tilde{X}\mapsto\mathrm{pr}_i(\tilde{X})$ is a bijection from the set of the connected $\tilde{X}$ with $d_M(\tilde{X})\le L$ that contain $\tilde{\square}$ onto the set of the small polymers of $T_i$ that contain $\mathrm{pr}_i(\tilde{\square})$.
For a connected admissible $\tilde{X}$ the following three conditions are equivalent: $\mathrm{pr}_1(\tilde{X})$ is small; $d_M(\tilde{X})\le L$; $\mathrm{pr}_2(\tilde{X})$ is small.
When they hold, $d_M(\mathrm{pr}_i\tilde{X})=d_M(\tilde{X})$ and $\operatorname{Vol}(\mathrm{pr}_i\tilde{X})=\operatorname{Vol}(\tilde{X})$.
\item[(e)] (Sub-objects.) Let $\tilde{Y}$ be admissible.
The map $\mathrm{pr}_i$ induces bijections
from the non-empty connected unions of $M$-cubes contained in $\tilde{Y}$ onto the polymers of $T_i$ contained in $\mathrm{pr}_i(\tilde{Y})$;
from the parameter cubes of each kind contained in $\tilde{Y}$ onto the parameter cubes of that kind contained in $\mathrm{pr}_i(\tilde{Y})$;
from the unit sites, the fermion sites and the basis elements contained in (over) $\tilde{Y}$ onto those contained in (over) $\mathrm{pr}_i(\tilde{Y})$.
The map $\mathrm{pr}_i$ is injective on $\tilde{Y}^{\sim(L+1)}$, which is a subset of $N(\tilde{Y})$.

The sets $P_0$, $\mathrm{Yp}$, $Y'$ are those of Definition~24.D\,(D0) (see \cite[Section~17.4.3]{QEDI}).
For a non-empty connected union $\tilde{X}$ of $M$-cubes of $\mathbb{L}_k$ and a set $\tilde{P}$ of parameter cubes of $\mathbb{L}_k$, the sets $P_0(\tilde{X})$, $\mathrm{Yp}(\tilde{X},\tilde{P})$, $Y'(\tilde{X},\tilde{P})$ are formed by the wording of Definition~24.D\,(D0), with the parameter cubes and the $M$-cubes of $\mathbb{L}_k$.
The set $P_0$ is formed with ``the kinds that occur'', and the assertions below hold for every choice of a set of kinds.
Let $\tilde{X}$, $\tilde{X}_1$, $\tilde{X}_2$ be connected non-empty unions of $M$-cubes contained in $\tilde{Y}$, and let $\tilde{P}$ be a set of parameter cubes that are contained in $\tilde{Y}$ and do not belong to $P_0(\tilde{X})$.
Then
\begin{gather*}
P_0(\mathrm{pr}_i\tilde{X})=\mathrm{pr}_i P_0(\tilde{X}),\\
Y'(\tilde{X},\tilde{P})\subset\tilde{Y}\quad\text{and}\quad\mathrm{pr}_i Y'(\tilde{X},\tilde{P})=Y'(\mathrm{pr}_i\tilde{X},\mathrm{pr}_i\tilde{P}),\\
\bar{\tilde{X}}^{L}\subset\tilde{Y}^{\sim(L-1)}\quad\text{and}\quad\mathrm{pr}_i\bar{\tilde{X}}^{L}=\overline{\mathrm{pr}_i\tilde{X}}^{L},
\end{gather*}
where $\bar{\tilde{X}}^{L}$ is the $LM$-closure of $\tilde{X}$.
Further, $\tilde{X}_1\cup\tilde{X}_2$ is connected if and only if $\mathrm{pr}_i\tilde{X}_1\cup\mathrm{pr}_i\tilde{X}_2$ is connected, so that incompatibility is preserved; and inclusions, unions, the relation ``meets'', common faces and links are preserved in both directions.
Finally, if $Y'(X,P)=\mathbf{Y}$, then $X$ and every cube of $P_0(X)\cup P$ are contained in $\mathbf{Y}$.
\item[(f)] If $\mathscr{E}^{(1)}\sim\mathscr{E}^{(2)}$ (families with slot), then the families $\mathscr{E}^{(i)}|_0$ are local families without slot, and $\mathscr{E}^{(1)}|_0\sim\mathscr{E}^{(2)}|_0$.
\item[(g)] (Without slot.) Let $F$ be a local family without slot on $T_i$, and let
\[
F^{\mathrm{c}}(X;\mathscr{A},\mathscr{C}):=F(X;\mathscr{A})\qquad\text{for }(\mathscr{A},\mathscr{C})\in\mathrm{Dm}^{\mathrm{o}}(T_i)\times\mathfrak{X}_k(T_i)
\]
be its extension that is constant in $\mathscr{C}$.
Then $F^{\mathrm{c}}$ is a local family in the sense of (D3), with $\mathrm{Dm}=\mathrm{Dm}^{\mathrm{o}}\times\mathfrak{X}_k$: it does not depend on $\mathscr{C}$, so that (I.k2) for $F^{\mathrm{c}}$ is (k2)$^{\mathrm{o}}$ for $F$, read on $\mathrm{Dm}^{\mathrm{o}}$ (for $\mathrm{Dm}^{\mathrm{o}}=\tilde{R}_k$ this is Proposition~48.1\,(b)).
Therefore (b) holds with $\mathscr{C}$ deleted: for a flat $X=\mathrm{pr}_i(\tilde{X})$ the element $F(X;\mathscr{A})=F^{\mathrm{c}}(X;\mathscr{A},0)$ lies in $\mathrm{Gr}_i(\tilde{X})$, and the last display of (D4) is meaningful.
One has $F^{(1)}\sim F^{(2)}$ if and only if $F^{(1)\mathrm{c}}\sim F^{(2)\mathrm{c}}$.
Therefore (c) holds with $\mathscr{C}$, $\mathscr{C}_i$ deleted and with $\mathrm{Dm}^{\mathrm{o}}$ in the place of $\mathrm{Dm}$: $F^{(1)}\sim F^{(2)}$ if and only if
\[
\mathrm{idn}\,F^{(1)}(\mathrm{pr}_1\tilde{X};\mathscr{A})=F^{(2)}(\mathrm{pr}_2\tilde{X};\mathrm{pj}^{*}\mathscr{A})
\]
for every connected admissible $\tilde{X}$ and every $\mathscr{A}\in\mathrm{Dm}^{\mathrm{o}}(T_1)$; and in that case $\mathrm{idn}\,F^{(1)}(\mathrm{pr}_1\tilde{X};\mathscr{A}_1)=F^{(2)}(\mathrm{pr}_2\tilde{X};\mathscr{A}_2)$ whenever $\mathscr{A}_i\in\mathrm{Dm}^{\mathrm{o}}(T_i)$ have equal lifts on the bonds that meet $\tilde{X}$.
\end{itemize}
\end{named}

\begin{named}{Remark}
The Definition above and Proposition~27.2 correspond to \cite[Definition~19.9(c)--(e)]{QEDI}, to the paragraph on families without the second connection that follows that definition in \cite[Section~19.4.1]{QEDI}, and to \cite[Lemma~19.10]{QEDI}.
They differ in the following points.
The generators $\psi((c,g),\,\cdot\,)$ are indexed by fermion sites, and the generators $\chi$ by the pairs of $\mathrm{Pzt}_k$, pairs with $c=c'$ included; at $k=0$ the generators $W$ carry the label $g$, and there is no separate family $\Psi$.
The kinds (I), (II) are those of Section~I.1.7; kind (II) means $x'=x^{(3)}$.
Clause (b) states in addition that the lifted pair is the shortest lift.
(D1) defines equal lifts on a set of sites for connections, for elements $\mathsf{Z}$ of $V_k$ and for parameter families.
The last sentence of \cite[Lemma~19.10(c)]{QEDI} is written in (e) as identities.
Clause (d) includes the admissibility of the connected $\tilde{X}$ with $d_M(\tilde{X})\le L$.
All unions of $M$-cubes are non-empty.
As to notation, the isomorphism $\mathfrak{i}$ of \cite[Definition~19.9(c)]{QEDI} is written $\mathrm{idn}$ here, because $\mathfrak{i}(A)=e_k\cdot|\mathrm{Im}\,A|_{\infty}$ by (I.1.6); the admissible set $\tilde{V}$ of \cite[Lemma~19.10(c)]{QEDI} is written $\tilde{Y}$.
\end{named}

\subsubsection*{Proposition 27.3 (primitives near a flat set)}

\begin{named}{Proposition 27.3}
Assume (TT) with $k<k_m$.
Let $\tilde{Y}$ be admissible, and put $Y_i:=\mathrm{pr}_i(\tilde{Y})$.
Let $s^{(1)}\in\mathrm{Pd}_k(T_1)$ vanish off $Y_1$, and let $s^{(2)}$ be the parameter family on $T_2$ given by
\[
s^{(2)}_{\mathrm{pr}_2(\widetilde{\mathrm{pc}})}:=s^{(1)}_{\mathrm{pr}_1(\widetilde{\mathrm{pc}})}
\]
for the parameter cubes $\widetilde{\mathrm{pc}}$ contained in $\tilde{Y}$, and by the value $0$ at every other cube.
Let $\mathrm{Op}$ be one of the families (p1)--(p3), $\Gamma_k$; let $j$ be an integer with $0\le j\le5d_U$; let $t$ be a row index of $\mathrm{Op}$ on $T_1$ that is contained in $Y_1^{\sim j}$, and let $\tilde{t}$ be its lift contained in $\tilde{Y}^{\sim j}$ (Proposition~27.1\,(d)).
Then the following holds for every $A^{(1)}\in\tilde{G}_k(T_1)$.
\begin{itemize}
\item[(a)] $\mathrm{Op}^{(1)}(s^{(1)},A^{(1)})(t,t')=0$ unless $t'$ is contained in $Y_1^{\sim(j+c_l)}$, and the row at $t$ depends on $A^{(1)}$ only through its values on the bonds contained in $Y_1^{\sim(j+d_U)}$.
\item[(b)] Let $A^{(2)}\in\tilde{G}_k(T_2)$ have the same lift as $A^{(1)}$ on the bonds contained in $\tilde{Y}^{\sim(j+d_U)}$.
Then
\[
\mathrm{Op}^{(1)}(s^{(1)},A^{(1)})(\tilde{t},\tilde{t}')=\mathrm{Op}^{(2)}(s^{(2)},A^{(2)})(\tilde{t},\tilde{t}')
\]
for every column index $\tilde{t}'$ contained in $N(\tilde{Y})$, and $\mathrm{Op}^{(2)}(s^{(2)},A^{(2)})(\tilde{t},t'')=0$ for the indices $t''$ of $T_2$ that are not contained in $Y_2^{\sim(j+c_l)}$.
\item[(c)] Let $\tilde{X}$ be a non-empty connected union of $M$-cubes contained in $\tilde{Y}$, let $\tilde{t}$, $\tilde{t}'$ be indices over $\tilde{X}$, and let $A^{(i)}\in\tilde{G}_k(T_i)$ have equal lifts on the bonds that meet $\tilde{X}$.
Then
\[
S^{(1)}_{k,y}(\mathrm{pr}_1\tilde{X},A^{(1)})(\tilde{t},\tilde{t}')=S^{(2)}_{k,y}(\mathrm{pr}_2\tilde{X},A^{(2)})(\tilde{t},\tilde{t}')\qquad\text{for all }y\ge0,
\]
and, for $k\ge1$,
\[
M_k^{(1)}(\mathrm{pr}_1\tilde{X},A^{(1)})(\tilde{t},\tilde{t}')=M_k^{(2)}(\mathrm{pr}_2\tilde{X},A^{(2)})(\tilde{t},\tilde{t}').
\]
\end{itemize}
Kernels at indices of the covering lattice are taken as in Section~G4.0.4.
At $k=0$ the row and column indices of (p3), $\Gamma_0$ and $S_{0,y}$ are fermion sites, and the kernels are those of (I.1.4).
By the definition of $\bar{H}_k$, clauses (a) and (b) hold for $\bar{H}_k(s'',\,\cdot\,)$ as well.
\end{named}

\begin{named}{Remark}
Proposition~27.3 corresponds to \cite[Lemma~19.11]{QEDI}.
The index sets and the kernels are those of Section~G4.0.4: the indices of spinors on the fine lattice are fermion sites; the kernel of an operator whose source is a space of spinors on fermion sites is taken by (I.1.4), while for an operator whose source is a space of spinors on unit sites the kernel is defined by the plain sum over the unit sites, $(O v)(t)=\sum_{y'}O(t,y')\,v(y')$.
The identity for $M_k$ is asserted for $k\ge1$ only, there being no family $M_0$.
The lift $\tilde{t}$ of the row index is the one given by Proposition~27.1\,(d).
\end{named}

\subsubsection*{Proposition 27.4 (exact operators, forms, homomorphisms and integrals near a flat set)}

\begin{named}{Proposition 27.4}
Assume (TT) with $k<k_m$.
Let $\tilde{Y}$ be admissible, let $\tilde{\mathbf{Y}}$ be a non-empty connected union of $M$-cubes contained in $\tilde{Y}$, put $\mathbf{Y}_i:=\mathrm{pr}_i(\tilde{\mathbf{Y}})$, and write $\mathrm{idn}:=\mathrm{idn}[\tilde{Y}]$.
\begin{itemize}
\item[(a)] (Operators.) Let $B_i\in\tilde{G}_k(T_i)$, $i=1,2$, have equal lifts on the bonds that meet $\tilde{\mathbf{Y}}$.
All indices are written in $N(\tilde{Y})$, in the following sense: the row of an operator $O$ on $T_i$ at a covering row index $t$ is written by the coefficients $\mathrm{ro}(t,t')$ if
\[
(O f)(t)=\sum_{t'}\mathrm{ro}(t,t')\,f(t')
\]
for every $f$ in the source space on $T_i$, the sum extending over the indices $t'$ contained in $N(\tilde{Y})$; by Proposition~27.1\,(b), the coefficients $\mathrm{ro}(t,t')$ are determined by $O$ and $t$.
Written in this way, the rows and kernels listed below coincide for $i=1$ and $i=2$, vanish at the columns of $T_i$ that have no lift in $N(\tilde{Y})$, and are the rows and kernels of the operators $\mathrm{Op}_{\mathbb{L}}$ of Section~I.1.9 on the covering lattice.
\begin{itemize}
\item[(a-1)] The rows $C(b',\,\cdot\,)$ at the unit bonds $b'$ contained in $\tilde{\mathbf{Y}}^{\sim(1+2d_U)}$ (no connection enters; range less than $2L$).
\item[(a-2)] The rows of $\delta_{\alpha,B_i}$ and of $\delta_{\alpha,-B_i}$ at the pairs that are admissible for polymers contained in $\mathbf{Y}_i$ (range $M$; they read $B_i$ on bonds that meet $\tilde{\mathbf{Y}}$).
\item[(a-3)] For $k\ge1$: the rows of $T^{(1)}(B_i)$, $T^{(2)}(B_i)$, $\bar{T}^{(1)}(B_i)$, $\bar{T}^{(2)}(B_i)$ of (I.5.4), formed with the polymer $\mathbf{Y}_i$ or with a polymer $\mathrm{pr}_i(\tilde{X})$, where $\tilde{X}$ is connected and contained in $\tilde{\mathbf{Y}}$, at the unit sites of this polymer (range less than $1$; they read $B_i$ on bonds with an end point in $\Delta_y$; the membership of $c_{3}$ in the polymer is the membership of its lift in the lift of the polymer).
\item[(a-4)] For $k\ge1$: the kernels of $Q_k(B_i)\cdot J^{*}$ and of $J\cdot Q_k^{T}(-B_i)$ (range less than $1$), and of $B_{k,y}(B_i)$, $y\ge0$ (range less than $L$), at the indices over $\tilde{\mathbf{Y}}$.
By the rows nos.~4--6 of the table of Proposition~25.1, the row of $Q_k(B_i)\cdot J^{*}$ at a unit site $u$ reads $B_i$ on bonds contained in $\Delta_u$;
the row of $J\cdot Q_k^{T}(-B_i)$ at a fermion site $x=(c,g)$ reads $B_i$ on bonds contained in $\Delta_{[c]}$;
and the row of $B_{k,y}(B_i)$ at a unit site $u$ vanishes outside the $L$-block $B(y')$ that contains $u$, and it reads $B_i$ on the fine bonds of the paths $\boldsymbol{\Gamma}^{\pi}(y',z)$, $z\in B(y')$, which are contained in $\mathrm{Bl}_{k+1}(y')$.
The sets $\Delta_u$, $\Delta_{[c]}$ and $\mathrm{Bl}_{k+1}(y')$ lie in the $M$-cube of the row index, so that every bond read in these rows is contained in $\tilde{\mathbf{Y}}$.
\item[(a-5)] For $k=0$: the kernel (I.1.4) of
\[
F_0(B_i)=D_{B_i}+\bar{m}\cdot\eta+b_{L}\cdot J\cdot P(B_i)\cdot J^{*}
\]
at the fermion sites $x$, $x'$ over $\tilde{\mathbf{Y}}$.
Its range is $L-1$ (the row no.~8 of the table of Proposition~25.1): it vanishes unless $x'\in\{x,x^{(0)},x^{(1)},x^{(2)},x^{(3)}\}$ or $c$, $c'$ lie in one $L$-block, and it reads $B_i$ on the bond $(c,c_{3})$ and inside the $L$-block of $c$.
\end{itemize}
\item[(b)] (Forms.) For an $M$-cube $\tilde{\square}$ of $\tilde{\mathbf{Y}}$,
\[
\mathrm{idn}\,\langle\bar{\psi},\psi\rangle_{\mathrm{pr}_1\tilde{\square}}=\langle\bar{\psi},\psi\rangle_{\mathrm{pr}_2\tilde{\square}} .
\]
For $k=0$, let $G^{(1)}[\mathrm{Zf}](\square)$ denote the display (I.5.5) with a complex bond function $\mathrm{Zf}$ in the place of $\mathscr{Z}(s')$.
If $\mathrm{Zf}_1$ on $T_1$ and $\mathrm{Zf}_2$ on $T_2$ have equal lifts on the bonds that meet $\tilde{\square}$, then
\[
\mathrm{idn}\,G^{(1)}[\mathrm{Zf}_1](\mathrm{pr}_1\tilde{\square})=G^{(1)}[\mathrm{Zf}_2](\mathrm{pr}_2\tilde{\square}) .
\]
\item[(c)] (Extraction numbers.) Let $\tilde{X}$ be non-empty and connected, with $d_M(\tilde{X})\le L$ and $\tilde{X}\subset\tilde{Y}$, and let $F_1$, $F_2$ be even elements of $\mathrm{Gr}_1(\tilde{X})$, $\mathrm{Gr}_2(\tilde{X})$ that contain $\psi(\xi)$ only for $\xi$ over $\tilde{X}$ and satisfy $\mathrm{idn}\,F_1=F_2$.
Then the numbers $(F_i)_{00}$ and the $4\times4$ matrices
\[
\int\mathrm{d}x\int\mathrm{d}x'\,(F_i)_{20}(x,x'),
\]
formed with all arguments over the lift $\tilde{X}$ ((I.AP) of Section~I.3.6, and Section~I.5.6), coincide for $i=1$ and $i=2$.
Here, for an element $F$ of the algebra, $F_{00}$ denotes the coefficient of the empty monomial of $F$, and $F_{20}$ denotes the antisymmetric kernel, with respect to the weight $\eta_k^3/48$, of the part of $F$ that has degree $2$ in $\psi$ and degree $0$ in each of the other families; for $F\in\mathrm{Gr}[\psi,\chi]$ these are the objects of Section~I.5.1.
\item[(d)] (Adjustment.) Let $\tilde{X}$ be as in (c), and let $\mathscr{T}_i\in(\tilde{R}_k+\mathfrak{X}_k)(T_i)$ have equal lifts on the bonds contained in $\tilde{X}$.
Then
\[
\mathrm{idn}\circ\boldsymbol{\Theta}_{\mathrm{pr}_1\tilde{X},\mathscr{T}_1}=\boldsymbol{\Theta}_{\mathrm{pr}_2\tilde{X},\mathscr{T}_2}\circ\mathrm{idn}
\]
on the subalgebra generated by the $\psi(\xi)$ with $\xi$ over $\tilde{X}$; here $\boldsymbol{\Theta}$ is the homomorphism (I.5.3), formed with the chain (I.SY6)$^{\mathrm{ch}}$.
\item[(e)] (Scaling.) Let $\tilde{U}$ be a non-empty union of $M$-cubes of $\mathbb{L}_{k+1}$ that is admissible at level $k+1$, that is, $\mathrm{diam}(\tilde{U})\le\ell_1/L-14d_U\cdot M$.
Then $L\tilde{U}$ is admissible at level $k$, and $N(L\tilde{U})\subset L\,N(\tilde{U})$.
The homomorphism
\[
S_L\colon\quad\psi(x)\mapsto L^{-1}\psi(x/L),\qquad\chi(x,x')\mapsto L^{-1-\alpha}\chi(x/L,x'/L)
\]
of Section~I.3.7 maps the part of $\mathrm{Gr}_i(L\tilde{U})$ generated by $\psi$ and $\chi$ into the algebra $\mathrm{Gr}_i(\tilde{U})$ of level $k+1$, and
\[
S_L\circ\mathrm{idn}[L\tilde{U}]=\mathrm{idn}[\tilde{U}]\circ S_L .
\]
\item[(f)] (Integrals.) For $F\in\mathrm{Gr}_1(\tilde{Y})$,
\[
\mathrm{idn}\int\mathrm{d}\mu_I(W)\,F=\int\mathrm{d}\mu_I(W)\,\mathrm{idn}\,F .
\]
By (I.B3), (I.B4), the integral of a monomial in the $W(u)$ is either $0$ or $\pm\prod\nu(\mathrm{mo})^{-1}$, the product extending over the modes $\mathrm{mo}$ whose pair $W(\mathrm{mo})$, $\bar{W}(\mathrm{mo})$ occurs in the monomial (Proposition~15\,(e)); here $\nu=1$ for $k\ge1$ and $\nu=1/48$ for $k=0$, the same numbers on both tori ((I.L0.5)).
In the coordinates $\mathsf{Z}(\Upsilon)$, the measure $\hat{\mu}$ is, on both tori, the product over the basis of copies of the standard normal law restricted to $[-p_{0,k},p_{0,k}]$ and normalised; and
\[
\delta\varepsilon^{0}_{k}=-d_V\cdot\log\left((2\pi)^{-\frac{1}{2}}\int_{-p_{0,k}}^{p_{0,k}}\mathrm{d}z\,\exp\left(-\frac{z^2}{2}\right)\right)
\]
is the same number on both tori.
\end{itemize}
\end{named}

\begin{named}{Remark}
Proposition~27.4 corresponds to the following parts of the proof of \cite[Proposition~19.12]{QEDI}, where they are steps of the argument and not a statement: the second half of Step~1 of the proof of part~(b); in Step~2, the first paragraph, on the operators that are given by one formula, and the comparison of their kernels on the two tori at the beginning of the second paragraph; point~3 of Step~3, on exact operators and forms, as a whole; and the sentences of Step~4 on $\hat{\mu}$ and on $\mathrm{d}\mu_I(W)$.
In both papers the coincidence of rows and kernels rests on the locality of the formulas: the clause (I.f4) of Section~I.1.9 agrees with the requirement of locality that is part of the notion of an operator given by one formula in \cite[Definition~5.5(2)]{QEDI}.
From the comparison of two tori in that definition, point~3 of Step~3 obtains the coincidence of rows also for connections with equal lifts that need not be of the form $(B,\mathrm{pj}^{*}B)$, and for polymers $\mathrm{pr}_2(\tilde{X})$, with $\tilde{X}\subset\tilde{\mathbf{Y}}$ connected, that need not be equal to $\mathrm{pj}^{-1}(\mathrm{pr}_1\tilde{X})$.
As in the list (a-4), (a-5), the operators $Q_k(B)$, $Q_k^{T}(-B)$, $P(B)$, $B_{k,y}(B)$ and, for $k=0$, $F_0(B)$ are given in Step~2 by one formula with a range, there the range $L$, together with the distances beyond which their entries vanish and the bonds on which they read $B$.
The differences are the following.
Here the coincidence is asserted for all $B_i\in\tilde{G}_k(T_i)$ with equal lifts on the bonds that meet $\tilde{\mathbf{Y}}$, there for the connections that occur in the proof; in particular, for the operators of the determinant, (a-4) and (a-5), Step~2 compares the two tori at the pair $(B,\mathrm{pj}^{*}B)$ only.
The index sets and the kernels are those of Section~I.1.9.
In (a-4) and (a-5) the cluster maps $J$, $J^{*}$ occur, as in $Q_k(B)\cdot J^{*}$, $J\cdot Q_k^{T}(-B)$ and $J\cdot P(B)\cdot J^{*}$, and a range is stated for each operator separately.
In the rows (I.5.4) and in the display (I.5.5) the hopping matrices $M_{i}(x)$ stand in the place of the matrices $P^{\boldsymbol{\sigma}}_{\nu}$ of the corresponding displays in \cite[Section~17.4.4]{QEDI}, and at $k=0$ the rows of $D_B$ are those of (I.O1).
The chain of the adjustment is the chain (I.SY6)$^{\mathrm{ch}}$, whose points are fermion sites that carry labels.
At level $0$ the modes in (f) have the weight $1/48$.
\end{named}

\subsubsection*{Proposition 27.5 (the objects of the step agree: the insertion into the displays)}

\begin{named}{Proposition 27.5}
Assume (TT) with $k<k_m$.
Let $m$ be a real number with $|m|\le w^{\mathrm{m}}_k$.
For $i=1,2$ let $\mathscr{E}^{(i)}\in K^{x}_{k}(\ell_i)$ satisfy $\left\|\mathscr{E}^{(i)}\right\|^{x}_{k}\le w^{\mathrm{E}}_k$, where $\mathscr{E}^{(i)}=0$ if $k=0$, and let $\mathscr{E}^{(1)}\sim\mathscr{E}^{(2)}$ in the sense of Proposition~27.2\,(D4).
Let $\tilde{Y}$ be an admissible set, and write $\mathrm{idn}:=\mathrm{idn}[\tilde{Y}]$.
Let $(A_1,\mathscr{C}_1)\in\tilde{R}_k^{-}(T_1)\times\mathfrak{X}_k(T_1)$, and put
\[
(A_2,\mathscr{C}_2):=\mathrm{pj}^{*}(A_1,\mathscr{C}_1),\qquad
\mathscr{T}_i:=A_i+\mathscr{C}_i\quad(i=1,2).
\]
On $T_i$ the step map is the one of Section~I.5.7, given by the displays (8.1)$^{\mathrm{ch}}$, (8.2), (8.3)$^{\mathrm{ch}}$, (8.4), (8.6)$^{\mathrm{ch}}$ of Definition~24.D together with (I.5.4) and (I.5.5), and formed with the system $\mathrm{Op}^{(i)}$.
Put
\[
m'^{(i)}:=m+m\left(\mathscr{E}^{(i)}\right),\qquad
\mathscr{E}^{r,(i)}:=R\,\mathscr{E}^{(i)}.
\]

For a non-empty connected union $\tilde{\mathbf{Y}}$ of $M$-cubes contained in $\tilde{Y}$ let $\mathbf{Y}_i:=\mathrm{pr}_i(\tilde{\mathbf{Y}})$, and denote by $\mathrm{Bs}$ the set of those basis elements of the covering lattice that are contained in $\tilde{\mathbf{Y}}$.
With every
\[
\mathrm{zval}=\left(\mathrm{zval}_{\Upsilon}\right)_{\Upsilon\in\mathrm{Bs}}\in[-p_{0,k},p_{0,k}]^{\mathrm{Bs}}
\]
associate the element $\mathsf{Z}^{(i)}\in\mathrm{Box}_k(T_i)$ whose coordinate at $\mathrm{pr}_i(\Upsilon)$ is $\mathrm{zval}_{\Upsilon}$ for every $\Upsilon\in\mathrm{Bs}$, and whose coordinates at all other basis elements of $T_i$ are $0$.
Let $\mathrm{sval}=(\mathrm{sval}_{\tilde{\mathrm{pc}}})$ be a family of complex numbers of modulus at most $M^{\alpha_0}$, indexed by the parameter cubes $\tilde{\mathrm{pc}}$ of the three kinds that are contained in $\tilde{\mathbf{Y}}$.
With it associate the parameter family $s^{(i)}=\left(s'^{(i)},s''^{(i)}\right)\in\mathrm{Pd}_k(T_i)$ that has the value $\mathrm{sval}_{\tilde{\mathrm{pc}}}$ at $\mathrm{pr}_i(\tilde{\mathrm{pc}})$ and the value $0$ at every other parameter cube of $T_i$.

Under these assumptions the following assertions hold.
\begin{itemize}
\item[(a)] (Extraction, adjustment.) One has $m'^{(1)}=m'^{(2)}$.
Moreover
\[
R\,\mathscr{E}^{(1)}\sim R\,\mathscr{E}^{(2)}\qquad\text{and}\qquad
\left(R\,\mathscr{E}^{(1)}\right)^{\mathrm{nat}}\sim\left(R\,\mathscr{E}^{(2)}\right)^{\mathrm{nat}},
\]
as local families on $\tilde{R}_k\times\mathfrak{X}_k$.
\item[(b)] (Determinant.) Let $\tilde{X}$ be a non-empty connected union of $M$-cubes contained in $\tilde{Y}$, and let $B_i\in\tilde{G}_k(T_i)$, $i=1,2$, have equal lifts on the bonds that meet $\tilde{X}$.
Then, for every $y\ge0$, the integrand at $y$ of $\acute{E}^{(1)}_k\left(\mathrm{pr}_1\tilde{X},B_1\right)$ coincides with the integrand at $y$ of $\acute{E}^{(2)}_k\left(\mathrm{pr}_2\tilde{X},B_2\right)$ (Section~I.5.7).
Here, for a polymer $X$ and a connection $B$ of one of the two tori, the traces are the following sums and integrals.
For $k\ge1$,
\begin{multline*}
\mathrm{Tr}\left[\gamma_5\cdot B_{k,y}(B)\cdot Q_k(B)\cdot J^{*}\cdot S_{k,y}(X,B)\cdot J\cdot Q_k^{T}(-B)\cdot B_{k,y}(B)\right]\\
=\sum_{z_1,z_2}\mathrm{tr}\left[\left(B_{k,y}(B)\cdot\gamma_5\cdot B_{k,y}(B)\right)(z_2,z_1)\cdot\left(Q_k(B)\cdot J^{*}\cdot S_{k,y}(X,B)\cdot J\cdot Q_k^{T}(-B)\right)(z_1,z_2)\right],
\end{multline*}
the sum extending over the unit sites $z_1$, $z_2$ of $X$.
For $k=0$,
\begin{gather*}
\mathrm{Tr}\left[F_0(B)\cdot S_{0,y}(X,B)\right]=\int_X\mathrm{d}x\int_X\mathrm{d}x'\,\mathrm{tr}\left[F_0(B)(x,x')\cdot S_{0,y}(X,B)(x',x)\right],\\
\mathrm{Tr}\left[\gamma_5\cdot S_{0,y}(X,B)\right]=\int_X\mathrm{d}x\,\mathrm{tr}\left[\gamma_5\cdot S_{0,y}(X,B)(x,x)\right].
\end{gather*}
Consequently
\begin{gather*}
\acute{E}^{(1)}_k\left(\mathrm{pr}_1\tilde{X},B_1\right)=\acute{E}^{(2)}_k\left(\mathrm{pr}_2\tilde{X},B_2\right),\qquad
E^{d,(1)}_k\left(\mathrm{pr}_1\tilde{X},B_1\right)=E^{d,(2)}_k\left(\mathrm{pr}_2\tilde{X},B_2\right),\\
{\delta E^{\mathrm{det}}}^{(1)}_k\left(\mathrm{pr}_1\tilde{X};A_1,\mathscr{C}_1\right)={\delta E^{\mathrm{det}}}^{(2)}_k\left(\mathrm{pr}_2\tilde{X};A_2,\mathscr{C}_2\right),\\
E_k^{\mathrm{det},(1)}\left(\mathrm{pr}_1\tilde{X},\mathscr{C}_1\right)=E_k^{\mathrm{det},(2)}\left(\mathrm{pr}_2\tilde{X},\mathscr{C}_2\right).
\end{gather*}
\item[(c)] Let $\tilde{\mathbf{Y}}$ be a non-empty connected union of $M$-cubes contained in $\tilde{Y}$, and let $\mathrm{zval}$ and $\mathrm{sval}$ be as above.
In (c3)--(c7) the letter $\tilde{X}$ stands for a non-empty connected union of $M$-cubes contained in $\tilde{\mathbf{Y}}$.
\begin{itemize}
\item[(c1)] (The field $\mathscr{Z}$.) For every fine bond $b$ contained in $\tilde{\mathbf{Y}}^{\sim(1+d_U)}$,
\[
\mathscr{Z}^{(1)}\left(s'^{(1)}\right)\left(\mathrm{pr}_1 b\right)=\mathscr{Z}^{(2)}\left(s'^{(2)}\right)\left(\mathrm{pr}_2 b\right),
\]
where on each torus the field $\mathscr{Z}(s')$ is formed with $\mathsf{Z}=\mathsf{Z}^{(i)}$:
\[
\mathscr{Z}(s')(b)=\sum_{b'}\sum_{\Upsilon,\Upsilon'}H_{k}(s')(b,b')\cdot C(b',\Upsilon)\cdot C_{k}^{1/2}(s')(\Upsilon,\Upsilon')\cdot\mathsf{Z}(\Upsilon').
\]
Written on the covering lattice, the terms of this sum that do not vanish have $b'$ contained in $\tilde{\mathbf{Y}}^{\sim(1+2d_U)}$, $\Upsilon$ contained in the $M$-cubes of the end points of $b'$, and $\Upsilon'$ contained in $\tilde{\mathbf{Y}}^{\sim(1+3d_U)}$, a subset of $N(\tilde{\mathbf{Y}})$.
\item[(c2)] (Connections.) One has $1<\mathrm{rad}_2\le\mathrm{rad}_1$ and $\mathrm{rad}_4>1$.
For $|t|\le\mathrm{rad}_1$ put
\[
B_i(t):=\mathscr{T}_i+t\cdot\mathscr{Z}^{(i)}\left(s'^{(i)}\right).
\]
Then $B_i(t)\in\tilde{G}_k(T_i)$ and $A_i+t\cdot\mathscr{Z}^{(i)}\left(s'^{(i)}\right)\in\tilde{R}_k(T_i)$.
The fields $B_1(t)$ and $B_2(t)$ have equal lifts on the bonds of $\tilde{\mathbf{Y}}^{\sim(1+d_U)}$, and so have $A_1+t\cdot\mathscr{Z}^{(1)}\left(s'^{(1)}\right)$ and $A_2+t\cdot\mathscr{Z}^{(2)}\left(s'^{(2)}\right)$.
\item[(c3)] (Rows of the fermionic primitives.) Consider the connections $B_i(t)$, $|t|\le\mathrm{rad}_1$, and $\mathscr{T}_i$.
Written with all indices in $N(\tilde{\mathbf{Y}})$, as in Proposition~27.4\,(a), the following rows and kernels coincide for $i=1$ and $i=2$:
the rows of $H^{(i)}_{k}\left(s''^{(i)},\,\cdot\,\right)$ and of $\bar{H}^{(i)}_{k}\left(s''^{(i)},\,\cdot\,\right)$ at these connections, at the fermion sites over the photon sites whose sup-distance from $\tilde{\mathbf{Y}}$ is at most $\eta_k$;
the rows of $\Gamma^{(i)}_k\left(s''^{(i)},\mathscr{T}_i\right)$ at the unit sites of $\tilde{\mathbf{Y}}$ (for $k=0$, at the fermion sites over $\tilde{\mathbf{Y}}$);
and, for $k\ge1$, the kernels of $M^{(i)}_k\left(\mathrm{pr}_i\tilde{X},B_i(t)\right)$.
These rows vanish at the columns that have no lift in $\tilde{\mathbf{Y}}$.
\item[(c4)] (Data.) Write $\mathscr{Z}^{(i)}:=\mathscr{Z}^{(i)}\left(s'^{(i)}\right)$.
For $|t|\le\mathrm{rad}_1$,
\[
\mathrm{idn}\,\mathscr{E}^{r,(1)}\left(\mathrm{pr}_1\tilde{X};A_1+t\cdot\mathscr{Z}^{(1)},\mathscr{C}_1\right)=\mathscr{E}^{r,(2)}\left(\mathrm{pr}_2\tilde{X};A_2+t\cdot\mathscr{Z}^{(2)},\mathscr{C}_2\right).
\]
\item[(c5)] (First localisation.) Let $j\in\{1,2,3\}$; for $j=3$, and for $j=1$ at $k=0$, let $\tilde{X}$ be a single $M$-cube.
The element $G^{(j),(i)}\left(s^{(i)},\mathrm{pr}_i\tilde{X}\right)$ belongs to $\mathrm{Gr}_i(\tilde{\mathbf{Y}})$: it is a polynomial in the generators $\psi$, $\chi$, $\Psi$ with indices in $N(\tilde{\mathbf{Y}})$.
For every $t$ on the contour, $\mathrm{idn}$ maps the integrand on $T_1$ to the integrand on $T_2$, and therefore
\[
\mathrm{idn}\,G^{(j),(1)}\left(s^{(1)},\mathrm{pr}_1\tilde{X}\right)=G^{(j),(2)}\left(s^{(2)},\mathrm{pr}_2\tilde{X}\right).
\]
Here, for $k\ge1$ and on each torus,
\begin{gather*}
G^{(1)}(s,X)=\oint_{\mathrm{rad}_1}\left\langle\bar{\Psi},M_k\left(X,\mathscr{T}+t\cdot\mathscr{Z}(s')\right)\Psi\right\rangle,\\
G^{(2)}(s,X)=\oint_{\mathrm{rad}_2}\mathscr{E}^{r}\left(X;A+t\cdot\mathscr{Z}(s'),\mathscr{C},\psi^{\#}+t\cdot J_k^{\#}\left(s'',\mathscr{T},\mathscr{Z}(s')\right)\Psi\right),\\
G^{(3)}(s,\square)=\oint_{\mathrm{rad}_2}(-m')\cdot\left\langle\bar{\psi}+t\cdot\bar{J}_k\left(s'',\mathscr{T},\mathscr{Z}(s')\right)\bar{\Psi},\ \psi+t\cdot J_k\left(s'',\mathscr{T},\mathscr{Z}(s')\right)\Psi\right\rangle_{\square},
\end{gather*}
and $J_k^{\#}\left(s'',\mathscr{T},\mathscr{Z}(s')\right)$ is formed with the two connections $B(1)$ and $B(0)$, where $B(t):=\mathscr{T}+t\cdot\mathscr{Z}(s')$.
For $k=0$, $G^{(1)}(s,\square)$ is the display (I.5.5), and $G^{(2)}=G^{(3)}=0$.
\item[(c6)] The identity of (c5) holds for every $\mathrm{sval}$ in the polydisk.
By Proposition~27.2\,(e), the pairs $(X,P)$ of $T_1$ with $Y'(X,P)=\mathbf{Y}_1$ correspond to the pairs $(X,P)$ of $T_2$ with $Y'(X,P)=\mathbf{Y}_2$.
It follows that
\[
\mathrm{idn}\,\mathrm{Loc}\left[G^{(j)}\right]^{(1)}(\mathbf{Y}_1)=\mathrm{Loc}\left[G^{(j)}\right]^{(2)}(\mathbf{Y}_2),
\]
and that, after the substitution (8.3a)$^{\mathrm{ch}}$ (for $k\ge1$; Proposition~27.4\,(a-3)),
\[
\mathrm{idn}\,E^{(j),x,(1)}(\mathbf{Y}_1)=E^{(j),x,(2)}(\mathbf{Y}_2)\qquad(j=1,2,3)
\]
at $\left(A_i,\mathscr{C}_i,\mathsf{Z}^{(i)}\right)$.
\item[(c7)] (Second localisation.) One has
\[
\mathrm{idn}\left(G^{(\le3)}+G^{(4)}+G^{(5)}\right)^{(1)}\left(s''^{(1)},\mathrm{pr}_1\tilde{X}\right)=\left(G^{(\le3)}+G^{(4)}+G^{(5)}\right)^{(2)}\left(s''^{(2)},\mathrm{pr}_2\tilde{X}\right),
\]
where the field
\[
W_k^{\#}(s'',\mathscr{T})=H_k^{\#}(s'',\mathscr{T})\cdot\tilde{\Gamma}_k(s'',\mathscr{T})\cdot W
\]
enters at the fermion sites over $\tilde{X}$ and at the pairs that are admissible for $\tilde{X}$.
Hence, for every $\mathrm{zval}$,
\[
\mathrm{idn}\,E^{(\le5),x,(1)}(\mathbf{Y}_1)=E^{(\le5),x,(2)}(\mathbf{Y}_2)
\]
at $\left(A_i,\mathscr{C}_i,\mathsf{Z}^{(i)}\right)$.
\end{itemize}
\item[(d)] (Cluster expansion.) Let $\tilde{X}$ be a non-empty connected union of $M$-cubes contained in $\tilde{Y}$, and let $\mathsf{Z}^{(i)}$ be formed with $\tilde{\mathbf{Y}}=\tilde{X}$.
Then
\begin{gather*}
\mathrm{idn}\,\mathrm{Km}^{(1)}\left(\mathrm{pr}_1\tilde{X}\right)=\mathrm{Km}^{(2)}\left(\mathrm{pr}_2\tilde{X}\right)\quad\text{at }\mathsf{Z}^{(i)},\qquad
\mathrm{idn}\,\mathrm{Hm}^{(1)}\left(\mathrm{pr}_1\tilde{X}\right)=\mathrm{Hm}^{(2)}\left(\mathrm{pr}_2\tilde{X}\right),\\
\mathrm{idn}\,E^{\#,x,(1)}_k\left(\mathrm{pr}_1\tilde{X};A_1,\mathscr{C}_1\right)=E^{\#,x,(2)}_k\left(\mathrm{pr}_2\tilde{X};A_2,\mathscr{C}_2\right),
\end{gather*}
and the series (8.4) converge in norm on both tori.
\end{itemize}
\end{named}

\emph{Conditions on the constants.}
Of the conditions on the constants listed in Proposition~14.2\,(b), the following enter.
The third inequality of (A1) gives $\mathrm{rad}_2\le\mathrm{rad}_1$, and (A5) gives $\mathrm{rad}_i\ge2$.
Of (A8) the first and the fourth inequality enter; at $v=-\log e_k$ they give
\[
e_k^{\varepsilon}\le\frac{1}{4},\qquad e_k^{\frac{1}{4}}+e_k^{\varepsilon}\le1 .
\]
The first of these gives $e_k<1$ and $0<1-e_k^{\varepsilon}<1$.
Of (S1) one needs $4\varepsilon\le\frac{1}{2}$ and the order of the three exponents,
\[
-\frac{3}{4}+7\varepsilon\le-\frac{1}{4}+3\varepsilon\le-\frac{1}{4}+6\varepsilon .
\]
Of (M4) one needs $M^{\alpha_0}>1$, so that $[0,1]^P$ lies in the interior of the polydisk.
Of (L5) and (M1) one needs that $L$ is odd, $L>1$, and that $M$ is a power of $L$ with $M>1$; these give the fact (g0) of \cite[Section~17.4.1]{QEDI}.
The inequality $L>1$ of (L5) gives moreover $b_{k}\ge b_{-}=b\cdot\left(1-L^{-1}\right)>0$, where $b=1$ by (C0).
Of (L5) one needs, besides $L$ odd and $L>1$, the inequality $L\ge5$ itself.
The inequality $L\ge5$ and the fourth inequality of (A8), together with $e_k^{\varepsilon}\le\frac{1}{4}$, enter through \cite[Lemma~5.17(a)]{QEDI}, that is, \cite[Lemma~17.4(a)]{QEDI}: the sets $\tilde{R}_k$ and $\tilde{R}_k+\mathfrak{X}_k$ are contained in $\tilde{G}_k$.
In the proof of that lemma the cube $\square^{\sim5}$ lies in $\square^{\sim L}$ for every $M$-cube $\square$, because $L\ge5$; one has $e_k<1$; and the fourth inequality of (A8) at $v=-\log e_k$, multiplied by $e_k^{-1+\varepsilon}$, gives
\[
e_k^{-\frac{3}{4}+\varepsilon}+e_k^{-1+2\varepsilon}\le e_k^{-1+\varepsilon} .
\]
These inclusions enter the data: by the fact (TT.3) of Chapter~27b, Section~0.1, the bond functions $\mathscr{T}_i$, $\mathscr{C}_i$ and $0$ lie in $\tilde{G}_k(T_i)$.
They enter also (c2), which asserts $B_i(t)\in\tilde{G}_k(T_i)$.
Finally, one needs $\gamma_{\mathrm{w}}>0$ and $\kappa_{\mathrm{w}}>0$; these hold by the rows $\gamma_{\mathrm{w}}$ and $\kappa_{\mathrm{w}}$ of Proposition~14.2\,(a), the second together with (C2) and (C0).
These two enter (b), through the existence of the integrals over $y$ in $\acute{E}_k$: the factor $\exp\left(-\gamma_{\mathrm{w}}\,d-\kappa_{\mathrm{w}}\,d_M(X)\right)$ of W7$^{\mathrm{ch}}$ (Section~G4.0.4), with $d$ the distance of the two indices, is at most $1$.

\begin{named}{Remark}
Proposition~27.5 corresponds to clause (b) of \cite[Proposition~19.12]{QEDI}, as it is obtained in Steps~1--4 of the proof given there, Step~3 with its six parts.
The following points differ.
The identities by which $\mathrm{idn}$ carries the integrand of $G^{(j),(1)}$ to the integrand of $G^{(j),(2)}$ are stated as clauses, (c5) and (c7).
They come with the form of each $G^{(j)}$ in the generators over the lift (the assertion of (c5) that $G^{(j),(i)}$ belongs to $\mathrm{Gr}_i(\tilde{\mathbf{Y}})$), with the list of the rows, parameter families and connections that enter ((c1)--(c3)), and with the inequalities $1<\mathrm{rad}_2\le\mathrm{rad}_1$ ((c2)).
The operators of the substitution (8.3a)$^{\mathrm{ch}}$ and the vertex of level $0$ are given by the displays (I.5.4) and (I.5.5).
In (b) the cluster maps $J$, $J^{*}$ and the matrix $\gamma_5$ occur, and the traces are taken by (I.ke2); at $k=0$ the sums over sites are the integrals $\int_X\mathrm{d}x$.
The rows of the family (p3) are taken at the fermion sites over the photon sites at distance at most $\eta_k$ from $\tilde{\mathbf{Y}}$, a set that contains the site $x^{(3)}$ of every site $x$ over $\tilde{\mathbf{Y}}$; and the rows of $\Gamma_0$ are taken at fermion sites.
\end{named}

\subsubsection*{Proposition 27.6 (the steps on two tori agree; extracted mass and energy coincide)}

\begin{named}{Proposition 27.6}
Assume (TT) with $0\le k<k_m$.
Let $m$ be a real number with $|m|\le w^{\mathrm{m}}_k$.
For $i=1,2$ let $\mathscr{E}^{(i)}\in K^{x}_{k}(\ell_i)$ with $\left\|\mathscr{E}^{(i)}\right\|^{x}_{k}\le w^{\mathrm{E}}_k$, where $\mathscr{E}^{(i)}=0$ if $k=0$, and $\mathscr{E}^{(1)}\sim\mathscr{E}^{(2)}$; and let $\mathfrak{P}^{(i)}_k$ be families of complex functions on $\mathfrak{X}_k(T_i)$ with $\mathfrak{P}^{(1)}_k\sim\mathfrak{P}^{(2)}_k$.
Then the following hold.
\begin{itemize}
\item[(a)] One has $\varepsilon\left(\mathscr{E}^{(1)}\right)=\varepsilon\left(\mathscr{E}^{(2)}\right)$ and $m\left(\mathscr{E}^{(1)}\right)=m\left(\mathscr{E}^{(2)}\right)$, and $\delta\varepsilon^{0}_{k}$ is the same number on both tori.
\item[(b)] For $i=1,2$ put
\[
\mathscr{E}^{(i)+}:=L R\,\mathscr{E}^{(i)}+L\,{\delta E^{\mathrm{det}}}^{(i)}_k+E^{*,x,(i)}_k\left(m,\mathscr{E}^{(i)}\right).
\]
Then $\mathscr{E}^{(1)+}\sim\mathscr{E}^{(2)+}$ at the level $k+1$, that is, for the tori $T_{k+1}(R)$ and $T_{k+1}(R_{\mathrm{top}})$ and as local families on $\tilde{R}_{k+1}\times\mathfrak{X}_{k+1}$.
\item[(c)] For $i=1,2$, for $Y\in D_{k+1}$ and for $\mathscr{C}'\in\mathfrak{X}_{k+1}$ of the corresponding torus of level $k+1$ put
\[
\mathfrak{P}^{(i)}_{k+1}(Y,\mathscr{C}'):=\sum_{X:\ \bar{X}^{L}=LY}\left(\mathfrak{P}^{(i)}_k\left(X,\mathscr{C}'_L\right)+E_k^{\mathrm{det},(i)}\left(X,\mathscr{C}'_L\right)\right).
\]
Then $\mathfrak{P}^{(1)}_{k+1}\sim\mathfrak{P}^{(2)}_{k+1}$.
\item[(d)] Let $E_{(i)}:=\mathscr{E}^{(i)}|_0$.
Then
\[
L R\,E_{(1)}+L\,E_k^{\mathrm{det},(1)}+E^{*,(1)}_k\left(m,E_{(1)}\right)\sim L R\,E_{(2)}+L\,E_k^{\mathrm{det},(2)}+E^{*,(2)}_k\left(m,E_{(2)}\right).
\]
\item[(e)] (Without slot.) The assertions (a) and (b) hold for $E_{(i)}\in K_k(\ell_i)$ with $\|E_{(i)}\|_k\le w^{\mathrm{E}}_k$, $E_{(i)}=0$ if $k=0$, and $E_{(1)}\sim E_{(2)}$, with $E_k^{\mathrm{det}}$ in the place of ${\delta E^{\mathrm{det}}}_k$ and $E^{*}_k$ in the place of $E^{*,x}_k$.
\end{itemize}
At the level $0$ the data are $(m,0)$.
Then (a) reads $\varepsilon\left(\mathscr{E}^{(i)}\right)=0$ and $m\left(\mathscr{E}^{(i)}\right)=0$ on both tori; in (b) only the terms $L\,{\delta E^{\mathrm{det}}}_0$ and $E^{*,x}_0(m,0)$ occur; and the generators $W$ carry the label and the weight $\frac{1}{48}$.
\end{named}

\begin{named}{Remark}
Clauses (a)--(d) of Proposition~27.6 correspond to \cite[Proposition~19.12]{QEDI}, formed here with the index sets of Part~I.
Clauses (a)--(c) are the clauses (a)--(c) of \cite[Proposition~19.12]{QEDI}, which are also \cite[Theorem~5.80(a)--(c)]{QEDI}; clause (d) is the assertion that closes the statement of \cite[Proposition~19.12]{QEDI}, the restriction to $\mathscr{C}=0$.
Clause (e), the form of (a) and (b) without slot, has no counterpart in \cite{QEDI}; on the forms without slot see the Remark after Proposition~48.2.
\end{named}

\subsection*{28. The flow; the tuned pair}

\subsubsection*{Proposition 28.1 (the flow on one torus)}

\begin{named}{Proposition 28.1}
Let (FL) hold, let $\left(m_{\mathrm{in}},\varepsilon^{\mathrm{v}}_{\mathrm{in}}\right)$ be a real pair, and let $0\le k\le k_{*}(R)$.
Assume that the numbers $m_i$, $i<k$, which the recursion of Section~G4.0.5 defines one after another, satisfy $|m_i|\le w^{\mathrm{m}}_i$.
Then the trajectory on $R$ exists up to the level $k$, so that the data of the level $k$ are defined, and the following hold.
The numbers $\varepsilon^{\mathrm{v},t}_k$ and $m_k$ are real.
One has $\mathscr{E}_k\in K^{x}_{k}(\ell_k)$ and $E_k\in K_k(\ell_k)$, and $\mathscr{E}_k|_0=E_k$.
For every $X\in D_k$, $\mathfrak{P}_k(X,\cdot)$ is a complex function on $\mathfrak{X}_k$ with $\mathfrak{P}_k(X,0)=0$.
Finally,
\[
\|E_k\|_k\le\|\mathscr{E}_k\|^{x}_{k}\le4\,e_k^{\varepsilon}\cdot w^{\mathrm{E}}_k\le w^{\mathrm{E}}_k .
\]
Consequently a real pair is controlled on $R$ if and only if, recursively, $|m_k|\le w^{\mathrm{m}}_k$ for $k=0,1,\dots,k_{*}(R)$.
At the level $0$ one has $E_0=0$, $\mathscr{E}_0=0$, $\mathfrak{P}_0=0$, the mass coupling is $m_{\mathrm{in}}$, and $\varepsilon^{\mathrm{v},t}_0=\varepsilon^{\mathrm{v}}_{\mathrm{in}}$ (see (I.L0.21)).
\end{named}

\emph{Constants.}
The constants that occur are $c_{\mathrm{fl}}$ of (C0), $L$ of (C3) and $\varepsilon$ of (C7).
The bounds involve $N$ only through the datum $e_k=e\cdot L^{\frac{k-N}{2}}$, and they involve neither $R$ nor the twist $\theta$ nor the pair (compare Proposition~24.1\,(i)).

\emph{Conditions on the constants.}
The following enter: $c_{\mathrm{fl}}\cdot L^{-\frac{3}{8}+6\bar{\varepsilon}}\le\frac{1}{16}$, from (C3); $e_0^{\varepsilon}\le\frac{1}{4}$, from (A8); (S1), of which $0<\varepsilon\le\bar{\varepsilon}$ is needed; and (L5), of which $L>1$ is needed.
Further, the conditions (A8), $L\ge5$, $L^{-\frac{3}{4}}\le\frac{1}{2}$ (a consequence of (L1), by Proposition~14.2\,(f)) and $\varepsilon\le\frac{1}{200}$ are hypotheses of \cite[Lemma~5.17]{QEDI} and \cite[Lemma~5.18]{QEDI}, on which the last sentence of Definition~24.D\,(D6) rests.
Finally, the row (e1) at $v=v_K$ supplies the data of (Sc1).

\begin{named}{Remark}
Proposition~28.1 corresponds to \cite[Proposition~19.13]{QEDI}, together with the inequalities at the head of \cite[Section~19.5]{QEDI}, and to \cite[Theorem~5.81(a)]{QEDI}.
The spaces $K_k$ and $K^{x}_{k}$ are those of Section~I.5.4, formed on $T^{\mathrm{f}}_k$; the numbers in the statement are those of \cite[Proposition~19.13]{QEDI}.
\end{named}

\subsubsection*{Proposition 28.2 (exactly one pair stays in the windows and ends at zero)}

\begin{named}{Proposition 28.2}
Let (FL) hold, and write $k_{*}:=k_{*}(R)$.
There is exactly one real pair $\left(m_{\mathrm{in}},\varepsilon^{\mathrm{v}}_{\mathrm{in}}\right)$ that is controlled on $R$ and satisfies
\[
m_{k_{*}}=0,\qquad\varepsilon^{\mathrm{v},t}_{k_{*}}=0 .
\]
For $0\le k\le k_{*}$ its trajectory on $R$ satisfies
\[
|m_k|\le\frac{w^{\mathrm{m}}_k}{8},\qquad
\varepsilon^{\mathrm{v},t}_k=-\sum_{i=k}^{k_{*}-1}L^{-3(i-k)}\cdot\left(\varepsilon(E_i)+\delta\varepsilon^{0}_{i}\right),\qquad
\left|\varepsilon^{\mathrm{v},t}_k\right|\le5\,c_{\mathrm{fl}}\cdot e_k^{\varepsilon}\cdot w^{\mathrm{E}}_k .
\]
\end{named}

\emph{Constants.}
The constants that occur are $c_{\mathrm{fl}}$ of (C0), $L$ of (C3) and $\varepsilon$ of (C7).

\emph{Conditions on the constants.}
The following enter: from (C3), the inequality $L^{-\frac{5}{8}}\le\frac{1}{4}$, which contains $L\ge4$, and the inequality $c_{\mathrm{fl}}\cdot L^{-\frac{3}{8}+6\bar{\varepsilon}}\le\frac{1}{16}$; from (C0), the inequality $c_{\mathrm{fl}}\ge1$; from (A8), the inequality $e_0^{\varepsilon}\le\frac{1}{4}$; and, of (S1), the range $0<\varepsilon\le\bar{\varepsilon}$.
Further, (fl) enters in the form $\mathrm{sm}_k=(c_{\mathrm{fl}}+2)\cdot e_k^{\varepsilon}\le\frac{1}{16}$ for each $k\le k_{*}$.
Under (FL) this holds by (Sc6), that is, by the rows (e1) at $v=v_K$ and (e2) at $v=v_{*}$; the rows (e1)--(e8) hold for all $v\ge-\log e_1$ by the second assertion of Proposition~14.3\,(f).
The row (e1) at $v=v_K$ also supplies the data of (Sc1).

\begin{named}{Remark}
Proposition~28.2 corresponds to \cite[Proposition~19.14]{QEDI} and to \cite[Theorem~5.81(b)]{QEDI}.
The fixed point is taken in the spaces $K_k$ without slot, with the extraction and the contraction in their forms without slot (Propositions~48.2 and 48.5; see the Remark after Proposition~48.2), and $m_{\mathrm{in}}$ is the term of index $k=0$ of the fixed point.
Exactly one real mass, the same on the $48$ sites over a photon site (see (I.E1)), and one vacuum energy are tuned; no other coupling is tuned.
\end{named}

\subsubsection*{Proposition 28.3 (one pair for all tori of the scheme)}

\begin{named}{Definition}[The tuned pair]
Let $\bar{m}\in(0,1]$, let $e>0$ satisfy $e_K=e\cdot h^{\frac{1}{2}}\le e_1$, and let $N\ge N_{\min}\left(R_{\mathrm{top}}(e),e\right)$.
The pair
\[
P_N(e)=\left(m_{0,N}(e),\varepsilon^{\mathrm{v}}_{0,N}(e)\right)
\]
is the pair of Proposition~28.2 on the torus $R_{\mathrm{top}}(e)$.
Its trajectory on $R_{\mathrm{top}}$ is denoted by $\left(\varepsilon^{\mathrm{v},a}_k,m^a_k,E^a_k\right)$, and its trajectory with slot by $\left(\mathscr{E}^a_k,\mathfrak{P}^a_k\right)$, where $0\le k\le k_{\mathrm{top}}$.
For $0\le N<N_{\min}\left(R_{\mathrm{top}}(e),e\right)$ one puts $m_{0,N}(e):=0$.
For every integer $N\ge0$ let
\[
\boldsymbol{\mu}_N:=\bar{m}\cdot\eta+m_{0,N}(e).
\]
\end{named}

\begin{named}{Proposition 28.3}
Let (FL) hold.
\begin{itemize}
\item[(a)] The pair $P_N(e)$ is defined and controlled on $R$; for $0\le k\le k_*(R)$ its trajectory on $R$ satisfies
\begin{equation}
m_k=m^a_k,\qquad
\varepsilon^{\mathrm{v},t}_k=\varepsilon^{\mathrm{v},a}_k,\qquad
\mathscr{E}_k\sim\mathscr{E}^a_k,\qquad
\mathfrak{P}_k\sim\mathfrak{P}^a_k,\qquad
E_k\sim E^a_k;
\tag{6.1}
\end{equation}
here $\sim$ is the relation of Proposition~27.2\,(D4) at the level $k$, and for $R=R_{\mathrm{top}}$ it is read as equality.
\item[(b)] For $k\le k_*(R)$ one has
\[
|m_k|\le\frac{w^{\mathrm{m}}_k}{8},\qquad
\|E_k\|_k\le\|\mathscr{E}_k\|^x_k\le4\,e_k^{\varepsilon}\cdot w^{\mathrm{E}}_k,
\]
and
\[
\left|\varepsilon^{\mathrm{v},t}_{k_*(R)}\right|\le5\,c_{\mathrm{fl}}\cdot e_*^{\varepsilon}\cdot w^{\mathrm{E}}_*,
\]
where $e_*$ and $w^{\mathrm{E}}_*$ are the numbers of the torus $R$.
\item[(c)] The pair $P_N(e)$ is real.
It is a function of the data $(N,e,\bar{m})$ of the Definition and of the scheme (Section~I.6.2: the constants of the lines (C0)--(C9) of Section~I.8, among them $\rho$ and the weights; the basis $\{e^{0}_{i}\}$; the systems of Proposition~26.5 on $T_k(R_{\mathrm{top}})$, $k<k_{\mathrm{top}}$).
No torus is among the arguments of $P_N(e)$.
In particular the torus $R$ of (FL) is not one of them: for given $(N,e,\bar{m})$ the pair $P_N(e)$ is the same for all $R\in\mathrm{Tor}(e)$ with $N\ge N_{\min}(R,e)$, the torus $R=R_{\mathrm{top}}(e)$ included.
The only torus that occurs in the definition of $P_N(e)$ is $R_{\mathrm{top}}(e)$, which is itself a function of $(e,\bar{m})$ and of the constants $L$, $M$, $r$.
Neither a twist, nor a probe, nor a number of the line (C10) formed with a probe enters the definition of $P_N(e)$.
\item[(d)] Assume only that $\bar{m}\in(0,1]$ and that $e>0$ satisfies $e_K\le e_1$; no torus and no integer $N$ as in (FL) are required.
Then for every integer $N\ge0$,
\[
\left|m_{0,N}(e)\right|\le\frac{w^{\mathrm{m}}_0}{8}<\frac{1}{8},
\qquad\text{where}\quad
w^{\mathrm{m}}_0=q^{\frac{3}{4}-8\varepsilon},\quad q=e\cdot L^{-\frac{N}{2}}.
\]
Thus $N\mapsto m_{0,N}(e)$ is a real sequence independent of any torus, that is, a mass assignment in the sense of Section~I.7.5.
\end{itemize}
The sign of $m_{0,N}(e)$, and with it the sign of $\boldsymbol{\mu}_N$, is left undetermined; every real lattice mass is admitted (I.K10).
\end{named}

\emph{Conditions on the constants.}
Of the conditions listed in Proposition~14.2, the following enter.
For the Definition one needs (M1), of which it is used that $M$ is a power of $L$, and the row (e1) at $v=v_K$, in the form
\[
v_K\ge8,\qquad
\left(1+\frac{\mathrm{lg}}{v_K}\right)^r\le\left(1+\frac{3\,\mathrm{lg}}{v_K}\right)^r<L
\]
(the rows (e1)--(e8) hold for all $v\ge-\log e_1$ by the second assertion of Proposition~14.3\,(f)).
Both are contained in (Sc0), and the second one supplies the data of (Sc1) and (Sc2).
For (d) one needs $e_1\le\frac{1}{3}$, which is part of (C9), and (S1), of which it is used that $\frac{3}{4}-8\varepsilon>0$.
Finally (L5) enters, through $L>1$.

\begin{named}{Remark}
The Definition and Proposition~28.3 correspond to \cite[Definition~5.82]{QEDI} together with the definition of $P_N(e)$ at the beginning of \cite[Section~19.6]{QEDI}, to \cite[Theorem~19.16]{QEDI} and \cite[Theorem~5.83]{QEDI}, and to the assertions on $m_{0,N}(e)$ in \cite[Theorem~21.3]{QEDI} and in Step~3 of its proof in \cite[Section~21.9]{QEDI}.
The scheme of clause (c) has the same three parts as the scheme of \cite[Theorem~19.16]{QEDI} and of \cite[Definition~5.5(10)]{QEDI}: constants, the basis $\{e^{0}_{i}\}$, and the systems on the tori $T_k(R_{\mathrm{top}})$, $k<k_{\mathrm{top}}$ (those of Proposition~26.5 here, the systems of primitives of \cite[Lemma~19.7]{QEDI} there).
The following points differ.
Here the parameter $\rho$ and the weights are among the constants of the scheme.
Clause (c) states explicitly that no probe, and no number of the line (C10) formed with a probe, enters the definition of $P_N(e)$, while \cite[Theorem~19.16]{QEDI} names the constants chosen in the stages 1--10 of the order of choice of \cite[Section~5.6.1]{QEDI} and states that $P_N(e)$ is independent of the torus $R\in\mathrm{Tor}(e)$ and of the twist $\theta$.
Clause (d) is the part of Step~3 in \cite[Section~21.9]{QEDI} that concerns $m_{0,N}(e)$, without the inequality $\boldsymbol{\mu}_N>-\frac{1}{8}$ drawn there.
The bound in (d) rests on the inequality $w^{\mathrm{m}}_0<1$, which holds for every integer $N\ge0$ because $e\le\frac{1}{3}$, $L>1$ and $\frac{3}{4}-8\varepsilon>0$.
For $N\ge N_{\min}\left(R_{\mathrm{top}}(e),e\right)$ this inequality is the third clause of (T1) in Section~I.8.4, taken for the torus $R_{\mathrm{top}}(e)$.
Finally, $m_{0,N}(e)$ is the first entry of $P_N(e)$ for every $N$ for which the pair is defined, that is, for every $N\ge N_{\min}\left(R_{\mathrm{top}}(e),e\right)$, whereas \cite[Theorem~21.3]{QEDI} puts $m_{0,N}(e)=0$ for all $N<N_0(e)$.
Since $N_{\min}\left(R_{\mathrm{top}}(e),e\right)\le N_0(e)$, the two sequences can differ only at the finitely many $N$ with $N_{\min}\left(R_{\mathrm{top}}(e),e\right)\le N<N_0(e)$;
at these $N$ the first entry of $P_N(e)$ is needed in Proposition~32.3\,(10.1), (10.2) on the tori $R>R^{*}(e)$ (Section~I.6.2).
\end{named}

\subsection*{29. The terminal integral}

\subsubsection*{Proposition 29 (on the small-field set the twisted measure with external connection is the terminal integral)}

\begin{named}{Definition}
Let (FL) hold, and let $\left(m_{\mathrm{in}},\varepsilon^{\mathrm{v}}_{\mathrm{in}}\right)$ be a real pair that is controlled on $R$, with its trajectories on $R$.
The index $*$ refers to the level $k_*(R)$.
For $\mathscr{C}\in\mathfrak{X}_{k_*}$ and $A\in\mathrm{Dom}_{k_*}$ put
\[
Y^x_*(A;\mathscr{C}):=y^x_{k_*}\left[\varepsilon^{\mathrm{v},t}_*,m_*,\mathscr{E}_*\right](A;\mathscr{C}),\qquad
Y_*(A):=Y^x_*(A;0),\qquad
D_*(\cdot):=D_{k_*}(\cdot),
\]
where $y^x_{k_*}$ is the function of Section~I.5.7, formed with $\mathfrak{P}_*:=\mathfrak{P}_{k_*}$.
Explicitly, let $\mathscr{A}:=H_{k_*}A$ and $\mathscr{T}:=\mathscr{A}+\mathscr{C}$, let $(\psi,\chi)=H^{\#}_{k_*}(\mathscr{T})\Psi$, and let the barred fields be those of Section~I.5.4.
Then
\begin{multline*}
Y^x_*(A;\mathscr{C})=\left(\det D_*(0)\right)^{-1}\cdot\exp\left\{-\varepsilon^{\mathrm{v},t}_*\cdot\ell_*^3+\sum_X\mathfrak{P}_*(X,\mathscr{C})\right\}\\
\cdot\int\mathrm{D}\Psi\,\exp\left\{-\left\langle\bar{\Psi},D_*(\mathscr{T})\Psi\right\rangle-m_*\cdot\langle\bar{\psi},\psi\rangle+\sum_X\mathscr{E}_*(X;\mathscr{A},\mathscr{C},\psi,\chi)\right\}.
\end{multline*}
Here $\int\mathrm{D}\Psi$ is the Berezin integral (I.3.5) with $\nu=1$ over the $4\ell_*^3$ pairs of generators $\Psi_{\beta}(y)$, $\bar{\Psi}_{\beta}(y)$, $y\in T^0_{k_*}$, $\beta=1,\dots,4$ (as $k_*\ge1$, the block field has four components).
The sums are over $X\in D_{k_*}$, and $\langle\bar{\psi},\psi\rangle$ is the form of Section~I.3.6 on $T^{\mathrm{f}}_{k_*}$, with $192$ components per photon site.
By Proposition~28.1 one has $\mathscr{E}_*|_0=E_*$ and $\mathfrak{P}_*(\cdot,0)=0$; hence
\begin{equation}
\begin{split}
Y_*(A)={}&\left(\det D_*(0)\right)^{-1}\\
&\cdot\int\mathrm{D}\Psi\,\exp\left\{-\left\langle\bar{\Psi},D_*(\mathscr{A})\Psi\right\rangle-m_*\cdot\langle\bar{\psi},\psi\rangle-\varepsilon^{\mathrm{v},t}_*\cdot\ell_*^3+\sum_XE_*(X;\mathscr{A},\psi,\chi)\right\}.
\end{split}
\tag{Y}
\end{equation}
An external connection is an $a\in C^{1}(n)\otimes\mathbb{C}$ such that $a^{[j]}\in\mathfrak{X}_j$, formed with the side $\ell_j(R)$, for all $j\le k_*$; for such $a$ put $\mathscr{C}_*:=a^{[k_*]}$.
The density $\mathrm{dens}^{\theta,a}_N$ and the measure $\nu^{\theta,a}$ are those of (I.7.2), formed with the pair $\left(m_{\mathrm{in}},\varepsilon^{\mathrm{v}}_{\mathrm{in}}\right)$.
\end{named}

\begin{named}{Proposition 29}
Let (FL) hold, let $\left(m_{\mathrm{in}},\varepsilon^{\mathrm{v}}_{\mathrm{in}}\right)$ be a real pair that is controlled on $R$, let $a$ be an external connection, and let $Y^x_*$, $\mathscr{C}_*=a^{[k_*]}$ and $\nu^{\theta,a}$ be as in the preceding Definition.
\begin{itemize}
\item[(a)] For $A\in\mathrm{Dom}_{k_*}$, with $\zeta=(\zeta_j)_{j<k_*}$,
\[
Y^x_*(A;\mathscr{C}_*)=\int\prod_{j<k_*}\mathrm{d}\mu_I(\zeta_j)\,\prod_{j<k_*}\mathbf{1}_{\mathrm{Box}_j}(\zeta_j)\cdot y^x_0\left[\varepsilon^{\mathrm{v}}_{\mathrm{in}},m_{\mathrm{in}},0\right]\left(\mathrm{Af}_0(A,\zeta);a\right),
\]
where
\[
y^x_0\left[\varepsilon^{\mathrm{v}}_{\mathrm{in}},m_{\mathrm{in}},0\right](A_0;a)=\exp\left(-\varepsilon^{\mathrm{v}}_{\mathrm{in}}\cdot n^3\right)\cdot\frac{\det\left(D_{A_0+a}+\bar{m}\cdot\eta+m_{\mathrm{in}}\right)}{\det\left(D_0+\bar{m}\cdot\eta\right)}.
\]
On the compact set $\prod_j\mathrm{Box}_j$ the integrand is a continuous function of $\zeta$.
\item[(b)] Let $\boldsymbol{\lambda}$ be a real function on $T^0_{k_*}$.
If $A$ and $A+d\boldsymbol{\lambda}$ lie in $\mathrm{Dom}_{k_*}$, then
\[
Y^x_*(A+d\boldsymbol{\lambda};\mathscr{C}_*)=Y^x_*(A;\mathscr{C}_*).
\]
\item[(c)] For every $\theta\in\mathbb{R}^3$ and every bounded Borel function $\Phi$ on ${V_{\mathrm{t}}}(T^0_{k_*})$,
\[
\nu^{\theta,a}\left(\mathbf{1}_{\mathfrak{G}}\cdot\Phi(A_*)\right)=\int\mathrm{d}\mathsf{m}_{*}(A)\,\Phi(A)\cdot\chi_*(A)\cdot Y^x_*\left(A+c_{\theta};\mathscr{C}_*\right).
\]
The function $A\mapsto\chi_*(A)\cdot Y^x_*\left(A+c_{\theta};\mathscr{C}_*\right)$ is bounded and Borel on ${V_{\mathrm{t}}}(T^0_{k_*})$, and on the set $\{\chi_*=1\}$ it is continuous in $(\theta,A)$.
\end{itemize}
\end{named}

No constant enters Proposition~29.
The twist enters only through the harmonic part $c_{\theta}$ of the first argument of $Y^x_*$, and $a$ enters only through $\mathscr{C}_*$.

\emph{Conditions on the constants.}
Of the conditions listed in Proposition~14.2, the following enter:
(A1) (first clause), (A2), (A8) (first and fourth clause), (A10) (third and fourth clause); (M3), of which it is used that $\gamma_{\mathrm{w}}>0$; (L5), (M1), (S1), (C1); and the range $0<\rho<1$ of (C0), for a bound of the determinant of level~$0$ of which only the finiteness is used.
All of them except the last are hypotheses of \cite[Proposition~5.38\,(i), (ii)]{QEDI}, of the results \cite[Propositions~5.37 and~5.32]{QEDI} on which it rests, and of \cite[Lemma~17.4(a)]{QEDI}; the last one is the only condition used directly in the proof of Proposition~29.

\begin{named}{Remark}
Proposition~29 corresponds to \cite[Theorem~19.17]{QEDI}, together with the representation of $Y_*$ displayed in \cite[Section~19.7]{QEDI}, and to \cite[Theorem~5.84]{QEDI}.
The following points differ.
The determinant of level~$0$ has the size $192\,n^3$, while the terminal integral is over $4\ell_*^3$ pairs of generators.
The gauge functions are functions of the photon sites $c$.
In (c) the functions $\Phi$ form a single class, the bounded Borel functions on ${V_{\mathrm{t}}}(T^0_{k_*})$, the space that carries $\mathrm{d}\mathsf{m}_{*}$; moreover (c) contains the boundedness and the Borel measurability of $A\mapsto\chi_*(A)\cdot Y^x_*\left(A+c_{\theta};\mathscr{C}_*\right)$.
Clause (a) also states that the integrand is continuous on the compact set $\prod_j\mathrm{Box}_j$.
The application of Fubini's theorem by which the one-step identity is iterated over the levels, in the proof of (a) as in that of \cite[Theorem~19.17(a)]{QEDI}, rests on this continuity, taken at each level $k\le k_*$ on $\prod_{j<k}\mathrm{Box}_j$.
The proof of (c), like that of \cite[Theorem~19.17(c)]{QEDI}, uses that $Q_{k_*}$ preserves means, that is, $\mathrm{mean}(Q_{k_*}A_0)=\mathrm{mean}(A_0)$ for $A_0\in C^{1}(n)$; this is the identity $\mathrm{mean}(Q_jA)=\mathrm{mean}(A)$ of \cite[Proposition~14.4(c)]{QEDI}, taken with $a=1$, $j=k_*$ and $l=n$.
\end{named}

\subsection*{G6. Conventions for the Propositions 38.m--47.m and 50}

\paragraph{G6.A1. Conventions.}
In the Propositions~38.m--47.m and~50 the following data are fixed throughout and are not repeated in the statements: $\rho\in(0,1)$, the standard weights, $L$, $M$ and all constants of Section~I.8 in their order; $\bar{m}\in(0,1]$; and $h$.
Lattice units are those of Sections~I.2.5 and~I.7.1, and physical units those of Section~I.7.5.
A cochain on a torus is a periodic cochain on $\mathbb{Z}^3$, and every $g_{\mathrm{w}}\in W_{\mathrm{af}}$ acts on cochains through $\mathbb{Z}^3$ (Section~I.4.7).
Hence no action depends on an induced map of the sites of a torus.

\paragraph{G6.A3. Notation.}
In the Propositions~38.m--47.m and~50 the following letters have a fixed meaning.
\begin{description}
\item[$\mathrm{hw}$] the width of a general coordinate cell structure; $\mathrm{hw}\in\frac{1}{2}\mathbb{Z}$, $\mathrm{hw}>0$.
It is the width $h$ of \cite[Definition~2.18]{QEDII}; for the cell structure $\mathcal{Q}_{*}$ it is $\mathrm{hc}$.
\item[$o=(o_1,o_2,o_3)$] the offsets of such a cell structure.
\item[$\mathrm{sv}=(\mathrm{sv}_1,\mathrm{sv}_2,\mathrm{sv}_3)$] numbers of cells per side.
\item[$\mathrm{ncs}$] the number of cells per side of a reference torus.
\item[$\check{j}$] an integer $\ge2$.
\item[$\mathrm{kr}$] a ratio of two sides, a rational number.
\item[$\mathrm{wq}$] a complex weight on $C^{1}(n)$.
\item[$\mathrm{Ypt}$] a finite set of points of $\mathbb{R}^3$.
\item[$\mathrm{ix}$] an index in $\{0,+,-\}$; the three tests are $g_{\mathrm{ix}}$, and their probes are $\beta_{\mathrm{ix}}$.
\item[$s_{B}$, $C_{3\mathrm{p}}$] two numbers of the Propositions~43.m (constants of the line (C10) of Section~I.8.2).
\item[$Z^{\mathrm{ph}}_n$, $\mathrm{Ld}(n)$] the photon numbers of \cite[Section~21.8.1 and Lemma~21.32]{QEDI}.
\end{description}
The numbers $\mathrm{Mw}_m$ and $\mathrm{Tw}_m$ stand in the places of $\dot{M}_m$ and $\boldsymbol{\Theta}^{\mathrm{af}}_m$ of \cite[Theorem~21.2]{QEDI}.
In Theorem~44\,(B0), (P8) the index $i$ numbers a sequence of cutoffs $(N_i,T_i)$, with
\[
\eta_i:=L^{-N_i},\qquad
n_i:=8h\cdot L^{T_i+N_i},\qquad
R_i:=8h\cdot L^{T_i};
\]
there $k$ is the number of the tests.
Finally, $\mathfrak{F}(n)$ is the class of the Borel functions $\Phi$ on $C^{2}(n)$ such that $|\Phi(u)|\le\mathrm{cf}\cdot\exp\left(\mathrm{cf}\cdot\|u\|\right)$ for some $\mathrm{cf}$.

\subsection*{38. The determinant on strips; the zero-flux sector; the shift identity}

\subsubsection*{Proposition 38.1. The determinant at complex connections}

\begin{named}{Proposition 38.1}
Let $\boldsymbol{\mu}$ be real, let $q>0$, let $n_\mu\ge1$, $\mu=1,2,3$, be integers, let $B\in C^{1}(n)\otimes\mathbb{C}$, and put
\[
|\operatorname{Im}B|_\infty:=\max_b\left|\operatorname{Im}B(b)\right|,\qquad
M_{\mathrm{ch}}:=\boldsymbol{\mu}+\frac{\rho}{2}\cdot\mathrm{sw}.
\]
Here, for sides $n_\mu\ge1$, $C^{j}(n)$ denotes the space of the real $j$-cochains on $\mathbb{Z}^3$ that are periodic under $\sum_\mu n_\mu\mathbb{Z}e_\mu$ (Section~G6.A1; these are the data of Proposition~46, whereas Section~I.7.1 has $n_\mu\ge4$), and the links of the torus are the $3|n|$ orbits of links.
The number $\operatorname{det}_n\left(D_B+\boldsymbol{\mu}\right)$ is the determinant of the linear map $D_B+\boldsymbol{\mu}$ of $\mathfrak{S}^{\mathrm{ch}}(1;n)$ (Proposition~46\,(e)).
\begin{itemize}
\item[(a)] $\operatorname{det}_n\left(D_B+\boldsymbol{\mu}\right)$ is a polynomial in the $6|n|$ numbers $\exp\left(iq\cdot B(b)\right)$, $\exp\left(-iq\cdot B(b)\right)$, where $b$ ranges over the links of the torus, and its degree in each of them is at most $32$.
Hence it is an entire function of $B\in\mathbb{C}^{3|n|}$.
(It is the value at $u_b=\exp\left(iq\cdot B(b)\right)$ of the determinant polynomial $P_{\mathbb{T}}$ of \cite[Definition~4.4]{QEDII}, a Laurent polynomial in the indeterminates $u_b$, formed with the mass $M=M_{\mathrm{ch}}$ for the general torus $\mathbb{T}=\left(\sum_\mu n_\mu\mathbb{Z}e_\mu,\ \chi\right)$, $\chi(n_\mu e_\mu)=-1$, in the sense of \cite[Definition~4.1]{QEDII};
for $n_\mu\ge4$ this is the coordinate torus $C(n_1,n_2,n_3)$ of \cite[Definition~2.11]{QEDII}.)
\item[(b)] $\operatorname{det}_n\left(D_{B+r_q\cdot m}+\boldsymbol{\mu}\right)=\operatorname{det}_n\left(D_B+\boldsymbol{\mu}\right)$ for every $m\in C^{1}(n;\mathbb{Z})$.
\item[(c)] $\operatorname{det}_n\left(D_{B+d\boldsymbol{\lambda}}+\boldsymbol{\mu}\right)=\operatorname{det}_n\left(D_B+\boldsymbol{\mu}\right)$ for every complex zero-cochain $\boldsymbol{\lambda}$ of the torus.
\item[(d)] One has
\[
\left|\operatorname{det}_n\left(D_B+\boldsymbol{\mu}\right)\right|\le\left(|M_{\mathrm{ch}}|+\frac{1+\rho}{2}\left(\mathrm{sw}-6+6\exp\left(q\cdot|\operatorname{Im}B|_\infty\right)\right)\right)^{192|n|}.
\]
\item[(e)] For $\theta\in\mathbb{C}^3$ and $\mu=1,2,3$,
\[
H_{2\pi e_\mu}=r_q\cdot m^{(\mu)}-d\boldsymbol{\lambda}^{(\mu)},
\]
where
\[
m^{(\mu)}(c,c+e_\nu):=\delta_{\mu\nu}\cdot\mathbf{1}\left[c_\mu\equiv n_\mu-1\pmod{n_\mu}\right],\qquad
\boldsymbol{\lambda}^{(\mu)}(c):=-\frac{r_q}{n_\mu}\cdot c_\mu\quad\text{for }c_\mu\in\{0,\dots,n_\mu-1\}.
\]
Consequently the function $\theta\mapsto\operatorname{det}_n\left(D_{B+H_{\theta}}+\boldsymbol{\mu}\right)$ is $2\pi$-periodic in each $\theta_\mu$.
\end{itemize}
\end{named}

\begin{named}{Remark}
Proposition~38.1 corresponds to \cite[Theorem~21.1(d1)]{QEDI}.
It corresponds further to three of the properties of the fermion weight stated in \cite[Section~10.1.2]{QEDI}: the invariance under $A\mapsto A+d\boldsymbol{\lambda}$; the invariance under $A\mapsto A+\frac{2\pi}{q}\cdot m$ for every integer one-cochain $m$; and the continuity and boundedness on real arguments, together with the entire extension that keeps both invariances and is bounded on each set $\left\{|\operatorname{Im}A|_\infty\le c\right\}$.
It also corresponds to the first of the three properties of the determinant $\Delta$ stated in \cite[Section~11.1.2]{QEDI}: $\Delta$ is entire, bounded on each set $\left\{|\operatorname{Im}B|_\infty\le c\right\}$, of period $2\pi$ in its value at each link, and invariant under $B\mapsto B+d\boldsymbol{\lambda}$.
Clause (d) corresponds to the bound of the determinant in \cite[Section~7.1.2]{QEDI}, and clause (e), together with (I.6.1), to the period $2\pi$ in $\theta$ in \cite[Lemma~7.7(a)]{QEDI}.
The following points differ.
The operator is the chamber operator: a photon link enters through the phases of $8+8$ oriented fermion bonds of type~$3$ with four components each, whence the degree $32$ in (a).
The boundedness on the sets $\left\{|\operatorname{Im}B|_\infty\le c\right\}$ is replaced by the explicit bound (d).
The relation in (e), which gives the period $2\pi$ in $\theta$ used in \cite[Lemma~11.2]{QEDI}, is part of the statement.
\end{named}

\subsection*{50. (W1)--(W6) and the invariances of the functionals, on rectangular tori of either parity}

In this subsection (W1)--(W6) refer to the six properties \cite[(10.3)--(10.8)]{QEDI} of the fermion weight, stated in \cite[Section~10.1.2]{QEDI}:
non-negativity at real arguments,
invariance under the addition of a coboundary,
invariance under the addition of $2\pi/q$ times an integer valued one-cochain,
continuity and boundedness at real arguments together with an entire extension that is bounded on every set on which the imaginary part is bounded by a constant,
evenness,
and invariance under the isometries of the lattice.

\subsubsection*{Proposition 50}

\begin{named}{Proposition 50}
Let $\boldsymbol{\mu}$ be real, let $q>0$ and $r_q=2\pi/q$, and let $n=(n_1,n_2,n_3)$ and $n'=(n'_1,n'_2,n'_3)$ be triples of integers $n_\mu\ge1$, $n'_\mu\ge1$, of either parity.
For $B\in C^1(n)\otimes\mathbb{C}$ the number $\operatorname{det}_n\left(D_B+\boldsymbol{\mu}\right)$ is the determinant on $\mathfrak{S}^{\mathrm{ch}}(1;n)$ (Proposition~46\,(e)).
Let $g_{\mathrm{w}}=(a,r)\in W_{\mathrm{af}}$ satisfy
\[
r\left(\sum_\mu n_\mu\mathbb{Z}e_\mu\right)=\sum_\mu n'_\mu\mathbb{Z}e_\mu ,
\]
and write $re_\mu=\epsilon_\mu\cdot e_{\mu'}$, where $\epsilon_\mu=\pm1$ and $\mu\mapsto\mu'$ is a permutation of the directions; then $n'_{\mu'}=n_\mu$.
The pull-back $g_{\mathrm{w}}^{*}$ maps $C^j(n')$ onto $C^j(n)$, and $g_{\mathrm{w}*}:=\left(g_{\mathrm{w}}^{-1}\right)^{*}$ maps $C^j(n)$ onto $C^j(n')$.
Both maps are orthogonal for the counting pairings and commute with $d$.
\begin{itemize}
\item[(a)] Let $B$ be real.
In the two eigenspaces of $\Sigma_3$ the operator $D_B+\boldsymbol{\mu}$ has the block form
\[
D_B+\boldsymbol{\mu}=\operatorname{diag}\left(D_B^{\boldsymbol{\sigma}}+\boldsymbol{\mu},\left(D_B^{\boldsymbol{\sigma}}+\boldsymbol{\mu}\right)^{*}\right),
\]
and therefore
\[
\operatorname{det}_n\left(D_B+\boldsymbol{\mu}\right)=\left|\det\left(D_B^{\boldsymbol{\sigma}}+\boldsymbol{\mu}\right)\right|^2\ge0 .
\]
Moreover $\operatorname{Re}\left(D_B+\boldsymbol{\mu}\right)=\boldsymbol{\mu}+\frac{\rho}{2}\cdot\nabla_B^{*}\nabla_B\ge\boldsymbol{\mu}$, and the norm of $D_B+\boldsymbol{\mu}$ on $\ell^2$ does not exceed $\left|M_{\mathrm{ch}}\right|+\frac{1+\rho}{2}\cdot\mathrm{sw}$.
In particular $\operatorname{det}_n\left(D_B+\boldsymbol{\mu}\right)\ge\boldsymbol{\mu}^{192|n|}>0$ whenever $\boldsymbol{\mu}>0$.
\item[(b)] By Proposition~38.1\,(a)--(d), the function $B\mapsto\operatorname{det}_n\left(D_B+\boldsymbol{\mu}\right)$ is entire.
It is invariant under $B\mapsto B+d\boldsymbol{\lambda}$ for every complex zero-cochain $\boldsymbol{\lambda}$ of the torus, and under $B\mapsto B+r_q\cdot m$ for every $m\in C^1(n;\mathbb{Z})$.
It is bounded and continuous on $C^1(n)$, and it is bounded on every set $\left\{B:\ \max_b\left|\operatorname{Im}B(b)\right|\le\mathrm{cf}\right\}$.
\item[(c)] $\operatorname{det}_n\left(D_{-B}+\boldsymbol{\mu}\right)=\operatorname{det}_n\left(D_B+\boldsymbol{\mu}\right)$ for every complex $B$.
\item[(d)] The map $U_{\mathrm{gw}}$, given by $\left(U_{\mathrm{gw}}f\right)(x):=S_{\mathrm{gw}}f\left(g_{\mathrm{w}}^{-1}x\right)$ as in Section~I.4.7, maps $\mathfrak{S}^{\mathrm{ch}}(1;n)$ onto $\mathfrak{S}^{\mathrm{ch}}(1;n')$, and for every complex $B$
\[
U_{\mathrm{gw}}\left(D_B+\boldsymbol{\mu}\right)U_{\mathrm{gw}}^{-1}=D_{g_{\mathrm{w}*}B}+\boldsymbol{\mu},\qquad\operatorname{det}_{n'}\left(D_{g_{\mathrm{w}*}B}+\boldsymbol{\mu}\right)=\operatorname{det}_n\left(D_B+\boldsymbol{\mu}\right).
\]
\end{itemize}
From now on $n_\mu\ge4$ and $n'_\mu\ge4$.
\begin{itemize}
\item[(e)] (Convergence and positivity.)
For $\Phi\in\mathfrak{F}(n)$ the sum and the integral in (I.7.1) converge absolutely, and
\[
\sum_{w\in\mathbb{Z}^3}\left|Z^{(w)}_n[\Phi]\right|\le Z_n\left[|\Phi|\right]<\infty .
\]
For every $w\in\mathbb{Z}^3$ one has
\[
0<Z^{(w)}_n[1]\le Z^{(0)}_n[1]\le Z_n[1]<\infty .
\]
Hence, for every real $\boldsymbol{\mu}$, the functionals $E^{\mathrm{af}}_n$ and $E^{\mathrm{av}}_n$ are expectations of probability laws of $F^{\mathrm{u}}$.
\item[(f)] For $\Phi\in\mathfrak{F}(n)$,
\[
Z_{n'}\left[\Phi\circ g_{\mathrm{w}}^{*}\right]=Z_n[\Phi],\qquad Z^{(0)}_{n'}\left[\Phi\circ g_{\mathrm{w}}^{*}\right]=Z^{(0)}_n[\Phi].
\]
\item[(g)] (Single sectors.)
For $w\in\mathbb{Z}^3$ set $w_{\nu\mu}:=-w_{\mu\nu}$, and define $g_{\mathrm{w}}\cdot w\in\mathbb{Z}^3$, extended by the same convention, through
\[
\left(g_{\mathrm{w}}\cdot w\right)_{\mu'\nu'}:=\epsilon_\mu\epsilon_\nu\cdot w_{\mu\nu}.
\]
Then $w\mapsto g_{\mathrm{w}}\cdot w$ is a bijection of $\mathbb{Z}^3$ with $g_{\mathrm{w}}\cdot0=0$, and $\left\|h_{g_{\mathrm{w}}\cdot w}\right\|=\left\|h_w\right\|$, the left side being formed on $n'$ and the right side on $n$.
For a closed $u'\in C^2(n')$ the flux of $g_{\mathrm{w}}^{*}u'$ is $w$ if and only if the flux of $u'$ is $g_{\mathrm{w}}\cdot w$.
Finally,
\[
Z^{(g_{\mathrm{w}}\cdot w)}_{n'}\left[\Phi\circ g_{\mathrm{w}}^{*}\right]=Z^{(w)}_n[\Phi]\qquad\text{for every }\Phi\in\mathfrak{F}(n).
\]
\item[(g$'$)] (Sign.)
$Z^{(-w)}_n\left[\Phi(-\,\cdot\,)\right]=Z^{(w)}_n[\Phi]$ for every $w\in\mathbb{Z}^3$ and every $\Phi\in\mathfrak{F}(n)$.
Consequently $Z_n$ and $Z^{(0)}_n$ are invariant under $F^{\mathrm{u}}\mapsto-F^{\mathrm{u}}$.
\item[(g$''$)] On a cubic torus $(n,n,n)$ the functionals $E^{\mathrm{af}}_n$ and $E^{\mathrm{av}}_n$ are invariant under $u\mapsto g_{\mathrm{w}}^{*}u$ for every $g_{\mathrm{w}}\in W_{\mathrm{af}}$, and under $u\mapsto-u$.
Invariance means that $E\left[\Phi\circ g_{\mathrm{w}}^{*}\right]=E[\Phi]$ and $E\left[\Phi(-\,\cdot\,)\right]=E[\Phi]$ for every $\Phi\in\mathfrak{F}(n)$, where $E$ is $E^{\mathrm{af}}_n$ or $E^{\mathrm{av}}_n$.
Equivalently, the image of the law of $F^{\mathrm{u}}$ under $u\mapsto g_{\mathrm{w}}^{*}u$, respectively under $u\mapsto-u$, is this law itself.
\item[(h)] (The free ratio.)
Let $n_\mu\ge1$, let $\mathbf{m}>0$, and let $a\in C^1(n)$ be a real one-cochain; in this clause the letter $a$ stands for this cochain, and the translation part of $g_{\mathrm{w}}=(a,r)$ is not written.
For $t\in\mathbb{C}$ put
\[
N_a(t):=\frac{\operatorname{det}_n\left(D_{ta}+\mathbf{m}\right)}{\operatorname{det}_n\left(D_0+\mathbf{m}\right)} .
\]
The denominator is positive.
The function $N_a$ is entire and even, it is real and positive on the real axis, with
\[
N_a(t)\ge\frac{\mathbf{m}^{192|n|}}{\operatorname{det}_n\left(D_0+\mathbf{m}\right)}\qquad(t\in\mathbb{R}),
\]
and $N_a(0)=1$; moreover $N_{g_{\mathrm{w}*}a}$, formed on $n'$, coincides with $N_a$.
In particular the numbers $c_m$, $a_{\mathrm{ch}}$, $b_{\mathrm{ch}}$ of Section~I.7.6 remain the same if $\hat{\mathfrak{a}}$ is replaced by $\hat{\mathfrak{a}}_{\mathfrak{c}}$, or if the placement $v$ is replaced by another vector of $\eta\mathbb{Z}^3$.
The last assertion concerns the data $0<\eta\le1$, $e>0$, $q=e\cdot\eta^{\frac{1}{2}}$, $\mathbf{m}=\bar{m}\cdot\eta$ and $a=\hat{\mathfrak{a}}/e$, for which $N_a$ is the free ratio $N_0$ of Section~I.7.6, formed on an arbitrary cubic torus $(n,n,n)$ with $n=R/\eta$ an integer $\ge1$;
here $\hat{\mathfrak{a}}$ is given, also for $R/\eta=1$, by the displays of Section~G5.1.3.
For the copies $\hat{\mathfrak{a}}_{\mathfrak{c}}$ the data are those of Section~I.6.4, namely $\eta=L^{-N}$, $e\cdot h^{\frac{1}{2}}\le e_1$ and $N\ge N_0(e)$, and the number $\mathrm{per}$ need not divide $n$.
\end{itemize}
\end{named}

\begin{named}{Remark}
Proposition~50 corresponds to \cite[Theorem~21.1\,(d2)--(d5)]{QEDI}, whose proofs are given in \cite[Section~21.3.2]{QEDI}.
It also corresponds to the properties (W1)--(W6) of \cite[Section~10.1.2]{QEDI}, to the absolute convergence, the positivity of $Z_n[1]$ and the invariances in \cite[Section~10.1.3]{QEDI}, and to the positivity of the sector sums in \cite[Section~10.1.4]{QEDI}; to the properties \cite[(11.2)--(11.4)]{QEDI} of the determinant in \cite[Section~11.1.2]{QEDI}; to the sign of the determinant and the convergence in \cite[Section~7.1.2]{QEDI}, which are here (a) and (e), together with Sections~I.7.2, I.7.3 and Propositions~38.1\,(d) and 31.1\,(a)--(c); to \cite[Lemma~7.7\,(b)]{QEDI}, which is here (d), (f) and (g$''$); and to \cite[Lemma~7.7\,(c)]{QEDI}, which is here (c), (g$'$) and (g$''$).
Clause by clause, (a) is the counterpart of (W1) and of \cite[Theorem~21.1\,(d2)]{QEDI}, (b) of (W2)--(W4) and of \cite[Theorem~21.1\,(d1)]{QEDI}, (c) of (W5) and of \cite[Theorem~21.1\,(d3)]{QEDI}, (d) of (W6) and of \cite[Theorem~21.1\,(d4)]{QEDI}, and (f), (g), (g$'$), (g$''$) of \cite[Theorem~21.1\,(d5)]{QEDI}.
The evenness (c) rests on the transposition identity (I.O5); its counterpart in the proof of \cite[Theorem~21.1\,(d3)]{QEDI} is the identity $\mathbf{C}\cdot D[B]^{T}\cdot\mathbf{C}^{-1}=D[-B]$ for the Wilson operator $D[B]$, with the constant matrix $\mathbf{C}$ of that proof.

The mathematical differences are as follows.
The Wilson determinant, taken on a space of dimension $4|n|$, is replaced by $\operatorname{det}_n\left(D_B+\boldsymbol{\mu}\right)$, taken on a space of dimension $192|n|$.
Every real lattice mass $\boldsymbol{\mu}$ is admitted, as required by (I.K10); the restrictions to masses $>-1$ in \cite[Section~10]{QEDI} and to masses $>-\frac{1}{2}$ in \cite[Section~11]{QEDI} and \cite[Section~21.7]{QEDI} do not occur.
In \cite[Section~21.1.4]{QEDI} the laws $E^{\mathrm{af}}_n$ and $E^{\mathrm{av}}_n$ are quotients whose denominators are shown to be positive for masses $>-1$ only; clause (e) gives this positivity for every real $\boldsymbol{\mu}$.
The determinant is considered on tori with sides $n_\mu\ge1$ and the functionals on tori with sides $n_\mu\ge4$, whereas \cite{QEDI} works with sides $n_\mu\ge2$.
The maps $g_{\mathrm{w}*}$, $g_{\mathrm{w}}^{*}$ and $p\mapsto g_{\mathrm{w}}p$ belong to the pair $(a,r)$ and are defined through $\mathbb{Z}^3$ (Section~G6.A1); on a torus with a side equal to $2$ they are not determined by the map that $g_{\mathrm{w}}$ induces on the sites, and under the condition $n_\mu,n'_\mu\ge4$ of (e)--(g$''$) this case does not arise.
The product formula for the free Wilson determinant in \cite[Section~10.1.3]{QEDI} has no counterpart here, and the positivity of the partition functions rests on \cite[Section~4]{QEDII}.
Clause (g), the rule for a single sector $w\ne0$, is the clause on the sectors of fixed flux in \cite[Theorem~21.1\,(d5)]{QEDI}.
In \cite[Definition~13.7\,(e)]{QEDII} it is stated that the elements of $W_{\mathrm{af}}$ permute the single flux sectors, the fluxes being relabelled by a permutation of their components and changes of sign, and that the harmonic-averaged law is invariant under the translations by $\mathbb{Z}^3$, under $W_{\mathrm{af}}$ on a cubic torus and under $F\mapsto-F$; these are the counterparts of (g), (g$'$) and (g$''$) in \cite{QEDII}.
The same relabelling of the fluxes accompanies the replacement of the time direction $e_1$ by $e_2$ or $e_3$ in \cite[Lemma~13.8\,(c)]{QEDII}.
Clause (h) collects the last sentence of \cite[Theorem~21.1\,(d5)]{QEDI} and \cite[Lemma~20.4\,(b)]{QEDI}; in particular the numbers $c_m$, $a_{\mathrm{ch}}$, $b_{\mathrm{ch}}$ formed with $\tau_v\beta$, $v\in\eta\mathbb{Z}^3$, coincide with those formed with $\beta$.
For $\boldsymbol{\sigma}^2$ the corresponding facts are Proposition~47.1\,(b5$'$) and Proposition~41\,(c).
\end{named}

\subsection*{38 (continued). The determinant on strips; the zero-flux sector; the shift identity}

\subsubsection*{Proposition 38.2. The zero-flux sector is the harmonic average}

\begin{named}{Proposition 38.2}
Let $q>0$, and let the sides $n_\mu\ge4$ be of either parity.
Let $C^1(n)={V_{\mathrm{g}}}\oplus{V_{\mathrm{h}}}\oplus{V_{\mathrm{t}}}$ be the decomposition into the exact cochains, the constant cochains and the orthogonal complement of both, and let $\sigma$ be the Lebesgue measure of ${V_{\mathrm{t}}}$.
Let $\Delta_0:=d^{*}d$ on $C^0(n)$, let $\operatorname{det}'\Delta_0$ be the product of the non-zero eigenvalues of $\Delta_0$, and put
\[
\mathrm{Nc}_n:=r_q^{|n|+2}\cdot\left(\operatorname{det}'\Delta_0\right)^{\frac{1}{2}} .
\]
\begin{itemize}
\item[(a)] Let $\mathrm{wq}\colon C^1(n)\to\mathbb{C}$ be bounded and continuous, with $\mathrm{wq}(A+d\boldsymbol{\lambda})=\mathrm{wq}(A)$ for every real zero-cochain $\boldsymbol{\lambda}$ and $\mathrm{wq}(A+r_q\cdot m)=\mathrm{wq}(A)$ for every $m\in C^1(n;\mathbb{Z})$.
Then for every $\Phi\in\mathfrak{F}(n)$
\begin{multline*}
\sum_{\substack{\mathfrak{z}\in C^2(n;\mathbb{Z}),\ d\mathfrak{z}=0\\ \text{flux }0}}\ \int_{[0,r_q)^{\mathrm{links}}}\mathrm{d}A\,\Phi\left(dA+r_q\mathfrak{z}\right)\cdot\exp\left(-\frac{1}{2}\left\|dA+r_q\mathfrak{z}\right\|^2\right)\cdot\mathrm{wq}(A)\\
=\mathrm{Nc}_n\int_{[-\pi,\pi]^3}\frac{\mathrm{d}\theta}{(2\pi)^3}\int_{V_{\mathrm{t}}}\mathrm{d}\sigma(A)\,\Phi(dA)\cdot\exp\left(-\frac{1}{2}\|dA\|^2\right)\cdot\mathrm{wq}\left(A+H_\theta\right),
\end{multline*}
and both sides converge absolutely.
\item[(b)] For every real $\boldsymbol{\mu}$ and every $\Phi\in\mathfrak{F}(n)$,
\[
Z^{(0)}_n[\Phi]=\mathrm{Nc}_n\int_{[-\pi,\pi]^3}\frac{\mathrm{d}\theta}{(2\pi)^3}\int_{V_{\mathrm{t}}}\mathrm{d}\sigma(A)\,\Phi(dA)\cdot\exp\left(-\frac{1}{2}\|dA\|^2\right)\cdot\det\left(D_{A+H_\theta}+\boldsymbol{\mu}\right),
\]
both sides converging absolutely; this is (a) for the weight $\mathrm{wq}(A)=\det\left(D_A+\boldsymbol{\mu}\right)$.
It follows that the law of $F^{\mathrm{u}}$ under $E^{\mathrm{av}}_n$ is the law of $dA$ under the probability measure on ${V_{\mathrm{t}}}$ that is proportional to
\[
\left(\int_{[-\pi,\pi]^3}\mathrm{d}\theta\,\exp\left(-\frac{1}{2}\|dA\|^2\right)\cdot\det\left(D_{A+H_\theta}+\boldsymbol{\mu}\right)\right)\mathrm{d}\sigma(A).
\]
\item[(c)] Let the torus be cubic, with $n\ge4$ sites per side, and let the objects of Section~I.7.3 be formed with a real pair $\left(m_{\mathrm{in}},\varepsilon^{\mathrm{v}}_{\mathrm{in}}\right)$ and with $\boldsymbol{\mu}=\bar{m}\cdot\eta+m_{\mathrm{in}}$.
Then $\bar{\mu}$ is the law of $F^{\mathrm{u}}$ under $E^{\mathrm{av}}_n$, and
\[
Z^{(0)}_n[1]=Z^{\mathrm{ph}}_n\cdot\mathrm{zr}_n(q,\boldsymbol{\mu}),
\]
where
\begin{gather*}
Z^{\mathrm{ph}}_n:=\mathrm{Nc}_n\int_{V_{\mathrm{t}}}\mathrm{d}\sigma(A)\,\exp\left(-\frac{1}{2}\|dA\|^2\right)=r_q^{n^3+2}\cdot(2\pi)^{n^3-1}\cdot\exp\left(-\frac{1}{2}\mathrm{Ld}(n)\right),\\
\mathrm{Ld}(n):=\log\operatorname{det}'\Delta_0=\sum_{k\in(\mathbb{Z}/n\mathbb{Z})^3,\ k\ne0}\log\left(\sum_\mu\left(2-2\cos\frac{2\pi k_\mu}{n}\right)\right).
\end{gather*}
Furthermore $\mathrm{zr}_n(q,\boldsymbol{\mu})>0$, and
\[
\bar{\nu}(1)=\frac{\exp\left(-\varepsilon^{\mathrm{v}}_{\mathrm{in}}\cdot n^3\right)\cdot\mathrm{zr}_n(q,\boldsymbol{\mu})}{\det\left(D_0+\bar{m}\cdot\eta\right)}\in(0,\infty).
\]
The number
\[
\mathrm{zV}_n:=\int_{V_{\mathrm{t}}}\mathrm{d}\sigma(A)\,\exp\left(-\frac{1}{2}\|dA\|^2\right)=\frac{(2\pi)^{n^3-1}}{\operatorname{det}'\Delta_0}
\]
satisfies $Z^{\mathrm{ph}}_n=\mathrm{Nc}_n\cdot\mathrm{zV}_n$, and
\[
\int_{V_{\mathrm{t}}}\mathrm{d}\sigma(A)\,G(dA)\cdot\exp\left(-\frac{1}{2}\|dA\|^2\right)=\mathrm{zV}_n\int\mathrm{d}\gamma_N\,G
\]
for every Borel function $G$ on $F_N$ that is non-negative or satisfies $|G(u)|\le\mathrm{cf}\cdot\exp\left(\mathrm{cf}\cdot\|u\|\right)$ for some $\mathrm{cf}$.
In (c) the number $q>0$ is arbitrary and is not tied to $\eta$, and every real $\boldsymbol{\mu}$ occurs, with $m_{\mathrm{in}}:=\boldsymbol{\mu}-\bar{m}\cdot\eta$; the assertions of (c) on $Z^{(0)}_n[1]$, $Z^{\mathrm{ph}}_n$, $\mathrm{zV}_n$ and $\mathrm{zr}_n(q,\boldsymbol{\mu})$ involve $(n,q,\boldsymbol{\mu})$ only.
\end{itemize}
\end{named}

\begin{named}{Remark}
Proposition~38.2 corresponds to \cite[Lemma~10.2\,(b)]{QEDI}, to \cite[Lemma~11.2]{QEDI} with the identity \cite[(11.6)]{QEDI}, to the identification of the harmonic-averaged law in \cite[Section~21.1.4]{QEDI}, and to \cite[Lemma~21.32\,(a)]{QEDI} together with the first identity of \cite[Lemma~21.32\,(b)]{QEDI}.
In \cite[Definition~13.7\,(b)]{QEDII} the harmonic-averaged law is defined as the law of the curvature conditioned on the flux $0$; clauses (b) and (c) express it through the average over the twist.

The mathematical differences are as follows.
Clause (a) is stated for a bounded continuous complex weight with the two invariances, and no positivity of the weight enters.
The proof of \cite[Lemma~11.2]{QEDI} writes the absolute convergence as the finiteness of a positive measure.
\cite[Lemma~11.2]{QEDI} is formulated for the determinant $\Delta$ of \cite[(11.1)]{QEDI}.
In the proof of the third identity of \cite[Lemma~11.7]{QEDI} it is shown that \cite[Lemma~11.2]{QEDI} holds when $\Delta$ is replaced by the weight $\Delta(\cdot+qa)$, which takes complex values when $a$ is complex, and that the constant stays the same, since it does not involve the weight; the identity is then taken at $\Phi=1$.
The argument is that the proof of \cite[Lemma~11.2]{QEDI} uses the weight only through the two invariances in \cite[(11.2)]{QEDI} and through the absolute convergence of the integral over $[0,r_q)^{\mathrm{links}}$ of the series over the closed integer valued two-cochains, and that this absolute convergence holds for $\Delta(\cdot+qa)$.
Indeed, by the bound in \cite[(11.2)]{QEDI}, applied with $c=\max_b\left|\operatorname{Im}(qa_b)\right|$, there is a number $M_a$ with $\left|\Delta(qA+qh+qa)\right|\le M_a$ for all real $A,h\in C^1(n)$;
hence the integral of the series of the moduli is at most $M_a$ times the same integral without the weight, which is finite by the bound on that series in \cite[Section~11.1.3]{QEDI}, a bound in which $\Delta$ does not enter.
In that no positivity of the weight is needed, clause (a) therefore agrees with this proof.
Clause (a) differs from this proof in that it is a statement on every bounded continuous complex weight with the two invariances, whereas in \cite{QEDI} the lemma is stated for $\Delta$ and the case of the weight $\Delta(\cdot+qa)$ is established within the proof of the third identity of \cite[Lemma~11.7]{QEDI}.
For sides $n_\mu\ge4$ the identity needed in that step is the case $\mathrm{wq}(A)=\Delta(qA+qa)$, $\Phi=1$ of (a).
The measures are specified: on the left of (a) the Lebesgue measure of $\mathbb{R}^{\mathrm{links}}$, on the right $\sigma$ and $\mathrm{d}\theta/(2\pi)^3$, with the explicit constant $\mathrm{Nc}_n$; in \cite[Lemma~10.2\,(b)]{QEDI} the sector of flux $0$ is written as an integral over $C^1_n/\Gamma_n$, and the constant $V_n$ that relates it to the twist average is not evaluated.
The fact that a closed integer valued one-cochain whose sums along the three coordinate loops vanish is $d\boldsymbol{\lambda}$ with an integer valued $\boldsymbol{\lambda}$ is derived from the decomposition \cite[(10.1)]{QEDI}.
Every real $\boldsymbol{\mu}$ is admitted, and the sides are $n_\mu\ge4$, of either parity.
Clause (c) identifies the law of $F^{\mathrm{u}}$ under $E^{\mathrm{av}}_n$ with the measure $\bar{\mu}$ of Section~I.7.3, and it contains the positivity $\mathrm{zr}_n(q,\boldsymbol{\mu})>0$, which the logarithms in the definition of $D^{\mathrm{mat}}$ in Section~I.7.6 require.
\end{named}

\subsection*{47. Probes and cells on the even tori}

\subsubsection*{Proposition 47.1. Lattice probes on every torus}

\begin{named}{Proposition 47.1}
Let $\beta$ be a probe (Section~I.7.6), $a=d^{*}\beta$ and $f=dd^{*}\beta$.
For $0<\eta\le1$ and a smooth function $\mathrm{fs}$ on $\mathbb{R}^3$ let
\[
\delta_{\eta,\mu}\mathrm{fs}(y):=\frac{\mathrm{fs}\left(y+\frac{\eta}{2}e_\mu\right)-\mathrm{fs}\left(y-\frac{\eta}{2}e_\mu\right)}{\eta},
\]
let $J_\eta$ be the cube average of Section~I.7.6, and put
\[
a'_\nu:=-\sum_\mu\delta_{\eta,\mu}\beta_{\mu\nu},\qquad\tilde{a}_\eta:=J_\eta a',\qquad f^{(\eta)}_{\mu\nu}:=\delta_{\eta,\mu}a'_\nu-\delta_{\eta,\nu}a'_\mu ;
\]
thus $\tilde{a}_{\eta,\nu}=\mathrm{gs}_\nu[\beta]$ in the notation of Section~G5.1.3.
For a smooth function, vector field or two-form $u$ on $\mathbb{R}^3$ with bounded derivatives, $\left|D^iu\right|_\infty$ denotes the largest of the sup norms of the derivatives of order exactly $i$ of the components of $u$, and $|u|_{C^i}:=\max_{j\le i}\left|D^ju\right|_\infty$.
For a two-form $u\in\mathcal{S}$,
\[
\|u\|_{H^s}^2:=\sum_{\mu<\nu}\int\frac{\mathrm{d}p}{(2\pi)^3}\,\left(1+|p|_2^2\right)^s\cdot\left|\hat{u}_{\mu\nu}(p)\right|^2,
\]
so that $\|u\|_2=\|u\|_{H^0}$.

\emph{Grade 0.}
Let $0<\eta\le1$, and let $R>D_\beta$ be such that $n:=R/\eta$ is an integer $\ge2$, of either parity.
(The condition $R>D_\beta$ says that $Q_\beta$ lies in the interior of a translate of $[0,R)^3$; it is invariant under $\beta\mapsto\tau_b\beta$, $b\in\mathbb{R}^3$.)
Let $u$ and $a_\eta$ be formed on $\Lambda_{\eta,R}$.
A function on $\mathbb{R}^3$ is evaluated at a point of the torus through its $R\mathbb{Z}^3$-periodisation, according to the rule of Section~I.7.6; because $R>D_\beta$, at most one term of the periodising sum differs from $0$.
\begin{itemize}
\item[(b1)] $a_{\eta,\nu}(c)=\tilde{a}_{\eta,\nu}\left(c+\frac{\eta}{2}e_\nu\right)$ and $\left(f^{(\eta)}\right)_\eta=d_\eta a_\eta$.
The cochain $a_\eta$ lies in the range of $d_\eta^{*}$; equivalently, it is orthogonal to the exact and to the constant one-cochains.
\item[(b2)] $\tilde{a}_\eta$ and $f^{(\eta)}$ are real and smooth, with supports in $\operatorname{supp}\beta+[-\eta,\eta]^3$.
Both end points of every bond of $\operatorname{supp}a_\eta$ lie, modulo $R\mathbb{Z}^3$, in $\operatorname{supp}\beta+\left[-\frac{3\eta}{2},\frac{3\eta}{2}\right]^3$, which is a subset of $Q_\beta$.
\item[(b2$^{\circ}$)] (The form without periodisation; of the data of grade 0 it uses only $0<\eta\le1$.)
Let
\[
u^{\circ}_{\mu\nu}(c):=\left(J_\eta\beta_{\mu\nu}\right)\left(c+\frac{\eta}{2}\left(e_\mu+e_\nu\right)\right),\qquad c\in\eta\mathbb{Z}^3,
\]
a two-cochain of finite support on $\eta\mathbb{Z}^3$, and let $a^{\circ}_\eta:=d_\eta^{*}u^{\circ}$, where $d_\eta^{*}$ is the adjoint of $d_\eta$ for the pairings $\langle\cdot,\cdot\rangle_\eta$ with the sums over $\eta\mathbb{Z}^3$.
Then $a^{\circ}_{\eta,\nu}(c)=\tilde{a}_{\eta,\nu}\left(c+\frac{\eta}{2}e_\nu\right)$ for all $c\in\eta\mathbb{Z}^3$, and both end points of every bond of $\operatorname{supp}a^{\circ}_\eta$ lie in $\operatorname{supp}\beta+\left[-\frac{3\eta}{2},\frac{3\eta}{2}\right]^3$, as points of $\mathbb{R}^3$.
For every $R>0$ such that $R/\eta$ is an integer $\ge2$, the cochain $a_\eta$ on $\Lambda_{\eta,R}$ is the $R\mathbb{Z}^3$-periodisation of $a^{\circ}_\eta$:
\[
a_{\eta,\nu}(c)=\sum_{z\in\mathbb{Z}^3}a^{\circ}_{\eta,\nu}(c+Rz).
\]
\item[(b3)] $\left\|f^{(\eta)}-f\right\|_{H^s}\le\eta^2\cdot\|\beta\|_{H^{s+4}}$ for every real $s\ge0$.
\item[(b4)] $\left|D^i\tilde{a}_\eta\right|_\infty\le2|\beta|_{C^{i+1}}$ and $\left|\tilde{a}_\eta-a\right|_{C^i}\le3\eta^2\cdot|\beta|_{C^{i+3}}$ for every integer $i\ge0$.
\item[(b0)] For every test form $f'\in\mathcal{S}$ with $\operatorname{supp}f'\subset y_0+\left(-\frac{R}{2},\frac{R}{2}\right)^3$, where $y_0\in\mathbb{R}^3$ is arbitrary, one has $\left|f'_\eta\right|_\eta\le\|f'\|_2$.
\end{itemize}

\emph{Grade 1.}
In addition let $v\in\mathbb{R}^3$, and let $\hat{\mathfrak{a}}:=\eta^{\frac{1}{2}}a_\eta$, $\boldsymbol{\sigma}^2:=\left\|d\hat{\mathfrak{a}}\right\|^2$ and $X:=\left\langle d\hat{\mathfrak{a}},F^{\mathrm{u}}\right\rangle$ be formed with $\tau_v\beta$ on $\Lambda_{\eta,R}$.
\begin{itemize}
\item[(b2$'$)] $X=X_{\eta,R}\left(\tau_vf^{(\eta)}\right)$ as functions of $F^{\mathrm{u}}\in C^2(n)$.
\item[(b5)] One has
\[
\boldsymbol{\sigma}\le\left\|f^{(\eta)}\right\|_2\le\|f\|_2+\|\beta\|_{H^4}=:s_f
\]
and
\[
\left|\boldsymbol{\sigma}-\|f\|_2\right|\le\eta^2\cdot\|\beta\|_{H^4}+21\eta\cdot|\beta|_{C^3}\cdot D_\beta^{\frac{3}{2}} .
\]
\item[(b5$'$)] If $v\in\eta\mathbb{Z}^3$, then $\hat{\mathfrak{a}}$ is the translate by $v/\eta$ of the cochain formed with $\beta$, and $\boldsymbol{\sigma}^2$ depends neither on $v$ nor on $R$.
\end{itemize}

\emph{Grade 2} ($\eta=L^{-N}$, $\beta$ fits, $v$ is a placement, $n\in\mathrm{per}\,\mathbb{Z}$ and $n\ge n_*$): Proposition~41.
\end{named}

\begin{named}{Remark}
Proposition~47.1 corresponds to \cite[Theorem~21.1\,(b)]{QEDI}, which is proved in \cite[Section~21.3.1]{QEDI}.
Clause (b0) is the assertion of \cite[Lemma~7.5]{QEDI} on tests supported in a period cube, the one that precedes the bound $p_1(h_\kappa)\le31\lambda_\kappa$; the same assertion is part of \cite[Lemma~12.5\,(c)]{QEDII}.

The comparison is as follows.
The hypotheses are arranged in three grades.
Clauses (b2$'$) and (b5) hold for the same sides $R$ as the last sentence of \cite[Theorem~21.1\,(b2)]{QEDI} and \cite[Theorem~21.1\,(b5)]{QEDI}, namely for every $R$ as in \cite[Theorem~21.1\,(b)]{QEDI}, and they differ from these two assertions in the other data.
In \cite{QEDI} both assertions presuppose integers $N\ge0$ and $T^{*}\ge1$ with $\eta=L^{-N}$, a probe $\beta$ that fits $\left(N,T^{*}\right)$, and a placement $v\in\eta\mathbb{Z}^3$ as in \cite[Section~21.1.6]{QEDI};
here, in grade~1, $\eta$ is any number with $0<\eta\le1$ for which $R/\eta$ is an integer $\ge2$, $v$ is any vector of $\mathbb{R}^3$, and $\beta$ need not fit.
\cite[Theorem~21.1\,(b)]{QEDI} assumes that $Q_\beta$ lies in the interior of a period cube; here this is the condition $R>D_\beta$, which is invariant under translations, so that it holds for $\tau_v\beta$ as soon as it holds for $\beta$.
In (b5) the cube that has to lie in the interior of a period cube is $v+Q_\beta$; the inclusion $v+Q_\beta\subset[\eta,\ell_c-\eta]^3$, which the hypothesis of \cite[Theorem~21.1\,(b5)]{QEDI} that $\beta$ fits provides, is not needed.
The number $n$ may be of either parity, and no fermion enters.
Clause (b2$'$) holds for every real $v$, because $\tau_v\beta$ is a probe and translations commute with $\delta_{\eta,\mu}$ and $J_\eta$; the condition $v\in\eta\mathbb{Z}^3$ enters in (b5$'$) only.
In \cite{QEDI} the probes are also described by the half side $\rho_\beta$, with $D_\beta=2\rho_\beta$; for instance, the condition that $\beta$ fits reads $2\rho_\beta+3\eta\le\ell_c$ there.
\end{named}

\subsection*{41. The probes live in the open cell}

\subsubsection*{Proposition 41}

\begin{named}{Proposition 41}
Let $e\cdot h^{\frac{1}{2}}\le e_1$, $N\ge N_0(e)$ and $\eta=L^{-N}$, and let $\mathrm{hc}$, $m_{\mathrm{p}}$, $\ell_c$ and $\mathcal{Q}_{*}$ be as in Section~I.6.4.
Let $\beta$ be a probe that fits $\left(N,T^{*}(e)\right)$, that is, $D_\beta+3\eta\le\ell_c$, and let $v$ be any placement, that is, $v\in\eta\mathbb{Z}^3$ and $\left|v+y_\beta-\frac{\ell_c}{2}(1,1,1)\right|\le\frac{\eta}{2}$.
Let $\hat{\mathfrak{a}}^{\circ}:=\eta^{\frac{1}{2}}a_\eta$ be the one-cochain on $\mathbb{Z}^3$, in lattice units, that the definitions of Section~I.7.6 give for $\tau_v\beta$ when no periodisation is made: here $u$ is the finitely supported two-cochain on $\eta\mathbb{Z}^3$ with
\[
u_{\mu\nu}(c)=\left(J_\eta\left(\tau_v\beta\right)_{\mu\nu}\right)\left(c+\frac{\eta}{2}\left(e_\mu+e_\nu\right)\right),
\]
and $d_\eta^{*}=\eta^{-1}d^{*}$ is the adjoint on the finitely supported cochains \cite[Definition~2.1]{QEDII}.
For a cubic torus with $n$ sites per side, where $n\in\mathrm{per}\,\mathbb{Z}$ and $n\ge n_*$, the cochain $\hat{\mathfrak{a}}$ is the periodisation of $\hat{\mathfrak{a}}^{\circ}$.
\begin{itemize}
\item[(a)] $v+Q_\beta\subset[\eta,\ell_c-\eta]^3$, and $m_{\mathrm{p}}\ge7$.
The cochain $\hat{\mathfrak{a}}^{\circ}$ has finite support and does not depend on $n$.
Both end points of every bond of $\operatorname{supp}\hat{\mathfrak{a}}^{\circ}$ have their lattice coordinates in $\left[\frac{2}{\eta}-\frac{1}{2},m_{\mathrm{p}}+1-\frac{2}{\eta}\right]^3$, and hence lie in $\left\{2,\dots,m_{\mathrm{p}}-1\right\}^3$.
\item[(b)] Every plaquette of $\operatorname{supp}d\hat{\mathfrak{a}}^{\circ}$ is contained, as a closed set, in the box
\[
\mathrm{Bx}_0:=\frac{\mathrm{hc}}{2}(1,1,1)+[-r_0,r_0]^3,\qquad r_0:=\frac{\mathrm{hc}}{2}-\frac{1}{\eta}\left(2-\frac{3\eta}{2}\right)>0,
\]
and its vertices lie in $\left\{1,\dots,m_{\mathrm{p}}\right\}^3$.
The box $\mathrm{Bx}_0$ lies in the open cell $Q^{\circ}_0=(0,\mathrm{hc})^3$, and its distance from the wall $\left\{y_\mu=j\cdot\mathrm{hc}\right\}$ equals
\[
\frac{2}{\eta}-\frac{3}{2}+\max(-j,j-1)\cdot\mathrm{hc} ;
\]
this is $\frac{2}{\eta}-\frac{3}{2}\ge\frac{1}{2}$ for the six walls of $Q^{\circ}_0$ ($j=0,1$), and it is larger for every other wall.
\item[(c)] Let $\mathfrak{c}\in\mathbb{Z}^3$.
The cochain $\hat{\mathfrak{a}}_{\mathfrak{c}}=\left(\mathsf{g}_{\mathfrak{c}}\right)_{*}\hat{\mathfrak{a}}$ is the periodisation of $\left(\mathsf{g}_{\mathfrak{c}}\right)_{*}\hat{\mathfrak{a}}^{\circ}$, and $\operatorname{supp}d\left(\mathsf{g}_{\mathfrak{c}}\right)_{*}\hat{\mathfrak{a}}^{\circ}$ lies in $\mathsf{g}_{\mathfrak{c}}\mathrm{Bx}_0$, a box in $Q^{\circ}_{\mathfrak{c}}$.
Every plaquette of $\mathbb{Z}^3$ in $\operatorname{supp}d\hat{\mathfrak{a}}_{\mathfrak{c}}$ lies in $\mathsf{g}_{\mathfrak{c}}\mathrm{Bx}_0+nz$ for exactly one $z\in\mathbb{Z}^3$, and $\mathsf{g}_{\mathfrak{c}}\mathrm{Bx}_0+nz=\mathsf{g}_{\mathfrak{c}+(2n/\mathrm{per})z}\mathrm{Bx}_0$ is the box of the same cell of the torus $n$.
The cochain $\hat{\mathfrak{a}}_{\mathfrak{c}}$ depends only on the cell of the torus, that is, on the class of $\mathfrak{c}$ modulo $(2n/\mathrm{per})\mathbb{Z}^3$.
For two cells $\mathfrak{c}$, $\mathfrak{c}'$ of the torus $n$,
\[
\left\langle d\hat{\mathfrak{a}}_{\mathfrak{c}},d\hat{\mathfrak{a}}_{\mathfrak{c}'}\right\rangle=
\begin{cases}
\boldsymbol{\sigma}^2 & \text{if }\mathfrak{c}=\mathfrak{c}',\\
0 & \text{otherwise},
\end{cases}
\]
and $\boldsymbol{\sigma}^2=\left\|d\hat{\mathfrak{a}}^{\circ}\right\|^2$ depends neither on $n$, nor on $\mathfrak{c}$, nor on the placement.
\item[(d)] For a box $\mathrm{Bx}$, that is, a closed cube $y+[-r,r]^3$ with $y\in\mathbb{R}^3$ and $r>0$, let $\mathrm{Ln}(\mathrm{Bx})$ and $\mathrm{Pq}(\mathrm{Bx})$ be the sets of the links and of the plaquettes of $\mathbb{Z}^3$ that are contained in $\mathrm{Bx}$;
they are the sets $B^{(1)}$ and $B^{(2)}$ of \cite[Definition~2.12]{QEDII} for $B=\mathrm{Bx}$.
For $\hat{A}\in\mathbb{R}^{\mathrm{Ln}(\mathrm{Bx}_0)}$ the value $(d\hat{A})(p)$ is defined for every $p\in\mathrm{Pq}(\mathrm{Bx}_0)$, because the four links of such a plaquette lie in $\mathrm{Bx}_0$, and
\[
\left\langle d\hat{\mathfrak{a}}^{\circ},d\hat{A}\right\rangle:=\sum_{p\in\mathrm{Pq}(\mathrm{Bx}_0)}\left(d\hat{\mathfrak{a}}^{\circ}\right)(p)\cdot\left(d\hat{A}\right)(p);
\]
by (b), every plaquette of $\operatorname{supp}d\hat{\mathfrak{a}}^{\circ}$ belongs to $\mathrm{Pq}(\mathrm{Bx}_0)$.
For every polynomial $\mathrm{Pol}$ with complex coefficients the map $\hat{A}\mapsto\mathrm{Pol}\left(\left\langle d\hat{\mathfrak{a}}^{\circ},d\hat{A}\right\rangle\right)$ on $\mathbb{R}^{\mathrm{Ln}(\mathrm{Bx}_0)}$ is a bosonic element of $\mathrm{Ad}(\mathrm{Bx}_0)$, hence of $\mathcal{A}\left(Q^{\circ}_0\right)$, and $n$ does not enter it.
The box $\mathrm{Bx}_0$ embeds in the torus with the period lattice $n\mathbb{Z}^3$, since $2r_0+2<\mathrm{per}\le n$, and the reading of this element on the torus $n$ through a local lift \cite[Definition~2.12]{QEDII} is the function $\mathrm{Pol}(X)$, $X=\left\langle d\hat{\mathfrak{a}},F^{\mathrm{u}}\right\rangle$, of $F^{\mathrm{u}}\in C^2(n)$.
In this sense $X$ and every polynomial in $X$ with complex coefficients are bosonic elements of $\mathrm{Ad}(\mathrm{Bx}_0)$, each being one and the same element for all sides $n$ admitted here.
For $e>0$, real numbers $c_m$ and $m\in\{2,4\}$, the function $O_m:=e^{-m}\cdot H_m(X)-c_m$ is real, and
\[
O_m^{[\mathfrak{c}]}=e^{-m}\cdot H_m\left(X_{\mathfrak{c}}\right)-c_m\qquad\text{for every }\mathfrak{c}\in\mathbb{Z}^3.
\]
\item[(e)] For $\mathfrak{c}\in2\mathbb{Z}^3$ the map $\mathsf{g}_{\mathfrak{c}}$ is the translation by $\mathrm{hc}\cdot\mathfrak{c}\in\mathrm{per}\,\mathbb{Z}^3$.
\end{itemize}
\end{named}

\begin{named}{Remark}
Proposition~41 corresponds to the second half of \cite[Theorem~21.1\,(b2)]{QEDI}, namely the assertions that every bond of the support of $\hat{\mathfrak{a}}$ has both end points in $\left\{1,\dots,m_{\mathrm{p}}-1\right\}^3$, that the support of $d\hat{\mathfrak{a}}$ is contained in $P^o(\mathfrak{c}_0)$, that the cochains $d\hat{\mathfrak{a}}_{\mathfrak{c}}$ of different cells have disjoint supports, and that $\left\|d\hat{\mathfrak{a}}_{\mathfrak{c}}\right\|^2=\boldsymbol{\sigma}^2$, with their proof in \cite[Section~21.3.1]{QEDI}; it also corresponds to a part of \cite[Lemma~20.3\,(d)]{QEDI}.

The mathematical differences are as follows.
The open cell takes the place of the closed one, and the distance $\frac{2}{\eta}-\frac{3}{2}$ of $\mathrm{Bx}_0$ from the walls of $Q^{\circ}_0$ is made explicit.
The condition that $\beta$ fits is the same as in \cite{QEDI}: no additional lattice spacing is needed, because it is $\operatorname{supp}\beta+[-2,2]^3$, and not merely $\operatorname{supp}\beta$, that lies in $Q_\beta$.
As in \cite[Theorem~21.1\,(b2)]{QEDI}, the spacing $\eta=L^{-N}$ and the requirement that the torus carry the cells are hypotheses:
there the assertions in question are made for $\eta=L^{-N}$ and for $R=k\cdot R^{*}$ with an integer $k\ge1$, where $R^{*}=h\cdot L^{T^{*}}$, that is, on the torus $k$ of \cite[Section~21.1.5]{QEDI}.
The difference lies in the torus: here it may be any cubic torus with $n$ sites per side, where $n\in\mathrm{per}\,\mathbb{Z}$ and $n\ge n_*$, of either parity, whereas there $R=k\cdot R^{*}$ with an integer $k\ge1$.
The assertions hold for every placement, since of the placement only the inequality $\left|v+y_\beta-\frac{\ell_c}{2}(1,1,1)\right|\le\frac{\eta}{2}$ enters.
\end{named}

\subsection*{38 (continued). The determinant on strips; the zero-flux sector; the shift identity}

\subsubsection*{Proposition~38.3. The shift identity}

\begin{named}{Proposition 38.3}
Let $\boldsymbol{\mu}$ be real and $q>0$, and let $n=(n_1,n_2,n_3)$ with integers $n_\mu\ge4$, each of either parity.
The pairing $\langle\,\cdot\,,\,\cdot\,\rangle$ is extended to complex cochains as a bilinear form.
For $a\in C^{1}(n)\otimes\mathbb{C}$ and $w\in\mathbb{Z}^3$ let $Z^{\mathrm{x},(w)}_n(a)$ and $\mathbf{Z}^{\mathrm{x}}_n(a)$ denote the expressions of Section~I.7.4.
\begin{itemize}
\item[(a)] The two expressions converge absolutely, locally uniformly in $a$, and they are entire functions of $a\in\mathbb{C}^{3|n|}$.
Moreover $\sum_{w}\left|Z^{\mathrm{x},(w)}_n(a)\right|<\infty$, locally uniformly in $a$.
\item[(a$'$)] (A further factor; the measures $\bar{\nu}^a$.)
Let $\Phi\in\mathfrak{F}(n)$, and let $Z^{\mathrm{x},(w)}_n(a;\Phi)$ and $\mathbf{Z}^{\mathrm{x}}_n(a;\Phi)$ be the expressions of Section~I.7.4 whose integrands carry the further factor $\Phi(F^{\mathrm{u}})$, respectively $\Phi(dA)$.
Both converge absolutely, locally uniformly in $a$, and both are entire functions of $a\in\mathbb{C}^{3|n|}$.
One has
\[
Z^{\mathrm{x},(0)}_n(a;\Phi)=\mathrm{Nc}_n\cdot\mathbf{Z}^{\mathrm{x}}_n(a;\Phi)\quad\text{for every complex }a,
\qquad
Z^{\mathrm{x},(w)}_n(0;\Phi)=Z^{(w)}_n[\Phi].
\]
Let now the torus be cubic with $n\ge4$, and take the objects of Section~I.7.3 for a real pair $(m_{\mathrm{in}},\varepsilon^{\mathrm{v}}_{\mathrm{in}})$, with $\boldsymbol{\mu}=\bar{m}\cdot\eta+m_{\mathrm{in}}$; for complex $a$ the measures $\nu^{\theta,a}$ and $\bar{\nu}^a$ are complex measures.
The integral
\[
\bar{\nu}^a(\Phi):=\int_{[-\pi,\pi]^3}\frac{\mathrm{d}\theta}{(2\pi)^3}\int\mathrm{d}\gamma_N\,\Phi\cdot\mathrm{dens}^{\theta,a}_N
\]
converges absolutely, and
\[
\bar{\nu}^a(\Phi)=\frac{\exp\left(-\varepsilon^{\mathrm{v}}_{\mathrm{in}}\cdot n^3\right)}{\det\left(D_0+\bar{m}\cdot\eta\right)\cdot\mathrm{zV}_n}\cdot\mathbf{Z}^{\mathrm{x}}_n(a;\Phi).
\]
Consequently the map $a\mapsto\bar{\nu}^a(\Phi)$ is entire on $\mathbb{C}^{3n^3}$, and so is $a\mapsto\bar{\nu}^a(\mathrm{Ev})$ for every Borel set $\mathrm{Ev}\subset F_N$.
Furthermore
\[
\left|\bar{\nu}^a(\mathrm{Ev})\right|\le\frac{\exp\left(-\varepsilon^{\mathrm{v}}_{\mathrm{in}}\cdot n^3\right)}{\det\left(D_0+\bar{m}\cdot\eta\right)}\cdot\left(|M_{\mathrm{ch}}|+\frac{1+\rho}{2}\left(\mathrm{sw}-6+6\exp\left(q\cdot|\operatorname{Im}a|_{\infty}\right)\right)\right)^{192n^3},
\]
and $t\mapsto\bar{\nu}^{a(t)}(\mathrm{Ev})$ is entire for every linear map $t\mapsto a(t)$ with values in $C^{1}(n)\otimes\mathbb{C}$; on $n=n_*$ this applies to $a(t)=e^{-1}\sum_{\mathfrak{c}}t_{\mathfrak{c}}\cdot\hat{\mathfrak{a}}_{\mathfrak{c}}$.
\item[(b)] For every $w\in\mathbb{Z}^3$ and every complex $a$,
\[
Z^{(w)}_n\left[\exp\langle da,F^{\mathrm{u}}\rangle\right]=\exp\left(\frac{1}{2}\langle da,da\rangle\right)\cdot Z^{\mathrm{x},(w)}_n(a).
\]
Summation over $w$ gives
\[
Z_n\left[\exp\langle da,F^{\mathrm{u}}\rangle\right]=\exp\left(\frac{1}{2}\langle da,da\rangle\right)\cdot\sum_{w}Z^{\mathrm{x},(w)}_n(a).
\]
\item[(c)] For every complex $a$ one has $Z^{\mathrm{x},(0)}_n(a)=\mathrm{Nc}_n\cdot\mathbf{Z}^{\mathrm{x}}_n(a)$; moreover $\mathbf{Z}^{\mathrm{x}}_n(0)>0$.
\item[(d)] Let $a^{(1)},\dots,a^{(r)}\in C^{1}(n)$ be finitely many real one-cochains, and put $X^{(j)}:=\langle da^{(j)},F^{\mathrm{u}}\rangle$.
For all $z\in\mathbb{C}^r$,
\[
E^{\mathrm{av}}_n\exp\left(\sum_{j}z_j\cdot X^{(j)}-\frac{1}{2}\sum_{i,j}z_i\cdot z_j\cdot\left\langle da^{(i)},da^{(j)}\right\rangle\right)=\frac{\mathbf{Z}^{\mathrm{x}}_n\left(\sum_{j}z_j\cdot a^{(j)}\right)}{\mathbf{Z}^{\mathrm{x}}_n(0)}.
\]
\item[(e)] (Wick powers on the reference torus.)
Consider the setting of Section~I.7.6: a probe that fits, an arbitrary placement, $n=n_*$, the lattice mass $\boldsymbol{\mu}=\boldsymbol{\mu}_N$ or any real $\boldsymbol{\mu}$, and the $N_*$ cochains $\hat{\mathfrak{a}}_{\mathfrak{c}}$.
Let $\bar{Z}^{\mathrm{ext}}$ and $\mathring{\Phi}$ be the displays of Section~I.7.6 formed with this $\boldsymbol{\mu}$, and let $\Theta^{(0)}_m$ be the display formed with this $\mathring{\Phi}$ and with the real number $c_m$ that occurs below.
Then
\[
\bar{Z}^{\mathrm{ext}}(t)=\mathrm{zV}_{n_*}^{-1}\cdot\mathbf{Z}^{\mathrm{x}}_{n_*}\left(e^{-1}\sum_{\mathfrak{c}}t_{\mathfrak{c}}\cdot\hat{\mathfrak{a}}_{\mathfrak{c}}\right),
\qquad
\bar{Z}^{\mathrm{ext}}(0)>0,
\]
and the function
\[
\mathring{\Phi}(t)=\frac{\mathbf{Z}^{\mathrm{x}}_{n_*}\left(e^{-1}\sum_{\mathfrak{c}}t_{\mathfrak{c}}\cdot\hat{\mathfrak{a}}_{\mathfrak{c}}\right)}{\mathbf{Z}^{\mathrm{x}}_{n_*}(0)}
\]
is entire in $t\in\mathbb{C}^{N_*}$.
For $z\in\mathbb{C}^{N_*}$,
\[
E^{\mathrm{av}}_{n_*}\exp\sum_{\mathfrak{c}}\left(z_{\mathfrak{c}}\cdot X_{\mathfrak{c}}-\frac{1}{2}z_{\mathfrak{c}}^2\cdot\boldsymbol{\sigma}^2\right)=\mathring{\Phi}(e\cdot z).
\]
Finally, for $m\in\{2,4\}$ and every real $c_m$,
\[
E^{\mathrm{av}}_{n_*}\prod_{\mathfrak{c}\in\mathrm{Cel}}\left(e^{-m}\cdot H_m(X_{\mathfrak{c}})-c_m\right)=\left(\prod_{\mathfrak{c}\in\mathrm{Cel}}\left(\partial_{t_{\mathfrak{c}}}^m-c_m\right)\mathring{\Phi}\right)(0),
\]
so that
\[
\left|E^{\mathrm{av}}_{n_*}\prod_{\mathfrak{c}\in\mathrm{Cel}}\left(e^{-m}\cdot H_m(X_{\mathfrak{c}})-c_m\right)\right|=\left(\Theta^{(0)}_m\right)^{N_*}.
\]
\end{itemize}
\end{named}

\begin{named}{Remark}
Proposition~38.3 corresponds to the assertion of \cite[Section~11.4.1]{QEDI} that the two functionals of an external field defined there are entire, to \cite[Lemma~11.7]{QEDI}, and to Step~3 of the proof of \cite[Theorem~21.2]{QEDI} in \cite[Section~21.8.2]{QEDI}, namely to the identity for the generating function of the cell variables under the averaged law on the reference torus $n_*$ and to the identity for the product of the centred Hermite polynomials over all cells that is derived from it.
The statements differ in the following points.
\begin{itemize}
\item[(1)] In (a) and (a$'$) the absolute convergence for complex $a$ and the analyticity in $a$ are part of the statement.
The convergence is locally uniform in $a$: for every bounded set of values of $a$, the moduli of the integrands and those of the terms of the series over $w$ admit an integrable, respectively summable, majorant that is the same for all $a$ in that set.
Such majorants are furnished by the two bounds (38.6) and (38.7) of Chapter~38, Lemma~2.A.
\item[(2)] For complex $a$, clause (c) is the case $\mathrm{wq}(A)=\det\left(D_{A+a}+\boldsymbol{\mu}\right)$ of Proposition~38.2\,(a), which admits bounded continuous complex weights with the two invariances.
\item[(3)] Clause (d) holds for arbitrary real one-cochains $a^{(1)},\dots,a^{(r)}$: the Gram matrix $\left(\langle da^{(i)},da^{(j)}\rangle\right)_{i,j}$ enters, and the supports of the $da^{(j)}$ need not be disjoint.
In Step~3 of the proof of \cite[Theorem~21.2]{QEDI} the cochains $d\hat{\mathfrak{a}}_{\mathfrak{c}}$ have pairwise disjoint supports, so that only the diagonal of this matrix occurs there.
\item[(4)] Clause (e) gives the modulus of $E^{\mathrm{av}}_{n_*}\prod_{\mathfrak{c}}\left(e^{-m}\cdot H_m(X_{\mathfrak{c}})-c_m\right)$ and no sign.
In \cite[Section~21.8.2]{QEDI} the corresponding number is shown to be non-negative by the chessboard estimate \cite[Proposition~21.22]{QEDI} for the zero-flux functional formed with the Wilson determinant.
The functional $E^{\mathrm{av}}_{n_*}$ of (e) is the zero-flux sector of Section~I.7.3, formed with the chamber determinant;
for this law no reflection positivity is available on tori with an odd number of sites per side, such as $n_*$ (Section~I.7.3).
\item[(5)] Clause (e) contains the facts that $\bar{Z}^{\mathrm{ext}}(0)>0$ and that $\mathring{\Phi}$ is entire, on which the definition of $\Theta^{(0)}_m$ rests (Section~I.7.6; \cite[Section~21.1.7]{QEDI}).
\item[(6)] Every real lattice mass $\boldsymbol{\mu}$ is admitted, and the sides may have either parity.
\end{itemize}
\end{named}

\subsection*{47 (continued). Probes and cells on the even tori}

\subsubsection*{Proposition~47.2. The cells on the even tori; the translation form}

\begin{named}{Proposition 47.2}
Let $e\cdot h^{\frac{1}{2}}\le e_1$, $N\ge N_0(e)$ and $T\ge T^{*}(e)$; put $n:=8h\cdot L^{T+N}$ and $R:=n\cdot\eta=8h\cdot L^{T}$, let $\boldsymbol{\mu}$ be real, and let $q=e\cdot\eta^{\frac{1}{2}}$.
In (b)--(e) let $\beta$ be a probe that fits $(N,T^{*}(e))$, let $v$ be a placement, and let $f=dd^{*}\beta$.
For $y\in(R^{*}(e)/L)\,\mathbb{Z}^3=\mathrm{per}\cdot\eta\,\mathbb{Z}^3$ put
\[
\mathfrak{c}(y):=\frac{2y}{\mathrm{per}\cdot\eta}\in2\mathbb{Z}^3.
\]
\begin{itemize}
\item[(a)] One has $n=\mathrm{kt}\cdot n_*$ with $\mathrm{kt}=8L^{T-T^{*}(e)}$; the number $\mathrm{per}$ divides $n$, and $n/\mathrm{per}=\mathrm{kt}\cdot L$.
The torus $C(n,n,n)$ carries the cell structure $\mathcal{Q}_{*}$; it has $2\mathrm{kt}\cdot L$ cells per side, hence $(2\mathrm{kt}\cdot L)^3$ cells, and $n\ge4$.
Moreover $\mathrm{kt}\ge L$ if and only if $T\ge T^{*}(e)+1$.
The walls $\{y_\mu=j\cdot\mathrm{hc}\}$ are site planes for even $j$ and link planes for odd $j$.
The two mirrors $\{y_\mu=0\}$ and $\{y_\mu=n/2\}$, which cut the torus into two halves, are both site planes and both walls, because $n/2=\mathrm{kt}\cdot L\cdot\mathrm{hc}$ and $\mathrm{kt}\cdot L$ is even.
\item[(b)] Let $y\in(R^{*}(e)/L)\,\mathbb{Z}^3$.
Then $\mathsf{g}_{\mathfrak{c}(y)}$ is the translation of $\mathbb{Z}^3$ by the vector $y/\eta\in\mathrm{per}\cdot\mathbb{Z}^3$.
For every function $O\in\mathcal{A}(Q^{\circ}_0)$ the mirror copy is $O^{[\mathfrak{c}(y)]}(u)=O\left(\mathsf{g}_{\mathfrak{c}(y)}^{*}u\right)$, without complex conjugation.
Furthermore
\[
X_{\mathfrak{c}(y)}=X_{\eta,R}\left(\tau_{v+y}f^{(\eta)}\right).
\]
\item[(c)] For $y,y'\in(R^{*}(e)/L)\,\mathbb{Z}^3$ the cells $\mathfrak{c}(y)$ and $\mathfrak{c}(y')$ of the torus coincide if and only if $y-y'\in R\,\mathbb{Z}^3$.
(For a finite set $\mathrm{Ypt}$, the inequality $R>2\max_{y\in\mathrm{Ypt}}|y|_{\infty}$ is a sufficient condition for the cells $\mathfrak{c}(y)$, $y\in\mathrm{Ypt}$, to be pairwise distinct.)
\item[(d)] For every finite subset $\mathrm{Ypt}$ of $(R^{*}(e)/L)\,\mathbb{Z}^3$ and every polynomial $\mathrm{Pol}$ in $\#|\mathrm{Ypt}|$ variables,
\[
E^{\mathrm{af}}_n\,\mathrm{Pol}\left(\left(X_{\eta,R}\left(\tau_{v+y}f^{(\eta)}\right)\right)_{y\in\mathrm{Ypt}}\right)=E^{\mathrm{af}}_n\,\mathrm{Pol}\left(\left(X_{\eta,R}\left(\tau_{y}f^{(\eta)}\right)\right)_{y\in\mathrm{Ypt}}\right).
\]
\item[(e)] The two-form $g_N:=f^{(\eta)}$ is real, lies in $C_c^{\infty}(\mathbb{R}^3;\Lambda^2)$, and is supported in $Q_\beta$, a subset of the fixed cube $[-r_g,r_g]^3$, where $r_g:=|y_\beta|_{\infty}+D_\beta/2$.
One has $\|g_N-f\|_2\le\eta^2\cdot\|\beta\|_{H^4}$, and $\boldsymbol{\sigma}^2\to\|f\|_2^2$ as $N\to\infty$, with the rate of Proposition~47.1\,(b5).
The number $(R^{*}(e)/L)/\eta=\mathrm{per}$ is an integer.
\end{itemize}
\end{named}

\begin{named}{Remark}
Proposition~47.2 corresponds to the cells of \cite[Section~21.1.5]{QEDI} on the torus $\mathrm{kt}$ for even $\mathrm{kt}$, and to the part of Step~6 of the proof in \cite[Section~21.10.1]{QEDI} that precedes the bound on the expectation of the product over $i=1,\dots,J$.
The statements differ in the following points.
\begin{itemize}
\item[(1)] The cells of \cite[Section~21.1.5]{QEDI} are formed here on tori with an even number of sites per side; nothing in the construction of Section~I.6.4 depends on the parity of $\mathrm{kt}$.
In particular the cells exist on the tori $n=8h\cdot L^{T+N}$, because $\mathrm{per}$ divides $n$ by (a).
\item[(2)] Clause (c) is an exact criterion for the coincidence of two cells; in Step~6 of the proof in \cite[Section~21.10.1]{QEDI} the cells are distinct under the sufficient condition $R>2\max_i|y_i|$.
\item[(3)] Clause (d) is the invariance under the lattice translation by $v$ that Step~6 of the proof in \cite[Section~21.10.1]{QEDI} invokes for the laws of \cite{QEDI}.
\end{itemize}
\end{named}

\subsection*{31. The zero-flux sector; the complement of the small-field set}

\subsubsection*{Proposition~31.1 (the twist-averaged measure is the zero-flux sector; sign, mass, Gaussian domination)}

\begin{named}{Proposition 31.1}
\emph{Data.}
All objects are formed on the lattice $\mathbb{Z}^3$ of spacing one (Section~I.2.5).
Let $\bar{m}\in(0,1]$, $e>0$, let $N\ge0$ be an integer, $\eta=L^{-N}$, $q=e\cdot\eta^{\frac{1}{2}}$ and $r_q=2\pi/q$.
Let $n\ge4$ be an integer of either parity; the torus is the cubic torus $\mathbb{Z}^3/n\mathbb{Z}^3$, and the spinors lie in $\mathfrak{S}^{\mathrm{ch}}(1;n)$.
Let $(m_{\mathrm{in}},\varepsilon^{\mathrm{v}}_{\mathrm{in}})$ be an arbitrary real pair, and $\boldsymbol{\mu}:=\bar{m}\cdot\eta+m_{\mathrm{in}}$.

\emph{Definitions.}
The objects $C^{j}(n)$, $V_{\mathrm{g}}$, $V_{\mathrm{h}}$, $V_{\mathrm{t}}$, $\sigma$, $F_N:=dC^{1}(n)$ with its standard Gaussian measure $\gamma_N$, $A(F^{\mathrm{u}})\in V_{\mathrm{t}}$, $H_\theta$, $\mathrm{dens}^{\theta,a}_N$, $\nu^{\theta,a}$, $\nu^{\theta}$, $\bar{\nu}^a$, $\bar{\nu}$ and $\mathrm{zr}_n(q,\mathbf{m})$ are given by the displays of Sections~I.7.1 and~I.7.3 (for $H_\theta$, see Section~I.6.6).
The functional $Z^{(0)}_n[\,\cdot\,]$, the expectation $E^{\mathrm{av}}_n$ and the number $\mathrm{Nc}_n$ are those of Section~I.7.3, formed with $\det\left(D_A+\boldsymbol{\mu}\right)$.
Let $\mathrm{PrF}$ be the orthogonal projection of $C^{2}(n)$ onto $F_N$.
The number
\[
\mathrm{zV}_n:=\int_{V_{\mathrm{t}}}\mathrm{d}\sigma(A)\,\exp\left(-\frac{1}{2}\|dA\|^2\right)
\]
lies in $(0,\infty)$, and
\[
\int_{V_{\mathrm{t}}}\mathrm{d}\sigma(A)\,G(dA)\cdot\exp\left(-\frac{1}{2}\|dA\|^2\right)=\mathrm{zV}_n\cdot\int\mathrm{d}\gamma_N\,G
\]
for every Borel function $G$ on $F_N$ that is non-negative or satisfies $|G(F^{\mathrm{u}})|\le\mathrm{cf}\cdot\exp\left(\mathrm{cf}\cdot\|F^{\mathrm{u}}\|\right)$ for some number $\mathrm{cf}$.

\emph{Statement.}
\begin{itemize}
\item[(a)] (Sign.)
Let $a$ be a real one-cochain on $\mathbb{Z}^3/n\mathbb{Z}^3$ and $z$ a real number.
Then $\operatorname{Re}\left(D_a+z\right)\ge z$ with respect to the scalar product $(\,\cdot\,,\,\cdot\,)^{\mathrm{ch}}$, and $\det\left(D_a+z\right)\ge0$; if $z>0$, then $\det\left(D_0+z\right)>0$.
Consequently, for real $a$ and all $\theta\in\mathbb{R}^3$, $\mathrm{dens}^{\theta,a}_N$ is a bounded, continuous, non-negative function of $(\theta,F^{\mathrm{u}})$, and $\nu^{\theta,a}$ is a finite positive measure on $F_N$, that is, a countably additive measure with values in $[0,\infty)$.
\item[(b)] (Identification.)
For every Borel function $G$ on $C^{2}(n)$ with $|G(F^{\mathrm{u}})|\le\mathrm{cf}\cdot\exp\left(\mathrm{cf}\cdot\|F^{\mathrm{u}}\|\right)$ for some $\mathrm{cf}$,
\[
\bar{\nu}(G)=\frac{\exp\left(-\varepsilon^{\mathrm{v}}_{\mathrm{in}}\cdot n^3\right)}{\det\left(D_0+\bar{m}\cdot\eta\right)\cdot\mathrm{Nc}_n\cdot\mathrm{zV}_n}\cdot Z^{(0)}_n[G].
\]
In particular
\[
\mathrm{zr}_n(q,\boldsymbol{\mu})=\frac{Z^{(0)}_n[1]}{\mathrm{Nc}_n\cdot\mathrm{zV}_n}=\exp\left(\varepsilon^{\mathrm{v}}_{\mathrm{in}}\cdot n^3\right)\cdot\det\left(D_0+\bar{m}\cdot\eta\right)\cdot\bar{\nu}(1).
\]
\item[(c)] (Mass.)
One has $0<\bar{\nu}(1)<\infty$.
Hence $\bar{\mu}:=\bar{\nu}/\bar{\nu}(1)$ is a probability measure on $F_N$.
It is the law of $F^{\mathrm{u}}$ under $E^{\mathrm{av}}_n=E^{\mathrm{af}}_n\left(\,\cdot\,|\,\mathrm{flux}=0\right)$ in the following sense: this law, a measure on $C^{2}(n)$, is the image of $\bar{\mu}$ under the inclusion of $F_N$ in $C^{2}(n)$.
The measure $\bar{\mu}$ depends on $(n,q,\boldsymbol{\mu})$ only, and not on $\varepsilon^{\mathrm{v}}_{\mathrm{in}}$.
\item[(d)] (Gaussian domination with the constant one.)
For every real two-cochain $\mathrm{gc}$ on $\mathbb{Z}^3/n\mathbb{Z}^3$,
\[
\bar{\mu}\left(\exp\langle\mathrm{gc},F^{\mathrm{u}}\rangle\right)\le\exp\left(\frac{1}{2}\|\mathrm{PrF}\,\mathrm{gc}\|^2\right).
\]
For $0\ne\phi\in F_N$ and $y\ge0$,
\[
\bar{\mu}\left\{\langle\phi,F^{\mathrm{u}}\rangle\ge y\right\}\le\exp\left(-\frac{y^2}{2\|\phi\|^2}\right).
\]
For $\phi\in F_N$ with $\|\phi\|\le1$ and $y\ge0$,
\[
\bar{\mu}\left\{\left|\langle\phi,F^{\mathrm{u}}\rangle\right|>y\right\}\le2\exp\left(-\frac{y^2}{2}\right).
\]
\end{itemize}
No constant occurs in (a)--(d), and no condition is imposed on the sign or on the size of $\boldsymbol{\mu}$ (I.K10).
The product $\mathrm{Nc}_n\cdot\mathrm{zV}_n$ is the number $Z^{\mathrm{ph}}_n$ of Proposition~38.2\,(c) and of \cite[Section~21.8.1]{QEDI}.
\end{named}

\begin{named}{Remark}
Proposition~31.1 corresponds to the following statements of \cite{QEDI}:
\cite[Lemma~20.4(a)]{QEDI} and \cite[Proposition~20.5]{QEDI};
the first three assertions of \cite[Lemma~19.23]{QEDI} and the two parts \emph{Positivity} and \emph{The averaged law} of its proof;
the last sentence of \cite[Theorem~6.1]{QEDI}, the Gaussian bound of \cite[Theorem~6.3]{QEDI} together with the positivity of the normalising integral asserted there, and \cite[Corollary~6.4(iii)]{QEDI}, in the form in which these are applied on the torus in the proof of \cite[Lemma~19.23]{QEDI};
the sentence of \cite[Section~5.2.1]{QEDI} which identifies the normalised zero-flux part of the sum with the harmonic-averaged law;
\cite[Section~7.1.2]{QEDI}, to which (a)--(c) correspond together with Sections~I.7.2, I.7.3 and Propositions~50\,(a), (e) and~38.1\,(d);
and \cite[Proposition~7.2]{QEDI} with the last sentence of \cite[Theorem~7.11]{QEDI}, to which (d) corresponds.
The statements differ in the following points.
\begin{itemize}
\item[(1)] The hypotheses $n\ge2$, $q\ne0$ and lattice mass $>-1$ of \cite[Section~6]{QEDI} stand here as $n\ge4$, $q>0$ and $\boldsymbol{\mu}$ real.
\cite[Lemma~19.23]{QEDI} assumes a pair that is controlled on $R$, which gives $|m_{\mathrm{in}}|\le w^{\mathrm{m}}_0<1$; here every real pair is admitted.
\item[(2)] The side $n$ may have either parity.
\item[(3)] The sign of the determinant in (a), the positivity of the mass in (c) and the Gaussian domination in (d) are not obtained through the slice factorisation of the Wilson determinant \cite[Lemma~6.8]{QEDI}, the Proca parent \cite[Section~6.1.5]{QEDI} and the limit of vanishing Proca mass \cite[Theorem~6.3]{QEDI}.
They follow from statements of \cite{QEDII} on the chamber operator:
the sign from \cite[Proposition~3.18]{QEDII};
the positivity of $Z^{(0)}_n[1]$ from \cite[Lemma~13.12(b)]{QEDII} together with \cite[Theorem~4.26]{QEDII}, or, when $\rho^2\ne1$, from \cite[Theorem~4.13(d)]{QEDII}, by which the term of the vanishing integer two-cochain is positive;
the Gaussian domination from \cite[Proposition~13.11]{QEDII}, which gives the bound with the constant one on every torus $C(n_1,n_2,n_3)$ with $n_\mu\ge4$, for $-1\le\rho\le1$ and every real value of the mass parameter.
\item[(4)] \cite[Definition~13.7]{QEDII} defines the averaged law as the law conditioned on the flux $0$, and does not identify it with the average over the twist $\theta$.
Here this identification is clause (b), with the constant
\[
\frac{\exp\left(-\varepsilon^{\mathrm{v}}_{\mathrm{in}}\cdot n^3\right)}{\det\left(D_0+\bar{m}\cdot\eta\right)\cdot\mathrm{Nc}_n\cdot\mathrm{zV}_n};
\]
for the Wilson determinant it is the last sentence of the paragraph on the all-flux law in \cite[Section~5.2.1]{QEDI}.
\item[(5)] That $\bar{\nu}(1)\ne0$ for an arbitrary real pair is clause (c).
\item[(6)] The determinant is that of a linear map of a space of dimension $192n^3$ (Proposition~46\,(e)); the Wilson determinant of \cite{QEDI} is that of a linear map of a space of dimension $4n^3$.
\end{itemize}
\end{named}

\subsubsection*{Proposition~31.2 (the complement of the small-field set)}

\begin{named}{Proposition 31.2}
Let (FL) hold, and let $(m_{\mathrm{in}},\varepsilon^{\mathrm{v}}_{\mathrm{in}})$ be an arbitrary real pair.
Let $n=R/\eta=\ell_0$, so that $n\ge L\cdot M\cdot r_0\ge4$ and $n$ is odd, and let $\bar{\mu}$ be the measure of Proposition~31.1.
Then
\begin{equation}
\bar{\mu}\left(F_N\setminus\mathfrak{G}\right)\le\varepsilon_3,
\qquad
\mathrm{d}\mathsf{m}_{*}\left(\chi_{*}=0\right)\le\varepsilon_4.
\tag{9.2}
\end{equation}
Constants: $C_0$ of the line (C6), of which only the inequality $C_0\ge1$ enters, and the integers $p_0$, $p$ of the line (C7).
Conditions on the constants: the first two clauses of (T1) (Proposition~14.2\,(c)).
\end{named}

\begin{named}{Remark}
Proposition~31.2 corresponds to the two bounds displayed in \cite[Lemma~19.23]{QEDI}, with the two parts of its proof that derive them, and to \cite[Theorem~5.85(b)]{QEDI}.
Here the real pair is arbitrary, whereas \cite[Lemma~19.23]{QEDI} assumes that it is controlled on $R$; the third clause $w^{\mathrm{m}}_0<1$ of (T1) is not required.
The numbers $4n_j$, $6\ell_*^3$, $28$ and $L^{-4i}$ that occur in the proof of \cite[Lemma~19.23]{QEDI} are the same here, since they come from the photon field alone.
\end{named}

\subsection*{32. Two-sided bound on the partition function}

\subsubsection*{Proposition~32.1 (terminal estimate)}

\begin{named}{Proposition 32.1}
\emph{Data.}
Assume (FL), and let $(m_{\mathrm{in}},\varepsilon^{\mathrm{v}}_{\mathrm{in}})$ be a real pair that is controlled on $R$.
Let $Y_*$ and $D_*(\,\cdot\,)$ be as in Proposition~29, and put $c_\theta:=c^{(k_*)}_\theta$.
Let $\mathrm{Rtw}_\theta$ be the number of Proposition~30.3 for
\[
(\ell,\mathfrak{N},\mu,\mathsf{b})=\left(\ell_*,L^{k_*},\mu_*,b_{k_*}\right);
\]
these values are admissible there, since $k_*\ge1$, so that $\mathfrak{N}\ge3$ is odd, and since $b_-\le b_{k_*}\le b_+$ and $0<\mu_*\le1$.
By Proposition~30.3\,(f),
\[
\mathrm{Rtw}_\theta=\frac{\det D_*(c_\theta)}{\det D_*(0)}>0.
\]

\emph{Statement.}
Let $\theta\in\mathbb{R}^3$, and let $A$ be a real bond function on $T^0_{k_*}$ with mean zero and $|dA|_{\infty}\le p_*$.
Then $A+c_\theta$ lies in $\mathrm{Dom}_{k_*}$, and the number
\[
\mathrm{vrem}(\theta,A):=\frac{Y_*\left(A+c_\theta\right)\cdot\exp\left(\varepsilon^{\mathrm{v},t}_{*}\cdot\ell_*^3\right)}{\mathrm{Rtw}_\theta}-1
\]
satisfies
\begin{equation}
\left|\mathrm{vrem}(\theta,A)\right|\le\varepsilon_1=\frac{1}{2}\,e_{*}^{\frac{1}{4}-9\varepsilon}.
\tag{9.1}
\end{equation}
On the set of these pairs $(\theta,A)$ the function $\mathrm{vrem}$ is continuous.

Constants: $C_0$ and $c_{\kappa}$ of the line (C6).
Conditions on the constants (Proposition~14.2): (T2) and (T3); in addition, as hypotheses of \cite[Proposition~5.38(i)]{QEDI}, \cite[Lemma~5.17]{QEDI}, \cite[Proposition~5.32]{QEDI} and (Sc4) (Section~G4.0.2), the first clause of (A1), (A2), the first and the fourth clause of (A8), the third and the fourth clause of (A10), $L\ge5$, $L\ge33$ (a consequence of (L1)), $C_{P}\le C_{w}$, $\gamma_{\mathrm{w}}\le\gamma_{P}$ and $0<\alpha<1$.
\end{named}

\begin{named}{Remark}
Proposition~32.1 corresponds to \cite[Proposition~19.22]{QEDI}, read with the setting at the head of \cite[Section~19.9]{QEDI}, and to \cite[Theorem~5.85(a)]{QEDI}; there the numbers $\mathrm{vrem}(\theta,A)$ and $\mathrm{Rtw}_\theta$ are denoted $\vartheta$ and $R_\theta$.
The statements differ in the following points.
\begin{itemize}
\item[(1)] Here $\mathrm{vrem}$ is defined as a function of $(\theta,A)$, and its continuity is part of the statement; the proof of \cite[Theorem~19.24]{QEDI} integrates the corresponding quantity $\vartheta(A)$ with respect to $\mathrm{d}\mathsf{m}_{*}(A)$, at fixed $\theta$, which presupposes its measurability.
\item[(2)] The constant $C_0$ is formed with the constants $C^{F}$ and $\gamma^{F}$ of the chamber operator (Section~I.8.2, lines (C4r) and (C6)).
\item[(3)] All numerical constants of \cite[Proposition~19.22]{QEDI} are the same, because the terminal integral extends over four components per site of $T^0_{k_*}$, as in \cite{QEDI}.
\end{itemize}
\end{named}

\subsubsection*{Proposition 32.2. The twist-averaged partition function on one torus}

\begin{named}{Proposition 32.2}
Let the flow data (FL) of Section~G4.0.1 be given, and let $\left(m_{\mathrm{in}},\varepsilon^{\mathrm{v}}_{\mathrm{in}}\right)$ be a real pair that is controlled on $R$ in the sense of Section~G4.0.5.
The measures $\nu^{\theta}$ and $\bar{\nu}$ of Section~I.7.3 are formed with this pair, $\mathfrak{G}$ is the small-field set of Section~G4.0.6, and $\varepsilon^{\mathrm{v},t}_{*}$ is the vacuum energy of the trajectory of the pair on $R$ at the stop level $k_{*}=k_{*}(R)$.
Put $n=R/\eta$ and $\mathrm{zr}_n:=\mathrm{zr}_n\left(q,\bar{m}\cdot\eta+m_{\mathrm{in}}\right)$, let $\mathrm{zr}^{\mathrm{free}}_n$ be the number of Section~I.7.3, and let $\mathrm{Rtw}_{\theta}$ be the number of Proposition~30.3 for the data $\left(\ell,\mathfrak{N},\mu,\mathsf{b}\right)=\left(\ell_*,L^{k_{*}},\mu_*,b_{k_{*}}\right)$, as in Proposition~32.1.
The numbers $\varepsilon_1,\dots,\varepsilon_4$ and the function $\varepsilon_5$ are those of Section~G4.0.3.
Then, for every $\theta\in\mathbb{R}^3$,
\begin{equation}
\begin{split}
\exp\left(-\varepsilon^{\mathrm{v},t}_{*}\cdot\ell_*^3\right)\cdot\mathrm{Rtw}_{\theta}\cdot\left(1-\varepsilon_1\right)\cdot\left(1-\varepsilon_4\right)
&\le\nu^{\theta}(\mathfrak{G})\\
&\le\exp\left(-\varepsilon^{\mathrm{v},t}_{*}\cdot\ell_*^3\right)\cdot\mathrm{Rtw}_{\theta}\cdot\left(1+\varepsilon_1\right).
\end{split}
\tag{7.1}
\end{equation}
Moreover,
\begin{equation}
\bar{\nu}(\mathfrak{G})\le\bar{\nu}(1)\le\frac{\bar{\nu}(\mathfrak{G})}{1-\varepsilon_3},
\tag{7.2}
\end{equation}
\begin{equation}
\begin{split}
\left|\log\mathrm{zr}_n-\log\mathrm{zr}^{\mathrm{free}}_n-\varepsilon^{\mathrm{v}}_{\mathrm{in}}\cdot n^3+\varepsilon^{\mathrm{v},t}_{*}\cdot\ell_*^3\right|
&\le2\,\varepsilon_1+\varepsilon_2+2\,\varepsilon_3+2\,\varepsilon_4\\
&\le2\,e_{*}^{\frac{1}{4}-9\varepsilon},
\end{split}
\tag{7.3}
\end{equation}
and
\begin{equation}
\left|\log\mathrm{zr}_n-\left(p^{0}_{\mathrm{ch}}(\bar{m}\cdot\eta)+\varepsilon^{\mathrm{v}}_{\mathrm{in}}\right)\cdot n^3+\varepsilon^{\mathrm{v},t}_{*}\cdot\ell_*^3\right|\le\varepsilon_5(\bar{m}\cdot R)+2\,e_{*}^{\frac{1}{4}-9\varepsilon}.
\tag{7.4}
\end{equation}
Conditions on the constants (Proposition~14.2): the third clause of (T2), and (T4), (T5), all three of Proposition~14.2\,(c) and of Section~G4.0.3;
of the condition (S1) of Proposition~14.2\,(b) only $\varepsilon>0$ and $\frac{1}{4}-9\varepsilon>0$ are used.
\end{named}

\begin{named}{Remark}
Proposition~32.2 corresponds to \cite[Theorem~19.24]{QEDI} and to \cite[Theorem~5.85(c)]{QEDI}; the estimates (7.1)--(7.4) stand in the place of the four estimates there.
The function $p^0$ of \cite[Theorem~5.85]{QEDI} is replaced by the function $p^{0}_{\mathrm{ch}}$ of Proposition~30.1, and the function $x\mapsto312\cdot\log\left(\left(1-\exp\left(-\frac{x}{4}\right)\right)^{-1}\right)$ by the function $\varepsilon_5(x)=C_{6}\cdot\log\left(\left(1-\exp\left(-\kappa_{(0)}\cdot x\right)\right)^{-1}\right)$ of Section~G4.0.3, formed with the constant $\kappa_{(0)}=\min\left\{\frac{\theta_{0}}{2},\frac{c_{\mathrm{el}}\cdot\theta_{0}}{12K_{\mathrm{Str}}}\right\}$ of line (C0r);
the number $\varepsilon_2=C_{\Sigma}\cdot\ell_*^{-1}\cdot\exp\left(-\gamma_{\Sigma}\cdot\ell_*\right)$ is formed with the constants $\gamma_{\Sigma}$ and $C_{\Sigma}$ of Proposition~30.3\,(e).
The numbers $\mathrm{zr}_n$, $\mathrm{zr}^{\mathrm{free}}_n$, $\mathrm{Rtw}_{\theta}$ stand in the place of $\rho_n$, $\rho^{\mathrm{free}}_n$, $R_{\theta}$ of \cite{QEDI}.
The function $A\mapsto\chi_{*}(A)\cdot\left(1+\mathrm{vrem}(\theta,A)\right)$, with $\mathrm{vrem}$ of Proposition~32.1, which is integrated against $\mathrm{d}\mathsf{m}_{*}$ in the twist average, is bounded and Borel on ${V_{\mathrm{t}}}\left(T^0_{k_{*}}\right)$, by Proposition~32.1 and Proposition~29\,(c).
\end{named}

\subsubsection*{Proposition 32.3. All tori of the scheme, one $p_F$, one mass, uniformly in $N$; the reference datum (R2)}

\begin{named}{Proposition 32.3}
Let $\bar{m}\in(0,1]$, and let $e>0$ satisfy $e_K=e\cdot h^{\frac{1}{2}}\le e_1$.
Let $R\in\mathrm{Tor}(e)$, let $N\ge N_{\min}(R,e)$ be an integer, and put $n=R/\eta$.
Let $P_N(e)=\left(m_{0,N}(e),\varepsilon^{\mathrm{v}}_{0,N}(e)\right)$ be the pair of Proposition~28.3, let
\[
\mathrm{zr}_n:=\mathrm{zr}_n\left(e\cdot\eta^{\frac{1}{2}},\bar{m}\cdot\eta+m_{0,N}(e)\right)
\]
be the value of the function $\mathrm{zr}_n(\cdot,\cdot)$ of Section~I.7.3, let $\mathrm{zr}^{\mathrm{free}}_n$ be the number of Section~I.7.3, and put
\[
p_F:=p^{0}_{\mathrm{ch}}(\bar{m}\cdot\eta)+\varepsilon^{\mathrm{v}}_{0,N}(e);
\]
this is a real number, which depends on $(N,e,\bar{m})$ and on the scheme, but not on $R$.
Then
\begin{equation}
\left|\log\mathrm{zr}_n-\log\mathrm{zr}^{\mathrm{free}}_n-\varepsilon^{\mathrm{v}}_{0,N}(e)\cdot n^3\right|\le3\,e_{*}^{\frac{1}{4}-9\varepsilon},
\tag{10.1}
\end{equation}
\begin{equation}
\left|\log\mathrm{zr}_n-p_F\cdot n^3\right|\le\varepsilon_P(R):=\varepsilon_5(\bar{m}\cdot R)+3\,e_{*}^{\frac{1}{4}-9\varepsilon}\le\varepsilon_5\left(\frac{1}{L}\right)+3\,e_K^{\frac{1}{4}-9\varepsilon};
\tag{10.2}
\end{equation}
here $e_{*}$ is the charge at the stop level of the torus $R$, so that $\varepsilon_P(R)$ is defined for every $R\in\mathrm{Tor}(e)$ with $N\ge N_{\min}(R,e)$.
Let moreover $N\ge N_0(e)=N_{\min}\left(R^{*}(e),e\right)$, put $n_*=R^{*}(e)/\eta$, and let $D^{\mathrm{mat}}=L^3\cdot\log\mathrm{zr}_{n_*}-\log\mathrm{zr}_{L\cdot n_*}$ be the number of Section~I.7.6.
Then
\begin{equation}
\begin{split}
\left|D^{\mathrm{mat}}\right|&\le L^3\cdot\varepsilon_P\left(R^{*}\right)+\varepsilon_P\left(L\cdot R^{*}\right)\\
&\le\left(L^3+1\right)\cdot\left(\varepsilon_5\left(\bar{m}\cdot R^{*}\right)+3\,e_K^{\frac{1}{4}-9\varepsilon}\right)\\
&\le C_{\mathrm{mat}}=\left(L^3+1\right)\cdot\left(\varepsilon_5\left(\frac{1}{L}\right)+1\right).
\end{split}
\tag{10.3}
\end{equation}
The right sides of (10.2) and (10.3) do not depend on $N$.
Their last members depend only on $(L,\rho)$, on the exponent $\varepsilon$ and on $e_K$: the exponent $\varepsilon$ occurs in the last member of (10.2), and $C_{\mathrm{mat}}$ is a function of $(L,\rho)$.
In particular, (10.3) contains the reference datum (R2) of Proposition~43.1: $D^{\mathrm{mat}}\le C_{\mathrm{mat}}$ for all $N\ge N_0(e)$.
The constants that enter are $C_{6}$ and $\kappa_{(0)}$ (line (C0r)), $c_{\mathrm{fl}}$ (line (C0)) and $C_{\mathrm{mat}}$, which is listed in line (C10) and is a function of $(L,\rho)$.
Conditions on the constants (Proposition~14.2): the second and the third clause of (T2) of Proposition~14.2\,(c), where the condition (e3) of Proposition~14.2\,(b) at $v=v_K$ gives $e_K^{\frac{1}{4}-9\varepsilon}\le\frac{1}{8}$;
of the condition (S1) of Proposition~14.2\,(b) only $\varepsilon>0$ and $\frac{1}{4}-9\varepsilon>0$ are used.
\end{named}

\begin{named}{Remark}
Proposition~32.3 corresponds to \cite[Theorem~19.25]{QEDI} and to \cite[Theorem~5.86]{QEDI}, with (10.1)--(10.3) in the place of their estimates, and to the bound $D^{\mathrm{mat}}\le C_{\mathrm{mat}}$ of \cite[Theorem~21.3, part~1]{QEDI}, which is derived in Step~5 of the proof of that theorem in \cite[Section~21.9]{QEDI}.
The function $p^0$ of \cite{QEDI} is replaced by $p^{0}_{\mathrm{ch}}$, and $\varepsilon_5$ by the function $\varepsilon_5$ of Section~G4.0.3;
the number $L^{-m_0-1}$ of \cite{QEDI} becomes $1/L$, because here $m_0=0$ and $h'=h$ (Section~I.6.1).
The last inequality of (10.3), which in \cite{QEDI} is obtained only in the proof of \cite[Theorem~21.3]{QEDI}, belongs to the statement, and all inequalities of (10.3) are stated for $\left|D^{\mathrm{mat}}\right|$.
\end{named}

\subsection*{49. Symmetry of the odd reference torus under the reflections in its walls}

\subsubsection*{Proposition 49}

\begin{named}{Proposition 49}
Let $\bar{m}\in(0,1]$, let $e>0$ satisfy $e_K=e\cdot h^{\frac{1}{2}}\le e_1$, let $N\ge N_0(e)$ be an integer, and put $\eta=L^{-N}$ and $q=e\cdot\eta^{\frac{1}{2}}$.
The numbers $n_*$, $\mathrm{per}=n_*/L=2m_{\mathrm{p}}+1$ and $\mathrm{hc}=\mathrm{per}/2$, the cell structure $\mathcal{Q}_{*}$, its walls $\left\{y_\mu=j\cdot\mathrm{hc}\right\}$ with $j\in\mathbb{Z}$ (site planes for even $j$, link planes for odd $j$), the open cells $Q^{\circ}_{\mathfrak{c}}$, the group $W(\mathcal{Q}_{*})$ and its elements $\mathsf{g}_{\mathfrak{c}}$ are those of Section~I.6.4.
Let $\mathrm{kt}\ge1$ be an integer and $n:=\mathrm{kt}\cdot n_*$; the cubic torus with $n$ sites per side is the torus kt.
For $\mathrm{kt}=1$ it is the odd reference torus $R^{*}(e)$, for $\mathrm{kt}=L$ the second reference torus; even values of $\mathrm{kt}$ are admitted in (a) and in (c)--(f).
The spinors are the elements of $\mathfrak{S}^{\mathrm{ch}}(1;n)$.
The cells of the torus kt are indexed by $\mathfrak{c}\in\left(\mathbb{Z}/(2\cdot\mathrm{kt}\cdot L)\mathbb{Z}\right)^3$; for $\mathrm{kt}=1$ their set is $\mathrm{Cel}$.
On cochains $(g_{\mathrm{w}})_*:=\left(g_{\mathrm{w}}^{-1}\right)^{*}$ (Section~I.4.7).
In the second half of (d), in (e) and in the last sentence of (f) a probe $\beta$ and a placement $v$ are given, and $\hat{\mathfrak{a}}$ is the lattice probe formed with them on the torus kt (Section~I.7.6);
the only property of $\hat{\mathfrak{a}}$ that matters is that it is a real $n\mathbb{Z}^3$-periodic one-cochain, and $\beta$ need not fit $\left(N,T^{*}(e)\right)$.
Then the following assertions hold.
\begin{itemize}
\item[(a)] (The group of the walls.)
For every integer $j$, even or odd, the reflection in the wall $\left\{y_\mu=j\cdot\mathrm{hc}\right\}$ is the map $y_\mu\mapsto j\cdot\mathrm{per}-y_\mu$, the other two coordinates being fixed, and it is an element of $\mathrm{Gc}=\mathbb{Z}^3\rtimes\{\pm1\}^3$.
For $\mathfrak{c}\in\mathbb{Z}^3$ one has $\mathsf{g}_{\mathfrak{c}}y=r_{\mathfrak{c}}y+a_{\mathfrak{c}}$, where
\[
r_{\mathfrak{c}}:=\operatorname{diag}\left((-1)^{\mathfrak{c}_1},(-1)^{\mathfrak{c}_2},(-1)^{\mathfrak{c}_3}\right)\in\{\pm1\}^3,\qquad
(a_{\mathfrak{c}})_\mu:=\begin{cases}\frac{\mathfrak{c}_\mu}{2}\cdot\mathrm{per}&\text{if }\mathfrak{c}_\mu\text{ is even},\\ \frac{\mathfrak{c}_\mu+1}{2}\cdot\mathrm{per}&\text{if }\mathfrak{c}_\mu\text{ is odd};\end{cases}
\]
thus $a_{\mathfrak{c}}\in\mathrm{per}\cdot\mathbb{Z}^3$, and
\[
W(\mathcal{Q}_{*})=\left\{\mathsf{g}_{\mathfrak{c}}:\mathfrak{c}\in\mathbb{Z}^3\right\}=\left\{(a,r):a\in\mathrm{per}\cdot\mathbb{Z}^3,\ r\in\{\pm1\}^3\right\}.
\]
The group $W(\mathcal{Q}_{*})$ contains the translations by $n\mathbb{Z}^3$;
the quotient $W(\mathcal{Q}_{*})/n\mathbb{Z}^3$ acts on $\mathbb{Z}^3/n\mathbb{Z}^3$, and it acts simply transitively on the $(2\cdot\mathrm{kt}\cdot L)^3$ cells of the torus kt.
For $g_{\mathrm{w}}\in W(\mathcal{Q}_{*})$ let $g_{\mathrm{w}}.\mathfrak{c}$ be the index with $g_{\mathrm{w}}Q^{\circ}_{\mathfrak{c}}=Q^{\circ}_{g_{\mathrm{w}}.\mathfrak{c}}$;
then $g_{\mathrm{w}}\mathsf{g}_{\mathfrak{c}}=\mathsf{g}_{g_{\mathrm{w}}.\mathfrak{c}}$, modulo $n\mathbb{Z}^3$ on the torus.
If $g_{\mathrm{w}}$ is the reflection in $\left\{y_\mu=j\cdot\mathrm{hc}\right\}$, then $(g_{\mathrm{w}}.\mathfrak{c})_\mu=2j-1-\mathfrak{c}_\mu$, and the other two entries of $\mathfrak{c}$ are unchanged.
On the fermion sites an element $g_{\mathrm{w}}=(a,r)$ acts by $g_{\mathrm{w}}(c,g)=(rc+a,rg)$ (Section~I.4.7);
this is a bijection of $\mathbb{Z}^3\times O_{\mathrm{cub}}$ with $g_{\mathrm{w}}(x+n\cdot e_\mu)=g_{\mathrm{w}}x+n\cdot r\,e_\mu$ and $\left(g_{\mathrm{w}}x\right)^{(i)}=g_{\mathrm{w}}\left(x^{(i)}\right)$, and a reflection fixes no fermion site.
For $S\in\mathrm{Sp}$ with $r(S)=r$ and $\tilde{g}:=(S,a)$, the map $f\mapsto\tilde{g}f$, $(\tilde{g}f)(x)=S\,f\left(\tilde{g}^{-1}x\right)$, is a linear bijection of $\mathfrak{S}^{\mathrm{ch}}(1;n)$ onto itself.
\item[(b)] (Levels; the torus $R^{*}(e)$, $\mathrm{kt}=1$.)
Let $k_{*}:=k_{*}\left(R^{*}(e)\right)$ and $0\le j\le k_{*}$.
Then $\ell_j(R^{*})$ lies in $\mathrm{Cl}$, so that the group $\mathbb{G}_j^{L}$ of the side $\ell_j(R^{*})$ is defined.
In the coordinates of level $j$ the map $\mathsf{g}_{\mathfrak{c}}$ reads $c\mapsto r_{\mathfrak{c}}c+L^{-j}\cdot a_{\mathfrak{c}}$, and $L^{-j}\cdot a_{\mathfrak{c}}\in L\cdot M\cdot\mathbb{Z}^3$.
Hence, for every $S\in\mathrm{Sp}$ with $r(S)=r_{\mathfrak{c}}$, the pair $\tilde{g}^{(j)}_{\mathfrak{c}}:=\left(S,L^{-j}\cdot a_{\mathfrak{c}}\right)$ is an element of $\mathbb{G}_j^{L}$ that induces $\mathsf{g}_{\mathfrak{c}}$;
the matrix $S$ can be chosen in $\mathrm{Sp}_3$, and then $\tilde{g}^{(j)}_{\mathfrak{c}}\in\mathbb{G}^{L}_{j,3}$.
For every complex one-cochain $a$ on $\mathbb{Z}^3/n_*\mathbb{Z}^3$,
\[
\left((\mathsf{g}_{\mathfrak{c}})_*a\right)^{[j]}=\tilde{g}^{(j)}_{\mathfrak{c}}\left(a^{[j]}\right).
\]
\item[(c)] (Determinant.)
Let $A$ be a complex one-cochain on $\mathbb{Z}^3/n\mathbb{Z}^3$, let $z$ be a complex number and $g_{\mathrm{w}}\in\mathrm{Gc}$.
Then, the determinants being taken on $\mathfrak{S}^{\mathrm{ch}}(1;n)$,
\begin{gather*}
\det\left(D_{(g_{\mathrm{w}})_*A}+z\right)=\det\left(D_A+z\right),\qquad
\det\left(D_{-A}+z\right)=\det\left(D_A+z\right),\\
\overline{\det\left(D_{\overline{A}}+\overline{z}\right)}=\det\left(D_A+z\right).
\end{gather*}
\item[(d)] (Free ratio.)
Let $\mathrm{ah}$ be a real one-cochain on $\mathbb{Z}^3/n\mathbb{Z}^3$, let $\mathbf{m}\in(0,1]$, and put
\[
N[\mathrm{ah}](t):=\frac{\det\left(D_{(t/e)\cdot\mathrm{ah}}+\mathbf{m}\right)}{\det\left(D_0+\mathbf{m}\right)}\qquad(t\in\mathbb{C}).
\]
With the charge $q=e\cdot\eta^{\frac{1}{2}}$ fixed above, $N[\mathrm{ah}]$ is the function $N_a$ of Proposition~50\,(h), formed on the cubic torus $(n,n,n)$ with the same $\mathbf{m}$ and with the real one-cochain $a=\mathrm{ah}/e$:
\[
N[\mathrm{ah}](t)=N_{\mathrm{ah}/e}(t)\qquad(t\in\mathbb{C}),
\]
since both sides are the quotient $\det\left(D_{(t/e)\cdot\mathrm{ah}}+\mathbf{m}\right)/\det\left(D_0+\mathbf{m}\right)$ of determinants on $\mathfrak{S}^{\mathrm{ch}}(1;n)$.
Then $N[\mathrm{ah}]$ is an entire function, $N[\mathrm{ah}](0)=1$, $N[\mathrm{ah}](-t)=N[\mathrm{ah}](t)$ and $\overline{N[\mathrm{ah}](\overline{t})}=N[\mathrm{ah}](t)$;
consequently $N[\mathrm{ah}]$ is real on the real axis, and there moreover $N[\mathrm{ah}](t)>0$.
Furthermore $N\left[(g_{\mathrm{w}})_*\mathrm{ah}\right]=N[\mathrm{ah}]$ for every $g_{\mathrm{w}}\in\mathrm{Gc}$.
In particular, let $\hat{\mathfrak{a}}$ be the lattice probe of Section~I.7.6 on the torus kt, and let $\mathbf{m}=\bar{m}\cdot\eta$.
For these data $N[\hat{\mathfrak{a}}]$ is the free ratio $N_0$ of Section~I.7.6, formed on the torus kt:
\[
N[\hat{\mathfrak{a}}](t)=N_0(t)=\frac{\det\left(D_{e^{-1}t\hat{\mathfrak{a}}}+\bar{m}\cdot\eta\right)}{\det\left(D_0+\bar{m}\cdot\eta\right)}\qquad(t\in\mathbb{C}).
\]
Hence $N_0$ does not change when $\hat{\mathfrak{a}}$ is replaced by a mirror copy $\hat{\mathfrak{a}}_{\mathfrak{c}}=(\mathsf{g}_{\mathfrak{c}})_*\hat{\mathfrak{a}}$;
the function $N_0$ is even and real, $c_m=0$ for $m=1,3$, and $c_2$, $c_4$ are real.
(The link phases of $D_{(t/e)\cdot\hat{\mathfrak{a}}}$ are $\exp\left(\pm i\,t\cdot\eta\cdot a_{\eta}(b)\right)$, so that the explicit factor $e$ cancels; but $a_{\eta}$ still depends on $e$, through the period $\mathrm{kt}\cdot R^{*}(e)$ and through the placement.)
\item[(e)] (Partition function with marks.)
Let $\boldsymbol{\mu}$ be a real number and let $\hat{\mathfrak{a}}_{\mathfrak{c}}$ be as in (d).
For a family $t=(t_{\mathfrak{c}})$ of complex numbers indexed by the cells $\mathfrak{c}$ of the torus kt and for $\theta\in\mathbb{R}^3$ put
\begin{gather*}
\mathrm{Zmk}^{\theta}(t):=\int\mathrm{d}\gamma_N(F^{\mathrm{u}})\,\det\left(D_{A(F^{\mathrm{u}})+H_{\theta}+e^{-1}\sum_{\mathfrak{c}}t_{\mathfrak{c}}\hat{\mathfrak{a}}_{\mathfrak{c}}}+\boldsymbol{\mu}\right),\\
\bar{Z}^{\mathrm{ext}}(t)=\int_{[-\pi,\pi]^3}\frac{\mathrm{d}\theta}{(2\pi)^3}\,\mathrm{Zmk}^{\theta}(t);
\end{gather*}
for $\mathrm{kt}=1$ and $\boldsymbol{\mu}=\boldsymbol{\mu}_N$ the function $\bar{Z}^{\mathrm{ext}}$ so defined is that of Section~I.7.6.
For $g_{\mathrm{w}}=(a,r)\in W(\mathcal{Q}_{*})$ let $g_{\mathrm{w}}t$ be the family with $(g_{\mathrm{w}}t)_{g_{\mathrm{w}}.\mathfrak{c}}:=t_{\mathfrak{c}}$.
Then
\[
\mathrm{Zmk}^{r\theta}(g_{\mathrm{w}}t)=\mathrm{Zmk}^{\theta}(t),\qquad
\mathrm{Zmk}^{-\theta}(-t)=\mathrm{Zmk}^{\theta}(t),\qquad
\overline{\mathrm{Zmk}^{\theta}(\overline{t})}=\mathrm{Zmk}^{\theta}(t);
\]
hence
\begin{gather*}
\bar{Z}^{\mathrm{ext}}(g_{\mathrm{w}}t)=\bar{Z}^{\mathrm{ext}}(t),\qquad
\bar{Z}^{\mathrm{ext}}(-t)=\bar{Z}^{\mathrm{ext}}(t),\qquad
\overline{\bar{Z}^{\mathrm{ext}}(\overline{t})}=\bar{Z}^{\mathrm{ext}}(t),\\
\bar{Z}^{\mathrm{ext}}(t)\ge0\ \text{ for real }t,\qquad
\bar{Z}^{\mathrm{ext}}(0)>0,
\end{gather*}
and the same three identities hold for $\mathring{\Phi}(t)=\bar{Z}^{\mathrm{ext}}(t)/\bar{Z}^{\mathrm{ext}}(0)$.
In particular, the function of one variable obtained from $\mathring{\Phi}$ by setting all marks except $t_{\mathfrak{c}}$ equal to $0$ is the same for all cells $\mathfrak{c}$, and the number
\[
\left.\prod_{\mathfrak{c}}\left(\partial_{t_{\mathfrak{c}}}^{m}-c_m\right)\mathring{\Phi}(t)\right|_{t=0}
\]
is real for every integer $m\ge0$ and every real number $c_m$;
the numbers $c_m$ of Section~I.7.6, on whichever torus they are formed, are real by (d).
\item[(f)] (Measures.)
For every real pair $\left(m_{\mathrm{in}},\varepsilon^{\mathrm{v}}_{\mathrm{in}}\right)$ the measures of Proposition~31.1 on the torus kt have the following properties.
For $g_{\mathrm{w}}=(a,r)\in\mathrm{Gc}$ the image of $\nu^{\theta}$ under $F^{\mathrm{u}}\mapsto(g_{\mathrm{w}})_*F^{\mathrm{u}}$ is $\nu^{r\theta}$, and its image under $F^{\mathrm{u}}\mapsto-F^{\mathrm{u}}$ is $\nu^{-\theta}$.
The measures $\bar{\nu}$ and $\bar{\mu}$ are invariant under $F^{\mathrm{u}}\mapsto(g_{\mathrm{w}})_*F^{\mathrm{u}}$ for $g_{\mathrm{w}}\in\mathrm{Gc}$, and under $F^{\mathrm{u}}\mapsto-F^{\mathrm{u}}$.
Finally $X_{\mathfrak{c}}=\left\langle d\hat{\mathfrak{a}}_{\mathfrak{c}},F^{\mathrm{u}}\right\rangle=\left\langle d\hat{\mathfrak{a}},\mathsf{g}_{\mathfrak{c}}^{*}F^{\mathrm{u}}\right\rangle$;
so the law of $X_{\mathfrak{c}}$ under $\bar{\mu}$ does not depend on $\mathfrak{c}$, and it is symmetric.
\end{itemize}
No constant occurs in Proposition~49.
Conditions on the constants (Proposition~14.2): in (b), the condition (G4.2.1) of Proposition~14.2\,(c), in the form that $r_*$ is a power of $L$ with $r_*\ge L^2$, and, through (Sc0), (Sc1), (Sc4) of Section~G4.0.2, the first and the third inequality of (e1) of Proposition~14.2\,(b) at $v=v_K$ and the inequality $L\ge13$, a consequence of (L1) of Proposition~14.2\,(b);
in all clauses, (L5), which gives that $\mathrm{per}$ is odd and $n\ge5$, and (M1), which gives that $T^{*}(e)$ is an integer $\ge1$, both of Proposition~14.2\,(b).
\end{named}

\begin{named}{Remark}
Proposition~49 corresponds to \cite[Lemma~20.4(b)]{QEDI}, together with \cite[Section~20.1.3]{QEDI} and \cite[Lemma~20.2(a),~(b)]{QEDI}.
In the proof of \cite[Lemma~20.4(b)]{QEDI} the link phases of $D_{(t/e)\cdot\hat{\mathfrak{a}}}$ on the torus $R^{*}$ are written, for real $t$, as
\[
\exp\left(\pm i\,q\cdot\frac{t}{e}\cdot\hat{\mathfrak{a}}(b)\right)=\exp\left(\pm i\,t\cdot\eta\cdot a_{\eta}(b)\right),
\]
since $q=e\cdot\eta^{\frac{1}{2}}$ and $\hat{\mathfrak{a}}=\eta^{\frac{1}{2}}\cdot a_{\eta}$ (\cite[Section~20.1.4]{QEDI});
the identity at the end of (d) is the same, on the torus kt.
The differences are the following.
The fermion sites $(c,g)$ carry a label $g\in O_{\mathrm{cub}}$; the group acts on them by $(c,g)\mapsto(rc+a,rg)$ and on spinors through the spin matrix $S$.
The determinants are taken on $\mathfrak{S}^{\mathrm{ch}}(1;n)$, a space of dimension $192\,n^3$.
Of the clauses (c), (d), (e), on the determinant, the free ratio and the partition function with marks, only (d) has a counterpart in \cite[Lemma~20.4(b)]{QEDI}.
All clauses except (b) are stated on every torus kt, $\mathrm{kt}\ge1$, of either parity, and not on the torus $R^{*}$ only.
The two determinant identities $\det\left(D_{-A}+z\right)=\det\left(D_A+z\right)$ and $\det\left(D_{(g_{\mathrm{w}})_*A}+z\right)=\det\left(D_A+z\right)$ follow for every complex $A$ from Section~I.4.4 and (I.O5), with no analytic continuation in $A$;
the proof of \cite[Lemma~20.4(b)]{QEDI} takes them from \cite[Proposition~5.47]{QEDI}, applied at the level $0$ in the scale-free setting.
The element $\mathsf{g}_{\mathfrak{c}}$ is given by the formula of (a), and its existence and uniqueness follow from the simple transitivity of the action on the cells;
in \cite[Section~20.1.3]{QEDI} it is introduced as the element of the group generated by the reflections that maps the cell with index $0$ onto the cell with index $\mathfrak{c}$.
Finally, (b) contains the fact that $\ell_j(R^{*})$ lies in $\mathrm{Cl}$, on which the definition of the group $\mathbb{G}_j^{L}$ in \cite[Lemma~20.2(b)]{QEDI} depends.
\end{named}

\subsection*{G5. Definitions for the Propositions 33.1--37.3}

\subsubsection*{G5.0.3. Notation}

In the Sections G5.0.3--G5.6.0 and in the Propositions 33.1--37.3, $\mathrm{sv}$ always denotes a real parameter or a variable of integration; in Proposition~39\,(e), in Proposition~40 and in Section~G6.A3 it always denotes a number of cells per side.

The notation of \cite[Section~20]{QEDI} is related to the notation used here as follows.
In \cite[Section~20.1.7]{QEDI} the numbers
\[
\left(b_-^{-1}+\mu_-^{-1}\right)^{\frac{1}{2}},\qquad b_-^{-1}+\frac{4}{\mu_-},\qquad 2\,b_-^{-1}+4\,\mathrm{e}\cdot\mu_-^{-1}\cdot\mathrm{ssum}\left(\frac{\theta_{\mathrm{CT}}}{2}\right)
\]
are denoted $g$, $g_1$ and $\boldsymbol{\alpha}$; here they are denoted $\mathrm{gGa}$, $\mathrm{gGa}_1$ and $\mathrm{alG}$.
The constants $c_1$, $C_V$, $C_W$ and $M_L=2+\exp\left(\mathbf{D}/2\right)$ of that section correspond to $c_\tau$, $\mathrm{CVj}$, $\mathrm{CWj}$ and $\mathrm{mxL}$, and its number $z_2$ equals $t_2/e$; the numbers $c_\tau$, $\mathrm{CVj}$, $\mathrm{CWj}$, $\mathrm{mxL}$ and $t_2$ are defined in Section~G5.1.8.
Where \cite[Section~20.3.4]{QEDI} uses the row sums $\boldsymbol{\alpha}_{\vartheta}$, $\boldsymbol{\alpha}_0$ of a covariance, both numbers are taken equal to $\mathrm{alG}$.
The radii $\rho$, $\rho'$ of \cite[Sections~20.1.4, 20.1.6, 20.3.1 and~20.3.5]{QEDI}, that is, the radius of a polydisc $\mathbb{D}(\rho)$ and the radius in the norm of a jet, are denoted $\mathrm{rad}$, $\mathrm{rad}'$.
The norms $\rho_h$, $N_h$ of even elements and the weights $w_{\boldsymbol{\delta}}$, $w_P$ of \cite[Section~20.3.4]{QEDI} are denoted $\mathrm{rw}_G$, $\mathrm{Nw}_G$, $\mathrm{wg}_{\delta}$ and $\mathrm{wg}_P$.
The half side $\rho_\beta$ of \cite[Section~20.1.4]{QEDI} is $D_\beta/2$.

The cells $c$, $c'$ of \cite[Section~20]{QEDI}, its cell $c_0$ and the unit vector $e_c$ of $\mathbb{C}^{\mathrm{Cel}}$ are denoted $\mathfrak{c}$, $\mathfrak{c}'$, $0$ and $\mathrm{ue}_{\mathfrak{c}}$.
The sets of cells $S$, $T$ of \cite[Sections~20.3.1 and~20.8]{QEDI} are denoted $\mathrm{Scl}$, $\mathrm{Scl}'$, and the elements $k_S$ of $\mathcal{J}$ of \cite[Section~20.8]{QEDI} are denoted $\mathrm{kj}_{\mathrm{Scl}}$.
The Taylor coefficients $n_i$ of $N_0$ in \cite[Lemma~20.23\,(e)]{QEDI} are denoted $\mathrm{nj}_i$, and the function $w_c$ of \cite[Proposition~20.24]{QEDI}, the difference $\mathrm{Lg}-\mathrm{lgN}$, is denoted $\mathrm{wdf}_{\mathfrak{c}}$.
The measure $\boldsymbol{\varphi}$ on $[-\pi,\pi]^3\times V_{\mathrm{t}}\left(T^0_*\right)$ of \cite[Section~20.7]{QEDI} is denoted $\mathrm{vm}$.
The terminal integral $\mathsf{Z}(t)$, the small-field part $\mathscr{Y}(t)$, the exponent $V$ with its parts $V_0$, $V_+$, and the covariance $\mathcal{G}$ of \cite[Sections~20.5 and~20.7]{QEDI} are denoted $\mathrm{Ztm}(t)$, $\mathrm{Ysf}(t)$, $\mathrm{Vj}$, $\mathrm{Vj}_0$, $\mathrm{Vj}_+$ and $\mathrm{Gcv}$.
In the notation used here the number $\mathbf{Z}(t)$ of \cite[Section~20.1.5]{QEDI} is $\bar{\nu}^{a(t)}(F_N)$.

The coefficient $\mathfrak{H}_A(\boldsymbol{\upsilon};u)$ of $\Psi(u)$ in a leg, \cite[Section~20.4]{QEDI}, is denoted $\mathrm{Leg}_A(\mathrm{up};u)$.
The function $E(y,y')=\exp\left(-\gamma_F\cdot d(y,y')\right)$ of \cite[Section~20.2.2]{QEDI} is denoted $\mathrm{Edy}(y,y')$.
The operator $\mathsf{D}(A)$ of that section, its inverse $S(A)$ and the H\"older quotient $\mathrm{Sa}(A)$ are denoted $\mathrm{Fst}(A)$, $S_*(A)$ and $\mathrm{Sa}_*(A)$; the operator $\mathrm{Fst}(A)$ is the operator $F_\infty(A)$ of Section~I.3.5 at the level $k_*$.
For the function $\sigma_A(y,x)$ of \cite[Section~20.2.2]{QEDI} one writes $\mathrm{phx}_A(y,c)$; at the level $k_*$ the phase of (I.3.3) is $\varphi_A(y,c)=\exp\left(i\,\mathrm{phx}_A(y,c)\right)$.
The sets of unit sites $P_B$, $P_B^+$, $P_c$, $P_c^+$ of \cite[Sections~20.1.4 and~20.4]{QEDI} are denoted $\mathrm{Ypr}_B$, $\mathrm{Ypr}_B^+$, $\mathrm{Ypr}_{\mathfrak{c}}$, $\mathrm{Ypr}_{\mathfrak{c}}^+$, and the domains $G_{\mathrm{w}}$, $G_{\mathrm{w}}^+$ of \cite[Section~20.4]{QEDI} are denoted $\mathrm{Gwk}$, $\mathrm{Gwk}^+$.
The scaled probe $\acute{g}_\nu$ of \cite[Section~20.1.4]{QEDI} is denoted $\mathrm{gsc}_\nu$.
The weight $a_2=C_0^{-1}\cdot e_*^{-\frac{1}{4}+\varepsilon}$ of \cite[Section~20.2.4]{QEDI} equals $C_0^{-1}\cdot h_{*,2}$; it is the number $\mathrm{aw}_2$ of Proposition~20.1.
The weight $a$ of \cite[Lemma~20.19]{QEDI} is denoted $\mathrm{aw}$, as in Proposition~15 and Proposition~20.1, and the real parameter $u$ of \cite[Sections~20.2.2 and~20.2.3]{QEDI} is denoted $\mathrm{sv}$.

The numbers $h_L$, $\vartheta$, $\boldsymbol{\theta}_{\mathrm{CT}}$, $\bar{\kappa}$ and $\boldsymbol{\lambda}_L$ of \cite[Section~20.1.7]{QEDI} are denoted $\mathrm{hL}$, $\mathrm{vt}$, $\theta_{\mathrm{CT}}$, $\bar{\kappa}_1$ and $\lambda_L$.
The number $\mathsf{p}$, the cube $\bar{Q}$ and the polydiscs $\mathbb{D}(\rho)$, $\mathbb{D}^{\circ}(\rho)$ of \cite[Section~20.1.4]{QEDI} are denoted $\mathrm{pp}$, $\mathrm{Qb}$, $\mathrm{DD}(\mathrm{rad})$ and $\mathrm{DDo}(\mathrm{rad})$.
In \cite[Section~20.2.2]{QEDI} one has $U_A=\exp\left(i\,\upsilon_A\right)$; the function $\upsilon_A$ is denoted $\mathrm{ups}_A$, and a leg index $\boldsymbol{\upsilon}$ of \cite[Section~20.4]{QEDI} is denoted $\mathrm{up}$.
The function $\mathfrak{g}$ of \cite[Sections~20.1.6 and~20.3.1]{QEDI} is denoted $\mathrm{gf}$, and an even element $\mathbf{F}$ of \cite[Section~20.3.4]{QEDI}, with its kernels $\mathbf{F}_m$, is denoted $\mathrm{Fv}$, with the kernels $\mathrm{Fv}_m$.
In the notation of Section~G5.1.3 the function $\mathring{g}^0_\nu$ of \cite[Section~20.1.4]{QEDI} is $\mathrm{gs}_\nu[\beta]$, formed with $\eta$, and the function $\mathring{g}_\nu$ of that section, which belongs to the placement $v$, is
\[
\mathring{g}_\nu(y)=\sum_{z\in\mathbb{Z}^3}\mathrm{gs}_\nu[\beta]\left(y-v-R^*(e)\cdot z\right)\qquad\left(y\in\mathbb{R}^3\right);
\]
accordingly $\mathrm{gs}^0_\nu=\mathrm{gs}_\nu[U_r\beta]$ is the function $\mathring{g}^0_\nu$ with $U_r\beta$ and $\eta'$ in place of $\beta$ and $\eta$.

In Section~G5.6.0 and in the Propositions 36 and 37.1--37.3 the following notation is used.
The continuum two-photon coefficient is denoted $\mathrm{pv}(f)$ and its density $\mathrm{vpo}_{\bar{m}}(\mathrm{ts})$; the external momentum in physical units is denoted $\mathrm{kph}$.
The lattice vacuum polarisation is denoted $\mathrm{Vpl}^{\eta,R}$, $\mathrm{Vpl}^{\eta}$, and the lattice transform of the probe $\mathrm{ahl}_{\eta,\nu}$.
The inverse of the symbol is denoted $\mathrm{Rsy}(p)$, and the derivative of the symbol $\mathrm{Vx}_\nu(p)$.
The vertices are denoted $\mathrm{Vd}_j$; where a vertex is normalised with the factor $1/k!$, it is written $\mathrm{Vd}_k/k!$.
The hop matrices in the fibre are denoted $\mathrm{Ehp}^{\pm}_\nu$, the loop integrand $\mathrm{lp}_{\mu\nu}(p,k)$, and
\[
\mathrm{G0}:=\left(D^{\mathrm{ch}}[0]+\bar{m}\cdot\eta\right)^{-1}.
\]
Finally, the parity involution of the fibre $V_{\mathrm{v}}$ is denoted $\mathrm{Pa}$; it is defined in Section~G5.6.0.

\subsubsection*{G5.1. Common setting of the Propositions 33.1--35.10}

\subsubsection*{G5.1.1. The standing setting (ST*)$^{\mathrm{ch}}$; scales}

\textbf{(ST*)$^{\mathrm{ch}}$.}
The constants are those of Section~I.8, lines (C0)--(C9), chosen in this order.
The data are: a mass $\bar{m}\in(0,1]$, which determines $h$ (Section~I.0.5); a charge $e>0$ with $e\cdot h^{\frac{1}{2}}\le e_1$; the torus of side
\[
R=R^*(e)=\frac{h\cdot M\cdot r_K(e)}{L}
\]
(Section~I.6.2); and an integer $N\ge N_0(e)=N_{\min}\left(R^*(e),e\right)$.
With these data
\[
\eta:=L^{-N},\qquad q:=e\cdot\eta^{\frac{1}{2}},\qquad n_*:=\frac{R^*(e)}{\eta},
\]
and $n_*$ is odd.
The pair is
\[
\left(m_{\mathrm{in}},\varepsilon^{\mathrm{v}}_{\mathrm{in}}\right)=P_N(e)=\left(m_{0,N}(e),\varepsilon^{\mathrm{v}}_{0,N}(e)\right)
\]
(Section~I.6.2), and $\boldsymbol{\mu}_N=\bar{m}\cdot\eta+m_{0,N}(e)$.
At the levels $j<k_*$ the primitives on $T_j\left(R^*(e)\right)$ are the projections (Section~I.6.3) of the systems of the scheme on $T_j\left(R_{\mathrm{top}}(e)\right)$.

Let $k_*:=k_*\left(R^*(e)\right)$.
A subscript $*$ indicates the value at $k=k_*$.
This applies to $a_*$, $e_*$, $v_*$, $r_*$, $p_*$, $p_{0*}:=p_{0,k_*}$, $\ell_*$, $\mu_*=\bar{m}\cdot a_*$, $b_*$, $\eta_*=L^{-k_*}$, $T_*$, $T^{\mathrm{f}}_*$, $T^0_*$, $\mathrm{Pz}_*$, to $D_*$, which without an argument is the set of the polymers of $T_*$, to
\[
h_*=\left(h_{*,1},h_{*,2}\right)=\left(e_*^{-\frac{1}{4}},\,e_*^{-\frac{1}{4}+\varepsilon}\right),
\]
and to $w^{\mathrm{m}}_*$, $w^{\mathrm{E}}_*=e_*^{\frac{1}{4}-7\varepsilon}$, $\tilde{R}_*$, $\tilde{G}_*$, $\mathfrak{X}_*$, $\mathrm{Dom}_*$.
The number $n_*$ is of another kind.
It is the number $R^*(e)/\eta=\ell_0$ of the photon sites per side (Section~I.6.4), whereas $n_k=\ell_k^3$ (Section~I.6.1); thus $n_{k_*}=\ell_*^3$.

The trajectories of $P_N(e)$ on $R^*(e)$ (Section~I.5.7, paragraph ``Flows'') are $\left(\varepsilon^{\mathrm{v},t}_k,m_k,E_k\right)$ and $\left(\mathscr{E}_k,\mathfrak{P}_k\right)$, $0\le k\le k_*$.
One writes $m_*:=m_{k_*}$, $\varepsilon^{\mathrm{v},t}_*:=\varepsilon^{\mathrm{v},t}_{k_*}$, $\mathscr{E}_*:=\mathscr{E}_{k_*}$ and $\mathfrak{P}_*:=\mathfrak{P}_{k_*}$.
Moreover
\[
\mu_-:=L^{-3}.
\]
In \cite[Section~20.1.2]{QEDI} this number is $L^{-m_0-3}$; here $m_0=0$ by Section~I.6.1.

No condition on the sign of $\boldsymbol{\mu}_N$ is imposed in Section~G5 or in the Propositions 33.1--37.3 (see (I.K10)): the propagators contain the running masses $\bar{m}\cdot a_j>0$, and $m_j$ enters as a coupling.

\subsubsection*{G5.1.2. Cells}

At the cutoff $N$, Section~I.6.4 gives
\[
\mathrm{per}=\frac{n_*}{L}=2m_{\mathrm{p}}+1,\qquad \mathrm{hc}=\frac{\mathrm{per}}{2},\qquad \ell_c=\mathrm{hc}\cdot\eta=\frac{R^*(e)}{2L},\qquad N_*=8L^3.
\]
The set of the $N_*$ cells of $C(n_*)$ is denoted $\mathrm{Cel}$; the cells are indexed by $\mathfrak{c}\in(\mathbb{Z}/2L)^3$, and $\mathsf{g}_{\mathfrak{c}}$ is the map of Section~I.6.4.
Physical coordinates are $\eta$ times the lattice coordinates of Section~I.6.4; in them the closure of the cell $Q^{\circ}_0$ is $[0,\ell_c]^3$.
For a lattice isometry $g_{\mathrm{w}}$ and a bond function $A$ (Section~I.4.7),
\[
\left(\left(g_{\mathrm{w}}\right)_*A\right)(c,c')=A\left(g_{\mathrm{w}}^{-1}c,\,g_{\mathrm{w}}^{-1}c'\right).
\]

\subsubsection*{G5.1.3. Probes, placements, the numbers $C_\beta$ and $\mathbf{D}$}

\emph{Probe.}
A probe is a real two-form $\beta\in C_c^\infty(\mathbb{R}^3;\Lambda^2)$ with $\beta\ne0$ and $\beta_{\nu\mu}=-\beta_{\mu\nu}$.
One puts $a:=d^*\beta$ and $f:=dd^*\beta$; the number $b(f)$ is that of Section~I.7.6.
For an integer $i\ge0$, $|\beta|_{C^i}$ is the largest of the sup norms of the derivatives of order at most $i$ of the components of $\beta$.
For a smooth function $\mathrm{fs}$ on $\mathbb{R}^3$ with bounded derivatives, $|D^i\mathrm{fs}|_\infty$ is the largest of the sup norms of the partial derivatives of order exactly $i$ of $\mathrm{fs}$, as in Proposition~47.1.
The least side of a closed coordinate cube that contains $\operatorname{supp}\beta$ is denoted $D_{\mathrm{be}}$, and $D_\beta$ and $Q_\beta=y_\beta+\left[-D_\beta/2,D_\beta/2\right]^3$ are those of Section~I.7.6.
One has $D_\beta=D_{\mathrm{be}}+4$, because a closed coordinate cube contains $\operatorname{supp}\beta+[-2,2]^3$ if and only if the concentric cube whose side is smaller by $4$ contains $\operatorname{supp}\beta$.
The forms $U_r\beta$, $r\in O_{\mathrm{cub}}$, and $\tau_v\beta$, $v\in\mathbb{R}^3$, are those of Section~I.7.5.

\emph{Lattice probe on a torus.}
Let $0<\eta\le1$, and let $R>0$ be such that $R/\eta$ is an integer $\ge2$.
For a smooth function $\mathrm{fs}$ on $\mathbb{R}^3$ put
\[
\delta_{\eta,\mu}\mathrm{fs}(y):=\frac{\mathrm{fs}\left(y+\frac{\eta}{2}e_\mu\right)-\mathrm{fs}\left(y-\frac{\eta}{2}e_\mu\right)}{\eta},
\]
and let $J_\eta$ be the cube average of Section~I.7.6.
For a probe $\beta'$ let
\begin{gather*}
\mathrm{gs}_\nu[\beta']:=-\sum_{\mu\ne\nu}\delta_{\eta,\mu}J_\eta\beta'_{\mu\nu},\\
u_{\mu\nu}(c):=\sum_{z\in\mathbb{Z}^3}\left(J_\eta\beta'_{\mu\nu}\right)\left(c+\frac{\eta}{2}\left(e_\mu+e_\nu\right)+R\cdot z\right),\qquad a_\eta[\beta']:=d_\eta^*u,
\end{gather*}
that is,
\[
a_\eta[\beta']_\nu(c)=-\eta^{-1}\sum_{\mu\ne\nu}\left(u_{\mu\nu}(c)-u_{\mu\nu}(c-\eta\cdot e_\mu)\right)\qquad\left(c\in\Lambda_{\eta,R}\right).
\]
Here $d_\eta=\eta^{-1}d$, and $d_\eta^*$ is its adjoint for $\langle\,,\rangle_\eta$.
The function $\mathrm{gs}_\nu[\beta']$ belongs to $C_c^\infty(\mathbb{R}^3)$ and depends on the mesh: it is formed with the number $\eta$, or $\eta'$, that is named where it occurs.
In lattice units one uses the one-cochain $\eta^{\frac{1}{2}}\cdot a_\eta[\beta']$ on $\mathbb{Z}^3/n\mathbb{Z}^3$, $n=R/\eta$, whose value on the bond $(c,c+e_\nu)$ is $\eta^{\frac{1}{2}}\cdot a_\eta[\beta']_\nu(\eta\cdot c)$.
This is the lattice probe of Section~I.7.6; neither a cell nor a placement enters its definition.

\emph{Placement} (setting (ST*)$^{\mathrm{ch}}$).
A vector $v\in\eta\mathbb{Z}^3$ is an \emph{admitted placement} of $\beta$ if, in physical coordinates,
\begin{equation}
\mathrm{Qb}:=v+Q_\beta\ \subset\ \left[\frac{\ell_c}{8},\frac{7\ell_c}{8}\right]^3.
\tag{G5.1.1}
\end{equation}
A \emph{centred placement} is a vector $v\in\eta\mathbb{Z}^3$ with
\[
\left|v+y_\beta-\frac{\ell_c}{2}(1,1,1)\right|\le\frac{\eta}{2}
\]
(Section~I.7.6); there may be several of them.
Every centred placement is admitted if $D_\beta/2+\eta/2\le\frac{3}{8}\ell_c$.
For every $v\in\eta\mathbb{Z}^3$ put
\begin{gather*}
a_\eta:=a_\eta[\tau_v\beta]\ \text{on}\ \Lambda_{\eta,R^*(e)},\qquad \hat{\mathfrak{a}}:=\eta^{\frac{1}{2}}\cdot a_\eta\ \text{(lattice units)},\qquad \boldsymbol{\sigma}^2:=\|d\hat{\mathfrak{a}}\|^2,\\
\hat{\mathfrak{a}}_{\mathfrak{c}}:=\left(\mathsf{g}_{\mathfrak{c}}\right)_*\hat{\mathfrak{a}},\qquad X_{\mathfrak{c}}:=\left\langle d\hat{\mathfrak{a}}_{\mathfrak{c}},F^{\mathrm{u}}\right\rangle,\qquad X:=X_0\qquad\left(\mathfrak{c}\in\mathrm{Cel}\right).
\end{gather*}

\emph{The three norms; $|a|_{\mathrm{ext}}$.}
For a bond function $A$ on a lattice of mesh $\eta''$, a torus $T_k$ or $\eta''\mathbb{Z}^3$, the numbers $|A|_\infty$, $|\partial A|_\infty$ and $|\delta_\alpha\partial A|_\infty$ are the three norms of Section~I.5.3 (that is, of \cite[Section~5.3.1]{QEDI}), formed with the mesh $\eta''$; the third of them is
\[
|\delta_\alpha\partial A|_\infty:=\sup_{\mu,\nu}\ \sup_{0<|c-c'|\le1}|c-c'|^{-\alpha}\cdot\left|\partial_\nu A_\mu(c')-\partial_\nu A_\mu(c)\right|.
\]
Further, $a^{[k]}(b):=L^{\frac{k}{2}}\cdot a\left(L^k\cdot b\right)$ (Section~I.6.6).
For a complex one-cochain $a$ on $\mathbb{Z}^3/n_*\mathbb{Z}^3$, $|a|_{\mathrm{ext}}$ denotes the sum of the three norms of $a^{[k_*]}$ on $T_*$.

\emph{The numbers $C_\beta$ and $\mathbf{D}$.}
They are explicit, and they depend neither on $e$ nor on a torus or a placement.
Let $\mathrm{Ind}_\beta$ be the set of the triples $(N',\mathrm{asc},r)$ formed by an integer $N'\ge0$, one of the two numbers $\mathrm{asc}\in\left\{h/L,\,h/L^2\right\}$ with $N'+\log_L\mathrm{asc}\ge1$, and $r\in\{\pm1\}^3$.
For such a triple put $\eta':=L^{-N'}$ and $\eta'':=\eta'/\mathrm{asc}$, let $\mathrm{gs}^0_\nu:=\mathrm{gs}_\nu[U_r\beta]$, formed with $\eta'$, and let
\begin{align*}
\mathrm{gsc}_\nu(c)&:=\mathrm{asc}^{\frac{1}{2}}\cdot\mathrm{gs}^0_\nu\left(\mathrm{asc}\cdot c+\frac{\eta'}{2}e_\nu\right)&&\left(c\in\eta''\mathbb{Z}^3\right),\\
\mathrm{w0}_\nu(c)&:=\eta'^{\frac{1}{2}}\cdot\mathrm{gs}^0_\nu\left(\eta'\cdot c+\frac{\eta'}{2}e_\nu\right)&&\left(c\in\mathbb{Z}^3\right);
\end{align*}
thus $\mathrm{gsc}$ is a bond function on $\eta''\mathbb{Z}^3$, and $\mathrm{w0}$ is a one-cochain on $\mathbb{Z}^3$.
Then
\begin{gather*}
C_\beta:=\sup_{\mathrm{Ind}_\beta}\left(|\mathrm{gsc}|_\infty+|\partial\,\mathrm{gsc}|_\infty+|\delta_\alpha\partial\,\mathrm{gsc}|_\infty\right),\qquad \mathbf{D}:=\sup_{\mathrm{Ind}_\beta}\|d\,\mathrm{w0}\|^2,\\
\mathrm{Cb}:=1+h^{\frac{1}{2}}\cdot C_\beta,
\end{gather*}
where $\|d\,\mathrm{w0}\|$ is the counting norm over the plaquettes of $\mathbb{Z}^3$.

\subsubsection*{G5.1.4. Marks, external connection, the sets of a probe}

Throughout this section the setting is (ST*)$^{\mathrm{ch}}$, $\beta$ is a probe, and the placement is admitted.
Put
\[
t_x:=\frac{e_*^{3\varepsilon}}{\mathrm{Cb}},\qquad t_x^+:=e_*^{3\varepsilon}\cdot\left(h^{\frac{1}{2}}\cdot C_\beta\right)^{-1},\qquad \mathbf{r}:=\frac{t_x}{4},\qquad q_x:=\frac{2}{t_x},
\]
with $t_x^+:=\infty$ if $C_\beta=0$; thus $t_x^+>t_x$.
For $t=\left(t_{\mathfrak{c}}\right)\in\mathbb{C}^{\mathrm{Cel}}$ let $\mathrm{ue}_{\mathfrak{c}}$ be the unit vector of the cell $\mathfrak{c}$, and put
\[
a(t):=\sum_{\mathfrak{c}}\frac{t_{\mathfrak{c}}}{e}\cdot\hat{\mathfrak{a}}_{\mathfrak{c}},\qquad \hat{c}_{\mathfrak{c}}:=\frac{1}{e}\cdot\hat{\mathfrak{a}}_{\mathfrak{c}}^{[k_*]},\qquad \mathscr{C}(t):=a(t)^{[k_*]}=\sum_{\mathfrak{c}}t_{\mathfrak{c}}\cdot\hat{c}_{\mathfrak{c}};
\]
each $\hat{c}_{\mathfrak{c}}$ is a real bond function on $T_*$.
The polydiscs are
\[
\mathrm{DD}(\mathrm{rad}):=\left\{t:\ \left|t_{\mathfrak{c}}\right|\le\mathrm{rad}\ \text{for all}\ \mathfrak{c}\right\},\qquad \mathrm{DDo}(\mathrm{rad}):=\left\{t:\ \left|t_{\mathfrak{c}}\right|<\mathrm{rad}\ \text{for all}\ \mathfrak{c}\right\}.
\]
To the cell $\mathfrak{c}$ belong the sets
\begin{align*}
\mathrm{Bd}_{\mathfrak{c}}&:=\left\{\text{bonds}\ b\ \text{of}\ T_*:\ \hat{c}_{\mathfrak{c}}(b)\ne0\right\},\\
\mathrm{Ypr}_{\mathfrak{c}}&:=\left\{y\in T^0_*:\ \Delta_y\ \text{contains an end point of a bond of}\ \mathrm{Bd}_{\mathfrak{c}}\right\},\\
\mathrm{Ypr}_{\mathfrak{c}}^+&:=\left\{y\in T^0_*:\ d\left(y,\mathrm{Ypr}_{\mathfrak{c}}\right)\le M+1\right\},
\end{align*}
and one puts
\[
\mathrm{pp}:=\left(\frac{D_\beta}{a_*}+2M+4\right)^3,\qquad \epsilon_{\mathfrak{c}}:=e_*\cdot\left|\hat{c}_{\mathfrak{c}}\right|_\infty,\qquad d_{\mathrm{sep}}:=\min_{\mathfrak{c}\ne\mathfrak{c}'}d\left(\mathrm{Ypr}_{\mathfrak{c}},\mathrm{Ypr}_{\mathfrak{c}'}\right).
\]

\emph{Distances of sets; empty sets.}
For a unit site $y$ and sets $Z$, $Z'$ of unit sites of $T^0_*$,
\[
d(y,Z):=d(Z,y):=\min_{w\in Z}d(y,w),\qquad d(Z,Z'):=\min_{w\in Z,\ w'\in Z'}d(w,w').
\]
The distance of a point or of a set to the empty set is $+\infty$, and an empty sum is $0$.
If $\hat{\mathfrak{a}}=0$, then all $\hat{c}_{\mathfrak{c}}$ vanish, all $\mathrm{Bd}_{\mathfrak{c}}$, $\mathrm{Ypr}_{\mathfrak{c}}$, $\mathrm{Ypr}_{\mathfrak{c}}^+$ are empty, $d_{\mathrm{sep}}=+\infty$, and $\mathscr{C}(t)=0$ for all $t$.
Finally $\exp(-\infty):=0$.

\subsubsection*{G5.1.5. Measures, partition functions, free ratio}

The objects $F_N$, $\gamma_N$, $A(F^{\mathrm{u}})$, $H_\theta$, $\mathrm{dens}^{\theta,a}_N$, $\nu^{\theta,a}$, $\nu^\theta$, $\bar{\nu}^a$, $\bar{\nu}$ and $\bar{\mu}$ are those of Section~I.7.3 on $C(n_*)$, formed with the pair $P_N(e)$.
The slice $\mathrm{Sl}$, the fields $A_j$, $A_*$, $\zeta_j$, the spaces $V_j$ with the bases $\{\Upsilon\}$, the law $\mathrm{d}\mathsf{m}_*$, the function $\chi_*(A)=\mathbf{1}\left[\,|dA|_\infty\le p_*\right]$ and the sets $\mathfrak{G}$, $\mathfrak{G}^-$ are the photon objects of \cite[Definition~5.11]{QEDI} (Section~I.5.7).
The space of the coclosed real bond functions on $T^0_*$ with mean zero is denoted $V_{\mathrm{t}}\left(T^0_*\right)$.
Further, $c_\theta:=c^{(k_*)}_\theta$ (Section~I.6.6) is the constant bond function on $T_*$ and on $T^0_*$ with the components $\theta_\mu/\left(e_*\cdot\ell_*\right)$, $\theta\in\mathbb{R}^3$; in Proposition~34.4, $\theta\in\mathbb{C}^3$.
Let
\begin{gather*}
\mathring{\Phi}(t):=\frac{\bar{\nu}^{a(t)}(F_N)}{\bar{\nu}(F_N)},\qquad N_0(t):=\frac{\det\left(D_{(t/e)\hat{\mathfrak{a}}}+\bar{m}\cdot\eta\right)}{\det\left(D_0+\bar{m}\cdot\eta\right)}\quad(t\in\mathbb{C}),\\
c_m:=N_0^{(m)}(0),\qquad \Theta^{(0)}_m:=\left|\left.\prod_{\mathfrak{c}}\left(\partial_{t_{\mathfrak{c}}}^m-c_m\right)\mathring{\Phi}(t)\right|_{t=0}\right|^{\frac{1}{N_*}}\qquad(m=1,2,3,4).
\end{gather*}
Here $\mathring{\Phi}(t)=\bar{Z}^{\mathrm{ext}}(t)/\bar{Z}^{\mathrm{ext}}(0)$, with $\bar{Z}^{\mathrm{ext}}$ of Section~I.7.6 formed with the placement at hand, and the determinants are those of operators on $\mathfrak{S}^{\mathrm{ch}}(1;n_*)$, of the size $192\,n_*^3$.
For $A\in\mathrm{Dom}_*$ and $\mathscr{C}\in\mathfrak{X}_*$ let
\[
Y^x_*(A;\mathscr{C}):=y^x_{k_*}\left[\varepsilon^{\mathrm{v},t}_*,m_*,\mathscr{E}_*\right](A;\mathscr{C})
\]
(Section~I.5.7, formed with $\mathfrak{P}_*$).

\subsubsection*{G5.1.6. Jets; M\"obius components; two norms on $\mathrm{Gr}[\Psi]$}

Let
\[
\mathcal{J}:=\mathbb{C}\left[t_{\mathfrak{c}}:\ \mathfrak{c}\in\mathrm{Cel}\right]\big/\left(t_{\mathfrak{c}}^5:\ \mathfrak{c}\in\mathrm{Cel}\right),
\]
with the norm
\[
\|P\|_{\mathrm{rad}}:=\sum_{\mathrm{ak}}\left|P_{\mathrm{ak}}\right|\cdot\mathrm{rad}^{|\mathrm{ak}|},
\]
the sum ranging over the multi-indices $\mathrm{ak}$ with $\mathrm{ak}_{\mathfrak{c}}\le4$; one has $\|P\cdot P'\|_{\mathrm{rad}}\le\|P\|_{\mathrm{rad}}\cdot\|P'\|_{\mathrm{rad}}$.
A function that is analytic near $t=0$ is identified with its image in $\mathcal{J}$, namely its Taylor polynomial with $\mathrm{ak}_{\mathfrak{c}}\le4$ for all $\mathfrak{c}$.
This identification is a homomorphism of rings, and it maps $\exp(\mathrm{gf})$ to the exponential series of the image of $\mathrm{gf}$.

For a complex function $\mathrm{gf}$ on $\mathrm{DD}(\mathrm{rad})$, a subset $\mathrm{Scl}'$ of $\mathrm{Cel}$ and $t\in\mathrm{DD}(\mathrm{rad})$, let $t^{\mathrm{Scl}'}$ be the point with the components $t_{\mathfrak{c}}$ for $\mathfrak{c}\in\mathrm{Scl}'$ and $0$ for the other cells.
For a subset $\mathrm{Scl}$ of $\mathrm{Cel}$ the M\"obius component is
\[
\mathrm{gf}_{\mathrm{Scl}}(t):=\sum_{\mathrm{Scl}'\subset\mathrm{Scl}}(-1)^{\#|\mathrm{Scl}\setminus\mathrm{Scl}'|}\cdot\mathrm{gf}\left(t^{\mathrm{Scl}'}\right).
\]

The Grassmann algebra generated by the generators $\Psi(u)$, $u=(y,\beta,\omega)$, $y\in T^0_*$ (weight $\nu=1$; Section~I.3.6), is denoted $\mathrm{Gr}[\Psi]$, and its norm with the one weight $\mathrm{aw}>0$ (Section~I.5.1) is denoted $|F|_{\mathrm{aw}}$.
Put $d(u,u'):=d(y,y')$ and, for a set $P$ of unit sites, $d(u,P):=d(P,u):=d(y,P)$ (Section~G5.1.4).
An even element without constant term is written
\[
\mathrm{Fv}=\sum_{m\ge2}\frac{1}{m!}\sum_{u_1,\dots,u_m}\mathrm{Fv}_m(u_1,\dots,u_m)\cdot\Psi(u_1)\cdots\Psi(u_m),
\]
with antisymmetric kernels $\mathrm{Fv}_m$.
A weight is a symmetric function $\mathrm{wg}\ge1$ of the finite families of indices $u$.
For a weight $\mathrm{wg}$, with the number $h_G$ of Section~G5.1.8, let
\begin{align*}
\mathrm{rw}_G(\mathrm{Fv};\mathrm{wg})&:=\sum_m\frac{h_G^m}{(m-1)!}\,\sup_{u_1}\sum_{u_2,\dots,u_m}\left|\mathrm{Fv}_m(u_1,\dots,u_m)\right|\cdot\mathrm{wg}(u_1,\dots,u_m),\\
\mathrm{Nw}_G(\mathrm{Fv};\mathrm{wg})&:=\sum_m\frac{h_G^m}{(m-1)!}\sum_{u_1,\dots,u_m}\left|\mathrm{Fv}_m(u_1,\dots,u_m)\right|\cdot\mathrm{wg}(u_1,\dots,u_m).
\end{align*}
The weights used are
\[
\mathrm{wg}_\delta(u_1,\dots,u_m):=\exp\left(\mathrm{vt}\cdot\max_{i,j}d(u_i,u_j)\right),\qquad \mathrm{wg}_P(u_1,\dots,u_m):=\exp\left(\mathrm{vt}\cdot\max_i d(u_i,P)\right),
\]
the latter for a non-empty set $P$ of unit sites, and the weight $1$, the constant function one.

\subsubsection*{G5.1.7. Operators at the stop level; majorants; the domains $\mathrm{Gwk}$, $\mathrm{Gwk}^+$; legs}

In this section the level is $k_*$, the side $\ell_*$, the mass $\mu_*$ and the charge $e_*$; $A$ is a complex bond function on $T_*$, and $x=(c,g)$ is a fermion site of $T^{\mathrm{f}}_*$.

\emph{Phases.}
The phases $u_A$, $U_A$ and the number $\mathfrak{i}(A)$ are defined in (I.1.6), and $\varphi_A(y,c)$ in (I.3.3).
Put
\begin{align*}
\mathrm{phx}_A(y,c)&:=e_*\cdot\eta_*\cdot\left(\tau_{k_*}A\right)(y,c)&&\left(c\in\Delta_y\right),\\
\mathrm{ups}_A(c,c')&:=e_*\cdot\eta_*\cdot(\tau A)(c,c')&&\left(|c-c'|\le M\right),
\end{align*}
so that $\varphi_A=\exp\left(i\,\mathrm{phx}_A\right)$ and $U_A=\exp\left(i\,\mathrm{ups}_A\right)$.
Both $\mathrm{phx}_A$ and $\mathrm{ups}_A$ are linear in $A$.
The function $\mathrm{phx}_A(y,\cdot)$ depends only on the values of $A$ on the bonds with both end points in $\Delta_y$, and $\left|\mathrm{phx}_A(y,c)\right|\le\frac{3}{2}\,e_*\cdot|A|_\infty$.
The number $\mathrm{ups}_A(c,c')$ depends only on the values of $A$ on the bonds in the lattice box spanned by $c$ and $c'$; moreover $\left|\mathrm{ups}_A(c,c')\right|\le3\,|c-c'|\cdot e_*\cdot|A|_\infty$ and $\mathrm{ups}_A(c,c)=0$.
These facts on the photon objects are proved in Chapter~33, Section~0.4, (Ph1)--(Ph3).

\emph{Operators} (Sections~I.2.3, I.3.3 and~I.3.5 at $k=k_*$, $\mu=\mu_*$).
Besides $D_A$ and the averaging operators $Q_*(A)J^*$, $JQ_*^T(-A)$, $JP_*(A)J^*$ one uses
\begin{align*}
\mathrm{Fst}(A)&:=D_A+\mu_*+b_*\cdot JP_*(A)J^*\qquad\text{on}\ \mathfrak{S}^{\mathrm{ch}}(\eta_*,\ell_*),\\
S_*(A)&:=\mathrm{Fst}(A)^{-1}\qquad\text{wherever the inverse exists},\\
\mathrm{Sa}_*(A)&:=\delta_{\alpha,A}\cdot S_*(A),\ \text{restricted to the short pairs}\ \left(d_X<1\right),\\
H_*(A)&:=b_*\cdot S_*(A)\cdot JQ_*^T(-A),\qquad \bar{H}_*(A):=\mathsf{C}^{-1}\cdot H_*(-A)\cdot\mathsf{C},\\
D_*(A)&:=b_*-b_*^2\cdot Q_*(A)J^*\cdot S_*(A)\cdot JQ_*^T(-A)=b_*-b_*\cdot Q_*(A)J^*\cdot H_*(A)\qquad\text{on}\ \mathfrak{S}(1,\ell_*),
\end{align*}
where $\delta_{\alpha,A}$ is the H\"older quotient (I.1.7).
The only hop that carries a phase has the matrix $M_3(x)=3\left(\rho-\gamma\cdot(ge_1)\right)$, and by (I.m6), (I.m7)
\[
\left|M_3(x)\cdot v\right|_c\le3\,(1+\rho)\cdot|v|_c\le6\,|v|_c,\qquad \left|M_3(x)\right|_2=3\,(1+\rho)\le6.
\]

\emph{Majorants.}
A spinor on $T^0_*$ is localised at the sites $y$, a spinor on $T^{\mathrm{f}}_*$ in the sets $\Delta_y\times O_{\mathrm{cub}}$, and a function of the pairs $(x,x')$ of $\mathrm{Pz}_*$ in the set of the first point; $\mathbf{1}_y$ denotes the corresponding indicator.
All sup norms are formed with $|\cdot|_c$ (Section~I.1.5).
For an operator $O$ with values in the functions of pairs, the majorant of $O$ for the short pairs is the majorant of the operator that maps $f$ to the restriction of $Of$ to the pairs $(x,x')$ of $\mathrm{Pz}_*$ with $d_X(x,x')<1$.
For a linear operator $O$ between such spaces and $\gamma\ge0$ let
\begin{align*}
|O|(y,y')&:=\sup\left\{\left|\mathbf{1}_y\cdot O\cdot\mathbf{1}_{y'}\cdot f\right|_\infty:\ |f|_\infty\le1\right\},\\
\|O\|_\gamma&:=\max\left\{\sup_y\sum_{y'}|O|(y,y')\cdot\exp\left(\gamma\cdot d(y,y')\right),\ \sup_{y'}\sum_y|O|(y,y')\cdot\exp\left(\gamma\cdot d(y,y')\right)\right\},
\end{align*}
and put
\[
\mathrm{Edy}(y,y'):=\exp\left(-\gamma^F\cdot d(y,y')\right).
\]
Two rules hold.
\begin{itemize}
\item[(M1)] $\left|O_1\cdot O_2\right|(y,y')\le\sum_z\left|O_1\right|(y,z)\cdot\left|O_2\right|(z,y')$ and $\left\|O_1\cdot O_2\right\|_\gamma\le\left\|O_1\right\|_\gamma\cdot\left\|O_2\right\|_\gamma$.
\item[(M2)] Let $O$ be an operator on $\mathfrak{S}(1,\ell_*)$ with the $4\times4$ blocks $O(y,y')$. Then $\left|O(y,y')\right|_c\le|O|(y,y')$; the $\ell^2$ operator norm satisfies $|O|_{\ell^2}\le2\,\|O\|_0$; and the trace norm satisfies $|O|_1\le4\sum_{y,y'}|O|(y,y')$.
\end{itemize}
Proofs of (M1) and (M2) are given in Chapter~33, Section~0.6; a proof of (M2) is also given in Chapter~34, (P3)\,(a) and (P4)\,(e).
They are not the conditions (M1), (M2) of Proposition~14.2\,(b).

For an analytic function $O(\cdot)$ of the connection with values in operators put
\begin{align*}
\left(\partial_BO\right)(A)&:=\left.\frac{\mathrm{d}}{\mathrm{d}\,\mathrm{sv}}O(A+\mathrm{sv}\cdot B)\right|_{\mathrm{sv}=0},\\
\left(\partial_B\partial_{B'}O\right)(A)&:=\left(\partial_B\left(\partial_{B'}O\right)\right)(A)=\left.\frac{\mathrm{d}}{\mathrm{d}\,\mathrm{sv}}\frac{\mathrm{d}}{\mathrm{d}\,\mathrm{sv}'}O\left(A+\mathrm{sv}\cdot B+\mathrm{sv}'\cdot B'\right)\right|_{\mathrm{sv}=\mathrm{sv}'=0};
\end{align*}
the second derivative is symmetric in $(B,B')$.
For a complex bond function $B$ on $T_*$ let
\begin{align*}
\epsilon_B&:=e_*\cdot|B|_\infty,\\
\mathrm{Ypr}_B&:=\left\{y\in T^0_*:\ \Delta_y\ \text{contains an end point of a bond}\ b\ \text{with}\ B(b)\ne0\right\},\\
\mathrm{Ypr}_B^+&:=\left\{y:\ d\left(y,\mathrm{Ypr}_B\right)\le M+1\right\}.
\end{align*}
For $B=\hat{c}_{\mathfrak{c}}$ these are $\epsilon_{\mathfrak{c}}$, $\mathrm{Ypr}_{\mathfrak{c}}$ and $\mathrm{Ypr}_{\mathfrak{c}}^+$.

\emph{The domains $\mathrm{Gwk}$ and $\mathrm{Gwk}^+$.}
\begin{align*}
\mathrm{Gwk}&:=\left\{A_r+\mathscr{C}:\ A_r\ \text{real and in}\ \tilde{R}_*,\ \mathscr{C}\in\mathfrak{X}_*\right\},\\
\mathrm{Gwk}^+&:=\left\{A+A_1:\ A\in\mathrm{Gwk},\ A_1\ \text{a complex bond function on}\ T_*\ \text{with}\ e_*\cdot|A_1|_\infty\le c_\tau\right\}.
\end{align*}

\emph{Legs.}
A leg index is either $\mathrm{up}=\xi=(x,\beta,\omega)$ with $x\in T^{\mathrm{f}}_*$, or $\mathrm{up}=\boldsymbol{\zeta}=(x,x',\beta,\omega)$ with $(x,x')\in\mathrm{Pz}_*$ (Section~I.3.6).
Put $\mathrm{pos}(\xi):=\{[x]\}$, and
\[
\mathrm{pos}(\boldsymbol{\zeta}):=\{[x]\}\ \text{if}\ d_X(x,x')<1,\qquad \mathrm{pos}(\boldsymbol{\zeta}):=\{[x],[x']\}\ \text{if}\ d_X(x,x')\ge1,
\]
where $[x]=[c]$ (I.POS).
For $A\in\tilde{G}_*$ and $u=(y',\beta',\omega')$, $\mathrm{Leg}_A(\mathrm{up};u)$ is the coefficient of $\Psi(u)$ in $\psi(\xi)$, respectively in $\chi(\boldsymbol{\zeta})$, at
\[
\psi=H_*(A)\Psi,\qquad \bar{\psi}=\bar{H}_*(A)\bar{\Psi},\qquad \chi=\delta_{\alpha,A}\psi,\qquad \bar{\chi}=\delta_{\alpha,-A}\bar{\psi};
\]
it vanishes for $\omega'\ne\omega$.
Put
\[
\left|\mathrm{Leg}_A\right|(\mathrm{up};y'):=\sum_{\beta'}\left|\mathrm{Leg}_A\left(\mathrm{up};(y',\beta',\omega)\right)\right|,
\]
and define the same quantity in the same way for the derivatives of $\mathrm{Leg}_A$ in $A$.

\subsubsection*{G5.1.8. Constants}

\emph{Constants of Section~I.8.2.}
The following constants of Section~I.8.2 are used: $K_0$, $\kappa_0$, $c_X$, $c_{\mathrm{rb}}$; $b=1$ and $b_-=1-1/L$; $\kappa$; $L$; $N_*$, $a_L$; $M$; $C^F\ge1$ and $\gamma^F\in(0,1]$ with $\gamma^F\cdot M\ge16$; $C_0\ge\max\left(1,C^F\right)$; $C_d$; $r$, $p_0$, $p$, $\varepsilon$, $s$.

\emph{The line (C6'), which follows (C6).}
Its entries are functions of the constants of the lines (C0)--(C6); they do not depend on $\bar{m}$, $\beta$, $e$, $N$, $\theta$.
The table also contains $\mu_-$, which stands in the line (C4).
Its last column gives the corresponding definition in \cite[Section~20.1.7]{QEDI}, where the constants $C^F$, $\gamma^F$ are written $C_F$, $\gamma_F$.

\begin{tabular}{p{0.2\textwidth}p{0.46\textwidth}p{0.26\textwidth}}
constant & definition & in \cite[Section~20.1.7]{QEDI}\\[2pt]
$\mu_-$ (line (C4)) & $L^{-3}$ & $L^{-m_0-3}$ (defined in \cite[Section~20.1.2]{QEDI})\\
$\mathrm{ssum}(\gamma)$ & $\left(12/\gamma\right)^3$ for $0<\gamma\le6$ & the same\\
$\mathbf{S}$; $\mathrm{hL}$ & $C^F\cdot\mathrm{ssum}\left(\gamma^F/2\right)$; $6\,\mathbf{S}$ & the same formulas\\
$C_K$ & $226\,\mathbf{S}^2$ & $100\,\mathbf{S}^2$\\
$C_K'$ & $1328\,\mathbf{S}^2$ & $400\,\mathbf{S}^2$\\
$C_H$ & $203\,M\cdot\left(C^F\right)^2$ & $100\,M\cdot C_F^2$\\
$\mathrm{CVj}$ & $10^3\,C_H^2\cdot\mathrm{ssum}\left(\gamma^F/2\right)^4\cdot K_0$ & $C_V$, the same formula\\
$\theta_{\mathrm{CT}}$; $\mathrm{vt}$ & $\mu_-/(12\,\mathrm{e})$; $\min\left\{\theta_{\mathrm{CT}}/2,\ \gamma^F/8,\ 1/M\right\}$ & the same\\
$\mathrm{gGa}$; $h_G$; $\mathrm{gGa}_1$ & $\left(b_-^{-1}+\mu_-^{-1}\right)^{\frac{1}{2}}$; $1+\mathrm{gGa}$; $b_-^{-1}+4/\mu_-$ & $g$, $h_G$, $g_1$\\
$\mathrm{alG}$ & $2\,b_-^{-1}+4\,\mathrm{e}\cdot\mu_-^{-1}\cdot\mathrm{ssum}\left(\theta_{\mathrm{CT}}/2\right)$ & $\boldsymbol{\alpha}$\\
$c_\tau$ & $\min\left\{1/8,\ \mu_-/48,\ 1/(292\,\mathbf{S})\right\}$ & $c_1$, with $1/(80\,\mathbf{S})$ in place of $1/(292\,\mathbf{S})$\\
$\mathrm{CWj}$ & $3\,\mathrm{CVj}\cdot h_G^2+2\,\mathrm{alG}\cdot\mathrm{CVj}^2\cdot h_G^4+4\,\mathrm{e}^6\cdot C_d\cdot M^3\cdot K_0$ & $C_W$, the same formula
\end{tabular}

For $0<\gamma\le6$, on every torus $(\mathbb{Z}/\ell)^3$ and for every site $y$,
\[
\sum_{y'}\exp\left(-\gamma\cdot d(y,y')\right)\le\mathrm{ssum}(\gamma).
\]
This is the bound $c(\gamma)\le(12/\gamma)^3$, $\gamma\le6$, of \cite[Section~13.1.6]{QEDI}; a proof is also given in the proof of Proposition~35.4\,(SF2).

\emph{The line (C10).}
After $\bar{m}$ and $\beta$ this line contains $D_{\mathrm{be}}$, $D_\beta$, $C_\beta$, $\mathbf{D}$, $\mathrm{Cb}$ (Section~G5.1.3) and
\[
C_{\mathrm{fr}}:=2592\,C_d\cdot K_0\cdot C_\beta^2\cdot\left(M^2\cdot D_\beta^3\cdot L^4\cdot h^{-2}+8\,M^5\cdot h\right).
\]

\emph{Quantities that depend on $e$.}
With the exception of $\mathrm{mxL}$, the following numbers depend on $e$, through $e$, $e_*$, $a_*$, $\ell_*$, $r_*$, $p_*$ and $p_{0*}$; by Proposition~35.1\,(c) they do not depend on $N$, with the exception of $\varpi_{\mathrm{sep}}$.
\begin{align*}
\tau_s&:=c_\tau\cdot\ell_*;\qquad \mathrm{pp},\ t_x,\ t_x^+,\ \mathbf{r},\ q_x\ \ \text{(Section~G5.1.4)};\qquad \bar{\kappa}_1:=6\,e_*^{\frac{1}{4}-9\varepsilon};\\
\bar{r}&:=4\,h_G^2\cdot C_K\cdot e_*^{2\varepsilon}+650\,K_0\cdot e_*^{\frac{1}{4}-7\varepsilon};\\
B_1&:=16\,\mathrm{gGa}_1\cdot C_K'\cdot\mathrm{pp}\cdot e_*^{2\varepsilon},\qquad B_{\mathrm{wr}}:=16\,\mathrm{e}^6\cdot C_d\cdot M^3\cdot K_0\cdot\mathrm{pp}\cdot\exp\left(-r_*\right),\\
B_0&:=C_{\mathrm{fr}}\cdot e_*^{4\varepsilon}+B_{\mathrm{wr}}+B_1;\\
\varpi_{\mathrm{sep}}&:=2^{N_*}\cdot\mathrm{CWj}\cdot\mathrm{pp}\cdot\exp\left(-\mathrm{vt}\cdot\left(d_{\mathrm{sep}}-4M-6\right)\right),\\
\varpi_{\mathrm{sep}}^+&:=2^{N_*}\cdot\mathrm{CWj}\cdot\mathrm{pp}\cdot\exp\left(-\mathrm{vt}\cdot\left(\frac{M\cdot r_*}{8L}-4M-7\right)\right);\\
\varpi_0&:=18\,\ell_*^3\cdot\exp\left(-\frac{p_{0*}}{2}\right)+6\,\ell_*^3\cdot\exp\left(-\frac{\left(p_*-1\right)^2}{2\,C_0^2}\right);\\
t_2&:=e\cdot\min\left\{N_*^{-\frac{1}{2}},\ \frac{1}{C_\beta},\ \frac{1}{2\sqrt{N_*\cdot\mathbf{D}}}\right\};\\
\varpi_L&:=\varpi_0^{\frac{1}{2}}\cdot\exp\left(\frac{\mathbf{D}}{2}\right),\qquad \mathrm{mxL}:=2+\exp\left(\frac{\mathbf{D}}{2}\right),\\
\lambda_L&:=\left(\varpi_L^{\frac{1}{3}}\cdot\mathrm{mxL}^{\frac{2}{3}}\right)^{\frac{1}{N_*}}\cdot5\cdot\left(\frac{2N_*\cdot\mathbf{r}}{t_2}\right)^4.
\end{align*}
The number $\varpi_{\mathrm{sep}}$ depends on $N$ and on the placement through $d_{\mathrm{sep}}$, and $\varpi_{\mathrm{sep}}=0$ if $\hat{\mathfrak{a}}=0$.
The number $\mathrm{mxL}$ is a function of the constant $\mathbf{D}$ of the line (C10) alone; hence it depends only on $\beta$, $h$ and $L$, and neither on $e$ nor on $N$, the torus or the placement.
It is listed here only because $\lambda_L$ is formed with it.
In the minimum that defines $t_2$ an entry with the denominator $0$ is read as $+\infty$.

The number $\mathrm{pp}$ takes two values, since $a_*\in\left\{h/L,\,h/L^2\right\}$, and
\[
(2M+4)^3\le\mathrm{pp}\le\left(\frac{D_\beta\cdot L^2}{h}+2M+4\right)^3.
\]

\subsubsection*{G5.1.9. The conditions (g1)--(g9); the number $e_{\sharp}(\beta)$}

The conditions (g1)--(g9) are the following inequalities.
\begin{align}
&5\,\mathrm{hL}\cdot h_{G}\le e_{*}^{-\frac{1}{4}+\varepsilon},\qquad
\left(4+\frac{M^3}{2}\right)\cdot e_{*}^{\varepsilon}\le1,\qquad
\mathrm{alG}\cdot\bar{r}\le\frac{1}{2};\tag{g1}\\
&\frac{M\cdot r_{*}}{8L}-1\ge6M+8;\tag{g2}\\
&4\,C_{0}\cdot\ell_*^3\cdot\mathrm{gGa}_{1}\le e_{*}^{-\frac{1}{4}+\varepsilon},\qquad
4\,\ell_*^3\cdot\left(2C_K+K_{0}+4\right)\le e_{*}^{-2\varepsilon},\qquad
e_{*}^{\frac{1}{4}-9\varepsilon}\le\frac{1}{8};\tag{g3}\\
&\begin{aligned}[t]
&48\,e_{*}^{2\varepsilon}\le\mu_-,\qquad
2\,\mathrm{gGa}_{1}\cdot C_K'\cdot e_{*}^{2\varepsilon}\le\frac{1}{2},\qquad
\tau_s\ge1,\\
&52\,B_1\cdot\exp(-\tau_s)\le e_{*}^{\frac{1}{4}-9\varepsilon},\qquad
B_0\le\frac{1}{4};
\end{aligned}\tag{g4}\\
&2\,\varpi_{\mathrm{sep}}^{+}\le\bar{\kappa}_1^{N_{*}};\tag{g5}\\
&N_{*}\cdot\bar{\kappa}_1^{1-s}\le\frac{1}{2};\tag{g6}\\
&\varpi_0\le\frac{1}{4},\qquad8\,N_{*}\cdot t_2\le t_x;\tag{g7}\\
&\lambda_L\le\bar{\kappa}_1^{s};\tag{g8}\\
&\frac{D_\beta}{2}+\frac{1}{2}\le\frac{3}{8}\,\ell_c,\qquad
D_\beta+3\le\ell_c\qquad
\left(\ell_c=\frac{h\cdot M\cdot r_K(e)}{2L^2}\right).\tag{g9}
\end{align}
Each of the conditions (g1)--(g9) is an inequality between numbers which are functions of $e$, of $\bar{m}$, of $C_{\beta}$, $\mathbf{D}$, $D_\beta$, and of the constants of the lines (C0)--(C9) and (C6') of Section~I.8.2.
The conditions (Z9.s1), (Z9.s2), (Z9.s3) of Section~I.8.4 are, in this order, the second condition in (g1), the first condition in (g3) and the first condition in (g4).
The two inequalities of (g9) appear in the proof of \cite[Lemma~20.33]{QEDI}.

Let $e_{\sharp}(\beta)$ be one half of the supremum of the numbers $e'\in\left(0,e_1\cdot h^{-\frac{1}{2}}\right]$ with the property that (g1)--(g9) hold for every $e\in(0,e']$;
if there is no such $e'$, put $e_{\sharp}(\beta):=0$.
The number $e_{\sharp}(\beta)$ is a constant of the line (C10) of Section~I.8.2.
By its definition, the conditions (g1)--(g9) hold for every $e\in\left(0,e_{\sharp}(\beta)\right]$.
For such a charge $e$ and every $N\ge N_0(e)$ one has $\eta\le1$, so that (g9) yields
\[
\frac{D_\beta}{2}+\frac{\eta}{2}\le\frac{3}{8}\,\ell_c,\qquad
D_\beta+3\eta\le\ell_c;
\]
consequently every centred placement is admitted (Section~G5.1.3), and $\beta$ fits $\left(N,T^{*}(e)\right)$ in the sense of Section~I.7.6.
The inequality $e_{\sharp}(\beta)>0$ is Proposition~35.10.

The following conditions enter the Propositions~33.1--35.10.
All entries of the list but the last are conditions of Section~I.8.4, and both kinds of conditions of Section~I.8.1 occur among them;
the four inequalities of the last entry are consequences of the lines (C2), (C3), (C5) of Section~I.8.2.
The condition (2.1), the first two clauses of (T1) and the condition (T3) are conditions on the data $(e,T,N)$, here at $R=R^{*}(e)$ and $N\ge N_0(e)$, and not conditions on the constants;
they hold whenever $e\cdot h^{\frac{1}{2}}\le e_1$, by the rows (e1) and (e4) of Proposition~14.2\,(b) (Proposition~14.2\,(c) and Proposition~35.1\,(e)).
The other entries, namely (A8), (A10), the conditions (A1) and (A2) with $C_{P}\le C_{w}$ and $\gamma_{\mathrm{w}}\le\gamma_{P}$, and the last entry, are conditions on the constants.
\begin{itemize}
\item The condition (A8); in particular the first, the third and the fourth of the four inequalities of its row in Proposition~14.2\,(b), namely
\[
e_0^{\varepsilon}\le\frac{1}{4},\qquad
3M\cdot e_0^{\varepsilon}\le1,\qquad
e_0^{\frac{1}{4}}+e_0^{\varepsilon}\le1.
\]
As $e_{*}\le e_0$ (Proposition~35.1\,(a)), the first two of them give $e_{*}^{2\varepsilon}\le\frac{1}{16}$ and $3M\cdot e_{*}^{2\varepsilon}\le\frac{1}{4}$.
The inequality $e_0^{\frac{1}{4}}+e_0^{\varepsilon}\le1$, the last one displayed and the fourth of the row of (A8), is the inequality on which the proof of \cite[Lemma~17.4(a)]{QEDI} rests.
\item The condition (2.1), that is, $r_{*}>31L+16$; it is the condition (G4.2.1) of Section~I.8.4.
\item The first two clauses of (T1): $v_{*}\ge8$ and $p_0\cdot v_{*}^{p_0-1}\ge28$.
\item The condition (T3): $4\,C_{0}\cdot\ell_*^2\cdot p_*<e_{*}^{-\varepsilon}$.
\item The condition (A10), that is, the four inequalities of its row in Proposition~14.2\,(b), which are imposed for all $v\ge-\log e_0$:
\begin{gather*}
v\ge9,\qquad
8M\cdot\left(\frac{v}{4}+4\right)\cdot\exp(-\varepsilon v)\le1,\\
8\,C_{B}\cdot v^{p}\le\exp\left(\left(\frac{3}{4}-3\varepsilon\right)\cdot v\right),\qquad
8\exp\left(-\left(\frac{3}{4}-\frac{7\varepsilon}{2}\right)\cdot v\right)\le1 .
\end{gather*}
Since $v_{*}\ge-\log e_0$, the third and the fourth of them hold at $v=v_{*}$, where they read
\[
8\,C_{B}\cdot p_*\le e_{*}^{-\frac{3}{4}+3\varepsilon},\qquad
8\,e_{*}^{\frac{3}{4}-\frac{7\varepsilon}{2}}\le1 ;
\]
these are the two displayed hypotheses of \cite[Proposition~15.30]{QEDI} at $e_k=e_{*}$.
Of the condition (A10), the Propositions~33.1--35.10 require only these two inequalities, in the clause (SF5) of Proposition~35.4.
The first two inequalities of (A10) are conditions of the clause (d4) of Proposition~13, a clause which these propositions do not require.
Of the conditions on the constants listed after Proposition~13, the clauses (d0)--(d3) of that proposition, on which the clause (SF4) of Proposition~35.4 rests, require the following:
$\rho<1$ (line (C0));
the conditions (L5), (M1), (Z2.s1) of Proposition~14.2\,(b);
the further conditions of Proposition~2.8, namely (M2), the part $M\ge108\pi\cdot C'_{V}\cdot a'$ of (M5), (Z2.s2) and $\alpha<1$ (the range of the line (C1)), which enter the property W1 and the clause (d0) through Proposition~2.8\,(a), (b), enter (d0) also through Propositions~12 and~10\,(e), (g), and enter (d3) through Proposition~10\,(b), (g);
and $\kappa_{\mathrm{w}}-1\ge\kappa+1$.
None of these conditions is part of (A10).
\item The conditions (A1) and (A2), together with $C_{P}\le C_{w}$ and $\gamma_{\mathrm{w}}\le\gamma_{P}$.
With $v=v_{*}$ in (A1) and (A2), these conditions and the two inequalities displayed at the end of the preceding entry are hypotheses of \cite[Proposition~15.30]{QEDI} at $e_k=e_{*}$.
The further hypotheses of that proposition, besides those which its clause~(ii) adds, are
\[
0<\varepsilon\le\bar{\varepsilon},\qquad
L^{\frac{3}{4}}\ge\frac{8}{3},\qquad
0<e_{*}<1 .
\]
The first of them holds by the condition (S1) of Proposition~14.2\,(b), since $\alpha<1$ by the line (C1) of Section~I.8.2;
the second holds by the condition (L5) of Proposition~14.2\,(b), since $L\ge5$ and $5^{3}>\left(\frac{8}{3}\right)^{4}$;
the third holds since $0<e_{*}\le e_0$ (Proposition~35.1\,(a)) and since the first inequality of (A8) gives $e_0<1$.
\item The inequalities
\begin{gather*}
\kappa\ge7\kappa_{0}+7,\qquad
L\cdot\left(\kappa-\kappa_{0}-1\right)\ge\kappa,\\
c_{\mathrm{rb}}\cdot L^3\cdot\exp\left(-\left(L\cdot\left(\kappa-\kappa_{0}-1\right)-\kappa\right)\right)\le\frac{1}{2},\qquad
M\ge L^2.
\end{gather*}
\end{itemize}
The bound in the clause (SF3) of Proposition~35.4 is the display (2.1);
this display is meant in the clause (A6) of that proposition and wherever (2.1) is named together with (SF3),
whereas the condition (2.1), as in Proposition~35.1\,(e), is the inequality $r_{*}>31L+16$ of the list above.
In the Propositions~33.1--35.10, (A6) and (A7) are the last two clauses of Proposition~35.4, and (N0)--(N3), attached to Proposition~35.5, are the clauses of that proposition.

\subsection*{35 (first part). Scales, cells, probes, the inputs at the stop level}

\subsubsection*{Proposition~35.1 (scales on the reference torus)}

\begin{named}{Proposition 35.1}
Let the constants be those of Section~I.8, let $\bar{m}\in(0,1]$, and let $e>0$ satisfy $e\cdot h^{\frac{1}{2}}\le e_1$.
Then $R^{*}(e)\in\mathrm{Tor}(e)$, and the number $N_0(e)=N_{\min}\left(R^{*}(e),e\right)$ is finite;
that is, the set of the integers $N_1$ in the definition of $N_{\min}$ in Section~I.6.2 is not empty.
Let $N\ge N_0(e)$ and $R=R^{*}(e)$.
Then the following hold.
\begin{itemize}
\item[(a)] One has $e\le\frac{1}{3}$ and $v_K\ge8$;
the number $T^{*}(e)$ is an integer with $T^{*}(e)\ge1$;
$k_{*}\in\{K-1,K-2\}$ and $k_{*}\ge1$;
$\ell_*=M\cdot r_{*}$, $r_{*}\in\{r_K,L\cdot r_K\}$ and $a_{*}\in\{h/L,h/L^2\}$;
$\mu_-<\mu_*\le1/L$;
\[
e\cdot\left(\frac{h}{L^2}\right)^{\frac{1}{2}}\le e_{*}=e\cdot a_{*}^{\frac{1}{2}}\le e_K=e\cdot h^{\frac{1}{2}}\le e_1\le e_0;
\]
$b_{-}\le b_*\le1$;
and $\eta_*\le1/L$.
\item[(b)] One has $v_K\le v_{*}\le v_K+\log L$ and $v_{*}^r\le r_{*}<L\cdot v_{*}^r$;
the number $r_{*}$ is a power of $L$, and $r_{*}\ge L^2$;
and $M\cdot v_{*}^r\le\ell_*<L\cdot M\cdot v_{*}^r$.
\item[(c)] The numbers $K-k_{*}$, $a_{*}$, $e_{*}$, $v_{*}$, $r_{*}$, $\ell_*$, $\mu_*$, $p_*=v_{*}^p$, $p_{0*}=v_{*}^{p_0}$, $R^{*}(e)$ and $\ell_c$ do not depend on $N\ge N_0(e)$;
the numbers $k_{*}$, $\eta_*$, $b_*$, $n_*$ and $\mathrm{per}$ depend on $N$.
For an arbitrary integer $N$, the condition $N\ge N_0(e)$ is equivalent to
\[
N\ge0\qquad\text{and}\qquad N+\log_L a_{*}\ge1,
\]
where $a_{*}$ denotes the number which, by the first sentence of (c), is the same for all $N\ge N_0(e)$, and which is therefore defined without reference to $N$;
moreover $k_{*}=N+\log_L a_{*}$.
\item[(d)] For $0\le j\le k_{*}$ the side $\ell_j$ lies in $\mathrm{Cl}$.
The data $\left(k_{*},\ell_*,\mu_*,e_{*}\right)$ are data of the scale-free setting of Section~I.1.8, with $k_{*}\ge1$.
For $0\le j<k_{*}$, the sides $\ell_j$ and $\ell_j/L=\ell_{j+1}$ lie in $\mathrm{Cl}$, $\ell_j\ge M\cdot r_j$, the running mass $\bar{m}\cdot a_j$ lies in $(0,1/L]$, and $e_{j+1}=L^{\frac{1}{2}}\cdot e_j\le e_0$;
at $j=0$ the mass is $\bar{m}\cdot\eta$ and the charge is $q$.
\item[(e)] The condition (2.1), the first two clauses of (T1) and the condition (T3), in the forms written out in Section~G5.1.9, hold.
\item[(f)] One has
\[
\ell_c=\frac{h\cdot M\cdot r_K}{2L^2}\ge\frac{h\cdot M\cdot v_K^r}{2L^2},\qquad
\mathrm{per}=L^{k_{*}-1}\cdot M\cdot r_{*}\ge L^4,\qquad
R^{*}(e)\ge h\cdot L,
\]
and $\bar{m}\cdot R^{*}(e)>r_K\ge v_K$, whence $\exp\left(-\bar{m}\cdot R^{*}(e)\right)<e\cdot h^{\frac{1}{2}}$.
\end{itemize}
\end{named}

\begin{named}{Remark}
Proposition~35.1 corresponds to \cite[Lemma~20.1(a)--(c)]{QEDI},
to the first sentence of the paragraph ``Independence of the bounds'' in the proof of \cite[Theorem~20.29]{QEDI},
and to Steps~2 and~4 of the proof of part~1 of \cite[Theorem~21.3]{QEDI} together with the estimate of $\exp\left(-\bar{m}\cdot R^{*}\right)$ in the proof of its part~2\,(c), both in \cite[Section~21.9]{QEDI}.
It differs from these statements in the following points.
Here $m_0=0$ and $h'=h$, and the upper end of the range of the masses is $1$ (Section~I.6.1);
hence $\mu_-=L^{-3}$ and $\mu_*\le1/L$.
The number $e_1$ is the number of the line (C9) of Section~I.8.2, formed with the constants of Section~I.8.
Clause~(c) includes the numbers $p_*$ and $p_{0*}$;
their independence of $N$ is not asserted in \cite[Proposition~19.1(c)]{QEDI}, and it follows from that of $e_{*}$, since both are functions of $e_{*}$.
The assertion of \cite[Lemma~20.1(c)]{QEDI} on the systems of primitives is contained in (SF4) of Proposition~35.4.
\end{named}

\subsubsection*{Proposition~35.2 (cells and the groups of the scheme)}

\begin{named}{Proposition 35.2}
Assume the setting (ST*)$^{\mathrm{ch}}$ of Section~G5.1.1, and let $\mathfrak{c}\in\mathrm{Cel}$.
\begin{itemize}
\item[(a)] One has $\mathsf{g}_{\mathfrak{c}}y=r_{\mathfrak{c}}y+a_{\mathfrak{c}}$, where $r_{\mathfrak{c}}\in\{\pm1\}^3$ has the diagonal entries $(r_{\mathfrak{c}})_{\mu\mu}=(-1)^{\mathfrak{c}_\mu}$, and where $a_{\mathfrak{c}}\in\mathrm{per}\cdot\mathbb{Z}^3$ has the components
\[
(a_{\mathfrak{c}})_\mu=i\cdot\mathrm{per}\quad\text{if }\mathfrak{c}_\mu=2i,\qquad
(a_{\mathfrak{c}})_\mu=(i+1)\cdot\mathrm{per}\quad\text{if }\mathfrak{c}_\mu=2i+1.
\]
Here the index $\mathfrak{c}$ is read through a representative in $\mathbb{Z}^3$.
Another representative changes $a_{\mathfrak{c}}$ by a vector of $n_*\mathbb{Z}^3$ and leaves $r_{\mathfrak{c}}$ unchanged;
thus $a_{\mathfrak{c}}$ is a class in $\mathrm{per}\cdot\mathbb{Z}^3/n_*\mathbb{Z}^3$, and in (b) the vector $L^{-j}\cdot a_{\mathfrak{c}}$ is a class in $L\cdot M\cdot\mathbb{Z}^3/\ell_j\mathbb{Z}^3$, which is the form in which the translation part enters the groups of Section~I.4.2.
\item[(b)] Let $0\le j\le k_{*}$.
In the coordinates of level $j$ the map $\mathsf{g}_{\mathfrak{c}}$ is $c\mapsto r_{\mathfrak{c}}c+L^{-j}\cdot a_{\mathfrak{c}}$, and $L^{-j}\cdot a_{\mathfrak{c}}\in L\cdot M\cdot\mathbb{Z}^3$.
For every $S\in\mathrm{Sp}$ with $r(S)=r_{\mathfrak{c}}$ (there are four such elements, and two of them lie in $\mathrm{Sp}_3$), the pair $\tilde{g}:=\left(S,L^{-j}\cdot a_{\mathfrak{c}}\right)$ is an element of $\mathbb{G}_j^{L}$.
It acts on photon sites, on cells and on polymers as $\mathsf{g}_{\mathfrak{c}}$;
on fermion sites it acts by
\[
\tilde{g}(c,g)=\left(r_{\mathfrak{c}}c+L^{-j}\cdot a_{\mathfrak{c}},\ r_{\mathfrak{c}}g\right)
\]
(see (I.4.2));
and $\tilde{g}A=(\mathsf{g}_{\mathfrak{c}})_*A$ for every bond function $A$ on $T_j$.
\end{itemize}
\end{named}

\begin{named}{Remark}
Proposition~35.2 corresponds to \cite[Lemma~20.2]{QEDI}.
In addition, it describes the action of $\tilde{g}$ on the fermion sites, and it gives the number of the elements $S\in\mathrm{Sp}$ with $r(S)=r_{\mathfrak{c}}$;
this number is four because the kernel of $r$ has four elements (Section~I.4.1).
\end{named}

\subsubsection*{Proposition~35.3 (probes)}

\begin{named}{Proposition 35.3}
Clause~(b1) concerns only a probe, a number $\eta\in(0,1]$ and a number $R>0$ such that $R/\eta$ is an integer $\ge2$;
clause~(c1) concerns only $\bar{m}$ and a probe $\beta$.
In the other clauses the setting (ST*)$^{\mathrm{ch}}$ is assumed, $\beta$ is a probe, and $v\in\eta\mathbb{Z}^3$;
from (b2) on, $v$ is an admitted placement.
A function on $\mathbb{R}^3$ is evaluated at a point of the torus through its $R^{*}(e)\mathbb{Z}^3$-periodisation.
\begin{itemize}
\item[(a)] The cochain $\hat{\mathfrak{a}}$ and all the cochains $\hat{\mathfrak{a}}_{\mathfrak{c}}$ lie in ${V_{\mathrm{t}}}$, the orthogonal complement of $\ker d$ in $C^1(n_*)$.
\item[(b1)] For every probe $\beta'$, every $\eta\in(0,1]$, every $R>0$ such that $R/\eta$ is an integer $\ge2$, and every $c\in\Lambda_{\eta,R}$,
\[
a_{\eta}[\beta']_\nu(c)=\sum_{z\in\mathbb{Z}^3}\mathrm{gs}_\nu[\beta']\left(c+\frac{\eta}{2}\cdot e_\nu+R\cdot z\right),
\]
a sum with finitely many non-zero terms.
In the setting (ST*)$^{\mathrm{ch}}$, with the rule for the points of the torus,
\[
a_{\eta,\nu}(c)=\mathrm{gs}_\nu[\tau_v\beta]\left(c+\frac{\eta}{2}\cdot e_\nu\right).
\]
For every probe $\beta'$, every $\eta\in(0,1]$ and every integer $i\ge0$ one has $\left|D^i\mathrm{gs}_\nu[\beta']\right|_\infty\le2\,|\beta'|_{C^{i+1}}$;
and $\mathrm{gs}_\nu[\beta'](y)\ne0$ only if the sup distance of $y$ from $\operatorname{supp}\beta'$ is at most $\eta$.
\item[(b2)] Both end points of every bond in the support of $\hat{\mathfrak{a}}$ lie in $\mathrm{Qb}$, in physical coordinates.
\item[(c1)] One has
\[
C_{\beta}\le2\,h^{\frac{1}{2}}\cdot\left(1+h+3h^2\right)\cdot|\beta|_{C^3},\qquad
\mathbf{D}\le48\,|\beta|_{C^2}^2\cdot\left(D_\beta+3\right)^3 .
\]
The numbers $C_{\beta}$ and $\mathbf{D}$ depend on $(\beta,h,L,\alpha)$ only.
\item[(c2)] The number $|\hat{\mathfrak{a}}|_{\mathrm{ext}}$ is the sum of the three norms of $\mathrm{gsc}$ at the triple $(N,a_{*},1)$ of $\mathrm{Ind}_\beta$, and $\boldsymbol{\sigma}^2$ is the number $\left\|d\,\mathrm{w0}\right\|^2$ at this triple.
Hence $|\hat{\mathfrak{a}}|_{\mathrm{ext}}\le C_{\beta}$ and $\boldsymbol{\sigma}^2\le\mathbf{D}$.
\item[(d)] The three norms of $\hat{\mathfrak{a}}_{\mathfrak{c}}^{[k_{*}]}$ are those of $\hat{\mathfrak{a}}^{[k_{*}]}$;
hence $\epsilon_{\mathfrak{c}}$ does not depend on $\mathfrak{c}$.
The supports of the cochains $d\hat{\mathfrak{a}}_{\mathfrak{c}}$, $\mathfrak{c}\in\mathrm{Cel}$, are pairwise disjoint, and $\left\|d\hat{\mathfrak{a}}_{\mathfrak{c}}\right\|^2=\boldsymbol{\sigma}^2$.
\item[(e)] One has
\[
\#\left|\mathrm{Ypr}_{\mathfrak{c}}^{+}\right|\le\mathrm{pp},\qquad
d_{\mathrm{sep}}\ge\frac{\ell_*}{8L}-1=\frac{M\cdot r_{*}}{8L}-1\ge4;
\]
and, for $\mathfrak{c}\ne\mathfrak{c}'$, every end point of a bond of $\mathrm{Bd}_{\mathfrak{c}}$ has distance at least $\frac{\ell_*}{8L}$ in $T_*$ from every end point of a bond of $\mathrm{Bd}_{\mathfrak{c}'}$.
\item[(f)] For every $t\in\mathbb{C}^{\mathrm{Cel}}$, each of the three norms of $\mathscr{C}(t)$ is the maximum over $\mathfrak{c}$ of $|t_{\mathfrak{c}}|$ times the same norm of $\hat{c}_{\mathfrak{c}}$.
Hence
\[
|a(t)|_{\mathrm{ext}}\le\max_{\mathfrak{c}}|t_{\mathfrak{c}}|\cdot\frac{C_{\beta}}{e},
\]
and $|a(t)|_{\mathrm{ext}}<e_{*}^{-1+3\varepsilon}$ for $t\in\mathrm{DDo}\left(t_x^{+}\right)$.
Moreover $t_x\cdot\epsilon_{\mathfrak{c}}<e_{*}^{3\varepsilon}$, so that $\epsilon_{\mathfrak{c}}\le q_x$.
\item[(g)] If $\hat{\mathfrak{a}}\ne0$, then every set $\mathrm{Bd}_{\mathfrak{c}}$, $\mathrm{Ypr}_{\mathfrak{c}}$, $\mathrm{Ypr}_{\mathfrak{c}}^{+}$ is non-empty, and $d_{\mathrm{sep}}<\infty$.
If $\hat{\mathfrak{a}}=0$, all these sets are empty.
\end{itemize}
\end{named}

\begin{named}{Remark}
Proposition~35.3 corresponds to \cite[Section~20.1.4]{QEDI}, including its paragraph on the constants $C_{\beta}$ and $\mathbf{D}$, and to \cite[Lemma~20.3]{QEDI}.
No fermionic object occurs in Proposition~35.3, and its mathematical content is that of these two places of \cite{QEDI}, with the following differences.
The numbers $C_{\beta}$ and $\mathbf{D}$ are defined as suprema over the index set $\mathrm{Ind}_\beta$ of Section~G5.1.3, without reference to the charge, to a torus or to a placement, and (c1) is stated without the setting (ST*)$^{\mathrm{ch}}$ and without a placement.
Clause~(b1) is stated for the functions $\mathrm{gs}_\nu[\beta']$ before periodisation, and for all probes $\beta'$.
The evaluation of functions on $\mathbb{R}^3$ at points of the torus through periodisation is part of the statement.
In (f) the hypothesis $d_{\mathrm{sep}}\ge4$ of \cite[Lemma~20.3(f)]{QEDI} is omitted, since it always holds by (e), because $M\ge L^2$ and $r_{*}\ge L^2$.
The first sentence of~(g) corresponds to the assertion of \cite[Lemma~20.20]{QEDI} that, if $\hat{\mathfrak{a}}\ne0$, the sets $\mathrm{Bd}_c$, $P_c$ and $P_c^{+}$ are non-empty for every $c\in\mathrm{Cel}$;
this assertion is proved in the second paragraph of the proof of that lemma, where it is also noted that then $d_{\mathrm{sep}}=\min_{c\ne c'}d\left(P_c,P_{c'}\right)$, so that $d_{\mathrm{sep}}<\infty$.
The sets $P_c$ and $P_c^{+}$ of \cite{QEDI} are the sets $\mathrm{Ypr}_{\mathfrak{c}}$ and $\mathrm{Ypr}_{\mathfrak{c}}^{+}$ of Section~G5.1.4, and for $\hat{\mathfrak{a}}=0$ the definition of $d_{\mathrm{sep}}$ in \cite[Section~20.1.4]{QEDI} puts $d_{\mathrm{sep}}:=\infty$.
The second sentence of~(g), by which all these sets are empty if $\hat{\mathfrak{a}}=0$, is not stated in \cite{QEDI}.
The assertion of \cite[Lemma~20.20]{QEDI} that corresponds to the first sentence of~(g) is made under the hypotheses of that lemma, among them the smallness conditions (sm1), (sm2) and (sm3) of \cite[Section~20.1.8]{QEDI};
of these three conditions, the second paragraph of the proof of the lemma, which proves the assertion, uses only the clause of (sm2) by which the placement is admitted.
Here the first sentence of~(g) is a clause of Proposition~35.3, and it assumes nothing beyond the setting (ST*)$^{\mathrm{ch}}$, a probe $\beta$ and an admitted placement $v\in\eta\mathbb{Z}^3$;
in particular, none of the conditions (g1)--(g9) of Section~G5.1.9, the first eight of which correspond to the smallness conditions (sm1)--(sm8) of \cite[Section~20.1.8]{QEDI}, is among its hypotheses.
The numbers $D_\beta/2$ and $h$ stand in the places of the numbers $\rho_\beta$ and $h'$ of \cite{QEDI}.
\end{named}

\subsubsection*{Proposition~35.4 (what holds at the stop level and below it)}

\begin{named}{Proposition 35.4}
Assume the setting (ST*)$^{\mathrm{ch}}$.
\begin{itemize}
\item[(SF1)] \emph{Numbers.}
The pair $P_N(e)$ is real and controlled on $R^{*}(e)$:
its trajectories (Section~G5.1.1) exist up to the level $k_{*}$, the numbers $\varepsilon^{\mathrm{v},t}_k$ and $m_k$ are real, and $|m_k|\le w^{\mathrm{m}}_k/8$ for $k\le k_{*}$.
In particular $|m_*|\le e_{*}^{\frac{3}{4}-8\varepsilon}/8$.
\item[(SF2)] \emph{Domains, operators.}
At every level $j\le k_{*}$ (at $j=0$ the domains being formed with $q$), the sets $G_j$ and $\tilde{R}_j+\mathfrak{X}_j$, in particular the set $\mathfrak{X}_j$, are contained in $\tilde{G}_j$;
the sets $\mathrm{cf}\cdot\tilde{R}_j$ ($0<\mathrm{cf}\le1$) and $\tilde{G}_j$ are invariant under the addition of real constants, under the map $A\mapsto-A$, and under real gauge transformations.
For every $A\in\tilde{G}_*$ the operator $\mathrm{Fst}(A)$ is a bijection of $\mathfrak{S}^{\mathrm{ch}}(\eta_*,\ell_*)$;
the map $A\mapsto S_*(A)$ is analytic on the open set $\tilde{G}_*$;
and for every $A\in\tilde{G}_*$ and all unit sites $y$, $y'$
\[
|S_*(A)|(y,y'),\quad
|\mathrm{Sa}_*(A)|(y,y'),\quad
|H_*(A)|(y,y'),\quad
\left|\delta_{\alpha,A}H_*(A)\right|(y,y')\quad\le\quad C^{F}\cdot\mathrm{Edy}(y,y'),
\]
the second and the fourth majorant being taken for the short pairs ($d_X<1$), \emph{the pairs inside a cluster included}.
For $A\in\tilde{G}_*$, for $(x,x')\in\mathrm{Pz}_*$ with $d_X(x,x')\ge1$ and for every $f\in\mathfrak{S}^{\mathrm{ch}}(\eta_*,\ell_*)$,
\[
\left|\left(\delta_{\alpha,A}f\right)(x,x')\right|_c\le\mathrm{e}\cdot|f(x')|_c+|f(x)|_c .
\]
The set $\mathrm{Gwk}$ is contained in $\tilde{R}_*+\mathfrak{X}_*$, hence in $\tilde{G}_*$, and it is invariant under $A\mapsto-A$.
For $A\in\mathrm{Gwk}$ one has
\[
\mathfrak{i}(A)<e_{*}^{2\varepsilon}\le\frac{1}{16},\qquad
3M\cdot\mathfrak{i}(A)<\frac{1}{4},\qquad
\left\|S_*(A)\right\|_{\gamma^{F}/2}\le\mathbf{S},\qquad
\left\|\mathrm{Sa}_*(A)\right\|_{\gamma^{F}/2}\le\mathbf{S};
\]
and $\mathfrak{i}<\frac{1}{4}$ on $\mathrm{Gwk}^{+}$.
\item[(SF3)] \emph{Activities with slot.}
One has $\mathscr{E}_*\in K^{x}_*$ (Section~I.5.4, read on $T^{\mathrm{f}}_*$), and
\[
\left\|\mathscr{E}_*\right\|^{x}_*\le4\,e_{*}^{\varepsilon}\cdot w^{\mathrm{E}}_* .
\]
For $X\in D_*$ and $(\mathscr{A},\mathscr{C})\in\tilde{R}_*\times\mathfrak{X}_*$, let $\mathscr{E}_{nm}\left(X;\mathscr{A},\mathscr{C};\xi_1,\dots,\xi_n;\boldsymbol{\zeta}_1,\dots,\boldsymbol{\zeta}_m\right)$ be the kernels of $\mathscr{E}_*(X;\mathscr{A},\mathscr{C})$ that are antisymmetric in $\xi_1,\dots,\xi_n$ and, separately, in $\boldsymbol{\zeta}_1,\dots,\boldsymbol{\zeta}_m$;
by Proposition~15\,(a), such kernels exist, and they are determined by $\mathscr{E}_*(X;\mathscr{A},\mathscr{C})$.
Membership in $K^{x}_*$ means in particular that these kernels are analytic in $(\mathscr{A},\mathscr{C})\in\tilde{R}_*\times\mathfrak{X}_*$;
that they vanish unless all $\xi_i$ lie over $X$ and all pairs $\boldsymbol{\zeta}_j$ are admissible for $X$ (the kinds (I), (II) of Section~I.1.7);
that they depend on $\mathscr{A}(b)$, $\mathscr{C}(b)$ only for the bonds $b$ that meet $X$ (a bond meets $X$ if it has an end point in $X$; Section~I.1.1);
and that they have equally many barred and unbarred arguments.
For these antisymmetric kernels, the left side of the display (2.1) below equals the norm $\left\|\mathscr{E}_*(X;\mathscr{A},\mathscr{C})\right\|_{h_*}$ of Section~I.5.1, with $h_*=\left(h_{*,1},h_{*,2}\right)$;
this is the case of equality in Section~I.5.1.
Hence the clause (k6) of Section~I.5.4 and the bound on $\left\|\mathscr{E}_*\right\|^{x}_*$ give, for every $X\in D_*$ and every $(\mathscr{A},\mathscr{C})\in\tilde{R}_*\times\mathfrak{X}_*$,
\begin{equation}
\sum_{n,m}\frac{h_{*,1}^n\cdot h_{*,2}^m}{n!\cdot m!}\cdot\left\|\mathscr{E}_{nm}(X;\mathscr{A},\mathscr{C})\right\|\le4\,e_{*}^{\varepsilon}\cdot w^{\mathrm{E}}_*\cdot\exp\left(-\kappa\cdot d_M(X)\right),\tag{2.1}
\end{equation}
where $\left\|\mathscr{E}_{nm}\right\|=\int\mathrm{d}\xi\,\mathrm{d}\boldsymbol{\zeta}\,\left|\mathscr{E}_{nm}\right|$ is formed with the weights $\eta_*^3/48$ and $\left(\eta_*^3/48\right)^2$ of Section~I.3.6.
\item[(SF4)] \emph{Determinants.}
For $0\le j<k_{*}$ the system of primitives on $T_j\left(R^{*}(e)\right)$ has the properties W1--W17 for $\mathrm{par}$;
these are the properties numbered 1--17 in \cite[Section~5.4.2]{QEDI}, with W1, W2, W5--W17 in the displayed form of Section~G1.11, and they are read as in Section~I.5.8;
at $j=0$, W7 is the bound with the factor $(1+y)^{-1}$, together with Borel measurability in $y$.
For $X\in D_j$ the functions $E^{d}_j(X,\cdot)$ of Section~I.5.7, formed with this system (for $j\ge1$ with $\acute{E}_j$, which is a trace on $\mathfrak{S}(1,\ell_j)$; \emph{at $j=0$ with $\acute{E}_0$}, which is formed with traces on $\mathfrak{S}^{\mathrm{ch}}(1,\ell_0)$), are analytic on $\tilde{G}_j$, and for $A,\tilde{A}\in\tilde{G}_j$:
\begin{itemize}
\item[(i)] $\left|E^{d}_j(X,A)\right|\le C_{d}\cdot M^3\cdot\exp\left(-(\kappa+1)\cdot d_M(X)\right)$;
\item[(ii)] $E^{d}_j(X,A)=E^{d}_j(X,\tilde{A})$ if $A$ and $\tilde{A}$ coincide on the bonds that meet $X$;
\item[(iii)] $E^{d}_j(X,A-d\boldsymbol{\lambda})=E^{d}_j(X,A)$ for every real function $\boldsymbol{\lambda}$ of the photon sites of $T_j$;
\item[(iv)] $E^{d}_j(X,-A)=E^{d}_j(X,A)$.
\end{itemize}
Moreover $E_j^{\mathrm{det}}(X,A)=E^{d}_j(X,A)-E^{d}_j(X,0)$, one has $\mathfrak{P}_0=0$, and
\[
\mathfrak{P}_{j+1}(Y,\mathscr{C})=\sum_{X:\ \bar{X}^{L}=LY}\left(\mathfrak{P}_j\left(X,\mathscr{C}_L\right)+E_j^{\mathrm{det}}\left(X,\mathscr{C}_L\right)\right)
\qquad\left(Y\in D_{j+1},\ \mathscr{C}\in\mathfrak{X}_{j+1}\right).
\]
\begin{itemize}
\item[(v)] (Polymers, at every level $j\le k_{*}$.)
For every $M$-cube $\square$ of $T_j$,
\[
\sum_{X\in D_j:\ X\supset\square}\exp\left(-\kappa_{0}\cdot d_M(X)\right)\le K_{0};
\]
any two points of a polymer $X\in D_j$ have distance at most $M\cdot\left(d_M(X)+2\right)$;
every block $\Delta_y$ lies in one $M$-cube.
For $j<k_{*}$, every $X\in D_j$ has exactly one $Y\in D_{j+1}$ with $\bar{X}^{L}=LY$, and then $d_M(X)\ge L\cdot d_M(Y)$;
and for every $Y\in D_{j+1}$
\[
\sum_{X\in D_j:\ \bar{X}^{L}=LY}\exp\left(-\kappa\cdot d_M(X)\right)\le c_{\mathrm{rb}}\cdot L^3\cdot\exp\left(-L\cdot\left(\kappa-\kappa_{0}-1\right)\cdot d_M(Y)\right).
\]
\item[(vi)] Let $a$ be a complex one-cochain on $\mathbb{Z}^3/n_*\mathbb{Z}^3$ with $a^{[j]}\in\mathfrak{X}_j$ for all $j\le k_{*}$.
Then $D_*\left(a^{[k_{*}]}\right)$ is defined, $\det D_*(0)\ne0$, and
\begin{equation}
\begin{split}
\frac{\det\left(D_a+\bar{m}\cdot\eta\right)}{\det\left(D_0+\bar{m}\cdot\eta\right)}&=\exp\left\{\sum_{X\in D_*}\mathfrak{P}_*\left(X,a^{[k_{*}]}\right)\right\}\cdot\frac{\det D_*\left(a^{[k_{*}]}\right)}{\det D_*(0)},\\
\sum_{X\in D_*}\mathfrak{P}_*\left(X,a^{[k_{*}]}\right)&=\sum_{j<k_{*}}\ \sum_{X\in D_j}E_j^{\mathrm{det}}\left(X,a^{[j]}\right).
\end{split}\tag{2.2}
\end{equation}
On the left the determinants have the size $192\,n_*^3$ and are taken on $\mathfrak{S}^{\mathrm{ch}}(1;n_*)$;
on the right they have the size $4\,\ell_*^3$ and are taken on $\mathfrak{S}(1,\ell_*)$;
the sum over $j$ contains the term $j=0$.
\end{itemize}
\item[(SF5)] \emph{Small-field set, block field; photon objects.}
The map $F^{\mathrm{u}}\mapsto\left(A_{*};\zeta_0,\dots,\zeta_{k_{*}-1}\right)$ is a linear bijection of $F_N$ onto ${V_{\mathrm{t}}}(T^0_*)\times\prod_{j<k_{*}}V_{j}$.
One has $\zeta_j(\Upsilon)=\left\langle\phi_{j,\Upsilon},F^{\mathrm{u}}\right\rangle$ with an orthonormal family $\left\{\phi_{j,\Upsilon}\right\}$ in $F_N$;
the space $V_{j}$ has at most $2n_j$ basis elements $\Upsilon$, where $n_j=\ell_*^3\cdot L^{3(k_{*}-j)}$;
and $(dA_{*})(p)=\left\langle\phi_p,F^{\mathrm{u}}\right\rangle$ with $\phi_p\in F_N$ and $\|\phi_p\|\le1$, for each of the $3\ell_*^3$ plaquettes $p$ of $T^0_*$.
The measure $\mathrm{d}\mathsf{m}_{*}$ is a non-degenerate centred Gaussian measure on ${V_{\mathrm{t}}}(T^0_*)$.
The set $\left\{A\in{V_{\mathrm{t}}}(T^0_*):\ \chi_{*}(A)=1\right\}$ is compact and convex, and $0$ is an interior point of it.
For every $A$ in this set there is a real function $\boldsymbol{\lambda}$ on $T^0_*$ such that, with $A^{g}:=A-d\boldsymbol{\lambda}$ and with the photon minimiser $\mathscr{A}^0:=H_*A^{g}$,
\begin{equation}
\begin{aligned}
&\left|A^{g}\right|_\infty\le2\,\ell_*^2\cdot p_*;\qquad
\mathscr{A}^0\ \text{is real and}\ \left|\mathscr{A}^0\right|_\infty<\frac{1}{2}\,e_{*}^{-\varepsilon};\\
&\text{for all }\theta\in\mathbb{R}^3\text{ and }\mathrm{sv}\in[0,1]:\quad
A+c_{\theta}\in\mathrm{Dom}_*,\quad
\mathrm{sv}\cdot A^{g}+c_{\theta}\in\mathrm{Dom}_*,\\
&\qquad H_*\left[\mathrm{sv}\cdot A^{g}+c_{\theta}\right]=\mathrm{sv}\cdot\mathscr{A}^0+c_{\theta}\in\frac{1}{4}\,\tilde{R}_* .
\end{aligned}\tag{2.3}
\end{equation}
Here $H_*[\,\cdot\,]$ is the photon minimiser $H_{k_{*}}$ applied to a bond function on $T^0_*$, and not the fermion operator $H_*(A)$ of Section~G5.1.7.
\item[(SF6)] \emph{Representation.}
Let $A\in\mathrm{Dom}_*$ and $\mathscr{C}\in\mathfrak{X}_*$, put $\mathscr{A}:=H_*A$ and $\mathscr{T}:=\mathscr{A}+\mathscr{C}$, and insert the fields
\[
\psi=H_*(\mathscr{T})\Psi,\qquad
\bar{\psi}=\bar{H}_*(\mathscr{T})\bar{\Psi},\qquad
\chi=\delta_{\alpha,\mathscr{T}}\psi,\qquad
\bar{\chi}=\delta_{\alpha,-\mathscr{T}}\bar{\psi}.
\]
Then
\begin{equation}
\begin{split}
Y^{x}_{*}(A;\mathscr{C})={}&\det D_*(0)^{-1}\cdot\exp\left\{-\varepsilon^{\mathrm{v},t}_*\cdot\ell_*^3+\sum_{X\in D_*}\mathfrak{P}_*(X,\mathscr{C})\right\}\\
&\cdot\int\mathrm{D}\Psi\,\exp\left\{-\left\langle\bar{\Psi},D_*(\mathscr{T})\Psi\right\rangle-m_*\cdot\left\langle\bar{\psi},\psi\right\rangle+\sum_{X\in D_*}\mathscr{E}_*\left(X;\mathscr{A},\mathscr{C},\psi,\chi\right)\right\},
\end{split}\tag{2.4}
\end{equation}
where $\left\langle\bar{\psi},\psi\right\rangle=\int_{T^{\mathrm{f}}_*}\mathrm{d}x\,\bar{\psi}(x)\cdot\psi(x)$, and where the Berezin integral $\int\mathrm{D}\Psi$ extends over the $4\ell_*^3$ pairs of generators, with the weight $\nu=1$.
Let $a$ be a complex one-cochain on $\mathbb{Z}^3/n_*\mathbb{Z}^3$ with $a^{[j]}\in\mathfrak{X}_j$ for $j\le k_{*}$.
Then $Y^{x}_{*}\left(A+d\boldsymbol{\lambda};a^{[k_{*}]}\right)=Y^{x}_{*}\left(A;a^{[k_{*}]}\right)$ whenever $A$ and $A+d\boldsymbol{\lambda}$ lie in $\mathrm{Dom}_*$, $\boldsymbol{\lambda}$ being a real function on $T^0_*$;
and for every $\theta\in\mathbb{R}^3$ and every bounded Borel function $\Phi$ of $A_{*}$,
\begin{equation}
\nu^{\theta,a}\left(\mathbf{1}_{\mathfrak{G}}\,\Phi\right)=\int\mathrm{d}\mathsf{m}_{*}(A)\,\Phi(A)\cdot\chi_{*}(A)\cdot Y^{x}_{*}\left(A+c_{\theta};a^{[k_{*}]}\right).\tag{2.5}
\end{equation}
\item[(DF5)] \emph{External connections.}
For \emph{every} complex one-cochain $a$ on $\mathbb{Z}^3/n_*\mathbb{Z}^3$ and every $0\le k\le k_{*}$, with $i:=k_{*}-k$,
\[
a^{[k]}(b)=L^{-\frac{i}{2}}\cdot a^{[k_{*}]}\left(b/L^i\right),\qquad
\left|a^{[k]}\right|_\infty=L^{-\frac{i}{2}}\cdot\left|a^{[k_{*}]}\right|_\infty,\qquad
\left|\partial a^{[k]}\right|_\infty=L^{-\frac{3i}{2}}\cdot\left|\partial a^{[k_{*}]}\right|_\infty,
\]
and, for $k<k_{*}$, $\left(a^{[k+1]}\right)_L=a^{[k]}$.
If $|a|_{\mathrm{ext}}<e_{*}^{-1+3\varepsilon}$, then $a^{[k]}\in\mathfrak{X}_k$ for all $k\le k_{*}$.
Consequently, for a probe $\beta$, an admitted placement and $t\in\mathrm{DDo}\left(t_x^{+}\right)$, one has $a(t)^{[j]}\in\mathfrak{X}_j$ for all $j\le k_{*}$, and $e_{*}\cdot\left|\mathscr{C}(t)\right|_\infty<e_{*}^{3\varepsilon}$.
\item[(A6)] \emph{Terminal bounds.}
Let $A\in\tilde{G}_*$, and let $\psi=H_*(A)\Psi$ and $\bar{\psi}=\bar{H}_*(A)\bar{\Psi}$.
Then
\[
\left|\left\langle\bar{\psi},\psi\right\rangle\right|_{\mathrm{aw}}\le4\,e_{*}^{-\frac{1}{2}+2\varepsilon}\cdot\ell_*^3\qquad\text{for }\mathrm{aw}=C_{0}^{-1}\cdot h_{*,2}.
\]
For $(\mathscr{A},\mathscr{C})\in\tilde{R}_*\times\mathfrak{X}_*$, $\mathscr{T}=\mathscr{A}+\mathscr{C}$ and $X\in D_*$, the norm $\left|\mathscr{E}_*\left(X;\mathscr{A},\mathscr{C},H_*^{\#}(\mathscr{T})\Psi\right)\right|_{\mathrm{aw}}$, at the same $\mathrm{aw}$, is at most the left side of (2.1).
\item[(A7)] \emph{Block covariance at real fields.}
Let $A$ be a real bond function on $T_*$.
Then
\[
\operatorname{Re}\left(f,\left(D_A+\mu_*\right)f\right)^{\mathrm{ch}}\ge\mu_*\cdot(f,f)^{\mathrm{ch}}\qquad\left(f\in\mathfrak{S}^{\mathrm{ch}}(\eta_*,\ell_*)\right);
\]
one has $\left(Q_*(A)J^{*}\right)^{*}=J\,Q_*^{T}(-A)$;
the operator $\mathrm{Fst}(A)$ is invertible;
the operator $D_*(A)$ is invertible, $\left|D_*(A)^{-1}\right|_{\ell^2}\le b_*^{-1}+\mu_*^{-1}$, and its $4\times4$ blocks satisfy, for all unit sites $y$, $y'$,
\[
\left|D_*(A)^{-1}(y,y')\right|_2\le b_*^{-1}\cdot\mathbf{1}[y=y']+\frac{2\mathrm{e}}{\mu_*}\cdot\exp\left(-\frac{\mu_*}{12\mathrm{e}}\cdot d(y,y')\right).
\]
\end{itemize}
\end{named}

\begin{named}{Remark}
Proposition~35.4 corresponds to the following statements of \cite{QEDI}:
the clauses (SF1)--(SF6) to the clauses (SL1)--(SL6) of \cite[Proposition~20.6]{QEDI};
(DF5) to \cite[Lemma~20.7]{QEDI} together with the clause (SL7) of \cite[Proposition~20.6]{QEDI};
(A6) and (A7) to the displays (Gr6) and (Gr7) of \cite[Section~20.2.4]{QEDI};
the last two sentences of (SF2) to the facts on the sets $\mathrm{Gwk}$ and $\mathrm{Gwk}^{+}$, which are denoted $G_{\mathrm{w}}$ and $G_{\mathrm{w}}^{+}$ in \cite{QEDI}, at the beginning of \cite[Section~20.4]{QEDI};
and the first sentence of (SF4) to the assertion of \cite[Lemma~20.1(c)]{QEDI} on the systems of primitives.
The differences are the following.
\begin{itemize}
\item[(1)] All fine objects are read on $T^{\mathrm{f}}_*$, with the integral $\int\mathrm{d}x$;
the short pairs are those with $d_X<1$, the pairs inside a cluster included;
the pairs of kind (II) are those with $x'=x^{(3)}$.
\item[(2)] The constants $(C_F,\gamma_F)$ of \cite{QEDI} are here the constants $(C^{F},\gamma^{F})$ of Proposition~14.1\,(i), and the norm is $|\cdot|_c$.
\item[(3)] Clause~(SF4) specifies that at $j=0$ the property W7 is taken with the factor $(1+y)^{-1}$, and that $\acute{E}_0$ is formed with traces on spaces of dimension $192\,n_0$;
(2.2) specifies the sizes of the determinants.
\item[(4)] Clause~(SF4)\,(v) contains the fact that every block $\Delta_y$ lies in one $M$-cube;
this is the fact (g0) of \cite[Section~17.4.1]{QEDI}.
\item[(5)] Clause~(SF5) contains the compactness and the convexity of the set $\{\chi_{*}=1\}$.
\item[(6)] Clause~(DF5) states the scaling relations for every complex one-cochain $a$.
\item[(7)] Clause~(A7) contains the invertibility of $\mathrm{Fst}(A)$, and it specifies the norm $|\cdot|_2$ of the blocks.
\item[(8)] The displays (Gr1)--(Gr5) of \cite[Section~20.2.4]{QEDI} have no counterpart here:
they are the rules (I.N1), (I.S6c)\,(i)--(iii), (I.N2), (I.N0) of Section~I.5.2 for the generators of Section~I.3.6 with the weight $\nu=1$, which hold for every finite index set.
\item[(9)] No hypothesis that the lattice mass be larger than $-1$ is made (I.K10).
\end{itemize}
\end{named}

\subsection*{33. The fermionic operators at the stop level as functions of the connection}

In Propositions 33.1--33.4 the setting is (ST*)$^{\mathrm{ch}}$, the notation is that of Section~G5.1.7, and the constants are those of Section~G5.1.8.
No probe enters.
Every bound below holds uniformly in $N$, in $b_{*}\in[b_{-},1]$, in $\mu_{*}\in\left(\mu_{-},\frac{1}{L}\right]$ and in the bond functions that occur in it.

\subsubsection*{Proposition~33.1 (elementary operators)}

\begin{named}{Proposition 33.1}
The operators
\[
D_A,\qquad Q_{*}(A)J^{*},\qquad JQ_{*}^{T}(-A),\qquad JP_{*}(A)J^{*},\qquad \mathrm{Fst}(A),\qquad \delta_{\alpha,A}
\]
are entire functions of the complex bond function $A$ on $T_{*}$.
Let $A$, $B$, $B'$ be complex bond functions on $T_{*}$ such that $\mathfrak{i}(A)\le\frac{1}{4}$.
\begin{itemize}
\item[(a)] For all unit sites $y$, $y'$,
\begin{gather*}
\left|Q_{*}(A)J^{*}\right|(y,y'),\ \left|JQ_{*}^{T}(-A)\right|(y,y')\le\frac{3}{2}\,\mathbf{1}[y=y'],\\
\left|\partial_B\left(Q_{*}(A)J^{*}\right)\right|(y,y'),\ \left|\partial_B\left(JQ_{*}^{T}(-A)\right)\right|(y,y')\le\frac{9}{4}\,\epsilon_B\cdot\mathbf{1}[y=y'\in\mathrm{Ypr}_B].
\end{gather*}
\item[(b)] Let $f\in\mathfrak{S}^{\mathrm{ch}}(\eta_{*},\ell_{*})$ and $x=(c,g)$, and put $c_{3}:=c+\eta_{*}\cdot g\cdot e_1$ and $x^{(3)}:=(c_{3},g\cdot s_3)$.
Then $\partial_BD_A$ consists of one term only:
\[
\left(\partial_BD_A\,f\right)(x)=-i\cdot e_{*}\cdot B(c,c_{3})\cdot M_{3}(x)\cdot u_A(c,c_{3})\cdot f\left(x^{(3)}\right).
\]
Its majorant satisfies $\left|\partial_BD_A\right|(z,z')\le4\,(1+\rho)\cdot\epsilon_B\le8\,\epsilon_B$ if $z=z'$ or if $z$, $z'$ are nearest neighbours in $T^0_{*}$, and $\left|\partial_BD_A\right|(z,z')=0$ for all other pairs.
Moreover $\left|\partial_B\left(JP_{*}(A)J^{*}\right)\right|(y,y')\le7\,\epsilon_B\cdot\mathbf{1}[y=y']$.
The first majorant vanishes unless $z$ and $z'$ lie in $\mathrm{Ypr}_B$, the second unless $y$ lies in $\mathrm{Ypr}_B$.
Let $\mathfrak{d}(z,z):=15$, let $\mathfrak{d}(z,z'):=8$ if $z$, $z'$ are nearest neighbours, and let $\mathfrak{d}(z,z'):=0$ for all other pairs.
It follows that
\[
\left|\partial_B\mathrm{Fst}(A)\right|(z,z')\le\epsilon_B\cdot\mathfrak{d}(z,z')\cdot\mathbf{1}[z,z'\in\mathrm{Ypr}_B],\qquad
\left\|\partial_B\mathrm{Fst}(A)\right\|_{\gamma}\le146\,\epsilon_B\quad\text{for }0\le\gamma\le1,
\]
and that for all unit sites $y$, $y'$
\begin{equation}\tag{4.1}
\begin{split}
&\sum_{z,z'}\mathrm{Edy}(y,z)\cdot\left|\partial_B\mathrm{Fst}(A)\right|(z,z')\cdot\mathrm{Edy}(z',y')\\
&\quad\le146\,\epsilon_B\sum_{z\in\mathrm{Ypr}_B}\mathrm{Edy}(y,z)\cdot\mathrm{Edy}(z,y').
\end{split}
\end{equation}
\item[(c)] Assume in addition $3M\cdot\mathfrak{i}(A)\le\frac{1}{4}$.
For $(x,x')\in\mathrm{Pz}_{*}$ and every $f$ one has $\left|U_A(x,x')^{\pm1}\right|\le\frac{4}{3}$ and
\[
\left|\left(\partial_B\delta_{\alpha,A}f\right)(x,x')\right|_c\le4\,d_X(x,x')^{1-\alpha}\cdot\epsilon_B\cdot|f(x')|_c\le4M\cdot\epsilon_B\cdot|f(x')|_c,
\]
and the left side is at most $4\,\epsilon_B\cdot|f(x')|_c$ if $d_X(x,x')<1$.
The left side vanishes for the pairs inside a cluster ($c=c'$).
It vanishes unless $d([x],\mathrm{Ypr}_B)\le M$ and $d([x'],\mathrm{Ypr}_B)\le M$; if $d_X(x,x')<1$, it vanishes unless $d([x],\mathrm{Ypr}_B)\le1$.
\item[(d)] (Additivity.)
If $d(\mathrm{Ypr}_B,\mathrm{Ypr}_{B'})\ge M+2$, then $\partial_B\partial_{B'}O(A)=0$ for each of the functions
\[
O=D_{(\cdot)},\qquad Q_{*}(\cdot)J^{*},\qquad JQ_{*}^{T}(-\,\cdot\,),\qquad JP_{*}(\cdot)J^{*},\qquad \mathrm{Fst}(\cdot),\qquad \delta_{\alpha,\cdot}.
\]
\end{itemize}
\end{named}

\begin{named}{Remark}
Proposition 33.1 corresponds to \cite[Lemma~20.14]{QEDI}, together with the formulas for $D_A$, $Q_{*}(A)$, $Q_{*}^{T}(-A)$, $\delta_{\alpha,A}$ and the statement on the rows of the projections $P^{\boldsymbol{\sigma}}_{\nu}$ in \cite[Section~20.2.2]{QEDI}.
Proposition~33.1 and these statements of \cite{QEDI} differ in the points (1)--(3) below; point (4) records that the condition on the pairs $(z,z')$ in (b) is the one of \cite[Lemma~20.14(b)]{QEDI}.
(1) In \cite[Section~20.2.2]{QEDI} the operator $D_A$ is built from six hops with the projections $P^{\boldsymbol{\sigma}}_{\nu}$, whose rows have norm one, and at most one of these hops leads into a given neighbouring cube; here there is a single hop, with the matrix $M_{3}(x)$, whose rows have norm $3\,(1+\rho)\le6$.
Hence the majorant of $\partial_BD_A$ is $8\,\epsilon_B$ on the diagonal, as in the proof of \cite[Lemma~20.14]{QEDI}, and it is $8\,\epsilon_B$ also between neighbouring cubes, where that proof has $\frac{4}{3}\,\epsilon_B$.
Accordingly the value $\mathfrak{d}(z,z')=8$ for neighbours takes the place of the number $\frac{3}{2}$ of \cite[Lemma~20.14(b)]{QEDI}, and $15+6\cdot8\cdot\mathrm{e}\le146$ takes the place of the constant $40$ there, also in (4.1).
(2) The averaging operators carry the factors $J^{*}$ and $J$, which have norm one for the sup norms by (I.J2); the identity $\int_{\Delta_y}\mathrm{d}x=1$ takes the place of $\eta_{*}^3\cdot\#|\Delta_y|=1$; the numbers $\frac{3}{2}$, $\frac{9}{4}$ and $7$ are those of \cite[Lemma~20.14]{QEDI}.
(3) In (c) the distance $d_X(x,x')$ takes the place of $|x-x'|$, and the assertion on the pairs inside a cluster has no counterpart in \cite[Lemma~20.14]{QEDI}.
(4) In \cite[Lemma~20.14(b)]{QEDI} the condition $|z-z'|_1=1$ concerns the $\ell^1$ distance of two unit sites of the torus; it says that $z$ and $z'$ are nearest neighbours in $T^0_{*}$, which is the condition used in (b).
The distances of pairs of fine sites in \cite[Section~20.2.2]{QEDI} are likewise distances on the torus.
On the present lattice the operator $D_A$ is given by (I.2.3), the operators $Q_{*}(A)J^{*}$ and $JQ_{*}^{T}(-A)$ by (I.3.4) and $\delta_{\alpha,A}$ by (I.1.7); the rows of the hopping matrices $M_i(x)$ are described by (I.m7).
These formulas take the place of the four displays of \cite[Section~20.2.2]{QEDI} for $D_A$, $Q_{*}(A)$, $Q_{*}^{T}(-A)$ and $\delta_{\alpha,A}$, and of the statement there on the rows of $P^{\boldsymbol{\sigma}}_{\nu}$.
\end{named}

The conditions on the constants that enter Proposition~33.1 are the following; they stand in the lines of Section~I.8.2 and in the rows of Proposition~14.2.
\begin{itemize}
\item $0<\rho<1$ (line (C0)), in the form $3\,(1+\rho)\le6$.
It gives the bound $4\,(1+\rho)\le8$ in (b), and with it the numbers $8$, $15$ and $146$ in (b) and in (4.1).
\item $0<\alpha<1$ (line (C1)).
It is used in (c): $d_X(x,x')^{1-\alpha}\le1$ if $d_X(x,x')<1$, and $d_X(x,x')^{1-\alpha}\le M$ if $1\le d_X(x,x')\le M$.
\item $L$ odd and $L>1$ (row (L5) of Proposition~14.2\,(b); line (C3)).
That $L$ is odd is used for the blocks $\Delta_y$, for the neighbour structure of $T^0_{*}$ and for the supports in $\mathrm{Ypr}_B$, in (a), (b) and (d); $L>1$ gives $\eta_{*}=L^{-k_{*}}<1$, which enters the bound $\left|u_A(c,c_{3})\right|\le\frac{4}{3}$ in (b).
\item $M$ an integer with $M\ge1$; this holds because $M$ is a power of $L$ (row (M1) of Proposition~14.2\,(b); line (C5)).
It is used in (c) and (d).
\item $0<\gamma^{F}\le1$ (line (C4r)).
It gives the factor $\mathrm{e}$ in (4.1).
\item $b=1$ (line (C0)), through $b_{*}\le1$ (Proposition~35.1\,(a)).
It gives the number $15$ in (b).
\end{itemize}
No condition on $e_0$, $e_1$ or $e_{\sharp}(\beta)$ enters, that is, none of (A8) and (g1)--(g9).
The smallness that the proof needs is carried by the inequalities $\mathfrak{i}(A)\le\frac{1}{4}$ and $3M\cdot\mathfrak{i}(A)\le\frac{1}{4}$, which are hypotheses on $A$, not conditions on the constants.

\subsubsection*{Proposition~33.2 ($S_{*}$, the legs and $D_{*}(\cdot)$ on $\mathrm{Gwk}$)}

\begin{named}{Proposition 33.2}
Let $A\in\mathrm{Gwk}$, and let $B$, $B'$ be complex bond functions on $T_{*}$; in (H2), (L3) and (D2) assume moreover $d(\mathrm{Ypr}_B,\mathrm{Ypr}_{B'})\ge M+2$.
Then the following bounds hold for every leg index $\mathrm{up}$ and all unit sites $y$, $y'$.
\begin{itemize}
\item[(H1)] The following holds, the second majorant being the majorant for the short pairs:
\begin{equation}\tag{4.2}
\begin{split}
&\left|\partial_BH_{*}(A)\right|(y,y'),\ \left|\delta_{\alpha,A}\cdot\partial_BH_{*}(A)\right|(y,y')\\
&\quad\le149\,(C^{F})^{2}\cdot\epsilon_B\sum_{z\in\mathrm{Ypr}_B}\mathrm{Edy}(y,z)\cdot\mathrm{Edy}(z,y').
\end{split}
\end{equation}
\item[(H2)] If $d(\mathrm{Ypr}_B,\mathrm{Ypr}_{B'})\ge M+2$, then, the second majorant again being the majorant for the short pairs,
\begin{equation}\tag{4.3}
\begin{split}
&\left|\partial_B\partial_{B'}H_{*}(A)\right|(y,y'),\ \left|\delta_{\alpha,A}\cdot\partial_B\partial_{B'}H_{*}(A)\right|(y,y')\\
&\quad\le21645\,(C^{F})^{3}\cdot\epsilon_B\cdot\epsilon_{B'}\sum_{z\in\mathrm{Ypr}_B}\ \sum_{z'\in\mathrm{Ypr}_{B'}}[\mathrm{Edy}(y,z)\cdot\mathrm{Edy}(z,z')\cdot\mathrm{Edy}(z',y')\\
&\qquad+\mathrm{Edy}(y,z')\cdot\mathrm{Edy}(z',z)\cdot\mathrm{Edy}(z,y')].
\end{split}
\end{equation}
\item[(L1)] One has
\[
\left|\mathrm{Leg}_A\right|(\mathrm{up};y')\le3\,C^{F}\sum_{w\in\mathrm{pos}(\mathrm{up})}\mathrm{Edy}(w,y');
\]
consequently $\left|\mathrm{Leg}_A\right|(\mathrm{up};y')\le\mathrm{hL}\cdot\exp\left(-\gamma^{F}\cdot d\left(\mathrm{pos}(\mathrm{up}),y'\right)\right)$, and
\[
\sum_{u}\left|\mathrm{Leg}_A(\mathrm{up};u)\right|\cdot\exp\left(\gamma\cdot d\left(\mathrm{pos}(\mathrm{up}),u\right)\right)\le\mathrm{hL}\qquad\text{for }0\le\gamma\le\frac{\gamma^{F}}{2}.
\]
\item[(L2)] One has
\[
\left|\partial_B\mathrm{Leg}_A\right|(\mathrm{up};y')\le C_H\cdot\epsilon_B\sum_{w\in\mathrm{pos}(\mathrm{up})}\ \sum_{z\in\mathrm{Ypr}_B^{+}}\mathrm{Edy}(w,z)\cdot\mathrm{Edy}(z,y'),
\]
and the bound holds with $149\,(C^{F})^{2}$ in place of $C_H$ if $\mathrm{up}=\xi$, with $160\,(C^{F})^{2}$ if $\mathrm{up}=\boldsymbol{\zeta}$ belongs to a short pair, and with $203M\cdot(C^{F})^{2}$ if $\mathrm{up}=\boldsymbol{\zeta}$ belongs to a long pair.
\item[(L3)] One has
\begin{multline*}
\left|\partial_B\partial_{B'}\mathrm{Leg}_A\right|(\mathrm{up};y')\le C_H^{2}\cdot\epsilon_B\cdot\epsilon_{B'}\sum_{w\in\mathrm{pos}(\mathrm{up})}\ \sum_{z\in\mathrm{Ypr}_B^{+}}\ \sum_{z'\in\mathrm{Ypr}_{B'}^{+}}\\
[\mathrm{Edy}(w,z)\cdot\mathrm{Edy}(z,z')\cdot\mathrm{Edy}(z',y')+\mathrm{Edy}(w,z')\cdot\mathrm{Edy}(z',z)\cdot\mathrm{Edy}(z,y')].
\end{multline*}
\item[(D1)] One has
\[
\left|\partial_BD_{*}(A)\right|(y,y')\le226\,(C^{F})^{2}\cdot\epsilon_B\sum_{z\in\mathrm{Ypr}_B}\mathrm{Edy}(y,z)\cdot\mathrm{Edy}(z,y');
\]
consequently
\begin{align*}
\left\|\partial_BD_{*}(A)\right\|_{\frac{\gamma^{F}}{2}}&\le C_K\cdot\epsilon_B,\\
\sum_{y,y'}\left|\partial_BD_{*}(A)\right|(y,y')&\le C_K\cdot\epsilon_B\cdot\#|\mathrm{Ypr}_B|.
\end{align*}
\item[(D2)] One has
\begin{multline*}
\left|\partial_B\partial_{B'}D_{*}(A)\right|(y,y')\le C_H^{2}\cdot\epsilon_B\cdot\epsilon_{B'}\sum_{z\in\mathrm{Ypr}_B}\ \sum_{z'\in\mathrm{Ypr}_{B'}}\\
[\mathrm{Edy}(y,z)\cdot\mathrm{Edy}(z,z')\cdot\mathrm{Edy}(z',y')+\mathrm{Edy}(y,z')\cdot\mathrm{Edy}(z',z)\cdot\mathrm{Edy}(z,y')].
\end{multline*}
\item[(D3)] If $A_0$ and $A_0+B$ lie in $\mathrm{Gwk}$ together with the segment joining them, then
\begin{align*}
\left\|D_{*}(A_0+B)-D_{*}(A_0)\right\|_{\frac{\gamma^{F}}{2}}&\le C_K\cdot\epsilon_B,\\
\sum_{y,y'}\left|D_{*}(A_0+B)-D_{*}(A_0)\right|(y,y')&\le C_K\cdot\epsilon_B\cdot\#|\mathrm{Ypr}_B|.
\end{align*}
\end{itemize}
The bounds (L1)--(L3) hold for barred as well as for unbarred legs.
\end{named}

\begin{named}{Remark}
Proposition 33.2 corresponds to \cite[Lemma~20.15]{QEDI} together with the paragraph on legs that precedes it in \cite[Section~20.4]{QEDI}, and clause (D3) corresponds to \cite[Corollary~20.16]{QEDI}.
The numerical constants of Proposition~33.2 arise from the constant $146$ of Proposition~33.1\,(b) through the following computations:
$146\,(C^{F})^{2}+\frac{9}{4}\,C^{F}\le149\,(C^{F})^{2}$;
$149+11=160$;
$4M\cdot C^{F}+\frac{4}{3}\cdot149\,(C^{F})^{2}\le203M\cdot(C^{F})^{2}$;
$146^2+329=21645$;
the coefficient of a chain in (L3) is at most $\left(4M\cdot149+11\cdot149+\frac{4}{3}\cdot21645\right)\cdot(C^{F})^{3}\le31095\,M\cdot(C^{F})^{3}\le C_H^{2}$;
the coefficient in (D1) is $\frac{9}{4}+\frac{3}{2}\cdot149\le226$;
the coefficient in (D2) is $\frac{9}{4}\cdot149+\frac{3}{2}\cdot21645\le32803$, and $32803\,(C^{F})^{3}\le C_H^{2}$.
The proof of \cite[Lemma~20.15]{QEDI} starts from the constant $40$ of \cite[Lemma~20.14(b)]{QEDI} in the place of $146$; the numbers in the same places are there $43$, $54$, $62M$ and $1600+90=1690$, the coefficients of a chain in (L3) add up to $2899$, and the coefficients in (D2) add up to $2632$, while the coefficient in (D1) is $67$ in \cite[Lemma~20.15]{QEDI}.
All differences between these numbers and those of Proposition~33.2 come from the constant $146$ in the place of $40$.
The bounds (H1) and (H2) occur in \cite{QEDI} only as displays within the proof of \cite[Lemma~20.15]{QEDI}.
The leg indices are formed on $T^{\mathrm{f}}_{*}$, and $\mathrm{pos}$ is defined through $[c]$ and $d_X$.
\end{named}

\subsubsection*{Proposition~33.3 (a larger imaginary part)}

\begin{named}{Proposition 33.3}
Let $A+A_1\in\mathrm{Gwk}^{+}$, where $A\in\mathrm{Gwk}$ and $e_{*}\cdot|A_1|_{\infty}\le c_{\tau}$.
Then $\mathrm{Fst}(A+A_1)$ is invertible, and
\[
\left\|S_{*}(A+A_1)\right\|_{\frac{\gamma^{F}}{2}}\le2\,\mathbf{S}.
\]
The functions $S_{*}$ and $D_{*}(\cdot)$, defined on $\mathrm{Gwk}^{+}$ by the formulas of Section~G5.1.7, are analytic in a neighbourhood of every point of $\mathrm{Gwk}^{+}$.
If $A'$ and $A'+B$ lie in $\mathrm{Gwk}^{+}$ together with the segment joining them, then
\begin{align*}
\left\|D_{*}(A'+B)-D_{*}(A')\right\|_{\frac{\gamma^{F}}{2}}&\le C_K'\cdot\epsilon_B,\\
\sum_{y,y'}\left|D_{*}(A'+B)-D_{*}(A')\right|(y,y')&\le C_K'\cdot\epsilon_B\cdot\#|\mathrm{Ypr}_B|.
\end{align*}
\end{named}

\begin{named}{Remark}
Proposition 33.3 corresponds to \cite[Lemma~20.17]{QEDI}.
The smallness $146\,c_{\tau}\cdot\mathbf{S}\le\frac{1}{2}$ takes the place of $40\,c_1\cdot\mathbf{S}\le\frac{1}{2}$, with the constant $c_1$ of \cite[Section~20.1.7]{QEDI}; that is, the entry $\frac{1}{292\,\mathbf{S}}$ in the definition of $c_{\tau}$ takes the place of the entry $\frac{1}{80\,\mathbf{S}}$ in that of $c_1$.
The bound $\left\|\partial_BD_{*}(\cdot)\right\|_{\frac{\gamma^{F}}{2}}\le\left(13.5\,\mathbf{S}+1314\,\mathbf{S}^2\right)\cdot\epsilon_B\le C_K'\cdot\epsilon_B$ takes the place of the bound with $13.5\,\mathbf{S}+360\,\mathbf{S}^2$ in the proof of \cite[Lemma~20.17]{QEDI}.
\end{named}

\subsubsection*{Proposition~33.4 (block covariances)}

\begin{named}{Proposition 33.4}
\begin{itemize}
\item[(a)] Let $A$ be a real bond function on $T_{*}$.
Then $\mathrm{Gcv}:=D_{*}(A)^{-1}$ exists, $|\mathrm{Gcv}|_{\ell^2}\le\mathrm{gGa}^2$, and the matrix $\mathrm{Gcv}$ on $\mathrm{Ix}:=T^0_{*}\times\{1,\dots,4\}$ satisfies
\begin{align*}
\sup_{(y,\beta)}\sum_{(y',\beta')}\left|\mathrm{Gcv}(y,\beta;y',\beta')\right|\cdot\exp\left(\mathrm{vt}\cdot d(y,y')\right)&\le\mathrm{alG},\\
\sup_{(y',\beta')}\sum_{(y,\beta)}\left|\mathrm{Gcv}(y,\beta;y',\beta')\right|\cdot\exp\left(\mathrm{vt}\cdot d(y,y')\right)&\le\mathrm{alG};
\end{align*}
that is, the larger of the two suprema is at most $\mathrm{alG}$.
\item[(b)] Let $A=A_r+A_1$, where $A_r$ is a real and $A_1$ a complex bond function on $T_{*}$ with $e_{*}\cdot|A_1|_{\infty}\le\frac{\mu_{*}}{24}$, and assume that $\mathrm{Fst}(A)$ is invertible.
Then the operator norm of $D_A-D_{A_r}$ on $L^2$, with the scalar product $(\cdot,\cdot)^{\mathrm{ch}}$, is at most $12\,e_{*}\cdot|A_1|_{\infty}\le\frac{\mu_{*}}{2}$; the operator $D_A+\mu_{*}$ is invertible, with $\left|\left(D_A+\mu_{*}\right)^{-1}\right|_{L^2}\le\frac{2}{\mu_{*}}$; the operator $D_{*}(A)$ is invertible, and
\begin{align*}
D_{*}(A)^{-1}&=b_{*}^{-1}+Q_{*}(A)J^{*}\cdot\left(D_A+\mu_{*}\right)^{-1}\cdot JQ_{*}^{T}(-A),\\
\left|D_{*}(A)^{-1}\right|_{\ell^2}&\le b_{*}^{-1}+\frac{4}{\mu_{*}}\le\mathrm{gGa}_1.
\end{align*}
Every row of the matrix $D_{*}(A)^{-1}$ has $\ell^1$ norm at most $4\,\ell_{*}^3\cdot\mathrm{gGa}_1$.
\end{itemize}
\end{named}

\begin{named}{Remark}
Proposition 33.4 has the content of \cite[Lemma~20.18]{QEDI}; the formula for $D_{*}(A)^{-1}$ and the two intermediate bounds in (b) appear in \cite{QEDI} within the proof of that lemma.
In the estimate of $D_A-D_{A_r}$ the six terms $P^{\boldsymbol{\sigma}}_{\nu}\cdot\left(m^{\boldsymbol{\sigma}}_{\nu}-1\right)\cdot T^{\boldsymbol{\sigma}}_{\nu}(A_r)$ of the proof of \cite[Lemma~20.18]{QEDI} correspond here to the single term of (I.O6),
\[
\left(\left(D_A-D_{A_r}\right)f\right)(x)=-\eta_{*}^{-1}\cdot M_{3}(x)\cdot\left(u_{A_1}(c,c_{3})-1\right)\cdot u_{A_r}(c,c_{3})\cdot f\left(x^{(3)}\right),
\]
where $|M_{3}(x)|_2\le6$, the map $x\mapsto x^{(3)}$ is a bijection that preserves $\int\mathrm{d}x$, and $\left|u_{A_1}-1\right|\le2\,\eta_{*}\cdot e_{*}\cdot|A_1|_{\infty}$; the resulting constant $12$ is the same as in \cite{QEDI}.
The Cauchy--Schwarz inequality is applied in $\Delta_y\times O_{\mathrm{cub}}$, a set of measure one.
The $L^2$ and $\ell^2$ norms are formed with $|\cdot|_2$ (Section~I.1.5).
The constant $\mathrm{alG}$ is formed from the bound of Proposition~35.4\,(A7), that is, of Proposition~19\,(e), for the norm $|\cdot|_2$ of the blocks (Section~I.1.5).
\end{named}

\subsection*{34. The interacting factor; the functions $\mathfrak{P}_j$; the free ratio; one cell and the twist}

\subsubsection*{G5.4.0. Data and definitions for the Propositions 34.1--34.4 and 35.5--35.10}

\paragraph{(DAT)}
The setting is (ST*)$^{\mathrm{ch}}$; $\beta$ is a probe, the charge satisfies $0<e\le e_{\sharp}(\beta)$ (Section~G5.1.9), and $v$ is an admitted placement (Section~G5.1.3).
Wherever a block field occurs, $\theta\in\mathbb{R}^3$, $A\in V_{\mathrm{t}}(T^0_{*})$ with $\chi_{*}(A)=1$, and $(A^{g},\mathscr{A}^0)$ is a choice as in Proposition~35.4\,(SF5).
Put
\[
\mathscr{A}:=\mathscr{A}^0+c_{\theta}=H_{*}\left[A^{g}+c_{\theta}\right],\qquad\mathscr{T}(t):=\mathscr{A}+\mathscr{C}(t),\qquad\mathrm{Gcv}:=D_{*}(\mathscr{A})^{-1};
\]
here $H_{*}[\,\cdot\,]$ is the photon minimiser, as in the display (2.3) of Proposition~35.4\,(SF5), so that $\mathscr{A}$ is real and lies in $\frac{1}{4}\tilde{R}_{*}$, and $\mathrm{Gcv}$ exists by Proposition~33.4\,(a).
For $t\in\mathrm{DDo}(t_x^{+})$ let
\begin{equation}\tag{G5.5.1}
\begin{split}
\mathrm{Ztm}(t):=\int\mathrm{D}\Psi\,\exp\Bigl\{&-\left\langle\bar{\Psi},D_{*}\left(\mathscr{T}(t)\right)\Psi\right\rangle-m_{*}\cdot\left\langle\bar{\psi},\psi\right\rangle\\
&+\sum_{X\in D_{*}}\mathscr{E}_{*}\left(X;\mathscr{A},\mathscr{C}(t),\psi,\chi\right)\Bigr\},
\end{split}
\end{equation}
where the inserted fields are formed with $\mathscr{T}(t)$, as in Proposition~35.4\,(SF6).

In the Propositions 34.1--34.4 and 35.5--35.10 a bound is called uniform if it depends neither on $N\ge N_0(e)$, nor on $\theta$, nor on $A$, nor on the choice of $(A^{g},\mathscr{A}^0)$, nor on the admitted placement, and if it is unchanged when $\beta$ is replaced by $U_r\beta$, $r\in\{\pm1\}^3$.

\subsubsection*{Proposition~34.2 (the functions $\mathfrak{P}_j$)}

\begin{named}{Proposition 34.2}
Assume (ST*)$^{\mathrm{ch}}$, and let $0\le j\le k_{*}$ and $X\in D_j$.
The function $\mathfrak{P}_j(X,\cdot)$ is analytic on $\mathfrak{X}_j$, it vanishes at $0$, and $\mathfrak{P}_j(X,\mathscr{C})=\mathfrak{P}_j(X,\tilde{\mathscr{C}})$ whenever $\mathscr{C},\tilde{\mathscr{C}}\in\mathfrak{X}_j$ coincide on the bonds that meet $X$.
If $d_M(X)\ge1$, then
\[
\left|\mathfrak{P}_j(X,\mathscr{C})\right|\le2\,C_{d}\cdot M^3\cdot\exp\left(-\kappa\cdot d_M(X)\right)\qquad\text{for all }\mathscr{C}\in\mathfrak{X}_j.
\]
Here $\mathfrak{P}_0=0$, and $\mathfrak{P}_1$ contains $E_0^{\mathrm{det}}$, which comes from the determinant of the first step, a determinant of size $192\,n_0$.
\end{named}

\begin{named}{Remark}
Proposition 34.2 corresponds to \cite[Lemma~20.22]{QEDI}.
The constant $C_{d}$ is that of Proposition~13; it contains the row sum (I.L0.12) and the factor $48$ of the level $0$.
The proof of Proposition~34.2 uses the scaling relations of Section~I.3.7.
For $0\le j<k_{*}$ and every complex bond function $\mathscr{C}$ on $T_{j+1}$ the bond function $\mathscr{C}_L$ on $T_j$ is given by $\mathscr{C}_L(b)=L^{-\frac{1}{2}}\cdot\mathscr{C}(b/L)$, and
\[
|\mathscr{C}_L|_{\infty}=L^{-\frac{1}{2}}\cdot|\mathscr{C}|_{\infty},\qquad|\partial\mathscr{C}_L|_{\infty}=L^{-\frac{3}{2}}\cdot|\partial\mathscr{C}|_{\infty}.
\]
The set $\mathfrak{X}_j$ is defined by the thresholds $e_j^{-1+2\varepsilon}$ for $|\mathscr{C}|_{\infty}$ and $e_j^{-1+3\varepsilon}$ for $|\partial\mathscr{C}|_{\infty}$ (Section~I.5.3); since $0<e_j<e_{j+1}$ and the exponents $-1+2\varepsilon$, $-1+3\varepsilon$ are negative, the relations show that $\mathscr{C}\mapsto\mathscr{C}_L$ maps $\mathfrak{X}_{j+1}$ into $\mathfrak{X}_j$.
Hence the recursion of Proposition~35.4\,(SF4) defines $\mathfrak{P}_{j+1}(Y,\cdot)$ on $\mathfrak{X}_{j+1}$ for every $Y\in D_{j+1}$.
The same relations are contained in Proposition~35.4\,(DF5), which holds for every complex one-cochain, and in the proof of \cite[Lemma~17.5]{QEDI}; clause (c) of \cite[Lemma~17.5]{QEDI} states their consequence that $\mathscr{C}_L\in\mathfrak{X}_j$ whenever $\mathscr{C}\in\mathfrak{X}_{j+1}$, with $L^{-\frac{3}{4}}$ times the thresholds.
\end{named}

The conditions on the constants that enter are $L\cdot\left(\kappa-\kappa_0-1\right)\ge\kappa$, $c_{\mathrm{rb}}\cdot L^3\cdot\exp\left(-\left(L\cdot\left(\kappa-\kappa_0-1\right)-\kappa\right)\right)\le\frac{1}{2}$, the first condition of (A8), $0<\varepsilon<\frac{1}{3}$ and $L>1$.

\subsubsection*{Proposition~34.1 (the interacting factor; the functions $\mathrm{Rt}$ and $\mathrm{Lg}$)}

\begin{named}{Proposition 34.1}
Assume the data (DAT) of Section~G5.4.0, and let $t\in\mathrm{DDo}(t_x^{+})$.
\begin{itemize}
\item[(a)] One has $\mathscr{C}(t)\in\mathfrak{X}_{*}$ and $e_{*}\cdot|\mathscr{C}(t)|_{\infty}<e_{*}^{3\varepsilon}$, and the fields
\[
\mathscr{T}(t),\qquad c_{\theta}+\mathscr{C}(t),\qquad\mathscr{A}+\mathrm{sv}\cdot\mathscr{C}(t),\qquad c_{\theta}+\mathrm{sv}\cdot\mathscr{A}^0+\mathscr{C}(t)\qquad(0\le\mathrm{sv}\le1)
\]
lie in $\mathrm{Gwk}$.
\item[(b)] The operator $D_{*}\left(c_{\theta}+\mathscr{C}(t)\right)$ is invertible, and the function
\[
\mathrm{ff}(t):=\frac{\mathrm{Ztm}(t)}{\det D_{*}\left(c_{\theta}+\mathscr{C}(t)\right)}
\]
is analytic on $\mathrm{DDo}(t_x^{+})$ and satisfies the uniform bound
\[
\left|\mathrm{ff}(t)-1\right|\le\frac{1}{2}\,e_{*}^{\frac{1}{4}-9\varepsilon}\le\frac{1}{16}.
\]
In particular $\mathrm{Ztm}(t)\ne0$.
\item[(c)] One has
\begin{align*}
Y^{x}_{*}\left(A+c_{\theta};\mathscr{C}(t)\right)&=Y^{x}_{*}\left(A^{g}+c_{\theta};\mathscr{C}(t)\right)\\
&=\det D_{*}(0)^{-1}\cdot\exp\left\{-\varepsilon^{\mathrm{v},t}_{*}\cdot\ell_{*}^3+\sum_{X\in D_{*}}\mathfrak{P}_{*}\left(X,\mathscr{C}(t)\right)\right\}\cdot\mathrm{Ztm}(t)\ne0.
\end{align*}
The function
\begin{equation}\tag{6.1}
\mathrm{Rt}(t;\theta,A):=\frac{Y^{x}_{*}\left(A+c_{\theta};\mathscr{C}(t)\right)}{Y^{x}_{*}\left(A+c_{\theta};0\right)}=\exp\left\{\sum_{X\in D_{*}}\mathfrak{P}_{*}\left(X,\mathscr{C}(t)\right)\right\}\cdot\frac{\mathrm{Ztm}(t)}{\mathrm{Ztm}(0)}
\end{equation}
does not depend on the choice of $(A^{g},\mathscr{A}^0)$; it is continuous in $(t,\theta,A)$ on $\mathrm{DDo}(t_x^{+})\times\mathbb{R}^3\times\left\{A\in V_{\mathrm{t}}(T^0_{*}):\chi_{*}(A)=1\right\}$, analytic in $t$, and without zeros.
\item[(d)] The function
\begin{equation}\tag{6.2}
\mathrm{Lg}(t;\theta,A):=\int_0^1\mathrm{d}\,\mathrm{sv}\,\sum_{\mathfrak{c}\in\mathrm{Cel}}t_{\mathfrak{c}}\cdot\frac{\left(\partial_{t_{\mathfrak{c}}}\mathrm{Rt}\right)(\mathrm{sv}\cdot t;\theta,A)}{\mathrm{Rt}(\mathrm{sv}\cdot t;\theta,A)}
\end{equation}
is analytic in $t$ on $\mathrm{DDo}(t_x^{+})$ and continuous in $(t,\theta,A)$; it satisfies $\mathrm{Lg}(0;\theta,A)=0$ and $\exp\mathrm{Lg}=\mathrm{Rt}$, and it is the only continuous function of $t$ on $\mathrm{DDo}(t_x^{+})$ with these two properties.
\end{itemize}
\end{named}

\begin{named}{Remark}
Proposition 34.1 corresponds to the first paragraph of \cite[Section~20.5.1]{QEDI}, including the definition of the fermion integral given there, which is (G5.5.1), to \cite[Lemma~20.21]{QEDI}, and to \cite[Section~20.6.1]{QEDI} up to the definition of $\mathrm{Lg}$.
The assertions that \cite{QEDI} makes in the text between these lemmas, namely the representation (6.1), the continuity of $\mathrm{Rt}$ and the uniqueness of $\mathrm{Lg}$, are part of the statement here.
The weight $a_2$ of \cite[Section~20.2.4]{QEDI} is $C_0^{-1}\cdot h_{*,2}$.
The pairing $\left\langle\bar{\psi},\psi\right\rangle$ and the kernels of $\mathscr{E}_{*}$ are those of $T^{\mathrm{f}}_{*}$, with the weights $\frac{1}{48}$; the numbers $4\,\ell_{*}^3$ and $K_{0}\cdot\left(\frac{\ell_{*}}{M}\right)^3$ in the estimates are those of \cite{QEDI} (Section~I.5.1).
\end{named}

The conditions on the constants that enter are the three inequalities of (g3) and the first inequality of (g4) of Section~G5.1.9, the first condition of (A8), $\kappa\ge\kappa_0$, $C_{0}\ge1$, $M\ge1$, $\mathrm{gGa}_1\ge1$ and $\varepsilon>0$.

\subsubsection*{Proposition~34.3 (the free ratio on the disk of radius $t_x$)}

\begin{named}{Proposition 34.3}
Assume the data (DAT) without a block field, and let $\mathfrak{c}\in\mathrm{Cel}$.
For $|t|<t_x^{+}$ and $\theta\in\mathbb{R}^3$ put
\begin{gather*}
\mathrm{gD}^{\theta}(t):=\operatorname{Tr}\log\left[1+D_{*}(c_{\theta})^{-1}\cdot\left(D_{*}\left(c_{\theta}+t\cdot\hat{c}_{\mathfrak{c}}\right)-D_{*}(c_{\theta})\right)\right],\\
\mathrm{lgN}(t):=\sum_{X\in D_{*}}\mathfrak{P}_{*}\left(X,t\cdot\hat{c}_{\mathfrak{c}}\right)+\mathrm{gD}^{0}(t),
\end{gather*}
the trace being taken on $\mathfrak{S}(1,\ell_{*})$.
\begin{itemize}
\item[(a)] For $|t|<t_x^{+}$ one has $N_0(t)=\exp\mathrm{lgN}(t)$.
\item[(b)] For $|t|<t_x^{+}$,
\[
\left|\sum_{X\in D_{*}}\mathfrak{P}_{*}\left(X,t\cdot\hat{c}_{\mathfrak{c}}\right)\right|\le C_{\mathrm{fr}}\cdot e_{*}^{-2\varepsilon}\cdot|t|^2+B_{\mathrm{wr}}.
\]
\item[(c)] The operator $D_{*}(c_{\theta})^{-1}\cdot\left(D_{*}\left(c_{\theta}+t\cdot\hat{c}_{\mathfrak{c}}\right)-D_{*}(c_{\theta})\right)$ has $\ell^2$ operator norm at most $\frac{1}{2}$, so that the logarithm is defined by its power series.
One has $\left|\mathrm{gD}^{\theta}(t)\right|\le B_1$ and
\[
\exp\mathrm{gD}^{\theta}(t)=\frac{\det D_{*}\left(c_{\theta}+t\cdot\hat{c}_{\mathfrak{c}}\right)}{\det D_{*}(c_{\theta})}.
\]
For every $\theta\in\mathbb{R}^3$ the function $t\mapsto\mathrm{gD}^{\theta}(t)$ is analytic on $|t|<t_x^{+}$, and $\mathrm{gD}^{\theta}(0)=0$.
\item[(d)] For $|t|\le t_x$ one has $\left|\mathrm{lgN}(t)\right|\le B_0\le\frac{1}{4}$.
The function $\mathrm{lgN}$ is analytic on $|t|<t_x^{+}$, even, and real at real points, and $\mathrm{lgN}(0)=0$.
It does not depend on $\mathfrak{c}$, whereas $\mathrm{gD}^{\theta}$ is formed with the cell $\mathfrak{c}$.
\item[(e)] Let $\phi_0(t)=\sum_{i\le4}\mathrm{nj}_i\cdot t^i$ be the four-jet of $N_0$ at $0$, and let $c_{m}=\phi_0^{(m)}(0)$.
Then $\mathrm{nj}_0=1$ and $\mathrm{nj}_1=\mathrm{nj}_3=0$, so that $c_{1}=c_{3}=0$; the coefficients $\mathrm{nj}_i$ are real, $\left|\mathrm{nj}_i\right|\le2\,B_0\cdot t_x^{-i}$ for $1\le i\le4$, and
\begin{equation}\tag{6.4}
\begin{split}
&|c_{2}|\le\frac{\mathbf{r}^{-2}}{16},\qquad|c_{4}|\le\frac{\mathbf{r}^{-4}}{16},\qquad\|\phi_0\|_{\mathbf{r}}\le2,\\
&\|\phi_0-1\|_{\mathbf{r}}\le\frac{1}{6},\qquad c_{4}-3\,c_{2}^2=b_{\mathrm{ch}}\left(\eta,R^{*}(e)\right).
\end{split}
\end{equation}
Here $\|\cdot\|_{\mathbf{r}}$ is the norm $\|\cdot\|_{\mathrm{rad}}$ of Section~G5.1.6 with $\mathrm{rad}=\mathbf{r}$, applied to polynomials of degree at most $4$ in the single variable $t$: the norm of $\sum_{i\le4}p_i\cdot t^i$ is $\sum_{i\le4}|p_i|\cdot\mathbf{r}^i$.
Thus $\|\phi_0\|_{\mathbf{r}}=\sum_{i\le4}\left|\mathrm{nj}_i\right|\cdot\mathbf{r}^i$ and $\|\phi_0-1\|_{\mathbf{r}}=\sum_{1\le i\le4}\left|\mathrm{nj}_i\right|\cdot\mathbf{r}^i$.
\end{itemize}
All bounds are uniform.
\end{named}

\begin{named}{Remark}
Proposition 34.3 corresponds to \cite[Lemma~20.23]{QEDI}, the lemma on the free Wilson ratio on a disc, whose assertions \cite[Theorem~5.91]{QEDI} collects.
Here $N_0$ is formed with the determinant of size $192\,n_{*}^3$, and $b_{\mathrm{ch}}\left(\eta,R^{*}(e)\right)$ takes the place of $b_{\eta,R^{*}}(0)$.
The sum in (b) is made of the contributions of the levels $0\le j<k_{*}$, the level $0$ included on the same footing as the others.
The constants $C_{d}$, $C_K$, $C_K'$, and with them $C_{\mathrm{fr}}$, $B_{\mathrm{wr}}$, $B_1$ and $B_0$, have other values than in \cite{QEDI}.
\end{named}

The conditions on the constants that enter are the second and the fifth inequality of (g4), the condition (2.1) of Section~G5.1.9, that is $r_{*}>31L+16$, the first condition of (A8), $0<\varepsilon\le\frac{1}{200}$, $\kappa-1\ge\kappa_0$, $L\ge3$, $L^{1-\varepsilon}\ge2$, $M\ge2$, $C_K\le C_K'$, $\mathrm{gGa}^2\le\mathrm{gGa}_1$, $\mathrm{Cb}\ge1$ and $\gamma^{F}\ge0$.
The inequality $\gamma^{F}\ge0$ holds since $0<\gamma^{F}\le1$ (line (C4r) of Section~I.8.2; Proposition~14.2\,(f)).
Of the condition (sm2) of \cite[Section~20.1.8]{QEDI}, which \cite[Lemma~20.23]{QEDI} assumes, only the clause that the placement is admitted is needed, and here this clause is part of the data (DAT).

\subsubsection*{Proposition~34.4 (one cell; the twisted free ratio)}

\begin{named}{Proposition 34.4}
Assume the data (DAT), and let $\mathfrak{c}\in\mathrm{Cel}$.
\begin{itemize}
\item[(a)] (Complex twist.)
Fix $t$ with $|t|<t_x^{+}$.
For $\theta\in\mathbb{C}^3$ with $\left|\operatorname{Im}\theta_{\mu}\right|\le\tau_s$ ($\mu=1,2,3$) the connections $c_{\theta}+\mathrm{sv}\cdot t\cdot\hat{c}_{\mathfrak{c}}$, $0\le\mathrm{sv}\le1$, lie in $\mathrm{Gwk}^{+}$, and $D_{*}(c_{\theta})$ is invertible with $\left|D_{*}(c_{\theta})^{-1}\right|_{\ell^2}\le\mathrm{gGa}_1$.
The function $\theta\mapsto\mathrm{gD}^{\theta}(t)$, defined by the formula of Proposition~34.3, is analytic in a neighbourhood of the closed strip $\left\{\theta\in\mathbb{C}^3:\left|\operatorname{Im}\theta_{\mu}\right|\le\tau_s,\ \mu=1,2,3\right\}$, it is $2\pi$-periodic in each $\theta_{\mu}$, and $\left|\mathrm{gD}^{\theta}(t)\right|\le B_1$ there.
Consequently
\[
\left|\mathrm{gD}^{\theta}(t)-\mathrm{gD}^{0}(t)\right|\le52\,B_1\cdot\exp(-\tau_s)\le e_{*}^{\frac{1}{4}-9\varepsilon}\qquad\text{for all }\theta\in\mathbb{R}^3.
\]
\item[(b)] Let $\theta\in\mathbb{R}^3$, and let $A$ satisfy $\chi_{*}(A)=1$.
For $|t|<t_x^{+}$ put, with the principal branch of the logarithm,
\[
\mathrm{wdf}_{\mathfrak{c}}(t):=\mathrm{Lg}\left(t\cdot\mathrm{ue}_{\mathfrak{c}};\theta,A\right)-\mathrm{lgN}(t)=\mathrm{gD}^{\theta}(t)-\mathrm{gD}^{0}(t)+\log\mathrm{ff}\left(t\cdot\mathrm{ue}_{\mathfrak{c}}\right)-\log\mathrm{ff}(0).
\]
This function is analytic, it vanishes at $0$, and it satisfies the uniform bound $\left|\mathrm{wdf}_{\mathfrak{c}}(t)\right|\le3\,e_{*}^{\frac{1}{4}-9\varepsilon}$.
\end{itemize}
\end{named}

\begin{named}{Remark}
Proposition 34.4 corresponds to \cite[Proposition~20.24]{QEDI}.
Clause (a), the free factor at a complex twist, is Step~2 of the proof of \cite[Proposition~20.24]{QEDI}.
With $A_r=c_{\operatorname{Re}\theta}$ and $A_1=i\cdot c_{\operatorname{Im}\theta}$ one has
\[
e_{*}\cdot|A_1|_{\infty}\le\frac{\tau_s}{\ell_{*}}=c_{\tau}\le\frac{\mu_{-}}{48}<\frac{\mu_{*}}{24}.
\]
In that step the same decomposition gives $e_{*}\cdot|A_1|_{\infty}\le c_1$, with the constant $c_1$ of \cite[Section~20.1.7]{QEDI} in the place of $c_{\tau}$; this implies the bound $c_1+e_{*}^{2\varepsilon}$ for $e_{*}\cdot|A_1|_{\infty}$ that is used there.
\end{named}

The conditions on the constants that enter are the third inequality of (g3), the first four inequalities of (g4), $c_{\tau}\le\frac{\mu_{-}}{48}$, $\gamma^{F}\ge0$, the first condition of (A8) and $\varepsilon>0$, and, through Proposition~34.1, the conditions (g3) and the first inequality of (g4).
Here $\gamma^{F}\ge0$ follows from $0<\gamma^{F}\le1$ (line (C4r) of Section~I.8.2; Proposition~14.2\,(f)).

\subsection*{Propositions 35.m, second part. Several cells; the complement of the small-field set; the representation; $\Theta^{(0)}_m$; $e_{\sharp}(\beta)$}

\subsubsection*{Proposition~35.5 (the terminal fermion integral as a function of the marks)}

Let the data (DAT) of Section~G5.4.0 be given, together with a block field.
For $X\in D_{*}$ and $t\in\mathrm{DDo}\left(t_x^{+}\right)$ put
\begin{equation*}
\mathrm{Ea}(X;t):=\mathscr{E}_{*}\left(X;\mathscr{A},\mathscr{C}(t),\,\cdot\,\right)-\mathbf{1}[X\text{ is an }M\text{-cube}]\cdot m_{*}\cdot\langle\bar{\psi},\psi\rangle_X\in\mathrm{Gr}[\psi,\chi],
\end{equation*}
and denote by $\mathrm{Ea}_{nm}(X;t;\,\cdot\,)$ its antisymmetric kernels.
Further let
\begin{align*}
\mathrm{Kq}(t)&:=D_{*}\left(\mathscr{T}(t)\right)-D_{*}(\mathscr{A}),\\
\mathrm{Vj}(t)&:=-\left\langle\bar{\Psi},\mathrm{Kq}(t)\Psi\right\rangle+\sum_{X\in D_{*}}\mathrm{Ea}(X;t)=:\mathrm{Vj}_0(t)+\mathrm{Vj}_+(t),\tag{5.2}
\end{align*}
where the activities in the sum are taken at $(\psi,\chi)=H_{*}^{\#}\left(\mathscr{T}(t)\right)\Psi$ and at the barred fields formed in the same way.
Here $\mathrm{Vj}_0(t)$ is the part of degree zero of $\mathrm{Vj}(t)$, namely
\begin{equation*}
\mathrm{Vj}_0(t)=\sum_{X}\mathscr{E}_{00}\left(X;\mathscr{A},\mathscr{C}(t)\right),
\end{equation*}
and $\mathrm{Vj}_+(t)=\mathrm{Vj}(t)-\mathrm{Vj}_0(t)$ is an even element of $\mathrm{Gr}[\Psi]$ without constant term.
For a cell $\mathfrak{c}$ write $\partial_{\mathfrak{c}}:=\partial/\partial t_{\mathfrak{c}}$;
on an element of a Grassmann algebra whose coefficients are analytic in $t$ this operator acts coefficient by coefficient.

\begin{named}{Proposition 35.5}
Assume the data (DAT) with a block field, and let $X\in D_{*}$.
\begin{itemize}
\item[(a)] For $0<\mathrm{aw}\le e_{*}^{-\frac{1}{4}+\varepsilon}$ and $t\in\mathrm{DD}(t_x)$,
\begin{equation*}
\sum_{n,m}\frac{\mathrm{aw}^{n+m}}{n!\,m!}\cdot\left\|\mathrm{Ea}_{nm}(X;t)\right\|\le w^{\mathrm{E}}_{*}\cdot\exp\left(-\kappa\cdot d_M(X)\right).
\end{equation*}
\item[(b)] For the same $\mathrm{aw}$, for $t\in\mathrm{DD}(t_x/2)$ and for cells $\mathfrak{c}\ne\mathfrak{c}'$,
\begin{align*}
\sum_{n,m}\frac{\mathrm{aw}^{n+m}}{n!\,m!}\cdot\left\|\partial_{\mathfrak{c}}\mathrm{Ea}_{nm}(X;t)\right\|&\le q_x\cdot w^{\mathrm{E}}_{*}\cdot\exp\left(-\kappa\cdot d_M(X)\right),\\
\sum_{n,m}\frac{\mathrm{aw}^{n+m}}{n!\,m!}\cdot\left\|\partial_{\mathfrak{c}}\partial_{\mathfrak{c}'}\mathrm{Ea}_{nm}(X;t)\right\|&\le q_x^2\cdot w^{\mathrm{E}}_{*}\cdot\exp\left(-\kappa\cdot d_M(X)\right).
\end{align*}
\item[(c)] $\partial_{\mathfrak{c}}\mathrm{Ea}_{nm}(X;t)=0$ unless $X$ contains a point of $\mathrm{Ypr}_{\mathfrak{c}}$.
\item[(d)] All coefficients of $\mathrm{Vj}(t)$ are analytic functions of $t$ on $\mathrm{DDo}\left(t_x^{+}\right)$, and
\begin{equation*}
\mathrm{Ztm}(t)=\det D_{*}(\mathscr{A})\cdot\exp\left(\mathrm{Vj}_0(t)\right)\cdot\int\mathrm{d}\mu_{\mathrm{Gcv}}\,\exp\left(\mathrm{Vj}_+(t)\right).\tag{5.3}
\end{equation*}
\item[(N1)] For $t\in\mathrm{DD}(t_x)$, without any restriction on $\hat{\mathfrak{a}}$,
\begin{equation*}
\mathrm{rw}_G\left(\mathrm{Vj}_+(t);\mathrm{wg}_{\delta}\right)\le\bar{r}.
\end{equation*}
\item[(N2)] If $\hat{\mathfrak{a}}\ne0$, then for $t\in\mathrm{DD}(t_x/2)$ and $\mathfrak{c}\in\mathrm{Cel}$
\begin{equation*}
\mathrm{Nw}_G\left(\partial_{\mathfrak{c}}\mathrm{Vj}_+(t);\mathrm{wg}_{\mathrm{Ypr}_{\mathfrak{c}}}\right)\le q_x\cdot\mathrm{CVj}\cdot h_{G}^2\cdot\mathrm{pp},\qquad
\mathrm{rw}_G\left(\partial_{\mathfrak{c}}\mathrm{Vj}_+(t);\mathrm{wg}_{\mathrm{Ypr}_{\mathfrak{c}}}\right)\le q_x\cdot\mathrm{CVj}\cdot h_{G}^2.
\end{equation*}
\item[(N3)] If $\hat{\mathfrak{a}}\ne0$, then for $t\in\mathrm{DD}(t_x/2)$ and cells $\mathfrak{c}\ne\mathfrak{c}'$
\begin{equation*}
\mathrm{Nw}_G\left(\partial_{\mathfrak{c}}\partial_{\mathfrak{c}'}\mathrm{Vj}_+(t);1\right)+\left|\partial_{\mathfrak{c}}\partial_{\mathfrak{c}'}\mathrm{Vj}_0(t)\right|\le q_x^2\cdot\mathrm{CVj}\cdot h_{G}^2\cdot\mathrm{pp}\cdot\exp\left(-\mathrm{vt}\cdot\left(d_{\mathrm{sep}}-4M-6\right)\right).
\end{equation*}
\item[(N0)] If $\hat{\mathfrak{a}}=0$, then $\mathscr{C}(t)=0$ and $\mathrm{Kq}(t)=0$, and $\mathrm{Ea}(X;t)$, $\mathrm{Vj}_0(t)$ and $\mathrm{Vj}_+(t)$ do not depend on $t$.
\end{itemize}
The bounds in (a), (b), (N1) and (N2) are uniform in the sense of Section~G5.4.0.
The right side of (N3) depends on the choices listed in Section~G5.4.0 only through $d_{\mathrm{sep}}$;
it is a decreasing function of $d_{\mathrm{sep}}$, and it does not exceed the uniform number
\begin{equation*}
q_x^2\cdot\mathrm{CVj}\cdot h_{G}^2\cdot\mathrm{pp}\cdot\exp\left(-\mathrm{vt}\cdot\left(\frac{M\cdot r_{*}}{8L}-4M-7\right)\right).
\end{equation*}
\end{named}

In connection with Proposition~35.5 the labels (N0)--(N3) denote its clauses; they are not the rules (I.N0)--(I.N6) of Section~I.5.2.

\begin{named}{Remark}
Proposition~35.5 corresponds to \cite[Lemmas~20.19 and~20.20]{QEDI}, together with the part of \cite[Section~20.5.1]{QEDI} that begins with the paragraph on the mass term in the activities, including there the definition of the interaction and the Gaussian form of the terminal integral;
the clauses (a)--(c) correspond to \cite[Lemma~20.19]{QEDI}, and (N1)--(N3) to the bounds (V1)--(V3) of \cite[Lemma~20.20]{QEDI}.
Compared with these places of \cite{QEDI}, the following holds.
In Step~2 of the proof of \cite[Lemma~20.20]{QEDI} the derivative of $D_{*}$ is estimated by the first inequality of (D$'$) in \cite[Lemma~20.15]{QEDI}, which carries the numeral $67$;
for the present operator the corresponding inequality is Proposition~33.2\,(D1), which carries the numeral $226$,
and the inequality $C_K\le C_H\cdot\mathrm{ssum}\left(\gamma^{F}/2\right)^2$ that is applied in that step holds for the values of Section~G5.1.8, because $226\le203\,M$.
In Step~3 of the proof of \cite[Lemma~20.20]{QEDI} the bound for each product given by the product rule contains one of two exponential factors, according to whether a single leg carries both derivatives or not;
the proof of Proposition~35.5 bounds every product with the sum of these two factors, and therefore has the numeral $256$ where the proof of \cite[Lemma~20.20]{QEDI} has $128$;
since $h_{G}\ge1$, one has $16\,h_{G}^2+256\le10^3\,h_{G}^2$.
Expressed through $\mathrm{CVj}$, $\bar{r}$ and $\mathrm{pp}$, the bounds have the same form as in \cite[Lemmas~20.19 and~20.20]{QEDI}.
The leg indices range over $T^{\mathrm{f}}_{*}$ and $\mathrm{Pz}_{*}$, with the weights of Section~I.3.6;
in Step~0\,(ii) of the proof of \cite[Lemma~20.20]{QEDI} a pair of kind~(II) is $\left(x,x^{(3)}\right)$, with $d_X\left(x,x^{(3)}\right)=\eta_{*}<1$, and $c$ or $c_{3}$ lies in $X$.
As in \cite[Section~20.5.1]{QEDI}, the norm of the mass term is $4\,M^3\cdot|m_{*}|\cdot\mathrm{aw}^2$ (Section~I.5.1).
As to the ranges of the clauses, Proposition~35.5 agrees with \cite[Lemma~20.20]{QEDI}:
(N1), like (V1), is asserted without restriction on $\hat{\mathfrak{a}}$, and (N2) and (N3), like (V2) and (V3), for $\hat{\mathfrak{a}}\ne0$ only;
(N0) is the counterpart of the last sentence of \cite[Lemma~20.20]{QEDI}, which settles the case $\hat{\mathfrak{a}}=0$.
Beyond that sentence, (N0) states that $\mathrm{Kq}(t)=0$ and that $\mathrm{Ea}(X;t)$ does not depend on $t$;
in \cite{QEDI} these two facts stand in the first paragraph of the proof of \cite[Lemma~20.20]{QEDI}.
The non-emptiness of the sets $\mathrm{Bd}_{\mathfrak{c}}$, $\mathrm{Ypr}_{\mathfrak{c}}$ and $\mathrm{Ypr}_{\mathfrak{c}}^{+}$ for $\hat{\mathfrak{a}}\ne0$, which \cite[Lemma~20.20]{QEDI} states before (V2), is Proposition~35.3\,(g).
As in \cite[Section~20.3.4]{QEDI} and in the last sentence of \cite[Section~20.5.1]{QEDI}, where the weight is denoted $w_P$, the weight $\mathrm{wg}_P$ of Section~G5.1.6 is defined for non-empty sets $P$ only;
for this reason (N2), like (V2), is restricted to $\hat{\mathfrak{a}}\ne0$.
Finally, the fact, needed in the proof of \cite[Lemma~20.20]{QEDI}, that every block $\Delta_y$ lies in one $M$-cube is Proposition~35.4\,(SF4)\,(v).
\end{named}

The proof of Proposition~35.5 requires the following conditions on the constants:
the first and the second inequality of (g1);
(g2);
the third inequality of (g3);
the first inequality of (A8), whence $e_{*}\le e_0<1$;
$\gamma^{F}\cdot M\ge16$;
$M\ge2$, $M$ odd;
$C^{F}\ge1$;
$\gamma^{F}\le1$;
$\mathrm{vt}\le\min\left(\gamma^{F}/8,1/M\right)$;
$\kappa-1\ge\kappa_{0}$;
$K_{0}\ge1$;
$\varepsilon>0$;
and, among the inequalities of Section~G5.9.3, the inequality $C_K\le C_H\cdot\mathrm{ssum}\left(\gamma^{F}/2\right)^2$, the two bounds
\begin{equation*}
8\,C_K+3\,K_{0}+1200\,C_H\cdot\mathrm{ssum}\left(\gamma^{F}/2\right)\cdot K_{0}\le\mathrm{CVj},\qquad
4\,C_K+82\,K_{0}+1200\,C_H\cdot\mathrm{ssum}\left(\gamma^{F}/2\right)^2\cdot K_{0}\le\mathrm{CVj},
\end{equation*}
and the inequality
\begin{equation*}
\left(16\,h_{G}^2+256\right)\cdot C_H^2\cdot\mathrm{ssum}\left(\gamma^{F}/2\right)^4\cdot K_{0}\le\mathrm{CVj}\cdot h_{G}^2.
\end{equation*}

\subsubsection*{Proposition~35.6 (several cells)}

\begin{named}{Proposition 35.6}
Assume the data (DAT) with a block field.
For every subset $\mathrm{Scl}$ of $\mathrm{Cel}$ with $\#|\mathrm{Scl}|\ge2$,
\begin{equation*}
\sup_{t\in\mathrm{DD}(t_x/2)}\left|\mathrm{Lg}_{\mathrm{Scl}}(t;\theta,A)\right|\le\varpi_{\mathrm{sep}}\le\varpi_{\mathrm{sep}}^{+};
\end{equation*}
these bounds are uniform, except that $\varpi_{\mathrm{sep}}$ depends on $d_{\mathrm{sep}}$.
If $\hat{\mathfrak{a}}=0$, then $\mathrm{Lg}=0$.
\end{named}

\begin{named}{Remark}
Proposition~35.6 corresponds to \cite[Proposition~20.25]{QEDI}, together with the display of \cite[Section~20.6.1]{QEDI} that introduces the function $W(t)$ and expresses the logarithm $\mathrm{Lg}(t)$ through it, and with the paragraph preceding that display.
The case $\hat{\mathfrak{a}}=0$ and the second inequality, which follows from Proposition~35.3\,(e), are not part of \cite[Proposition~20.25]{QEDI}.
In the application of \cite[Lemma~20.11]{QEDI} the data $\left(C,g,h_G,\boldsymbol{\alpha}_{\vartheta}\right)$ of that lemma are $\left(\mathrm{Gcv},\mathrm{gGa},h_{G},\mathrm{alG}\right)$.
\end{named}

Of the conditions on the constants, the proof of Proposition~35.6 requires
the third inequality of (g1);
(g2);
$\kappa-\kappa_{0}\ge2$;
$\mathrm{vt}\le1/M$;
and the inequality
\begin{equation*}
2\,\mathrm{CVj}\cdot h_{G}^2+2\,\mathrm{alG}\cdot\mathrm{CVj}^2\cdot h_{G}^4+2\,\mathrm{e}^4\cdot C_{d}\cdot M^3\cdot K_{0}\le\mathrm{CWj}
\end{equation*}
of Section~G5.9.3.

\subsubsection*{Proposition~35.7 (entire functions; the complement of the small-field set; the measure $\mathrm{vm}$; the remainder)}

In the setting (DAT) put, for $t\in\mathbb{C}^{\mathrm{Cel}}$,
\begin{equation*}
\mathrm{Ysf}(t):=\bar{\nu}^{a(t)}(\mathfrak{G}),\qquad\mathrm{zL}(t):=\bar{\nu}^{a(t)}\left(F_N\setminus\mathfrak{G}\right),
\end{equation*}
and let $\mathrm{vm}$ be the measure on $[-\pi,\pi]^3\times{V_{\mathrm{t}}}\left(T^0_{*}\right)$ given by
\begin{equation*}
\mathrm{vm}(\mathrm{d}\theta\,\mathrm{d}A):=\bar{\nu}(F_N)^{-1}\cdot\chi_{*}(A)\cdot Y^{x}_{*}\left(A+c_{\theta};0\right)\cdot(2\pi)^{-3}\,\mathrm{d}\theta\,\mathrm{d}\mathsf{m}_{*}(A).
\end{equation*}
Here, and in the display (2.5) of Proposition~35.4\,(SF6), a product $\chi_{*}(A)\cdot Y^{x}_{*}\left(A+c_{\theta};\,\cdot\,\right)$ is read as $0$ where $\chi_{*}(A)=0$:
the function $Y^{x}_{*}$ is defined on $\mathrm{Dom}_{*}$ (Section~G5.1.5), and by Proposition~35.4\,(SF5) the field $A+c_{\theta}$ lies in $\mathrm{Dom}_{*}$ wherever $\chi_{*}(A)=1$.
Functions that are defined only where $\chi_{*}(A)=1$, such as $\mathrm{Lg}$ and $\mathrm{kj}_{\mathrm{Scl}}$, are integrated over $[-\pi,\pi]^3\times\{\chi_{*}=1\}$;
the complement of this set is a null set for $\mathrm{vm}$.

\begin{named}{Proposition 35.7}
Assume the data (DAT).
\begin{itemize}
\item[(a)] For every $F^{\mathrm{u}}\in F_N$ and every $\theta\in\mathbb{R}^3$ the function $t\mapsto\mathrm{dens}^{\theta,a(t)}_N\left(F^{\mathrm{u}}\right)$ is entire on $\mathbb{C}^{\mathrm{Cel}}$;
on every compact set of values of $t$ it is bounded uniformly in $\left(F^{\mathrm{u}},\theta\right)$, and it is continuous in $\left(t,F^{\mathrm{u}},\theta\right)$.
The function $N_0$ is entire on $\mathbb{C}$, and for every Borel set $\mathrm{Ev}\subset F_N$ the function $t\mapsto\bar{\nu}^{a(t)}(\mathrm{Ev})$ is entire on $\mathbb{C}^{\mathrm{Cel}}$.
Moreover $0<\bar{\nu}(F_N)<\infty$.
\item[(b)] $\bar{\mu}\left(F_N\setminus\mathfrak{G}^{-}\right)\le\varpi_0$.
\item[(c)] For every $t\in\mathrm{DDo}\left(t_x^{+}\right)$ the map $(\theta,A)\mapsto Y^{x}_{*}\left(A+c_{\theta};\mathscr{C}(t)\right)$ is continuous on $\mathbb{R}^3\times\{\chi_{*}=1\}$;
consequently the density of $\mathrm{vm}$ is a bounded Borel function.
The measure $\mathrm{vm}$ is positive; its total mass is $\bar{\mu}(\mathfrak{G})$, and $\bar{\mu}(\mathfrak{G})\in[1-\varpi_0,1]$.
\item[(d)] For $t\in\mathrm{DDo}\left(t_x^{+}\right)$,
\begin{equation*}
\frac{\mathrm{Ysf}(t)}{\bar{\nu}(F_N)}=\int\mathrm{vm}(\mathrm{d}\theta\,\mathrm{d}A)\,\exp\mathrm{Lg}(t;\theta,A).
\end{equation*}
\item[(e)] For every real $t\in\mathrm{DDo}\left(t_x^{+}\right)$, every $\theta\in\mathbb{R}^3$ and every $A$ with $\chi_{*}(A)=1$ the number $\mathrm{Lg}(t;\theta,A)$ is real.
\item[(f)] For every real $t\in\mathrm{DD}(t_2)$,
\begin{equation*}
0\le\frac{\mathrm{zL}(t)}{\bar{\nu}(F_N)}\le\varpi_L.
\end{equation*}
\item[(g)] For all $z\in\mathbb{C}^{\mathrm{Cel}}$ and all $t\in\mathbb{C}^{\mathrm{Cel}}$,
\begin{equation*}
\mathring{\Phi}(e\cdot z)=\bar{\mu}\left(\exp\sum_{\mathfrak{c}}\left(z_{\mathfrak{c}}\cdot X_{\mathfrak{c}}-\frac{z_{\mathfrak{c}}^2\cdot\boldsymbol{\sigma}^2}{2}\right)\right),\qquad
\left|\mathring{\Phi}(t)\right|\le\exp\left(\sum_{\mathfrak{c}}\frac{\left(\operatorname{Im}t_{\mathfrak{c}}\right)^2\cdot\boldsymbol{\sigma}^2}{2e^2}\right);
\end{equation*}
for $t\in\mathrm{DD}(t_2)$ one has $\left|\exp\mathrm{Lg}(t;\theta,A)\right|\le2$ and $\left|\mathrm{Ysf}(t)\right|\le2\,\bar{\nu}(F_N)$.
\item[(h)] For $t\in\mathrm{DD}\left(\frac{t_2}{2N_{*}}\right)$,
\begin{equation*}
\left|\frac{\mathrm{zL}(t)}{\bar{\nu}(F_N)}\right|\le\varpi_L^{\frac{1}{3}}\cdot\mathrm{mxL}^{\frac{2}{3}}.
\end{equation*}
\item[(i)] The image $\mathrm{Lr}$ of $\mathrm{zL}/\bar{\nu}(F_N)$ in $\mathcal{J}$ has real coefficients, and $\|\mathrm{Lr}\|_{\mathbf{r}}\le\lambda_L^{N_{*}}$.
\end{itemize}
\end{named}

\begin{named}{Remark}
Proposition~35.7 corresponds to the paragraph of \cite[Section~20.1.5]{QEDI} that immediately precedes \cite[Lemma~20.4]{QEDI}, to \cite[Lemma~20.26]{QEDI}, to the definitions that precede \cite[Lemma~20.27]{QEDI}, to \cite[Lemma~20.27]{QEDI} and to \cite[Proposition~20.28]{QEDI}, with the same numerical constants;
clause (a) corresponds to that paragraph, (b) to \cite[Lemma~20.26]{QEDI}, (c)--(e) to \cite[Lemma~20.27]{QEDI}, and (f)--(i) to \cite[Proposition~20.28]{QEDI}.
It differs from them in three points.
First, in (a) the entries of the matrix of the operator are entire functions of the numbers $A(b)$, by (I.O1), and only the hop of type~3 carries a phase;
this takes the place of the description in \cite[Section~20.1.5]{QEDI} by polynomials in the link phases with coefficients of modulus at most $3$.
Analyticity in several variables is obtained from analyticity in each variable together with continuity, by Osgood's lemma, or from Cauchy's formula under the integral.
Second, the inequality $\bar{\nu}(F_N)>0$ and the Gaussian domination of $\bar{\mu}$ are those of Proposition~31.1, which hold for every real $\boldsymbol{\mu}_N$;
the condition $\boldsymbol{\mu}_N>-1$ on the lattice mass, on which the proofs of \cite[Lemma~20.4(a)]{QEDI} and \cite[Proposition~20.5]{QEDI} rest for $\bar{\nu}(F_N)>0$ and for the Gaussian domination, is not needed here.
Third, the continuity of the map $(\theta,A)\mapsto Y^{x}_{*}\left(A+c_{\theta};\mathscr{C}(t)\right)$ on $\mathbb{R}^3\times\{\chi_{*}=1\}$ and the fact that the density of $\mathrm{vm}$ is a bounded Borel function are part of the statement, in clause (c);
in \cite{QEDI} they stand in the paragraph that precedes \cite[Lemma~20.27]{QEDI}, where the set $\{\chi_{*}=1\}$ is shown to be compact, and in the proof of \cite[Lemma~20.27(b)]{QEDI}.
The positivity of $\mathrm{vm}$ follows from Tonelli's theorem, as in the proof of \cite[Lemma~20.27(a)]{QEDI}, or else from a pointwise positivity of the kind stated in (e).
\end{named}

The conditions on the constants that enter the proof of Proposition~35.7 are
the first two clauses of (T1);
the third inequality of (g3);
(g5), (g6) and (g7);
$0<s<1$;
$C_0\ge1$;
$\rho\ge0$, which is part of $\rho\in(0,1)$ (line (C0)) and enters the proof of clause (a) through (I.m6), in the form $|M_{i}(x)|_2=\frac{w_i}{2}\,(1+\rho)$;
the standard weights (line (C0)), which enter the proof of clause (a) through the closure relation (I.1.1), by way of (I.m1) and the second identity of (I.2.5), through the inequalities $w_i>0$, by (I.m6), and through $w_3=6$ and $\mathrm{sw}=w_0+w_1+w_2+w_3$, in the bound of the norm of the operator $D_B+\boldsymbol{\mu}_N$ and in the resulting bound of $\mathrm{dens}^{\theta,a(t)}_N$;
$\bar{m}\in(0,1]$ (line (C10)), a datum of the setting (ST*)$^{\mathrm{ch}}$ of Section~G5.1.1, with both of its bounds:
$\bar{m}>0$ gives $\bar{m}\,\eta>0$, which is used for $\det\left(D_0+\bar{m}\,\eta\right)>0$ (Proposition~31.1\,(a)),
and $\bar{m}\le1$ gives $h\ge1$, $h$ being an integral power of $L$ (Section~I.0.5);
$L\ge5$, whence $\sum_{i\ge1}L^{-4i}\le1$, and, together with $h\ge1$ and with the facts $T^{*}(e)\ge1$ and $N\ge0$ on the scales (Section~I.6.2), $n_*=h\cdot L^{T^{*}(e)+N}\ge L\ge5$;
that $p_0$ and $p$ are integers $\ge1$;
and $N_{*}\ge8$.

\subsubsection*{Proposition~35.8 (the four-jet representation)}

The data (DAT) being given, the following elements of $\mathcal{J}$ are defined for $\theta\in\mathbb{R}^3$ and $A$ with $\chi_{*}(A)=1$:
\begin{equation*}
\mathrm{kj}_{\{\mathfrak{c}\}}(\theta,A;\,\cdot\,):=\exp\left(\mathrm{wdf}_{\mathfrak{c}}(t_{\mathfrak{c}})\right)-1,\qquad
\mathrm{kj}_{\mathrm{Scl}}(\theta,A;\,\cdot\,):=\exp\left(\mathrm{Lg}_{\mathrm{Scl}}\right)-1\quad\left(\#|\mathrm{Scl}|\ge2\right).
\end{equation*}

\begin{named}{Proposition 35.8}
Assume the data (DAT).
For every $N\ge N_0(e)$ one has, in $\mathcal{J}$,
\begin{equation*}
\mathring{\Phi}(t)=\prod_{\mathfrak{c}\in\mathrm{Cel}}\phi_0(t_{\mathfrak{c}})\cdot\int\mathrm{vm}(\mathrm{d}\theta\,\mathrm{d}A)\,\prod_{\emptyset\ne\mathrm{Scl}\subset\mathrm{Cel}}\left(1+\mathrm{kj}_{\mathrm{Scl}}(\theta,A;t)\right)+\mathrm{Lr}(t),\tag{8.1}
\end{equation*}
where the integral of an element of $\mathcal{J}$ whose coefficients are bounded Borel functions of $(\theta,A)$ is formed coefficient by coefficient, and where the following hold.
\begin{itemize}
\item[(a)] $\phi_0$ is the four-jet of the free ratio $N_0$;
it is even, its coefficients are real, $\phi_0(0)=1$, and the relations (6.4) of Proposition~34.3\,(e) hold.
\item[(b)] $\mathrm{vm}$ is a positive measure of total mass $\bar{\mu}(\mathfrak{G})\in[1-\varpi_0,1]$.
\item[(c)] The coefficients of $\mathrm{kj}_{\mathrm{Scl}}(\theta,A;\,\cdot\,)$ are real and depend continuously on $(\theta,A)$ on $\mathbb{R}^3\times\{\chi_{*}=1\}$;
every monomial of $\mathrm{kj}_{\mathrm{Scl}}(\theta,A;\,\cdot\,)$ contains each variable $t_{\mathfrak{c}}$ with $\mathfrak{c}\in\mathrm{Scl}$, and no other variable.
\item[(d)] For every cell $\mathfrak{c}$,
\begin{equation*}
\sup_{\theta,A}\left\|\mathrm{kj}_{\{\mathfrak{c}\}}\right\|_{\mathbf{r}}\le\bar{\kappa}_1=6\,e_{*}^{\frac{1}{4}-9\varepsilon}.
\end{equation*}
\item[(e)] For $\#|\mathrm{Scl}|\ge2$,
\begin{equation*}
\sup_{\theta,A}\left\|\mathrm{kj}_{\mathrm{Scl}}\right\|_{\mathbf{r}}\le2\,\varpi_{\mathrm{sep}}\le2\,\varpi_{\mathrm{sep}}^{+}\le\bar{\kappa}_1^{N_{*}}.
\end{equation*}
\item[(f)] $\|\mathrm{Lr}\|_{\mathbf{r}}\le\lambda_L^{N_{*}}$, and the coefficients of $\mathrm{Lr}$ are real.
\item[(g)] $\mathbf{r}=\frac{t_x}{4}=\frac{e_{*}^{3\varepsilon}}{4\,\mathrm{Cb}}$.
\end{itemize}
In (d) and (e) the suprema extend over $\theta\in\mathbb{R}^3$ and over the $A$ with $\chi_{*}(A)=1$.
The bounds in (6.4), the numbers $\varpi_0$, $\bar{\kappa}_1$, $\varpi_{\mathrm{sep}}^{+}$, $\bar{\kappa}_1^{N_{*}}$, $\lambda_L^{N_{*}}$ and the number $\mathbf{r}$ do not depend on $N$, on $\theta$, on the admitted placement, nor on a replacement of $\beta$ by $U_r\beta$ with $r\in\{\pm1\}^3$.
\end{named}

\begin{named}{Remark}
Proposition~35.8 corresponds to \cite[Theorem~20.29]{QEDI}, together with the definitions that precede it in \cite[Section~20.8]{QEDI};
the same representation is stated in \cite[Theorem~5.90]{QEDI}.
Here $\phi_0$ is the four-jet of the free ratio $N_0$ of Section~G5.1.5, formed with the determinant of size $192\,n_*^3$, in the place of the free Wilson ratio.
In place of the hypotheses (sm1)--(sm7) of \cite[Theorem~20.29]{QEDI}, which are smallness conditions of \cite[Section~20.1.8]{QEDI}, stand the conditions (g1)--(g7) of Section~G5.1.9;
these hold because $e\le e_{\sharp}(\beta)$, the definition of the number $e_{\sharp}(\beta)$ containing (g1)--(g9).
The elements $\mathrm{kj}_{\mathrm{Scl}}$ are defined only where $\chi_{*}(A)=1$.
That $p_{*}$ and $p_{0*}$ do not depend on $N$ is Proposition~35.1\,(c).
Clause (e) contains the bound by $\varpi_{\mathrm{sep}}^{+}$ of Section~G5.1.8, which does not depend on the placement.
\end{named}

The proof of Proposition~35.8 requires, of the conditions on the constants, the third inequality of (g3), the condition (g5) and the second inequality of (g7).

\subsubsection*{Proposition~35.9 (the numbers $\Theta^{(0)}_m$; the reference datum (R3))}

\begin{named}{Proposition 35.9}
Assume the data (DAT).
For $m\in\{1,2,3,4\}$ let $c_{m}=\phi_0^{(m)}(0)$; thus $c_1=c_{3}=0$.
\begin{itemize}
\item[(a)] One has
\begin{equation*}
\Theta^{(0)}_m\le m!\cdot\mathbf{r}^{-m}\cdot\left(2\,\mathrm{e}\cdot\bar{\kappa}_1^s+\lambda_L\right)\le8\,m!\cdot\mathbf{r}^{-m}\cdot\bar{\kappa}_1^s.
\end{equation*}
\item[(b)] Let $\mathfrak{e}=s\cdot\left(\frac{1}{4}-9\varepsilon\right)-12\varepsilon$ (Section~I.8.5).
Then
\begin{align*}
\Theta^{(0)}_4&\le3\cdot10^5\cdot\mathrm{Cb}^4\cdot e_{*}^{\mathfrak{e}}\le3\cdot10^5\cdot\mathrm{Cb}^4\cdot h^{\frac{1}{8}}\cdot e^{\mathfrak{e}},\\
\left(1+|c_2|\right)\cdot\Theta^{(0)}_2&\le3\cdot10^4\cdot\mathrm{Cb}^4\cdot e_{*}^{\mathfrak{e}}\le3\cdot10^4\cdot\mathrm{Cb}^4\cdot h^{\frac{1}{8}}\cdot e^{\mathfrak{e}}.
\end{align*}
\item[(c)] One has
\begin{equation*}
|c_2|+|c_4|^{\frac{1}{2}}\le8\,\mathrm{Cb}^2\cdot e_{*}^{-6\varepsilon}\le8\,\mathrm{Cb}^2\cdot\left(\frac{L^2}{h}\right)^{3\varepsilon}\cdot e^{-6\varepsilon}.
\end{equation*}
\end{itemize}
In particular the proposition holds for every centred placement;
for such a placement $\mathring{\Phi}$, $N_0$, $c_{m}$ and $\Theta^{(0)}_m$ ($m=2,4$) are the objects of Section~I.7.6.
\end{named}

\begin{named}{Remark}
Proposition~35.9 corresponds to \cite[Proposition~20.31]{QEDI} and \cite[Corollary~20.32]{QEDI}, which are combined in \cite[Theorem~5.92]{QEDI}, and to the assertions (a) and (b) in the proof of \cite[Theorem~21.3]{QEDI} in \cite[Section~21.9]{QEDI}.
The hypothesis $\lambda_L\le\bar{\kappa}_1^s$ of \cite[Proposition~20.31]{QEDI} and \cite[Corollary~20.32]{QEDI} is the condition (g8), which holds because $e\le e_{\sharp}(\beta)$.
Clause (b) contains in addition the bounds in powers of $e$, obtained from Proposition~35.1\,(a) with $0<\mathfrak{e}\le\frac{1}{4}$ and $h\ge1$;
so (b) is the reference datum (R3) of Section~I.7.6, with the constant $3\cdot10^5\cdot\mathrm{Cb}^4\cdot h^{\frac{1}{8}}$.
Clause (c) is (R4) in the form of \cite[Corollary~20.32]{QEDI}, with the exponent $6\varepsilon$;
the form of (R4) with $\mathfrak{e}_2=0$ (Section~I.8.5) is Proposition~37.3\,(b).
The numerals $3\cdot10^5$, $3\cdot10^4$, $8$ and $2\,\mathrm{e}$ are those of \cite[Proposition~20.31]{QEDI} and \cite[Corollary~20.32]{QEDI};
in (b) and (c) the factor in front of each power of $e_{*}$ is the product of one of these numerals with a power of $\mathrm{Cb}$, and contains no other constant.
For these bounds, what differs from \cite[Proposition~20.31]{QEDI} and \cite[Corollary~20.32]{QEDI} is the range $e\le e_{\sharp}(\beta)$ of the charge.
\end{named}

The conditions on the constants that enter the proof of Proposition~35.9 are
(g5), through Proposition~35.8\,(e);
(g6);
(g8);
for the last sentence, the first inequality of (g9);
$0<s<1$;
$\sigma=\mathfrak{e}-\frac{2a_{L}}{N_{*}}>0$, which is (S2) of Proposition~14.2\,(b) and gives $\mathfrak{e}>0$ in (b);
$\varepsilon>0$, from (S1) of Proposition~14.2\,(b), whence, with $\mathfrak{e}>0$ and $s>0$, also $\frac{1}{4}-9\varepsilon>0$;
$e_1\le\frac{1}{3}$ (Proposition~14.2\,(f));
the condition (L5) of Proposition~14.2\,(b), in the forms $L\ge1$ (so that $\eta\le1$), $N_{*}=8L^3\ge2$ and $a_{L}>0$;
and $h\ge1$, which holds because $\bar{m}\le1$ (Section~I.0.5).

\subsubsection*{Proposition~35.10 ($e_{\sharp}(\beta)>0$)}

\begin{named}{Proposition 35.10}
Let the constants of the lines (C0)--(C9) and (C6$'$) be fixed, let $\bar{m}\in(0,1]$, and let $\beta$ be a probe.
Then $e_{\sharp}(\beta)>0$.
This number depends only on these constants (among them $\rho$), on $\bar{m}$, and on $\beta$ through $\left(C_{\beta},\mathbf{D},D_\beta\right)$;
it does not depend on $N$, on a torus, on a placement or on $\theta$.
For $0<e\le e_{\sharp}(\beta)$ and every $N\ge N_0(e)$ the following hold:
$e\cdot h^{\frac{1}{2}}\le e_1$;
every centred placement is admitted;
$D_\beta+3\eta\le\ell_c$, that is, $\beta$ fits $\left(N,T^{*}(e)\right)$;
$R^{*}(e)\ge D_{\mathrm{be}}+1$;
and (g1)--(g9) hold.
\end{named}

\begin{named}{Remark}
Proposition~35.10 corresponds to \cite[Lemma~20.33]{QEDI}, which is stated also as \cite[Theorem~5.93]{QEDI}, together with the smallness conditions (sm1)--(sm8) of \cite[Section~20.1.8]{QEDI}, whose place is taken by the conditions (g1)--(g9).
Here $e_{\sharp}(\beta)$ is a definite number, and its definition (Section~G5.1.9) contains (g1)--(g9).
The conditions do not refer to a placement:
the number $M\cdot r_{*}/(8L)-1$ stands in the place of $d_{\mathrm{sep}}$ in (g2) and, through $\varpi_{\mathrm{sep}}^{+}$, in (g5).
The inequality $R^{*}(e)\ge D_{\mathrm{be}}+1$ has no counterpart in \cite[Lemma~20.33]{QEDI};
here $R^{*}(e)=2L\cdot\ell_c$.
The finiteness of $C_{\beta}$ and $\mathbf{D}$ is Proposition~35.3\,(c1), which presupposes neither a placement nor a pair $(e,N)$.
\end{named}

The proof of Proposition~35.10 rests on the following conditions on the constants:
$r\ge2$;
$p_0\ge2$, which follows from $p_0>3r+2$;
$p>2p_0$;
$0<\varepsilon\le\bar{\varepsilon}=\frac{1}{200}$ (the lines (C0), (C7) and (S1) of Proposition~14.2\,(b));
$0<s<1$;
the conditions (L5) and (M1) of Proposition~14.2\,(b), that is, $L$ is odd and $L\ge5$, so that $L>1$, and $M$ is a power of $L$ with $M\ge L^2$;
$0<e_1\le\frac{1}{3}$ (Propositions~14.2\,(f) and~14.3\,(f));
the finiteness and the positivity of the constants of the lines (C0), (C4), (C6) and (C6$'$) that occur in (g1)--(g9), those of the line (C6) being $C_{0}$ and $C_{d}$,
which follow from the ranges
\begin{equation*}
C^{F}\ge1,\qquad0<\gamma^{F}\le1,\qquad C_{0}\ge\max\left(1,C^{F}\right),\qquad K_{0}\ge1,\qquad C_{d}<\infty
\end{equation*}
of (G5.T-1) of Section~G5.7 (a consequence of Proposition~14.1 and of Proposition~14.2\,(a), (f)),
from the inequality $C_{w}\ge C^{F}$ of Proposition~14.2\,(f),
and from the paragraph ``Ranges'' of Proposition~14.2\,(f), by which $C^{F}$, $C_{0}$ and $C_{w}$ are finite and $C_{d}$ is positive,
so that $1\le C_{0}<\infty$ and $0<C_{d}<\infty$;
$h\ge1$, which holds because $\bar{m}\le1$, and the integrality of $\log_L h$ (Section~I.0.5);
and, through Proposition~35.1, the conditions (e1) of Proposition~14.2\,(b).

\subsection*{Propositions 37.m and 36. The Bloch symbol; the two-photon coefficient; the four-photon coefficient}

Propositions~37.1, 36 and~37.2 are statements about the free field.
They involve neither the charge nor a cell nor the flow, and the scale factor $L$ enters them only through the sequence of spacings $L^{-N}$ in the last clause of Proposition~37.2\,(f);
the number $q>0$ in Proposition~36\,(h) is arbitrary.
Proposition~37.3 specialises them to the reference torus.

\subsubsection*{G5.6.0. Setting and definitions of the free-field statements}

\paragraph{The setting (FF).}
Let $\rho\in(0,1)$, with the standard weights of Section~I.0.5;
let $\bar{m}\in(0,1]$;
let $\beta$ be a probe (Section~G5.1.3), with $a=d^{*}\beta$ and $f=dd^{*}\beta$;
let $\eta\in(0,1]$, and let $R>0$ be such that $n:=R/\eta$ is an integer $\ge2$.
Lattice units are used throughout, on the cubic torus $C(n,n,n)$ with the spinors $\mathfrak{S}^{\mathrm{ch}}(1;n)$ of Proposition~46, a space of dimension $192\,n^3$.
A one-cochain on $\mathbb{Z}^3/n\mathbb{Z}^3$ is an $n\mathbb{Z}^3$-periodic cochain on $\mathbb{Z}^3$, as in Proposition~46;
thus for $n=2$ the bonds $(c,c+e_\nu)$ and $(c,c-e_\nu)$ are different bonds.
The lattice probe $a_{\eta}[\beta]$ is the one of Section~G5.1.3, formed with $\beta$ itself and not with a translate of $\beta$.
Put
\begin{gather*}
B_{t}(c,c+e_\nu):=t\cdot\eta\cdot a_{\eta}[\beta]_\nu(\eta\cdot c)\quad(t\in\mathbb{C}),\qquad
N_0(t):=\frac{\det\left(D^{\mathrm{ch}}[B_{t}]+\bar{m}\cdot\eta\right)}{\det\left(D^{\mathrm{ch}}[0]+\bar{m}\cdot\eta\right)},\qquad
c_{m}:=N_0^{(m)}(0),\\
a_{\mathrm{ch}}(\eta,R):=\left.\frac{\mathrm{d}^2}{\mathrm{d}t^2}\log N_0(t)\right|_{t=0},\qquad
b_{\mathrm{ch}}(\eta,R):=\left.\frac{\mathrm{d}^4}{\mathrm{d}t^4}\log N_0(t)\right|_{t=0}.
\end{gather*}
Here $D^{\mathrm{ch}}$ is the operator (I.2.3) in lattice units, by Proposition~46\,(b).
With $\hat{\mathfrak{a}}=\eta^{\frac{1}{2}}\cdot a_{\eta}[\beta]$, any $e>0$ and the charge $q=e\cdot\eta^{\frac{1}{2}}$ of level $0$ (Section~I.2.5) one has $B_{t}=q\cdot(t/e)\cdot\hat{\mathfrak{a}}$,
and $N_0$ is the free ratio of Section~I.7.6 on $C(n,n,n)$, formed with $\beta$;
this is Proposition~46\,(g+) at $v=0$.
Where the probe has to be named, one writes $a_{\mathrm{ch}}(\eta,R)[\beta]$, $b_{\mathrm{ch}}(\eta,R)[\beta]$ and $N_0[\beta]$.

\paragraph{Fibre.}
The space ${V_{\mathrm{v}}}$ of Section~I.2.6 carries the inner product
\begin{equation*}
(u',u):=\frac{1}{48}\sum_{g}\overline{u'(g)}\cdot u(g);
\end{equation*}
$|O|_2$ is the operator norm on ${V_{\mathrm{v}}}$, and $\mathrm{tr}$ is the trace on ${V_{\mathrm{v}}}$, a space of dimension $192$;
the maps $J\colon\mathbb{C}^4\to{V_{\mathrm{v}}}$ and $J^{*}\colon{V_{\mathrm{v}}}\to\mathbb{C}^4$ are those of (I.3.1).
For a number $\mathbf{m}\ge0$ and $p\in\mathbb{C}^3$ let $D^{(p)}$ be the symbol (I.2.6) at $\boldsymbol{\mu}=\mathbf{m}$, and $D_0^{(p)}$ the symbol at $\boldsymbol{\mu}=0$.
Moreover let
\begin{gather*}
(R_iu)(g):=u(gs_i)\quad(i\le2),\qquad
\left(R_3(p)u\right)(g):=\exp\left(i\,p\cdot(ge_1)\right)\cdot u(gs_3),\qquad
\hat{d}_i(g):=-g\hat n_i,\\
\mathrm{Rsy}(p):=\left(D^{(p)}\right)^{-1},\qquad
\mathrm{Vx}_\nu(p):=\frac{\partial D^{(p)}}{\partial p_\nu},\\
\left(\mathrm{Ehp}^{\pm}_\nu u\right)(g):=\mathbf{1}[ge_1=\pm e_\nu]\cdot3\left(\rho\mp\gamma_\nu\right)u(gs_3),\qquad
(\mathrm{Pa}\,u)(g):=\gamma_5\,u(-g).
\end{gather*}
The operator $\mathrm{Rsy}(p)$ is defined wherever the inverse exists;
for $\mathbf{m}>0$ and real $p$ it exists by Proposition~37.1\,(a).

\paragraph{Continuum.}
For $\mathrm{kph}\in\mathbb{R}^3$ let
\begin{equation*}
\hat{a}(\mathrm{kph}):=\int\mathrm{d}y\,a(y)\cdot\exp\left(-i\,\mathrm{kph}\cdot y\right);
\end{equation*}
$S_{\bar{m}}$ is the operator of Section~I.2.6, and $b(f)$ the number of Section~I.7.6.
Put
\begin{gather*}
\mathrm{vpo}_{\bar{m}}(\mathrm{ts}):=\frac{1}{\pi}\int_0^1\mathrm{d}\,\mathrm{sv}\;\mathrm{sv}\cdot(1-\mathrm{sv})\cdot\left(\bar{m}^2+\mathrm{sv}\cdot(1-\mathrm{sv})\cdot\mathrm{ts}\right)^{-\frac{1}{2}}\quad(\mathrm{ts}\ge0),\\
\mathrm{pv}(f):=-\int_{\mathbb{R}^3}\frac{\mathrm{d}^3\mathrm{kph}}{(2\pi)^3}\,|\mathrm{kph}|_2^2\cdot\mathrm{vpo}_{\bar{m}}\left(|\mathrm{kph}|_2^2\right)\cdot\left|\hat{a}(\mathrm{kph})\right|_2^2\quad(\le0).
\end{gather*}
Let $L^2_{\mathrm{ap}}\left(T_R;\mathbb{C}^4\right)$ be the space of the locally square integrable functions $\psi\colon\mathbb{R}^3\to\mathbb{C}^4$ with $\psi\left(y+R\cdot e_\mu\right)=-\psi(y)$, normed by the $L^2$ norm over a period cube,
and let $S^R:=\left(\bar{m}+\gamma\cdot\partial\right)^{-1}$ on this space, that is, the multiplier $\left(\bar{m}+i\gamma\cdot p\right)^{-1}$ on the momenta $p\in\frac{2\pi}{R}\left(\mathbb{Z}+\frac{1}{2}\right)^3$.
For $R\ge D_{\mathrm{be}}+1$ put $a^R:=\sum_{z\in\mathbb{Z}^3}a\left(\,\cdot\,+R\cdot z\right)$, the translates of $\operatorname{supp}\beta$ being then pairwise disjoint, and
\begin{equation*}
b_{R}(f):=-6\operatorname{Tr}_{L^2_{\mathrm{ap}}\left(T_R;\mathbb{C}^4\right)}\left[\left(S^R\,i\gamma\cdot a^R\right)^4\right].
\end{equation*}

\subsubsection*{Proposition~37.1 (the Bloch symbol and its inverse)}

\begin{named}{Proposition 37.1}
Let $\rho\in(0,1)$; the objects of the fibre are those of Section~G5.6.0.
\begin{itemize}
\item[(a)] For real $p$ the symbol without mass splits as
\[
D_0^{(p)}=\frac{\rho}{2}\,(\nabla^{*}\nabla)(p)+\frac{1}{2}\,\mathrm{Gh}(p),
\]
where
\begin{gather*}
(\nabla^{*}\nabla)(p):=\sum_{i\le2}w_i\cdot\left(1-R_i\right)+w_3\cdot\left(1-R_3(p)\right),\\
\left(\mathrm{Gh}(p)u\right)(g):=\sum_{i\le2}w_i\cdot\left(\gamma\cdot\hat{d}_i(g)\right)\cdot u(gs_i)+w_3\cdot\left(\gamma\cdot\hat{d}_3(g)\right)\cdot\exp\left(ip\cdot(ge_1)\right)\cdot u(gs_3).
\end{gather*}
Here $(\nabla^{*}\nabla)(p)=(\nabla^{*}\nabla)(p)^{*}\ge0$ and $\mathrm{Gh}(p)^{*}=-\mathrm{Gh}(p)$.
In particular $\operatorname{Re}\,\left(u,D^{(p)}u\right)\ge\mathbf{m}\cdot(u,u)$ for every $u\in{V_{\mathrm{v}}}$.
\item[(b)] For real $p\notin2\pi\mathbb{Z}^3$ the kernel of $(\nabla^{*}\nabla)(p)$ is $\{0\}$.
At $p=0$ one has $\ker(\nabla^{*}\nabla)(0)=J\,\mathbb{C}^4$, and the smallest non-zero eigenvalue of $(\nabla^{*}\nabla)(0)$ is the number $\lambda_1$ of (C0r).
\item[(c)] Both $D_0^{(0)}J=0$ and $J^{*}D_0^{(0)}=0$ hold, and for every $p$
\[
J^{*}D_0^{(p)}J=i\,\gamma\cdot\sin p+\rho\sum_\mu\left(1-\cos p_\mu\right).
\]
\item[(d)] There is a matrix $D_{0,\mathrm{loc}}$, independent of $p$, with
\[
D_0^{(p)}=D_{0,\mathrm{loc}}-\sum_\nu\left(\exp(ip_\nu)\cdot\mathrm{Ehp}^{+}_{\nu}+\exp(-ip_\nu)\cdot\mathrm{Ehp}^{-}_{\nu}\right).
\]
For all complex numbers $z^{\pm}_\nu$,
\[
\left|\sum_\nu\left(z^{+}_\nu\cdot\mathrm{Ehp}^{+}_{\nu}+z^{-}_\nu\cdot\mathrm{Ehp}^{-}_{\nu}\right)\right|_2\le3(1+\rho)\cdot\max_{\nu,\pm}\left|z^{\pm}_\nu\right|,
\qquad
J^{*}\mathrm{Ehp}^{\pm}_{\nu}J=\frac{1}{2}\left(\rho\mp\gamma_\nu\right).
\]
Consequently
\[
\mathrm{Vx}_\nu(p)=-i\left(\exp(ip_\nu)\cdot\mathrm{Ehp}^{+}_{\nu}-\exp(-ip_\nu)\cdot\mathrm{Ehp}^{-}_{\nu}\right),\qquad
\partial_\nu^2D^{(p)}=\exp(ip_\nu)\cdot\mathrm{Ehp}^{+}_{\nu}+\exp(-ip_\nu)\cdot\mathrm{Ehp}^{-}_{\nu},
\]
the mixed second derivatives of $D^{(p)}$ vanish, and $J^{*}\mathrm{Vx}_\nu(0)J=i\gamma_\nu$.
For real $p,p'$ one has $\left|\partial^{\mathrm{ak}}D^{(p)}\right|_2\le6$ for every multi-index $\mathrm{ak}\ne0$, and $\left|D^{(p)}-D^{(p')}\right|_2\le6\,|p-p'|_\infty$.
\item[(e)] (Parity.)
The map $\mathrm{Pa}$ is a unitary involution of ${V_{\mathrm{v}}}$, and
\[
\mathrm{Pa}\,D^{(p)}\,\mathrm{Pa}=D^{(-p)},\qquad\mathrm{Pa}\,\mathrm{Vx}_\nu(p)\,\mathrm{Pa}=-\mathrm{Vx}_\nu(-p).
\]
\item[(f)] (One light species, slope one, no additive mass.)
There is a number $C_E=C_E(\rho)<\infty$ such that for all $\mathbf{m}>0$ and all $p\in[-\pi,\pi]^3$ the matrix $D^{(p)}$ is invertible and
\[
\mathrm{Rsy}(p)=J\left(\mathbf{m}+i\gamma\cdot p\right)^{-1}J^{*}+\mathrm{Er}(p,\mathbf{m}),\qquad\left|\mathrm{Er}(p,\mathbf{m})\right|_2\le C_E.
\]
It follows that $\left|\mathrm{Rsy}(p)\right|_2\le C_E+\left(\mathbf{m}^2+|p|_2^2\right)^{-\frac{1}{2}}$, and, when $\mathbf{m}\le1$,
\[
\left|\mathrm{Rsy}(p)\right|_2\le\frac{\sqrt{2}+7C_E}{|p|_2+\mathbf{m}}.
\]
\item[(g)] For every multi-index $\mathrm{ak}$ there is $\mathrm{Csy}_{\mathrm{ak}}(\rho)<\infty$ such that
\[
\left|\partial^{\mathrm{ak}}\mathrm{Rsy}(p)\right|_2\le\mathrm{Csy}_{\mathrm{ak}}(\rho)\cdot\left(\mathrm{dist}_2\left(p,2\pi\mathbb{Z}^3\right)+\mathbf{m}\right)^{-1-|\mathrm{ak}|}
\]
for all real $p$ and all $0<\mathbf{m}\le1$; here $\mathrm{dist}_2$ denotes the Euclidean distance.
\item[(h)] (Scaling limit.)
Let $\bar{m}\in(0,1]$, let $\eta_j\in(0,1]$ with $\eta_j\to0$, and let $\mathrm{kph}_j\to\mathrm{kph}\in\mathbb{R}^3$.
At the index $j$ let $\mathrm{Rsy}$ be formed with $\mathbf{m}=\bar{m}\cdot\eta_j$.
Then
\[
\eta_j\cdot\mathrm{Rsy}\left(\eta_j\cdot\mathrm{kph}_j\right)\to J\left(\bar{m}+i\gamma\cdot\mathrm{kph}\right)^{-1}J^{*},
\]
and, for every multi-index with $|\mathrm{ak}|\le2$,
\[
\eta_j^{1+|\mathrm{ak}|}\cdot\left(\partial^{\mathrm{ak}}\mathrm{Rsy}\right)\left(\eta_j\cdot\mathrm{kph}_j\right)\to J\,\partial^{\mathrm{ak}}\left(\bar{m}+i\gamma\cdot\mathrm{kph}\right)^{-1}J^{*}.
\]
\end{itemize}
The constants $C_E(\rho)$ and $\mathrm{Csy}_{\mathrm{ak}}(\rho)$ obtained in the proof contain the minimum of the lowest eigenvalue of $(\nabla^{*}\nabla)(p)$ over the set $\left\{p\in[-\pi,\pi]^3:\ |p|_2\ge\delta(\rho)\right\}$, where
\[
\delta(\rho):=\min\left\{\frac{\lambda_1}{12},\ \left(2+\frac{288}{\rho\cdot\lambda_1}\right)^{-1}\right\}>0;
\]
this minimum comes from a compactness argument, and the constants are not explicit.
Every number $C_E(\rho)$ for which (f) holds satisfies $C_E(\rho)\ge\frac{2}{\rho\cdot\lambda_1}$, and the numbers $\mathrm{Csy}_{\mathrm{ak}}(\rho)$ obtained in the proof tend to infinity as $\rho\to0$.
\end{named}

\begin{named}{Remark}
Proposition~37.1 corresponds to the first and the third bound of \cite[Lemma~21.8]{QEDI};
the second bound there, $\left|(m+i\gamma\cdot p)^{-1}\right|\le\sqrt{2}/(m+|p|)$, concerns the continuum and remains as it stands.
In \cite[Lemma~21.8]{QEDI} the symbol $m+\mathrm{Wl}(p)+i\gamma\cdot\operatorname{sn}(p)$ of the Wilson operator is a normal $4\times4$ matrix whose inverse is known in closed form;
$D^{(p)}$ is a $192\times192$ matrix for which no such form is available, and (f) and (g) take the place of the two bounds of \cite[Lemma~21.8]{QEDI}.
At $\mathbf{m}=\bar{m}\cdot\eta$ with $\bar{m},\eta\in(0,1]$ and at $p=\eta\cdot\mathrm{kph}$, clause (f) gives
\[
\left|\eta\cdot\mathrm{Rsy}\left(\eta\cdot\mathrm{kph}\right)-J\left(\bar{m}+i\gamma\cdot\mathrm{kph}\right)^{-1}J^{*}\right|_2\le C_E\cdot\eta\qquad\left(|\mathrm{kph}|_\infty\le\frac{\pi}{\eta}\right).
\]
This is the third bound of \cite[Lemma~21.8]{QEDI}, read with $m=\bar{m}$ and $p=\mathrm{kph}$, with the factor $\eta$ in the place of $15\,\eta^{\frac{1}{2}}\cdot(m+|p|)^{-\frac{1}{2}}$, apart from the geometric mean that \cite[Lemma~21.8]{QEDI} takes.
Clauses (a) and (c) are the forms in the fibre of (I.2.4), (I.O2) and (I.O8).
\end{named}

\subsubsection*{Proposition~36 (the two-photon coefficient)}

For Proposition~36 the setting is (G5.FF), and the symbols of the fibre of Section~G5.6.0 are formed with the mass $\mathbf{m}=\bar{m}\cdot\eta$.
On $\eta\mathbb{Z}^3$, without periodisation, the lattice probe is
\[
\mathrm{a0}_{\eta,\nu}(y):=\mathrm{gs}_\nu[\beta]\left(y+\frac{\eta}{2}e_\nu\right)\qquad\left(y\in\eta\mathbb{Z}^3\right),
\]
and its transform is the finite sum
\[
\mathrm{ahl}_{\eta,\nu}(\mathrm{kph}):=\eta^3\sum_{y\in\eta\mathbb{Z}^3}\mathrm{a0}_{\eta,\nu}(y)\cdot\exp\left(-i\,\mathrm{kph}\cdot\left(y+\frac{\eta}{2}e_\nu\right)\right)\qquad\left(\mathrm{kph}\in\mathbb{R}^3\right).
\]
Let
\[
\mathrm{Kx}_{\eta,R}:=\left\{\frac{2\pi}{R}\,j:\ j\in\mathbb{Z}^3,\ -\frac{n}{2}<j_\nu\le\frac{n}{2}\ \text{for}\ \nu=1,2,3\right\},
\]
and let $P_{n,0}$ be the set of momenta of Section~I.2.6.
For $p,p'\in\mathbb{R}^3$ put
\[
\mathrm{lp}_{\mu\nu}(p,p'):=\delta_{\mu\nu}\operatorname{tr}\left[\mathrm{Rsy}(p)\,\partial_\nu^2D^{(p)}\right]-\operatorname{tr}\left[\mathrm{Rsy}(p)\,\mathrm{Vx}_\mu\left(p+\frac{p'}{2}\right)\mathrm{Rsy}(p+p')\,\mathrm{Vx}_\nu\left(p+\frac{p'}{2}\right)\right],
\]
and define the lattice vacuum polarisation on the torus and on the infinite lattice by
\[
\mathrm{Vpl}^{\eta,R}_{\mu\nu}(\mathrm{kph}):=\eta^{-1}\cdot n^{-3}\sum_{p\in P_{n,0}}\mathrm{lp}_{\mu\nu}\left(p,\eta\cdot\mathrm{kph}\right),\qquad
\mathrm{Vpl}^{\eta}_{\mu\nu}(\mathrm{kph}):=\eta^{-1}\int_{[-\pi,\pi]^3}\frac{\mathrm{d}^3p}{(2\pi)^3}\,\mathrm{lp}_{\mu\nu}\left(p,\eta\cdot\mathrm{kph}\right).
\]

\begin{named}{Proposition 36}
Assume the setting (G5.FF).
\begin{itemize}
\item[(a)] The function $N_0$ is entire in $t$, with $N_0(0)=1$ and $N_0(-t)=N_0(t)$, and $N_0(t)>0$ for real $t$.
Hence $c_1=c_3=0$, and the numbers $a_{\mathrm{ch}}(\eta,R)=c_2$ and $b_{\mathrm{ch}}(\eta,R)=c_4-3c_2^2$ are real.
\item[(b)] (Exact identity.)
\[
a_{\mathrm{ch}}(\eta,R)=R^{-3}\sum_{\mathrm{kph}\in\mathrm{Kx}_{\eta,R}}\ \sum_{\mu,\nu}\overline{\mathrm{ahl}_{\eta,\mu}(\mathrm{kph})}\cdot\mathrm{ahl}_{\eta,\nu}(\mathrm{kph})\cdot\mathrm{Vpl}^{\eta,R}_{\mu\nu}(\mathrm{kph}).
\]
\item[(c)] There is $\mathrm{Cpv}_0=\mathrm{Cpv}_0(\rho)<\infty$ such that, for all $\bar{m}\in(0,1]$, all real $\mathrm{kph}$, all $\eta\in(0,1]$ and all $R$ as in (G5.FF),
\begin{gather*}
\left|\mathrm{Vpl}^{\eta,R}_{\mu\nu}(\mathrm{kph})-\mathrm{Vpl}^{\eta}_{\mu\nu}(\mathrm{kph})\right|\le\mathrm{Cpv}_0\cdot\bar{m}^{-3}\cdot R^{-4},\\
\left|\mathrm{Vpl}^{\eta}_{\mu\nu}(\mathrm{kph})\right|\le\frac{\mathrm{Cpv}_0\cdot|\mathrm{kph}|_2^2}{\bar{m}},\qquad
\left|\operatorname{grad}\mathrm{Vpl}^{\eta}_{\mu\nu}(\mathrm{kph})\right|_2\le\frac{\mathrm{Cpv}_0\cdot|\mathrm{kph}|_2}{\bar{m}}.
\end{gather*}
\item[(d)] As $\eta\to0$, locally uniformly in $\mathrm{kph}$,
\[
\mathrm{Vpl}^{\eta}_{\mu\nu}(\mathrm{kph})\to-\left(|\mathrm{kph}|_2^2\cdot\delta_{\mu\nu}-\mathrm{kph}_\mu\cdot\mathrm{kph}_\nu\right)\cdot\mathrm{vpo}_{\bar{m}}\left(|\mathrm{kph}|_2^2\right);
\]
no counterterm is subtracted, and the mass is not tuned.
\item[(e)] For every integer $i\ge0$ there is $\mathrm{Cah}_i(\beta)<\infty$ such that
\[
\left|\mathrm{ahl}_{\eta,\nu}(\mathrm{kph})\right|\le\mathrm{Cah}_i(\beta)\cdot\left(1+|\mathrm{kph}|_2\right)^{-i}
\]
for all $\eta\in(0,1]$ and all $\mathrm{kph}$ with $|\mathrm{kph}|_\infty\le\frac{3\pi}{2\eta}$.
Moreover $\mathrm{ahl}_{\eta,\nu}(\mathrm{kph})\to\hat{a}_\nu(\mathrm{kph})$ as $\eta\to0$, and $\sum_\nu\mathrm{kph}_\nu\cdot\hat{a}_\nu(\mathrm{kph})=0$.
\item[(f)] There are constants $\mathrm{Cpv}_1=\mathrm{Cpv}_1(\rho,\beta)<\infty$ and $\mathrm{Cpv}_2=\mathrm{Cpv}_2(\rho,\bar{m},\beta)<\infty$, and a bounded function $\mathrm{epv}\colon(0,1]\to[0,\infty)$ which depends on $(\rho,\bar{m},\beta)$ but not on $R$ and satisfies $\mathrm{epv}(\eta)\to0$ as $\eta\to0$, such that
\[
\left|a_{\mathrm{ch}}(\eta,R)-\mathrm{pv}(f)\right|\le\mathrm{Cpv}_1\cdot\bar{m}^{-3}\cdot R^{-4}+\mathrm{Cpv}_2\cdot R^{-1}+\mathrm{epv}(\eta)
\]
for all $\eta\in(0,1]$ and all $R\ge1$ for which $R/\eta$ is an integer $\ge2$.
In particular $a_{\mathrm{ch}}(\eta,R)\to\mathrm{pv}(f)$ as $\eta\to0$ and $R\to\infty$ in any manner.
The support of $\beta$ need not lie in the interior of a cube of side $R$, since the cochain $a_\eta[\beta]$ of Section~G5.1.3 is formed with the periodised probe.
\item[(g)] The number
\[
\mathrm{Cpv}_\beta:=\sup\left\{\left|a_{\mathrm{ch}}(\eta,R)\right|:\ 0<\eta\le1,\ R\ge1,\ R/\eta\ \text{an integer}\ \ge2\right\}
\]
is finite, and $\mathrm{Cpv}_\beta\le|\mathrm{pv}(f)|+\mathrm{Cpv}_1\cdot\bar{m}^{-3}+\mathrm{Cpv}_2+\sup\mathrm{epv}$.
It depends on $(\rho,\bar{m},\beta)$.
\item[(h)] (Invariance.)
For every $q>0$, every complex one-cochain $A$ on $\mathbb{Z}^3/n\mathbb{Z}^3$, every real number $\boldsymbol{\mu}$ and every $g_{\mathrm{w}}\in W_{\mathrm{af}}$,
\[
\det\left(D_{(g_{\mathrm{w}})_{*}A}+\boldsymbol{\mu}\right)=\det\left(D_A+\boldsymbol{\mu}\right)\qquad\text{on }\mathfrak{S}^{\mathrm{ch}}(1;n).
\]
For $v\in\eta\mathbb{Z}^3$ one has $\eta^{\frac{1}{2}}a_\eta[\tau_v\beta]=(g_{\mathrm{w}})_{*}\left(\eta^{\frac{1}{2}}a_\eta[\beta]\right)$, where $g_{\mathrm{w}}$ is the translation by $v/\eta$.
Consequently $N_0$, the numbers $c_m$, $a_{\mathrm{ch}}(\eta,R)$ and $b_{\mathrm{ch}}(\eta,R)$ remain the same when $\beta$ is replaced by $\tau_v\beta$ with $v\in\eta\mathbb{Z}^3$, and when the cochain $\eta^{\frac{1}{2}}a_\eta[\beta]$ is replaced by its image under $(g_{\mathrm{w}})_{*}$ with $g_{\mathrm{w}}\in W_{\mathrm{af}}$.
\end{itemize}
\end{named}

\begin{named}{Remark}
Proposition~36 has no counterpart in \cite{QEDI}.
It yields (R4) of Proposition~43.1 with $\mathfrak{e}_2=0$, where \cite{QEDI} has the third bound of \cite[Corollary~20.32]{QEDI} and \cite[Theorem~21.3, Part~2(b)]{QEDI}, with $\mathfrak{e}_2=6\varepsilon$.
Clause (h) takes the place of the last sentence of \cite[Theorem~21.1(d5)]{QEDI} and of the sentence on the independence of the placement in \cite[Section~21.1.6]{QEDI}.
\end{named}

\subsubsection*{Proposition~37.2 (the four-photon coefficient)}

In the setting (G5.FF) let
\[
\mathrm{G0}:=\left(D^{\mathrm{ch}}[0]+\bar{m}\cdot\eta\right)^{-1}\quad\text{on }\mathfrak{S}^{\mathrm{ch}}(1;n),\qquad
\mathrm{Vd}_j:=\left.\frac{\mathrm{d}^j}{\mathrm{d}t^j}D^{\mathrm{ch}}[B_t]\right|_{t=0}\quad(j\ge1);
\]
explicitly, with $B_1$ the value of $B_t$ at $t=1$, so that $B_t=t\,B_1$ by the definition of $B_t$ in (G5.FF), one has for $\psi\in\mathfrak{S}^{\mathrm{ch}}(1;n)$
\[
\left(\mathrm{Vd}_j\psi\right)(c,g)=-M_3(c,g)\left(i\,B_1(c,c+ge_1)\right)^j\psi(c+ge_1,gs_3).
\]
The Schatten class of order four is denoted $I_4$.

\begin{named}{Proposition 37.2}
Assume the setting (G5.FF) with $R\ge D_{\mathrm{be}}+1$.
\begin{itemize}
\item[(a)] The operator $S_{\bar{m}}\,i\gamma\cdot a$ on $L^2(\mathbb{R}^3;\mathbb{C}^4)$ and the operator $S^{R}\,i\gamma\cdot a^{R}$ on $L^2_{\mathrm{ap}}(T_R;\mathbb{C}^4)$ belong to $I_4$.
Hence $b(f)$ and $b_{R}(f)$ are defined, and both are real.
\item[(b)] (Exact identity.)
\begin{multline*}
b_{\mathrm{ch}}(\eta,R)=\operatorname{Tr}\left(\mathrm{G0}\,\mathrm{Vd}_4\right)-4\operatorname{Tr}\left(\mathrm{G0}\,\mathrm{Vd}_1\,\mathrm{G0}\,\mathrm{Vd}_3\right)-3\operatorname{Tr}\left(\left(\mathrm{G0}\,\mathrm{Vd}_2\right)^2\right)\\
+12\operatorname{Tr}\left(\left(\mathrm{G0}\,\mathrm{Vd}_1\right)^2\mathrm{G0}\,\mathrm{Vd}_2\right)-6\operatorname{Tr}\left(\left(\mathrm{G0}\,\mathrm{Vd}_1\right)^4\right).
\end{multline*}
\item[(c)] (Contact traces.)
There is $C_{\mathrm{ct}}=C_{\mathrm{ct}}(\rho,\beta)$ such that
\begin{multline*}
\left|\operatorname{Tr}\left(\mathrm{G0}\,\mathrm{Vd}_4\right)\right|+4\left|\operatorname{Tr}\left(\mathrm{G0}\,\mathrm{Vd}_1\,\mathrm{G0}\,\mathrm{Vd}_3\right)\right|+3\left|\operatorname{Tr}\left(\left(\mathrm{G0}\,\mathrm{Vd}_2\right)^2\right)\right|+12\left|\operatorname{Tr}\left(\left(\mathrm{G0}\,\mathrm{Vd}_1\right)^2\mathrm{G0}\,\mathrm{Vd}_2\right)\right|\\
\le C_{\mathrm{ct}}\cdot\eta\cdot\left(1+\log\left(1+\frac{5}{\eta\cdot\bar{m}}\right)\right).
\end{multline*}
\item[(d)] There is $C_{1}^{\mathrm{N2}}=C_{1}^{\mathrm{N2}}(\rho,\bar{m},\beta)<\infty$ such that
\[
\left|b_{\mathrm{ch}}(\eta,R)-b_{R}(f)\right|\le C_{1}^{\mathrm{N2}}\cdot\eta^{\frac{1}{4}}.
\]
\item[(e)] There is $C_{2}^{\mathrm{N2}}=C_{2}^{\mathrm{N2}}(\bar{m},\beta)<\infty$ such that
\[
\left|b_{R}(f)-b(f)\right|\le C_{2}^{\mathrm{N2}}\cdot\exp\left(-\bar{m}\left(R-D_{\mathrm{be}}\right)\right).
\]
\item[(f)] Consequently
\[
\left|b_{\mathrm{ch}}(\eta,R)-b(f)\right|\le C_{1}^{\mathrm{N2}}\cdot\eta^{\frac{1}{4}}+C_{2}^{\mathrm{N2}}\cdot\exp\left(-\bar{m}\left(R-D_{\mathrm{be}}\right)\right),\qquad
\left|b_{\mathrm{ch}}(\eta,R)\right|\le|b(f)|+C_{1}^{\mathrm{N2}}+C_{2}^{\mathrm{N2}};
\]
$b_{\mathrm{ch}}(\eta,R)\to b(f)$ as $\eta\to0$ and $R\to\infty$ in any manner;
and if $R\cdot L^N$ is an integer for all sufficiently large $N$, then $b_{R}(f)=\lim_{N\to\infty}b_{\mathrm{ch}}\left(L^{-N},R\right)$.
\item[(g)] The constants $C_{\mathrm{ct}}$, $C_{1}^{\mathrm{N2}}$ and $C_{2}^{\mathrm{N2}}$ depend neither on $\eta$, nor on $R$, nor on the position of $\operatorname{supp}\beta$: they are the same for all $\tau_v\beta$, $v\in\mathbb{R}^3$.
The parameter $\rho$ enters $C_{\mathrm{ct}}$ and $C_{1}^{\mathrm{N2}}$ only through $C_E(\rho)$.
\end{itemize}
\end{named}

\begin{named}{Remark}
Proposition~37.2 corresponds to \cite[Theorem~21.1(c2)]{QEDI} with its proof in \cite[Section~21.4.3]{QEDI};
on the range $R\ge D_{\mathrm{be}}+1$ it restates \cite[Theorem~21.1(c1), (c3)]{QEDI} and the definition of $b_{R}(f)$ in \cite[Section~21.1.9]{QEDI}.
The mathematical differences are these.
\begin{itemize}
\item[(1)] The rate is $\eta^{\frac{1}{4}}$ instead of $\eta^{\frac{1}{2}}$ (Section~I.8.5), and $C_{1}^{\mathrm{N2}}$ depends also on $\rho$.
\item[(2)] $D_{\mathrm{be}}$ stands in the place of $D_\beta$.
The range $R\ge D_{\mathrm{be}}+1$ is larger than the range $R\ge D_\beta+1=D_{\mathrm{be}}+5$, and
\[
\exp\left(-\bar{m}\left(R-D_{\mathrm{be}}\right)\right)=\exp\left(-4\bar{m}\right)\cdot\exp\left(-\bar{m}\left(R-D_\beta\right)\right);
\]
so (e) contains \cite[Theorem~21.1(c3)]{QEDI}, with the constant $C_{2}^{\mathrm{N2}}$.
The argument of \cite[Section~21.4.4]{QEDI} is a second proof of (e) on its range, once two of its ingredients are replaced:
clause (a) stands in the place of \cite[Theorem~21.1(c1)]{QEDI},
and the definition of $S^{R}$ in Section~G5.6.0, as the multiplier $\left(\bar{m}+i\gamma\cdot p\right)^{-1}$ on the momenta $p\in\frac{2\pi}{R}\left(\mathbb{Z}+\frac{1}{2}\right)^3$, stands in the place of the orthonormal basis $e_p\otimes\epsilon$, $p\in P_R$, of $L^2_{\mathrm{ap}}(T_R;\mathbb{C}^4)$ that \cite[Section~21.4.4]{QEDI} takes from \cite[Section~21.4.3]{QEDI}, and on which it identifies $S^{R}$ with the convolution with the integral kernel of $S_{\bar{m}}$.
\item[(3)] $b_{R}(f)$ is defined exactly on the range on which (a) is proved.
In \cite[Section~21.1.9]{QEDI} the definition admits every $R>0$ for which the support of $a$ lies in the interior of a period cube, and on this whole range the operator $\left(S^{R}\,i\gamma\cdot a^{R}\right)^4$ is of trace class by \cite[Lemma~21.5]{QEDI}, which bounds the Hilbert--Schmidt norm of $S^{R}\,i\gamma\cdot a^{R}\,S^{R}$ for every $R>0$ and shows that $b_{R}(f)$ is real wherever it is defined;
in \cite[Theorem~21.1(c1)]{QEDI} the Hilbert--Schmidt property is stated again for $R\ge D_\beta+1$.
The two definitions differ only in their range: here it is $R\ge D_{\mathrm{be}}+1$, in \cite[Section~21.1.9]{QEDI} it is the wider range just described.
For $R\ge D_{\mathrm{be}}+1$ they agree:
the support of $a$ lies in that of $\beta$, hence in a closed coordinate cube of side $D_{\mathrm{be}}<R$, and so in the interior of a period cube;
and both definitions give $-6\operatorname{Tr}\left[\left(S^{R}\,i\gamma\cdot a^{R}\right)^4\right]$, with the same operator $S^{R}$ and the $R$-periodisation $a^{R}$ of $a$.
\item[(4)] The traces in (b) are traces on a space of dimension $192\,n^3$, and the vertices contain $M_3$ in the place of the projections $\frac{1}{2}\left(1\mp\gamma_\nu\right)$.
The identity (b) is the expression of $b_{\eta,R}(0)$ through the operators $T_1,\dots,T_4$ in \cite[Section~21.4.3]{QEDI}, where the vertex $V_r$ contains the factor $1/r!$.
\item[(5)] The bound (c) is sharper than the bound $C\cdot\eta^{\frac{1}{2}}$ for the four traces in the brackets of that expression in \cite[Section~21.4.3]{QEDI}.
\end{itemize}
\end{named}

\subsubsection*{Proposition~37.3 (the free coefficients on the reference torus: (R4) with $\mathfrak{e}_2=0$, (R5))}

\begin{named}{Proposition 37.3}
Let the constants be those of Section~I.8, let $\bar{m}\in(0,1]$, let $\beta$ be a probe, and let $e>0$ satisfy $e\cdot h^{\frac{1}{2}}\le e_1$ and $R^{*}(e)\ge D_{\mathrm{be}}+1$; the last inequality holds for $e\le e_{\sharp}(\beta)$ (Proposition~35.10).
Let $N\ge N_0(e)$ and $\eta=L^{-N}$.
Let $v\in\eta\mathbb{Z}^3$ be arbitrary, let $\hat{\mathfrak{a}}=\eta^{\frac{1}{2}}a_\eta[\tau_v\beta]$ on $C(n_*)$, let $\mathfrak{c}\in\mathrm{Cel}$, let $N_0$ be formed with $\hat{\mathfrak{a}}$ or with $\hat{\mathfrak{a}}_{\mathfrak{c}}$, and let $c_m=N_0^{(m)}(0)$.
Then the following hold.
\begin{itemize}
\item[(a)] $N_0$ depends neither on $v$ nor on $\mathfrak{c}$, and $c_1=c_3=0$.
Moreover $c_2=a_{\mathrm{ch}}\left(\eta,R^{*}(e)\right)[\beta]$ and $c_4-3c_2^2=b_{\mathrm{ch}}\left(\eta,R^{*}(e)\right)[\beta]$, the numbers formed with $\beta$ itself.
\item[(b)] $|c_2|\le\mathrm{Cpv}_\beta$ and $|c_4|\le|b(f)|+C_{1}^{\mathrm{N2}}+C_{2}^{\mathrm{N2}}+3\,\mathrm{Cpv}_\beta^2$.
Hence, for all $e$ and $N$ as above,
\[
|c_2|+|c_4|^{\frac{1}{2}}\le\mathrm{Cpv}_\beta+\left(|b(f)|+C_{1}^{\mathrm{N2}}+C_{2}^{\mathrm{N2}}+3\,\mathrm{Cpv}_\beta^2\right)^{\frac{1}{2}};
\]
this is (R4) of Proposition~43.1 with $\mathfrak{e}_2=0$.
\item[(c)] One has
\[
\left|c_4-3c_2^2-b(f)\right|\le C_{1}^{\mathrm{N2}}\cdot\eta^{\frac{1}{4}}+C_{2}^{\mathrm{N2}}\cdot\exp\left(-\bar{m}\left(R^{*}(e)-D_{\mathrm{be}}\right)\right)\le C_{1}^{\mathrm{N2}}\cdot\eta^{\frac{1}{4}}+C_{2}^{\mathrm{N2}}\cdot\exp\left(\bar{m}\cdot D_{\mathrm{be}}\right)\cdot h^{\frac{1}{2}}\cdot e;
\]
this is (R5) of Proposition~43.1 with the exponent $\frac{1}{4}$.
\item[(d)] $\left|c_2-\mathrm{pv}(f)\right|\le\mathrm{Cpv}_1\cdot\bar{m}^{-3}\cdot R^{*}(e)^{-4}+\mathrm{Cpv}_2\cdot R^{*}(e)^{-1}+\mathrm{epv}(\eta)$.
\end{itemize}
\end{named}

\begin{named}{Remark}
Proposition~37.3 corresponds to \cite[Theorem~21.3, Part~2(b), (c)]{QEDI}, together with the corresponding parts of its proof in \cite[Section~21.9]{QEDI}.
Here $\mathfrak{e}_2=0$ in the place of $6\varepsilon$, and $\eta^{\frac{1}{4}}$ in the place of $\eta^{\frac{1}{2}}$.
The bound (b) is derived from the one-loop statement Proposition~36, not from the renormalisation group form of the free determinant;
and the equality of the numbers formed with $\beta$ and with $\tau_v\beta$ is clause (a).
Clause (d) has no counterpart in \cite{QEDI}.
\end{named}

\subsection*{39. Chessboard estimate and lengthening for coordinate cells of half-integer width}

\subsubsection*{Proposition~39}

\begin{named}{Definition 39.0}[Coordinate cell structures]
The following extends \cite[Definitions~2.18 and~2.19]{QEDII} for the frame $\mathfrak{F}_{\mathrm{c}}$.
A coordinate cell structure is $\mathcal{Q}=(\mathfrak{F}_{\mathrm{c}};o_1,o_2,o_3;\mathrm{hw})$ with $o_\mu\in\frac{1}{2}\mathbb{Z}$ and $\mathrm{hw}\in\frac{1}{2}\mathbb{Z}$, $\mathrm{hw}>0$.
Its walls are the planes $(y_\mu,o_\mu+j\cdot\mathrm{hw})$, that is, $\left\{y_\mu=o_\mu+j\cdot\mathrm{hw}\right\}$, for $\mu=1,2,3$ and $j\in\mathbb{Z}$;
its open cells are
\[
Q^{\circ}_{\mathfrak{c}}(\mathcal{Q}):=\left\{y\in\mathbb{R}^3:\ o_\mu+\mathfrak{c}_\mu\cdot\mathrm{hw}<y_\mu<o_\mu+(\mathfrak{c}_\mu+1)\cdot\mathrm{hw},\ \mu=1,2,3\right\}\qquad\left(\mathfrak{c}\in\mathbb{Z}^3\right).
\]
$W(\mathcal{Q})$ denotes the group generated by the reflections $\theta_{y_\mu,o_\mu+j\cdot\mathrm{hw}}$.
A torus $C(n)$, $n_\mu\ge4$, carries $\mathcal{Q}$ if $n_\mu\in2\,\mathrm{hw}\,\mathbb{Z}$ for $\mu=1,2,3$;
its cells are the orbits of the cells under $\sum_{\mu} n_{\mu}\mathbb{Z}\,e_{\mu}$.
The phrases ``for every cell $\mathfrak{c}$ of $C(n)$'' and ``over the cells $\mathfrak{c}'$ of $C(n)$'' refer to a set of representatives $\mathfrak{c}\in\mathbb{Z}^3$ of these orbits.
As an element of the space $\mathcal{A}(\mathbb{R}^3)$ of \cite[Definition~2.12]{QEDII}, which contains every $\mathcal{A}\left(Q^{\circ}_{\mathfrak{c}}(\mathcal{Q})\right)$ and the products of elements of these spaces over distinct cells, a product over the cells depends on this set;
the function of the configuration on $C(n)$ that it defines does not, and neither does its value under $Z_n$ \cite[Lemma~6.4(c)]{QEDII}.
For $O\in\mathcal{A}\left(Q^{\circ}_{\mathfrak{c}}(\mathcal{Q})\right)$, the algebra of \cite[Definition~2.12]{QEDII} for the open set $Q^{\circ}_{\mathfrak{c}}$, and for a cell $\mathfrak{c}'$ put
\[
O^{[\mathfrak{c}']}:=O^{g_{\mathrm{w}}},
\]
where $g_{\mathrm{w}}$ is the element of $W(\mathcal{Q})$ with $g_{\mathrm{w}}Q^{\circ}_{\mathfrak{c}}=Q^{\circ}_{\mathfrak{c}'}$ (Section~I.4.8);
$O^{\sharp}_n$ denotes the product of the $O^{[\mathfrak{c}']}$ over the cells $\mathfrak{c}'$ of $C(n)$.
The structure $\mathcal{Q}_{*}$ of Section~I.6.4 is the case $o=0$, $\mathrm{hw}=\mathrm{hc}$; for $\mathfrak{c}=0$ the element $g_{\mathrm{w}}$ is then $\mathsf{g}_{\mathfrak{c}'}$.
Let $\mathcal{A}(C(n))$ be the span of the products of $\mathcal{A}(\mathbb{R}^3)$ that are read on $C(n)$ \cite[Definition~5.11]{QEDII};
it is the set of the elements of $\mathcal{A}(\mathbb{R}^3)$ that are read on $C(n)$, and it contains every element of $\mathcal{A}(\mathbb{R}^3)$ that embeds in $C(n)$.
On $\mathcal{A}(C(n))$ the functional $Z_n[\,\cdot\,]$ is the functional $Z_T$ of \cite[Definition~2.13(a)]{QEDII} with $T=C(n)$, the mass parameter of \cite[Definition~2.10]{QEDII} being $M_{\mathrm{ch}}$ (Section~I.7.2).
\end{named}

\begin{named}{Proposition 39}
Let $\boldsymbol{\mu}$ be real, let $q>0$, and let $\mathcal{Q}$ be as in Definition~39.0.
\begin{itemize}
\item[(a)] (Soundness of the definition.)
Every wall is a mirror in the sense of Section~I.4.8, because $o_\mu+j\cdot\mathrm{hw}$ lies in $\operatorname{Adm}(y_\mu)=\frac{1}{2}\mathbb{Z}$.
For a wall $(y_\mu,\mathrm{po})$ the product $\theta_{y_\mu,\mathrm{po}+\mathrm{hw}}\,\theta_{y_\mu,\mathrm{po}}$ is the translation by $2\,\mathrm{hw}\,e_\mu$, a vector of $\mathbb{Z}^3$.
The group $W(\mathcal{Q})$ is the product of three commuting infinite dihedral groups, and it acts simply transitively on the cells.
If $C(n)$ carries $\mathcal{Q}$, then $\sum_{\mu} n_{\mu}\mathbb{Z}\,e_{\mu}$ is contained in the subgroup $(2\,\mathrm{hw}\,\mathbb{Z})^3$ of the translations in $W(\mathcal{Q})$, the quotient acts simply transitively on the cells of $C(n)$, and the number of these cells is
\[
\mathrm{nc}(n):=\prod_\mu\frac{n_\mu}{\mathrm{hw}},
\]
a product of three even integers.
\item[(b)] (Reading.)
Let $C(n)$ carry $\mathcal{Q}$, and let $O_{\mathfrak{c}}\in\mathcal{A}\left(Q^{\circ}_{\mathfrak{c}}(\mathcal{Q})\right)$ for every cell $\mathfrak{c}$ of $C(n)$.
Each $O_{\mathfrak{c}}$ is a finite sum of products \cite[Definition~2.12]{QEDII}, and for every single product that arises when $\prod_{\mathfrak{c}}O_{\mathfrak{c}}$ is expanded, the boxes of the factors and all their translates by $\sum_{\mu} n_{\mu}\mathbb{Z}\,e_{\mu}$ are pairwise disjoint.
The product $\prod_{\mathfrak{c}}O_{\mathfrak{c}}$ is read on $C(n)$ in the sense of \cite[Definition~2.13(a)]{QEDII}, and $Z_n\left[\prod_{\mathfrak{c}}O_{\mathfrak{c}}\right]$ converges absolutely.
An element of $\mathcal{A}\left(Q^{\circ}_{0}(\mathcal{Q})\right)$ is an integrand on $\mathbb{R}^3$: one and the same $O$ is read on all tori that carry $\mathcal{Q}$.
\item[(c)] (Positivity; chessboard estimate.)
In the situation of (b),
\[
Z_n\left[\left(O_{\mathfrak{c}}\right)^{\sharp}_n\right]\ge0\quad\text{for every }\mathfrak{c},\qquad
\left|Z_n\left[\prod_{\mathfrak{c}}O_{\mathfrak{c}}\right]\right|\le\prod_{\mathfrak{c}}Z_n\left[\left(O_{\mathfrak{c}}\right)^{\sharp}_n\right]^{\frac{1}{\mathrm{nc}(n)}}.
\]
\item[(d)] (One norm for all copies.)
For $O\in\mathcal{A}\left(Q^{\circ}_{0}(\mathcal{Q})\right)$ and every cell $\mathfrak{c}$ one has $\left(O^{[\mathfrak{c}]}\right)^{[\mathfrak{c}']}=O^{[\mathfrak{c}']}$ for all $\mathfrak{c}'$;
hence $\left(O^{[\mathfrak{c}]}\right)^{\sharp}_n=O^{\sharp}_n$.
\item[(e)] (Lengthening.)
For $\mathrm{sv}=(\mathrm{sv}_1,\mathrm{sv}_2,\mathrm{sv}_3)$ with even positive integers $\mathrm{sv}_\mu$ such that $\mathrm{hw}\cdot\mathrm{sv}_\mu\ge4$, and for $O\in\mathcal{A}\left(Q^{\circ}_{0}(\mathcal{Q})\right)$, put
\[
\mathcal{Z}_{\mathrm{sv}}(O):=Z_{\mathrm{hw}\cdot\mathrm{sv}}\left[O^{\sharp}_{\mathrm{hw}\cdot\mathrm{sv}}\right];
\]
here $\mathrm{hw}\cdot\mathrm{sv}_\mu$ is an integer.
Fix $\mu$ and the two other entries of $\mathrm{sv}$.
On the set $\left\{\mathrm{sv}_\mu\ \text{even}:\ \mathrm{hw}\cdot\mathrm{sv}_\mu\ge4\right\}$ the function $\mathrm{sv}_\mu\mapsto\mathcal{Z}_{\mathrm{sv}}(O)^{\frac{1}{\mathrm{sv}_\mu}}$ is non-increasing;
if $\mathcal{Z}_{\mathrm{sv}}(O)>0$ for one $\mathrm{sv}_\mu$ of this set, then $\mathcal{Z}_{\mathrm{sv}}(O)>0$ for all of them, and $\mathrm{sv}_\mu\mapsto\log\mathcal{Z}_{\mathrm{sv}}(O)$ is convex on this set.
\item[(f)] (Normalised form on cubic tori.)
For $n\in2\,\mathrm{hw}\,\mathbb{Z}$ with $n\ge4$, and for $O_{\mathfrak{c}}$ as in (b) on the cubic torus $C(n,n,n)$,
\[
E^{\mathrm{af}}_n\left[\left(O_{\mathfrak{c}}\right)^{\sharp}_n\right]\ge0,\qquad
\left|E^{\mathrm{af}}_n\prod_{\mathfrak{c}}O_{\mathfrak{c}}\right|\le\prod_{\mathfrak{c}}E^{\mathrm{af}}_n\left[\left(O_{\mathfrak{c}}\right)^{\sharp}_n\right]^{\left(\frac{\mathrm{hw}}{n}\right)^3}.
\]
For a function $O$ of the curvature, $O^{[\mathfrak{c}']}(u)=O\left(g_{\mathrm{w}}^{*}u\right)$ if $g_{\mathrm{w}}$ is even, and $O^{[\mathfrak{c}']}(u)=\overline{O\left(g_{\mathrm{w}}^{*}u\right)}$ if $g_{\mathrm{w}}$ is odd.
\end{itemize}
\end{named}

\begin{named}{Remark}
In \cite{QEDII} the width of a cell structure is a positive integer.
The points (I.h1)--(I.h5) of Section~I.6.4 extend statements of \cite{QEDII} to widths in $\frac{1}{2}\mathbb{Z}$ for the frame $\mathfrak{F}_{\mathrm{c}}$.
With these extensions, the clauses (a)--(f) of Proposition~39 are, in this order:
\cite[Definition~2.18]{QEDII} together with the paragraph on the groups of the walls in \cite[Definition~6.13]{QEDII};
\cite[Lemma~6.5(d)]{QEDII} together with \cite[Definition~6.13]{QEDII};
\cite[Theorem~6.16(a)]{QEDII};
\cite[Definition~6.13]{QEDII} together with the action of the group $W$ on even integrands (the paragraph that follows \cite[Definition~2.12]{QEDII}) and \cite[Proposition~3.16(f$''$)]{QEDII};
\cite[Proposition~6.23(a)]{QEDII};
and \cite[Proposition~12.10(f)]{QEDII} with its proof, which divides \cite[Theorem~6.16(a)]{QEDII} by $Z_T[1]$, on every cubic torus that carries $\mathcal{Q}$ and not only for the sides $n=8h\cdot L^{T+N}$ of \cite[Definition~12.1]{QEDII}.
The proofs of \cite[Definition~6.7, Lemma~6.5(a), (b), (d), Proposition~6.8, Definition~6.13, Lemma~6.14]{QEDII}, of \cite[Proposition~6.15]{QEDII} with $T_0=T_1=T_2=T_3=C(n)$, and of \cite[Theorem~6.16(a), Lemma~6.17, Proposition~6.23]{QEDII} hold without change for the frame $\mathfrak{F}_{\mathrm{c}}$, the families $\ell=y_\mu$, the tori $T=C(n)$ and widths in $\frac{1}{2}\mathbb{Z}$.
In particular, for $e\cdot h^{\frac{1}{2}}\le e_1$, $N\ge N_0(e)$ and $T\ge T^{*}(e)$, the cells of $\mathcal{Q}_{*}$ exist on the tori with the sides $n=8h\cdot L^{T+N}$:
there $n=\mathrm{kt}\cdot n_*$ with $\mathrm{kt}=8L^{T-T^{*}(e)}$, so that $n/\mathrm{per}=\mathrm{kt}\cdot L$ is an integer and $\mathrm{per}$ divides $n$ (Section~I.6.5).
For the diagonal frames nothing changes.
The two diagonal families of such a frame have $\operatorname{Adm}=\mathbb{Z}$, while its third family is an axis family, with $\operatorname{Adm}=\frac{1}{2}\mathbb{Z}$ and positions in $\frac{1}{2}\mathbb{Z}$;
as one width serves all three families, the walls of the two diagonal families are mirrors only if the width is an integer, and so the width remains an integer.

In the proofs just listed the width enters at the following places only.
\begin{itemize}
\item[(1)] The walls are mirrors; this needs $o_\mu,\mathrm{hw}\in\frac{1}{2}\mathbb{Z}$.
\item[(2)] The vector $\tau:=2\,\mathrm{hw}\,e_\mu$ lies in $\mathbb{Z}^3$; this needs $2\,\mathrm{hw}\in\mathbb{Z}$.
The operator $V_\tau$ of \cite[Definition~5.27]{QEDII} maps the mirror space $\mathrm{Hs}(\mathrm{po})$ of that definition to $\mathrm{Hs}(\mathrm{po}+2\,\mathrm{hw})$, the space of a plane of the same kind.
\item[(3)] \cite[Proposition~6.8]{QEDII} is formulated for every height $h'$ with $\mathrm{po}_0+j\cdot h'\in\operatorname{Adm}(\ell)$ for all $j$.
The slab operators $\mathrm{Cs}[G]\colon\mathrm{Hs}(\mathrm{po}_{j-1})\to\mathrm{Hs}(\mathrm{po}_j)$ of \cite[Definition~5.28]{QEDII} go from the space of a site plane to the space of a link plane or back, and \cite[Lemma~6.3]{QEDII} is H\"older's inequality for operators between different Hilbert spaces.
\item[(4)] In \cite[Lemma~6.5(d)]{QEDII} the conditions are $2r<\mathrm{hw}$ for a box of half side $r$ in a cell and $|\ell(\gamma)|\ge2\,\mathrm{hw}$: inequalities between real numbers.
\item[(5)] The shifts $\mathrm{sh}(v)$, that is, the cell indices of translations in \cite[Definition~6.13]{QEDII}, are defined for $v\in\mathbb{Z}^3$ with $v_\mu\in\mathrm{hw}\,\mathbb{Z}$;
for $\mathrm{hw}=\mathrm{per}/2$ with odd $\mathrm{per}$ this means $v\in\mathrm{per}\,\mathbb{Z}^3$, and $\mathrm{sh}(v)$ is then even.
\item[(6)] The volume $|Q|=\mathrm{hw}^3$ of a cell enters only through the ratio $|n|/\mathrm{nc}(n)$.
\item[(7)] $\mathrm{hw}\cdot\mathrm{sv}_\mu$ is an integer, because $\mathrm{sv}_\mu$ is even.
\item[(8)] $n_\mu/(2\,\mathrm{hw})$ is an integer, by the definition of ``carries'':
the number $n_\mu/\mathrm{hw}$ of the layers is even, and $n_\mu\cdot e_\mu=\frac{n_\mu}{2\,\mathrm{hw}}\cdot\tau$ lies in the dihedral group of the walls.
\end{itemize}
\end{named}

\begin{named}{Remark 39.R}
Let $n_2,n_3\ge4$, put $v_s:=n_2\cdot n_3$, and let $\mathrm{Fk}'$, $\mathrm{Fk}$, $\mathrm{Qs}(A)\colon\mathrm{Fk}'\to\mathrm{Fk}$, $\mathrm{Wq}(a)\in\operatorname{End}\mathrm{Fk}$ and $\mathrm{Qc}$ be the objects of \cite[Definition~2.27]{QEDII} for the cross-section $(n_2,n_3)$, with $e_1$ as the time direction.
In the notation of that definition, $\mathrm{Fk}=\mathcal{F}$, $\mathrm{Fk}'=\mathcal{F}'$, $\mathrm{Qs}(A)=\widehat{\mathfrak{N}}(A)$ and $\mathrm{Wq}(a)=u(a)$;
for a subset $I$ of the set $\mathcal{J}^{\mathrm{l}}(\Lambda_{\mathrm{s}})$ of the $64v_s$ modes of the link plane, $\mathrm{Qc}_I$ is the charge $\mathfrak{q}(I)\in\mathbb{Z}^{\Lambda_{\mathrm{s}}}$ of the paragraph ``Coordinate planes'' of that definition,
so that $\mathrm{Qc}_I(y)$ is the number of the modes of $I$ over the site $y$ that belong to the groups 1 or 2, minus the number of the modes of $I$ over $y$ that belong to the groups 3 or 4.
One has $\dim\mathrm{Fk}'=\dim\mathrm{Fk}=2^{64v_s}$.
For an arbitrary unitary map $\mathrm{Uf}\colon\mathrm{Fk}\to\mathrm{Fk}'$ put
\[
\mathrm{Qf}(A):=\mathrm{Qs}(A)\,\mathrm{Uf}\in\operatorname{End}\mathrm{Fk},\qquad s:=1,\qquad\mathrm{Gmf}:=\mathrm{Wq}.
\]
Let $n_1\ge1$, put $n:=(n_1,n_2,n_3)$, and let $A$ be a real one-cochain on $\mathbb{Z}^3$ that is periodic under $\sum_\mu n_\mu\mathbb{Z}\,e_\mu$.
For $t\in\mathbb{Z}$ let $A_t$ be the family of the values of $A$ on the links of $\mathbb{Z}^3$ that lie in the plane $\{y_1=t\}$, read as a one-cochain of the slice $\Lambda_{\mathrm{s}}=\mathbb{Z}^2/(n_2\mathbb{Z}\times n_3\mathbb{Z})$,
and let $a_t$ be the function on the sites of the slice whose value at $y$ is the value of $A$ on the link from $(t,y)$ to $(t+1,y)$;
these are the slice variables $A_j$, $a_j$ of \cite[Definition~2.26]{QEDII}, and $A_{t+n_1}=A_t$, $a_{t+n_1}=a_t$.
Then, and in the same way with $e_2$ or $e_3$ as the time direction,
\[
\det\nolimits_n\left(D_A+\boldsymbol{\mu}\right)=\operatorname{tr}_{\mathrm{Fk}}\prod_{t=n_1-1,\dots,0}\left[s\,\mathrm{Qf}(A_{t+1})^{*}\,\mathrm{Gmf}(a_t)\,\mathrm{Qf}(A_t)\,s\right];
\]
here, for $n_1<4$ as for $n_1\ge4$, $\det\nolimits_n$ is the determinant on $\mathfrak{S}^{\mathrm{ch}}(1;n)$ (Proposition~46\,(e), Proposition~38.1).
The display rests on the identity
\[
\det\nolimits_n\left(D_A+\boldsymbol{\mu}\right)=\operatorname{tr}_{\mathrm{Fk}'}\prod_{t=n_1-1,\dots,0}\left(\mathrm{Qs}(A_{t+1})^{*}\,\mathrm{Wq}(a_t)\,\mathrm{Qs}(A_t)\right),\qquad A_{n_1}=A_0,
\]
in which, as in the display above, the later times stand to the left;
for $n_1\ge4$ it is the display of \cite[Proposition~5.33]{QEDII} with $\Upsilon_j=u(a_j)$, together with the last sentence of that proposition, and for every $n_1\ge1$ it is Lemma~8.1 of Chapter~39.
The display follows from this identity, because $\mathrm{Uf}\,\mathrm{Uf}^{*}=1$ cancels between neighbouring factors and $\operatorname{tr}_{\mathrm{Fk}}\left(\mathrm{Uf}^{*}Y\,\mathrm{Uf}\right)=\operatorname{tr}_{\mathrm{Fk}'}Y$.
The map $\mathrm{Qf}$ is bounded, continuous and a function of the link phases.
With the commuting self-adjoint charges $Q^{\mathrm{ch}}_y:=\operatorname{diag}\left(\mathrm{Qc}_I(y)\right)_{I\subset\mathcal{J}^{\mathrm{l}}(\Lambda_{\mathrm{s}})}$, diagonal in the orthonormal basis $(e_I)$ of $\mathrm{Fk}$, $y$ a site of the slice, whose spectra lie in $\{-32,\dots,32\}$,
\[
\mathrm{Gmf}(a)=\exp\left(iq\sum_ya_y\cdot Q^{\mathrm{ch}}_y\right)
\]
is a unitary representation, and $\mathrm{Gmf}(a)=1$ if all $a_y$ lie in $r_q\mathbb{Z}$.
All these data are formed with $(n_2,n_3)$, with $(q,\rho,M_{\mathrm{ch}})$ and with $\mathrm{Uf}$, and not with $n_1$.
They are therefore slice data for the cross-section $(n_2,n_3)$, in the sense of \cite[Section~11.2.1]{QEDI}, of the function $A\mapsto\det\nolimits_n\left(D_A+\boldsymbol{\mu}\right)$ of the chamber operator:
they satisfy the conditions (SD1), (SD2) stated there, where in (SD2) this function stands in the place of the Wilson determinant $\Delta(qA)$.
Equivalently, they have the properties (a), (b) of \cite[Lemma~21.4]{QEDI} and the properties (FK), (F1)--(F3) of \cite[Section~10.2.1]{QEDI}, where in (b) and in (FK) the same function stands in the place of the Wilson determinant.

\emph{Closed cells.}
Let $\mathrm{hw}\ge\frac{3}{2}$, and let $C(n)$ carry $\mathcal{Q}$.
For $\mathfrak{c}\in\mathbb{Z}^3$ let $\mathrm{Pqp}(\mathfrak{c})$ be the set of the plaquettes of $\mathbb{Z}^3$ that are contained in the closure of $Q^{\circ}_{\mathfrak{c}}(\mathcal{Q})$ and in no wall (it is the set $P^{+}(\mathfrak{c}_c)$ of \cite[Definition~21.18]{QEDI}), and let $\mathrm{Fcl}(\mathfrak{c})$ be the set of the Borel functions $\Phi$ on $\mathbb{R}^{\mathrm{Pqp}(\mathfrak{c})}$ with
\[
|\Phi(u)|\le\mathrm{Cq}\cdot\exp\left(\left(\frac{1}{2}-\epsilon\right)\sum_{p\in\mathrm{Pqp}(\mathfrak{c})}u_p^2\right)
\]
for some $\mathrm{Cq}<\infty$ and some $\epsilon\in(0,\frac{1}{2}]$.
With $g_{\mathrm{w}}$ as in Definition~39.0 put $\Phi^{[\mathfrak{c}']}(u):=\Phi\left(g_{\mathrm{w}}^{*}u\right)$ if $g_{\mathrm{w}}$ is even and $\Phi^{[\mathfrak{c}']}(u):=\overline{\Phi\left(g_{\mathrm{w}}^{*}u\right)}$ if $g_{\mathrm{w}}$ is odd, and let $\Phi^{\sharp}_n$ be the product of the $\Phi^{[\mathfrak{c}']}$ over the cells $\mathfrak{c}'$ of $C(n)$.
Then $Z_n\left[\prod_{\mathfrak{c}}\Phi_{\mathfrak{c}}\right]$ converges absolutely, and the clauses (c), (d), (e), (f) of Proposition~39 hold with $\Phi_{\mathfrak{c}}\in\mathrm{Fcl}(\mathfrak{c})$ in the place of $O_{\mathfrak{c}}\in\mathcal{A}\left(Q^{\circ}_{\mathfrak{c}}(\mathcal{Q})\right)$.
This contains the closed-cell form of \cite[Propositions~21.22 and~21.24(a)]{QEDI} for the all-flux functional, and it gives a second proof of (c) and (e) for functions of the curvature.
Of the fermion, this proof requires, besides the slice data for the time direction $e_1$, only the covariance of the determinant under the two cyclic permutations of the axes (Proposition~50\,(d), in the form of \cite[Lemma~5.46]{QEDII}), for (c) and for (e) in the directions $2$ and $3$, and the inequality $0<Z_n[1]$ (Proposition~50\,(e)), for (f).
\end{named}

\begin{named}{Remark}
Definition~39.0, Proposition~39 and Remark~39.R correspond to \cite[Definition~21.18]{QEDI}, \cite[Lemma~21.20]{QEDI}, \cite[Proposition~21.22]{QEDI} together with the reduction to bounded functions at the beginning of \cite[Section~21.7.4]{QEDI}, and \cite[Proposition~21.24(a)]{QEDI}.
They differ from these statements as follows.
\begin{itemize}
\item[(1)] The contents of the cells are read in the open cell, that is, in $\mathcal{A}\left(Q^{\circ}_{\mathfrak{c}}\right)$; the sets $P^{o}(\mathfrak{c}_c)$, $\mathfrak{A}(\mathfrak{c}_c)$ and $\mathfrak{A}^{+}$ of \cite{QEDI} do not occur.
\item[(2)] The contents may be unbounded (the class $\mathrm{Ad}$: growth at most $\exp\left(\left(\frac{1}{2}-\epsilon\right)\sum F_p^2\right)$ in pairwise disjoint boxes), and they may be Grassmann valued and even;
the reduction to bounded functions at the beginning of \cite[Section~21.7.4]{QEDI} does not enter.
\item[(3)] Only the all-flux functional occurs: there is no Bloch angle and no flux sector, and nothing is stated about $Z^{(0)}_n$.
\item[(4)] Clause (e) compares any two even numbers of cells in one direction, not only $N_\mu$ and $k\,N_\mu$ cells as in \cite[Proposition~21.24(a)]{QEDI}.
\item[(5)] Every real $\boldsymbol{\mu}$ is admitted, whereas in \cite{QEDI} the lattice mass satisfies $\boldsymbol{\mu}>-\frac{1}{2}$;
the condition $m_{\mathrm{p}}\ge2$ of \cite{QEDI} becomes $\mathrm{hw}\cdot\mathrm{sv}_\mu\ge4$.
\item[(6)] Offsets $o_\mu\in\frac{1}{2}\mathbb{Z}$ are admitted.
This contains \cite[Remark~21.27]{QEDI}, which has integer offsets; the half-integer offsets are those of \cite[Definition~2.18]{QEDII}.
\item[(7)] In the proofs of \cite[Propositions~21.22 and~21.24]{QEDI} the assertions are proved for the direction singled out as time and carried to the two other directions by the isometry of the torus that interchanges two axes ((Q2) of \cite[Section~21.7.1]{QEDI});
here each of the three directions is treated by the slab operators of the family of planes orthogonal to it, as in \cite[Proposition~6.15]{QEDII}.
\end{itemize}
\end{named}

\subsection*{40. Transfer to all tori}

\subsubsection*{Proposition~40}

\begin{named}{Proposition 40}
Let $\boldsymbol{\mu}$ be a real number and let $q>0$.
Let $\mathcal{Q}=(\mathfrak{F}_{\mathrm{c}};o;\mathrm{hw})$ be a coordinate cell structure as in Definition~39.0,
let $\mathrm{ncs}\ge2$ be an even integer with $n_r:=\mathrm{hw}\cdot\mathrm{ncs}\ge4$,
put $N_r:=\mathrm{ncs}^3$,
and let $\check{j}\ge2$ be an integer.
For every even integer $\mathrm{sv}$ with $\mathrm{hw}\cdot\mathrm{sv}\ge4$, the torus $\mathrm{hw}\cdot\mathrm{sv}$ is the cubic torus with $\mathrm{hw}\cdot\mathrm{sv}$ sites per side;
it has $\mathrm{sv}^3$ cells.
On it one writes
\[
Z_{[\mathrm{sv}]}:=Z_{\mathrm{hw}\cdot\mathrm{sv}}[1],\qquad
E_{[\mathrm{sv}]}:=E^{\mathrm{af}}_{\mathrm{hw}\cdot\mathrm{sv}},\qquad
\mathcal{Z}_{[\mathrm{sv}]}(O):=Z_{\mathrm{hw}\cdot\mathrm{sv}}\left[O^{\sharp}_{\mathrm{hw}\cdot\mathrm{sv}}\right]\quad\left(O\in\mathcal{A}\left(Q^{\circ}_{0}(\mathcal{Q})\right)\right),
\]
together with $\mathrm{kr}:=\mathrm{sv}/\mathrm{ncs}$.
Define further
\[
\|O\|_{*}:=\left(\frac{\mathcal{Z}_{[\mathrm{ncs}]}(O)}{Z_{[\mathrm{ncs}]}}\right)^{\frac{1}{N_r}}=\left(E_{[\mathrm{ncs}]}O^{\sharp}\right)^{\frac{1}{N_r}},\qquad
D_{\check{j}}:=\check{j}^3\cdot\log Z_{[\mathrm{ncs}]}-\log Z_{[\check{j}\cdot\mathrm{ncs}]},\qquad
\delta_{\check{j}}:=\frac{D_{\check{j}}}{(\check{j}-1)^3\cdot N_r}.
\]
\begin{itemize}
\item[(a)] One has $D_{\check{j}}\ge0$.
\item[(b)] For every even integer $\mathrm{sv}\ge\mathrm{ncs}$ and every $O\in\mathcal{A}\left(Q^{\circ}_{0}(\mathcal{Q})\right)$,
\[
\mathcal{Z}_{[\mathrm{sv}]}(O)\le\mathcal{Z}_{[\mathrm{ncs}]}(O)^{\mathrm{kr}^3}.
\]
In particular, $\log Z_{[\mathrm{sv}]}\le\mathrm{kr}^3\cdot\log Z_{[\mathrm{ncs}]}$.
\item[(c)] If $\mathrm{sv}=\mathrm{ncs}$, or if $\mathrm{sv}$ is an even integer with $\mathrm{sv}\ge\check{j}\cdot\mathrm{ncs}$, then
\[
\log Z_{[\mathrm{sv}]}\ge\mathrm{kr}^3\cdot\log Z_{[\mathrm{ncs}]}-\left(\frac{\mathrm{kr}-1}{\check{j}-1}\right)^3\cdot D_{\check{j}}\ge\mathrm{kr}^3\cdot\left(\log Z_{[\mathrm{ncs}]}-\frac{D_{\check{j}}}{(\check{j}-1)^3}\right).
\]
\item[(d)] Let $\mathrm{sv}=\mathrm{ncs}$, or let $\mathrm{sv}$ be an even integer with $\mathrm{sv}\ge\check{j}\cdot\mathrm{ncs}$.
Then, for every set $\mathrm{Scl}$ of distinct cells of the torus $\mathrm{hw}\cdot\mathrm{sv}$ and for all $O_{\mathfrak{c}}\in\mathcal{A}\left(Q^{\circ}_{0}(\mathcal{Q})\right)$, $\mathfrak{c}\in\mathrm{Scl}$,
\[
\left|E_{[\mathrm{sv}]}\prod_{\mathfrak{c}\in\mathrm{Scl}}O_{\mathfrak{c}}^{[\mathfrak{c}]}\right|\le\prod_{\mathfrak{c}\in\mathrm{Scl}}\exp\left(\delta_{\check{j}}\right)\cdot\|O_{\mathfrak{c}}\|_{*}.
\]
\item[(e)] (The cell structure $\mathcal{Q}_{*}$.)
Consider the setting of Section~I.6.4, that is, $e\cdot h^{\frac{1}{2}}\le e_1$, $N\ge N_0(e)$, $q=e\cdot\eta^{\frac{1}{2}}$, and $\boldsymbol{\mu}$ an arbitrary real number.
Choose $\mathcal{Q}=\mathcal{Q}_{*}$, $\mathrm{hw}=\mathrm{hc}$, $\mathrm{ncs}=2L$ and $\check{j}=L$.
Then $n_r=n_*$, $N_r=N_{*}$, $D_{\check{j}}=D_L$, $\delta_{\check{j}}=\delta$ and $\|O\|_{*}=\left(E^{\mathrm{af}}_{n_*}O^{\sharp}\right)^{\frac{1}{N_{*}}}$.
Moreover $D_L\ge0$, and (b)--(d) hold on the cubic tori with $n$ sites per side for $n=n_*$ and for every $n\in\mathrm{per}\,\mathbb{Z}$ with $n\ge L\cdot n_*$;
for these tori $\mathrm{kr}=n/n_*$ and $\mathrm{sv}=2n/\mathrm{per}$.
They are the tori $\mathrm{kt}$ with $\mathrm{kt}=1$ or with $\mathrm{kt}\ge L$ an integer,
among them the second reference torus $\mathrm{kt}=L$ and the even tori $n=8h\cdot L^{T+N}$ with $T\ge T^{*}(e)+1$, for which $\mathrm{kt}=8L^{T-T^{*}(e)}$;
and, besides these, tori whose side is not a multiple of $n_*$.
\item[(c+)] For every even integer $\mathrm{sv}$ with $\mathrm{ncs}\le\mathrm{sv}\le\check{j}\cdot\mathrm{ncs}$,
\[
\log Z_{[\mathrm{sv}]}\ge\mathrm{kr}^3\cdot\left(\log Z_{[\mathrm{ncs}]}-\frac{D_{\check{j}}}{\check{j}^3}\right)\ge\mathrm{kr}^3\cdot\left(\log Z_{[\mathrm{ncs}]}-\frac{D_{\check{j}}}{(\check{j}-1)^3}\right).
\]
\item[(d+)] The inequality of (d) holds for every even integer $\mathrm{sv}\ge\mathrm{ncs}$, for every set $\mathrm{Scl}$ of distinct cells of the torus $\mathrm{hw}\cdot\mathrm{sv}$ and for all $O_{\mathfrak{c}}\in\mathcal{A}\left(Q^{\circ}_{0}(\mathcal{Q})\right)$.
\item[(e+)] In the setting of (e), the inequality of (d) holds on the cubic torus with $n$ sites per side for every $n\in\mathrm{per}\,\mathbb{Z}$ with $n\ge n_*$.
These tori include all tori $\mathrm{kt}$ with $\mathrm{kt}\ge1$ an integer, and the even tori $n=8h\cdot L^{T+N}$ with $T\ge T^{*}(e)$.
\end{itemize}
\end{named}

\begin{named}{Remark}
Proposition~40 corresponds to the parameters and reference numbers of \cite[Section~21.7.5]{QEDI}, to \cite[Theorem~21.25]{QEDI} and to \cite[Remark~21.27]{QEDI}.
It differs from these statements in the following points.
The ratio $\mathrm{kr}$ need not be an integer:
the side of the torus has only to be a multiple of $2\mathrm{hw}$ and at least $\check{j}\cdot n_r$, and in (c+), (d+), (e+) at least $n_r$.
In \cite[Theorem~21.25]{QEDI} the side of the torus is $\mathrm{kr}\cdot n_r$ with an integer $\mathrm{kr}\in\{1\}\cup\{\check{j},\check{j}+1,\dots\}$;
for the cells of \cite[Remark~21.27]{QEDI} the transfer inequality is therefore asserted there neither on tori whose side is not a multiple of $n_*$, nor on the tori $\mathrm{kt}$ with $1<\mathrm{kt}<L$.
By (d+) and (e+), the inequality of (d) holds on tori of both kinds whose side lies in $2\mathrm{hw}\,\mathbb{Z}$ (for $\mathcal{Q}_{*}$: in $\mathrm{per}\,\mathbb{Z}$) and is at least $n_r$;
Proposition~42.2, on the other hand, is stated for $n=n_*$ and for $n\in\mathrm{per}\,\mathbb{Z}$ with $n\ge L\cdot n_*$ only.
Cubic tori whose side lies outside $\mathrm{per}\,\mathbb{Z}$ do not carry $\mathcal{Q}_{*}$ (Section~I.6.4).
Clause~(d) is the special case of (d+) in which $\mathrm{sv}=\mathrm{ncs}$ or $\mathrm{sv}\ge\check{j}\cdot\mathrm{ncs}$.
The first inequality of (c) has the same hypothesis on $\mathrm{sv}$;
for $\mathrm{ncs}<\mathrm{sv}<\check{j}\cdot\mathrm{ncs}$ the lower bound for $\log Z_{[\mathrm{sv}]}$ is that of (c+).
The contents of the cells are elements of the algebra $\mathcal{A}\left(Q^{\circ}_{0}(\mathcal{Q})\right)$ of the open cell (Proposition~39), in place of the closed-cell class $\mathfrak{A}^{+}$ of \cite[Theorem~21.25]{QEDI};
by Proposition~39\,(b) one and the same element is read on all tori that carry $\mathcal{Q}$.
The lattice mass $\boldsymbol{\mu}$ is an arbitrary real number, whereas \cite[Section~21.7.5]{QEDI} assumes $\boldsymbol{\mu}>-\frac{1}{2}$.
Only the all-flux law occurs; no transfer statement for the harmonic-averaged law is made.
\end{named}

\subsection*{42. The defect of the two reference tori; the transfer inequality for Wick powers}

\subsubsection*{Proposition~42.1. The three parts of the defect}

\begin{named}{Proposition 42.1}
Let $n\ge4$ be an integer, and consider the cubic torus with $n$ sites per side, of either parity.
Let $q>0$, let $\boldsymbol{\mu}$ be real, and let $\check{j}\ge2$ be an integer.
Let $Z^{\mathrm{ph}}_n$ and $\mathrm{Ld}(n)$ be as in Proposition~38.2\,(c), put $\mathrm{zr}_n:=\mathrm{zr}_n(q,\boldsymbol{\mu})$, and let
\[
\varsigma_n:=\sum_{w\in\mathbb{Z}^3,\ w\ne0}\mathfrak{q}_w(n),\qquad\mathfrak{q}_w(n)=\exp\left(-\frac{2\pi^2\cdot|w|_2^2}{q^2\cdot n}\right),
\]
\[
G_{\check{j}}:=\left(1+\frac{\check{j}}{3.2}\right)\cdot\frac{2\check{j}^2}{2.56},\qquad
K_{\check{j}}:=\left(1+\check{j}+\check{j}^2\right)\cdot\left(\log\check{j}+\check{j}\cdot G_{\check{j}}\right).
\]
\begin{itemize}
\item[(a)] One has
\[
Z^{(0)}_n[1]=Z^{\mathrm{ph}}_n\cdot\mathrm{zr}_n,\qquad Z_n[1]=\frac{Z^{\mathrm{ph}}_n\cdot\mathrm{zr}_n}{P_n(0)};
\]
moreover $P_n(w)\le\mathfrak{q}_w(n)$ for every $w\in\mathbb{Z}^3$, and $1-\varsigma_n\le P_n(0)\le1$.
\item[(b)] If $q^2\cdot n\le1$, then
\[
\varsigma_n\le7\exp\left(-\frac{2\pi^2}{q^2\cdot n}\right)\le\frac{1}{8}.
\]
\item[(c)] One has
\[
\check{j}^3\cdot\log Z_n[1]-\log Z_{\check{j}\cdot n}[1]=D^{\mathrm{ph}}+D^{\mathrm{mat}}_{\check{j}}+D^{\mathrm{fl}},
\]
where
\begin{align*}
D^{\mathrm{ph}}&:=\left(\check{j}^3-1\right)\cdot\log\frac{2\pi}{q^2}+\frac{1}{2}\left(\mathrm{Ld}(\check{j}\cdot n)-\check{j}^3\cdot\mathrm{Ld}(n)\right)\\
&\phantom{:}\le\left(\check{j}^3-1\right)\cdot\log\frac{2\pi}{q^2\cdot n}+K_{\check{j}},\\
D^{\mathrm{mat}}_{\check{j}}&:=\check{j}^3\cdot\log\mathrm{zr}_n-\log\mathrm{zr}_{\check{j}\cdot n},\\
D^{\mathrm{fl}}&:=-\check{j}^3\cdot\log P_n(0)+\log P_{\check{j}\cdot n}(0);
\end{align*}
and if $\varsigma_n<1$, then
\[
D^{\mathrm{fl}}\le-\check{j}^3\cdot\log\left(1-\varsigma_n\right).
\]
\end{itemize}
\end{named}

\begin{named}{Remark}
Proposition~42.1 corresponds to \cite[Lemma~21.32]{QEDI}.
The bound for $D^{\mathrm{ph}}$ is equivalent to \cite[Lemma~21.31]{QEDI}, an inequality for the Laplacian of the cubic torus into which no fermion enters, and which holds here as stated.
The numbers $\mathrm{zr}_n$, $Z_n[1]$ and $P_n(w)$ are formed with the determinant $\det\left(D_A+\boldsymbol{\mu}\right)$ of Section~I.7.2 in place of the Wilson determinant of \cite[Lemma~21.32]{QEDI}.
The lattice mass $\boldsymbol{\mu}$ is any real number and $n\ge4$, while \cite[Lemma~21.32]{QEDI} assumes $\boldsymbol{\mu}>-\frac{1}{2}$ and $n\ge2$.
The first identity of (a) is Proposition~38.2\,(c);
the bound $P_n(w)\le\mathfrak{q}_w(n)$ is \cite[Lemma~13.12(a)]{QEDII}, which takes the place of \cite[Corollary~10.8(d)]{QEDI}.
\end{named}

\subsubsection*{Proposition~42.2. The transfer inequality}

\begin{named}{Proposition 42.2}
Let $0<e\le e_G$, $N\ge N_0(e)$, $\eta=L^{-N}$, $q=e\cdot\eta^{\frac{1}{2}}$ and $R^{*}=R^{*}(e)$.
Let $n_*$, $\mathrm{per}$, $\mathrm{hc}$, $\mathcal{Q}_{*}$, $N_{*}=8L^3$ and $a_L$ be as in Sections~I.6.4 and~I.7.6,
and let $K_L$ be as in Proposition~14.2\,(a);
it is the number $K_{\check{j}}$ of Proposition~42.1 at $\check{j}=L$.
Let $\boldsymbol{\mu}$ be an arbitrary real number;
the numbers $\mathrm{zr}_n=\mathrm{zr}_n(q,\boldsymbol{\mu})$, $D^{\mathrm{mat}}=L^3\cdot\log\mathrm{zr}_{n_*}-\log\mathrm{zr}_{L\cdot n_*}$, $D_L$ and $\delta$ are formed with this $\boldsymbol{\mu}$.
Let $\beta$ be a probe that fits $\left(N,T^{*}(e)\right)$, let $v$ be an arbitrary placement,
and let $\hat{\mathfrak{a}}$, $\boldsymbol{\sigma}^2$, $X$, $X_{\mathfrak{c}}$, $H_2$, $H_4$ and $\mathring{\Phi}$ be as in Section~I.7.6, where $\mathring{\Phi}$ is formed with $\boldsymbol{\mu}$.
Let $c_2$ and $c_4$ be arbitrary real numbers, and let
\[
\Theta^{(0)}_m:=\left|\left(\prod_{\mathfrak{c}\in\mathrm{Cel}}\left(\partial_{t_{\mathfrak{c}}}^m-c_m\right)\mathring{\Phi}\right)(0)\right|^{\frac{1}{N_{*}}}\qquad(m=2,4)
\]
be the display of Section~I.7.6 formed with these numbers $c_m$, and not with $N_0^{(m)}(0)$.
For $m=2,4$ put
\begin{gather*}
\mathrm{Mw}_4:=600\,N_{*}^2\cdot\boldsymbol{\sigma}^4\cdot e^{-4}+|c_4|,\qquad
\mathrm{Mw}_2:=12\,N_{*}\cdot\boldsymbol{\sigma}^2\cdot e^{-2}+|c_2|,\\
\mathrm{Tw}_m:=\Theta^{(0)}_m+2e^6\cdot\mathrm{Mw}_m,\qquad
\mathsf{u}_m:=\exp(\delta)\cdot\mathrm{Tw}_m.
\end{gather*}
Let $n=n_*$, or let $n\in\mathrm{per}\,\mathbb{Z}$ with $n\ge L\cdot n_*$,
and let $\kappa_2$ and $\kappa_4$ be the second and the fourth cumulant of $X=\langle d\hat{\mathfrak{a}},F^{\mathrm{u}}\rangle$ under $E^{\mathrm{af}}_n$.
Then
\begin{align}
\left|e^{-2}\cdot\left(\kappa_2-\boldsymbol{\sigma}^2\right)-c_2\right|&\le\mathsf{u}_2,\tag{2.1}\\
\left|e^{-4}\cdot\kappa_4-\left(c_4-3c_2^2\right)\right|&\le\mathsf{u}_4+3\mathsf{u}_2\cdot\left(2|c_2|+\mathsf{u}_2\right),\tag{2.2}\\
0\le\delta&\le\frac{a_L}{N_{*}}\cdot\log\frac{2\pi}{e^2\cdot R^{*}}+\frac{K_L+D^{\mathrm{mat}}+1}{(L-1)^3\cdot N_{*}}.\tag{2.3}
\end{align}
For every set $\mathrm{Scl}$ of distinct cells of the torus $n$ and for $m=2,4$,
\begin{equation}
\left|E^{\mathrm{af}}_n\prod_{\mathfrak{c}\in\mathrm{Scl}}\left(e^{-m}\cdot H_m(X_{\mathfrak{c}})-c_m\right)\right|\le\mathsf{u}_m^{\#|\mathrm{Scl}|}.\tag{2.4}
\end{equation}
Finally, with $O_m:=e^{-m}\cdot H_m(X)-c_m$,
\begin{equation}
\left(E^{\mathrm{af}}_{n_*}O_m^{\sharp}\right)^{\frac{1}{N_{*}}}\le\mathrm{Tw}_m;\tag{2.4$'$}
\end{equation}
the left side of (2.4$'$) is the reference seminorm $\|O_m\|_{*}$ of Proposition~40\,(e).
\end{named}

\begin{named}{Remark}
Proposition~42.2 corresponds to \cite[Theorem~21.2]{QEDI} for the all-flux law, together with Steps~1--4 of its proof in \cite[Section~21.8.2]{QEDI}.
In \cite[Theorem~21.2]{QEDI} the conditions $0<e\le\frac{1}{3}$ and $6N_{*}\cdot e^2\cdot R^{*}\cdot\log\left(e^{-1}\right)\le\pi^2$ are hypotheses.
Here they are not:
for $e\le e_G$ both hold by the Propositions~14.m, the second being the condition (W3) of Proposition~14.2\,(d);
they are conditions on the constants $e_1$ and $e_G$.
The lattice mass is an arbitrary real number, whereas \cite[Theorem~21.2]{QEDI} assumes $\boldsymbol{\mu}_N>-\frac{1}{2}$.
More tori are admitted:
$n=n_*$, or any multiple $n$ of $\mathrm{per}$ with $n\ge L\cdot n_*$, of either parity;
in \cite[Theorem~21.2]{QEDI} the tori are those with $\mathrm{kt}\cdot n_*$ sites per side, $\mathrm{kt}=1$ or $\mathrm{kt}\ge L$.
Of the two cases $\mathrm{nat}\in\{\mathrm{af},\mathrm{av}\}$ of \cite[Theorem~21.2]{QEDI}, only $\mathrm{nat}=\mathrm{af}$ has a counterpart:
the cumulants in (2.1), (2.2) and the expectations in (2.4), (2.4$'$) are taken under the all-flux law, and none of these bounds is formulated under the harmonic-averaged law $E^{\mathrm{av}}_n$.
Accordingly the numbers $\boldsymbol{\Gamma}^{\mathrm{av}}$, $\boldsymbol{\Theta}^{\mathrm{av}}_m$, $Y$, $\check{X}$ of \cite[Theorem~21.2]{QEDI}, the fifth step of its proof in \cite[Section~21.8.2]{QEDI} and \cite[Proposition~21.28]{QEDI} have no counterpart here.
The harmonic-averaged law enters only on the reference torus $n_*$, through the number $E^{\mathrm{av}}_{n_*}O_m^{\sharp}$.
In the fourth step of the proof in \cite[Section~21.8.2]{QEDI} the bound for $E^{\mathrm{af}}_{n_*}O_m^{\sharp}$ is derived from the equality $E^{\mathrm{av}}_{n_*}O_m^{\sharp}=\left(\Theta^{(0)}_m\right)^{N_{*}}\ge0$ of the third step there.
Without any condition on the sign of $E^{\mathrm{av}}_{n_*}O_m^{\sharp}$ one has
\[
E^{\mathrm{af}}_{n_*}O_m^{\sharp}=P_{n_*}(0)\cdot E^{\mathrm{av}}_{n_*}O_m^{\sharp}+E^{\mathrm{af}}_{n_*}\left[O_m^{\sharp}\cdot\mathbf{1}[\mathrm{flux}\ne0]\right]\le\left|E^{\mathrm{av}}_{n_*}O_m^{\sharp}\right|+\varsigma_{n_*}^{\frac{1}{2}}\cdot\left(E^{\mathrm{af}}_{n_*}\left|O_m^{\sharp}\right|^2\right)^{\frac{1}{2}},
\]
and $\left|E^{\mathrm{av}}_{n_*}O_m^{\sharp}\right|=\left(\Theta^{(0)}_m\right)^{N_{*}}$ by Proposition~38.3\,(e).
In \cite[Theorem~21.2]{QEDI} the numbers $\exp(\delta)$, $\mathrm{Tw}_m$ and $\mathrm{Mw}_m$ are denoted $\boldsymbol{\Gamma}^{\mathrm{af}}$, $\boldsymbol{\Theta}^{\mathrm{af}}_m$ and $\dot{M}_m$.
Proposition~42.2 holds for every placement $v$, and $\hat{\mathfrak{a}}$, $X$ and the $X_{\mathfrak{c}}$ are formed on each of the tori $n$ named in it.
The numerals $600$, $12$, $2e^6$ and the constant $K_L$ in the statement of Proposition~42.2 are those of the statement of \cite[Theorem~21.2]{QEDI}.
The numeral $7$ is that of Proposition~42.1\,(b), equal to the numeral of \cite[Lemma~21.32(c)]{QEDI};
it does not occur in the statement of Proposition~42.2.
It enters the proof of (2.3), where it bounds $\varsigma_{n_*}$ and with it $D^{\mathrm{fl}}$, as in Step~1 of the proof of \cite[Theorem~21.2]{QEDI} in \cite[Section~21.8.2]{QEDI},
and the proof of the factor $2e^6$ in $\mathrm{Tw}_m$.
These numbers come from Gaussian domination with the constant one and from the arithmetic of the photon field.
\end{named}

\subsection*{43. Reference data; the ultraviolet proposition}

\subsubsection*{Proposition~43.1. Reference data}

\begin{named}{Proposition 43.1}
The clauses (R0)--(R2) below hold for every $e>0$ with $e\cdot h^{\frac{1}{2}}\le e_1$, and the clause (R6) does not contain $e$.
In the assertion that follows (R6) the two constants $e_{\sharp}(\beta)$ and $C^{G}_{\beta}$ are fixed first, at the line (C10) of Section~I.8.2,
and the charge $e$ is chosen after them, at the line (C11).
The numbers $T^{*}(e)$, $R^{*}(e)$, $N_0(e)$, $m_{0,N}(e)$ and the lattice mass $\boldsymbol{\mu}_N=\bar{m}\cdot\eta+m_{0,N}(e)$ are those of Sections~I.6.2 and~I.2.5, and
\[
C_{\mathrm{mat}}:=\left(L^3+1\right)\cdot\left(\varepsilon_5\left(\frac{1}{L}\right)+1\right).
\]
\begin{itemize}
\item[(R0)] For every integer $N\ge0$ the number $m_{0,N}(e)$ is real, and $\left|m_{0,N}(e)\right|<\frac{1}{8}$.
These numbers depend on $(N,e,\bar{m},\rho)$ and on the scheme (Section~I.6.2); they depend neither on a torus nor on a twist.
Hence $N\mapsto m_{0,N}(e)$ is a mass assignment in the sense of Section~I.7.5.
\item[(R1)] One has $e\le\frac{1}{3}$.
The number $T^{*}(e)$ is an integer with $T^{*}(e)\ge1$, one has $R^{*}(e)=h\cdot L^{T^{*}(e)}$, and both $R^{*}(e)$ and $L\cdot R^{*}(e)$ belong to $\mathrm{Tor}(e)$.
The length $\ell_c=\frac{R^{*}(e)}{2L}$ tends to $\infty$ as $e\to0$, and
\[
\exp\left(-\bar{m}\cdot R^{*}(e)\right)\le e\cdot h^{\frac{1}{2}}.
\]
If $e\le e_G$, then the inequality (W3) of Proposition~14.2\,(d) holds:
\[
6N_{*}\cdot e^2\cdot R^{*}(e)\cdot\log\frac{1}{e}\le\pi^2.
\]
\item[(R2)] For every $N\ge N_0(e)$, with $q=e\cdot\eta^{\frac{1}{2}}$ and $\boldsymbol{\mu}=\boldsymbol{\mu}_N$, one has
\[
D^{\mathrm{mat}}\le C_{\mathrm{mat}}.
\]
\item[(R6)] The exponents satisfy
\[
0<\mathfrak{e}\le\frac{1}{4},\qquad0\le\mathfrak{e}_2<1,\qquad0<\sigma=\mathfrak{e}-\frac{2a_L}{N_{*}}\le\frac{1}{4}.
\]
\end{itemize}
For every probe $\beta$ there are numbers $e_{\sharp}(\beta)\in\left(0,e_1\cdot h^{-\frac{1}{2}}\right]$ and $C^{G}_{\beta}<\infty$,
which depend only on the constants of the lines (C0)--(C9), on $\bar{m}$ and on $\beta$, with the following property.
Let $0<e\le\min\left(e_G,e_{\sharp}(\beta)\right)$, let $N\ge N_0(e)$, and let $v$ be any placement;
put $\boldsymbol{\mu}=\boldsymbol{\mu}_N$ and $c_m:=N_0^{(m)}(0)$, where $N_0$ is the free ratio on $C(n_*)$ (Section~I.7.6).
Then:
\begin{itemize}
\item[(R3)] the probe $\beta$ fits $\left(N,T^{*}(e)\right)$, and
\[
\Theta^{(0)}_4\le C^{G}_{\beta}\cdot e^{\mathfrak{e}},\qquad\left(1+|c_2|\right)\cdot\Theta^{(0)}_2\le C^{G}_{\beta}\cdot e^{\mathfrak{e}};
\]
\item[(R4)] the numbers $c_2$ and $c_4$ are real, they do not depend on the placement, and
\[
|c_2|+|c_4|^{\frac{1}{2}}\le C^{G}_{\beta}\cdot e^{-\mathfrak{e}_2};
\]
\item[(R5)] one has $c_4-3c_2^2=b_{\mathrm{ch}}\left(\eta,R^{*}(e)\right)$, and
\[
\left|c_4-3c_2^2-b\left(d\,d^{*}\beta\right)\right|\le C^{G}_{\beta}\cdot\left(e+\eta^{\frac{1}{4}}\right).
\]
\end{itemize}
The constant is
\[
C^{G}_{\beta}:=\max\left\{C_{35}(\beta),\ C_{\mathrm{R4}},\ C_{1}^{\mathrm{N2}},\ C_{2}^{\mathrm{N2}}\cdot\exp\left(\bar{m}\cdot D_{\mathrm{be}}\right)\cdot h^{\frac{1}{2}}\right\},
\]
where
\begin{align*}
C_{35}(\beta)&:=3\cdot10^5\cdot\mathrm{Cb}^4\cdot h^{\frac{1}{8}},\\
C_{\mathrm{R4}}&:=\mathrm{Cpv}_{\beta}+\left(\left|b\left(d\,d^{*}\beta\right)\right|+C_{1}^{\mathrm{N2}}+C_{2}^{\mathrm{N2}}+3\,\mathrm{Cpv}_{\beta}^2\right)^{\frac{1}{2}}.
\end{align*}
The first of these is the constant of Proposition~35.9\,(b).
The second is the constant of Proposition~37.3\,(b);
since $R^{*}(e)\ge D_{\mathrm{be}}+1$, one has $\exp\left(-\bar{m}\cdot\left(R^{*}(e)-D_{\mathrm{be}}\right)\right)\le1$.
The numbers $\mathrm{Cb}$, $\mathrm{Cpv}_{\beta}$, $C_{1}^{\mathrm{N2}}$, $C_{2}^{\mathrm{N2}}$ and $D_{\mathrm{be}}$ are those of the probe $\beta$ in Section~G5.1.3, Proposition~36\,(g) and Proposition~37.2\,(d), (e).
\end{named}

\begin{named}{Remark}
Proposition~43.1 corresponds to \cite[Theorem~21.3]{QEDI}, restated as \cite[Theorem~5.95]{QEDI}, together with its proof in \cite[Section~21.9]{QEDI}.
The differences are the following.
The numbers $D^{\mathrm{mat}}$ and $\Theta^{(0)}_m$ and the free ratio $N_0$ are formed with the determinant of Section~I.7.2 in place of the Wilson determinant.
The numbers $h'$ and $m_0$ of \cite[Theorem~21.3]{QEDI} do not occur, since $h'=h$ and $m_0=0$;
accordingly $C_{\mathrm{mat}}$ contains $\varepsilon_5\left(\frac{1}{L}\right)$, formed with the function $\varepsilon_5$ of Section~I.8.2, which takes the place of the function $\varepsilon_5$ of \cite[Section~5.6.2]{QEDI}.
In (R5) the rate in the spacing is $\eta^{\frac{1}{4}}$, where \cite[Theorem~21.3]{QEDI} has $\eta^{\frac{1}{2}}$.
Here $\mathfrak{e}_2=0$ (Section~I.8.5), whereas $\mathfrak{e}_2=6\varepsilon$ in \cite[Theorem~5.95]{QEDI};
and (R4) is obtained from Proposition~36 and the Propositions~37.m, not from the renormalisation group representation of the free determinant.
That $N\mapsto m_{0,N}(e)$ is a mass assignment does not require $\boldsymbol{\mu}_N>-\frac{1}{2}$,
and none of the clauses requires a sign of $\boldsymbol{\mu}_N$ (I.K10).
The dependence on $\rho$ is stated in (R0).
In Section~I.6.2 the number $m_{0,N}(e)$ is the first entry of $P_N(e)$ for every $N\ge N_{\min}(R_{\mathrm{top}}(e),e)$ and $0$ for $0\le N<N_{\min}(R_{\mathrm{top}}(e),e)$,
whereas \cite[Theorem~21.3]{QEDI} puts $m_{0,N}(e)=0$ for every $N<N_0(e)$;
since $N_{\min}(R_{\mathrm{top}}(e),e)\le N_0(e)$, the two rules differ only for the finitely many $N$ with $N_{\min}(R_{\mathrm{top}}(e),e)\le N<N_0(e)$,
and among the clauses of Proposition~43.1 only (R0) concerns these $N$.
Where \cite[Theorem~21.3]{QEDI} takes the placement of \cite[Section~21.1.6]{QEDI}, the clauses (R3)--(R5) hold for every placement;
that $c_m$ does not depend on the placement is Proposition~50\,(h).
In the definition of $C^{G}_{\beta}$ the exponent contains $D_{\mathrm{be}}$, the least side of a closed coordinate cube that contains $\operatorname{supp}\beta$;
one has $D_{\mathrm{be}}=D_{\beta}-4$.
\end{named}

\subsubsection*{Proposition~43.2. The ultraviolet proposition}

\emph{Data.}
Let $\beta$ be a probe, and let $r_{\beta}>0$ be a number such that $\operatorname{supp}\beta$ lies in a ball of radius $r_{\beta}$.
This number is a datum in addition to $(\beta,x_0)$;
it enters only $d_0$ and the inequality displayed below, and none of the constants below depends on it.
Put
\[
a=d^{*}\beta,\qquad\|a\|_{\infty}:=\sup_{y\in\mathbb{R}^3}|a(y)|_2,\qquad f=d\,d^{*}\beta,
\]
and assume $b(f)>0$;
such probes exist by \cite[Theorem~21.1(a)]{QEDI}.
Let $x_0\in\mathbb{R}^3$ be such that $d_0:=|x_0|_2-2r_{\beta}>0$ and such that the condition of \cite[Section~2.8]{QEDI} holds:
\[
\frac{6}{\bar{m}^2}\cdot\|a\|_{\infty}^4\cdot\left(\frac{4\pi\cdot r_{\beta}^3}{3}\right)^2\cdot\left(\bar{m}^2+\left(\bar{m}+\frac{1}{d_0}\right)^2\right)\cdot\frac{\exp\left(-2\bar{m}\cdot d_0\right)}{4\pi^2\cdot d_0^2}\le\frac{b(f)}{12}.
\]
Let $\beta_0:=\beta$ and $\beta_{\pm}:=\beta\pm\tau_{x_0}\beta$; these are three probes.
Put $g_{\mathrm{ix}}:=d\,d^{*}\beta_{\mathrm{ix}}$ for $\mathrm{ix}\in\{0,+,-\}$, so that $g_0=f$ and $g_{\pm}=f\pm\tau_{x_0}f$.
For $0<\eta\le1$ let $g^{(\eta)}_{\mathrm{ix}}$ be the two-form $f^{(\eta)}$ defined at the beginning of Proposition~47.1, formed with the probe $\beta_{\mathrm{ix}}$ and not with the test $g_{\mathrm{ix}}$;
below, $\eta=L^{-N}$.

\emph{Constants.}
The following constants belong to the line (C10) and are fixed in the order in which they are listed;
each of them depends only on the preceding ones and on $(\bar{m},\beta,x_0)$.
\begin{align*}
s_{B}&:=\max_{\mathrm{ix}}\left(\|g_{\mathrm{ix}}\|_2+\|\beta_{\mathrm{ix}}\|_{H^4}\right),\qquad C_{3\mathrm{p}}:=\max_{\mathrm{ix}}C^{G}_{\beta_{\mathrm{ix}}},\\
C_{\Gamma}^{(16)}&:=(2\pi)^{\frac{a_L}{N_{*}}}\cdot\exp\left(\frac{K_{L}+C_{\mathrm{mat}}+1}{(L-1)^3\cdot N_{*}}\right),
\end{align*}
\begin{multline*}
C_{U}:=\max\left\{C_{\Gamma}^{(16)}\cdot\left(C_{3\mathrm{p}}+1200N_{*}^2\cdot s_{B}^4+2C_{3\mathrm{p}}^2\right),\right.\\
\left.C_{\Gamma}^{(16)}\cdot\left(C_{3\mathrm{p}}+2\left(12N_{*}\cdot s_{B}^2+C_{3\mathrm{p}}\right)\cdot\left(1+C_{3\mathrm{p}}\right)\right),\ C_{3\mathrm{p}},\ 60\,s_{B}^4\right\},
\end{multline*}
\[
C_{\bar{\kappa}}:=2\left(8C_{U}+3C_{U}^2\right),\qquad C'_{\bar{\kappa}}:=8C_{\bar{\kappa}}+3C_{\bar{\kappa}}^2;
\]
finally, $e_2$ is the largest number such that
\[
e_2\le e_G,\qquad e_2\le e_{\sharp}(\beta_{\mathrm{ix}})\quad\text{for }\mathrm{ix}\in\{0,+,-\},\qquad\frac{5}{2}\,C'_{\bar{\kappa}}\cdot e_2^{\sigma}\le\frac{b(f)}{4}.
\]

\begin{named}{Proposition 43.2}
Let $e\in(0,e_2]$, and let $N\ge N_0(e)$ and $T\ge T^{*}(e)+1$ be integers.
Put
\[
\eta=L^{-N},\qquad n=8h\cdot L^{T+N},\qquad R=8h\cdot L^{T},\qquad q=e\cdot\eta^{\frac{1}{2}},\qquad\boldsymbol{\mu}=\boldsymbol{\mu}_N=\bar{m}\cdot\eta+m_{0,N}(e),
\]
and consider the law $E^{\mathrm{af}}_n$, under which $X_{\eta,R}(\,\cdot\,)$ has the law given by $P_{N,T}$.
For $\mathrm{ix}\in\{0,+,-\}$ let $v$ be a placement of $\beta_{\mathrm{ix}}$, and let
\[
c_m=c_m(N,e,\beta_{\mathrm{ix}}):=N_0^{(m)}(0),\qquad\boldsymbol{\sigma}^2=\boldsymbol{\sigma}^2(N,\beta_{\mathrm{ix}}),\qquad\mathsf{u}_m=\mathsf{u}_m(N,e,\beta_{\mathrm{ix}};v)\qquad(m=2,4)
\]
be the numbers of Proposition~42.2;
here $\boldsymbol{\sigma}^2=\|d\hat{\mathfrak{a}}\|^2$ is the number of Section~I.7.6, formed with $\tau_v\beta_{\mathrm{ix}}$, and not the square of the exponent $\sigma$ of Section~I.8.5;
through $\Theta^{(0)}_m$ the number $\mathsf{u}_m$ may depend on the placement $v$.
All the following clauses hold for every placement $v$ of $\beta_{\mathrm{ix}}$.
\begin{itemize}
\item[(U-a)] One has
\begin{multline*}
\left|\kappa_4\left(X_{\eta,R}(g_{\mathrm{ix}})\right)-e^4\cdot b(g_{\mathrm{ix}})\right|\le\frac{C_{\bar{\kappa}}}{2}\cdot e^{4+\sigma}+C_{U}\cdot\left(\eta^2+e^4\cdot\eta^{\frac{1}{4}}\right)\\
\le C_{\bar{\kappa}}\cdot\left(e^{4+\sigma}+\eta^{\frac{1}{4}}\right).
\end{multline*}
\item[(U-a$'$)] One has
\begin{gather*}
\left|\kappa_4\left(X_{\eta,R}(g_{\mathrm{ix}})\right)-e^4\cdot b_{\mathrm{ch}}\left(\eta,R^{*}(e)\right)[\beta_{\mathrm{ix}}]\right|\le\left(7C_{U}+3C_{U}^2\right)\cdot e^{4+\sigma}+C_{U}\cdot\eta^2,\\
\left|\kappa_2\left(X_{\eta,R}(g_{\mathrm{ix}})\right)-\boldsymbol{\sigma}^2-e^2\cdot c_2\right|\le e^2\cdot\mathsf{u}_2+2s_{B}^2\cdot\eta^2.
\end{gather*}
\item[(U-b)] For every finite subset $\mathrm{Ypt}$ of $\frac{R^{*}(e)}{L}\,\mathbb{Z}^3$ whose points are pairwise incongruent modulo $R\mathbb{Z}^3$, and for $m=2,4$,
\[
\left|E^{\mathrm{af}}_n\prod_{y\in\mathrm{Ypt}}\mathrm{Pw}_m\left(X_{\eta,R}\left(\tau_y\,g^{(\eta)}_{\mathrm{ix}}\right);\boldsymbol{\sigma}^2,c_m\right)\right|\le\mathsf{u}_m^{\#|\mathrm{Ypt}|}.
\]
\item[(U-c)] One has
\begin{gather*}
\mathsf{u}_4\le C_{U}\cdot e^{\sigma},\qquad\left(1+|c_2|\right)\cdot\mathsf{u}_2\le C_{U}\cdot e^{\sigma},\\
|c_2|\le C_{U}\cdot e^{-\mathfrak{e}_2},\qquad\left|c_4-3c_2^2-b(g_{\mathrm{ix}})\right|\le C_{U}\cdot\left(e^{\sigma}+\eta^{\frac{1}{4}}\right).
\end{gather*}
\end{itemize}
In consequence, the following holds.
The exponents $\sigma$ and $\mathfrak{e}_2$ of Section~I.8.5 satisfy $\sigma>0$ and $\mathfrak{e}_2\in[0,1)$;
the constants above satisfy $C_{\bar{\kappa}}<\infty$, $e_2\in(0,1]$ and $\frac{5}{2}\,C'_{\bar{\kappa}}\cdot e_2^{\sigma}\le\frac{b(f)}{4}$.
Let $e\in(0,e_2]$.
Then $N\mapsto m_{0,N}(e)$ is a mass assignment, which does not depend on $T$;
$N_0(e)$ and $T^{*}(e)+1$ are integers;
and the number $\frac{R^{*}(e)}{L}>0$ satisfies $\frac{R^{*}(e)}{L\cdot\eta}\in\mathbb{Z}$ for $N\ge N_0(e)$.
For $N\ge N_0(e)$ and $\mathrm{ix}\in\{0,+,-\}$ fix one placement $v(N,e,\mathrm{ix})$ of $\beta_{\mathrm{ix}}$, and form $c_m$ and $\mathsf{u}_m$, $m=2,4$, with it.
Then $c_m$ is real, $\mathsf{u}_m\ge0$ and $\boldsymbol{\sigma}^2\ge0$;
the test $g_N:=g^{(\eta)}_{\mathrm{ix}}$ is real, lies in $C_c^{\infty}(\mathbb{R}^3;\Lambda^2)$ and is supported in the cube $[-r_g,r_g]^3$, where
\[
r_g:=|y_{\beta_{\mathrm{ix}}}|_{\infty}+\frac{D_{\beta_{\mathrm{ix}}}}{2}
\]
does not depend on $N$;
and $\|g_N-g_{\mathrm{ix}}\|_2\to0$ and $\boldsymbol{\sigma}^2\to\|g_{\mathrm{ix}}\|_2^2$ as $N\to\infty$.
Finally, let $N\ge N_0(e)$ and $T\ge T^{*}(e)+1$.
Then
\[
\left|\kappa_4\left(X_{\eta,R}(g_{\mathrm{ix}})\right)-e^4\cdot b(g_{\mathrm{ix}})\right|\le C_{\bar{\kappa}}\cdot\left(e^{4+\sigma}+\eta^{\frac{1}{4}}\right);
\]
for every integer $J\ge1$, all pairwise distinct points $y_1,\dots,y_J\in\frac{R^{*}(e)}{L}\,\mathbb{Z}^3$ with $R>2\max_i|y_i|_{\infty}+2r_g$, and $m=2,4$,
\[
\left|E^{\mathrm{af}}_n\prod_{i=1}^{J}\mathrm{Pw}_m\left(X_{\eta,R}\left(\tau_{y_i}g_N\right);\boldsymbol{\sigma}^2,c_m\right)\right|\le\mathsf{u}_m^{J};
\]
and
\begin{gather*}
\mathsf{u}_4\le C_{\bar{\kappa}}\cdot e^{\sigma},\qquad\left(1+|c_2|\right)\cdot\mathsf{u}_2\le C_{\bar{\kappa}}\cdot e^{\sigma},\\
|c_2|\le C_{\bar{\kappa}}\cdot e^{-\mathfrak{e}_2},\qquad\left|c_4-3c_2^2-b(g_{\mathrm{ix}})\right|\le C_{\bar{\kappa}}\cdot\left(e^{\sigma}+\eta^{\frac{1}{4}}\right).
\end{gather*}
\end{named}

\begin{named}{Remark}
Proposition~43.2 corresponds to Steps~0--3 and to the beginning of Step~6 of the proof in \cite[Section~21.10.1]{QEDI}, with the inequalities displayed there,
to \cite[Theorem~3.1]{QEDI}, Part~II\,(ii),
and to the four conditions on $e_2$ in \cite[Section~5.6.2]{QEDI}.
The differences are the following.
The rate in the spacing is $\eta^{\frac{1}{4}}$ in place of $\eta^{\frac{1}{2}}$.
The law is the all-flux law on the even tori $n=8h\cdot L^{T+N}$, with $T\ge T^{*}(e)+1$,
whereas \cite[Theorem~3.1]{QEDI}, Part~II\,(ii), takes $T\ge T^{*}(e)$ on the odd tori $n=h\cdot L^{T+N}$.
The constants are explicit functions of $C_{3\mathrm{p}}$, $s_{B}$ and $C_{\Gamma}^{(16)}$, whereas the proof in \cite[Section~21.10.1]{QEDI} uses unspecified constants $C$.
The third condition on $e_2$ in \cite[Section~5.6.2]{QEDI}, the constant $C_Y$ and the number $\check{X}$ of Step~1 of \cite[Section~21.10.1]{QEDI} concern the averaged law only;
nothing corresponds to them here.
The last condition in the definition of $e_2$ is stated with $C'_{\bar{\kappa}}$;
it implies
\[
\frac{5}{2}\,C_{\bar{\kappa}}\cdot e_2^{\sigma}\le\frac{b(f)}{4}\qquad\text{and}\qquad C'_{\bar{\kappa}}\cdot e_2^{\sigma}\le\frac{b(f)}{10},
\]
the first of which is the fourth condition on $e_2$ in \cite[Section~5.6.2]{QEDI}.
In (U-b) the factors are indexed by the points of a finite set, and the points are only required to be pairwise incongruent modulo $R\mathbb{Z}^3$;
this requirement is weaker than the one of the last paragraph, that the points be pairwise distinct with $R>2\max_i|y_i|_{\infty}+2r_g$,
and it is exactly the condition under which the corresponding cells are distinct (Proposition~47.2\,(c)).
The clause (U-a$'$) contains the first inequality of \cite[Theorem~3.1]{QEDI}, Part~II\,(ii), and the bounds on the second cumulant among the inequalities displayed in Step~3 of \cite[Section~21.10.1]{QEDI}.
Finally, (U-c) holds with $C_{U}\le C_{\bar{\kappa}}$, so that the bounds of the last paragraph hold a fortiori.
\end{named}

\subsection*{Notation for laws on $\mathcal{S}'$, used in the Propositions~45.m and in Theorem~44}

The notation below agrees with that of \cite[Sections~8.1.1--8.1.6, 7.5 and~7.6]{QEDI} and of \cite[Definitions~2.38 and~14.61]{QEDII};
the space denoted $Q$ in \cite[Section~8]{QEDI} is written $\mathcal{S}'$ here.
Throughout, $\mathrm{lw}$ denotes a law, $\mathrm{po}$ a real number, the position of a plane, $\hat{n}$ a unit normal vector, and $\mathrm{Gi}$ a Euclidean map of $\mathbb{R}^3$.
\begin{itemize}
\item[(n1)] $\mathcal{S}=\mathcal{S}(\mathbb{R}^3;\Lambda^2)$ is the \emph{real} Schwartz space and $\mathcal{S}'$ its dual;
$X(f)(\omega):=\omega(f)$;
$\Sigma$ is the cylinder $\sigma$-algebra, and for an open set $O$
\[
\Sigma_O:=\sigma\left(X(f):\ f\in C_c^{\infty}(O;\Lambda^2)\ \text{real}\right).
\]
A law is a probability measure $\mathrm{lw}$ on $(\mathcal{S}',\Sigma)$.
One puts $\langle\mathrm{Ga},\mathrm{Gb}\rangle_{\mathrm{lw}}:=\mathrm{lw}\left(\overline{\mathrm{Ga}}\cdot\mathrm{Gb}\right)$,
and $L^2(\Sigma_O,\mathrm{lw})$ denotes the closed subspace of $L^2(\mathrm{lw})$ formed by the classes that have a $\Sigma_O$-measurable representative.
\item[(n2)] For a Euclidean map $\mathrm{Gi}\colon y\mapsto\mathrm{Rl}\,y+a$ with $\mathrm{Rl}\in\mathrm{O}(3)$ put
\[
(\mathrm{Gi}\cdot f)(y):=(\wedge^2\mathrm{Rl})\,f\left(\mathrm{Gi}^{-1}y\right),\qquad(\mathrm{Gi}\,\omega)(f):=\omega\left(\mathrm{Gi}^{-1}\cdot f\right),\qquad\left(\widehat{\mathrm{Gi}}\,\mathrm{Ga}\right)(\omega):=\mathrm{Ga}\left(\mathrm{Gi}^{-1}\omega\right),
\]
and $U(a):=\widehat{\tau_a}$.
Let $r_{\hat{n},\mathrm{po}}$ be the reflection in the plane $\{\hat{n}\cdot y=\mathrm{po}\}$, and $\theta_j:=r_{e_j,0}$.
For a law $\mathrm{lw}$ the image law is
\[
(\mathrm{Gi}_{*}\mathrm{lw})(\mathrm{Ev}):=\mathrm{lw}\left(\left\{\omega:\mathrm{Gi}\,\omega\in\mathrm{Ev}\right\}\right),\qquad\mathrm{Ev}\in\Sigma;
\]
here $X(f)(\mathrm{Gi}\,\omega)=X\left(\mathrm{Gi}^{-1}\cdot f\right)(\omega)$ for $f\in\mathcal{S}$, and the $X(f)$ generate $\Sigma$, so that $\omega\mapsto\mathrm{Gi}\,\omega$ is $\Sigma$-measurable.
The law $\mathrm{lw}$ is \emph{invariant} under $\mathrm{Gi}$ if $\mathrm{Gi}_{*}\mathrm{lw}=\mathrm{lw}$.
It is invariant under $\mathbb{R}^3$ (under $\mathbb{R}e_j$) if it is invariant under $\tau_a\colon y\mapsto y+a$ for all $a\in\mathbb{R}^3$ (for all $a\in\mathbb{R}e_j$);
it is invariant under $O_{\mathrm{cub}}$ (under $\{\pm1\}^3$, under $\mathrm{O}(3)$) if it is invariant under $y\mapsto\mathrm{Rl}\,y$ for every $\mathrm{Rl}$ in this group;
and it is invariant under $\omega\mapsto-\omega$ if $\mathrm{lw}\left(\left\{\omega:-\omega\in\mathrm{Ev}\right\}\right)=\mathrm{lw}(\mathrm{Ev})$ for all $\mathrm{Ev}\in\Sigma$.
\item[(n3)] (G): $\mathrm{lw}\left(\exp\left(X(f)\right)\right)\le\exp\left(\frac{1}{2}\|f\|_2^2\right)$ for every real $f\in\mathcal{S}$.
(dF): $\mathrm{lw}\left(\left\{dF=0\text{ in }\mathcal{S}'\right\}\right)=1$.
The law $\mathrm{lw}$ is \emph{reflection positive} at $\{\hat{n}\cdot y=\mathrm{po}\}$ if
\[
\sum_{p,q}\overline{c_p}\cdot c_q\cdot\mathrm{lw}\left(\exp\left(i\left(X(f_q)-X\left(r_{\hat{n},\mathrm{po}}\cdot f_p\right)\right)\right)\right)\ge0
\]
for all finite families of numbers $c_p\in\mathbb{C}$ and of real $f_p\in C_c^{\infty}\left(\{\hat{n}\cdot y>\mathrm{po}\};\Lambda^2\right)$.
(H): the law has (G), is invariant under $\mathbb{R}^3$ and under $O_{\mathrm{cub}}$, and is reflection positive at $\{y_j=0\}$ for $j=1,2,3$.
$\mathcal{C}_9$ is the class of the laws that are invariant under $\mathbb{R}^3$ and that, for each plane of the nine families
\[
\{y_i=\mathrm{po}\}\quad(i=1,2,3),\qquad\{y_i-y_l=\mathrm{po}\},\quad\{y_i+y_l=\mathrm{po}\}\quad(1\le i<l\le3),
\]
in each of which $\mathrm{po}$ ranges over $\mathbb{R}$, are invariant under the reflection in this plane and reflection positive at it.
The law $\mathrm{lw}$ is \emph{$\mathbb{R}^3$-ergodic} if it is invariant under $\mathbb{R}^3$ and every $\mathrm{Ga}\in L^2(\mathrm{lw})$ with $U(a)\mathrm{Ga}=\mathrm{Ga}$ for all $a$ is constant.
$P^{123}_{\mathrm{lw}}$ and $P^{j}_{\mathrm{lw}}$ denote the orthogonal projections of $L^2(\mathrm{lw})$ onto the vectors that are invariant under all $U(a)$ with $a\in\mathbb{R}^3$, respectively with $a\in\mathbb{R}e_j$.
\item[(n4)] For a law that is invariant under $\mathbb{R}^3$ and under $\theta_j$ and reflection positive at $\{y_j=0\}$,
\[
\left(\mathcal{H}_{j},\ [\,\cdot\,]_j,\ \mathsf{H}_{j},\ V_{j}(y)\ (y\perp e_j),\ \Omega_{j}=[1]_j\right)
\]
is the reconstruction of \cite[Proposition~8.19]{QEDI};
for $i\ne j$, $\mathsf{P}_{i}$ is the self-adjoint generator of $u\mapsto V_{j}(u\cdot e_i)$.
\item[(n5)] $\mathrm{Sw}_n:=\mathcal{S}\left(\mathbb{R}^{3n};(\Lambda^2)^{\otimes n}\right)$, complex valued, and $\mathrm{Sw}_0:=\mathbb{C}$;
\[
q_m(\Psi):=\max_{|\alpha|\le m}\sup_{y}\left(1+|y|_2^2\right)^{\frac{m}{2}}\cdot\left|\partial^{\alpha}\Psi(y)\right|,
\]
the maximum being formed also over the components.
For a law with (G), $X$ extends to $L^2(\mathbb{R}^3;\Lambda^2)$, and $X^{(n)}\colon\mathrm{Sw}_n\to L^2(\mathrm{lw})$ is the continuous linear map with $X^{(n)}(f_1\otimes\dots\otimes f_n)=X(f_1)\cdots X(f_n)$ (Proposition~45.5);
\[
S_n(\Psi):=\mathrm{lw}\left(X^{(n)}(\Psi)\right),\qquad S_0:=1,\qquad S_n(f_1,\dots,f_n):=\mathrm{lw}\left(X(f_1)\cdots X(f_n)\right).
\]
The joint cumulants are
\[
\kappa_n(f_1,\dots,f_n):=\sum_{\mathrm{pn}}(-1)^{\#|\mathrm{pn}|-1}\cdot\left(\#|\mathrm{pn}|-1\right)!\cdot\prod_{\mathrm{bk}\in\mathrm{pn}}\mathrm{lw}\left(\prod_{a\in\mathrm{bk}}X(f_a)\right),
\]
where the sum extends over the partitions $\mathrm{pn}$ of $\{1,\dots,n\}$ into non-empty blocks $\mathrm{bk}$, and $\#|\mathrm{pn}|$ is the number of the blocks;
for a law with (G) this is the derivative $\partial_{z_1}\cdots\partial_{z_n}\log\mathrm{lw}\left(\exp\left(\sum_az_a\cdot X(f_a)\right)\right)$ at $z=0$, the number that is bounded in \cite[Proposition~7.10(7)]{QEDI}.
Further $\kappa_4(f'):=\kappa_4(f',f',f',f')$,
\[
\mathrm{Ext}_{x_0}(f'):=\kappa_4(f')-3\kappa_4\left(f',f',\tau_{x_0}f',\tau_{x_0}f'\right),\qquad\Theta^{\hat{n}}_m\Psi:=\overline{\left(r_{\hat{n},0}\,\cdot\right)^{\otimes m}\Psi}.
\]
An upper index $\mu$ or $\xi$ names the law with which a quantity is formed.
$\Pi_{\mathrm{c}}$ is the projection of $L^2(\mathbb{R}^3;\Lambda^2)$ that acts on $\widehat{\ast f}$ by the multiplier
\[
\Pi(p):=1-\frac{p\cdot p^{T}}{|p|_2^2},\qquad\text{where }\ast f:=(f_{23},-f_{13},f_{12}).
\]
\item[(n6)] A law is a \emph{mixture} of laws with a given property if it is equal to $\int\mathrm{d}\,\mathrm{mx}(\mathfrak{l})\,\gamma_{\mathfrak{l}}$ for a probability space $(\boldsymbol{\Lambda},\mathrm{mx})$ and laws $\gamma_{\mathfrak{l}}$, $\mathfrak{l}\in\boldsymbol{\Lambda}$, which all have the property, the map $\mathfrak{l}\mapsto\gamma_{\mathfrak{l}}(\mathrm{Ev})$ being measurable for every $\mathrm{Ev}\in\Sigma$ \cite[Section~21.1.8]{QEDI}.
A law is \emph{Gaussian} if every $X(f')$, $f'\in\mathcal{S}$, is a normal random variable, possibly degenerate \cite[Section~9.3]{QEDI}.
\item[(n7)] Let $\mu$ be a law that is invariant under $\mathbb{R}^3$.
The set of the laws carries the $\sigma$-algebra generated by the maps $\mathrm{lw}\mapsto\mathrm{lw}(\mathrm{Ev})$, $\mathrm{Ev}\in\Sigma$.
Then $\nu$ denotes the unique probability measure on this set that is carried by the $\mathbb{R}^3$-ergodic laws and satisfies
\[
\mu=\int\mathrm{d}\nu(\mathrm{lw})\,\mathrm{lw};
\]
its points are written $\mu_{\xi}$ (\cite[Theorem~8.1(a)]{QEDI} with the index set $\{1,2,3\}$).
The words ``for $\nu$-almost every $\xi$'' mean: for all $\xi$ in a set of $\nu$-measure $1$ that belongs to this $\sigma$-algebra; the measure $\nu$ is not completed.
\item[(n8)] $L^2(\mathbb{R}^3;\Lambda^2)$ is the \emph{real} Hilbert space, $\Lambda^2$ being real \cite[Section~8.1.1]{QEDI};
``real $f\in L^2$'' means $f\in L^2(\mathbb{R}^3;\Lambda^2)$;
the complexification is written $L^2(\mathbb{R}^3;\Lambda^2)\otimes\mathbb{C}$, and $X$ is extended to it complex linearly.
For $\mathrm{Rl}\in\mathrm{O}(3)$, $\wedge^2\mathrm{Rl}$ is the linear map of $\Lambda^2$ with $(\wedge^2\mathrm{Rl})(\mathrm{va}\wedge\mathrm{vb})=\mathrm{Rl}\,\mathrm{va}\wedge\mathrm{Rl}\,\mathrm{vb}$ for $\mathrm{va},\mathrm{vb}\in\mathbb{R}^3$.
One puts
\[
\left\{dF=0\text{ in }\mathcal{S}'\right\}:=\left\{\omega\in\mathcal{S}':\ \omega\left(d^{*}\mathrm{gs}\right)=0\text{ for every real }\mathrm{gs}\in\mathcal{S}(\mathbb{R}^3)\right\},
\]
with $d^{*}\mathrm{gs}$ as defined at the beginning of the Propositions~45.m below;
this set belongs to $\Sigma$ \cite[Section~8.1.6]{QEDI}.
With the Hermite basis $(h_{\kappa})$ and the eigenvalues $\lambda_{\kappa}$ of the Propositions~45.m below, put
\[
\|\omega\|_s^2:=\sum_{\kappa}\lambda_{\kappa}^{2s}\cdot\omega(h_{\kappa})^2,\qquad H_{s}:=\left\{\omega\in\mathcal{S}':\|\omega\|_s<\infty\right\},
\]
as in \cite[Section~7.1.5]{QEDI}, applied to the real dual $\mathcal{S}'$.
For every real $s$ the set $H_{s}$ belongs to $\Sigma$, and the trace of $\Sigma$ on $H_{s}$ is the Borel $\sigma$-algebra of the Hilbert space $H_{s}$ \cite[Section~7.1.5]{QEDI};
a Borel probability measure $\mu$ on $H_{s}$ is regarded as the law $\mathrm{Ev}\mapsto\mu(\mathrm{Ev}\cap H_{s})$ on $(\mathcal{S}',\Sigma)$.
In this way the limit points, the elements of $\mathcal{V}_N$ and the volume-first limit points of Section~I.7.5 ($s=-4$) are laws.
For a Euclidean map $\mathrm{Gi}$ with linear part $\mathrm{Rl}$,
\[
\left(\left(\mathrm{Gi}\,\cdot\right)^{\otimes n}\Psi\right)(y_1,\dots,y_n):=(\wedge^2\mathrm{Rl})^{\otimes n}\,\Psi\left(\mathrm{Gi}^{-1}y_1,\dots,\mathrm{Gi}^{-1}y_n\right),\qquad\tau_b^{\otimes n}:=\left(\tau_b\,\cdot\right)^{\otimes n}.
\]
For $n\ge1$ and $\Psi\in\mathrm{Sw}_n$, $Q_{\Psi}$ is the function on $\mathcal{S}'$ defined by the series of Proposition~45.2\,(2).
For $n=0$ one puts, as in \cite[Section~7.6]{QEDI},
\[
Q_{\Psi_0}:=\Psi_0,\qquad S_0(\Psi_0):=\Psi_0,\qquad\Theta^{\hat{n}}_0\Psi_0:=\overline{\Psi_0},\qquad\Psi_0\otimes\Psi:=\Psi_0\cdot\Psi,
\]
so that $Q_{\Psi_0}$ is the constant function $\Psi_0$ on $\mathcal{S}'$;
in addition $\Psi\otimes\Psi_0:=\Psi_0\cdot\Psi$ and, for every Euclidean map $\mathrm{Gi}$, $\left(\mathrm{Gi}\,\cdot\right)^{\otimes0}\Psi_0:=\Psi_0$, so that $\tau_b^{\otimes0}\Psi_0=\Psi_0$.
The cumulant $\kappa_n(f_1,\dots,f_n)$ is given by the display of (n5);
this display and the identity stated after it hold also for $f_i\in L^2(\mathbb{R}^3;\Lambda^2)\otimes\mathbb{C}$ (Chapter~45-1, Part~F, Lemma~F.1).
In the Propositions~45.m and in Theorem~44 the letter $w$ denotes an element of $O_{\mathrm{cub}}$, as in \cite[Section~8.1.5]{QEDI}, and as $r$ does in (L3) and in Proposition~45.1\,(6);
for $w\in O_{\mathrm{cub}}$, $w_{*}\nu$ is the image of $\nu$ under the measurable map $\mathrm{lw}'\mapsto w_{*}\mathrm{lw}'$ of the set of the laws.
On the reconstruction space one has $\left\langle[\mathrm{Ga}]_j,[\mathrm{Gb}]_j\right\rangle:=\left\langle\widehat{\theta_j}\,\mathrm{Ga},\mathrm{Gb}\right\rangle_{\mathrm{lw}}$, which is antilinear in the first argument.
For a self-adjoint operator $H\ge0$ and $t\ge0$, $\exp(-tH)$ is the operator given by the functional calculus for the bounded function $\lambda\mapsto\exp\left(-t\max(\lambda,0)\right)$.
\end{itemize}

\subsection*{45. Statements on single laws}

The notation is that of (n1)--(n8).
For a real $\mathrm{gs}\in\mathcal{S}(\mathbb{R}^3)$, $d^{*}\mathrm{gs}$ is the two-form with $\ast\left(d^{*}\mathrm{gs}\right)=-\operatorname{grad}\mathrm{gs}$;
for a complex $\mathrm{gs}=\mathrm{gs}'+i\,\mathrm{gs}''$ with real $\mathrm{gs}',\mathrm{gs}''\in\mathcal{S}(\mathbb{R}^3)$ one puts $d^{*}\mathrm{gs}:=d^{*}\mathrm{gs}'+i\,d^{*}\mathrm{gs}''$.
For real $f\in\mathcal{S}$ let $\Phi_{\mathrm{lw}}(f):=\mathrm{lw}\left(\exp\left(i\,X(f)\right)\right)$.
The family $(h_{\kappa})_{\kappa}$ is the Hermite basis of $L^2(\mathbb{R}^3;\Lambda^2)$, and the $\lambda_{\kappa}$ are its eigenvalues \cite[Section~7.1.5]{QEDI};
$\mathrm{ZH}_s:=\sum_{\kappa}\lambda_{\kappa}^{-s}$ is the number $\mathsf{Z}_s$ of \cite[Section~7.1.5]{QEDI};
and
\[
\sigma_n:=2\cdot2^{\frac{n}{2}}\cdot\Gamma\left(\frac{n}{2}+1\right)\le2\left(2^n\cdot n!\right)^{\frac{1}{2}}
\]
are the numbers of \cite[Proposition~7.10]{QEDI}.
For a law $\mathrm{lw}$ the following four conditions are considered.
\begin{itemize}
\item[(L1)] $\mathrm{lw}\left(\exp\left(X(f)\right)\right)\le\exp\left(\frac{1}{2}\|f\|_2^2\right)$ for real $f\in\mathcal{S}$; this is (G).
\item[(L2)] $X\left(d^{*}\mathrm{gs}\right)=0$ $\mathrm{lw}$-almost surely, for every real $\mathrm{gs}\in C_c^{\infty}(\mathbb{R}^3)$.
\item[(L3)] $\Phi_{\mathrm{lw}}(\tau_b\cdot f)=\Phi_{\mathrm{lw}}(f)=\Phi_{\mathrm{lw}}(r\cdot f)$ for real $f\in\mathcal{S}$, $b\in\mathbb{R}^3$ and $r\in O_{\mathrm{cub}}$.
\item[(L4)] For $j=1,2,3$, for $(c_1,\dots,c_s)\in\mathbb{C}^s$ and for real $f_1,\dots,f_s\in C_c^{\infty}\left(\{y_j>0\};\Lambda^2\right)$,
\[
\sum_{p,q}\overline{c_p}\cdot c_q\cdot\Phi_{\mathrm{lw}}\left(f_q-\theta_j\cdot f_p\right)\ge0;
\]
this is reflection positivity at $\{y_j=0\}$.
\end{itemize}

\subsubsection*{Proposition~45.6. The ergodic decomposition}

\begin{named}{Proposition 45.6}
Let $\mathrm{lw}$ be a law that is invariant under $\mathbb{R}^3$, and let $I$ be the $\sigma$-algebra of the events $\mathrm{Ev}\in\Sigma$ with $\tau_a\mathrm{Ev}=\mathrm{Ev}$ for all $a\in\mathbb{R}^3$.
\begin{itemize}
\item[(a)] There is a probability kernel $(\omega,\mathrm{Ev})\mapsto\mathrm{lw}_{\omega}(\mathrm{Ev})$ from $(\mathcal{S}',I)$ to $(\mathcal{S}',\Sigma)$ such that, for each $\mathrm{Ev}\in\Sigma$,
\[
\mathrm{lw}_{\omega}(\mathrm{Ev})=E_{\mathrm{lw}}\left[\mathbf{1}_{\mathrm{Ev}}\mid I\right](\omega)\qquad\text{for }\mathrm{lw}\text{-almost every }\omega.
\]
Two such kernels agree for $\mathrm{lw}$-almost every $\omega$.
\item[(a$'$)] Let $\Phi\ge0$ be measurable, or let $\Phi\in L^1(\mathrm{lw})$.
Then $\omega\mapsto\mathrm{lw}_{\omega}(\Phi)$ is $I$-measurable and defined for $\mathrm{lw}$-almost every $\omega$,
one has $\mathrm{lw}_{\omega}(\Phi)=E_{\mathrm{lw}}\left[\Phi\mid I\right](\omega)$ for $\mathrm{lw}$-almost every $\omega$, and
\[
\mathrm{lw}(\Phi)=\int\mathrm{d}\,\mathrm{lw}(\omega)\,\mathrm{lw}_{\omega}(\Phi).
\]
For a measurable $\Phi\ge0$ that is not $\mathrm{lw}$-integrable, $E_{\mathrm{lw}}\left[\Phi\mid I\right]$ is the conditional expectation with values in $[0,\infty]$,
that is, an $I$-measurable function $\mathrm{Gu}\colon\mathcal{S}'\to[0,\infty]$ with
\[
\int_{\mathrm{Ie}}\mathrm{d}\,\mathrm{lw}\,\mathrm{Gu}=\int_{\mathrm{Ie}}\mathrm{d}\,\mathrm{lw}\,\Phi\qquad\text{for all }\mathrm{Ie}\in I;
\]
it is unique up to $\mathrm{lw}$-null sets.
\item[(a$''$)] The image $\nu$ of $\mathrm{lw}$ under $\omega\mapsto\mathrm{lw}_{\omega}$ is the only probability measure on the laws that is carried by the $\mathbb{R}^3$-ergodic laws and represents $\mathrm{lw}$.
For a property that defines a measurable set of laws, ``for $\nu$-almost every $\xi$'' and ``for $\mathrm{lw}$-almost every $\omega$, for $\mathrm{lw}_{\omega}$'' mean the same.
If $\mathrm{lw}$ is $\mathbb{R}^3$-ergodic, then $\nu$ is the point mass at $\mathrm{lw}$.
\item[(b)] If $\mathrm{lw}$ is invariant under $\theta_j$ and reflection positive at $\{y_j=0\}$, then $\nu$-almost every $\mathrm{lw}_{\xi}$ is invariant under $\theta_j$ and reflection positive at every plane $\{y_j=\mathrm{po}\}$.
\item[(c)] If $\mathrm{lw}$ has (G), then $\nu$-almost every $\mathrm{lw}_{\xi}$ has (G), for all real $f\in\mathcal{S}$ simultaneously.
\item[(c$'$)] If $\mathrm{Ev}\in\Sigma$ and $\mathrm{lw}(\mathrm{Ev})=1$, then $\mathrm{lw}_{\xi}(\mathrm{Ev})=1$ for $\nu$-almost every $\xi$.
\item[(d0)] For every $\Sigma$-measurable $\Phi\colon\mathcal{S}'\to[0,\infty]$ the map $\xi\mapsto\mathrm{lw}_{\xi}(\Phi)=\xi(\Phi)$, from the laws to $[0,\infty]$, is measurable for the $\sigma$-algebra of (n7).
For every $\Sigma$-measurable complex function $\Phi$ the set $\left\{\xi:\xi\left(|\Phi|\right)<\infty\right\}$ belongs to this $\sigma$-algebra, and $\xi\mapsto\xi(\Phi)$ is measurable on it.
Every set $\{\xi_0\}$ consisting of one point belongs to it.
\item[(d)] If $\mathrm{Ga}\in L^2(\mathrm{lw})$ is real and $\left|E_{\mathrm{lw}}\left[\mathrm{Ga}\mid I\right]\right|\le\mathrm{cf}$ holds $\mathrm{lw}$-almost surely, then $\left|\mathrm{lw}_{\xi}(\mathrm{Ga})\right|\le\mathrm{cf}$ for $\nu$-almost every $\xi$.
Here $\mathrm{lw}_{\xi}(\mathrm{Ga})$ is formed with a real $\Sigma$-measurable representative of the class $\mathrm{Ga}$;
this representative is $\mathrm{lw}_{\xi}$-integrable for $\nu$-almost every $\xi$, and two representatives give the same number for $\nu$-almost every $\xi$.
\end{itemize}
\end{named}

\begin{named}{Remark}
Proposition~45.6 corresponds to \cite[Theorem~8.1]{QEDI} for the index set $\{1,2,3\}$, together with the identity between the kernel and the conditional expectation, for non-negative measurable and for integrable functions, that is displayed in its proof in \cite[Section~8.5]{QEDI};
to \cite[Proposition~9.7(b)]{QEDI} in the form in which it is applied in Step~6 of the proof in \cite[Section~21.10.1]{QEDI};
and, through its clause (c), to the first assertion of \cite[Corollary~9.10]{QEDI}.
The identity between the kernel and the conditional expectation is the first numbered display of the proof of \cite[Theorem~8.1]{QEDI} in \cite[Section~8.5]{QEDI}, obtained there by the monotone class theorem, and \cite[Theorem~8.1(a)]{QEDI} is its case of indicator functions.
Here that display is clause (a$'$), and its case of indicator functions is the first sentence of clause (a);
the last sentence of (a) adds that two kernels with the property of (a) agree for $\mathrm{lw}$-almost every $\omega$.
Clause (a$''$) contains the assertions of \cite[Theorem~8.1(a)]{QEDI} on $\nu$, among them the uniqueness of $\nu$, which is the last sentence of \cite[Theorem~8.1(a)]{QEDI};
it contains in addition the equivalence of ``for $\nu$-almost every $\xi$'' and ``for $\mathrm{lw}$-almost every $\omega$, for $\mathrm{lw}_{\omega}$'', which is derived in the same proof from the formula $\nu(B)=\mathrm{lw}\left(\left\{\omega:\mathrm{lw}_{\omega}\in B\right\}\right)$ for the measurable sets $B$ of laws.
Clause (d) is the step by which a bound on many points reaches the components;
it requires the measurability of $\xi\mapsto\mathrm{lw}_{\xi}(\mathrm{Ga})$ on the set of the laws, which is (d0);
the measurability of $\omega\mapsto\mathrm{lw}_{\omega}(\Phi)$ with respect to $I$ is contained in (a$'$).
The last sentence of (a$''$) is the fact used for (E4) and (E8) in Part~II of \cite[Theorem~1.4]{QEDII}.
The $\sigma$-algebra $I$ consists of the strictly invariant events, as in \cite[Section~8.1.4]{QEDI};
\cite[Lemma~9.4]{QEDI} uses instead the events that are invariant up to null sets of the law.
The two conditional expectations are the same element of $L^2(\mathrm{lw})$:
each of them is the orthogonal projection onto $\left\{\mathrm{Ga}:U(a)\mathrm{Ga}=\mathrm{Ga}\text{ for all }a\right\}$ (\cite[Lemma~8.13(iv)]{QEDI} and \cite[Lemma~9.4(b)]{QEDI}).
\end{named}

\subsubsection*{Proposition 45.1. Regularity, closedness, invariance, reflection positivity, determinacy}

\begin{named}{Proposition 45.1}
Let $\mathrm{lw}$ be a probability measure on $(\mathcal{S}',\Sigma)$ that satisfies (L1)--(L4).
Then the following assertions hold.
\begin{itemize}
\item[(1)] The map $f\mapsto X(f)$, $f\in\mathcal{S}$, has exactly one extension to a bounded linear map $X\colon L^2(\mathbb{R}^3;\Lambda^2)\to L^2(\mathrm{lw})$.
For every $f$ in the real space $L^2(\mathbb{R}^3;\Lambda^2)$ the element $X(f)$ is a class of real functions, and
\[
\mathrm{lw}\left(\exp X(f)\right)\le\exp\left(\frac{1}{2}\|f\|_2^2\right).
\]
All the conclusions (1)--(7) of \cite[Proposition~7.10]{QEDI} hold for the probability space $(\mathcal{S}',\Sigma,\mathrm{lw})$, the space $L^2(\mathbb{R}^3;\Lambda^2)$ with the norm $\|\cdot\|_2$, and this map $X$.
\item[(3)] Let $s>\frac{3}{2}$.
Then
\[
\mathrm{lw}\left(\|\omega\|_{-s}^2\right)\le\mathrm{ZH}_{2s},\qquad
\mathrm{lw}\left(\exp\left(t\cdot\|\omega\|_{-s}^2\right)\right)\le\left(1-2t\cdot\mathrm{ZH}_{2s}\right)^{-\frac{1}{2}}\quad\text{for }0\le2t\cdot\mathrm{ZH}_{2s}<1 .
\]
In particular $\mathrm{lw}(H_{-s})=1$.
\item[(4)] One has $\mathrm{lw}\{dF=0\text{ in }\mathcal{S}'\}=1$, and $X(f)=X(\Pi_{\mathrm{c}}f)$ almost surely for every $f\in L^2(\mathbb{R}^3;\Lambda^2)$.
For every $f\in L^2(\mathbb{R}^3;\Lambda^2)$,
\[
\mathrm{lw}\left(\exp X(f)\right)\le\exp\left(\frac{1}{2}\langle f,\Pi_{\mathrm{c}}\,f\rangle\right),\qquad
0\le\mathrm{lw}\left(X(f)^2\right)\le\langle f,\Pi_{\mathrm{c}}\,f\rangle\le\|f\|_2^2 .
\]
\item[(5)] $(\tau_b)_*\mathrm{lw}=\mathrm{lw}$ for every $b\in\mathbb{R}^3$, and $b\mapsto U(b)$ is a strongly continuous unitary representation of $\mathbb{R}^3$ on $L^2(\mathrm{lw})$.
\item[(6)] $r_*\mathrm{lw}=\mathrm{lw}$ for every $r\in O_{\mathrm{cub}}$.
\item[(7)] Let $j\in\{1,2,3\}$ and $\mathrm{po}\in\mathbb{R}$; then $(r_{e_j,\mathrm{po}})_*\mathrm{lw}=\mathrm{lw}$, and
\[
\left\langle\hat{r}_{e_j,\mathrm{po}}\mathrm{Ga},\mathrm{Ga}\right\rangle_{\mathrm{lw}}\ge0\qquad\text{for every }\mathrm{Ga}\in L^2\left(\Sigma_{\{y_j>\mathrm{po}\}},\mathrm{lw}\right).
\]
For every $f\in L^2(\mathbb{R}^3;\Lambda^2)$ that vanishes almost everywhere on $\{y_j<\mathrm{po}\}$ the class $X(f)$ has a $\Sigma_{\{y_j>\mathrm{po}\}}$-measurable representative, so that the inequality applies to every polynomial and to every exponential in such functions $X(f)$.
\item[(8)] Let $\mathrm{lw}'$ be a probability measure on $(\mathcal{S}',\Sigma)$ with the following property: for every $n$ and all real $f_1,\dots,f_n\in C_c^{\infty}(\mathbb{R}^3;\Lambda^2)$ the product $X(f_1)\cdots X(f_n)$ is $\mathrm{lw}'$-integrable, and
\[
\mathrm{lw}'\left(X(f_1)\cdots X(f_n)\right)=\mathrm{lw}\left(X(f_1)\cdots X(f_n)\right).
\]
Then $\mathrm{lw}'=\mathrm{lw}$.
\item[(9)] Suppose in addition that $\Phi_{\mathrm{lw}}(-f)=\Phi_{\mathrm{lw}}(f)$ for every real $f\in\mathcal{S}$.
Then $\mathrm{lw}$ is invariant under $\omega\mapsto-\omega$, and $S_n=0$ for every odd $n$.
\item[(10)] (Converse statements.)
Let $\mathrm{lw}$ be an arbitrary law; (L1)--(L4) are not assumed here.
The three conditions (L1), (L3), (L4) together form the hypothesis (H).
The condition (L3) is equivalent to the invariance of $\mathrm{lw}$ under $\mathbb{R}^3$ and under $O_{\mathrm{cub}}$, and (L2) is equivalent to (dF).
Hence every law with (H) and (dF) satisfies (L1)--(L4).
\item[(11)] (Dependence on the hypotheses.)
Let $\mathrm{lw}$ be an arbitrary law.
The assertions (1), (3) and (8) hold under (L1); (4) holds under (L1) and (L2); (5) under (L1) and the first equality in (L3); (6) under the second equality in (L3); (7) under (L1), (L4), the first equality in (L3) for $b\in\mathbb{R}e_j$ and the second equality in (L3) for $r=\theta_j$.
The second sentence of (7), in the form ``$X(f)$ has a $\Sigma_{\{y_j>\mathrm{po}\}}$-measurable representative, and every polynomial and every exponential in such functions $X(f)$ lies in $L^2(\Sigma_{\{y_j>\mathrm{po}\}},\mathrm{lw})$'', holds under (L1) alone.
Finally, (9) holds under (L1) and $\Phi_{\mathrm{lw}}(-f)=\Phi_{\mathrm{lw}}(f)$.
\end{itemize}
\end{named}

\begin{named}{Remark}
Proposition~45.1 corresponds to the assertions (1) and (3)--(9) of \cite[Theorem~7.15]{QEDI}, which are made there for the limit points of admissible families, together with the part of \cite[Proposition~7.19]{QEDI} that concerns them.
The assertion (2) of \cite[Theorem~7.15]{QEDI} is a statement on the defining sequence of a limit point and has no counterpart here.
For a limit point of the laws $P_{N,T}$ of Section~I.7.5, this assertion and the conditions (L1)--(L4) are established in the proof of \cite[Proposition~12.12]{QEDII}, the former in its step (4), the latter in its steps (3), (5), (6) and (7); that proposition applies to the family $(P_{N,T})$ by \cite[Proposition~12.10]{QEDII}.
The mathematical differences are the following.
The hypothesis bears on one law with (L1)--(L4); neither a defining sequence nor the space $H_{-4}$ enters.
In \cite[Proposition~7.19]{QEDI} the corresponding assertions for a law with (L1)--(L4) are derived from the proofs of \cite[Theorem~7.15]{QEDI}.
In (7) the invariance of $\mathrm{lw}$ under $r_{e_j,\mathrm{po}}$ for $\mathrm{po}\ne0$ is part of the assertion, since it is needed for $\hat{r}_{e_j,\mathrm{po}}$ to act on classes; it follows from (5), (6) and the identity $r_{e_j,\mathrm{po}}=\tau_{\mathrm{po}\cdot e_j}\,\theta_j\,\tau_{-\mathrm{po}\cdot e_j}$.
In (8) the integrability of the products under $\mathrm{lw}'$ is part of the hypothesis.
The reflected form in (7) is written with the linear map $\hat{r}$ and the sesquilinear form $\langle\,\cdot\,,\cdot\,\rangle_{\mathrm{lw}}$, as in \cite[Section~8.1.2]{QEDI}, whereas \cite[Theorem~7.15(7)]{QEDI} writes $\int\mathrm{d}\mu\,(\Theta_{\nu,c}A)\cdot A$ with the antilinear map $\Theta_{\nu,c}$; the two expressions are equal (Section~I.4.8).
\end{named}

\subsubsection*{Proposition 45.2. Schwinger functions}

\begin{named}{Proposition 45.2}
Let $\mathrm{lw}$ be as in Proposition~45.1, and put
\[
S_n(f_1,\dots,f_n):=\mathrm{lw}\left(X(f_1)\cdots X(f_n)\right)\qquad\left(f_1,\dots,f_n\in L^2(\mathbb{R}^3;\Lambda^2)\otimes\mathbb{C}\right).
\]
\begin{itemize}
\item[(1)] One has $S_0=1$ and $S_1=0$; every $S_n$ is multilinear and symmetric, and it is real at real arguments.
Moreover
\begin{align*}
\left|S_n(f_1,\dots,f_n)\right|&\le2\cdot\left(2^n\cdot n!\right)^{\frac{1}{2}}\cdot\prod_{i=1}^n\|f_i\|_2\qquad\text{for real }f_1,\dots,f_n,\\
\left|S_n(f_1,\dots,f_n)\right|&\le2\cdot\left(4^n\cdot n!\right)^{\frac{1}{2}}\cdot\prod_{i=1}^n\|f_i\|_2\qquad\text{for complex }f_1,\dots,f_n.
\end{align*}
\item[(2)] For an $n$-tuple $\mathrm{kk}=(\kappa_1,\dots,\kappa_n)$ of indices of the Hermite basis $(h_\kappa)_\kappa$ let $h_{\mathrm{kk}}:=h_{\kappa_1}\otimes\dots\otimes h_{\kappa_n}$, and for $\Psi\in\mathrm{Sw}_n$ let $c_{\mathrm{kk}}(\Psi):=\langle h_{\mathrm{kk}},\Psi\rangle_{L^2}$.
For every $\omega\in\mathcal{S}'$ the series, extended over all such $n$-tuples,
\[
Q_\Psi(\omega):=\sum_{\mathrm{kk}}c_{\mathrm{kk}}(\Psi)\cdot\prod_{i=1}^n\omega(h_{\kappa_i})
\]
converges absolutely, and $\Psi\mapsto Q_\Psi(\omega)$ is the only continuous linear functional on $\mathrm{Sw}_n$ whose value at $f_1\otimes\dots\otimes f_n$ is $\prod_{i=1}^n\omega(f_i)$.
One has
\[
\|Q_\Psi\|_{L^2(\mathrm{lw})}\le\sigma_{2n}^{\frac{1}{2}}\cdot\sum_{\mathrm{kk}}\left|c_{\mathrm{kk}}(\Psi)\right| .
\]
The map $\Psi\mapsto S_n(\Psi):=\mathrm{lw}(Q_\Psi)$ is the only continuous linear functional on $\mathrm{Sw}_n$ with $S_n(f_1\otimes\dots\otimes f_n)=S_n(f_1,\dots,f_n)$, and
\[
\left|S_n(\Psi)\right|\le2\cdot66^n\cdot(n!)^{\frac{1}{2}}\cdot q_{6n}(\Psi)\qquad(\Psi\in\mathrm{Sw}_n).
\]
The map that assigns to $\Psi$ the class of $Q_\Psi$ is the only continuous linear map $\mathrm{Sw}_n\to L^2(\mathrm{lw})$ whose value at $f_1\otimes\dots\otimes f_n$ is $X(f_1)\cdots X(f_n)$.
It is therefore the map $X^{(n)}$ of (n5), and $Q_\Psi=X^{(n)}(\Psi)$ in $L^2(\mathrm{lw})$.
\item[(4)] For $b\in\mathbb{R}^3$ and $r\in O_{\mathrm{cub}}$,
\[
S_n\left(\tau_b^{\otimes n}\Psi\right)=S_n(\Psi)=S_n\left((r\cdot)^{\otimes n}\Psi\right).
\]
Moreover $S_n\left(d^{*}\mathrm{gs}\otimes\Psi'\right)=0$ for every $\mathrm{gs}\in\mathcal{S}(\mathbb{R}^3)$, real or complex, and every $\Psi'\in\mathrm{Sw}_{n-1}$.
\item[(5)] Let $j\in\{1,2,3\}$.
Let $\Psi_0\in\mathbb{C}$ and $\Psi_n\in\mathrm{Sw}_n$, $n\ge1$, be given, only finitely many of them different from $0$, such that every $\Psi_n$ is supported in $\{(y_a)_j\ge0\text{ for all }a\}$.
Then
\[
\sum_{m,n}S_{m+n}\left(\Theta^{e_j}_m\Psi_m\otimes\Psi_n\right)\ge0 .
\]
\item[(6)] (Rules, valid at every $\omega\in\mathcal{S}'$.)
One has $Q_f(\omega)=\omega(f)$ for $f\in\mathrm{Sw}_1$, and $Q_{\Psi\otimes\Phi}(\omega)=Q_\Psi(\omega)\cdot Q_\Phi(\omega)$ for $\Psi\in\mathrm{Sw}_m$, $\Phi\in\mathrm{Sw}_n$, $m,n\ge0$.
Further $\overline{Q_\Psi(\omega)}=Q_{\overline{\Psi}}(\omega)$; $Q_{(\mathrm{Gi}\cdot)^{\otimes n}\Psi}=\widehat{\mathrm{Gi}}\,Q_\Psi$ for every Euclidean map $\mathrm{Gi}$; $Q_\Psi(-\omega)=(-1)^n\cdot Q_\Psi(\omega)$; and $\overline{\hat{\theta}_jQ_\Psi}=Q_{\Theta^{e_j}_n\Psi}$.
If $\Phi_{\mathrm{lw}}(-f)=\Phi_{\mathrm{lw}}(f)$ for every real $f\in\mathcal{S}$, then $S_n(\Psi)=0$ for every odd $n$ and every $\Psi\in\mathrm{Sw}_n$.
\item[(7)] (Support.)
If $\Psi\in\mathrm{Sw}_n$ is supported in $\{(y_a)_j\ge0\text{ for all }a\}$, then the class of $Q_\Psi$ lies in $L^2(\Sigma_{\{y_j>0\}},\mathrm{lw})$.
\item[(8)] (Dependence on the hypotheses.)
Let $\mathrm{lw}$ be an arbitrary law.
The assertions (1), (2) and (7) hold under (L1).
The rules of (6) involve no law, and the last sentence of (6) holds under (L1).
The assertion (4) holds under the hypotheses that Proposition~45.1\,(4), (5), (6) require, and (5) under those that Proposition~45.1\,(7) requires; these hypotheses are listed in Proposition~45.1\,(11).
\end{itemize}
\end{named}

\begin{named}{Remark}
Proposition~45.2 corresponds to the assertions (1), (2), (4) and (5) of \cite[Theorem~7.17]{QEDI}, made there for the limit points of admissible families, together with the corresponding part of \cite[Proposition~7.19]{QEDI}.
The assertion (3) of \cite[Theorem~7.17]{QEDI} is a statement on limit points: it concerns the lattice laws along the defining sequence of a limit point, and it has no counterpart here.
For a limit point of the laws $P_{N,T}$ of Section~I.7.5 it is contained in \cite[Proposition~12.12(b)]{QEDII}, and it is proved in step (9) of the proof of that proposition.
The hypothesis bears on one law with (L1)--(L4).
The tensor power $\omega^{\otimes n}$ is not presupposed: $Q_\Psi(\omega)$ is defined by the series of (2) and characterised by the property stated there.
The bound on $|S_n(\Psi)|$ in (2) is displayed as an inequality; it is the bound of \cite[Theorem~7.17(2)]{QEDI}, where it is identified with the linear growth condition of Osterwalder and Schrader \cite{OsterwalderSchrader1975} for the index $s=6$, and it is the literal analogue of that condition for tests on $\mathbb{R}^{3n}$ with values in $(\Lambda^2)^{\otimes n}$.
The last sentence of (2) relates $Q_\Psi$ to the map $X^{(n)}$ of Proposition~45.5.
\end{named}

\subsubsection*{Proposition 45.3. Transfer semigroups}

\begin{named}{Proposition 45.3}
Let $\mathrm{lw}$ be as in Proposition~45.1, and let $j\in\{1,2,3\}$.
Put $\mathrm{Ep}_j:=L^2(\Sigma_{\{y_j>0\}},\mathrm{lw})$ and
\[
\mathrm{bf}(\mathrm{Ga},\mathrm{Gb}):=\left\langle\hat{\theta}_j\mathrm{Ga},\mathrm{Gb}\right\rangle_{\mathrm{lw}}\qquad(\mathrm{Ga},\mathrm{Gb}\in\mathrm{Ep}_j).
\]
Let $\mathcal{H}_j$ be the completion of $\mathrm{Ep}_j$ modulo $\{\mathrm{Ga}\in\mathrm{Ep}_j:\mathrm{bf}(\mathrm{Ga},\mathrm{Ga})=0\}$, let $[\mathrm{Ga}]_j$ denote the class of $\mathrm{Ga}$, and put $\Omega_j:=[1]_j$.
\begin{itemize}
\item[(1)] For $t\ge0$ the assignment $[\mathrm{Ga}]_j\mapsto[U(t\cdot e_j)\mathrm{Ga}]_j$ defines a strongly continuous semigroup of self-adjoint contractions $\exp(-t\cdot\mathsf{H}_j)$, where $\mathsf{H}_j=\mathsf{H}_j^{*}\ge0$.
Moreover $\|\Omega_j\|=1$ and $\mathsf{H}_j\Omega_j=0$.
\item[(2)] For $y\in e_j^{\perp}$ the formula $V_j(y)[\mathrm{Ga}]_j:=[U(y)\mathrm{Ga}]_j$ defines a strongly continuous unitary representation of $e_j^{\perp}$, so that $V_j(y)V_j(y')=V_j(y+y')$; it commutes with $\exp(-t\cdot\mathsf{H}_j)$ and fixes $\Omega_j$.
\item[(3)] Let $m,n\ge0$ and $t\ge0$, and let $\Psi\in\mathrm{Sw}_m$ and $\Phi\in\mathrm{Sw}_n$ be supported in $\{(y_a)_j\ge0\text{ for all }a\}$.
Then
\[
S_{m+n}\left(\Theta^{e_j}_m\Psi\otimes\tau_{t\cdot e_j}^{\otimes n}\Phi\right)=\left\langle[Q_\Psi]_j,\exp(-t\cdot\mathsf{H}_j)[Q_\Phi]_j\right\rangle,
\]
and this number converges, as $t\to\infty$, to $\left\langle[Q_\Psi]_j,E_0[Q_\Phi]_j\right\rangle$, where $E_0$ is the orthogonal projection onto $\ker\mathsf{H}_j$.
\item[(4)] $\left(\mathcal{H}_j,[\,\cdot\,]_j,\mathsf{H}_j,V_j,\Omega_j\right)$ is the reconstruction of \cite[Proposition~8.19]{QEDI} for the law $\mathrm{lw}$.
\item[(5)] (Dependence on the hypotheses.)
For the direction $j$, the assertions (1)--(4) hold for every law with (L1), (L4), the first equality in (L3), and the second equality in (L3) for $r=\theta_j$.
\end{itemize}
\end{named}

\begin{named}{Remark}
Proposition~45.3 corresponds to \cite[Corollary~7.18]{QEDI}, stated there for the limit points of admissible families, together with the corresponding part of \cite[Proposition~7.19]{QEDI}.
The hypothesis bears on one law with (L1)--(L4).
The assertion (2) gives a single strongly continuous unitary representation of the plane $e_j^{\perp}$, whereas \cite[Corollary~7.18(2)]{QEDI} gives one unitary group for each coordinate direction orthogonal to $e_j$; the operators $V_j(y)$ for all $y\in e_j^{\perp}$ occur in (P7) of Theorem~44 and in \cite[Definition~2.38]{QEDII}.
The assertion (4) identifies the construction of (1) and (2) with the reconstruction of \cite[Proposition~8.19]{QEDI}.
\end{named}

\subsubsection*{Proposition 45.4. Soft laws}

\begin{named}{Proposition 45.4}
A law $\mathrm{lw}$ is called \emph{soft} if it satisfies (H) and (dF) and is invariant under $\omega\mapsto-\omega$.
Let $\mathrm{lw}$ be soft, and let $\nu$ be the measure of its $\mathbb{R}^3$-ergodic decomposition (n7).
\begin{itemize}
\item[(a)] The law $\mathrm{lw}$ satisfies (L1)--(L4), and $\Phi_{\mathrm{lw}}$ is even; hence Propositions~45.1--45.3 hold for $\mathrm{lw}$.
These are the properties (P1)--(P7) of Theorem~44\,(B1).
\item[(b)] The following six statements hold, each with its own null set:
\begin{itemize}
\item[$\bullet$] for $\nu$-almost every $\xi$ the law $\mathrm{lw}_\xi$ is invariant and ergodic under $\mathbb{R}^3$;
\item[$\bullet$] for $\nu$-almost every $\xi$ the law $\mathrm{lw}_\xi$ has (G);
\item[$\bullet$] for $\nu$-almost every $\xi$ the law $\mathrm{lw}_\xi$ has (dF);
\item[$\bullet$] for each $j\in\{1,2,3\}$: for $\nu$-almost every $\xi$ the law $\mathrm{lw}_\xi$ is invariant under $\theta_j$ and reflection positive at every plane $\{y_j=\mathrm{po}\}$.
\end{itemize}
\item[(c)] Let $\xi$ be any point at which the six properties of (b) hold.
Then $\mathrm{lw}_\xi$ lies in the class $\mathcal{C}$ of \cite[Section~8.1.3]{QEDI}, has (G) and (dF), is $\mathbb{R}^3$-ergodic and is invariant under $\{\pm1\}^3$.
Consequently \cite[Theorem~8.3(i)--(iv)]{QEDI} and \cite[Proposition~8.19]{QEDI} hold for the single law $\mathrm{lw}_\xi$, and Proposition~45.5 applies to it.
\item[(d)] For every fixed event $\mathrm{Ev}\in\Sigma$ with $\mathrm{lw}(\mathrm{Ev})=1$ one has $\mathrm{lw}_\xi(\mathrm{Ev})=1$ for $\nu$-almost every $\xi$, the null set depending on $\mathrm{Ev}$.
Moreover $w_*\nu=\nu$ for every $w\in O_{\mathrm{cub}}$; this is a statement on $\nu$, not on single components.
\end{itemize}
\end{named}

\begin{named}{Remark}
Proposition~45.4 corresponds to \cite[Lemma~21.36]{QEDI}, by which every soft law has the properties (P1)--(P7) and the property (c) of \cite[Theorem~3.1, Part~I]{QEDI}, and to \cite[Theorem~3.1, Part~I\,(c)]{QEDI}.
In the latter the assertions on the components $\mu_\xi$ of a limit point $\mu$ carry quantifiers of three kinds, and so do the properties 1--6 of \cite[Corollary~8.9]{QEDI}.
First, the assertions at the head of Part~I\,(c), namely invariance and ergodicity under $\mathbb{R}^3$, invariance under $\{\pm1\}^3$, (G), (dF) and reflection positivity at every plane $\{y_j=\mathrm{po}\}$, together with the properties 1--4 and the first clause of property~6, by which $w_*\mu_\xi$ has the properties 1--4 and depends on $w\in O_{\mathrm{cub}}$ only through its class modulo $\{\pm1\}^3$, hold for every $\xi$ of one set of full $\nu$-measure.
In \cite[Corollary~8.9]{QEDI} the properties 1--4 and the first clause of property~6 hold for every $\xi$ in the set of the laws of the class $\mathcal{C}$ that have (G) and are $\mathbb{R}^3$-ergodic, a set of full $\nu$-measure.
Secondly, property~5, by which $\mu_\xi(\mathrm{Ev})=1$, and the second clause of property~6, by which $(w_*\mu_\xi)(\mathrm{Ev})=1$, hold for every event $\mathrm{Ev}\in\Sigma$ with $\mu(\mathrm{Ev})=1$ and $\nu$-almost every $\xi$, the null set depending on $\mathrm{Ev}$.
Thirdly, the last clause of property~6 is the identity $w_*\nu=\nu$, a statement on $\nu$.
In Proposition~45.4 the assertions (b)--(d) keep three kinds of statements apart: almost sure statements, each with its own null set, in (b); consequences for a single law, in (c); a statement whose null set depends on the event, and a statement on $\nu$, in (d).
In this respect Proposition~45.4\,(b)--(d) agrees with \cite[Theorem~3.1, Part~I\,(c)]{QEDI} and with \cite[Corollary~8.9]{QEDI}.
Among the differences are the following.
Proposition~45.4 is stated for an arbitrary soft law, whereas Part~I\,(c) of \cite[Theorem~3.1]{QEDI} is stated for the limit points.
Its assertion (b) lists six almost sure statements one by one, each with its own null set, and (c) concerns any point $\xi$ at which the six hold, whereas \cite{QEDI} names one set of full $\nu$-measure; the two forms are equivalent, since a finite union of null sets is a null set.
The assertion (c) obtains the properties of the single law $\mathrm{lw}_\xi$ by applying \cite[Theorem~8.3(i)--(iv)]{QEDI} and \cite[Proposition~8.19]{QEDI} to it.
Of property~6, Proposition~45.4 contains the last clause, as the identity $w_*\nu=\nu$ in (d), and it makes no assertion on the laws $w_*\mathrm{lw}_\xi$, $w\in O_{\mathrm{cub}}$.
The first and the second clause of property~6 follow from (b)--(d), as follows.
Let $w\in O_{\mathrm{cub}}$, and let $\mathrm{Ev}\in\Sigma$ with $\mathrm{lw}(\mathrm{Ev})=1$.
The set $w^{-1}\mathrm{Ev}:=\{\omega:w\,\omega\in\mathrm{Ev}\}$ belongs to $\Sigma$ by (n2), and $\mathrm{lw}\left(w^{-1}\mathrm{Ev}\right)=(w_*\mathrm{lw})(\mathrm{Ev})=1$ by (H).
The first sentence of (d), applied to the event $w^{-1}\mathrm{Ev}$, gives
\[
\left(w_*\mathrm{lw}_\xi\right)(\mathrm{Ev})=\mathrm{lw}_\xi\left(w^{-1}\mathrm{Ev}\right)=1\qquad\text{for $\nu$-almost every }\xi,
\]
the null set depending on $\mathrm{Ev}$; this is the second clause of property~6.
At every point $\xi$ at which the six statements of (b) hold, the law $\mathrm{lw}_\xi$ is invariant under $\{\pm1\}^3$ by (c), so that $(w\varepsilon)_*\mathrm{lw}_\xi=w_*\left(\varepsilon_*\mathrm{lw}_\xi\right)=w_*\mathrm{lw}_\xi$ for $w\in O_{\mathrm{cub}}$ and $\varepsilon\in\{\pm1\}^3$; thus $w_*\mathrm{lw}_\xi$ is determined by the class of $w$ modulo $\{\pm1\}^3$.
For the properties 1--4 of the laws $w_*\mathrm{lw}_\xi$, let $\Xi'$ be a set of $\nu$-measure $1$ in the $\sigma$-algebra of (n7) at every point of which the six statements of (b) hold.
For $w\in O_{\mathrm{cub}}$ the point $w_*\xi$ of the set of the laws is the law $w_*\mathrm{lw}_\xi$, and the identity $w_*\nu=\nu$ states that $\xi\mapsto w_*\xi$ preserves $\nu$.
Hence the set $\{\xi:w_*\xi\in\Xi'\}$ has $\nu$-measure $1$ for each $w\in O_{\mathrm{cub}}$, and so has the intersection $\Xi''$ of these finitely many sets; it is contained in $\Xi'$.
For $\xi\in\Xi''$ and $w\in O_{\mathrm{cub}}$ the six statements of (b) and the assertions of (c) hold at the point $w_*\xi$: the law $w_*\mathrm{lw}_\xi$ lies in $\mathcal{C}$, has (G) and (dF), is $\mathbb{R}^3$-ergodic, invariant under $\{\pm1\}^3$ and reflection positive at every plane $\{y_j=\mathrm{po}\}$, and \cite[Theorem~8.3(i)--(iv)]{QEDI} and \cite[Proposition~8.19]{QEDI} hold for it; these give its properties 1--4.
Thus the first clause of property~6 holds for every $\xi$ of the set $\Xi''$, of full $\nu$-measure, on which the six statements of (b) hold as well; this is the form of \cite[Theorem~3.1, Part~I\,(c)]{QEDI}.
When $\nu$ has no atoms, the null set in property~5 of \cite[Corollary~8.9]{QEDI} cannot be chosen independently of the event.
\end{named}

\subsubsection*{Proposition 45.5. Tensor tests}

\begin{named}{Proposition 45.5}
Let $\mathrm{lw}$ be a law with (G).
\begin{itemize}
\item[(a)] For real $f_1,\dots,f_n\in\mathcal{S}$,
\[
\left\|X(f_1)\cdots X(f_n)\right\|_{L^2(\mathrm{lw})}\le\left(2n\cdot2^n\cdot\Gamma(n)\right)^{\frac{1}{2}}\cdot\prod_{p=1}^n\|f_p\|_2 .
\]
There is exactly one continuous linear map $X^{(n)}\colon\mathrm{Sw}_n\to L^2(\mathrm{lw})$ with $X^{(n)}(f_1\otimes\dots\otimes f_n)=X(f_1)\cdots X(f_n)$.
For $\Psi_0\in\mathrm{Sw}_0=\mathbb{C}$ put $X^{(0)}(\Psi_0):=\Psi_0$, and put $S_n(\Psi):=\mathrm{lw}\left(X^{(n)}(\Psi)\right)$.
\item[(b)] For $\Phi\in\mathrm{Sw}_m$ and $\Psi\in\mathrm{Sw}_n$ one has $X^{(m)}(\Phi)\cdot X^{(n)}(\Psi)=X^{(m+n)}(\Phi\otimes\Psi)$ in $L^1(\mathrm{lw})$, and $\overline{X^{(n)}(\Psi)}=X^{(n)}(\overline{\Psi})$.
For every Euclidean map $\mathrm{Gi}$ with $\mathrm{Gi}_*\mathrm{lw}=\mathrm{lw}$,
\[
\widehat{\mathrm{Gi}}\,X^{(n)}(\Psi)=X^{(n)}\left((\mathrm{Gi}\cdot)^{\otimes n}\Psi\right).
\]
\item[(c)] Let $O\subset\mathbb{R}^3$ be open, and let $\Psi\in\mathrm{Sw}_n$ have compact support contained in $O^n$.
Then $X^{(n)}(\Psi)$ lies in $L^2(\Sigma_O,\mathrm{lw})$.
Let $\hat{n}$ be a unit vector and $\mathrm{po}\in\mathbb{R}$.
If $\Psi\in\mathrm{Sw}_n$ is supported in $\{\hat{n}\cdot y_a\ge\mathrm{po}\text{ for all }a\}$, then $X^{(n)}(\Psi)$ lies in $L^2(\Sigma_{\{\hat{n}\cdot y>\mathrm{po}\}},\mathrm{lw})$.
\item[(d)] Let $\hat{n}$ be a unit vector, and suppose that $\mathrm{lw}$ is invariant under $r_{\hat{n},0}$ and reflection positive at $\{\hat{n}\cdot y=0\}$.
Then
\[
\left\langle\hat{r}_{\hat{n},0}\mathrm{Ga},\mathrm{Ga}\right\rangle_{\mathrm{lw}}\ge0\qquad\text{for every }\mathrm{Ga}\in L^2\left(\Sigma_{\{\hat{n}\cdot y>0\}},\mathrm{lw}\right),
\]
and
\[
\sum_{m,n}S_{m+n}\left(\Theta^{\hat{n}}_m\Psi_m\otimes\Psi_n\right)\ge0
\]
for every finite sequence of elements $\Psi_n\in\mathrm{Sw}_n$ that are supported in $\{\hat{n}\cdot y_a\ge0\text{ for all }a\}$.
\end{itemize}
\end{named}

\begin{named}{Remark}
Proposition~45.5 corresponds to \cite[Lemma~12.22]{QEDII}: the assertion (a) to \cite[Lemma~12.22(a), (b)]{QEDII}, the assertion (b) to \cite[Lemma~12.22(c), (d)]{QEDII}, the first clause of (c) to \cite[Lemma~12.22(e)]{QEDII}, and the second clause of (c) to \cite[Lemma~12.22(f)]{QEDII}, which is stated there only for $\hat{n}=e_j$ and $\mathrm{po}=0$.
The assertions (c) and (d) are stated here for an arbitrary unit normal $\hat{n}$, not only for the coordinate vectors $e_j$.
In the proof of (c) the arguments are shifted by $\hat{n}/k$; this uses the continuity of the translations on $\mathrm{Sw}_n$ and the continuity of $X^{(n)}$, but no invariance of $\mathrm{lw}$.
The first inequality of (d) is \cite[Lemma~8.16(3)]{QEDI}, transported to the plane $\{\hat{n}\cdot y=0\}$.
The last clause of (b) and the first clause of (c) are contained in \cite[Lemma~12.22(d), (e)]{QEDII}; the proof of \cite[Theorem~1.4, Part~II]{QEDII} uses the former for the Euclidean invariance (E1) and the latter for the rate of clustering in (E4).
\end{named}

\subsubsection*{Proposition 45.7. One further mirror is inherited by the components}

\begin{named}{Proposition 45.7}
Let $\mathrm{lw}$ lie in the class $\mathcal{C}$ of \cite[Section~8.1.3]{QEDI}; that is, for $j=1,2,3$ the law $\mathrm{lw}$ is invariant under $\mathbb{R}^3$ and under $\theta_j$, and reflection positive at $\{y_j=0\}$.
Let $\nu$ be as in Proposition~45.6, and let $\hat{n}$ be a unit vector such that $(r_{\hat{n},0})_*\mathrm{lw}=\mathrm{lw}$ and $\mathrm{lw}$ is reflection positive at $\{\hat{n}\cdot y=0\}$.
Then the following holds.
\begin{itemize}
\item[(a)] For $\nu$-almost every $\xi$ the law $\mathrm{lw}_\xi$ is invariant under $r_{\hat{n},0}$, and therefore, being invariant under $\mathbb{R}^3$, under $r_{\hat{n},\mathrm{po}}$ for every real $\mathrm{po}$.
\item[(b)] For $\nu$-almost every $\xi$ the law $\mathrm{lw}_\xi$ is reflection positive at every plane $\{\hat{n}\cdot y=\mathrm{po}\}$, $\mathrm{po}\in\mathbb{R}$.
\item[(c)] As $|s|\to\infty$, the operators $U(s\cdot\hat{n})$ converge weakly on $L^2(\mathrm{lw})$ to the orthogonal projection onto $\{\mathrm{Ga}:U(s\cdot\hat{n})\mathrm{Ga}=\mathrm{Ga}\text{ for all }s\in\mathbb{R}\}$.
\item[(d)] For an arbitrary unit vector $\hat{n}$ let $\mathrm{Ls}_{\hat{n}}$ be the set of the laws that are invariant under $\mathbb{R}^3$ and under $r_{\hat{n},0}$ and reflection positive at $\{\hat{n}\cdot y=0\}$.
Then $\mathrm{Ls}_{\hat{n}}$ belongs to the $\sigma$-algebra of (n7).
Every law in $\mathrm{Ls}_{\hat{n}}$ is invariant under $r_{\hat{n},\mathrm{po}}$ and reflection positive at $\{\hat{n}\cdot y=\mathrm{po}\}$, for every real $\mathrm{po}$.
Thus $\mathrm{Ls}_{\hat{n}}$ is the set of all laws that are invariant under $\mathbb{R}^3$, invariant under $r_{\hat{n},\mathrm{po}}$ for every real $\mathrm{po}$, and reflection positive at every plane $\{\hat{n}\cdot y=\mathrm{po}\}$.
Under the hypotheses of the proposition, $\nu(\mathrm{Ls}_{\hat{n}})=1$ for the unit vector $\hat{n}$ given there.
\end{itemize}
\end{named}

\begin{named}{Remark}
Proposition~45.7 corresponds to \cite[Theorem~8.8(v)]{QEDI} together with its proof.
The proof of \cite[Proposition~12.33(a)]{QEDII} applies \cite[Theorem~8.8(v)]{QEDI} to the six unit vectors $(e_i\pm e_l)/\sqrt{2}$, $1\le i<l\le3$;
for these six vectors the assertions (a) and (b) give six of the nine statements (c2) of Theorem~44\,(D).
The proof of \cite[Theorem~8.8(v)]{QEDI} transports $\mathrm{lw}$ by a rotation $\mathrm{gr}\in O(3)$ with $\mathrm{gr}\,e_1=\hat{n}$: the law $\mathrm{lw}':=(\mathrm{gr}^{-1})_*\mathrm{lw}$ is invariant under $\mathbb{R}^3$ and under $\theta_1$, and reflection positive at $\{y_1=0\}$.
Of the hypothesis $\mathrm{lw}\in\mathcal{C}$ that proof uses only the invariance under $\mathbb{R}^3$.
The assertion (d) adds the measurability of the sets of laws that occur in it; it follows by transport, since the map $\mathrm{lw}'\mapsto\mathrm{gr}_*\mathrm{lw}'$ is bimeasurable.
\end{named}

\subsection*{44. The main theorem}

\subsubsection*{Theorem 44}

\paragraph{Data.}
The following data are fixed:
$\rho\in(0,1)$, the standard weights and the constants of Section~I.8, chosen in their order;
$\bar{m}\in(0,1]$ and $h$;
the probe $\beta$, the point $x_0$, the test $f$, the probes $\beta_{\mathrm{ix}}$ and the tests $g_{\mathrm{ix}}$ of the Data of Proposition~43.2, so that $b(f)>0$ and $x_0$ satisfies the inequality of \cite[Section~2.8]{QEDI};
the exponents $\sigma$ and $\mathfrak{e}_2$ of Section~I.8.5;
the constants $C_U$, $C_{\bar{\kappa}}$, $C'_{\bar{\kappa}}$ and the number $e_2$ of Proposition~43.2;
and the absolute number
\[
c_{\mathrm{cl}}:=\frac{96\cdot\left(1+\sqrt{2}\right)\cdot\left(128+64\cdot\left(\frac{16}{\pi}+\frac{64}{\pi^2}\right)\right)}{4\pi^2},
\]
on which no condition is imposed.
In the notation of \cite[Theorem~12.39]{QEDII} it is the constant $C_0=3\cdot C_3$, where $C_3=32\cdot\left(1+\sqrt{2}\right)\cdot C_2$ and $C_2=\left(128+64\cdot\left(\frac{16}{\pi}+\frac{64}{\pi^2}\right)\right)/(4\pi^2)$.
Here $C_0$, $C_3$ and $C_2$ are the letters of \cite[Theorem~12.39]{QEDII}:
the numbers $C_3$ and $C_2$ are not constants of Section~I.8,
and $C_0$, $C_3$ are not the constants that bear these names in the lines (C6) and (C4r) of Section~I.8.2.

\paragraph{Laws.}
Let $e\in(0,e_2]$.
The numbers $T^{*}(e)$, $R^{*}(e)$, $N_0(e)$ and the real numbers $m_{0,N}(e)$ are those of Section~I.6.2.
For integers $N,T\ge0$ let $P_{N,T}$ be the law of Section~I.7.5 on the even torus with $n=8h\cdot L^{T+N}$ sites per side, formed with $q=e\cdot\eta^{\frac{1}{2}}$ and the lattice mass $\boldsymbol{\mu}_N=\bar{m}\cdot\eta+m_{0,N}(e)$, that is, with the mass parameter
\[
M_{\mathrm{ch}}=\bar{m}\cdot\eta+m_{0,N}(e)+\frac{\rho}{2}\cdot\mathrm{sw}.
\]
Limit points, the sets $\mathcal{V}_N$ and volume-first limit points are those of Section~I.7.5.
A defining sequence of a limit point $\mu$ is a sequence $(N_i,T_i)$ with $N_i\to\infty$, $T_i\to\infty$ and $P_{N_i,T_i}\to\mu$ weakly on $H_{-4}$; along it one puts $\eta_i:=L^{-N_i}$, $n_i:=8h\cdot L^{T_i+N_i}$ and $R_i:=8h\cdot L^{T_i}$.

\begin{named}{Theorem 44}
For every integer $N\ge0$ and every $e$ with $0<e\cdot h^{\frac{1}{2}}\le e_1$ the number $m_{0,N}(e)$ is real and satisfies $|m_{0,N}(e)|<\frac{1}{8}$; it depends on $N$, $e$, $\bar{m}$, $\rho$ and on the scheme (Section~I.6.2), but not on $T$.
For every $e\in(0,e_2]$ the following assertions hold.
\begin{itemize}
\item[(A)] (Cutoffs.)
For all $N\ge N_0(e)$, all $T\ge T^{*}(e)+1$ and every $\mathrm{ix}\in\{0,+,-\}$ the assertions (U-a), (U-a$'$), (U-b) and (U-c) of Proposition~43.2 hold.
In particular, with $\eta=L^{-N}$ and $R=8h\cdot L^{T}$,
\[
\left|\kappa_4\left(X_{\eta,R}(g_{\mathrm{ix}})\right)-e^4\cdot b(g_{\mathrm{ix}})\right|\le\frac{C_{\bar{\kappa}}}{2}\cdot e^{4+\sigma}+C_{\bar{\kappa}}\cdot\left(\eta^2+e^4\cdot\eta^{\frac{1}{4}}\right),
\]
uniformly in $T$.
\item[(B)] (Every limit point.)
\begin{itemize}
\item[(B0)] One has
\[
\sup_{N,T}\int\mathrm{d}P_{N,T}(\omega)\,\|\omega\|_{-3}^2\le61 .
\]
The family $(P_{N,T})$ is tight on $H_{-4}$, and every sequence $(N_i,T_i)$ with $N_i\to\infty$ and $T_i\to\infty$ has a subsequence along which $P_{N_i,T_i}$ converges weakly on $H_{-4}$.
\end{itemize}
Let $\mu$ be a limit point.
\begin{itemize}
\item[(B1)] The law $\mu$ has the properties (P1)--(P8):
\begin{itemize}
\item[(P1)] For real $f'\in\mathcal{S}$,
\[
\mu\left(\exp X(f')\right)\le\exp\left(\frac{1}{2}\langle f',\Pi_{\mathrm{c}}\,f'\rangle\right)\le\exp\left(\frac{1}{2}\|f'\|_2^2\right).
\]
The map $f'\mapsto X(f')$ extends to a bounded linear map $L^2(\mathbb{R}^3;\Lambda^2)\to L^2(\mu)$, for which the same bound holds.
Moreover $\left|\mu\left(\exp\left(z\cdot X(f')\right)\right)\right|\le\exp\left(\frac{1}{2}(\operatorname{Re}z)^2\cdot\|f'\|_2^2\right)$, and $\mu\left(\exp\left(z\cdot X(f')\right)\right)\ne0$ whenever $|z|\cdot\|f'\|_2<0.9005$.
Let $n\ge1$ be an integer, and let $f_1,\dots,f_n$ lie in $L^2(\mathbb{R}^3;\Lambda^2)$ or in its complexification $L^2(\mathbb{R}^3;\Lambda^2)\otimes\mathbb{C}$, on which $X$ is extended complex linearly as in (n8) and $\|f_i\|_2^2=\|\operatorname{Re}f_i\|_2^2+\|\operatorname{Im}f_i\|_2^2$.
On the set of the $z=(z_1,\dots,z_n)\in\mathbb{C}^n$ with
\[
\sum_{i=1}^n|z_i|\cdot\|f_i\|_2\le\frac{2}{5}
\]
one has $\left|\mu\left(\exp\sum_{i=1}^nz_i\cdot X(f_i)\right)-1\right|\le\frac{1}{2}$, so that the principal branch of the logarithm is defined there;
the function $z\mapsto\log\mu\left(\exp\sum_{i=1}^nz_i\cdot X(f_i)\right)$, formed with this branch, is analytic, and of modulus at most $\log2$, on this set.
\item[(P2)] For real $f_1,\dots,f_n\in L^2(\mathbb{R}^3;\Lambda^2)$,
\begin{align*}
\left|S_n(f_1,\dots,f_n)\right|&\le2\cdot\left(2^n\cdot n!\right)^{\frac{1}{2}}\cdot\prod_{i=1}^n\|f_i\|_2,\\
\left|\kappa_n(f_1,\dots,f_n)\right|&\le(\log2)\cdot\left(\frac{5\mathrm{e}}{2}\right)^n\cdot n!\cdot\prod_{i=1}^n\|f_i\|_2 .
\end{align*}
The functional $S_n$ has exactly one extension to a continuous linear functional on $\mathrm{Sw}_n$, and $\left|S_n(\Psi)\right|\le2\cdot66^n\cdot(n!)^{\frac{1}{2}}\cdot q_{6n}(\Psi)$.
\item[(P3)] One has $S_0=1$; each $S_n$ is symmetric and real at real tests; $\mu$ is invariant under $\omega\mapsto-\omega$; and $S_n=0$ for odd $n$.
\item[(P4)] The law $\mu$ is the only law with the moments $S_n$ on tests in $C_c^{\infty}(\mathbb{R}^3;\Lambda^2)$, in the sense of Proposition~45.1\,(8); $\mu(H_{-s})=1$ for $s>\frac{3}{2}$; and $\mu$ has (dF).
\item[(P5)] The law $\mu$ is invariant under $\mathbb{R}^3$, which acts on $L^2(\mu)$ by a strongly continuous unitary group, and under $O_{\mathrm{cub}}$, which acts on two-forms.
\item[(P6)] For $j\in\{1,2,3\}$ and real $\mathrm{po}$,
\[
\left\langle\hat{r}_{e_j,\mathrm{po}}\mathrm{Ga},\mathrm{Ga}\right\rangle_\mu\ge0\qquad\text{for every }\mathrm{Ga}\in L^2\left(\Sigma_{\{y_j>\mathrm{po}\}},\mu\right);
\]
and $\sum_{m,n}S_{m+n}\left(\Theta^{e_j}_m\Psi_m\otimes\Psi_n\right)\ge0$ for every finite sequence of elements $\Psi_n\in\mathrm{Sw}_n$ supported in $\{(y_a)_j\ge0\text{ for all }a\}$.
\item[(P7)] For $j\in\{1,2,3\}$ the reconstruction (n4) applies to $\mu$, with $\mathsf{H}_j=\mathsf{H}_j^{*}\ge0$, $\mathsf{H}_j\Omega_j=0$ and $\|\Omega_j\|=1$, and $V_j$ is a strongly continuous unitary representation of $e_j^{\perp}$ that commutes with $\exp(-t\cdot\mathsf{H}_j)$ and fixes $\Omega_j$.
\item[(P8)] Let $f_1,\dots,f_k\in\mathcal{S}$.
Along a defining sequence of $\mu$ the functions
\[
z\mapsto E^{\mathrm{af}}_{n_i}\exp\left(\sum_{a=1}^kz_a\cdot X_{\eta_i,R_i}(f_a)\right)
\]
converge to $z\mapsto\mu\left(\exp\sum_{a=1}^kz_a\cdot X(f_a)\right)$, uniformly on the compact subsets of $\mathbb{C}^k$, and all moments and all cumulants converge.
\end{itemize}
\item[(B2)] For $\mathrm{ix}\in\{0,+,-\}$,
\[
\left|\kappa_4^{\mu}(g_{\mathrm{ix}})-e^4\cdot b(g_{\mathrm{ix}})\right|\le C'_{\bar{\kappa}}\cdot e^{4+\sigma},\qquad
\left|S_2^{\mu}(g_{\mathrm{ix}},g_{\mathrm{ix}})-\|g_{\mathrm{ix}}\|_2^2\right|\le C'_{\bar{\kappa}}\cdot e^{2-\mathfrak{e}_2}.
\]
\item[(B3)] One has $\kappa_4^{\mu}(f)\ge\frac{1}{2}e^4\cdot b(f)>0$ and $\mathrm{Ext}^{\mu}_{x_0}(f)\ge\frac{1}{2}e^4\cdot b(f)$.
The law $\mu$ is not Gaussian, and it is not a mixture of translation invariant Gaussian laws with (G).
\item[(B4)] (Many points.)
For $\mathrm{ix}\in\{0,+,-\}$ and $m\in\{2,4\}$ there are real numbers $c^{\mathrm{inf}}_{m,\mathrm{ix}}$ and numbers $\mathsf{u}^{\mathrm{inf}}_{m,\mathrm{ix}}\ge0$ with
\begin{gather*}
\mathsf{u}^{\mathrm{inf}}_{4,\mathrm{ix}}\le C_U\cdot e^{\sigma},\qquad
\left(1+\left|c^{\mathrm{inf}}_{2,\mathrm{ix}}\right|\right)\cdot\mathsf{u}^{\mathrm{inf}}_{2,\mathrm{ix}}\le C_U\cdot e^{\sigma},\\
\left|c^{\mathrm{inf}}_{2,\mathrm{ix}}\right|\le C_U\cdot e^{-\mathfrak{e}_2},\qquad
\left|c^{\mathrm{inf}}_{4,\mathrm{ix}}-3\left(c^{\mathrm{inf}}_{2,\mathrm{ix}}\right)^2-b(g_{\mathrm{ix}})\right|\le C_U\cdot e^{\sigma},
\end{gather*}
such that, with
\[
\mathrm{Wk}_{m,\mathrm{ix}}:=\mathrm{Pw}_m\left(X(g_{\mathrm{ix}});\|g_{\mathrm{ix}}\|_2^2,c^{\mathrm{inf}}_{m,\mathrm{ix}}\right),
\]
one has
\[
\left|\mu\left(\prod_{y\in\mathrm{Ypt}}U(y)\,\mathrm{Wk}_{m,\mathrm{ix}}\right)\right|\le\left(\mathsf{u}^{\mathrm{inf}}_{m,\mathrm{ix}}\right)^{\#|\mathrm{Ypt}|}
\]
for every finite subset $\mathrm{Ypt}$ of $\left(R^{*}(e)/L\right)\mathbb{Z}^3$.
One family $\left(c^{\mathrm{inf}}_{m,\mathrm{ix}},\mathsf{u}^{\mathrm{inf}}_{m,\mathrm{ix}}\right)$ with these properties is fixed once and for all; it is not determined by $\mu$.
The assertions (B5) and (D) refer to this family and to its functions $\mathrm{Wk}_{m,\mathrm{ix}}$; they hold for every such family, and the set $\Xi_\mu$ of (D) depends on the family.
\item[(B5)] (Components: nine almost sure statements, each with its own null set.)
\begin{itemize}
\item[(c1)] Three statements: for $\nu$-almost every $\xi$ the law $\mu_\xi$ is invariant and ergodic under $\mathbb{R}^3$; for $\nu$-almost every $\xi$ the law $\mu_\xi$ has (G), for all real $f'\in\mathcal{S}$ simultaneously; for $\nu$-almost every $\xi$ the law $\mu_\xi$ has (dF).
\item[(c3)] Six statements: for each pair $(m,\mathrm{ix})$, for $\nu$-almost every $\xi$ one has $\left|\mu_\xi(\mathrm{Wk}_{m,\mathrm{ix}})\right|\le\mathsf{u}^{\mathrm{inf}}_{m,\mathrm{ix}}$.
\end{itemize}
For every $\xi$ at which these nine statements hold, and for $\mathrm{ix}\in\{0,+,-\}$,
\begin{gather*}
\left|\kappa_4^{\xi}(g_{\mathrm{ix}})-e^4\cdot b(g_{\mathrm{ix}})\right|\le C'_{\bar{\kappa}}\cdot e^{4+\sigma},\qquad
\left|S_2^{\xi}(g_{\mathrm{ix}},g_{\mathrm{ix}})-\|g_{\mathrm{ix}}\|_2^2\right|\le C'_{\bar{\kappa}}\cdot e^{2-\mathfrak{e}_2},\\
\mathrm{Ext}^{\xi}_{x_0}(f)\ge\frac{1}{2}e^4\cdot b(f),\qquad
\kappa_4^{\xi}(f)\ge\frac{1}{2}e^4\cdot b(f)>0 .
\end{gather*}
\end{itemize}
\item[(C)] (Volume first.)
For every $N$ the set $\mathcal{V}_N$ is not empty; volume-first limit points exist; and every volume-first limit point is a limit point, so that (B) holds for it.
Let $\mu$ be a volume-first limit point.
\begin{itemize}
\item[(C1)] The law $\mu$ lies in $\mathcal{C}_9$.
\item[(C2)] $\mathrm{Gi}_*\mu=\mu$ for every Euclidean map $\mathrm{Gi}$ of $\mathbb{R}^3$, that is, for every element of $\mathbb{R}^3\rtimes O(3)$; and $S_n(\mathrm{Gi}\cdot f_1,\dots,\mathrm{Gi}\cdot f_n)=S_n(f_1,\dots,f_n)$ for every $n$ and all $f_1,\dots,f_n\in L^2(\mathbb{R}^3;\Lambda^2)$.
\item[(C3)] For every unit vector $\hat{n}$ and every real $\mathrm{po}$,
\[
\left\langle\hat{r}_{\hat{n},\mathrm{po}}\mathrm{Ga},\mathrm{Ga}\right\rangle_\mu\ge0\qquad\text{for every }\mathrm{Ga}\in L^2\left(\Sigma_{\{\hat{n}\cdot y>\mathrm{po}\}},\mu\right).
\]
\item[(C4)] For $j=1,2,3$ the joint spectral measure of $\left(\mathsf{H}_j,(\mathsf{P}_i)_{i\ne j}\right)$ is carried by the set $\left\{(E,p)\in[0,\infty)\times\mathbb{R}^2:E\ge|p|_2\right\}$.
\item[(C5)] $P^j_\mu=P^{123}_\mu$ for $j=1,2,3$.
\item[(C6)] For every $\mathrm{vp}>0$, all $\mathrm{Ga},\mathrm{Gb}\in L^2\left(\Sigma_{(-\mathrm{vp},\mathrm{vp})^3},\mu\right)$ and all $a\in\mathbb{R}^3$ with $|a|_\infty\ge4\,\mathrm{vp}$,
\[
\left|\langle\mathrm{Ga},U(a)\mathrm{Gb}\rangle_\mu-\langle\mathrm{Ga},P^{123}_\mu\mathrm{Gb}\rangle_\mu\right|\le c_{\mathrm{cl}}\cdot\|\mathrm{Ga}\|_2\cdot\|\mathrm{Gb}\|_2\cdot\frac{\mathrm{vp}}{|a|_\infty}.
\]
\item[(C7)] There is exactly one positive measure $\mathrm{kl}_\mu$ on $[0,\infty)$, of mass at most $1$, such that
\[
S_2(f',f'')=(2\pi)^{-3}\int\mathrm{d}p\,\overline{\hat{b}'(p)}\cdot\Pi(p)\,\hat{b}''(p)\cdot\int\mathrm{d}\,\mathrm{kl}_\mu(t)\,\frac{|p|_2^2}{|p|_2^2+t^2}
\]
for all real $f',f''\in L^2(\mathbb{R}^3;\Lambda^2)$, where $b'=\ast f'$ and $b''=\ast f''$.
\item[(C8)] The functionals $S_n$ on $\mathrm{Sw}_n$, coinciding arguments included, satisfy (E0$'$), (E1), (E2) and (E3) below.
\end{itemize}
\item[(D)] (Components of a volume-first limit point; $3+9+6$ almost sure statements.)
Besides (c1) and (c3) the following nine statements hold.
\begin{itemize}
\item[(c2)] Nine statements, one for each of the nine families of planes, namely the three families $\{y_i=\mathrm{po}\}$ with $i\in\{1,2,3\}$ and the six families $\{y_i-y_l=\mathrm{po}\}$ and $\{y_i+y_l=\mathrm{po}\}$ with $1\le i<l\le3$, where in each family $\mathrm{po}$ ranges over $\mathbb{R}$: for $\nu$-almost every $\xi$ the law $\mu_\xi$ is invariant under the reflections in all planes of the family and reflection positive at all of them.
\end{itemize}
Let $\Xi_\mu$ be the set of the $\xi$ at which the $18$ statements (c1), (c2), (c3) hold.
Then $\nu(\Xi_\mu)=1$, and for every $\xi\in\Xi_\mu$ the single law $\mu_\xi$ has the following properties.
\begin{itemize}
\item[(D1)] It is a law on closed two-forms with (G); it satisfies (H) and (dF) and lies in $\mathcal{C}_9$; and it is $\mathbb{R}^3$-ergodic.
\item[(D2)] It has the properties (C2), (C3), (C4) and (C7), with $\mu_\xi$ in place of $\mu$.
\item[(D3)] $\ker\mathsf{H}_j(\mu_\xi)=\mathbb{C}\,\Omega_j$ for $j=1,2,3$.
\item[(D4)] For every $\mathrm{vp}>0$, all $\mathrm{Ga},\mathrm{Gb}\in L^2\left(\Sigma_{(-\mathrm{vp},\mathrm{vp})^3},\mu_\xi\right)$ and all $a\in\mathbb{R}^3$ with $|a|_\infty\ge4\,\mathrm{vp}$,
\[
\left|\mu_\xi\left(\overline{\mathrm{Ga}}\cdot U(a)\mathrm{Gb}\right)-\overline{\mu_\xi(\mathrm{Ga})}\cdot\mu_\xi(\mathrm{Gb})\right|\le c_{\mathrm{cl}}\cdot\|\mathrm{Ga}\|_2\cdot\|\mathrm{Gb}\|_2\cdot\frac{\mathrm{vp}}{|a|_\infty};
\]
moreover $\mu_\xi\left(\overline{\mathrm{Ga}}\cdot U(a)\mathrm{Gb}\right)\to\overline{\mu_\xi(\mathrm{Ga})}\cdot\mu_\xi(\mathrm{Gb})$ as $|a|\to\infty$ in any manner, for all $\mathrm{Ga},\mathrm{Gb}\in L^2(\mu_\xi)$.
\item[(D5)] It satisfies (E0$'$), (E1), (E2), (E3) and (E4) below.
\item[(D6)] The four inequalities at the end of (B5) hold for it, and $\kappa_4^{\xi}(\mathrm{Gi}\cdot f)=\kappa_4^{\xi}(f)\ge\frac{1}{2}e^4\cdot b(f)>0$ for every Euclidean map $\mathrm{Gi}$.
Hence $\mu_\xi$ is not Gaussian, and it is not the point mass at $\omega=0$.
\end{itemize}
\end{itemize}
\emph{The five properties.}
Let $\mathrm{lw}$ be a law with (G), and let $S_n=S_n^{\mathrm{lw}}$.
\begin{itemize}
\item[(E0$'$)] $S_0=1$, and $\left|S_n(\Psi)\right|\le2\cdot66^n\cdot(n!)^{\frac{1}{2}}\cdot q_{6n}(\Psi)$ for all $n\ge1$ and $\Psi\in\mathrm{Sw}_n$.
(Here $2\cdot66^n\cdot(n!)^{\frac{1}{2}}\le2\exp(2178)\cdot n!$.)
\item[(E1)] $S_n\left((\mathrm{Gi}\cdot)^{\otimes n}\Psi\right)=S_n(\Psi)$ for every Euclidean map $\mathrm{Gi}$, reflections included.
\item[(E2)] For every unit vector $\hat{n}$ and every finite sequence $\Psi_0\in\mathbb{C}$, $\Psi_n\in\mathrm{Sw}_n$ of elements supported in $\{\hat{n}\cdot y_a\ge0\text{ for all }a\}$,
\[
\sum_{m,n}S_{m+n}\left(\Theta^{\hat{n}}_m\Psi_m\otimes\Psi_n\right)\ge0 .
\]
\item[(E3)] $S_n$ is invariant under the permutations of its $n$ arguments together with their tensor indices.
\item[(E4)] $S_{m+n}\left(\Psi\otimes\tau_a^{\otimes n}\Phi\right)\to S_m(\Psi)\cdot S_n(\Phi)$ as $|a|\to\infty$ in any manner, for $\Psi\in\mathrm{Sw}_m$ and $\Phi\in\mathrm{Sw}_n$.
\end{itemize}
\end{named}

\begin{named}{Remark}
The inequality (E0$'$) bounds $|S_n(\Psi)|$ by the $n$-th term of a sequence times a norm of $\Psi$ of index $ns$; here $s=6$, the norms are the $q_m$, and the sequence is $2\cdot66^n\cdot(n!)^{\frac{1}{2}}$, which is at most $\alpha\cdot(n!)^{\beta}$ with $\alpha=2\exp(2178)$ and $\beta=1$.
In this remark $s$, $\alpha$, $\beta$ have only the meaning just given; in particular $s$ is neither the exponent $s$ of the line (C7) of Section~I.8.2 nor the index of $H_{-s}$ in (P4).
In this form (E0$'$) is the literal analogue, for tests on $\mathbb{R}^{3n}$ with values in $(\Lambda^2)^{\otimes n}$, of the linear growth condition of Osterwalder and Schrader \cite{OsterwalderSchrader1975}, with $s=6$ and the constants $2\cdot66^n\cdot(n!)^{\frac{1}{2}}$; that paper formulates its conditions for scalar tests on $\mathbb{R}^{4n}$.
No assertion of Theorem~44 depends on more than the inequality (E0$'$) as displayed.
\end{named}

\begin{named}{Remark}
Theorem~44 corresponds to Part~II of \cite[Theorem~3.1]{QEDI} and to \cite[Corollary~3.2]{QEDI}, together with the parts of their proofs in \cite[Section~21.10.1]{QEDI} and \cite[Section~21.10.2]{QEDI} that concern the limit points and their components, and to Part~I of \cite[Theorem~3.1]{QEDI}, read for the laws $P_{N,T}$ of Section~I.7.5.
That every volume-first limit point is a limit point follows by the argument of \cite[Remark~21.17]{QEDI}, which proves \cite[Theorem~3.1, Part~I\,(d)]{QEDI}.
Theorem~44 compares with \cite{QEDI} as follows.
\begin{itemize}
\item[(1)] The laws are formed with the determinant of the chamber operator in place of the Wilson determinant; only the all-flux law occurs, on the even tori with $8h\cdot L^{T+N}$ sites per side; and (A) requires $T\ge T^{*}(e)+1$.
\item[(2)] The bound in (A) contains $\eta^{\frac{1}{4}}$ where \cite[Theorem~3.1(ii)]{QEDI} has $\eta^{\frac{1}{2}}$.
\item[(3)] The following assertions of (C) and (D) have no counterpart in the assertions of \cite[Theorem~3.1]{QEDI} and \cite[Corollary~3.2]{QEDI} on the limit points and their components:
reflection positivity at the families of planes $\{y_i-y_l=\mathrm{po}\}$ and $\{y_i+y_l=\mathrm{po}\}$, and with it the class $\mathcal{C}_9$ in (C1) and (D1) and the statements (c2) on these families;
the invariance of the single component under $O_{\mathrm{cub}}$, by which it satisfies (H) in (D1);
the clauses (C2)--(C7) and (D2);
the unique vacuum (D3) of the $\mathbb{R}^3$-ergodic components;
the rate in (D4);
(E1);
(E2) for unit vectors $\hat{n}\notin\{e_1,e_2,e_3\}$;
the clustering in (E4) of the Schwinger functions of every order as $|a|\to\infty$ in any manner;
and the Euclidean invariance of $\kappa_4^{\xi}$ in (D6).
Invariance and reflection positivity at planes other than the planes $\{y_j=\mathrm{po}\}$, the invariance under $O(3)$, the identities $P^j=P^{123}$ of (C5), the unique vacuum of the $\mathbb{R}^3$-ergodic components and the clustering as $|a|\to\infty$ in any manner occur in \cite[Theorem~8.8]{QEDI}, in particular in its parts (ii) and (iv)--(vi), where they follow from hypotheses on the law that \cite{QEDI} does not establish for its limit points;
the round cone, the rate of clustering, the K\"all\'en--Lehmann form and, for Schwinger functions, the conditions (E1) and (E2) at every plane are not stated in \cite{QEDI}.
Among the other assertions of (C) and (D), the third assertion at the head of (C) corresponds to the first assertion of Part~I\,(d) of \cite[Theorem~3.1]{QEDI};
for the statements (c2) on the families $\{y_i=\mathrm{po}\}$ and for the set $\Xi_\mu$ see (4);
(E0$'$) and (E3) in (C8) correspond to (P2) and (P3) of \cite[Theorem~3.1]{QEDI}, and (E0$'$) in (D5) corresponds to the property (P2) of $\mu_\xi$ in \cite[Corollary~3.2(a)]{QEDI}, see (5) and (6);
the closedness, (G), (dF) and the ergodicity in (D1) are the statements (c1) at the point $\xi$, which correspond to assertions at the head of Part~I\,(c) of \cite[Theorem~3.1]{QEDI} and in \cite[Corollary~3.2(a)]{QEDI};
and the four inequalities in (D6), which are those of (B5), and the non-Gaussianity of $\mu_\xi$ correspond to Part~II\,(vi) and (vii) of \cite[Theorem~3.1]{QEDI} and to the end of \cite[Corollary~3.2(a)]{QEDI}.
The clauses (D3), (E4) and the second sentence of (D4) are stronger forms of the properties~3 and~4 of Part~I\,(c) of \cite[Theorem~3.1]{QEDI}, in which $\dim\ker\mathsf{H}_j\in\{1,\infty\}$, the moments cluster as $\min_l|a_l|\to\infty$, and the two-point function clusters as $|a|\to\infty$ in any manner;
the equality $\ker\mathsf{H}_j=\mathbb{C}\,\Omega_j$ is asserted in \cite[Corollary~3.2(b)]{QEDI} for the components along one axis, see (7).
\item[(4)] In Theorem~44, as in Part~I\,(c) of \cite[Theorem~3.1]{QEDI}, there are almost sure statements on the components, a set of full $\nu$-measure on which they hold together, and properties of the single law $\mu_\xi$ for every $\xi$ of that set.
The differences concern which assertions are made, the listing of the almost sure statements one by one, and the two sets that are named.
In Part~I\,(c) of \cite[Theorem~3.1]{QEDI} the assertions at its head and the properties 1--4 stand under one quantifier ``for $\nu$-almost every $\xi$'', and the set of full measure on which they hold together is denoted $\Xi_\mu$; it is not the set $\Xi_\mu$ of (D), which is defined by the $18$ statements (c1), (c2), (c3) and depends on the family fixed in (B4).
There property~5 and the second clause of property~6 hold, for each event $\mathrm{Ev}\in\Sigma$ with $\mu(\mathrm{Ev})=1$, outside a $\nu$-null set that depends on $\mathrm{Ev}$, and the last clause of property~6 is a statement on $\nu$.
In \cite[Corollary~3.2(a)]{QEDI} the properties of the components stand under one quantifier ``for $\nu$-almost every $\xi$''.
Of the assertions of Part~I\,(c), Theorem~44 contains for every limit point $\mu$ the hypothesis (H), by (P1), (P5) and (P6), and, for the components, the three statements (c1): invariance and ergodicity under $\mathbb{R}^3$, (G), and (dF).
For a volume-first limit point it contains in addition the invariance of the components under $\{\pm1\}^3$ and their reflection positivity at the planes $\{y_j=\mathrm{po}\}$, through the three families $\{y_i=\mathrm{po}\}$ of (c2), together with the properties (D1)--(D6) of the single components.
For an arbitrary limit point $\mu$, which is soft by (B1), the assertions of Part~I\,(c) on the components other than those of (c1) follow from Proposition~45.4\,(b)--(d), applied to $\mu$.
In Theorem~44 the almost sure statements are listed one by one, each with its own null set: the nine statements (c1) and (c3) stand in (B5) and are made for every limit point, and the nine statements (c2) stand in (D) and are made for a volume-first limit point; there are $18$ in all.
Every other assertion on the components is made pointwise, as a property of the single law $\mu_\xi$, for every $\xi$ in a set named in the theorem, and there are two such sets.
The four inequalities at the end of (B5) hold, $\mu$ being any limit point, for every $\xi$ at which the nine statements (c1) and (c3) hold.
The properties (D1)--(D6) hold, $\mu$ being a volume-first limit point, for every $\xi$ in the set $\Xi_\mu$, of full $\nu$-measure, at which all $18$ statements hold.
Since $\Xi_\mu$ is contained in the first of these two sets, the four inequalities hold at every $\xi\in\Xi_\mu$, as (D6) states.
\item[(5)] In (E0$'$) the inequality, the value $S_0=1$ and the factorial growth of the constants are displayed; \cite[Theorem~3.1, (P2)]{QEDI} states the same bound in the form of the linear growth condition of Osterwalder and Schrader.
\item[(6)] Let $\xi$ be a point at which the second statement of (c1) holds, so that the component $\mu_\xi$ has (G), that is, (L1).
For such a component the bound on $S_n(f_1,\dots,f_n)$ in (P2), and the extension of $S_n$ to $\mathrm{Sw}_n$ with its bound by $q_{6n}$, are Proposition~45.2\,(1) and (2) for the law $\mu_\xi$; these hold under (L1) by Proposition~45.2\,(8).
The bound on $\kappa_n(f_1,\dots,f_n)$ in (P2) is Proposition~45.1\,(1) for the law $\mu_\xi$, namely conclusion (7) of \cite[Proposition~7.10]{QEDI}, which holds under (L1) by Proposition~45.1\,(11); here $\kappa_n$ is the cumulant of (n5), which for a law with (G) is the derivative $\partial_{z_1}\cdots\partial_{z_n}\log\mu_\xi\left(\exp\sum_{a=1}^nz_a\cdot X(f_a)\right)$ at $z=0$.
In \cite{QEDI} these bounds are obtained inside the proof of \cite[Corollary~3.2]{QEDI}.
\item[(7)] The assertion (v) of Part~II of \cite[Theorem~3.1]{QEDI}, which concerns a set of components of positive $\nu$-measure, and \cite[Corollary~3.2(b)]{QEDI}, which concerns the components along one axis, have no counterpart in Theorem~44.
The assertion (B5) is stronger than (v); and for a volume-first limit point the decomposition along an axis is, by (C5), the $\mathbb{R}^3$-ergodic one.
\item[(8)] \cite[Corollary~3.3]{QEDI}, on the zero-harmonic law, has no counterpart.
\item[(9)] The bound is $|m_{0,N}(e)|<\frac{1}{8}$; Part~II of \cite[Theorem~3.1]{QEDI} states $|m_{0,N}(e)|<\frac{1}{2}$, and \cite[Theorem~21.3]{QEDI} states the bound $\frac{1}{8}$.
\item[(10)] The step from the many-point bound (B4) to the components uses three clauses of Proposition~45.6:
the identity $\mathrm{lw}_\omega(\Phi)=E\left[\Phi\mid I\right](\omega)$ for integrable $\Phi$, which is (a$'$);
the measurability of $\xi\mapsto\mu_\xi(\mathrm{Wk}_{m,\mathrm{ix}})$ on the measurable set $\left\{\xi:\mu_\xi\left(|\mathrm{Wk}_{m,\mathrm{ix}}|\right)<\infty\right\}$, which is the second sentence of (d0);
and the step from the bound on $E_\mu\left[\mathrm{Wk}_{m,\mathrm{ix}}\mid I\right]$ to the same bound for $\nu$-almost every component, which is (d).
In \cite{QEDI} the identity for integrable $\Phi$ is the display in the proof of \cite[Theorem~8.1]{QEDI} and in Step~1 of the proof of \cite[Proposition~9.7(a)]{QEDI}, and \cite[Theorem~8.1(a)]{QEDI} is its case of indicator functions.
The proof of \cite[Proposition~9.7(b)]{QEDI} contains two measurability statements:
the map $\xi\mapsto\int\mathrm{d}\mu_\xi\,\Phi$ is measurable for every measurable $\Phi$ with values in $[0,\infty]$;
and, for every measurable complex $\Phi$, in particular for a signed integrable $\Phi$ such as $\mathrm{Wk}_{m,\mathrm{ix}}$, the set $\left\{\xi:\int\mathrm{d}\mu_\xi\,|\Phi|<\infty\right\}$ is measurable and $\xi\mapsto\int\mathrm{d}\mu_\xi\,\Phi$ is measurable on it.
The second of them is the second sentence of Proposition~45.6\,(d0).
In the proof of \cite[Corollary~3.2(b)]{QEDI} in \cite[Section~21.10.2]{QEDI}, on the set of the laws on $\mathcal{S}'$ with its $\sigma$-algebra, the first statement is taken from the first paragraph of the proof of \cite[Lemma~8.17]{QEDI}, the second is derived from it for real $\Phi$, and both are applied to the decomposition along one axis.
The clause (d) corresponds to the application of \cite[Proposition~9.7(b)]{QEDI} in Step~6 of the proof in \cite[Section~21.10.1]{QEDI}, by which the bound on the conditional expectation yields the same bound for $\nu$-almost every component.
Thus each of the clauses (a$'$), (d0) and (d) has its counterpart in \cite{QEDI}; there they stand inside proofs, whereas here they are clauses of the statement of Proposition~45.6.
\end{itemize}
\end{named}

\section*{Part III. Annexes}

\subsection*{Annex A. The forms in which clauses proved elsewhere are used}

Each of the following lists belongs to one range of Propositions.
It states clauses of Part~I, of Propositions outside that range and of \cite{QEDI}, \cite{QEDII}, in the form in which the proofs of the Propositions of the range need them, and it names further clauses that these proofs use in the form in which they are stated.
A clause that the proofs of a range use and that does not appear in its list is needed in the form in which it is stated.
The list of Section~G1.Z1, for Propositions~1.1--14.4, contains in addition entries on which no proof of these Propositions rests (see the first paragraph of Section~G1.Z1).

\subsubsection*{G1.Z1. The forms needed for Propositions 1.1--14.4}

The entries of this list are of two kinds.
The entries of the first kind are clauses on which proofs of Propositions~1.1--14.4 rest.
They are (G1.T-1)--(G1.T-4);
(G1.T-9), that is, Propositions~30.1\,(c) and~30.3\,(e), which the proof of Proposition~14.2 uses and which stand before the Propositions~14.m in Part~II;
and, of (G1.T-11), only Propositions~35.3\,(c1), 35.10, 36\,(g) and~37.2\,(a), (d), (e), which the proof of Proposition~14.3\,(g) uses, and Propositions~37.1\,(f), (g) and~36\,(c), which the proof of Proposition~14.4\,(a) uses for the letters $C_E$, $\mathrm{Csy}_{\mathrm{ak}}$, $\mathrm{Cpv}_0$.
These uses of later Propositions are among the exceptions to the order of Part~II listed in Section~IV.1, and the order in which the clauses concerned are proved is given in Section~IV.1\,(B14) and~(B16).
The entries of the second kind are those on which no proof of Propositions~1.1--14.4 rests:
(G1.T-5)--(G1.T-8), (G1.T-10), (G1.T-12) and the other parts of (G1.T-11);
among the latter, $C^{G}_{\beta}$, $C_{\mathrm{mat}}$, $s_{B}$, $C_{3\mathrm{p}}$, $C_{\Gamma}^{(16)}$, $C_{U}$, $C_{\bar{\kappa}}$ and $C'_{\bar{\kappa}}$ are given by their formulas in Proposition~14.2\,(a), and $e_2$ is given by Proposition~14.3\,(g).
The entries of the second kind show with which constants, numerals and conditions the Propositions named at the heads of (G1.T-5)--(G1.T-8), (G1.T-10) and (G1.T-11) are stated, that is, where the rows of Proposition~14.2\,(a), (b) whose reference or whose indication ``Used in'' is one of these Propositions come from.
In the proofs of Propositions~14.2--14.4 these rows are read as definitions (Section~IV.1\,(B1), and Section~IV.1, paragraph~D; Chapter~14b, Step~0\,(i), and the paragraph ``Other Propositions'' of Chapter~14b).
That these Propositions hold with these constants, numerals and conditions is proved in their own chapters, in the order of the proofs given in Section~IV.1.
All of them stand after the Propositions~14.m in Part~II, except Proposition~15, named at the head of (G1.T-6), and Proposition~16.1, named at the head of (G1.T-7).
These two stand before Proposition~14.1 in Part~II and belong to the entry (B11$'$) of Section~IV.1, which comes after (B1)--(B5) and before (B12);
in the sense of Section~IV.1 they quote no statement, and of the Propositions~14.m they take only rows of Proposition~14.2\,(a), (b), as definitions and as conditions on the choice of the constants, and lines of Proposition~14.2\,(f), which are parts of (B1) and (B5).

\paragraph{(G1.T-1)}
From Part~I the following are used in the form in which they are stated:
(I.s1)--(I.s4), (I.s6);
(I.c1)--(I.c6);
(I.m1)--(I.m7);
(I.O1), (I.O2), (I.O4)--(I.O8);
(I.J1)--(I.J4), (I.J6);
(I.Q1)--(I.Q3);
(I.ke1)--(I.ke5);
(I.a1)--(I.a4);
Sections~I.4.1, I.4.4 and I.4.5;
(I.POS);
(I.1.4), (I.1.7), (I.2.3), (I.2.4), (I.2.6), (I.3.2)--(I.3.4), (I.4.1), (I.5.4);
Sections~I.1.2 and I.1.9;
Sections~I.3.4, I.3.5, I.3.7 and I.3.9;
Section~I.5.4 with (I.k1)--(I.k7), and Sections~I.5.6, I.5.7, I.5.8;
Section~I.6.3;
(I.L0.1)--(I.L0.18) and (I.L0.26).

\paragraph{(G1.T-2)}
Polymers: \cite[Proposition~5.58]{QEDI}, whose proof is \cite[Theorem~13.2]{QEDI}.
Let $\mathrm{pr}$ denote the projection of $\mathbb{R}^3$ onto $\mathbb{R}^3/\ell\mathbb{Z}^3$.
In (t2) and (t8), $\kappa_0$ is the constant $54\log72$ of the line (C0) of Section~I.8.2.
The following clauses are needed.
\begin{itemize}
\item[(t1)] Every $X\in D_k$ satisfies $|X|_M\le c_X\cdot\left(d_M(X)+1\right)$.
\item[(t2)] For every $M$-cube $\square$ of $T_k$,
\[
\sum_{X\in D_k,\ X\supseteq\square}\exp\left(-\kappa_0\cdot d_M(X)\right)\le K_0 .
\]
\item[(t8)] Let $\ell$ and $\ell/L$ lie in $\mathrm{Cl}$, let $Y\in D_{k+1}$, and let $\mathrm{kq}\ge\kappa_0$.
Then
\[
\sum_{X:\ \bar{X}^{L}=LY}\exp\left(-\mathrm{kq}\cdot d_M(X)\right)\le c_{\mathrm{rb}}\cdot L^3\cdot\exp\left(-L\cdot\left(\mathrm{kq}-\kappa_0-1\right)\cdot d_M(Y)\right).
\]
\item[(t7)] The assertion on $\mathbb{R}^3$: any two points of a finite union $\tilde{X}$ of $M$-cubes of $\mathbb{R}^3$ that is connected through faces have sup distance at most $M\cdot\left(d_M(\tilde{X})+2\right)$, where $d_M(\tilde{X})$ is formed with nets in $\mathbb{R}^3$.
\item[(t10)] The first assertion: if $X\in D_k$ satisfies $2M\cdot\left(d_M(X)+2\right)<\ell$, then $X=\mathrm{pr}(\tilde{X})$ for such a union $\tilde{X}$ on which $\mathrm{pr}$ is injective, and $d_M(\tilde{X})=d_M(X)$.
\end{itemize}

\paragraph{(G1.T-3)}
Domains: \cite[Lemmas~5.18--5.20]{QEDI}, \cite[Section~17.3]{QEDI} and the assertion (g0) of \cite[Section~17.4.1]{QEDI}.
Let $\ell,\ell/L\in\mathrm{Cl}$ and $e_{k+1}=L^{\frac{1}{2}}\cdot e_k$.
Then $\mathscr{A}\in\tilde{R}_{k+1}$ implies $\mathscr{A}_L\in L^{-\frac{3}{4}}\cdot\tilde{R}_k$, which is a subset of $\tilde{R}_k$;
and $\mathscr{C}\in\mathfrak{X}_{k+1}$ implies $\mathscr{C}_L\in\mathfrak{X}_k$.
The set $\tilde{R}_k$ is invariant under the addition of real constants.
At every level whose side lies in $\mathrm{Cl}$ and whose charge satisfies $e_k\le e_0$, the sets $\tilde{R}_k$ and $\mathfrak{X}_k$ are open (\cite[Lemma~17.3]{QEDI}), and they are invariant under $\mathbb{G}_k$, under $A\mapsto-A$ and under $A\mapsto\overline{A}$ (\cite[Lemma~17.6]{QEDI}).
The map $\mathrm{pj}^{*}$, for the tori $T_1$, $T_2$ of \cite[Definition~5.13]{QEDI}, maps $\tilde{G}_k(T_1)$ into $\tilde{G}_k(T_2)$.
By (g0), the unit blocks, the $L$-blocks (together with the unit blocks of their sites) and the $M$-cubes are nested, and the $M$-cubes are nested in the $LM$-cubes.

\paragraph{(G1.T-4)}
Spaces of activities (Sections~I.5.4 and I.5.6; Propositions~16.1--18.6).
The spaces $K_{k+1}$ and $K^{x}_{k+1}$ are needed with the clauses
(I.k1) (analyticity on $\tilde{R}_{k+1}$, resp.\ on $\tilde{R}_{k+1}\times\mathfrak{X}_{k+1}$),
(I.k2) (dependence on the bonds that meet $Y$),
(I.k3) (invariance under $\mathbb{G}_{k+1}$),
(I.k4) (gauge invariance in each slot on the bonds at distance at most $1$ from $Y$),
(I.k5) and (I.k7) (charge conjugation and reality),
and (I.k6) (finiteness of the norm $\sup\exp\left(\kappa\cdot d_M(Y)\right)\cdot\left|\,\cdots\right|$).
For activities without fermion legs,
\[
(L E)(Y;\mathscr{A},\mathscr{C})=\sum_{X:\ \bar{X}^{L}=LY}E\left(X;\mathscr{A}_L,\mathscr{C}_L\right).
\]
From Proposition~16.1, for a side $\ell\in\mathrm{Cl}$:
\begin{itemize}
\item[(a1)] for $\tilde{g}\in\mathbb{G}_k$ the map $X\mapsto\tilde{g}X$ is a bijection of $D_k$, and the same holds at the level $k+1$ if $\ell/L$ lies in $\mathrm{Cl}$;
moreover $\overline{\tilde{g}X}^{L}=\tilde{g}\left(\bar{X}^{L}\right)$ for $\tilde{g}\in\mathbb{G}_k^{L}$;
\item[(d1)] the map $\tilde{g}=(S,a)\mapsto\tilde{g}':=(S,a/L)$ is an isomorphism of $\mathbb{G}_k^{L}$ onto $\mathbb{G}_{k+1}$;
\item[(d2)] $L\left(\tilde{g}'Y\right)=\tilde{g}\left(LY\right)$;
\item[(d3)] $\left(\tilde{g}'A\right)_L=\tilde{g}\left(A_L\right)$;
\item[(d4)] $\left(d\Lambda\right)_L=d\left(\Lambda_L\right)$.
\end{itemize}

\paragraph{(G1.T-5)}
Projection of systems of primitives (Propositions~26.1--26.6).
Here W1--W17 are the properties of systems of primitives of \cite[Section~5.4.2]{QEDI}, in the form of Section~G1.11 and of the table of Section~G4.0.4.
The projection of a system with the properties W1--W17 for $\mathrm{par}^{0}$ has the properties W1--W17 for $\mathrm{par}$, with
\[
C_w=64\,K_0\cdot C_w^{(1)} .
\]
The factor $64\,K_0$ comes from $\gamma_{\mathrm{w}}\cdot M\ge1$, $K_0\ge1$, $\kappa_{\mathrm{w}}\ge0$ and the clause (t2) of \cite[Proposition~5.58]{QEDI}.
The proof uses in addition that $L$ is odd, that $M$ is a power of $L$, $M\ge L^2$, $M>4L$, $M\ge M^{P}$, $0<\alpha_0\le\frac{1}{4}$, and
\[
C_w^{(1)}\ge4\exp\left(2\gamma_{\mathrm{w}}\right)\cdot M^{27\alpha_0} ;
\]
it uses moreover, through Proposition~11\,(b) at the property W1 and through Proposition~8\,(b) in Proposition~26.3\,(c) and in W17, the conditions of these two Propositions, namely (L5), (M1), (M2), (M5), (Z2.s1), (Z2.s2) and $\alpha<1$
(Section~G4.A.6, first row).
All the conditions named above, except (Z2.s1) and (Z2.s2), are conditions or definitions of the lines (C0)--(C6) of Section~I.8.2.
The conditions (Z2.s1) and (Z2.s2) are conditions on $e_0$ and stand in the line (C8); they contain the number $\varepsilon$ of the line (C7), and $e_0$ is chosen after $C_w$.
None of the conditions named above contains $C_w$;
in particular, the letters of (Z2.s1) are $L$ and $\varepsilon$, those of (Z2.s2) are $C'_V$, $a'$ and $\varepsilon$, and the number $M^{F}$ of (M5) is formed with $L$, $C'_V$, $a'$ and $\kappa'$, of which $C'_V$ and $\kappa'$ are constants of the line (C4) and $a'$ is a constant of the line (C4r).
Hence $C_w$ is the number $64\,K_0\cdot C_w^{(1)}$ given by its formula in the line (C6).

\paragraph{(G1.T-6)}
The step (Propositions~15, 23.1, 23.2 and 24.1--24.3).
The weight $\mathsf{w}_W=\max\nu^{-\frac{1}{2}}$ of (I.5.2) equals $1$ for $k\ge1$, and $\mathsf{w}_W=\sqrt{48}\le7$ at $k=0$.
The bound $\mathsf{w}_W\le7$, which holds at every level, is what enters the constant $C_W=28\,C_w^2\cdot c_{\gamma}$ and the condition (A6) below; both are used at all levels.
The step requires the conditions (A1)--(A5), (A7), (A8) with the numerals of \cite[Section~5.6.2]{QEDI}, which are those of the rows (A1)--(A5), (A7), (A8) of Proposition~14.2\,(b), and the condition (A6) in the form
\[
112\,C_w^2\cdot c_{\gamma}\cdot e_k^{2\varepsilon}\le3 ,
\]
whereas the condition (A6) of \cite[Section~5.6.2]{QEDI} carries the numeral $16$ in place of $112$.
The second part of (A3), $96\,M^3\cdot e_k^{\frac{1}{4}}\le1$, serves the level $0$, together with the bound $12\,M^3\cdot q^{\frac{1}{2}-\frac{\varepsilon}{2}}$ on $\left\|G^{(1)}(s,\square)\right\|_{h_0/2}$ (Proposition~23.2).
The constant $c_{\gamma}$ is given by its formula of Section~I.5.8.

\paragraph{(G1.T-7)}
Propositions~16.1--18.6.
The constant $c_{\mathrm{fl}}=61\,c_X\cdot K_0$ is an absolute number.
The conditions (A9) of Section~I.8.4 and (L1), (S1) of Proposition~14.2\,(b) are required in the form in which they are stated;
the condition (L2) of Proposition~14.2\,(b) is required only in Propositions~28.1 and 28.2.

\paragraph{(G1.T-8)}
Propositions~20.1--20.3.
Proposition~20.1, which is the form of \cite[Proposition~5.64]{QEDI} for the present operator, holds with the constant $4b_{+}+16\,C^{F}\cdot c\left(\gamma^{F}\right)$.
Hence $C_0$ is given by its formula in the line (C6), which is that of \cite[Definition~5.97]{QEDI}:
\[
C_0=\max\left\{C_P\cdot c_{\mathrm{bd}}\left(\gamma_P\right),\ C^{F},\ 4b_{+}+16\,C^{F}\cdot c\left(\gamma^{F}\right)\right\}.
\]

\paragraph{(G1.T-9)}
Propositions~30.1--30.3.
There is a constant $C_6=C_6(\rho)$ such that the function
\[
\varepsilon_5(x)=C_6\cdot\log\left(\left(1-\exp\left(-\kappa_{(0)}\cdot x\right)\right)^{-1}\right),
\]
where $\kappa_{(0)}=\min\left\{\frac{\theta_{0}}{2},\ \frac{c_{\mathrm{el}}\cdot\theta_{0}}{12\,K_{\mathrm{Str}}}\right\}$ is the constant of the line (C0r) of Section~I.8.2, bounds the quantities that \cite{QEDI} bounds by
\[
312\,\log\left(\left(1-\exp\left(-\frac{x}{4}\right)\right)^{-1}\right),
\]
the function $\varepsilon_5$ of \cite[Section~5.6.2]{QEDI}.
The numbers $\gamma_{\Sigma}>0$ and $C_{\Sigma}<\infty$ depend only on $(b_{-},b_{+},\rho)$;
they depend neither on $\ell$, nor on $\mathfrak{N}$, nor on $\theta$, nor on $\mathsf{b}\in[b_{-},b_{+}]$.

\paragraph{(G1.T-10)}
Propositions~33.1--33.4.
These Propositions, which correspond to \cite[Lemmas~20.14, 20.15, 20.17 and~20.18]{QEDI}, hold with
\[
C_K=226\,\mathbf{S}^2,\qquad C_K'=1328\,\mathbf{S}^2,\qquad C_H=203\,M\cdot\left(C^{F}\right)^2,\qquad c_{\tau}=\min\left\{\frac{1}{8},\ \frac{\mu_-}{48},\ \frac{1}{292\,\mathbf{S}}\right\}.
\]
The other constants of \cite[Section~20.1.7]{QEDI} are given by the formulas stated there.

\paragraph{(G1.T-11)}
Propositions~32.1--32.3, 35.3\,(c1), 35.10, 36, 37.1, 37.2, 43.1 and~43.2.
The constants of the line (C0r), namely $\mathrm{Csy}_{\mathrm{ak}}$, $\mathrm{Cpv}_0$, $C_E$, and those of the line (C10), namely all the others below, exist, and each depends only on the quantities that the table of Proposition~14.2\,(a) names for it:
\begin{itemize}
\item $D_\beta$, $C_{\beta}$, $\mathbf{D}$;
\item $C_{1}^{\mathrm{N2}}$, $C_{2}^{\mathrm{N2}}$, which depend on $\rho$, $\bar{m}$ and the probe, but not on its position;
\item $\mathrm{Csy}_{\mathrm{ak}}$, $\mathrm{Cpv}_0$, $C_E$, which depend on $\rho$;
\item $\mathrm{Cah}_i$, $\mathrm{Cpv}_1$, $\mathrm{Cpv}_2$, $\mathrm{epv}$, $\mathrm{Cpv}_\beta$, $C_{\mathrm{ct}}$, which depend on $\rho$, $\bar{m}$ and the probe;
\item $e_{\sharp}>0$, with the conditions (g1)--(g9) of Proposition~14.2\,(d) for $e\le e_{\sharp}$;
\item $C^{G}_{\beta}$;
\item $C_{\mathrm{mat}}$;
\item $s_{B}$, $C_{3\mathrm{p}}$, $C_{\Gamma}^{(16)}$, $C_{U}$;
\item $C_{\bar{\kappa}}$, $C'_{\bar{\kappa}}$;
\item $e_2$.
\end{itemize}

\paragraph{(G1.T-12)}
All Propositions; the rule on the conditions.
The proofs of all Propositions follow the rule that is stated with the conditions of Proposition~14.2\,(b) and that is a convention for these proofs (Chapter~14b, end of Step~9):
apart from the ranges of Proposition~14.2\,(a) and the conditions of Proposition~14.2\,(b), no proof requires a condition on $\rho$, $\alpha$, $\kappa$, $L$ or $M$, and no condition bounds $M$ from above;
every further condition that a proof imposes requires that $e_0$ (resp.\ $e_1$, $e_{\sharp}$, $e_2$) be small, depending on constants chosen earlier, and it stands in the list of Proposition~14.2.
The rule is discharged in the sections ``Conditions on the constants'' of the chapters, which name, for every inequality between the constants used in a proof, the row of Proposition~14.2\,(a), (b), (d) from which it comes or the line of Proposition~14.2\,(c), (f) under which it stands (Section~IV.1, paragraph~E).
The condition (Pb) of Proposition~14.2\,(d) is of another type: it restricts the choice of the probe, of $r_\beta$ and of $x_0$, and it belongs to the data of Proposition~43.2 and of Theorem~44.

\subsubsection*{G2.10.1. The forms needed for Propositions 15--20.3 and 48.1--48.6}

\paragraph{(G2.T1)}
From the Propositions~2.m, in the form of \cite[Proposition~5.42]{QEDI} for the present operator.
Data (G2.S-3).
For every $A\in\tilde{G}_k$ the operator
\[
F_{\infty}(A)=D_A+\mu+b_k\cdot J\cdot P_k(A)\cdot J^{*}
\]
is bijective on $\mathfrak{S}^{\mathrm{ch}}(\eta_k,\ell)$, and $A\mapsto S_k(A)$ is analytic on $\tilde{G}_k$.

\paragraph{(G2.T2)}
From Proposition~5: the bound on $M_k$, in the form of \cite[Proposition~5.44]{QEDI}.
Data (G2.S-3).
For $A\in\tilde{G}_k$ and unit sites $y$, $y'$,
\[
\left|M_k(A)(y,y')\right|_c\le C^{F}\cdot\exp\left(-\gamma^{F}\cdot d(y,y')\right).
\]

\paragraph{(G2.T3)}
From Proposition~6: the first clause of the form of \cite[Proposition~5.45]{QEDI}.
Data (G2.S-3).
For $A\in\tilde{G}_k$ and $v\in\mathfrak{S}(1,\ell)$ one has $\left|H_k(A)v\right|_{\infty}\le C^{F}\cdot|v|_{\infty}$ and
\[
\sup_{(x,x')\in\mathrm{Pz}_k}\left|\left(\delta_{\alpha,A}H_k(A)v\right)(x,x')\right|_c\le C^{F}\cdot|v|_{\infty},
\]
the supremum extending over all pairs of $\mathrm{Pz}_k$, the pairs inside a cluster ($d_X=\eta_k$, no phase) and the long pairs included.
The same two bounds hold for $\bar{H}_k(A)$, with $\delta_{\alpha,-A}$ in place of $\delta_{\alpha,A}$.
All the norms are formed with $|\cdot|_c$.

\paragraph{(G2.T4)}
From Proposition~13\,(d0), in the form of \cite[Proposition~5.48]{QEDI}.
Data: a level $j\ge0$, $\ell_j\in\mathrm{Cl}$, $\mu_j\in(0,1]$, $e_j\in(0,e_0]$, and a family (p5) with the properties W1, W2, W7$^{\mathrm{ch}}$, W12--W14 for $\mathrm{par}$.
For all $A\in\tilde{G}_j$,
\[
\delta Z_j(A)=\exp\left\{c^{\mathrm{d}}_{j}\cdot n_j+\sum_{X\in D_j}\acute{E}_j(X,A)\right\}\quad\left(c^{\mathrm{d}}_{0}=0\right),
\qquad
\frac{\delta Z_j(A)}{\delta Z_j(0)}=\exp\left\{\sum_{X\in D_j}E_j^{\mathrm{det}}(X,A)\right\}.
\]

\paragraph{(G2.T5)}
From Proposition~13\,(d5).
In the setting of one step, $L E_k^{\mathrm{det}}$ lies in $K_{k+1}$, and $L\,{\delta E^{\mathrm{det}}}_k$ lies in $K^{x}_{k+1}$.

\paragraph{(G2.T6)}
From Proposition~24.1\,(a), the form of \cite[Theorem~5.71\,(a)]{QEDI}.
For every $(m,\mathscr{E})\in\mathrm{Ball}^{x}_{k}(4)$, with $\mathscr{E}=0$ if $k=0$, the activity $E^{*,x}_k(m,\mathscr{E})$ lies in $K^{x,C}_{k+1}$;
for real $m$ and $\mathscr{E}\in K^{x}_k$ it lies in $K^{x}_{k+1}$.
Moreover $0\le\delta\varepsilon^{0}_{k}$.

\paragraph{(G2.T7)}
From Proposition~24.1\,(b), the form of \cite[Theorem~5.71\,(b)]{QEDI}.
For every $(m,\mathscr{E})\in\mathrm{Ball}^{x}_{k}(4)$,
\[
E^{*,x}_k(m,\mathscr{E})|_0=E^{*}_k\left(m,\mathscr{E}|_0\right),
\]
and $E^{*}_k$ is defined on $\mathrm{Ball}_{k}(4)$, with values in $K^{C}_{k+1}$.

\paragraph{(G2.T8)}
From \cite[Proposition~19.1\,(c), (e), (f)]{QEDI}, the clause (f) for $N\ge N_{\min}(R,e)$;
and from (G4.Sc6) of Section~G4.0.2, that is, Proposition~14.2\,(c), (f) together with the second assertion of Proposition~14.3\,(f), of which the row (e1) of Proposition~14.2\,(b) is needed, at $v=v_K$ and at $v=v_{*}\ge v_K$.
These stand in the place of \cite[Proposition~19.2]{QEDI}, and they give the hypotheses of \cite[Proposition~19.1]{QEDI}:
\[
v_K\ge8,\qquad\left(1+\frac{\mathrm{lg}}{v_K}\right)^{r}\le\left(1+\frac{3\,\mathrm{lg}}{v_K}\right)^{r}<L,
\]
and the condition (2.1) of Section~G4.0.3, $r_{*}>31L+16$;
the number $e_1$ is that of Proposition~14.2\,(a), line (C9).
Under (G2.S-4) one has $1\le k_{*}(R)\le K$;
for $j\le k_{*}(R)$,
\[
\ell_j\in\mathrm{Cl},\qquad e_j\le e_K\le e_1\le e_0<1,\qquad 0<\mu_j\le\mu_K\le1;
\]
and for $j<k_{*}(R)$ moreover
\[
\frac{\ell_j}{L}=\ell_{j+1}\in\mathrm{Cl},\qquad \ell_j\ge M\cdot r_j,\qquad \mu_j\le\frac{1}{L}.
\]

\paragraph{(G2.T9)}
From Part~I, facts that Part~I states together with their argument, or together with the statement from which they follow:
(I.1.1), (I.2.4), (I.s1)--(I.s5), (I.c1)--(I.c6), (I.Sch), (I.4.1), (I.m2)--(I.m4), (I.m6), (I.m7), (I.O2), (I.O6), (I.O9), (I.J1)--(I.J4), (I.Q1), (I.Q2), (I.ke1)--(I.ke3);
Section~I.1.1 (faces; the sentences on lifts);
Section~I.1.2 (the $s_i$ lie in $O_{\mathrm{cub}}$, which has $48$ elements);
Section~I.1.4 (the sets $Z\times O_{\mathrm{cub}}$ form a partition of $\mathbb{L}^{\mathrm{f}}_k$, and $\Delta_y\times O_{\mathrm{cub}}$ has $48\,\eta_k^{-3}$ sites);
Section~I.1.9 (a translation of $\tilde{U}$ by $\ell\cdot w$ multiplies every generator by $(-1)^{w_1+w_2+w_3}$);
Section~I.3.4 ($b_{-}\le b_k\le b_{+}$);
Section~I.3.7 (for two of its sentences Part~I names the Propositions from which they follow: Proposition~18.3\,(e) for the division by $L$ of admissible pairs, and Proposition~16.1\,(d) for what the definition of $S_L$ requires of $\ell/L$);
Section~I.5.4 (the assertion on $\mathscr{E}|_0$, which Part~I states with Proposition~48.1\,(a), the statement that stands behind it; see Section~IV.1, paragraph ``Sentences of Part~I behind which a statement stands'', (2)).
The relations (I.B1)--(I.B4) are proved in Proposition~15\,(e), and the assertions (I.a1), (I.a2), (I.a4) in Proposition~16.1\,(a2), (d).

\paragraph{(G2.T10)}
From \cite{QEDI}, in the form in which they are stated there:
\begin{itemize}
\item the conventions of \cite[Section~13.1.7]{QEDI} that define $d_M$ by nets and the set $\bar{X}^{L}$, and the convention there on the value of a cut-off function at a bond;
\item the clauses (t0), (t1), (t2), (t5), (t7), (t8), (t10) of \cite[Proposition~5.58]{QEDI}, that is, of \cite[Theorem~13.2]{QEDI};
\item \cite[Lemmas~13.9 and~13.10]{QEDI}, with the properties (f1)--(f3) of the cut-off functions of \cite[Section~13.6]{QEDI};
and the two facts of linear algebra established at the beginning of the proof of \cite[Theorem~13.18]{QEDI}: an operator on a finite-dimensional complex Hilbert space whose real part is bounded below by a number $\vartheta>0$ is invertible, its inverse has norm at most $\vartheta^{-1}$ and a non-negative real part; and the criterion printed there for the invertibility of the sum of an invertible operator and a positive multiple of a product of two linear maps, with the formulas for the inverse;
\item \cite[Lemma~17.3]{QEDI} (the domains are open; $\operatorname{Im}A=\operatorname{Im}A_1$), \cite[Lemma~17.4]{QEDI}, which is \cite[Lemma~5.17]{QEDI}, and \cite[Lemma~17.6]{QEDI}, which is \cite[Lemma~5.19]{QEDI};
\item from \cite[Section~5]{QEDI}: the pairing of cochains and $\hat{\mu}$ of \cite[Section~5.2.5]{QEDI}, and the meaning of the word ``reads'' fixed in \cite[Definition~5.5\,(1)]{QEDI}.
\end{itemize}

\subsubsection*{G3.0.4. The forms needed for Propositions 21.1--25.3}

\paragraph{(G3.T-2a)}
Proposition~2.8\,(a), (b), the form of \cite[Theorem~16.16\,(a), (b)]{QEDI} for the present operator.
For every $k\ge0$, $\ell\in\mathrm{Cl}$, $\mu\in(0,1]$, $e_k\in(0,e_0]$, $y\in Y_{k}$ and $A\in\tilde{G}_k$, the operator $F_{y}(A)$ is bijective on $\mathfrak{S}^{\mathrm{ch}}(\eta_k,\ell)$, and $A\mapsto S_{y}(A)$ is analytic on $\tilde{G}_k$.
It is needed at $y=\infty$ and at $y=0$ for $k\ge1$, and at $y=0$ for $k=0$.

\paragraph{(G3.T-9)}
Proposition~9\,(d), the form of \cite[Proposition~16.27\,(d)]{QEDI}.
Let $k\ge1$, $\ell\in\mathrm{Cl}$, $\mu\in(0,1]$, $e_k\in(0,e_0]$ and $B\in\tilde{R}_k+\mathfrak{X}_k$.
For every polymer $\mathbf{Y}$ and $i=1,2$, the rows of $T^{(i)}(B,\mathbf{Y})$ and of $\bar{T}^{(i)}(B,\mathbf{Y})$ at the unit sites of $\mathbf{Y}$ satisfy $|\cdot|_{1,\infty}\le C_{T}$.
Here $|\cdot|_{1,\infty}$ is the norm of (I.ke3) and Section~I.5.2 (for $i=2$: the sum of the moduli of the coefficients of the values $F(x,x^{(3)})_{\beta}$), and $C_{T}$ is the constant of the line (C4) of Section~I.8.2.
The operators $T^{(i)}$ of Proposition~9 are those of the display (23.1), which is (I.5.4).

\paragraph{(G3.T-13)}
Proposition~13\,(d0), (d5), (d6), the form of \cite[Proposition~5.48]{QEDI}.
Assume the setting (ST1)$^{\mathrm{ch}}$, and let the family (p5) of the given system have the properties W1, W2, W7 (at $k=0$: the kernel bound with the factor $(1+y)^{-1}$, and Borel dependence on $y$), W12, W13, and W14 for $\mathbb{G}^{L}_{k,3}$, for $\mathrm{par}$.
Then the following hold.
\begin{itemize}
\item[(d0)] For $B\in\tilde{G}_k$ the integrals in $\acute{E}_k(X,B)$ converge absolutely, and
\[
\delta Z_k(B)=\exp\left\{c^{\mathrm{d}}_{k}\cdot n_k+\sum_{X}\acute{E}_k(X,B)\right\},\qquad
\frac{\delta Z_k(B)}{\delta Z_k(0)}=\exp\left\{\sum_{X\in D_k}E_k^{\mathrm{det}}(X,B)\right\}.
\]
\item[(d5)] $L E_k^{\mathrm{det}}$ lies in $K_{k+1}$, and $L\,{\delta E^{\mathrm{det}}}_k$ lies in $K^{x}_{k+1}$, on the side $\ell/L$.
\item[(d6)] $\left\|L E_k^{\mathrm{det}}\right\|_{k+1}\le e_{k+1}^{\frac{1}{4}-2\varepsilon}$ and $\left\|L\,{\delta E^{\mathrm{det}}}_k\right\|^{x}_{k+1}\le e_{k+1}^{\frac{1}{4}-2\varepsilon}$.
\end{itemize}

\paragraph{(G3.T-15)}
Proposition~15\,(a)--(d), (f)--(k), with Section~I.5.2.
The rule (N3)$^{\mathrm{w}}$, which is (I.5.2), with $\mathsf{w}_W=\max\nu^{-\frac{1}{2}}$;
the rules (N0), (N1), (N2) (with $|\cdot|_{1,\infty}$ the norm for the sup norms, (I.ke3)), (N4), (N4$'$), (N5) and (N6), for arbitrary finite index sets with weights.

\paragraph{(G3.T-16)}
Proposition~16.3\,(a)--(d), the extraction, in the form of \cite[Proposition~5.59]{QEDI}.
In the setting (G3.S22), with $k\ge1$, let $\mathscr{E}\in K^{x,C}_k$.
\begin{itemize}
\item[(i)] $\varepsilon(\mathscr{E})$ and $m(\mathscr{E})$ do not depend on $\square$, and $m(\mathscr{E})$ is a multiple of the unit matrix $1$;
both are complex linear in $\mathscr{E}$, and real for $\mathscr{E}\in K^{x}_k$;
$\varepsilon(\mathscr{E})=\varepsilon\left(\mathscr{E}|_0\right)$ and $m(\mathscr{E})=m\left(\mathscr{E}|_0\right)$.
\item[(ii)] $\left|\varepsilon(\mathscr{E})\right|\le c_{\mathrm{fl}}\cdot\|\mathscr{E}\|^{x}_k$ and $\left|m(\mathscr{E})\right|\le c_{\mathrm{fl}}\cdot e_k^{\frac{1}{2}}\cdot\|\mathscr{E}\|^{x}_k$.
\item[(iii)] For all $(\mathscr{A},\mathscr{C})\in\tilde{R}_k\times\mathfrak{X}_k$, in $\mathrm{Gr}[\psi,\chi]$,
\[
\sum_X\mathscr{E}(X;\mathscr{A},\mathscr{C},\,\cdot\,)=-\varepsilon(\mathscr{E})\cdot n_k-m(\mathscr{E})\cdot\langle\bar{\psi},\psi\rangle+\sum_X\left(R\mathscr{E}\right)(X;\mathscr{A},\mathscr{C},\,\cdot\,).
\]
\item[(iv)] $R$ is linear and maps $K^{x,C}_k$ and $K^{x}_k$ into themselves;
\[
\left\|\left(R\mathscr{E}\right)(X;\mathscr{A},\mathscr{C})\right\|_{h_k}\le\left\|\mathscr{E}(X;\mathscr{A},\mathscr{C})\right\|_{h_k}+\left\|\mathscr{E}(X;0,0)\right\|_{h_k};
\]
and $\left(R\mathscr{E}\right)|_0=R\left(\mathscr{E}|_0\right)$.
\end{itemize}

\paragraph{(G3.T-17)}
Proposition~17.2\,(a), the adjustment, in the form of \cite[Proposition~5.60\,(a)]{QEDI}.
Let $E\in K^{x,C}_k$, let $X$ be a small polymer, let $(\mathscr{A},\mathscr{C})\in\tilde{R}_k\times\mathfrak{X}_k$ and $\mathscr{T}=\mathscr{A}+\mathscr{C}$.
Then $E^{\mathrm{nat}}(X;\mathscr{A},\mathscr{C},\,\cdot\,)$ and $E(X;\mathscr{A},\mathscr{C},\,\cdot\,)$ have the same image under every homomorphism $\mathrm{ev}$ of $\mathrm{Gr}[\psi,\chi]$ with
\[
\mathrm{ev}(\chi)=\delta_{\alpha,\mathscr{T}}\,\mathrm{ev}(\psi),\qquad\mathrm{ev}(\bar{\chi})=\delta_{\alpha,-\mathscr{T}}\,\mathrm{ev}(\bar{\psi}).
\]

\paragraph{(G3.T-18)}
Proposition~18.5\,(a)--(c), the contraction, in the form of \cite[Proposition~5.62]{QEDI}.
In the setting (G3.S22), $L R$ maps $K^{x,C}_k$ into $K^{x,C}_{k+1}$ and $K^{x}_k$ into $K^{x}_{k+1}$;
\[
\left\|L R\mathscr{E}\right\|^{x}_{k+1}\le c_{\mathrm{fl}}\cdot L^{-\frac{1}{4}+2\varepsilon}\cdot\|\mathscr{E}\|^{x}_k ;
\]
and $\left(L R\mathscr{E}\right)|_0=L R\left(\mathscr{E}|_0\right)$.

\paragraph{(G3.T-48)}
The Definition in Proposition~48.1, with the clauses (k1)$^{\mathrm{o}}$--(k7)$^{\mathrm{o}}$, and Propositions~48.2\,(a)--(d), 48.3\,(a), 48.5\,(a)--(c).
The assertions (G3.T-16), (G3.T-17) and (G3.T-18) hold for the families without slot, with $K_k^{C}$, $K_k$, $\|\cdot\|_k$ and $E(X;0)$, and with the same constants.

\paragraph{(G3.T-W)}
The hypothesis on the primitives in Propositions~24.1--24.3 consists of the properties W1--W8, W10--W17 of \cite[Section~5.4.2]{QEDI} for $\mathrm{par}$, read by the rules of Section~I.5.8;
they are stated in full in Section~G1.11 and in the table of Section~G4.0.4.
These properties are established in Proposition~7 ($k\ge1$) and Proposition~8 ($k=0$) for the fermionic families (p3)--(p5),
in Propositions~26.1--26.6, with the parameters $\mathrm{par}$, for the systems of primitives formed from the families (p1)--(p5) and for their projections,
and in \cite[Theorem~15.35 and Proposition~15.36]{QEDI} for the photon families (p1), (p2).
At the places where a fermion index occurs they are needed in the following form.
\begin{description}
\item[W1.] $H_k(1)=H_k$, $C_k^{1/2}(1)=C_k^{1/2}$, $H_k(1,A)=H_k(A)$, $\Gamma_k(1,A)=\Gamma_k(A)$, and $\sum_XM_k(X,A)=M_k(A)$ ($k\ge1$);
at $k=0$: $H_0(1)=1$, $H_0(1,A)=1$ on $\mathfrak{S}^{\mathrm{ch}}(1,\ell)$, and $\Gamma_0(1,A)=\Gamma_0(A)$.
\item[W5.] For $k\ge1$:
\[
\left|H_k(s'',A)v\right|_{\infty}\le C_w\cdot|v|_{\infty},\qquad\left|\left(\delta_{\alpha,A}H_k(s'',A)v\right)(x,x')\right|_c\le C_w\cdot|v|_{\infty}
\]
at all short pairs, $d_X<1$, the pairs inside a cluster included.
For $k=0$:
\[
\left|H_0(s'',A)f\right|_{\infty}\le4\exp\left(2\gamma_{\mathrm{w}}\right)\cdot M^{27\alpha_0}\cdot|f|_{\infty}\le C_w\cdot|f|_{\infty}.
\]
\item[W6.] For $k\ge1$:
\[
\left|\Gamma_k(s'',A)(y,y')\right|_c\le C_w\cdot\exp\left(-\gamma_{\mathrm{w}}\cdot d(y,y')\right),\qquad
\left|M_k(X,A)(y,y')\right|_c\le C_w\cdot\exp\left(-\gamma_{\mathrm{w}}\cdot d(y,y')-\kappa_{\mathrm{w}}\cdot d_M(X)\right).
\]
For $k=0$: $\left|\Gamma_0(s'',A)(x,x')\right|_c\le C_w\cdot\exp\left(-\gamma_{\mathrm{w}}\cdot d(x,x')\right)$ for the kernel (I.1.4).
Consequence: $\left|\Gamma_k(s'',A)\right|_{1,\infty}\le C_w\cdot c_{\gamma}$ for every $k\ge0$.
In the proof of Proposition~24.1, W6 enters for $\Gamma_k$ only through this row sum, as in \cite[Lemma~17.11\,(q4)]{QEDI}, and for $M_k$ through the clause (q2) of that lemma.
\item[W8, W10, W11, W12.] A fermion site $x$ is an index of the kind $\pi_k$, with $\operatorname{supp}x=\{c\}$ and $K_x=K_c$.
\item[W13.] For a real function $\boldsymbol{\lambda}$ on the sites of $T_k$:
\[
H_k(s'',A-d\boldsymbol{\lambda})=\exp\left(\mathrm{i}\,e_k\cdot\boldsymbol{\lambda}\right)\cdot H_k(s'',A)\cdot\exp\left(-\mathrm{i}\,e_k\cdot\boldsymbol{\lambda}^{\mathrm{u}}\right);
\]
\[
O(\,\cdot\,,A-d\boldsymbol{\lambda})=\exp\left(\mathrm{i}\,e_k\cdot\boldsymbol{\lambda}^{\mathrm{u}}\right)\cdot O(\,\cdot\,,A)\cdot\exp\left(-\mathrm{i}\,e_k\cdot\boldsymbol{\lambda}^{\mathrm{u}}\right)\quad\text{for (p4)};
\]
at $k=0$ all phases are read at $c$.
\item[W14.] $\mathrm{Op}(\tilde{g}s,\tilde{g}A)\cdot\tilde{g}=\tilde{g}\cdot\mathrm{Op}(s,A)$ and $\mathrm{Op}(\tilde{g}X,\tilde{g}A)\cdot\tilde{g}=\tilde{g}\cdot\mathrm{Op}(X,A)$ for $\tilde{g}\in\mathbb{G}_k^{L}$, with the action (I.4.2);
here $\mathrm{Op}(s,A)$ is $H_k(s'',A)$, $\Gamma_k(s'',A)$ or a member of the photon families (p1), (p2), and $\mathrm{Op}(X,A)$ is $M_k(X,A)$.
For the family (p5), $S_{k,y}(X,A)$, the relation is required for $\tilde{g}\in\mathbb{G}^{L}_{k,3}$ only, and it is needed only through (G3.T-13).
\item[W15.] $\Gamma_k(s'',A)^{T}=\mathsf{C}^{-1}\cdot\Gamma_k(s'',-A)\cdot\mathsf{C}$ and $M_k(X,A)^{T}=\mathsf{C}^{-1}\cdot M_k(X,-A)\cdot\mathsf{C}$, the transposes being formed for the pairings of Section~I.1.6 (at $k=0$: $\langle\cdot,\cdot\rangle^{\mathrm{ch}}$).
\item[W16.] $\overline{O\left(\overline{s},\overline{A}\right)}=\varkappa^{-1}\cdot O(s,-A)\cdot\varkappa$ for (p3), (p4); and (p1), (p2) are real for real $s$.
\item[W2, W3, W4.] As stated.
W7 is needed only through (G3.T-13).
\end{description}

From \cite{QEDI} the following are needed in the form in which they are stated there:
(LA2) and (LA3) of \cite[Section~14.1.4]{QEDI}, \cite[Proposition~14.8]{QEDI}, and \cite[Theorem~14.15]{QEDI} (see \cite[Proposition~5.52]{QEDI});
\cite[Propositions~15.21, 15.26 and~15.30\,(ii)]{QEDI}, that is, \cite[Propositions~5.35, 5.36 and~5.38\,(ii)]{QEDI};
\cite[Theorem~13.2]{QEDI}, that is, \cite[Proposition~5.58]{QEDI};
\cite[Lemmas~17.3--17.6]{QEDI} (the last three of which contain \cite[Lemmas~5.17--5.19]{QEDI}), \cite[Lemmas~17.12 and~17.13]{QEDI}, and (g0)--(g4) of \cite[Section~17.4.1]{QEDI}.
The proof of \cite[Lemma~17.17]{QEDI}, a statement on $\mathrm{Gr}[\psi,\chi,W]$ and $\mathrm{d}\mu_I(W)$, rests on \cite[Proposition~5.57]{QEDI} and on \cite[Corollary~17.15]{QEDI}, in whose place stand here Proposition~15 and Proposition~24.3.
That statement is stated and proved for all $k\ge0$ as Lemma~5.C of Chapter~24, in the proof of Proposition~24.1\,(a).

\subsubsection*{G4.A. The forms needed for Propositions 26.1--32.3 and 49}

\subsubsection*{G4.A.1. From Propositions 1.1--14.4}

\paragraph{Proposition 1.2.}
Let $0\le\mathbf{m}\le1$; lattice units are used.
\begin{itemize}
\item[(S1)] The heavy block $D^{\mathrm{hh}}(p)$ of the symbol (I.2.6) at $\boldsymbol{\mu}=\mathbf{m}$ is invertible on $\left\{|\operatorname{Im}p_\mu|\le\theta_{0}\right\}$, and the norm of its inverse is at most $2/\hat{g}$.
\item[(S3)] $\mathrm{Dsc}_{\mathbf{m}}(p)$ is holomorphic and $2\pi\mathbb{Z}^3$-periodic on $\left\{|\operatorname{Im}p_\mu|\le\theta_{0}\right\}$, and $\left|\mathrm{Dsc}_{\mathbf{m}}(p)\right|_2\le K_{\mathrm{Str}}$ there, uniformly in $\mathbf{m}\in[0,1]$.
\item[(S7)] Let $\kappa_{(0)}=\min\left\{\frac{\theta_{0}}{2},\ \frac{c_{\mathrm{el}}\cdot\theta_{0}}{12\,K_{\mathrm{Str}}}\right\}$ be the constant of the line (C0r) of Section~I.8.2.
For $\mathbf{m}>0$ and $p=p'+\mathrm{i}\,p''$ with $p'\in[-\pi,\pi]^3$ and $|p''_\mu|\le\kappa_{(0)}\cdot\mathbf{m}$, the Schur symbol $\mathrm{Dsc}_{\mathbf{m}}(p)$ is invertible, and
\[
\left|\partial^{\mathrm{ak}}\mathrm{Dsc}_{\mathbf{m}}(p)^{-1}\right|\le A'_{\mathrm{ak}}\cdot\left(|p'|_2^2+\mathbf{m}^2\right)^{-\frac{1+|\mathrm{ak}|}{2}},\qquad A'_{\mathrm{ak}}=A'_{\mathrm{ak}}(\rho).
\]
\item[(S8)] $\det D^{(p)}=\det D^{\mathrm{hh}}(p)\cdot\det\mathrm{Dsc}_{\mathbf{m}}(p)$.
\item[(S9)] Both factors are positive for real $p$ and $\mathbf{m}>0$.
\end{itemize}
Further, \cite[Lemma~12.4]{QEDI} on the Fourier coefficients of symbols of homogeneous type, which involves no operator.

\paragraph{Proposition 14.1\,(ii)}
(Proposition~7 and Proposition~8\,(a), (b); the properties W1, W2, W5--W17$^{\mathrm{ch}}$ of the table of Section~G1.11).
For all data (G4.SF), the families (p3), (p4), (p5) of level $k\ge1$ (Proposition~7), and $\Gamma_0(s'',A)$, $S_{0,y}(X,A)$ on $\mathfrak{S}^{\mathrm{ch}}(1,\ell)$ (Proposition~8), have the properties of the table of Section~G4.0.4 for $\left(\alpha_0,\gamma^{F},\kappa_{\mathrm{w}}+\kappa_0;C^{F}\right)$, where $\kappa_0$ is the constant $54\log72$ of the line (C0) of Section~I.8.2.
The diagonal family $H_0(s'',A)$ has W1, W2, W8--W11, W13, W14, W16, and W5 with the constant
\[
M^{8\alpha_0}\le4\exp(2\gamma)\cdot M^{27\alpha_0}\qquad\text{for every }\gamma\ge0
\]
(Proposition~8\,(b)).
W14 holds for (p5) with $\mathbb{G}^{L}_{k,3}$;
W7 holds at $k=0$ with the factor $(1+y)^{-1}$ and with Borel dependence on $y$.

\paragraph{Proposition 14.1, its paragraph on the constants, and the first clause of Proposition 14.1\,(i)}
(the form of \cite[Proposition~5.42]{QEDI}; Proposition~2.8\,(a), (b) at $y=\infty$).
One has $0<\gamma^{F}\le\frac{1}{6L}$.
For data (G4.SF) with $k\ge1$ and every $A\in\tilde{G}_k$ the operator $D_A+\mu+b_k\cdot J\cdot P_k(A)\cdot J^{*}$ is bijective on $\mathfrak{S}^{\mathrm{ch}}(\eta_k,\ell)$;
hence $S_k(A)$, $H_k(A)$, $\bar{H}_k(A)$ and $D_k(A)$ exist (Section~I.3.5; Section~I.5.7, paragraph ``Settings'').

\paragraph{Proposition 11\,(b)}
(the form of \cite[Proposition~5.49]{QEDI}).
Let $\ell_2=\mathrm{Jr}\cdot\ell_1$ with $\mathrm{Jr}$ odd and $\ell_1,\ell_2\in\mathrm{Cl}$, let $\mu\in(0,1]$, $e_k\in(0,e_0]$ and $A\in\tilde{G}_k(T_1)$, the upper index $(i)$ marking the operator on the torus $T_i$ of side $\ell_i$.
Then
\[
O^{(2)}\left(\mathrm{pj}^{*}A\right)\cdot\mathrm{pj}^{*}=\mathrm{pj}^{*}\cdot O^{(1)}(A)
\]
for $O(A)=H_k(A)$, $\Gamma_k(A)$, $M_k(A)$, $S_{k,y}(A)$ with $y\ge0$ ($k\ge1$), and for $O(A)=\Gamma_0(A)$, $S_{0,y}(A)$ ($k=0$).

\paragraph{Proposition 13\,(d1), (d5), (d6)}
(the form of \cite[Proposition~5.48]{QEDI}).
Assume (G4.OS); in the data of Section~G1.1 the statement is made under (G1.SF+), and its clause (d1) under (G1.SF).
For every family (p5) with W1, W2, W7, W12--W14$^{\mathrm{ch}}$ for $\mathrm{par}$ and every $k\ge0$:
\begin{itemize}
\item[(d1)] $\acute{E}_k(X,\cdot\,)$ and $E^{d}_k(X,\cdot\,)$ are analytic functions on $\tilde{G}_k$;
\item[(d5)] $L E_k^{\mathrm{det}}$ lies in $K_{k+1}$, and $L\,{\delta E^{\mathrm{det}}}_k$ lies in $K^{x}_{k+1}$;
\item[(d6)] their norms are at most $e_{k+1}^{\frac{1}{4}-2\varepsilon}$.
\end{itemize}

\paragraph{Proposition 14.2.}
The constants of Section~I.8.2, with:
$e_1$ defined by the conditions (e1)--(e8), which are those of \cite[Section~19.1.2]{QEDI} (Section~G4.0.3);
$e_1\le e_0$;
$M\ge M^{F}$ at $\kappa'=\kappa_{\mathrm{w}}+\kappa_0$, with the constant $\kappa_0=54\log72$ of the line (C0);
\[
C_0=\max\left\{C_P\cdot c_{\mathrm{bd}}\left(\gamma_P\right),\ C^{F},\ 4b_{+}+16\,C^{F}\cdot c\left(\gamma^{F}\right)\right\}\ge1;
\]
$c_{\kappa}\le K_0$;
and $\gamma_{\Sigma}$, $C_{\Sigma}$ in the line (C6), chosen before $e_0$ and $e_1$.
The conditions on the constants that the proofs of Propositions~26.1--32.3 and 49 require are listed in Section~G4.A.6.
Further:
Proposition~14.3\,(f) (the conditions (e1)--(e8) hold for all $v\ge-\log e_1$);
Proposition~14.2\,(f) ($C_P\le C_w$, $\gamma_{\mathrm{w}}\le\gamma_P$, $C_0\ge1$, $e_1\le e_0<1$, $L\ge2400^8$, $0<\alpha<1$);
Proposition~14.2\,(e);
and Proposition~14.4\,(a), the rule (I.U4)$^{\mathrm{ch}}$ in the form of a statement.

\subsubsection*{G4.A.2. From Propositions 15--20.3 and 48.1--48.6}

\paragraph{Proposition 15\,(a), (b), (d), (e), (k), (l), (n); Section I.5.2.}
The rules (N0) (the antisymmetric kernels are determined by the element), (N1), (N2), (N6) and (S6c)\,(i)--(iii);
the relations (I.B1)--(I.B4).
At $k\ge1$ they are needed with $\nu=1$; at $k=0$ with $\nu=\frac{1}{48}$, where
\[
\int\mathrm{d}\mu_I\,W_{\beta}(x)\,\bar{W}_{\beta'}(x')=48\cdot\mathbf{1}\left[x=x',\ \beta=\beta'\right],
\]
the same number on every torus.

\paragraph{Proposition 16.3\,(a), (d), (f)}
(the extraction, in the form of \cite[Proposition~5.59]{QEDI}).
For $\mathscr{E}\in K^{x,C}_k$: $\varepsilon(\mathscr{E})$ and $m(\mathscr{E})$ do not depend on $\square$;
$m(\mathscr{E})$ is a multiple of the unit matrix $1$;
both are linear, and real on $K^{x}_k$;
$\varepsilon(\mathscr{E})=\varepsilon\left(\mathscr{E}|_0\right)$ and $m(\mathscr{E})=m\left(\mathscr{E}|_0\right)$;
and $R$ maps $K^{x}_k$ into itself.
The quantities $\mathrm{al}_{0}$, $\mathrm{al}_{2}$ are formed with all arguments over one lift of the small polymer (Section~I.5.6, (I.AP)).

\paragraph{Propositions 17.1 and 17.2\,(c), (f)}
(the adjustment, in the form of \cite[Proposition~5.60]{QEDI}, with (I.SY6)$^{\mathrm{ch}}$).
$E^{\mathrm{nat}}$ has (I.k2);
$\boldsymbol{\Theta}_{X,\mathscr{T}}$ is the homomorphism (I.5.3);
all pairs of a chain are of kind (I) for $X$;
and the transports read $\mathscr{T}$ on the bonds contained in $X$.

\paragraph{Proposition 18.5\,(f)}
(the contraction, in the form of \cite[Proposition~5.62]{QEDI}).
In the sense of Section~I.1.9, for the sides $\ell$ with $\ell,\ell/L\in\mathrm{Cl}$:
$S_L$ is given on the generators indexed by $\mathbb{L}^{\mathrm{f}}_k$ and by $\mathrm{Pzt}_k$ by the one formula (G2.r5) of Section~G2.0.5, that is, by Section~I.3.7, of range $0$;
and
\[
\mathrm{cov}_{\ell/L}\circ S_L=S_L\circ\mathrm{cov}_{\ell}
\]
on the subalgebra of $\mathrm{Grc}_k(\tilde{U})$ generated by $\psi$, $\chi$, as soon as the sup-diameter of $\tilde{U}$ is smaller than $\ell-\eta_k$.

\paragraph{Proposition 19\,(a), (d)}
(the form of \cite[Proposition~5.65]{QEDI}).
Let the data (G4.SF) hold, and let $A$ be a real bond function on $T_k$.
If $k\ge1$, then by Proposition~19\,(d) the operator $D_k(A)$ is invertible, with
\[
D_k(A)^{-1}=b_k^{-1}+Q_k(A)\cdot J^{*}\cdot\left(D_A+\mu\right)^{-1}\cdot J\cdot Q_k^{T}(-A),
\]
and the operator norm of $D_k(A)^{-1}$ is at most $b_k^{-1}+\mu^{-1}$.
For every $k\ge0$, Proposition~19\,(a) gives $\operatorname{Re}\left(D_A+\mu\right)\ge\mu$;
hence $D_A+\mu$ is invertible and $\det\left(D_A+\mu\right)\ne0$.
For $k\ge1$ this provides the inverse $\left(D_A+\mu\right)^{-1}$ in the display above.
At $k=0$ one has $\eta_0=1$, the charge is $q$, the space is $\mathfrak{S}^{\mathrm{ch}}(1,\ell)$, and $D_0(A)=D_A+\mu$ (Section~I.3.9);
so Proposition~19\,(a) gives $\det D_0(0)\ne0$, which in the flow, where $\mu=\bar{m}\cdot\eta$, reads $\det\left(D_0+\bar{m}\cdot\eta\right)\ne0$.
This form at $k=0$ is the one needed in the proof of Proposition~29 (Chapter~29, (P2.1); Section~I.5.7, paragraph ``Settings'');
there Proposition~19\,(d) is needed at $A=0$ for $1\le k\le k_{*}$, and in Chapter~32, (32.10), it is needed at the level $k_{*}$ with the real bond function $A=c_\theta$.

\paragraph{Proposition 20.1}
(the terminal bounds, in the form of \cite[Proposition~5.64]{QEDI}).
Let (G4.SF) hold with $k\ge1$, and put
\[
\mathrm{aw}_1:=e_k^{-\frac{1}{4}},\qquad\mathrm{aw}_2:=C_0^{-1}\cdot e_k^{-\frac{1}{4}+\varepsilon}.
\]
\begin{itemize}
\item[(a)] For $A\in\tilde{G}_k$: $\left|\langle\bar{\Psi},D_k(A)\Psi\rangle\right|_{\mathrm{aw}_1}\le C_0\cdot e_k^{-\frac{1}{2}}\cdot\ell^3$, and the entries of $D_k(A)$ are analytic functions on $\tilde{G}_k$.
\item[(b)] $\left|\langle\bar{\psi},\psi\rangle\right|_{\mathrm{aw}_2}\le4\,e_k^{-\frac{1}{2}}\cdot\ell^3$ at $\psi=H_k(A)\Psi$, $\bar{\psi}=\bar{H}_k(A)\bar{\Psi}$.
\item[(c)] For $(\mathscr{A},\mathscr{C})\in\tilde{R}_k\times\mathfrak{X}_k$, the number $\left|\mathscr{E}\left(X;\mathscr{A},\mathscr{C},H_k^{\#}(\mathscr{T})\Psi\right)\right|_{\mathrm{aw}_2}$ is at most the norm of $\mathscr{E}(X;\mathscr{A},\mathscr{C})$ with the weights $h_{k}$.
\end{itemize}

\paragraph{Propositions 48.1\,(a), (d), 48.2\,(a), (b), 48.5\,(a), (b) and 48.6.}
For $\mathscr{E}\in K^{x}_k$ the family $\mathscr{E}|_0$ lies in $K_k$, and $\left\|\mathscr{E}|_0\right\|_k\le\|\mathscr{E}\|^{x}_k$, the assertion stated in Section~I.5.4.
The space $K_k$ is a real Banach space.
The forms without slot hold with the same constants:
$\varepsilon\left(E\right)$ and $m\left(E\right)$ are linear, and real on $K_k$, with
\[
\left|\varepsilon\left(E\right)\right|\le c_{\mathrm{fl}}\cdot\|E\|_k,\qquad\left|m\left(E\right)\right|\le c_{\mathrm{fl}}\cdot e_k^{\frac{1}{2}}\cdot\|E\|_k;
\]
$L R$ is a linear map of $K_k$ into $K_{k+1}$, of norm at most $c_{\mathrm{fl}}\cdot L^{-\frac{1}{4}+2\varepsilon}$;
and the relation $E_k=\mathscr{E}_k|_0$ holds along every trajectory.
Further:
the Definition in Proposition~48.1, with the clause (k2)$^{\mathrm{o}}$, and Proposition~48.1\,(b): the extension of a family that is constant in $\mathscr{C}$ has (I.k2) if and only if the family has (k2)$^{\mathrm{o}}$;
Proposition~48.1\,(a) at the level $k+1$: $\mathscr{E}|_0$ lies in $K_{k+1}$ for $\mathscr{E}\in K^{x}_{k+1}$;
Proposition~48.2\,(d), by which $R$ maps $K_k$ into itself, and Proposition~48.2\,(f);
Proposition~48.3\,(c), by which $E^{\mathrm{nat}}$ has (k2)$^{\mathrm{o}}$, and Proposition~48.3\,(f);
and Proposition~48.5\,(e).

\subsubsection*{G4.A.3. From the Propositions 21.1--25.3}

\begin{description}
\item[Proposition 22.2\,(c)]
For every real pair $(m_{\mathrm{in}},\varepsilon^{\mathrm{v}}_{\mathrm{in}})$, every $A_0\in\mathrm{Dom}_0$ and every $a\in\mathfrak{X}_0$ one has, wherever $y^{x}_0$ is defined,
\[
y^{x}_0[\varepsilon^{\mathrm{v}}_{\mathrm{in}},m_{\mathrm{in}},0](A_0;a)=\exp\left(-\varepsilon^{\mathrm{v}}_{\mathrm{in}}\cdot n_0\right)\cdot\frac{\det\left(D_{A_0+a}+\bar{m}\eta+m_{\mathrm{in}}\right)}{\det\left(D_0+\bar{m}\eta\right)}.
\]
The Berezin identity on which this formula rests holds for every complex bond function and for all complex $m$ and $\varepsilon^{\mathrm{v}}$.
The determinants are those of linear maps of $\mathfrak{S}^{\mathrm{ch}}(1;n)$, and $n_0=n^3$.
(Used in Proposition~29.)

\item[Propositions 23.1\,(b), (d) and 23.2\,(c)]
The displays (I.5.4) and (I.5.5) have the following locality.
Let $y$ be a unit site of $\mathbf{Y}$.
The rows of $T^{(1)}(B)$ and of $T^{(2)}(B)$ at $y$ contain $\psi$ at fermion sites over $\mathbf{Y}$ and $\chi$ at the pairs $(x,x^{(3)})$ with $c\in\Delta_y$ and $c_3\notin\mathbf{Y}$, and they read $B$ on bonds that have an end in $\Delta_y$.
The display (I.5.5) contains $\psi$ and $\bar{\psi}$ at the sites $x$ over $\square$, $\chi$ and $\bar{\chi}$ at the pairs $(x,x^{(3)})$, and $\mathscr{Z}(s')$ on the bonds $(c,c_3)$ with $c\in\square$.
(Used in Proposition~27.4\,(a-3), (b).)

\item[Propositions 24.1\,(a)--(e), (i) and 24.3]
These are Theorem~Z1$^{\mathrm{ch}}$; the step map is that of Definition~24.D.
They are used under (OS), which contains the setting (ST1)$^{\mathrm{ch}}$ of Section~G3.4.0 (see Section~I.1.8), at every level $k\ge0$, in the following forms.
Each form is named after the clause of Z1$^{\mathrm{ch}}$ to which it belongs: Z1.a--Z1.e and Z1.i are the forms of the clauses (a)--(e) and (i), that is, of Proposition~24.1\,(a)--(e) and~(i), and Z1.g is the form of the second half of the clause (g).
\begin{itemize}
\item[Z1.a]
For $(m,\mathscr{E})\in\mathrm{Ball}^{x}_k(4)$, $A\in\tilde{R}_k^{-}$ and $\mathscr{C}\in\mathfrak{X}_k$ the series (8.4) of Definition~24.D converge in norm, and $0\le\delta\varepsilon^{0}_{k}\le e_k^7$.
\item[Z1.b]
$E^{*}_k$ maps the real data of $\mathrm{Ball}_{k}(4)$ into $K_{k+1}$, with norm at most $e_k^{\frac{1}{4}-6\varepsilon}$; moreover $E^{*,x}_k(m,\mathscr{E})|_0=E^{*}_k\left(m,\mathscr{E}|_0\right)$.
\item[Z1.c]
Let $\varepsilon^{\mathrm{v}}$ and $m$ be real with $|m|\le w^{\mathrm{m}}_k$, let $\mathscr{E}\in K^{x}_{k}$ with $\|\mathscr{E}\|^{x}_k\le w^{\mathrm{E}}_k$ ($\mathscr{E}=0$ if $k=0$), let $\mathfrak{P}_k$ be an arbitrary family of complex functions on $\mathfrak{X}_k$, let $A'\in\mathrm{Dom}_{k+1}$, with arbitrary harmonic part, and let $\mathscr{C}'\in\mathfrak{X}_{k+1}$.
Put
\begin{gather*}
\left((\varepsilon^{\mathrm{v}})^{+},m^{+},\mathscr{E}^{+}\right)=\left(L^3\cdot\left(\varepsilon^{\mathrm{v}}+\varepsilon(\mathscr{E})+\delta\varepsilon^{0}_{k}\right),\ L\cdot\left(m+m(\mathscr{E})\right),\ L\,R\,\mathscr{E}+L\,{\delta E^{\mathrm{det}}}_k+E^{*,x}_k(m,\mathscr{E})\right),\\
\mathfrak{P}_{k+1}(Y,\mathscr{C}')=\sum_{X:\,\bar{X}^{L}=LY}\left(\mathfrak{P}_k(X,\mathscr{C}'_L)+E_k^{\mathrm{det}}(X,\mathscr{C}'_L)\right).
\end{gather*}
Then the identity (2.4x)$^{\mathrm{ch}}$ holds in the form
\[
y^{x}_{k+1}[(\varepsilon^{\mathrm{v}})^{+},m^{+},\mathscr{E}^{+}](A';\mathscr{C}')=\int_{\mathrm{Box}_{k}}\mathrm{d}\mu_I(\zeta)\,y^{x}_k[\varepsilon^{\mathrm{v}},m,\mathscr{E}]\left(A^{\min}_{k}(A')+\mathrm{Sq}_{k}\zeta;\mathscr{C}'_L\right).
\]
\item[Z1.d]
In the situation of Z1.c, $\mathscr{E}^{+}$ lies in $K^{x}_{k+1}$, one has $\mathscr{E}^{+}|_0=L\,R\left(\mathscr{E}|_0\right)+L\,E_k^{\mathrm{det}}+E^{*}_k\left(m,\mathscr{E}|_0\right)$, and
\[
\left\|\mathscr{E}^{+}\right\|^{x}_{k+1}\le c_{\mathrm{fl}}\cdot L^{-\frac{1}{4}+2\varepsilon}\cdot\|\mathscr{E}\|^{x}_k+e_{k+1}^{\frac{1}{4}-2\varepsilon}+e_k^{\frac{1}{4}-6\varepsilon}.
\]
\item[Z1.e]
Of this clause the second sentence is used, that is, the bound for $E^{*}_k$ (in the reading (o) of Definition~24.D\,(D7)) for $(m_1,E_1)$, $(m_2,E_2)$ in $\mathrm{Ball}_{k}(1)$, with the norms $\|\cdot\|_k$ and $\|\cdot\|_{k+1}$, together with the closing parenthesis on $k=0$ as far as it concerns this bound: for $k=0$ one has $E_1=E_2=0$, and the second quotient is absent.
It is used in the form of the Lipschitz bound
\[
\left\|E^{*}_k(m_1,E_1)-E^{*}_k(m_2,E_2)\right\|_{k+1}\le e_k^{\frac{1}{4}-6\varepsilon}\cdot\left(\frac{|m_1-m_2|}{w^{\mathrm{m}}_k}+\frac{\|E_1-E_2\|_k}{w^{\mathrm{E}}_k}\right)\qquad\left((m_i,E_i)\in\mathrm{Ball}_{k}(1),\ i=1,2\right).
\]
\item[Z1.i]
The bounds of Z1.a--Z1.e hold with the numbers printed there, for all data (OS) and for every system of primitives with the properties named; apart from numerals, the norms of the data and $d_M(X)$, they contain the datum $e_k$ and constants only.
\end{itemize}
The clauses Z1.a--Z1.i are used in Propositions~28.1 and~28.2.
\begin{itemize}
\item[Z1.g]
Of this clause the second half is used; it consists of Proposition~24.3\,(adm), (q1), (al), (ga1)--(ga4), (de), (ze), which take the place of \cite[Proposition~17.24\,(a), (c), (d), Corollary~17.15 and Lemma~17.11\,(q1)]{QEDI}.
It is used in the following form.
Every term $(X,P)$ of the display (8.2) of Definition~24.D with $Y'(X,P)=\mathbf{Y}$ has $X$ and all cubes of $P_0(X)\cup P$ contained in $\mathbf{Y}$.
The primitives enter the objects attached to $\mathbf{Y}$ at parameter families that vanish outside the parameter cubes contained in $\mathbf{Y}$, through the rows named in Proposition~27.5\,(c1), (c3), and at the connections $\mathscr{T}+t\cdot\mathscr{Z}(s')$ with $|t|\le\mathrm{rad}_1$ and $\mathscr{T}$.
The elements $E^{(j),x}(\mathbf{Y})$ and $E^{(\le5),x}(\mathbf{Y})$ depend on $\mathsf{Z}$ through the coordinates $\mathsf{Z}(\Upsilon)$ with $\Upsilon$ contained in $\mathbf{Y}$, and on $A$, $\mathscr{C}$ through the bonds that meet $\mathbf{Y}$.
For $s\in\mathrm{Pd}_k$, $\mathsf{Z}\in\mathrm{Box}_{k}$, $A\in\tilde{R}_k^{-}$, $\mathscr{C}\in\mathfrak{X}_k$ and $|t|\le\mathrm{rad}_1$ one has
\[
A+t\cdot\mathscr{Z}(s')\in\tilde{R}_k,\qquad\mathscr{T}+t\cdot\mathscr{Z}(s')\in\tilde{G}_k.
\]
Further, the last sentence of Proposition~24.3\,(q5) is used: every function to which $\mathrm{Loc}$ is applied in (D4) of Definition~24.D satisfies the hypothesis of (D3) there.
For Proposition~27.6\,(e), Proposition~24.1\,(b) and the last sentence of Proposition~24.3 are used, that is, the same statements in the reading (o) of Definition~24.D\,(D7).
\end{itemize}
(The forms of this entry are used in Propositions~27.5, 27.6, 28.1, 28.2 and~29.)

\item[Propositions 25.1, 25.2 and 25.3\,(a), (b)]
These are used in the sense of Section~I.1.9, for the class of the sides $\ell\in\mathrm{Cl}$, in the following form.
\begin{itemize}
\item[(i)]
The operators $C$ (with the range $\mathrm{rg}<2L$); $\delta_{\alpha,B}$ and $\delta_{\alpha,-B}$ ($\mathrm{rg}=M$); $T^{(j)}(B)$ and $\bar{T}^{(j)}(B)$, $j=1,2$, which depend on a union of $M$-cubes ($\mathrm{rg}<1$); $Q_k(B)\,J^{*}$ and $J\,Q_k^{T}(-B)$ ($\mathrm{rg}<1$); $B_{k,y}(B)$ ($\mathrm{rg}<L$); and $F_0(B)$ ($\mathrm{rg}<L$) are given by one formula on $\mathbb{L}^{\mathrm{f}}_k$, with the properties (f1)--(f4) of Section~I.1.9.
\item[(ii)]
The form $\langle\bar{\psi},\psi\rangle_{\square}$, the form (I.5.5) as a function of a bond function $\mathrm{Zf}$, and the homomorphism $\boldsymbol{\Theta}_{X,\mathscr{T}}$ for small $X$ are objects attached to a polymer that are given by one formula.
The homomorphism $S_L$ is not attached to a polymer.
It is given by one formula, and it is indexed by a finite set $\tilde{U}$ of sites of $\mathbb{L}_k$, which stands in the place of $\tilde{X}$; one has
\[
\mathrm{cov}_{\ell/L}\circ S_L=S_L\circ\mathrm{cov}_{\ell}\quad\text{if }\operatorname{sup-diameter}(\tilde{U})+L\cdot M<\ell,\qquad T_{v/L}\circ S_L=S_L\circ T_v\quad\left(v\in LM\mathbb{Z}^3\right).
\]
This is Proposition~25.2\,(o7), with Section~I.1.9.
\end{itemize}
The operators of (i) in the fourth, the fifth and the sixth place, namely $Q_k(B)\,J^{*}$ and $J\,Q_k^{T}(-B)$, $B_{k,y}(B)$, and $F_0(B)$, occur in the determinant of Section~I.5.7, and not in the displays (8.1)$^{\mathrm{ch}}$--(8.4) of Definition~24.D.
In \cite{QEDI} the corresponding operators enter through Step~2 of the proof of \cite[Proposition~19.12]{QEDI}, and not through \cite[Proposition~17.24\,(b)]{QEDI}.
(Used in Proposition~27.4.)
\end{description}

\subsubsection*{G4.A.4. From the Propositions 38.1--47.2 and 50}

\begin{description}
\item[Propositions 38.2\,(b), (c) (last display), 38.1\,(a), (d) and 38.3\,(a)]
\begin{itemize}
\item[(i)]
Let $n=(n_1,n_2,n_3)$ with $n_\mu\ge4$, let $\boldsymbol{\mu}$ be real, and let $\Phi$ be a Borel function with $|\Phi|\le\mathrm{cf}\cdot\exp\left(\mathrm{cf}\cdot\|F^{\mathrm{u}}\|\right)$.
Then
\[
Z^{(0)}_n[\Phi]=\mathrm{Nc}_n\int_{[-\pi,\pi]^3}\frac{\mathrm{d}\theta}{(2\pi)^3}\int_{{V_{\mathrm{t}}}}\mathrm{d}\sigma(A)\,\Phi(dA)\cdot\exp\left(-\frac{\|dA\|^2}{2}\right)\cdot\det\left(D_{A+H_{\theta}}+\boldsymbol{\mu}\right),
\]
and both sides converge absolutely.
\item[(ii)]
Let $\bar{Z}^{\mathrm{ext}}$ be the function of Proposition~49\,(e) on the cubic torus with $n=\mathrm{kt}\cdot n_*$; for $\mathrm{kt}=1$ and $\boldsymbol{\mu}=\boldsymbol{\mu}_N$ it is the function of Section~I.7.6.
By the last display of Proposition~38.2\,(c), which holds on every cubic torus with $n\ge4$ and for every real $\boldsymbol{\mu}$ and expresses integrals against the standard Gaussian measure of $F_N$ through $\exp\left(-\frac{\|dA\|^2}{2}\right)\mathrm{d}\sigma(A)$ on ${V_{\mathrm{t}}}$ (equivalently, by the definitions of Proposition~31.1), one has
\[
\bar{Z}^{\mathrm{ext}}(t)=\mathrm{zV}_n^{-1}\cdot\mathbf{Z}^{\mathrm{x}}_n\left(a(t)\right),\qquad a(t)=e^{-1}\sum_{\mathfrak{c}}t_{\mathfrak{c}}\cdot\hat{\mathfrak{a}}_{\mathfrak{c}}.
\]
The function $\bar{Z}^{\mathrm{ext}}$ is an entire function of the marks $t$.
\end{itemize}
(Used in Propositions~31.1\,(b) and~49\,(e).)

\item[Proposition 46\,(b), (e), (f)]
The operator of \cite[Definition~2.10]{QEDII} with the mass parameter $M_{\mathrm{ch}}$ is $D_A+\boldsymbol{\mu}$.
The determinant $\det D_A^{T}$ of \cite[Definition~2.11]{QEDII} for $T=C(n)$ is $\det\left(D_A+\boldsymbol{\mu}\right)$, the determinant of a linear map of $\mathfrak{S}^{\mathrm{ch}}(1;n)$.
Let $n_\mu\ge4$, and let $\Phi$ be a Borel function of the curvature $F^{\mathrm{u}}$ with $|\Phi(F^{\mathrm{u}})|\le\mathrm{cf}\cdot\exp\left(\mathrm{cf}\cdot\|F^{\mathrm{u}}\|\right)$ for some number $\mathrm{cf}$.
Then the number $Z_T[\Phi]$ of \cite[Definition~2.13]{QEDII}, formed for $T=C(n)$ with the mass parameter $M_{\mathrm{ch}}$, equals the number $Z_n[\Phi]$.
(Used in Proposition~31.1; the first sentence of~(e) is used in Proposition~32.2.)
\end{description}

\subsubsection*{G4.A.5. Results of the companion papers that are used as they are stated}

The following results of \cite{QEDI} are used in the form in which they are stated there.
\cite[Section~19.2]{QEDI}, with the exception of \cite[Proposition~19.2]{QEDI}, in whose place stands (G4.Sc6) of Section~G4.0.2, that is Proposition~14.2\,(c), (f).
The Hodge decomposition \cite[(10.1)]{QEDI} of \cite[Section~10.1.1]{QEDI}.
The definitions of $d_M$ and of $\bar{X}^{L}$ in \cite[Section~13.1.7]{QEDI}, and \cite[Section~13.3]{QEDI}.
The assertions (t2), (t5), (t7) and (t10) of \cite[Theorem~13.2]{QEDI}, the last one in its converse part, for non-empty $\tilde{X}$.
The Hodge decomposition (Hodge) of \cite[Section~14.1.1]{QEDI}; \cite[Proposition~14.4\,(c) and Proposition~14.8\,(ii)]{QEDI}; \cite[Section~14.4.1]{QEDI} and \cite[Lemma~14.13\,(a), (b)]{QEDI}, with the definition of the map $C$ given in \cite[Section~14.4.1]{QEDI}; \cite[Section~14.4.2]{QEDI} with \cite[Proposition~14.14\,(a), (b)]{QEDI}; \cite[Theorem~14.15\,(a)]{QEDI}; \cite[Section~14.5.1]{QEDI}, with its definition of the fields $A_j$; \cite[Theorem~14.16\,(a)]{QEDI}; the first paragraph of \cite[Section~14.5.2]{QEDI}; \cite[Theorem~14.17]{QEDI}; and \cite[Corollary~14.18]{QEDI} without its last sentence, the one that begins with ``In particular''.
The paragraph on kernels in \cite[Section~15.1.1]{QEDI}; \cite[Proposition~15.5]{QEDI}; \cite[Theorem~15.14]{QEDI}, with the constants $C_P$, $\gamma_P$, $c_P$ fixed in \cite[Section~15.7]{QEDI} and with the closing paragraph of \cite[Section~15.7]{QEDI} on these constants; the bounds on sums of exponentials in \cite[Section~15.1.1]{QEDI}; \cite[Corollary~15.20]{QEDI}; \cite[Proposition~15.21]{QEDI}, with its definition of the map $C$; \cite[Propositions~15.28, 15.29 and~15.30]{QEDI}; the definitions of \cite[Sections~15.10.2 and~15.10.3]{QEDI}, among them the number $M^{P}$ and the families $H_k(s')$, $C_k^{1/2}(s')$; \cite[Theorem~15.35 and Proposition~15.36]{QEDI}.
\cite[Lemmas~17.3--17.7]{QEDI}; the facts (g0)--(g4) of \cite[Section~17.4.1]{QEDI}; and the sets $P_0(X)$, $Y^{\mathrm{p}}(X,P)$, $Y'(X,P)$ with the map $\mathrm{Loc}$ of \cite[Section~17.4.3]{QEDI}.

The following results of \cite{QEDII} are used in the form in which they are stated there:
\cite[Definitions~2.1 and~2.6]{QEDII};
the action $G\mapsto G^{g}$ of \cite[Section~2.3]{QEDII} on the functions of the curvature;
\cite[Definitions~2.10, 2.11 and~2.13]{QEDII};
\cite[Proposition~3.18]{QEDII};
\cite[Theorem~4.13\,(d) and Theorem~4.26]{QEDII};
\cite[Lemma~6.26]{QEDII};
\cite[Definition~13.7 and Proposition~13.11]{QEDII};
and \cite[Lemma~13.12\,(b)]{QEDII}, together with the Bloch representation of \cite[Lemma~13.8]{QEDII} for the time directions $e_1$ and $e_2$, on which the last two results rest.
For the time direction $e_2$ the representation is \cite[Lemma~13.8\,(c)]{QEDII}; it also follows from the case of the time direction $e_1$ by \cite[Lemma~6.26]{QEDII}.

\subsubsection*{G5.7. Forms used in the proofs of the Propositions 33.1--37.3}

Each entry names the statement from which the form is taken, the form that is needed, in the notation of Section~G5.1, and the places where it is used.
The forms are clauses of Proposition~35.4 unless they are printed here.
The properties W1--W17 are those of \cite[Section~5.4.2]{QEDI}, read as in Section~I.5.8.

\begin{description}
\item[(G5.T-1)]
\emph{Taken from} the Propositions~14.m, with Sections~I.8.2 and~I.8.4.
\emph{Form:}
$C^{F}\ge1$; $0<\gamma^{F}\le1$; $\gamma^{F}M\ge16$; $C_{0}\ge\max\left(1,C^{F}\right)$; $M\ge L^2$; $K_{0}\ge1$; $C_{d}<\infty$, one constant for all levels $j\ge0$; (A8) and (A10); the three consequences of (C2), (C3) listed in Section~G5.1.9; $e_1>0$, defined by the conditions (e1)--(e8) of \cite[Section~19.1.2]{QEDI} formed with the constants of Section~I.8.2, and $e_1\le e_0$.
The pair $(C^{F},\gamma^{F})$ consists of the constants of FE1$^{\mathrm{ch}}$--FE4$^{\mathrm{ch}}$ in the norm $|\cdot|_c$ (Proposition~14.1\,(i)).
\emph{Used in} proofs of the Propositions~33.1--35.10, the range of Section~G5.1.9: its conditions and constants enter these proofs at the places named in the sections ``Conditions on the constants'' of Chapters~33, 34, 35a, 35b and~35c.
None of them enters Propositions~36, 37.1 and~37.2, and Proposition~37.3 involves the form only through Propositions~35.1\,(f) and~35.2\,(a) and through its datum $e\cdot h^{\frac{1}{2}}\le e_1$.

\item[(G5.T-2)]
\emph{Taken from} Proposition~14.1\,(i) with Proposition~2.8: FE1$^{\mathrm{ch}}$ at $k=k_{*}$, as stated in Proposition~14.1\,(i), in the scale-free setting $k\ge1$, $\ell\in\mathrm{Cl}$, $\mu\in(0,1]$, $e_k\le e_0$.
\emph{Form:}
for $A\in\tilde{G}_*$, the operator $\mathrm{Fst}(A)$ is bijective on $\mathfrak{S}^{\mathrm{ch}}(\eta_*,\ell_*)$; $S_*(A)$ is analytic on $\tilde{G}_*$; the majorants of $S_*(A)$ and of $\mathrm{Sa}_*(A)$ at $(y,y')$ are at most $C^{F}\cdot\mathrm{Edy}(y,y')$, the second one over all short pairs ($d_X<1$), the pairs inside a cluster included.
\emph{Used in} Proposition~35.4\,(SF2) and in Propositions~33.2, 33.3.

\item[(G5.T-3)]
\emph{Taken from} Proposition~6 (FE4$^{\mathrm{ch}}$ at $k=k_{*}$).
\emph{Form:}
for $A\in\tilde{G}_*$, the majorants of $H_*(A)$ and of $\delta_{\alpha,A}H_*(A)$ at $(y,y')$, the latter for the short pairs, clusters included, are at most $C^{F}\cdot\mathrm{Edy}(y,y')$; the sup norm of $H^{\#}_*(A)v$ is at most $C^{F}$ times that of $v$, the supremum extending over the sites and over all pairs of $\mathrm{Pz}_*$; the same holds for $\bar{H}_*(A)$ with $\delta_{\alpha,-A}$.
\emph{Used in} Proposition~35.4\,(SF2), and through (G5.T-6).

\item[(G5.T-4)]
\emph{Taken from} (I.Lc2)$^{\mathrm{ch}}$ and from \cite[Lemmas~17.3, 17.4 and~17.6]{QEDI} at the levels $j\le k_{*}$.
\emph{Form:}
the modulus of $U_B(x,x')^{\pm1}$ is at most $\mathrm{e}$ on $\tilde{G}_*$; the inclusions and invariances of Proposition~35.4\,(SF2) hold; $G_j$, $\mathfrak{X}_j$, $\mathrm{cf}\cdot\tilde{R}_j$ and $\tilde{G}_j$ are open.
\emph{Used in} Proposition~35.4\,(SF2); \cite[Lemma~17.3]{QEDI} is used also in Propositions~34.1, 34.2 and~34.3.

\item[(G5.T-5)]
\emph{Taken from} Proposition~19\,(a)--(e) (AC8$^{\mathrm{ch}}$) at $k=k_{*}$ and $\mu=\mu_{*}$; the hypothesis $\mu\le12\mathrm{e}$ of Proposition~19\,(e) holds, since $\mu_{*}\le1$ by (G4.Sc3).
\emph{Form:}
clause (A7) of Proposition~35.4 in its full wording: for every real bond function $A$ on $T_*$, not only for $A\in\tilde{G}_*$; with the invertibility of $\mathrm{Fst}(A)$; with the blocks in the norm $|\cdot|_2$; and with the numbers $\frac{2\mathrm{e}}{\mu_*}$ and $\frac{\mu_*}{12\mathrm{e}}$.
\emph{Used in} Propositions~33.4\,(a) and~34.3\,(c).

\item[(G5.T-6)]
\emph{Taken from} the Propositions~20.m (AC7$^{\mathrm{ch}}$\,(b), (c) at $k=k_{*}$).
\emph{Form:}
clause (A6) of Proposition~35.4, with the weight $C_0^{-1}h_{*,2}$ and the number $4\,e_*^{-\frac{1}{2}+2\varepsilon}\cdot\ell_*^3$.
\emph{Used in} Proposition~34.1.

\item[(G5.T-7)]
\emph{Taken from} the Propositions~20.m (AC9$^{\mathrm{ch}}$).
\emph{Form:}
clause (SF4)\,(vi) of Proposition~35.4 with its display (2.2), for every complex $a$ with $a^{[j]}\in\mathfrak{X}_j$ ($j\le k_{*}$), the first step included:
\[
\det D_1(B')=\frac{b^{4n_1}\cdot\det D_0(B)}{\delta Z_0(B)}\qquad\text{(I.L0.10)}.
\]
The identity $\delta Z_j(A)/\delta Z_j(0)=\exp\sum_{X}E_j^{\mathrm{det}}(X,A)$ of Proposition~13\,(d0) enters only through Proposition~20.3\,(c).
\emph{Used in} Proposition~34.3\,(a), (b).

\item[(G5.T-8)]
\emph{Taken from} Proposition~13 (FE7$^{\mathrm{ch}}$\,(d1)--(d3)), for every system with W1, W2, W7, W12--W14 for $\mathrm{par}$, at every level $j\ge0$.
\emph{Form:}
Proposition~35.4\,(SF4)\,(i)--(iv): analyticity on $\tilde{G}_j$; the bound $C_{d}\cdot M^3\cdot\exp\left(-(\kappa+1)\cdot d_M(X)\right)$ for $E^{d}_j$; dependence on $A$ through the bonds that meet $X$; invariance under $A\mapsto A-d\boldsymbol{\lambda}$, $\boldsymbol{\lambda}$ a real function of the photon sites; evenness.
At $j=0$ the form is used with the function $\acute{E}_0$ of Section~I.5.7.
\emph{Used in} Proposition~35.4\,(SF4) and in Propositions~34.2, 34.3\,(b), 35.6.

\item[(G5.T-9)]
\emph{Taken from} the Propositions~26.m (W18$^{\mathrm{ch}}$), with Propositions~7 and~8.
\emph{Form:}
the projected systems on $T_j(R^{*}(e))$, $j<k_{*}$, have W1--W17 for $\mathrm{par}$ (Section~I.5.8).
\emph{Used in} Proposition~35.4\,(SF4).

\item[(G5.T-10)]
\emph{Taken from} Proposition~10 (FE6$^{\mathrm{ch}}$\,(a), for real constants only).
\emph{Form:}
for $A\in\tilde{G}_*$ and a real constant bond function $\mathrm{Acl}$ with $e_*\cdot\ell_*\cdot\mathrm{Acl}_\mu\in2\pi\mathbb{Z}$,
\[
D_*\left(A-\mathrm{Acl}\right)=\mathrm{uf}^0\,D_*(A)\left(\mathrm{uf}^0\right)^{-1},\qquad\mathrm{uf}^0(y)=\exp\left(i\,e_*\cdot\mathrm{Acl}\cdot y\right);
\]
$\mathrm{uf}^0$ is a diagonal matrix of phases that does not depend on $A$.
\emph{Used in} Proposition~34.4\,(a).

\item[(G5.T-11)]
\emph{Taken from} the Propositions~31.m.
\emph{Form:}
on $C(n_*)$, for the pair $P_N(e)$ and every real $\boldsymbol{\mu}_N$:
(i) $\mathrm{dens}^{\theta,a}_N\ge0$ for real $a$;
(ii) $0<\bar{\nu}(F_N)<\infty$;
(iii) (Gaussian domination) for every real two-cochain $\mathrm{gc}$ on $\mathbb{Z}^3/n_*\mathbb{Z}^3$,
\[
\bar{\mu}\left(\exp\langle\mathrm{gc},F^{\mathrm{u}}\rangle\right)\le\exp\left(\frac{\|\mathrm{PrF}\,\mathrm{gc}\|^2}{2}\right),
\]
where $\mathrm{PrF}$ is the orthogonal projection of $C^2(n_*)$ onto $F_N$; consequently, for $\phi\in F_N$ of norm at most $1$ and $y\ge0$,
\[
\bar{\mu}\left\{\left|\langle\phi,F^{\mathrm{u}}\rangle\right|>y\right\}\le2\exp\left(-\frac{y^2}{2}\right).
\]
\emph{Used in} Proposition~35.7.

\item[(G5.T-12)]
\emph{Taken from} the Propositions~28.m (the chamber forms of \cite[Theorem~5.81\,(a)]{QEDI} and of \cite[Theorem~5.83]{QEDI}).
\emph{Form:}
the clauses (SF1), (SF3) of Proposition~35.4: $P_N(e)$ is real and controlled on $R^{*}(e)$, and $|m_k|\le w^{\mathrm{m}}_k/8$; $\mathscr{E}_*\in K^{x}_*$, with norm at most $4\,e_*^{\varepsilon}\cdot w^{\mathrm{E}}_*$.
The identity $\mathfrak{P}_k(X,0)=0$, which holds by Proposition~28.1, enters the Propositions~34.m and~35.m through Proposition~34.2, where it is derived from (SF4).
\emph{Used in} Propositions~34.1, 35.5 and~35.7.

\item[(G5.T-13)]
\emph{Taken from} Proposition~29 (the chamber form of \cite[Theorem~5.84\,(b), (c)]{QEDI}).
\emph{Form:}
clause (SF6) of Proposition~35.4: the gauge invariance of $Y^{x}_*$ in $\mathrm{Dom}_*$, and the display (2.5) for every complex $a$ with $a^{[j]}\in\mathfrak{X}_j$, every $\theta\in\mathbb{R}^3$ and every bounded Borel function $\Phi$.
\emph{Used in} Propositions~34.1\,(c) and~35.7\,(c)--(e).

\item[(G5.T-14)]
\emph{Taken from} Proposition~49\,(d).
\emph{Form:}
on $C(n_*)$ the function $N_0$ is entire, even and real on the real axis, $N_0(0)=1$, and $N_0$ does not change if $\hat{\mathfrak{a}}$ is replaced by $(g_{\mathrm{w}})_*\hat{\mathfrak{a}}$ with $g_{\mathrm{w}}\in\mathrm{Gc}$, that is, by $\hat{\mathfrak{a}}_{\mathfrak{c}}$ with $\mathfrak{c}\in\mathrm{Cel}$, or by the cochain formed with another $v\in\eta\mathbb{Z}^3$.
(This is contained in Proposition~36\,(a), (h) together with Proposition~35.2\,(a).)
\emph{Used in} Proposition~34.3\,(a), (d), (e).

\item[(G5.T-15)]
\emph{Taken from} Proposition~46\,(b), (g+).
\emph{Form:}
$D^{\mathrm{ch}}[qA]+\boldsymbol{\mu}=D_A+\boldsymbol{\mu}$; $B_t=q\cdot(t/e)\cdot\hat{\mathfrak{a}}$; and the quotient $\det\left(D^{\mathrm{ch}}[B_t]+\bar{m}\eta\right)/\det\left(D^{\mathrm{ch}}[0]+\bar{m}\eta\right)$ is $N_0(t)$; all of this for every $\eta\in(0,1]$ and every $R$ with $R/\eta$ an integer $\ge2$.
\emph{Used in} Section~G5.6.0.

\item[(G5.T-16)]
\emph{Taken from} Proposition~15 and from Sections~I.3.6, I.5.2 at $\nu=1$.
\emph{Form:}
(I.N0), (I.N1), (I.N2), (I.S6c)\,(i)--(iii) and (I.B1), for the index set $T^0_*\times\{1,\dots,4\}$ and for every finite index set.
\emph{Used in} Propositions~34.1, 35.5, 35.6, and in \cite[Lemma~20.9\,(a)]{QEDI}.

\item[(G5.T-17)]
\emph{Taken from} \cite{QEDI}, in the form in which the results are stated there.
\emph{Form:}
\cite[Proposition~19.1\,(a)--(c), (e)--(g)]{QEDI};
the assertions (t2), (t5), (t7), (t8) of \cite[Theorem~13.2]{QEDI};
\cite[Lemma~17.5\,(c)]{QEDI} and the fact (g0) of \cite[Section~17.4.1]{QEDI};
\cite[Section~14.5.1]{QEDI}, the first paragraph of \cite[Section~14.5.2]{QEDI}, \cite[Theorem~14.16\,(a) and Theorem~14.17]{QEDI};
the Hodge decomposition of \cite[Section~15.1.1]{QEDI}, \cite[Proposition~15.5]{QEDI}, \cite[Theorem~15.14]{QEDI} with the constants $C_P$, $\gamma_P$, $c_P$ fixed in \cite[Section~15.7]{QEDI} and the closing paragraph of \cite[Section~15.7]{QEDI} on these constants, the bounds on sums of exponentials in \cite[Section~15.1.1]{QEDI}, \cite[Proposition~15.21\,(ii), Proposition~15.30\,(i) and Corollary~15.20]{QEDI};
\cite[Section~7.1.3]{QEDI}, for the translations $\tau_b$ (in Proposition~36\,(h));
\cite[Lemmas~20.8--20.13]{QEDI}.
\emph{Used in} Proposition~35.1, Proposition~35.4\,(SF4)\,(v), (SF5), the Propositions~34.m, and the Propositions~35.5--35.10.
\end{description}

\subsubsection*{G6.Z1. Forms used in the proofs of the Propositions 38.1--47.2 and 50 and of Theorem 44}

\begin{description}
\item[(G6.T-14)]
\emph{Taken from} Sections~I.0.5, I.6.1, I.6.2 and~I.8.2; from \cite[Proposition~19.1\,(b), (c), (d)]{QEDI}; from the first display of (G4.Sc0); and from (G4.Sc6), that is Proposition~14.2\,(b), row (e1) (its first and its third inequality, at $v=v_K$), and (c), (f), with the second assertion of Proposition~14.3\,(f); these stand in the place of the first two sentences of \cite[Proposition~19.2]{QEDI}, that is, of that proposition without its assertions on the condition (W3).
These give (i)--(iv).
Of these facts, Proposition~35.1 contains the following: before its clauses, $R^{*}(e)\in\mathrm{Tor}(e)$ and the finiteness of $N_0(e)$; in (a), $e\le\frac{1}{3}$, $v_K\ge8$, and that $T^{*}(e)$ is an integer $\ge1$; in (f), the lower bound on $\ell_c$, and $\bar{m}\cdot R^{*}(e)>r_K\ge v_K$ with its consequence for $\exp\left(-\bar{m}\cdot R^{*}(e)\right)$; the clauses (a) and (f) are stated for $N\ge N_0(e)$.
The two inequalities of the display in (i), the assertions on $R_{\mathrm{top}}(e)$ and on $N_{\min}\left(R_{\mathrm{top}}(e),e\right)$ in (ii), and the inequality $\bar{m}\cdot R^{*}(e)\ge\bar{m}\cdot h>\frac{1}{L}$ with the bound on $\varepsilon_5$ in (iv) come only from the sections and statements named above.
The clauses (v)--(vii) are taken from Proposition~14.2\,(b), rows (e3), (e8), from Proposition~14.2\,(d), conditions (Z8.w) and (W3), from Proposition~14.2\,(f), and, for (vii), also from the lines (C5) and (C7) of Section~I.8.2;
the value $\mathfrak{e}_2=0$ in (vii) is the one fixed in Section~I.8.5, which Proposition~14.2\,(a) repeats.
The clauses concern the arithmetic of the scales and conditions on the constants.
In \cite{QEDI} the corresponding statements are \cite[Propositions~19.1 and~19.2]{QEDI} and Steps~1, 2, 4 and~6 of the proof of Part~1 of \cite[Theorem~21.3]{QEDI} in \cite[Section~21.9]{QEDI}, read with $h$ in place of $h'$ and with $0$ in place of $m_0$.

\emph{Form.}
Let $e>0$ with $e\cdot h^{\frac{1}{2}}\le e_1$.
\begin{itemize}
\item[(i)]
$e\le\frac{1}{3}$, $v_K\ge8$, and
\[
\left(1+\frac{\mathrm{lg}}{v_K}\right)^{r}\le\left(1+\frac{3\,\mathrm{lg}}{v_K}\right)^{r}<L.
\]
The last inequality is the third inequality of the row (e1) of Proposition~14.2\,(b) at $v=v_K\ge-\log e_1$, by the second assertion of Proposition~14.3\,(f); the first one holds because $\mathrm{lg}>0$, $v_K>0$ and $r\ge1$.
The inequalities $v_K\ge1$ and $\left(1+\frac{\mathrm{lg}}{v_K}\right)^{r}\le L$ are the hypotheses of \cite[Proposition~19.1]{QEDI}.
\item[(ii)]
$T^{*}(e)=\log_L\left(M\cdot r_K(e)\right)-1$ is an integer $\ge1$; $R^{*}(e)=h\cdot L^{T^{*}(e)}$ and $L\cdot R^{*}(e)=R_{\mathrm{top}}(e)$ lie in $\mathrm{Tor}(e)$; and
\[
N_{\min}\left(R_{\mathrm{top}}(e),e\right)\le N_0(e)=N_{\min}\left(R^{*}(e),e\right)<\infty.
\]
\item[(iii)]
$\ell_c=\frac{R^{*}(e)}{2L}\ge\frac{h\cdot M\cdot v_K^{r}}{2L^2}$ with $v_K=-\log\left(e\cdot h^{\frac{1}{2}}\right)$; hence $\ell_c\to\infty$ as $e\to0$.
\item[(iv)]
$\bar{m}\cdot R^{*}(e)\ge r_K(e)\ge v_K$, so that $\exp\left(-\bar{m}\cdot R^{*}(e)\right)\le e\cdot h^{\frac{1}{2}}$; and $\bar{m}\cdot R^{*}(e)\ge\bar{m}\cdot h>\frac{1}{L}$, so that $\varepsilon_5\left(\bar{m}\cdot R^{*}(e)\right)\le\varepsilon_5\left(\frac{1}{L}\right)$.
\item[(v)]
If $e\le e_G$, then (W3) holds:
\[
6\,N_{*}\cdot e^2\cdot R^{*}(e)\cdot\log\frac{1}{e}\le\pi^2.
\]
\item[(vi)]
$3\left(e\cdot h^{\frac{1}{2}}\right)^{\frac{1}{4}-9\varepsilon}\le1$.
\item[(vii)]
(C3): $L>8$.
(C5): $M\ge L^2$.
(C7): $0<s<1$ and
\[
\sigma=s\left(\frac{1}{4}-9\varepsilon\right)-12\varepsilon-\frac{2a_L}{N_{*}}>0,
\]
hence $0<\mathfrak{e}\le\frac{1}{4}$.
Further $\mathfrak{e}_2=0$ (Section~I.8.5 and Proposition~14.2\,(a)), so that $0\le\mathfrak{e}_2<1$, which is the clause on $\mathfrak{e}_2$ of (R6) in Proposition~43.1.
\item[(viii)]
(From Proposition~14.2\,(a), rows $G_L$ and $K_L$.)
\[
G_L=\left(1+\frac{L}{3.2}\right)\cdot\frac{2L^2}{2.56},\qquad K_L=\left(1+L+L^2\right)\left(\log L+L\cdot G_L\right);
\]
these are the numbers $G_{\check{j}}$, $K_{\check{j}}$ of Proposition~42.1 at $\check{j}=L$.
\end{itemize}

\item[(G6.T-28)]
\emph{Taken from} Proposition~28.3\,(c), (d).
\emph{Form:}
for $e\cdot h^{\frac{1}{2}}\le e_1$ and $N\ge N_{\min}\left(R_{\mathrm{top}}(e),e\right)$ the pair
\[
P_N(e)=\left(m_{0,N}(e),\varepsilon^{\mathrm{v}}_{0,N}(e)\right)
\]
is defined and real.
It depends on $(N,e,\bar{m},\rho)$ and on the scheme of Section~I.6.2; it depends neither on $R\in\mathrm{Tor}(e)$ nor on $\theta$.
Moreover
\[
\left|m_{0,N}(e)\right|\le\frac{w^{\mathrm{m}}_0}{8}=\frac{q^{\frac{3}{4}-8\varepsilon}}{8},\qquad q=e\cdot\eta^{\frac{1}{2}}\le\frac{1}{3}.
\]

\item[(G6.T-32)]
\emph{Taken from} Proposition~32.3, display (10.3).
\emph{Form:}
for $e\cdot h^{\frac{1}{2}}\le e_1$ and all $N\ge N_0(e)$, with $n_*=R^{*}(e)/\eta$, $q=e\cdot\eta^{\frac{1}{2}}$ and $\mathrm{zr}_n=\mathrm{zr}_n\left(q,\bar{m}\eta+m_{0,N}(e)\right)$ of Section~I.7.3,
\[
D^{\mathrm{mat}}=L^3\log\mathrm{zr}_{n_*}-\log\mathrm{zr}_{Ln_*}\le\left(L^3+1\right)\left(\varepsilon_5\left(\bar{m}\cdot R^{*}(e)\right)+3\left(e\cdot h^{\frac{1}{2}}\right)^{\frac{1}{4}-9\varepsilon}\right).
\]
The bound presupposes $\mathrm{zr}_n>0$, which is Proposition~38.2\,(c).

\item[(G6.T-35)]
\emph{Taken from} Proposition~35.9\,(b) with the last sentence of Proposition~35.9, from Proposition~35.10, and from Proposition~35.3\,(c1), which gives $C_\beta<\infty$ and hence $C_{35}(\beta)<\infty$.
\emph{Form:}
for every probe $\beta$ there are numbers
\[
e_{\sharp}(\beta)\in\left(0,e_1\cdot h^{-\frac{1}{2}}\right],\qquad C_{35}(\beta)=3\cdot10^{5}\cdot\mathrm{Cb}^{4}\cdot h^{\frac{1}{8}}<\infty,
\]
which depend only on the constants of the lines (C0)--(C9), on $\bar{m}$ and on $\beta$, such that the following holds for all $0<e\le e_{\sharp}(\beta)$, all $N\ge N_0(e)$ and every placement $v$ of Section~I.7.6.
First, $D_\beta+3\eta\le\ell_c$.
Second, let $\boldsymbol{\mu}=\boldsymbol{\mu}_N=\bar{m}\eta+m_{0,N}(e)$, let $\mathring{\Phi}$ and $\Theta^{(0)}_m$ be as in Section~I.7.6 on $C(n_*)$, and let $c_{m}:=N_0^{(m)}(0)$; then
\[
\Theta^{(0)}_4\le C_{35}(\beta)\cdot e^{\mathfrak{e}},\qquad\left(1+|c_{2}|\right)\cdot\Theta^{(0)}_2\le C_{35}(\beta)\cdot e^{\mathfrak{e}}.
\]
In \cite{QEDI} the corresponding statements are \cite[Corollary~20.32 and Lemma~20.33]{QEDI}, in whose place Propositions~35.9 and~35.10 stand.
The inequalities
\[
e\cdot\left(\frac{h}{L^2}\right)^{\frac{1}{2}}\le e_{*}\le e\cdot h^{\frac{1}{2}}
\]
are Proposition~35.1\,(a).

\item[(G6.T-36)]
\emph{Taken from} Proposition~36\,(a), (g), (h).
\emph{Form:}
for every probe $\beta$ there is a number $\mathrm{Cpv}_\beta=\mathrm{Cpv}_\beta(\rho,\bar{m},\beta)<\infty$ with
\[
\left|a_{\mathrm{ch}}(\eta,R)\right|\le\mathrm{Cpv}_\beta
\]
for all $0<\eta\le1$ and all $R\ge1$ such that $R/\eta$ is an integer $\ge2$ and $R>D_\beta$.
Moreover $a_{\mathrm{ch}}(\eta,R)=N_0''(0)$, where $N_0$ is formed with $\tau_v\beta$ for an arbitrary $v\in\eta\mathbb{Z}^3$.

\item[(G6.T-37)]
\emph{Taken from} Proposition~37.2\,(a), (d)--(g).
\emph{Form:}
let $\beta$ be a probe, and let $D_{\mathrm{be}}$ be the least side of a closed coordinate cube that contains $\operatorname{supp}\beta$, so that $D_{\mathrm{be}}=D_\beta-4$.
There are numbers $C_{1}^{\mathrm{N2}}=C_{1}^{\mathrm{N2}}(\rho,\bar{m},\beta)<\infty$ and $C_{2}^{\mathrm{N2}}=C_{2}^{\mathrm{N2}}(\bar{m},\beta)<\infty$ such that for all $0<\eta\le1$ and all $R\ge D_{\mathrm{be}}+1$ with $R/\eta$ an integer $\ge2$
\[
\left|b_{\mathrm{ch}}(\eta,R)-b(dd^{*}\beta)\right|\le C_{1}^{\mathrm{N2}}\cdot\eta^{\frac{1}{4}}+C_{2}^{\mathrm{N2}}\cdot\exp\left(-\bar{m}\cdot\left(R-D_{\mathrm{be}}\right)\right).
\]
The numbers $C_{1}^{\mathrm{N2}}$, $C_{2}^{\mathrm{N2}}$ depend neither on $\eta$ nor on $R$ nor on the position of $\operatorname{supp}\beta$.
The number $b(dd^{*}\beta)$ is real (Proposition~37.2\,(a); see also \cite[Lemma~21.5]{QEDI}).
\end{description}

The following results of \cite{QEDI} are used in the form in which they are stated there.
\cite[Theorem~21.1\,(a)]{QEDI}, the existence of a probe $\beta$ with $b(f)>0$; \cite[Lemmas~21.5, 21.31 and~21.33]{QEDI}.
\cite[Proposition~19.1\,(b)--(d)]{QEDI}.
The Hodge decomposition \cite[(10.1)]{QEDI} and the statement \cite[(10.2)]{QEDI} on fluxes and integer coboundaries, both in \cite[Section~10.1.1]{QEDI}, and \cite[Lemma~10.6]{QEDI}.
\cite[Proposition~7.10]{QEDI}; the first clause of \cite[Theorem~7.12\,(2)]{QEDI}, with the last sentence of its proof; \cite[Lemmas~7.14 and~7.16]{QEDI}; \cite[Sections~7.1.1, 7.1.3 and~7.1.5]{QEDI}.
\cite[Sections~8.1.1--8.1.6]{QEDI}, with the exception of the last paragraph of \cite[Section~8.1.5]{QEDI}, which concerns the limit points of the lattice laws and in whose place Theorem~44\,(B1) stands; \cite[Lemmas~8.11--8.18 and~8.21]{QEDI}; \cite[Theorem~8.3]{QEDI}; \cite[Propositions~8.19 and~8.20]{QEDI}.
\cite[Section~9.2]{QEDI}, with the exception of its last sentence, which concerns the limit points of the lattice laws and in whose place Theorem~44\,(B2), (B5) stand; \cite[Sections~9.3 and~9.7.1]{QEDI}; the inequality between $\|a\|_\infty$, $r$, $m$, $d_0$ and $b(f)$ among the hypotheses of \cite[Theorem~9.21]{QEDI}; \cite[Lemma~9.1\,(a), (c)]{QEDI}; \cite[Definition~9.5]{QEDI}; and \cite[Theorem~9.21\,(i)]{QEDI}.
Further, \cite[Remark~21.17]{QEDI} is used in the reading of Section~I.9.2: without its last sentence, which concerns \cite[Theorem~3.1, Part~I\,(d)]{QEDI}, and for an arbitrary family of Borel probability measures on $H_{-4}$ indexed by the cutoffs $(N,T)$.
In this reading, an iterated weak limit of such a family, taken in either order, is a weak limit of the family along a sequence of cutoffs $(N_i,T_i)$ with $N_i\to\infty$ and $T_i\to\infty$.
The argument uses only the Prokhorov metric on $H_{-4}$, which is a separable complete metric space.

The following results of \cite{QEDII} are used in the form in which they are stated there.
\cite[Definitions~2.1--2.13 and~2.15--2.19]{QEDII}, \cite[Definition~2.21]{QEDII}, \cite[Definitions~2.26--2.28]{QEDII} and \cite[Definitions~2.36--2.38]{QEDII}.
The third clause of \cite[Proposition~3.6]{QEDII}, on the equivariance of the projection $\pi$; \cite[Propositions~3.16 and~3.18]{QEDII}; \cite[Lemmas~3.10--3.15]{QEDII}.
\cite[Definitions~4.1, 4.3 and~4.4]{QEDII} with \cite[Lemma~4.2]{QEDII}; \cite[Lemmas~4.6 and~4.7]{QEDII}; \cite[Theorem~4.13]{QEDII}; \cite[Proposition~4.23]{QEDII}; \cite[Theorems~4.25 and~4.26]{QEDII}.
\cite[Proposition~3.20]{QEDII}; \cite[Definition~2.35]{QEDII}.
\cite[Definitions~5.2, 5.7, 5.9, 5.11 and~5.12]{QEDII} and \cite[Lemmas~5.3, 5.13, 5.14 and~5.15]{QEDII}, the last one with its bound for sums of Gaussian terms and its facts on Hilbert--Schmidt kernels; \cite[Proposition~5.17]{QEDII}; \cite[Lemma~5.4]{QEDII}; \cite[Example~5.6]{QEDII}; \cite[Lemma~5.18]{QEDII}; \cite[Lemmas~5.8 and~5.19]{QEDII} and \cite[Definitions~5.20 and~5.21]{QEDII}, with \cite[Lemma~5.22]{QEDII}; \cite[Corollary~5.30]{QEDII}; \cite[Proposition~5.33, Lemma~5.35, Corollary~5.41 and Proposition~5.43]{QEDII}; the third clause of \cite[Lemma~5.40]{QEDII}; \cite[Lemma~5.46]{QEDII} with \cite[Proposition~5.47]{QEDII}.
\cite[Lemmas~6.3, 6.4, 6.5, 6.17, 6.18 and~6.26]{QEDII}; \cite[Definition~6.7]{QEDII}; \cite[Proposition~6.8]{QEDII}; \cite[Theorem~6.16\,(a)]{QEDII}; \cite[Proposition~6.19\,(a), (b)]{QEDII}; \cite[Proposition~6.23\,(a)]{QEDII}; \cite[Theorem~6.27]{QEDII}, read with (I.h1)--(I.h4).
\cite[Theorem~1.2]{QEDII}.
\cite[Definition~12.4]{QEDII}, for the norms $p_R$; \cite[Lemma~12.5\,(c)]{QEDII}; \cite[Lemma~12.9]{QEDII}; \cite[Propositions~12.10, 12.13, 12.15 and~12.16]{QEDII}; \cite[Proposition~12.12]{QEDII}, with the steps (1), (2), (3) and (5)--(8) of its proof; \cite[Lemma~12.18\,(b), (c)]{QEDII}; \cite[Lemma~12.25]{QEDII}; \cite[Propositions~12.32 and~12.33]{QEDII} and \cite[Theorems~12.39 and~1.3]{QEDII}, read with (I.h5).
\cite[Definition~13.7]{QEDII}; the identity of \cite[Lemma~13.8\,(a)]{QEDII}; \cite[Lemma~13.12]{QEDII}.
\cite[Lemma~14.12\,(b)]{QEDII}; \cite[Theorem~14.26\,(c)]{QEDII}; \cite[Theorem~1.4\,(a), (b), (d), (f)]{QEDII}.

\subsection*{Annex B. Conditions on the constants for the list of Proposition 14.2}

The conditions collected in this annex restrict the choice of the constants, which is made in the order of Section~I.8.2; the constants are chosen so that all of them hold.
The complete list is that of Proposition~14.2.
This annex contains no list for the Propositions~1.1--13.
For Propositions~1.1--1.4, 1.5--1.7, 4, 10 and~12 the conditions on the constants that enter them are named in the sections ``Conditions on the constants'' of Chapters~1a, 1b, 4, 10 and~12, in this order;
for Propositions~3, 5, 6, 7, 8, 9, 11 and~13 they are named in the paragraphs on the conditions attached to these statements and in the sections ``Conditions on the constants'' of Chapters~3, 5, 6, 7, 8, 9, 11 and~13.
For the Propositions~2.m they are named in the paragraph that precedes Proposition~2.1 and in the sentences on the conditions that follow Propositions~2.5, 2.6, 2.7 and~2.8;
the last two of these sentences name, besides the conditions of that paragraph, the range $\alpha<1$ of line (C1) of Section~I.8.2.
The sections ``Conditions on the constants'' of Chapters~2a and~2b name further conditions that enter the proofs of these statements:
for Propositions~2.1--2.4 the ranges $0<\rho<1$, $b>0$ and $\alpha<1$, and the closure relation (I.1.1) of the standard weights;
for Propositions~2.5--2.8 the ranges $b>0$ and $\varepsilon>0$.
The sections ``Conditions on the constants'' of Chapters~3, 5, 6, 7, 8, 11 and~13 name further conditions that enter the proofs of Propositions~3, 5, 6, 7, 8, 11 and~13:
for Proposition~3, $b>0$ and $\varepsilon>0$;
for Proposition~5, $b>0$;
for Proposition~6, the value $b=1$;
for Proposition~7, $\varepsilon>0$ and $b>0$;
for Proposition~8, $\varepsilon>0$;
for Proposition~11, $b>0$;
for Proposition~13, the standard weights $w=\left(6,6\sqrt{2},6\sqrt{2},6\right)$ with their closure relation (I.1.1), $b>0$ and $\gamma_{\mathrm{w}}\ge0$.
For Propositions~3, 7 and~8 these sections name, moreover, the range $\alpha<1$ of line (C1), which is a condition of Propositions~2.7 and~2.8 and enters these proofs through them;
in the proof of Proposition~3 it enters through the second bound of Proposition~2.7\,(e) and through Proposition~2.8.
The section ``Conditions on the constants'' of Chapter~9 names, among the conditions that enter the proof of Proposition~9, $\rho\le1$ (from the range $0<\rho<1$ of line (C0)), the closure relation (I.1.1) of the standard weights together with the value $w_3=6$, and the inequalities $b>0$ and $\alpha>0$.
The places of use that Proposition~14.2\,(b) gives for each condition do not form a complete list.

\subsubsection*{G2.9. Conditions used in the Propositions 15--20.3 and 48.1--48.6}

Apart from the last two entries, all of the following are conditions of Sections~I.8.2 and~I.8.4.
Apart from the first two entries, the range of $\rho$ and the standard weights, which have no counterpart in \cite{QEDI} (Section~I.8.2, line (C0)), and apart from the last two entries, they are conditions of \cite[Sections~5.6.1 and~5.6.2]{QEDI}.
The second to last entry concerns the class $\mathrm{Cl}$ of the sides (Section~I.1.1), which is a datum of the statements and not a condition on a constant.
The last entry is the inequality $48\le49$, a numerical fact of level $0$, where $\nu=1/48$ and $\mathsf{w}_W=\sqrt{48}$; it has no line in Section~I.8.2 and no counterpart in \cite{QEDI}.
None of them contains $\rho$, a weight $w_i$ or a letter of the lines (C0r), (C4r), (C6), with three exceptions: the range of $\rho$ itself (first entry), the standard weights (second entry), and the entries marked $\dagger$.
The constants $M$, $e_0$ and $e_1$, which occur in the entries (C5), (A8), (A9), (C9) and, through $\mathrm{Cl}$, in the entry on the sides, are fermionic constants in the sense of Proposition~14.4: they depend on $\rho$ through their choice (Proposition~14.4\,(c)).
Each entry gives the condition, its line in Section~I.8.2, the statements in which it is used, and, where the proofs need only a weaker form of the condition, that form.

\begin{description}
\item[(C0)]
$0<\rho<1$; only $0\le\rho\le1$ is needed.
Used in Propositions~19\,(a), (e) and~48.1\,(c).

\item[(C0)]
The standard weights $w=\left(6,6\sqrt{2},6\sqrt{2},6\right)$; only $w_i>0$, the closure relation (I.1.1) and $w_3=6$ are needed.
Used in Proposition~19\,(a), (e), where the weights are data, and in Proposition~48.1\,(c), through $w_3=6$.

\item[(C0)]
The numbers
\[
c_X=54,\qquad\kappa_0=54\log72,\qquad K_0=\frac{\exp(\kappa_0)}{36},\qquad c_{\mathrm{rb}}=c_X\cdot K_0,\qquad c_{\mathrm{fl}}=61\,c_X\cdot K_0,
\]
and $b=1$.
Used in Proposition~16.3\,(b), (d), in the Propositions~17.m, in Proposition~18.5, in Proposition~48.5\,(b) (through $61\,c_{\mathrm{rb}}=c_{\mathrm{fl}}$), in Proposition~19, in Proposition~20.1, in Proposition~20.2 (through $b>0$ and $b_k\le b_{+}=b$), and in Propositions~48.2\,(b), (d) and~48.4\,(b).

\item[(C1), (C7)]
$0<\alpha<1$; $1-\varepsilon<\alpha$; $\varepsilon\le\bar{\varepsilon}=\frac{1}{200}$.
Used in Proposition~17.1\,(a), (c3), in Propositions~18.1, 18.2, in Proposition~18.5\,(b), (d), and in Propositions~48.1\,(c), 48.4\,(b), 48.5\,(b), (d); the inequality $\varepsilon>0$ is used in Propositions~20.1, 20.3\,(a) and~48.3\,(d).

\item[(C2)]
$\kappa\ge7\kappa_0+7$.
The proofs need $\kappa\ge\kappa_0$, and in Propositions~18.5\,(b) and~48.5\,(b) also $\kappa-1\ge\kappa_0$, which serves to apply the assertion (t8) of \cite[Theorem~13.2]{QEDI} at $\kappa-1$ to the large polymers.
Used in Propositions~16.3\,(b), 18.5\,(b), 48.1\,(d) (through $\kappa\ge0$), 48.2\,(b) and~48.5\,(b).

\item[(C3), first condition]
$L$ is odd, and $L^{\frac{1}{8}-\frac{3\bar{\varepsilon}}{2}}\ge2400$.
This condition implies
\begin{gather*}
L\ge5,\qquad L^{\frac{1}{4}}\ge2,\qquad L^{\frac{3}{4}}\ge110,\qquad L^3\exp(-L)\le L^{-\frac{1}{4}},\\
L^{\frac{1}{4}-\frac{3\varepsilon}{2}}\ge13.6,\qquad L^{-\frac{7}{4}}+4L^{-\frac{7}{8}+\frac{\varepsilon}{2}}<1,\qquad L^{-\frac{3}{4}}\le\frac{1}{2}.
\end{gather*}
Used in Proposition~18.1\,(c), in Proposition~18.2\,(ii), (iii), in Proposition~18.5\,(b), (d), and in Propositions~48.4\,(b), 48.5\,(b), (d).
Of this condition only the part ``$L$ odd and $L\ge5$'' enters Proposition~15\,(n) (for $\eta_k\le1$), the Propositions~16.1--16.3, 17.1, 17.2, 18.1\,(b), 18.3\,(e), 18.4\,(a), and the Propositions~48.1--48.3.
The inequality $L\ge5$ enters, moreover, Propositions~20.1\,(c) and~48.6\,(a), (b), through \cite[Lemma~17.4\,(a)]{QEDI}, whose proof uses it in the form of the inclusion $\square^{\sim5}\subset\square^{\sim L}$ for every $M$-cube $\square$.
Of the part ``$L$ odd and $L\ge5$'', only the oddness of $L$ enters Propositions~19\,(b), 20.2, 18.5, 18.6\,(e), 20.3\,(c) and~48.5.
Besides it, the consequences $L^{\frac{1}{4}}\ge2$, $L^{\frac{3}{4}}\ge110$ and $L^3\exp(-L)\le L^{-\frac{1}{4}}$ displayed above enter Propositions~18.5\,(b), (d) and~48.5\,(b), (d), and the inequality $L>1$, a consequence of $L\ge5$, enters Propositions~19\,(c)--(e) and~20.2, for $b_{L}>0$, $b_k>0$ and $\eta_k\le1$.

\item[(C3), third condition]
$L\left(\kappa-6\kappa_0-4\right)-1\ge\kappa$; only $L\left(\kappa-\kappa_0-2\right)\ge\kappa$ is needed.
Used in Propositions~18.5\,(b) and~48.5\,(b).

\item[(C4r)\,$\dagger$]
$\gamma^{F}>0$; only $c(\gamma^{F})<\infty$ is needed.
Used in Proposition~20.1.

\item[(C5)]
$M$ is a power of $L$, $M$ is odd, and $M\ge L^2$.
Used in Proposition~15\,(n) (through $1\le M$), in Propositions~16.1--16.3, 17.1, 17.2, 18.1\,(c), 18.2, 18.3\,(e), 18.4\,(a), in Propositions~48.1\,(c), 48.2\,(b), (e), (f), 48.3\,(f) and~48.4;
in Propositions~18.5\,(f), 18.6\,(a), (b) and~48.5\,(f), through the oddness of $M\cdot L^k$;
in Proposition~48.5\,(e), through $\eta_k\le1<M$;
and in Proposition~20.2, through $\frac{L-1}{2}\le M$.

\item[(C6)\,$\dagger$]
$C_0\ge C^{F}$ and $C_0\ge4b_{+}+16\,C^{F}\cdot c(\gamma^{F})$.
Used in Proposition~20.1.

\item[(A8)]
$e_0^{\varepsilon}\le\frac{1}{4}$.
Used in Propositions~17.1\,(a), 18.1\,(c), in Proposition~20.1 (in its data and in (b), (c)), and in Propositions~20.3\,(a), 48.1\,(c).
The fourth inequality of (A8), which at $v=-\log e_0$ reads $e_0^{\frac{1}{4}}+e_0^{\varepsilon}\le1$, enters Propositions~20.1\,(c) and~48.6 through \cite[Lemma~17.4\,(a)]{QEDI}.

\item[(A9)]
\[
160\,L^{1+\frac{3\varepsilon}{2}}\cdot M\cdot e_0^{\varepsilon}\le1,\qquad90\,c_X\cdot L\cdot M\cdot e_0^{\varepsilon}\le1,\qquad9\,M\cdot L^{\varepsilon}\cdot e_0^{2\varepsilon}\le1.
\]
Used in Propositions~17.1\,(c2), 17.2\,(d), 18.2, 48.3\,(d) and~48.4\,(b).

\item[(C9)]
$e_1\le e_0$.
Used in Propositions~20.3 and~48.6\,(b).

\item[(C9)]
$e_1\le\frac{1}{3}$; only $e_k<1$ is needed, so that $v_k>0$, $p_{0,k}>0$ and $\mu_I(\mathrm{Box}_{k})>0$.
Used in Proposition~48.6\,(b).

\item[(C9); (C3); (C5)]
The condition (e1), of which only
\[
v\ge8,\qquad v^{r}>31L+16,\qquad\left(1+\frac{3\,\mathrm{lg}}{v}\right)^{r}<L\qquad\text{for all }v\ge-\log e_1
\]
is needed; further $L\ge33$, which follows from (L1); and $M$ a power of $L$.
These are used, through \cite[Proposition~19.1\,(e), (f)]{QEDI} and through (G4.Sc6), in Propositions~20.3 and~48.6\,(b).

\item[Class of sides]
$\ell\in\mathrm{Cl}$, a datum of the statements.
The proofs need that $\ell$ is a multiple of $LM$ and that $\ell\ge13\,LM$; the oddness of $\ell/(LM)$ enters only Propositions~16.1\,(d) and~20.2\,(Scaling), through the fact that $\ell/L$ is an odd multiple of $M$.
Used in all the statements of this list.

\item[The inequality $48\le49$]
At level $0$ one has $\nu=1/48$ and $\mathsf{w}_W=\sqrt{48}$, and $48\le49$ gives $\mathsf{w}_W\le7$.
Used in Proposition~15\,(n); through the bound $\mathsf{w}_W\le7$ it enters the constant $C_W$ and the condition (A6), see (I.L0.19).
\end{description}

No condition of this list bounds a constant by a constant that is chosen later.
Among these conditions, the restrictions on $\kappa$, $L$ and $M$ are the following: (C2); the first and the third condition of (C3); and (C5), that is, $M$ is a power of $L$ (hence odd, since $L$ is odd) and $M\ge L^2$.

\subsubsection*{G3.6.1. Conditions on the constants in the Propositions 21.1--25.3}

The conditions collected here are conditions on the choice of the constants in the order of Section~I.8.2; the statements are asserted for constants so chosen.
Their labels are those of Section~I.8.4, where they are attached to the conditions of \cite[Section~5.6.2]{QEDI}, and the inequalities themselves are stated in Proposition~14.2\,(b)--(d); the lines (C0)--(C11) are those of Section~I.8.2.
The conditions are arranged in two layers.
For the statements named in each entry below, the first layer consists of the conditions that are used in the arguments of these statements themselves, and the second of the conditions that enter through the statements applied in these arguments; the labels (G3.T-$n$) refer to Section~G3.0.4.
The complete list of the conditions is that of the Propositions~14.m.

\begin{description}
\item[Propositions 21.1--21.3, 21.4\,(a)--(c), (d1), (d3), and 21.5.]
First layer:
$L$ odd, $L>1$;
$M$ a power of $L$ with $L/2<M$;
$b>0$;
$\rho\ge0$ and $w_i>0$ in Proposition~21.1\,(e) and in Proposition~21.5;
the closure relation (I.1.1) in Proposition~21.1\,(c), (e), and hence in Proposition~21.5;
$\rho\ne0$, $\rho^2\ne1$ and $w_3\ne0$ in the last assertion of Proposition~21.4\,(a) at $k=0$: at $B=0$, if $F_{0}(0)$ is invertible, then $\Psi^{c}$ is constant on clusters only if $\Psi'=0$, and $\bar{\Psi}^{c}$ only if $\bar{\Psi}'=0$.
All of these are contained in the lines (C0), (C3), (C5).
Second layer: none.

\item[Proposition 21.4\,(d2).]
First layer: none.
Second layer, through (G3.T-2a):
$L\ge5$ odd;
$M$ a power of $L$ with $M>4L$;
$M\ge108\pi\cdot C'_V\cdot a'$, which is part of (Z2.m);
(Z2.s1), and (Z2.s2) formed with $C'_V$ and $a'$;
$0\le\alpha<1$ (the conditions stated with Proposition~2.7 include $\alpha<1$, and $a'$ contains the factor $(1-\alpha)^{-1}$).

\item[Propositions 22.1, 22.2\,(a), (c), and 22.4.]
First layer:
$b>0$ in Proposition~22.2\,(a);
$\rho\ge0$ and $w_i>0$, in the form that the weights are real and $w_i\ge0$, together with the closure relation (I.1.1), in Proposition~22.2\,(c);
$M\ge1$, and $L$, $M$ odd, in Proposition~22.4\,(b), so that by the fact (g0) of \cite[Section~17.4.1]{QEDI} the $M$-cubes form a partition of $T_0$.
Second layer:
for Proposition~22.4\,(a), the second layer of Proposition~22.3 at $k=0$, through Proposition~22.3\,(a);
for Propositions~22.1, 22.2\,(a), (c) and~22.4\,(b), none.

\item[Propositions 22.2\,(b) and 22.3.]
First layer:
$b>0$;
$L>1$, which gives $e_k<e_{k+1}\le e_0$ and $\mu<1$;
$e_0<1$ (the first inequality of (A8)).
Second layer.
Through \cite[Lemmas~5.17--5.19]{QEDI}: (A8), $L\ge5$, $L^{-\frac{3}{4}}\le\frac{1}{2}$ and $\varepsilon\le\frac{1}{200}$.
Through \cite[Proposition~5.38(ii)]{QEDI}:
(A1) and (A2);
the inequalities
\[
8\,C_B\cdot p_k\le e_k^{-\frac{3}{4}+3\varepsilon},\qquad 8\,e_k^{\frac{3}{4}-\frac{7\varepsilon}{2}}\le1
\]
for all $e_k\le e_0$, which are part of (A10) in the form of Section~I.8.4;
$e_k<1$ (by (A8));
$\varepsilon\le\frac{1}{200}$;
$L^{\frac{3}{4}}\ge\frac{8}{3}$;
$\gamma_{\mathrm{w}}\le\gamma^{P}\le\gamma_{P}$;
$C_{w}\ge C^{P}\ge C_{P}$;
$\gamma_{\mathrm{w}}>0$;
$0\le\alpha<1$, a hypothesis of \cite[Theorem~15.14]{QEDI}, that is, of \cite[Proposition~5.32]{QEDI} ($C_{P}$ contains the factor $(1-\alpha)^{-1}$).
The conditions $M\ge M^{P}$ and $\alpha_{0}\le\alpha^{P}$, which are quantifiers of \cite[Theorem~15.37]{QEDI}, the theorem that collects the properties stated in \cite[Propositions~5.31--5.41]{QEDI}, enter only in \cite[Section~15.10]{QEDI}, and \cite[Propositions~15.30, 15.29, 15.24 and Theorem~15.14]{QEDI} are proved without them.
Through (G3.T-2a): the conditions listed in the entry of Proposition~21.4\,(d2).
Through Propositions~21.1\,(d), (f), 21.4\,(a), (b), (c) and~21.5\,(b), (c): the conditions of the first entry of this list other than $\rho\ne0$, $\rho^2\ne1$ and $w_3\ne0$, that is,
$L$ odd, $L>1$;
$M$ a power of $L$ with $L/2<M$;
$b>0$;
$\rho\ge0$ and $w_i>0$;
the closure relation (I.1.1).
Through (G3.T-16), for $k\ge1$: the conditions of the Propositions~16.1--16.3.

\item[Propositions 23.1\,(a)--(f) and 23.2\,(a)--(d).]
First layer:
$L$, $M$ odd (for the fact (g0) of \cite[Section~17.4.1]{QEDI});
$M\ge1$;
$b>0$ and $L>1$, so that $b_{k}>0$;
$0<\rho<1$, in the proof of Proposition~23.2\,(c), where the diagonal entries $3\rho$ and $3\,(\rho\pm1)$ of $M_3$ do not vanish, and in Proposition~23.2\,(d), where the sums of the moduli of the entries of $M_3$ are at most $12\,(1+\rho)$, and $12\,(1+\rho)\le24$;
$0<\varepsilon\le2$, in the second sentence of Proposition~23.2\,(d); this is contained in the line (C7).
Second layer, through Proposition~21.1\,(c), in the second assertion of Proposition~23.1\,(e):
the closure relation (I.1.1) of the standard weights (line (C0)), on which Proposition~21.1\,(c) rests;
by Proposition~21.1\,(c), $F_\infty(-B)$ is invertible whenever $F_\infty(B)$ is, so that $\bar{H}_k(B)=\mathsf{C}^{-1}H_k(-B)\,\mathsf{C}$ is defined.

\item[Proposition 23.1\,(g).]
First layer: none.
Second layer, through (G3.T-9), that is, through Proposition~9\,(d):
(Z2.s4), which gives $\mathfrak{i}(B)\le1$ and $e_k\cdot\sup|dB|\le1$ for $B\in\tilde{R}_k+\mathfrak{X}_k$;
(L5);
(M1), which together with (L5) makes $M$ odd with $M>4L$;
(S1), that is, $0<\varepsilon\le\bar{\varepsilon}$ for the $\varepsilon$ fixed in the line (C7); this is among the data of \cite[Section~16.2.1]{QEDI} under which the display (G3) of \cite[Section~16.2.7]{QEDI} is stated;
$0\le\rho\le1$, so that $3\,(1+\rho)\le6$; this is contained in the condition $0<\rho<1$ of the line (C0);
the standard weights, through the closure relation (I.1.1) and the value $w_3=6$;
$\alpha\in(0,1)$, by the line (C1), in two forms: as $\alpha\ge0$, which gives $\eta_k^{\alpha}\le1$, and as $0<\alpha<1$, which is among the data of \cite[Section~16.2.1]{QEDI};
$b=1>0$ and $L>1$, so that $b_{k}$ is defined and positive and $b_{-}>0$;
$b_{-}\le b_{k}$.

\item[Proposition 23.2\,(e).]
First layer:
the second part of (A3), that is,
\[
96\,M^3\cdot e_k^{\frac{1}{4}}\le1;
\]
$\varepsilon\ge0$, which follows from $0<\varepsilon\le2$ and is thus contained in the line (C7), in the step (e-2) of the proof, where $\frac{1}{4}-\frac{\varepsilon}{2}\ge\frac{1}{4}-5\varepsilon$ is used;
$L\ge1$, by which the inequality $L^{\frac{1}{2}}\cdot q\le e_0$ of (ST1)$^{\mathrm{ch}}$ at $k=0$ gives $q\le e_0$, so that the conditions on $e_0$ hold at $v=-\log q$.
Second layer, through Proposition~23.2\,(c), (d): those conditions of the entry of Propositions~23.1\,(a)--(f) and~23.2\,(a)--(d) that enter these two clauses, namely
$L$, $M$ odd;
$M\ge1$;
$0<\rho<1$;
$0<\varepsilon\le2$.
Second layer, through Proposition~24.3\,(q1).
For the assertion of (q1) that Proposition~23.2\,(e) uses, the bound $\frac{1}{2}\,e_k^{-\frac{\varepsilon}{2}}$ on the three norms of $\mathscr{Z}(s')$:
the first inequality of (A1), that is, $p_{0,k}\le e_k^{-\frac{\varepsilon}{4}}$;
(A2);
$C_{w}\ge4\exp\left(2\gamma_{\mathrm{w}}\right)\cdot M^{27\alpha_{0}}$ (by the definition of $C_{w}^{(1)}$), which this bound needs at the level $0$, the only level of Proposition~23.2\,(e): by W17 the diagonal family $H_0(s')$ has the property W3 with the constant $4\exp\left(2\gamma_{\mathrm{w}}\right)\cdot M^{27\alpha_{0}}$, hence with $C_w$;
$L$ odd, $L\ge3$ and $M$ a power of $L$, the setting of \cite[Proposition~5.35]{QEDI}, that is, of \cite[Proposition~15.21]{QEDI}, from which the fact (P2) of Chapter~24 is taken.
For the two other assertions of (q1), which Proposition~23.2\,(e) does not use:
the third inequality of (A1), that is, $e_k^{\frac{5\varepsilon}{2}}\le\frac{1}{2}$;
(A5), together with $M\ge1$;
$\varepsilon\le\frac{1}{200}$ in the form $\frac{1}{2}-4\varepsilon>0$;
$\varepsilon\ge0$;
$e_k\le1$, by the first inequality of (A8);
these enter the inequalities between the radii and between the thresholds;
moreover, the first and the fourth inequality of (A8), together with $L\ge5$, as it stands, as hypotheses of \cite[Lemma~5.17(a)]{QEDI}, that is, of \cite[Lemma~17.4(a)]{QEDI}, by which the set $\tilde{R}_k+\mathfrak{X}_k$ is contained in $\tilde{G}_k$.

\item[Propositions 24.1--24.3.]
First layer:
$L$ odd, $M$ a power of $L$;
$L\ge4^8$, through its consequences $L\ge5$ and $L\ge4$;
$L^{-\frac{3}{4}}\le\frac{1}{2}$;
$\kappa\ge7\kappa_{0}+7$;
$L\cdot\left(\kappa-6\kappa_{0}-4\right)-1\ge\kappa$;
$\kappa_{\mathrm{w}}=\kappa+\kappa_{0}+3$;
$\kappa_{0}\ge c_{X}\ge1$ and $K_{0}\ge1$ (by the values in the line (C0));
$M>4L$;
$M^{\alpha_{0}}\ge1+\exp(\hat{\kappa})$;
$C_{w}\ge4\exp\left(2\gamma_{\mathrm{w}}\right)\cdot M^{27\alpha_{0}}$ (by the definition of $C_{w}^{(1)}$);
$C_{w}\ge1$ and $c_{\gamma}\ge1$;
$0<1-\alpha<\varepsilon\le\frac{1}{200}$;
(A1), with its five inequalities;
(A2);
(A3), both parts;
(A4);
(A5);
(A6), with the numeral $112$;
(A7);
(A8);
$3M\cdot e_0^{\varepsilon}\le1$, for (I.Lc2)$^{\mathrm{ch}}$;
$\ell$, $\ell/L\in\mathrm{Cl}$, which is a datum of (ST1)$^{\mathrm{ch}}$, used through the facts that $M$ divides $\ell$, that $\ell/(LM)$ is an integer, and that $(2L+1)\cdot M<\ell$ and $2M<\ell/L$.
Second layer.
Through (G3.T-16), (G3.T-17), (G3.T-18) and (G3.T-48):
(A9);
the first inequality of (C3),
\[
L^{\frac{1}{8}-\frac{3\bar{\varepsilon}}{2}}\ge2400;
\]
$1-\varepsilon<\alpha<1$.
Through (G3.T-13):
$\rho<1$ (line (C0));
(Z2.s1);
(Z2.s4);
(Z2.s5), with $C_{d}$ formed with the constant $C_{w}$ of $\mathrm{par}$;
the inequalities
\[
v\ge9,\qquad 8\,M\cdot\left(\frac{v}{4}+4\right)\cdot\exp\left(-\varepsilon v\right)\le1
\]
for $v\ge-\log e_0$;
$r\ge2$;
$L\cdot\left(\kappa-\kappa_{0}-1\right)\ge\kappa$;
$\kappa\ge\kappa_{0}$, $\varepsilon\le\frac{1}{8}$ and $\kappa_{\mathrm{w}}-1\ge\kappa+1$, which follow from $\kappa\ge7\kappa_{0}+7$, $\varepsilon\le\frac{1}{200}$ and $\kappa_{\mathrm{w}}=\kappa+\kappa_{0}+3$;
$\ell\ge M\cdot r_k$, which is a datum.
Through (G3.T-9): the conditions of the entry of Proposition~23.1\,(g).
Through (G3.T-2a): the conditions of the entry of Proposition~21.4\,(d2).
Through \cite[Proposition~5.38(ii)]{QEDI}, that is, \cite[Proposition~15.30(ii)]{QEDI}: the conditions listed under this proposition in the entry of Propositions~22.2\,(b) and~22.3, from (A1) and (A2) to $0\le\alpha<1$.
Through \cite[Propositions~5.35 and~5.36]{QEDI}, that is, \cite[Propositions~15.21 and~15.26]{QEDI}, from which the facts (P1)--(P4) of Chapter~24 are taken:
$L$ odd and $L\ge3$ (the setting of \cite[Section~15]{QEDI});
$M$ a power of $L$;
$p_{0,k}\ge1$, the hypothesis of the bound on $\delta\varepsilon^{0}_{k}$ in (P4); since $p_{0,k}=v_k^{p_0}$ with $v_k=-\log e_k$, it holds because $v_k\ge1$ and $p_0>0$:
the first follows from $e_k\le e_0\le\exp(-1)$, which is (Z3.s1) of Section~I.8.4 and is contained in the first inequality of (A8) together with $\varepsilon\le\frac{1}{200}$;
the second follows from the range $p_0>3r+2>0$ of the line (C7).
Through \cite[Theorem~15.37]{QEDI}, by which Chapter~24 obtains \cite[Propositions~5.31, 5.32 and~5.34--5.36]{QEDI} for the constants of Section~I.8.2:
$M\ge M^{P}$ ((M6)) and $\alpha_{0}\le\alpha^{P}$ (by $\alpha_{0}=\min\left(\alpha^{F},\alpha^{P}\right)$, line (C4));
these are quantifiers of that theorem, they hold for the constants, and the proofs of \cite[Theorem~15.14 and Propositions~15.21, 15.24, 15.26, 15.29, 15.30]{QEDI} do not use them.
Through Propositions~22.2\,(b), 22.3, 22.4\,(a), 23.1 and~23.2: the conditions of the entries of these statements above (for Proposition~22.4\,(a), those of Proposition~22.3 at $k=0$), among them $\rho\ge0$, $w_i>0$, the closure relation (I.1.1) and $0<\rho<1$.

\item[Propositions 25.1--25.3.]
First layer:
$L$, $M$ odd;
$L>1$, in the rows nos.~7 and~8 of the table of Proposition~25.1 (where $\eta_k\le L-\eta_k$ and $1\le L-1$ are used) and for $b_{k}>0$;
$M$ a power of $L$;
$L-1<M$;
$b>0$, in the rows nos.~6, 7, 9, 10 and in (K3) of Proposition~25.1;
the closure relation (I.1.1), through (I.m1), in the row no.~3 and in (K4) of Proposition~25.1.
Proposition~25.2\,(o10), and the part of Proposition~25.3\,(b5) that rests on it, are stated for $p_{0,k}>0$;
this is a restriction on the datum $e_k$; it holds for $e_k<1$, hence for every $e_k\in(0,e_0]$, because $e_0<1$ by the first inequality of (A8), or by the first inequality of (A10).
Second layer: \cite[Proposition~5.58\,(t7), (t10)]{QEDI}.
\end{description}

\begin{named}{Remark}
\begin{itemize}
\item[(1)] Among the conditions of the first layer, (A6) is the only one whose numeral differs from that of the condition with the same label in \cite[Section~5.6.2]{QEDI}: the factor $16$ there is $112$ here.
\item[(2)] The parameter $\rho$ enters the first layer through its range and through the constants $\gamma_{\mathrm{w}}$, $M$, $C_{w}$, $c_{\gamma}$; the constant $\gamma_{\mathrm{w}}$ stands explicitly in the lower bound on $C_{w}$ in the entry of the Propositions~24.1--24.3.
Through its range it enters as $\rho\ge0$ in Propositions~21.1\,(e), 21.5 and~22.2\,(c);
as $\rho\ne0$, $\rho^2\ne1$ in the last assertion of Proposition~21.4\,(a) at $k=0$, and in the proof of Proposition~23.2\,(c), where the diagonal entries $3\rho$ and $3\,(\rho\pm1)$ of $M_3$ do not vanish;
and as $0<\rho<1$ in Proposition~23.2\,(d), where the sums of the moduli of the entries of $M_3$ are at most $12\,(1+\rho)$, and $12\,(1+\rho)\le24$.
It enters $M$ through the lower bounds (M3), (M5) on $M$ (Proposition~14.4\,(c)), and $c_{\gamma}$ through $\gamma_{\mathrm{w}}$.
The constant $C_{T}$ does not contain $\rho$ (Section~I.8.2).
\item[(3)] No condition on $\kappa$, $L$ or $M$ in the list of Section~G3.6.1 contains $C_{w}$, $C_{T}$, $\varepsilon$ or $e_0$.
Where an inequality between $L$ and $\varepsilon$ is needed, it follows from the form of (C3) in which $\bar{\varepsilon}$ stands, together with $\varepsilon\le\bar{\varepsilon}$.
\item[(4)] The weight $\mathsf{w}_W=\sqrt{48}\le7$ occurs at $k=0$ only, since $\mathsf{w}_W=1$ for $k\ge1$; the numeral $7$ is the bound for it that enters $C_{W}$.
The constant $C_{W}=28\,C_{w}^2\cdot c_{\gamma}$ is used at all levels (I.L0.19).
\end{itemize}
\end{named}

\subsubsection*{G4.A.6. Conditions on the constants in the Propositions 26.1--32.3 and 49}

Each entry of the following list names a condition, the line of Section~I.8.2 that imposes it (the mark $e_1$ means that the condition follows from the rows (e1)--(e8) of Proposition~14.2\,(b), by which $e_1$ is chosen), and the statements in which it is used.
The facts (Sc0)--(Sc6) are those of Section~G4.0.2.

\begin{itemize}
\item
Condition:
$L$ odd;
$M$ a power of $L$ with $M\ge L^2$, $M>4L$, $\gamma_{\mathrm{w}}\cdot M\ge1$ and $M\ge M^{P}$;
$0<\alpha_{0}\le\frac{1}{4}$ and $\kappa_{\mathrm{w}}\ge0$;
$K_{0}\ge1$;
$C_{w}^{(1)}\ge4\exp\left(2\gamma_{\mathrm{w}}\right)\cdot M^{27\alpha_{0}}$;
the conditions of Proposition~11, which are those of Proposition~2.8:
(L5), that is, $L\ge5$ besides $L$ odd;
(M1);
(M2);
the part $M\ge108\pi\cdot C'_{V}\cdot a'$ of (M5);
(Z2.s1) and (Z2.s2);
$\alpha<1$;
the conditions of Proposition~8, which are those of Proposition~3, that is, the conditions of Proposition~2.8 with (M5) in full.
Lines: (C3), (C5), (C4), (C0), (C6); for the conditions of Propositions~11 and~8, (C3) for (L5), (C5) for (M1), (M2) and (M5), (C8) for (Z2.s1) and (Z2.s2), and (C1) for $\alpha<1$.
Used in Proposition~26.3:
the conditions of Proposition~11 in clause~(b), at W1, through Proposition~11\,(b);
the conditions of Proposition~8 in clause~(c), and through it at W17 in clause~(b), through Proposition~8\,(b).
Of these conditions, $L$ odd, $M$ a power of $L$ and $M>4L$ are used also in Propositions~26.2 and~26.4.

\item
Condition:
$M\ge M^{F}$ at $\kappa'=\kappa_{\mathrm{w}}+\kappa_{0}$;
$M\ge M^{P}$;
(Z2.s1)--(Z2.s4);
the hypotheses $0\le\alpha<1$, $L$ odd, $L\ge3$, $0<\alpha_{0}\le\frac{1}{4}$, $M$ a power of $L$ and $M\ge L^2$ of \cite[Theorem~15.35]{QEDI}.
Lines: (C5), (C8) for the first three; (C1), (C3), (C4), (C5) for the hypotheses.
Used in Proposition~26.5.

\item
Condition: (2.1), that is, $r_{*}>31L+16$.
It is a consequence of $e_K\le e_1$.
Used in Proposition~26.6 (through (Sc4)), in Propositions~27.1 and~27.2\,(d), in Proposition~27.3 (where $2M\le\ell_1-14\,d_U\cdot M$ is used), and in Proposition~49\,(b).

\item
Condition:
the third clause of (A1);
(A5);
the first clause of (A8), and in the domains also its fourth clause;
(S1);
(M4);
(L5) and (M1);
$\gamma_{\mathrm{w}}>0$ and $\kappa_{\mathrm{w}}>0$ (rows of Proposition~14.2\,(a)).
Lines: (C8) for (A1), (A5), (A8); (C7) for (S1); (C5) for (M4); (C3), (C5) for (L5), (M1); (C4r), (C4) for $\gamma_{\mathrm{w}}$, $\kappa_{\mathrm{w}}$.
Used in Proposition~27.5:
in (c2);
(A8) also in the domains, through the fact (TT.3) of Chapter~27b, Section~0.1: its first clause, which gives $e_k^{\varepsilon}\le\frac{1}{4}$, and, since (TT.3) applies \cite[Lemma~5.17(a)]{QEDI}, that is, \cite[Lemma~17.4(a)]{QEDI}, its fourth clause together with $L\ge5$, as it stands;
(M4) in (c6);
(L5) and (M1) through the fact (g0) of \cite[Section~17.4.1]{QEDI} and through $b_{-}>0$;
$\gamma_{\mathrm{w}}$ and $\kappa_{\mathrm{w}}$ in (b), where the factor $\exp\left(-\gamma_{\mathrm{w}}d-\kappa_{\mathrm{w}}d_M(X)\right)$ of W7$^{\mathrm{ch}}$ (Section~G4.0.4) is at most $1$.

\item
Condition:
\[
c_{\mathrm{fl}}\cdot L^{-\frac{3}{8}+6\bar{\varepsilon}}\le\frac{1}{16},\qquad L^{-\frac{5}{8}}\le\frac{1}{4};
\]
$c_{\mathrm{fl}}\ge1$;
$e_0^{\varepsilon}\le\frac{1}{4}$.
Lines: (C3), (C0), and (A8).
Used in Propositions~28.1 and~28.2.

\item
Condition: (fl), in the form $\mathrm{sm}_k\le\frac{1}{16}$.
Line: $e_1$.
Used in Proposition~28.2.

\item
Condition: the first two clauses of (T1).
Line: $e_1$.
Used in Proposition~31.2.

\item
Condition: $e_1\le\frac{1}{3}$.
It gives $w^{\mathrm{m}}_0<1$ for all $N\ge0$; the third clause of (T1) is the case $N\ge N_{\min}(R,e)$ of this.
Line: (C9).
Used in Proposition~28.3\,(d).

\item
Condition: (T2) and (T3).
Line: $e_1$.
Used in Proposition~32.1; (T2) is used also in Propositions~32.2 and~32.3.

\item
Condition: (T4) and (T5).
Line: $e_1$.
Used in Proposition~32.2.

\item
Condition:
the first clause of (A1);
(A2);
the first and the fourth clause of (A8);
the third and the fourth clause of (A10);
$\gamma_{\mathrm{w}}\cdot M\ge16$, through its consequence $\gamma_{\mathrm{w}}>0$;
$C_{P}\le C_{w}$ and $\gamma_{\mathrm{w}}\le\gamma_{P}$ (by the definitions);
$L$ odd, $L\ge5$;
$M$ a power of $L$, $M\ge L^2$;
$0<\varepsilon\le\bar{\varepsilon}$;
$0\le\alpha<1$.
Lines: (C8) for the clauses of (A1), (A2), (A8), (A10); (C5) for $\gamma_{\mathrm{w}}\cdot M\ge16$; (C4), (C4r), (C6) for $C_{P}\le C_{w}$ and $\gamma_{\mathrm{w}}\le\gamma_{P}$; (C3) for $L$; (C5) for $M$; (C7) for $\varepsilon$; (C1) for $\alpha$.
Used in Propositions~29 and~32.1, as hypotheses of \cite[Proposition~5.38(i), (ii)]{QEDI}, of \cite[Propositions~5.37 and~5.32]{QEDI}, on which it rests, and of \cite[Lemma~5.17(a)]{QEDI}, that is, of \cite[Lemma~17.4(a)]{QEDI}.

\item
Condition:
(e1);
$L\ge33$ (from (L1));
$M$ a power of $L$.
Lines: $e_1$; (C3); (C5).
Used through (Sc0)--(Sc6) in:
Proposition~26.6;
Propositions~27.1--27.6, through the standing data (TT) of the Propositions~27.m:
there (Sc2) gives $k_m\le k_{\mathrm{top}}$, and (Sc3), together with (Sc1) and with the condition (2.1), which (Sc6) supplies, gives $\ell_1\ge M\cdot r_{*}>(31L+16)\cdot M$;
in this way the first and the third inequality of (e1) enter at $v=v_K$, where, together with $r>0$, they give the hypotheses $v_K\ge1$ and $\left(1+\frac{\mathrm{lg}}{v_K}\right)^r\le L$ of \cite[Proposition~19.1]{QEDI}, the second inequality of (e1) enters at $v=v_{*}$, and the condition that $M$ is a power of $L$ enters through the proof of \cite[Proposition~19.1]{QEDI};
moreover, the consequence $L\ge13$ of $L\ge33$ enters through Proposition~26.6\,(d), which supplies in (TT) the data of Proposition~26.1, and through (Sc4), on which that clause rests;
in Proposition~27.4\,(f), in addition, the first inequality of (e1) gives, through (Sc3) and (Sc6), $v_k\ge v_K\ge8$, so that $p_{0,k}>0$;
Proposition~28.1, through (Sc1);
Proposition~28.2, through (Sc1) and (Sc6);
Proposition~28.3, through (Sc0)--(Sc2);
Proposition~30.3\,(f), in its closing part ``In the flow'';
Proposition~31.2;
Propositions~32.1--32.3;
Proposition~49\,(b), through (Sc0), (Sc1), (Sc4), with the first and the third inequality of (e1) at $v=v_K$ and with the consequence $L\ge13$ of $L\ge33$;
and further in Propositions~20.3 and~48.6\,(b).

\item
Condition:
$r$ an integer $\ge2$, in the form $r>0$, that is, $r\ge1$;
$p_0$ an integer $\ge1$, which follows from $p_0>3r+2$.
Line: (C7).
Used at every place of the preceding entry, among them Proposition~27.4\,(f) and the Propositions~27.1--27.6, in which (Sc2) and (Sc3) enter through the standing data (TT):
the condition $r\ge1$ inside each of (Sc0)--(Sc5), and so in every clause of \cite[Proposition~19.1]{QEDI}, namely in the step
\[
\left(1+\frac{\mathrm{lg}}{v_K}\right)^r\le\left(1+\frac{3\,\mathrm{lg}}{v_K}\right)^r
\]
from the third inequality of (e1) to the hypothesis $\left(1+\frac{\mathrm{lg}}{v_K}\right)^r\le L$, under which all clauses of \cite[Proposition~19.1]{QEDI} stand, and in the clauses (a), (b) of that proposition, from which its clauses (c)--(g) are derived;
the condition $r\ge1$ also in the assertion of (Sc6) that the conditions of Section~G4.0.3 hold, which is derived with the clauses (c), (e) of \cite[Proposition~19.1]{QEDI}; the first display of (Sc6) does not require it;
the condition $p_0\ge1$ in Proposition~31.2, in a binomial expansion.

\item
Condition:
$0<\rho<1$;
the standard weights, through $w_i>0$, the closure relation (I.1.1) and $w_3=6$.
Line: (C0).
Used in:
Proposition~29, where $0<\rho<1$ enters a bound of the determinant of level $0$, through the finiteness of this bound;
Propositions~30.1--30.3 ($\rho>0$, $\rho\le1$; $w_i>0$; (I.1.1); $w_3=6$);
Proposition~31.1 ($\rho\ge0$, $\rho^2\ne1$, $-1\le\rho\le1$; $w_i>0$; (I.1.1));
Proposition~49\,(c)--(f) ($\rho\ge0$, $\rho^2\ne1$; $w_i>0$; (I.1.1); and, through Proposition~38.1\,(d), the value $w_3=6$, with $w_0+w_1+w_2=\mathrm{sw}-6$).

\item
Condition:
$L$ odd and $L\ge5$, that is, (L5);
$M$ a power of $L$ and $M\ge L^2$, that is, (M1).
Lines: (C3); (C5).
Used also in Propositions~27.1--27.4, 27.6 and~31.2, and in all clauses of Proposition~49 ($\mathrm{per}$ is odd and $n\ge5$; $T^{*}(e)$ is an integer $\ge1$).
The condition $L\ge5$ is used in Proposition~28.1 as it stands, together with the first and the fourth clause of (A8), as a hypothesis of \cite[Lemma~5.17(a)]{QEDI}, that is, of \cite[Lemma~17.4(a)]{QEDI}, on which, together with \cite[Lemma~5.18]{QEDI}, the last sentence of Definition~24.D\,(D6) rests;
it is used further in the form $L\ge4$, which follows also from $L^{-\frac{5}{8}}\le\frac{1}{4}$, in Proposition~28.2, and in the form $L>1$ in Propositions~28.1--28.3.
The conditions $L$ odd and $L\ge5$ alone are used in Proposition~30.3\,(f);
in the form $L\ge1$, which gives $\eta=L^{-N}\le1$, in Proposition~30.2\,(b);
in Proposition~26.1, in the form $L$ odd, $L\ge3$, which is the setting of \cite[Proposition~5.35]{QEDI};
and in Proposition~32.2, where $\mathfrak{N}=L^{k_{*}}$ is odd and $\ge3$.
The condition $L$ odd alone is used in Proposition~26.6\,(d), where $\mathrm{Jr}=R_{\mathrm{top}}/R$ is odd.
The condition that $M$ is a power of $L$ alone is used in Proposition~28.3.
The condition $M\ge L^2$ is used also in Proposition~32.3: it gives $T^{*}(e)\ge1$, so that $R^{*}(e)$ and $R_{\mathrm{top}}(e)$ lie in $\mathrm{Tor}(e)$.

\item
Condition: $L^{-\frac{3}{4}}\le\frac{1}{2}$ (from (L1)).
Line: (C3).
Used through \cite[Lemma~5.18]{QEDI}, that is, \cite[Lemma~17.5]{QEDI}, in Propositions~27.6 and~28.1.

\item
Condition: (S1), in the form $0<\varepsilon\le\bar{\varepsilon}$.
Line: (C7).
Used in Proposition~27.6, in the form $\varepsilon\le\frac{1}{200}$, as a hypothesis of \cite[Lemma~17.5(a), (c)]{QEDI}, of which \cite[Lemma~5.18]{QEDI} is a special case;
and in Propositions~28.1--28.3, 32.2 and~32.3.

\item
Condition: the first clause of (A8); $\alpha>0$.
Lines: (C8); (C1).
Used in Proposition~27.6, as hypotheses of \cite[Lemma~17.5]{QEDI}, of which \cite[Lemma~5.18]{QEDI} is a special case:
the first clause of (A8), through $e_k^{\varepsilon}\le\frac{1}{4}$, together with $L^{-\frac{3}{4}}\le\frac{1}{2}$, in \cite[Lemma~17.5(b)]{QEDI}, the inclusion $L^{-\frac{3}{4}}\cdot\tilde{R}_k\subset\tilde{R}^{-}_k$;
$\alpha>0$ in \cite[Lemma~17.5(a)]{QEDI}.
The first clause of (A8) enters Proposition~27.6 also through the fact (TT.3) of Chapter~27b, Section~0.1, where it gives $e_k^{\varepsilon}\le\frac{1}{4}$ and $e_k<1$;
since (TT.3) applies \cite[Lemma~5.17(a)]{QEDI}, that is, \cite[Lemma~17.4(a)]{QEDI}, the fourth clause of (A8) and $L\ge5$, as it stands, enter Proposition~27.6 as well.

\item
Condition: $e_1\le e_0$.
Line: (C9).
Used in Propositions~26.6 and~32.1.
\end{itemize}

In the Propositions~26.1--32.3 and~49 no condition on $\kappa$, $L$ or $M$ arises beyond those of the lines (C2), (C3), (C5).
The conditions on $e_0$, and the rows (e1)--(e8) of Proposition~14.2\,(b) from which the conditions marked $e_1$ follow, contain neither $N$ nor $\theta$ nor the pair $(m_{\mathrm{in}},\varepsilon^{\mathrm{v}}_{\mathrm{in}})$; by contrast, the condition (fl) and the number $w^{\mathrm{m}}_0$ in (T1) contain $N$ (Section~G4.0.3).
The constants $\gamma_{\Sigma}$, $C_{\Sigma}$, $C_{6}$ and the constant $\kappa_{(0)}$ of the line (C0r) are chosen before $e_0$.

\subsubsection*{G5.9. Order of the constants and conditions in the Propositions 33.1--37.3}

\paragraph{G5.9.1. Order and dependence.}
The following list gives, line by line of Section~I.8.2, the constants and what each of them may depend on.

\begin{description}
\item[(C0r).]
$C_E(\rho)$; $\mathrm{Csy}_{\mathrm{ak}}(\rho)$ for $|\mathrm{ak}|\le4$; $\mathrm{Cpv}_0(\rho)$.
They depend on $\rho$; they are obtained by compactness and are not explicit.
\item[(C6$'$), after (C6).]
$\mu_{-}$ (a constant of the line (C4)), $\mathrm{ssum}$, $\mathbf{S}$, $\mathrm{hL}$, $C_K$, $C_K'$, $C_H$, $\mathrm{CVj}$, $\theta_{\mathrm{CT}}$, $\mathrm{vt}$, $\mathrm{gGa}$, $h_{G}$, $\mathrm{gGa}_{1}$, $\mathrm{alG}$, $c_{\tau}$, $\mathrm{CWj}$ of Section~G5.1.8.
They may depend on $b$, $L$, $M$, $K_{0}$, $C^{F}$, $\gamma^{F}$, $C_{d}$, and through these on $\rho$; they do not depend on $\bar{m}$, $\beta$, $e$, $N$.
\item[(C10), after $\bar{m}$ and $\beta$.]
$h$;
$D_{\mathrm{be}}$ and $D_\beta=D_{\mathrm{be}}+4$;
$C_{\beta}$, $\mathbf{D}$, $\mathrm{Cb}$;
$C_{\mathrm{fr}}$;
$\mathrm{Cah}_i(\beta)$;
$\mathrm{Cpv}_1(\rho,\beta)$, $\mathrm{Cpv}_2(\rho,\bar{m},\beta)$, $\mathrm{epv}(\,\cdot\,)$, $\mathrm{Cpv}_{\beta}$;
$C_{\mathrm{ct}}(\rho,\beta)$, $C_{1}^{\mathrm{N2}}(\rho,\bar{m},\beta)$, $C_{2}^{\mathrm{N2}}(\bar{m},\beta)$.
Their dependence is the one indicated.
\item[(C10), next.]
$e_{\sharp}(\beta)$.
It may depend on the constants of the lines (C0)--(C9) and (C6$'$), on $\bar{m}$, and on $C_{\beta}$, $\mathbf{D}$, $D_\beta$.
\item[(C11).]
First $e$; then the quantities of Section~G5.1.8 that depend on $e$, with the exception of $\varpi_{\mathrm{sep}}$; then $N$; then the placement, $\theta$ and $A$.
The number $\varpi_{\mathrm{sep}}$ of Section~G5.1.8, which contains the number $d_{\mathrm{sep}}$ of Section~G5.1.4, is formed after $N$ and the placement.
\end{description}

\paragraph{G5.9.2. Conditions.}
On $e_0$ (line (C8)) the condition (A8) enters directly, through the first, the third and the fourth of its four inequalities.
Through \cite[Proposition~15.30]{QEDI}, that is, in Proposition~35.4\,(SF5) only, there enter the third and the fourth inequality of (A10), at $v=v_{*}$, as the two displayed hypotheses of \cite[Proposition~15.30]{QEDI}, and the conditions (A1) and (A2), together with $C_{P}\le C_{w}$ and $\gamma_{\mathrm{w}}\le\gamma_{P}$ (Proposition~14.2\,(f)).
The first two inequalities of (A10) are not needed in the Propositions~33.1--37.3.
On $e_1$ (line (C9)): the condition (2.1), that is, $r_{*}>31L+16$, the first two clauses of (T1), and (T3); they hold by the rows (e1) and (e4) of Proposition~14.2\,(b), which are the conditions of \cite[Section~19.1.2]{QEDI} formed with the constants of Section~I.8.2.
On $L$ and $\kappa$: the three consequences of (C2) and (C3) listed in Section~G5.1.9, namely
\[
\kappa\ge7\kappa_{0}+7,\qquad L\cdot\left(\kappa-\kappa_{0}-1\right)\ge\kappa,\qquad c_{\mathrm{rb}}\cdot L^3\cdot\exp\left(-\left(L\cdot\left(\kappa-\kappa_{0}-1\right)-\kappa\right)\right)\le\frac{1}{2}.
\]
On $M$: $M\ge L^2$, and $\gamma_{\mathrm{w}}\cdot M\ge16$ together with $\gamma_{\mathrm{w}}\le\gamma^{F}$.
Besides these, the following are used, as ranges of Proposition~14.2\,(a), as rows of Proposition~14.2\,(b), or as consequences stated in Proposition~14.2\,(f):
$L$ odd and $L\ge5$, that is, (L5);
$L\ge13$ (from (L1));
$M$ odd;
$e_1\le\frac{1}{3}$;
of the line (C7), $r\ge2$, $p_0\ge2$, $p>2p_0$, $0<s<1$, $\sigma>0$ (that is, (S2)) and $0<\varepsilon\le\frac{1}{200}$ (that is, (S1));
$0<\alpha<1$ (line (C1)),
with both of its bounds in Proposition~33.1\,(c) (there $0<1-\alpha<1$ gives $d_X^{1-\alpha}\le1$ for $d_X<1$ and $d_X^{1-\alpha}\le M^{1-\alpha}\le M$ for $1\le d_X\le M$),
in the form $\alpha>0$ in Proposition~33.2 (there $d_X^{-\alpha}\le1$ for the long pairs),
in the form $\alpha<1$ in Proposition~35.3\,(c1) (there $|c-c'|\le|c-c'|^{\alpha}$ for $0<|c-c'|\le1$),
and in the form $0\le\alpha<1$ in Proposition~35.4\,(SF5), through \cite[Theorem~15.14]{QEDI};
$C_{0}\ge1$ and $C_{0}\ge C_{P}\cdot c_{\mathrm{bd}}(\gamma_{P})$, by the definition of $C_{0}$.
Further, the arguments use these conditions of the lines (C0) and (C10):
$0<\rho<1$ (line (C0)), in the forms $3\,(1+\rho)\le6$ and $\rho\ge0$ in Propositions~33.1\,(b) and~33.4\,(b), in the form $\rho\ge0$ in Proposition~35.7\,(a), where $|M_i(x)|_2=\frac{w_i}{2}\,(1+\rho)$ by (I.m6), in the form $\rho>0$ in Proposition~36, where $\rho\in(0,1)$ enters also through Proposition~37.1, in the forms $\rho\ge0$, $\rho>0$, $\rho\le1$ in Propositions~37.1 and~37.2, and, through Propositions~36 and~37.2, in Proposition~37.3;
the standard weights (line (C0)):
through the closure relation (I.1.1), (I.m1) and $w_3=6$ in Propositions~36, 37.1 and~37.2, and in Proposition~37.1 also through their values $w=(6,6\sqrt{2},6\sqrt{2},6)$ themselves, which determine the gap $\lambda_1$;
in Proposition~35.7\,(a) through the closure relation, by way of (I.m1) and (I.2.5), through $w_i>0$ and through $w_3=6$;
in Propositions~33.1\,(b) and~33.4\,(b) through $w_3=6$, in $M_3(x)=3\,\left(\rho-\gamma\cdot(ge_1)\right)$;
$\bar{m}\in(0,1]$ (line (C10)), in Propositions~36, 37.1 and~37.2, and, in the form $\bar{m}\le1$, that is, $h\ge1$, in Propositions~35.9, 35.10 and~37.3, in Proposition~35.1\,(a), where $e=e_K\cdot h^{-\frac{1}{2}}\le e_K$, and in Proposition~35.7\,(a), where $n_*=h\cdot L^{T^{*}(e)+N}\ge L\ge5$, so that Proposition~31.1 applies.
The arguments use, moreover, the bounds $K_{0}\ge1$, $C^{F}\ge1$ and $0<\gamma^{F}\le1$, which hold by (G5.T-1) of Section~G5.7 and by Section~G5.1.8, and the bound $c_{\tau}\le\frac{1}{8}$, which holds by the definition of $c_{\tau}$ in the line (C6$'$).
The Propositions~33.1--37.3 impose no new condition on $\kappa$, $L$, $M$, $\varepsilon$, $s$, $e_0$ or $e_1$.

On $e_{\sharp}(\beta)$ (line (C10)): the conditions (g1)--(g9) of Section~G5.1.9.
The conditions (g1)--(g8) correspond to the smallness conditions (sm1)--(sm8) of \cite[Section~20.1.8]{QEDI}, and compared with these the following holds.
The constants $C_K$, $C_K'$, $C_H$, $c_{\tau}$, $C^{F}$, $\gamma^{F}$, $C_{0}$, $C_{d}$ and $M$ have other values;
of these, $C_H$ enters the conditions through $\mathrm{CVj}$ and $\mathrm{CWj}$, and $c_{\tau}$ through $\tau_s=c_{\tau}\cdot\ell_*$.
The numerals $5$, $4$, $\frac{1}{2}$, $48$, $52$, $8$, $2$ in the conditions are the same.
The conditions (g2) and (g5) are written with $M\cdot r_{*}/(8L)-1$ in the place of $d_{\mathrm{sep}}$.
The condition (g8) has the same form as (sm8).
Finally, (g9) has no counterpart among (sm1)--(sm8).

\paragraph{G5.9.3. Inequalities between the constants that the proofs use.}
For the values of Section~G5.1.8 the following inequalities hold; the reason is given after each of them.
\begin{itemize}
\item $C_K\le C_K'$, because $226\le1328$.
\item $C_K\le C_H\cdot\mathrm{ssum}(\gamma^{F}/2)^2$: both sides contain the factor $\mathrm{ssum}(\gamma^{F}/2)^2$, and $226\,(C^{F})^2\le203\,M\cdot(C^{F})^2$ because $M\ge2$.
\item $149\,(C^{F})^2\le160\,(C^{F})^2\le C_H$, because $160\le203\,M$.
\item $C_H^2=41209\,M^2\cdot(C^{F})^4$, and this number is at least $31095\,M\cdot(C^{F})^3$, at least $32803\,(C^{F})^3$ and at least $21645\,(C^{F})^3$, because $M\ge1$ and $C^{F}\ge1$.
\item $\mathbf{S}\ge1$, because $C^{F}\ge1$ and $\mathrm{ssum}(\gamma^{F}/2)=(24/\gamma^{F})^3\ge1$, the latter since $\gamma^{F}\le1$.
\item $13.5\,\mathbf{S}+1314\,\mathbf{S}^2\le C_K'$, because $\mathbf{S}\ge1$.
\item $146\,c_{\tau}\cdot\mathbf{S}\le\frac{1}{2}$, because $c_{\tau}\le\frac{1}{292\,\mathbf{S}}$.
\item The two inequalities
\begin{gather*}
8\,C_K+3\,K_{0}+1200\,C_H\cdot\mathrm{ssum}(\gamma^{F}/2)\cdot K_{0}\le\mathrm{CVj},\\
4\,C_K+82\,K_{0}+1200\,C_H\cdot\mathrm{ssum}(\gamma^{F}/2)^2\cdot K_{0}\le\mathrm{CVj},
\end{gather*}
because $C_H\ge2$, $K_{0}\ge1$ and $\mathrm{ssum}(\gamma^{F}/2)\ge1$.
\item $\left(16\,h_{G}^2+256\right)\cdot C_H^2\cdot\mathrm{ssum}(\gamma^{F}/2)^4\cdot K_{0}\le\mathrm{CVj}\cdot h_{G}^2$, because $h_{G}\ge1$ and $16+256\le10^3$.
\item $2\,\mathrm{CVj}\cdot h_{G}^2+2\,\mathrm{alG}\cdot\mathrm{CVj}^2\cdot h_{G}^4+2\,\mathrm{e}^4\cdot C_{d}\cdot M^3\cdot K_{0}\le\mathrm{CWj}$, by the definition of $\mathrm{CWj}$, since $2\le3$ and $2\,\mathrm{e}^4\le4\,\mathrm{e}^6$.
\item $\mathrm{gGa}^2\le\mathrm{gGa}_{1}$, because $\mathrm{gGa}^2=b_{-}^{-1}+\mu_{-}^{-1}$ and $\mathrm{gGa}_{1}=b_{-}^{-1}+4\,\mu_{-}^{-1}$.
\item $\mathrm{gGa}_{1}\ge1$, because $\mathrm{gGa}_{1}\ge b_{-}^{-1}$ and $b_{-}<1$.
\item $6\,C^{F}\le\mathrm{hL}$, because $\mathrm{hL}=6\,\mathbf{S}=6\,C^{F}\cdot\mathrm{ssum}(\gamma^{F}/2)$ and $\mathrm{ssum}(\gamma^{F}/2)\ge1$.
\item $\theta_{\mathrm{CT}}-\mathrm{vt}\ge\frac{\theta_{\mathrm{CT}}}{2}$ and $\gamma^{F}-\mathrm{vt}\ge\frac{7\gamma^{F}}{8}$, because $\mathrm{vt}\le\frac{\theta_{\mathrm{CT}}}{2}$ and $\mathrm{vt}\le\frac{\gamma^{F}}{8}$.
\item $c_{\tau}\le\frac{\mu_{-}}{48}<\frac{\mu_{*}}{24}$, by the definition of $c_{\tau}$ and because $\mu_{-}<\mu_{*}$; here $\mu_{*}$ is formed in the setting of Section~G5.1.1, in which $\mu_{-}<\mu_{*}$ by Proposition~14.2\,(f).
\end{itemize}

\subsubsection*{G6.Z4. Conditions on the constants in the Propositions 38.1--47.2 and 50 and in Theorem 44}

Each entry of the following list names a condition, the constant on which it is imposed together with its line of Section~I.8.2, and the statements in which it is used.
The label (G6.T-14) refers to Section~G6.Z1; (W3), (Z8.w), (Z10.s1), (Z10.s2), (Z10.s4) and (Pb) are the conditions of Proposition~14.2\,(d).

\begin{itemize}
\item
Condition: $L>8$.
Constant: $L$, line (C3).
Used in Proposition~47.2\,(a): there $\mathrm{kt}=8\,L^{T-T^{*}(e)}\ge L$ holds if and only if $T\ge T^{*}(e)+1$, and $L>8$ is needed for the direction ``only if''; Proposition~40\,(e) asserts the direction ``if'', which needs only $8L\ge L$.

\item
Condition: $L$ odd.
Constant: $L$, line (C3).
Used in Proposition~41 and in Proposition~47.2\,(a): $\mathrm{per}$ is odd, $\mathrm{hc}=m_{\mathrm{p}}+\frac{1}{2}$, and the walls with odd $j$ are link planes.
Used also in Proposition~43.2 and in Theorem~44, as a datum of \cite[Definition~2.36]{QEDII}, together with $L\ge3$.

\item
Condition: $L\ge5$.
Constant: $L$, line (C3).
Used in Proposition~38.3\,(e), where $n_*\ge5\ge4$;
in Proposition~40\,(e), where $n_*\ge L\ge4$;
in Proposition~42.2, where $N_{*}=8L^3\ge1000$, $10\,\mathrm{hc}\le2L\cdot\mathrm{hc}$ and $n_*\ge L\ge4$;
in Proposition~43.1, where $n_*\ge4$.

\item
Condition: $M$ a power of $L$, and $M\ge L^2$.
Constant: $M$, line (C5).
Used in (G6.T-14)\,(ii), (iv), for $T^{*}(e)\ge1$ and $\bar{m}\cdot R^{*}(e)\ge r_K(e)$;
and in Propositions~38.3\,(e), 41, 47.2\,(a), 40\,(e), 42.2, 43.1 and~50\,(h), where $T^{*}(e)$ is an integer $\ge1$ and $\mathrm{per}=n_*/L$ is an integer.

\item
Condition: $r\ge2$.
Constant: $r$, line (C7).
Used in Proposition~42.2, in the form $r>0$, so that $v_K^r>1$, and in Proposition~43.1\,(R1), where $v_K^r\le r_K<L\cdot v_K^r$ and $r_K\ge v_K$.

\item
Condition:
$0<s<1$ and
\[
\sigma=s\cdot\left(\frac{1}{4}-9\varepsilon\right)-12\varepsilon-\frac{2a_{L}}{N_{*}}>0;
\]
the condition (S1) in the form $\varepsilon>0$, which follows from $\alpha<1$ and $1-\alpha<\varepsilon$;
the value $\mathfrak{e}_2=0$ of Section~I.8.5 (the closing paragraph of Proposition~14.2\,(a), after the line (C11)), on which the clause $0\le\mathfrak{e}_2<1$ of (R6) rests.
Constants: $\alpha$, $\varepsilon$ and $s$, lines (C1), (C7); the exponent $\mathfrak{e}_2$ is fixed in Section~I.8.5.
Used in Proposition~43.1\,(R6), where $\varepsilon>0$ gives $\mathfrak{e}<\frac{s}{4}<\frac{1}{4}$, and hence $\sigma<\frac{1}{4}$.
Through (R6), that is, through $\mathfrak{e}\le\frac{1}{4}$ and $\sigma\le\frac{1}{4}$, the conditions $0<s<1$ and $\varepsilon>0$ are used also in Proposition~43.2 (Chapter~43, Section~2.2, Steps~2 and~3);
the condition $\sigma>0$ is used also in Proposition~43.2, through (R6), and in Theorem~44.
The value $\mathfrak{e}_2=0$ is used also in Proposition~43.1\,(R4), where it gives $C^{G}_{\beta}=C^{G}_{\beta}\cdot e^{-\mathfrak{e}_2}$, so that $C_{\mathrm{R4}}\le C^{G}_{\beta}\cdot e^{-\mathfrak{e}_2}$;
in Proposition~43.2, through (R6), in the form $0\le\mathfrak{e}_2<1$: in Chapter~43, Section~2.2, Step~2, where $6-2\mathfrak{e}_2>4$ and $e^{-\mathfrak{e}_2}\le e^{-1}$ for the charge $0<e\le1$, and in (PU-1), which proves the assertion $\mathfrak{e}_2\in[0,1)$ of Proposition~43.2;
and in Theorem~44, in the form $\mathfrak{e}_2\ge0$, which gives $1\le e^{-\mathfrak{e}_2}$ for the charges $0<e\le1$.

\item
Condition:
$e_1\le\frac{1}{3}$;
$v_K\ge8$;
$\left(1+\frac{3\,\mathrm{lg}}{v_K}\right)^r<L$;
$3\,e_1^{\frac{1}{4}-9\varepsilon}\le1$
(the rows (e1) and (e3) of Proposition~14.2\,(b)).
Constant: $e_1$, line (C9).
Used in Proposition~43.1\,(R1), where the third inequality gives that $N_0(e)$ is finite, and (R2); and in Proposition~42.2.
The condition $e_1\le\frac{1}{3}$ is used also in Theorem~44, where it gives $e\le1$.

\item
Condition: (W3), that is,
\[
6\,N_{*}\cdot e^2\cdot R^{*}(e)\cdot\log\frac{1}{e}\le\pi^2
\]
for all $e\le e_G$; it follows from the row (e8) of Proposition~14.2\,(b) together with (Z8.w), that is, $e\cdot h^{\frac{1}{2}}\cdot\log h\le1$.
Constants: $e_1$, line (C9), and $e_G$, line (C10).
Used in Proposition~42.2, where it gives $D^{\mathrm{fl}}\le1$ and $\varsigma^{\frac{1}{2N_{*}}}\le2\,e^6$.

\item
Condition: (Z10.s1), that is, $e_2\le e_G$.
Constant: $e_2$, line (C10).
Used in Proposition~43.2 and in Theorem~44.

\item
Condition: (Z10.s2), that is, $e_2\le e_{\sharp}(\beta_{\mathrm{ix}})$ for $\mathrm{ix}=0,+,-$.
Constant: $e_2$, line (C10).
Used in Proposition~43.2, where the three probes fit and (R3)--(R5) of Proposition~43.1 hold, and in Theorem~44, where the probes fit.

\item
Condition: (Z10.s4), in the form
\[
\frac{5}{2}\,C'_{\bar{\kappa}}\cdot e_2^{\sigma}\le\frac{b(f)}{4},\qquad C'_{\bar{\kappa}}=8\,C_{\bar{\kappa}}+3\,C_{\bar{\kappa}}^2.
\]
Constant: $e_2$, line (C10), after $C_{\bar{\kappa}}$.
Used in Proposition~43.2 (Chapter~43, Section~2.2, Step~8, (PU-1)) and in Theorem~44\,(B3), (B5), (D6).

\item
Condition: the inequality on $x_0$ of \cite[Section~2.8]{QEDI}, which is part of (Pb):
\[
\frac{6}{\bar{m}^2}\cdot\|a\|_{\infty}^4\cdot\left(\frac{4\pi\cdot r_\beta^3}{3}\right)^2\cdot\left(\bar{m}^2+\left(\bar{m}+\frac{1}{d_0}\right)^2\right)\cdot\frac{\exp\left(-2\bar{m}\cdot d_0\right)}{4\pi^2\cdot d_0^2}\le\frac{b(f)}{12},
\]
with $a=d^{*}\beta$, $\|a\|_{\infty}=\sup_y|a(y)|_2$, the radius $r_\beta$ and $d_0=|x_0|_2-2r_\beta>0$ as in (Pb).
Constant: $x_0$, line (C10), after $\beta$.
Used in Theorem~44\,(B3), (B5), for $\mathrm{Ext}\ge\frac{e^4\cdot b(f)}{2}$.

\item
Condition: $\rho\in(0,1)$; the standard weights; $\bar{m}\in(0,1]$.
Constants: $\rho$ and the weights, line (C0); $\bar{m}$, line (C10).
Used in Propositions~43.1, 43.2 and in Theorem~44, as follows.
The inequality $\rho>0$ enters (0.M) of Chapter~43, where it gives $\kappa_{(0)}>0$ and that $C_{\mathrm{mat}}$ is a real number.
In the proof of Proposition~43.2, at (5.2) of Chapter~43, Section~2.2, the condition $\rho\in(0,1)$ is used in the form $-1\le\rho\le1$, that is, $\rho$ lies in the range for which \cite[Theorem~6.27]{QEDII} is stated.
The number $\rho\in(0,1)$ is a datum of \cite[Definition~2.36]{QEDII} in Proposition~43.2 (Chapter~43, Section~2.2, Step~8, and (PU-7)) and in Theorem~44 (Chapter~44, Section~0.3), where it is also a datum of \cite[Theorems~1.2, 1.3 and~1.4]{QEDII};
in the form $0<\rho<1$ it enters \cite[Sections~8--10]{QEDII}, through Theorem~44\,(C), (D).
The standard weights are a datum of \cite[Definition~2.36]{QEDII} at the same places;
at (5.2) of Chapter~43 they fulfil, by $w_i>0$, the hypothesis $w_0,\dots,w_3>0$ of \cite[Theorem~6.27]{QEDII};
and $w_3=6$ enters (0.M) of Chapter~43.
The range $\bar{m}\in(0,1]$ enters (0.H) of Chapter~43, that is, $h\ge1$ and $\frac{1}{L}<\bar{m}\cdot h\le1$, in Propositions~43.1 and~43.2;
in the form $h\ge1$ it enters Theorem~44 (Chapter~44, Section~0.2, and $R=8h\cdot L^{T}\ge1$ in the proof of Lemma~2.1\,(b) of Chapter~44);
and it is a datum of \cite[Definition~2.36]{QEDII}.
\end{itemize}

The Propositions~38.1--47.2 and~50 and Theorem~44 add no condition on $\alpha$, $\kappa$, $L$, $M$, $\varepsilon$, $s$, $e_0$ or $e_1$.
Of the four conditions (Z10.s1), (Z10.s2), (Z10.s3), (Z10.s4) on $e_2$ in \cite[Section~5.6.2]{QEDI}, the line (C10) imposes (Z10.s1), (Z10.s2) and (Z10.s4).
Apart from the entries above (for $L$ and $M$ in Propositions~38.3\,(e), 41, 47.2\,(a), 40\,(e), 42.2 and~50\,(h); for $r$, $e_1$ and $e_G$ in Proposition~42.2), the Propositions~38.1--38.3, 39, 40, 41, 42.1, 42.2, 46, 47.1, 47.2 and~50 contain no condition on a constant besides these:
$\rho\in(0,1)$ (in the forms $0\le\rho$ in $\operatorname{Re}\left(D_B+\boldsymbol{\mu}\right)\ge\boldsymbol{\mu}$; $|\rho|\le1$ in \cite[Sections~5, 6 and~13]{QEDII}, in the last through \cite[Lemmas~13.8 and~13.12]{QEDII}; $\rho^2\ne1$ in \cite[Theorem~4.13]{QEDII});
the standard weights (through $w_i>0$, the closure relation (I.1.1) and $w_0=w_3=6$);
$\bar{m}\in(0,1]$ (through $\bar{m}\cdot\eta>0$ and the fact that $h$ is a power of $L$ with an exponent $\ge0$);
and the sides $n_\mu\ge4$.

\section*{Part IV. The order of the proofs}

\subsection*{IV.1. Dependence}

Let $X$ and $Y$ be any two of the Propositions and Theorem~44;
Proposition~46 stands in Part~I, before all the others, which stand in Part~II.
We say that $X$ quotes $Y$ if the proof of $X$ in Part~V refers to $Y$, by its number, as a statement on which it rests.
The references of the proofs of Propositions~24.3 and~24.2 to Section~5.4 of Chapter~24 count as quotations as well:
that section proves the convergence of the series (G3.8.4), that is, the first sentences of Proposition~24.1\,(a), (b), and these references take exactly these sentences;
they give the pairs from Proposition~24.3 and from Proposition~24.2 to Proposition~24.1, which point backward.
A reference to a form of Annex~A is a reference to the statements from which the form is taken;
they are named with the form, at its head or after the words ``Taken from''.
In this way the proof of Proposition~43.1 quotes Proposition~36\,(a), (g), (h) and Proposition~37.2\,(a), (d)--(g), through (G6.T-36) and (G6.T-37) (Chapter~43, Step~5).
Where the proof of Proposition~14.2 refers to (G1.T-11) (Chapter~14b, Step~3, row~28), it concerns only the letters on which $C_E$, $\mathrm{Csy}_{\mathrm{ak}}$, $\mathrm{Cpv}_0$ depend;
of the statements named with that form it reads only Propositions~37.1 and~36, which give these constants, and only as definitions, as (B1) below describes, so that it quotes none of them.
The form (G5.T-1) of Section~G5.7 is referred to by its name in the proofs of Propositions~35.5--35.8 (Chapter~35b, (n3)) and of Proposition~35.10 (Chapter~35c, paragraph~3.0 and Section~4.2).
These proofs take of it only
\[
C^{F}\ge1,\quad 0<\gamma^{F}\le1,\quad \gamma^{F}\cdot M\ge16,\quad C_{0}\ge\max\left(1,C^{F}\right),\quad M\ge L^2,\quad K_{0}\ge1,\quad C_{d}<\infty,\quad e_1>0,
\]
that is, ranges of the constants and facts listed under E\,(i), (iii) below (Proposition~14.2\,(a), (b), (f); Proposition~14.3\,(f)), together with the constants $C^{F}$, $\gamma^{F}$ of Proposition~14.1\,(i), named in the paragraph on the constants of Proposition~14.1.
Such a reference is a use of ranges of the constants and of facts of E, and no quotation of the Propositions~14.m:
the pairs from Proposition~35.10 to Propositions~14.1, 14.2 and~14.3, which (B13) states, come from the references by number in its proof,
and the form gives no pair from Propositions~35.5--35.8.
Neither proof takes the entry of the form by which $C_{d}$ is one constant for all levels $j\ge0$, the only entry behind which Proposition~14.4\,(a) stands;
so neither Proposition~35.10 nor any of the Propositions~35.5--35.8 quotes Proposition~14.4.
The indications ``Used in'' of the forms of Annex~A belong to no proof and give no pair.
A circle of the relation is a sequence $X_0,X_1,\dots,X_n=X_0$ of statements, $n\ge 2$, in which $X_{i-1}$ quotes $X_i$ for $i=1,\dots,n$;
the two pairs from Proposition~43.1 named above point backward and lie on no circle.
References and uses of the following four kinds do not belong to the relation.
\begin{itemize}
\item The places of the kinds listed in D below;
among them is the reading of the Definitions in Proposition~27.1 by the proofs of Proposition~18.5\,(f) and Proposition~48.5\,(e).
\item The uses of the inequalities between the constants and of the facts listed in E below;
among them is what the closing part ``In the flow'' of Proposition~30.3\,(f) takes of the Propositions~14.m through (Sc6) and (Sc4) of Section~G4.0.2 (see the paragraph on the statements below).
Such a use is by itself no quotation of Proposition~14.2 or of Proposition~14.3, also where the proof names the clause from which it is taken.
For the statements that stand after Proposition~14.2 in Part~II, the relation contains the pairs to Propositions~14.2 and~14.3 that are stated in B, in E and in the paragraph on the statements below;
it contains no such pair from a statement that stands before Proposition~14.2.
Further pairs to Proposition~14.2 or to Proposition~14.3 that such uses by statements standing after Proposition~14.2 could give close no circle of the relation with the ten pairs of the paragraph on the statements, also when several of them are included:
Proposition~14.2 quotes, directly or through other statements, only statements that stand before it;
of the statements other than itself that stand after Proposition~14.2, Proposition~14.3 quotes, directly or through other statements, only Propositions~35.1, 35.2, 35.3, 35.10, 36, 37.1 and~37.2,
and of these only Proposition~35.10 uses an assertion of Proposition~14.3, through the pair from Proposition~35.10 to Proposition~14.3, which B contains.
So the statements that lie on a circle of that relation do not depend on which of these further pairs are included.
\item The uses of the sentences of Part~I treated at the end of this section.
Such a use is by itself no quotation of the statement behind the sentence;
a statement whose proof takes the sentence may quote that statement in addition, by its number.
\item The references between Propositions~26.1 and~26.3 that C below treats:
the reference of the Definition in Proposition~26.1, and of the step (1.9) of its verification in Chapter~26a, to the first sentence of Proposition~26.3\,(a);
and the use, in the proof of Proposition~26.3, of the data, the notation, the facts and the Definition of Proposition~26.1 through the steps (1.0), (1.2), (1.3), (1.4), (1.5), (1.7) and (1.9) of that verification.
\end{itemize}
Every use that a proof of Part~V makes of another statement is a quotation, a use of one of these four kinds, or a use of one of the following three kinds.
\begin{itemize}
\item A lemma of a later chapter.
The proof of Proposition~16.1 uses Lemma~A and (Un) of Chapter~17, Section~0.5;
their proofs rest on no statement, so that Proposition~16.1 quotes no statement.
The proof of the closing sentence of Proposition~2.8 (Chapter~2b, end of the proof of Proposition~2.8) uses the inequalities $64a'\le608a'\le C^{F}$, and it derives them there from the row $C^{F}$ of the line (C6) of Proposition~14.2\,(a), through the third entry of the maximum that defines $C^{F}$:
since $\mathrm{e}\ge1$, $\max(1,b)\ge1$ and $a'>0$,
\[
C^{F}\ge 2\cdot304\,\mathrm{e}\cdot\max(1,b)\cdot a'\ge608a'\ge64a'.
\]
The same inequalities stand in Lemma~1\,(K2) of Chapter~14a, a lemma of the later chapter that contains the proof of Proposition~14.1;
the proof of Proposition~2.8 names that lemma only as a second place at which they stand, and it takes nothing from Chapter~14a.
The assertion (K2) follows from the rows $C_1$, $C^{F}$ and $\alpha_0$ of Proposition~14.2\,(a) together with (M1) and (L5), and it rests on no statement.
What the closing sentence of Proposition~2.8 uses is therefore a part of (B1) below, and not of (B12).
Hence Proposition~2.8 does not quote Proposition~14.1, no forward pair arises, and the pair from Proposition~14.1 to Proposition~2.8, which (B12) contains, closes no circle.
Between clauses no use points forward either:
the closing sentence of Proposition~2.8, which belongs to (B11$'$), uses a part of (B1),
and the proof of Proposition~14.1\,(i) takes of Proposition~2.8 only the clauses (a)--(d), and not the closing sentence.
\item A sentence of a section of Part~II whose number carries one of the prefixes G1--G6, in a chapter of definitions or next to the Propositions that use it
(Sections~G1.1 and~G2.0.4; (FL) in Section~G4.0.1; (DAT) in Section~G5.4.0, under the heading~34):
it repeats a sentence of Part~I treated in (6) or (7) at the end of this section, or it names statements that stand before the statements whose proofs use it.
\item A step or a section of the proof of another statement of the same chapter.
Let the proof of $X$ use a step or a section of the proof of $Y$.
If the proof of $X$ quotes $Y$ by its number, then, whichever way the use points, it gives, read as a dependence, only the pair from $X$ to $Y$, which the quotation gives already.
Pointing backward, this is the case for the proof of Proposition~26.3, which uses the step (b1) of the proof of Proposition~26.2 (Chapter~26a, Section~3.7, (lk2)), a consequence of Proposition~26.2\,(a), and which quotes Proposition~26.2;
the proof of Proposition~26.2 uses neither Proposition~26.1 nor Proposition~26.3.
Pointing forward, it is the case in Chapter~24.
The proofs of Propositions~24.1 (Chapter~24, Sections~5 and~9) and~24.2 (Chapter~24, Section~8) use the displays (24.2), (24.3), (24.5)--(24.8) and Sections~2.2, 3.1--3.12 and~6 of Chapter~24, which belong to the proof of Proposition~24.3,
and Propositions~24.1 and~24.2 quote Proposition~24.3; these are two of the ten pairs of the paragraph on the statements below.
Each display and each section so used is established at one of the entries (A-2), (A-4), (A-6) of A, before the step at which it is used, so that the order A resolves these uses.
If the use points backward and the proof of $X$ does not quote $Y$, then the use, read as a dependence, gives the pair from $X$ to $Y$, which points backward.
This pair closes a circle only if $Y$ reaches $X$, directly or through other statements.
Every pair of the relation other than the ten pairs points backward, and so does every further pair of this kind;
hence a statement reaches a statement that stands after it only through one of the ten pairs.
Two instances are these.
The proof of Proposition~22.4\,(b) uses the display (P.3) of the proof of Proposition~22.2\,(c);
paragraph~4.6 of the proof of Proposition~35.8 (Chapter~35b) uses paragraph~1.10 of the proof of Proposition~35.5, of which it takes only that $C_\beta$, $\mathbf{D}$ and $D_\beta$ do not change when $\beta$ is replaced by $U_{r'}\beta$ with $r'\in\{\pm1\}^3$;
neither proof refers to the earlier statement by its number.
Proposition~22.2 reaches none of the Propositions~23.2, 24.1, 24.2, 14.3 and~14.4, at which the ten pairs start, and hence no statement that stands after it.
The ten pairs end at Propositions~24.2, 24.3, 35.3, 35.10, 36, 37.1 and~37.2, none of which reaches Proposition~35.8;
so Proposition~35.5 does not reach Proposition~35.8.
Neither of the two pairs, therefore, closes a circle.
Moreover, (P.3) rests on (I.s1) alone, and the fact taken from paragraph~1.10 rests on Sections~I.7.5 and~G5.1.3 alone, so that what is used rests on no statement.
The following uses of this kind have reasons of their own.
The proof of Proposition~26.3 uses the steps (1.0), (1.2), (1.3), (1.4), (1.5), (1.7) and (1.9) of the verification of Proposition~26.1 in Chapter~26a (see C).
The proof of Proposition~35.9 uses Sections~3.0 and~3.1 of Chapter~35c, which belong to the proof of Proposition~35.10 (see D).
The proof of Proposition~18.3\,(e) (Chapter~18a, Section~3.4) uses, pointing forward, the paragraph ``The objects'' at the beginning of the proof of Proposition~18.4 (Chapter~18a, Section~4);
that paragraph rests on Proposition~16.1\,(d), on (sf1), (sf4) and on the definitions of Sections~I.1.1, I.1.3, I.1.7, I.3.6 and of (G2.r5), and on no clause of Propositions~18.3 and~18.4.
The proofs of Propositions~24.3 and~24.2 use Section~5.4 of Chapter~24;
these uses are the quotations of the first sentences of Proposition~24.1\,(a), (b) described above, and A below resolves them, since (A-5) comes before (A-6).
\end{itemize}

\paragraph{The statements.}
Part~II prints the statements in an order in which every statement quotes earlier statements only, with the following ten exceptions:
\begin{itemize}
\item[] Proposition~23.2 quotes Proposition~24.3;
\item[] Proposition~24.1 quotes Propositions~24.2 and~24.3;
\item[] Proposition~24.2 quotes Proposition~24.3;
\item[] Proposition~14.3 quotes Propositions~35.3, 35.10, 36 and~37.2;
\item[] Proposition~14.4 quotes Propositions~36 and~37.1.
\end{itemize}
Besides these, the relation contains the pairs from Propositions~24.2 and~24.3 to Proposition~24.1, given by the references to Section~5.4 of Chapter~24 (see above),
and the pairs from Propositions~24.1, 24.2 and~24.3 to Proposition~23.2;
all of them point backward.
When these ten pairs are included, the statements that lie on a circle are those of two sets,
the set of Propositions~23.2, 24.1, 24.2, 24.3, joined by the circles
\[
24.1\to24.2\to24.1,\qquad 24.1\to24.3\to24.1,\qquad 23.2\to24.3\to23.2,
\]
and the set of Propositions~14.3 and~35.10,
and two statements lie on a common circle if and only if they belong to the same one of these sets.
If the references between Propositions~26.1 and~26.3 treated in C belonged to the relation, they would add an eleventh forward pair, from Proposition~26.1 to Proposition~26.3, and a third such set, formed by Propositions~26.1 and~26.3;
the order (C-1)--(C-4) of C below resolves it.

The closing part ``In the flow'' of Proposition~30.3\,(f) takes from (Sc6) the inequality $v_K\ge 8$ and the condition (2.1), which belong to Proposition~14.2\,(f) and~(c),
and (e1) at $v=v_K$, which is given by the second assertion of Proposition~14.3\,(f);
through (Sc4), that is, through the proof of \cite[Proposition~19.1(f)]{QEDI}, it takes the inequality $L\ge 33$, which Proposition~14.2\,(f) derives from (L1), and the fact that $M$ is a power of $L$, which is the row (M1) of Proposition~14.2\,(b).
By E below this use is no quotation of the Propositions~14.m.
Read as a dependence of Proposition~30.3 on Propositions~14.2 and~14.3, it gives two further forward pairs, from Proposition~30.3 to Proposition~14.2 and from Proposition~30.3 to Proposition~14.3.
Then the Propositions
\[
14.3\to35.3\to35.1\to14.2\to30.3\to14.3
\]
form a circle, and in the place of the set formed by Propositions~14.3 and~35.10 there is one set of statements any two of which lie on a common circle, formed by the nine Propositions~30.3, 14.2, 14.3, 35.1, 35.2, 35.3, 35.10, 36, 37.2;
the order (B0)--(B14) of B below resolves it.
Propositions~35.1, 35.2, 35.3, 36 and~37.2 belong to this set through the pair from Proposition~35.1 to Proposition~14.2.
This pair belongs to the relation (see E),
although all that the proof of Proposition~35.1 takes from Proposition~14.2 are rows of (a) and (b), read as definitions and as conditions on the constants, among them the definition of $e_1$ in the line (C9), which Step~1 of that proof names (Chapter~35a), and the inequality $L\ge2400^8$ of (f) (see (B9)), that is, definitions and inequalities of E\,(i);
the same definition of $e_1$ is taken, and named, without a pair by the proof of Proposition~49 (Chapter~49, Section~2.1), and it is taken without a pair by the closing part ``In the flow'' of Proposition~30.3\,(f) (Chapter~30, Section~4.5).
If the pair from Proposition~35.1 to Proposition~14.2 is left out, the set formed by Propositions~30.3, 14.2, 14.3 and~35.10 takes the place of the set of nine,
and the order (B0)--(B14) resolves it as well.
Read as a dependence on Proposition~14.2 alone (Chapter~30, Section~4.5, derives (e1) at $v=v_K$ itself: see (B5$'$)), it leaves the two sets as they are and adds a third one, formed by Propositions~30.3 and~14.2;
the order (B0)--(B5$'$) resolves it.

The proofs of Propositions~30.3, 20.3, 48.6, 26.6, 27.1--27.4, 28.2, 31.2, 32.1--32.3, 49 and~43.1 use (Sc6), or \cite[Proposition~19.2]{QEDI} in whose place (Sc6) stands.
The proof of Proposition~28.3 takes the same assertion through the first display of (Sc0),
and the proof of Proposition~28.1 through (Sc1), which Section~G4.0.2 states under (S-4) alone,
that is, with the hypotheses $v_K\ge 1$ and $(1+\mathrm{lg}/v_K)^r\le L$ of \cite[Proposition~19.1]{QEDI} derived from (S-4).
Read as dependences on Propositions~14.2 and~14.3, these uses add, besides the two forward pairs from Proposition~30.3, only backward pairs, which start at Propositions~27.1--27.4, 28.1--28.3, 31.2 and~49;
the other statements named here quote Propositions~14.2 and~14.3 already (see E).
The statements that lie on a circle are then those of the set of Propositions~23.2, 24.1, 24.2, 24.3 and of the set of nine statements described above,
or, without the pair from Proposition~35.1 to Proposition~14.2, of the set formed by Propositions~30.3, 14.2, 14.3 and~35.10
(for the case in which the uses (6) of the paragraph on the sentences of Part~I are read as dependences as well, see the end of this section).

Among the ten pairs, those from Proposition~14.3 to Propositions~35.3, 36, 37.2 and those from Proposition~14.4 to Propositions~36, 37.1 point forward and close no circle:
the statements that Propositions~35.3, 36, 37.1 and~37.2 quote, directly or through other statements, include neither Proposition~14.3 nor Proposition~14.4 nor Proposition~35.10.
With the two pairs from Proposition~30.3, the statements that Propositions~35.3, 36 and~37.2 quote, directly or through other statements, include Propositions~14.3 and~35.10, through the pair from Proposition~35.1 to Proposition~14.2;
without that pair they include neither.
Also with these pairs no circle between clauses arises.
Indeed, Proposition~14.2 takes from Proposition~30.3 only the clause (e) and the definitions of its constants, that is, a part of (B0), which does not use the closing part ``In the flow'' of (f);
the entry (B5$'$) comes after (B4);
and no clause of the Propositions~14.m, 30.m, 35.m, 36 and~37.m uses the closing part ``In the flow'' of Proposition~30.3\,(f);
of the later statements, at most Proposition~32.1 uses it (see (B5$'$)).

\paragraph{The clauses.}
The ten exceptions are resolved by the orders of clauses A and B below, and the references between Propositions~26.1 and~26.3 by the order C.
In each of these orders every entry uses only entries that precede it, with the two refinements stated in (B0) and (B11$'$):
the entries of Lemma~1 of Chapter~14b that (B0) names, which are a part of (B1), are established before (B0),
and the part of (B15) that (B11$'$) names is established before (B11$'$).
The dependences stated in A--C are those of the proofs in the chapters named with them.
Every forward pair among the ten falls under A or under B;
C concerns a reference to a later statement that stands in the Definition of Proposition~26.1 and in its verification.

\paragraph{A. The Propositions 23.m and 24.m.}
This order is stated also at the end of Definition~24.D, and the proofs in Chapters~23 and~24 follow it (Chapter~24, Section~0.2).
\begin{itemize}
\item[(A-1)] Proposition~23.1\,(a)--(g) and Proposition~23.2\,(a)--(d);
they use no clause of the Propositions~24.m.
\item[(A-2)] Proposition~24.3\,(adm), (q1);
they use no clause of the Propositions~23.m.
\item[(A-3)] Proposition~23.2\,(e),
from Proposition~23.2\,(c), (d), Proposition~24.3\,(q1), the property W2 of the family (p2) and the second part of the condition (A3) of Proposition~14.2\,(b).
\item[(A-4)] Proposition~24.3\,(q5), (al), (be), (ga), (de), (ze), (et);
among these, (ze), the fourth sentence of (al) and the last sentence of (ga4) are established at this stage only for the finite sums and integrals and for the terms of the series (G3.8.4) of Definition~24.D.
\item[(A-5)] Proposition~24.1\,(a), first sentence with its display, and Proposition~24.1\,(b), first sentence (likewise in the reading (o)):
the series (G3.8.4) converge.
\item[(A-6)] Proposition~24.3\,(ze), the fourth sentence of (al) and the last sentence of (ga4), for $E^{\#,x}_k(m,\mathscr{E})(\mathbf{Y})$;
the closing sentences of Proposition~24.3, on the two other parts $(R\mathscr{E})^{\mathrm{nat}}(X)$ and ${\delta E^{\mathrm{det}}}_k(X)$ of $\mathscr{E}^{+}(Y)$, into which no clause of the Propositions~23.m and~24.m enters
(their proof, in Chapter~24, Section~6, rests only on the display of Definition~24.D that gives $\mathscr{E}^{+}(Y)$, on Sections~I.5.6 and~I.5.7 and on W12);
Proposition~24.2.
\item[(A-7)] Proposition~24.1\,(f), (g);
the remaining sentences of (a);
(b) from its second sentence on;
the first clause of (d);
(c);
the rest of (d);
(e);
(i).
\end{itemize}
The reading (o) of each clause of Proposition~24.3 is established at the step of that clause, (A-2), (A-4) or (A-6) (Chapter~24, Section~6, paragraph on the reading (o)).

\paragraph{B. The Propositions 14.m together with the Propositions 30.m, 35.m, 36 and 37.m.}
The dependences of the entries are established in Chapter~30, Sections~5.2 and~6;
in Chapter~14b, in the order of the proofs of Propositions~14.2 and~14.3 at the beginning of Section~1, (o1)--(o6), and in Section~2\,(g), on the numbers that are given;
in Chapter~35c, Sections~3.0, 3.1 and~3.4;
and in Chapter~14c, Section~1.3, on the letters $C_E$, $\mathrm{Csy}_{\mathrm{ak}}$, $\mathrm{Cpv}_0$.
\begin{itemize}
\item[(B00)] Proposition~46; Propositions~1.1--1.7.
Proposition~46 quotes no statement;
each of the Propositions~1.1--1.7 quotes only earlier ones among them.
Of the Propositions~14.m they read rows of Proposition~14.2\,(a), as definitions and with the ranges stated there, and the condition (L5).
\item[(B0)] Propositions~30.1 and~30.2; Proposition~30.3\,(a)--(e), (g), and (f) without its closing part ``In the flow''.
The signs and ranges of constants that their proofs use are entries of Lemma~1 of Chapter~14b for the following rows of Proposition~14.2\,(a), the rows being numbered as in Chapter~14b, Section~0:
the rows 1, 2, 10, 12, 15, 19--21;
the row 33, that is, $L$, whose entry ($L$ is an odd integer, $L\ge5$) the proof of Proposition~30.2\,(b) uses in the form $L\ge1$, so that $0<\eta=L^{-N}\le1$, and which Proposition~30.3 uses as well;
and, for Proposition~30.3 only, also the rows 29, 35, 36.
That lemma is a part of (B1); it is proved, and used, row by row.
The entries just named, and those of the earlier rows on which they rest, quote no statement other than Proposition~1.2\,(S3), (S6), (S7);
they are established before Proposition~30.1,
and they do not rest on the entry of the row 25, which quotes Proposition~30.1\,(c).
The statements of this entry stand before the Propositions~14.m and contain no letter of the lines (C7)--(C9) (Chapter~30, Sections~5.2 and~6).
\item[(B1)] Proposition~14.2: (a) as a list of definitions, (b), the paragraph that follows the conditions of (b) without its second sentence (Chapter~14b, Step~9), the first two sentences of (e), and (f), with three parts left out:
the assertions of (f) that begin with $K_0\ge 1$, namely the last three rows of the display in the paragraph ``Ranges'' of (f) and what follows that display down to $\mu_-<\mu_*$, the sentence ``If $e\cdot h^{\frac{1}{2}}\le e_1$'' included;
the last inequality of Annex~B, Section~G5.9.3;
and the inequalities that are named in the indications ``Used in'' of (b) and that (f) does not state.
This is (o1) of Chapter~14b, together with (f1)--(f8), (f10), (f11) there.
It quotes Propositions~1.2, 1.3\,(b), 1.5, 1.6, 30.1, 30.3 and~46, and no later statement.
Propositions~1.3\,(b) and~46 enter only through the entry of Lemma~1 of Chapter~14b for the row 26, and both belong to (B00).
The entries of (a) whose reference is a statement that stands later are read as definitions.
\item[(B2)] Proposition~14.3\,(a)--(e).
This is (o2) of Chapter~14b; it rests on (B1).
\item[(B3)] Proposition~14.2\,(c), first sentence (Chapter~14b, paragraph~5.1, display (5.1)):
if the set in the definition of $e_1$ is not empty, then (e1)--(e8) hold for all $v\ge-\log e_1$.
This is (o3) of Chapter~14b.
The second assertion of Proposition~14.3\,(f) states the same conclusion again, without the condition, and its proof quotes Proposition~14.2\,(c) for it.
\item[(B4)] Proposition~14.3\,(f): $e_1>0$, and (e1)--(e8) hold for all $v\ge-\log e_1$.
This is (o4) of Chapter~14b; it rests on (B1)--(B3).
\item[(B5)] The rest of Proposition~14.2\,(c), from its second sentence on (Chapter~14b, paragraphs~5.2 and~5.3);
Proposition~14.2\,(d) without the assertion, made in its sentence on (g2) and in its sentence on $\varpi_{\mathrm{sep}}^{+}$, that $\frac{M\cdot r_{*}}{8L}-1$ is a lower bound of $d_{\mathrm{sep}}$;
the third and the fourth sentence of Proposition~14.2\,(e);
the three parts left out in (B1).
This is (o5) of Chapter~14b, together with (f9), (f12) there.
Each of these sentences that contains $e_1$ either carries the datum $e>0$ with $e\cdot h^{\frac{1}{2}}\le e_1$ or is void without such an $e$;
under this datum the set in the definition of $e_1$ is not empty, and (B3) applies.
Nothing of (B4) is used: the proof of Proposition~14.2 uses no clause of Proposition~14.3 (Chapter~14b, Step~0\,(i), and (o5) there).
Where Steps~5 and~6 of Chapter~14b name Proposition~14.3\,(f), in a bracket in paragraph~5.1 and in a conditional clause in paragraph~6.1, no step rests on it:
paragraph~5.1 treats both cases of the set in the definition of $e_1$, and the condition $e_1>0$ of paragraph~6.1 holds under the datum.
Proposition~14.2\,(d) names the setting (ST*)$^{\mathrm{ch}}$ of Section~G5.1.1 only in order to fix the meaning of the index $*$ and the numbers $\mathrm{pp}$ and $d_{\mathrm{sep}}$ of Section~G5.1.4.
Of this setting, (d) and its proof (Chapter~14b, Step~6, in particular paragraph~6.3, and paragraphs~7.1, 7.3 and~7.4) involve only $\bar{m}$, the charge $e>0$ with $e\cdot h^{\frac{1}{2}}\le e_1$, the side $R=R^{*}(e)$ and the integer $N\ge N_0(e)$, that is, (S-4) at $R^{*}(e)$;
they use it through the scales of Sections~I.6.1 and~I.6.2 and through (Sc1) and (Sc5) of Section~G4.0.2, that is, \cite[Proposition~19.1(c)]{QEDI} and \cite[Proposition~19.1(g)]{QEDI}.
The probe enters $\mathrm{pp}$ and $d_{\mathrm{sep}}$ as well, and the placement enters $d_{\mathrm{sep}}$.
Neither the pair $P_N(e)$ of Proposition~28.3 nor a primitive (Propositions~26.1, 26.5, 26.6\,(e)) is involved;
hence (B5) rests neither on Proposition~28.3 nor on Propositions~26.1, 26.5 and~26.6\,(e), and not on (B12), on which Propositions~26.5, 26.6 and~28.3 rest.
Two sentences of Proposition~14.2\,(d) name the lower bound $\frac{M\cdot r_{*}}{8L}-1$ of $d_{\mathrm{sep}}$ of Proposition~35.3\,(e): the sentence on (g2) and the sentence on $\varpi_{\mathrm{sep}}^{+}$.
In both, the assertion that $d_{\mathrm{sep}}\ge\frac{M\cdot r_{*}}{8L}-1$ for every admitted placement is an assertion of Proposition~35.3\,(e) and belongs to (B9).
Of the sentence on $\varpi_{\mathrm{sep}}^{+}$, the entry (B5) contains three things: the formula of $\varpi_{\mathrm{sep}}^{+}$, the identity
\[
\left(\frac{M\cdot r_{*}}{8L}-1\right)-4M-6=\frac{M\cdot r_{*}}{8L}-4M-7,
\]
and the comparison with the bound of \cite[Theorem~20.29]{QEDI}.
No step of the proof of Proposition~14.2 depends on the lower bound:
paragraph~7.4 of Chapter~14b uses only that (g2) and (g5) are written with $\frac{M\cdot r_{*}}{8L}-1$ and with $\varpi_{\mathrm{sep}}^{+}$, without $d_{\mathrm{sep}}$,
and of Proposition~35.3\,(e) the proof uses only which letters it contains (Chapter~14b, the paragraph ``Other Propositions'' that follows the order (o1)--(o6) of the proofs).
\item[(B5$'$)] Proposition~30.3\,(f), its closing part ``In the flow''.
It is stated under the datum $e\cdot h^{\frac{1}{2}}\le e_1$;
the definition of $e_1$ contains the constants $\gamma_\Sigma$, $C_\Sigma$ of Proposition~30.3\,(e).
It takes from (Sc6) the inequality $v_K\ge 8$ and the condition (2.1), which belong to (B5),
and the inequality $(1+3\,\mathrm{lg}/v_K)^r<L$, which is (e1) at $v=v_K$, a case of the sentence of (B3) that Proposition~14.3\,(f), that is (B4), asserts once more;
through (Sc4), that is, through the proof of \cite[Proposition~19.1(f)]{QEDI}, it takes the inequality $L\ge 33$ (Proposition~14.2\,(f), from (L1)) and the fact that $M$ is a power of $L$ (the row (M1) of Proposition~14.2\,(b)), which belong to (B1).
No clause of the Propositions~30.m and no clause of the Propositions~14.m uses this part.
Of the later statements, at most Proposition~32.1 uses it, through its citation of Proposition~30.3\,(c), (f) at the data of the flow (Chapter~32, (32.6)), whose admissibility it verifies again;
Proposition~32.2 cites only the clauses (c), (d), (e) of Proposition~30.3.
Nothing of (B4) is needed beyond (B3) under the datum:
Chapter~30, Section~4.5 (paragraph ``In the flow''), takes $v_K\ge 8$ from Proposition~14.2\,(f) and (2.1) from Proposition~14.2\,(c),
and it derives (e1) at $v=v_K$ from the definition of $e_1$ and the datum $e>0$ with $e\cdot h^{\frac{1}{2}}\le e_1$, by the argument of (B3).
The three facts $v_K\ge 8$, (2.1) and (e1) at $v=v_K$ follow also, as in Step~1 of the proof of Proposition~35.1, from the definition of $e_1$, the datum $e>0$ with $e\cdot h^{\frac{1}{2}}\le e_1$ and the row (e1).
\item[(B6)] Proposition~37.1.
It quotes no statement.
\item[(B7)] Proposition~35.1.
\item[(B8)] Proposition~35.2.
\item[(B9)] Proposition~35.3.
The assertion of the sentences on (g2) and on $\varpi_{\mathrm{sep}}^{+}$ in Proposition~14.2\,(d) that $\frac{M\cdot r_{*}}{8L}-1$ is a lower bound of $d_{\mathrm{sep}}$ belongs to this entry, as an assertion of Proposition~35.3\,(e) (see (B5)).
Of the Propositions~14.m, the entries (B7)--(B9) take three things.
First, rows of Proposition~14.2\,(a), read as definitions and with the ranges stated there: $b=1$, $\alpha<1$, $r\ge 2$, and the definition of $e_1$ in the row~118 (the rows numbered as in Chapter~14b, Section~0).
Second, the conditions (L1), (L5), (M1), (S1) and the rows (e1), (e4) of Proposition~14.2\,(b).
Third, the inequality $L\ge 2400^8$ of Proposition~14.2\,(f), which they use only in the form $L\ge 13$, for the class $\mathrm{Cl}$.
All three are parts of (B1);
the inequality $e_1>0$ follows from the datum $e\cdot h^{\frac{1}{2}}\le e_1$.
Proposition~35.2, the clauses of Proposition~35.3 other than (b1) and (c1), and the middle sentence of Proposition~35.3\,(b1), which is its first sentence at $R=R^{*}(e)$ and $\beta'=\tau_v\beta$, assume the setting (ST*)$^{\mathrm{ch}}$ of Section~G5.1.1, which is (FL) at $R^{*}(e)$ with the pair $P_N(e)$ and the primitives.
Of this setting, these clauses and their proofs (Chapter~35a, Sections~2 and~3) involve only $\bar{m}$, $e>0$ with $e\cdot h^{\frac{1}{2}}\le e_1$, $R=R^{*}(e)$ and $N\ge N_0(e)$, that is, (S-4) at $R^{*}(e)$;
for the middle sentence of (b1), the parts of the setting involved are the scales $\eta=L^{-N}\le1$ and $n_*=R^{*}(e)/\eta=2L\cdot\mathrm{hc}\ge2$, which its proof takes from Proposition~35.1\,(a), (b) (Chapter~35a, Section~3, (ps6) and (Eq3.5)).
These clauses and their proofs involve neither the pair $P_N(e)$ of Proposition~28.3 nor a primitive (Propositions~26.1, 26.5, 26.6\,(e));
hence they rest neither on Proposition~28.3 nor on Propositions~26.1, 26.5 and~26.6\,(e), and in particular not on Proposition~14.1, that is, on (B12), on which Propositions~26.5, 26.6 and~28.3 rest.
The later entries of B take from (B7)--(B9) only Proposition~35.1, Proposition~35.3\,(c1), and the first and the last sentence of Proposition~35.3\,(b1);
it is these that are stated without the setting (ST*)$^{\mathrm{ch}}$.
\item[(B10)] Proposition~36.
\item[(B11)] Proposition~37.2.
The entries (B10), (B11) rest on Proposition~46, Proposition~37.1 and the first and the last sentence of Proposition~35.3\,(b1).
\item[(B11$'$)] Propositions~2.1--2.8, 3--8, 15, 16.1, 9--13, in this order, which is the order of Part~II.
Each of them quotes only statements among the Propositions~1.1--1.7 and earlier statements of this entry
(Proposition~13 quotes Proposition~16.1\,(a1), (d1)--(d4); Propositions~15 and~16.1 quote no statement).
None of them quotes a statement among the Propositions~14.m, 30.m, 35.m, 36 and~37.m;
Propositions~5 and~8 refer to the ordered list of Proposition~14.2 only for the definitions of $C_1$ and $C^F$ (see D).
Of the Propositions~14.m they take rows of Proposition~14.2\,(a), (b), as definitions and as conditions on the choice of the constants, and lines of Proposition~14.2\,(f):
parts of (B1) and (B5).
Of this kind are the inequalities $64a'\le608a'\le C^{F}$ that the proof of the closing sentence of Proposition~2.8 uses:
that proof derives them in Chapter~2b from the row $C^{F}$ of Proposition~14.2\,(a), so that they are a part of (B1), and not of (B12)
(see the kind ``A lemma of a later chapter'' in the first paragraph of this section),
and the order needs no refinement besides the two stated in (B0) and in the present entry.
Through Section~I.8.1, Propositions~3 and~8 take in addition Proposition~14.4\,(a) for the constants in their bounds;
this is a part of (B15), which rests on (B00)--(B5) only and can be established before the present entry
(see (6) in the paragraph on the sentences of Part~I below).
No clause of Proposition~14.2 or of Proposition~14.3 quotes them.
\item[(B12)] Proposition~14.1.
It stands before Proposition~14.2 and quotes Propositions~1.6, 2.8, 3\,(a) and~4--13;
of Proposition~14.2 it reads the lists (a) and (b) as definitions.
\item[(B13)] Proposition~35.10: $e_\sharp(\beta')>0$.
It rests on Proposition~35.1\,(a), (b), (c), (f), Proposition~35.3\,(c1), Proposition~14.3\,(f), Proposition~14.2\,(a), (f) and Proposition~14.1,
that is, on (B1), (B4), (B5), (B7), (B9), (B12).
The assertions of Proposition~35.10 for $0<e\le e_\sharp(\beta)$, that every centred placement is admitted, that $\beta$ fits $(N,T^{*}(e))$ and that $R^{*}(e)\ge D_{\mathrm{be}}+1$, are proved in the setting (ST*)$^{\mathrm{ch}}$, with Lemma~1 and paragraph~3.6 of Chapter~35c.
Of this setting, Lemma~1, paragraph~3.6 and the notions of these assertions (centred and admitted placement, Section~G5.1.3; a probe that fits, Section~I.7.6) involve only $\bar{m}$, $e>0$ with $e\cdot h^{\frac{1}{2}}\le e_1$, $R^{*}(e)$ and $N\ge N_0(e)$, that is, (S-4) at $R^{*}(e)$,
through $\eta=L^{-N}\le1$ and $\ell_c=R^{*}(e)/(2L)$, together with the numbers $D_\beta$ and $y_\beta$ of the probe;
they involve neither the pair $P_N(e)$ nor a primitive.
Hence Proposition~35.10 does not rest on Propositions~26.5, 26.6\,(e) and~28.3.
The assertion $e_\sharp(\beta')>0$ is proved in Chapter~35c, Sections~3.0--3.4, without the setting.
\item[(B14)] Proposition~14.3\,(g): $e_G>0$, $e_2>0$, and admissible pairs $(\beta,x_0)$ exist.
This is (o6) of Chapter~14b.
It rests on (B4), (B5), Proposition~35.10 and, for $C'_{\bar{\kappa}}<\infty$, on Proposition~35.3\,(c1), Proposition~36\,(g) and Proposition~37.2\,(a), (d), (e).
\item[(B15)] Proposition~14.4\,(b), (c), (d), and Proposition~14.4\,(a) for all letters other than $C_E$, $\mathrm{Csy}_{\mathrm{ak}}$, $\mathrm{Cpv}_0$.
It rests on Proposition~14.2, Proposition~14.3\,(c), (f), and Propositions~1.2, 1.5, 30.1, 30.3.
This is what Propositions~32.1--32.3, the only statements that quote Proposition~14.4, take from it;
they stand before Propositions~36 and~37.1.
What is taken from Proposition~14.4\,(a) through Section~I.8.1 belongs to this entry as well
(see (6) in the paragraph on the sentences of Part~I below).
\item[(B16)] Proposition~14.4\,(a) for the three letters $C_E$, $\mathrm{Csy}_{\mathrm{ak}}$, $\mathrm{Cpv}_0$.
It comes after Proposition~37.1\,(f), (g) and Proposition~36\,(c).
The three letters occur in no row up to $e_1$ other than their own rows, and in no condition of Proposition~14.2\,(b).
\end{itemize}

\paragraph{C. Propositions 26.1 and 26.3.}
The Definition in Proposition~26.1 rests on the first sentence of Proposition~26.3\,(a), to which it refers, and so does the step (1.9) of its verification in Chapter~26a;
the Definition states also that this sentence concerns $\widetilde{\mathrm{Op}}$ only.
The proof of Proposition~26.3 uses the data, the notation, the facts and the Definition of Proposition~26.1 through the steps (1.0), (1.2), (1.3), (1.4), (1.5), (1.7) and (1.9) of that verification;
the step (1.1) enters through (1.9) and, without its number, in the proof of Proposition~26.3\,(e) (Chapter~26a, Section~3.3, (iii)).
The steps (1.0)--(1.7) verify the data, the notation and the facts; they use neither the Definition nor any clause of Proposition~26.3.
The step (1.8) uses the definition of $H^{\mathrm{pj}}_k$, and the step (1.9) is the verification of the Definition.
Of the steps (1.0)--(1.7), the first sentence of Proposition~26.3\,(a) takes only (1.4) and (1.5) (Chapter~26a, beginning of Section~3).
The clauses are proved in the following order (Chapter~26a, step (1.9) of the verification of Proposition~26.1, and Section~3.2).
\begin{itemize}
\item[(C-1)] Proposition~26.1, the notation and the facts stated with it, with the steps (1.0)--(1.7) of its verification.
\item[(C-2)] Proposition~26.3\,(a), first sentence;
it concerns $\widetilde{\mathrm{Op}}$ only and does not use the projection.
\item[(C-3)] The Definition in Proposition~26.1, with the steps (1.8) and (1.9) of its verification.
\item[(C-4)] The rest of Proposition~26.3.
\end{itemize}

\paragraph{D. Later statements that are named without being quoted.}
Apart from the cases of A--C, a statement names a later statement only at places of the following kinds, and none of these places is a quotation.
\begin{itemize}
\item The statements that give the constants of Proposition~14.2\,(a), or to which these constants belong, at the places of (a) where they are named:
the references of the entries;
the statements named in the definitions of the entries whose constants they give
($C_E$, $\mathrm{Csy}_{\mathrm{ak}}$, $\mathrm{Cpv}_0$: Propositions~37.1\,(f), (g) and~36\,(c);
$\mathrm{Cah}_i$, $\mathrm{Cpv}_1$, $\mathrm{Cpv}_2$, $\mathrm{epv}$, $\mathrm{Cpv}_{\beta}$: Proposition~36\,(e)--(g);
$C_{\mathrm{ct}}$, $C_1^{\mathrm{N2}}$, $C_2^{\mathrm{N2}}$: Proposition~37.2\,(c)--(e));
and the sentence of the preamble of (a) on a letter of a later entry in the property that defines a constant, which names Proposition~36\,(c) for $\mathrm{Cpv}_0$.
Likewise the indications ``Used in'' of Proposition~14.2\,(b), and the closing paragraph of Proposition~14.4, on the terms used in its clauses (a)--(d).
\item The other places of Proposition~14.2 at which a later statement is named:
after the conditions of (b), the second sentence of the paragraph that follows them, by which each condition on the constants that the proof of a Proposition adds is of the type ``$e_0$ (respectively $e_1$, $e_{\sharp}$, $e_2$) small depending on earlier constants'' and stands in the list, with its reference to (G1.T-12), a form headed ``All Propositions'';
in the preamble of (a), the sentence by which the sets in which $L$, $M$, $(r,p_0,p,\varepsilon,s)$ and $e_0$ are chosen are not empty and $e_1$, $e_G$, $e_2$ are positive, by Proposition~14.3,
and the sentence by which the constants of the entries $C_E$, $\mathrm{Csy}_{\mathrm{ak}}$, $\mathrm{Cpv}_0$ and of the line (C10) from $D_{\mathrm{be}}$ on exist, by the statements named as their references and, for $e_2$, by Proposition~14.3\,(g);
in (a), the entry $e_\sharp$, which is positive by Proposition~35.10,
the definition in the entry $\beta,x_0$, which names Proposition~43.2, besides the reference of the entry, as the statement in whose data the probes $\beta_{\pm}=\beta\pm\tau_{x_0}\beta$ and the tests $g_{\pm}=f\pm\tau_{x_0}f$, $f=dd^{*}\beta$, are formed,
the entry $s_B$, which is formed with the data of Proposition~43.2 and with the norm of Proposition~47.1 that the entry states,
and the sentence at the end of (a) on the absolute number $c_{\mathrm{cl}}$ of the data of Theorem~44;
in (d), the sentence that names the statements that need the conditions of the line (C10), namely Propositions~34.1, 34.3, 34.4, Propositions~35.5--35.10, Propositions~42.2, 43.1, 43.2 and Theorem~44,
and the three sentences that follow it, which name the places at which these conditions enter those statements: Propositions~34.1, 34.3, 34.4 and~35.5--35.10, Proposition~42.2 with the Remark that follows it, Proposition~43.1\,(R1), Proposition~43.2, Theorem~44 and Section~G6.Z4;
like the indications ``Used in'' of (b), these three sentences are indications of place and no quotations, Proposition~14.2 rests on none of the statements that they name, and they are no part of what Proposition~14.2 asserts;
the comparison with Proposition~35.1\,(c) in the sentence by which $r_{*}$, $v_{*}$, $p_{*}$ and $p_{0*}$ do not depend on $N\ge N_0(e)$, which (d) derives from \cite[Proposition~19.1(c)]{QEDI};
the sentence by which $\varpi_{\mathrm{sep}}^{+}$ arises from $\varpi_{\mathrm{sep}}$ when $d_{\mathrm{sep}}$ is replaced by its lower bound $\frac{M\cdot r_{*}}{8L}-1$ of Proposition~35.3\,(e);
the sentence on (g2), by which $d_{\mathrm{sep}}\ge\frac{M\cdot r_{*}}{8L}-1$ by Proposition~35.3\,(e);
and the sentences on (Pb) and on the radius $r_\beta$, by which $(\beta,x_0)$ and $r_\beta$ are data of Proposition~43.2;
in (e), the radius $r_\beta$, a datum of Proposition~43.2;
in (f), the bracket by which $C_E$, $\mathrm{Csy}_{\mathrm{ak}}$, $\mathrm{Cpv}_0$ are given by Propositions~37.1 and~36.
At none of these places does Proposition~14.2 rest on the later statement.
The second sentence of the paragraph after (b) is a rule for the proofs of all Propositions, and no part of what Proposition~14.2 asserts or of what Chapter~14b proves:
it is shown proof by proof in the sections ``Conditions on the constants'' (see E);
no step of the proofs of Propositions~14.2 and~14.3 rests on it (Chapter~14b, end of Step~9);
and the statements that quote Proposition~14.2 take rows, displays and lines of it, not this sentence.
Hence the references of Proposition~14.2 and of its proof to (G1.T-12) are no quotations and close no circle, as for (G1.T-11) above.
Of the other statements, the proof of Proposition~14.2 rests only on Propositions~1.2\,(S3), (S6), (S7), 1.3\,(b), 1.5, 1.6, 30.1\,(c), 30.3\,(e) and~46
(the setting (ST*)$^{\mathrm{ch}}$ named in (d) enters only as (S-4) at $R^{*}(e)$, see (B5));
of the later statements named in (a) and in the sentences on $\varpi_{\mathrm{sep}}^{+}$ and on (g2) it uses only which letters the property or the formula contains;
the assertions that the sentences above assign to later statements are no part of what Proposition~14.2 asserts;
and the independence of $r_{*}$, $v_{*}$, $p_{*}$, $p_{0*}$ of $N$, as well as the absence of a placement from (g2) and (g5), are proved without Propositions~35.1 and~35.3
(Chapter~14b, Step~0\,(i), (iv), paragraphs~7.1 and~7.4, and the paragraph ``Other Propositions'' that follows the order (o1)--(o6) of the proofs).
\item Definitions of constants.
Such a definition may be used before the place where it is stated, since the constants are those of Section~I.8 in the form given by Proposition~14.2 (head of Part~II):
Proposition~5\,(e), Proposition~6\,(f) and Proposition~8 name $C_1$ and $C^F$;
the constants of Proposition~14.1 include the row $M^F$ of Proposition~14.2\,(a).
\item The conditions on the constants that are stated with Propositions~3 and~13 for Proposition~14.2.
\item Definitions without slot:
Proposition~17.2\,(e) and Proposition~18.5\,(c) name the maps of Propositions~48.3 and~48.5, whose definitions stand in Section~I.5.11.
\item Definitions that are stated inside a later statement.
The last sentence of Proposition~18.5\,(f), the sentence on the two margins, concerns in its second half, after the semicolon, the admissible sets of the Propositions~27.m:
for the two sides $\ell=\ell_1$ and $\ell=\ell_2$ of the standing data (TT) of the Propositions~27.m, the two margins hold also for the sets that are admissible in the sense of the Definitions in Proposition~27.1.
With it, the second sentence of Proposition~48.5\,(e) concerns these sets as well.
The proofs (Chapter~18b, (7.7); Chapter~48b, Section~3.6\,(e-iii)\,(5)) read the standing data (TT) of the Propositions~27.m and the Definitions in Proposition~27.1,
\[
\mathrm{diam}(\tilde{X})\le \ell_1-14\,d_U\cdot M,\qquad d_U=2L+1,\qquad \ell_1\le\ell_2,
\]
as they are stated, and no assertion of the Propositions~27.m.
\item Remarks on wording, on notation and on margins:
Proposition~46\,(g+) (Proposition~36 is stated on the same range);
Proposition~20.2 (the list (S21) stated with the Propositions~21.m);
Proposition~47.1 (the notation $\tilde{a}_{\eta,\nu}=\mathrm{gs}_{\nu}[\beta]$ of Section~G5.1.3);
Proposition~18.5\,(f) (the margin stated in Proposition~25.2\,(o7)).
\item Proposition~45.4 names Theorem~44\,(B1) and the hypotheses of Proposition~45.5.
\item Remarks that name the statements in which a clause is applied
(Propositions~2.8 and~47.1;
the Remark ``The projected families'' that follows Proposition~14.1, for the clause (iii) of Proposition~14.1;
and Propositions~32.3 and~37.3\,(b), (c), which name the reference data (R2), (R4), (R5) of Proposition~43.1).
The closing sentence of Proposition~2.8 names Proposition~14.1\,(i) as the place at which the clauses (a)--(c) of Proposition~2.8 are applied;
it is proved in Chapter~2b with the inequalities $64a'\le608a'\le C^{F}$, which are derived there from the row $C^{F}$ of Proposition~14.2\,(a)
(see the kind ``A lemma of a later chapter'' in the first paragraph of this section).
\item The Remark ``The projected families'' that follows Proposition~14.1, and the last sentence of Chapter~14a, Section~2.5, at the end of the proof of Proposition~14.1\,(iii).
The Remark names Propositions~26.1, 26.3\,(b), 26.5 and~26.6\,(b), (e), and the sentence names Propositions~26.6\,(b) and~26.3\,(b), as the ground of a fact:
the primitives of the flow at the levels $k<k_{*}(R)$, projections included, are data (OS), hence (SF+) with $\ell\ge M\cdot r_k$, and their family (p5) has W1, W2, W7, W12--W14 for $\mathrm{par}$, so that the assertions of Proposition~14.1\,(iii) hold for it.
Neither place belongs to Proposition~14.1 or to a step of its proof:
the Remark stands after the statement,
and the sentence states that the Remark is no part of (iii) and that no step of the proof of (iii) rests on it.
The data of Proposition~14.1 are (SF) and, in (iii), a family (p5) with the properties required in Proposition~13;
neither $R$, nor a flow, nor a projection is among them.
The fact is established after Proposition~26.6, in Chapter~26b, Remark ``How (a)--(d) are applied'', from Propositions~26.6\,(b) and~26.3\,(b).
No step of Chapter~14a and no later proof rests on the Remark or on the sentence:
the statements in which, by the Remark, Proposition~13 is applied to the projected families apply it directly to the family (p5) of their data (for instance through (G3.T-13) and (G5.T-8) of Annex~A);
Proposition~26.5 takes from Proposition~14.1 its clause (ii) and constants;
and Proposition~35.10 takes from it only $C^{F}\ge1$ and $0<\gamma^{F}\le1$, through (G5.T-1).
Hence there is no pair from Proposition~14.1 to Proposition~26.3 or to Proposition~26.6, and (B12) and the sets of statements described above stand.
If the sentence of Chapter~14a, Section~2.5, is read as a dependence of the proof of Proposition~14.1\,(iii), it gives two further forward pairs, from Proposition~14.1 to Propositions~26.3 and~26.6.
Through the circles
\[
14.1\to26.6\to26.5\to14.1,\qquad 14.3\to35.10\to14.1\to26.6\to14.3,
\]
the set formed by Propositions~14.3 and~35.10 then grows to the set of Propositions~14.1, 14.3, 35.10, 26.5 and~26.6,
and the following order of the clauses resolves it:
(B12) without the sentence; Proposition~26.5; the order C; Proposition~26.6\,(c), (d), (a), (e), (b); then the sentence;
these entries stand between (B11$'$) and (B13).
\item A datum whose meaning is proved later.
That each of (g1)--(g9) is a condition on the number $e$ alone,
so that $e_\sharp(\beta)$, and with it the datum $0<e\le e_\sharp(\beta)$ of (DAT) (Propositions~34.1, 34.3, 34.4 and~35.5--35.9), is defined,
is stated in Sections~G5.1.8 and~G5.1.9 with a reference to Proposition~35.1\,(c),
and it is proved in Chapter~35c, Sections~3.0 and~3.1, within the proof of Proposition~35.10;
Lemma~0 of Chapter~35c, and through it the proof of Proposition~35.9, refers forward to these two sections.
No circle arises.
Besides definitions (Sections~I.0.5, I.6.1, I.6.2, I.6.4 and~I.8.2; Sections~G5.1.3, G5.1.4, G5.1.8 and~G5.1.9)
and ranges of the constants ((C0), (C7), (C9); the form (G5.T-1) of Section~G5.7; Proposition~14.2\,(a), (f)),
Sections~3.0 and~3.1 of Chapter~35c quote only Proposition~35.1\,(a), (b), (c), (f) and Proposition~35.3\,(c1);
these stand before Proposition~34.1, and their data contain no $e_\sharp(\beta)$;
and Sections~3.0 and~3.1 use neither Lemma~0 nor Section~2 of Chapter~35c.
\item Orders of the clauses inside one statement, stated in the chapter of its proof:
for Proposition~1.5 (Chapter~1b, Sections~1.4--1.6: (b), (d); (a), whose proof uses the integrability of the integrand defining $G_{\mathbf{m}}$ shown in Section~1.4, and (c); (e), (f)),
Proposition~5 (Chapter~5: (a) without its last assertion; (b); (c); (d); (e); the last assertion of (a), on data (SF) with $k=0$, in Section~6),
Proposition~13 (Chapter~13), Proposition~18.5 (Chapter~18b), Proposition~21.1 (Chapter~21a),
Proposition~24.2 (Chapter~24, Sections~8.5--8.8: (a), (b), (d), (c)),
Proposition~26.3 (Chapter~26a, Sections~3.2--3.7: (a), (e), (d), (c), (b)),
Proposition~26.6 (Chapter~26b, paragraphs~2.5--2.7: (c); (d); (a); (e), whose proof uses (d) and the systems constructed for (a); (b)),
Propositions~30.1 and~30.3 (Chapter~30, Section~5.2), Proposition~34.3 (Chapter~34),
Proposition~35.8 (Chapter~35b, Section~4: (a), (b), (f), (g) in paragraph~4.2; (c), (d), (e) in paragraphs~4.3--4.5, without the identity (8.1); the identity (8.1), whose proof in paragraph~4.1 uses (c), (d), (e) for the integral in it),
Proposition~36 (Chapter~36),
Proposition~37.2 (Chapter~37, Section~2: (b) in paragraph~2.2, then (a) in paragraph~2.6),
Proposition~38.3 (Chapter~38, Sections~3.7 and~3.8: (a$'$) after (e)),
Proposition~40 (Chapter~40, Sections~4--6: (b); (a), from the second assertion of (b); (c)),
Proposition~46 (Chapter~46, Section~10: (g+); then (g), the case $v=0$ of (g+)),
Proposition~48.5 (Chapter~48b) and Proposition~50 (Chapter~50).
Apart from those of Propositions~26.3 and~30.3, none of these orders involves another statement;
the order of Proposition~26.3 involves Proposition~26.1 as described in C.
Apart from Proposition~24.2, none of these statements lies on a circle of the relation ``quotes'', the ten pairs included;
Proposition~26.3 lies on no circle because the references between Propositions~26.1 and~26.3 do not belong to the relation.
Proposition~24.2 lies on a circle, and the order of its clauses lies inside the entry (A-6) of A.
The order of Proposition~30.3 involves Proposition~14.2, through a use that is no quotation:
(B1) quotes Proposition~30.3\,(e), and the closing part ``In the flow'' of Proposition~30.3\,(f) takes (B1), (B3) and (B5)
(see the paragraph on the statements, and (B0)--(B5$'$)).
In five of these orders a clause is divided:
\begin{itemize}
\item[] for Proposition~5: (a) without its last assertion; (b)--(e); the last assertion of (a) (Chapter~5, Section~6);
\item[] for Proposition~13: (d1) with the first assertion of (d0); (d2); (d3); the rest of (d0); (d4)--(d6) (Chapter~13, the order of its Sections~3--9);
\item[] for Proposition~50: (a)--(d); the convergence in (e); (f), (g), (g$'$); the rest of (e), because the argument of \cite[Lemma~13.12(b)]{QEDII} used there needs (g); (g$''$), (h) (Chapter~50, order of the proof);
\item[] for Proposition~30.3: (B0), then (B5$'$);
\item[] for Proposition~26.3: the first sentence of (a), which is (C-2); then the Definition in Proposition~26.1, which is (C-3); then the rest of (a), by which the projection is well defined and has the kernels (8.0), within (C-4) (Chapter~26a, Section~3.2).
\end{itemize}
\end{itemize}

\paragraph{E. Inequalities between the constants.}
The following facts are available to every proof, through the conditions on the constants stated with its statement, through the introduction to the Propositions with the same first number, through Annex~B, or, for Propositions~35.5--35.8 and~35.10, through the form (G5.T-1) of Annex~A;
a use of them is no quotation of Proposition~14.2, of Proposition~14.3 or of Proposition~35.10.
\begin{itemize}
\item[(i)] the inequalities between the constants that are rows of Proposition~14.2\,(a), (b), or that are the displays (Z10.s1), (Z10.s2), (Z10.s4) and (Pb) of Proposition~14.2\,(d), or that Proposition~14.2\,(f) derives from these;
\item[(ii)] under the datum $e>0$ with $e\cdot h^{\frac{1}{2}}\le e_1$:
the conditions of Proposition~14.2\,(c) and the sentence ``If $e\cdot h^{\frac{1}{2}}\le e_1$'' of Proposition~14.2\,(f), also in the forms (Sc0), (Sc6) of Section~G4.0.2;
\item[(iii)] $e_0>0$, $e_1>0$, and the second assertion of Proposition~14.3\,(f), namely that (e1)--(e8) hold for all $v\ge-\log e_1$ (Proposition~14.3\,(e), (f));
\item[(iv)] under the datum $0<e\le e_G$: $e\cdot h^{\frac{1}{2}}\le e_1$ and (Z8.w), by the definition of $e_G$ (Proposition~14.2\,(a), row~121),
and (W3), which the first paragraph of Proposition~14.2\,(d), a part of (B5), derives from (Z8.w), $e\cdot h^{\frac{1}{2}}\le e_1$ and (e8);
under the datum $0<e\le e_\sharp(\beta)$ of (DAT): $e\cdot h^{\frac{1}{2}}\le e_1$ and (g1)--(g9), by the definition of $e_\sharp(\beta)$ as one half of a supremum (Section~G5.1.9; Section~I.8.1\,(ii)).
That these conditions contain of the data only $e$ is stated in Proposition~14.2\,(d), a part of (B5) (see also D, on a datum whose meaning is proved later).
Under these two data the sets in the definitions of $e_G$ and of $e_\sharp(\beta)$ are not empty, so that nothing of Proposition~35.10, that is (B13), and nothing of Proposition~14.3\,(g), that is (B14), is used.
\end{itemize}
The inequalities (i) are conditions on the choice of the constants, and the constants are chosen so that they hold (Section~I.8.1; head of Part~II).
The facts (ii)--(iv) follow from them, from the definitions of $e_1$, $e_G$ and $e_\sharp(\beta)$, and from the data: (B2)--(B5).
That every inequality used in a proof is a row of Proposition~14.2\,(a), (b) or a display of Proposition~14.2\,(d), or follows from these,
is shown in the section ``Conditions on the constants'' of the chapter of that proof, which names the row, the display or the line of Proposition~14.2\,(c), (f);
this is the rule stated in the second sentence of the paragraph that follows the conditions of Proposition~14.2\,(b) (see D).
The inequalities $64a'\le608a'\le C^{F}$ of the proof of the closing sentence of Proposition~2.8, which Section~5 of Chapter~2b names with the row $C^{F}$ of the line (C6) of Proposition~14.2\,(a), stand also in Lemma~1\,(K2) of Chapter~14a
(see the kind ``A lemma of a later chapter'' in the first paragraph of this section).
That these rows and displays can be satisfied is Proposition~14.3, and for (g1)--(g9) its clause (g) together with Proposition~35.10.
Propositions~14.2 and~14.3 settle this for all statements at once,
including Proposition~46, Propositions~1.1--1.7, 2.1--2.8, 3--13, 15, 19, 30.1--30.3, Proposition~16.1 and Proposition~14.1, which stand before them;
of the statements that use the inequalities, Propositions~14.2 and~14.3 quote only those named under (B1) and (B14), row by row in the order of the constants (Proposition~14.2\,(e));
these give constants, except Proposition~46, which enters through the entry for the row 26.

The proofs in Chapters~2b, 5, 7, 13, 22, 24, 27a, 28, 49, 35b, 33, 34, 40 and~44 use Proposition~14.2\,(c) or (f) in this way,
mostly in their sections on the conditions on the constants (within a proof: Chapter~24, Section~0.3; Chapter~28, Section~0.4; Chapter~33, Section~0.1),
and the statements proved there do not quote Proposition~14.2.
Chapter~47 uses Proposition~14.2\,(f) in this way, while Proposition~47.2 quotes Proposition~14.2, but not its clause (f).
Chapters~20, 48b, 26b and~43 use Proposition~14.2\,(c), (f), and these clauses are quoted by Propositions~20.3, 48.6, 26.6 and~43.1.
The proofs in Chapters~34, 35b, 35c (the proof of Proposition~35.9) and~44 use displays of Proposition~14.2\,(d) in the ways (i) and (iv), and the statements proved there do not quote Proposition~14.2;
Chapters~42 and~43 use them as well, and Propositions~42.2 and~43.1 quote Proposition~14.2.

The proof of Proposition~31.2 takes (e1) at $v=v_K$ through (Sc6) (Chapter~31, Section~2.1);
behind this assertion of (Sc6) stands the second assertion of Proposition~14.3\,(f), and Proposition~31.2 does not quote Proposition~14.3.
Chapter~49, Section~2.1, does not take (e1) at $v=v_K$ from that assertion:
like Chapter~30, Section~4.5, it derives (e1) at $v=v_K$ from the definition of $e_1$ (Proposition~14.2\,(a), line (C9)) and the datum $e>0$ with $e\cdot h^{\frac{1}{2}}\le e_1$
(the first assertion of Proposition~14.2\,(c), which it names, is a further statement of the fact that it derives);
Proposition~49 quotes no statement among the Propositions~14.m.
Chapters~48b and~43 use this assertion as well, and Propositions~48.6, 43.1 and~43.2 quote Proposition~14.3\,(f).
Chapter~28 uses it at three places (Section~0.4\,(u6); Section~3.0; Section~4, nos.~7 and~10),
and Propositions~28.2 and~28.3 quote it, also in their conditions on the constants.
Chapter~20, Section~3.1; Chapter~26b, paragraphs~2.2 and~2.10; Chapter~32, Section~0.2; and Chapter~42, Step~0, parts~0.1 and~0.2, use Proposition~14.3\,(f), in particular its second assertion, by its number,
and Propositions~20.3, 26.6, 32.1--32.3 and~42.2 quote Proposition~14.3\,(f).
These pairs point backward, and Proposition~14.3 quotes none of the statements from which they start, directly or through other statements;
so they lie on no circle.
Chapters~18b and~48a prove the consequences of the rows that they need, without Proposition~14.2\,(f)
(Chapter~18b, Lemma~1: from (L1), (L3) and the range of $\kappa$; Chapter~48a, Section~4: from (L1)).

Among the chapters that use Proposition~14.2\,(c), (f) or Proposition~14.3\,(f), only Chapter~30 proves statements that Proposition~14.2 or Proposition~14.3\,(a)--(f) quote, namely Proposition~30.1\,(c) and Proposition~30.3\,(e).
It uses $L\ge 2400^8$, $v_K\ge 8$, (2.1) and (e1) at $v=v_K$ only in the closing part ``In the flow'' of Proposition~30.3\,(f),
and the Propositions~14.m enter the proof of Proposition~30.3 only for this part: (B0), (B5$'$).
The remaining chapters among them whose statements stand before the Propositions~14.m, namely Chapters~2b, 5, 7 and~13,
use consequences of the definitions ($\gamma^F\cdot L\le 1/6$, $a'\ge 1$, $\kappa_{\mathrm{w}}-1\ge\kappa+1$),
of the ranges of $\alpha$ and $\kappa$ ($\alpha<1$, $\kappa\ge\kappa_0$)
and of the conditions (M5), (L3) ($M\ge 108\,\pi\cdot C'_V\cdot a'$, $L\,(\kappa-\kappa_0-1)\ge\kappa$),
all of which belong to (B1);
no statement proved in them is quoted by Proposition~14.2 or Proposition~14.3.
Propositions~14.3, 14.4, 20.3, 48.6, 26.6, 29, 47.2, 31.2, 32.1, 32.2, 32.3, 35.1, 35.4, 35.10, 42.2 and~43.1 quote Proposition~14.2;
the reference to Proposition~14.2 for the closing part ``In the flow'' of Proposition~30.3\,(f) is no quotation (see the paragraph on the statements).

\paragraph{Sentences of Part I behind which a statement stands.}
A proof may also use a statement without naming it, by taking from Part~I a sentence behind which that statement stands.
Besides Proposition~46, Part~I contains sentences that are, by their own wording, clauses of a statement or rest on one;
a proof that takes such a sentence from its place in Part~I uses the statement behind it.
\begin{itemize}
\item[(1)] The statement behind the sentence quotes no statement;
so a use of the sentence closes no circle, wherever the statement using it stands.
This is the case for
Section~I.5.2 (the rules (I.N0), (I.N1), (I.N2), (I.N4), (I.N4$'$), (I.N5), (I.N6), (I.S6c)\,(i), (iii), (iv): Proposition~15\,(a), (b), (d), (g), (h), (j), (k), (l); (I.N3)$^{\mathrm{w}}$ with (I.5.2): Proposition~15\,(f));
for Sections~I.3.7, I.4.3 and the last sentence of Section~I.5.11 (what the homomorphism $S_L$ and the automorphism $a_{\tilde{g}}$ need of the side and of the vector): Proposition~16.1\,(b), (d);
for Section~I.4.5 (each of the four maps is compatible with (I.AP), sends every generator to a generator times a number of modulus one, and preserves the weights $\nu$, the bidegree and the norms of Section~I.5.1, on the four families $\psi$, $\chi$, $\Psi$, $W$): Proposition~16.1\,(b);
and for Section~I.1.5 (the bound of $D_k(A)^{-1}(y,y')$ in $|\cdot|_2$) together with $\det D_k(0)\neq 0$ in the paragraph ``Settings'' of Section~I.5.7: Proposition~19\,(e), (a), (d).
\item[(2)] Section~I.5.4, the sentence on $\mathscr{E}|_0$ after (I.k7) ($\mathscr{E}|_0$ lies in $K_k^C$, and $\|\mathscr{E}|_0\|_k\le\|\mathscr{E}\|^x_k$):
Proposition~48.1\,(a), which quotes only Proposition~15\,(a) and Proposition~16.1\,(a1).
Proposition~18.5\,(c) stands before it; it takes the sentence through Section~I.5.4, and its proof gives the reasons (Chapter~18b, proof of (c)).
Propositions~24.1 and~28.1 quote Proposition~48.1\,(a) beside the sentence.
\item[(3)] Section~I.5.7, paragraph ``Settings'' (under (SF) the operators $H_k(\mathscr{T})$, $\bar{H}_k(\mathscr{T})$, $D_k(\mathscr{T})$ exist):
Proposition~14.1\,(i).
It stands before every statement that forms $I^x_k$, $\mathrm{Act}^x_k$, $\mathscr{F}^x_k$, $y^x_k$.
For the Propositions~22.m the fact is proved again (Section~G3.2.0: Proposition~21.4\,(d2), Proposition~21.5\,(b));
Proposition~24.1 takes it through Section~I.5.7 (Chapter~24, Section~9.6\,(0), at the level $k+1$);
Proposition~29 quotes Proposition~14.1\,(i) and Proposition~19.
\item[(4)] Section~I.6.2 ($N_{\min}(R,e)$ exists; $1\le k_*\le K$; $\ell_k(R)$ lies in $\mathrm{Cl}$):
(Sc0), (Sc1), (Sc4) of Section~G4.0.2, that is, \cite[Proposition~19.1]{QEDI} together with \cite[Proposition~19.2]{QEDI} read with the constants and the number $e_1$ of Section~I.8.2.
In this reading \cite[Proposition~19.2]{QEDI} consists of Proposition~14.2\,(c), the condition (W3) of Proposition~14.2\,(d) together with (Z8.w), the sentence ``If $e\cdot h^{\frac{1}{2}}\le e_1$'' of Proposition~14.2\,(f), and the second assertion of Proposition~14.3\,(f);
apart from (W3), it is stated again in the first display of (Sc0) and in (Sc6) of Section~G4.0.2.
This use falls under E\,(ii), (iii);
the only statement taking it that stands before Proposition~14.2 is Proposition~30.3, in the closing part ``In the flow'' of its clause (f): (B5$'$).
(For $R=R^{*}(e)$, Proposition~35.1 proves the existence of $N_{\min}(R,e)$ again.)
\item[(5)] Section~I.6.4, (I.h1)--(I.h5) (the five replacements in \cite{QEDII}):
Proposition~39, which quotes Propositions~46 and~50 and, through Proposition~50, Proposition~38.1.
Propositions~41 and~47.2 stand before it and read the definitions \cite[Definitions~2.18 and~2.19]{QEDII} with (I.h1);
the assertions of \cite{QEDII} in which the widening is made are quoted by Proposition~39 alone.
\item[(6)] Section~I.8.1, (I.U4)$^{\mathrm{ch}}$, second sentence:
Proposition~14.4\,(a).
Propositions~32.1--32.3, which stand after it, quote it, and their proofs take (I.U4)$^{\mathrm{ch}}$ as well.
Proposition~14.1\,(iv), which stands before it, is proved without it.
Propositions~3 and~8, which stand before it, use Section~I.8.1 with Section~G1.1 and take the sentence for the numbers in their bounds, which are constants of the lines (C4), (C4r), (C5), (C6) (Chapter~3, Section~4; Chapter~8, Section~4).
For these letters Proposition~14.4\,(a) is a part of (B15), which rests on (B00)--(B5) only and quotes no statement of (B11$'$);
it can be established before (B11$'$), so that no circle between clauses arises.
Chapter~33 (end of Section~2) and Chapter~35a, Section~4.10, take the sentence after Proposition~14.4,
while Propositions~33.2 and~35.4 do not quote Proposition~14.4.
The sentence is taken after Proposition~14.4, without a reference to Proposition~14.4 by its number, also in Chapter~20, paragraph~3.6, in the proof of Proposition~20.3;
in Chapter~48a, Section~5, for Propositions~48.1--48.4;
in Chapter~27a, paragraph ``Uniformity'', in the proof of Proposition~27.1;
and in Chapter~27b, paragraph ``Range of validity'', in the proof of Proposition~27.5.
(Of (I.U4)$^{\mathrm{ch}}$, Proposition~43.1 takes the first sentence, ``a constant is a number'';
the assertion behind it, Proposition~14.2\,(f), paragraph ``Ranges'', is quoted by Proposition~43.1.)
\item[(7)] References.
The sentence says where a fact is proved or an object is formed, and the proofs that take it stand after the statement named:
\begin{itemize}
\item[] Section~I.1.8 (the primitives of (FL): Proposition~26.5; they are defined under the four conditions of (FL), so that (FL) and (S-4) are the same condition: Proposition~26.6\,(e));
\item[] Section~I.2.9, the Remark after Proposition~46 (the independence of the placement in clause (g+) of Proposition~46: Proposition~50\,(h));
\item[] Section~I.3.5 (existence of the inverses: Propositions~2.8 and~4);
\item[] Section~I.3.7 (division by $L$ carries a pair that is admissible for $X$, with $X$ contained in $LY$, to a pair of the same kind that is admissible for $Y$: Proposition~18.3\,(e); its proof, Chapter~18a, Section~3.4, takes from Section~I.3.7 the definitions only, and Chapter~24, (Sc2), derives the fact again);
\item[] Section~I.4.4 (covariances of the composed operators: Proposition~10);
\item[] Section~I.4.6 (a single $\mathrm{al}_2(E,X)$ need not be a multiple of $1$: the Remark after Proposition~16.2; the seven light-heavy couplings are allowed by the symmetry: Proposition~18.6\,(f); neither sentence is used in a proof);
\item[] Section~I.5.3 ((I.Lc2)$^{\mathrm{ch}}$: Proposition~6\,(b), by the wording of that clause; the proofs of Propositions~24.1 and~24.3 take (I.Lc2)$^{\mathrm{ch}}$ without quoting Proposition~6, which is used in Chapter~24, Section~5.1, while Proposition~35.4 takes (I.Lc2)$^{\mathrm{ch}}$ and quotes Proposition~6; DF5: the clause (DF5) of Proposition~35.4);
\item[] Section~I.5.6 ($m({E})$ is a multiple of $1$: Proposition~16.3\,(a); without slot: Proposition~48.2\,(a));
\item[] Section~I.5.7 (the step map is defined on $\mathrm{Ball}^x_k(4)$: Proposition~24.1\,(a); the data of a flow are (OS), (SF): Proposition~26.6\,(b), (c); level $0$ integrates to the determinant: Proposition~22.2\,(c));
\item[] Section~I.5.9 (the facts on the vertex: Proposition~23.2; Proposition~22.4\,(b) proves the identity again);
\item[] Section~I.5.10 ((I.L0.13): the bound of Proposition~1.7, and, at the levels $k\ge 1$, the input at zero field, Proposition~1.6);
\item[] Section~I.5.11 (in the transition to global fields, $\mathscr{C}=0$ gives $\mathscr{C}^{b}=0$: Proposition~18.2\,(o); the invariant $E_k=\mathscr{E}_k|_0$: Proposition~48.6\,(b));
\item[] Section~I.6.1 (the mass range: Proposition~14.1\,(v));
\item[] Section~I.6.2 (the tuned pair: Proposition~28.3; $m_{0,N}(e)$ is real, of modulus $<1/8$ and independent of $T$: Proposition~28.3\,(d));
\item[] Section~I.6.3 (the projection of a system of primitives: the Definition in Proposition~26.1, with the display (8.0) of Proposition~26.3\,(a));
\item[] Section~I.6.5 (the range of the transfer: Propositions~40, 42.1 and~42.2);
\item[] Sections~I.7.3 and~I.7.4 (identification, shift identity: Proposition~38.2\,(b), Proposition~38.3\,(b), (c));
\item[] Section~I.7.6 ($\boldsymbol{\sigma}^2$ and $N_0$ do not depend on the placement, and $N_0$ does not depend on the cell: Proposition~47.1\,(b5$'$), Proposition~50\,(h), Proposition~49\,(d), Proposition~36\,(h); the probe lies in the open cell: Proposition~41; (R0)--(R6): Proposition~43.1; together with Proposition~42.2 they give the ultraviolet proposition, Proposition~43.2);
\item[] Sections~I.8.1, I.8.2 and~I.8.4 (the complete list of the conditions, and the fact that their order contains no circle: Proposition~14.2, with its clause (e); admissible constants and admissible pairs $(\beta,x_0)$ exist: Proposition~14.3; how the proofs use these two statements is described in E above; some of the fermionic constants that carry $\rho$ degenerate as $\rho\to 0$: Proposition~14.4\,(d), a sentence of Section~I.8.1 that is used in no proof);
\item[] Section~I.8.5 (the bound of the fluctuation field: Proposition~24.3\,(q1); Propositions~23.2, 24.1 and~24.2, which stand before it, quote this clause: A);
\item[] Section~I.9.1 (the continuum data of the chamber operator: Propositions~37.1\,(f) and~37.2\,(f); the relations that take the place of the hypothesis for the flow: Proposition~26.6; the property DF5: the clause (DF5) of Proposition~35.4).
\end{itemize}
The remaining places at which Part~I names a statement are of the kind (7) or of the kinds listed in D:
remarks on wording, on letters and on the statements in which a clause is applied;
the statements that Section~I.8.2 names beside the constants they give;
and the places of Section~I.9, which sets the statements of \cite{QEDI} and \cite{QEDII} beside the Propositions that correspond to them or take their places, and which repeats (4) for \cite[Proposition~19.2]{QEDI}.
\end{itemize}

Consequently the uses through Part~I that point forward in the order of Part~II are
the uses of Proposition~14.4 by Propositions~3 and~8,
the use of Proposition~48.1 by Proposition~18.5,
the uses of Proposition~39 by Propositions~41 and~47.2,
and (B5$'$);
the setting (ST*)$^{\mathrm{ch}}$, which Proposition~14.2\,(d) names, gives no further one (see below).
Also when these uses are read as dependences, neither Proposition~48.1 nor Proposition~39 quotes, directly or through other statements, a statement that uses it:
Proposition~48.1 reaches only Propositions~15 and~16.1,
and Proposition~39 reaches only Propositions~46, 38.1 and~50.
If the use of Proposition~14.4 by Propositions~48.1--48.4 through Section~I.8.1, which points backward (see (6)), is read as a dependence as well,
Proposition~48.1 reaches in addition Proposition~14.4 and the statements that Proposition~14.4 reaches,
which lie among Propositions~46, 1.m, 2.m, 3--13, 15, 16.1, 30.m, 14.1--14.3, 35.1--35.3, 35.10, 36, 37.1 and~37.2;
none of Propositions~18.5, 20.3, 27.1, 27.5 and~48.m is among them.
Proposition~14.4 reaches Propositions~3 and~8 through
\[
\text{Proposition }14.3\to\text{Proposition }35.10\to\text{Proposition }14.1,
\]
that is, through the clause Proposition~14.3\,(g).
With the two uses of (6), the set formed by Propositions~14.3 and~35.10 is enlarged to the set of the Propositions~3, 7, 8, 14.1, 14.3, 14.4 and~35.10, any two of which lie on a common circle
(Proposition~14.1 quotes Propositions~3, 7 and~8, and Proposition~7 quotes Proposition~3);
the order of the clauses resolves it: of Proposition~14.3, the entry (B15) takes only (c) and (f).
The uses that point backward
(among them, without a reference by number: the use of Propositions~14.1 and~19 by Proposition~24.1; the use of Proposition~6 by Propositions~24.1 and~24.3; the use of Proposition~14.4 by Propositions~20.3, 27.1, 27.5, 33.2, 35.4 and~48.1--48.4)
start at statements that the statement behind the sentence does not quote, directly or through other statements,
with one exception.
The statements behind the sentences of Sections~I.1.8 and~I.6.2 on (FL) and on the setting (ST*)$^{\mathrm{ch}}$, Propositions~26.6 and~28.3, quote Proposition~14.3\,(f), and through Proposition~14.3 they reach Propositions~35.2, 35.3 and~35.10.
Proposition~35.2, the clauses of Proposition~35.3 other than (b1) and (c1), and the middle sentence of Proposition~35.3\,(b1) assume that setting;
Proposition~35.10 has no setting among its data, and only the proof of its assertions for $0<e\le e_\sharp(\beta)$ (Chapter~35c, Lemma~1 and paragraph~3.6) is made in the setting, together with the notions of placement used in these assertions.
These three statements make no use of the two sentences, since what they and their proofs involve of the setting is only (S-4) at $R^{*}(e)$, and neither the pair $P_N(e)$ nor a primitive ((B9), (B13)).
Propositions~26.6 and~28.3 reach Proposition~14.2 as well, and Proposition~14.2 stands before them, so that a use of the two sentences by Proposition~14.2 would point forward.
Proposition~14.2\,(d) names the setting (ST*)$^{\mathrm{ch}}$, but only to fix the meaning of the index $*$ and the numbers $\mathrm{pp}$ and $d_{\mathrm{sep}}$;
what (d) and its proof involve of the setting is again only (S-4) at $R^{*}(e)$, and neither the pair $P_N(e)$ nor a primitive (see (B5)).
So Proposition~14.2 makes no use of the two sentences either.
Hence no circle between clauses arises,
and apart from the cases of Proposition~14.4 and of (B5$'$) the sets of statements that lie on a common circle remain as they are
(for what changes if the use in (B5$'$) is read as a dependence, see the paragraph on the statements).
If the two uses of (6) and the use in (B5$'$), as a dependence on Propositions~14.2 and~14.3, are all read as dependences, the set of seven and the set of nine, which share Propositions~14.3 and~35.10, form one set,
of the fourteen Propositions~3, 7, 8, 14.1, 14.2, 14.3, 14.4, 30.3, 35.1, 35.2, 35.3, 35.10, 36 and~37.2 (the pair from Proposition~35.1 to Proposition~14.2 included), and no other statement joins it;
the order B with its two refinements stated in (B0) and (B11$'$) resolves it.
If the use in (B5$'$) is read as a dependence on Proposition~14.2 alone, the set of seven and the set formed by Propositions~30.3 and~14.2 remain separate,
since Propositions~14.2 and~30.3 reach, besides each other, only Propositions~46, 1.m and~30.m.

\section*{Part V. The proofs}

\section*{Chapter 46. The operator dictionary: proof of Proposition~46}

\subsection*{0. Conventions of this chapter}

The displays of \cite{QEDI} and of \cite{QEDII} that the clauses of Proposition~46 name are read with the data of Proposition~46, as its preamble prescribes.

In this chapter an element $z$ of $\mathbb{Z}^3$ is the index of the period $\mathrm{nz}$ defined in 2.3, and $z(\cdot)$ denotes a map with values in $\mathbb{Z}^3$.
Apart from the placement $v$ of Proposition~46\,(g+), the letter $v$ denotes a vector of $\mathbb{Z}^3$ that is added to a site (2.2, 7.1, 7.4, 7.6) or a vector of $\mathbb{R}^3$ (2.1, 5.2, 8.1).
The numbers $v_1$, $v_2$, $v_3$ are the coordinates of $v_0$, except in the display (46.1) of 2.1 and in the proof of (I.c1) in 5.2, where $v_j$ and $v_\mu$ are the coordinates of the vector $v$ named there.
The numbers $n_1$, $n_2$, $n_3$ are the sides of a torus, except in 4.3 and 4.4, where $n_0,\dots,n_3$ are the vectors $\hat n_0,\dots,\hat n_3$.
Indices of the four spin components are written $\beta$, $\beta'$.
The transpose ${}^{T}$ is applied to vectors of $\mathbb{R}^3$ and to $2\times2$ and $4\times4$ matrices only.
Besides the curvature, the letters $F$, $F_0,\dots,F_3$ and $\mathrm{Top}$ denote elements of a Grassmann algebra (Section~I.3.6 and (I.3.5)), and $\mathrm{Dm}$ denotes a complex matrix.

When a display of \cite{QEDI} or of \cite{QEDII} is reproduced in order to be compared, it keeps the letters of the paper from which it is taken, except for bound variables (points, sites, vectors, indices, and the function to which an operator or an integral is applied) and except where the contrary is stated.
In particular, the function on which an operator acts is written $\mathrm{fX}$ when it is defined on the set $X$ of the fermion sites of \cite{QEDII}, and $f$ when it is defined on $\mathbb{Z}^3\times O_{\mathrm{cub}}$; the letters $\psi$, $\bar{\psi}$ are reserved for the Grassmann generators of 7.6.

The following letters have the meaning given to them in \cite{QEDII}.
$W$ is the group generated by the reflections in all mirrors \cite[Definition~2.2]{QEDII}.
The fundamental chamber $\mathrm{Chm}_{0}=\left\{y\in\mathbb{R}^3:\ 0<y_3<y_2<y_1<\tfrac12\right\}$ is written $C_0$ in \cite[Definition~2.3]{QEDII}; its walls are written $H_0,\dots,H_3$ here and $\Pi_0,\dots,\Pi_3$ there, and $\mathrm{Ar}_i$ denotes the area of the face $H_i\cap\overline{\mathrm{Chm}_{0}}$.
$X=Wv_0$ is the set of the fermion sites of \cite[Definition~2.4]{QEDII} (written $\mathbb{X}$ there), $x^{(i)}$ is the neighbour of a site $x$ of $X$ across the wall of type $i$, $\pi$ is the map to the nearest photon site, and $\mathrm{dlt}_i(x):=x^{(i)}-x$.
Where a statement of \cite{QEDII} is reported, $x$, $x'$ and $x^{(i)}$ are sites of $X$, that is, points of $\mathbb{R}^3$, and $\mathrm{gb}$ is the linear part of the element of $W$ that maps $v_0$ to $x$ (written $\vec g$ in \cite{QEDII}).
The parameter $M$ of \cite[Definition~2.10]{QEDII} is the mass parameter, written $M_{\mathrm{ch}}$ here; $m$ is the lattice mass and $h(\cdot,\cdot)$ the hopping matrix of the same definition (where the hopping matrix is written $w(x,x')$).
The operator of \cite[Definition~2.10]{QEDII} on the spinors of a torus $\mathbb{T}$ is written $D^{\mathbb{T}}_A$ in \cite[Definition~2.11]{QEDII}; its upper index is the torus, not a transpose.
The numbers $d_i=\operatorname{dist}(v_0,H_i)$, $i=0,\dots,3$, are those of \cite[Lemma~3.3]{QEDII}; in particular $d_0=v_3$.
In the paragraph of 5.3 on \cite[Section~6.3.2]{QEDI}, $w$ is the unitary matrix chosen there and $P^{\pm}_{\nu}$ are the matrices $\left(1\pm\gamma_\nu\right)/2$ of \cite[Section~6.1.3]{QEDI}; $s$ is the number of the spin components in \cite[Section~6.1.2]{QEDI}.
The operators $\mathcal{D}_A$, $D^{\nabla}_A$ and $D^{\mathrm{ch}}[B]$ are defined in 4.2, 4.3 and 4.4.

\subsection*{1. The statement}

The statement is Proposition~46 (Section~I.2.9).
This chapter contains the proof of its clauses (a)--(g+).
Clause (h) is proved in Section~I.2.9, from the clauses (a), (b), (e) and the display (46.7b).

\subsection*{2. Notation of the proofs}

\subsubsection*{2.1. The point group}

The columns of a matrix of $O_{\mathrm{cub}}=\mathrm{O}(3)\cap\mathrm{GL}(3,\mathbb{Z})$ are vectors of $\mathbb{Z}^3$ of Euclidean length one, and they are pairwise orthogonal.
The only vectors of $\mathbb{Z}^3$ of Euclidean length one are $\pm e_j$.
Consequently every $g\in O_{\mathrm{cub}}$ satisfies
\begin{equation}
g\,e_j=\epsilon_j\cdot e_{\mathrm{sp}(j)}\quad(j=1,2,3),\qquad\text{hence}\qquad\left(g\,v\right)_{\mathrm{sp}(j)}=\epsilon_j\cdot v_j\quad\left(v\in\mathbb{R}^3\right),
\tag{46.1}
\end{equation}
with a permutation $\mathrm{sp}$ of $\{1,2,3\}$ and signs $\epsilon_j\in\{1,-1\}$.
Every pair $(\mathrm{sp},\epsilon)$ occurs, so that $O_{\mathrm{cub}}$ has $6\cdot8=48$ elements.

\subsubsection*{2.2. The displays of Part~I in lattice units}

At $k=0$, with $\eta_0=1$ and with $q$ in the place of $e_k$, Sections~I.1.1, I.1.2, I.1.3 and the displays (I.1.6), (I.2.1), (I.2.3) take the following form:
\begin{gather*}
\hat n_0=e_3,\qquad\hat n_1=\frac{e_2-e_3}{\sqrt2},\qquad\hat n_2=\frac{e_1-e_2}{\sqrt2},\qquad\hat n_3=-e_1;\qquad s_i=1-2\,\hat n_i\hat n_i^{T};\\
s_i^2=1,\qquad s_i\hat n_i=-\hat n_i,\qquad s_3e_1=-e_1;\\
w=\left(6,6\sqrt2,6\sqrt2,6\right),\qquad\mathrm{sw}=12+12\sqrt2;\qquad\sum_iw_i\cdot\hat n_i=6e_3+6\left(e_2-e_3\right)+6\left(e_1-e_2\right)-6e_1=0\qquad\text{[(I.1.1)]};\\
x=(c,g)\in\mathbb{Z}^3\times O_{\mathrm{cub}}:\qquad x^{(i)}=\left(c,gs_i\right)\quad(i=0,1,2),\qquad x^{(3)}=\left(c+ge_1,gs_3\right),\qquad\hat d_i(x)=-g\hat n_i\qquad\text{[(I.1.2)]};\\
M_i(x)=\frac{w_i}{2}\cdot\left(\rho-\gamma\cdot\hat d_i(x)\right)=\frac{w_i}{2}\cdot\left(\rho+\gamma\cdot\left(g\hat n_i\right)\right)\qquad\text{[(I.2.1)]}.
\end{gather*}

A one-cochain $A$ on $\mathbb{Z}^3$ is a complex function on the ordered pairs $(c,c')$ of nearest neighbours with $A(c',c)=-A(c,c')$; one writes $A_\mu(c)=A(c,c+e_\mu)$ and puts $A(c,c):=0$ (Section~I.1.1, in agreement with \cite[Definition~2.1]{QEDII}).
By (I.1.6), $U_A(x,x')=\exp\left(iq\cdot(\tau A)(c,c')\right)$, where
\[
(\tau A)(c,c')=\frac16\sum_{\pi}A\left(\boldsymbol{\Gamma}^{\pi}(c,c')\right),\qquad(\tau A)(c,c)=0 .
\]
If $c$ and $c'$ are nearest neighbours, each of the six paths $\boldsymbol{\Gamma}^{\pi}(c,c')$ consists of the single bond $(c,c')$, and therefore $(\tau A)(c,c')=\frac16\cdot6\cdot A(c,c')=A(c,c')$.
It follows that
\begin{gather}
U_A\left(x,x^{(i)}\right)=1\quad(i=0,1,2),\qquad U_A\left(x,x^{(3)}\right)=\exp\left(iq\cdot A\left(c,c+ge_1\right)\right),\tag{46.2}\\
\left(D_Af\right)(x)=-\sum_{i=0}^{3}M_i(x)\cdot\left[U_A\left(x,x^{(i)}\right)\cdot f\left(x^{(i)}\right)-f(x)\right]\qquad\text{for every }f\colon\mathbb{Z}^3\times O_{\mathrm{cub}}\to\mathbb{C}^4\qquad\text{[(I.2.3)]}.\tag{46.3}
\end{gather}
For $v\in\mathbb{Z}^3$ let $(c,g)+v:=(c+v,g)$.
Then
\begin{equation}
\left(x^{(i)}\right)^{(i)}=x,\qquad\hat d_i\left(x^{(i)}\right)=-\hat d_i(x),\qquad(x+v)^{(i)}=x^{(i)}+v.
\tag{46.4}
\end{equation}
Indeed, for $i\le2$ one has $\left(c,gs_is_i\right)=(c,g)$, and for $i=3$
\[
\left(c+ge_1+gs_3e_1,\ gs_3s_3\right)=\left(c+ge_1-ge_1,\ g\right)=(c,g);
\]
moreover $\hat d_i\left(x^{(i)}\right)=-gs_i\hat n_i=g\hat n_i$; in the third relation $v$ is added to the first entry on both sides.

\subsubsection*{2.3. The torus}

For $z\in\mathbb{Z}^3$ put
\[
\mathrm{nz}:=\sum_{\mu}z_\mu n_\mu e_\mu,\qquad\boldsymbol{\chi}(z):=(-1)^{z_1+z_2+z_3}
\]
(Section~I.6.3), and $\mathrm{Ga}_n:=\left\{\mathrm{nz}:\ z\in\mathbb{Z}^3\right\}=\sum_{\mu}n_{\mu}\mathbb{Z}\,e_{\mu}$.
A function $f\colon\mathbb{Z}^3\times O_{\mathrm{cub}}\to\mathbb{C}^4$ belongs to $\mathfrak{S}^{\mathrm{ch}}(1;n)$ if and only if $f(c+\mathrm{nz},g)=\boldsymbol{\chi}(z)\cdot f(c,g)$ for all $c$, $g$ and $z$.
In fact the three conditions of the definition are the cases $z=e_\mu$, and the condition in the direction $\mu$, applied $|z_\mu|$ times (to $c$ or to $c-n_\mu e_\mu$), produces the factor $\prod_\mu(-1)^{|z_\mu|}=\boldsymbol{\chi}(z)$.
Let
\[
\mathrm{Fb}:=\prod_{\mu}\left\{0,\dots,n_\mu-1\right\},\qquad\mathrm{Yf}:=\mathrm{Fb}\times O_{\mathrm{cub}},\qquad\text{so that}\quad\#|\mathrm{Yf}|=48\,|n| .
\]
Division with remainder in each coordinate shows that every $c\in\mathbb{Z}^3$ can be written as $c=c^0+\mathrm{nz}$ with exactly one pair $c^0\in\mathrm{Fb}$, $z\in\mathbb{Z}^3$.

All sites $x^{(i)}$ and all bonds $(c,c+ge_1)$ are formed in $\mathbb{Z}^3$, and this is how they are read for all sides $n_\mu\ge1$.
For $n_\mu=1,2$ two different bonds of $\mathbb{Z}^3$ may have the same end points modulo $\mathrm{Ga}_n$; the displays are never reduced modulo $\mathrm{Ga}_n$.
The rule of Section~I.1.1, by which on the torus the paths are formed on a lift with $|c-c'|<\ell/2$, does not apply to nearest neighbours in a direction $\mu$ with $n_\mu\le2$: for $n_\mu=2$ no representative of $1$ modulo $2$ has modulus less than $1$, and for $n_\mu=1$ the two ends are the same point of the torus.
For $n_\mu\ge3$ the lift with $|c-c'|=1<n_\mu/2$ is the only one.

\emph{Spinors with four components.}
The space $\mathfrak{S}(1;n)$ consists of the functions $\xi\colon\mathbb{Z}^3\to\mathbb{C}^4$ with $\xi(c+n_\mu e_\mu)=-\xi(c)$ for $\mu=1,2,3$.
By (I.3.1),
\[
(J\xi)(c,g)=\xi(c),\qquad\left(J^{*}f\right)(c)=\frac{1}{48}\sum_{g\in O_{\mathrm{cub}}}f(c,g).
\]
For $\xi\in\mathfrak{S}(1;n)$ one has $(J\xi)(c+n_\mu e_\mu,g)=\xi(c+n_\mu e_\mu)=-(J\xi)(c,g)$.
For $f\in\mathfrak{S}^{\mathrm{ch}}(1;n)$ one has
\[
\left(J^{*}f\right)(c+n_\mu e_\mu)=\frac{1}{48}\sum_{g}f(c+n_\mu e_\mu,g)=-\left(J^{*}f\right)(c).
\]
Hence $J$ maps $\mathfrak{S}(1;n)$ into $\mathfrak{S}^{\mathrm{ch}}(1;n)$, and $J^{*}$ maps $\mathfrak{S}^{\mathrm{ch}}(1;n)$ into $\mathfrak{S}(1;n)$.
Over a photon site $c$, an element of $\mathfrak{S}^{\mathrm{ch}}(1;n)$ has the $48\cdot4=192$ components $f_\beta(c,g)$, and an element of $\mathfrak{S}(1;n)$ has the four components $\xi_\beta(c)$.

\subsection*{3. Proof of (a)}

\subsubsection*{3.0. The group}

In the letters of Section~I.4.8, \cite[Definition~2.2]{QEDII} reads as follows.
$\mathcal{L}$ is the set of the nine linear forms $\ell=y_\mu$ $(\mu=1,2,3)$ and $\ell=y_\mu-y_\nu$, $\ell=y_\mu+y_\nu$ $(\mu<\nu)$ on $\mathbb{R}^3$, and $\ell(y)=n_\ell\cdot y$; the squared Euclidean length of $n_\ell$ is written $|n_\ell|_2^{2}$ (the index $2$ is explained in 3.4).
The admissible positions are $\operatorname{Adm}(\ell)=\frac12\mathbb{Z}$ for $\ell=y_\mu$ and $\operatorname{Adm}(\ell)=\mathbb{Z}$ for the six other forms.
For $\mathrm{po}\in\operatorname{Adm}(\ell)$,
\[
\theta_{\ell,\mathrm{po}}(y)=y-\frac{2}{|n_\ell|_2^{2}}\cdot\left(\ell(y)-\mathrm{po}\right)\cdot n_\ell ,
\]
and $W$ is the group generated by all $\theta_{\ell,\mathrm{po}}$.

\begin{named}{Lemma 3.0}
$W=W_{\mathrm{af}}=\left\{y\mapsto ry+a:\ a\in\mathbb{Z}^3,\ r\in O_{\mathrm{cub}}\right\}$.
\end{named}

\begin{proof}
Write $(a,r)$ for the map $y\mapsto ry+a$.
The set $W_{\mathrm{af}}$ is a group, because
\[
(a',r')(a,r)=\left(a'+r'a,\ r'r\right),\qquad(a,r)^{-1}=\left(-r^{-1}a,\ r^{-1}\right).
\]
One has
\[
\theta_{\ell,\mathrm{po}}(y)=r_\ell\,y+\frac{2\,\mathrm{po}}{|n_\ell|_2^{2}}\cdot n_\ell,\qquad r_\ell:=1-\frac{2}{|n_\ell|_2^{2}}\cdot n_\ell n_\ell^{T}.
\]
Let $\ell=y_\mu$.
Then $n_\ell=e_\mu$ and $|n_\ell|_2^{2}=1$; the matrix $r_\ell=1-2e_\mu e_\mu^{T}$ changes the sign of the coordinate $\mu$, and $\frac{2\,\mathrm{po}}{|n_\ell|_2^{2}}\cdot n_\ell=2\,\mathrm{po}\cdot e_\mu$ lies in $\mathbb{Z}^3$, since $2\,\mathrm{po}$ is an integer.
Let $\ell=y_\mu-\epsilon y_\nu$ with $\epsilon=\pm1$ and $\mu<\nu$.
Then $n_\ell=e_\mu-\epsilon e_\nu$ and $|n_\ell|_2^{2}=2$;
\[
r_\ell e_\mu=e_\mu-\left(e_\mu-\epsilon e_\nu\right)=\epsilon e_\nu,\qquad r_\ell e_\nu=e_\nu+\epsilon\left(e_\mu-\epsilon e_\nu\right)=\epsilon e_\mu,\qquad r_\ell e_{\mathfrak{l}}=e_{\mathfrak{l}}
\]
for the third index $\mathfrak{l}$, and $\frac{2\,\mathrm{po}}{|n_\ell|_2^{2}}\cdot n_\ell=\mathrm{po}\cdot n_\ell$ lies in $\mathbb{Z}^3$.
So every generator of $W$ lies in $W_{\mathrm{af}}$.
Conversely,
\[
\theta_{y_\mu,\frac12}\left(\theta_{y_\mu,0}(y)\right)=\theta_{y_\mu,\frac12}\left(y-2y_\mu e_\mu\right)=y-2y_\mu e_\mu-2\left(-y_\mu-\tfrac12\right)e_\mu=y+e_\mu ,
\]
and $W$ contains the translations by $\mathbb{Z}^3$.
It contains $r_\ell=\theta_{\ell,0}$ for the nine forms, among them the three changes of sign of one coordinate and, for $\epsilon=1$, the three transpositions of two coordinates.
By (46.1), an element $g$ of $O_{\mathrm{cub}}$ is the product of the permutation matrix $e_j\mapsto e_{\mathrm{sp}(j)}$, which is a product of transpositions, and of the diagonal matrix $e_j\mapsto\epsilon_je_j$, which is a product of changes of sign of one coordinate.
Thus $W$ contains $O_{\mathrm{cub}}$ and $\mathbb{Z}^3$, hence $W_{\mathrm{af}}$.
\end{proof}

\subsubsection*{3.1. The map $\mathrm{geo}$ is a bijection; the map $\pi$}

Let $v_0=(v_1,v_2,v_3)$ with $0<v_3<v_2<v_1<\frac12$.
By (46.1), $\left(g\,v_0\right)_{\mathrm{sp}(j)}=\epsilon_j\cdot v_j$, and therefore $\left|g\,v_0\right|_\infty=v_1<\frac12$.

\emph{Injectivity.}
Suppose that $c+g\,v_0=c'+g'v_0$.
Then $c-c'=g'v_0-g\,v_0$ is a vector of $\mathbb{Z}^3$ of sup norm at most $2v_1<1$, so that $c=c'$ and $g\,v_0=g'v_0$.
The moduli $v_1$, $v_2$, $v_3$ are pairwise different and non-zero.
Hence the vector $g\,v_0$ determines, for each $j$, the index $\mathrm{sp}(j)$ (the coordinate whose modulus is $v_j$) and the sign $\epsilon_j$ (the sign of that coordinate), and (46.1) gives $g=g'$.

\emph{Surjectivity.}
By Lemma~3.0,
\[
X=W_{\mathrm{af}}\,v_0=\left\{rv_0+a:\ a\in\mathbb{Z}^3,\ r\in O_{\mathrm{cub}}\right\}=\left\{\mathrm{geo}(a,r):\ a\in\mathbb{Z}^3,\ r\in O_{\mathrm{cub}}\right\}.
\]

\emph{The map $\pi$.}
In \cite[Definition~2.4]{QEDII} the map $\pi\colon X\to\mathbb{Z}^3$ is defined by $\pi(c+g\,v_0):=c$ for $c\in\mathbb{Z}^3$ and $g\in O_{\mathrm{cub}}$ (the nearest photon site).
By injectivity the point $c+g\,v_0$ determines $c$, and the definition says that $\pi\left(\mathrm{geo}(c,g)\right)=c$.
For $c'\in\mathbb{Z}^3$ with $c'\neq c$,
\[
\left|\mathrm{geo}(c,g)-c'\right|_\infty\ge\left|c-c'\right|_\infty-\left|g\,v_0\right|_\infty\ge1-v_1>\frac12>v_1=\left|\mathrm{geo}(c,g)-c\right|_\infty .
\]

\subsubsection*{3.2. The two actions}

Let $g_{\mathrm{w}}=(a,r)$ and $g_{\mathrm{w}}'=(a',r')$ be elements of $W_{\mathrm{af}}$, and put $g_{\mathrm{w}}(c,g):=(rc+a,\ rg)$.
Then
\[
g_{\mathrm{w}}'\left(g_{\mathrm{w}}(c,g)\right)=\left(r'(rc+a)+a',\ r'rg\right)=\left(a'+r'a,\ r'r\right)(c,g)=\left(g_{\mathrm{w}}'g_{\mathrm{w}}\right)(c,g),
\]
so that this is an action of $W_{\mathrm{af}}$ on $\mathbb{Z}^3\times O_{\mathrm{cub}}$.
Moreover
\[
g_{\mathrm{w}}\left(\mathrm{geo}(c,g)\right)=r\left(c+g\,v_0\right)+a=(rc+a)+(rg)\,v_0=\mathrm{geo}\left(g_{\mathrm{w}}(c,g)\right).
\]
The element $\mathrm{gt}:=(c,g)$ of $W_{\mathrm{af}}$, that is, $\mathrm{gt}\,y=g\,y+c$, satisfies $\mathrm{geo}(c,g)=\mathrm{gt}\,v_0$, and by 3.1 it is the only element of $W_{\mathrm{af}}$ with this property.

\subsubsection*{3.3. Neighbours}

By \cite[Definition~2.3]{QEDII} the walls of $\mathrm{Chm}_{0}$ are
\[
H_0=\left\{y_3=0\right\},\qquad H_1=\left\{y_2=y_3\right\},\qquad H_2=\left\{y_1=y_2\right\},\qquad H_3=\left\{y_1=\tfrac12\right\};
\]
here $\mathrm{st}_i$ denotes the reflection of $\mathbb{R}^3$ in $H_i$, which is written $s_i$ in \cite{QEDII}.
By \cite[Definition~2.4]{QEDII}, for a site $\mathrm{gt}\,v_0$ of $X$, with $\mathrm{gt}\in W$ unique, and for $i\in\{0,1,2,3\}$, the site $\mathrm{gt}\,\mathrm{st}_iv_0$ is the neighbour of $\mathrm{gt}\,v_0$ across the wall of type $i$.

For a unit vector $\hat n$ and a real number $\mathrm{pw}$, the reflection in the plane $\left\{y:\ \hat n\cdot y=\mathrm{pw}\right\}$ is the map $y\mapsto y-2\left(\hat n\cdot y-\mathrm{pw}\right)\hat n$.
Let $\mathrm{pw}_0=\mathrm{pw}_1=\mathrm{pw}_2:=0$ and $\mathrm{pw}_3:=-\frac12$.
Since
\[
\hat n_0\cdot y=y_3,\qquad\hat n_1\cdot y=\frac{y_2-y_3}{\sqrt2},\qquad\hat n_2\cdot y=\frac{y_1-y_2}{\sqrt2},\qquad\hat n_3\cdot y=-y_1,
\]
one has $H_i=\left\{y:\ \hat n_i\cdot y=\mathrm{pw}_i\right\}$, and therefore
\[
\mathrm{st}_i\,y=y-2\left(\hat n_i\cdot y\right)\hat n_i=s_i\,y\quad(i=0,1,2),\qquad\mathrm{st}_3\,y=y-2\left(-y_1+\tfrac12\right)\left(-e_1\right)=\left(y-2y_1e_1\right)+e_1=s_3\,y+e_1 .
\]
Consequently, for $x=(c,g)$ and $\mathrm{gt}=(c,g)$,
\begin{align*}
\mathrm{gt}\,\mathrm{st}_iv_0&=c+gs_iv_0=\mathrm{geo}\left(c,gs_i\right)=\mathrm{geo}\left(x^{(i)}\right)\qquad(i=0,1,2),\\
\mathrm{gt}\,\mathrm{st}_3v_0&=c+g\left(s_3v_0+e_1\right)=\left(c+ge_1\right)+\left(gs_3\right)v_0=\mathrm{geo}\left(c+ge_1,gs_3\right)=\mathrm{geo}\left(x^{(3)}\right).
\end{align*}

\subsubsection*{3.4. Directions}

Let $\mathrm{dw}_i:=\hat n_i\cdot v_0-\mathrm{pw}_i$, that is,
\[
\mathrm{dw}_0=v_3,\qquad\mathrm{dw}_1=\frac{v_2-v_3}{\sqrt2},\qquad\mathrm{dw}_2=\frac{v_1-v_2}{\sqrt2},\qquad\mathrm{dw}_3=\frac12-v_1 .
\]
All four numbers are positive.
The formula for the reflection gives $\mathrm{st}_iv_0-v_0=-2\,\mathrm{dw}_i\cdot\hat n_i$.
Hence
\[
\mathrm{geo}\left(x^{(i)}\right)-\mathrm{geo}(x)=\mathrm{gt}\,\mathrm{st}_iv_0-\mathrm{gt}\,v_0=g\left(\mathrm{st}_iv_0-v_0\right)=-2\,\mathrm{dw}_i\cdot g\hat n_i,\qquad\left|\mathrm{geo}\left(x^{(i)}\right)-\mathrm{geo}(x)\right|_2=2\,\mathrm{dw}_i,
\]
because $g$ is orthogonal and $|\hat n_i|_2=1$, and
\[
\frac{\mathrm{geo}\left(x^{(i)}\right)-\mathrm{geo}(x)}{\left|\mathrm{geo}\left(x^{(i)}\right)-\mathrm{geo}(x)\right|_2}=-g\hat n_i=\hat d_i(x).
\]

\emph{The modulus in \cite[Definitions~2.4 and~2.10]{QEDII}.}
In \cite[Definition~2.4]{QEDII} the direction of the bond of type $i$ at $x$ is $\hat d_i(x)=\mathrm{dlt}_i(x)/\left|\mathrm{dlt}_i(x)\right|$, and in \cite[Definition~2.10]{QEDII} the direction of a bond is $\hat d=(x'-x)/|x'-x|$: the modulus carries no index.
In \cite[Definitions~2.2, 2.4 and~2.10]{QEDII} the moduli without index of vectors of $\mathbb{R}^3$ are $|n_\ell|$, $\left|x^{(i)}-x\right|=\left|\mathrm{dlt}_i(x)\right|$ and $|x'-x|$; the other bars without index in \cite[Definitions~2.3 and~2.10]{QEDII} denote the areas $\mathrm{Ar}_i$ of the faces $H_i\cap\overline{\mathrm{Chm}_{0}}$, the volume of $\mathrm{Chm}_{0}$ and, in the form $\#|\cdot|$, cardinalities.
Each of these three moduli is the Euclidean norm.
Indeed, \cite[Definition~2.2]{QEDII} has $|n_\ell|^2\in\{1,2\}$ and calls $n_\ell/|n_\ell|$ a unit normal, whereas for the diagonal families $n_\ell=e_\mu\mp e_\nu$ one has $|n_\ell|_2^2=2$ and $|n_\ell|_\infty^2=1$.
Further, \cite[Lemma~3.3]{QEDII} states $|\mathrm{dlt}_i|=2d_i$ with $d_i=\operatorname{dist}(v_0,H_i)$,
\[
d_0=v_3,\qquad d_1=\frac{v_2-v_3}{\sqrt2},\qquad d_2=\frac{v_1-v_2}{\sqrt2},\qquad d_3=\frac12-v_1,
\]
and $\hat d_i(x)=-\mathrm{gb}\,\hat n_i$.
Here $d_i=\mathrm{dw}_i$, and $|\mathrm{dlt}_i|_2=2\,\mathrm{dw}_i$ by the computation at the beginning of this subsection.
On the other hand, for $i=1,2$ the vector $g\hat n_i$ has two entries of modulus $1/\sqrt2$ and one entry $0$, so that
\[
\left|\mathrm{dlt}_i\right|_\infty=\frac{2\,\mathrm{dw}_i}{\sqrt2}=\sqrt2\,\mathrm{dw}_i\neq2\,\mathrm{dw}_i,\qquad\frac{\mathrm{dlt}_i}{\left|\mathrm{dlt}_i\right|_\infty}=-\sqrt2\,g\hat n_i .
\]
Thus the two identities of \cite[Lemma~3.3]{QEDII} hold for $|\cdot|_2$ and not for $|\cdot|_\infty$.
For the hopping matrix $h(\cdot,\cdot)$ of \cite[Definition~2.10]{QEDII}, \cite[Remark~3.5]{QEDII} gives
\[
h\left(x,x^{(i)}\right)=\frac{w_i}{2}\cdot\left(\rho+\gamma\cdot\left(\mathrm{gb}\,\hat n_i\right)\right),
\]
that is, $\hat d=-\mathrm{gb}\,\hat n_i$ on the bond $\left\{x,x^{(i)}\right\}$.
So the quotient $\left(\mathrm{geo}\left(x^{(i)}\right)-\mathrm{geo}(x)\right)/\left|\mathrm{geo}\left(x^{(i)}\right)-\mathrm{geo}(x)\right|_2$ computed above is the vector $\hat d_i$ of \cite[Definition~2.4]{QEDII} at the site $\mathrm{geo}(x)$, and it is the vector $\hat d$ of \cite[Definition~2.10]{QEDII} on the bond $\left\{\mathrm{geo}(x),\mathrm{geo}\left(x^{(i)}\right)\right\}$.

\subsubsection*{3.5. The last sentence of (a)}

By 3.1--3.4, under the bijection $\mathrm{geo}$ the set $X$ corresponds to $\mathbb{Z}^3\times O_{\mathrm{cub}}$, the map ``neighbour across the wall of type $i$'' to the map $x\mapsto x^{(i)}$ of (I.1.2), the bond $\left\{\mathrm{geo}\,x,\mathrm{geo}\,x^{(i)}\right\}$ of type $i$ to the pair $\left\{x,x^{(i)}\right\}$ of type $i$, the map $\pi$ to $(c,g)\mapsto c$, the vector $\hat d_i$ to $-g\hat n_i$, and the action of $W_{\mathrm{af}}$ on $X$ to the action $g_{\mathrm{w}}(c,g)=(rc+a,rg)$.
None of these expressions contains $v_0$.
A bond has only one type: $x^{(0)}$, $x^{(1)}$, $x^{(2)}$ are pairwise different, since $s_0$, $s_1$, $s_2$ are; $x^{(3)}$ lies over $c+ge_1\neq c$; and, seen from its other end, the bond $\left\{x,x^{(i)}\right\}$ is $\left\{x^{(i)},\left(x^{(i)}\right)^{(i)}\right\}$ by (46.4).
The formulas of 3.3 and 3.4 are those of the table in \cite[Lemma~3.3]{QEDII} (where $d_i=\mathrm{dw}_i$ and $\mathrm{dlt}_i(x)=-2d_i\cdot\mathrm{gb}\,\hat n_i$) and of \cite[Remark~3.5]{QEDII}.
This proves Proposition~46\,(a).

\subsection*{4. Proof of (b)}

\subsubsection*{4.1. The identity (I.m1)}

By (I.1.1),
\[
\sum_iM_i(x)=\frac{\rho}{2}\cdot\sum_iw_i+\frac12\,\gamma\cdot\left(g\sum_iw_i\cdot\hat n_i\right)=\frac{\rho}{2}\cdot\mathrm{sw}.
\]
Since $M_{\mathrm{ch}}=\boldsymbol{\mu}+\frac{\rho}{2}\cdot\mathrm{sw}$, it follows that
\begin{equation}
M_{\mathrm{ch}}-\sum_{i=0}^{3}M_i(x)=\boldsymbol{\mu}\qquad\text{for every }x.
\tag{46.5}
\end{equation}

\subsubsection*{4.2. The operator of \cite[Definition~2.10]{QEDII}}

Let $\mathcal{D}_A$ and $\mathcal{U}_A$ denote the operator $D_A$ and the parallel transport $U_A$ of \cite[Definition~2.10]{QEDII}, formed with the Wilson parameter $\rho$, the weights $w$, the mass parameter $M_{\mathrm{ch}}$ in the place of the parameter $M$ of that definition, and the charge $q$.
For a function $\mathrm{fX}\colon X\to\mathbb{C}^4$ and a site $\mathrm{xa}$ of $X$ with the neighbours $\mathrm{xa}^{(i)}$ of \cite[Definition~2.4]{QEDII} one has
\[
\left(\mathcal{D}_A\mathrm{fX}\right)(\mathrm{xa})=M_{\mathrm{ch}}\cdot\mathrm{fX}(\mathrm{xa})-\sum_{i=0}^{3}h\left(\mathrm{xa},\mathrm{xa}^{(i)}\right)\cdot\mathcal{U}_A\left(\mathrm{xa},\mathrm{xa}^{(i)}\right)\cdot\mathrm{fX}\left(\mathrm{xa}^{(i)}\right),
\]
where, for a bond $\{\mathrm{xa},\mathrm{xb}\}$ of type $i$,
\[
h(\mathrm{xa},\mathrm{xb})=\frac{w_i}{2}\cdot\left(\rho-\gamma\cdot\hat d\right),\qquad\hat d=\frac{\mathrm{xb}-\mathrm{xa}}{\left|\mathrm{xb}-\mathrm{xa}\right|_2},\qquad\mathcal{U}_A(\mathrm{xa},\mathrm{xb})=\exp\left(iq\cdot A\left(\pi\,\mathrm{xa},\pi\,\mathrm{xb}\right)\right).
\]
By 3.4 the modulus $|x'-x|$ of \cite[Definition~2.10]{QEDII} is the Euclidean norm, which is written $|\cdot|_2$ here.
The standard weights of \cite[Definition~2.10]{QEDII}, $w_i=48\,\mathrm{Ar}_i$ with the areas $\mathrm{Ar}=\left(1,\sqrt2,\sqrt2,1\right)/8$ of \cite[Definition~2.3]{QEDII}, are $\left(6,6\sqrt2,6\sqrt2,6\right)$.

Let $\mathrm{xa}=\mathrm{geo}(x)$ with $x=(c,g)$, and let $f:=\mathrm{fX}\circ\mathrm{geo}$.
By (a), $\mathrm{xa}^{(i)}=\mathrm{geo}\left(x^{(i)}\right)$ and $\hat d=\hat d_i(x)$, whence
\[
h\left(\mathrm{xa},\mathrm{xa}^{(i)}\right)=\frac{w_i}{2}\cdot\left(\rho-\gamma\cdot\hat d_i(x)\right)=M_i(x);
\]
moreover $\pi\,\mathrm{xa}=c$, $\pi\,\mathrm{xa}^{(i)}=c$ for $i\le2$, and $\pi\,\mathrm{xa}^{(3)}=c+ge_1$.
Therefore $\mathcal{U}_A\left(\mathrm{xa},\mathrm{xa}^{(i)}\right)=\exp\left(iq\cdot A(c,c)\right)=1$ for $i\le2$, because $A(c,c)=0$, and $\mathcal{U}_A\left(\mathrm{xa},\mathrm{xa}^{(3)}\right)=\exp\left(iq\cdot A\left(c,c+ge_1\right)\right)$; these are the numbers $U_A\left(x,x^{(i)}\right)$ of (46.2).
Hence
\begin{align*}
\left(\mathcal{D}_A\mathrm{fX}\right)(\mathrm{geo}\,x)
&=M_{\mathrm{ch}}\cdot f(x)-\sum_iM_i(x)\cdot U_A\left(x,x^{(i)}\right)\cdot f\left(x^{(i)}\right)\qquad\text{[the form (I.2.5)]}\\
&=\left[M_{\mathrm{ch}}-\sum_iM_i(x)\right]\cdot f(x)-\sum_iM_i(x)\cdot\left[U_A\left(x,x^{(i)}\right)\cdot f\left(x^{(i)}\right)-f(x)\right]\\
&=\boldsymbol{\mu}\cdot f(x)+\left(D_Af\right)(x)\qquad\text{by (46.5) and (46.3)}.
\end{align*}

\subsubsection*{4.3. The operator $D_A+\boldsymbol{\mu}$ written with covariant differences}

On the sites $(c,g)\in\mathbb{Z}^3\times O_{\mathrm{cub}}$ let $n_i:=\hat n_i$ $(i=0,\dots,3)$; the reflection $s_i\in O_{\mathrm{cub}}$ in the plane orthogonal to $n_i$ is then $1-2\,\hat n_i\hat n_i^{T}$, as in 2.2.
Let $x^{(i)}$ and $\hat d_i(x)$ be the sites and the vectors of the display (I.1.2), let $w=\left(6,6\sqrt2,6\sqrt2,6\right)$ be the standard weights, and let the phases be $U\left(x,x^{(3)}\right):=\exp\left(iq\cdot A\left(c,c+ge_1\right)\right)$ and $U\left(x,x^{(i)}\right):=1$ for $i\le2$.
The covariant differences are
\[
\left(\nabla_if\right)(x):=U\left(x,x^{(i)}\right)\cdot f\left(x^{(i)}\right)-f(x),
\]
and, for three matrices $\gamma_1$, $\gamma_2$, $\gamma_3$ acting on $\mathbb{C}^4$ and a number $m$, the operator with covariant differences is
\[
D^{\nabla}_A:=m-\sum_{i=0}^{3}\frac{w_i}{2}\cdot\left(\rho-\gamma\cdot\hat d_i(x)\right)\cdot\nabla_i .
\]
With the matrices $\gamma_\mu$ of Section~I.2.1 and with $m=\boldsymbol{\mu}$, (46.2) and (46.3) give
\[
\left(D^{\nabla}_Af\right)(x)=\boldsymbol{\mu}\cdot f(x)-\sum_iM_i(x)\cdot\left[U_A\left(x,x^{(i)}\right)\cdot f\left(x^{(i)}\right)-f(x)\right]=\boldsymbol{\mu}\cdot f(x)+\left(D_Af\right)(x),
\]
that is, $D^{\nabla}_A=D_A+\boldsymbol{\mu}$.

\subsubsection*{4.4. The operator $D^{\mathrm{ch}}[B]$}

Let $\gamma_\mu=\operatorname{diag}\left(\sigma_\mu,-\sigma_\mu\right)$ be the matrices of Section~I.2.1, let $n_i=\hat n_i$ for $i\le2$, let $s_0$, $s_1$, $s_2$ be the reflections in the planes orthogonal to $n_0$, $n_1$, $n_2$ and $s_3$ the reflection $e_1\mapsto-e_1$, and let $w$ be the standard weights.
For a function $B$ on the oriented nearest-neighbour bonds of $\mathbb{Z}^3$ with $B(c',c)=-B(c,c')$, the operator $D^{\mathrm{ch}}[B]$ on the functions $f\colon\mathbb{Z}^3\times O_{\mathrm{cub}}\to\mathbb{C}^4$ is defined by
\begin{multline*}
\left(D^{\mathrm{ch}}[B]f\right)(c,g)=\sum_{i=0}^{2}\frac{w_i}{2}\cdot\left(\rho+\gamma\cdot\left(gn_i\right)\right)\cdot\left(f(c,g)-f\left(c,gs_i\right)\right)\\
+\frac{w_3}{2}\cdot\left(\rho-\gamma\cdot\left(ge_1\right)\right)\cdot\left(f(c,g)-\exp\left(iB\left(c,c+ge_1\right)\right)\cdot f\left(c+ge_1,gs_3\right)\right).
\end{multline*}
With $n_3:=-e_1$ the four terms take one form: $D^{\mathrm{ch}}[B]$ is the sum over $i=0,\dots,3$ of $\frac{w_i}{2}\cdot\left(\rho+\gamma\cdot\left(gn_i\right)\right)$, applied to $f(c,g)-f\left(c,gs_i\right)$ for $i\le2$ and to $f(c,g)-\exp\left(iB\left(c,c+ge_1\right)\right)\cdot f\left(c+ge_1,gs_3\right)$ for $i=3$.

For $i\le2$ one has $\frac{w_i}{2}\cdot\left(\rho+\gamma\cdot\left(g\hat n_i\right)\right)=M_i(x)$ and $f\left(c,gs_i\right)=f\left(x^{(i)}\right)=U_A\left(x,x^{(i)}\right)\cdot f\left(x^{(i)}\right)$.
In the last term, $\hat n_3=-e_1$ gives
\[
\frac{w_3}{2}\cdot\left(\rho-\gamma\cdot\left(ge_1\right)\right)=\frac{w_3}{2}\cdot\left(\rho+\gamma\cdot\left(g\hat n_3\right)\right)=M_3(x);
\]
further $f\left(c+ge_1,gs_3\right)=f\left(x^{(3)}\right)$, and for $B=qA$ one has $\exp\left(iB\left(c,c+ge_1\right)\right)=U_A\left(x,x^{(3)}\right)$.
Therefore
\begin{equation}
\left(D^{\mathrm{ch}}[qA]f\right)(x)=\sum_{i=0}^{3}M_i(x)\cdot\left(f(x)-U_A\left(x,x^{(i)}\right)\cdot f\left(x^{(i)}\right)\right)=\left(D_Af\right)(x),\qquad D^{\mathrm{ch}}[qA]+\boldsymbol{\mu}=D_A+\boldsymbol{\mu}.
\tag{46.6}
\end{equation}

\subsubsection*{4.5. Level $k$}

Let $k\ge0$, $\eta_k=L^{-k}$ and $e_k>0$, let $A$ be a complex bond function on $\mathbb{L}_k=\eta_k\mathbb{Z}^3$, and let $f\colon\mathbb{L}^{\mathrm{f}}_k\to\mathbb{C}^4$.
For $c\in\mathbb{L}_k$ put $\mathrm{cq}:=c/\eta_k\in\mathbb{Z}^3$, and let
\[
\mathrm{fq}(\mathrm{cq},g):=f\left(\eta_k\cdot\mathrm{cq},g\right),\qquad B\left(\mathrm{cq},\mathrm{cq}'\right):=e_k\cdot\eta_k\cdot A\left(\eta_k\cdot\mathrm{cq},\eta_k\cdot\mathrm{cq}'\right)
\]
for nearest neighbours $\mathrm{cq}$, $\mathrm{cq}'$ of $\mathbb{Z}^3$.
At level $k$ the display (I.1.2) gives $x^{(i)}=\left(c,gs_i\right)$ for $i\le2$ and $x^{(3)}=\left(c+\eta_k\cdot ge_1,gs_3\right)$; hence $f\left(x^{(i)}\right)=\mathrm{fq}\left(\mathrm{cq},gs_i\right)$ for $i\le2$ and $f\left(x^{(3)}\right)=\mathrm{fq}\left(\mathrm{cq}+ge_1,gs_3\right)$.
By (I.2.3), $U_A\left(x,x^{(i)}\right)=1$ for $i\le2$, and
\[
U_A\left(x,x^{(3)}\right)=\exp\left(i\,e_k\cdot\eta_k\cdot A\left(c,c+\eta_k\cdot ge_1\right)\right)=\exp\left(iB\left(\mathrm{cq},\mathrm{cq}+ge_1\right)\right).
\]
The matrix $M_i(x)=\frac{w_i}{2}\cdot\left(\rho+\gamma\cdot\left(g\hat n_i\right)\right)$ depends on the label only.
By the definition of $D^{\mathrm{ch}}$ in 4.4 and by $\hat n_3=-e_1$,
\begin{align*}
\left(D^{\mathrm{ch}}[B]\,\mathrm{fq}\right)(\mathrm{cq},g)
&=\sum_{i=0}^{2}M_i(x)\cdot\left(\mathrm{fq}(\mathrm{cq},g)-\mathrm{fq}\left(\mathrm{cq},gs_i\right)\right)\\
&\qquad+M_3(x)\cdot\left(\mathrm{fq}(\mathrm{cq},g)-\exp\left(iB\left(\mathrm{cq},\mathrm{cq}+ge_1\right)\right)\cdot\mathrm{fq}\left(\mathrm{cq}+ge_1,gs_3\right)\right)\\
&=\sum_{i=0}^{3}M_i(x)\cdot\left(f(x)-U_A\left(x,x^{(i)}\right)\cdot f\left(x^{(i)}\right)\right),
\end{align*}
and (I.2.3) gives
\[
\left(D_Af\right)(c,g)=-\eta_k^{-1}\sum_{i=0}^{3}M_i(x)\cdot\left[U_A\left(x,x^{(i)}\right)\cdot f\left(x^{(i)}\right)-f(x)\right]=\eta_k^{-1}\cdot\left(D^{\mathrm{ch}}[B]\,\mathrm{fq}\right)(\mathrm{cq},g).
\]

At level $k$ the operator (I.2.3) also has the form of 4.3.
On the functions on $\mathbb{L}_k\times O_{\mathrm{cub}}$ consider the operator $\eta_k^{-1}\cdot D^{\nabla}_A$, where $D^{\nabla}_A$ is taken with $m=0$ and with the phase $U\left(x,x^{(3)}\right)=\exp\left(i\,e_k\cdot\eta_k\cdot A\left(c,c+\eta_k\cdot ge_1\right)\right)$, and is read in the variable $c/\eta_k$.
Applied to $\mathrm{fq}$ at the site $(\mathrm{cq},g)$ of $\mathbb{Z}^3\times O_{\mathrm{cub}}$, with the phases $1$ for $i\le2$ and $\exp\left(iB\left(\mathrm{cq},\mathrm{cq}+ge_1\right)\right)$ for $i=3$, the operator $D^{\nabla}_A$ with $m=0$ gives
\[
-\sum_{i=0}^{2}M_i(x)\cdot\left(\mathrm{fq}\left(\mathrm{cq},gs_i\right)-\mathrm{fq}(\mathrm{cq},g)\right)-M_3(x)\cdot\left(\exp\left(iB\left(\mathrm{cq},\mathrm{cq}+ge_1\right)\right)\cdot\mathrm{fq}\left(\mathrm{cq}+ge_1,gs_3\right)-\mathrm{fq}(\mathrm{cq},g)\right),
\]
which is the first right-hand side in the computation of $\left(D^{\mathrm{ch}}[B]\,\mathrm{fq}\right)(\mathrm{cq},g)$ above.
Multiplied by $\eta_k^{-1}$, it is therefore $\left(D_Af\right)(c,g)$.
Hence the operator $\eta_k^{-1}\cdot D^{\nabla}_A$, formed in this way, is the operator (I.2.3), and it coincides with $\eta_k^{-1}\cdot D^{\mathrm{ch}}[B]$ in the variable $c/\eta_k$, where $B\left(c/\eta_k,c'/\eta_k\right)=e_k\cdot\eta_k\cdot A(c,c')$.
This proves Proposition~46\,(b).

\subsection*{5. Proof of (c)}

\subsubsection*{5.1. Rules}

For $2\times2$ matrices $\mathrm{ua}$ and $\mathrm{ub}$, the Kronecker product $\mathrm{ua}\otimes\mathrm{ub}$ is the $4\times4$ matrix whose $2\times2$ blocks are the matrices $\mathrm{ua}_{ij}\cdot\mathrm{ub}$, $i,j=1,2$.
Multiplication of block matrices gives
\begin{equation}
\begin{aligned}
&\left(\mathrm{ua}\otimes\mathrm{ub}\right)\cdot\left(\mathrm{ua}'\otimes\mathrm{ub}'\right)=\left(\mathrm{ua}\cdot\mathrm{ua}'\right)\otimes\left(\mathrm{ub}\cdot\mathrm{ub}'\right),\qquad\left(\mathrm{ua}\otimes\mathrm{ub}\right)^T=\mathrm{ua}^T\otimes\mathrm{ub}^T,\\
&\overline{\mathrm{ua}\otimes\mathrm{ub}}=\overline{\mathrm{ua}}\otimes\overline{\mathrm{ub}},\qquad\left(\mathrm{ua}\otimes\mathrm{ub}\right)^*=\mathrm{ua}^*\otimes\mathrm{ub}^*.
\end{aligned}
\tag{K}
\end{equation}
The Pauli matrices are
\[
\sigma_1=\left(\begin{array}{cc}0&1\\1&0\end{array}\right),\qquad
\sigma_2=\left(\begin{array}{cc}0&-i\\i&0\end{array}\right),\qquad
\sigma_3=\left(\begin{array}{cc}1&0\\0&-1\end{array}\right).
\]
Multiplying these matrices one finds, for $j,l\in\{1,2,3\}$,
\begin{align}
&\sigma_j^*=\sigma_j,\qquad\sigma_j^2=1_2,\qquad\sigma_j\cdot\sigma_l=-\sigma_l\cdot\sigma_j\quad(j\neq l),\notag\\
&\sigma_1\cdot\sigma_2=i\sigma_3,\qquad\sigma_2\cdot\sigma_3=i\sigma_1,\qquad\sigma_3\cdot\sigma_1=i\sigma_2;\tag{Pa1}\\
&\sigma_j^T=\overline{\sigma_j}=\mathrm{tp}_j\cdot\sigma_j,\qquad\text{where }\mathrm{tp}_1=\mathrm{tp}_3=1,\ \mathrm{tp}_2=-1;\tag{Pa2}\\
&\sigma_2\cdot\sigma_j\cdot\sigma_2=-\mathrm{tp}_j\cdot\sigma_j=-\sigma_j^T=-\overline{\sigma_j}.\tag{Pa3}
\end{align}
For (Pa3): if $j=2$, the left side is $\sigma_2^3=\sigma_2$;
if $j=1$ or $j=3$, it is $\sigma_2\cdot\sigma_j\cdot\sigma_2=-\sigma_j\cdot\sigma_2^2=-\sigma_j$.

The eight matrices of Section~I.2.1 are
\begin{gather*}
\gamma_\mu=\sigma_3\otimes\sigma_\mu=\operatorname{diag}\left(\sigma_\mu,-\sigma_\mu\right)\quad(\mu=1,2,3),\qquad
\gamma_5=\sigma_1\otimes1_2=\left(\begin{array}{cc}0&1_2\\1_2&0\end{array}\right),\qquad
\Sigma_3=\sigma_3\otimes1_2=\operatorname{diag}(1,1,-1,-1),\\
\mathsf{C}=i\gamma_2=\sigma_3\otimes\left(i\sigma_2\right),\qquad
\varkappa=\gamma_5\cdot\mathsf{C},\qquad
\mathrm{Cc}=1_2\otimes\sigma_2=\operatorname{diag}\left(\sigma_2,\sigma_2\right).
\end{gather*}
Since $\sigma_1\cdot\sigma_3=-\sigma_3\cdot\sigma_1=-i\sigma_2$, one has
\[
\varkappa=\left(\sigma_1\otimes1_2\right)\cdot\left(i\sigma_3\otimes\sigma_2\right)=i\left(\sigma_1\cdot\sigma_3\right)\otimes\sigma_2=i\left(-i\sigma_2\right)\otimes\sigma_2=\sigma_2\otimes\sigma_2.
\]
\emph{Entries.}
As $i\sigma_2=\left(\begin{array}{cc}0&1\\-1&0\end{array}\right)$,
\[
\mathsf{C}=\operatorname{diag}\left(i\sigma_2,-i\sigma_2\right)=\left(\begin{array}{cccc}0&1&0&0\\-1&0&0&0\\0&0&0&-1\\0&0&1&0\end{array}\right).
\]
Further,
\[
\sigma_2\otimes\sigma_2=\left(\begin{array}{cc}0&-i\sigma_2\\i\sigma_2&0\end{array}\right),\qquad
-i\sigma_2=\left(\begin{array}{cc}0&-1\\1&0\end{array}\right).
\]
So the rows of $\varkappa$ are $(0,0,0,-1)$, $(0,0,1,0)$, $(0,1,0,0)$ and $(-1,0,0,0)$:
$\varkappa$ is the antidiagonal matrix whose entries, from the upper right to the lower left corner, are $-1,1,1,-1$.

\subsubsection*{5.2. (c1)--(c6)}

\begin{itemize}
\item[(c1)]
By (K) and (Pa1), the matrices $\gamma_\mu$, $\gamma_5$ and $\Sigma_3$ are Hermitian.
Moreover,
\begin{gather*}
\gamma_\mu\cdot\gamma_\nu+\gamma_\nu\cdot\gamma_\mu=\sigma_3^2\otimes\left(\sigma_\mu\cdot\sigma_\nu+\sigma_\nu\cdot\sigma_\mu\right)=1_2\otimes\left(2\delta_{\mu\nu}\cdot1_2\right),\\
\gamma_5\cdot\gamma_\mu=\left(\sigma_1\cdot\sigma_3\right)\otimes\sigma_\mu=-\left(\sigma_3\cdot\sigma_1\right)\otimes\sigma_\mu=-\gamma_\mu\cdot\gamma_5,\qquad
\gamma_5^2=\sigma_1^2\otimes1_2=1,\qquad
\Sigma_3^2=\sigma_3^2\otimes1_2=1.
\end{gather*}
For real $v$ it follows that
\[
\left(\gamma\cdot v\right)^2=\frac12\sum_{\mu,\nu}v_\mu\cdot v_\nu\cdot\left(\gamma_\mu\cdot\gamma_\nu+\gamma_\nu\cdot\gamma_\mu\right)=|v|_2^2.
\]
\item[(c2)]
One has
\begin{gather*}
\gamma_1\cdot\gamma_2\cdot\gamma_3=\sigma_3^3\otimes\left(\sigma_1\cdot\sigma_2\cdot\sigma_3\right)=\sigma_3\otimes\left(\left(i\sigma_3\right)\cdot\sigma_3\right)=i\sigma_3\otimes1_2=i\Sigma_3,\qquad
\Sigma_3\cdot\gamma_\mu=\sigma_3^2\otimes\sigma_\mu=\gamma_\mu\cdot\Sigma_3,\\
\Sigma_3\cdot\gamma_5=\left(\sigma_3\cdot\sigma_1\right)\otimes1_2=-\left(\sigma_1\cdot\sigma_3\right)\otimes1_2=-\gamma_5\cdot\Sigma_3,\qquad
i\gamma_5\cdot\Sigma_3=i\left(\sigma_1\cdot\sigma_3\right)\otimes1_2=i\left(-i\sigma_2\right)\otimes1_2=\sigma_2\otimes1_2.
\end{gather*}
\item[(c3)]
The entries of $\mathsf{C}$ show that $\mathsf{C}$ is real.
Next, $\mathsf{C}^T=\sigma_3^T\otimes\left(i\sigma_2\right)^T=\sigma_3\otimes\left(-i\sigma_2\right)=-\mathsf{C}$, and $\mathsf{C}^2=i^2\gamma_2^2=-1$, so that $\mathsf{C}^{-1}=-\mathsf{C}$.
By (Pa3) and $\sigma_3^T=\sigma_3$,
\[
\mathsf{C}^{-1}\cdot\gamma_\mu\cdot\mathsf{C}=\left(-i\gamma_2\right)\cdot\gamma_\mu\cdot\left(i\gamma_2\right)=\gamma_2\cdot\gamma_\mu\cdot\gamma_2=\sigma_3^3\otimes\left(\sigma_2\cdot\sigma_\mu\cdot\sigma_2\right)=\sigma_3\otimes\left(-\sigma_\mu^T\right)=-\left(\sigma_3\otimes\sigma_\mu\right)^T=-\gamma_\mu^T.
\]
Finally, $\mathsf{C}^{-1}\cdot\gamma_5\cdot\mathsf{C}=\gamma_2\cdot\gamma_5\cdot\gamma_2=-\gamma_5\cdot\gamma_2^2=-\gamma_5=-\gamma_5^T$, since $\gamma_5=\sigma_1\otimes1_2$ is symmetric.
\item[(c4)]
The entries of $\varkappa$ show that $\varkappa$ is real.
Next, $\varkappa^T=\sigma_2^T\otimes\sigma_2^T=\left(-\sigma_2\right)\otimes\left(-\sigma_2\right)=\varkappa$ and $\varkappa^2=\sigma_2^2\otimes\sigma_2^2=1$.
By (Pa3),
\[
\varkappa^{-1}\cdot\gamma_\mu\cdot\varkappa=\left(\sigma_2\cdot\sigma_3\cdot\sigma_2\right)\otimes\left(\sigma_2\cdot\sigma_\mu\cdot\sigma_2\right)=\left(-\sigma_3\right)\otimes\left(-\overline{\sigma_\mu}\right)=\overline{\sigma_3\otimes\sigma_\mu}=\overline{\gamma_\mu},
\]
and $\varkappa^{-1}\cdot\gamma_5\cdot\varkappa=\left(\sigma_2\cdot\sigma_1\cdot\sigma_2\right)\otimes\sigma_2^2=-\sigma_1\otimes1_2=-\gamma_5=-\overline{\gamma_5}$.
As $\mathsf{C}=i\gamma_2$ anticommutes with $\gamma_5$, one has $\varkappa\cdot\mathsf{C}=\gamma_5\cdot\mathsf{C}^2=-\gamma_5$ and $\mathsf{C}\cdot\varkappa=\mathsf{C}\cdot\gamma_5\cdot\mathsf{C}=-\gamma_5\cdot\mathsf{C}^2=\gamma_5$;
hence $\varkappa\cdot\mathsf{C}=-\mathsf{C}\cdot\varkappa$.
\item[(c5)]
One has $i\Sigma_3\cdot\mathrm{Cc}=i\left(\sigma_3\otimes1_2\right)\cdot\left(1_2\otimes\sigma_2\right)=i\sigma_3\otimes\sigma_2=\mathsf{C}$, and $\mathrm{Cc}^2=1_2\otimes\sigma_2^2=1$, so that $\mathrm{Cc}^{-1}=\mathrm{Cc}$.
By (Pa2) and (Pa3), $\sigma_2\cdot\sigma_\mu^T\cdot\sigma_2=\mathrm{tp}_\mu\cdot\sigma_2\cdot\sigma_\mu\cdot\sigma_2=-\mathrm{tp}_\mu^2\cdot\sigma_\mu=-\sigma_\mu$.
Therefore
\[
\mathrm{Cc}\cdot\gamma_\mu^T\cdot\mathrm{Cc}^{-1}=\sigma_3^T\otimes\left(\sigma_2\cdot\sigma_\mu^T\cdot\sigma_2\right)=-\sigma_3\otimes\sigma_\mu=-\gamma_\mu,
\qquad
\mathrm{Cc}\cdot\gamma_5^T\cdot\mathrm{Cc}^{-1}=\sigma_1^T\otimes\left(\sigma_2\cdot1_2\cdot\sigma_2\right)=\sigma_1\otimes1_2=+\gamma_5.
\]
\item[(c6)]
Each of the matrices $1_2$, $\sigma_1$, $\sigma_2$, $\sigma_3$ has exactly one non-zero entry in each row and in each column, and this entry has modulus one.
The entries of $\mathrm{ua}\otimes\mathrm{ub}$ are the products $\mathrm{ua}_{ij}\cdot\mathrm{ub}_{i'j'}$;
hence $\mathrm{ua}\otimes\mathrm{ub}$ and $i\,\mathrm{ua}\otimes\mathrm{ub}$ have the same property whenever $\mathrm{ua}$ and $\mathrm{ub}$ have it.
Each of the eight matrices is of one of these two forms, with $\mathrm{ua},\mathrm{ub}\in\{1_2,\sigma_1,\sigma_2,\sigma_3\}$.
The matrices $\gamma_3=\sigma_3\otimes\sigma_3$ and $\Sigma_3=\sigma_3\otimes1_2$ are Kronecker products of diagonal matrices, hence diagonal.
\end{itemize}

\subsubsection*{5.3. The rows of the table}

The rows are those of the table of Section~I.2.1.
In \cite{QEDI} and \cite{QEDII} the matrices of these rows are defined as follows.

\paragraph{Row $\gamma_1,\gamma_2,\gamma_3$.}
In \cite[Section~5.2.4]{QEDI}, and again in \cite[Sections~13.4.1 and~14.6.1]{QEDI}, one has $\gamma_\mu=\sigma_3\otimes\sigma_{\mu+1}$ for $\mu=0,1,2$;
these are the matrices $\gamma_{\mu+1}$ of Section~I.2.1.
In \cite[Section~2.2]{QEDI}, in \cite[Section~6.1.2]{QEDI} (the four-component representation) and in \cite[Sections~7.1.2, 9.7.1, 10.1.2 and~21.1.1]{QEDI}, one has $\gamma_\mu=\operatorname{diag}\left(\sigma_\mu,-\sigma_\mu\right)$ for $\mu=1,2,3$.
In \cite[Definition~2.7]{QEDII}, $\gamma_\mu=\operatorname{diag}\left(\sigma_\mu,-\sigma_\mu\right)$.

\paragraph{Row $\gamma_5$.}
In \cite[Section~5.2.4]{QEDI}, and again in \cite[Sections~13.4.1 and~14.6.1]{QEDI}, the matrix $\gamma_3$ of that numbering is $\sigma_1\otimes1$.
In \cite[Section~21.1.1]{QEDI},
\[
\hat{\Gamma}=\left(\begin{array}{cc}0&1_2\\1_2&0\end{array}\right)=\sigma_1\otimes1_2.
\]
In \cite[Section~7.1.2]{QEDI}, $\gamma_4=\tau_1\otimes1$, where $\tau_1$ is the first Pauli matrix acting on the two blocks.
In \cite[Definition~2.7]{QEDII},
\[
\gamma_5=\left(\begin{array}{cc}0&1\\1&0\end{array}\right)\otimes1_2.
\]
All of these matrices are $\sigma_1\otimes1_2$.

\paragraph{Row $\Sigma_3$.}
Formed with the matrices of \cite[Section~5.2.4]{QEDI}, the matrix $-i\gamma_0\cdot\gamma_1\cdot\gamma_2$ is, in the numbering of Section~I.2.1, the matrix $-i\gamma_1\cdot\gamma_2\cdot\gamma_3$, and $-i\gamma_1\cdot\gamma_2\cdot\gamma_3=-i\left(i\Sigma_3\right)=\Sigma_3$ by (c2).
The other entries of the row name no matrix.
The two blocks of $\gamma$ in \cite{QEDII}, that is, the two summands of the decomposition $\mathbb{C}^4=\mathbb{C}^2\oplus\mathbb{C}^2$ of \cite[Definition~2.7]{QEDII}, are the two eigenspaces of $\Sigma_3=\operatorname{diag}(1,1,-1,-1)$;
on them $\gamma_\mu=\operatorname{diag}\left(\sigma_\mu,-\sigma_\mu\right)$ acts as $\sigma_\mu$ and as $-\sigma_\mu$, respectively.

\paragraph{Row $\mathsf{C}$.}
In \cite[Section~5.2.4]{QEDI}, and again in \cite[Sections~13.4.1 and~14.6.1]{QEDI}, $\mathsf{C}=i\gamma_1$ in the numbering $0,1,2$ of the directions;
this is $i\gamma_2$ in the numbering of Section~I.2.1.

\paragraph{Row $\varkappa$.}
In \cite[Section~5.2.4]{QEDI}, and again in \cite[Section~13.4.1]{QEDI}, $\varkappa=\gamma_3\cdot\mathsf{C}$, where $\gamma_3=\sigma_1\otimes1$ is the matrix of that numbering;
thus $\varkappa=\gamma_5\cdot\mathsf{C}$.

\paragraph{Row $\mathrm{Cc}$.}
In the proof of \cite[Theorem~21.1(d3)]{QEDI} in \cite[Section~21.3.2]{QEDI}, on which the identity $\mathfrak{w}(-A)=\mathfrak{w}(A)$ of \cite[Section~10.1.2]{QEDI} rests, and in \cite[Definition~2.7]{QEDII}:
\[
\mathrm{Cc}=\operatorname{diag}\left(\sigma_2,\sigma_2\right)=1_2\otimes\sigma_2.
\]

\paragraph{Repetitions inside statements and proofs.}
At each of the following places \cite{QEDI} repeats one of the definitions above.
For $\gamma$:
the proof of \cite[Lemma~6.6(v)]{QEDI}, with $\gamma_\mu=\operatorname{diag}\left(\sigma_\mu,-\sigma_\mu\right)$ in the four-component representation;
the proof of \cite[Lemma~7.7(b)]{QEDI}, with $\gamma_\nu=\tau_3\otimes\sigma_\nu$;
the proof of \cite[Lemma~7.7(c)]{QEDI}, with $\gamma_\nu=\operatorname{diag}\left(\sigma_\nu,-\sigma_\nu\right)$;
the proof of the lemma on the Clifford relations in \cite[Section~17.1.4]{QEDI}, with $\gamma_\nu=\sigma_3\otimes\sigma_{\nu+1}$ for $\nu\le2$, as in \cite[Section~5.2.4]{QEDI};
the proofs of \cite[Lemma~19.23]{QEDI} and of \cite[Lemma~20.4(a)]{QEDI}, with $\gamma_\nu=\sigma_3\otimes\sigma_{\nu+1}=\operatorname{diag}\left(\sigma_{\nu+1},-\sigma_{\nu+1}\right)$, as in \cite[Section~5.2.4]{QEDI};
\cite[Theorem~21.1(d2)]{QEDI} and its proof in \cite[Section~21.3.2]{QEDI}, with the decomposition $\mathbb{C}^4=\mathbb{C}^2\oplus\mathbb{C}^2$ in which $\gamma_\mu=\operatorname{diag}\left(\sigma_\mu,-\sigma_\mu\right)$, as in \cite[Section~21.1.1]{QEDI};
the proof of \cite[Lemma~21.15]{QEDI}, with $\gamma\cdot w=\operatorname{diag}\left(\sigma\cdot w,-\sigma\cdot w\right)$ for $w\in\mathbb{R}^3$, where $\sigma\cdot w=\sum_\mu w_\mu\cdot\sigma_\mu$, as in \cite[Section~21.1.1]{QEDI}.
For $\gamma_5$:
the proof of \cite[Lemma~7.7(b)]{QEDI}, with $\gamma_4=\tau_1\otimes1$;
the proof of \cite[Lemma~9.20(a)]{QEDI}, with $\hat{\Gamma}=\left(\begin{array}{cc}0&1\\1&0\end{array}\right)$ in blocks of size two;
the proof of \cite[Lemma~16.30]{QEDI} and the proof of the lemma on the Clifford relations in \cite[Section~17.1.4]{QEDI}, with $\gamma_3=\sigma_1\otimes1$, as in \cite[Section~5.2.4]{QEDI}.
For $\mathsf{C}$: the proofs of \cite[Proposition~16.21]{QEDI} and of \cite[Proposition~16.27]{QEDI}, and the proof of the lemma on the Clifford relations in \cite[Section~17.1.4]{QEDI}, with $\mathsf{C}=i\gamma_1$, as in \cite[Section~5.2.4]{QEDI}.
For $\varkappa$: the proof of the lemma on the Clifford relations in \cite[Section~17.1.4]{QEDI}, with $\varkappa=\gamma_3\cdot\mathsf{C}$, as in \cite[Section~5.2.4]{QEDI}.
For $\mathrm{Cc}$: the proof of \cite[Lemma~7.7(c)]{QEDI}.
In that proof the letter $\mathsf{C}$ denotes the operator that acts at every site by the matrix $\operatorname{diag}\left(\sigma_2,\sigma_2\right)=\mathrm{Cc}$.
This is not the matrix $\mathsf{C}=i\gamma_1$ of \cite[Section~5.2.4]{QEDI}, which is $i\gamma_2$ in the numbering of Section~I.2.1;
the two differ by the factor $i\Sigma_3$, since $i\Sigma_3\cdot\mathrm{Cc}=\mathsf{C}$ by (c5).
Hence this place belongs to the row $\mathrm{Cc}$.

These are the places of \cite{QEDI} and \cite{QEDII} that belong to the table of Section~I.2.1 and at which the matrices of the lattice model are defined;
inside its proofs, \cite{QEDII} repeats $\gamma_\mu=\operatorname{diag}\left(\sigma_\mu,-\sigma_\mu\right)$ of \cite[Definition~2.7]{QEDII}.
At each of these places the matrix is the matrix in the first column of its row.
This proves the assertion of (c) on the places at which the matrices of the lattice model are defined.

\paragraph{The places at which no matrix is defined.}
The first sentence of \cite[Section~6.1.2]{QEDI} asks for Hermitian matrices $\gamma_\mu$ with $\gamma_\mu\cdot\gamma_\nu+\gamma_\nu\cdot\gamma_\mu=2\delta_{\mu\nu}$.
Convention~3 of \cite[Section~12.1]{QEDI} asks for the same at the size four, and in addition for a Hermitian matrix $\hat{\Gamma}$ with $\hat{\Gamma}^2=1$.
By (c1), the matrices $\gamma_1$, $\gamma_2$, $\gamma_3$ and $\gamma_5$ satisfy these relations.

\cite[Section~11.1.2]{QEDI} asks for $\gamma_\mu=\operatorname{diag}\left(\sigma_\mu,-\sigma_\mu\right)$, $\mu=0,1,2$, where $\sigma_0,\sigma_1,\sigma_2$ are the three Pauli matrices in an order that is not specified.
If the matrix $\sigma_\mu$ of that section is taken to be $\sigma_{\mu+1}$, these are the matrices $\gamma_{\mu+1}$ of Section~I.2.1.

In \cite[Section~6.3.2]{QEDI}, in the paragraph on the spinor basis adapted to the time direction, $\nu$ is the direction fixed there as time and $i$ ranges over the two other directions.
A unitary matrix $w$ is chosen there with $w\cdot\gamma_\nu\cdot w^*=\operatorname{diag}\left(1_{\frac{s}{2}},-1_{\frac{s}{2}}\right)$;
then
\[
w\cdot\gamma_i\cdot w^*=\left(\begin{array}{cc}0&\sigma_i\\\sigma_i^*&0\end{array}\right)
\]
where, for each of the two indices $i$, the letter $\sigma_i$ names the upper right block of $w\cdot\gamma_i\cdot w^*$, a matrix of size $s/2$, and does not denote the Pauli matrix with the same index.
From there on $\gamma_\nu$ and $\gamma_i$ are assumed to have this form.
This change of the basis of $\mathbb{C}^s$ is kept from that paragraph on, through the rest of \cite[Section~6]{QEDI};
so in the proof of \cite[Lemma~6.8]{QEDI} one has $P^{+}_{\nu}=\operatorname{diag}(1,0)$ and $P^{-}_{\nu}=\operatorname{diag}(0,1)$, where $P^{\pm}_{\nu}=\frac12\left(1\pm\gamma_\nu\right)$ are the matrices of the display that defines the Wilson operator in \cite[Section~6.1.3]{QEDI}.
For $s=4$ one has $w\cdot\gamma_\nu\cdot w^*=\operatorname{diag}(1,1,-1,-1)=\Sigma_3$.
This is none of the matrices $\gamma_1,\gamma_2,\gamma_3$:
the matrices $\gamma_1$ and $\gamma_2$ are not diagonal, and $\gamma_3=\sigma_3\otimes\sigma_3=\operatorname{diag}(1,-1,-1,1)$.
Hence, from this paragraph of \cite[Section~6.3.2]{QEDI} on, the letters $\gamma_\nu$, $\gamma_i$ of \cite[Section~6]{QEDI} do not denote the matrices of the row $\gamma_1,\gamma_2,\gamma_3$ of the table.
The triple $\left(w\cdot\gamma_\mu\cdot w^*\right)_{\mu}$ is unitarily conjugate to the triple of \cite[Section~6.1.2]{QEDI}, which in the four-component representation is the triple of Section~I.2.1 (row $\gamma_1,\gamma_2,\gamma_3$ above).

The values that \cite{QEDI} inserts or names are the following.
In \cite[Section~6.1.2]{QEDI}, the third sentence gives the triple above, and the same section states that, for $\mu=1,2,3$, the matrix $\gamma_\mu$ of the four-component representation is the matrix $\gamma_{\mu-1}$ of \cite[Section~5.2.4]{QEDI}, where the directions are numbered $0,1,2$.
In the proof of \cite[Proposition~16.6]{QEDI}, $\hat{\Gamma}=\gamma_3$, and the matrices $\gamma_1,\gamma_2,\gamma_3$ of Convention~3 of \cite[Section~12.1]{QEDI} are taken to be $\gamma_0,\gamma_1,\gamma_2$;
in the proof of \cite[Lemma~16.7]{QEDI}, $\hat{\Gamma}=\gamma_3$;
in both proofs $\gamma_0,\dots,\gamma_3$ are the matrices of \cite[Section~5.2.4]{QEDI}.
In \cite[Section~11.1.2]{QEDI}, after the display that defines $D[B]$, the operator $D[B]$ is identified, after a renumbering of the directions, with the operator of the display in \cite[Section~21.1.4]{QEDI}.
The hop in the direction $\mu$ of the display of \cite[Section~11.1.2]{QEDI} carries $\frac12\left(1-\gamma_\mu\right)$, with the matrix $\gamma_\mu$ of that section;
the hop in the direction $e_\nu$ of the display of \cite[Section~21.1.4]{QEDI} carries $\frac12\left(1-\gamma_\nu\right)$, with $\gamma_\nu=\operatorname{diag}\left(\sigma_\nu,-\sigma_\nu\right)$.
With the numbering $\mu\mapsto\mu+1$ of the directions (Section~I.1.1), the two displays agree hop by hop if and only if the matrix $\sigma_\mu$ of \cite[Section~11.1.2]{QEDI} is $\sigma_{\mu+1}$.
Up to the numbering of the directions, the matrices of \cite[Sections~5.2.4, 6.1.2 and~11.1.2]{QEDI} and those of \cite[Section~21.1.1]{QEDI} are the same matrices.
\qed

\subsection*{6. Proof of (d)}

In \cite[Definition~2.10]{QEDII} the lattice mass is $m:=M-\frac{\rho}{2}\cdot\sum_iw_i$.
For $M=M_{\mathrm{ch}}=\boldsymbol{\mu}+\frac{\rho}{2}\cdot\mathrm{sw}$ and $\sum_iw_i=\mathrm{sw}$ this gives $m=\boldsymbol{\mu}$.
In \cite[Definition~2.36]{QEDII} one has $\eta=L^{-N}$, $q=e\cdot\eta^{\frac12}$, $\boldsymbol{\mu}_N:=\bar{m}\cdot\eta+m_{0,N}$ (written $m_N$ there), $M=\boldsymbol{\mu}_N+\frac{\rho}{2}\cdot\sum_iw_i$, and the lattice mass $m$ of the chamber operator equals $\boldsymbol{\mu}_N$.
In \cite[Section~21.1.4]{QEDI}, in the paragraph on the laws $P^{\mathrm{nat}}_{N,T}$, one has $q=e\cdot\eta^{\frac12}$ and $\boldsymbol{\mu}_N:=\bar{m}\cdot L^{-N}+m_{0,N}$, and the law at the cutoff is formed with $\boldsymbol{\mu}=\boldsymbol{\mu}_N$.
Let $\boldsymbol{\mu}=\bar{m}\cdot\eta+m_{\mathrm{in}}$ with $m_{\mathrm{in}}=m_{0,N}$.
Then $\boldsymbol{\mu}_N=\boldsymbol{\mu}$ in both papers, and the number $M$ of \cite[Definition~2.36]{QEDII} is
\[
\boldsymbol{\mu}+\frac{\rho}{2}\cdot\mathrm{sw}=M_{\mathrm{ch}}=\bar{m}\cdot\eta+m_{\mathrm{in}}+\frac{\rho}{2}\cdot\sum_iw_i.
\]
By (b), the operator of \cite[Definition~2.10]{QEDII} with these parameters is $D_A+\boldsymbol{\mu}$.
\qed

\subsection*{7. Proof of (e)}

\subsubsection*{7.1. The map}

Let $f\in\mathfrak{S}^{\mathrm{ch}}(1;n)$, $z\in\mathbb{Z}^3$ and $v=\mathrm{nz}$.
By (46.4), $(x+v)^{(i)}=x^{(i)}+v$.
Since $M_i(x)$ depends on the label only, $M_i(x+v)=M_i(x)$;
since $A$ is $\mathrm{Ga}_n$-periodic, (46.2) gives $U_A\left(x+v,x^{(i)}+v\right)=U_A\left(x,x^{(i)}\right)$.
Hence, by (46.3),
\[
\left(D_Af\right)(x+v)=-\sum_iM_i(x)\cdot\left[U_A\left(x,x^{(i)}\right)\cdot f\left(x^{(i)}+v\right)-f(x+v)\right]=\boldsymbol{\chi}(z)\cdot\left(D_Af\right)(x).
\]
Thus $D_A$, and therefore $D_A+\boldsymbol{\mu}$, is a linear map of $\mathfrak{S}^{\mathrm{ch}}(1;n)$ into itself.

\subsubsection*{7.2. Dimension, basis, matrix}

By 2.3, the restriction $f\mapsto f|_{\mathrm{Yf}}$ is a linear bijection of $\mathfrak{S}^{\mathrm{ch}}(1;n)$ onto the space of the functions $\mathrm{Yf}\to\mathbb{C}^4$.
Its inverse sends such a function $f$ to the function given by $f\left(c^0+\mathrm{nz},g\right):=\boldsymbol{\chi}(z)\cdot f\left(c^0,g\right)$;
this is well defined because the pair $\left(c^0,z\right)$ is unique, and it is antiperiodic because $\boldsymbol{\chi}(z+z')=\boldsymbol{\chi}(z)\cdot\boldsymbol{\chi}(z')$.
Consequently
\[
\mathrm{dq}:=\dim\mathfrak{S}^{\mathrm{ch}}(1;n)=4\cdot48\cdot|n|=192|n|.
\]
For $x'\in\mathrm{Yf}$ and $\beta'\in\{1,\dots,4\}$ let $\mathrm{bs}_{x',\beta'}$ be the element of $\mathfrak{S}^{\mathrm{ch}}(1;n)$ with
\[
\left(\mathrm{bs}_{x',\beta'}(x)\right)_\beta=\mathbf{1}[x=x']\cdot\mathbf{1}[\beta=\beta']\qquad(x\in\mathrm{Yf}).
\]
These $\mathrm{dq}$ elements form a basis.
The matrix $\mathrm{Dm}$ of a linear map $O$ of $\mathfrak{S}^{\mathrm{ch}}(1;n)$ in this basis is determined by
\begin{equation}
(Of)_\beta(x)=\sum_{x'\in\mathrm{Yf}}\sum_{\beta'=1}^{4}\mathrm{Dm}_{(x,\beta),(x',\beta')}\cdot f_{\beta'}(x')\qquad\left(x\in\mathrm{Yf},\ \beta=1,\dots,4\right).
\tag{46.7}
\end{equation}

\subsubsection*{7.3. The determinant of a linear map; basis, pairing and $v_0$}

Put $\det O:=\det\mathrm{Dm}$.
In another basis the matrix of $O$ is $\mathrm{Bc}^{-1}\cdot\mathrm{Dm}\cdot\mathrm{Bc}$ with an invertible matrix $\mathrm{Bc}$, and $\det\left(\mathrm{Bc}^{-1}\cdot\mathrm{Dm}\cdot\mathrm{Bc}\right)=\det\mathrm{Dm}$;
so the number does not depend on the basis.
Its definition involves no pairing.
Neither the display (46.3), nor the number $\boldsymbol{\mu}$, nor the space $\mathfrak{S}^{\mathrm{ch}}(1;n)$ involves $v_0$.

\subsubsection*{7.4. The three identifications}

\begin{itemize}
\item[(i)]
\emph{The tori of \cite[Definition~2.11]{QEDII}.}
A torus is a pair $T=(\Gamma,\chi)$ formed by a lattice of periods $\Gamma$ and a character $\chi\colon\Gamma\to\{1,-1\}$.
The spinors of $T$ are the functions $\mathrm{fX}\colon X\to\mathbb{C}^4$ on the set $X$ of the fermion sites such that $\mathrm{fX}(\mathrm{xa}+v)=\chi(v)\cdot\mathrm{fX}(\mathrm{xa})$ for all $\mathrm{xa}\in X$ and $v\in\Gamma$;
these are the fields on $T$ of \cite[Definition~2.11]{QEDII}, described in the same words in \cite[Definitions~3.9 and~4.1]{QEDII}.
The operator $D^{T}_A$ is the operator of \cite[Definition~2.10]{QEDII} on the spinors of $T$;
for $M=M_{\mathrm{ch}}$ it is the operator $\mathcal{D}_A$ of 4.2 on these spinors.
The torus $C(n)$ has $\Gamma=\mathrm{Ga}_n$ and $\chi(n_\mu e_\mu)=-1$, that is, $\chi(\mathrm{nz})=\boldsymbol{\chi}(z)$.
For $v\in\mathbb{Z}^3$ one has $\mathrm{geo}(c,g)+v=\mathrm{geo}(c+v,g)$.
By (a) and 2.3 it follows that $\mathrm{fX}\mapsto\mathrm{fX}\circ\mathrm{geo}$ is a linear bijection of the space of the spinors of $C(n)$ onto $\mathfrak{S}^{\mathrm{ch}}(1;n)$, and by 4.2, with $M=M_{\mathrm{ch}}$,
\[
\left(D^{T}_A\mathrm{fX}\right)\circ\mathrm{geo}=\left(\mathcal{D}_A\mathrm{fX}\right)\circ\mathrm{geo}=\left(D_A+\boldsymbol{\mu}\right)\left(\mathrm{fX}\circ\mathrm{geo}\right).
\]
The preimages of the elements $\mathrm{bs}_{x',\beta'}$ under this bijection form a basis of the spinors of $C(n)$, and in this basis $D^{T}_A$ has the matrix of $D_A+\boldsymbol{\mu}$.
Therefore $\det D^{T}_A=\det\left(D_A+\boldsymbol{\mu}\right)$ for every $v_0\in\mathrm{Chm}_0$.
If $n_\mu\ge4$ for $\mu=1,2,3$, then $C(n)$ is a torus in the sense of \cite[Definitions~2.6 and~2.11]{QEDII}.
If $\min_\mu n_\mu\le3$, then \cite[Definition~2.11]{QEDII} defines no torus $C(n)$:
a period lattice $\Gamma$ has $\mathrm{sys}(\Gamma)\ge4$ by \cite[Definition~2.6]{QEDII}, and the condition $n_\mu\ge4$ is part of the definition of $C(n)$ in \cite[Definition~2.11]{QEDII}.
For these $n$ the same words define $T=C(n)$, its spinors and $D^{T}_A$ in \cite[Definition~4.1 and Proposition~4.23]{QEDII}, where $\Gamma$ is any subgroup of $\mathbb{Z}^3$ of rank three and $n_\mu\ge1$.
\cite[Definition~4.1]{QEDII} names a real $\Gamma$-periodic $A$;
for complex $A$ the display is that of \cite[Definition~2.10]{QEDII}, which admits one-cochains on $\mathbb{Z}^3$ with real or complex values, on the same spinors.
In the argument above, $\Gamma$ enters only through $\Gamma=\mathrm{Ga}_n$, and $A$ only through its $\mathrm{Ga}_n$-periodicity (7.1);
so its conclusions hold for these $n$ and for every complex $A$.
\item[(ii)]
\emph{The determinant $\mathrm{Dt}^{\mathrm{ch}}$.}
Section~I.2.7 gives the formula
\[
\mathrm{Dt}^{\mathrm{ch}}(B)=\det\left(D^{\mathrm{ch}}[B]+\boldsymbol{\mu}\right),
\]
with the operator $D^{\mathrm{ch}}[B]$ of Proposition~46\,(b), which is the operator of 4.4.
The setting in which this determinant is taken, a torus with antiperiodic spinors, with four spin components and with $\gamma_\mu=\operatorname{diag}\left(\sigma_\mu,-\sigma_\mu\right)$, is the one that the preamble of Proposition~46 prescribes, together with Sections~I.1.5 and~I.2.1.
The operator $D^{\mathrm{ch}}[B]$ is the operator of 4.3 taken with $m=0$ and with the phases $\exp(iB)$, that is, with $\exp\left(iB\left(c,c+ge_1\right)\right)$ in the place of $U\left(x,x^{(3)}\right)$.
Indeed, by (I.2.1) the coefficient of $\nabla_i$ in 4.3 is $-M_i(x)$, and by the computation in 4.4 the coefficients in the display of $D^{\mathrm{ch}}[B]$ are $M_0(x),\dots,M_3(x)$;
so in both operators the term of index $i$ is $M_i(x)$ applied to $f(x)-f\left(x^{(i)}\right)$ for $i\le2$ and to $f(x)-\exp\left(iB\left(c,c+ge_1\right)\right)\cdot f\left(x^{(3)}\right)$ for $i=3$.
For $B=qA$ this operator is $D_A$, by 4.3 at $m=0$ and by (46.6).
Let $B$ be a $\mathrm{Ga}_n$-periodic complex one-cochain.
Since $q>0$, one has $B=qA$ with the $\mathrm{Ga}_n$-periodic cochain $A:=q^{-1}B$, and by (46.6) and 7.1 the operator $D^{\mathrm{ch}}[B]+\boldsymbol{\mu}=D_A+\boldsymbol{\mu}$ maps $\mathfrak{S}^{\mathrm{ch}}(1;n)$, the space of the spinors on the torus with the sides $n$, into itself;
the number $\mathrm{Dt}^{\mathrm{ch}}(B)$ is the determinant, in the sense of 7.3, of the restriction of $D^{\mathrm{ch}}[B]+\boldsymbol{\mu}$ to this space.
So $\mathrm{Dt}^{\mathrm{ch}}(qA)=\det\left(D_A+\boldsymbol{\mu}\right)$.
\item[(iii)]
\emph{The free ratio.}
The left side of the identity in Proposition~46\,(g+) is the quotient
\[
N(t)=\frac{\det\left(D^{\mathrm{ch}}[B_t]+\bar{m}\cdot\eta\right)}{\det\left(D^{\mathrm{ch}}[0]+\bar{m}\cdot\eta\right)}
\]
of 10.5, formed on the space
\[
\left\{f:\ f(c+ne_\mu,g)=-f(c,g),\ \mu=1,2,3\right\}=\mathfrak{S}^{\mathrm{ch}}(1;(n,n,n)).
\]
Its numerator is $\det\left(D^{\mathrm{ch}}[B_t]+\bar{m}\cdot\eta\right)$ on this space.
By (46.6) it is $\det\left(D_A+\boldsymbol{\mu}\right)$ for the sides $(n,n,n)$, for $\boldsymbol{\mu}=\bar{m}\cdot\eta$ and for the cochain $A$ with $qA=B_t$.
\end{itemize}

\subsubsection*{7.5. The rule for the Berezin integral}

\begin{named}{Lemma 7.5}
Let $\mathrm{zt}_j,\mathrm{zb}_j$, $j=1,\dots,\mathrm{dq}$, be the generators of a Grassmann algebra over $\mathbb{C}$, and put $\mathrm{Top}:=\prod_{j=1}^{\mathrm{dq}}\mathrm{zt}_j\cdot\mathrm{zb}_j$.
For every complex $\mathrm{dq}\times\mathrm{dq}$ matrix $\mathrm{Dm}$ the coefficient of $\mathrm{Top}$ in
\[
\exp\left(-\sum_{j,j'}\mathrm{zb}_j\cdot\mathrm{Dm}_{jj'}\cdot\mathrm{zt}_{j'}\right)
\]
is $\det\mathrm{Dm}$.
\end{named}

\begin{proof}
Put
\[
\mathrm{en}_{jj'}:=-\mathrm{zb}_j\cdot\mathrm{Dm}_{jj'}\cdot\mathrm{zt}_{j'}=\mathrm{Dm}_{jj'}\cdot\mathrm{zt}_{j'}\cdot\mathrm{zb}_j.
\]
The elements $\mathrm{en}_{jj'}$ are even, so they commute with one another;
a product of two of them with the same first index or with the same second index vanishes, because $\mathrm{zb}_j^2=0$ and $\mathrm{zt}_{j'}^2=0$.
For every integer $\mathrm{ii}\ge0$ the power $\left(\sum_{j,j'}\mathrm{en}_{jj'}\right)^{\mathrm{ii}}$ is homogeneous of degree $2\,\mathrm{ii}$.
Hence the part of degree $2\,\mathrm{dq}$ of $\exp\left(\sum_{j,j'}\mathrm{en}_{jj'}\right)$ is
\[
\frac{1}{\mathrm{dq}!}\left(\sum_{j,j'}\mathrm{en}_{jj'}\right)^{\mathrm{dq}}=\frac{1}{\mathrm{dq}!}\sum\mathrm{en}_{j_1j'_1}\cdots\mathrm{en}_{j_{\mathrm{dq}}j'_{\mathrm{dq}}},
\]
where the sum extends over all sequences $\left((j_1,j'_1),\dots,(j_{\mathrm{dq}},j'_{\mathrm{dq}})\right)$.
A term vanishes unless $j_1,\dots,j_{\mathrm{dq}}$ are pairwise different and $j'_1,\dots,j'_{\mathrm{dq}}$ are pairwise different.
In that case there is exactly one permutation $\mathrm{pe}$ of $\{1,\dots,\mathrm{dq}\}$ with $j'_{\mathrm{ii}}=\mathrm{pe}(j_{\mathrm{ii}})$ for $\mathrm{ii}=1,\dots,\mathrm{dq}$, and, since the factors commute, the term equals $\prod_{j=1}^{\mathrm{dq}}\mathrm{en}_{j,\mathrm{pe}(j)}$.
Each $\mathrm{pe}$ arises from exactly $\mathrm{dq}!$ sequences, namely from the orderings of the $\mathrm{dq}$ factors.
So the part of degree $2\,\mathrm{dq}$ equals
\[
\sum_{\mathrm{pe}}\left(\prod_{j}\mathrm{Dm}_{j,\mathrm{pe}(j)}\right)\cdot\prod_{j=1}^{\mathrm{dq}}\left(\mathrm{zt}_{\mathrm{pe}(j)}\cdot\mathrm{zb}_j\right).
\]
\emph{Sign.}
Let $\mathrm{oa}_1,\dots,\mathrm{oa}_{\mathrm{dq}},\mathrm{ob}_1,\dots,\mathrm{ob}_{\mathrm{dq}}$ be odd elements that anticommute with one another.
Moving $\mathrm{oa}_j$ to the left across $\mathrm{ob}_1,\dots,\mathrm{ob}_{j-1}$ produces $j-1$ changes of sign, and $\sum_j(j-1)=\frac{\mathrm{dq}(\mathrm{dq}-1)}{2}$;
hence
\[
\prod_{j=1}^{\mathrm{dq}}\left(\mathrm{oa}_j\cdot\mathrm{ob}_j\right)=(-1)^{\frac{\mathrm{dq}(\mathrm{dq}-1)}{2}}\cdot\mathrm{oa}_1\cdots\mathrm{oa}_{\mathrm{dq}}\cdot\mathrm{ob}_1\cdots\mathrm{ob}_{\mathrm{dq}}.
\]
Apply this with $\mathrm{oa}_j=\mathrm{zt}_{\mathrm{pe}(j)}$ and $\mathrm{ob}_j=\mathrm{zb}_j$;
then use $\mathrm{zt}_{\mathrm{pe}(1)}\cdots\mathrm{zt}_{\mathrm{pe}(\mathrm{dq})}=\operatorname{sgn}(\mathrm{pe})\cdot\mathrm{zt}_1\cdots\mathrm{zt}_{\mathrm{dq}}$;
then apply the same rule with $\mathrm{oa}_j=\mathrm{zt}_j$ and $\mathrm{ob}_j=\mathrm{zb}_j$, read from right to left.
The two factors $(-1)^{\frac{\mathrm{dq}(\mathrm{dq}-1)}{2}}$ cancel, and
\[
\prod_j\left(\mathrm{zt}_{\mathrm{pe}(j)}\cdot\mathrm{zb}_j\right)=\operatorname{sgn}(\mathrm{pe})\cdot\mathrm{Top}.
\]
Therefore the coefficient of $\mathrm{Top}$ is $\sum_{\mathrm{pe}}\operatorname{sgn}(\mathrm{pe})\cdot\prod_j\mathrm{Dm}_{j,\mathrm{pe}(j)}=\det\mathrm{Dm}$.
\end{proof}

\subsubsection*{7.6. The two integrals}

\paragraph{Generators.}
The generators are $\psi_\beta(x)$ and $\bar{\psi}_\beta(x)$ with $x\in\mathrm{Yf}$ and $\beta=1,\dots,4$;
these are $\mathrm{dq}$ pairs.
At the other points of $\mathbb{Z}^3\times O_{\mathrm{cub}}$ put
\[
\psi_\beta\left(c^0+\mathrm{nz},g\right):=\boldsymbol{\chi}(z)\cdot\psi_\beta\left(c^0,g\right),\qquad
\bar{\psi}_\beta\left(c^0+\mathrm{nz},g\right):=\boldsymbol{\chi}(z)\cdot\bar{\psi}_\beta\left(c^0,g\right)
\]
((I.AP); \cite[Definition~2.11]{QEDII}).
In \cite[Definitions~2.8 and~2.11]{QEDII} the generators are indexed by the points $\mathrm{geo}(x)$ of $X$ and the components by $\alpha$;
by (a) this is an indexing by $x$, and $\mathrm{Yf}$ corresponds to a set of representatives of $X_T=X/\mathrm{Ga}_n$.

\paragraph{The Berezin integral.}
For a finite set $Y$ of fermion sites, the Berezin integral of \cite[Definition~2.8]{QEDII} is
\[
\int_Y:=\int\prod_{x:\ \mathrm{geo}(x)\in Y}\ \prod_\beta\left(\mathrm{d}\bar{\psi}_\beta(x)\,\mathrm{d}\psi_\beta(x)\right),
\qquad\text{normalised by}\qquad
\int\mathrm{d}\bar{\psi}_\beta(x)\,\mathrm{d}\psi_\beta(x)\,\psi_\beta(x)\cdot\bar{\psi}_\beta(x)=1;
\]
it is a product of even pairs, whose order is immaterial.
Write $\psi=\psi_\beta(x)$ and $\bar{\psi}=\bar{\psi}_\beta(x)$, and let $F=F_0+\psi\cdot F_1+\bar{\psi}\cdot F_2+\psi\cdot\bar{\psi}\cdot F_3$ with elements $F_0,\dots,F_3$ that do not contain these two generators;
the integral over this pair maps $F$ to $F_3$.
The elements $\psi_\beta(x)\cdot\bar{\psi}_\beta(x)$ are even and commute with all elements.
So the integral over all pairs, that is, $\int_Y$ for the set $Y:=\mathrm{geo}(\mathrm{Yf})$, maps $F$ to the coefficient of
\[
\mathrm{Top}=\prod_{x\in\mathrm{Yf}}\prod_\beta\psi_\beta(x)\cdot\bar{\psi}_\beta(x)
\]
in $F$.
(This is (I.3.5) with $\nu=1$.)
In 7.6 and 7.7, $Y$ denotes this set, the image under $\mathrm{geo}$ of the fundamental domain $\mathrm{Yf}$ of $\mathrm{Ga}_n$, as in Proposition~46\,(e).

\paragraph{The forms.}
For a linear map $O$ of $\mathfrak{S}^{\mathrm{ch}}(1;n)$, the element $(O\psi)_\beta(x)$ is given by (46.7) with the generators in the place of the values:
\[
(O\psi)_\beta(x):=\sum_{x'\in\mathrm{Yf}}\sum_{\beta'}\mathrm{Dm}_{(x,\beta),(x',\beta')}\cdot\psi_{\beta'}(x').
\]
With $\bar{\psi}(x)\cdot(O\psi)(x)=\sum_\beta\bar{\psi}_\beta(x)\cdot(O\psi)_\beta(x)$ one has
\begin{align*}
\sum_{x\in\mathrm{Yf}}\bar{\psi}(x)\cdot(O\psi)(x)&=\sum_{(x,\beta),(x',\beta')}\bar{\psi}_\beta(x)\cdot\mathrm{Dm}_{(x,\beta),(x',\beta')}\cdot\psi_{\beta'}(x'),\\
\left\langle\bar{\psi},O\psi\right\rangle^{\mathrm{ch}}=\int\mathrm{d}x\,\bar{\psi}(x)\cdot(O\psi)(x)&=\frac{1}{48}\sum_{c\in\mathrm{Fb}}\sum_{g\in O_{\mathrm{cub}}}\bar{\psi}(c,g)\cdot(O\psi)(c,g)\\
&=\sum_{(x,\beta),(x',\beta')}\bar{\psi}_\beta(x)\cdot\left(\frac{\mathrm{Dm}}{48}\right)_{(x,\beta),(x',\beta')}\cdot\psi_{\beta'}(x'),
\end{align*}
by (I.1.3) at $\eta_0=1$, with $c$ ranging over the fundamental domain $\mathrm{Fb}$ of $\mathrm{Ga}_n$: each of the $48$ sites of a cluster has the weight $1/48$ (Section~I.3.6).
Lemma~7.5, applied to $\mathrm{Dm}$ and to $\mathrm{Dm}/48$, which are matrices of size $\mathrm{dq}=192|n|$, gives
\[
\int_Y\exp\left(-\sum_x\bar{\psi}(x)\cdot(O\psi)(x)\right)=\det\mathrm{Dm}=\det O,\qquad
\int_Y\exp\left(-\left\langle\bar{\psi},O\psi\right\rangle^{\mathrm{ch}}\right)=\det\left(\frac{\mathrm{Dm}}{48}\right)=48^{-192|n|}\cdot\det O.
\]

\paragraph{Maps given by a display.}
Let $O$ be the restriction to $\mathfrak{S}^{\mathrm{ch}}(1;n)$ of an operator on the functions $f\colon\mathbb{Z}^3\times O_{\mathrm{cub}}\to\mathbb{C}^4$ of the form $(Of)(x)=\sum_{x''}\mathrm{Kv}(x,x'')\cdot f(x'')$, with $4\times4$ blocks $\mathrm{Kv}(x,x'')$ and, for each $x$, a finite sum over points $x''$ of $\mathbb{Z}^3\times O_{\mathrm{cub}}$;
assume that this operator maps $\mathfrak{S}^{\mathrm{ch}}(1;n)$ into itself.
By 2.3, every $x''$ is $x'+\mathrm{nz}$ with exactly one pair $x'\in\mathrm{Yf}$, $z\in\mathbb{Z}^3$.
For $f\in\mathfrak{S}^{\mathrm{ch}}(1;n)$ one has $f(x'+\mathrm{nz})=\boldsymbol{\chi}(z)\cdot f(x')$;
hence
\[
(Of)(x)=\sum_{x'\in\mathrm{Yf}}\left[\sum_{z\in\mathbb{Z}^3}\boldsymbol{\chi}(z)\cdot\mathrm{Kv}(x,x'+\mathrm{nz})\right]\cdot f(x')\qquad(x\in\mathrm{Yf}).
\]
The values $f(x')$, $x'\in\mathrm{Yf}$, are arbitrary (7.2).
Comparison with (46.7) shows that the $4\times4$ block of $\mathrm{Dm}$ at $(x,x')$ is $\sum_z\boldsymbol{\chi}(z)\cdot\mathrm{Kv}(x,x'+\mathrm{nz})$.
The generators obey the same rule, $\psi(x'+\mathrm{nz})=\boldsymbol{\chi}(z)\cdot\psi(x')$.
So the display, with the generators in the place of the values, reads
\begin{equation}
\sum_{x''}\mathrm{Kv}(x,x'')\cdot\psi(x'')=\sum_{x'\in\mathrm{Yf}}\left[\sum_z\boldsymbol{\chi}(z)\cdot\mathrm{Kv}(x,x'+\mathrm{nz})\right]\cdot\psi(x')=(O\psi)(x)\qquad(x\in\mathrm{Yf}).
\tag{46.7a}
\end{equation}
For $O=D_A+\boldsymbol{\mu}$ (7.1) the display is the first line of the computation in 4.2:
\[
\mathrm{Kv}(x,x)=M_{\mathrm{ch}},\qquad
\mathrm{Kv}\left(x,x^{(i)}\right)=-M_i(x)\cdot U_A\left(x,x^{(i)}\right)\quad(i=0,\dots,3),\qquad
\mathrm{Kv}(x,x'')=0\ \text{ for all other }x''.
\]
Here the five points $x,x^{(0)},\dots,x^{(3)}$ are pairwise different, because their labels are $g,gs_0,\dots,gs_3$ and $1,s_0,\dots,s_3$ are five different matrices.
In \cite[Definitions~2.10 and~2.11]{QEDII} the element $\left(\mathcal{D}_A\psi\right)(\mathrm{xa})$ is formed by the display of \cite[Definition~2.10]{QEDII}, that is, by the display of $\mathcal{D}_A$ in 4.2, with the generators $\psi(\mathrm{xa})$, $\psi\left(\mathrm{xa}^{(i)}\right)$ at the sites of $X$, where $\mathrm{xa}^{(i)}$ is the neighbour of $\mathrm{xa}$ across the wall of type $i$ (\cite[Definition~2.4]{QEDII}), and with the rule $\psi(\mathrm{xa}+v)=\chi(v)\cdot\psi(\mathrm{xa})$, $v\in\Gamma$.
Let $T=C(n)$, $M=M_{\mathrm{ch}}$ and $\mathrm{xa}=\mathrm{geo}(x)$, and name the generator at $\mathrm{geo}(x)$ by $\psi(x)$.
If $n_\mu\ge4$ for $\mu=1,2,3$, the Grassmann algebra $\mathrm{Gr}_T$, the element $S_T(A)$ and the integral $\int_T$ are those of \cite[Definition~2.11]{QEDII}.
If $\min_\mu n_\mu\le3$, then $C(n)$ is not a torus of \cite[Definition~2.11]{QEDII} (7.4\,(i));
for these $n$ the objects $\mathrm{Gr}_T$, $S_T(A)$ and $\int_T$ on the general torus $C(n)$ of \cite[Definition~4.1 and Proposition~4.23]{QEDII} are defined by the words of \cite[Definition~2.11]{QEDII}.
The computation that follows uses of $T$ only that $\Gamma=\mathrm{Ga}_n$ and $\chi(\mathrm{nz})=\boldsymbol{\chi}(z)$, and of $A$ only that it is $\mathrm{Ga}_n$-periodic;
so it holds for every $n$.
By the identifications of 4.2, namely
\[
\mathrm{xa}^{(i)}=\mathrm{geo}\left(x^{(i)}\right),\qquad h\left(\mathrm{xa},\mathrm{xa}^{(i)}\right)=M_i(x),\qquad\mathcal{U}_A\left(\mathrm{xa},\mathrm{xa}^{(i)}\right)=U_A\left(x,x^{(i)}\right),
\]
and with $\psi(x'+\mathrm{nz})=\boldsymbol{\chi}(z)\cdot\psi(x')$ (by 7.4\,(i), $\chi(\mathrm{nz})=\boldsymbol{\chi}(z)$ and $\mathrm{geo}(x')+\mathrm{nz}=\mathrm{geo}(x'+\mathrm{nz})$), the element $\left(\mathcal{D}_A\psi\right)(\mathrm{xa})$ is
\[
M_{\mathrm{ch}}\cdot\psi(x)-\sum_iM_i(x)\cdot U_A\left(x,x^{(i)}\right)\cdot\psi\left(x^{(i)}\right).
\]
By (46.7a) it is $\left(\left(D_A+\boldsymbol{\mu}\right)\psi\right)(x)$.
The set $Y=\mathrm{geo}(\mathrm{Yf})$ is a set of representatives of $X_T$.
Hence, with $\mathrm{xa}$ ranging over a set of representatives of $X_T$,
\begin{equation}
S_T(A)=\sum_{\mathrm{xa}\in X_T}\bar{\psi}(\mathrm{xa})\cdot\left(\mathcal{D}_A\psi\right)(\mathrm{xa})=\sum_{x\in\mathrm{Yf}}\bar{\psi}(x)\cdot\left(\left(D_A+\boldsymbol{\mu}\right)\psi\right)(x):
\tag{46.7b}
\end{equation}
the element $S_T(A)$ of \cite[Definition~2.11]{QEDII} is the counting form of the map $D_A+\boldsymbol{\mu}$.
The integral $\int_T$ of \cite[Definition~2.11]{QEDII} is the Berezin integral over the generators at a set of representatives of $X_T$, here $\int_Y$.
By the first of the two integrals above,
\[
\int_T\exp\left(-S_T(A)\right)=\det\left(D_A+\boldsymbol{\mu}\right),
\]
and this is $\det D^{T}_A$ by 7.4\,(i).

\paragraph{Other fundamental domains.}
By 2.3, every set of representatives of $\left(\mathbb{Z}^3\times O_{\mathrm{cub}}\right)/\mathrm{Ga}_n$ has the form
\[
\mathrm{Yf}'=\left\{x+n\,z(x):\ x\in\mathrm{Yf}\right\}
\]
with a map $z(\cdot)\colon\mathrm{Yf}\to\mathbb{Z}^3$, where $n\,z(x)$ is the vector $\mathrm{nz}$ for $z=z(x)$.
The generators at the points of $\mathrm{Yf}'$ are $\boldsymbol{\chi}(z(x))\cdot\psi_\beta(x)$ and $\boldsymbol{\chi}(z(x))\cdot\bar{\psi}_\beta(x)$;
so the product $\psi_\beta(x')\cdot\bar{\psi}_\beta(x')$ at $x'=x+n\,z(x)$ is $\boldsymbol{\chi}(z(x))^2\cdot\psi_\beta(x)\cdot\bar{\psi}_\beta(x)$.
Thus $\mathrm{Top}$ is the same element, and the Berezin integral over the generators at $\mathrm{geo}(\mathrm{Yf}')$, the coefficient of $\mathrm{Top}$ in the part of highest degree, which is a multiple of $\mathrm{Top}$, is the same functional as $\int_Y$.
The basis that belongs to $\mathrm{Yf}'$ consists of the elements $\boldsymbol{\chi}(z(x'))\cdot\mathrm{bs}_{x',\beta'}$;
the matrix of $O$ in this basis has the entries $\boldsymbol{\chi}(z(x))\cdot\mathrm{Dm}_{(x,\beta),(x',\beta')}\cdot\boldsymbol{\chi}(z(x'))$, and the element $(O\psi)_\beta(x+n\,z(x))$ formed with it is
\[
\sum_{x',\beta'}\boldsymbol{\chi}(z(x))\cdot\mathrm{Dm}_{(x,\beta),(x',\beta')}\cdot\boldsymbol{\chi}(z(x'))^2\cdot\psi_{\beta'}(x')=\boldsymbol{\chi}(z(x))\cdot(O\psi)_\beta(x).
\]
Hence $\bar{\psi}(x')\cdot(O\psi)(x')$ takes the same value at $x'=x+n\,z(x)$ as at $x'=x$, and both forms are the same elements.

\subsubsection*{7.7. $\mathrm{D}\Psi$; the constant}

The functional $\mathrm{D}\Psi:=48^{192|n|}\cdot\int_Y$ satisfies, by 7.6,
\[
\int\mathrm{D}\Psi\,\exp\left(-\left\langle\bar{\Psi},O\Psi\right\rangle^{\mathrm{ch}}\right)=\det O
\]
for every $O$.
If a multiple $\mathrm{cz}\cdot\int_Y$ has this property, then $\mathrm{cz}\cdot48^{-192|n|}=1$ (take $O=1$).
$\mathrm{D}\Psi$ is the integral (I.3.5) for the index set $\left\{(x,\beta):\ x\in\mathrm{Yf},\ \beta=1,\dots,4\right\}$ with the weights $\nu=1/48$:
\[
\prod_{(x,\beta)}\nu^{-1}=48^{192|n|}.
\]
The number $48^{192|n|}$ depends on $|n|$ only; it involves neither $A$ nor $\boldsymbol{\mu}$.
For linear maps $O_1$, $O_2$ of $\mathfrak{S}^{\mathrm{ch}}(1;n)$ with $\det O_2\neq0$, the quotient of the integrals of $\exp\left(-\left\langle\bar{\psi},O_1\psi\right\rangle^{\mathrm{ch}}\right)$ and of $\exp\left(-\left\langle\bar{\psi},O_2\psi\right\rangle^{\mathrm{ch}}\right)$ is $\det O_1/\det O_2$, whether it is formed with $\int_Y$ or with $\mathrm{D}\Psi$, and it equals the quotient of the integrals $\int_Y$ of the two counting forms.
The density (I.7.2) is $\exp\left(-\varepsilon^{\mathrm{v}}_{\mathrm{in}}\cdot n^3\right)$ times such a quotient, with
\[
O_1=D_{A(F^{\mathrm{u}})+H_\theta+a}+\bar{m}\cdot\eta+m_{\mathrm{in}},\qquad O_2=D_0+\bar{m}\cdot\eta.
\]
For $\rho\ge0$ one has $\det\left(D_0+\bar{m}\cdot\eta\right)>0$: this determinant is $\ge0$ by Lemma~8.1\,(c) and $\neq0$ by Lemma~8.1\,(d), applied with $A=0$ and $\boldsymbol{\mu}=\bar{m}\cdot\eta>0$.
For $\rho<0$ the number $48^{192|n|}$ is a common factor of the two integrals, and the quotient is taken where the denominator does not vanish.
The expectations $E^{\mathrm{af}}_n=Z_n[\,\cdot\,]/Z_n[1]$ and $E^{\mathrm{av}}_n=Z^{(0)}_n[\,\cdot\,]/Z^{(0)}_n[1]$ and the measure $\bar{\mu}=\bar{\nu}/\bar{\nu}(1)$ are quotients of two expressions, each of which is an integral of a determinant on $\mathfrak{S}^{\mathrm{ch}}(1;n)$ against a measure that does not involve the fermion;
a common factor of all these determinants cancels.
\qed

\subsection*{8. Proof of (f)}

\subsubsection*{8.1. Adjoint, norm, sign, real part}

On $\mathfrak{S}^{\mathrm{ch}}(1;n)$ consider the Hermitian form and the norm
\[
(f',f):=\sum_{x\in\mathrm{Yf}}\overline{f'(x)}\cdot f(x),\qquad\|f\|^2:=(f,f).
\]
The symbol $\|\cdot\|$ also denotes the corresponding operator norm, and ${}^{*}$ denotes the adjoint with respect to $(\cdot,\cdot)$.
The pairing $(\cdot,\cdot)^{\mathrm{ch}}$ of Section~I.1.6 is $\frac{1}{48}$ times this form, so that the two give the same adjoints and the same operator norms.
Since $\boldsymbol{\chi}(z)^2=1$, the summand is $\mathrm{Ga}_n$-periodic; the sum may therefore extend over any set of representatives of $(\mathbb{Z}^3\times O_{\mathrm{cub}})/\mathrm{Ga}_n$.
Let $R_i(A)$ be the operator of (I.2.2), and let $\mathrm{Mu}_i$ be the multiplication by $M_i$:
\[
\left(R_i(A)f\right)(x):=U_A\left(x,x^{(i)}\right)f\left(x^{(i)}\right),\qquad\left(\mathrm{Mu}_if\right)(x):=M_i(x)f(x).
\]

\begin{named}{Lemma 8.1}
Let $n_\mu\ge1$, let $\rho$ and $\boldsymbol{\mu}$ be real, let $A$ be real and $\mathrm{Ga}_n$-periodic, and let $O:=D_A+\boldsymbol{\mu}$ on $\mathfrak{S}^{\mathrm{ch}}(1;n)$.
Then
\[
O=M_{\mathrm{ch}}-\sum_{i=0}^{3}\mathrm{Mu}_i\,R_i(A),
\]
and the following hold.
\begin{itemize}
\item[(a)] $R_i(A)$ maps $\mathfrak{S}^{\mathrm{ch}}(1;n)$ into itself, $R_i(A)^2=1$ and $R_i(A)^{*}=R_i(A)$; moreover $O^{*}=\gamma_5\,O\,\gamma_5$.
\item[(b)] One has
\[
\|O\|\le\mathrm{Kd}:=\left|M_{\mathrm{ch}}\right|+\frac{1+|\rho|}{2}\cdot\mathrm{sw},\qquad\left|\det O\right|\le\mathrm{Kd}^{192|n|}.
\]
\item[(c)] $\det O\ge0$.
\item[(d)] One has
\[
O+O^{*}=2\boldsymbol{\mu}+\rho\sum_{i=0}^{3}w_i\left(1-R_i(A)\right),\qquad\left(f,\left(1-R_i(A)\right)f\right)=\frac{1}{2}\left\|f-R_i(A)f\right\|^2\ge0.
\]
Consequently, for $\rho\ge0$ one has $\operatorname{Re}(f,Of)\ge\boldsymbol{\mu}\cdot\|f\|^2$ for every $f$, and $\det O\ne0$ if $\rho\ge0$ and $\boldsymbol{\mu}>0$.
\end{itemize}
\end{named}

\begin{proof}
The identity $O=M_{\mathrm{ch}}-\sum_i\mathrm{Mu}_iR_i(A)$ is the first line of the computation in 4.2.
That computation shows, for $f\in\mathfrak{S}^{\mathrm{ch}}(1;n)$,
\[
M_{\mathrm{ch}}f(x)-\sum_{i=0}^{3}M_i(x)\,U_A\left(x,x^{(i)}\right)f\left(x^{(i)}\right)=\boldsymbol{\mu}f(x)+(D_Af)(x),
\]
and the left side is $\left(M_{\mathrm{ch}}f-\sum_i\mathrm{Mu}_iR_i(A)f\right)(x)$.

(a)
Let $z\in\mathbb{Z}^3$.
By the three facts recorded in 7.1,
\[
\left(R_i(A)f\right)(x+\mathrm{nz})=U_A\left(x,x^{(i)}\right)f\left(x^{(i)}+\mathrm{nz}\right)=\boldsymbol{\chi}(z)\left(R_i(A)f\right)(x),
\]
so that $R_i(A)$ maps $\mathfrak{S}^{\mathrm{ch}}(1;n)$ into itself.
By (46.4),
\[
\left(R_i(A)^2f\right)(x)=U_A\left(x,x^{(i)}\right)U_A\left(x^{(i)},\left(x^{(i)}\right)^{(i)}\right)f\left(\left(x^{(i)}\right)^{(i)}\right)=U_A\left(x,x^{(i)}\right)U_A\left(x^{(i)},x\right)f(x).
\]
For $i\le2$ both factors equal $1$.
For $i=3$, the display (46.2), applied at the site $x^{(3)}=(c+ge_1,gs_3)$, gives
\[
U_A\left(x^{(3)},x\right)=\exp\left(iqA\left(c+ge_1,c+ge_1+gs_3e_1\right)\right)=\exp\left(iqA\left(c+ge_1,c\right)\right)=U_A\left(x,x^{(3)}\right)^{-1}.
\]
Hence $R_i(A)^2=1$.
As $A$ is real, $\left|U_A\left(x,x^{(i)}\right)\right|=1$.
By (46.4), $x\mapsto x^{(i)}$ is a bijection of $\mathbb{Z}^3\times O_{\mathrm{cub}}$ which commutes with the translations by $\mathrm{Ga}_n$; it therefore maps $\mathrm{Yf}$ onto a set of representatives, and
\[
\left\|R_i(A)f\right\|^2=\sum_{x\in\mathrm{Yf}}\left|f\left(x^{(i)}\right)\right|_2^2=\|f\|^2.
\]
A linear isometry of a finite-dimensional space is unitary; thus $R_i(A)^{*}=R_i(A)^{-1}=R_i(A)$.

The matrices $M_i(x)$ are Hermitian, because $\rho$ and $w_i$ are real and, by (c1), $\gamma\cdot v$ is Hermitian for real $v$.
Hence $\mathrm{Mu}_i^{*}=\mathrm{Mu}_i$ and $\left(\mathrm{Mu}_iR_i(A)\right)^{*}=R_i(A)\,\mathrm{Mu}_i$.
Now $\left(R_i(A)\,\mathrm{Mu}_if\right)(x)=U_A\left(x,x^{(i)}\right)M_i\left(x^{(i)}\right)f\left(x^{(i)}\right)$, and, by (46.4) and by (c1) in the form $\gamma_5\gamma_\mu\gamma_5=-\gamma_\mu$,
\[
M_i\left(x^{(i)}\right)=\frac{w_i}{2}\left(\rho-\gamma\cdot\hat{d}_i\left(x^{(i)}\right)\right)=\frac{w_i}{2}\left(\rho+\gamma\cdot\hat{d}_i(x)\right)=\gamma_5M_i(x)\gamma_5.
\]
The constant matrix $\gamma_5$ commutes with $R_i(A)$.
Therefore
\[
\left(\mathrm{Mu}_iR_i(A)\right)^{*}=\gamma_5\,\mathrm{Mu}_i\,\gamma_5\,R_i(A)=\gamma_5\,\mathrm{Mu}_iR_i(A)\,\gamma_5 .
\]
Since $M_{\mathrm{ch}}$ is real and $\gamma_5^2=1$, it follows that
\[
O^{*}=M_{\mathrm{ch}}-\sum_{i=0}^{3}\gamma_5\,\mathrm{Mu}_iR_i(A)\,\gamma_5=\gamma_5\,O\,\gamma_5 .
\]

(b)
By (c1), $\gamma\cdot\hat{d}_i(x)$ is Hermitian and its square is $\left|\hat{d}_i(x)\right|_2^2=1$, so that its eigenvalues lie in $\{1,-1\}$.
The eigenvalues of the Hermitian matrix $M_i(x)$ therefore lie in $\left\{\frac{w_i}{2}(\rho-1),\frac{w_i}{2}(\rho+1)\right\}$, and $\left|M_i(x)\right|_2\le\frac{w_i}{2}\left(1+|\rho|\right)$.
Hence $\left\|\mathrm{Mu}_i\right\|\le\frac{w_i}{2}\left(1+|\rho|\right)$; together with $\left\|R_i(A)\right\|=1$ this gives
\[
\|O\|\le\left|M_{\mathrm{ch}}\right|+\sum_{i=0}^{3}\frac{w_i}{2}\left(1+|\rho|\right)=\mathrm{Kd}.
\]
The number $\det O$ is the product of the $192|n|$ eigenvalues of $O$, each of which has modulus at most $\|O\|$.

(c)
By (c2), $\Sigma_3$ commutes with the matrices $\gamma_\mu$, hence with every $\mathrm{Mu}_i$; being a constant matrix, it commutes with $R_i(A)$.
So $O\Sigma_3=\Sigma_3O$.
As $\Sigma_3$ is Hermitian and $\Sigma_3^2=1$, the space $\mathfrak{S}^{\mathrm{ch}}(1;n)$ is the orthogonal sum of
\[
V_{\pm}:=\left\{f:\ \Sigma_3f=\pm f\right\},
\]
and each of $O$, $O^{*}$ maps $V_+$ and $V_-$ into themselves.
With $O_{\pm}:=O|_{V_{\pm}}$ one has
\[
\det O=\det O_+\cdot\det O_-,\qquad O^{*}|_{V_-}=\left(O_-\right)^{*}.
\]
By (c1) and (c2), $\gamma_5$ is unitary, $\gamma_5^2=1$ and $\gamma_5\Sigma_3=-\Sigma_3\gamma_5$; consequently $\gamma_5$ maps $V_-$ onto $V_+$ and $V_+$ onto $V_-$.
By (a), $\left(O_-\right)^{*}=\gamma_5\,O_+\,\gamma_5$ on $V_-$, an operator conjugate to $O_+$.
Hence $\overline{\det O_-}=\det\left(O_-\right)^{*}=\det O_+$, and
\[
\det O=\det O_+\cdot\overline{\det O_+}=\left|\det O_+\right|^2\ge0.
\]

(d)
By (a),
\[
O+O^{*}=2M_{\mathrm{ch}}-\sum_{i=0}^{3}\left(\mathrm{Mu}_i+\gamma_5\,\mathrm{Mu}_i\,\gamma_5\right)R_i(A),
\]
and
\[
M_i(x)+\gamma_5M_i(x)\gamma_5=\frac{w_i}{2}\left(\rho-\gamma\cdot\hat{d}_i(x)\right)+\frac{w_i}{2}\left(\rho+\gamma\cdot\hat{d}_i(x)\right)=w_i\rho .
\]
Inserting $2M_{\mathrm{ch}}=2\boldsymbol{\mu}+\rho\sum_iw_i$ one obtains $O+O^{*}=2\boldsymbol{\mu}+\rho\sum_iw_i\left(1-R_i(A)\right)$.
Let $\mathrm{Rh}:=R_i(A)$, which is Hermitian and unitary.
Then
\[
\left\|f-\mathrm{Rh}f\right\|^2=\|f\|^2-2\left(f,\mathrm{Rh}f\right)+\left\|\mathrm{Rh}f\right\|^2=2\left(f,(1-\mathrm{Rh})f\right).
\]
For $\rho\ge0$ it follows that $2\operatorname{Re}(f,Of)=\left(f,\left(O+O^{*}\right)f\right)\ge2\boldsymbol{\mu}\cdot\|f\|^2$.
If moreover $\boldsymbol{\mu}>0$, then $Of=0$ implies $f=0$.
\end{proof}

\subsubsection*{8.2. Convergence}

The torus $\mathbb{Z}^3/\mathrm{Ga}_n$ has the $|n|$ sites $c\in\mathrm{Fb}$, the $3|n|$ links $(c,\mu)$ and the $3|n|$ plaquettes $(c;\mu\nu)$, $\mu<\nu$.
The spaces $C^{j}(n)$, the counting norm $\|\cdot\|$ and the coboundary $d$ are those of Sections~I.7.1 and~I.1.1.

\begin{named}{Lemma 8.2}
Let $n_\mu\ge4$ (the range in which Section~I.7.1 defines $C^{j}(n)$; the condition enters only through this definition), let $\rho$ and $\boldsymbol{\mu}$ be real, and let $\Phi$ be a Borel function on $C^{2}(n)$ with $|\Phi\left(F\right)|\le\mathrm{cf}\cdot\exp\left(\mathrm{cf}\cdot\|F\|\right)$.
Then
\begin{multline*}
\sum_{\mathfrak{z}\in C^{2}(n;\mathbb{Z})}\int_{[0,r_q)^{\mathrm{links}}}\mathrm{d}A\,\exp\left(-\frac{1}{2}\left\|dA+r_q\mathfrak{z}\right\|^2\right)\cdot\left|\det\left(D_A+\boldsymbol{\mu}\right)\right|\cdot\left|\Phi\left(dA+r_q\mathfrak{z}\right)\right|\\
\le\mathrm{cf}\cdot\exp\left(\mathrm{cf}^2\right)\cdot\mathrm{Kd}^{192|n|}\cdot\left(r_q\cdot\mathrm{Cth}(1/4)\right)^{3|n|},
\end{multline*}
where
\[
\mathrm{Cth}(1/4):=\sup_{\mathrm{xr}\in\mathbb{R}}\sum_{j\in\mathbb{Z}}\exp\left(-\frac{1}{4}\left(\mathrm{xr}+r_qj\right)^2\right)\le2\sum_{j\ge0}\exp\left(-\frac{1}{4}r_q^2j^2\right)<\infty .
\]
\end{named}

For $\varepsilon>0$ let $\mathrm{Cth}(\varepsilon):=\sup_{y\in\mathbb{R}}\sum_{k\in\mathbb{Z}}\exp\left(-\varepsilon\left(y+r_qk\right)^2\right)$.
This is the number $C_q(\varepsilon)$ of \cite[Lemma~3.14]{QEDII}, and the number of Lemma~8.2 is its value at $\varepsilon=1/4$.
The number $C_q$ without argument of \cite[Lemma~3.13]{QEDII} and of \cite[Theorem~4.13(c)]{QEDII} is $\mathrm{Cth}(1/2)$.

\begin{proof}
The integrand is a Borel function of $A$.
Indeed, by (46.2) and (46.3) the entries of the matrix of $D_A+\boldsymbol{\mu}$ are polynomials in the numbers $\exp\left(\pm iqA(b)\right)$, $b$ a link, so that $\det\left(D_A+\boldsymbol{\mu}\right)$ is continuous in $A$; and $A\mapsto dA+r_q\mathfrak{z}$ is affine.
For real $\mathrm{xr}$ one has
\[
\mathrm{cf}\cdot\mathrm{xr}-\frac{1}{2}\mathrm{xr}^2=\mathrm{cf}^2-\frac{1}{4}\mathrm{xr}^2-\left(\frac{1}{2}\mathrm{xr}-\mathrm{cf}\right)^2\le\mathrm{cf}^2-\frac{1}{4}\mathrm{xr}^2 .
\]
With $\mathrm{xr}=\|F\|$ this gives
\[
\exp\left(-\frac{1}{2}\|F\|^2\right)\cdot\left|\Phi\left(F\right)\right|\le\mathrm{cf}\cdot\exp\left(\mathrm{cf}^2\right)\cdot\exp\left(-\frac{1}{4}\|F\|^2\right)=\mathrm{cf}\cdot\exp\left(\mathrm{cf}^2\right)\cdot\prod_p\exp\left(-\frac{1}{4}F_p^2\right),
\]
the product extending over the $3|n|$ plaquettes $p$.
For fixed $A$ the sum over $\mathfrak{z}\in C^{2}(n;\mathbb{Z})=\mathbb{Z}^{3|n|}$ factorises:
\[
\sum_{\mathfrak{z}}\prod_p\exp\left(-\frac{1}{4}\left((dA)_p+r_q\mathfrak{z}_p\right)^2\right)=\prod_p\sum_{j\in\mathbb{Z}}\exp\left(-\frac{1}{4}\left((dA)_p+r_qj\right)^2\right)\le\mathrm{Cth}(1/4)^{3|n|}.
\]
All terms being non-negative, the sum over $\mathfrak{z}$ may be taken under the integral.
The box $[0,r_q)^{\mathrm{links}}$ has the volume $r_q^{3|n|}$, and $\left|\det\left(D_A+\boldsymbol{\mu}\right)\right|\le\mathrm{Kd}^{192|n|}$ by Lemma~8.1\,(b).
This proves the inequality.

\emph{The number $\mathrm{Cth}(1/4)$.}
The sum is periodic in $\mathrm{xr}$ with the period $r_q$; let $0\le\mathrm{xr}<r_q$.
For $j\ge0$ one has $\left(\mathrm{xr}+r_qj\right)^2\ge r_q^2j^2$.
For $j\le-1$ one has $\mathrm{xr}+r_qj<r_q(j+1)\le0$, and hence $\left(\mathrm{xr}+r_qj\right)^2\ge r_q^2\left(|j|-1\right)^2$.
So each of the two parts of the sum, over $j\ge0$ and over $j\le-1$, is at most $\sum_{j\ge0}\exp\left(-\frac{1}{4}r_q^2j^2\right)$.
\end{proof}

\subsubsection*{8.3. The displays}

Let $n_\mu\ge4$, let $\Phi$ be as in (f), and, as in (f), let $T=C(n)$ and $M=M_{\mathrm{ch}}$.
Thus $T$ is the torus with the period lattice $\mathrm{Ga}_n$ and the character $\chi$ given by $\chi(n_\mu e_\mu)=-1$, $\mu=1,2,3$ (7.4\,(i)).
The display (I.7.1) reads
\[
Z_n[\Phi]=\sum_{\substack{\mathfrak{z}\in C^{2}(n;\mathbb{Z})\\ d\mathfrak{z}=0}}\int_{[0,r_q)^{\mathrm{links}}}\mathrm{d}A\,\exp\left(-\frac{1}{2}\left\|F^{\mathrm{u}}\right\|^2\right)\cdot\det\left(D_A+\boldsymbol{\mu}\right)\cdot\Phi\left(F^{\mathrm{u}}\right),\qquad F^{\mathrm{u}}=dA+r_q\mathfrak{z}.
\]
In \cite[Section~5.2.1]{QEDI} the display of $Z_n[G]$ is written for a cubic torus with an odd side, and in \cite[Definition~10.1]{QEDI}, to which \cite[Section~5.2.1]{QEDI} refers, for the sides $n_1,n_2,n_3$.
It is the same display, with $2\pi/q$ written for $r_q$ and with the Wilson determinant in place of $\det\left(D_A+\boldsymbol{\mu}\right)$.
The class of the functions $G$ admitted there is given by $|G\left(F\right)|\le\mathrm{cf}'\cdot\exp\left(\mathrm{cf}''\cdot\|F\|\right)$; it is the class of (f), with $\mathrm{cf}:=\max\left(\mathrm{cf}',\mathrm{cf}''\right)$.

In \cite[Section~21.1.4]{QEDI} the functional is
\[
\omega_n(\Phi)=\int_{[0,r_q)^{\mathrm{links}}}\mathrm{d}A\sum_{\substack{\mathfrak{z}\in C^{2}(n;\mathbb{Z})\\ d\mathfrak{z}=0}}\Phi\left(dA+r_q\mathfrak{z}\right)\cdot\exp\left(-\frac{1}{2}\left\|dA+r_q\mathfrak{z}\right\|^2\right)\cdot\mathrm{Dt}(qA),
\]
$\mathrm{Dt}(qA)$ being the Wilson determinant.

By \cite[Definition~2.6]{QEDII},
\[
\int_{\mathrm{Conf}(\mathrm{Ga}_n)}\mathrm{d}(A,\mathfrak{z})\,G'(A,\mathfrak{z}):=\sum_{\mathfrak{z}:\ d\mathfrak{z}=0}\ \int_{[0,r_q)^{\mathrm{links}}}\mathrm{d}A\,G'(A,\mathfrak{z}),\qquad F:=dA+r_q\mathfrak{z}.
\]
By \cite[Definition~2.13]{QEDII}, together with \cite[Definition~5.11]{QEDII} for the class $\mathrm{Ig}(T)$,
\[
Z_T[G]=\sum_{\substack{\mathfrak{z}\in C^{2}(\mathrm{Ga}_n;\mathbb{Z})\\ d\mathfrak{z}=0}}\int_{[0,r_q)^{\mathrm{links\ of\ }T}}\mathrm{d}A\,\exp\left(-\frac{1}{2}\|F\|^2\right)\cdot\int_T\exp\left(-S_T(A)\right)\cdot G(A,\mathfrak{z}),
\]
where $S_T(A)=\sum_{x\in X_T}\bar{\psi}(x)\,(D_A\psi)(x)$, with the operator $D_A$ of \cite[Definition~2.10]{QEDII}, which contains the mass parameter $M$, and where $\int_T$ is the Berezin integral of \cite[Definition~2.8]{QEDII} over $X_T=X/\mathrm{Ga}_n$.
The class $\mathrm{Ig}(T)$ consists of the Borel maps $G$ on $\left\{(A,\mathfrak{z}):\ d\mathfrak{z}=0\right\}$ with values in the Grassmann algebra of $T$ such that $G\left(A+r_qm,\mathfrak{z}-dm\right)=G(A,\mathfrak{z})$ for $m\in C^{1}(\mathrm{Ga}_n;\mathbb{Z})$ and such that the coefficients of $G(A,\mathfrak{z})$ are bounded by $\mathrm{Cq}\cdot\exp\left(\left(\frac{1}{2}-\varepsilon_G\right)\|F\|^2\right)$ for some $\mathrm{Cq}$ and some $\varepsilon_G\in\left(0,\frac{1}{2}\right]$.

\emph{Common objects.}
In all these displays the cochains are those of $\mathbb{Z}^3/\mathrm{Ga}_n$, the norm is the counting norm, and the coboundary is given by
\begin{align*}
(dA)_{\mu\nu}(c)&=A_\mu(c)+A_\nu(c+e_\mu)-A_\mu(c+e_\nu)-A_\nu(c),\\
(d\mathfrak{z})(c)&=\left[\mathfrak{z}_{23}(c+e_1)-\mathfrak{z}_{23}(c)\right]-\left[\mathfrak{z}_{13}(c+e_2)-\mathfrak{z}_{13}(c)\right]+\left[\mathfrak{z}_{12}(c+e_3)-\mathfrak{z}_{12}(c)\right];
\end{align*}
these two formulas stand alike in \cite[Definition~2.1]{QEDII}, in \cite[Section~10.1.1]{QEDI} and in Section~I.1.1.

\emph{Comparison with the functionals $Z_n[G]$ and $\omega_n$.}
Once the Wilson determinant is replaced by $\det\left(D_A+\boldsymbol{\mu}\right)$, the displays of \cite[Section~5.2.1]{QEDI} and of \cite[Definition~10.1]{QEDI} become (I.7.1).
The display of \cite[Section~21.1.4]{QEDI} differs from (I.7.1) only in the order of the sum and the integral.
By Lemma~8.2 the integrand is absolutely summable and integrable over $\left\{\mathfrak{z}\ \text{closed}\right\}\times[0,r_q)^{\mathrm{links}}$ with respect to the product of the counting measure and the Lebesgue measure, and by Fubini's theorem the two orders give the same number.
Hence $\omega_n(\Phi)$, formed with $\det\left(D_A+\boldsymbol{\mu}\right)$ in place of the Wilson determinant, equals $Z_n[\Phi]$.

\emph{Comparison with $Z_T$.}
The function $G(A,\mathfrak{z}):=\Phi\left(dA+r_q\mathfrak{z}\right)$ is a scalar Borel function.
It has the required invariance, since $d\left(A+r_qm\right)+r_q\left(\mathfrak{z}-dm\right)=dA+r_q\mathfrak{z}$.
Since $\mathrm{cf}\cdot\mathrm{xr}\le\mathrm{cf}^2+\frac{1}{4}\mathrm{xr}^2$ for real $\mathrm{xr}$,
\[
\mathrm{cf}\cdot\exp\left(\mathrm{cf}\cdot\|F\|\right)\le\mathrm{cf}\cdot\exp\left(\mathrm{cf}^2\right)\cdot\exp\left(\frac{1}{4}\|F\|^2\right).
\]
So $G$ belongs to $\mathrm{Ig}(T)$, with $\mathrm{Cq}=\mathrm{cf}\cdot\exp\left(\mathrm{cf}^2\right)$ and $\varepsilon_G=\frac{1}{4}$.
The Berezin integral is linear and $G$ is a scalar, so that
\[
\int_T\exp\left(-S_T(A)\right)\cdot G=\Phi\left(F\right)\cdot\int_T\exp\left(-S_T(A)\right).
\]
By (46.7b), with $M=M_{\mathrm{ch}}$, the element $S_T(A)$ is the counting form of the map $D_A+\boldsymbol{\mu}$, and by the first of the two integrals of 7.6
\[
\int_T\exp\left(-S_T(A)\right)=\det\left(D_A+\boldsymbol{\mu}\right),
\]
which is $\det D_A^T$ by 7.4\,(i).
Therefore $Z_T[\Phi]$ is the right side of (I.7.1).

\emph{The values of $\rho$.}
In its second paragraph \cite[Definition~2.13]{QEDII} defines $Z_T$ on the class $\mathrm{Ig}(T)$ for $-1\le\rho\le1$, and \cite[Definition~5.11]{QEDII} stands under the hypothesis $-1\le\rho\le1$ made throughout \cite[Section~5]{QEDII}.
For every other real $\rho$, the number $Z_T[\Phi]$ of (f) is the same display, read with the data of Proposition~46, as the preamble of Proposition~46 prescribes.
This display has a meaning for every real $\rho$.
Indeed, by (46.7b) and the first of the two integrals of 7.6 the Berezin integral in it equals $\Phi\left(F\right)\cdot\det\left(D_A+\boldsymbol{\mu}\right)$, as shown above, and by Lemma~8.2, which holds for every real $\rho$, the sum over $\mathfrak{z}$ and the integral over $A$ converge absolutely.
For a bounded $\Phi$ and for $\Phi\ge0$, \cite[Definition~2.13]{QEDII} defines $Z_T[\Phi]$ for every real $\rho$, and its value is
\[
\int_{\mathrm{Conf}(\mathrm{Ga}_n)}\mathrm{d}(A,\mathfrak{z})\,\exp\left(-\frac{1}{2}\|F\|^2\right)\cdot\det D_A^T\cdot\Phi\left(F\right);
\]
by the definition of the measure $\mathrm{d}(A,\mathfrak{z})$ and by 7.4\,(i), the display gives the same number.
This proves (f).

\subsection*{9. Proof of (f+)}

The setting is that of (f).
For $\mu<\nu$ let $\mathfrak{l}$ be the third direction.
For $j\in\mathbb{Z}/n_{\mathfrak{l}}\mathbb{Z}$ let $\Pi_{\mu\nu}(j)$ be the set of the $n_\mu n_\nu$ plaquettes $(c;\mu\nu)$ with $c_{\mathfrak{l}}=j$, and for a two-cochain $u$ put
\[
\mathrm{Su}_{\mu\nu}(u;j):=\sum_{p\in\Pi_{\mu\nu}(j)}u_p .
\]

\subsubsection*{9.1. Exact cochains}

Let $A\in C^{1}(n)$.
Then
\[
\mathrm{Su}_{\mu\nu}(dA;j)=\sum_{c:\ c_{\mathfrak{l}}=j}\left[A_\mu(c)-A_\mu(c+e_\nu)\right]+\sum_{c:\ c_{\mathfrak{l}}=j}\left[A_\nu(c+e_\mu)-A_\nu(c)\right]=0,
\]
because $c\mapsto c+e_\nu$ and $c\mapsto c+e_\mu$ are bijections of $\left\{c\in\mathbb{Z}^3/\mathrm{Ga}_n:\ c_{\mathfrak{l}}=j\right\}$.
Consequently, for $F=dA+r_q\mathfrak{z}$,
\[
\mathrm{Su}_{\mu\nu}(F;j)=r_q\cdot\mathrm{Su}_{\mu\nu}(\mathfrak{z};j).
\]

\subsubsection*{9.2. Closed cochains}

Let $u$ be a two-cochain with $du=0$, and sum
\[
(du)(c)=\left[u_{23}(c+e_1)-u_{23}(c)\right]-\left[u_{13}(c+e_2)-u_{13}(c)\right]+\left[u_{12}(c+e_3)-u_{12}(c)\right]
\]
over the $c$ with $c_{\mathfrak{l}}=j$.
The sum of a bracket with a shift $e_{\mu'}$, $\mu'\ne\mathfrak{l}$, vanishes, since $c\mapsto c+e_{\mu'}$ is a bijection of the set $\left\{c_{\mathfrak{l}}=j\right\}$.
Only the bracket with the shift $e_{\mathfrak{l}}$ remains, and one obtains
\begin{align*}
\mathfrak{l}=3:&\qquad0=\mathrm{Su}_{12}(u;j+1)-\mathrm{Su}_{12}(u;j),\\
\mathfrak{l}=2:&\qquad0=-\left[\mathrm{Su}_{13}(u;j+1)-\mathrm{Su}_{13}(u;j)\right],\\
\mathfrak{l}=1:&\qquad0=\mathrm{Su}_{23}(u;j+1)-\mathrm{Su}_{23}(u;j).
\end{align*}
Hence $\mathrm{Su}_{\mu\nu}(u;j)$ does not depend on $j$.
For a closed $\mathfrak{z}\in C^{2}(n;\mathbb{Z})$ the integer
\[
w_{\mu\nu}(\mathfrak{z}):=\mathrm{Su}_{\mu\nu}(\mathfrak{z};j)
\]
is therefore the same whether it is formed over one plane, as in Section~I.7.1 and in \cite[Definition~13.7]{QEDII}, over the plane $c_{\mathfrak{l}}=0$, as in \cite[Definition~11.1]{QEDI}, or over a coordinate two-torus in the directions $\mu$, $\nu$, as in \cite[(10.2)]{QEDI}.

\emph{Order of the components.}
In Section~I.7.1 and in \cite[Definition~13.7]{QEDII} the flux is the triple $w=\left(w_{23},w_{13},w_{12}\right)$; \cite[Definition~13.7]{QEDII} writes it $k=\left(k_{23},k_{13},k_{12}\right)$.
The flux of \cite[Section~21.1.4]{QEDI} is that of \cite[Definition~11.1]{QEDI}, where the directions are numbered $0,1,2$ and the triple is $\left(\mathfrak{n}_{12},\mathfrak{n}_{02},\mathfrak{n}_{01}\right)$; with the renumbering $\mu\mapsto\mu+1$ of the directions (Section~I.1.1) this is $\left(w_{23},w_{13},w_{12}\right)$.
In \cite[(10.2)]{QEDI} and in \cite[Section~10.1.4]{QEDI} the triple is written $w=\left(w_{12},w_{13},w_{23}\right)$.
The components are labelled by the pairs $\mu<\nu$ (Section~I.7.1), and in \cite[(10.2)]{QEDI}, as in Section~I.7.1, $w_{\mu\nu}$ is the sum of $\mathfrak{z}$ over a coordinate two-torus in the directions $\mu$, $\nu$.
So only the order in which the three components are listed differs, and $Z_{n,w}$ and $h_w$ of \cite[Section~10.1.4]{QEDI} are formed with the same components $w_{\mu\nu}$ as the objects of 9.3 and 9.6.

\subsubsection*{9.3. The sectors}

By \cite[Definition~13.7(a)]{QEDII},
\[
\Phi_w\left(F\right):=\prod_{\mu<\nu}\mathbf{1}\left[\sum_{p\in\Pi_{\mu\nu}}F_p=r_qw_{\mu\nu}\right],\qquad Z^{(w)}_T[G]:=Z_T\left[\Phi_wG\right],
\]
where $\Pi_{\mu\nu}$ is the set of the plaquettes $(c;\mu\nu)$ of one plane, that is, $\Pi_{\mu\nu}(j)$ for one $j$.
By 9.1 and 9.2, $\Phi_w\left(F\right)=\mathbf{1}\left[w(\mathfrak{z})=w\right]$ for $F=dA+r_q\mathfrak{z}$ with closed $\mathfrak{z}$.
Thus the indicator $\mathbf{1}[\text{flux}=w]$ of Section~I.7.3 is the Borel function $\Phi_w$ of $F^{\mathrm{u}}$.
As $0\le\Phi_w\le1$, the function $\Phi_w\Phi$ satisfies the bound of (f) with the same $\mathrm{cf}$, and (f), applied to $\Phi_w\Phi$, gives
\[
Z^{(w)}_T[\Phi]=Z_T\left[\Phi_w\Phi\right]=Z_n\left[\Phi_w\Phi\right]=Z_n\left[\mathbf{1}[\text{flux}=w]\,\Phi\right]=Z^{(w)}_n[\Phi].
\]
By the display of $\omega_n$ in \cite[Section~21.1.4]{QEDI} and the sentence that follows it, $\omega^{(w)}_n(\Phi)$ is the part of the sum in $\omega_n(\Phi)$ with $w(\mathfrak{z})=w$.
By the sentence after the display of $\mathfrak{q}_w(n)$ in the same section, it is $Z_{n,w}[\Phi]$ of \cite[Section~10.1.4]{QEDI}, the part of the sum with $w(\nu)=w$, where $\nu$ denotes the closed integer two-cochain of \cite[Definition~10.1]{QEDI}.
So
\[
\omega^{(w)}_n(\Phi)=\int_{[0,r_q)^{\mathrm{links}}}\mathrm{d}A\sum_{\mathfrak{z}\ \text{closed}}\mathbf{1}\left[w(\mathfrak{z})=w\right]\cdot\Phi\left(F\right)\cdot\exp\left(-\frac{1}{2}\|F\|^2\right)\cdot\mathrm{Dt}(qA)=\omega_n\left(\Phi_w\Phi\right).
\]
Formed with $\det\left(D_A+\boldsymbol{\mu}\right)$ in place of the Wilson determinant, this is $Z_n\left[\Phi_w\Phi\right]$, by 8.3.

\subsubsection*{9.4. $P$}

By \cite[Definition~13.7(a)]{QEDII} the weight of the sector $w$ is $P_T(w):=Z^{(w)}_T[1]/Z_T[1]$, and by Section~I.7.3, $P_n(w):=Z^{(w)}_n[1]/Z_n[1]$.
The numerators are equal by 9.3 with $\Phi=1$, and the denominators are equal by (f).
The denominator lies in $(0,\infty)$ by \cite[Theorem~4.26(a)]{QEDII}: for real $\rho$ and $M$, $w_i>0$, $q>0$ and $v_0\in C_0$, one has $0<Z_T[1]<\infty$ for every torus $T=(\Gamma,\chi)$ in the sense of \cite[Definition~2.11]{QEDII}, whatever its period lattice and its character with values $\pm1$, the trivial character included.
Here $C_0$ is the chamber $\mathrm{Chm}_0$ of \cite[Definition~2.3]{QEDII}; $\rho$ and $M=M_{\mathrm{ch}}$ are real; and $T=C(n)$ with $n_\mu\ge4$ is a torus in the sense of \cite[Definition~2.11]{QEDII}.

\subsubsection*{9.5. The averaged law}

By \cite[Definition~13.7(b)]{QEDII}, $\mu^{\mathrm{av}}_T$ is the law of $\widetilde{F}$ under $Z^{(0)}_T/Z^{(0)}_T[1]$, where $\widetilde{F}$ is the periodic extension of $F$ to the plaquettes of $\mathbb{Z}^3$ (\cite[Definition~2.6]{QEDII}).
By Section~I.7.3, $E^{\mathrm{av}}_n:=Z^{(0)}_n[\,\cdot\,]/Z^{(0)}_n[1]$.

\emph{The denominator.}
By 9.3 with $w=0$ and $\Phi=1$, the number $Z^{(0)}_T[1]=Z^{(0)}_n[1]$ is the sum, over the closed $\mathfrak{z}\in C^{2}(n;\mathbb{Z})$ of flux $0$, of the numbers
\[
\int_{[0,r_q)^{\mathrm{links}}}\mathrm{d}A\,\exp\left(-\frac{1}{2}\left\|dA+r_q\mathfrak{z}\right\|^2\right)\cdot\det\left(D_A+\boldsymbol{\mu}\right).
\]
Each of these numbers is non-negative by Lemma~8.1\,(c), and their sum is finite by Lemma~8.2.
The cochain $\mathfrak{z}=0$ is closed and has flux $0$.
By 7.4\,(i), for $T=C(n)$ and $M=M_{\mathrm{ch}}$, its term is
\[
\int_{[0,r_q)^{\mathrm{links}}}\mathrm{d}A\,\exp\left(-\frac{1}{2}\|dA\|^2\right)\cdot\det D_A^T ,
\]
and this number is positive.
For $\rho^2\ne1$ this is \cite[Theorem~4.13(d)]{QEDII} at $\mathfrak{z}=0$.
The data of that theorem are a real $\rho$ with $\rho^2\ne1$, a real $M$, $w_i>0$, $q>0$, a subgroup $\Gamma$ of $\mathbb{Z}^3$ of rank three, a unitary character $\chi$ of $\Gamma$ and $T=(\Gamma,\chi)$; its clause (d) states that, for every closed $\mathfrak{z}\in C^{2}(\Gamma;\mathbb{Z})$, the number
\[
\int_{[0,r_q)^{\mathrm{Ln}(T)}}\mathrm{d}A\,\exp\left(-\frac{1}{2}\left\|dA+r_q\mathfrak{z}\right\|^2\right)\cdot\det D_A^T
\]
is positive, $\mathrm{Ln}(T)$ being the set of the links of $T$.
For $\rho^2=1$ it is \cite[Theorem~4.25(b)]{QEDII}, stated for $-1\le\rho\le1$ and the other data as in \cite[Theorem~4.13]{QEDII}: every flux term of $Z_T[1]$ is positive.
The two cases cover every real $\rho$.
Hence $0<Z^{(0)}_T[1]<\infty$.
The same conclusion follows from \cite[Theorem~4.26(b)]{QEDII}, by which every flux term of $Z_T[1]$, in particular the term of $\mathfrak{z}=0$, is positive; thus $\mu^{\mathrm{av}}_T$ is defined.

\emph{The law.}
By 9.3 with $w=0$, the two functionals agree on the bounded Borel functions of $F=F^{\mathrm{u}}$.
By Lemma~8.1\,(c) and Lemma~8.2,
\[
\Phi_0\left(F^{\mathrm{u}}\right)\cdot\exp\left(-\frac{1}{2}\left\|F^{\mathrm{u}}\right\|^2\right)\cdot\det\left(D_A+\boldsymbol{\mu}\right)\mathrm{d}A,
\]
summed over the closed $\mathfrak{z}$, is a finite positive measure on $\left\{\mathfrak{z}\ \text{closed}\right\}\times[0,r_q)^{\mathrm{links}}$, and $Z^{(0)}_n[\,\cdot\,]$ is the integral against its image under $(A,\mathfrak{z})\mapsto F^{\mathrm{u}}$.
Therefore $E^{\mathrm{av}}_n$ is the expectation of a probability law of $F^{\mathrm{u}}$ on $C^{2}(n)$, and $\mu^{\mathrm{av}}_T$ is the image of this law under $F^{\mathrm{u}}\mapsto\widetilde{F}$.

\subsubsection*{9.6. $h_w$}

The cochain $h_w$ is defined in the same way in \cite[Definition~13.7(d)]{QEDII}, where it is written $\mathfrak{h}_k$, and in Section~I.7.1: it is the two-cochain on $\mathbb{Z}^3/\mathrm{Ga}_n$ with the constant components $r_qw_{\mu\nu}/(n_\mu n_\nu)$.
In \cite[Section~10.1.4]{QEDI} it has the same components, written $2\pi w_{\mu\nu}/(qn_\mu n_\nu)$.
For each pair $\mu<\nu$ there are $|n|$ plaquettes $(c;\mu\nu)$, and $|n|=n_\mu n_\nu n_{\mathfrak{l}}$.
Hence
\[
\left\|h_w\right\|^2=\sum_{\mu<\nu}|n|\cdot\frac{r_q^2w_{\mu\nu}^2}{\left(n_\mu n_\nu\right)^2}=r_q^2\sum_{\mu<\nu}w_{\mu\nu}^2\cdot\frac{n_{\mathfrak{l}}}{n_\mu n_\nu}=\frac{r_q^2}{|n|}\sum_{\mu<\nu}n_{\mathfrak{l}}^2w_{\mu\nu}^2,
\]
since $n_{\mathfrak{l}}/(n_\mu n_\nu)=n_{\mathfrak{l}}^2/|n|$; \cite[Definition~13.7(d)]{QEDII} writes this norm in the first of these forms.
It follows that
\[
\exp\left(-\frac{1}{2}\left\|h_w\right\|^2\right)=\exp\left(-\frac{r_q^2}{2|n|}\sum_{\mu<\nu}n_{\mathfrak{l}}^2w_{\mu\nu}^2\right),
\]
which is the display of $\mathfrak{q}_w(n)$ in \cite[Section~21.1.4]{QEDI}; the sentence after that display states $\mathfrak{q}_w(n)=\exp\left(-\frac{1}{2}\left\|h_w\right\|^2\right)$, with $h_w$ of \cite[Section~10.1.4]{QEDI}.
This completes the proof of (f+).

\subsection*{10. Proof of (g+) and of (g)}

\subsubsection*{10.1. Data; the two-cochain $u$}

Let $0<\eta\le1$, and let $R>0$ be such that $n=R/\eta$ is an integer $\ge2$.
Let $\beta$ be a probe, that is, a real two-form in $C_c^\infty(\mathbb{R}^3;\Lambda^2)$ with $\beta\ne0$ and $\beta_{\nu\mu}=-\beta_{\mu\nu}$.
Let $v\in\eta\mathbb{Z}^3$; then $\tau_v\beta$ is again such a two-form.
Let $e>0$, $q=e\eta^{\frac{1}{2}}$ and $0<\bar{m}\le1$.
As in Section~I.7.6 and in \cite[Section~21.1.6]{QEDI}, $J_\eta$ is the average over a cube of side $\eta$,
\[
\left(J_\eta\mathrm{fs}\right)(y)=\eta^{-3}\int_{[-\eta/2,\eta/2]^3}\mathrm{d}y'\,\mathrm{fs}\left(y+y'\right),
\]
for a function $\mathrm{fs}$ on $\mathbb{R}^3$.
For $c\in\eta\mathbb{Z}^3$ and $\mu,\nu\in\{1,2,3\}$ put
\begin{equation}
u_{\mu\nu}(c):=\sum_{z\in\mathbb{Z}^3}\left(J_\eta\left(\tau_v\beta\right)_{\mu\nu}\right)\left(c+\frac{\eta}{2}\left(e_\mu+e_\nu\right)+Rz\right).\tag{46.8}
\end{equation}
This is the cochain $u$ of Section~I.7.6, which is that of \cite[Section~21.1.6]{QEDI}, formed with $\tau_v\beta$.
There $u_{\mu\nu}(c)$ is $\left(J_\eta\left(\tau_v\beta\right)_{\mu\nu}\right)\left(c+(\eta/2)(e_\mu+e_\nu)\right)$, periodised, and the periodisation of a function $\mathrm{fs}$ on $\mathbb{R}^3$ is $\mathrm{fs}^R(y)=\sum_z\mathrm{fs}(y+Rz)$ (Sections~I.7.5 and~I.7.6).
The sum in (46.8) is finite, because $J_\eta\left(\tau_v\beta\right)_{\mu\nu}$ has compact support.
Moreover $u_{\mu\nu}$ is $R\mathbb{Z}^3$-periodic, $u_{\nu\mu}=-u_{\mu\nu}$ and $u_{\mu\mu}=0$.
The two-cochain $u$ on $\Lambda_{\eta,R}$ takes the value $u_{\mu\nu}(c)$ on the plaquette $(c;\mu\nu)$, $\mu<\nu$.

\subsubsection*{10.2. $d_\eta^{*}$}

By \cite[Section~7.1.1]{QEDI} and Section~I.7.5, the pairing on the cochains of each degree of $\Lambda_{\eta,R}$ is $\langle u,u'\rangle_\eta=\eta^3\sum u\,u'$; the coboundary $d$ is that of the lattice of mesh one, $d_\eta=\eta^{-1}d$, and $d_\eta^{*}$ is the adjoint of $d_\eta$ with respect to $\langle\cdot,\cdot\rangle_\eta$.
These definitions involve only the lattice $\Lambda_{\eta,R}=\eta\mathbb{Z}^3/R\mathbb{Z}^3$, its cells and the number $\eta$, and they are read with the data of 10.1.
(\cite[Section~21.1.3]{QEDI} takes the definitions of \cite[Section~7.1.1]{QEDI} for every $\eta\in(0,1]$ and every $R\ge1$ with $R/\eta$ an integer $\ge2$; Section~I.7.5 takes them for $n$ of either parity; and the preamble of Proposition~46 takes them for the data of the clause, here for every $R>0$ with this property.)
Let $d^{*}$ be the adjoint of $d$ with respect to the counting pairing $\langle\cdot,\cdot\rangle$.
Then
\[
\left\langle d_\eta A,u\right\rangle_\eta=\eta^3\eta^{-1}\langle dA,u\rangle=\eta^3\left\langle A,\eta^{-1}d^{*}u\right\rangle=\left\langle A,\eta^{-1}d^{*}u\right\rangle_\eta,
\]
and therefore $d_\eta^{*}=\eta^{-1}d^{*}$.

\emph{A formula for $d^{*}$.}
On $\Lambda_{\eta,R}$ one has $(dA)_{\mu\nu}(c)=A_\mu(c)+A_\nu(c+\eta e_\mu)-A_\mu(c+\eta e_\nu)-A_\nu(c)$, and the shifts $c\mapsto c+\eta e_\mu$ are bijections of $\Lambda_{\eta,R}$, for every $n\ge1$.
Hence
\begin{align*}
\langle dA,u\rangle&=\sum_c\sum_{\mu<\nu}u_{\mu\nu}(c)\left[A_\mu(c)-A_\mu(c+\eta e_\nu)\right]+\sum_c\sum_{\mu<\nu}u_{\mu\nu}(c)\left[A_\nu(c+\eta e_\mu)-A_\nu(c)\right]\\
&=\sum_c\sum_{\mu<\nu}A_\mu(c)\left[u_{\mu\nu}(c)-u_{\mu\nu}(c-\eta e_\nu)\right]-\sum_c\sum_{\mu<\nu}A_\nu(c)\left[u_{\mu\nu}(c)-u_{\mu\nu}(c-\eta e_\mu)\right].
\end{align*}
In the coefficient of $A_{\mathfrak{l}}(c)$, the first sum contributes the terms with $\mu=\mathfrak{l}<\nu$,
\[
+\left[u_{\mathfrak{l}\nu}(c)-u_{\mathfrak{l}\nu}(c-\eta e_\nu)\right]=-\left[u_{\nu\mathfrak{l}}(c)-u_{\nu\mathfrak{l}}(c-\eta e_\nu)\right],
\]
and the second sum contributes the terms with $\mu<\nu=\mathfrak{l}$, namely $-\left[u_{\mu\mathfrak{l}}(c)-u_{\mu\mathfrak{l}}(c-\eta e_\mu)\right]$.
Together, since $u_{\mathfrak{l}\mathfrak{l}}=0$,
\begin{equation}
\left(d^{*}u\right)_{\mathfrak{l}}(c)=-\sum_{\mu=1}^{3}\left[u_{\mu\mathfrak{l}}(c)-u_{\mu\mathfrak{l}}(c-\eta e_\mu)\right].\tag{46.9}
\end{equation}

\subsubsection*{10.3. $a_\eta$}

By 10.2,
\[
\left(a_\eta\right)_\nu(c)=\left(d_\eta^{*}u\right)_\nu(c)=-\eta^{-1}\sum_{\mu}\left(u_{\mu\nu}(c)-u_{\mu\nu}(c-\eta e_\mu)\right)=:a_{\eta,\nu}(c).
\]
Thus the one-cochain $a_\eta$ of Section~I.7.6, formed with $\tau_v\beta$, is the function $a_{\eta,\nu}$ given by this explicit display.

\subsubsection*{10.4. $B_t$}

Through the identification $c\mapsto\eta c$ of $\mathbb{Z}^3/n\mathbb{Z}^3$ with $\Lambda_{\eta,R}$ (Section~I.7.5), the cochain $\hat{\mathfrak{a}}=\eta^{\frac{1}{2}}a_\eta$ is read in lattice units:
\[
\hat{\mathfrak{a}}(c,c+e_\nu)=\eta^{\frac{1}{2}}a_{\eta,\nu}(\eta c)\quad\left(c\in\mathbb{Z}^3\right),\qquad\hat{\mathfrak{a}}(c+e_\nu,c)=-\hat{\mathfrak{a}}(c,c+e_\nu);
\]
it is $n\mathbb{Z}^3$-periodic.
Let
\[
B_t(c,c+e_\nu):=t\eta\cdot a_{\eta,\nu}(\eta c),
\]
and, as for the functions $B$ of the oriented bonds in 4.4, $B_t(c',c):=-B_t(c,c')$.
Then, for every $e>0$,
\[
q\cdot\frac{t}{e}\cdot\hat{\mathfrak{a}}(c,c+e_\nu)=e\eta^{\frac{1}{2}}\cdot\frac{t}{e}\cdot\eta^{\frac{1}{2}}a_{\eta,\nu}(\eta c)=t\eta\cdot a_{\eta,\nu}(\eta c)=B_t(c,c+e_\nu).
\]

\subsubsection*{10.5. $N(t)$}

Let
\[
N(t):=\frac{\det\left(D^{\mathrm{ch}}[B_t]+\bar{m}\eta\right)}{\det\left(D^{\mathrm{ch}}[0]+\bar{m}\eta\right)}
\]
on the space $\left\{f:\ f(c+ne_\mu,g)=-f(c,g),\ \mu=1,2,3\right\}=\mathfrak{S}^{\mathrm{ch}}(1;(n,n,n))$, with the operator $D^{\mathrm{ch}}[B]$ of 4.4.
Apply (46.6) with $A=(t/e)\hat{\mathfrak{a}}$, for which $qA=B_t$ by 10.4, and with $A=0$, in both cases with $\boldsymbol{\mu}=\bar{m}\eta$.
This gives
\[
D^{\mathrm{ch}}[B_t]+\bar{m}\eta=D_{(t/e)\hat{\mathfrak{a}}}+\bar{m}\eta,\qquad D^{\mathrm{ch}}[0]+\bar{m}\eta=D_0+\bar{m}\eta
\]
as maps of $\mathfrak{S}^{\mathrm{ch}}(1;(n,n,n))$ into itself (7.1).
So the numerator of $N(t)$ is the numerator of $N_0(t)$ of Section~I.7.6 on the torus $C(n,n,n)$, and the two denominators are equal.
(For $n=2,3$ the name $C(n,n,n)$ is read as in 7.4\,(i): it is the lattice $n\mathbb{Z}^3$ with $\chi(ne_\mu)=-1$, as in \cite[Proposition~4.23]{QEDII}; what enters the two determinants is the space $\mathfrak{S}^{\mathrm{ch}}(1;(n,n,n))$.)
These two equalities hold for every real $\rho$.
For $\rho\ge0$ the denominator is positive.
Indeed, by Lemma~8.1\,(d) with $A=0$ and $\boldsymbol{\mu}=\bar{m}\eta>0$, one has $\operatorname{Re}\left(f,\left(D_0+\bar{m}\eta\right)f\right)\ge\bar{m}\eta\cdot\|f\|^2$ and the determinant does not vanish; it is non-negative by Lemma~8.1\,(c).
Then $N(t)=N_0(t)$.
For $\rho<0$ the two quotients are read where the denominator does not vanish.
Since (46.6) holds for complex $A$, the identity holds for complex $t$, with $D^{\mathrm{ch}}[B]$ read for complex $B$.

\subsubsection*{10.6. The range of $R$ in (g+); clause (g)}

By Section~I.7.6, $D_{\mathrm{be}}$ is the least side of a closed coordinate cube that contains $\operatorname{supp}\beta$; it does not change when $\beta$ is replaced by $\tau_v\beta$.
The condition that $\operatorname{supp}\tau_v\beta$ lie in the interior of a cube of side $R$, which (g+) does not impose, reads $D_{\mathrm{be}}<R$ when the cube is a coordinate cube $y+[0,R]^3$, that is, a period cell of $R\mathbb{Z}^3$.
Indeed, if $D_{\mathrm{be}}<R$, the closed coordinate cube of side $D_{\mathrm{be}}$ that contains $\operatorname{supp}\tau_v\beta$ lies in the interior of the concentric coordinate cube of side $R$.
Conversely, a compact subset of the interior of a coordinate cube of side $R$ lies in a closed concentric coordinate cube of smaller side; hence $D_{\mathrm{be}}<R$.

For cubes in arbitrary position the converse fails.
Let $R<\mathrm{ls}<\sqrt{3}R$.
The segment $\left\{\mathrm{xr}\,e_1:\ \mathrm{xr}\in\mathbb{R},\ |\mathrm{xr}|\le\mathrm{ls}/2\right\}$ lies in the interior of the cube of side $R$ with centre $0$ and a diagonal parallel to $e_1$: this diagonal has the length $\sqrt{3}R$, and its points other than its two end points are interior points of the cube.
The segment is compact, so that a neighbourhood of it lies in the interior as well.
If $\operatorname{supp}\tau_v\beta$ lies in this neighbourhood and contains the two end points of the segment, then $D_{\mathrm{be}}\ge\mathrm{ls}>R$.

The displays (46.8), 10.3, 10.4 and the display of $N(t)$ in 10.5 have a meaning for every $R>0$ with $R/\eta$ an integer $\ge2$, and the arguments of 10.1--10.5 hold without the condition $D_{\mathrm{be}}<R$; for $R\le D_{\mathrm{be}}$ these displays are the definition of $B_t$ and of $N(t)$.
(The setting (FF) of Proposition~36, Section~G5.6.0, has $R>0$ with $R/\eta$ an integer $\ge2$.)

Clause (g) is the case $v=0$.
The lattice probe of \cite[Section~21.1.6]{QEDI} formed with $\beta$ itself is given by (46.8) and 10.3 with $\tau_0\beta=\beta$.
(The definitions of $J_\eta$, $u$ and $a_\eta$ of 10.1--10.3 stand in the first paragraph of \cite[Section~21.1.6]{QEDI} with its displays, and the display in the paragraph on the placement of the probe is $\hat{\mathfrak{a}}=\eta^{\frac{1}{2}}a_\eta$.)
This completes the proof of (g+) and of (g).

\subsection*{11. Remarks on the clauses}

The clauses (a)--(g+) of Proposition~46 consist for the most part of identities between objects defined by displays.
Among their assertions that are not of this kind are the following:
\begin{itemize}
\item[(i)] in (a), the map $\mathrm{geo}$ is a bijection (3.1);
\item[(ii)] for $\rho\ge0$ the free denominator $\det\left(D_0+\bar{m}\eta\right)$ is positive, and $\operatorname{Re}\left(D_0+\bar{m}\eta\right)\ge\bar{m}\eta>0$; (g+) asserts this for (g), for (g+) and for the density named in (e) (7.7, 10.5);
\item[(iii)] the two normalisations of (f+) satisfy $0<Z_T[1]<\infty$ and $0<Z^{(0)}_T[1]<\infty$, so that the quotients $P_T(w)$, $P_n(w)$ and the laws $\mu^{\mathrm{av}}_T$, $E^{\mathrm{av}}_n$ are defined (9.4, 9.5);
\item[(iv)] in (e), the map $D_A+\boldsymbol{\mu}$ sends $\mathfrak{S}^{\mathrm{ch}}(1;n)$ into itself, this space has the dimension $192|n|$, and $\det\left(D_A+\boldsymbol{\mu}\right)$ depends on no basis, on no pairing and not on $v_0$ (7.1--7.3);
\item[(v)] in (f) and (f+), the sum over $\mathfrak{z}$ and the integral over $A$ converge absolutely, so that the displays have a meaning for every real $\rho$ and the two orders of summation and integration give the same number (Lemma~8.2; 8.3); moreover $\det\left(D_A+\boldsymbol{\mu}\right)\ge0$ for real $A$ (Lemma~8.1\,(c)), by which $E^{\mathrm{av}}_n$ is the expectation of a probability law (9.5);
\item[(vi)] in (g+), the sum in (46.8) is finite, so that the displays that define $B_t$ and the two quotients have a meaning for every $R>0$ with $R/\eta$ an integer $\ge2$, in particular for $R\le D_{\mathrm{be}}$ (10.1, 10.6); and the condition that $\operatorname{supp}\tau_v\beta$ lie in the interior of a coordinate cube of side $R$ is equivalent to $D_{\mathrm{be}}<R$ (10.6).
\end{itemize}
Each of the clauses (a)--(g+) holds for all values of the data named in it, namely for every real $\rho$, every real $\boldsymbol{\mu}$, every $q>0$ and every $v_0\in\mathrm{Chm}_0$;
for all sides $n_\mu\ge1$ in (a)--(e) and $n_\mu\ge4$ in (f), (f+);
for every complex periodic $A$ in (b), (e), (g), (g+), whereas in (f), (f+) the cochain $A$ is the real variable of integration;
at every level $k\ge0$ and for every mesh $\eta_k$ in the last sentence of (b);
and for every $\eta\in(0,1]$ and all $R$, $e$, $\bar{m}$, $\beta$, $v$, $t$ in (g), (g+).
In (g) and (g+) the numbers $q$ and $e$ are not independent data: they are tied by $q=e\eta^{\frac{1}{2}}$, and the lattice mass there is $\bar{m}\eta$.
For $\rho<0$ the quotients in (g), (g+) and in the last sentence of (e) are taken where the denominator $\det\left(D_0+\bar{m}\eta\right)$ does not vanish.

\emph{Level $0$.}
The clauses (a)--(g+) are statements in lattice units, that is, at level $0$, except for the last sentence of (b), which holds at every level $k\ge0$ (4.5).
At level $0$ the block spinor has $192$ components per photon site and a fermion site has the weight $1/48$; in 7.6 this weight gives the factor $48^{-192|n|}$.

\subsection*{12. Conditions on the constants}

No constant is chosen in this chapter.
Proposition~46 is stated for every real $\rho$ and for the standard weights, and the restriction $n_\mu\ge4$ in (f) and (f+) is a hypothesis of these clauses on the torus.
The proofs of this chapter use no condition of Sections~I.8.2 and~I.8.4 and of Proposition~14.2 except those of the following two entries of the list of the constants:
the first row of the line (C0) of Section~I.8.2, which fixes $\rho\in(0,1)$, the weights $w$ and $\mathrm{sw}$;
and the entry $\bar{m}\in(0,1]$ (Section~I.0.5; the line (C10) of Section~I.8.2, which is the row~[119] of Proposition~14.2\,(a)).
Of $\rho$ and the weights the proofs use only the following properties, each of which holds under the hypotheses of the assertion in which it enters:
\begin{itemize}
\item the values $w=\left(6,6\sqrt{2},6\sqrt{2},6\right)$ of the weights, through (I.1.1), in (46.5), and in the comparisons of 4.2--4.4 with the weights of \cite[Definition~2.10]{QEDII} and of the operators of 4.3 and 4.4;
\item $w_i>0$, in the weaker form $w_i\ge0$ in Lemma~8.1\,(b), (d), and as a hypothesis of \cite[Theorems~4.26, 4.13 and~4.25]{QEDII} in 9.4 and 9.5;
\item that $\rho$ is real: directly in Lemma~8.1\,(a) and (b); through (a) in Lemma~8.1\,(c) and in the first display of Lemma~8.1\,(d); through Lemma~8.1\,(b) in Lemma~8.2, and through Lemma~8.2 in 8.3 and 9.5; as a hypothesis of \cite[Theorem~4.26]{QEDII} in 9.4 and 9.5 and of \cite[Theorem~4.13]{QEDII} in 9.5, whereas \cite[Theorem~4.25]{QEDII}, applied in 9.5, has it in the form $-1\le\rho\le1$; and in the head of \cite[Definition~2.13]{QEDII}, applied in 8.3;
\item $\rho\ge0$, in Lemma~8.1\,(d), through the inequality $\operatorname{Re}\left(f,\left(D_A+\boldsymbol{\mu}\right)f\right)\ge\boldsymbol{\mu}\cdot\|f\|^2$ for real $\mathrm{Ga}_n$-periodic $A$; this part of Lemma~8.1\,(d) is applied only in 7.7 and 10.5, to the denominator $\det\left(D_0+\bar{m}\eta\right)$, for which (g+) assumes $\rho\ge0$.
\end{itemize}
The hypothesis $\rho^2\ne1$ of \cite[Theorem~4.13]{QEDII} and the range $-1\le\rho\le1$ in which \cite[Definitions~2.13 and~5.11]{QEDII} define $Z_T$ on the class $\mathrm{Ig}(T)$ give no condition on $\rho$.
In 9.5 the case $\rho^2=1$ is covered by \cite[Theorem~4.25]{QEDII}.
For real $\rho$ outside $[-1,1]$, the number $Z_T[\Phi]$ of (f) is the display of 8.3, read with the data of Proposition~46, and this display converges absolutely by Lemma~8.2 (8.3).

Of the entry $\bar{m}\in(0,1]$ only the half $\bar{m}>0$ is used, in the form $\boldsymbol{\mu}=\bar{m}\eta>0$ (with $\eta>0$), as the hypothesis $\boldsymbol{\mu}>0$ of the last assertion of Lemma~8.1\,(d).
It enters in 7.7 and 10.5 only, that is, only through the positivity of the denominator $\det\left(D_0+\bar{m}\eta\right)$ ((g+); Section~11\,(ii)).
The bound $\bar{m}\le1$ is not used, and the equalities of (d), (e), (g) and (g+) use no property of $\bar{m}$.

\section*{Chapter 1a. The free operator, first part (light and heavy part, Schur symbol, coercivity, Poincar\'e inequality): proofs of Propositions 1.1--1.4}

\subsection*{0. Conventions of this chapter; the setting (GS); eleven facts}

All four proofs take place in one setting, the data (GS) of Proposition~1.1\,(f).
These data contain the data of Proposition~1.1\,(a)--(e), the spaces of Section~G1.4 and the fibres of the Bloch decomposition.
The displays of this chapter are numbered (1a.$n$).
Throughout, $\sum_a$ denotes the sum over the six vectors $a\in\{\pm e_1,\pm e_2,\pm e_3\}$.

\subsubsection*{0.0. Labels and letters}

In the proofs of this chapter, (s1)--(s4), (m1)--(m7) and (c1)--(c6) denote the facts (I.s1)--(I.s4), (I.m1)--(I.m7) and (I.c1)--(I.c6) of Sections~I.1.3, I.2.2 and I.2.1, and (S1)--(S9) denote the clauses of Proposition~1.2, in all sections.
The objects $R_i$, $R_3(p)$, $D_0^{(p)}$, $\hat{d}_i(g)$ and $(\nabla^{*}\nabla)(p)$, which are defined in 0.2\,(In3), have the meaning that they have in Section~G5.6.0 and in Proposition~37.1\,(a).
As in Proposition~1.2, $\mathrm{tw}(p):=\sum_\mu(1-\cos p_\mu)$ and, for real $p$, $\mathrm{rp}(p):=\left(|p|_2^2+\mathbf{m}^2\right)^{\frac{1}{2}}$;
$\mathrm{ak}$ denotes a multi-index;
$T_v$ is the translation $(T_vf)(t):=f(t-v)$ of Section~I.1.9.
Where they occur as bound variables, $s$ denotes a real number, $F$ a function on the sites and $z$ a site of $\mathrm{me}\,\mathbb{Z}^3$.
The constants $A_{\mathrm{ak}}$, $A'_{\mathrm{ak}}$ of (S6), (S7) carry a multi-index;
$A_0$ in 2.6 and $A'_0$ in 2.7 are the members with $\mathrm{ak}=0$.
The two numbers $A_1$, $A_2$ of Proposition~1.3\,(b) carry a numeral and are not members of these families.
The connections $A_0$ and $A_1$ of Section~I.5.3, the two summands of an element $A_0+A_1$ of $\tilde{G}_k$, do not occur in this chapter.

\subsubsection*{0.1. Data and definitions}

\paragraph{(GS)}
The data $(\mathrm{me},\mathrm{Lab},R,\theta,e_k,A)$ are the data (GS) of Proposition~1.1\,(f):
a mesh $\mathrm{me}>0$;
a translate $\mathrm{Lab}=t+\mathrm{me}\,\mathbb{Z}^3$ of $\mathrm{me}\,\mathbb{Z}^3$;
a positive integer multiple $R$ of $\mathrm{me}$, or $R=\infty$;
a vector $\theta\in\mathbb{R}^3$;
a number $e_k>0$;
a complex function $A$ on the oriented bonds of $\mathrm{Lab}$ with $A(c',c)=-A(c,c')$ and, if $R<\infty$, with $A(c+R\,e_\mu,c'+R\,e_\mu)=A(c,c')$.
The oriented bonds of $\mathrm{Lab}$ are the pairs $(c,c')$ of points of $\mathrm{Lab}$ with $|c-c'|_2=\mathrm{me}$.
The vector $t$ enters no formula.
Put
\begin{equation*}
u_A(c,c'):=\exp(i\,e_k\cdot\mathrm{me}\cdot A(c,c')).
\end{equation*}
Then $u_A(c',c)=u_A(c,c')^{-1}$, the function $u_A$ is $R\mathbb{Z}^3$-periodic, and $|u_A(c,c')|=1$ if $A$ is real.

\emph{Sites.}
A site is a pair $x=(c,g)\in\mathrm{Lab}\times O_{\mathrm{cub}}$.
As in (I.1.2), (I.2.1) and (I.1.6), with $\mathrm{me}$, $\mathrm{Lab}$ in the places of $\eta_k$, $\mathbb{L}_k$, put
\begin{gather*}
x^{(i)}:=(c,gs_i)\quad(i=0,1,2),\qquad
x^{(3)}:=(c+\mathrm{me}\cdot ge_1,\,gs_3),\qquad
c_{3}:=c+\mathrm{me}\cdot ge_1,\\
\hat{d}_i(x):=-g\hat n_i,\qquad
M_{i}(x):=\frac{w_i}{2}\left(\rho-\gamma\cdot\hat{d}_i(x)\right),\\
U_A(x,x^{(i)}):=1\quad(i=0,1,2),\qquad
U_A(x,x^{(3)}):=u_A(c,c_{3}).
\end{gather*}

\emph{The facts of Sections~I.1.3, I.2.1 and I.2.2 in the setting (GS).}
Of the facts (I.s1)--(I.s4), (I.m1), (I.m2), (I.m5), (I.m6), (I.m7), (I.c1), (I.c2), (I.c3), only two contain the mesh or the lattice:
(s1), through the identity $(x^{(3)})^{(3)}=x$ and through the sets on which $x\mapsto x^{(i)}$ is a bijection (see \emph{Congruence} below),
and (s4), through the oriented photon bond $(c,c+\eta_k a)$;
in the setting (GS), (s4) is proved as (GS8).
All of them hold in (GS), with $\mathrm{me}$, $\mathrm{Lab}$ in the places of $\eta_k$, $\mathbb{L}_k$;
for (s1)--(s4) this is one of the assertions of Proposition~1.1\,(f).
Each of them follows in one line from the definitions of Sections~I.1.2 and I.2.1 and from the displays above;
of (c1)--(c3) and (m1)--(m7) only the following parts are derived.
\begin{itemize}
\item[(s1)] $(x^{(i)})^{(i)}=(c,gs_is_i)=x$ for $i\le2$, because $s_i^2=1$;
$(x^{(3)})^{(3)}=\left(c+\mathrm{me}\cdot ge_1+\mathrm{me}\cdot(gs_3)e_1,\,gs_3s_3\right)=x$, because $s_3e_1=-e_1$.
\item[(s2)] $\hat{d}_i(x^{(i)})=-gs_i\hat n_i=g\hat n_i=-\hat{d}_i(x)$, because $s_i\hat n_i=-\hat n_i$.
\item[(s3)] $\sum_iw_i\hat{d}_i(x)=-g\sum_iw_i\hat n_i=-g\left[6e_3+6(e_2-e_3)+6(e_1-e_2)-6e_1\right]=0$.
\item[(s4)] This is (GS8) below.
\item[(c1)] $\gamma_\mu=\sigma_3\otimes\sigma_\mu$ and $\gamma_5=\sigma_1\otimes1_2$ are Hermitian;
$\gamma_\mu\gamma_\nu+\gamma_\nu\gamma_\mu=1_2\otimes(\sigma_\mu\sigma_\nu+\sigma_\nu\sigma_\mu)=2\delta_{\mu\nu}$;
$\gamma_5\gamma_\mu=(\sigma_1\sigma_3)\otimes\sigma_\mu=-\gamma_\mu\gamma_5$;
$\gamma_5^2=1$.
Consequently $(\gamma\cdot v)^2=|v|_2^2$ for real $v$.
\item[(c2)] In part: $\Sigma_3=\sigma_3\otimes1_2$ commutes with $\gamma_\mu=\sigma_3\otimes\sigma_\mu$.
\item[(c3)] $\mathsf{C}=i\gamma_2$; hence $\mathsf{C}^2=-1$ and $\mathsf{C}^{-1}=-\mathsf{C}$.
Moreover $\gamma_\mu^T=\gamma_\mu$ for $\mu=1,3$, and $\gamma_2^T=-\gamma_2$ because $\sigma_2^T=-\sigma_2$.
Next, $\mathsf{C}^{-1}\gamma_\mu\mathsf{C}=\gamma_2\gamma_\mu\gamma_2$, which equals $-\gamma_\mu$ for $\mu=1,3$ and $\gamma_2$ for $\mu=2$.
Therefore $-\gamma_\mu^T=\mathsf{C}^{-1}\gamma_\mu\mathsf{C}$;
a second conjugation with $\mathsf{C}$, together with $\mathsf{C}^2=-1$, gives $\mathsf{C}^{-1}\gamma_\mu^T\mathsf{C}=-\gamma_\mu$.
\item[(m1)] $\sum_iM_{i}(x)=\frac{\rho}{2}\,\mathrm{sw}-\frac{1}{2}\,\gamma\cdot\left(\sum_iw_i\hat{d}_i(x)\right)=\frac{\rho}{2}\,\mathrm{sw}$ by (s3);
$M_{i}(x^{(i)})=\frac{w_i}{2}\left(\rho+\gamma\cdot\hat{d}_i(x)\right)$ by (s2);
hence $\sum_iM_{i}(x^{(i)})=\frac{\rho}{2}\,\mathrm{sw}=\sum_iM_{i}(x)$.
\item[(m2)] $\mathsf{C}^{-1}M_{i}(x)^T\mathsf{C}=\frac{w_i}{2}\left(\rho+\gamma\cdot\hat{d}_i(x)\right)=M_{i}(x^{(i)})$, by (c3).
\item[(m5)] In part: $\Sigma_3\,M_{i}(x)=M_{i}(x)\,\Sigma_3$, by (c2).
\item[(m6)] $\hat{d}_i(x)$ is a real unit vector, because $g$ is orthogonal.
Hence $\gamma\cdot\hat{d}_i(x)$ is Hermitian, with square $1$ and trace $0$ (note that $\mathrm{tr}\,\gamma_\mu=\mathrm{tr}\,\sigma_3\cdot\mathrm{tr}\,\sigma_\mu=0$);
its eigenvalues are $+1$ and $-1$, each of multiplicity two.
The eigenvalues of $M_{i}(x)$ are $\frac{w_i}{2}(\rho\mp1)$, and $|M_{i}(x)|_2=\frac{w_i}{2}\max\left(|\rho-1|,\rho+1\right)=\frac{w_i}{2}(1+\rho)$, because $\rho>0$.
\item[(m7)] For $M_3$: Section~3, proof of (c).
\end{itemize}
By (s1), $U_A(x^{(i)},x)=U_A(x,x^{(i)})^{-1}$ for $i=0,\dots,3$;
for $i=3$ this reads $U_A(x^{(3)},(x^{(3)})^{(3)})=u_A(c_{3},c)$.

\emph{Congruence.}
Two sites, two bonds $\{x,x^{(i)}\}$, or two photon bonds $\{c,c'\}$ are congruent if a translation by a vector of $R\mathbb{Z}^3$ maps the one onto the other;
for $R=\infty$ this means that they are equal.
``Modulo $R\mathbb{Z}^3$'' means: over one representative of each class of congruent objects.
The map $x\mapsto x^{(i)}$ commutes with the translations;
by (s1) it is an involution of the set of the sites and of the set of their classes.
The site $x^{(i)}$ is never congruent to $x$, because the labels $gs_i$ and $g$ differ.

\emph{Spaces, pairings.}
Let
\begin{gather*}
\mathfrak{S}^{\mathrm{ch}}:=\left\{f:\mathrm{Lab}\times O_{\mathrm{cub}}\to\mathbb{C}^4\,:\,f(c+R\,e_\mu,g)=-\exp(i\theta_\mu)\,f(c,g)\right\},\\
\mathfrak{S}:=\left\{\xi:\mathrm{Lab}\to\mathbb{C}^4\,:\,\xi(c+R\,e_\mu)=-\exp(i\theta_\mu)\,\xi(c)\right\},
\end{gather*}
with no condition if $R=\infty$, and
\begin{gather*}
\int\mathrm{d}x\,F(x):=\frac{\mathrm{me}^3}{48}\sum_{c\bmod R\mathbb{Z}^3}\ \sum_{g\in O_{\mathrm{cub}}}F(c,g),\qquad
(f',f)^{\mathrm{ch}}:=\int\mathrm{d}x\,\overline{f'(x)}\cdot f(x),\qquad
\|f\|_{L^2}:=\left((f,f)^{\mathrm{ch}}\right)^{\frac{1}{2}},\\
(\xi',\xi):=\mathrm{me}^3\sum_{c\bmod R\mathbb{Z}^3}\overline{\xi'(c)}\cdot\xi(c),\qquad
\|\xi\|_{L^2}:=(\xi,\xi)^{\frac{1}{2}}.
\end{gather*}
The summands are $R\mathbb{Z}^3$-periodic, because $\left|-\exp(i\theta_\mu)\right|=1$.
$\mathrm{HH}$ and $\mathrm{HH}_4$ denote the sets of the elements of $\mathfrak{S}^{\mathrm{ch}}$ and of $\mathfrak{S}$ with finite norm (all elements, if $R<\infty$);
they are Hilbert spaces.
$\|O\|_{L^2}$ denotes the operator norm.

\emph{Operators.}
For $f\in\mathfrak{S}^{\mathrm{ch}}$, $\xi\in\mathfrak{S}$ and $a\in\{\pm e_\mu\}$ let
\begin{gather*}
(R_i(A)f)(x):=U_A(x,x^{(i)})\,f(x^{(i)}),\qquad
(D_Af)(x):=-\mathrm{me}^{-1}\sum_{i=0}^{3}M_{i}(x)\left[U_A(x,x^{(i)})\,f(x^{(i)})-f(x)\right],\\
\nabla_A^{*}\nabla_A:=\sum_{i=0}^{3}w_i\left(1-R_i(A)\right),\qquad
(J_Af)(x):=\sum_{i=0}^{3}w_i\,\gamma\cdot\hat{d}_i(x)\,U_A(x,x^{(i)})\,f(x^{(i)}),\\
(J\xi)(c,g):=\xi(c),\qquad
(J^{*}f)(c):=\frac{1}{48}\sum_{g\in O_{\mathrm{cub}}}f(c,g),\qquad
\mathrm{Plt}:=JJ^{*},\qquad
\mathrm{Phv}:=1-\mathrm{Plt},\\
(\mathrm{dc}_{A,a}\xi)(c):=\mathrm{me}^{-1}\left(u_A(c,c+\mathrm{me}\cdot a)\,\xi(c+\mathrm{me}\cdot a)-\xi(c)\right),\qquad
\|\mathrm{dc}_A\xi\|^2:=\sum_{\mu=1}^{3}\|\mathrm{dc}_{A,e_\mu}\xi\|_{L^2}^2,\\
D^{\mathrm{ll}}_{A}:=J^{*}D_AJ,\qquad
D^{\mathrm{lh}}_{A}:=J^{*}D_A\,\mathrm{Phv},\qquad
D^{\mathrm{hl}}_{A}:=\mathrm{Phv}\,D_AJ,\qquad
D^{\mathrm{hh}}_{A}:=\mathrm{Phv}\,D_A\,\mathrm{Phv}.
\end{gather*}
All these maps respect the conditions of quasi-periodicity:
$x\mapsto x^{(i)}$ commutes with the translations, $M_{i}(x)$ and $\hat{d}_i(x)$ depend on the label only, and $u_A$ is periodic.
They are defined on the whole of $\mathfrak{S}^{\mathrm{ch}}$ and of $\mathfrak{S}$, without any condition of summability.

\subsubsection*{0.2. Instances (the three instances of Proposition~1.1\,(f))}

\begin{itemize}
\item[(In1)] \emph{The data of Proposition~1.1\,(a)--(e).}
Here $\mathrm{me}=\eta_k$, and either $\mathrm{Lab}$ and $R$ are as in Section~G1.2 ($R\ge3L$ a multiple of $L$), with $\theta=(0,0,0)$ if $\boldsymbol{\sigma}=-1$ and $\theta=(\pi,\pi,\pi)$ if $\boldsymbol{\sigma}=+1$,
or $\mathrm{Lab}=\mathbb{L}_k$ and $R=\infty$.
The objects of 0.1 are then those of Section~G1.2, of (I.2.2)--(I.2.4), (I.3.1), (I.J4), (I.J6) and of the data of Proposition~1.1.
$R$ is a multiple of $\mathrm{me}$, since $R/\eta_k=RL^k$ is an integer;
and $-\exp(i\theta_\mu)=\boldsymbol{\sigma}$ in both cases.
The level $k$ enters through $\eta_k$ only.
\emph{Level $0$:} the formulas are read with $\eta_0=1$, by (I.L0.1);
$e_k>0$ is one of the data of Proposition~1.1 at every level;
$\xi$ has $4$ components and $f$ has $192$ components per photon site, as at every level.
Nothing else is different at level~$0$.
\item[(In2)] \emph{Section~G1.4.}
$\mathrm{Lab}=\mathrm{me}\,\mathbb{Z}^3$ and $A=0$.
Then $\mathfrak{S}^{\mathrm{ch}}=\mathfrak{S}^{\mathrm{ch}}(\mathrm{me},R;\theta)$, $\mathrm{HH}=\mathrm{HH}(\mathrm{me},R;\theta)$, $D_0=D$, $\nabla_0^{*}\nabla_0=\nabla^{*}\nabla$ and $R_i(0)=R_i$, where $D$, $\nabla^{*}\nabla$ and $R_i$ are the operators of Section~G1.4.
\item[(In3)] \emph{The fibre $p$, $p\in\mathbb{R}^3$.}
$\mathrm{me}=1$, $\mathrm{Lab}=\mathbb{Z}^3$, $R=1$, $A=0$ and $\theta_\mu=p_\mu-\pi$.
The conditions of quasi-periodicity become $f(c+e_\mu,g)=\exp(ip_\mu)\,f(c,g)$ and $\xi(c+e_\mu)=\exp(ip_\mu)\,\xi(c)$.
Hence $f\in\mathfrak{S}^{\mathrm{ch}}$ if and only if $f(c,g)=\exp(ip\cdot c)\,u(g)$ with $u:=f(0,\cdot\,)\in{V_{\mathrm{v}}}$,
and $\xi\in\mathfrak{S}$ if and only if $\xi(c)=\exp(ip\cdot c)\,v$ with $v:=\xi(0)\in\mathbb{C}^4$.
Modulo $R\mathbb{Z}^3$ there is one photon site;
therefore $(f',f)^{\mathrm{ch}}=\frac{1}{48}\sum_g\overline{u'(g)}\cdot u(g)=:(u',u)$, $\|u\|:=(u,u)^{\frac{1}{2}}$, and $(\xi',\xi)=\overline{v'}\cdot v$.
For all $p\in\mathbb{C}^3$ define
\begin{gather*}
(R_iu)(g):=u(gs_i)\quad(i=0,1,2),\qquad
(R_3(p)u)(g):=\exp\left(ip\cdot(ge_1)\right)\,u(gs_3),\qquad
M_{i}(g):=M_{i}(c,g),\qquad
\hat{d}_i(g):=-g\hat n_i,\\
\left(D_0^{(p)}u\right)(g):=-\sum_{i\le2}M_{i}(g)\left(u(gs_i)-u(g)\right)-M_3(g)\left(\exp\left(ip\cdot(ge_1)\right)\,u(gs_3)-u(g)\right),\\
(\nabla^{*}\nabla)(p):=\sum_{i\le2}w_i(1-R_i)+w_3\left(1-R_3(p)\right),\qquad
(Jv)(g):=v,\qquad
J^{*}u:=\frac{1}{48}\sum_gu(g).
\end{gather*}
Under the identification $f\leftrightarrow u$, $\xi\leftrightarrow v$, the operator $R_i(0)$ becomes $R_i$ ($i\le2$),
$R_3(0)$ becomes $R_3(p)$ (because $f(x^{(3)})=\exp\left(ip\cdot(c+ge_1)\right)\,u(gs_3)$),
$D_0$ becomes $D_0^{(p)}$,
$\nabla_0^{*}\nabla_0$ becomes $(\nabla^{*}\nabla)(p)$,
and $J$, $J^{*}$, $\mathrm{Plt}$, $\mathrm{Phv}$ become the maps with the same names on $\mathbb{C}^4$ and ${V_{\mathrm{v}}}$.
The symbol $D^{(p)}$ of (I.2.6) at $\boldsymbol{\mu}=\mathbf{m}$ is $\mathbf{m}+D_0^{(p)}$.
Further $(\mathrm{dc}_{0,a}\xi)(c)=\exp(ip\cdot c)\left(\exp(ip\cdot a)-1\right)v$, so that
\begin{equation}\tag{1a.1}
\|\mathrm{dc}_0\xi\|^2=\sum_{\mu=1}^{3}\left|\exp(ip_\mu)-1\right|^2\cdot|v|_2^2=2\,\mathrm{tw}(p)\cdot|v|_2^2\qquad\left(\left|\exp(is)-1\right|^2=2-2\cos s\right).
\end{equation}
\end{itemize}

\subsubsection*{0.3. Facts}

The facts (GS2)--(GS8) and (GS10) below are assertions of Proposition~1.1\,(f);
(GS1), (GS9) and (GS11) are auxiliary.

\begin{named}{(GS1)}[complex $A$]
$R_i(A)^2=1$.
\end{named}
\begin{proof}
By (s1) and the identity $U_A(x^{(i)},x)=U_A(x,x^{(i)})^{-1}$ of 0.1,
$(R_i(A)R_i(A)f)(x)=U_A(x,x^{(i)})\,U_A(x^{(i)},x)\,f\left((x^{(i)})^{(i)}\right)=f(x)$.
\end{proof}

\begin{named}{(GS2)}[complex $A$; Proposition~1.1\,(f): (I.2.4) in (GS)]
\begin{equation*}
D_A=\mathrm{me}^{-1}\left[\frac{\rho}{2}\,\nabla_A^{*}\nabla_A+\frac{1}{2}\,J_A\right].
\end{equation*}
\end{named}
\begin{proof}
For each $i$,
\begin{equation*}
-M_{i}(x)\left[U_A(x,x^{(i)})f(x^{(i)})-f(x)\right]=\frac{\rho}{2}\,w_i\left(f-R_i(A)f\right)(x)+\frac{1}{2}\,w_i\,\gamma\cdot\hat{d}_i(x)\,(R_i(A)f)(x)-\frac{1}{2}\,w_i\,\gamma\cdot\hat{d}_i(x)\,f(x).
\end{equation*}
In the sum over $i$ the last term drops out, since $\sum_iw_i\hat{d}_i(x)=0$ by (s3).
\end{proof}

\begin{named}{(GS3)}[real $A$; Proposition~1.1\,(f)]
$R_i(A)$ is a unitary and self-adjoint operator on $\mathrm{HH}$.
\end{named}
\begin{proof}
One has $|U_A|=1$, and $x\mapsto x^{(i)}$ is a bijection of the classes of sites;
hence $\|R_i(A)f\|_{L^2}^2=\int\mathrm{d}x\,|f(x^{(i)})|_2^2=\|f\|_{L^2}^2$.
A map that preserves the norm and satisfies $R_i(A)^2=1$ (GS1) is unitary, with $R_i(A)^{*}=R_i(A)^{-1}=R_i(A)$.
\end{proof}

\begin{named}{(GS4)}[real $A$, $f\in\mathrm{HH}$; the energy; Proposition~1.1\,(f)]
For $i=0,\dots,3$,
\begin{equation}\tag{1a.2}
\begin{split}
\left(f,(1-R_i(A))f\right)^{\mathrm{ch}}
&=\frac{1}{2}\,\left\|(1-R_i(A))f\right\|_{L^2}^2
=\frac{1}{2}\cdot\frac{\mathrm{me}^3}{48}\sum_{x\bmod R\mathbb{Z}^3}\left|U_A(x,x^{(i)})f(x^{(i)})-f(x)\right|_2^2\\
&=\frac{\mathrm{me}^3}{48}\sum_{\{x,x^{(i)}\}\bmod R\mathbb{Z}^3}\left|U_A(x,x^{(i)})f(x^{(i)})-f(x)\right|_2^2,
\end{split}
\end{equation}
where the last sum extends over the unoriented bonds of type $i$.
Consequently $\nabla_A^{*}\nabla_A=(\nabla_A^{*}\nabla_A)^{*}\ge0$, and $(f,\nabla_A^{*}\nabla_Af)^{\mathrm{ch}}$ is the sum over $i$ of $w_i$ times (1a.2);
this is the formula of (I.O2).
\end{named}
\begin{proof}
By (GS3), $(f,R_i(A)f)^{\mathrm{ch}}$ is real, and $\|(1-R_i(A))f\|^2=2\|f\|^2-2(f,R_i(A)f)^{\mathrm{ch}}$.
The classes of sites fall into the pairs $\{x,x^{(i)}\}$ (0.1), and the summand takes the same value at $x$ and at $x^{(i)}$:
\begin{equation*}
\left|U_A(x^{(i)},x)f(x)-f(x^{(i)})\right|_2=\left|U_A(x^{(i)},x)\right|\cdot\left|f(x)-U_A(x,x^{(i)})f(x^{(i)})\right|_2.
\end{equation*}
\end{proof}

\begin{named}{(GS5)}[real $A$; Proposition~1.1\,(f)]
On $\mathrm{HH}$,
\begin{equation*}
J_A^{*}=-J_A,\qquad
\mathrm{Re}\,(f,D_Af)^{\mathrm{ch}}=\frac{\rho}{2\,\mathrm{me}}\,(f,\nabla_A^{*}\nabla_Af)^{\mathrm{ch}},\qquad
D_A^{*}=\gamma_5D_A\gamma_5,\qquad
\|D_A\|_{L^2}\le\frac{(1+\rho)\,\mathrm{sw}}{\mathrm{me}}.
\end{equation*}
\end{named}
\begin{proof}
Let $\mathrm{Mx}_i$ be the multiplication by the Hermitian matrix $\gamma\cdot\hat{d}_i(x)$.
Then $J_A=\sum_iw_i\,\mathrm{Mx}_iR_i(A)$, and by (s2)
\begin{equation*}
(R_i(A)\,\mathrm{Mx}_if)(x)=U_A(x,x^{(i)})\,\gamma\cdot\hat{d}_i(x^{(i)})\,f(x^{(i)})=-(\mathrm{Mx}_iR_i(A)f)(x).
\end{equation*}
So $(\mathrm{Mx}_iR_i(A))^{*}=R_i(A)\,\mathrm{Mx}_i=-\mathrm{Mx}_iR_i(A)$.
Together with (GS2) and (GS4) this gives $D_A^{*}=\mathrm{me}^{-1}\left[\frac{\rho}{2}\,\nabla_A^{*}\nabla_A-\frac{1}{2}\,J_A\right]$, and the formula for the real part follows.
The matrix $\gamma_5$ commutes with $R_i(A)$ and anticommutes with $\mathrm{Mx}_i$ (c1), and $\gamma_5^2=1$;
hence $\gamma_5\nabla_A^{*}\nabla_A\gamma_5=\nabla_A^{*}\nabla_A$ and $\gamma_5J_A\gamma_5=-J_A$.
For the norm: $|M_{i}(x)|_2=\frac{w_i}{2}(1+\rho)$ by (m6), so
\begin{equation*}
\|D_Af\|\le\mathrm{me}^{-1}\sum_i\frac{w_i}{2}(1+\rho)\left(\|R_i(A)f\|+\|f\|\right)=\mathrm{me}^{-1}(1+\rho)\,\mathrm{sw}\,\|f\|.
\end{equation*}
\end{proof}

\begin{named}{(GS6)}[cluster maps; Proposition~1.1\,(f)]
One has $J^{*}J=1$;
$(f,J\xi)^{\mathrm{ch}}=(J^{*}f,\xi)$ for $f\in\mathrm{HH}$ and $\xi\in\mathrm{HH}_4$;
$\|J\xi\|_{L^2}=\|\xi\|_{L^2}$;
$\mathrm{Plt}$ and $\mathrm{Phv}$ are orthogonal projections of $\mathrm{HH}$;
$(\mathrm{Phv}\,f)(c,g)=f(c,g)-(J^{*}f)(c)$, $\sum_g(\mathrm{Phv}\,f)(c,g)=0$, $J^{*}\mathrm{Phv}=0$, $\mathrm{Phv}\,J=0$;
and for every $c$
\begin{equation}\tag{1a.3}
\frac{1}{48}\sum_g|f(c,g)|_2^2=|(J^{*}f)(c)|_2^2+\frac{1}{48}\sum_g|(\mathrm{Phv}\,f)(c,g)|_2^2.
\end{equation}
\end{named}
\begin{proof}
First, $\|J\xi\|_{L^2}^2=\mathrm{me}^3\sum_c\frac{1}{48}\sum_g|\xi(c)|_2^2=\|\xi\|_{L^2}^2$.
The identity (1a.3) concerns a single photon site and requires no summability;
it is verified at the end of this proof.
Multiplying it by $\mathrm{me}^3$ and summing over $c$ gives $\|J^{*}f\|_{L^2}^2+\|\mathrm{Phv}\,f\|_{L^2}^2=\|f\|_{L^2}^2$.
Hence $J$ maps $\mathrm{HH}_4$ into $\mathrm{HH}$, while $J^{*}$ and $\mathrm{Phv}$ map $\mathrm{HH}$ into $\mathrm{HH}_4$ and $\mathrm{HH}$, respectively.
Next,
\begin{equation*}
(f,J\xi)^{\mathrm{ch}}=\mathrm{me}^3\sum_c\frac{1}{48}\sum_g\overline{f(c,g)}\cdot\xi(c)=\mathrm{me}^3\sum_c\overline{(J^{*}f)(c)}\cdot\xi(c);
\end{equation*}
so $J^{*}$ is the adjoint of $J$.
Since moreover $J^{*}J=1$, the operator $\mathrm{Plt}=JJ^{*}$ is self-adjoint and idempotent.
Finally, (1a.3): $f(c,g)=(J^{*}f)(c)+(\mathrm{Phv}\,f)(c,g)$, and the mean over $g$ of the cross term vanishes.
\end{proof}

\begin{named}{(GS7)}[the cluster Laplacian; Proposition~1.1\,(f): the bounds on $\lambda_{\mathrm{int}}$]
On ${V_{\mathrm{v}}}$, with the inner product $(u',u)$ of (In3), put $\mathrm{Lcl}:=\sum_{i\le2}w_i(1-R_i)$.
Then
\begin{equation}\tag{1a.4}
\begin{aligned}
&(u,\mathrm{Lcl}\,u)=\sum_{i\le2}\frac{w_i}{2}\cdot\frac{1}{48}\sum_g\left|u(gs_i)-u(g)\right|_2^2,\qquad
\mathrm{Lcl}\,J=0,\\
&(u,\mathrm{Lcl}\,u)\ge\lambda_{\mathrm{int}}\,\|\mathrm{Phv}\,u\|^2\quad\text{for all }u\in{V_{\mathrm{v}}},\qquad
0<\lambda_{\mathrm{int}}\le12+24\sqrt{2}<48.
\end{aligned}
\end{equation}
\end{named}
\begin{proof}
The map $g\mapsto gs_i$ is an involution of $O_{\mathrm{cub}}$.
Hence $R_i$ is a unitary operator on ${V_{\mathrm{v}}}$ with $R_i^2=1$, thus self-adjoint;
$(u,R_iu)$ is real, and $\|(1-R_i)u\|^2=2\|u\|^2-2(u,R_iu)=2\left(u,(1-R_i)u\right)$.
This is the first formula.
Since $(R_iJv)(g)=v$, one has $\mathrm{Lcl}\,J=0$.
On each of the four components of $u$, $\mathrm{Lcl}$ acts by the real symmetric $48\times48$ matrix $\sum_{i\le2}w_i(1-R_i)$ on $\ell^2(O_{\mathrm{cub}})$ of (I.O8).
If $(u,\mathrm{Lcl}\,u)=0$, then $u(gs_i)=u(g)$ for all $g$ and $i=0,1,2$;
so $u$ is constant on the orbits of the right multiplication by the group generated by $s_0$, $s_1$, $s_2$.
This group is $O_{\mathrm{cub}}$:
$s_1$ and $s_2$ are the transpositions $(23)$ and $(12)$, which generate the six permutation matrices;
the conjugates of $s_0=\mathrm{diag}(1,1,-1)$ by these are $\mathrm{diag}(1,-1,1)$ and $\mathrm{diag}(-1,1,1)$, and the three diagonal matrices generate $\{\pm1\}^3$;
every signed permutation matrix is the product of an element of $\{\pm1\}^3$ and a permutation matrix.
Hence the kernel of the matrix consists of the constants, and its smallest non-zero eigenvalue $\lambda_{\mathrm{int}}$ is positive.
By the spectral theorem for real symmetric matrices \cite[Section~4.1]{HornJohnson2013} there is an orthonormal basis of eigenvectors, and here the orthogonal projection onto the kernel is the mean over $O_{\mathrm{cub}}$;
therefore $(u_\beta,\mathrm{Lcl}\,u_\beta)\ge\lambda_{\mathrm{int}}\,\|u_\beta-J^{*}u_\beta\|^2$ for each component $u_\beta$.
Summation over $\beta$ gives the inequality.
An eigenvalue is at most the norm:
$\lambda_{\mathrm{int}}\le\sum_{i\le2}w_i\,\|1-R_i\|\le2(w_0+w_1+w_2)=12+24\sqrt{2}$, and $\sqrt{2}<\frac{3}{2}$.
\end{proof}

\begin{named}{(GS8)}[octets; Proposition~1.1\,(f): (I.s4) in (GS)]
For $c\in\mathrm{Lab}$ and $a\in\{\pm e_\mu\}$ put $\mathrm{Oc}_a(c):=\{(c,g):ge_1=a\}$.
Then $\mathrm{Oc}_a(c)$ has $8$ elements;
the six sets $\mathrm{Oc}_a(c)$ partition the cluster of $c$;
$x\mapsto x^{(3)}$ maps $\mathrm{Oc}_a(c)$ onto $\mathrm{Oc}_{-a}(c+\mathrm{me}\cdot a)$;
and the $8$ bonds $\{x,x^{(3)}\}$, $x\in\mathrm{Oc}_a(c)$, are the bonds of type $3$ that join the clusters of $c$ and $c+\mathrm{me}\cdot a$.
This is (s4).
\end{named}
\begin{proof}
The vector $ge_1$ is the first column of $g$.
A signed permutation matrix whose first column is a given vector $\pm e_\mu$ is determined by a signed permutation of the two other coordinates, and there are $2\cdot4=8$ of these.
Moreover $(gs_3)e_1=-ge_1$.
Since $ge_1\in\{\pm e_\mu\}$ for every $g$, the six sets partition the cluster.
For $x=(c,g)\in\mathrm{Oc}_a(c)$ the site $x^{(3)}=(c+\mathrm{me}\cdot a,\,gs_3)$ lies in $\mathrm{Oc}_{-a}(c+\mathrm{me}\cdot a)$.
The same computation shows that $x\mapsto x^{(3)}$ maps $\mathrm{Oc}_{-a}(c+\mathrm{me}\cdot a)$ into $\mathrm{Oc}_a(c)$;
as $(x^{(3)})^{(3)}=x$ by (s1), it is a bijection between the two octets.
A bond of type $3$ that joins the clusters of $c$ and $c+\mathrm{me}\cdot a$ has one end $x$ over $c$, and $c_{3}(x)=c+\mathrm{me}\cdot a$ means $ge_1=a$.
\end{proof}

\begin{named}{(GS9)}[covariant differences; complex $A$, $\xi\in\mathfrak{S}$]
For $c\in\mathrm{Lab}$, $a\in\{\pm e_\mu\}$ and $c':=c+\mathrm{me}\cdot a$,
\begin{equation}\tag{1a.5}
(\mathrm{dc}_{A,-a}\xi)(c')=-u_A(c',c)\,(\mathrm{dc}_{A,a}\xi)(c).
\end{equation}
Consequently, for real $A$, $|(\mathrm{dc}_{A,-a}\xi)(c')|_2=|(\mathrm{dc}_{A,a}\xi)(c)|_2$;
this number is a function of the unoriented photon bond $\{c,c'\}$;
every unoriented photon bond is $\{c,c+\mathrm{me}\cdot e_\mu\}$ for exactly one pair $(c,\mu)$;
and $\sum_a\|\mathrm{dc}_{A,a}\xi\|_{L^2}^2=2\,\|\mathrm{dc}_A\xi\|^2$ for $\xi\in\mathrm{HH}_4$.
\end{named}
\begin{proof}
Since $u_A(c',c)\,u_A(c,c')=1$,
\begin{equation*}
(\mathrm{dc}_{A,-a}\xi)(c')=\mathrm{me}^{-1}\left(u_A(c',c)\,\xi(c)-\xi(c')\right)=-u_A(c',c)\,\mathrm{me}^{-1}\left(u_A(c,c')\,\xi(c')-\xi(c)\right).
\end{equation*}
The translation by $\mathrm{me}\cdot e_\mu$ permutes the classes of photon sites;
hence $\|\mathrm{dc}_{A,-e_\mu}\xi\|_{L^2}=\|\mathrm{dc}_{A,e_\mu}\xi\|_{L^2}$.
\end{proof}

\begin{named}{(GS10)}[upper bound of the energy; Proposition~1.1\,(f); real $A$, $f\in\mathrm{HH}$, $\xi:=J^{*}f$, $\mathrm{hp}:=\mathrm{Phv}\,f$]
\begin{equation}\tag{1a.6}
(f,\nabla_A^{*}\nabla_Af)^{\mathrm{ch}}\le3\,\mathrm{me}^2\cdot\|\mathrm{dc}_A\xi\|^2+\left(30+24\sqrt{2}\right)\|\mathrm{hp}\|_{L^2}^2.
\end{equation}
\end{named}
\begin{proof}
Write $f=J\xi+\mathrm{hp}$.
For $i\le2$ one has $f(x^{(i)})-f(x)=\mathrm{hp}(x^{(i)})-\mathrm{hp}(x)$ and $|\mathrm{hp}(x^{(i)})-\mathrm{hp}(x)|_2^2\le2|\mathrm{hp}(x^{(i)})|_2^2+2|\mathrm{hp}(x)|_2^2$;
by the second expression in (1a.2) and the bijectivity of $x\mapsto x^{(i)}$,
\begin{equation*}
\left(f,(1-R_i(A))f\right)^{\mathrm{ch}}\le\frac{1}{2}\cdot\frac{\mathrm{me}^3}{48}\cdot4\sum_x|\mathrm{hp}(x)|_2^2=2\,\|\mathrm{hp}\|^2.
\end{equation*}
For $i=3$, $x=(c,g)$ and $a=ge_1$,
\begin{equation*}
U_A(x,x^{(3)})f(x^{(3)})-f(x)=\mathrm{me}\cdot(\mathrm{dc}_{A,a}\xi)(c)+u_A(c,c_{3})\,\mathrm{hp}(x^{(3)})-\mathrm{hp}(x),
\end{equation*}
whose square is at most $3\,\mathrm{me}^2\cdot|(\mathrm{dc}_{A,a}\xi)(c)|_2^2+3|\mathrm{hp}(x^{(3)})|_2^2+3|\mathrm{hp}(x)|_2^2$.
By (GS8) each $a$ occurs for $8$ labels; hence
\begin{align*}
\left(f,(1-R_3(A))f\right)^{\mathrm{ch}}
&\le\frac{1}{2}\cdot\frac{\mathrm{me}^3}{48}\left[3\,\mathrm{me}^2\cdot8\sum_c\sum_a|(\mathrm{dc}_{A,a}\xi)(c)|_2^2+6\sum_x|\mathrm{hp}(x)|_2^2\right]\\
&=\frac{\mathrm{me}^2}{4}\sum_a\|\mathrm{dc}_{A,a}\xi\|_{L^2}^2+3\,\|\mathrm{hp}\|^2
=\frac{\mathrm{me}^2}{2}\,\|\mathrm{dc}_A\xi\|^2+3\,\|\mathrm{hp}\|^2,
\end{align*}
the last step by (GS9).
Multiplication by the weights gives (1a.6):
$w_3\cdot\frac{\mathrm{me}^2}{2}=3\,\mathrm{me}^2$ and $3w_3+2(w_0+w_1+w_2)=18+12+24\sqrt{2}$.
\end{proof}

\begin{named}{(GS11)}[operators with positive real part]
Let $O$ be a bounded operator on a Hilbert space with $\mathrm{Re}\,(f,Of)\ge\delta\,\|f\|^2$ for all $f$, where $\delta>0$.
Then $O$ is bijective and $\|O^{-1}\|\le1/\delta$.
\end{named}
\begin{proof}
One has $\delta\,\|f\|^2\le|(f,Of)|\le\|f\|\,\|Of\|$;
so $\|Of\|\ge\delta\,\|f\|$, $O$ is injective, and its range is closed.
Since $\mathrm{Re}\,(f,O^{*}f)=\mathrm{Re}\,(f,Of)$, the operator $O^{*}$ is injective too.
The orthogonal complement of the range of $O$ is the kernel of $O^{*}$, that is $\{0\}$;
and a closed subspace whose orthogonal complement is $\{0\}$ is the whole space, by the projection theorem \cite[Chapter~4]{Rudin1987}.
\end{proof}

\subsection*{1. Proposition G; Proposition 1.1}

Proposition~G states, for all data (GS), the clauses (a)--(d) of Proposition~1.1 with $\mathrm{me}$ in the place of $\eta_k$ (in (d) with the two forms of the set $E$), and the clause (e) without the maps $\mathrm{mS}$ of Section~G1.2.
In the instance (In1) it gives Proposition~1.1\,(a)--(d), and of Proposition~1.1\,(e) only the part on the multiplication by functions of $c$, the translations of $\mathrm{Lab}$ and the constant spin matrices.
The commutation of $J$, $J^{*}$, $\mathrm{Plt}$, $\mathrm{Phv}$ with the maps $\mathrm{mS}$ is not an instance of Proposition~G;
it is proved by a separate computation in 1.3.
Together with the facts of Section~0, Proposition~G gives Proposition~1.1\,(f) (see 1.3).

\subsubsection*{1.1. Constants; local quantities}

\emph{Constants.}
The constants $\lambda_{\mathrm{int}}$, $\hat{g}$, $c_{\mathrm{hl}}$, $c_{*}$ are those of Section~I.8.2, line (C0r), with the formulas of (I.O8) and of Proposition~1.1\,(b), (c), (d):
$\lambda_{\mathrm{int}}$ is the smallest non-zero eigenvalue of $\sum_{i\le2}w_i(1-R_i)$ on $\ell^2(O_{\mathrm{cub}})$,
$\hat{g}=\rho\,\lambda_{\mathrm{int}}/2$,
$c_{\mathrm{hl}}=\sqrt{3}\,(1+\rho)$
and $c_{*}=\min(w_3,\lambda_{\mathrm{int}})/18$.
The number $\lambda_{\mathrm{int}}$ is treated in (GS7).

\emph{Local quantities.}
Let the data be (GS) with real $A$.
A set $E$ of photon sites and a function $f$ are taken in one of the following two readings.
\begin{itemize}
\item[(Es1)] $E$ is an $R\mathbb{Z}^3$-periodic subset of $\mathrm{Lab}$ (an arbitrary subset if $R=\infty$), $f\in\mathrm{HH}$, and all sums over sites and bonds are taken modulo $R\mathbb{Z}^3$.
($E$ is a set of photon sites of the torus.)
\item[(Es2)] $E$ is a finite subset of $\mathrm{Lab}$, $f$ is an arbitrary function $E\times O_{\mathrm{cub}}\to\mathbb{C}^4$, and all sums extend over the sites and bonds of $\mathrm{Lab}$ themselves.
\end{itemize}
A site $x=(c,g)$ lies over $E$ if $c\in E$.
In both readings $J^{*}f$ and $\mathrm{Phv}\,f$ are formed over the sites of $E$ by the displays of 0.1:
$(J^{*}f)(c)=\frac{1}{48}\sum_gf(c,g)$ and $(\mathrm{Phv}\,f)(c,g)=f(c,g)-(J^{*}f)(c)$ for $c\in E$.
With $\xi:=J^{*}f$ put
\begin{align*}
\|\mathrm{dc}_A\xi\|_E^2&:=\mathrm{me}^3\sum\left|(\mathrm{dc}_{A,a}\xi)(c)\right|_2^2,\\
(f,\nabla_A^{*}\nabla_Af)_E&:=\frac{\mathrm{me}^3}{48}\sum_{i=0}^{3}w_i\sum\left|U_A(x,x^{(i)})f(x^{(i)})-f(x)\right|_2^2,\\
\|f\|_{L^2,E}^2&:=\frac{\mathrm{me}^3}{48}\sum_{x\text{ over }E}|f(x)|_2^2.
\end{align*}
In the first line the sum extends over the unoriented photon bonds $\{c,c+\mathrm{me}\cdot a\}$ with both ends in $E$;
by (GS9) the summand is a function of the unoriented bond.
In the second line the inner sum extends over the unoriented bonds $\{x,x^{(i)}\}$ of type $i$ with both ends over $E$;
as in the proof of (GS4), the summand takes the same value at $x$ and at $x^{(i)}$.

The identity (1a.5), the equality $\left|U_A(x^{(i)},x)f(x)-f(x^{(i)})\right|_2=\left|f(x)-U_A(x,x^{(i)})f(x^{(i)})\right|_2$ from the proof of (GS4), and the identity (1a.3) involve only the values of $\xi$, respectively of $f$, at the two ends of one bond, respectively over one photon site;
hence they hold in the reading (Es2) as well.
In the reading (Es1) the summands are $R\mathbb{Z}^3$-periodic:
a translation by $R\,e_\mu$ multiplies $f$, $\xi$ and $\mathrm{Phv}\,f$ by the number $-\exp(i\theta_\mu)$ of modulus one, and $u_A$ is periodic.
The three quantities depend only on the values of $f$ on $E\times O_{\mathrm{cub}}$.
One has
\begin{equation*}
(f,\nabla_A^{*}\nabla_Af)_E=(f,\nabla_A^{*}\nabla_Af)_{E,3}+(f,\nabla_A^{*}\nabla_Af)_{E,\mathrm{int}},
\end{equation*}
the contributions of the bonds of type $3$ and of the bonds of the types $0$, $1$, $2$.

In the reading (Es1), by (1a.2), $(f,\nabla_A^{*}\nabla_Af)_E$ is the contribution to $(f,\nabla_A^{*}\nabla_Af)^{\mathrm{ch}}$ of the bonds with both ends over $E$.
For $E=\mathrm{Lab}$ in the reading (Es1),
\begin{equation*}
\|\mathrm{dc}_A\xi\|_E=\|\mathrm{dc}_A\xi\|\ \text{ by (GS9)},\qquad
(f,\nabla_A^{*}\nabla_Af)_E=(f,\nabla_A^{*}\nabla_Af)^{\mathrm{ch}}\ \text{ by (GS4)},\qquad
\|f\|_{L^2,E}=\|f\|_{L^2}.
\end{equation*}

\subsubsection*{1.2. The proposition}

The next proposition is the part of Proposition~1.1\,(f) that concerns the clauses (a)--(e):
the clauses (a)--(d) with $\mathrm{me}$ in the place of $\eta_k$,
the clause (e) without the maps $\mathrm{mS}$,
and the two forms of the set $E$ in (d).

\begin{named}{Proposition G}
Let the data be (GS).
\begin{itemize}
\item[(a)] For complex $A$ and $\xi\in\mathfrak{S}$,
\[
(J^{*}D_AJ\xi)(c)=-\mathrm{me}^{-1}\sum_a\frac{1}{2}\left(\rho-\gamma\cdot a\right)\left[u_A(c,c+\mathrm{me}\,a)\,\xi(c+\mathrm{me}\,a)-\xi(c)\right],
\]
where $a$ ranges over the six vectors $\pm e_1,\pm e_2,\pm e_3$.
\item[(b)] For complex $A$ and $\xi\in\mathfrak{S}$,
\[
(D_AJ\xi)(c,g)=-M_3(x)\,(\mathrm{dc}_{A,ge_1}\xi)(c),\qquad x=(c,g).
\]
For real $A$, $\xi\in\mathrm{HH}_4$ and $f\in\mathrm{HH}$,
\[
\|D^{\mathrm{hl}}_A\xi\|_{L^2}\le c_{\mathrm{hl}}\,\|\mathrm{dc}_A\xi\|,
\qquad
\left|(\xi,D^{\mathrm{lh}}_Af)\right|\le c_{\mathrm{hl}}\,\|\mathrm{dc}_A\xi\|\cdot\|f\|_{L^2}.
\]
\item[(c)] For real $A$ and $f=\mathrm{Phv}\,f$ in $\mathrm{HH}$,
\[
\mathrm{Re}\,(f,D_Af)^{\mathrm{ch}}\ge\frac{\hat g}{\mathrm{me}}\,(f,f)^{\mathrm{ch}}.
\]
For every real $t>-\hat g/\mathrm{me}$ the operator $D^{\mathrm{hh}}_A+t$ is invertible on the range of $\mathrm{Phv}$ in $\mathrm{HH}$, and
\[
\left\|\left(D^{\mathrm{hh}}_A+t\right)^{-1}\right\|_{L^2}\le\left(t+\hat g/\mathrm{me}\right)^{-1}.
\]
\item[(d)] Let $A$ be real, let $(E,f)$ be taken in the reading (Es1) or in the reading (Es2), and put $\xi:=J^{*}f$.
Then
\[
\mathrm{me}^{-2}\,(f,\nabla_A^{*}\nabla_Af)_E\ge c_{*}\,\|\mathrm{dc}_A\xi\|_E^2,
\qquad
(f,\nabla_A^{*}\nabla_Af)_E\ge\lambda_{\mathrm{int}}\,\|\mathrm{Phv}\,f\|_{L^2,E}^2.
\]
In particular, taking $E=\mathrm{Lab}$ in the reading (Es1), one has for every $f\in\mathrm{HH}$
\[
\mathrm{me}^{-2}\,(f,\nabla_A^{*}\nabla_Af)^{\mathrm{ch}}\ge c_{*}\,\|\mathrm{dc}_A\xi\|^2,
\qquad
(f,\nabla_A^{*}\nabla_Af)^{\mathrm{ch}}\ge\lambda_{\mathrm{int}}\,\|\mathrm{Phv}\,f\|_{L^2}^2.
\]
\item[(e)] On arbitrary functions $\xi:\mathrm{Lab}\to\mathbb{C}^4$ and $f:\mathrm{Lab}\times O_{\mathrm{cub}}\to\mathbb{C}^4$ the maps $J$, $J^{*}$, $\mathrm{Plt}$, $\mathrm{Phv}$ commute with the multiplication by functions of $c$, with the translations by the vectors of $\mathrm{me}\,\mathbb{Z}^3$, and with constant $4\times4$ matrices.
\end{itemize}
The numbers $c_{\mathrm{hl}}$, $\hat g$, $c_{*}$, $\lambda_{\mathrm{int}}$ are functions of $\rho$ and of the weights.
They depend neither on the data $(\mathrm{me},\mathrm{Lab},R,\theta,e_k,A)$ nor on $E$.
\end{named}

\begin{proof}
\emph{(a) and the first identity of (b); complex $A$, $\xi\in\mathfrak{S}$.}
Let $x=(c,g)$.
If $i\in\{0,1,2\}$, the site $x^{(i)}$ lies over $c$ and $U_A(x,x^{(i)})=1$;
hence $U_A(x,x^{(i)})\,(J\xi)(x^{(i)})-(J\xi)(x)=\xi(c)-\xi(c)=0$.
For $i=3$ one has $(J\xi)(x^{(3)})=\xi(c_3)$ with $c_3=c+\mathrm{me}\,ge_1$, and $U_A(x,x^{(3)})=u_A(c,c_3)$.
Therefore
\begin{equation}
\begin{aligned}
(D_AJ\xi)(c,g)&=-\mathrm{me}^{-1}M_3(x)\left[u_A(c,c+\mathrm{me}\,ge_1)\,\xi(c+\mathrm{me}\,ge_1)-\xi(c)\right]\\
&=-M_3(x)\,(\mathrm{dc}_{A,ge_1}\xi)(c),
\qquad
M_3(x)=3\left(\rho-\gamma\cdot(ge_1)\right).
\end{aligned}
\tag{1a.7}
\end{equation}
This is the first identity of (b).
To average (1a.7) over the $48$ labels, group together the labels $g$ with $ge_1=a$;
by (GS8) each of the six vectors $a$ arises from exactly $8$ labels, and so
\begin{align*}
(J^{*}D_AJ\xi)(c)&=-\frac{1}{48}\sum_a8\cdot3\left(\rho-\gamma\cdot a\right)(\mathrm{dc}_{A,a}\xi)(c)\\
&=-\sum_a\frac{1}{2}\left(\rho-\gamma\cdot a\right)\mathrm{me}^{-1}\left[u_A(c,c+\mathrm{me}\,a)\,\xi(c+\mathrm{me}\,a)-\xi(c)\right],
\end{align*}
which is (a).
Both identities hold pointwise, between functions on $\mathrm{Lab}\times O_{\mathrm{cub}}$, resp.\ on $\mathrm{Lab}$, and no summability enters;
for $R=\infty$ and unbounded $\mathrm{Im}\,A$ the function $D_AJ\xi$ need not be square summable.

\emph{(b), the two bounds; real $A$, $\xi\in\mathrm{HH}_4$, $f\in\mathrm{HH}$.}
By (1a.7), by $|M_3(x)|_2=3(1+\rho)$ (see (I.m6)), by (GS8) and by (GS9),
\begin{equation}
\begin{aligned}
\|D_AJ\xi\|_{L^2}^2&=\mathrm{me}^3\,\frac{1}{48}\sum_c\sum_g\left|M_3(x)\,(\mathrm{dc}_{A,ge_1}\xi)(c)\right|_2^2\\
&\le9(1+\rho)^2\,\mathrm{me}^3\,\frac{8}{48}\sum_c\sum_a\left|(\mathrm{dc}_{A,a}\xi)(c)\right|_2^2\\
&=9(1+\rho)^2\cdot\frac{1}{6}\cdot2\,\|\mathrm{dc}_A\xi\|^2=c_{\mathrm{hl}}^2\,\|\mathrm{dc}_A\xi\|^2.
\end{aligned}
\tag{1a.8}
\end{equation}
As $\mathrm{Phv}$ is an orthogonal projection (GS6),
$\|D^{\mathrm{hl}}_A\xi\|_{L^2}=\|\mathrm{Phv}\,D_AJ\xi\|_{L^2}\le\|D_AJ\xi\|_{L^2}\le c_{\mathrm{hl}}\,\|\mathrm{dc}_A\xi\|$.
For the second bound one uses (GS6), (GS5) and the fact that the constant matrix $\gamma_5$ commutes with $J$:
\[
(\xi,D^{\mathrm{lh}}_Af)=(\xi,J^{*}D_A\,\mathrm{Phv}\,f)=(J\xi,D_A\,\mathrm{Phv}\,f)^{\mathrm{ch}}=(D_A^{*}J\xi,\mathrm{Phv}\,f)^{\mathrm{ch}}=(\gamma_5D_AJ\gamma_5\xi,\mathrm{Phv}\,f)^{\mathrm{ch}}.
\]
Since $\gamma_5$ is unitary and commutes with $\mathrm{dc}_{A,a}$, (1a.8) yields
$\|\gamma_5D_AJ\gamma_5\xi\|_{L^2}=\|D_AJ(\gamma_5\xi)\|_{L^2}\le c_{\mathrm{hl}}\,\|\mathrm{dc}_A(\gamma_5\xi)\|=c_{\mathrm{hl}}\,\|\mathrm{dc}_A\xi\|$,
and therefore
$\left|(\xi,D^{\mathrm{lh}}_Af)\right|\le c_{\mathrm{hl}}\,\|\mathrm{dc}_A\xi\|\cdot\|\mathrm{Phv}\,f\|_{L^2}\le c_{\mathrm{hl}}\,\|\mathrm{dc}_A\xi\|\cdot\|f\|_{L^2}$.

\emph{The internal part of the energy; the second inequality of (d); real $A$, $(E,f)$ in the reading (Es1) or (Es2).}
For a photon site $c\in E$ let $f(c,\cdot\,)\in{V_{\mathrm{v}}}$ be the function $g\mapsto f(c,g)$.
For $i\le2$ the unoriented bonds of type $i$ over $c$ are the $24$ pairs $\{(c,g),(c,gs_i)\}$;
on each of them $U_A=1$, and a sum over the $48$ labels $g$ meets each of them twice.
By (1a.4), the contribution of the site $c$ to $(f,\nabla_A^{*}\nabla_Af)_{E,\mathrm{int}}$ is therefore
\begin{equation}
\begin{split}
\frac{\mathrm{me}^3}{48}\sum_{i\le2}w_i\,\frac{1}{2}\sum_g\left|f(c,gs_i)-f(c,g)\right|_2^2
&=\mathrm{me}^3\left(f(c,\cdot\,),\mathrm{Lcl}\,f(c,\cdot\,)\right)\\
&\ge\lambda_{\mathrm{int}}\cdot\mathrm{me}^3\,\frac{1}{48}\sum_g\left|(\mathrm{Phv}\,f)(c,g)\right|_2^2;
\end{split}
\tag{1a.9}
\end{equation}
here $\mathrm{Phv}$ on ${V_{\mathrm{v}}}$ is the map $\mathrm{Phv}$ of Section~0.1 over the site $c$.
A bond of type $0$, $1$ or $2$ lies over a single photon site.
Summation of (1a.9) over $c\in E$ therefore gives
\begin{equation}
(f,\nabla_A^{*}\nabla_Af)_E\ge(f,\nabla_A^{*}\nabla_Af)_{E,\mathrm{int}}\ge\lambda_{\mathrm{int}}\,\|\mathrm{Phv}\,f\|_{L^2,E}^2.
\tag{1a.10}
\end{equation}

\emph{(c); real $A$, $f=\mathrm{Phv}\,f\in\mathrm{HH}$.}
By (GS5), and by (1a.10) for $E=\mathrm{Lab}$ in the reading (Es1), in which $(f,\nabla_A^{*}\nabla_Af)_E=(f,\nabla_A^{*}\nabla_Af)^{\mathrm{ch}}$ (Section~1.1),
\[
\mathrm{Re}\,(f,D_Af)^{\mathrm{ch}}=\frac{\rho}{2\,\mathrm{me}}\,(f,\nabla_A^{*}\nabla_Af)^{\mathrm{ch}}\ge\frac{\rho}{2\,\mathrm{me}}\,\lambda_{\mathrm{int}}\,\|\mathrm{Phv}\,f\|_{L^2}^2=\frac{\hat g}{\mathrm{me}}\,(f,f)^{\mathrm{ch}}.
\]
The range of $\mathrm{Phv}$ is closed in $\mathrm{HH}$.
The operator $D^{\mathrm{hh}}_A=\mathrm{Phv}\,D_A\,\mathrm{Phv}$ maps this range into itself and is bounded (GS5),
and for $f=\mathrm{Phv}\,f$ one has $(f,D^{\mathrm{hh}}_Af)^{\mathrm{ch}}=(\mathrm{Phv}\,f,D_Af)^{\mathrm{ch}}=(f,D_Af)^{\mathrm{ch}}$.
Consequently $\mathrm{Re}\,\left(f,(D^{\mathrm{hh}}_A+t)f\right)^{\mathrm{ch}}\ge\left(t+\hat g/\mathrm{me}\right)(f,f)^{\mathrm{ch}}$ on the range of $\mathrm{Phv}$,
and (GS11) with $\delta=t+\hat g/\mathrm{me}>0$ yields the invertibility together with $\left\|\left(D^{\mathrm{hh}}_A+t\right)^{-1}\right\|_{L^2}\le\left(t+\hat g/\mathrm{me}\right)^{-1}$.

\emph{(d), the first inequality; real $A$, $(E,f)$ in the reading (Es1) or (Es2), $\xi:=J^{*}f$, $\mathrm{hp}:=\mathrm{Phv}\,f$.}
Let $\{c,c'\}$ be an unoriented photon bond with both ends in $E$, and $c'=c+\mathrm{me}\,a$.
By (GS8) the map $x\mapsto x^{(3)}$ takes $\mathrm{Oc}_a(c)$ onto $\mathrm{Oc}_{-a}(c')$;
consequently $\frac{1}{8}\sum_{x\in\mathrm{Oc}_a(c)}f(x^{(3)})=\frac{1}{8}\sum_{x'\in\mathrm{Oc}_{-a}(c')}f(x')$.
Let $\mathrm{ub}:=u_A(c,c')$; this number equals $U_A(x,x^{(3)})$ for every $x\in\mathrm{Oc}_a(c)$.
Then
\begin{multline*}
\mathrm{ub}\,\xi(c')-\xi(c)=\frac{1}{8}\sum_{x\in\mathrm{Oc}_a(c)}\left[\mathrm{ub}\,f(x^{(3)})-f(x)\right]\\
+\mathrm{ub}\,\frac{1}{8}\sum_{x'\in\mathrm{Oc}_{-a}(c')}\left[\xi(c')-f(x')\right]-\frac{1}{8}\sum_{x\in\mathrm{Oc}_a(c)}\left[\xi(c)-f(x)\right],
\end{multline*}
because on the right side the two means of $f$ over $\mathrm{Oc}_{-a}(c')$ cancel, as do the two means of $f$ over $\mathrm{Oc}_a(c)$.
For $x$ over $c$ one has $\xi(c)-f(x)=-\mathrm{hp}(x)$, and for $x'$ over $c'$ one has $\xi(c')-f(x')=-\mathrm{hp}(x')$ (Section~1.1).
By the inequality $|\mathrm{vv}_1+\mathrm{vv}_2+\mathrm{vv}_3|_2^2\le3\sum_j|\mathrm{vv}_j|_2^2$,
by the Cauchy--Schwarz inequality in the form $\left|\frac{1}{8}\sum_{j=1}^{8}\mathrm{vv}_j\right|_2^2\le\frac{1}{8}\sum_j|\mathrm{vv}_j|_2^2$ (for $\mathrm{vv}_j\in\mathbb{C}^4$),
and by $|\mathrm{ub}|=1$,
\begin{equation}
\begin{split}
\mathrm{me}^2\left|(\mathrm{dc}_{A,a}\xi)(c)\right|_2^2\le\frac{3}{8}\Biggl[&\sum_{x\in\mathrm{Oc}_a(c)}\left|U_A(x,x^{(3)})\,f(x^{(3)})-f(x)\right|_2^2\\
&+\sum_{x'\in\mathrm{Oc}_{-a}(c')}\left|\mathrm{hp}(x')\right|_2^2+\sum_{x\in\mathrm{Oc}_a(c)}\left|\mathrm{hp}(x)\right|_2^2\Biggr].
\end{split}
\tag{1a.11}
\end{equation}
Now sum (1a.11) over the unoriented photon bonds with both ends in $E$;
in the reading (Es1) the sum extends over their classes, and in (i), (ii) below the words distinct and disjoint refer to classes.
\begin{itemize}
\item[(i)] For $x\in\mathrm{Oc}_a(c)$ the $8$ bonds $\{x,x^{(3)}\}$ are of type $3$ and have both ends over $E$; by (GS8), distinct photon bonds carry distinct bonds of type $3$.
Hence the sum of the first terms does not exceed the sum over all bonds of type $3$ with both ends over $E$, and the latter equals $\frac{48}{\mathrm{me}^3w_3}\,(f,\nabla_A^{*}\nabla_Af)_{E,3}$.
\item[(ii)] The second and third terms are the sums of $|\mathrm{hp}|_2^2$ over the two octets $\mathrm{Oc}_{a''}(c'')$ such that $c''$ is one end of the bond and $c''+\mathrm{me}\,a''$ the other.
The bond is determined by the pair $(c'',a'')$.
Consequently each octet occurs at most once in the whole sum, and distinct octets are disjoint subsets of $E\times O_{\mathrm{cub}}$ (GS8).
Hence the sum of these terms does not exceed $\sum_{x\text{ over }E}|\mathrm{hp}(x)|_2^2=\frac{48}{\mathrm{me}^3}\,\|\mathrm{hp}\|_{L^2,E}^2$.
\end{itemize}
Multiplying by $\mathrm{me}^3$ one obtains
\begin{equation}
\begin{aligned}
\mathrm{me}^2\,\|\mathrm{dc}_A\xi\|_E^2
&\le\frac{3}{8}\cdot48\left[w_3^{-1}(f,\nabla_A^{*}\nabla_Af)_{E,3}+\|\mathrm{hp}\|_{L^2,E}^2\right]\\
&\le18\left[w_3^{-1}(f,\nabla_A^{*}\nabla_Af)_{E,3}+\lambda_{\mathrm{int}}^{-1}(f,\nabla_A^{*}\nabla_Af)_{E,\mathrm{int}}\right]\\
&\le18\max\left(\frac{1}{w_3},\frac{1}{\lambda_{\mathrm{int}}}\right)(f,\nabla_A^{*}\nabla_Af)_E=c_{*}^{-1}\,(f,\nabla_A^{*}\nabla_Af)_E,
\end{aligned}
\tag{1a.12}
\end{equation}
where the second step rests on (1a.10).
This proves the first inequality of (d).
For $E=\mathrm{Lab}$ in the reading (Es1) the three local quantities reduce to $\|\mathrm{dc}_A\xi\|^2$, $(f,\nabla_A^{*}\nabla_Af)^{\mathrm{ch}}$ and $\|\cdot\|_{L^2}^2$ (Section~1.1);
so (1a.12) and (1a.10) become the last two inequalities of (d).

\emph{(e).}
Let $\xi$ be a function on $\mathrm{Lab}$, $f$ a function on $\mathrm{Lab}\times O_{\mathrm{cub}}$, $\psi$ a function of $c$, $v\in\mathrm{me}\,\mathbb{Z}^3$, and $\mathrm{Mx}$ a constant $4\times4$ matrix;
put $(T_vf)(c,g):=f(c-v,g)$ and $(T_v\xi)(c):=\xi(c-v)$.
Then
\begin{gather*}
(J\psi\xi)(c,g)=\psi(c)\,\xi(c)=(\psi J\xi)(c,g),
\qquad
(J^{*}\psi f)(c)=\frac{1}{48}\sum_g\psi(c)\,f(c,g)=\psi(c)\,(J^{*}f)(c),\\
(JT_v\xi)(c,g)=\xi(c-v)=(T_vJ\xi)(c,g),
\qquad
(J^{*}T_vf)(c)=\frac{1}{48}\sum_gf(c-v,g)=(T_vJ^{*}f)(c),
\end{gather*}
and, by linearity, $J\,\mathrm{Mx}\,\xi=\mathrm{Mx}\,J\xi$ and $J^{*}\mathrm{Mx}\,f=\mathrm{Mx}\,J^{*}f$.
Every map commuting with $J$ and $J^{*}$ commutes with $\mathrm{Plt}=JJ^{*}$ and with $\mathrm{Phv}=1-\mathrm{Plt}$.

\emph{The constants.}
By Section~1.1, the numbers $c_{\mathrm{hl}}$, $\hat g$, $c_{*}$, $\lambda_{\mathrm{int}}$ are functions of $\rho$ and of the weights.
The arguments above involve the data $(\mathrm{me},\mathrm{Lab},R,\theta,e_k,A)$ only through the properties stated in (GS) and, in the clauses for real $A$, through the equality $|u_A|=1$;
they involve $E$ only through the properties stated in (Es1), (Es2).
Hence none of the four numbers depends on these data or on $E$.
\end{proof}

\subsubsection*{1.3. Proposition~1.1}

The standing data are those fixed for the Propositions~1.1--1.7.

\begin{proof}[Proof of Proposition~1.1]
\emph{The data of (a)--(e).}
These data form the instance (In1) of (GS), with $\mathrm{me}=\eta_k$.
On a block torus $\mathbb{T}=(t,R,\boldsymbol{\sigma})$ the spinors $f$, $\xi$ of Proposition~1.1 range over $\mathfrak{S}^{\mathrm{ch}}$, $\mathfrak{S}$, which coincide with $\mathrm{HH}$, $\mathrm{HH}_4$ since $R<\infty$.
When $\mathrm{Lab}=\mathbb{L}_k$ carries no periodicity ($R=\infty$), the spinors of Proposition~1.1 are the square summable functions, i.e.\ the elements of $\mathrm{HH}$, $\mathrm{HH}_4$.
The connections admitted in Proposition~1.1 are exactly the functions $A$ of (GS).
The objects $u_A$, $D_A$, $\nabla_A^{*}\nabla_A$, $J$, $J^{*}$, $\mathrm{Plt}$, $\mathrm{Phv}$, the four blocks, $\mathrm{dc}_{A,a}$, $\|\mathrm{dc}_A\xi\|$, the pairings and the norms of Proposition~1.1 are the objects of Section~0.1 (Section~0.2, (In1)),
and $c_{\mathrm{hl}}$, $\hat g$, $c_{*}$ are the numbers of Section~1.1.

\emph{(a).}
This is Proposition~G\,(a);
at $\mathrm{me}=\eta_k$ the right side of its identity is $(D^{W,\rho}_A\xi)(c)$ of (I.3.2).

\emph{(b), (c).}
These are Proposition~G\,(b), (c), word for word, with $\mathrm{me}=\eta_k$.

\emph{(d).}
The global inequalities of (d) are the third and the fourth inequality of Proposition~G\,(d).
A set of photon sites of $\mathbb{T}$ is the image of an $R\mathbb{Z}^3$-periodic subset $E$ of $\mathrm{Lab}$, and a set of photon sites of $\mathbb{L}_k$ is a subset $E$ of $\mathrm{Lab}$;
together with $f\in\mathrm{HH}$ this is the reading (Es1).
The three quantities of Section~1.1 coincide with those of Proposition~1.1\,(d):
$\|\mathrm{dc}_A\xi\|_E^2$ and $\|\cdot\|_{L^2,E}$ by their displays,
and $(f,\nabla_A^{*}\nabla_Af)_E$ because, by (1a.2), which is the formula of (I.O2), it is the part of $(f,\nabla_A^{*}\nabla_Af)^{\mathrm{ch}}$ coming from the bonds $\{x,x^{(i)}\}$ with both ends over $E$.
Hence the local inequalities of (d) are the first two inequalities of Proposition~G\,(d).

\emph{(e).}
The translations of $\mathrm{Lab}$ are those by the vectors of $\eta_k\mathbb{Z}^3$.
For the multiplication by functions of $c$, for these translations and for constant spin matrices the assertion is Proposition~G\,(e).
Let $\mathrm{mph}(c)=rc+\mathrm{tv}$, with $r\in\{\pm1\}^3$, be a morphism $\mathbb{T}\to\mathbb{T}'$, and let $\mathrm{mS}=(S,\mathrm{mph})$ be as in Section~G1.2.
The map $\mathrm{mS}$ takes the functions on $\mathrm{Lab}\times O_{\mathrm{cub}}$ to the functions on $\mathrm{Lab}'\times O_{\mathrm{cub}}$, where $\mathrm{Lab}'=\mathrm{mph}(\mathrm{Lab})$ is the lattice of $\mathbb{T}'$,
and it acts on the functions on $\mathrm{Lab}$ as well; for $c\in\mathrm{Lab}'$ and $g\in O_{\mathrm{cub}}$,
\[
(\mathrm{mS}\,f)(c,g)=S\,f(\mathrm{mph}^{-1}c,r^{-1}g),
\qquad
(\mathrm{mS}\,\xi)(c)=S\,\xi(\mathrm{mph}^{-1}c).
\]
On both lattices the letters $J$, $J^{*}$, $\mathrm{Plt}$, $\mathrm{Phv}$ denote the maps of Section~0.1.
Then
\[
(J\,\mathrm{mS}\,\xi)(c,g)=S\,\xi(\mathrm{mph}^{-1}c)=S\,(J\xi)(\mathrm{mph}^{-1}c,r^{-1}g)=(\mathrm{mS}\,J\xi)(c,g),
\]
and, since $g\mapsto r^{-1}g$ is a bijection of $O_{\mathrm{cub}}$ (the group $\{\pm1\}^3$ is a subgroup of $O_{\mathrm{cub}}$),
\[
(J^{*}\mathrm{mS}\,f)(c)=\frac{1}{48}\sum_gS\,f(\mathrm{mph}^{-1}c,r^{-1}g)=S\,\frac{1}{48}\sum_{g'}f(\mathrm{mph}^{-1}c,g')=(\mathrm{mS}\,J^{*}f)(c).
\]
This computation is that of (I.a3).
Of $\mathrm{mph}$ it requires only that $c\mapsto rc+\mathrm{tv}$ be a bijection of $\mathrm{Lab}$ onto $\mathrm{Lab}'$ and that $r$ lie in $O_{\mathrm{cub}}$.
Thus $\mathrm{mS}$ commutes with $J$ and with $J^{*}$, hence also with $\mathrm{Plt}=JJ^{*}$ and with $\mathrm{Phv}=1-\mathrm{Plt}$.

\emph{(f).}
The data (GS), the spinors and the objects named in (f) are those of Section~0.1.
Indeed, the displays of Section~0.1 are the displays (I.1.2), (I.2.1)--(I.2.4), (I.3.1), (I.J4), (I.J6), those of Section~I.1.6 and those of the data of Proposition~1.1, with $\mathrm{me}$, $\mathrm{Lab}$ in the places of $\eta_k$, $\mathbb{L}_k$;
the counting modulo $R\mathbb{Z}^3$ is the congruence defined in Section~0.1;
and the square summable $f$, $\xi$ are the elements of $\mathrm{HH}$, $\mathrm{HH}_4$.
The individual assertions of (f) are obtained as follows.
\begin{itemize}
\item The facts (I.s1)--(I.s3) in (GS): Section~0.1 (for the bijections in (I.s1): the congruence of Section~0.1).
The fact (I.s4): the display of $x^{(i)}$ in Section~0.1 ($x^{(i)}$ lies over $c$ for $i\le2$, and $x^{(3)}$ lies over $c_3=c+\mathrm{me}\,ge_1$, a nearest neighbour of $c$), together with (GS8).
The identity (I.2.4): (GS2).
\item The clauses (a)--(d) with $\mathrm{me}$ in the place of $\eta_k$, and (e) without the maps $\mathrm{mS}$: Proposition~G\,(a)--(e), with the numbers $\lambda_{\mathrm{int}}$, $\hat g$, $c_{\mathrm{hl}}$, $c_{*}$ of Section~1.1.
With $\mathrm{me}$ in the place of $\eta_k$, the right side of Proposition~G\,(a) is the display (I.3.2).
The readings (Es1), (Es2) of Section~1.1 are the two forms of the set $E$ in the local form of (d).
\item The facts for real $A$.
The operator $R_i(A)$ is unitary and self-adjoint: (GS3).
The sum over bonds, and $\nabla_A^{*}\nabla_A=(\nabla_A^{*}\nabla_A)^{*}\ge0$: (GS4) with (1a.2).
The relation $J_A^{*}=-J_A$, the real part, $D_A^{*}=\gamma_5D_A\gamma_5$ and the norm of $D_A$: (GS5).
The relation $J^{*}J=1$, the adjoint, the isometry, the two orthogonal projections and the identity over one photon site: (GS6) with (1a.3).
The upper bound of the energy: (GS10), that is, (1a.6).
\item The bounds $0<\lambda_{\mathrm{int}}\le12+24\sqrt{2}<48$: (GS7), that is, (1a.4).
\item The three instances: Section~0.2, (In1), (In2), (In3).
In (In3) the operator $D_0$ corresponds to $D_0^{(p)}$, which is $D^{(p)}$ of (I.2.6) at $\boldsymbol{\mu}=0$, and $\|\mathrm{dc}_0\xi\|^2=2\,\mathrm{tw}(p)\,|v|_2^2$ is (1a.1).
\end{itemize}

\emph{The constants; uniformity.}
By the last assertion of Proposition~G, the numbers $c_{\mathrm{hl}}$, $\hat g$, $c_{*}$, $\lambda_{\mathrm{int}}$ depend neither on $(\mathrm{me},\mathrm{Lab},R,\theta,e_k,A)$ nor on $E$;
this is the assertion of Proposition~1.1 on the constants for the data of (f).
In the instance (In1) these data are built from $(k,t,R,\boldsymbol{\sigma},e_k,A)$.
Hence no constant depends on $A$, $e_k$, $k$, or on the block torus $\mathbb{T}=(t,R,\boldsymbol{\sigma})$.
Level $0$ corresponds to $\mathrm{me}=\eta_0=1$ (Section~0.2, (In1)).
\end{proof}

\subsection*{2. Proposition~1.2}

Let the data be those of Proposition~1.2.
In the instance (In3) of (GS), the fibre $p$ of Proposition~1.1\,(f), the clauses (b), (c), (d) of Proposition~G are the clauses (b), (c), (d) of Proposition~1.1 on the torus of one photon site with the character $c\mapsto\exp(ip\cdot c)$.

\subsubsection*{2.0. The fibre}

\emph{Norms.}
On $\mathbb{C}^4$ the norm is $|v|_2$, and on ${V_{\mathrm{v}}}$ it is $\|u\|$ (Section~0.2, (In3)).
For a linear map $O$ between two of the spaces $\mathbb{C}^4$, ${V_{\mathrm{v}}}$, $\ker J^{*}$, the operator norm is denoted by $|O|$; for a $4\times4$ matrix it is $|O|_2$.
By (GS6), $J$ is an isometry with adjoint $J^{*}$, one has $|J^{*}|\le1$, $|\mathrm{Phv}|\le1$ and $\ker J^{*}=\mathrm{ran}\,\mathrm{Phv}$,
and ${V_{\mathrm{v}}}=\mathrm{ran}\,J\oplus\ker J^{*}$ is an orthogonal sum:
$u=Jv+\mathrm{hp}$ with $v=J^{*}u$ and $\mathrm{hp}=\mathrm{Phv}\,u$, and $\|u\|^2=|v|_2^2+\|\mathrm{hp}\|^2$.

\emph{The mass in the blocks.}
One has $D^{(p)}=\mathbf{m}+D_0^{(p)}$.
Since $J^{*}J=1$, $J^{*}\mathrm{Phv}=0$ and $\mathrm{Phv}\,J=0$,
\begin{gather*}
D^{\mathrm{ll}}(p)=\mathbf{m}+J^{*}D_0^{(p)}J,\qquad
D^{\mathrm{lh}}(p)=J^{*}D_0^{(p)}\,\mathrm{Phv},\qquad
D^{\mathrm{hl}}(p)=\mathrm{Phv}\,D_0^{(p)}J,\\
D^{\mathrm{hh}}(p)=\mathbf{m}+\mathrm{Phv}\,D_0^{(p)}\,\mathrm{Phv}\quad\text{on }\ker J^{*}.
\end{gather*}

\emph{Real $p$.}
By Section~0.2, (In3), the operators $D_0^{(p)}$, $(\nabla^{*}\nabla)(p)$, $J$, $J^{*}$, $\mathrm{Phv}$ on ${V_{\mathrm{v}}}$, $\mathbb{C}^4$ are the operators $D_0$, $\nabla_0^{*}\nabla_0$, $J$, $J^{*}$, $\mathrm{Phv}$ of the instance (In3) of (GS),
and $\|\mathrm{dc}_0\xi\|^2=2\,\mathrm{tw}(p)\,|v|_2^2$ by (1a.1).
Hence, for real $p$, $u\in{V_{\mathrm{v}}}$ and $v\in\mathbb{C}^4$, the following hold.
\begin{itemize}
\item[(Fb1)] By (GS4) and (GS5),
\[
\mathrm{Re}\,(u,D^{(p)}u)=\mathbf{m}\,\|u\|^2+\frac{\rho}{2}\,\left(u,(\nabla^{*}\nabla)(p)\,u\right),
\qquad
\left(u,(\nabla^{*}\nabla)(p)\,u\right)\ge0,
\]
\[
\left(D^{(p)}\right)^{*}=\gamma_5D^{(p)}\gamma_5,
\qquad
|D^{(p)}|\le\mathbf{m}+(1+\rho)\,\mathrm{sw}.
\]
\item[(Fb2)] By Proposition~G\,(b) and (1a.1),
\[
\|D^{\mathrm{hl}}(p)v\|\le c_{\mathrm{hl}}\left(2\,\mathrm{tw}(p)\right)^{\frac{1}{2}}|v|_2,
\qquad
\left|(v,D^{\mathrm{lh}}(p)u)\right|\le c_{\mathrm{hl}}\left(2\,\mathrm{tw}(p)\right)^{\frac{1}{2}}|v|_2\cdot\|u\|.
\]
\item[(Fb3)] By Proposition~G\,(c) with $\mathrm{me}=1$,
\[
\mathrm{Re}\,(u,D^{(p)}u)\ge(\mathbf{m}+\hat g)\,\|u\|^2\qquad\text{for }u\in\ker J^{*}.
\]
\item[(Fb4)] By the third inequality of Proposition~G\,(d) with $\mathrm{me}=1$, and by (1a.1),
\[
\left(u,(\nabla^{*}\nabla)(p)\,u\right)\ge2c_{*}\cdot\mathrm{tw}(p)\,|J^{*}u|_2^2.
\]
\end{itemize}

\emph{Complex $p$.}
For complex $p$, the symbols $\sin p_\mu$, $\cos p_\mu$ and $\mathrm{tw}(p)$ denote the entire functions of $p$.
\begin{itemize}
\item[(Fb5)] For $p,\tilde p\in\mathbb{C}^3$,
\[
\left(\left(D^{(p)}-D^{(\tilde p)}\right)u\right)(g)=-M_3(g)\left(\exp\left(ip\cdot(ge_1)\right)-\exp\left(i\tilde p\cdot(ge_1)\right)\right)u(gs_3),
\]
and
\[
\left|D^{(p)}-D^{(\tilde p)}\right|\le3(1+\rho)\max_a\left|\exp(ip\cdot a)-\exp(i\tilde p\cdot a)\right|.
\]
Indeed, $p$ enters neither the terms with $i\le2$ nor the mass term.
Moreover $|M_3(g)|_2=3(1+\rho)$ by (I.m6), and $g\mapsto gs_3$ is a bijection of $O_{\mathrm{cub}}$; hence
\[
\left\|\left(D^{(p)}-D^{(\tilde p)}\right)u\right\|^2\le\left(3(1+\rho)\max_a\left|\exp(ip\cdot a)-\exp(i\tilde p\cdot a)\right|\right)^2\frac{1}{48}\sum_g|u(gs_3)|_2^2.
\]
\item[(Fb6)] For $p=p'+ip''$ with $p',p''\in\mathbb{R}^3$, and for $a=\pm e_\mu$,
\[
\left|\exp(ip\cdot a)-\exp(ip'\cdot a)\right|=\left|\exp(-p''\cdot a)-1\right|\le\exp(|p''_\mu|)-1.
\]
Indeed, $\exp(ip\cdot a)=\exp(ip'\cdot a)\exp(-p''\cdot a)$, $|\exp(ip'\cdot a)|=1$ and $p''\cdot a=\pm p''_\mu$;
moreover $1-\exp(-s)\le\exp(s)-1$ for $s\ge0$, because $\exp(s)+\exp(-s)\ge2$.
\item[(Fb7)] For $p\in\mathbb{C}^3$ and $v\in\mathbb{C}^4$,
\[
(D^{(p)}Jv)(g)=\mathbf{m}\,v-M_3(g)\left(\exp\left(ip\cdot(ge_1)\right)-1\right)v,
\qquad
D^{\mathrm{ll}}(p)=\mathbf{m}+i\gamma\cdot\sin p+\rho\,\mathrm{tw}(p).
\]
Indeed, constants are annihilated by the terms with $i\le2$ of (I.2.6).
Further, $M_3(g)=3\left(\rho-\gamma\cdot(ge_1)\right)$, and each $a$ equals $ge_1$ for exactly $8$ labels $g$ (GS8); therefore
\[
J^{*}D^{(p)}Jv=\mathbf{m}\,v-\sum_a\frac{1}{2}\left(\rho-\gamma\cdot a\right)\left(\exp(ip\cdot a)-1\right)v=\mathbf{m}\,v-\sum_\mu\left[\rho\left(\cos p_\mu-1\right)-i\gamma_\mu\sin p_\mu\right]v.
\]
\end{itemize}

\emph{Strips.}
Let $\mathrm{Stp}:=\{p\in\mathbb{C}^3:|\mathrm{Im}\,p_\mu|\le\theta_0,\ \mu=1,2,3\}$.
By the definition of $\theta_0$ in Proposition~1.2, namely $3(1+\rho)\left(\exp(\theta_0)-1\right)=\hat g/2$, one has $\theta_0=\log\left(1+\frac{\hat g}{6(1+\rho)}\right)$.
Put
\[
\theta_0':=\log\left(1+\frac{\hat g}{3(1+\rho)}\right)>\theta_0,
\qquad
\mathrm{Stp}':=\{p\in\mathbb{C}^3:|\mathrm{Im}\,p_\mu|<\theta_0'\};
\]
$\mathrm{Stp}'$ is an open set containing $\mathrm{Stp}$.
Holomorphy on $\mathrm{Stp}$ is always established in the form of holomorphy on $\mathrm{Stp}'$.

\emph{Entire entries.}
Fix a basis of ${V_{\mathrm{v}}}$ adapted to the decomposition $\mathrm{ran}\,J\oplus\ker J^{*}$.
In this basis the entries of $D^{(p)}$ are polynomials in $\exp(\pm ip_\mu)$; they are entire and $2\pi\mathbb{Z}^3$-periodic.
The same holds for the entries of $D^{\mathrm{ll}}(p)$, $D^{\mathrm{lh}}(p)$, $D^{\mathrm{hl}}(p)$, $D^{\mathrm{hh}}(p)$.

\emph{Two tools.}
\begin{itemize}
\item[(Cy)] (Cauchy's inequalities.)
Let $\mathrm{hf}$ be a function with values in the $4\times4$ matrices, holomorphic in a neighbourhood of the closed polydisc $\{\mathrm{pz}\in\mathbb{C}^3:|\mathrm{pz}_\mu-p_\mu|\le\mathrm{rd}\}$, and let $|\mathrm{hf}(\mathrm{pz})|_2\le\mathrm{Kq}$ there.
Then $|\partial^{\mathrm{ak}}\mathrm{hf}(p)|_2\le\mathrm{ak}!\cdot\mathrm{Kq}\cdot\mathrm{rd}^{-|\mathrm{ak}|}$ for every multi-index $\mathrm{ak}$.
\end{itemize}
\begin{proof}[Proof of (Cy)]
Fix unit vectors $v$, $v'$;
the function $\mathrm{pz}\mapsto(v',\mathrm{hf}(\mathrm{pz})v)$ is holomorphic and bounded in modulus by $\mathrm{Kq}$.
By Cauchy's integral formula for polydiscs, differentiated under the integral sign \cite[Section~2.2]{Hormander1990},
\[
\partial^{\mathrm{ak}}(v',\mathrm{hf}\,v)(p)=\mathrm{ak}!\,(2\pi i)^{-3}\int_{|\mathrm{pz}_1-p_1|=\mathrm{rd}}\mathrm{d}\mathrm{pz}_1\int_{|\mathrm{pz}_2-p_2|=\mathrm{rd}}\mathrm{d}\mathrm{pz}_2\int_{|\mathrm{pz}_3-p_3|=\mathrm{rd}}\mathrm{d}\mathrm{pz}_3\,(v',\mathrm{hf}(\mathrm{pz})v)\prod_\mu(\mathrm{pz}_\mu-p_\mu)^{-\mathrm{ak}_\mu-1},
\]
where the three integrations extend over the circles $|\mathrm{pz}_\mu-p_\mu|=\mathrm{rd}$, each of length $2\pi\,\mathrm{rd}$.
This gives $|\partial^{\mathrm{ak}}(v',\mathrm{hf}\,v)(p)|\le\mathrm{ak}!\cdot\mathrm{Kq}\cdot\mathrm{rd}^{-|\mathrm{ak}|}$,
and taking the supremum over $v$, $v'$ yields the norm.
\end{proof}
\begin{itemize}
\item[(Li)] (Derivatives of an inverse.)
Let $\mathrm{Mx}$ be a holomorphic function on an open subset of $\mathbb{C}^3$ with values in the $4\times4$ matrices, and $p$ a point at which $\mathrm{Mx}(p)$ is invertible.
Let $\mathrm{cf}>0$, $0<\mathrm{rd}\le\mathrm{rd}_{\max}$ and $d_{\mathrm{ak}}\ge0$ ($\mathrm{ak}\ne0$) be numbers with
\[
|\mathrm{Mx}(p)^{-1}|_2\le(\mathrm{cf}\cdot\mathrm{rd})^{-1},
\qquad
|\partial^{\mathrm{ak}}\mathrm{Mx}(p)|_2\le d_{\mathrm{ak}}.
\]
Define recursively
\[
\mathrm{Bq}_0:=\frac{1}{\mathrm{cf}},
\qquad
\mathrm{Bq}_{\mathrm{ak}}:=\frac{1}{\mathrm{cf}}\sum_{0\ne\mathrm{ak}'\le\mathrm{ak}}\binom{\mathrm{ak}}{\mathrm{ak}'}d_{\mathrm{ak}'}\cdot\mathrm{Bq}_{\mathrm{ak}-\mathrm{ak}'}\cdot\mathrm{rd}_{\max}^{|\mathrm{ak}'|-1}.
\]
Then $\mathrm{Mx}^{-1}$ is holomorphic near $p$, and $|\partial^{\mathrm{ak}}(\mathrm{Mx}^{-1})(p)|_2\le\mathrm{Bq}_{\mathrm{ak}}\cdot\mathrm{rd}^{-1-|\mathrm{ak}|}$ for all $\mathrm{ak}$.
\end{itemize}
\begin{proof}[Proof of (Li)]
The determinant of $\mathrm{Mx}$ does not vanish near $p$, so $\mathrm{Mx}^{-1}$ is holomorphic there by Cramer's rule.
For $\mathrm{ak}\ne0$ the Leibniz rule, applied to $\mathrm{Mx}\cdot\mathrm{Mx}^{-1}=1$, gives
\[
\partial^{\mathrm{ak}}(\mathrm{Mx}^{-1})=-\mathrm{Mx}^{-1}\sum_{0\ne\mathrm{ak}'\le\mathrm{ak}}\binom{\mathrm{ak}}{\mathrm{ak}'}\left(\partial^{\mathrm{ak}'}\mathrm{Mx}\right)\left(\partial^{\mathrm{ak}-\mathrm{ak}'}(\mathrm{Mx}^{-1})\right).
\]
The bound is proved by induction on $|\mathrm{ak}|$ at the point $p$; for $\mathrm{ak}=0$ it is the hypothesis $|\mathrm{Mx}(p)^{-1}|_2\le\mathrm{Bq}_0\cdot\mathrm{rd}^{-1}$.
For $\mathrm{ak}\ne0$ the term with index $\mathrm{ak}'$ is at most
\[
(\mathrm{cf}\cdot\mathrm{rd})^{-1}\binom{\mathrm{ak}}{\mathrm{ak}'}d_{\mathrm{ak}'}\cdot\mathrm{Bq}_{\mathrm{ak}-\mathrm{ak}'}\cdot\mathrm{rd}^{-1-|\mathrm{ak}|+|\mathrm{ak}'|}
=\frac{1}{\mathrm{cf}}\binom{\mathrm{ak}}{\mathrm{ak}'}d_{\mathrm{ak}'}\cdot\mathrm{Bq}_{\mathrm{ak}-\mathrm{ak}'}\cdot\mathrm{rd}^{|\mathrm{ak}'|-1}\cdot\mathrm{rd}^{-1-|\mathrm{ak}|},
\]
and $\mathrm{rd}^{|\mathrm{ak}'|-1}\le\mathrm{rd}_{\max}^{|\mathrm{ak}'|-1}$ because $|\mathrm{ak}'|\ge1$.
Summing over $\mathrm{ak}'$ gives the bound $\mathrm{Bq}_{\mathrm{ak}}\cdot\mathrm{rd}^{-1-|\mathrm{ak}|}$.
\end{proof}
For a holomorphic function, the derivatives of its restriction to $\mathbb{R}^3$ in real directions coincide with its complex derivatives;
so (Cy) and (Li) apply to the derivatives in (S6).

\subsubsection*{2.1. (S1)}

Let $p=p'+ip''\in\mathrm{Stp}'$, and put
\[
\delta(p):=3(1+\rho)\left(\exp\left(\max_\mu|p''_\mu|\right)-1\right).
\]
By the definitions of $\theta_0'$ and $\theta_0$, one has $\delta(p)<\hat g$ on $\mathrm{Stp}'$ and $\delta(p)\le\hat g/2$ on $\mathrm{Stp}$;
moreover $\delta(p)=0$ for real $p$, and $|D^{(p)}-D^{(p')}|\le\delta(p)$ by (Fb5) and (Fb6).
For $u\in\ker J^{*}$ one has $(u,D^{\mathrm{hh}}(p)u)=(\mathrm{Phv}\,u,D^{(p)}u)=(u,D^{(p)}u)$, and so, by (Fb3),
\begin{equation}
\mathrm{Re}\,(u,D^{\mathrm{hh}}(p)u)\ge\mathrm{Re}\,(u,D^{(p')}u)-\delta(p)\,\|u\|^2\ge\left(\mathbf{m}+\hat g-\delta(p)\right)\|u\|^2.
\tag{1a.13}
\end{equation}
By (GS11) on the space $\ker J^{*}$, the operator $D^{\mathrm{hh}}(p)$ is invertible for all $p\in\mathrm{Stp}'$, with
\[
\left|D^{\mathrm{hh}}(p)^{-1}\right|\le\left(\mathbf{m}+\hat g-\delta(p)\right)^{-1}.
\]
This bound is at most $2/\hat g$ on $\mathrm{Stp}$ (because $\mathbf{m}\ge0$), and it equals $(\hat g+\mathbf{m})^{-1}$ for real $p$.
Put $\mathrm{Hi}(p):=D^{\mathrm{hh}}(p)^{-1}$.
As $\det D^{\mathrm{hh}}(p)$ does not vanish on $\mathrm{Stp}'$ and the entries of $D^{\mathrm{hh}}(p)$ are entire functions, Cramer's rule shows that $\mathrm{Hi}$ is holomorphic on $\mathrm{Stp}'$;
it is $2\pi\mathbb{Z}^3$-periodic.

\subsubsection*{2.2. (S2)}

\begin{proof}[Proof of Proposition~1.2\,(S2)]
Let $p\in\mathbb{R}^3$.
The first inequality of (Fb2) is the bound
\begin{equation*}
|D^{\mathrm{hl}}(p)|\le c_{\mathrm{hl}}\cdot\left(2\,\mathrm{tw}(p)\right)^{\frac{1}{2}}
\end{equation*}
for the operator norm.
For $u\in{V_{\mathrm{v}}}$ the second inequality of (Fb2) yields
\begin{equation*}
|D^{\mathrm{lh}}(p)\,u|_2=\sup\left\{|(v,D^{\mathrm{lh}}(p)\,u)|\,:\,v\in\mathbb{C}^4,\ |v|_2=1\right\}\le c_{\mathrm{hl}}\cdot\left(2\,\mathrm{tw}(p)\right)^{\frac{1}{2}}\cdot\|u\|.
\end{equation*}
At $p=0$ one has $\mathrm{tw}(0)=0$, and therefore $D^{\mathrm{lh}}(0)=0$ and $D^{\mathrm{hl}}(0)=0$.
Finally, $1-\cos s\le s^2/2$ for real $s$ gives $2\,\mathrm{tw}(p)\le|p|_2^2$.
\end{proof}

\subsubsection*{2.3. (S3)}

\begin{proof}[Proof of Proposition~1.2\,(S3)]
The matrix
\begin{equation*}
\mathrm{Dsc}_{\mathbf{m}}(p)=D^{\mathrm{ll}}(p)-D^{\mathrm{lh}}(p)\,\mathrm{Hi}(p)\,D^{\mathrm{hl}}(p)
\end{equation*}
is built by products and sums from functions that are holomorphic on $\mathrm{Stp}'$ and $2\pi\mathbb{Z}^3$-periodic (2.0, 2.1);
so it has both properties.

\emph{The bound.}
Let $p=p'+ip''\in\mathrm{Stp}$ with $p',p''\in\mathbb{R}^3$.
From (Fb1), (Fb5), (Fb6) and $\mathbf{m}\le1$,
\begin{equation*}
|D^{(p)}|\le|D^{(p')}|+\frac{\hat{g}}{2}\le K_D:=1+(1+\rho)\,\mathrm{sw}+\frac{\hat{g}}{2}.
\end{equation*}
The blocks $D^{\mathrm{ll}}(p)$, $D^{\mathrm{lh}}(p)$, $D^{\mathrm{hl}}(p)$ arise from $D^{(p)}$ by composition with the maps $J$, $J^{*}$, $\mathrm{Phv}$, whose norms are at most one;
their norms are therefore at most $K_D$.
Together with (S1) this gives
\begin{equation*}
|\mathrm{Dsc}_{\mathbf{m}}(p)|_2\le K_D+\frac{2K_D^2}{\hat{g}}\qquad\text{for all }p\in\mathrm{Stp}\text{ and all }\mathbf{m}\in[0,1].
\end{equation*}
Consequently $K_{\mathrm{Str}}\le K_D+2K_D^2/\hat{g}<\infty$, and this bound is a function of $\rho$.

\emph{Real $p$.}
For $p\in\mathbb{R}^3$, (Fb7) gives
\begin{equation*}
\mathrm{Dsc}_{\mathbf{m}}(p)=\mathbf{m}+i\gamma\cdot\sin p+\mathrm{Er}_{\mathbf{m}}(p),\qquad
\mathrm{Er}_{\mathbf{m}}(p):=\rho\cdot\mathrm{tw}(p)-D^{\mathrm{lh}}(p)\,\mathrm{Hi}(p)\,D^{\mathrm{hl}}(p),
\end{equation*}
and (S1), (S2) give
\begin{equation*}
|\mathrm{Er}_{\mathbf{m}}(p)|_2\le\rho\cdot\mathrm{tw}(p)+c_{\mathrm{hl}}^2\cdot\frac{2\,\mathrm{tw}(p)}{\hat{g}+\mathbf{m}}\le\left(\rho+\frac{2c_{\mathrm{hl}}^2}{\hat{g}}\right)\mathrm{tw}(p)=K_{\rho}\cdot\mathrm{tw}(p).
\end{equation*}
In particular $\mathrm{Dsc}_{\mathbf{m}}(0)=\mathbf{m}$;
hence $K_{\mathrm{Str}}\ge|\mathrm{Dsc}_1(0)|_2=1$.
\end{proof}

\subsubsection*{2.4. (S4)}

\begin{proof}[Proof of Proposition~1.2\,(S4)]
Let $p\in\mathbb{R}^3$ and $v\in\mathbb{C}^4$, and put
\begin{equation*}
u:=Jv-\mathrm{Hi}(p)\,D^{\mathrm{hl}}(p)\,v.
\end{equation*}
Its two parts are $J^{*}u=v$ and $\mathrm{Phv}\,u=-\mathrm{Hi}(p)\,D^{\mathrm{hl}}(p)\,v$.
Hence
\begin{align*}
J^{*}D^{(p)}u&=D^{\mathrm{ll}}(p)\,v-D^{\mathrm{lh}}(p)\,\mathrm{Hi}(p)\,D^{\mathrm{hl}}(p)\,v=\mathrm{Dsc}_{\mathbf{m}}(p)\,v,\\
\mathrm{Phv}\,D^{(p)}u&=D^{\mathrm{hl}}(p)\,v-D^{\mathrm{hh}}(p)\,\mathrm{Hi}(p)\,D^{\mathrm{hl}}(p)\,v=0,
\end{align*}
that is, $D^{(p)}u=J\,\mathrm{Dsc}_{\mathbf{m}}(p)\,v$, and
\begin{equation*}
(u,D^{(p)}u)=(u,J\,\mathrm{Dsc}_{\mathbf{m}}(p)\,v)=(J^{*}u,\mathrm{Dsc}_{\mathbf{m}}(p)\,v)=(v,\mathrm{Dsc}_{\mathbf{m}}(p)\,v).
\end{equation*}
Since $\|u\|^2\ge|v|_2^2$ and $\mathbf{m}\ge0$, (Fb1) and (Fb4) give
\begin{align*}
\mathrm{Re}\,(v,\mathrm{Dsc}_{\mathbf{m}}(p)\,v)
&=\mathbf{m}\cdot\|u\|^2+\frac{\rho}{2}\,(u,(\nabla^{*}\nabla)(p)\,u)\\
&\ge\mathbf{m}\cdot|v|_2^2+\frac{\rho}{2}\cdot2c_{*}\cdot\mathrm{tw}(p)\cdot|v|_2^2=\left(\mathbf{m}+\rho\cdot c_{*}\cdot\mathrm{tw}(p)\right)|v|_2^2.
\end{align*}
\end{proof}

\subsubsection*{2.5. (S5)}

Three inequalities between the constants come first;
(i) and (ii) establish the inequalities $K_{\rho}>1$, $t_{0}\le1/4$, $c_{*}\le1/3$ that are asserted in brackets in the list of constants of Proposition~1.2.
\begin{itemize}
\item[(i)] By $(1+\rho)^2\ge4\rho$ and (1a.4),
\begin{equation*}
K_{\rho}\ge\frac{2c_{\mathrm{hl}}^2}{\hat{g}}=\frac{12(1+\rho)^2}{\rho\cdot\lambda_{\mathrm{int}}}\ge\frac{48}{\lambda_{\mathrm{int}}}>1;
\end{equation*}
hence $t_{0}=1/(4K_{\rho}^2)\le1/4$.
\item[(ii)] $c_{*}\le w_3/18=1/3$.
\item[(iii)] $1/(\sqrt{2}\,\pi)>0.22$, because $\sqrt{2}\,\pi<1.4143\cdot3.1416<4.4432<1/0.22$.
\end{itemize}

\begin{proof}[Proof of Proposition~1.2\,(S5)]
Let $p\in[-\pi,\pi]^3$ and $t_\mu:=1-\cos p_\mu\in[0,2]$, so that $\mathrm{tw}(p)=\sum_\mu t_\mu$ and $\sin^2p_\mu=t_\mu(2-t_\mu)$.
The matrix $\gamma\cdot\sin p$ is Hermitian, with $(\gamma\cdot\sin p)^2=|\sin p|_2^2$ (c1);
therefore
\begin{equation}\tag{1a.14}
|(\mathbf{m}+i\gamma\cdot\sin p)\,v|_2^2=\left(v,(\mathbf{m}-i\gamma\cdot\sin p)(\mathbf{m}+i\gamma\cdot\sin p)\,v\right)=\left(\mathbf{m}^2+|\sin p|_2^2\right)|v|_2^2.
\end{equation}

\emph{The case $\mathrm{tw}(p)\le t_{0}$.}
Each $t_\mu$ is at most $t_{0}\le1$, so that $|\sin p|_2^2=\sum_\mu t_\mu(2-t_\mu)\ge\mathrm{tw}(p)$.
By (S3),
\begin{equation*}
|\mathrm{Er}_{\mathbf{m}}(p)|_2\le K_{\rho}\cdot\mathrm{tw}(p)\le K_{\rho}\cdot t_{0}^{\frac{1}{2}}\cdot\mathrm{tw}(p)^{\frac{1}{2}}=\frac{1}{2}\,\mathrm{tw}(p)^{\frac{1}{2}}\le\frac{1}{2}\,|\sin p|_2\le\frac{1}{2}\left(\mathbf{m}^2+|\sin p|_2^2\right)^{\frac{1}{2}},
\end{equation*}
and together with (1a.14)
\begin{equation}\tag{1a.15}
|\mathrm{Dsc}_{\mathbf{m}}(p)\,v|_2\ge\frac{1}{2}\left(\mathbf{m}^2+|\sin p|_2^2\right)^{\frac{1}{2}}|v|_2\ge\frac{1}{2}\left(\mathbf{m}^2+\mathrm{tw}(p)\right)^{\frac{1}{2}}|v|_2.
\end{equation}
For $|s|\le\pi$ one has $1-\cos s=2\sin^2(s/2)\ge2s^2/\pi^2$;
indeed $\sin$ is concave on $[0,\pi/2]$, so that $\sin(s/2)\ge s/\pi$ for $0\le s\le\pi$.
Hence $\mathrm{tw}(p)\ge(2/\pi^2)\,|p|_2^2$, and, since $2/\pi^2<1$, also $\mathbf{m}^2+\mathrm{tw}(p)\ge(2/\pi^2)\,\mathrm{rp}(p)^2$.
It follows that
\begin{equation*}
|\mathrm{Dsc}_{\mathbf{m}}(p)\,v|_2\ge\frac{1}{\sqrt{2}\,\pi}\,\mathrm{rp}(p)\cdot|v|_2\ge0.22\,\mathrm{rp}(p)\cdot|v|_2\ge c_{\mathrm{el}}\cdot\mathrm{rp}(p)\cdot|v|_2.
\end{equation*}

\emph{The case $\mathrm{tw}(p)\ge t_{0}$.}
By (S4),
\begin{equation*}
|\mathrm{Dsc}_{\mathbf{m}}(p)\,v|_2\cdot|v|_2\ge\mathrm{Re}\,(v,\mathrm{Dsc}_{\mathbf{m}}(p)\,v)\ge\left(\mathbf{m}+\rho\cdot c_{*}\cdot t_{0}\right)|v|_2^2.
\end{equation*}
On the other hand $\mathrm{rp}(p)\le\mathbf{m}+|p|_2\le\mathbf{m}+\pi\sqrt{3}$, and $c_{\mathrm{el}}\le0.22<1$, $c_{\mathrm{el}}\cdot\pi\sqrt{3}\le\rho\cdot c_{*}\cdot t_{0}$;
so $c_{\mathrm{el}}\cdot\mathrm{rp}(p)\le\mathbf{m}+\rho\cdot c_{*}\cdot t_{0}$, and the inequality of (S5) holds in this case too.
\end{proof}

\subsubsection*{2.6. (S6)}

\begin{proof}[Proof of Proposition~1.2\,(S6)]
Let $p\in[-\pi,\pi]^3\setminus\{0\}$.
Then $\mathrm{rp}(p)>0$, and by (S5) the matrix $\mathrm{Dsc}_{\mathbf{m}}(p)$ is invertible, with $|\mathrm{Dsc}_{\mathbf{m}}(p)^{-1}|_2\le1/(c_{\mathrm{el}}\cdot\mathrm{rp}(p))$.
The closed polydisc of radius $\theta_{0}$ about the real point $p$ is contained in $\mathrm{Stp}$, where $|\mathrm{Dsc}_{\mathbf{m}}|_2\le K_{\mathrm{Str}}$, and $\mathrm{Dsc}_{\mathbf{m}}$ is holomorphic on the neighbourhood $\mathrm{Stp}'$ of $\mathrm{Stp}$.
By (Cy),
\begin{equation*}
|\partial^{\mathrm{ak}}\mathrm{Dsc}_{\mathbf{m}}(p)|_2\le d_{\mathrm{ak}}:=\mathrm{ak}!\cdot K_{\mathrm{Str}}\cdot\theta_{0}^{-|\mathrm{ak}|}.
\end{equation*}
Moreover $\mathrm{rp}(p)\le\mathrm{rp}_{\mathrm{max}}:=\left(3\pi^2+1\right)^{\frac{1}{2}}$.
Thus (Li) applies with $\mathrm{Mx}=\mathrm{Dsc}_{\mathbf{m}}$, $\mathrm{cf}=c_{\mathrm{el}}$, $\mathrm{rd}=\mathrm{rp}(p)$, $\mathrm{rd}_{\mathrm{max}}=\mathrm{rp}_{\mathrm{max}}$, and it gives (S6) with
\begin{equation*}
A_0:=\frac{1}{c_{\mathrm{el}}},\qquad
A_{\mathrm{ak}}:=\frac{1}{c_{\mathrm{el}}}\sum_{0\ne\mathrm{ak}'\le\mathrm{ak}}\binom{\mathrm{ak}}{\mathrm{ak}'}\cdot\mathrm{ak}'!\cdot K_{\mathrm{Str}}\cdot\theta_{0}^{-|\mathrm{ak}'|}\cdot A_{\mathrm{ak}-\mathrm{ak}'}\cdot\mathrm{rp}_{\mathrm{max}}^{|\mathrm{ak}'|-1}.
\end{equation*}
These numbers are formed with $c_{\mathrm{el}}$, $K_{\mathrm{Str}}$, $\theta_{0}$;
they are therefore functions of $\rho$.
\end{proof}

\subsubsection*{2.7. (S7)}

\begin{proof}[Proof of Proposition~1.2\,(S7)]
Let $\mathbf{m}>0$ and $p=p'+ip''$ with $p'\in[-\pi,\pi]^3$ and $|p''_\mu|\le\kappa_{(0)}\cdot\mathbf{m}$.
As $\kappa_{(0)}\le\theta_{0}/2$ and $\mathbf{m}\le1$, every point $\mathrm{pz}$ of the segment $\mathrm{pz}(s):=p'+is\,p''$, $0\le s\le1$, satisfies $|\mathrm{Im}\,\mathrm{pz}_\mu|\le\theta_{0}/2$, and the closed polydisc of radius $\theta_{0}/2$ about $\mathrm{pz}$ lies in $\mathrm{Stp}$.
By (Cy),
\begin{equation}\tag{1a.16}
|\partial^{\mathrm{ak}}\mathrm{Dsc}_{\mathbf{m}}(\mathrm{pz})|_2\le d'_{\mathrm{ak}}:=\mathrm{ak}!\cdot K_{\mathrm{Str}}\cdot(\theta_{0}/2)^{-|\mathrm{ak}|}\qquad\text{for all }\mathrm{pz}\text{ with }|\mathrm{Im}\,\mathrm{pz}_\mu|\le\theta_{0}/2.
\end{equation}
Since
\begin{equation*}
\mathrm{Dsc}_{\mathbf{m}}(p)-\mathrm{Dsc}_{\mathbf{m}}(p')=\int_0^1\mathrm{d}s\,\sum_\mu ip''_\mu\,(\partial_\mu\mathrm{Dsc}_{\mathbf{m}})(\mathrm{pz}(s)),
\end{equation*}
the bound (1a.16) with $|\mathrm{ak}|=1$ and the inequality $\kappa_{(0)}\le c_{\mathrm{el}}\cdot\theta_{0}/(12K_{\mathrm{Str}})$ give
\begin{equation*}
|\mathrm{Dsc}_{\mathbf{m}}(p)-\mathrm{Dsc}_{\mathbf{m}}(p')|_2\le3\kappa_{(0)}\cdot\mathbf{m}\cdot\frac{2K_{\mathrm{Str}}}{\theta_{0}}\le\frac{c_{\mathrm{el}}}{2}\,\mathbf{m}\le\frac{c_{\mathrm{el}}}{2}\,\mathrm{rp}(p').
\end{equation*}
Combined with (S5) at $p'$ this yields
\begin{equation*}
|\mathrm{Dsc}_{\mathbf{m}}(p)\,v|_2\ge c_{\mathrm{el}}\cdot\mathrm{rp}(p')\cdot|v|_2-\frac{c_{\mathrm{el}}}{2}\,\mathrm{rp}(p')\cdot|v|_2=\frac{c_{\mathrm{el}}}{2}\,\mathrm{rp}(p')\cdot|v|_2,
\end{equation*}
which is the first assertion; moreover $\mathrm{rp}(p')\ge\mathbf{m}>0$.
So $\mathrm{Dsc}_{\mathbf{m}}(p)$ is invertible, and $|\mathrm{Dsc}_{\mathbf{m}}(p)^{-1}|_2\le2/(c_{\mathrm{el}}\cdot\mathrm{rp}(p'))$.
Now (Li), applied with $\mathrm{Mx}=\mathrm{Dsc}_{\mathbf{m}}$ on $\mathrm{Stp}'$, $\mathrm{cf}=c_{\mathrm{el}}/2$, $\mathrm{rd}=\mathrm{rp}(p')$, $\mathrm{rd}_{\mathrm{max}}=\mathrm{rp}_{\mathrm{max}}$ and the numbers $d'_{\mathrm{ak}}$ of (1a.16), gives the second assertion with
\begin{equation*}
A'_0:=\frac{2}{c_{\mathrm{el}}},\qquad
A'_{\mathrm{ak}}:=\frac{2}{c_{\mathrm{el}}}\sum_{0\ne\mathrm{ak}'\le\mathrm{ak}}\binom{\mathrm{ak}}{\mathrm{ak}'}\cdot\mathrm{ak}'!\cdot K_{\mathrm{Str}}\cdot(\theta_{0}/2)^{-|\mathrm{ak}'|}\cdot A'_{\mathrm{ak}-\mathrm{ak}'}\cdot\mathrm{rp}_{\mathrm{max}}^{|\mathrm{ak}'|-1}.
\end{equation*}
\end{proof}

\subsubsection*{2.8. (S8)}

\begin{proof}[Proof of Proposition~1.2\,(S8)]
Let $p\in\mathrm{Stp}$; the argument holds on all of $\mathrm{Stp}'$.
The map $(v,\mathrm{hp})\mapsto Jv+\mathrm{hp}$ from $\mathbb{C}^4\oplus\ker J^{*}$ to ${V_{\mathrm{v}}}$ is unitary (2.0).
In this decomposition $D^{(p)}$ is the block matrix
\begin{equation*}
\left(\begin{array}{cc}D^{\mathrm{ll}}&D^{\mathrm{lh}}\\D^{\mathrm{hl}}&D^{\mathrm{hh}}\end{array}\right),
\end{equation*}
with $D^{\mathrm{lh}}$ restricted to $\ker J^{*}$:
for $u=Jv+\mathrm{hp}$ one has $J^{*}D^{(p)}u=D^{\mathrm{ll}}v+D^{\mathrm{lh}}\mathrm{hp}$ and $\mathrm{Phv}\,D^{(p)}u=D^{\mathrm{hl}}v+D^{\mathrm{hh}}\mathrm{hp}$.
With $\mathrm{Hi}=\mathrm{Hi}(p)$ of (S1), multiplying out the three factors gives
\begin{equation}\tag{1a.17}
\begin{aligned}
&\left(\begin{array}{cc}1&D^{\mathrm{lh}}\mathrm{Hi}\\0&1\end{array}\right)
\left(\begin{array}{cc}\mathrm{Dsc}_{\mathbf{m}}&0\\0&D^{\mathrm{hh}}\end{array}\right)
\left(\begin{array}{cc}1&0\\\mathrm{Hi}\,D^{\mathrm{hl}}&1\end{array}\right)
=\left(\begin{array}{cc}1&D^{\mathrm{lh}}\mathrm{Hi}\\0&1\end{array}\right)
\left(\begin{array}{cc}\mathrm{Dsc}_{\mathbf{m}}&0\\D^{\mathrm{hl}}&D^{\mathrm{hh}}\end{array}\right)\\
&\qquad=\left(\begin{array}{cc}\mathrm{Dsc}_{\mathbf{m}}+D^{\mathrm{lh}}\mathrm{Hi}\,D^{\mathrm{hl}}&D^{\mathrm{lh}}\\D^{\mathrm{hl}}&D^{\mathrm{hh}}\end{array}\right)=D^{(p)}.
\end{aligned}
\end{equation}
The two outer factors are triangular, with identity blocks on the diagonal, and have determinant one.
Hence
\begin{equation*}
\det D^{(p)}=\det\mathrm{Dsc}_{\mathbf{m}}(p)\cdot\det D^{\mathrm{hh}}(p)
\end{equation*}
at every point of $\mathrm{Stp}$, whether $\mathrm{Dsc}_{\mathbf{m}}(p)$ is invertible or not.
If $\mathrm{Dsc}_{\mathbf{m}}(p)$ is invertible and $\mathrm{Sc}:=\mathrm{Dsc}_{\mathbf{m}}(p)^{-1}$, the inverse of (1a.17) is
\begin{equation*}
\begin{aligned}
&\left(\begin{array}{cc}1&0\\-\mathrm{Hi}\,D^{\mathrm{hl}}&1\end{array}\right)
\left(\begin{array}{cc}\mathrm{Sc}&0\\0&\mathrm{Hi}\end{array}\right)
\left(\begin{array}{cc}1&-D^{\mathrm{lh}}\mathrm{Hi}\\0&1\end{array}\right)
=\left(\begin{array}{cc}\mathrm{Sc}&0\\-\mathrm{Hi}\,D^{\mathrm{hl}}\mathrm{Sc}&\mathrm{Hi}\end{array}\right)
\left(\begin{array}{cc}1&-D^{\mathrm{lh}}\mathrm{Hi}\\0&1\end{array}\right)\\
&\qquad=\left(\begin{array}{cc}\mathrm{Sc}&-\mathrm{Sc}\,D^{\mathrm{lh}}\mathrm{Hi}\\-\mathrm{Hi}\,D^{\mathrm{hl}}\mathrm{Sc}&\mathrm{Hi}+\mathrm{Hi}\,D^{\mathrm{hl}}\mathrm{Sc}\,D^{\mathrm{lh}}\mathrm{Hi}\end{array}\right).
\end{aligned}
\end{equation*}
\end{proof}

\subsubsection*{2.9. (S9)}

\begin{proof}[Proof of Proposition~1.2\,(S9)]
Let $p\in\mathbb{R}^3$.

\emph{Splitting.}
For $v\in\mathbb{C}^4$ write $v=(v^+,v^-)\in\mathbb{C}^2\oplus\mathbb{C}^2$.
By Section~I.2.1,
\begin{equation*}
\Sigma_3=\mathrm{diag}(1_2,-1_2),\qquad
\gamma_\mu=\mathrm{diag}(\sigma_\mu,-\sigma_\mu),\qquad
\gamma_5(v^+,v^-)=(v^-,v^+).
\end{equation*}
Let $V_{\mathrm{v},2}$ be the space of the functions $O_{\mathrm{cub}}\to\mathbb{C}^2$, with the pairing $\frac{1}{48}\sum_g\overline{u'(g)}\cdot u(g)$, and let $\mathrm{Kh}:=\{u\in V_{\mathrm{v},2}\,:\,\sum_gu(g)=0\}$.
Then ${V_{\mathrm{v}}}=V_{\mathrm{v},2}\oplus V_{\mathrm{v},2}$, $\ker J^{*}=\mathrm{Kh}\oplus\mathrm{Kh}$ and $\mathbb{C}^4=\mathbb{C}^2\oplus\mathbb{C}^2$ are the decompositions into the eigenspaces of $\Sigma_3$ for the eigenvalues $+1$ and $-1$, each eigenspace being identified with $V_{\mathrm{v},2}$, $\mathrm{Kh}$, $\mathbb{C}^2$ by the coordinates.
Let $\mathrm{Vq}$ be one of the three spaces $V_{\mathrm{v},2}$, $\mathrm{Kh}$, $\mathbb{C}^2$, and $O$ an operator on $\mathrm{Vq}\oplus\mathrm{Vq}$;
on $\mathrm{Vq}\oplus\mathrm{Vq}$ the matrices $\Sigma_3$ and $\gamma_5$ act by $\Sigma_3\,(u^+,u^-)=(u^+,-u^-)$ and $\gamma_5(u^+,u^-)=(u^-,u^+)$.
\begin{itemize}
\item[(i)] $O$ commutes with $\Sigma_3$ if and only if $O=\mathrm{diag}(O^+,O^-)$ with operators $O^+$, $O^-$ on $\mathrm{Vq}$.
\item[(ii)] In that case $\gamma_5O\gamma_5=\mathrm{diag}(O^-,O^+)$ and $O^{*}=\mathrm{diag}\left((O^+)^{*},(O^-)^{*}\right)$.
\item[(iii)] Consequently, if $O$ commutes with $\Sigma_3$ and $O^{*}=\gamma_5O\gamma_5$, then $O^-=(O^+)^{*}$, $O=\mathrm{diag}\left(O^+,(O^+)^{*}\right)$, and
\begin{equation*}
\det O=\det O^+\cdot\overline{\det O^+}=|\det O^+|^2\ge0.
\end{equation*}
\end{itemize}

\emph{The three operators.}
The matrix $\Sigma_3$ commutes with $M_{i}(g)$ (m5) and with $R_i$, $R_3(p)$, which act on the labels only;
so it commutes with $D^{(p)}$.
$\Sigma_3$ and $\gamma_5$ are constant spin matrices and commute with $J$, $J^{*}$, $\mathrm{Phv}$ (Proposition~G\,(e)).
Therefore $\Sigma_3$ commutes with the four blocks, with $\mathrm{Hi}(p)$ and with $\mathrm{Dsc}_{\mathbf{m}}(p)$.
By (Fb1), $(D^{(p)})^{*}=\gamma_5D^{(p)}\gamma_5$;
consequently
\begin{align*}
D^{\mathrm{hh}}(p)^{*}&=\mathrm{Phv}\,(D^{(p)})^{*}\,\mathrm{Phv}=\gamma_5D^{\mathrm{hh}}(p)\gamma_5,&
D^{\mathrm{ll}}(p)^{*}&=J^{*}(D^{(p)})^{*}J=\gamma_5D^{\mathrm{ll}}(p)\gamma_5,\\
D^{\mathrm{lh}}(p)^{*}&=\mathrm{Phv}\,(D^{(p)})^{*}J=\gamma_5D^{\mathrm{hl}}(p)\gamma_5,&
D^{\mathrm{hl}}(p)^{*}&=J^{*}(D^{(p)})^{*}\,\mathrm{Phv}=\gamma_5D^{\mathrm{lh}}(p)\gamma_5,
\end{align*}
and
\begin{align*}
\mathrm{Hi}(p)^{*}&=\left(D^{\mathrm{hh}}(p)^{*}\right)^{-1}=\gamma_5\,\mathrm{Hi}(p)\,\gamma_5,\\
\mathrm{Dsc}_{\mathbf{m}}(p)^{*}&=D^{\mathrm{ll}}(p)^{*}-D^{\mathrm{hl}}(p)^{*}\,\mathrm{Hi}(p)^{*}\,D^{\mathrm{lh}}(p)^{*}\\
&=\gamma_5\left[D^{\mathrm{ll}}(p)-D^{\mathrm{lh}}(p)\,\mathrm{Hi}(p)\,D^{\mathrm{hl}}(p)\right]\gamma_5=\gamma_5\,\mathrm{Dsc}_{\mathbf{m}}(p)\,\gamma_5.
\end{align*}
By (iii), with $\mathrm{Hh}(p):=D^{\mathrm{hh}}(p)^{+}$ on $\mathrm{Kh}$ and $\mathrm{Ds}(p):=\mathrm{Dsc}_{\mathbf{m}}(p)^{+}$ on $\mathbb{C}^2$,
\begin{gather*}
D^{\mathrm{hh}}(p)=\mathrm{diag}\left(\mathrm{Hh}(p),\mathrm{Hh}(p)^{*}\right),\qquad
\mathrm{Dsc}_{\mathbf{m}}(p)=\mathrm{diag}\left(\mathrm{Ds}(p),\mathrm{Ds}(p)^{*}\right),\\
\det D^{\mathrm{hh}}(p)=|\det\mathrm{Hh}(p)|^2,\qquad
\det\mathrm{Dsc}_{\mathbf{m}}(p)=|\det\mathrm{Ds}(p)|^2\ge0.
\end{gather*}

\emph{Strict signs.}
By (S1) the operator $D^{\mathrm{hh}}(p)$ is invertible;
hence $\det D^{\mathrm{hh}}(p)>0$ for all real $p$ and all $\mathbf{m}\in[0,1]$.
Let $\tilde{p}\in[-\pi,\pi]^3$ be congruent to $p$ modulo $2\pi\mathbb{Z}^3$;
by the periodicity in (S3), $\mathrm{Dsc}_{\mathbf{m}}(p)=\mathrm{Dsc}_{\mathbf{m}}(\tilde{p})$.
If $\mathbf{m}>0$ or $p\notin2\pi\mathbb{Z}^3$, then $\mathrm{rp}(\tilde{p})>0$, and $\mathrm{Dsc}_{\mathbf{m}}(\tilde{p})$ is injective by (S5);
so $\det\mathrm{Dsc}_{\mathbf{m}}(p)>0$.
If $\mathbf{m}=0$ and $p\in2\pi\mathbb{Z}^3$, then $\mathrm{Dsc}_0(p)=\mathrm{Dsc}_0(0)=0$ (2.3), and the determinant vanishes.
\end{proof}

\begin{proof}[Proof of the uniformity in $\mathbf{m}$ stated in Proposition~1.2]
The numbers $\theta_{0}$, $\theta_{0}'$, $K_{\rho}$, $t_{0}$, $c_{\mathrm{el}}$, $K_D$, $K_{\mathrm{Str}}$, $\kappa_{(0)}$, $A_{\mathrm{ak}}$, $A'_{\mathrm{ak}}$ are formed with $\rho$ and the weights;
none of them contains $\mathbf{m}$.
Every bound in (S1)--(S9) has been proved for an arbitrary $\mathbf{m}\in[0,1]$ with these numbers;
hence all bounds of Proposition~1.2 are uniform in $\mathbf{m}\in[0,1]$.
\end{proof}

\begin{proof}[Proof of the last sentence of (I.O8)]
At $\boldsymbol{\mu}=0$ the symbol $D^{(0)}$ is $D_0^{(0)}$ (0.2\,(In3)).
By (Fb7) at $p=0$ one has $D_0^{(0)}J=0$, and (Fb1) gives
\begin{equation*}
J^{*}D_0^{(0)}=\left(\left(D_0^{(0)}\right)^{*}J\right)^{*}=\left(\gamma_5\,D_0^{(0)}J\,\gamma_5\right)^{*}=0.
\end{equation*}
\end{proof}

\subsection*{3. Proposition 1.3}

\begin{proof}[Proof of Proposition~1.3]
The data are the instance (In2) of (GS):
$\mathrm{Lab}=\mathrm{me}\,\mathbb{Z}^3$, $A=0$, $D=D_0$ and $\nabla^{*}\nabla=\nabla_0^{*}\nabla_0$.

\emph{(a).}
This is (GS5) at $A=0$.

\emph{(c).}
Let $\psi$ be a function on $\mathrm{me}\,\mathbb{Z}^3$ and $u$ a function on $\mathrm{me}\,\mathbb{Z}^3\times O_{\mathrm{cub}}$, let $x=(c,g)$, and write $\pi(x^{(i)})$ for the photon site of $x^{(i)}$.
By the definition of $D$ in Section~G1.4,
\begin{align*}
(D(\psi u))(x)-\psi(c)\,(Du)(x)
&=-\mathrm{me}^{-1}\sum_iM_{i}(x)\left[\psi(\pi(x^{(i)}))\,u(x^{(i)})-\psi(c)\,u(x)\right]\\
&\qquad+\mathrm{me}^{-1}\sum_iM_{i}(x)\left[\psi(c)\,u(x^{(i)})-\psi(c)\,u(x)\right]\\
&=-\mathrm{me}^{-1}\sum_iM_{i}(x)\left(\psi(\pi(x^{(i)}))-\psi(c)\right)u(x^{(i)}).
\end{align*}
Since $\pi(x^{(i)})=c$ for $i=0,1,2$ and $\pi(x^{(3)})=c_{3}$, only the term $i=3$ remains, and this is the identity of (c);
neither periodicity nor summability has been used.
The map $x\mapsto x^{(3)}$ is a bijection by (s1).
Further $c_{3}-c=\mathrm{me}\cdot ge_1$ with $ge_1\in\{\pm e_\mu\}$, so that $|c_{3}-c|=\mathrm{me}$;
and $|M_3(x)|_2=3(1+\rho)$ by (m6).
If $ge_1=\pm e_\mu$, then $M_3(x)=3(\rho\mp\gamma_\mu)$.
For $\mu=1,2$ each row of $\gamma_\mu$ has one entry of modulus one, off the diagonal, and no other non-zero entry;
so every row sum of moduli of $M_3(x)$ equals $3(\rho+1)$.
For $\mu=3$ one has $\gamma_3=\mathrm{diag}(1,-1,-1,1)$, and the row sums are $3|\rho\mp1|$, the larger of which is $3(1+\rho)$, because $\rho>0$.
Hence $|M_3(x)|_c=3(1+\rho)$, which is (m7), and $3(1+\rho)\le6$ because $\rho<1$.

\emph{(d).}
The identity $D^{*}=\gamma_5D\gamma_5$ on $\mathrm{HH}(\mathrm{me},R;\theta)$ is (GS5) at $A=0$.
For the transpose, the pairing is $\langle f',f\rangle^{\mathrm{ch}}=\int\mathrm{d}x\,f'(x)\cdot f(x)$ (Section~G1.4).
In the case $R<\infty$ let $f'\in\mathfrak{S}^{\mathrm{ch}}(\mathrm{me},R;-\theta)$ and $f\in\mathfrak{S}^{\mathrm{ch}}(\mathrm{me},R;\theta)$.
The function $f'(x)\cdot f(x)$ is $R\mathbb{Z}^3$-periodic, because $\left(-\mathrm{e}^{-i\theta_\mu}\right)\left(-\mathrm{e}^{i\theta_\mu}\right)=1$;
the same holds for all products below, since $D$ maps each of the two spaces into itself (0.1).
In the case $R=\infty$, $\mathfrak{S}^{\mathrm{ch}}(\mathrm{me},\infty;-\theta)$ is the set of all functions on $\mathrm{me}\,\mathbb{Z}^3\times O_{\mathrm{cub}}$;
let $f'$ be one of them and $f$ a function with finite support.
Then $Df$ has finite support, and all sums below are finite.
In both cases one uses $\mathrm{vv}'\cdot(\mathrm{Mx}\,\mathrm{vv})=(\mathrm{Mx}^T\mathrm{vv}')\cdot\mathrm{vv}$ for a $4\times4$ matrix $\mathrm{Mx}$ and $\mathrm{vv},\mathrm{vv}'\in\mathbb{C}^4$, and, in the terms that contain $f(x^{(i)})$, the substitution $x\mapsto x^{(i)}$, a bijection of the classes of sites with $(x^{(i)})^{(i)}=x$;
this gives
\begin{align*}
\langle f',Df\rangle^{\mathrm{ch}}
&=-\mathrm{me}^{-1}\sum_i\int\mathrm{d}x\,f'(x)\cdot M_{i}(x)\left[f(x^{(i)})-f(x)\right]\\
&=-\mathrm{me}^{-1}\sum_i\int\mathrm{d}x\,\left[M_{i}(x^{(i)})^Tf'(x^{(i)})-M_{i}(x)^Tf'(x)\right]\cdot f(x).
\end{align*}
The admitted functions $f$ separate the functions $f'$.
Indeed, for $R<\infty$ the function $\overline{f'}$ lies in $\mathfrak{S}^{\mathrm{ch}}(\mathrm{me},R;\theta)$, and $\langle f',\overline{f'}\rangle^{\mathrm{ch}}=\|f'\|_{L^2}^2$;
for $R=\infty$ and the function $f$ that equals the $\beta$-th standard basis vector of $\mathbb{C}^4$ at one site $x$ and vanishes at all other sites, $\langle f',f\rangle^{\mathrm{ch}}=(\mathrm{me}^3/48)\cdot f'_\beta(x)$.
So the function $D^Tf'$ is determined by $\langle D^Tf',f\rangle^{\mathrm{ch}}=\langle f',Df\rangle^{\mathrm{ch}}$ for all admitted $f$, and
\begin{equation*}
(D^Tf')(x)=-\mathrm{me}^{-1}\sum_i\left[M_{i}(x^{(i)})^Tf'(x^{(i)})-M_{i}(x)^Tf'(x)\right],
\end{equation*}
a function in $\mathfrak{S}^{\mathrm{ch}}(\mathrm{me},R;-\theta)$.
For $R=\infty$ and $f',f\in\mathrm{HH}$ all sums converge absolutely and the same identity holds;
hence the restriction of $D^T$ to $\mathrm{HH}$ is the transpose of $D$ for the pairing of $\mathrm{HH}$ with $\mathrm{HH}$.
By (m2) at $x$ and at $x^{(i)}$,
\begin{equation*}
M_{i}(x)^T=\mathsf{C}\,M_{i}(x^{(i)})\,\mathsf{C}^{-1},\qquad
M_{i}(x^{(i)})^T=\mathsf{C}\,M_{i}(x)\,\mathsf{C}^{-1},
\end{equation*}
and by (m1), $\sum_iM_{i}(x^{(i)})=\sum_iM_{i}(x)$.
Therefore
\begin{align*}
(D^Tf')(x)
&=-\mathrm{me}^{-1}\,\mathsf{C}\sum_i\left[M_{i}(x)\,\mathsf{C}^{-1}f'(x^{(i)})-M_{i}(x^{(i)})\,\mathsf{C}^{-1}f'(x)\right]\\
&=-\mathrm{me}^{-1}\,\mathsf{C}\sum_iM_{i}(x)\left[(\mathsf{C}^{-1}f')(x^{(i)})-(\mathsf{C}^{-1}f')(x)\right]=(\mathsf{C}\,D\,\mathsf{C}^{-1}f')(x).
\end{align*}
Since $\mathsf{C}^{-1}=-\mathsf{C}$ (c3), one has $\mathsf{C}\,D\,\mathsf{C}^{-1}=(-\mathsf{C}^{-1})\,D\,(-\mathsf{C})=\mathsf{C}^{-1}D\,\mathsf{C}$.

\emph{(b).}
The proof has three steps: a lower bound for the Schur symbol, an inequality in each fibre, and the Fourier decomposition.

\emph{Step 1.}
For all $p\in\mathbb{R}^3$ and $v\in\mathbb{C}^4$,
\begin{equation}\tag{1a.18}
|\mathrm{Dsc}_0(p)\,v|_2\ge c_{\mathrm{lo}}\cdot\mathrm{tw}(p)^{\frac{1}{2}}\cdot|v|_2,\qquad
c_{\mathrm{lo}}:=\frac{\rho\cdot c_{*}}{2K_{\rho}}.
\end{equation}
Let $t_\mu:=1-\cos p_\mu\in[0,2]$.
Suppose first that $\mathrm{tw}(p)\le t_{0}$.
Then $t_\mu\le t_{0}\le1$ ($t_{0}\le1/4$ by 2.5\,(i)), so that $|\sin p|_2^2=\sum_\mu t_\mu(2-t_\mu)\ge\mathrm{tw}(p)$;
further $|i\gamma\cdot\sin p\,v|_2=|\sin p|_2\cdot|v|_2$ by (1a.14) at $\mathbf{m}=0$, whose proof holds for every real $p$.
By (S3) at $\mathbf{m}=0$,
\begin{align*}
|\mathrm{Dsc}_0(p)\,v|_2
&\ge|\sin p|_2\cdot|v|_2-K_{\rho}\cdot\mathrm{tw}(p)\cdot|v|_2\\
&\ge\mathrm{tw}(p)^{\frac{1}{2}}\cdot|v|_2-K_{\rho}\cdot t_{0}^{\frac{1}{2}}\cdot\mathrm{tw}(p)^{\frac{1}{2}}\cdot|v|_2=\frac{1}{2}\,\mathrm{tw}(p)^{\frac{1}{2}}\cdot|v|_2,
\end{align*}
and $\frac{1}{2}\ge c_{\mathrm{lo}}$, because $\rho\cdot c_{*}\le c_{*}\le1/3<1<K_{\rho}$ ($\rho<1$; 2.5\,(i), (ii)).
Suppose next that $\mathrm{tw}(p)\ge t_{0}$.
Then (S4) at $\mathbf{m}=0$ gives
\begin{align*}
|\mathrm{Dsc}_0(p)\,v|_2\cdot|v|_2
&\ge\mathrm{Re}\,(v,\mathrm{Dsc}_0(p)\,v)\ge\rho\cdot c_{*}\cdot\mathrm{tw}(p)\cdot|v|_2^2\\
&\ge\rho\cdot c_{*}\cdot t_{0}^{\frac{1}{2}}\cdot\mathrm{tw}(p)^{\frac{1}{2}}\cdot|v|_2^2=c_{\mathrm{lo}}\cdot\mathrm{tw}(p)^{\frac{1}{2}}\cdot|v|_2^2,
\end{align*}
since $t_{0}^{\frac{1}{2}}=1/(2K_{\rho})$.

\emph{Step 2.}
For all $p\in\mathbb{R}^3$ and $u\in{V_{\mathrm{v}}}$,
\begin{equation}\tag{1a.19}
\left\|D_0^{(p)}u\right\|^2\ge c_{b}\cdot(u,(\nabla^{*}\nabla)(p)\,u).
\end{equation}
The notation is that of Proposition~1.2 at $\mathbf{m}=0$:
the symbol $D^{(p)}$ of Proposition~1.2 is then $D_0^{(p)}$, the blocks $D^{\mathrm{ll}}$, $D^{\mathrm{lh}}$, $D^{\mathrm{hl}}$, $D^{\mathrm{hh}}$ are taken at $p$, and $\mathrm{Hi}=(D^{\mathrm{hh}})^{-1}$.
Let
\begin{equation*}
v:=J^{*}u,\qquad
\mathrm{hp}:=\mathrm{Phv}\,u,\qquad
\mathrm{fs}:=D_0^{(p)}u,\qquad
\mathrm{fs}_{\mathrm{l}}:=J^{*}\mathrm{fs},\qquad
\mathrm{fs}_{\mathrm{h}}:=\mathrm{Phv}\,\mathrm{fs}.
\end{equation*}
Then $\|\mathrm{fs}\|^2=|\mathrm{fs}_{\mathrm{l}}|_2^2+\|\mathrm{fs}_{\mathrm{h}}\|^2$ (2.0), and
\begin{equation*}
\mathrm{fs}_{\mathrm{l}}=D^{\mathrm{ll}}v+D^{\mathrm{lh}}\mathrm{hp},\qquad
\mathrm{fs}_{\mathrm{h}}=D^{\mathrm{hl}}v+D^{\mathrm{hh}}\mathrm{hp}.
\end{equation*}
The second equation gives $\mathrm{hp}=\mathrm{Hi}\left(\mathrm{fs}_{\mathrm{h}}-D^{\mathrm{hl}}v\right)$;
inserted into the first, it gives $\mathrm{Dsc}_0(p)\,v=\mathrm{fs}_{\mathrm{l}}-D^{\mathrm{lh}}\mathrm{Hi}\,\mathrm{fs}_{\mathrm{h}}$.
By (S2) and $\mathrm{tw}(p)\le6$,
\begin{equation*}
|D^{\mathrm{lh}}|\le c_{\mathrm{hl}}\cdot\left(2\,\mathrm{tw}(p)\right)^{\frac{1}{2}}\le c_{\mathrm{hl}}\sqrt{12}=6(1+\rho),
\end{equation*}
and by (S1) at $\mathbf{m}=0$, $|\mathrm{Hi}|\le1/\hat{g}$.
Together with (1a.18) this gives
\begin{align*}
c_{\mathrm{lo}}\cdot\mathrm{tw}(p)^{\frac{1}{2}}\cdot|v|_2
&\le|\mathrm{Dsc}_0(p)\,v|_2\le|\mathrm{fs}_{\mathrm{l}}|_2+\frac{6(1+\rho)}{\hat{g}}\,\|\mathrm{fs}_{\mathrm{h}}\|\le\left(1+\frac{6(1+\rho)}{\hat{g}}\right)\|\mathrm{fs}\|,
\end{align*}
that is,
\begin{align}
\left(2\,\mathrm{tw}(p)\right)^{\frac{1}{2}}|v|_2
&\le\frac{\sqrt{2}}{c_{\mathrm{lo}}}\left(1+\frac{6(1+\rho)}{\hat{g}}\right)\|\mathrm{fs}\|=A_1\,\|\mathrm{fs}\|,\tag{1a.20}\\
\|\mathrm{hp}\|
&\le\frac{1}{\hat{g}}\left(\|\mathrm{fs}_{\mathrm{h}}\|+c_{\mathrm{hl}}\cdot\left(2\,\mathrm{tw}(p)\right)^{\frac{1}{2}}|v|_2\right)\le\frac{1}{\hat{g}}\left(1+c_{\mathrm{hl}}\cdot A_1\right)\|\mathrm{fs}\|=A_2\,\|\mathrm{fs}\|.\tag{1a.21}
\end{align}
Here $\sqrt{2}/c_{\mathrm{lo}}=2\sqrt{2}\,K_{\rho}/(\rho\cdot c_{*})$, $|\mathrm{fs}_{\mathrm{l}}|_2\le\|\mathrm{fs}\|$ and $\|\mathrm{fs}_{\mathrm{h}}\|\le\|\mathrm{fs}\|$.
By (GS10) in the instance (In3), together with (1a.1),
\begin{align*}
(u,(\nabla^{*}\nabla)(p)\,u)
&\le3\cdot2\,\mathrm{tw}(p)\cdot|v|_2^2+\left(30+24\sqrt{2}\right)\|\mathrm{hp}\|^2\\
&\le\left[3A_1^2+\left(30+24\sqrt{2}\right)A_2^2\right]\|\mathrm{fs}\|^2=c_{b}^{-1}\left\|D_0^{(p)}u\right\|^2.
\end{align*}

\emph{Step 3.}
Let $f\in\mathrm{HH}(\mathrm{me},R;\theta)$.

\emph{The case $R<\infty$.}
Put $n:=R/\mathrm{me}$, and let $P_{n,\theta}$ be the set of the $n^3$ momenta $p$ with $p_\mu=(2\pi j_\mu+\pi+\theta_\mu)/n$, $j\in\{0,\dots,n-1\}^3$ (Section~I.2.6).
For $p\in P_{n,\theta}$ the function $\epsilon_p(c):=\exp(ip\cdot c/\mathrm{me})$ on $\mathrm{me}\,\mathbb{Z}^3$ satisfies
\begin{equation*}
\epsilon_p(c+R\,e_\mu)=\mathrm{e}^{ip_\mu n}\,\epsilon_p(c)=-\mathrm{e}^{i\theta_\mu}\,\epsilon_p(c).
\end{equation*}
The space of the scalar functions on $\mathrm{me}\,\mathbb{Z}^3$ with this property has dimension $n^3$.
For $p,\tilde{p}\in P_{n,\theta}$, with the sum over the period cell $c=\mathrm{me}\,j$, $j\in\{0,\dots,n-1\}^3$,
\begin{equation*}
\sum_c\overline{\epsilon_{\tilde{p}}(c)}\,\epsilon_p(c)=\prod_\mu\sum_{j=0}^{n-1}\mathrm{e}^{i(p_\mu-\tilde{p}_\mu)j}=n^3\,\mathbf{1}[p=\tilde{p}],
\end{equation*}
because $p_\mu-\tilde{p}_\mu$ lies in $(2\pi/n)\mathbb{Z}$ and has modulus less than $2\pi$.
So the functions $\epsilon_p$ form an orthogonal basis of this space.
Consequently, if $\hat{f}$ is defined by the first line of the following display, the second line holds:
\begin{equation}\tag{1a.22}
\begin{aligned}
&\hat{f}(p)(g):=n^{-3}\sum_{c\bmod R\mathbb{Z}^3}\exp(-ip\cdot c/\mathrm{me})\,f(c,g)\in{V_{\mathrm{v}}},\\
&f(c,g)=\sum_{p\in P_{n,\theta}}\exp(ip\cdot c/\mathrm{me})\,\hat{f}(p)(g),\qquad
(f',f)^{\mathrm{ch}}=R^3\sum_{p\in P_{n,\theta}}\left(\hat{f}'(p),\hat{f}(p)\right),
\end{aligned}
\end{equation}
where $\mathrm{me}^3\cdot n^3=R^3$ has been used.
On a function $\exp(ip\cdot c/\mathrm{me})\,u(g)$ the hops act by $u(g)\mapsto u(gs_i)$ for $i\le2$, and by
\begin{equation*}
\exp\left(ip\cdot(c+\mathrm{me}\cdot ge_1)/\mathrm{me}\right)\,u(gs_3)=\exp(ip\cdot c/\mathrm{me})\,(R_3(p)u)(g)
\end{equation*}
for $i=3$.
Hence
\begin{equation}\tag{1a.23}
\widehat{Df}(p)=\mathrm{me}^{-1}\,D_0^{(p)}\hat{f}(p),\qquad
\widehat{\nabla^{*}\nabla f}(p)=(\nabla^{*}\nabla)(p)\,\hat{f}(p).
\end{equation}

\emph{The case $R=\infty$.}
Here $\mathrm{HH}=\ell^2(\mathrm{me}\,\mathbb{Z}^3\times O_{\mathrm{cub}};\mathbb{C}^4)$.
Put $\hat{f}(p)(g):=\sum_c\exp(-ip\cdot c/\mathrm{me})\,f(c,g)$ for $p\in[-\pi,\pi]^3$, the series converging in $L^2([-\pi,\pi]^3)$.
Parseval's identity for Fourier series in three variables \cite[Chapter~VII, Section~1]{SteinWeiss1971}, applied to each of the $192$ components $(g,\beta)$, gives
\begin{equation*}
(f',f)^{\mathrm{ch}}=\mathrm{me}^3\int_{[-\pi,\pi]^3}\frac{\mathrm{d}^3p}{(2\pi)^3}\,\left(\hat{f}'(p),\hat{f}(p)\right).
\end{equation*}
The transform of $c\mapsto f(c+\mathrm{me}\cdot a,g')$ is $\mathrm{e}^{ip\cdot a}\hat{f}(p)(g')$.
So the transform of $(R_3f)(c,g)=f(c+\mathrm{me}\cdot ge_1,gs_3)$ is $\mathrm{e}^{ip\cdot(ge_1)}\hat{f}(p)(gs_3)=(R_3(p)\hat{f}(p))(g)$, that of $(R_if)(c,g)=f(c,gs_i)$ is $(R_i\hat{f}(p))(g)$, and (1a.23) holds for almost every $p$.

\emph{Conclusion.}
Let $\int\mathrm{dm}(p):=R^3\sum_{p\in P_{n,\theta}}$ in the first case and $\int\mathrm{dm}(p):=\mathrm{me}^3\int_{[-\pi,\pi]^3}\mathrm{d}^3p/(2\pi)^3$ in the second.
In both cases (1a.22), respectively Parseval's identity, together with (1a.23) and with (1a.19), all momenta being real, gives
\begin{align*}
\|Df\|_{L^2}^2
&=\mathrm{me}^{-2}\int\mathrm{dm}(p)\,\left\|D_0^{(p)}\hat{f}(p)\right\|^2\\
&\ge c_{b}\cdot\mathrm{me}^{-2}\int\mathrm{dm}(p)\,\left(\hat{f}(p),(\nabla^{*}\nabla)(p)\,\hat{f}(p)\right)=c_{b}\cdot\mathrm{me}^{-2}\cdot(f,\nabla^{*}\nabla f)^{\mathrm{ch}}.
\end{align*}
For $f=u$ this is (b).

\emph{Constants; uniformity.}
The constant $c_{b}$ is formed with $\rho$, $c_{*}$, $K_{\rho}$, $\hat{g}$, $c_{\mathrm{hl}}$.
The mesh $\mathrm{me}>0$ is arbitrary, and no constant and no step depends on $\mathrm{me}$, $R$ or $\theta$.
\end{proof}

\subsection*{4. Proposition 1.4}

\emph{The objects of Section~G1.4 on a cube.}
Let $E=\{z+\mathrm{me}\cdot(j_1,j_2,j_3)\,:\,1\le j_\mu\le n\}$ with some $z\in\mathrm{me}\,\mathbb{Z}^3$.
The quantities $P_Eu$, $\|u\|_{L^2,E}$ and $(u,\nabla^{*}\nabla u)_E$ are those of the paragraph \emph{Blocks} of Section~G1.4, whose last sentence extends them to cubes of $n^3$ sites;
one has $\#|E\times O_{\mathrm{cub}}|=48n^3$.
Moreover $\|u\|_{L^2,E}^2$ and $(u,\nabla^{*}\nabla u)_E$ are the quantities $\|f\|_{L^2,E}^2$ and $(f,\nabla_A^{*}\nabla_Af)_E$ of 1.1 for $f=u$ and $A=0$, in the reading (Es2).
For a function $\xi:E\to\mathbb{C}^4$ the norm is $\|\xi\|_{L^2,E}^2:=\mathrm{me}^3\sum_{c\in E}|\xi(c)|_2^2$;
$P_E\xi$ is the mean of $\xi$ over the $n^3$ sites of $E$, taken as a constant function on $E$;
and, by 1.1 at $A=0$,
\begin{equation*}
\|\mathrm{dc}_0\xi\|_E^2=\mathrm{me}^3\sum_{\mu}\ \sum_{c\,:\,c,\,c+\mathrm{me}\cdot e_\mu\in E}\mathrm{me}^{-2}\,\left|\xi(c+\mathrm{me}\cdot e_\mu)-\xi(c)\right|_2^2.
\end{equation*}

\begin{named}{Lemma P}
For every function $\xi:E\to\mathbb{C}^4$,
\begin{equation*}
\|\xi-P_E\xi\|_{L^2,E}^2\le\frac{\ell_E^2}{8}\,\|\mathrm{dc}_0\xi\|_E^2.
\end{equation*}
\end{named}

\begin{named}{Remark}
For the blocks of \cite[Section~12.1]{QEDI}, whose sides contain a power of $L$ sites, Lemma~P is \cite[Lemma~12.3]{QEDI}.
Here $n\ge1$ is an arbitrary integer;
the proof is the same.
\end{named}

\begin{proof}[Proof of Lemma~P]
\emph{One dimension.}
Let $\mathrm{vv}_1,\dots,\mathrm{vv}_n\in\mathbb{C}^4$ and $\mathrm{vvb}:=\frac{1}{n}\sum_i\mathrm{vv}_i$.
One has
\begin{equation*}
\sum_{i,j}|\mathrm{vv}_i-\mathrm{vv}_j|_2^2=2n\sum_i|\mathrm{vv}_i|_2^2-2\left|\sum_i\mathrm{vv}_i\right|_2^2=2n\sum_i|\mathrm{vv}_i-\mathrm{vvb}|_2^2,
\end{equation*}
and $\mathrm{vv}_j-\mathrm{vv}_i=\sum_{i\le\mathrm{ib}<j}(\mathrm{vv}_{\mathrm{ib}+1}-\mathrm{vv}_{\mathrm{ib}})$ for $i<j$.
With the Cauchy--Schwarz inequality it follows that
\begin{align*}
\sum_i|\mathrm{vv}_i-\mathrm{vvb}|_2^2
&=\frac{1}{n}\sum_{i<j}|\mathrm{vv}_i-\mathrm{vv}_j|_2^2
\le\frac{1}{n}\sum_{i<j}(j-i)\sum_{i\le\mathrm{ib}<j}|\mathrm{vv}_{\mathrm{ib}+1}-\mathrm{vv}_{\mathrm{ib}}|_2^2\\
&=\frac{1}{n}\sum_{\mathrm{ib}=1}^{n-1}|\mathrm{vv}_{\mathrm{ib}+1}-\mathrm{vv}_{\mathrm{ib}}|_2^2\sum_{i=1}^{\mathrm{ib}}\ \sum_{j=\mathrm{ib}+1}^{n}(j-i).
\end{align*}
The inner double sum equals
\begin{align*}
\sum_{i=1}^{\mathrm{ib}}\ \sum_{j=\mathrm{ib}+1}^{n}(j-i)
&=\frac{\mathrm{ib}\cdot\left[n(n+1)-\mathrm{ib}\cdot(\mathrm{ib}+1)\right]}{2}-\frac{(n-\mathrm{ib})\cdot\mathrm{ib}\cdot(\mathrm{ib}+1)}{2}\\
&=\frac{\mathrm{ib}}{2}\left[n^2-n\cdot\mathrm{ib}\right]=\frac{\mathrm{ib}\cdot(n-\mathrm{ib})\cdot n}{2}\le\frac{n^3}{8}.
\end{align*}
Therefore
\begin{equation}\tag{1a.24}
\sum_i|\mathrm{vv}_i-\mathrm{vvb}|_2^2\le\frac{n^2}{8}\sum_{\mathrm{ib}=1}^{n-1}|\mathrm{vv}_{\mathrm{ib}+1}-\mathrm{vv}_{\mathrm{ib}}|_2^2.
\end{equation}

\emph{Three dimensions.}
For $\mu=1,2,3$ let $\mathrm{Avl}_\mu$ be the operator on the functions on $E$ that replaces $\xi(c)$ by the mean of $\xi$ over the $n$ sites of the line of $E$ through $c$ in the direction $e_\mu$.
The operators $\mathrm{Avl}_1$, $\mathrm{Avl}_2$, $\mathrm{Avl}_3$ are commuting orthogonal projections on the Hilbert space of the functions $E\to\mathbb{C}^4$ with the norm $\|\cdot\|_{L^2,E}$, being means over the fibres of a product;
moreover $P_E=\mathrm{Avl}_1\mathrm{Avl}_2\mathrm{Avl}_3$, and
\begin{equation*}
1-P_E=(1-\mathrm{Avl}_1)+\mathrm{Avl}_1(1-\mathrm{Avl}_2)+\mathrm{Avl}_1\mathrm{Avl}_2(1-\mathrm{Avl}_3).
\end{equation*}
The three terms have mutually orthogonal ranges:
the second and the third term begin with $\mathrm{Avl}_1$, and $(1-\mathrm{Avl}_1)\,\mathrm{Avl}_1=0$;
further
\begin{equation*}
\left(\mathrm{Avl}_1(1-\mathrm{Avl}_2)\right)^{*}\mathrm{Avl}_1\mathrm{Avl}_2(1-\mathrm{Avl}_3)=\mathrm{Avl}_1(1-\mathrm{Avl}_2)\,\mathrm{Avl}_2(1-\mathrm{Avl}_3)=0.
\end{equation*}
As projections do not increase the norm,
\begin{align*}
\|\xi-P_E\xi\|_{L^2,E}^2
&=\|(1-\mathrm{Avl}_1)\xi\|^2+\|\mathrm{Avl}_1(1-\mathrm{Avl}_2)\xi\|^2+\|\mathrm{Avl}_1\mathrm{Avl}_2(1-\mathrm{Avl}_3)\xi\|^2\\
&\le\sum_\mu\|(1-\mathrm{Avl}_\mu)\xi\|_{L^2,E}^2.
\end{align*}
Applying (1a.24) on each line in the direction $e_\mu$ and summing over these lines, one obtains
\begin{equation*}
\|(1-\mathrm{Avl}_\mu)\xi\|_{L^2,E}^2\le\frac{n^2}{8}\cdot\mathrm{me}^3\sum_{c\,:\,c,\,c+\mathrm{me}\cdot e_\mu\in E}\left|\xi(c+\mathrm{me}\cdot e_\mu)-\xi(c)\right|_2^2.
\end{equation*}
The sum of the right sides over $\mu$ is $(n^2\cdot\mathrm{me}^2/8)\,\|\mathrm{dc}_0\xi\|_E^2$, and $n\cdot\mathrm{me}=\ell_E$.
\end{proof}

\begin{proof}[Proof of Proposition~1.4]
Let $\xi:=J^{*}u$, a function on $E$, and $\mathrm{hp}:=u-J\xi=\mathrm{Phv}\,u$, a function on $E\times O_{\mathrm{cub}}$.
The mean of $u$ over $E\times O_{\mathrm{cub}}$ is
\begin{equation*}
\frac{1}{n^3}\sum_{c\in E}\frac{1}{48}\sum_gu(c,g)=\frac{1}{n^3}\sum_{c\in E}\xi(c);
\end{equation*}
hence $P_Eu=J\,P_E\xi$, and
\begin{equation*}
u-P_Eu=J(\xi-P_E\xi)+\mathrm{hp}.
\end{equation*}
At each $c\in E$ the cluster mean of $u-P_Eu$ is $\xi(c)-P_E\xi$ and its heavy part is $\mathrm{hp}$;
applying (1a.3) to $u-P_Eu$ at each $c\in E$, multiplying by $\mathrm{me}^3$ and summing over $c\in E$, one finds
\begin{equation}\tag{1a.25}
\|u-P_Eu\|_{L^2,E}^2=\|\xi-P_E\xi\|_{L^2,E}^2+\|\mathrm{hp}\|_{L^2,E}^2.
\end{equation}
Proposition~G\,(d) in the instance (In2), that is with $A=0$, and in the reading (Es2) is the local form of Proposition~1.1\,(d) for a finite set $E$ of sites of $\mathrm{me}\,\mathbb{Z}^3$ and a function on $E\times O_{\mathrm{cub}}$.
Its first inequality, for $f=u$, and Lemma~P give
\begin{equation*}
\|\xi-P_E\xi\|_{L^2,E}^2\le\frac{\ell_E^2}{8}\,\|\mathrm{dc}_0\xi\|_E^2\le\frac{\ell_E^2}{8c_{*}}\cdot\mathrm{me}^{-2}\cdot(u,\nabla^{*}\nabla u)_E.
\end{equation*}
Its second inequality and $\ell_E=n\cdot\mathrm{me}\ge\mathrm{me}$ give
\begin{equation*}
\|\mathrm{hp}\|_{L^2,E}^2\le\lambda_{\mathrm{int}}^{-1}\cdot(u,\nabla^{*}\nabla u)_E\le\lambda_{\mathrm{int}}^{-1}\cdot\ell_E^2\cdot\mathrm{me}^{-2}\cdot(u,\nabla^{*}\nabla u)_E.
\end{equation*}
Inserting the sum of the two bounds into (1a.25) gives the assertion, with $C_{\mathrm{Po}}=1/(8c_{*})+1/\lambda_{\mathrm{int}}$.
This constant is formed with the weights only, since $c_{*}$ and $\lambda_{\mathrm{int}}$ do not contain $\rho$;
it depends neither on $\mathrm{me}$ nor on $n$ or $z$.
\end{proof}

\subsection*{5. Conditions on the constants}

Proposition~G and Propositions~1.1--1.4 hold for every choice of the constants of Section~I.8, in the form of Proposition~14.2\,(a), that has the following properties.

\medskip
\noindent
\begin{tabular}{p{0.17\textwidth}p{0.57\textwidth}p{0.14\textwidth}}
\emph{Property} & \emph{Where it enters} & \emph{Line of Section~I.8.2}\\
\hline
$\rho>0$ &
(m6): $|M_{i}(x)|_2=(w_i/2)(1+\rho)$, from $|\rho-1|\le1+\rho$;
through (m6): the norm of $D_A$ in (GS5), (Fb1), $K_D$;
$c_{\mathrm{hl}}$ in (1a.8), hence Proposition~G\,(b), (S2), $K_{\rho}$, $A_1$, $A_2$;
(Fb5), $\delta(p)$, hence $\theta_{0}$, $\theta_{0}'$.
Proposition~1.3\,(c): $|M_3(x)|_2$ and $|M_3(x)|_c$ (the row sums $3(\rho+1)$; the largest of the numbers $3|\rho\mp1|$).
An inequality is multiplied by $\rho$: Proposition~G\,(c); 2.4\,(S4); $\rho\cdot\mathrm{tw}(p)\ge0$ in 2.3; $K_{\rho}\ge2c_{\mathrm{hl}}^2/\hat{g}$ in 2.5\,(i).
$\hat{g}>0$: (S1); $\theta_{0}>0$; $\theta_{0}<\theta_{0}'$; division by $\hat{g}$ in $K_{\rho}$, $2/\hat{g}$, $A_1$, $A_2$.
$c_{\mathrm{el}}>0$, $c_{\mathrm{lo}}>0$: division in $A_{\mathrm{ak}}$, $A'_{\mathrm{ak}}$, $A_1$; $\kappa_{(0)}>0$ &
(C0), ``$\rho\in(0,1)$''\\
\hline
$\rho<1$ &
$3(1+\rho)\le6$ in Proposition~1.3\,(c); $\rho\cdot c_{*}\le c_{*}$ in Section~3, Step~1 (for $c_{\mathrm{lo}}\le1/2$) &
(C0)\\
\hline
$w=(6,6\sqrt{2},6\sqrt{2},6)$, with $\sum_iw_i\hat n_i=0$ (I.1.1) &
(GS2) through (s3), hence (GS5) and the real part everywhere;
Proposition~1.3\,(d) through (m1);
the numerals $3=w_3/2$, $1/2=(8/48)(w_3/2)$, $18=(3/8)\cdot48$, $30+24\sqrt{2}=3w_3+2(w_0+w_1+w_2)$;
$2(w_0+w_1+w_2)=12+24\sqrt{2}<48$ in (GS7);
$w_3/18=1/3$ in 2.5\,(ii);
$\mathrm{sw}$ in (GS5), $K_D$ &
(C0)\\
\hline
the definitions of $\lambda_{\mathrm{int}}$, $\hat{g}$, $c_{\mathrm{hl}}$, $c_{*}$, $\theta_{0}$, $K_{\rho}$, $t_{0}$, $c_{\mathrm{el}}$, $K_{\mathrm{Str}}$, $\kappa_{(0)}$, $c_{b}$ (with $A_1$, $A_2$), $C_{\mathrm{Po}}$ &
as stated in Propositions~1.1--1.4 &
(C0r)\\
\hline
$L$ is an integer &
0.2\,(In1): $R/\eta_k=RL^k$ is an integer, so that the side of a block torus is a multiple of the mesh &
(C3); Section~I.0.5\\
\hline
\end{tabular}
\medskip

In particular, these propositions hold for all values of $\alpha$, $\kappa$, $M$, $\varepsilon$, $e_0$ and of the constants chosen after them, and for every integer $L$.

The following relations between the constants enter the proofs;
each of them follows from the definitions and is proved at the place indicated.
\begin{itemize}
\item $0<\lambda_{\mathrm{int}}\le12+24\sqrt{2}<48$ (GS7).
\item $K_{\rho}>1$, $t_{0}\le1/4$, $c_{*}\le1/3$, $1/(\sqrt{2}\,\pi)>0.22$ (2.5).
\item $\theta_{0}<\theta_{0}'$ (2.0).
\item $1\le K_{\mathrm{Str}}\le K_D+2K_D^2/\hat{g}$, where $K_D=1+(1+\rho)\,\mathrm{sw}+\hat{g}/2$ (2.3).
\item $\kappa_{(0)}\le\theta_{0}/2$ and $6\kappa_{(0)}\cdot K_{\mathrm{Str}}/\theta_{0}\le c_{\mathrm{el}}/2$ (definition of $\kappa_{(0)}$; 2.7).
\item $c_{\mathrm{lo}}=\rho\cdot c_{*}/(2K_{\rho})\le1/2$ (Section~3, Step~1).
\end{itemize}

The constants $A_{\mathrm{ak}}$, $A'_{\mathrm{ak}}$ whose existence is asserted in Proposition~1.2\,(S6), (S7) are given by the recursions of 2.6 and 2.7;
they are functions of $(c_{\mathrm{el}},K_{\mathrm{Str}},\theta_{0})$, hence of $\rho$.
They are constants of line (C0r) of Section~I.8.2, and in the order of Proposition~14.2\,(a) they are chosen after $\kappa_{(0)}$.

\section*{Chapter 1b. The free Green function; the blocked operator at zero field; the first step: proofs of Propositions 1.5--1.7}

\subsection*{0. Conventions of this chapter: standing data; notation; four facts at an arbitrary mesh; letters and labels}

\subsubsection*{0.1. Standing data}

The standing data are those of the Propositions~1.m:
the Wilson parameter $\rho\in(0,1)$,
the standard weights $w=(w_0,w_1,w_2,w_3)$,
the spin table of Section~I.2.1,
and the numbers $b$, $\alpha$, $L$ of Section~I.8.

\subsubsection*{0.2. Notation}

The notation is that of Sections~I.1--I.3 and of Sections~G1.3 and~G1.4.
In addition, the following conventions are recalled or introduced.

\begin{description}
\item[(n1)]
For $v\in\mathbb{C}^4$, $|v|_2$ denotes the Euclidean norm, and for a $4\times4$ matrix $m$, $|m|_2$ denotes the operator norm associated with it.
For $n\in\mathbb{R}^3$, a difference of photon sites, $|n|=|n|_\infty$ is the maximum norm and $|n|_1$ is the $\ell^1$ norm.
The modulus of a momentum is its Euclidean norm; it is always written $|p|_2$.

\item[(n2)]
At a mesh $\mathrm{me}>0$ the measure on the fermion sites is
\[
\int\mathrm{d}x\,(\cdot)=\mathrm{me}^3\cdot\frac{1}{48}\sum_c\sum_{g\in O_{\mathrm{cub}}}(\cdot).
\]
For a set $Z$ of photon sites, $\int_Z\mathrm{d}x$ is the integral over $Z\times O_{\mathrm{cub}}$, and
\[
|Z|:=\int_Z\mathrm{d}x=\mathrm{me}^3\cdot\#|Z|,
\]
each of the $48$ sites of a cluster having the weight $1/48$ relative to $\mathrm{me}^3$.
Further,
\begin{gather*}
\|u\|_{L^2,Z}^2:=\int_Z\mathrm{d}x\,|u(x)|_2^2,\qquad
|u|_{\mathrm{su},Z}:=\sup\left\{|u(x)|_2:\ x\text{ over }Z\right\},\qquad
\|u\|_{L^1}:=\int\mathrm{d}x\,|u(x)|_2,
\end{gather*}
and for a single site $x$ one sets $\|u\|_{L^2,\{x\}}^2:=\frac{\mathrm{me}^3}{48}\,|u(x)|_2^2$.

\item[(n3)]
On ${V_{\mathrm{v}}}=\{u:O_{\mathrm{cub}}\to\mathbb{C}^4\}$ let
\[
(u',u):=\frac{1}{48}\sum_g\overline{u'(g)}\cdot u(g),\qquad\|u\|:=(u,u)^{\frac{1}{2}}.
\]
For a linear map $O$ between two of the spaces $(\mathbb{C}^4,|\cdot|_2)$, $({V_{\mathrm{v}}},\|\cdot\|)$ and $(\ker J^{*},\|\cdot\|)$, $\|O\|$ denotes its operator norm.
On a single cluster, $(Jv)(g)=v$ and $J^{*}u=\frac{1}{48}\sum_gu(g)$.
By (I.J1) and (I.J4), $J^{*}J=1$, $J$ is an isometry, and $\mathrm{Plt}=JJ^{*}$ and $\mathrm{Phv}=1-\mathrm{Plt}$ are orthogonal projections of ${V_{\mathrm{v}}}$.
Consequently $\|J\|=\|J^{*}\|=1$ and $\|\mathrm{Phv}\|\le1$.

\item[(n4)]
The letters $\mathrm{ak}$, $\mathrm{ak}'$ denote multi-indices in $\{0,1,2,\dots\}^3$, and $|\mathrm{ak}|:=\mathrm{ak}_1+\mathrm{ak}_2+\mathrm{ak}_3$.
The relation $\mathrm{ak}'\le\mathrm{ak}$ is meant componentwise; moreover
\[
\binom{\mathrm{ak}}{\mathrm{ak}'}:=\prod_\mu\binom{\mathrm{ak}_\mu}{\mathrm{ak}'_\mu},\qquad
\partial^{\mathrm{ak}}:=\prod_\mu\left(\frac{\partial}{\partial p_\mu}\right)^{\mathrm{ak}_\mu}.
\]
For two families $\mathrm{fa}=(\mathrm{fa}_{\mathrm{ak}})$ and $\mathrm{fb}=(\mathrm{fb}_{\mathrm{ak}})$ of numbers $\ge0$, indexed by the multi-indices with $|\mathrm{ak}|\le4$, let
\[
(\mathrm{fa}*\mathrm{fb})_{\mathrm{ak}}:=\sum_{\mathrm{ak}'\le\mathrm{ak}}\binom{\mathrm{ak}}{\mathrm{ak}'}\,\mathrm{fa}_{\mathrm{ak}'}\cdot\mathrm{fb}_{\mathrm{ak}-\mathrm{ak}'}.
\]
This product is commutative and associative:
it is the product of the formal power series in three indeterminates with the coefficients $\mathrm{fa}_{\mathrm{ak}}/\mathrm{ak}!$, where $\mathrm{ak}!:=\prod_\mu\mathrm{ak}_\mu!$, truncated at the degree $4$.

\item[(n5)]
Let $\mathrm{rp}_{\max}:=(3\pi^2+1)^{\frac{1}{2}}$.
For $p'\in[-\pi,\pi]^3$ and $\mathbf{m}\in[0,1]$, with $\mathrm{rp}(p')=\left(|p'|_2^2+\mathbf{m}^2\right)^{\frac{1}{2}}$ as in Proposition~1.2, one has
\[
|p'|_2\le\mathrm{rp}(p')\le\mathrm{rp}_{\max}\qquad\text{and}\qquad\mathbf{m}\le\mathrm{rp}(p').
\]

\item[(n6)]
The following facts of Part~I are recalled.
The map $x\mapsto x^{(i)}$ is a bijection of $\mathrm{me}\,\mathbb{Z}^3\times O_{\mathrm{cub}}$ with $\left(x^{(i)}\right)^{(i)}=x$ and $x^{(i)}\ne x$, and it commutes with the translations of $c$, by (I.s1).
One has $\pi\left(x^{(i)}\right)=c$ for $i=0,1,2$, and $|c_3-c|=\mathrm{me}$, by (I.1.2) and (I.s4).
The matrix $M_i(x)$ depends on the label only, and $|M_i(x)|_2=\frac{w_i}{2}(1+\rho)$; in particular $|M_3(x)|_2=3(1+\rho)\le6$, by (I.m6).
\end{description}

\subsubsection*{0.3. Two statements of \cite{QEDI}; classical facts}

\begin{named}{Lemma}[{\cite[Lemma~12.4]{QEDI}}]
Let $\mathrm{kk}<3$ be real, let $N>3-\mathrm{kk}$ be an integer, and let $\mathrm{As}\ge0$.
Let $\mathrm{sy}$ be a $2\pi\mathbb{Z}^3$-periodic function on $\mathbb{R}^3\setminus2\pi\mathbb{Z}^3$ with values in the matrices, of class $C^N$, such that
\[
\left|\partial^{\mathrm{ak}}\mathrm{sy}(p)\right|\le\mathrm{As}\cdot|p|_2^{-\mathrm{kk}-|\mathrm{ak}|}
\]
for all $p\in[-\pi,\pi]^3\setminus\{0\}$ and all $\mathrm{ak}$ with $|\mathrm{ak}|\le N$.
Then $\mathrm{sy}$ is integrable over $[-\pi,\pi]^3$, and
\[
\mathrm{sy}^{v}(n):=\int_{[-\pi,\pi]^3}\frac{\mathrm{d}^3p}{(2\pi)^3}\,\exp(ip\cdot n)\,\mathrm{sy}(p)
\]
satisfies
\[
\left|\mathrm{sy}^{v}(n)\right|\le c(\mathrm{kk},N)\cdot\mathrm{As}\cdot(1+|n|)^{-(3-\mathrm{kk})}\qquad\text{for all }n\in\mathbb{Z}^3,
\]
with a number $c(\mathrm{kk},N)$ that depends on $(\mathrm{kk},N)$ only.
\end{named}

In this lemma, and wherever a pair $(\mathrm{kk},N)$ is named, $N$ is the integer of \cite[Lemma~12.4]{QEDI} and not the index of the spacing.
Matrices carry the operator norm $|\cdot|_2$, which is the norm of matrices in the conventions of \cite[Section~12.1]{QEDI}.
In \cite[Lemma~12.4]{QEDI} the modulus $|p|$ of $p\in\mathbb{R}^3$ is the Euclidean norm $|p|_2$ of Section~I.1.1, as stated at the beginning of \cite[Section~12.4]{QEDI}.

\begin{named}{Lemma}[{\cite[Lemma~12.7(a)]{QEDI}}]
Let $\mathrm{me}>0$ and $\mathrm{vr}\ge\mathrm{me}$.
Then, the sums extending over $z\in\mathrm{me}\,\mathbb{Z}^3$,
\[
\sum_{|z|\le\mathrm{vr}}\mathrm{me}^3\,(\mathrm{me}+|z|)^{-2}\le27\,\mathrm{vr},\qquad
\sum_{|z|\le\mathrm{vr}}\mathrm{me}^3\,(\mathrm{me}+|z|)^{-2-\alpha}\le\frac{27\,\mathrm{vr}^{1-\alpha}}{1-\alpha}.
\]
\end{named}

\begin{proof}
Let $\mathrm{tq}$ be $2$ or $2+\alpha$.
For an integer $j\ge1$ there are $(2j+1)^3-(2j-1)^3=24j^2+2\le26j^2$ points $z\in\mathrm{me}\,\mathbb{Z}^3$ with $|z|=j\cdot\mathrm{me}$, and $(\mathrm{me}+j\cdot\mathrm{me})^{-\mathrm{tq}}\le(j\cdot\mathrm{me})^{-\mathrm{tq}}$.
Let $\mathrm{jm}:=\lfloor\mathrm{vr}/\mathrm{me}\rfloor$; then $\mathrm{jm}\ge1$ and $\mathrm{jm}\cdot\mathrm{me}\le\mathrm{vr}$.
Since the point $z=0$ contributes $\mathrm{me}^{3-\mathrm{tq}}$, the sum $\sum_{|z|\le\mathrm{vr}}\mathrm{me}^3\,(\mathrm{me}+|z|)^{-\mathrm{tq}}$ is at most
\[
\mathrm{me}^{3-\mathrm{tq}}\left(1+26\sum_{j=1}^{\mathrm{jm}}j^{2-\mathrm{tq}}\right).
\]
For $\mathrm{tq}=2$ this equals $\mathrm{me}\,(1+26\,\mathrm{jm})\le27\,\mathrm{jm}\cdot\mathrm{me}\le27\,\mathrm{vr}$, since $1\le\mathrm{jm}$.
For $\mathrm{tq}=2+\alpha$ one has $\sum_{j\le\mathrm{jm}}j^{-\alpha}\le\mathrm{jm}^{1-\alpha}/(1-\alpha)$ and $1\le\mathrm{jm}^{1-\alpha}/(1-\alpha)$;
so the bound is at most $27\,(\mathrm{jm}\cdot\mathrm{me})^{1-\alpha}/(1-\alpha)\le27\,\mathrm{vr}^{1-\alpha}/(1-\alpha)$.
\end{proof}

The following classical facts are recalled.
Leibniz' rule; Cramer's rule; Cauchy's theorem for a rectangle \cite[Chapter~4, Section~1.4]{Ahlfors1979};
the length function of a finite reflection group \cite[Sections~1.6 and~1.7]{Humphreys1990};
the Schwarz inequality \cite[Theorem~4.2]{Rudin1987};
monotone convergence for series with terms $\ge0$, that is, \cite[Theorem~1.26]{Rudin1987} for the counting measure;
and two facts on a bounded operator $O$ of a Hilbert space \cite[Chapter~IV, Theorem~5.2, and Chapter~III, Section~5.5]{Kato1995}:
if $O$ is bounded below, that is, $\|Ou\|\ge t_1\|u\|$ for all $u$ with a number $t_1>0$, then the range of $O$ is closed;
and the orthogonal complement of the range of $O$ is the kernel of $O^{*}$.

\subsubsection*{0.4. Paths at the mesh $\mathrm{me}$}

\begin{named}{(s6)$_{\mathrm{me}}$}
Let $x=(c,g)$ and $\tilde x=(\tilde c,\tilde g)$ be two different sites of $\mathrm{me}\,\mathbb{Z}^3\times O_{\mathrm{cub}}$.
They are joined by a sequence $x=x^0,x^1,\dots,x^{\mathrm{nb}}=\tilde x$ with $x^{j+1}=\left(x^j\right)^{(i_j)}$, $i_j\in\{0,1,2,3\}$, and
\[
\mathrm{nb}\le39\,\frac{d_X(x,\tilde x)}{\mathrm{me}},
\]
such that all photon sites $\pi(x^j)$ lie in the lattice box spanned by $c$ and $\tilde c$, that is, in the set of the $c'\in\mathrm{me}\,\mathbb{Z}^3$ for which $c'_\mu$ lies between $c_\mu$ and $\tilde c_\mu$ for $\mu=1,2,3$.
\end{named}

This is (I.s6) at $k=0$, where $\eta_0=1$, transported by the bijection $(c,g)\mapsto(c/\mathrm{me},g)$ of $\mathrm{me}\,\mathbb{Z}^3\times O_{\mathrm{cub}}$ onto $\mathbb{L}^{\mathrm{f}}_0$:
by Section~G1.4 and (I.1.2) this map takes $x^{(i)}$ to $(x/\mathrm{me})^{(i)}$;
it takes boxes to boxes, and $d_X/\mathrm{me}$ to the distance $d_X$ of $\mathbb{L}^{\mathrm{f}}_0$.
If the two sites lie in one cluster ($c=\tilde c$, $d_X=\mathrm{me}$), then all $x^j$ lie over $c$, and all $i_j$ belong to $\{0,1,2\}$.

\begin{proof}[Proof of (I.s6)]
Lattice units are used.
The vectors $\hat n_2$, $\hat n_1$, $\hat n_0$ are positive multiples of $e_1-e_2$, $e_2-e_3$, $e_3$.
These three vectors form a system of simple roots of the root system whose positive roots are the nine vectors $e_\mu$ and $e_\mu\pm e_\nu$ ($\mu<\nu$);
the reflection group of this root system is $O_{\mathrm{cub}}$, generated by $s_0$, $s_1$, $s_2$ (Section~I.1.2).
The least number of factors $s_0$, $s_1$, $s_2$ in a product that represents a given element of $O_{\mathrm{cub}}$ is the number of positive roots that this element maps to negative roots \cite[Sections~1.6 and~1.7]{Humphreys1990};
hence it is at most $9$.
It follows that two sites $(c,g)$ and $(c,g')$ over one photon site are joined by at most $9$ bonds of the types $0$, $1$, $2$:
one writes $g'=g\,s_{i_1}s_{i_2}\cdots$ with at most $9$ factors and uses $(c,g\,s_i)=(c,g)^{(i)}$ for $i\le2$.
A path of nearest neighbours in $\mathbb{Z}^3$ from $c$ to $\tilde c$ that is monotone in each coordinate has $|c-\tilde c|_1\le3\,|c-\tilde c|$ steps, and it stays in the box.
For a step in the direction $a\in\{\pm e_\mu\}$, one moves inside the cluster to a site with $g\,e_1=a$, using at most $9$ bonds, and then follows the bond of type $3$ of this site, which ends over $c+a$.
At the end one moves inside the cluster of $\tilde c$ to $\tilde x$, using at most $9$ bonds.
In total this gives at most $10\cdot3\,|c-\tilde c|+9\le39\,d_X(x,\tilde x)$ bonds, because $d_X(x,\tilde x)=\max\left(|c-\tilde c|,1\right)\ge1$.
\end{proof}

As a consequence, for all $j$,
\[
|\pi(x^j)-c|\le|c-\tilde c|\le d_X(x,\tilde x)\qquad\text{and}\qquad|\pi(x^j)-\tilde c|\le d_X(x,\tilde x).
\]

\subsubsection*{0.5. Blocks, period cells, energy and the norm of $D$ at an arbitrary mesh}

In this subsection $\mathrm{me}>0$ is an arbitrary mesh, $j\ge0$ an integer and $\theta\in\mathbb{R}^3$.
The blocks of level $j$, the quantities $(u,\nabla^{*}\nabla u)_E$, the operators $D$, $\nabla^{*}\nabla$, $R_i$ and the measure $\int\mathrm{d}x$ are those of Section~G1.4 at the mesh $\mathrm{me}$;
moreover $(\cdot,\cdot):=(\cdot,\cdot)^{\mathrm{ch}}$ and $\|\cdot\|:=\|\cdot\|_{L^2}$.

\begin{named}{(me1)}[Blocks]
The blocks of level $j$ form a partition of $\mathrm{me}\,\mathbb{Z}^3$.
A block of level $j$ has $L^{3j}$ photon sites, and any two of its points $c$, $c'$ satisfy $|c-c'|\le(L^j-1)\,\mathrm{me}$.
A block of level $j+1$ is the union of the $L^3$ blocks of level $j$ whose centres lie in it.
\end{named}

\begin{proof}
The number $L^j$ is odd.
Hence every $c/\mathrm{me}\in\mathbb{Z}^3$ has exactly one nearest point in $L^j\mathbb{Z}^3$, at the maximum distance $\le(L^j-1)/2$.
Therefore the blocks
\[
\left\{c:\ |c-z|<\frac{L^j\cdot\mathrm{me}}{2}\right\}=\left\{c:\ |c-z|\le\frac{(L^j-1)\,\mathrm{me}}{2}\right\},\qquad z\in L^j\cdot\mathrm{me}\,\mathbb{Z}^3,
\]
form a partition of $\mathrm{me}\,\mathbb{Z}^3$, and on each block every coordinate takes exactly $L^j$ values.
Let $z'\in L^{j+1}\cdot\mathrm{me}\,\mathbb{Z}^3$.
The points of $L^j\cdot\mathrm{me}\,\mathbb{Z}^3$ in the block of $z'$ are the $L^3$ points $z=z'+L^j\cdot\mathrm{me}\cdot n$ with $n\in\mathbb{Z}^3$, $|n|\le(L-1)/2$.
If $|c-z|\le(L^j-1)\,\mathrm{me}/2$ for one of them, then
\[
|c-z'|\le\frac{(L^j-1)\,\mathrm{me}}{2}+L^j\cdot\mathrm{me}\cdot\frac{L-1}{2}=\frac{(L^{j+1}-1)\,\mathrm{me}}{2}.
\]
The blocks of these $L^3$ points are disjoint, have $L^{3j}$ sites each, and are contained in the block of $z'$, which has $L^{3(j+1)}$ sites;
so their union is the block of $z'$.
\end{proof}

\begin{named}{(me2)}[Period cells]
Let $R\in L^j\cdot\mathrm{me}\,\mathbb{N}$.
For $z^{\mathrm{o}}\in L^j\cdot\mathrm{me}\,\mathbb{Z}^3$, the union of the blocks of level $j$ with centres in
\[
z^{\mathrm{o}}+\left\{0,\,L^j\cdot\mathrm{me},\,2L^j\cdot\mathrm{me},\,\dots,\,R-L^j\cdot\mathrm{me}\right\}^3
\]
is a period cell of $R\mathbb{Z}^3$ in $\mathrm{me}\,\mathbb{Z}^3$, and every block of level $j$ belongs to such a cell (take its centre as $z^{\mathrm{o}}$).
\end{named}

\begin{proof}
By (me1), every $c\in\mathrm{me}\,\mathbb{Z}^3$ lies in exactly one block of level $j$.
For $n\in\mathbb{Z}^3$ the vector $Rn$ lies in $L^j\cdot\mathrm{me}\,\mathbb{Z}^3$, and the translation by $Rn$ maps the block with the centre $z$ onto the block with the centre $z+Rn$.
Since $R/(L^j\cdot\mathrm{me})$ is an integer $\ge1$, every class of $L^j\cdot\mathrm{me}\,\mathbb{Z}^3$ modulo $R\mathbb{Z}^3$ has exactly one element in $z^{\mathrm{o}}+\{0,L^j\cdot\mathrm{me},\dots,R-L^j\cdot\mathrm{me}\}^3$.
Hence for every $c$ there is exactly one $n\in\mathbb{Z}^3$ such that $c+Rn$ lies in the union.
\end{proof}

For $u,u'\in\mathfrak{S}^{\mathrm{ch}}(\mathrm{me},R;\theta)$ the functions $\overline{u'(x)}\cdot u(x)$ and $|u(x)|_2$ are $R\mathbb{Z}^3$-periodic, because the character $-\exp(i\theta_\mu)$ has modulus one.
Hence, with $(u',u)_E:=\int_E\mathrm{d}x\,\overline{u'(x)}\cdot u(x)$,
\[
(u',u)=\sum_E(u',u)_E,\qquad\|u\|^2=\sum_E\|u\|_{L^2,E}^2,
\]
where $E$ ranges over the blocks of level $j$ of one such cell; the sums do not depend on the cell.
If $R=\infty$, the two identities hold for $u,u'\in\mathrm{HH}(\mathrm{me},\infty;\theta)$, with $E$ ranging over all blocks of level $j$, by (me1).

\begin{named}{(me3)}[Energy]
Let $R\in L^j\cdot\mathrm{me}\,\mathbb{N}\cup\{\infty\}$ and $u\in\mathrm{HH}(\mathrm{me},R;\theta)$.
Then, with the sum $\sum_E$ of (me2),
\[
(u,\nabla^{*}\nabla u)=\frac{1}{2}\sum_{i=0}^{3}w_i\int\mathrm{d}x\,\left|u\left(x^{(i)}\right)-u(x)\right|_2^2\ge\sum_E(u,\nabla^{*}\nabla u)_E.
\]
\end{named}

\begin{proof}
If $R=\infty$, read $\mathrm{me}\,\mathbb{Z}^3$ in place of $\mathrm{me}\,\mathbb{Z}^3/R\mathbb{Z}^3$ and in place of the cell, and disregard the periodicity in $R\mathbb{Z}^3$.
By (n6), $x\mapsto x^{(i)}$ is a bijection of $\mathrm{me}\,\mathbb{Z}^3\times O_{\mathrm{cub}}$ with $\left(x^{(i)}\right)^{(i)}=x$ that commutes with the translations of $c$.
So $R_i$ maps $\mathfrak{S}^{\mathrm{ch}}(\mathrm{me},R;\theta)$ into itself, and $R_i^2=1$.
All sites have the same weight $\mathrm{me}^3/48$, the function $|u|_2^2$ is $R\mathbb{Z}^3$-periodic, and $x\mapsto x^{(i)}$ induces a bijection of $(\mathrm{me}\,\mathbb{Z}^3/R\mathbb{Z}^3)\times O_{\mathrm{cub}}$;
hence $\int\mathrm{d}x\,\left|u\left(x^{(i)}\right)\right|_2^2=\int\mathrm{d}x\,|u(x)|_2^2$, that is, $R_i$ is unitary on $\mathrm{HH}(\mathrm{me},R;\theta)$.
It follows that $R_i=R_i^{*}$, that $(u,R_iu)$ is real, and that
\[
\|(1-R_i)u\|^2=2\|u\|^2-2\,(u,R_iu)=2\,\left(u,(1-R_i)u\right).
\]
Multiplied by $w_i/2$ and summed over $i$, this is the equality, since $\nabla^{*}\nabla=\sum_iw_i(1-R_i)$.
The function $x\mapsto\left|u\left(x^{(i)}\right)-u(x)\right|_2^2$ is $\ge0$ and $R\mathbb{Z}^3$-periodic.
Integrate it over a cell of (me2), and keep, for every block $E$ of the cell, only the sites $x$ over $E$ for which $x^{(i)}$ lies over $E$ as well.
In the integral over these $x$, every unoriented bond $\{x,x^{(i)}\}$ of type $i$ with both ends over $E$ occurs twice, once as $x$ and once as $x^{(i)}$, with the weight $\mathrm{me}^3/48$;
bonds of different types are different pairs of sites, because $x^{(0)}$, $x^{(1)}$, $x^{(2)}$ are three different sites over $c$, and $x^{(3)}$ lies over $c_3\ne c$.
Therefore
\[
\frac{1}{2}\sum_iw_i\int_{x,\,x^{(i)}\text{ over }E}\mathrm{d}x\,\left|u\left(x^{(i)}\right)-u(x)\right|_2^2=(u,\nabla^{*}\nabla u)_E
\]
in the sense of Section~G1.4, and the inequality follows.
\end{proof}

\begin{named}{(me4)}[Norm of $D$]
Let $R\in\mathrm{me}\,\mathbb{N}\cup\{\infty\}$.
Then $\|Du\|\le\mathrm{me}^{-1}(1+\rho)\cdot\mathrm{sw}\cdot\|u\|$ for $u\in\mathrm{HH}(\mathrm{me},R;\theta)$.
\end{named}

\begin{proof}
One has $D=-\mathrm{me}^{-1}\sum_iM_i(\cdot)(R_i-1)$, and $|M_i(x)|_2=\frac{w_i}{2}(1+\rho)$ by (n6).
By the proof of (me3), at $j=0$, the operator $R_i$ is unitary on $\mathrm{HH}(\mathrm{me},R;\theta)$, so that $\|R_i-1\|\le2$.
Finally $\sum_i\frac{w_i}{2}(1+\rho)\cdot2=(1+\rho)\cdot\mathrm{sw}$.
\end{proof}

\subsubsection*{0.6. Letters and labels}

In this chapter $p'$ and $p''$ denote the real and the imaginary part of a complex momentum $p=p'+ip''$, as in Proposition~1.2\,(S7).
The letter $\mathrm{kk}$ denotes both the real number of \cite[Lemma~12.4]{QEDI} (subsection~0.3) and the integer that indexes a class $\mathrm{Sy}(\mathrm{kk};\mathrm{fa})$ (with $\mathrm{kk}_1$, $\mathrm{kk}_2$ in (Sy1)).
In a pair $(\mathrm{kk},N)$ and in \cite[Lemma~12.4]{QEDI}, $N$ is the integer of that lemma, not the index of the spacing $\eta=L^{-N}$;
$c(\mathrm{kk},N)$, with two arguments, is the number of that lemma, and it is not the lattice sum $c(\gamma)$ of Section~G1.3, which has one argument.
In Section~2, $\operatorname{Pen}$ is the operator $\operatorname{Pen}_y(A)$ at $A=0$, formed on $\mathrm{HH}(\eta_k,R;\theta)$ by the display of Section~I.3.5 and written without index and argument (subsection~2.1), and $\operatorname{Pen}_E$ is its part over an $L$-block $E$:
$(\operatorname{Pen}u)(x)=(\operatorname{Pen}_Eu)(x)$ for $x$ over $E$, by (b3).
The letter $\mathrm{rad}$ denotes a radius.
Fermion sites are written $x$, $x'$, $x''$, $\tilde x$, $x^j$, and photon sites $c$, $c'$, $c''$, $\tilde c$, $c^j$, $c^{\mathrm{o}}$, $c^E$;
the letters $z$, $z'$, $z^{\mathrm{o}}$ denote centres of blocks (in Proposition~1.6\,(c), unit sites), except in the statement of \cite[Lemma~12.7(a)]{QEDI} in subsection~0.3, where $z$ is the variable of summation;
$y'$, which occurs in subsection~3.1 only, is a site of $L\mathbb{Z}^3$, and $y$ always denotes the resolvent parameter.
In Section~1 the letters (a)--(f) without further reference are the clauses of Proposition~1.5;
in Section~2, outside Lemma~6.4, the letters (a)--(d) without further reference are the clauses of Proposition~1.6.
The labels (Pn1) and (Pn3) denote the conditions (Pen1) and (Pen3) of \cite[Definition~12.8]{QEDI}, and (Pn2) denotes the first of the two inequalities (Pen2) there.

\subsection*{1. Proposition 1.5}

This section proves Proposition~1.5.

\subsubsection*{1.1. Kernels on ${V_{\mathrm{v}}}$}

For a linear map $O$ of ${V_{\mathrm{v}}}$, the kernel $[O](g,g')$, a $4\times4$ matrix, is defined by
\[
(Ou)(g)=\frac{1}{48}\sum_{g'}[O](g,g')\,u(g').
\]
For $g'\in O_{\mathrm{cub}}$ and $v\in\mathbb{C}^4$ let $\delta_{g',v}\in{V_{\mathrm{v}}}$ be the function $\delta_{g',v}(g):=48\cdot\mathbf{1}[g=g']\,v$.

\begin{description}
\item[(K0)]
$[O](g,g')\,v=(O\delta_{g',v})(g)$, and $\|\delta_{g',v}\|^2=\frac{1}{48}\cdot48^2\,|v|_2^2=48\,|v|_2^2$.

\item[(K1)]
$|[O](g,g')|_2\le48\,\|O\|$.
Indeed, by (K0),
\[
\frac{1}{48}\,\left|[O](g,g')\,v\right|_2^2\le\|O\delta_{g',v}\|^2\le48\,\|O\|^2\,|v|_2^2.
\]

\item[(K2)]
$[J\,m\,J^{*}](g,g')=m$ for every $4\times4$ matrix $m$ and all $g$, $g'$, because $(J\,m\,J^{*}u)(g)=m\,\frac{1}{48}\sum_{g'}u(g')$.

\item[(K3)]
$[1](g,g')=48\cdot\mathbf{1}[g=g']\,1_4$.

\item[(K4)]
For two linear maps $O_1$, $O_2$ of ${V_{\mathrm{v}}}$, the function $g\mapsto[O_1O_2](g,g')\,v$ is the image under $O_1$ of the function $g\mapsto[O_2](g,g')\,v$;
this follows from (K0), applied twice.

\item[(K5)]
The map $O\mapsto[O](g,g')$ is linear.
Hence, if $O(p)$ depends differentiably on $p$, then $\partial^{\mathrm{ak}}[O(p)](g,g')=[\partial^{\mathrm{ak}}O(p)](g,g')$.
\end{description}

\subsubsection*{1.2. The symbol and its blocks}

Let $D^{(p)}$ be the matrix (I.2.6) at $\boldsymbol{\mu}=\mathbf{m}$, and let $D^{\mathrm{ll}}(p)$, $D^{\mathrm{lh}}(p)$, $D^{\mathrm{hl}}(p)$, $D^{\mathrm{hh}}(p)$ and $\mathrm{Dsc}_{\mathbf{m}}(p)$ be formed from it as in Proposition~1.2.
Put
\[
\mathrm{Hi}(p):=D^{\mathrm{hh}}(p)^{-1},\qquad\mathrm{Sc}(p):=\mathrm{Dsc}_{\mathbf{m}}(p)^{-1}.
\]

\begin{description}
\item[(y1)]
In (I.2.6) the variable $p$ occurs in the term of type $3$ only.
For $p,\tilde p\in\mathbb{C}^3$, $u\in{V_{\mathrm{v}}}$ and $|\mathrm{ak}|\ge1$,
\begin{gather*}
\left(\left(D^{(p)}-D^{(\tilde p)}\right)u\right)(g)=-M_3(g)\left(\exp\left(ip\cdot(g\,e_1)\right)-\exp\left(i\tilde p\cdot(g\,e_1)\right)\right)u(g\,s_3),\\
\left(\partial^{\mathrm{ak}}D^{(p)}u\right)(g)=-M_3(g)\,i^{|\mathrm{ak}|}\,(g\,e_1)^{\mathrm{ak}}\exp\left(ip\cdot(g\,e_1)\right)u(g\,s_3),
\end{gather*}
where $(g\,e_1)^{\mathrm{ak}}:=\prod_\mu\left((g\,e_1)_\mu\right)^{\mathrm{ak}_\mu}\in\{0,1,-1\}$, with $0^0:=1$.
The map $g\mapsto g\,s_3$ is a bijection of $O_{\mathrm{cub}}$, $|M_3(g)|_2=3(1+\rho)$, and $g\,e_1$ is one of the six vectors $a\in\{\pm e_\mu\}$.
Summing $\frac{1}{48}|\cdot|_2^2$ over $g$, one obtains
\begin{gather*}
\left\|D^{(p)}-D^{(\tilde p)}\right\|\le3(1+\rho)\max_a\left|\exp(ip\cdot a)-\exp(i\tilde p\cdot a)\right|,\\
\left\|\partial^{\mathrm{ak}}D^{(p)}\right\|\le3(1+\rho)\exp\left(\max_\mu|\operatorname{Im}p_\mu|\right)\qquad(|\mathrm{ak}|\ge1).
\end{gather*}
The maps $J$, $J^{*}$, $\mathrm{Phv}$ do not depend on $p$ and have norm $\le1$.
Therefore both bounds hold also with $D^{\mathrm{lh}}(p)=J^{*}D^{(p)}\,\mathrm{Phv}$, $D^{\mathrm{hl}}(p)=\mathrm{Phv}\,D^{(p)}J$ and $D^{\mathrm{hh}}(p)=\mathrm{Phv}\,D^{(p)}\,\mathrm{Phv}$ in place of $D^{(p)}$.

\item[(y2)]
$\exp(\theta_0)<3$.
Consequently, for $0\le\mathbf{m}\le1$,
\[
\exp(\kappa_{(0)}\cdot\mathbf{m})\le\exp(\theta_0/2)<\sqrt{3}.
\]
Indeed, $\lambda_{\mathrm{int}}$ is an eigenvalue of $\sum_{i\le2}w_i(1-R_i)$ on $\ell^2(O_{\mathrm{cub}})$, where $(R_iu)(g)=u(g\,s_i)$ (I.O8), and the $R_i$ are unitary;
so $\lambda_{\mathrm{int}}\le2(w_0+w_1+w_2)=12+24\sqrt{2}<46$.
By the definition of $\theta_0$ in the list of constants of Proposition~1.2 and by that of $\hat{g}$ in Section~I.8.2, line (C0r),
\[
\exp(\theta_0)-1=\frac{\hat{g}}{6(1+\rho)}=\frac{\rho\cdot\lambda_{\mathrm{int}}}{12(1+\rho)}<\frac{46}{24}<2,
\]
because $\rho/(1+\rho)<1/2$.
Finally, $\kappa_{(0)}\le\theta_0/2$ by the definition of $\kappa_{(0)}$ in the list of constants of Proposition~1.2.

\item[(y3)]
(Holomorphy and periodicity.)
The entries of $D^{(p)}$ are entire and $2\pi\mathbb{Z}^3$-periodic (Section~I.2.6).
The set of the $p\in\mathbb{C}^3$ at which $D^{\mathrm{hh}}(p)$ is invertible on $\ker J^{*}$ is open and $2\pi\mathbb{Z}^3$-periodic;
by Proposition~1.2\,(S1) it contains the closed strip $\{|\operatorname{Im}p_\mu|\le\theta_0\}$.
On this set $\mathrm{Hi}$ is holomorphic, by Cramer's rule, and periodic, and so is $\mathrm{Dsc}_{\mathbf{m}}=D^{\mathrm{ll}}-D^{\mathrm{lh}}\,\mathrm{Hi}\,D^{\mathrm{hl}}$.
On the open subset where $\mathrm{Dsc}_{\mathbf{m}}(p)$ is invertible, $\mathrm{Sc}$ is holomorphic and periodic.
This subset contains $\mathbb{R}^3\setminus2\pi\mathbb{Z}^3$, by Proposition~1.2\,(S5), in which $\mathrm{rp}(p)>0$ for $p\in[-\pi,\pi]^3\setminus\{0\}$, and by periodicity.
If $\mathbf{m}>0$, it also contains all $p$ with $|\operatorname{Im}p_\mu|\le\kappa_{(0)}\cdot\mathbf{m}$, $\mu=1,2,3$, by Proposition~1.2\,(S7) and periodicity;
here $\kappa_{(0)}\cdot\mathbf{m}\le\theta_0/2$.

\item[(y4)]
For a function that is holomorphic near the point $p'+ip''$ ($p',p''\in\mathbb{R}^3$), the derivatives with respect to the real variables $p'$ are the complex derivatives $\partial^{\mathrm{ak}}$ at $p'+ip''$.

\item[(y5)]
(The inverse of the symbol.)
Let $p$ be a point of the strip at which $\mathrm{Dsc}_{\mathbf{m}}(p)$ is invertible.
Identify the block $\operatorname{ran}J$ with $\mathbb{C}^4$ by means of $J$.
Then Proposition~1.2\,(S8), written as an identity between linear maps of ${V_{\mathrm{v}}}$, reads
\begin{gather*}
\left(D^{(p)}\right)^{-1}=J\,\mathrm{Sc}\,J^{*}+\mathrm{Hp}_1(p)+\mathrm{Hp}_2(p)+\mathrm{Hp}_3(p)+\mathrm{Hp}_4(p),\\
\mathrm{Hp}_1:=-J\,\mathrm{Sc}\,D^{\mathrm{lh}}\,\mathrm{Hi}\,\mathrm{Phv},\qquad
\mathrm{Hp}_2:=-\mathrm{Hi}\,D^{\mathrm{hl}}\,\mathrm{Sc}\,J^{*},\\
\mathrm{Hp}_3:=\mathrm{Hi}\,D^{\mathrm{hl}}\,\mathrm{Sc}\,D^{\mathrm{lh}}\,\mathrm{Hi}\,\mathrm{Phv},\qquad
\mathrm{Hp}_4:=\mathrm{Hi}\,\mathrm{Phv},
\end{gather*}
where $\mathrm{Hi}\,\mathrm{Phv}$ is the map $u\mapsto\mathrm{Hi}(\mathrm{Phv}\,u)$ of ${V_{\mathrm{v}}}$ into $\ker J^{*}$, a subspace of ${V_{\mathrm{v}}}$.
\end{description}

\subsubsection*{1.3. Classes of symbols}

\begin{named}{Definition}
Let $\mathbf{m}\in[0,1]$, and let $p''\in\mathbb{R}^3$ satisfy $|p''_\mu|\le\kappa_{(0)}\cdot\mathbf{m}$ for $\mu=1,2,3$ (so that $p''=0$ if $\mathbf{m}=0$).
Let $\mathrm{kk}$ be an integer, and let $\mathrm{fa}=(\mathrm{fa}_{\mathrm{ak}})_{|\mathrm{ak}|\le4}$ be numbers $\ge0$.
A function $O$ of $p$ whose values are complex numbers, or linear maps between two of the spaces $\mathbb{C}^4$, ${V_{\mathrm{v}}}$, $\ker J^{*}$, belongs to $\mathrm{Sy}(\mathrm{kk};\mathrm{fa})$ at $(\mathbf{m},p'')$ if it is holomorphic in a neighbourhood of $\{p'+ip'':\ p'\in\mathbb{R}^3\setminus2\pi\mathbb{Z}^3\}$, $2\pi\mathbb{Z}^3$-periodic, and
\[
\left\|\partial^{\mathrm{ak}}O(p'+ip'')\right\|\le\mathrm{fa}_{\mathrm{ak}}\cdot\mathrm{rp}(p')^{-\mathrm{kk}-|\mathrm{ak}|}
\]
for all $p'\in[-\pi,\pi]^3\setminus\{0\}$ and all $\mathrm{ak}$ with $|\mathrm{ak}|\le4$.
\end{named}

\begin{description}
\item[(Sy1)]
(Products.)
Let $O_1$ belong to $\mathrm{Sy}(\mathrm{kk}_1;\mathrm{fa})$ and $O_2$ to $\mathrm{Sy}(\mathrm{kk}_2;\mathrm{fb})$ at $(\mathbf{m},p'')$, and let the product $O_1O_2$ be defined.
Then $O_1O_2$ belongs to $\mathrm{Sy}(\mathrm{kk}_1+\mathrm{kk}_2;\mathrm{fa}*\mathrm{fb})$ at $(\mathbf{m},p'')$.
Indeed, by Leibniz' rule, in which the order of the factors is kept,
\[
\partial^{\mathrm{ak}}(O_1O_2)=\sum_{\mathrm{ak}'\le\mathrm{ak}}\binom{\mathrm{ak}}{\mathrm{ak}'}\left(\partial^{\mathrm{ak}'}O_1\right)\left(\partial^{\mathrm{ak}-\mathrm{ak}'}O_2\right),
\]
and the norm of the term with the index $\mathrm{ak}'$ is at most $\mathrm{fa}_{\mathrm{ak}'}\cdot\mathrm{rp}^{-\mathrm{kk}_1-|\mathrm{ak}'|}\cdot\mathrm{fb}_{\mathrm{ak}-\mathrm{ak}'}\cdot\mathrm{rp}^{-\mathrm{kk}_2-|\mathrm{ak}|+|\mathrm{ak}'|}$.

\item[(Sy2)]
Multiplication from the left or from the right by one of the constant maps $J$, $J^{*}$, $\mathrm{Phv}$ does not lead out of $\mathrm{Sy}(\mathrm{kk};\mathrm{fa})$, because these maps do not depend on $p$ and have norm $\le1$, by (n3).
\end{description}

\begin{named}{Lemma 5.1}[The factors]
For $|\mathrm{ak}|\le4$ define
\begin{gather*}
\mathrm{fS}_{\mathrm{ak}}:=\max\left(A_{\mathrm{ak}},A'_{\mathrm{ak}}\right),\\
\mathrm{fL}_0:=c_{\mathrm{hl}}\,(1+3\kappa_{(0)}),\qquad
\mathrm{fL}_{\mathrm{ak}}:=3\,c_{\mathrm{hl}}\cdot\mathrm{rp}_{\max}^{|\mathrm{ak}|-1}\quad(|\mathrm{ak}|\ge1),\\
\mathrm{hq}_0:=\frac{2}{\hat{g}},\qquad
\mathrm{hq}_{\mathrm{ak}}:=\frac{6\,c_{\mathrm{hl}}}{\hat{g}}\sum_{\mathrm{ak}'\le\mathrm{ak},\ \mathrm{ak}'\ne0}\binom{\mathrm{ak}}{\mathrm{ak}'}\,\mathrm{hq}_{\mathrm{ak}-\mathrm{ak}'}\quad(|\mathrm{ak}|\ge1),\qquad
\mathrm{fH}_{\mathrm{ak}}:=\mathrm{hq}_{\mathrm{ak}}\cdot\mathrm{rp}_{\max}^{|\mathrm{ak}|},
\end{gather*}
where $A_{\mathrm{ak}}$, $A'_{\mathrm{ak}}$ are the numbers of Proposition~1.2\,(S6), (S7), and $\mathrm{hq}_{\mathrm{ak}}$ is defined by recursion on $|\mathrm{ak}|$;
let $\mathrm{fS}:=(\mathrm{fS}_{\mathrm{ak}})$, $\mathrm{fL}:=(\mathrm{fL}_{\mathrm{ak}})$, $\mathrm{fH}:=(\mathrm{fH}_{\mathrm{ak}})$.
For $a\in\{\pm e_\mu\}$ let $\mathrm{fd}_a(p):=\exp(ip\cdot a)-1$.
Then for every $\mathbf{m}\in[0,1]$ and every $p''$ with $|p''_\mu|\le\kappa_{(0)}\cdot\mathbf{m}$, at $(\mathbf{m},p'')$:
\begin{itemize}
\item[(i)] $\mathrm{Sc}$ belongs to $\mathrm{Sy}(1;\mathrm{fS})$;
\item[(ii)] $D^{\mathrm{lh}}$, $D^{\mathrm{hl}}$ and the six functions $\mathrm{fd}_a$ belong to $\mathrm{Sy}(-1;\mathrm{fL})$;
\item[(iii)] $\mathrm{Hi}$ belongs to $\mathrm{Sy}(0;\mathrm{fH})$.
\end{itemize}
The numbers $\mathrm{fS}_{\mathrm{ak}}$, $\mathrm{fL}_{\mathrm{ak}}$, $\mathrm{fH}_{\mathrm{ak}}$ depend on $\rho$ only.
\end{named}

\begin{proof}
Holomorphy and periodicity hold by (y3); the functions $\mathrm{fd}_a$ are entire and periodic.
Derivatives are read by (y4).
Let $p=p'+ip''$ with $p'\in[-\pi,\pi]^3\setminus\{0\}$.

(i) For $p''=0$ this is Proposition~1.2\,(S6).
For $p''\ne0$ one has $\mathbf{m}>0$, and the assertion is the second bound of Proposition~1.2\,(S7).

(ii) For every $a\in\{\pm e_\mu\}$, by (y2),
\[
\left|\exp(ip\cdot a)-\exp(ip'\cdot a)\right|=\left|\exp(-p''\cdot a)-1\right|\le|p''\cdot a|\cdot\exp\left(|p''\cdot a|\right)\le\sqrt{3}\,\kappa_{(0)}\cdot\mathbf{m}.
\]
Let $O$ be one of the two functions $D^{\mathrm{lh}}$, $D^{\mathrm{hl}}$.
By Proposition~1.2\,(S2), $\|O(p')\|\le c_{\mathrm{hl}}\cdot|p'|_2$.
With (y1), the identity $3(1+\rho)\sqrt{3}=3\,c_{\mathrm{hl}}$ and (n5),
\[
\|O(p)\|\le\|O(p')\|+\|O(p)-O(p')\|\le c_{\mathrm{hl}}\cdot|p'|_2+3\,c_{\mathrm{hl}}\cdot\kappa_{(0)}\cdot\mathbf{m}\le c_{\mathrm{hl}}\,(1+3\kappa_{(0)})\,\mathrm{rp}(p')=\mathrm{fL}_0\cdot\mathrm{rp}(p').
\]
For $|\mathrm{ak}|\ge1$, by (y1), (y2) and (n5),
\[
\left\|\partial^{\mathrm{ak}}O(p)\right\|\le3(1+\rho)\exp(\kappa_{(0)}\cdot\mathbf{m})\le3\,c_{\mathrm{hl}}=3\,c_{\mathrm{hl}}\cdot\mathrm{rp}(p')^{|\mathrm{ak}|-1}\cdot\mathrm{rp}(p')^{1-|\mathrm{ak}|}\le\mathrm{fL}_{\mathrm{ak}}\cdot\mathrm{rp}(p')^{1-|\mathrm{ak}|}.
\]
Since $c_{\mathrm{hl}}=\sqrt{3}(1+\rho)\ge1$,
\[
|\mathrm{fd}_a(p)|\le\left|\exp(ip'\cdot a)-1\right|+\sqrt{3}\,\kappa_{(0)}\cdot\mathbf{m}\le|p'|_2+\sqrt{3}\,\kappa_{(0)}\cdot\mathbf{m}\le\left(1+\sqrt{3}\,\kappa_{(0)}\right)\mathrm{rp}(p')\le\mathrm{fL}_0\cdot\mathrm{rp}(p').
\]
For $|\mathrm{ak}|\ge1$,
\[
\left|\partial^{\mathrm{ak}}\mathrm{fd}_a(p)\right|\le\left|\exp(ip\cdot a)\right|\le\exp(\kappa_{(0)}\cdot\mathbf{m})\le\sqrt{3}\le3\,c_{\mathrm{hl}}=3\,c_{\mathrm{hl}}\cdot\mathrm{rp}(p')^{|\mathrm{ak}|-1}\cdot\mathrm{rp}(p')^{1-|\mathrm{ak}|}\le\mathrm{fL}_{\mathrm{ak}}\cdot\mathrm{rp}(p')^{1-|\mathrm{ak}|}.
\]

(iii) One has $|\operatorname{Im}p_\mu|\le\kappa_{(0)}\cdot\mathbf{m}\le\theta_0/2$.
By Proposition~1.2\,(S1), $\|\mathrm{Hi}(p)\|\le2/\hat{g}=\mathrm{hq}_0$.
Differentiating $D^{\mathrm{hh}}\,\mathrm{Hi}=1$ (on $\ker J^{*}$) one obtains, for $|\mathrm{ak}|\ge1$,
\[
\partial^{\mathrm{ak}}\mathrm{Hi}=-\mathrm{Hi}\sum_{\mathrm{ak}'\le\mathrm{ak},\ \mathrm{ak}'\ne0}\binom{\mathrm{ak}}{\mathrm{ak}'}\left(\partial^{\mathrm{ak}'}D^{\mathrm{hh}}\right)\left(\partial^{\mathrm{ak}-\mathrm{ak}'}\mathrm{Hi}\right).
\]
By (y1) and (y2), $\left\|\partial^{\mathrm{ak}'}D^{\mathrm{hh}}(p)\right\|\le3\,c_{\mathrm{hl}}$.
By induction on $|\mathrm{ak}|$,
\[
\left\|\partial^{\mathrm{ak}}\mathrm{Hi}(p)\right\|\le\frac{2}{\hat{g}}\cdot3\,c_{\mathrm{hl}}\sum_{\mathrm{ak}'\le\mathrm{ak},\ \mathrm{ak}'\ne0}\binom{\mathrm{ak}}{\mathrm{ak}'}\,\mathrm{hq}_{\mathrm{ak}-\mathrm{ak}'}=\mathrm{hq}_{\mathrm{ak}}.
\]
Finally $\mathrm{hq}_{\mathrm{ak}}=\mathrm{hq}_{\mathrm{ak}}\cdot\mathrm{rp}(p')^{|\mathrm{ak}|}\cdot\mathrm{rp}(p')^{-|\mathrm{ak}|}\le\mathrm{fH}_{\mathrm{ak}}\cdot\mathrm{rp}(p')^{-|\mathrm{ak}|}$.

The numbers $\hat{g}$, $c_{\mathrm{hl}}$, $\theta_0$, $\kappa_{(0)}$, $A_{\mathrm{ak}}$, $A'_{\mathrm{ak}}$ are functions of $\rho$ (Section~I.8.2, line (C0r); Proposition~1.2).
\end{proof}

\begin{named}{Lemma 5.2}[Fourier coefficients of a member of a class]
Let $O$ belong to $\mathrm{Sy}(\mathrm{kk};\mathrm{fa})$ at $(\mathbf{m},p'')$, with $\mathrm{kk}\in\{1,0,-1\}$, and let the values of $O$ be linear maps of ${V_{\mathrm{v}}}$, or $4\times4$ matrices.
Put
\[
\mathrm{As}:=\begin{cases}
\max_{|\mathrm{ak}|\le3}\mathrm{fa}_{\mathrm{ak}} & \text{if }\mathrm{kk}=1,\\
\max_{|\mathrm{ak}|\le4}\mathrm{fa}_{\mathrm{ak}} & \text{if }\mathrm{kk}=0,\\
\mathrm{rp}_{\max}\cdot\max_{|\mathrm{ak}|\le4}\mathrm{fa}_{\mathrm{ak}} & \text{if }\mathrm{kk}=-1,
\end{cases}
\qquad
\mathrm{kk}^{+}:=\max(\mathrm{kk},0).
\]
Then $p'\mapsto O(p'+ip'')$ is integrable over $[-\pi,\pi]^3$.
If the values of $O$ are linear maps of ${V_{\mathrm{v}}}$, then for all $n\in\mathbb{Z}^3$ and $g,g'\in O_{\mathrm{cub}}$
\[
\left|\int_{[-\pi,\pi]^3}\frac{\mathrm{d}^3p'}{(2\pi)^3}\,\exp(ip'\cdot n)\,[O(p'+ip'')](g,g')\right|_2\le48\,c(\mathrm{kk}^{+},4-\mathrm{kk}^{+})\cdot\mathrm{As}\cdot(1+|n|)^{-(3-\mathrm{kk}^{+})}.
\]
If the values of $O$ are $4\times4$ matrices, then for all $n\in\mathbb{Z}^3$
\[
\left|\int_{[-\pi,\pi]^3}\frac{\mathrm{d}^3p'}{(2\pi)^3}\,\exp(ip'\cdot n)\,O(p'+ip'')\right|_2\le c(\mathrm{kk}^{+},4-\mathrm{kk}^{+})\cdot\mathrm{As}\cdot(1+|n|)^{-(3-\mathrm{kk}^{+})}.
\]
\end{named}

\begin{proof}
If the values of $O$ are linear maps of ${V_{\mathrm{v}}}$, let $\mathrm{sy}(p'):=[O(p'+ip'')](g,g')$ and $t_1:=48$;
if they are $4\times4$ matrices, let $\mathrm{sy}(p'):=O(p'+ip'')$ and $t_1:=1$.
Then $\mathrm{sy}$ is a $2\pi\mathbb{Z}^3$-periodic function of class $C^\infty$ on $\mathbb{R}^3\setminus2\pi\mathbb{Z}^3$ with values in the $4\times4$ matrices.
By (y4), and for the maps of ${V_{\mathrm{v}}}$ by (K5) and (K1),
\[
\left|\partial^{\mathrm{ak}}\mathrm{sy}(p')\right|_2\le t_1\cdot\mathrm{fa}_{\mathrm{ak}}\cdot\mathrm{rp}(p')^{-\mathrm{kk}-|\mathrm{ak}|}\qquad\text{on }[-\pi,\pi]^3\setminus\{0\}.
\]
If $\mathrm{kk}$ is $1$ or $0$, the exponent $-\mathrm{kk}-|\mathrm{ak}|$ is $\le0$, and $\mathrm{rp}(p')\ge|p'|_2$ gives $\mathrm{rp}(p')^{-\mathrm{kk}-|\mathrm{ak}|}\le|p'|_2^{-\mathrm{kk}-|\mathrm{ak}|}$.
If $\mathrm{kk}=-1$, then $\mathrm{rp}(p')^{1-|\mathrm{ak}|}\le\mathrm{rp}_{\max}\cdot\mathrm{rp}(p')^{-|\mathrm{ak}|}\le\mathrm{rp}_{\max}\cdot|p'|_2^{-|\mathrm{ak}|}$.
Hence $\mathrm{sy}$ satisfies the hypothesis of \cite[Lemma~12.4]{QEDI} (subsection~0.3), with the pair $(\mathrm{kk}^{+},4-\mathrm{kk}^{+})$ in place of the pair $(\mathrm{kk},N)$ there and with the number $t_1\cdot\mathrm{As}$ in place of $\mathrm{As}$;
here $\mathrm{kk}$ is the integer of the class $\mathrm{Sy}(\mathrm{kk};\mathrm{fa})$, so that the pairs are $(1,3)$, $(0,4)$, $(0,4)$ for the classes with $\mathrm{kk}=1,0,-1$.
By \cite[Lemma~12.4]{QEDI}, $\mathrm{sy}$ is integrable, and the two displays hold;
a function with values in the linear maps of ${V_{\mathrm{v}}}$ is integrable if its $48^2$ kernels are.
\end{proof}

\begin{named}{Lemma 5.3}[Shift of the contour]
Let $\mathrm{ti}>0$, and let $O$ be a $2\pi\mathbb{Z}^3$-periodic function with values in a finite dimensional space that is holomorphic in a neighbourhood of $\{p\in\mathbb{C}^3:\ |\operatorname{Im}p_\mu|\le\mathrm{ti},\ \mu=1,2,3\}$.
Let $n\in\mathbb{Z}^3\setminus\{0\}$, let $\mu$ be a direction with $|n_\mu|=|n|$, and let $p'':=\operatorname{sgn}(n_\mu)\,\mathrm{ti}\cdot e_\mu$.
Then
\[
\int_{[-\pi,\pi]^3}\mathrm{d}^3p\,\exp(ip\cdot n)\,O(p)=\exp(-\mathrm{ti}\cdot|n|)\int_{[-\pi,\pi]^3}\mathrm{d}^3p'\,\exp(ip'\cdot n)\,O(p'+ip'').
\]
\end{named}

\begin{proof}
Fix the two real coordinates $p_\nu$, $\nu\ne\mu$.
The function $p_\mu\mapsto\exp(ip\cdot n)\,O(p)$ is holomorphic in a neighbourhood of the closed rectangle with the corners $-\pi$, $\pi$, $\pi+i\operatorname{sgn}(n_\mu)\,\mathrm{ti}$, $-\pi+i\operatorname{sgn}(n_\mu)\,\mathrm{ti}$, and it has the period $2\pi$, because $n_\mu$ is an integer.
By Cauchy's theorem for a rectangle \cite[Chapter~4, Section~1.4]{Ahlfors1979} its integral over the boundary of the rectangle vanishes, and the integrals over the two vertical sides cancel by periodicity.
Hence $\int_{-\pi}^{\pi}\mathrm{d}p_\mu\,\exp(ip\cdot n)\,O(p)$ equals the integral of the same function along $p_\mu=p'_\mu+i\operatorname{sgn}(n_\mu)\,\mathrm{ti}$, $p'_\mu\in[-\pi,\pi]$, on which $\exp(i\,p_\mu\cdot n_\mu)=\exp(i\,p'_\mu\cdot n_\mu)\exp(-\mathrm{ti}\cdot|n_\mu|)$.
The integrands are continuous, and an integration over the two other coordinates gives the assertion.
\end{proof}

\subsubsection*{1.4. The constant $C_G$; proof of (b) and (d)}

\paragraph{The constant.}
Let $\mathrm{fS}$, $\mathrm{fL}$, $\mathrm{fH}$ be the families of Lemma~5.1, and let $c(1,3)$ and $c(0,4)$ be the numbers $c(\mathrm{kk},N)$ of \cite[Lemma~12.4]{QEDI} for the pairs $(\mathrm{kk},N)=(1,3)$ and $(\mathrm{kk},N)=(0,4)$.
Define
\begin{align*}
\mathrm{As}_1&:=\max_{|\mathrm{ak}|\le3}\mathrm{fS}_{\mathrm{ak}},\qquad
\mathrm{As}_2:=\max_{|\mathrm{ak}|\le4}\left(\mathrm{fL}*\mathrm{fS}\right)_{\mathrm{ak}},\qquad
\mathrm{As}_3:=\max_{|\mathrm{ak}|\le4}\left(\mathrm{fS}*\mathrm{fL}*\mathrm{fH}\right)_{\mathrm{ak}},\\
\mathrm{As}_4&:=\mathrm{rp}_{\mathrm{max}}\cdot\max_{|\mathrm{ak}|\le4}\left(\mathrm{fH}*\mathrm{fL}*\mathrm{fS}*\mathrm{fL}*\mathrm{fH}\right)_{\mathrm{ak}},\qquad
\mathrm{As}_5:=\frac{2}{\hat{g}}\cdot\exp\left(\frac{\theta_0}{2}\right)\cdot\left(\frac{6}{\mathrm{e}\cdot\theta_0}\right)^{3},
\end{align*}
and
\[
C_G:=\max\left\{c(1,3)\cdot\mathrm{As}_1,\ \ c(0,4)\cdot\mathrm{As}_2,\ \ 48\left[c(0,4)\cdot\left(2\,\mathrm{As}_3+\mathrm{As}_4\right)+\mathrm{As}_5\right]\right\}.
\]
By Lemma~5.1, and because $c(1,3)$ and $c(0,4)$ are absolute numbers, $C_G$ is a function of $\rho$ alone.

\begin{proof}[Proof of Proposition~1.5\,(b) and (d)]
\emph{The integral defining $G_{\mathbf{m}}$.}
Let $\mathbf{m}\in[0,1]$, and let $p''$ be as in the Definition of subsection~1.3.
Combining (y5) with Lemma~5.1, (Sy1) and (Sy2), one obtains at $(\mathbf{m},p'')$, for each $a\in\{\pm e_1,\pm e_2,\pm e_3\}$,
\begin{equation}\tag{V.1}
\begin{aligned}
&J\,\mathrm{Sc}\,J^{*}\in\mathrm{Sy}(1;\mathrm{fS}),\qquad
\mathrm{Hp}_1,\ \mathrm{Hp}_2\in\mathrm{Sy}(0;\mathrm{fS}*\mathrm{fL}*\mathrm{fH}),\qquad
\mathrm{Hp}_3\in\mathrm{Sy}(-1;\mathrm{fH}*\mathrm{fL}*\mathrm{fS}*\mathrm{fL}*\mathrm{fH}),\\
&\mathrm{Hp}_4\in\mathrm{Sy}(0;\mathrm{fH}),\qquad
\mathrm{fd}_a\,\mathrm{Sc}\in\mathrm{Sy}(0;\mathrm{fL}*\mathrm{fS}).
\end{aligned}
\end{equation}
Lemma~5.2, applied with $p''=0$, shows that each of these functions is integrable over $[-\pi,\pi]^3$.
Consequently $G_{\mathbf{m}}(x,x')$ and $\mathrm{Sc}^{v}_{\mathbf{m}}(n)$ are well defined for every $\mathbf{m}\in[0,1]$.
Put $n:=c-c'$.
By (y5) and (K2),
\begin{equation}\tag{V.2}
G_{\mathbf{m}}(x,x')=\mathrm{Sc}^{v}_{\mathbf{m}}(n)+G^{h}_{\mathbf{m}}(x,x'),\qquad
G^{h}_{\mathbf{m}}(x,x'):=\sum_{j=1}^{4}\int_{[-\pi,\pi]^3}\frac{\mathrm{d}^3p}{(2\pi)^3}\,\mathrm{e}^{ip\cdot n}\left[\mathrm{Hp}_j(p)\right](g,g'),
\end{equation}
and this is the decomposition asserted in (b).

\emph{The three bounds of (b), and those of (d).}
Let $\mathbf{m}\in[0,1]$ and $n\in\mathbb{Z}^3$.
The two cases are treated simultaneously, with the following choices of a vector $p''\in\mathbb{R}^3$ and of a number $\mathrm{fac}$:
\begin{description}
\item[Case (b)] $p'':=0$ and $\mathrm{fac}:=1$;
\item[Case (d)] $\mathbf{m}>0$ and $n\ne0$; $\mu$ is a direction with $|n_\mu|=|n|$, $p'':=\operatorname{sgn}(n_\mu)\,\kappa_{(0)}\cdot\mathbf{m}\,e_\mu$ and $\mathrm{fac}:=\exp\left(-\kappa_{(0)}\cdot\mathbf{m}\,|n|\right)$.
\end{description}
In case (d) the five functions $\mathrm{Sc}$, $\mathrm{fd}_a\,\mathrm{Sc}$, $\mathrm{Hp}_1$, $\mathrm{Hp}_2$, $\mathrm{Hp}_3$ are $2\pi\mathbb{Z}^3$-periodic and holomorphic in a neighbourhood of the set $\left\{p\in\mathbb{C}^3:|\operatorname{Im}p_\nu|\le\kappa_{(0)}\cdot\mathbf{m},\ \nu=1,2,3\right\}$, by (y3).
Lemma~5.3 with $\mathrm{ti}=\kappa_{(0)}\cdot\mathbf{m}$ therefore yields, in both cases and for $O$ any one of these five functions,
\begin{equation}\tag{V.3}
\int_{[-\pi,\pi]^3}\frac{\mathrm{d}^3p}{(2\pi)^3}\,\mathrm{e}^{ip\cdot n}\,O(p)=\mathrm{fac}\cdot\int_{[-\pi,\pi]^3}\frac{\mathrm{d}^3p'}{(2\pi)^3}\,\mathrm{e}^{ip'\cdot n}\,O(p'+ip'');
\end{equation}
in case (b) the two sides coincide trivially.
\begin{itemize}
\item[(1)] By (V.3), Lemma~5.1\,(i) and Lemma~5.2, the latter for $4\times4$ matrices and with $\mathrm{kk}=1$,
\[
\left|\mathrm{Sc}^{v}_{\mathbf{m}}(n)\right|_2\le c(1,3)\cdot\mathrm{As}_1\cdot\mathrm{fac}\cdot(1+|n|)^{-2}\le C_G\cdot\mathrm{fac}\cdot(1+|n|)^{-2}.
\]
\item[(2)] Let $a\in\{\pm e_1,\pm e_2,\pm e_3\}$.
Then
\[
\mathrm{Sc}^{v}_{\mathbf{m}}(n+a)-\mathrm{Sc}^{v}_{\mathbf{m}}(n)=\int_{[-\pi,\pi]^3}\frac{\mathrm{d}^3p}{(2\pi)^3}\,\mathrm{e}^{ip\cdot n}\,\mathrm{fd}_a(p)\,\mathrm{Sc}(p),
\]
and (V.3), (V.1) and Lemma~5.2, the latter for $4\times4$ matrices and with $\mathrm{kk}=0$, give
\[
\left|\mathrm{Sc}^{v}_{\mathbf{m}}(n+a)-\mathrm{Sc}^{v}_{\mathbf{m}}(n)\right|_2\le c(0,4)\cdot\mathrm{As}_2\cdot\mathrm{fac}\cdot(1+|n|)^{-3}\le C_G\cdot\mathrm{fac}\cdot(1+|n|)^{-3}.
\]
\item[(3)] By (V.3), (V.1) and Lemma~5.2, with $\mathrm{kk}=0$ for $j=1,2$ and with $\mathrm{kk}=-1$ for $j=3$, the sum of the terms $j=1,2,3$ of (V.2) has $|\cdot|_2$-norm at most
\[
48\,c(0,4)\cdot\left(2\,\mathrm{As}_3+\mathrm{As}_4\right)\cdot\mathrm{fac}\cdot(1+|n|)^{-3}.
\]
\item[(4)] There remains the term $j=4$.
The function $\mathrm{Hp}_4=\mathrm{Hi}\,\mathrm{Phv}$ is $2\pi\mathbb{Z}^3$-periodic and holomorphic in a neighbourhood of the strip $\left\{p\in\mathbb{C}^3:|\operatorname{Im}p_\nu|\le\theta_0,\ \nu=1,2,3\right\}$, and $\|\mathrm{Hp}_4(p)\|\le2/\hat{g}$ on this strip; this follows from (y3) and Proposition~1.2\,(S1).
For $n\ne0$, Lemma~5.3 with $\mathrm{ti}=\theta_0$ and (K1) give
\[
\left|\int_{[-\pi,\pi]^3}\frac{\mathrm{d}^3p}{(2\pi)^3}\,\mathrm{e}^{ip\cdot n}\left[\mathrm{Hp}_4(p)\right](g,g')\right|_2\le48\cdot\frac{2}{\hat{g}}\cdot\exp\left(-\theta_0|n|\right),
\]
and for $n=0$ the same bound follows from (K1) without any shift of the contour.
Write
\[
\exp\left(-\theta_0|n|\right)=\exp\left(-\frac{\theta_0}{2}|n|\right)\cdot\exp\left(-\frac{\theta_0}{2}|n|\right).
\]
The first factor does not exceed $\mathrm{fac}$, because $\kappa_{(0)}\cdot\mathbf{m}\le\theta_0/2$.
For the second factor observe that, for every real $\mathrm{ts}\ge0$,
\[
(1+\mathrm{ts})^3\exp\left(-\frac{\theta_0}{2}\,\mathrm{ts}\right)=\exp\left(\frac{\theta_0}{2}\right)(1+\mathrm{ts})^3\exp\left(-\frac{\theta_0}{2}\,(1+\mathrm{ts})\right)\le\exp\left(\frac{\theta_0}{2}\right)\left(\frac{6}{\mathrm{e}\cdot\theta_0}\right)^{3},
\]
since $\mathrm{tq}^3\exp\left(-\frac{\theta_0}{2}\,\mathrm{tq}\right)\le\left(\frac{6}{\mathrm{e}\cdot\theta_0}\right)^{3}$ for all real $\mathrm{tq}\ge0$, the maximum being attained at $\mathrm{tq}=6/\theta_0$.
Hence the term $j=4$ is at most $48\,\mathrm{As}_5\cdot\mathrm{fac}\cdot(1+|n|)^{-3}$.
\end{itemize}
Together, (3) and (4) give
\[
\left|G^{h}_{\mathbf{m}}(x,x')\right|_2\le C_G\cdot\mathrm{fac}\cdot(1+|c-c'|)^{-3}.
\]
In case (b) the estimates obtained are the three bounds of (b).
In case (d) they are the same three bounds multiplied by $\exp(-\kappa_{(0)}\cdot\mathbf{m}\,|n|)$; for $n=0$ this factor equals $1$, and the bounds are those of (b).

\emph{The two consequences, in (b) and in (d).}
With $n=c-c'$, let $\mathrm{fac}(n):=1$ in (b) and $\mathrm{fac}(n):=\exp(-\kappa_{(0)}\cdot\mathbf{m}\,|n|)$ in (d).
Then
\[
\left|G_{\mathbf{m}}(x,x')\right|_2\le C_G\cdot\mathrm{fac}(n)\cdot\left[(1+|n|)^{-2}+(1+|n|)^{-3}\right]\le2\,C_G\cdot\mathrm{fac}(n)\cdot(1+|n|)^{-2}.
\]
Let $i\in\{0,1,2\}$.
The site $x^{(i)}=(c,gs_i)$ lies over $c$; hence the term $\mathrm{Sc}^{v}_{\mathbf{m}}(n)$ of (V.2) has the same value at $x^{(i)}$ and at $x$, and it cancels in the difference:
\[
\left|G_{\mathbf{m}}(x^{(i)},x')-G_{\mathbf{m}}(x,x')\right|_2=\left|G^{h}_{\mathbf{m}}(x^{(i)},x')-G^{h}_{\mathbf{m}}(x,x')\right|_2\le2\,C_G\cdot\mathrm{fac}(n)\cdot(1+|n|)^{-3}.
\]
Let $i=3$.
Then $x^{(3)}=(c+a,gs_3)$ with $a=ge_1\in\{\pm e_1,\pm e_2,\pm e_3\}$, and
\[
G_{\mathbf{m}}(x^{(3)},x')-G_{\mathbf{m}}(x,x')=\left[\mathrm{Sc}^{v}_{\mathbf{m}}(n+a)-\mathrm{Sc}^{v}_{\mathbf{m}}(n)\right]+G^{h}_{\mathbf{m}}(x^{(3)},x')-G^{h}_{\mathbf{m}}(x,x').
\]
One has $1+|n+a|\ge\frac{1}{2}(1+|n|)$: for $n=0$ the left side equals $2$, and for $|n|\ge1$ it is at least $|n|\ge\frac{1}{2}(1+|n|)$.
Therefore $(1+|n+a|)^{-3}\le8\,(1+|n|)^{-3}$.
Moreover $\mathrm{fac}(n+a)\le\exp(\kappa_{(0)}\cdot\mathbf{m})\cdot\mathrm{fac}(n)\le\sqrt3\cdot\mathrm{fac}(n)$ by (y2).
It follows that
\[
\left|G_{\mathbf{m}}(x^{(3)},x')-G_{\mathbf{m}}(x,x')\right|_2\le\left(1+8\sqrt3+1\right)C_G\cdot\mathrm{fac}(n)\cdot(1+|n|)^{-3}\le20\,C_G\cdot\mathrm{fac}(n)\cdot(1+|n|)^{-3}.
\]
This proves (b) and (d).
One constant serves for all $\mathbf{m}\in[0,1]$, since the families $\mathrm{fS}$, $\mathrm{fL}$, $\mathrm{fH}$ do not depend on $\mathbf{m}$.
\end{proof}

\subsubsection*{1.5. Proof of (a) and (c)}

\begin{proof}[Proof of Proposition~1.5\,(a)]
Fix $x'=(c',g')$ and $v\in\mathbb{C}^4$.
For $p\in[-\pi,\pi]^3\setminus\{0\}$ let $u_p\in{V_{\mathrm{v}}}$ be the function
\[
u_p(g):=\left[\left(D^{(p)}\right)^{-1}\right](g,g')\,v.
\]
By (I.2.6) the operator $D_0+\mathbf{m}$ sends the function $(c,g)\mapsto\mathrm{e}^{ip\cdot c}u_p(g)$ to the function $(c,g)\mapsto\mathrm{e}^{ip\cdot c}\left(D^{(p)}u_p\right)(g)$, and by (K4) and (K3)
\[
\left(D^{(p)}u_p\right)(g)=\left[D^{(p)}\left(D^{(p)}\right)^{-1}\right](g,g')\,v=48\cdot\mathbf{1}[g=g']\,v.
\]
The value of $(D_0+\mathbf{m})f$ at $x$ is a linear combination of the five values $f(x),f(x^{(0)}),\dots,f(x^{(3)})$, with $4\times4$ matrices as coefficients that do not depend on $p$.
The integrand defining $G_{\mathbf{m}}$ is integrable (subsection~1.4, the paragraph on the integral defining $G_{\mathbf{m}}$).
Hence $D_0+\mathbf{m}$ may be applied under the integral sign, and
\[
\left((D_0+\mathbf{m})\,G_{\mathbf{m}}(\cdot,x')\,v\right)(x)=\int_{[-\pi,\pi]^3}\frac{\mathrm{d}^3p}{(2\pi)^3}\,\mathrm{e}^{ip\cdot(c-c')}\,48\cdot\mathbf{1}[g=g']\,v=48\cdot\mathbf{1}[c=c']\,\mathbf{1}[g=g']\,v.
\]
\end{proof}

\begin{proof}[Proof of Proposition~1.5\,(c)]
Let $v$ be finitely supported, and let $\mathrm{fs}:=(D_0+\mathbf{m})v$, which is finitely supported as well.
The functions
\[
\mathrm{vh}(p)(g):=\sum_c\mathrm{e}^{-ip\cdot c}\,v(c,g),\qquad\mathrm{fh}(p)(g):=\sum_c\mathrm{e}^{-ip\cdot c}\,\mathrm{fs}(c,g)
\]
are trigonometric polynomials with values in ${V_{\mathrm{v}}}$.
One has $v(c,g)=\int_{[-\pi,\pi]^3}\frac{\mathrm{d}^3p}{(2\pi)^3}\,\mathrm{e}^{ip\cdot c}\,\mathrm{vh}(p)(g)$; applying $D_0+\mathbf{m}$ under the integral and using (I.2.6), one obtains
\[
\mathrm{fs}(c,g)=\int_{[-\pi,\pi]^3}\frac{\mathrm{d}^3p}{(2\pi)^3}\,\mathrm{e}^{ip\cdot c}\left(D^{(p)}\mathrm{vh}(p)\right)(g).
\]
Since two trigonometric polynomials with the same Fourier coefficients are equal, $\mathrm{fh}(p)=D^{(p)}\mathrm{vh}(p)$.
As the sum over $x'$ is finite,
\begin{align*}
\frac{1}{48}\sum_{x'}G_{\mathbf{m}}(x,x')\,\mathrm{fs}(x')
&=\int_{[-\pi,\pi]^3}\frac{\mathrm{d}^3p}{(2\pi)^3}\,\mathrm{e}^{ip\cdot c}\,\frac{1}{48}\sum_{g'}\left[\left(D^{(p)}\right)^{-1}\right](g,g')\sum_{c'}\mathrm{e}^{-ip\cdot c'}\,\mathrm{fs}(c',g')\\
&=\int_{[-\pi,\pi]^3}\frac{\mathrm{d}^3p}{(2\pi)^3}\,\mathrm{e}^{ip\cdot c}\left(\left(D^{(p)}\right)^{-1}\mathrm{fh}(p)\right)(g)\\
&=\int_{[-\pi,\pi]^3}\frac{\mathrm{d}^3p}{(2\pi)^3}\,\mathrm{e}^{ip\cdot c}\,\mathrm{vh}(p)(g)=v(c,g),
\end{align*}
where the third equality rests on $\left(D^{(p)}\right)^{-1}D^{(p)}=1$ for $p\in[-\pi,\pi]^3\setminus\{0\}$.
\end{proof}

\subsubsection*{1.6. Proof of (e) and (f)}

\begin{proof}[Proof of Proposition~1.5\,(e)]
Let $\mathrm{me}>0$ and $0\le\mu\le1/\mathrm{me}$, and put $\mathbf{m}:=\mu\cdot\mathrm{me}\in[0,1]$.
To a function $\mathrm{Fl}$ on $\mathbb{Z}^3\times O_{\mathrm{cub}}$ associate the function $f(x):=\mathrm{Fl}(x/\mathrm{me})$ on $\mathrm{me}\,\mathbb{Z}^3\times O_{\mathrm{cub}}$.
Since $(x/\mathrm{me})^{(i)}=x^{(i)}/\mathrm{me}$ and $M_i(x/\mathrm{me})=M_i(x)$ (Section~G1.4, (I.1.2), (I.2.1)),
\begin{equation}\tag{V.4}
\begin{aligned}
(Df)(x)&=-\mathrm{me}^{-1}\sum_iM_i(x)\left[\mathrm{Fl}\left(x^{(i)}/\mathrm{me}\right)-\mathrm{Fl}(x/\mathrm{me})\right]=\mathrm{me}^{-1}\left(D_0\mathrm{Fl}\right)(x/\mathrm{me}),\\
\left((D+\mu)f\right)(x)&=\mathrm{me}^{-1}\left((D_0+\mathbf{m})\mathrm{Fl}\right)(x/\mathrm{me}).
\end{aligned}
\end{equation}
Taking $\mathrm{Fl}=G_{\mathbf{m}}(\cdot,x'/\mathrm{me})\,v$ and using (a), one finds
\[
\left((D+\mu)\,G(\cdot,x')\,v\right)(x)=\mathrm{me}^{-2}\,\mathrm{me}^{-1}\,48\cdot\mathbf{1}[x/\mathrm{me}=x'/\mathrm{me}]\,v=48\,\mathrm{me}^{-3}\,\mathbf{1}[x=x']\,v,
\]
which is the kernel of the identity with respect to $\int\mathrm{d}x=\mathrm{me}^3\,\frac{1}{48}\sum$, by (I.ke1).
By (b),
\[
|G(x,x')|_2\le\mathrm{me}^{-2}\cdot2\,C_G\left(1+\frac{|c-c'|}{\mathrm{me}}\right)^{-2}=2\,C_G\left(\mathrm{me}+|c-c'|\right)^{-2},
\]
and, for $i=0,\dots,3$,
\[
\left|G(x^{(i)},x')-G(x,x')\right|_2\le\mathrm{me}^{-2}\cdot20\,C_G\left(1+\frac{|c-c'|}{\mathrm{me}}\right)^{-3}=20\,C_G\cdot\mathrm{me}\cdot\left(\mathrm{me}+|c-c'|\right)^{-3}.
\]
Finally, let $v$ be finitely supported on $\mathrm{me}\,\mathbb{Z}^3\times O_{\mathrm{cub}}$, and let $\mathrm{Fl}$ be the function on $\mathbb{Z}^3\times O_{\mathrm{cub}}$ with $\mathrm{Fl}(x/\mathrm{me})=v(x)$.
By (c) and (V.4), the sums running over $x'\in\mathrm{me}\,\mathbb{Z}^3\times O_{\mathrm{cub}}$,
\begin{align*}
v(x)=\mathrm{Fl}(x/\mathrm{me})
&=\frac{1}{48}\sum_{x'}G_{\mathbf{m}}(x/\mathrm{me},x'/\mathrm{me})\,\mathrm{me}\left((D+\mu)v\right)(x')\\
&=\mathrm{me}^3\,\frac{1}{48}\sum_{x'}\mathrm{me}^{-2}\,G_{\mathbf{m}}(x/\mathrm{me},x'/\mathrm{me})\left((D+\mu)v\right)(x')
=\int\mathrm{d}x'\,G(x,x')\left((D+\mu)v\right)(x').
\end{align*}
\end{proof}

\begin{proof}[Proof of Proposition~1.5\,(f)]
Let $\alpha\in[0,1]$ be arbitrary.
Only two properties of $\alpha$ enter the argument: $\alpha\ge0$, through the inequality $\mathrm{ts}^{\alpha}\le(3\,\mathrm{hd})^{\alpha}$ for $0<\mathrm{ts}\le3\,\mathrm{hd}$, and $\alpha\le1$, through $3^{\alpha}\le3$ and through $(\mathrm{hd}/\mathrm{rn})^{1-\alpha}\le1$ for $\mathrm{hd}<\mathrm{rn}$.
Let $x\ne\tilde{x}$, and put
\[
\mathrm{hd}:=d_X(x,\tilde{x})=\max\left(|c-\tilde{c}|,\mathrm{me}\right),\qquad\mathrm{rn}:=\mathrm{me}+\min\left(|c-c'|,|\tilde{c}-c'|\right).
\]

\emph{Case $\mathrm{hd}\ge\mathrm{rn}/2$.}
As $|c-\tilde{c}|\le\mathrm{hd}$, both $\mathrm{me}+|c-c'|$ and $\mathrm{me}+|\tilde{c}-c'|$ are at most $\mathrm{rn}+\mathrm{hd}\le3\,\mathrm{hd}$.
By (e), therefore,
\[
|G(x,x')|_2\le2\,C_G\left(\mathrm{me}+|c-c'|\right)^{-2-\alpha}\left(\mathrm{me}+|c-c'|\right)^{\alpha}\le2\cdot3^{\alpha}\,C_G\cdot\mathrm{hd}^{\alpha}\left(\mathrm{me}+|c-c'|\right)^{-2-\alpha},
\]
and in the same way
\[
|G(\tilde{x},x')|_2\le2\,C_G\left(\mathrm{me}+|\tilde{c}-c'|\right)^{-2-\alpha}\left(\mathrm{me}+|\tilde{c}-c'|\right)^{\alpha}\le2\cdot3^{\alpha}\,C_G\cdot\mathrm{hd}^{\alpha}\left(\mathrm{me}+|\tilde{c}-c'|\right)^{-2-\alpha}.
\]
Since $2\cdot3^{\alpha}\le6\le6240$, the bound of (f) follows from $|G(x,x')-G(\tilde{x},x')|_2\le|G(x,x')|_2+|G(\tilde{x},x')|_2$.

\emph{Case $\mathrm{hd}<\mathrm{rn}/2$.}
Let $x=x^0,x^1,\dots,x^{\mathrm{nb}}=\tilde{x}$ be the sequence of (s6)$_{\mathrm{me}}$, so that $\mathrm{nb}\le39\,\mathrm{hd}/\mathrm{me}$, and let $c^j:=\pi(x^j)$.
By subsection~0.4, $|c^j-c|\le\mathrm{hd}$ and $|c^j-\tilde{c}|\le\mathrm{hd}$; hence $\mathrm{me}+|c^j-c'|\ge\mathrm{rn}-\mathrm{hd}>\mathrm{rn}/2$.
Applied to the site $x^j$ and to its neighbour $x^{j+1}=(x^j)^{(i_j)}$, (e) gives
\[
\left|G(x^{j+1},x')-G(x^j,x')\right|_2\le20\,C_G\cdot\mathrm{me}\cdot\left(\frac{\mathrm{rn}}{2}\right)^{-3}.
\]
Summing over $j$ one obtains
\[
\left|G(x,x')-G(\tilde{x},x')\right|_2\le39\,\frac{\mathrm{hd}}{\mathrm{me}}\cdot20\,C_G\cdot\mathrm{me}\cdot8\,\mathrm{rn}^{-3}=6240\,C_G\cdot\mathrm{hd}^{\alpha}\cdot\mathrm{rn}^{-2-\alpha}\left(\frac{\mathrm{hd}}{\mathrm{rn}}\right)^{1-\alpha}\le6240\,C_G\cdot\mathrm{hd}^{\alpha}\cdot\mathrm{rn}^{-2-\alpha},
\]
and $\mathrm{rn}^{-2-\alpha}$ is the larger of the two terms in the bracket of (f).

For a pair inside a cluster, that is $c=\tilde{c}$ and $\mathrm{hd}=\mathrm{me}$, the two terms of the bracket are equal; the first case then means $\mathrm{me}+|c-c'|\le2\,\mathrm{me}$, and in the second case the sequence consists of bonds of the types $0$, $1$, $2$ over $c$.
All points occurring in (e) and (f) lie in $\mathrm{me}\,\mathbb{Z}^3\times O_{\mathrm{cub}}$, and $|c-c'|$ and $d_X$ are the distances of this lattice.
\end{proof}

\begin{named}{Remark}
One constant $C_G$, a function of $\rho$ alone, serves throughout Proposition~1.5: for every $\mathbf{m}\in[0,1]$ (subsection~1.4), and in (e) and (f) for every $\mathrm{me}>0$, every $\mu\in[0,1/\mathrm{me}]$ and every $\alpha\in[0,1]$ (proof of (f)).
The proposition concerns the free operator, without a connection; since neither the charge nor a side nor an index of the spacing enters it, its bounds hold with the same constants for all values of these three quantities.
The clauses (a)--(d) are statements in lattice units, that is, on $\mathbb{L}^{\mathrm{f}}_0$ (level $0$, $\eta_0=1$, with $48$ sites of weight $1/48$ over each photon site).
For $\mathrm{me}=\eta_k$ the clauses (e) and (f) are the corresponding statements on $\mathbb{L}^{\mathrm{f}}_k$; for $k=0$ they reduce to (a)--(c) with $\mu=\mathbf{m}$, and in (f) one then has $\mathrm{hd}\ge1$.
\end{named}

\subsection*{2. Proposition 1.6}

This section is devoted to the proof of Proposition~1.6.

\subsubsection*{2.1. Setting: blocks, period cells, projections, energy}

Throughout this section $k\ge1$ is an integer, $\mathrm{me}=\eta_k=L^{-k}$, $R\in L\mathbb{N}\cup\{\infty\}$ with $\mathbb{N}=\{1,2,\dots\}$, $\theta\in\mathbb{R}^3$, $\mu\in[0,1]$ and $y\in\mathbb{R}\cup\{\infty\}$.
Let $\mathrm{HH}:=\mathrm{HH}(\eta_k,R;\theta)$; on $\mathrm{HH}$ the pairing $(\cdot,\cdot)$ stands for $(\cdot,\cdot)^{\mathrm{ch}}$ and the norm $\|\cdot\|$ for $\|\cdot\|_{L^2}$.
The operators $D$, $\nabla^{*}\nabla$ and $R_i$ are those of Section~G1.4 at the mesh $\eta_k$.
Let
\[
\operatorname{Pen}:=\boldsymbol{\alpha}_y\,JP_k(0)J^{*}+\boldsymbol{\beta}_y\,JP_{k+1}(0)J^{*},\qquad\text{so that}\qquad F_y(0)=D+\mu+\operatorname{Pen},
\]
and recall the convention $-\infty:=\infty$ of Section~G1.9.
\begin{itemize}
\item[(b1)] \emph{Blocks.}
A \emph{unit block} is a block $\Delta_z$ of level $k$, with $z\in\mathbb{Z}^3$ and side $L^k\mathrm{me}=1$; an \emph{$L$-block} is a block $E$ of level $k+1$, with centre in $L\mathbb{Z}^3$ and side $L^{k+1}\mathrm{me}=L$.
By (me1) with $j=k$ and with $j=k+1$, the unit blocks form a partition of $\mathrm{me}\,\mathbb{Z}^3$, and so do the $L$-blocks.
The block $\Delta_z$ contains $L^{3k}$ photon sites, the set $\Delta_z\times O_{\mathrm{cub}}$ contains $48\,L^{3k}$ sites, and $|\Delta_z|=1$.
Further $|E|=L^3$, and $E$ is the union of the $L^3$ unit blocks $\Delta_z$ with $z\in E$.
Any two points $c$, $c'$ of one $L$-block satisfy $|c-c'|\le L-\mathrm{me}<L$ and $|c-c'|_1<3L$.
\item[(b2)] \emph{Period cells.}
Either $R$ is a multiple of $L=L^{k+1}\mathrm{me}$, or $R=\infty$.
In this section $\sum_E$ always denotes the sum over the $L$-blocks of one period cell of (me2) with $j=k+1$, or over all $L$-blocks if $R=\infty$.
With $(u',u)_E:=\int_E\mathrm{d}x\,\overline{u'(x)}\cdot u(x)$, (me2) gives, for $u,u'\in\mathrm{HH}$,
\[
(u',u)=\sum_E(u',u)_E,\qquad\|u\|^2=\sum_E\|u\|_{L^2,E}^2;
\]
the sums do not depend on the choice of the cell, and every $L$-block of the covering lattice $\mathrm{me}\,\mathbb{Z}^3$ belongs to such a cell.
\item[(b3)] \emph{Projections.}
For an $L$-block $E$ and $u\in L^2(E\times O_{\mathrm{cub}};\mathbb{C}^4)$ put
\[
(P_Eu)(x)=L^{-3}\int_E\mathrm{d}x'\,u(x')\quad(x\text{ over }E),\qquad
(P'_Eu)(x):=\int_{\Delta_z}\mathrm{d}x'\,u(x')=\mathrm{me}^3\sum_{c'\in\Delta_z}\frac{1}{48}\sum_{g'\in O_{\mathrm{cub}}}u(c',g')
\]
for $x$ over a unit block $\Delta_z\subset E$.
Thus $P_E$ is the mean over the $48\,L^{3(k+1)}$ sites of $E\times O_{\mathrm{cub}}$ (Section~G1.4), and $P'_E$ is the mean over the $48\,L^{3k}$ sites of a unit block times the cluster.
Both are orthogonal projections of $L^2(E\times O_{\mathrm{cub}};\mathbb{C}^4)$, onto the functions that are constant on $E\times O_{\mathrm{cub}}$, respectively on each $\Delta_z\times O_{\mathrm{cub}}$.
They commute with constant spin matrices, and $P_EP'_E=P'_EP_E=P_E$ by (b1).
Consequently $P'_E-P_E$ and $P_E$ are orthogonal projections whose ranges are orthogonal to each other.
Their kernels with respect to $\int\mathrm{d}x$ are
\[
P_E(x,x')=L^{-3}\,1_4,\qquad P'_E(x,x')=\mathbf{1}[c\text{ and }c'\text{ lie in one unit block}]\,1_4.
\]
Let $Z$ range over the unit blocks, or over the $L$-blocks.
Consider the map that replaces an element $u$ of $\mathrm{HH}$, over each $Z$, by its mean over $Z\times O_{\mathrm{cub}}$.
It preserves the character, because the translations by $R\,e_\mu$ permute the blocks $Z$.
The mean over $Z\times O_{\mathrm{cub}}$ is an orthogonal projection of $L^2(Z\times O_{\mathrm{cub}};\mathbb{C}^4)$; by (me2) with $j=k$, respectively $j=k+1$, the map therefore does not increase the norm of $\mathrm{HH}$, is idempotent and is Hermitian.
Its range consists of the elements of $\mathrm{HH}$ that are constant on each $Z\times O_{\mathrm{cub}}$.
So the map is the orthogonal projection $JP_k(0)J^{*}$, respectively $JP_{k+1}(0)J^{*}$, of the data of Proposition~1.6: for $u\in\mathrm{HH}$ and $x$ over $E$,
\[
\left(JP_k(0)J^{*}u\right)(x)=(P'_Eu)(x),\qquad\left(JP_{k+1}(0)J^{*}u\right)(x)=(P_Eu)(x),
\]
where on the right $u$ is restricted to $E\times O_{\mathrm{cub}}$.
Hence
\[
(\operatorname{Pen}u)(x)=(\operatorname{Pen}_Eu)(x)\quad\text{for }x\text{ over }E,\qquad
\operatorname{Pen}_E:=\boldsymbol{\alpha}_yP'_E+\boldsymbol{\beta}_yP_E=\boldsymbol{\alpha}_y\left(P'_E-P_E\right)+\left(\boldsymbol{\alpha}_y+\boldsymbol{\beta}_y\right)P_E.
\]
\item[(b4)] \emph{Energy.}
For $u\in\mathrm{HH}$, (me3) with $j=k+1$ gives
\[
(u,\nabla^{*}\nabla u)=\frac{1}{2}\sum_{i=0}^{3}w_i\int\mathrm{d}x\,\left|u(x^{(i)})-u(x)\right|_2^2\ge\sum_E(u,\nabla^{*}\nabla u)_E.
\]
\item[(b5)] \emph{Boundedness.}
By (me4), $\|D\|\le\mathrm{me}^{-1}(1+\rho)\cdot\mathrm{sw}$.
Together with Lemma~6.1\,(Pn3$'$) below, it follows that $F_y(0)$ is a bounded operator of $\mathrm{HH}$.
\end{itemize}

\subsubsection*{2.2. The penalty}

\begin{named}{Lemma 6.1}
Let $k\ge1$, $y\in\mathbb{R}\cup\{\infty\}$ and $v\in\mathbb{C}^4$; here $v'\cdot v$ is the bilinear product on $\mathbb{C}^4$, so that $\overline{v}\cdot mv$ is the Hermitian form of a $4\times4$ matrix $m$.
Then:
\begin{itemize}
\item[(al1)] $|\boldsymbol{\alpha}_yv|_2^2=b_k\operatorname{Re}\left(\overline{v}\cdot\boldsymbol{\alpha}_yv\right)$;
\item[(al2)] $\left|(\boldsymbol{\alpha}_y+\boldsymbol{\beta}_y)v\right|_2^2\le b_k\operatorname{Re}\left(\overline{v}\cdot(\boldsymbol{\alpha}_y+\boldsymbol{\beta}_y)v\right)$;
\item[(al3)] $\operatorname{Re}\left(\overline{v}\cdot(\boldsymbol{\alpha}_y+\boldsymbol{\beta}_y)v\right)\ge\frac{b_-}{L}\,|v|_2^2$;
\item[(al4)] $|\boldsymbol{\alpha}_y|_2\le b$ and $|\boldsymbol{\beta}_y|_2\le b_L$;
\item[(al5)] $\boldsymbol{\alpha}_y^{*}=\boldsymbol{\alpha}_{-y}$ and $\boldsymbol{\beta}_y^{*}=\boldsymbol{\beta}_{-y}$; $\boldsymbol{\alpha}_y^{T}=\mathsf{C}^{-1}\boldsymbol{\alpha}_{-y}\mathsf{C}$ and $\boldsymbol{\beta}_y^{T}=\mathsf{C}^{-1}\boldsymbol{\beta}_{-y}\mathsf{C}$; and $\boldsymbol{\alpha}_y$, $\boldsymbol{\beta}_y$ commute with $\gamma_5$.
\end{itemize}
For every $L$-block $E$ and every $u\in L^2(E\times O_{\mathrm{cub}};\mathbb{C}^4)$, moreover:
\begin{itemize}
\item[(Pn1)] $\operatorname{Re}(u,\operatorname{Pen}_Eu)_E\ge\frac{b_-}{L}\,\|P_Eu\|_{L^2,E}^2$;
\item[(Pn2)] $\|\operatorname{Pen}_Eu\|_{L^2,E}^2\le b\operatorname{Re}(u,\operatorname{Pen}_Eu)_E$;
\item[(Pn3)] $(\operatorname{Pen}_Eu)(x)=\int_E\mathrm{d}x'\,K_E(x,x')\,u(x')$ with $K_E(x,x')=\boldsymbol{\alpha}_y\,\mathbf{1}[c,c'\text{ in one unit block}]+\boldsymbol{\beta}_y\,L^{-3}$, and $|K_E(x,x')|_2\le2b$;
\item[(Pn3$'$)] $\|\operatorname{Pen}_Eu\|_{L^2,E}\le(b+b_L)\,\|u\|_{L^2,E}$.
\end{itemize}
\end{named}

\begin{named}{Remark}
The inequalities (Pn1) and (Pn3) are the conditions (Pen1) and (Pen3) of the requirement (B3) in \cite[Definition~12.8]{QEDI}, read on $L^2(E\times O_{\mathrm{cub}};\mathbb{C}^4)$ with $\ell_E=L$, and (Pn2) is the first of the two inequalities of (Pen2) there.
The constants are those of \cite[Lemma~12.10]{QEDI}, $c_1=b_-$, $C_1=L\cdot b$ and $C_2=2L^4\cdot b$, for which $b_-\cdot\ell_E^{-1}=b_-/L$, $L\cdot b\cdot\ell_E^{-1}=b$ and $2L^4\cdot b\cdot\ell_E^{-4}=2b$.
The second inequality of (Pen2), $\|\operatorname{Pen}_E^{*}u\|_{L^2,E}^2\le b\operatorname{Re}(u,\operatorname{Pen}_Eu)_E$, is (Pn2) for $-y$: by (al5), $\operatorname{Pen}_E^{*}$ is the operator $\operatorname{Pen}_E$ formed with $-y$, and $\operatorname{Re}(u,\operatorname{Pen}_E^{*}u)_E=\operatorname{Re}(u,\operatorname{Pen}_Eu)_E$.
\end{named}

\begin{proof}
By (I.c1), $\gamma_5$ is Hermitian and $\gamma_5^2=1$.
Hence $\mathbb{C}^4$ is the orthogonal sum of the eigenspaces of $\gamma_5$ for the eigenvalues $\mathrm{sn}\in\{1,-1\}$.
The matrices $\boldsymbol{\alpha}_y$ and $\boldsymbol{\beta}_y$ of Section~I.3.5 are rational functions of $\gamma_5$; they commute with $\gamma_5$, and on the eigenspace of $\mathrm{sn}$ they act as multiplication by numbers $\mathrm{ev}_1$ and $\mathrm{ev}_3$, respectively.
Let $\mathrm{ev}_2:=\mathrm{ev}_1+\mathrm{ev}_3$.
For real $y$ put $\mathrm{ys}:=\mathrm{sn}\cdot y$; then
\[
\mathrm{ev}_1=\frac{b_k\cdot i\,\mathrm{ys}}{b_k+i\,\mathrm{ys}}=\frac{b_k\cdot\mathrm{ys}^2+i\,b_k^2\cdot\mathrm{ys}}{b_k^2+\mathrm{ys}^2},\qquad
\mathrm{ev}_3=\frac{b_k^2\cdot b_L}{\left(b_k+b_L+i\,\mathrm{ys}\right)\left(b_k+i\,\mathrm{ys}\right)},\qquad
\mathrm{ev}_2=\frac{b_k\left(b_L+i\,\mathrm{ys}\right)}{b_k+b_L+i\,\mathrm{ys}},
\]
the last formula because $b_k\cdot i\,\mathrm{ys}\left(b_k+b_L+i\,\mathrm{ys}\right)+b_k^2\cdot b_L=b_k\left(b_k+i\,\mathrm{ys}\right)\left(b_L+i\,\mathrm{ys}\right)$.
It follows that
\begin{gather*}
|\mathrm{ev}_1|^2=\frac{b_k^2\cdot\mathrm{ys}^2}{b_k^2+\mathrm{ys}^2}=b_k\operatorname{Re}\mathrm{ev}_1\le b_k^2,\qquad
|\mathrm{ev}_3|\le\frac{b_k^2\cdot b_L}{\left(b_k+b_L\right)b_k}\le b_L,\\
\operatorname{Re}\mathrm{ev}_2=b_k\cdot\frac{b_L\left(b_k+b_L\right)+\mathrm{ys}^2}{\left(b_k+b_L\right)^2+\mathrm{ys}^2}\ge\frac{b_k\cdot b_L}{b_k+b_L},\qquad
|\mathrm{ev}_2|^2=\frac{b_k^2\left(b_L^2+\mathrm{ys}^2\right)}{\left(b_k+b_L\right)^2+\mathrm{ys}^2}\le b_k\operatorname{Re}\mathrm{ev}_2;
\end{gather*}
the lower bound for $\operatorname{Re}\mathrm{ev}_2$ holds because, for $0<t_1\le t_2$, the quotient $(t_1+\mathrm{ts})/(t_2+\mathrm{ts})$ is non-decreasing in $\mathrm{ts}\ge0$, and the bound for $|\mathrm{ev}_2|^2$ holds because $b_L^2\le b_L\left(b_k+b_L\right)$.
For $y=\infty$ one has $\mathrm{ev}_1=\mathrm{ev}_2=b_k$ and $\mathrm{ev}_3=0$, so that $|\mathrm{ev}_1|^2=b_k\operatorname{Re}\mathrm{ev}_1=b_k^2$, $|\mathrm{ev}_3|=0\le b_L$, $\operatorname{Re}\mathrm{ev}_2=b_k\ge b_k\cdot b_L/(b_k+b_L)$ and $|\mathrm{ev}_2|^2=b_k\operatorname{Re}\mathrm{ev}_2$.
By Section~I.3.4, $b_k\cdot b_L/(b_k+b_L)=b_{k+1}/L\ge b_-/L$ and $b_k\le b$.
Decomposing $v$ into its components in the two eigenspaces, one obtains (al1)--(al4).

For (al5), note that $\gamma_5^{*}=\gamma_5$; hence the adjoint of a rational function of $i\gamma_5y$ with real coefficients is the same function of $-i\gamma_5y$.
By (I.c3), $\gamma_5^{T}=\mathsf{C}^{-1}(-\gamma_5)\mathsf{C}$; hence the transpose of a rational function of $i\gamma_5y$ is $\mathsf{C}^{-1}$ times the same function of $-i\gamma_5y$ times $\mathsf{C}$.
For $y=\infty$ one has $\boldsymbol{\alpha}_y=b_k$ and $\boldsymbol{\beta}_y=0$.

Now let $u_1:=(P'_E-P_E)u$ and $u_2:=P_Eu$.
By (b3), $\boldsymbol{\alpha}_yu_1$ lies in the range of $P'_E-P_E$ and $(\boldsymbol{\alpha}_y+\boldsymbol{\beta}_y)u_2$ in the range of $P_E$, and these two ranges are orthogonal to each other and to $(1-P'_E)u$.
Therefore
\[
(u,\operatorname{Pen}_Eu)_E=(u_1,\boldsymbol{\alpha}_yu_1)_E+\left(u_2,(\boldsymbol{\alpha}_y+\boldsymbol{\beta}_y)u_2\right)_E,\qquad
\|\operatorname{Pen}_Eu\|_{L^2,E}^2=\|\boldsymbol{\alpha}_yu_1\|_{L^2,E}^2+\left\|(\boldsymbol{\alpha}_y+\boldsymbol{\beta}_y)u_2\right\|_{L^2,E}^2.
\]
Applying (al1)--(al3) at each site and integrating over $E$, one obtains $\operatorname{Re}(u_1,\boldsymbol{\alpha}_yu_1)_E\ge0$ and $\operatorname{Re}\left(u_2,(\boldsymbol{\alpha}_y+\boldsymbol{\beta}_y)u_2\right)_E\ge\frac{b_-}{L}\,\|u_2\|_{L^2,E}^2$, which give (Pn1), and $\|\operatorname{Pen}_Eu\|_{L^2,E}^2\le b_k\operatorname{Re}(u,\operatorname{Pen}_Eu)_E$, which together with $b_k\le b$ gives (Pn2).
(Pn3) follows from the kernels in (b3), from (al4) and from $b+b_L\cdot L^{-3}\le2b$.
(Pn3$'$) follows from (al4) and from $\|P'_E\|,\|P_E\|\le1$.
\end{proof}

\subsubsection*{2.3. The $L^2$ bound; proof of (a) and (b)}

\begin{named}{Lemma 6.2}
For every $u\in\mathrm{HH}$,
\[
\|u\|\le L\cdot C_I\cdot\|F_y(0)u\|.
\]
\end{named}

\begin{proof}
Let $f:=F_y(0)u$ and $\mathrm{Xr}:=\operatorname{Re}(u,f)$, so that $\mathrm{Xr}\le\|u\|\,\|f\|$.
By Proposition~1.3\,(a) and by (b2), (b3) of subsection~2.1,
\begin{equation}\tag{VI.1}
\mathrm{Xr}=\frac{\rho}{2\,\mathrm{me}}\,(u,\nabla^{*}\nabla u)+\mu\,\|u\|^2+\sum_E\operatorname{Re}(u,\operatorname{Pen}_Eu)_E,
\end{equation}
and the three terms on the right are non-negative, by (b4), by $\mu\ge0$ and by (Pn1).
Each of them is therefore at most $\mathrm{Xr}$, with the following consequences:
\begin{itemize}
\item[(i)] $\sum_E\|P_Eu\|_{L^2,E}^2\le\frac{L}{b_-}\,\mathrm{Xr}$, by (Pn1);
\item[(ii)] $\|\operatorname{Pen}u\|^2=\sum_E\|\operatorname{Pen}_Eu\|_{L^2,E}^2\le b\cdot\mathrm{Xr}$, by (Pn2);
\item[(iii)] $\|\mu u\|^2=\mu\left(\mu\,\|u\|^2\right)\le\mu\,\mathrm{Xr}\le\mathrm{Xr}$, because $\mu\le1$.
\end{itemize}
Since $Du=f-\mu u-\operatorname{Pen}u$, (ii), (iii) and the identity $1+b^{\frac{1}{2}}=K_1\cdot L^{-\frac{1}{2}}$ give $\|Du\|\le\|f\|+K_1\cdot L^{-\frac{1}{2}}\cdot\mathrm{Xr}^{\frac{1}{2}}$.
Proposition~1.3\,(b) then yields
\begin{equation}\tag{VI.2}
c_b\cdot\mathrm{me}^{-2}\,(u,\nabla^{*}\nabla u)\le\|Du\|^2\le2\,\|f\|^2+2K_1^2\cdot L^{-1}\cdot\mathrm{Xr}\le2\,\|f\|^2+2K_1^2\cdot L^{-1}\,\|u\|\,\|f\|.
\end{equation}
An $L$-block $E$ is a cube of $(L/\mathrm{me})^3$ sites with $\ell_E=L$.
By Proposition~1.4 and by (b4) of subsection~2.1, therefore,
\[
\sum_E\|u-P_Eu\|_{L^2,E}^2\le C_{\mathrm{Po}}\cdot L^2\cdot\mathrm{me}^{-2}\sum_E(u,\nabla^{*}\nabla u)_E\le C_{\mathrm{Po}}\cdot L^2\cdot\mathrm{me}^{-2}\,(u,\nabla^{*}\nabla u).
\]
Combining this with (b2), (i) and (VI.2), one arrives at
\[
\|u\|^2=\sum_E\left(\|P_Eu\|_{L^2,E}^2+\|u-P_Eu\|_{L^2,E}^2\right)\le L\left(\frac{1}{b_-}+\frac{2C_{\mathrm{Po}}\cdot K_1^2}{c_b}\right)\|u\|\,\|f\|+\frac{2C_{\mathrm{Po}}}{c_b}\,L^2\,\|f\|^2.
\]
Now an inequality $\|u\|^2\le t_1\|u\|\,\|f\|+t_2\|f\|^2$ with numbers $t_1,t_2\ge0$ implies $\|u\|\le\left(t_1+t_2^{\frac{1}{2}}\right)\|f\|$.
Indeed, otherwise $\|u\|>t_2^{\frac{1}{2}}\|f\|$ and
\[
\|u\|^2>\left(t_1+t_2^{\frac{1}{2}}\right)\|f\|\,\|u\|\ge t_1\|u\|\,\|f\|+t_2^{\frac{1}{2}}\|f\|\cdot t_2^{\frac{1}{2}}\|f\|,
\]
which is impossible.
Hence
\[
\|u\|\le L\left[\frac{1}{b_-}+\frac{2C_{\mathrm{Po}}\cdot K_1^2}{c_b}+\left(\frac{2C_{\mathrm{Po}}}{c_b}\right)^{\frac{1}{2}}\right]\|f\|=L\cdot C_I\cdot\|f\|.
\]
\end{proof}

\begin{proof}[Proof of Proposition~1.6\,(a)]
\emph{Adjoint.}
On $\mathrm{HH}$ one has $D^{*}=\gamma_5D\gamma_5$ by Proposition~1.3\,(d); the number $\mu$ is real; and $JP_k(0)J^{*}$, $JP_{k+1}(0)J^{*}$ are orthogonal projections that commute with spin matrices.
By (al5),
\[
F_y(0)^{*}=\gamma_5D\gamma_5+\mu+\boldsymbol{\alpha}_{-y}\,JP_k(0)J^{*}+\boldsymbol{\beta}_{-y}\,JP_{k+1}(0)J^{*}=\gamma_5\,F_{-y}(0)\,\gamma_5,
\]
because $\gamma_5^2=1$ and $\gamma_5$ commutes with $\boldsymbol{\alpha}_{-y}$, $\boldsymbol{\beta}_{-y}$ and with the two projections.

\emph{Transpose.}
Complex conjugation maps $\mathrm{HH}(\eta_k,R;-\theta)$ onto $\mathrm{HH}(\eta_k,R;\theta)$, and $\langle f',f\rangle^{\mathrm{ch}}=(\overline{f'},f)$ for $f'\in\mathrm{HH}(\eta_k,R;-\theta)$ and $f\in\mathrm{HH}$.
Every bounded operator $O$ of $\mathrm{HH}$ therefore has the transpose $O^{T}f':=\overline{O^{*}\,\overline{f'}}$, a bounded operator of $\mathrm{HH}(\eta_k,R;-\theta)$ with $\langle O^{T}f',f\rangle^{\mathrm{ch}}=\langle f',Of\rangle^{\mathrm{ch}}$; moreover $(O_1O_2)^{T}=O_2^{T}O_1^{T}$.
By Proposition~1.3\,(d), $D^{T}=\mathsf{C}^{-1}D\,\mathsf{C}$.
Let $\mathrm{Pb}$ be one of the two projections.
By (b2) and (b3), $\langle f',\mathrm{Pb}f\rangle^{\mathrm{ch}}$ is the sum, over the unit blocks (respectively the $L$-blocks) $Z$ of a period cell, of
\[
|Z|^{-1}\left(\int_Z\mathrm{d}x\,f'(x)\right)\cdot\left(\int_Z\mathrm{d}x\,f(x)\right),
\]
an expression symmetric in $f'$ and $f$.
Hence $\mathrm{Pb}^{T}$ is the projection given by the same formula on $\mathrm{HH}(\eta_k,R;-\theta)$.
For a constant matrix $m$ one has $\langle f',m\,\mathrm{Pb}f\rangle^{\mathrm{ch}}=\langle m^{T}f',\mathrm{Pb}f\rangle^{\mathrm{ch}}$, whence $(m\,\mathrm{Pb})^{T}=\mathrm{Pb}\,m^{T}=m^{T}\,\mathrm{Pb}$.
With (al5) it follows that, on $\mathrm{HH}(\eta_k,R;-\theta)$,
\[
F_y(0)^{T}=\mathsf{C}^{-1}D\,\mathsf{C}+\mu+\mathsf{C}^{-1}\boldsymbol{\alpha}_{-y}\mathsf{C}\,JP_k(0)J^{*}+\mathsf{C}^{-1}\boldsymbol{\beta}_{-y}\mathsf{C}\,JP_{k+1}(0)J^{*}=\mathsf{C}^{-1}F_{-y}(0)\,\mathsf{C}.
\]

\emph{Bijectivity.}
The operator $F_y(0)$ is bounded by (b5).
By Lemma~6.2 it is injective, and its range is closed, by the classical facts recalled in subsection~0.3: a bounded operator that is bounded below has a closed range, and the orthogonal complement of the range of a bounded operator is the kernel of its adjoint.
The matrix $\gamma_5$ acts as a unitary operator of $\mathrm{HH}$; so Lemma~6.2 for $-y$ shows that $F_y(0)^{*}=\gamma_5F_{-y}(0)\gamma_5$ is injective.
The orthogonal complement of the range of $F_y(0)$, which is the kernel of $F_y(0)^{*}$, is therefore $\{0\}$.
Hence $F_y(0)$ is a bijection of $\mathrm{HH}$.
(If $R<\infty$, then $\mathrm{HH}$ has the finite dimension $192\,(R/\mathrm{me})^3$, and injectivity alone suffices.)
By Lemma~6.2 the inverse is bounded.

\emph{Inverses.}
Let $S_y(0):=F_y(0)^{-1}$, as in Section~I.3.5.
From $S_y(0)F_y(0)=F_y(0)S_y(0)=1$ one obtains $F_y(0)^{*}S_y(0)^{*}=S_y(0)^{*}F_y(0)^{*}=1$ and $F_y(0)^{T}S_y(0)^{T}=S_y(0)^{T}F_y(0)^{T}=1$.
Therefore $\left(F_y(0)^{*}\right)^{-1}=\left(F_y(0)^{-1}\right)^{*}$ and $\left(F_y(0)^{T}\right)^{-1}=\left(F_y(0)^{-1}\right)^{T}$.
\end{proof}

\begin{proof}[Proof of Proposition~1.6\,(b)]
Apply Lemma~6.2 to $u=F_y(0)^{-1}f$.
\end{proof}

\subsubsection*{2.4. Exponential weights}

\begin{named}{Lemma 6.3}
Let $\psi$ be a real, bounded, $R\mathbb{Z}^3$-periodic function of $c\in\mathrm{me}\,\mathbb{Z}^3$ such that
\[
|\psi(c)-\psi(c')|\le a_0\cdot\frac{|c-c'|_1}{L}\qquad\text{for all }c,c'.
\]
Then every $f\in\mathrm{HH}$ satisfies
\[
\left\|\exp(\psi)\,F_y(0)^{-1}f\right\|\le2C_I\cdot L\left\|\exp(\psi)\,f\right\|.
\]
\end{named}

\begin{proof}
Since $\psi$ is real, bounded and $R\mathbb{Z}^3$-periodic, the multiplications by $\exp(\psi)$ and by $\exp(-\psi)$ are bounded operators of $\mathrm{HH}$.
Put
\[
\mathrm{Rm}:=\exp(\psi)\,F_y(0)\exp(-\psi)-F_y(0)=\exp(\psi)\,[D,\exp(-\psi)]+\left(\exp(\psi)\operatorname{Pen}\exp(-\psi)-\operatorname{Pen}\right);
\]
the term $\mu$ commutes with the multiplications and cancels.
Fix $\mathrm{vv}\in\mathrm{HH}$.

\emph{The hopping term.}
Proposition~1.3\,(c), applied to the function $\exp(-\psi)$, gives
\[
\left(\exp(\psi)\,[D,\exp(-\psi)]\,\mathrm{vv}\right)(x)=-\mathrm{me}^{-1}\cdot M_3(x)\left(\exp\left(\psi(c)-\psi(c_3)\right)-1\right)\mathrm{vv}\left(x^{(3)}\right).
\]
Here $|c-c_3|_1=\mathrm{me}$, so that $|\psi(c)-\psi(c_3)|\le a_0\cdot\mathrm{me}/L\le1$ (recall $a_0\le1/3$ and $\mathrm{me}\le1/L$); moreover $|\exp(\mathrm{ts})-1|\le\mathrm{e}\,|\mathrm{ts}|$ for real $\mathrm{ts}$ with $|\mathrm{ts}|\le1$.
As $|M_3(x)|_2=3(1+\rho)$, the value of the left side at $x$ has the norm $|\cdot|_2$ at most $3(1+\rho)\cdot\mathrm{e}\cdot(a_0/L)\left|\mathrm{vv}\left(x^{(3)}\right)\right|_2$.
The map $x\mapsto x^{(3)}$ is a bijection which commutes with the translations, and all sites carry the same weight $\mathrm{me}^3/48$; hence $\int\mathrm{d}x\,\left|\mathrm{vv}\left(x^{(3)}\right)\right|_2^2=\|\mathrm{vv}\|^2$, and
\[
\left\|\exp(\psi)\,[D,\exp(-\psi)]\,\mathrm{vv}\right\|\le3(1+\rho)\cdot\mathrm{e}\cdot\frac{a_0}{L}\,\|\mathrm{vv}\|.
\]

\emph{The penalty term.}
Let $x$ lie over an $L$-block $E$.
By (Pn3),
\[
\left(\left(\exp(\psi)\operatorname{Pen}\exp(-\psi)-\operatorname{Pen}\right)\mathrm{vv}\right)(x)=\int_E\mathrm{d}x'\,K_E(x,x')\left(\exp\left(\psi(c)-\psi(c')\right)-1\right)\mathrm{vv}(x').
\]
For $c,c'\in E$ one has $|c-c'|_1<3L$ by (b1), whence $|\psi(c)-\psi(c')|\le3a_0\le1$.
Together with the bound $|K_E(x,x')|_2\le2b$ of (Pn3) and with the Schwarz inequality ($|E|=L^3$), this shows that the value of the left side at $x$ has the norm $|\cdot|_2$ at most $2b\cdot3\mathrm{e}\cdot a_0\cdot L^{\frac32}\,\|\mathrm{vv}\|_{L^2,E}$.
It follows that
\[
\left\|\left(\exp(\psi)\operatorname{Pen}\exp(-\psi)-\operatorname{Pen}\right)\mathrm{vv}\right\|_{L^2,E}\le L^{\frac32}\cdot6b\cdot\mathrm{e}\cdot a_0\cdot L^{\frac32}\,\|\mathrm{vv}\|_{L^2,E}=6L^4\cdot b\cdot\mathrm{e}\cdot\frac{a_0}{L}\,\|\mathrm{vv}\|_{L^2,E}.
\]
Summing the squares over the $L$-blocks $E$ of a period cell, by (b2), one obtains
\[
\left\|\left(\exp(\psi)\operatorname{Pen}\exp(-\psi)-\operatorname{Pen}\right)\mathrm{vv}\right\|\le6L^4\cdot b\cdot\mathrm{e}\cdot\frac{a_0}{L}\,\|\mathrm{vv}\|.
\]

\emph{Conclusion.}
Since $c_{\mathrm{wt}}=\mathrm{e}\left(3(1+\rho)+6L^4\cdot b\right)$, the two bounds add up to
\[
\|\mathrm{Rm}\,\mathrm{vv}\|\le c_{\mathrm{wt}}\cdot\frac{a_0}{L}\,\|\mathrm{vv}\|\qquad\text{for every }\mathrm{vv}\in\mathrm{HH}.
\]
Now choose $\mathrm{vv}:=\exp(\psi)\,F_y(0)^{-1}f$.
Then
\[
F_y(0)\,\mathrm{vv}=\exp(\psi)\,F_y(0)\exp(-\psi)\,\mathrm{vv}-\mathrm{Rm}\,\mathrm{vv}=\exp(\psi)\,f-\mathrm{Rm}\,\mathrm{vv}.
\]
Lemma~6.2 and the inequality $C_I\cdot c_{\mathrm{wt}}\cdot a_0\le\frac12$, which holds by the definition of $a_0$, give
\[
\|\mathrm{vv}\|\le L\cdot C_I\left\|\exp(\psi)\,f\right\|+C_I\cdot c_{\mathrm{wt}}\cdot a_0\,\|\mathrm{vv}\|\le L\cdot C_I\left\|\exp(\psi)\,f\right\|+\frac12\,\|\mathrm{vv}\|.
\]
As $\|\mathrm{vv}\|$ is finite, this yields $\|\mathrm{vv}\|\le2C_I\cdot L\left\|\exp(\psi)\,f\right\|$.
\end{proof}

\subsubsection*{2.5. Interior regularity, uniformly in the mesh}

In Lemma~6.4 the mesh $\mathrm{me}>0$ is arbitrary, $D$ is the operator of Section~G1.4 acting on functions on $\mathrm{me}\,\mathbb{Z}^3\times O_{\mathrm{cub}}$, $G$ is the function of Proposition~1.5\,(e), and $C_R$ is the constant of Proposition~1.6:
\[
C_R=\frac{18\,000\,000\,C_G}{1-\alpha}.
\]

\begin{named}{Lemma 6.4}
Let $c^{\mathrm{o}}\in\mathrm{me}\,\mathbb{Z}^3$, $\mathrm{rad}\ge8\,\mathrm{me}$ and $0\le\mu\le1/\mathrm{me}$, and put
\[
\mathrm{Bo}:=\left\{c\in\mathrm{me}\,\mathbb{Z}^3:|c-c^{\mathrm{o}}|<\mathrm{rad}\right\},\qquad\mathrm{Bo}^+:=\left\{|c-c^{\mathrm{o}}|<\mathrm{rad}+\mathrm{me}\right\},\qquad\mathrm{Bo}':=\left\{|c-c^{\mathrm{o}}|\le\frac{\mathrm{rad}}{4}\right\}.
\]
Let $u$ be a function on $\mathrm{Bo}^+\times O_{\mathrm{cub}}$ and $\mathrm{fs}$ a function on $\mathrm{Bo}\times O_{\mathrm{cub}}$, both with values in $\mathbb{C}^4$, such that $\left((D+\mu)u\right)(x)=\mathrm{fs}(x)$ for all $x$ over $\mathrm{Bo}$, and put
\[
\mathrm{NN}:=\mathrm{rad}\cdot|\mathrm{fs}|_{\mathrm{su},\mathrm{Bo}}+\mathrm{rad}^{-\frac32}\,\|u\|_{L^2,\mathrm{Bo}^+}.
\]
Then for all sites $x\ne\tilde x$ over $\mathrm{Bo}'$:
\begin{itemize}
\item[(a)] $|u(x)|_2\le C_R\cdot\mathrm{NN}$;
\item[(b)] $|u(x)-u(\tilde x)|_2\le C_R\left(\dfrac{d_X(x,\tilde x)}{\mathrm{rad}}\right)^{\alpha}\mathrm{NN}$.
\end{itemize}
(For $x$ over $\mathrm{Bo}$ the four neighbours $x^{(i)}$ lie over $\mathrm{Bo}^+$, so that $(Du)(x)$ is defined.)
\end{named}

\begin{proof}
Let
\[
\mathrm{ct}(c):=\min\left\{1,\max\left\{0,2-\frac{2|c-c^{\mathrm{o}}|}{\mathrm{rad}}\right\}\right\}.
\]
Then $\mathrm{ct}=1$ for $|c-c^{\mathrm{o}}|\le\mathrm{rad}/2$, $\mathrm{ct}=0$ for $|c-c^{\mathrm{o}}|\ge\mathrm{rad}$, and $|\mathrm{ct}(c)-\mathrm{ct}(c')|\le2|c-c'|/\mathrm{rad}$.
Extend $u$ by zero off $\mathrm{Bo}^+\times O_{\mathrm{cub}}$ and let $v:=\mathrm{ct}\cdot u$, a finitely supported function on $\mathrm{me}\,\mathbb{Z}^3\times O_{\mathrm{cub}}$.
At every site
\[
(D+\mu)v=\mathrm{ct}\cdot(D+\mu)u+[D,\mathrm{ct}]\,u.
\]
If $x$ lies over $\mathrm{Bo}$, its neighbours lie over $\mathrm{Bo}^+$ and $\left((D+\mu)u\right)(x)=\mathrm{fs}(x)$; if $x$ does not lie over $\mathrm{Bo}$, then $\mathrm{ct}(c)=0$.
Hence
\[
(D+\mu)v=\mathrm{ct}\cdot\mathrm{fs}+\mathrm{kf},\qquad\mathrm{kf}:=[D,\mathrm{ct}]\,u,
\]
where $\mathrm{ct}\cdot\mathrm{fs}:=0$ off $\mathrm{Bo}\times O_{\mathrm{cub}}$.
By Proposition~1.5\,(e),
\[
v=\mathrm{Gc}[\mathrm{ct}\cdot\mathrm{fs}]+\mathrm{Gc}[\mathrm{kf}],
\]
where, for a finitely supported function $\mathrm{vv}$,
\[
\mathrm{Gc}[\mathrm{vv}](x):=\int\mathrm{d}x'\,G(x,x')\,\mathrm{vv}(x')=\mathrm{me}^3\sum_{c'}\frac1{48}\sum_{g'\in O_{\mathrm{cub}}}G\left(x,(c',g')\right)\mathrm{vv}(c',g').
\]

\emph{The function $\mathrm{kf}$.}
By Proposition~1.3\,(c),
\[
\mathrm{kf}(x)=-\mathrm{me}^{-1}\cdot M_3(x)\left(\mathrm{ct}(c_3)-\mathrm{ct}(c)\right)u\left(x^{(3)}\right),
\]
so that
\[
|\mathrm{kf}(x)|_2\le\mathrm{me}^{-1}\cdot3(1+\rho)\cdot\frac{2\,\mathrm{me}}{\mathrm{rad}}\left|u\left(x^{(3)}\right)\right|_2\le\frac{12}{\mathrm{rad}}\left|u\left(x^{(3)}\right)\right|_2.
\]
If $\mathrm{ct}(c_3)\ne\mathrm{ct}(c)$, then $\max\left(|c-c^{\mathrm{o}}|,|c_3-c^{\mathrm{o}}|\right)>\mathrm{rad}/2$; as $|c_3-c|=\mathrm{me}$, it follows that $\mathrm{kf}(x)\ne0$ implies $|c-c^{\mathrm{o}}|>\mathrm{rad}/2-\mathrm{me}$.
The map $x\mapsto x^{(3)}$ is a bijection, all sites carry the same weight $\mathrm{me}^3/48$, and $u$ vanishes off $\mathrm{Bo}^+\times O_{\mathrm{cub}}$; therefore $\int\mathrm{d}x\,\left|u\left(x^{(3)}\right)\right|_2=\int_{\mathrm{Bo}^+}\mathrm{d}x'\,|u(x')|_2$.
The Schwarz inequality gives
\[
\|\mathrm{kf}\|_{L^1}\le\frac{12}{\mathrm{rad}}\int_{\mathrm{Bo}^+}\mathrm{d}x'\,|u(x')|_2\le\frac{12}{\mathrm{rad}}\,|\mathrm{Bo}^+|^{\frac12}\,\|u\|_{L^2,\mathrm{Bo}^+}\le45\,\mathrm{rad}^{\frac12}\,\|u\|_{L^2,\mathrm{Bo}^+};
\]
indeed, an open interval of length $2(\mathrm{rad}+\mathrm{me})$ contains at most $2(\mathrm{rad}+\mathrm{me})/\mathrm{me}+1$ points of $\mathrm{me}\,\mathbb{Z}$, so that $|\mathrm{Bo}^+|\le(2\,\mathrm{rad}+3\,\mathrm{me})^3\le(2.4\,\mathrm{rad})^3$ (as $\mathrm{me}\le\mathrm{rad}/8$), and $12\cdot2.4^{\frac32}\le45$.

\emph{The term $\mathrm{Gc}[\mathrm{kf}]$.}
Let $x'=(c',g')$ be a site over $\mathrm{Bo}'$ and $x''=(c'',g'')$ a site with $\mathrm{kf}(x'')\ne0$.
Then $|c'-c''|>\mathrm{rad}/2-\mathrm{me}-\mathrm{rad}/4\ge\mathrm{rad}/8$.
By Proposition~1.5\,(e),
\[
|G(x',x'')|_2\le2C_G\left(\frac{\mathrm{rad}}{8}\right)^{-2},\qquad\left|G\left(x'^{(i)},x''\right)-G(x',x'')\right|_2\le20C_G\cdot\mathrm{me}\left(\frac{\mathrm{rad}}{8}\right)^{-3}\quad(i=0,\dots,3).
\]
Consequently
\[
|\mathrm{Gc}[\mathrm{kf}](x')|_2\le128\,C_G\cdot\mathrm{rad}^{-2}\,\|\mathrm{kf}\|_{L^1}\le5760\,C_G\cdot\mathrm{rad}^{-\frac32}\,\|u\|_{L^2,\mathrm{Bo}^+},
\]
\[
\left|\mathrm{Gc}[\mathrm{kf}]\left(x'^{(i)}\right)-\mathrm{Gc}[\mathrm{kf}](x')\right|_2\le10240\,C_G\cdot\mathrm{me}\cdot\mathrm{rad}^{-3}\,\|\mathrm{kf}\|_{L^1}\le460800\,C_G\cdot\mathrm{me}\cdot\mathrm{rad}^{-\frac52}\,\|u\|_{L^2,\mathrm{Bo}^+}.
\]
Let $\mathrm{hd}:=d_X(x,\tilde x)$, and let $x=x^0,\dots,x^{\mathrm{nb}}=\tilde x$, with $x^{j+1}=\left(x^j\right)^{(i_j)}$, be the sequence of (s6)$_{\mathrm{me}}$; thus $\mathrm{nb}\le39\,\mathrm{hd}/\mathrm{me}$.
The box spanned by $c$ and $\tilde c$ lies in the cube $\mathrm{Bo}'$, so that all $x^j$ lie over $\mathrm{Bo}'$.
Summing the second bound, taken at $x'=x^j$ and $i=i_j$, over $j=0,\dots,\mathrm{nb}-1$, one finds
\[
|\mathrm{Gc}[\mathrm{kf}](x)-\mathrm{Gc}[\mathrm{kf}](\tilde x)|_2\le39\cdot460800\,C_G\cdot\mathrm{hd}\cdot\mathrm{rad}^{-\frac52}\,\|u\|_{L^2,\mathrm{Bo}^+}=17\,971\,200\,C_G\cdot\frac{\mathrm{hd}}{\mathrm{rad}}\cdot\mathrm{rad}^{-\frac32}\,\|u\|_{L^2,\mathrm{Bo}^+},
\]
and $\mathrm{hd}/\mathrm{rad}\le(\mathrm{hd}/\mathrm{rad})^{\alpha}$, because $\mathrm{hd}=\max\left(|c-\tilde c|,\mathrm{me}\right)\le\mathrm{rad}/2$.

\emph{The term $\mathrm{Gc}[\mathrm{ct}\cdot\mathrm{fs}]$.}
For $x$ over $\mathrm{Bo}'$ and $c'\in\mathrm{Bo}$ one has $|c-c'|<5\,\mathrm{rad}/4$; moreover $0\le\mathrm{ct}\le1$.
By Proposition~1.5\,(e) the norm $|\cdot|_2$ of the integrand $G\left(x,(c',g')\right)\mathrm{ct}(c')\,\mathrm{fs}(c',g')$ is at most $2C_G\left(\mathrm{me}+|c-c'|\right)^{-2}|\mathrm{fs}|_{\mathrm{su},\mathrm{Bo}}$.
This majorant does not depend on $g'$, so that the mean over the $48$ labels $g'$, each with the weight $1/48$, contributes the factor one.
The first bound of \cite[Lemma~12.7(a)]{QEDI} (see subsection~0.3), applied with the radius $5\,\mathrm{rad}/4\ge\mathrm{me}$, then gives
\[
|\mathrm{Gc}[\mathrm{ct}\cdot\mathrm{fs}](x)|_2\le|\mathrm{fs}|_{\mathrm{su},\mathrm{Bo}}\cdot\mathrm{me}^3\sum_{|c'-c|\le5\mathrm{rad}/4}2C_G\left(\mathrm{me}+|c-c'|\right)^{-2}\le2\cdot27\cdot\frac54\,C_G\cdot\mathrm{rad}\cdot|\mathrm{fs}|_{\mathrm{su},\mathrm{Bo}}=67.5\,C_G\cdot\mathrm{rad}\cdot|\mathrm{fs}|_{\mathrm{su},\mathrm{Bo}}.
\]
By Proposition~1.5\,(f), and by the second bound of \cite[Lemma~12.7(a)]{QEDI}, applied with the radius $5\,\mathrm{rad}/4$ to the sums centred at $c$ and at $\tilde c$ (both points of $\mathrm{Bo}'$),
\begin{align*}
&\left|\mathrm{Gc}[\mathrm{ct}\cdot\mathrm{fs}](x)-\mathrm{Gc}[\mathrm{ct}\cdot\mathrm{fs}](\tilde x)\right|_2\\
&\qquad\le6240\,C_G\cdot\mathrm{hd}^{\alpha}\cdot2\cdot27\left(\frac{5\,\mathrm{rad}}{4}\right)^{1-\alpha}(1-\alpha)^{-1}\,|\mathrm{fs}|_{\mathrm{su},\mathrm{Bo}}\\
&\qquad\le\frac{421\,200\,C_G}{1-\alpha}\left(\frac{\mathrm{hd}}{\mathrm{rad}}\right)^{\alpha}\mathrm{rad}\cdot|\mathrm{fs}|_{\mathrm{su},\mathrm{Bo}}.
\end{align*}

\emph{Conclusion.}
Since $\mathrm{ct}=1$ on $\mathrm{Bo}'$, one has $u=v=\mathrm{Gc}[\mathrm{ct}\cdot\mathrm{fs}]+\mathrm{Gc}[\mathrm{kf}]$ at $x$ and at $\tilde x$.
Adding the bounds, (a) holds with the constant $\max(67.5,5760)\,C_G$, and (b) with the constant $\max\left(421\,200/(1-\alpha),17\,971\,200\right)C_G$; both numbers are at most $C_R$.
\end{proof}

\subsubsection*{2.6. From block to block}

With the two abbreviations $C_3^{\mathrm{n}}$ and $C_3''$, the constants $C_3$ and $C_5$ of Proposition~1.6 read
\[
C_3^{\mathrm{n}}:=\frac{\mathrm{e}^{\frac13}}{4}+\left(L^4\cdot b+46\right)\mathrm{e}^{1.3}\,C_I,\qquad C_3'':=\max\left\{C_R\cdot C_3^{\mathrm{n}},\,14L^3\cdot C_I\right\},
\]
\[
C_3=\left(1+(1+\mathrm{e})L^2\right)C_3'',\qquad C_5=\mathrm{e}^2L\cdot C_3.
\]
For $c\in\mathrm{me}\,\mathbb{Z}^3$ and an $L$-block $E_1$ let
\[
d_1(c,E_1):=\min\left\{\frac{|c-c'|_1}{L}:c'\in E_1+R\mathbb{Z}^3\right\}\quad(c'\in E_1\text{ if }R=\infty),\qquad d_1(x,E_1):=d_1(c,E_1).
\]
This function of $c$ is $R\mathbb{Z}^3$-periodic, vanishes on $E_1+R\mathbb{Z}^3$, and satisfies $|d_1(c,E_1)-d_1(c',E_1)|\le|c-c'|_1/L$.

\begin{named}{Lemma 6.5}
Let $E_1$ be an $L$-block, let $f\in\mathfrak{S}^{\mathrm{ch}}(\eta_k,R;\theta)$ be bounded and supported over $E_1+R\mathbb{Z}^3$, and let $u:=F_y(0)^{-1}f$.
Then for every site $x$ of $\mathrm{me}\,\mathbb{Z}^3\times O_{\mathrm{cub}}$ and every site $\tilde x$ with $0<d_X(x,\tilde x)\le L$,
\[
|u(x)|_2+L^{\alpha}\,d_X(x,\tilde x)^{-\alpha}\,|u(x)-u(\tilde x)|_2\le C_3\cdot L\exp\left(-a_0\,d_1(x,E_1)\right)|f|_{\mathrm{su}}.
\]
\end{named}

\begin{proof}
Both sides are homogeneous in $f$; let $|f|_{\mathrm{su}}=1$.
Then $f$ lies in $\mathrm{HH}$, and $\|f\|\le|E_1|^{\frac12}=L^{\frac32}$, because a period cell contains one translate of $E_1$.

\emph{Step 0 (weighted $L^2$ bound).}
For $\mathrm{La}>0$ the function $\psi:=a_0\min\left(d_1(\cdot,E_1),\mathrm{La}\right)$ satisfies the hypotheses of Lemma~6.3 and vanishes on the support of $f$.
Hence $\|\exp(\psi)\,u\|\le2C_I\cdot L\,\|f\|\le\mathrm{Uw}:=2C_I\cdot L^{\frac52}$.
Letting $\mathrm{La}\to\infty$ (monotone convergence; subsection~0.3), one obtains
\begin{equation}
\left\|\exp\left(a_0\,d_1(\cdot,E_1)\right)u\right\|\le\mathrm{Uw}.\tag{VI.3}
\end{equation}
The integrand of the norm on the left is $R\mathbb{Z}^3$-periodic.
Therefore, by (b2), every subset $Z$ of an $L$-block of the covering lattice satisfies
\[
\|u\|_{L^2,Z}\le\exp\left(-a_0\min_{c'\in Z}d_1(c',E_1)\right)\mathrm{Uw}.
\]

\emph{Step 1 (sup bound, $k=1$).}
Here $\mathrm{me}=1/L$.
By (n2), and by the consequence of (VI.3) just obtained, taken with $Z=\{c\}$,
\begin{align*}
|u(x)|_2&=48^{\frac12}\,\mathrm{me}^{-\frac32}\,\|u\|_{L^2,\{x\}}\le7L^{\frac32}\,\|u\|_{L^2,\{c\}}\\
&\le7L^{\frac32}\cdot2C_I\cdot L^{\frac52}\exp\left(-a_0\,d_1(x,E_1)\right)=14L^3\cdot C_I\cdot L\exp\left(-a_0\,d_1(x,E_1)\right).
\end{align*}

\emph{Step 2 ($k\ge2$).}
Here $\mathrm{me}\le L^{-2}$.
Let $\mathrm{rad}:=L/4$ and $c^{\mathrm{o}}:=c=\pi(x)$, and let $\mathrm{Bo}$, $\mathrm{Bo}^+$, $\mathrm{Bo}'$ be as in Lemma~6.4.
One has $\mathrm{rad}\ge8\,\mathrm{me}$, because $L^{k+1}\ge125\ge32$, and $\mu\cdot\mathrm{me}\le1$.
On the whole lattice
\[
(D+\mu)u=\mathrm{fs}:=f-\operatorname{Pen}u.
\]
The cube $\mathrm{Bo}^+$ has the side $2(\mathrm{rad}+\mathrm{me})<L$ in each direction; hence it meets at most two $L$-blocks in each direction, and at most $8$ $L$-blocks $E''$ altogether.
A point $c'$ of such a block satisfies $|c'-c|<(\mathrm{rad}+\mathrm{me})+L$, by (b1); so $|c'-c|_1/L<3.75+3\,\mathrm{me}/L\le3.9$, whence $d_1(c',E_1)\ge d_1(x,E_1)-3.9$; moreover $3.9\,a_0\le1.3$.
Step~0 therefore gives
\[
\|u\|_{L^2,E''}\le\mathrm{e}^{1.3}\exp\left(-a_0\,d_1(x,E_1)\right)\mathrm{Uw},\qquad\|u\|_{L^2,\mathrm{Bo}^+}\le8^{\frac12}\,\mathrm{e}^{1.3}\exp\left(-a_0\,d_1(x,E_1)\right)\mathrm{Uw}.
\]
By (Pn3) and the Schwarz inequality, $|\operatorname{Pen}u|_{\mathrm{su},E''}\le2b\cdot L^{\frac32}\,\|u\|_{L^2,E''}$.
The function $f$ vanishes over $\mathrm{Bo}$ unless $E_1+R\mathbb{Z}^3$ meets $\mathrm{Bo}$; in that case $d_1(x,E_1)<3\,\mathrm{rad}/L=3/4$, and $1\le\mathrm{e}^{\frac13}\exp\left(-a_0\,d_1(x,E_1)\right)$.
Consequently
\[
\begin{aligned}
\mathrm{NN}&=\mathrm{rad}\cdot|\mathrm{fs}|_{\mathrm{su},\mathrm{Bo}}+\mathrm{rad}^{-\frac32}\,\|u\|_{L^2,\mathrm{Bo}^+}\\
&\le\left[\frac{L}{4}\,\mathrm{e}^{\frac13}+\frac{L}{4}\cdot2b\cdot L^{\frac32}\,\mathrm{e}^{1.3}\cdot2C_I\cdot L^{\frac52}+8L^{-\frac32}\cdot8^{\frac12}\,\mathrm{e}^{1.3}\cdot2C_I\cdot L^{\frac52}\right]\exp\left(-a_0\,d_1(x,E_1)\right)\\
&\le C_3^{\mathrm{n}}\cdot L\exp\left(-a_0\,d_1(x,E_1)\right),
\end{aligned}
\]
since $\frac{L}{4}\cdot2b\cdot L^{\frac32}\cdot2C_I\cdot L^{\frac52}=L^4\cdot b\cdot C_I\cdot L$ and $8\cdot8^{\frac12}\cdot2=45.25\ldots\le46$.
Lemma~6.4 now yields
\[
|u(x)|_2\le C_R\cdot C_3^{\mathrm{n}}\cdot L\exp\left(-a_0\,d_1(x,E_1)\right),
\]
\[
|u(x)-u(\tilde x)|_2\le C_R\left(\frac{4\,d_X(x,\tilde x)}{L}\right)^{\alpha}C_3^{\mathrm{n}}\cdot L\exp\left(-a_0\,d_1(x,E_1)\right)\qquad\text{if }|\tilde c-c|\le\frac{L}{16}.
\]

\emph{Step 3 (sup bound, all $k$).}
By Steps~1 and~2, for every $k\ge1$ and for every site $x$,
\[
|u(x)|_2\le C_3''\cdot L\exp\left(-a_0\,d_1(x,E_1)\right).
\]

\emph{Step 4 (H\"older quotient).}
Let $\mathrm{hd}:=d_X(x,\tilde x)\in(0,L]$.
If $k\ge2$ and $\mathrm{hd}\le L/16$ (this case contains the pairs inside a cluster, for which $\mathrm{hd}=\mathrm{me}\le L^{-2}$), then $|\tilde c-c|\le L/16$, and Step~2 gives
\[
L^{\alpha}\,\mathrm{hd}^{-\alpha}\,|u(x)-u(\tilde x)|_2\le4^{\alpha}\,C_R\cdot C_3^{\mathrm{n}}\cdot L\exp\left(-a_0\,d_1(x,E_1)\right)\le4\,C_3''\cdot L\exp\left(-a_0\,d_1(x,E_1)\right).
\]
Otherwise either $k=1$, and then $\mathrm{hd}\ge\mathrm{me}=L/L^2$ (for the pairs inside a cluster $\mathrm{hd}=\mathrm{me}$), or $\mathrm{hd}>L/16$; in both cases $(L/\mathrm{hd})^{\alpha}\le\max(L^2,16)=L^2$.
Since $|c-\tilde c|_1\le3\,\mathrm{hd}\le3L$, one has $d_1(\tilde x,E_1)\ge d_1(x,E_1)-3$, and $3a_0\le1$.
Step~3, applied at $x$ and at $\tilde x$, gives
\[
L^{\alpha}\,\mathrm{hd}^{-\alpha}\,|u(x)-u(\tilde x)|_2\le L^2(1+\mathrm{e})\,C_3''\cdot L\exp\left(-a_0\,d_1(x,E_1)\right).
\]
As $4\le(1+\mathrm{e})L^2$, Steps~3 and~4 give the assertion with $C_3=\left(1+(1+\mathrm{e})L^2\right)C_3''$.
\end{proof}

\subsubsection*{2.7. Proof of Proposition 1.6\,(c), (d); the constants}

\begin{proof}[Proof of Proposition~1.6\,(c)]
Let $E_1$ be the $L$-block that contains $\Delta_{z'}$ (see (b1)).
The function $f$ is supported over $\Delta_{z'}+R\mathbb{Z}^3$, a subset of $E_1+R\mathbb{Z}^3$.
Lemma~6.5 therefore applies to $v=F_y(0)^{-1}f$, to $x$ over $\Delta_z$ and to $\tilde x$ with $0<d_X(x,\tilde x)\le1\le L$; as $L^{\alpha}\ge1$, it gives
\[
|v(x)|_2+d_X(x,\tilde x)^{-\alpha}\,|v(x)-v(\tilde x)|_2\le C_3\cdot L\exp\left(-a_0\,d_1(x,E_1)\right)|f|_{\mathrm{su}}.
\]
Let $c''=c^E+Rn$, with $c^E\in E_1$ and $n\in\mathbb{Z}^3$ ($n=0$ if $R=\infty$), be a point with $d_1(x,E_1)=|c-c''|_1/L$.
Then $|z-c|<1/2$, and $|c^E-z'|<L$ by (b1), because $z'$ lies in $\Delta_{z'}$, a subset of $E_1$.
Hence the sup distance of $z$ and $z'$ modulo $R\mathbb{Z}^3$ satisfies
\[
|z-z'|\le|z-c|+|c-c''|+\left|c''-(z'+Rn)\right|<\frac12+|c-c''|+L,
\]
and therefore
\begin{align*}
d_1(x,E_1)&\ge\frac{|c-c''|}{L}>\frac1L\left(|z-z'|-L-\frac12\right),\\
\exp\left(-a_0\,d_1(x,E_1)\right)&\le\exp\left(a_0\left(1+\frac{1}{2L}\right)\right)\exp\left(-\frac{a_0}{L}\,|z-z'|\right)\\
&\le\mathrm{e}^2\exp\left(-\gamma_{(5)}\,|z-z'|\right).
\end{align*}
This is (c), with $C_5=\mathrm{e}^2L\cdot C_3$ and $\gamma_{(5)}=a_0/L$.
\end{proof}

\begin{proof}[Proof of Proposition~1.6\,(d)]
By (a),
\[
\left(F_y(0)^{-1}\right)^*=\left(\gamma_5\,F_{-y}(0)\,\gamma_5\right)^{-1}=\gamma_5\,F_{-y}(0)^{-1}\gamma_5\quad\text{on }\mathrm{HH}(\eta_k,R;\theta),
\]
\[
\left(F_y(0)^{-1}\right)^T=\mathsf{C}^{-1}F_{-y}(0)^{-1}\mathsf{C}\quad\text{on }\mathrm{HH}(\eta_k,R;-\theta).
\]
The matrix $\gamma_5$ is Hermitian with $\gamma_5^2=1$, and $\mathsf{C}$ is real with $\mathsf{C}^T=\mathsf{C}^{-1}$ ((I.c1), (I.c3)); so both are unitary matrices.
They act at each site, preserve $|\cdot|_2$ and supports, and commute with differences of values.
Hence (c), taken for the parameter $-y$ and the character $\theta$, resp.\ $-\theta$, gives the bound of (c) for $\left(F_y(0)^{-1}\right)^*f$, with $f$ as in (c), and for $\left(F_y(0)^{-1}\right)^Tf'$, with $f'\in\mathfrak{S}^{\mathrm{ch}}(\eta_k,R;-\theta)$ bounded and supported over the orbit of $\Delta_{z'}$.
\end{proof}

\begin{proof}[Proof of the assertions of Proposition~1.6 on its constants]
One has $0<a_0\le1/3$, hence $0<\gamma_{(5)}\le1/(3L)$.
The inequality $C_5\ge1$ follows from the formulas of $c_*$, $K_\rho$, $C_{\mathrm{Po}}$ (Section~I.8.2, line~(C0r)) and of $A_1$, $c_b$ (Proposition~1.3\,(b)).
Indeed, $c_*\le w_3/18=1/3$ gives $C_{\mathrm{Po}}\ge1/(8c_*)\ge3/8$; $K_\rho\ge\rho$ gives $A_1\ge2\sqrt2/c_*\ge6\sqrt2$, hence $1/c_b\ge3A_1^2\ge216$ by Proposition~1.3\,(b); so $C_I\ge\left(2C_{\mathrm{Po}}/c_b\right)^{\frac12}\ge1$, and $C_5\ge\mathrm{e}^2L\cdot14L^3\cdot C_I\ge1$.

The numbers $K_1$, $C_I$, $c_{\mathrm{wt}}$, $a_0$, $\gamma_{(5)}$ are given by formulas in $(b,L,\rho)$ and in $C_{\mathrm{Po}}$, $c_b$, which are functions of $\rho$ (Propositions~1.3, 1.4); $C_R$ is given by $C_G$, a function of $\rho$ (Proposition~1.5), and by $\alpha$; and $C_3^{\mathrm{n}}$, $C_3''$, $C_3$, $C_5$ are given by $(b,L,C_I,C_R)$.
None of the Lemmas~6.1--6.5 has a constant depending on $k$, $R$, $\theta$, $\mu\in[0,1]$ or $y\in\mathbb{R}\cup\{\infty\}$.
The level enters only through $\mathrm{me}=L^{-k}\le1/L$ and $b_-\le b_k\le b$; the period $R$ only through $R\in L\mathbb{N}$; the character $\theta$ not at all, since only $|u|_2$ and products $\overline{u'}\cdot u$ are integrated; $\mu$ only through $0\le\mu\le1$; and $y$ only through (al1)--(al5).
Hence $C_5$, $\gamma_{(5)}$, $C_I$ depend on $(b,\alpha,L,\rho)$ only;
in particular, they depend neither on a connection, nor on the charge, nor on the index of the spacing, nor on the side $R$ of the torus.
\end{proof}

\paragraph{Level $0$.}
Proposition~1.6 has no instance at $k=0$: the number $b_0$ and the matrices $\boldsymbol{\alpha}_y$, $\boldsymbol{\beta}_y$ are not defined there (Sections~I.3.4, I.3.9), and there is no short pair.
The operator $F_y(0)$ of level $0$ is the subject of Proposition~1.7.

\subsection*{3. Proposition 1.7}

\subsubsection*{3.1. Proof}

\begin{proof}[Proof of Proposition~1.7]
\emph{Setting.}
In this section $(\cdot,\cdot):=(\cdot,\cdot)^{\mathrm{ch}}$ and $\|\cdot\|:=\|\cdot\|_{L^2}$ on $\mathrm{HH}(\mathrm{me},R;\theta)$, and
\[
\mathrm{HH}_{\mathrm{Wd}}:=\left\{u\in\mathrm{HH}(\mathrm{me},R;\theta):u=0\text{ off }\mathrm{Wd}\times O_{\mathrm{cub}}\right\},
\]
a closed subspace.
The blocks $E$ of level one are the cubes $\left\{c\in\mathrm{me}\,\mathbb{Z}^3:|c-z|<\ell_1/2\right\}$, $z\in\ell_1\mathbb{Z}^3$.
By (me1) with $j=1$ they partition $\mathrm{me}\,\mathbb{Z}^3$; each of them has $L^3$ photon sites, $E\times O_{\mathrm{cub}}$ has $48L^3$ sites, and $|E|=\ell_1^3$.
Hence a block of level one is either contained in $\mathrm{Wd}$ or disjoint from $\mathrm{Wd}$.
The period $R$ is a multiple of $\ell_1=L\cdot\mathrm{me}$, or $R=\infty$.
Fix a period cell of (me2) with $j=1$ (the whole lattice $\mathrm{me}\,\mathbb{Z}^3$ if $R=\infty$).
In pairings and norms, $\sum_E$ denotes the sum over those blocks of level one of this cell which are contained in $\mathrm{Wd}$.
The elements of $\mathrm{HH}_{\mathrm{Wd}}$ vanish over the other blocks of the cell; therefore, by (me2), for $u,u'\in\mathrm{HH}_{\mathrm{Wd}}$,
\[
(u',u)=\sum_E(u',u)_E,\qquad\|u\|^2=\sum_E\|u\|_{L^2,E}^2.
\]
In the definition of $\mathrm{Tst}$, on the other hand, the sum extends over all blocks $E$ of level one contained in $\mathrm{Wd}$, and $P_Eu$ is continued by zero off $E\times O_{\mathrm{cub}}$:
\[
\left(\sum_Eb_E\,\ell_1^{-1}P_Eu\right)(x)=
\begin{cases}
b_E\,\ell_1^{-1}(P_Eu)(x)&\text{for $x$ over a block $E$ contained in $\mathrm{Wd}$},\\
0&\text{for $x$ off $\mathrm{Wd}\times O_{\mathrm{cub}}$}.
\end{cases}
\]

\emph{$\mathrm{Tst}$ maps $\mathrm{HH}_{\mathrm{Wd}}$ into itself and is bounded.}
The operators $\mathbf{1}_{\mathrm{Wd}}$, $D$, $\sum_Eb_E\,\ell_1^{-1}P_E$ and $\gamma_5$ commute with the translations by $Re_\mu$, because $\mathrm{Wd}$ and $E\mapsto b_E$ are $R\mathbb{Z}^3$-periodic; so they preserve the character.
Each of the three terms of $\mathrm{Tst}\,u$ vanishes off $\mathrm{Wd}\times O_{\mathrm{cub}}$.
By (me4), $\|D\|\le\mathrm{me}^{-1}(1+\rho)\,\mathrm{sw}$; the $P_E$ are orthogonal projections of $L^2(E\times O_{\mathrm{cub}};\mathbb{C}^4)$, $0<b_E\le b$, and $|\gamma_5|_2=1$.
Therefore
\[
\|\mathrm{Tst}\|\le\mathrm{me}^{-1}(1+\rho)\,\mathrm{sw}+\mu+\frac{b}{\ell_1}+|y|.
\]

\emph{The real part.}
Let $u\in\mathrm{HH}_{\mathrm{Wd}}$.
\begin{itemize}
\item[(1)] One has $\left(u,\mathbf{1}_{\mathrm{Wd}}(D+\mu)u\right)=\left(\mathbf{1}_{\mathrm{Wd}}u,(D+\mu)u\right)=(u,Du)+\mu\|u\|^2$, and $\mu\|u\|^2\ge0$.
\item[(2)] By Proposition~1.3\,(a), $\operatorname{Re}(u,Du)=\frac{\rho}{2\,\mathrm{me}}\,(u,\nabla^{*}\nabla u)$.
By (me3) with $j=1$, $(u,\nabla^{*}\nabla u)$ is at least the sum of $(u,\nabla^{*}\nabla u)_E$ over all blocks $E$ of level one of the cell.
The blocks of the cell that are not contained in $\mathrm{Wd}$ contribute $(u,\nabla^{*}\nabla u)_E=0$, since $u$ vanishes over them.
Thus $(u,\nabla^{*}\nabla u)\ge\sum_E(u,\nabla^{*}\nabla u)_E$.
(The terms omitted in (me3) belong to the bonds of type~$3$ that join two blocks of level one; the bonds leaving $\mathrm{Wd}$ are among them.)
\item[(3)] Proposition~1.4, applied to the cube $E$ of $L^3$ sites, with $\ell_E=\ell_1$, gives $(u,\nabla^{*}\nabla u)_E\ge\left(C_{\mathrm{Po}}\cdot\ell_1^2\cdot\mathrm{me}^{-2}\right)^{-1}\|u-P_Eu\|_{L^2,E}^2$.
Hence, as $\mathrm{me}/\ell_1=1/L$,
\[
\operatorname{Re}(u,Du)\ge\sum_E\frac{\rho\cdot\mathrm{me}}{2C_{\mathrm{Po}}\cdot\ell_1^2}\,\|u-P_Eu\|_{L^2,E}^2=\sum_E\frac{\rho}{2C_{\mathrm{Po}}\cdot L}\,\ell_1^{-1}\,\|u-P_Eu\|_{L^2,E}^2.
\]
\item[(4)] The function $\sum_Eb_E\,\ell_1^{-1}P_Eu$ lies in $\mathrm{HH}_{\mathrm{Wd}}$ and coincides with $b_E\,\ell_1^{-1}P_Eu$ over each block $E$ contained in $\mathrm{Wd}$.
By the setting,
\[
\left(u,\sum_Eb_E\,\ell_1^{-1}P_Eu\right)=\sum_Eb_E\,\ell_1^{-1}\,(u,P_Eu)_E,
\]
where on the left the sum extends over all blocks contained in $\mathrm{Wd}$, and on the right over those of the cell.
As $P_E$ is an orthogonal projection of $L^2(E\times O_{\mathrm{cub}};\mathbb{C}^4)$, one has $(u,P_Eu)_E=\|P_Eu\|_{L^2,E}^2$; moreover $b_E\ge b_-=b(1-1/L)$.
\item[(5)] The matrix $\gamma_5$ is Hermitian ((I.c1)) and $y$ is real; so $(u,i\gamma_5y\,u)=iy\,(u,\gamma_5u)$ is purely imaginary.
\end{itemize}
Adding (1)--(5), and using $\|u\|_{L^2,E}^2=\|P_Eu\|_{L^2,E}^2+\|u-P_Eu\|_{L^2,E}^2$, one obtains
\[
\operatorname{Re}(u,\mathrm{Tst}\,u)\ge\sum_E\ell_1^{-1}\left[\frac{\rho}{2C_{\mathrm{Po}}\cdot L}\,\|u-P_Eu\|_{L^2,E}^2+b\left(1-\frac1L\right)\|P_Eu\|_{L^2,E}^2\right]\ge c_{\mathrm{re}}\cdot\ell_1^{-1}\sum_E\|u\|_{L^2,E}^2=c_{\mathrm{re}}\cdot\ell_1^{-1}\,(u,u).
\]

\emph{Bijectivity.}
From $\|\mathrm{Tst}\,u\|\,\|u\|\ge\operatorname{Re}(u,\mathrm{Tst}\,u)\ge c_{\mathrm{re}}\cdot\ell_1^{-1}\|u\|^2$ it follows that $\|\mathrm{Tst}\,u\|\ge c_{\mathrm{re}}\cdot\ell_1^{-1}\|u\|$.
So $\mathrm{Tst}$ is injective, and its range is closed in $\mathrm{HH}_{\mathrm{Wd}}$ (by the classical facts of subsection~0.3).
For the adjoint $\mathrm{Tst}^*$ of $\mathrm{Tst}$ in $\mathrm{HH}_{\mathrm{Wd}}$ one has $\operatorname{Re}(u,\mathrm{Tst}^*u)=\operatorname{Re}(\mathrm{Tst}\,u,u)=\operatorname{Re}(u,\mathrm{Tst}\,u)\ge c_{\mathrm{re}}\cdot\ell_1^{-1}\|u\|^2$; so $\mathrm{Tst}^*$ is injective, and the orthogonal complement of the range of $\mathrm{Tst}$ in $\mathrm{HH}_{\mathrm{Wd}}$ is $\{0\}$.
Hence $\mathrm{Tst}$ is a bijection of $\mathrm{HH}_{\mathrm{Wd}}$, and $\|\mathrm{Tst}^{-1}\|_{L^2}\le\ell_1/c_{\mathrm{re}}$.

\emph{Level $0$.}
Let $\mathrm{me}=1$, $\mathrm{Wd}=\mathbb{Z}^3$ and $b_E=b$.
Then $\ell_1=L$, and the blocks of level one are the $L$-blocks $B(y')=\left\{c\in\mathbb{Z}^3:|c-y'|<L/2\right\}$, $y'\in L\mathbb{Z}^3$ (Section~I.1.1).
By (I.3.1) and Section~I.3.3 at $A=0$, where all phases are equal to $1$, one has for $c\in B(y')$
\[
\left(J\,P(0)\,J^{*}f\right)(c,g)=\left(Q^{T}(0)\,Q(0)\,J^{*}f\right)(c)=L^{-3}\sum_{c'\in B(y')}\frac1{48}\sum_{g'\in O_{\mathrm{cub}}}f(c',g')=(P_Ef)(c,g),\qquad E=B(y'),
\]
that is, the mean over the $48L^3$ sites of $B(y')\times O_{\mathrm{cub}}$; on one block its range consists of the constants, which form a space of dimension $4$ in a space of dimension $192L^3$.
Hence $b\,\ell_1^{-1}\sum_EP_E=b_L\,J\,P(0)\,J^{*}$, $\mathbf{1}_{\mathrm{Wd}}=1$, and
\[
\mathrm{Tst}=D_0+\mu+b_L\,J\,P(0)\,J^{*}+i\gamma_5y=F_y(0),
\]
the operator of Section~I.3.9 at $A=0$, on $\mathrm{HH}(1,R;\theta)$; for $\theta=0$ and $R=\ell$ this space is $\mathfrak{S}^{\mathrm{ch}}(1,\ell)$.
This is the assertion of Proposition~1.7 on level $0$.
\end{proof}

\paragraph{Uniformity.}
The constant $c_{\mathrm{re}}$ depends on $(b,L,\rho)$ only, since $C_{\mathrm{Po}}$ is a function of the weights (Proposition~1.4).
In particular it depends neither on $\mathrm{me}$, $R$, $\theta$, $\mathrm{Wd}$, nor on the numbers $b_E\in[b_-,b_+]$, $\mu\ge0$ and $y\in\mathbb{R}$, nor on a connection, the charge or the index of the spacing.

\subsection*{Conditions on the constants}

Propositions~1.5, 1.6 and~1.7 hold for all values of the constants that satisfy the following conditions.

\begin{tabular}{p{0.26\textwidth}p{0.42\textwidth}p{0.24\textwidth}}
\hline
condition & where it is used & where it is imposed\\
\hline
$\rho\in(0,1)$; the standard weights $w$ & everywhere (through Propositions~1.2--1.4; the bound $|M_3(x)|_2=3(1+\rho)\le6$; the inequality $\rho/(1+\rho)<1/2$ in (y2)) & Section~I.8.2, line~(C0); Proposition~14.2\,(a), line~(C0)\\
$b>0$ ($b:=1$) & Lemma~6.1; Proposition~1.7 & Section~I.8.2, line~(C0); Proposition~14.2\,(a), line~(C0)\\
the definitions of $\lambda_{\mathrm{int}}$, $\hat{g}$, $\theta_0$, $\kappa_{(0)}$, $c_{\mathrm{hl}}$, $c_{*}$, $K_{\rho}$, $C_{\mathrm{Po}}$, $c_b$ (with $A_1$) & (y2) ($\lambda_{\mathrm{int}}\le2(w_0+w_1+w_2)$; the formulas of $\hat{g}$ and $\theta_0$; $\kappa_{(0)}\le\theta_0/2$), and through (y2) Section~1; Lemma~5.1 ($c_{\mathrm{hl}}=\sqrt{3}\,(1+\rho)$); the proof of $C_5\ge1$ ($c_{*}$, $K_{\rho}$, $C_{\mathrm{Po}}$, $A_1$, $c_b$); the formulas of $C_I$ ($C_{\mathrm{Po}}$, $c_b$) and of $c_{\mathrm{re}}$ ($C_{\mathrm{Po}}$) & Section~I.8.2, line~(C0r); Proposition~14.2\,(a), line~(C0r)\\
$\alpha\in(0,1)$ & Proposition~1.5\,(f) (there $0\le\alpha\le1$ suffices, as stated in the data of Proposition~1.5); the second bound of \cite[Lemma~12.7(a)]{QEDI} and Lemma~6.4 ($\alpha<1$); Lemma~6.5 ($0\le\alpha\le1$) & Section~I.8.2, line~(C1): $\alpha\in(1-\bar{\varepsilon},1)$\\
$L$ odd, $L\ge5$ & (me1) ($L$ odd; through (me1) in Sections~2 and~3); Lemma~6.5 ($L\ge5$, through $L^3\ge32$, $3\,\mathrm{me}/L\le0.15$ and $16\le L^2$); the consequence $L>1$, through $b_->0$ in Lemmas~6.1 and~6.2 (so that $C_I$ is finite) and through $c_{\mathrm{re}}>0$ in Proposition~1.7 & Proposition~14.2\,(b), condition~(L5); Section~I.8.2, line~(C3)\\
the definitions $b_L=b/L$, $b_-=b(1-1/L)$, $b_+=b$; the relations $b_-\le b_k\le b$ and $b_k\cdot b_L/(b_k+b_L)=b_{k+1}/L$ & Lemma~6.1; Lemma~6.2 and the formula of $C_I$ ($b_-$); the uniformity in $k$ of the constants of Proposition~1.6 ($b_-\le b_k\le b$); Proposition~1.7 ($b_E\le b$, $b_E\ge b_-$, and $b/\ell_1=b_L$ at level~$0$) & Section~I.8.2, line~(C4); Proposition~14.2\,(a), line~(C4); Section~I.3.4\\
the definitions of $K_1$, $C_I$, $c_{\mathrm{wt}}$, $a_0$, $\gamma_{(5)}$, $C_R$, $C_3$, $C_5=\mathrm{e}^2L\cdot C_3$, $c_{\mathrm{re}}$ & Propositions~1.6, 1.7 & Section~I.8.2, line~(C4r); Proposition~14.2\,(a), lines~(C4r)\\
\hline
\end{tabular}

All these conditions are among those imposed in Section~I.8.2 and in Proposition~14.2\,(a), (b), or follow from them and from the definition of the numbers $b_k$ in Section~I.3.4.
Hence Propositions~1.5--1.7 hold for the constants chosen there, and they hold without any condition on $M$, $\varepsilon$, $e_0$ or on a constant chosen later.

The constant $C_G$ of Proposition~1.5, which also enters the lists of constants of Section~I.8.2 and of Proposition~14.2\,(a), is the number
\[
C_G=\max\left\{c(1,3)\,\mathrm{As}_1,\;c(0,4)\,\mathrm{As}_2,\;48\left[c(0,4)\left(2\,\mathrm{As}_3+\mathrm{As}_4\right)+\mathrm{As}_5\right]\right\}
\]
of subsection~1.4.
It is formed with $\hat{g}$, $c_{\mathrm{hl}}$, $\theta_0$, $\kappa_{(0)}$, with the numbers $A_{\mathrm{ak}}$, $A'_{\mathrm{ak}}$ ($|\mathrm{ak}|\le4$) whose existence is asserted by Proposition~1.2\,(S6), (S7), and with the absolute numbers $c(1,3)$, $c(0,4)$ whose existence is asserted by \cite[Lemma~12.4]{QEDI}; hence it is a function of $\rho$.
The constants $C_R$, $C_3$, $C_5$ are given by the formulas of Proposition~1.6, which are also the formulas of Proposition~14.2\,(a);
Lemma~6.4 holds with the constant $C_R$ so defined (subsection~2.5), Lemma~6.5 with $C_3$ (subsection~2.6), and Proposition~1.6\,(c) with $C_5$ (subsection~2.7).
In the order in which Proposition~14.2\,(a) chooses the constants, $C_G$ stands in line~(C0r), after $\theta_0$, $\hat{g}$, $c_{\mathrm{hl}}$, $\kappa_{(0)}$ and the numbers $A_{\mathrm{ak}}$, $A'_{\mathrm{ak}}$ of Proposition~1.2, and $C_R$, $C_3$, $C_5$ stand in line~(C4r), after $\alpha$, $L$, $C_I$.

\section*{Chapter 2a. The elementary operators on block tori; the inverse of $F_{y}$ at zero and at small fields: proofs of Propositions 2.1--2.4}

\subsection*{0. Conventions of this chapter; setting and tools}

The conventions of Part~I and those of Sections~G1.1--G1.4, G1.9 and G1.10 of Part~II are in force.
In this chapter the letter $\eta$ denotes the mesh $\eta_k=L^{-k}$ of the level $k$ (Section~I.1.1), and not the spacing $L^{-N}$ of Section~I.0.5.

Throughout Sections~0--4 the following data are fixed, and $\mu$, $y$, $\gamma$, $A$, $B$ are the quantities of Propositions~2.1--2.4:
an odd integer $L\ge5$;
numbers $b>0$ and $b_{L}=\frac{b}{L}$;
the Wilson parameter $\rho\in(0,1)$ and the standard weights, with $\mathrm{sw}=12+12\sqrt{2}$;
a number $\alpha\in(0,1)$;
a block torus $\mathbb{T}=(t,R,\boldsymbol{\sigma})$ of level $k\ge0$ (Section~G1.2), with the mesh $\eta:=\eta_k=L^{-k}$;
and a number $e_k>0$.
Sites are points of the covering lattices $\mathrm{Lab}$, $\mathrm{Lab}_j$ and $\mathrm{Lab}\times O_{\mathrm{cub}}$, not classes modulo $R\mathbb{Z}^3$, unless classes are named explicitly.
The letter $n$ denotes a vector of $\mathbb{Z}^3$, and $\boldsymbol{\sigma}^n:=\boldsymbol{\sigma}^{n_1+n_2+n_3}$;
$\mathrm{vv}$ denotes a vector of $\mathbb{C}^4$, and $\mathrm{Mx}$ a constant $4\times4$ matrix.
A scalar is an $R\mathbb{Z}^3$-periodic function of $c$ (Section~G1.2).
On $\mathbb{C}^4$ and on the $4\times4$ matrices the norm $|\cdot|_2$ is used throughout;
the sup norms are $|\cdot|_{\mathrm{su}}$ and $\|\cdot\|_{\mathrm{su}}$, and the $L^2$ norms are $\|\cdot\|_{L^2}$ (Section~G1.3).
The letters $R_i$, $E$ and $P_E$ have the meaning of Section~G1.4.
In Section~3\,(b) the symbols $\mathrm{Rw}$, $\mathrm{Rw}^{\mathrm{hop}}$, $\mathrm{Rw}^{\mathrm{pen}}$ denote the operators defined there, not the row sums $\mathrm{Rw}_z(\cdot)$ of Section~G1.3.

\subsubsection*{0.1. Blocks}

For $c\in\mathrm{Lab}$ and $0\le j\le k+1$ let $\mathrm{cen}_j(c)$ be the point $z\in\mathrm{Lab}_j$ with $c\in\mathrm{Bl}_j(z)$;
it exists and is unique, because the blocks of one side form a partition of $\mathrm{Lab}$ (Section~G1.2).
Thus $\mathrm{cen}_0(c)=c$, the unit site of $c$ is $[c]:=\mathrm{cen}_k(c)$, and $\mathrm{cen}_{k+1}(c)$ is the site of the $L$-lattice whose block contains $c$.

\begin{description}
\item[(bl1)]
The block $\mathrm{Bl}_j(z)$ consists of $L^{3j}$ sites.
Hence $\#|\Delta_z|=\eta^{-3}$, $\#|\mathrm{Bl}_{k+1}(z')|=L^3\eta^{-3}$, and $B(z')$ consists of $L^3$ unit sites.
Since $\int\mathrm{d}x=\eta^3\cdot\frac{1}{48}\sum_c\sum_{g\in O_{\mathrm{cub}}}$, it follows that
\[
\int_{\Delta_z}\mathrm{d}x=\eta^3\cdot\frac{1}{48}\cdot\eta^{-3}\cdot48=1,\qquad\int_{\mathrm{Bl}_{k+1}(z')}\mathrm{d}x=L^3.
\]
For $k=0$ one has $\Delta_z=\{z\}$, and $\Delta_z\times O_{\mathrm{cub}}$ is the cluster of $z$, which consists of $48$ sites of weight $\frac{1}{48}$.

\item[(bl2)]
As $\mathrm{Lab}_j\subset\mathrm{Lab}$, the vector $c-\mathrm{cen}_j(c)$ lies in $\eta\mathbb{Z}^3$, and its sup norm is less than $\frac{L^j\eta}{2}$.
Since $L^j$ is odd, it follows that
\[
|c-\mathrm{cen}_j(c)|\le\frac{\left(L^j-1\right)\eta}{2};\qquad\text{in particular}\quad|c-[c]|\le\frac{1-\eta}{2}\qquad(j=k).
\]

\item[(bl3)]
For $0\le j\le k$ one has
\[
|\mathrm{cen}_{j+1}(c)-\mathrm{cen}_j(c)|\le\frac{\left(L^{j+1}-L^j\right)\eta}{2}.
\]
Indeed, by the nesting of the blocks, $\mathrm{Bl}_j(\mathrm{cen}_j(c))\subset\mathrm{Bl}_{j+1}(\mathrm{cen}_{j+1}(c))$.
Let $\mu$ be a coordinate in which the sup norm $|\mathrm{cen}_j(c)-\mathrm{cen}_{j+1}(c)|$ is attained.
The site of $\mathrm{Bl}_j(\mathrm{cen}_j(c))$ whose coordinate $\mu$ is farthest from that of $\mathrm{cen}_{j+1}(c)$ differs from $\mathrm{cen}_{j+1}(c)$ in this coordinate by $|\mathrm{cen}_j(c)-\mathrm{cen}_{j+1}(c)|+\frac{\left(L^j-1\right)\eta}{2}$;
as this site lies in $\mathrm{Bl}_{j+1}(\mathrm{cen}_{j+1}(c))$, the number is at most $\frac{\left(L^{j+1}-1\right)\eta}{2}$ by (bl2).

\item[(bl4)]
Two sites $c$, $c'$ of one block $\mathrm{Bl}_j(z)$ satisfy $|c-c'|\le\left(L^j-1\right)\eta$ and $|c-c'|_1\le3\left(L^j-1\right)\eta$.
Consequently $|c-c'|<1$ and $|c-c'|_1<3$ for two sites of one unit block,
$|c-c'|<L$ and $|c-c'|_1<3L$ for two sites of one block of side $L$,
and $|z_1-z_2|\le L-1$ for two unit sites $z_1$, $z_2$ of one set $B(z')$, because $z_1-z_2\in\mathbb{Z}^3$.

\item[(bl5)]
If $c'=c+\eta a$ with $a\in\{\pm e_\mu\}$, then $[c']=[c]$ or $[c']=[c]+a$.
Indeed, $[c']-[c]$ lies in $\mathbb{Z}^3$.
By (bl2) each coordinate $\nu\ne\mu$ of this vector has modulus at most $\frac{1-\eta}{2}+\frac{1-\eta}{2}<1$, hence vanishes;
its coordinate $\mu$ differs from $\eta a_\mu$ by at most $1-\eta$, hence lies in $\{0,a_\mu\}$.

\item[(bl6)]
If $|c-c'|<1$, then $|[c]-[c']|\le1$: this sup norm is less than $\frac12+1+\frac12$, and $[c]-[c']\in\mathbb{Z}^3$.

\item[(bl7)]
One has $\mathrm{cen}_j(c+Rn)=\mathrm{cen}_j(c)+Rn$, and $d(z,z+a)\le|a|$ for $a\in\mathbb{Z}^3$.

\item[(bl8)]
The path $\boldsymbol{\Gamma}^{\pi}(c,c')$ consists of $\frac{|c-c'|_1}{\eta}\le\frac{3\,|c-c'|}{\eta}$ bonds, all of which lie in the lattice box spanned by $c$ and $c'$.
A block is a lattice box.
Therefore every bond of a path $\boldsymbol{\Gamma}^{\pi}$ joining two sites of a block is contained in this block.

\item[(bl9)]
(Period cells.)
The block $\mathrm{Bl}_{k+1}(z')$ is the disjoint union of the unit blocks $\Delta_z$, $z\in B(z')$.
Indeed, the unit blocks form a partition of $\mathrm{Lab}$, each of them is contained in one block of side $L$ (Section~G1.2), and $z\in\Delta_z$;
hence $\Delta_z\subset\mathrm{Bl}_{k+1}(z')$ if and only if $z\in\mathrm{Bl}_{k+1}(z')$.

Fix a period cell of $\mathrm{Lab}_{k+1}$.
The union of the blocks $\mathrm{Bl}_{k+1}(z')$, $z'$ in this cell, is a period cell of $\mathrm{Lab}$:
for $c\in\mathrm{Lab}$ and $n\in\mathbb{Z}^3$ the site $c+Rn$ lies in this union if and only if $\mathrm{cen}_{k+1}(c+Rn)=\mathrm{cen}_{k+1}(c)+Rn$ (by (bl7)) lies in the cell of $\mathrm{Lab}_{k+1}$, and this is the case for exactly one $n$.
Since $\mathrm{Lab}^0$ is invariant under $R\mathbb{Z}^3$ (Section~G1.2), the intersection of this period cell of $\mathrm{Lab}$ with $\mathrm{Lab}^0$, which is the union of the sets $B(z')=\mathrm{Lab}^0\cap\mathrm{Bl}_{k+1}(z')$, is a period cell of $\mathrm{Lab}^0$;
and the period cell of $\mathrm{Lab}$ is the union of the blocks $\Delta_z$, $z$ in this period cell of $\mathrm{Lab}^0$.
For these three cells, which have $\left(\frac{R}{L}\right)^3$, $R^3$ and $R^3\eta^{-3}$ sites, one has
\[
\sum_c=\sum_z\sum_{c\in\Delta_z}=\sum_{z'}\sum_{c\in\mathrm{Bl}_{k+1}(z')},\qquad\sum_z=\sum_{z'}\sum_{z\in B(z')}.
\]
For $k=0$: $\mathrm{Lab}^0=\mathrm{Lab}$, $\Delta_z=\{z\}$ and $B(z')=\mathrm{Bl}_1(z')$.
\end{description}

\subsubsection*{0.2. Spinors, operators, kernels}

Let a linear operator be defined on all functions on $\mathrm{Lab}\times O_{\mathrm{cub}}$ (on $\mathrm{Lab}^0$, on $\mathrm{Lab}_{k+1}$) by a formula that commutes with the translations by $R\mathbb{Z}^3$.
Then it maps spinors to spinors.
All operators of Section~0.4 are of this kind, because $A$ is $R\mathbb{Z}^3$-periodic, (bl7) holds, $x\mapsto x^{(i)}$ commutes with the translations (I.s1), and $M_i(x)$ depends on the label only.

For $x'\in\mathrm{Lab}\times O_{\mathrm{cub}}$ and $\mathrm{vv}\in\mathbb{C}^4$ let $\mathrm{dsp}[x',\mathrm{vv}]$ be the spinor on $\mathbb{T}$ with
\[
\mathrm{dsp}[x',\mathrm{vv}](x'+Rn)=\boldsymbol{\sigma}^n\cdot48\,\eta^{-3}\cdot\mathrm{vv}\qquad(n\in\mathbb{Z}^3)
\]
that vanishes off the orbit of $x'$.
Every spinor $f$ on $\mathbb{T}$ can be written as
\[
f=\sum_{x'}\frac{\eta^3}{48}\cdot\mathrm{dsp}[x',f(x')],
\]
where $x'$ ranges over one period cell.
Accordingly, for a linear operator $O$ on the spinors on $\mathbb{T}$ with values in any one of the spaces of Section~G1.3 one has the following.

\begin{description}
\item[(kn1)]
Put $O(t,x')\,\mathrm{vv}:=\left(O\,\mathrm{dsp}[x',\mathrm{vv}]\right)(t)$.
This is a kernel of $O$ in the sense of Section~G1.2, and it is the only one: for every period cell,
\[
(Of)(t)=\eta^3\cdot\frac{1}{48}\sum_{x'}O(t,x')\cdot f(x')=\int\mathrm{d}x'\,O(t,x')\cdot f(x').
\]
The relation $\mathrm{dsp}[x'+Re_\mu,\mathrm{vv}]=\boldsymbol{\sigma}\cdot\mathrm{dsp}[x',\mathrm{vv}]$ gives $O(t,x'+Re_\mu)=\boldsymbol{\sigma}\cdot O(t,x')$;
and since the column $t\mapsto O(t,x')\,\mathrm{vv}$ is an element of the target space, $O(t+Re_\mu,x')=\boldsymbol{\sigma}\cdot O(t,x')$.

\item[(kn2)]
For scalars $\psi$, $\psi'$ one has $(\psi\,O\,\psi')(t,x')=\psi(t)\cdot O(t,x')\cdot\psi'(c')$, because $\psi'\,\mathrm{dsp}[x',\mathrm{vv}]=\psi'(c')\cdot\mathrm{dsp}[x',\mathrm{vv}]$.

\item[(kn3)]
$(O_1O_2)(t,x')\,\mathrm{vv}=\left(O_1\left[O_2(\cdot,x')\,\mathrm{vv}\right]\right)(t)$.

\item[(kn4)]
(This is (Nk5) of Section~G1.3; level $0$.)
For an operator $O$ on the spinors on a block torus of level $0$,
\[
|O(x,x')|_2\le48\,\|O\|_{L^2}.
\]
Indeed, $\eta=1$.
The spinor $\mathrm{dsp}[x',\mathrm{vv}]$ is carried by one site of a period cell, so that $\|\mathrm{dsp}[x',\mathrm{vv}]\|_{L^2}^2=\frac{1}{48}\,|48\,\mathrm{vv}|_2^2=48\,|\mathrm{vv}|_2^2$.
For every spinor $u$ and every site $x$ one has $\frac{1}{48}\,|u(x)|_2^2\le\|u\|_{L^2}^2$.
Hence
\[
|O(x,x')\,\mathrm{vv}|_2=\left|\left(O\,\mathrm{dsp}[x',\mathrm{vv}]\right)(x)\right|_2\le\sqrt{48}\,\|O\|_{L^2}\cdot\sqrt{48}\,|\mathrm{vv}|_2.
\]

\item[(kn5)]
For an $R\mathbb{Z}^3$-periodic function $\mathrm{Fn}$ on $\mathrm{Lab}\times O_{\mathrm{cub}}$ and $i=0,\dots,3$,
\[
\int\mathrm{d}x\,\mathrm{Fn}\left(x^{(i)}\right)=\int\mathrm{d}x\,\mathrm{Fn}(x).
\]
Indeed, $x\mapsto x^{(i)}$ is a bijection of $\mathrm{Lab}\times O_{\mathrm{cub}}$ that commutes with the translations (I.s1);
hence it permutes the classes modulo $R\mathbb{Z}^3$, and all sites have the weight $\frac{\eta^3}{48}$.
The products $f'(x)\cdot f(x)$ and $\overline{f'(x)}\cdot f(x)$ of two spinors are periodic, because $\boldsymbol{\sigma}^2=1$.
\end{description}

\subsubsection*{0.3. Phases}

For a connection $A$ on $\mathbb{T}$ the phases of (I.1.6) and (I.3.3) take the following form in the notation of Section~G1.2:
\begin{align*}
u_A(c,c')&=\exp\left\{i\cdot e_k\cdot\eta\cdot A(c,c')\right\}\qquad\text{on bonds},\\
U_A(c,c')&=\exp\left\{i\cdot e_k\cdot\eta\cdot(\tau A)(c,c')\right\},\qquad U_A(c,c)=1,\\
(\tau_kA)(z,c)&=\sum_{j<k}(\tau A)\left(\mathrm{cen}_{j+1}(c),\mathrm{cen}_j(c)\right)\qquad\text{for }c\in\Delta_z,\\
\varphi_A(z,c)&=\exp\left\{i\cdot e_k\cdot\eta\cdot(\tau_kA)(z,c)\right\}\qquad\text{for }c\in\Delta_z,\\
\varphi^{+}_A(z',c)&=U_A(z',[c])\cdot\varphi_A([c],c)\qquad\text{for }c\in\mathrm{Bl}_{k+1}(z').
\end{align*}
For $k=0$ one has $\tau_0A=0$, $\varphi_A(z,z)=1$, $\varphi^{+}_A(z',c)=U_A(z',c)$ and $\mathrm{Bl}_1(z')=B(z')$.

\begin{description}
\item[(ph0)]
Each of the four phases has the form $\exp\left\{i\cdot e_k\cdot\eta\cdot\mathrm{lf}(A)\right\}$, where $\mathrm{lf}$ is a linear functional of $A$ with real coefficients.
If $\mathrm{pf}_A$ denotes any one of them, then
\begin{gather*}
\mathrm{pf}_{A+B}=\mathrm{pf}_A\cdot\mathrm{pf}_B,\qquad\mathrm{pf}_0=1,\qquad\mathrm{pf}_{-A}=\mathrm{pf}_A^{-1},\\
\overline{\mathrm{pf}_{\overline{A}}}=\mathrm{pf}_{-A},\qquad|\mathrm{pf}_A|=1\ \text{for real }A,\qquad u_A(c',c)=u_{-A}(c,c').
\end{gather*}
The phases do not change when all their arguments are translated by $Rn$.
By (bl8) and the nesting of the blocks, $\varphi_A(z,c)$ depends on $A$ only through the values $A(b)$ at bonds $b$ contained in $\Delta_z$,
and $U_A(z',z)$ ($z\in B(z')$) and $\varphi^{+}_A(z',c)$ depend on $A$ only through the values $A(b)$ at bonds $b$ contained in $\mathrm{Bl}_{k+1}(z')$.
\end{description}

\begin{named}{Lemma 0.1}
Let $z\in\mathrm{Lab}^0$ and $z'\in\mathrm{Lab}_{k+1}$, and let $c$, $c'$ be sites at which the respective phases are defined (Section~0.3).
\begin{itemize}
\item[(a)]
One has
\begin{gather*}
\left|u_A(c,c')^{\pm1}\right|\le\exp\left(\eta\cdot\mathfrak{i}(A)\right),\qquad\left|\varphi_A(z,c)^{\pm1}\right|\le\exp\left(\frac{3}{2}\,\mathfrak{i}(A)\right),\\
\left|\varphi^{+}_A(z',c)^{\pm1}\right|\le\exp\left(\frac{3}{2}\,L\cdot\mathfrak{i}(A)\right),\qquad\left|U_A(c,c')^{\pm1}\right|\le\exp\left(3\,|c-c'|\cdot\mathfrak{i}(A)\right).
\end{gather*}
\item[(b)]
If $e_k\cdot|B|_\infty\le\delta$, then
\begin{gather*}
\left|u_B(c,c')-1\right|\le\eta\cdot\delta\cdot\exp(\eta\cdot\delta),\qquad\left|\varphi_B(z,c)^{-1}\cdot\varphi_B(z,c')-1\right|\le3\delta\cdot\exp(3\delta),\\
\left|\varphi^{+}_B(z',c)^{-1}\cdot\varphi^{+}_B(z',c')-1\right|\le3L\cdot\delta\cdot\exp(3L\cdot\delta),\\
\left|U_B(c,c')-1\right|\le3\,|c-c'|\cdot\delta\cdot\exp\left(3\,|c-c'|\cdot\delta\right).
\end{gather*}
\end{itemize}
\end{named}

\begin{proof}
Let $B$ be a connection.
By (bl8), $\eta\cdot\left|B\left(\boldsymbol{\Gamma}^{\pi}(c,c')\right)\right|\le3\,|c-c'|\cdot|B|_\infty$, and the mean $(\tau B)(c,c')$ over the six orders $\pi$ obeys the same bound.
By (bl3),
\[
\eta\cdot\left|(\tau_kB)(z,c)\right|\le3\sum_{j<k}\left(L^{j+1}-L^j\right)\frac{\eta}{2}\cdot|B|_\infty=\frac{3}{2}\left(L^k-1\right)\eta\cdot|B|_\infty\le\frac{3}{2}\,|B|_\infty.
\]
For $z\in B(z')$ one has $z-z'\in\mathbb{Z}^3$ and $|z-z'|<\frac{L}{2}$, hence $|z-z'|\le\frac{L-1}{2}$, and therefore
\[
\eta\cdot\left|(\tau B)(z',z)\right|\le\frac{3}{2}\,(L-1)\cdot|B|_\infty.
\]
The sum of the last two bounds is $\frac{3}{2}\,L\cdot|B|_\infty$.
Consequently the functionals $\mathrm{lf}$ of (ph0) satisfy
\[
\eta\cdot|\mathrm{lf}(B)|\le\eta\cdot|B|_\infty,\qquad\frac{3}{2}\,|B|_\infty,\qquad\frac{3}{2}\,L\cdot|B|_\infty,\qquad3\,|c-c'|\cdot|B|_\infty
\]
for $u$, $\varphi$, $\varphi^{+}$ and $U$, respectively.

(a) For a complex number $\mathrm{cx}$ one has $\left|\exp(i\cdot\mathrm{cx})^{\pm1}\right|\le\exp|\operatorname{Im}\mathrm{cx}|$, and $\operatorname{Im}\mathrm{lf}(A)=\mathrm{lf}(\operatorname{Im}A)$;
the bounds above, applied to the connection $\operatorname{Im}A$, give (a).

(b) For a complex number $\mathrm{cx}$ one has $|\exp(\mathrm{cx})-1|\le|\mathrm{cx}|\cdot\exp|\mathrm{cx}|$.
The two quotients in (b) are $\exp\left\{i\cdot e_k\cdot\eta\cdot\left(\mathrm{lf}'(B)-\mathrm{lf}(B)\right)\right\}$ with two functionals $\mathrm{lf}$, $\mathrm{lf}'$ of the same kind;
their exponents are therefore bounded in modulus by twice the bounds above, that is, by $3\delta$ and by $3L\cdot\delta$.
\end{proof}

\begin{named}{Remark}
Lemma~0.1 is \cite[Lemma~16.2]{QEDI}.
There the charge ranges over $e_k\in(0,e_0]$; here every $e_k>0$ is admitted.
\end{named}

\subsubsection*{0.4. The elementary operators in explicit form}

Let $x=(c,g)$, and put
\[
c_3=c+\eta\cdot ge_1,\qquad x^{(i)}=(c,gs_i)\quad(i\le2),\qquad x^{(3)}=(c_3,gs_3),\qquad z=[c],\qquad z'=\mathrm{cen}_{k+1}(c).
\]
The hopping matrices (I.2.1) are
\[
M_i(x)=\frac{w_i}{2}\left(\rho+\gamma\cdot(g\hat n_i)\right),\qquad M_i\left(x^{(i)}\right)=\frac{w_i}{2}\left(\rho-\gamma\cdot(g\hat n_i)\right),
\]
since the label of $x^{(i)}$ is $gs_i$ and $s_i\hat n_i=-\hat n_i$; the vector $g\hat n_i$ is a real unit vector.
The following five facts are (I.m1), (I.m2), (I.m3), (I.m5) and (I.m6); the reason is indicated with each.

\begin{description}
\item[(hm1)]
$\sum_{i=0}^{3}M_i(x)=\frac{\rho}{2}\,\mathrm{sw}$, because $\sum_iw_i\hat n_i=0$ (I.1.1).
\item[(hm2)]
$M_i\left(x^{(i)}\right)^T=\mathsf{C}^{-1}\cdot M_i(x)\cdot\mathsf{C}$, because $\gamma_\mu^T=-\mathsf{C}^{-1}\cdot\gamma_\mu\cdot\mathsf{C}$ (I.c3).
\item[(hm3)]
$\overline{M_i(x)}=\varkappa^{-1}\cdot M_i(x)\cdot\varkappa$, because $\overline{\gamma_\mu}=\varkappa^{-1}\cdot\gamma_\mu\cdot\varkappa$ (I.c4).
\item[(hm4)]
$M_i(x)^*=M_i(x)$ and $\gamma_5\cdot M_i(x)\cdot\gamma_5=M_i\left(x^{(i)}\right)$, because $\gamma_\mu$ is Hermitian and $\gamma_5\cdot\gamma_\mu\cdot\gamma_5=-\gamma_\mu$ (I.c1).
\item[(hm5)]
$|M_i(x)|_2=\frac{w_i}{2}\,(1+\rho)$, and $|M_3(x)|_2=3\,(1+\rho)<6$.
Indeed, $\gamma\cdot(g\hat n_i)$ is Hermitian with square $1$ (I.c1); hence $M_i(x)$ is Hermitian with the eigenvalues $\frac{w_i}{2}\,(\rho\pm1)$.
\end{description}

\begin{description}
\item[(op1)]
By (I.2.3) and (hm1),
\begin{align*}
(D_Af)(c,g)&=-\eta^{-1}\left\{\sum_{i=0}^{2}M_i(x)\left[f(c,gs_i)-f(c,g)\right]+M_3(x)\left[u_A(c,c_3)\cdot f(c_3,gs_3)-f(c,g)\right]\right\}\\
&=\eta^{-1}\,\frac{\rho}{2}\,\mathrm{sw}\cdot f(x)-\eta^{-1}\sum_{i=0}^{3}M_i(x)\cdot U_A\left(x,x^{(i)}\right)\cdot f\left(x^{(i)}\right).
\end{align*}

\item[(op2)]
The sites $x^{(0)}$, $x^{(1)}$, $x^{(2)}$ lie over $c$.
Hence, for connections $A$, $B$, a function $\psi$ of $c$ and a function $\mathrm{uf}$ of $c$ without zeros,
\begin{align*}
\left((D_A-D_B)f\right)(x)&=-\eta^{-1}M_3(x)\cdot(u_A-u_B)(c,c_3)\cdot f\left(x^{(3)}\right),\\
\left([D_A,\psi]f\right)(x)&=-\eta^{-1}M_3(x)\cdot u_A(c,c_3)\cdot\left(\psi(c_3)-\psi(c)\right)\cdot f\left(x^{(3)}\right),\\
\left(\left(\mathrm{uf}\,D_A\,\mathrm{uf}^{-1}-D_A\right)f\right)(x)&=-\eta^{-1}M_3(x)\cdot u_A(c,c_3)\cdot\left(\frac{\mathrm{uf}(c)}{\mathrm{uf}(c_3)}-1\right)\cdot f\left(x^{(3)}\right).
\end{align*}
The first two identities are (I.O6).

\item[(op3)]
For $k\ge1$,
\begin{gather*}
\left(Q_k(A)J^{*}f\right)(z)=\eta^3\sum_{c\in\Delta_z}\varphi_A(z,c)\cdot\frac{1}{48}\sum_{g\in O_{\mathrm{cub}}}f(c,g)=\int_{\Delta_z}\mathrm{d}x\,\varphi_A(z,c)\cdot f(x),\\
\left(JQ_k^{T}(-A)v\right)(c,g)=\varphi_A([c],c)^{-1}\cdot v([c])\qquad\text{for each of the }48\text{ labels }g.
\end{gather*}

\item[(op4)]
One has
\[
\left(Q(A)v\right)(z')=L^{-3}\sum_{z\in B(z')}U_A(z',z)\cdot v(z),\qquad\left(Q^{T}(-A)v'\right)(z)=U_A(z',z)^{-1}\cdot v'(z')\quad\left(z\in B(z')\right).
\]

\item[(op5)]
For $k\ge1$,
\[
\left(JP_k(A)J^{*}f\right)(x)=\varphi_A(z,c)^{-1}\int_{\Delta_z}\mathrm{d}\tilde x\,\varphi_A(z,\tilde c)\cdot f(\tilde x).
\]

\item[(op6)]
For $k\ge1$,
\begin{align*}
\left(Q(A)Q_k(A)J^{*}f\right)(z')&=L^{-3}\int_{\mathrm{Bl}_{k+1}(z')}\mathrm{d}\tilde x\,\varphi^{+}_A(z',\tilde c)\cdot f(\tilde x),\\
\left(JQ_k^{T}(-A)Q^{T}(-A)v'\right)(x)&=\varphi^{+}_A(z',c)^{-1}\cdot v'(z'),\\
\left(JP_{k+1}(A)J^{*}f\right)(x)&=\varphi^{+}_A(z',c)^{-1}\cdot L^{-3}\int_{\mathrm{Bl}_{k+1}(z')}\mathrm{d}\tilde x\,\varphi^{+}_A(z',\tilde c)\cdot f(\tilde x).
\end{align*}
These follow from (op3) and (op4):
$\mathrm{Bl}_{k+1}(z')$ is the disjoint union of the blocks $\Delta_z$, $z\in B(z')$ (bl9);
$\varphi^{+}_A(z',\tilde c)=U_A(z',[\tilde c])\cdot\varphi_A([\tilde c],\tilde c)$;
and $P_{k+1}(A)=Q_k^{T}(-A)\,Q^{T}(-A)\,Q(A)\,Q_k(A)$.

\item[(op7)]
For $k=0$,
\begin{align*}
\left(Q(A)J^{*}f\right)(z')&=L^{-3}\sum_{c\in B(z')}U_A(z',c)\cdot\frac{1}{48}\sum_{g\in O_{\mathrm{cub}}}f(c,g)=L^{-3}\int_{B(z')}\mathrm{d}x\,U_A(z',c)\cdot f(x),\\
\left(JQ^{T}(-A)v'\right)(c,g)&=U_A(z',c)^{-1}\cdot v'(z'),\\
\left(JP(A)J^{*}f\right)(x)&=U_A(z',c)^{-1}\cdot L^{-3}\int_{B(z')}\mathrm{d}\tilde x\,U_A(z',\tilde c)\cdot f(\tilde x).
\end{align*}
These are the displays of (op6) at $k=0$ (Section~0.3); moreover $Q_0=P_0=1$ ((I.3.4), Section~I.3.9).

\item[(op8)]
The operator $F_y(A)$ is
\begin{align*}
F_y(A)&=D_A+\mu+\boldsymbol{\alpha}_{y}\cdot JP_k(A)J^{*}+\boldsymbol{\beta}_{y}\cdot JP_{k+1}(A)J^{*}\qquad\left(k\ge1,\ y\in\mathbb{R}\cup\{\infty\}\right),\\
F_y(A)&=D_A+\mu+b_{L}\cdot JP(A)J^{*}+i\gamma_5y\qquad(k=0,\ y\in\mathbb{R}).
\end{align*}

\item[(op9)]
(The kernel of $F_y(A)$.)
The set $\mathrm{Bl}_{k+1}(z')\cup\{c_3\}$ has sup-diameter at most $L<R$; hence it lies in a period cell.
For $\tilde x$ over this cell, with the kernel $48\,\eta^{-3}\cdot\mathbf{1}[\tilde x=x]\cdot\mathbf{1}_4$ of the identity (I.ke1), one has for $k\ge1$
\begin{align*}
F_y(A)(x,\tilde x)={}&48\,\eta^{-3}\left\{\left(\eta^{-1}\,\frac{\rho}{2}\,\mathrm{sw}+\mu\right)\mathbf{1}[\tilde x=x]-\eta^{-1}\sum_{i=0}^{3}M_i(x)\cdot U_A\left(x,x^{(i)}\right)\cdot\mathbf{1}\left[\tilde x=x^{(i)}\right]\right\}\\
&+\boldsymbol{\alpha}_{y}\cdot\varphi_A(z,c)^{-1}\cdot\varphi_A(z,\tilde c)\cdot\mathbf{1}[\tilde c\in\Delta_z]\\
&+\boldsymbol{\beta}_{y}\cdot L^{-3}\cdot\varphi^{+}_A(z',c)^{-1}\cdot\varphi^{+}_A(z',\tilde c)\cdot\mathbf{1}\left[\tilde c\in\mathrm{Bl}_{k+1}(z')\right],
\end{align*}
and for $k=0$
\begin{align*}
F_y(A)(x,\tilde x)={}&48\left\{\left(\frac{\rho}{2}\,\mathrm{sw}+\mu+i\gamma_5y\right)\mathbf{1}[\tilde x=x]-\sum_{i=0}^{3}M_i(x)\cdot U_A\left(x,x^{(i)}\right)\cdot\mathbf{1}\left[\tilde x=x^{(i)}\right]\right\}\\
&+b_{L}\cdot L^{-3}\cdot U_A(z',c)^{-1}\cdot U_A(z',\tilde c)\cdot\mathbf{1}\left[\tilde c\in B(z')\right].
\end{align*}
Consequently $F_y(A)(x,\tilde x)=0$ unless, for some $n$, the site $\tilde x+Rn$ is one of $x$, $x^{(0)}$, $x^{(1)}$, $x^{(2)}$, $x^{(3)}$, or $c$ and $\tilde c+Rn$ lie in one block of side $L$.
\end{description}

\subsubsection*{0.5. The matrices of the penalty}

Let $\mathrm{Mx}=\mathrm{Mx}'+i\,\mathrm{Mx}''\cdot\gamma_5$ with real numbers $\mathrm{Mx}'$, $\mathrm{Mx}''$, and put $\mathrm{Mx}^{\mathrm{fl}}:=\mathrm{Mx}'-i\,\mathrm{Mx}''\cdot\gamma_5$.
From (I.c1), (I.c3), (I.c4) and Section~I.4.1 one obtains the following.

\begin{description}
\item[(mx1)]
$\mathrm{Mx}^T=\mathsf{C}^{-1}\cdot\mathrm{Mx}^{\mathrm{fl}}\cdot\mathsf{C}$, because $\gamma_5^T=-\mathsf{C}^{-1}\cdot\gamma_5\cdot\mathsf{C}$.
\item[(mx2)]
$\overline{\mathrm{Mx}}=\varkappa^{-1}\cdot\mathrm{Mx}\cdot\varkappa$, because $\overline{i\gamma_5}=-i\,\overline{\gamma_5}=i\,\varkappa^{-1}\cdot\gamma_5\cdot\varkappa$.
\item[(mx3)]
$\mathrm{Mx}^*=\mathrm{Mx}^{\mathrm{fl}}=\gamma_5\cdot\mathrm{Mx}^{\mathrm{fl}}\cdot\gamma_5$, because $\gamma_5$ is Hermitian and $\gamma_5^2=1$.
\item[(mx4)]
$S\cdot\mathrm{Mx}\cdot S^{-1}=\mathrm{Mx}$ for $S\in\mathrm{Sp}_3$, and $S\cdot\mathrm{Mx}\cdot S^{-1}=\mathrm{Mx}^{\mathrm{fl}}$ for $S\in\mathrm{Sp}\setminus\mathrm{Sp}_3$, because such an $S$ anticommutes with $\gamma_5$.
\item[(mx5)]
$|\mathrm{Mx}|_2=\left((\mathrm{Mx}')^2+(\mathrm{Mx}'')^2\right)^{\frac{1}{2}}$, because $\mathrm{Mx}$ is normal with the eigenvalues $\mathrm{Mx}'\pm i\,\mathrm{Mx}''$.
\item[(mx6)]
For $k\ge1$ and real $y$,
\[
\left(b_k+i\gamma_5y\right)^{-1}=\frac{b_k-i\gamma_5y}{b_k^2+y^2},
\]
a matrix of norm $\left(b_k^2+y^2\right)^{-\frac{1}{2}}\le\frac{1}{b_k}$ by (mx5); the same holds with $b_k+b_{L}$ in place of $b_k$.
Hence the matrices of Section~I.3.5 are
\begin{align*}
\boldsymbol{\alpha}_{y}&=\frac{b_k\cdot y^2+i\,b_k^2\cdot y\cdot\gamma_5}{b_k^2+y^2},\\
\boldsymbol{\beta}_{y}&=\frac{b_k^2\cdot b_{L}\cdot\left[(b_k+b_{L})\cdot b_k-y^2-i\,(2b_k+b_{L})\cdot y\cdot\gamma_5\right]}{\left[(b_k+b_{L})^2+y^2\right]\cdot\left[b_k^2+y^2\right]}.
\end{align*}
Thus $\boldsymbol{\alpha}_{y}$, $\boldsymbol{\beta}_{y}$ ($k\ge1$) and $i\gamma_5y$ ($k=0$) have the form $\mathrm{Mx}'+i\,\mathrm{Mx}''\cdot\gamma_5$, and
\[
(\boldsymbol{\alpha}_{y})^{\mathrm{fl}}=\boldsymbol{\alpha}_{-y},\qquad(\boldsymbol{\beta}_{y})^{\mathrm{fl}}=\boldsymbol{\beta}_{-y},\qquad(i\gamma_5y)^{\mathrm{fl}}=i\gamma_5\cdot(-y);
\]
moreover $\boldsymbol{\alpha}_{\infty}=b_k$ and $\boldsymbol{\beta}_{\infty}=0$ are real numbers, and $-\infty:=\infty$.
If $y$ is neutral (Section~G1.9), then $\mathrm{Mx}''=0$ and $\mathrm{Mx}^{\mathrm{fl}}=\mathrm{Mx}$ for these matrices.
\item[(mx7)]
(These are the bounds of \cite[Section~16.2.5]{QEDI}.)
For all $y\in\mathbb{R}\cup\{\infty\}$,
\[
|\boldsymbol{\alpha}_{y}|_2\le b_k\le b,\qquad|\boldsymbol{\beta}_{y}|_2\le\frac{b_k\cdot b_{L}}{b_k+b_{L}}\le b_{L}.
\]
Indeed, for real $y$ one has $|\boldsymbol{\alpha}_{y}|_2=b_k\cdot|y|\cdot\left(b_k^2+y^2\right)^{-\frac{1}{2}}$ and $|\boldsymbol{\beta}_{y}|_2\le\frac{b_k^2\cdot b_{L}}{(b_k+b_{L})\cdot b_k}$, by (mx5) and (mx6);
for $y=\infty$ the bounds follow from $\boldsymbol{\alpha}_{\infty}=b_k$ and $\boldsymbol{\beta}_{\infty}=0$;
and $b_k=\frac{b\left(1-L^{-1}\right)}{1-L^{-k}}\le b$.
\end{description}

\subsubsection*{0.6. The norms $N_{\gamma}$ and $N^{\mathrm{ke}}_{\gamma}$ of Section~G1.3: proofs of their properties}

In each of the three spaces of Section~G1.3 (spinors on $\mathbb{T}$, spinors on the unit lattice, functions of short pairs) one has, with $z$ ranging over $\mathrm{Lab}^0$ modulo $R\mathbb{Z}^3$,
\[
\sum_z\mathbf{1}_{\Delta_z}=1,\qquad\mathbf{1}_{\Delta_z}\mathbf{1}_{\Delta_{z'}}=0\quad(z\ne z'),\qquad|F|_{\mathrm{su}}=\sup_z\left|\mathbf{1}_{\Delta_z}F\right|_{\mathrm{su}}.
\]
The sums below are finite.
By the sign rules the moduli $|f(x)|_2$, $|v(z)|_2$ and $|F(x,x')|_2$ do not change under the translations by $R\mathbb{Z}^3$ (for a pair, both points being translated);
hence a supremum over the orbit of $\Delta_z\times O_{\mathrm{cub}}$ equals the supremum over $\Delta_z\times O_{\mathrm{cub}}$, and in the estimates of $\left\|\mathbf{1}_{\Delta_z}O\mathbf{1}_{\Delta_{z'}}\right\|_{\mathrm{su}}$ the row index is taken in $\Delta_z$ itself.
The properties (N1)--(N3) and (Nk1)--(Nk5) of Section~G1.3 are proved as follows.

\begin{description}
\item[(N1)]
$\left|\mathbf{1}_{\Delta_z}Of\right|_{\mathrm{su}}\le\sum_{z'}\left\|\mathbf{1}_{\Delta_z}O\mathbf{1}_{\Delta_{z'}}\right\|_{\mathrm{su}}\cdot|f|_{\mathrm{su}}\le N_{\gamma}(O)\cdot|f|_{\mathrm{su}}$;
the second assertion of (N1) retains one term of the sum that defines $N_{\gamma}$.
\item[(N2)]
The first assertion is the triangle inequality.
The identity
\[
\mathbf{1}_{\Delta_z}O_1O_2\mathbf{1}_{\Delta_{z'}}=\sum_{z''}\left(\mathbf{1}_{\Delta_z}O_1\mathbf{1}_{\Delta_{z''}}\right)\cdot\left(\mathbf{1}_{\Delta_{z''}}O_2\mathbf{1}_{\Delta_{z'}}\right)
\]
and the inequality $d(z,z')\le d(z,z'')+d(z'',z')$ give
\begin{align*}
\sum_{z'}\exp\left(\gamma\cdot d(z,z')\right)\cdot\left\|\mathbf{1}_{\Delta_z}O_1O_2\mathbf{1}_{\Delta_{z'}}\right\|_{\mathrm{su}}&\le\sum_{z''}\exp\left(\gamma\cdot d(z,z'')\right)\cdot\left\|\mathbf{1}_{\Delta_z}O_1\mathbf{1}_{\Delta_{z''}}\right\|_{\mathrm{su}}\cdot N_{\gamma}(O_2)\\
&\le N_{\gamma}(O_1)\cdot N_{\gamma}(O_2).
\end{align*}
\item[(N3)]
A scalar $\psi$ commutes with $\mathbf{1}_{\Delta_z}$, since both multiply by functions of the site (on functions of pairs, of the first point), and $|\psi F|_{\mathrm{su}}\le|\psi|_\infty\cdot|F|_{\mathrm{su}}$ in each of the three spaces.
Hence $\mathbf{1}_{\Delta_z}\psi\mathbf{1}_{\Delta_{z'}}=0$ for $z\ne z'$.
Therefore, for each $z$, the sum that defines $N_{\gamma}(\psi)$ reduces to its term $z'=z$, in which $d(z,z)=0$, and
\[
N_{\gamma}(\psi)=\sup_z\left\|\mathbf{1}_{\Delta_z}\psi\mathbf{1}_{\Delta_z}\right\|_{\mathrm{su}}.
\]
Consider first the spinors on $\mathbb{T}$, which at $k=0$ are the spinors on the unit lattice.
Since $|\psi(c)|$ and $|f(c,g)|_2$ do not change under the translations by $R\mathbb{Z}^3$, every spinor $f$ satisfies
\[
\left|\mathbf{1}_{\Delta_z}\psi\mathbf{1}_{\Delta_z}f\right|_{\mathrm{su}}=\sup\left\{|\psi(c)|\cdot|f(c,g)|_2:\ c\in\Delta_z,\ g\in O_{\mathrm{cub}}\right\}\le\sup_{c\in\Delta_z}|\psi(c)|\cdot|f|_{\mathrm{su}}.
\]
Equality holds for the spinor $f=\mathrm{dsp}[(c,g),\mathrm{vv}]$ of Section~0.2, where $\mathrm{vv}\ne0$ and $c\in\Delta_z$ is a point at which $|\psi|$ attains its maximum over the finite set $\Delta_z$:
this spinor vanishes off the orbit of $(c,g)$, so that $\mathbf{1}_{\Delta_z}\psi\mathbf{1}_{\Delta_z}f=\psi(c)\cdot f$.
Thus $\left\|\mathbf{1}_{\Delta_z}\psi\mathbf{1}_{\Delta_z}\right\|_{\mathrm{su}}=\sup_{c\in\Delta_z}|\psi(c)|$.
The blocks $\Delta_z$, $z\in\mathrm{Lab}^0$, form a partition of $\mathrm{Lab}$ (Section~G1.2), and $\psi$ is periodic;
it follows that $N_{\gamma}(\psi)=|\psi|_\infty$ on the spinors on $\mathbb{T}$.
Consider next the spinors on the unit lattice at $k\ge1$, on which $\psi$ acts at a unit site $z$ by its value $\psi(z)$.
There the operator $\mathbf{1}_{\Delta_z}\psi\mathbf{1}_{\Delta_z}$ multiplies the values on the class of $z$ by the number $\psi(z)$ and replaces the values on the other classes by $0$.
Its norm is $|\psi(z)|$:
the inequality $\left|\mathbf{1}_{\Delta_z}\psi\mathbf{1}_{\Delta_z}v\right|_{\mathrm{su}}\le|\psi(z)|\cdot|v|_{\mathrm{su}}$ holds for every spinor $v$ on the unit lattice, with equality when $v\ne0$ vanishes outside the class of $z$.
Consequently, on the spinors on the unit lattice at $k\ge1$,
\[
N_{\gamma}(\psi)=\sup_{z\in\mathrm{Lab}^0}|\psi(z)|\le|\psi|_\infty,
\]
where the inequality holds because $\mathrm{Lab}^0\subset\mathrm{Lab}$.
Finally, let $O$ be an operator between any two of the three spaces, the functions of short pairs being admitted as values.
Then $\mathbf{1}_{\Delta_z}\psi O\mathbf{1}_{\Delta_{z'}}=\psi\,\mathbf{1}_{\Delta_z}O\mathbf{1}_{\Delta_{z'}}$, and therefore
\[
\left\|\mathbf{1}_{\Delta_z}\psi O\mathbf{1}_{\Delta_{z'}}\right\|_{\mathrm{su}}\le|\psi|_\infty\cdot\left\|\mathbf{1}_{\Delta_z}O\mathbf{1}_{\Delta_{z'}}\right\|_{\mathrm{su}};
\]
multiplying by $\exp\left(\gamma\cdot d(z,z')\right)$ and summing over $z'$, one obtains $N_{\gamma}(\psi O)\le|\psi|_\infty\cdot N_{\gamma}(O)$.
In particular $N_{\gamma}(1)=1$ on the spinors on $\mathbb{T}$ and on the spinors on the unit lattice.
\end{description}

Now let the level be $0$, let $O$, $O_1$, $O_2$ be operators on the spinors on $\mathbb{T}$, and recall that $\Delta_z\times O_{\mathrm{cub}}$ is the cluster of $z$.

\begin{description}
\item[(Nk1)]
For $x$ over $z$,
\[
\left|\left(\mathbf{1}_{\Delta_z}O\mathbf{1}_{\Delta_{z'}}f\right)(x)\right|_2=\left|\frac{1}{48}\sum_{g'}O\left(x,(z',g')\right)\cdot f(z',g')\right|_2\le\sup\left\{|O(x,x')|_2:\ x'\text{ over }z'\right\}\cdot|f|_{\mathrm{su}},
\]
the sum consisting of $48$ terms of weight $\frac{1}{48}$.
\item[(Nk2)]
Let $x$ lie over $z$ and $x'$ over $z'$, and put $u:=O_2(\cdot,x')\,\mathrm{vv}$, a spinor.
By (kn3),
\[
\left|(O_1O_2)(x,x')\,\mathrm{vv}\right|_2=\left|(O_1u)(x)\right|_2\le\sum_{z''}\left\|\mathbf{1}_{\Delta_z}O_1\mathbf{1}_{\Delta_{z''}}\right\|_{\mathrm{su}}\cdot\sup\left\{|O_2(x'',x')|_2:\ x''\text{ over }z''\right\}\cdot|\mathrm{vv}|_2.
\]
Together with $d(z,z')\le d(z,z'')+d(z'',z')$ this gives
\begin{align*}
&\sum_{z'}\exp\left(\gamma\cdot d(z,z')\right)\cdot\sup\left\{|(O_1O_2)(x,x')|_2:\ x\text{ over }z,\ x'\text{ over }z'\right\}\\
&\quad\le\sum_{z''}\exp\left(\gamma\cdot d(z,z'')\right)\cdot\left\|\mathbf{1}_{\Delta_z}O_1\mathbf{1}_{\Delta_{z''}}\right\|_{\mathrm{su}}\\
&\qquad\cdot\sum_{z'}\exp\left(\gamma\cdot d(z'',z')\right)\cdot\sup\left\{|O_2(x'',x')|_2:\ x''\text{ over }z'',\ x'\text{ over }z'\right\}\\
&\quad\le N_{\gamma}(O_1)\cdot N^{\mathrm{ke}}_{\gamma}(O_2).
\end{align*}
\item[(Nk3)]
This follows from (kn2).
\item[(Nk4)]
This retains one term of the sum that defines $N^{\mathrm{ke}}_{\gamma}$.
\item[(Nk5)]
This is (kn4).
\item[(cg)]
The number $c(\gamma)=\sum_{n\in\mathbb{Z}^3}\exp\left(-\gamma\cdot|n|\right)$ is finite for $\gamma>0$, since at most $(2m+1)^3$ points $n$ satisfy $|n|=m$;
it is at least $1$ (the term $n=0$), and it does not increase in $\gamma$.
On $\mathbb{T}$ one has $\sum_{z'}\exp\left(-\gamma\cdot d(z,z')\right)\le c(\gamma)$:
every class $z'$ has a representative with $|z'-z|=d(z,z')$, different classes have different representatives, and $z'-z\in\mathbb{Z}^3$.
\end{description}

\subsubsection*{0.7. The Neumann series}

The operators on the spinors on $\mathbb{T}$ form a finite-dimensional algebra.
On it $N_{\gamma}$ is a norm (it is definite, because $N_{\gamma}(O)\ge\|O\|_{\mathrm{su}}$ by (N1)), and $N_{\gamma}(O_1O_2)\le N_{\gamma}(O_1)\cdot N_{\gamma}(O_2)$, $N_{\gamma}(1)=1$.
If $N_{\gamma}(O)\le\frac{1}{2}$, then the series $\sum_{n\ge0}O^n$ converges in $N_{\gamma}$, its sum satisfies $N_{\gamma}\left(\sum_{n\ge0}O^n\right)\le2$, and it is a left and a right inverse of $1-O$.
Indeed, $N_{\gamma}(O^n)\le2^{-n}$ and $\sum_n2^{-n}=2$; moreover
\[
(1-O)\sum_{n<m}O^n=\sum_{n<m}O^n\,(1-O)=1-O^m,
\]
and $N_{\gamma}(O^m)\to0$ as $m\to\infty$.

\subsubsection*{0.8. The constants of the Propositions 2.m}

The constants $C'_{V}$, $c_{\mathrm{re}}$, $a_{z}$, $k_{1}$, $C_{00}$, $C_z$, $\gamma_{z}$, $\gamma^{F}$, $a'$ are those defined before Proposition~2.1;
they are listed in Section~I.8.2, lines (C4), (C4r), and their formulas are those of Proposition~14.2\,(a), lines (C4), (C4r):
\begin{gather*}
C'_{V}:=2\,\mathrm{e}\cdot(42+6b),\qquad a_{z}:=\min\left\{\frac{1}{3L},\ \frac{c_{\mathrm{re}}}{2\,\mathrm{e}\cdot L\cdot(6+3b)}\right\},\qquad k_{1}:=60+b,\\
C_{00}:=48\cdot\frac{2L}{c_{\mathrm{re}}}\cdot\left(1+2k_{1}\right),\qquad C_z:=\max\left(C_{5},C_{00}\right),\\
\gamma_{z}:=\min\left(\gamma_{(5)},a_{z}\right),\qquad\gamma^{F}:=\frac{\gamma_{z}}{2},\qquad a':=C_z\cdot c\left(\gamma^{F}\right).
\end{gather*}
Here $C_{5}\ge1$ and $\gamma_{(5)}=\frac{a_{0}}{L}>0$ are the constants of Proposition~1.6, and $c_{\mathrm{re}}>0$ is the constant of Proposition~1.7.
In particular, (co1), (co3) and (co4) below give the inequalities $\gamma^{F}L\le\frac{1}{6}<\log2$, $a'\ge1$ and $C'_{V}\ge24\,\mathrm{e}$ stated with these definitions.

\begin{description}
\item[(co1)]
$0<\gamma_{z}\le a_{z}\le\frac{1}{3L}$.
Hence $0<\gamma^{F}\le\frac{1}{6L}$, $\gamma^{F}L\le\frac{1}{6}<\log2$ (as $\exp\left(\frac{1}{6}\right)<2$), and $\gamma^{F}\le\frac{1}{30}$.
\item[(co2)]
$\gamma_{(5)}-\gamma^{F}\ge\gamma_{z}-\frac{\gamma_{z}}{2}=\gamma^{F}$, and $a_{z}-\gamma^{F}\ge\gamma^{F}$.
\item[(co3)]
$a'=C_z\cdot c\left(\gamma^{F}\right)$ is finite and at least $1$, because $C_z\ge C_{5}\ge1$ and by (cg).
\item[(co4)]
$C'_{V}=2\,\mathrm{e}\cdot(42+6b)\ge84\,\mathrm{e}\ge24\,\mathrm{e}$.
\item[(co5)]
$c_{\mathrm{re}}=\min\left\{b\left(1-\frac{1}{L}\right),\ \frac{\rho}{2C_{\mathrm{Po}}L}\right\}$ (Proposition~1.7), where $C_{\mathrm{Po}}=\frac{1}{8c_{*}}+\frac{1}{\lambda_{\mathrm{int}}}$, $c_{*}=\frac{\min\left(w_3,\lambda_{\mathrm{int}}\right)}{18}$, $\lambda_{\mathrm{int}}>0$ and $w_3=6$ (Section~I.8.2, lines (C0), (C0r)).
Hence $c_{\mathrm{re}}\le b\left(1-\frac{1}{L}\right)<b$.
Moreover $c_{*}\le\frac{w_3}{18}=\frac{1}{3}$ and $C_{\mathrm{Po}}\ge\frac{1}{8c_{*}}\ge\frac{3}{8}$; therefore
\begin{gather*}
c_{\mathrm{re}}\le\frac{\rho}{2C_{\mathrm{Po}}L}\le\frac{4\rho}{3L}<\frac{4}{3L},\qquad\frac{2L}{c_{\mathrm{re}}}>\frac{3}{2}\,L^2>37\quad(L\ge5),\\
\frac{C_{00}}{48}=\frac{2L}{c_{\mathrm{re}}}\cdot\left(1+2k_{1}\right)>3.
\end{gather*}
\item[(co6)]
$12\sqrt{2}\le17$, because $288\le289$; hence $\mathrm{sw}\le29$ and $(1+\rho)\cdot\mathrm{sw}<58$.
\end{description}

\subsection*{1. Perturbations, commutators, H\"older quotient of a product}

\begin{proof}[Proof of Proposition~2.1]
\textbf{(a)}
By (op8), the difference $V:=F_y(B)-F_y(0)$ contains neither the number $\mu$ nor, at $k=0$, the matrix $i\gamma_5y$.
For $k\ge1$ one has
\begin{gather*}
V=V^{\mathrm{hop}}+V^{(k)}+V^{(k+1)},\qquad V^{\mathrm{hop}}:=D_B-D_0,\\
V^{(k)}:=\boldsymbol{\alpha}_y\left(JP_k(B)J^{*}-JP_k(0)J^{*}\right),\qquad V^{(k+1)}:=\boldsymbol{\beta}_y\left(JP_{k+1}(B)J^{*}-JP_{k+1}(0)J^{*}\right),
\end{gather*}
and for $k=0$
\[
V=V^{\mathrm{hop}}+V^{(1)},\qquad V^{(1)}:=b_L\left(JP(B)J^{*}-JP(0)J^{*}\right).
\]
Fix $z_1\in\mathrm{Lab}^0$, a site $x=(c,g)$ over $\Delta_{z_1}$ and a spinor $f$, and put $z':=\mathrm{cen}_{k+1}(c)=\mathrm{cen}_{k+1}(z_1)$.
By (ph0), every phase formed with the connection $0$ equals $1$.

\emph{The operator $V^{\mathrm{hop}}$, all $k$.}
By (op2),
\[
\left(V^{\mathrm{hop}}f\right)(x)=-\eta^{-1}\cdot M_3(x)\cdot\left(u_B(c,c_3)-1\right)\cdot f\left(x^{(3)}\right).
\]
From (hm5) and Lemma~0.1\,(b) one obtains
\[
\left|\left(V^{\mathrm{hop}}f\right)(x)\right|_2\le\eta^{-1}\cdot6\cdot\eta\delta\cdot\exp(\eta\delta)\cdot\left|f\left(x^{(3)}\right)\right|_2=6\delta\cdot\exp(\eta\delta)\cdot\left|f\left(x^{(3)}\right)\right|_2 .
\]
According to (bl5), the site $x^{(3)}$ lies over $\Delta_{z_1}$ or over $\Delta_{z_1+a}$, where $a=ge_1\in\{\pm e_1,\pm e_2,\pm e_3\}$.
Consequently $\mathbf{1}_{\Delta_{z_1}}V^{\mathrm{hop}}\mathbf{1}_{\Delta_{z_2}}=0$ unless $z_2$ is congruent modulo $R\mathbb{Z}^3$ to one of the seven sites $z_1$, $z_1\pm e_\mu$.
For each of these seven classes one has $\left\|\mathbf{1}_{\Delta_{z_1}}V^{\mathrm{hop}}\mathbf{1}_{\Delta_{z_2}}\right\|_{\mathrm{su}}\le6\delta\cdot\exp(\eta\delta)$, and $d(z_1,z_1\pm e_\mu)\le1$ by (bl7).
Therefore
\[
\sum_{z_2}\exp\left(\gamma\,d(z_1,z_2)\right)\left\|\mathbf{1}_{\Delta_{z_1}}V^{\mathrm{hop}}\mathbf{1}_{\Delta_{z_2}}\right\|_{\mathrm{su}}\le\left(1+6\exp(\gamma)\right)\cdot6\delta\cdot\exp(\eta\delta)\le42\exp(\gamma)\cdot\exp(\delta)\cdot\delta .
\]

\emph{The operator $V^{(k)}$, $k\ge1$.}
By (op5),
\[
\left(V^{(k)}f\right)(x)=\boldsymbol{\alpha}_y\int_{\Delta_{z_1}}\mathrm{d}\tilde x\,\left[\varphi_B(z_1,c)^{-1}\varphi_B(z_1,\tilde c)-1\right]f(\tilde x),
\]
where $\tilde x=(\tilde c,\tilde g)$ and the integral stands for $\eta^3\sum_{\tilde c\in\Delta_{z_1}}\frac{1}{48}\sum_{\tilde g\in O_{\mathrm{cub}}}$.
By (mx7), Lemma~0.1\,(b) and $\int_{\Delta_{z_1}}\mathrm{d}\tilde x=1$ (bl1),
\[
\left|\left(V^{(k)}f\right)(x)\right|_2\le b\cdot3\delta\cdot\exp(3\delta)\cdot\left|\mathbf{1}_{\Delta_{z_1}}f\right|_{\mathrm{su}} .
\]
Hence $\mathbf{1}_{\Delta_{z_1}}V^{(k)}\mathbf{1}_{\Delta_{z_2}}=0$ whenever $z_2$ is not congruent to $z_1$, and
\[
\sum_{z_2}\exp\left(\gamma\,d(z_1,z_2)\right)\left\|\mathbf{1}_{\Delta_{z_1}}V^{(k)}\mathbf{1}_{\Delta_{z_2}}\right\|_{\mathrm{su}}\le3b\cdot\delta\cdot\exp(3\delta).
\]

\emph{The operator $V^{(k+1)}$, $k\ge1$.}
The block $\mathrm{Bl}_{k+1}(z')$ is the disjoint union of the $L^3$ unit blocks $\Delta_{z_2}$, $z_2\in B(z')$, by (bl9).
So (op6) gives
\[
\left(V^{(k+1)}f\right)(x)=\boldsymbol{\beta}_y\,L^{-3}\sum_{z_2\in B(z')}\int_{\Delta_{z_2}}\mathrm{d}\tilde x\,\left[\varphi^{+}_B(z',c)^{-1}\varphi^{+}_B(z',\tilde c)-1\right]f(\tilde x).
\]
It follows that $\mathbf{1}_{\Delta_{z_1}}V^{(k+1)}\mathbf{1}_{\Delta_{z_2}}=0$ unless $z_2$ is congruent to a site of $B(z')$.
For $z_2\in B(z')$, (mx7), Lemma~0.1\,(b) and (bl1) yield
\[
\left\|\mathbf{1}_{\Delta_{z_1}}V^{(k+1)}\mathbf{1}_{\Delta_{z_2}}\right\|_{\mathrm{su}}\le b_L\cdot L^{-3}\cdot3L\delta\cdot\exp(3L\delta),
\]
and $d(z_1,z_2)\le L-1$ by (bl4), (bl7).
Therefore
\begin{align*}
&\sum_{z_2}\exp\left(\gamma\,d(z_1,z_2)\right)\left\|\mathbf{1}_{\Delta_{z_1}}V^{(k+1)}\mathbf{1}_{\Delta_{z_2}}\right\|_{\mathrm{su}}\\
&\qquad\le L^3\exp\left(\gamma(L-1)\right)\cdot b_LL^{-3}\cdot3L\delta\cdot\exp(3L\delta)\\
&\qquad=3b\cdot\delta\cdot\exp\left(\gamma(L-1)\right)\cdot\exp(3L\delta).
\end{align*}

\emph{The sum, $k\ge1$.}
Since $\gamma L\le\log2$ and $3L\delta\le1$, each of the numbers $\exp(\gamma)$, $1$, $\exp\left(\gamma(L-1)\right)$ is at most $\exp(\gamma L)\le2$, and each of the numbers $\exp(\delta)$, $\exp(3\delta)$, $\exp(3L\delta)$ is at most $\exp(3L\delta)\le\mathrm{e}$.
With (N2) one arrives at
\[
N_\gamma(V)\le2\mathrm{e}\cdot(42+3b+3b)\cdot\delta=C'_V\,\delta .
\]

\emph{Level $0$.}
Here $\Delta_{z_1}=\{z_1\}$, $c=z_1$ and $\eta=1$.
The site $x^{(3)}$ lies over $z_1+ge_1\ne z_1$; so the term $z_2=z_1$ does not occur in the sum for $V^{\mathrm{hop}}$, and that sum is at most $36\exp(\gamma)\cdot\exp(\delta)\cdot\delta$.
By (op7),
\[
\left(V^{(1)}f\right)(x)=b_L\,L^{-3}\sum_{z_2\in B(z')}\left[U_B(z',z_1)^{-1}U_B(z',z_2)-1\right]\frac{1}{48}\sum_{\tilde g\in O_{\mathrm{cub}}}f(z_2,\tilde g),
\]
and $U_B(z',\cdot\,)=\varphi^{+}_B(z',\cdot\,)$ (Section~0.3).
The mean over the $48$ sites of the cluster of $z_2$ has the norm $|\cdot|_2$ at most $\left|\mathbf{1}_{\Delta_{z_2}}f\right|_{\mathrm{su}}$.
Hence $\left\|\mathbf{1}_{\Delta_{z_1}}V^{(1)}\mathbf{1}_{\Delta_{z_2}}\right\|_{\mathrm{su}}\le b_LL^{-3}\cdot3L\delta\cdot\exp(3L\delta)$ for the $L^3$ sites $z_2$ of $B(z')$, the operator vanishes for every other class, and the corresponding sum is at most $3b\cdot\delta\cdot\exp\left(\gamma(L-1)\right)\cdot\exp(3L\delta)$.
Altogether
\[
N_\gamma(V)\le2\mathrm{e}\cdot(36+3b)\cdot\delta\le C'_V\,\delta .
\]

\textbf{(b)}
The number $\mu$ and the constant matrices $\boldsymbol{\alpha}_y$, $\boldsymbol{\beta}_y$, $i\gamma_5y$ commute with the multiplication by $\psi$.
Therefore
\begin{align*}
[F_y(A),\psi]&=[D_A,\psi]+\boldsymbol{\alpha}_y\left[JP_k(A)J^{*},\psi\right]+\boldsymbol{\beta}_y\left[JP_{k+1}(A)J^{*},\psi\right]&&(k\ge1),\\
[F_y(A),\psi]&=[D_A,\psi]+b_L\left[JP(A)J^{*},\psi\right]&&(k=0).
\end{align*}
Let $z_1$, $x$, $z'$, $f$ be as in (a).
By (op2), (op5), (op6),
\begin{align*}
\left([D_A,\psi]f\right)(x)&=-\eta^{-1}\cdot M_3(x)\cdot u_A(c,c_3)\cdot\left(\psi(c_3)-\psi(c)\right)\cdot f\left(x^{(3)}\right),\\
\left(\left[JP_k(A)J^{*},\psi\right]f\right)(x)&=\varphi_A(z_1,c)^{-1}\int_{\Delta_{z_1}}\mathrm{d}\tilde x\,\varphi_A(z_1,\tilde c)\left(\psi(\tilde c)-\psi(c)\right)f(\tilde x),\\
\left(\left[JP_{k+1}(A)J^{*},\psi\right]f\right)(x)&=\varphi^{+}_A(z',c)^{-1}L^{-3}\sum_{z_2\in B(z')}\int_{\Delta_{z_2}}\mathrm{d}\tilde x\,\varphi^{+}_A(z',\tilde c)\left(\psi(\tilde c)-\psi(c)\right)f(\tilde x).
\end{align*}
One has $|c_3-c|=|c_3-c|_1=\eta\le L$, whence $|\psi(c_3)-\psi(c)|\le\mathrm{Lip}\cdot\eta$.
For $\tilde c\in\Delta_{z_1}$ one has $|\tilde c-c|<1$ and $|\psi(\tilde c)-\psi(c)|\le3\,\mathrm{Lip}$; for $\tilde c\in\mathrm{Bl}_{k+1}(z')$ one has $|\tilde c-c|<L$ and $|\psi(\tilde c)-\psi(c)|\le3L\cdot\mathrm{Lip}$ (bl4).
Lemma~0.1\,(a) gives
\[
|u_A|\le\exp\left(\eta\,\mathfrak{i}(A)\right),\qquad\left|\varphi_A(z_1,c)^{-1}\varphi_A(z_1,\tilde c)\right|\le\exp\left(3\,\mathfrak{i}(A)\right),\qquad\left|\varphi^{+}_A(z',c)^{-1}\varphi^{+}_A(z',\tilde c)\right|\le\exp\left(3L\,\mathfrak{i}(A)\right).
\]
Together with (hm5), (mx7), (bl1), (bl5), and since $\mathrm{Bl}_{k+1}(z')$ is the disjoint union of the blocks $\Delta_{z_2}$, $z_2\in B(z')$, this implies
\begin{align*}
\left\|\mathbf{1}_{\Delta_{z_1}}[D_A,\psi]\mathbf{1}_{\Delta_{z_2}}\right\|_{\mathrm{su}}&\le6\,\mathrm{Lip}\cdot\exp\left(\eta\,\mathfrak{i}(A)\right)&&\text{for the seven classes }z_1,\ z_1\pm e_\mu,\\
\left\|\mathbf{1}_{\Delta_{z_1}}\boldsymbol{\alpha}_y\left[JP_k(A)J^{*},\psi\right]\mathbf{1}_{\Delta_{z_2}}\right\|_{\mathrm{su}}&\le3b\cdot\mathrm{Lip}\cdot\exp\left(3\,\mathfrak{i}(A)\right)&&\text{for }z_2=z_1,\\
\left\|\mathbf{1}_{\Delta_{z_1}}\boldsymbol{\beta}_y\left[JP_{k+1}(A)J^{*},\psi\right]\mathbf{1}_{\Delta_{z_2}}\right\|_{\mathrm{su}}&\le b_LL^{-3}\cdot3L\cdot\mathrm{Lip}\cdot\exp\left(3L\,\mathfrak{i}(A)\right)&&\text{for the }L^3\text{ classes of }B(z'),
\end{align*}
while each of the three operators vanishes for the remaining classes $z_2$.
Moreover $d(z_1,z_1\pm e_\mu)\le1$, and $d(z_1,z_2)\le L-1$ for $z_2\in B(z')$, by (bl4), (bl7).
By (N2), and since $\exp(\gamma L)\le2$ and $\exp\left(3L\,\mathfrak{i}(A)\right)\le\mathrm{e}$,
\begin{align*}
N_\gamma\left([F_y(A),\psi]\right)&\le\left(1+6\exp(\gamma)\right)\cdot6\,\mathrm{Lip}\cdot\exp\left(\eta\,\mathfrak{i}(A)\right)+3b\cdot\mathrm{Lip}\cdot\exp\left(3\,\mathfrak{i}(A)\right)\\
&\qquad+L^3\exp\left(\gamma(L-1)\right)\cdot b_LL^{-3}\cdot3L\cdot\mathrm{Lip}\cdot\exp\left(3L\,\mathfrak{i}(A)\right)\\
&\le\exp(\gamma L)\cdot\exp\left(3L\,\mathfrak{i}(A)\right)\cdot(42+3b+3b)\cdot\mathrm{Lip}\le2\mathrm{e}\cdot(42+6b)\cdot\mathrm{Lip}=C'_V\,\mathrm{Lip}.
\end{align*}

\emph{Level $0$.}
Now $c=z_1$, and $x^{(3)}$ lies over $z_1+ge_1\ne z_1$.
From (op7),
\[
\left(\left[JP(A)J^{*},\psi\right]f\right)(x)=U_A(z',z_1)^{-1}L^{-3}\sum_{z_2\in B(z')}U_A(z',z_2)\left(\psi(z_2)-\psi(z_1)\right)\frac{1}{48}\sum_{\tilde g\in O_{\mathrm{cub}}}f(z_2,\tilde g),
\]
with $U_A(z',\cdot\,)=\varphi^{+}_A(z',\cdot\,)$.
Consequently $\left\|\mathbf{1}_{\Delta_{z_1}}[D_A,\psi]\mathbf{1}_{\Delta_{z_2}}\right\|_{\mathrm{su}}\le6\,\mathrm{Lip}\cdot\exp\left(\mathfrak{i}(A)\right)$ for the six classes $z_1\pm e_\mu$ and $=0$ for all others, and $\left\|\mathbf{1}_{\Delta_{z_1}}b_L\left[JP(A)J^{*},\psi\right]\mathbf{1}_{\Delta_{z_2}}\right\|_{\mathrm{su}}\le b_LL^{-3}\cdot3L\cdot\mathrm{Lip}\cdot\exp\left(3L\,\mathfrak{i}(A)\right)$ for the $L^3$ classes of $B(z')$ and $=0$ for all others.
Hence
\[
N_\gamma\left([F_y(A),\psi]\right)\le36\exp(\gamma)\cdot\mathrm{Lip}\cdot\exp\left(\mathfrak{i}(A)\right)+3b\cdot\mathrm{Lip}\cdot\exp\left(\gamma(L-1)\right)\cdot\exp\left(3L\,\mathfrak{i}(A)\right)\le2\mathrm{e}\cdot(36+3b)\cdot\mathrm{Lip}\le C'_V\,\mathrm{Lip}.
\]

\emph{The kernel.}
By (kn2),
\[
[F_y(A),\psi](x,x')=F_y(A)(x,x')\cdot\left(\psi(c')-\psi(c)\right),
\]
which vanishes if $\psi(c')=\psi(c)$.
By (op9), $F_y(A)(x,x')=0$ unless there is an $n\in\mathbb{Z}^3$ such that either $x'+Rn\in\left\{x,x^{(0)},x^{(1)},x^{(2)}\right\}$, or $x'+Rn=x^{(3)}$, or $c$ and $c'+Rn$ lie in one block of side $L$.
In the first case $c'+Rn=c$, and then $\psi(c')=\psi(c)$, since $\psi$ is $R\mathbb{Z}^3$-periodic.
So the kernel vanishes unless $\psi(c)\ne\psi(c')$ and, modulo $R\mathbb{Z}^3$, either $x'=x^{(3)}$ or $c$ and $c'$ lie in one block of side $L$.

\textbf{(c)}
By (I.1.7),
\begin{align*}
\left(\delta_{\alpha,A}(\psi f)\right)(x,x')&=d_X(x,x')^{-\alpha}\left(U_A(x,x')\,\psi(c')\,f(x')-\psi(c)\,f(x)\right)\\
&=\psi(c)\cdot d_X(x,x')^{-\alpha}\left(U_A(x,x')\,f(x')-f(x)\right)+d_X(x,x')^{-\alpha}\left(\psi(c')-\psi(c)\right)U_A(x,x')\,f(x'),\\
\left(\delta_{\alpha,B}f-\delta_{\alpha,0}f\right)(x,x')&=d_X(x,x')^{-\alpha}\left(U_B(x,x')\,f(x')-f(x)\right)-d_X(x,x')^{-\alpha}\left(f(x')-f(x)\right)\\
&=d_X(x,x')^{-\alpha}\left(U_B(x,x')-1\right)f(x').
\end{align*}
For a pair inside a cluster one has $c'=c$; then $\psi(c')-\psi(c)=0$ and $U_B(x,x')=U_B(c,c)=1$.

\emph{The operator $E_{\mathrm{wq}}O$ (the last assertion of (c)).}
Let $\mathrm{wq}$ be a function on the short pairs with values in $\mathbb{C}$, or in the complex $4\times4$ matrices, such that $\mathrm{wq}(x+Re_\mu,x'+Re_\mu)=\mathrm{wq}(x,x')$ and $|\mathrm{wq}(x,x')|\le\mathrm{wq}_0$, respectively $|\mathrm{wq}(x,x')|_2\le\mathrm{wq}_0$.
Then $\left|\mathrm{wq}(x,x')\,\mathrm{vv}\right|_2\le\mathrm{wq}_0\,|\mathrm{vv}|_2$ for every $\mathrm{vv}\in\mathbb{C}^4$.
Moreover $E_{\mathrm{wq}}Of$ is a function of short pairs in the sense of Section~G1.3, because
\[
\left(E_{\mathrm{wq}}Of\right)(x+Re_\mu,x'+Re_\mu)=\mathrm{wq}(x,x')\,(Of)(x'+Re_\mu)=\boldsymbol{\sigma}\,\left(E_{\mathrm{wq}}Of\right)(x,x').
\]
At $k=0$ there is no short pair, and the assertion is void.
Let $k\ge1$, and let $(x,x')$ be a short pair with $x$ over $\Delta_z$.
Then $|c-c'|\le d_X(x,x')<1$, so that $x'$ lies over $\Delta_{\tilde z}$ with $\tilde z=z+a$, $a\in\mathbb{Z}^3$, $|a|\le1$ (bl6).
For every element $\mathrm{fs}$ of the space on which $O$ acts,
\begin{align*}
\left|\left(\mathbf{1}_{\Delta_z}E_{\mathrm{wq}}O\mathbf{1}_{\Delta_{z'}}\mathrm{fs}\right)(x,x')\right|_2&=\left|\mathrm{wq}(x,x')\left(\mathbf{1}_{\Delta_{\tilde z}}O\mathbf{1}_{\Delta_{z'}}\mathrm{fs}\right)(x')\right|_2\le\mathrm{wq}_0\left|\left(\mathbf{1}_{\Delta_{\tilde z}}O\mathbf{1}_{\Delta_{z'}}\mathrm{fs}\right)(x')\right|_2\\
&\le\mathrm{wq}_0\sum_{|a|\le1}\left\|\mathbf{1}_{\Delta_{z+a}}O\mathbf{1}_{\Delta_{z'}}\right\|_{\mathrm{su}}\,|\mathrm{fs}|_{\mathrm{su}}.
\end{align*}
By (bl7), $\exp\left(\gamma\,d(z,z')\right)\le\exp(\gamma)\cdot\exp\left(\gamma\,d(z+a,z')\right)$.
Summing over $z'$ one obtains
\begin{align*}
\sum_{z'}\exp\left(\gamma\,d(z,z')\right)\left\|\mathbf{1}_{\Delta_z}E_{\mathrm{wq}}O\mathbf{1}_{\Delta_{z'}}\right\|_{\mathrm{su}}&\le\mathrm{wq}_0\cdot\exp(\gamma)\sum_{|a|\le1}\sum_{z'}\exp\left(\gamma\,d(z+a,z')\right)\left\|\mathbf{1}_{\Delta_{z+a}}O\mathbf{1}_{\Delta_{z'}}\right\|_{\mathrm{su}}\\
&\le27\exp(\gamma)\cdot\mathrm{wq}_0\cdot N_\gamma(O).
\end{align*}

\emph{Uniformity.}
The constant $C'_V$ depends on $b$ alone.
None of the bounds in (a)--(c) depends on $k$, $t$, $R$, $\boldsymbol{\sigma}$, $e_k$, $\mu$ or $y$, nor on $A$, $B$, $\psi$, $f$ otherwise than through the numbers $\delta$, $\mathrm{Lip}$, $\mathrm{wq}_0$, $\gamma$ that appear in it.
\end{proof}

\subsection*{2. Algebra, rows and covariances of the elementary operators}

\begin{proof}[Proof of Proposition~2.2]
\emph{Level $0$ (the data of Proposition~2.2).}
By (I.3.4) and Section~I.3.9 one has $Q_0=P_0=1$; the unit lattice is $\mathrm{Lab}$, its spinors are the spinors on $\mathbb{T}$, and $\mathrm{uf}^0=\mathrm{uf}$.
The relations that contain $Q_k(A)J^{*}$, $JQ_k^T(-A)$, or one of the four-component operators $Q(A)$, $Q^T(-A)$, $P(A)$, are relations for $k\ge1$.
At $k=0$ the operator ``$JP_{k+1}(A)J^{*}$'' is $JP(A)J^{*}$, and the operators under consideration are $D_A$, $Q(A)J^{*}$, $JQ^T(-A)$, $JP(A)J^{*}$, $\operatorname{Pen}_y(A)=b_L\,JP(A)J^{*}+i\gamma_5y$ and $F_y(A)$.
At $k=0$ the operator $JP_k(A)J^{*}$ is read as $P_0=1$.
What (f0) asserts about it is then the following: $1\cdot1=1$; the kernel of $1$ is $48\cdot\mathbf{1}[\tilde x=x]\,\mathbf{1}_4$ ((I.ke1) at $\eta=1$, with $\tilde x$ ranging over a period cell that contains $x$), a multiple of $\mathbf{1}_4$ that couples $x$ to $x$ only and does not contain $A$; $1$ commutes with every constant $4\times4$ matrix; and the two products with $JP(A)J^{*}$ are computed in (f0) below.
The relations of (a)--(e) that contain this operator become
\[
1=\mathrm{uf}\,1\,\mathrm{uf}^{-1},\qquad1\,\mathrm{mS}=\mathrm{mS}\,1,\qquad1^T=1,\qquad\overline{1}=\varkappa^{-1}\,1\,\varkappa,\qquad1^{*}=1=1^2.
\]

\emph{Pairings} (Section~I.1.6, Section~G1.2; every sum extends over one period cell).
On the spinors on $\mathbb{T}$, on $\mathrm{Lab}^0$ ($k\ge1$) and on $\mathrm{Lab}_{k+1}$ the bilinear pairings are
\[
\langle f',f\rangle^{\mathrm{ch}}=\eta^3\sum_c\frac{1}{48}\sum_gf'(c,g)\cdot f(c,g),\qquad\langle v'',v\rangle_1=\sum_zv''(z)\cdot v(z),\qquad\langle v'',v'\rangle_L=L^3\sum_{z'}v''(z')\cdot v'(z'),
\]
and the sesquilinear pairings $(\cdot,\cdot)$ carry the complex conjugation on the first factor.
Transposes and adjoints are taken with respect to these pairings.
The summands are $R\mathbb{Z}^3$-periodic because $\boldsymbol{\sigma}^2=1$, so that no pairing depends on the choice of its period cell (Section~I.1.6).
In (c) and (e) the sums of the pairings extend over the three period cells of (bl9).

\textbf{(f0)}
\emph{Products.}
Let $k\ge1$.
By (op3) and (bl1),
\[
\left(Q_k(A)J^{*}\,JQ_k^T(-A)\,v\right)(z)=\eta^3\sum_{c\in\Delta_z}\varphi_A(z,c)\cdot\frac{1}{48}\sum_{g\in O_{\mathrm{cub}}}\varphi_A(z,c)^{-1}\,v(z)=\eta^3\cdot\eta^{-3}\cdot\frac{48}{48}\,v(z)=v(z),
\]
and by (op4)
\[
\left(Q(A)\,Q^T(-A)\,v'\right)(z')=L^{-3}\sum_{z\in B(z')}U_A(z',z)\,U_A(z',z)^{-1}\,v'(z')=v'(z').
\]
Let $k=0$.
By (op7),
\[
\left(Q(A)J^{*}\,JQ^T(-A)\,v'\right)(z')=L^{-3}\sum_{c\in B(z')}U_A(z',c)\cdot\frac{1}{48}\sum_{g\in O_{\mathrm{cub}}}U_A(z',c)^{-1}\,v'(z')=v'(z').
\]

\emph{Idempotents.}
For $k\ge1$ one has (Section~I.3.3)
\begin{align*}
JP_k(A)J^{*}&=\left(JQ_k^T(-A)\right)\left(Q_k(A)J^{*}\right),\\
P(A)&=Q^T(-A)\,Q(A),\\
JP_{k+1}(A)J^{*}&=\left(JQ_k^T(-A)\right)P(A)\left(Q_k(A)J^{*}\right).
\end{align*}
The two products just computed give
\begin{align*}
\left(JP_k(A)J^{*}\right)^2&=\left(JQ_k^T(-A)\right)\left[\left(Q_k(A)J^{*}\right)\left(JQ_k^T(-A)\right)\right]\left(Q_k(A)J^{*}\right)=JP_k(A)J^{*},\\
P(A)^2&=Q^T(-A)\left[Q(A)\,Q^T(-A)\right]Q(A)=P(A),\\
\left(JP_{k+1}(A)J^{*}\right)^2&=\left(JQ_k^T(-A)\right)P(A)\left[\left(Q_k(A)J^{*}\right)\left(JQ_k^T(-A)\right)\right]P(A)\left(Q_k(A)J^{*}\right)\\
&=\left(JQ_k^T(-A)\right)P(A)^2\left(Q_k(A)J^{*}\right)=JP_{k+1}(A)J^{*},\\
\left(JP_k(A)J^{*}\right)\left(JP_{k+1}(A)J^{*}\right)&=\left(JQ_k^T(-A)\right)\left[\left(Q_k(A)J^{*}\right)\left(JQ_k^T(-A)\right)\right]P(A)\left(Q_k(A)J^{*}\right)=JP_{k+1}(A)J^{*},\\
\left(JP_{k+1}(A)J^{*}\right)\left(JP_k(A)J^{*}\right)&=\left(JQ_k^T(-A)\right)P(A)\left[\left(Q_k(A)J^{*}\right)\left(JQ_k^T(-A)\right)\right]\left(Q_k(A)J^{*}\right)=JP_{k+1}(A)J^{*}.
\end{align*}
For $k=0$ one has $JP(A)J^{*}=\left(JQ^T(-A)\right)\left(Q(A)J^{*}\right)$, and therefore
\[
\left(JP(A)J^{*}\right)^2=\left(JQ^T(-A)\right)\left[\left(Q(A)J^{*}\right)\left(JQ^T(-A)\right)\right]\left(Q(A)J^{*}\right)=JP(A)J^{*};
\]
since $P_0=1$, both products $P_0\left(JP(A)J^{*}\right)$ and $\left(JP(A)J^{*}\right)P_0$ equal $JP(A)J^{*}$.

\emph{Kernels.}
Let $z=[c]$ and $z'=\mathrm{cen}_{k+1}(c)$ (for the kernel of $P(A)$ let $z'=\mathrm{cen}_{k+1}(z_1)$), and let $\tilde x$, respectively $z_2$, range over a period cell that contains the block in question.
By (op5), (op6), (op7), the kernels with respect to $\int\mathrm{d}\tilde x$ are
\begin{align*}
\left(JP_k(A)J^{*}\right)(x,\tilde x)&=\varphi_A(z,c)^{-1}\varphi_A(z,\tilde c)\cdot\mathbf{1}[\tilde c\in\Delta_z]\,\mathbf{1}_4,\\
\left(JP_{k+1}(A)J^{*}\right)(x,\tilde x)&=L^{-3}\,\varphi^{+}_A(z',c)^{-1}\varphi^{+}_A(z',\tilde c)\cdot\mathbf{1}\left[\tilde c\in\mathrm{Bl}_{k+1}(z')\right]\mathbf{1}_4,\\
\left(JP(A)J^{*}\right)(x,\tilde x)&=L^{-3}\,U_A(z',c)^{-1}U_A(z',\tilde c)\cdot\mathbf{1}\left[\tilde c\in B(z')\right]\mathbf{1}_4\qquad(k=0),
\end{align*}
and by (op4) the kernel of $P(A)$ with respect to $\sum_{z_2}$ is
\[
P(A)(z_1,z_2)=L^{-3}\,U_A(z',z_1)^{-1}U_A(z',z_2)\cdot\mathbf{1}\left[z_2\in B(z')\right]\mathbf{1}_4 .
\]
All of them are multiples of $\mathbf{1}_4$.

\emph{Constant matrices.}
(sc)\enspace By (op3), (op4), (op7), each of the six operators $Q_k(A)J^{*}$, $JQ_k^T(-A)$, $Q(A)$, $Q^T(-A)$, $Q(A)J^{*}$, $JQ^T(-A)$ multiplies the values of the spinors by numbers and adds the results.
Hence these six operators and the four idempotents commute with every constant $4\times4$ matrix.

\emph{Traces.}
For two linear maps $O_1$, $O_2$ in opposite directions between two finite-dimensional spaces, $\operatorname{Tr}(O_1O_2)=\operatorname{Tr}(O_2O_1)$.
The spinors on the $L$-lattice form a space of dimension $4(R/L)^3$.
Consequently
\begin{align*}
\operatorname{Tr}P(A)&=\operatorname{Tr}Q^T(-A)\,Q(A)=\operatorname{Tr}Q(A)\,Q^T(-A)=\operatorname{Tr}1=4(R/L)^3&&(k\ge1),\\
\operatorname{Tr}JP(A)J^{*}&=\operatorname{Tr}\left(JQ^T(-A)\right)\left(Q(A)J^{*}\right)=\operatorname{Tr}\left(Q(A)J^{*}\right)\left(JQ^T(-A)\right)=4(R/L)^3&&(k=0).
\end{align*}

\emph{Rows.}
The kernels displayed above show that $JP_k(A)J^{*}$ couples $x$ only to the sites over $\Delta_z$, the unit block of $c$ (and to their translates by $R\mathbb{Z}^3$, with the sign $\boldsymbol{\sigma}^n$), and that $P(A)$, $JP_{k+1}(A)J^{*}$ and $JP(A)J^{*}$ couple only sites in (respectively over) the block of side $L$.
By (ph0), the phases in the row at $x$ contain $A(b)$ only for bonds $b$ contained in this block.
The row of $D_A$ is read off (op1).

\textbf{(a)}
\emph{Step 1: $U_{\mathrm{Acl}}(c,c')=\mathrm{uf}(c')/\mathrm{uf}(c)$ for all $c,c'\in\mathrm{Lab}$.}
Let $\mathrm{Wb}$ be a lattice box of $\mathrm{Lab}$ that contains $c$ and $c'$.
For two orders $\pi$, $\pi'$, the path $\boldsymbol{\Gamma}^{\pi}(c,c')$ followed by the path $\boldsymbol{\Gamma}^{\pi'}(c,c')$ traversed in the opposite direction is a closed lattice path in $\mathrm{Wb}$ (bl8).
The last assertion of \cite[Lemma~16.26]{QEDI} bounds the sum of a connection over a closed lattice path of $n$ bonds in a box by $2n$ times the largest side length of the box times the supremum of the modulus of its curvature over the plaquettes of the box.
Since $d\mathrm{Acl}=0$, it applies with $\sup_p\left|(d\mathrm{Acl})(p)\right|=0$, and the sum of $\mathrm{Acl}$ over the closed path above vanishes.
Thus $\mathrm{Acl}\left(\boldsymbol{\Gamma}^{\pi}(c,c')\right)$ does not depend on $\pi$, and it equals $(\tau\mathrm{Acl})(c,c')$.
Multiplying the relations $\mathrm{uf}(c'')=\mathrm{uf}(\tilde c)\exp\left\{ie_k\eta\,\mathrm{Acl}(\tilde c,c'')\right\}$ along the bonds $(\tilde c,c'')$ of one path $\boldsymbol{\Gamma}^{\pi}(c,c')$, one obtains
\[
\mathrm{uf}(c')=\mathrm{uf}(c)\exp\left\{ie_k\eta\,\mathrm{Acl}\left(\boldsymbol{\Gamma}^{\pi}(c,c')\right)\right\}=\mathrm{uf}(c)\,U_{\mathrm{Acl}}(c,c').
\]

\emph{Step 2: phases.}
By (ph0), $\mathrm{pf}_{A-\mathrm{Acl}}=\mathrm{pf}_A\,\mathrm{pf}_{\mathrm{Acl}}^{-1}$.
Step 1, together with the telescoping product
\[
\prod_{j<k}\frac{\mathrm{uf}\left(\mathrm{cen}_{j+1}(c)\right)}{\mathrm{uf}\left(\mathrm{cen}_j(c)\right)}=\frac{\mathrm{uf}(z)}{\mathrm{uf}(c)}\qquad(c\in\Delta_z),
\]
yields
\begin{equation}\tag{ga1}
u_{A-\mathrm{Acl}}(c,c_3)=u_A(c,c_3)\,\frac{\mathrm{uf}(c)}{\mathrm{uf}(c_3)},\qquad U_{A-\mathrm{Acl}}(c,c')=U_A(c,c')\,\frac{\mathrm{uf}(c)}{\mathrm{uf}(c')},\qquad\varphi_{A-\mathrm{Acl}}(z,c)=\varphi_A(z,c)\,\frac{\mathrm{uf}(z)}{\mathrm{uf}(c)}.
\end{equation}

\emph{Step 3: operators.}
Since $\mathrm{uf}$ is periodic, $\mathrm{uf}\,f$ is a spinor on $\mathbb{T}$; likewise $\mathrm{uf}^0v$ and $\mathrm{uf}^1v'$ are spinors.

The operator $D$: by (op2) and (ga1),
\begin{align*}
\left(\left(\mathrm{uf}\,D_A\,\mathrm{uf}^{-1}-D_A\right)f\right)(x)&=-\eta^{-1}\cdot M_3(x)\cdot u_A(c,c_3)\cdot\left(\frac{\mathrm{uf}(c)}{\mathrm{uf}(c_3)}-1\right)\cdot f\left(x^{(3)}\right)\\
&=-\eta^{-1}\cdot M_3(x)\cdot\left(u_{A-\mathrm{Acl}}-u_A\right)(c,c_3)\cdot f\left(x^{(3)}\right)=\left(\left(D_{A-\mathrm{Acl}}-D_A\right)f\right)(x).
\end{align*}
Let $k\ge1$.
By (op3), (op4) and (ga1),
\begin{align*}
\left(Q_k(A-\mathrm{Acl})J^{*}f\right)(z)&=\int_{\Delta_z}\mathrm{d}x\,\varphi_A(z,c)\,\mathrm{uf}(z)\,\mathrm{uf}(c)^{-1}f(x)=\mathrm{uf}(z)\left(Q_k(A)J^{*}\,\mathrm{uf}^{-1}f\right)(z),\\
\left(JQ_k^T\left(-(A-\mathrm{Acl})\right)v\right)(x)&=\varphi_A(z,c)^{-1}\,\mathrm{uf}(c)\,\mathrm{uf}(z)^{-1}\,v(z)=\mathrm{uf}(c)\left(JQ_k^T(-A)\,(\mathrm{uf}^0)^{-1}v\right)(x),\\
\left(Q(A-\mathrm{Acl})\,v\right)(z')&=L^{-3}\sum_{z\in B(z')}U_A(z',z)\,\mathrm{uf}(z')\,\mathrm{uf}(z)^{-1}\,v(z)=\mathrm{uf}(z')\left(Q(A)\,(\mathrm{uf}^0)^{-1}v\right)(z'),\\
\left(Q^T\left(-(A-\mathrm{Acl})\right)v'\right)(z)&=U_A(z',z)^{-1}\,\mathrm{uf}(z)\,\mathrm{uf}(z')^{-1}\,v'(z')=\mathrm{uf}(z)\left(Q^T(-A)\,(\mathrm{uf}^1)^{-1}v'\right)(z).
\end{align*}
Composing these four relations one finds
\begin{align*}
JP_k(A-\mathrm{Acl})J^{*}&=\mathrm{uf}\,JQ_k^T(-A)\,(\mathrm{uf}^0)^{-1}\,\mathrm{uf}^0\,Q_k(A)J^{*}\,\mathrm{uf}^{-1}=\mathrm{uf}\,JP_k(A)J^{*}\,\mathrm{uf}^{-1},\\
P(A-\mathrm{Acl})&=\mathrm{uf}^0\,Q^T(-A)\,(\mathrm{uf}^1)^{-1}\,\mathrm{uf}^1\,Q(A)\,(\mathrm{uf}^0)^{-1}=\mathrm{uf}^0\,P(A)\,(\mathrm{uf}^0)^{-1},\\
JP_{k+1}(A-\mathrm{Acl})J^{*}&=\mathrm{uf}\,JQ_k^T(-A)\,(\mathrm{uf}^0)^{-1}\,\mathrm{uf}^0\,P(A)\,(\mathrm{uf}^0)^{-1}\,\mathrm{uf}^0\,Q_k(A)J^{*}\,\mathrm{uf}^{-1}=\mathrm{uf}\,JP_{k+1}(A)J^{*}\,\mathrm{uf}^{-1}.
\end{align*}
The matrices $\boldsymbol{\alpha}_y$, $\boldsymbol{\beta}_y$ and the number $\mu$ commute with $\mathrm{uf}$; hence $\operatorname{Pen}_y(A-\mathrm{Acl})=\mathrm{uf}\operatorname{Pen}_y(A)\,\mathrm{uf}^{-1}$ and $F_y(A-\mathrm{Acl})=\mathrm{uf}\,F_y(A)\,\mathrm{uf}^{-1}$.

Let $k=0$.
By (op7) and (ga1),
\begin{align*}
\left(Q(A-\mathrm{Acl})J^{*}f\right)(z')&=L^{-3}\int_{B(z')}\mathrm{d}x\,U_A(z',c)\,\mathrm{uf}(z')\,\mathrm{uf}(c)^{-1}f(x)=\mathrm{uf}(z')\left(Q(A)J^{*}\,\mathrm{uf}^{-1}f\right)(z'),\\
\left(JQ^T\left(-(A-\mathrm{Acl})\right)v'\right)(x)&=U_A(z',c)^{-1}\,\mathrm{uf}(c)\,\mathrm{uf}(z')^{-1}\,v'(z')=\mathrm{uf}(c)\left(JQ^T(-A)\,(\mathrm{uf}^1)^{-1}v'\right)(x),
\end{align*}
and the product of the two is $JP(A-\mathrm{Acl})J^{*}=\mathrm{uf}\,JP(A)J^{*}\,\mathrm{uf}^{-1}$.
As $i\gamma_5y$ and $\mu$ commute with $\mathrm{uf}$,
\begin{gather*}
\operatorname{Pen}_y(A-\mathrm{Acl})=b_L\,\mathrm{uf}\,JP(A)J^{*}\,\mathrm{uf}^{-1}+i\gamma_5y=\mathrm{uf}\operatorname{Pen}_y(A)\,\mathrm{uf}^{-1},\\
F_y(A-\mathrm{Acl})=\mathrm{uf}\,D_A\,\mathrm{uf}^{-1}+\mu+\mathrm{uf}\operatorname{Pen}_y(A)\,\mathrm{uf}^{-1}=\mathrm{uf}\,F_y(A)\,\mathrm{uf}^{-1}.
\end{gather*}

The H\"older quotient: let $(x,x')$ be an ordered pair of distinct points of $\mathrm{Lab}\times O_{\mathrm{cub}}$ (among these are the short pairs and, on $T_k$, the pairs of $\mathrm{Pzt}_k$).
By (ga1),
\[
\left(\delta_{\alpha,A-\mathrm{Acl}}(\mathrm{uf}\,f)\right)(x,x')=d_X(x,x')^{-\alpha}\left(U_A(x,x')\,\mathrm{uf}(c)\,\mathrm{uf}(c')^{-1}\,\mathrm{uf}(c')\,f(x')-\mathrm{uf}(c)\,f(x)\right)=\mathrm{uf}(c)\left(\delta_{\alpha,A}f\right)(x,x').
\]
For a pair inside a cluster, $c'=c$ and both phases equal $1$.

\emph{The example.}
Let $\boldsymbol{\lambda}$ be a real scalar.
On every bond $(c,c')$ one has $\mathrm{e}^{ie_k\boldsymbol{\lambda}(c')}=\mathrm{e}^{ie_k\boldsymbol{\lambda}(c)}\exp\left\{ie_k\eta\,(d\boldsymbol{\lambda})(c,c')\right\}$; moreover $d\,d\boldsymbol{\lambda}=0$ and $\left|\mathrm{e}^{ie_k\boldsymbol{\lambda}}\right|=1$.

\textbf{(b)}
Write $\mathrm{mph}(c)=rc+\mathrm{tv}$ with $r=\operatorname{diag}(r_1,r_2,r_3)\in\{\pm1\}^3$; then $r^{-1}=r$.
On the spinors on the unit lattice and on the $L$-lattice, $\mathrm{mS}$ acts by $(\mathrm{mS}\,v)(z)=S\,v(\mathrm{mph}^{-1}z)$ (Section~G1.2).

(mo1)\enspace The map $\mathrm{mph}$ is a bijection $\mathrm{Lab}_j\to\mathrm{Lab}'_j$ for $0\le j\le k+1$, with $\mathrm{mph}(c+Rn)=\mathrm{mph}(c)+Rrn$ and $|\mathrm{mph}\,c-\mathrm{mph}\,c'|=|c-c'|$.
Hence $\mathrm{mph}\left(\mathrm{Bl}_j(z)\right)=\mathrm{Bl}'_j(\mathrm{mph}\,z)$, $\mathrm{cen}'_j(\mathrm{mph}\,c)=\mathrm{mph}\left(\mathrm{cen}_j(c)\right)$ and $\mathrm{mph}\left(B(z')\right)=B'(\mathrm{mph}\,z')$.
Since $\boldsymbol{\sigma}^{rn}=\boldsymbol{\sigma}^n$, the functions $\mathrm{mS}\,f$ and $\mathrm{mS}\,v$ are spinors on $\mathbb{T}'$, and $\mathrm{mS}\,A$ is a connection on $\mathbb{T}'$.

(mo2)\enspace For $a\in\{\pm e_\mu\}$ one has $\mathrm{mph}(c+\eta a)=\mathrm{mph}(c)+\eta\,ra$ and $ra=r_\mu a$; thus bonds are mapped onto bonds.
The coordinate $\mu$ of $\mathrm{mph}\,c$ depends only on the coordinate $\mu$ of $c$.
Therefore $\mathrm{mph}$ maps $\boldsymbol{\Gamma}^{\pi}(c,c')$ bond by bond, in the order and with the orientation of the path, onto $\boldsymbol{\Gamma}^{\pi}(\mathrm{mph}\,c,\mathrm{mph}\,c')$, with the same $\pi$ (compare the last paragraph of \cite[Definition~16.1]{QEDI}).
As $(\mathrm{mS}\,A)(\mathrm{mph}\,\tilde c,\mathrm{mph}\,c'')=A(\tilde c,c'')$ on bonds, it follows that $(\mathrm{mS}\,A)\left(\boldsymbol{\Gamma}^{\pi}(\mathrm{mph}\,c,\mathrm{mph}\,c')\right)=A\left(\boldsymbol{\Gamma}^{\pi}(c,c')\right)$ and $(\tau\,\mathrm{mS}\,A)(\mathrm{mph}\,c,\mathrm{mph}\,c')=(\tau A)(c,c')$.
Together with (mo1) this gives
\begin{equation}\tag{mo5}
\begin{aligned}
&u_{\mathrm{mS}\,A}(\mathrm{mph}\,c,\mathrm{mph}\,c_3)=u_A(c,c_3),&&U_{\mathrm{mS}\,A}(\mathrm{mph}\,c,\mathrm{mph}\,c')=U_A(c,c'),\\
&\varphi_{\mathrm{mS}\,A}(\mathrm{mph}\,z,\mathrm{mph}\,c)=\varphi_A(z,c),&&\varphi^{+}_{\mathrm{mS}\,A}(\mathrm{mph}\,z',\mathrm{mph}\,c)=\varphi^{+}_A(z',c).
\end{aligned}
\end{equation}

(mo3)\enspace On fermion sites, $\mathrm{mph}(c,g)=(\mathrm{mph}\,c,rg)$ defines a bijection $\mathrm{Lab}\times O_{\mathrm{cub}}\to\mathrm{Lab}'\times O_{\mathrm{cub}}$ (the left multiplication by $r$ permutes $O_{\mathrm{cub}}$), and this bijection preserves the weight $\eta^3/48$.
One has $(\mathrm{mph}\,x)^{(i)}=(\mathrm{mph}\,c,rgs_i)=\mathrm{mph}\left(x^{(i)}\right)$ for $i\le2$, and
\[
(\mathrm{mph}\,x)^{(3)}=\left(\mathrm{mph}\,c+\eta\,rge_1,\ rgs_3\right)=\mathrm{mph}\left(c+\eta\,ge_1,\ gs_3\right)=\mathrm{mph}\left(x^{(3)}\right).
\]
Moreover $\hat d_i(\mathrm{mph}\,x)=-rg\hat n_i=r\,\hat d_i(x)$, and $d_X(\mathrm{mph}\,x,\mathrm{mph}\,x')=d_X(x,x')$ by (I.1.5), since $|\mathrm{mph}\,c-\mathrm{mph}\,c'|=|c-c'|$.

(mo4)\enspace By (I.4.1) and (mo3),
\[
S\,M_i(x)\,S^{-1}=\frac{w_i}{2}\left(\rho-S\,\gamma\cdot\hat d_i(x)\,S^{-1}\right)=\frac{w_i}{2}\left(\rho-\gamma\cdot\left(r\,\hat d_i(x)\right)\right)=M_i(\mathrm{mph}\,x).
\]
This is the relation (I.m4) for a morphism.

The operator $D$: by (op1), (mo3), (mo4), (mo5), and since $(\mathrm{mS}\,f)(\mathrm{mph}\,\tilde x)=S\,f(\tilde x)$,
\begin{align*}
&\left(D_{\mathrm{mS}\,A}\,\mathrm{mS}\,f\right)(\mathrm{mph}\,x)\\
&\qquad=-\eta^{-1}\sum_{i=0}^{3}M_i(\mathrm{mph}\,x)\left[U_{\mathrm{mS}\,A}\left(\mathrm{mph}\,x,\mathrm{mph}\,x^{(i)}\right)S\,f\left(x^{(i)}\right)-S\,f(x)\right]\\
&\qquad=-\eta^{-1}\sum_{i=0}^{3}S\,M_i(x)\,S^{-1}S\left[U_A\left(x,x^{(i)}\right)f\left(x^{(i)}\right)-f(x)\right]\\
&\qquad=S\,(D_Af)(x)=\left(\mathrm{mS}\,D_Af\right)(\mathrm{mph}\,x).
\end{align*}
Let $k\ge1$.
In (op3), at the site $\mathrm{mph}\,z$ of $\mathbb{T}'$, substitute $\tilde x=\mathrm{mph}\,x$; by (mo1) and (mo3), $x$ then ranges over $\Delta_z\times O_{\mathrm{cub}}$, with the same weights.
By (mo5),
\begin{align*}
\left(Q_k(\mathrm{mS}\,A)J^{*}\,\mathrm{mS}\,f\right)(\mathrm{mph}\,z)&=\int_{\Delta_z}\mathrm{d}x\,\varphi_A(z,c)\,S\,f(x)=S\left(Q_k(A)J^{*}f\right)(z)=\left(\mathrm{mS}\,Q_k(A)J^{*}f\right)(\mathrm{mph}\,z),\\
\left(JQ_k^T(-\mathrm{mS}\,A)\,\mathrm{mS}\,v\right)(\mathrm{mph}\,x)&=\varphi_A(z,c)^{-1}\,S\,v(z)=\left(\mathrm{mS}\,JQ_k^T(-A)\,v\right)(\mathrm{mph}\,x),\\
\left(Q(\mathrm{mS}\,A)\,\mathrm{mS}\,v\right)(\mathrm{mph}\,z')&=L^{-3}\sum_{z\in B(z')}U_A(z',z)\,S\,v(z)=\left(\mathrm{mS}\,Q(A)\,v\right)(\mathrm{mph}\,z'),\\
\left(Q^T(-\mathrm{mS}\,A)\,\mathrm{mS}\,v'\right)(\mathrm{mph}\,z)&=U_A(z',z)^{-1}\,S\,v'(z')=\left(\mathrm{mS}\,Q^T(-A)\,v'\right)(\mathrm{mph}\,z).
\end{align*}
Composing these four relations:
\begin{align*}
JP_k(\mathrm{mS}\,A)J^{*}\,\mathrm{mS}&=JQ_k^T(-\mathrm{mS}\,A)\,Q_k(\mathrm{mS}\,A)J^{*}\,\mathrm{mS}=JQ_k^T(-\mathrm{mS}\,A)\,\mathrm{mS}\,Q_k(A)J^{*}=\mathrm{mS}\,JP_k(A)J^{*},\\
P(\mathrm{mS}\,A)\,\mathrm{mS}&=Q^T(-\mathrm{mS}\,A)\,Q(\mathrm{mS}\,A)\,\mathrm{mS}=Q^T(-\mathrm{mS}\,A)\,\mathrm{mS}\,Q(A)=\mathrm{mS}\,P(A),\\
JP_{k+1}(\mathrm{mS}\,A)J^{*}\,\mathrm{mS}&=JQ_k^T(-\mathrm{mS}\,A)\,P(\mathrm{mS}\,A)\,Q_k(\mathrm{mS}\,A)J^{*}\,\mathrm{mS}=JQ_k^T(-\mathrm{mS}\,A)\,P(\mathrm{mS}\,A)\,\mathrm{mS}\,Q_k(A)J^{*}\\
&=JQ_k^T(-\mathrm{mS}\,A)\,\mathrm{mS}\,P(A)\,Q_k(A)J^{*}=\mathrm{mS}\,JP_{k+1}(A)J^{*}.
\end{align*}
Let $k=0$.
By (op7), (mo1), (mo3), (mo5),
\begin{align*}
\left(Q(\mathrm{mS}\,A)J^{*}\,\mathrm{mS}\,f\right)(\mathrm{mph}\,z')&=L^{-3}\int_{B(z')}\mathrm{d}x\,U_A(z',c)\,S\,f(x)=\left(\mathrm{mS}\,Q(A)J^{*}f\right)(\mathrm{mph}\,z'),\\
\left(JQ^T(-\mathrm{mS}\,A)\,\mathrm{mS}\,v'\right)(\mathrm{mph}\,x)&=U_A(z',c)^{-1}\,S\,v'(z')=\left(\mathrm{mS}\,JQ^T(-A)\,v'\right)(\mathrm{mph}\,x),
\end{align*}
and for their product
\[
JP(\mathrm{mS}\,A)J^{*}\,\mathrm{mS}=JQ^T(-\mathrm{mS}\,A)\,Q(\mathrm{mS}\,A)J^{*}\,\mathrm{mS}=JQ^T(-\mathrm{mS}\,A)\,\mathrm{mS}\,Q(A)J^{*}=\mathrm{mS}\,JP(A)J^{*}.
\]
The H\"older quotient: for every ordered pair of distinct points of $\mathrm{Lab}\times O_{\mathrm{cub}}$, (mo3) and (mo5) give
\[
\left(\delta_{\alpha,\mathrm{mS}\,A}\,\mathrm{mS}\,f\right)(\mathrm{mph}\,x,\mathrm{mph}\,x')=d_X(x,x')^{-\alpha}\left(U_A(x,x')\,S\,f(x')-S\,f(x)\right)=S\left(\delta_{\alpha,A}f\right)(x,x').
\]
The operator $F_y$: every constant matrix $\mathrm{Mx}$ satisfies $\mathrm{Mx}\,\mathrm{mS}=\mathrm{mS}\left(S^{-1}\mathrm{Mx}\,S\right)$.
Apply (mx4) and (mx6) to $S^{-1}$, which lies in $\mathrm{Sp}_3$ if and only if $S$ does:
\[
S^{-1}\boldsymbol{\alpha}_{y'}S=\boldsymbol{\alpha}_y,\qquad S^{-1}\boldsymbol{\beta}_{y'}S=\boldsymbol{\beta}_y,\qquad S^{-1}\left(i\gamma_5y'\right)S=i\gamma_5y,
\]
where $y'=y$ for $S\in\mathrm{Sp}_3$ and $y'=-y$ for $S\notin\mathrm{Sp}_3$.
If $y$ is neutral, the three matrices are real numbers and $-y=y$, so that $y'=y$ in both cases.
Hence, for $k\ge1$,
\begin{align*}
F_{y'}(\mathrm{mS}\,A)\,\mathrm{mS}&=D_{\mathrm{mS}\,A}\,\mathrm{mS}+\mu\,\mathrm{mS}+\boldsymbol{\alpha}_{y'}\,JP_k(\mathrm{mS}\,A)J^{*}\,\mathrm{mS}+\boldsymbol{\beta}_{y'}\,JP_{k+1}(\mathrm{mS}\,A)J^{*}\,\mathrm{mS}\\
&=\mathrm{mS}\,D_A+\mathrm{mS}\,\mu+\boldsymbol{\alpha}_{y'}\,\mathrm{mS}\,JP_k(A)J^{*}+\boldsymbol{\beta}_{y'}\,\mathrm{mS}\,JP_{k+1}(A)J^{*}=\mathrm{mS}\,F_y(A),
\end{align*}
and for $k=0$
\[
F_{y'}(\mathrm{mS}\,A)\,\mathrm{mS}=\mathrm{mS}\,D_A+\mathrm{mS}\,\mu+b_L\,\mathrm{mS}\,JP(A)J^{*}+i\gamma_5y'\,\mathrm{mS}=\mathrm{mS}\,F_y(A).
\]

\textbf{(c)}
\emph{The operator $D_A$.}
Write $\langle f',D_Af\rangle^{\mathrm{ch}}$ by means of (op1), apply (kn5) to the $i$-th integral (with $x$ replaced by $x^{(i)}$, and $\left(x^{(i)}\right)^{(i)}=x$), and use $\mathrm{vv}'\cdot(\mathrm{Mx}\,\mathrm{vv})=(\mathrm{Mx}^T\mathrm{vv}')\cdot\mathrm{vv}$:
\begin{align*}
\langle f',D_Af\rangle^{\mathrm{ch}}&=\eta^{-1}\,\frac{\rho}{2}\,\mathrm{sw}\int\mathrm{d}x\,f'(x)\cdot f(x)-\eta^{-1}\sum_{i=0}^{3}\int\mathrm{d}x\,f'(x)\cdot M_i(x)\,U_A\left(x,x^{(i)}\right)f\left(x^{(i)}\right)\\
&=\eta^{-1}\,\frac{\rho}{2}\,\mathrm{sw}\int\mathrm{d}x\,f'(x)\cdot f(x)-\eta^{-1}\sum_{i=0}^{3}\int\mathrm{d}x\,\left[U_A\left(x^{(i)},x\right)M_i\left(x^{(i)}\right)^Tf'\left(x^{(i)}\right)\right]\cdot f(x).
\end{align*}
Here $U_A(x^{(i)},x)=1=U_{-A}(x,x^{(i)})$ for $i\le2$, and $U_A(x^{(3)},x)=u_A(c_3,c)=u_{-A}(c,c_3)=U_{-A}(x,x^{(3)})$ by (ph0); moreover $M_i(x^{(i)})^T=\mathsf{C}^{-1}M_i(x)\,\mathsf{C}$ by (hm2).
Therefore
\[
\left(D_A^Tf'\right)(x)=\eta^{-1}\,\frac{\rho}{2}\,\mathrm{sw}\,f'(x)-\eta^{-1}\sum_{i=0}^{3}\mathsf{C}^{-1}M_i(x)\,\mathsf{C}\,U_{-A}\left(x,x^{(i)}\right)f'\left(x^{(i)}\right)=\left(\mathsf{C}^{-1}D_{-A}\,\mathsf{C}f'\right)(x),
\]
the last equality being the second form of (op1) for $D_{-A}$ (the scalar $\frac{\rho}{2}\,\mathrm{sw}$ commutes with $\mathsf{C}$).

\emph{Averages.}
Let $k\ge1$.
By Section~I.3.3 and (ph0), $\left(Q_k^T(A)v\right)(c)=\varphi_{-A}(z,c)^{-1}v(z)=\varphi_A(z,c)\,v(z)$; with (op3),
\begin{align*}
\langle v,Q_k(A)J^{*}f\rangle_1&=\sum_zv(z)\cdot\eta^3\sum_{c\in\Delta_z}\varphi_A(z,c)\,\frac{1}{48}\sum_gf(c,g)\\
&=\eta^3\sum_c\frac{1}{48}\sum_g\left[\varphi_A([c],c)\,v([c])\right]\cdot f(c,g)=\langle JQ_k^T(A)\,v,f\rangle^{\mathrm{ch}}.
\end{align*}
By (op4) and $\left(Q^T(A)v'\right)(z)=U_A(z',z)\,v'(z')$,
\[
\langle v',Q(A)\,v\rangle_L=L^3\sum_{z'}v'(z')\cdot L^{-3}\sum_{z\in B(z')}U_A(z',z)\,v(z)=\sum_z\left[U_A(z',z)\,v'(z')\right]\cdot v(z)=\langle Q^T(A)\,v',v\rangle_1 .
\]
Let $k=0$.
By (op7),
\begin{align*}
\langle v',Q(A)J^{*}f\rangle_L&=L^3\sum_{z'}v'(z')\cdot L^{-3}\sum_{c\in B(z')}U_A(z',c)\,\frac{1}{48}\sum_gf(c,g)\\
&=\sum_c\frac{1}{48}\sum_g\left[U_A(z',c)\,v'(z')\right]\cdot f(c,g)=\langle JQ^T(A)\,v',f\rangle^{\mathrm{ch}}.
\end{align*}

\emph{Idempotents.}
Recall that $(O_1O_2)^T=O_2^TO_1^T$ and $(O^T)^T=O$.
The relations obtained for the averages, at $A$ and at $-A$, give $\left(JQ_k^T(-A)\right)^T=Q_k(-A)J^{*}$, $Q^T(-A)^T=Q(-A)$ and $\left(JQ^T(-A)\right)^T=Q(-A)J^{*}$.
Hence
\begin{align*}
\left(JP_k(A)J^{*}\right)^T&=\left(Q_k(A)J^{*}\right)^T\left(JQ_k^T(-A)\right)^T=JQ_k^T(A)\,Q_k(-A)J^{*}=JP_k(-A)J^{*},\\
P(A)^T&=Q(A)^T\,Q^T(-A)^T=Q^T(A)\,Q(-A)=P(-A),\\
\left(JP_{k+1}(A)J^{*}\right)^T&=\left(Q_k(A)J^{*}\right)^TP(A)^T\left(JQ_k^T(-A)\right)^T=JQ_k^T(A)\,P(-A)\,Q_k(-A)J^{*}=JP_{k+1}(-A)J^{*},
\end{align*}
and for $k=0$
\[
\left(JP(A)J^{*}\right)^T=\left(Q(A)J^{*}\right)^T\left(JQ^T(-A)\right)^T=JQ^T(A)\,Q(-A)J^{*}=JP(-A)J^{*}.
\]

\emph{The operator $F_y$.}
The transpose of the multiplication by a constant matrix $\mathrm{Mx}$ is the multiplication by $\mathrm{Mx}^T$.
By (mx1) and (mx6),
\[
\boldsymbol{\alpha}_y^T=\mathsf{C}^{-1}\boldsymbol{\alpha}_{-y}\mathsf{C},\qquad\boldsymbol{\beta}_y^T=\mathsf{C}^{-1}\boldsymbol{\beta}_{-y}\mathsf{C},\qquad\left(i\gamma_5y\right)^T=\mathsf{C}^{-1}\left(i\gamma_5(-y)\right)\mathsf{C}.
\]
With (sc),
\begin{align*}
\left(\boldsymbol{\alpha}_y\,JP_k(A)J^{*}\right)^T&=JP_k(-A)J^{*}\,\mathsf{C}^{-1}\boldsymbol{\alpha}_{-y}\mathsf{C}=\mathsf{C}^{-1}\boldsymbol{\alpha}_{-y}\,JP_k(-A)J^{*}\,\mathsf{C},\\
\left(\boldsymbol{\beta}_y\,JP_{k+1}(A)J^{*}\right)^T&=JP_{k+1}(-A)J^{*}\,\mathsf{C}^{-1}\boldsymbol{\beta}_{-y}\mathsf{C}=\mathsf{C}^{-1}\boldsymbol{\beta}_{-y}\,JP_{k+1}(-A)J^{*}\,\mathsf{C},\\
\left(b_L\,JP(A)J^{*}\right)^T&=b_L\,JP(-A)J^{*}=\mathsf{C}^{-1}b_L\,JP(-A)J^{*}\,\mathsf{C}.
\end{align*}
Adding $D_A^T=\mathsf{C}^{-1}D_{-A}\mathsf{C}$ and $\mu$, one obtains $F_y(A)^T=\mathsf{C}^{-1}F_{-y}(-A)\,\mathsf{C}$, for $k\ge1$ and for $k=0$.

\textbf{(d)}
For an operator $O$ let $\overline{O}$ denote the operator $f\mapsto\overline{O\,\overline{f}}$; together with $f$, also $\overline{f}$ is a spinor, as $\boldsymbol{\sigma}$ is real.
One has $\overline{O_1O_2}=\overline{O_1}\;\overline{O_2}$, and the conjugate of the multiplication by $\mathrm{Mx}$ is the multiplication by $\overline{\mathrm{Mx}}$.

The operator $D$: by the second form of (op1), by (hm3), and by $\overline{U_{\overline{A}}(x,x^{(i)})}=U_{-A}(x,x^{(i)})$ (ph0),
\[
\left(\overline{D_{\overline{A}}}\,f\right)(x)=\eta^{-1}\,\frac{\rho}{2}\,\mathrm{sw}\,f(x)-\eta^{-1}\sum_{i=0}^{3}\varkappa^{-1}M_i(x)\,\varkappa\,U_{-A}\left(x,x^{(i)}\right)f\left(x^{(i)}\right)=\left(\varkappa^{-1}D_{-A}\,\varkappa f\right)(x).
\]
The averages: in (op3), (op4), (op7) every coefficient is a positive number times a phase $\mathrm{pf}_A^{\pm1}$, and $\overline{\mathrm{pf}_{\overline{A}}^{\pm1}}=\mathrm{pf}_{-A}^{\pm1}$ by (ph0).
Hence $\overline{O(\overline{A})}=O(-A)$ for $O=Q_kJ^{*}$, $JQ_k^T(-\,\cdot\,)$, $Q$, $Q^T(-\,\cdot\,)$, and, at $k=0$, for $O=QJ^{*}$, $JQ^T(-\,\cdot\,)$; consequently the same holds for their products $JP_kJ^{*}$, $P$, $JP_{k+1}J^{*}$, $JPJ^{*}$.
By (sc), $O(-A)=\varkappa^{-1}O(-A)\,\varkappa$.

The operator $F_y$: by (mx2) and (mx6), $\overline{\boldsymbol{\alpha}_y}=\varkappa^{-1}\boldsymbol{\alpha}_y\varkappa$, $\overline{\boldsymbol{\beta}_y}=\varkappa^{-1}\boldsymbol{\beta}_y\varkappa$ and $\overline{i\gamma_5y}=\varkappa^{-1}\left(i\gamma_5y\right)\varkappa$, while the numbers $\mu$ and $b_L$ are real.
With (sc),
\begin{align*}
\overline{\boldsymbol{\alpha}_y\,JP_k(\overline{A})J^{*}}&=\varkappa^{-1}\boldsymbol{\alpha}_y\varkappa\,JP_k(-A)J^{*}=\varkappa^{-1}\boldsymbol{\alpha}_y\,JP_k(-A)J^{*}\,\varkappa,\\
\overline{\boldsymbol{\beta}_y\,JP_{k+1}(\overline{A})J^{*}}&=\varkappa^{-1}\boldsymbol{\beta}_y\varkappa\,JP_{k+1}(-A)J^{*}=\varkappa^{-1}\boldsymbol{\beta}_y\,JP_{k+1}(-A)J^{*}\,\varkappa,\\
\overline{b_L\,JP(\overline{A})J^{*}}&=b_L\,JP(-A)J^{*}=\varkappa^{-1}b_L\,JP(-A)J^{*}\,\varkappa .
\end{align*}
So $\overline{F_y(\overline{A})}=\varkappa^{-1}F_y(-A)\,\varkappa$, with the same $y$, for $k\ge1$ and for $k=0$.

\textbf{(e)}
Let $A$ be real.
Then $|\mathrm{pf}_A|=1$ and $\overline{\mathrm{pf}_A}=\mathrm{pf}_A^{-1}=\mathrm{pf}_{-A}$ (ph0).

The operator $D_A$: by (op1), by (kn5) applied to the $i$-th integral, and by $\overline{\mathrm{vv}'}\cdot(\mathrm{Mx}\,\mathrm{vv})=\overline{\mathrm{Mx}^{*}\mathrm{vv}'}\cdot\mathrm{vv}$,
\[
(f',D_Af)^{\mathrm{ch}}=\eta^{-1}\,\frac{\rho}{2}\,\mathrm{sw}\int\mathrm{d}x\,\overline{f'(x)}\cdot f(x)-\eta^{-1}\sum_{i=0}^{3}\int\mathrm{d}x\,\overline{\left[\overline{U_A\left(x^{(i)},x\right)}\,M_i\left(x^{(i)}\right)^{*}f'\left(x^{(i)}\right)\right]}\cdot f(x).
\]
Here $\overline{U_A(x^{(i)},x)}=U_A(x^{(i)},x)^{-1}=U_A(x,x^{(i)})$: for $i\le2$ both sides equal $1$, and for $i=3$ one has $u_A(c_3,c)^{-1}=u_A(c,c_3)$.
Moreover $M_i(x^{(i)})^{*}=M_i(x^{(i)})=\gamma_5M_i(x)\gamma_5$ by (hm4).
Hence, as $\gamma_5^2=1$,
\[
\left(D_A^{*}f'\right)(x)=\eta^{-1}\,\frac{\rho}{2}\,\mathrm{sw}\,f'(x)-\eta^{-1}\sum_{i=0}^{3}\gamma_5M_i(x)\gamma_5\,U_A\left(x,x^{(i)}\right)f'\left(x^{(i)}\right)=\left(\gamma_5D_A\gamma_5f'\right)(x).
\]
The averages: for $k\ge1$,
\begin{align*}
(v,Q_k(A)J^{*}f)_1&=\sum_z\overline{v(z)}\cdot\eta^3\sum_{c\in\Delta_z}\varphi_A(z,c)\,\frac{1}{48}\sum_gf(c,g)\\
&=\eta^3\sum_c\frac{1}{48}\sum_g\overline{\varphi_A([c],c)^{-1}v([c])}\cdot f(c,g)=\left(JQ_k^T(-A)\,v,f\right)^{\mathrm{ch}},\\
(v',Q(A)\,v)_L&=\sum_z\overline{U_A(z',z)^{-1}v'(z')}\cdot v(z)=\left(Q^T(-A)\,v',v\right)_1 ,
\end{align*}
and for $k=0$
\[
(v',Q(A)J^{*}f)_L=\sum_c\frac{1}{48}\sum_g\overline{U_A(z',c)^{-1}v'(z')}\cdot f(c,g)=\left(JQ^T(-A)\,v',f\right)^{\mathrm{ch}},
\]
that is, $\left(Q(A)J^{*}\right)^{*}=JQ^T(-A)$.

The projections:
\begin{align*}
\left(JP_k(A)J^{*}\right)^{*}&=\left(Q_k(A)J^{*}\right)^{*}\left(JQ_k^T(-A)\right)^{*}=JQ_k^T(-A)\,Q_k(A)J^{*}=JP_k(A)J^{*},\\
P(A)^{*}&=Q(A)^{*}\,Q^T(-A)^{*}=Q^T(-A)\,Q(A)=P(A),\\
\left(JP_{k+1}(A)J^{*}\right)^{*}&=JQ_k^T(-A)\,P(A)^{*}\,Q_k(A)J^{*}=JP_{k+1}(A)J^{*},\\
\left(JP(A)J^{*}\right)^{*}&=JQ^T(-A)\,Q(A)J^{*}=JP(A)J^{*}.
\end{align*}
An idempotent of (f0) that is self-adjoint is an orthogonal projection.

The operator $F_y$: by (mx3) and (mx6),
\[
\boldsymbol{\alpha}_y^{*}=\boldsymbol{\alpha}_{-y}=\gamma_5\boldsymbol{\alpha}_{-y}\gamma_5,\qquad\boldsymbol{\beta}_y^{*}=\boldsymbol{\beta}_{-y}=\gamma_5\boldsymbol{\beta}_{-y}\gamma_5,\qquad\left(i\gamma_5y\right)^{*}=i\gamma_5(-y)=\gamma_5\left(i\gamma_5(-y)\right)\gamma_5 .
\]
With (sc) and the projections just treated,
\begin{align*}
\left(\boldsymbol{\alpha}_y\,JP_k(A)J^{*}\right)^{*}&=JP_k(A)J^{*}\,\boldsymbol{\alpha}_{-y}=\gamma_5\boldsymbol{\alpha}_{-y}\,JP_k(A)J^{*}\,\gamma_5,\\
\left(\boldsymbol{\beta}_y\,JP_{k+1}(A)J^{*}\right)^{*}&=JP_{k+1}(A)J^{*}\,\boldsymbol{\beta}_{-y}=\gamma_5\boldsymbol{\beta}_{-y}\,JP_{k+1}(A)J^{*}\,\gamma_5,\\
\left(b_L\,JP(A)J^{*}\right)^{*}&=b_L\,JP(A)J^{*}=\gamma_5\,b_L\,JP(A)J^{*}\,\gamma_5,
\end{align*}
and $\mu$ is real.
Adding $D_A^{*}=\gamma_5D_A\gamma_5$, one obtains $F_y(A)^{*}=\gamma_5F_{-y}(A)\,\gamma_5$, for $k\ge1$ and for $k=0$.

All the relations above are identities.
They hold for every block torus, every $e_k>0$, every $\mu\ge0$, every $y\in Y_k$ and every complex connection (in (e): every real connection).
\end{proof}

\subsection*{3. Zero field on every block torus}

\begin{proof}[Proof of Proposition~2.3]
\emph{The translation $\mathrm{sft}$.}
Suppose that $t=(\eta/2)(1,1,1)$, and let $\mathbb{T}_0:=(0,R,\boldsymbol{\sigma})$.
For a function $f$ on $\mathrm{Lab}\times O_{\mathrm{cub}}$ define the function $\mathrm{sft}\,f$ on $\eta\mathbb{Z}^3\times O_{\mathrm{cub}}$ by
\[
(\mathrm{sft}\,f)(c,g):=f\left(c+\frac{L}{2}(1,1,1),\,g\right)\qquad\left(c\in\eta\mathbb{Z}^3\right).
\]
Because $L=L^{k+1}\eta$, every $j$ with $0\le j\le k+1$ satisfies
\[
t_j-\frac{L}{2}(1,1,1)=L^j\eta\cdot\frac{1-L^{k+1-j}}{2}\,(1,1,1),
\]
and the integer $1-L^{k+1-j}$ is even.
Hence, for each of these $j$, the map $c\mapsto c-\frac{L}{2}(1,1,1)$ carries $\mathrm{Lab}_j$ onto $L^j\eta\mathbb{Z}^3$, and it carries the block $\mathrm{Bl}_j(z)$ onto the block of $\mathbb{T}_0$ with centre $z-\frac{L}{2}(1,1,1)$.
The translation moves $c$ only; the label $g$ stays fixed.
This has the following consequences.
\begin{itemize}
\item[(sf1)]
The map $\mathrm{sft}$ is a bijection of the spinors on $\mathbb{T}$ onto the spinors on $\mathbb{T}_0$, and
\[
|\mathrm{sft}\,f|_{\mathrm{su}}=|f|_{\mathrm{su}},\qquad\|\mathrm{sft}\,f\|_{L^2}=\|f\|_{L^2},\qquad\left(\mathrm{sft}\,f',\mathrm{sft}\,f\right)^{\mathrm{ch}}=(f',f)^{\mathrm{ch}},
\]
since the translate of a period cell is a period cell.
A spinor $f$ is supported over the orbit of $\Delta_z$ if and only if $\mathrm{sft}\,f$ is supported over the orbit of $\Delta_{z-\frac{L}{2}(1,1,1)}$.
The distances $d$ and $d_X$, and the relation $\tilde{x}=x^{(i)}$, are invariant under the translation.
\item[(sf2)]
One has $\mathrm{sft}\,F_y(0)=F_y(0)\,\mathrm{sft}$, where the operator $F_y(0)$ on the left-hand side is that of $\mathbb{T}$ and the one on the right-hand side is that of $\mathbb{T}_0$.
Indeed, at $A=0$ the expression (op1) contains $c$ only through $c_3-c=\eta g e_1$.
At $A=0$ all phases are equal to $1$; so, by (op5)--(op7), each of $JP_k(0)J^{*}f$, $JP_{k+1}(0)J^{*}f$ and $JP(0)J^{*}f$ is equal, on the product of a block with $O_{\mathrm{cub}}$, to the mean of $f$ over this product, and the translation maps blocks onto blocks.
Finally, $\mu$, $\boldsymbol{\alpha}_{y}$, $\boldsymbol{\beta}_{y}$ and $i\gamma_5y$ are constants.
\item[(sf3)]
The spinors on $\mathbb{T}_0$ are the elements of $\mathfrak{S}^{\mathrm{ch}}(\eta,R;\theta)=\mathrm{HH}(\eta,R;\theta)$ of Section~G1.4 with $\mathrm{me}=\eta$, where $\theta=0$ if $\boldsymbol{\sigma}=-1$ and $\theta=(\pi,\pi,\pi)$ if $\boldsymbol{\sigma}=+1$, so that $-\exp(i\theta_\mu)=\boldsymbol{\sigma}$; and $R$ is a multiple of $L$.
On $\mathbb{T}_0$ one has $\mathrm{Lab}_j=L^j\eta\mathbb{Z}^3$, and the blocks $\mathrm{Bl}_j(z)$ of $\mathbb{T}_0$ are the blocks of level $j$ of Section~G1.4.
The measure $\int\mathrm{d}x$, the pairing $(\cdot,\cdot)^{\mathrm{ch}}$, the distance $d_X$ and the sup distance modulo $R\mathbb{Z}^3$ are those of Section~G1.4, and $D_0$ is the operator $D$ of Section~G1.4.
\end{itemize}

\emph{(a)}
Let $k\ge1$, and consider first the torus $\mathbb{T}_0$.
By (op5), (op6) at $A=0$ and by (bl1), the spinor $JP_k(0)J^{*}f$ is constant on $\Delta_z\times O_{\mathrm{cub}}$, with the value
\[
\int_{\Delta_z}\mathrm{d}x\,f(x)=\eta^3\sum_{c\in\Delta_z}\frac{1}{48}\sum_{g\in O_{\mathrm{cub}}}f(c,g),
\]
which is the mean of $f$ over the $48\eta^{-3}$ sites of $\Delta_z\times O_{\mathrm{cub}}$.
Likewise, $JP_{k+1}(0)J^{*}f$ is constant on $\mathrm{Bl}_{k+1}(z')\times O_{\mathrm{cub}}$, with the value
\[
L^{-3}\int_{\mathrm{Bl}_{k+1}(z')}\mathrm{d}x\,f(x),
\]
the mean of $f$ over the $48L^3\eta^{-3}$ sites of this set.
All sites carry the same weight $\eta^3/48$.
Therefore the orthogonal projection, with respect to $(\cdot,\cdot)^{\mathrm{ch}}$, onto the spinors that are constant on every set $E\times O_{\mathrm{cub}}$, $E$ ranging over the blocks of one side, replaces $f$ on each set $E\times O_{\mathrm{cub}}$ by its mean over that set.
Consequently $JP_k(0)J^{*}$ and $JP_{k+1}(0)J^{*}$ are the two projections of Proposition~1.6, and, by (sf3), the operator $F_y(0)$ of (op8) on $\mathbb{T}_0$ is the operator $F_y(0)$ of Proposition~1.6 on $\mathrm{HH}(\eta,R;\theta)$.
By Proposition~1.6\,(a) it is bijective.
Proposition~1.6\,(c), applied with $(z,z')=(z_1,z_2)$, yields the inequality asserted in (a); note that every spinor on a torus is bounded.
This proves (a) for $t=0$.

Let now $t=(\eta/2)(1,1,1)$.
By (sf1) and (sf2), $F_y(0)$ is bijective on the spinors on $\mathbb{T}$, and $\mathrm{sft}\,S_y(0)f=S_y(0)\,\mathrm{sft}\,f$.
Put $v=S_y(0)f$.
The values $v(x)$ and $v(\tilde{x})$ are the values of $\mathrm{sft}\,v$ at the translated points;
the spinor $\mathrm{sft}\,f$ is supported over the orbit of $\Delta_{z_2-\frac{L}{2}(1,1,1)}$, and $|\mathrm{sft}\,f|_{\mathrm{su}}=|f|_{\mathrm{su}}$;
and the numbers $d(z_1,z_2)$, $d_X(x,\tilde{x})$ are not changed by the translation.
The inequality of (a) on $\mathbb{T}$ is therefore the inequality of (a) on $\mathbb{T}_0$ at the translated points, which has been proved.

\emph{(b)}
Let $k=0$, so that $\eta=1$, and write
\[
\mathrm{Fz}:=F_y(0)=D_0+\mu+b_L\cdot JP(0)J^{*}+i\gamma_5y.
\]

\emph{Step 1 (real part).}
The first claim is the inequality
\begin{equation}\tag{re1}
\operatorname{Re}\,(u,\mathrm{Fz}\,u)^{\mathrm{ch}}\ge\frac{c_{\mathrm{re}}}{L}\,(u,u)^{\mathrm{ch}}\qquad\text{for every spinor $u$ on $\mathbb{T}$.}
\end{equation}
On $\mathbb{T}_0$, the blocks of level one of Section~G1.4 at $\mathrm{me}=1$ are the sets $E=B(z')$, $z'\in L\mathbb{Z}^3$, and $\ell_1=L$.
By (op7) at $A=0$, the spinor $JP(0)J^{*}u$ is constant on $E\times O_{\mathrm{cub}}$, with the value
\[
L^{-3}\sum_{c\in E}\frac{1}{48}\sum_{g\in O_{\mathrm{cub}}}u(c,g),
\]
the mean of $u$ over the $48L^3$ sites of $E\times O_{\mathrm{cub}}$; there it coincides with $P_Eu$.
Hence $b_L\cdot JP(0)J^{*}=\sum_Eb\cdot\ell_1^{-1}P_E$, and $\mathrm{Fz}$ is the operator $\mathrm{Tst}$ of Proposition~1.7 with $\mathrm{me}=1$, $\mathrm{Wd}=\mathbb{Z}^3$ and $b_E=b$.
Proposition~1.7 gives (re1) on $\mathbb{T}_0$.
On $\mathbb{T}$ with $t\ne0$, the relations (sf1), (sf2) and the inequality (re1) on $\mathbb{T}_0$ give
\begin{align*}
\operatorname{Re}\,(u,\mathrm{Fz}\,u)^{\mathrm{ch}}&=\operatorname{Re}\,\left(\mathrm{sft}\,u,\mathrm{sft}\,\mathrm{Fz}\,u\right)^{\mathrm{ch}}=\operatorname{Re}\,\left(\mathrm{sft}\,u,\mathrm{Fz}\,\mathrm{sft}\,u\right)^{\mathrm{ch}}\\
&\ge\frac{c_{\mathrm{re}}}{L}\,\left(\mathrm{sft}\,u,\mathrm{sft}\,u\right)^{\mathrm{ch}}=\frac{c_{\mathrm{re}}}{L}\,(u,u)^{\mathrm{ch}}.
\end{align*}

\emph{Step 2 (weights).}
Let $\mathrm{wf}$ be a real scalar such that
\[
\left|\mathrm{wf}(c)-\mathrm{wf}(\tilde{c})\right|\le a_z\cdot|c-\tilde{c}|_1\qquad\text{for all }c,\tilde{c}\in\mathrm{Lab},
\]
and put
\[
\mathrm{Fw}:=\exp(\mathrm{wf})\,\mathrm{Fz}\,\exp(-\mathrm{wf}),\qquad\mathrm{Rw}:=\mathrm{Fw}-\mathrm{Fz}.
\]
The terms $\mu$ and $i\gamma_5y$ commute with $\exp(\mathrm{wf})$.
The third identity of (op2), taken with $\mathrm{uf}=\exp(\mathrm{wf})$ and $A=0$, together with (op7), gives $\mathrm{Rw}=\mathrm{Rw}^{\mathrm{hop}}+\mathrm{Rw}^{\mathrm{pen}}$, where
\begin{align*}
\left(\mathrm{Rw}^{\mathrm{hop}}u\right)(x)&=-M_3(x)\left(\exp\left(\mathrm{wf}(c)-\mathrm{wf}(c_3)\right)-1\right)u\left(x^{(3)}\right),\\
\left(\mathrm{Rw}^{\mathrm{pen}}u\right)(x)&=b_L\cdot L^{-3}\sum_{\tilde{c}\in B(z')}\left(\exp\left(\mathrm{wf}(c)-\mathrm{wf}(\tilde{c})\right)-1\right)\frac{1}{48}\sum_{\tilde{g}\in O_{\mathrm{cub}}}u(\tilde{c},\tilde{g}),
\end{align*}
where $z'=\mathrm{cen}_1(c)$.
For a complex number $\mathrm{cx}$ one has $\left|\exp(\mathrm{cx})-1\right|\le|\mathrm{cx}|\cdot\exp|\mathrm{cx}|$, and hence $\left|\exp(\mathrm{cx})-1\right|\le|\mathrm{cx}|\cdot\mathrm{e}$ if $|\mathrm{cx}|\le1$.
Moreover $\left|\mathrm{wf}(c)-\mathrm{wf}(c_3)\right|\le a_z\le1$.
Together with (hm5), and with (kn5) for $i=3$, this gives
\[
\left\|\mathrm{Rw}^{\mathrm{hop}}u\right\|_{L^2}^2\le\left(6a_z\cdot\mathrm{e}\right)^2\int\mathrm{d}x\,\left|u\left(x^{(3)}\right)\right|_2^2=\left(6a_z\cdot\mathrm{e}\right)^2\|u\|_{L^2}^2.
\]
If $c$ and $\tilde{c}$ lie in one block $B(z')$, then (bl4) and $a_z\le1/(3L)$ give $\left|\mathrm{wf}(c)-\mathrm{wf}(\tilde{c})\right|\le3(L-1)\cdot a_z\le1$, and therefore
\[
\left|\exp\left(\mathrm{wf}(c)-\mathrm{wf}(\tilde{c})\right)-1\right|\le3L\cdot a_z\cdot\mathrm{e}.
\]
The set $B(z')\times O_{\mathrm{cub}}$ consists of $48L^3$ sites of weight $1/48$, so that its total mass is $L^3$.
The Cauchy--Schwarz inequality on this set gives
\begin{align*}
\left|\left(\mathrm{Rw}^{\mathrm{pen}}u\right)(x)\right|_2&\le b_L\cdot L^{-3}\cdot3L\cdot a_z\cdot\mathrm{e}\int_{B(z')}\mathrm{d}\tilde{x}\,|u(\tilde{x})|_2\\
&\le b_L\cdot L^{-3}\cdot3L\cdot a_z\cdot\mathrm{e}\cdot L^{\frac{3}{2}}\,\|u\|_{L^2,B(z')}.
\end{align*}
Integrating the square over the sites $x$ that lie over $B(z')$, which form a set of mass $L^3$, one obtains
\[
\left\|\mathrm{Rw}^{\mathrm{pen}}u\right\|_{L^2,B(z')}\le b_L\cdot3L\cdot a_z\cdot\mathrm{e}\cdot\|u\|_{L^2,B(z')}=3b\cdot a_z\cdot\mathrm{e}\cdot\|u\|_{L^2,B(z')}.
\]
By (bl9), the union of the $(R/L)^3$ blocks $B(z')=\mathrm{Bl}_1(z')$, $z'$ in a period cell of $\mathrm{Lab}_1$, is a period cell of $\mathrm{Lab}$; adding the squares over these blocks one finds $\left\|\mathrm{Rw}^{\mathrm{pen}}\right\|_{L^2}\le3b\cdot a_z\cdot\mathrm{e}$.
The definition of $a_z$ now gives
\begin{equation}\tag{re2}
\|\mathrm{Rw}\|_{L^2}\le(6+3b)\cdot\mathrm{e}\cdot a_z\le\frac{c_{\mathrm{re}}}{2L}.
\end{equation}
From (re1) and (re2) it follows that
\[
\operatorname{Re}\,(u,\mathrm{Fw}\,u)^{\mathrm{ch}}\ge\left(\frac{c_{\mathrm{re}}}{L}-\frac{c_{\mathrm{re}}}{2L}\right)\|u\|_{L^2}^2,
\]
and hence that $\frac{c_{\mathrm{re}}}{2L}\,\|u\|_{L^2}^2\le\|u\|_{L^2}\cdot\|\mathrm{Fw}\,u\|_{L^2}$.
Thus $\mathrm{Fw}$ is injective, hence bijective because the dimension is finite, and $\left\|\mathrm{Fw}^{-1}\right\|_{L^2}\le2L/c_{\mathrm{re}}$.
For $\mathrm{wf}=0$ this shows that $\mathrm{Fz}$ is bijective.

\emph{Step 3 (the factor $(1+|y|)^{-1}$).}
The operator $R_i$, $(R_iu)(x):=u\left(x^{(i)}\right)$, is unitary on $L^2$ by (kn5), and $D_0=\sum_{i=0}^{3}M_i(\cdot)\,(1-R_i)$ by (op1).
With (hm5) and (co6) this gives
\[
\|D_0\|_{L^2}\le\sum_{i}\frac{w_i}{2}(1+\rho)\cdot2=(1+\rho)\cdot\mathrm{sw}<58.
\]
On each block $B(z')$, the Cauchy--Schwarz inequality on a set of mass $L^3$ gives
\[
\left\|JP(0)J^{*}u\right\|_{L^2,B(z')}^2=L^3\left|L^{-3}\int_{B(z')}\mathrm{d}\tilde{x}\,u(\tilde{x})\right|_2^2\le\|u\|_{L^2,B(z')}^2,
\]
and the sum over the blocks of the period cell of (bl9) gives $\left\|JP(0)J^{*}\right\|_{L^2}\le1$.
Also $\mu\le1$.
Since $\mathrm{Fw}-i\gamma_5y=D_0+\mu+b_L\cdot JP(0)J^{*}+\mathrm{Rw}$, the bound (re2), the relations (co5) and the inequality $L\ge5$ give
\[
\left\|\mathrm{Fw}-i\gamma_5y\right\|_{L^2}\le58+1+\frac{b}{L}+\frac{b}{2L}\le59+\frac{3}{10}\,b\le k_1.
\]
As $\gamma_5$ is unitary, $\|\mathrm{Fw}\,u\|_{L^2}\ge\left(|y|-k_1\right)\|u\|_{L^2}$, and consequently
\[
\left\|\mathrm{Fw}^{-1}\right\|_{L^2}\le\frac{1}{|y|-k_1}\le\frac{2}{|y|}\qquad\text{for }|y|\ge2k_1.
\]
If $|y|\le2k_1$, then
\[
\left\|\mathrm{Fw}^{-1}\right\|_{L^2}\le\frac{2L}{c_{\mathrm{re}}}\le\frac{2L}{c_{\mathrm{re}}}\cdot\frac{1+2k_1}{1+|y|}.
\]
If $|y|\ge2k_1$, then, by (co5),
\begin{align*}
\left\|\mathrm{Fw}^{-1}\right\|_{L^2}&\le\frac{2(1+|y|)/|y|}{1+|y|}\le\frac{2+1/k_1}{1+|y|}\\
&\le\frac{3}{1+|y|}\le\frac{2L}{c_{\mathrm{re}}}\cdot\frac{1+2k_1}{1+|y|}.
\end{align*}
Hence, for every real $y$,
\begin{equation}\tag{re3}
\left\|\mathrm{Fw}^{-1}\right\|_{L^2}\le\frac{C_{00}}{48}\cdot(1+|y|)^{-1}.
\end{equation}

\emph{Step 4 (kernel).}
Fix $x'=(c',g')$, and put
\[
d_{\mathrm{l1}}(c,c'):=\min_{n\in\mathbb{Z}^3}\left|c-c'-Rn\right|_1.
\]
As a function of $c$ it is $R\mathbb{Z}^3$-periodic; moreover $d_{\mathrm{l1}}(c',c')=0$, $\left|d_{\mathrm{l1}}(c,c')-d_{\mathrm{l1}}(\tilde{c},c')\right|\le|c-\tilde{c}|_1$, and $d_{\mathrm{l1}}(c,c')\ge d(c,c')$.
Choose $\mathrm{wf}:=a_z\cdot d_{\mathrm{l1}}(\cdot,c')$ in Step~2.
Then $S_y(0)=\mathrm{Fz}^{-1}=\exp(-\mathrm{wf})\,\mathrm{Fw}^{-1}\exp(\mathrm{wf})$, and (kn2), (kn4), (re3) give
\begin{align*}
\left|S_y(0)(x,x')\right|_2&=\exp\left(-\mathrm{wf}(c)+\mathrm{wf}(c')\right)\cdot\left|\mathrm{Fw}^{-1}(x,x')\right|_2\\
&\le\exp\left(-a_z\cdot d_{\mathrm{l1}}(c,c')\right)\cdot48\cdot\frac{C_{00}}{48}\cdot(1+|y|)^{-1}\\
&\le C_{00}\cdot(1+|y|)^{-1}\cdot\exp\left(-a_z\cdot d(x,x')\right).
\end{align*}

\emph{(c)}
Let $k\ge1$, so that $w_k=1$.
For a spinor $f$, the spinor $\mathbf{1}_{\Delta_{z_2}}f$ is supported over the orbit of $\Delta_{z_2}$, and $\left|\mathbf{1}_{\Delta_{z_2}}f\right|_{\mathrm{su}}\le|f|_{\mathrm{su}}$.
The moduli of a spinor and of a function of short pairs are invariant under the translations by $R\mathbb{Z}^3$, both points of a pair being translated; therefore a supremum over the orbit of $\Delta_{z_1}$ is the supremum over $\Delta_{z_1}$.
For a short pair $(x,\tilde{x})$ one has
\[
\left(\delta_{\alpha,0}v\right)(x,\tilde{x})=d_X(x,\tilde{x})^{-\alpha}\left(v(\tilde{x})-v(x)\right),\qquad0<d_X(x,\tilde{x})<1,
\]
with $d_X=\eta$ for the pairs inside a cluster.
By (a), therefore,
\begin{align*}
\left\|\mathbf{1}_{\Delta_{z_1}}S_y(0)\mathbf{1}_{\Delta_{z_2}}\right\|_{\mathrm{su}}&\le C_5\cdot\exp\left(-\gamma_{(5)}\cdot d(z_1,z_2)\right),\\
\left\|\mathbf{1}_{\Delta_{z_1}}\delta_{\alpha,0}S_y(0)\mathbf{1}_{\Delta_{z_2}}\right\|_{\mathrm{su}}&\le C_5\cdot\exp\left(-\gamma_{(5)}\cdot d(z_1,z_2)\right).
\end{align*}
By (co2) and (cg),
\begin{align*}
\sum_{z_2}\exp\left(\gamma^{F}\cdot d(z_1,z_2)\right)\cdot C_5\cdot\exp\left(-\gamma_{(5)}\cdot d(z_1,z_2)\right)&\le C_5\cdot c\left(\gamma_{(5)}-\gamma^{F}\right)\\
&\le C_5\cdot c(\gamma^{F})\le C_z\cdot c(\gamma^{F})=a'.
\end{align*}
This bounds both $N_{\gamma^{F}}(S_y(0))$ and $N_{\gamma^{F}}(\delta_{\alpha,0}S_y(0))$ by $a'$.

Let $k=0$.
By (b), (co2) and (cg),
\begin{align*}
N^{\mathrm{ke}}_{\gamma^{F}}(S_y(0))&\le\sup_z\sum_{z'}\exp\left(\gamma^{F}\cdot d(z,z')\right)\cdot C_{00}\cdot(1+|y|)^{-1}\cdot\exp\left(-a_z\cdot d(z,z')\right)\\
&\le C_{00}\cdot c\left(a_z-\gamma^{F}\right)\cdot w_0(y)\le C_z\cdot c(\gamma^{F})\cdot w_0(y)=a'\cdot w_0(y),
\end{align*}
and (Nk1) gives $N_{\gamma^{F}}(S_y(0))\le a'\cdot w_0(y)$.
There is no short pair at $k=0$, since $d_X\ge\eta=1$.

\emph{Uniformity.}
By Proposition~1.6 the constants $C_5$ and $\gamma_{(5)}$ depend only on $(b,\alpha,L,\rho)$, and by Proposition~1.7 the constant $c_{\mathrm{re}}$ depends only on $(b,L,\rho)$; hence $a_z$, $k_1$, $C_{00}$, $C_z$, $\gamma_z$, $\gamma^{F}$ and $a'$ depend only on $(b,\alpha,L,\rho)$.
No constant depends on $k$, $t$, $R$, $\boldsymbol{\sigma}$, on $\mu\in[0,1]$, or on $y$.
The bounds depend on $y$ only through the displayed factor $(1+|y|)^{-1}=w_0(y)$ at $k=0$.
\end{proof}

\subsection*{4. Small fields on a block torus}

\begin{proof}[Proof of Proposition~2.4]
Write $N:=N_{\gamma^{F}}$.
By (co1), $0\le\gamma^{F}L\le\log2$; so Proposition~2.1 holds with $\gamma=\gamma^{F}$.
Moreover, $w_k(y)\le1$ for all $k\ge0$ and all $y\in Y_k$.

\emph{(a)}
Since $3L\delta\le1$, Proposition~2.1\,(a) gives $N(V)\le C'_{V}\cdot\delta$, and Proposition~2.3\,(c) gives $N(S_y(0))\le a'\cdot w_k(y)\le a'$.
By (N2) and $2C'_{V}\cdot a'\cdot\delta\le1$,
\begin{equation}\tag{ne1}
N\left(VS_y(0)\right)\le C'_{V}\cdot\delta\cdot a'\le\frac{1}{2},\qquad N\left(S_y(0)V\right)\le\frac{1}{2}.
\end{equation}
Section~0.7, applied with $O=-VS_y(0)$ and with $O=-S_y(0)V$, shows that the series
\[
\mathrm{Xr}:=\sum_{n\ge0}\left(-VS_y(0)\right)^n,\qquad\mathrm{Xl}:=\sum_{n\ge0}\left(-S_y(0)V\right)^n
\]
converge in $N$, and that
\begin{gather*}
N(\mathrm{Xr})\le2,\qquad N(\mathrm{Xl})\le2,\\
\left(1+VS_y(0)\right)\mathrm{Xr}=\mathrm{Xr}\left(1+VS_y(0)\right)=1,\qquad\left(1+S_y(0)V\right)\mathrm{Xl}=\mathrm{Xl}\left(1+S_y(0)V\right)=1.
\end{gather*}
The operator $F_y(0)$ is bijective by Proposition~2.3\,(a), (b), and
\[
F_y(B)=F_y(0)+V=\left(1+VS_y(0)\right)F_y(0)=F_y(0)\left(1+S_y(0)V\right).
\]
In each of these two products both factors are bijective.
Hence $F_y(B)$ is bijective, and
\begin{align*}
S_y(B)&=S_y(0)\left(1+VS_y(0)\right)^{-1}=S_y(0)\,\mathrm{Xr}=\sum_{n\ge0}S_y(0)\left(-VS_y(0)\right)^n,\\
S_y(B)&=\left(1+S_y(0)V\right)^{-1}S_y(0)=\mathrm{Xl}\,S_y(0)=\sum_{n\ge0}\left(-S_y(0)V\right)^nS_y(0).
\end{align*}
By (N2) and (ne1), the $n$-th term of each of the two series has $N$-norm at most $a'\cdot w_k(y)\cdot2^{-n}$, for every connection $B$ with $e_k|B|_{\infty}\le\delta$; thus the convergence is uniform in $B$.

\emph{(b)}
One has $N(S_y(B))=N\left(S_y(0)\,\mathrm{Xr}\right)\le N(S_y(0))\cdot N(\mathrm{Xr})\le2a'\cdot w_k(y)$.

\emph{(c)}
At $k=0$ there is no short pair.
Let $k\ge1$.
Since $\delta_{\alpha,0}S_y(B)=\left(\delta_{\alpha,0}S_y(0)\right)\mathrm{Xr}$, (N2) and Proposition~2.3\,(c) give $N\left(\delta_{\alpha,0}S_y(B)\right)\le2a'$.
The second identity of Proposition~2.1\,(c), applied to the spinor $S_y(B)f$, gives
\[
\delta_{\alpha,B}S_y(B)-\delta_{\alpha,0}S_y(B)=E_{\mathrm{wq}}S_y(B),\qquad\mathrm{wq}(x,x'):=d_X(x,x')^{-\alpha}\left(U_B(x,x')-1\right);
\]
here $\mathrm{wq}$ is a complex function of the short pairs which, by (ph0), is invariant under the translation of both points by $Rn$.
Inside a cluster one has $c'=c$, $d_X=\eta$ and $U_B=1$, hence $\mathrm{wq}=0$; on these pairs both $\delta_{\alpha,B}S_y(B)f$ and $\delta_{\alpha,0}S_y(B)f$ are equal to $\eta^{-\alpha}\left((S_y(B)f)(x')-(S_y(B)f)(x)\right)$.
For a short pair with $c\ne c'$ one has $d_X(x,x')=|c-c'|<1$, and Lemma~0.1\,(b) gives
\begin{align*}
|\mathrm{wq}(x,x')|&\le3\delta\cdot\exp(3\delta)\cdot|c-c'|^{1-\alpha}\le3\delta\cdot\exp(3\delta)\le L^{-1}\cdot\exp\left(\frac{1}{L}\right)\\
&\le\frac{1}{5}\cdot\exp\left(\frac{1}{5}\right)\le\frac{1}{5}\cdot\frac{5}{4}=\frac{1}{4},
\end{align*}
because $3\delta\le1/L$, $L\ge5$, and $\exp(\mathrm{cx})\le1/(1-\mathrm{cx})$ for $0\le\mathrm{cx}<1$.
Proposition~2.1\,(c) with $\mathrm{wq}_0=1/4$, part (b), and $\exp(\gamma^{F})\le\exp(1/30)\le30/29$ give
\begin{align*}
N\left(E_{\mathrm{wq}}S_y(B)\right)&\le27\cdot\exp(\gamma^{F})\cdot\frac{1}{4}\cdot2a'\cdot w_k(y)\\
&\le\frac{27}{4}\cdot\frac{30}{29}\cdot2a'\le8\cdot2a'=16a'.
\end{align*}
By (N2), $N\left(\delta_{\alpha,B}S_y(B)\right)\le2a'+16a'=18a'$.

\emph{(d)}
Let $k=0$.
Then $S_y(B)=\mathrm{Xl}\,S_y(0)$, and (Nk2) together with Proposition~2.3\,(c) gives
\[
N^{\mathrm{ke}}_{\gamma^{F}}(S_y(B))\le N(\mathrm{Xl})\cdot N^{\mathrm{ke}}_{\gamma^{F}}(S_y(0))\le2a'\cdot w_0(y).
\]
The kernel is measured in the last factor $S_y(0)$ alone; the factor $48$ enters through the constant $C_{00}$ contained in $a'$.

\emph{Uniformity.}
The bounds contain only $a'$ and $w_k(y)$, and $a'$, $\gamma^{F}$, $C'_{V}$ depend only on $(b,\alpha,L,\rho)$.
No constant depends on $k$, $t$, $R$, $\boldsymbol{\sigma}$, $e_k$, on $\mu\in[0,1]$, or on $y$, $\delta$, $B$.
\end{proof}

\subsection*{5. Conditions on the constants}

The proofs of Sections~0--4 use the following conditions on the constants, together with the definitions of constants (rows marked ``definitions'') with which their numerical bounds are computed.
For each entry, the last column names the place in Section~I.8.2 or in Proposition~14.2 where it is imposed or defined;
hence the constants chosen there satisfy all of these conditions.

\noindent
\begin{tabular}{p{0.34\textwidth}p{0.42\textwidth}p{0.14\textwidth}}
\hline
condition & enters in & line\\
\hline
$\rho\in(0,1)$
& (hm5), $|M_3(x)|_2=3(1+\rho)<6$ (Proposition~2.1\,(a), (b), Proposition~2.3\,(b)); (co6), $(1+\rho)\cdot\mathrm{sw}<58$ (Proposition~2.3\,(b)); (co5), $4\rho/(3L)<4/(3L)$; and $c_{\mathrm{re}}>0$
& (C0)\\
$w_3=6$, $\mathrm{sw}=12+12\sqrt{2}$ (definitions)
& (hm5), $|M_3(x)|_2=(w_3/2)(1+\rho)$; (co5), $w_3/18=1/3$; (co6), $\mathrm{sw}\le29$
& (C0)\\
the standard weights, with $\sum_iw_i\hat n_i=0$
& (hm1), which is (I.m1): the diagonal term of $D_A$ is the scalar $(\rho/2)\cdot\mathrm{sw}/\eta_k$ (Proposition~2.2\,(c), (d), (e))
& (C0); display (I.1.1)\\
$b>0$ (with $b:=1$)
& (mx6), (mx7) (through $b_k>0$); $C'_{V}$; $k_1$; $c_{\mathrm{re}}>0$
& (C0)\\
$\alpha<1$
& Proposition~2.4\,(c), $|c-c'|^{1-\alpha}\le1$ on short pairs
& (C1)\\
$L$ odd
& Section~G1.2 ($\mathrm{Lab}_{j+1}\subset\mathrm{Lab}_j$; nesting of the blocks); (bl2), (bl3); the translation $\mathrm{sft}$ ($1-L^{k+1-j}$ is even)
& (C3), (L5)\\
$L\ge5$
& Step~3 of the proof of Proposition~2.3\,(b), $3b/(2L)\le(3/10)\,b$; (co1), $\gamma^{F}\le1/30$; (co5), $(3/2)L^2>37$; Proposition~2.4\,(c), $L^{-1}\cdot\exp(1/L)\le1/4$
& (C3), (L5)\\
$C'_{V}=2\mathrm{e}\cdot(42+6b)$
& Proposition~2.1\,(a), (b)
& (C4)\\
$C_5\ge1$ and $\gamma_{(5)}>0$ (Proposition~1.6); $c_{\mathrm{re}}>0$ (Proposition~1.7)
& (co1), (co3); Proposition~2.3
& (C4r)\\
$c_{\mathrm{re}}=\min\left\{b(1-1/L),\,\rho/(2C_{\mathrm{Po}}L)\right\}$, $C_{\mathrm{Po}}=1/(8c_{*})+1/\lambda_{\mathrm{int}}$, $c_{*}=\min(w_3,\lambda_{\mathrm{int}})/18$, $\lambda_{\mathrm{int}}>0$ (definitions); and, as consequences of them, $c_{\mathrm{re}}\le b(1-1/L)$, $c_{\mathrm{re}}\le\rho/(2C_{\mathrm{Po}}L)$, $C_{\mathrm{Po}}\ge1/(8c_{*})$, $c_{*}\le w_3/18$
& (co5), and through (co5) in Step~3 of the proof of Proposition~2.3\,(b): the $L^2$ norm of $\mathrm{Rw}$ is at most $c_{\mathrm{re}}/(2L)\le b/(2L)$, and $3\le C_{00}/48$
& (C4r), (C0r)\\
$a_z\le1/(3L)$ and $a_z\le c_{\mathrm{re}}/\left(2\mathrm{e}\cdot L\cdot(6+3b)\right)$
& (co1); Step~2 of the proof of Proposition~2.3\,(b), namely (re2)
& (C4r), and the definition of $a_z$\\
$k_1=60+b$ and $C_{00}=48\cdot(2L/c_{\mathrm{re}})\cdot(1+2k_1)$
& Steps~3 and~4 of the proof of Proposition~2.3\,(b)
& (C4r)\\
$C_z=\max(C_5,C_{00})$, $\gamma_z=\min\left(\gamma_{(5)},a_z\right)$, $\gamma^{F}=\gamma_z/2$, $a'=C_z\cdot c(\gamma^{F})$
& Proposition~2.3\,(c); Proposition~2.4
& (C4r)\\
\hline
\end{tabular}

The inequalities $0\le\gamma L\le\log2$, $3L\delta\le1$, $3L\cdot\mathfrak{i}(A)\le1$ and $2C'_{V}\cdot a'\cdot\delta\le1$ are hypotheses of Propositions~2.1--2.4 on their data $\gamma$, $\delta$ and $A$.
The conditions stated before Proposition~2.1, which stand in Proposition~14.2\,(b), are (L5), (M1), (M2), the inequality $M\ge108\pi\cdot C'_{V}\cdot a'$ contained in (M5), (Z2.s1) and (Z2.s2).
Of these six, Propositions~2.1--2.4 require only (L5);
the other five are conditions on the constants $M$ and $e_0$, which are chosen after the constants of line (C4r).
For $\delta=2e_k^{\varepsilon}$ the two conditions on $\delta$ in Proposition~2.4 are the conditions (Z2.s1) and (Z2.s2).

\section*{Chapter 2b. Local inverses on mirror tori, parametrix and walks, the inverse of $F_y(A)$ on $\tilde{G}_k$: proofs of Propositions 2.5--2.8}

\subsection*{0. Conventions of this chapter; preliminaries for the four proofs}

Throughout this chapter $\mathrm{1v}:=(1,1,1)$, and for $n\in\mathbb{Z}^3$ one writes $\mathrm{pa}(n):=n_1+n_2+n_3$, so that $\mathrm{pa}(n+n')=\mathrm{pa}(n)+\mathrm{pa}(n')$.
A prime or an index $n$ attached to $\mathrm{res}$, $\mathrm{ap}$, $\mathrm{ap2}$, $\mathrm{ext}$, $\mathrm{rs}$, $\mathrm{Fo}$, $\mathrm{emb}$, $o$, $\tilde{W}$ indicates the object of the same name formed with another cube or with another lift, and $\tilde{z}'$ is the lift of $z'$.

\subsubsection*{0.1. Data and constants}

The data are those of (SF) in Section~G1.1; in Proposition~2.5 the number $e_k$ may be any positive number.
By (L5) and (M1), $L$ is odd with $L\ge5$, and $M$ is a power of $L$ with $M\ge L^2$; consequently $M/L$ is a power of $L$.
By the definition of $\mathrm{Cl}$, the side $\ell$ is a multiple of $LM$ with $\ell\ge13LM$.
The constants $C'_V$, $\gamma^{F}$, $a'$ are those of the lines (C4), (C4r) of Proposition~14.2\,(a), namely
\begin{gather*}
C'_V=2\mathrm{e}\cdot\left(42+6b\right),\qquad\gamma_{z}=\min\left(\gamma_{(5)},a_{z}\right),\qquad\gamma^{F}=\frac{\gamma_{z}}{2},\\
C_{z}=\max\left(C_5,C_{00}\right),\qquad a'=C_{z}\cdot c\left(\gamma^{F}\right).
\end{gather*}
They satisfy
\begin{gather*}
\gamma^{F}L\le\frac{1}{6}<\log2,\qquad\exp\left(\gamma^{F}\right)\le\exp\left(\frac{1}{30}\right)<1.04,\qquad a'\ge1,\\
C'_V\ge84\mathrm{e}\ge24\mathrm{e},\qquad M\ge108\pi\cdot C'_V\cdot a'\ge108\cdot24\pi\mathrm{e}>1600.\tag*{[P.1]}
\end{gather*}
Indeed, $\gamma^{F}\le\gamma_{(5)}/2$, while $0<\gamma_{(5)}\le1/(3L)$ by the list of constants of Proposition~1.6 and $L\ge5$; this gives the first two inequalities.
Next, $C_{z}\ge C_5\ge1$ by Proposition~1.6, and $c(\gamma^{F})\ge1$, since the term $n=0$ of the sum that defines $c(\gamma^{F})$ equals $1$; hence $a'\ge1$.
The bound on $C'_V$ holds because $b=1>0$.
Finally, the first bound on $M$ is contained in (M5), and $108\cdot24\cdot\pi\cdot\mathrm{e}>108\cdot24\cdot3\cdot2.7=20995.2>1600$.

For $e_k\in(0,e_0]$ the number $\delta_k:=2e_k^{\varepsilon}$ satisfies
\begin{equation*}
3L\cdot\delta_k\le1\qquad\text{and}\qquad2C'_V\cdot a'\cdot\delta_k\le1.\tag*{[P.2]}
\end{equation*}
In fact $3L\cdot\delta_k=6L\cdot e_k^{\varepsilon}\le6L\cdot e_0^{\varepsilon}\le1$ by (Z2.s1), and $2C'_V\cdot a'\cdot\delta_k=4C'_V\cdot a'\cdot e_k^{\varepsilon}\le1$ by (Z2.s2).
Thus $\delta_k$ is admissible as the number $\delta$ of Proposition~2.4.

The following inequalities hold:
\begin{gather*}
\frac{3M}{4}+L<\frac{3M}{2}-\eta_k,\qquad\frac{3M}{2}+L<2M,\qquad\frac{3M}{2}+3<2M,\qquad\frac{3M}{2}+2L+1<3M,\\
3M<\frac{\ell}{2}.\tag*{[P.3]}
\end{gather*}
The first four are consequences of (M2), that is of $M>4L$: one has $3M/4>3L\ge L+1\ge L+\eta_k$; next $L<M/4$; next $M>20$; and $2L+1<M/2+1<3M/2$.
The last one follows from $\ell\ge13LM\ge65M$.

\subsubsection*{0.2. The operator $F_y$ on a lattice without periodicity; its rows}

\emph{Blocked lattices.}
Let $\mathrm{Lb}$ be one of the two lattices $\mathbb{L}_k=\eta_k\mathbb{Z}^3$ (shift $t=0$) and $\mathrm{Lab}^W=(\eta_k/2)\mathrm{1v}+\eta_k\mathbb{Z}^3$.
For $0\le j\le k+1$ put $\mathrm{Lb}_j:=t_j+L^j\eta_k\mathbb{Z}^3$, with $t_j=0$ in the first case and $t_j=(L^j\eta_k/2)\mathrm{1v}$ in the second.
As in Sections~I.1.1 and~G1.2, the blocks are
\begin{equation*}
\mathrm{Bl}_j(z)=\left\{c\in\mathrm{Lb}:|c-z|<\frac{L^j\eta_k}{2}\right\}\qquad\left(z\in\mathrm{Lb}_j,\ 0\le j\le k+1\right),
\end{equation*}
the unit blocks are $\Delta_z=\mathrm{Bl}_k(z)$, and $B(z')=\mathrm{Lb}_k\cap\mathrm{Bl}_{k+1}(z')$.
Every block is a lattice box, that is, the set of the points of $\mathrm{Lb}$ that lie in a product of three intervals.
Since $L$ is odd, the blocks of one side partition $\mathrm{Lb}$, and blocks of different sides are nested.

\emph{The operator.}
Let $\tilde{A}$ be a complex function on the oriented nearest-neighbour bonds of $\mathrm{Lb}$ with $\tilde{A}(c',c)=-\tilde{A}(c,c')$; no periodicity is assumed.
Let $e_k>0$, $\mu\ge0$ and $y\in Y_k$.
The phases $u_{\tilde{A}}$, $U_{\tilde{A}}$, $\varphi_{\tilde{A}}$, $\varphi^{+}_{\tilde{A}}$ are formed by (I.1.6) and (I.3.3).
For an arbitrary function $f\colon\mathrm{Lb}\times O_{\mathrm{cub}}\to\mathbb{C}^4$ the display (I.2.3) and the displays of Sections~I.3.3, I.3.5 ($k\ge1$) and I.3.9 ($k=0$) are finite sums.
To write them out, let $x=(c,g)$, $c_3=c+\eta_k\,ge_1$, let $z$ be the point of $\mathrm{Lb}_k$ with $c\in\Delta_z$, and let $z'$ be the point of $\mathrm{Lb}_{k+1}$ with $c\in\mathrm{Bl}_{k+1}(z')$.
For $k\ge1$ one obtains
\begin{align*}
\left(F_y(\tilde{A})f\right)(c,g)={}&-\eta_k^{-1}\sum_{i=0}^{2}M_i(x)\left[f(c,gs_i)-f(c,g)\right]\\
&-\eta_k^{-1}M_3(x)\left[u_{\tilde{A}}(c,c_3)\,f(c_3,gs_3)-f(c,g)\right]+\mu f(c,g)\\
&+\boldsymbol{\alpha}_y\,\varphi_{\tilde{A}}(z,c)^{-1}\,\frac{\eta_k^3}{48}\sum_{c'\in\Delta_z}\ \sum_{g'\in O_{\mathrm{cub}}}\varphi_{\tilde{A}}(z,c')\,f(c',g')\\
&+\boldsymbol{\beta}_y\,\varphi^{+}_{\tilde{A}}(z',c)^{-1}\,L^{-3}\,\frac{\eta_k^3}{48}\sum_{c'\in\mathrm{Bl}_{k+1}(z')}\ \sum_{g'\in O_{\mathrm{cub}}}\varphi^{+}_{\tilde{A}}(z',c')\,f(c',g'),\tag*{[P.4]}
\end{align*}
and for $k=0$, where $\eta_0=1$ and $\Delta_z=\{c\}$,
\begin{align*}
\left(F_y(\tilde{A})f\right)(c,g)={}&-\sum_{i=0}^{2}M_i(x)\left[f(c,gs_i)-f(c,g)\right]\\
&-M_3(x)\left[u_{\tilde{A}}(c,c_3)\,f(c_3,gs_3)-f(c,g)\right]+\mu f(c,g)+iy\,\gamma_5f(c,g)\\
&+b_L\,U_{\tilde{A}}(z',c)^{-1}\,L^{-3}\,\frac{1}{48}\sum_{c'\in B(z')}\ \sum_{g'\in O_{\mathrm{cub}}}U_{\tilde{A}}(z',c')\,f(c',g').\tag*{[P.5]}
\end{align*}
The fourth term of [P.4] is $\boldsymbol{\alpha}_y\,JP_k(\tilde{A})J^{*}f$, because $\left(Q_k(\tilde{A})J^{*}f\right)(z)=\int_{\Delta_z}\mathrm{d}x'\,\varphi_{\tilde{A}}(z,c')f(x')$ and $\left(JQ_k^{T}(-\tilde{A})v\right)(x)=\varphi_{\tilde{A}}(z,c)^{-1}v(z)$.
The fifth term is $\boldsymbol{\beta}_y\,JQ_k^{T}(-\tilde{A})Q^{T}(-\tilde{A})Q(\tilde{A})Q_k(\tilde{A})J^{*}f$, because
\begin{align*}
\left(Q(\tilde{A})Q_k(\tilde{A})J^{*}f\right)(z')&=L^{-3}\sum_{z''\in B(z')}U_{\tilde{A}}(z',z'')\int_{\Delta_{z''}}\mathrm{d}x'\,\varphi_{\tilde{A}}(z'',c')f(x')\\
&=L^{-3}\int_{\mathrm{Bl}_{k+1}(z')}\mathrm{d}x'\,\varphi^{+}_{\tilde{A}}(z',c')f(x')
\end{align*}
and $\varphi_{\tilde{A}}(z,c)^{-1}\,U_{\tilde{A}}(z',z)^{-1}=\varphi^{+}_{\tilde{A}}(z',c)^{-1}$.
The last term of [P.5] is $b_L\,JQ^{T}(-\tilde{A})Q(\tilde{A})J^{*}f$.
Here $M_i(x)=(w_i/2)\left(\rho+\gamma\cdot(g\hat{n}_i)\right)$ depends on the label only, and $\boldsymbol{\alpha}_y$, $\boldsymbol{\beta}_y$ are the constant matrices of Section~I.3.5.

\begin{named}{(F1)}[What a row reads]
The value $\left(F_y(\tilde{A})f\right)(c,g)$ reads $f$ at the four sites $(c,g)$, $(c,gs_0)$, $(c,gs_1)$, $(c,gs_2)$ over $c$, at the single site $(c_3,gs_3)$ over the nearest neighbour $c_3$ of $c$, and at the $48(L/\eta_k)^3$ sites over the block $\mathrm{Bl}_{k+1}(z')$ of side $L$ that contains $c$.
It reads $\tilde{A}$ on the bond $(c,c_3)$ and on bonds whose two end points lie in $\mathrm{Bl}_{k+1}(z')$.
\end{named}

\begin{proof}
The sites are read off from [P.4] and [P.5].
The exponent of $\varphi_{\tilde{A}}(z,c')$ is $ie_k\eta_k\sum_{j<k}(\tau\tilde{A})(y_{j+1},y_j)$, where $y_j$ denotes the centre of the block of side $L^j\eta_k$ that contains $c'$.
The points $y_j$ and $y_{j+1}$ both lie in $\mathrm{Bl}_{j+1}(y_{j+1})$, and, the blocks being nested, this block is contained in $\Delta_z$ and in $\mathrm{Bl}_{k+1}(z')$.
A path $\boldsymbol{\Gamma}^{\pi}(y_{j+1},y_j)$ changes one coordinate at a time, and does so monotonically; therefore it stays in the lattice box spanned by its end points, hence in the lattice box $\mathrm{Bl}_{j+1}(y_{j+1})$.
In the same way the six paths $\boldsymbol{\Gamma}^{\pi}(z',z'')$, $z''\in B(z')$, which enter $U_{\tilde{A}}(z',z'')$, stay in the lattice box spanned by $z'$ and $z''$, and this box is contained in the lattice box $\mathrm{Bl}_{k+1}(z')$.
\end{proof}

\begin{named}{(F2)}[Translations]
Let $\mathrm{Lb}$, $\mathrm{Lb}'$ be blocked lattices, and let $v\in\mathbb{R}^3$ satisfy $\mathrm{Lb}_j+v=\mathrm{Lb}'_j$ for $0\le j\le k+1$.
Put $(T_vf)(c,g):=f(c-v,g)$ and $(T_v\tilde{A})(c,c'):=\tilde{A}(c-v,c'-v)$.
Then $T_vF_y(\tilde{A})=F_y(T_v\tilde{A})\,T_v$, and
\begin{equation*}
\left(\delta_{\alpha,T_v\tilde{A}}T_vf\right)(x+v,x'+v)=\left(\delta_{\alpha,\tilde{A}}f\right)(x,x')
\end{equation*}
for all pairs $x\ne x'$ with $d_X(x,x')\le M$.
\end{named}

\begin{proof}
One has $(x+v)^{(i)}=x^{(i)}+v$ and $M_i(x+v)=M_i(x)$ by (I.1.2) and (I.2.1), the label not being moved.
Moreover $\mathrm{Bl}_j(z)+v=\mathrm{Bl}_j(z+v)$, so that blocks are mapped to blocks and centres to centres.
Since $\boldsymbol{\Gamma}^{\pi}(c+v,c'+v)=\boldsymbol{\Gamma}^{\pi}(c,c')+v$, one has $(\tau T_v\tilde{A})(c+v,c'+v)=(\tau\tilde{A})(c,c')$; consequently $u$, $U$, $\varphi$, $\varphi^{+}$ formed with $T_v\tilde{A}$ at the translated arguments coincide with those formed with $\tilde{A}$ at the original arguments.
Also $d_X(x+v,x'+v)=d_X(x,x')$.
Inserting these facts in [P.4], [P.5] and in (I.1.7) one obtains $\left(F_y(T_v\tilde{A})T_vf\right)(x+v)=\left(F_y(\tilde{A})f\right)(x)$ and the relation for $\delta_{\alpha}$.
\end{proof}

\begin{named}{(F3)}[Tori]
Let $\mathbb{T}=(t,R,\boldsymbol{\sigma})$ be a block torus with the lattice $\mathrm{Lb}$, and let $A$ be a connection on $\mathbb{T}$.
A spinor on $\mathbb{T}$ is a function $f$ on $\mathrm{Lb}\times O_{\mathrm{cub}}$ with $f(c+Re_{\mu},g)=\boldsymbol{\sigma}f(c,g)$, the connection $A$ is an $R\mathbb{Z}^3$-periodic bond function on $\mathrm{Lb}$, and $F_y(A)$ on $\mathbb{T}$ is formed by the displays listed in Section~G1.2.
Hence $F_y(A)f$ on $\mathbb{T}$ is the function given by [P.4], [P.5] with $\tilde{A}=A$.
Since $R$ is a multiple of $L=L^{k+1}\eta_k$, (F2) with $v=Re_{\mu}$ shows that this function is again a spinor on $\mathbb{T}$.
This applies in particular to $T_k=(0,\ell,-1)$ and to $\mathbb{T}^{W}$.
\end{named}

\begin{named}{(F4)}[Entries]
Consider the basis of the functions that are equal to a unit vector of $\mathbb{C}^4$ at one site of a period cell and are continued by the sign rule.
In this basis each entry of the matrix of $F_y(A)$ on $\mathbb{T}$ is a sum of terms of the following kinds: constants;
\begin{gather*}
-\eta_k^{-1}\left(M_3(x)\right)_{\beta\beta'}\exp\left\{ie_k\eta_k\,A(c,c_3)\right\};\\
\left(\boldsymbol{\alpha}_y\right)_{\beta\beta'}\frac{\eta_k^3}{48}\exp\left\{ie_k\eta_k\left[(\tau_kA)(z,c')-(\tau_kA)(z,c)\right]\right\};\\
\left(\boldsymbol{\beta}_y\right)_{\beta\beta'}L^{-3}\frac{\eta_k^3}{48}\exp\left\{ie_k\eta_k\left[(\tau A)(z',z'')+(\tau_kA)(z'',c')-(\tau A)(z',z)-(\tau_kA)(z,c)\right]\right\},
\end{gather*}
where $z''$ is the unit site with $c'\in\Delta_{z''}$, so that the last exponent is that of $\varphi^{+}_A(z',c)^{-1}\varphi^{+}_A(z',c')$; and, at $k=0$, in the place of the last two kinds,
\begin{equation*}
b_L\,L^{-3}\,\frac{1}{48}\exp\left\{ie_k\left[(\tau A)(z',c')-(\tau A)(z',c)\right]\right\}\mathbf{1}[\beta=\beta']\qquad\text{and}\qquad iy\left(\gamma_5\right)_{\beta\beta'}.
\end{equation*}
Each term is multiplied by $\boldsymbol{\sigma}^{\mathrm{pa}(n)}$ if the site that the row reads is $x'+Rn$, $x'$ being the site of the period cell.
The exponents are linear forms in the numbers $A(b)$; hence every entry is an entire function of the numbers $A(b)$.
For real $y$ the matrices $\boldsymbol{\alpha}_y$, $\boldsymbol{\beta}_y$ depend continuously on $y$: the matrix $\gamma_5$ is Hermitian with $\gamma_5^2=1$, and on its eigenspaces $b_k+i\gamma_5y$ and $b_k+b_L+i\gamma_5y$ are the numbers $b_k\pm iy$ and $b_k+b_L\pm iy$, none of which vanishes.
Hence every entry is continuous in $(y,A)$, $y\in\mathbb{R}$.
\end{named}

\subsubsection*{0.3. Row sums: the facts (N1)--(N3), (Nk1)--(Nk4) of Section~G1.3}

Two settings occur.
In the first, $\mathbb{T}$ is a block torus, with its unit sites taken modulo $R\mathbb{Z}^3$ and with the distance $d$.
In the second, one considers the set $\tilde{W}$ of Section~G1.5 together with its unit sites, the points of $\mathbb{Z}^3$ in $\tilde{W}$, and the distance $d(z,z'):=|z-z'|$.
In both settings the sets $\Delta_z\times O_{\mathrm{cub}}$ (on $\mathbb{T}$: their orbits) partition the fermion sites.
On $\mathbb{T}$ this is part of Section~G1.2.
For $\tilde{W}$ it holds because $\tilde{W}$ is a union of blocks of side $L$ by [Q.1]\,(a) in Section~0.4, hence, the blocks being nested, a union of the unit blocks $\Delta_z$, $z\in\mathbb{Z}^3\cap\tilde{W}$.
Therefore $\sum_z\mathbf{1}_{\Delta_z}=1$ on spinors (over $\tilde{W}$: on the functions on $\tilde{W}\times O_{\mathrm{cub}}$), and $d$ satisfies the triangle inequality.
Let $\gamma\ge0$.
For an operator $O$ between the spaces of Section~G1.3 and a unit site $z$, the row sums of Section~G1.3 are formed with the unit sites $z'$ and the distance $d$ of the setting:
\begin{gather*}
\mathrm{Rw}_z(O)=\sum_{z'}\exp\left(\gamma d(z,z')\right)\left\|\mathbf{1}_{\Delta_z}O\mathbf{1}_{\Delta_{z'}}\right\|_{\mathrm{su}},\qquad N_{\gamma}(O)=\sup_z\mathrm{Rw}_z(O),\\
\mathrm{Rk}_z(O)=\sum_{z'}\exp\left(\gamma d(z,z')\right)\sup\left\{|O(x,x')|_2:x\text{ over }z,\ x'\text{ over }z'\right\},\\
N^{\mathrm{ke}}_{\gamma}(O)=\sup_z\mathrm{Rk}_z(O),\tag*{[P.6]}
\end{gather*}
the quantities of the second and third lines being formed at $k=0$ only.
Over $\tilde{W}$ these are the norms of Section~G1.6.
The kernels are those of Section~G1.2 on $\mathbb{T}$ and those of Section~G1.6 over $\tilde{W}$, where the kernel is taken with respect to $\int_{\tilde{W}}\mathrm{d}x$.

\begin{named}{(Rw1)--(Rw5)}
\begin{itemize}
\item[(Rw1)] $\mathrm{Rw}_z(O_1+O_2)\le\mathrm{Rw}_z(O_1)+\mathrm{Rw}_z(O_2)$.
\item[(Rw2)] If $O_2$ takes values in spinors, then $\mathrm{Rw}_z(O_1O_2)\le\mathrm{Rw}_z(O_1)\cdot N_{\gamma}(O_2)$.
\item[(Rw3)] Let $\psi$, $\psi'$ be functions of $c$: on $\mathbb{T}$, scalars in the sense of Section~G1.2, that is, $R\mathbb{Z}^3$-periodic functions; over $\tilde{W}$, arbitrary functions on $\tilde{W}$.
Then $\mathrm{Rw}_z(\psi O\psi')\le\sup_{\Delta_z}|\psi|\cdot|\psi'|_{\infty}\cdot\mathrm{Rw}_z(O)$, also when $O$ takes values in the functions of short pairs; in particular $\mathrm{Rw}_z(\psi O)=0$ if $\psi$ vanishes on $\Delta_z$.
Moreover, for the multiplication by $\psi$ on the spinors on $\mathbb{T}$, which at $k=0$ are the spinors on the unit lattice, and over $\tilde{W}$ on the functions on $\tilde{W}\times O_{\mathrm{cub}}$, one has $N_{\gamma}(\psi)=|\psi|_{\infty}$;
for the multiplication by $\psi$ on the spinors on the unit lattice of $\mathbb{T}$ at $k\ge1$, where $\psi$ acts at a unit site $z$ by its value $\psi(z)$, one has $N_{\gamma}(\psi)=\sup_z|\psi(z)|\le|\psi|_{\infty}$, the supremum being taken over the unit sites $z$.
Also $N_{\gamma}(1)=1$.
\item[(Rw4)] $\|O\|_{\mathrm{su}}\le N_{\gamma}(O)$, and $\left\|\mathbf{1}_{\Delta_z}O\mathbf{1}_{\Delta_{z'}}\right\|_{\mathrm{su}}\le\exp\left(-\gamma d(z,z')\right)N_{\gamma}(O)$.
\item[(Rw5)] ($k=0$.) $\mathrm{Rw}_z(O)\le\mathrm{Rk}_z(O)$; $\mathrm{Rk}_z(O_1+O_2)\le\mathrm{Rk}_z(O_1)+\mathrm{Rk}_z(O_2)$; $\mathrm{Rk}_z(O_1O_2)\le\mathrm{Rw}_z(O_1)\cdot N^{\mathrm{ke}}_{\gamma}(O_2)$; $\mathrm{Rk}_z(\psi O\psi')\le|\psi(z)|\cdot|\psi'|_{\infty}\cdot\mathrm{Rk}_z(O)$ for $\psi$, $\psi'$ as in (Rw3); and $|O(x,x')|_2\le\exp\left(-\gamma d(x,x')\right)N^{\mathrm{ke}}_{\gamma}(O)$.
\end{itemize}
\end{named}

\begin{proof}
(Rw1) is the triangle inequality for $\|\cdot\|_{\mathrm{su}}$.

(Rw2).
One has
\begin{gather*}
\mathbf{1}_{\Delta_z}O_1O_2\mathbf{1}_{\Delta_{z'}}=\sum_{z''}\left(\mathbf{1}_{\Delta_z}O_1\mathbf{1}_{\Delta_{z''}}\right)\left(\mathbf{1}_{\Delta_{z''}}O_2\mathbf{1}_{\Delta_{z'}}\right),\\
\exp\left(\gamma d(z,z')\right)\le\exp\left(\gamma d(z,z'')\right)\exp\left(\gamma d(z'',z')\right).
\end{gather*}
Summing first over $z'$ and then over $z''$ one finds
\begin{equation*}
\mathrm{Rw}_z(O_1O_2)\le\sum_{z''}\exp\left(\gamma d(z,z'')\right)\left\|\mathbf{1}_{\Delta_z}O_1\mathbf{1}_{\Delta_{z''}}\right\|_{\mathrm{su}}\mathrm{Rw}_{z''}(O_2),
\end{equation*}
and $\mathrm{Rw}_{z''}(O_2)\le N_{\gamma}(O_2)$.

(Rw3).
The functions $\psi$ and $\psi'$ commute with the indicators: on functions of pairs $\psi$ multiplies by its value at the first point, and $\mathbf{1}_{\Delta_z}$ restricts the first point (Section~G1.3).
Hence
\begin{equation*}
\mathbf{1}_{\Delta_z}\psi O\psi'\mathbf{1}_{\Delta_{z'}}=\left(\psi\mathbf{1}_{\Delta_z}\right)\left(\mathbf{1}_{\Delta_z}O\mathbf{1}_{\Delta_{z'}}\right)\psi',
\end{equation*}
where the multiplication by $\psi\mathbf{1}_{\Delta_z}$ has a norm at most $\sup_{\Delta_z}|\psi|$ and the multiplication by $\psi'$ a norm at most $|\psi'|_{\infty}$.
On $\mathbb{T}$ this uses the periodicity of $\psi$ and $\psi'$: their products with spinors are spinors, the product of $\psi$ with a function of short pairs is a function of short pairs, and on the spinors on $\mathbb{T}$ the operator $\mathbf{1}_{\Delta_z}$ is the indicator of the orbit of $\Delta_z\times O_{\mathrm{cub}}$ (on functions of short pairs it restricts the first point to this orbit), on which $|\psi|$ takes the same values as on $\Delta_z$.
On the spinors on the unit lattice at $k\ge1$ the operator $\mathbf{1}_{\Delta_z}$ is the indicator of the class of $z$ (Section~G1.3), and $\psi$ acts there by the number $\psi(z)$;
so the multiplication by $\psi\mathbf{1}_{\Delta_z}$ has the norm $|\psi(z)|$, which is at most $\sup_{\Delta_z}|\psi|$ because $z\in\Delta_z$.
For $O=1$ one has $\mathbf{1}_{\Delta_z}\psi\mathbf{1}_{\Delta_{z'}}=0$ whenever $z\ne z'$, so that $\mathrm{Rw}_z(\psi)$ is the norm of the multiplication by $\psi\mathbf{1}_{\Delta_z}$:
on the spinors on $\mathbb{T}$ and over $\tilde{W}$ one has $\mathrm{Rw}_z(\psi)=\sup_{\Delta_z}|\psi|$, and on the spinors on the unit lattice at $k\ge1$ one has $\mathrm{Rw}_z(\psi)=|\psi(z)|$.
Taking the supremum over $z$ one obtains the values of $N_{\gamma}(\psi)$ given in (Rw3); for the first one it is used that the blocks $\Delta_z$ cover the sites of $\mathbb{T}$, respectively of $\tilde{W}$.
For $\psi=1$ both values are equal to $1$.

(Rw4).
For every $f$,
\begin{equation*}
\left|\mathbf{1}_{\Delta_z}Of\right|_{\mathrm{su}}\le\sum_{z'}\left\|\mathbf{1}_{\Delta_z}O\mathbf{1}_{\Delta_{z'}}\right\|_{\mathrm{su}}|f|_{\mathrm{su}}\le\mathrm{Rw}_z(O)\,|f|_{\mathrm{su}},
\end{equation*}
and $|Of|_{\mathrm{su}}$ is the supremum over $z$ of the left side; this gives the first inequality.
The second one is obtained by keeping a single term of the sum $\mathrm{Rw}_z(O)$.

(Rw5).
At $k=0$ the set $\Delta_z\times O_{\mathrm{cub}}$ is the cluster over $z$, which consists of $48$ sites of weight $1/48$.
For $x$ over $z$,
\begin{align*}
\left|\left(O\mathbf{1}_{\Delta_{z'}}f\right)(x)\right|_2&=\left|\frac{1}{48}\sum_{g'}O\left(x,(z',g')\right)f(z',g')\right|_2\\
&\le\sup\left\{|O(x,x')|_2:x'\text{ over }z'\right\}|f|_{\mathrm{su}};
\end{align*}
on $\mathbb{T}$ the representative $z'$ is taken in the period cell, and the modulus does not depend on this choice.
This is the first inequality; the second is the triangle inequality for $|\cdot|_2$.
For the third, let $\mathrm{vv}\in\mathbb{C}^4$ and fix $x'$.
The column $t\mapsto O_2(t,x')\mathrm{vv}$ is a spinor (by Section~G1.2, $O_2(t+Re_{\mu},x')=\boldsymbol{\sigma}O_2(t,x')$; over $\tilde{W}$ it is a function on $\tilde{W}\times O_{\mathrm{cub}}$), and by linearity
\begin{equation*}
(O_1O_2f)(x)=\int\mathrm{d}x'\,\left(O_1\left[O_2(\cdot,x')f(x')\right]\right)(x);
\end{equation*}
therefore $(O_1O_2)(x,x')\mathrm{vv}=\left(O_1\left[O_2(\cdot,x')\mathrm{vv}\right]\right)(x)$.
For $x$ over $z$ and $x'$ over $z'$ this gives
\begin{equation*}
\left|(O_1O_2)(x,x')\mathrm{vv}\right|_2\le\sum_{z''}\left\|\mathbf{1}_{\Delta_z}O_1\mathbf{1}_{\Delta_{z''}}\right\|_{\mathrm{su}}\sup\left\{|O_2(t,x')|_2:t\text{ over }z''\right\}|\mathrm{vv}|_2.
\end{equation*}
Multiplying by $\exp\left(\gamma d(z,z')\right)\le\exp\left(\gamma d(z,z'')\right)\exp\left(\gamma d(z'',z')\right)$ and summing over $z'$ one arrives at the third inequality.
The fourth follows from $(\psi O\psi')(x,x')=\psi(c)\,O(x,x')\,\psi'(c')$ (Section~G1.2).
The fifth is one term of the sum $\mathrm{Rk}_z(O)$, namely the term with $z=c$ and $z'=c'$.
\end{proof}

Taking suprema over $z$, one sees that (Rw1)--(Rw4) contain (N1)--(N3) and that (Rw5) contains (Nk1)--(Nk4).
By (Rw1), (Rw2) and (Rw3), $N_{\gamma}$ is a submultiplicative norm with $N_{\gamma}(1)=1$ on the operators on the finite-dimensional space of the spinors on $\mathbb{T}$; it is definite, because $N_{\gamma}(O)=0$ implies $O=\sum_{z,z'}\mathbf{1}_{\Delta_z}O\mathbf{1}_{\Delta_{z'}}=0$.

\subsubsection*{0.4. Photon statements}

The following facts concern photon objects only, and they are proved in \cite{QEDI}, where the coordinate directions carry the indices $0,1,2$.
They are stated here in the notation of this paper, with the directions $\mu=1,2,3$.
A \emph{mirror plane} is a plane $\{u\in\mathbb{R}^3:u_{\mu}=3Mn\}$ with $\mu\in\{1,2,3\}$ and $n\in\mathbb{Z}$.
The mirror planes divide the mirror torus $\mathbb{T}^{W}$ into the eight \emph{chambers of $\mathbb{T}^{W}$},
\begin{equation*}
\left\{u\in\mathrm{Lab}^W:u_{\mu}\in\left(3Ms_{\mu},3M(s_{\mu}+1)\right)+6M\mathbb{Z}\ \text{ for }\mu=1,2,3\right\}\qquad\left(s\in\{0,1\}^3\right);
\end{equation*}
each of them is invariant under the translations by $6M\mathbb{Z}^3$ and is regarded as a set of sites of $\mathbb{T}^{W}$, and $\mathrm{Ch}$ is the chamber of $\mathbb{T}^{W}$ with $s=(0,0,0)$ \cite[Definition~16.9]{QEDI}.

\begin{named}{[Q.1]}
\begin{itemize}
\item[(a)] For $0\le j\le k+1$ the number $3M/(L^j\eta_k)$ is an odd integer.
No site of $\mathrm{Lab}^W$ lies on a mirror plane $\{u_{\mu}\in3M\mathbb{Z}\}$, and every block of $\mathbb{T}^{W}$ lies in one chamber of $\mathbb{T}^{W}$.
The map $\mathrm{emb}$ is a bijection of $\tilde{W}$ onto $\mathrm{Ch}$ (modulo $6M\mathbb{Z}^3$); it maps nearest neighbours to nearest neighbours, and the blocks of $T_k$ contained in $\tilde{W}$ onto the blocks of $\mathbb{T}^{W}$ contained in $\mathrm{Ch}$, centres to centres.
The set $\tilde{W}$ is a union of blocks of side $L$.
\item[(b)] The map $\mathrm{Fo}$ maps $\mathrm{Lab}^W$ onto $\tilde{W}$, it is $6M\mathbb{Z}^3$-periodic, and $\mathrm{Fo}(\mathrm{emb}\,c)=c$ for $c\in\tilde{W}$.
If $u$, $u'$ are nearest neighbours in $\mathrm{Lab}^W$, then either they lie in one chamber of $\mathbb{T}^{W}$, and then $\mathrm{Fo}\,u$, $\mathrm{Fo}\,u'$ are nearest neighbours in $\tilde{W}$, or they are separated by a mirror plane, and then $\mathrm{Fo}\,u=\mathrm{Fo}\,u'$.
\item[(c)] For unit sites $z$, $z'$ in $\tilde{W}$ the distance of $\mathrm{emb}\,z$ and $\mathrm{emb}\,z'$ on $\mathbb{T}^{W}$ is $|z-z'|$; this is also the distance of $z$ and $z'$ on $T_k$.
\end{itemize}
\end{named}
This is \cite[Lemma~16.10]{QEDI}.

\begin{named}{[Q.2]}
The function $A^W$ is a connection on $\mathbb{T}^{W}$; the map $A\mapsto A^W$ is complex linear; $(\overline{A})^W=\overline{A^W}$; $|A^W|_{\infty}=|A|_{\infty}$, the supremum on the right being formed with the bonds of $\tilde{W}$; and $(d\boldsymbol{\lambda})^W=d(\boldsymbol{\lambda}^W)$.
The set $\mathrm{Gloc}(\delta)$ of Section~G1.5 is open: if $A=B+d\boldsymbol{\lambda}$ with $e_k\cdot|B|_{\infty}<\delta$, then $A+A'=(B+A')+d\boldsymbol{\lambda}$ with $e_k\cdot|B+A'|_{\infty}<\delta$ for every $A'$ with $e_k\cdot|A'|_{\infty}<\delta-e_k\cdot|B|_{\infty}$.
\end{named}
These facts are established in the paragraphs that follow \cite[Definition~16.11]{QEDI}.

\begin{named}{[Q.3]}
\begin{itemize}
\item[(G1)] For $A\in\tilde{G}_k$ and every $M$-cube $\square$ there are a real function $\boldsymbol{\lambda}$ on the sites of $\square^{\sim1}$ and a complex function $B$ on the bonds of $\square^{\sim1}$ such that $A=B+d\boldsymbol{\lambda}$ on the bonds of $\square^{\sim1}$ and $e_k\cdot|B|_{\infty}<2e_k^{\varepsilon}$; moreover $\mathfrak{i}(A)<e_k^{\varepsilon}$.
\item[(G2)] The set $\tilde{G}_k$ is open in the space of the connections on $T_k$, and it is invariant under the maps $A\mapsto A-d\boldsymbol{\lambda}$ ($\boldsymbol{\lambda}$ a real scalar on $T_k$, that is, a real $\ell\mathbb{Z}^3$-periodic function of $c$), $A\mapsto-A$, $A\mapsto\overline{A}$ and $A\mapsto\tilde{g}A$ ($\tilde{g}\in\mathbb{G}_k$).
\end{itemize}
The diameter of $\tilde{W}$ in the sup norm, increased by $\eta_k$, equals $3M<\ell$; hence the bonds of $\square^{\sim1}$ are the images of the bonds of $\tilde{W}$ (Section~I.1.1, lifts).
Accordingly, (G1) says that the restriction of $A$ to the bonds of $\tilde{W}$ lies in $\mathrm{Gloc}(2e_k^{\varepsilon})$ (Section~G1.6).
\end{named}
The properties (G1) and (G2) are proved in \cite[Section~16.2.7]{QEDI}.

\begin{named}{[Q.4]}
$\left|U_A(c,c')^{\pm1}\right|\le\exp\left(3|c-c'|\cdot\mathfrak{i}(A)\right)$.
\end{named}
This is the last inequality of \cite[Lemma~16.2(a)]{QEDI}.

\begin{named}{[Q.5]}
By its definition in Section~G1.7, $h_{\square}$ is a real $\ell\mathbb{Z}^3$-periodic function of $c$ with $0\le h_{\square}\le1$, and the second relation of (6.1) in Section~G1.7 says that $h_{\square}(c)\ne0$ only if $|c+\ell n-z_{\square}|<3M/4$ for some $n\in\mathbb{Z}^3$.
The last sentence of (6.1), by which a unit block meets $\operatorname{supp}h_{\square}$ for at most $8$ cubes $\square$, rests on the following observation: for a fixed lift of $z$, a cube $\square$ for which $\operatorname{supp}h_{\square}$ meets $\Delta_z$ has a lift of its centre with $|z_{\mu}-z_{\square,\mu}|<3M/4+1/2$ for $\mu=1,2,3$, and an open interval of length $3M/2+1<2M$ contains at most two multiples of $M$.
\end{named}
The relations (6.1) are those of \cite[Section~16.7.1]{QEDI}.

\subsection*{1. Proposition 2.5 (local inverses on mirror tori)}

\paragraph{Proof of Proposition~2.5.}
The triple $\mathbb{T}^{W}=\left((\eta_k/2)\mathrm{1v},6M,+1\right)$ is a block torus of level $k$, since $6M=6(M/L)L$ is an integer multiple of $L$ and $6M\ge3L$.

\paragraph{Step 0 (the translation by $o$).}
Recall that $o=z_{\square}-(3M/2)\mathrm{1v}$, and let $0\le j\le k+1$.
The point $z_{\square}$ lies in $M\mathbb{Z}^3$, which is contained in $L^j\eta_k\mathbb{Z}^3$, because $M/(L^j\eta_k)=(M/L)L^{k+1-j}$ is an integer for $j\le k+1$.
The number $3M/2$ lies in $L^j\eta_k(\mathbb{Z}+1/2)$, because $3M/(L^j\eta_k)$ is odd by [Q.1]\,(a).
Hence
\begin{equation*}
L^j\eta_k\mathbb{Z}^3-o=\frac{L^j\eta_k}{2}\mathrm{1v}+L^j\eta_k\mathbb{Z}^3=\mathrm{Lab}^W_j\qquad(0\le j\le k+1).\tag*{[5.1]}
\end{equation*}
Consequently (F2) applies with $\mathrm{Lb}=\mathrm{Lab}^W$, $\mathrm{Lb}'=\mathbb{L}_k$ and $v=o$.
The map $T_o$, given by $(T_ov)(c,g)=v(c-o,g)$, takes functions on $\mathrm{Lab}^W\times O_{\mathrm{cub}}$ to functions on $\mathbb{L}^{\mathrm{f}}_k$ and bond functions on $\mathrm{Lab}^W$ to bond functions on $\mathbb{L}_k$; it is injective.
The translation $c\mapsto c-o$ maps $\tilde{W}$ onto
\begin{equation*}
\mathrm{Ch}_0:=\left\{u\in\mathrm{Lab}^W:0<u_{\mu}<3M,\ \mu=1,2,3\right\},
\end{equation*}
and $\mathrm{Ch}$ is the disjoint union of the sets $\mathrm{Ch}_0+6Mn$, $n\in\mathbb{Z}^3$.
Since $\mathrm{fo}(t)=t$ for $0\le t\le3M$ and $\mathrm{fo}$ is $6M$-periodic,
\begin{equation*}
\mathrm{Fo}(c-o+6Mn)=c\qquad(c\in\tilde{W},\ n\in\mathbb{Z}^3).\tag*{[5.2]}
\end{equation*}
For $x=(c,g)$ over $\tilde{W}$ put $\mathrm{emb}\,x:=(c-o,g)$, the representative of $\mathrm{emb}\,c$ in $\mathrm{Ch}_0$ with the label unchanged.
Spinors on $\mathbb{T}^{W}$ are $6M\mathbb{Z}^3$-periodic, so that $(\mathrm{rs}\,v)(x)=v(\mathrm{emb}\,x)$.
Let $\mathrm{res}$ denote the restriction of a function on $\mathbb{L}^{\mathrm{f}}_k$ to $\tilde{W}\times O_{\mathrm{cub}}$, let $\mathrm{cz}$ denote the continuation by zero of a function on $\tilde{W}\times O_{\mathrm{cub}}$ to $\mathbb{L}^{\mathrm{f}}_k$, and put $A^{o}:=T_oA^W$.
Then
\begin{gather*}
A^{o}\text{ is a }6M\mathbb{Z}^3\text{-periodic antisymmetric function on the bonds of }\mathbb{L}_k,\\
\text{and }A^{o}=A\text{ on the bonds of }\tilde{W};\tag*{[5.3]}\\
\mathrm{rs}\,v=\mathrm{res}\,T_ov,\qquad T_o\,\mathrm{ext}\,f=\sum_{n\in\mathbb{Z}^3}T_{6Mn}\,\mathrm{cz}\,f,\qquad\mathrm{rs}\cdot\mathrm{ext}=1;\tag*{[5.4]}\\
T_oF^{W}_y(A)v=F_y(A^{o})\,T_ov\quad\text{for every spinor }v\text{ on }\mathbb{T}^{W};\tag*{[5.5]}
\end{gather*}
in the sum of [5.4] the supports $\tilde{W}+6Mn$ of the terms are pairwise disjoint.

To prove [5.3], let $(c,c')$ be a bond of $\tilde{W}$; then $A^{o}(c,c')=A^W(c-o,c'-o)=A\left(\mathrm{Fo}(c-o),\mathrm{Fo}(c'-o)\right)=A(c,c')$ by [5.2].
Periodicity and antisymmetry follow from [Q.2].
In [5.4] the first relation is the definition of $\mathrm{rs}$.
For the second, $(T_o\,\mathrm{ext}\,f)(c,g)=(\mathrm{ext}\,f)(c-o,g)$ is equal to $f\left(\mathrm{Fo}(c-o),g\right)$ if $c-o\in\mathrm{Ch}$, and to $0$ otherwise.
The point $c-o$ lies in $\mathrm{Ch}$ if and only if $c-6Mn\in\tilde{W}$ for some $n\in\mathbb{Z}^3$, and this $n$ is then unique; moreover $\mathrm{Fo}(c-o)=c-6Mn$ by [5.2].
So the value is $(\mathrm{cz}\,f)(c-6Mn,g)=(T_{6Mn}\,\mathrm{cz}\,f)(c,g)$ if $c-o\in\mathrm{Ch}$, and otherwise every term of the sum vanishes at $(c,g)$.
The third relation is obtained by applying $\mathrm{res}$ to the second.
Finally, [5.5] follows from (F3) on $\mathbb{T}^{W}$ together with (F2).

\paragraph{Step 1 (reduction to a connection $B^W$ with $e_k\cdot|B^W|_{\infty}<\delta$ on the whole mirror torus).}
Let $A=B+d\boldsymbol{\lambda}$ with $\boldsymbol{\lambda}$ real on $\tilde{W}$ and $e_k\cdot|B|_{\infty}<\delta$.
By [Q.2], $A^W=B^W+d(\boldsymbol{\lambda}^W)$ and $e_k\cdot|B^W|_{\infty}=e_k\cdot|B|_{\infty}<\delta$; by [Q.1]\,(b), $\boldsymbol{\lambda}^W=\boldsymbol{\lambda}\circ\mathrm{Fo}$ is a real scalar on $\mathbb{T}^{W}$.
Put $\mathrm{uf}:=\exp\left(ie_k\boldsymbol{\lambda}^W\right)$, a $6M\mathbb{Z}^3$-periodic function of modulus one.
Proposition~2.2\,(a) on $\mathbb{T}^{W}$, applied to the connection $A^W$, to $\mathrm{Acl}:=d(\boldsymbol{\lambda}^W)$ and to this function $\mathrm{uf}$ (the example given there), yields
\begin{gather*}
F_y\left(A^W-d\boldsymbol{\lambda}^W\right)=\mathrm{uf}\cdot F_y(A^W)\cdot\mathrm{uf}^{-1},\\
\left(\delta_{\alpha,A^W-d\boldsymbol{\lambda}^W}(\mathrm{uf}\,f)\right)(t,t')=\mathrm{uf}(t)\left(\delta_{\alpha,A^W}f\right)(t,t'),
\end{gather*}
where $\mathrm{uf}(t)$ is the value of $\mathrm{uf}$ at the photon site of $t$.
(At $k=0$ the relation for $F_y=D+\mu+b_LJPJ^{*}+iy\gamma_5$ is the sum of the relations given there for $D_A$ and for $JP(A)J^{*}$, since $\mu$ and $iy\gamma_5$ commute with the scalar function $\mathrm{uf}$.)
As $A^W-d\boldsymbol{\lambda}^W=B^W$, this means
\begin{gather*}
F^{W}_y(A)=\mathrm{uf}^{-1}\cdot F_y(B^W)\cdot\mathrm{uf},\\
\left(\delta_{\alpha,A^W}(\mathrm{uf}^{-1}v)\right)(t,t')=\mathrm{uf}(t)^{-1}\left(\delta_{\alpha,B^W}v\right)(t,t')\tag*{[5.6]}
\end{gather*}
for every spinor $v$ on $\mathbb{T}^{W}$ and every short pair $(t,t')$.

\paragraph{(a)}
Proposition~2.4\,(a) on $\mathbb{T}^{W}$, with the data $e_k>0$, $\mu\in(0,1]$, $y\in Y_k$, $\delta$ and the connection $B^W$, for which $e_k\cdot|B^W|_{\infty}\le\delta$, shows that $F_y(B^W)$ is bijective.
By [5.6], $F^{W}_y(A)$ is bijective as well, and
\begin{gather*}
F^{W}_y(A)^{-1}=\mathrm{uf}^{-1}\cdot S_y(B^W)\cdot\mathrm{uf},\\
\mathrm{Sloc}_{\square,y}(A)=\mathrm{rs}\cdot\mathrm{uf}^{-1}\cdot S_y(B^W)\cdot\mathrm{uf}\cdot\mathrm{ext}.\tag*{[5.7]}
\end{gather*}
By (F4), the entries of the matrix of $F^{W}_y(A)=F_y(A^W)$ are entire functions of the numbers $A^W(b')$, and these are linear in the numbers $A(b)$, $b$ a bond of $\tilde{W}$, by [Q.2]; for real $y$ the entries are continuous in $(y,A)$.
The determinant does not vanish on $\mathrm{Gloc}(\delta)$, for any $y\in Y_k$.
By Cramer's rule \cite[Section~0.8]{HornJohnson2013}, the entries of the inverse are polynomials in the entries, divided by the determinant.
Therefore $A\mapsto F^{W}_y(A)^{-1}$ is analytic on $\mathrm{Gloc}(\delta)$, which is open by [Q.2], and continuous on $\mathbb{R}\times\mathrm{Gloc}(\delta)$.
Since $\mathrm{rs}$ and $\mathrm{ext}$ do not depend on $(y,A)$, the same holds for $\mathrm{Sloc}_{\square,y}(A)=\mathrm{rs}\cdot F^{W}_y(A)^{-1}\cdot\mathrm{ext}$.

\paragraph{(b)}
\emph{Localisation.}
By [Q.1]\,(a), $\tilde{W}$ is a union of blocks of side $L$, hence, the blocks being nested, a union of unit blocks.
Let $z_1$, $z_2$ be unit sites of $\tilde{W}$.
By [5.1], $\Delta_{z_1}-o=\Delta_{z_1-o}$ and $\Delta_{z_2}-o=\Delta_{z_2-o}$ are unit blocks of $\mathbb{T}^{W}$, and they are contained in $\mathrm{Ch}_0$.
If $(x,x')$ is a short pair with both points over $\tilde{W}$, $x=(c,g)$, $x'=(c',g')$, then $(\mathrm{emb}\,x,\mathrm{emb}\,x')$ is a short pair of $\mathbb{T}^{W}$, since $d_X(\mathrm{emb}\,x,\mathrm{emb}\,x')=\max\left(|c-c'|,\eta_k\right)=d_X(x,x')$.
For a function $\mathrm{Gp}$ of short pairs on $\mathbb{T}^{W}$ let $(\mathrm{rs2}\,\mathrm{Gp})(x,x'):=\mathrm{Gp}(\mathrm{emb}\,x,\mathrm{emb}\,x')$; this is a function on the short pairs with both points over $\tilde{W}$.
Then
\begin{gather*}
\mathbf{1}_{\Delta_{z_1}}\cdot\mathrm{rs}=\mathrm{rs}\cdot\mathbf{1}_{\Delta_{z_1-o}},\qquad\mathbf{1}_{\Delta_{z_1}}\cdot\mathrm{rs2}=\mathrm{rs2}\cdot\mathbf{1}_{\Delta_{z_1-o}},\qquad\mathrm{ext}\cdot\mathbf{1}_{\Delta_{z_2}}=\mathbf{1}_{\Delta_{z_2-o}}\cdot\mathrm{ext};\\
|\mathrm{rs}\,v|_{\mathrm{su}}\le|v|_{\mathrm{su}},\qquad|\mathrm{rs2}\,\mathrm{Gp}|_{\mathrm{su}}\le|\mathrm{Gp}|_{\mathrm{su}},\qquad|\mathrm{ext}\,f|_{\mathrm{su}}=|f|_{\mathrm{su}}.\tag*{[5.8]}
\end{gather*}

\emph{Proof of [5.8].}
For $c\in\tilde{W}$ the point $c-o$ lies in $\mathrm{Ch}_0$.
Both $\mathrm{Ch}_0$ and $\Delta_{z_1-o}$ lie in the open cube $(0,3M)^3$; hence $c-o$ belongs to the orbit $\Delta_{z_1-o}+6M\mathbb{Z}^3$ if and only if it belongs to $\Delta_{z_1-o}$, that is, if and only if $c\in\Delta_{z_1}$.
Therefore $(\mathrm{rs}\,\mathbf{1}_{\Delta_{z_1-o}}v)(c,g)=\mathbf{1}[c\in\Delta_{z_1}]\,v(c-o,g)$ and, the indicator restricting the first point of a pair (Section~G1.3), $(\mathrm{rs2}\,\mathbf{1}_{\Delta_{z_1-o}}\mathrm{Gp})(x,x')=\mathbf{1}[c\in\Delta_{z_1}]\,\mathrm{Gp}(\mathrm{emb}\,x,\mathrm{emb}\,x')$.
In the third relation both sides vanish at $(u,g)$ with $u\notin\mathrm{Ch}$, and at $u=c-o+6Mn$, $c\in\tilde{W}$, both are equal to $\mathbf{1}[c\in\Delta_{z_2}]\,f(c,g)$, by [5.2].
The values of $\mathrm{rs}\,v$ and of $\mathrm{rs2}\,\mathrm{Gp}$ are values of $v$ and of $\mathrm{Gp}$, while the values of $\mathrm{ext}\,f$ are the values of $f$ and $0$; this gives the three bounds.

Let $\mathrm{Rs}$ stand for $\mathrm{rs}$, respectively for $\mathrm{rs2}$, and let $O^W$ be an operator on the spinors of $\mathbb{T}^{W}$ with values in the spinors, respectively in the functions of short pairs on $\mathbb{T}^{W}$.
By [5.8],
\begin{align*}
\left\|\mathbf{1}_{\Delta_{z_1}}\cdot\mathrm{Rs}\cdot O^W\cdot\mathrm{ext}\cdot\mathbf{1}_{\Delta_{z_2}}\right\|_{\mathrm{su}}&=\left\|\mathrm{Rs}\cdot\left(\mathbf{1}_{\Delta_{z_1-o}}O^W\mathbf{1}_{\Delta_{z_2-o}}\right)\cdot\mathrm{ext}\right\|_{\mathrm{su}}\\
&\le\left\|\mathbf{1}_{\Delta_{z_1-o}}O^W\mathbf{1}_{\Delta_{z_2-o}}\right\|_{\mathrm{su}}.
\end{align*}
By [Q.1]\,(c) the distance of $z_1-o$ and $z_2-o$ on $\mathbb{T}^{W}$ is $|z_1-z_2|$, and $z\mapsto z-o$ is injective from the unit sites of $\tilde{W}$ to the unit sites of $\mathbb{T}^{W}$ modulo $6M\mathbb{Z}^3$.
Hence the sum $\mathrm{Rw}_{z_1}(\mathrm{Rs}\cdot O^W\cdot\mathrm{ext})$ of [P.6], formed over $\tilde{W}$, is bounded by the part of the sum $\mathrm{Rw}_{z_1-o}(O^W)$, formed on $\mathbb{T}^{W}$, that consists of the terms with the indices $z_2-o$, $z_2\in\tilde{W}$:
\begin{gather*}
\mathrm{Rw}_{z_1}\left(\mathrm{Rs}\cdot O^W\cdot\mathrm{ext}\right)\le\mathrm{Rw}_{z_1-o}\left(O^W\right),\\
N_{\gamma}\left(\mathrm{Rs}\cdot O^W\cdot\mathrm{ext}\right)\le N_{\gamma}\left(O^W\right)\quad\text{for every }\gamma\ge0.\tag*{[5.9]}
\end{gather*}
The functions $\mathrm{uf}$ and $\mathrm{uf}^{-1}$ are functions of $c$ of modulus one; they commute with the indicators (on functions of pairs they multiply by the value at the first point), and they leave $|\cdot|_{\mathrm{su}}$ unchanged.
Hence, for an operator $O^W$ as above and unit sites $u$, $u'$ of $\mathbb{T}^{W}$,
\begin{gather*}
\left\|\mathbf{1}_{\Delta_u}\cdot\mathrm{uf}^{-1}\cdot O^W\cdot\mathrm{uf}\cdot\mathbf{1}_{\Delta_{u'}}\right\|_{\mathrm{su}}=\left\|\mathrm{uf}^{-1}\cdot\left(\mathbf{1}_{\Delta_u}O^W\mathbf{1}_{\Delta_{u'}}\right)\cdot\mathrm{uf}\right\|_{\mathrm{su}}=\left\|\mathbf{1}_{\Delta_u}O^W\mathbf{1}_{\Delta_{u'}}\right\|_{\mathrm{su}},\\
N_{\gamma}\left(\mathrm{uf}^{-1}\cdot O^W\cdot\mathrm{uf}\right)=N_{\gamma}\left(O^W\right).\tag*{[5.10]}
\end{gather*}

\emph{First bound.}
By [5.7], by [5.9] with $\mathrm{Rs}=\mathrm{rs}$, by [5.10] and by Proposition~2.4\,(b),
\begin{equation*}
N_{\gamma^{F}}\left(\mathrm{Sloc}_{\square,y}(A)\right)\le N_{\gamma^{F}}\left(\mathrm{uf}^{-1}\cdot S_y(B^W)\cdot\mathrm{uf}\right)=N_{\gamma^{F}}\left(S_y(B^W)\right)\le2a'\cdot w_k(y).
\end{equation*}

\emph{Second bound ($k\ge1$).}
Let $f_0$ be a function on $\tilde{W}\times O_{\mathrm{cub}}$, and let $(x,x')$ be a short pair with both points over $\tilde{W}$.
The quantity $(\delta_{\alpha,A}f_0)(x,x')=d_X(x,x')^{-\alpha}\left(U_A(x,x')f_0(x')-f_0(x)\right)$ is given by the display (I.1.7): the six paths $\boldsymbol{\Gamma}^{\pi}(c,c')$ that enter $U_A(x,x')=\exp\left\{ie_k\eta_k(\tau A)(c,c')\right\}$ stay in the lattice box spanned by $c$ and $c'$, which is contained in the lattice box $\tilde{W}$, so that $U_A(x,x')$ involves $A$ only on bonds of $\tilde{W}$.
On these bonds $A^{o}=A$ by [5.3]; hence $U_{A^{o}}(x,x')=U_A(x,x')$.
(For $c=c'$, that is, for a pair inside a cluster, both are equal to $1$.)
Let $v$ be a spinor on $\mathbb{T}^{W}$.
By (F2) for the translation by $o$, and then by the relation $\mathrm{rs}\,v=\mathrm{res}\,T_ov$ of [5.4],
\begin{align*}
\left(\delta_{\alpha,A^W}v\right)(\mathrm{emb}\,x,\mathrm{emb}\,x')&=\left(\delta_{\alpha,A^{o}}T_ov\right)(x,x')\\
&=d_X(x,x')^{-\alpha}\left(U_A(x,x')\,(\mathrm{rs}\,v)(x')-(\mathrm{rs}\,v)(x)\right)\\
&=\left(\delta_{\alpha,A}\,\mathrm{rs}\,v\right)(x,x'),\qquad\text{that is,}\qquad\delta_{\alpha,A}\cdot\mathrm{rs}=\mathrm{rs2}\cdot\delta_{\alpha,A^W}.\tag*{[5.11]}
\end{align*}
By [5.7], [5.11] and [5.6],
\begin{align*}
\delta_{\alpha,A}\,\mathrm{Sloc}_{\square,y}(A)&=\mathrm{rs2}\cdot\delta_{\alpha,A^W}\cdot\mathrm{uf}^{-1}\cdot S_y(B^W)\cdot\mathrm{uf}\cdot\mathrm{ext}\\
&=\mathrm{rs2}\cdot\mathrm{uf}^{-1}\cdot\left[\delta_{\alpha,B^W}S_y(B^W)\right]\cdot\mathrm{uf}\cdot\mathrm{ext}.\tag*{[5.12]}
\end{align*}
By [5.9] with $\mathrm{Rs}=\mathrm{rs2}$, by [5.10] and by Proposition~2.4\,(c),
\begin{equation*}
N_{\gamma^{F}}\left(\delta_{\alpha,A}\,\mathrm{Sloc}_{\square,y}(A)\right)\le N_{\gamma^{F}}\left(\delta_{\alpha,B^W}S_y(B^W)\right)\le18a'.
\end{equation*}

\emph{Level $0$.}
Here $d_X\ge1$ for all pairs, so there is no short pair and the second bound is void.
The unit sites of $\tilde{W}$ are the sites of $\tilde{W}$, and $\Delta_z\times O_{\mathrm{cub}}$ is the cluster over $z$, consisting of $48$ sites of weight $1/48$.
Let $O^W$ be an operator on the spinors of $\mathbb{T}^{W}$ with the kernel $O^W(t,t')$ of Section~G1.2.
Choose the period cell $\{u\in\mathrm{Lab}^W:0<u_{\mu}<6M\}$; its intersection with $\mathrm{Ch}$ is $\mathrm{Ch}_0$.
For $f$ on $\tilde{W}\times O_{\mathrm{cub}}$ and $x$ over $\tilde{W}$, [5.2] gives
\begin{align*}
\left(\mathrm{rs}\cdot O^W\cdot\mathrm{ext}\,f\right)(x)&=\frac{1}{48}\sum_{u'\in\mathrm{Ch}_0}\ \sum_{g'\in O_{\mathrm{cub}}}O^W\left(\mathrm{emb}\,x,(u',g')\right)f(\mathrm{Fo}\,u',g')\\
&=\frac{1}{48}\sum_{c'\in\tilde{W}}\ \sum_{g'\in O_{\mathrm{cub}}}O^W\left(\mathrm{emb}\,x,\mathrm{emb}(c',g')\right)f(c',g').
\end{align*}
Hence the kernel of $\mathrm{rs}\cdot O^W\cdot\mathrm{ext}$ with respect to $\int_{\tilde{W}}\mathrm{d}x$ (Section~G1.6) is
\begin{equation*}
\left(\mathrm{rs}\cdot O^W\cdot\mathrm{ext}\right)(x,x')=O^W(\mathrm{emb}\,x,\mathrm{emb}\,x')\qquad(x,\ x'\text{ over }\tilde{W}).\tag*{[5.13]}
\end{equation*}
For sites $z_1$, $z_2$ of $\tilde{W}$ it follows that
\begin{align*}
&\sup\left\{\left|\left(\mathrm{rs}\cdot O^W\cdot\mathrm{ext}\right)(x,x')\right|_2:x\text{ over }z_1,\ x'\text{ over }z_2\right\}\\
&\qquad=\sup\left\{\left|O^W(t,t')\right|_2:t\text{ over }z_1-o,\ t'\text{ over }z_2-o\right\}.
\end{align*}
Since $d(z_1-o,z_2-o)=|z_1-z_2|$ and $z\mapsto z-o$ is injective, $\mathrm{Rk}_{z_1}(\mathrm{rs}\cdot O^W\cdot\mathrm{ext})\le\mathrm{Rk}_{z_1-o}(O^W)$, and therefore $N^{\mathrm{ke}}_{\gamma}(\mathrm{rs}\cdot O^W\cdot\mathrm{ext})\le N^{\mathrm{ke}}_{\gamma}(O^W)$.
The kernel of $\mathrm{uf}^{-1}\cdot S_y(B^W)\cdot\mathrm{uf}$ is $\mathrm{uf}(t)^{-1}\cdot S_y(B^W)(t,t')\cdot\mathrm{uf}(t')$ (Section~G1.2), and its modulus $|\cdot|_2$ equals that of $S_y(B^W)(t,t')$.
Together with [5.7] and Proposition~2.4\,(d) this yields
\begin{equation*}
N^{\mathrm{ke}}_{\gamma^{F}}\left(\mathrm{Sloc}_{\square,y}(A)\right)\le N^{\mathrm{ke}}_{\gamma^{F}}\left(S_y(B^W)\right)\le2a'\cdot w_0(y).
\end{equation*}

\paragraph{(c)}
Let $\tilde{A}$ be an antisymmetric function on the bonds of $\mathbb{L}_k$ that coincides with $A$ on the bonds of $\tilde{W}$, and let $F_y(\tilde{A})$ be the operator of Section~0.2 on the functions on $\mathbb{L}^{\mathrm{f}}_k$.

\emph{First assertion.}
Let $x=(c,g)$ with $c\in W^{-}$.
Then $|c_3-z_{\square}|\le|c-z_{\square}|+\eta_k<3M/2$; thus $c_3\in\tilde{W}$, and $(c,c_3)$ is a bond of $\tilde{W}$.
The block of side $L$ that contains $c$ is contained in $\tilde{W}$, because $\tilde{W}$ is a union of blocks of side $L$ by [Q.1]\,(a).
By (F1) the row at $x$ reads $f$ at $(c,g)$, $(c,gs_0)$, $(c,gs_1)$, $(c,gs_2)$, $(c_3,gs_3)$ and over this block, that is, over $\tilde{W}$, and it reads $\tilde{A}$ on bonds of $\tilde{W}$ only.

\emph{Second assertion.}
Let $f$ be a function on $\tilde{W}\times O_{\mathrm{cub}}$, and put $v':=F^{W}_y(A)^{-1}\mathrm{ext}\,f$ and
\begin{equation*}
v:=\mathrm{cz}\cdot\mathrm{Sloc}_{\square,y}(A)f=\mathrm{cz}\cdot\mathrm{rs}\,v'=\mathrm{cz}\cdot\mathrm{res}\,T_ov',
\end{equation*}
the last equality by [5.4].
The functions $v$ and $T_ov'$ coincide over $\tilde{W}$, and $\tilde{A}$ and $A^{o}$ coincide on the bonds of $\tilde{W}$ by [5.3].
By the first assertion, [5.5] and [5.4], for $x$ over $W^{-}$,
\begin{equation*}
\left(F_y(\tilde{A})v\right)(x)=\left(F_y(A^{o})T_ov'\right)(x)=\left(T_oF^{W}_y(A)v'\right)(x)=\left(T_o\,\mathrm{ext}\,f\right)(x)=f(x).
\end{equation*}

\emph{Third assertion.}
Let $\mathrm{fs}$ be a function on $\mathbb{L}^{\mathrm{f}}_k$ that vanishes off $W^{-}\times O_{\mathrm{cub}}$, and let $x=(c,g)$.

\emph{Case $c\notin\tilde{W}$.}
The block of side $L$ that contains $c$ is disjoint from $\tilde{W}$, since the blocks of side $L$ partition $\mathbb{L}_k$ and $\tilde{W}$ is a union of such blocks; so the sums over blocks in [P.4] and [P.5] vanish.
The function $\mathrm{fs}$ vanishes at the four sites over $c$.
What remains is
\begin{equation*}
\left(F_y(\tilde{A})\mathrm{fs}\right)(x)=-\eta_k^{-1}M_3(x)\,u_{\tilde{A}}(c,c_3)\,\mathrm{fs}(c_3,gs_3).
\end{equation*}
If $c_3$ were in $W^{-}$, then $|c-z_{\square}|\le|c_3-z_{\square}|+\eta_k<3M/2$, that is, $c\in\tilde{W}$.
Hence $\mathrm{fs}(c_3,gs_3)=0$, and $\left(F_y(\tilde{A})\mathrm{fs}\right)(x)=0$.

\emph{Case $c\in\tilde{W}$.}
The block of side $L$ that contains $c$ lies in $\tilde{W}$, so the phases in the sums over blocks read $\tilde{A}$ on bonds of $\tilde{W}$ only.
If $c_3\in\tilde{W}$, then $(c,c_3)$ is a bond of $\tilde{W}$.
If $c_3\notin\tilde{W}$, then $\mathrm{fs}(c_3,gs_3)=0$, and the product $u_{\tilde{A}}(c,c_3)\,\mathrm{fs}(c_3,gs_3)$ vanishes for every $\tilde{A}$.

Consequently $F_y(\tilde{A})\mathrm{fs}$ vanishes off $\tilde{W}\times O_{\mathrm{cub}}$, and it is one and the same function for all continuations $\tilde{A}$ of $A$; in particular

\begin{equation}
F_y(\tilde{A})\,\mathrm{fs}=F_y(A^{o})\,\mathrm{fs}=\mathrm{cz}\,\mathrm{res}\,F_y(\tilde{A})\,\mathrm{fs}.
\tag*{[5.14]}
\end{equation}
For $n\in\mathbb{Z}^3$ and every $j\le k+1$ the vector $6Mn$ lies in $L^j\eta_k\mathbb{Z}^3$, and $T_{6Mn}A^{o}=A^{o}$;
by (F2), the operator $F_y(A^{o})$ therefore commutes with $T_{6Mn}$.
The rows of $F_y(A^{o})$ have finite support, and the sums in [5.4] are locally finite.
Applying in turn [5.5], then [5.4] to $\mathrm{res}\,\mathrm{fs}$ (note that $\mathrm{cz}\,\mathrm{res}\,\mathrm{fs}=\mathrm{fs}$), then [5.14], and finally [5.4] to $\mathrm{res}\,F_y(\tilde{A})\,\mathrm{fs}$, one obtains
\begin{align*}
T_o\,F^{W}_{y}(A)\,\mathrm{ext}\,\mathrm{res}\,\mathrm{fs}
&=F_y(A^{o})\sum_{n\in\mathbb{Z}^3}T_{6Mn}\,\mathrm{fs}
=\sum_{n\in\mathbb{Z}^3}T_{6Mn}\,F_y(A^{o})\,\mathrm{fs}\\
&=\sum_{n\in\mathbb{Z}^3}T_{6Mn}\,\mathrm{cz}\,\mathrm{res}\,F_y(\tilde{A})\,\mathrm{fs}
=T_o\,\mathrm{ext}\,\mathrm{res}\,F_y(\tilde{A})\,\mathrm{fs}.
\end{align*}
Since $T_o$ is injective, it follows that $F^{W}_{y}(A)\,\mathrm{ext}\,\mathrm{res}\,\mathrm{fs}=\mathrm{ext}\,\mathrm{res}\,F_y(\tilde{A})\,\mathrm{fs}$.
Applying $\mathrm{rs}\,F^{W}_{y}(A)^{-1}$ to both sides and using $\mathrm{rs}\,\mathrm{ext}=1$, one arrives at
\[
\mathrm{res}\,\mathrm{fs}=\mathrm{Sloc}_{\square,y}(A)\,\mathrm{res}\,F_y(\tilde{A})\,\mathrm{fs}.
\]

\textbf{(d)}
The function $A-d\boldsymbol{\lambda}$ belongs to $\mathrm{Gloc}(\delta)$, and $(A-d\boldsymbol{\lambda})^{W}=A^{W}-d(\boldsymbol{\lambda}^{W})$ by [Q.2].
Put $\mathrm{uf}_1:=\exp\left(ie_k\boldsymbol{\lambda}^{W}\right)$.
Proposition~2.2\,(a) on $\mathbb{T}^{W}$ gives $F_y\left(A^{W}-d\boldsymbol{\lambda}^{W}\right)=\mathrm{uf}_1\,F_y(A^{W})\,\mathrm{uf}_1^{-1}$, and hence
\[
F^{W}_{y}(A-d\boldsymbol{\lambda})^{-1}=\mathrm{uf}_1\,F^{W}_{y}(A)^{-1}\,\mathrm{uf}_1^{-1}.
\]
For a spinor $v$ on $\mathbb{T}^{W}$ and a function $f$ on $\tilde{W}\times O_{\mathrm{cub}}$, [5.2] gives
\begin{align*}
(\mathrm{rs}\,\mathrm{uf}_1v)(c,g)&=\exp\left(ie_k\boldsymbol{\lambda}\left(\mathrm{Fo}(c-o)\right)\right)v(c-o,g)=\exp\left(ie_k\boldsymbol{\lambda}(c)\right)(\mathrm{rs}\,v)(c,g),\\
(\mathrm{uf}_1^{-1}\,\mathrm{ext}\,f)(u,g)&=\mathbf{1}[u\in\mathrm{Ch}]\,\exp\left(-ie_k\boldsymbol{\lambda}(\mathrm{Fo}\,u)\right)f(\mathrm{Fo}\,u,g)=\left(\mathrm{ext}\,\exp\left(-ie_k\boldsymbol{\lambda}\right)f\right)(u,g).
\end{align*}
Consequently
\begin{align*}
\mathrm{Sloc}_{\square,y}(A-d\boldsymbol{\lambda})&=\mathrm{rs}\,\mathrm{uf}_1\,F^{W}_{y}(A)^{-1}\,\mathrm{uf}_1^{-1}\,\mathrm{ext}\\
&=\exp\left(ie_k\boldsymbol{\lambda}\right)\mathrm{Sloc}_{\square,y}(A)\exp\left(-ie_k\boldsymbol{\lambda}\right).
\end{align*}

\textbf{(e)}
Let $\tilde{g}=(S,a)$, where $a\in M\mathbb{Z}^3$ is a representative of the translation part, and put $r:=r(S)\in\{\pm1\}^3$, so that $r^{-1}=r$, and $\mathrm{mph}(c):=rc+a$.
The map $\mathrm{mph}$ sends $\mathbb{L}_k$ onto itself and preserves the distance $|c-c'|$.
Hence $\tilde{g}\square$ is the $M$-cube whose centre has the lift $z_{\tilde{g}\square}:=rz_{\square}+a\in M\mathbb{Z}^3$, and for this lift
\[
\tilde{W}':=\left\{c:\ |c-z_{\tilde{g}\square}|<\tfrac{3M}{2}\right\}=\mathrm{mph}(\tilde{W})=r\tilde{W}+a,\qquad o':=z_{\tilde{g}\square}-\tfrac{3M}{2}(1,1,1).
\]
The maps $\mathrm{Fo}'$, $\mathrm{emb}'$, $\mathrm{ext}'$ and $\mathrm{rs}'$ are formed with $o'$; the mirror torus is the same block torus $\mathbb{T}^{W}$.
For $c\in\tilde{W}'$ put $(\tilde{g}\boldsymbol{\lambda})(c):=\boldsymbol{\lambda}(\mathrm{mph}^{-1}c)$.
Then
\[
\left(d(\tilde{g}\boldsymbol{\lambda})\right)(c,c')=\frac{\boldsymbol{\lambda}(\mathrm{mph}^{-1}c')-\boldsymbol{\lambda}(\mathrm{mph}^{-1}c)}{\eta_k}=(\tilde{g}\,d\boldsymbol{\lambda})(c,c'),
\]
so that $\tilde{g}A=\tilde{g}B+d(\tilde{g}\boldsymbol{\lambda})$, with $\tilde{g}\boldsymbol{\lambda}$ real and $|\tilde{g}B|_\infty=|B|_\infty$;
thus $\tilde{g}A$ lies in the set $\mathrm{Gloc}(\delta)$ of the cube $\tilde{g}\square$.
Define
\[
\mathrm{mph}^{W}(u):=r\left(u-\tfrac{3M}{2}(1,1,1)\right)+\tfrac{3M}{2}(1,1,1);
\]
in the coordinate $\mu$ this is the map $u_\mu\mapsto u_\mu$ if $r_\mu=1$, and $u_\mu\mapsto3M-u_\mu$ if $r_\mu=-1$.
For $0\le j\le k+1$ the map $\mathrm{mph}^{W}$ sends $\mathrm{Lab}^{W}_{j}=L^j\eta_k\left(\mathbb{Z}+\tfrac12\right)^3$ onto itself.
Indeed, $3M$ is an odd multiple of $L^j\eta_k$ by [Q.1]\,(a), and for $n\in\mathbb{Z}$
\[
3M-L^j\eta_k\left(n+\frac12\right)=L^j\eta_k\left(\frac{3M}{L^j\eta_k}-1-n+\frac12\right).
\]
Hence $\mathrm{mph}^{W}$ is a morphism of $\mathbb{T}^{W}$ onto itself with linear part $r$ (Section~G1.2).
On fermion sites it acts by $\mathrm{mph}^{W}(u,g)=(\mathrm{mph}^{W}u,rg)$; let $\mathrm{mS}^{W}:=(S,\mathrm{mph}^{W})$.
The map $\mathrm{mph}^{W}$ sends $\mathrm{Ch}$ onto $\mathrm{Ch}$, because $0<u_\mu<3M$ holds if and only if $0<3M-u_\mu<3M$ (modulo $6M$).
For every $c\in\mathbb{R}^3$,
\begin{equation}
\mathrm{mph}^{W}(c-o)=r(c-z_{\square})+\tfrac{3M}{2}(1,1,1)=\mathrm{mph}(c)-o'.
\tag*{[5.15]}
\end{equation}
Every real $t$ satisfies $\operatorname{dist}(t,6M\mathbb{Z})+\operatorname{dist}(t,3M+6M\mathbb{Z})=3M$;
hence $\mathrm{fo}(3M-t)=\operatorname{dist}(t-3M,6M\mathbb{Z})=3M-\mathrm{fo}(t)$.
With the notation $\mathrm{fo}(u):=\left(\mathrm{fo}(u_1),\mathrm{fo}(u_2),\mathrm{fo}(u_3)\right)$ this says that $\mathrm{fo}(\mathrm{mph}^{W}u)=\mathrm{mph}^{W}\left(\mathrm{fo}(u)\right)$.
Taking $c=o+\mathrm{fo}(u)=\mathrm{Fo}\,u$ in [5.15], one obtains
\begin{equation}
\mathrm{Fo}'(\mathrm{mph}^{W}u)=o'+\mathrm{mph}^{W}\left(\mathrm{fo}(u)\right)=\mathrm{mph}(\mathrm{Fo}\,u).
\tag*{[5.16]}
\end{equation}
Now let $f$ be a function on $\tilde{W}\times O_{\mathrm{cub}}$, $v$ a spinor on $\mathbb{T}^{W}$, and $u,u'$ nearest neighbours in $\mathrm{Lab}^{W}$;
recall that $\mathrm{mph}$ and $\mathrm{mph}^{W}$ are bijections and that $(\mathrm{mph}^{W})^{-1}=\mathrm{mph}^{W}$.
If $\mathrm{Fo}'u\neq\mathrm{Fo}'u'$, then
\begin{align*}
(\tilde{g}A)^{W'}(u,u')&=(\tilde{g}A)(\mathrm{Fo}'u,\mathrm{Fo}'u')=A\left(\mathrm{mph}^{-1}\mathrm{Fo}'u,\mathrm{mph}^{-1}\mathrm{Fo}'u'\right)\\
&=A\left(\mathrm{Fo}\,(\mathrm{mph}^{W})^{-1}u,\mathrm{Fo}\,(\mathrm{mph}^{W})^{-1}u'\right)=(\mathrm{mS}^{W}A^{W})(u,u');
\end{align*}
if $\mathrm{Fo}'u=\mathrm{Fo}'u'$, which by [5.16] is equivalent to $\mathrm{Fo}\,(\mathrm{mph}^{W})^{-1}u=\mathrm{Fo}\,(\mathrm{mph}^{W})^{-1}u'$, both sides vanish.
Further,
\begin{align*}
(\mathrm{ext}'\,\tilde{g}f)(u,g)&=\mathbf{1}[u\in\mathrm{Ch}]\,Sf\left(\mathrm{mph}^{-1}\mathrm{Fo}'u,rg\right)\\
&=\mathbf{1}\left[(\mathrm{mph}^{W})^{-1}u\in\mathrm{Ch}\right]Sf\left(\mathrm{Fo}\,(\mathrm{mph}^{W})^{-1}u,rg\right)\\
&=S(\mathrm{ext}\,f)\left((\mathrm{mph}^{W})^{-1}u,rg\right)=(\mathrm{mS}^{W}\,\mathrm{ext}\,f)(u,g),
\end{align*}
and, for $c\in\tilde{W}'$, by [5.15],
\begin{align*}
(\mathrm{rs}'\,\mathrm{mS}^{W}v)(c,g)&=Sv\left((\mathrm{mph}^{W})^{-1}(c-o'),rg\right)=Sv\left(\mathrm{mph}^{-1}c-o,rg\right)\\
&=S(\mathrm{rs}\,v)\left(\mathrm{mph}^{-1}c,rg\right)=(\tilde{g}\,\mathrm{rs}\,v)(c,g).
\end{align*}
In operator form,
\begin{equation}
(\tilde{g}A)^{W'}=\mathrm{mS}^{W}(A^{W}),\qquad\mathrm{ext}'\,\tilde{g}=\mathrm{mS}^{W}\,\mathrm{ext},\qquad\mathrm{rs}'\,\mathrm{mS}^{W}=\tilde{g}\,\mathrm{rs}.
\tag*{[5.17]}
\end{equation}
Proposition~2.2\,(b), applied to the morphism $\mathrm{mph}^{W}$ of $\mathbb{T}^{W}$, gives $F_{y'}(\mathrm{mS}^{W}A^{W})\,\mathrm{mS}^{W}=\mathrm{mS}^{W}F_y(A^{W})$, with $y'$ as in (e).
Both operators $F$ in this identity are bijective by (a), whence $F_{y'}(\mathrm{mS}^{W}A^{W})^{-1}\,\mathrm{mS}^{W}=\mathrm{mS}^{W}F_y(A^{W})^{-1}$.
Together with [5.17] this yields
\begin{align*}
\mathrm{Sloc}_{\tilde{g}\square,y'}(\tilde{g}A)\,\tilde{g}
&=\mathrm{rs}'\,F_{y'}\left((\tilde{g}A)^{W'}\right)^{-1}\mathrm{ext}'\,\tilde{g}
=\mathrm{rs}'\,F_{y'}\left(\mathrm{mS}^{W}A^{W}\right)^{-1}\mathrm{mS}^{W}\,\mathrm{ext}\\
&=\mathrm{rs}'\,\mathrm{mS}^{W}F_y(A^{W})^{-1}\,\mathrm{ext}
=\tilde{g}\,\mathrm{rs}\,F_y(A^{W})^{-1}\,\mathrm{ext}
=\tilde{g}\,\mathrm{Sloc}_{\square,y}(A).
\end{align*}

\textbf{(f)}
The functions $-A=-B+d(-\boldsymbol{\lambda})$ and $\overline{A}=\overline{B}+d\boldsymbol{\lambda}$ lie in $\mathrm{Gloc}(\delta)$, and $(-A)^{W}=-A^{W}$, $(\overline{A})^{W}=\overline{A^{W}}$ by [Q.2].
Let $v$ be a spinor on $\mathbb{T}^{W}$ and $f$ a function on $\tilde{W}\times O_{\mathrm{cub}}$.
Computing the pairing over the period cell $\{0<u_\mu<6M\}$ and using [5.2], one finds
\begin{align*}
\langle v,\mathrm{ext}\,f\rangle^{\mathrm{ch}}
&=\frac{\eta_k^3}{48}\sum_{u\in\mathrm{Ch}_0}\sum_{g\in O_{\mathrm{cub}}}v(u,g)\cdot f(\mathrm{Fo}\,u,g)\\
&=\frac{\eta_k^3}{48}\sum_{c\in\tilde{W}}\sum_{g\in O_{\mathrm{cub}}}v(c-o,g)\cdot f(c,g)
=\int_{\tilde{W}}\mathrm{d}x\,(\mathrm{rs}\,v)\cdot f.
\end{align*}
Hence $\mathrm{ext}^{T}=\mathrm{rs}$ and $\mathrm{rs}^{T}=\mathrm{ext}$.
The maps $\mathrm{rs}$ and $\mathrm{ext}$ do not act on the spin index and have real coefficients;
they therefore commute with constant $4\times4$ matrices, and $\overline{\mathrm{rs}}=\mathrm{rs}$, $\overline{\mathrm{ext}}=\mathrm{ext}$.
By Proposition~2.2\,(c) on $\mathbb{T}^{W}$ one has $F_y(A^{W})^{T}=\mathsf{C}^{-1}F_{-y}(-A^{W})\,\mathsf{C}$, and the transpose of an inverse is the inverse of the transpose.
Therefore
\begin{align*}
\mathrm{Sloc}_{\square,y}(A)^{T}&=\mathrm{ext}^{T}\left(F_y(A^{W})^{T}\right)^{-1}\mathrm{rs}^{T}=\mathrm{rs}\,\mathsf{C}^{-1}F_{-y}\left((-A)^{W}\right)^{-1}\mathsf{C}\,\mathrm{ext}\\
&=\mathsf{C}^{-1}\,\mathrm{Sloc}_{\square,-y}(-A)\,\mathsf{C}.
\end{align*}
Complex conjugation of operators satisfies $\overline{O_1O_2}=\overline{O_1}\;\overline{O_2}$ and $\overline{O^{-1}}=\left(\overline{O}\right)^{-1}$.
Proposition~2.2\,(d) on $\mathbb{T}^{W}$, at the connection $A^{W}$, gives $\overline{F_y\left(\overline{A^{W}}\right)}=\varkappa^{-1}F_y(-A^{W})\,\varkappa$.
It follows that
\begin{align*}
\overline{\mathrm{Sloc}_{\square,y}(\overline{A})}&=\mathrm{rs}\left(\overline{F_y\left(\overline{A^{W}}\right)}\right)^{-1}\mathrm{ext}=\mathrm{rs}\,\varkappa^{-1}F_y\left((-A)^{W}\right)^{-1}\varkappa\,\mathrm{ext}\\
&=\varkappa^{-1}\,\mathrm{Sloc}_{\square,y}(-A)\,\varkappa.
\end{align*}

\textbf{(g)}
Replace the lift $z_{\square}$ by $z_{\square}+\ell n$.
Then $\tilde{W}$, $W^{-}$ and $o$ are replaced by $\tilde{W}+\ell n$, $W^{-}+\ell n$ and $o+\ell n$, the map $\mathrm{Fo}$ by $\mathrm{Fo}+\ell n$, and $\mathrm{emb}\,c$ by $\mathrm{emb}(c-\ell n)$.
For an $\ell\mathbb{Z}^3$-periodic $A$ the folded connection formed with the lift $z_{\square}+\ell n$ takes at $(u,u')$ the value
\[
A(\mathrm{Fo}\,u+\ell n,\mathrm{Fo}\,u'+\ell n)=A(\mathrm{Fo}\,u,\mathrm{Fo}\,u')=A^{W}(u,u')
\]
(and the value $0$ where $\mathrm{Fo}\,u=\mathrm{Fo}\,u'$), so that $F^{W}_{y}(A)$ is unchanged.
The maps formed with the new lift are
\begin{align*}
(\mathrm{ext}_nf)(u,g)&=\mathbf{1}[u\in\mathrm{Ch}]\,f(\mathrm{Fo}\,u+\ell n,g)=(\mathrm{ext}\,T_{-\ell n}f)(u,g),\\
(\mathrm{rs}_nv)(c,g)&=v\left(\mathrm{emb}(c-\ell n),g\right)=(T_{\ell n}\,\mathrm{rs}\,v)(c,g).
\end{align*}
Hence the local inverse formed with the lift $z_{\square}+\ell n$ equals $T_{\ell n}\,\mathrm{Sloc}_{\square,y}(A)\,T_{-\ell n}$, the translate of $\mathrm{Sloc}_{\square,y}(A)$.

\emph{Uniformity.}
Besides the function $w_k$, the bounds of (b) contain only the constants $a'$ and $\gamma^{F}$ of the line (C4r).
Proposition~2.4 holds with these constants on every block torus of every level.
Hence the constants $2a'$, $18a'$ and $\gamma^{F}$ in (b) do not depend on $k$, $\ell$, $\mu\in(0,1]$, $e_k>0$, $\delta$, $y$, $\square$, the lift $z_{\square}$, or $A\in\mathrm{Gloc}(\delta)$;
the level $k$ and the parameter $y$ enter the bounds of (b) only through the factor $w_k(y)$.
\qed

\subsection*{2. Proposition 2.6 (operators over $\tilde{W}$ as operators on $T_k$: four compatibilities)}

\begin{proof}[Proof of Proposition~2.6]
\emph{Step 0 (three maps).}
Suppose that $c$ lies in $\tilde{W}$ and in $\tilde{W}+\ell n$ for some $n\in\mathbb{Z}^3$.
Then
\[
\ell\,|n|_\infty=|\ell n|\le|c-z_{\square}|+|c-\ell n-z_{\square}|<3M<\ell,
\]
and therefore $n=0$.
This proves
\begin{equation}
\begin{aligned}
&\text{the sets }\tilde{W}+\ell n,\ n\in\mathbb{Z}^3,\text{ are pairwise disjoint, and}\\
&\text{the covering map }\mathbb{L}_k\to T_k\text{ is injective on }\tilde{W}.
\end{aligned}
\tag*{[6.1]}
\end{equation}
For $f\in\mathfrak{S}^{\mathrm{ch}}(\eta_k,\ell)$ let $\mathrm{res}\,f$ be its restriction to $\tilde{W}\times O_{\mathrm{cub}}$.
For a function $f_0$ on $\tilde{W}\times O_{\mathrm{cub}}$ define the function $\mathrm{ap}\,f_0$ on $\mathbb{L}^{\mathrm{f}}_k$ by
\[
(\mathrm{ap}\,f_0)(x+\ell n):=(-1)^{\mathrm{pa}(n)}f_0(x)\qquad(x\text{ over }\tilde{W},\ n\in\mathbb{Z}^3),
\]
and by $\mathrm{ap}\,f_0:=0$ off the orbit of $\tilde{W}$.
For a function $\mathrm{Gq}$ on the short pairs with both points over $\tilde{W}$ define
\[
(\mathrm{ap2}\,\mathrm{Gq})(x+\ell n,x'+\ell n):=(-1)^{\mathrm{pa}(n)}\,\mathrm{Gq}(x,x'),
\]
and let $\mathrm{ap2}\,\mathrm{Gq}$ vanish on the short pairs that do not lie over a single translate of $\tilde{W}$.
Both definitions are unambiguous by [6.1]: a site lies over at most one translate of $\tilde{W}$, and so does a pair.
The function $\mathrm{ap}\,f_0$ belongs to $\mathfrak{S}^{\mathrm{ch}}(\eta_k,\ell)$, and $\mathrm{ap2}\,\mathrm{Gq}$ is a function of short pairs on $T_k$ (Section~G1.3).
In this notation the definition of Section~G1.6 takes the form
\begin{equation}
\begin{aligned}
&E_{\square}(O)=\mathrm{ap}\,O\,\mathrm{res},\quad\text{resp.}\quad E_{\square}(O)=\mathrm{ap2}\,O\,\mathrm{res}\\
&\quad(\text{the latter for }O\text{ with values in the functions on the short pairs over }\tilde{W}),
\end{aligned}
\tag*{[6.2]}
\end{equation}
and the three maps satisfy
\begin{equation}
\begin{aligned}
&\mathrm{res}\,\mathrm{ap}=1;\qquad|\mathrm{res}\,f|_{\mathrm{su}}\le|f|_{\mathrm{su}};\qquad|\mathrm{ap}\,f_0|_{\mathrm{su}}=|f_0|_{\mathrm{su}},\quad|\mathrm{ap2}\,\mathrm{Gq}|_{\mathrm{su}}=|\mathrm{Gq}|_{\mathrm{su}};\\
&\mathrm{res},\ \mathrm{ap},\ \mathrm{ap2}\ \text{are linear}.
\end{aligned}
\tag*{[6.3]}
\end{equation}

\textbf{(a)} \emph{Linearity and values.}
By [6.2] and [6.3], $E_{\square}(O)$ is linear, with values in $\mathfrak{S}^{\mathrm{ch}}(\eta_k,\ell)$, resp.\ in the functions of short pairs, and $O\mapsto E_{\square}(O)$ is linear.

\emph{Change of the lift.}
Let $n\in\mathbb{Z}^3$, let $\mathrm{res}_n$, $\mathrm{ap}_n$, $\mathrm{ap2}_n$ be the maps formed with $\tilde{W}+\ell n$, and let $O_n:=T_{\ell n}\,O\,T_{-\ell n}$ be the translate of $O$, an operator on the functions on $(\tilde{W}+\ell n)\times O_{\mathrm{cub}}$.
For $f\in\mathfrak{S}^{\mathrm{ch}}(\eta_k,\ell)$ and $x'$ over $\tilde{W}$ one has $(T_{-\ell n}\,\mathrm{res}_nf)(x')=f(x'+\ell n)=(-1)^{\mathrm{pa}(n)}f(x')$;
thus $T_{-\ell n}\,\mathrm{res}_n=(-1)^{\mathrm{pa}(n)}\,\mathrm{res}$.
For $f_0$ on $\tilde{W}\times O_{\mathrm{cub}}$, $x$ over $\tilde{W}$ and $n'\in\mathbb{Z}^3$,
\begin{align*}
(\mathrm{ap}_nT_{\ell n}f_0)(x+\ell n+\ell n')&=(-1)^{\mathrm{pa}(n')}f_0(x)=(-1)^{\mathrm{pa}(n)}(-1)^{\mathrm{pa}(n+n')}f_0(x)\\
&=(-1)^{\mathrm{pa}(n)}(\mathrm{ap}\,f_0)(x+\ell n+\ell n');
\end{align*}
thus $\mathrm{ap}_nT_{\ell n}=(-1)^{\mathrm{pa}(n)}\,\mathrm{ap}$.
For $\mathrm{Gq}$ on the short pairs over $\tilde{W}$, with $(T_{\ell n}\mathrm{Gq})(x+\ell n,x'+\ell n):=\mathrm{Gq}(x,x')$, one has
\begin{align*}
(\mathrm{ap2}_nT_{\ell n}\mathrm{Gq})(x+\ell n+\ell n',x'+\ell n+\ell n')&=(-1)^{\mathrm{pa}(n')}\,\mathrm{Gq}(x,x')\\
&=(-1)^{\mathrm{pa}(n)}(\mathrm{ap2}\,\mathrm{Gq})(x+\ell n+\ell n',x'+\ell n+\ell n'),
\end{align*}
and both sides vanish on the pairs that lie over no translate of $\tilde{W}$;
thus $\mathrm{ap2}_nT_{\ell n}=(-1)^{\mathrm{pa}(n)}\,\mathrm{ap2}$.
Consequently
\[
\mathrm{ap}_n\,O_n\,\mathrm{res}_n=\mathrm{ap}_n\,T_{\ell n}\,O\,T_{-\ell n}\,\mathrm{res}_n=(-1)^{2\mathrm{pa}(n)}\,\mathrm{ap}\,O\,\mathrm{res}=E_{\square}(O),
\]
and the same holds with $\mathrm{ap2}_n$, $\mathrm{ap2}$ in place of $\mathrm{ap}_n$, $\mathrm{ap}$ when $O$ takes values in functions of pairs.
In particular, let $A$ be $\ell\mathbb{Z}^3$-periodic, and let $A|_{\tilde{W}}$ denote the restriction of $A$ to the bonds of $\tilde{W}$.
Then the operator $E_{\square}\left(\mathrm{Sloc}_{\square,y}(A|_{\tilde{W}})\right)$ does not depend on the lift $z_{\square}$, since by Proposition~2.5\,(g) a change of the lift turns the local inverse into its translate.
This operator is the operator $\mathrm{Sloc}_{\square,y}(A)$ on $T_k$ of Section~G1.6.

\emph{Products.}
If $O'$ takes values in the functions on $\tilde{W}\times O_{\mathrm{cub}}$, then [6.3] gives
\[
E_{\square}(O)\,E_{\square}(O')=\mathrm{ap}\,O\,\mathrm{res}\,\mathrm{ap}\,O'\,\mathrm{res}=\mathrm{ap}\,O\,O'\,\mathrm{res}=E_{\square}(OO'),
\]
with $\mathrm{ap2}$ in the first place if $O$ takes values in functions of pairs.

\emph{Norms.}
By [Q.1]\,(a) the set $\tilde{W}$ is a union of blocks of side $L$, hence, the blocks being nested, a union of unit blocks.
Therefore, for a unit site $z$ of $T_k$, either the orbit of $\Delta_z$ is contained in the orbit of $\tilde{W}$, in which case $z$ has exactly one lift $\tilde{z}$ in $\tilde{W}$ by [6.1], or the orbit of $\Delta_z$ is disjoint from the orbit of $\tilde{W}$.
In the first case $\mathbf{1}_{\Delta_z}\mathrm{ap}=\mathrm{ap}\,\mathbf{1}_{\Delta_{\tilde{z}}}$ and $\mathbf{1}_{\Delta_z}\mathrm{ap2}=\mathrm{ap2}\,\mathbf{1}_{\Delta_{\tilde{z}}}$ (the indicator restricts the first point), and $\mathrm{res}\,\mathbf{1}_{\Delta_z}=\mathbf{1}_{\Delta_{\tilde{z}}}\mathrm{res}$.
In the second case $\mathbf{1}_{\Delta_z}\mathrm{ap}=0$ and $\mathbf{1}_{\Delta_z}\mathrm{ap2}=0$ (a pair whose first point does not lie over the orbit of $\tilde{W}$ lies over no translate of $\tilde{W}$), and $\mathrm{res}\,\mathbf{1}_{\Delta_z}=0$.
It follows that
\begin{equation}
\mathbf{1}_{\Delta_z}\,E_{\square}(O)\,\mathbf{1}_{\Delta_{z'}}=
\begin{cases}
E_{\square}\left(\mathbf{1}_{\Delta_{\tilde{z}}}\,O\,\mathbf{1}_{\Delta_{\tilde{z}'}}\right)&\text{if $z$ and $z'$ have lifts $\tilde{z}$, $\tilde{z}'$ in $\tilde{W}$},\\
0&\text{otherwise}.
\end{cases}
\tag*{[6.4]}
\end{equation}
For every operator $O''$ as in Section~G1.6 one has $\|E_{\square}(O'')\|_{\mathrm{su}}=\|O''\|_{\mathrm{su}}$.
Indeed, with $\mathrm{Ap}=\mathrm{ap}$, resp.\ $\mathrm{Ap}=\mathrm{ap2}$, [6.3] gives
\[
|\mathrm{Ap}\,O''\,\mathrm{res}\,f|_{\mathrm{su}}=|O''\,\mathrm{res}\,f|_{\mathrm{su}}\le\|O''\|_{\mathrm{su}}\cdot|\mathrm{res}\,f|_{\mathrm{su}}\le\|O''\|_{\mathrm{su}}\cdot|f|_{\mathrm{su}};
\]
conversely, for $f_0$ on $\tilde{W}\times O_{\mathrm{cub}}$ the spinor $f:=\mathrm{ap}\,f_0$ satisfies $\mathrm{res}\,f=f_0$, $|f|_{\mathrm{su}}=|f_0|_{\mathrm{su}}$ and $|E_{\square}(O'')f|_{\mathrm{su}}=|O''f_0|_{\mathrm{su}}$, again by [6.3].
The distance of $z$ and $z'$ on $T_k$ is $|\tilde{z}-\tilde{z}'|$ by [Q.1]\,(c);
here one uses $|\tilde{z}-\tilde{z}'|\le3M-1\le\ell/2$, the last inequality because $3M<\ell/2$ by [P.3].
Inserting this and [6.4] into [P.6], one obtains $\mathrm{Rw}_z(E_{\square}(O))=\mathrm{Rw}_{\tilde{z}}(O)$ if $z$ has the lift $\tilde{z}$ in $\tilde{W}$, and $\mathrm{Rw}_z(E_{\square}(O))=0$ otherwise.
Hence $N_{\gamma}(E_{\square}(O))=N_{\gamma}(O)$ for every $\gamma\ge0$.

\emph{Level $0$; kernels.}
Let $k=0$, and let $O$ act on the functions on $\tilde{W}\times O_{\mathrm{cub}}$, with the kernel $O(x,x')$ taken with respect to $\int_{\tilde{W}}\mathrm{d}x$ (Section~G1.6).
By [6.1] there is a fundamental domain of $\ell\mathbb{Z}^3$ in $\mathbb{L}_0$ that contains $\tilde{W}$.
Computing with this domain, one has for $x$ over $\tilde{W}$
\[
(E_{\square}(O)f)(x)=\frac{1}{48}\sum_{c'\in\tilde{W}}\sum_{g'\in O_{\mathrm{cub}}}O(x,x')\,f(x').
\]
By the sign rule of Section~G1.2, the kernel (I.1.4) of $E_{\square}(O)$ is therefore given by
\begin{equation}
\begin{aligned}
&E_{\square}(O)(x+\ell n,x'+\ell n')=(-1)^{\mathrm{pa}(n)+\mathrm{pa}(n')}\,O(x,x')\qquad(x,x'\text{ over }\tilde{W};\ n,n'\in\mathbb{Z}^3),\\
&E_{\square}(O)(t,t')=0\quad\text{if $t$ or $t'$ does not lie over the orbit of $\tilde{W}$}.
\end{aligned}
\tag*{[6.5]}
\end{equation}
Consequently, if $z$ and $z'$ have the lifts $\tilde{z}$ and $\tilde{z}'$ in $\tilde{W}$, then
\[
\sup\left\{|E_{\square}(O)(t,t')|_2:\ t\text{ over }z,\ t'\text{ over }z'\right\}=\sup\left\{|O(x,x')|_2:\ x\text{ over }\tilde{z},\ x'\text{ over }\tilde{z}'\right\},
\]
and otherwise the left side vanishes.
Since $d(z,z')=|\tilde{z}-\tilde{z}'|$, it follows that $\mathrm{Rk}_z(E_{\square}(O))=\mathrm{Rk}_{\tilde{z}}(O)$, resp.\ $\mathrm{Rk}_z(E_{\square}(O))=0$, and $N^{\mathrm{ke}}_{\gamma}(E_{\square}(O))=N^{\mathrm{ke}}_{\gamma}(O)$.

\textbf{(b)}
The function $\mathrm{uf}$ is $\ell\mathbb{Z}^3$-periodic and has no zeros;
hence $\mathrm{uf}^{\pm1}f\in\mathfrak{S}^{\mathrm{ch}}(\eta_k,\ell)$ for $f\in\mathfrak{S}^{\mathrm{ch}}(\eta_k,\ell)$.
Writing $\mathrm{uf}$ also for the restriction of $\mathrm{uf}$ to $\tilde{W}$, one has
\[
\mathrm{res}\,\mathrm{uf}^{-1}=\mathrm{uf}^{-1}\,\mathrm{res},\qquad\mathrm{uf}\,\mathrm{ap}=\mathrm{ap}\,\mathrm{uf},\qquad\mathrm{uf}\,\mathrm{ap2}=\mathrm{ap2}\,\mathrm{uf};
\]
for the last identity note that $\mathrm{uf}$ acts on functions of pairs through its value at the first point, and that $\mathrm{uf}(c+\ell n)=\mathrm{uf}(c)$.
Therefore
\[
\mathrm{uf}\,E_{\square}(O)\,\mathrm{uf}^{-1}=\mathrm{uf}\,\mathrm{ap}\,O\,\mathrm{res}\,\mathrm{uf}^{-1}=\mathrm{ap}\,\mathrm{uf}\,O\,\mathrm{uf}^{-1}\,\mathrm{res}=E_{\square}\left(\mathrm{uf}\,O\,\mathrm{uf}^{-1}\right).
\]
A constant invertible $4\times4$ matrix $\mathrm{Mx}$ acts on the values of the functions and commutes with $\mathrm{res}$, $\mathrm{ap}$ and $\mathrm{ap2}$;
hence $\mathrm{Mx}\,E_{\square}(O)\,\mathrm{Mx}^{-1}=E_{\square}\left(\mathrm{Mx}\,O\,\mathrm{Mx}^{-1}\right)$.
Finally, for a bond $(c,c')$ of $\tilde{W}$ the value
\[
(A-d\boldsymbol{\lambda})(c,c')=A(c,c')-\frac{\boldsymbol{\lambda}(c')-\boldsymbol{\lambda}(c)}{\eta_k}
\]
involves $\boldsymbol{\lambda}$ only at the two points $c$, $c'$ of $\tilde{W}$;
it equals $\left(A|_{\tilde{W}}-d(\boldsymbol{\lambda}|_{\tilde{W}})\right)(c,c')$, where $\boldsymbol{\lambda}|_{\tilde{W}}$ is the restriction of $\boldsymbol{\lambda}$ to $\tilde{W}$.

\textbf{(c)}
Let $\tilde{g}=(S,a)$ with a representative $a\in M\mathbb{Z}^3$, and $r:=r(S)$.
By (I.4.2), $\tilde{g}c=rc+a$, $\tilde{g}(c,g)=(rc+a,rg)$ and $\tilde{g}^{-1}(c,g)=\left(r(c-a),rg\right)$.
The set
\[
\tilde{g}\tilde{W}=r\tilde{W}+a=\left\{c:\ \left|c-(rz_{\square}+a)\right|<\tfrac{3M}{2}\right\}
\]
is the set $\tilde{W}$ of the cube $\tilde{g}\square$ for the lift $rz_{\square}+a$ of its centre;
the maps $\mathrm{res}'$, $\mathrm{ap}'$, $\mathrm{ap2}'$ and $E_{\tilde{g}\square}$ are formed with it.
The element $\tilde{g}$ maps functions on $\tilde{W}\times O_{\mathrm{cub}}$ to functions on $\tilde{g}\tilde{W}\times O_{\mathrm{cub}}$, and functions on the short pairs over $\tilde{W}$ to functions on the short pairs over $\tilde{g}\tilde{W}$, by
\[
(\tilde{g}f_0)(x):=Sf_0(\tilde{g}^{-1}x),\qquad(\tilde{g}\,\mathrm{Gq})(x,x'):=S\,\mathrm{Gq}(\tilde{g}^{-1}x,\tilde{g}^{-1}x');
\]
for the second definition one uses that $\tilde{g}$ preserves $d_X$ (I.a1).
For $f\in\mathfrak{S}^{\mathrm{ch}}(\eta_k,\ell)$ and $x$ over $\tilde{g}\tilde{W}$ one has $(\mathrm{res}'\,\tilde{g}f)(x)=Sf(\tilde{g}^{-1}x)$, and $\tilde{g}^{-1}x$ lies over $\tilde{W}$;
hence $\mathrm{res}'\,\tilde{g}=\tilde{g}\,\mathrm{res}$, which is the second assertion of (c).
Next let $f_0$ be a function on $\tilde{W}\times O_{\mathrm{cub}}$, $x$ a site over $\tilde{g}\tilde{W}$, and $n\in\mathbb{Z}^3$.
Then $\tilde{g}^{-1}(x+\ell n)=\tilde{g}^{-1}x+\ell rn$, and $\mathrm{pa}(rn)=\sum_\mu r_\mu n_\mu$ has the parity of $\mathrm{pa}(n)$.
Hence
\[
(\tilde{g}\,\mathrm{ap}\,f_0)(x+\ell n)=S(\mathrm{ap}\,f_0)(\tilde{g}^{-1}x+\ell rn)=(-1)^{\mathrm{pa}(n)}Sf_0(\tilde{g}^{-1}x)=(\mathrm{ap}'\,\tilde{g}f_0)(x+\ell n).
\]
If $t$ does not lie over the orbit of $\tilde{g}\tilde{W}$, then $\tilde{g}^{-1}t$ does not lie over the orbit of $\tilde{W}$, and both sides vanish at $t$.
Thus $\tilde{g}\,\mathrm{ap}=\mathrm{ap}'\,\tilde{g}$.
Similarly, for $\mathrm{Gq}$ on the short pairs over $\tilde{W}$ and a short pair $(x,x')$ with both points over $\tilde{g}\tilde{W}$,
\begin{align*}
(\tilde{g}\,\mathrm{ap2}\,\mathrm{Gq})(x+\ell n,x'+\ell n)&=S(\mathrm{ap2}\,\mathrm{Gq})(\tilde{g}^{-1}x+\ell rn,\tilde{g}^{-1}x'+\ell rn)\\
&=(-1)^{\mathrm{pa}(n)}S\,\mathrm{Gq}(\tilde{g}^{-1}x,\tilde{g}^{-1}x')=(\mathrm{ap2}'\,\tilde{g}\,\mathrm{Gq})(x+\ell n,x'+\ell n).
\end{align*}
A short pair lies over a translate of $\tilde{g}\tilde{W}$ if and only if its image under $\tilde{g}^{-1}$ lies over a translate of $\tilde{W}$;
so both sides vanish on the remaining short pairs, and $\tilde{g}\,\mathrm{ap2}=\mathrm{ap2}'\,\tilde{g}$.
From $\mathrm{res}'\,\tilde{g}=\tilde{g}\,\mathrm{res}$ one gets $\mathrm{res}\,\tilde{g}^{-1}=\tilde{g}^{-1}\,\mathrm{res}'$, and therefore
\[
\tilde{g}\,E_{\square}(O)\,\tilde{g}^{-1}=\tilde{g}\,\mathrm{ap}\,O\,\mathrm{res}\,\tilde{g}^{-1}=\mathrm{ap}'\left(\tilde{g}\,O\,\tilde{g}^{-1}\right)\mathrm{res}'=E_{\tilde{g}\square}\left(\tilde{g}\,O\,\tilde{g}^{-1}\right),
\]
with $\mathrm{ap2}$, $\mathrm{ap2}'$ in place of $\mathrm{ap}$, $\mathrm{ap}'$ if $O$ takes values in functions of pairs.
(If $a$ is replaced by $a+\ell n$, the action of $\tilde{g}$ on $\mathfrak{S}^{\mathrm{ch}}(\eta_k,\ell)$ is multiplied by $(-1)^{\mathrm{pa}(n)}$, which leaves the left side unchanged;
on the right side the lift and the operator are translated by $\ell n$, which leaves $E_{\tilde{g}\square}(\cdot)$ unchanged by (a).)
Finally, if $(c,c')$ is a bond of $\tilde{g}\tilde{W}$, then $(\tilde{g}^{-1}c,\tilde{g}^{-1}c')$ is a bond of $\tilde{W}$ and $(\tilde{g}A)(c,c')=A(\tilde{g}^{-1}c,\tilde{g}^{-1}c')$;
thus the restriction of $\tilde{g}A$ to the bonds of $\tilde{g}\tilde{W}$ is $\tilde{g}\left(A|_{\tilde{W}}\right)$.

\textbf{(d)}
Let $O$ take values in the functions on $\tilde{W}\times O_{\mathrm{cub}}$.
For $f'\in\mathfrak{S}^{\mathrm{ch}}(\eta_k,\ell)$ and $f_0$ on $\tilde{W}\times O_{\mathrm{cub}}$ the product $f'\cdot(\mathrm{ap}\,f_0)$ is $\ell\mathbb{Z}^3$-periodic and vanishes off the orbit of $\tilde{W}$.
Evaluating the pairing on a fundamental domain that contains $\tilde{W}$, which exists by [6.1] (by Section~I.1.6 the choice of the domain is immaterial), one obtains
\[
\langle f',\mathrm{ap}\,f_0\rangle^{\mathrm{ch}}=\frac{\eta_k^3}{48}\sum_{c\in\tilde{W}}\sum_{g\in O_{\mathrm{cub}}}f'(c,g)\cdot f_0(c,g)=\int_{\tilde{W}}\mathrm{d}x\,(\mathrm{res}\,f')\cdot f_0.
\]
Thus $\mathrm{ap}^{T}=\mathrm{res}$ and $\mathrm{res}^{T}=\mathrm{ap}$ with respect to the pairing $\langle\cdot,\cdot\rangle^{\mathrm{ch}}$ on $T_k$ and the pairing $\int_{\tilde{W}}\mathrm{d}x\,f'\cdot f$, and
\[
E_{\square}(O)^{T}=\mathrm{res}^{T}\,O^{T}\,\mathrm{ap}^{T}=\mathrm{ap}\,O^{T}\,\mathrm{res}=E_{\square}(O^{T}).
\]

\textbf{(e)}
The maps $\mathrm{res}$, $\mathrm{ap}$ and $\mathrm{ap2}$ have the coefficients $0$, $1$, $-1$;
hence $\overline{\mathrm{res}}=\mathrm{res}$, $\overline{\mathrm{ap}}=\mathrm{ap}$, $\overline{\mathrm{ap2}}=\mathrm{ap2}$, and
\[
\overline{E_{\square}(O)}=\left(\overline{\mathrm{ap}}\right)\left(\overline{O}\right)\left(\overline{\mathrm{res}}\right)=E_{\square}\left(\overline{O}\right).
\]
\end{proof}

\begin{named}{Corollary 6.A}[the relations of Proposition~2.5\,(d)--(f) on $T_k$]
Let the data (SF) be given, let $A\in\tilde{G}_k$ and $y\in Y_{k}$, and let $\mathrm{Sloc}_{\square,y}(A)$ be the operator on $T_k$ of Section~G1.6.
For every real function $\boldsymbol{\lambda}$ of $c$ on $T_k$, every $\tilde{g}=(S,a)\in\mathbb{G}_k$, and $y'$ as in Proposition~2.5\,(e),
\begin{gather*}
\mathrm{Sloc}_{\square,y}(A-d\boldsymbol{\lambda})=\exp\left(ie_k\boldsymbol{\lambda}\right)\mathrm{Sloc}_{\square,y}(A)\exp\left(-ie_k\boldsymbol{\lambda}\right),\\
\mathrm{Sloc}_{\tilde{g}\square,y'}(\tilde{g}A)\,\tilde{g}=\tilde{g}\,\mathrm{Sloc}_{\square,y}(A),\\
\mathrm{Sloc}_{\square,y}(A)^{T}=\mathsf{C}^{-1}\,\mathrm{Sloc}_{\square,-y}(-A)\,\mathsf{C},\\
\overline{\mathrm{Sloc}_{\square,y}(\overline{A})}=\varkappa^{-1}\,\mathrm{Sloc}_{\square,y}(-A)\,\varkappa.
\end{gather*}
\end{named}

\begin{proof}
By [Q.3]\,(G2) the connections $A-d\boldsymbol{\lambda}$, $\tilde{g}A$, $-A$ and $\overline{A}$ lie in $\tilde{G}_k$;
by [Q.3]\,(G1) their restrictions lie in $\mathrm{Gloc}(\delta_k)$, and by [P.2] the number $\delta_k$ satisfies the conditions imposed on $\delta$ in Proposition~2.5.
For the first relation, apply $E_{\square}$ to Proposition~2.5\,(d) and use Proposition~2.6\,(b) with $\mathrm{uf}=\exp\left(ie_k\boldsymbol{\lambda}\right)$.
For the second, write Proposition~2.5\,(e) in the form $\mathrm{Sloc}_{\tilde{g}\square,y'}(\tilde{g}A)=\tilde{g}\,\mathrm{Sloc}_{\square,y}(A)\,\tilde{g}^{-1}$ over $\tilde{g}\tilde{W}$, apply $E_{\tilde{g}\square}$, and use Proposition~2.6\,(c).
For the last two, apply $E_{\square}$ to Proposition~2.5\,(f) and use Proposition~2.6\,(d), (e), together with Proposition~2.6\,(b) for $\mathrm{Mx}=\mathsf{C}^{-1}$ and $\mathrm{Mx}=\varkappa^{-1}$.
\end{proof}

\subsection*{3. Proposition 2.7 (parametrix; sums over walks)}

\textit{Proof of Proposition~2.7.}
The conditions (L5), (M1), (M2), (M5), (Z2.s1) and (Z2.s2) on the constants enter the proof through [P.1]--[P.3].
Throughout the proof $N:=N_{\gamma^{F}}$, that is, $N$ is the localised norm of Section~G1.3 on $T_k$ with $\gamma=\gamma^{F}$;
accordingly, the row sums $\mathrm{Rw}_z$ of [P.6], and the inequalities (Rw1)--(Rw3) for them, are taken with $\gamma=\gamma^{F}$.

\emph{Step 0 (setting).}
The letter $A$ also denotes the $\ell\mathbb{Z}^3$-periodic bond function on $\mathbb{L}_k$ given by the connection.
By (F3), for $f\in\mathfrak{S}^{\mathrm{ch}}(\eta_k,\ell)$ the spinor $Ff$ is the function $F_y(A)f$ on $\mathbb{L}^{\mathrm{f}}_k$ of Section~0.2;
by (F2), $F_y(A)$ commutes with the translations $T_{\ell n}$, since $\ell n\in L\mathbb{Z}^3$ and $T_{\ell n}A=A$.
The restriction $A|_{\tilde{W}}$ of $A$ to the bonds of $\tilde{W}$ lies in $\mathrm{Gloc}(\delta_k)$ by [Q.3], and $\delta_k$ is a number $\delta$ as in Proposition~2.4 by [P.2].
Hence Proposition~2.5 applies to $A|_{\tilde{W}}$, and in Proposition~2.5\,(c) one may take the continuation $\tilde{A}:=A$.
In the notation of [6.2], $\mathrm{Sloc}_{\square}=\mathrm{ap}\,\mathrm{Sloc}_{\square,y}(A|_{\tilde{W}})\,\mathrm{res}$.
Proposition~2.6\,(a) and Proposition~2.5\,(b) give
\begin{equation}
\begin{aligned}
&N(\mathrm{Sloc}_{\square})\le2a'\cdot w_k(y);\qquad N\left(E_{\square}\left(\delta_{\alpha,A}\mathrm{Sloc}_{\square,y}(A|_{\tilde{W}})\right)\right)\le18a';\\
&\text{at }k=0:\quad N^{\mathrm{ke}}_{\gamma^{F}}(\mathrm{Sloc}_{\square})\le2a'\cdot w_0(y).
\end{aligned}
\tag*{[7.1]}
\end{equation}
By [Q.3], $\mathfrak{i}(A)<e_k^{\varepsilon}$; together with (Z2.s1) this gives
\begin{equation}
3L\cdot\mathfrak{i}(A)\le3L\cdot e_0^{\varepsilon}\le\frac12,\qquad\exp\left(3\,\mathfrak{i}(A)\right)\le\mathrm{e}.
\tag*{[7.2]}
\end{equation}
Moreover, $T_k$ is the block torus $(0,\ell,-1)$, since $\ell\ge3L$ is a multiple of $L$;
$\gamma^{F}L\le\log2$ by [P.1];
$3L\cdot\mathfrak{i}(A)\le1$ by [7.2];
and $h_{\square}$ is a real scalar on $T_k$ by [Q.5], with $|h_{\square}(c)-h_{\square}(c')|\le(\pi/M)\,|c-c'|_1$ by the third relation of (6.1) in Section~G1.7.
Therefore the hypotheses of Proposition~2.1\,(b) are fulfilled on $T_k$ with $\gamma=\gamma^{F}$, $\psi=h_{\square}$ and $\mathrm{Lip}=\pi/M$, and those of Proposition~2.1\,(c) with $\gamma=\gamma^{F}$.
On the lift $\tilde{W}$, [Q.5] and [P.3] give
\[
\operatorname{supp}h_{\square}\cap\tilde{W}\subset\left\{c:\ |c-z_{\square}|<\tfrac{3M}{4}\right\}\subset W^{-}.
\]

\textbf{(c)}
One has $\mathrm{Cm}_{\square}=-[F,h_{\square}]=h_{\square}F-Fh_{\square}$.
By the product rule for kernels (Section~G1.2),
\[
\mathrm{Cm}_{\square}(x,x')=h_{\square}(c)\cdot F(x,x')-F(x,x')\cdot h_{\square}(c')=\left(h_{\square}(c)-h_{\square}(c')\right)F(x,x').
\]
To write $\mathrm{Cm}_{\square}$ out, insert [P.4] and [P.5].
The terms of $F$ that stay over $c$, namely the three internal hops $i=0,1,2$, the diagonal terms, $\mu$ and, at $k=0$, $iy\gamma_5$, commute with $h_{\square}$ and cancel;
the hop of type $3$ and the penalty remain.
For $x=(c,g)$, and with $z$, $z'$ as in Section~0.2, this gives
\begin{equation}
\begin{aligned}
(k\ge1)\quad(\mathrm{Cm}_{\square}f)(c,g)={}&\eta_k^{-1}M_{3}(x)\,u_A(c,c_3)\left(h_{\square}(c_3)-h_{\square}(c)\right)f(c_3,gs_3)\\
&+\boldsymbol{\alpha}_{y}\,\varphi_A(z,c)^{-1}\,\frac{\eta_k^3}{48}\sum_{c'\in\Delta_z}\sum_{g'\in O_{\mathrm{cub}}}\varphi_A(z,c')\left(h_{\square}(c)-h_{\square}(c')\right)f(c',g')\\
&+\boldsymbol{\beta}_{y}\,\varphi^{+}_A(z',c)^{-1}L^{-3}\,\frac{\eta_k^3}{48}\\
&\qquad\cdot\sum_{c'\in\mathrm{Bl}_{k+1}(z')}\sum_{g'\in O_{\mathrm{cub}}}\varphi^{+}_A(z',c')\left(h_{\square}(c)-h_{\square}(c')\right)f(c',g');\\
(k=0)\quad(\mathrm{Cm}_{\square}f)(c,g)={}&M_{3}(x)\,u_A(c,c_3)\left(h_{\square}(c_3)-h_{\square}(c)\right)f(c_3,gs_3)\\
&+b_{L}\,U_A(z',c)^{-1}L^{-3}\,\frac{1}{48}\sum_{c'\in B(z')}\sum_{g'\in O_{\mathrm{cub}}}U_A(z',c')\left(h_{\square}(c)-h_{\square}(c')\right)f(c',g').
\end{aligned}
\tag*{[7.3]}
\end{equation}
Next, if $\mathrm{Cm}_{\square}(x,x')\neq0$, then $c$ and $c'$ have representatives modulo $\ell\mathbb{Z}^3$ at distance less than $3M/4+L$ from $z_{\square}$, and these representatives lie in $W^{-}$.
To prove this, let $\mathrm{Cm}_{\square}(x,x')\neq0$.
By [7.3] (this is also the clause on the kernel in Proposition~2.1\,(b)), $h_{\square}(c)\neq h_{\square}(c')$, and there is $n''\in\mathbb{Z}^3$ such that either $x'+\ell n''=x^{(3)}$, or $c$ and $c'+\ell n''$ lie in one block of side $L$.
In the first case $|c-(c'+\ell n'')|=\eta_k<L$;
in the second $|c-(c'+\ell n'')|\le L-\eta_k<L$.
At least one of the numbers $h_{\square}(c)$, $h_{\square}(c')$ is different from $0$.
If $h_{\square}(c)\neq0$, then by [Q.5] there is $n$ with $|c+\ell n-z_{\square}|<3M/4$, and then $|c'+\ell n''+\ell n-z_{\square}|<3M/4+L$.
If $h_{\square}(c')\neq0$, there is $n'$ with $|c'+\ell n'-z_{\square}|<3M/4$, and then $|c+\ell(n'-n'')-z_{\square}|<3M/4+L$.
In either case $c$ and $c'$ have representatives at distance less than $3M/4+L$ from $z_{\square}$.
Since $3M/4+L<3M/2-\eta_k$ by [P.3], these representatives lie in $W^{-}$.
In particular, the rows of $\mathrm{Cm}_{\square}$ at the sites $x$ that do not lie over the orbit of $\left\{c:\ |c-z_{\square}|<3M/4+L\right\}$ vanish for every $A$.
For $x$ over $W^{-}$ one has $(\mathrm{Cm}_{\square}f)(x)=h_{\square}(c)\cdot(Ff)(x)-\left(F(h_{\square}f)\right)(x)$, and the row of $F=F_y(A)$ at $x$ involves $A$ only on the bonds of $\tilde{W}$ (first assertion of Proposition~2.5\,(c)).
As $\mathrm{Cm}_{\square}f$ is antiperiodic, the rows over $W^{-}+\ell n$ coincide, up to the sign $(-1)^{\mathrm{pa}(n)}$, with the rows over $W^{-}$.
Consequently $\mathrm{Cm}_{\square}(A)=\mathrm{Cm}_{\square}(A')$ whenever two connections $A,A'\in\tilde{G}_k$ coincide on the bonds of $\tilde{W}$;
that is, $\mathrm{Cm}_{\square}$ reads $A$ only on the bonds of $\tilde{W}$.
Let now $\square'$ be a second $M$-cube, and suppose that $h_{\square'}(c)\cdot\mathrm{Cm}_{\square}(x,x')\neq0$ or $\mathrm{Cm}_{\square}(x',x)\cdot h_{\square'}(c)\neq0$.
Then $d(c,z_{\square'})<3M/4$ by [Q.5], and $d(c,z_{\square})<3M/4+L$ by what has just been shown;
hence $d(z_{\square},z_{\square'})<3M/2+L<2M$ by [P.3].
As $z_{\square}-z_{\square'}\in M\mathbb{Z}^3$ and $\ell\in M\mathbb{Z}$, the distance $d(z_{\square},z_{\square'})$ is a multiple of $M$, hence at most $M$.
It remains to note that the kernels of $h_{\square'}\mathrm{Cm}_{\square}$ and $\mathrm{Cm}_{\square}h_{\square'}$ are $h_{\square'}(c)\cdot\mathrm{Cm}_{\square}(x,x')$ and $\mathrm{Cm}_{\square}(x,x')\cdot h_{\square'}(c')$.

\textbf{(a)}
Let $f\in\mathfrak{S}^{\mathrm{ch}}(\eta_k,\ell)$, and put $\mathrm{vl}:=\mathrm{Sloc}_{\square,y}(A|_{\tilde{W}})\,\mathrm{res}(h_{\square}f)$ and $v:=\mathrm{Sloc}_{\square}h_{\square}f=\mathrm{ap}\,\mathrm{vl}$.
Then
\[
F(h_{\square}v)=h_{\square}Fv+[F,h_{\square}]v=h_{\square}Fv-\mathrm{Cm}_{\square}v.
\]
Let $x=(c,g)$ lie over $\tilde{W}$.
If $h_{\square}(c)=0$, then $(h_{\square}Fv)(x)=0=h_{\square}(c)^2f(x)$.
If $h_{\square}(c)\neq0$, then $c\in W^{-}$;
the row of $F_y(A)$ at $x$ involves the function only over $\tilde{W}$ (first assertion of Proposition~2.5\,(c)), the functions $v$ and $\mathrm{cz}\,\mathrm{vl}$ coincide over $\tilde{W}$, and the second assertion of Proposition~2.5\,(c) yields
\[
(Fv)(x)=(F_y(A)\,\mathrm{cz}\,\mathrm{vl})(x)=\left(\mathrm{res}(h_{\square}f)\right)(x)=h_{\square}(c)\cdot f(x).
\]
Thus $h_{\square}Fv=h_{\square}^2f$ over $\tilde{W}$.
Both functions belong to $\mathfrak{S}^{\mathrm{ch}}(\eta_k,\ell)$ and vanish off the orbit of $\tilde{W}$, where $h_{\square}$ vanishes;
hence they are equal, and
\[
F\,h_{\square}\mathrm{Sloc}_{\square}h_{\square}f=h_{\square}^2f-\mathrm{Cm}_{\square}\mathrm{Sloc}_{\square}h_{\square}f.
\]

\textbf{(b)}
Let $f\in\mathfrak{S}^{\mathrm{ch}}(\eta_k,\ell)$.
Since $h_{\square}Ff=F(h_{\square}f)+\mathrm{Cm}_{\square}f$, one has
\[
h_{\square}\mathrm{Sloc}_{\square}h_{\square}Ff=h_{\square}\mathrm{Sloc}_{\square}F(h_{\square}f)+h_{\square}\mathrm{Sloc}_{\square}\mathrm{Cm}_{\square}f,
\]
and it suffices to prove that $\mathrm{Sloc}_{\square}F(h_{\square}f)=h_{\square}f$.
Put $\mathrm{fs}:=\mathrm{cz}\,\mathrm{res}(h_{\square}f)$;
by Step~0 this is a function on $\mathbb{L}^{\mathrm{f}}_k$ that vanishes off $W^{-}\times O_{\mathrm{cub}}$.
The spinor $h_{\square}f$ is antiperiodic and vanishes off the orbit of $\tilde{W}$, so that
\[
h_{\square}f=\mathrm{ap}\,\mathrm{res}(h_{\square}f)=\sum_{n\in\mathbb{Z}^3}(-1)^{\mathrm{pa}(n)}\,T_{\ell n}\,\mathrm{fs},
\]
a locally finite sum.
The rows of $F_y(A)$ have finite support, and $F_y(A)$ commutes with $T_{\ell n}$;
hence
\[
F(h_{\square}f)=\sum_{n\in\mathbb{Z}^3}(-1)^{\mathrm{pa}(n)}\,T_{\ell n}\,F_y(A)\,\mathrm{fs}.
\]
By the third assertion of Proposition~2.5\,(c), $F_y(A)\,\mathrm{fs}$ vanishes off $\tilde{W}\times O_{\mathrm{cub}}$;
so $T_{\ell n}F_y(A)\,\mathrm{fs}$ vanishes off $(\tilde{W}+\ell n)\times O_{\mathrm{cub}}$, which for $n\neq0$ is disjoint from $\tilde{W}\times O_{\mathrm{cub}}$ by [6.1].
Consequently $\mathrm{res}\,F(h_{\square}f)=\mathrm{res}\,F_y(A)\,\mathrm{fs}$, and the third assertion of Proposition~2.5\,(c) gives
\[
\mathrm{res}\,\mathrm{Sloc}_{\square}F(h_{\square}f)=\mathrm{Sloc}_{\square,y}(A|_{\tilde{W}})\,\mathrm{res}\,F_y(A)\,\mathrm{fs}=\mathrm{res}\,\mathrm{fs}=\mathrm{res}(h_{\square}f).
\]
Both $\mathrm{Sloc}_{\square}F(h_{\square}f)$ and $h_{\square}f$ are of the form $\mathrm{ap}(\cdot)$, and $\mathrm{res}\,\mathrm{ap}=1$;
hence they are equal.

\textbf{(d)} \emph{The bound for $N(\mathrm{Cm}_{\square})$.}
Proposition~2.1\,(b) with $\psi=h_{\square}$ and $\mathrm{Lip}=\pi/M$ (Step~0) gives
\[
N(\mathrm{Cm}_{\square})=N\left([F,h_{\square}]\right)\le\frac{\pi\cdot C'_V}{M}.
\]

\emph{Products.}
By (Rw3), $N(h_{\square})\le1$.
With (Rw2), (Rw3), [7.1] and $w_k(y)\le1$ one obtains
\[
N\left(h_{\square}\mathrm{Sloc}_{\square}h_{\square}\right)\le N(\mathrm{Sloc}_{\square})\le2a'\cdot w_k(y)
\]
and
\begin{multline*}
N\left(\mathrm{Cm}_{\square}\mathrm{Sloc}_{\square}h_{\square}\right),\ N\left(h_{\square}\mathrm{Sloc}_{\square}\mathrm{Cm}_{\square}\right),\ N\left(\mathrm{Cm}_{\square}h_{\square'}\mathrm{Sloc}_{\square'}\right)\\
\le\frac{\pi\cdot C'_V}{M}\cdot2a'\cdot w_k(y)\le\frac{2\pi\cdot C'_V\cdot a'}{M}.
\end{multline*}

\emph{The bound for $N\left(\delta_{\alpha,A}h_{\square}\mathrm{Sloc}_{\square}\right)$.}
Let $k\ge1$;
at $k=0$ there is no short pair, and the operator is $0$.
The operator in question is the composition on $T_k$ (Section~G1.6) that sends $f$ to the restriction of $\delta_{\alpha,A}(h_{\square}\mathrm{Sloc}_{\square}f)$ to the short pairs.
It vanishes on every short pair $(x,x')$ that does not lie over a single translate of $\tilde{W}$.
Indeed, if $h_{\square}(c)\neq0$ or $h_{\square}(c')\neq0$, one of the two points has a representative at distance less than $3M/4$ from $z_{\square}$, the other point, translated by the same vector, lies at distance less than $3M/4+1<3M/2$, and so the pair lies over a translate of $\tilde{W}$.
Hence the composition equals $E_{\square}$ applied to the operator $\delta_{\alpha,A}h_{\square}\mathrm{Sloc}_{\square,y}(A|_{\tilde{W}})$ over $\tilde{W}$.
Now let $f\in\mathfrak{S}^{\mathrm{ch}}(\eta_k,\ell)$ and $v:=\mathrm{Sloc}_{\square}f$.
By Proposition~2.1\,(c), for every short pair $(x,x')$,
\begin{equation}
\begin{aligned}
\left(\delta_{\alpha,A}(h_{\square}v)\right)(x,x')&=h_{\square}(c)\cdot\left(\delta_{\alpha,A}v\right)(x,x')+\mathrm{wq}(x,x')\,v(x'),\\
\mathrm{wq}(x,x')&:=d_X(x,x')^{-\alpha}\left(h_{\square}(c')-h_{\square}(c)\right)U_A(x,x').
\end{aligned}
\tag*{[7.4]}
\end{equation}

\emph{The first term of [7.4].}
Let $h_{\square}(c)\neq0$.
After a translation of the pair by a vector of $\ell\mathbb{Z}^3$, which multiplies both sides of the identity below by the same sign, one has $|c-z_{\square}|<3M/4$;
as $|c-c'|\le d_X(x,x')<1$, also $|c'-z_{\square}|<3M/4+1<3M/2$, and both points lie over $\tilde{W}$.
There $v=\mathrm{Sloc}_{\square,y}(A|_{\tilde{W}})\,\mathrm{res}\,f$, the six paths of the pair lie in the lattice box $\tilde{W}$, and
\begin{align*}
\left(\delta_{\alpha,A}v\right)(x,x')&=\left(\delta_{\alpha,A}\mathrm{Sloc}_{\square,y}(A|_{\tilde{W}})\,\mathrm{res}\,f\right)(x,x')\\
&=\left(E_{\square}\left(\delta_{\alpha,A}\mathrm{Sloc}_{\square,y}(A|_{\tilde{W}})\right)f\right)(x,x').
\end{align*}
So the first term of [7.4] is $h_{\square}\,E_{\square}\left(\delta_{\alpha,A}\mathrm{Sloc}_{\square,y}(A|_{\tilde{W}})\right)f$, that is, the operator denoted $\delta_{\alpha,A}\mathrm{Sloc}_{\square}$ in Section~G1.6, multiplied by $h_{\square}$ at the first point.
By (Rw3) and [7.1] its norm $N$ is at most $18a'$.

\emph{The second term of [7.4].}
In the notation of Proposition~2.1\,(c) it is $E_{\mathrm{wq}}\mathrm{Sloc}_{\square}f$.
The function $\mathrm{wq}$ is complex valued on the short pairs and satisfies $\mathrm{wq}(x+\ell e_\mu,x'+\ell e_\mu)=\mathrm{wq}(x,x')$:
the distance $d_X$ does not change, the six paths are translated together with the pair, and $h_{\square}$ and $A$ are $\ell\mathbb{Z}^3$-periodic.
For a pair inside a cluster one has $c'=c$ and $\mathrm{wq}=0$.
For $c'\neq c$ one has $d_X(x,x')=|c-c'|$, because $|c-c'|\ge\eta_k$, and $|c-c'|_1\le3|c-c'|$;
by the third relation of (6.1) in Section~G1.7, by [Q.4] and by [7.2], and since $|c-c'|<1$ and $\alpha<1$,
\begin{align*}
|\mathrm{wq}(x,x')|&\le|c-c'|^{-\alpha}\cdot\frac{\pi}{M}\cdot3|c-c'|\cdot\exp\left(3|c-c'|\cdot\mathfrak{i}(A)\right)\\
&\le\frac{3\pi}{M}\cdot|c-c'|^{1-\alpha}\cdot\exp\left(3\,\mathfrak{i}(A)\right)\le\frac{3\pi\cdot\mathrm{e}}{M}.
\end{align*}
Proposition~2.1\,(c) and [7.1] now give
\[
N\left(E_{\mathrm{wq}}\mathrm{Sloc}_{\square}\right)\le27\exp(\gamma^{F})\cdot\frac{3\pi\cdot\mathrm{e}}{M}\cdot2a'=\frac{162\pi\cdot\mathrm{e}\cdot\exp(\gamma^{F})\cdot a'}{M}.
\]
Since $162\pi\cdot\mathrm{e}\cdot\exp(\gamma^{F})<162\cdot3.1416\cdot2.7183\cdot1.04<1439<1600\le M$ by [P.1], it follows that $N\left(E_{\mathrm{wq}}\mathrm{Sloc}_{\square}\right)\le a'$.
By (Rw1),
\[
N\left(\delta_{\alpha,A}h_{\square}\mathrm{Sloc}_{\square}\right)\le18a'+a'=19a'.
\]

\textbf{(e)}
Let $n\ge0$, let $\mathrm{wk}=(\square_0,\dots,\square_n)\in\mathrm{Wk}_n$, and let $z$ be a unit site of $T_k$.
Write
\begin{align*}
S_{y,\mathrm{wk}}(A)&=O_0O_1\cdots O_n,\\
&\qquad O_0:=h_0\,\mathrm{Sloc}_0\,h_0,\quad O_i:=\mathrm{Cm}_i\,\mathrm{Sloc}_i\,h_i\quad(1\le i\le n);\\
S'_{y,\mathrm{wk}}(A)&=O'_nO'_{n-1}\cdots O'_0\,h_0\quad(n\ge1),\\
&\qquad O'_n:=h_n\,\mathrm{Sloc}_n,\quad O'_{i-1}:=-\mathrm{Cm}_i\,h_{i-1}\,\mathrm{Sloc}_{i-1}\quad(1\le i\le n);\\
S'_{y,\mathrm{wk}}(A)&=h_0\,\mathrm{Sloc}_0\,h_0=S_{y,\mathrm{wk}}(A)\quad(n=0).
\end{align*}
By (d), $N(O_0)\le2a'\cdot w_k(y)$ and $N(O_i),\,N(O'_{i-1})\le2\pi\cdot C'_V\cdot a'/M$ for $1\le i\le n$.
By (Rw3) and [7.1], $N(O'_n)\le N(\mathrm{Sloc}_n)\le2a'\cdot w_k(y)$, and $N(h_0)\le1$.
Applying (Rw2) $n$ times ($n+1$ times for $S'$) and (Rw3), one obtains
\begin{equation}
\begin{aligned}
&\mathrm{Rw}_z\left(S_{y,\mathrm{wk}}(A)\right)\le\mathrm{Rw}_z(O_0)\prod_{i=1}^{n}N(O_i)\le2a'\cdot w_k(y)\cdot\left(\frac{2\pi\cdot C'_V\cdot a'}{M}\right)^{n},\\
&\qquad\text{and }\mathrm{Rw}_z\left(S_{y,\mathrm{wk}}(A)\right)=0\text{ unless }\Delta_z\text{ meets }\operatorname{supp}h_0;\\
&\mathrm{Rw}_z\left(S'_{y,\mathrm{wk}}(A)\right)\le\mathrm{Rw}_z(O'_n)\prod_{i=1}^{n}N(O'_{i-1})\cdot N(h_0)\le2a'\cdot w_k(y)\cdot\left(\frac{2\pi\cdot C'_V\cdot a'}{M}\right)^{n},\\
&\qquad\text{and }\mathrm{Rw}_z\left(S'_{y,\mathrm{wk}}(A)\right)=0\text{ unless }\Delta_z\text{ meets }\operatorname{supp}h_n.
\end{aligned}
\tag*{[7.5]}
\end{equation}

\emph{Counting.}
By [Q.5], the support $\operatorname{supp}h_{\square}$ meets $\Delta_z$ for at most $8$ cubes $\square$.
For a fixed cube $\square$, at most $27$ cubes $\square'$ satisfy $d(z_{\square},z_{\square'})\le M$, since every coordinate of $z_{\square'}-z_{\square}$ is congruent to $-M$, $0$ or $M$ modulo $\ell$.
Choosing $\square_0$ first and then $\square_1,\dots,\square_n$ in turn, one sees that at most $8\cdot 27^n$ walks of $\mathrm{Wk}_n$ have a cube $\square_0$ for which $\operatorname{supp}h_0$ meets $\Delta_z$.
Choosing $\square_n$ first and then $\square_{n-1},\dots,\square_0$, one sees in the same way that at most $8\cdot 27^n$ walks of $\mathrm{Wk}_n$ have a cube $\square_n$ for which $\operatorname{supp}h_n$ meets $\Delta_z$.
Every sum below is a finite sum of non-negative terms.
With [7.5] it follows that
\begin{align*}
&\sum_{z'}\exp\left(\gamma^{F}d(z,z')\right)\sum_{\mathrm{wk}\in\mathrm{Wk}_n}\mathrm{wp}^{n+1}\left\|\mathbf{1}_{\Delta_z}S^{\#}_{\mathrm{wk}}\mathbf{1}_{\Delta_{z'}}\right\|_{\mathrm{su}}
=\sum_{\mathrm{wk}\in\mathrm{Wk}_n}\mathrm{wp}^{n+1}\cdot\mathrm{Rw}_z\left(S^{\#}_{\mathrm{wk}}\right)\\
&\qquad\le 8\cdot 27^n\cdot\mathrm{wp}^{n+1}\cdot 2a'\cdot w_{k}(y)\cdot\left(\frac{2\pi C'_{V}a'}{M}\right)^{n}
=16\,a'\cdot\mathrm{wp}\cdot w_{k}(y)\cdot\mathrm{qq}(\mathrm{wp})^{n}.
\end{align*}
Taking the supremum over $z$ gives the first bound of (e).

\emph{The second bound} ($k\ge1$; at $k=0$ the left-hand side vanishes).
By Section~G1.6, $\delta_{\alpha,A}S^{\#}_{\mathrm{wk}}$ denotes the composition on $T_k$, restricted to the short pairs; thus
\begin{equation*}
\delta_{\alpha,A}S_{y,\mathrm{wk}}(A)=\left(\delta_{\alpha,A}h_0\,\mathrm{Sloc}_0\right)h_0\,O_1\cdots O_n,
\qquad
\delta_{\alpha,A}S'_{y,\mathrm{wk}}(A)=\left(\delta_{\alpha,A}h_n\,\mathrm{Sloc}_n\right)O'_{n-1}\cdots O'_0\,h_0 .
\end{equation*}
By (d), the first factor satisfies $N\left(\delta_{\alpha,A}h_{\square}\,\mathrm{Sloc}_{\square}\right)\le19a'$, with $\square=\square_0$, resp.\ $\square=\square_n$; together with (Rw2) and (Rw3) this gives
\begin{equation*}
\mathrm{Rw}_z\left(\delta_{\alpha,A}S^{\#}_{\mathrm{wk}}\right)\le 19\,a'\cdot\left(\frac{2\pi C'_{V}a'}{M}\right)^{n}.
\end{equation*}
Let $\square$ be the cube at the left end of the product, that is $\square_0$, resp.\ $\square_n$.
The indicator $\mathbf{1}_{\Delta_z}$ acts on the first factor, so $\mathrm{Rw}_z\left(\delta_{\alpha,A}S^{\#}_{\mathrm{wk}}\right)=0$ unless $\mathbf{1}_{\Delta_z}\delta_{\alpha,A}h_{\square}\mathrm{Sloc}_{\square}\neq0$; assume the latter.
The difference quotient involves the spinor $h_{\square}\mathrm{Sloc}_{\square}f$ only at the two points of the pair.
Hence there is a short pair $(x,x')$, $x=(c,g)$, $x'=(c',g')$, with $c\in\Delta_z$ and with $h_{\square}(c)\neq0$ or $h_{\square}(c')\neq0$.
Since $|c'-z|\le|c'-c|+|c-z|<3/2$, the set $\operatorname{supp}h_{\square}$ contains a point at distance less than $3/2$ from $z$, and consequently $d(z,z_{\square})<3M/4+3/2$.
Choose a lift of $z$.
Every cube of this kind then possesses a lift of its centre in $M\mathbb{Z}^3$ with $|z_\mu-z_{\square,\mu}|<3M/4+3/2$ for $\mu=1,2,3$.
By [P.3], an open interval of length $3M/2+3<2M$ contains at most two multiples of $M$.
So at most $8$ cubes can stand at the left end, and the counting above gives
\begin{equation*}
\sum_{\mathrm{wk}\in\mathrm{Wk}_n}\mathrm{wp}^{n+1}\cdot\mathrm{Rw}_z\left(\delta_{\alpha,A}S^{\#}_{\mathrm{wk}}\right)
\le 8\cdot 27^n\cdot\mathrm{wp}^{n+1}\cdot 19\,a'\cdot\left(\frac{2\pi C'_{V}a'}{M}\right)^{n}
=152\,a'\cdot\mathrm{wp}\cdot\mathrm{qq}(\mathrm{wp})^{n}.
\end{equation*}
The supremum over $z$ yields the second bound of (e).

\textbf{(f)} ($k=0$.)
At level $0$ the unit blocks are the sites, and $\Delta_z\times O_{\mathrm{cub}}$ is the cluster over $z$.
The quantity of the statement, summed over $z'$ with the weight $\exp\left(\gamma^{F}d(z,z')\right)$, is therefore the number $\mathrm{Rk}_z\left(S^{\#}_{\mathrm{wk}}\right)$ of [P.6].
In each product exactly one factor is estimated through its kernel, namely the last local inverse.
For $n=0$, (Rw5), [7.1] and (d) give
\begin{equation*}
\mathrm{Rk}_z\left(h_0\,\mathrm{Sloc}_0\,h_0\right)\le|h_0(z)|\cdot N^{\mathrm{ke}}_{\gamma^{F}}(\mathrm{Sloc}_0)\le 2a'\cdot w_0(y),
\qquad
\mathrm{Rk}_z\left(h_0\,\mathrm{Sloc}_0\,h_0\right)=0\ \text{ if } h_0(z)=0 .
\end{equation*}
For $n\ge1$ the same three facts give
\begin{align*}
N^{\mathrm{ke}}_{\gamma^{F}}(O_n)
&=N^{\mathrm{ke}}_{\gamma^{F}}\left(\mathrm{Cm}_n\left(\mathrm{Sloc}_n\,h_n\right)\right)
\le N(\mathrm{Cm}_n)\cdot N^{\mathrm{ke}}_{\gamma^{F}}\left(\mathrm{Sloc}_n\,h_n\right)
\le\frac{\pi C'_{V}}{M}\cdot 2a'\cdot w_0(y)\le\frac{2\pi C'_{V}a'}{M},\\
N^{\mathrm{ke}}_{\gamma^{F}}\left(O'_0\,h_0\right)
&=N^{\mathrm{ke}}_{\gamma^{F}}\left(\left(\mathrm{Cm}_1\,h_0\right)\left(\mathrm{Sloc}_0\,h_0\right)\right)
\le N(\mathrm{Cm}_1)\cdot N(h_0)\cdot N^{\mathrm{ke}}_{\gamma^{F}}\left(\mathrm{Sloc}_0\,h_0\right)
\le\frac{2\pi C'_{V}a'}{M},\\
\mathrm{Rk}_z\left(S_{y,\mathrm{wk}}(A)\right)
&\le\mathrm{Rw}_z\left(O_0\cdots O_{n-1}\right)\cdot N^{\mathrm{ke}}_{\gamma^{F}}(O_n)
\le 2a'\cdot w_0(y)\cdot\left(\frac{2\pi C'_{V}a'}{M}\right)^{n-1}\cdot\frac{2\pi C'_{V}a'}{M},\\
\mathrm{Rk}_z\left(S'_{y,\mathrm{wk}}(A)\right)
&\le\mathrm{Rw}_z\left(O'_n\cdots O'_1\right)\cdot N^{\mathrm{ke}}_{\gamma^{F}}\left(O'_0\,h_0\right)
\le 2a'\cdot w_0(y)\cdot\left(\frac{2\pi C'_{V}a'}{M}\right)^{n-1}\cdot\frac{2\pi C'_{V}a'}{M}.
\end{align*}
Moreover, by (Rw2) and (Rw3), $\mathrm{Rw}_z\left(O_0\cdots O_{n-1}\right)=0$ unless $h_0(z)\neq0$, and $\mathrm{Rw}_z\left(O'_n\cdots O'_1\right)=0$ unless $h_n(z)\neq0$.
The counting of (e) now yields
\begin{multline*}
\sum_{z'}\exp\left(\gamma^{F}d(z,z')\right)\sum_{\mathrm{wk}\in\mathrm{Wk}_n}\mathrm{wp}^{n+1}\sup\left\{\left|S^{\#}_{\mathrm{wk}}(x,x')\right|_2:\ x\text{ over }z,\ x'\text{ over }z'\right\}\\
=\sum_{\mathrm{wk}\in\mathrm{Wk}_n}\mathrm{wp}^{n+1}\cdot\mathrm{Rk}_z\left(S^{\#}_{\mathrm{wk}}\right)
\le 16\,a'\cdot\mathrm{wp}\cdot w_0(y)\cdot\mathrm{qq}(\mathrm{wp})^{n}.
\end{multline*}

\emph{Level $0$.}
At $k=0$ the operator is $F=D_A+\mu+b_{L}\,J\,P(A)\,J^{*}+iy\gamma_5$ on $\mathfrak{S}^{\mathrm{ch}}(1,\ell)$, with $y\in\mathbb{R}$.
The unit block of a site is its cluster, $\mathrm{Cm}_{\square}$ is given by the second display of [7.3], and the weight appearing in (d), (e) and (f) is $w_0(y)=(1+|y|)^{-1}$.
There is no short pair; accordingly the last bound of (d) and the second bound of (e) are void, and (f) is a clause for $k=0$ alone.
The proofs of (a)--(c), of every bound of (d) except the last one, and of the first bound of (e) are the same at all levels.

\emph{Uniformity.}
The right-hand sides of (d)--(f) involve only $C'_{V}$, $a'$, $\gamma^{F}$, $M$, $\mathrm{wp}$ and $w_k(y)$.
The bounds on which the argument rests hold with these constants for all data (SF) of Section~G1.1, all $y\in Y_k$ and all $A\in\tilde{G}_k$, and the numbers $8$ and $27$ are geometric counts.
Hence nothing depends on $k$, $\ell$, $\mu$, $e_k$, $A$, $\square$ or $z$, nor on $y$ except through the factor $w_k(y)$.
\qed

\subsection*{4. Proposition 2.8 (FE1$^{\mathrm{ch}}$, FE2$^{\mathrm{ch}}$; the level-$0$ part of FE3$^{\mathrm{ch}}$)}

\begin{proof}[Proof of Proposition~2.8]
Two classical facts enter.
The first is the Neumann series: if $N$ is a submultiplicative norm on the operators of a finite-dimensional space and $O$ is an operator with $N(O)<1$, then $1-O$ is invertible and $(1-O)^{-1}=\sum_{n\ge0}O^n$, the series converging in $N$ \cite[Section~5.6]{HornJohnson2013}.
The second is that a uniform limit of analytic functions is analytic \cite[Theorem~10.28]{Rudin1987}; for functions of several variables see \cite[Chapter~I, Section~A]{GunningRossi1965}.
Throughout the proof $N:=N_{\gamma^{F}}$ and $F:=F_y(A)$.
By (M5) and [P.1] one has $\mathrm{qq}(1)=54\pi C'_{V}a'/M\le1/2$.
The number of $M$-cubes of $T_k$ is $(\ell/M)^3$.

\textbf{(a)} \emph{A right inverse.}
Define
\begin{equation*}
\mathrm{Spx}:=\sum_{\square}h_{\square}\,\mathrm{Sloc}_{\square}\,h_{\square},
\qquad
\mathrm{Rr}:=\sum_{\square}\mathrm{Cm}_{\square}\,\mathrm{Sloc}_{\square}\,h_{\square},
\end{equation*}
where $\square$ ranges over the $M$-cubes of $T_k$.
Summation of Proposition~2.7\,(a) over $\square$, together with the first relation $\sum_{\square}h_{\square}^2=1$ of (6.1) in Section~G1.7, in which the sum extends over all $M$-cubes of $T_k$ \cite[Section~16.7.1]{QEDI}, gives
\begin{equation}
F\,\mathrm{Spx}=1-\mathrm{Rr}.
\tag*{[8.1]}
\end{equation}
Consider a unit site $z$.
If $\mathbf{1}_{\Delta_z}\mathrm{Cm}_{\square}\neq0$, Proposition~2.7\,(c) provides a point $c\in\Delta_z$ with $d(c,z_{\square})<3M/4+L$, whence $d(z,z_{\square})<3M/4+L+1/2$.
Choose a lift of $z$.
Every such cube has a lift of its centre in $M\mathbb{Z}^3$ with $|z_\mu-z_{\square,\mu}|<3M/4+L+1/2$ for $\mu=1,2,3$.
An open interval of length $3M/2+2L+1<3M$ (see [P.3]) contains at most three multiples of $M$.
Therefore at most $27$ cubes contribute to $\mathrm{Rw}_z(\mathrm{Rr})$, and (Rw1) together with Proposition~2.7\,(d) gives
\begin{equation*}
\mathrm{Rw}_z(\mathrm{Rr})\le 27\cdot\frac{2\pi C'_{V}a'}{M}=\mathrm{qq}(1).
\end{equation*}
Thus $N(\mathrm{Rr})\le\mathrm{qq}(1)\le1/2$.
By Section~0.3, $N$ is a submultiplicative norm with $N(1)=1$.
The Neumann series shows that $1-\mathrm{Rr}$ is invertible, and by [8.1] one has $F\,\mathrm{Spx}\,(1-\mathrm{Rr})^{-1}=1$.
So $F$, a linear map of the finite-dimensional space $\mathfrak{S}^{\mathrm{ch}}(\eta_k,\ell)$ (of dimension $192\,(\ell/\eta_k)^3$) into itself, possesses a right inverse.
It is therefore bijective, and
\begin{equation}
S_y(A)=\mathrm{Spx}\,(1-\mathrm{Rr})^{-1}=\sum_{n\ge0}\mathrm{Spx}\,\mathrm{Rr}^{n},
\tag*{[8.2]}
\end{equation}
where the series converges in $N$.
As in Section~G1.8 write $h_i:=h_{\square_i}$, $\mathrm{Sloc}_i:=\mathrm{Sloc}_{\square_i}$, $\mathrm{Cm}_i:=\mathrm{Cm}_{\square_i}$.
Expanding the finite sums over the cubes one obtains
\begin{equation*}
\mathrm{Spx}\,\mathrm{Rr}^{n}=\sum_{\square_0,\dots,\square_n}\left(h_0\,\mathrm{Sloc}_0\,h_0\right)\left(\mathrm{Cm}_1\,\mathrm{Sloc}_1\,h_1\right)\cdots\left(\mathrm{Cm}_n\,\mathrm{Sloc}_n\,h_n\right),
\end{equation*}
the sum extending over all sequences $(\square_0,\dots,\square_n)$ of $M$-cubes.
Each product contains the factors $h_{i-1}\mathrm{Cm}_i$, $1\le i\le n$, and by Proposition~2.7\,(c) these vanish unless $d(z_{\square_{i-1}},z_{\square_i})\le M$.
Only walks contribute, and
\begin{equation}
\mathrm{Spx}\,\mathrm{Rr}^{n}=\sum_{\mathrm{wk}\in\mathrm{Wk}_n}S_{y,\mathrm{wk}}(A).
\tag*{[8.3]}
\end{equation}

\emph{A left inverse.}
Define $\mathrm{Rr}':=-\sum_{\square}h_{\square}\,\mathrm{Sloc}_{\square}\,\mathrm{Cm}_{\square}$.
Summing Proposition~2.7\,(b) over $\square$ yields $\mathrm{Spx}\,F=1-\mathrm{Rr}'$.
By [Q.5], $\mathbf{1}_{\Delta_z}h_{\square}\neq0$ holds for at most $8$ cubes.
Hence $\mathrm{Rw}_z(\mathrm{Rr}')\le 8\cdot 2\pi C'_{V}a'/M\le\mathrm{qq}(1)$, so that $N(\mathrm{Rr}')\le1/2$ and $(1-\mathrm{Rr}')^{-1}\,\mathrm{Spx}\,F=1$.
As $F$ is bijective, it follows that
\begin{equation*}
S_y(A)=(1-\mathrm{Rr}')^{-1}\,\mathrm{Spx}=\sum_{n\ge0}\mathrm{Rr}'^{\,n}\,\mathrm{Spx},
\end{equation*}
with convergence in $N$.
In the expansion, denote the cubes of the $n$ factors $\mathrm{Rr}'$, read from the left to the right, by $\square_n,\square_{n-1},\dots,\square_1$, and let $\square_0$ be the cube of $\mathrm{Spx}$.
Then
\begin{align*}
\mathrm{Rr}'^{\,n}\,\mathrm{Spx}
&=\sum\left(-h_n\,\mathrm{Sloc}_n\,\mathrm{Cm}_n\right)\left(-h_{n-1}\,\mathrm{Sloc}_{n-1}\,\mathrm{Cm}_{n-1}\right)\cdots\left(-h_1\,\mathrm{Sloc}_1\,\mathrm{Cm}_1\right)\left(h_0\,\mathrm{Sloc}_0\,h_0\right)\\
&=\sum\left(h_n\,\mathrm{Sloc}_n\right)\left(-\mathrm{Cm}_n\,h_{n-1}\,\mathrm{Sloc}_{n-1}\right)\cdots\left(-\mathrm{Cm}_1\,h_0\,\mathrm{Sloc}_0\right)h_0 ,
\end{align*}
where the two lines differ only in the position of the brackets, the order of the factors being the same.
Each product contains the factors $\mathrm{Cm}_i\,h_{i-1}$, which vanish unless $d(z_{\square_i},z_{\square_{i-1}})\le M$.
Hence
\begin{equation}
\mathrm{Rr}'^{\,n}\,\mathrm{Spx}=\sum_{\mathrm{wk}\in\mathrm{Wk}_n}S'_{y,\mathrm{wk}}(A).
\tag*{[8.4]}
\end{equation}

\emph{Absolute and uniform convergence.}
Consider a walk $\mathrm{wk}$ of length $n$ and a connection $A\in\tilde{G}_k$.
By Section~G1.8, $S_{y,\mathrm{wk}}(A)$ is the product of $h_0\,\mathrm{Sloc}_0\,h_0$ and of the $n$ operators $\mathrm{Cm}_i\,\mathrm{Sloc}_i\,h_i$, $1\le i\le n$.
Moving the brackets of $S'_{y,\mathrm{wk}}(A)$ as in the display preceding [8.4], one has
\begin{equation*}
S'_{y,\mathrm{wk}}(A)=(-1)^{n}\left(h_n\,\mathrm{Sloc}_n\,\mathrm{Cm}_n\right)\left(h_{n-1}\,\mathrm{Sloc}_{n-1}\,\mathrm{Cm}_{n-1}\right)\cdots\left(h_1\,\mathrm{Sloc}_1\,\mathrm{Cm}_1\right)\left(h_0\,\mathrm{Sloc}_0\,h_0\right).
\end{equation*}
For every cube $\square$, Proposition~2.7\,(d) gives $N\left(h_{\square}\,\mathrm{Sloc}_{\square}\,h_{\square}\right)\le 2a'\cdot w_k(y)$, and it bounds both $N\left(\mathrm{Cm}_{\square}\,\mathrm{Sloc}_{\square}\,h_{\square}\right)$ and $N\left(h_{\square}\,\mathrm{Sloc}_{\square}\,\mathrm{Cm}_{\square}\right)$ by $2\pi C'_{V}a'/M$.
The inequality (Rw2) yields submultiplicativity, $N(O_1O_2)\le N(O_1)\cdot N(O_2)$.
Consequently
\begin{equation*}
N\left(S^{\#}_{\mathrm{wk}}\right)\le 2a'\cdot w_k(y)\cdot\left(\frac{2\pi C'_{V}a'}{M}\right)^{n}.
\end{equation*}
Moreover $\#|\mathrm{Wk}_n|\le(\ell/M)^3\cdot 27^n$: the cube $\square_0$ is one of the $(\ell/M)^3$ cubes, and for a given cube $\square$ at most $27$ cubes $\square'$ satisfy $d(z_{\square},z_{\square'})\le M$, every coordinate of $z_{\square'}-z_{\square}$ being congruent to $-M$, $0$ or $M$ modulo $\ell$.
Since $27\cdot 2\pi C'_{V}a'/M=\mathrm{qq}(1)$, one obtains
\begin{equation}
\begin{aligned}
&\sum_{\mathrm{wk}\in\mathrm{Wk}}N\left(S^{\#}_{y,\mathrm{wk}}(A)\right)\le(\ell/M)^3\cdot 2a'\cdot w_k(y)\sum_{n\ge0}\mathrm{qq}(1)^{n}\le 4a'\cdot(\ell/M)^3\cdot w_k(y),\\
&\sum_{\mathrm{wk}:\,|\mathrm{wk}|\ge j}N\left(S^{\#}_{y,\mathrm{wk}}(A)\right)\le 4a'\cdot(\ell/M)^3\cdot 2^{-j}\qquad(j\ge0\ \text{an integer}),
\end{aligned}
\tag*{[8.5]}
\end{equation}
for all $A\in\tilde{G}_k$ and all $y\in Y_k$.
The two series therefore converge absolutely in $N$ (criterion (Nc) of Section~G1.3), uniformly on $\tilde{G}_k$ and in $y$, and by [8.2]--[8.4] the sum of each of them is $S_y(A)$.
In localised form, Proposition~2.7\,(e) with $\mathrm{wp}=1$ provides a bound that does not involve $\ell$:
\begin{equation}
\begin{split}
&\sup_z\sum_{z'}\exp\left(\gamma^{F}d(z,z')\right)\sum_{\mathrm{wk}\in\mathrm{Wk}}\left\|\mathbf{1}_{\Delta_z}S^{\#}_{\mathrm{wk}}\mathbf{1}_{\Delta_{z'}}\right\|_{\mathrm{su}}\\
&\qquad\le\sum_{n\ge0}16\,a'\cdot w_k(y)\cdot\mathrm{qq}(1)^{n}\le 32\,a'\cdot w_k(y).
\end{split}
\tag*{[8.6]}
\end{equation}

\textbf{(b)}
By [Q.3]\,(G2), $\tilde{G}_k$ is an open subset of the finite-dimensional complex vector space of the connections on $T_k$.
Fix a walk $\mathrm{wk}$.
The map that sends $A$ to its restriction to the bonds of the set $\tilde{W}$ of $\square_i$ is linear and maps $\tilde{G}_k$ into $\mathrm{Gloc}(\delta_k)$ ([Q.3]\,(G1)).
The map $\mathrm{Sloc}_{\square_i,y}(\cdot)$ is analytic on $\mathrm{Gloc}(\delta_k)$ by Proposition~2.5\,(a) and [P.2], and $O\mapsto E_{\square_i}(O)$ is linear by Proposition~2.6\,(a).
Hence $A\mapsto\mathrm{Sloc}_i$ is analytic on $\tilde{G}_k$.
By (F4), the entries of $\mathrm{Cm}_i=h_i\,F_y(A)-F_y(A)\,h_i$ are entire functions of the numbers $A(b)$, and $h_i$ does not depend on $A$.
Being a finite product of such factors, $S_{y,\mathrm{wk}}(A)$ is analytic on $\tilde{G}_k$.
By [8.5] and the inequality $\|O\|_{\mathrm{su}}\le N(O)$ of (Rw4), the series $\sum_{\mathrm{wk}}S_{y,\mathrm{wk}}(A)$ converges uniformly on $\tilde{G}_k$, entry by entry.
Consequently the sum $S_y(A)$ is analytic on $\tilde{G}_k$.

\textbf{(c)}
Combining (a), (Rw1), which extends to the convergent series because $\|\cdot\|_{\mathrm{su}}$ is continuous, and [8.6], one finds for every unit site $z$
\begin{equation*}
\mathrm{Rw}_z\left(S_y(A)\right)\le\sum_{\mathrm{wk}\in\mathrm{Wk}}\mathrm{Rw}_z\left(S_{y,\mathrm{wk}}(A)\right)\le 32\,a'\cdot w_k(y);
\end{equation*}
hence $N\left(S_y(A)\right)\le 32\,a'\cdot w_k(y)$.
Restricted to the short pairs, $\delta_{\alpha,A}$ is a linear map on the finite-dimensional space $\mathfrak{S}^{\mathrm{ch}}(\eta_k,\ell)$.
For the compositions on $T_k$ (Section~G1.6) it follows that $\delta_{\alpha,A}S_y(A)=\sum_{\mathrm{wk}\in\mathrm{Wk}}\delta_{\alpha,A}S_{y,\mathrm{wk}}(A)$, and the second bound of Proposition~2.7\,(e) with $\mathrm{wp}=1$ gives
\begin{equation*}
\mathrm{Rw}_z\left(\delta_{\alpha,A}S_y(A)\right)\le\sum_{n\ge0}152\,a'\cdot\mathrm{qq}(1)^{n}\le 304\,a'
\end{equation*}
for every unit site $z$, that is $N\left(\delta_{\alpha,A}S_y(A)\right)\le304\,a'$.
The two inequalities displayed in (c) are the instances $O=S_y(A)$ and $O=\delta_{\alpha,A}S_y(A)$ of the inequality $\left\|\mathbf{1}_{\Delta_z}O\,\mathbf{1}_{\Delta_{z'}}\right\|_{\mathrm{su}}\le\exp\left(-\gamma^{F}d(z,z')\right)\cdot N(O)$ of (Rw4).
The short pairs are the pairs with $d_X<1$; for $k\ge1$ they include the $48\cdot47$ ordered pairs inside each cluster, for which $d_X=\eta_k$ and no phase occurs.
When $k=0$ the second inequality is empty, there being no short pairs.

\textbf{(d)} ($k=0$, $Y_0=\mathbb{R}$.)
The kernel (I.1.4) of an operator depends linearly on the operator and, the dimension being finite, continuously; hence $S_{0,y}(A)(x,x')=\sum_{\mathrm{wk}\in\mathrm{Wk}}S_{y,\mathrm{wk}}(A)(x,x')$.
By (Rw5) and Proposition~2.7\,(f) with $\mathrm{wp}=1$, every site $z$ of $T^0_0$ satisfies
\begin{equation*}
\mathrm{Rk}_z\left(S_{0,y}(A)\right)\le\sum_{n\ge0}\sum_{\mathrm{wk}\in\mathrm{Wk}_n}\mathrm{Rk}_z\left(S_{y,\mathrm{wk}}(A)\right)\le\sum_{n\ge0}16\,a'\cdot w_0(y)\cdot\mathrm{qq}(1)^{n}\le 32\,a'\cdot w_0(y).
\end{equation*}
Therefore $N^{\mathrm{ke}}_{\gamma^{F}}\left(S_{0,y}(A)\right)\le 32\,a'\cdot(1+|y|)^{-1}$, and the last inequality of (Rw5) gives
\begin{equation*}
\left|S_{0,y}(A)(x,x')\right|_2\le 32\,a'\cdot(1+|y|)^{-1}\cdot\exp\left(-\gamma^{F}d(x,x')\right)
\qquad\text{for all fermion sites } x,x' \text{ and all real } y .
\end{equation*}
At $y=0$ one has $S_{0,0}(A)=F_0(A)^{-1}=\Gamma_0(A)=S^{0}_{1}(A)$ (Section~I.3.9).
By (a) with $k=0$ and $y=0$, the operator $F_0(A)=D_A+\mu+b_{L}\,J\,P(A)\,J^{*}$ is bijective; hence $\delta Z_0(A)=\det F_0(A)\neq0$.
In accordance with (I.ke5), the factor $48$ of the kernel of the identity enters this bound only once, through $C_{00}$ in $a'$: in each product a single factor, a local inverse, is estimated in the kernel norm, and Proposition~2.5\,(b) bounds it by means of Proposition~2.4\,(d).

\emph{The final assertion of Proposition~2.8.}
Let $k\ge1$.
By Section~I.3.5, $S_k(A)=S_{\infty}(A)$, $S_{k,y}(A)=S_y(A)$ for real $y$, and $S^{0}_{k+1}(A)=S_0(A)$.
Accordingly (a)--(c) at $y=\infty$ yield FE1$^{\mathrm{ch}}$, and (a)--(c) at real $y$ yield FE2$^{\mathrm{ch}}$, once the norm $|\cdot|_{\mathrm{su}}$ is replaced by $|\cdot|_{\infty}$ at the cost of the factor $2$ (Sections~I.1.5 and G1.3); this puts $64a'$ and $608a'$ in the place of $32a'$ and $304a'$.
Here FE1$^{\mathrm{ch}}$ and FE2$^{\mathrm{ch}}$ are the clauses of Proposition~14.1\,(i), which are stated in the norm $|\cdot|_{\infty}$ with the constant $C^{F}$; they follow from the bounds above because
\begin{equation*}
64a'\le608a'\le C^{F}.
\end{equation*}
Indeed, the row of $C^{F}$ in the line (C6) of Proposition~14.2\,(a) reads
\begin{equation*}
C^{F}=2\max\left\{C_1,\ C_{\mathrm{By}},\ 304\,\mathrm{e}\cdot\max(1,b)\cdot a',\ C_{\Gamma},\ \mathrm{e}\cdot b^{2}\cdot(32a')^{2}\right\},
\end{equation*}
and the third entry of the maximum gives $C^{F}\ge608\,\mathrm{e}\cdot\max(1,b)\cdot a'\ge608a'$, since $\mathrm{e}\ge1$, $\max(1,b)\ge1$ and $a'>0$; the first inequality holds because $a'>0$.
The same inequalities are contained in Chapter~14a, Lemma~1\,(K2), which the proof of Proposition~14.1\,(i) applies.

\emph{Uniformity.}
The numbers $32a'$, $304a'$ and $\gamma^{F}$ are constants of the line (C4r).
The bounds of (c), (d) and [8.6] hold for all data (SF), that is for every level $k\ge0$, every side $\ell\in\mathrm{Cl}$, every $\mu\in(0,1]$ and every $e_k\in(0,e_0]$, and for all $y\in Y_k$ and $A\in\tilde{G}_k$.
The cutoff, that is the integer $N$ of Section~I.0.5 with $\eta=L^{-N}$ (in this proof the letter $N$ stands for the norm $N_{\gamma^{F}}$), the torus and the charge $e$ enter a flow only through such data; the bounds are therefore uniform in them.
Of the bounds above, only [8.5] involves $\ell$.
\end{proof}

\subsection*{5. Conditions on the constants}

The proofs of Propositions~2.5--2.8 use the following conditions on the constants; for each of them the table gives the place where it is stated and the places where it enters.

\noindent
\begin{tabular}{p{0.25\textwidth}p{0.23\textwidth}p{0.42\textwidth}}
\hline
condition & stated in & enters in\\
\hline
$b=1>0$
& Section~I.8.2, line (C0); Proposition~14.2\,(a)
& [P.1] ($C'_{V}\ge24\,\mathrm{e}$); (F4) ($b_k>0$)\\
$\alpha<1$
& Section~I.8.2, line (C1); Proposition~14.2\,(a), (f)
& Proposition~2.7\,(d): $|c-c'|^{1-\alpha}\le1$ on short pairs\\
$\varepsilon>0$
& Section~I.8.2, line (C7); Proposition~14.2\,(a)
& $e_k^{\varepsilon}\le e_0^{\varepsilon}$ for $e_k\le e_0$, in [P.2] and [7.2]\\
(L5): $L$ odd, $L\ge5$
& Proposition~14.2\,(b)
& $L$ odd: nested blocks (Section~0.2), [Q.1], [5.1], Proposition~2.5\,(e).
$L\ge5$: $\exp\left(\gamma^{F}\right)\le\exp(1/30)$ in [P.1]; $M>20$ and $\ell\ge65M$ in the proof of [P.3]\\
(M1): $M$ a power of $L$ ($M\ge L^2$)
& Proposition~14.2\,(b)
& $6M$ is a multiple of $L$ and $6M\ge3L$ ($\mathbb{T}^{W}$ is a block torus); [Q.1]\,(a); [5.1]; $M\mathbb{Z}^3\subset L^j\eta_k\mathbb{Z}^3$\\
(M2): $M>4L$
& Proposition~14.2\,(b)
& [P.3]: Proposition~2.7\,(c) ($3M/4+L<3M/2-\eta_k$; $3M/2+L<2M$), Proposition~2.7\,(e) ($3M/2+3<2M$), Proposition~2.8\,(a) ($3M/2+2L+1<3M$); through [P.3] also [Q.5] ($3M/2+1<2M$, which gives the count $8$ used in Proposition~2.7\,(e), (f) and Proposition~2.8\,(a)), Step~0 of the proof of Proposition~2.7 ($3M/4\le3M/2-\eta_k$, so that $\operatorname{supp}h_{\square}\cap\tilde{W}\subset W^{-}$, used in Proposition~2.7\,(a), (b)) and Proposition~2.7\,(d) ($3M/4+1<3M/2$)\\
(M5): $M\ge M^{F}$, in particular $M\ge108\pi C'_{V}a'$
& Proposition~14.2\,(b), (f)
& $M\ge1600$ in Proposition~2.7\,(d) (the number $19$); $\mathrm{qq}(1)\le1/2$ in Proposition~2.8\\
(Z2.s1): $6L\,e_0^{\varepsilon}\le1$
& Proposition~14.2\,(b); Section~I.8.4
& [P.2] ($3L\delta_k\le1$); [7.2] ($3L\cdot\mathfrak{i}(A)\le1$, $\exp\left(3\,\mathfrak{i}(A)\right)\le\mathrm{e}$)\\
(Z2.s2): $4C'_{V}a'\,e_0^{\varepsilon}\le1$
& Proposition~14.2\,(b); Section~I.8.4
& [P.2] ($2C'_{V}a'\,\delta_k\le1$)\\
the definitions $C'_{V}=2\mathrm{e}\,(42+6b)$, $\gamma^{F}=\min\left(\gamma_{(5)},a_{z}\right)/2$, $a'=\max\left(C_{5},C_{00}\right)\cdot c(\gamma^{F})$; $0<\gamma_{(5)}\le1/(3L)$, $C_{5}\ge1$
& Proposition~14.2\,(a), lines (C4), (C4r); Proposition~14.2\,(f); Proposition~1.6
& [P.1]\\
$64a'\le608a'\le C^{F}$
& Proposition~14.2\,(a), line (C6): the row of $C^{F}$, whose third entry gives $C^{F}\ge608\,\mathrm{e}\cdot\max(1,b)\cdot a'$
& the final assertion of Proposition~2.8 (FE1$^{\mathrm{ch}}$, FE2$^{\mathrm{ch}}$ in the norm $|\cdot|_{\infty}$, with $64a'$ and $608a'$ in the place of $32a'$ and $304a'$)\\
\hline
\end{tabular}

The three inequalities $3M/2+1<2M$, $3M/4\le3M/2-\eta_k$ and $3M/4+1<3M/2$ named last in the row (M2) follow from [P.3], the first from its third inequality and the other two from its first.
Each of them holds as soon as $M>2$ (for the second one recall that $\eta_k\le1$), hence already under (L5) and (M1), which give $M\ge L^2\ge25$.

The data are subject to the following requirements.
The side $\ell\in\mathrm{Cl}$, a multiple of $LM$ with $\ell\ge13LM$, gives $3M<\ell/2$ ([P.3]; this enters [6.1], [Q.1]\,(c) and [Q.3]) and makes $T_k$ a block torus.
Further $\mu\in(0,1]$ and $e_k\in(0,e_0]$; in Proposition~2.5 only $e_k>0$ is required, the number $\delta$ being a datum there.

Of the conditions (L5), (M1), (M2), (M5), (Z2.s1) and (Z2.s2), only (L5) and (M1) enter the proofs of Propositions~2.5 and~2.6.
Besides these, the proof of Proposition~2.5\,(a) uses $b>0$ through (F4) ($b_k>0$), for the continuity in $(y,A)$ at real $y$, and the proof of Proposition~2.5 uses the two conditions on the datum $\delta$.
The conditions (M2), (M5), (Z2.s1) and (Z2.s2) enter in Propositions~2.7 and~2.8, and Corollary~6.A rests on (Z2.s1) and (Z2.s2) through [P.2].

\section*{Chapter 3. Weakened operators and polymer propagators on the fine lattice: proof of Proposition 3}

\subsection*{Proof}

\subsubsection*{0. Conventions of this chapter: setting; operators, norms, series, kernels}

The labels (G1) and (G2) refer to two properties of the domain $\tilde{G}_k$ of complex connections that are proved under these labels in \cite[Section~16.2.7]{QEDI}.
By (G1), for every $A\in\tilde{G}_k$ and every $M$-cube $\square$ there are a real function $\boldsymbol{\lambda}$ on the sites of $\square^{\sim1}$ and a complex function $B$ on the bonds of $\square^{\sim1}$ such that $A=B+d\boldsymbol{\lambda}$ on these bonds and $e_k\cdot|B|_{\infty}<2e_k^{\varepsilon}$.
By (G2), $\tilde{G}_k$ is an open subset of the space of connections on $T_k$, and together with $A$ also $A-d\boldsymbol{\lambda}$, $-A$, $\overline{A}$ and $\tilde{g}A$ belong to $\tilde{G}_k$, for every real function $\boldsymbol{\lambda}$ of $c$ and every $\tilde{g}\in\mathbb{G}_k$.

Throughout the chapter $(k,\ell,\mu,e_k)$ is a quadruple of data (SF), $y\in Y_k$ and $A\in\tilde{G}_k$.
The torus $T_k$ is the block torus $(0,\ell,-1)$ of Section~G1.2, and $N:=N_{\gamma^{F}}$ is the norm of Section~G1.3 on this torus.
The letters $z$, $z'$ stand for sites of $T^0_k$; there are $n_k=\ell^3$ of them.
Put $\mathbf{1}_z:=\mathbf{1}_{\Delta_z}$.
On spinors $\mathbf{1}_z$ is the multiplication by the indicator of the orbit of $\Delta_z\times O_{\mathrm{cub}}$ under the translations by $\ell\mathbb{Z}^3$; on functions of short pairs it is the restriction to those pairs whose first point lies over this orbit (Section~G1.3, (I.POS)).
For $k=0$ one has $\Delta_z=\{z\}$, and $\mathbf{1}_z$ is the indicator of the cluster $\{z\}\times O_{\mathrm{cub}}$, whose $48$ sites carry the weight $\frac{1}{48}$ each, so that the cluster has measure $1$.

Denote by $\mathrm{Lop}$ the space of linear operators on $\mathfrak{S}^{\mathrm{ch}}(\eta_k,\ell)$, and by $\mathrm{Lop}^{\mathrm{sp}}$ the space of linear maps from $\mathfrak{S}^{\mathrm{ch}}(\eta_k,\ell)$ to the functions of short pairs.
By Section~G1.3 the short pairs are the ordered pairs $(x,x')$ of distinct points of $\mathbb{L}^{\mathrm{f}}_k$ with $d_X(x,x')<1$; for $k\ge1$ the $48\cdot47$ ordered pairs inside each cluster, for which $d_X=\eta_k$, are among them.
One has $\dim\mathfrak{S}^{\mathrm{ch}}(\eta_k,\ell)=192\left(\frac{\ell}{\eta_k}\right)^3$ (Section~I.1.5).
Let $k\ge1$.
Then $\eta_k<1$, and $\frac{1}{\eta_k}$ is an integer.
A point $x=(c,g)$ is the first point of exactly
\[
48\left(\frac{2}{\eta_k}-1\right)^3-1
\]
short pairs; their second points are the points $x'=(c',g')\ne x$ with $|c'-c|<1$, equivalently with $|c'-c|\le1-\eta_k$.
Each short pair has exactly one translate by a vector of $\ell\mathbb{Z}^3$ whose first point lies over a given fundamental domain of $\ell\mathbb{Z}^3$ in $\mathbb{L}_k$.
A function $F$ of short pairs takes values in $\mathbb{C}^4$ and satisfies $F(x+\ell e_{\mu},x'+\ell e_{\mu})=-F(x,x')$ (Section~G1.3).
It is therefore determined by its values at the short pairs whose first point lies over this domain, and these values can be prescribed arbitrarily.
Hence the functions of short pairs form a space of dimension
\[
4\cdot48\cdot\left(\frac{\ell}{\eta_k}\right)^3\cdot\left(48\left(\frac{2}{\eta_k}-1\right)^3-1\right).
\]
Consequently $\mathrm{Lop}$ and $\mathrm{Lop}^{\mathrm{sp}}$ are of finite dimension.
For $k=0$ there is no short pair, and $\mathrm{Lop}^{\mathrm{sp}}=\{0\}$.
For $O\in\mathrm{Lop}$, the element $\delta_{\alpha,A}O$ of $\mathrm{Lop}^{\mathrm{sp}}$ is the composition of $O$ with the linear map $f\mapsto\delta_{\alpha,A}f$ of (I.1.7), restricted to the short pairs (Section~G1.6).

\paragraph{0.1. $N$ is a norm.}
The blocks of one side partition the lattice (Section~G1.2); hence each photon site lies in exactly one unit block.
It follows that $\sum_z\mathbf{1}_z=1$ on spinors as well as on functions of short pairs, and that every $O$ in $\mathrm{Lop}$ or in $\mathrm{Lop}^{\mathrm{sp}}$ can be written as
\[
O=\sum_{z,z'}\mathbf{1}_z\cdot O\cdot\mathbf{1}_{z'}.
\]
Therefore $N(O)=0$ implies $O=0$.
Moreover $N(t\cdot O)=|t|\cdot N(O)$ for complex $t$, and $N(O_1+O_2)\le N(O_1)+N(O_2)$ by (G1.N2) of Section~G1.3.
Since a supremum of $n_k$ non-negative numbers does not exceed their sum, one has
\begin{equation}
N(O)\le\sum_z\sum_{z'}\exp\left(\gamma^{F}\cdot d(z,z')\right)\cdot\left\|\mathbf{1}_z\cdot O\cdot\mathbf{1}_{z'}\right\|_{\mathrm{su}}\qquad\left(O\in\mathrm{Lop}\ \text{or}\ O\in\mathrm{Lop}^{\mathrm{sp}}\right).
\tag{3-0.1}
\end{equation}

\paragraph{0.2. Series.}
By (G1.Nc) of Section~G1.3, a series $\sum_jO_j$ in $\mathrm{Lop}$ or in $\mathrm{Lop}^{\mathrm{sp}}$ converges absolutely in $N$ if $\sum_jN(O_j)<\infty$.
A normed space of finite dimension is complete.
Hence such a series converges, its sum does not depend on the order of the terms, and the terms may be collected in groups.
Every real-linear map between spaces of finite dimension is continuous, so that it may be applied to a convergent series term by term.
Maps of this kind are $O\mapsto O_1\cdot O\cdot O_2$ (with $O_1$, $O_2$ fixed), $O\mapsto\mathbf{1}_z\cdot O\cdot\mathbf{1}_{z'}$, $O\mapsto\delta_{\alpha,A}O$, $O\mapsto O(x,x')$ (paragraph~0.3), $O\mapsto O^T$ and $O\mapsto\overline{O}$ (paragraph~0.4).

\paragraph{0.3. Kernels.}
For $O\in\mathrm{Lop}$ and $x,x'\in\mathbb{L}^{\mathrm{f}}_k$, the $4\times4$ matrix $O(x,x')$ is the kernel (I.1.4), taken with the rules of Section~G1.2 for $R=\ell$ and $\boldsymbol{\sigma}=-1$.
Thus
\[
(Of)(x)=\int\mathrm{d}x'\,O(x,x')\cdot f(x')=\eta_k^3\cdot\frac{1}{48}\sum_{c'}\sum_{g'\in O_{\mathrm{cub}}}O\left(x,(c',g')\right)\cdot f(c',g'),
\]
where $c'$ ranges over a fundamental domain of $\ell\mathbb{Z}^3$;
\[
O(x+\ell e_{\mu},x')=O(x,x'+\ell e_{\mu})=-O(x,x');
\]
and $(\psi\cdot O\cdot\psi')(x,x')=\psi(c)\cdot O(x,x')\cdot\psi'(c')$ for $\ell\mathbb{Z}^3$-periodic functions $\psi$, $\psi'$ of position (the scalars of Section~G1.2).
For $x'=(c',g')\in\mathbb{L}^{\mathrm{f}}_k$ and $v\in\mathbb{C}^4$ let $e_{x',v}$ be the element of $\mathfrak{S}^{\mathrm{ch}}(\eta_k,\ell)$ that takes the value $(-1)^{n_1+n_2+n_3}v$ at the sites $x'+\ell n$, $n\in\mathbb{Z}^3$, and the value $0$ at all other sites.
If the fundamental domain is chosen so as to contain $c'$, the sum above has only one non-vanishing term, and
\begin{equation}
(Oe_{x',v})(x)=\eta_k^3\cdot\frac{1}{48}\cdot O(x,x')\cdot v,\qquad\text{that is,}\qquad O(x,x')\cdot v=48\eta_k^{-3}\cdot(Oe_{x',v})(x).
\tag{3-0.2}
\end{equation}
In particular, $O\mapsto O(x,x')$ is linear.
The row of $O$ at $x$ is the function $x'\mapsto O(x,x')$; two operators $O$, $\tilde{O}$ have the same row at $x$ if and only if $(Of)(x)=(\tilde{O}f)(x)$ for all $f$.
The column of $O$ at $x'$ is the function $x\mapsto O(x,x')$.
The modulus of $O(x,x')$, and the question whether $O(x,x')$ vanishes, depend only on the classes of $x$ and $x'$ modulo $\ell\mathbb{Z}^3$.
For a set $Z$ of sites of $T_k$, the words ``$x$ lies over $Z$'' mean that the class of $c=\pi(x)$ belongs to $Z$ ((I.POS), Section~G1.2).

\paragraph{0.4. Transposes and conjugates.}
The transpose ${}^T$ is taken with respect to the bilinear pairing
\[
\langle f',f\rangle^{\mathrm{ch}}=\eta_k^3\cdot\frac{1}{48}\sum_c\sum_{g\in O_{\mathrm{cub}}}f'(c,g)\cdot f(c,g)
\]
of Section~I.1.6; here $f'(c,g)\cdot f(c,g)$ is the bilinear dot product of the two vectors $f'(c,g)$, $f(c,g)$ of $\mathbb{C}^4$.
One has $(O_1\cdot O_2)^T=O_2^T\cdot O_1^T$ and $(O^T)^T=O$, and, the pairing being bilinear, $(t\cdot O)^T=t\cdot O^T$ for complex $t$.
For an $\ell\mathbb{Z}^3$-periodic function $\psi$ of $c$ the identity $f'(x)\cdot\left(\psi(c)f(x)\right)=\left(\psi(c)f'(x)\right)\cdot f(x)$ holds at every site; so $\psi^T=\psi$.
For a constant $4\times4$ matrix $\mathrm{Mx}$ one has $f'(x)\cdot\left(\mathrm{Mx}\,f(x)\right)=\left(\mathrm{Mx}^Tf'(x)\right)\cdot f(x)$; so the transpose of the operator $\mathrm{Mx}$ is the operator given by the transposed matrix $\mathrm{Mx}^T$.
The conjugate $\overline{O}$ is the operator $f\mapsto\overline{O\overline{f}}$ ((I.O5)).
One has $\overline{O_1\cdot O_2}=\overline{O_1}\cdot\overline{O_2}$ and $\overline{t\cdot O}=\overline{t}\cdot\overline{O}$ for complex $t$, $\overline{\psi}=\psi$ for a real function $\psi$, and $\overline{\mathrm{Mx}}$ is the matrix whose entries are the complex conjugates of those of $\mathrm{Mx}$.
Functions of $c$ commute with constant matrices.

\subsubsection*{1. Cubes, walks, weights}

\paragraph{1.1. Cubes.}
Since $\ell\in\mathrm{Cl}$, the side $\ell$ is a multiple of $M$.
As $M$ is a power of $L$ and $L$ is odd, $\frac{M}{\eta_k}$ is an odd integer.
Hence no site of $\mathbb{L}_k$ has a coordinate at distance $\frac{M}{2}$ from a point of $M\mathbb{Z}$, and the $M$-cubes
\[
\left\{c\in T_k:\ d(c,z)<\frac{M}{2}\right\},\qquad z\in M\mathbb{Z}^3/\ell\mathbb{Z}^3,
\]
form a partition of $T_k$ (Section~I.1.1, ``Faces'', and Section~I.1.4).
Two consequences follow.
\begin{itemize}
\item[(i)] An $M$-cube that has a site in common with a union $W$ of $M$-cubes is one of the cubes of $W$.
\item[(ii)] For unions $W$, $W'$ of $M$-cubes, $W\subset W'$ holds if and only if every cube of $W$ is a cube of $W'$.
\end{itemize}
Two $M$-cubes have a common face if and only if, for one direction $\mu$, their centres differ by $Me_{\mu}$ or by $-Me_{\mu}$ modulo $\ell\mathbb{Z}^3$.
Let $\square$ be an $M$-cube with the centre $z_{\square}$.
Then $\square^{\sim1}$ is the union of the cubes centred at the points $z_{\square}+Mn$, $n\in\{-1,0,1\}^3$.
These $27$ cubes are pairwise different, because $3M<\ell$ ($\ell\in\mathrm{Cl}$), and
\begin{equation}
\square^{\sim1}=\left\{c\in T_k:\ d(c,z_{\square})<\frac{3M}{2}\right\}.
\tag{3-1.1}
\end{equation}
Indeed, a site $c$ with $d(c,z_{\square})<\frac{3M}{2}$ has a lift $\tilde{c}$ with $|\tilde{c}_{\mu}-z_{\square,\mu}|<\frac{3M}{2}$ for $\mu=1,2,3$.
No coordinate of $\tilde{c}$ lies at distance $\frac{M}{2}$ from a point of $M\mathbb{Z}$.
So $\tilde{c}_{\mu}$ lies in exactly one of the three open intervals of length $M$ centred at the points $z_{\square,\mu}+Mn_{\mu}$, $n_{\mu}\in\{-1,0,1\}$, and $c$ lies in the cube centred at $z_{\square}+Mn$.
Conversely, every site of the latter cube has a lift $\tilde{c}$ with $|\tilde{c}_{\mu}-z_{\square,\mu}|<M+\frac{M}{2}$.
The set $\tilde{W}$ of Section~G1.5 is a lift of $\square^{\sim1}$, and the covering map is injective on $\tilde{W}$, because the sup-diameter of $\tilde{W}$ is smaller than $3M$ and $3M+\eta_k<\ell$.
A bond of $\tilde{W}$, that is, a bond with both ends in $\tilde{W}$, is mapped onto a bond of $T_k$ with both ends in $\square^{\sim1}$.

\paragraph{1.2. $X_{\mathrm{wk}}$ is a polymer, as asserted in Proposition~3\,(d).}
Let $\mathrm{wk}=(\square_0,\ldots,\square_n)\in\mathrm{Wk}_n$ and $z_i:=z_{\square_i}$.
For $n'=(n'_1,n'_2,n'_3)\in\{-1,0,1\}^3$, the cube of $\square_i^{\sim1}$ centred at $z_i+Mn'$ is joined to $\square_i$ by the cubes centred at
\[
z_i+M(n'_1,n'_2,n'_3),\qquad z_i+M(0,n'_2,n'_3),\qquad z_i+M(0,0,n'_3),\qquad z_i;
\]
these cubes belong to $\square_i^{\sim1}$, and any two consecutive ones either coincide or have a common face.
So $\square_i^{\sim1}$ is connected through faces.
The points $z_{i-1}$ and $z_i$ lie in $M\mathbb{Z}^3$, and $d(z_{i-1},z_i)\le M$.
Hence $z_i=z_{i-1}+Mn'$ modulo $\ell\mathbb{Z}^3$ for some $n'\in\{-1,0,1\}^3$, and the cube $\square_i$ belongs both to $\square_{i-1}^{\sim1}$ and to $\square_i^{\sim1}$.
Induction on $i$ shows that the union of $\square_0^{\sim1},\ldots,\square_i^{\sim1}$ is connected through faces.
Therefore $X_{\mathrm{wk}}$ is a non-empty union of $M$-cubes that is connected through faces, that is, $X_{\mathrm{wk}}\in D_k$.

\paragraph{1.3. Numbers: the bounds $|X_{\mathrm{wk}}|_M\le27(n+1)$ and $d_M(X_{\mathrm{wk}})\le27(n+1)$ of Proposition~3\,(d).}
The set $X_{\mathrm{wk}}$ is a union of $n+1$ sets of $27$ cubes; hence $|X_{\mathrm{wk}}|_M\le27(n+1)$.
For every $X\in D_k$ one has $d_M(X)\le|X|_M-1$.
To see this, consider the graph whose vertices are the cubes of $X$ and whose edges are the pairs of cubes with a common face.
This graph is connected; so it has a spanning tree, with $|X|_M-1$ edges.
With each edge associate the image in $\mathbb{R}^3/\ell\mathbb{Z}^3$ of the closed segment of $\mathbb{R}^3$ that joins a lift $\tilde{z}$ of the centre of one of its two cubes to the lift $\tilde{z}\pm Me_{\mu}$ of the centre of the other; its Euclidean length is $M$.
The tree being connected, the union of these $|X|_M-1$ segments (the centre of the cube, if $|X|_M=1$) is a net in the sense of \cite[Section~5.2.3]{QEDI} (see Section~I.1.1).
This net contains the centre of every cube of $X$, and the lengths of its segments add up to $(|X|_M-1)M$.
Since $M\cdot d_M(X)$ is the infimum of the lengths of the nets that meet every closed $M$-cube of $X$, the inequality follows.
Consequently
\begin{equation}
d_M(X_{\mathrm{wk}})\le27(n+1).
\tag{3-1.2}
\end{equation}
The set $\mathrm{Wk}_n$ is finite, with $\#|\mathrm{Wk}_n|=\left(\frac{\ell}{M}\right)^3\cdot27^n$.
The set $D_k$ is finite.

\paragraph{1.4. Weights.}
Put
\[
\mathrm{wp}_s:=M^{27\alpha_0},\qquad\mathrm{wp}_X:=\exp(27\kappa').
\]
Since $M\ge1$ and $\kappa'=\kappa_{\mathrm{w}}+\kappa_0>0$, one has $M^{\alpha_0}\ge1$, $\mathrm{wp}_s\ge1$ and $\mathrm{wp}_X\ge1$.
By paragraph~1.1\,(ii), the cubes of $\pi_k$ that are contained in $X_{\mathrm{wk}}$ are exactly the cubes of $X_{\mathrm{wk}}$.
Therefore, for $s\in\mathrm{Pd}_k$, the product $s_{\mathrm{wk}}$ of Section~G1.9 has $|X_{\mathrm{wk}}|_M$ factors, each of modulus at most $M^{\alpha_0}$ (definition of $\mathrm{Pd}_k$, Section~I.1.1).
Together with paragraph~1.3 this gives
\begin{equation}
|s_{\mathrm{wk}}|\le M^{\alpha_0\cdot|X_{\mathrm{wk}}|_M}\le M^{27\alpha_0(n+1)}=\mathrm{wp}_s^{n+1}\qquad\left(\mathrm{wk}\in\mathrm{Wk}_n,\ s\in\mathrm{Pd}_k\right).
\tag{3-1.3}
\end{equation}
As $\alpha_0=\frac{1}{54}$, one has $\mathrm{wp}_s=M^{\frac{1}{2}}$.
By condition (M5), $M\ge\left(108\pi\cdot C'_{V}\cdot a'\right)^2$, that is, $M^{\frac{1}{2}}\ge108\pi\cdot C'_{V}\cdot a'$; hence
\[
M=M^{\frac{1}{2}}\cdot M^{\frac{1}{2}}\ge108\pi\cdot C'_{V}\cdot a'\cdot\mathrm{wp}_s.
\]
Condition (M5) also gives
\[
M\ge108\pi\cdot C'_{V}\cdot a'\cdot\exp(27\kappa')=108\pi\cdot C'_{V}\cdot a'\cdot\mathrm{wp}_X.
\]
With $\mathrm{qq}(\mathrm{wp})=\frac{54\pi\cdot C'_{V}\cdot a'\cdot\mathrm{wp}}{M}$ (Proposition~2.7\,(e)) it follows that
\begin{equation}
\mathrm{qq}(\mathrm{wp}_s)\le\frac{1}{2},\qquad\mathrm{qq}(\mathrm{wp}_X)\le\frac{1}{2}.
\tag{3-1.4}
\end{equation}

\subsubsection*{2. The three kinds of factors}

\paragraph{2.1. The factors.}
Let $\square$ be an $M$-cube, and let $z_{\square}$, $\tilde{W}$, $W^{-}$ be as in Section~G1.5.
Write $A|_{\square}$ for the restriction, to the bonds of $\tilde{W}$, of the $\ell\mathbb{Z}^3$-periodic bond function on $\mathbb{L}_k$ that $A$ defines.
The map $A\mapsto A|_{\square}$ is complex linear, and
\[
(-A)|_{\square}=-\left(A|_{\square}\right),\qquad\overline{A}|_{\square}=\overline{A|_{\square}}.
\]
Put $\delta_k:=2e_k^{\varepsilon}$.
By Section~G1.6, that is, by (G1) of \cite[Section~16.2.7]{QEDI}, the restriction $A|_{\square}$ lies in $\mathrm{Gloc}(\delta_k)$.
Since $e_k\le e_0$, condition (Z2.s1) gives $3L\cdot\delta_k=6L\cdot e_k^{\varepsilon}\le1$, and condition (Z2.s2) gives $2C'_{V}\cdot a'\cdot\delta_k=4C'_{V}\cdot a'\cdot e_k^{\varepsilon}\le1$.
Thus $\delta_k$ satisfies the two conditions $3L\cdot\delta\le1$ and $2C'_{V}\cdot a'\cdot\delta\le1$ on the number $\delta$ in the data of Proposition~2.4, and Proposition~2.5 applies at $A|_{\square}$ with $\delta=\delta_k$.
The factors are the following three elements of $\mathrm{Lop}$:
\begin{itemize}
\item $h_{\square}$, the multiplication by the function $h_{\square}(c)$ of Section~G1.7, which is $\ell\mathbb{Z}^3$-periodic and real, takes the same value at the $48$ sites over $c$, and is scalar in the spin index;
\item $\mathrm{Sloc}_{\square,y}(A)=\mathrm{Ecb}\left(\mathrm{Sloc}_{\square,y}(A|_{\square})\right)$ (Section~G1.6);
\item $\mathrm{Cm}_{\square,y}(A)=-\left[F_y(A),h_{\square}\right]$ (Section~G1.7).
\end{itemize}
The operator $\mathrm{Sloc}_{\square,y}(A)$ does not depend on the choice of the lift $z_{\square}$.
Indeed, if $z_{\square}$ is replaced by $z_{\square}+\ell n$, then $\tilde{W}$ is replaced by $\tilde{W}+\ell n$ and $\mathrm{Sloc}_{\square,y}(A|_{\square})$ by its translate (Proposition~2.5\,(g)), and $\mathrm{Ecb}$, formed with $\tilde{W}+\ell n$ and applied to this translate, gives the same operator on $\mathfrak{S}^{\mathrm{ch}}(\eta_k,\ell)$ (Proposition~2.6\,(a)).

\paragraph{2.2. The commutator, written out.}
Let $\psi$ be an $\ell\mathbb{Z}^3$-periodic function of $c$, let $f\in\mathfrak{S}^{\mathrm{ch}}(\eta_k,\ell)$ and $x=(c,g)\in\mathbb{L}^{\mathrm{f}}_k$, let $z:=[c]$ be the unit site with $c\in\Delta_z$, and let $z^{+}$ be the site of $L\mathbb{Z}^3$ with $z\in B(z^{+})$.
All sites are sites of the covering lattices $\mathbb{L}_k$, $\mathbb{Z}^3$, $L\mathbb{Z}^3$, and $A$ stands for its $\ell\mathbb{Z}^3$-periodic extension.
In the integrals below $x'=(c',g')$, and
\[
\int_Z\mathrm{d}x'=\eta_k^3\cdot\frac{1}{48}\sum_{c'\in Z}\sum_{g'\in O_{\mathrm{cub}}}.
\]
Let $k\ge1$.
By (I.2.3), (I.3.3), Section~I.3.3 and Section~I.3.5,
\begin{equation*}
\begin{split}
(F_y(A)f)(x)&=-\eta_k^{-1}\sum_{i=0}^{3}M_i(x)\cdot\left[U_A(x,x^{(i)})\cdot f(x^{(i)})-f(x)\right]+\mu\cdot f(x)\\
&\quad+\boldsymbol{\alpha}_{y}\int_{\Delta_z}\mathrm{d}x'\,\varphi_A(z,c)^{-1}\cdot\varphi_A(z,c')\cdot f(x')\\
&\quad+\boldsymbol{\beta}_{y}\cdot L^{-3}\int_{\mathrm{Bl}_{k+1}(z^{+})}\mathrm{d}x'\,\varphi^{+}_A(z^{+},c)^{-1}\cdot\varphi^{+}_A(z^{+},c')\cdot f(x').
\end{split}
\end{equation*}
In fact, the identities
\[
\left(Q_k(A)\cdot J^{*}f\right)(z)=\int_{\Delta_z}\mathrm{d}x'\,\varphi_A(z,c')\cdot f(x'),\qquad\left(J\cdot Q_k^{T}(-A)v\right)(x)=\varphi_A(z,c)^{-1}\cdot v(z)
\]
of Section~I.3.3 give the term with $J\cdot P_k(A)\cdot J^{*}$, and the identities
\begin{gather*}
\left(P(A)v\right)(z)=U_A(z^{+},z)^{-1}\cdot L^{-3}\sum_{z_1\in B(z^{+})}U_A(z^{+},z_1)\cdot v(z_1),\\
\varphi^{+}_A(z^{+},c')=U_A(z^{+},z_1)\cdot\varphi_A(z_1,c')\qquad\left(c'\in\Delta_{z_1}\right)
\end{gather*}
(Section~I.3.3 and (I.3.3)) give the term with $J\cdot P_{k+1}(A)\cdot J^{*}=J\cdot Q_k^{T}(-A)\cdot P(A)\cdot Q_k(A)\cdot J^{*}$.
The sites $x^{(0)}$, $x^{(1)}$, $x^{(2)}$ lie over $c$, the site $x^{(3)}$ lies over $c_3$, and $U_A(x,x^{(3)})=u_A(c,c_3)$ ((I.1.2), (I.1.6)).
Replacing $f$ by $\psi f$ and subtracting $\psi(c)\cdot(F_y(A)f)(x)$, one obtains for $k\ge1$
\begin{equation}
\begin{split}
\left([F_y(A),\psi]f\right)(x)&=-\eta_k^{-1}\cdot M_3(x)\cdot u_A(c,c_3)\cdot\left(\psi(c_3)-\psi(c)\right)\cdot f(x^{(3)})\\
&\quad+\boldsymbol{\alpha}_{y}\int_{\Delta_z}\mathrm{d}x'\,\varphi_A(z,c)^{-1}\cdot\varphi_A(z,c')\cdot\left(\psi(c')-\psi(c)\right)\cdot f(x')\\
&\quad+\boldsymbol{\beta}_{y}\cdot L^{-3}\int_{\mathrm{Bl}_{k+1}(z^{+})}\mathrm{d}x'\,\varphi^{+}_A(z^{+},c)^{-1}\cdot\varphi^{+}_A(z^{+},c')\cdot\left(\psi(c')-\psi(c)\right)\cdot f(x').
\end{split}
\tag{3-2.1}
\end{equation}
Now let $k=0$; then $\eta_0=1$ and $z=c$.
By (I.2.3), (I.3.1), Section~I.3.3 and Section~I.3.9, together with
\begin{gather*}
\left(J\cdot P(A)\cdot J^{*}f\right)(x)=\left(P(A)\cdot J^{*}f\right)(c),\\
\left(P(A)v\right)(c)=U_A(z^{+},c)^{-1}\cdot L^{-3}\sum_{c'\in B(z^{+})}U_A(z^{+},c')\cdot v(c'),\\
\left(J^{*}f\right)(c')=\frac{1}{48}\sum_{g'\in O_{\mathrm{cub}}}f(c',g'),
\end{gather*}
one has for $k=0$
\begin{equation*}
\begin{split}
(F_y(A)f)(x)&=-\sum_{i=0}^{3}M_i(x)\cdot\left[U_A(x,x^{(i)})\cdot f(x^{(i)})-f(x)\right]+\mu\cdot f(x)+i\gamma_5y\cdot f(x)\\
&\quad+b_{L}\cdot L^{-3}\sum_{c'\in B(z^{+})}U_A(z^{+},c)^{-1}\cdot U_A(z^{+},c')\cdot\frac{1}{48}\sum_{g'\in O_{\mathrm{cub}}}f(c',g').
\end{split}
\end{equation*}
Replacing $f$ by $\psi f$ and subtracting $\psi(c)\cdot(F_y(A)f)(x)$ in the same way, one obtains for $k=0$
\begin{equation}
\begin{split}
\left([F_y(A),\psi]f\right)(x)&=-M_3(x)\cdot u_A(c,c_3)\cdot\left(\psi(c_3)-\psi(c)\right)\cdot f(x^{(3)})\\
&\quad+b_{L}\cdot L^{-3}\sum_{c'\in B(z^{+})}U_A(z^{+},c)^{-1}\cdot U_A(z^{+},c')\cdot\left(\psi(c')-\psi(c)\right)\\
&\qquad\cdot\frac{1}{48}\sum_{g'\in O_{\mathrm{cub}}}f(c',g').
\end{split}
\tag{3-2.1$_0$}
\end{equation}
In the commutator the three hops inside the cluster, the diagonal term of $D_A$, the mass $\mu$ and, for $k=0$, the term $i\gamma_5y$ have cancelled.
The numbers
\[
u_A(c,c_3),\qquad\varphi_A(z,c)^{-1}\cdot\varphi_A(z,c'),\qquad\varphi^{+}_A(z^{+},c)^{-1}\cdot\varphi^{+}_A(z^{+},c'),\qquad U_A(z^{+},c)^{-1}\cdot U_A(z^{+},c')
\]
are exponentials of $i\,e_k\eta_k$ times linear forms in the numbers $A(b)$ (Section~I.1.1, (I.1.6), (I.3.3)).
The matrices $M_3(x)$, $\boldsymbol{\alpha}_{y}$, $\boldsymbol{\beta}_{y}$ and the number $b_{L}$ do not contain $A$.
The parameter $y$ enters (3-2.1) only through $\boldsymbol{\alpha}_{y}$ and $\boldsymbol{\beta}_{y}$, and it does not enter (3-2.1$_0$).
For real $t>0$ and real $y$ one has
\[
(t+i\gamma_5y)(t-i\gamma_5y)=t^2+y^2,
\]
because $\gamma_5^2=1$ ((I.c1)); hence $(t+i\gamma_5y)^{-1}=\frac{t-i\gamma_5y}{t^2+y^2}$.
This applies to $t=b_k$ and to $t=b_k+b_{L}$, which are both positive: $b_k\ge b_{-}=b\left(1-\frac{1}{L}\right)>0$ by Section~I.3.4, and $b_{L}=\frac{b}{L}>0$ by Section~I.0.5.
Consequently $\boldsymbol{\alpha}_{y}$ and $\boldsymbol{\beta}_{y}$ (Section~I.3.5) are continuous functions of $y\in\mathbb{R}$.

\paragraph{2.3. Dependence on $A$.}
The function $h_{\square}$ does not contain $A$.
By its definition, $\mathrm{Sloc}_{\square,y}(A)$ is a function of $A|_{\square}$.
The operator $\mathrm{Cm}_{\square,y}(A)$ reads $A$ on the bonds of $\tilde{W}$ only (Proposition~2.7\,(c)); thus it, too, is a function of $A|_{\square}$.
In terms of (3-2.1) and (3-2.1$_0$) with $\psi=h_{\square}$ this reads as follows.
By Proposition~2.7\,(c) the row at $x$ vanishes unless $c$ has a lift $\tilde{c}$ with $|\tilde{c}-z_{\square}|<\frac{3M}{4}+L$.
Such a row contains $A$ on the bond $(c,c_3)$ and on the bonds of the paths that enter $\varphi_A$, $\varphi^{+}_A$, $U_A$; these paths join sites of the block of side $L$ that contains $c$.
Combined with the last sentence of paragraph~1.1, this shows: if $A,\tilde{A}\in\tilde{G}_k$ coincide on the bonds of $T_k$ with both ends in $\square^{\sim1}$, then
\[
\mathrm{Sloc}_{\square,y}(A)=\mathrm{Sloc}_{\square,y}(\tilde{A})\qquad\text{and}\qquad\mathrm{Cm}_{\square,y}(A)=\mathrm{Cm}_{\square,y}(\tilde{A}).
\]

\paragraph{2.4. Analyticity in $A$, continuity in $y$.}
The connections on $T_k$ form a complex vector space of dimension $3\left(\frac{\ell}{\eta_k}\right)^3$, and $\tilde{G}_k$ is an open subset of this space ((G2) of \cite[Section~16.2.7]{QEDI}).
A function with values in $\mathrm{Lop}$ is called analytic if its entries with respect to a fixed basis are holomorphic functions of the complex variables on which it depends.
The map $A\mapsto A|_{\square}$ is linear and maps $\tilde{G}_k$ into $\mathrm{Gloc}(\delta_k)$ (paragraph~2.1).
The map $\mathrm{Sloc}_{\square,y}(\,\cdot\,)$ is analytic on $\mathrm{Gloc}(\delta_k)$ for every $y\in Y_k$ and continuous in $(y,\,\cdot\,)$ on $\mathbb{R}\times\mathrm{Gloc}(\delta_k)$ (Proposition~2.5\,(a)).
The map $O\mapsto\mathrm{Ecb}(O)$ is linear (Proposition~2.6\,(a)).
Hence $A\mapsto\mathrm{Sloc}_{\square,y}(A)$ is analytic on $\tilde{G}_k$ for every $y\in Y_k$, and $(y,A)\mapsto\mathrm{Sloc}_{\square,y}(A)$ is continuous on $\mathbb{R}\times\tilde{G}_k$.
By (3-0.2), together with (3-2.1) and (3-2.1$_0$) for $\psi=h_{\square}$, the entries of $\mathrm{Cm}_{\square,y}(A)$ are entire functions of $A$ for every $y\in Y_k$, and continuous functions of $(y,A)\in\mathbb{R}\times\tilde{G}_k$.

\paragraph{2.5. Gauge.}
Let $\boldsymbol{\lambda}$ be a real function of $c$ on $T_k$, and put $\mathrm{uf}:=\exp\left(i e_k\cdot\boldsymbol{\lambda}\right)$; it is an $\ell\mathbb{Z}^3$-periodic function of $c$ with $|\mathrm{uf}|=1$.
By (G2) of \cite[Section~16.2.7]{QEDI}, $A-d\boldsymbol{\lambda}$ lies in $\tilde{G}_k$.
Proposition~2.2\,(a) on $T_k$ with $\mathrm{Acl}=d\boldsymbol{\lambda}$ (the example given there) yields
\[
F_y(A-d\boldsymbol{\lambda})=\mathrm{uf}\cdot F_y(A)\cdot\mathrm{uf}^{-1}.
\]
The multiplication by $\mathrm{uf}$ commutes with $h_{\square}$.
By Proposition~2.6\,(b), $(A-d\boldsymbol{\lambda})|_{\square}=A|_{\square}-d\left(\boldsymbol{\lambda}|_{\tilde{W}}\right)$.
By Proposition~2.5\,(d), $\mathrm{Sloc}_{\square,y}\left(A|_{\square}-d\left(\boldsymbol{\lambda}|_{\tilde{W}}\right)\right)=\mathrm{uf}\cdot\mathrm{Sloc}_{\square,y}(A|_{\square})\cdot\mathrm{uf}^{-1}$ as operators over $\tilde{W}$.
By Proposition~2.6\,(b), $\mathrm{Ecb}\left(\mathrm{uf}\cdot O\cdot\mathrm{uf}^{-1}\right)=\mathrm{uf}\cdot\mathrm{Ecb}(O)\cdot\mathrm{uf}^{-1}$.
It follows that
\begin{equation}
\begin{split}
&h_{\square}=\mathrm{uf}\cdot h_{\square}\cdot\mathrm{uf}^{-1},\qquad\mathrm{Sloc}_{\square,y}(A-d\boldsymbol{\lambda})=\mathrm{uf}\cdot\mathrm{Sloc}_{\square,y}(A)\cdot\mathrm{uf}^{-1},\\
&\mathrm{Cm}_{\square,y}(A-d\boldsymbol{\lambda})=-\left[\mathrm{uf}\cdot F_y(A)\cdot\mathrm{uf}^{-1},h_{\square}\right]=\mathrm{uf}\cdot\mathrm{Cm}_{\square,y}(A)\cdot\mathrm{uf}^{-1}.
\end{split}
\tag{3-2.2}
\end{equation}

\paragraph{2.6. Lattice.}
Let $\tilde{g}=(S,a)\in\mathbb{G}_k$ and $r:=r(S)$, with $a$ taken in $M\mathbb{Z}^3$ (Section~I.4.3).
Let $y'=y$ if $S\in\mathrm{Sp}_3$ or if $y$ is neutral (Section~G1.9), and $y'=-y$ otherwise.
The map $c\mapsto\tilde{g}c=rc+a$, with $r$ diagonal with entries $\pm1$, is an isometry for the sup norm; it maps $\mathbb{L}_k$ and $M\mathbb{Z}^3$ onto themselves, and classes modulo $\ell\mathbb{Z}^3$ onto classes modulo $\ell\mathbb{Z}^3$, since $r\ell\mathbb{Z}^3=\ell\mathbb{Z}^3$.
Therefore it acts on $T_k$ and preserves $d$.
It maps the $M$-cube centred at $z$ onto the $M$-cube centred at $\tilde{g}z$, and cubes with a common face onto cubes with a common face; by (3-1.1), it maps $\square^{\sim1}$ onto $(\tilde{g}\square)^{\sim1}$.
By (G2) of \cite[Section~16.2.7]{QEDI}, $\tilde{g}A$ lies in $\tilde{G}_k$.
The map $\tilde{g}$ is a bijection of $\mathfrak{S}^{\mathrm{ch}}(\eta_k,\ell)$, given by $(\tilde{g}f)(x)=S\cdot f(\tilde{g}^{-1}x)$ with the action (I.4.2) on the sites, and $\pi(\tilde{g}^{-1}x)=\tilde{g}^{-1}c$ ((I.a1)).
For $f\in\mathfrak{S}^{\mathrm{ch}}(\eta_k,\ell)$, Section~G1.7\,(6.1) at the point $\tilde{g}^{-1}c$ gives
\[
\left(\tilde{g}(h_{\square}f)\right)(x)=S\cdot h_{\square}(\tilde{g}^{-1}c)\cdot f(\tilde{g}^{-1}x)=h_{\tilde{g}\square}(c)\cdot(\tilde{g}f)(x).
\]
The element $\tilde{g}$ induces the morphism $c\mapsto rc+a$ of $T_k$ onto itself, with $\mathrm{mS}=\tilde{g}$ (Section~G1.2); so Proposition~2.2\,(b) on $T_k$ gives
\[
F_{y'}(\tilde{g}A)\cdot\tilde{g}=\tilde{g}\cdot F_y(A).
\]
The set $\tilde{g}\tilde{W}:=r\tilde{W}+a$ is a lift of $(\tilde{g}\square)^{\sim1}$ with the centre $\tilde{g}z_{\square}\in M\mathbb{Z}^3$; by paragraph~2.1 the operator $\mathrm{Sloc}_{\tilde{g}\square,y'}(\tilde{g}A)\in\mathrm{Lop}$ may be formed with this lift.
One has $(\tilde{g}A)|_{\tilde{g}\square}=\tilde{g}\left(A|_{\square}\right)$ (Proposition~2.6\,(c)).
As operators over $\tilde{g}\tilde{W}$, $\mathrm{Sloc}_{\tilde{g}\square,y'}\left(\tilde{g}(A|_{\square})\right)=\tilde{g}\cdot\mathrm{Sloc}_{\square,y}(A|_{\square})\cdot\tilde{g}^{-1}$ (Proposition~2.5\,(e)).
Finally, $E_{\tilde{g}\square}\left(\tilde{g}\cdot O\cdot\tilde{g}^{-1}\right)=\tilde{g}\cdot\mathrm{Ecb}(O)\cdot\tilde{g}^{-1}$ (Proposition~2.6\,(c)).
Hence
\begin{equation}
\begin{split}
h_{\tilde{g}\square}\cdot\tilde{g}&=\tilde{g}\cdot h_{\square},\qquad\mathrm{Sloc}_{\tilde{g}\square,y'}(\tilde{g}A)\cdot\tilde{g}=\tilde{g}\cdot\mathrm{Sloc}_{\square,y}(A),\\
\mathrm{Cm}_{\tilde{g}\square,y'}(\tilde{g}A)\cdot\tilde{g}&=-F_{y'}(\tilde{g}A)\cdot h_{\tilde{g}\square}\cdot\tilde{g}+h_{\tilde{g}\square}\cdot F_{y'}(\tilde{g}A)\cdot\tilde{g}\\
&=-\tilde{g}\cdot F_y(A)\cdot h_{\square}+\tilde{g}\cdot h_{\square}\cdot F_y(A)=\tilde{g}\cdot\mathrm{Cm}_{\square,y}(A).
\end{split}
\tag{3-2.3}
\end{equation}
If $a$ is replaced by $a+\ell n$, the operator $\tilde{g}$ on $\mathfrak{S}^{\mathrm{ch}}(\eta_k,\ell)$ is multiplied by $(-1)^{n_1+n_2+n_3}$, on both sides of each of these relations.

\paragraph{2.7. Transposition.}
The connection $-A$ lies in $\tilde{G}_k$ ((G2) of \cite[Section~16.2.7]{QEDI}), and $-y\in Y_k$, with $-\infty:=\infty$ (Section~G1.9).
One has $h_{\square}^T=h_{\square}$ (paragraph~0.4).
By Proposition~2.2\,(c) on $T_k$,
\[
F_y(A)^T=\mathsf{C}^{-1}\cdot F_{-y}(-A)\cdot\mathsf{C}.
\]
For $O\in\mathrm{Lop}$,
\[
\left[O,h_{\square}\right]^T=h_{\square}\cdot O^T-O^T\cdot h_{\square}=-\left[O^T,h_{\square}\right].
\]
The operator $\mathrm{Sloc}_{\square,y}(A|_{\square})$ takes values in the functions on $\tilde{W}\times O_{\mathrm{cub}}$ (Section~G1.5).
Therefore Proposition~2.6\,(d), Proposition~2.5\,(f), the identity $(-A)|_{\square}=-\left(A|_{\square}\right)$ and Proposition~2.6\,(b) with $\mathrm{Mx}=\mathsf{C}^{-1}$ give
\begin{equation*}
\begin{split}
\mathrm{Ecb}\left(\mathrm{Sloc}_{\square,y}(A|_{\square})\right)^T&=\mathrm{Ecb}\left(\mathrm{Sloc}_{\square,y}(A|_{\square})^T\right)=\mathrm{Ecb}\left(\mathsf{C}^{-1}\cdot\mathrm{Sloc}_{\square,-y}\left((-A)|_{\square}\right)\cdot\mathsf{C}\right)\\
&=\mathsf{C}^{-1}\cdot\mathrm{Ecb}\left(\mathrm{Sloc}_{\square,-y}\left((-A)|_{\square}\right)\right)\cdot\mathsf{C}.
\end{split}
\end{equation*}
Since $\mathsf{C}$ commutes with $h_{\square}$, it follows that
\begin{equation}
\begin{split}
&h_{\square}^T=\mathsf{C}^{-1}\cdot h_{\square}\cdot\mathsf{C},\qquad\mathrm{Sloc}_{\square,y}(A)^T=\mathsf{C}^{-1}\cdot\mathrm{Sloc}_{\square,-y}(-A)\cdot\mathsf{C},\\
&\mathrm{Cm}_{\square,y}(A)^T=\left[F_y(A)^T,h_{\square}\right]=\mathsf{C}^{-1}\cdot\left[F_{-y}(-A),h_{\square}\right]\cdot\mathsf{C}=-\mathsf{C}^{-1}\cdot\mathrm{Cm}_{\square,-y}(-A)\cdot\mathsf{C}.
\end{split}
\tag{3-2.4}
\end{equation}

\paragraph{2.8. Conjugation.}
The connections $\overline{A}$ and $-A$ lie in $\tilde{G}_k$ ((G2) of \cite[Section~16.2.7]{QEDI}).
The function $h_{\square}$ is real.
By Proposition~2.2\,(d) on $T_k$,
\[
\overline{F_y(\overline{A})}=\varkappa^{-1}\cdot F_y(-A)\cdot\varkappa.
\]
The identity $\overline{A}|_{\square}=\overline{A|_{\square}}$, Proposition~2.6\,(e), Proposition~2.5\,(f) and Proposition~2.6\,(b) with $\mathrm{Mx}=\varkappa^{-1}$ give
\begin{equation*}
\begin{split}
\overline{\mathrm{Ecb}\left(\mathrm{Sloc}_{\square,y}\left(\overline{A|_{\square}}\right)\right)}&=\mathrm{Ecb}\left(\overline{\mathrm{Sloc}_{\square,y}\left(\overline{A|_{\square}}\right)}\right)=\mathrm{Ecb}\left(\varkappa^{-1}\cdot\mathrm{Sloc}_{\square,y}\left((-A)|_{\square}\right)\cdot\varkappa\right)\\
&=\varkappa^{-1}\cdot\mathrm{Ecb}\left(\mathrm{Sloc}_{\square,y}\left((-A)|_{\square}\right)\right)\cdot\varkappa.
\end{split}
\end{equation*}
Since $\varkappa$ commutes with $h_{\square}$, it follows that
\begin{equation}
\begin{split}
&\overline{h_{\square}}=\varkappa^{-1}\cdot h_{\square}\cdot\varkappa,\qquad\overline{\mathrm{Sloc}_{\square,y}(\overline{A})}=\varkappa^{-1}\cdot\mathrm{Sloc}_{\square,y}(-A)\cdot\varkappa,\\
&\overline{\mathrm{Cm}_{\square,y}(\overline{A})}=-\left[\overline{F_y(\overline{A})},h_{\square}\right]=\varkappa^{-1}\cdot\mathrm{Cm}_{\square,y}(-A)\cdot\varkappa.
\end{split}
\tag{3-2.5}
\end{equation}

\paragraph{2.9. The words.}
Let $\mathrm{wk}=(\square_0,\ldots,\square_n)\in\mathrm{Wk}_n$, and let $h_i$, $\mathrm{Sloc}_i$, $\mathrm{Cm}_i$ be as in Section~G1.8.
Written without brackets, the two products of Section~G1.8 are
\begin{align}
S_{y,\mathrm{wk}}(A)&=h_0\cdot\mathrm{Sloc}_0\cdot h_0\cdot\mathrm{Cm}_1\cdot\mathrm{Sloc}_1\cdot h_1\cdot\mathrm{Cm}_2\cdot\mathrm{Sloc}_2\cdot h_2\cdots\mathrm{Cm}_n\cdot\mathrm{Sloc}_n\cdot h_n,\tag{3-2.6}\\
S'_{y,\mathrm{wk}}(A)&=(-1)^n\,h_n\cdot\mathrm{Sloc}_n\cdot\mathrm{Cm}_n\cdot h_{n-1}\cdot\mathrm{Sloc}_{n-1}\cdot\mathrm{Cm}_{n-1}\cdots\notag\\
&\qquad\cdots h_1\cdot\mathrm{Sloc}_1\cdot\mathrm{Cm}_1\cdot h_0\cdot\mathrm{Sloc}_0\cdot h_0.\tag{3-2.7}
\end{align}
For $n=0$ both are equal to $h_0\cdot\mathrm{Sloc}_0\cdot h_0$.
The product (3-2.6) has $h_0$ at its left end and $h_n$ at its right end; the product (3-2.7) has $h_n$ at its left end and $h_0$ at its right end.
Finally, $S^{\mathrm{sy}}_{y,\mathrm{wk}}(A)=\frac{1}{2}\left(S_{y,\mathrm{wk}}(A)+S'_{y,\mathrm{wk}}(A)\right)$ (Section~G1.9).

\subsubsection*{3. Proof of the clauses}

Throughout this section $\sum_{\mathrm{wk}}$ is an abbreviation for $\sum_{n\ge0}\sum_{\mathrm{wk}\in\mathrm{Wk}_n}$.
By paragraph~1.3 each of the inner sums has finitely many terms.

\begin{proof}[Proof of Proposition~3\,(a)]
\emph{Three inequalities at a fixed length.}
Let $\mathrm{wp}$ be either of the two numbers $\mathrm{wp}_s$, $\mathrm{wp}_X$.
Paragraph~1.4 gives $\mathrm{wp}\ge1$ and $\mathrm{qq}(\mathrm{wp})\le\frac{1}{2}$, whence $\mathrm{qq}(\mathrm{wp})^n\le2^{-n}$.
Consequently Proposition~2.7\,(e) yields, for every $n\ge0$, every unit site $z$, and for $S^{\#}=S$ as well as for $S^{\#}=S'$,
\begin{align}
\sum_{z'}\exp\left(\gamma^{F}\cdot d(z,z')\right)\sum_{\mathrm{wk}\in\mathrm{Wk}_n}\mathrm{wp}^{n+1}\left\|\mathbf{1}_zS^{\#}_{y,\mathrm{wk}}(A)\mathbf{1}_{z'}\right\|_{\mathrm{su}}&\le16a'\cdot\mathrm{wp}\cdot w_k(y)\cdot2^{-n},\tag{3-3.1}\\
\sum_{z'}\exp\left(\gamma^{F}\cdot d(z,z')\right)\sum_{\mathrm{wk}\in\mathrm{Wk}_n}\mathrm{wp}^{n+1}\left\|\mathbf{1}_z\delta_{\alpha,A}S^{\#}_{y,\mathrm{wk}}(A)\mathbf{1}_{z'}\right\|_{\mathrm{su}}&\le152a'\cdot\mathrm{wp}\cdot2^{-n},\tag{3-3.2}
\end{align}
and, at $k=0$, Proposition~2.7\,(f) yields
\begin{equation}
\begin{split}
&\sum_{z'}\exp\left(\gamma^{F}\cdot d(z,z')\right)\sum_{\mathrm{wk}\in\mathrm{Wk}_n}\mathrm{wp}^{n+1}\sup\left\{\left|S^{\#}_{y,\mathrm{wk}}(A)(x,x')\right|_2:\ x\text{ over }z,\ x'\text{ over }z'\right\}\\
&\qquad\le16a'\cdot\mathrm{wp}\cdot w_0(y)\cdot2^{-n}.
\end{split}
\tag{3-3.3}
\end{equation}
For fixed unit sites $z$, $z'$ each of the three maps
\[
O\mapsto\left\|\mathbf{1}_zO\mathbf{1}_{z'}\right\|_{\mathrm{su}},\qquad
O\mapsto\left\|\mathbf{1}_z\delta_{\alpha,A}O\mathbf{1}_{z'}\right\|_{\mathrm{su}},\qquad
O\mapsto\sup\left\{|O(x,x')|_2:\ x\text{ over }z,\ x'\text{ over }z'\right\}
\]
is a seminorm on $\mathrm{Lop}$.
Its value at $S^{\mathrm{sy}}_{y,\mathrm{wk}}(A)=\frac{1}{2}\left(S_{y,\mathrm{wk}}(A)+S'_{y,\mathrm{wk}}(A)\right)$ is therefore at most the mean of its values at $S_{y,\mathrm{wk}}(A)$ and at $S'_{y,\mathrm{wk}}(A)$.
It follows that (3-3.1), (3-3.2) and (3-3.3) hold for $S^{\#}=S^{\mathrm{sy}}$ too.
At $k=0$ one has $\mathrm{Lop}^{\mathrm{sp}}=\{0\}$, and (3-3.2) is void.

\emph{Convergence.}
Let $s\in\mathrm{Pd}_k$, and let $\mathrm{nlb}\ge0$ be an integer.
All terms below being non-negative, the summation over $z$ may be performed last; so (3-1.3), (3-0.1), together with (3-3.1), (3-3.2) for $\mathrm{wp}=\mathrm{wp}_s$, give
\begin{align}
&\sum_{n\ge\mathrm{nlb}}\sum_{\mathrm{wk}\in\mathrm{Wk}_n}N\left(s_{\mathrm{wk}}\cdot S^{\#}_{y,\mathrm{wk}}(A)\right)\notag\\
&\qquad\le\sum_z\sum_{n\ge\mathrm{nlb}}\sum_{z'}\exp\left(\gamma^{F}\cdot d(z,z')\right)\sum_{\mathrm{wk}\in\mathrm{Wk}_n}\mathrm{wp}_s^{n+1}\left\|\mathbf{1}_zS^{\#}_{y,\mathrm{wk}}(A)\mathbf{1}_{z'}\right\|_{\mathrm{su}}\notag\\
&\qquad\le32n_k\cdot a'\cdot\mathrm{wp}_s\cdot w_k(y)\cdot2^{-\mathrm{nlb}},\tag{3-3.4}\\
&\sum_{n\ge\mathrm{nlb}}\sum_{\mathrm{wk}\in\mathrm{Wk}_n}N\left(s_{\mathrm{wk}}\cdot\delta_{\alpha,A}S^{\#}_{y,\mathrm{wk}}(A)\right)\notag\\
&\qquad\le\sum_z\sum_{n\ge\mathrm{nlb}}\sum_{z'}\exp\left(\gamma^{F}\cdot d(z,z')\right)\sum_{\mathrm{wk}\in\mathrm{Wk}_n}\mathrm{wp}_s^{n+1}\left\|\mathbf{1}_z\delta_{\alpha,A}S^{\#}_{y,\mathrm{wk}}(A)\mathbf{1}_{z'}\right\|_{\mathrm{su}}\notag\\
&\qquad\le304n_k\cdot a'\cdot\mathrm{wp}_s\cdot2^{-\mathrm{nlb}}.\tag{3-3.5}
\end{align}
Here $w_k(y)\le1$.
For $\mathrm{nlb}=0$ the left sides are finite, so that, by paragraph~0.2, the series $\sum_{\mathrm{wk}}s_{\mathrm{wk}}\cdot S^{\#}_{y,\mathrm{wk}}(A)$ converges absolutely in $N$.
The right side of (3-3.4) involves neither $s$ nor $A$, and, once $w_k(y)\le1$ is used, not $y$ either; it tends to $0$ as $\mathrm{nlb}\to\infty$.
Hence the convergence is uniform on $\mathrm{Pd}_k\times\tilde{G}_k$, and uniform in $y\in Y_k$ as well.
For $S^{\#}=S$ and for $S^{\#}=S'$ the sum of the series is $S^{\#}_y(s,A)$, by Section~G1.9; for $S^{\#}=S^{\mathrm{sy}}$ it is $\frac{1}{2}\left(S_y(s,A)+S'_y(s,A)\right)=S^{\mathrm{sy}}_y(s,A)$.
The family $s=1$, that is, $s_{\mathrm{pc}}=1$ for every $\mathrm{pc}$, belongs to $\mathrm{Pd}_k$ because $M^{\alpha_0}\ge1$, and all its coefficients are $s_{\mathrm{wk}}=1$.
Every term of the series that defines $S^{\#}_y(X,A)$ is a term of the series at $s=1$.
Therefore (3-3.4) holds for this series as well; it converges absolutely in $N$, uniformly on $\tilde{G}_k$ and in $y$, and for $S^{\#}=S^{\mathrm{sy}}$ its sum is $\frac{1}{2}\left(S_y(X,A)+S'_y(X,A)\right)=S^{\mathrm{sy}}_y(X,A)$.
Finally, (3-3.5) and paragraph~0.2 show that
\[
\delta_{\alpha,A}S^{\#}_y(s,A)=\sum_{\mathrm{wk}}s_{\mathrm{wk}}\cdot\delta_{\alpha,A}S^{\#}_{y,\mathrm{wk}}(A),
\]
the series converging absolutely in $N$.

\emph{The first and the second bound.}
Paragraph~0.2 gives
\[
\mathbf{1}_zS^{\#}_y(s,A)\mathbf{1}_{z'}=\sum_{\mathrm{wk}}s_{\mathrm{wk}}\cdot\mathbf{1}_zS^{\#}_{y,\mathrm{wk}}(A)\mathbf{1}_{z'}.
\]
Using (3-1.3) and (3-3.1) with $\mathrm{wp}=\mathrm{wp}_s$, one finds for every unit site $z$
\begin{align*}
\sum_{z'}\exp\left(\gamma^{F}\cdot d(z,z')\right)\left\|\mathbf{1}_zS^{\#}_y(s,A)\mathbf{1}_{z'}\right\|_{\mathrm{su}}
&\le\sum_{n\ge0}\sum_{z'}\exp\left(\gamma^{F}\cdot d(z,z')\right)\sum_{\mathrm{wk}\in\mathrm{Wk}_n}\mathrm{wp}_s^{n+1}\left\|\mathbf{1}_zS^{\#}_{y,\mathrm{wk}}(A)\mathbf{1}_{z'}\right\|_{\mathrm{su}}\\
&\le\sum_{n\ge0}16a'\cdot\mathrm{wp}_s\cdot w_k(y)\cdot2^{-n}=32a'\cdot\mathrm{wp}_s\cdot w_k(y).
\end{align*}
The supremum over $z$ of the left side is $N\left(S^{\#}_y(s,A)\right)$; since $\mathrm{wp}_s=M^{27\alpha_0}$, this gives $N\left(S^{\#}_y(s,A)\right)\le32a'\cdot M^{27\alpha_0}\cdot w_k(y)$.
In the same manner, the identity at the end of the preceding paragraph and paragraph~0.2 give
\[
\mathbf{1}_z\delta_{\alpha,A}S^{\#}_y(s,A)\mathbf{1}_{z'}=\sum_{\mathrm{wk}}s_{\mathrm{wk}}\cdot\mathbf{1}_z\delta_{\alpha,A}S^{\#}_{y,\mathrm{wk}}(A)\mathbf{1}_{z'},
\]
and (3-1.3) together with (3-3.2) at $\mathrm{wp}=\mathrm{wp}_s$ yields, for every $z$,
\begin{align*}
\sum_{z'}\exp\left(\gamma^{F}\cdot d(z,z')\right)\left\|\mathbf{1}_z\delta_{\alpha,A}S^{\#}_y(s,A)\mathbf{1}_{z'}\right\|_{\mathrm{su}}
&\le\sum_{n\ge0}\sum_{z'}\exp\left(\gamma^{F}\cdot d(z,z')\right)\sum_{\mathrm{wk}\in\mathrm{Wk}_n}\mathrm{wp}_s^{n+1}\left\|\mathbf{1}_z\delta_{\alpha,A}S^{\#}_{y,\mathrm{wk}}(A)\mathbf{1}_{z'}\right\|_{\mathrm{su}}\\
&\le\sum_{n\ge0}152a'\cdot\mathrm{wp}_s\cdot2^{-n}=304a'\cdot\mathrm{wp}_s.
\end{align*}
Taking the supremum over $z$ one obtains $N\left(\delta_{\alpha,A}S^{\#}_y(s,A)\right)\le304a'\cdot M^{27\alpha_0}$.

\emph{The third bound.}
Let $X\in D_k$.
If $\mathrm{wk}\in\mathrm{Wk}_n$ and $X_{\mathrm{wk}}=X$, then (3-1.2) gives $d_M(X)\le27(n+1)$, which may be written
\begin{equation}
1\le\exp\left(-\kappa'\cdot d_M(X)\right)\cdot\exp\left(27\kappa'\cdot(n+1)\right)=\exp\left(-\kappa'\cdot d_M(X)\right)\cdot\mathrm{wp}_X^{n+1}.\tag{3-3.6}
\end{equation}
By paragraph~0.2,
\[
\mathbf{1}_zS^{\#}_y(X,A)\mathbf{1}_{z'}=\sum_{\mathrm{wk}:\,X_{\mathrm{wk}}=X}\mathbf{1}_zS^{\#}_{y,\mathrm{wk}}(A)\mathbf{1}_{z'}.
\]
Let $z$, $z'$ be unit sites.
Apply the triangle inequality and (3-3.6); enlarge the sum over the walks with $X_{\mathrm{wk}}=X$ to the sum over all walks; and, for each $n$, bound the term belonging to $z'$ by the whole sum over the unit sites in (3-3.1) with $\mathrm{wp}=\mathrm{wp}_X$, all of whose terms are non-negative.
The result is
\begin{align*}
\exp\left(\gamma^{F}\cdot d(z,z')+\kappa'\cdot d_M(X)\right)\left\|\mathbf{1}_zS^{\#}_y(X,A)\mathbf{1}_{z'}\right\|_{\mathrm{su}}
&\le\sum_{n\ge0}\sum_{\mathrm{wk}\in\mathrm{Wk}_n}\mathrm{wp}_X^{n+1}\exp\left(\gamma^{F}\cdot d(z,z')\right)\left\|\mathbf{1}_zS^{\#}_{y,\mathrm{wk}}(A)\mathbf{1}_{z'}\right\|_{\mathrm{su}}\\
&\le\sum_{n\ge0}16a'\cdot\mathrm{wp}_X\cdot w_k(y)\cdot2^{-n}=32a'\cdot\exp(27\kappa')\cdot w_k(y).
\end{align*}
If no walk has $X_{\mathrm{wk}}=X$, then $S^{\#}_y(X,A)=0$.

\emph{The kernels at level~$0$.}
Let $k=0$, and let $x=(c,g)$, $x'=(c',g')$ be fermion sites.
Then $w_0(y)=\left(1+|y|\right)^{-1}$, the unit sites are the photon sites, and $d(x,x')=d(c,c')$ by (I.POS).
By paragraphs~0.2 and~0.3 the kernel of the sum of a series is the sum of the kernels of its terms.
Using (3-1.3) and (3-3.3) with $\mathrm{wp}=\mathrm{wp}_s$ and $z=c$, and retaining only the term $z'=c'$, one obtains
\begin{align*}
\exp\left(\gamma^{F}\cdot d(c,c')\right)\left|S^{\#}_y(s,A)(x,x')\right|_2
&\le\sum_{n\ge0}\exp\left(\gamma^{F}\cdot d(c,c')\right)\sum_{\mathrm{wk}\in\mathrm{Wk}_n}\mathrm{wp}_s^{n+1}\left|S^{\#}_{y,\mathrm{wk}}(A)(x,x')\right|_2\\
&\le\sum_{n\ge0}16a'\cdot\mathrm{wp}_s\cdot w_0(y)\cdot2^{-n}=32a'\cdot M^{27\alpha_0}\cdot\left(1+|y|\right)^{-1}.
\end{align*}
For the polymer family one starts from $S^{\#}_y(X,A)(x,x')=\sum_{\mathrm{wk}:\,X_{\mathrm{wk}}=X}S^{\#}_{y,\mathrm{wk}}(A)(x,x')$, applies (3-3.6), enlarges the sum over the walks with $X_{\mathrm{wk}}=X$ to the sum over all walks, and uses (3-3.3) with $\mathrm{wp}=\mathrm{wp}_X$ and $z=c$, retaining only the term $z'=c'$:
\begin{align*}
\exp\left(\gamma^{F}\cdot d(c,c')+\kappa'\cdot d_M(X)\right)\left|S^{\#}_y(X,A)(x,x')\right|_2
&\le\sum_{n\ge0}\exp\left(\gamma^{F}\cdot d(c,c')\right)\sum_{\mathrm{wk}\in\mathrm{Wk}_n}\mathrm{wp}_X^{n+1}\left|S^{\#}_{y,\mathrm{wk}}(A)(x,x')\right|_2\\
&\le\sum_{n\ge0}16a'\cdot\mathrm{wp}_X\cdot w_0(y)\cdot2^{-n}=32a'\cdot\exp(27\kappa')\cdot\left(1+|y|\right)^{-1}.
\end{align*}
These two inequalities are the kernel bounds at $k=0$ stated in (a).
\end{proof}

\begin{proof}[Proof of Proposition~3\,(b)]
Proposition~2.8\,(a) gives
\[
S_y(A)=\sum_{\mathrm{wk}}S_{y,\mathrm{wk}}(A)=\sum_{\mathrm{wk}}S'_{y,\mathrm{wk}}(A).
\]
At $s=1$ every coefficient $s_{\mathrm{wk}}$ equals $1$.
Hence $S_y(1,A)=S'_y(1,A)=S_y(A)$, and $S^{\mathrm{sy}}_y(1,A)=\frac{1}{2}\left(S_y(A)+S_y(A)\right)=S_y(A)$.
To each walk $\mathrm{wk}$ belongs exactly one set $X_{\mathrm{wk}}$, and $X_{\mathrm{wk}}\in D_k$ by paragraph~1.2; the set $D_k$ is finite by paragraph~1.3.
Grouping the terms of the absolutely convergent series $\sum_{\mathrm{wk}}S^{\#}_{y,\mathrm{wk}}(A)$ according to the value of $X_{\mathrm{wk}}$, as paragraph~0.2 permits, one obtains
\[
\sum_{X\in D_k}S^{\#}_y(X,A)=\sum_{X\in D_k}\ \sum_{\mathrm{wk}:\,X_{\mathrm{wk}}=X}S^{\#}_{y,\mathrm{wk}}(A)=\sum_{\mathrm{wk}}S^{\#}_{y,\mathrm{wk}}(A)=S^{\#}_y(1,A)=S_y(A).\qedhere
\]
\end{proof}

\begin{proof}[Proof of Proposition~3\,(c)]
According to paragraph~2.4 and to (3-2.6), (3-2.7), each of $S_{y,\mathrm{wk}}(A)$ and $S'_{y,\mathrm{wk}}(A)$ is a finite product, and $S^{\mathrm{sy}}_{y,\mathrm{wk}}(A)$ is the mean of two such products, of factors that are analytic in $A$ on $\tilde{G}_k$ for every $y\in Y_k$ and continuous in $(y,A)$ on $\mathbb{R}\times\tilde{G}_k$.
These operators therefore have the same two properties.
The coefficient $s_{\mathrm{wk}}$ is a monomial in the variables $s_{\mathrm{pc}}$.
For $\mathrm{nlb}\ge1$ put
\begin{align*}
\mathrm{Pt}_{\mathrm{nlb}}(y,s,A)&:=\sum_{n<\mathrm{nlb}}\sum_{\mathrm{wk}\in\mathrm{Wk}_n}s_{\mathrm{wk}}\cdot S^{\#}_{y,\mathrm{wk}}(A),\\
\mathrm{Pt}_{\mathrm{nlb}}(y,X,A)&:=\sum_{n<\mathrm{nlb}}\sum_{\mathrm{wk}\in\mathrm{Wk}_n,\ X_{\mathrm{wk}}=X}S^{\#}_{y,\mathrm{wk}}(A);
\end{align*}
by paragraph~1.3 these are finite sums.
Hence, for every $y\in Y_k$, the function $\mathrm{Pt}_{\mathrm{nlb}}(y,\cdot,\cdot)$ is analytic on $\mathbb{C}^{\pi_k}\times\tilde{G}_k$ and $\mathrm{Pt}_{\mathrm{nlb}}(y,X,\cdot)$ is analytic on $\tilde{G}_k$; moreover the first is a continuous function of $(y,s,A)\in\mathbb{R}\times\mathbb{C}^{\pi_k}\times\tilde{G}_k$, and the second a continuous function of $(y,A)\in\mathbb{R}\times\tilde{G}_k$.
By (3-3.4),
\[
N\left(S^{\#}_y(s,A)-\mathrm{Pt}_{\mathrm{nlb}}(y,s,A)\right)\le32n_k\cdot a'\cdot\mathrm{wp}_s\cdot2^{-\mathrm{nlb}}\qquad\text{for all }(y,s,A)\in Y_k\times\mathrm{Pd}_k\times\tilde{G}_k,
\]
and by (3-3.4) for the series with $X$ (proof of (a))
\[
N\left(S^{\#}_y(X,A)-\mathrm{Pt}_{\mathrm{nlb}}(y,X,A)\right)\le32n_k\cdot a'\cdot\mathrm{wp}_s\cdot2^{-\mathrm{nlb}}\qquad\text{for all }(y,A)\in Y_k\times\tilde{G}_k.
\]
Thus the partial sums converge in the norm $N$, uniformly on $Y_k\times\mathrm{Pd}_k\times\tilde{G}_k$, resp. on $Y_k\times\tilde{G}_k$; the dimension being finite, they converge uniformly in every entry.

A uniform limit of continuous functions is continuous.
Therefore $(s,A)\mapsto S^{\#}_y(s,A)$ is continuous on $\mathrm{Pd}_k\times\tilde{G}_k$ for every $y\in Y_k$, and the maps $(y,s,A)\mapsto S^{\#}_y(s,A)$ and $(y,A)\mapsto S^{\#}_y(X,A)$ are continuous on $\mathbb{R}\times\mathrm{Pd}_k\times\tilde{G}_k$ and on $\mathbb{R}\times\tilde{G}_k$, respectively.
In particular all of them depend continuously on $y\in\mathbb{R}$.

On an open subset of $\mathbb{C}^d$ a uniform limit of holomorphic functions is holomorphic.
Indeed, the limit is continuous; it is holomorphic in each variable separately \cite[Theorem~10.28]{Rudin1987}; and a continuous function that is holomorphic in each variable separately is holomorphic, by Osgood's lemma \cite[Chapter~I, Section~A]{GunningRossi1965}.
This fact is applied on three open sets.
\begin{itemize}
\item[(1)] Since $\tilde{G}_k$ is open (paragraph~2.4), the interior of $\mathrm{Pd}_k\times\tilde{G}_k$ is
\[
\left\{s:\ |s_{\mathrm{pc}}|<M^{\alpha_0}\text{ for all }\mathrm{pc}\in\pi_k\right\}\times\tilde{G}_k,
\]
and on this set $S^{\#}_y(s,A)$ is analytic in $(s,A)$.
\item[(2)] Fix $s\in\mathrm{Pd}_k$, whether or not $|s_{\mathrm{pc}}|=M^{\alpha_0}$ for some $\mathrm{pc}$.
The functions $A\mapsto\mathrm{Pt}_{\mathrm{nlb}}(y,s,A)$ are analytic on $\tilde{G}_k$ and converge uniformly there; hence $A\mapsto S^{\#}_y(s,A)$ is analytic on $\tilde{G}_k$.
\item[(3)] The functions $A\mapsto\mathrm{Pt}_{\mathrm{nlb}}(y,X,A)$ are analytic on $\tilde{G}_k$ and converge uniformly there; hence $A\mapsto S^{\#}_y(X,A)$ is analytic on $\tilde{G}_k$.
\end{itemize}
\end{proof}

\begin{proof}[Proof of Proposition~3\,(d)]
Let $\mathrm{wk}=(\square_0,\dots,\square_n)\in\mathrm{Wk}_n$, and write $z_i:=z_{\square_i}$.

\emph{The set $X_{\mathrm{wk}}$.}
That $X_{\mathrm{wk}}$ is a polymer was shown in paragraph~1.2.
The bound $|X_{\mathrm{wk}}|_M\le27(n+1)$ is in paragraph~1.3, and $d_M(X_{\mathrm{wk}})\le27(n+1)$ is (3-1.2).

\emph{Supports.}
Paragraph~2.9 and the product rule of paragraph~0.3 provide operators $O,O'\in\mathrm{Lop}$ such that, for $x=(c,g)$ and $x'=(c',g')$,
\begin{equation}
S_{y,\mathrm{wk}}(A)(x,x')=h_0(c)\cdot O(x,x')\cdot h_n(c'),\qquad S'_{y,\mathrm{wk}}(A)(x,x')=h_n(c)\cdot O'(x,x')\cdot h_0(c').\tag{3-3.7}
\end{equation}
By display (6.1) of Section~G1.7, $h_i(c)\ne0$ implies $d(c,z_i)<\frac{3M}{4}<\frac{3M}{2}$, hence $c\in\square_i^{\sim1}$ by (3-1.1); and $\square_i^{\sim1}\subset X_{\mathrm{wk}}$.
Both kernels in (3-3.7) therefore vanish unless $c$ and $c'$ belong to $X_{\mathrm{wk}}$, that is, unless $x$ and $x'$ lie over $X_{\mathrm{wk}}$.
The kernel of $S^{\mathrm{sy}}_{y,\mathrm{wk}}(A)$ is the mean of these two kernels (paragraph~0.3); so it, too, vanishes unless $x$ and $x'$ lie over $X_{\mathrm{wk}}$.

\emph{Dependence on $A$.}
Let $A,\tilde{A}\in\tilde{G}_k$ coincide on the bonds with both end points in $X_{\mathrm{wk}}$.
Then, for every $i$, they coincide on the bonds with both end points in $\square_i^{\sim1}$, and paragraph~2.3 shows that the factors $h_i$, $\mathrm{Sloc}_i$, $\mathrm{Cm}_i$ of (3-2.6), (3-2.7) are the same for $A$ and for $\tilde{A}$.
Consequently
\[
S_{y,\mathrm{wk}}(A)=S_{y,\mathrm{wk}}(\tilde{A}),\qquad S'_{y,\mathrm{wk}}(A)=S'_{y,\mathrm{wk}}(\tilde{A}),\qquad S^{\mathrm{sy}}_{y,\mathrm{wk}}(A)=S^{\mathrm{sy}}_{y,\mathrm{wk}}(\tilde{A}).
\]

\emph{The cubes that meet $K_c$.}
Suppose that the entry at $(x,x')$ of $S_{y,\mathrm{wk}}(A)$, of $S'_{y,\mathrm{wk}}(A)$ or of $S^{\mathrm{sy}}_{y,\mathrm{wk}}(A)$ is different from $0$; in the third case one of the two other entries is different from $0$.
By (3-3.7) there are $i,j\in\{0,n\}$ with $h_i(c)\ne0$ and $h_j(c')\ne0$.
Let $(c'',i'')$ be either of the two pairs $(c,i)$, $(c',j)$, so that $h_{i''}(c'')\ne0$.
By display (6.1) of Section~G1.7, $c''$ has a lift $\tilde{c}''$ with $|\tilde{c}''-z_{i''}|<\frac{3M}{4}$, where $z_{i''}$ is the lift of Section~G1.5.
Every site $c'''$ of $K_{c''}=\left\{c''':\ d(c''',c'')\le2L\right\}$ has a lift $\tilde{c}'''$ with $|\tilde{c}'''-\tilde{c}''|\le2L$, whence
\[
|\tilde{c}'''-z_{i''}|<\frac{3M}{4}+2L<\frac{3M}{4}+\frac{M}{2}<\frac{3M}{2},
\]
because $M>4L$ by (M2).
By (3-1.1) it follows that $K_{c''}\subset\square_{i''}^{\sim1}$.
An $M$-cube that meets $K_{c''}$ thus has a site in common with $\square_{i''}^{\sim1}$; by paragraph~1.1\,(i) it is one of the cubes of $\square_{i''}^{\sim1}$, and hence it is contained in $X_{\mathrm{wk}}$.
For $(c'',i'')=(c,i)$ this is the assertion on the row argument, and for $(c'',i'')=(c',j)$ the assertion on the column argument.
\end{proof}

\begin{proof}[Proof of Proposition~3\,(e)]
Let $\boldsymbol{\lambda}$ and $\mathrm{uf}$ be as in paragraph~2.5.
According to (3-2.2), every factor $O(A)$ of the words (3-2.6), (3-2.7), where $O$ is one of $h_i$, $\mathrm{Sloc}_i$, $\mathrm{Cm}_i$, satisfies
\[
O(A-d\boldsymbol{\lambda})=\mathrm{uf}\cdot O(A)\cdot\mathrm{uf}^{-1}.
\]
In the products the inner factors $\mathrm{uf}^{-1}\cdot\mathrm{uf}$ cancel, and
\[
S_{y,\mathrm{wk}}(A-d\boldsymbol{\lambda})=\mathrm{uf}\cdot S_{y,\mathrm{wk}}(A)\cdot\mathrm{uf}^{-1},\qquad S'_{y,\mathrm{wk}}(A-d\boldsymbol{\lambda})=\mathrm{uf}\cdot S'_{y,\mathrm{wk}}(A)\cdot\mathrm{uf}^{-1};
\]
the mean of these two relations is the relation for $S^{\mathrm{sy}}_{y,\mathrm{wk}}$.
Neither $s_{\mathrm{wk}}$ nor $X_{\mathrm{wk}}$ depends on $A$.
Paragraph~0.2 now gives
\begin{align*}
S^{\#}_y(s,A-d\boldsymbol{\lambda})&=\sum_{\mathrm{wk}}s_{\mathrm{wk}}\cdot\mathrm{uf}\cdot S^{\#}_{y,\mathrm{wk}}(A)\cdot\mathrm{uf}^{-1}=\mathrm{uf}\cdot S^{\#}_y(s,A)\cdot\mathrm{uf}^{-1},\\
S^{\#}_y(X,A-d\boldsymbol{\lambda})&=\sum_{\mathrm{wk}:\,X_{\mathrm{wk}}=X}\mathrm{uf}\cdot S^{\#}_{y,\mathrm{wk}}(A)\cdot\mathrm{uf}^{-1}=\mathrm{uf}\cdot S^{\#}_y(X,A)\cdot\mathrm{uf}^{-1},
\end{align*}
where $\mathrm{uf}=\exp\left(i\,e_k\cdot\boldsymbol{\lambda}\right)$.
\end{proof}

\begin{proof}[Proof of Proposition~3\,(f)]
Let $\tilde{g}$, $r$, $y'$ be as in paragraph~2.6.

\emph{Walks.}
The centre of $\tilde{g}\square_i$ is $\tilde{g}z_i$, and $d(\tilde{g}z_{i-1},\tilde{g}z_i)=d(z_{i-1},z_i)\le M$.
So $\tilde{g}\,\mathrm{wk}$ is a walk of the same length as $\mathrm{wk}$, and $\mathrm{wk}\mapsto\tilde{g}\,\mathrm{wk}$ is a bijection of $\mathrm{Wk}_n$ onto itself, with the inverse $\mathrm{wk}\mapsto\tilde{g}^{-1}\mathrm{wk}$.
Paragraph~2.6 gives $(\tilde{g}\square_i)^{\sim1}=\tilde{g}(\square_i^{\sim1})$, so that
\[
X_{\tilde{g}\,\mathrm{wk}}=\bigcup_i\tilde{g}(\square_i^{\sim1})=\tilde{g}X_{\mathrm{wk}}.
\]
The cubes of $\pi_k$ contained in $\tilde{g}X_{\mathrm{wk}}$ are the cubes $\tilde{g}\,\mathrm{pc}$ with $\mathrm{pc}\in\pi_k$ contained in $X_{\mathrm{wk}}$; therefore
\[
(\tilde{g}s)_{\tilde{g}\,\mathrm{wk}}=\prod\left\{(\tilde{g}s)_{\tilde{g}\,\mathrm{pc}}:\ \mathrm{pc}\subset X_{\mathrm{wk}}\right\}=\prod\left\{s_{\mathrm{pc}}:\ \mathrm{pc}\subset X_{\mathrm{wk}}\right\}=s_{\mathrm{wk}}.
\]
Moreover $\tilde{g}s\in\mathrm{Pd}_k$, since the moduli are merely permuted, and $\tilde{g}X\in D_k$, since common faces are preserved.

\emph{Words.}
The words (3-2.6) and (3-2.7) for the data $(y',\tilde{g}\,\mathrm{wk},\tilde{g}A)$ consist of the factors $h_{\tilde{g}\square_i}$, $\mathrm{Sloc}_{\tilde{g}\square_i,y'}(\tilde{g}A)$, $\mathrm{Cm}_{\tilde{g}\square_i,y'}(\tilde{g}A)$, arranged in the orders displayed there, and they carry the same sign $(-1)^n$.
Let the operator $\tilde{g}$ stand at the right end of either word.
By (3-2.3) it can be moved to the left through the factors, one at a time, and each factor is thereby changed into the corresponding factor $h_i$, $\mathrm{Sloc}_i$, $\mathrm{Cm}_i$ for $(y,\mathrm{wk},A)$.
This gives
\[
S_{y',\tilde{g}\,\mathrm{wk}}(\tilde{g}A)\cdot\tilde{g}=\tilde{g}\cdot S_{y,\mathrm{wk}}(A),\qquad S'_{y',\tilde{g}\,\mathrm{wk}}(\tilde{g}A)\cdot\tilde{g}=\tilde{g}\cdot S'_{y,\mathrm{wk}}(A),
\]
and the mean of the two relations is $S^{\mathrm{sy}}_{y',\tilde{g}\,\mathrm{wk}}(\tilde{g}A)\cdot\tilde{g}=\tilde{g}\cdot S^{\mathrm{sy}}_{y,\mathrm{wk}}(A)$.

\emph{Sums.}
The map $\mathrm{wk}\mapsto\tilde{g}\,\mathrm{wk}$ is a bijection of $\mathrm{Wk}$, and $X_{\tilde{g}\,\mathrm{wk}}=\tilde{g}X$ holds if and only if $X_{\mathrm{wk}}=X$.
Paragraph~0.2 therefore gives
\begin{align*}
S^{\#}_{y'}(\tilde{g}s,\tilde{g}A)\cdot\tilde{g}&=\sum_{\mathrm{wk}}(\tilde{g}s)_{\tilde{g}\,\mathrm{wk}}\cdot S^{\#}_{y',\tilde{g}\,\mathrm{wk}}(\tilde{g}A)\cdot\tilde{g}=\sum_{\mathrm{wk}}s_{\mathrm{wk}}\cdot\tilde{g}\cdot S^{\#}_{y,\mathrm{wk}}(A)=\tilde{g}\cdot S^{\#}_y(s,A),\\
S^{\#}_{y'}(\tilde{g}X,\tilde{g}A)\cdot\tilde{g}&=\sum_{\mathrm{wk}:\,X_{\mathrm{wk}}=X}S^{\#}_{y',\tilde{g}\,\mathrm{wk}}(\tilde{g}A)\cdot\tilde{g}=\sum_{\mathrm{wk}:\,X_{\mathrm{wk}}=X}\tilde{g}\cdot S^{\#}_{y,\mathrm{wk}}(A)=\tilde{g}\cdot S^{\#}_y(X,A).\qedhere
\end{align*}
\end{proof}

\begin{proof}[Proof of Proposition~3\,(g)]
Write $\widehat{\mathrm{Sloc}}_i:=\mathrm{Sloc}_{\square_i,-y}(-A)$ and $\widehat{\mathrm{Cm}}_i:=\mathrm{Cm}_{\square_i,-y}(-A)$.
Transposition reverses the order of the factors of (3-2.6):
\[
S_{y,\mathrm{wk}}(A)^T=h_n^T\cdot\mathrm{Sloc}_n^T\cdot\mathrm{Cm}_n^T\cdot h_{n-1}^T\cdot\mathrm{Sloc}_{n-1}^T\cdot\mathrm{Cm}_{n-1}^T\cdots h_1^T\cdot\mathrm{Sloc}_1^T\cdot\mathrm{Cm}_1^T\cdot h_0^T\cdot\mathrm{Sloc}_0^T\cdot h_0^T.
\]
Insert (3-2.4).
The inner products $\mathsf{C}\cdot\mathsf{C}^{-1}$ cancel, and each of the $n$ factors $\mathrm{Cm}_i^T$ contributes one sign; hence
\begin{align*}
S_{y,\mathrm{wk}}(A)^T&=\mathsf{C}^{-1}\cdot(-1)^n\,h_n\cdot\widehat{\mathrm{Sloc}}_n\cdot\widehat{\mathrm{Cm}}_n\cdot h_{n-1}\cdot\widehat{\mathrm{Sloc}}_{n-1}\cdot\widehat{\mathrm{Cm}}_{n-1}\cdots h_1\cdot\widehat{\mathrm{Sloc}}_1\cdot\widehat{\mathrm{Cm}}_1\cdot h_0\cdot\widehat{\mathrm{Sloc}}_0\cdot h_0\cdot\mathsf{C}\\
&=\mathsf{C}^{-1}\cdot S'_{-y,\mathrm{wk}}(-A)\cdot\mathsf{C},
\end{align*}
by (3-2.7) at $(-y,-A)$.
This is the first relation; it holds for all $(y,A)\in Y_k\times\tilde{G}_k$.
At the point $(-y,-A)$ it reads $S_{-y,\mathrm{wk}}(-A)^T=\mathsf{C}^{-1}\cdot S'_{y,\mathrm{wk}}(A)\cdot\mathsf{C}$.
Transposing this identity, with paragraph~0.4 and with $\mathsf{C}^T=\mathsf{C}^{-1}$, $(\mathsf{C}^{-1})^T=\mathsf{C}$ from (I.c3), one obtains
\[
S_{-y,\mathrm{wk}}(-A)=\mathsf{C}^T\cdot S'_{y,\mathrm{wk}}(A)^T\cdot(\mathsf{C}^{-1})^T=\mathsf{C}^{-1}\cdot S'_{y,\mathrm{wk}}(A)^T\cdot\mathsf{C}.
\]
Therefore
\[
S'_{y,\mathrm{wk}}(A)^T=\mathsf{C}\cdot S_{-y,\mathrm{wk}}(-A)\cdot\mathsf{C}^{-1}=\mathsf{C}^{-1}\cdot S_{-y,\mathrm{wk}}(-A)\cdot\mathsf{C},
\]
since $\mathsf{C}=-\mathsf{C}^{-1}$ by (I.c3) and the two signs cancel.
This is the second relation.

Now let $y$ be neutral, that is, $y=0$, or $k\ge1$ and $y=\infty$.
Then $-y=y$, and the mean of the two relations is
\[
S^{\mathrm{sy}}_{y,\mathrm{wk}}(A)^T=\mathsf{C}^{-1}\cdot S^{\mathrm{sy}}_{y,\mathrm{wk}}(-A)\cdot\mathsf{C}.
\]
Neither $s_{\mathrm{wk}}$ nor $X_{\mathrm{wk}}$ depends on $A$ or on the choice between $S$ and $S'$.
By paragraphs~0.2 and~0.4,
\begin{align*}
S^{\mathrm{sy}}_y(s,A)^T&=\sum_{\mathrm{wk}}s_{\mathrm{wk}}\cdot\mathsf{C}^{-1}\cdot S^{\mathrm{sy}}_{y,\mathrm{wk}}(-A)\cdot\mathsf{C}=\mathsf{C}^{-1}\cdot S^{\mathrm{sy}}_y(s,-A)\cdot\mathsf{C},\\
S^{\mathrm{sy}}_y(X,A)^T&=\sum_{\mathrm{wk}:\,X_{\mathrm{wk}}=X}\mathsf{C}^{-1}\cdot S^{\mathrm{sy}}_{y,\mathrm{wk}}(-A)\cdot\mathsf{C}=\mathsf{C}^{-1}\cdot S^{\mathrm{sy}}_y(X,-A)\cdot\mathsf{C}.\qedhere
\end{align*}
\end{proof}

\begin{proof}[Proof of Proposition~3\,(h)]
By paragraph~0.4 the conjugate of a product is the product of the conjugates, and $(-1)^n$ is real.
According to (3-2.5), every factor $O(A)$ of the words (3-2.6), (3-2.7), with $O$ one of $h_i$, $\mathrm{Sloc}_i$, $\mathrm{Cm}_i$, satisfies
\[
\overline{O(\overline{A})}=\varkappa^{-1}\cdot O(-A)\cdot\varkappa,
\]
with the same $y$.
The inner products $\varkappa\cdot\varkappa^{-1}$ cancel, and
\[
\overline{S_{y,\mathrm{wk}}(\overline{A})}=\varkappa^{-1}\cdot S_{y,\mathrm{wk}}(-A)\cdot\varkappa,\qquad\overline{S'_{y,\mathrm{wk}}(\overline{A})}=\varkappa^{-1}\cdot S'_{y,\mathrm{wk}}(-A)\cdot\varkappa;
\]
the mean of the two is the relation for $S^{\mathrm{sy}}$.
The family $\overline{s}:=\left(\overline{s_{\mathrm{pc}}}\right)$ lies in $\mathrm{Pd}_k$, and $(\overline{s})_{\mathrm{wk}}=\overline{s_{\mathrm{wk}}}$.
Paragraphs~0.2 and~0.4 give
\begin{align*}
\overline{S^{\#}_y(\overline{s},\overline{A})}&=\sum_{\mathrm{wk}}\overline{(\overline{s})_{\mathrm{wk}}}\cdot\overline{S^{\#}_{y,\mathrm{wk}}(\overline{A})}=\sum_{\mathrm{wk}}s_{\mathrm{wk}}\cdot\varkappa^{-1}\cdot S^{\#}_{y,\mathrm{wk}}(-A)\cdot\varkappa=\varkappa^{-1}\cdot S^{\#}_y(s,-A)\cdot\varkappa,\\
\overline{S^{\#}_y(X,\overline{A})}&=\sum_{\mathrm{wk}:\,X_{\mathrm{wk}}=X}\varkappa^{-1}\cdot S^{\#}_{y,\mathrm{wk}}(-A)\cdot\varkappa=\varkappa^{-1}\cdot S^{\#}_y(X,-A)\cdot\varkappa.\qedhere
\end{align*}
\end{proof}

\begin{proof}[Proof of Proposition~3\,(i)]
Two cubes of $\pi_k$ are linked if they have a common face.
A link-component of $\mathscr{Y}$ is a class of the equivalence relation on $\mathscr{Y}$ generated by this relation, and a link-component of $Y$ is the union of the cubes of a link-component of $\mathscr{Y}$ (\cite[Definition~5.8]{QEDI}; Section~I.1.1).
Since the $M$-cubes are pairwise disjoint (paragraph~1.1), distinct link-components of $Y$ are disjoint sets of sites, and $Y$ is their union.

\emph{The coefficients $s_{\mathrm{wk}}$.}
Let $\mathrm{wk}$ be a walk such that $X_{\mathrm{wk}}$ is not contained in $Y$.
By paragraph~1.1\,(ii) some cube $\mathrm{pc}$ of $X_{\mathrm{wk}}$ does not belong to $\mathscr{Y}$; then $s_{\mathrm{pc}}=0$ is one of the factors of $s_{\mathrm{wk}}$, and $s_{\mathrm{wk}}=0$.
Let now $X_{\mathrm{wk}}\subset Y$.
By paragraph~1.1\,(ii) the cubes of $X_{\mathrm{wk}}$ belong to $\mathscr{Y}$, and by paragraph~1.2 any two of them are joined by a chain of cubes of $X_{\mathrm{wk}}$ in which consecutive cubes have a common face.
So they lie in one link-component of $\mathscr{Y}$, and $X_{\mathrm{wk}}$ lies in one link-component of $Y$.

\emph{Rows.}
Let $Y_1$ be a link-component of $Y$, let $x=(c,g)$ lie over $Y_1$, and let $\tilde{s}$, $\tilde{A}$ be as in Proposition~3\,(i).
Since $\tilde{s}$ vanishes outside $\mathscr{Y}$, the preceding paragraph applies to $\tilde{s}$ as well.
The walks fall into three classes.
\begin{itemize}
\item[(1)] $X_{\mathrm{wk}}$ is not contained in $Y$.
Then $s_{\mathrm{wk}}=0=\tilde{s}_{\mathrm{wk}}$.
\item[(2)] $X_{\mathrm{wk}}$ is contained in a link-component of $Y$ different from $Y_1$.
Then $c\notin X_{\mathrm{wk}}$, and by (d) the rows at $x$ of $S^{\#}_{y,\mathrm{wk}}(A)$ and of $S^{\#}_{y,\mathrm{wk}}(\tilde{A})$ vanish.
\item[(3)] $X_{\mathrm{wk}}\subset Y_1$.
By (d), $S^{\#}_{y,\mathrm{wk}}(A)(x,x')=0$ unless $x'$ lies over $X_{\mathrm{wk}}$, hence over $Y_1$.
The cubes of $X_{\mathrm{wk}}$ are cubes of $Y_1$, so that $\tilde{s}_{\mathrm{wk}}=s_{\mathrm{wk}}$.
Every bond with both end points in $X_{\mathrm{wk}}$ meets $Y_1$; hence $\tilde{A}$ and $A$ coincide on these bonds, and $S^{\#}_{y,\mathrm{wk}}(\tilde{A})=S^{\#}_{y,\mathrm{wk}}(A)$ by (d).
\end{itemize}
By paragraphs~0.2 and~0.3,
\[
S^{\#}_y(s,A)(x,x')=\sum_{\mathrm{wk}}s_{\mathrm{wk}}\cdot S^{\#}_{y,\mathrm{wk}}(A)(x,x'),
\]
and only the walks of class~(3) contribute.
Consequently $S^{\#}_y(s,A)(x,x')=0$ unless $x'$ lies over $Y_1$, and for all $x'$
\[
S^{\#}_y(s,A)(x,x')=\sum_{\mathrm{wk}:\,X_{\mathrm{wk}}\subset Y_1}s_{\mathrm{wk}}\cdot S^{\#}_{y,\mathrm{wk}}(A)(x,x')=\sum_{\mathrm{wk}:\,X_{\mathrm{wk}}\subset Y_1}\tilde{s}_{\mathrm{wk}}\cdot S^{\#}_{y,\mathrm{wk}}(\tilde{A})(x,x')=S^{\#}_y(\tilde{s},\tilde{A})(x,x').
\]

\emph{Columns.}
Let $x'=(c',g')$ lie over $Y_1$, and let $\tilde{s}$, $\tilde{A}$ be as in Proposition~3\,(i).
Again the walks fall into three classes.
\begin{itemize}
\item[(1)] $X_{\mathrm{wk}}$ is not contained in $Y$.
Then $s_{\mathrm{wk}}=0=\tilde{s}_{\mathrm{wk}}$.
\item[(2)] $X_{\mathrm{wk}}$ is contained in a link-component of $Y$ different from $Y_1$.
Then $c'\notin X_{\mathrm{wk}}$, and by (d) the columns at $x'$ of $S^{\#}_{y,\mathrm{wk}}(A)$ and of $S^{\#}_{y,\mathrm{wk}}(\tilde{A})$ vanish.
\item[(3)] $X_{\mathrm{wk}}\subset Y_1$.
By (d), $S^{\#}_{y,\mathrm{wk}}(A)(x,x')=0$ unless $x$ lies over $X_{\mathrm{wk}}$, hence over $Y_1$.
The cubes of $X_{\mathrm{wk}}$ are cubes of $Y_1$, so that $\tilde{s}_{\mathrm{wk}}=s_{\mathrm{wk}}$.
Every bond with both end points in $X_{\mathrm{wk}}$ meets $Y_1$; hence $\tilde{A}$ and $A$ coincide on these bonds, and $S^{\#}_{y,\mathrm{wk}}(\tilde{A})=S^{\#}_{y,\mathrm{wk}}(A)$ by (d).
\end{itemize}
Only the walks of class~(3) contribute to $\sum_{\mathrm{wk}}s_{\mathrm{wk}}\cdot S^{\#}_{y,\mathrm{wk}}(A)(x,x')$ and to $\sum_{\mathrm{wk}}\tilde{s}_{\mathrm{wk}}\cdot S^{\#}_{y,\mathrm{wk}}(\tilde{A})(x,x')$, and their contributions to the two sums are equal.
Therefore $S^{\#}_y(s,A)(x,x')=0$ unless $x$ lies over $Y_1$, and $S^{\#}_y(s,A)(x,x')=S^{\#}_y(\tilde{s},\tilde{A})(x,x')$ for all $x$.

\emph{The cubes that meet $K_c$.}
Suppose that the row of $S^{\#}_y(s,A)$ at $x=(c,g)$ is different from $0$.
Then $s_{\mathrm{wk}}\cdot S^{\#}_{y,\mathrm{wk}}(A)(x,x')\ne0$ for some walk $\mathrm{wk}$ and some $x'$.
By the paragraph on the coefficients $s_{\mathrm{wk}}$, $X_{\mathrm{wk}}\subset Y$.
By the last assertion of (d), every cube of $\pi_k$ that meets $K_c$ is contained in $X_{\mathrm{wk}}$, hence in $Y$, and by paragraph~1.1\,(ii) it belongs to $\mathscr{Y}$.
The cube of $\pi_k$ that contains $c$ meets $K_c$, because $c\in K_c$; so $c\in Y$.
Hence the row at $x$ vanishes unless $\mathscr{Y}$ contains all cubes of $\pi_k$ that meet $K_c$, and in particular it vanishes if $c\notin Y$.
(If the column of $S^{\#}_y(s,A)$ at $x'=(c',g')$ is different from $0$, then $s_{\mathrm{wk}}\cdot S^{\#}_{y,\mathrm{wk}}(A)(x,x')\ne0$ for some $\mathrm{wk}$ and some $x$; then $X_{\mathrm{wk}}\subset Y$, and the last assertion of (d), applied to the column argument, shows that every cube of $\pi_k$ that meets $K_{c'}$ is contained in $X_{\mathrm{wk}}$; so it belongs to $\mathscr{Y}$.)
\end{proof}

\subsubsection*{4. Uniformity}

The numbers that occur in the bounds of (a) are
\[
32a'\cdot M^{27\alpha_0},\qquad304a'\cdot M^{27\alpha_0},\qquad32a'\cdot\exp(27\kappa'),\qquad\gamma^{F},\qquad\kappa'.
\]
They are built from $a'$, $\gamma^{F}$ (line (C4r) of Proposition~14.2\,(a)), $\alpha_0$, $\kappa'$ (line (C4)) and $M$ (line (C5)).
None of them depends on the level $k$, the side $\ell$, the mass $\mu$ or the number $e_k$, nor on $A$, $y$, $s$, $X$, $z$, $z'$, $x$, $x'$ (Section~I.8.1, (I.U4)$^{\mathrm{ch}}$; Section~G1.1).
The cutoff $N$, the torus and the charge $e$ enter Proposition~3 only through the data $(k,\ell,\mu,e_k)$; hence the bounds are uniform in them as well.
The dependence on $y$ enters through $w_k(y)$, with $w_k(y)=1$ for $k\ge1$ and $w_0(y)=\left(1+|y|\right)^{-1}$.
The number $n_k=\ell^3$ enters only (3-3.4), (3-3.5) and the tail estimates in the proof of (c), which are derived from (3-3.4), that is, the proof that, for given data (SF), the series converge absolutely and uniformly on $\mathrm{Pd}_k\times\tilde{G}_k$.
The bounds of Proposition~3 are independent of $n_k$.
Moreover, the partial sums $\mathrm{Pt}_{\mathrm{nlb}}$ converge in $N$ at a rate independent of $n_k$: restricting the sum over $n$ in the computation of the first bound of (a) to $n\ge\mathrm{nlb}$, one obtains
\[
N\left(S^{\#}_y(s,A)-\mathrm{Pt}_{\mathrm{nlb}}(y,s,A)\right)\le\sum_{n\ge\mathrm{nlb}}16a'\cdot\mathrm{wp}_s\cdot w_k(y)\cdot2^{-n}=32a'\cdot\mathrm{wp}_s\cdot w_k(y)\cdot2^{-\mathrm{nlb}}.
\]
No constant occurs in the clauses (b)--(i).
At level~$0$, Proposition~3 and its proof differ from those at the levels $k\ge1$ in the following points: $Y_0=\mathbb{R}$; the weight is $w_0(y)$; the penalty is (3-2.1$_0$) in place of (3-2.1); there are no short pairs, so that (3-3.2), (3-3.5) and the second bound of (a) are void; $y=0$ is the only neutral value; and (a) contains in addition the kernel bounds, which rest on (3-3.3).

\subsection*{Conditions on the constants}
\leavevmode\par

\noindent
\begin{tabular}{p{0.27\textwidth}p{0.35\textwidth}p{0.27\textwidth}}
\hline
condition & stated in & enters in\\
\hline
$b=1$ (in the form $b>0$) & Proposition~14.2\,(a), line (C0); Section~I.0.5; Section~I.8.2, line (C0) & paragraph~2.2: $b_k>0$, $b_L>0$\\
\hline
$L$ odd, $L\ge5$ & (L5) of Proposition~14.2\,(b); Section~I.0.5; Section~I.8.2, line (C3) & paragraph~1.1 ($M/\eta_k$ is odd); paragraph~0 ($\eta_k<1$ for $k\ge1$); paragraph~2.2 ($b_->0$)\\
\hline
$M$ a power of $L$ (hence $M\ge1$) & (M1) of Proposition~14.2\,(b); Section~I.0.5; Section~I.8.2, line (C5) & paragraph~1.1; paragraph~1.4 ($M^{\alpha_0}\ge1$, $\mathrm{wp}_s\ge1$; the family $s=1$ lies in $\mathrm{Pd}_k$)\\
\hline
$M>4L$ & (M2) of Proposition~14.2\,(b); Section~I.8.2, line (C5) & proof of (d): $\frac{3M}{4}+2L<\frac{3M}{2}$\\
\hline
$M\ge\left(108\pi\cdot C'_V\cdot a'\right)^2$ & (M5) of Proposition~14.2\,(b), second entry of $M^{F}$ (Proposition~14.2\,(a), line (C4r)); Section~I.8.2, line (C5) & (3-1.4): $\mathrm{qq}(\mathrm{wp}_s)\le\frac{1}{2}$\\
\hline
$M\ge108\pi\cdot C'_V\cdot a'\cdot\exp(27\kappa')$ & (M5) of Proposition~14.2\,(b), third entry of $M^{F}$; Section~I.8.2, line (C5) & (3-1.4): $\mathrm{qq}(\mathrm{wp}_X)\le\frac{1}{2}$\\
\hline
$\alpha_0=\frac{1}{54}$ & Proposition~14.2\,(a), line (C4) (Section~I.8.2, line (C4), lists $\alpha_0$) & paragraph~1.4: $\mathrm{wp}_s=M^{\frac{1}{2}}$\\
\hline
$\kappa'=\kappa_{\mathrm{w}}+\kappa_0>0$ & Proposition~14.2\,(a), line (C4); Section~I.8.2, line (C4) & paragraph~1.4 ($\mathrm{wp}_X\ge1$); (3-3.6)\\
\hline
$6L\cdot e_0^{\varepsilon}\le1$ & (Z2.s1) of Proposition~14.2\,(b); Section~I.8.2, line (C8); Section~I.8.4 & paragraph~2.1: Proposition~2.5 at $\delta=2e_k^{\varepsilon}$\\
\hline
$4C'_V\cdot a'\cdot e_0^{\varepsilon}\le1$ & (Z2.s2) of Proposition~14.2\,(b); Section~I.8.2, line (C8); Section~I.8.4 & paragraph~2.1: Proposition~2.5 at $\delta=2e_k^{\varepsilon}$\\
\hline
$\varepsilon>0$ & Section~I.0.5; Section~I.8.2, line (C7); Proposition~14.2\,(a), line (C7), and (b), (S1) ($\varepsilon>1-\alpha>0$) & paragraph~2.1: $e_k^{\varepsilon}\le e_0^{\varepsilon}$ for $e_k\le e_0$; through Proposition~2.7\\
\hline
the conditions of the Propositions~2.m: (L5), (M1), (M2), $M\ge108\pi\cdot C'_V\cdot a'$ (contained in (M5)), (Z2.s1), (Z2.s2) & the paragraph preceding Proposition~2.1; Proposition~14.2\,(b) & through Propositions~2.2, 2.5, 2.6, 2.7, 2.8\\
\hline
$\alpha<1$ & Proposition~14.2\,(a), line (C1) ($\alpha\in(1-\bar{\varepsilon},1)$); Section~I.8.2, line (C1) & only through Proposition~2.7\,(e), second bound ((3-3.2), (3-3.5), the second bound of (a)), and through Proposition~2.8 (conditions of Propositions~2.7 and~2.8)\\
\hline
\end{tabular}

\medskip
All these conditions are rows of Proposition~14.2\,(a), (b).
Two further facts come from the data (SF) and not from the constants.
First, $\ell\in\mathrm{Cl}$ implies that $\ell$ is a multiple of $M$ and that $\ell\ge13L\cdot M$; these facts enter in the forms $3M<\ell$ and $3M+\eta_k<\ell$ (paragraph~1.1).
Second, $e_k\le e_0$ (paragraph~2.1).

\section*{Chapter 4. Resolvent algebra between the fine and the unit lattice: proof of Proposition 4}

\subsection*{0. Conventions of this chapter}

\noindent\textbf{Abbreviations.}
Following Proposition~4 (and Section~G1.10 for $\mathrm{Qk}$ and $\mathrm{Qt}$), let
\[
\mathrm{Qk}:=Q_k(A)\cdot J^{*},\qquad
\mathrm{Qt}:=J\cdot Q_k^{T}(-A),\qquad
P:=P(A);
\]
explicit formulas for these three operators are given in (0.1) below.

\medskip
\noindent\textbf{Notation.}
Throughout the chapter the data $(k,\ell,\mu,e_k)$, $A$ and $y$ are fixed:
$k\ge1$, $\ell\in\mathrm{Cl}$, $\mu\in(0,1]$ and $e_k\in(0,e_0]$ (Section~I.1.8), $A\in\tilde{G}_k$, and $y$ is a real number.
The letter $y$ is the resolvent parameter;
for this reason unit sites are written $z$, $z''$, and sites of the $L$-lattice are written $z'$ (Section~I.1.1).
The letter $\mu$ is the mass, and a direction is written $\nu$.
Further, $\mathrm{Gy}:=i\gamma_5\cdot y$, a $4\times4$ matrix ($\gamma_5$: Section~I.2.1);
$\mathrm{bz}$ denotes a real number $>0$ (later $b_k$ or $b_k+b_L$);
$\mathrm{Mx}$ denotes an arbitrary constant $4\times4$ matrix;
$\mathrm{od}$ denotes an odd positive integer (later a power of $L$);
and
\begin{gather*}
\mathrm{Bq}:=\left(b_k+\mathrm{Gy}\right)^{-1}\cdot(1-P)+\left(b_k+b_L+\mathrm{Gy}\right)^{-1}\cdot P,\qquad
\mathrm{Zq}:=b_k^2\cdot\mathrm{Qt}\cdot\mathrm{Bq}\cdot\mathrm{Qk},\\
\mathrm{Si}:=S_{\infty}(A),\qquad
\mathrm{Sy}:=S_{y}(A),\\
\mathrm{Fu}:=D_k(A)+b_L\cdot P+\mathrm{Gy},\qquad
\mathrm{Gq}:=\mathrm{Bq}+b_k^2\cdot\mathrm{Bq}\cdot\mathrm{Qk}\cdot\mathrm{Sy}\cdot\mathrm{Qt}\cdot\mathrm{Bq}.
\end{gather*}
The operators $\mathrm{Bq}$ and $\mathrm{Zq}$ are formed in Step~2, once Step~1 has shown that the two inverses of $4\times4$ matrices occurring in $\mathrm{Bq}$ exist;
$\mathrm{Si}$, $\mathrm{Sy}$, $\mathrm{Fu}$ and $\mathrm{Gq}$ are formed in Step~3, once the inverses $S_{\infty}(A)$ and $S_{y}(A)$ are known to exist.

\subsection*{Proposition 4}

\begin{proof}[Proof of Proposition~4]
\textbf{Step 0} (blocks, phases, spaces; the operators $\mathrm{Qk}$, $\mathrm{Qt}$, $P$).

\medskip
\noindent\textbf{(0.0)} \emph{Blocks and phases.}
\begin{itemize}
\item[(i)] Every photon site $c$ of $\mathbb{L}_k$ lies in exactly one unit block $\Delta_z$, $z\in\mathbb{Z}^3$.
Every unit site $z$ lies in exactly one $L$-block $B(z')$, $z'\in L\mathbb{Z}^3$.
One has $\#|B(z')|=L^3$.
\item[(ii)] The side $\ell$ is an integer multiple of $L$.
For $\nu=1,2,3$ the translation by $\ell e_{\nu}$ maps each of the lattices $L^j\eta_k\mathbb{Z}^3$, $0\le j\le k+1$, onto itself (among them $\mathbb{L}_k$ for $j=0$, $\mathbb{Z}^3$ for $j=k$ and $L\mathbb{Z}^3$ for $j=k+1$), and
\[
\Delta_{z+\ell e_{\nu}}=\Delta_z+\ell e_{\nu},\qquad
B(z'+\ell e_{\nu})=B(z')+\ell e_{\nu}.
\]
\item[(iii)] For $c\in\Delta_z$ and $z\in B(z')$ the coefficients $\varphi_A(z,c)$ and $U_A(z',z)$ are non-zero complex numbers, not matrices, and
\[
\varphi_A(z+\ell e_{\nu},c+\ell e_{\nu})=\varphi_A(z,c),\qquad
U_A(z'+\ell e_{\nu},z+\ell e_{\nu})=U_A(z',z).
\]
\end{itemize}

\noindent\emph{Proof of (0.0).}
(i) By Section~I.1.1 one has $\mathbb{L}_k=\eta_k\mathbb{Z}^3$ with $\eta_k=L^{-k}$, and
\[
\Delta_z=\left\{c\in\mathbb{L}_k:|c-z|<\tfrac{1}{2}\right\},\qquad
B(z')=\left\{z\in\mathbb{Z}^3:|z-z'|<\tfrac{L}{2}\right\},
\]
where $|\cdot|$ is the sup norm.
Let $\mathrm{od}$ be an odd positive integer.
The integers of modulus less than $\mathrm{od}/2$ are the $\mathrm{od}$ integers $-(\mathrm{od}-1)/2,\dots,(\mathrm{od}-1)/2$, one from each residue class modulo $\mathrm{od}$.
Consequently every integer differs by less than $\mathrm{od}/2$ from exactly one multiple of $\mathrm{od}$.
Since $L$ is odd (Section~I.0.5), every power of $L$ is odd.
Applied with $\mathrm{od}=L^k$ to the three coordinates of $c/\eta_k\in\mathbb{Z}^3$, this gives exactly one $z\in\mathbb{Z}^3$ with $\left|c/\eta_k-L^k\cdot z\right|<L^k/2$, that is, with $|c-z|<1/2$.
Applied with $\mathrm{od}=L$ to the three coordinates of $z$, it gives exactly one $z'\in L\mathbb{Z}^3$ with $|z-z'|<L/2$;
moreover, for fixed $z'$ each coordinate of $z-z'$ takes $L$ values.
With $\mathrm{od}=L^j$ the same holds for the blocks of side $L^j\eta_k$, $0\le j<k$, which occur in (iii): $c$ lies in exactly one of them.
(In accordance with this, no photon site lies on a face of a block, by Section~I.1.1, \emph{Faces}, and the blocks of one side partition the lattice, by Section~G1.2.)

(ii) The condition $\ell\in\mathrm{Cl}$ means that $\ell/(LM)$ is an odd integer $\ge13$ (Section~I.1.1), and $M$ is a power of $L$ (Section~I.0.5).
Hence $\ell$ is an integer multiple of $L$, and also of $L^j\eta_k=L^{j-k}$ for every $j\le k+1$.
The blocks are defined through distances, and distances are invariant under translations.

(iii) The connection $A$ is a complex bond function on $T_k$ (Section~I.5.3), that is, a complex function on the bonds of $\mathbb{L}_k$ that takes the same value on a bond and on its translates by $\ell\mathbb{Z}^3$ (Section~I.1.1).
By (I.3.3) and (I.1.6),
\[
\varphi_A(z,c)=\exp\left(i e_k\cdot\eta_k\cdot(\tau_k A)(z,c)\right),\qquad
U_A(z',z)=\exp\left(i e_k\cdot\eta_k\cdot(\tau A)(z',z)\right)
\]
(here $|z-z'|<L/2\le M$).
The exponential function has no zeros.
By Sections~I.1.1 and I.3.2, $(\tau A)(z',z)$ is the mean, over the six orders $\pi$, of the sums of $A$ over the bonds of the rectilinear paths $\boldsymbol{\Gamma}^{\pi}(z',z)$,
and $(\tau_k A)(z,c)$ is the sum of the numbers $(\tau A)(y_{j+1},y_j)$, $0\le j<k$;
here, in the notation of Section~I.3.2, $y_j\in L^j\eta_k\mathbb{Z}^3$ is the centre of the block of side $L^j\eta_k$ that contains $c$, so that $y_0=c$ and $y_k=z$
(these $y_j$ are sites, not values of the resolvent parameter).
By (ii) the translation by $\ell e_{\nu}$ maps these centres and these paths onto the centres and the paths that belong to $(z+\ell e_{\nu},c+\ell e_{\nu})$, respectively to $(z'+\ell e_{\nu},z+\ell e_{\nu})$,
and $A$ takes the same values on them.

\medskip
\noindent\emph{Spaces.}
By Section~I.1.5 and the table of Section~I.3.5,
\begin{gather*}
\mathrm{Qk}:\mathfrak{S}^{\mathrm{ch}}(\eta_k,\ell)\to\mathfrak{S}(1,\ell),\qquad
\mathrm{Qt}:\mathfrak{S}(1,\ell)\to\mathfrak{S}^{\mathrm{ch}}(\eta_k,\ell),\qquad
P:\mathfrak{S}(1,\ell)\to\mathfrak{S}(1,\ell),\\
Q(A):\mathfrak{S}(1,\ell)\to\mathfrak{S}(L,\ell),\qquad
Q^{T}(-A):\mathfrak{S}(L,\ell)\to\mathfrak{S}(1,\ell),
\end{gather*}
and $D_A$, $J\cdot P_k(A)\cdot J^{*}$, $J\cdot P_{k+1}(A)\cdot J^{*}$, $\operatorname{Pen}_{y}(A)$, $F_{y}(A)$, $F_{\infty}(A)$ map $\mathfrak{S}^{\mathrm{ch}}(\eta_k,\ell)$ into itself.
A constant $4\times4$ matrix $\mathrm{Mx}$ acts on these spaces site by site:
$(\mathrm{Mx}\,v)(z):=\mathrm{Mx}\cdot v(z)$ and $(\mathrm{Mx}\,f)(c,g):=\mathrm{Mx}\cdot f(c,g)$.
A scalar is identified with the corresponding multiple of the identity map, and $1$ denotes the identity map of the space in question.
With these identifications every sum and every product below is a linear map whose sources and targets match:
products of the form $\mathrm{Qt}\cdot(\,\cdot\,)\cdot\mathrm{Qk}$ map $\mathfrak{S}^{\mathrm{ch}}(\eta_k,\ell)$ into itself, and products of the form $\mathrm{Qk}\cdot(\,\cdot\,)\cdot\mathrm{Qt}$ map $\mathfrak{S}(1,\ell)$ into itself.

\medskip
\noindent\emph{Formulas.}
Let $f\in\mathfrak{S}^{\mathrm{ch}}(\eta_k,\ell)$ and $v\in\mathfrak{S}(1,\ell)$;
let $z$ be a unit site, $c\in\Delta_z$ and $g\in O_{\mathrm{cub}}$, so that $(c,g)$ is a fermion site (Sections~I.1.2, I.1.3);
let $z'$ be the site of $L\mathbb{Z}^3$ with $z\in B(z')$ (see (0.0)\,(i)).
By (I.3.1) and Section~I.3.3,
\begin{equation}
\tag{0.1}
\begin{aligned}
(\mathrm{Qk}\,f)(z)&=\left(Q_k(A)\cdot J^{*}f\right)(z)=\eta_k^3\sum_{c'\in\Delta_z}\varphi_A(z,c')\cdot\frac{1}{48}\sum_{g'\in O_{\mathrm{cub}}}f(c',g'),\\
(\mathrm{Qt}\,v)(c,g)&=\left(J\cdot Q_k^{T}(-A)\,v\right)(c,g)=\varphi_A(z,c)^{-1}\cdot v(z),\\
(P\,v)(z)&=\left(Q^{T}(-A)\cdot Q(A)\,v\right)(z)=U_A(z',z)^{-1}\cdot L^{-3}\sum_{z''\in B(z')}U_A(z',z'')\cdot v(z'').
\end{aligned}
\end{equation}
By (0.0)\,(i), the second and the third line define $(\mathrm{Qt}\,v)(c,g)$ at every fermion site $(c,g)$ and $(P\,v)(z)$ at every unit site $z$.
By (0.0)\,(iii), the coefficients $\varphi_A(z,c')$, $\varphi_A(z,c)^{-1}$, $U_A(z',z)^{-1}$, $U_A(z',z'')$ are non-zero complex numbers.
The value $(\mathrm{Qt}\,v)(c,g)$ does not depend on the label $g$.

The boundary condition, that is, the statement that the three operators take values in the spaces named in the table of Section~I.3.5, is checked as follows.
By (0.0)\,(ii), the sums over $\Delta_{z+\ell e_{\nu}}$ and over $B(z'+\ell e_{\nu})$ are the sums over $c'+\ell e_{\nu}$, $c'\in\Delta_z$, and over $z''+\ell e_{\nu}$, $z''\in B(z')$.
With (0.0)\,(iii) and the antiperiodicity $f(c'+\ell e_{\nu},g')=-f(c',g')$, $v(z+\ell e_{\nu})=-v(z)$, one obtains
\begin{align*}
(\mathrm{Qk}\,f)(z+\ell e_{\nu})&=\eta_k^3\sum_{c'\in\Delta_z}\varphi_A(z+\ell e_{\nu},c'+\ell e_{\nu})\cdot\frac{1}{48}\sum_{g'\in O_{\mathrm{cub}}}f(c'+\ell e_{\nu},g')=-(\mathrm{Qk}\,f)(z),\\
(\mathrm{Qt}\,v)(c+\ell e_{\nu},g)&=\varphi_A(z+\ell e_{\nu},c+\ell e_{\nu})^{-1}\cdot v(z+\ell e_{\nu})=-(\mathrm{Qt}\,v)(c,g),\\
(P\,v)(z+\ell e_{\nu})&=U_A(z'+\ell e_{\nu},z+\ell e_{\nu})^{-1}\cdot L^{-3}\sum_{z''\in B(z')}U_A(z'+\ell e_{\nu},z''+\ell e_{\nu})\cdot v(z''+\ell e_{\nu})\\
&=-(P\,v)(z).
\end{align*}

\medskip
\noindent\textbf{(0.2)} For every constant $4\times4$ matrix $\mathrm{Mx}$,
\begin{equation}
\tag{0.2}
\mathrm{Qk}\cdot\mathrm{Mx}=\mathrm{Mx}\cdot\mathrm{Qk},\qquad
\mathrm{Qt}\cdot\mathrm{Mx}=\mathrm{Mx}\cdot\mathrm{Qt},\qquad
P\cdot\mathrm{Mx}=\mathrm{Mx}\cdot P.
\end{equation}
Indeed, each line of (0.1) is a linear combination, with numerical coefficients, of values of $f$, respectively of $v$, and $\mathrm{Mx}$ acts linearly on $\mathbb{C}^4$:
\begin{align*}
(\mathrm{Qk}\cdot\mathrm{Mx}\,f)(z)&=\eta_k^3\sum_{c'\in\Delta_z}\varphi_A(z,c')\cdot\frac{1}{48}\sum_{g'\in O_{\mathrm{cub}}}\mathrm{Mx}\cdot f(c',g')=\mathrm{Mx}\cdot(\mathrm{Qk}\,f)(z),\\
(\mathrm{Qt}\cdot\mathrm{Mx}\,v)(c,g)&=\varphi_A(z,c)^{-1}\cdot\mathrm{Mx}\cdot v(z)=\mathrm{Mx}\cdot(\mathrm{Qt}\,v)(c,g),\\
(P\cdot\mathrm{Mx}\,v)(z)&=U_A(z',z)^{-1}\cdot L^{-3}\sum_{z''\in B(z')}U_A(z',z'')\cdot\mathrm{Mx}\cdot v(z'')=\mathrm{Mx}\cdot(P\,v)(z).
\end{align*}
The three relations are also contained in Proposition~2.2\,(f0), applied on the block torus $T_k$ as in the proof of (0.3) below:
in its assertion on constant $4\times4$ matrices, $P(A)$ is one of the idempotents, and $Q_k(A)\cdot J^{*}$, $J\cdot Q_k^{T}(-A)$ are two of the six operators.

\medskip
\noindent\textbf{(0.3)} One has
\begin{equation}
\tag{0.3}
P^2=P;\qquad\text{hence}\qquad(1-P)^2=1-P,\qquad P\cdot(1-P)=(1-P)\cdot P=0.
\end{equation}
\emph{Proof of (0.3).}
By Proposition~2.2\,(f0), applied to the block torus $T_k=(0,\ell,-1)$ of Section~G1.2, whose spinors on the unit lattice are the elements of $\mathfrak{S}(1,\ell)$, and to the complex connection $A$ on it, $P(A)$ is idempotent.
The same follows directly from the formulas of Section~I.3.3.
Let $v'\in\mathfrak{S}(L,\ell)$ and $z'\in L\mathbb{Z}^3$.
Since every $z\in B(z')$ lies in no other $L$-block and $\#|B(z')|=L^3$ (see (0.0)\,(i)),
\begin{align*}
\left(Q(A)\cdot Q^{T}(-A)\,v'\right)(z')
&=L^{-3}\sum_{z\in B(z')}U_A(z',z)\cdot\left(Q^{T}(-A)\,v'\right)(z)\\
&=L^{-3}\sum_{z\in B(z')}U_A(z',z)\cdot U_A(z',z)^{-1}\cdot v'(z')=v'(z').
\end{align*}
Thus $Q(A)\cdot Q^{T}(-A)=1$ on $\mathfrak{S}(L,\ell)$, and
\[
P^2=Q^{T}(-A)\cdot\left[Q(A)\cdot Q^{T}(-A)\right]\cdot Q(A)=Q^{T}(-A)\cdot Q(A)=P.
\]
The remaining three relations follow by multiplying out.

\medskip
\noindent\textbf{(0.4)} One has
\begin{equation}
\tag{0.4}
\mathrm{Qt}\cdot\mathrm{Qk}=J\cdot P_k(A)\cdot J^{*},\qquad
\mathrm{Qt}\cdot P\cdot\mathrm{Qk}=J\cdot P_{k+1}(A)\cdot J^{*}.
\end{equation}
Indeed, by the associativity of composition,
$\mathrm{Qt}\cdot\mathrm{Qk}=J\cdot Q_k^{T}(-A)\cdot Q_k(A)\cdot J^{*}$ and $\mathrm{Qt}\cdot P\cdot\mathrm{Qk}=J\cdot Q_k^{T}(-A)\cdot P(A)\cdot Q_k(A)\cdot J^{*}$,
and $P_k(A):=Q_k^{T}(-A)\cdot Q_k(A)$, $P_{k+1}(A):=Q_k^{T}(-A)\cdot P(A)\cdot Q_k(A)$ are the definitions of Section~I.3.3.
Written out with (0.1), for $c\in\Delta_z$, $z\in B(z')$ and $g\in O_{\mathrm{cub}}$,
\begin{align*}
(\mathrm{Qt}\cdot\mathrm{Qk}\,f)(c,g)
&=\varphi_A(z,c)^{-1}\cdot\eta_k^3\sum_{c'\in\Delta_z}\varphi_A(z,c')\cdot\frac{1}{48}\sum_{g'\in O_{\mathrm{cub}}}f(c',g'),\\
(\mathrm{Qt}\cdot P\cdot\mathrm{Qk}\,f)(c,g)
&=\varphi_A(z,c)^{-1}\cdot U_A(z',z)^{-1}\cdot L^{-3}\sum_{z''\in B(z')}U_A(z',z'')\\
&\qquad\cdot\eta_k^3\sum_{c'\in\Delta_{z''}}\varphi_A(z'',c')\cdot\frac{1}{48}\sum_{g'\in O_{\mathrm{cub}}}f(c',g').
\end{align*}

\medskip
\noindent\textbf{Step 1} (the $4\times4$ matrices).
By (I.c1), $\gamma_5^2=1$.
Hence $\mathrm{Gy}^2=i^2\cdot y^2\cdot\gamma_5^2=-y^2$ (times $1_4$).
For every real number $\mathrm{bz}>0$,
\[
\left(\mathrm{bz}+\mathrm{Gy}\right)\cdot\left(\mathrm{bz}-\mathrm{Gy}\right)
=\left(\mathrm{bz}-\mathrm{Gy}\right)\cdot\left(\mathrm{bz}+\mathrm{Gy}\right)
=\mathrm{bz}^2-\mathrm{Gy}^2=\mathrm{bz}^2+y^2>0.
\]
Therefore $\mathrm{bz}+\mathrm{Gy}$ is invertible, and
\begin{equation}
\tag{1.1}
\left(\mathrm{bz}+\mathrm{Gy}\right)^{-1}
=\left(\mathrm{bz}^2+y^2\right)^{-1}\cdot\left(\mathrm{bz}-\mathrm{Gy}\right)
=\left(\mathrm{bz}^2+y^2\right)^{-1}\cdot\left(\mathrm{bz}-i\gamma_5\cdot y\right).
\end{equation}
The two cases needed are $\mathrm{bz}=b_k$ and $\mathrm{bz}=b_k+b_L$.
Both numbers are positive: $b_k\ge b_{-}=b\cdot(1-1/L)>0$ and $b_L=b/L>0$ (Sections~I.3.4 and I.0.5: $b=1$, $L\ge5$).
The matrices $\mathrm{Gy}$, $b_k+\mathrm{Gy}$, $b_k+b_L+\mathrm{Gy}$ and, by (1.1), the two inverses $\left(b_k+\mathrm{Gy}\right)^{-1}$, $\left(b_k+b_L+\mathrm{Gy}\right)^{-1}$ are polynomials in the single matrix $\gamma_5$.
Consequently any two of these five matrices commute, and by (0.2) each of them commutes with $\mathrm{Qk}$, $\mathrm{Qt}$ and $P$.
In particular the matrices of Section~I.3.5,
\[
\boldsymbol{\alpha}_{y}=b_k\cdot\mathrm{Gy}\cdot\left(b_k+\mathrm{Gy}\right)^{-1},\qquad
\boldsymbol{\beta}_{y}=b_k^2\cdot b_L\cdot\left(b_k+b_L+\mathrm{Gy}\right)^{-1}\cdot\left(b_k+\mathrm{Gy}\right)^{-1},
\]
are defined, and they commute with $\mathrm{Qk}$, $\mathrm{Qt}$, $P$.
They satisfy two relations.
The first is
\begin{equation}
\tag{1.2}
b_k-b_k^2\cdot\left(b_k+\mathrm{Gy}\right)^{-1}
=b_k\cdot\left[\left(b_k+\mathrm{Gy}\right)-b_k\right]\cdot\left(b_k+\mathrm{Gy}\right)^{-1}
=b_k\cdot\mathrm{Gy}\cdot\left(b_k+\mathrm{Gy}\right)^{-1}
=\boldsymbol{\alpha}_{y}.
\end{equation}
For the second, one computes
\begin{align*}
b_k-b_k^2\cdot\left(b_k+b_L+\mathrm{Gy}\right)^{-1}
&=b_k\cdot\left[\left(b_k+b_L+\mathrm{Gy}\right)-b_k\right]\cdot\left(b_k+b_L+\mathrm{Gy}\right)^{-1}\\
&=b_k\cdot\left(b_L+\mathrm{Gy}\right)\cdot\left(b_k+b_L+\mathrm{Gy}\right)^{-1},\\
b_k\cdot\left(b_L+\mathrm{Gy}\right)\cdot\left(b_k+b_L+\mathrm{Gy}\right)^{-1}-\boldsymbol{\alpha}_{y}
&=b_k\cdot\left[\left(b_L+\mathrm{Gy}\right)\cdot\left(b_k+\mathrm{Gy}\right)-\mathrm{Gy}\cdot\left(b_k+b_L+\mathrm{Gy}\right)\right]\\
&\qquad\cdot\left(b_k+b_L+\mathrm{Gy}\right)^{-1}\cdot\left(b_k+\mathrm{Gy}\right)^{-1},
\end{align*}
and
\begin{multline*}
\left(b_L+\mathrm{Gy}\right)\cdot\left(b_k+\mathrm{Gy}\right)-\mathrm{Gy}\cdot\left(b_k+b_L+\mathrm{Gy}\right)\\
=b_L\cdot b_k+b_L\cdot\mathrm{Gy}+b_k\cdot\mathrm{Gy}+\mathrm{Gy}^2-b_k\cdot\mathrm{Gy}-b_L\cdot\mathrm{Gy}-\mathrm{Gy}^2
=b_k\cdot b_L,
\end{multline*}
so that
\begin{equation}
\tag{1.3}
b_k-b_k^2\cdot\left(b_k+b_L+\mathrm{Gy}\right)^{-1}
=\boldsymbol{\alpha}_{y}+b_k^2\cdot b_L\cdot\left(b_k+b_L+\mathrm{Gy}\right)^{-1}\cdot\left(b_k+\mathrm{Gy}\right)^{-1}
=\boldsymbol{\alpha}_{y}+\boldsymbol{\beta}_{y}.
\end{equation}

\medskip
\noindent\textbf{Step 2} (clause (a)).
\emph{The first identity of (a).}
Since $\left(b_k+\mathrm{Gy}\right)\cdot(1-P)+\left(b_k+b_L+\mathrm{Gy}\right)\cdot P=\left(b_k+\mathrm{Gy}\right)+b_L\cdot P$, one has
\begin{equation}
\tag{2.1}
b_k+b_L\cdot P+\mathrm{Gy}=\left(b_k+\mathrm{Gy}\right)\cdot(1-P)+\left(b_k+b_L+\mathrm{Gy}\right)\cdot P.
\end{equation}
Multiply (2.1) by $\mathrm{Bq}=\left(b_k+\mathrm{Gy}\right)^{-1}\cdot(1-P)+\left(b_k+b_L+\mathrm{Gy}\right)^{-1}\cdot P$, once from the right and once from the left.
By Step~1 all the matrices involved commute with $P$.
By (0.3) the two mixed terms vanish, as they contain the factor $(1-P)\cdot P=0$, respectively $P\cdot(1-P)=0$, and
\begin{align*}
\left(b_k+b_L\cdot P+\mathrm{Gy}\right)\cdot\mathrm{Bq}
&=\left(b_k+\mathrm{Gy}\right)\cdot\left(b_k+\mathrm{Gy}\right)^{-1}\cdot(1-P)^2\\
&\qquad+\left(b_k+b_L+\mathrm{Gy}\right)\cdot\left(b_k+b_L+\mathrm{Gy}\right)^{-1}\cdot P^2\\
&=(1-P)+P=1,\\
\mathrm{Bq}\cdot\left(b_k+b_L\cdot P+\mathrm{Gy}\right)
&=\left(b_k+\mathrm{Gy}\right)^{-1}\cdot\left(b_k+\mathrm{Gy}\right)\cdot(1-P)^2\\
&\qquad+\left(b_k+b_L+\mathrm{Gy}\right)^{-1}\cdot\left(b_k+b_L+\mathrm{Gy}\right)\cdot P^2\\
&=(1-P)+P=1.
\end{align*}
Hence $b_k+b_L\cdot P+i\gamma_5\cdot y$ is invertible on $\mathfrak{S}(1,\ell)$, and its inverse, which is the operator $B_{k,y}(A)$ of Section~I.3.5, equals $\mathrm{Bq}$:
\begin{equation}
\tag{2.2}
B_{k,y}(A)=\mathrm{Bq},\qquad
\mathrm{Bq}^{-1}=b_k+b_L\cdot P+\mathrm{Gy}.
\end{equation}
(By (1.1),
$B_{k,y}(A)=\left(b_k^2+y^2\right)^{-1}\cdot\left(b_k-i\gamma_5\cdot y\right)\cdot(1-P)+\left(\left(b_k+b_L\right)^2+y^2\right)^{-1}\cdot\left(b_k+b_L-i\gamma_5\cdot y\right)\cdot P$.)

\emph{The second identity of (a).}
Writing $b_k=b_k\cdot(1-P)+b_k\cdot P$ and using (1.2) and (1.3), one obtains
\begin{equation}
\tag{2.3}
\begin{aligned}
b_k-b_k^2\cdot\mathrm{Bq}
&=\left[b_k-b_k^2\cdot\left(b_k+\mathrm{Gy}\right)^{-1}\right]\cdot(1-P)+\left[b_k-b_k^2\cdot\left(b_k+b_L+\mathrm{Gy}\right)^{-1}\right]\cdot P\\
&=\boldsymbol{\alpha}_{y}\cdot(1-P)+\left(\boldsymbol{\alpha}_{y}+\boldsymbol{\beta}_{y}\right)\cdot P
=\boldsymbol{\alpha}_{y}+\boldsymbol{\beta}_{y}\cdot P.
\end{aligned}
\end{equation}

\emph{The third identity of (a).}
Multiply (2.3) by $\mathrm{Qt}$ from the left and by $\mathrm{Qk}$ from the right.
The left-hand side becomes $b_k\cdot\mathrm{Qt}\cdot\mathrm{Qk}-b_k^2\cdot\mathrm{Qt}\cdot\mathrm{Bq}\cdot\mathrm{Qk}$, and $\mathrm{Qt}\cdot\mathrm{Qk}=J\cdot P_k(A)\cdot J^{*}$ by (0.4).
On the right-hand side, $\boldsymbol{\alpha}_{y}$ and $\boldsymbol{\beta}_{y}$ commute with $\mathrm{Qt}$ (Step~1), and (0.4) is applied to both terms:
\begin{equation}
\tag{2.4}
\begin{aligned}
b_k\cdot J\cdot P_k(A)\cdot J^{*}-b_k^2\cdot\mathrm{Qt}\cdot\mathrm{Bq}\cdot\mathrm{Qk}
&=\mathrm{Qt}\cdot\left(\boldsymbol{\alpha}_{y}+\boldsymbol{\beta}_{y}\cdot P\right)\cdot\mathrm{Qk}
=\boldsymbol{\alpha}_{y}\cdot\mathrm{Qt}\cdot\mathrm{Qk}+\boldsymbol{\beta}_{y}\cdot\mathrm{Qt}\cdot P\cdot\mathrm{Qk}\\
&=\boldsymbol{\alpha}_{y}\cdot J\cdot P_k(A)\cdot J^{*}+\boldsymbol{\beta}_{y}\cdot J\cdot P_{k+1}(A)\cdot J^{*}
=\operatorname{Pen}_{y}(A).
\end{aligned}
\end{equation}
By Section~I.3.5, $\boldsymbol{\alpha}_{\infty}=b_k$ and $\boldsymbol{\beta}_{\infty}=0$;
hence $\operatorname{Pen}_{\infty}(A)=b_k\cdot J\cdot P_k(A)\cdot J^{*}$ and $F_{\infty}(A)=D_A+\mu+b_k\cdot J\cdot P_k(A)\cdot J^{*}$.
Adding $D_A+\mu$ to both sides of (2.4) gives
\begin{equation}
\tag{2.5}
F_{\infty}(A)-b_k^2\cdot\mathrm{Qt}\cdot\mathrm{Bq}\cdot\mathrm{Qk}=D_A+\mu+\operatorname{Pen}_{y}(A)=F_{y}(A).
\end{equation}
As $\mathrm{Bq}=B_{k,y}(A)$, the identities (2.2), (2.3) and (2.5) are the three assertions of clause (a).

\medskip
\noindent\textbf{Step 3} (clause (b)).
For $k\ge1$ one has $Y_{k}=\mathbb{R}\cup\{\infty\}$ (Section~I.3.5).
By Proposition~2.8\,(a), applied with the parameter $\infty$ and with the real parameter $y$, the maps $F_{\infty}(A)$ and $F_{y}(A)$ are bijective on $\mathfrak{S}^{\mathrm{ch}}(\eta_k,\ell)$.
Hence $\mathrm{Si}=S_{\infty}(A)=F_{\infty}(A)^{-1}=S_k(A)$ and $\mathrm{Sy}=S_{y}(A)=F_{y}(A)^{-1}$ exist, and with $\mathrm{Si}$ the operators of Section~I.3.5
\begin{gather*}
H_{k}(A)=b_k\cdot\mathrm{Si}\cdot\mathrm{Qt},\qquad
\mathfrak{H}_{k}(A)=b_k\cdot\mathrm{Qk}\cdot\mathrm{Si},\\
D_k(A)=b_k-b_k^2\cdot\mathrm{Qk}\cdot\mathrm{Si}\cdot\mathrm{Qt},\qquad
M_k(A)=b_k-D_k(A)=b_k^2\cdot\mathrm{Qk}\cdot\mathrm{Si}\cdot\mathrm{Qt}
\end{gather*}
are defined.
By (2.5) and the definition of $\mathrm{Zq}$,
\begin{equation}
\tag{3.1}
F_{\infty}(A)-F_{y}(A)=\mathrm{Zq}.
\end{equation}
Multiplying (3.1) by $\mathrm{Si}$ from the left and by $\mathrm{Sy}$ from the right, one obtains on the left-hand side $\mathrm{Si}\cdot F_{\infty}(A)\cdot\mathrm{Sy}-\mathrm{Si}\cdot F_{y}(A)\cdot\mathrm{Sy}=\mathrm{Sy}-\mathrm{Si}$.
Multiplying (3.1) by $\mathrm{Sy}$ from the left and by $\mathrm{Si}$ from the right, one obtains on the left-hand side $\mathrm{Sy}\cdot F_{\infty}(A)\cdot\mathrm{Si}-\mathrm{Sy}\cdot F_{y}(A)\cdot\mathrm{Si}=\mathrm{Sy}-\mathrm{Si}$.
Consequently
\begin{equation}
\tag{3.2}
\mathrm{Sy}-\mathrm{Si}=\mathrm{Si}\cdot\mathrm{Zq}\cdot\mathrm{Sy},\qquad
\mathrm{Sy}-\mathrm{Si}=\mathrm{Sy}\cdot\mathrm{Zq}\cdot\mathrm{Si}.
\end{equation}
By the definition of $D_k(A)$ and by (2.2),
\begin{equation}
\tag{3.3}
\mathrm{Fu}=b_k-b_k^2\cdot\mathrm{Qk}\cdot\mathrm{Si}\cdot\mathrm{Qt}+b_L\cdot P+\mathrm{Gy}=\mathrm{Bq}^{-1}-b_k^2\cdot\mathrm{Qk}\cdot\mathrm{Si}\cdot\mathrm{Qt}.
\end{equation}
Multiplying out, and recognising the factor $b_k^2\cdot\mathrm{Qt}\cdot\mathrm{Bq}\cdot\mathrm{Qk}=\mathrm{Zq}$ in the last of the four terms that each product produces, one finds
\begin{align*}
\mathrm{Fu}\cdot\mathrm{Gq}
&=\left(\mathrm{Bq}^{-1}-b_k^2\cdot\mathrm{Qk}\cdot\mathrm{Si}\cdot\mathrm{Qt}\right)\cdot\left(\mathrm{Bq}+b_k^2\cdot\mathrm{Bq}\cdot\mathrm{Qk}\cdot\mathrm{Sy}\cdot\mathrm{Qt}\cdot\mathrm{Bq}\right)\\
&=1+b_k^2\cdot\mathrm{Qk}\cdot\mathrm{Sy}\cdot\mathrm{Qt}\cdot\mathrm{Bq}-b_k^2\cdot\mathrm{Qk}\cdot\mathrm{Si}\cdot\mathrm{Qt}\cdot\mathrm{Bq}\\
&\qquad-b_k^2\cdot\mathrm{Qk}\cdot\mathrm{Si}\cdot\left(b_k^2\cdot\mathrm{Qt}\cdot\mathrm{Bq}\cdot\mathrm{Qk}\right)\cdot\mathrm{Sy}\cdot\mathrm{Qt}\cdot\mathrm{Bq}\\
&=1+b_k^2\cdot\mathrm{Qk}\cdot\left[\mathrm{Sy}-\mathrm{Si}-\mathrm{Si}\cdot\mathrm{Zq}\cdot\mathrm{Sy}\right]\cdot\mathrm{Qt}\cdot\mathrm{Bq}=1
\end{align*}
by the first relation (3.2), and
\begin{align*}
\mathrm{Gq}\cdot\mathrm{Fu}
&=\left(\mathrm{Bq}+b_k^2\cdot\mathrm{Bq}\cdot\mathrm{Qk}\cdot\mathrm{Sy}\cdot\mathrm{Qt}\cdot\mathrm{Bq}\right)\cdot\left(\mathrm{Bq}^{-1}-b_k^2\cdot\mathrm{Qk}\cdot\mathrm{Si}\cdot\mathrm{Qt}\right)\\
&=1-b_k^2\cdot\mathrm{Bq}\cdot\mathrm{Qk}\cdot\mathrm{Si}\cdot\mathrm{Qt}+b_k^2\cdot\mathrm{Bq}\cdot\mathrm{Qk}\cdot\mathrm{Sy}\cdot\mathrm{Qt}\\
&\qquad-b_k^2\cdot\mathrm{Bq}\cdot\mathrm{Qk}\cdot\mathrm{Sy}\cdot\left(b_k^2\cdot\mathrm{Qt}\cdot\mathrm{Bq}\cdot\mathrm{Qk}\right)\cdot\mathrm{Si}\cdot\mathrm{Qt}\\
&=1+b_k^2\cdot\mathrm{Bq}\cdot\mathrm{Qk}\cdot\left[\mathrm{Sy}-\mathrm{Si}-\mathrm{Sy}\cdot\mathrm{Zq}\cdot\mathrm{Si}\right]\cdot\mathrm{Qt}=1
\end{align*}
by the second relation (3.2).
So $\mathrm{Fu}=D_k(A)+b_L\cdot P(A)+i\gamma_5\cdot y$ is a bijection of $\mathfrak{S}(1,\ell)$ onto itself, with the inverse $\mathrm{Gq}$;
this inverse is the operator $\Gamma_{k,y}(A)$ of Section~I.3.5.
With $\mathrm{Bq}=B_{k,y}(A)$ and $\mathrm{Sy}=S_{y}(A)$,
\begin{equation}
\tag{3.4}
\Gamma_{k,y}(A)=\mathrm{Gq}=B_{k,y}(A)+b_k^2\cdot B_{k,y}(A)\cdot\mathrm{Qk}\cdot S_{y}(A)\cdot\mathrm{Qt}\cdot B_{k,y}(A).
\end{equation}

\medskip
\noindent\textbf{Step 4} (clause (c)).
By (3.4), by $H_{k}(A)=b_k\cdot\mathrm{Si}\cdot\mathrm{Qt}$, $\mathfrak{H}_{k}(A)=b_k\cdot\mathrm{Qk}\cdot\mathrm{Si}$, and by the definition of $\mathrm{Zq}$,
\begin{align*}
H_{k}(A)\cdot\Gamma_{k,y}(A)\cdot\mathfrak{H}_{k}(A)
&=b_k\cdot\mathrm{Si}\cdot\mathrm{Qt}\cdot\left(\mathrm{Bq}+b_k^2\cdot\mathrm{Bq}\cdot\mathrm{Qk}\cdot\mathrm{Sy}\cdot\mathrm{Qt}\cdot\mathrm{Bq}\right)\cdot b_k\cdot\mathrm{Qk}\cdot\mathrm{Si}\\
&=\mathrm{Si}\cdot\left(b_k^2\cdot\mathrm{Qt}\cdot\mathrm{Bq}\cdot\mathrm{Qk}\right)\cdot\mathrm{Si}\\
&\qquad+\mathrm{Si}\cdot\left(b_k^2\cdot\mathrm{Qt}\cdot\mathrm{Bq}\cdot\mathrm{Qk}\right)\cdot\mathrm{Sy}\cdot\left(b_k^2\cdot\mathrm{Qt}\cdot\mathrm{Bq}\cdot\mathrm{Qk}\right)\cdot\mathrm{Si}\\
&=\mathrm{Si}\cdot\mathrm{Zq}\cdot\mathrm{Si}+\mathrm{Si}\cdot\mathrm{Zq}\cdot\mathrm{Sy}\cdot\mathrm{Zq}\cdot\mathrm{Si}
=\mathrm{Si}\cdot\mathrm{Zq}\cdot\left(\mathrm{Si}+\mathrm{Sy}\cdot\mathrm{Zq}\cdot\mathrm{Si}\right).
\end{align*}
By the second relation (3.2), $\mathrm{Si}+\mathrm{Sy}\cdot\mathrm{Zq}\cdot\mathrm{Si}=\mathrm{Sy}$;
by the first one, $\mathrm{Si}\cdot\mathrm{Zq}\cdot\mathrm{Sy}=\mathrm{Sy}-\mathrm{Si}$.
Therefore $H_{k}(A)\cdot\Gamma_{k,y}(A)\cdot\mathfrak{H}_{k}(A)=\mathrm{Sy}-\mathrm{Si}$, that is,
\[
S_{y}(A)=S_k(A)+H_{k}(A)\cdot\Gamma_{k,y}(A)\cdot\mathfrak{H}_{k}(A).
\]

\medskip
\noindent\textbf{Step 5} (clause (d)).
Since Steps~1--3 are valid for every real $y$, they may be applied at $y=0$.
There $\mathrm{Gy}=0$ and $\mathrm{Fu}=D_k(A)+b_L\cdot P(A)$, which by Step~3 is a bijection of $\mathfrak{S}(1,\ell)$ onto itself.
Its inverse, the operator $\Gamma_k(A)$ of Section~I.3.5, therefore exists:
$\Gamma_k(A)=\left(D_k(A)+b_L\cdot P(A)\right)^{-1}$, and it is $\Gamma_{k,y}(A)$ at $y=0$.
By the definitions of Section~I.3.5, $B_{k,y}(A)$ at $y=0$ is $B_{k}(A)$, and $S_{y}(A)$ at $y=0$ is $S^{0}_{k+1}(A)$.
Thus (3.4) at $y=0$ reads
\[
\Gamma_k(A)=B_{k}(A)+b_k^2\cdot B_{k}(A)\cdot\mathrm{Qk}\cdot S^{0}_{k+1}(A)\cdot\mathrm{Qt}\cdot B_{k}(A).
\]
The space $\mathfrak{S}(1,\ell)$ has the finite dimension $4\ell^3=4n_k$ (Sections~I.1.5, I.1.1), and a bijective linear map of a finite-dimensional space into itself has a non-zero determinant; hence
\[
\delta Z_k(A)=\det\left(D_k(A)+b_L\cdot P(A)\right)\ne0.
\]
\end{proof}

\begin{named}{Remark 1}[A more general setting]
Let $\ell$ be an integer multiple of $L$ with $\ell\ge3L$, so that $T_k$ is a block torus in the sense of Section~G1.2, the setting of Proposition~2.2\,(f0);
let $A$ be an arbitrary complex bond function on $T_k$, let $e_k$ and $\mu$ be arbitrary numbers with $e_k>0$ and $\mu\ge0$, as in Proposition~2.2, and let $y$ be real.
For these data the argument of Steps~0--2 applies verbatim.
Hence clause~(a) holds for every complex bond function $A$ on $T_k$ and every real $y$.
If, in addition, $F_{\infty}(A)$ and $F_{y}(A)$ are bijective on $\mathfrak{S}^{\mathrm{ch}}(\eta_k,\ell)$, the arguments of Steps~3 and~4 give clauses~(b) and~(c);
if $F_{\infty}(A)$ and $F_{y}(A)$ at $y=0$ are bijective, Step~5 gives clause~(d).
The same conclusions hold when $D_A$ is replaced by an arbitrary linear map of $\mathfrak{S}^{\mathrm{ch}}(\eta_k,\ell)$ into itself, $F_{\infty}(A)$, $F_{y}(A)$ and the operators derived from them being formed with that map.
For the data of Proposition~4, that is, for $A\in\tilde{G}_k$, $\mu\in(0,1]$, $e_k\le e_0$ and $\ell\in\mathrm{Cl}$, the maps $F_{\infty}(A)$ and $F_{y}(A)$ are bijective for every real $y$ by Proposition~2.8\,(a).
\end{named}

\begin{named}{Remark 2}[Level 0]
Proposition~4 carries the datum $k\ge1$, and none of its clauses has a counterpart of the same form at $k=0$:
$b_0$, $S_k(A)=S_{\infty}(A)$ at $k=0$, $M_0(A)$, $B_{0,y}(A)$ and $\Gamma_{0,y}(A)$ are not defined (Section~I.3.9; $Y_{0}=\mathbb{R}$ does not contain $\infty$).
At $k=0$ the fine lattice and the unit lattice coincide, and by the definitions of Section~I.3.9
\begin{gather*}
F_{y}(A)\text{ at }y=0\text{ is }F_0(A)=D_A+\mu+b_L\cdot J\cdot P(A)\cdot J^{*}\text{ on }\mathfrak{S}^{\mathrm{ch}}(1,\ell),\\
\Gamma_0(A)=S^{0}_{1}(A)=F_0(A)^{-1},\qquad
\delta Z_0(A)=\det F_0(A).
\end{gather*}
In the place of (d), the following holds for all data (SF) with $k=0$ and all $A\in\tilde{G}_0$.
The map $F_0(A)$ is bijective on $\mathfrak{S}^{\mathrm{ch}}(1,\ell)$ by Proposition~2.8\,(a) at $k=0$ and $y=0$ (here $\eta_0=1$, and $0$ lies in $Y_{0}=\mathbb{R}$);
the space $\mathfrak{S}^{\mathrm{ch}}(1,\ell)$ has the finite dimension $192\ell^3$ (Section~I.1.5);
hence $\Gamma_0(A)$ exists, and $\delta Z_0(A)\ne0$ by the fact on determinants used at the end of Step~5.
These two conclusions are the last assertion of Proposition~2.8\,(d).
\end{named}

\begin{named}{Remark 3}[Range of validity]
The clauses of Proposition~4 are exact identities between linear maps and assertions of invertibility, each for a single datum $(k,\ell,\mu,e_k,A,y)$;
they contain no constants and no bounds.
They hold on the whole range of data on which Proposition~2.8\,(a) holds: all (SF) with $k\ge1$, all $A\in\tilde{G}_k$ and all real $y$.
\end{named}

\subsubsection*{Conditions on the constants}
\leavevmode\par\nopagebreak
\noindent
\begin{tabular}{|p{0.28\textwidth}|p{0.32\textwidth}|p{0.28\textwidth}|}
\hline
condition & where it enters & where it stands in the ordered list\\
\hline
$b>0$ and $L>1$ (they give $b_k\ge b_{-}>0$ and $b_L>0$)
& Step~1
& Section~I.8.2, that is, Proposition~14.2\,(a): (C0) $b:=1$; (C3) $L$, with $L\ge5$ from (L5); (C4) $b_L$, $b_{-}$\\
\hline
$L$ odd (every photon site lies in exactly one $\Delta_z$, every unit site in exactly one $B(z')$; $\#|B(z')|=L^3$)
& Step~0: (0.0)\,(i); through it the formulas (0.1) for $\mathrm{Qt}=J\cdot Q_k^{T}(-A)$ and for $P=Q^{T}(-A)\cdot Q(A)$, the phases $\varphi_A$, and (0.3)
& Proposition~14.2\,(b): (L5), in the line (C3)\\
\hline
$M$ a power of $L$, $M\ge L^2$ (what enters: $M$ is an integer, so that a side $\ell\in\mathrm{Cl}$ is an integer multiple of $L$ and $\ell\ge3L$; and $L/2\le M$)
& Step~0: (0.0)\,(ii), (iii); $T_k$ is a block torus (Section~G1.2)
& Proposition~14.2\,(b): (M1), in the line (C5)\\
\hline
$L\ge5$; (M2) $M>4L$; $M\ge108\pi\cdot C'_{V}\cdot a'$ (part of (M5)); (Z2.s1) $6L\cdot e_0^{\varepsilon}\le1$; (Z2.s2) $4C'_{V}\cdot a'\cdot e_0^{\varepsilon}\le1$; $\alpha<1$
& Step~3, by way of Proposition~2.8\,(a), and through Step~3 also Steps~4 and~5: the first five stand in the list of conditions at the head of the Propositions~2.m, and $\alpha<1$ is among the conditions of Proposition~2.8, through Proposition~2.7\,(d), (e) ($a'$ contains $1/(1-\alpha)$)
& Proposition~14.2\,(b): (L5) in the line (C3); (M2), (M5) in the line (C5); (Z2.s1), (Z2.s2) in the line (C8); Proposition~14.2\,(a): $\alpha\in(1-\bar{\varepsilon},1)$, line (C1)\\
\hline
\end{tabular}

\medskip
Consequently, Proposition~4 holds for every choice of the constants that satisfies the conditions of Proposition~14.2\,(a), (b).

\section*{Chapter 5. Kernels of the unit-lattice operators at the levels $k\ge1$ (FE3$^{\mathrm{ch}}$): proof of Proposition~5}

\subsection*{Proof of Proposition~5}

\subsubsection*{0. Conventions of this chapter; preliminaries}

\paragraph{0.0. Conventions.}
Throughout this chapter the data (SF) of Section~G1.1 have $k\ge1$;
the one exception is the last paragraph of Section~6, which treats clause (a) at $k=0$.
The connection $A$ belongs to $\tilde{G}_k$ (Section~I.5.3), $\eta_k=L^{-k}$, and, as in Section~G1.10 and in Proposition~4,
\[
\mathrm{Qk}:=Q_k(A)\cdot J^{*},\qquad\mathrm{Qt}:=J\cdot Q_k^{T}(-A),\qquad P:=P(A).
\]
In this chapter $y$ denotes a real number or, where this is said, an element of $Y_k=\mathbb{R}\cup\{\infty\}$;
$y'$ denotes a point of $L\mathbb{Z}^3$;
$z,z',z'',z_1,z_2$ denote points of $\mathbb{Z}^3$, regarded as lifts of sites of $T^0_k$;
$n,n',n'',\tilde{n}$ denote vectors of $\mathbb{Z}^3$, and $(-1)^n:=(-1)^{n_1+n_2+n_3}$ for $n=(n_1,n_2,n_3)$;
$\mathrm{vv}$ denotes a vector of $\mathbb{C}^4$, and $\mathrm{Mx}$ a $4\times4$ matrix.
All operators are linear.

\paragraph{0.1. $L$-blocks on the torus.}
By Section~I.1.1 the number $\ell/(LM)$ is an odd integer $\ge13$, and $M$ is an integer, namely a power of $L$ (Section~I.0.5).
Consequently $\ell\mathbb{Z}^3\subset L\mathbb{Z}^3$ and $L-1<\ell$.
Since $L$ is odd (Section~I.0.5), every integer $\mathrm{ti}$ has exactly one representation $\mathrm{ti}=L\,\mathrm{ti}'+\mathrm{ti}''$ with integers $\mathrm{ti}'$, $\mathrm{ti}''$ such that $|\mathrm{ti}''|\le(L-1)/2$;
moreover, for the integer vector $z-y'$ the conditions $|z-y'|<L/2$ and $|z-y'|\le(L-1)/2$ are equivalent.
It follows that
\begin{equation}\tag{0.1}
\begin{gathered}
\text{each $z$ lies in exactly one block }B(y')=\left\{z:|z-y'|<\frac{L}{2}\right\}=\left\{z:|z-y'|\le\frac{L-1}{2}\right\};\\
\#|B(y')|=L^3;\qquad|z-z''|\le L-1\ \text{ for }z,z''\in B(y');\qquad B(y')+\ell n=B(y'+\ell n).
\end{gathered}
\end{equation}

\begin{named}{Definition}
Points $z,z'$ of $\mathbb{Z}^3$ \emph{lie in one $L$-block} (on the torus), in symbols $z\sim z'$, if for some $n\in\mathbb{Z}^3$ the points $z$ and $z'+\ell n$ lie in one block $B(y')$.
One puts $\mathbf{1}[\text{one $L$-block}]:=\mathbf{1}[z\sim z']$;
further, $\mathbf{1}[z=z']:=1$ if $z-z'\in\ell\mathbb{Z}^3$, and $\mathbf{1}[z=z']:=0$ otherwise.
\end{named}

The relation $\sim$ has the following four properties.

(Lb1) Whether $z\sim z'$ holds depends only on the classes of $z$ and $z'$ modulo $\ell\mathbb{Z}^3$, and the relation is symmetric.
Indeed, suppose that $z$ and $z'+\ell n$ lie in $B(y')$.
By the last relation in (0.1), the points $z+\ell n'$ and $(z'+\ell n'')+\ell\,(n+n'-n'')$ lie in $B(y'+\ell n')$ for all $n',n''\in\mathbb{Z}^3$;
and the points $z'$ and $z-\ell n$ lie in $B(y'-\ell n)$.
If $z-z'\in\ell\mathbb{Z}^3$, then $z\sim z'$;
in other words, $\mathbf{1}[z=z']\le\mathbf{1}[\text{one $L$-block}]$.

(Lb2) The vector $n$ in the definition is unique:
if $z'+\ell n$ and $z'+\ell\tilde{n}$ lie in one block $B(y')$, then $\ell\cdot|n-\tilde{n}|\le L-1<\ell$ by (0.1), and therefore $n=\tilde{n}$.

(Lb3) Let $\mathrm{fd}$ be a fundamental domain of $\ell\mathbb{Z}^3$ in $\mathbb{Z}^3$.
For every $z$ there are exactly $L^3$ points $z'\in\mathrm{fd}$ with $z\sim z'$,
and for every $z'$ there are exactly $L^3$ points $z\in\mathrm{fd}$ with $z\sim z'$.
To see this, let $B(y')$ be the block of $z$, and assign to each $z''\in B(y')$ the point of $\mathrm{fd}$ that is congruent to $z''$ modulo $\ell\mathbb{Z}^3$.
This map is injective by (Lb2), because two congruent points $z''$ and $z''+\ell n$ of $B(y')$ have $n=0$;
and its image is the set $\{z'\in\mathrm{fd}:z\sim z'\}$.
This is the first assertion.
By (Lb1), $z\sim z'$ is equivalent to $z'\sim z$;
so the second assertion is the first one with the letters $z$ and $z'$ interchanged.

(Lb4) If $z\sim z'$, then $d(z,z')\le|z-(z'+\ell n)|\le L-1$.
Here $d(z,z')=\min_n|z-z'-\ell n|$ is the sup distance on the torus (Section~I.1.1);
it satisfies the triangle inequality.

\paragraph{0.2. Kernels of operators on $\mathfrak{S}(1,\ell)$.}
Given $z'$ and $\mathrm{vv}$, let $\mathrm{dsp}(z',\mathrm{vv})$ be the function on $\mathbb{Z}^3$ whose value at the point $z'+\ell n$ ($n\in\mathbb{Z}^3$) is $(-1)^n\,\mathrm{vv}$ and whose value at every other point is $0$;
the points $z'+\ell n$ are pairwise distinct.
This function belongs to $\mathfrak{S}(1,\ell)$ (Section~I.1.5):
the point $z'+\ell\,(n+e_\mu)$ carries the value $-(-1)^n\,\mathrm{vv}$, and the set of the remaining points is invariant under $\ell\mathbb{Z}^3$.
It depends linearly on $\mathrm{vv}$, and $\mathrm{dsp}(z'+\ell e_\mu,\mathrm{vv})=-\mathrm{dsp}(z',\mathrm{vv})$.
For an operator $O$ on $\mathfrak{S}(1,\ell)$ the $4\times4$ matrix $O(z,z')$ is defined by
\begin{equation}\tag{0.2}
O(z,z')\cdot\mathrm{vv}:=\left(O\,\mathrm{dsp}(z',\mathrm{vv})\right)(z).
\end{equation}

(Kn1) For every $v\in\mathfrak{S}(1,\ell)$ and every fundamental domain $\mathrm{fd}$ one has
\[
v=\sum_{z'\in\mathrm{fd}}\mathrm{dsp}\left(z',v(z')\right);
\]
indeed, at the point $z'_0+\ell n$ with $z'_0\in\mathrm{fd}$ both sides take the value $(-1)^n\,v(z'_0)$.
Hence $(O\,v)(z)=\sum_{z'\in\mathrm{fd}}O(z,z')\cdot v(z')$:
the matrix $O(z,z')$ is the kernel of Section~I.1.6 and of the head of Section~G1.11 for operators on $\mathfrak{S}(1,\ell)$, whatever the choice of $\mathrm{fd}$.
No other matrix function has this property, as one sees by inserting $v=\mathrm{dsp}(z',\mathrm{vv})$ and choosing a fundamental domain that contains $z'$.

(Kn2) One has $O(z+\ell e_\mu,z')=O(z,z'+\ell e_\mu)=-O(z,z')$,
since $O\,\mathrm{dsp}(z',\mathrm{vv})$ belongs to $\mathfrak{S}(1,\ell)$, $\mathrm{dsp}(z'+\ell e_\mu,\mathrm{vv})=-\mathrm{dsp}(z',\mathrm{vv})$, and $O$ is linear.
Therefore $|O(z,z')|_2$, $|O(z,z')|_c$, and the vanishing of $O(z,z')$ depend only on the classes of $z$ and $z'$ modulo $\ell\mathbb{Z}^3$ (Section~G1.2 with $R=\ell$, $\boldsymbol{\sigma}=-1$).
The row sums $\sum_{z'\in\mathrm{fd}}|O(z,z')|_2$ and the column sums $\sum_{z\in\mathrm{fd}}|O(z,z')|_2$ do not depend on $\mathrm{fd}$.

(Kn3) One has
\[
(O_1\cdot O_2)(z,z')=\sum_{z''\in\mathrm{fd}}O_1(z,z'')\cdot O_2(z'',z').
\]
This is (Kn1) for $O_1$ and for $v:=O_2\,\mathrm{dsp}(z',\mathrm{vv})$, whose values are $v(z'')=O_2(z'',z')\cdot\mathrm{vv}$.

(Kn4) Let a constant $4\times4$ matrix $\mathrm{Mx}$ act on $\mathfrak{S}(1,\ell)$ at each site.
Then
\[
\mathrm{Mx}(z,z')=\mathrm{sz}(z,z')\cdot\mathrm{Mx}\qquad\text{and}\qquad(\mathrm{Mx}\cdot O)(z,z')=\mathrm{Mx}\cdot O(z,z'),
\]
where $\mathrm{sz}(z,z'):=(-1)^n$ if $z=z'+\ell n$, and $\mathrm{sz}(z,z'):=0$ if $z-z'\notin\ell\mathbb{Z}^3$;
note that $|\mathrm{sz}(z,z')|=\mathbf{1}[z=z']$.
Indeed,
\begin{gather*}
\left(\mathrm{Mx}\,\mathrm{dsp}(z',\mathrm{vv})\right)(z)=\mathrm{Mx}\cdot\left(\mathrm{dsp}(z',\mathrm{vv})(z)\right)=\mathrm{sz}(z,z')\cdot\mathrm{Mx}\cdot\mathrm{vv},\\
\left(\mathrm{Mx}\cdot O\,\mathrm{dsp}(z',\mathrm{vv})\right)(z)=\mathrm{Mx}\cdot\left(O\,\mathrm{dsp}(z',\mathrm{vv})\right)(z).
\end{gather*}

\paragraph{0.3. The imaginary part of $A$.}
The last assertion of the display (G1) in \cite[Section~16.2.7]{QEDI} gives
\[
\mathfrak{i}(A)=e_k\cdot|\operatorname{Im}A|_\infty<e_k^{\varepsilon}.
\]
In fact $A=A_0+A_1$ with $A_0$ real and $|A_1|_\infty<e_k^{-1+\varepsilon}$ (Section~I.5.3), so that $\operatorname{Im}A=\operatorname{Im}A_1$;
the quantity $\mathfrak{i}(A)$ is defined in (I.1.6).
The condition (Z2.s1) is imposed for all values $\ge-\log e_0$ of its variable (Proposition~14.2\,(b));
at the value $v_k=-\log e_k\ge-\log e_0$ (Section~G1.1) it reads $6L\cdot e_k^{\varepsilon}\le1$.
Hence
\begin{equation}\tag{0.3}
\mathfrak{i}(A)<\frac{1}{6L}.
\end{equation}

\paragraph{0.4. Phases.}
Recall that $T_k$ is the block torus $(0,\ell,-1)$ of Section~G1.2.
For the connection $A$ on $T_k$, and with $c$, $c'$ written for photon sites, \cite[Lemma~16.2(a)]{QEDI} gives
\[
\left|\varphi_A(z,c)^{\pm1}\right|\le\exp\left(\frac{3}{2}\,\mathfrak{i}(A)\right)\quad(c\in\Delta_z),\qquad
\left|U_A(c,c')^{\pm1}\right|\le\exp\left(3\,|c-c'|\cdot\mathfrak{i}(A)\right).
\]
Both bounds rest on the following facts.
For a complex number $\mathrm{wz}$ one has $\left|\exp(i\,\mathrm{wz})\right|=\exp\left(-\operatorname{Im}\mathrm{wz}\right)$.
A path $\boldsymbol{\Gamma}^{\pi}(c,c')$ has $|c-c'|_1/\eta_k\le3\,|c-c'|/\eta_k$ bonds;
hence $e_k\cdot\eta_k\cdot\left|\operatorname{Im}(\tau A)(c,c')\right|\le3\,|c-c'|\cdot\mathfrak{i}(A)$.
The centres $c=y_0,y_1,\dots,y_k=z$ of Section~I.3.2 satisfy $|y_{j+1}-y_j|\le\left(L^{j+1}-L^j\right)\eta_k/2$.
To see this, note that $y_j-y_{j+1}\in L^j\eta_k\mathbb{Z}^3$ and $|y_j-y_{j+1}|\le|y_j-c|+|c-y_{j+1}|<(L+1)\,L^j\eta_k/2$;
since $L$ is odd, $(L+1)/2$ is an integer, and so $|y_j-y_{j+1}|\le\left(\frac{L+1}{2}-1\right)L^j\eta_k$.
Consequently
\[
e_k\cdot\eta_k\cdot\left|\operatorname{Im}(\tau_kA)(z,c)\right|\le\sum_{j<k}\frac{3}{2}\left(L^{j+1}-L^j\right)\eta_k\cdot\mathfrak{i}(A)=\frac{3}{2}\left(L^k-1\right)\eta_k\cdot\mathfrak{i}(A)\le\frac{3}{2}\,\mathfrak{i}(A).
\]
For $z\in B(y')$ one has $|y'-z|\le(L-1)/2\le M$ (Section~I.0.5: $M\ge L^2$), so that $U_A(y',z)$ is defined by (I.1.6) and Section~I.3.2.
Together with (0.3) this yields
\begin{equation}\tag{0.4}
\begin{gathered}
\left|U_A(y',z)^{\pm1}\right|\le\exp\left(\frac{3}{2}\,(L-1)\cdot\mathfrak{i}(A)\right),\\
\exp\left(3\,(L-1)\cdot\mathfrak{i}(A)\right)\le\exp\left(\frac{1}{2}\right)\le\mathrm{e},\qquad
\exp\left(3\,\mathfrak{i}(A)\right)\le\exp\left(\frac{1}{2L}\right)\le\mathrm{e}.
\end{gathered}
\end{equation}

\paragraph{0.5. The operators $\mathrm{Qk}$ and $\mathrm{Qt}$ in the Euclidean sup norms.}
By (I.3.1) and Section~I.3.3, for $f\in\mathfrak{S}^{\mathrm{ch}}(\eta_k,\ell)$ and $v\in\mathfrak{S}(1,\ell)$,
\begin{equation}\tag{0.5}
\begin{gathered}
(\mathrm{Qk}\,f)(z)=\eta_k^3\sum_{c\in\Delta_z}\varphi_A(z,c)\cdot\frac{1}{48}\sum_{g\in O_{\mathrm{cub}}}f(c,g),\\
(\mathrm{Qt}\,v)(c,g)=\varphi_A(z,c)^{-1}\cdot v(z)\qquad\left(c\in\Delta_z,\ g\in O_{\mathrm{cub}}\right).
\end{gathered}
\end{equation}
The block $\Delta_z=\{c\in\mathbb{L}_k:|c-z|<1/2\}=\{z+\eta_k n:n\in\mathbb{Z}^3,\ |n|\le(L^k-1)/2\}$ (recall that $L^k$ is odd) has $L^{3k}=\eta_k^{-3}$ points,
and $O_{\mathrm{cub}}$ has $48$ elements (Section~I.1.2).
Thus $(\mathrm{Qk}\,f)(z)$ is the arithmetic mean of the $48\,\eta_k^{-3}$ vectors $\varphi_A(z,c)\cdot f(c,g)$, $(c,g)\in\Delta_z\times O_{\mathrm{cub}}$.
The operator $\mathbf{1}_{\Delta_z}$ is the multiplication by the indicator of the orbit of $\Delta_z\times O_{\mathrm{cub}}$ under $\ell\mathbb{Z}^3$ (Section~G1.3),
and ``$x$ over $\Delta_z$'' means $x=(c,g)$ with $c\in\Delta_z$ (Section~I.1.3, (I.POS)).
The first bound of Section~0.4 now gives two facts.

(Qs1) One has $(\mathrm{Qk}\,f)(z)=\left(\mathrm{Qk}\cdot\mathbf{1}_{\Delta_z}f\right)(z)$ and
\[
\left|(\mathrm{Qk}\,f)(z)\right|_2\le\exp\left(\frac{3}{2}\,\mathfrak{i}(A)\right)\cdot\sup\left\{|f(x)|_2:x\text{ over }\Delta_z\right\}.
\]
Indeed, by (0.5) the vector $(\mathrm{Qk}\,f)(z)$ involves the values of $f$ over $\Delta_z$ only;
the norm of an arithmetic mean of vectors is at most the largest of their norms;
and $\left|\varphi_A(z,c)\cdot f(c,g)\right|_2=\left|\varphi_A(z,c)\right|\cdot|f(c,g)|_2$.

(Qs2) One has $\mathrm{Qt}\,\mathrm{dsp}(z',\mathrm{vv})=\mathbf{1}_{\Delta_{z'}}\mathrm{Qt}\,\mathrm{dsp}(z',\mathrm{vv})$ and
\[
\left|\mathrm{Qt}\,\mathrm{dsp}(z',\mathrm{vv})\right|_{\mathrm{su}}\le\exp\left(\frac{3}{2}\,\mathfrak{i}(A)\right)\cdot|\mathrm{vv}|_2.
\]
Indeed, by (0.5) the value at $(c,g)$, $c\in\Delta_{z''}$, is $\varphi_A(z'',c)^{-1}\cdot\mathrm{dsp}(z',\mathrm{vv})(z'')$.
It vanishes unless $z''=z'+\ell n$ for some $n$, and $\Delta_{z'+\ell n}=\Delta_{z'}+\ell n$;
where it does not vanish, its norm is $\left|\varphi_A(z'',c)^{-1}\right|\cdot|\mathrm{vv}|_2$.

The facts (Qs1), (Qs2) are (I.J2), (I.Q3) for the Euclidean norm $|\cdot|_2$ on $\mathbb{C}^4$,
and the two arguments above use of $|\cdot|_2$ only that it is a norm on $\mathbb{C}^4$.
They therefore hold with the norms $|\cdot|_c$, $|\cdot|_\infty$ of Section~I.1.5 in place of $|\cdot|_2$, $|\cdot|_{\mathrm{su}}$:
\begin{equation}\tag{0.5$'$}
\begin{gathered}
\left|(\mathrm{Qk}\,f)(z)\right|_c\le\exp\left(\frac{3}{2}\,\mathfrak{i}(A)\right)\cdot\sup\left\{|f(x)|_c:x\text{ over }\Delta_z\right\},\\
\left|\mathrm{Qt}\,\mathrm{dsp}(z',\mathrm{vv})\right|_\infty\le\exp\left(\frac{3}{2}\,\mathfrak{i}(A)\right)\cdot|\mathrm{vv}|_c.
\end{gathered}
\end{equation}

\paragraph{0.6. Norms of $4\times4$ blocks.}
For a $4\times4$ matrix $\mathrm{Mx}$ and $\mathrm{vv}\in\mathbb{C}^4$, Section~I.1.5 gives
\[
|\mathrm{Mx}\cdot\mathrm{vv}|_c\le|\mathrm{Mx}\cdot\mathrm{vv}|_2\le|\mathrm{Mx}|_2\cdot|\mathrm{vv}|_2\le2\,|\mathrm{Mx}|_2\cdot|\mathrm{vv}|_c,
\]
and hence $|\mathrm{Mx}|_c\le2\,|\mathrm{Mx}|_2$.
For a complex number $\mathrm{wz}$ one has $|\mathrm{wz}\cdot\mathbf{1}_4|_c=|\mathrm{wz}\cdot\mathbf{1}_4|_2=|\mathrm{wz}|$.

\paragraph{0.7. The rate; the numbers $b_k$.}
The constants defined before the Propositions~2.m (see Proposition~14.2\,(a), line (C4r)) satisfy
\[
\gamma^{F}=\gamma_{z}/2,\qquad\gamma_{z}=\min\left(\gamma_{(5)},a_{z}\right),\qquad\gamma_{(5)}>0,\qquad0<a_{z}\le\frac{1}{3L}.
\]
Hence $0<\gamma^{F}\le1/(6L)$, and $0<\gamma^{F}L\le1/6<\log2$, because $\exp(1/6)<\exp(1/2)<2$ as $\mathrm{e}<4$.
It follows that
\begin{equation}\tag{0.7}
\exp\left(\gamma^{F}(L-1)\right)\le\exp\left(\gamma^{F}L\right)\le2,\qquad\exp\left(2\gamma^{F}(L-1)\right)\le\exp\left(2\gamma^{F}L\right)\le4.
\end{equation}
By Sections~I.0.5 and~I.3.4 one has $b>0$ and $L>1$;
so $b_{-}=b\cdot\left(1-L^{-1}\right)>0$ and $b_{L}=b/L>0$, and $b_k=b_{-}/(1-L^{-k})$ with $1-L^{-1}\le1-L^{-k}<1$ for $k\ge1$.
Therefore $0<b_{-}\le b_k\le b_{-}/(1-L^{-1})=b=b_{+}$.

\paragraph{0.8. The rules (N1), (N2) on $T_k$.}
Let $\gamma\ge0$, and let $O$, $O_1$, $O_2$ be operators on $\mathfrak{S}^{\mathrm{ch}}(\eta_k,\ell)$.
In $N_{\gamma}$ (Section~G1.3, formed on the block torus $T_k$) the sites range over the classes of $\mathbb{Z}^3$ modulo $\ell\mathbb{Z}^3$, that is, over $\mathrm{fd}$.
By Section~I.1.4 the sets $\Delta_{z''}\times O_{\mathrm{cub}}$, $z''\in\mathbb{Z}^3$, form a partition of $\mathbb{L}^{\mathrm{f}}_k$, and $\Delta_{z''+\ell n}=\Delta_{z''}+\ell n$.
Hence the orbits under $\ell\mathbb{Z}^3$ of the sets $\Delta_{z''}\times O_{\mathrm{cub}}$ with $z''\in\mathrm{fd}$ form a partition of $\mathbb{L}^{\mathrm{f}}_k$; that is,
\[
\sum_{z''\in\mathrm{fd}}\mathbf{1}_{\Delta_{z''}}=1\qquad\text{and}\qquad\mathbf{1}_{\Delta_{z''}}\cdot\mathbf{1}_{\Delta_{z''}}=\mathbf{1}_{\Delta_{z''}}.
\]

\emph{(N1), second part.}
The number $\exp\left(\gamma\cdot d(z,z')\right)\cdot\left\|\mathbf{1}_{\Delta_z}\cdot O\cdot\mathbf{1}_{\Delta_{z'}}\right\|_{\mathrm{su}}$ is one of the terms of the sum
\[
\sum_{z''\in\mathrm{fd}}\exp\left(\gamma\cdot d(z,z'')\right)\cdot\left\|\mathbf{1}_{\Delta_z}\cdot O\cdot\mathbf{1}_{\Delta_{z''}}\right\|_{\mathrm{su}},
\]
whose terms are non-negative and which is at most $N_{\gamma}(O)$.
Therefore $\left\|\mathbf{1}_{\Delta_z}\cdot O\cdot\mathbf{1}_{\Delta_{z'}}\right\|_{\mathrm{su}}\le\exp\left(-\gamma\cdot d(z,z')\right)\cdot N_{\gamma}(O)$.

\emph{(N2), product.}
One has
\[
\mathbf{1}_{\Delta_z}\cdot O_1\cdot O_2\cdot\mathbf{1}_{\Delta_{z'}}=\sum_{z''\in\mathrm{fd}}\left(\mathbf{1}_{\Delta_z}\cdot O_1\cdot\mathbf{1}_{\Delta_{z''}}\right)\cdot\left(\mathbf{1}_{\Delta_{z''}}\cdot O_2\cdot\mathbf{1}_{\Delta_{z'}}\right)
\]
and $\exp\left(\gamma\cdot d(z,z')\right)\le\exp\left(\gamma\cdot d(z,z'')\right)\cdot\exp\left(\gamma\cdot d(z'',z')\right)$.
Being an operator norm, $\|\cdot\|_{\mathrm{su}}$ is subadditive and submultiplicative.
Hence, for every $z$,
\begin{align*}
\sum_{z'\in\mathrm{fd}}\exp\left(\gamma\cdot d(z,z')\right)\cdot\left\|\mathbf{1}_{\Delta_z}\cdot O_1\cdot O_2\cdot\mathbf{1}_{\Delta_{z'}}\right\|_{\mathrm{su}}
&\le\sum_{z''\in\mathrm{fd}}\exp\left(\gamma\cdot d(z,z'')\right)\cdot\left\|\mathbf{1}_{\Delta_z}\cdot O_1\cdot\mathbf{1}_{\Delta_{z''}}\right\|_{\mathrm{su}}\\
&\qquad\cdot\left[\sum_{z'\in\mathrm{fd}}\exp\left(\gamma\cdot d(z'',z')\right)\cdot\left\|\mathbf{1}_{\Delta_{z''}}\cdot O_2\cdot\mathbf{1}_{\Delta_{z'}}\right\|_{\mathrm{su}}\right]\\
&\le N_{\gamma}(O_1)\cdot N_{\gamma}(O_2),
\end{align*}
that is, $N_{\gamma}(O_1\cdot O_2)\le N_{\gamma}(O_1)\cdot N_{\gamma}(O_2)$.

\subsubsection*{1. Clause (a)}

Let $z\in B(y')$ and $v\in\mathfrak{S}(1,\ell)$.
By Section~I.3.3,
\begin{equation}\tag{1.1}
\begin{aligned}
(P\,v)(z)&=\left(Q^{T}(-A)\cdot Q(A)\,v\right)(z)=U_A(y',z)^{-1}\cdot\left(Q(A)\,v\right)(y')\\
&=L^{-3}\sum_{z''\in B(y')}U_A(y',z)^{-1}\cdot U_A(y',z'')\cdot v(z'').
\end{aligned}
\end{equation}
Now take $v:=\mathrm{dsp}(z',\mathrm{vv})$.
Then $v(z'')=0$ unless $z''=z'+\ell n$ for some $n$, and by (Lb2) at most one of these points lies in $B(y')$.
With $n$ the vector for which $z'+\ell n\in B(y')$, if there is one, (0.2) gives
\begin{equation}\tag{1.2}
\begin{gathered}
P(z,z')=(-1)^n\cdot L^{-3}\cdot U_A(y',z)^{-1}\cdot U_A(y',z'+\ell n)\cdot\mathbf{1}_4\quad\text{if $z\in B(y')$ and $z'+\ell n\in B(y')$};\\
P(z,z')=0\quad\text{if $z\sim z'$ does not hold.}
\end{gathered}
\end{equation}
For $n=0$, that is, for points $z$, $z'$ of $\mathbb{Z}^3$ that lie in one block $B(y')$, this is the display of clause (a):
the display holds identically for such points, and with $z'+\ell n$ in place of $z'$ the kernel acquires the sign $(-1)^n$.
On the torus (Section~G1.2 and the head of Section~G1.11), sites $z$, $z'$ lie in one block exactly when $z\sim z'$.
If $z\sim z'$ does not hold, no translate $z'+\ell n$ lies in the block of $z$, and (1.2) gives $P(z,z')=0$, as asserted in clause (a).
The kernel is a scalar multiple of $\mathbf{1}_4$;
hence $|P(z,z')|_2=|P(z,z')|_c$ is the modulus of this scalar (Section~0.6), which is denoted by $|P(z,z')|$.
Applying (0.4) to the two points $z$ and $z'+\ell n$ of $B(y')$ one obtains
\begin{equation}\tag{1.3}
\begin{aligned}
|P(z,z')|&\le L^{-3}\cdot\exp\left(\frac{3}{2}\,(L-1)\cdot\mathfrak{i}(A)\right)^2\cdot\mathbf{1}[\text{one $L$-block}]\\
&=L^{-3}\cdot\exp\left(3\,(L-1)\cdot\mathfrak{i}(A)\right)\cdot\mathbf{1}[\text{one $L$-block}]\\
&\le\mathrm{e}\cdot L^{-3}\cdot\mathbf{1}[\text{one $L$-block}].
\end{aligned}
\end{equation}
The last assertion of clause (a), which concerns data (SF) with $k=0$, is proved in the last paragraph of Section~6 (clause (a) at $k=0$).

\subsubsection*{2. Clause (b)}

(b1) For $v'\in\mathfrak{S}(L,\ell)$, Section~I.3.3 gives
\[
\left(Q(A)\cdot Q^{T}(-A)\,v'\right)(y')=L^{-3}\sum_{z\in B(y')}U_A(y',z)\cdot U_A(y',z)^{-1}\cdot v'(y')=v'(y'),
\]
because $B(y')$ has $L^3$ points.
Thus $Q(A)\cdot Q^{T}(-A)=1$, and
\[
P^2=Q^{T}(-A)\cdot\left[Q(A)\cdot Q^{T}(-A)\right]\cdot Q(A)=P.
\]
By (1.1) the coefficients of $P$ are scalars;
consequently $P\cdot\mathrm{Mx}=\mathrm{Mx}\cdot P$ for every constant $4\times4$ matrix $\mathrm{Mx}$.

(b2) Let $\mathrm{tb}>0$, and let $y$ be real.
By (I.c1), $\gamma_5^{*}=\gamma_5$ and $\gamma_5^2=1$.
Hence
\[
\left(\mathrm{tb}+i\gamma_5\cdot y\right)^{*}\cdot\left(\mathrm{tb}+i\gamma_5\cdot y\right)=\left(\mathrm{tb}-i\gamma_5\cdot y\right)\cdot\left(\mathrm{tb}+i\gamma_5\cdot y\right)=\mathrm{tb}^2+y^2,
\]
that is, $\left|\left(\mathrm{tb}+i\gamma_5\cdot y\right)\cdot\mathrm{vv}\right|_2=\left(\mathrm{tb}^2+y^2\right)^{\frac{1}{2}}\cdot|\mathrm{vv}|_2$ for all $\mathrm{vv}$.
So $\mathrm{tb}+i\gamma_5\cdot y$ is invertible, and $\left|\left(\mathrm{tb}+i\gamma_5\cdot y\right)^{-1}\right|_2=\left(\mathrm{tb}^2+y^2\right)^{-\frac{1}{2}}$.
Since $b_k\ge b_{-}>0$ and $b_{L}>0$ (Section~0.7), one may put
\begin{equation}\tag{2.1}
\begin{gathered}
\mathrm{rb}_1:=\left(b_k+i\gamma_5\cdot y\right)^{-1},\qquad\mathrm{rb}_2:=\left(b_k+b_{L}+i\gamma_5\cdot y\right)^{-1};\\
\left|\mathrm{rb}_1\right|_2=\left(b_k^2+y^2\right)^{-\frac{1}{2}},\qquad
\left|\mathrm{rb}_2\right|_2=\left(\left(b_k+b_{L}\right)^2+y^2\right)^{-\frac{1}{2}}\le\left(b_k^2+y^2\right)^{-\frac{1}{2}}.
\end{gathered}
\end{equation}

(b3) One has
\[
b_k+b_{L}\cdot P+i\gamma_5\cdot y=\left(b_k+i\gamma_5\cdot y\right)\cdot(1-P)+\left(b_k+b_{L}+i\gamma_5\cdot y\right)\cdot P.
\]
By (b1), $(1-P)^2=1-P$ and $P\cdot(1-P)=(1-P)\cdot P=0$, and $P$ commutes with $\gamma_5$, $\mathrm{rb}_1$ and $\mathrm{rb}_2$.
Therefore
\[
\left[\left(b_k+i\gamma_5\cdot y\right)\cdot(1-P)+\left(b_k+b_{L}+i\gamma_5\cdot y\right)\cdot P\right]\cdot\left[\mathrm{rb}_1\cdot(1-P)+\mathrm{rb}_2\cdot P\right]=(1-P)+P=1,
\]
and, by the same relations,
\[
\left[\mathrm{rb}_1\cdot(1-P)+\mathrm{rb}_2\cdot P\right]\cdot\left[\left(b_k+i\gamma_5\cdot y\right)\cdot(1-P)+\left(b_k+b_{L}+i\gamma_5\cdot y\right)\cdot P\right]=(1-P)+P=1.
\]
In each of the two products both mixed terms contain the factor $(1-P)\cdot P=0$ or $P\cdot(1-P)=0$.
The two remaining terms are $\left(b_k+i\gamma_5\cdot y\right)\cdot\mathrm{rb}_1\cdot(1-P)^2=1-P$ and $\left(b_k+b_{L}+i\gamma_5\cdot y\right)\cdot\mathrm{rb}_2\cdot P^2=P$ in the first product,
and $\mathrm{rb}_1\cdot\left(b_k+i\gamma_5\cdot y\right)\cdot(1-P)^2=1-P$ and $\mathrm{rb}_2\cdot\left(b_k+b_{L}+i\gamma_5\cdot y\right)\cdot P^2=P$ in the second.
Hence $b_k+b_{L}\cdot P+i\gamma_5\cdot y$ is invertible, and the definition in Section~I.3.5 gives
\begin{equation}\tag{2.2}
B_{k,y}(A)=\mathrm{rb}_1\cdot(1-P)+\mathrm{rb}_2\cdot P=\mathrm{rb}_1+\left(\mathrm{rb}_2-\mathrm{rb}_1\right)\cdot P.
\end{equation}

(b4) By (2.2) and (Kn4),
\[
B_{k,y}(A)(z,z')=\mathrm{sz}(z,z')\cdot\mathrm{rb}_1+\left(\mathrm{rb}_2-\mathrm{rb}_1\right)\cdot P(z,z').
\]
The first term vanishes unless $z-z'\in\ell\mathbb{Z}^3$, and the second unless $z\sim z'$ (see (1.2));
by (Lb1) it follows that $B_{k,y}(A)(z,z')=0$ unless $z\sim z'$.
By (2.1) and (1.3),
\begin{equation}\tag{2.3}
\begin{aligned}
\left|B_{k,y}(A)(z,z')\right|_2&\le\left|\mathrm{rb}_1\right|_2\cdot\mathbf{1}[z=z']+\left(\left|\mathrm{rb}_1\right|_2+\left|\mathrm{rb}_2\right|_2\right)\cdot|P(z,z')|\\
&\le\left(b_k^2+y^2\right)^{-\frac{1}{2}}\cdot\left(\mathbf{1}[z=z']+2\,\mathrm{e}\cdot L^{-3}\cdot\mathbf{1}[\text{one $L$-block}]\right).
\end{aligned}
\end{equation}
This contains the first inequality of clause (b).

(b5) One has $2\left(1+y^2\right)-\left(1+|y|\right)^2=\left(1-|y|\right)^2\ge0$.
Further, $1+y^2\le\max\left(1,b_k^{-2}\right)\cdot\left(b_k^2+y^2\right)$:
if $b_k\ge1$, then $1+y^2\le b_k^2+y^2$;
if $b_k<1$, then $b_k^{-2}\cdot\left(b_k^2+y^2\right)=1+y^2/b_k^2\ge1+y^2$.
Hence $\left(1+|y|\right)^2\le2\max\left(1,b_k^{-2}\right)\cdot\left(b_k^2+y^2\right)$, and, since $b_k\ge b_{-}$,
\begin{equation}\tag{2.4}
\left(b_k^2+y^2\right)^{-\frac{1}{2}}\le\sqrt{2}\cdot\max\left(1,\frac{1}{b_k}\right)\cdot\left(1+|y|\right)^{-1}\le\sqrt{2}\cdot\max\left(1,\frac{1}{b_{-}}\right)\cdot\left(1+|y|\right)^{-1}.
\end{equation}
Together with $\mathbf{1}[z=z']+2\,\mathrm{e}\cdot L^{-3}\le1+2\,\mathrm{e}$ and the definition of $C_{\mathrm{By}}$ in Proposition~5, this is the second inequality of clause (b):
\[
\left(b_k^2+y^2\right)^{-\frac{1}{2}}\cdot\left(\mathbf{1}[z=z']+2\,\mathrm{e}\cdot L^{-3}\right)\le C_{\mathrm{By}}\cdot\left(1+|y|\right)^{-1}.
\]

(b6) Fix $z$.
Exactly one point $z'\in\mathrm{fd}$ has $\mathbf{1}[z=z']=1$, and by (Lb3) exactly $L^3$ points $z'\in\mathrm{fd}$ have $\mathbf{1}[\text{one $L$-block}]=1$.
So (2.3) gives
\begin{equation}\tag{2.5}
\sum_{z'\in\mathrm{fd}}\left|B_{k,y}(A)(z,z')\right|_2\le\left(b_k^2+y^2\right)^{-\frac{1}{2}}\cdot\left(1+2\,\mathrm{e}\cdot L^{-3}\cdot L^3\right)=(1+2\,\mathrm{e})\cdot\left(b_k^2+y^2\right)^{-\frac{1}{2}}.
\end{equation}
Now fix $z'$.
Exactly one point $z\in\mathrm{fd}$ has $\mathbf{1}[z=z']=1$, and by (Lb3) exactly $L^3$ points $z\in\mathrm{fd}$ have $\mathbf{1}[\text{one $L$-block}]=1$.
Again by (2.3),
\[
\sum_{z\in\mathrm{fd}}\left|B_{k,y}(A)(z,z')\right|_2\le\left(b_k^2+y^2\right)^{-\frac{1}{2}}\cdot\left(1+2\,\mathrm{e}\cdot L^{-3}\cdot L^3\right)=(1+2\,\mathrm{e})\cdot\left(b_k^2+y^2\right)^{-\frac{1}{2}}.
\]

\subsubsection*{3. Clause (c)}

Let $O$ be an operator on $\mathfrak{S}^{\mathrm{ch}}(\eta_k,\ell)$;
then $\mathrm{Qk}\cdot O\cdot\mathrm{Qt}$ is an operator on $\mathfrak{S}(1,\ell)$.
By (0.2), (Qs2) and (Qs1),
\begin{align*}
\left(\mathrm{Qk}\cdot O\cdot\mathrm{Qt}\right)(z,z')\cdot\mathrm{vv}&=\left(\mathrm{Qk}\cdot O\cdot\mathrm{Qt}\,\mathrm{dsp}(z',\mathrm{vv})\right)(z)=\left(\mathrm{Qk}\cdot\left[\mathbf{1}_{\Delta_z}\cdot O\cdot\mathbf{1}_{\Delta_{z'}}\right]\cdot\mathrm{Qt}\,\mathrm{dsp}(z',\mathrm{vv})\right)(z),\\
\left|\left(\mathrm{Qk}\cdot O\cdot\mathrm{Qt}\right)(z,z')\cdot\mathrm{vv}\right|_2
&\le\exp\left(\frac{3}{2}\,\mathfrak{i}(A)\right)\cdot\left|\left[\mathbf{1}_{\Delta_z}\cdot O\cdot\mathbf{1}_{\Delta_{z'}}\right]\cdot\mathrm{Qt}\,\mathrm{dsp}(z',\mathrm{vv})\right|_{\mathrm{su}}\\
&\le\exp\left(\frac{3}{2}\,\mathfrak{i}(A)\right)\cdot\left\|\mathbf{1}_{\Delta_z}\cdot O\cdot\mathbf{1}_{\Delta_{z'}}\right\|_{\mathrm{su}}\cdot\left|\mathrm{Qt}\,\mathrm{dsp}(z',\mathrm{vv})\right|_{\mathrm{su}}\\
&\le\exp\left(3\,\mathfrak{i}(A)\right)\cdot\left\|\mathbf{1}_{\Delta_z}\cdot O\cdot\mathbf{1}_{\Delta_{z'}}\right\|_{\mathrm{su}}\cdot|\mathrm{vv}|_2.
\end{align*}
By (0.4), $\exp\left(3\,\mathfrak{i}(A)\right)\le\mathrm{e}$.
Taking the supremum over all $\mathrm{vv}$ with $|\mathrm{vv}|_2\le1$ one arrives at clause (c).

\subsubsection*{4. Clause (d)}

For $y\in Y_k$ the operator $S_y(A)=F_y(A)^{-1}$ on $\mathfrak{S}^{\mathrm{ch}}(\eta_k,\ell)$ exists (Proposition~2.8\,(a)),
and $w_k(y)=1$ (Section~I.3.5, $k\ge1$).
The norm $N_{\gamma^{F}}$ is formed on the block torus $T_k$ (Section~G1.3).

(d1) For $y\in Y_k$ put $\mathrm{Ku}_y:=\mathrm{Qk}\cdot S_y(A)\cdot\mathrm{Qt}$.
By Proposition~2.8\,(c),
\[
\left\|\mathbf{1}_{\Delta_z}\cdot S_y(A)\cdot\mathbf{1}_{\Delta_{z'}}\right\|_{\mathrm{su}}\le32\,a'\cdot\exp\left(-\gamma^{F}\cdot d(z,z')\right).
\]
Clause (c) with $O=S_y(A)$ therefore gives
\begin{equation}\tag{4.1}
\left|\mathrm{Ku}_y(z,z')\right|_2\le32\,\mathrm{e}\cdot a'\cdot\exp\left(-\gamma^{F}\cdot d(z,z')\right)\qquad(y\in Y_k).
\end{equation}
By Section~I.3.5, $M_k(A)=b_k-D_k(A)=b_k^2\cdot\mathrm{Qk}\cdot S_k(A)\cdot\mathrm{Qt}=b_k^2\cdot\mathrm{Ku}_\infty$, because $S_k(A)=S_\infty(A)$.
Since $0<b_k\le b$ (Section~0.7), it follows that
\[
\left|M_k(A)(z,z')\right|_2\le32\,\mathrm{e}\cdot b^2\cdot a'\cdot\exp\left(-\gamma^{F}\cdot d(z,z')\right).
\]

(d2) By Section~I.3.5, $\dot{M}_k(A)=b_k^2\cdot\mathrm{Qk}\cdot S_k(A)^2\cdot\mathrm{Qt}$.
By (N2) and Proposition~2.8\,(c) at $y=\infty$,
\[
N_{\gamma^{F}}\left(S_k(A)^2\right)\le N_{\gamma^{F}}\left(S_k(A)\right)^2\le(32\,a')^2,
\]
and (N1) gives $\left\|\mathbf{1}_{\Delta_z}\cdot S_k(A)^2\cdot\mathbf{1}_{\Delta_{z'}}\right\|_{\mathrm{su}}\le(32\,a')^2\cdot\exp\left(-\gamma^{F}\cdot d(z,z')\right)$;
here (N1), (N2) are the rules of Section~G1.3, see also Section~0.8.
Clause (c) with $O=S_k(A)^2$, together with $b_k\le b$, yields
\[
\left|\dot{M}_k(A)(z,z')\right|_2\le\mathrm{e}\cdot b^2\cdot(32\,a')^2\cdot\exp\left(-\gamma^{F}\cdot d(z,z')\right).
\]

(d3) \emph{Products with $B_{k,y}(A)$ on both sides.}
Let $\mathrm{Ku}$ be an operator on $\mathfrak{S}(1,\ell)$, and let $C_{\mathrm{Ku}}\ge0$ be a number such that
$\left|\mathrm{Ku}(z_1,z_2)\right|_2\le C_{\mathrm{Ku}}\cdot\exp\left(-\gamma^{F}\cdot d(z_1,z_2)\right)$ for all $z_1$, $z_2$.
Then, for every real $y$ and all $z$, $z'$,
\begin{equation}\tag{4.2}
\left|\left(B_{k,y}(A)\cdot\mathrm{Ku}\cdot B_{k,y}(A)\right)(z,z')\right|_2\le4\,(1+2\,\mathrm{e})^2\cdot\left(b_k^2+y^2\right)^{-1}\cdot C_{\mathrm{Ku}}\cdot\exp\left(-\gamma^{F}\cdot d(z,z')\right).
\end{equation}
\begin{proof}[Proof of (4.2)]
Two applications of (Kn3) give
\[
\left(B_{k,y}(A)\cdot\mathrm{Ku}\cdot B_{k,y}(A)\right)(z,z')=\sum_{z_1\in\mathrm{fd}}\sum_{z_2\in\mathrm{fd}}B_{k,y}(A)(z,z_1)\cdot\mathrm{Ku}(z_1,z_2)\cdot B_{k,y}(A)(z_2,z').
\]
By (b4) a term vanishes unless $z\sim z_1$ and $z_2\sim z'$.
For such $z_1$, $z_2$, (Lb4) and the triangle inequality give
\[
d(z,z')\le d(z,z_1)+d(z_1,z_2)+d(z_2,z')\le2\,(L-1)+d(z_1,z_2),
\]
and hence, by (0.7),
\[
\exp\left(-\gamma^{F}\cdot d(z_1,z_2)\right)\le\exp\left(2\gamma^{F}(L-1)\right)\cdot\exp\left(-\gamma^{F}\cdot d(z,z')\right)\le4\exp\left(-\gamma^{F}\cdot d(z,z')\right).
\]
Since $|\cdot|_2$ is submultiplicative, it follows that
\begin{multline*}
\left|\left(B_{k,y}(A)\cdot\mathrm{Ku}\cdot B_{k,y}(A)\right)(z,z')\right|_2\\
\le4\,C_{\mathrm{Ku}}\cdot\exp\left(-\gamma^{F}\cdot d(z,z')\right)\cdot\left[\sum_{z_1\in\mathrm{fd}}\left|B_{k,y}(A)(z,z_1)\right|_2\right]\cdot\left[\sum_{z_2\in\mathrm{fd}}\left|B_{k,y}(A)(z_2,z')\right|_2\right].
\end{multline*}
By (b6) the first bracket, a row sum, and the second bracket, a column sum, are each at most $(1+2\,\mathrm{e})\cdot\left(b_k^2+y^2\right)^{-\frac{1}{2}}$.
\end{proof}

(d4) Let $y$ be real.
By Proposition~4\,(b) the operator $\Gamma_{k,y}(A)=\left(D_k(A)+b_{L}\cdot P(A)+i\gamma_5\cdot y\right)^{-1}$ exists, and
\[
\Gamma_{k,y}(A)-B_{k,y}(A)=b_k^2\cdot B_{k,y}(A)\cdot\mathrm{Qk}\cdot S_y(A)\cdot\mathrm{Qt}\cdot B_{k,y}(A)=b_k^2\cdot B_{k,y}(A)\cdot\mathrm{Ku}_y\cdot B_{k,y}(A).
\]
By (4.1) the hypothesis of (d3) is satisfied by $\mathrm{Ku}=\mathrm{Ku}_y$ with $C_{\mathrm{Ku}}=32\,\mathrm{e}\cdot a'$, and (4.2) gives
\begin{align*}
\left|\left(\Gamma_{k,y}(A)-B_{k,y}(A)\right)(z,z')\right|_2
&\le b_k^2\cdot4\,(1+2\,\mathrm{e})^2\cdot\left(b_k^2+y^2\right)^{-1}\cdot32\,\mathrm{e}\cdot a'\cdot\exp\left(-\gamma^{F}\cdot d(z,z')\right)\\
&=128\,\mathrm{e}\cdot(1+2\,\mathrm{e})^2\cdot a'\cdot b_k^2\cdot\left(b_k^2+y^2\right)^{-1}\cdot\exp\left(-\gamma^{F}\cdot d(z,z')\right).
\end{align*}

(d5) By Section~I.3.5, $\Gamma_k(A)=\Gamma_{k,0}(A)$ and $B_k(A)=B_{k,0}(A)$.
From (2.3) at $y=0$, together with $b_k\ge b_{-}$, with $\mathbf{1}[z=z']\le\mathbf{1}[\text{one $L$-block}]$ (Lb1) and with $L^{-3}\le1$, and from (Lb4) and (0.7), one obtains
\begin{equation}\tag{4.3}
\begin{gathered}
\left|B_k(A)(z,z')\right|_2\le\frac{1+2\,\mathrm{e}}{b_{-}}\cdot\mathbf{1}[\text{one $L$-block}],\\
\mathbf{1}[\text{one $L$-block}]\le\exp\left(\gamma^{F}(L-1)\right)\cdot\exp\left(-\gamma^{F}\cdot d(z,z')\right)\le2\exp\left(-\gamma^{F}\cdot d(z,z')\right).
\end{gathered}
\end{equation}
By (d4) at $y=0$, where $b_k^2\cdot\left(b_k^2+y^2\right)^{-1}=1$, and by the definition of $C_{\Gamma}$ in Proposition~5,
\begin{align*}
\left|\Gamma_k(A)(z,z')\right|_2&\le\left|B_k(A)(z,z')\right|_2+\left|\left(\Gamma_{k,0}(A)-B_{k,0}(A)\right)(z,z')\right|_2\\
&\le\left[\frac{2\,(1+2\,\mathrm{e})}{b_{-}}+128\,\mathrm{e}\cdot(1+2\,\mathrm{e})^2\cdot a'\right]\cdot\exp\left(-\gamma^{F}\cdot d(z,z')\right)
=C_{\Gamma}\cdot\exp\left(-\gamma^{F}\cdot d(z,z')\right).
\end{align*}

\subsubsection*{5. Clause (e)}

(e1) \emph{The bounds in $|\cdot|_c$.}
By Section~0.6 one has $|\mathrm{Mx}|_c\le 2\,|\mathrm{Mx}|_2$ for each of the $4\times4$ blocks $\mathrm{Mx}$ that occur in (b), (c) and (d),
and $|P(A)(z,z')|_c=|P(A)(z,z')|$ is the modulus of the number in front of $\mathbf{1}_4$.
By (Lb1),
\[
\mathbf{1}[z=z']+2\,\mathrm{e}\cdot L^{-3}\cdot\mathbf{1}[\text{one $L$-block}]
=\left(\mathbf{1}[z=z']+2\,\mathrm{e}\cdot L^{-3}\right)\cdot\mathbf{1}[\text{one $L$-block}].
\]
Therefore (1.3), (2.3), (2.4) and (2.5) give, for every real $y$ and all unit sites $z$, $z'$,
\begin{align*}
\left|P(A)(z,z')\right|_c&\le\mathrm{e}\cdot L^{-3}\le 2\,\mathrm{e}\cdot L^{-3},\\
\left|B_{k,y}(A)(z,z')\right|_c&\le 2\left(b_k^2+y^2\right)^{-\frac{1}{2}}\cdot\left(\mathbf{1}[z=z']+2\,\mathrm{e}\cdot L^{-3}\right)\cdot\mathbf{1}[\text{one $L$-block}]\\
&\le 2\,C_{\mathrm{By}}\cdot(1+|y|)^{-1}\cdot\mathbf{1}[\text{one $L$-block}],
\end{align*}
and the row sums and the column sums of $\left|B_{k,y}(A)(z,z')\right|_c$ are at most $2\,(1+2\,\mathrm{e})\left(b_k^2+y^2\right)^{-\frac{1}{2}}$.
In the same way (c) and (d) give, for every operator $O$ on $\mathfrak{S}^{\mathrm{ch}}(\eta_k,\ell)$, every real $y$ and all $z$, $z'$,
\begin{align*}
\left|\left(\mathrm{Qk}\cdot O\cdot\mathrm{Qt}\right)(z,z')\right|_c&\le 2\,\mathrm{e}\cdot\left\|\mathbf{1}_{\Delta_z}\cdot O\cdot\mathbf{1}_{\Delta_{z'}}\right\|_{\mathrm{su}},\\
\left|M_k(A)(z,z')\right|_c&\le 64\,\mathrm{e}\cdot b^2\cdot a'\cdot\exp\left(-\gamma^{F}\cdot d(z,z')\right),\\
\left|\dot{M}_k(A)(z,z')\right|_c&\le 2\,\mathrm{e}\cdot b^2\cdot(32\,a')^2\cdot\exp\left(-\gamma^{F}\cdot d(z,z')\right),\\
\left|\left(\Gamma_{k,y}(A)-B_{k,y}(A)\right)(z,z')\right|_c&\le 256\,\mathrm{e}\cdot(1+2\,\mathrm{e})^2\cdot a'\\
&\qquad\cdot b_k^2\cdot\left(b_k^2+y^2\right)^{-1}\cdot\exp\left(-\gamma^{F}\cdot d(z,z')\right),\\
\left|\Gamma_k(A)(z,z')\right|_c&\le 2\,C_{\Gamma}\cdot\exp\left(-\gamma^{F}\cdot d(z,z')\right).
\end{align*}

(e2) \emph{Comparison with $C^{F}$.}
By Proposition~14.2\,(a), line (C6),
\[
C^{F}=2\max\left\{C_{1},\ C_{\mathrm{By}},\ 304\,\mathrm{e}\cdot\max(1,b)\cdot a',\ C_{\Gamma},\ \mathrm{e}\cdot b^2\cdot(32\,a')^2\right\},
\]
where
\begin{multline*}
C_{1}=\max\left\{304\,\mathrm{e}\cdot\max(1,b)\cdot a'\cdot M^{27\alpha_{0}},\ C_{\Gamma}\cdot M^{27\alpha_{0}},\right.\\
\left.32\,\mathrm{e}\cdot\max\left(1,b^2\right)\cdot a'\cdot\exp\left(27\kappa'\right)\right\}.
\end{multline*}
The doubled constants of (e1) are the eight numbers
\begin{gather*}
2\,\mathrm{e}\cdot L^{-3},\qquad 2\,C_{\mathrm{By}},\qquad \frac{2\,(1+2\,\mathrm{e})}{b_{-}},\qquad 2\,\mathrm{e},\\
64\,\mathrm{e}\cdot b^2\cdot a',\qquad 2\,\mathrm{e}\cdot b^2\cdot(32\,a')^2,\qquad 256\,\mathrm{e}\cdot(1+2\,\mathrm{e})^2\cdot a',\qquad 2\,C_{\Gamma}.
\end{gather*}
Here $2\,(1+2\,\mathrm{e})/b_{-}$ bounds $2\,(1+2\,\mathrm{e})\left(b_k^2+y^2\right)^{-\frac{1}{2}}$, because $b_k\ge b_{-}$;
and $256\,\mathrm{e}\cdot(1+2\,\mathrm{e})^2\cdot a'$ stands with the factor $b_k^2\cdot\left(b_k^2+y^2\right)^{-1}\le 1$.
Now
\[
\mathrm{e}\cdot L^{-3}\le\mathrm{e}\le 1+2\,\mathrm{e}\le C_{\mathrm{By}}
\qquad\text{and}\qquad
\frac{1+2\,\mathrm{e}}{b_{-}}\le C_{\mathrm{By}},
\]
because $\sqrt{2}\cdot\max\left(1,1/b_{-}\right)\ge\max\left(1,1/b_{-}\right)$;
\[
32\,\mathrm{e}\cdot b^2\cdot a'\le 32\,\mathrm{e}\cdot\max\left(1,b^2\right)\cdot a'\cdot\exp\left(27\kappa'\right)\le C_{1},
\]
because $a'>0$ and $\kappa'=\kappa_{\mathrm{w}}+\kappa_{0}>0$;
\[
128\,\mathrm{e}\cdot(1+2\,\mathrm{e})^2\cdot a'\le C_{\Gamma},
\]
because $(1+2\,\mathrm{e})/b_{-}>0$.
Moreover $C_{\mathrm{By}}$, $C_{\Gamma}$, $\mathrm{e}\cdot b^2\cdot(32\,a')^2$ and $C_{1}$ are entries of the maximum that defines $C^{F}$.
Hence each of the eight numbers is at most $C^{F}$.

(e3) \emph{The label FE3$^{\mathrm{ch}}$.}
From (e1) and (e2),
\begin{gather*}
\left|B_{k,y}(A)(z,z')\right|_c\le C^{F}\cdot(1+|y|)^{-1}\cdot\mathbf{1}[\text{one $L$-block}],\\
\left|M_k(A)(z,z')\right|_c,\ \left|\dot{M}_k(A)(z,z')\right|_c,\ \left|\Gamma_k(A)(z,z')\right|_c\le C^{F}\cdot\exp\left(-\gamma^{F}\cdot d(z,z')\right).
\end{gather*}
By Section~I.3.5, $D_k(A)=b_k-M_k(A)$.
By (Kn4), $D_k(A)(z,z')=\mathrm{sz}(z,z')\cdot b_k\mathbf{1}_4-M_k(A)(z,z')$, and $|\mathrm{sz}(z,z')|=\mathbf{1}[z=z']$.
Since $|b_k\mathbf{1}_4|_c=b_k\le b_{+}$ (Sections~0.6 and~0.7), it follows that
\[
\left|D_k(A)(z,z')\right|_c\le b_{+}\cdot\mathbf{1}[z=z']+C^{F}\cdot\exp\left(-\gamma^{F}\cdot d(z,z')\right).
\]

(e4) \emph{Clause (c) with the operator norm for $|\cdot|_\infty$ on the right side.}
Let $\|\cdot\|_\infty$ denote the operator norm for $|\cdot|_\infty$; it is the norm $|\cdot|_{1,\infty}$ of (I.ke3).
Let $O$ be an operator on $\mathfrak{S}^{\mathrm{ch}}(\eta_k,\ell)$ and $\mathrm{vv}\in\mathbb{C}^4$.
By the first line of the display of Section~3 and by (0.5$'$),
\begin{align*}
&\left|\left(\mathrm{Qk}\cdot O\cdot\mathrm{Qt}\right)(z,z')\cdot\mathrm{vv}\right|_c\\
&\qquad\le\exp\left(\frac{3}{2}\,\mathfrak{i}(A)\right)\cdot\left|\left[\mathbf{1}_{\Delta_z}\cdot O\cdot\mathbf{1}_{\Delta_{z'}}\right]\cdot\mathrm{Qt}\,\mathrm{dsp}(z',\mathrm{vv})\right|_\infty\\
&\qquad\le\exp\left(\frac{3}{2}\,\mathfrak{i}(A)\right)\cdot\left\|\mathbf{1}_{\Delta_z}\cdot O\cdot\mathbf{1}_{\Delta_{z'}}\right\|_\infty\cdot\left|\mathrm{Qt}\,\mathrm{dsp}(z',\mathrm{vv})\right|_\infty\\
&\qquad\le\exp\left(3\,\mathfrak{i}(A)\right)\cdot\left\|\mathbf{1}_{\Delta_z}\cdot O\cdot\mathbf{1}_{\Delta_{z'}}\right\|_\infty\cdot|\mathrm{vv}|_c.
\end{align*}
With (0.4), and taking the supremum over $|\mathrm{vv}|_c\le 1$, one obtains
\[
\left|\left(\mathrm{Qk}\cdot O\cdot\mathrm{Qt}\right)(z,z')\right|_c\le\mathrm{e}\cdot\left\|\mathbf{1}_{\Delta_z}\cdot O\cdot\mathbf{1}_{\Delta_{z'}}\right\|_\infty .
\]
By (e2), $\mathrm{e}\le 2\,\mathrm{e}\le C^{F}$.

\subsubsection*{6. Uniformity in the data; level 0}

The numbers that occur in the bounds of Proposition~5 are
\begin{gather*}
\mathrm{e},\qquad L^{-3},\qquad C_{\mathrm{By}},\qquad C_{\Gamma},\qquad 32\,\mathrm{e}\cdot b^2\cdot a',\\
\mathrm{e}\cdot b^2\cdot(32\,a')^2,\qquad 128\,\mathrm{e}\cdot(1+2\,\mathrm{e})^2\cdot a',\qquad \gamma^{F},\qquad C^{F},\qquad b_{+}
\end{gather*}
and, in the last sentence of (a), $b_{L}$ and $b$.
By Proposition~14.2\,(a) they are functions of $b$, $L$, $a'$, $\gamma^{F}$ (lines (C0), (C3), (C4), (C4r)) and, in the case of $C^{F}$, also of $M$, $\alpha_{0}$, $\kappa'$ (lines (C4), (C5), (C6));
thus none of them depends on $k$, $\ell$, $\mu$, $e_k$, $A$, $y$, $z$, $z'$ or $O$.
The number $b_k$ occurs in (b) and in the third bound of (d) as the function of $(b,L,k)$ defined in Section~I.3.4;
wherever a constant is asserted, $b_k$ has been replaced by means of $b_{-}\le b_k\le b_{+}=b$, which holds for all $k\ge1$ (Section~0.7).

The level $k$ enters through $b_k$;
through the number $48\,\eta_k^{-3}$ of the terms of the mean (0.5);
through the phases $\varphi_A(z,c)$ and $U_A(y',z)$, where a path has $|c-c'|_1/\eta_k$ bonds, each with the factor $\eta_k$, and the $k$ links of the chain of centres add up to
\[
\frac{3}{2}\left(L^k-1\right)\cdot\eta_k\cdot\mathfrak{i}(A)=\frac{3}{2}\left(1-\eta_k\right)\cdot\mathfrak{i}(A)\le\frac{3}{2}\,\mathfrak{i}(A),
\]
so that the bounds $\exp\left(\frac{3}{2}\,\mathfrak{i}(A)\right)$ and $\exp\left(\frac{3}{2}\,(L-1)\cdot\mathfrak{i}(A)\right)$ of Section~0.4 do not contain $k$;
through $w_k(y)=1$;
and through Propositions~2.8 and~4, whose constants do not depend on $k$.

The side $\ell$ appears through the relations $\ell\in L\mathbb{Z}$ and $L-1<\ell$ (see (0.1)), through $\ell\ge 3L$, and, in Propositions~2.8 and~4, through $\ell\in\mathrm{Cl}$.
Here $\ell\ge 13\,L\cdot M\ge 3L$ by Sections~I.1.1 and~I.0.5;
together with $\ell\in L\mathbb{Z}$ this means that $T_k=(0,\ell,-1)$ is a block torus (Section~G1.2).
The bounds of \cite[Lemma~16.2(a)]{QEDI}, applied in Section~0.4, are stated for block tori,
and so are the definitions of Section~G1.2 (kernels, used in (Kn2) and in Section~1)
and of Section~G1.3 (the localisations $\mathbf{1}_{\Delta_z}$, used in Section~0.5, and the norms $N_{\gamma}$, used in Sections~0.8 and~4).

In Sections~0--5 the mass $\mu$ enters only by way of the operators $S_{y}(A)$, $y\in Y_{k}$, of Section~I.3.5;
in particular $M_k(A)$, $\dot{M}_k(A)$, $D_k(A)$, $\Gamma_{k,y}(A)$ and $\Gamma_k(A)$ depend on $\mu$, but only through the operators $S_{y}(A)$ from which they are formed.
Uniformity in $\mu$ therefore comes from Proposition~2.8\,(c), whose bounds hold for all $\mu\in(0,1]$.
The charge $e_k$ enters through the inequality $e_k\le e_0$, in (0.3), and through Propositions~2.8 and~4.
The bound (4.1) holds for all $y\in Y_{k}$ with one constant,
while in (b) and in the third bound of (d) the dependence on $y$ is the explicit one.
It follows that the bounds hold, with the same constants, for every quadruple $(k,\ell,\mu,e_k)$ that satisfies (SF) with $k\ge1$,
whatever the torus, the cutoff $N$ and the charge $e$ from which the quadruple is formed (Sections~I.1.8 and~I.6.1).

Averages over the $48$ sites of a cluster enter only through the operators of (0.5):
$\mathrm{Qk}$ contains the mean $\frac{1}{48}\sum_{g\in O_{\mathrm{cub}}}$, and $\mathrm{Qt}\,v$ is constant on each cluster for every $v\in\mathfrak{S}(1,\ell)$.
Both operators have norm at most $\exp\left(\frac{3}{2}\,\mathfrak{i}(A)\right)$.
All kernels of the statement at the levels $k\ge1$ are $4\times4$ blocks at pairs of unit sites, with the weight $1$ of Section~I.1.6.

At level $0$ the objects $b_0$, $M_0(A)$, $B_{0,y}(A)$, $\Gamma_{0,y}(A)$ are not defined (Section~I.3.9, (I.L0.9)), and the unit lattice carries $192$ components per site.
For $k=0$ Proposition~5 asserts only the last sentence of (a), which is proved in the next paragraph.
The bounds of FE3$^{\mathrm{ch}}$ at level $0$ are bounds of the kernels (I.1.4) of $S_{0,y}(A)$ and $\Gamma_0(A)$, namely Proposition~2.8\,(d).

\emph{Clause (a) at $k=0$ ($P(A)$ at level $0$).}
Sections~0.1--0.4, 0.6 and~1 do not use $k\ge1$:
Sections~0.1, 0.2 and~0.6 contain neither $k$ nor $A$;
Section~0.3 uses $e_k\le e_0$ and $A\in\tilde{G}_k$;
Sections~0.4 and~1 contain $k$ only through $\eta_k$
(at $k=0$ one has $\eta_0=1$ and $\Delta_z=\{z\}$, the chain of centres has no link, and $\varphi_A(z,z)=1$).
Hence, for data (SF) with $k=0$ and $A\in\tilde{G}_0$, the kernel of the operator $P(A)$ on $\mathfrak{S}(1,\ell)$ (Section~I.3.3) is given by (1.2) and satisfies (1.3);
that is, (a) holds for these data.
At $k=0$ one has $\mathbb{L}_0=\mathbb{Z}^3$.
Let $c$, $c'$ be points of $\mathbb{Z}^3$, $x=(c,g)$, $x'=(c',g')$, and $f\in\mathfrak{S}^{\mathrm{ch}}(1,\ell)$.
By (I.3.1), by (Kn1), and by (I.1.3) with $\eta_0=1$,
\begin{align*}
\left(JP(A)J^{*}f\right)(c,g)&=\left(P(A)J^{*}f\right)(c)
=\sum_{c'\in\mathrm{fd}}P(A)(c,c')\,\frac{1}{48}\sum_{g'\in O_{\mathrm{cub}}}f(c',g')\\
&=\int\mathrm{d}x'\,P(A)(c,c')\,f(x').
\end{align*}
Together with (Kn2) this says that the kernel (I.1.4) of the operator $JP(A)J^{*}$ of (I.3.4) at $(x,x')$ is $P(A)(c,c')$; this is (I.J5).
This kernel does not depend on $g$, $g'$, and it is a scalar multiple $\mathrm{wz}\,\mathbf{1}_4$ of $\mathbf{1}_4$,
for which $\sum_{\beta'}\left|\left(\mathrm{wz}\,\mathbf{1}_4\right)_{\beta\beta'}\right|=|\mathrm{wz}|$ for every $\beta$,
that is, $\left|\mathrm{wz}\,\mathbf{1}_4\right|_c=|\mathrm{wz}|$ (Section~0.6).
By (1.3) and (Lb3),
\begin{align*}
\sup_x\int\mathrm{d}x'\,\left|b_{L}\left(JP(A)J^{*}\right)(x,x')\right|_c
&=b_{L}\cdot\sup_c\sum_{c'\in\mathrm{fd}}\frac{1}{48}\sum_{g'\in O_{\mathrm{cub}}}\left|P(A)(c,c')\right|\\
&=b_{L}\cdot\sup_c\sum_{c'\in\mathrm{fd}}\left|P(A)(c,c')\right|\\
&\le b_{L}\cdot L^3\cdot\mathrm{e}\cdot L^{-3}=b_{L}\cdot\mathrm{e}\le b\cdot\mathrm{e}.
\end{align*}
This is the last term of (I.L0.12).
\qed

\subsection*{Conditions on the constants}

Besides the conditions under which Propositions~2.8 and~4 hold, the proof of Proposition~5 rests on the following properties of the constants.

\medskip
\noindent
\begin{tabular}{p{0.29\textwidth}p{0.29\textwidth}p{0.29\textwidth}}
Condition & Used in & Where it is stated\\
\hline
(Z2.s1): $6L\cdot\exp\left(-\varepsilon\cdot v\right)\le 1$ for all $v\ge-\log e_0$
& (0.3), hence (0.4), (1.3), (2.3), Section~3, (e4), Section~6 (clause (a) at $k=0$)
& Proposition~14.2\,(b), row (Z2.s1); a condition on $e_0$, line (C8) of Section~I.8.2\\
\hline
$L$ odd
& (0.1); Section~0.4 (chain of centres); Section~0.5 ($L^k$ odd)
& Proposition~14.2\,(b), (L5); line (C3); range of $L$ in Section~I.0.5\\
\hline
$b>0$ and $L>1$ (so that $b_{-}=b\left(1-1/L\right)>0$ and $b_{L}=b/L>0$)
& Section~0.7; (b2), (2.1), (2.4), (4.3); finiteness of $C_{\mathrm{By}}$ and $C_{\Gamma}$; Section~6, clause (a) at $k=0$ ($b_{L}\le b$)
& $b=1$: line (C0); (L5): $L\ge5$; $b_{-}$, $b_{L}$: line (C4)\\
\hline
$M$ an integer; $(L-1)/2\le M$
& Section~0.1 ($\ell\in L\mathbb{Z}$); Section~0.4 ($U_A(y',z)$ is defined)
& (M1): $M$ a power of $L$, $M\ge L^2$; line (C5); range of $M$ in Section~I.0.5\\
\hline
$0<\gamma^{F}\cdot L\le\log 2$
& (0.7), hence (4.2), (4.3)
& a consequence of the definitions $a_{z}\le 1/(3L)$, $\gamma_{z}\le a_{z}$, $\gamma^{F}=\gamma_{z}/2$ of line (C4r) (Section~0.7); stated in Proposition~14.2\,(f)\\
\hline
$\kappa'>0$; $a'>0$
& (e2)
& consequences of $\kappa'=\kappa_{\mathrm{w}}+\kappa_{0}=\kappa+2\kappa_{0}+3$, lines (C0), (C2), (C4), and of $a'=C_z\cdot c(\gamma^{F})\ge1$, line (C4r), Proposition~14.2\,(f)\\
\hline
$b_{-}\le b_k\le b_{+}=b$
& (2.1), (2.4), (d1), (d2), (4.3), (e2), (e3)
& a consequence of the definition of $b_k$ in Section~I.3.4 (Section~0.7)\\
\hline
\end{tabular}

\medskip
Consequently Proposition~5 holds whenever the constants have the properties of the table and satisfy (L5), (M1), (M2), $M\ge 108\,\pi\cdot C'_{V}\cdot a'$ (part of (M5)), (Z2.s1), (Z2.s2) and $\alpha<1$;
indeed, Proposition~2.8 holds under these conditions, by way of Proposition~2.7, and so does Proposition~4, which rests on Proposition~2.8\,(a).
Under them the side $\ell\in\mathrm{Cl}$ of the data satisfies $\ell\in L\mathbb{Z}$, $L-1<\ell$ and $\ell\ge 3L$ (see (0.1) and Section~6), so that $T_k$ is a block torus (Section~G1.2).

\section*{Chapter 6. The field $\psi=H_k(A)\Psi$ and all its H\"older quotients (FE4$^{\mathrm{ch}}$): proof of Proposition~6}

\subsection*{Proof}

\subsubsection*{0. Conventions of this chapter; auxiliary facts}

Throughout this chapter $(k,\ell,\mu,e_k)$ are data (SF) in the sense of Section~I.1.8, and $A\in\tilde{G}_k$ (Section~I.5.3).
The constants are those of Section~I.8, as fixed in Proposition~14.2.
On $\mathbb{C}^4$ a modulus $|\cdot|$ without index stands for either of the norms $|\cdot|_2$, $|\cdot|_{c}$, in an inequality that is valid for both;
accordingly, $\sup|f|$ then means $|f|_{\mathrm{su}}$, respectively $|f|_{\infty}$ (Sections~G1.3, I.1.5).
As in Section~I.1.1, the letters $y$, $y'$ denote unit sites, that is, points of $\mathbb{Z}^3$;
the letter $y$ denotes the resolvent parameter of Section~I.3.5 only in the expression $y=\infty$, as an element of $Y_k$, and as index or argument in $F_y$, $S_y$, $w_k(y)$.

In this chapter $t$ denotes a row index, as in (I.1.4) and Section~G1.11;
$m$ denotes a $4\times4$ block, as in Section~I.1.5;
$n$ denotes a point of $\mathbb{Z}^3$, as in Sections~G1.3 and~G1.6;
$F$ denotes a function of pairs in the sense of Section~I.1.7;
$\mathrm{Mx}$ denotes a constant $4\times4$ matrix, $\mathrm{vv}$ an element of $\mathbb{C}^4$, and $\mathrm{cf}$ a number.
The point $\tilde{z}'$ and the spinor $\mathrm{vz}$ are defined in Section~0.2, the set $\mathrm{Dfu}$ in Section~0.8;
$\mathrm{psf}$ denotes a field that is defined anew in each of Sections~1, 2.2 and~3;
the labels (pin), (psh), (plg) of the three kinds of pairs are introduced in Section~0.3.
All other letters have the meaning given to them in Proposition~6, in Part~I and in Sections~G1.1--G1.11.

\paragraph{0.1. Functions on the covering lattices; localisation.}
The connection $A$ is identified with its $\ell\mathbb{Z}^3$-periodic extension to the bonds of $\mathbb{L}_k$.
The elements of $\mathfrak{S}(1,\ell)$ are functions on $\mathbb{Z}^3$, and the elements of $\mathfrak{S}^{\mathrm{ch}}(\eta_k,\ell)$ are functions on $\mathbb{L}^{\mathrm{f}}_k=\mathbb{L}_k\times O_{\mathrm{cub}}$;
both change sign under the translation by $\ell e_{\mu}$ (Section~I.1.5).
A function of pairs is a function $F$ on $\mathrm{Pzt}_k$ with $F(x+\ell e_{\mu},x'+\ell e_{\mu})=-F(x,x')$ (Section~I.1.7).
Consequently the moduli $|v(y)|$, $|f(x)|$, $|F(x,x')|$ are $\ell\mathbb{Z}^3$-periodic,
and each supremum below has the same value over the covering index set as over the torus
($\mathbb{Z}^3$ and $T^0_k$; $\mathbb{L}^{\mathrm{f}}_k$ and $T^{\mathrm{f}}_k$; $\mathrm{Pzt}_k$ and $\mathrm{Pz}_k$).

A site $z$ of $T^0_k$ is a class of $\mathbb{Z}^3$ modulo $\ell\mathbb{Z}^3$.
For $c\in\mathbb{L}_k$ let $[c]\in\mathbb{Z}^3$ be the unit site with $c\in\Delta_{[c]}$;
the blocks $\Delta_y$, $y\in\mathbb{Z}^3$, form a partition of $\mathbb{L}_k$ (Sections~I.1.1, I.1.4).
A fermion site $x=(c,g)$ is said to lie over $\Delta_z$ if $[c]$ belongs to the class $z$ (I.POS).
By Section~G1.3 the localisations act as follows.
\begin{itemize}
\item The operator $\mathbf{1}_{z}$ multiplies $v\in\mathfrak{S}(1,\ell)$ by the indicator of the class $z$.
\item The operator $\mathbf{1}_{\Delta_z}$ multiplies $f\in\mathfrak{S}^{\mathrm{ch}}(\eta_k,\ell)$ by the indicator of the set $\left\{(c,g):[c]\in z,\ g\in O_{\mathrm{cub}}\right\}$, which contains all $48$ labels,
and it multiplies a function of pairs by the indicator of $\left\{(x,x'):[c]\in z\right\}$, where $c=\pi(x)$ is the photon site of the \emph{first} point.
\end{itemize}
These sets are invariant under $\ell\mathbb{Z}^3$, so that the products are again antiperiodic.

\paragraph{0.2. The spinor $\mathbf{1}_{z'}\mathrm{vv}$.}
Let $z'$ be a site of $T^0_k$, let $\tilde{z}'\in\mathbb{Z}^3$ be a representative of $z'$, and let $\mathrm{vv}\in\mathbb{C}^4$.
The spinor $\mathbf{1}_{z'}\mathrm{vv}$ of Section~G1.3 (localisation), formed with the representative $\tilde{z}'$, is the element $\mathrm{vz}$ of $\mathfrak{S}(1,\ell)$ given by
\begin{equation}\tag{0.1}
\mathrm{vz}(\tilde{z}'+\ell n):=(-1)^{n_1+n_2+n_3}\,\mathrm{vv}\quad(n\in\mathbb{Z}^3),\qquad\mathrm{vz}(y):=0\quad\text{for $y$ outside the class $z'$}.
\end{equation}
Another choice of the representative multiplies $\mathrm{vz}$ by $+1$ or by $-1$;
in what follows only moduli of images of $\mathrm{vz}$ under linear maps occur, and these do not depend on the choice.
One has $|\mathrm{vz}|_{\mathrm{su}}=|\mathrm{vv}|_2$, $|\mathrm{vz}|_{\infty}=|\mathrm{vv}|_{c}$ and $\mathbf{1}_{z'}\mathrm{vz}=\mathrm{vz}$.
Every $v\in\mathfrak{S}(1,\ell)$ satisfies $v(\tilde{z}'+\ell n)=(-1)^{n_1+n_2+n_3}v(\tilde{z}')$, and therefore $\mathbf{1}_{z'}v=\mathbf{1}_{z'}\left(v(\tilde{z}')\right)$.
It follows that, for every linear map $O$ on $\mathfrak{S}(1,\ell)$ with values in spinors or in functions of pairs,
\begin{equation}\tag{0.2}
\begin{split}
\left\|O\,\mathbf{1}_{z'}\right\|_{\mathrm{su}}
&=\sup\left\{\left|O\,\mathbf{1}_{z'}v\right|_{\mathrm{su}}:\ v\in\mathfrak{S}(1,\ell),\ |v|_{\mathrm{su}}\le1\right\}\\
&=\sup\left\{\left|O\,\mathbf{1}_{z'}\mathrm{vv}\right|_{\mathrm{su}}:\ \mathrm{vv}\in\mathbb{C}^4,\ |\mathrm{vv}|_2\le1\right\},
\end{split}
\end{equation}
and the same holds with $|\cdot|_{\infty}$, $|\cdot|_{c}$ in place of $|\cdot|_{\mathrm{su}}$, $|\cdot|_2$ (on the left, the operator norm is then the one for the norms $|\cdot|_{\infty}$).
Indeed, let $|\cdot|$ be one of the two norms on $\mathbb{C}^4$.
If $\sup|v|\le1$, then $O\,\mathbf{1}_{z'}v=O\,\mathbf{1}_{z'}\mathrm{vv}$ with $\mathrm{vv}:=v(\tilde{z}')$, and $|\mathrm{vv}|\le1$.
If $|\mathrm{vv}|\le1$, then $v:=\mathbf{1}_{z'}\mathrm{vv}$ satisfies $\sup|v|=|\mathrm{vv}|\le1$ and $\mathbf{1}_{z'}v=v$, whence $O\,\mathbf{1}_{z'}v=O\,\mathbf{1}_{z'}\mathrm{vv}$.
Hence the two sets of numbers whose suprema occur in (0.2) coincide.

\paragraph{0.3. Pairs (Section~I.1.7).}
Let $(x,x')\in\mathrm{Pzt}_k$, $x=(c,g)$, $x'=(c',g')$.
Then $x\ne x'$ and $d_X(x,x')=\max\left(|c-c'|,\eta_k\right)\le M$;
consequently $|c-c'|\le M<\ell/2$, since $\ell\ge13L\cdot M$ for $\ell\in\mathrm{Cl}$ (Section~I.1.1).
Two distinct points of $\mathbb{L}_k=\eta_k\mathbb{Z}^3$ have sup distance at least $\eta_k$.
Therefore exactly one of the following three cases occurs.
\begin{itemize}
\item[(pin)] $c=c'$ and $g\ne g'$, so that the pair lies inside a cluster;
then $d_X=\eta_k$ and $U_A(x,x')=U_A(c,c)=1$ by (I.1.6);
over every photon site there are $48\cdot47$ such ordered pairs.
\item[(psh)] $0<|c-c'|<1$; then $d_X=|c-c'|\in[\eta_k,1)$.
\item[(plg)] $1\le|c-c'|\le M$; then $d_X=|c-c'|\in[1,M]$.
\end{itemize}
\emph{The case $k\ge1$.}
Here $\eta_k\le1/L<1$;
the short pairs ($d_X<1$) are the pairs (pin) and (psh), and the long pairs ($1\le d_X\le M$) are the pairs (plg).

\emph{The case $k=0$.}
Here $\eta_0=1$ and $\mathbb{L}_0=\mathbb{Z}^3$;
the case (psh) does not occur, and in the case (pin) one has $d_X=1\le M$.
Every pair of $\mathrm{Pzt}_0$, whether (pin) or (plg), is long, and there is no short pair.

\paragraph{0.4. Numbers.}
For every $k\ge0$ one has
\[
\mathfrak{i}(A)=e_k\cdot\left|\operatorname{Im}A\right|_{\infty}<e_k^{\varepsilon}
\]
by \cite[Section~16.2.7]{QEDI}:
the connection is of the form $A=A_0+A_1$ with $A_0$ real and $|A_1|_{\infty}<e_k^{-1+\varepsilon}$, so that $\operatorname{Im}A=\operatorname{Im}A_1$.
The conditions (Z2.s1) and (Z2.s3) are imposed at every $e_k\in(0,e_0]$:
by Proposition~14.2\,(b) they are required at all values $v_k=-\log e_k\ge-\log e_0$, and $\exp\left(-\varepsilon\cdot v_k\right)=e_k^{\varepsilon}$.
In this form they read
\[
6L\cdot e_k^{\varepsilon}\le1,\qquad3M\cdot e_k^{\varepsilon}\le1.
\]
Hence
\begin{equation}\tag{0.3}
\frac{3}{2}\,\mathfrak{i}(A)<\frac{1}{4L}<1,\qquad3M\cdot\mathfrak{i}(A)<1.
\end{equation}
For $k\ge1$ one has
\[
b_k=\frac{b\cdot\left(1-L^{-1}\right)}{1-L^{-k}}\le b,
\]
because $1-L^{-k}\ge1-L^{-1}$ (Section~I.3.4).
Moreover $\mathrm{e}<3$; thus
\begin{equation}\tag{0.4}
1+\mathrm{e}<4,\qquad32\left(1+\mathrm{e}\right)<128<304.
\end{equation}

\paragraph{0.5. Phases.}
Here $T_k$ is the block torus $(0,\ell,-1)$ of level $k$ (Section~G1.2).
For every $k\ge0$ and every connection $A$ on $T_k$, \cite[Lemma~16.2(a)]{QEDI} gives
\begin{gather*}
\left|\varphi_A(y,c)^{\pm1}\right|\le\exp\left(\frac{3}{2}\,\mathfrak{i}(A)\right)\qquad(y\in\mathbb{Z}^3,\ c\in\Delta_y),\\
\left|U_A(c,c')^{\pm1}\right|\le\exp\left(3\,|c-c'|\cdot\mathfrak{i}(A)\right)\qquad(c,c'\in\mathbb{L}_k,\ |c-c'|<\ell/2).
\end{gather*}
These two bounds follow from the definitions, as follows.
The functionals $\tau$ and $\tau_k$ are sums over bonds with real coefficients (Section~I.1.1, (I.3.3));
hence $\operatorname{Im}(\tau A)=\tau(\operatorname{Im}A)$ and $\operatorname{Im}(\tau_kA)=\tau_k(\operatorname{Im}A)$, and (I.1.6), (I.3.3) give
\begin{gather*}
\left|U_A(c,c')^{\pm1}\right|=\exp\left(\mp e_k\cdot\eta_k\cdot(\tau\operatorname{Im}A)(c,c')\right),\\
\left|\varphi_A(y,c)^{\pm1}\right|=\exp\left(\mp e_k\cdot\eta_k\cdot(\tau_k\operatorname{Im}A)(y,c)\right).
\end{gather*}
Each of the six paths $\boldsymbol{\Gamma}^{\pi}(c,c')$ consists of $|c-c'|_1/\eta_k\le3\,|c-c'|/\eta_k$ bonds;
therefore $e_k\cdot\eta_k\cdot\left|(\tau\operatorname{Im}A)(c,c')\right|\le3\,|c-c'|\cdot\mathfrak{i}(A)$.
The centres $y_j$ of Section~I.3.2 satisfy $|y_{j+1}-y_j|\le\left(L^{j+1}-L^j\right)\cdot\eta_k/2$:
indeed, $y_j$ lies in the block of side $L^{j+1}\eta_k$ with centre $y_{j+1}$, the difference $y_j-y_{j+1}$ lies in $L^j\eta_k\mathbb{Z}^3$, and $L$ is odd.
Consequently
\begin{align*}
e_k\cdot\eta_k\cdot\left|(\tau_k\operatorname{Im}A)(y,c)\right|
&\le\sum_{j<k}3\,|y_{j+1}-y_j|\cdot\mathfrak{i}(A)
\le\sum_{j<k}\frac{3}{2}\left(L^{j+1}-L^j\right)\cdot\eta_k\cdot\mathfrak{i}(A)\\
&=\frac{3}{2}\left(1-L^{-k}\right)\cdot\mathfrak{i}(A)\le\frac{3}{2}\,\mathfrak{i}(A).
\end{align*}

For a pair $(x,x')$ of $\mathrm{Pzt}_k$ the phase $U_A(x,x')=U_A(c,c')$ is formed with the two points $c$, $c'$ of $\mathbb{L}_k$ ((I.1.6), Section~I.1.8), and $|c-c'|\le M<\ell/2$ by Section~0.3.
Together with (0.3) this gives
\begin{equation}\tag{0.5}
\left|\varphi_A(y,c)^{\pm1}\right|\le\mathrm{e}\quad(c\in\Delta_y),\qquad
\left|U_A(x,x')^{\pm1}\right|\le\exp\left(3M\cdot\mathfrak{i}(A)\right)\le\mathrm{e}
\end{equation}
for \emph{every} pair $(x,x')$ of $\mathrm{Pzt}_k$, short or long, and for every $k\ge0$.
The second bound in (0.5) is the first assertion of Proposition~6\,(b); see Section~2.1.

\paragraph{0.6. The operator $J\,Q_k^{T}(-A)$ ($k\ge1$).}
By (I.3.1) and Section~I.3.3, for $v\in\mathfrak{S}(1,\ell)$ and $x=(c,g)$,
\begin{equation}\tag{0.6}
\left(J\,Q_k^{T}(-A)\,v\right)(c,g)=\varphi_A([c],c)^{-1}\cdot v([c]),
\end{equation}
that is, the same vector at the $48$ sites of the cluster of $c$, multiplied by a number.
By (0.5), $\left|\left(J\,Q_k^{T}(-A)\,v\right)(x)\right|\le\mathrm{e}\cdot|v([c])|$ for both norms on $\mathbb{C}^4$, and therefore
\begin{equation}\tag{0.7}
\left|J\,Q_k^{T}(-A)\,v\right|_{\mathrm{su}}\le\mathrm{e}\cdot|v|_{\mathrm{su}},\qquad\left|J\,Q_k^{T}(-A)\,v\right|_{\infty}\le\mathrm{e}\cdot|v|_{\infty};
\end{equation}
this is (I.Q3), (I.J2) combined with $\exp\left(\frac{3}{2}\,\mathfrak{i}(A)\right)\le\mathrm{e}$.
The value $(\mathbf{1}_{z'}v)([c])$ equals $v([c])$ if $[c]$ lies in the class $z'$, and $0$ otherwise.
By (0.6) it follows that
\begin{equation}\tag{0.8}
J\,Q_k^{T}(-A)\,\mathbf{1}_{z'}=\mathbf{1}_{\Delta_{z'}}\,J\,Q_k^{T}(-A).
\end{equation}

\paragraph{0.7. Consequences of Proposition~2.8 ($k\ge1$).}
The point $y=\infty$ belongs to $Y_k$, and $w_k(\infty)=1$ (Section~I.3.5).
By Proposition~2.8\,(a) at $y=\infty$, the operator $F_{\infty}(A)=D_A+\mu+b_k\cdot J\,P_k(A)\,J^{*}$ is bijective on $\mathfrak{S}^{\mathrm{ch}}(\eta_k,\ell)$.
Hence $S_k(A)=S_{\infty}(A)=F_{\infty}(A)^{-1}$ exists, and the operator of Section~I.3.5
\begin{equation}\tag{0.9}
H_k(A)=b_k\cdot S_k(A)\,J\,Q_k^{T}(-A)\colon\mathfrak{S}(1,\ell)\to\mathfrak{S}^{\mathrm{ch}}(\eta_k,\ell)
\end{equation}
is defined.
Proposition~2.8\,(c) at $y=\infty$ states that, for all sites $z$, $z'$ of $T^0_k$ and every $f\in\mathfrak{S}^{\mathrm{ch}}(\eta_k,\ell)$,
\begin{align}
\left|\mathbf{1}_{\Delta_z}S_k(A)\,\mathbf{1}_{\Delta_{z'}}f\right|_{\mathrm{su}}&\le32a'\cdot\exp\left(-\gamma^{F}\cdot d(z,z')\right)\cdot|f|_{\mathrm{su}},\tag{0.10}\\
\left|\mathbf{1}_{\Delta_z}\delta_{\alpha,A}S_k(A)\,\mathbf{1}_{\Delta_{z'}}f\right|_{\mathrm{su}}&\le304a'\cdot\exp\left(-\gamma^{F}\cdot d(z,z')\right)\cdot|f|_{\mathrm{su}},\tag{0.11}
\end{align}
where $\delta_{\alpha,A}S_k(A)$ denotes the composition, that is, the map $\delta_{\alpha,A}$ of (I.1.7) applied to the spinor $S_k(A)\,\mathbf{1}_{\Delta_{z'}}f$ (the convention at the end of Section~G1.6) and restricted to the short pairs (Section~G1.3);
the supremum in (0.11) extends over the short pairs whose first point lies over $\Delta_z$, including the pairs (pin) inside a cluster.
Moreover $N_{\gamma^{F}}\left(S_k(A)\right)\le32a'$ and $N_{\gamma^{F}}\left(\delta_{\alpha,A}S_k(A)\right)\le304a'$.
By the inequality (N1) of Section~G1.3, namely $\|O\|_{\mathrm{su}}\le N_{\gamma}(O)$, every $f\in\mathfrak{S}^{\mathrm{ch}}(\eta_k,\ell)$ satisfies
\begin{equation}\tag{0.12}
\begin{split}
&\left|S_k(A)f\right|_{\mathrm{su}}\le32a'\cdot|f|_{\mathrm{su}},\\
&\left|\left(\delta_{\alpha,A}S_k(A)f\right)(x,x')\right|_2\le304a'\cdot|f|_{\mathrm{su}}\quad\text{for every short pair $(x,x')$}.
\end{split}
\end{equation}

\paragraph{0.8. Row sums.}
Let $O$ be a linear map from $\mathfrak{S}(1,\ell)$ to $\mathfrak{S}^{\mathrm{ch}}(\eta_k,\ell)$ or to the functions of pairs,
and let $t$ be a row index as in (I.1.4): $t=x\in\mathbb{L}^{\mathrm{f}}_k$ in the first case, $t=(x,x')\in\mathrm{Pzt}_k$ in the second.
Let $\mathrm{Dfu}$ be a fundamental domain of $\ell\mathbb{Z}^3$ in $\mathbb{Z}^3$; it consists of $\ell^3$ unit sites.
The kernel of $O$ is the function $O(t,y')$ with values in the $4\times4$ matrices defined by $O(t,y')\,\mathrm{vv}:=\left(O\,\mathbf{1}_{y'}\mathrm{vv}\right)(t)$,
where $\mathbf{1}_{y'}\mathrm{vv}$ is formed with the representative $y'\in\mathbb{Z}^3$.
By (0.1), $\mathbf{1}_{y'+\ell e_{\mu}}\mathrm{vv}=-\mathbf{1}_{y'}\mathrm{vv}$, so that $O(t,y'+\ell e_{\mu})=-O(t,y')$;
and $(O\,v)(t)=\sum_{y'\in\mathrm{Dfu}}O(t,y')\,v(y')$ (Section~I.1.6), because $v=\sum_{y'\in\mathrm{Dfu}}\mathbf{1}_{y'}\left(v(y')\right)$.
By (I.ke3), for the source $\mathfrak{S}(1,\ell)$,
\begin{equation}\tag{0.13}
|O|_{1,\infty}=\sup_t\max_{\beta}\sum_{y'\in\mathrm{Dfu}}\sum_{\beta'=1}^{4}\left|O(t,y')_{\beta\beta'}\right|;
\end{equation}
every unit site carries the weight $1$ (Sections~I.1.6, I.3.6), and (0.13) contains neither a factor $48$ nor a factor $1/48$.
The moduli of the entries are $\ell\mathbb{Z}^3$-periodic in $y'$; hence the number (0.13) does not depend on the choice of $\mathrm{Dfu}$.
\begin{itemize}
\item[(i)] $|O|_{1,\infty}=\sup\left\{|O\,v|_{\infty}:\ v\in\mathfrak{S}(1,\ell),\ |v|_{\infty}\le1\right\}$ (I.ke3).
\item[(ii)] If $\|O\|_{\mathrm{su}}\le\mathrm{cf}$ for a number $\mathrm{cf}$, then $|O|_{1,\infty}\le2\,\mathrm{cf}$.
\end{itemize}
To prove (i), let first $|v|_{\infty}\le1$; then
\[
\left|(O\,v)(t)_{\beta}\right|\le\sum_{y'\in\mathrm{Dfu}}\sum_{\beta'}\left|O(t,y')_{\beta\beta'}\right|\cdot\left|v(y')_{\beta'}\right|\le|O|_{1,\infty}.
\]
Conversely, fix $(t,\beta)$.
The values of an element of $\mathfrak{S}(1,\ell)$ on $\mathrm{Dfu}$ can be prescribed arbitrarily.
For $y'\in\mathrm{Dfu}$ put
\[
v(y')_{\beta'}:=\frac{\overline{O(t,y')_{\beta\beta'}}}{\left|O(t,y')_{\beta\beta'}\right|}\quad\text{where this entry is not zero},\qquad v(y')_{\beta'}:=0\quad\text{otherwise}.
\]
Then $|v|_{\infty}\le1$ and $(O\,v)(t)_{\beta}=\sum_{y'\in\mathrm{Dfu}}\sum_{\beta'}\left|O(t,y')_{\beta\beta'}\right|$.
To prove (ii), recall that $|\mathrm{vv}|_{c}\le|\mathrm{vv}|_2\le2\,|\mathrm{vv}|_{c}$ on $\mathbb{C}^4$ (Section~I.1.5);
hence $|O\,v|_{\infty}\le|O\,v|_{\mathrm{su}}\le\mathrm{cf}\cdot|v|_{\mathrm{su}}\le2\,\mathrm{cf}\cdot|v|_{\infty}$ for every $v\in\mathfrak{S}(1,\ell)$, and (i) applies.

\subsubsection*{1. Clause (a) ($k\ge1$)}

Let $z$, $z'$ be sites of $T^0_k$ and $v\in\mathfrak{S}(1,\ell)$, and put
\[
f:=J\,Q_k^{T}(-A)\,\mathbf{1}_{z'}v\in\mathfrak{S}^{\mathrm{ch}}(\eta_k,\ell).
\]
By (0.8), $f=\mathbf{1}_{\Delta_{z'}}J\,Q_k^{T}(-A)\,v$, and hence $f=\mathbf{1}_{\Delta_{z'}}f$.
By (0.7), $|f|_{\mathrm{su}}\le\mathrm{e}\cdot|\mathbf{1}_{z'}v|_{\mathrm{su}}\le\mathrm{e}\cdot|v|_{\mathrm{su}}$.
From (0.9), and since $\delta_{\alpha,A}$ is linear (I.1.7), one obtains
\begin{equation}\tag{1.1}
\begin{split}
H_k(A)\,\mathbf{1}_{z'}v&=b_k\cdot S_k(A)\,\mathbf{1}_{\Delta_{z'}}f,\\
\delta_{\alpha,A}H_k(A)\,\mathbf{1}_{z'}v&=b_k\cdot \delta_{\alpha,A}S_k(A)\,\mathbf{1}_{\Delta_{z'}}f.
\end{split}
\end{equation}
By (0.10), (0.11) and $b_k\le b$,
\begin{equation}\tag{1.2}
\begin{split}
\left|\mathbf{1}_{\Delta_z}H_k(A)\,\mathbf{1}_{z'}v\right|_{\mathrm{su}}
&\le b_k\cdot32a'\cdot\exp\left(-\gamma^{F}\cdot d(z,z')\right)\cdot|f|_{\mathrm{su}}\\
&\le32\,\mathrm{e}\cdot b\cdot a'\cdot\exp\left(-\gamma^{F}\cdot d(z,z')\right)\cdot|v|_{\mathrm{su}},
\end{split}
\end{equation}
\begin{equation}\tag{1.3}
\begin{split}
\left|\mathbf{1}_{\Delta_z}\delta_{\alpha,A}H_k(A)\,\mathbf{1}_{z'}v\right|_{\mathrm{su}}
&\le b_k\cdot304a'\cdot\exp\left(-\gamma^{F}\cdot d(z,z')\right)\cdot|f|_{\mathrm{su}}\\
&\le304\,\mathrm{e}\cdot b\cdot a'\cdot\exp\left(-\gamma^{F}\cdot d(z,z')\right)\cdot|v|_{\mathrm{su}},
\end{split}
\end{equation}
where the supremum in (1.3) extends over the short pairs, (pin) and (psh), whose first point lies over $\Delta_z$.
The suprema of the left sides of (1.2) and (1.3) over all $v$ with $|v|_{\mathrm{su}}\le1$ are the two operator norms in Proposition~6\,(a),
so that (1.2) and (1.3) give the two bounds of Proposition~6\,(a).

\emph{Pairs inside a cluster.}
Let $\mathrm{psf}:=H_k(A)\,\mathbf{1}_{z'}v$, let $c$ be a photon site with $[c]\in z$, and let $g\ne g'$ be elements of $O_{\mathrm{cub}}$.
For the pair $\left((c,g),(c,g')\right)$ one has $d_X=\eta_k$ and $U_A=1$ (Section~0.3, case (pin));
hence $\left(\delta_{\alpha,A}\mathrm{psf}\right)\left((c,g),(c,g')\right)=\eta_k^{-\alpha}\cdot\left(\mathrm{psf}(c,g')-\mathrm{psf}(c,g)\right)$,
and (1.3) yields, for each of the $48\cdot47$ ordered pairs of labels,
\begin{equation}\tag{1.4}
\left|\mathrm{psf}(c,g')-\mathrm{psf}(c,g)\right|_2\le304\,\mathrm{e}\cdot b\cdot a'\cdot\eta_k^{\alpha}\cdot\exp\left(-\gamma^{F}\cdot d(z,z')\right)\cdot|v|_{\mathrm{su}}.
\end{equation}

\subsubsection*{2. Clause (b)}

\paragraph{2.1. First part (every $k\ge0$).}
Let $(x,x')\in\mathrm{Pzt}_k$ and $f\in\mathfrak{S}^{\mathrm{ch}}(\eta_k,\ell)$.
By (0.5), $\left|U_A(x,x')^{\pm1}\right|\le\mathrm{e}$;
this holds for every pair of $\mathrm{Pzt}_k$, short or long, and it is the first assertion of Proposition~6\,(b).
Now let the pair be long.
Then $d_X(x,x')\ge1$ and $\alpha>0$, so that $d_X(x,x')^{-\alpha}\le1$.
For both norms on $\mathbb{C}^4$, (I.1.7) gives
\begin{equation}\tag{2.1}
\begin{split}
\left|\left(\delta_{\alpha,A}f\right)(x,x')\right|
&=d_X(x,x')^{-\alpha}\cdot\left|U_A(x,x')\cdot f(x')-f(x)\right|
\le\left|U_A(x,x')\right|\cdot|f(x')|+|f(x)|\\
&\le\mathrm{e}\cdot|f(x')|+|f(x)|\le\left(1+\mathrm{e}\right)\cdot\sup|f|\le4\sup|f|,
\end{split}
\end{equation}
the last inequality by (0.4).
With $|\cdot|_2$ and $|f|_{\mathrm{su}}$ this is the display in Proposition~6\,(b).
With $|\cdot|_{c}$ and $|f|_{\infty}$, and together with $\left|U_A(x,x')^{\pm1}\right|\le\mathrm{e}$ for every pair, it is (I.Lc2)$^{\mathrm{ch}}$ for $B=A\in\tilde{G}_k$.

\emph{Level $0$.}
Here $\eta_0=1$, and by Section~0.3 every pair of $\mathrm{Pzt}_0$ is long;
hence (2.1) holds for \emph{all} pairs of $\mathrm{Pzt}_0$.
The bounds (0.3), (0.5) and the steps of (2.1) are valid for every $k\ge0$, including $k=0$.

\paragraph{2.2. Second part ($k\ge1$).}
Let $z$, $z'$, $z''$ be sites of $T^0_k$ and $\mathrm{vv}\in\mathbb{C}^4$, put $\mathrm{vz}:=\mathbf{1}_{z'}\mathrm{vv}$ as in (0.1) and $\mathrm{psf}:=H_k(A)\,\mathrm{vz}$,
and let $(x,x')$ be a long pair with $x$ over $\Delta_z$ and $x'$ over $\Delta_{z''}$.
By Section~0.2, $\mathrm{vz}=\mathbf{1}_{z'}\mathrm{vz}$ and $|\mathrm{vz}|_{\mathrm{su}}=|\mathrm{vv}|_2$.
Applying (1.2) with $v=\mathrm{vz}$, once for the block $\Delta_z$ and once for the block $\Delta_{z''}$, one obtains
\begin{gather*}
|\mathrm{psf}(x)|_2\le\left|\mathbf{1}_{\Delta_z}H_k(A)\,\mathbf{1}_{z'}\mathrm{vz}\right|_{\mathrm{su}}\le32\,\mathrm{e}\cdot b\cdot a'\cdot\exp\left(-\gamma^{F}\cdot d(z,z')\right)\cdot|\mathrm{vv}|_2,\\
|\mathrm{psf}(x')|_2\le\left|\mathbf{1}_{\Delta_{z''}}H_k(A)\,\mathbf{1}_{z'}\mathrm{vz}\right|_{\mathrm{su}}\le32\,\mathrm{e}\cdot b\cdot a'\cdot\exp\left(-\gamma^{F}\cdot d(z'',z')\right)\cdot|\mathrm{vv}|_2.
\end{gather*}
Both exponentials are at most $\exp\left(-\gamma^{F}\cdot\min\left(d(z,z'),d(z'',z')\right)\right)$.
By (2.1) for $f=\mathrm{psf}$ and by (0.4),
\begin{equation}\tag{2.2}
\begin{split}
\left|\left(\delta_{\alpha,A}H_k(A)\,\mathbf{1}_{z'}\mathrm{vv}\right)(x,x')\right|_2
&\le\mathrm{e}\cdot|\mathrm{psf}(x')|_2+|\mathrm{psf}(x)|_2\\
&\le\left(1+\mathrm{e}\right)\cdot32\,\mathrm{e}\cdot b\cdot a'\cdot\exp\left(-\gamma^{F}\cdot\min\left(d(z,z'),d(z'',z')\right)\right)\cdot|\mathrm{vv}|_2\\
&\le128\,\mathrm{e}\cdot b\cdot a'\cdot\exp\left(-\gamma^{F}\cdot\min\left(d(z,z'),d(z'',z')\right)\right)\cdot|\mathrm{vv}|_2.
\end{split}
\end{equation}

\subsubsection*{3. Clause (c) ($k\ge1$)}

Let $v\in\mathfrak{S}(1,\ell)$, and put $f:=J\,Q_k^{T}(-A)\,v$ and $\mathrm{psf}:=H_k(A)\,v=b_k\cdot S_k(A)f$.
By (0.7), $|f|_{\mathrm{su}}\le\mathrm{e}\cdot|v|_{\mathrm{su}}$.
The first bound in (0.12) and $b_k\le b$ give
\begin{equation}\tag{3.1}
\left|H_k(A)\,v\right|_{\mathrm{su}}=b_k\cdot\left|S_k(A)f\right|_{\mathrm{su}}\le b\cdot32a'\cdot\mathrm{e}\cdot|v|_{\mathrm{su}}=32\,\mathrm{e}\cdot b\cdot a'\cdot|v|_{\mathrm{su}}.
\end{equation}
Let $(x,x')\in\mathrm{Pzt}_k$.
By Section~0.3 the pair is short or long.

\emph{Short pairs, (pin) and (psh).}
One has $\delta_{\alpha,A}\mathrm{psf}=b_k\cdot \delta_{\alpha,A}S_k(A)f$, and the second bound in (0.12) together with $b_k\le b$ yields
\begin{equation}\tag{3.2}
\left|\left(\delta_{\alpha,A}H_k(A)\,v\right)(x,x')\right|_2\le b\cdot304a'\cdot\mathrm{e}\cdot|v|_{\mathrm{su}}=304\,\mathrm{e}\cdot b\cdot a'\cdot|v|_{\mathrm{su}}.
\end{equation}

\emph{Long pairs, (plg).}
By (2.1) for $f=\mathrm{psf}$, by (3.1) and by (0.4),
\begin{equation}\tag{3.3}
\begin{split}
\left|\left(\delta_{\alpha,A}H_k(A)\,v\right)(x,x')\right|_2
&\le\left(1+\mathrm{e}\right)\cdot|\mathrm{psf}|_{\mathrm{su}}\le\left(1+\mathrm{e}\right)\cdot32\,\mathrm{e}\cdot b\cdot a'\cdot|v|_{\mathrm{su}}\\
&\le128\,\mathrm{e}\cdot b\cdot a'\cdot|v|_{\mathrm{su}}\le304\,\mathrm{e}\cdot b\cdot a'\cdot|v|_{\mathrm{su}}.
\end{split}
\end{equation}
The bounds (3.1)--(3.3) constitute Proposition~6\,(c).
Written out for the three kinds of pairs, with $x=(c,g)$ and $x'=(c',g')$, they take the following form;
in the cases (pin) and (psh) one multiplies (3.2) by $d_X^{\alpha}$, and in the case (plg) one uses the triangle inequality, (0.5), (3.1) and (0.4).
\begin{itemize}
\item[(pin)] For the $48\cdot47$ ordered pairs over each photon site, without phase:
\[
\left|\mathrm{psf}(c,g')-\mathrm{psf}(c,g)\right|_2\le304\,\mathrm{e}\cdot b\cdot a'\cdot\eta_k^{\alpha}\cdot|v|_{\mathrm{su}}.
\]
\item[(psh)] For $0<|c-c'|<1$ and all $48^2$ pairs of labels:
\[
\left|U_A(c,c')\cdot\mathrm{psf}(c',g')-\mathrm{psf}(c,g)\right|_2\le304\,\mathrm{e}\cdot b\cdot a'\cdot|c-c'|^{\alpha}\cdot|v|_{\mathrm{su}}.
\]
\item[(plg)] For $1\le|c-c'|\le M$ and all $48^2$ pairs of labels:
\[
\left|U_A(c,c')\cdot\mathrm{psf}(c',g')-\mathrm{psf}(c,g)\right|_2\le128\,\mathrm{e}\cdot b\cdot a'\cdot|v|_{\mathrm{su}}.
\]
\end{itemize}

\subsubsection*{4. Clause (d) (barred fields: $k\ge1$; in the first part of (b): every $k\ge0$)}

\paragraph{4.1. The matrix $\mathsf{C}$.}
By the entries of $\mathsf{C}$ listed in Section~I.2.1, for $\mathrm{vv}=(\mathrm{vv}_1,\mathrm{vv}_2,\mathrm{vv}_3,\mathrm{vv}_4)\in\mathbb{C}^4$ one has
\[
\mathsf{C}\,\mathrm{vv}=\left(\mathrm{vv}_2,-\mathrm{vv}_1,-\mathrm{vv}_4,\mathrm{vv}_3\right),\qquad
\mathsf{C}^{-1}\mathrm{vv}=-\mathsf{C}\,\mathrm{vv}=\left(-\mathrm{vv}_2,\mathrm{vv}_1,\mathrm{vv}_4,-\mathrm{vv}_3\right),
\]
in accordance with $\mathsf{C}^{-1}=\mathsf{C}^{T}=-\mathsf{C}$ (I.c3).
Hence
\begin{equation}\tag{4.1}
\left|\mathsf{C}\,\mathrm{vv}\right|_2=\left|\mathsf{C}^{-1}\mathrm{vv}\right|_2=|\mathrm{vv}|_2,\qquad\left|\mathsf{C}\,\mathrm{vv}\right|_{c}=\left|\mathsf{C}^{-1}\mathrm{vv}\right|_{c}=|\mathrm{vv}|_{c}.
\end{equation}
A constant $4\times4$ matrix $\mathrm{Mx}$ acts on spinors and on functions of pairs by
\[
(\mathrm{Mx}\,f)(x):=\mathrm{Mx}\,f(x),\qquad(\mathrm{Mx}\,v)(y):=\mathrm{Mx}\,v(y),\qquad(\mathrm{Mx}\,F)(x,x'):=\mathrm{Mx}\,F(x,x').
\]
This action preserves the antiperiodicity, and it commutes with $\mathbf{1}_{z'}$ and with $\mathbf{1}_{\Delta_z}$, which are multiplications by numbers.
For every bond function $B$ the quantities $U_B(x,x')$ and $d_X(x,x')^{-\alpha}$ are numbers; therefore (I.1.7) gives
\begin{equation}\tag{4.2}
\delta_{\alpha,B}\left(\mathrm{Mx}\,f\right)=\mathrm{Mx}\,\delta_{\alpha,B}f.
\end{equation}
By (4.1), $\mathsf{C}$ and $\mathsf{C}^{-1}$ are bijective isometries of $\mathfrak{S}(1,\ell)$, of $\mathfrak{S}^{\mathrm{ch}}(\eta_k,\ell)$ and of the space of functions of pairs, for $|\cdot|_{\mathrm{su}}$ as well as for $|\cdot|_{\infty}$.

\paragraph{4.2. Reduction to the field $-A$.}
For every $k\ge0$ the connection $-A$ lies in $\tilde{G}_k$ \cite[Section~16.2.7]{QEDI}.
Hence Sections~0--3 apply with $-A$ in place of $A$ (``at $-A$'').
In particular, for $k\ge1$ the operator $H_k(-A)$ is defined, and $\bar{H}_k(A)=\mathsf{C}^{-1}H_k(-A)\,\mathsf{C}\colon\mathfrak{S}(1,\ell)\to\mathfrak{S}^{\mathrm{ch}}(\eta_k,\ell)$ (Section~I.3.5).
By (4.2) with $B=-A$ and $\mathrm{Mx}=\mathsf{C}^{-1}$, and since $\mathsf{C}$ and $\mathsf{C}^{-1}$ commute with $\mathbf{1}_{z'}$ and $\mathbf{1}_{\Delta_z}$ (Section~4.1),
\begin{equation}\tag{4.3}
\begin{split}
\mathbf{1}_{\Delta_z}\bar{H}_k(A)\,\mathbf{1}_{z'}&=\mathsf{C}^{-1}\left[\mathbf{1}_{\Delta_z}H_k(-A)\,\mathbf{1}_{z'}\right]\mathsf{C},\\
\mathbf{1}_{\Delta_z}\delta_{\alpha,-A}\bar{H}_k(A)\,\mathbf{1}_{z'}&=\mathsf{C}^{-1}\left[\mathbf{1}_{\Delta_z}\delta_{\alpha,-A}H_k(-A)\,\mathbf{1}_{z'}\right]\mathsf{C},
\end{split}
\end{equation}
\begin{equation}\tag{4.4}
\begin{split}
&\bar{H}_k(A)\,v=\mathsf{C}^{-1}H_k(-A)\left(\mathsf{C}\,v\right),\qquad\delta_{\alpha,-A}\bar{H}_k(A)\,v=\mathsf{C}^{-1}\delta_{\alpha,-A}H_k(-A)\left(\mathsf{C}\,v\right),\\
&\mathsf{C}\left(\mathbf{1}_{z'}\mathrm{vv}\right)=\mathbf{1}_{z'}\left(\mathsf{C}\,\mathrm{vv}\right),
\end{split}
\end{equation}
where in the last identity both sides are formed with the same representative $\tilde{z}'$ (see (0.1)).

\paragraph{4.3. The bounds.}
\emph{Proposition~6\,(a) for the barred operators ($k\ge1$).}
By (4.3), by Section~4.1 (composition with bijective isometries leaves the operator norm unchanged) and by (1.2), (1.3) at $-A$,
\begin{align*}
\left\|\mathbf{1}_{\Delta_z}\bar{H}_k(A)\,\mathbf{1}_{z'}\right\|_{\mathrm{su}}&=\left\|\mathbf{1}_{\Delta_z}H_k(-A)\,\mathbf{1}_{z'}\right\|_{\mathrm{su}}\le32\,\mathrm{e}\cdot b\cdot a'\cdot\exp\left(-\gamma^{F}\cdot d(z,z')\right),\\
\left\|\mathbf{1}_{\Delta_z}\delta_{\alpha,-A}\bar{H}_k(A)\,\mathbf{1}_{z'}\right\|_{\mathrm{su}}&=\left\|\mathbf{1}_{\Delta_z}\delta_{\alpha,-A}H_k(-A)\,\mathbf{1}_{z'}\right\|_{\mathrm{su}}\le304\,\mathrm{e}\cdot b\cdot a'\cdot\exp\left(-\gamma^{F}\cdot d(z,z')\right),
\end{align*}
the second bound on the short pairs, the pairs inside a cluster included.

\emph{Proposition~6\,(b), first part (every $k\ge0$).}
By (0.5) and (2.1) at $-A$, one has $\left|U_{-A}(x,x')^{\pm1}\right|\le\mathrm{e}$ for every pair of $\mathrm{Pzt}_k$, and for a long pair
\[
\left|\left(\delta_{\alpha,-A}f\right)(x,x')\right|\le\mathrm{e}\cdot|f(x')|+|f(x)|\le4\sup|f|,
\]
for both norms on $\mathbb{C}^4$.

\emph{Proposition~6\,(b), second part ($k\ge1$).}
By (4.4), (4.1) and (2.2) at $-A$, with $\mathsf{C}\,\mathrm{vv}$ in place of $\mathrm{vv}$ (note that $|\mathsf{C}\,\mathrm{vv}|_2=|\mathrm{vv}|_2$), one has, for a long pair with $x$ over $\Delta_z$ and $x'$ over $\Delta_{z''}$,
\begin{equation*}
\begin{split}
\left|\left(\delta_{\alpha,-A}\bar{H}_k(A)\,\mathbf{1}_{z'}\mathrm{vv}\right)(x,x')\right|_2
&=\left|\left(\delta_{\alpha,-A}H_k(-A)\,\mathbf{1}_{z'}\left(\mathsf{C}\,\mathrm{vv}\right)\right)(x,x')\right|_2\\
&\le128\,\mathrm{e}\cdot b\cdot a'\cdot\exp\left(-\gamma^{F}\cdot\min\left(d(z,z'),d(z'',z')\right)\right)\cdot|\mathrm{vv}|_2.
\end{split}
\end{equation*}

\emph{Proposition~6\,(c) for the barred operators ($k\ge1$).}
By (4.4), (4.1) and (3.1)--(3.3) at $-A$, with $\mathsf{C}\,v$ in place of $v$ (note that $|\mathsf{C}\,v|_{\mathrm{su}}=|v|_{\mathrm{su}}$),
\begin{align*}
\left|\bar{H}_k(A)\,v\right|_{\mathrm{su}}&=\left|H_k(-A)\left(\mathsf{C}\,v\right)\right|_{\mathrm{su}}\le32\,\mathrm{e}\cdot b\cdot a'\cdot|v|_{\mathrm{su}},\\
\sup_{(x,x')\in\mathrm{Pzt}_k}\left|\left(\delta_{\alpha,-A}\bar{H}_k(A)\,v\right)(x,x')\right|_2&=\sup_{(x,x')\in\mathrm{Pzt}_k}\left|\left(\delta_{\alpha,-A}H_k(-A)\left(\mathsf{C}\,v\right)\right)(x,x')\right|_2\\
&\le304\,\mathrm{e}\cdot b\cdot a'\cdot|v|_{\mathrm{su}},
\end{align*}
where the supremum extends over all pairs (pin), (psh), (plg).

\subsubsection*{5. Clause (e) (level 0)}

Let $k=0$.
Then $\eta_0=1$ and $\mathbb{L}^{\mathrm{f}}_0=\mathbb{Z}^3\times O_{\mathrm{cub}}$;
the fine lattice coincides with the unit lattice, and the fields of level $0$ are the elements of $\mathfrak{S}^{\mathrm{ch}}(1,\ell)$, with $192$ components over every photon site (Section~I.3.9).
\begin{itemize}
\item[(1)]
By definition, $H_0(A)=\bar{H}_0(A)=1$ on $\mathfrak{S}^{\mathrm{ch}}(1,\ell)$ (Section~I.3.9).
This agrees with the formula by which the barred operator is formed at the levels $k\ge1$:
\[
\mathsf{C}^{-1}H_0(-A)\,\mathsf{C}=\mathsf{C}^{-1}\mathsf{C}=1.
\]
\item[(2)]
At level $0$ there is no short pair:
for every pair of $\mathrm{Pzt}_0$ one has $d_X(x,x')=\max\left(|c-c'|,1\right)\ge1$ (Section~0.3).
\item[(3)]
Let $x=(c,g)$ and $x'=(c,g')$ with $g\neq g'$, so that the pair lies inside a cluster.
Then $d_X(x,x')=1$ and $U_A(x,x')=U_A(c,c)=1$, and (I.1.7) gives
\[
\left(\delta_{\alpha,A}f\right)(x,x')=1^{-\alpha}\left(f(x')-f(x)\right)=f(x')-f(x).
\]
As $U_{-A}(c,c)=1$ as well, one has likewise $\left(\delta_{\alpha,-A}f\right)(x,x')=f(x')-f(x)$.
\item[(4)]
By (1), $H_0^{\#}(A)f=\left(f,\delta_{\alpha,A}f\right)$, and $\left|H_0^{\#}(A)f\right|_\infty$ is the larger of the two suprema that occur in (5.1) below.
Here $H_0^{\#}(A)$ and $\left|H_0^{\#}(A)f\right|_\infty$ are defined by Section~I.3.5, which applies at $k=0$ by the first sentence of Section~I.5.10, with $H_0(A)=1$ and with the source $\mathfrak{S}^{\mathrm{ch}}(1,\ell)$ of Section~I.3.9;
this is the quantity that is bounded in (I.L0.15).
By (2), every pair of $\mathrm{Pzt}_0$ is long.
The estimate (2.1), taken with the norm $|\cdot|_c$, is valid at $k=0$ (Section~2.1, paragraph on level~$0$), and it yields
\begin{equation}
\begin{split}
\left|H_0^{\#}(A)f\right|_\infty
&=\max\left\{|f|_\infty,\ \sup_{(x,x')\in\mathrm{Pzt}_0}\left|\left(\delta_{\alpha,A}f\right)(x,x')\right|_c\right\}\\
&\le\max\{1,1+\mathrm{e}\}\cdot|f|_\infty=(1+\mathrm{e})\cdot|f|_\infty\le4|f|_\infty.
\end{split}
\tag{5.1}
\end{equation}
For the pairs of (3) the bound is $2|f|_\infty$.
\item[(5)]
The connection $-A$ lies in $\tilde{G}_0$ \cite[Section~16.2.7, (G2)]{QEDI}.
By (1),
\[
\left(\bar{H}_0(A)f,\ \delta_{\alpha,-A}\bar{H}_0(A)f\right)=\left(f,\ \delta_{\alpha,-A}f\right),
\]
and (2.1) with the norm $|\cdot|_c$ at $-A$, which holds for all pairs of $\mathrm{Pzt}_0$ (Section~4.3, (b), first part), gives
\begin{equation}
\begin{split}
\max\left\{\left|\bar{H}_0(A)f\right|_\infty,\ \sup_{(x,x')\in\mathrm{Pzt}_0}\left|\left(\delta_{\alpha,-A}\bar{H}_0(A)f\right)(x,x')\right|_c\right\}
&\le(1+\mathrm{e})\cdot|f|_\infty\\
&\le4|f|_\infty.
\end{split}
\tag{5.2}
\end{equation}
\end{itemize}
The points (1)--(5) are the assertions of (e);
the last sentence of (e) only restricts (a) and the second part of (b) to the levels $k\ge1$.

\subsubsection*{6. Clause (f) (the label FE4$^{\mathrm{ch}}$, in the norms of Section~I.1.5)}

\paragraph{6.1. Inequalities between constants.}
\emph{The bound $a'\ge1$.}
By the definitions of the constants (Proposition~14.2\,(a), lines (C4r)),
\[
a'=C_z\cdot c(\gamma^{F}),\qquad
C_z=\max\left(C_{5},C_{00}\right)\ge C_{00}=48\cdot\frac{2L}{c_{\mathrm{re}}}\cdot\left(1+2k_{1}\right),
\]
where $k_1=60+b=61$, because $b=1$, and
\[
c_{\mathrm{re}}=\min\left\{b\left(1-\frac{1}{L}\right),\ \frac{\rho}{2C_{\mathrm{Po}}\cdot L}\right\}.
\]
The first entry of this minimum is $1-1/L$, a number in $(0,1)$, since $L\ge5$.
The second entry is positive.
Indeed, $\rho>0$ (line (C0));
by the lines (C0r), $C_{\mathrm{Po}}=1/(8c_*)+1/\lambda_{\mathrm{int}}$ and $c_*=\min\left(w_3,\lambda_{\mathrm{int}}\right)/18$, where $w_3=6$ (Section~I.0.5) and $0<\lambda_{\mathrm{int}}<\infty$, $\lambda_{\mathrm{int}}$ being a positive eigenvalue (I.O8);
hence $0<c_*<\infty$ and $0<C_{\mathrm{Po}}<\infty$.
Consequently $0<c_{\mathrm{re}}\le1-1/L<1$, so that $2L/c_{\mathrm{re}}>2L$ and
\[
C_{00}>48\cdot2L\cdot123>1.
\]
Moreover
\[
c(\gamma^{F})=\sum_{n\in\mathbb{Z}^3}\exp\left(-\gamma^{F}\cdot|n|\right)\ge1
\]
(Section~G1.3), the term $n=0$ alone being equal to $1$.
It follows that $a'\ge C_{00}>1$.

\emph{The constant $C^{F}$.}
By its definition (Proposition~14.1; Proposition~14.2\,(a), line (C6)),
\begin{align*}
C^{F}&=2\max\left\{C_{1},\ C_{\mathrm{By}},\ 304\,\mathrm{e}\cdot\max(1,b)\cdot a',\ C_{\Gamma},\ \mathrm{e}\cdot b^2\cdot\left(32a'\right)^2\right\}\\
&\ge608\,\mathrm{e}\cdot\max(1,b)\cdot a'.
\end{align*}
Since $a'\ge1>0$, $b>0$, $\max(1,b)\ge b$ and $\max(1,b)\ge1$, one obtains
\begin{align}
64\,\mathrm{e}\cdot b\cdot a'
&\le128\,\mathrm{e}\cdot b\cdot a'\le256\,\mathrm{e}\cdot b\cdot a'\le608\,\mathrm{e}\cdot b\cdot a'
\notag\\
&\le608\,\mathrm{e}\cdot\max(1,b)\cdot a'\le C^{F},
\tag{6.1}\\
4&\le608\,\mathrm{e}\le608\,\mathrm{e}\cdot\max(1,b)\cdot a'\le C^{F}.
\tag{6.2}
\end{align}

\paragraph{6.2. Global bounds for $k\ge1$.}
The operator $H_k(A)$ is a linear map from $\mathfrak{S}(1,\ell)$ to $\mathfrak{S}^{\mathrm{ch}}(\eta_k,\ell)$, and $\delta_{\alpha,A}H_k(A)$ is a linear map from $\mathfrak{S}(1,\ell)$ to the functions of pairs.
The rows of the first are indexed by $(x,\beta)$ with $x\in\mathbb{L}^{\mathrm{f}}_k$;
the rows of the second by the triples $(x,x',\beta)$ with $(x,x')\in\mathrm{Pzt}_k$, where \emph{all} pairs occur, those of the kinds (pin), (psh) and (plg).
Clause (c) gives
\[
\left\|H_k(A)\right\|_{\mathrm{su}}\le32\,\mathrm{e}\cdot b\cdot a',\qquad
\left\|\delta_{\alpha,A}H_k(A)\right\|_{\mathrm{su}}\le304\,\mathrm{e}\cdot b\cdot a',
\]
and Section~0.8\,(ii) together with (6.1) yields
\begin{equation}
\begin{aligned}
\left|H_k(A)\right|_{1,\infty}&\le64\,\mathrm{e}\cdot b\cdot a'\le C^{F},\\
\left|\delta_{\alpha,A}H_k(A)\right|_{1,\infty}&\le608\,\mathrm{e}\cdot b\cdot a'\le C^{F}.
\end{aligned}
\tag{6.3}
\end{equation}
By Section~0.8\,(i), the bounds (6.3) mean that for every $v\in\mathfrak{S}(1,\ell)$
\begin{gather*}
\left|H_k(A)v\right|_\infty\le C^{F}\cdot|v|_\infty,\\
\sup_{(x,x')\in\mathrm{Pzt}_k}\left|\left(\delta_{\alpha,A}H_k(A)v\right)(x,x')\right|_c\le C^{F}\cdot|v|_\infty.
\end{gather*}
A translation of the row index by an element of $\ell\mathbb{Z}^3$ leaves the moduli of the entries of the row unchanged (Section~0.1);
the row sums are therefore functions on $T^{\mathrm{f}}_k$, respectively on $\mathrm{Pz}_k=\mathrm{Pzt}_k/\ell\mathbb{Z}^3$.
For a union $Z$ of $M$-cubes, the pairs that are admissible for $Z$ form a subset $\mathrm{Pz}_k(Z)$ of $\mathrm{Pz}_k$ (Section~I.1.7), and the supremum of the row sums over a subset of the rows does not exceed their supremum over all rows.
Hence, for every such $Z$, the second bound of (6.3) holds for the operator $v\mapsto\delta_{\alpha,A}H_k(A)v$ with its rows restricted to $\mathrm{Pz}_k(Z)$.

\paragraph{6.3. Localised bounds for $k\ge1$.}
Let $z,z',z''$ be sites of $T^0_k$ and $\mathrm{vv}\in\mathbb{C}^4$.
One has $\left|\mathbf{1}_{z'}\mathrm{vv}\right|_{\mathrm{su}}=|\mathrm{vv}|_2\le2|\mathrm{vv}|_c$, and $|F|_\infty\le|F|_{\mathrm{su}}$ for every spinor and every function of pairs $F$ (Section~I.1.5).
Inserting $v=\mathbf{1}_{z'}\mathrm{vv}$ into (1.2) and (1.3), and using (6.1), one obtains
\begin{align}
\left|\mathbf{1}_{\Delta_z}H_k(A)\mathbf{1}_{z'}\mathrm{vv}\right|_\infty
&\le32\,\mathrm{e}\cdot b\cdot a'\cdot\exp\left(-\gamma^{F}\cdot d(z,z')\right)\cdot|\mathrm{vv}|_2
\notag\\
&\le64\,\mathrm{e}\cdot b\cdot a'\cdot\exp\left(-\gamma^{F}\cdot d(z,z')\right)\cdot|\mathrm{vv}|_c
\notag\\
&\le C^{F}\cdot\exp\left(-\gamma^{F}\cdot d(z,z')\right)\cdot|\mathrm{vv}|_c,
\tag{6.4}\\
\left|\mathbf{1}_{\Delta_z}\delta_{\alpha,A}H_k(A)\mathbf{1}_{z'}\mathrm{vv}\right|_\infty
&\le304\,\mathrm{e}\cdot b\cdot a'\cdot\exp\left(-\gamma^{F}\cdot d(z,z')\right)\cdot|\mathrm{vv}|_2
\notag\\
&\le608\,\mathrm{e}\cdot b\cdot a'\cdot\exp\left(-\gamma^{F}\cdot d(z,z')\right)\cdot|\mathrm{vv}|_c
\notag\\
&\le C^{F}\cdot\exp\left(-\gamma^{F}\cdot d(z,z')\right)\cdot|\mathrm{vv}|_c;
\tag{6.5}
\end{align}
in (6.5) the supremum ranges over the short pairs whose first point lies over $\Delta_z$, the pairs inside a cluster included.
By (0.2), read with the norms $|\cdot|_\infty$ and $|\cdot|_c$, the bounds (6.4) and (6.5) state that the operators $\mathbf{1}_{\Delta_z}H_k(A)\mathbf{1}_{z'}$ and $\mathbf{1}_{\Delta_z}\delta_{\alpha,A}H_k(A)\mathbf{1}_{z'}$ (short pairs) have norms at most $C^{F}\cdot\exp\left(-\gamma^{F}\cdot d(z,z')\right)$ for the sup norms $|\cdot|_\infty$.
Now let $(x,x')$ be a long pair with $x$ over $\Delta_z$ and $x'$ over $\Delta_{z''}$.
Then (2.2) and (6.1) give
\begin{equation}
\begin{split}
\left|\left(\delta_{\alpha,A}H_k(A)\mathbf{1}_{z'}\mathrm{vv}\right)(x,x')\right|_c
&\le128\,\mathrm{e}\cdot b\cdot a'\cdot\exp\left(-\gamma^{F}\cdot\min\left\{d(z,z'),d(z'',z')\right\}\right)\cdot|\mathrm{vv}|_2\\
&\le256\,\mathrm{e}\cdot b\cdot a'\cdot\exp\left(-\gamma^{F}\cdot\min\left\{d(z,z'),d(z'',z')\right\}\right)\cdot|\mathrm{vv}|_c\\
&\le C^{F}\cdot\exp\left(-\gamma^{F}\cdot\min\left\{d(z,z'),d(z'',z')\right\}\right)\cdot|\mathrm{vv}|_c.
\end{split}
\tag{6.6}
\end{equation}
With the Euclidean norm on both sides, (2.2) and the inequality $128\,\mathrm{e}\cdot b\cdot a'\le C^{F}$ of (6.1) give the same bound:
\[
\left|\left(\delta_{\alpha,A}H_k(A)\mathbf{1}_{z'}\mathrm{vv}\right)(x,x')\right|_2\le C^{F}\cdot\exp\left(-\gamma^{F}\cdot\min\left\{d(z,z'),d(z'',z')\right\}\right)\cdot|\mathrm{vv}|_2.
\]

\emph{Kernels.}
The kernels are those of Section~0.8;
for a $4\times4$ block $m$, $|m|_c$ is its norm as an operator for $|\cdot|_c$, that is, the largest row sum of the moduli of its entries (Section~I.1.5).
Since $O(t,z')\mathrm{vv}=\left(O\mathbf{1}_{z'}\mathrm{vv}\right)(t)$, the bounds (6.4)--(6.6) amount to
\begin{equation}
\begin{aligned}
\left|H_k(A)(x,z')\right|_c&\le C^{F}\cdot\exp\left(-\gamma^{F}\cdot d([x],z')\right),\\
\left|\left(\delta_{\alpha,A}H_k(A)\right)\left((x,x'),z'\right)\right|_c&\le
\begin{cases}
C^{F}\cdot\exp\left(-\gamma^{F}\cdot d([x],z')\right) & \text{for short pairs},\\
C^{F}\cdot\exp\left(-\gamma^{F}\cdot\min\left\{d([x],z'),d([x'],z')\right\}\right) & \text{for long pairs}.
\end{cases}
\end{aligned}
\tag{6.7}
\end{equation}

\paragraph{6.4. Barred operators for $k\ge1$.}
The connection $-A$ lies in $\tilde{G}_k$ \cite[Section~16.2.7, (G2)]{QEDI};
so (6.3)--(6.6) hold at $-A$.
By (4.1), $\mathsf{C}$ and $\mathsf{C}^{-1}$ preserve the norm $|\cdot|_c$ on $\mathbb{C}^4$, and hence the norm $|\cdot|_\infty$ on $\mathfrak{S}(1,\ell)$, on $\mathfrak{S}^{\mathrm{ch}}(\eta_k,\ell)$ and on the functions of pairs.

\emph{Global bounds.}
Let $v\in\mathfrak{S}(1,\ell)$.
By (4.4), and by (6.3) together with Section~0.8\,(i) at $-A$, applied to the spinor $\mathsf{C}v$,
\begin{align}
\left|\bar{H}_k(A)v\right|_\infty
&=\left|\mathsf{C}^{-1}H_k(-A)\left(\mathsf{C}v\right)\right|_\infty
=\left|H_k(-A)\left(\mathsf{C}v\right)\right|_\infty
\notag\\
&\le C^{F}\cdot|\mathsf{C}v|_\infty=C^{F}\cdot|v|_\infty,
\tag{6.8}\\
\sup_{(x,x')\in\mathrm{Pzt}_k}\left|\left(\delta_{\alpha,-A}\bar{H}_k(A)v\right)(x,x')\right|_c
&=\sup_{(x,x')\in\mathrm{Pzt}_k}\left|\left(\delta_{\alpha,-A}H_k(-A)\left(\mathsf{C}v\right)\right)(x,x')\right|_c
\notag\\
&\le C^{F}\cdot|\mathsf{C}v|_\infty=C^{F}\cdot|v|_\infty.
\tag{6.9}
\end{align}
Section~0.8\,(i), applied to $O=\bar{H}_k(A)$ and to $O=\delta_{\alpha,-A}\bar{H}_k(A)$, turns (6.8) and (6.9) into
\begin{equation}
\left|\bar{H}_k(A)\right|_{1,\infty}\le C^{F},\qquad
\left|\delta_{\alpha,-A}\bar{H}_k(A)\right|_{1,\infty}\le C^{F},
\tag{6.10}
\end{equation}
where the rows are all fermion sites, respectively \emph{all} pairs of $\mathrm{Pzt}_k$, of the kinds (pin), (psh) and (plg).
For a union $Z$ of $M$-cubes, the rows at the pairs of $\mathrm{Pz}_k(Z)$ form a subset of the rows of the second operator;
hence the restriction of this operator to $\mathrm{Pz}_k(Z)$ satisfies the second bound of (6.10).

\emph{Localised bounds.}
Let $z,z',z''$ be sites of $T^0_k$ and $\mathrm{vv}\in\mathbb{C}^4$.
By (4.4),
\begin{gather*}
\bar{H}_k(A)\mathbf{1}_{z'}\mathrm{vv}=\mathsf{C}^{-1}H_k(-A)\mathbf{1}_{z'}\left(\mathsf{C}\mathrm{vv}\right),\\
\delta_{\alpha,-A}\bar{H}_k(A)\mathbf{1}_{z'}\mathrm{vv}=\mathsf{C}^{-1}\delta_{\alpha,-A}H_k(-A)\mathbf{1}_{z'}\left(\mathsf{C}\mathrm{vv}\right).
\end{gather*}
The multiplication by $\mathbf{1}_{\Delta_z}$ commutes with $\mathsf{C}^{-1}$ (Section~4.1), and $\left|\mathsf{C}\mathrm{vv}\right|_c=|\mathrm{vv}|_c$.
Applying (6.4), (6.5) and (6.6) at $-A$, with $\mathsf{C}\mathrm{vv}$ in place of $\mathrm{vv}$, one finds
\begin{align}
\left|\mathbf{1}_{\Delta_z}\bar{H}_k(A)\mathbf{1}_{z'}\mathrm{vv}\right|_\infty
&=\left|\mathbf{1}_{\Delta_z}H_k(-A)\mathbf{1}_{z'}\left(\mathsf{C}\mathrm{vv}\right)\right|_\infty
\notag\\
&\le C^{F}\cdot\exp\left(-\gamma^{F}\cdot d(z,z')\right)\cdot|\mathrm{vv}|_c,
\tag{6.11}\\
\left|\mathbf{1}_{\Delta_z}\delta_{\alpha,-A}\bar{H}_k(A)\mathbf{1}_{z'}\mathrm{vv}\right|_\infty
&=\left|\mathbf{1}_{\Delta_z}\delta_{\alpha,-A}H_k(-A)\mathbf{1}_{z'}\left(\mathsf{C}\mathrm{vv}\right)\right|_\infty
\notag\\
&\le C^{F}\cdot\exp\left(-\gamma^{F}\cdot d(z,z')\right)\cdot|\mathrm{vv}|_c,
\tag{6.12}
\end{align}
where in (6.12) the supremum ranges over the short pairs whose first point lies over $\Delta_z$, the pairs inside a cluster included;
and, for a long pair $(x,x')$ with $x$ over $\Delta_z$ and $x'$ over $\Delta_{z''}$,
\begin{equation}
\begin{split}
\left|\left(\delta_{\alpha,-A}\bar{H}_k(A)\mathbf{1}_{z'}\mathrm{vv}\right)(x,x')\right|_c
&=\left|\left(\delta_{\alpha,-A}H_k(-A)\mathbf{1}_{z'}\left(\mathsf{C}\mathrm{vv}\right)\right)(x,x')\right|_c\\
&\le C^{F}\cdot\exp\left(-\gamma^{F}\cdot\min\left\{d(z,z'),d(z'',z')\right\}\right)\cdot|\mathrm{vv}|_c.
\end{split}
\tag{6.13}
\end{equation}
With the Euclidean norm on both sides, the display of Section~4.3, (b), second part, and the inequality $128\,\mathrm{e}\cdot b\cdot a'\le C^{F}$ of (6.1) give
\[
\left|\left(\delta_{\alpha,-A}\bar{H}_k(A)\mathbf{1}_{z'}\mathrm{vv}\right)(x,x')\right|_2\le C^{F}\cdot\exp\left(-\gamma^{F}\cdot\min\left\{d(z,z'),d(z'',z')\right\}\right)\cdot|\mathrm{vv}|_2.
\]
Since $O(t,z')\mathrm{vv}=\left(O\mathbf{1}_{z'}\mathrm{vv}\right)(t)$, the bounds (6.11)--(6.13) show that the three bounds of (6.7) hold for the kernels $\bar{H}_k(A)(x,z')$ and $\left(\delta_{\alpha,-A}\bar{H}_k(A)\right)\left((x,x'),z'\right)$.

\paragraph{6.5. The level $k=0$.}
At $k=0$ the label FE4$^{\mathrm{ch}}$ consists of (5.1) and (5.2), and $4\le C^{F}$ by (6.2).
In terms of (I.ke3) these bounds read as follows.
The kernels are those of (I.1.4), and the measure is that of (I.1.3) at $\eta_0=1$,
\[
\int\mathrm{d}x''=\frac{1}{48}\sum_{c''}\sum_{g''},
\]
where $x''=(c'',g'')$ ranges over a fundamental domain;
equalities between sites are understood modulo $\ell\mathbb{Z}^3$, with the sign of (I.1.4), which does not affect a modulus.
\begin{itemize}
\item
The kernel of $H_0(A)=1$ is $48\cdot\mathbf{1}[x''=x]\,\mathbf{1}_4$ (I.ke1).
Its row sum is $\frac{1}{48}\cdot48=1$, so that $\left|H_0(A)\right|_{1,\infty}=1$.
\item
In the row $(x,x')$, the kernel of $\delta_{\alpha,A}H_0(A)$ is
\[
d_X(x,x')^{-\alpha}\left(U_A(x,x')\cdot48\cdot\mathbf{1}[x''=x']-48\cdot\mathbf{1}[x''=x]\right)\mathbf{1}_4.
\]
The sites $x$ and $x'$ are distinct modulo $\ell\mathbb{Z}^3$ (Section~0.3: either $0<|c-c'|\le M<\ell$, or $c=c'$ and $g\neq g'$).
Hence the row sum equals $d_X(x,x')^{-\alpha}\left(\left|U_A(x,x')\right|+1\right)\cdot\frac{1}{48}\cdot48$, which is at most $1+\mathrm{e}$ by (0.5) and $d_X\ge1$.
Thus $\left|\delta_{\alpha,A}H_0(A)\right|_{1,\infty}\le1+\mathrm{e}\le4\le C^{F}$.
\item
The kernel of $\bar{H}_0(A)=1$ is $48\cdot\mathbf{1}[x''=x]\,\mathbf{1}_4$, with the row sum $\frac{1}{48}\cdot48=1$.
In the row $(x,x')$, the kernel of $\delta_{\alpha,-A}\bar{H}_0(A)$ is
\[
d_X(x,x')^{-\alpha}\left(U_{-A}(x,x')\cdot48\cdot\mathbf{1}[x''=x']-48\cdot\mathbf{1}[x''=x]\right)\mathbf{1}_4,
\]
with the row sum $d_X(x,x')^{-\alpha}\left(\left|U_{-A}(x,x')\right|+1\right)\cdot\frac{1}{48}\cdot48\le1+\mathrm{e}$;
indeed, $U_{-A}(x,x')=U_A(x,x')^{-1}$ by (I.1.6), $\tau$ being linear, and by (0.5) this number has modulus at most $\mathrm{e}$.
Thus $\left|\bar{H}_0(A)\right|_{1,\infty}=1$ and $\left|\delta_{\alpha,-A}\bar{H}_0(A)\right|_{1,\infty}\le1+\mathrm{e}\le4\le C^{F}$.
\end{itemize}
The weight $\frac{1}{48}$ of the measure cancels the factor $48$ in the kernel of the identity, and no factor $48$ remains.

\subsubsection*{7. Uniformity}

The bounds contain the numbers $32\,\mathrm{e}\cdot b\cdot a'$, $304\,\mathrm{e}\cdot b\cdot a'$, $128\,\mathrm{e}\cdot b\cdot a'$, $\gamma^{F}$, $1+\mathrm{e}$, $4$ and $C^{F}$.
Here $b=1$;
$a'$ and $\gamma^{F}$ are constants of the lines (C4r) of Proposition~14.2\,(a), functions of $(b,\alpha,L,\rho)$;
and $C^{F}$ is a constant of the line (C6).
None of these numbers depends on $k$, $\ell$, $\mu$, $e_k$, $A$, on a site, on a pair or on a spinor.
The facts on which the proof rests hold with these constants for all data (SF) and all $A\in\tilde{G}_k$:
(0.10)--(0.12) by Proposition~2.8\,(c), for every $k\ge1$, $\ell\in\mathrm{Cl}$, $\mu\in(0,1]$ and $e_k\in(0,e_0]$;
(0.3) by (Z2.s1) and (Z2.s3), which are imposed for every $e_k\in(0,e_0]$;
$b_k\le b$ for every $k\ge1$;
and $M<\ell/2$ for every $\ell\in\mathrm{Cl}$.
The mass $\mu$ enters only through $S_k(A)$, that is, through Proposition~2.8.
Proposition~6 involves no further datum:
a cutoff $N$, a charge $e$, a torus or a twist can enter only through the quadruple $(k,\ell,\mu,e_k)$ and through the element $A$ of $\tilde{G}_k$.
Therefore (a)--(f) hold with the stated constants for every level $k$
($k\ge1$ in (a), in the second part of (b), in (c), in the barred forms of these three clauses in (d), and in (f) up to its last sentence;
$k\ge0$ in the first part of (b) and in its barred form in (d);
$k=0$ in (e) and in the last sentence of (f)),
for every side $\ell\in\mathrm{Cl}$, every $\mu\in(0,1]$, every $e_k\in(0,e_0]$ and every $A\in\tilde{G}_k$.
\qed

\subsection*{Conditions on the constants}

\begin{description}
\item[(Z2.s1)]
The inequality $6L\cdot e_k^{\varepsilon}\le1$ for all $e_k\in(0,e_0]$.
It is imposed in Proposition~14.2\,(b), as a condition on $e_0$ (line (C8) of Section~I.8.2).
It gives the first inequality of (0.3), and through it the bound $\exp\left(\frac{3}{2}\mathfrak{i}(A)\right)\le\mathrm{e}$ in (0.5) and (0.7), on which the clauses (a), (b) (second part), (c), (d) and (f) rest.
This bound already follows from the consequence $\frac{3}{2}e_k^{\varepsilon}\le1$ of (Z2.s1), since $\mathfrak{i}(A)<e_k^{\varepsilon}$ (Section~0.4).
\item[(Z2.s3)]
The inequality $3M\cdot e_k^{\varepsilon}\le1$ for all $e_k\in(0,e_0]$.
It is imposed in Proposition~14.2\,(b), as a condition on $e_0$ (line (C8));
it is the third inequality of (A8) and the fifth of (A1).
It gives the second inequality of (0.3), and through it the bound $\left|U_A(x,x')^{\pm1}\right|\le\mathrm{e}$ in (0.5), on which the clauses (b), (c) (long pairs), (d), (e) and (f) rest.
\item[$\alpha>0$]
This is part of the range of the line (C1).
It gives $d_X^{-\alpha}\le1$ for $d_X\ge1$ in (2.1).
\item[$b=1$; $L$ odd, $L\ge5$]
These are imposed in the line (C0) and in (L5).
They give $b_k\le b$ (which needs $L>1$), as well as $k_1=61$ and $0<b(1-1/L)<1$ in Section~6.1.
That $L$ is odd gives the partition of $\mathbb{L}_k$ into the unit blocks (Section~0.1, by Section~I.1.1) and enters the argument of Section~0.5 that bounds the phases.
\item[$M\ge1$]
This follows from (M1) ($M\ge L^2$).
It enters only in Section~0.3, where at $k=0$ it places the pairs inside a cluster, for which $d_X=1$, in $\mathrm{Pzt}_0$.
\item[$\rho>0$]
This is part of the range $\rho\in(0,1)$ in the line (C0).
It gives $c_{\mathrm{re}}>0$ in Section~6.1, hence $a'\ge1$, and through it (6.1) and $4\le C^{F}$.
\item[The conditions that enter through Proposition~2.8\,(a), (c)]
These are (L5), (M1), (M2), the inequality $M\ge108\pi C'_{V}\cdot a'$ (contained in (M5)), (Z2.s1), (Z2.s2), and $\alpha<1$ (through Proposition~2.7\,(d), (e)).
Apart from $\alpha<1$, they are imposed in Proposition~14.2\,(b);
$\alpha<1$ is part of the range of the line (C1), Proposition~14.2\,(a).
The whole list is that of the conditions of Proposition~2.8.
They enter through (0.9)--(0.12).
\end{description}
By Proposition~14.2\,(a), (b), the constants satisfy all these conditions.
The relation $\ell\in\mathrm{Cl}$ is a datum and not a condition on the constants;
it gives $M<\ell/2$ in Sections~0.3 and~0.5.

\section*{Chapter 7. The localised operator families at the levels $k\ge1$ (W1, W2, W5--W16): proofs of Proposition~7}

The statement proved in this chapter is Proposition~7 of Part~II.
The properties W1, W2, W5--W16, the setting in which they are formulated, and the paragraph ``Barred family'' are those of Section~G1.11.

\subsection*{0. Conventions of this chapter}

\subsubsection*{0.1. Data and notation}

Throughout the chapter $(k,\ell,\mu,e_k)$ are data (SF) with $k\ge1$; the data (SF) are those of Section~I.1.8, which Section~G1.1 repeats.
The letter $s=(s_{\mathrm{pc}})_{\mathrm{pc}\in\pi_k}$ denotes a point of $\mathrm{Pd}_k$, the letters $A$ and $\tilde{A}$ denote points of $\tilde{G}_k$, and $X$ denotes a polymer of $D_k$.
The resolvent parameter is $y\in Y_k=\mathbb{R}\cup\{\infty\}$, and $w_k(y)=1$ for every $y\in Y_k$ (Section~I.3.5).
As in Proposition~3, one writes $N:=N_{\gamma^F}$ for the localised norm of Section~G1.3; moreover
\[
w_s:=M^{27\alpha_0}.
\]
The letters $z$, $z'$, $z_1$, $z_2$ denote unit sites, $y'$ a site of the $L$-lattice, $c$, $c'$ photon sites of $\mathbb{L}_k$, and $x=(c,g)$, $x'=(c',g')$ fermion sites.
Further, $\mathrm{vv}$ is a vector of $\mathbb{C}^4$, $v$ an element of $\mathfrak{S}(1,\ell)$, and $f$ an element of $\mathfrak{S}^{\mathrm{ch}}(\eta_k,\ell)$.
The torus $T_k$ is the block torus $(0,\ell,-1)$ of Section~G1.2, and Proposition~2.2 is applied with this block torus.
The symbol $S^{\#}$ stands for $S$, $S'$ or $S^{\mathrm{sy}}$ (Section~G1.9).
The inequalities (N1) are those of Section~G1.3, and (I.POS) is the convention of Section~I.1.4 on where positions are read.
The properties (G1) and (G2) of the domain $\tilde{G}_k$ are the displays (G1) and (G2) of \cite[Section~16.2.7]{QEDI}, and (g0) is the fact (g0) of \cite[Section~17.4.1]{QEDI}.
The conditions (Z2.s1), (L5), (M1) on the constants are those of Proposition~14.2\,(b).

\subsubsection*{0.2. Level $0$}

The data of Proposition~7 have $k\ge1$.
At $k=0$ the objects $b_0$, $B_{0,y}(A)$ and $M_0(X,A)$ are not defined, by (I.L0.9), and there is no short pair;
the families of level~$0$ are the subject of Proposition~8.

\subsubsection*{0.3. Indices on the torus}

Kernels are functions on the covering lattices which satisfy
\[
O(t+\ell e_\mu,t')=O(t,t'+\ell e_\mu)=-O(t,t')
\]
(Sections~G1.2 and G1.11).
Consequently the modulus of $O(t,t')$, and the question whether $O(t,t')$ vanishes, depend only on the classes of $t$ and $t'$ modulo $\ell\mathbb{Z}^3$.
A relation between indices, such as ``$z=z'$'', ``$z$ and $z'$ lie in one $L$-block'' or ``$x$ lies over $Y_1$'', is a relation between sites of the torus:
it is said to hold if it holds for $t$ and $t'+\ell n$ for some $n\in\mathbb{Z}^3$.
A subset of $T_k$ (an $M$-cube, a polymer, the sets $Y$, $Y_1$ below) is identified with its $\ell\mathbb{Z}^3$-periodic preimage in $\mathbb{L}_k$.
The blocks $\Delta_z$ and $\mathrm{Bl}_{k+1}(y')$, where $z\in\mathbb{Z}^3$ and $y'\in L\mathbb{Z}^3$, are the subsets of $\mathbb{L}_k$, and $B(y')$ is the subset of $\mathbb{Z}^3$, defined in Section~I.1.1.
In a sum over an intermediate index this index ranges over one representative of each class modulo $\ell\mathbb{Z}^3$;
the summand $O_1(t,t_1)\cdot O_2(t_1,t')$ does not depend on the choice of the representative, since the two signs cancel.

\subsection*{1. Preliminaries}

\subsubsection*{1.1. Numbers and norms}

\begin{itemize}
\item[(a)] One has $b_-\le b_k\le b$ (Section~I.3.4).
\item[(b)] One has
\[
\mathfrak{i}(A)<e_k^{\varepsilon}\le e_0^{\varepsilon}\le\frac{1}{6L}:
\]
the first inequality is (G1) of \cite[Section~16.2.7]{QEDI}, the second follows from $e_k\le e_0$, which is part of the data (SF), and the third is the condition (Z2.s1).
Consequently
\[
\exp\left(\frac{3}{2}\,\mathfrak{i}(A)\right)\le\exp\left(\frac{1}{4L}\right)\le\mathrm{e}.
\]
\item[(c)] One has $\gamma^FL\le\frac{1}{6}<\log2$; this is among the consequences of the definitions of the constants stated before Proposition~2.1.
Hence
\[
\exp\left(\gamma^FL\right)\le2,\qquad\exp\left(2\gamma^FL\right)\le4.
\]
\item[(d)] $w_s\ge1$, because $M\ge1$ and $\alpha_0>0$.
\item[(e)] For $\mathrm{vv}\in\mathbb{C}^4$ one has $|\mathrm{vv}|_c\le|\mathrm{vv}|_2\le2\,|\mathrm{vv}|_c$ (Section~I.1.5).
Hence $|F|_\infty\le|F|_{\mathrm{su}}\le2\,|F|_\infty$ for every $\mathbb{C}^4$-valued function $F$ of sites or of pairs.
For a $4\times4$ matrix $m$,
\[
|m\,\mathrm{vv}|_c\le|m\,\mathrm{vv}|_2\le|m|_2\cdot|\mathrm{vv}|_2\le2\,|m|_2\cdot|\mathrm{vv}|_c,
\]
that is, $|m|_c\le2\,|m|_2$.
If a linear operator $O$ between two of the spaces of Section~G1.3 satisfies $|OF|_{\mathrm{su}}\le C\cdot|F|_{\mathrm{su}}$ for all $F$, then
\[
|OF|_\infty\le|OF|_{\mathrm{su}}\le C\cdot|F|_{\mathrm{su}}\le2C\cdot|F|_\infty.
\]
\end{itemize}

\subsubsection*{1.2. Unit blocks, blocks of side $L$, $M$-cubes, link-components (photon sites)}

The fact (g0) of \cite[Section~17.4.1]{QEDI} holds at every level $k\ge0$ under the standing conditions of \cite[Section~17.4.1]{QEDI}, namely $L$ odd, $M$ a power of $L$, $M>4L$ and $\ell\in\mathrm{Cl}$.
It states: the numbers $L$ and $M$ are odd;
the $M$-cubes
\[
\left\{c:\ |c-\mathrm{mc}|<\frac{M}{2}\right\},\qquad\mathrm{mc}\in M\mathbb{Z}^3,
\]
form a partition of the set of the sites of $T_k$ and of the set of the sites of $T^0_k$;
and every unit block $\Delta_z$ and every $L$-block $B(y')$ is contained in one $M$-cube.
By \cite[Definition~16.1]{QEDI}, the blocks of one side form a partition of $\mathbb{L}_k$, and every block is the union of the $L^3$ blocks of the next smaller side whose centres it contains.
In particular, $\mathrm{Bl}_{k+1}(y')$ is the union of the blocks $\Delta_z$, $z\in B(y')$.
For a unit site $z$ let $\square(z)$ be the $M$-cube that contains $z$, and let $y'(z)$ be the site of the $L$-lattice with $z\in B(y'(z))$.

\begin{itemize}
\item[(a)] The sets $\Delta_z$, $B(y'(z))$ and $\mathrm{Bl}_{k+1}(y'(z))$ are contained in $\square(z)$.
\begin{proof}
The $M$-cube that contains $B(y'(z))$ contains $z$, and two different $M$-cubes have no site in common;
therefore this cube is $\square(z)$.
Let $z_1\in B(y'(z))$.
The $M$-cube that contains $\Delta_{z_1}$ contains $z_1$, which is a site of $\square(z)$; therefore it is $\square(z)$.
The block $\mathrm{Bl}_{k+1}(y'(z))$ is the union of these blocks $\Delta_{z_1}$, and $\Delta_z$ is one of them.

The same conclusion can be seen in coordinates.
Since $M$ is a power of $L$, $M\ge L^2$ and $L$ is odd, the quotient $M/L$ is an odd integer.
Hence every element of $L\mathbb{Z}$ can be written as $Mi+Lj$ with integers $i$, $j$ such that $|j|\le\frac{1}{2}\left(\frac{M}{L}-1\right)$.
It follows that for every $y'\in L\mathbb{Z}^3$ there is a point $\mathrm{mc}\in M\mathbb{Z}^3$ with $|y'-\mathrm{mc}|\le\frac{M-L}{2}$, and then $|c-y'|<\frac{L}{2}$ implies $|c-\mathrm{mc}|<\frac{M}{2}$.
\end{proof}
\item[(b)] Let $Z$ be a union of $M$-cubes.
If $z$ lies in $Z$, then $\square(z)$, and with it $\Delta_z$, $B(y'(z))$ and $\mathrm{Bl}_{k+1}(y'(z))$, is contained in $Z$.
If $z$ does not lie in $Z$, then $\Delta_z$ does not meet $Z$.
\begin{proof}
Suppose that a site of $\square(z)$ lies in $Z$.
This site lies in one of the $M$-cubes of which $Z$ is the union;
that cube has a site in common with $\square(z)$, so it is $\square(z)$, and $\square(z)\subset Z$.
Applied to the site $z$, this observation together with (a) gives the first assertion.
If a site of $\Delta_z$ lay in $Z$, then, since $\Delta_z\subset\square(z)$ by (a), the cube $\square(z)$ would be contained in $Z$, and $z$ would lie in $Z$.
\end{proof}
\item[(c)] (Link-components.) Two cubes of $\pi_k$ are linked if they have a common face;
the link-components of a subset $\mathscr{Y}$ of $\pi_k$ are the classes of the equivalence relation on $\mathscr{Y}$ generated by this relation (Section~I.1.1; \cite[Definition~5.8]{QEDI}).
A link-component of the union $Y$ of $\mathscr{Y}$ is the union of the cubes of one link-component of $\mathscr{Y}$ (Section~I.1.1).
Therefore the link-components of $Y$ are unions of $M$-cubes, they are pairwise disjoint as sets of sites, and their union is $Y$.
Consequently (b) applies to $Z=Y_1$ for every link-component $Y_1$ of $Y$.
Moreover, an index $t$ whose support consists of a single site (a fermion site $x=(c,g)$, for which $\operatorname{supp}x=\{c\}$ by (I.POS), or a unit site $z$) is contained in a link-component of $Y$ if and only if its site lies in $Y$.
\end{itemize}

\subsubsection*{1.3. The bonds that the phases contain}

In this subsection $y_0,\dots,y_k$ denote sites, as in Section~I.3.2.

\begin{itemize}
\item[(a)] Let $c\in\Delta_z$.
Then $\left(\tau_kA\right)(z,c)$ is a linear form in the numbers $A(b)$, where $b$ ranges over the bonds contained in $\Delta_z$.
Consequently
\[
\varphi_A(z,c)=\exp\left(i\,e_k\eta_k\cdot\left(\tau_kA\right)(z,c)\right)
\]
is an entire function of $A$, and $\varphi_A(z,c)=\varphi_{\tilde{A}}(z,c)$ for any two connections $A$, $\tilde{A}$ that coincide on the bonds contained in $\Delta_z$.
\begin{proof}
Let $c=y_0,y_1,\dots,y_k=z$ be the centres of the blocks of the sides $\eta_k,L\eta_k,\dots,1$ that contain $c$.
The blocks of one side form a partition of $\mathbb{L}_k$, and every block is a union of blocks of the next smaller side;
so the block of side $L^j\eta_k$ that contains $c$ is contained in the block of side $L^{j+1}\eta_k$ that contains $c$.
Hence all these blocks are contained in $\Delta_z$, and $y_0,\dots,y_k$ lie in $\Delta_z$.
One has
\[
\left(\tau_kA\right)(z,c)=\sum_{j<k}\frac{1}{6}\sum_{\pi}A\left(\boldsymbol{\Gamma}^{\pi}(y_{j+1},y_j)\right).
\]
The path $\boldsymbol{\Gamma}^{\pi}(y_{j+1},y_j)$ changes the coordinates one after the other, each one monotonically, from its value at $y_{j+1}$ to its value at $y_j$;
so for every site $c''$ of the path and every $\mu$ the number $c''_\mu$ lies between $y_{j+1,\mu}$ and $y_{j,\mu}$.
The block
\[
\Delta_z=\left\{c''\in\mathbb{L}_k:\ \left|c''_\mu-z_\mu\right|<\frac{1}{2},\ \mu=1,2,3\right\}
\]
is a product of three intervals, and it contains $y_{j+1}$ and $y_j$.
Therefore all sites of the path lie in $\Delta_z$, and its bonds are contained in $\Delta_z$.
\end{proof}
\item[(b)] Let $z\in B(y')$.
Then $\left(\tau A\right)(y',z)$ is a linear form in the numbers $A(b)$, where $b$ ranges over the bonds contained in $\mathrm{Bl}_{k+1}(y')$.
Consequently $U_A(y',z)$ is an entire function of $A$, and $U_A(y',z)=U_{\tilde{A}}(y',z)$ if $A$ and $\tilde{A}$ coincide on the bonds contained in $\mathrm{Bl}_{k+1}(y')$.
\begin{proof}
One has
\[
\left(\tau A\right)(y',z)=\frac{1}{6}\sum_{\pi}A\left(\boldsymbol{\Gamma}^{\pi}(y',z)\right),
\]
where $A$ is summed over the fine bonds of the paths (Section~I.3.2).
The path $\boldsymbol{\Gamma}^{\pi}(y',z)$ changes the coordinates one after the other, each one monotonically, from its value at $y'$ to its value at $z$;
so every site $c''$ of the path has $c''_\mu$ between $y'_\mu$ and $z_\mu$ for $\mu=1,2,3$.
The block
\[
\mathrm{Bl}_{k+1}(y')=\left\{c''\in\mathbb{L}_k:\ \left|c''_\mu-y'_\mu\right|<\frac{L}{2},\ \mu=1,2,3\right\}
\]
is a product of three intervals, and it contains $y'$ and $z$.
Therefore all sites of the path lie in $\mathrm{Bl}_{k+1}(y')$, and its bonds are contained in $\mathrm{Bl}_{k+1}(y')$.
\end{proof}
\end{itemize}

\subsubsection*{1.4. The operators $Q_k(A)\cdot J^{*}$, $J\cdot Q_k^T(-A)$, $P(A)$, $B_k(A)$}

\begin{itemize}
\item[(a)] (Displays; Section~I.3.3, (I.3.1), (I.1.3).)
For $f\in\mathfrak{S}^{\mathrm{ch}}(\eta_k,\ell)$ and $v\in\mathfrak{S}(1,\ell)$ one has
\begin{gather*}
\left(Q_k(A)\cdot J^{*}f\right)(z)=\int_{\Delta_z}\mathrm{d}x\,\varphi_A(z,c)\cdot f(x)=\eta_k^3\sum_{c\in\Delta_z}\varphi_A(z,c)\cdot\frac{1}{48}\sum_{g\in O_{\mathrm{cub}}}f(c,g),\\
\left(J\cdot Q_k^T(-A)\,v\right)(c,g)=\varphi_A(z,c)^{-1}\cdot v(z)\qquad\text{for }c\in\Delta_z\text{ and each of the }48\text{ labels }g\in O_{\mathrm{cub}}.
\end{gather*}
The coefficients are scalars.
Hence both operators commute with every constant $4\times4$ matrix, and by 1.3\,(a) their coefficients are entire functions of $A$.
\item[(b)] (Norm.) By \cite[Lemma~16.2(a)]{QEDI}, $\left|\varphi_A(z,c)^{\pm1}\right|\le\exp\left(\frac{3}{2}\,\mathfrak{i}(A)\right)$.
Together with (a) and 1.1\,(b) this gives
\begin{equation}
\left|J\cdot Q_k^T(-A)\,v\right|_{\mathrm{su}}\le\exp\left(\frac{3}{2}\,\mathfrak{i}(A)\right)\cdot|v|_{\mathrm{su}}\le\mathrm{e}\cdot|v|_{\mathrm{su}}.
\tag{1.1}
\end{equation}
This is (I.Q3) for the Euclidean sup norms: $J$ has norm one by (I.J2), and the phase is a scalar.
\item[(c)] ($P(A)$.) By Section~I.3.3, for $z\in B(y')$,
\[
\left(P(A)\,v\right)(z)=\left(Q^T(-A)\cdot Q(A)\,v\right)(z)=U_A(y',z)^{-1}\cdot L^{-3}\sum_{z'\in B(y')}U_A(y',z')\cdot v(z').
\]
Therefore
\[
P(A)(z,z')=
\begin{cases}
L^{-3}\cdot U_A(y',z)^{-1}\cdot U_A(y',z')\cdot\mathbf{1}_4&\text{if }z,z'\text{ lie in one }L\text{-block }B(y'),\\
0&\text{otherwise}.
\end{cases}
\]
The formula is written for points $z$, $z'$ of $\mathbb{Z}^3$ that lie in one block $B(y')$ of $\mathbb{Z}^3$;
at $(z,z'+\ell n)$ the kernel is multiplied by $(-1)^{n_1+n_2+n_3}$, by 0.3.
\item[(d)] ($B_k(A)$.) The operator $B_k(A)=B_{k,0}(A)$ is defined in Section~I.3.5 as $\left(b_k+b_L\cdot P(A)\right)^{-1}$.
By Proposition~2.2\,(f0), $P(A)^2=P(A)$.
Hence, writing $P:=P(A)$,
\begin{align*}
&\left(b_k+b_L\cdot P\right)\cdot\left(b_k^{-1}\cdot\left(1-P\right)+\left(b_k+b_L\right)^{-1}\cdot P\right)\\
&\qquad=\left(1-P\right)+\frac{b_L}{b_k}\cdot\left(P-P^2\right)+\frac{b_k\cdot P+b_L\cdot P^2}{b_k+b_L}=\left(1-P\right)+P=1.
\end{align*}
Since $\mathfrak{S}(1,\ell)$ has finite dimension, $b_k+b_L\cdot P(A)$ is invertible, and
\begin{equation}
\begin{aligned}
B_k(A)&=b_k^{-1}\cdot\left(1-P(A)\right)+\left(b_k+b_L\right)^{-1}\cdot P(A),\\
B_k(A)(z,z')&=b_k^{-1}\cdot\mathbf{1}[z=z']\cdot\mathbf{1}_4-\frac{b_L}{b_k\left(b_k+b_L\right)}\cdot P(A)(z,z').
\end{aligned}
\tag{1.2}
\end{equation}
Here $\mathbf{1}[z=z']\cdot\mathbf{1}_4$ is the kernel of the identity.
It is written for $z=z'$ in $\mathbb{Z}^3$; at $(z,z+\ell n)$, $n\in\mathbb{Z}^3$, it is multiplied by $(-1)^{n_1+n_2+n_3}$, as in (c), by 0.3.
From (1.2), (c) and 1.3\,(b) one obtains:
\begin{itemize}
\item[(d1)] the kernel of $B_k(A)$ consists of scalar multiples of $\mathbf{1}_4$, so that $B_k(A)$ commutes with every constant $4\times4$ matrix;
\item[(d2)] $B_k(A)(z,z')=0$ unless $z$ and $z'$ lie in one $L$-block;
\item[(d3)] for $z,z'\in B(y')$ the entry $B_k(A)(z,z')$ is an entire function of $A$ which contains only the numbers $A(b)$ with $b$ contained in $\mathrm{Bl}_{k+1}(y')$.
\end{itemize}
\begin{itemize}
\item[(d4)] For a linear operator $O$ on $\mathfrak{S}(1,\ell)$ put
\[
\mathrm{bfun}(O):=b_k^{-1}\cdot\left(1-O\right)+\left(b_k+b_L\right)^{-1}\cdot O.
\]
By (1.2), $B_k(A)=\mathrm{bfun}\left(P(A)\right)$ for every $A\in\tilde{G}_k$.
For every $O$ and every invertible linear operator $\tilde{O}$ on $\mathfrak{S}(1,\ell)$,
\[
\mathrm{bfun}\left(\tilde{O}\cdot O\cdot\tilde{O}^{-1}\right)=\tilde{O}\cdot\mathrm{bfun}(O)\cdot\tilde{O}^{-1},\qquad
\mathrm{bfun}(O)^T=\mathrm{bfun}\left(O^T\right),\qquad
\overline{\mathrm{bfun}(O)}=\mathrm{bfun}\left(\overline{O}\right);
\]
here the transpose is taken for the pairing $\langle\cdot,\cdot\rangle_1$ of Section~I.1.6, and $\overline{O}$ is the operator $v\mapsto\overline{O\,\overline{v}}$, as in (I.O5).
Indeed, $\tilde{O}\cdot1\cdot\tilde{O}^{-1}=1$, $1^T=1$ and $\overline{1}=1$;
each of the maps $O\mapsto\tilde{O}\cdot O\cdot\tilde{O}^{-1}$, $O\mapsto O^T$, $O\mapsto\overline{O}$ is additive;
for every real number $\mathrm{cf}$ one has $\left(\mathrm{cf}\cdot O\right)^T=\mathrm{cf}\cdot O^T$, $\tilde{O}\cdot\left(\mathrm{cf}\cdot O\right)\cdot\tilde{O}^{-1}=\mathrm{cf}\cdot\tilde{O}\cdot O\cdot\tilde{O}^{-1}$ and $\overline{\mathrm{cf}\cdot O}=\mathrm{cf}\cdot\overline{O}$;
and the numbers $b_k^{-1}$ and $\left(b_k+b_L\right)^{-1}$ are real.
\end{itemize}
\end{itemize}

\subsubsection*{1.5. Kernels of the composed operators}

Let $O$ be a linear map from $\mathfrak{S}(1,\ell)$ to $\mathfrak{S}(1,\ell)$ or to $\mathfrak{S}^{\mathrm{ch}}(\eta_k,\ell)$.
Its kernel $O(t,z')$, $z'\in\mathbb{Z}^3$, is the $4\times4$ matrix with $O(t,z')\,\mathrm{vv}=\left(O\,\mathrm{te}\right)(t)$, where $\mathrm{te}\in\mathfrak{S}(1,\ell)$ is the spinor with
\[
\mathrm{te}(z'+\ell n)=(-1)^{n_1+n_2+n_3}\,\mathrm{vv}\qquad\left(n\in\mathbb{Z}^3\right)
\]
and $\mathrm{te}=0$ at all other sites.
The spinor $\mathrm{te}$ changes its sign if $z'$ is replaced by $z'+\ell e_\mu$;
so $O(t,z'+\ell e_\mu)=-O(t,z')$, which is the rule of 0.3.
Since $O\,\mathrm{te}$ is an element of $\mathfrak{S}(1,\ell)$ or of $\mathfrak{S}^{\mathrm{ch}}(\eta_k,\ell)$, one has $\left(O\,\mathrm{te}\right)(t+\ell e_\mu)=-\left(O\,\mathrm{te}\right)(t)$, that is, $O(t+\ell e_\mu,z')=-O(t,z')$.
Every $v\in\mathfrak{S}(1,\ell)$ is the sum of the spinors $\mathrm{te}$ formed with $z'$ and $\mathrm{vv}=v(z')$, where $z'$ ranges over a fundamental domain of $\ell\mathbb{Z}^3$ in $\mathbb{Z}^3$.
Hence $\left(O\,v\right)(t)=\sum_{z'}O(t,z')\,v(z')$, with $z'$ in this domain; this is the kernel of Section~G1.11.
The kernel of an operator with the source $\mathfrak{S}^{\mathrm{ch}}(\eta_k,\ell)$ is that of (I.1.4).

Let $O$ be a linear operator on $\mathfrak{S}^{\mathrm{ch}}(\eta_k,\ell)$ with the kernel $O(x,x')$.
Then
\begin{equation}
\begin{aligned}
\left(O\cdot J\cdot Q_k^T(-A)\right)(x,z')&=\int_{\Delta_{z'}}\mathrm{d}x'\,O(x,x')\cdot\varphi_A(z',c')^{-1}\\
&=\frac{\eta_k^3}{48}\sum_{c'\in\Delta_{z'}}\sum_{g'\in O_{\mathrm{cub}}}\varphi_A(z',c')^{-1}\cdot O\left(x,(c',g')\right),
\end{aligned}
\tag{K1}
\end{equation}
\begin{equation}
\left(Q_k(A)\cdot J^{*}\cdot O\cdot J\cdot Q_k^T(-A)\right)(z,z')=\int_{\Delta_z}\mathrm{d}x\int_{\Delta_{z'}}\mathrm{d}x'\,\varphi_A(z,c)\cdot O(x,x')\cdot\varphi_A(z',c')^{-1},
\tag{K2}
\end{equation}
where each of the two integrals is $\eta_k^3$ times the sum over the photon sites of the unit block times $\frac{1}{48}$ times the sum over the $48$ labels.
For operators $O_1$, $O_2$ on $\mathfrak{S}(1,\ell)$,
\begin{equation}
\left(O_1\cdot O_2\right)(z,z')=\sum_{z_1}O_1(z,z_1)\cdot O_2(z_1,z').
\tag{K3}
\end{equation}
Here $z$, $z'$ are points of $\mathbb{Z}^3$, and $\Delta_z$, $\Delta_{z'}$ are the unit blocks in $\mathbb{L}_k$.
\begin{proof}
Let $\mathrm{te}$ be the spinor formed with $z'$ and $\mathrm{vv}$.

(K1) Let $\mathrm{Fu}^0$ be a fundamental domain of $\ell\mathbb{Z}^3$ in $\mathbb{Z}^3$, and let $\mathrm{Fu}$ be the union of the blocks $\Delta_{z_1}$, $z_1\in\mathrm{Fu}^0$.
The unit blocks form a partition of $\mathbb{L}_k$, and $\Delta_{z_1+\ell n}=\Delta_{z_1}+\ell n$;
so $\mathrm{Fu}$ is a fundamental domain of $\ell\mathbb{Z}^3$ in $\mathbb{L}_k$.
By (I.1.4) and 1.4\,(a),
\begin{align*}
\left(O\cdot J\cdot Q_k^T(-A)\,v\right)(x)&=\int_{\mathrm{Fu}}\mathrm{d}x'\,O(x,x')\cdot\left(J\cdot Q_k^T(-A)\,v\right)(x')\\
&=\sum_{z_1\in\mathrm{Fu}^0}\left[\int_{\Delta_{z_1}}\mathrm{d}x'\,O(x,x')\cdot\varphi_A(z_1,c')^{-1}\right]v(z_1).
\end{align*}
For $v=\mathrm{te}$, and with $\mathrm{Fu}^0$ chosen so that it contains $z'$, the sum reduces to the term $z_1=z'$, which is the assertion.

(K2) Apply the first display of 1.4\,(a) to $f=O\cdot J\cdot Q_k^T(-A)\,\mathrm{te}$, for which $f(x)=\left(O\cdot J\cdot Q_k^T(-A)\right)(x,z')\,\mathrm{vv}$, and insert (K1).

(K3) One has
\[
\left(O_1\cdot O_2\,\mathrm{te}\right)(z)=\sum_{z_1}O_1(z,z_1)\cdot\left(O_2\,\mathrm{te}\right)(z_1)=\sum_{z_1}O_1(z,z_1)\cdot O_2(z_1,z')\,\mathrm{vv}.
\qedhere
\]
\end{proof}

\subsubsection*{1.6. The families are defined}

By Proposition~3\,(a) the series of Section~G1.9 that define $S_\infty(s,A)$, $S^{\mathrm{sy}}_0(s,A)$, $S^{\mathrm{sy}}_\infty(X,A)$ and $S_y(X,A)$ converge absolutely in $N$.
The operator $B_k(A)$ exists by 1.4\,(d).
By (G2) of \cite[Section~16.2.7]{QEDI}, the connection $-A$ lies in $\tilde{G}_k$.
Therefore the operators
\[
H_k(s,A),\ \bar{H}_k(s,A)\colon\mathfrak{S}(1,\ell)\to\mathfrak{S}^{\mathrm{ch}}(\eta_k,\ell),
\]
the operators $\Gamma_k(s,A)$, $M_k(X,A)$ on $\mathfrak{S}(1,\ell)$ and the operator $S_{k,y}(X,A)$ on $\mathfrak{S}^{\mathrm{ch}}(\eta_k,\ell)$ are defined for all $s\in\mathrm{Pd}_k$, $A\in\tilde{G}_k$, $X\in D_k$ and $y\ge0$;
they are linear operators between the spaces that Section~G1.11 assigns to the families (p3), (p4), (p5).

\subsection*{2. W1}

\begin{proof}[Proof of W1 (Section~G1.11)]
By Proposition~3\,(b), for every $y\in Y_k$,
\[
S^{\#}_y(1,A)=S_y(A)\qquad\text{and}\qquad\sum_{X\in D_k}S^{\#}_y(X,A)=S_y(A);
\]
the set $D_k$ is finite, since $T_k$ has $(\ell/M)^3$ $M$-cubes.
By Section~I.3.5,
\[
S_\infty(A)=S_k(A),\qquad S_0(A)=S^0_{k+1}(A),\qquad S_y(A)=S_{k,y}(A).
\]
\begin{itemize}
\item[(i)] By the definition of $H_k(A)$ in Section~I.3.5,
\[
H_k(1,A)=b_k\cdot S_\infty(1,A)\cdot J\cdot Q_k^T(-A)=b_k\cdot S_k(A)\cdot J\cdot Q_k^T(-A)=H_k(A).
\]
\item[(ii)] By Proposition~4\,(d),
\begin{align*}
\Gamma_k(1,A)&=B_k(A)+b_k^2\cdot B_k(A)\cdot Q_k(A)\cdot J^{*}\cdot S^{\mathrm{sy}}_0(1,A)\cdot J\cdot Q_k^T(-A)\cdot B_k(A)\\
&=B_k(A)+b_k^2\cdot B_k(A)\cdot Q_k(A)\cdot J^{*}\cdot S^0_{k+1}(A)\cdot J\cdot Q_k^T(-A)\cdot B_k(A)=\Gamma_k(A).
\end{align*}
\item[(iii)] By the definitions of $D_k(A)$ and $M_k(A)$ in Section~I.3.5,
\begin{align*}
\sum_XM_k(X,A)&=b_k^2\cdot Q_k(A)\cdot J^{*}\cdot\left[\sum_XS^{\mathrm{sy}}_\infty(X,A)\right]\cdot J\cdot Q_k^T(-A)\\
&=b_k^2\cdot Q_k(A)\cdot J^{*}\cdot S_k(A)\cdot J\cdot Q_k^T(-A)=b_k-D_k(A)=M_k(A).
\end{align*}
\item[(iv)] For $y\ge0$,
\[
\sum_XS_{k,y}(X,A)=\sum_XS_y(X,A)=S_y(A)=S_{k,y}(A).
\]
\end{itemize}
\end{proof}

\subsection*{3. W2}

\begin{proof}[Proof of W2 (Section~G1.11)]
Let $\mathrm{Vc}$ be the complex vector space of the connections on $T_k$, and let $\mathbb{C}^{\pi_k}$ be the space of the families $(s_{\mathrm{pc}})$; both spaces have finite dimension.
By (G2) of \cite[Section~16.2.7]{QEDI}, the set $\tilde{G}_k$ is open in $\mathrm{Vc}$;
so the interior of $\mathrm{Pd}_k\times\tilde{G}_k$ is
\[
\left\{s:\ |s_{\mathrm{pc}}|<M^{\alpha_0}\text{ for all }\mathrm{pc}\in\pi_k\right\}\times\tilde{G}_k.
\]
A function with values in the linear maps between two of the spaces $\mathfrak{S}(1,\ell)$, $\mathfrak{S}^{\mathrm{ch}}(\eta_k,\ell)$, which have finite dimension, is continuous (analytic) if and only if the entries of its kernel are continuous (analytic);
sums and products of continuous (analytic) functions of this kind are continuous (analytic).

By 1.4\,(a) and (d3), the remaining entries of $B_k(A)$ being $0$ by (d2), the kernels of $Q_k(A)\cdot J^{*}$, $J\cdot Q_k^T(-A)$ and $B_k(A)$ are entire functions of $A\in\mathrm{Vc}$, and they do not depend on $s$.
By Proposition~3\,(c), $S_\infty(s,A)$ and $S^{\mathrm{sy}}_0(s,A)$ are continuous on $\mathrm{Pd}_k\times\tilde{G}_k$, analytic in $(s,A)$ in the interior, and analytic in $A\in\tilde{G}_k$ for every fixed $s\in\mathrm{Pd}_k$, the boundary points included;
and $S^{\mathrm{sy}}_\infty(X,\cdot)$ and $S_y(X,\cdot)$, $y\ge0$, are analytic on $\tilde{G}_k$.

The operators
\begin{gather*}
H_k(s,A)=b_k\cdot S_\infty(s,A)\cdot J\cdot Q_k^T(-A),\\
\Gamma_k(s,A)=B_k(A)+b_k^2\cdot B_k(A)\cdot Q_k(A)\cdot J^{*}\cdot S^{\mathrm{sy}}_0(s,A)\cdot J\cdot Q_k^T(-A)\cdot B_k(A)
\end{gather*}
are sums of products of these functions; hence they have the three properties just listed.
The operators
\[
M_k(X,A)=b_k^2\cdot Q_k(A)\cdot J^{*}\cdot S^{\mathrm{sy}}_\infty(X,A)\cdot J\cdot Q_k^T(-A)\qquad\text{and}\qquad S_{k,y}(X,A)=S_y(X,A)
\]
are analytic on $\tilde{G}_k$.
\end{proof}

\subsection*{4. W5}

\begin{proof}[Proof of W5 (Section~G1.11)]
Proposition~3\,(a) at $y=\infty$, where $w_k=1$, and (N1) give
\[
\left\|S_\infty(s,A)\right\|_{\mathrm{su}}\le N\left(S_\infty(s,A)\right)\le32a'\cdot w_s,\qquad
\left\|\delta_{\alpha,A}S_\infty(s,A)\right\|_{\mathrm{su}}\le N\left(\delta_{\alpha,A}S_\infty(s,A)\right)\le304a'\cdot w_s.
\]
Here $\delta_{\alpha,A}S_\infty(s,A)$ is the composition of the operator $S_\infty(s,A)$ on $\mathfrak{S}^{\mathrm{ch}}(\eta_k,\ell)$ with the H\"older quotient (I.1.7) (last sentence of Section~G1.6), restricted to the short pairs of Section~G1.3, that is, to the ordered pairs $(x,x')$ of distinct points of $\mathbb{L}^{\mathrm{f}}_k$ with $d_X(x,x')<1$.
Since $k\ge1$, one has $\eta_k<1$;
hence the $48\cdot47$ ordered pairs inside the cluster of a photon site $c$ (those with $c'=c$, $g'\neq g$, $d_X=\eta_k$) are short pairs.
At such a pair, by Section~I.1.8,
\[
\left(\delta_{\alpha,A}H_k(s,A)\,v\right)\left((c,g),(c,g')\right)=\eta_k^{-\alpha}\left[\left(H_k(s,A)\,v\right)(c,g')-\left(H_k(s,A)\,v\right)(c,g)\right],
\]
without a phase.

One has
\[
H_k(s,A)\,v=b_k\cdot S_\infty(s,A)\left[J\cdot Q_k^T(-A)\,v\right],\qquad
\delta_{\alpha,A}H_k(s,A)\,v=b_k\cdot\left[\delta_{\alpha,A}S_\infty(s,A)\right]\left[J\cdot Q_k^T(-A)\,v\right].
\]
With (1.1) and $b_k\le b$ it follows that
\begin{align}
\left|H_k(s,A)\,v\right|_{\mathrm{su}}&\le b_k\cdot32a'\cdot w_s\cdot\mathrm{e}\cdot|v|_{\mathrm{su}}\le32\,\mathrm{e}\cdot b\cdot a'\cdot M^{27\alpha_0}\cdot|v|_{\mathrm{su}},\tag{4.1}\\
\left|\delta_{\alpha,A}H_k(s,A)\,v\right|_{\mathrm{su}}&\le b_k\cdot304a'\cdot w_s\cdot\mathrm{e}\cdot|v|_{\mathrm{su}}\le304\,\mathrm{e}\cdot b\cdot a'\cdot M^{27\alpha_0}\cdot|v|_{\mathrm{su}},\tag{4.2}
\end{align}
the supremum in (4.2) extending over the short pairs.
These are the two bounds in the Euclidean norms that Proposition~7 asserts;
both constants are at most $304\,\mathrm{e}\cdot\max(1,b)\cdot a'\cdot M^{27\alpha_0}\le C_1$.
By 1.1\,(e),
\[
\left|H_k(s,A)\,v\right|_\infty\le2C_1\cdot|v|_\infty,\qquad
\left|\left(\delta_{\alpha,A}H_k(s,A)\,v\right)(x,x')\right|_c\le2C_1\cdot|v|_\infty
\]
for all short pairs $(x,x')$, the pairs inside a cluster included.
\end{proof}

\subsection*{5. W6}

The property W6 is stated in Section~G1.11; for $k\ge1$ its indices are unit sites.

\paragraph*{The kernel of $\Gamma_k(s,A)$.}
Put
\[
G:=Q_k(A)J^{*}\cdot S^{\mathrm{sy}}_0(s,A)\cdot JQ_k^{T}(-A),
\]
an operator on $\mathfrak{S}(1,\ell)$.
Proposition~5\,(c), the inequality (N1) of Section~G1.3 and Proposition~3\,(a), taken at $y=0$, where $w_k=1$, give
\begin{equation}
\begin{split}
\left|G(z_1,z_2)\right|_2
&\le\mathrm{e}\cdot\left\|\mathbf{1}_{\Delta_{z_1}}S^{\mathrm{sy}}_0(s,A)\,\mathbf{1}_{\Delta_{z_2}}\right\|_{\mathrm{su}}
\le\mathrm{e}\cdot\exp\left(-\gamma^{F}d(z_1,z_2)\right)\cdot N\left(S^{\mathrm{sy}}_0(s,A)\right)\\
&\le 32\mathrm{e}\cdot a'\cdot w_s\cdot\exp\left(-\gamma^{F}d(z_1,z_2)\right).
\end{split}
\tag{5.1}
\end{equation}
Proposition~5\,(b), at the value $0$ of the resolvent parameter, where $B_{k,0}(A)=B_k(A)$, states that $B_k(A)(z,z')=0$ unless $z$ and $z'$ lie in one $L$-block, and that
\begin{equation}
\begin{aligned}
&\left|B_k(A)(z,z')\right|_2\le b_k^{-1}\cdot\left(\mathbf{1}[z=z']+2\mathrm{e}\cdot L^{-3}\right),\\
&\sum_{z_1}\left|B_k(A)(z,z_1)\right|_2\le\frac{1+2\mathrm{e}}{b_k},\qquad
\sum_{z_2}\left|B_k(A)(z_2,z')\right|_2\le\frac{1+2\mathrm{e}}{b_k}.
\end{aligned}
\tag{5.2}
\end{equation}
By (K3),
\[
\left(B_k(A)\cdot G\cdot B_k(A)\right)(z,z')=\sum_{z_1,z_2}B_k(A)(z,z_1)\cdot G(z_1,z_2)\cdot B_k(A)(z_2,z').
\]
In a non-vanishing term of this sum the site $z_1$ lies in the $L$-block of $z$, and $z_2$ lies in the $L$-block of $z'$.
Two sites of one $L$-block are at distance at most $L-1$ from each other.
Consequently
\[
d(z,z')\le d(z,z_1)+d(z_1,z_2)+d(z_2,z')\le d(z_1,z_2)+2L,
\]
and 1.1\,(c) yields
\[
\exp\left(-\gamma^{F}d(z_1,z_2)\right)\le\exp\left(2\gamma^{F}L\right)\cdot\exp\left(-\gamma^{F}d(z,z')\right)\le 4\exp\left(-\gamma^{F}d(z,z')\right).
\]
The norm $|\cdot|_2$ of $4\times4$ matrices is submultiplicative.
Together with (5.1) and (5.2) one obtains
\begin{equation}
\begin{split}
b_k^2\cdot\left|\left(B_k(A)\cdot G\cdot B_k(A)\right)(z,z')\right|_2
&\le b_k^2\cdot 32\mathrm{e}\cdot a'\cdot w_s\cdot 4\exp\left(-\gamma^{F}d(z,z')\right)\\
&\qquad\cdot\left[\sum_{z_1}\left|B_k(A)(z,z_1)\right|_2\right]\cdot\left[\sum_{z_2}\left|B_k(A)(z_2,z')\right|_2\right]\\
&\le 128\mathrm{e}\cdot(1+2\mathrm{e})^2\cdot a'\cdot w_s\cdot\exp\left(-\gamma^{F}d(z,z')\right);
\end{split}
\tag{5.3}
\end{equation}
in the last step the factor $b_k^2$ cancels against the two factors $b_k^{-1}$ that (5.2) supplies.

The first term of $\Gamma_k(s,A)$ is $B_k(A)$.
By (5.2) and $b_k\ge b_-$ its entries satisfy $\left|B_k(A)(z,z')\right|_2\le\frac{1+2\mathrm{e}}{b_-}$, and $B_k(A)(z,z')$ vanishes unless $z$ and $z'$ lie in one $L$-block.
In the latter case $d(z,z')\le L$, so that
\[
1\le\exp\left(\gamma^{F}L\right)\cdot\exp\left(-\gamma^{F}d(z,z')\right)\le 2\exp\left(-\gamma^{F}d(z,z')\right).
\]
Therefore
\begin{equation}
\left|B_k(A)(z,z')\right|_2\le\frac{2(1+2\mathrm{e})}{b_-}\cdot\exp\left(-\gamma^{F}d(z,z')\right).
\tag{5.4}
\end{equation}
The kernel of $\Gamma_k(s,A)=B_k(A)+b_k^2\cdot B_k(A)\cdot G\cdot B_k(A)$ is the sum of the kernels estimated in (5.4) and (5.3).
Since $w_s\ge1$, and since $C_{\Gamma}=\frac{2(1+2\mathrm{e})}{b_-}+128\mathrm{e}\cdot(1+2\mathrm{e})^2\cdot a'$ by its definition in Proposition~5, it follows that
\begin{equation}
\begin{split}
\left|\Gamma_k(s,A)(z,z')\right|_2
&\le\left[\frac{2(1+2\mathrm{e})}{b_-}+128\mathrm{e}\cdot(1+2\mathrm{e})^2\cdot a'\cdot w_s\right]\cdot\exp\left(-\gamma^{F}d(z,z')\right)\\
&\le C_{\Gamma}\cdot M^{27\alpha_0}\cdot\exp\left(-\gamma^{F}d(z,z')\right).
\end{split}
\tag{5.5}
\end{equation}

\paragraph*{The kernel of $M_k(X,A)$.}
Proposition~5\,(c), the third bound of Proposition~3\,(a) at $y=\infty$, where $w_k=1$, and the inequality $b_k\le b$ give
\begin{equation}
\begin{split}
\left|M_k(X,A)(z,z')\right|_2
&\le b_k^2\cdot\mathrm{e}\cdot\left\|\mathbf{1}_{\Delta_z}S^{\mathrm{sy}}_{\infty}(X,A)\,\mathbf{1}_{\Delta_{z'}}\right\|_{\mathrm{su}}\\
&\le 32\mathrm{e}\cdot b^2\cdot a'\cdot\exp(27\kappa')\cdot\exp\left(-\gamma^{F}d(z,z')-\kappa'\cdot d_M(X)\right).
\end{split}
\tag{5.6}
\end{equation}
The inequalities (5.5) and (5.6) are the bounds in the Euclidean norm that Proposition~7 asserts for $\Gamma_k(s,A)$ and $M_k(X,A)$.
By the definition of $C_1$ in Proposition~7,
\[
C_{\Gamma}\cdot M^{27\alpha_0}\le C_1,\qquad
32\mathrm{e}\cdot b^2\cdot a'\cdot\exp(27\kappa')\le 32\mathrm{e}\cdot\max\left(1,b^2\right)\cdot a'\cdot\exp(27\kappa')\le C_1.
\]
As $|m|_c\le 2|m|_2$ for every $4\times4$ matrix $m$ (1.1\,(e)), one concludes that
\[
\left|\Gamma_k(s,A)(z,z')\right|_c\le 2C_1\exp\left(-\gamma^{F}d(z,z')\right),\qquad
\left|M_k(X,A)(z,z')\right|_c\le 2C_1\exp\left(-\gamma^{F}d(z,z')-\kappa'\cdot d_M(X)\right).
\]

\subsection*{6. W7}

The property W7 is stated in Section~G1.11.
Since $k\ge1$, it is its clause for $k\ge1$ that has to be proved.
Let $y\in[0,\infty)$; then $w_k(y)=1$, and $S_{k,y}(X,A)=S_y(X,A)$.
The third bound of Proposition~3\,(a), for $S^{\#}=S$, gives
\begin{equation}
\left|\mathbf{1}_{\Delta_z}S_{k,y}(X,A)\,\mathbf{1}_{\Delta_{z'}}f\right|_{\mathrm{su}}\le 32a'\cdot\exp(27\kappa')\cdot\exp\left(-\gamma^{F}d(z,z')-\kappa'\cdot d_M(X)\right)\cdot|f|_{\mathrm{su}},
\tag{6.1}
\end{equation}
and $32a'\cdot\exp(27\kappa')\le 32\mathrm{e}\cdot\max\left(1,b^2\right)\cdot a'\cdot\exp(27\kappa')\le C_1$.
By 1.1\,(e) it follows that
\[
\left|\mathbf{1}_{\Delta_z}S_{k,y}(X,A)\,\mathbf{1}_{\Delta_{z'}}f\right|_{\infty}\le 2C_1\exp\left(-\gamma^{F}d(z,z')-\kappa'\cdot d_M(X)\right)\cdot|f|_{\infty}\qquad\text{for all }y\ge0.
\]
Here $\mathbf{1}_{\Delta_z}$ is the multiplication by the indicator function of the orbit of $\Delta_z\times O_{\mathrm{cub}}$ (Section~G1.3 and (I.POS)).

\subsection*{7. The setting of W8--W11; three facts}

The setting of Sections~8--10, denoted (Loc), is the following.
\begin{itemize}
\item[(Loc)] $\mathscr{Y}$ is a subset of $\pi_k$, and $Y$ is its union; $s\in\mathrm{Pd}_k$ satisfies $s_{\mathrm{pc}}=0$ for every $\mathrm{pc}\in\pi_k$ outside $\mathscr{Y}$; $Y_1$ is a link-component of $Y$ (1.2\,(c)); $\tilde s\in\mathrm{Pd}_k$ vanishes outside $\mathscr{Y}$ and coincides with $s$ on the cubes of $Y_1$; $\tilde A\in\tilde{G}_k$ coincides with $A$ on the bonds that meet $Y_1$.
\end{itemize}
The pair $(\tilde s,\tilde A)$ fulfils the hypotheses imposed on $(s,A)$, because $\tilde s$ vanishes outside $\mathscr{Y}$.
Hence every assertion below of the form ``the entry vanishes unless \dots'' holds for $(\tilde s,\tilde A)$ as well.

\begin{named}{Fact (Fa1)}
Let $S^{\#}_y$ stand either for $S_{\infty}$, that is $S^{\#}=S$ at $y=\infty$, or for $S^{\mathrm{sy}}_0$, that is $S^{\#}=S^{\mathrm{sy}}$ at $y=0$.
Then the following holds.
\begin{itemize}
\item \emph{Rows.} If $x$ lies over $Y_1$, then $S^{\#}_y(s,A)(x,x')=0$ unless $x'$ lies over $Y_1$, and $S^{\#}_y(s,A)(x,\cdot\,)=S^{\#}_y(\tilde s,\tilde A)(x,\cdot\,)$.
\item \emph{Columns.} If $x'$ lies over $Y_1$, then $S^{\#}_y(s,A)(x,x')=0$ unless $x$ lies over $Y_1$, and $S^{\#}_y(s,A)(\cdot\,,x')=S^{\#}_y(\tilde s,\tilde A)(\cdot\,,x')$.
\item \emph{Third clause.} The row of $S^{\#}_y(s,A)$ at $x=(c,g)$ vanishes unless $\mathscr{Y}$ contains all cubes of $\pi_k$ that meet $K_{c}$; in particular it vanishes if $c\notin Y$.
\end{itemize}
\end{named}
These are the assertions of Proposition~3\,(i) for the two choices of $(S^{\#},y)$.

\begin{named}{Fact (Fa2)}
Let $z\in Y_1$.
Then $\Delta_z$, $B(y'(z))$ and $\mathrm{Bl}_{k+1}(y'(z))$ are contained in $Y_1$; one has $\varphi_{\tilde A}(z,c)=\varphi_A(z,c)$ for every $c\in\Delta_z$; and $B_k(\tilde A)(z,z_1)=B_k(A)(z,z_1)$, $B_k(\tilde A)(z_1,z)=B_k(A)(z_1,z)$ for every unit site $z_1$.
If $z\notin Y_1$, then no fermion site over $\Delta_z$ lies over $Y_1$.
\end{named}
\begin{proof}
By 1.2\,(c) the set $Y_1$ is a union of $M$-cubes.
So 1.2\,(b) applies to $Z=Y_1$; together with (I.POS) it gives the three inclusions and the last assertion.
A bond contained in $\Delta_z$ or in $\mathrm{Bl}_{k+1}(y'(z))$ has both end points in $Y_1$; hence it meets $Y_1$ (Section~I.1.1), and $\tilde A=A$ on this bond.
The equality of the phases is therefore a consequence of 1.3\,(a).
If $z_1\notin B(y'(z))$, the four entries $B_k(\tilde A)(z,z_1)$, $B_k(A)(z,z_1)$, $B_k(\tilde A)(z_1,z)$, $B_k(A)(z_1,z)$ vanish by 1.4\,(d2).
If $z_1\in B(y'(z))$, they coincide in pairs by 1.4\,(d3).
\end{proof}

\begin{named}{Fact (Fa3)}
Let $G(s,A):=Q_k(A)J^{*}\cdot S^{\mathrm{sy}}_0(s,A)\cdot JQ_k^{T}(-A)$.
By (K2), with $x=(c,g)$ and $x'=(c',g')$,
\begin{equation}
G(s,A)(z_1,z_2)=\int_{\Delta_{z_1}}\mathrm{d}x\int_{\Delta_{z_2}}\mathrm{d}x'\,\varphi_A(z_1,c)\cdot S^{\mathrm{sy}}_0(s,A)(x,x')\cdot\varphi_A(z_2,c')^{-1}.
\tag{7.1}
\end{equation}
The following holds.
\begin{itemize}
\item \emph{Rows.} If $z_1\in Y_1$, then $G(s,A)(z_1,z_2)=0$ unless $z_2\in Y_1$, and $G(s,A)(z_1,z_2)=G(\tilde s,\tilde A)(z_1,z_2)$ for every $z_2$.
\item \emph{Columns.} If $z_2\in Y_1$, then $G(s,A)(z_1,z_2)=0$ unless $z_1\in Y_1$, and $G(s,A)(z_1,z_2)=G(\tilde s,\tilde A)(z_1,z_2)$ for every $z_1$.
\end{itemize}
\end{named}
\begin{proof}
Rows.
Let $z_1\in Y_1$.
In (7.1) every $x$ lies over $\Delta_{z_1}$, hence over $Y_1$ by (Fa2).
If $z_2\notin Y_1$, then no $x'$ over $\Delta_{z_2}$ lies over $Y_1$ (Fa2), and the rows clause of (Fa1) shows that the integrand vanishes identically; the same reasoning applies to $(\tilde s,\tilde A)$.
If $z_2\in Y_1$, then (Fa2) gives $\varphi_{\tilde A}(z_1,c)=\varphi_A(z_1,c)$ and $\varphi_{\tilde A}(z_2,c')=\varphi_A(z_2,c')$, while the rows clause of (Fa1) gives $S^{\mathrm{sy}}_0(\tilde s,\tilde A)(x,x')=S^{\mathrm{sy}}_0(s,A)(x,x')$; thus the two integrands are equal.

Columns.
Let $z_2\in Y_1$.
In (7.1) every $x'$ lies over $\Delta_{z_2}$, hence over $Y_1$.
If $z_1\notin Y_1$, then no $x$ over $\Delta_{z_1}$ lies over $Y_1$, and the integrand vanishes identically by the columns clause of (Fa1); likewise for $(\tilde s,\tilde A)$.
If $z_1\in Y_1$, the phases coincide by (Fa2), and $S^{\mathrm{sy}}_0(\tilde s,\tilde A)(x,x')=S^{\mathrm{sy}}_0(s,A)(x,x')$ by the columns clause of (Fa1).
\end{proof}

\subsection*{8. W8 (rows)}

The property W8 is stated in Section~G1.11.
Let the data be as in (Loc).

\paragraph*{The case $\mathrm{Op}=H_k$.}
The row index is a fermion site $x=(c,g)$ with $c\in Y_1$, and the column index is a unit site $z'$.
For every fermion site $x$, over $Y_1$ or not, and every unit site $z'$, (K1) gives
\begin{equation}
\begin{split}
H_k(s,A)(x,z')&=b_k\int_{\Delta_{z'}}\mathrm{d}x'\,S_{\infty}(s,A)(x,x')\cdot\varphi_A(z',c')^{-1}\\
&=b_k\cdot\frac{\eta_k^3}{48}\sum_{c'\in\Delta_{z'}}\sum_{g'\in O_{\mathrm{cub}}}\varphi_A(z',c')^{-1}\cdot S_{\infty}(s,A)\left(x,(c',g')\right).
\end{split}
\tag{8.1}
\end{equation}
If $z'\notin Y_1$, no $x'$ over $\Delta_{z'}$ lies over $Y_1$ (Fa2), and each term of (8.1) vanishes by the rows clause of (Fa1).
So $H_k(s,A)(x,z')=0$ unless $z'$ is contained in $Y_1$, and the same is true for $(\tilde s,\tilde A)$.
If $z'\in Y_1$, then $\varphi_{\tilde A}(z',c')=\varphi_A(z',c')$ for $c'\in\Delta_{z'}$ by (Fa2), and $S_{\infty}(\tilde s,\tilde A)(x,x')=S_{\infty}(s,A)(x,x')$ by (Fa1); hence (8.1) takes the same value for $(s,A)$ and for $(\tilde s,\tilde A)$.
It follows that $H_k(s,A)(x,\cdot\,)=H_k(\tilde s,\tilde A)(x,\cdot\,)$.

\paragraph*{The case $\mathrm{Op}=\Gamma_k$.}
The row index is a unit site $z\in Y_1$.
With $G(s,A)$ as in (Fa3), (K3) gives for all unit sites $z,z'$, whether $z\in Y_1$ or not,
\begin{equation}
\Gamma_k(s,A)(z,z')=B_k(A)(z,z')+b_k^2\sum_{z_1,z_2}B_k(A)(z,z_1)\cdot G(s,A)(z_1,z_2)\cdot B_k(A)(z_2,z').
\tag{8.2}
\end{equation}
If $B_k(A)(z,z')\ne0$, then $z'\in B(y'(z))$ by 1.4\,(d2), and $B(y'(z))\subset Y_1$ by (Fa2).
If a term of the double sum does not vanish, then first $z_1\in B(y'(z))\subset Y_1$; next $z_2\in Y_1$ by the rows clause of (Fa3); finally $z'\in B(y'(z_2))$, and this set is contained in $Y_1$ by (Fa2).
Thus $\Gamma_k(s,A)(z,z')=0$ unless $z'$ is contained in $Y_1$, and the same is true for $(\tilde s,\tilde A)$.

Now let $z'\in Y_1$.
For $(s,A)$ as well as for $(\tilde s,\tilde A)$, only the terms of (8.2) with $z_1\in Y_1$ and $z_2\in Y_1$ contribute.
For these terms
\[
B_k(\tilde A)(z,z_1)=B_k(A)(z,z_1),\qquad G(\tilde s,\tilde A)(z_1,z_2)=G(s,A)(z_1,z_2),\qquad B_k(\tilde A)(z_2,z')=B_k(A)(z_2,z'),
\]
by (Fa2) with $z\in Y_1$, by the rows clause of (Fa3) with $z_1\in Y_1$, and by (Fa2) with $z'\in Y_1$, respectively; moreover $B_k(\tilde A)(z,z')=B_k(A)(z,z')$.
Hence $\Gamma_k(s,A)(z,\cdot\,)=\Gamma_k(\tilde s,\tilde A)(z,\cdot\,)$.

\subsection*{9. W9 (columns)}

The property W9 is stated in Section~G1.11.
Let the data be as in (Loc).

\paragraph*{The case $\mathrm{Op}=H_k$.}
The column index is a unit site $z'\in Y_1$.
Every $x'$ in (8.1) lies over $\Delta_{z'}$, hence over $Y_1$ (Fa2).
By the columns clause of (Fa1), $S_{\infty}(s,A)(x,x')=0$ unless $x$ lies over $Y_1$; so $H_k(s,A)(x,z')=0$ unless $x$ is contained in $Y_1$.
For every fermion site $x$ one has $S_{\infty}(\tilde s,\tilde A)(x,x')=S_{\infty}(s,A)(x,x')$ for all $x'$ over $\Delta_{z'}$, by the columns clause of (Fa1), and $\varphi_{\tilde A}(z',c')=\varphi_A(z',c')$ by (Fa2).
Hence (8.1) takes the same value for $(s,A)$ and for $(\tilde s,\tilde A)$, that is, $H_k(s,A)(\cdot\,,z')=H_k(\tilde s,\tilde A)(\cdot\,,z')$.

\paragraph*{The case $\mathrm{Op}=\Gamma_k$.}
The column index is a unit site $z'\in Y_1$.
If $B_k(A)(z,z')\ne0$, then $z\in B(y'(z'))$, and this set is contained in $Y_1$ (1.4\,(d2) and (Fa2)).
If a term of the double sum in (8.2) does not vanish, then first $z_2\in B(y'(z'))\subset Y_1$; next $z_1\in Y_1$ by the columns clause of (Fa3); finally $z\in B(y'(z_1))$, a set contained in $Y_1$.
Thus $\Gamma_k(s,A)(z,z')=0$ unless $z$ is contained in $Y_1$, and the same is true for $(\tilde s,\tilde A)$.

Now let $z\in Y_1$.
For $(s,A)$ as well as for $(\tilde s,\tilde A)$, only the terms of (8.2) with $z_1\in Y_1$ and $z_2\in Y_1$ contribute, and for them
\[
B_k(\tilde A)(z,z_1)=B_k(A)(z,z_1),\qquad G(\tilde s,\tilde A)(z_1,z_2)=G(s,A)(z_1,z_2),\qquad B_k(\tilde A)(z_2,z')=B_k(A)(z_2,z'),
\]
by (Fa2) with $z\in Y_1$, by the columns clause of (Fa3) with $z_2\in Y_1$, and by (Fa2) with $z'\in Y_1$, respectively; moreover $B_k(\tilde A)(z,z')=B_k(A)(z,z')$.
Hence $\Gamma_k(s,A)(\cdot\,,z')=\Gamma_k(\tilde s,\tilde A)(\cdot\,,z')$.

\subsection*{10. W10 and W11}

The properties W10 and W11 are stated in Section~G1.11.
Of the setting (Loc) only $\mathscr{Y}$, $Y$ and $s$ enter.
Let $x=(c,g)$.
By (8.1) the row $H_k(s,A)(x,\cdot\,)$ is built from the row $S_{\infty}(s,A)(x,\cdot\,)$; so it vanishes whenever the latter vanishes.

\paragraph*{W11.}
By the third clause of (Fa1), the row $S_{\infty}(s,A)(x,\cdot\,)$ vanishes unless $\mathscr{Y}$ contains all cubes of $\pi_k$ that meet $K_{c}$; and $K_{x}=K_{c}$ by (I.POS).

\paragraph*{W10.}
By 1.2\,(c), the site $x$ is contained in no link-component of $Y$ if and only if $c\notin Y$.
In this case the row $S_{\infty}(s,A)(x,\cdot\,)$ vanishes by the third clause of (Fa1).
(Indeed, $c\in K_{c}$; so the $M$-cube of $c$ meets $K_{c}$, and this cube does not belong to $\mathscr{Y}$.)

\subsection*{11. W12}

The property W12 is stated in Section~G1.11.
Let $S^{\#}$ be $S$ or $S^{\mathrm{sy}}$, and let $y\in Y_{k}$.
By Section~G1.9,
\[
S^{\#}_y(X,A)=\sum_{\mathrm{wk}\,:\,X_{\mathrm{wk}}=X}S^{\#}_{y,\mathrm{wk}}(A),
\]
and the series converges absolutely with respect to $N(\cdot)$ (Proposition~3\,(a)).
The space $\mathfrak{S}^{\mathrm{ch}}(\eta_k,\ell)$ has finite dimension, and for fixed $(x,x')$ the entries of the kernel $O(x,x')$ are linear functionals of $O$; therefore the kernel of the sum is the sum of the kernels.
By Proposition~3\,(d) the kernels of $S_{y,\mathrm{wk}}(A)$ and of $S'_{y,\mathrm{wk}}(A)$, and hence the kernel of $S^{\mathrm{sy}}_{y,\mathrm{wk}}(A)=\frac{1}{2}\left(S_{y,\mathrm{wk}}(A)+S'_{y,\mathrm{wk}}(A)\right)$, vanish unless both arguments lie over $X_{\mathrm{wk}}=X$.
So
\begin{equation}
S^{\#}_y(X,A)(x,x')=0\qquad\text{unless $x$ and $x'$ lie over $X$.}
\tag{11.1}
\end{equation}
Let $A,\tilde A\in\tilde{G}_k$ coincide on the bonds that meet $X$.
A bond with both end points in $X$ meets $X$.
Hence $S^{\#}_{y,\mathrm{wk}}(A)=S^{\#}_{y,\mathrm{wk}}(\tilde A)$ for every walk with $X_{\mathrm{wk}}=X$ (Proposition~3\,(d)), and
\begin{equation}
S^{\#}_y(X,A)=S^{\#}_y(X,\tilde A).
\tag{11.2}
\end{equation}

\paragraph*{The family $S_{k,y}(X,A)$, $y\ge0$.}
For $S^{\#}=S$ the relations (11.1) and (11.2) are the assertions of W12.

\paragraph*{The family $M_k(X,A)$.}
By (K2),
\begin{equation}
M_k(X,A)(z,z')=b_k^2\int_{\Delta_z}\mathrm{d}x\int_{\Delta_{z'}}\mathrm{d}x'\,\varphi_A(z,c)\cdot S^{\mathrm{sy}}_{\infty}(X,A)(x,x')\cdot\varphi_A(z',c')^{-1}.
\tag{11.3}
\end{equation}
The polymer $X$ is a union of $M$-cubes.
If $z\notin X$, then $\Delta_z$ does not meet $X$ (1.2\,(b)), no $x$ over $\Delta_z$ lies over $X$, and the integrand vanishes by (11.1).
If $z'\notin X$, then $\Delta_{z'}$ does not meet $X$ (1.2\,(b)), no $x'$ over $\Delta_{z'}$ lies over $X$, and the integrand vanishes by (11.1).
Thus $M_k(X,A)(z,z')=0$ unless $z\in X$ and $z'\in X$, and the same is true for $\tilde A$.
If $z,z'\in X$, then $\Delta_z$ and $\Delta_{z'}$ are contained in $X$ (1.2\,(b)), and the bonds contained in them meet $X$.
Hence $\varphi_{\tilde A}(z,c)=\varphi_A(z,c)$ and $\varphi_{\tilde A}(z',c')=\varphi_A(z',c')$ by 1.3\,(a), and together with (11.2) the integrands of (11.3) for $A$ and for $\tilde A$ coincide.
Therefore $M_k(X,A)=M_k(X,\tilde A)$.

\subsection*{12. W13}

The property W13 is stated in Section~G1.11; since $k\ge1$, its clauses for $k\ge1$ are the ones to be proved.
Let $\boldsymbol{\lambda}$ be a real function of $c$ on $T_k$, that is, an $\ell\mathbb{Z}^3$-periodic function on $\mathbb{L}_k$, and let $\boldsymbol{\lambda}^{\mathrm{u}}$ be its restriction to $T^0_k$.
Put
\[
\mathrm{uf}:=\exp\left(ie_k\boldsymbol{\lambda}\right),\qquad\mathrm{uf}^{0}:=\exp\left(ie_k\boldsymbol{\lambda}^{\mathrm{u}}\right).
\]
By (I.POS), $\mathrm{uf}$ acts on $\mathfrak{S}^{\mathrm{ch}}(\eta_k,\ell)$ by $(\mathrm{uf}\cdot f)(c,g)=\mathrm{uf}(c)\cdot f(c,g)$, with one and the same phase on the $48$ sites over $c$.
The connection $A-d\boldsymbol{\lambda}$ lies in $\tilde{G}_k$ by the property (G2) of \cite[Section~16.2.7]{QEDI}, which states in particular that $A-d\boldsymbol{\lambda}$, $-A$, $\overline{A}$ and $\tilde gA$ belong to $\tilde{G}_k$ for every $A\in\tilde{G}_k$, every real scalar $\boldsymbol{\lambda}$ and every $\tilde g\in\mathbb{G}_k$.

The connection $\mathrm{Acl}:=d\boldsymbol{\lambda}$ is real, and $d\,\mathrm{Acl}=0$, since
\[
(dd\boldsymbol{\lambda})_{\mu\nu}=\partial_\mu\partial_\nu\boldsymbol{\lambda}-\partial_\nu\partial_\mu\boldsymbol{\lambda}=0,
\]
the two difference quotients commuting.
The function $\mathrm{uf}$ is $\ell\mathbb{Z}^3$-periodic with $|\mathrm{uf}|=1$, and on every bond $(c,c')$
\[
\mathrm{uf}(c')=\mathrm{uf}(c)\cdot\exp\left(ie_k\eta_k\cdot\mathrm{Acl}(c,c')\right),
\]
because $\eta_k\cdot(d\boldsymbol{\lambda})(c,c')=\boldsymbol{\lambda}(c')-\boldsymbol{\lambda}(c)$.
This is the example mentioned in Proposition~2.2\,(a).
Proposition~2.2\,(a) on $T_k$ therefore gives
\begin{gather*}
JQ_k^{T}\left(-(A-d\boldsymbol{\lambda})\right)=\mathrm{uf}\cdot JQ_k^{T}(-A)\cdot\left(\mathrm{uf}^{0}\right)^{-1},\qquad
Q_k(A-d\boldsymbol{\lambda})J^{*}=\mathrm{uf}^{0}\cdot Q_k(A)J^{*}\cdot\mathrm{uf}^{-1},\\
P(A-d\boldsymbol{\lambda})=\mathrm{uf}^{0}\cdot P(A)\cdot\left(\mathrm{uf}^{0}\right)^{-1},
\end{gather*}
and 1.4\,(d4) with $\tilde O=\mathrm{uf}^{0}$ gives
\[
B_k(A-d\boldsymbol{\lambda})=\mathrm{bfun}\left(\mathrm{uf}^{0}\cdot P(A)\cdot\left(\mathrm{uf}^{0}\right)^{-1}\right)=\mathrm{uf}^{0}\cdot B_k(A)\cdot\left(\mathrm{uf}^{0}\right)^{-1}.
\]
By Proposition~3\,(e),
\[
S^{\#}_y(s,A-d\boldsymbol{\lambda})=\mathrm{uf}\cdot S^{\#}_y(s,A)\cdot\mathrm{uf}^{-1},\qquad
S^{\#}_y(X,A-d\boldsymbol{\lambda})=\mathrm{uf}\cdot S^{\#}_y(X,A)\cdot\mathrm{uf}^{-1}.
\]
Consequently
\begin{align*}
H_k(s,A-d\boldsymbol{\lambda})
&=b_k\cdot\left[\mathrm{uf}\cdot S_{\infty}(s,A)\cdot\mathrm{uf}^{-1}\right]\cdot\left[\mathrm{uf}\cdot JQ_k^{T}(-A)\cdot\left(\mathrm{uf}^{0}\right)^{-1}\right]
=\mathrm{uf}\cdot H_k(s,A)\cdot\left(\mathrm{uf}^{0}\right)^{-1},\\
\Gamma_k(s,A-d\boldsymbol{\lambda})
&=\mathrm{uf}^{0}\cdot B_k(A)\cdot\left(\mathrm{uf}^{0}\right)^{-1}
+b_k^2\cdot\left[\mathrm{uf}^{0}\cdot B_k(A)\cdot\left(\mathrm{uf}^{0}\right)^{-1}\right]\cdot\left[\mathrm{uf}^{0}\cdot Q_k(A)J^{*}\cdot\mathrm{uf}^{-1}\right]\\
&\qquad\cdot\left[\mathrm{uf}\cdot S^{\mathrm{sy}}_0(s,A)\cdot\mathrm{uf}^{-1}\right]\cdot\left[\mathrm{uf}\cdot JQ_k^{T}(-A)\cdot\left(\mathrm{uf}^{0}\right)^{-1}\right]\cdot\left[\mathrm{uf}^{0}\cdot B_k(A)\cdot\left(\mathrm{uf}^{0}\right)^{-1}\right]\\
&=\mathrm{uf}^{0}\cdot\Gamma_k(s,A)\cdot\left(\mathrm{uf}^{0}\right)^{-1},\\
M_k(X,A-d\boldsymbol{\lambda})
&=b_k^2\cdot\left[\mathrm{uf}^{0}\cdot Q_k(A)J^{*}\cdot\mathrm{uf}^{-1}\right]\cdot\left[\mathrm{uf}\cdot S^{\mathrm{sy}}_{\infty}(X,A)\cdot\mathrm{uf}^{-1}\right]\cdot\left[\mathrm{uf}\cdot JQ_k^{T}(-A)\cdot\left(\mathrm{uf}^{0}\right)^{-1}\right]\\
&=\mathrm{uf}^{0}\cdot M_k(X,A)\cdot\left(\mathrm{uf}^{0}\right)^{-1},\\
S_{k,y}(X,A-d\boldsymbol{\lambda})
&=\mathrm{uf}\cdot S_{k,y}(X,A)\cdot\mathrm{uf}^{-1}\qquad(y\ge0).
\end{align*}

\subsection*{13. W14}

The property W14 is stated in Section~G1.11.
The relations proved below hold for all $\tilde g\in\mathbb{G}_k$ for the families (p3), (p4), and for all $\tilde g\in\mathbb{G}_{k,3}$ for the family (p5), as Proposition~7 asserts.
They contain W14, since $\mathbb{G}_k^{L}$ is a subgroup of $\mathbb{G}_k$ and $\mathbb{G}^{L}_{k,3}$ is a subgroup of $\mathbb{G}_{k,3}$ (Section~I.4.2).

Let $\tilde g=(S,a)\in\mathbb{G}_k$, put $r:=r(S)$, and represent $a$ by a vector of $M\mathbb{Z}^3$.
By Section~I.4.3, $\tilde g$ acts on $\mathfrak{S}^{\mathrm{ch}}(\eta_k,\ell)$ by $(\tilde gf)(c,g)=S\cdot f\left(r(c-a),rg\right)$ and on $\mathfrak{S}(1,\ell)$ by $(\tilde gv)(z)=S\cdot v\left(r(z-a)\right)$.
If $a$ is replaced by $a+\ell n$ with $n\in\mathbb{Z}^3$, both actions are multiplied by $(-1)^{n_1+n_2+n_3}$; hence each of the relations below holds for all representatives once it holds for one of them.
Both actions are invertible, because $S$ is invertible and $c\mapsto r(c-a)$, $g\mapsto rg$ are bijections.
By the property (G2) of \cite[Section~16.2.7]{QEDI}, $\tilde gA\in\tilde{G}_k$.
Further, $\tilde gs\in\mathrm{Pd}_k$, because $\left|(\tilde gs)_{\mathrm{pc}}\right|=\left|s_{\tilde g^{-1}\mathrm{pc}}\right|\le M^{\alpha_0}$.
Finally $\tilde gX\in D_k$, because the map $c\mapsto rc+a$, where $r$ is diagonal with entries $\pm1$ and $a\in M\mathbb{Z}^3$, maps $M$-cubes onto $M$-cubes and common faces onto common faces.

By Section~G1.2, $\tilde g$ is the action $\mathrm{mS}$ of the morphism $c\mapsto rc+a$ of the block torus $T_k$ onto itself.
Proposition~2.2\,(b) yields
\[
Q_k(\tilde gA)J^{*}\cdot\tilde g=\tilde g\cdot Q_k(A)J^{*},\qquad
JQ_k^{T}(-\tilde gA)\cdot\tilde g=\tilde g\cdot JQ_k^{T}(-A),\qquad
P(\tilde gA)\cdot\tilde g=\tilde g\cdot P(A),
\]
and 1.4\,(d4) with $\tilde O=\tilde g$ on $\mathfrak{S}(1,\ell)$ yields $B_k(\tilde gA)=\mathrm{bfun}\left(\tilde g\cdot P(A)\cdot\tilde g^{-1}\right)=\tilde g\cdot B_k(A)\cdot\tilde g^{-1}$, that is, $B_k(\tilde gA)\cdot\tilde g=\tilde g\cdot B_k(A)$.
The values $y=\infty$ and $y=0$ are neutral (Section~G1.9).
Hence in Proposition~3\,(f) the parameter on the left side is $y$ itself, and
\[
S_{\infty}(\tilde gs,\tilde gA)\cdot\tilde g=\tilde g\cdot S_{\infty}(s,A),\qquad
S^{\mathrm{sy}}_0(\tilde gs,\tilde gA)\cdot\tilde g=\tilde g\cdot S^{\mathrm{sy}}_0(s,A),\qquad
S^{\mathrm{sy}}_{\infty}(\tilde gX,\tilde gA)\cdot\tilde g=\tilde g\cdot S^{\mathrm{sy}}_{\infty}(X,A).
\]
In each of the following products, $\tilde g$ is moved by these relations from the right end to the left end, across one factor after another; for $\Gamma_k$ the factors are crossed in the order $B_k$, $JQ_k^{T}$, $S^{\mathrm{sy}}_0$, $Q_kJ^{*}$, $B_k$:
\begin{align*}
H_k(\tilde gs,\tilde gA)\cdot\tilde g
&=b_k\cdot S_{\infty}(\tilde gs,\tilde gA)\cdot JQ_k^{T}(-\tilde gA)\cdot\tilde g
=b_k\cdot S_{\infty}(\tilde gs,\tilde gA)\cdot\tilde g\cdot JQ_k^{T}(-A)\\
&=\tilde g\cdot b_k\cdot S_{\infty}(s,A)\cdot JQ_k^{T}(-A)=\tilde g\cdot H_k(s,A),\\
\Gamma_k(\tilde gs,\tilde gA)\cdot\tilde g
&=B_k(\tilde gA)\cdot\tilde g+b_k^2\cdot B_k(\tilde gA)\cdot Q_k(\tilde gA)J^{*}\cdot S^{\mathrm{sy}}_0(\tilde gs,\tilde gA)\cdot JQ_k^{T}(-\tilde gA)\cdot B_k(\tilde gA)\cdot\tilde g\\
&=\tilde g\cdot B_k(A)+b_k^2\cdot\tilde g\cdot B_k(A)\cdot Q_k(A)J^{*}\cdot S^{\mathrm{sy}}_0(s,A)\cdot JQ_k^{T}(-A)\cdot B_k(A)
=\tilde g\cdot\Gamma_k(s,A),\\
M_k(\tilde gX,\tilde gA)\cdot\tilde g
&=b_k^2\cdot Q_k(\tilde gA)J^{*}\cdot S^{\mathrm{sy}}_{\infty}(\tilde gX,\tilde gA)\cdot JQ_k^{T}(-\tilde gA)\cdot\tilde g\\
&=\tilde g\cdot b_k^2\cdot Q_k(A)J^{*}\cdot S^{\mathrm{sy}}_{\infty}(X,A)\cdot JQ_k^{T}(-A)=\tilde g\cdot M_k(X,A).
\end{align*}

\paragraph*{The family (p5).}
Let $S\in\mathrm{Sp}_3$ and $y\ge0$.
As $S\in\mathrm{Sp}_3$, the parameter on the left side of Proposition~3\,(f) is again $y$ itself, and
\[
S_{k,y}(\tilde gX,\tilde gA)\cdot\tilde g=S_y(\tilde gX,\tilde gA)\cdot\tilde g=\tilde g\cdot S_y(X,A)=\tilde g\cdot S_{k,y}(X,A).
\]

\subsection*{14. W15 and W16}

\paragraph*{W15.}
The property W15 is stated in Section~G1.11.
By the property (G2) of \cite[Section~16.2.7]{QEDI}, $-A\in\tilde{G}_k$.
Transposes are taken with respect to $\langle\,,\rangle^{\mathrm{ch}}$ on $\mathfrak{S}^{\mathrm{ch}}(\eta_k,\ell)$ and with respect to $\langle\,,\rangle_1$ on $\mathfrak{S}(1,\ell)$ (Section~I.1.6); one has $(O_1O_2)^{T}=O_2^{T}O_1^{T}$ and $\left(O^{T}\right)^{T}=O$.
Proposition~2.2\,(c), applied at $A$ and at $-A$, gives
\begin{gather*}
\left(Q_k(A)J^{*}\right)^{T}=JQ_k^{T}(A)=JQ_k^{T}\left(-(-A)\right),\qquad
\left(JQ_k^{T}(-A)\right)^{T}=\left(\left(Q_k(-A)J^{*}\right)^{T}\right)^{T}=Q_k(-A)J^{*},\\
P(A)^{T}=P(-A),
\end{gather*}
and 1.4\,(d4) gives
\[
B_k(A)^{T}=\mathrm{bfun}\left(P(A)\right)^{T}=\mathrm{bfun}\left(P(A)^{T}\right)=\mathrm{bfun}\left(P(-A)\right)=B_k(-A).
\]
By Proposition~3\,(g) at the neutral values $y=0$ and $y=\infty$,
\[
S^{\mathrm{sy}}_0(s,A)^{T}=\mathsf{C}^{-1}\cdot S^{\mathrm{sy}}_0(s,-A)\cdot\mathsf{C},\qquad
S^{\mathrm{sy}}_{\infty}(X,A)^{T}=\mathsf{C}^{-1}\cdot S^{\mathrm{sy}}_{\infty}(X,-A)\cdot\mathsf{C}.
\]
The operators $Q_k(-A)J^{*}$, $JQ_k^{T}(A)$ and $B_k(-A)$ commute with $\mathsf{C}$ (1.4\,(a) and (d1)).
Hence
\begin{align*}
M_k(X,A)^{T}
&=b_k^2\cdot\left(JQ_k^{T}(-A)\right)^{T}\cdot S^{\mathrm{sy}}_{\infty}(X,A)^{T}\cdot\left(Q_k(A)J^{*}\right)^{T}\\
&=b_k^2\cdot Q_k(-A)J^{*}\cdot\mathsf{C}^{-1}\cdot S^{\mathrm{sy}}_{\infty}(X,-A)\cdot\mathsf{C}\cdot JQ_k^{T}\left(-(-A)\right)
=\mathsf{C}^{-1}\cdot M_k(X,-A)\cdot\mathsf{C},\\
\Gamma_k(s,A)^{T}
&=B_k(A)^{T}+b_k^2\cdot B_k(A)^{T}\cdot\left(JQ_k^{T}(-A)\right)^{T}\cdot S^{\mathrm{sy}}_0(s,A)^{T}\cdot\left(Q_k(A)J^{*}\right)^{T}\cdot B_k(A)^{T}\\
&=B_k(-A)+b_k^2\cdot B_k(-A)\cdot Q_k(-A)J^{*}\cdot\mathsf{C}^{-1}\cdot S^{\mathrm{sy}}_0(s,-A)\cdot\mathsf{C}\cdot JQ_k^{T}\left(-(-A)\right)\cdot B_k(-A)\\
&=\mathsf{C}^{-1}\cdot\Gamma_k(s,-A)\cdot\mathsf{C}.
\end{align*}

\paragraph*{W16.}
The property W16 is stated in Section~G1.11.
By the property (G2) of \cite[Section~16.2.7]{QEDI}, $\overline{A}$ and $-A$ lie in $\tilde{G}_k$; moreover $\overline{s}\in\mathrm{Pd}_k$.
For a linear operator $O$ between spaces of spinors, $\overline{O}$ denotes the operator $F\mapsto\overline{O\,\overline{F}}$ ((I.O5)).
One has $\overline{O_1O_2}=\overline{O_1}\cdot\overline{O_2}$, $\overline{O_1+O_2}=\overline{O_1}+\overline{O_2}$, and $\overline{\mathrm{cf}\cdot O}=\mathrm{cf}\cdot\overline{O}$ for every real number $\mathrm{cf}$.
Proposition~2.2\,(d), applied to the operator-valued functions $A\mapsto Q_k(A)J^{*}$, $A\mapsto JQ_k^{T}(-A)$ and $A\mapsto P(A)$, gives
\[
\overline{Q_k(\overline{A})J^{*}}=\varkappa^{-1}\cdot Q_k(-A)J^{*}\cdot\varkappa,\qquad
\overline{JQ_k^{T}(-\overline{A})}=\varkappa^{-1}\cdot JQ_k^{T}(A)\cdot\varkappa,\qquad
\overline{P(\overline{A})}=\varkappa^{-1}\cdot P(-A)\cdot\varkappa,
\]
and 1.4\,(d4) with $\tilde O=\varkappa^{-1}$ gives
\[
\overline{B_k(\overline{A})}=\overline{\mathrm{bfun}\left(P(\overline{A})\right)}=\mathrm{bfun}\left(\overline{P(\overline{A})}\right)=\mathrm{bfun}\left(\varkappa^{-1}\cdot P(-A)\cdot\varkappa\right)=\varkappa^{-1}\cdot\mathrm{bfun}\left(P(-A)\right)\cdot\varkappa=\varkappa^{-1}\cdot B_k(-A)\cdot\varkappa.
\]
By Proposition~3\,(h),
\[
\overline{S^{\#}_y(\overline{s},\overline{A})}=\varkappa^{-1}\cdot S^{\#}_y(s,-A)\cdot\varkappa,\qquad
\overline{S^{\#}_y(X,\overline{A})}=\varkappa^{-1}\cdot S^{\#}_y(X,-A)\cdot\varkappa.
\]
The number $b_k$ is real.
Hence, the inner factors $\varkappa\cdot\varkappa^{-1}$ cancelling,
\begin{align*}
\overline{H_k(\overline{s},\overline{A})}
&=b_k\cdot\left[\varkappa^{-1}\cdot S_{\infty}(s,-A)\cdot\varkappa\right]\cdot\left[\varkappa^{-1}\cdot JQ_k^{T}(A)\cdot\varkappa\right]\\
&=\varkappa^{-1}\cdot b_k\cdot S_{\infty}(s,-A)\cdot JQ_k^{T}\left(-(-A)\right)\cdot\varkappa=\varkappa^{-1}\cdot H_k(s,-A)\cdot\varkappa,\\
\overline{\Gamma_k(\overline{s},\overline{A})}
&=\varkappa^{-1}\cdot B_k(-A)\cdot\varkappa
+b_k^2\cdot\left[\varkappa^{-1}\cdot B_k(-A)\cdot\varkappa\right]\cdot\left[\varkappa^{-1}\cdot Q_k(-A)J^{*}\cdot\varkappa\right]\cdot\left[\varkappa^{-1}\cdot S^{\mathrm{sy}}_0(s,-A)\cdot\varkappa\right]\\
&\qquad\cdot\left[\varkappa^{-1}\cdot JQ_k^{T}(A)\cdot\varkappa\right]\cdot\left[\varkappa^{-1}\cdot B_k(-A)\cdot\varkappa\right]
=\varkappa^{-1}\cdot\Gamma_k(s,-A)\cdot\varkappa,\\
\overline{M_k(X,\overline{A})}
&=b_k^2\cdot\left[\varkappa^{-1}\cdot Q_k(-A)J^{*}\cdot\varkappa\right]\cdot\left[\varkappa^{-1}\cdot S^{\mathrm{sy}}_{\infty}(X,-A)\cdot\varkappa\right]\cdot\left[\varkappa^{-1}\cdot JQ_k^{T}(A)\cdot\varkappa\right]\\
&=\varkappa^{-1}\cdot M_k(X,-A)\cdot\varkappa.
\end{align*}

\subsection*{15. Continuity in $y$}

By Proposition~3\,(c), the operators $S^{\#}_y(X,A)$ depend continuously on $y\in\mathbb{R}$.
Hence $y\mapsto S_{k,y}(X,A)=S_y(X,A)$ is continuous on $[0,\infty)$, and a continuous function is a Borel function.
This is the requirement ``Borel in $y$'' in the definition of the family (p5) (Section~G1.11).

\subsection*{16. The constants; uniformity}

The bounds (4.1), (4.2), (5.5), (5.6) and (6.1) contain the numbers $32$, $304$, $128$, $2$, $\mathrm{e}$ and the constants $b$, $b_{-}$, $a'$, $M$, $\alpha_{0}$, $\kappa'$, $\gamma^{F}$, which belong to the lines (C0)--(C5) of Section~I.8.2; the constants $C_{\Gamma}$ and $C_{1}$ are formed from them (line (C6)).
The bounds of Propositions~3 and~5 that were applied hold with these constants for all data (SF) with $k\ge1$ and for all $s\in\mathrm{Pd}_k$, $A\in\tilde{G}_k$, $X\in D_k$: Proposition~3\,(a) for all $y\in Y_{k}$; Proposition~5\,(b) at $y=0$, the value at which it enters Section~5; and Proposition~5\,(c), which contains no $y$.
Apart from these bounds, the data entered Sections~4--6 only through $b_{-}\le b_k\le b$ (1.1\,(a)), the cancellation $b_k^2\cdot b_k^{-2}=1$ in (5.3), the inequality $\mathfrak{i}(A)<e_k^{\varepsilon}\le e_0^{\varepsilon}$ (1.1\,(b)), and $w_k(y)=1$.
Therefore $C_1$ and $2C_1$ depend neither on the level $k\ge1$, nor on the side $\ell\in\mathrm{Cl}$, the mass $\mu\in(0,1]$, the number $e_k\in(0,e_0]$, nor on $s$, $A$, $X$, $y$.
(In a flow, the index of the spacing, the physical side of the torus, the charge $e$ and the mass $\bar{m}$ enter the objects of a level only through the quadruple $(k,\ell,\mu,e_k)=(k,\ell_k,\mu_k,e_k)$; see Sections~I.1.8 and~I.6.1.)

By Sections~2--15, the families (p3), (p4), (p5) have the properties W1, W2, W5--W16 of Section~G1.11 for the set of parameters $(\alpha_0,\gamma^{F},\kappa';2C_1)$.
In the Euclidean norms the properties W5--W7 hold with the constant $C_1$, in the forms (4.1), (4.2), (5.5), (5.6), (6.1).
W14 holds for all $\tilde g\in\mathbb{G}_k$, and for the family (p5) for all $\tilde g\in\mathbb{G}_{k,3}$.
Finally, $y\mapsto S_{k,y}(X,A)$ is continuous on $[0,\infty)$.
This completes the proof of Proposition~7.
\qed

\subsection*{17. Remark: the barred family}

This section proves the paragraph ``Barred family'' of Section~G1.11.
Its assertions are consequences of Sections~2--14 for the operator $\bar{H}_k(s,A)=\mathsf{C}^{-1}\cdot H_k(s,-A)\cdot\mathsf{C}$.

The matrix $\mathsf{C}$ is real and monomial with entries $\pm1$, and $\mathsf{C}^{-1}=\mathsf{C}^{T}$ ((I.c3), (I.c6)).
Hence $\left|\mathsf{C}^{\pm1}\mathrm{vv}\right|_2=|\mathrm{vv}|_2$ and $\left|\mathsf{C}^{\pm1}\mathrm{vv}\right|_c=|\mathrm{vv}|_c$ for every $\mathrm{vv}\in\mathbb{C}^4$.
The phase in (I.1.7) is a scalar; so $\delta_{\alpha,-A}$ commutes with constant matrices, and
\[
\delta_{\alpha,-A}\bar{H}_k(s,A)=\mathsf{C}^{-1}\cdot\left[\delta_{\alpha,-A}H_k(s,-A)\right]\cdot\mathsf{C}.
\]
Applied at the point $-A$ of $\tilde{G}_k$, the results of Sections~2--14 give the following.
\begin{itemize}
\item By (4.1) and (4.2),
\[
\left|\bar{H}_k(s,A)v\right|_{\mathrm{su}}\le 32\mathrm{e}\cdot b\cdot a'\cdot M^{27\alpha_0}\cdot|v|_{\mathrm{su}},\qquad
\left|\delta_{\alpha,-A}\bar{H}_k(s,A)v\right|_{\mathrm{su}}\le 304\mathrm{e}\cdot b\cdot a'\cdot M^{27\alpha_0}\cdot|v|_{\mathrm{su}}
\]
(short pairs), and the same holds in the norms $|\cdot|_{\infty}$, $|\cdot|_c$ with the constant $2C_1$.
\item W1 gives $\bar{H}_k(1,A)=\mathsf{C}^{-1}\cdot H_k(-A)\cdot\mathsf{C}=\bar{H}_k(A)$ (Section~I.3.5).
W2 holds for $\bar{H}_k(s,A)$, because $A\mapsto-A$ is linear and maps $\tilde{G}_k$ onto itself.
\item W8--W11 hold for $\bar{H}_k(s,A)$: its kernel is $\mathsf{C}^{-1}\cdot H_k(s,-A)(x,z')\cdot\mathsf{C}$, and $\tilde A=A$ on a set of bonds if and only if $-\tilde A=-A$ there.
\item W13 at $(-A,-\boldsymbol{\lambda})$ gives
\[
\bar{H}_k(s,A-d\boldsymbol{\lambda})=\mathsf{C}^{-1}\cdot H_k\left(s,-A-d(-\boldsymbol{\lambda})\right)\cdot\mathsf{C}=\exp\left(-ie_k\boldsymbol{\lambda}\right)\cdot\bar{H}_k(s,A)\cdot\exp\left(ie_k\boldsymbol{\lambda}^{\mathrm{u}}\right).
\]
\item W14 gives $\bar{H}_k(\tilde gs,\tilde gA)\cdot\mathrm{glb}=\mathrm{glb}\cdot\bar{H}_k(s,A)$ for $\tilde g=(S,a)\in\mathbb{G}_k$, where $(\mathrm{glb}\,F)(t):=(S^{-1})^{T}F(\tilde g^{-1}t)$ (Section~G1.10).
Indeed, $\tilde g(-A)=-\tilde gA$, and $\mathsf{C}^{-1}\cdot\tilde g\cdot\mathsf{C}=\pm\mathrm{glb}$ on both spaces with the same sign, because $(S^{-1})^{T}=\pm\mathsf{C}^{-1}\cdot S\cdot\mathsf{C}$ (Section~I.4.1); therefore
\[
\bar{H}_k(\tilde gs,\tilde gA)\cdot\left(\mathsf{C}^{-1}\cdot\tilde g\cdot\mathsf{C}\right)=\mathsf{C}^{-1}\cdot H_k\left(\tilde gs,\tilde g(-A)\right)\cdot\tilde g\cdot\mathsf{C}=\mathsf{C}^{-1}\cdot\tilde g\cdot H_k(s,-A)\cdot\mathsf{C}=\left(\mathsf{C}^{-1}\cdot\tilde g\cdot\mathsf{C}\right)\cdot\bar{H}_k(s,A).
\]
\item W16 at the point $-A$ gives
\[
\overline{\bar{H}_k(\overline{s},\overline{A})}=\mathsf{C}^{-1}\cdot\left[\overline{H_k\left(\overline{s},\overline{-A}\right)}\right]\cdot\mathsf{C}=\mathsf{C}^{-1}\cdot\varkappa^{-1}\cdot H_k(s,A)\cdot\varkappa\cdot\mathsf{C}=\varkappa^{-1}\cdot\mathsf{C}^{-1}\cdot H_k(s,A)\cdot\mathsf{C}\cdot\varkappa=\varkappa^{-1}\cdot\bar{H}_k(s,-A)\cdot\varkappa;
\]
here $\mathsf{C}$ is real ((I.c3)), and $\varkappa\cdot\mathsf{C}=-\mathsf{C}\cdot\varkappa$ ((I.c4)) gives $\mathsf{C}^{-1}\cdot\varkappa^{-1}=-\varkappa^{-1}\cdot\mathsf{C}^{-1}$, so that the two signs cancel.
\end{itemize}
Apart from the two Euclidean bounds, the derivations above rest only on the properties W1, W2, W5, W8--W11, W13, W14, W16 of the family $H_k(s,A)$ at the point $-A$.
They therefore apply to every family (p3) that has these properties for a set of parameters $(\alpha_0,\gamma,\mathrm{kq};C)$, with $C$ in place of $2C_1$, and for those $\tilde g$ for which the family has W14; this is the paragraph ``Barred family'' of Section~G1.11.
Section~G1.11 states that paragraph for all data (SF), $k=0$ included.
Apart from the two Euclidean bounds, the derivations above do not use $k\ge1$, and at $k=0$ they are read as follows.
\begin{itemize}
\item The domain and the target space of $H_0(s,A)$ and of $\bar{H}_0(s,A)$ are $\mathfrak{S}^{\mathrm{ch}}(1,\ell)$ (Section~G1.11); accordingly $v$ is an element of $\mathfrak{S}^{\mathrm{ch}}(1,\ell)$, the column index is a fermion site $x'$ in place of $z'$, the kernel is the one of (I.1.4), and $\tilde g$ and $\mathrm{glb}$ act on both sides by their formulas for fermion sites.
\item $T^0_0=T_0$ (Section~I.3.9), so that $\boldsymbol{\lambda}^{\mathrm{u}}=\boldsymbol{\lambda}$, and $\exp\left(ie_k\boldsymbol{\lambda}^{\mathrm{u}}\right)$ acts by (I.POS).
\item $H_0(-A)=1$ and $\bar{H}_0(A)=1$ by Section~I.3.9, which takes the place of Section~I.3.5; hence $\mathsf{C}^{-1}\cdot H_0(-A)\cdot\mathsf{C}=\mathsf{C}^{-1}\cdot\mathsf{C}=1=\bar{H}_0(A)$.
\item There is no short pair (Section~I.1.7), so that the part of W5 that concerns $\delta_{\alpha,-A}$ is void.
\end{itemize}

\subsection*{Conditions on the constants}

The arguments of this chapter use the following conditions.

\medskip
\noindent
\begin{tabular}{p{0.25\textwidth}p{0.33\textwidth}p{0.29\textwidth}}
\hline
condition & where it stands & where it is used\\
\hline
$L$ odd & (L5) of Proposition~14.2\,(b); Section~I.8.2, line (C3) & 1.2 ((g0) of \cite[Section~17.4.1]{QEDI}; $M/L$ odd), 1.3, 1.5 (the blocks partition $\mathbb{L}_k$)\\
$M$ a power of $L$, $M\ge L^2$ & (M1) of Proposition~14.2\,(b); Section~I.8.2, line (C5) & 1.2, where it gives, with $L$ odd, that $M/L$ is an odd integer; 1.1\,(d), where it gives $M\ge1$\\
$M>4L$ & (M2) of Proposition~14.2\,(b); Section~I.8.2, line (C5); it is among the conditions of Proposition~3 & 1.2: one of the conditions under which (g0) is stated in \cite[Section~17.4.1]{QEDI}\\
$6L\cdot e_0^{\varepsilon}\le1$ & (Z2.s1) of Proposition~14.2\,(b); Section~I.8.2, line (C8); Section~I.8.4 & 1.1\,(b), hence (1.1), (4.1), (4.2)\\
$\varepsilon>0$ & Section~I.0.5; Section~I.8.2, line (C7), and Proposition~14.2\,(a), line (C7): $\varepsilon\in(1-\alpha,\bar{\varepsilon}]$; (S1) of Proposition~14.2\,(b): $1-\alpha<\varepsilon\le\bar{\varepsilon}$; together with $\alpha<1$ of line (C1) & 1.1\,(b): $e_k^{\varepsilon}\le e_0^{\varepsilon}$ for $e_k\le e_0$\\
$\gamma^{F}L\le\frac{1}{6}<\log2$ & the list of constants at the head of the Propositions~2.m; Proposition~14.2\,(f); it follows from $\gamma^{F}=\frac{1}{2}\gamma_{z}$ and $\gamma_{z}\le a_{z}\le\frac{1}{3L}$, lines (C4r) of Proposition~14.2\,(a) & 1.1\,(c), hence (5.3), (5.4)\\
$\alpha_0>0$ & Proposition~14.2\,(a), line (C4) & 1.1\,(d), hence (5.5)\\
$b>0$, $L>1$ (so $b_{-}>0$; $b_{-}\le b_k\le b$; $b_{L}>0$) & Section~I.3.4; Proposition~14.2\,(a), lines (C0), (C4) & 1.4\,(d), (4.1), (4.2), (5.4), (5.6)\\
\hline
\end{tabular}
\medskip

The side $\ell\in\mathrm{Cl}$ is a datum; it enters through (g0) of \cite[Section~17.4.1]{QEDI}, by way of the fact that $M$ divides $\ell$.

Further conditions enter through the Propositions that are applied.
These are the conditions of Proposition~3, namely those listed at the head of the Propositions~2.m, that is, (L5); (M1); (M2) $M>4L$; $M\ge108\pi\cdot C'_{V}\cdot a'$; (Z2.s1); (Z2.s2) $4C'_{V}\cdot a'\cdot e_0^{\varepsilon}\le1$; together with (M5) in full; and besides them the range $\alpha<1$ (Proposition~14.2\,(a), line (C1): $\alpha\in(1-\bar{\varepsilon},1)$; Section~I.8.2, line (C1)), which is among the conditions of Propositions~2.7 and~2.8.
Proposition~3 depends on the range $\alpha<1$ through Propositions~2.7\,(e) and~2.8\,(a), and Proposition~4\,(d) through Proposition~2.8\,(a); in particular, the bound (4.2) depends on it through the second bound of Proposition~2.7\,(e).
Clauses (b) and (c) of Proposition~5 hold under (Z2.s1), (L5), (M1); Proposition~4\,(d) and Proposition~2.2 require no condition besides those named above.
All the conditions named in this section are among those of Proposition~14.2; hence the results of this chapter hold whenever the constants satisfy the conditions of Proposition~14.2.

\section*{Chapter 8. The localised operator families at level 0 (W17, fermionic part): proof of Proposition~8}

\subsection*{0. Conventions of this chapter; the setting at level 0}

\paragraph{0.1. Data.}
Throughout the chapter the data are (SF) of Section~G1.1 with $k=0$:
a side $\ell\in\mathrm{Cl}$, a mass $\mu\in(0,1]$, and a number $e_k\in(0,e_0]$, which is the fourth entry of the quadruple;
here $e_0$ denotes the ceiling (Section~I.2.5).
One has $\eta_0=1$, $T_0=T^0_0=\mathbb{Z}^3/\ell\mathbb{Z}^3$ and $T^{\mathrm{f}}_0=T_0\times O_{\mathrm{cub}}$,
so that above every photon site $c$ there are the $48$ fermion sites $x=(c,g)$, $g\in O_{\mathrm{cub}}$.
The covering map $\mathbb{Z}^3\to T_0$ is denoted by $\mathrm{pr}$.
The condition $\ell\in\mathrm{Cl}$ means that $\ell/(L\cdot M)$ is an odd integer $\ge13$;
consequently $M$ divides $\ell$, and $\ell\ge13L\cdot M>2M$.
The set $\tilde{G}_0$ of Section~I.5.3 is an open subset of the space of the connections on $T_0$,
and it is invariant under the maps $A\mapsto A-d\boldsymbol{\lambda}$ with $\boldsymbol{\lambda}$ real, $A\mapsto-A$, $A\mapsto\overline{A}$, and $A\mapsto\tilde{g}A$ with $\tilde{g}\in\mathbb{G}_0$;
this is the display (G2) in \cite[Section~16.2.7]{QEDI}.

\paragraph{0.2. The space.}
Recall that
\[
\mathfrak{S}^{\mathrm{ch}}(1,\ell)=\left\{f\colon\mathbb{Z}^3\times O_{\mathrm{cub}}\to\mathbb{C}^4:\ f(c+\ell\cdot e_\mu,g)=-f(c,g)\right\};
\]
its dimension is $4\cdot48\cdot\ell^3=192\ell^3$.
By Section~I.3.9 and (I.L0.3), at level $0$ this one space is at the same time the space of the fine spinors and the space of the spinors on the unit lattice.
Denote by $\mathrm{Lin}$ the complex vector space of the linear maps of $\mathfrak{S}^{\mathrm{ch}}(1,\ell)$ into itself;
it has finite dimension.
By (I.L0.16), the fermionic part of a system of level $0$ consists of families with values in $\mathrm{Lin}$,
namely $\Gamma_0(s,A)$, $S_{0,y}(X,A)$ and the diagonal family $H_0(s,A)$.

\paragraph{0.3. Measure and kernels.}
For $\eta_0=1$ the definitions (I.1.3) and (I.1.4) take the following form.
The measure is $\int\mathrm{d}x=\frac{1}{48}\sum_c\sum_{g\in O_{\mathrm{cub}}}$, where $c$ ranges over a fundamental domain of $\ell\mathbb{Z}^3$,
and the kernel $O(x,x')$ of an element $O$ of $\mathrm{Lin}$ is the $4\times4$ matrix, defined for $x,x'\in\mathbb{Z}^3\times O_{\mathrm{cub}}$, such that
\begin{equation}
\begin{split}
(Of)(x)&=\frac{1}{48}\sum_{c'}\sum_{g'\in O_{\mathrm{cub}}}O\left(x,(c',g')\right)\cdot f(c',g'),\\
O(x+\ell\cdot e_\mu,x')&=O(x,x'+\ell\cdot e_\mu)=-O(x,x'),
\end{split}
\tag{8.0.1}
\end{equation}
with $c'$ ranging over a fundamental domain of $\ell\mathbb{Z}^3$.
This is the kernel of Section~G1.2 for $(R,\boldsymbol{\sigma})=(\ell,-1)$.
The sum does not depend on the choice of the fundamental domain, since the change of sign of $O(x,\cdot\,)$ and that of $f$ cancel.
Whether $O(x,x')$ vanishes, and its modulus, depend only on the classes of $x$ and $x'$ modulo $\ell\mathbb{Z}^3$;
accordingly, an assertion that $O(x,x')=0$ unless $x$ and $x'$ satisfy a certain relation is understood for these classes (Sections~G1.2, G1.11).
By (I.ke1), the multiplication by a complex function $\psi$ of $c\in T_0$ has the kernel
\begin{equation}
48\cdot\mathbf{1}[x=x']\cdot\psi(c)\cdot\mathbf{1}_4\qquad(x,x'\text{ in one fundamental domain}).
\tag{8.0.2}
\end{equation}
In view of the sign rule in (8.0.1), this means for arbitrary $x,x'\in\mathbb{Z}^3\times O_{\mathrm{cub}}$:
the kernel equals $(-1)^{n_1+n_2+n_3}\cdot48\cdot\psi(c)\cdot\mathbf{1}_4$ at $x'=x+\ell\cdot n$, $n\in\mathbb{Z}^3$,
and it vanishes at every $x'$ outside the orbit $x+\ell\mathbb{Z}^3$.

\paragraph{0.4. Positions.}
By (I.POS), the site $x=(c,g)$ lies over a set $Z$ of photon sites if and only if $c\in Z$.
Moreover $\operatorname{supp}x=\{c\}$, so that the phrase ``the index $x$ is contained in $Z$'' (Section~G1.11, W8--W10) has the same meaning.
Further, $d(x,x')=d(c,c')$ and
\[
K_x=K_c=\left\{c'\in T_0:\ d(c',c)\le2L\right\}.
\]
The unit block of a site $z$ of $T_0$ is $\Delta_z=\{z\}$, and $\mathbf{1}_{\Delta_z}$ is the indicator of the orbit of the cluster $\{z\}\times O_{\mathrm{cub}}$.

\paragraph{0.5. Exact operators (Section~I.3.9).}
For $y\in Y_0=\mathbb{R}$ and $A\in\tilde{G}_0$,
\begin{gather*}
F_y(A)=D_A+\mu+b_L\cdot J\cdot P(A)\cdot J^{*}+i\gamma_5\cdot y=F_0(A)+i\gamma_5\cdot y,\\
w_0(y)=(1+|y|)^{-1},\qquad H_0(A)=1.
\end{gather*}
For all $y\in\mathbb{R}$ and $A\in\tilde{G}_0$ the operator $F_y(A)$ is a bijection of $\mathfrak{S}^{\mathrm{ch}}(1,\ell)$ onto itself, by Proposition~2.8\,(a).
Hence the inverses that occur in the following definitions of Section~I.3.9 exist:
\[
S_y(A)=F_y(A)^{-1},\qquad S_{0,y}(A)=S_y(A),\qquad\Gamma_0(A)=S^0_1(A)=F_0(A)^{-1}=S_{0,0}(A);
\]
the last operator is $S_y(A)$ at $y=0$.
For $y\ne0$ the operator $F_y(A)$ contains $\gamma_5$;
thus $y=0$ is the only neutral value of $y$ (Section~G1.9).

\paragraph{0.6. The families of Section~G1.9.}
Let $\mathrm{wk}=(\square_0,\dots,\square_n)$ be a walk of $M$-cubes of $T_0$ (Section~G1.8).
The set $X_{\mathrm{wk}}$, the union of the sets $\square_i^{\sim1}$, is a polymer (Proposition~3\,(d));
$S_{y,\mathrm{wk}}(A)$ and $S'_{y,\mathrm{wk}}(A)$ are the two products of Section~G1.8;
and
\[
s_{\mathrm{wk}}=\prod\left\{s_{\mathrm{pc}}:\ \mathrm{pc}\in\pi_0,\ \mathrm{pc}\subset X_{\mathrm{wk}}\right\}.
\]
For all $y\in\mathbb{R}$ and $A\in\tilde{G}_0$ the factors of the two products are elements of $\mathrm{Lin}$.
Indeed, $h_\square$ is a multiplication (Section~G1.7);
$\mathrm{Cm}_{\square,y}(A)=-\left[F_y(A),h_\square\right]$;
and the local inverse $\mathrm{Sloc}_{\square,y}(A)$ exists and is an element of $\mathrm{Lin}$, by Sections~G1.5, G1.6 and Propositions~2.5\,(a), 2.6\,(a).
Here Proposition~2.5\,(a) is applied with $\delta=2e_k^{\varepsilon}$, which is admissible there, since
\begin{align*}
3L\cdot\delta&=6L\cdot e_k^{\varepsilon}\le6L\cdot e_0^{\varepsilon}\le1&&\text{by (Z2.s1)},\\
2C'_V\cdot a'\cdot\delta&=4C'_V\cdot a'\cdot e_k^{\varepsilon}\le4C'_V\cdot a'\cdot e_0^{\varepsilon}\le1&&\text{by (Z2.s2)}.
\end{align*}
By the definitions in Section~G1.9 and in Proposition~8,
\begin{equation}
\begin{split}
\Gamma_0(s,A)&=\frac{1}{2}\sum_{\mathrm{wk}\in\mathrm{Wk}}s_{\mathrm{wk}}\cdot\left(S_{y,\mathrm{wk}}(A)+S'_{y,\mathrm{wk}}(A)\right)\quad\text{at }y=0,\\
S_{0,y}(X,A)&=\sum_{\mathrm{wk}:\,X_{\mathrm{wk}}=X}S_{y,\mathrm{wk}}(A)\quad(y\ge0).
\end{split}
\tag{8.0.3}
\end{equation}
By Proposition~3\,(a), the series $S_y(s,A)$, $S'_y(s,A)$ and $S_y(X,A)$ converge absolutely in the norm $N_{\gamma^F}$ of Section~G1.3,
for $s\in\mathrm{Pd}_0$, $A\in\tilde{G}_0$, $X\in D_0$ and $y\in\mathbb{R}$.
As $N_{\gamma^F}$ is a norm on the finite-dimensional space $\mathrm{Lin}$ (Lemma~8.B below and its proof),
the series (8.0.3) define elements of $\mathrm{Lin}$.

\paragraph{0.7. $M$-cubes and parameter cubes.}
The notions are those of Section~I.1.1 and of \cite[Section~5.2.3 and Definition~5.8]{QEDI}.
By (M1) and (L5), $M$ is a power of the odd integer $L$.
Hence $M$ is odd, $(M-1)/2$ is an integer, and an integer has modulus $<M/2$ if and only if its modulus is at most $(M-1)/2$.
For $z\in M\mathbb{Z}^3$ the $M$-cube of $\mathbb{Z}^3$ with centre $z$ is
\[
\tilde{\square}(z):=\left\{c'\in\mathbb{Z}^3:\ |c'_\mu-z_\mu|\le\frac{M-1}{2},\ \mu=1,2,3\right\}.
\]
The sets of $M$ consecutive integers whose middle elements are the points of $M\mathbb{Z}$ form a partition of $\mathbb{Z}$;
therefore the cubes $\tilde{\square}(z)$, $z\in M\mathbb{Z}^3$, form a partition of $\mathbb{Z}^3$.
Since $M$ divides $\ell$ (subsection~0.1), for every $n\in\mathbb{Z}^3$ the set $\tilde{\square}(z)+\ell\cdot n=\tilde{\square}(z+\ell\cdot n)$ is again one of these cubes.
The $M$-cubes of $T_0$ are the sets $\mathrm{pr}\,\tilde{\square}(z)$, $z\in M\mathbb{Z}^3$.
\begin{itemize}
\item[(l1)] If $\mathrm{pr}\,\tilde{\square}(z)$ and $\mathrm{pr}\,\tilde{\square}(z')$ have a site in common, then $z'-z\in\ell\mathbb{Z}^3$ and $\mathrm{pr}\,\tilde{\square}(z')=\mathrm{pr}\,\tilde{\square}(z)$.
Indeed, there are $c'\in\tilde{\square}(z)$ and $n\in\mathbb{Z}^3$ with $c'+\ell\cdot n\in\tilde{\square}(z')$.
The point $c'+\ell\cdot n$ belongs also to $\tilde{\square}(z+\ell\cdot n)$.
Since the cubes of $\mathbb{Z}^3$ form a partition of $\mathbb{Z}^3$, it follows that $z'=z+\ell\cdot n$, and
$\mathrm{pr}\,\tilde{\square}(z')=\mathrm{pr}\left(\tilde{\square}(z)+\ell\cdot n\right)=\mathrm{pr}\,\tilde{\square}(z)$.
\end{itemize}
Every site of $T_0$ is of the form $\mathrm{pr}\,c'$ with $c'$ in some cube $\tilde{\square}(z)$.
Consequently, the $M$-cubes form a partition of the set of the sites of $T_0$ (``the $M$-cube of $c$'' is the one that contains the site $c$);
$\mathrm{pr}$ maps the cube of $\mathbb{Z}^3$ that contains a point $c'$ onto the $M$-cube of $\mathrm{pr}\,c'$;
and $\mathrm{pr}\,\tilde{\square}(z)=\mathrm{pr}\,\tilde{\square}(z')$ if and only if $z-z'\in\ell\mathbb{Z}^3$.

Let $\pi_0$ be the set of the $M$-cubes of $T_0$, and
\[
\mathrm{Pd}_0=\left\{s=(s_{\mathrm{pc}})_{\mathrm{pc}\in\pi_0}:\ |s_{\mathrm{pc}}|\le M^{\alpha_0}\right\}.
\]
Two cubes of $\pi_0$ are linked if they have a common face.
A subset of $\pi_0$ is link-connected if any two of its cubes are joined by a chain of cubes of the subset in which consecutive cubes are linked.
The link-components of a subset $\mathscr{Y}$ of $\pi_0$ are its maximal link-connected subsets;
a link-component of the union $Y$ of $\mathscr{Y}$ is the union of the cubes of a link-component of $\mathscr{Y}$.
By (l1), the cube of $\mathscr{Y}$ that contains a given site of $Y$ is the $M$-cube of this site.
Two further consequences follow.
\begin{itemize}
\item[(l2)] The link-components of $\mathscr{Y}$ are the classes of the equivalence relation ``joined by a chain of linked cubes of $\mathscr{Y}$''; hence they form a partition of $\mathscr{Y}$.
Together with (l1): $Y$ is the disjoint union of its link-components.
\item[(l3)] The union of two link-connected subsets of $\mathscr{Y}$ that have a cube in common is link-connected.
Hence, by maximality, a link-connected subset of $\mathscr{Y}$ that contains a cube of a link-component of $\mathscr{Y}$ is contained in this link-component.
\end{itemize}
In the definition of $H_0(s,A)$ in Proposition~8, and in W11, ``$\mathrm{pc}$ meets $K_c$'' means that the $M$-cube $\mathrm{pc}$ and the set $K_c$ have a photon site of $T_0$ in common;
this is the meaning of ``meet'' in \cite[Section~5.2.3]{QEDI}, paragraph ``Indices'', to which Section~I.1.1 refers.

\paragraph{0.8. Notation.}
In this chapter $\mathrm{Lin}$ is the space of subsection~0.2 and $\mathrm{pr}$ the covering map of subsection~0.1;
$\mathrm{pr}_1$, $\mathrm{pr}_2$ denote the same map for the sides $\ell_1$, $\ell_2$ (Section~4 below).
The cubes $\tilde{\square}(z)$ are defined in subsection~0.7.
The objects $\mathrm{Cub}(c)$, $\tilde{c}$, $\tilde{o}$, $\tilde{K}$, $\mathrm{jv}$, $\mathrm{jv}'$, $\mathrm{dv}$, $\mathrm{Nb}_\mu$ are those of Lemma~8.C and of its proof;
$\tilde{c}'$ denotes a lift of $c'$;
$\mathrm{Cub}_1$, $\mathrm{Cub}_2$ are the sets $\mathrm{Cub}$ on the tori of the sides $\ell_1$, $\ell_2$ (Section~4).
The function $\mathrm{fa}_s(c)$ is defined in subsection~3.1, and $\mathscr{Y}_1$ is the link-component of $\mathscr{Y}$ whose union is $Y_1$ (Section~3).
The symbols $z_\square$ (Section~G1.10), $\mathscr{Y}$, $Y$, $Y_1$, $\tilde{s}$, $\tilde{A}$ (Section~G1.11, W8), $\delta$ (Proposition~2.4),
the functions of position $\psi$, $\psi'$ ((I.POS), Section~G1.2), and $\mathrm{pj}$, $\mathrm{pj}^{*}$, $\ell_1$, $\ell_2$ (Section~I.6.3) have the meanings given at these places.
The number $(8.n.m)$ is carried by the $m$-th numbered display of Section~$n$,
and a number $n.m$ without further indication denotes subsection~$n.m$ of this chapter.

\subsection*{1. Three lemmas}

\begin{named}{Lemma 8.A}
For every $4\times4$ matrix $m$ one has $|m|_c\le2|m|_2$.
\end{named}

\begin{proof}
Let $v\in\mathbb{C}^4$.
By Section~I.1.5,
\[
|mv|_c\le|mv|_2\le|m|_2\cdot|v|_2\le2|m|_2\cdot|v|_c.
\]
\end{proof}

\begin{named}{Lemma 8.B}
Let $O\in\mathrm{Lin}$ and $\gamma\ge0$, and let $x$ lie over $z$ and $x'$ over $z'$.
Then
\[
|O(x,x')|_2\le48\left\|\mathbf{1}_{\Delta_z}\cdot O\cdot\mathbf{1}_{\Delta_{z'}}\right\|_{\mathrm{su}}\le48N_\gamma(O).
\]
Consequently, if $O_n\in\mathrm{Lin}$ and $\sum_nN_\gamma(O_n)<\infty$, then $\sum_nO_n$ converges in $\mathrm{Lin}$, and
\[
\left(\sum_nO_n\right)(x,x')=\sum_nO_n(x,x')\qquad\text{for all }x,x',
\]
the series on the right converging absolutely.
\end{named}

\begin{proof}
Let $v\in\mathbb{C}^4$, and let $f$ be the function on $\mathbb{Z}^3\times O_{\mathrm{cub}}$ given by $f(x'+\ell\cdot n):=(-1)^{n_1+n_2+n_3}\cdot v$ for $n\in\mathbb{Z}^3$ and by $f:=0$ at all other points.
Then $f\in\mathfrak{S}^{\mathrm{ch}}(1,\ell)$, $f=\mathbf{1}_{\Delta_{z'}}\cdot f$, and $|f|_{\mathrm{su}}=|v|_2$.
Of the $48\ell^3$ terms of the sum in (8.0.1), all vanish except the one at the point $x'+\ell\cdot n'$ of the orbit of $x'$ that lies over the fundamental domain.
By the sign rule in (8.0.1) and by the definition of $f$, this term equals
\[
\frac{1}{48}(-1)^{n'_1+n'_2+n'_3}\cdot O(x,x')\cdot(-1)^{n'_1+n'_2+n'_3}\cdot v=\frac{1}{48}O(x,x')\cdot v.
\]
Hence $(Of)(x)=\frac{1}{48}O(x,x')\cdot v$, and
\[
\left|O(x,x')\cdot v\right|_2=48\left|\left(\mathbf{1}_{\Delta_z}\cdot O\cdot\mathbf{1}_{\Delta_{z'}}\cdot f\right)(x)\right|_2\le48\left\|\mathbf{1}_{\Delta_z}\cdot O\cdot\mathbf{1}_{\Delta_{z'}}\right\|_{\mathrm{su}}\cdot|v|_2.
\]
As $v$ is arbitrary, the first inequality follows.
The second one holds because every term of the sum that defines $N_\gamma(O)$ is non-negative and $\exp\left(\gamma\cdot d(z,z')\right)\ge1$.
The map $N_\gamma$ is a norm on $\mathrm{Lin}$:
it is a supremum of finite sums of seminorms, and $N_\gamma(O)=0$ implies $O(x,x')=0$ for all $x,x'$ by the first part, hence $O=0$ by (8.0.1).
Since $\mathrm{Lin}$ has finite dimension, it is complete, and therefore the series converges in $\mathrm{Lin}$.
The map $O\mapsto O(x,x')$ is linear and bounded by $48N_\gamma(O)$;
hence it may be applied term by term.
\end{proof}

\begin{named}{Lemma 8.C}[The cubes that meet $K_c$]
For a site $c$ of $T_0$ put
\[
\mathrm{Cub}(c):=\left\{\mathrm{pc}\in\pi_0:\ \mathrm{pc}\text{ and }K_c\text{ have a site in common}\right\}.
\]
Then the following hold.
\begin{itemize}
\item[(i)] The $M$-cube of $c$ belongs to $\mathrm{Cub}(c)$.
\item[(ii)] $\mathrm{Cub}(c)$ has at most $8$ elements.
\item[(iii)] $\mathrm{Cub}(c)$ is link-connected.
\item[(iv)] $\mathrm{Cub}(\tilde{g}c)=\left\{\tilde{g}\,\mathrm{pc}:\ \mathrm{pc}\in\mathrm{Cub}(c)\right\}$ for every $\tilde{g}\in\mathbb{G}_0$.
\item[(v)] Let $\tilde{c}\in\mathbb{Z}^3$ be a lift of $c$, let $z_\square\in M\mathbb{Z}^3$ be the centre of the cube of $\mathbb{Z}^3$ that contains $\tilde{c}$, put $\tilde{o}:=\tilde{c}-z_\square$, and for $\mu=1,2,3$ put
\[
\mathrm{Nb}_\mu:=\begin{cases}
\{0,1\} & \text{if }\tilde{o}_\mu\ge\frac{M+1}{2}-2L,\\
\{-1,0\} & \text{if }\tilde{o}_\mu\le-\left(\frac{M+1}{2}-2L\right),\\
\{0\} & \text{otherwise}.
\end{cases}
\]
Then $\frac{M+1}{2}-2L\ge1$, so that the first two cases exclude each other,
and $\mathrm{jv}\mapsto\mathrm{pr}\,\tilde{\square}(z_\square+M\cdot\mathrm{jv})$ is a bijection of $\mathrm{Nb}_1\times\mathrm{Nb}_2\times\mathrm{Nb}_3$ onto $\mathrm{Cub}(c)$.
\end{itemize}
\end{named}

\begin{proof}
Let $\tilde{c}$, $z_\square$ and $\tilde{o}$ be as in (v);
thus $|\tilde{o}_\mu|\le\frac{M-1}{2}$.
Put $\tilde{K}:=\left\{c'\in\mathbb{Z}^3:\ |c'-\tilde{c}|\le2L\right\}$.
Since $d(c',c)$ is the least sup norm of a difference of lifts, one has $\mathrm{pr}\,\tilde{K}=K_c$.

\emph{Step 1: $\mathrm{Cub}(c)=\left\{\mathrm{pr}\,\tilde{\square}(z):\ z\in M\mathbb{Z}^3,\ \tilde{\square}(z)\text{ meets }\tilde{K}\right\}$.}
Let $\mathrm{pc}\in\mathrm{Cub}(c)$, let $c'$ be a common site of $\mathrm{pc}$ and $K_c$, and let $\tilde{c}'\in\tilde{K}$ be a lift of $c'$.
The cube $\tilde{\square}(z)$ that contains $\tilde{c}'$ meets $\tilde{K}$, and by (l1) the set $\mathrm{pr}\,\tilde{\square}(z)$ is the $M$-cube of $c'$, that is, $\mathrm{pc}$.
Conversely, $\mathrm{pr}$ maps a common point of $\tilde{\square}(z)$ and $\tilde{K}$ to a common site of $\mathrm{pr}\,\tilde{\square}(z)$ and $K_c$.

\emph{Step 2.}
Both $\tilde{\square}(z)$ and $\tilde{K}$ are products of three intervals of integers:
in the direction $\mu$ the first has the midpoint $z_\mu$ and the half-width $\frac{M-1}{2}$, the second the midpoint $\tilde{c}_\mu$ and the half-width $2L$, and all four numbers are integers.
The intersection of the two products is the product of the three intersections;
hence $\tilde{\square}(z)$ meets $\tilde{K}$ if and only if in each of the three directions the two intervals have an integer in common.
Two such intervals have an integer in common if and only if the distance of their midpoints is at most the sum of their half-widths.
For $z=z_\square+M\cdot\mathrm{jv}$ with $\mathrm{jv}\in\mathbb{Z}^3$ the distance of the midpoints in the direction $\mu$ is $|z_\mu-\tilde{c}_\mu|=|M\cdot\mathrm{jv}_\mu-\tilde{o}_\mu|$.
Therefore
\begin{equation}
\tilde{\square}(z_\square+M\cdot\mathrm{jv})\text{ meets }\tilde{K}\quad\text{if and only if}\quad|M\cdot\mathrm{jv}_\mu-\tilde{o}_\mu|\le\frac{M-1}{2}+2L\quad\text{for }\mu=1,2,3.
\tag{8.1.1}
\end{equation}
By (M2) one has $4L<M$;
as $4L$ is even and $M$ is odd, $4L\le M-1$, that is, $2L\le\frac{M-1}{2}$, and the right-hand side of (8.1.1) is at most $M-1$.
Fix $\mu$.
\begin{itemize}
\item If $\mathrm{jv}_\mu=0$, then $|\tilde{o}_\mu|\le\frac{M-1}{2}$, and the condition holds.
\item If $|\mathrm{jv}_\mu|\ge2$, then $|M\cdot\mathrm{jv}_\mu-\tilde{o}_\mu|\ge2M-\frac{M-1}{2}>M-1$, and the condition fails.
\item If $\mathrm{jv}_\mu=1$, the condition reads $M-\tilde{o}_\mu\le\frac{M-1}{2}+2L$, that is, $\tilde{o}_\mu\ge\frac{M+1}{2}-2L$;
moreover $\frac{M+1}{2}-2L\ge\frac{M+1}{2}-\frac{M-1}{2}=1$.
\item If $\mathrm{jv}_\mu=-1$, the condition reads $M+\tilde{o}_\mu\le\frac{M-1}{2}+2L$, that is, $\tilde{o}_\mu\le-\left(\frac{M+1}{2}-2L\right)\le-1$.
\end{itemize}
The last two cases exclude each other.
Hence the set of the integers $\mathrm{jv}_\mu$ that satisfy the condition is the set $\mathrm{Nb}_\mu$ of (v);
it is $\{0\}$, $\{0,1\}$ or $\{-1,0\}$.
By Step~1 and (8.1.1),
\begin{equation}
\mathrm{Cub}(c)=\left\{\mathrm{pr}\,\tilde{\square}(z_\square+M\cdot\mathrm{jv}):\ \mathrm{jv}\in\mathrm{Nb}_1\times\mathrm{Nb}_2\times\mathrm{Nb}_3\right\}.
\tag{8.1.2}
\end{equation}
For $z=z_\square+M\cdot\mathrm{jv}$ with $\mathrm{jv}\in\mathbb{Z}^3$, the closed real cubes $\left\{y\in\mathbb{R}^3:\ |y-z|\le M/2\right\}$ and $\left\{y\in\mathbb{R}^3:\ |y-\tilde{c}|\le2L\right\}$ meet if and only if $|M\cdot\mathrm{jv}_\mu-\tilde{o}_\mu|\le M/2+2L$ for $\mu=1,2,3$,
that is, if and only if the inequalities of (8.1.1) hold with $M/2+2L$ in the place of $\frac{M-1}{2}+2L$.
Since $|M\cdot\mathrm{jv}_\mu-\tilde{o}_\mu|$ is an integer and $M$ is odd, this is the case if and only if $\tilde{\square}(z)$ meets $\tilde{K}$;
thus the closed real cubes define the same set of cubes.

\emph{(v).}
The map is onto by (8.1.2).
Let $\mathrm{jv}$, $\mathrm{jv}'$ be elements of the product with $\mathrm{pr}\,\tilde{\square}(z_\square+M\cdot\mathrm{jv})=\mathrm{pr}\,\tilde{\square}(z_\square+M\cdot\mathrm{jv}')$, and put $\mathrm{dv}:=\mathrm{jv}-\mathrm{jv}'$.
By (l1), $M\cdot\mathrm{dv}\in\ell\mathbb{Z}^3$.
Each $\mathrm{Nb}_\mu$ consists of one integer or of two consecutive integers;
hence $\mathrm{dv}_\mu\in\{-1,0,1\}$ and $|M\cdot\mathrm{dv}_\mu|\le M<\ell$ (subsection~0.1).
A multiple of $\ell$ of modulus $<\ell$ is $0$.
Hence $\mathrm{dv}=0$.

\emph{(i).}
The point $\mathrm{jv}=0$ belongs to the product, and $\mathrm{pr}\,\tilde{\square}(z_\square)$ is the $M$-cube of $c$ by (l1).

\emph{(ii).}
The product has at most $2^3=8$ elements.

\emph{(iii).}
Since each $\mathrm{Nb}_\mu$ consists of one integer or of two consecutive integers, any two elements $\mathrm{jv}$, $\mathrm{jv}'$ of the product are joined inside the product by at most three steps, each of which changes one coordinate by $+1$ or by $-1$.
It remains to show that for $z\in M\mathbb{Z}^3$ the $M$-cubes $\mathrm{pr}\,\tilde{\square}(z)$ and $\mathrm{pr}\,\tilde{\square}(z\pm M\cdot e_\mu)$ of $T_0$ are linked.
The closed cubes $\left\{y\in\mathbb{R}^3:\ |y-z|\le M/2\right\}$ and $\left\{y\in\mathbb{R}^3:\ |y-z\mp M\cdot e_\mu|\le M/2\right\}$, whose lattice points form $\tilde{\square}(z)$ and $\tilde{\square}(z\pm M\cdot e_\mu)$, have the common face
\[
\left\{y:\ y_\mu=z_\mu\pm M/2,\ |y_\nu-z_\nu|\le M/2\text{ for }\nu\ne\mu\right\}.
\]
The union of $\tilde{\square}(z)$ and $\tilde{\square}(z\pm M\cdot e_\mu)$ has the sup-diameter $2M-1$:
in the direction $\mu$ its coordinates fill an interval of $2M$ consecutive integers, in the two other directions an interval of $M$ consecutive integers;
and $2M-1+\eta_0=2M<\ell$ (subsection~0.1).
Hence this union is a lift of its image in the sense of the paragraph on lifts in Section~I.1.1, and $\mathrm{pr}$ is injective on it.
By that paragraph, a lift of a union of $M$-cubes is a union of $M$-cubes of $\mathbb{L}_0=\mathbb{Z}^3$ with the same common faces;
consequently the two images are two distinct $M$-cubes of $T_0$ with the common face of their lifts, that is, they are linked.
In the terms of Section~I.1.1: the centres $\mathrm{pr}\,z$ and $\mathrm{pr}(z\pm M\cdot e_\mu)$ of the two cubes differ by $\pm M\cdot e_\mu$ modulo $\ell\mathbb{Z}^3$.

\emph{(iv).}
Let $\tilde{g}=(S,a)$, $r:=r(S)\in\{\pm1\}^3$, $a\in M\mathbb{Z}^3$.
The matrix $r$ is diagonal with entries $\pm1$.
Hence the map $c'\mapsto rc'+a$ of $\mathbb{Z}^3$ is a bijection, preserves the sup norm of differences, maps the orbits of $\ell\mathbb{Z}^3$ onto such orbits (as $r\ell\mathbb{Z}^3=\ell\mathbb{Z}^3$), and maps $\tilde{\square}(z)$ onto $\tilde{\square}(rz+a)$, where $rz+a\in M\mathbb{Z}^3$.
Therefore the induced map $c'\mapsto\tilde{g}c'$ is a bijection of $T_0$ that preserves $d$ and permutes the $M$-cubes of $T_0$.
It follows that $\tilde{g}K_c=K_{\tilde{g}c}$, and that $\mathrm{pc}$ has a site in common with $K_c$ if and only if $\tilde{g}\,\mathrm{pc}$ has a site in common with $K_{\tilde{g}c}$.
\end{proof}

\subsection*{2. Proof of (a)}

\begin{proof}[Proof of Proposition~8\,(a)]
Of the properties named in (a), those that concern the two families are the following (Section~G1.11):
for $\Gamma_0(s,A)$, the properties W1, W2, W6, W8, W9, W13, W14, W15, W16;
for $S_{0,y}(X,A)$, the properties W1, W2, W7, W12, W13, W14.
Throughout, $s\in\mathrm{Pd}_0$, $A\in\tilde{G}_0$, $X\in D_0$ and $y\ge0$;
since $[0,\infty)\subset Y_0=\mathbb{R}$, Proposition~3 applies to every such $y$.

\textbf{W1.}
Proposition~3\,(b) with $S^{\#}=S^{\mathrm{sy}}$ and $y=0$, together with subsection~0.5, gives
\[
\Gamma_0(1,A)=S^{\mathrm{sy}}_0(1,A)=S_{0,0}(A)=F_0(A)^{-1}=\Gamma_0(A).
\]
Proposition~3\,(b) with $S^{\#}=S$, together with subsection~0.5, gives
\[
\sum_{X\in D_0}S_{0,y}(X,A)=\sum_{X\in D_0}S_y(X,A)=S_y(A)=F_y(A)^{-1}=\left(F_0(A)+i\gamma_5\cdot y\right)^{-1}=S_{0,y}(A).
\]
The inverses exist by Proposition~2.8\,(a), as noted in subsection~0.5.

\textbf{W2; continuity in $y$.}
By Proposition~3\,(c) with $S^{\#}=S^{\mathrm{sy}}$ and $y=0$, the family $\Gamma_0(s,A)=S^{\mathrm{sy}}_0(s,A)$ is continuous on $\mathrm{Pd}_0\times\tilde{G}_0$, analytic in $(s,A)$ in the interior, and analytic in $A\in\tilde{G}_0$ for every fixed $s\in\mathrm{Pd}_0$, the boundary included.
By Proposition~3\,(c) with $S^{\#}=S$, the map $S_{0,y}(X,\cdot\,)=S_y(X,\cdot\,)$ is analytic on $\tilde{G}_0$ for every $y\ge0$,
and $y\mapsto S_y(X,A)$ is continuous on $\mathbb{R}$, hence on $[0,\infty)$.
A continuous map of $[0,\infty)$ into the finite-dimensional space $\mathrm{Lin}$ is Borel;
equivalently, by Lemma~8.B, every entry of every kernel $S_{0,y}(X,A)(x,x')$ is a continuous, hence Borel, function of $y$.

\textbf{W6.}
At $k=0$ this is a bound of the kernel (8.0.1) at pairs of fermion sites, see (I.L0.17).
The second display of Proposition~3\,(a) with $S^{\#}=S^{\mathrm{sy}}$ and $y=0$, where $(1+|y|)^{-1}=1$, gives
\[
\left|\Gamma_0(s,A)(x,x')\right|_2\le32a'\cdot M^{27\alpha_0}\cdot\exp\left(-\gamma^F\cdot d(x,x')\right)\qquad\text{for all fermion sites }x,x'.
\]
One has $C_1\ge304\mathrm{e}\cdot\max(1,b)\cdot a'\cdot M^{27\alpha_0}\ge32a'\cdot M^{27\alpha_0}$, because $304\mathrm{e}\ge32$ and $\max(1,b)\ge1$.
With Lemma~8.A it follows that
\[
\left|\Gamma_0(s,A)(x,x')\right|_c\le64a'\cdot M^{27\alpha_0}\cdot\exp\left(-\gamma^F\cdot d(x,x')\right)\le2C_1\cdot\exp\left(-\gamma^F\cdot d(x,x')\right).
\]

\textbf{W7.}
At $k=0$ this is a bound of the kernel (8.0.1) with the factor $(1+y)^{-1}$, see (I.L0.17).
The second display of Proposition~3\,(a) with $S^{\#}=S$ gives, since $|y|=y$ for $y\ge0$,
\[
\left|S_{0,y}(X,A)(x,x')\right|_2\le32a'\cdot\exp(27\kappa')\cdot(1+y)^{-1}\cdot\exp\left(-\gamma^F\cdot d(x,x')-\kappa'\cdot d_M(X)\right).
\]
One has $C_1\ge32\mathrm{e}\cdot\max(1,b^2)\cdot a'\cdot\exp(27\kappa')\ge32a'\cdot\exp(27\kappa')$.
With Lemma~8.A it follows that
\begin{align*}
\left|S_{0,y}(X,A)(x,x')\right|_c&\le64a'\cdot\exp(27\kappa')\cdot(1+y)^{-1}\cdot\exp\left(-\gamma^F\cdot d(x,x')-\kappa'\cdot d_M(X)\right)\\
&\le2C_1\cdot(1+y)^{-1}\cdot\exp\left(-\gamma^F\cdot d(x,x')-\kappa'\cdot d_M(X)\right).
\end{align*}
The rates are $\gamma=\gamma^F$ and $\mathrm{kq}=\kappa'$.
The factor $48$ of the weight in (8.0.1) is contained in $a'$, through $C_{00}$ ((I.L0.13), (I.ke5)).

\textbf{W8, W9.}
Let $\mathscr{Y}$ be a subset of $\pi_0$ and $Y$ its union, let $s\in\mathrm{Pd}_0$ satisfy $s_{\mathrm{pc}}=0$ for $\mathrm{pc}\notin\mathscr{Y}$, and let $Y_1$ be a link-component of $Y$.
Let $\tilde{s}\in\mathrm{Pd}_0$ be a family that vanishes outside $\mathscr{Y}$ and coincides with $s$ on the cubes of $Y_1$,
and let $\tilde{A}\in\tilde{G}_0$ be a connection that coincides with $A$ on the bonds that meet $Y_1$.
Proposition~3\,(i) with $S^{\#}=S^{\mathrm{sy}}$ and $y=0$ asserts:
\begin{itemize}
\item if $x$ lies over $Y_1$, then $\Gamma_0(s,A)(x,x')=0$ unless $x'$ lies over $Y_1$, and $\Gamma_0(s,A)(x,\cdot\,)=\Gamma_0(\tilde{s},\tilde{A})(x,\cdot\,)$;
\item if $x'$ lies over $Y_1$, then $\Gamma_0(s,A)(x,x')=0$ unless $x$ lies over $Y_1$, and $\Gamma_0(s,A)(\cdot\,,x')=\Gamma_0(\tilde{s},\tilde{A})(\cdot\,,x')$.
\end{itemize}
By subsection~0.4, ``lies over $Y_1$'' means the same as ``is contained in $Y_1$''.
The first assertion is W8, the second is W9, for $\mathrm{Op}=\Gamma_0$.

\textbf{W12.}
Let $\mathrm{wk}$ be a walk with $X_{\mathrm{wk}}=X$.
By Proposition~3\,(d), the kernel $S_{y,\mathrm{wk}}(A)(x,x')$ vanishes unless $x$ and $x'$ lie over $X$.
By (8.0.3), the absolute convergence in $N_{\gamma^F}$ (Proposition~3\,(a)) and Lemma~8.B,
\[
S_{0,y}(X,A)(x,x')=\sum_{\mathrm{wk}:\,X_{\mathrm{wk}}=X}S_{y,\mathrm{wk}}(A)(x,x'),
\]
which therefore vanishes unless $x$ and $x'$ lie over $X$.
Now let $A,\tilde{A}\in\tilde{G}_0$ coincide on the bonds that meet $X$.
A bond with both end points in $X$ has an end point in $X$, that is, it meets $X$ (Section~I.1.1).
Hence $A$ and $\tilde{A}$ coincide on the bonds with both end points in $X_{\mathrm{wk}}=X$,
and $S_{y,\mathrm{wk}}(A)=S_{y,\mathrm{wk}}(\tilde{A})$ by Proposition~3\,(d), for every walk of the sum (8.0.3).
It follows that $S_{0,y}(X,A)=S_{0,y}(X,\tilde{A})$.

\textbf{W13.}
Let $\boldsymbol{\lambda}$ be a real function of $c$ on $T_0$.
Then $A-d\boldsymbol{\lambda}\in\tilde{G}_0$ by \cite[Section~16.2.7, (G2)]{QEDI}.
At level $0$ one has $T^0_0=T_0$ and $\boldsymbol{\lambda}^{u}=\boldsymbol{\lambda}$,
and by (I.POS) the phase acts as
\[
\left(\exp\left(ie_k\cdot\boldsymbol{\lambda}\right)\cdot f\right)(c,g)=\exp\left(ie_k\cdot\boldsymbol{\lambda}(c)\right)\cdot f(c,g),
\]
with the same number on the $48$ sites over $c$.
Proposition~3\,(e) with $S^{\#}=S^{\mathrm{sy}}$ and $y=0$, respectively with $S^{\#}=S$, gives
\begin{gather*}
\Gamma_0(s,A-d\boldsymbol{\lambda})=\exp\left(ie_k\cdot\boldsymbol{\lambda}\right)\cdot\Gamma_0(s,A)\cdot\exp\left(-ie_k\cdot\boldsymbol{\lambda}\right),\\
S_{0,y}(X,A-d\boldsymbol{\lambda})=\exp\left(ie_k\cdot\boldsymbol{\lambda}\right)\cdot S_{0,y}(X,A)\cdot\exp\left(-ie_k\cdot\boldsymbol{\lambda}\right).
\end{gather*}

\textbf{W14.}
Let $\tilde{g}=(S,a)\in\mathbb{G}_0$, with the actions of Section~I.4.3:
\[
\tilde{g}(c,g)=\left(r(S)c+a,\ r(S)g\right),\qquad(\tilde{g}f)(x)=S\cdot f(\tilde{g}^{-1}x),\qquad(\tilde{g}s)_{\mathrm{pc}}=s_{\tilde{g}^{-1}\mathrm{pc}}.
\]
Then $\tilde{g}s\in\mathrm{Pd}_0$, its entries being a permutation of those of $s$ (Lemma~8.C, proof of (iv));
$\tilde{g}A\in\tilde{G}_0$ by \cite[Section~16.2.7, (G2)]{QEDI};
and $\tilde{g}X\in D_0$.
For the last assertion, note that the affine map of $\mathbb{R}^3$ with the linear part $r(S)$ and the translation part $a$ maps the orbits of $\ell\mathbb{Z}^3$ onto such orbits, and maps the closed cubes of side $M$ with centres in $M\mathbb{Z}^3$, together with their faces, onto such cubes and faces;
hence $\tilde{g}$ maps $M$-cubes of $T_0$ with a common face onto $M$-cubes with a common face.
In the terms of Section~I.1.1: $r(S)$ is diagonal with entries $\pm1$, so that centres that differ by $\pm M\cdot e_\mu$ modulo $\ell\mathbb{Z}^3$ are mapped to centres that differ by $\pm M\cdot e_\mu$ modulo $\ell\mathbb{Z}^3$.
In Proposition~3\,(f) one has $y'=y$ if $S\in\mathrm{Sp}_3$ or if $y$ is neutral.
The value $y=0$ is neutral (subsection~0.5).
Hence for every $\tilde{g}\in\mathbb{G}_0$, Proposition~3\,(f) with $S^{\#}=S^{\mathrm{sy}}$ and $y=y'=0$ gives
\[
\Gamma_0(\tilde{g}s,\tilde{g}A)\cdot\tilde{g}=\tilde{g}\cdot\Gamma_0(s,A).
\]
For $\tilde{g}\in\mathbb{G}_{0,3}$, that is, for $S\in\mathrm{Sp}_3$, and for every $y\ge0$, Proposition~3\,(f) with $S^{\#}=S$ and $y'=y$ gives
\[
S_{0,y}(\tilde{g}X,\tilde{g}A)\cdot\tilde{g}=\tilde{g}\cdot S_{0,y}(X,A).
\]
Since $\mathbb{G}_0^{L}$ is a subgroup of $\mathbb{G}_0$ and $\mathbb{G}^{L}_{0,3}$ is a subgroup of $\mathbb{G}_{0,3}$ (Section~I.4.2), W14 holds in the form of Section~G1.11.

\textbf{W15.}
By \cite[Section~16.2.7, (G2)]{QEDI}, $-A\in\tilde{G}_0$;
and $y=0$ is neutral.
The second sentence of Proposition~3\,(g), at $y=0$, gives
\[
\Gamma_0(s,A)^{T}=S^{\mathrm{sy}}_0(s,A)^{T}=\mathsf{C}^{-1}\cdot S^{\mathrm{sy}}_0(s,-A)\cdot\mathsf{C}=\mathsf{C}^{-1}\cdot\Gamma_0(s,-A)\cdot\mathsf{C},
\]
the transpose being taken with respect to $\langle f',f\rangle^{\mathrm{ch}}=\frac{1}{48}\sum_c\sum_gf'(c,g)\cdot f(c,g)$.

\textbf{W16.}
One has $\overline{s}\in\mathrm{Pd}_0$, and $\overline{A}$, $-A$ lie in $\tilde{G}_0$ by \cite[Section~16.2.7, (G2)]{QEDI}.
Proposition~3\,(h) with $S^{\#}=S^{\mathrm{sy}}$ and $y=0$ gives
\[
\overline{\Gamma_0(\overline{s},\overline{A})}=\varkappa^{-1}\cdot\Gamma_0(s,-A)\cdot\varkappa.
\]

This proves (a):
the radius of $\mathrm{Pd}_0$ is $M^{\alpha_0}$, the rates are $\gamma^F$ and $\kappa'$, and the constant is $2C_1$.
\end{proof}

\begin{named}{Remark}[The identities W13--W16 in kernels]
Let $O\in\mathrm{Lin}$, let $\psi$, $\psi'$ be functions of $c$, and let $\tilde{g}=(S,a)\in\mathbb{G}_0$.
From (8.0.1) one obtains the following.
\begin{itemize}
\item $(\psi\cdot O\cdot\psi')(x,x')=\psi(c)\cdot O(x,x')\cdot\psi'(c')$ (Section~G1.2).
\item $\overline{O}(x,x')=\overline{O(x,x')}$, where $\overline{O}$ is the operator $f\mapsto\overline{O\overline{f}}$ ((I.O5)); this holds because the weight $\frac{1}{48}$ is real.
\item $O^{T}(x,x')=O(x',x)^{T}$, the transpose of the $4\times4$ matrix, because
\[
\langle f',Of\rangle^{\mathrm{ch}}=\int\mathrm{d}x\int\mathrm{d}x'\,f'(x)\cdot O(x,x')\cdot f(x')=\int\mathrm{d}x'\left(\int\mathrm{d}x\,O(x,x')^{T}\cdot f'(x)\right)\cdot f(x').
\]
\item $(\tilde{g}\cdot O\cdot\tilde{g}^{-1})(x,x')=S\cdot O(\tilde{g}^{-1}x,\tilde{g}^{-1}x')\cdot S^{-1}$, because $x\mapsto\tilde{g}x$ preserves $\int\mathrm{d}x$ (Section~I.4.3, (I.a1)).
\end{itemize}
Hence, with $r:=r(S)$ and $\tilde{g}^{-1}(c,g)=\left(r(c-a),rg\right)$,
\begin{gather*}
\Gamma_0(s,A-d\boldsymbol{\lambda})(x,x')=\exp\left(ie_k\cdot\left(\boldsymbol{\lambda}(c)-\boldsymbol{\lambda}(c')\right)\right)\cdot\Gamma_0(s,A)(x,x'),\\
\Gamma_0(\tilde{g}s,\tilde{g}A)(x,x')=S\cdot\Gamma_0(s,A)(\tilde{g}^{-1}x,\tilde{g}^{-1}x')\cdot S^{-1},\\
\Gamma_0(s,A)(x',x)^{T}=\mathsf{C}^{-1}\cdot\Gamma_0(s,-A)(x,x')\cdot\mathsf{C},\\
\overline{\Gamma_0(\overline{s},\overline{A})(x,x')}=\varkappa^{-1}\cdot\Gamma_0(s,-A)(x,x')\cdot\varkappa.
\end{gather*}
The first two of these identities hold also for $S_{0,y}(X,A)$, with $X$ and $\tilde{g}X$ in the places of $s$ and $\tilde{g}s$;
the second of the two holds for $S\in\mathrm{Sp}_3$.
\end{named}

\subsection*{3. Proof of (b)}

\begin{proof}[Proof of Proposition~8\,(b)]
\textbf{3.1. The factor.}
For $s\in\mathrm{Pd}_0$ and a site $c$ of $T_0$ let
\[
\mathrm{fa}_s(c):=\prod_{\mathrm{pc}\in\mathrm{Cub}(c)}s_{\mathrm{pc}},
\]
where $\mathrm{Cub}(c)$ is the set of Lemma~8.C.
The definition of $H_0(s,A)$ in Proposition~8, the meaning of ``meets'' fixed in subsection~0.7 and the equality $K_x=K_c$ of (I.POS) give
\begin{equation}
\left(H_0(s,A)f\right)(c,g)=\mathrm{fa}_s(c)\cdot f(c,g)\qquad\left(f\in\mathfrak{S}^{\mathrm{ch}}(1,\ell),\ g\in O_{\mathrm{cub}}\right).
\tag{8.3.1}
\end{equation}
Since $\mathrm{fa}_s$ is a function on $T_0$, it is an $\ell\mathbb{Z}^3$-periodic function of $c\in\mathbb{Z}^3$;
consequently $H_0(s,A)f$ is antiperiodic, and $H_0(s,A)$ belongs to $\mathrm{Lin}$.
The connection $A$ does not occur on the right-hand side of (8.3.1).
The number $\mathrm{fa}_s(c)$ does not involve $g$, so that it takes the same value at the $48$ sites over $c$;
and it multiplies the four components of $f(c,g)$ by one and the same factor, so that $H_0(s,A)$ acts in the spin indices as a multiple of $\mathbf{1}_4$.
Applying (8.0.2) to the function $\mathrm{fa}_s$ one obtains the kernel
\begin{equation}
H_0(s,A)(x,x')=48\cdot\mathbf{1}[x=x']\cdot\mathrm{fa}_s(c)\cdot\mathbf{1}_4\qquad\left(x=(c,g),\ x'\ \text{in one fundamental domain}\right).
\tag{8.3.2}
\end{equation}
In the row at $x$, and likewise in the column at $x$, at most one entry can be different from zero, namely the entry at $(x,x)$ (modulo $\ell\mathbb{Z}^3$);
it equals $48\,\mathrm{fa}_s(c)\cdot\mathbf{1}_4$.

The connections $-A$, $\overline{A}$, $A-d\boldsymbol{\lambda}$ with a real $\boldsymbol{\lambda}$, and $\tilde{g}A$ with $\tilde{g}\in\mathbb{G}_0$, which occur below as arguments, lie in $\tilde{G}_0$ by the invariance property (G2) of \cite[Section~16.2.7]{QEDI}.
The families $\overline{s}$ and $\tilde{g}s$ lie in $\mathrm{Pd}_0$:
the entries of $\overline{s}$ have the same moduli as those of $s$, and the entries of $\tilde{g}s$ are those of $s$ in another order.

The matrices $\mathsf{C}$ and $\varkappa$ are constant and invertible (Section~I.2.1);
hence every multiple of $\mathbf{1}_4$ commutes with $\mathsf{C}$.
Together with $H_0(s,-A)=H_0(s,A)$ this gives
\[
\bar{H}_0(s,A)=\mathsf{C}^{-1}\cdot H_0(s,-A)\cdot\mathsf{C}=H_0(s,-A)=H_0(s,A).
\]

\textbf{W1.}
For $s=1$ every factor of the product equals $1$.
Thus $\mathrm{fa}_1(c)=1$, and $H_0(1,A)=1=H_0(A)$ by subsection~0.5.

\textbf{W2.}
By (8.3.1) and Lemma~8.C\,(ii), each entry of $H_0(s,A)$ is, as a function of $(s,A)$, a monomial of degree at most $8$ in the finitely many variables $s_{\mathrm{pc}}$, and it is constant in $A$.
It follows that $H_0(s,A)$ is continuous on $\mathrm{Pd}_0\times\tilde{G}_0$ and analytic in $(s,A)$ in the interior,
and that for every fixed $s\in\mathrm{Pd}_0$, the boundary points of $\mathrm{Pd}_0$ included, it is analytic in $A\in\tilde{G}_0$.

\textbf{W5.}
One has $|s_{\mathrm{pc}}|\le M^{\alpha_0}$ and $M^{\alpha_0}\ge1$, because $M\ge1$ and $\alpha_0>0$;
moreover $\#\left|\mathrm{Cub}(c)\right|\le8$ by Lemma~8.C\,(ii).
Hence $\left|\mathrm{fa}_s(c)\right|\le M^{8\alpha_0}$, and (8.3.1) yields for $f\in\mathfrak{S}^{\mathrm{ch}}(1,\ell)$
\[
\left|H_0(s,A)f\right|_{\infty}=\sup_{(c,g)}\left|\mathrm{fa}_s(c)\right|\cdot\left|f(c,g)\right|_c\le M^{8\alpha_0}\cdot|f|_{\infty}.
\]
For every $\gamma\ge0$ one has
\[
M^{8\alpha_0}\le M^{27\alpha_0}\le4\exp(2\gamma)\cdot M^{27\alpha_0},
\]
since $M\ge1$, $8\alpha_0\le27\alpha_0$ and $4\exp(2\gamma)\ge1$.
The second part of W5 is a bound at the short pairs, that is, at the pairs with $d_X(x,x')<1$.
At level $0$, (I.1.5) gives
\[
d_X(x,x')=\max\left(|c-c'|,\eta_0\right)\ge1\qquad\text{for all }x\ne x';
\]
in particular each of the $48\cdot47$ ordered pairs inside a cluster has $d_X=1$.
There is therefore no short pair, and the second part of W5 is void.

\textbf{Locality.}
Let $\mathscr{Y}$ be a subset of $\pi_0$, let $Y$ be its union, and let $s\in\mathrm{Pd}_0$ satisfy $s_{\mathrm{pc}}=0$ for every $\mathrm{pc}\notin\mathscr{Y}$.

\textup{$(*)$} If $\mathrm{Cub}(c)$ is not contained in $\mathscr{Y}$, then $\mathrm{fa}_s(c)=0$, because one of the factors is $s_{\mathrm{pc}}$ with $\mathrm{pc}\notin\mathscr{Y}$.

\textbf{W11.}
According to (8.3.2), the row of $H_0(s,A)$ at $x=(c,g)$ is zero as soon as $\mathrm{fa}_s(c)=0$.
By $(*)$ this row therefore vanishes unless $\mathscr{Y}$ contains $\mathrm{Cub}(c)$, the set of all cubes of $\pi_0$ that meet $K_x=K_c$.

\textbf{W10.}
Let $x=(c,g)$ be contained in no link-component of $Y$.
By (l2), $Y$ is the union of its link-components;
so $c$ is not in $Y$, and the $M$-cube of $c$, which contains $c$, is not a cube of $\mathscr{Y}$.
This cube is an element of $\mathrm{Cub}(c)$ by Lemma~8.C\,(i).
Now $(*)$ gives $\mathrm{fa}_s(c)=0$, and the row at $x$ vanishes.

\textbf{W8, W9.}
Let $Y_1$ be a link-component of $Y$ and $\mathscr{Y}_1$ the link-component of $\mathscr{Y}$ whose union is $Y_1$.
Let $\tilde{s}\in\mathrm{Pd}_0$ be a family that vanishes outside $\mathscr{Y}$ and coincides with $s$ on the cubes of $Y_1$, and let $\tilde{A}\in\tilde{G}_0$.
Let $x=(c,g)$ be contained in $Y_1$, that is, $c\in Y_1$.

By (8.3.2), $H_0(s,A)(x,x')=0$ unless $x'=x$ modulo $\ell\mathbb{Z}^3$, and in that case $x'$ is contained in $Y_1$.
Interchanging the two arguments in (8.3.2), one finds for $x'=(c',g')$
\[
H_0(s,A)(x',x)=48\cdot\mathbf{1}[x'=x]\cdot\mathrm{fa}_s(c')\cdot\mathbf{1}_4,
\]
which is zero unless $x'=x$ modulo $\ell\mathbb{Z}^3$; and then $x'$ is contained in $Y_1$.

Consequently the row and the column of $H_0(s,A)$ at $x$ are determined by the number $\mathrm{fa}_s(c)$, and those of $H_0\left(\tilde{s},\tilde{A}\right)$ by the number $\mathrm{fa}_{\tilde{s}}(c)$.
It remains to show that $\mathrm{fa}_s(c)=\mathrm{fa}_{\tilde{s}}(c)$.

Case~1: $\mathrm{Cub}(c)$ is not contained in $\mathscr{Y}$.
Both $s$ and $\tilde{s}$ vanish outside $\mathscr{Y}$, and $(*)$, applied to $s$ and to $\tilde{s}$, gives $\mathrm{fa}_s(c)=0=\mathrm{fa}_{\tilde{s}}(c)$.

Case~2: $\mathrm{Cub}(c)$ is contained in $\mathscr{Y}$.
Since $c\in Y_1$, the site $c$ lies in a cube of $\mathscr{Y}_1$, and by (l1) this cube is the $M$-cube of $c$.
By Lemma~8.C\,(i) and~(iii), $\mathrm{Cub}(c)$ is a link-connected subset of $\mathscr{Y}$ that contains the $M$-cube of $c$;
by (l3) it is therefore contained in $\mathscr{Y}_1$.
The cubes of $\mathscr{Y}_1$ are the cubes of $Y_1$, and on these $\tilde{s}=s$.
So the products that define $\mathrm{fa}_s(c)$ and $\mathrm{fa}_{\tilde{s}}(c)$ have the same factors.

It follows that $H_0(s,A)(x,\cdot)=H_0\left(\tilde{s},\tilde{A}\right)(x,\cdot)$, which is W8, and that $H_0(s,A)(\cdot,x)=H_0\left(\tilde{s},\tilde{A}\right)(\cdot,x)$, which is W9.

\textbf{W13.}
Let $\boldsymbol{\lambda}$ be a real function of $c$ on $T_0$.
As $T^0_0=T_0$, one has $\boldsymbol{\lambda}^u=\boldsymbol{\lambda}$.
Since $H_0$ does not depend on the connection, $H_0(s,A-d\boldsymbol{\lambda})=H_0(s,A)$.
The multiplications by the scalar functions $\mathrm{fa}_s(c)$ and $\exp\left(\pm ie_k\cdot\boldsymbol{\lambda}(c)\right)$ commute with each other.
Therefore
\[
\exp\left(ie_k\cdot\boldsymbol{\lambda}\right)\cdot H_0(s,A)\cdot\exp\left(-ie_k\cdot\boldsymbol{\lambda}^u\right)=H_0(s,A)=H_0(s,A-d\boldsymbol{\lambda}).
\]

\textbf{W14, for all $\tilde{g}\in\mathbb{G}_0$.}
Let $\tilde{g}=(S,a)\in\mathbb{G}_0$, $f\in\mathfrak{S}^{\mathrm{ch}}(1,\ell)$ and $x=(c,g)$.
By (I.a1), $\pi\left(\tilde{g}^{-1}x\right)=\tilde{g}^{-1}c$.
From (8.3.1) and $\left(\tilde{g}f\right)(x)=S\cdot f\left(\tilde{g}^{-1}x\right)$ one obtains
\begin{align*}
\left(H_0\left(\tilde{g}s,\tilde{g}A\right)\cdot\tilde{g}f\right)(x)&=\mathrm{fa}_{\tilde{g}s}(c)\cdot S\cdot f\left(\tilde{g}^{-1}x\right),\\
\left(\tilde{g}\cdot H_0(s,A)f\right)(x)&=S\cdot\left(H_0(s,A)f\right)\left(\tilde{g}^{-1}x\right)=\mathrm{fa}_s\left(\tilde{g}^{-1}c\right)\cdot S\cdot f\left(\tilde{g}^{-1}x\right).
\end{align*}
Since $\left(\tilde{g}s\right)_{\mathrm{pc}}=s_{\tilde{g}^{-1}\mathrm{pc}}$, Lemma~8.C\,(iv), applied to $\tilde{g}^{-1}$, gives
\[
\mathrm{fa}_{\tilde{g}s}(c)=\prod_{\mathrm{pc}\in\mathrm{Cub}(c)}s_{\tilde{g}^{-1}\mathrm{pc}}=\prod_{\mathrm{pc}'\in\mathrm{Cub}\left(\tilde{g}^{-1}c\right)}s_{\mathrm{pc}'}=\mathrm{fa}_s\left(\tilde{g}^{-1}c\right).
\]
Hence $H_0\left(\tilde{g}s,\tilde{g}A\right)\cdot\tilde{g}=\tilde{g}\cdot H_0(s,A)$.
If the representative $a\in M\mathbb{Z}^3$ is changed by an element of $\ell\mathbb{Z}^3$, the operator $\tilde{g}$ is multiplied by the same sign on both sides of this identity.
As $\mathbb{G}_0^{L}$ is a subgroup of $\mathbb{G}_0$, the identity holds in particular for all $\tilde{g}\in\mathbb{G}_0^{L}$.

\textbf{W16.}
Let $f\in\mathfrak{S}^{\mathrm{ch}}(1,\ell)$.
By (I.O5), $\overline{O}$ denotes the operator $f\mapsto\overline{O\overline{f}}$; hence
\[
\left(\overline{H_0\left(\overline{s},\overline{A}\right)}f\right)(c,g)=\overline{\mathrm{fa}_{\overline{s}}(c)\cdot\overline{f(c,g)}}=\overline{\mathrm{fa}_{\overline{s}}(c)}\cdot f(c,g)=\mathrm{fa}_s(c)\cdot f(c,g),
\]
because the complex conjugate of $\prod_{\mathrm{pc}\in\mathrm{Cub}(c)}\overline{s_{\mathrm{pc}}}$ is $\prod_{\mathrm{pc}\in\mathrm{Cub}(c)}s_{\mathrm{pc}}$.
Every multiple of $\mathbf{1}_4$ commutes with $\varkappa$, and $H_0(s,-A)=H_0(s,A)$;
so $\varkappa^{-1}\cdot H_0(s,-A)\cdot\varkappa=H_0(s,A)$.
Comparing the two identities,
\[
\overline{H_0\left(\overline{s},\overline{A}\right)}=\varkappa^{-1}\cdot H_0(s,-A)\cdot\varkappa.
\]

The last assertion of (b), namely that $H_0(s,A)$ does not depend on $A$, is scalar in the spin indices and is the same at the $48$ sites of a cluster, was established in~3.1.
This completes the proof of (b).
\end{proof}

\subsection*{4. Uniformity; two sides}

\textbf{Uniformity.}
Proposition~8 is a statement on the single level $k=0$.
The numbers that enter its bounds are $\alpha_0$, $\gamma^{F}$, $\kappa'=\kappa_{\mathrm{w}}+\kappa_0$, $M$ and
\[
C_1=\max\left\{304\,\mathrm{e}\cdot\max(1,b)\cdot a'\cdot M^{27\alpha_0},\ C_{\Gamma}\cdot M^{27\alpha_0},\ 32\,\mathrm{e}\cdot\max(1,b^2)\cdot a'\cdot\exp(27\kappa')\right\}.
\]
They are constants of the lines (C4), (C4r), (C5), (C6) of Proposition~14.2\,(a), in the sense of (I.U4)$^{\mathrm{ch}}$ (Section~I.8.1, Section~G1.1);
none of them depends on $\ell$, $\mu$, $e_k$, $s$, $A$, $X$ or $y$.
Proposition~3 holds for all data (SF) with these constants, and apart from $L$ and $M$ the only numbers that appear in Lemmas~8.A--8.C are the numerals $2$, $48$ and $8$.
Consequently the bounds W5, W6 and W7 hold with the same constants for every side $\ell\in\mathrm{Cl}$, every mass $\mu\in(0,1]$ and every $e_k\in(0,e_0]$, and for all $s\in\mathrm{Pd}_0$, $A\in\tilde{G}_0$, $X\in D_0$, $y\ge0$.
For a flow, the cutoff $N$, the torus and the charge $e$ reach level $0$ only through the three data
\[
\ell=R/\eta,\qquad\mu=\bar{m}\cdot\eta,\qquad e_k=q=e\cdot\eta^{\frac{1}{2}}
\]
(Section~I.6.1, (I.L0.2));
hence no constant depends on them.
In subsection~0.7 and in Sections~1 and~3 the side $\ell$ enters only through the two facts that $M$ divides $\ell$ and that $2M<\ell$.

\textbf{Two sides.}
In Lemma~8.C\,(v), the sets $\mathrm{Nb}_\mu$ and the $M$-cubes $\tilde{\square}\left(z_{\square}+M\cdot\mathrm{jv}\right)$ of $\mathbb{Z}^3$ are formed from the lift $\tilde{c}$ and from $L$ and $M$ alone.
The side $\ell$ thus enters $\mathrm{Cub}(c)$, and with it $\mathrm{fa}_s(c)$, only through the covering map $\mathrm{pr}$:
\begin{equation}
\mathrm{fa}_s\left(\mathrm{pr}\,\tilde{c}\right)=\prod_{\mathrm{jv}\in\mathrm{Nb}_1\times\mathrm{Nb}_2\times\mathrm{Nb}_3}s_{\mathrm{pr}\,\tilde{\square}\left(z_{\square}+M\cdot\mathrm{jv}\right)},
\tag{8.4.1}
\end{equation}
where the cubes $\mathrm{pr}\,\tilde{\square}\left(z_{\square}+M\cdot\mathrm{jv}\right)$, $\mathrm{jv}\in\mathrm{Nb}_1\times\mathrm{Nb}_2\times\mathrm{Nb}_3$, are pairwise different.

Let $\ell_1,\ell_2\in\mathrm{Cl}$ be sides with $\ell_2\in\ell_1\mathbb{Z}$.
The quotient $\ell_2/\ell_1$ is odd, because $\ell_1/(L\cdot M)$ and $\ell_2/(L\cdot M)$ are odd.
Let $\mathrm{pr}_i:\mathbb{Z}^3\to\mathbb{Z}^3/\ell_i\mathbb{Z}^3$ ($i=1,2$) be the covering maps, let $\mathrm{pj}:\mathbb{Z}^3/\ell_2\mathbb{Z}^3\to\mathbb{Z}^3/\ell_1\mathbb{Z}^3$ be the covering map with $\mathrm{pj}\circ\mathrm{pr}_2=\mathrm{pr}_1$ (Section~I.6.3),
and let $\mathrm{Cub}_i(\cdot)$ denote the sets of Lemma~8.C on the torus of side $\ell_i$.
For $z\in M\mathbb{Z}^3$ one has $\mathrm{pj}\,\mathrm{pr}_2\,\tilde{\square}(z)=\mathrm{pr}_1\,\tilde{\square}(z)$;
so $\mathrm{pj}$ maps $M$-cubes onto $M$-cubes.
Fix $\tilde{c}\in\mathbb{Z}^3$.
Lemma~8.C\,(v), applied on each of the two tori with this same $\tilde{c}$, shows that for $i=1,2$ the map $\mathrm{jv}\mapsto\mathrm{pr}_i\,\tilde{\square}\left(z_{\square}+M\cdot\mathrm{jv}\right)$ is a bijection of one and the same set $\mathrm{Nb}_1\times\mathrm{Nb}_2\times\mathrm{Nb}_3$ onto $\mathrm{Cub}_i\left(\mathrm{pr}_i\,\tilde{c}\right)$.
As $\mathrm{pj}$ maps $\mathrm{pr}_2\,\tilde{\square}\left(z_{\square}+M\cdot\mathrm{jv}\right)$ to $\mathrm{pr}_1\,\tilde{\square}\left(z_{\square}+M\cdot\mathrm{jv}\right)$ for every $\mathrm{jv}$, the map $\mathrm{pc}\mapsto\mathrm{pj}\,\mathrm{pc}$ is a bijection of $\mathrm{Cub}_2\left(\mathrm{pr}_2\,\tilde{c}\right)$ onto $\mathrm{Cub}_1\left(\mathrm{pr}_1\,\tilde{c}\right)$.
Let $s=(s_{\mathrm{pc}})$ be a family on the $M$-cubes of the torus of side $\ell_1$, and let $\mathrm{pj}^{*}s$ be the family on the $M$-cubes of the torus of side $\ell_2$ given by $\left(\mathrm{pj}^{*}s\right)_{\mathrm{pc}}:=s_{\mathrm{pj}\,\mathrm{pc}}$.
Then (8.4.1) yields
\begin{equation}
\begin{split}
\mathrm{fa}_{\mathrm{pj}^{*}s}\left(\mathrm{pr}_2\,\tilde{c}\right)
&=\prod_{\mathrm{jv}\in\mathrm{Nb}_1\times\mathrm{Nb}_2\times\mathrm{Nb}_3}s_{\mathrm{pj}\,\mathrm{pr}_2\,\tilde{\square}\left(z_{\square}+M\cdot\mathrm{jv}\right)}\\
&=\prod_{\mathrm{jv}\in\mathrm{Nb}_1\times\mathrm{Nb}_2\times\mathrm{Nb}_3}s_{\mathrm{pr}_1\,\tilde{\square}\left(z_{\square}+M\cdot\mathrm{jv}\right)}
=\mathrm{fa}_s\left(\mathrm{pr}_1\,\tilde{c}\right).
\end{split}
\tag{8.4.2}
\end{equation}

\subsection*{5. Conditions on the constants}

The proofs of this chapter use the following conditions on the constants.

\noindent
\begin{tabular}{|p{0.27\textwidth}|p{0.33\textwidth}|p{0.28\textwidth}|}
\hline
Condition & Where it is used & Place in Section~I.8.2 / Proposition~14.2\\
\hline
(M2): $M>4L$
& proof of Lemma~8.C, Step~2 ($2L\le(M-1)/2$); through it, in (b): W2 (degree), W5, W8, W9
& Proposition~14.2\,(b), (M2); line (C5)\\
\hline
(M1), (L5): $M$ is a power of $L$, and $L$ is odd
& subsection~0.7 ($M$ is an odd integer); W5 ($M\ge1$)
& Proposition~14.2\,(b), (M1), (L5); lines (C5), (C3); also Section~I.0.5 and Section~I.1.1 (Faces)\\
\hline
$\alpha_0=1/54>0$
& W5 ($M^{\alpha_0}\ge1$; $8\alpha_0\le27\alpha_0$)
& Proposition~14.2\,(a), line (C4)\\
\hline
$C_1\ge304\,\mathrm{e}\cdot\max(1,b)\cdot a'\cdot M^{27\alpha_0}$, \ $C_1\ge32\,\mathrm{e}\cdot\max(1,b^2)\cdot a'\cdot\exp(27\kappa')$
& W6, W7
& definition of $C_1$: Proposition~14.2\,(a), line (C6)\\
\hline
(Z2.s1): $6L\cdot e_0^{\varepsilon}\le1$; \ (Z2.s2): $4C'_{V}\cdot a'\cdot e_0^{\varepsilon}\le1$
& subsection~0.6 ($\delta=2e_k^{\varepsilon}$ is admissible in Proposition~2.5\,(a))
& Proposition~14.2\,(b), under these labels; line (C8); both are conditions of Proposition~3\\
\hline
$\varepsilon>0$
& subsection~0.6 ($e_k^{\varepsilon}\le e_0^{\varepsilon}$ for $e_k\in(0,e_0]$, in both lines of the display)
& Proposition~14.2\,(a), line (C7): $\varepsilon\in(1-\alpha,\bar{\varepsilon}]$, with $\alpha<1$ of line (C1); Proposition~14.2\,(b), (S1); also Section~I.0.5\\
\hline
the conditions of Proposition~3: (L5); (M1); (M2); (M5) in full; (Z2.s1); (Z2.s2); the range $\alpha<1$, which is among the conditions of Propositions~2.7 and~2.8
& through Proposition~3\,(a)--(i), Proposition~2.8\,(a), Proposition~2.5\,(a) and Proposition~2.6\,(a)
& Proposition~14.2\,(b), under these labels; lines (C3), (C5), (C8); $\alpha<1$: the range of $\alpha$ in Proposition~14.2\,(a), line (C1)\\
\hline
\end{tabular}

\medskip
All these conditions are imposed in Proposition~14.2:
(L5), (M1), (M2), (M5), (Z2.s1), (Z2.s2) and (S1) are conditions of its part~(b),
and the value $\alpha_0=1/54$, the two lower bounds for $C_1$, the range $\alpha<1$ of line (C1) and the range $\varepsilon\in(1-\alpha,\bar{\varepsilon}]$ of line (C7), which gives $\varepsilon>0$,
are fixed by the definitions and ranges of its part~(a).

The side $\ell$ is a datum of Proposition~8, not a constant:
it ranges over the class $\mathrm{Cl}$ of the sides (Section~I.1.1), for which the definitions of Part~I are made.
Besides the two facts that $M$ divides $\ell$ and that $2M<\ell$, which are all that subsection~0.7 and Sections~1 and~3 need of $\ell$,
the membership $\ell\in\mathrm{Cl}$ enters subsections~0.5 and~0.6 and Section~2 through Propositions~2.5\,(a), 2.6\,(a), 2.8\,(a) and~3\,(a)--(i), which are stated for $\ell\in\mathrm{Cl}$.
These Propositions rest on consequences of $\ell\in\mathrm{Cl}$, among them $3M<\ell$ and $3M-1\le\ell/2$ (Proposition~2.6), $3M<\ell/2$ and the fact that $T_k$ is a block torus (Chapter~2b), and $3M+\eta_k<\ell$ (Chapter~3).
In Section~4 one needs, in addition, that $\ell_1/(L\cdot M)$ and $\ell_2/(L\cdot M)$ are odd integers;
hence the quotient $\ell_2/\ell_1$, which is the number $\mathrm{Jr}$ of Section~I.6.3, is odd, as that section requires.

\section*{Chapter 15. Grassmann algebras, norms, Gaussian integration on an index set with weights: proofs of Proposition 15}

The statement proved in this chapter is Proposition~15 of Part~II.

\subsection*{Proof}

Section~0 fixes the conventions of the chapter and establishes seven facts on Grassmann algebras; the clauses (a)--(n) are then proved in their order.
By 0.1\,(iii), 0.6\,(iii) and the paragraph ``What $m(J)$ depends on'' in (e), the norms, $\int\mathrm{D}W$, $\mathrm{d}\mu_I(W)$ and $\mathrm{d}\mu_{\mathrm{Cov}}$ do not depend on the orders, and $\mathrm{sn}(J)$ and $m(J)$ depend only on the order of the elements of $J$ among themselves, as stated in (D2).
For an empty set $\mathrm{Mo}$, (D2) states that $\mathrm{Gr}[\phi,W]=\mathrm{Gr}[\phi]$, $\mathrm{Top}=1$ and $\int\mathrm{D}W$ is the identity, and that the assertions of Proposition~15 hold with an arbitrary number $\mathsf{w}_W>0$, except the sentence of (f) according to which $\mathsf{w}_W$ cannot be lowered;
this follows from 0.6\,(iv) and from the section ``The empty set of modes'' following (m).
The factors of $\mathrm{Top}$ in (D3) are even, so that their order is immaterial, as is seen at the beginning of 0.6;
and the integral defined in (D3) is the integral of (I.3.5), by 0.6\,(v).

\subsubsection*{0. Conventions of this chapter; seven facts}

\paragraph{0.0. Labels and letters.}
In this chapter (D1)--(D5) and (a)--(n) are the data and the clauses of Proposition~15, and (G2.0.1) is the display in (G2.r7) of Section~G2.0.5.
The labels 0.0--0.7, (b1)--(b4), (d$'$), (f$'$), (l$'$), (n1)--(n8), (1)--(3), (1$'$)--(3$'$) and the displays (15.1)--(15.31) refer to this chapter.

The following notation is used throughout:
for a finite set $I$, $\#|I|$ is the number of its elements, and $\cup$, $\times$, $\setminus$ are, as usual, the union, the Cartesian product and the difference of sets;
$\mathrm{Gen}$ is a set of generators, $\mathrm{ge}$ a generator and $\widehat{\mathrm{ge}}$ its image under a homomorphism;
$\mathrm{pm}$ is a permutation, $\mathrm{id}$ the identity permutation, and $\mathrm{sgn}$ the sign of a permutation (in (15.16) the sign of a sequence);
$\mathrm{dg}$, $\mathrm{dg}'$ are degrees, and $\mathrm{hm}$ is a homomorphism.
In Step~1 of (i) and in (j) the symbols $\mathrm{ia}$, $\mathrm{ib}$ are indices of places;
in (j), $\mathrm{nt}$ is a number of points, $\mathrm{Vs}$ a normed space and $\mathrm{fc}$ a function on the circle;
$\mathrm{fdo}$ and $\widetilde{\mathrm{fdo}}$ are fundamental domains, and $\mathrm{sz}$ is a sign in (n5).
Moreover,
\[
\mathrm{ww}(\mathrm{mo}):=W(\mathrm{mo})\cdot\bar{W}(\mathrm{mo}),\qquad
\mathrm{Ew}:=\exp\left(-\langle\bar{W},W\rangle_\nu\right),\qquad
\mathrm{Iw}({F}):=\int\mathrm{d}\mu_I(W)\,F,
\]
and for $J\subset\mathrm{Mo}\times\{0,1\}$ the set $\mathrm{md}(J)$ consists of the modes of $J$, that is, of the $\mathrm{mo}$ with $(\mathrm{mo},0)\in J$ or $(\mathrm{mo},1)\in J$.
The letters $J$, $J'$ denote subsets of $\mathrm{Mo}\times\{0,1\}$;
the letters $p$, $q$, $r$, $n$, $s$, $t$, $v$, $f$, the kernels $K_n$, the moments $m(J)$, the families of generators $\theta_j$, the numbers $C_{ij}$ and the homomorphisms $s_U$ have the meaning that they have in Proposition~15.
In 0.1, 0.2, 0.4, (b3), (b4) and in Step~1 of (i), however, $n$ is a natural number (a number of factors, an exponent, or half the size of a matrix), and in 0.4 the $\theta_j$ are anticommuting elements.

\paragraph{0.1. Generators, monomials, signs.}
Let $\mathrm{Gen}$ be a finite totally ordered set.
The Grassmann algebra over $\mathbb{C}$ with the independent generators $\mathrm{ge}\in\mathrm{Gen}$ is the associative complex algebra with identity element that is generated by the elements $\mathrm{ge}\in\mathrm{Gen}$, in which
\begin{equation}
\mathrm{ge}\cdot\mathrm{ge}'+\mathrm{ge}'\cdot\mathrm{ge}=0\quad\text{for all }\mathrm{ge},\mathrm{ge}'\in\mathrm{Gen}\qquad\left(\text{in particular }\mathrm{ge}\cdot\mathrm{ge}=0\right),
\tag{15.1}
\end{equation}
and in which the $2^{\#|\mathrm{Gen}|}$ ordered monomials $\mathrm{ge}(I)$, $I\subset\mathrm{Gen}$, form a basis of the algebra as a complex vector space;
here $\mathrm{ge}(I)$ denotes the product of the elements of $I$ in increasing order, and $\mathrm{ge}(\emptyset)=1$.
The set of generators of $\mathrm{Gr}[\phi_1,\dots,\phi_p]$ is $\mathrm{Gen}=\{\phi_i(t):1\le i\le p,\ t\in\mathrm{Ind}_i\}$, ordered first by $i$ and then by $t$.
In $\mathrm{Gr}[\phi,W]:=\mathrm{Gr}[\phi_1,\dots,\phi_p,W]$ the generators $W(u)$, $u\in\mathrm{Mo}\times\{0,1\}$, come after these, in the order fixed in (D2).
The ordered monomials are thus $\phi_1(I_1)\cdots\phi_p(I_p)$, respectively $\phi_1(I_1)\cdots\phi_p(I_p)\cdot W(J)$, as in (D1), (D3).
With $I=(I_1,\dots,I_p)$ the following abbreviations are used:
\[
\phi(I):=\phi_1(I_1)\cdots\phi_p(I_p),\qquad
\#|I|:=\#|I_1|+\dots+\#|I_p|,\qquad
\mathsf{w}^I:=\mathsf{w}_1^{\#|I_1|}\cdots\mathsf{w}_p^{\#|I_p|},\qquad
\nu(I):=\prod_i\nu_i(I_i).
\]
The $p$-tuple $I$ is identified with the subset $\{\phi_i(t):t\in I_i,\ 1\le i\le p\}$ of $\mathrm{Gen}$; then $\phi(I)=\mathrm{ge}(I)$.
A \emph{monomial of} $F$ is an ordered monomial whose coefficient in $F$ is different from zero.
\begin{itemize}
\item[(i)] Let $\mathrm{ge}_1,\dots,\mathrm{ge}_n\in\mathrm{Gen}$.
If two of them coincide, then $\mathrm{ge}_1\cdots\mathrm{ge}_n=0$.
If they are pairwise distinct, then $\mathrm{ge}_{\mathrm{pm}(1)}\cdots\mathrm{ge}_{\mathrm{pm}(n)}=\mathrm{sgn}(\mathrm{pm})\cdot\mathrm{ge}_1\cdots\mathrm{ge}_n$ for every permutation $\mathrm{pm}$ of $\{1,\dots,n\}$.
Indeed, by (15.1) the interchange of two neighbouring factors multiplies a product by $-1$.
Every permutation $\mathrm{pm}$ is a product of interchanges of neighbours, and the number of these interchanges has the parity of $\mathrm{pm}$.
If two factors coincide, interchanges of neighbours bring them next to each other, and $\mathrm{ge}\cdot\mathrm{ge}=0$.
\item[(ii)] For subsets $I$, $I'$ of $\mathrm{Gen}$ one has $\mathrm{ge}(I)\cdot\mathrm{ge}(I')=0$ if $I$ and $I'$ meet;
otherwise $\mathrm{ge}(I)\cdot\mathrm{ge}(I')$ equals $+\mathrm{ge}(I\cup I')$ or $-\mathrm{ge}(I\cup I')$, with a sign that depends on $(I,I')$ alone.
This is a consequence of (i).
\item[(iii)] When the order of $\mathrm{Gen}$ is replaced by another order, each ordered monomial gets multiplied by $1$ or by $-1$, by (i).
The moduli of the coefficients $c_F$, and hence all norms (I.5.1), are therefore independent of the orders.
\end{itemize}

\paragraph{0.2. Degree and parity.}
An element $F$ is called homogeneous of degree $\mathrm{dg}$ if every monomial $\mathrm{ge}(I)$ of $F$ has $\#|I|=\mathrm{dg}$, and even (odd) if every monomial of $F$ has even (odd) $\#|I|$.
The element $0$ is homogeneous of every degree, and it is both even and odd.
\begin{itemize}
\item[(i)] If two elements are homogeneous of the degrees $\mathrm{dg}$ and $\mathrm{dg}'$, their product is homogeneous of degree $\mathrm{dg}+\mathrm{dg}'$, by 0.1\,(ii).
So the even elements form a subalgebra that contains $1$, and a product of $n$ odd elements is even if $n$ is even and odd if $n$ is odd.
\item[(ii)] For subsets $I$, $I'$ of $\mathrm{Gen}$,
\[
\mathrm{ge}(I')\cdot\mathrm{ge}(I)=(-1)^{\#|I|\cdot\#|I'|}\cdot\mathrm{ge}(I)\cdot\mathrm{ge}(I').
\]
Both sides vanish if $I$ and $I'$ meet.
If they are disjoint, the $\#|I'|$ factors of $\mathrm{ge}(I')$ are moved, one after the other, past the $\#|I|$ factors of $\mathrm{ge}(I)$, and this takes $\#|I|\cdot\#|I'|$ interchanges of neighbours.
\item[(iii)] From (ii) and bilinearity one obtains: an even element commutes with every element; two odd elements $F$, $G$ anticommute, $F\cdot G=-G\cdot F$; and the square of an odd element is zero, since $F\cdot F=-F\cdot F$.
\end{itemize}

\paragraph{0.3. Homomorphisms.}
Homomorphisms are always understood as homomorphisms of complex algebras with identity element; in particular they map $1$ to $1$.

\begin{named}{Lemma 0.3}
Let $\mathrm{Gr}$ and $\mathrm{Gr}'$ be Grassmann algebras, let $\mathrm{Gen}$ be the set of generators of $\mathrm{Gr}$, and suppose that for every $\mathrm{ge}\in\mathrm{Gen}$ an odd element $\widehat{\mathrm{ge}}$ of $\mathrm{Gr}'$ is given.
Then there exists exactly one homomorphism $\mathrm{hm}$ of $\mathrm{Gr}$ to $\mathrm{Gr}'$ with $\mathrm{hm}(\mathrm{ge})=\widehat{\mathrm{ge}}$ for all $\mathrm{ge}\in\mathrm{Gen}$.
It maps even elements to even elements and odd elements to odd elements.
\end{named}

\begin{proof}
Since the elements $\mathrm{ge}$ generate $\mathrm{Gr}$, there is at most one such homomorphism.
To construct it, let $\mathrm{hm}$ be the linear map with $\mathrm{hm}(1):=1$ that sends $\mathrm{ge}(I)$ to the product of the elements $\widehat{\mathrm{ge}}$, $\mathrm{ge}\in I$, taken in the increasing order of the $\mathrm{ge}$.
According to 0.2\,(iii) the elements $\widehat{\mathrm{ge}}$ satisfy the relations (15.1):
\[
\widehat{\mathrm{ge}}\cdot\widehat{\mathrm{ge}'}+\widehat{\mathrm{ge}'}\cdot\widehat{\mathrm{ge}}=0\qquad\text{for all }\mathrm{ge},\mathrm{ge}'\in\mathrm{Gen}.
\]
The facts 0.1\,(i), (ii) were deduced from (15.1) alone, so they hold for products of the elements $\widehat{\mathrm{ge}}$, with the same signs.
Consequently, if $I$ and $I'$ meet, then $\mathrm{hm}(\mathrm{ge}(I))\cdot\mathrm{hm}(\mathrm{ge}(I'))=0=\mathrm{hm}(\mathrm{ge}(I)\cdot\mathrm{ge}(I'))$;
if they are disjoint, then $\mathrm{hm}(\mathrm{ge}(I))\cdot\mathrm{hm}(\mathrm{ge}(I'))=\pm\mathrm{hm}(\mathrm{ge}(I\cup I'))$ with the sign of 0.1\,(ii), and this is $\mathrm{hm}(\mathrm{ge}(I)\cdot\mathrm{ge}(I'))$.
By bilinearity, $\mathrm{hm}(F\cdot G)=\mathrm{hm}({F})\cdot\mathrm{hm}(G)$ for all $F,G\in\mathrm{Gr}$.
The statement on parity follows from 0.2\,(i).
\end{proof}

In every substitution that occurs in Proposition~15 the image of a generator is a linear combination of generators, and therefore an odd element.
Lemma~0.3 thus defines the following maps and elements:
the map $F\mapsto F'$ in (d);
the homomorphism $s_U$ in (h);
in (j), the element $F(\phi+t\theta)$, the image of $F$ under $\phi_i(s)\mapsto\phi_i(s)+t\cdot\theta_i(s)$;
in (l), the element $f(\bar{\Psi},\mathrm{Cov}\,\Psi)$, and the element $E\left(H^{(1)}\Psi,H^{(2)}\Psi\right)$, the image of $E$ under $\psi(s)\mapsto(H^{(1)}\Psi)(s)$, $\chi(s)\mapsto(H^{(2)}\Psi)(s)$, where $(H^{(i)}\Psi)(s)$ are the elements of (D2).

\paragraph{0.4. The Leibniz formula.}
The following lemma is the formula (cl5) of \cite[Section~13]{QEDI}, with the Grassmann algebra replaced by an arbitrary associative complex algebra.

\begin{named}{Lemma 0.4}[Leibniz formula]
Let $\theta_1,\dots,\theta_n$ be elements of an associative complex algebra such that $\theta_i\cdot\theta_j+\theta_j\cdot\theta_i=0$ for all $i$, $j$, the case $i=j$ included, and let $\mathrm{Mx}$ be a complex $n\times n$ matrix.
Then
\begin{equation}
\prod_{i=1}^{n}\left(\sum_{j=1}^{n}\mathrm{Mx}_{ij}\cdot\theta_j\right)=\det\mathrm{Mx}\cdot\theta_1\cdots\theta_n,
\tag{15.2}
\end{equation}
where the factors of the product on the left are taken in the order $i=1,\dots,n$.
\end{named}

\begin{proof}
Expanding the product, the left side becomes the sum, over all maps $i\mapsto j(i)$ of $\{1,\dots,n\}$ into itself, of the terms $\mathrm{Mx}_{1j(1)}\cdots\mathrm{Mx}_{nj(n)}\cdot\theta_{j(1)}\cdots\theta_{j(n)}$.
The argument of 0.1\,(i) shows that a term with $j(i)=j(i')$ for some $i\ne i'$ vanishes, and that $\theta_{j(1)}\cdots\theta_{j(n)}=\mathrm{sgn}(j)\cdot\theta_1\cdots\theta_n$ when $j$ is a permutation.
It remains to note that the sum of $\mathrm{sgn}(j)\cdot\prod_i\mathrm{Mx}_{ij(i)}$ over all permutations $j$ is $\det\mathrm{Mx}$.
\end{proof}

\paragraph{0.5. Linear maps of the functions on $\mathrm{Mo}$.}
For a linear map $O$ with the kernel $O(\mathrm{mo},\mathrm{mo}')$ put
\begin{equation}
[O]_{\mathrm{mo}\,\mathrm{mo}'}:=O(\mathrm{mo},\mathrm{mo}')\cdot\nu(\mathrm{mo}'),\qquad\text{so that}\qquad
(Of)(\mathrm{mo})=\sum_{\mathrm{mo}'}[O]_{\mathrm{mo}\,\mathrm{mo}'}\cdot f(\mathrm{mo}').
\tag{15.3}
\end{equation}
\begin{itemize}
\item[(i)] The matrix of $O$ in the basis formed by the indicator functions of the points of $\mathrm{Mo}$ is $[O]$.
Consequently every linear map has exactly one kernel, $\det O=\det[O]$, $[O_1O_2]=[O_1]\cdot[O_2]$, and $[1]$ is the identity matrix; in other words, the kernel of the identity map is $\mathbf{1}[\mathrm{mo}=\mathrm{mo}']\cdot\nu(\mathrm{mo})^{-1}$.
\item[(ii)] The transpose $O^T$ is the linear map with $\langle O^Tf',f\rangle_\nu=\langle f',Of\rangle_\nu$ for all $f$, $f'$; there is only one such map, because all $\nu(\mathrm{mo})$ are positive.
Written out, the two sides read
\begin{align*}
\langle f',Of\rangle_\nu&=\sum_{\mathrm{mo},\mathrm{mo}'}\nu(\mathrm{mo})\cdot\nu(\mathrm{mo}')\cdot O(\mathrm{mo},\mathrm{mo}')\cdot f'(\mathrm{mo})\cdot f(\mathrm{mo}'),\\
\langle O^Tf',f\rangle_\nu&=\sum_{\mathrm{mo},\mathrm{mo}'}\nu(\mathrm{mo}')\cdot\nu(\mathrm{mo})\cdot O^T(\mathrm{mo}',\mathrm{mo})\cdot f'(\mathrm{mo})\cdot f(\mathrm{mo}').
\end{align*}
Since the two sums coincide for all $f$, $f'$ and all weights are positive, it follows that
\begin{equation}
O^T(\mathrm{mo}',\mathrm{mo})=O(\mathrm{mo},\mathrm{mo}'),\qquad
[O^T]_{\mathrm{mo}'\,\mathrm{mo}}=\nu(\mathrm{mo}')^{-1}\cdot[O]_{\mathrm{mo}\,\mathrm{mo}'}\cdot\nu(\mathrm{mo}).
\tag{15.4}
\end{equation}
Thus $[O^T]=\mathrm{diag}(\nu)^{-1}\cdot[O]^t\cdot\mathrm{diag}(\nu)$, where $[O]^t$ is the transposed matrix and $\mathrm{diag}(\nu)$ the diagonal matrix with the entries $\nu(\mathrm{mo})$.
Hence $\det O^T=\det O$ and $(O_1O_2)^T=O_2^TO_1^T$, and every invertible $U$ satisfies $(U^{-1})^T=(U^T)^{-1}$ and $\det(U^{-1})^T=(\det U)^{-1}$.
\item[(iii)] Let $\mathrm{el}=(\mathrm{el}(\mathrm{mo}))$ and $\mathrm{el}'=(\mathrm{el}'(\mathrm{mo}))$, $\mathrm{mo}\in\mathrm{Mo}$, be families of elements of a Grassmann algebra, and let $(O\,\mathrm{el})(\mathrm{mo})$ and $\langle\mathrm{el}',\mathrm{el}\rangle_\nu$ be formed as in (D2).
By (15.3), $(O\,\mathrm{el})(\mathrm{mo})=\sum_{\mathrm{mo}'}[O]_{\mathrm{mo}\,\mathrm{mo}'}\cdot\mathrm{el}(\mathrm{mo}')$, and by (i), $O_1(O_2\,\mathrm{el})=(O_1O_2)\,\mathrm{el}$.
From (15.3) and (15.4) one obtains, without moving any factor,
\begin{equation}
\langle O^T\mathrm{el}',\mathrm{el}\rangle_\nu=\sum_{\mathrm{mo},\mathrm{mo}'}\nu(\mathrm{mo})\cdot\nu(\mathrm{mo}')\cdot O(\mathrm{mo},\mathrm{mo}')\cdot\mathrm{el}'(\mathrm{mo})\cdot\mathrm{el}(\mathrm{mo}')=\langle\mathrm{el}',O\,\mathrm{el}\rangle_\nu.
\tag{15.5}
\end{equation}
A homomorphism $\mathrm{hm}$ with $\mathrm{hm}(W(\mathrm{mo}))=\mathrm{el}(\mathrm{mo})$ and $\mathrm{hm}(\bar{W}(\mathrm{mo}))=\mathrm{el}'(\mathrm{mo})$ for all $\mathrm{mo}$ sends $\langle\bar{W},OW\rangle_\nu$ to $\langle\mathrm{el}',O\,\mathrm{el}\rangle_\nu$.
\end{itemize}

\paragraph{0.6. The map $\int\mathrm{D}W$.}
Let $J_f:=\mathrm{Mo}\times\{0,1\}$.
The factors $W(\mathrm{mo})\cdot\bar{W}(\mathrm{mo})$ of $\mathrm{Top}$ are even by 0.2\,(i), and so they commute with one another by 0.2\,(iii); for this reason $\mathrm{Top}$ is independent of the order of its factors.
In the product $\mathrm{Top}$ each of the $2\cdot\#|\mathrm{Mo}|$ generators $W(u)$, $u\in J_f$, occurs exactly once.
Hence, by 0.1\,(i), $W(J_f)=\mathrm{sn}_f\cdot\mathrm{Top}$ for some $\mathrm{sn}_f\in\{1,-1\}$.
Since the elements $W(J)$, $J\subset J_f$, form a basis of $\mathrm{Gr}[W]$, the coefficient of $\mathrm{Top}$ in $W(J)$ equals $\mathrm{sn}_f$ for $J=J_f$ and $0$ for $J\ne J_f$, and (D3) takes the form
\begin{equation}
\int\mathrm{D}W\,F=\mathrm{sn}_f\cdot\left(\prod_{\mathrm{mo}}\nu(\mathrm{mo})^{-1}\right)\cdot\sum_Ic_F(I,J_f)\cdot\phi(I)\qquad\text{for }F=\sum_{I,J}c_F(I,J)\cdot\phi(I)\cdot W(J).
\tag{15.6}
\end{equation}
\begin{itemize}
\item[(i)] It follows from (15.6) that $\int\mathrm{D}W\,F=0$ whenever no monomial of $F$ contains all the generators $W(u)$, $u\in J_f$.
Moreover $\int\mathrm{D}W\,\mathrm{Top}=\prod_{\mathrm{mo}}\nu(\mathrm{mo})^{-1}$, because $\mathrm{Top}=\mathrm{sn}_f\cdot W(J_f)$ and $\mathrm{sn}_f\cdot\mathrm{sn}_f=1$.
\item[(ii)] For $F_0\in\mathrm{Gr}[\phi]$ and $F\in\mathrm{Gr}[\phi,W]$ one has $\int\mathrm{D}W\,(F_0\cdot F)=F_0\cdot\int\mathrm{D}W\,F$.
Indeed, both sides are bilinear in $(F_0,F)$.
Let $F_0=\phi(I')$ and $F=\phi(I)\cdot W(J)$.
By 0.1\,(ii) the product $\phi(I')\cdot\phi(I)$ is $0$ or $\pm$ an ordered monomial of $\mathrm{Gr}[\phi]$, so that each side equals $\phi(I')\cdot\phi(I)$ times $\prod_{\mathrm{mo}}\nu(\mathrm{mo})^{-1}$ times the coefficient of $\mathrm{Top}$ in $W(J)$.
\item[(iii)] The map $\int\mathrm{D}W$ does not depend on the orders of the index sets.
To see this, let $\mathrm{tc}(G)$, for $G\in\mathrm{Gr}[W]$, be the number such that $G-\mathrm{tc}(G)\cdot\mathrm{Top}$ is a linear combination of products of fewer than $2\cdot\#|\mathrm{Mo}|$ generators.
Such a number exists: $G$ is a linear combination of the $W(J)$, one has $W(J_f)=\mathrm{sn}_f\cdot\mathrm{Top}$, and for $J\ne J_f$ the element $W(J)$ is a product of fewer than $2\cdot\#|\mathrm{Mo}|$ generators.
It is unique: by 0.1\,(i) a product of fewer than $2\cdot\#|\mathrm{Mo}|$ generators $W(u)$ is $0$ or $\pm W(J)$ with $J\ne J_f$, and $\mathrm{Top}=\mathrm{sn}_f\cdot W(J_f)$ is linearly independent of these basis elements.
The definition of $\mathrm{tc}(G)$ involves no order, and for every choice of the orders $\mathrm{tc}(W(J))$ is the coefficient of $\mathrm{Top}$ in $W(J)$.
By (ii), therefore, $\int\mathrm{D}W$ is the linear map with
\[
\int\mathrm{D}W\,(F_0\cdot G)=F_0\cdot\left(\prod_{\mathrm{mo}}\nu(\mathrm{mo})^{-1}\right)\cdot\mathrm{tc}(G)\qquad\text{for }F_0\in\mathrm{Gr}[\phi],\ G\in\mathrm{Gr}[W],
\]
and the elements $F_0\cdot G$ span $\mathrm{Gr}[\phi,W]$.
Consequently $\mathrm{d}\mu_I(W)$ and $\mathrm{d}\mu_{\mathrm{Cov}}(W)$ do not depend on the orders either.
\item[(iv)] \emph{The empty set of modes} (the objects that occur in the assertion of (D2) on an empty $\mathrm{Mo}$).
Let $\mathrm{Mo}$ be empty.
Then there is no generator $W(u)$, and $\mathrm{Gr}[\phi,W]=\mathrm{Gr}[\phi]$.
The set $J_f$ is empty, it is the only subset $J$, and $W(J_f)=1$.
Both $\mathrm{Top}$ and $\prod_{\mathrm{mo}}\nu(\mathrm{mo})^{-1}$ are empty products, equal to $1$, and $\mathrm{sn}_f=1$.
Thus (15.6) becomes $\int\mathrm{D}W\,F=F$, that is, $\int\mathrm{D}W$ is the identity of $\mathrm{Gr}[\phi]$.
The space of the functions on $\mathrm{Mo}$ is $\{0\}$; its only linear map is the identity, which is invertible and has the determinant $1$ (the sum over the permutations that defines a determinant consists here of one term, an empty product).
Every form $\langle\mathrm{el}',\mathrm{el}\rangle_\nu$ is an empty sum, hence $0$, and $\exp0=1$.
It follows that $\int\mathrm{d}\mu_I(W)\,F=\int\mathrm{d}\mu_{\mathrm{Cov}}(W)\,F=F$ for all $F\in\mathrm{Gr}[\phi]$.
\item[(v)] \emph{Comparison with (I.3.5)} (the assertion of (D3) that the integral defined there is the integral of (I.3.5)).
The subalgebra of $\mathrm{Gr}[\phi,W]$ generated by the generators outside the family $W$ is $\mathrm{Gr}[\phi]$.
The degree of $\phi(I)\cdot W(J)$ in the family is $\#|J|$.
A term of lower degree in the family is a linear combination of products of generators in which fewer than $2\cdot\#|\mathrm{Mo}|$ factors are generators $W(u)$; this definition involves no order.
By 0.1\,(i) such a product is $0$ or $\pm$ an ordered monomial $\phi(I)\cdot W(J)$ with $J\ne J_f$, and conversely every $\phi(I)\cdot W(J)$ with $J\ne J_f$ is such a product.
The terms of lower degree are therefore exactly the linear combinations of the $\phi(I)\cdot W(J)$ with $J\ne J_f$.
For $F=\sum_{I,J}c_F(I,J)\cdot\phi(I)\cdot W(J)$ let
\[
F_0:=\mathrm{sn}_f\cdot\sum_Ic_F(I,J_f)\cdot\phi(I),
\]
an element of $\mathrm{Gr}[\phi]$.
Since $\mathrm{sn}_f\cdot\mathrm{Top}=W(J_f)$,
\[
F_0\cdot\mathrm{Top}=\sum_Ic_F(I,J_f)\cdot\phi(I)\cdot W(J_f),\qquad
F-F_0\cdot\mathrm{Top}=\sum_I\sum_{J\ne J_f}c_F(I,J)\cdot\phi(I)\cdot W(J).
\]
Hence $F=F_0\cdot\mathrm{Top}+(\text{a term of lower degree in the family})$, and $F_0$ is the only element of $\mathrm{Gr}[\phi]$ with this property.
In fact, the difference $G$ of two such elements lies in $\mathrm{Gr}[\phi]$, and $G\cdot\mathrm{Top}=\mathrm{sn}_f\cdot\sum_Ic_G(I)\cdot\phi(I)\cdot W(J_f)$ is a linear combination of the $\phi(I)\cdot W(J)$ with $J\ne J_f$; as the $\phi(I)\cdot W(J)$ form a basis, all $c_G(I)$ vanish.
So $F_0$ is the element $[F]_{\mathrm{Top}}$ of Section~I.3.6, and $\left(\prod_{\mathrm{mo}}\nu(\mathrm{mo})^{-1}\right)\cdot[F]_{\mathrm{Top}}$ is the right side of (15.6).
The element $\mathrm{Top}$ is even by 0.2\,(i); hence $\mathrm{Top}\cdot F_0=F_0\cdot\mathrm{Top}$ by 0.2\,(iii), and the coefficient on the left is the same element.
Now let $\mathrm{Mo}=\{(t,\beta)\}$ and $\nu(t,\beta):=\nu(t)$.
Then $W(\mathrm{mo})=W_\beta(t)$, $\bar{W}(\mathrm{mo})=\bar{W}_\beta(t)$ and $O(\mathrm{mo},\mathrm{mo}')=O(t,t')_{\beta\beta'}$ for $\mathrm{mo}=(t,\beta)$, $\mathrm{mo}'=(t',\beta')$.
With these identifications the objects $\mathrm{Top}$, $\langle\bar{W},W\rangle_\nu$, the kernel of a linear map $O$ and $\langle\bar{W},OW\rangle_\nu$ of Section~I.3.6 coincide, term by term, with those of (D2), (D3).
Therefore the map $\int\mathrm{D}W$ of (I.3.5) is the map $\int\mathrm{D}W$ of (D3), and $\mathrm{d}\mu_I(W)=\exp\left(-\langle\bar{W},W\rangle_\nu\right)\mathrm{D}W$ of (I.3.5) is the functional $F\mapsto\int\mathrm{D}W\,\exp\left(-\langle\bar{W},W\rangle_\nu\right)\cdot F$ of (D3); the exponential is even by (b3), so that by 0.2\,(iii) its position beside $F$ is immaterial.
Finally, Section~I.5.2 defines $\int\mathrm{d}\mu_{\mathrm{Cov}}(W)\,F$ and $\mathsf{w}_W$ by the formulas of (D3).
\end{itemize}

\begin{named}{Lemma 0.7}[$\int\mathrm{D}W$ under linear substitutions]
Let $O_1$, $O_2$ be linear maps of the functions on $\mathrm{Mo}$, invertible or not, and let $s_{O_1,O_2}$ be the homomorphism of $\mathrm{Gr}[\phi,W]$ to itself with
\[
W(\mathrm{mo})\mapsto(O_1W)(\mathrm{mo}),\qquad\bar{W}(\mathrm{mo})\mapsto(O_2\bar{W})(\mathrm{mo}),\qquad\phi_i(t)\mapsto\phi_i(t).
\]
Then
\begin{equation}
\int\mathrm{D}W\,s_{O_1,O_2}({F})=\det O_1\cdot\det O_2\cdot\int\mathrm{D}W\,F\qquad\text{for all }F\in\mathrm{Gr}[\phi,W].
\tag{15.7}
\end{equation}
\end{named}

\begin{proof}
For subsets $\mathrm{Mo}'$, $\mathrm{Mo}''$ of $\mathrm{Mo}$, $W(\mathrm{Mo}'\times\{0\})$ is the ordered product of the $W(\mathrm{mo})$, $\mathrm{mo}\in\mathrm{Mo}'$, and $W(\mathrm{Mo}''\times\{1\})$ is the ordered product of the $\bar{W}(\mathrm{mo})$, $\mathrm{mo}\in\mathrm{Mo}''$.
By 0.1\,(i), $W(J)=\pm W(\mathrm{Mo}'\times\{0\})\cdot W(\mathrm{Mo}''\times\{1\})$ with $\mathrm{Mo}'=\{\mathrm{mo}:(\mathrm{mo},0)\in J\}$ and $\mathrm{Mo}''=\{\mathrm{mo}:(\mathrm{mo},1)\in J\}$.
Both sides of (15.7) are linear in $F$, so it suffices to take $F=\phi(I)\cdot W(\mathrm{Mo}'\times\{0\})\cdot W(\mathrm{Mo}''\times\{1\})$.
For such $F$,
\[
s_{O_1,O_2}({F})=\phi(I)\cdot\prod_{\mathrm{mo}\in\mathrm{Mo}'}(O_1W)(\mathrm{mo})\cdot\prod_{\mathrm{mo}\in\mathrm{Mo}''}(O_2\bar{W})(\mathrm{mo}),
\]
with the factors in the order of the corresponding factors of $F$.
The first product is a linear combination of products of $\#|\mathrm{Mo}'|$ generators $W(\mathrm{mo}_1)$, the second a linear combination of products of $\#|\mathrm{Mo}''|$ generators $\bar{W}(\mathrm{mo}_1)$.
If $\mathrm{Mo}'\ne\mathrm{Mo}$, then every monomial of $s_{O_1,O_2}({F})$, and $F$ itself, contains fewer than $\#|\mathrm{Mo}|$ generators $W(\mathrm{mo})$, and both sides of (15.7) vanish by 0.6\,(i).
If $\mathrm{Mo}''\ne\mathrm{Mo}$, then every monomial of $s_{O_1,O_2}({F})$, and $F$ itself, contains fewer than $\#|\mathrm{Mo}|$ generators $\bar{W}(\mathrm{mo})$, and again both sides of (15.7) vanish by 0.6\,(i).
There remains the case $\mathrm{Mo}'=\mathrm{Mo}''=\mathrm{Mo}$.
Applying (15.2) with the generators $W(\mathrm{mo})$, in their order, in the place of the $\theta_j$ and with $\mathrm{Mx}=[O_1]$, rows and columns arranged in this order, one finds
\[
\prod_{\mathrm{mo}}(O_1W)(\mathrm{mo})=\det[O_1]\cdot W(\mathrm{Mo}\times\{0\});
\]
applying (15.2) with the generators $\bar{W}(\mathrm{mo})$, in their order, in the place of the $\theta_j$ and with $\mathrm{Mx}=[O_2]$, rows and columns arranged in that order, one finds
\[
\prod_{\mathrm{mo}}(O_2\bar{W})(\mathrm{mo})=\det[O_2]\cdot W(\mathrm{Mo}\times\{1\}).
\]
(A determinant is unchanged when its rows and its columns are permuted by one and the same permutation.)
Hence $s_{O_1,O_2}({F})=\det O_1\cdot\det O_2\cdot F$, by 0.5\,(i).
\end{proof}

\subsubsection*{(a) Kernels}

\begin{proof}[Proof of Proposition~15\,(a)]
Fix $n=(n_1,\dots,n_p)$.
A point of $\mathrm{Ind}_1^{n_1}\times\dots\times\mathrm{Ind}_p^{n_p}$ is written $t=(t^i_j)$, $1\le i\le p$, $1\le j\le n_i$, and one puts
\[
\nu(t):=\prod_{i,j}\nu_i(t^i_j),\qquad
\phi(t):=\phi_1(t^1_1)\cdots\phi_1(t^1_{n_1})\cdot\phi_2(t^2_1)\cdots\phi_p(t^p_{n_p}).
\]
Since $\int\mathrm{d}t$ is the sum over $t$ with the weight $\nu(t)$, the term of $F$ with the index $n$ is
\[
F_n=\left(\prod_i\frac{1}{n_i!}\right)\cdot\sum_tK_n(t)\cdot\nu(t)\cdot\phi(t),\qquad\text{and}\qquad
\int\mathrm{d}t\,|K_n|=\sum_t|K_n(t)|\cdot\nu(t).
\]
The point $t$ is called \emph{injective} if $t^i_j\ne t^i_{j'}$ for all $i$ and all $j\ne j'$.
By 0.1\,(i), $\phi(t)=0$ for every $t$ that is not injective.
If $n_i>\#|\mathrm{Ind}_i|$ for some $i$, no $t$ is injective, and $F_n=0$; consequently only finitely many $F_n$ are different from $0$, and $F=\sum_nF_n$ is well defined.

The injective points $t$ are in one-to-one correspondence with the pairs $(I,\mathrm{pm})$ in which $I=(I_1,\dots,I_p)$ satisfies $\#|I_i|=n_i$ and $\mathrm{pm}=(\mathrm{pm}_1,\dots,\mathrm{pm}_p)$ consists of permutations $\mathrm{pm}_i$ of $\{1,\dots,n_i\}$:
if $s^i_1<\dots<s^i_{n_i}$ are the elements of $I_i$, the point $t=:I^{\mathrm{pm}}$ has the coordinates $t^i_j=s^i_{\mathrm{pm}_i(j)}$.
Applying 0.1\,(i) to each of the $p$ groups of neighbouring factors one gets
\[
\phi(I^{\mathrm{pm}})=\mathrm{sgn}(\mathrm{pm})\cdot\phi(I),\qquad\mathrm{sgn}(\mathrm{pm}):=\prod_i\mathrm{sgn}(\mathrm{pm}_i),
\]
and moreover $\nu(I^{\mathrm{pm}})=\nu(I)$.
An ordered monomial $\phi(I)$ with $\#|I_i|=n_i$ for all $i$ occurs in no term other than $F_n$.
Hence, for such $I$,
\begin{equation}
c_F(I)=\left(\prod_i\frac{1}{n_i!}\right)\cdot\nu(I)\cdot\sum_{\mathrm{pm}}\mathrm{sgn}(\mathrm{pm})\cdot K_n(I^{\mathrm{pm}}).
\tag{15.8}
\end{equation}
The triangle inequality, together with the fact that the terms added in the last step are non-negative, yields
\begin{equation}
\begin{split}
\sum_{I:\,\#|I_i|=n_i}\left|c_F(I)\right|
&\le\left(\prod_i\frac{1}{n_i!}\right)\cdot\sum_{I,\mathrm{pm}}\left|K_n(I^{\mathrm{pm}})\right|\cdot\nu(I^{\mathrm{pm}})\\
&=\left(\prod_i\frac{1}{n_i!}\right)\cdot\sum_{t\text{ injective}}\left|K_n(t)\right|\cdot\nu(t)
\le\left(\prod_i\frac{1}{n_i!}\right)\cdot\int\mathrm{d}t\,|K_n|.
\end{split}
\tag{15.9}
\end{equation}
Multiplying by $\prod_i\mathsf{w}_i^{n_i}$ and summing over $n$ one arrives at
\[
\|F\|_{\mathsf{w}}\le\sum_n\prod_i\frac{\mathsf{w}_i^{n_i}}{n_i!}\int\mathrm{d}t\,|K_n|,
\]
where the right side is a sum of non-negative numbers and may be $+\infty$.

\emph{Antisymmetric kernels.}
Suppose now that every $K_n$ is antisymmetric in each group of arguments.
Then $K_n(I^{\mathrm{pm}})=\mathrm{sgn}(\mathrm{pm})\cdot K_n(I^{\mathrm{id}})$, $\mathrm{id}$ denoting the identity, and $K_n(t)=0$ for every $t$ that is not injective, since interchanging two equal arguments of one group gives $K_n(t)=-K_n(t)$.
The number $K_n(I^{\mathrm{id}})$ is the value of $K_n$ at the increasing arguments, and it is written $K_n(I_1;\dots;I_p)$.
All $\prod_in_i!$ terms of the sum in (15.8) are then equal to $K_n(I_1;\dots;I_p)$, and therefore
\[
c_F(I_1,\dots,I_p)=K_n(I_1;\dots;I_p)\cdot\prod_i\nu_i(I_i).
\]
So $|c_F(I)|=\left(\prod_i\frac{1}{n_i!}\right)\cdot\sum_{\mathrm{pm}}\left|K_n(I^{\mathrm{pm}})\right|\cdot\nu(I)$, which means that the first inequality in (15.9) is an equality.
The second inequality in (15.9) is an equality as well, since $K_n$ vanishes at all points $t$ that are not injective.
This proves the equality of the norms.

\emph{Determination.}
The last display shows that $F$ determines $K_n$ at the increasing arguments (because $\nu(I)>0$); by antisymmetry it determines $K_n$ at all injective arguments, and $K_n=0$ at the remaining points.
Conversely, let $F\in\mathrm{Gr}[\phi]$ be arbitrary, and define functions $K_n$ by these three rules, that is, $K_n(I^{\mathrm{pm}}):=\mathrm{sgn}(\mathrm{pm})\cdot\nu(I)^{-1}\cdot c_F(I)$ and $K_n:=0$ at all other points.
These functions are antisymmetric in each group, and by (15.8) they represent $F$.

\emph{A generator.}
For $F=\phi_i(t)$ the coefficient $c_F(I)$ equals $1$ if $I_i=\{t\}$ and $I_j=\emptyset$ for all $j\ne i$, and $0$ otherwise; hence $\|\phi_i(t)\|_{\mathsf{w}}=\mathsf{w}_i$.
The definition (I.5.1) is formed with the coefficients $c_F$ and the numbers $\mathsf{w}_i$ only, and contains none of the weights $\nu_i$.
\end{proof}

\subsubsection*{(b) Product; the even elements; the exponential}

\begin{proof}[Proof of Proposition~15\,(b)]
\textup{(b1)}
Let $F$ and $G$ lie in $\mathrm{Gr}[\phi]=\mathrm{Gr}[\phi_1,\dots,\phi_p]$.
Expanding both factors in the basis of ordered monomials, one has
\[
F\cdot G=\sum_{I,I'}c_F(I)\cdot c_G(I')\cdot\phi(I)\cdot\phi(I').
\]
According to 0.1\,(ii), the product $\phi(I)\cdot\phi(I')$ vanishes as soon as $I_i\cap I'_i\neq\emptyset$ for some $i$.
In the opposite case it equals $\pm\phi(I'')$, where $I''_i=I_i\cup I'_i$ for every $i$, and then $\mathsf{w}^{I''}=\mathsf{w}^{I}\cdot\mathsf{w}^{I'}$.
Consequently
\[
\left|c_{FG}(I'')\right|\le\sum\left|c_F(I)\right|\cdot\left|c_G(I')\right|,
\]
the sum extending over those pairs $(I,I')$ for which, for every $i$, the sets $I_i$ and $I'_i$ are disjoint and $I_i\cup I'_i=I''_i$;
moreover, a given pair $(I,I')$ contributes to at most one $I''$.
It follows that
\[
\|F\cdot G\|_{\mathsf{w}}\le\sum_{I,I'}\left|c_F(I)\right|\cdot\left|c_G(I')\right|\cdot\mathsf{w}^{I}\cdot\mathsf{w}^{I'}=\|F\|_{\mathsf{w}}\cdot\|G\|_{\mathsf{w}}.
\]

\textup{(b2)}
The space $\mathrm{Gr}[\phi]$ has finite dimension, and $\|\cdot\|_{\mathsf{w}}$ is a norm on it, namely the $\ell^1$ norm of the coefficients with respect to the positive weights $\mathsf{w}^I$.
Hence $\mathrm{Gr}[\phi]$ is complete, and $\|1\|_{\mathsf{w}}=1$.
The even elements form a linear subspace, which is therefore closed;
by 0.2\,(i) this subspace is a subalgebra containing $1$, and by 0.2\,(iii) it is commutative.
Together with (b1) one concludes that the even elements form a commutative Banach algebra with identity element $1$.

\textup{(b3)}
For even $F$ define
\[
\exp F:=\sum_{n\ge0}\frac{F^n}{n!}.
\]
By (b1), $\|F^n\|_{\mathsf{w}}\le\|F\|_{\mathsf{w}}^n$.
The series therefore converges absolutely, its sum is even, and
\[
\|\exp F-1\|_{\mathsf{w}}\le\sum_{n\ge1}\frac{\|F\|_{\mathsf{w}}^n}{n!}=\exp\left(\|F\|_{\mathsf{w}}\right)-1.
\]

\textup{(b4)}
The exponential has the following three properties.
Let $F$ and $G$ be even.
\begin{itemize}
\item[(i)] $\exp(F+G)=\exp F\cdot\exp G$.
Indeed, $F$ and $G$ commute; the binomial formula and the Cauchy product of two absolutely convergent series give the identity.
\item[(ii)] If $F\cdot F=0$, then $\exp F=1+F$.
\item[(iii)] If $\mathrm{hm}$ is a homomorphism as in 0.3, then $\mathrm{hm}(\exp F)=\exp\left(\mathrm{hm}({F})\right)$.
Indeed, $\mathrm{hm}$ is a linear map between spaces of finite dimension, hence continuous;
furthermore $\mathrm{hm}(F^n)=\mathrm{hm}({F})^n$, and $\mathrm{hm}({F})$ is even.
\end{itemize}
\end{proof}

\subsubsection*{(c) Row sums}

\begin{proof}[Proof of Proposition~15\,(c)]
Let $f$ be a function on $\mathrm{Ind}$ with $|f|_\infty\le1$.
For every $t\in\mathrm{Ind}'$,
\[
\left|(Of)(t)\right|=\left|\sum_sO(t,s)\cdot f(s)\cdot\nu(s)\right|\le\sum_s\left|O(t,s)\right|\cdot\nu(s)\cdot\left|f(s)\right|\le\sum_s\left|O(t,s)\right|\cdot\nu(s)\le|O|_{1,\infty}.
\]
Hence
\[
\sup\left\{|Of|_\infty:\ |f|_\infty\le1\right\}\le|O|_{1,\infty}.
\]
For the reverse inequality fix $t\in\mathrm{Ind}'$ and define
\[
f_t(s):=\frac{\overline{O(t,s)}}{\left|O(t,s)\right|}\quad\text{if }O(t,s)\neq0,\qquad f_t(s):=0\quad\text{otherwise}.
\]
Then $|f_t|_\infty\le1$ and $(Of_t)(t)=\sum_s\left|O(t,s)\right|\cdot\nu(s)$, so that
\[
\sup\left\{|Of|_\infty:\ |f|_\infty\le1\right\}\ge\left|(Of_t)(t)\right|=\sum_s\left|O(t,s)\right|\cdot\nu(s).
\]
Taking the supremum over $t$ gives the inequality $\ge|O|_{1,\infty}$.
\end{proof}

\begin{named}{Remark}
The number $\sum_s\left|O(t,s)\right|\cdot\nu(s)$ is the sum of the moduli of the coefficients $O(t,s)\cdot\nu(s)$ with which the values $f(s)$ enter $(Of)(t)$.
Thus $|O|_{1,\infty}$ is determined by the linear map $O$ alone;
the measure $\nu$ intervenes only through the way in which the map is represented by a kernel.
\end{named}

\subsubsection*{(d) Linear substitution}

\begin{proof}[Proof of Proposition~15\,(d)]
The image of a generator is
\[
\phi_i(t)'=\sum_j\sum_s\left(O_{ij}(t,s)\cdot\nu'_j(s)\right)\cdot\theta_j(s),
\]
a linear combination of generators of $\mathrm{Gr}[\theta]$.
It is odd, and by 0.3 the homomorphism $F\mapsto F'$ exists and is unique.
Distinct pairs $(j,s)$ give distinct ordered monomials $\theta_j(s)$;
the coefficients of $\phi_i(t)'$ in the basis of $\mathrm{Gr}[\theta]$ are therefore the numbers $O_{ij}(t,s)\cdot\nu'_j(s)$, standing at the monomials $\theta_j(s)$.
From (I.5.1) and (D4) one obtains
\[
\left\|\phi_i(t)'\right\|_{\mathsf{w}'}=\sum_j\mathsf{w}'_j\sum_s\left|O_{ij}(t,s)\right|\cdot\nu'_j(s)\le\sum_j\mathsf{w}'_j\cdot|O_{ij}|_{1,\infty}\le\sum_jC_{ij}\cdot\mathsf{w}'_j\le\mathsf{w}_i.
\]
The image of $\phi(I)$ is the product of the images of its $\#|I|$ factors, so that by (b1) its norm is at most $\mathsf{w}^I$.
Since $F'=\sum_Ic_F(I)\cdot\phi(I)'$, it follows that
\[
\|F'\|_{\mathsf{w}'}\le\sum_I\left|c_F(I)\right|\cdot\mathsf{w}^I=\|F\|_{\mathsf{w}}.
\]

\emph{(d$'$) Subalgebras and the rows that occur: the last assertion of (d).}
The index sets in (d) are arbitrary.
For every $i$ let $\mathrm{Ind}^0_i$ be a subset of $\mathrm{Ind}_i$, and let $F$ lie in the subalgebra generated by the $\phi_i(t)$, $t\in\mathrm{Ind}^0_i$.
This subalgebra is the Grassmann algebra on the index sets $\mathrm{Ind}^0_i$, and the coefficients and the norm of $F$ are the same in the two algebras.
Hence (d) holds for such $F$ when $|O_{ij}|_{1,\infty}$ is replaced by the supremum of the row sums over the rows $t\in\mathrm{Ind}^0_i$ only.
\end{proof}

\subsubsection*{(e) Moments; the determinant}

\begin{proof}[Proof of Proposition~15\,(e)]
\emph{Notation.}
Put $\mathrm{ww}(\mathrm{mo}):=W(\mathrm{mo})\cdot\bar{W}(\mathrm{mo})$.
For a subset $\mathrm{Mo}'$ of $\mathrm{Mo}$ let
\[
\mathrm{ww}(\mathrm{Mo}'):=\prod_{\mathrm{mo}\in\mathrm{Mo}'}\mathrm{ww}(\mathrm{mo}),\qquad\nu(\mathrm{Mo}'):=\prod_{\mathrm{mo}\in\mathrm{Mo}'}\nu(\mathrm{mo}),
\]
both products being equal to $1$ for $\mathrm{Mo}'=\emptyset$.
Further let
\[
\mathrm{Ew}:=\exp\left(-\langle\bar{W},W\rangle_\nu\right),\qquad\mathrm{Iw}({F}):=\int\mathrm{d}\mu_I(W)\,F=\int\mathrm{D}W\,\mathrm{Ew}\cdot F,
\]
and, for a subset $J$ of $J_f$,
\[
\mathrm{md}(J):=\left\{\mathrm{mo}\in\mathrm{Mo}:\ (\mathrm{mo},0)\in J\ \text{or}\ (\mathrm{mo},1)\in J\right\}.
\]
Every $\mathrm{ww}(\mathrm{mo})$ is even, and so commutes with every element (0.2\,(iii));
in particular the order of the factors in $\mathrm{ww}(\mathrm{Mo}')$ does not matter.
By 0.1\,(i), $\mathrm{ww}(\mathrm{mo})\cdot\mathrm{ww}(\mathrm{mo})=0$.
Finally, $\mathrm{Top}=\mathrm{ww}(\mathrm{Mo})$.

\emph{Step 1 (the element $\mathrm{Ew}$).}
By (15.1), $\bar{W}(\mathrm{mo})\cdot W(\mathrm{mo})=-\mathrm{ww}(\mathrm{mo})$.
Hence $-\langle\bar{W},W\rangle_\nu=\sum_{\mathrm{mo}}\nu(\mathrm{mo})\cdot\mathrm{ww}(\mathrm{mo})$, and (b4)\,(i), (ii) yield
\begin{equation}
\begin{split}
\mathrm{Ew}&=\prod_{\mathrm{mo}}\exp\left(\nu(\mathrm{mo})\cdot\mathrm{ww}(\mathrm{mo})\right)=\prod_{\mathrm{mo}}\left(1+\nu(\mathrm{mo})\cdot\mathrm{ww}(\mathrm{mo})\right)\\
&=\sum_{\mathrm{Mo}'\subset\mathrm{Mo}}\nu(\mathrm{Mo}')\cdot\mathrm{ww}(\mathrm{Mo}'),
\end{split}
\tag{15.10}
\end{equation}
where the sum extends over all subsets $\mathrm{Mo}'$ of $\mathrm{Mo}$.
In particular, $\mathrm{Ew}$ is even.

\emph{Step 2 (the products that reach $\mathrm{Top}$).}
Let $J$ be a subset of $J_f$ and $\mathrm{Mo}'$ a subset of $\mathrm{Mo}$.
If $\mathrm{Mo}'$ meets $\mathrm{md}(J)$, some generator occurs twice in $\mathrm{ww}(\mathrm{Mo}')\cdot W(J)$, and this product vanishes by 0.1\,(i).
If $\mathrm{Mo}'\cap\mathrm{md}(J)=\emptyset$, then $J\subset(\mathrm{Mo}\setminus\mathrm{Mo}')\times\{0,1\}$, and $\mathrm{ww}(\mathrm{Mo}')\cdot W(J)$ is $\pm$ the ordered monomial with the index set $(\mathrm{Mo}'\times\{0,1\})\cup J$;
this index set equals $J_f$ if and only if $J=(\mathrm{Mo}\setminus\mathrm{Mo}')\times\{0,1\}$.

\emph{Step 3 ($J$ not paired).}
Every set of the form $\mathrm{Mo}''\times\{0,1\}$ is paired.
Let $J$ be not paired.
By (15.10) and Step 2, no monomial of $\mathrm{Ew}\cdot W(J)$ contains all the generators $W(u)$, $u\in J_f$;
by 0.6\,(i) it follows that $m(J)=\int\mathrm{D}W\,\mathrm{Ew}\cdot W(J)=0$.

\emph{Step 4 ($J$ paired).}
In this case $J=\mathrm{Mo}''\times\{0,1\}$ with $\mathrm{Mo}'':=\mathrm{md}(J)=\left\{\mathrm{mo}:\ (\mathrm{mo},0)\in J\right\}$.
Each generator $W(u)$, $u\in J$, occurs exactly once in the product $\mathrm{ww}(\mathrm{Mo}'')$.
By 0.1\,(i),
\[
W(J)=\mathrm{sn}(J)\cdot\mathrm{ww}(\mathrm{Mo}'')=\mathrm{sn}(J)\cdot\prod_{(\mathrm{mo},0)\in J}W(\mathrm{mo})\cdot\bar{W}(\mathrm{mo}),\qquad\mathrm{sn}(J)\in\{1,-1\}.
\]
According to Step 2, among the terms of (15.10) only the term with $\mathrm{Mo}'=\mathrm{Mo}\setminus\mathrm{Mo}''$ contributes to $\int\mathrm{D}W\,\mathrm{Ew}\cdot W(J)$.
Since $\mathrm{ww}(\mathrm{Mo}\setminus\mathrm{Mo}'')\cdot\mathrm{ww}(\mathrm{Mo}'')=\mathrm{Top}$, 0.6\,(i) gives
\begin{equation}
\begin{split}
m(J)&=\mathrm{sn}(J)\cdot\nu(\mathrm{Mo}\setminus\mathrm{Mo}'')\cdot\int\mathrm{D}W\,\mathrm{Top}=\mathrm{sn}(J)\cdot\frac{\nu(\mathrm{Mo}\setminus\mathrm{Mo}'')}{\nu(\mathrm{Mo})}\\
&=\mathrm{sn}(J)\cdot\nu(\mathrm{Mo}'')^{-1}=\mathrm{sn}(J)\cdot\prod_{(\mathrm{mo},0)\in J}\nu(\mathrm{mo})^{-1}.
\end{split}
\tag{15.11}
\end{equation}
As $\mathrm{ww}(\mathrm{Mo}'')=\mathrm{sn}(J)\cdot W(J)$ and $\mathrm{sn}(J)\cdot\mathrm{sn}(J)=1$, this may be restated as
\begin{equation}
\mathrm{Iw}\left(\mathrm{ww}(\mathrm{Mo}'')\right)=\nu(\mathrm{Mo}'')^{-1}\qquad\text{for every subset }\mathrm{Mo}''\text{ of }\mathrm{Mo}.
\tag{15.12}
\end{equation}

\emph{Reality and modulus.}
The numbers $\nu(\mathrm{mo})$ are real, and so are all the $m(J)$.
Let $J$ be paired.
Each mode $\mathrm{mo}\in\mathrm{Mo}''$ contributes to $J$ the two elements $(\mathrm{mo},0)$ and $(\mathrm{mo},1)$, each with $\nu(u)^{-\frac{1}{2}}=\nu(\mathrm{mo})^{-\frac{1}{2}}$;
therefore
\[
\nu^{-\frac{1}{2}}(J)=\prod_{\mathrm{mo}\in\mathrm{Mo}''}\nu(\mathrm{mo})^{-1}=\left|m(J)\right|.
\]

\emph{What $m(J)$ depends on.}
The sign $\mathrm{sn}(J)$ is the sign of the permutation that takes the increasing arrangement of the elements of $J$ to an arrangement of the form $(\mathrm{mo}_1,0),(\mathrm{mo}_1,1),(\mathrm{mo}_2,0),(\mathrm{mo}_2,1),\dots$.
In view of (15.11), $m(J)$ is thus determined by the order of the elements of $J$ among themselves and by the weights of the modes of $J$;
the other modes of $\mathrm{Mo}$ play no role.

\emph{Step 5 (special cases).}
The empty set $J$ is paired, with $W(J)=1$ and $\mathrm{sn}(J)=1$;
so (15.11) gives $\int\mathrm{d}\mu_I(W)=1$.
Let $\mathrm{mo}\neq\mathrm{mo}'$.
None of the sets $\left\{(\mathrm{mo},0),(\mathrm{mo}',1)\right\}$, $\left\{(\mathrm{mo},0),(\mathrm{mo}',0)\right\}$, $\left\{(\mathrm{mo},1),(\mathrm{mo}',1)\right\}$ is paired, and the products $W(\mathrm{mo})\cdot\bar{W}(\mathrm{mo}')$, $W(\mathrm{mo})\cdot W(\mathrm{mo}')$, $\bar{W}(\mathrm{mo})\cdot\bar{W}(\mathrm{mo}')$ are $\pm$ the corresponding $W(J)$;
their integrals vanish by Step 3.
For $\mathrm{mo}=\mathrm{mo}'$ one has $W(\mathrm{mo})\cdot W(\mathrm{mo})=\bar{W}(\mathrm{mo})\cdot\bar{W}(\mathrm{mo})=0$, whereas (15.12) gives
\[
\mathrm{Iw}\left(W(\mathrm{mo})\cdot\bar{W}(\mathrm{mo})\right)=\mathrm{Iw}\left(\mathrm{ww}(\mathrm{mo})\right)=\frac{1}{\nu(\mathrm{mo})}.
\]

\emph{Step 6 (the determinant).}
Let $O$ be a linear map of the functions on $\mathrm{Mo}$, and let $\mathrm{hm}:=s_{O,1}$ be the homomorphism of Lemma~0.7 with $O_1=O$ and $O_2=1$.
By 0.5\,(iii), $\mathrm{hm}\left(\langle\bar{W},W\rangle_\nu\right)=\langle\bar{W},OW\rangle_\nu$;
by (b4)\,(iii), $\mathrm{hm}(\mathrm{Ew})=\exp\left(-\langle\bar{W},OW\rangle_\nu\right)$.
Now (15.7) and Step 5 give
\begin{equation*}
\begin{split}
\int\mathrm{D}W\,\exp\left(-\langle\bar{W},OW\rangle_\nu\right)&=\int\mathrm{D}W\,\mathrm{hm}(\mathrm{Ew})=\det O\cdot\det1\cdot\int\mathrm{D}W\,\mathrm{Ew}\\
&=\det O\cdot\int\mathrm{d}\mu_I(W)=\det O.
\end{split}
\end{equation*}
No invertibility of $O$ has been used.

The assertions (I.B1)--(I.B4) of Section~I.3.6 are among these statements:
(I.B1) is Step~6, Step~5 contains (I.B2) and (I.B3), and Steps~3 and~4 give (I.B4).
\end{proof}

\subsubsection*{(f) Integration}

\begin{proof}[Proof of Proposition~15\,(f)]
\emph{The formula for the integral.}
Let $F_0\in\mathrm{Gr}[\phi]$ and $G\in\mathrm{Gr}[\phi,W]$.
The even element $\mathrm{Ew}$ commutes with $F_0$ (0.2\,(iii)), and 0.6\,(ii) applies;
therefore
\[
\mathrm{Iw}(F_0\cdot G)=\int\mathrm{D}W\,\mathrm{Ew}\cdot F_0\cdot G=\int\mathrm{D}W\,F_0\cdot\mathrm{Ew}\cdot G=F_0\cdot\int\mathrm{D}W\,\mathrm{Ew}\cdot G=F_0\cdot\mathrm{Iw}({G}).
\]
Taking $F_0=\phi(I)$ and $G=W(J)$ and using linearity, one obtains for $F=\sum_{I,J}c_F(I,J)\cdot\phi(I)\cdot W(J)$
\begin{equation}
\int\mathrm{d}\mu_I(W)\,F=\sum_I\left(\sum_Jc_F(I,J)\cdot m(J)\right)\cdot\phi(I).
\tag{15.13}
\end{equation}

\emph{The two inequalities.}
By (e), $\left|m(J)\right|=\nu^{-\frac{1}{2}}(J)$ for paired $J$, and $m(J)=0$ otherwise.
Moreover $\nu(u)^{-\frac{1}{2}}\le\mathsf{w}_W$ for every $u$, so that $\nu^{-\frac{1}{2}}(J)\le\mathsf{w}_W^{\#|J|}$.
Consequently
\begin{equation}
\begin{split}
\left\|\int\mathrm{d}\mu_I(W)\,F\right\|_{(\mathsf{w}_1,\dots,\mathsf{w}_p)}&=\sum_I\left|\sum_Jc_F(I,J)\cdot m(J)\right|\cdot\mathsf{w}^I\le\sum_{I,J}\left|c_F(I,J)\right|\cdot\mathsf{w}^I\cdot\nu^{-\frac{1}{2}}(J)\\
&\le\sum_{I,J}\left|c_F(I,J)\right|\cdot\mathsf{w}^I\cdot\mathsf{w}_W^{\#|J|}=\|F\|_{(\mathsf{w}_1,\dots,\mathsf{w}_p,\mathsf{w}_W)}.
\end{split}
\tag{15.14}
\end{equation}
The outer members of (15.14) form the inequality (I.5.2).

\emph{The weight cannot be lowered.}
Let $\mathrm{mo}$ be a mode with $\nu(\mathrm{mo})^{-\frac{1}{2}}=\mathsf{w}_W$, and put $J:=\left\{(\mathrm{mo},0),(\mathrm{mo},1)\right\}$ and $F:=W(\mathrm{mo})\cdot\bar{W}(\mathrm{mo})=\mathrm{sn}(J)\cdot W(J)$.
The only non-zero coefficient of $F$ is $c_F(I,J)=\mathrm{sn}(J)$, with $I$ empty.
The three members of the chain of inequalities in (f) are
\[
\left|\int\mathrm{d}\mu_I(W)\,F\right|=\frac{1}{\nu(\mathrm{mo})}=\mathsf{w}_W^2\quad\text{by (15.12)},\qquad\nu^{-\frac{1}{2}}(J)=\nu(\mathrm{mo})^{-1}=\mathsf{w}_W^2,\qquad\|F\|_{(\mathsf{w},\mathsf{w}_W)}=\mathsf{w}_W^2.
\]
If a weight $\mathsf{w}'_W<\mathsf{w}_W$ is put in the place of $\mathsf{w}_W$, the right-hand side becomes $(\mathsf{w}'_W)^2<\mathsf{w}_W^2=\left|\int\mathrm{d}\mu_I(W)\,F\right|$.

\emph{(f$'$) The modes that do not occur: the last assertion of (f).}
Let $\mathrm{Mo}^0$ be a non-empty subset of $\mathrm{Mo}$ with the weights $\nu(\mathrm{mo})$, $\mathrm{mo}\in\mathrm{Mo}^0$, and let $\mathrm{Mo}^0\times\{0,1\}$ carry the order induced by the order of $\mathrm{Mo}\times\{0,1\}$.
By 0.1\,(ii), the linear span of the ordered monomials $\phi(I)\cdot W(J)$ with $J\subset\mathrm{Mo}^0\times\{0,1\}$ is a subalgebra of $\mathrm{Gr}[\phi,W]$.
It contains the $\phi_i(t)$ and the $W(u)$, $u\in\mathrm{Mo}^0\times\{0,1\}$, and the monomials in question are products of these elements;
hence it is the subalgebra generated by them.
It is the Grassmann algebra of (D2) formed with $\mathrm{Mo}^0$ in the place of $\mathrm{Mo}$, and the coefficients $c_F(I,J)$ of its elements are the same in both algebras.
Let $J\subset\mathrm{Mo}^0\times\{0,1\}$.
The number $m(J)$ formed with $\mathrm{Mo}^0$ coincides with the number $m(J)$ formed with $\mathrm{Mo}$:
by (e) both vanish if $J$ is not paired, and otherwise both equal $\mathrm{sn}(J)\cdot\prod_{(\mathrm{mo},0)\in J}\nu(\mathrm{mo})^{-1}$ with the same sign $\mathrm{sn}(J)$, because $W(J)$ and the elements $\mathrm{ww}(\mathrm{mo})$, $\mathrm{mo}\in\mathrm{md}(J)$, are the same in both algebras.
By (15.13) it follows that, for every $F$ in the subalgebra, $\int\mathrm{d}\mu_I(W)\,F$ is one and the same element of $\mathrm{Gr}[\phi]$, whether the integral is formed with $\mathrm{Mo}$ or with $\mathrm{Mo}^0$.

Let now $\mathrm{Mo}^0$ be empty.
The only subset $J$ of $\mathrm{Mo}^0\times\{0,1\}$ is the empty set, and the subalgebra is $\mathrm{Gr}[\phi]$.
Formed with $\mathrm{Mo}$:
$m(J)=1$ for the empty set $J$ ((e), Step~5), and for $F\in\mathrm{Gr}[\phi]$ only the coefficients $c_F(I,J)$ with empty $J$ can differ from $0$, so that (15.13) reads $\int\mathrm{d}\mu_I(W)\,F=F$.
Formed with the empty set of modes:
$\int\mathrm{d}\mu_I(W)\,F=F$ by 0.6\,(iv), and in particular $m(J)=1$.
\end{proof}

\subsubsection*{(g) Factorisation}

\begin{proof}[Proof of Proposition~15\,(g)]
Both sides are bilinear in $(F_1,F_2)$.
The hypothesis states that $\mathrm{md}(J)$ and $\mathrm{md}(J')$ are disjoint whenever $\phi(I)\cdot W(J)$ is a monomial of $F_1$ and $\phi(I')\cdot W(J')$ is a monomial of $F_2$.
It therefore suffices to treat $F_1=\phi(I)\cdot W(J)$ and $F_2=\phi(I')\cdot W(J')$ with $\mathrm{md}(J)\cap\mathrm{md}(J')=\emptyset$.
By (15.13),
\[
\mathrm{Iw}(F_1)\cdot\mathrm{Iw}(F_2)=m(J)\cdot m(J')\cdot\phi(I)\cdot\phi(I').
\]
By 0.2\,(ii), $W(J)\cdot\phi(I')=(-1)^{\#|J|\cdot\#|I'|}\,\phi(I')\cdot W(J)$.
Hence $F_1\cdot F_2=(-1)^{\#|J|\cdot\#|I'|}\,\phi(I)\cdot\phi(I')\cdot W(J)\cdot W(J')$, and the identity $\mathrm{Iw}(F_0\cdot G)=F_0\cdot\mathrm{Iw}({G})$ from the beginning of the proof of (f) gives
\[
\mathrm{Iw}(F_1\cdot F_2)=(-1)^{\#|J|\cdot\#|I'|}\,\phi(I)\cdot\phi(I')\cdot\mathrm{Iw}\left(W(J)\cdot W(J')\right).
\]
The sets $J$ and $J'$ are disjoint, and $W(J)\cdot W(J')=\pm W(J\cup J')$ by 0.1\,(ii).
The set $J\cup J'$ is paired if and only if $J$ and $J'$ are both paired.
Indeed, for $\mathrm{mo}\in\mathrm{md}(J)$ an element $(\mathrm{mo},\omega)$ lies in $J\cup J'$ exactly when it lies in $J$, because $\mathrm{mo}\notin\mathrm{md}(J')$;
for $\mathrm{mo}\in\mathrm{md}(J')$ it lies in $J\cup J'$ exactly when it lies in $J'$, because $\mathrm{mo}\notin\mathrm{md}(J)$;
and a mode outside $\mathrm{md}(J)\cup\mathrm{md}(J')$ occurs in none of the three sets.

\emph{Case 1: $J$ or $J'$ is not paired.}
Then $\mathrm{Iw}\left(W(J)\cdot W(J')\right)=\pm m(J\cup J')=0$, and $m(J)\cdot m(J')=0$ by (e).
Both sides vanish.

\emph{Case 2: $J$ and $J'$ are paired.}
Then $\#|J|$ is even, and the sign $(-1)^{\#|J|\cdot\#|I'|}$ is $1$.
By (e), Step~4,
\[
W(J)\cdot W(J')=\mathrm{sn}(J)\cdot\mathrm{sn}(J')\cdot\mathrm{ww}\left(\mathrm{md}(J)\right)\cdot\mathrm{ww}\left(\mathrm{md}(J')\right)=\mathrm{sn}(J)\cdot\mathrm{sn}(J')\cdot\mathrm{ww}\left(\mathrm{md}(J)\cup\mathrm{md}(J')\right).
\]
By (15.12) and the disjointness,
\[
\mathrm{Iw}\left(\mathrm{ww}\left(\mathrm{md}(J)\cup\mathrm{md}(J')\right)\right)=\nu\left(\mathrm{md}(J)\right)^{-1}\cdot\nu\left(\mathrm{md}(J')\right)^{-1}.
\]
Together with (15.11) this gives $\mathrm{Iw}\left(W(J)\cdot W(J')\right)=m(J)\cdot m(J')$.
Therefore
\[
\mathrm{Iw}(F_1\cdot F_2)=m(J)\cdot m(J')\cdot\phi(I)\cdot\phi(I')=\mathrm{Iw}(F_1)\cdot\mathrm{Iw}(F_2).
\]

\emph{The last assertion of (g): disjoint sets of points in the scheme.}
Let $\mathrm{mo}=(t,\beta)$.
If a mode $(t,\beta)$ occurs in a monomial of $F_1$ and in a monomial of $F_2$, the point $t$ belongs to both sets of points.
Disjoint sets of points therefore imply the hypothesis.
\end{proof}

\subsubsection*{(h) Invariance}

\begin{proof}[Proof of Proposition~15\,(h)]
In the notation of Lemma~0.7 one has $s_U=s_{O_1,O_2}$ with $O_1=U$ and $O_2=(U^{-1})^{T}$.

\emph{The form.}
Apply (15.5) with $O=U^{-1}$, $\mathrm{el}'=\bar{W}$ and $\mathrm{el}=UW$, and then 0.5\,(iii):
\begin{equation}
\left\langle(U^{-1})^{T}\bar{W},UW\right\rangle_\nu=\left\langle\bar{W},U^{-1}(UW)\right\rangle_\nu=\left\langle\bar{W},(U^{-1}U)W\right\rangle_\nu=\langle\bar{W},W\rangle_\nu.
\tag{15.15}
\end{equation}

\emph{The map $\int\mathrm{D}W$.}
By (15.7) and 0.5\,(ii),
\[
\int\mathrm{D}W\,s_U({F})=\det U\cdot\det(U^{-1})^{T}\cdot\int\mathrm{D}W\,F=\det U\cdot(\det U)^{-1}\cdot\int\mathrm{D}W\,F=\int\mathrm{D}W\,F.
\]

\emph{The functional $\mathrm{d}\mu_I$.}
By 0.5\,(iii), $s_U\left(\langle\bar{W},W\rangle_\nu\right)$ is the left-hand side of (15.15);
hence $s_U\left(\langle\bar{W},W\rangle_\nu\right)=\langle\bar{W},W\rangle_\nu$, and $s_U(\mathrm{Ew})=\mathrm{Ew}$ by (b4)\,(iii).
Consequently
\begin{equation*}
\begin{split}
\int\mathrm{d}\mu_I(W)\,s_U({F})&=\int\mathrm{D}W\,\mathrm{Ew}\cdot s_U({F})=\int\mathrm{D}W\,s_U(\mathrm{Ew})\cdot s_U({F})=\int\mathrm{D}W\,s_U(\mathrm{Ew}\cdot F)\\
&=\int\mathrm{D}W\,\mathrm{Ew}\cdot F=\int\mathrm{d}\mu_I(W)\,F.
\end{split}
\end{equation*}
\end{proof}

\subsubsection*{(i) Wick's theorem}

\begin{proof}[Proof of Proposition~15\,(i)]
By the identity $\mathrm{Iw}(F_0\cdot G)=F_0\cdot\mathrm{Iw}({G})$ from the beginning of the proof of (f),
\[
\int\mathrm{d}\mu_I(W)\,F_0\cdot v_1\cdots v_r=F_0\cdot\mathrm{Iw}(v_1\cdots v_r).
\]
The assertion is thus a statement about the numbers $\mathrm{Iw}(v_1\cdots v_r)$.
Write
\[
v_j=\sum_{u\in J_f}\mathrm{kf}_j(u)\cdot W(u)
\]
with complex numbers $\mathrm{kf}_j(u)$.
The $v_j$ are odd;
by 0.2\,(iii), $v_i\cdot v_j=-v_j\cdot v_i$ and $v_i\cdot v_i=0$.
Hence the matrix with the entries $\mathrm{Mx}_{ij}:=\mathrm{Iw}(v_i\cdot v_j)$, $1\le i,j\le r$, is antisymmetric;
it is the matrix of (i).

\emph{Step 1 (the Pfaffian and its expansion along the first row).}
For an antisymmetric complex $r\times r$ matrix $\mathrm{Mx}$ with $r=2n\ge2$ let
\begin{equation}
\mathrm{Pfaff}(\mathrm{Mx}):=\sum\mathrm{sgn}(i_1,j_1,i_2,j_2,\dots,i_n,j_n)\cdot\mathrm{Mx}_{i_1j_1}\cdot\mathrm{Mx}_{i_2j_2}\cdots\mathrm{Mx}_{i_nj_n},
\tag{15.16}
\end{equation}
where the sum is over the partitions of $\{1,\dots,r\}$ into $n$ pairs, each partition being written exactly once, in the form $\{i_1,j_1\},\dots,\{i_n,j_n\}$ with $i_{\mathrm{ia}}<j_{\mathrm{ia}}$ for $\mathrm{ia}=1,\dots,n$ and $i_1<i_2<\dots<i_n$;
here $\mathrm{sgn}(\cdot)$ is $-1$ raised to the number of inversions of the sequence, an inversion being a pair of places $\mathrm{ia}<\mathrm{ib}$ such that the entry at the place $\mathrm{ia}$ is larger than the entry at the place $\mathrm{ib}$.
For $r=0$ let $\mathrm{Pfaff}:=1$.
In every term $i_1=1$.
Fix $j:=j_1\in\{2,\dots,r\}$.
The other pairs $\{i_2,j_2\},\dots,\{i_n,j_n\}$ then range over the partitions of $\{2,\dots,r\}\setminus\{j\}$ into $n-1$ pairs, each partition written once, with $i_{\mathrm{ia}}<j_{\mathrm{ia}}$ for $\mathrm{ia}=2,\dots,n$ and $i_2<\dots<i_n$.
In the sequence $(1,j,i_2,j_2,\dots)$ the entry $1$ belongs to no inversion;
the entry $j$ belongs exactly to the inversions with the $j-2$ numbers $2,\dots,j-1$, all of which stand to its right;
the remaining inversions are those of $(i_2,j_2,\dots,i_n,j_n)$.
Thus $\mathrm{sgn}(1,j,i_2,j_2,\dots)=(-1)^{j-2}\,\mathrm{sgn}(i_2,j_2,\dots,i_n,j_n)$, and
\begin{equation}
\mathrm{Pfaff}(\mathrm{Mx})=\sum_{j=2}^{r}(-1)^j\,\mathrm{Mx}_{1j}\cdot\mathrm{Pfaff}\left(\mathrm{Mx}^{(1j)}\right),
\tag{15.17}
\end{equation}
where $\mathrm{Mx}^{(1j)}$ is the $(r-2)\times(r-2)$ matrix that remains after deletion of the rows and the columns $1$ and $j$, the others keeping their order.
(The increasing bijection of $\{2,\dots,r\}\setminus\{j\}$ onto $\{1,\dots,r-2\}$ maps the partitions, written as above, onto the partitions of $\{1,\dots,r-2\}$ written as in (15.16), and it preserves the inversions of the sequences.)

\emph{Step 2 (the map $\mathrm{der}_u$).}
Fix $u\in J_f$.
Let $\mathrm{Gr}^{(u)}$ be the linear span of the ordered monomials $\phi(I)\cdot W(J)$ with $u\notin J$.
By 0.1\,(ii) it is a subalgebra, and it contains all generators other than $W(u)$.
By 0.1\,(i), an ordered monomial $\phi(I)\cdot W(J)$ with $u\in J$ equals $\pm W(u)\cdot\phi(I)\cdot W(J\setminus\{u\})$.
Hence the elements $\phi(I)\cdot W(J')$ and $W(u)\cdot\phi(I)\cdot W(J')$, with $u\notin J'$, form a basis of $\mathrm{Gr}[\phi,W]$, and every $G\in\mathrm{Gr}[\phi,W]$ can be written in exactly one way as
\[
G=G'+W(u)\cdot G''\quad\text{with }G',G''\in\mathrm{Gr}^{(u)};\qquad\mathrm{der}_uG:=G''.
\]
The map $\mathrm{der}_u$ is linear, and $\mathrm{der}_u1=0$.
Since $\mathrm{der}_uG$ lies in $\mathrm{Gr}^{(u)}$, none of its monomials contains $W(u)$, and 0.6\,(i) gives
\begin{equation}
\int\mathrm{D}W\,\mathrm{der}_uG=0\qquad\text{for all }G\in\mathrm{Gr}[\phi,W].
\tag{15.18}
\end{equation}
Let $v=\sum_{u'}\mathrm{kf}(u')\cdot W(u')$.
Then
\begin{equation}
\mathrm{der}_u(v\cdot G)=\mathrm{kf}(u)\cdot G-v\cdot\mathrm{der}_uG\qquad\text{for all }G\in\mathrm{Gr}[\phi,W].
\tag{15.19}
\end{equation}
For the proof put $v':=v-\mathrm{kf}(u)\cdot W(u)$, an odd element of $\mathrm{Gr}^{(u)}$, and write $G=G'+W(u)\cdot G''$.
Since $W(u)\cdot W(u)=0$ and $v'\cdot W(u)=-W(u)\cdot v'$,
\[
v\cdot G=\mathrm{kf}(u)\cdot W(u)\cdot G'+v'\cdot G'+v'\cdot W(u)\cdot G''=v'\cdot G'+W(u)\cdot\left(\mathrm{kf}(u)\cdot G'-v'\cdot G''\right),
\]
and $v'\cdot G'$ and $\mathrm{kf}(u)\cdot G'-v'\cdot G''$ lie in $\mathrm{Gr}^{(u)}$.
So $\mathrm{der}_u(v\cdot G)=\mathrm{kf}(u)\cdot G'-v'\cdot G''$.
On the other hand,
\[
\mathrm{kf}(u)\cdot G-v\cdot G''=\mathrm{kf}(u)\cdot G'+\mathrm{kf}(u)\cdot W(u)\cdot G''-\mathrm{kf}(u)\cdot W(u)\cdot G''-v'\cdot G''=\mathrm{kf}(u)\cdot G'-v'\cdot G''.
\]
For $r\ge1$, induction on the number of factors, starting from (15.19) and $\mathrm{der}_u1=0$, gives
\begin{equation}
\mathrm{der}_u(v_2\cdot v_3\cdots v_r)=\sum_{j=2}^{r}(-1)^j\,\mathrm{kf}_j(u)\cdot v_2\cdots\widehat{v_j}\cdots v_r
\tag{15.20}
\end{equation}
(for $r=1$ both sides are $0$), where $\widehat{v_j}$ indicates that the factor $v_j$ is omitted.
Indeed, suppose that (15.20) holds for the $r-2$ factors $v_3,\dots,v_r$, that is,
\[
\mathrm{der}_u(v_3\cdots v_r)=\sum_{j=3}^{r}(-1)^{j-1}\,\mathrm{kf}_j(u)\cdot v_3\cdots\widehat{v_j}\cdots v_r.
\]
Then (15.19) gives
\[
\mathrm{der}_u\left(v_2\cdot(v_3\cdots v_r)\right)=\mathrm{kf}_2(u)\cdot v_3\cdots v_r-v_2\cdot\sum_{j=3}^{r}(-1)^{j-1}\,\mathrm{kf}_j(u)\cdot v_3\cdots\widehat{v_j}\cdots v_r,
\]
which is the right-hand side of (15.20).

\emph{Step 3 ($\mathrm{der}_u$ and $\mathrm{Ew}$).}
For every mode $\mathrm{mo}$ and every $G\in\mathrm{Gr}[\phi,W]$,
\begin{equation}
\begin{split}
\mathrm{der}_{(\mathrm{mo},1)}(\mathrm{Ew}\cdot G)&=\mathrm{Ew}\cdot\mathrm{der}_{(\mathrm{mo},1)}G-\nu(\mathrm{mo})\cdot W(\mathrm{mo})\cdot\mathrm{Ew}\cdot G,\\
\mathrm{der}_{(\mathrm{mo},0)}(\mathrm{Ew}\cdot G)&=\mathrm{Ew}\cdot\mathrm{der}_{(\mathrm{mo},0)}G+\nu(\mathrm{mo})\cdot\bar{W}(\mathrm{mo})\cdot\mathrm{Ew}\cdot G.
\end{split}
\tag{15.21}
\end{equation}
To prove this, note that by (15.10)
\[
\mathrm{Ew}=\left(1+\nu(\mathrm{mo})\cdot W(\mathrm{mo})\cdot\bar{W}(\mathrm{mo})\right)\cdot\mathrm{Ew}',\qquad\mathrm{Ew}':=\prod_{\mathrm{mo}'\neq\mathrm{mo}}\left(1+\nu(\mathrm{mo}')\cdot\mathrm{ww}(\mathrm{mo}')\right),
\]
where $\mathrm{Ew}'$ is even and contains neither $W(\mathrm{mo})$ nor $\bar{W}(\mathrm{mo})$.

For the first formula let $u=(\mathrm{mo},1)$, so that $W(u)=\bar{W}(\mathrm{mo})$, and write $G=G'+\bar{W}(\mathrm{mo})\cdot G''$.
Since $\bar{W}(\mathrm{mo})\cdot\bar{W}(\mathrm{mo})=0$ and $W(\mathrm{mo})\cdot\bar{W}(\mathrm{mo})=-\bar{W}(\mathrm{mo})\cdot W(\mathrm{mo})$,
\begin{equation*}
\begin{split}
\mathrm{Ew}\cdot G&=\mathrm{Ew}'\cdot\left(G'+\bar{W}(\mathrm{mo})\cdot G''+\nu(\mathrm{mo})\cdot W(\mathrm{mo})\cdot\bar{W}(\mathrm{mo})\cdot G'\right)\\
&=\mathrm{Ew}'\cdot G'+\bar{W}(\mathrm{mo})\cdot\mathrm{Ew}'\cdot\left(G''-\nu(\mathrm{mo})\cdot W(\mathrm{mo})\cdot G'\right).
\end{split}
\end{equation*}
The elements $\mathrm{Ew}'\cdot G'$ and $\mathrm{Ew}'\cdot\left(G''-\nu(\mathrm{mo})\cdot W(\mathrm{mo})\cdot G'\right)$ lie in $\mathrm{Gr}^{(u)}$.
Hence $\mathrm{der}_u(\mathrm{Ew}\cdot G)=\mathrm{Ew}'\cdot\left(G''-\nu(\mathrm{mo})\cdot W(\mathrm{mo})\cdot G'\right)$.
As $W(\mathrm{mo})$ commutes with $\mathrm{Ew}'$ and $W(\mathrm{mo})\cdot W(\mathrm{mo})=0$, the right-hand side of the first formula equals
\begin{equation*}
\begin{split}
&\mathrm{Ew}'\cdot\left(1+\nu(\mathrm{mo})\cdot W(\mathrm{mo})\cdot\bar{W}(\mathrm{mo})\right)\cdot G''-\nu(\mathrm{mo})\cdot\mathrm{Ew}'\cdot W(\mathrm{mo})\cdot\left(G'+\bar{W}(\mathrm{mo})\cdot G''\right)\\
&\qquad=\mathrm{Ew}'\cdot\left(G''-\nu(\mathrm{mo})\cdot W(\mathrm{mo})\cdot G'\right).
\end{split}
\end{equation*}

For the second formula let $u=(\mathrm{mo},0)$, so that $W(u)=W(\mathrm{mo})$, and write $G=G'+W(\mathrm{mo})\cdot G''$.
Since $W(\mathrm{mo})\cdot\bar{W}(\mathrm{mo})\cdot W(\mathrm{mo})=-W(\mathrm{mo})\cdot W(\mathrm{mo})\cdot\bar{W}(\mathrm{mo})=0$,
\begin{equation*}
\begin{split}
\mathrm{Ew}\cdot G&=\mathrm{Ew}'\cdot\left(G'+W(\mathrm{mo})\cdot G''+\nu(\mathrm{mo})\cdot W(\mathrm{mo})\cdot\bar{W}(\mathrm{mo})\cdot G'\right)\\
&=\mathrm{Ew}'\cdot G'+W(\mathrm{mo})\cdot\mathrm{Ew}'\cdot\left(G''+\nu(\mathrm{mo})\cdot\bar{W}(\mathrm{mo})\cdot G'\right),
\end{split}
\end{equation*}
so that $\mathrm{der}_u(\mathrm{Ew}\cdot G)=\mathrm{Ew}'\cdot\left(G''+\nu(\mathrm{mo})\cdot\bar{W}(\mathrm{mo})\cdot G'\right)$.
As $\bar{W}(\mathrm{mo})\cdot W(\mathrm{mo})\cdot\bar{W}(\mathrm{mo})=0$ and $W(\mathrm{mo})\cdot\bar{W}(\mathrm{mo})+\bar{W}(\mathrm{mo})\cdot W(\mathrm{mo})=0$, the right-hand side of the second formula equals
\begin{equation*}
\begin{split}
&\mathrm{Ew}'\cdot\left(1+\nu(\mathrm{mo})\cdot W(\mathrm{mo})\cdot\bar{W}(\mathrm{mo})\right)\cdot G''+\nu(\mathrm{mo})\cdot\mathrm{Ew}'\cdot\bar{W}(\mathrm{mo})\cdot\left(G'+W(\mathrm{mo})\cdot G''\right)\\
&\qquad=\mathrm{Ew}'\cdot\left(G''+\nu(\mathrm{mo})\cdot\bar{W}(\mathrm{mo})\cdot G'\right).
\end{split}
\end{equation*}

\emph{Step 4 (integration by parts).}
Apply $\int\mathrm{D}W$ to the two formulas (15.21).
The left-hand sides give $0$ by (15.18), and $\mathrm{Ew}$ commutes with $W(\mathrm{mo})$ and $\bar{W}(\mathrm{mo})$.
Hence, for all $G\in\mathrm{Gr}[\phi,W]$,
\begin{equation}
\begin{split}
\mathrm{Iw}\left(W(\mathrm{mo})\cdot G\right)&=\nu(\mathrm{mo})^{-1}\cdot\mathrm{Iw}\left(\mathrm{der}_{(\mathrm{mo},1)}G\right),\\
\mathrm{Iw}\left(\bar{W}(\mathrm{mo})\cdot G\right)&=-\nu(\mathrm{mo})^{-1}\cdot\mathrm{Iw}\left(\mathrm{der}_{(\mathrm{mo},0)}G\right).
\end{split}
\tag{15.22}
\end{equation}

\emph{Step 5 (the recursion).}
For $r\ge1$ and $v_1,\dots,v_r$ in the linear span of the $W(u)$,
\begin{equation}
\mathrm{Iw}(v_1\cdot v_2\cdots v_r)=\sum_{j=2}^{r}(-1)^j\,\mathrm{Iw}(v_1\cdot v_j)\cdot\mathrm{Iw}\left(v_2\cdots\widehat{v_j}\cdots v_r\right)
\tag{15.23}
\end{equation}
(the sum is $0$ for $r=1$; an empty product is $1$).
Both sides are linear in $v_1$;
it suffices therefore to take $v_1=W(\mathrm{mo})$ or $v_1=\bar{W}(\mathrm{mo})$.
For $v_1=W(\mathrm{mo})$, the formulas (15.22) and (15.20) with $u=(\mathrm{mo},1)$ give
\[
\mathrm{Iw}\left(W(\mathrm{mo})\cdot v_2\cdots v_r\right)=\sum_{j=2}^{r}(-1)^j\,\nu(\mathrm{mo})^{-1}\cdot\mathrm{kf}_j(\mathrm{mo},1)\cdot\mathrm{Iw}\left(v_2\cdots\widehat{v_j}\cdots v_r\right).
\]
For the single factor $v_j$, the formulas (15.22) and (15.19) with $G=1$, together with $\mathrm{der}_u1=0$ and $\mathrm{Iw}(1)=1$ (see (e)), give
\[
\mathrm{Iw}\left(W(\mathrm{mo})\cdot v_j\right)=\nu(\mathrm{mo})^{-1}\cdot\mathrm{Iw}\left(\mathrm{der}_{(\mathrm{mo},1)}v_j\right)=\nu(\mathrm{mo})^{-1}\cdot\mathrm{kf}_j(\mathrm{mo},1).
\]
The last two formulas together are (15.23).
For $v_1=\bar{W}(\mathrm{mo})$, the second formula of (15.22), the formula (15.20) and the formula (15.19) with $G=1$, all with $u=(\mathrm{mo},0)$, give
\[
\mathrm{Iw}\left(\bar{W}(\mathrm{mo})\cdot v_2\cdots v_r\right)=-\sum_{j=2}^{r}(-1)^j\,\nu(\mathrm{mo})^{-1}\cdot\mathrm{kf}_j(\mathrm{mo},0)\cdot\mathrm{Iw}\left(v_2\cdots\widehat{v_j}\cdots v_r\right)
\]
and
\[
\mathrm{Iw}\left(\bar{W}(\mathrm{mo})\cdot v_j\right)=-\nu(\mathrm{mo})^{-1}\cdot\mathrm{kf}_j(\mathrm{mo},0);
\]
together they are (15.23).

\emph{Step 6 (induction on $r$).}
The claim is that $\mathrm{Iw}(v_1\cdots v_r)=0$ for odd $r$ and $\mathrm{Iw}(v_1\cdots v_r)=\mathrm{Pfaff}(\mathrm{Mx})$ for even $r$.
For $r=0$: $\mathrm{Iw}(1)=1=\mathrm{Pfaff}$.
For $r=1$, (15.23) gives $0$.
Let $r\ge2$, and let the claim hold for $r-2$ factors.
The $r-2$ elements $v_2,\dots,v_r$ other than $v_j$, taken in this order, have the matrix of pair integrals $\mathrm{Mx}^{(1j)}$.
If $r$ is odd, every term of (15.23) vanishes.
If $r$ is even, (15.23) reads
\[
\mathrm{Iw}(v_1\cdots v_r)=\sum_{j=2}^{r}(-1)^j\,\mathrm{Mx}_{1j}\cdot\mathrm{Pfaff}\left(\mathrm{Mx}^{(1j)}\right),
\]
which is $\mathrm{Pfaff}(\mathrm{Mx})$ by (15.17).
\end{proof}

\subsubsection*{(j) Contours}

\begin{proof}[Proof of Proposition~15\,(j)]
Let $\mathrm{Vs}$ be the finite-dimensional normed space in which $f$ takes its values.
That $f$ is analytic means that its components with respect to a basis of $\mathrm{Vs}$ are analytic functions, and integrals of $f$ are formed component by component.

\emph{An inequality.}
For every continuous function $\mathrm{fc}$ from the circle $\{|t|=\mathrm{rad}\}$ to $\mathrm{Vs}$ one has
\[
\left\|\oint_{|t|=\mathrm{rad}}\mathrm{d}t\,\mathrm{fc}(t)\right\|\le2\pi\cdot\mathrm{rad}\cdot\sup_{|t|=\mathrm{rad}}\left\|\mathrm{fc}(t)\right\|.
\]
Indeed, let $t_1,\dots,t_{\mathrm{nt}}$ be $\mathrm{nt}$ points of the circle in counterclockwise order, and put $t_{\mathrm{nt}+1}:=t_1$.
The integral is the limit of the sums
\[
\sum_{\mathrm{ia}=1}^{\mathrm{nt}}\mathrm{fc}(t_{\mathrm{ia}})\cdot\left(t_{\mathrm{ia}+1}-t_{\mathrm{ia}}\right),
\]
formed for such subdivisions of the circle, as the longest arc between neighbouring points tends to $0$.
The norm of each of these sums is at most
\[
\sup\|\mathrm{fc}\|\cdot\sum_{\mathrm{ia}=1}^{\mathrm{nt}}\left|t_{\mathrm{ia}+1}-t_{\mathrm{ia}}\right|\le\sup\|\mathrm{fc}\|\cdot2\pi\cdot\mathrm{rad},
\]
since no chord is longer than its arc and the $\mathrm{nt}$ arcs together make up the circle.
The norm being continuous, the same bound holds for the limit.

\emph{The values.}
For $t\notin\{0,1\}$,
\begin{equation}\tag{15.24}
\frac{1}{t(t-1)}=\frac{1}{t-1}-\frac{1}{t},\qquad\frac{1}{t^2(t-1)}=\frac{1}{t-1}-\frac{1}{t}-\frac{1}{t^2},
\end{equation}
the second identity because $t^2-t(t-1)-(t-1)=1$.
By (15.24),
\begin{align*}
2\pi i\cdot\oint_{\mathrm{rad}}f&=\oint_{|t|=\mathrm{rad}}\mathrm{d}t\,\left(\frac{f(t)}{t-1}-\frac{f(t)}{t}\right),\\
2\pi i\cdot\oint^{(2)}_{\mathrm{rad}}f&=\oint_{|t|=\mathrm{rad}}\mathrm{d}t\,\left(\frac{f(t)}{t-1}-\frac{f(t)}{t}-\frac{f(t)}{t^2}\right).
\end{align*}
As $\mathrm{rad}>1$, the points $0$ and $1$ lie inside the circle, and the components of $f$ are analytic in a neighbourhood of the closed disk.
The residue of $f(t)/(t-1)$ at $t=1$ is $f(1)$, that of $f(t)/t$ at $t=0$ is $f(0)$, and that of $f(t)/t^2$ at $t=0$ is $f'(0)$, the coefficient of $t$ in the Taylor series of $f$ at $0$; the closed disk contains no other singularities.
The residue theorem, applied to the components and to the circle oriented counterclockwise (in the form \cite[Section~13, (cl1)]{QEDI}), therefore gives
\[
\oint_{\mathrm{rad}}f=f(1)-f(0),\qquad\oint^{(2)}_{\mathrm{rad}}f=f(1)-f(0)-f'(0).
\]

\emph{The bounds.}
On the circle, $|t|=\mathrm{rad}$ and $|t-1|\ge\mathrm{rad}-1>0$.
Hence, with the suprema over $|t|=\mathrm{rad}$,
\[
\left\|\frac{f(t)}{t(t-1)}\right\|\le\frac{\sup\|f\|}{\mathrm{rad}\,(\mathrm{rad}-1)},\qquad\left\|\frac{f(t)}{t^2(t-1)}\right\|\le\frac{\sup\|f\|}{\mathrm{rad}^2\,(\mathrm{rad}-1)}.
\]
Taking into account the factor $1/(2\pi)$ and the inequality above, one obtains
\[
\left\|\oint_{\mathrm{rad}}f\right\|\le\frac{\sup\|f\|}{\mathrm{rad}-1},\qquad\left\|\oint^{(2)}_{\mathrm{rad}}f\right\|\le\frac{\sup\|f\|}{\mathrm{rad}\,(\mathrm{rad}-1)}.
\]

\emph{The shift.}
Let $\theta=(\theta_1,\dots,\theta_p)$, where $\theta_i$ is a family of generators on $(\mathrm{Ind}_i,\nu_i)$.
For a complex number $t$ let
\[
f(t):=F(\phi+t\cdot\theta)\in\mathrm{Gr}[\phi,\theta]
\]
be the image of $F$ under the homomorphism $\phi_i(s)\mapsto\phi_i(s)+t\cdot\theta_i(s)$, which exists by 0.3.
The image of $\phi(I)$ is a product of $\#|I|$ factors, each of which is a polynomial of degree one in $t$.
Thus $f$ is a polynomial in $t$ with values in the finite-dimensional space $\mathrm{Gr}[\phi,\theta]$, and $f(1)=F(\phi+\theta)$, $f(0)=F(\phi)$.
The homomorphism has the form considered in (d), with the $2p$ families $\phi_1,\dots,\phi_p,\theta_1,\dots,\theta_p$ in place of $\theta_1,\dots,\theta_q$ and with the following operators: the identity from the functions on the index set of $\phi_i$ to the functions on $\mathrm{Ind}_i$; $t$ times the identity from the functions on the index set of $\theta_i$ to the functions on $\mathrm{Ind}_i$; the zero operator for all other families.
The identity on the functions on $(\mathrm{Ind}_i,\nu_i)$ has the kernel $\mathbf{1}[s=s']/\nu_i(s')$, whose row sums are
\[
\sum_{s'}\frac{\mathbf{1}[s=s']}{\nu_i(s')}\cdot\nu_i(s')=1.
\]
The numbers $C_{ij}$ of (d) may therefore be taken to be $1$, $|t|$ and $0$.
For $|t|=\mathrm{rad}$ the condition of (d) reads $\mathsf{w}'_i+\mathrm{rad}\cdot\mathsf{w}''_i\le\mathsf{w}_i$, and (d) gives $\|f(t)\|_{(\mathsf{w}',\mathsf{w}'')}\le\|F\|_{\mathsf{w}}$.
By the first part of the proof it follows that
\[
\left\|F(\phi+\theta)-F(\phi)\right\|_{(\mathsf{w}',\mathsf{w}'')}=\left\|\oint_{\mathrm{rad}}f\right\|_{(\mathsf{w}',\mathsf{w}'')}\le\frac{\|F\|_{\mathsf{w}}}{\mathrm{rad}-1}.
\]
\end{proof}

\subsubsection*{(k) Monotonicity, degree}

\begin{proof}[Proof of Proposition~15\,(k)]
Suppose that $0<\mathsf{w}'_i\le\mathsf{w}_i$ for all $i$.
Then $\mathsf{w}'^{I}\le\mathsf{w}^{I}$ for every $I$, whence
\[
\|F\|_{\mathsf{w}'}=\sum_I|c_F(I)|\cdot\mathsf{w}'^{I}\le\sum_I|c_F(I)|\cdot\mathsf{w}^{I}=\|F\|_{\mathsf{w}}.
\]
Suppose in addition that $\#|I_i|\ge r_i$ for every monomial $\phi(I)$ of $F$.
For these $I$, since $0<\mathsf{w}'_i/\mathsf{w}_i\le1$,
\[
\mathsf{w}'^{I}=\mathsf{w}^{I}\cdot\prod_i\left(\frac{\mathsf{w}'_i}{\mathsf{w}_i}\right)^{\#|I_i|}\le\mathsf{w}^{I}\cdot\prod_i\left(\frac{\mathsf{w}'_i}{\mathsf{w}_i}\right)^{r_i}.
\]
Multiplication by $|c_F(I)|$ and summation over $I$ yield the second bound.
\end{proof}

\subsubsection*{(l) The rules for one family $\Psi$}

\begin{proof}[Proof of Proposition~15\,(l)]
Throughout, $\int\mathrm{D}\Psi$, $\mathrm{d}\mu_I(\Psi)$ and $\mathrm{d}\mu_{\mathrm{Cov}}(\Psi)$ denote the objects of (D3) for the family $\Psi$ on $\mathrm{Mo}\times\{0,1\}$, as in (D5).
\begin{itemize}
\item[(i)] This is (b1) for $p=1$, $\mathrm{Ind}_1=\mathrm{Mo}\times\{0,1\}$ and $\mathsf{w}_1=\mathrm{aw}$.
\item[(ii)] This is (f) for $p=0$ and $W=\Psi$: here $\mathrm{Gr}[\phi]=\mathbb{C}$, the norm with the empty weight vector is the modulus, and $\|F\|_{(\mathsf{w}_W)}=|F|_{\mathsf{w}_W}$.
\item[(iii)] Let $\mathrm{hm}:=s_{\mathrm{Cov},1}$ be the homomorphism $\Psi(\mathrm{mo})\mapsto(\mathrm{Cov}\,\Psi)(\mathrm{mo})$, $\bar\Psi(\mathrm{mo})\mapsto\bar\Psi(\mathrm{mo})$, that is, the homomorphism of Lemma~0.7 with $O_1=\mathrm{Cov}$ and $O_2=1$; thus $f(\bar\Psi,\mathrm{Cov}\,\Psi)=\mathrm{hm}(f)$.
By 0.5\,(iii),
\[
\mathrm{hm}\left(\langle\bar\Psi,\mathrm{Cov}^{-1}\Psi\rangle_\nu\right)=\langle\bar\Psi,\mathrm{Cov}^{-1}(\mathrm{Cov}\,\Psi)\rangle_\nu=\langle\bar\Psi,\Psi\rangle_\nu,
\]
and by (b4)\,(iii),
\[
\mathrm{hm}\left(\exp\left(-\langle\bar\Psi,\mathrm{Cov}^{-1}\Psi\rangle_\nu\right)\right)=\exp\left(-\langle\bar\Psi,\Psi\rangle_\nu\right).
\]
Hence, by (15.7),
\begin{align*}
\int\mathrm{d}\mu_I\,\mathrm{hm}(f)&=\int\mathrm{D}\Psi\,\exp\left(-\langle\bar\Psi,\Psi\rangle_\nu\right)\cdot\mathrm{hm}(f)\\
&=\int\mathrm{D}\Psi\,\mathrm{hm}\left(\exp\left(-\langle\bar\Psi,\mathrm{Cov}^{-1}\Psi\rangle_\nu\right)\cdot f\right)\\
&=\det\mathrm{Cov}\cdot\int\mathrm{D}\Psi\,\exp\left(-\langle\bar\Psi,\mathrm{Cov}^{-1}\Psi\rangle_\nu\right)\cdot f\\
&=\int\mathrm{d}\mu_{\mathrm{Cov}}\,f.
\end{align*}

\emph{Total mass.}
By the definition of $\mathrm{d}\mu_{\mathrm{Cov}}$ and by the last sentence of (e), applied with $O=\mathrm{Cov}^{-1}$,
\[
\int\mathrm{d}\mu_{\mathrm{Cov}}=\det\mathrm{Cov}\cdot\det\left(\mathrm{Cov}^{-1}\right)=1.
\]

\emph{Two-point function.}
By the identity just proved, by (15.3) and by (e),
\begin{align*}
\int\mathrm{d}\mu_{\mathrm{Cov}}\,\Psi(\mathrm{mo})\cdot\bar\Psi(\mathrm{mo}')&=\int\mathrm{d}\mu_I\,(\mathrm{Cov}\,\Psi)(\mathrm{mo})\cdot\bar\Psi(\mathrm{mo}')\\
&=\sum_{\mathrm{mo}''}\mathrm{Cov}(\mathrm{mo},\mathrm{mo}'')\cdot\nu(\mathrm{mo}'')\cdot\frac{\mathbf{1}[\mathrm{mo}''=\mathrm{mo}']}{\nu(\mathrm{mo}')}=\mathrm{Cov}(\mathrm{mo},\mathrm{mo}').
\end{align*}
The weight $\nu(\mathrm{mo}')$ in the sum cancels against the moment $1/\nu(\mathrm{mo}')$, so that what remains is the kernel of $\mathrm{Cov}$, and not the entry $[\mathrm{Cov}]_{\mathrm{mo}\,\mathrm{mo}'}=\mathrm{Cov}(\mathrm{mo},\mathrm{mo}')\cdot\nu(\mathrm{mo}')$ of its matrix.

\emph{Last identity.}
By the definition of $\mathrm{d}\mu_{\mathrm{Cov}}$,
\begin{align*}
&\det\left(\mathrm{Cov}^{-1}\right)\cdot\int\mathrm{d}\mu_{\mathrm{Cov}}\,F\\
&\qquad=\det\left(\mathrm{Cov}^{-1}\right)\cdot\det\mathrm{Cov}\cdot\int\mathrm{D}\Psi\,\exp\left(-\langle\bar\Psi,\mathrm{Cov}^{-1}\Psi\rangle_\nu\right)\cdot F\\
&\qquad=\int\mathrm{D}\Psi\,\exp\left(-\langle\bar\Psi,\mathrm{Cov}^{-1}\Psi\rangle_\nu\right)\cdot F.
\end{align*}
\item[(iv)] This is (d) with $p=2$, $(\phi_1,\phi_2)=(\psi,\chi)$, $q=1$, $\theta_1=\Psi$ on the index set $\mathrm{Mo}\times\{0,1\}$ with the weights $\nu(\mathrm{mo},\omega)=\nu(\mathrm{mo})$, $O_{11}=H^{(1)}$, $O_{21}=H^{(2)}$, $C_{11}=C_{21}=\mathrm{Cr}$ and $\mathsf{w}'_1=\mathrm{aw}$; the condition $\sum_jC_{ij}\,\mathsf{w}'_j\le\mathsf{w}_i$ of (d) reads $\mathrm{Cr}\cdot\mathrm{aw}\le\mathsf{w}_1$ and $\mathrm{Cr}\cdot\mathrm{aw}\le\mathsf{w}_2$.
\end{itemize}

\emph{(l$'$) The last sentence of (iii): further generators.}
Let $f$ and $F$ lie in $\mathrm{Gr}[\phi,\Psi]$, and let $\mathrm{hm}=s_{\mathrm{Cov},1}$ fix the generators $\phi_i(t)$.
Lemma~0.7 holds in $\mathrm{Gr}[\phi,\Psi]$, and 0.5\,(iii) and (b4)\,(iii) hold in every Grassmann algebra.
Besides these three facts, the computation in the proof of the first identity of (iii) uses only that $\mathrm{hm}$ is a homomorphism and the definitions of $\mathrm{d}\mu_I$ and $\mathrm{d}\mu_{\mathrm{Cov}}$, while the proof of the last identity uses only the definition of $\mathrm{d}\mu_{\mathrm{Cov}}$.
Consequently both identities hold in $\mathrm{Gr}[\phi]$:
\begin{gather*}
\int\mathrm{d}\mu_{\mathrm{Cov}}\,f=\int\mathrm{d}\mu_I\,\mathrm{hm}(f),\\
\int\mathrm{D}\Psi\,\exp\left(-\langle\bar\Psi,\mathrm{Cov}^{-1}\Psi\rangle_\nu\right)\cdot F=\det\left(\mathrm{Cov}^{-1}\right)\cdot\int\mathrm{d}\mu_{\mathrm{Cov}}\,F
\end{gather*}
for all $f,F\in\mathrm{Gr}[\phi,\Psi]$.
\end{proof}

\subsubsection*{(m) Reduction to the weight one}

\begin{proof}[Proof of Proposition~15\,(m)]
Since $\nu(\mathrm{mo},0)^{\frac{1}{2}}\cdot\nu(\mathrm{mo},1)^{\frac{1}{2}}=\nu(\mathrm{mo})$,
\[
\sum_{\mathrm{mo}}\overline{\mathrm{Wr}}(\mathrm{mo})\cdot\mathrm{Wr}(\mathrm{mo})=\sum_{\mathrm{mo}}\nu(\mathrm{mo})\cdot\bar W(\mathrm{mo})\cdot W(\mathrm{mo})=\langle\bar W,W\rangle_\nu.
\]
Each $\mathrm{Wr}(u)$ is a non-zero multiple of $W(u)$.
The $\mathrm{Wr}(u)$ satisfy (15.1), and if $\mathrm{Wr}(J)$ denotes the ordered product of the $\mathrm{Wr}(u)$, $u\in J$, then
\[
\mathrm{Wr}(J)=\nu^{\frac{1}{2}}(J)\cdot W(J),\qquad\nu^{\frac{1}{2}}(J):=\prod_{u\in J}\nu(u)^{\frac{1}{2}}>0.
\]
Hence the $\phi(I)\cdot\mathrm{Wr}(J)$ form a basis: the $\phi_i(t)$ and the $\mathrm{Wr}(u)$ are independent generators of the same algebra.
Put
\[
\mathrm{Top}_r:=\prod_{\mathrm{mo}}\mathrm{Wr}(\mathrm{mo})\cdot\overline{\mathrm{Wr}}(\mathrm{mo})=\left(\prod_{\mathrm{mo}}\nu(\mathrm{mo})\right)\cdot\mathrm{Top};
\]
note also that $\nu^{\frac{1}{2}}(J_f)=\prod_{\mathrm{mo}}\nu(\mathrm{mo})$.
From $W(J_f)=\mathrm{sn}_f\cdot\mathrm{Top}$ one obtains $\mathrm{Wr}(J_f)=\mathrm{sn}_f\cdot\mathrm{Top}_r$.
The coefficient of $\mathrm{Top}_r$ in $\mathrm{Wr}(J)$ is therefore $\mathrm{sn}_f$ for $J=J_f$ and $0$ otherwise, exactly as for $\mathrm{Top}$ and $W(J)$.
By (D3) and linearity,
\begin{align*}
\int\mathrm{D}W\,\phi(I)\cdot\mathrm{Wr}(J)&=\nu^{\frac{1}{2}}(J)\cdot\phi(I)\cdot\left(\prod_{\mathrm{mo}}\nu(\mathrm{mo})^{-1}\right)\\
&\qquad\cdot\left(\text{the coefficient of }\mathrm{Top}\text{ in }W(J)\right)\\
&=\phi(I)\cdot\left(\text{the coefficient of }\mathrm{Top}_r\text{ in }\mathrm{Wr}(J)\right),
\end{align*}
because both sides vanish for $J\ne J_f$, and $\nu^{\frac{1}{2}}(J_f)\cdot\prod_{\mathrm{mo}}\nu(\mathrm{mo})^{-1}=1$.
This is the assertion: $\int\mathrm{D}W$ is the linear map that sends $\phi(I)\cdot\mathrm{Wr}(J)$ to $\phi(I)$ times the coefficient of $\prod_{\mathrm{mo}}\mathrm{Wr}(\mathrm{mo})\cdot\overline{\mathrm{Wr}}(\mathrm{mo})$ in $\mathrm{Wr}(J)$.

In \cite[Section~13.2]{QEDI} the following objects are defined for generators $W_\beta(y)$, $\bar W_\beta(y)$:
\[
\langle\bar W,W\rangle:=\sum_{y,\beta}\bar W_\beta(y)\cdot W_\beta(y),\qquad\mathrm{Top}:=\prod_{(y,\beta)}W_\beta(y)\cdot\bar W_\beta(y);
\]
$\int\mathrm{D}W$ is the linear map of $\mathrm{Gr}[\phi,W]$ to $\mathrm{Gr}[\phi]$ that sends $\phi(I)\cdot W(J)$ to $\phi(I)$ times the coefficient of $\mathrm{Top}$ in $W(J)$; and
\[
\int\mathrm{d}\mu_I(W)\,F:=\int\mathrm{D}W\,\exp\left(-\langle\bar W,W\rangle\right)\cdot F.
\]
So the map described above is the map $\int\mathrm{D}W$ of \cite[Section~13.2]{QEDI}, read with the generators $\mathrm{Wr}$ and with the index set $\mathrm{Mo}$ in place of $\{(y,\beta)\}$.
Together with the first display of this proof it follows that
\[
\int\mathrm{d}\mu_I(W)\,F=\int\mathrm{D}W\,\exp\left(-\sum_{\mathrm{mo}}\overline{\mathrm{Wr}}(\mathrm{mo})\cdot\mathrm{Wr}(\mathrm{mo})\right)\cdot F
\]
is the functional $\mathrm{d}\mu_I$ of \cite[Section~13.2]{QEDI} in the generators $\mathrm{Wr}$.
\end{proof}

(The coefficients of $F$ with respect to the basis $\phi(I)\cdot\mathrm{Wr}(J)$ are $c_F(I,J)\cdot\nu^{-\frac{1}{2}}(J)$.
Hence the middle term of the display in (f) is the norm of $F$ formed with the generators $(\phi,\mathrm{Wr})$ and the weights $(\mathsf{w}_1,\dots,\mathsf{w}_p,1)$.)

\subsubsection*{The empty set of modes}

\begin{proof}[Proof of the assertion of (D2) on an empty set $\mathrm{Mo}$]
Let $\mathrm{Mo}$ be empty, and let $\mathsf{w}_W>0$ be any number.
The objects are those of 0.6\,(iv): $\mathrm{Gr}[\phi,W]=\mathrm{Gr}[\phi]$; $\mathrm{Top}=1$; the maps $\int\mathrm{D}W$, $F\mapsto\int\mathrm{d}\mu_I(W)\,F$ and $F\mapsto\int\mathrm{d}\mu_{\mathrm{Cov}}(W)\,F$ are the identity of $\mathrm{Gr}[\phi]$; the only subset $J$ is the empty set, with $W(J)=1$ and $c_F(I,J)=c_F(I)$; the only linear map $O$ is the identity, with $\det O=1$; and all forms $\langle\cdot,\cdot\rangle_\nu$ vanish.
For a family $\Psi$ on $\mathrm{Mo}\times\{0,1\}$ as in (D5) one has $\mathrm{Gr}[\Psi]=\mathbb{C}$, $\mathrm{Gr}[\phi,\Psi]=\mathrm{Gr}[\phi]$, and $|F|_{\mathrm{aw}}=|F|$ for every $\mathrm{aw}>0$.
The clauses (a)--(d), (j) and (k) do not involve $\mathrm{Mo}$, and in (n) the set $\mathrm{Mo}$ is not empty.
The other clauses take the following form.
\begin{itemize}
\item[(e)] The empty set $J$ is paired.
$W(J)=1$ and $m(J)=\int\mathrm{d}\mu_I(W)\,1=1$ are $\mathrm{sn}(J)$ times an empty product, with $\mathrm{sn}(J)=1$; $m(J)$ is real; $\nu^{-\frac{1}{2}}(J)=1=|m(J)|$; and $\int\mathrm{d}\mu_I(W)=1$.
The identities for the products of two generators are void, there being no mode.
Finally, $\int\mathrm{D}W\,\exp\left(-\langle\bar W,OW\rangle_\nu\right)=\int\mathrm{D}W\,1=1=\det O$.
\item[(f)] The first display reads $F=F$.
In the inequalities $\nu^{-\frac{1}{2}}(J)=1=\mathsf{w}_W^{\#|J|}$, so that all three terms are equal to $\|F\|_{(\mathsf{w}_1,\dots,\mathsf{w}_p)}$.
The only subset $\mathrm{Mo}^0$ is the empty set $\mathrm{Mo}$ itself, and $\int\mathrm{d}\mu_I(W)\,F=F$.
(The assertion that $\mathsf{w}_W$ cannot be lowered concerns a mode; it is the exception made in (D2).)
\item[(g)] Both sides are equal to $F_1\cdot F_2$.
\item[(h)] $U$ is the identity, and $s_U$ is the identity of $\mathrm{Gr}[\phi]$; the three assertions read $0=0$, $F=F$ and $F=F$.
\item[(i)] The linear span of the $W(u)$ is $\{0\}$.
For $r=0$ both sides are equal to $F_0$.
For $r\ge1$ all $v_j$ vanish and the integral is $0$; for even $r\ge2$ the matrix of the clause is $0$, and so is its Pfaffian, since every term of (15.16) contains the factor $\mathrm{Mx}_{i_1j_1}$.
\item[(l)] (i) $|F\cdot G|=|F|\cdot|G|$.
(ii) $\left|\int\mathrm{d}\mu_I(\Psi)\,F\right|=|F|=|F|_{\mathsf{w}_W}$.
(iii) $\mathrm{Cov}$ is the identity, and so is the homomorphism.
The first and the last identity read $f=f$ and $F=F$, in $\mathbb{C}$ and in $\mathrm{Gr}[\phi,\Psi]=\mathrm{Gr}[\phi]$; $\int\mathrm{d}\mu_{\mathrm{Cov}}=1$; the identity for $\Psi(\mathrm{mo})\cdot\bar\Psi(\mathrm{mo}')$ is void.
(iv) $(H^{(i)}\Psi)(s)$ is an empty sum, equal to $0$.
Thus $E(H^{(1)}\Psi,H^{(2)}\Psi)$ is the image of $E$ under $\psi(s)\mapsto0$, $\chi(s)\mapsto0$ (0.3; the element $0$ is odd).
This image is the coefficient of the monomial $1$ in $E$, and its modulus is one of the terms of the sum (I.5.1) for $\|E\|_{(\mathsf{w}_1,\mathsf{w}_2)}$.
\item[(m)] There is no $\mathrm{Wr}(u)$, and both sides of the first identity are empty sums.
Moreover $\prod_{\mathrm{mo}}\mathrm{Wr}(\mathrm{mo})\cdot\overline{\mathrm{Wr}}(\mathrm{mo})=1=\mathrm{Top}$, and the identity of $\mathrm{Gr}[\phi]$ sends $\phi(I)$ to $\phi(I)$ times the coefficient of $1$ in $1$.
The definitions of \cite[Section~13.2]{QEDI} recalled in the proof of (m), read with the empty index set, give the identity of $\mathrm{Gr}[\phi]$ both for $\int\mathrm{D}W$ and for $\mathrm{d}\mu_I(W)$.
\end{itemize}
\end{proof}

\subsubsection*{(n) Values in the scheme}

\begin{proof}[Proof of Proposition~15\,(n)]
Let $k\ge0$ be an integer and $\ell$ a side; in (n2)--(n5) only the fact that $\ell$ is a positive integer is used.
Let $\mathrm{fdo}$ be a fundamental domain of $\ell\mathbb{Z}^3$ in $\mathbb{Z}^3$, that is, a subset that contains exactly one point of every class modulo $\ell\mathbb{Z}^3$; then $\#|\mathrm{fdo}|=\ell^3=n_k$.
By (I.AP) in Section~I.3.6 the independent generators are those whose point (for a pair: whose first point) lies over a fundamental domain; it is shown in (n5) that nothing depends on the choice of this domain.
The total orders required in (D1) and (D2) are chosen arbitrarily; by 0.1\,(iii) and 0.6\,(iii) neither the norms nor the integrals depend on them.

\textbf{(n1) $\psi$ and $\chi$, every $k\ge0$.}
This is (D1) with $p=2$.
The first index set is $\mathrm{Ind}_1=\{\xi\}$, where $\xi=(c,g,\beta,\omega)$ with $c$ in a fundamental domain of $\ell\mathbb{Z}^3$ in $\mathbb{L}_k$ (it consists of $(\ell/\eta_k)^3$ photon sites), $g\in O_{\mathrm{cub}}$ ($48$ labels), $\beta\in\{1,\dots,4\}$ and $\omega\in\{0,1\}$; the weights are $\nu_1(\xi)=\eta_k^3/48$.
The second index set is $\mathrm{Ind}_2=\{\boldsymbol{\zeta}\}$, where $\boldsymbol{\zeta}=(x,x',\beta,\omega)$ with $(x,x')\in\mathrm{Pz}_k$; the weights are $\nu_2(\boldsymbol{\zeta})=(\eta_k^3/48)^2$.
Thus
\[
\int\mathrm{d}\xi\,(\cdot)=\sum_c\sum_{g\in O_{\mathrm{cub}}}\sum_\beta\sum_\omega\frac{\eta_k^3}{48}\,(\cdot),
\]
and by (a) the norm (I.5.1) with $\mathsf{w}=(\mathsf{w}_1,\mathsf{w}_2)$ is, for antisymmetric kernels,
\[
\sum_{n,m}\frac{\mathsf{w}_1^n\cdot\mathsf{w}_2^m}{n!\,m!}\int\mathrm{d}\xi\,\mathrm{d}\boldsymbol{\zeta}\,|E_{nm}|,
\]
which is the norm of Section~I.5.1.
The weights $1/48$ stand in the integrals and not in (I.5.1).
For example, by Section~I.1.6 one has $\int_{\square}\mathrm{d}x=\operatorname{Vol}(\square)=M^3$ for an $M$-cube $\square$, and the cluster of a photon site has the measure $48\cdot\eta_k^3/48=\eta_k^3$; hence $\square$ contains $(M/\eta_k)^3$ photon sites.
Consequently
\[
\langle\bar\psi,\psi\rangle_{\square}=\sum_{c\in\square}\sum_{g\in O_{\mathrm{cub}}}\sum_\beta\frac{\eta_k^3}{48}\cdot\bar\psi_\beta(c,g)\cdot\psi_\beta(c,g)
\]
has $(M/\eta_k)^3\cdot48\cdot4$ monomials, each with a coefficient of modulus $\eta_k^3/48$, and its norm is
\[
\left(\frac{M}{\eta_k}\right)^3\cdot48\cdot4\cdot\frac{\eta_k^3}{48}\cdot\mathsf{w}_1^2=4M^3\cdot\mathsf{w}_1^2,
\]
the value stated in Section~I.5.1.

\textbf{(n2) $\Psi$ and $W$, levels $k\ge1$ (the paragraph on the levels $k\ge1$ in (n)).}
Here $\mathrm{Mo}=\mathrm{fdo}\times\{1,\dots,4\}$, $\mathrm{mo}=(y,\beta)$, $\nu(\mathrm{mo})=1$ and $\#|\mathrm{Mo}|=4\,n_k$.
Then $\mathsf{w}_W=\max_{\mathrm{mo}}1^{-\frac{1}{2}}=1$ and $\prod_{\mathrm{mo}}\nu(\mathrm{mo})^{-1}=1$, and by (15.11) $m(J)=\mathrm{sn}(J)$ for paired $J$.
Restriction to $\mathrm{fdo}$ is a linear bijection of $\mathfrak{S}(1,\ell)$ onto the functions on $\mathrm{Mo}$: an antiperiodic function is determined by its values on $\mathrm{fdo}$, and every function on $\mathrm{fdo}\times\{1,\dots,4\}$ extends by
\[
v(y+\ell z)=(-1)^{z_1+z_2+z_3}\cdot v(y),\qquad z\in\mathbb{Z}^3.
\]
Under this bijection $\langle v',v\rangle_\nu=\sum_{y\in\mathrm{fdo}}v'(y)\cdot v(y)=\langle v',v\rangle_1$.
The kernel of (D2) is $O(\mathrm{mo},\mathrm{mo}')=O(y,y')_{\beta\beta'}$, where $(Ov)(y)=\sum_{y'\in\mathrm{fdo}}O(y,y')\cdot v(y')$ (Section~I.1.6).
The form
\[
\langle\bar\Psi,O\Psi\rangle=\sum_{y,y'\in\mathrm{fdo}}\bar\Psi(y)\cdot O(y,y')\cdot\Psi(y')
\]
of Section~I.3.6 is $\langle\bar\Psi,O\Psi\rangle_\nu$.
A linear map $O$ of $\mathfrak{S}(1,\ell)$ and the corresponding linear map of the functions on $\mathrm{Mo}$ have the same determinant, because the determinant of a linear map does not change under conjugation with a linear bijection.
Accordingly, for $y,y'\in\mathrm{fdo}$, the clauses (e) and (l)\,(iii) read
\begin{gather*}
\int\mathrm{d}\mu_I\,W_\beta(y)\cdot\bar W_{\beta'}(y')=\mathbf{1}[y=y',\,\beta=\beta'],\\
\int\mathrm{D}\Psi\,\exp\left(-\langle\bar\Psi,O\Psi\rangle\right)=\det O\qquad\text{for every linear map $O$ of $\mathfrak{S}(1,\ell)$},\\
\int\mathrm{d}\mu_{\mathrm{Cov}}\,\Psi(y)\cdot\bar\Psi(y')^T=\mathrm{Cov}(y,y'),
\end{gather*}
the last being the $4\times4$ matrix with the entries $\int\mathrm{d}\mu_{\mathrm{Cov}}\,\Psi_\beta(y)\cdot\bar\Psi_{\beta'}(y')=\mathrm{Cov}(\mathrm{mo},\mathrm{mo}')$.
(The family $\Psi'$ of the fifth row of the table in Section~I.3.6 is a family of this kind on $T^0_{k+1}$, for every $k\ge0$, the positive integer $\ell/L$ standing in place of $\ell$.)

\textbf{(n3) Level $0$: the objects.}
Here $\mathrm{Mo}=\mathrm{fdo}\times O_{\mathrm{cub}}\times\{1,\dots,4\}$ and $\mathrm{mo}=(x,\beta)=(c,g,\beta)$ with $x=(c,g)\in T^{\mathrm{f}}_0$; moreover $\nu(\mathrm{mo})=1/48$ for all $\mathrm{mo}$, and $\#|\mathrm{Mo}|=48\cdot4\cdot\ell^3=192\,n_0$, so that there are $192$ modes over every photon site ($48$ sites of the cluster, $4$ spin components).
Restriction to $\mathrm{fdo}\times O_{\mathrm{cub}}$ is a linear bijection of $\mathfrak{S}^{\mathrm{ch}}(1,\ell)$ onto the functions on $\mathrm{Mo}$: an antiperiodic function is determined by its values on $\mathrm{fdo}\times O_{\mathrm{cub}}$, every function on $\mathrm{fdo}\times O_{\mathrm{cub}}\times\{1,\dots,4\}$ extends by
\[
f(c+\ell z,g)=(-1)^{z_1+z_2+z_3}\cdot f(c,g),\qquad z\in\mathbb{Z}^3,
\]
and both spaces have the dimension $192\,\ell^3$ (Section~I.1.5).
With $\eta_0=1$ the objects of (D2) and (D3) are as follows.
By (I.1.3) and Section~I.1.6,
\begin{align*}
\langle f',f\rangle_\nu&=\sum_{c\in\mathrm{fdo}}\sum_{g\in O_{\mathrm{cub}}}\frac{1}{48}\sum_{\beta=1}^{4}f'_\beta(c,g)\cdot f_\beta(c,g)\\
&=\int\mathrm{d}x\,f'(x)\cdot f(x)=\langle f',f\rangle^{\mathrm{ch}}.
\end{align*}
A linear map $O$ acts by
\begin{align*}
(Of)_\beta(x)&=\sum_{c'\in\mathrm{fdo}}\sum_{g'\in O_{\mathrm{cub}}}\frac{1}{48}\sum_{\beta'=1}^{4}O\left(x,(c',g')\right)_{\beta\beta'}\cdot f_{\beta'}(c',g')\\
&=\left(\int\mathrm{d}x'\,O(x,x')\cdot f(x')\right)_\beta,
\end{align*}
so that $O(\mathrm{mo},\mathrm{mo}')=O(x,x')_{\beta\beta'}$ is the kernel (I.1.4).
By Section~I.3.6 (the forms at $k=0$),
\begin{align*}
\langle\bar W,OW\rangle_\nu&=\sum_{c\in\mathrm{fdo}}\sum_{g\in O_{\mathrm{cub}}}\frac{1}{48}\sum_\beta\bar W_\beta(c,g)\cdot(OW)_\beta(c,g)\\
&=\int\mathrm{d}x\,\bar W(x)\cdot(OW)(x)=\langle\bar W,OW\rangle^{\mathrm{ch}}.
\end{align*}
Further,
\[
\mathrm{Top}=\prod_{c\in\mathrm{fdo}}\prod_{g\in O_{\mathrm{cub}}}\prod_\beta W_\beta(c,g)\cdot\bar W_\beta(c,g)
\]
is a product of $192\,n_0$ factors, and $\prod_{\mathrm{mo}}\nu(\mathrm{mo})^{-1}=48^{192n_0}$.
By (15.10),
\[
\mathrm{Ew}=\prod_{c,g,\beta}\left(1+\frac{1}{48}\cdot W_\beta(c,g)\cdot\bar W_\beta(c,g)\right).
\]
The weight of the family is
\[
\mathsf{w}_W=\max_{\mathrm{mo}}\left(\frac{1}{48}\right)^{-\frac{1}{2}}=\sqrt{48},
\]
and $\sqrt{48}\le7$ because $48\le49=7^2$.
Finally, a paired $J$ has $\#|J|/2$ modes, so that (15.11) gives
\begin{equation}\tag{15.25}
m(J)=\mathrm{sn}(J)\cdot48^{\frac{\#|J|}{2}},\quad\left|m(J)\right|=\left(\sqrt{48}\right)^{\#|J|}\quad\text{for paired }J,\qquad m(J)=0\quad\text{otherwise}.
\end{equation}
The family $\Psi$ of level $0$ has the same index set and the same weights (Section~I.3.6); hence all the identities just listed hold with $\Psi$, $\bar\Psi$ in place of $W$, $\bar W$, in accordance with (D5).
In the last assertion of (g) the points $t$ are the fermion sites $x=(c,g)$: the $48$ sites of a cluster are $48$ different points, each carrying $4$ modes.
In (h) and in (l)\,(iii) the maps $U$ and $\mathrm{Cov}$ are linear maps of $\mathfrak{S}^{\mathrm{ch}}(1,\ell)$, and ${}^T$ is the transpose with respect to $\langle\cdot,\cdot\rangle^{\mathrm{ch}}$; by (15.4) the kernel (I.1.4) of $O^T$ is
\[
O^T(x',x)_{\beta'\beta}=O(x,x')_{\beta\beta'},
\]
that is, the $4\times4$ block $O^T(x',x)$ is the transpose of the block $O(x,x')$.

\textbf{(n4) Level $0$: the three identities.}
\emph{First.}
By (e), $\int\mathrm{d}\mu_I\,W(\mathrm{mo})\cdot\bar W(\mathrm{mo}')=\mathbf{1}[\mathrm{mo}=\mathrm{mo}']/\nu(\mathrm{mo})$; with $\mathrm{mo}=(x,\beta)$, $\mathrm{mo}'=(x',\beta')$ and $\nu=1/48$ this becomes
\[
\int\mathrm{d}\mu_I\,W_\beta(x)\cdot\bar W_{\beta'}(x')=48\cdot\mathbf{1}[x=x',\,\beta=\beta'].
\]
\emph{Second.}
Through the bijection of (n3) a linear map $O$ of $\mathfrak{S}^{\mathrm{ch}}(1,\ell)$ is a linear map of the functions on $\mathrm{Mo}$ with the same determinant (the determinant of a linear map does not change under conjugation with a linear bijection), and $\langle\bar\Psi,O\Psi\rangle^{\mathrm{ch}}=\langle\bar\Psi,O\Psi\rangle_\nu$ by (n3).
Hence the last sentence of (e) gives
\[
\int\mathrm{D}\Psi\,\exp\left(-\langle\bar\Psi,O\Psi\rangle^{\mathrm{ch}}\right)=\det O;
\]
here $O$ need not be invertible.
\emph{Third.}
$\int\mathrm{d}\mu_{\mathrm{Cov}}\,\Psi(x)\cdot\bar\Psi(x')^T$ is the $4\times4$ matrix with the entries $\int\mathrm{d}\mu_{\mathrm{Cov}}\,\Psi_\beta(x)\cdot\bar\Psi_{\beta'}(x')$.
By (l)\,(iii) this entry is $\mathrm{Cov}(\mathrm{mo},\mathrm{mo}')=\mathrm{Cov}(x,x')_{\beta\beta'}$, which by (n3) is the entry of the kernel (I.1.4).

\textbf{(n5) The fundamental domain; the covering lattices (every $k\ge0$).}
Let $\widetilde{\mathrm{fdo}}$ be a second fundamental domain of $\ell\mathbb{Z}^3$ in $\mathbb{Z}^3$.
By (I.AP) a generator changes its sign when its point is translated by $\ell e_\mu$; read from right to left, (I.AP) states the same for the translation by $-\ell e_\mu$.
Hence a translation of the point by $\ell z$, $z\in\mathbb{Z}^3$, multiplies a generator by $(-1)^{z_1+z_2+z_3}$.

\emph{Level $0$: the generators.}
For $c\in\mathrm{fdo}$ let $c+\ell z(c)$, with $z(c)\in\mathbb{Z}^3$, be the point of $\widetilde{\mathrm{fdo}}$ in the class of $c$, and put $\mathrm{sz}(c):=(-1)^{z_1(c)+z_2(c)+z_3(c)}$.
The map $c\mapsto c+\ell z(c)$ is a bijection of $\mathrm{fdo}$ onto $\widetilde{\mathrm{fdo}}$, and
\begin{equation}\tag{15.26}
W\left(c+\ell z(c),g,\beta,\omega\right)=\mathrm{sz}(c)\cdot W(c,g,\beta,\omega)\qquad(\omega=0\text{ and }\omega=1)
\end{equation}
($c\in\mathrm{fdo}$, $g\in O_{\mathrm{cub}}$, $\beta\in\{1,\dots,4\}$).
Thus the generators that belong to $\widetilde{\mathrm{fdo}}$ are the generators that belong to $\mathrm{fdo}$, multiplied by signs, and the sign is the \emph{same} for $W(\mathrm{mo})$ and $\bar W(\mathrm{mo})$, $\mathrm{mo}=(c,g,\beta)$.
This has the following consequences.
\begin{itemize}
\item[(1)] A product of generators of one choice is $\pm$ the product of the corresponding generators of the other choice.
So an ordered monomial of one choice is $\pm$ an ordered monomial of the other choice (0.1\,(i)), and the moduli of the coefficients of an element, hence all norms (I.5.1), are the same for both choices.
\item[(2)] One has
\begin{align*}
W\left(c+\ell z(c),g,\beta,0\right)\cdot W\left(c+\ell z(c),g,\beta,1\right)&=\mathrm{sz}(c)^2\cdot W(\mathrm{mo})\cdot\bar W(\mathrm{mo})\\
&=\mathrm{ww}(\mathrm{mo}).
\end{align*}
Therefore $\mathrm{Top}=\prod_{\mathrm{mo}}\mathrm{ww}(\mathrm{mo})$, $\langle\bar W,W\rangle_\nu=-\sum_{\mathrm{mo}}\frac{1}{48}\cdot\mathrm{ww}(\mathrm{mo})$ and $\mathrm{Ew}=\exp\left(-\langle\bar W,W\rangle_\nu\right)$ are the same elements for both choices.
\item[(3)] By (1) the linear span of the products of fewer than $2\,\#|\mathrm{Mo}|=384\,n_0$ generators $W(u)$ is the same for both choices, and by (2) so is $\mathrm{Top}$.
Hence the number $\mathrm{tc}(G)$ of 0.6\,(iii) is the same.
The factor $\prod_{\mathrm{mo}}\nu(\mathrm{mo})^{-1}$ is $48^{192n_0}$ for both choices.
By the description of $\int\mathrm{D}W$ in 0.6\,(iii), $\int\mathrm{D}W$ is the same map; together with (2), $\mathrm{d}\mu_I(W)$ is the same functional.
\end{itemize}
By (I.AP), (15.26) holds for $\Psi$ as well; so (1)--(3) hold with $\Psi$, $\bar\Psi$ in place of $W$, $\bar W$.

\emph{Level $0$: kernels, forms, determinants.}
Let $O$ be a linear map of $\mathfrak{S}^{\mathrm{ch}}(1,\ell)$.
Its kernel (I.1.4) is a function of $x,x'\in\mathbb{L}^{\mathrm{f}}_0$ with
\begin{equation}\tag{15.27}
O(x,x'+\ell e_\mu)=-O(x,x'),\qquad O(x+\ell e_\mu,x')=-O(x,x')\qquad(\mu=1,2,3).
\end{equation}
The first relation is part of (I.1.4).
To prove the second, recall from (I.1.4) that for every $x\in\mathbb{L}^{\mathrm{f}}_0$ and every $f\in\mathfrak{S}^{\mathrm{ch}}(1,\ell)$
\[
(Of)_\beta(x)=\sum_{c'\in\mathrm{fdo}}\sum_{g'\in O_{\mathrm{cub}}}\frac{1}{48}\sum_{\beta'}O\left(x,(c',g')\right)_{\beta\beta'}\cdot f_{\beta'}(c',g').
\]
The values $f_{\beta'}(c',g')$, $c'\in\mathrm{fdo}$, are arbitrary complex numbers (n3).
So the numbers $O(x,x')_{\beta\beta'}$, $x'$ over $\mathrm{fdo}$, are determined by the linear functional $f\mapsto(Of)_\beta(x)$: they are $48$ times its coefficients.
Now $Of$ is antiperiodic, $(Of)(x+\ell e_\mu)=-(Of)(x)$ for every $f$.
Hence $O(x+\ell e_\mu,x')=-O(x,x')$ for $x'$ over $\mathrm{fdo}$, and then, by the first relation, for all $x'\in\mathbb{L}^{\mathrm{f}}_0$.

By the first relation of (15.27), $O(x,x')\cdot f(x')$ does not change under $x'\mapsto x'+\ell e_\mu$.
So (I.1.4) holds with one and the same function $O(x,x')$ for every fundamental domain of integration: the kernel does not depend on the choice.
By the first relation and (I.AP), $O(x,x')\cdot\Psi(x')$ does not change under $x'\mapsto x'+\ell e_\mu$.
Therefore
\[
(O\Psi)(x):=\int\mathrm{d}x'\,O(x,x')\cdot\Psi(x')
\]
is, for every $x\in\mathbb{L}^{\mathrm{f}}_0$, the same element for both choices, and $(O\Psi)(x+\ell e_\mu)=-(O\Psi)(x)$ by the second relation.
With (I.AP) it follows that $\bar\Psi(x)\cdot(O\Psi)(x)$ does not change under $x\mapsto x+\ell e_\mu$.
So $\langle\bar\Psi,O\Psi\rangle^{\mathrm{ch}}=\int\mathrm{d}x\,\bar\Psi(x)\cdot(O\Psi)(x)$ is the same element for both choices.
The number $\det O$ is the determinant of a linear map of the vector space $\mathfrak{S}^{\mathrm{ch}}(1,\ell)$; its definition involves no fundamental domain, and by (n4) it is the determinant of the corresponding map of the functions on $\mathrm{Mo}$ for each of the two choices.
Together with (3) this shows that
\[
\int\mathrm{D}\Psi\,\exp\left(-\langle\bar\Psi,O\Psi\rangle^{\mathrm{ch}}\right)\cdot F
\]
and, for invertible $\mathrm{Cov}$,
\[
\int\mathrm{d}\mu_{\mathrm{Cov}}(\Psi)\,F=\det\mathrm{Cov}\cdot\int\mathrm{D}\Psi\,\exp\left(-\langle\bar\Psi,\mathrm{Cov}^{-1}\Psi\rangle^{\mathrm{ch}}\right)\cdot F
\]
do not depend on the choice.

\emph{Level $0$: arguments on the covering lattice.}
Let $x=(c+\ell z,g)$ and $x'=(c'+\ell z',g')$ with $c,c'\in\mathrm{fdo}$ and $z,z'\in\mathbb{Z}^3$.
By (I.AP) and the first identity of (n4),
\begin{equation}\tag{15.28}
\begin{split}
\int\mathrm{d}\mu_I\,W_\beta(x)\cdot\bar W_{\beta'}(x')&=(-1)^{z_1+z_2+z_3+z'_1+z'_2+z'_3}\\
&\qquad\cdot48\cdot\mathbf{1}\left[(c,g)=(c',g'),\,\beta=\beta'\right].
\end{split}
\end{equation}
The right side is the kernel (I.1.4) of the identity, which is $48\cdot\mathbf{1}[x=x']\cdot1_4$ for $x,x'$ over $\mathrm{fdo}$ by (I.ke1), continued by (15.27).
The third identity of (n4),
\[
\int\mathrm{d}\mu_{\mathrm{Cov}}\,\Psi(x)\cdot\bar\Psi(x')^T=\mathrm{Cov}(x,x'),
\]
holds for all $x,x'\in\mathbb{L}^{\mathrm{f}}_0$: for $x,x'$ over $\mathrm{fdo}$ it is (n4), and both sides change their sign under $x\mapsto x+\ell e_\mu$ and under $x'\mapsto x'+\ell e_\mu$, the left side by (I.AP), the right side by (15.27).

\emph{Levels $k\ge1$: the generators.}
For $y\in\mathrm{fdo}$ let $y+\ell z(y)$, with $z(y)\in\mathbb{Z}^3$, be the point of $\widetilde{\mathrm{fdo}}$ in the class of $y$, and put $\mathrm{sz}(y):=(-1)^{z_1(y)+z_2(y)+z_3(y)}$.
The map $y\mapsto y+\ell z(y)$ is a bijection of $\mathrm{fdo}$ onto $\widetilde{\mathrm{fdo}}$, and
\begin{equation}\tag{15.29}
W\left(y+\ell z(y),\beta,\omega\right)=\mathrm{sz}(y)\cdot W(y,\beta,\omega)\qquad(\omega=0\text{ and }\omega=1)
\end{equation}
($y\in\mathrm{fdo}$, $\beta\in\{1,\dots,4\}$).
The sign is the same for $W(\mathrm{mo})$ and $\bar W(\mathrm{mo})$, $\mathrm{mo}=(y,\beta)$.
The consequences are these.
\begin{itemize}
\item[(1$'$)] A product of generators of one choice equals $\pm$ the product of the corresponding generators of the other choice.
Hence an ordered monomial of one choice is $\pm$ an ordered monomial of the other choice (0.1\,(i)), and the moduli of the coefficients of an element, and with them all norms (I.5.1), coincide for the two choices.
\item[(2$'$)] One has
\begin{align*}
W\left(y+\ell z(y),\beta,0\right)\cdot W\left(y+\ell z(y),\beta,1\right)&=\mathrm{sz}(y)^2\cdot W(\mathrm{mo})\cdot\bar W(\mathrm{mo})\\
&=\mathrm{ww}(\mathrm{mo}).
\end{align*}
So $\mathrm{Top}=\prod_{\mathrm{mo}}\mathrm{ww}(\mathrm{mo})$, $\langle\bar W,W\rangle_\nu=-\sum_{\mathrm{mo}}\mathrm{ww}(\mathrm{mo})$ (here $\nu=1$) and $\mathrm{Ew}=\exp\left(-\langle\bar W,W\rangle_\nu\right)$ are the same elements for both choices.
\item[(3$'$)] By (1$'$) the linear span of the products of fewer than $2\,\#|\mathrm{Mo}|=8\,n_k$ generators $W(u)$ is the same for both choices, and by (2$'$) the element $\mathrm{Top}$ is the same.
So the number $\mathrm{tc}(G)$ of 0.6\,(iii) is the same.
The factor $\prod_{\mathrm{mo}}\nu(\mathrm{mo})^{-1}$ equals $1$ for both choices.
By the description of $\int\mathrm{D}W$ in 0.6\,(iii), $\int\mathrm{D}W$ is the same map; together with (2$'$), $\mathrm{d}\mu_I(W)$ is the same functional.
\end{itemize}
By (I.AP), (15.29) holds for $\Psi$ as well; so (1$'$)--(3$'$) hold with $\Psi$, $\bar\Psi$ in place of $W$, $\bar W$.

\emph{Levels $k\ge1$: kernels, forms, determinants.}
Let $O$ be a linear map of $\mathfrak{S}(1,\ell)$.
Its kernel is a function of $y,y'\in\mathbb{Z}^3$ with values in the $4\times4$ matrices such that $(Ov)(y)=\sum_{y'\in\mathrm{fdo}}O(y,y')\cdot v(y')$ for all $y\in\mathbb{Z}^3$ and all $v\in\mathfrak{S}(1,\ell)$ (Section~I.1.6), and
\begin{equation}\tag{15.30}
O(y,y'+\ell e_\mu)=-O(y,y'),\qquad O(y+\ell e_\mu,y')=-O(y,y')\qquad(\mu=1,2,3).
\end{equation}
The first relation is part of Section~I.1.6.
For the second, observe that the values $v_{\beta'}(y')$, $y'\in\mathrm{fdo}$, are arbitrary complex numbers (n2).
So the numbers $O(y,y')_{\beta\beta'}$, $y'\in\mathrm{fdo}$, are determined by the linear functional $v\mapsto(Ov)_\beta(y)$: they are its coefficients.
Since $Ov$ is antiperiodic, $(Ov)(y+\ell e_\mu)=-(Ov)(y)$ for every $v$.
Hence $O(y+\ell e_\mu,y')=-O(y,y')$ for $y'\in\mathrm{fdo}$, and then, by the first relation, for all $y'\in\mathbb{Z}^3$.

By the first relation of (15.30), $O(y,y')\cdot v(y')$ does not change under $y'\mapsto y'+\ell e_\mu$.
Hence $(Ov)(y)=\sum_{y'}O(y,y')\cdot v(y')$ holds with one and the same function $O(y,y')$ for every fundamental domain of summation: the kernel does not depend on the choice.
By (15.30) and (I.AP), $\bar\Psi(y)\cdot O(y,y')\cdot\Psi(y')$ does not change under $y\mapsto y+\ell e_\mu$ or under $y'\mapsto y'+\ell e_\mu$.
So
\[
\langle\bar\Psi,O\Psi\rangle=\sum_{y,y'}\bar\Psi(y)\cdot O(y,y')\cdot\Psi(y')
\]
is the same element for both choices.
The number $\det O$ is the determinant of a linear map of the vector space $\mathfrak{S}(1,\ell)$; its definition involves no fundamental domain, and by (n2) it is the determinant of the corresponding map of the functions on $\mathrm{Mo}$ for each of the two choices.
With (3$'$) it follows that
\[
\int\mathrm{D}\Psi\,\exp\left(-\langle\bar\Psi,O\Psi\rangle\right)\cdot F
\]
and, for invertible $\mathrm{Cov}$,
\[
\int\mathrm{d}\mu_{\mathrm{Cov}}(\Psi)\,F=\det\mathrm{Cov}\cdot\int\mathrm{D}\Psi\,\exp\left(-\langle\bar\Psi,\mathrm{Cov}^{-1}\Psi\rangle\right)\cdot F
\]
do not depend on the choice.

\emph{Levels $k\ge1$: arguments on the covering lattice.}
For $y,y'\in\mathrm{fdo}$ and $z,z'\in\mathbb{Z}^3$, by (I.AP) and (n2),
\begin{equation}\tag{15.31}
\begin{split}
\int\mathrm{d}\mu_I\,W_\beta(y+\ell z)\cdot\bar W_{\beta'}(y'+\ell z')&=(-1)^{z_1+z_2+z_3+z'_1+z'_2+z'_3}\\
&\qquad\cdot\mathbf{1}[y=y',\,\beta=\beta'].
\end{split}
\end{equation}
The number $1$ stands where (15.28) has the number $48$, because $\nu=1$.
The right side is the kernel of the identity of $\mathfrak{S}(1,\ell)$, which is $\mathbf{1}[y=y']\cdot1_4$ for $y,y'\in\mathrm{fdo}$ (0.5\,(i) with $\nu=1$), continued by (15.30).
The identity
\[
\int\mathrm{d}\mu_{\mathrm{Cov}}\,\Psi(y)\cdot\bar\Psi(y')^T=\mathrm{Cov}(y,y')
\]
holds for all $y,y'\in\mathbb{Z}^3$: for $y,y'\in\mathrm{fdo}$ it is (n2), and both sides change their sign under $y\mapsto y+\ell e_\mu$ and under $y'\mapsto y'+\ell e_\mu$, the left side by (I.AP), the right side by (15.30).

\emph{$\psi$ and $\chi$, every level $k\ge0$.}
By (I.AP) a change of the fundamental domain of $\ell\mathbb{Z}^3$ in $\mathbb{L}_k$ multiplies every generator $\psi(\xi)$, $\chi(\boldsymbol{\zeta})$ by $1$ or by $-1$ (for $\chi(\boldsymbol{\zeta})$ both points of the pair are translated by one and the same vector of $\ell\mathbb{Z}^3$).
So an ordered monomial of one choice is $\pm$ an ordered monomial of the other choice (0.1\,(i)), and the moduli of the coefficients and the norms of (n1) are the same for all choices.

\textbf{(n6) The rows.}
For the index sets of (n1)--(n3), the operator $H^{(2)}$ of (l)\,(iv) and the operators $O_{2j}$ of an application of (d) in which $\chi$ is substituted map into the functions on $\{\boldsymbol{\zeta}\}$.
By (D4) the norm of such an operator is the supremum of the row sums over \emph{all} $\boldsymbol{\zeta}=(x,x',\beta,\omega)$, $(x,x')\in\mathrm{Pz}_k$.
By Section~I.1.7, $\mathrm{Pz}_k$ contains, for every photon site $c$, the $48\cdot47$ ordered pairs $\left((c,g),(c,g')\right)$, $g\ne g'$, inside the cluster of $c$ (for them $d_X=\eta_k\le M$).
The proof of (d) uses the bound on the row $t$ for every generator $\phi_i(t)$ that occurs in $F$, and an element of $\mathrm{Gr}[\psi,\chi]$ may contain every $\chi(\boldsymbol{\zeta})$.
So the hypothesis of (d) and of (l)\,(iv) is a bound on all rows, the rows at the pairs inside a cluster included; fewer rows suffice only in the form (d$'$).
For $\chi\mapsto\delta_{\alpha,B}H\Psi$ the row at $(x,x',\beta,0)$ is the linear functional
\[
v\mapsto d_X(x,x')^{-\alpha}\cdot\left(U_B(x,x')\cdot(Hv)_\beta(x')-(Hv)_\beta(x)\right)
\]
by (I.1.7); at a pair inside a cluster it is the functional $v\mapsto\eta_k^{-\alpha}\cdot\left((Hv)_\beta(c,g')-(Hv)_\beta(c,g)\right)$, without phase.
(Let $\bar H$ be a second linear map with the source and the target of $H$, that is, of $\mathfrak{S}(1,\ell)$ (for $k=0$: of $\mathfrak{S}^{\mathrm{ch}}(1,\ell)$) to $\mathfrak{S}^{\mathrm{ch}}(\eta_k,\ell)$;
like $H$ it is arbitrary, it is not determined by $H$, and it is the map inserted for the barred field.
In the applications, for a connection $A$, one has $H=H_{k}(A)$, $\bar H=\bar{H}_{k}(A)=\mathsf{C}^{-1}\cdot H_{k}(-A)\cdot\mathsf{C}$ and $B=A$ (Sections~I.3.5 and~I.5.4); at $k=0$ both maps are $1$ (Section~I.3.9).
For $\bar\chi\mapsto\delta_{\alpha,-B}\bar H\bar\Psi$ the row at $(x,x',\beta,1)$ is the linear functional
\[
v\mapsto d_X(x,x')^{-\alpha}\cdot\left(U_{-B}(x,x')\cdot(\bar Hv)_\beta(x')-(\bar Hv)_\beta(x)\right),
\]
by (I.1.7) for barred fields; at a pair inside a cluster it is $v\mapsto\eta_k^{-\alpha}\cdot\left((\bar Hv)_\beta(c,g')-(\bar Hv)_\beta(c,g)\right)$.
A row with $\omega=0$ involves only the columns with $\omega=0$, and a row with $\omega=1$ only those with $\omega=1$.)
At level $0$ one has $\eta_0=1$, so that the factor $\eta_k^{-\alpha}$ equals $1$; for $H=1$ (the operator $H_0(A)$ of Section~I.3.9) the row at a pair inside a cluster is $v\mapsto v_\beta(c,g')-v_\beta(c,g)$.
Its kernel with respect to $\nu=1/48$ has the two entries $48$ and $-48$, and its row sum is
\[
48\cdot\frac{1}{48}+48\cdot\frac{1}{48}=2:
\]
the factor $48$ stands in the kernel and not in the row sum.
By (c), $|\delta_{\alpha,B}H|_{1,\infty}\le\mathrm{Cr}$ is equivalent to
\[
\left|(\delta_{\alpha,B}Hv)(x,x')\right|_c\le\mathrm{Cr}\cdot|v|_\infty\qquad\text{for all $v$ and all pairs of $\mathrm{Pz}_k$};
\]
here $|f|_\infty$ of (D4), over an index set of the form $\{(t,\beta)\}$, $\beta\in\{1,\dots,4\}$, is $\sup_t\max_\beta|f_\beta(t)|$, the sup norm formed with the norm $|\cdot|_c$ of Section~I.1.5, and $t$ is $y$, resp.\ $x$, for $v$, and the pair $(x,x')$ for $\delta_{\alpha,B}Hv$.

\textbf{(n7) The two assertions of (G2.r7) in Section~G2.0.5.}
\emph{The sum over the bidegrees.}
Let $F\in\mathrm{Gr}[\psi,\chi]$.
Every ordered monomial $\psi(I_1)\cdot\chi(I_2)$ has exactly one bidegree $(\#|I_1|,\#|I_2|)$, and $F^{[n,m]}$ is the sum of the terms $c_F(I_1,I_2)\cdot\psi(I_1)\cdot\chi(I_2)$ of $F$ with $(\#|I_1|,\#|I_2|)=(n,m)$.
So the sum (I.5.1) for $F$ is the sum over $(n,m)$ of the sums (I.5.1) for the $F^{[n,m]}$:
\[
\|F\|_{\mathsf{w}}=\sum_{n,m}\left\|F^{[n,m]}\right\|_{\mathsf{w}}.
\]
\emph{The display (G2.0.1).}
Let $E\in\mathrm{Gr}[\psi,\chi]$ be balanced in the sense of (G2.r8) in Section~G2.0.5, and let $E_{20}(\xi_1,\xi_2)$ be the antisymmetric kernel of $E^{[2,0]}$ given by (a), so that
\[
E^{[2,0]}=\frac{1}{2}\int\mathrm{d}\xi_1\,\mathrm{d}\xi_2\,E_{20}(\xi_1,\xi_2)\cdot\psi(\xi_1)\cdot\psi(\xi_2).
\]
By (a), $E_{20}(\xi_1,\xi_2)\cdot\nu_1(\xi_1)\cdot\nu_1(\xi_2)$ is, for $\xi_1<\xi_2$, the coefficient of $\psi(\xi_1)\cdot\psi(\xi_2)$ in $E$; it vanishes if $\omega_1=\omega_2$, because $E$ is balanced.
Hence only $(\omega_1,\omega_2)=(1,0)$ and $(\omega_1,\omega_2)=(0,1)$ contribute.
Exchanging the names $\xi_1$ and $\xi_2$ carries the second part into the first, since
\[
E_{20}(\xi_2,\xi_1)\cdot\psi(\xi_2)\cdot\psi(\xi_1)=E_{20}(\xi_1,\xi_2)\cdot\psi(\xi_1)\cdot\psi(\xi_2)
\]
(two changes of sign).
With $\sum_{\xi:\,\omega\text{ fixed}}\nu_1(\xi)\,(\cdot)=\int\mathrm{d}x\,\sum_\beta(\cdot)$ one obtains
\begin{align*}
E^{[2,0]}&=\int\mathrm{d}x\int\mathrm{d}x'\,\sum_{\beta,\beta'}E_{20}\left((x,\beta,1),(x',\beta',0)\right)\cdot\bar\psi_\beta(x)\cdot\psi_{\beta'}(x')\\
&=\int\mathrm{d}x\int\mathrm{d}x'\,\bar\psi(x)^T\cdot E_{20}(x,x')\cdot\psi(x'),\\
\left\|E^{[2,0]}\right\|_{\mathsf{w}}&=\frac{\mathsf{w}_1^2}{2}\int\mathrm{d}\xi_1\,\mathrm{d}\xi_2\,|E_{20}|=\frac{\mathsf{w}_1^2}{2}\cdot\|E_{20}\|\\
&=\frac{\mathsf{w}_1^2}{2}\cdot2\sum_{\beta,\beta'}\int\mathrm{d}x\int\mathrm{d}x'\,\left|E_{20}(x,x')_{\beta\beta'}\right|;
\end{align*}
the last expression for $E^{[2,0]}$ is written in the matrix notation of Section~I.5.1, and the first equality for the norm is the case of equality in (a).

\textbf{(n8) The number $7$ (the display of (f) with $7$ in place of $\mathsf{w}_W$ at every level).}
By (k), $\|F\|_{(\mathsf{w}_1,\dots,\mathsf{w}_p,\mathsf{w}_W)}\le\|F\|_{(\mathsf{w}_1,\dots,\mathsf{w}_p,\mathsf{w}'_W)}$ for every $\mathsf{w}'_W\ge\mathsf{w}_W$.
Hence, by (15.14),
\begin{align*}
\left\|\int\mathrm{d}\mu_I(W)\,F\right\|_{(\mathsf{w}_1,\dots,\mathsf{w}_p)}&\le\sum_{I,J}|c_F(I,J)|\cdot\mathsf{w}^I\cdot\nu^{-\frac{1}{2}}(J)\\
&\le\|F\|_{(\mathsf{w}_1,\dots,\mathsf{w}_p,7)}
\end{align*}
at level $0$, where $\sqrt{48}\le7$, and at the levels $k\ge1$, where $1\le7$.
The assertion of (f) that $\mathsf{w}_W$ cannot be lowered is excepted.
\end{proof}

\subsubsection*{Uniformity}

The clauses (a)--(m) contain no constant of Section~I.8: the index sets, the weights $\nu_i$ and $\nu$, the weight vectors and the maps are arbitrary data, and apart from the norms of the data the bounds contain only the explicit factors
\[
\frac{1}{\mathrm{rad}-1},\qquad\frac{1}{\mathrm{rad}\,(\mathrm{rad}-1)},\qquad\prod_i\left(\frac{\mathsf{w}'_i}{\mathsf{w}_i}\right)^{r_i}.
\]
The numbers of (n) are the following.
First, the weights $\eta_k^3/48$ of the indices $\xi$ and $(\eta_k^3/48)^2$ of the indices $\boldsymbol{\zeta}$, which depend on $k$; they stand in the integrals by which kernels are written, and they enter no norm (I.5.1), no factor of a bound and no moment (see (a), the remark following the proof of (c), and (n1)).
Second, $\nu=1$, $\mathsf{w}_W=1$ and the number $1$ of the two-point function for $k\ge1$.
Third, $\nu=1/48$, $\mathsf{w}_W=\sqrt{48}$ and the number $48$ of the two-point function for $k=0$.
Fourth, the number $7$, at every level; see (n8).
By (15.11) and (15.25) the moments are, for paired $J$, equal to $\mathrm{sn}(J)$ for $k\ge1$ and to $\mathrm{sn}(J)\cdot48^{\frac{\#|J|}{2}}$ for $k=0$, and they vanish for the other $J$; they depend on the order of the elements of $J$ among themselves and on $\#|J|$.
By (f$'$) the integral $\int\mathrm{d}\mu_I(W)\,F$ does not change if modes that occur in no monomial of $F$ are added to $\mathrm{Mo}$ or removed from it.
Consequently the factors in the bounds of (a)--(m), the numbers $\nu(\mathrm{mo})$, $\mathsf{w}_W$ and $7$, the moments $m(J)$, the numbers $1$, resp.\ $48$, of the two-point functions and, for a fixed $F$, the element $\int\mathrm{d}\mu_I(W)\,F$ depend neither on the level within $k\ge1$, nor on the side $\ell$ (the torus), nor on $N$, the charge, a mass, a twist or a connection.

\subsection*{Conditions on the constants}

The assertions of this chapter hold under the conditions on the constants of Section~I.8 and of Proposition~14.2.
The clauses (a)--(m) are free of the constants of Section~I.8.2 and of Proposition~14.2, and they hold for all values of these constants.

In (n) the constants $L$ (in $\eta_k=L^{-k}$), $M$ (in $\mathrm{Pz}_k$) and $\alpha$ (in $\delta_{\alpha,B}$) enter through the definitions of Part~I recalled there.
Among them, (n) requires only the inequality $\eta_k\le M$, in (n6), by which the pairs inside a cluster belong to $\mathrm{Pz}_k$ (Section~I.1.7).
It holds because $\eta_k\le1\le M$, by $L\ge5$ (Section~I.0.5; line (C3)) and $M\ge L^2$ (Section~I.0.5; line (C5)).

The bound $\sqrt{48}\le7$ of (n3) follows from $48\le49=7^2$ and holds for all values of the constants.
Through the bound $\mathsf{w}_W\le7$ it enters the constant $C_W:=28\,C_w^2\cdot c_\gamma$ of line (C6) and the condition (A6), $112\,C_w^2\cdot c_\gamma\cdot e_k^{2\varepsilon}\le3$, of line (C8); see (I.L0.19), Section~I.8.4 and Proposition~14.2\,(b).

\section*{Chapter 16. The group lemma, the low kernels, the extraction: proofs of Propositions 16.1--16.3}

\subsection*{0. Conventions of this chapter; arithmetic of the sides; four lemmas}

\subsubsection*{0.1. Notation}

Throughout the chapter $k\ge0$ is an integer, $\ell\in\mathrm{Cl}$, and $\eta:=\eta_k=L^{-k}$.
The symbol $\mathrm{pr}$ stands for each of the covering maps
\[
\mathbb{R}^3\to\mathbb{R}^3/\ell\mathbb{Z}^3,\qquad\mathbb{L}_k\to T_k,\qquad\mathbb{L}^{\mathrm{f}}_k\to T^{\mathrm{f}}_k,\qquad\mathrm{Pzt}_k\to\mathrm{Pz}_k,
\]
where on fermion sites $\mathrm{pr}(c,g):=(\mathrm{pr}\,c,g)$.
For $n\in\mathbb{Z}^3$ let $\boldsymbol{\chi}(n):=(-1)^{n_1+n_2+n_3}$ (Section~I.6.3).
If $\xi=(x,\beta,\omega)$ and $\boldsymbol{\zeta}=(x,x',\beta,\omega)$ are indices and $v$ is a vector, one sets
\[
\xi+v:=(x+v,\beta,\omega),\qquad\boldsymbol{\zeta}+v:=(x+v,x'+v,\beta,\omega).
\]

In this chapter the letter $\eta$ always denotes the mesh $\eta_k$ of the level $k$ (Section~I.1.1), as in \cite[Section~13.1.1]{QEDI}, and not the spacing $L^{-N}$ of Section~I.0.5.
In (a1.5) the letter $w$ denotes a point of the lattice $\mathrm{hm}\,\mathbb{Z}^3=L^{-i}\mathbb{Z}^3$, and in (a2.6) a site of the six paths.
The letters $p$, $p'$ denote points of $\mathbb{R}^3$; in the closed $M$-cube $\left\{p:d(p,z)\le\frac{M}{2}\right\}$ of the continuum torus, $p$ is a point of that torus.
The letter $y$ denotes a site, as in Proposition~16.1\,(a1).
The letter $t$ denotes a real number with $0\le t\le1$, the parameter of a segment (paragraph on nets below, and (a1.6)).
Plaquettes are written $(c;\mu\nu)$ ((b.1)), and the indices of the generators are written $\xi$, $\boldsymbol{\zeta}$, $u$.
The letters $\mathsf{Z}$, $\mathsf{Z}'$ denote finite subsets of $\mathbb{L}_k$, in particular lifts, as in Section~I.1.1 and in Proposition~16.1\,(a2), and $a'$ is the translation part of an element $(S,a')$ of $\mathbb{G}_{k+1}$, as in Proposition~16.1\,(d).

\paragraph{The element $\tilde{g}$.}
Let $\tilde{g}=(S,a)\in\mathbb{G}_k$, and let $r:=r(S)=\mathrm{diag}(r_1,r_2,r_3)$ with $r_\mu\in\{1,-1\}$, so that $r^{-1}=r$ (Section~I.4.1).
The translation part, a class in $M\mathbb{Z}^3/\ell\mathbb{Z}^3$, is represented by a vector of $M\mathbb{Z}^3$, which is again denoted by $a$ (Section~I.4.3).
The element $\tilde{g}$ acts on $\mathbb{R}^3$ by $p\mapsto rp+a$ and on $\mathbb{L}^{\mathrm{f}}_k$ by $(c,g)\mapsto(rc+a,rg)$; accordingly
\[
\tilde{g}^{-1}p=r(p-a),\qquad\tilde{g}^{-1}(c,g)=\left(r(c-a),rg\right).
\]
When $a$ is replaced by $a+\ell n$ with $n\in\mathbb{Z}^3$, the point $\tilde{g}p$ is replaced by $\tilde{g}p+\ell n$, and $\tilde{g}^{-1}p$ by $\tilde{g}^{-1}p-\ell\cdot rn$.

\paragraph{Cubes.}
An $M$-cube of $\mathbb{L}_k$ is a set
\[
\tilde{\square}=\left\{c\in\mathbb{L}_k:|c-\tilde{z}|<\frac{M}{2}\right\},\qquad\tilde{z}\in M\mathbb{Z}^3,
\]
whose centre is $\tilde{z}$; the closed $M$-cube of $\mathbb{R}^3$ with the same centre is $\left\{p\in\mathbb{R}^3:|p-\tilde{z}|\le\frac{M}{2}\right\}$.
An $M$-cube of $T_k$ is a set $\square=\left\{c\in T_k:d(c,z)<\frac{M}{2}\right\}$ with $z\in M\mathbb{Z}^3/\ell\mathbb{Z}^3$, and the closed $M$-cube of the continuum torus with the same centre is $\left\{p:d(p,z)\le\frac{M}{2}\right\}$.
Two $M$-cubes have a common face if their centres differ by $M\cdot e_\mu$ or by $-M\cdot e_\mu$ for one direction $\mu$, on the torus modulo $\ell\mathbb{Z}^3$ (Section~I.1.1).
A union of $M$-cubes is identified with the set of the sites that lie in it (G2.r1).
Let $\tilde{X}$ be a non-empty finite union of $M$-cubes of $\mathbb{L}_k$ that is connected through faces.
Then $d_M(\tilde{X})$ is the number $d_M$ of \cite[Section~13.1.7]{QEDI}, formed with closed $M$-cubes of $\mathbb{R}^3$ and with nets in $\mathbb{R}^3$, as in the assertion (t10) of \cite[Theorem~13.2]{QEDI}:
$M\cdot d_M(\tilde{X})$ is the infimum of the lengths of the nets in $\mathbb{R}^3$ that meet every closed $M$-cube of $\mathbb{R}^3$ whose centre is the centre of a cube of $\tilde{X}$;
here a net in $\mathbb{R}^3$ is a connected finite union of closed straight segments of $\mathbb{R}^3$, a point counting as a segment, and its length is the sum of the Euclidean lengths of these segments.
Further, $|\tilde{X}|_M$ denotes the number of the cubes of $\tilde{X}$, and $\operatorname{Vol}(\tilde{X}):=M^3\cdot|\tilde{X}|_M$.

\paragraph{Nets.}
In the definition of $d_M$ on the torus \cite[Section~13.1.7]{QEDI} the segments of the torus, their lengths and the length of a net are those of Section~I.1.1:
a closed straight segment of $\mathbb{R}^3/\ell\mathbb{Z}^3$ is the image of a closed straight segment of $\mathbb{R}^3$ and has the Euclidean length of that segment, and the length of a net refers to a writing of the net as a finite union of segments.
For $p,p'\in\mathbb{R}^3$,
\[
\mathrm{seg}=\left\{p+t(p'-p):0\le t\le1\right\}
\]
denotes a closed straight segment of $\mathbb{R}^3$, of length $|p-p'|_2$, and $\mathrm{pr}(\mathrm{seg})$ denotes its image in the torus.
The infimum in the definition of $d_M$ extends over the nets together with their writings as finite unions of segments.

\paragraph{Enlargements.}
For an integer $i\ge0$ and a union $X$ of $M$-cubes, $X^{\sim i}$ is the set of Section~I.1.1: the union of the $M$-cubes whose centre has sup-distance at most $i\cdot M$ from the centre of an $M$-cube of $X$.
For an $M$-cube $\square$ of $T_k$ with the centre $z$ one has
\begin{equation}\tag{enl}
\square^{\sim i}=\left\{c\in T_k:d(c,z)<\left(i+\frac{1}{2}\right)\cdot M\right\}.
\end{equation}

\begin{proof}[Proof of (enl)]
Let $c$ lie in the $M$-cube with the centre $z'$, and let $d(z',z)\le i\cdot M$.
Then $d(c,z)\le d(c,z')+d(z',z)<\frac{M}{2}+i\cdot M$.
Conversely, let $d(c,z)<\left(i+\frac{1}{2}\right)\cdot M$.
First, $c$ lies in an $M$-cube.
Indeed, let $\tilde{c}$ be a point of $\mathbb{L}_k$ over $c$.
For every direction $\mu$ some multiple $\tilde{z}'_\mu$ of $M$ satisfies $|\tilde{c}_\mu-\tilde{z}'_\mu|\le\frac{M}{2}$, and equality is impossible by (Ar2) of Section~0.2 below.
Hence $|\tilde{c}-\tilde{z}'|<\frac{M}{2}$ for the point $\tilde{z}':=(\tilde{z}'_1,\tilde{z}'_2,\tilde{z}'_3)$ of $M\mathbb{Z}^3$.
The point $z':=\mathrm{pr}\,\tilde{z}'$ is the centre of this cube, and $d(c,z')<\frac{M}{2}$.
Consequently $d(z',z)\le d(z',c)+d(c,z)<(i+1)\cdot M$.
Since $\ell\mathbb{Z}^3\subset M\mathbb{Z}^3$ by (Ar1) of Section~0.2, every vector in the class of $z'-z$ modulo $\ell\mathbb{Z}^3$ belongs to $M\mathbb{Z}^3$, and the sup norm of a vector of $M\mathbb{Z}^3$ is an integral multiple of $M$.
So $d(z',z)$, the least of these sup norms, is an integral multiple of $M$, and therefore $d(z',z)\le i\cdot M$.
Thus the $M$-cube that contains $c$ is one of the cubes of $\square^{\sim i}$.
\end{proof}

In particular,
\[
\square^{\dagger}=\square^{\sim L}=\left\{c:d(c,z)<\left(L+\frac{1}{2}\right)\cdot M\right\}
\]
has the side $(2L+1)\cdot M$, and $\square^{\sim5}=\left\{c:d(c,z)<\left(5+\frac{1}{2}\right)\cdot M\right\}$ has the side $11M$, in accordance with \cite[Definition~5.7]{QEDI}.

\paragraph{Integrals over subsets of the covering lattice.}
For a finite subset $\mathsf{Z}$ of $\mathbb{L}_k$ and a function $f$ on $\mathbb{L}^{\mathrm{f}}_k$ with values in a vector space, the integral over $\mathsf{Z}$ is that of Proposition~16.1\,(a2):
\[
\int_{\mathsf{Z}}\mathrm{d}x\,f(x)=\eta^3\cdot\frac{1}{48}\sum_{c\in\mathsf{Z}}\sum_{g\in O_{\mathrm{cub}}}f(c,g).
\]
If $\mathsf{Z}$ is a lift of a set $Z$ of sites of $T_k$ and $f$ is $\ell\mathbb{Z}^3$-periodic, this is the integral $\int_Z\mathrm{d}x\,f(x)$ of Section~I.1.6.

\paragraph{Labels.}
The definitions (G2.r1)--(G2.r8) and the display (G2.0.1) are those of Section~G2.0.5, and the settings (G2.S-1), (G2.S-2) are those of Section~G2.0.4.
The displays (2.1) and (2.2) are those of Proposition~16.1\,(b) and of Proposition~16.3\,(b).
The labels (I.AP), (I.Sch), (I.ke1) and (I.c2)--(I.c6) refer to Sections~I.3.6, I.4.1, I.1.6 and I.2.1, respectively; (k1)--(k7) are the clauses of the definition in Section~I.5.4; (C0)--(C5) are lines of the ordered list of Section~I.8.2; and (A8) is the condition of that name in Section~I.8.4.
Finally, (t0), (t2), (t7), (t10) denote the assertions of \cite[Theorem~13.2]{QEDI} that carry these marks.

\subsubsection*{0.2. Arithmetic}

\begin{itemize}
\item[(Ar1)]
$L$ is an odd integer, and $M$ is a power of $L$ with $M\ge L^2$ (Section~I.8.2, lines (C3) and (C5); rows (L5) and (M1) of Proposition~14.2\,(b)); moreover $\frac{\ell}{L\cdot M}$ is an odd integer $\ge13$.
Consequently, in the list
\[
\ell\mathbb{Z}^3,\qquad LM\mathbb{Z}^3,\qquad M\mathbb{Z}^3,\qquad L\mathbb{Z}^3,\qquad\mathbb{Z}^3,\qquad\mathbb{L}_k
\]
each group is contained in the next one; in particular $M\mathbb{Z}^3\subset L^{-i}\mathbb{Z}^3$ for $-1\le i\le k$.
\item[(Ar2)]
Let $\mathrm{sid}$ be one of the numbers $L^{-i}$ with $-1\le i\le k$, $M$, $L\cdot M$, or $(2i+1)\cdot M$ with an integer $i\ge0$.
Then $\frac{\mathrm{sid}}{\eta}=\mathrm{sid}\cdot L^k$ is an odd integer.
Hence, for every $z\in\mathbb{L}_k$,
\[
\left\{c\in\mathbb{L}_k:|c-z|<\frac{\mathrm{sid}}{2}\right\}=\left\{c\in\mathbb{L}_k:|c-z|\le\frac{\mathrm{sid}-\eta}{2}\right\},
\]
no site $c$ of $\mathbb{L}_k$ satisfies $|c_\mu-z_\mu|=\frac{\mathrm{sid}}{2}$, this set is a lattice box with $\frac{\mathrm{sid}}{\eta}$ sites in each direction, and its sup-diameter is $\mathrm{sid}-\eta$.
\item[(Ar3)]
As $\eta\le1\le M\le L\cdot M$ and $13L\cdot M\le\ell$,
\begin{gather*}
2M<\ell,\qquad(2L+1)\cdot M+\eta\le4L\cdot M<\ell,\qquad11M+\eta\le12M<\ell,\\
2M\cdot(L+2)\le6L\cdot M<\ell,\qquad M\cdot(L+2)+M+2\eta\le6L\cdot M<\ell.
\end{gather*}
\end{itemize}

\subsubsection*{0.3. Lifts (Section~I.1.1)}

\begin{named}{Lemma 0.1}
Let $Z$ be a set of sites of $T_k$, and let $\mathsf{Z}$ be a lift of $Z$, that is, a subset of $\mathbb{L}_k$ with $\mathrm{pr}(\mathsf{Z})=Z$ and $\operatorname{sup-diameter}(\mathsf{Z})+\eta<\ell$.
\begin{itemize}
\item[(Lft1)] The map $\mathrm{pr}$ is injective on $\mathsf{Z}$.
\item[(Lft2)] Let $\mathrm{sd}\ge0$ be such that $\operatorname{sup-diameter}(\mathsf{Z})+\mathrm{sd}<\ell$.
If $\tilde{c},\tilde{c}'\in\mathsf{Z}$ and $d(\mathrm{pr}\,\tilde{c},\mathrm{pr}\,\tilde{c}')\le\mathrm{sd}$, then $\tilde{c}-\tilde{c}'$ is the only vector of sup norm at most $\mathrm{sd}$ in its class modulo $\ell\mathbb{Z}^3$, and $|\tilde{c}-\tilde{c}'|=d(\mathrm{pr}\,\tilde{c},\mathrm{pr}\,\tilde{c}')$.
\item[(Lft3)] If $Z$ is connected through nearest neighbours, then every lift of $Z$ has the form $\mathsf{Z}+\ell n$ with some $n\in\mathbb{Z}^3$.
\item[(Lft4)] If $Z$ is a union of $M$-cubes of $T_k$, then $\mathsf{Z}$ is a union of $M$-cubes of $\mathbb{L}_k$, exactly one of them lying over each $M$-cube of $Z$; the centre of the cube over $\square$ lies over the centre of $\square$; and two of these cubes have a common face if and only if their images have a common face.
\end{itemize}
\end{named}

\begin{proof}
(Lft1)
Two distinct points of $\mathsf{Z}$ with the same image differ by a vector $\ell n$ with $n\ne0$.
The sup norm of such a vector is at least $\ell$, which exceeds the sup-diameter of $\mathsf{Z}$.

(Lft2)
The number $d(\mathrm{pr}\,\tilde{c},\mathrm{pr}\,\tilde{c}')$ is the least sup norm of a vector in the class of $\tilde{c}-\tilde{c}'$; so this class contains a vector $v$ with $|v|\le\mathrm{sd}$.
Write $\tilde{c}-\tilde{c}'=v+\ell n$.
If $n\ne0$, then
\[
|\tilde{c}-\tilde{c}'|\ge\ell-|v|\ge\ell-\mathrm{sd}>\operatorname{sup-diameter}(\mathsf{Z}),
\]
which is impossible.
Therefore $v=\tilde{c}-\tilde{c}'$.
This conclusion holds for every such $v$, in particular for a vector $v$ of the class with $|v|=d(\mathrm{pr}\,\tilde{c},\mathrm{pr}\,\tilde{c}')$.

(Lft3)
Let $\mathsf{Z}'$ be another lift of $Z$.
For $c\in Z$ let $\tilde{c}\in\mathsf{Z}$ and $\tilde{c}'\in\mathsf{Z}'$ be the points over $c$, which are unique by (Lft1), and write $\tilde{c}'-\tilde{c}=\ell n(c)$.
Let $c^a$, $c^b$ be nearest neighbours in $Z$.
Applied with $\mathrm{sd}=\eta$ to $\mathsf{Z}$ and to $\mathsf{Z}'$, (Lft2) shows that $\tilde{c}^a-\tilde{c}^b$ and $\tilde{c}^{\prime a}-\tilde{c}^{\prime b}$ both coincide with the only vector of sup norm at most $\eta$ in the class of $c^a-c^b$.
Hence $n(c^a)=n(c^b)$, and $n$ is constant on $Z$.

(Lft4)
Let $\square$ be an $M$-cube of $Z$ with the centre $z$.
By (Ar1), $z$ is a site of $\square$.
Let $\tilde{z}$ be the point of $\mathsf{Z}$ over $z$; it belongs to $M\mathbb{Z}^3$, because $\ell\mathbb{Z}^3\subset M\mathbb{Z}^3$.
The points of $\mathsf{Z}$ over $\square$ form a lift of $\square$.
The $M$-cube $\tilde{\square}$ of $\mathbb{L}_k$ with the centre $\tilde{z}$ is mapped onto $\square$, since $d(c,z)$ is the least of the numbers $|\tilde{c}-\tilde{z}|$, $\tilde{c}$ ranging over the points of $\mathbb{L}_k$ over $c$; by (Ar2) its sup-diameter is $M-\eta$, and $M<\ell$; so $\tilde{\square}$ is a lift of $\square$ as well.
Being a lattice box, $\square$ is connected through nearest neighbours.
By (Lft3) the two lifts of $\square$ differ by a vector $\ell n$.
Both contain $\tilde{z}$, and by (Lft1) neither contains another point over $z$; hence $n=0$.

Faces.
Let $\square_1$, $\square_2$ be cubes of $Z$ with the centres $z_1$ and $z_2=z_1+M\cdot e_\mu$ (modulo $\ell\mathbb{Z}^3$).
By (Ar2), $\frac{M-\eta}{2}\in\eta\mathbb{Z}$, so that
\[
c^a:=z_1+\frac{M-\eta}{2}\cdot e_\mu\in\square_1,\qquad c^b:=c^a+\eta\cdot e_\mu=z_2-\frac{M-\eta}{2}\cdot e_\mu\in\square_2
\]
are nearest neighbours in $Z$.
By the first part of the proof, the points of $\mathsf{Z}$ over them are $\tilde{z}_1+\frac{M-\eta}{2}\cdot e_\mu$ and $\tilde{z}_2-\frac{M-\eta}{2}\cdot e_\mu$, and by (Lft2) with $\mathrm{sd}=\eta$ the second minus the first equals $\eta\cdot e_\mu$.
Hence $\tilde{z}_2-\tilde{z}_1=M\cdot e_\mu$.
Conversely, if $\tilde{z}_2-\tilde{z}_1=\pm M\cdot e_\mu$, then $z_2-z_1=\pm M\cdot e_\mu$ modulo $\ell\mathbb{Z}^3$.
\end{proof}

\begin{named}{Lemma 0.2}[Polymers and their lifts]
\begin{itemize}
\item[(a)] A polymer of $T_k$, and likewise a non-empty finite union of $M$-cubes of $\mathbb{L}_k$ that is connected through faces, is connected through nearest neighbours as a set of sites.
\item[(b)] Let $X$ be a polymer of $T_k$ with a lift $\tilde{X}$.
Then every lift of $X$ has the form $\tilde{X}+\ell n$; $\tilde{X}$ is a non-empty finite union of $M$-cubes of $\mathbb{L}_k$ that is connected through faces; and $|\tilde{X}|_M=|X|_M$, $\operatorname{Vol}(\tilde{X})=\operatorname{Vol}(X)$.
\item[(c)] Let $X$ be a small polymer.
Then $X$ has a lift $\tilde{X}$, and
\[
\operatorname{sup-diameter}(\tilde{X})\le M\cdot(L+2),\qquad d_M(\tilde{X})=d_M(X);
\]
moreover $\tilde{X}$ is the set of the sites of $\mathbb{L}_k$ in the union of $M$-cubes of $\mathbb{R}^3$ that (t10) attaches to $X$.
\item[(d)] Let $X$ be a polymer of $T_k$ with a lift $\tilde{X}$.
Then $X$ is small if and only if $d_M(\tilde{X})\le L$.
\end{itemize}
\end{named}

In particular the facts stated in (G2.r3) hold: they are contained in (b), (c), (d) together with the fifth inequality of (Ar3).

\begin{proof}
(a)
Every $M$-cube is a lattice box.
If two $M$-cubes have a common face, then the sites $c^a$, $c^b$ of the proof of (Lft4) are nearest neighbours, and each of the two cubes contains one of them.

(b)
By (a), $X$ is connected through nearest neighbours; so every lift of $X$ has the form $\tilde{X}+\ell n$, by (Lft3).
By (Lft4), $\tilde{X}$ is a union of $M$-cubes of $\mathbb{L}_k$, exactly one of them lying over each $M$-cube of $X$, and two of them have a common face if and only if their images have one.
Consequently $\tilde{X}$ is non-empty, finite and connected through faces, $|\tilde{X}|_M=|X|_M$, and $\operatorname{Vol}(\tilde{X})=M^3\cdot|\tilde{X}|_M=\operatorname{Vol}(X)$.

(c)
As $d_M(X)\le L$, (Ar3) gives
\[
2M\cdot\left(d_M(X)+2\right)\le2M\cdot(L+2)<\ell.
\]
By (t10), $X=\mathrm{pr}(\mathrm{Xr})$ for a finite union $\mathrm{Xr}$ of $M$-cubes of $\mathbb{R}^3$ (centred at points of $M\mathbb{Z}^3$) that is connected through faces; $\mathrm{pr}$ is injective on $\mathrm{Xr}$, and $d_M(\mathrm{Xr})=d_M(X)$.
By (t7) in $\mathbb{R}^3$, any two points of the closed cubes of $\mathrm{Xr}$ have sup-distance at most $M\cdot\left(d_M(\mathrm{Xr})+2\right)\le M\cdot(L+2)$.
Let $\tilde{X}$ be the set of the sites of $\mathbb{L}_k$ that lie in $\mathrm{Xr}$.
By (Ar2), $\tilde{X}$ is the union of the $M$-cubes of $\mathbb{L}_k$ whose centres are the centres of the cubes of $\mathrm{Xr}$; so the closed $M$-cubes of $\mathbb{R}^3$ that enter the definition of $d_M(\tilde{X})$ in Section~0.1 are the closed cubes of $\mathrm{Xr}$, and $d_M(\tilde{X})=d_M(\mathrm{Xr})=d_M(X)$.
By the fifth inequality of (Ar3),
\[
\operatorname{sup-diameter}(\tilde{X})+\eta\le M\cdot(L+2)+\eta<\ell.
\]
The map $\mathrm{pr}$ sends the sites of $\tilde{X}$ to sites of $X$.
Conversely, a site $c$ of $X$ equals $\mathrm{pr}(p)$ for some $p\in\mathrm{Xr}$, and $p$ lies in the class $c$, which by (Ar1) is a subset of $\mathbb{L}_k$; thus $p\in\tilde{X}$.
Hence $\mathrm{pr}(\tilde{X})=X$, and $\tilde{X}$ is a lift of $X$.

(d)
By (b), $\tilde{X}$ is a non-empty finite union of $M$-cubes of $\mathbb{L}_k$ that is connected through faces, so that $d_M(\tilde{X})$ is defined.
Suppose $d_M(\tilde{X})\le L$.
Then $2M\cdot\left(d_M(\tilde{X})+2\right)<\ell$ by (Ar3), and the converse part of (t10), applied to the non-empty union $\tilde{X}$, yields
\[
d_M(X)=d_M(\mathrm{pr}\,\tilde{X})=d_M(\tilde{X})\le L.
\]
Suppose now that $X$ is small.
Then (c) provides a lift $\tilde{X}'$ with $d_M(\tilde{X}')=d_M(X)\le L$, and $\tilde{X}=\tilde{X}'+\ell n$ by (b).
A translation of $\mathbb{R}^3$ maps nets to nets of the same length and the closed cubes of $\tilde{X}'$ onto the closed cubes of $\tilde{X}$.
Therefore $d_M(\tilde{X})=d_M(\tilde{X}')\le L$.
\end{proof}

\subsubsection*{0.4. Generators, kernels, quadratic elements (Sections~I.3.6, I.5.1; Proposition~15\,(a))}

Let $\mathrm{Fdm}$ be a fundamental domain of $\ell\mathbb{Z}^3$ in $\mathbb{L}_k$.
The independent generators of $\mathrm{Gr}[\psi,\chi]$ are the elements $\psi(\xi)$, $\xi=(x,\beta,\omega)$ with $x=(c,g)$ and $c\in\mathrm{Fdm}$, and the elements $\chi(\boldsymbol{\zeta})$, $\boldsymbol{\zeta}=(x,x',\beta,\omega)$ with $(x,x')\in\mathrm{Pzt}_k$ and $c\in\mathrm{Fdm}$.
They are the two families of Proposition~15\,(D1), with the weights
\[
\nu(\xi)=\frac{\eta^3}{48},\qquad\nu(\boldsymbol{\zeta})=\left(\frac{\eta^3}{48}\right)^2,
\]
and with arbitrary total orders of the two index sets.
At the remaining indices of the covering index sets the generators are given by (I.AP):
\begin{equation}\tag{0.a}
\psi(\xi+\ell n)=\boldsymbol{\chi}(n)\,\psi(\xi),\qquad\chi(\boldsymbol{\zeta}+\ell n)=\boldsymbol{\chi}(n)\,\chi(\boldsymbol{\zeta})\qquad(n\in\mathbb{Z}^3).
\end{equation}
Replacing $\mathrm{Fdm}$ or the orders by other choices multiplies every basis monomial $\psi(I_1)\chi(I_2)$ by a sign.
Hence the norms (I.5.1), the bidegree, and the properties of being even and of being balanced are independent of these choices.

A kernel $K_{nm}(\xi_1,\dots,\xi_n;\boldsymbol{\zeta}_1,\dots,\boldsymbol{\zeta}_m)$ is a function on the covering index sets that is multiplied by $\boldsymbol{\chi}(v)$ when one argument $\xi_i$ is replaced by $\xi_i+\ell v$, $v\in\mathbb{Z}^3$, the other arguments being kept fixed, and that behaves in the same way in each argument $\boldsymbol{\zeta}_j$.
For such a kernel the product
\[
K_{nm}(\xi_1,\dots,\xi_n;\boldsymbol{\zeta}_1,\dots,\boldsymbol{\zeta}_m)\,\psi(\xi_1)\cdots\psi(\xi_n)\,\chi(\boldsymbol{\zeta}_1)\cdots\chi(\boldsymbol{\zeta}_m)
\]
is $\ell\mathbb{Z}^3$-periodic in each of its $n+m$ arguments separately.
Consequently, in $\int\mathrm{d}\xi\,\mathrm{d}\boldsymbol{\zeta}$ each argument may be summed over an arbitrary set of representatives of its classes modulo $\ell\mathbb{Z}^3$, the set of representatives being chosen for each argument separately.
The antisymmetric kernels $F_{nm}$ of an element $F$ are determined by $F$ (Proposition~15\,(a)).

A matrix kernel is a function $K$ on $\mathbb{L}^{\mathrm{f}}_k\times\mathbb{L}^{\mathrm{f}}_k$ with values in the complex $4\times4$ matrices such that
\[
K(x+\ell n,x')=K(x,x'+\ell n)=\boldsymbol{\chi}(n)\,K(x,x')\qquad\text{for all }n\in\mathbb{Z}^3.
\]
The matrix kernels are precisely the functions $K$ of Proposition~16.2\,(v).
Indeed, the condition required there is the case $n=e_\mu$; evaluated at $x-\ell e_\mu$, resp. at $x'-\ell e_\mu$, it yields the case $n=-e_\mu$; and an arbitrary $n$ is a sum of $|n_1|+|n_2|+|n_3|$ vectors $\pm e_\mu$, a number of the same parity as $n_1+n_2+n_3$.
For a matrix kernel $K$ the element $\mathrm{Bil}[K]$ is that of Proposition~16.2\,(v); written out,
\begin{align*}
\mathrm{Bil}[K]&=\int\mathrm{d}x\int\mathrm{d}x'\,\bar{\psi}(x)^TK(x,x')\psi(x')\\
&=\left(\frac{\eta^3}{48}\right)^2\sum_{c,c'\in\mathrm{Fdm}}\ \sum_{g,g'\in O_{\mathrm{cub}}}\ \sum_{\beta,\beta'}K\left((c,g),(c',g')\right)_{\beta\beta'}\cdot\bar{\psi}_{\beta}(c,g)\,\psi_{\beta'}(c',g').
\end{align*}

\begin{named}{Lemma 0.3}
Let $K$ be a matrix kernel.
\begin{itemize}
\item[(a)] $\mathrm{Bil}[K]$ is balanced and of bidegree $(2,0)$, and its antisymmetric kernel $F_{20}$ is given by
\begin{equation}\tag{0.b}
\begin{aligned}
&F_{20}\left((x,\beta,1),(x',\beta',0)\right)=K(x,x')_{\beta\beta'},\qquad F_{20}\left((x',\beta',0),(x,\beta,1)\right)=-K(x,x')_{\beta\beta'},\\
&F_{20}(\xi,\xi')=0\quad\text{if }\xi\text{ and }\xi'\text{ have the same bar index.}
\end{aligned}
\end{equation}
Thus $K$ is the matrix $F_{20}(x,x')$ of Section~I.5.1, and $\mathrm{Bil}[K]=\mathrm{Bil}[K']$ implies $K=K'$.
\item[(b)] Every balanced element $F$ of bidegree $(2,0)$ equals $\mathrm{Bil}[K]$ for the matrix kernel $K$ with
\[
K(x,x')_{\beta\beta'}:=F_{20}\left((x,\beta,1),(x',\beta',0)\right).
\]
\item[(c)] For every weight vector $\mathsf{w}=(\mathsf{w}_1,\mathsf{w}_2)$,
\[
\left\|\mathrm{Bil}[K]\right\|_{\mathsf{w}}=\frac{\mathsf{w}_1^2}{2}\,\|F_{20}\|=\mathsf{w}_1^2\sum_{\beta,\beta'}\int\mathrm{d}x\int\mathrm{d}x'\,\left|K(x,x')_{\beta\beta'}\right|.
\]
\end{itemize}
\end{named}

In particular the display (G2.0.1) holds, and so do the assertions of Proposition~16.2\,(v) that precede its four transformation formulas.

\begin{proof}
(a)
By the written-out form of $\mathrm{Bil}[K]$, each of its monomials is the product of one barred and one unbarred generator of the family $\psi$; hence $\mathrm{Bil}[K]$ is balanced and of bidegree $(2,0)$.
Let $G(\xi,\xi')$ be the function defined by the right-hand sides of (0.b).
It is antisymmetric, and it is a kernel.
In
\[
\frac{1}{2}\int\mathrm{d}\xi\int\mathrm{d}\xi'\,G(\xi,\xi')\psi(\xi)\psi(\xi')
\]
the arguments with the bar indices $(1,0)$ contribute $\frac{1}{2}\mathrm{Bil}[K]$.
The arguments with the bar indices $(0,1)$ contribute
\[
\frac{1}{2}\sum_{\beta,\beta'}\int\mathrm{d}x\int\mathrm{d}x'\,\left(-K(x,x')_{\beta\beta'}\right)\psi_{\beta'}(x')\bar{\psi}_{\beta}(x)=\frac{1}{2}\mathrm{Bil}[K],
\]
because $\psi_{\beta'}(x')\bar{\psi}_{\beta}(x)=-\bar{\psi}_{\beta}(x)\psi_{\beta'}(x')$.
Thus $G$ is an antisymmetric kernel of $\mathrm{Bil}[K]$, and $G=F_{20}$ by Proposition~15\,(a).

(b)
By Proposition~15\,(a), $F_{20}(\xi,\xi')\,\nu(\xi)\,\nu(\xi')$ is the coefficient $c_F(\{\xi,\xi'\})$ whenever $\xi<\xi'$.
Since $F$ is balanced, this coefficient vanishes when the two bar indices coincide.
As $F_{20}$ is a kernel, the function $K$ defined in (b) is a matrix kernel, and $F_{20}$ has the form (0.b) with this $K$.
By the computation in (a),
\[
F=\frac{1}{2}\int\mathrm{d}\xi\int\mathrm{d}\xi'\,F_{20}(\xi,\xi')\psi(\xi)\psi(\xi')=\mathrm{Bil}[K].
\]

(c)
By the case of equality in Proposition~15\,(a),
\[
\|F\|_{\mathsf{w}}=\frac{\mathsf{w}_1^2}{2!}\int\mathrm{d}\xi\int\mathrm{d}\xi'\,\left|F_{20}(\xi,\xi')\right|
\]
for $F=\mathrm{Bil}[K]$, and by (0.b) the double integral equals $2\sum_{\beta,\beta'}\int\mathrm{d}x\int\mathrm{d}x'\,\left|K(x,x')_{\beta\beta'}\right|$.
\end{proof}

\subsubsection*{0.5. Elements with the first clause of (k2), written over a lift}

For a finite subset $\tilde{U}$ of $\mathbb{L}_k$ let $\mathrm{Grc}_k(\tilde{U})$ and, if $\operatorname{sup-diameter}(\tilde{U})+M<\ell$, the homomorphism $\mathrm{cov}_\ell$ be as in Section~I.1.9, with the generators $\psi$, $\chi$ only.
An index over $\tilde{U}$ is either an index $\xi=(x,\beta,\omega)$ with $c\in\tilde{U}$, or an index $\boldsymbol{\zeta}=(x,x',\beta,\omega)$ with $(x,x')\in\mathrm{Pzt}_k$ and $c,c'\in\tilde{U}$.
A datum on $\tilde{U}$ is a family $\tilde{K}=(\tilde{K}_{nm})$ of complex functions $\tilde{K}_{nm}(\xi_1,\dots,\xi_n;\boldsymbol{\zeta}_1,\dots,\boldsymbol{\zeta}_m)$ of indices over $\tilde{U}$ that are antisymmetric in the $\xi$ and, separately, in the $\boldsymbol{\zeta}$.
Put
\begin{multline*}
\mathrm{El}(\tilde{K}):=\sum_{n,m}\frac{1}{n!\,m!}\sum_{\xi_1,\dots,\xi_n,\boldsymbol{\zeta}_1,\dots,\boldsymbol{\zeta}_m}\left(\frac{\eta^3}{48}\right)^{n+2m}\\
\cdot\tilde{K}_{nm}(\xi_1,\dots,\xi_n;\boldsymbol{\zeta}_1,\dots,\boldsymbol{\zeta}_m)\,\psi(\xi_1)\cdots\psi(\xi_n)\,\chi(\boldsymbol{\zeta}_1)\cdots\chi(\boldsymbol{\zeta}_m)\in\mathrm{Grc}_k(\tilde{U}),
\end{multline*}
where the inner sum extends over the indices over $\tilde{U}$.
This is a finite sum: there are finitely many indices over $\tilde{U}$, and a product with a repeated generator vanishes.
For an element $F$ of $\mathrm{Gr}[\psi,\chi]$ of the torus of side $\ell$, the datum of $F$ on $\tilde{U}$ consists of the restrictions of the antisymmetric kernels $F_{nm}$ to the indices over $\tilde{U}$; for kernels this is the composition of the data with the covering map in the sense of Section~I.1.9.
For a set $\tilde{X}$ of sites of $\mathbb{L}_k$ let
\[
\tilde{X}^{+\eta}:=\left\{c\in\mathbb{L}_k:\mathrm{dist}(c,\tilde{X})\le\eta\right\},
\]
where $\mathrm{dist}$ is the sup-distance.

\begin{named}{Lemma 0.4}
Let $X$ be a polymer of $T_k$ with a lift $\tilde{X}$ such that $\operatorname{sup-diameter}(\tilde{X})+2\eta+M<\ell$, and put $\tilde{U}:=\tilde{X}^{+\eta}$.
Let $F\in\mathrm{Gr}[\psi,\chi]$ be an element whose antisymmetric kernels vanish unless all $\xi_i$ lie over $X$ and all pairs $\boldsymbol{\zeta}_j$ are admissible for $X$ (Section~I.1.7).
Then $\operatorname{sup-diameter}(\tilde{U})+M<\ell$, and
\[
F=\mathrm{cov}_\ell\left(\mathrm{El}(\tilde{K})\right),
\]
where $\tilde{K}$ is the datum of $F$ on $\tilde{U}$.
\end{named}

\begin{proof}
One has $\operatorname{sup-diameter}(\tilde{U})\le\operatorname{sup-diameter}(\tilde{X})+2\eta$.
Therefore $\operatorname{sup-diameter}(\tilde{U})+M<\ell$, and also $\operatorname{sup-diameter}(\tilde{U})+\eta<\ell$, since $\eta\le M$.
So $\tilde{U}$ is a lift of $\mathrm{pr}(\tilde{U})$, the map $\mathrm{pr}$ is injective on $\tilde{U}$ by (Lft1), and (Lft2) with $\mathrm{sd}=M$ applies to $\tilde{U}$ and to its subset $\tilde{X}$.

(1)
Let a class of indices $\xi$ modulo $\ell\mathbb{Z}^3$ lie over $X$.
It has exactly one member over $\tilde{U}$, and this member lies over $\tilde{X}$: the set $\tilde{X}$ contains a point over $\mathrm{pr}\,c$, and $\mathrm{pr}$ is injective on $\tilde{U}$.
A member over $\tilde{U}$ of a class that does not lie over $X$ contributes $0$ to $\mathrm{El}(\tilde{K})$, because the kernel vanishes there.

(2)
A pair of $\mathrm{Pz}_k$ that is admissible for $X$ has exactly one representative $(x,x')\in\mathrm{Pzt}_k$ with $c,c'\in\tilde{U}$.
Uniqueness: two representatives differ by a vector $\ell n$, and $\mathrm{pr}$ is injective on $\tilde{U}$.
Existence for a pair of kind (I): here the photon sites of the pair lie in $X$.
Let $(x,x')$ be a representative, so that $|c-c'|\le M$, and let $c^a$, $c^b$ be the points of $\tilde{X}$ over $\mathrm{pr}\,c$ and $\mathrm{pr}\,c'$.
Since $d(\mathrm{pr}\,c,\mathrm{pr}\,c')\le M$, (Lft2) with $\mathrm{sd}=M$ shows that $c^a-c^b$ is the only vector of sup norm at most $M$ in its class; hence $c-c'=c^a-c^b$.
As $c^a-c\in\ell\mathbb{Z}^3$, the pair $(x,x')+(c^a-c)$ is a representative over $\tilde{X}$.
This includes the pairs inside a cluster, for which $c=c'$, $c^a=c^b$ and $g\ne g'$.
Existence for a pair of kind (II): here $x'=x^{(3)}$, and $\mathrm{pr}\,c$ or $\mathrm{pr}\,c_3$ lies in $X$.
Take the representative for which this point lies in $\tilde{X}$.
The other one of the two points $c$ and $c_3=c+\eta\cdot ge_1$ has sup-distance $\eta$ from it, and so it lies in $\tilde{U}$.
Finally, a pair index over $\tilde{U}$ whose class is not admissible for $X$ contributes $0$ to $\mathrm{El}(\tilde{K})$, since by the hypothesis on $F$ the kernel vanishes there.

(3)
By Section~0.4, each argument of $\int\mathrm{d}\xi\,\mathrm{d}\boldsymbol{\zeta}$ may be summed over any set of representatives of its classes.
For the classes described in (1) and (2) take their members over $\tilde{U}$; the other classes contribute $0$.
By (1) and (2), the indices over $\tilde{U}$ whose classes are not of these two kinds contribute $0$ as well.
Hence $F$ is the sum that defines $\mathrm{El}(\tilde{K})$, with the generators of the torus at the same indices, in the sense of (I.AP), in place of the generators of $\mathrm{Grc}_k(\tilde{U})$.
That is, $F=\mathrm{cov}_\ell\left(\mathrm{El}(\tilde{K})\right)$.
\end{proof}

\subsection*{1. Proposition 16.1}

\subsubsection*{Proof of Proposition 16.1}

Two properties of $\tilde{g}=(S,a)$ underlie the whole argument: the matrix $r$ is diagonal with the entries $\pm1$, and the vector $a$ lies in $M\mathbb{Z}^3$.
Recall also the universal property of a Grassmann algebra: anticommuting images of independent generators define exactly one homomorphism, and two multiplicative maps which are both linear or both antilinear and which agree on the generators coincide (Chapter~17, Section~0.5, Lemma~A and (Un)).

\paragraph{(a1)}
\textbf{(a1.1) Isometry.}
For $v\in\mathbb{R}^3$ one has $(r\cdot v)_{\mu}=r_{\mu}\cdot v_{\mu}$, and therefore $|(r\cdot v)_{\mu}|=|v_{\mu}|$.
The sup norm, the $l^1$ norm and the Euclidean norm of $v$ are functions of $(|v_1|,|v_2|,|v_3|)$ alone.
As $\tilde{g}p-\tilde{g}p'=r\cdot(p-p')$, it follows that $|\tilde{g}p-\tilde{g}p'|=|p-p'|$ for each of these three norms.
Moreover, for $n\in\mathbb{Z}^3$,
\[
\tilde{g}(p+\ell\cdot n)=r\cdot p+\ell\cdot r\cdot n+a=\tilde{g}p+\ell\cdot r\cdot n,
\]
where $r\cdot n\in\mathbb{Z}^3$.

\textbf{(a1.2) Tori and the distance.}
Fix $i$ with $-1\le i\le k$.
One has $r\cdot L^{-i}\mathbb{Z}^3=L^{-i}\mathbb{Z}^3$, and $a\in M\mathbb{Z}^3\subset L^{-i}\mathbb{Z}^3$ by (Ar1).
Thus $\tilde{g}$ and $\tilde{g}^{-1}$ map $L^{-i}\mathbb{Z}^3$ into itself, so that $\tilde{g}$ is a bijection of $L^{-i}\mathbb{Z}^3$.
By (a1.1), $\tilde{g}$ and $\tilde{g}^{-1}$ carry classes modulo $\ell\mathbb{Z}^3$ onto classes.
Hence $\tilde{g}$ induces a bijection of $L^{-i}\mathbb{Z}^3/\ell\mathbb{Z}^3$ and a bijection of $\mathbb{R}^3/\ell\mathbb{Z}^3$, both denoted again by $\tilde{g}$, and $\mathrm{pr}\circ\tilde{g}=\tilde{g}\circ\mathrm{pr}$.
The induced maps are independent of the vector chosen for $a$ in its class.
For points $p$, $p'$ of $\mathbb{R}^3$,
\begin{align*}
d\left(\tilde{g}\,\mathrm{pr}\,p,\tilde{g}\,\mathrm{pr}\,p'\right)
&=\min_{n\in\mathbb{Z}^3}\left|r\cdot(p-p')+\ell\cdot n\right|
=\min_{n\in\mathbb{Z}^3}\left|r\cdot\left(p-p'+\ell\cdot r\cdot n\right)\right|\\
&=\min_{n\in\mathbb{Z}^3}\left|p-p'+\ell\cdot r\cdot n\right|
=d\left(\mathrm{pr}\,p,\mathrm{pr}\,p'\right),
\end{align*}
since $n\mapsto r\cdot n$ is a bijection of $\mathbb{Z}^3$.

\textbf{(a1.3) Lifts.}
Let $Z$ be a set of sites of $T_k$, and let $\mathsf{Z}$ be a lift of $Z$.
By (a1.2), $\tilde{g}\mathsf{Z}$ is a subset of $\mathbb{L}_k$ with $\mathrm{pr}(\tilde{g}\mathsf{Z})=\tilde{g}\,\mathrm{pr}(\mathsf{Z})=\tilde{g}Z$, and by (a1.1) the sup-diameter of $\tilde{g}\mathsf{Z}$ equals that of $\mathsf{Z}$.
Therefore $\tilde{g}\mathsf{Z}$ is a lift of $\tilde{g}Z$.

Let now $y$, $c$ be sites with $d(y,c)<\ell/2$.
A \emph{shortest lift} of $(y,c)$ is a pair $(\tilde{y},\tilde{c})$ of points over $y$ and $c$ with $|\tilde{y}-\tilde{c}|=d(y,c)$.
A shortest lift exists, and it is unique up to a common translation by a vector of $\ell\mathbb{Z}^3$.
Indeed, if $(\tilde{y},\tilde{c}+\ell\cdot n)$ with $n\ne0$ were another shortest lift with the same first point, then
\[
\left|\tilde{y}-\tilde{c}-\ell\cdot n\right|\ge\ell-|\tilde{y}-\tilde{c}|>\frac{\ell}{2}>d(y,c).
\]
The pair $(\tilde{g}\tilde{y},\tilde{g}\tilde{c})$ lies over $(\tilde{g}y,\tilde{g}c)$, and (a1.1), (a1.2) give
\[
|\tilde{g}\tilde{y}-\tilde{g}\tilde{c}|=|\tilde{y}-\tilde{c}|=d(y,c)=d(\tilde{g}y,\tilde{g}c).
\]
Consequently $(\tilde{g}\tilde{y},\tilde{g}\tilde{c})$ is a shortest lift of $(\tilde{g}y,\tilde{g}c)$, and $d(\tilde{g}y,\tilde{g}c)<\ell/2$.

\textbf{(a1.4) Sets defined by the distance from a point.}
Let $z\in\mathbb{R}^3$, and let $\mathrm{tq}>0$ be a number.
By (a1.1),
\[
\tilde{g}\left\{p:|p-z|<\mathrm{tq}\right\}=\left\{p:|p-\tilde{g}z|<\mathrm{tq}\right\};
\]
the same identity holds with $\le$ in place of $<$, and for the intersections of these sets with any lattice that $\tilde{g}$ maps onto itself; by (a1.2) it holds on the torus as well, with the distance $d$.
This yields the following facts.
\begin{itemize}
\item The block $\{c:|c-z|<L^{-i}/2\}$ with the centre $z\in L^{-i}\mathbb{Z}^3$ is mapped onto the block centred at $\tilde{g}z$, and $\tilde{g}z\in L^{-i}\mathbb{Z}^3$.
The block $\mathrm{Bl}_j(z)$, $0\le j\le k+1$, $z\in L^{j}\eta\mathbb{Z}^3$, is the case $i=k-j$ with $c\in\mathbb{L}_k$; $\Delta_y$ is $\mathrm{Bl}_k(y)$; $B(y')$ is the case $i=-1$ with $c\in\mathbb{Z}^3$.
\item $\tilde{g}(M\mathbb{Z}^3)=M\mathbb{Z}^3$, because $r\cdot M\mathbb{Z}^3=M\mathbb{Z}^3$ and $a\in M\mathbb{Z}^3$.
So the $M$-cube with the centre $z$ is mapped onto the $M$-cube with the centre $\tilde{g}z$, the closed $M$-cube onto the closed $M$-cube, and $\tilde{g}$ is a bijection of the set of the $M$-cubes.
The cubes of $\pi_k$ are the $M$-cubes.
The cube $\{c:|c-z|\le M\}$ of $\pi'_k$ is mapped onto the cube of $\pi'_k$ with the centre $\tilde{g}z$.
\item $\tilde{g}K_z=\{c:|c-\tilde{g}z|\le2L\}=K_{\tilde{g}z}$.
\end{itemize}

\textbf{(a1.5) Paths.}
Choose $i$ with $-1\le i\le k$, and write $\mathrm{hm}:=L^{-i}$ for the mesh of the lattice $L^{-i}\mathbb{Z}^3$.
The paths of Section~I.1.1 are the case $i=k$, $\mathrm{hm}=\eta$; the paths of the unit lattice are the case $i=0$.
The covering lattice is treated first.
Let $y,c\in\mathrm{hm}\,\mathbb{Z}^3$, and let $\pi=(\pi_1,\pi_2,\pi_3)$ be an order of $\{1,2,3\}$.
For $w\in\mathrm{hm}\,\mathbb{Z}^3$, a direction $\mu$ and $\mathrm{tq}\in\mathrm{hm}\,\mathbb{Z}$, the \emph{straight path} from $w$ to $w+\mathrm{tq}\cdot e_{\mu}$ is the sequence of the oriented bonds
\[
\left(w+(j-1)\cdot\mathrm{sn}\cdot\mathrm{hm}\cdot e_{\mu},\ w+j\cdot\mathrm{sn}\cdot\mathrm{hm}\cdot e_{\mu}\right),\qquad j=1,\dots,\frac{|\mathrm{tq}|}{\mathrm{hm}},
\]
where $\mathrm{sn}$ denotes the sign of $\mathrm{tq}$; for $\mathrm{tq}=0$ the sequence is empty.
Put $\mathrm{wc}_0:=y$ and $\mathrm{wc}_j:=\mathrm{wc}_{j-1}+(c-y)_{\pi_j}\cdot e_{\pi_j}$ for $j=1,2,3$; then $\mathrm{wc}_3=c$.
The path $\boldsymbol{\Gamma}^{\pi}(y,c)$ consists of the straight path from $\mathrm{wc}_0$ to $\mathrm{wc}_1$, followed by the straight paths from $\mathrm{wc}_1$ to $\mathrm{wc}_2$ and from $\mathrm{wc}_2$ to $\mathrm{wc}_3$ (Section~I.1.1).
Since $r$ is diagonal,
\[
\tilde{g}\left(w+j\cdot\mathrm{sn}\cdot\mathrm{hm}\cdot e_{\mu}\right)=\tilde{g}w+j\cdot(\mathrm{sn}\cdot r_{\mu})\cdot\mathrm{hm}\cdot e_{\mu};
\]
thus $\tilde{g}$ maps the straight path from $w$ to $w+\mathrm{tq}\cdot e_{\mu}$, bond by bond and in the same succession, onto the straight path from $\tilde{g}w$ to $\tilde{g}w+r_{\mu}\cdot\mathrm{tq}\cdot e_{\mu}$.
Furthermore
\[
\tilde{g}\,\mathrm{wc}_j=\tilde{g}\,\mathrm{wc}_{j-1}+r_{\pi_j}\cdot(c-y)_{\pi_j}\cdot e_{\pi_j}=\tilde{g}\,\mathrm{wc}_{j-1}+(\tilde{g}c-\tilde{g}y)_{\pi_j}\cdot e_{\pi_j}.
\]
Therefore $\tilde{g}\,\mathrm{wc}_0,\dots,\tilde{g}\,\mathrm{wc}_3$ are the corners of $\boldsymbol{\Gamma}^{\pi}(\tilde{g}y,\tilde{g}c)$, with the same order $\pi$, and
\begin{equation}\tag{eq.1}
\tilde{g}\,\boldsymbol{\Gamma}^{\pi}(y,c)=\boldsymbol{\Gamma}^{\pi}(\tilde{g}y,\tilde{g}c)\qquad\text{as sequences of oriented bonds,}
\end{equation}
where $\tilde{g}(c^a,c^b):=(\tilde{g}c^a,\tilde{g}c^b)$.
On the torus, for $d(y,c)<\ell/2$, the path $\boldsymbol{\Gamma}^{\pi}(y,c)$ is the image under $\mathrm{pr}$ of $\boldsymbol{\Gamma}^{\pi}(\tilde{y},\tilde{c})$, where $(\tilde{y},\tilde{c})$ is a shortest lift; translating the lift by $\ell\cdot n$ translates the path by $\ell\cdot n$, so the image is independent of the shortest lift chosen.
From (a1.3), from (eq.1) on the covering lattice and from $\mathrm{pr}\circ\tilde{g}=\tilde{g}\circ\mathrm{pr}$ one concludes that (eq.1) holds on the torus.

\textbf{(a1.6) Segments, nets, $d_M$.}
One has $\tilde{g}\left(p+t\cdot(p'-p)\right)=\tilde{g}p+t\cdot(\tilde{g}p'-\tilde{g}p)$.
Hence a closed straight segment $\mathrm{seg}$ of $\mathbb{R}^3$ is mapped onto a closed straight segment of the same Euclidean length (a1.1), and $\tilde{g}\,\mathrm{pr}(\mathrm{seg})=\mathrm{pr}(\tilde{g}\,\mathrm{seg})$.
On the continuum torus the induced map is an isometry for $d$ whose inverse is induced by $\tilde{g}^{-1}$; it is therefore a homeomorphism and carries connected sets onto connected sets.
It follows that a net, written as a union of segments, is mapped onto a net, written as the union of the image segments, and that both nets have the same length.
By (a1.4), a net meets the closed $M$-cube $\square$ if and only if its image meets the closed $M$-cube $\tilde{g}\square$.

\textbf{(a1.7) Faces, polymers.}
Let $z$, $z'$ be the centres of two $M$-cubes.
Then $\tilde{g}z-\tilde{g}z'=r\cdot(z-z')$, and for every direction $\mu$ the matrix $r$ maps the set $\{M\cdot e_{\mu},-M\cdot e_{\mu}\}+\ell\mathbb{Z}^3$ onto itself.
Hence two $M$-cubes have a common face if and only if their images have one.
Consequently $\tilde{g}$ maps the non-empty unions of $M$-cubes that are connected through faces onto unions of the same kind; that is, polymers go onto polymers.
As $\tilde{g}$ is a bijection of the set of the $M$-cubes, $|\tilde{g}X|_M=|X|_M$.
By (a1.6), applied to $\tilde{g}$ and to $\tilde{g}^{-1}$, the nets meeting every closed $M$-cube of $X$ correspond, with equal lengths, to the nets meeting every closed $M$-cube of $\tilde{g}X$.
Therefore $d_M(\tilde{g}X)=d_M(X)$, and $X$ is small if and only if $\tilde{g}X$ is small.

\textbf{(a1.8) Enlargements.}
An $M$-cube with the centre $z'$ belongs to $X^{\sim i}$ if and only if $d(z',z)\le i\cdot M$ for the centre $z$ of some cube of $X$.
Now $d(\tilde{g}z',\tilde{g}z)=d(z',z)$, and $\tilde{g}z$ ranges over the centres of the cubes of $\tilde{g}X$ as $z$ ranges over the centres of the cubes of $X$.
It follows that $\tilde{g}\left(X^{\sim i}\right)=(\tilde{g}X)^{\sim i}$; for $i=L$ and $X=\square$ this reads $\tilde{g}\left(\square^{\dagger}\right)=(\tilde{g}\square)^{\dagger}$.

\textbf{(a1.9) Linked cubes.}
Cubes of $\pi_k$ are linked when they share a face, and (a1.7) applies to them.
Cubes of $\pi'_k$ are linked when they have a common site of $T_k$; since $\tilde{g}$ is a bijection of $T_k$, a site $c$ lies in both cubes if and only if $\tilde{g}c$ lies in both images.

\textbf{(a1.10) Chains.}
By a chain of $M$-cubes of $X$ from $\square$ to $\square'$ one means a sequence $(\square_0,\dots,\square_n)$ of $M$-cubes of $X$ with $\square_0=\square$ and $\square_n=\square'$ in which any two consecutive cubes have a common face.
By (a1.7), $(\tilde{g}\square_0,\dots,\tilde{g}\square_n)$ is a chain of $\tilde{g}X$ from $\tilde{g}\square$ to $\tilde{g}\square'$ with the same $n$, and $\tilde{g}^{-1}$ maps the chains of $\tilde{g}X$ from $\tilde{g}\square$ to $\tilde{g}\square'$ to chains of $X$ from $\square$ to $\square'$.
The least $n$ is therefore the same in both cases, and $(\square_i)\mapsto(\tilde{g}\square_i)$ is a bijection of the set of the shortest chains of $X$ from $\square$ to $\square'$ onto the set of the shortest chains of $\tilde{g}X$ from $\tilde{g}\square$ to $\tilde{g}\square'$.

\textbf{(a1.11) The step group.}
Let $\tilde{g}\in\mathbb{G}_k^{L}$, that is, $a\in LM\mathbb{Z}^3$.
Then $\tilde{g}(LM\mathbb{Z}^3)=LM\mathbb{Z}^3$, and by (a1.4) the $LM$-cube $\left\{c:|c-z|<\frac{L\cdot M}{2}\right\}$, $z\in LM\mathbb{Z}^3$, is mapped onto the $LM$-cube with the centre $\tilde{g}z$, and the cube $\{c:|c-z|\le L\cdot M\}$ of $\pi''_k$ onto the cube of $\pi''_k$ with the centre $\tilde{g}z$; two cubes of $\pi''_k$ have a common site if and only if their images have one.
An $LM$-cube contains the $M$-cube $\square$ if and only if its image contains $\tilde{g}\square$.
By the definition of $\bar{X}^{L}$ in \cite[Section~13.1.7]{QEDI}, $\bar{X}^{L}$ is the union of the $LM$-cubes that contain an $M$-cube of $X$; hence $\tilde{g}\left(\bar{X}^{L}\right)=\overline{\tilde{g}X}^{L}$.

Let now $\tilde{g}=(S,a)\in\mathbb{G}_k$ with $a\notin LM\mathbb{Z}^3$; this is a condition on the class of $a$, because $\ell\mathbb{Z}^3\subset LM\mathbb{Z}^3$.
Let $\mathrm{pc}$ be the cube of $\pi''_k$ with the centre $z$.
Then
\[
\tilde{g}^{-1}\mathrm{pc}=\left\{c:\left|c-r\cdot(z-a)\right|\le L\cdot M\right\},\qquad r\cdot(z-a)\in M\mathbb{Z}^3\setminus LM\mathbb{Z}^3.
\]
A set $\{c\in T_k:d(c,z')\le L\cdot M\}$, $z'$ a site, determines $z'$.
Indeed, let $\tilde{z}'$ be a point of $\mathbb{L}_k$ over $z'$; the set has the lift $\{\tilde{c}:|\tilde{c}-\tilde{z}'|\le L\cdot M\}$, a lattice box with the midpoint $\tilde{z}'$ and the sup-diameter $2L\cdot M$, where $2L\cdot M+\eta<\ell$ by the second inequality of (Ar3); being the image of a lattice box, the set is connected through nearest neighbours, and its lifts are therefore the translates of this box by $\ell\mathbb{Z}^3$ (Lft3).
Hence $\tilde{g}^{-1}\mathrm{pc}$ is not a cube of $\pi''_k$, and $s_{\tilde{g}^{-1}\mathrm{pc}}$ has no meaning.
For $\tilde{g}\in\mathbb{G}_k^{L}$, on the other hand, $\tilde{g}^{-1}$ lies in $\mathbb{G}_k^{L}$ as well, and $\tilde{g}^{-1}\mathrm{pc}$ is a cube of $\pi''_k$.

\paragraph{(a2)}
\textbf{(a2.1) Bijection.}
The matrix $r$ belongs to $\{\pm1\}^3$, a subgroup of $O_{\mathrm{cub}}$ (Section~I.1.2); hence $r\cdot g\in O_{\mathrm{cub}}$, and $g\mapsto r\cdot g$ is a bijection of the $48$ labels which is its own inverse, because $r^2=1$.
Combined with (a1.2) for $i=k$, this shows that $x=(c,g)\mapsto(r\cdot c+a,r\cdot g)$ is a bijection of $\mathbb{L}^{\mathrm{f}}_k=\mathbb{L}_k\times O_{\mathrm{cub}}$, with the inverse $(c,g)\mapsto\left(r\cdot(c-a),r\cdot g\right)$.
One has
\[
\tilde{g}(x+\ell\cdot n)=\left(r\cdot c+\ell\cdot r\cdot n+a,\ r\cdot g\right)=\tilde{g}x+\ell\cdot r\cdot n,
\]
so that $\tilde{g}$ induces a bijection of $T^{\mathrm{f}}_k$.
Moreover $\pi(\tilde{g}x)=r\cdot c+a=\tilde{g}\pi(x)$, and $\tilde{g}\left(\{c\}\times O_{\mathrm{cub}}\right)=\{\tilde{g}c\}\times O_{\mathrm{cub}}$.

\textbf{(a2.2) Neighbours.}
With the neighbours $x^{(i)}$ of (I.1.2) one has, for $i=0,1,2$,
\[
\tilde{g}\left(x^{(i)}\right)=\left(r\cdot c+a,\ r\cdot g\cdot s_i\right)=(\tilde{g}x)^{(i)},
\]
and, for $i=3$,
\begin{align*}
\tilde{g}\left(x^{(3)}\right)&=\left(r\cdot(c+\eta\cdot g\cdot e_1)+a,\ r\cdot g\cdot s_3\right)\\
&=\left((r\cdot c+a)+\eta\cdot(r\cdot g)\cdot e_1,\ (r\cdot g)\cdot s_3\right)=(\tilde{g}x)^{(3)}.
\end{align*}
Further, $\hat{d}_i(\tilde{g}x)=-(r\cdot g)\cdot\hat{n}_i=r\cdot\left(-g\cdot\hat{n}_i\right)=r\cdot\hat{d}_i(x)$.
In these computations the left multiplication by $r$ commutes with the right multiplication by $s_i$.

\textbf{(a2.3) Measure: the three substitution rules of Proposition~16.1\,(a2).}
Let $f$ be an $\ell\mathbb{Z}^3$-periodic function on $\mathbb{L}^{\mathrm{f}}_k$ with values in a vector space; it is a function on $T^{\mathrm{f}}_k$.
One has
\[
\int\mathrm{d}x\,f(x)=\frac{\eta^3}{48}\sum_{x\in T^{\mathrm{f}}_k}f(x),
\]
and $x\mapsto\tilde{g}x$ is a bijection of $T^{\mathrm{f}}_k$ by (a2.1); therefore $\int\mathrm{d}x\,f(\tilde{g}x)=\int\mathrm{d}x\,f(x)$.
For a set $Z$ of sites of $T_k$ the relation $\pi(\tilde{g}x)=\tilde{g}\pi(x)$ gives $\mathbf{1}_{\tilde{g}Z}(\tilde{g}x)=\mathbf{1}_Z(x)$, whence
\[
\int_{\tilde{g}Z}\mathrm{d}x\,f(x)=\int\mathrm{d}x\,\mathbf{1}_{\tilde{g}Z}(x)\cdot f(x)=\int\mathrm{d}x\,\mathbf{1}_{\tilde{g}Z}(\tilde{g}x)\cdot f(\tilde{g}x)=\int_Z\mathrm{d}x\,f(\tilde{g}x).
\]
On the covering lattice, let $\mathsf{Z}$ be a finite subset of $\mathbb{L}_k$ and $f$ an \emph{arbitrary} function on $\mathbb{L}^{\mathrm{f}}_k$.
The map $\tilde{g}$ carries $\mathsf{Z}\times O_{\mathrm{cub}}$ bijectively onto $(\tilde{g}\mathsf{Z})\times O_{\mathrm{cub}}$.
So, for the integrals
\[
\int_{\mathsf{Z}}\mathrm{d}x\,f(x):=\frac{\eta^3}{48}\sum_{c\in\mathsf{Z}}\ \sum_{g\in O_{\mathrm{cub}}}f(c,g)
\]
of Proposition~16.1\,(a2), one obtains
\begin{equation}\tag{eq.2}
\int_{\tilde{g}\mathsf{Z}}\mathrm{d}x\,f(x)=\int_{\mathsf{Z}}\mathrm{d}x\,f(\tilde{g}x).
\end{equation}

\textbf{(a2.4) The distance $d_X$; pairs.}
Let $x$ and $x'$ be distinct points of $\mathbb{L}^{\mathrm{f}}_k$.
Then $\tilde{g}x\ne\tilde{g}x'$ and $|\pi(\tilde{g}x)-\pi(\tilde{g}x')|=|r\cdot(c-c')|=|c-c'|$, so that
\[
d_X(\tilde{g}x,\tilde{g}x')=\max\left(|c-c'|,\eta\right)=d_X(x,x').
\]
Hence $(x,x')\mapsto(\tilde{g}x,\tilde{g}x')$ maps $\mathrm{Pzt}_k$ into $\mathrm{Pzt}_k$; as the same is true of $\tilde{g}^{-1}$, this map is a bijection of $\mathrm{Pzt}_k$.
From
\[
\left(\tilde{g}(x+\ell\cdot n),\tilde{g}(x'+\ell\cdot n)\right)=\left(\tilde{g}x+\ell\cdot r\cdot n,\ \tilde{g}x'+\ell\cdot r\cdot n\right)
\]
it follows that classes modulo $\ell\mathbb{Z}^3$ are mapped onto classes; so a bijection of $\mathrm{Pz}_k$ is induced.

\textbf{(a2.5) Sorts.}
For a pair inside a cluster ($c=c'$, $g\ne g'$) one has $\tilde{g}c=\tilde{g}c'$ and $r\cdot g\ne r\cdot g'$.
For short pairs ($d_X<1$) and for long pairs ($1\le d_X\le M$) the assertion follows from (a2.4).
A bond of type $i$, $(x,x^{(i)})$, is sent to $(\tilde{g}x,(\tilde{g}x)^{(i)})$ by (a2.2).

\textbf{(a2.6) Kinds.}
Let $Z$ be a union of $M$-cubes of $T_k$, and let $(x,x')\in\mathrm{Pzt}_k$.
By (Ar3), $|c-c'|\le M<\ell/2$, and every other vector of the class of $c-c'$ modulo $\ell\mathbb{Z}^3$ has sup norm at least $\ell-M>M$; hence $(c,c')$ is a shortest lift of $(\mathrm{pr}\,c,\mathrm{pr}\,c')$.
The class of $(x,x')$ is of kind (I) for $Z$ if and only if $\mathrm{pr}\,w\in Z$ for every site $w$ of the six paths $\boldsymbol{\Gamma}^{\pi}(c,c')$; for $c=c'$ the condition is $\mathrm{pr}\,c\in Z$, and in any case the sites $c$, $c'$ belong to the paths.
By (eq.1) the sites of $\boldsymbol{\Gamma}^{\pi}(\tilde{g}c,\tilde{g}c')$ are the points $\tilde{g}w$, and $\mathrm{pr}\,\tilde{g}w=\tilde{g}\,\mathrm{pr}\,w$ lies in $\tilde{g}Z$ if and only if $\mathrm{pr}\,w$ lies in $Z$.
So $(x,x')$ is of kind (I) for $Z$ if and only if $(\tilde{g}x,\tilde{g}x')$ is of kind (I) for $\tilde{g}Z$.
Kind (II) for $Z$ means that $x'=x^{(3)}$ and that $\mathrm{pr}\,c$ or $\mathrm{pr}\,c_3$ lies in $Z$.
By (a2.2), $x'=x^{(3)}$ holds if and only if $\tilde{g}x'=(\tilde{g}x)^{(3)}$, and $\pi\left((\tilde{g}x)^{(3)}\right)=\tilde{g}c_3$.
So $(x,x')$ is of kind (II) for $Z$ if and only if $(\tilde{g}x,\tilde{g}x')$ is of kind (II) for $\tilde{g}Z$.
Together with (a2.4) this shows that the pairs admissible for $Z$, of either kind, are mapped \emph{onto} the pairs admissible for $\tilde{g}Z$, of the same kind.

\textbf{(a2.7) Chain points.}
According to Section~I.5.6 the chain points are formed over one lift of $X$, and a small polymer has a lift by Lemma~0.2\,(c).
Let $X$ be small and let $\tilde{X}$ be any lift of $X$; let $x=(c,g)$ and $\tilde{x}=(\tilde{c},\breve{g})$ be points over $\tilde{X}$.
Let $(\square_0,\dots,\square_n)$ be a shortest chain of $M$-cubes of $X$ from the cube of $\mathrm{pr}\,c$ to the cube of $\mathrm{pr}\,\tilde{c}$, let $\tilde{\square}_i$ be the cube of $\tilde{X}$ over $\square_i$, and let $z_i$ be its centre (Lft4).
The chain points (Section~I.5.6) are
\[
p_0=x,\qquad p_{i+1}=(z_i,g)\quad(0\le i\le n),\qquad p_{n+2}=\tilde{x}.
\]
They are determined by $X$, $x$, $\tilde{x}$ and the chain, since $\tilde{X}$ is the only lift of $X$ containing $c$.
Indeed, by Lemma~0.2\,(b) every lift of $X$ has the form $\tilde{X}+\ell\cdot n'$ with $n'\in\mathbb{Z}^3$; if $c\in\tilde{X}+\ell\cdot n'$, then $c$ and $c-\ell\cdot n'$ are points of $\tilde{X}$ with the same image, and (Lft1) gives $n'=0$.

The polymer $\tilde{g}X$ is small by (a1.7), and $\tilde{g}\tilde{X}$ is a lift of $\tilde{g}X$ by (a1.3).
The points $\tilde{g}x=(\tilde{g}c,r\cdot g)$ and $\tilde{g}\tilde{x}=(\tilde{g}\tilde{c},r\cdot\breve{g})$ lie over $\tilde{g}\tilde{X}$, and by (a1.10) the sequence $(\tilde{g}\square_i)$ is a shortest chain of $\tilde{g}X$ from the cube of $\mathrm{pr}\,\tilde{g}c=\tilde{g}\,\mathrm{pr}\,c$ to the cube of $\mathrm{pr}\,\tilde{g}\tilde{c}=\tilde{g}\,\mathrm{pr}\,\tilde{c}$.
By (a1.4), $\tilde{g}\tilde{\square}_i$ is the $M$-cube of $\mathbb{L}_k$ with the centre $\tilde{g}z_i$; it is contained in $\tilde{g}\tilde{X}$, and $\mathrm{pr}(\tilde{g}\tilde{\square}_i)=\tilde{g}\square_i$, so it is the cube of $\tilde{g}\tilde{X}$ over $\tilde{g}\square_i$ (Lft4).
Consequently
\[
\tilde{g}(z_i,g)=(\tilde{g}z_i,r\cdot g)=\left(\text{the centre of the cube of }\tilde{g}\tilde{X}\text{ over }\tilde{g}\square_i,\ r\cdot g\right),
\]
and the chain points of $\left(\tilde{g}X,\tilde{g}x,\tilde{g}\tilde{x},(\tilde{g}\square_i)\right)$, formed over the lift $\tilde{g}\tilde{X}$, are $\tilde{g}x$, then the points $(\tilde{g}z_i,r\cdot g)=\tilde{g}(z_i,g)$ for $0\le i\le n$, whose label $r\cdot g$ is the label of $\tilde{g}x$, and finally $\tilde{g}\tilde{x}$.

\paragraph{(b)}
\textbf{(b.1) Norms and pairings.}
Consider one of the lattices $L^{-i}\mathbb{Z}^3/\ell\mathbb{Z}^3$, $-1\le i\le k$.
By an oriented bond one means an ordered pair $b=(c,c')$ with $c'-c=\pm L^{-i}\cdot e_{\mu}$; put $\tilde{g}b:=(\tilde{g}c,\tilde{g}c')$.
Since $\tilde{g}c'-\tilde{g}c=\pm r_{\mu}\cdot L^{-i}\cdot e_{\mu}$, the map $b\mapsto\tilde{g}b$ is a bijection of the set of the oriented bonds, with the inverse given by $\tilde{g}^{-1}$, and it commutes with the reversal $(c,c')\mapsto(c',c)$; it therefore induces a bijection of the set of the unoriented bonds.
By definition $(\tilde{g}A)(\tilde{g}b)=A(b)$, and $\tilde{g}A$ is odd under reversal.
For bond functions $A$, $B$ the numbers $|A(b)|$ and $A(b)\cdot B(b)$ are unchanged when $b$ is reversed.
The sup norm is the supremum of the numbers $|A(b)|$; the $l^1$ norm and the square of the $l^2$ norm are a weight, the same for all bonds, times the sum of $|A(b)|$, resp.\ of $|A(b)|^2$, over the unoriented bonds; and the pairing
\[
\langle A,B\rangle=\sum_c\sum_{\mu}(\text{weight})\cdot A_{\mu}(c)\cdot B_{\mu}(c)
\]
of \cite[Section~5.2.5]{QEDI} is the sum of $(\text{weight})\cdot A(b)\cdot B(b)$ over the unoriented bonds.
All these quantities are consequently invariant under $(A,B)\mapsto(\tilde{g}A,\tilde{g}B)$.
An oriented plaquette is the cyclically ordered quadruple of the corners in Section~I.1.1, or the same quadruple with the reversed cyclic order.
The map $\tilde{g}$ sends the quadruple of $(c;\mu\nu)$ to
\begin{multline*}
\Bigl(\tilde{g}c,\ \tilde{g}c+r_{\mu}\cdot L^{-i}\cdot e_{\mu},\\
\tilde{g}c+r_{\mu}\cdot L^{-i}\cdot e_{\mu}+r_{\nu}\cdot L^{-i}\cdot e_{\nu},\ \tilde{g}c+r_{\nu}\cdot L^{-i}\cdot e_{\nu}\Bigr),
\end{multline*}
which is an oriented plaquette in the plane $\mu\nu$.
So $\tilde{g}$ permutes the oriented plaquettes and commutes with the reversal, and the preceding argument applies to the functions of the oriented plaquettes that are odd under reversal.

\textbf{(b.2) The three norms.}
Let $Z$ be a set of sites of $T_k$, the case $Z=T_k$ being admitted, and let $A$ be a function on the bonds inside $Z$ that is odd under reversal; on the bonds inside $\tilde{g}Z$ put $(\tilde{g}A)(b):=A(\tilde{g}^{-1}b)$.
A bond lies inside $\tilde{g}Z$ if and only if its image under $\tilde{g}^{-1}$ lies inside $Z$.
Define $\mathrm{sft}_{\mu}:=\eta\cdot e_{\mu}$ if $r_{\mu}=-1$, and $\mathrm{sft}_{\mu}:=0$ if $r_{\mu}=1$.
For $c\in T_k$ write $z:=\tilde{g}^{-1}c$; then $\tilde{g}^{-1}(c+\eta\cdot e_{\mu})=z+r_{\mu}\cdot\eta\cdot e_{\mu}$.
If $r_{\mu}=1$, then $(\tilde{g}A)_{\mu}(c)=A(z,z+\eta\cdot e_{\mu})=A_{\mu}(z)$.
If $r_{\mu}=-1$, then $(\tilde{g}A)_{\mu}(c)=A(z,z-\eta\cdot e_{\mu})=-A(z-\eta\cdot e_{\mu},z)=-A_{\mu}(z-\eta\cdot e_{\mu})$.
In both cases
\begin{equation}\tag{eq.3}
(\tilde{g}A)_{\mu}(c)=r_{\mu}\cdot A_{\mu}\left(\tilde{g}^{-1}c-\mathrm{sft}_{\mu}\right).
\end{equation}
Evaluating (eq.3) at $c+\eta\cdot e_{\nu}$ and at $c$, one obtains
\[
\left(\partial_{\nu}(\tilde{g}A)_{\mu}\right)(c)=\eta^{-1}\cdot r_{\mu}\cdot\left[A_{\mu}\left(z+r_{\nu}\cdot\eta\cdot e_{\nu}-\mathrm{sft}_{\mu}\right)-A_{\mu}\left(z-\mathrm{sft}_{\mu}\right)\right].
\]
For $r_{\nu}=1$ the right side equals $r_{\mu}\cdot(\partial_{\nu}A_{\mu})(z-\mathrm{sft}_{\mu})$; for $r_{\nu}=-1$ it equals $-r_{\mu}\cdot(\partial_{\nu}A_{\mu})(z-\mathrm{sft}_{\mu}-\eta\cdot e_{\nu})$.
Hence, for all $\mu$ and $\nu$, the case $\mu=\nu$ included,
\begin{equation}\tag{eq.4}
\left(\partial_{\nu}(\tilde{g}A)_{\mu}\right)(c)=r_{\mu}\cdot r_{\nu}\cdot\left(\partial_{\nu}A_{\mu}\right)\left(\tilde{g}^{-1}c-\mathrm{sft}_{\mu}-\mathrm{sft}_{\nu}\right).
\end{equation}
The sites of the quotient on the left, in the sense of (G2.r4), are the points $c+i\cdot\eta\cdot e_{\mu}+j\cdot\eta\cdot e_{\nu}$ with $i,j\in\{0,1\}$.
Under $\tilde{g}^{-1}$ they go to
\[
z+i\cdot r_{\mu}\cdot\eta\cdot e_{\mu}+j\cdot r_{\nu}\cdot\eta\cdot e_{\nu}=\left(z-\mathrm{sft}_{\mu}-\mathrm{sft}_{\nu}\right)+i'\cdot\eta\cdot e_{\mu}+j'\cdot\eta\cdot e_{\nu},
\]
where $i':=i$ if $r_{\mu}=1$, $i':=1-i$ if $r_{\mu}=-1$, and $j'$ is defined in the same way with $r_{\nu}$; indeed $\mathrm{sft}_{\mu}+i\cdot r_{\mu}\cdot\eta\cdot e_{\mu}=i'\cdot\eta\cdot e_{\mu}$.
Thus $\tilde{g}^{-1}$ maps the set of the sites of the quotient $\left(\partial_{\nu}(\tilde{g}A)_{\mu}\right)(c)$ onto the set of the sites of the quotient $(\partial_{\nu}A_{\mu})(\tilde{g}^{-1}c-\mathrm{sft}_{\mu}-\mathrm{sft}_{\nu})$.
It follows that the first quotient lies inside $\tilde{g}Z$ if and only if the second lies inside $Z$, and that, for fixed $(\mu,\nu)$, the map $c\mapsto\tilde{g}^{-1}c-\mathrm{sft}_{\mu}-\mathrm{sft}_{\nu}$ is a bijection between the quotients inside $\tilde{g}Z$ and the quotients inside $Z$; by (eq.4), corresponding quotients have equal moduli.
Consider a H\"older difference with the base points $c^a$, $c^b$, $0<d(c^a,c^b)\le1$, and with the directions $(\mu,\nu)$.
The shift is the same at the two base points, so that by (a1.2)
\[
d\left(\tilde{g}^{-1}c^a-\mathrm{sft}_{\mu}-\mathrm{sft}_{\nu},\ \tilde{g}^{-1}c^b-\mathrm{sft}_{\mu}-\mathrm{sft}_{\nu}\right)=d(c^a,c^b),
\]
and by (eq.4) the modulus of the difference of the two quotients does not change.
Together with (b.1), applied to the bonds inside $Z$, this shows that each of the three suprema $|\tilde{g}A|_{\infty}$, $|\partial\tilde{g}A|_{\infty}$, $|\delta_{\alpha}\partial\tilde{g}A|_{\infty}$ on $\tilde{g}Z$ is the supremum of the same set of numbers as the corresponding supremum for $A$ on $Z$.

\textbf{(b.3) $\tau$, $d$, $U_B$.}
Let $A$ be a bond function on one of the lattices of (b.1), let $\mathrm{hm}$ be its mesh, and let $y$, $c$ be sites of this lattice with $d(y,c)<\ell/2$.
By (eq.1) and $(\tilde{g}A)(\tilde{g}b)=A(b)$,
\[
(\tilde{g}A)\left(\boldsymbol{\Gamma}^{\pi}(\tilde{g}y,\tilde{g}c)\right)=\sum_{b\in\boldsymbol{\Gamma}^{\pi}(y,c)}(\tilde{g}A)(\tilde{g}b)=A\left(\boldsymbol{\Gamma}^{\pi}(y,c)\right);
\]
averaging over the six orders gives $(\tau\,\tilde{g}A)(\tilde{g}y,\tilde{g}c)=(\tau A)(y,c)$, both sides being $0$ when $y=c$.
For every oriented bond $(c,c')$ of the lattice, $(d\boldsymbol{\lambda})(c,c')=\mathrm{hm}^{-1}\cdot\left(\boldsymbol{\lambda}(c')-\boldsymbol{\lambda}(c)\right)$: for $c'=c+\mathrm{hm}\cdot e_{\mu}$ this is $\partial_{\mu}\boldsymbol{\lambda}(c)$, and the right side is odd under reversal.
Therefore
\[
\left(\tilde{g}(d\boldsymbol{\lambda})\right)(c,c')=\mathrm{hm}^{-1}\cdot\left(\boldsymbol{\lambda}(\tilde{g}^{-1}c')-\boldsymbol{\lambda}(\tilde{g}^{-1}c)\right)=\left(d(\boldsymbol{\lambda}\circ\tilde{g}^{-1})\right)(c,c').
\]
Let $(x,x')\in\mathrm{Pzt}_k$, and let $B$ be a bond function on $T_k$, so that $\mathrm{hm}=\eta$.
By (a2.6), $(c,c')$ is a shortest lift, and (I.1.6) gives
\[
U_{\tilde{g}B}(\tilde{g}x,\tilde{g}x')=\exp\left\{i\,e_k\cdot\eta\cdot(\tau\,\tilde{g}B)(\tilde{g}c,\tilde{g}c')\right\}=\exp\left\{i\,e_k\cdot\eta\cdot(\tau B)(c,c')\right\}=U_B(x,x');
\]
for $c=c'$ both sides equal $1$.

\textbf{(b.4) Monomial matrices.}
A $4\times4$ matrix $\mathrm{Mx}$ is called monomial when there are a permutation $\mathrm{pe}=\mathrm{pe}_{\mathrm{Mx}}$ of $\{1,\dots,4\}$ and numbers $\mathrm{um}_{\beta}$ of modulus one such that $\mathrm{Mx}_{\beta\beta'}=\mathrm{um}_{\beta}\cdot\mathbf{1}[\beta'=\mathrm{pe}(\beta)]$.
Then $(\mathrm{Mx}\,v)_{\beta}=\mathrm{um}_{\beta}\cdot v_{\mathrm{pe}(\beta)}$, and so $|\mathrm{Mx}\,v|_c=\max_{\beta}|v_{\mathrm{pe}(\beta)}|=|v|_c$.
The class of monomial matrices is closed under products, inverses and transposition.
By (I.c6) the matrices $\gamma_1$, $\gamma_2$, $\gamma_3$, $\gamma_5$, $\mathsf{C}$, $\varkappa$ are monomial; consequently every $S\in\mathrm{Sp}$ is monomial, and so are $(S^{-1})^T$, $\mathsf{C}$, $-\mathsf{C}$ and $\varkappa$.
The norm $|f|_{\infty}$ is the supremum of $|f(x)|_c$ over the sites, resp.\ of $|F(x,x')|_c$ over the pairs; it does not change when every value is multiplied by one and the same monomial matrix.

\textbf{(b.5) The four maps.}
A homomorphism of algebras, or an antilinear homomorphism of rings, of $\mathrm{Gr}[\psi,\chi]$ into itself is given by prescribing odd elements as the images of the independent generators of Section~0.4.
The rules of Section~I.4.5 are written for all indices of the covering index sets.
They are called \emph{compatible with} (I.AP) when their right sides satisfy (0.a); in that case the rules hold at all indices as soon as they hold at the indices of the independent generators.

For $a_{\tilde{g}}$: by (a2.1), $\tilde{g}^{-1}(x+\ell\cdot n)=\tilde{g}^{-1}x+\ell\cdot r\cdot n$, and $\boldsymbol{\chi}(r\cdot n)=\boldsymbol{\chi}(n)$, since $r_{\mu}\cdot n_{\mu}$ and $n_{\mu}$ have the same parity.
So (0.a) yields
\begin{align*}
S\cdot\psi\left(\tilde{g}^{-1}(x+\ell\cdot n)\right)&=\boldsymbol{\chi}(n)\cdot S\cdot\psi(\tilde{g}^{-1}x),\\
(S^{-1})^T\cdot\bar{\psi}\left(\tilde{g}^{-1}(x+\ell\cdot n)\right)&=\boldsymbol{\chi}(n)\cdot(S^{-1})^T\cdot\bar{\psi}(\tilde{g}^{-1}x),
\end{align*}
and, $(\tilde{g}^{-1}x,\tilde{g}^{-1}x')$ being a pair of $\mathrm{Pzt}_k$ by (a2.4),
\begin{align*}
S\cdot\chi\left(\tilde{g}^{-1}(x+\ell\cdot n),\tilde{g}^{-1}(x'+\ell\cdot n)\right)&=S\cdot\chi\left(\tilde{g}^{-1}x+\ell\cdot r\cdot n,\ \tilde{g}^{-1}x'+\ell\cdot r\cdot n\right)\\
&=\boldsymbol{\chi}(n)\cdot S\cdot\chi(\tilde{g}^{-1}x,\tilde{g}^{-1}x'),\\
(S^{-1})^T\cdot\bar{\chi}\left(\tilde{g}^{-1}(x+\ell\cdot n),\tilde{g}^{-1}(x'+\ell\cdot n)\right)&=\boldsymbol{\chi}(n)\cdot(S^{-1})^T\cdot\bar{\chi}(\tilde{g}^{-1}x,\tilde{g}^{-1}x').
\end{align*}
For $G_{\Lambda}$: $\Lambda$ is a function on $T_k$, so $\Lambda(c+\ell\cdot n)=\Lambda(c)$.
For $\mathrm{Cg}$ and $\mathbf{K}$: these maps do not move the sites.

The images of the generators are as follows, $\mathrm{pe}$ and $\mathrm{um}$ being those of (b.4) for the matrix named in each case:
\begin{align*}
a_{\tilde{g}}\psi(x,\beta,\omega)&=\mathrm{um}_{\beta}\cdot\psi\left(\tilde{g}^{-1}x,\mathrm{pe}(\beta),\omega\right),\\
a_{\tilde{g}}\chi(x,x',\beta,\omega)&=\mathrm{um}_{\beta}\cdot\chi\left(\tilde{g}^{-1}x,\tilde{g}^{-1}x',\mathrm{pe}(\beta),\omega\right),
\end{align*}
with the matrix $S$ for $\omega=0$ and the matrix $(S^{-1})^T$ for $\omega=1$;
\begin{align*}
G_{\Lambda}\psi(x,\beta,\omega)&=\exp\left(\pm i\,e_k\cdot\Lambda(c)\right)\cdot\psi(x,\beta,\omega),\\
G_{\Lambda}\chi(x,x',\beta,\omega)&=\exp\left(\pm i\,e_k\cdot\Lambda(c)\right)\cdot\chi(x,x',\beta,\omega),
\end{align*}
with the sign $+$ for $\omega=0$ and $-$ for $\omega=1$, where $c=\pi(x)$ is the photon site of the first point;
\begin{align*}
\mathrm{Cg}\,\psi(x,\beta,0)&=\mathrm{um}_{\beta}\cdot\psi(x,\mathrm{pe}(\beta),1),&\mathrm{Cg}\,\psi(x,\beta,1)&=-\mathrm{um}_{\beta}\cdot\psi(x,\mathrm{pe}(\beta),0),\\
\mathrm{Cg}\,\chi(x,x',\beta,0)&=\mathrm{um}_{\beta}\cdot\chi(x,x',\mathrm{pe}(\beta),1),&\mathrm{Cg}\,\chi(x,x',\beta,1)&=-\mathrm{um}_{\beta}\cdot\chi(x,x',\mathrm{pe}(\beta),0),
\end{align*}
with the matrix $\mathsf{C}$;
\begin{align*}
\mathbf{K}\psi(x,\beta,\omega)&=\mathrm{um}_{\beta}\cdot\psi(x,\mathrm{pe}(\beta),\omega),&\mathbf{K}\chi(x,x',\beta,\omega)&=\mathrm{um}_{\beta}\cdot\chi(x,x',\mathrm{pe}(\beta),\omega),
\end{align*}
with the matrix $\varkappa$.
The number $e_k\cdot\Lambda(c)$ is real.
Every generator is therefore sent to a generator of the same family ($\psi$ to $\psi$, $\chi$ to $\chi$) times a number of modulus one; if the index on the right is not the index of an independent generator, (0.a) contributes one more sign.
The maps of the classes of indices,
\begin{itemize}
\item $(x,\beta,\omega)\mapsto\left(\tilde{g}^{-1}x,\mathrm{pe}(\beta),\omega\right)$ and $(x,x',\beta,\omega)\mapsto\left(\tilde{g}^{-1}x,\tilde{g}^{-1}x',\mathrm{pe}(\beta),\omega\right)$ for $a_{\tilde{g}}$,
\item the identity for $G_{\Lambda}$,
\item $(x,\beta,\omega)\mapsto(x,\mathrm{pe}(\beta),1-\omega)$ and $(x,x',\beta,\omega)\mapsto(x,x',\mathrm{pe}(\beta),1-\omega)$ for $\mathrm{Cg}$,
\item $(x,\beta,\omega)\mapsto(x,\mathrm{pe}(\beta),\omega)$ and $(x,x',\beta,\omega)\mapsto(x,x',\mathrm{pe}(\beta),\omega)$ for $\mathbf{K}$,
\end{itemize}
are bijections: $\mathrm{pe}$ is a permutation, $x\mapsto\tilde{g}^{-1}x$ is a bijection of $T^{\mathrm{f}}_k$ which involves the permutation $g\mapsto r\cdot g$ of the labels (a2.1), and $(x,x')\mapsto(\tilde{g}^{-1}x,\tilde{g}^{-1}x')$ is a bijection of $\mathrm{Pz}_k$ (a2.4).
Hence each of the four maps sends the set of the independent generators bijectively onto itself, up to factors of modulus one, and is an automorphism ($\mathbf{K}$: an antilinear automorphism of rings).
The weights $\nu(\xi)=\frac{\eta^3}{48}$ and $\nu(\boldsymbol{\zeta})=\left(\frac{\eta^3}{48}\right)^2$ are constant on each family.
A product of $n$ generators $\psi$ and $m$ generators $\chi$ is mapped to a product of the same kind; thus the bidegree is preserved.

\emph{Proof of (2.1).}
Let $\mathrm{Tm}$ be one of the four maps.
By what has been shown, there exist a bijection $(I_1,I_2)\mapsto(I_1',I_2')$ of the set of the pairs of subsets of the two index sets, with $\#|I_1'|=\#|I_1|$ and $\#|I_2'|=\#|I_2|$, and numbers $\mathrm{um}(I_1,I_2)$ of modulus one, namely products of the factors above and of the sign that puts the images into increasing order, such that
\[
\mathrm{Tm}\left(\psi(I_1)\chi(I_2)\right)=\mathrm{um}(I_1,I_2)\cdot\psi(I_1')\chi(I_2').
\]
Hence $c_{\mathrm{Tm}F}(I_1',I_2')=\mathrm{um}(I_1,I_2)\cdot c_F(I_1,I_2)$, with $\overline{c_F(I_1,I_2)}$ in place of $c_F(I_1,I_2)$ for $\mathbf{K}$, and by (I.5.1)
\[
\|\mathrm{Tm}F\|_{\mathsf{w}}=\sum_{I_1,I_2}\left|c_F(I_1,I_2)\right|\cdot\mathsf{w}_1^{\#|I_1'|}\cdot\mathsf{w}_2^{\#|I_2'|}=\|F\|_{\mathsf{w}}.
\]

\emph{The families $\Psi$, $W$: the paragraph of Proposition~16.1\,(b) following (2.1).}
The index sets and the weights are those of Section~I.3.6.
Assume $k\ge1$.
The covering index set of $\Psi$ and of $W$ consists of the $u=(y,\beta,\omega)$ with $y\in\mathbb{Z}^3$; put $u+v:=(y+v,\beta,\omega)$; the weight is $\nu(u)=1$.
By (I.AP), as in (0.a),
\[
\Psi(u+\ell\cdot n)=\boldsymbol{\chi}(n)\cdot\Psi(u),\qquad W(u+\ell\cdot n)=\boldsymbol{\chi}(n)\cdot W(u)\qquad(n\in\mathbb{Z}^3),
\]
and the independent generators are those with $y\in\mathrm{Fdm}\cap\mathbb{Z}^3$; this is a set of representatives of $\mathbb{Z}^3/\ell\mathbb{Z}^3=T^0_k$, because $\mathbb{Z}^3$ is a union of classes modulo $\ell\mathbb{Z}^3$ (Ar1).
In the rules of Section~I.4.5 the variable $y$ takes the place of $x$.
By (a1.2) with $i=0$, $\tilde{g}^{-1}$ is a bijection of $\mathbb{Z}^3$ and induces a bijection of $T^0_k$; by Section~0.1,
\[
\tilde{g}^{-1}(y+\ell\cdot n)=r\cdot(y+\ell\cdot n-a)=\tilde{g}^{-1}y+\ell\cdot r\cdot n;
\]
and $\boldsymbol{\chi}(r\cdot n)=\boldsymbol{\chi}(n)$.
Hence
\begin{align*}
S\cdot\Psi\left(\tilde{g}^{-1}(y+\ell\cdot n)\right)&=\boldsymbol{\chi}(n)\cdot S\cdot\Psi(\tilde{g}^{-1}y),\\
(S^{-1})^T\cdot\bar{\Psi}\left(\tilde{g}^{-1}(y+\ell\cdot n)\right)&=\boldsymbol{\chi}(n)\cdot(S^{-1})^T\cdot\bar{\Psi}(\tilde{g}^{-1}y),
\end{align*}
and likewise for $W$.
For $G_{\Lambda}$: $\Lambda(y+\ell\cdot n)=\Lambda(y)$, because $\mathbb{Z}^3\subset\mathbb{L}_k$ (Ar1) and $\Lambda$ is a function on $T_k$.
For $\mathrm{Cg}$ and $\mathbf{K}$: the sites are not moved.
The right sides of the four rules therefore satisfy the two relations displayed above, that is, the rules are compatible with (I.AP).
With $\mathrm{pe}$, $\mathrm{um}$ of (b.4) for the same matrices as before,
\begin{align*}
a_{\tilde{g}}\Psi(y,\beta,\omega)&=\mathrm{um}_{\beta}\cdot\Psi\left(\tilde{g}^{-1}y,\mathrm{pe}(\beta),\omega\right),\\
G_{\Lambda}\Psi(y,\beta,\omega)&=\exp\left(\pm i\,e_k\cdot\Lambda(y)\right)\cdot\Psi(y,\beta,\omega),\\
\mathrm{Cg}\,\Psi(y,\beta,0)&=\mathrm{um}_{\beta}\cdot\Psi(y,\mathrm{pe}(\beta),1),\qquad\mathrm{Cg}\,\Psi(y,\beta,1)=-\mathrm{um}_{\beta}\cdot\Psi(y,\mathrm{pe}(\beta),0),\\
\mathbf{K}\Psi(y,\beta,\omega)&=\mathrm{um}_{\beta}\cdot\Psi(y,\mathrm{pe}(\beta),\omega),
\end{align*}
and the same with $W$ in place of $\Psi$.
The maps of the classes of indices, namely $(y,\beta,\omega)\mapsto\left(\tilde{g}^{-1}y,\mathrm{pe}(\beta),\omega\right)$, the identity, $(y,\beta,\omega)\mapsto(y,\mathrm{pe}(\beta),1-\omega)$ and $(y,\beta,\omega)\mapsto(y,\mathrm{pe}(\beta),\omega)$, are bijections.
For $k=0$ the covering index set of $\Psi$ and of $W$ is that of $\psi$, the rules are those of $\psi$, and $\nu=\frac{1}{48}=\frac{\eta^3}{48}$; so everything established for the family $\psi$ holds for $\Psi$ and for $W$.

Let $\mathrm{Gr}[\psi,\chi,\Psi,W]$ denote the Grassmann algebra whose independent generators are those of the four families.
For this algebra, too, prescribing odd elements as the images of the independent generators defines a homomorphism; hence each of the four rules defines a homomorphism ($\mathbf{K}$: an antilinear homomorphism of rings).
This homomorphism sends the set of the independent generators of each family bijectively onto itself, up to factors of modulus one; it is therefore an automorphism, the weights $\nu$ are constant on each family, and the degree in each family is preserved.
As in the proof of (2.1), now with four index sets,
\[
\mathrm{Tm}\left(\psi(I_1)\chi(I_2)\Psi(I_3)W(I_4)\right)=\mathrm{um}(I_1,\dots,I_4)\cdot\psi(I_1')\chi(I_2')\Psi(I_3')W(I_4'),
\]
with $\#|I_i'|=\#|I_i|$, $\left|\mathrm{um}(I_1,\dots,I_4)\right|=1$ and a bijection $(I_1,\dots,I_4)\mapsto(I_1',\dots,I_4')$; since each family is mapped to itself, only the order inside each of the four products has to be restored, and this costs a sign.
Hence $\|\mathrm{Tm}F\|_{\mathsf{w}}=\|F\|_{\mathsf{w}}$ for all $F\in\mathrm{Gr}[\psi,\chi,\Psi,W]$ and all $\mathsf{w}=(\mathsf{w}_1,\dots,\mathsf{w}_4)$.
Each of these arguments treats the families one at a time, and a family whose generators are fixed has the identity as its map of indices and the factors $1$; therefore the assertions hold for a map that acts by its rule on some of the families and fixes the generators of the others.
In (b.6) the two relations for $\Psi$, $W$ yield the same factor $\boldsymbol{\chi}(n)$ as (0.a).

\textbf{(b.6) The vector $a$ in its class: the last paragraph of Proposition~16.1\,(b).}
If $a$ is replaced by $a+\ell\cdot n$, then, by Section~0.1 and (0.a), $a_{\tilde{g}}$ is replaced by its composition with the automorphism that multiplies every generator by $\boldsymbol{\chi}(n)$.
The even elements are left fixed by this automorphism.
In Section~0.1 the letter $a$ stands for an arbitrary vector of $M\mathbb{Z}^3$ in the class, and none of the proofs of Proposition~16.1\,(a1)--(d) selects a particular one; hence all assertions of Proposition~16.1\,(a1)--(d) hold for every vector of the class, in (d) for every vector $a$ of $LM\mathbb{Z}^3$ in the class, with the vector $a/L$ for $\tilde{g}'$.

\paragraph{(c)}
Let $\mathrm{Tm}$ be one of the three maps
\[
A\mapsto\tilde{g}A,\qquad A\mapsto-A,\qquad A\mapsto\overline{A}
\]
of the space of the complex bond functions on $T_k$.
Put $\mathrm{Tm}\,Z:=\tilde{g}Z$ for the first of them and $\mathrm{Tm}\,Z:=Z$ for the other two, and let $\mathrm{Tm}$ act by the same rule on the functions on the bonds inside a set $Z$ of sites; the image is a function on the bonds inside $\mathrm{Tm}\,Z$.
The map $\mathrm{Tm}$ is real-linear and bijective, its inverse is a map of the same kind, and it sends real functions to real functions.
The three norms of $\mathrm{Tm}\,A$ on $\mathrm{Tm}\,Z$ coincide with those of $A$ on $Z$:
for $\tilde{g}$ this is (b.2);
for $-A$ and for $\overline{A}$ every value on a bond, every difference quotient and every H\"older difference is multiplied by $-1$, respectively replaced by its complex conjugate.

\textbf{(c.1)}
The sets $R_k$, $G_k$ and $\mathfrak{X}_k$ are defined by strict upper bounds on norms taken among $|A|_{\infty}$, $|\partial A|_{\infty}$, $|\delta_{\alpha}\partial A|_{\infty}$ on $T_k$ \cite[Definition~5.16]{QEDI}.
Hence $\mathrm{Tm}$ maps each of them onto itself.
For a real number $\mathrm{tq}$, the set $\mathrm{tq}\cdot R_k$ of \cite[Definition~5.16\,(a)]{QEDI} consists of the bond functions on $T_k$ whose three norms are smaller than $\mathrm{tq}$ times the thresholds of $R_k$;
thus
\[
\mathrm{tq}\cdot R_k=\left\{\mathrm{tq}\cdot A:\ A\in R_k\right\}\quad\text{if }\mathrm{tq}>0,\qquad\mathrm{tq}\cdot R_k=\emptyset\quad\text{if }\mathrm{tq}\le0.
\]
Since $\mathrm{Tm}$ preserves the three norms, it maps $\mathrm{tq}\cdot R_k$ onto itself for every real $\mathrm{tq}$; this applies in particular to $\mathrm{cf}\cdot R_k$ with $0<\mathrm{cf}\le1$.
In $\tilde{R}_k^{-}$ the factor is $\mathrm{tq}=1-e_k^{\varepsilon}$.
If $e_k\ge1$, this factor is $\le0$, so that $\left(1-e_k^{\varepsilon}\right)\cdot R_k$ is empty;
then $\tilde{R}_k^{-}$ is empty, and the assertions of Proposition~16.1\,(c) on $\tilde{R}_k^{-}$ hold vacuously.

\textbf{(c.2)}
Let $A=A_0+A_1$ be an element of $\mathrm{cf}\cdot\tilde{R}_k$.
Thus $A_1\in\mathrm{cf}\cdot R_k$, the function $A_0$ is real, and for every $M$-cube $\square$ there are a real function $\boldsymbol{\lambda}_{\square}$ on the sites of $\square^{\dagger}$ and a real function $A'_{\square}$ on the bonds inside $\square^{\dagger}$ such that $A_0=A'_{\square}+d\boldsymbol{\lambda}_{\square}$ on the bonds inside $\square^{\dagger}$ and the three norms of $A'_{\square}$ on $\square^{\dagger}$ are smaller than $\mathrm{cf}$ times the thresholds;
as in \cite[Definition~5.16]{QEDI}, this representation is local, the functions $\boldsymbol{\lambda}_{\square}$ and $A'_{\square}$ being defined on $\square^{\dagger}$ only.
For a bond $(c,c')$ inside $\square^{\dagger}$ the value
\[
(d\boldsymbol{\lambda}_{\square})(c,c')=\eta^{-1}\cdot\left(\boldsymbol{\lambda}_{\square}(c')-\boldsymbol{\lambda}_{\square}(c)\right)
\]
involves values of $\boldsymbol{\lambda}_{\square}$ at sites of $\square^{\dagger}$ only.
Therefore the set $\mathrm{cf}\cdot\tilde{R}_k$ remains the same if $\boldsymbol{\lambda}_{\square}$ is required to be a real function on all sites of $T_k$: such a function may be restricted to $\square^{\dagger}$, and a function on the sites of $\square^{\dagger}$ may be extended by $0$.
The same holds for the sets $\square^{\sim5}$ in the definition of $\tilde{G}_k$.
Now $\mathrm{Tm}\,A=\mathrm{Tm}\,A_0+\mathrm{Tm}\,A_1$, where $\mathrm{Tm}\,A_1\in\mathrm{cf}\cdot R_k$ by (c.1) and $\mathrm{Tm}\,A_0$ is real.

Consider first $\tilde{g}$.
Let $\square$ be an $M$-cube and $\square':=\tilde{g}^{-1}\square$, which is an $M$-cube by (a1.4).
By (a1.8), $\tilde{g}\left(\square'^{\dagger}\right)=\square^{\dagger}$, so that $\boldsymbol{\lambda}_{\square'}\circ\tilde{g}^{-1}$ is a real function on the sites of $\square^{\dagger}$.
The computation of (b.3) involves only the values of $\boldsymbol{\lambda}_{\square'}$ at the two end points of $\tilde{g}^{-1}b$; applied to the bonds $b$ inside $\square^{\dagger}$, it gives there
\[
\tilde{g}A_0=\tilde{g}A'_{\square'}+d\left(\boldsymbol{\lambda}_{\square'}\circ\tilde{g}^{-1}\right),
\]
and by (b.2) the three norms of $\tilde{g}A'_{\square'}$ on $\square^{\dagger}$ are those of $A'_{\square'}$ on $\square'^{\dagger}$.
For $A\mapsto-A$ one has
\[
-A_0=\left(-A'_{\square}\right)+d\left(-\boldsymbol{\lambda}_{\square}\right),
\]
and for $A\mapsto\overline{A}$ one has $\overline{A}=A_0+\overline{A_1}$.
In each case $\mathrm{Tm}\,A$ lies in $\mathrm{cf}\cdot\tilde{R}_k$.

In the case of $\tilde{R}_k^{-}$ one has $A=A_0+A_1$ with $A_1\in\left(1-e_k^{\varepsilon}\right)\cdot R_k$ and with $A_0$ as above for $\mathrm{cf}=1$.
By (c.1), $\mathrm{Tm}\,A_1\in\left(1-e_k^{\varepsilon}\right)\cdot R_k$.
The formulas just obtained for $\tilde{g}A_0$, $-A_0$ and $\overline{A}$, read with $\mathrm{cf}=1$, show that $\mathrm{Tm}\,A_0$ has the property required of the real part.
So $\mathrm{Tm}\,A\in\tilde{R}_k^{-}$.

In the case of $\tilde{G}_k$ one has $A=A_0+A_1$ with $A_1\in G_k$ and $A_0$ real, and for every $M$-cube $\square$
\[
A_0=A'_{\square}+d\boldsymbol{\lambda}_{\square}\qquad\text{on the bonds inside }\square^{\sim5},
\]
where $\boldsymbol{\lambda}_{\square}$ is a real function on the sites of $\square^{\sim5}$, $A'_{\square}$ is a real function on the bonds inside $\square^{\sim5}$, and $|A'_{\square}|_{\infty}<e_k^{-1+\varepsilon}$ on $\square^{\sim5}$.
By (c.1), $\mathrm{Tm}\,A_1\in G_k$.
For $\square'=\tilde{g}^{-1}\square$ one has $\tilde{g}\left(\square'^{\sim5}\right)=\square^{\sim5}$ ((a1.8) with $i=5$); on the bonds inside $\square^{\sim5}$
\[
\tilde{g}A_0=\tilde{g}A'_{\square'}+d\left(\boldsymbol{\lambda}_{\square'}\circ\tilde{g}^{-1}\right),
\]
and by (b.2) the number $|\tilde{g}A'_{\square'}|_{\infty}$ on $\square^{\sim5}$ equals $|A'_{\square'}|_{\infty}$ on $\square'^{\sim5}$.
Further $-A_0=\left(-A'_{\square}\right)+d\left(-\boldsymbol{\lambda}_{\square}\right)$ and $\overline{A}=A_0+\overline{A_1}$.
So $\mathrm{Tm}\,A\in\tilde{G}_k$.

\textbf{(c.3) Gauge transformations.}
Let $\boldsymbol{\lambda}$ be a real function on $T_k$.
Then
\[
A-d\boldsymbol{\lambda}=\left(A_0-d\boldsymbol{\lambda}\right)+A_1,
\]
the function $A_0-d\boldsymbol{\lambda}$ is real, and on the bonds inside $\square^{\dagger}$ (respectively $\square^{\sim5}$)
\[
A_0-d\boldsymbol{\lambda}=A'_{\square}+d\left(\boldsymbol{\lambda}_{\square}-\boldsymbol{\lambda}\right),
\]
where $\boldsymbol{\lambda}_{\square}-\boldsymbol{\lambda}$ is formed on the sites of $\square^{\dagger}$ (respectively $\square^{\sim5}$).
Neither $A'_{\square}$ nor $A_1$ is changed.
Hence $A-d\boldsymbol{\lambda}$ belongs to each of the sets $\mathrm{cf}\cdot\tilde{R}_k$, $\tilde{R}_k^{-}$, $\tilde{G}_k$ to which $A$ belongs.

\textbf{(c.4) Real constants.}
Let $\mathrm{kc}\in\mathbb{R}^3$, and let $\mathrm{Ac}$ be the bond function with $(\mathrm{Ac})_{\mu}(c)=\mathrm{kc}_{\mu}$ for all $c$ and $\mu$.
Then $A+\mathrm{Ac}=\left(A_0+\mathrm{Ac}\right)+A_1$, and $A_0+\mathrm{Ac}$ is real.
Let $\square$ be an $M$-cube with the centre $z$.
Put $Z:=\square^{\dagger}$ and $\mathrm{sid}:=(2L+1)\cdot M$ (respectively $Z:=\square^{\sim5}$ and $\mathrm{sid}:=11M$), let $\tilde{z}$ be a point of $M\mathbb{Z}^3$ over $z$, and let
\[
\mathsf{Z}:=\left\{\tilde{c}\in\mathbb{L}_k:\ |\tilde{c}-\tilde{z}|<\frac{\mathrm{sid}}{2}\right\}.
\]
By (enl), $Z=\left\{c\in T_k:\ d(c,z)<\frac{\mathrm{sid}}{2}\right\}$, and $d(c,z)$ is the least of the numbers $|\tilde{c}-\tilde{z}|$, $\tilde{c}$ ranging over the points of $\mathbb{L}_k$ over $c$; hence $\mathrm{pr}(\mathsf{Z})=Z$.
By (Ar2), $\operatorname{sup-diameter}(\mathsf{Z})=\mathrm{sid}-\eta$, and by (Ar3), $\mathrm{sid}+\eta<\ell$.
So $\mathsf{Z}$ is a lift of $Z$, with $\operatorname{sup-diameter}(\mathsf{Z})+\eta<\ell$.
For $c\in Z$ let $\tilde{c}$ be the point of $\mathsf{Z}$ over $c$ (Lft1), and put
\[
\boldsymbol{\lambda}_{\mathrm{kc}}(c):=\mathrm{kc}\cdot\left(\tilde{c}-\tilde{z}\right);
\]
this is a real function on the sites of $Z$.
For a bond $(c,c+\eta\cdot e_{\mu})$ inside $Z$, (Lft2) with $\mathrm{sd}=\eta$ shows that the points of $\mathsf{Z}$ over $c+\eta\cdot e_{\mu}$ and over $c$ differ by $\eta\cdot e_{\mu}$.
Consequently
\[
\left(d\boldsymbol{\lambda}_{\mathrm{kc}}\right)_{\mu}(c)=\eta^{-1}\cdot\left(\mathrm{kc}\cdot\eta e_{\mu}\right)=\mathrm{kc}_{\mu},
\]
that is, $\mathrm{Ac}=d\boldsymbol{\lambda}_{\mathrm{kc}}$ on the bonds inside $Z$, and there
\[
A_0+\mathrm{Ac}=A'_{\square}+d\left(\boldsymbol{\lambda}_{\square}+\boldsymbol{\lambda}_{\mathrm{kc}}\right),
\]
where $\boldsymbol{\lambda}_{\square}+\boldsymbol{\lambda}_{\mathrm{kc}}$ is a real function on the sites of $Z$ (c.2).
Neither $A'_{\square}$ nor $A_1$ is changed; so $A+\mathrm{Ac}$ belongs to each of the sets $\mathrm{cf}\cdot\tilde{R}_k$, $\tilde{R}_k^{-}$, $\tilde{G}_k$ to which $A$ belongs.

\paragraph{(d)}
\textbf{(d.0) The side $\ell/L$; the objects of level $k+1$.}
The number $j_{\ell}:=\ell/(L\cdot M)$ is an odd integer $\ge13$.
Hence $\ell/L=j_{\ell}\cdot M$ is an odd multiple of $M$, and $\ell/L\ge13M>2M$.
The objects of level $k+1$ that are named in Proposition~16.1\,(d) are given by the following formulas; after each of them the property of the side that it uses is stated.
\begin{itemize}
\item
The torus $T_{k+1}=\mathbb{L}_{k+1}/(\ell/L)\mathbb{Z}^3$ and $T^{\mathrm{f}}_{k+1}=T_{k+1}\times O_{\mathrm{cub}}$ (Sections~I.1.1 and~I.1.3).
Here $(\ell/L)\mathbb{Z}^3$ is a subgroup of $\mathbb{L}_{k+1}=\eta_{k+1}\mathbb{Z}^3$; what is used is that the side is an integer.
\item
The lattices $L^{-i}\mathbb{Z}^3/(\ell/L)\mathbb{Z}^3$, $-1\le i\le k+1$ (for $i=k+1$ this is $T_{k+1}$), their oriented bonds $(c,c')$ with $c'-c=\pm L^{-i}e_{\mu}$, and their bond functions, that is, the complex functions of the oriented bonds that are odd under reversal (Sections~I.1.1 and~I.4.3).
Here $(\ell/L)\mathbb{Z}^3$ is a subgroup of $M\mathbb{Z}^3$, and $M\mathbb{Z}^3$ is contained in $L\mathbb{Z}^3$, hence in every $L^{-i}\mathbb{Z}^3$ (Ar1); what is used is that the side is a multiple of $M$.
The six vectors $\pm L^{-i}e_{\mu}$ are pairwise different modulo $(\ell/L)\mathbb{Z}^3$; for this, $2L^{-i}\le2L\le2M$ (Ar1) and $2M<\ell/L$ are used.
\item
The $M$-cubes $\left\{c:\ d(c,z)<\frac{M}{2}\right\}$, $z\in M\mathbb{Z}^3/(\ell/L)\mathbb{Z}^3$.
Here $(\ell/L)\mathbb{Z}^3$ is a subgroup of $M\mathbb{Z}^3$; what is used is that the side is a multiple of $M$.
(The number $M/\eta_{k+1}$ is odd, so that no site lies on a face.)
\item
The set $\mathrm{Pzt}_{k+1}$ (Section~I.1.7), whose definition does not involve the side, and $\mathrm{Pz}_{k+1}=\mathrm{Pzt}_{k+1}/(\ell/L)\mathbb{Z}^3$.
A pair of distinct sites of $T^{\mathrm{f}}_{k+1}$ whose photon sites have distance at most $M$ has exactly one class of lifts in $\mathrm{Pzt}_{k+1}$: if $|c-c'|\le M$, then every other vector of the class of $c-c'$ has sup norm at least $\ell/L-M>M$.
What is used is $2M<\ell/L$.
\item
The generators $\psi(\xi)$, $\chi(\boldsymbol{\zeta})$ of level $k+1$, with (I.AP) (Section~I.3.6).
Their index sets are the two sets just named, and the signs in (I.AP) belong to the translations by $(\ell/L)\cdot e_{\mu}$.
\item
The spaces of spinors (Section~I.1.5)
\begin{align*}
\mathfrak{S}^{\mathrm{ch}}(\eta_{k+1},\ell/L)&=\left\{f':\mathbb{L}^{\mathrm{f}}_{k+1}\to\mathbb{C}^4:\ f'\left(x+\tfrac{\ell}{L}e_{\mu}\right)=-f'(x)\right\},\\
\mathfrak{S}(1,\ell/L)&=\left\{f':\mathbb{Z}^3\to\mathbb{C}^4:\ f'\left(y+\tfrac{\ell}{L}e_{\mu}\right)=-f'(y)\right\}.
\end{align*}
Here $(\ell/L)\mathbb{Z}^3$ is a subgroup of $\mathbb{L}_{k+1}$ and of $\mathbb{Z}^3$; what is used is that the side is an integer.
As for the matrix kernels in Section~0.4, one has for $n\in\mathbb{Z}^3$
\[
f'\left(x+\tfrac{\ell}{L}n\right)=\boldsymbol{\chi}(n)\cdot f'(x),\qquad f'\left(y+\tfrac{\ell}{L}n\right)=\boldsymbol{\chi}(n)\cdot f'(y).
\]
\item
The group
\[
\mathbb{G}_{k+1}=\left\{(S,a'):\ S\in\mathrm{Sp},\ a'\in M\mathbb{Z}^3/(\ell/L)\mathbb{Z}^3\right\}
\]
with the product of Section~I.4.2.
Here $(\ell/L)\mathbb{Z}^3$ is a subgroup of $M\mathbb{Z}^3$ that is mapped onto itself by every $r(S)$.
\item
The homomorphism $S_L$ of (G2.r5), with
\[
S_L\psi(x,\beta,\omega)=L^{-1}\cdot\psi(x/L,\beta,\omega),\qquad S_L\chi(x,x',\beta,\omega)=L^{-1-\alpha}\cdot\chi(x/L,x'/L,\beta,\omega).
\]
The map $x\mapsto x/L=(c/L,g)$ is a bijection of $\mathbb{L}^{\mathrm{f}}_k$ onto $\mathbb{L}^{\mathrm{f}}_{k+1}$, and $(x+\ell n)/L=x/L+(\ell/L)n$.
For $(x,x')\in\mathrm{Pzt}_k$ the points $x/L$ and $x'/L$ are distinct, and, since $\eta_{k+1}=\eta/L$, the number $d_X(x/L,x'/L)$ formed at level $k+1$ is
\[
\max\left(\frac{|c-c'|}{L},\eta_{k+1}\right)=\frac{d_X(x,x')}{L}\le\frac{M}{L}\le M;
\]
so $(x/L,x'/L)\in\mathrm{Pzt}_{k+1}$.
If $(x,x')$ is translated by $\ell n$, the right sides $L^{-1}\cdot\psi(x/L,\beta,\omega)$ and $L^{-1-\alpha}\cdot\chi(x/L,x'/L,\beta,\omega)$ of (G2.r5) are multiplied by $\boldsymbol{\chi}(n)$, by (I.AP) at level $k+1$ on the side $\ell/L$.
So the rule (G2.r5) is compatible with (I.AP) in the sense of (b.5), and there is exactly one homomorphism $S_L$ that satisfies (G2.r5) at all indices.
\end{itemize}
The arguments of (a1.1), (a1.2), (a2.1), (a2.2) and (a2.4), the part of (b.1) on the oriented bonds and on $\tilde{g}A$, the formula of (b.3) for $(d\boldsymbol{\lambda})(c,c')$, and the arguments of (b.4), (b.5) and (b.6) use only one property of the side: the translation part is a class of $M\mathbb{Z}^3$ modulo the side times $\mathbb{Z}^3$, that is, the side is a multiple of $M$.
Hence they hold at level $k+1$ on the side $\ell/L$, for every element $(S,a')$ of $\mathbb{G}_{k+1}$.
Such an element induces bijections of the lattices $L^{-i}\mathbb{Z}^3/(\ell/L)\mathbb{Z}^3$, $-1\le i\le k+1$ (among them $T_{k+1}$), of $T^{\mathrm{f}}_{k+1}$ and of $\mathrm{Pz}_{k+1}$, and the four maps $a_{(S,a')}$, $G_{\Lambda}$, $\mathrm{Cg}$, $\mathbf{K}$ of level $k+1$ are automorphisms, antilinear in the case of $\mathbf{K}$, that are compatible with (I.AP).

\emph{The actions of $\tilde{g}'=(S,a')\in\mathbb{G}_{k+1}$ on bond functions and on spinors} (Section~I.4.3; $a'$ is a vector of $M\mathbb{Z}^3$, and $r=r(S)$).
For an oriented bond $b=(c,c')$ of $L^{-i}\mathbb{Z}^3/(\ell/L)\mathbb{Z}^3$,
\[
\tilde{g}'c'-\tilde{g}'c=r(c'-c)=\pm r_{\mu}L^{-i}e_{\mu}\qquad\text{modulo }(\ell/L)\mathbb{Z}^3.
\]
Hence $\tilde{g}'b:=(\tilde{g}'c,\tilde{g}'c')$ is an oriented bond, and $b\mapsto\tilde{g}'b$ is a bijection of the set of the oriented bonds, with the inverse $b\mapsto\tilde{g}'^{-1}b$, which commutes with the reversal; therefore $(\tilde{g}'A)(b):=A(\tilde{g}'^{-1}b)$ is a bond function on this lattice.
For $f'\in\mathfrak{S}^{\mathrm{ch}}(\eta_{k+1},\ell/L)$ put $(\tilde{g}'f')(x):=Sf'(\tilde{g}'^{-1}x)$.
By (a2.1), $\tilde{g}'^{-1}\left(x+\frac{\ell}{L}n\right)=\tilde{g}'^{-1}x+\frac{\ell}{L}rn$, and $\boldsymbol{\chi}(rn)=\boldsymbol{\chi}(n)$ (b.5); hence
\[
(\tilde{g}'f')\left(x+\tfrac{\ell}{L}n\right)=Sf'\left(\tilde{g}'^{-1}x+\tfrac{\ell}{L}rn\right)=\boldsymbol{\chi}(n)\cdot(\tilde{g}'f')(x),
\]
and $\tilde{g}'f'\in\mathfrak{S}^{\mathrm{ch}}(\eta_{k+1},\ell/L)$.
For $f'\in\mathfrak{S}(1,\ell/L)$ put $(\tilde{g}'f')(y):=Sf'(\tilde{g}'^{-1}y)$; the same computation applies with $y\in\mathbb{Z}^3$ in place of $x$, since $\tilde{g}'^{-1}$ maps $\mathbb{Z}^3$ onto itself ((a1.2) with $i=0$).
Of the side, only the fact that it is a multiple of $M$ is used.

The parenthetical sentence of Proposition~16.1\,(d) names the objects of level $k+1$ for which a side in $\mathrm{Cl}$ is required; they are the following.
For the polymers $D_{k+1}$ one uses \cite[Proposition~5.58 and Theorem~13.2]{QEDI}, which are stated for sides in $\mathrm{Cl}$.
The domains of level $k+1$ are defined with the sets $\square^{\dagger}$, and the step (c.4), performed at level $k+1$, uses
\[
(2L+1)\cdot M+\eta_{k+1}<\ell/L,
\]
which does not follow from $\ell/L\ge13M$.
The reblocked families $LE'$ of Proposition~16.3\,(e) are indexed by $D_{k+1}$.
None of these objects occurs in (d1)--(d4).

In (d1)--(d4), $\tilde{g}=(S,a)\in\mathbb{G}_k^{L}$ with a vector $a$ of $LM\mathbb{Z}^3$, and $\tilde{g}':=(S,a/L)$ with the vector $a/L$ of $M\mathbb{Z}^3$; for both, $r=r(S)$.

\textbf{(d1)}
The map $a\mapsto a/L$ is an isomorphism of the group $LM\mathbb{Z}^3$ onto $M\mathbb{Z}^3$ that maps $\ell\mathbb{Z}^3$ onto $(\ell/L)\mathbb{Z}^3$ and commutes with every $r(S)$.
So it induces an isomorphism of $LM\mathbb{Z}^3/\ell\mathbb{Z}^3$ onto $M\mathbb{Z}^3/(\ell/L)\mathbb{Z}^3$.
For $(S,a)$ and $(S_1,a_1)$ in $\mathbb{G}_k^{L}$,
\[
\left((S,a)(S_1,a_1)\right)'=\left(SS_1,\frac{a+r(S)a_1}{L}\right)=\left(SS_1,\frac{a}{L}+r(S)\frac{a_1}{L}\right)=(S,a)'(S_1,a_1)'.
\]
Hence $\tilde{g}\mapsto\tilde{g}'$ is a bijective homomorphism of $\mathbb{G}_k^{L}$ onto $\mathbb{G}_{k+1}$.

\textbf{(d2)}
One has
\[
\frac{\tilde{g}c}{L}=\frac{rc+a}{L}=r\frac{c}{L}+\frac{a}{L}=\tilde{g}'\frac{c}{L},\qquad\frac{\tilde{g}x}{L}=\left(\frac{rc+a}{L},rg\right)=\tilde{g}'\left(\frac{c}{L},g\right)=\tilde{g}'\frac{x}{L}.
\]
For a set $Y$ of sites of $T_{k+1}$ let $LY:=\left\{c\in T_k:\ c/L\in Y\right\}$.
A site $c$ lies in $\tilde{g}(LY)$ if and only if $(\tilde{g}^{-1}c)/L\in Y$.
By (d1), $(\tilde{g}^{-1})'=(S^{-1},-ra/L)=(\tilde{g}')^{-1}$, so that $(\tilde{g}^{-1}c)/L=\tilde{g}'^{-1}(c/L)$.
Hence $c$ lies in $\tilde{g}(LY)$ if and only if $c/L$ lies in $\tilde{g}'Y$, that is, if and only if $c$ lies in $L(\tilde{g}'Y)$.

\textbf{(d3)}
Let $0\le i\le k+1$, let $A$ be a bond function on $L^{-i}\mathbb{Z}^3/(\ell/L)\mathbb{Z}^3$, and let $b=(c,c')$ be a bond of $L^{-(i-1)}\mathbb{Z}^3/\ell\mathbb{Z}^3$; then $b/L:=(c/L,c'/L)$ is a bond of the first lattice.
Since $-1\le i-1\le k$, the element $\tilde{g}$ acts on the bond functions of the second lattice, and by (d.0) the element $\tilde{g}'$ acts on those of the first lattice and on the two spaces of spinors below.
By (d2), applied to $\tilde{g}^{-1}$,
\begin{align*}
(\tilde{g}'A)_L(b)&=L^{-\frac{1}{2}}\cdot(\tilde{g}'A)(b/L)=L^{-\frac{1}{2}}\cdot A\left(\tilde{g}'^{-1}(b/L)\right)=L^{-\frac{1}{2}}\cdot A\left((\tilde{g}^{-1}b)/L\right)\\
&=A_L(\tilde{g}^{-1}b)=\left(\tilde{g}(A_L)\right)(b).
\end{align*}
For $f'\in\mathfrak{S}^{\mathrm{ch}}(\eta_{k+1},\ell/L)$,
\[
(\tilde{g}'f')_L(x)=L^{-1}\cdot Sf'\left(\tilde{g}'^{-1}(x/L)\right)=L^{-1}\cdot Sf'\left((\tilde{g}^{-1}x)/L\right)=Sf'_L(\tilde{g}^{-1}x)=\left(\tilde{g}(f'_L)\right)(x).
\]
For $f'\in\mathfrak{S}(1,\ell/L)$, with $f'_L(y)=L^{-1}\cdot f'(y/L)$ for $y\in L\mathbb{Z}^3$, the same computation applies with $y$ in place of $x$; here $\tilde{g}$ maps $L\mathbb{Z}^3$ onto itself (a1.2).

\textbf{(d4)}
In each of the relations both sides are homomorphisms of algebras (in the relation with $\mathbf{K}$: antilinear homomorphisms of rings) from $\mathrm{Gr}[\psi,\chi]$ of level $k$ to $\mathrm{Gr}[\psi,\chi]$ of level $k+1$.
So it suffices to compare them on the generators.
By (d2),
\begin{align*}
S_La_{\tilde{g}}\psi(x)&=S_L\left(S\psi(\tilde{g}^{-1}x)\right)=L^{-1}\cdot S\psi\left((\tilde{g}^{-1}x)/L\right)=L^{-1}\cdot S\psi\left(\tilde{g}'^{-1}(x/L)\right)\\
&=a_{\tilde{g}'}\left(L^{-1}\cdot\psi(x/L)\right)=a_{\tilde{g}'}S_L\psi(x),\\
S_La_{\tilde{g}}\bar{\psi}(x)&=L^{-1}\cdot\left(S^{-1}\right)^T\bar{\psi}\left((\tilde{g}^{-1}x)/L\right)=L^{-1}\cdot\left(S^{-1}\right)^T\bar{\psi}\left(\tilde{g}'^{-1}(x/L)\right)=a_{\tilde{g}'}S_L\bar{\psi}(x),\\
S_La_{\tilde{g}}\chi(x,x')&=L^{-1-\alpha}\cdot S\chi\left((\tilde{g}^{-1}x)/L,(\tilde{g}^{-1}x')/L\right)\\
&=L^{-1-\alpha}\cdot S\chi\left(\tilde{g}'^{-1}(x/L),\tilde{g}'^{-1}(x'/L)\right)=a_{\tilde{g}'}S_L\chi(x,x'),\\
S_La_{\tilde{g}}\bar{\chi}(x,x')&=L^{-1-\alpha}\cdot\left(S^{-1}\right)^T\bar{\chi}\left(\tilde{g}'^{-1}(x/L),\tilde{g}'^{-1}(x'/L)\right)=a_{\tilde{g}'}S_L\bar{\chi}(x,x').
\end{align*}
Further,
\begin{align*}
S_L\mathrm{Cg}\,\psi(x)&=S_L\left(\mathsf{C}\bar{\psi}(x)\right)=L^{-1}\cdot\mathsf{C}\bar{\psi}(x/L)=\mathrm{Cg}\left(L^{-1}\cdot\psi(x/L)\right)=\mathrm{Cg}\,S_L\psi(x),\\
S_L\mathrm{Cg}\,\bar{\psi}(x)&=-L^{-1}\cdot\mathsf{C}\psi(x/L)=\mathrm{Cg}\,S_L\bar{\psi}(x),\\
S_L\mathrm{Cg}\,\chi(x,x')&=L^{-1-\alpha}\cdot\mathsf{C}\bar{\chi}(x/L,x'/L)=\mathrm{Cg}\,S_L\chi(x,x'),\\
S_L\mathrm{Cg}\,\bar{\chi}(x,x')&=-L^{-1-\alpha}\cdot\mathsf{C}\chi(x/L,x'/L)=\mathrm{Cg}\,S_L\bar{\chi}(x,x').
\end{align*}
The numbers $L^{-1}$ and $L^{-1-\alpha}$ are real, so that $\mathbf{K}$ does not change them; hence
\begin{align*}
S_L\mathbf{K}\psi(x)&=L^{-1}\cdot\varkappa\psi(x/L)=\mathbf{K}\left(L^{-1}\cdot\psi(x/L)\right)=\mathbf{K}S_L\psi(x),\\
S_L\mathbf{K}\bar{\psi}(x)&=L^{-1}\cdot\varkappa\bar{\psi}(x/L)=\mathbf{K}S_L\bar{\psi}(x),\\
S_L\mathbf{K}\chi(x,x')&=L^{-1-\alpha}\cdot\varkappa\chi(x/L,x'/L)=\mathbf{K}S_L\chi(x,x'),\\
S_L\mathbf{K}\bar{\chi}(x,x')&=L^{-1-\alpha}\cdot\varkappa\bar{\chi}(x/L,x'/L)=\mathbf{K}S_L\bar{\chi}(x,x').
\end{align*}
Let $\Lambda$ be a real function on the sites of $T_{k+1}$, and $\Lambda_L(c)=L^{\frac{1}{2}}\cdot\Lambda(c/L)$, a real function on $T_k$.
Then
\[
e_k\cdot\Lambda_L(c)=e_kL^{\frac{1}{2}}\cdot\Lambda(c/L)=e_{k+1}\cdot\Lambda(c/L).
\]
Since the first point of the pair $(x/L,x'/L)$ has the photon site $c/L$, it follows that
\begin{align*}
G_{\Lambda}S_L\psi(x)&=L^{-1}\cdot\exp\left(ie_{k+1}\cdot\Lambda(c/L)\right)\psi(x/L)=S_L\left(\exp\left(ie_k\cdot\Lambda_L(c)\right)\psi(x)\right)=S_LG_{\Lambda_L}\psi(x),\\
G_{\Lambda}S_L\bar{\psi}(x)&=L^{-1}\cdot\exp\left(-ie_{k+1}\cdot\Lambda(c/L)\right)\bar{\psi}(x/L)=S_LG_{\Lambda_L}\bar{\psi}(x),\\
G_{\Lambda}S_L\chi(x,x')&=L^{-1-\alpha}\cdot\exp\left(ie_{k+1}\cdot\Lambda(c/L)\right)\chi(x/L,x'/L)=S_LG_{\Lambda_L}\chi(x,x'),\\
G_{\Lambda}S_L\bar{\chi}(x,x')&=L^{-1-\alpha}\cdot\exp\left(-ie_{k+1}\cdot\Lambda(c/L)\right)\bar{\chi}(x/L,x'/L)=S_LG_{\Lambda_L}\bar{\chi}(x,x').
\end{align*}
For a bond $(c,c')$ of $T_k$, with $\eta_{k+1}=\eta/L$,
\begin{align*}
(d\Lambda)_L(c,c')&=L^{-\frac{1}{2}}\cdot\eta_{k+1}^{-1}\cdot\left(\Lambda(c'/L)-\Lambda(c/L)\right)=\eta^{-1}L^{\frac{1}{2}}\cdot\left(\Lambda(c'/L)-\Lambda(c/L)\right)\\
&=\left(d(\Lambda_L)\right)(c,c').
\end{align*}

\textbf{Level~$0$.}
At $k=0$ one has $\eta=1$ and $\mathbb{L}_0=\mathbb{Z}^3$, the lattices of (a1) are $\mathbb{Z}^3$ and $L\mathbb{Z}^3$ ($i=0$ and $i=-1$), $\Delta_y=\{y\}$, $T^{\mathrm{f}}_0=T^0_0\times O_{\mathrm{cub}}$, and the number $e_k$ is $q$.
Every pair has $d_X\ge1$: there is no short pair (the assertion of (a2.5) on short pairs is void), and a pair inside a cluster has $d_X=1$.
The weights are $\nu(\xi)=\frac{1}{48}$ and $\nu(\boldsymbol{\zeta})=\left(\frac{1}{48}\right)^2$.
In (d), $S_L$ maps the algebra of level $0$ to that of level $1$, with $x/L\in L^{-1}\mathbb{Z}^3\times O_{\mathrm{cub}}$.
With these values every argument above applies at $k=0$ as well; for the families $\Psi$, $W$ the case $k=0$ is treated separately in (b.5).

No constant enters these arguments, and the assertions of Proposition~16.1 hold for every $k\ge0$, every $\ell\in\mathrm{Cl}$, every $e_k>0$ and every $\tilde{g}$.
\qed

\subsection*{2. Proposition 16.2}

\subsubsection*{Proof of Proposition 16.2}

\begin{proof}
\textbf{Notation.}
Put
\[
K_X(x,x'):=\mathscr{E}_{20}(X;0,0;x,x'),
\]
the matrix of Section~I.5.1; it is a matrix kernel in the sense of Section~0.4.
The element $\mathscr{E}(X;0,0)$ is balanced, hence so is its part of bidegree $(2,0)$, and Lemma~0.3\,(b) gives
\[
\mathscr{E}^{[2,0]}(X;0,0)=\mathrm{Bil}[K_X].
\]
For a small polymer $X$ with a lift $\tilde{X}$ (Lemma~0.2\,(c)) the matrix $\mathrm{al}_2(\mathscr{E},X)$ of Section~I.5.6, written out, is
\begin{align}
\mathrm{al}_2(\mathscr{E},X)&=\operatorname{Vol}(X)^{-1}\int_{\tilde{X}}\mathrm{d}x\int_{\tilde{X}}\mathrm{d}x'\,K_X(x,x')\notag\\
&=\operatorname{Vol}(X)^{-1}\cdot\left(\frac{\eta^3}{48}\right)^2\sum_{c,c'\in\tilde{X}}\ \sum_{g,g'\in O_{\mathrm{cub}}}K_X\left((c,g),(c',g')\right).\tag{eq.5}
\end{align}

\textbf{Step 1 (the four rules of Proposition~16.2\,(v)).}
Let $K$ be a matrix kernel, $\tilde{g}=(S,a)\in\mathbb{G}_k$, and $\Lambda$ a real function on $T_k$, and let $K^{\tilde{g}}$, $K^{\Lambda}$, $K^{\mathrm{Cg}}$, $K^{\mathbf{K}}$ be as in Proposition~16.2\,(v).
The four rules displayed there, the rules for $a_{\tilde{g}}$, $G_{\Lambda}$, $\mathrm{Cg}$ and $\mathbf{K}$, are denoted as follows:
\begin{align}
a_{\tilde{g}}\mathrm{Bil}[K]&=\mathrm{Bil}[K^{\tilde{g}}],\tag{Ru-g}\\
G_{\Lambda}\mathrm{Bil}[K]&=\mathrm{Bil}[K^{\Lambda}],\tag{Ru-G}\\
\mathrm{Cg}\,\mathrm{Bil}[K]&=\mathrm{Bil}[K^{\mathrm{Cg}}],\tag{Ru-C}\\
\mathbf{K}\,\mathrm{Bil}[K]&=\mathrm{Bil}[K^{\mathbf{K}}].\tag{Ru-K}
\end{align}
These rules hold, and $K^{\tilde{g}}$, $K^{\Lambda}$, $K^{\mathrm{Cg}}$, $K^{\mathbf{K}}$ are matrix kernels.

\emph{The rule} (Ru-g).
The map $a_{\tilde{g}}$ is a homomorphism, and $a_{\tilde{g}}\bar{\psi}(\tilde{x})=\left(S^{-1}\right)^T\bar{\psi}(\tilde{g}^{-1}\tilde{x})$ (Section~I.4.5), so that $\left(a_{\tilde{g}}\bar{\psi}(\tilde{x})\right)^T=\bar{\psi}(\tilde{g}^{-1}\tilde{x})^TS^{-1}$; moreover $a_{\tilde{g}}\psi(\tilde{x}')=S\psi(\tilde{g}^{-1}\tilde{x}')$.
Hence
\[
a_{\tilde{g}}\mathrm{Bil}[K]=\int\mathrm{d}\tilde{x}\int\mathrm{d}\tilde{x}'\,f(\tilde{x},\tilde{x}'),\qquad f(\tilde{x},\tilde{x}'):=\bar{\psi}(\tilde{g}^{-1}\tilde{x})^TS^{-1}K(\tilde{x},\tilde{x}')S\psi(\tilde{g}^{-1}\tilde{x}').
\]
The function $f$ is $\ell\mathbb{Z}^3$-periodic in $\tilde{x}$ and in $\tilde{x}'$ separately.
Indeed, $\tilde{g}^{-1}(\tilde{x}+\ell n)=\tilde{g}^{-1}\tilde{x}+\ell rn$.
Under $\tilde{x}\mapsto\tilde{x}+\ell n$ the kernel is multiplied by $\boldsymbol{\chi}(n)$, and
\[
\bar{\psi}(\tilde{g}^{-1}\tilde{x}+\ell rn)=\boldsymbol{\chi}(rn)\cdot\bar{\psi}(\tilde{g}^{-1}\tilde{x})=\boldsymbol{\chi}(n)\cdot\bar{\psi}(\tilde{g}^{-1}\tilde{x});
\]
under $\tilde{x}'\mapsto\tilde{x}'+\ell n$ the kernel is multiplied by $\boldsymbol{\chi}(n)$, and $\psi(\tilde{g}^{-1}\tilde{x}'+\ell rn)=\boldsymbol{\chi}(n)\cdot\psi(\tilde{g}^{-1}\tilde{x}')$.
By the invariance of the measure, Proposition~16.1\,(a2), applied first in $\tilde{x}$ and then in $\tilde{x}'$,
\[
\int\mathrm{d}\tilde{x}\int\mathrm{d}\tilde{x}'\,f(\tilde{x},\tilde{x}')=\int\mathrm{d}x\int\mathrm{d}x'\,f(\tilde{g}x,\tilde{g}x')=\int\mathrm{d}x\int\mathrm{d}x'\,\bar{\psi}(x)^TS^{-1}K(\tilde{g}x,\tilde{g}x')S\psi(x'),
\]
which is $\mathrm{Bil}[K^{\tilde{g}}]$.
Moreover
\begin{align*}
K^{\tilde{g}}(x+\ell n,x')&=S^{-1}K(\tilde{g}x+\ell rn,\tilde{g}x')S=\boldsymbol{\chi}(rn)\cdot K^{\tilde{g}}(x,x')=\boldsymbol{\chi}(n)\cdot K^{\tilde{g}}(x,x'),\\
K^{\tilde{g}}(x,x'+\ell n)&=S^{-1}K(\tilde{g}x,\tilde{g}x'+\ell rn)S=\boldsymbol{\chi}(n)\cdot K^{\tilde{g}}(x,x').
\end{align*}

\emph{The rule} (Ru-G).
One has $G_{\Lambda}\bar{\psi}(x)=\exp\left(-ie_k\cdot\Lambda(c)\right)\bar{\psi}(x)$ and $G_{\Lambda}\psi(x')=\exp\left(ie_k\cdot\Lambda(c')\right)\psi(x')$, which gives (Ru-G).
Since $\Lambda$ is $\ell\mathbb{Z}^3$-periodic, the factor $\exp\left\{ie_k\cdot\left(\Lambda(c')-\Lambda(c)\right)\right\}$ is periodic in each argument, and $K^{\Lambda}$ is a matrix kernel.

\emph{The rule} (Ru-C).
One has $\mathrm{Cg}\,\bar{\psi}(x)=-\mathsf{C}\psi(x)$ and $\mathrm{Cg}\,\psi(x')=\mathsf{C}\bar{\psi}(x')$.
Hence
\begin{align*}
\mathrm{Cg}\,\mathrm{Bil}[K]&=-\int\mathrm{d}x\int\mathrm{d}x'\,\psi(x)^T\mathsf{C}^TK(x,x')\mathsf{C}\bar{\psi}(x')\\
&=-\sum_{\beta,\beta'}\int\mathrm{d}x\int\mathrm{d}x'\,\left(\mathsf{C}^TK(x,x')\mathsf{C}\right)_{\beta\beta'}\psi_{\beta}(x)\bar{\psi}_{\beta'}(x')\\
&=\sum_{\beta,\beta'}\int\mathrm{d}x\int\mathrm{d}x'\,\bar{\psi}_{\beta'}(x')\left(\mathsf{C}^TK(x,x')\mathsf{C}\right)_{\beta\beta'}\psi_{\beta}(x)\\
&=\int\mathrm{d}x\int\mathrm{d}x'\,\bar{\psi}(x')^T\mathsf{C}^TK(x,x')^T\mathsf{C}\psi(x),
\end{align*}
because $\psi_{\beta}(x)\bar{\psi}_{\beta'}(x')=-\bar{\psi}_{\beta'}(x')\psi_{\beta}(x)$ and $\left(\mathsf{C}^TK\mathsf{C}\right)^T=\mathsf{C}^TK^T\mathsf{C}$.
Exchanging the names of $x$ and $x'$ and using $\mathsf{C}^T=\mathsf{C}^{-1}$ (I.c3) one obtains (Ru-C).
Moreover
\begin{align*}
K^{\mathrm{Cg}}(x+\ell n,x')&=\mathsf{C}^{-1}K(x',x+\ell n)^T\mathsf{C}=\boldsymbol{\chi}(n)\cdot K^{\mathrm{Cg}}(x,x'),\\
K^{\mathrm{Cg}}(x,x'+\ell n)&=\mathsf{C}^{-1}K(x'+\ell n,x)^T\mathsf{C}=\boldsymbol{\chi}(n)\cdot K^{\mathrm{Cg}}(x,x').
\end{align*}

\emph{The rule} (Ru-K).
The map $\mathbf{K}$ is antilinear and multiplicative, the weights $\left(\eta^3/48\right)^2$ are real, and $\mathbf{K}\bar{\psi}(x)=\varkappa\bar{\psi}(x)$, $\mathbf{K}\psi(x')=\varkappa\psi(x')$.
Hence
\[
\mathbf{K}\,\mathrm{Bil}[K]=\int\mathrm{d}x\int\mathrm{d}x'\,\bar{\psi}(x)^T\varkappa^T\,\overline{K(x,x')}\,\varkappa\psi(x'),
\]
and $\varkappa^T=\varkappa^{-1}$ (I.c4); this is (Ru-K).
Since $\boldsymbol{\chi}(n)$ is real, $K^{\mathbf{K}}$ is a matrix kernel.

\textbf{(o)}
The number $\mathrm{al}_0(\mathscr{E},X)=\operatorname{Vol}(X)^{-1}\cdot\mathscr{E}_{00}(X;0,0)$ involves no argument.
The polymer $X$ has a lift $\tilde{X}$, and every lift of $X$ has the form $\tilde{X}+\ell n$ (Lemma~0.2\,(b), (c)).
Since
\[
K_X(x+\ell n,x'+\ell n)=\boldsymbol{\chi}(n)^2\cdot K_X(x,x')=K_X(x,x'),
\]
one has
\[
\int_{\tilde{X}+\ell n}\mathrm{d}x\int_{\tilde{X}+\ell n}\mathrm{d}x'\,K_X(x,x')=\int_{\tilde{X}}\mathrm{d}x\int_{\tilde{X}}\mathrm{d}x'\,K_X(x+\ell n,x'+\ell n)=\int_{\tilde{X}}\mathrm{d}x\int_{\tilde{X}}\mathrm{d}x'\,K_X(x,x').
\]
So $\mathrm{al}_2(\mathscr{E},X)$ does not depend on the lift either.

\textbf{(i)}
Since $\tilde{g}0=0$, the clause (k3) of the definition in Section~I.5.4 reads at $(\mathscr{A},\mathscr{C})=(0,0)$
\[
a_{\tilde{g}}\left[\mathscr{E}(\tilde{g}X;0,0)\right]=\mathscr{E}(X;0,0).
\]
(The element $\mathscr{E}(\tilde{g}X;0,0)$ is even; by (b.6) the left side does not depend on the vector $a$ chosen in its class.)
The map $a_{\tilde{g}}$ preserves the bidegree, so that the parts of bidegree $(0,0)$ and $(2,0)$ satisfy the relation separately.
On numbers $a_{\tilde{g}}$ is the identity: $\mathscr{E}_{00}(\tilde{g}X;0,0)=\mathscr{E}_{00}(X;0,0)$.
The polymer $\tilde{g}X$ is small, and $\operatorname{Vol}(\tilde{g}X)=M^3\cdot|\tilde{g}X|_M=\operatorname{Vol}(X)$ (Proposition~16.1\,(a1)).
Hence $\mathrm{al}_0(\mathscr{E},\tilde{g}X)=\mathrm{al}_0(\mathscr{E},X)$.
In bidegree $(2,0)$ the relation reads $a_{\tilde{g}}\mathrm{Bil}[K_{\tilde{g}X}]=\mathrm{Bil}[K_X]$, and (Ru-g) together with Lemma~0.3\,(a) gives
\begin{equation}
S^{-1}K_{\tilde{g}X}(\tilde{g}x,\tilde{g}x')S=K_X(x,x')\qquad\text{for all }x,x'\in\mathbb{L}^{\mathrm{f}}_k.\tag{eq.6}
\end{equation}
The set $\tilde{g}\tilde{X}$ is a lift of $\tilde{g}X$ (Proposition~16.1\,(a1)); by (o) it may be used in the definition of $\mathrm{al}_2(\mathscr{E},\tilde{g}X)$.
By (eq.2), applied in both variables, and by (eq.6),
\[
\int_{\tilde{g}\tilde{X}}\mathrm{d}\tilde{x}\int_{\tilde{g}\tilde{X}}\mathrm{d}\tilde{x}'\,K_{\tilde{g}X}(\tilde{x},\tilde{x}')=\int_{\tilde{X}}\mathrm{d}x\int_{\tilde{X}}\mathrm{d}x'\,K_{\tilde{g}X}(\tilde{g}x,\tilde{g}x')=S\left[\int_{\tilde{X}}\mathrm{d}x\int_{\tilde{X}}\mathrm{d}x'\,K_X(x,x')\right]S^{-1}.
\]
(In the sums (eq.5) the substitution is $(c,g)\mapsto(rc+a,rg)$, $(c',g')\mapsto(rc'+a,rg')$; it permutes the $48$ labels of each of the two clusters.)
Division by $\operatorname{Vol}(\tilde{g}X)=\operatorname{Vol}(X)$ gives
\[
\mathrm{al}_2(\mathscr{E},\tilde{g}X)=S\,\mathrm{al}_2(\mathscr{E},X)\,S^{-1}.
\]

\textbf{(ii)}
The clause (k5) at $(0,0)$ reads $\mathrm{Cg}\left[\mathscr{E}(X;0,0)\right]=\mathscr{E}(X;0,0)$.
The map $\mathrm{Cg}$ preserves the bidegree.
In bidegree $(2,0)$, (Ru-C) and Lemma~0.3\,(a) give
\[
K_X(x,x')=\mathsf{C}^{-1}K_X(x',x)^T\mathsf{C}.
\]
Integrate in $x$ and in $x'$ over $\tilde{X}$; after an exchange of the names of the two variables,
\[
\int_{\tilde{X}}\mathrm{d}x\int_{\tilde{X}}\mathrm{d}x'\,K_X(x',x)^T=\left(\int_{\tilde{X}}\mathrm{d}x\int_{\tilde{X}}\mathrm{d}x'\,K_X(x,x')\right)^T.
\]
Division by $\operatorname{Vol}(X)$ gives $\mathrm{al}_2(\mathscr{E},X)=\mathsf{C}^{-1}\mathrm{al}_2(\mathscr{E},X)^T\mathsf{C}$.

\textbf{(iii)}
The clause (k7) at $(0,0)$ reads $\mathbf{K}\left[\mathscr{E}(X;0,0)\right]=\mathscr{E}(X;0,0)$.
The map $\mathbf{K}$ preserves the bidegree and acts on numbers by complex conjugation.
Hence $\overline{\mathscr{E}_{00}(X;0,0)}=\mathscr{E}_{00}(X;0,0)$, and, $\operatorname{Vol}(X)$ being real, $\mathrm{al}_0(\mathscr{E},X)$ is real.
In bidegree $(2,0)$, (Ru-K) and Lemma~0.3\,(a) give
\[
K_X(x,x')=\varkappa^{-1}\,\overline{K_X(x,x')}\,\varkappa.
\]
Integration in $x$ and in $x'$ over $\tilde{X}$, with real weights, and division by $\operatorname{Vol}(X)$ give $\mathrm{al}_2(\mathscr{E},X)=\varkappa^{-1}\,\overline{\mathrm{al}_2(\mathscr{E},X)}\,\varkappa$.

\textbf{(iv)}
Let $Z$ be a union of $M$-cubes of $T_k$ and $\mathrm{Mx}$ a $4\times4$ matrix.
The function $\bar{\psi}(x)^T\mathrm{Mx}\,\psi(x)$ is $\ell\mathbb{Z}^3$-periodic, and
\[
\int_Z\mathrm{d}x\,\bar{\psi}(x)^T\mathrm{Mx}\,\psi(x)=\frac{\eta^3}{48}\sum_{c\in Z}\ \sum_{g\in O_{\mathrm{cub}}}\ \sum_{\beta,\beta'}\mathrm{Mx}_{\beta\beta'}\,\bar{\psi}_{\beta}(c,g)\,\psi_{\beta'}(c,g).
\]
Put
\[
\mathrm{Kd}_Z[\mathrm{Mx}](x,x'):=48\,\eta^{-3}\cdot\mathbf{1}[\mathrm{pr}\,c\in Z]\cdot\left(\sum_{n\in\mathbb{Z}^3}\boldsymbol{\chi}(n)\,\mathbf{1}[x'=x+\ell n]\right)\mathrm{Mx}.
\]
This is a matrix kernel: the replacement of $x$ by $x+\ell n'$, or of $x'$ by $x'+\ell n'$, multiplies the sum by $\boldsymbol{\chi}(n')$ (substitute $n\mapsto n\mp n'$, and note $\boldsymbol{\chi}(-n')=\boldsymbol{\chi}(n')$).
If $c$ and $c'$ lie in one set on which $\mathrm{pr}$ is injective (a fundamental domain, or a lift), only the term $n=0$ contributes, and there
\[
\mathrm{Kd}_Z[\mathrm{Mx}](x,x')=48\,\eta^{-3}\cdot\mathbf{1}[x=x']\cdot\mathbf{1}[\mathrm{pr}\,c\in Z]\cdot\mathrm{Mx};
\]
this is the kernel $48\,\eta_k^{-3}\cdot\mathbf{1}[x=x'\text{ over }Z]\cdot\mathrm{Mx}$ of Proposition~16.2\,(iv).
For $x$ and $x'$ over $\mathrm{Fdm}$,
\[
\int\mathrm{d}x'\,\mathrm{Kd}_Z[\mathrm{Mx}](x,x')\psi(x')=\frac{\eta^3}{48}\cdot48\,\eta^{-3}\cdot\mathbf{1}[\mathrm{pr}\,c\in Z]\cdot\mathrm{Mx}\,\psi(x)=\mathbf{1}_Z(x)\,\mathrm{Mx}\,\psi(x),
\]
which is (I.ke1).
So $\mathrm{Bil}\left[\mathrm{Kd}_Z[\mathrm{Mx}]\right]=\int_Z\mathrm{d}x\,\bar{\psi}(x)^T\mathrm{Mx}\,\psi(x)$, and by Lemma~0.3\,(a) the function $\mathrm{Kd}_Z[\mathrm{Mx}]$ is the kernel of this element.

\emph{Integral.}
With both arguments over one fundamental domain (over $\mathsf{Z}$ if $Z$ has a lift $\mathsf{Z}$),
\[
\int\mathrm{d}x\int\mathrm{d}x'\,\mathrm{Kd}_Z[\mathrm{Mx}](x,x')=\int_Z\mathrm{d}x\,\mathrm{Mx}=\operatorname{Vol}(Z)\cdot\mathrm{Mx},
\]
because
\[
\int_Z\mathrm{d}x=\eta^3\cdot\frac{1}{48}\cdot48\cdot\left(\frac{M}{\eta}\right)^3\cdot|Z|_M=M^3\cdot|Z|_M=\operatorname{Vol}(Z);
\]
here an $M$-cube has $(M/\eta)^3$ photon sites by (Ar2), each of them with $48$ fermion sites.

\emph{Norm.}
By Lemma~0.3\,(c),
\[
\left\|\int_Z\mathrm{d}x\,\bar{\psi}(x)^T\mathrm{Mx}\,\psi(x)\right\|_{\mathsf{w}}=\mathsf{w}_1^2\sum_{\beta,\beta'}\int\mathrm{d}x\int\mathrm{d}x'\,\left|\mathrm{Kd}_Z[\mathrm{Mx}](x,x')_{\beta\beta'}\right|=\mathsf{w}_1^2\cdot\operatorname{Vol}(Z)\sum_{\beta,\beta'}\left|\mathrm{Mx}_{\beta\beta'}\right|.
\]

\emph{The four relations.}
For $a_{\tilde{g}}$: with the $\ell\mathbb{Z}^3$-periodic function $f(x):=\bar{\psi}(x)^T\mathrm{Mx}\,\psi(x)$,
\[
\left(a_{\tilde{g}}\bar{\psi}(\tilde{x})\right)^TS\,\mathrm{Mx}\,S^{-1}\left(a_{\tilde{g}}\psi(\tilde{x})\right)=\bar{\psi}(\tilde{g}^{-1}\tilde{x})^TS^{-1}S\,\mathrm{Mx}\,S^{-1}S\psi(\tilde{g}^{-1}\tilde{x})=f(\tilde{g}^{-1}\tilde{x}),
\]
and by Proposition~16.1\,(a2)
\[
a_{\tilde{g}}\left[\int_{\tilde{g}Z}\mathrm{d}\tilde{x}\,\bar{\psi}(\tilde{x})^TS\,\mathrm{Mx}\,S^{-1}\psi(\tilde{x})\right]=\int_{\tilde{g}Z}\mathrm{d}\tilde{x}\,f(\tilde{g}^{-1}\tilde{x})=\int_Z\mathrm{d}x\,f(x)=\int_Z\mathrm{d}x\,\bar{\psi}(x)^T\mathrm{Mx}\,\psi(x).
\]
For $G_{\Lambda}$: at every site $x=(c,g)$, for each of the $48$ labels, the two phases cancel,
\[
G_{\Lambda}\left(\bar{\psi}_{\beta}(x)\psi_{\beta'}(x)\right)=\exp\left(-ie_k\cdot\Lambda(c)\right)\exp\left(ie_k\cdot\Lambda(c)\right)\bar{\psi}_{\beta}(x)\psi_{\beta'}(x)=\bar{\psi}_{\beta}(x)\psi_{\beta'}(x),
\]
so that $G_{\Lambda}$ fixes $\int_Z\mathrm{d}x\,\bar{\psi}(x)^T\mathrm{Mx}\,\psi(x)$.
For $\mathrm{Cg}$: apply (Ru-C) to $K=\mathrm{Kd}_Z[\mathrm{Mx}]$.
The scalar factor of $\mathrm{Kd}_Z[\mathrm{Mx}](x',x)$ is
\[
48\,\eta^{-3}\cdot\mathbf{1}[\mathrm{pr}\,c'\in Z]\cdot\sum_{n}\boldsymbol{\chi}(n)\,\mathbf{1}[x=x'+\ell n],
\]
and it equals the scalar factor of $\mathrm{Kd}_Z[\mathrm{Mx}](x,x')$ (substitute $n\mapsto-n$, and note that $\mathrm{pr}\,c=\mathrm{pr}\,c'$ wherever a term does not vanish).
Hence $K^{\mathrm{Cg}}=\mathrm{Kd}_Z\left[\mathsf{C}^{-1}\mathrm{Mx}^T\mathsf{C}\right]$, and
\[
\mathrm{Cg}\left[\int_Z\mathrm{d}x\,\bar{\psi}(x)^T\mathrm{Mx}\,\psi(x)\right]=\int_Z\mathrm{d}x\,\bar{\psi}(x)^T\mathsf{C}^{-1}\mathrm{Mx}^T\mathsf{C}\,\psi(x).
\]
For $\mathbf{K}$: apply (Ru-K); the scalar factor is real, so that $K^{\mathbf{K}}=\mathrm{Kd}_Z\left[\varkappa^{-1}\overline{\mathrm{Mx}}\,\varkappa\right]$, and
\[
\mathbf{K}\left[\int_Z\mathrm{d}x\,\bar{\psi}(x)^T\mathrm{Mx}\,\psi(x)\right]=\int_Z\mathrm{d}x\,\bar{\psi}(x)^T\varkappa^{-1}\overline{\mathrm{Mx}}\,\varkappa\,\psi(x).
\]

\textbf{(v)}
The functions $K$ of Proposition~16.2\,(v) are the matrix kernels of Section~0.4.
That $\mathrm{Bil}[K]$ is balanced and of bidegree $(2,0)$, that its kernel in the matrix notation of Section~I.5.1 is $K$, and that $\mathrm{Bil}[K]=\mathrm{Bil}[K']$ implies $K=K'$, is Lemma~0.3\,(a).
That every balanced element of bidegree $(2,0)$ has the form $\mathrm{Bil}[K]$ is Lemma~0.3\,(b), and the norm of $\mathrm{Bil}[K]$ is given by Lemma~0.3\,(c).
The four rules, for $\tilde{g}=(S,a)\in\mathbb{G}_k$ and a real function $\Lambda$ of $c$ (a function on $T_k$), and the antiperiodicity of $K^{\tilde{g}}$, $K^{\Lambda}$, $K^{\mathrm{Cg}}$, $K^{\mathbf{K}}$ in each argument, are the content of Step~1.

\textbf{Level~$0$.}
At $k=0$ one has $\eta=1$ and $\int\mathrm{d}x=\frac{1}{48}\sum_c\sum_g$; further
\[
\mathrm{al}_2(\mathscr{E},X)=\operatorname{Vol}(X)^{-1}\cdot\left(\frac{1}{48}\right)^2\sum_{c,c'\in\tilde{X}}\ \sum_{g,g'}K_X\left((c,g),(c',g')\right),
\]
the kernel of (iv) is $48\cdot\mathbf{1}[x=x'\text{ over }Z]\cdot\mathrm{Mx}$, and the number $e_k$ is $q$.
With these values every argument above applies at $k=0$ as well, and none of them involves a constant.
\end{proof}

\begin{proof}[Verification of the Remark after Proposition~16.2]
There are a family $E\in K^{x}_{k}$ and a small polymer $X$ such that $\mathrm{al}_2(E,X)=i\,\Sigma_3$, a matrix that is not a multiple of $1$.
To construct them, let $z\in M\mathbb{Z}^3$ and put
\[
z_0:=z,\qquad z_1:=z+M\cdot e_1,\qquad z_2:=z_1+M\cdot e_2,\qquad z_3:=z_2+M\cdot e_3;
\]
let $\square_i$ be the $M$-cube with the centre $z_i$, and let $X$ be the union of the four cubes $\square_0,\dots,\square_3$.
One has $z_{i+1}-z_i=M\cdot e_{i+1}$, whereas the differences $z_2-z_0$, $z_3-z_1$, $z_3-z_0$ have two or three entries $M$ and are not congruent to a vector $\pm M\cdot e_{\mu}$ modulo $\ell\mathbb{Z}^3$, because $2M<\ell$.
(The four centres are distinct modulo $\ell\mathbb{Z}^3$, since their differences are non-zero vectors of sup norm $M<\ell$.)
Hence the pairs of cubes of $X$ with a common face are exactly the pairs $(\square_i,\square_{i+1})$, and $X$ is a polymer.
By \cite[Theorem~13.2\,(t0)]{QEDI}, and since $L\ge5$ by Proposition~14.2\,(L5),
\[
d_M(X)\le|X|_M-1=3\le L;
\]
so $X$ is small.

(1)
If $\tilde{g}=(S,a)\in\mathbb{G}_k$ and $\tilde{g}X=X$, then $r(S)=1$.
Indeed, by (a1.7) the element $\tilde{g}$ permutes the four cubes of $X$ and preserves common faces; hence either $\tilde{g}\square_i=\square_i$ for all $i$, or $\tilde{g}\square_i=\square_{3-i}$ for all $i$.
Moreover $\tilde{g}z_{i+1}-\tilde{g}z_i=r_{i+1}M\cdot e_{i+1}$.
In the first case $r_{i+1}M\cdot e_{i+1}\equiv M\cdot e_{i+1}$ modulo $\ell\mathbb{Z}^3$, and therefore $r_{i+1}=1$ for $i=0,1,2$, because $2M<\ell$.
In the second case $r_1M\cdot e_1\equiv z_2-z_3=-M\cdot e_3$ modulo $\ell\mathbb{Z}^3$, which is impossible.

(2)
For $\tilde{g}=(S,a)\in\mathbb{G}_k$ and for all $(\mathscr{A},\mathscr{C})$ put
\[
E(\tilde{g}X):=i\det\left(r(S)\right)\cdot\int_{\tilde{g}X}\mathrm{d}x\,\bar{\psi}(x)^T\,\Sigma_3\,\psi(x),
\]
and $E(X'):=0$ for the polymers $X'$ outside the orbit of $X$.
This is well defined: if $\tilde{g}X=\tilde{g}_1X$ with $\tilde{g}_1=(S_1,a_1)$, then $r(S^{-1}S_1)=1$ by (1), so that $r(S)=r(S_1)$.

(3)
From $\gamma_{\nu}S=r_{\nu}(S)S\gamma_{\nu}$ one obtains $S^{-1}\gamma_{\nu}S=r_{\nu}(S)\gamma_{\nu}$; together with $i\,\Sigma_3=\gamma_1\gamma_2\gamma_3$ (I.c2) this gives
\[
S^{-1}\,\Sigma_3\,S=\det\left(r(S)\right)\cdot\Sigma_3.
\]
Further $\Sigma_3^T=\Sigma_3$, and $\mathsf{C}=i\gamma_2$ commutes with $\Sigma_3$ (I.c2); hence $\mathsf{C}^{-1}\Sigma_3^T\mathsf{C}=\Sigma_3$.
Finally $\varkappa=\gamma_5\mathsf{C}$, and $\gamma_5$ anticommutes with $\Sigma_3$; hence $\varkappa^{-1}\Sigma_3\,\varkappa=-\Sigma_3$, and therefore
\[
\varkappa^{-1}\,\overline{i\,\Sigma_3}\,\varkappa=i\,\Sigma_3.
\]

(4)
(k1), (k2), (k6): the elements do not depend on $(\mathscr{A},\mathscr{C})$, they contain $\psi$ over the polymer only, and there are finitely many polymers.
(k3): for $\tilde{g}_1=(S_1,a_1)\in\mathbb{G}_k$, the relation for $a_{\tilde{g}_1}$ in (iv), with $\mathrm{Mx}=S_1^{-1}\,\Sigma_3\,S_1=\det\left(r(S_1)\right)\cdot\Sigma_3$, gives
\[
a_{\tilde{g}_1}\left[E(\tilde{g}_1\tilde{g}X)\right]=i\det\left(r(S_1S)\right)\det\left(r(S_1)\right)\cdot\int_{\tilde{g}X}\mathrm{d}x\,\bar{\psi}(x)^T\,\Sigma_3\,\psi(x)=E(\tilde{g}X),
\]
and for the polymers outside the orbit of $X$ both sides of (k3) vanish.
(k4), (k5), (k7): these follow from (iv) and (3).
The elements are even and balanced.
So $E\in K^{x}_{k}$, and $\mathrm{al}_2(E,X)=i\,\Sigma_3$ by (iv).
\end{proof}

\subsection*{3. Proposition 16.3}

\begin{proof}[Proof of Proposition~16.3]
\textbf{Preliminaries.}
(pre1) The pair $(0,0)$ belongs to $\tilde{R}_k\times\mathfrak{X}_k$ (Section~I.5.3).
Indeed, $0=A_0+A_1$ with $A_1=0$ and $A_0=0=A'+d\boldsymbol{\lambda}$, where $A'=0$ and $\boldsymbol{\lambda}=0$;
all the norms that occur vanish, while the thresholds are positive numbers.
Consequently $\mathscr{E}(X;0,0)$ is defined for every $X\in D_k$, and the clause (k6) of Section~I.5.4 yields
\begin{equation}
\left\|\mathscr{E}(X;0,0)\right\|_{h_k}\le\|\mathscr{E}\|^x_k\cdot\exp\left(-\kappa\cdot d_M(X)\right).
\tag{eq.7}
\end{equation}

(pre2) The family $\mathscr{E}$ is balanced, and by the clause (k2) of Section~I.5.4 its kernels $\mathscr{E}_{20}(X;0,0;x,x')$ vanish unless $x$ and $x'$ lie over $X$.
Thus $\mathscr{E}$ is a family as in Proposition~16.2.
Moreover the clauses (k3) and (k5) hold at $(0,0)$, and (k7) holds there whenever $\mathscr{E}\in K^x_k$.

(pre3) The set $D_k$ is finite.
For an $M$-cube $\square$ let $\mathrm{Pl}(\square)$ denote the set of the small polymers that contain $\square$.
As in the proof of Proposition~16.2, put $K_X:=\mathscr{E}_{20}(X;0,0;\,\cdot\,,\,\cdot\,)$.

(pre4) By (I.Sch), a $4\times4$ matrix $\mathrm{Mx}$ that commutes with $\gamma_1$, $\gamma_2$, $\gamma_3$, $\gamma_5$ is a multiple of the identity matrix $\mathbf{1}$.
This is seen as follows.
Such a matrix commutes with $\gamma_3=\operatorname{diag}(1,-1,-1,1)$ and, by (I.c2), with $\Sigma_3=-i\cdot\gamma_1\cdot\gamma_2\cdot\gamma_3=\operatorname{diag}(1,1,-1,-1)$.
The four pairs of corresponding diagonal entries, $(1,1)$, $(-1,1)$, $(-1,-1)$, $(1,-1)$, are pairwise distinct;
hence $\mathrm{Mx}=\operatorname{diag}(d_1,\dots,d_4)$.
The non-zero entries of $\gamma_1$ stand at the places $(1,2)$, $(2,1)$, $(3,4)$, $(4,3)$, which gives $d_1=d_2$ and $d_3=d_4$;
those of $\gamma_5$ stand at the places $(1,3)$, $(2,4)$, $(3,1)$, $(4,2)$, which gives $d_1=d_3$.

\paragraph{(a)}
\textbf{(a.1) Independence of the cube.}
Let $\square$, $\square'$ be $M$-cubes with the centres $z$, $z'$, and let $a\in M\mathbb{Z}^3$ be a vector in the class $z'-z$.
The element $\tilde{g}:=(\mathbf{1},a)$ lies in $\mathbb{G}_k$;
its matrix $r$ is the identity, $\tilde{g}c=c+a$ and $\tilde{g}(c,g)=(c+a,g)$, so that no label is moved, and $\tilde{g}\square=\square'$.
By Proposition~16.1\,(a1) the map $X\mapsto\tilde{g}X=X+a$ sends small polymers to small polymers, and $X$ contains $\square$ if and only if $X+a$ contains $\square'$;
its inverse is $X\mapsto X-a$.
Hence $X\mapsto X+a$ is a bijection of $\mathrm{Pl}(\square)$ onto $\mathrm{Pl}(\square')$.
Proposition~16.2\,(i) with $S=\mathbf{1}$ gives
\[
\mathrm{al}_{0}(\mathscr{E},X+a)=\mathrm{al}_{0}(\mathscr{E},X),\qquad\mathrm{al}_{2}(\mathscr{E},X+a)=\mathrm{al}_{2}(\mathscr{E},X),
\]
and summation over $\mathrm{Pl}(\square)$ yields $\varepsilon(\mathscr{E},\square')=\varepsilon(\mathscr{E},\square)$ and $m(\mathscr{E},\square')=m(\mathscr{E},\square)$.

\textbf{(a.2) A multiple of the identity.}
Let $z\in M\mathbb{Z}^3$ be the centre of $\square$, let $S$ be one of the matrices $\gamma_1$, $\gamma_2$, $\gamma_3$, $\gamma_5$, put $r:=r(S)$, and let $\tilde{g}_z=(S,(1-r)z)$ be the element of Section~I.4.2.
The diagonal entries of $1-r$ are $0$ or $2$;
so $(1-r)z$ lies in $2M\mathbb{Z}^3$, which is contained in $M\mathbb{Z}^3$, and $\tilde{g}_z\in\mathbb{G}_k$.
One has $\tilde{g}_zc=rc+(1-r)z=z+r(c-z)$, in particular $\tilde{g}_zz=z$;
therefore $\tilde{g}_z\square=\square$ by Proposition~16.1\,(a1).
On the fermion sites $\tilde{g}_z(c,g)=\left(z+r(c-z),rg\right)$:
the labels are multiplied by $r$, which is the rotation by $\pi$ about $e_\mu$ if $S=\gamma_\mu$, while $r=-1$ if $S=\gamma_5$.
The map $X\mapsto\tilde{g}_zX$ is a bijection of $\mathrm{Pl}(\square)$ onto itself;
its inverse is $X\mapsto\tilde{g}_z^{-1}X$, and $\tilde{g}_z^{-1}$ fixes $\square$ as well.
Proposition~16.2\,(i) now gives
\begin{align*}
-m(\mathscr{E},\square)&=\sum_{X\in\mathrm{Pl}(\square)}\mathrm{al}_{2}(\mathscr{E},\tilde{g}_zX)\\
&=S\cdot\left[\sum_{X\in\mathrm{Pl}(\square)}\mathrm{al}_{2}(\mathscr{E},X)\right]\cdot S^{-1}=-S\cdot m(\mathscr{E},\square)\cdot S^{-1}.
\end{align*}
Thus $m(\mathscr{E},\square)$ commutes with $\gamma_1$, $\gamma_2$, $\gamma_3$, $\gamma_5$, and (I.Sch) gives $m(\mathscr{E},\square)=m(\mathscr{E})\cdot\mathbf{1}$ with a complex number $m(\mathscr{E})$;
by (a.1) this number does not depend on $\square$.
Put $\varepsilon(\mathscr{E}):=\varepsilon(\mathscr{E},\square)$.

\textbf{(a.3) Linearity; restriction.}
The antisymmetric kernels of an element depend linearly on the element:
by Proposition~15\,(a) they are its coefficients divided by the weights.
Hence $\mathscr{E}\mapsto\mathscr{E}_{00}(X;0,0)$ and $\mathscr{E}\mapsto K_X$ are complex linear, and so are $\mathrm{al}_{0}(\,\cdot\,,X)$, $\mathrm{al}_{2}(\,\cdot\,,X)$, $\varepsilon(\,\cdot\,)$ and $m(\,\cdot\,)$ on $K^{x,C}_{k}$.
For the family without slot $\mathscr{E}|_0(X;\mathscr{A})=\mathscr{E}(X;\mathscr{A},0)$ the numbers $\mathrm{al}_{0}$, $\mathrm{al}_{2}$, $\varepsilon$, $m$ are formed with $\mathscr{E}|_0(X;0)$ (Section~I.5.11), and $\mathscr{E}|_0(X;0)=\mathscr{E}(X;0,0)$.
Therefore $\mathrm{al}_{0}(\mathscr{E}|_0,X)=\mathrm{al}_{0}(\mathscr{E},X)$ and $\mathrm{al}_{2}(\mathscr{E}|_0,X)=\mathrm{al}_{2}(\mathscr{E},X)$ for every small $X$;
consequently $\varepsilon(\mathscr{E}|_0,\square)=\varepsilon(\mathscr{E},\square)$ and $m(\mathscr{E}|_0,\square)=m(\mathscr{E},\square)$ for every $M$-cube $\square$.
By (a.1) and (a.2) the right-hand sides equal $\varepsilon(\mathscr{E})$ and $m(\mathscr{E})\cdot\mathbf{1}$, whatever the cube $\square$.
It follows that $\varepsilon(\mathscr{E}|_0)=\varepsilon(\mathscr{E})$ and $m(\mathscr{E}|_0)=m(\mathscr{E})$.

\textbf{(a.4) Reality.}
Let $\mathscr{E}\in K^{x}_{k}$.
Each number $\mathrm{al}_{0}(\mathscr{E},X)$ is real by Proposition~16.2\,(iii), and so is $\varepsilon(\mathscr{E})$.
Summing the relation $\mathrm{al}_{2}(\mathscr{E},X)=\varkappa^{-1}\cdot\overline{\mathrm{al}_{2}(\mathscr{E},X)}\cdot\varkappa$ of Proposition~16.2\,(iii) over $\mathrm{Pl}(\square)$, one finds
\[
m(\mathscr{E})\cdot\mathbf{1}=\varkappa^{-1}\cdot\overline{m(\mathscr{E})}\cdot\mathbf{1}\cdot\varkappa=\overline{m(\mathscr{E})}\cdot\mathbf{1},
\]
so that $m(\mathscr{E})$ is real.

\paragraph{(b)}
\textbf{The inequalities (2.2).}
Let $X$ be small and let $\tilde{X}$ be a lift of $X$.
By (G2.r7), that is, by Proposition~15\,(n),
\[
\operatorname{Vol}(X)\cdot\left|\mathrm{al}_{0}(\mathscr{E},X)\right|=\left|\mathscr{E}_{00}(X;0,0)\right|=\left\|\mathscr{E}^{[0,0]}(X;0,0)\right\|_{h_k}\le\left\|\mathscr{E}(X;0,0)\right\|_{h_k}.
\]
The function $\left|K_X(x,x')_{\beta\beta'}\right|$ is $\ell\mathbb{Z}^3$-periodic in each argument and vanishes unless $x$ and $x'$ lie over $X$;
its integral over $\tilde{X}\times\tilde{X}$ is therefore its integral over the torus.
With (eq.5), Lemma~0.3\,(c) for $\mathsf{w}=h_k$, and $h_{k,1}^{-2}=e_k^{\frac{1}{2}}$ one obtains
\begin{align*}
\operatorname{Vol}(X)\cdot\sum_{\beta,\beta'}\left|\mathrm{al}_{2}(\mathscr{E},X)_{\beta\beta'}\right|
&\le\sum_{\beta,\beta'}\int_{\tilde{X}}\mathrm{d}x\int_{\tilde{X}}\mathrm{d}x'\,\left|K_X(x,x')_{\beta\beta'}\right|\\
&=h_{k,1}^{-2}\cdot\left\|\mathscr{E}^{[2,0]}(X;0,0)\right\|_{h_k}=e_k^{\frac{1}{2}}\cdot\left\|\mathscr{E}^{[2,0]}(X;0,0)\right\|_{h_k}.
\end{align*}
These are the two inequalities (2.2) of Proposition~16.3\,(b).

\textbf{The two bounds.}
One has $\operatorname{Vol}(X)=M^3\cdot|X|_M\ge M^3$.
From (2.2), (G2.r7) and (eq.7),
\begin{align*}
\left|\mathrm{al}_{0}(\mathscr{E},X)\right|&\le M^{-3}\cdot\|\mathscr{E}\|^x_k\cdot\exp\left(-\kappa\cdot d_M(X)\right),\\
\left|\mathrm{al}_{2}(\mathscr{E},X)_{11}\right|&\le M^{-3}\cdot e_k^{\frac{1}{2}}\cdot\|\mathscr{E}\|^x_k\cdot\exp\left(-\kappa\cdot d_M(X)\right).
\end{align*}
Since $\kappa\ge\kappa_0$ by the condition (C2) of Section~I.8.2, \cite[Theorem~13.2\,(t2)]{QEDI} gives
\[
\sum_{X\in\mathrm{Pl}(\square)}\exp\left(-\kappa\cdot d_M(X)\right)\le\sum_{X\in D_k:\,X\supset\square}\exp\left(-\kappa_0\cdot d_M(X)\right)\le K_0.
\]
Now $\varepsilon(\mathscr{E})=-\sum_{X\in\mathrm{Pl}(\square)}\mathrm{al}_{0}(\mathscr{E},X)$, and $m(\mathscr{E})$ is the entry $(1,1)$ of the matrix $m(\mathscr{E},\square)=-\sum_{X\in\mathrm{Pl}(\square)}\mathrm{al}_{2}(\mathscr{E},X)$.
Hence
\[
\left|\varepsilon(\mathscr{E})\right|\le K_0\cdot M^{-3}\cdot\|\mathscr{E}\|^x_k,\qquad\left|m(\mathscr{E})\right|\le K_0\cdot M^{-3}\cdot e_k^{\frac{1}{2}}\cdot\|\mathscr{E}\|^x_k,
\]
and $K_0\cdot M^{-3}\le K_0\le61\cdot c_{X}\cdot K_0=c_{\mathrm{fl}}$, because $M\ge1$ and $c_{X}=54$.

\paragraph{(c)}
Fix $(\mathscr{A},\mathscr{C})\in\tilde{R}_k\times\mathfrak{X}_k$.
The definition of $R$ in Section~I.5.6 gives
\begin{multline*}
\sum_{X\in D_k}\mathscr{E}(X;\mathscr{A},\mathscr{C})-\sum_{X\in D_k}(R\mathscr{E})(X;\mathscr{A},\mathscr{C})\\
=\sum_{X\ \mathrm{small}}\left[\mathrm{al}_{0}(\mathscr{E},X)\cdot\operatorname{Vol}(X)+\int_X\mathrm{d}x\,\bar{\psi}(x)^T\cdot\mathrm{al}_{2}(\mathscr{E},X)\cdot\psi(x)\right].
\end{multline*}
A polymer $X$ is the disjoint union of its $M$-cubes, and $X\times O_{\mathrm{cub}}$ is the disjoint union of the sets $\square\times O_{\mathrm{cub}}$, $\square\subset X$ (Section~I.1.4);
hence $\operatorname{Vol}(X)=\sum_{\square\subset X}M^3$ and $\int_X=\sum_{\square\subset X}\int_\square$.
All the sums are finite, and $\sum_{X\ \mathrm{small}}\sum_{\square\subset X}=\sum_{\square}\sum_{X\in\mathrm{Pl}(\square)}$.
The right-hand side is therefore
\begin{multline*}
\sum_{\square}\left[M^3\cdot\sum_{X\in\mathrm{Pl}(\square)}\mathrm{al}_{0}(\mathscr{E},X)+\int_\square\mathrm{d}x\,\bar{\psi}(x)^T\cdot\left(\sum_{X\in\mathrm{Pl}(\square)}\mathrm{al}_{2}(\mathscr{E},X)\right)\cdot\psi(x)\right]\\
=\sum_{\square}\left[-\varepsilon(\mathscr{E},\square)\cdot M^3-\int_\square\mathrm{d}x\,\bar{\psi}(x)^T\cdot m(\mathscr{E},\square)\cdot\psi(x)\right].
\end{multline*}
The torus $T_k$ consists of $\frac{\ell^3}{M^3}$ $M$-cubes, and the sets $\square\times O_{\mathrm{cub}}$ partition $T^{\mathrm{f}}_k$.
By (a), the last expression is thus equal to
\begin{align*}
-\varepsilon(\mathscr{E})\cdot M^3\cdot\frac{\ell^3}{M^3}-m(\mathscr{E})\cdot\sum_{\square}\int_\square\mathrm{d}x\,\bar{\psi}(x)\cdot\psi(x)
&=-\varepsilon(\mathscr{E})\cdot\ell^3-m(\mathscr{E})\cdot\langle\bar{\psi},\psi\rangle,
\end{align*}
which is the identity of (c), with $n_k=\ell^3$.
In coordinates,
\[
m(\mathscr{E})\cdot\langle\bar{\psi},\psi\rangle=m(\mathscr{E})\cdot\eta^3\cdot\frac{1}{48}\sum_{c\in T_k}\sum_{g\in O_{\mathrm{cub}}}\sum_{\beta=1}^{4}\bar{\psi}_\beta(c,g)\cdot\psi_\beta(c,g):
\]
each of the $192$ components over a photon site carries the same coefficient.

\paragraph{(d)}
\textbf{(d.1) The kernels of $R\mathscr{E}$.}
Let $X$ be small.
By Proposition~16.2\,(iv), the subtracted element $\mathrm{al}_{0}(\mathscr{E},X)\cdot\operatorname{Vol}(X)+\int_X\mathrm{d}x\,\bar{\psi}^T\cdot\mathrm{al}_{2}(\mathscr{E},X)\cdot\psi$ consists of the part $\mathscr{E}_{00}(X;0,0)$ of bidegree $(0,0)$ and of a part of bidegree $(2,0)$ with the kernel $\mathrm{Kd}_X\left[\mathrm{al}_{2}(\mathscr{E},X)\right]$;
it has no other part.
Hence
\begin{align*}
(R\mathscr{E})_{00}(X;\mathscr{A},\mathscr{C})&=\mathscr{E}_{00}(X;\mathscr{A},\mathscr{C})-\mathscr{E}_{00}(X;0,0),\\
(R\mathscr{E})_{20}(X;\mathscr{A},\mathscr{C};x,x')&=\mathscr{E}_{20}(X;\mathscr{A},\mathscr{C};x,x')-\mathrm{Kd}_X\left[\mathrm{al}_{2}(\mathscr{E},X)\right](x,x'),
\end{align*}
and $(R\mathscr{E})_{nm}=\mathscr{E}_{nm}$ for every other bidegree $(n,m)$.
By (a.3), $R\mathscr{E}$ depends linearly on $\mathscr{E}$.

\textbf{(d.2) Normalisation.}
By (d.1), $(R\mathscr{E})_{00}(X;0,0)=0$.
With both arguments over a lift $\tilde{X}$, Proposition~16.2\,(iv) and (eq.5) give
\[
\int_{\tilde{X}}\mathrm{d}x\int_{\tilde{X}}\mathrm{d}x'\,(R\mathscr{E})_{20}(X;0,0;x,x')=\operatorname{Vol}(X)\cdot\mathrm{al}_{2}(\mathscr{E},X)-\operatorname{Vol}(X)\cdot\mathrm{al}_{2}(\mathscr{E},X)=0.
\]

\textbf{(d.3) The bound.}
Let $X$ be small.
By Proposition~16.2\,(iv) (the norm, with $\mathsf{w}=h_k$) and by (2.2),
\begin{align*}
\left\|\int_X\mathrm{d}x\,\bar{\psi}^T\cdot\mathrm{al}_{2}(\mathscr{E},X)\cdot\psi\right\|_{h_k}
&=h_{k,1}^2\cdot\operatorname{Vol}(X)\cdot\sum_{\beta,\beta'}\left|\mathrm{al}_{2}(\mathscr{E},X)_{\beta\beta'}\right|\\
&\le\left\|\mathscr{E}^{[2,0]}(X;0,0)\right\|_{h_k},
\end{align*}
and $\left|\mathrm{al}_{0}(\mathscr{E},X)\cdot\operatorname{Vol}(X)\right|=\left\|\mathscr{E}^{[0,0]}(X;0,0)\right\|_{h_k}$.
The norm is additive over the bidegrees (G2.r7);
so the two subtracted terms together have norm at most $\left\|\mathscr{E}(X;0,0)\right\|_{h_k}$, and
\[
\left\|(R\mathscr{E})(X;\mathscr{A},\mathscr{C})\right\|_{h_k}\le\left\|\mathscr{E}(X;\mathscr{A},\mathscr{C})\right\|_{h_k}+\left\|\mathscr{E}(X;0,0)\right\|_{h_k};
\]
for large $X$ the second term may even be omitted.
Multiplying by $\exp\left(\kappa\cdot d_M(X)\right)$ and applying (k6) at $(\mathscr{A},\mathscr{C})$ and at $(0,0)$ (see (pre1)), one obtains $\|R\mathscr{E}\|^x_k\le2\|\mathscr{E}\|^x_k$;
moreover $2\le c_{\mathrm{fl}}$.

\textbf{(d.4) The clauses of Section~I.5.4.}
For large $X$ one has $(R\mathscr{E})(X)=\mathscr{E}(X)$, and by Proposition~16.1\,(a1) the polymer $\tilde{g}X$ is large if and only if $X$ is.
Let $X$ be small and put $\mathrm{Mx}:=\mathrm{al}_{2}(\mathscr{E},X)$.
The subtracted terms are a number and an element with one barred and one unbarred leg;
hence $(R\mathscr{E})(X)$ is even and balanced.

(k1), (k2): the subtracted terms do not depend on $(\mathscr{A},\mathscr{C})$, and the kernel $\mathrm{Kd}_X[\mathrm{Mx}]$ vanishes unless both arguments lie over $X$.

(k3): by Proposition~16.1\,(c), $(\tilde{g}\mathscr{A},\tilde{g}\mathscr{C})$ lies in $\tilde{R}_k\times\mathfrak{X}_k$.
By (k3) for $\mathscr{E}$, by Proposition~16.2\,(i) and by the first relation of Proposition~16.2\,(iv),
\begin{align*}
a_{\tilde{g}}\left[(R\mathscr{E})(\tilde{g}X;\tilde{g}\mathscr{A},\tilde{g}\mathscr{C})\right]
&=\mathscr{E}(X;\mathscr{A},\mathscr{C})-\mathrm{al}_{0}(\mathscr{E},\tilde{g}X)\cdot\operatorname{Vol}(\tilde{g}X)\\
&\qquad-a_{\tilde{g}}\left[\int_{\tilde{g}X}\mathrm{d}x\,\bar{\psi}^T\cdot S\cdot\mathrm{Mx}\cdot S^{-1}\cdot\psi\right]\\
&=\mathscr{E}(X;\mathscr{A},\mathscr{C})-\mathrm{al}_{0}(\mathscr{E},X)\cdot\operatorname{Vol}(X)-\int_X\mathrm{d}x\,\bar{\psi}^T\cdot\mathrm{Mx}\cdot\psi,
\end{align*}
and the last line is $(R\mathscr{E})(X;\mathscr{A},\mathscr{C})$.

(k4): let the pairs $(\mathscr{A},\mathscr{C})$, $(\mathscr{A}^1,\mathscr{C}^1)$ and the functions $\boldsymbol{\lambda}$, $\boldsymbol{\lambda}'$ be as in (k4).
The automorphism $G_{\boldsymbol{\lambda}+\boldsymbol{\lambda}'}$ fixes the numbers and, by the second relation of Proposition~16.2\,(iv), the element $\int_X\mathrm{d}x\,\bar{\psi}^T\cdot\mathrm{Mx}\cdot\psi$;
together with (k4) for $\mathscr{E}$ this gives
\begin{align*}
G_{\boldsymbol{\lambda}+\boldsymbol{\lambda}'}\left[(R\mathscr{E})(X;\mathscr{A}^1,\mathscr{C}^1)\right]
&=\mathscr{E}(X;\mathscr{A},\mathscr{C})-\mathrm{al}_{0}(\mathscr{E},X)\cdot\operatorname{Vol}(X)-\int_X\mathrm{d}x\,\bar{\psi}^T\cdot\mathrm{Mx}\cdot\psi\\
&=(R\mathscr{E})(X;\mathscr{A},\mathscr{C}).
\end{align*}

(k5): by Proposition~16.1\,(c), $(-\mathscr{A},-\mathscr{C})$ lies in the domain.
The automorphism $\mathrm{Cg}$ fixes the numbers, and
\[
\mathrm{Cg}\left[\int_X\mathrm{d}x\,\bar{\psi}^T\cdot\mathrm{Mx}\cdot\psi\right]=\int_X\mathrm{d}x\,\bar{\psi}^T\cdot\mathsf{C}^{-1}\cdot\mathrm{Mx}^T\cdot\mathsf{C}\cdot\psi=\int_X\mathrm{d}x\,\bar{\psi}^T\cdot\mathrm{Mx}\cdot\psi
\]
by Proposition~16.2\,(iv) and (ii);
with (k5) for $\mathscr{E}$ it follows that $\mathrm{Cg}\left[(R\mathscr{E})(X;-\mathscr{A},-\mathscr{C})\right]=(R\mathscr{E})(X;\mathscr{A},\mathscr{C})$.

(k6): this is (d.3).

(k7), for $\mathscr{E}\in K^x_k$: by Proposition~16.1\,(c), $(-\overline{\mathscr{A}},-\overline{\mathscr{C}})$ lies in the domain.
One has $\mathbf{K}\left[\mathrm{al}_{0}(\mathscr{E},X)\cdot\operatorname{Vol}(X)\right]=\overline{\mathrm{al}_{0}(\mathscr{E},X)}\cdot\operatorname{Vol}(X)=\mathrm{al}_{0}(\mathscr{E},X)\cdot\operatorname{Vol}(X)$, and
\[
\mathbf{K}\left[\int_X\mathrm{d}x\,\bar{\psi}^T\cdot\mathrm{Mx}\cdot\psi\right]=\int_X\mathrm{d}x\,\bar{\psi}^T\cdot\varkappa^{-1}\cdot\overline{\mathrm{Mx}}\cdot\varkappa\cdot\psi=\int_X\mathrm{d}x\,\bar{\psi}^T\cdot\mathrm{Mx}\cdot\psi
\]
by Proposition~16.2\,(iv) and (iii);
with (k7) for $\mathscr{E}$ it follows that $\mathbf{K}\left[(R\mathscr{E})(X;-\overline{\mathscr{A}},-\overline{\mathscr{C}})\right]=(R\mathscr{E})(X;\mathscr{A},\mathscr{C})$.

Consequently $R$ maps $K^{x,C}_{k}$ into itself and $K^{x}_{k}$ into itself.

\textbf{(d.5) Restriction.}
For a family $E$ without slot, $(RE)(X;\mathscr{A})$ equals $E(X;\mathscr{A})-\mathrm{al}_{0}(E,X)\cdot\operatorname{Vol}(X)-\int_X\mathrm{d}x\,\bar{\psi}^T\cdot\mathrm{al}_{2}(E,X)\cdot\psi$ if $X$ is small, and $E(X;\mathscr{A})$ if $X$ is large (Section~I.5.11).
By (a.3), for small $X$,
\begin{align*}
(R\mathscr{E})(X;\mathscr{A},0)
&=\mathscr{E}|_0(X;\mathscr{A})-\mathrm{al}_{0}(\mathscr{E}|_0,X)\cdot\operatorname{Vol}(X)-\int_X\mathrm{d}x\,\bar{\psi}^T\cdot\mathrm{al}_{2}(\mathscr{E}|_0,X)\cdot\psi\\
&=\left(R(\mathscr{E}|_0)\right)(X;\mathscr{A}),
\end{align*}
and for large $X$ both sides equal $\mathscr{E}|_0(X;\mathscr{A})$.
Thus $(R\mathscr{E})|_0=R(\mathscr{E}|_0)$.

\textbf{(d.6) The clause (k2) alone.}
Let $E$ be a family of level $k$ in the sense of (G2.r8).
Its antisymmetric kernels at $(0,0)$ (see (pre1)) define $\mathrm{al}_{0}(E,X)$, the matrix $E_{20}(X;0,0;x,x')$ of Section~I.5.1, the matrix $\mathrm{al}_{2}(E,X)$ and the family $RE$;
$\mathrm{al}_{2}(E,X)$ does not depend on the lift, since the argument for Proposition~16.2\,(o) applies to every kernel with two arguments.
The kernels of $(RE)(X)$ differ from those of $E(X)$ by a number and by terms of bidegree $(2,0)$ which are constant in $(\mathscr{A},\mathscr{C})$ and vanish unless both arguments lie over $X$ (Proposition~16.2\,(iv), Lemma~0.3\,(a)).
Hence, if $E$ has (k2), then so has $RE$.

\paragraph{(e)}
The setting is (G2.S-2): $\ell/L\in\mathrm{Cl}$ and $e_{k+1}=L^{\frac{1}{2}}\cdot e_k\le e_0$.
Let $(\mathscr{A},\mathscr{C})\in\tilde{R}_{k+1}\times\mathfrak{X}_{k+1}$.
By Proposition~18.1\,(b), $\mathscr{C}_L$ lies in $L^{-\frac{3}{4}}\mathfrak{X}_k=\left\{L^{-\frac{3}{4}}\cdot\mathscr{C}':\mathscr{C}'\in\mathfrak{X}_k\right\}$.
This set is contained in $\mathfrak{X}_k$:
by Section~I.5.3, $\mathfrak{X}_k$ is the set of the complex bond functions $\mathscr{C}'$ with $|\mathscr{C}'|_\infty<e_k^{-1+2\varepsilon}$ and $|\partial\mathscr{C}'|_\infty<e_k^{-1+3\varepsilon}$;
the two norms of $L^{-\frac{3}{4}}\cdot\mathscr{C}'$ are $L^{-\frac{3}{4}}$ times those of $\mathscr{C}'$, and $L^{-\frac{3}{4}}\le1$.
By Proposition~18.1\,(c), $\mathscr{A}_L$ lies in $L^{-\frac{3}{4}}\tilde{R}_k$, which is contained in $\tilde{R}_k$.
So $(\mathscr{A}_L,\mathscr{C}_L)\in\tilde{R}_k\times\mathfrak{X}_k$, and $E'(X;\mathscr{A}_L,\mathscr{C}_L)$ is defined;
the sum over the $X\in D_k$ with $\bar{X}^{L}=LY$ is finite, and $(LE')(Y;\mathscr{A},\mathscr{C})$ is an element of $\mathrm{Gr}[\psi,\chi]$ of level $k+1$.

\textbf{(k3).}
Let $\tilde{g}'=(S,a')\in\mathbb{G}_{k+1}$, with a vector $a'\in M\mathbb{Z}^3$, and let $\tilde{g}:=(S,La')\in\mathbb{G}_k^{L}$, with the vector $La'$;
this is the pair of Proposition~16.1\,(d1).
If the elements are not even, (k3) is a relation that refers to a vector of the class (step (b.6) of the proof of Proposition~16.1), and it is asserted for every vector of the class (Section~I.4.3);
the computation below applies (k3) of $E'$ with the vector $La'$ and yields (k3) of $LE'$ with the vector $a'$.
By Proposition~16.1\,(a1) at level $k+1$, $\tilde{g}'Y$ lies in $D_{k+1}$, and by Proposition~16.1\,(c) at level $k+1$ (side $\ell/L\in\mathrm{Cl}$), $(\tilde{g}'\mathscr{A},\tilde{g}'\mathscr{C})$ lies in $\tilde{R}_{k+1}\times\mathfrak{X}_{k+1}$.
By Proposition~16.1\,(a1) and (d2), $\overline{\tilde{g}X}^{L}=\tilde{g}\left(\bar{X}^{L}\right)$ and $\tilde{g}(LY)=L(\tilde{g}'Y)$.
Hence $\bar{X}^{L}=LY$ holds if and only if $\overline{\tilde{g}X}^{L}=L(\tilde{g}'Y)$, and $X\mapsto\tilde{g}X$ is a bijection of $\left\{X\in D_k:\bar{X}^{L}=LY\right\}$ onto $\left\{X\in D_k:\bar{X}^{L}=L(\tilde{g}'Y)\right\}$.
Using Proposition~16.1\,(d3) (that is, $(\tilde{g}'\mathscr{A})_L=\tilde{g}(\mathscr{A}_L)$, and the same for $\mathscr{C}$), Proposition~16.1\,(d4) (that is, $a_{\tilde{g}'}S_L=S_La_{\tilde{g}}$) and (k3) of $E'$ for $\tilde{g}\in\mathbb{G}_k^{L}$, one obtains
\begin{align*}
a_{\tilde{g}'}\left[(LE')(\tilde{g}'Y;\tilde{g}'\mathscr{A},\tilde{g}'\mathscr{C})\right]
&=\sum_{X:\,\bar{X}^{L}=LY}a_{\tilde{g}'}S_L\left[E'\left(\tilde{g}X;\tilde{g}(\mathscr{A}_L),\tilde{g}(\mathscr{C}_L)\right)\right]\\
&=\sum_{X}S_La_{\tilde{g}}\left[E'\left(\tilde{g}X;\tilde{g}(\mathscr{A}_L),\tilde{g}(\mathscr{C}_L)\right)\right]\\
&=\sum_{X}S_L\left[E'(X;\mathscr{A}_L,\mathscr{C}_L)\right]=(LE')(Y;\mathscr{A},\mathscr{C}).
\end{align*}

\textbf{(k4).}
Let $(\mathscr{A},\mathscr{C})$ and $(\mathscr{A}^1,\mathscr{C}^1)$ lie in $\tilde{R}_{k+1}\times\mathfrak{X}_{k+1}$, let $\boldsymbol{\lambda}$, $\boldsymbol{\lambda}'$ be real functions on $T_{k+1}$, and let $\mathscr{A}^1=\mathscr{A}-d\boldsymbol{\lambda}$, $\mathscr{C}^1=\mathscr{C}-d\boldsymbol{\lambda}'$ on the bonds of $T_{k+1}$ at distance $\le1$ from $Y$.
Let $\bar{X}^{L}=LY$;
then $X\subset LY$.
Let $b$ be a bond of $T_k$ at distance $\le1$ from $X$.
Both end points of $b$ have sup-distance at most $1$ from $LY$;
hence both end points of $b/L$ have sup-distance at most $\frac{1}{L}\le1$ from $Y$.
Since $(d\Lambda)_L=d(\Lambda_L)$ by Proposition~16.1\,(d4), with $\Lambda_L(c)=L^{\frac{1}{2}}\cdot\Lambda(c/L)$ as in (G2.r5), it follows that
\[
\mathscr{A}^1_L(b)=L^{-\frac{1}{2}}\cdot\mathscr{A}^1(b/L)=L^{-\frac{1}{2}}\cdot\left(\mathscr{A}(b/L)-(d\boldsymbol{\lambda})(b/L)\right)=\mathscr{A}_L(b)-\left(d(\boldsymbol{\lambda}_L)\right)(b),
\]
and likewise $\mathscr{C}^1_L(b)=\mathscr{C}_L(b)-\left(d(\boldsymbol{\lambda}'_L)\right)(b)$.
The functions $\boldsymbol{\lambda}_L$, $\boldsymbol{\lambda}'_L$ are real functions on $T_k$, and both pairs $(\mathscr{A}_L,\mathscr{C}_L)$, $(\mathscr{A}^1_L,\mathscr{C}^1_L)$ lie in $\tilde{R}_k\times\mathfrak{X}_k$.
By (k4) of $E'$,
\[
G_{\boldsymbol{\lambda}_L+\boldsymbol{\lambda}'_L}\left[E'(X;\mathscr{A}^1_L,\mathscr{C}^1_L)\right]=E'(X;\mathscr{A}_L,\mathscr{C}_L).
\]
Moreover $(\boldsymbol{\lambda}+\boldsymbol{\lambda}')_L=\boldsymbol{\lambda}_L+\boldsymbol{\lambda}'_L$, and $G_{\Lambda}S_L=S_LG_{\Lambda_L}$ by Proposition~16.1\,(d4);
therefore
\begin{align*}
G_{\boldsymbol{\lambda}+\boldsymbol{\lambda}'}\left[(LE')(Y;\mathscr{A}^1,\mathscr{C}^1)\right]
&=\sum_{X}S_LG_{\boldsymbol{\lambda}_L+\boldsymbol{\lambda}'_L}\left[E'(X;\mathscr{A}^1_L,\mathscr{C}^1_L)\right]\\
&=\sum_{X}S_L\left[E'(X;\mathscr{A}_L,\mathscr{C}_L)\right]=(LE')(Y;\mathscr{A},\mathscr{C}).
\end{align*}

\textbf{(k5), (k7).}
By Proposition~16.1\,(c) at level $k+1$, the pairs $(-\mathscr{A},-\mathscr{C})$ and $(-\overline{\mathscr{A}},-\overline{\mathscr{C}})$ lie in $\tilde{R}_{k+1}\times\mathfrak{X}_{k+1}$.
As $L^{-\frac{1}{2}}$ is real, $(-\mathscr{A})_L=-\mathscr{A}_L$ and $(\overline{\mathscr{A}})_L=\overline{\mathscr{A}_L}$, and the same holds for $\mathscr{C}$.
With $\mathrm{Cg}\,S_L=S_L\,\mathrm{Cg}$ and $\mathbf{K}S_L=S_L\mathbf{K}$ (Proposition~16.1\,(d4)) one finds
\begin{align*}
\mathrm{Cg}\left[(LE')(Y;-\mathscr{A},-\mathscr{C})\right]&=\sum_{X}S_L\,\mathrm{Cg}\left[E'(X;-\mathscr{A}_L,-\mathscr{C}_L)\right]=\sum_{X}S_L\left[E'(X;\mathscr{A}_L,\mathscr{C}_L)\right],\\
\mathbf{K}\left[(LE')(Y;-\overline{\mathscr{A}},-\overline{\mathscr{C}})\right]&=\sum_{X}S_L\mathbf{K}\left[E'(X;-\overline{\mathscr{A}_L},-\overline{\mathscr{C}_L})\right]=\sum_{X}S_L\left[E'(X;\mathscr{A}_L,\mathscr{C}_L)\right],
\end{align*}
where the second equalities use (k5), resp.\ (k7), of $E'$;
the right-hand sides equal $(LE')(Y;\mathscr{A},\mathscr{C})$.

\paragraph{(f)}
\textbf{(f1).}
The definition of $\mathrm{Sm}(\tilde{\square})$ involves $M$, $L$, $\mathbb{L}_k$ and nets in $\mathbb{R}^3$;
the side $\ell$ does not enter it.
For $v\in M\mathbb{Z}^3$, the translation by $v$ maps $M$-cubes of $\mathbb{L}_k$ onto $M$-cubes, common faces onto common faces, closed cubes onto closed cubes, and nets in $\mathbb{R}^3$ onto nets of the same length;
hence $d_M(\tilde{X}+v)=d_M(\tilde{X})$, and $\mathrm{Sm}(\tilde{\square}+v)=\mathrm{Sm}(\tilde{\square})+v$.

\emph{The image lies in the set named.}
Let $\tilde{X}\in\mathrm{Sm}(\tilde{\square})$.
By (Ar3), $2M\cdot\left(d_M(\tilde{X})+2\right)\le2M\cdot(L+2)<\ell$.
Since $\tilde{X}$ is not empty, the converse part of \cite[Theorem~13.2\,(t10)]{QEDI} applies:
$\mathrm{pr}$ is injective on $\tilde{X}$, $\mathrm{pr}(\tilde{X})$ is a polymer, and $d_M(\mathrm{pr}\,\tilde{X})=d_M(\tilde{X})\le L$.
(That assertion concerns the union $\mathrm{Xr}$ of the $M$-cubes of $\mathbb{R}^3$ whose centres are the centres of the cubes of $\tilde{X}$;
the set of the sites of the polymer $\mathrm{pr}(\mathrm{Xr})$ is $\mathrm{pr}(\tilde{X})$, by the argument that concludes the proof of Lemma~0.2\,(c).)
Moreover $\mathrm{pr}(\tilde{X})$ contains $\mathrm{pr}(\tilde{\square})$.

\emph{Surjectivity.}
Let $X$ be a small polymer that contains $\mathrm{pr}(\tilde{\square})$, and let $\tilde{X}'$ be a lift of $X$ with $d_M(\tilde{X}')=d_M(X)\le L$ (Lemma~0.2\,(c));
by Lemma~0.2\,(b), $\tilde{X}'$ is a non-empty finite union of $M$-cubes of $\mathbb{L}_k$ that is connected through faces.
The cube of $\tilde{X}'$ over $\mathrm{pr}(\tilde{\square})$ (Lft4) has its centre over the centre of $\mathrm{pr}(\tilde{\square})$;
it is therefore $\tilde{\square}+\ell n$ for some $n\in\mathbb{Z}^3$.
The set $\tilde{X}:=\tilde{X}'-\ell n$ belongs to $\mathrm{Sm}(\tilde{\square})$, and $\mathrm{pr}(\tilde{X})=X$.

\emph{Injectivity.}
Let $\tilde{X}_1,\tilde{X}_2\in\mathrm{Sm}(\tilde{\square})$ with $\mathrm{pr}(\tilde{X}_1)=\mathrm{pr}(\tilde{X}_2)=:X$.
By \cite[Theorem~13.2\,(t7)]{QEDI} in $\mathbb{R}^3$, the sup-diameter of $\tilde{X}_i$, increased by $\eta$, is at most $M\cdot(L+2)+\eta$, which is smaller than $\ell$ by the fifth inequality of (Ar3);
so both sets are lifts of $X$.
By Lemma~0.2\,(b), $\tilde{X}_2=\tilde{X}_1+\ell n$ with some $n\in\mathbb{Z}^3$.
Then $\tilde{\square}$ and $\tilde{\square}-\ell n$ are cubes of $\tilde{X}_1$ with the same image;
since $\mathrm{pr}$ is injective on $\tilde{X}_1$ (Lft1), $n=0$.

\emph{$\tilde{X}$ is the lift of $\mathrm{pr}(\tilde{X})$.}
It is a lift, as has just been shown, and every lift of $\mathrm{pr}(\tilde{X})$ is a translate of it by a vector of $\ell\mathbb{Z}^3$ (Lemma~0.2\,(b));
by Lemma~0.2\,(c) it is, up to such a translation, the set of the sites in the union of cubes of \cite[Theorem~13.2\,(t10)]{QEDI}.

\textbf{(f2).}
Let $\mathrm{rg}:=\eta$.

\emph{The maps.}
Let $\tilde{X}$ be a non-empty finite union of $M$-cubes of $\mathbb{L}_k$ that is connected through faces, let $\tilde{U}:=\tilde{X}^{+\eta}$, and let $\tilde{K}^0$, $\tilde{K}$ be data on $\tilde{U}$ (Section~0.5).
Denote by $\tilde{K}^0_{20}(x,x')$ the $4\times4$ matrix with the entries $\tilde{K}^0_{20}\left((x,\beta,1),(x',\beta',0)\right)$.
Define
\begin{align*}
\Phi_0(\tilde{X},\tilde{K}^0)&:=\operatorname{Vol}(\tilde{X})^{-1}\cdot\tilde{K}^0_{00},\\
\Phi_2(\tilde{X},\tilde{K}^0)&:=\operatorname{Vol}(\tilde{X})^{-1}\int_{\tilde{X}}\mathrm{d}x\int_{\tilde{X}}\mathrm{d}x'\,\tilde{K}^0_{20}(x,x')\\
&\phantom{:}=\operatorname{Vol}(\tilde{X})^{-1}\cdot\left(\frac{\eta^3}{48}\right)^2\sum_{c,c'\in\tilde{X}}\sum_{g,g'\in O_{\mathrm{cub}}}\tilde{K}^0_{20}\left((c,g),(c',g')\right),
\end{align*}
and
\[
\Phi_R(\tilde{X},\tilde{K}^0,\tilde{K}):=
\begin{cases}
\mathrm{El}(\tilde{K})-\tilde{K}^0_{00}-\int_{\tilde{X}}\mathrm{d}x\,\bar{\psi}(x)^T\cdot\Phi_2(\tilde{X},\tilde{K}^0)\cdot\psi(x) & \text{if } d_M(\tilde{X})\le L,\\
\mathrm{El}(\tilde{K}) & \text{otherwise.}
\end{cases}
\]
Here $\Phi_0$ is a number, $\Phi_2$ a matrix, and $\Phi_R$ an element of $\mathrm{Grc}_k(\tilde{U})$.
The three maps involve $\eta$, $M$ and $L$, and no side $\ell$.
With $K=\tilde{K}^0$ and $F=\mathrm{El}(\tilde{K})$ they are the maps displayed in Proposition~16.3\,(f2).

\emph{Translations.}
Let $v\in M\mathbb{Z}^3$;
this set contains $LM\mathbb{Z}^3$.
The family $(T_v\tilde{K})_{nm}(\xi;\boldsymbol{\zeta}):=\tilde{K}_{nm}(\xi-v;\boldsymbol{\zeta}-v)$ is a datum on $\tilde{U}+v=(\tilde{X}+v)^{+\eta}$, and the rules $T_v\psi(\xi):=\psi(\xi+v)$, $T_v\chi(\boldsymbol{\zeta}):=\chi(\boldsymbol{\zeta}+v)$ define the isomorphism $T_v:\mathrm{Grc}_k(\tilde{U})\to\mathrm{Grc}_k(\tilde{U}+v)$.
By (f1), $\operatorname{Vol}(\tilde{X}+v)=\operatorname{Vol}(\tilde{X})$ and $d_M(\tilde{X}+v)=d_M(\tilde{X})$.
Substituting $\xi+v$, $\boldsymbol{\zeta}+v$, $x+v$ for the variables of summation, one obtains, for every $4\times4$ matrix $\mathrm{Mx}$,
\begin{gather*}
\mathrm{El}(T_v\tilde{K})=T_v\,\mathrm{El}(\tilde{K}),\\
\int_{\tilde{X}+v}\mathrm{d}x\int_{\tilde{X}+v}\mathrm{d}x'\,(T_v\tilde{K}^0)_{20}(x,x')=\int_{\tilde{X}}\mathrm{d}x\int_{\tilde{X}}\mathrm{d}x'\,\tilde{K}^0_{20}(x,x'),\\
\int_{\tilde{X}+v}\mathrm{d}x\,\bar{\psi}(x)^T\cdot\mathrm{Mx}\cdot\psi(x)=T_v\int_{\tilde{X}}\mathrm{d}x\,\bar{\psi}(x)^T\cdot\mathrm{Mx}\cdot\psi(x).
\end{gather*}
Therefore
\begin{gather*}
\Phi_0(\tilde{X}+v,T_v\tilde{K}^0)=\Phi_0(\tilde{X},\tilde{K}^0),\qquad\Phi_2(\tilde{X}+v,T_v\tilde{K}^0)=\Phi_2(\tilde{X},\tilde{K}^0),\\
\Phi_R(\tilde{X}+v,T_v\tilde{K}^0,T_v\tilde{K})=T_v\,\Phi_R(\tilde{X},\tilde{K}^0,\tilde{K}).
\end{gather*}

\emph{The torus.}
Let $\ell\in\mathrm{Cl}$, let $X$ be a polymer of $T_k$ with a lift $\tilde{X}$ whose sup-diameter, increased by $2\eta+M$, is smaller than $\ell$ (every small polymer has such a lift, by Lemma~0.2\,(c) and (Ar3)), let $(\mathscr{A},\mathscr{C})\in\tilde{R}_k\times\mathfrak{X}_k$, and let $\tilde{K}^0$, $\tilde{K}$ be the data of $\mathscr{E}(X;0,0)$ and of $\mathscr{E}(X;\mathscr{A},\mathscr{C})$ on $\tilde{U}$:
the antisymmetric kernels, which are functions on the covering index sets (Section~0.4), taken at the indices over $\tilde{U}$.
By Lemma~0.2\,(b), $\tilde{X}$ is a union of the kind on which the maps are defined, and $\operatorname{Vol}(\tilde{X})=\operatorname{Vol}(X)$.
By (k2) and Lemma~0.4,
\[
\mathscr{E}(X;\mathscr{A},\mathscr{C})=\mathrm{cov}_\ell\left(\mathrm{El}(\tilde{K})\right).
\]
If $X$ is small, then $\mathrm{al}_{0}(\mathscr{E},X)=\operatorname{Vol}(X)^{-1}\cdot\mathscr{E}_{00}(X;0,0)=\Phi_0(\tilde{X},\tilde{K}^0)$, and $\mathrm{al}_{2}(\mathscr{E},X)=\Phi_2(\tilde{X},\tilde{K}^0)$ by (eq.5).
For every matrix $\mathrm{Mx}$,
\[
\mathrm{cov}_\ell\left(\int_{\tilde{X}}\mathrm{d}x\,\bar{\psi}(x)^T\cdot\mathrm{Mx}\cdot\psi(x)\right)=\int_X\mathrm{d}x\,\bar{\psi}(x)^T\cdot\mathrm{Mx}\cdot\psi(x),
\]
because $\bar{\psi}(x)^T\cdot\mathrm{Mx}\cdot\psi(x)$ is $\ell\mathbb{Z}^3$-periodic on the torus and $\tilde{X}\times O_{\mathrm{cub}}$ contains exactly one point over each site over $X$ (Lft1);
moreover $\mathrm{cov}_\ell$ fixes the numbers.
By Lemma~0.2\,(d), $X$ is small if and only if $d_M(\tilde{X})\le L$.
Hence, in both cases of the definitions of $R$ and of $\Phi_R$,
\[
(R\mathscr{E})(X;\mathscr{A},\mathscr{C})=\mathrm{cov}_\ell\left(\Phi_R(\tilde{X},\tilde{K}^0,\tilde{K})\right).
\]
These are the three requirements of Section~I.1.9 (one map for all sides $\ell$; covariance under the translations by $LM\mathbb{Z}^3$; the value on the torus), with $\mathrm{rg}=\eta$ and with the data named above.
The identities hold for every lift $\tilde{X}$ of $X$.

\textbf{(f3).}
Let $T_k(\ell_i):=\mathbb{L}_k/\ell_i\mathbb{Z}^3$ be the torus $T_k$ formed with the side $\ell_i$, and let $\mathrm{pr}_i:\mathbb{L}_k\to T_k(\ell_i)$ be its covering map.
Let $\tilde{X}\in\mathrm{Sm}(\tilde{\square})$ and $X_i:=\mathrm{pr}_i\tilde{X}$ for $i=1,2$.
By (f1), $X_i$ is a small polymer of $T_k(\ell_i)$ with the lift $\tilde{X}$, and $\operatorname{Vol}(X_1)=\operatorname{Vol}(\tilde{X})=\operatorname{Vol}(X_2)$.
By the definitions, with (eq.5) for the lift $\tilde{X}$ on both tori,
\begin{align*}
\mathrm{al}_{0}(\mathscr{E}^{(i)},X_i)&=\operatorname{Vol}(\tilde{X})^{-1}\cdot\mathscr{E}^{(i)}_{00}(X_i;0,0),\\
\mathrm{al}_{2}(\mathscr{E}^{(i)},X_i)&=\operatorname{Vol}(\tilde{X})^{-1}\int_{\tilde{X}}\mathrm{d}x\int_{\tilde{X}}\mathrm{d}x'\,\mathscr{E}^{(i)}_{20}(X_i;0,0;x,x').
\end{align*}
By the hypothesis, the right-hand sides coincide for $i=1$ and $i=2$.
By (f1), $\tilde{X}\mapsto\mathrm{pr}_i\tilde{X}$ is a bijection of $\mathrm{Sm}(\tilde{\square})$ onto the set of the small polymers of $T_k(\ell_i)$ that contain $\mathrm{pr}_i\tilde{\square}$.
Hence
\begin{align*}
\varepsilon(\mathscr{E}^{(i)},\mathrm{pr}_i\tilde{\square})&=-\sum_{\tilde{X}\in\mathrm{Sm}(\tilde{\square})}\mathrm{al}_{0}(\mathscr{E}^{(i)},\mathrm{pr}_i\tilde{X}),\\
m(\mathscr{E}^{(i)},\mathrm{pr}_i\tilde{\square})&=-\sum_{\tilde{X}\in\mathrm{Sm}(\tilde{\square})}\mathrm{al}_{2}(\mathscr{E}^{(i)},\mathrm{pr}_i\tilde{X}),
\end{align*}
and for $i=1,2$ these are sums over one and the same index set with equal terms.

\textbf{Level 0.}
At $k=0$ one has $\eta=1$, the charge $e_k$ is the number $q$, and $h_{0,1}^{-2}=q^{\frac{1}{2}}$;
moreover $\int\mathrm{d}x=\frac{1}{48}\sum_{c}\sum_{g}$ and $\langle\bar{\psi},\psi\rangle=\frac{1}{48}\sum_{c,g,\beta}\bar{\psi}_\beta(c,g)\cdot\psi_\beta(c,g)$;
the kernel of the subtracted term is $48\cdot\mathbf{1}[x=x'\text{ over }X]\cdot\mathrm{al}_{2}(\mathscr{E},X)$;
$\tilde{U}=\tilde{X}^{+1}$;
the pairs inside a cluster have $d_X=1$ and are of kind (I) (step (2) of the proof of Lemma~0.4).
No step of the argument requires $k\ge1$.
For the family $\mathscr{E}=0$ (I.L0.21) one has $\mathrm{al}_{0}=0$, $\mathrm{al}_{2}=0$, $\varepsilon(0)=0$, $m(0)=0$ and $R\,0=0$.

\textbf{Uniformity.}
The constants in (b) and (d) are $K_0\cdot M^{-3}$, $c_{\mathrm{fl}}$ and $2$;
here $K_0$ and $c_{\mathrm{fl}}$ are the numbers of the line (C0) and $M$ is the constant of the line (C5) of Section~I.8.2.
None of them depends on $k$, $\ell$, $e_k$, $N$, the torus, $\mathscr{A}$ or $\mathscr{C}$.
\end{proof}

\subsection*{4. Conditions on the constants}

The proofs of Sections~0--3 rest on the following conditions on the constants.

\noindent
\begin{tabular}{p{0.19\textwidth}p{0.17\textwidth}p{0.24\textwidth}p{0.26\textwidth}}
\hline
condition & line of Section~I.8.2; label of Proposition~14.2 & consequence & places in the proofs\\
\hline
$L$ is an odd integer and $L\ge5$ & (C3); Proposition~14.2\,(L5) & $L$ odd, $L\ge1$; $L\ge3$ in the verification of the Remark following Proposition~16.2 & (Ar1), (Ar2): no site lies on a face; the sup-diameters of the cubes; (enl); (Ar3); (d.0): $L^{-i}\le L$ and $\frac{M}{L}\le M$; Proposition~16.3\,(e): $\frac{1}{L}\le1$ and $L^{-\frac{3}{4}}\le1$\\
\hline
$M$ is a power of $L$, $M\ge L^2$ & (C5); Proposition~14.2\,(M1) & $M$ is an odd integer and a multiple of $L$; together with $L\ge1$ (row (L5), line (C3)): $M\ge1$ and $L\le M$ & (Ar1): $M\mathbb{Z}^3\subset L\mathbb{Z}^3$, hence $M\mathbb{Z}^3\subset L^{-i}\mathbb{Z}^3$ for $-1\le i\le k$ ((a1.2)); (Ar2), (enl), (Ar3); $M\ge1$ in (Ar3) ($\eta\le1\le M\le L\cdot M$), in Lemma~0.4 ($\eta\le M$) and in Proposition~16.3\,(b) ($K_0\cdot M^{-3}\le K_0$); $L\le M$ in (d.0) ($2L\le2M$)\\
\hline
$c_{X}=54$, $\kappa_0=54\log72$, $K_0=\frac{\exp(\kappa_0)}{36}$, $c_{\mathrm{fl}}=61\cdot c_{X}\cdot K_0$ & (C0) & the bound of \cite[Theorem~13.2\,(t2)]{QEDI} with these numbers; $K_0\le c_{\mathrm{fl}}$; $2\le c_{\mathrm{fl}}$ & Proposition~16.3\,(b), (d)\\
\hline
$\kappa\ge7\kappa_0+7$ & (C2) & $\kappa\ge\kappa_0$ & Proposition~16.3\,(b)\\
\hline
$\varepsilon>0$ & the range of Section~I.0.5; it follows from $1-\alpha<\varepsilon$ of (C7), Proposition~14.2\,(S1), together with $\alpha<1$ of (C1) & $\varepsilon\ge0$; hence $1-e_k^{\varepsilon}\le0$ if $e_k\ge1$ & step (c.1) of the proof of Proposition~16.1, that is, the parenthesis of Proposition~16.1\,(c)\\
\hline
\end{tabular}

\medskip
The exponent $h_{k,1}=e_k^{-\frac{1}{4}}$ is fixed in Sections~I.5.1 and~I.8.5;
its value enters (2.2).

The side $\ell$ lies in $\mathrm{Cl}$;
the proofs rest on the facts that $\ell$ is a multiple of $L\cdot M$ and that $\ell\ge13L\cdot M$, and on the oddness of $\ell/(L\cdot M)$, which enters through the assertion of Proposition~16.1\,(d) that $\ell/L$ is an odd multiple of $M$.
In Proposition~16.1 the charge is a number $e_k>0$, and in Propositions~16.2 and~16.3 a number $e_k\in(0,e_0]$;
of the two conditions $e_k>0$ and $e_k\le e_0$, the proofs of Proposition~16.2 and of Proposition~16.3\,(a)--(d), (f) rest on the first alone.
The proof of Proposition~16.3\,(e) rests on both.
There $e_k\le e_0$ follows from $e_{k+1}=L^{\frac{1}{2}}\cdot e_k\le e_0$ and $L\ge1$;
it is part of the setting (G2.S-2) under which Proposition~18.1 is stated, and it is the form in which the condition (A8) enters, through the inequality $e_k^{\varepsilon}\le\frac{1}{4}$, labelled (nm5) in Chapter~18a, in the last sentence of the proof of Proposition~18.1\,(c).
In Proposition~16.3\,(e) the hypotheses $\ell/L\in\mathrm{Cl}$ and $e_{k+1}=L^{\frac{1}{2}}\cdot e_k\le e_0$ enter through Proposition~18.1\,(b), (c) and through Proposition~16.1\,(a1), (c) at level $k+1$;
through Proposition~18.1\,(b), (c), the assertion (e) rests in addition on the conditions of Proposition~18.1, namely $\alpha>0$, $\varepsilon\le\frac{1}{200}$, $L$ odd with $L^{-\frac{3}{4}}\le\frac{1}{2}$ (line (C3)), $M$ a power of $L$ (line (C5)), and (A8) of Section~I.8.4;
the lines of Section~I.8.2 and the rows of Proposition~14.2 that impose these conditions are named in Chapter~18a, Section~5.

\section*{Chapter 9. The operator $T_k(A)$ and its two pieces $T^{(1)}$, $T^{(2)}$ (FE5$^{\mathrm{ch}}$): proof of Proposition~9}

\subsection*{0. Conventions of this chapter}

This chapter contains the proof of Proposition~9, whose clauses are denoted (a)--(e).
The data are those of (SF), that is, a quadruple $(k,\ell,\mu,e_k)$ as in Section~G1.1 (equivalently, Section~I.1.8).
In Sections~1--6 it is assumed that $k\ge1$.
The labels (L5), (M1), (M2), (M5), (Z2.s1), (Z2.s2), (Z2.s4) denote rows of Proposition~14.2\,(b), and (G2), (G3) denote the two properties of the domain $\tilde{G}_k$ that carry these labels in \cite[Section~16.2.7]{QEDI}.

In this chapter the following letters have a local meaning.
$\mathrm{Ef}$, $\mathrm{Ef}_1$, $\mathrm{Ef}_2$, $\mathrm{Ef}_3$ are sets of fermion sites;
$\mathrm{Fa}_a(y)$ is a face of $\Delta_y$;
$\mathrm{ce}_j(c)$ denotes centres of blocks;
$\mathrm{jm}(c,a)$ is an index and $\mathrm{vz}(c,a)$ an exponent;
$\mathrm{Rw}_0,\dots,\mathrm{Rw}_3$ are the parts of a row, and they are not to be confused with the sums $\mathrm{Rw}_z(O)$ of Section~G1.3;
$\mathrm{sn}(S)$ is a sign, $\mathrm{Cn}$ a number and $\mathrm{Fn}$ a function;
$\mathrm{Rep}$, $\mathrm{coef}$, $\mathrm{mo}$ occur in Section~1.5, the sums $\mathrm{Su}_{\mathrm{in}}$, $\mathrm{Su}_1$, $\mathrm{Su}_2$ in Section~4.2, and the orders $\mathrm{od}_i$, $\mathrm{od}'_i$ of the three coordinates in Section~4.4.
In Section~4.4 the letters $\mathbb{W}$, $s_W$ and $\boldsymbol{\ell}$ have the meaning that they have in \cite[Lemma~16.26]{QEDI}: $\mathbb{W}$ is a box, a product of three intervals each of length at most $s_W$, and $\boldsymbol{\ell}$ is a closed lattice path in $\mathbb{W}$;
the bold letter $\boldsymbol{\ell}$ denotes a path and is distinct from the side $\ell$ of the torus.
The letters $n$, $n'$, $n''$ denote integers or vectors of $\mathbb{Z}^3$, as in Sections~G1.2 and G1.3.
The letter $t$ denotes an index of the source of an operator, that is, an element of the index set $\mathbb{L}^{\mathrm{f}}_k$ of (I.i4) or of the index set $\mathrm{Pzt}_k$ of (I.i5) in Section~I.1.9;
thus $t$ is a column index, which (I.1.4) writes $x'$ and (I.f4) writes $t'$, whereas the row index, written $t$ there, is here a unit site $y$.
In Section~4.4, however, $t$ is a complex number.
In Sections~1.2 and 4 the letter $a$ is a direction, $a\in\{\pm e_\mu:\mu=1,2,3\}$, as in (I.s4); in (9.P4) and in Sections~6.2 and 6.5 it is the translation part of an element $\tilde{g}=(S,a)$.
The symbol $T_k$ has two meanings, which are distinguished by the argument.
With a connection as its argument ($T_k(A)$, $T_k(-A)$, $T_k(\tilde{g}A)$, \dots) it is the operator of Section~I.3.5;
without an argument it is the torus $\mathbb{L}_k/\ell\mathbb{Z}^3$ of Section~I.1.1, also as a value of $\mathbf{Y}$ and as a lower index, as in $\mathbf{Y}=T_k$ and $T^{(1)}_{T_k}(A)$ (see (9.11) and Section~8.4).

\subsection*{1. Reading on the covering lattices; blocks, phases, spin matrices, norms}

\subsubsection*{1.1. Reading on the covering lattices}

A connection $A$ on $T_k$ is a complex bond function (Section~I.1.1).
It is identified with its $\ell\mathbb{Z}^3$-periodic extension to the bonds of $\mathbb{L}_k$.
A union $\mathbf{Y}$ of $M$-cubes of $T_k$ is identified with the $\ell\mathbb{Z}^3$-periodic set of the sites of $\mathbb{L}_k$ over it;
this set is a union of $M$-cubes $\{c\in\mathbb{L}_k:\ |c-z|<M/2\}$, $z\in M\mathbb{Z}^3$.
An element $\psi$ of $\mathfrak{S}^{\mathrm{ch}}(\eta_k,\ell)$ is a function on $\mathbb{L}^{\mathrm{f}}_k$ with $\psi(x+\ell e_\mu)=-\psi(x)$;
a function $\chi$ of pairs is a function on $\mathrm{Pzt}_k$ with $\chi(x+\ell e_\mu,x'+\ell e_\mu)=-\chi(x,x')$ (Section~I.1.7);
an element $v$ of $\mathfrak{S}(1,\ell)$ is a function on $\mathbb{Z}^3$ with $v(y+\ell e_\mu)=-v(y)$.
A unit site $y$ is a point of $\mathbb{Z}^3$, and $\Delta_y=\{c\in\mathbb{L}_k:\ |c-y|<1/2\}$.
Writing $x=(c,g)$, one has by (I.1.3) and (I.POS), for every function $\mathrm{Fn}$ on $\Delta_y\times O_{\mathrm{cub}}$,
\begin{equation}\tag{9.1}
\int_{\Delta_y}\mathrm{d}x\,\mathrm{Fn}(x)=\frac{\eta_k^3}{48}\sum_{c\in\Delta_y}\ \sum_{g\in O_{\mathrm{cub}}}\mathrm{Fn}(c,g).
\end{equation}
By (9.B1) below, $\Delta_y$ has $\eta_k^{-3}$ sites; hence $\int_{\Delta_y}\mathrm{d}x=1$.

The displays (I.5.4) contain the connection only through the numbers $\varphi_A(y,c)$ and $U_A(x,x^{(i)})$, which (I.1.6) and (I.3.3) define for every connection.
Accordingly, $T^{(i)}(A)=T^{(i)}_{\mathbf{Y}}(A)$ and $\bar{T}^{(i)}(A)=\bar{T}^{(i)}_{\mathbf{Y}}(A)$ denote the operators (I.5.4) for every connection $A$ and every union $\mathbf{Y}$ of $M$-cubes, with zero rows at the unit sites outside $\mathbf{Y}$ (Section~I.5.9).

\paragraph{(9.W) Values in $\mathfrak{S}(1,\ell)$.}
Let $n\in\mathbb{Z}^3$.
Then $\Delta_{y+\ell n}=\Delta_y+\ell n$;
$(x+\ell n)^{(i)}=x^{(i)}+\ell n$ by (I.s1);
$M_i(x+\ell n)=M_i(x)$, since $M_i$ depends on the label only (I.2.1);
$\varphi_A(y+\ell n,c+\ell n)=\varphi_A(y,c)$ and $U_A(x+\ell n,x^{(i)}+\ell n)=U_A(x,x^{(i)})$, since $A$ is periodic and blocks, centres and paths are translated;
and $c_3(x+\ell n)$ lies in $\mathbf{Y}$ if and only if $c_3(x)$ does.
Together with $\psi(x+\ell e_\mu)=-\psi(x)$ and $\chi(x+\ell e_\mu,x^{(3)}+\ell e_\mu)=-\chi(x,x^{(3)})$, the displays (I.5.4) yield
\begin{equation*}
\left(T^{(1)}(A)\psi\right)(y+\ell e_\mu)=-\left(T^{(1)}(A)\psi\right)(y),\qquad
\left(T^{(2)}(A)\chi\right)(y+\ell e_\mu)=-\left(T^{(2)}(A)\chi\right)(y).
\end{equation*}

\paragraph{Terminology.}
The phrase ``$(Of)(y)$ contains $f$ at the indices of a set'' means that $(Of)(y)$ is a linear combination, with $4\times4$ matrices as coefficients, of the values of $f$ at the indices of that set.
The phrase ``the row reads $A$ on the bonds of a set'' means that these coefficients are functions of the numbers $A(b)$, $b$ in that set.

\subsubsection*{1.2. Blocks ($L$ and $M$ are odd)}

\paragraph{(9.B1)}
Let $0\le j\le k$ and $z\in L^j\eta_k\mathbb{Z}^3$.
Since $L^j$ is odd,
\begin{equation*}
\mathrm{Bl}_j(z)=z+\eta_k\left\{n\in\mathbb{Z}^3:\ |n|\le\frac{L^j-1}{2}\right\},
\end{equation*}
a product of three intervals of $\eta_k\mathbb{Z}$, each of which has $L^j$ sites and the length $(L^j-1)\,\eta_k$.
The blocks $\mathrm{Bl}_j(z)$, $z\in L^j\eta_k\mathbb{Z}^3$, form a partition of $\mathbb{L}_k$, because every integer can be written in exactly one way as $L^jn'+n''$ with integers $n'$, $n''$ and $|n''|\le\frac{L^j-1}{2}$.
For a site $c$ let $\mathrm{ce}_j(c)$ denote the centre of the block of side $L^j\eta_k$ that contains $c$;
thus $\mathrm{ce}_0(c)=c$ and $\mathrm{ce}_k(c)=[c]$.
The unit block $\Delta_y=\mathrm{Bl}_k(y)$ has $L^{3k}=\eta_k^{-3}$ sites.

\paragraph{(9.B2)}
Let $0\le j<k$, $z\in L^{j+1}\eta_k\mathbb{Z}^3$ and $z'\in L^j\eta_k\mathbb{Z}^3$, and write $z'-z=L^j\eta_k\,n'$ with $n'\in\mathbb{Z}^3$.
If $|n'|\le\frac{L-1}{2}$, then $\mathrm{Bl}_j(z')\subset\mathrm{Bl}_{j+1}(z)$, because every $c\in\mathrm{Bl}_j(z')$ satisfies
\begin{equation*}
|c-z|\le\frac{(L^j-1)\,\eta_k}{2}+\frac{L^j\eta_k\,(L-1)}{2}=\frac{(L^{j+1}-1)\,\eta_k}{2}.
\end{equation*}
These $L^3$ blocks are pairwise disjoint, and together they have $L^3\cdot L^{3j}$ sites, as many as $\mathrm{Bl}_{j+1}(z)$.
Hence their union is $\mathrm{Bl}_{j+1}(z)$, and by (9.B1) a block $\mathrm{Bl}_j(z')$ with $|n'|>\frac{L-1}{2}$ does not meet $\mathrm{Bl}_{j+1}(z)$.

\paragraph{(9.B3)}
For every site $c$ and every $0\le j<k$ the blocks $\mathrm{Bl}_j(\mathrm{ce}_j(c))$ and $\mathrm{Bl}_{j+1}(\mathrm{ce}_{j+1}(c))$ have the point $c$ in common.
By (9.B2), therefore,
\begin{equation}\tag{9.2}
\mathrm{Bl}_j\left(\mathrm{ce}_j(c)\right)\subset\mathrm{Bl}_{j+1}\left(\mathrm{ce}_{j+1}(c)\right),\qquad
\left|\mathrm{ce}_{j+1}(c)-\mathrm{ce}_j(c)\right|\le\frac{(L^{j+1}-L^j)\,\eta_k}{2}.
\end{equation}
By induction, $\mathrm{Bl}_j(\mathrm{ce}_j(c))\subset\mathrm{Bl}_i(\mathrm{ce}_i(c))$ for $j\le i\le k$.
It follows that:
\begin{itemize}
\item[(i)] if $\mathrm{ce}_j(c)=\mathrm{ce}_j(c')$, then $c'$ lies in $\mathrm{Bl}_j(\mathrm{ce}_j(c))$, which is contained in $\mathrm{Bl}_i(\mathrm{ce}_i(c))$, so that $\mathrm{ce}_i(c')=\mathrm{ce}_i(c)$ for $j\le i\le k$;
\item[(ii)] for $c\in\Delta_y$ all centres $\mathrm{ce}_j(c)$, $0\le j\le k$, lie in $\Delta_y$, and $\mathrm{ce}_k(c)=y$;
\item[(iii)] by (9.B2) and induction on $k-j$, the unit block $\Delta_y=\mathrm{Bl}_k(y)$ is the disjoint union of $L^{3(k-j)}$ blocks of side $L^j\eta_k$, $0\le j\le k$.
\end{itemize}
In Section~I.3.2 the centres $\mathrm{ce}_j(c)$ are denoted $y_j$.

\paragraph{(9.B4)}
Let $y$ be a unit site and $z\in M\mathbb{Z}^3$ with $|y-z|\le\frac{M-1}{2}$; recall that $M$ is an odd integer.
Every $c\in\Delta_y$ satisfies
\begin{equation*}
|c-z|\le\frac{1-\eta_k}{2}+\frac{M-1}{2}<\frac{M}{2}.
\end{equation*}
So $\Delta_y$ is contained in the $M$-cube that contains $y$.
Consequently, if $y$ is a unit site of a union $\mathbf{Y}$ of $M$-cubes, then $\Delta_y\subset\mathbf{Y}$.

\paragraph{(9.B5)}
For a direction $a$ put
\begin{equation*}
\mathrm{Fa}_a(y):=\left\{c\in\Delta_y:\ c+\eta_ka\notin\Delta_y\right\}.
\end{equation*}
By (9.B1), $\mathrm{Fa}_a(y)=\left\{c\in\Delta_y:\ a\cdot(c-y)=\frac{1-\eta_k}{2}\right\}$, and this set has $\eta_k^{-2}$ sites.

\paragraph{(9.B6)}
For a direction $a$ and a site $c\in\Delta_y$ with $c+\eta_ka\in\Delta_y$ put
\begin{equation*}
\mathrm{jm}(c,a):=\min\left\{j\in\{1,\dots,k\}:\ \mathrm{ce}_j(c)=\mathrm{ce}_j(c+\eta_ka)\right\};
\end{equation*}
the set on the right contains $k$ by (9.B3)\,(ii).
For $1\le j\le k$ one has
\begin{equation}\tag{9.3}
\#\left|\left\{c\in\Delta_y:\ c+\eta_ka\in\Delta_y,\ \mathrm{jm}(c,a)=j\right\}\right|\le\eta_k^{-3}\cdot L^{-(j-1)}.
\end{equation}
For $j=1$ the right side is the number of the sites of $\Delta_y$.
Let $j\ge2$ and $\mathrm{jm}(c,a)=j$.
Then $\mathrm{ce}_{j-1}(c+\eta_ka)\ne\mathrm{ce}_{j-1}(c)$, that is, $c+\eta_ka$ does not lie in the block $\mathrm{Bl}_{j-1}(\mathrm{ce}_{j-1}(c))$;
by (9.B1) this means $a\cdot\left(c-\mathrm{ce}_{j-1}(c)\right)=\frac{(L^{j-1}-1)\,\eta_k}{2}$.
In each block $\mathrm{Bl}_{j-1}(z)$ the equation $a\cdot(c-z)=\frac{(L^{j-1}-1)\,\eta_k}{2}$ holds for $L^{2(j-1)}$ sites $c$, and by (9.B3)\,(iii) the unit block $\Delta_y$ consists of $L^{3(k-j+1)}$ such blocks.
The number on the left of (9.3) is therefore at most $L^{3(k-j+1)}\cdot L^{2(j-1)}=L^{3k}\cdot L^{-(j-1)}$.

\subsubsection*{1.3. Phases}

For $c\in\Delta_y$, the definitions (I.3.3) and (I.1.6), written with the centres of (9.B1), read
\begin{equation}\tag{9.4}
\begin{aligned}
&(\tau_kA)(y,c)=\sum_{j=0}^{k-1}(\tau A)\left(\mathrm{ce}_{j+1}(c),\mathrm{ce}_j(c)\right),\qquad
\varphi_A(y,c)=\exp\left\{\mathrm{i}\,e_k\cdot\eta_k\cdot(\tau_kA)(y,c)\right\},\\
&u_A(c,c')=\exp\left\{\mathrm{i}\,e_k\cdot\eta_k\cdot A(c,c')\right\}\ \ \text{(bonds)},\qquad
U_A(x,x^{(i)})=1\ \ (i\le2),\qquad
U_A(x,x^{(3)})=u_A(c,c_3).
\end{aligned}
\end{equation}

\paragraph{(9.P1) Paths.}
The path $\boldsymbol{\Gamma}^{\pi}(c,c')$ has $\eta_k^{-1}|c-c'|_1$ bonds, and its sites lie in the box
\begin{equation*}
\left\{c''\in\mathbb{L}_k:\ c''_\mu\ \text{lies between}\ c_\mu\ \text{and}\ c'_\mu,\ \mu=1,2,3\right\}.
\end{equation*}
A block is a product of intervals; hence it contains the paths $\boldsymbol{\Gamma}^{\pi}$ between any two of its sites.
By (9.2), for $c\in\Delta_y$ and $0\le j<k$ the six paths $\boldsymbol{\Gamma}^{\pi}\left(\mathrm{ce}_{j+1}(c),\mathrm{ce}_j(c)\right)$ lie in $\mathrm{Bl}_{j+1}(\mathrm{ce}_{j+1}(c))$, which is contained in $\Delta_y$, and each of them has at most $\frac{3}{2}\,(L^{j+1}-L^j)$ bonds.
Consequently $\varphi_A(y,c)$ reads $A$ only on bonds contained in $\Delta_y$.

\paragraph{(9.P2) Moduli.}
The statement \cite[Lemma~16.2(a)]{QEDI} concerns the photon sites and bonds of a block torus, and it applies to $T_k$, which is the block torus $(0,\ell,-1)$ of \cite[Definition~16.1]{QEDI}; there a photon site is written $x$, here $c$.
It gives, for every bond $(c,c')$ and every $c\in\Delta_y$,
\begin{equation*}
\left|u_A(c,c')^{\pm1}\right|\le\exp\left(\eta_k\cdot\mathfrak{i}(A)\right)\le\exp\left(\mathfrak{i}(A)\right),\qquad
\left|\varphi_A(y,c)^{\pm1}\right|\le\exp\left(\frac{3}{2}\,\mathfrak{i}(A)\right).
\end{equation*}

\paragraph{(9.P3) Gauge.}
Let $\boldsymbol{\lambda}$ be a function of $c$ on $T_k$.
On every bond one has $(d\boldsymbol{\lambda})(c,c')=\frac{\boldsymbol{\lambda}(c')-\boldsymbol{\lambda}(c)}{\eta_k}$ (Section~I.1.1; both sides are odd under the reversal of the bond).
Hence $\eta_k\cdot(d\boldsymbol{\lambda})(\boldsymbol{\Gamma})=\boldsymbol{\lambda}(c')-\boldsymbol{\lambda}(c)$ for every path $\boldsymbol{\Gamma}$ from $c$ to $c'$, $\eta_k\cdot(\tau\,d\boldsymbol{\lambda})(c,c')=\boldsymbol{\lambda}(c')-\boldsymbol{\lambda}(c)$, and
\begin{equation*}
\eta_k\cdot(\tau_k\,d\boldsymbol{\lambda})(y,c)=\sum_{j<k}\left(\boldsymbol{\lambda}\left(\mathrm{ce}_j(c)\right)-\boldsymbol{\lambda}\left(\mathrm{ce}_{j+1}(c)\right)\right)=\boldsymbol{\lambda}(c)-\boldsymbol{\lambda}(y).
\end{equation*}
The exponents in (9.4) are linear in $A$.
Therefore
\begin{equation}\tag{9.5}
\varphi_{A-d\boldsymbol{\lambda}}(y,c)\cdot\mathrm{e}^{\mathrm{i}e_k\boldsymbol{\lambda}(c)}=\mathrm{e}^{\mathrm{i}e_k\boldsymbol{\lambda}(y)}\cdot\varphi_A(y,c),\qquad
u_{A-d\boldsymbol{\lambda}}(c,c')\cdot\mathrm{e}^{\mathrm{i}e_k\boldsymbol{\lambda}(c')}=\mathrm{e}^{\mathrm{i}e_k\boldsymbol{\lambda}(c)}\cdot u_A(c,c').
\end{equation}

\paragraph{(9.P4) Lattice.}
Let $\tilde{g}=(S,a)\in\mathbb{G}_k$, with $a$ taken in $M\mathbb{Z}^3$, let $r:=r(S)\in\{\pm1\}^3$ (a diagonal matrix with the entries $\pm1$), and $\tilde{g}c=rc+a$ (Section~I.4.3).
The map $\tilde{g}$ maps each of the lattices $L^j\eta_k\mathbb{Z}^3$, $0\le j\le k$, onto itself, because $M\mathbb{Z}^3\subset\mathbb{Z}^3\subset L^j\eta_k\mathbb{Z}^3=L^{-(k-j)}\mathbb{Z}^3$;
it maps $M\mathbb{Z}^3$ onto itself, and $|\tilde{g}c-\tilde{g}c'|=|c-c'|$.
So $\tilde{g}\,\mathrm{Bl}_j(z)=\mathrm{Bl}_j(\tilde{g}z)$, $\tilde{g}\Delta_y=\Delta_{\tilde{g}y}$ and $\mathrm{ce}_j(\tilde{g}c)=\tilde{g}\,\mathrm{ce}_j(c)$;
moreover $\tilde{g}$ maps $M$-cubes onto $M$-cubes, unions of $M$-cubes onto unions of $M$-cubes, and the unit sites of $\mathbf{Y}$ onto the unit sites of $\tilde{g}\mathbf{Y}$.
Since $r$ is diagonal, the coordinate $\mu$ of $\tilde{g}c$ is a function of the coordinate $\mu$ of $c$;
so $\tilde{g}$ maps the path $\boldsymbol{\Gamma}^{\pi}(c,c')$, bond by bond, onto $\boldsymbol{\Gamma}^{\pi}(\tilde{g}c,\tilde{g}c')$, with the same order $\pi$.
By Section~I.4.3, $(\tilde{g}A)(\tilde{g}c,\tilde{g}c')=A(c,c')$, and $\tilde{g}A$ is $\ell\mathbb{Z}^3$-periodic.
Hence $(\tilde{g}A)\left(\boldsymbol{\Gamma}^{\pi}(\tilde{g}c,\tilde{g}c')\right)=A\left(\boldsymbol{\Gamma}^{\pi}(c,c')\right)$, $(\tau\tilde{g}A)(\tilde{g}c,\tilde{g}c')=(\tau A)(c,c')$, $(\tau_k\tilde{g}A)(\tilde{g}y,\tilde{g}c)=(\tau_kA)(y,c)$, and
\begin{equation}\tag{9.6}
\varphi_{\tilde{g}A}(\tilde{g}y,\tilde{g}c)=\varphi_A(y,c),\qquad
u_{\tilde{g}A}(\tilde{g}c,\tilde{g}c')=u_A(c,c').
\end{equation}

\paragraph{(9.P5) Conjugation, sign.}
The numbers $e_k$ and $\eta_k$ are real, and the linear forms $\tau$, $\tau_k$ have real coefficients, so that $(\tau_k\overline{A})(y,c)=\overline{(\tau_kA)(y,c)}$.
Hence
\begin{equation*}
\overline{\mathrm{i}\,e_k\cdot\eta_k\cdot(\tau_k\overline{A})(y,c)}=-\mathrm{i}\,e_k\cdot\eta_k\cdot(\tau_kA)(y,c),
\end{equation*}
and, for a bond, $\overline{\mathrm{i}\,e_k\cdot\eta_k\cdot\overline{A}(c,c')}=-\mathrm{i}\,e_k\cdot\eta_k\cdot A(c,c')$.
The exponents are linear in $A$.
It follows that
\begin{equation}\tag{9.7}
\overline{\varphi_{\overline{A}}(y,c)}=\varphi_{-A}(y,c)=\varphi_A(y,c)^{-1},\qquad
\overline{u_{\overline{A}}(c,c')}=u_{-A}(c,c')=u_A(c,c')^{-1}.
\end{equation}

\subsubsection*{1.4. Spin matrices}

\paragraph{(9.Sm1)}
By (I.m1), $M_i(x^{(i)})-M_i(x)=w_i\,\gamma\cdot\hat{d}_i(x)$ for $i=0,\dots,3$.
By (I.s3),
\begin{equation*}
\sum_{i=0}^{3}w_i\,\gamma\cdot\hat{d}_i(x)=\gamma\cdot\left(\sum_iw_i\,\hat{d}_i(x)\right)=0,
\qquad\text{that is,}\qquad
\sum_{i\le2}w_i\,\gamma\cdot\hat{d}_i(x)=-w_3\,\gamma\cdot\hat{d}_3(x).
\end{equation*}

\paragraph{(9.Sm2)}
One has $|M_3(x)|_2=3(1+\rho)\le6$ by (I.m6), and $|M_3(x)|_c=3(1+\rho)\le6$ by (I.m7).
Moreover $\left|w_3\,\gamma\cdot\hat{d}_3(x)\right|_2=6$:
indeed $\hat{d}_3(x)=ge_1$ is a real unit vector, $\gamma\cdot v$ is Hermitian with $(\gamma\cdot v)^2=|v|_2^2$ for real $v$ (I.c1), so that $|\gamma\cdot v|_2=|v|_2$, and $w_3=6$.

\paragraph{(9.Sm3)}
One has $\overline{M_i(x)}=\varkappa^{-1}M_i(x)\,\varkappa$ by (I.m3), and $M_i(\tilde{g}x)\,S=S\,M_i(x)$ for $\tilde{g}=(S,a)\in\mathbb{G}_k$ by (I.m4).

\paragraph{(9.Sm4)}
For every $S\in\mathrm{Sp}$ there is a sign $\mathrm{sn}(S)\in\{1,-1\}$ with
\begin{equation*}
\left(S^{-1}\right)^{T}=\mathrm{sn}(S)\cdot\mathsf{C}^{-1}S\,\mathsf{C}.
\end{equation*}
\begin{proof}
The maps $S\mapsto\left(S^{-1}\right)^{T}$ and $S\mapsto\mathsf{C}^{-1}S\,\mathsf{C}$ are homomorphisms of $\mathrm{Sp}$ into $\mathrm{GL}(4,\mathbb{C})$.
On the generators one has $\gamma_\mu^{-1}=\gamma_\mu$ and $\gamma_5^{-1}=\gamma_5$ by (I.c1), and $\gamma_\mu^{T}=-\mathsf{C}^{-1}\gamma_\mu\mathsf{C}$, $\gamma_5^{T}=-\mathsf{C}^{-1}\gamma_5\mathsf{C}$ by (I.c3).
Both homomorphisms map $-1$ to $-1$.
Hence $\mathrm{sn}(S)=(-1)^n$ if $S$ is $+$ or $-$ a product of $n$ generators.
\end{proof}

\paragraph{(9.Sm5)}
One has $\mathsf{C}^{-1}M_i(x)\,\mathsf{C}=M_i(x^{(i)})^{T}$.
\begin{proof}
The transpose of (I.m2), $\mathsf{C}^{-1}M_i(x)^{T}\mathsf{C}=M_i(x^{(i)})$, is $\mathsf{C}^{T}M_i(x)\left(\mathsf{C}^{-1}\right)^{T}=M_i(x^{(i)})^{T}$, and $\mathsf{C}^{T}=\mathsf{C}^{-1}$, $\left(\mathsf{C}^{-1}\right)^{T}=\mathsf{C}$ by (I.c3).
\end{proof}

\subsubsection*{1.5. Norms}

\paragraph{(9.No1)}
On $\mathbb{C}^4$ one has $|v|_c\le|v|_2\le2|v|_c$ (Section~I.1.5); hence $|f|_\infty\le|f|_{\mathrm{su}}\le2|f|_\infty$ (Section~G1.3).

\paragraph{(9.No2)}
Let $O$ be a linear operator from $\mathfrak{S}^{\mathrm{ch}}(\eta_k,\ell)$, or from the functions of pairs, to $\mathfrak{S}(1,\ell)$, and let $\mathrm{Cn}$ be a number such that $|(Of)(y)|_2\le\mathrm{Cn}\cdot|f|_{\mathrm{su}}$ for all $f$ and all unit sites $y$.
Then $|O|_{1,\infty}\le2\,\mathrm{Cn}$.
\begin{proof}
Fix a unit site $y$ and $\beta\in\{1,\dots,4\}$, and fix a set $\mathrm{Rep}$ of representatives modulo $\ell\mathbb{Z}^3$ of the indices $t$ of the source (fermion sites, or pairs of $\mathrm{Pzt}_k$, in which case both points are shifted).
Then
\begin{equation*}
(Of)(y)_\beta=\sum_{t\in\mathrm{Rep}}\ \sum_{\beta'}\mathrm{coef}(t,\beta')\cdot f(t)_{\beta'}
\end{equation*}
with complex numbers $\mathrm{coef}(t,\beta')$.
The sum that defines $|O|_{1,\infty}$ in the row $(y,\beta)$ is $\sum_{t,\beta'}|\mathrm{coef}(t,\beta')|$.
Indeed, for the source $\mathfrak{S}^{\mathrm{ch}}(\eta_k,\ell)$, (I.1.4) and (I.1.3) give $\mathrm{coef}(x',\beta')=\frac{\eta_k^3}{48}\,O(y,x')_{\beta\beta'}$, so that this sum equals $\sum_{\beta'}\int\mathrm{d}x'\,\left|O(y,x')_{\beta\beta'}\right|$, as in (I.ke3);
for the functions of pairs it is the definition in Section~I.5.2.
For $t\in\mathrm{Rep}$ put
\begin{equation*}
f(t)_{\beta'}:=
\begin{cases}
\dfrac{\overline{\mathrm{coef}(t,\beta')}}{|\mathrm{coef}(t,\beta')|} & \text{if }\mathrm{coef}(t,\beta')\ne0,\\[1ex]
0 & \text{otherwise},
\end{cases}
\end{equation*}
and extend $f$ antiperiodically.
Then $|f|_\infty\le1$, hence $|f|_{\mathrm{su}}\le2$ by (9.No1), and
\begin{equation*}
\sum_{t,\beta'}|\mathrm{coef}(t,\beta')|=(Of)(y)_\beta\le|(Of)(y)|_2\le2\,\mathrm{Cn}.
\end{equation*}
\end{proof}

\paragraph{(9.No3)}
Let $\mathrm{mo}$ be a monomial $4\times4$ matrix, that is, a matrix with one non-zero entry, of modulus one, in each row and in each column.
For $O$ as in (9.No2),
\begin{equation*}
\left|\mathrm{mo}^{-1}\cdot O\cdot\mathrm{mo}\right|_{1,\infty}=|O|_{1,\infty}.
\end{equation*}
\begin{proof}
For a unit site $y$ let $\mathrm{coef}(t)$ be the $4\times4$ matrix of the coefficients of $f(t)$ in $(Of)(y)$, that is, the matrix whose row $\beta$ consists of the numbers $\mathrm{coef}(t,\beta')$ of (9.No2).
The corresponding matrix of $\mathrm{mo}^{-1}\cdot O\cdot\mathrm{mo}$ is $\mathrm{mo}^{-1}\,\mathrm{coef}(t)\,\mathrm{mo}$, and $\mathrm{mo}^{-1}$ is monomial as well.
Multiplication by monomial matrices from the left and from the right permutes the rows and the columns of $\mathrm{coef}(t)$, by permutations that do not depend on $t$, and multiplies the entries by numbers of modulus one.
Hence the four row sums $\sum_{t,\beta'}\left|\mathrm{coef}(t)_{\beta\beta'}\right|$, $\beta=1,\dots,4$, are permuted.
\end{proof}

\subsection*{2. Proof of (a)}

\begin{proof}[Proof of Proposition~9\,(a)]
Let $k\ge1$ and $A\in\tilde{G}_k$.
By Section~I.3.5, $\boldsymbol{\alpha}_\infty=b_k$ and $\boldsymbol{\beta}_\infty=0$, so that $F_\infty(A)=D_A+\mu+b_k\cdot JP_k(A)J^{*}$.
By Proposition~2.8\,(a), applied at the element $\infty$ of $Y_k$, the operator $F_\infty(A)$ is bijective on $\mathfrak{S}^{\mathrm{ch}}(\eta_k,\ell)$, and $S_k(A)=F_\infty(A)^{-1}$.
The definitions of Section~I.3.5 therefore give
\begin{equation*}
T_k(A)\cdot H_k(A)=b_k^{-1}\cdot Q_k(A)J^{*}\cdot F_\infty(A)\cdot b_k\cdot F_\infty(A)^{-1}\cdot JQ_k^{T}(-A)=\left(Q_k(A)J^{*}\right)\left(JQ_k^{T}(-A)\right)=1
\qquad\text{on }\mathfrak{S}(1,\ell),
\end{equation*}
where the last equality is (I.Q1).
By the property (G2) of \cite[Section~16.2.7]{QEDI}, $-A$ lies in $\tilde{G}_k$;
so the relation also holds at $-A$, and with $\bar{T}_k(A)=\mathsf{C}^{-1}T_k(-A)\,\mathsf{C}$ and $\bar{H}_k(A)=\mathsf{C}^{-1}H_k(-A)\,\mathsf{C}$ (Section~I.3.5) one obtains
\begin{equation*}
\bar{T}_k(A)\cdot\bar{H}_k(A)=\mathsf{C}^{-1}T_k(-A)\,\mathsf{C}\cdot\mathsf{C}^{-1}H_k(-A)\,\mathsf{C}=\mathsf{C}^{-1}\cdot T_k(-A)\cdot H_k(-A)\cdot\mathsf{C}=\mathsf{C}^{-1}\mathsf{C}=1
\qquad\text{on }\mathfrak{S}(1,\ell).
\end{equation*}
The case $k=0$ is treated in Section~7.
\end{proof}

\subsection*{3. Proof of (b)}

\begin{proof}[Proof of Proposition~9\,(b)]
Let $k\ge1$ and let $A$ be a connection.

\emph{The display of $T_k(A)$.}
By Section~I.3.3, $P_k(A)=Q_k^{T}(-A)\,Q_k(A)$.
By (I.Q1),
\begin{equation*}
Q_k(A)J^{*}\cdot JP_k(A)J^{*}=\left(Q_k(A)J^{*}\cdot JQ_k^{T}(-A)\right)Q_k(A)J^{*}=Q_k(A)J^{*}.
\end{equation*}
Hence
\begin{equation*}
T_k(A)=b_k^{-1}\cdot Q_k(A)J^{*}\left(D_A+\mu+b_k\cdot JP_k(A)J^{*}\right)=Q_k(A)J^{*}\left[b_k^{-1}D_A+\frac{\mu}{b_k}+1\right],
\end{equation*}
which is the second form of $T_k(A)$ in Section~I.3.5.
By Section~I.3.3, $\left(Q_k(A)J^{*}f\right)(y)=\int_{\Delta_y}\mathrm{d}x\,\varphi_A(y,c)\cdot f(x)$.
Inserting $f=b_k^{-1}D_A\psi+\left(1+\frac{\mu}{b_k}\right)\psi$ and the display (I.2.3) of $D_A$, one finds
\begin{equation}\tag{9.8}
\begin{split}
\left(T_k(A)\psi\right)(y)=\int_{\Delta_y}\mathrm{d}x\,\varphi_A(y,c)\biggl\{&\left(1+\frac{\mu}{b_k}\right)\psi(x)\\
&-(b_k\cdot\eta_k)^{-1}\sum_{i=0}^{3}M_i(x)\left[U_A(x,x^{(i)})\cdot\psi(x^{(i)})-\psi(x)\right]\biggr\},
\end{split}
\end{equation}
where $U_A(x,x^{(i)})=1$ for $i\le2$ by (9.4).
This is the first display of (b).

\emph{The two pieces.}
Let $\mathbf{Y}$ be a union of $M$-cubes (a polymer is such a union: Section~I.1.1), and let $y$ be a unit site of $\mathbf{Y}$.
By (9.B4), $\Delta_y\subset\mathbf{Y}$.
Let $x=(c,g)$ lie over $\Delta_y$.
For $i\le2$ the site $x^{(i)}=(c,gs_i)$ lies over $c$ by (I.1.2), hence over $\mathbf{Y}$.
Thus in (I.5.4) the condition that $x^{(i)}$ lie over $\mathbf{Y}$ is satisfied for $i=0,1,2$, while for $i=3$ it reads $c_3\in\mathbf{Y}$.
So
\begin{equation}\tag{9.9}
\begin{split}
\left(T^{(1)}(A)\psi\right)(y)=\int_{\Delta_y}\mathrm{d}x\,\varphi_A(y,c)\biggl\{&\left(1+\frac{\mu}{b_k}\right)\psi(x)-(b_k\cdot\eta_k)^{-1}\biggl[\sum_{i\le2}M_i(x)\left(\psi(x^{(i)})-\psi(x)\right)\\
&+\mathbf{1}[c_3\in\mathbf{Y}]\cdot M_3(x)\left(u_A(c,c_3)\cdot\psi(x^{(3)})-\psi(x)\right)\biggr]\biggr\},
\end{split}
\end{equation}
\begin{equation}\tag{9.10}
\begin{split}
&\left(T_k(A)\psi\right)(y)-\left(T^{(1)}(A)\psi\right)(y)\\
&\qquad=-(b_k\cdot\eta_k)^{-1}\int_{\Delta_y}\mathrm{d}x\,\mathbf{1}[c_3\notin\mathbf{Y}]\cdot\varphi_A(y,c)\cdot M_3(x)\left(u_A(c,c_3)\cdot\psi(x^{(3)})-\psi(x)\right).
\end{split}
\end{equation}
Since $c_3=c+\eta_k\,ge_1\ne c$, the pair $(x,x^{(3)})$ is an ordered pair of distinct points of $\mathbb{L}^{\mathrm{f}}_k$ with $d_X(x,x^{(3)})=\max\left(|c-c_3|,\eta_k\right)=\eta_k\le M$, that is, an element of $\mathrm{Pzt}_k$;
and by (I.1.7) and (I.1.6),
\begin{equation*}
\eta_k^{\alpha}\cdot\left(\delta_{\alpha,A}\psi\right)(x,x^{(3)})=u_A(c,c_3)\cdot\psi(x^{(3)})-\psi(x).
\end{equation*}
Suppose that $\chi(x,x^{(3)})=\left(\delta_{\alpha,A}\psi\right)(x,x^{(3)})$ for the sites $x$ over $\Delta_y$ with $c_3\notin\mathbf{Y}$.
Then the right side of (9.10) equals
\begin{equation*}
-b_k^{-1}\cdot\eta_k^{\alpha-1}\int_{\Delta_y}\mathrm{d}x\,\mathbf{1}[c_3\notin\mathbf{Y}]\cdot\varphi_A(y,c)\cdot M_3(x)\,\chi(x,x^{(3)})=\left(T^{(2)}(A)\chi\right)(y)
\end{equation*}
by (I.5.4).
This is the identity $T^{(1)}(A)\psi+T^{(2)}(A)\chi=T_k(A)\psi$ at $y$.

\emph{What the rows contain.}
By (9.9), $\left(T^{(1)}(A)\psi\right)(y)$ contains $\psi$ at the sites $x$ and $x^{(i)}$ ($i\le2$) over $\Delta_y$ and at the sites $x^{(3)}$ with $x$ over $\Delta_y$ and $c_3\in\mathbf{Y}$.
As $\Delta_y\subset\mathbf{Y}$, all these sites lie over $\mathbf{Y}$.
The differences $\psi(x^{(i)})-\psi(x)$, $i\le2$, which belong to the pairs $(x,x^{(i)})$ inside the cluster of $c$, appear in $T^{(1)}(A)$, written with $\psi$, for every $\mathbf{Y}$;
the operator $T^{(2)}(A)$ contains no pair inside a cluster.
By (I.5.4), $\left(T^{(2)}(A)\chi\right)(y)$ contains $\chi$ only at the pairs $(x,x^{(3)})$ with $x$ over $\Delta_y$ and $c_3\notin\mathbf{Y}$.
The second point of such a pair is $x^{(3)}$, and $c\in\mathbf{Y}$;
so the pair is of kind (II) for $\mathbf{Y}$ (Section~I.1.7), hence admissible for $\mathbf{Y}$, and $d_X(x,x^{(3)})=\eta_k$.

\emph{What the rows read.}
The coefficients in (9.9) and in (I.5.4) are products of real numbers, of the matrices $M_i(x)$, of the numbers $\varphi_A(y,c)$, $c\in\Delta_y$, and, in (9.9), of the numbers $u_A(c,c_3)$, $c\in\Delta_y$.
By (9.P1), $\varphi_A(y,c)$ reads $A$ on bonds contained in $\Delta_y$;
and $u_A(c,c_3)$ reads $A$ on the bond $(c,c_3)$, whose end point $c$ lies in $\Delta_y$.
\end{proof}

\subsection*{4. Proof of (c)}

Let $k\ge1$, let $A$ be a connection, $\mathbf{Y}$ a union of $M$-cubes and $y$ a unit site of $\mathbf{Y}$, and put $\mathfrak{i}:=\mathfrak{i}(A)$ and $\mathrm{qd}:=e_k\cdot|dA|_\infty$.
The whole torus $T_k$ is a union of $M$-cubes, since $\ell$ is a multiple of $M$.
For $\mathbf{Y}=T_k$ the indicator in (9.9) equals $1$, and (9.8) gives
\begin{equation}\tag{9.11}
T_k(A)=T^{(1)}_{\mathbf{Y}}(A)\qquad\text{for }\mathbf{Y}=T_k.
\end{equation}
It therefore suffices to bound $T^{(1)}(A)$ and $T^{(2)}(A)$ for every $\mathbf{Y}$.

\subsubsection*{4.1. The sets}

Throughout 4.1--4.5 the unit site $y$ of $\mathbf{Y}$ is fixed.
A fermion site is written $x=(c,g)$, with $c=\pi(x)$ and $c_3=c_3(x)$.
Put
\begin{gather*}
\mathrm{Ef}:=\Delta_y\times O_{\mathrm{cub}},\qquad
\mathrm{Ef}_1:=\left\{x\in\mathrm{Ef}:c_3\in\Delta_y\right\},\\
\mathrm{Ef}_2:=\left\{x\in\mathrm{Ef}:c_3\in\mathbf{Y},\ c_3\notin\Delta_y\right\},\qquad
\mathrm{Ef}_3:=\left\{x\in\mathrm{Ef}:c_3\notin\mathbf{Y}\right\}.
\end{gather*}
Since $\Delta_y$ is contained in $\mathbf{Y}$ (9.B4), the set $\mathrm{Ef}$ is the disjoint union of $\mathrm{Ef}_1$, $\mathrm{Ef}_2$ and $\mathrm{Ef}_3$.
Moreover $\#|\mathrm{Ef}|=48\,\eta_k^{-3}$.
These sets have the following three properties.
\begin{itemize}
\item[(i)] For $i\le2$ the map $x\mapsto x^{(i)}$ is a bijection of $\mathrm{Ef}$ onto itself which leaves $c$ unchanged: $\pi(x^{(i)})=c$ by (I.1.2), and $(x^{(i)})^{(i)}=x$ by (I.s1).
\item[(ii)] The map $x\mapsto x^{(3)}$ is a bijection of $\mathrm{Ef}_1$ onto itself, with $\pi(x^{(3)})=c_3(x)$ and $c_3(x^{(3)})=c$.
Indeed, for $x\in\mathrm{Ef}_1$ the site $x^{(3)}$ lies over $c_3\in\Delta_y$, and $(x^{(3)})^{(3)}=x$ (I.s1) lies over $c\in\Delta_y$.
\item[(iii)] Each $g\in O_{\mathrm{cub}}$ maps $e_1$ to a direction $a\in\{\pm e_\mu:\mu=1,2,3\}$.
By (I.s4), for every $c$ and every direction $a$ there are exactly $8$ labels $g$ with $g\,e_1=a$; in other words, exactly $8$ sites $x$ over $c$ satisfy $c_3(x)=c+\eta_k a$ and $\hat{d}_3(x)=a$.
Together with (9.B5) this gives
\begin{equation}\tag{9.12}
\begin{aligned}
&\mathrm{Ef}_2\cup\mathrm{Ef}_3=\mathrm{Ef}\setminus\mathrm{Ef}_1=\bigsqcup_{a}\left\{(c,g):c\in\mathrm{Fa}_a(y),\ g\,e_1=a\right\},\\
&\#|\mathrm{Ef}_2\cup\mathrm{Ef}_3|=6\cdot\eta_k^{-2}\cdot8=48\,\eta_k^{-2},
\end{aligned}
\end{equation}
where $a$ ranges over the six directions and the union is disjoint.
\end{itemize}

\subsubsection*{4.2. The row after summation by parts}

\begin{named}{Lemma 9.A}
For every $\psi\in\mathfrak{S}^{\mathrm{ch}}(\eta_k,\ell)$ one has $\left(T^{(1)}(A)\psi\right)(y)=\mathrm{Rw}_0+\mathrm{Rw}_1+\mathrm{Rw}_2+\mathrm{Rw}_3$, where
\begin{equation}\tag{9.13}
\begin{aligned}
\mathrm{Rw}_0&:=\left(1+\frac{\mu}{b_k}\right)\frac{\eta_k^3}{48}\sum_{x\in\mathrm{Ef}}\varphi_A(y,c)\,\psi(x),\\
\mathrm{Rw}_1&:=b_k^{-1}\cdot\frac{\eta_k^2}{48}\sum_{x\in\mathrm{Ef}_2\cup\mathrm{Ef}_3}\varphi_A(y,c)\,w_3\,\gamma\cdot\hat{d}_3(x)\,\psi(x),\\
\mathrm{Rw}_2&:=-b_k^{-1}\cdot\frac{\eta_k^2}{48}\sum_{x\in\mathrm{Ef}_2}\varphi_A(y,c)\,M_3(x)\left[u_A(c,c_3)\,\psi(x^{(3)})-\psi(x)\right],\\
\mathrm{Rw}_3&:=-b_k^{-1}\cdot\frac{\eta_k^2}{48}\sum_{x\in\mathrm{Ef}_1}\left[\varphi_A(y,c_3)\,u_A(c_3,c)-\varphi_A(y,c)\right]M_3(x^{(3)})\,\psi(x).
\end{aligned}
\end{equation}
\end{named}

\begin{proof}
From (9.9) and (9.1), and since $(b_k\eta_k)^{-1}\cdot\frac{\eta_k^3}{48}=b_k^{-1}\cdot\frac{\eta_k^2}{48}$, one obtains
\begin{gather*}
\left(T^{(1)}(A)\psi\right)(y)=\mathrm{Rw}_0-b_k^{-1}\cdot\frac{\eta_k^2}{48}\left(\mathrm{Su}_{\mathrm{in}}+\mathrm{Su}_1+\mathrm{Su}_2\right),\\
\mathrm{Su}_{\mathrm{in}}:=\sum_{x\in\mathrm{Ef}}\sum_{i\le2}\varphi_A(y,c)\,M_i(x)\left[\psi(x^{(i)})-\psi(x)\right],\\
\mathrm{Su}_n:=\sum_{x\in\mathrm{Ef}_n}\varphi_A(y,c)\,M_3(x)\left[u_A(c,c_3)\,\psi(x^{(3)})-\psi(x)\right]\qquad(n=1,2).
\end{gather*}

\emph{The sum $\mathrm{Su}_{\mathrm{in}}$ (the $3\cdot48$ differences inside each cluster).}
Fix $i\le2$ and substitute $x\mapsto x^{(i)}$ in $\sum_{x\in\mathrm{Ef}}\varphi_A(y,c)\,M_i(x)\,\psi(x^{(i)})$.
By 4.1\,(i) the new summation variable again ranges over $\mathrm{Ef}$ and the photon site $c$ does not change, while $M_i(x)$ becomes $M_i(x^{(i)})$ and $\psi(x^{(i)})$ becomes $\psi(x)$.
So this sum equals $\sum_{x\in\mathrm{Ef}}\varphi_A(y,c)\,M_i(x^{(i)})\,\psi(x)$, and (9.Sm1) yields
\begin{equation}\tag{9.14}
\begin{aligned}
\mathrm{Su}_{\mathrm{in}}&=\sum_{x\in\mathrm{Ef}}\varphi_A(y,c)\sum_{i\le2}\left(M_i(x^{(i)})-M_i(x)\right)\psi(x)
=\sum_{x\in\mathrm{Ef}}\varphi_A(y,c)\sum_{i\le2}w_i\,\gamma\cdot\hat{d}_i(x)\,\psi(x)\\
&=-\sum_{x\in\mathrm{Ef}}\varphi_A(y,c)\,w_3\,\gamma\cdot\hat{d}_3(x)\,\psi(x).
\end{aligned}
\end{equation}

\emph{The sum $\mathrm{Su}_1$.}
Substitute $x\mapsto x^{(3)}$ in $\sum_{x\in\mathrm{Ef}_1}\varphi_A(y,c)\,M_3(x)\,u_A(c,c_3)\,\psi(x^{(3)})$.
By 4.1\,(ii) the new summation variable again ranges over $\mathrm{Ef}_1$; the site $c$ becomes $c_3$, the site $c_3$ becomes $c$, $M_3(x)$ becomes $M_3(x^{(3)})$, and $\psi(x^{(3)})$ becomes $\psi(x)$.
Hence
\begin{equation}\tag{9.15}
\begin{aligned}
\mathrm{Su}_1&=\sum_{x\in\mathrm{Ef}_1}\left\{\varphi_A(y,c_3)\,u_A(c_3,c)\,M_3(x^{(3)})-\varphi_A(y,c)\,M_3(x)\right\}\psi(x)\\
&=\sum_{x\in\mathrm{Ef}_1}\left[\varphi_A(y,c_3)\,u_A(c_3,c)-\varphi_A(y,c)\right]M_3(x^{(3)})\,\psi(x)\\
&\qquad+\sum_{x\in\mathrm{Ef}_1}\varphi_A(y,c)\left(M_3(x^{(3)})-M_3(x)\right)\psi(x),
\end{aligned}
\end{equation}
and $M_3(x^{(3)})-M_3(x)=w_3\,\gamma\cdot\hat{d}_3(x)$ by (9.Sm1).
In $\mathrm{Su}_{\mathrm{in}}+\mathrm{Su}_1$ the parts over $\mathrm{Ef}_1$ of the sums of $\varphi_A(y,c)\,w_3\,\gamma\cdot\hat{d}_3(x)\,\psi(x)$ cancel, and
\begin{align*}
\mathrm{Su}_{\mathrm{in}}+\mathrm{Su}_1=&-\sum_{x\in\mathrm{Ef}\setminus\mathrm{Ef}_1}\varphi_A(y,c)\,w_3\,\gamma\cdot\hat{d}_3(x)\,\psi(x)\\
&+\sum_{x\in\mathrm{Ef}_1}\left[\varphi_A(y,c_3)\,u_A(c_3,c)-\varphi_A(y,c)\right]M_3(x^{(3)})\,\psi(x).
\end{align*}
Multiplication by $-b_k^{-1}\cdot\frac{\eta_k^2}{48}$ turns the right side into $\mathrm{Rw}_1+\mathrm{Rw}_3$, because $\mathrm{Ef}\setminus\mathrm{Ef}_1=\mathrm{Ef}_2\cup\mathrm{Ef}_3$.
Finally $-b_k^{-1}\cdot\frac{\eta_k^2}{48}\,\mathrm{Su}_2=\mathrm{Rw}_2$.
\end{proof}

\subsubsection*{4.3. The terms $\mathrm{Rw}_0$, $\mathrm{Rw}_1$, $\mathrm{Rw}_2$}

By (9.P2) one has $|\varphi_A(y,c)|\le\exp\left(\frac{3}{2}\mathfrak{i}\right)$ and $|u_A(c,c_3)|\le\exp(\mathfrak{i})$.
Consequently $\left|u_A(c,c_3)\,\psi(x^{(3)})-\psi(x)\right|_2\le\left(\exp(\mathfrak{i})+1\right)|\psi|_{\mathrm{su}}$.
Using in addition (9.Sm2), (9.12), $\#|\mathrm{Ef}|=48\,\eta_k^{-3}$ and $\#|\mathrm{Ef}_2|\le48\,\eta_k^{-2}$, one finds
\begin{equation}\tag{9.16}
\begin{aligned}
|\mathrm{Rw}_0|_2&\le\left(1+\frac{\mu}{b_k}\right)\frac{\eta_k^3}{48}\cdot48\,\eta_k^{-3}\exp\left(\frac{3}{2}\mathfrak{i}\right)|\psi|_{\mathrm{su}}
=\left(1+\frac{\mu}{b_k}\right)\exp\left(\frac{3}{2}\mathfrak{i}\right)|\psi|_{\mathrm{su}},\\
|\mathrm{Rw}_1|_2&\le b_k^{-1}\cdot\frac{\eta_k^2}{48}\cdot48\,\eta_k^{-2}\exp\left(\frac{3}{2}\mathfrak{i}\right)\cdot6\,|\psi|_{\mathrm{su}}
=6\,b_k^{-1}\exp\left(\frac{3}{2}\mathfrak{i}\right)|\psi|_{\mathrm{su}},\\
|\mathrm{Rw}_2|_2&\le b_k^{-1}\cdot\frac{\eta_k^2}{48}\cdot48\,\eta_k^{-2}\exp\left(\frac{3}{2}\mathfrak{i}\right)\cdot6\left(\exp(\mathfrak{i})+1\right)|\psi|_{\mathrm{su}}\\
&=6\left(\exp(\mathfrak{i})+1\right)b_k^{-1}\exp\left(\frac{3}{2}\mathfrak{i}\right)|\psi|_{\mathrm{su}}.
\end{aligned}
\end{equation}

\subsubsection*{4.4. The term $\mathrm{Rw}_3$: the defect of the phases}

Let $a$ be a direction, and let $c$ and $c':=c+\eta_k a$ be two sites of $\Delta_y$.
By (9.4),
\begin{gather*}
\varphi_A(y,c')\,u_A(c',c)-\varphi_A(y,c)=\varphi_A(y,c)\left(\exp\left(\mathrm{vz}(c,a)\right)-1\right),\\
\mathrm{vz}(c,a):=i\,e_k\cdot\eta_k\left[(\tau_kA)(y,c')+A(c',c)-(\tau_kA)(y,c)\right].
\end{gather*}

\emph{Step 1 (the bracket is a mean over closed paths).}
Put $j:=\mathrm{jm}(c,a)$ (9.B6).
By (9.B3)\,(i), $\mathrm{ce}_i(c')=\mathrm{ce}_i(c)$ for $j\le i\le k$.
Hence in $(\tau_kA)(y,c')-(\tau_kA)(y,c)$ the terms of the sums (9.4) whose summation index is at least $j$ cancel, and
\begin{equation*}
(\tau_kA)(y,c')+A(c',c)-(\tau_kA)(y,c)=\sum_{i<j}(\tau A)\left(\mathrm{ce}_{i+1}(c'),\mathrm{ce}_i(c')\right)+A(c',c)-\sum_{i<j}(\tau A)\left(\mathrm{ce}_{i+1}(c),\mathrm{ce}_i(c)\right).
\end{equation*}
Let $\mathrm{od}_0,\dots,\mathrm{od}_{j-1}$ and $\mathrm{od}'_0,\dots,\mathrm{od}'_{j-1}$ be orders of the three coordinates.
Let $\boldsymbol{\ell}$ be the lattice path which traverses first the paths
\begin{equation*}
\boldsymbol{\Gamma}^{\mathrm{od}'_{j-1}}\left(\mathrm{ce}_j(c'),\mathrm{ce}_{j-1}(c')\right),\ \dots,\ \boldsymbol{\Gamma}^{\mathrm{od}'_0}\left(\mathrm{ce}_1(c'),c'\right),
\end{equation*}
then the bond $(c',c)$, and then the reversed paths of
\begin{equation*}
\boldsymbol{\Gamma}^{\mathrm{od}_0}\left(\mathrm{ce}_1(c),c\right),\ \dots,\ \boldsymbol{\Gamma}^{\mathrm{od}_{j-1}}\left(\mathrm{ce}_j(c),\mathrm{ce}_{j-1}(c)\right).
\end{equation*}
It begins at $\mathrm{ce}_j(c')$ and ends at $\mathrm{ce}_j(c)=\mathrm{ce}_j(c')$, so it is closed.
As $A$ is odd under the reversal of bonds,
\begin{equation*}
A(\boldsymbol{\ell})=\sum_{i<j}A\left(\boldsymbol{\Gamma}^{\mathrm{od}'_i}\left(\mathrm{ce}_{i+1}(c'),\mathrm{ce}_i(c')\right)\right)+A(c',c)-\sum_{i<j}A\left(\boldsymbol{\Gamma}^{\mathrm{od}_i}\left(\mathrm{ce}_{i+1}(c),\mathrm{ce}_i(c)\right)\right).
\end{equation*}
Each term on the right depends on at most one of the $2j$ orders.
Moreover $(\tau A)(\tilde{c},\tilde{c}')$ is the mean of $A\left(\boldsymbol{\Gamma}^{\pi}(\tilde{c},\tilde{c}')\right)$ over the six orders $\pi$; for $\tilde{c}=\tilde{c}'$ the path is empty and $(\tau A)(\tilde{c},\tilde{c})=0$ (Section~I.1.1).
Therefore the bracket in the definition of $\mathrm{vz}(c,a)$ is the mean of $A(\boldsymbol{\ell})$ over the $6^{2j}$ choices of the orders:
\begin{equation*}
\mathrm{vz}(c,a)=i\,e_k\cdot\eta_k\cdot6^{-2j}\sum_{\mathrm{od},\,\mathrm{od}'}A(\boldsymbol{\ell}),
\end{equation*}
the sum extending over all choices of $\mathrm{od}_0,\dots,\mathrm{od}_{j-1},\mathrm{od}'_0,\dots,\mathrm{od}'_{j-1}$.

\emph{Step 2 (one closed path).}
Put $\mathbb{W}:=\mathrm{Bl}_j\left(\mathrm{ce}_j(c)\right)$.
By (9.B1), $\mathbb{W}$ is a product of three intervals of $\eta_k\mathbb{Z}$, each of length $(L^j-1)\cdot\eta_k\le s_W:=L^j\cdot\eta_k$.
By (9.B3), $\mathbb{W}$ contains $c$, $c'$ and all the centres $\mathrm{ce}_i(c)$, $\mathrm{ce}_i(c')$, $i\le j$; by (9.P1) it contains all the paths of which $\boldsymbol{\ell}$ is composed, and the bond $(c',c)$.
By (9.P1) the number of bonds of $\boldsymbol{\ell}$ is at most
\begin{equation*}
2\sum_{i<j}\frac{3\left(L^{i+1}-L^i\right)}{2}+1=3\left(L^j-1\right)+1\le3L^j.
\end{equation*}
The statement of \cite[Lemma~16.26]{QEDI} concerns a complex antisymmetric function on the bonds of a box.
It is used here with the coordinates $1,2,3$ in the places of the coordinates $0,1,2$ of that lemma, so that its rectilinear paths change the coordinates in the order $1,2,3$; its $(dA)(p)$, which is $\eta_k^{-1}$ times the sum of $A$ over the boundary of the plaquette $p$, is $(dA)_{\mu\nu}(c)$ of Section~I.1.1.
Applied to the restriction of $A$ to the bonds of $\mathbb{W}$ and to the closed path $\boldsymbol{\ell}$, which has at most $3L^j$ bonds, it gives
\begin{equation*}
|A(\boldsymbol{\ell})|\le2\cdot3L^j\cdot L^j\cdot\eta_k\cdot|dA|_\infty=6L^{2j}\cdot\eta_k\cdot|dA|_\infty.
\end{equation*}
Since $\eta_k^2=L^{-2k}$ and $j\le k$, and since $|\exp(t)-1|\le|t|\exp|t|$ for every complex number $t$, it follows that
\begin{equation}\tag{9.17}
\begin{aligned}
&|\mathrm{vz}(c,a)|\le e_k\cdot\eta_k\cdot6L^{2j}\cdot\eta_k\cdot|dA|_\infty=6L^{2(j-k)}\cdot\mathrm{qd}\le6\,\mathrm{qd},\\
&\left|\exp\left(\mathrm{vz}(c,a)\right)-1\right|\le6L^{2(j-k)}\cdot\mathrm{qd}\cdot\exp(6\,\mathrm{qd})\qquad\left(j=\mathrm{jm}(c,a)\right).
\end{aligned}
\end{equation}

\emph{Step 3 (summation over $c$).}
By (9.17), (9.3) and $\eta_k=L^{-k}$, every direction $a$ satisfies
\begin{equation}\tag{9.18}
\begin{aligned}
\eta_k^2\sum_{\substack{c\in\Delta_y\\c+\eta_ka\in\Delta_y}}\left|\exp\left(\mathrm{vz}(c,a)\right)-1\right|
&\le\sum_{j=1}^{k}\eta_k^2\cdot\eta_k^{-3}\cdot L^{-(j-1)}\cdot6L^{2(j-k)}\cdot\mathrm{qd}\cdot\exp(6\,\mathrm{qd})\\
&=6\,\mathrm{qd}\cdot\exp(6\,\mathrm{qd})\cdot L^{1-k}\sum_{j=1}^{k}L^j\le9L\cdot\mathrm{qd}\cdot\exp(6\,\mathrm{qd}),
\end{aligned}
\end{equation}
because $L^k\cdot L^{1-j}\cdot L^{2j-2k}=L^{1-k}\cdot L^j$, $\sum_{j=1}^{k}L^j\le\frac{L^{k+1}}{L-1}$, and $\frac{L^2}{L-1}\le\frac{3L}{2}$ for $L\ge3$.

\emph{Step 4.}
By 4.1\,(iii), $\mathrm{Ef}_1$ is the disjoint union, over the six directions $a$, of the sets $\left\{(c,g):c\in\Delta_y,\ c+\eta_ka\in\Delta_y,\ g\,e_1=a\right\}$, with $8$ labels for each pair $(c,a)$; for such an $x$ one has $c_3=c+\eta_ka$.
In the following display the sums over $c$ extend over the sites $c\in\Delta_y$ with $c+\eta_ka\in\Delta_y$.
By (9.P2), (9.Sm2) and (9.18),
\begin{equation}\tag{9.19}
\begin{aligned}
|\mathrm{Rw}_3|_2&\le b_k^{-1}\cdot\frac{1}{48}\sum_a8\cdot\eta_k^2\sum_c\exp\left(\frac{3}{2}\mathfrak{i}\right)\left|\exp\left(\mathrm{vz}(c,a)\right)-1\right|\cdot6\,|\psi|_{\mathrm{su}}\\
&\le b_k^{-1}\cdot\frac{8\cdot6}{48}\cdot6\cdot9L\cdot\mathrm{qd}\cdot\exp(6\,\mathrm{qd})\exp\left(\frac{3}{2}\mathfrak{i}\right)|\psi|_{\mathrm{su}}\\
&=54L\cdot\mathrm{qd}\cdot\exp(6\,\mathrm{qd})\,b_k^{-1}\exp\left(\frac{3}{2}\mathfrak{i}\right)|\psi|_{\mathrm{su}}.
\end{aligned}
\end{equation}

\subsubsection*{4.5. Conclusion}

Lemma 9.A, (9.16), (9.19) and $b_k\ge b_{-}$ (Section~I.3.4) show that at every unit site $y$ of $\mathbf{Y}$
\begin{align*}
\left|\left(T^{(1)}(A)\psi\right)(y)\right|_2&\le\exp\left(\frac{3}{2}\mathfrak{i}\right)\left[1+\frac{\mu}{b_k}+b_k^{-1}\left(6+6\exp(\mathfrak{i})+6+54L\cdot\mathrm{qd}\cdot\exp(6\,\mathrm{qd})\right)\right]|\psi|_{\mathrm{su}}\\
&\le\exp\left(\frac{3}{2}\mathfrak{i}\right)\left[b_{-}^{-1}\left(54L\cdot\mathrm{qd}\cdot\exp(6\,\mathrm{qd})+12+6\exp(\mathfrak{i})+\mu\right)+1\right]|\psi|_{\mathrm{su}}.
\end{align*}
This is the asserted bound on $\|T^{(1)}(A)\|_{\mathrm{su}}$.
By (9.11) it contains the bound on $\|T_k(A)\|_{\mathrm{su}}$, at all unit sites $y$.
For $T^{(2)}(A)$, the display (I.5.4) together with (9.1), (9.P2), (9.Sm2) and $\#|\mathrm{Ef}_3|\le48\,\eta_k^{-2}$ (9.12) gives
\begin{equation*}
\left|\left(T^{(2)}(A)\chi\right)(y)\right|_2\le b_k^{-1}\cdot\eta_k^{\alpha-1}\cdot\frac{\eta_k^3}{48}\cdot48\,\eta_k^{-2}\exp\left(\frac{3}{2}\mathfrak{i}\right)\cdot6\,|\chi|_{\mathrm{su}}\le6\,b_{-}^{-1}\exp\left(\frac{3}{2}\mathfrak{i}\right)\eta_k^{\alpha}\,|\chi|_{\mathrm{su}}.
\end{equation*}
This proves (c).

\subsection*{5. Proof of (d)}

\begin{proof}
Let $k\ge1$ and $A\in\tilde{R}_k+\mathfrak{X}_k$ (Section~I.5.3; \cite[Definition~5.16]{QEDI}).

\emph{The numbers $\mathfrak{i}$ and $\mathrm{qd}$.}
The display (G3) of \cite[Section~16.2.7]{QEDI} is a statement on bond functions on $T_k$ for the data $(k,\ell,e_k)$ of (SF), under the condition (Z2.s4) on $e_0$.
For $A\in\tilde{R}_k+\mathfrak{X}_k$ it gives
\begin{equation*}
\mathfrak{i}(A)\le e_k^{\frac{1}{4}+\varepsilon}+e_k^{2\varepsilon}\le1,\qquad e_k\cdot|dA|_\infty\le4e_k^{\frac{1}{4}+2\varepsilon}+2e_k^{3\varepsilon}\le1.
\end{equation*}
Thus $\mathfrak{i}\le1$ and $\mathrm{qd}\le1$.

\emph{The unbarred operators.}
By (SF), $\mu\le1$.
Moreover $\eta_k^{\alpha}\le1$, and $\mathrm{e}\ge2$, so that $12\le6\,\mathrm{e}$ and $12+6\,\mathrm{e}\le12\,\mathrm{e}$.
The right sides of (c) are increasing functions of $\mathfrak{i}$, $\mathrm{qd}$ and $\mu$.
Hence
\begin{align*}
\|T^{(1)}(A)\|_{\mathrm{su}},\ \|T_k(A)\|_{\mathrm{su}}&\le\mathrm{e}^{\frac{3}{2}}\left[b_{-}^{-1}\left(54L\cdot\mathrm{e}^6+12+6\,\mathrm{e}+1\right)+1\right]\\
&\le\mathrm{e}^{\frac{3}{2}}\left[b_{-}^{-1}\left(54L\cdot\mathrm{e}^6+12\,\mathrm{e}+1\right)+1\right]=\frac{C_T}{2},\\
\|T^{(2)}(A)\|_{\mathrm{su}}&\le6\,b_{-}^{-1}\,\mathrm{e}^{\frac{3}{2}}\le\frac{C_T}{2}.
\end{align*}
At the unit sites outside $\mathbf{Y}$ the rows vanish.
So the hypothesis of (9.No2) holds with $\mathrm{Cn}=C_T/2$ for $O=T^{(1)}(A)$, $T^{(2)}(A)$, $T_k(A)$, and
\begin{equation*}
|T^{(1)}(A)|_{1,\infty},\ |T^{(2)}(A)|_{1,\infty},\ |T_k(A)|_{1,\infty}\le C_T.
\end{equation*}
For $T^{(2)}(A)$ the indices with a non-zero coefficient in the row at $y$ are the pairs $(x,x^{(3)})$, $x\in\mathrm{Ef}_3$.
They are pairwise incongruent modulo $\ell\mathbb{Z}^3$, because their first points lie over $\Delta_y$.
The row sum at $(y,\beta)$ is therefore the sum of the moduli of the coefficients of the values $\chi(x,x^{(3)})_{\beta'}$, as stated in (d).
The same number can be read off (I.5.4): with (9.1), (9.P2), (9.12) and $|M_3(x)|_{\mathrm{c}}\le6$ (9.Sm2) it equals the left side of
\begin{equation*}
b_k^{-1}\cdot\eta_k^{\alpha-1}\cdot\frac{\eta_k^3}{48}\sum_{x\in\mathrm{Ef}_3}|\varphi_A(y,c)|\sum_{\beta'}\left|M_3(x)_{\beta\beta'}\right|\le6\,b_{-}^{-1}\exp\left(\frac{3}{2}\mathfrak{i}\right)\eta_k^{\alpha}.
\end{equation*}

\emph{The barred operators.}
The sets $R_k$ and $\mathfrak{X}_k$ are defined by bounds on norms of the bond function, so they are invariant under $A\mapsto-A$.
Let $A_0$ be real, and suppose that for an $M$-cube $\square$ one has $A_0=A'+d\boldsymbol{\lambda}$ on $\square^{\dagger}$ with real $\boldsymbol{\lambda}$ and $A'$, the three norms of $A'$ lying below the thresholds.
Then $-A_0=(-A')+d(-\boldsymbol{\lambda})$, and $-A'$ has the same norms.
Therefore $\tilde{R}_k$, and with it $\tilde{R}_k+\mathfrak{X}_k$, is invariant under $A\mapsto-A$, and the bounds just proved hold at $-A$.
By (I.5.4) and Section~I.3.5,
\begin{equation*}
\bar{T}^{(i)}(A)=\mathsf{C}^{-1}T^{(i)}(-A)\,\mathsf{C},\qquad\bar{T}_k(A)=\mathsf{C}^{-1}T_k(-A)\,\mathsf{C},
\end{equation*}
and $\mathsf{C}$ is monomial (I.c6).
By (9.No3),
\begin{equation*}
|\bar{T}^{(i)}(A)|_{1,\infty}=|T^{(i)}(-A)|_{1,\infty}\le C_T\quad(i=1,2),\qquad|\bar{T}_k(A)|_{1,\infty}=|T_k(-A)|_{1,\infty}\le C_T.\qedhere
\end{equation*}
\end{proof}

\subsection*{6. Proof of (e)}

Let $k\ge1$, let $A$ be a connection, and let $\mathbf{Y}$ be a union of $M$-cubes.
All the relations to be proved are identities between elements of $\mathfrak{S}(1,\ell)$.
At a unit site outside $\mathbf{Y}$ (in 6.2 and 6.5: outside $\tilde{g}\mathbf{Y}$) both sides vanish.
Indeed, the rows of $T^{(i)}$ and of $\bar{T}^{(i)}$ vanish there, multiplication by a function of $y$ or by a constant matrix maps $0$ to $0$, and $(\tilde{g}v)(y')=S\,v(\tilde{g}^{-1}y')$, $(\mathrm{glb}\,v)(y')=(S^{-1})^T\,v(\tilde{g}^{-1}y')$ with $\tilde{g}^{-1}y'$ outside $\mathbf{Y}$.
It therefore suffices to compare the values at the unit sites $y$ of $\mathbf{Y}$ (respectively at the sites $\tilde{g}y$ of $\tilde{g}\mathbf{Y}$).
The set of the pairs $(x,i)$ with $x$ over $\Delta_y$ and $x^{(i)}$ over $\mathbf{Y}$ does not depend on the connection.

\subsubsection*{6.1. Gauge transformations}

Let $\boldsymbol{\lambda}$ be a real function of $c$ on $T_k$; as in (e), $\boldsymbol{\lambda}^{\mathrm{u}}$ is its restriction to $T^0_k$ and $\chi^{\boldsymbol{\lambda}}(x,x')=\mathrm{e}^{ie_k\boldsymbol{\lambda}(c)}\chi(x,x')$.
Then $\mathrm{e}^{ie_k\boldsymbol{\lambda}}\psi$ lies in $\mathfrak{S}^{\mathrm{ch}}(\eta_k,\ell)$, and $\chi^{\boldsymbol{\lambda}}$ is a function of pairs, since $\boldsymbol{\lambda}$ is $\ell\mathbb{Z}^3$-periodic.
The display (9.9), formed with $A-d\boldsymbol{\lambda}$ and $\mathrm{e}^{ie_k\boldsymbol{\lambda}}\psi$, contains three kinds of terms.
By (9.5), and because $x^{(i)}$ lies over $c$ for $i\le2$, they satisfy
\begin{align*}
&\varphi_{A-d\boldsymbol{\lambda}}(y,c)\,\mathrm{e}^{ie_k\boldsymbol{\lambda}(c)}\psi(x)=\mathrm{e}^{ie_k\boldsymbol{\lambda}(y)}\varphi_A(y,c)\,\psi(x),\\
&\varphi_{A-d\boldsymbol{\lambda}}(y,c)\,M_i(x)\left[\mathrm{e}^{ie_k\boldsymbol{\lambda}(c)}\psi(x^{(i)})-\mathrm{e}^{ie_k\boldsymbol{\lambda}(c)}\psi(x)\right]\\
&\qquad=\mathrm{e}^{ie_k\boldsymbol{\lambda}(y)}\varphi_A(y,c)\,M_i(x)\left[\psi(x^{(i)})-\psi(x)\right]\qquad(i\le2),\\
&\varphi_{A-d\boldsymbol{\lambda}}(y,c)\,M_3(x)\left[u_{A-d\boldsymbol{\lambda}}(c,c_3)\,\mathrm{e}^{ie_k\boldsymbol{\lambda}(c_3)}\psi(x^{(3)})-\mathrm{e}^{ie_k\boldsymbol{\lambda}(c)}\psi(x)\right]\\
&\qquad=\varphi_{A-d\boldsymbol{\lambda}}(y,c)\,\mathrm{e}^{ie_k\boldsymbol{\lambda}(c)}M_3(x)\left[u_A(c,c_3)\,\psi(x^{(3)})-\psi(x)\right]\\
&\qquad=\mathrm{e}^{ie_k\boldsymbol{\lambda}(y)}\varphi_A(y,c)\,M_3(x)\left[u_A(c,c_3)\,\psi(x^{(3)})-\psi(x)\right].
\end{align*}
Likewise, in (I.5.4) for $T^{(2)}$,
\begin{equation*}
\varphi_{A-d\boldsymbol{\lambda}}(y,c)\,M_3(x)\,\mathrm{e}^{ie_k\boldsymbol{\lambda}(c)}\chi(x,x^{(3)})=\mathrm{e}^{ie_k\boldsymbol{\lambda}(y)}\varphi_A(y,c)\,M_3(x)\,\chi(x,x^{(3)}).
\end{equation*}
Every term is thus multiplied by $\mathrm{e}^{ie_k\boldsymbol{\lambda}(y)}=\mathrm{e}^{ie_k\boldsymbol{\lambda}^{\mathrm{u}}(y)}$, and
\begin{equation}\tag{9.20}
T^{(1)}(A-d\boldsymbol{\lambda})\,\mathrm{e}^{ie_k\boldsymbol{\lambda}}\psi=\mathrm{e}^{ie_k\boldsymbol{\lambda}^{\mathrm{u}}}T^{(1)}(A)\psi,\qquad
T^{(2)}(A-d\boldsymbol{\lambda})\,\chi^{\boldsymbol{\lambda}}=\mathrm{e}^{ie_k\boldsymbol{\lambda}^{\mathrm{u}}}T^{(2)}(A)\chi.
\end{equation}

\subsubsection*{6.2. The lattice group}

Let $\tilde{g}=(S,a)\in\mathbb{G}_k$, with $a$ chosen in $M\mathbb{Z}^3$, and let $r=r(S)$.
Then $\tilde{g}x=(rc+a,rg)$ (I.4.2), and $(\tilde{g}\psi)(\tilde{g}x)=S\,\psi(x)$, $(\tilde{g}\chi)(\tilde{g}x,\tilde{g}x')=S\,\chi(x,x')$, $(\tilde{g}v)(\tilde{g}y)=S\,v(y)$ (Section~I.4.3).
Let $y$ be a unit site of $\mathbf{Y}$.
By (9.P4), $\tilde{g}y$ is a unit site of the union $\tilde{g}\mathbf{Y}$ of $M$-cubes, and $\tilde{g}\Delta_y=\Delta_{\tilde{g}y}$; moreover $g\mapsto rg$ is a bijection of $O_{\mathrm{cub}}$.
Hence $x\mapsto\tilde{g}x$ is a bijection of $\Delta_y\times O_{\mathrm{cub}}$ onto $\Delta_{\tilde{g}y}\times O_{\mathrm{cub}}$ (I.a1).
In the sum (9.1) for $\left(T^{(1)}_{\tilde{g}\mathbf{Y}}(\tilde{g}A)\,\tilde{g}\psi\right)(\tilde{g}y)$, formed by (I.5.4), write the summation variable as $\tilde{g}x$ with $x$ over $\Delta_y$.
Then:
\begin{itemize}
\item $(\tilde{g}\psi)(\tilde{g}x)=S\,\psi(x)$; $(\tilde{g}x)^{(i)}=\tilde{g}(x^{(i)})$ (I.a2), so that $(\tilde{g}\psi)\left((\tilde{g}x)^{(i)}\right)=S\,\psi(x^{(i)})$; $\pi\left((\tilde{g}x)^{(i)}\right)=\tilde{g}\,\pi(x^{(i)})$ (I.a1), which lies in $\tilde{g}\mathbf{Y}$ if and only if $\pi(x^{(i)})$ lies in $\mathbf{Y}$;
\item $M_i(\tilde{g}x)\,S=S\,M_i(x)$ (9.Sm3); $\varphi_{\tilde{g}A}(\tilde{g}y,\tilde{g}c)=\varphi_A(y,c)$ (9.6); $U_{\tilde{g}A}\left(\tilde{g}x,(\tilde{g}x)^{(i)}\right)=1=U_A(x,x^{(i)})$ for $i\le2$, and $U_{\tilde{g}A}\left(\tilde{g}x,(\tilde{g}x)^{(3)}\right)=u_{\tilde{g}A}(\tilde{g}c,\tilde{g}c_3)=u_A(c,c_3)$ (9.6).
\end{itemize}
Consequently
\begin{align*}
\left(T^{(1)}_{\tilde{g}\mathbf{Y}}(\tilde{g}A)\,\tilde{g}\psi\right)(\tilde{g}y)
&=\frac{\eta_k^3}{48}\sum_{x\text{ over }\Delta_y}\varphi_A(y,c)\Big\{\left(1+\frac{\mu}{b_k}\right)S\,\psi(x)\\
&\qquad-(b_k\eta_k)^{-1}\sum_{i:\,x^{(i)}\text{ over }\mathbf{Y}}S\,M_i(x)\left[U_A(x,x^{(i)})\,\psi(x^{(i)})-\psi(x)\right]\Big\}\\
&=S\left(T^{(1)}_{\mathbf{Y}}(A)\psi\right)(y)=\left(\tilde{g}\,T^{(1)}_{\mathbf{Y}}(A)\psi\right)(\tilde{g}y),\\
\left(T^{(2)}_{\tilde{g}\mathbf{Y}}(\tilde{g}A)\,\tilde{g}\chi\right)(\tilde{g}y)
&=-b_k^{-1}\cdot\eta_k^{\alpha-1}\cdot\frac{\eta_k^3}{48}\sum_{x\text{ over }\Delta_y}\mathbf{1}[\,c_3\notin\mathbf{Y}\,]\,\varphi_A(y,c)\,S\,M_3(x)\,\chi(x,x^{(3)})\\
&=S\left(T^{(2)}_{\mathbf{Y}}(A)\chi\right)(y)=\left(\tilde{g}\,T^{(2)}_{\mathbf{Y}}(A)\chi\right)(\tilde{g}y),
\end{align*}
where $(\tilde{g}\chi)\left(\tilde{g}x,(\tilde{g}x)^{(3)}\right)=(\tilde{g}\chi)\left(\tilde{g}x,\tilde{g}x^{(3)}\right)=S\,\chi(x,x^{(3)})$ has been used.
Every unit site of $\tilde{g}\mathbf{Y}$ is of the form $\tilde{g}y$.
Hence
\begin{equation}\tag{9.21}
T^{(i)}_{\tilde{g}\mathbf{Y}}(\tilde{g}A)\,\tilde{g}=\tilde{g}\,T^{(i)}_{\mathbf{Y}}(A)\qquad(i=1,2).
\end{equation}
If $a$ is replaced by $a+\ell n$, $n\in\mathbb{Z}^3$, the operator $\tilde{g}$ is multiplied by $(-1)^{n_1+n_2+n_3}$ on each of the three spaces, whereas $\tilde{g}A$ and $\tilde{g}\mathbf{Y}$ do not change; both sides of (9.21) contain $\tilde{g}$ once.

\subsubsection*{6.3. Conjugation}

For an operator $O$, the operator $\overline{O}$ is $f\mapsto\overline{O\,\overline{f}}$ (I.O5).
The numbers $\mu$, $b_k$, $\eta_k$, $\alpha$ are real.
By (9.7) and (9.Sm3), the terms of (9.9) and of (I.5.4) satisfy
\begin{align*}
\overline{\varphi_{\overline{A}}(y,c)\,\overline{\psi(x)}}&=\varphi_{-A}(y,c)\,\psi(x)=\varkappa^{-1}\varphi_{-A}(y,c)\,\varkappa\,\psi(x),\\
\overline{\varphi_{\overline{A}}(y,c)\,M_i(x)\left(U_{\overline{A}}(x,x^{(i)})\,\overline{\psi(x^{(i)})}-\overline{\psi(x)}\right)}
&=\varphi_{-A}(y,c)\,\varkappa^{-1}M_i(x)\,\varkappa\left(U_{-A}(x,x^{(i)})\,\psi(x^{(i)})-\psi(x)\right)\\
&=\varkappa^{-1}\varphi_{-A}(y,c)\,M_i(x)\left(U_{-A}(x,x^{(i)})\,\varkappa\,\psi(x^{(i)})-\varkappa\,\psi(x)\right),\\
\overline{\varphi_{\overline{A}}(y,c)\,M_3(x)\,\overline{\chi(x,x^{(3)})}}&=\varkappa^{-1}\varphi_{-A}(y,c)\,M_3(x)\,\varkappa\,\chi(x,x^{(3)}).
\end{align*}
Summation yields
\begin{equation}\tag{9.22}
\overline{T^{(i)}(\overline{A})}=\varkappa^{-1}T^{(i)}(-A)\,\varkappa\qquad(i=1,2).
\end{equation}

\subsubsection*{6.4. Barred operators: gauge transformations}

One has
\begin{equation*}
\bar{T}^{(i)}(A-d\boldsymbol{\lambda})=\mathsf{C}^{-1}T^{(i)}(-A+d\boldsymbol{\lambda})\,\mathsf{C}=\mathsf{C}^{-1}T^{(i)}\left((-A)-d(-\boldsymbol{\lambda})\right)\mathsf{C}.
\end{equation*}
For the connection $-A$ and the real function $-\boldsymbol{\lambda}$ the relations (9.20) read
\begin{equation*}
T^{(1)}\left((-A)-d(-\boldsymbol{\lambda})\right)\mathrm{e}^{-ie_k\boldsymbol{\lambda}}=\mathrm{e}^{-ie_k\boldsymbol{\lambda}^{\mathrm{u}}}T^{(1)}(-A),\qquad
T^{(2)}\left((-A)-d(-\boldsymbol{\lambda})\right)\chi^{-\boldsymbol{\lambda}}=\mathrm{e}^{-ie_k\boldsymbol{\lambda}^{\mathrm{u}}}T^{(2)}(-A)\chi.
\end{equation*}
The constant matrix $\mathsf{C}$ commutes with the multiplication by functions of $c$ and of $y$, and with $\chi\mapsto\chi^{-\boldsymbol{\lambda}}$.
Therefore
\begin{align*}
\bar{T}^{(1)}(A-d\boldsymbol{\lambda})\,\mathrm{e}^{-ie_k\boldsymbol{\lambda}}
&=\mathsf{C}^{-1}T^{(1)}\left((-A)-d(-\boldsymbol{\lambda})\right)\mathrm{e}^{-ie_k\boldsymbol{\lambda}}\,\mathsf{C}\\
&=\mathrm{e}^{-ie_k\boldsymbol{\lambda}^{\mathrm{u}}}\mathsf{C}^{-1}T^{(1)}(-A)\,\mathsf{C}=\mathrm{e}^{-ie_k\boldsymbol{\lambda}^{\mathrm{u}}}\bar{T}^{(1)}(A),\\
\bar{T}^{(2)}(A-d\boldsymbol{\lambda})\,\chi^{-\boldsymbol{\lambda}}
&=\mathsf{C}^{-1}T^{(2)}\left((-A)-d(-\boldsymbol{\lambda})\right)(\mathsf{C}\chi)^{-\boldsymbol{\lambda}}\\
&=\mathsf{C}^{-1}\mathrm{e}^{-ie_k\boldsymbol{\lambda}^{\mathrm{u}}}T^{(2)}(-A)\,\mathsf{C}\chi=\mathrm{e}^{-ie_k\boldsymbol{\lambda}^{\mathrm{u}}}\bar{T}^{(2)}(A)\chi.
\end{align*}

\subsubsection*{6.5. Barred operators: the lattice group}

The action $\mathrm{glb}$ is given by $(\mathrm{glb}\,f)(x)=(S^{-1})^Tf(\tilde{g}^{-1}x)$, and by the same matrix on the two other spaces: $(\mathrm{glb}\,F)(x,x')=(S^{-1})^TF(\tilde{g}^{-1}x,\tilde{g}^{-1}x')$ and $(\mathrm{glb}\,v)(y)=(S^{-1})^Tv(\tilde{g}^{-1}y)$ (Section~G1.10).
By (9.Sm4), $\mathrm{glb}=\mathrm{sn}(S)\,\mathsf{C}^{-1}\tilde{g}\,\mathsf{C}$ on each of the three spaces, with one sign $\mathrm{sn}(S)$.
On bond functions $\tilde{g}$ is linear: $-\tilde{g}A=\tilde{g}(-A)$.
With (9.21) at the connection $-A$ it follows that
\begin{align*}
\bar{T}^{(i)}_{\tilde{g}\mathbf{Y}}(\tilde{g}A)\,\mathrm{glb}
&=\mathsf{C}^{-1}T^{(i)}_{\tilde{g}\mathbf{Y}}\left(\tilde{g}(-A)\right)\mathsf{C}\cdot\mathrm{sn}(S)\,\mathsf{C}^{-1}\tilde{g}\,\mathsf{C}
=\mathrm{sn}(S)\,\mathsf{C}^{-1}T^{(i)}_{\tilde{g}\mathbf{Y}}\left(\tilde{g}(-A)\right)\tilde{g}\,\mathsf{C}\\
&=\mathrm{sn}(S)\,\mathsf{C}^{-1}\tilde{g}\,T^{(i)}_{\mathbf{Y}}(-A)\,\mathsf{C}
=\left(\mathrm{sn}(S)\,\mathsf{C}^{-1}\tilde{g}\,\mathsf{C}\right)\left(\mathsf{C}^{-1}T^{(i)}_{\mathbf{Y}}(-A)\,\mathsf{C}\right)
=\mathrm{glb}\,\bar{T}^{(i)}_{\mathbf{Y}}(A).
\end{align*}

\subsubsection*{6.6. Barred operators: conjugation}

The matrix $\mathsf{C}$ is real (I.c3); hence $\overline{\mathsf{C}^{-1}O\,\mathsf{C}}=\mathsf{C}^{-1}\,\overline{O}\,\mathsf{C}$ for every operator $O$.
Moreover $-\overline{A}=\overline{-A}$.
By (9.22) at the connection $-A$,
\begin{equation*}
\overline{\bar{T}^{(i)}(\overline{A})}=\mathsf{C}^{-1}\left[\overline{T^{(i)}\left(\overline{-A}\right)}\right]\mathsf{C}=\mathsf{C}^{-1}\varkappa^{-1}T^{(i)}(A)\,\varkappa\,\mathsf{C}.
\end{equation*}
By (I.c4), $\varkappa\,\mathsf{C}=-\mathsf{C}\,\varkappa$; taking inverses, $\mathsf{C}^{-1}\varkappa^{-1}=-\varkappa^{-1}\mathsf{C}^{-1}$.
The two signs cancel, and
\begin{equation*}
\overline{\bar{T}^{(i)}(\overline{A})}=\varkappa^{-1}\mathsf{C}^{-1}T^{(i)}(A)\,\mathsf{C}\,\varkappa=\varkappa^{-1}\bar{T}^{(i)}(-A)\,\varkappa,
\end{equation*}
because $\bar{T}^{(i)}(-A)=\mathsf{C}^{-1}T^{(i)}(A)\,\mathsf{C}$.
This completes the proof of (e).

\subsection*{7. Level 0}

At $k=0$ the operators of Proposition~9 are defined by Section~I.3.9, that is by (I.L0.7): for every connection $A$,
\begin{equation*}
T_0(A)=\bar{T}_0(A)=H_0(A)=\bar{H}_0(A):=1\qquad\text{on }\mathfrak{S}^{\mathrm{ch}}(1,\ell).
\end{equation*}
Hence $T_0(A)H_0(A)=1$ and $\bar{T}_0(A)\bar{H}_0(A)=1$ on $\mathfrak{S}^{\mathrm{ch}}(1,\ell)$, a space with $192$ components per unit site; the space $\mathfrak{S}(1,\ell)$ does not occur.
No inverse occurs, and no condition on $A$, on $e_k$ or on the constants is needed.
The formulas $\bar{H}_k(A)=\mathsf{C}^{-1}H_k(-A)\,\mathsf{C}$ and $\bar{T}_k(A)=\mathsf{C}^{-1}T_k(-A)\,\mathsf{C}$ of Section~I.3.5 give the same operators at $k=0$, since $\mathsf{C}^{-1}1\,\mathsf{C}=1$.
This proves the last sentence of (a).

The number $b_0$ and the operators $T^{(1)}$, $T^{(2)}$ are not defined at $k=0$ (I.L0.9).
The clauses (b)--(e) are statements for $k\ge1$ and have no part at level $0$.

\subsection*{8. Uniformity; conditions for the individual clauses; two supplements}

\subsubsection*{8.1. Uniformity}

\begin{named}{Remark}
The right sides of (c) are functions of $(b_{-},L,\mu,\mathfrak{i},\mathrm{qd},\eta_k^{\alpha})$ alone, and $C_T$ is a function of $(b_{-},L)$ alone.
Indeed, the level $k$ enters the proof of (c) only in the following ways:
through $b_k\ge b_{-}$;
through the cardinalities $\#|\mathrm{Ef}|=48\,\eta_k^{-3}$ and $\#|\mathrm{Ef}\setminus\mathrm{Ef}_1|=48\,\eta_k^{-2}$, which are compensated by the weights $\frac{\eta_k^3}{48}$ and $\frac{\eta_k^2}{48}$;
through the bounds (9.P2), whose right sides $\exp(\mathfrak{i})$ and $\exp\left(\frac{3}{2}\mathfrak{i}\right)$ do not contain $k$;
through the sum (9.18), which is bounded by $9L\cdot\mathrm{qd}\cdot\exp(6\,\mathrm{qd})$ for every $k\ge1$;
and, for $T^{(2)}$, through the factor $\eta_k^{\alpha}$ of the bound.
The side $\ell$ does not enter the bounds, since a row at $y$ contains sites, pairs and bonds within the sup distance $\frac{1+\eta_k}{2}$ of $y$.
The charge $e_k$ enters only through $\mathfrak{i}$ and $\mathrm{qd}$, and the same holds for $A$; the parameter $\rho$ enters only through $3(1+\rho)\le6$; neither $\mathbf{Y}$ nor $y$ enters.
In particular, the bounds depend neither on a cutoff $N$ nor on a torus $R$: neither occurs among the data (SF).
\end{named}

\subsubsection*{8.2. Conditions for the individual clauses}

\begin{named}{Remark}
Each clause of Proposition~9, and each relation of Sections~6.1--6.6, holds under the conditions listed for it below, by the proofs given above.
\begin{description}
\item[(a)] $F_{\infty}(A)^{-1}$ and $F_{\infty}(-A)^{-1}$ exist, and (I.Q1) holds.
\item[(b), 6.1, 6.2, 6.4, 6.5] $k\ge1$; $L$ and $M$ are odd; $\ell$ is a multiple of $M$; $L>1$, hence $L\ge3$, and $b\ne0$, so that the number $b_k$ of Section~I.3.4 is defined and different from $0$.
The number $\mu$ occurs only in the scalar coefficient $1+\frac{\mu}{b_k}$ of $\psi(x)$: the operators $T_k(A)$ and $T^{(1)}(A)$ are affine in $\mu$, and $T^{(2)}(A)$ does not contain $\mu$.
The proofs of (b), 6.1, 6.2, 6.4 and 6.5 use no property of this scalar, neither an inverse, nor a conjugate, nor a modulus.
Hence $\mu$ may be any complex number.
The charge $e_k$ may be any positive number.
In 6.1 and 6.4, $\boldsymbol{\lambda}$ may be any complex function on $\mathbb{L}_k$ for which $d\boldsymbol{\lambda}$ and $\mathrm{e}^{ie_k\boldsymbol{\lambda}}$ are $\ell\mathbb{Z}^3$-periodic: (9.P3) involves only differences of $\boldsymbol{\lambda}$ along paths, and the periodicity of $\mathrm{e}^{ie_k\boldsymbol{\lambda}}$ is needed for $\mathrm{e}^{ie_k\boldsymbol{\lambda}}\psi\in\mathfrak{S}^{\mathrm{ch}}(\eta_k,\ell)$.
\item[6.3, 6.6] The conditions of (b), and $\mu$ real.
\item[(c)] The conditions of (b) and, in addition, $0\le\rho\le1$ (this enters (9.Sm2) through (I.m6): the equality $|M_3(x)|_2=3(1+\rho)$ requires $\rho\ge0$, the bound $3(1+\rho)\le6$ requires $\rho\le1$); the closure relation (I.s3); the value $w_3=6$ (9.Sm2); $L\ge3$; $b>0$, so that $b_{-}>0$; $b_k\ge b_{-}$; $\mu\ge0$.
The bounds (9.P2) hold for every $e_k>0$ under the conditions of (b).
Indeed, by (9.4), $\left|u_A(c,c')^{\pm1}\right|=\exp\left(\mp e_k\cdot\eta_k\cdot\operatorname{Im}A(c,c')\right)\le\exp\left(\eta_k\cdot\mathfrak{i}(A)\right)$, because $\mathfrak{i}(A)=e_k\cdot|\operatorname{Im}A|_{\infty}$ by (I.1.6).
By (9.4) and (9.P1), $(\tau_kA)(y,c)$ is a sum over $j<k$ of means, with real coefficients, of the values of $A$ on paths with at most $\frac{3}{2}\left(L^{j+1}-L^j\right)$ bonds; hence $\operatorname{Im}(\tau_kA)(y,c)=(\tau_k\operatorname{Im}A)(y,c)$, and
\begin{equation*}
e_k\cdot\eta_k\cdot\left|\operatorname{Im}(\tau_kA)(y,c)\right|\le\frac{3}{2}\left(L^k-1\right)\cdot\eta_k\cdot\mathfrak{i}(A)=\frac{3}{2}\left(1-\eta_k\right)\cdot\mathfrak{i}(A),
\end{equation*}
so that $\left|\varphi_A(y,c)^{\pm1}\right|=\exp\left(\mp e_k\cdot\eta_k\cdot\operatorname{Im}(\tau_kA)(y,c)\right)\le\exp\left(\frac{3}{2}\,\mathfrak{i}(A)\right)$.
This uses only that $L$ is odd, that $k\ge1$ and that $e_k$ is a positive real number.
The inequality of \cite[Lemma~16.26]{QEDI}, which is applied in 4.4, Step~2, on the box $\mathbb{W}$ of $\eta_k\mathbb{Z}^3$, imposes no condition on $L$, because it does not depend on the mesh.
On a box of mesh $\eta$ the number $\eta\cdot(dA)(p)$ is the sum of $A$ over the boundary of the plaquette $p$, and the inequality takes the form
\begin{equation*}
|A(\boldsymbol{\ell})|\le2n\cdot\frac{s_W}{\eta}\cdot\sup_p\left|\eta\cdot(dA)(p)\right|,
\end{equation*}
where $n$ is the number of bonds of the closed path $\boldsymbol{\ell}$; besides $A(\boldsymbol{\ell})$, $n$ and these plaquette sums, only the ratio $\frac{s_W}{\eta}$ occurs in it.
The data of \cite[Section~16.2.1]{QEDI} admit the level $k=0$, at which the mesh equals $1$, and for the mesh $1$ no other datum of that section occurs in the statement of the lemma; so the lemma holds on every box of $\mathbb{Z}^3$, whatever $L$ is.
The dilation $x\mapsto\eta^{-1}x$ maps a box of $\eta\mathbb{Z}^3$ whose intervals have length at most $s_W$ onto a box of $\mathbb{Z}^3$ whose intervals have length at most $\frac{s_W}{\eta}$, and it maps bonds, plaquettes and closed paths with $n$ bonds onto bonds, plaquettes and closed paths with $n$ bonds.
The function $A$, transported by the dilation, has the same values on corresponding bonds and the same sums over the boundaries of corresponding plaquettes.
Therefore the case of the mesh $1$ gives the inequality on every box of $\eta\mathbb{Z}^3$, for every $\eta>0$, and in particular for $\eta=\eta_k=L^{-k}$ with $L\ge3$; no condition $L\ge5$ enters (c).
For complex $\mu$ the proof gives (c) with $|\mu|$ in the place of $\mu$.
\item[(d)] The conditions of (c) and, in addition, $\mu\le1$, $e_k\le e_0$ with (Z2.s4), and $\alpha\ge0$, which gives $\eta_k^{\alpha}\le1$ in Section~5.
Further, $\ell$, $L$, $M$, $\alpha$ and $\varepsilon$ are as required for the sets $\tilde{R}_k$ and $\mathfrak{X}_k$ of \cite[Definition~5.16]{QEDI} and for the display (G3) of \cite[Section~16.2.7]{QEDI}.
The cubes $\square^{\dagger}$ and the third norm $|\delta_{\alpha}\partial A|_{\infty}$ enter the definition of $\tilde{R}_k$ only;
the set $\mathfrak{X}_k$ is defined by the two bounds $e_k^{-1+2\varepsilon}$ and $e_k^{-1+3\varepsilon}$ on the sup norm of a bond function and on the sup norm of its difference quotients, and of $\ell$, $L$, $M$, $\alpha$, $\varepsilon$ these bounds contain only $\varepsilon$.
The requirements are:
$\ell\in\mathrm{Cl}$ (Section~I.1.1);
$L$ odd with $L\ge5$, which is (L5);
$M$ a power of $L$ with $M>4L$, which follows from (M1) together with (L5);
$\alpha\in(0,1)$, which holds by the line (C1) of Section~I.8.2, $\alpha\in(1-\bar{\varepsilon},1)$;
and $\varepsilon$ as fixed in Section~I.8, so that $0<\varepsilon\le\bar{\varepsilon}$.
Thus, for (d), neither the side $\ell$ nor $L$, $M$, $\alpha$, $\varepsilon$ ranges beyond (SF) and Section~I.8.
\end{description}
\end{named}

\subsubsection*{8.3. Supplement 1: the displays of the barred operators and the barred form of (b)}

By (I.5.4), (9.7) and (9.Sm5), for every connection $A$, every union $\mathbf{Y}$ of $M$-cubes and every unit site $y$ of $\mathbf{Y}$,
\begin{equation}\tag{9.23}
\begin{aligned}
\left(\bar{T}^{(1)}(A)f\right)(y)&=\int_{\Delta_y}\mathrm{d}x\,\varphi_A(y,c)^{-1}\Big\{\left(1+\frac{\mu}{b_k}\right)f(x)\\
&\qquad-(b_k\eta_k)^{-1}\sum_{i:\,x^{(i)}\text{ over }\mathbf{Y}}M_i(x^{(i)})^T\left[U_A(x,x^{(i)})^{-1}f(x^{(i)})-f(x)\right]\Big\},\\
\left(\bar{T}^{(2)}(A)F\right)(y)&=-b_k^{-1}\cdot\eta_k^{\alpha-1}\int_{\Delta_y}\mathrm{d}x\,\mathbf{1}[\,c_3\notin\mathbf{Y}\,]\,\varphi_A(y,c)^{-1}M_3(x^{(3)})^TF(x,x^{(3)}).
\end{aligned}
\end{equation}
Thus the rows of $\bar{T}^{(i)}(A)$ contain $f$, $F$ at the same sites and pairs as the rows of $T^{(i)}(A)$, and they read $A$ on the same bonds.
Suppose that $F(x,x^{(3)})=\left(\delta_{\alpha,-A}f\right)(x,x^{(3)})$ for the $x$ over $\Delta_y$ with $c_3\notin\mathbf{Y}$.
Then $\bar{T}^{(1)}(A)f+\bar{T}^{(2)}(A)F=\bar{T}_k(A)f$ at $y$.
To see this, apply (b) at the connection $-A$ to $\psi:=\mathsf{C}f$ and $\chi:=\mathsf{C}F$ (the operator $\delta_{\alpha,-A}$ commutes with the constant matrix $\mathsf{C}$), and multiply by $\mathsf{C}^{-1}$.

\subsubsection*{8.4. Supplement 2: $T_k$ and $\bar{T}_k$}

The torus satisfies $\tilde{g}T_k=T_k$.
By (9.11), the relations of (e) with $\mathbf{Y}=T_k$ are relations for $T_k(A)=T^{(1)}_{T_k}(A)$ and $\bar{T}_k(A)=\bar{T}^{(1)}_{T_k}(A)$.
So, for every connection $A$,
\begin{gather*}
T_k(A-d\boldsymbol{\lambda})\,\mathrm{e}^{ie_k\boldsymbol{\lambda}}=\mathrm{e}^{ie_k\boldsymbol{\lambda}^{\mathrm{u}}}T_k(A),\qquad
T_k(\tilde{g}A)\,\tilde{g}=\tilde{g}\,T_k(A),\qquad
\overline{T_k(\overline{A})}=\varkappa^{-1}T_k(-A)\,\varkappa,\\
\bar{T}_k(A-d\boldsymbol{\lambda})\,\mathrm{e}^{-ie_k\boldsymbol{\lambda}}=\mathrm{e}^{-ie_k\boldsymbol{\lambda}^{\mathrm{u}}}\bar{T}_k(A),\qquad
\bar{T}_k(\tilde{g}A)\,\mathrm{glb}=\mathrm{glb}\,\bar{T}_k(A),\qquad
\overline{\bar{T}_k(\overline{A})}=\varkappa^{-1}\bar{T}_k(-A)\,\varkappa.
\end{gather*}

\subsection*{9. Conditions on the constants}

The following three tables list the conditions on the constants and the bounds on the data that enter the arguments of this chapter, the places where they enter, and where they are stated.

\noindent
\begin{tabular}{p{0.31\textwidth}p{0.33\textwidth}p{0.27\textwidth}}
\hline
Condition & Where it is used & Where it is stated\\
\hline
$L$ odd
& (9.B1)--(9.B3), (9.B5), (9.B6): blocks with centres on the lattice, nested
& (L5) of Proposition~14.2\,(b); (C3) of Section~I.8.2\\
$L\ge3$ (contained in $L\ge5$)
& (9.18): $\frac{L^2}{L-1}\le\frac{3L}{2}$
& (L5) of Proposition~14.2\,(b); (C3) of Section~I.8.2\\
$M$ an odd integer ($M$ a power of $L$, $L$ odd)
& (9.B4): $\Delta_y$ is contained in the $M$-cube of $y$
& (M1) with (L5) of Proposition~14.2\,(b); (C5) of Section~I.8.2; Section~I.1.1 (faces)\\
$\rho\in(0,1)$
& (9.Sm2): $3(1+\rho)\le6$
& (C0) of Section~I.8.2\\
the standard weights $w$: the closure relation (I.1.1); the value $w_3=6$
& (9.Sm1), through (I.s3); (9.Sm2): the numbers $3(1+\rho)=\frac{w_3}{2}(1+\rho)$ and $6=w_3$
& (C0) of Section~I.8.2; (I.BARE)\\
$b=1>0$ and $L>1$: $b_k$ is defined and positive, and $b_{-}>0$
& every display of Sections~2--6 and~8.3 that contains $b_k^{-1}$ or $\frac{\mu}{b_k}$; $b_{-}^{-1}$ in 4.5, in Section~5 and in $C_T$
& (C0) of Section~I.8.2 ($b=1$); (L5) of Proposition~14.2\,(b) and (C3) of Section~I.8.2 ($L>1$); (C4) of Section~I.8.2 and Section~I.3.4 ($b_k$, $b_{-}$)\\
$b_{-}=b\left(1-\frac{1}{L}\right)\le b_k$
& 4.5; Section 5
& Section~I.3.4; (C4) of Section~I.8.2\\
$\alpha>0$
& Section 5: $\eta_k^{\alpha}\le1$
& (C1) of Section~I.8.2\\
(Z2.s4): for $v\ge-\log e_0$, $\exp(-\varepsilon v)\le\frac{1}{4}$; $\exp\left(-\frac{v}{4}\right)+\exp(-\varepsilon v)\le1$; $4\exp\left(-\left(\frac{1}{4}+2\varepsilon\right)v\right)+2\exp(-3\varepsilon v)\le1$
& Section 5, through the display (G3) of \cite[Section~16.2.7]{QEDI}
& (Z2.s4) of Proposition~14.2\,(b); (C8) of Section~I.8.2; Section~I.8.4\\
\hline
\end{tabular}

\medskip
\noindent
\begin{tabular}{p{0.31\textwidth}p{0.33\textwidth}p{0.27\textwidth}}
\hline
Condition & Where it is used & Where it is stated\\
\hline
the data on the constants of \cite[Section~16.2.1]{QEDI}, under which (G2), (G3), \cite[Lemma~16.2(a)]{QEDI} and \cite[Lemma~16.26]{QEDI} are stated: $L$ odd, $L\ge5$; $b>0$; $\alpha\in(0,1)$; $M$ a power of $L$ with $M>4L$; $0<\varepsilon\le\bar{\varepsilon}$
& Section~2: (G2); (9.P2): \cite[Lemma~16.2(a)]{QEDI}; 4.4, Step~2: \cite[Lemma~16.26]{QEDI}, whose inequality holds for every mesh by 8.2\,(c); Section~5: (G3)
& (L5) of Proposition~14.2\,(b); (C0) of Section~I.8.2 ($b=1$); (C1) of Section~I.8.2 ($1-\bar{\varepsilon}<\alpha<1$); (M1) with (L5), or (M2), of Proposition~14.2\,(b); (S1) of Proposition~14.2\,(b) and (C7) of Section~I.8.2 ($1-\alpha<\varepsilon\le\bar{\varepsilon}$)\\
$\ell\in\mathrm{Cl}$, that is, $\ell/(LM)$ is an odd integer $\ge13$; hence $\ell$ is a multiple of $M$ and of $L$, and $\ell\ge3L$
& $\ell$ a multiple of $M$: 1.1 (a union of $M$-cubes of $T_k$ lifts to a union of $M$-cubes of $\mathbb{L}_k$) and Section~4 ($T_k$ is a union of $M$-cubes); $\ell$ a multiple of $L$ with $\ell\ge3L$: (9.P2) ($T_k$ is the block torus $(0,\ell,-1)$ of \cite[Definition~16.1]{QEDI}); $\ell\in\mathrm{Cl}$: Section~2 (Proposition~2.8\,(a) and (G2)) and Section~5 ((G3) and \cite[Definition~5.16]{QEDI})
& the data (SF), Section~I.1.8; the definition of $\mathrm{Cl}$ in Section~I.1.1\\
\hline
\end{tabular}

\medskip
\noindent
\begin{tabular}{p{0.31\textwidth}p{0.33\textwidth}p{0.27\textwidth}}
\hline
Condition & Where it is used & Where it is stated\\
\hline
$\mu>0$
& as $\mu\ge0$: (9.16), 4.5 and Section~5; as $\mu$ real: 6.3 and 6.6
& the data (SF), Section~I.1.8\\
$\mu\le1$; $e_k\le e_0$
& Section~5 ($\mu\le1$; (G3)); Section~2, through Proposition~2.8\,(a) and (G2); (9.P2), where $e_k\le e_0$ enters through \cite[Lemma~16.2(a)]{QEDI}, although by 8.2\,(c) the bounds (9.P2) hold for every $e_k>0$
& the data (SF), Section~I.1.8\\
the conditions of Proposition~2.8\,(a): (L5), (M1), (M2), $M\ge108\pi\cdot C'_{V}\cdot a'$ (in (M5)), (Z2.s1), (Z2.s2); $\alpha<1$ (through Proposition~2.7\,(d), (e))
& Section 2, through Proposition~2.8\,(a)
& Proposition~14.2\,(b); for $\alpha<1$: (C1) of Section~I.8.2 and Proposition~14.2\,(a)\\
${\begin{aligned}[t]C_T=2\,\mathrm{e}^{\frac{3}{2}}\bigl[&b_{-}^{-1}\bigl(54L\cdot\mathrm{e}^6\\&+12\,\mathrm{e}+1\bigr)+1\bigr]\end{aligned}}$
& definition; Section 5, with the numerical inequality $\mathrm{e}\ge2$, which gives $12+6\,\mathrm{e}\le12\,\mathrm{e}$
& (C4) of Section~I.8.2; Proposition~14.2\,(a)\\
\hline
\end{tabular}

\section*{Chapter 10. Covariances of all exact operators (FE6$^{\mathrm{ch}}$): proofs of Proposition~10}

\subsection*{0. Conventions of this chapter}

In this chapter the letters $y$, $y'$, $\tilde{y}$ always denote the resolvent parameter.
Unit sites are therefore written $z$, and sites of the $L$-lattice $z'$; both pairs of letters denote these sites in Section~I.1.1.

The set $\tilde{G}_k$ of Section~I.5.3 is the set of \cite[Definition~5.16\,(e)]{QEDI}.
In its definition, $\boldsymbol{\lambda}$ is a function on the sites of $\square^{\sim 5}$, $A'$ is a function on the bonds contained in $\square^{\sim 5}$, and the relation $A_0=A'+d\boldsymbol{\lambda}$ is required on these bonds.

\subsection*{1. The statement}

Proposition~10 of Part~II consists of the clauses (a)--(f), which concern the levels $k\ge1$, and of clause~(g), which gives (a)--(f) at the level $k=0$.
The clauses (a)--(f) deal, in this order, with the replacement of $A$ by $A-\mathrm{Acl}$ with a real connection $\mathrm{Acl}$ satisfying $d\,\mathrm{Acl}=0$, the maps $\tilde{g}\in\mathbb{G}_k$, transposition, complex conjugation, the number $\delta Z_k$, and the adjoints at real $A$.

\subsection*{2. Operations on operators; rules that are used throughout}

The quadruple $(k,\ell,\mu,e_k)$ consists of data (SF) (Section~I.1.8), and every object below is formed with this single quadruple.
The relations of Proposition~10 are identities between objects that depend on the constants $\rho$, $w_i$ (through the matrices $M_i(x)$), $b$, $L$, $M$, $\alpha$ and $\varepsilon$ (the last one only through $\tilde{G}_k$), on the numbers $\ell$, $\mu$, $e_k$, $\eta_k=L^{-k}$, $b_k$ that belong to the data, and on the parameter $y$.
The conditions on the constants and the properties of the data that the proofs of Sections~2--6 require, directly or through the propositions applied there, are collected in Section~8.
Every relation below is an identity, and it is proved for every quadruple (SF), every $A\in\tilde{G}_k$ and every $y\in Y_k$;
the cutoff $N$, the torus and the charge $e$ enter the objects of Proposition~10 only through this quadruple (Section~I.1.8).

\subsubsection*{2.1. Spaces, pairings, transposes, conjugates, adjoints}

The spinor spaces are $\mathfrak{S}^{\mathrm{ch}}(\eta_k,\ell)$, $\mathfrak{S}(1,\ell)$ and $\mathfrak{S}(L,\ell)$ (Section~I.1.5).
They carry the bilinear pairings of Section~I.1.6, in which the sums extend over a fundamental domain of $\ell\mathbb{Z}^3$:
\begin{gather*}
\langle f',f\rangle^{\mathrm{ch}}=\eta_k^3\cdot\frac{1}{48}\sum_{c}\sum_{g\in O_{\mathrm{cub}}}f'(c,g)\cdot f(c,g),\\
\langle v',v\rangle_1=\sum_{z}v'(z)\cdot v(z),\qquad
\langle v',v\rangle_L=L^3\sum_{z'}v'(z')\cdot v(z').
\end{gather*}
Each of these pairings is symmetric and non-degenerate (the Gram matrix is computed below).
For a spinor space $\mathrm{Sv}$ the pairing that belongs to it is written $\langle\,\cdot\,,\cdot\,\rangle_{\mathrm{Sv}}$.
An \emph{operator} is a $\mathbb{C}$-linear map between two of the three spinor spaces, or from $\mathfrak{S}^{\mathrm{ch}}(\eta_k,\ell)$ to the functions of pairs (Section~I.1.7).
In a relation between operators, a symbol that acts on several spaces (a constant matrix, $\mathrm{uf}$, $\tilde{g}$, $\mathrm{glb}$, $\mathrm{Cx}$, $\mathrm{Kf}$) denotes its action on the space on which it stands.

For an operator $O\colon\mathrm{Sv}_1\to\mathrm{Sv}_2$ between spinor spaces, the transpose $O^T\colon\mathrm{Sv}_2\to\mathrm{Sv}_1$, the conjugate $\overline{O}$ and the adjoint $O^{*}\colon\mathrm{Sv}_2\to\mathrm{Sv}_1$ are defined by
\begin{gather*}
\langle O^Tf_2,f_1\rangle_{\mathrm{Sv}_1}=\langle f_2,Of_1\rangle_{\mathrm{Sv}_2}\qquad(f_1\in\mathrm{Sv}_1,\ f_2\in\mathrm{Sv}_2);\\
(\mathrm{Cx}\,f)(t):=\overline{f(t)},\qquad\overline{O}:=\mathrm{Cx}\cdot O\cdot\mathrm{Cx};\\
(f',f)_{\mathrm{Sv}}:=\langle\mathrm{Cx}\,f',f\rangle_{\mathrm{Sv}},\qquad(O^{*}f_2,f_1)_{\mathrm{Sv}_1}=(f_2,Of_1)_{\mathrm{Sv}_2}.
\end{gather*}
Here $\overline{O}$ is the map $f\mapsto\overline{O\,\overline{f}}$ of (I.O5), and $(\,\cdot\,,\cdot\,)_{\mathrm{Sv}}$ are the sesquilinear pairings $(\,\cdot\,,\cdot\,)^{\mathrm{ch}}$, $(\,\cdot\,,\cdot\,)_1$, $(\,\cdot\,,\cdot\,)_L$ of Section~I.1.6.
The complex conjugation $\mathrm{Cx}$ of functions maps each of the spaces onto itself, the space of the functions of pairs included, because the conditions $f(t+\ell e_\mu)=-f(t)$ have real coefficients.
It is antilinear, and $\mathrm{Cx}\cdot\mathrm{Cx}=1$.
For an operator $O$ with values in the functions of pairs one also puts $\overline{O}:=\mathrm{Cx}\cdot O\cdot\mathrm{Cx}$.

\emph{Bases; existence and uniqueness of $O^T$ and $O^{*}$.}
The three spinor spaces have finite dimension (Section~I.1.5).
For a spinor space $\mathrm{Sv}$, let $(\mathrm{bs}_j)$ be the basis formed by the functions that take a unit vector of $\mathbb{C}^4$ as value at one point of a fundamental domain, vanish at the other points of the fundamental domain, and are continued antiperiodically.
Then $\mathrm{Cx}\,\mathrm{bs}_j=\mathrm{bs}_j$, and
\[
\langle\mathrm{bs}_i,\mathrm{bs}_j\rangle_{\mathrm{Sv}}=\nu_{\mathrm{Sv}}\cdot\mathbf{1}[i=j]
\]
with one number $\nu_{\mathrm{Sv}}>0$, namely $\nu_{\mathrm{Sv}}=\frac{\eta_k^3}{48}$, $1$, $L^3$ for the three spaces.
Hence $\langle\,\cdot\,,\cdot\,\rangle_{\mathrm{Sv}}$ is non-degenerate, and so is $(\,\cdot\,,\cdot\,)_{\mathrm{Sv}}=\langle\mathrm{Cx}\,\cdot\,,\cdot\,\rangle_{\mathrm{Sv}}$.
Consequently the relation that defines $O^T$ (resp.\ $O^{*}$) determines the element $O^Tf_2$ (resp.\ $O^{*}f_2$) of $\mathrm{Sv}_1$ for every $f_2\in\mathrm{Sv}_2$, and there is at most one $O^T$ and at most one $O^{*}$.
Let $(\mathrm{bs}_j)$ and $(\mathrm{bs}'_i)$ be the bases of $\mathrm{Sv}_1$ and of $\mathrm{Sv}_2$, and write $O\,\mathrm{bs}_j=\sum_i o_{ij}\,\mathrm{bs}'_i$.
The linear map $\mathrm{Sv}_2\to\mathrm{Sv}_1$ that sends $\mathrm{bs}'_i$ to $\frac{\nu_{\mathrm{Sv}_2}}{\nu_{\mathrm{Sv}_1}}\sum_j o_{ij}\,\mathrm{bs}_j$ satisfies the relation that defines $O^T$ at $f_2=\mathrm{bs}'_i$, $f_1=\mathrm{bs}_j$, since both sides are then equal to $\nu_{\mathrm{Sv}_2}\cdot o_{ij}$.
Both sides of this relation are bilinear in $(f_2,f_1)$; so it holds for all $f_1$ and $f_2$.
Therefore $O^T$ exists and is linear, and in the bases $(\mathrm{bs}'_i)$, $(\mathrm{bs}_j)$ its matrix is $\frac{\nu_{\mathrm{Sv}_2}}{\nu_{\mathrm{Sv}_1}}$ times the transposed matrix of $O$.
The linear map $\mathrm{Cx}\cdot O^T\cdot\mathrm{Cx}$ satisfies the relation that defines $O^{*}$ (the computation is given in the proof of (Ru1)); hence $O^{*}$ exists and is linear.

\begin{named}{(Ru1)}
The three operations are additive.
For operators $O_1$, $O_2$ that can be composed and a complex number $\mathrm{cf}$,
\begin{gather*}
\left(O_1\cdot O_2\right)^T=O_2^T\cdot O_1^T,\qquad\left(O^T\right)^T=O,\qquad\left(\mathrm{cf}\cdot O\right)^T=\mathrm{cf}\cdot O^T;\\
\left(O_1\cdot O_2\right)^{*}=O_2^{*}\cdot O_1^{*},\qquad\left(O^{*}\right)^{*}=O,\qquad\left(\mathrm{cf}\cdot O\right)^{*}=\overline{\mathrm{cf}}\cdot O^{*};\\
\overline{O_1\cdot O_2}=\overline{O_1}\cdot\overline{O_2},\qquad\overline{\mathrm{cf}\cdot O}=\overline{\mathrm{cf}}\cdot\overline{O};\qquad O^{*}=\overline{O^T}.
\end{gather*}
If $O$ is bijective, then
\[
\left(O^{-1}\right)^T=\left(O^T\right)^{-1},\qquad\left(O^{-1}\right)^{*}=\left(O^{*}\right)^{-1},\qquad\overline{O^{-1}}=\left(\overline{O}\right)^{-1}.
\]
\end{named}

\begin{proof}
The relation that defines $O^T$ is linear in $O$, and $O^T$ is unique; this gives the additivity of the transposition and $\left(\mathrm{cf}\cdot O\right)^T=\mathrm{cf}\cdot O^T$.
The conjugation is additive because $\mathrm{Cx}$ is additive.
Next,
\[
\left\langle\left(O_1\cdot O_2\right)^Tf,f'\right\rangle=\left\langle f,O_1\cdot O_2\,f'\right\rangle=\left\langle O_1^Tf,O_2\,f'\right\rangle=\left\langle O_2^T\cdot O_1^Tf,f'\right\rangle,
\]
and $\left(O^T\right)^T=O$ because the pairings are symmetric.
Since $1^T=1$, one has $O^T\cdot\left(O^{-1}\right)^T=\left(O^{-1}\cdot O\right)^T=1$ and $\left(O^{-1}\right)^T\cdot O^T=\left(O\cdot O^{-1}\right)^T=1$.
Further, $\mathrm{Cx}\cdot O_1\cdot O_2\cdot\mathrm{Cx}=\left(\mathrm{Cx}\cdot O_1\cdot\mathrm{Cx}\right)\cdot\left(\mathrm{Cx}\cdot O_2\cdot\mathrm{Cx}\right)$ because $\mathrm{Cx}\cdot\mathrm{Cx}=1$; moreover $\mathrm{Cx}\cdot\mathrm{cf}=\overline{\mathrm{cf}}\cdot\mathrm{Cx}$, and $\overline{O}\cdot\overline{O^{-1}}=\overline{O\cdot O^{-1}}=1$.
The identity $O^{*}=\overline{O^T}$ follows from
\begin{align*}
(f_2,Of_1)_{\mathrm{Sv}_2}&=\langle\mathrm{Cx}\,f_2,Of_1\rangle_{\mathrm{Sv}_2}=\langle O^T\,\mathrm{Cx}\,f_2,f_1\rangle_{\mathrm{Sv}_1}\\
&=\left\langle\mathrm{Cx}\left(\mathrm{Cx}\cdot O^T\cdot\mathrm{Cx}\right)f_2,f_1\right\rangle_{\mathrm{Sv}_1}=\left(\overline{O^T}f_2,f_1\right)_{\mathrm{Sv}_1}.
\end{align*}
The rules for the adjoint follow from those for the transposition and for the conjugation.
First, $\left(O_1\cdot O_2\right)^{*}=\overline{O_2^T\cdot O_1^T}=O_2^{*}\cdot O_1^{*}$.
Second, $\overline{O^T}=\left(\overline{O}\right)^T$: the identity $\langle\mathrm{Cx}\,f',\mathrm{Cx}\,f\rangle_{\mathrm{Sv}}=\overline{\langle f',f\rangle_{\mathrm{Sv}}}$ gives
\begin{multline*}
\left\langle\left(\overline{O}\right)^Tf_2,f_1\right\rangle_{\mathrm{Sv}_1}=\left\langle f_2,\mathrm{Cx}\cdot O\cdot\mathrm{Cx}\,f_1\right\rangle_{\mathrm{Sv}_2}=\overline{\left\langle\mathrm{Cx}\,f_2,O\cdot\mathrm{Cx}\,f_1\right\rangle_{\mathrm{Sv}_2}}\\
=\overline{\left\langle O^T\cdot\mathrm{Cx}\,f_2,\mathrm{Cx}\,f_1\right\rangle_{\mathrm{Sv}_1}}=\left\langle\mathrm{Cx}\cdot O^T\cdot\mathrm{Cx}\,f_2,f_1\right\rangle_{\mathrm{Sv}_1}.
\end{multline*}
Hence
\begin{gather*}
\left(O^{*}\right)^{*}=\overline{\left(\overline{O^T}\right)^T}=\overline{\overline{\left(O^T\right)^T}}=O,\qquad\left(\mathrm{cf}\cdot O\right)^{*}=\overline{\mathrm{cf}\cdot O^T}=\overline{\mathrm{cf}}\cdot O^{*},\\
\left(O^{-1}\right)^{*}=\overline{\left(O^T\right)^{-1}}=\left(O^{*}\right)^{-1}.
\end{gather*}
\end{proof}

\begin{named}{(Ru2)}
A constant $4\times4$ matrix $\mathrm{Mx}$ acts on each of the spaces by $(\mathrm{Mx}\,f)(t):=\mathrm{Mx}\cdot f(t)$.
As an operator on a spinor space, $\mathrm{Mx}$ has the transposed matrix $\mathrm{Mx}^T$ as its transpose; $\overline{\mathrm{Mx}}$ is the complex conjugate matrix, and $\mathrm{Mx}^{*}$ is the Hermitian conjugate matrix.
\end{named}

\begin{proof}
Each pairing is a sum of terms
\[
(\text{weight})\cdot f'(t)\cdot\left(\mathrm{Mx}\cdot f(t)\right)=(\text{weight})\cdot\left(\mathrm{Mx}^T\cdot f'(t)\right)\cdot f(t),
\]
the weight being $\frac{\eta_k^3}{48}$, $1$ or $L^3$.
Further, $\mathrm{Cx}\cdot\mathrm{Mx}\cdot\mathrm{Cx}\,f=\overline{\mathrm{Mx}\cdot\overline{f}}=\overline{\mathrm{Mx}}\cdot f$.
Finally, $\mathrm{Mx}^{*}=\overline{\mathrm{Mx}^T}$ by (Ru1).
\end{proof}

\begin{named}{(Ru3)}
An $\ell\mathbb{Z}^3$-periodic function $\psi$ of $c$ acts on $\mathfrak{S}^{\mathrm{ch}}(\eta_k,\ell)$ by $(\psi f)(c,g)=\psi(c)\cdot f(c,g)$, with the same factor at the $48$ sites over $c$ (I.POS); on $\mathfrak{S}(1,\ell)$ and on $\mathfrak{S}(L,\ell)$ through its restrictions to $\mathbb{Z}^3$ and to $L\mathbb{Z}^3$; and on functions of pairs by $(\psi F)(x,x')=\psi(c)\cdot F(x,x')$ (Section~G1.3).
It maps each space into itself and commutes with the constant matrices.
If $\psi$ has no zero, it is bijective, with the inverse $\psi^{-1}$.
\end{named}

\begin{named}{(Ru4)}[Determinants]
Let $O$ be an operator of a spinor space $\mathrm{Sv}$ into itself.
Then:
\begin{itemize}
\item[(i)] $\det\left(\mathrm{Vm}\cdot O\cdot\mathrm{Vm}^{-1}\right)=\det O$ for every $\mathbb{C}$-linear bijection $\mathrm{Vm}$ of $\mathrm{Sv}$;
\item[(ii)] $\det\overline{O}=\overline{\det O}$;
\item[(iii)] $\det\left(O^T\right)=\det O$.
\end{itemize}
\end{named}

\begin{proof}
(i) is the multiplicativity of the determinant.
Let $(\mathrm{bs}_j)$ be the basis of $\mathrm{Sv}$ defined above under ``Bases'', and write $O\,\mathrm{bs}_j=\sum_i o_{ij}\,\mathrm{bs}_i$.
Since $\mathrm{Cx}\,\mathrm{bs}_j=\mathrm{bs}_j$, one has $\overline{O}\,\mathrm{bs}_j=\mathrm{Cx}\cdot O\,\mathrm{bs}_j=\sum_i\overline{o_{ij}}\,\mathrm{bs}_i$; this is (ii).
By the same paragraph with $\mathrm{Sv}_1=\mathrm{Sv}_2=\mathrm{Sv}$, where the quotient of the two numbers $\nu$ is $1$, the matrix of $O^T$ in this basis is the transposed matrix of $O$; this is (iii).
\end{proof}

\begin{named}{(Ru5)}[Inverses of intertwined maps]
Let $O$, $\tilde{O}$ and $\mathrm{Vm}$ be bijective maps of a set onto itself with $\tilde{O}\cdot\mathrm{Vm}=\mathrm{Vm}\cdot O$.
Then $\tilde{O}^{-1}\cdot\mathrm{Vm}=\mathrm{Vm}\cdot O^{-1}$.
\end{named}

\begin{proof}
Multiply the relation $\tilde{O}\cdot\mathrm{Vm}=\mathrm{Vm}\cdot O$ by $\tilde{O}^{-1}$ from the left and by $O^{-1}$ from the right.
\end{proof}

\subsubsection*{2.2. Operators with scalar coefficients}

An operator $O$ has \emph{scalar coefficients} if
\[
(Of)(t)_\beta=\sum_{t'}o(t,t')\cdot f(t')_\beta\qquad(\beta=1,\dots,4)
\]
with numbers $o(t,t')$ that do not depend on $\beta$.

\begin{named}{(Sca)}
An operator with scalar coefficients commutes with every constant matrix: $O\cdot\mathrm{Mx}=\mathrm{Mx}\cdot O$, where $\mathrm{Mx}$ on the right acts on the source and $\mathrm{Mx}$ on the left acts on the target.
Sums and products of operators with scalar coefficients have scalar coefficients.
\end{named}

By the displays (I.3.1), those of Section~I.3.3 and (I.1.7), the following operators have scalar coefficients for every complex connection $A$:
\begin{gather*}
\left(Q_k(A)\cdot J^{*}f\right)(z)=\int_{\Delta_z}\mathrm{d}x\,\varphi_A(z,c)\cdot f(x)=\eta_k^3\cdot\frac{1}{48}\sum_{c\in\Delta_z}\sum_{g\in O_{\mathrm{cub}}}\varphi_A(z,c)\cdot f(c,g),\\
\left(J\cdot Q_k^{T}(-A)\,v\right)(c,g)=\varphi_A(z,c)^{-1}\cdot v(z)\qquad\left(c\in\Delta_z,\ \text{all }48\text{ labels }g\right),\\
\left(Q(A)\,v\right)(z')=L^{-3}\sum_{z\in B(z')}U_A(z',z)\cdot v(z),\\
\left(Q^{T}(-A)\,v'\right)(z)=U_A(z',z)^{-1}\cdot v'(z')\qquad\left(z\in B(z')\right),\\
\left(\delta_{\alpha,A}f\right)(x,x')=d_X(x,x')^{-\alpha}\cdot\left(U_A(x,x')\cdot f(x')-f(x)\right)\qquad\left((x,x')\in\mathrm{Pzt}_k\right),
\end{gather*}
with $U_A(x,x')=1$ for a pair inside a cluster; the multiplication by a function $\psi$ of $c$ has scalar coefficients too.
Hence the products (definitions of Section~I.3.3)
\begin{gather*}
J\cdot P_k(A)\cdot J^{*}=\left(J\cdot Q_k^{T}(-A)\right)\cdot\left(Q_k(A)\cdot J^{*}\right),\qquad P(A)=Q^{T}(-A)\cdot Q(A),\\
J\cdot P_{k+1}(A)\cdot J^{*}=\left(J\cdot Q_k^{T}(-A)\right)\cdot P(A)\cdot\left(Q_k(A)\cdot J^{*}\right)
\end{gather*}
have scalar coefficients as well, and so do, at $k=0$ (by (I.3.4)), the operators
\begin{gather*}
\left(Q(A)\cdot J^{*}f\right)(z')=L^{-3}\sum_{z\in B(z')}U_A(z',z)\cdot\frac{1}{48}\sum_{g\in O_{\mathrm{cub}}}f(z,g),\\
\left(J\cdot Q^{T}(-A)\,v'\right)(z,g)=U_A(z',z)^{-1}\cdot v'(z')\qquad\left(z\in B(z'),\ \text{all }48\text{ labels }g\right),\\
J\cdot P(A)\cdot J^{*}=\left(J\cdot Q^{T}(-A)\right)\cdot\left(Q(A)\cdot J^{*}\right).
\end{gather*}
The operators $D_A$ and $F_y(A)$ do not have scalar coefficients: they contain the matrices $M_i(x)$ (see (I.2.3)).
For neutral $y$ (Section~G1.9) the operator $\operatorname{Pen}_y(A)$ has scalar coefficients, being a real multiple of one of the products above:
\[
\operatorname{Pen}_{\infty}(A)=b_k\cdot J\cdot P_k(A)\cdot J^{*},\qquad\operatorname{Pen}_0(A)=\frac{b_{k+1}}{L}\cdot J\cdot P_{k+1}(A)\cdot J^{*}\qquad(k\ge1),
\]
since $\boldsymbol{\alpha}_{\infty}=b_k$, $\boldsymbol{\beta}_{\infty}=0$, $\boldsymbol{\alpha}_0=0$, $\boldsymbol{\beta}_0=\frac{b_{k+1}}{L}$ (Section~I.3.5), and
\[
\operatorname{Pen}_0(A)=b_L\cdot J\cdot P(A)\cdot J^{*}\qquad(k=0;\ \text{Section~I.3.9}).
\]
For the other $y$ the operator $\operatorname{Pen}_y(A)$ contains $\gamma_5$.

\subsubsection*{2.3. The maps $\tilde{g}$, $\mathrm{glb}$, $\mathrm{Kf}$; conjugates of phases}

\emph{The map $\tilde{g}$.}
Let $\tilde{g}=(S,a)\in\mathbb{G}_k$ and $r:=r(S)=\operatorname{diag}(r_1,r_2,r_3)$ (Section~I.4.1).
As in Section~I.4.3, $a$ is taken in $M\mathbb{Z}^3$, with the same representative on all spaces.
The map acts by
\begin{gather*}
(\tilde{g}f)(x)=S\cdot f\left(\tilde{g}^{-1}x\right),\qquad\tilde{g}^{-1}(c,g)=\left(r(c-a),rg\right);\\
(\tilde{g}v)(z)=S\cdot v\left(r(z-a)\right)\quad\text{on }\mathfrak{S}(1,\ell)\text{ and on }\mathfrak{S}(L,\ell);\\
(\tilde{g}F)(x,x')=S\cdot F\left(\tilde{g}^{-1}x,\tilde{g}^{-1}x'\right)\quad\text{on the functions of pairs};
\end{gather*}
on $\mathfrak{S}(L,\ell)$ this formula makes sense because $M\mathbb{Z}^3$ is contained in $L\mathbb{Z}^3$, $M$ being a power of $L$ (Section~I.0.5).
The map $\tilde{g}$ sends each space into itself.
Indeed,
\[
(\tilde{g}f)(c+\ell e_\mu,g)=S\cdot f\left(r(c-a)+\ell r_\mu e_\mu,rg\right)=-(\tilde{g}f)(c,g),
\]
because $f$ changes its sign under the shift by $+\ell e_\mu$ and under the shift by $-\ell e_\mu$; in the same way $(\tilde{g}v)(z+\ell e_\mu)=S\cdot v\left(r(z-a)+\ell r_\mu e_\mu\right)=-(\tilde{g}v)(z)$ on $\mathfrak{S}(1,\ell)$ and on $\mathfrak{S}(L,\ell)$, and on the functions of pairs (Section~I.1.7)
\[
(\tilde{g}F)(x+\ell e_\mu,x'+\ell e_\mu)=S\cdot F\left(\tilde{g}^{-1}x+\ell r_\mu e_\mu,\tilde{g}^{-1}x'+\ell r_\mu e_\mu\right)=-(\tilde{g}F)(x,x').
\]
The map $\tilde{g}$ is $\mathbb{C}$-linear and bijective.
In fact, $x\mapsto\tilde{g}x$ is a bijection of the sites, and of $\mathrm{Pzt}_k$, with the inverse $x\mapsto\tilde{g}^{-1}x$ ((I.a1)); the map $f\mapsto S^{-1}\cdot f(\tilde{g}\,\cdot\,)$ sends each space into itself, because
\[
S^{-1}\cdot f\left(\tilde{g}(c+\ell e_\mu,g)\right)=S^{-1}\cdot f\left(\tilde{g}(c,g)+\ell r_\mu e_\mu\right)=-S^{-1}\cdot f\left(\tilde{g}(c,g)\right),
\]
and in the same way on the other spaces; and this map is the inverse of $\tilde{g}$.

\emph{The representative of $a$.}
If $a$ is replaced by $a+\ell n$ with $n\in\mathbb{Z}^3$, then $\tilde{g}x$ is replaced by $\tilde{g}x+\ell n$ and $\tilde{g}^{-1}x$ by $\tilde{g}^{-1}x-\ell rn$, and $rn$ has the entries $\pm n_\mu$.
Hence $\tilde{g}$, and $\mathrm{glb}$ (defined below), are multiplied on every space by $(-1)^{n_1+n_2+n_3}$, whereas $\tilde{g}A$ does not change, $A$ being $\ell\mathbb{Z}^3$-periodic.
In every relation of Proposition~10\,(b) between operators each of the two sides contains the map $\tilde{g}$ (resp.\ $\mathrm{glb}$) exactly once as a factor, and contains $\tilde{g}$ otherwise only through $\tilde{g}A$; here the last relation of Proposition~10\,(b) is taken in the form $\delta_{\alpha,\tilde{g}A}\cdot\tilde{g}=\tilde{g}\cdot\delta_{\alpha,A}$.
In the form in which that relation is stated in Proposition~10\,(b),
\[
\left(\delta_{\alpha,\tilde{g}A}\,\tilde{g}f\right)(\tilde{g}x,\tilde{g}x')=S\cdot\left(\delta_{\alpha,A}f\right)(x,x'),
\]
the right side does not contain $\tilde{g}$, and on the left side the factor $(-1)^{n_1+n_2+n_3}$ of $\tilde{g}f$ is cancelled by the factor $(-1)^{n_1+n_2+n_3}$ that the function of pairs $\delta_{\alpha,\tilde{g}A}\,\tilde{g}f$ acquires under the shift of $(\tilde{g}x,\tilde{g}x')$ by $\ell n$ (Section~I.1.7).
The relation $\delta Z_k(\tilde{g}A)=\delta Z_k(A)$ contains $\tilde{g}$ only through $\tilde{g}A$.
So the relations of Proposition~10\,(b) do not depend on the representative.
(The proofs in 5.2 and 6.2 hold for every representative $a\in M\mathbb{Z}^3$.)

By the last paragraph of Section~G1.2, $T_k$ is the block torus $(0,\ell,-1)$, the map $c\mapsto rc+a$ is a morphism of $T_k$ onto itself, and $\tilde{g}$ is the map $\mathrm{mS}$ of Proposition~2.2\,(b).

The map $A\mapsto\tilde{g}A$, $(\tilde{g}A)(c,c')=A\left(\tilde{g}^{-1}c,\tilde{g}^{-1}c'\right)$, is $\mathbb{C}$-linear and commutes with the complex conjugation:
\[
\tilde{g}(-A)=-\tilde{g}A,\qquad\tilde{g}\left(\overline{A}\right)=\overline{\tilde{g}A}.
\]

\emph{The map $\mathrm{glb}$} (Section~G1.10):
\[
(\mathrm{glb}\,f)(x):=\left(S^{-1}\right)^T\cdot f\left(\tilde{g}^{-1}x\right).
\]

\begin{named}{(Sn1)}
For $S\in\mathrm{Sp}$ there is a sign $\mathrm{sgn}(S)\in\{1,-1\}$ with $\left(S^{-1}\right)^T=\mathrm{sgn}(S)\cdot\mathsf{C}^{-1}\cdot S\cdot\mathsf{C}$ (Section~I.4.1).
Hence, on each of the spaces,
\[
\mathrm{glb}=\mathrm{sgn}(S)\cdot\mathsf{C}^{-1}\cdot\tilde{g}\cdot\mathsf{C}.
\]
\end{named}

\begin{proof}
By (I.c1) and (I.c3), each $S'$ among the four generators $\gamma_1$, $\gamma_2$, $\gamma_3$, $\gamma_5$ satisfies $S'^{-1}=S'$ and $S'^T=-\mathsf{C}^{-1}\cdot S'\cdot\mathsf{C}$, so that $\left(S'^{-1}\right)^T=-\mathsf{C}^{-1}\cdot S'\cdot\mathsf{C}$.
The maps $S\mapsto\left(S^{-1}\right)^T$ and $S\mapsto\mathsf{C}^{-1}\cdot S\cdot\mathsf{C}$ are multiplicative on the invertible matrices and fix $-1$.
Therefore, for $S=\pm S'_1\cdots S'_j$, a product of $j$ generators, $\left(S^{-1}\right)^T=(-1)^j\cdot\mathsf{C}^{-1}\cdot S\cdot\mathsf{C}$.
For the second assertion note that $\left(\mathsf{C}^{-1}\cdot\tilde{g}\cdot\mathsf{C}\,f\right)(x)=\mathsf{C}^{-1}\cdot S\cdot\mathsf{C}\cdot f\left(\tilde{g}^{-1}x\right)$.
\end{proof}

\begin{named}{(Sn2)}
One has $S\cdot\gamma_5\cdot S^{-1}=\gamma_5$ for $S\in\mathrm{Sp}_3$, and $S\cdot\gamma_5\cdot S^{-1}=-\gamma_5$ for $S\in\mathrm{Sp}\setminus\mathrm{Sp}_3$ (Section~I.4.1).
\end{named}

\emph{The map $\mathrm{Kf}$.}
On each of the spaces let $\mathrm{Kf}:=\varkappa\cdot\mathrm{Cx}$, that is, $(\mathrm{Kf}\,f)(t)=\varkappa\cdot\overline{f(t)}$.
The map $\mathrm{Kf}$ is additive, antilinear and bijective, with $\mathrm{Kf}^{-1}=\mathrm{Cx}\cdot\varkappa^{-1}$, and it commutes with the multiplication by real numbers.

\begin{named}{(Kf1)}
For operators $O_1$, $O_2$ with the same source and the same target,
\[
\overline{O_1}=\varkappa^{-1}\cdot O_2\cdot\varkappa\qquad\text{if and only if}\qquad O_2\cdot\mathrm{Kf}=\mathrm{Kf}\cdot O_1.
\]
\end{named}

\begin{proof}
By definition $\overline{O_1}=\mathrm{Cx}\cdot O_1\cdot\mathrm{Cx}$.
Multiplying $\mathrm{Cx}\cdot O_1\cdot\mathrm{Cx}=\varkappa^{-1}\cdot O_2\cdot\varkappa$ by $\varkappa$ from the left and by $\mathrm{Cx}$ from the right gives $\varkappa\cdot\mathrm{Cx}\cdot O_1=O_2\cdot\varkappa\cdot\mathrm{Cx}$.
Both multiplications are by bijections, so the step can be reversed.
\end{proof}

\begin{named}{(Kf2)}
$\left(i\gamma_5y\right)\cdot\mathrm{Kf}=\mathrm{Kf}\cdot\left(i\gamma_5y\right)$ for real $y$.
\end{named}

\begin{proof}
One has
\begin{align*}
\mathrm{Kf}\left(i\gamma_5y\,f\right)&=\varkappa\cdot\overline{i}\cdot\overline{\gamma_5}\cdot y\cdot\overline{f}=-iy\cdot\left(\varkappa\cdot\overline{\gamma_5}\right)\cdot\overline{f}\\
&=-iy\cdot\left(-\gamma_5\cdot\varkappa\right)\cdot\overline{f}=i\gamma_5y\cdot\mathrm{Kf}\,f;
\end{align*}
here $\varkappa\cdot\overline{\gamma_5}=-\gamma_5\cdot\varkappa$ is (I.c4), $\overline{\gamma_5}=-\varkappa^{-1}\cdot\gamma_5\cdot\varkappa$.
\end{proof}

\begin{named}{(Kf3)}
$\det\left(\mathrm{Kf}\cdot O\cdot\mathrm{Kf}^{-1}\right)=\overline{\det O}$ for every operator $O$ of a spinor space into itself.
\end{named}

\begin{proof}
One has $\mathrm{Kf}\cdot O\cdot\mathrm{Kf}^{-1}=\varkappa\cdot\left(\mathrm{Cx}\cdot O\cdot\mathrm{Cx}\right)\cdot\varkappa^{-1}=\varkappa\cdot\overline{O}\cdot\varkappa^{-1}$; now apply (Ru4)\,(i) and (ii).
\end{proof}

\begin{named}{(Kf4)}
For every operator $O$ between spinor spaces,
\[
\varkappa^{-1}\cdot\mathsf{C}^{-1}\cdot O\cdot\mathsf{C}\cdot\varkappa=\mathsf{C}^{-1}\cdot\varkappa^{-1}\cdot O\cdot\varkappa\cdot\mathsf{C}.
\]
\end{named}

\begin{proof}
By (I.c4), $\varkappa\cdot\mathsf{C}=-\mathsf{C}\cdot\varkappa$.
So $\mathsf{C}\cdot\varkappa=-\varkappa\cdot\mathsf{C}$ and $\varkappa^{-1}\cdot\mathsf{C}^{-1}=\left(\mathsf{C}\cdot\varkappa\right)^{-1}=-\mathsf{C}^{-1}\cdot\varkappa^{-1}$; the two signs cancel.
\end{proof}

\emph{Conjugates of phases.}
For every complex connection $A$, all sites for which $U$ is defined, and $c\in\Delta_z$,
\begin{equation}
\overline{U_{\overline{A}}}=U_{-A},\qquad\overline{\varphi_{\overline{A}}(z,c)}=\varphi_{-A}(z,c).
\tag{2.1}
\end{equation}

\begin{proof}
The maps $\tau$ and $\tau_k$ are linear with real coefficients (Section~I.1.1 and (I.3.3)), and $e_k$, $\eta_k$ are real; so
\[
\overline{\exp\left\{ie_k\eta_k\cdot\left(\tau\overline{A}\right)\right\}}=\exp\left\{-ie_k\eta_k\cdot(\tau A)\right\}.
\]
\end{proof}

\subsection*{3. The domain; existence of the inverses; the composed operators as words}

\subsubsection*{3.1. Closed real connections}

\begin{named}{Lemma 1}
Let $k\ge0$, and let $\mathrm{Acl}$ be a real connection on $T_k$ with $d\,\mathrm{Acl}=0$.
\begin{itemize}
\item[(i)] For every $M$-cube $\square$ there is a real function $\mathrm{lamc}$ on the sites of $\square^{\sim 5}$ with $\mathrm{Acl}=d\,\mathrm{lamc}$ on the bonds of $\square^{\sim 5}$, that is, on the bonds contained in $\square^{\sim 5}$, those with both end points in $\square^{\sim 5}$ (Section~I.1.1).
\item[(ii)] If $A\in\tilde{G}_k$, then $A-\mathrm{Acl}\in\tilde{G}_k$.
\end{itemize}
\end{named}

\begin{proof}
(i) Let $z_{\square}\in M\mathbb{Z}^3$ be a lift of the centre of $\square$, and put
\[
\mathrm{Wb}:=\left\{c\in\mathbb{L}_k:\ |c-z_{\square}|<\frac{11M}{2}\right\}.
\]
The set $\square^{\sim 5}$, which is $\square$ enlarged by five layers of $M$-cubes (Section~I.1.1), is the image of $\mathrm{Wb}$ in $T_k$.
The number $\frac{11M}{\eta_k}=11M\cdot L^k$ is odd, because $L$ is odd and $M$ is a power of $L$ (Section~I.0.5).
Hence $\mathrm{Wb}$ is the product of the three intervals
\[
\left\{c_\mu\in\eta_k\mathbb{Z}:\ |c_\mu-z_{\square,\mu}|\le\frac{11M-\eta_k}{2}\right\},\qquad\mu=1,2,3,
\]
each of length $11M-\eta_k$, and
\[
\text{sup-diameter}(\mathrm{Wb})+\eta_k=11M<13L\cdot M\le\ell\qquad(\ell\in\mathrm{Cl}).
\]
So $\mathrm{Wb}$ is a lift of $\square^{\sim 5}$ in the sense of Section~I.1.1: the covering map is injective on $\mathrm{Wb}$, and the bonds of $\square^{\sim 5}$ are the images of the bonds of $\mathrm{Wb}$.
Let $\mathrm{Acl}$ also denote the $\ell\mathbb{Z}^3$-periodic extension of $\mathrm{Acl}$ to the bonds of $\mathbb{L}_k$.
This extension is real and antisymmetric, and $(d\,\mathrm{Acl})(p)=0$ for every plaquette $p$ of $\mathbb{L}_k$, since $d$ is formed with differences along the bonds of $p$ and therefore commutes with the periodic extension; in particular this holds for the plaquettes with all corners in $\mathrm{Wb}$.
Now apply \cite[Lemma~16.26]{QEDI} to the restriction of $\mathrm{Acl}$ to the bonds of $\mathrm{Wb}$, with $\eta=\eta_k$, the three intervals of $\eta_k\mathbb{Z}$ above, $s_W:=11M$, the fixed site $o:=z_{\square}$, and, for $c\in\mathrm{Wb}$, the rectilinear path $\boldsymbol{\Gamma}(z_{\square},c)$ that brings the coordinates to their final values in the order $1,2,3$ (in the numbering of the directions used in that lemma, the order $0,1,2$).
The lemma says that the function
\[
\boldsymbol{\lambda}(c):=\eta_k\cdot\mathrm{Acl}\left(\boldsymbol{\Gamma}(z_{\square},c)\right)\qquad(c\in\mathrm{Wb})
\]
is real and that
\[
\left|\left(\mathrm{Acl}-d\boldsymbol{\lambda}\right)(b)\right|\le2s_W\cdot\sup_p\left|(d\,\mathrm{Acl})(p)\right|=0\qquad\text{for every bond }b\text{ of }\mathrm{Wb}.
\]
Let $\mathrm{lamc}$ be the function $\boldsymbol{\lambda}$ transported to $\square^{\sim 5}$ by the covering map.
A bond $(c,c')$ of $\square^{\sim 5}$ is the image of a bond $(\tilde{c},\tilde{c}')$ of $\mathrm{Wb}$, and therefore
\begin{align*}
(d\,\mathrm{lamc})(c,c')&=\frac{\mathrm{lamc}(c')-\mathrm{lamc}(c)}{\eta_k}=\frac{\boldsymbol{\lambda}(\tilde{c}')-\boldsymbol{\lambda}(\tilde{c})}{\eta_k}\\
&=(d\boldsymbol{\lambda})(\tilde{c},\tilde{c}')=\mathrm{Acl}(\tilde{c},\tilde{c}')=\mathrm{Acl}(c,c').
\end{align*}

(ii) By Section~I.5.3, $A=A_0+A_1$ with $A_1\in G_k$ and $A_0$ real, and for every $M$-cube $\square$ one has $A_0=A'+d\boldsymbol{\lambda}$ on $\square^{\sim 5}$ with real $\boldsymbol{\lambda}$ and $A'$ and with $|A'|_{\infty}<e_k^{-1+\varepsilon}$ on $\square^{\sim 5}$.
Here $\boldsymbol{\lambda}$ is a function on the sites of $\square^{\sim 5}$, $A'$ is a function on the bonds of $\square^{\sim 5}$, and the relation holds on these bonds: as in \cite[Definition~5.16]{QEDI}, the real functions $\boldsymbol{\lambda}$ and $A'$ of such a representation are defined on the enlarged cube only.
Then $A-\mathrm{Acl}=(A_0-\mathrm{Acl})+A_1$, the bond function $A_0-\mathrm{Acl}$ is real, and by (i)
\[
A_0-\mathrm{Acl}=A'+d\left(\boldsymbol{\lambda}-\mathrm{lamc}\right)\qquad\text{on }\square^{\sim 5},
\]
with the same $A'$ and the real function $\boldsymbol{\lambda}-\mathrm{lamc}$.
These are the conditions of Section~I.5.3 for $A-\mathrm{Acl}$.
\end{proof}

\subsubsection*{3.2. Invariances of the domain}

\begin{named}{(Dm)}
Let $A\in\tilde{G}_k$ and $\tilde{g}\in\mathbb{G}_k$, and let $\mathrm{Acl}$ be a real connection on $T_k$ with $d\mathrm{Acl}=0$.
Then each of the connections
\[
-A,\qquad\overline{A},\qquad\tilde{g}A,\qquad A-\mathrm{Acl}
\]
belongs to $\tilde{G}_k$.
Consequently $-\overline{A}$, $-\tilde{g}A=\tilde{g}(-A)$ and $-(A-\mathrm{Acl})=(-A)-(-\mathrm{Acl})$ belong to $\tilde{G}_k$ as well.
\end{named}

\begin{proof}
For $-A$, $\overline{A}$ and $\tilde{g}A$ the assertion is contained in the invariance property (G2) of \cite[Section~16.2.7]{QEDI}.
In the notation of the proof of Lemma~1\,(ii) it rests on the following decompositions.
One has $-A=(-A_0)+(-A_1)$, where $-A_0=(-A')+d(-\boldsymbol{\lambda})$ on $\square^{\sim 5}$.
Next, $\overline{A}=A_0+\overline{A_1}$.
Finally, $\tilde{g}A=\tilde{g}A_0+\tilde{g}A_1$, where $|\tilde{g}A_1|_{\infty}=|A_1|_{\infty}$ and
\[
\tilde{g}A_0=\tilde{g}A'+d\left(\boldsymbol{\lambda}\circ\tilde{g}^{-1}\right)
\qquad\text{on }\tilde{g}\left(\square^{\sim 5}\right)=\left(\tilde{g}\square\right)^{\sim 5},
\]
and every $M$-cube is of the form $\tilde{g}\square$.
For $A-\mathrm{Acl}$ the assertion is Lemma~1\,(ii).
The remaining memberships follow by applying the first one to the connections $\overline{A}$, $\tilde{g}A$ and $A-\mathrm{Acl}$.
\end{proof}

\subsubsection*{3.3. Existence of the inverses}

\begin{named}{(Iv)}
Let $B\in\tilde{G}_k$.
\begin{itemize}
\item[(i)] Let $k\ge0$.
For every $y\in Y_{k}$ the operator $F_{y}(B)$ is a bijection of $\mathfrak{S}^{\mathrm{ch}}(\eta_k,\ell)$ onto itself (Proposition~2.8\,(a)).
\item[(ii)] Let $k\ge1$ and let $y$ be real.
The operators
\[
b_{k}+b_{L}\cdot P(B)+i\gamma_5\cdot y
\qquad\text{and}\qquad
D_k(B)+b_{L}\cdot P(B)+i\gamma_5\cdot y
\]
are bijections of $\mathfrak{S}(1,\ell)$ onto itself (Proposition~4\,(a), (b); for $y=0$ also Proposition~4\,(d)).
\end{itemize}
\end{named}

Consequently every operator of Section~I.3.5 (for $k\ge1$) and of Section~I.3.9 (for $k=0$) is defined at each of the connections named in (Dm).

\subsubsection*{3.4. The composed operators of a level $k\ge1$}

Let $k\ge1$.
By Section~I.3.5 the operators of level $k$ are the following words, where $S_k=S_{\infty}$, $S^{0}_{k+1}=S_{0}$, $S_{k,y}=S_{y}$, $B_{k}=B_{k,0}$, $\Gamma_k=\Gamma_{k,0}$, and $-\infty:=\infty$ (Section~G1.9):
\begin{equation}\tag{3.1}
\begin{aligned}
&F_{y}(A)=D_A+\mu+\operatorname{Pen}_{y}(A),\qquad
S_{y}(A)=F_{y}(A)^{-1},\\
&F_{\infty}(A)=D_A+\mu+b_{k}\cdot J\cdot P_k(A)\cdot J^{*},\\
&H_{k}(A)=b_{k}\cdot S_{\infty}(A)\cdot J\cdot Q_k^{T}(-A),\qquad
\mathfrak{H}_{k}(A)=b_{k}\cdot Q_k(A)\cdot J^{*}\cdot S_{\infty}(A),\\
&T_k(A)=b_{k}^{-1}\cdot Q_k(A)\cdot J^{*}\cdot F_{\infty}(A),\\
&M_k(A)=b_{k}-D_k(A)=b_{k}^{2}\cdot Q_k(A)\cdot J^{*}\cdot S_{\infty}(A)\cdot J\cdot Q_k^{T}(-A),\qquad
D_k(A)=b_{k}-M_k(A),\\
&\dot{M}_k(A)=b_{k}^{2}\cdot Q_k(A)\cdot J^{*}\cdot S_{\infty}(A)^{2}\cdot J\cdot Q_k^{T}(-A),\\
&B_{k,y}(A)=\left(b_{k}+b_{L}\cdot P(A)+i\gamma_5\cdot y\right)^{-1},\qquad
\Gamma_{k,y}(A)=\left(D_k(A)+b_{L}\cdot P(A)+i\gamma_5\cdot y\right)^{-1},\\
&\delta Z_k(A)=\det\left(D_k(A)+b_{L}\cdot P(A)\right),\\
&\bar{H}_{k}(A)=\mathsf{C}^{-1}\cdot H_{k}(-A)\cdot\mathsf{C},\qquad
\bar{T}_k(A)=\mathsf{C}^{-1}\cdot T_k(-A)\cdot\mathsf{C}.
\end{aligned}
\end{equation}
The numbers $\mu$, $b_{k}$ and $b_{L}$ are real, and $b_{k}>0$ ((SF); Sections~I.3.4 and I.0.5).
Sources and targets are those of the table in Section~I.3.5.
The operators $D_A$, $\operatorname{Pen}_{y}$, $F_{y}$, $S_{y}$, $J\cdot P_k\cdot J^{*}$ and $J\cdot P_{k+1}\cdot J^{*}$ act on $\mathfrak{S}^{\mathrm{ch}}(\eta_k,\ell)$;
the operators $H_{k}$, $\bar{H}_{k}$ and $J\cdot Q_k^{T}(-\,\cdot\,)$ map $\mathfrak{S}(1,\ell)$ into $\mathfrak{S}^{\mathrm{ch}}(\eta_k,\ell)$;
the operators $T_k$, $\bar{T}_k$, $\mathfrak{H}_{k}$ and $Q_k\cdot J^{*}$ map $\mathfrak{S}^{\mathrm{ch}}(\eta_k,\ell)$ into $\mathfrak{S}(1,\ell)$;
the operators $D_k$, $M_k$, $\dot{M}_k$, $B_{k,y}$, $\Gamma_{k,y}$ and $P$ act on $\mathfrak{S}(1,\ell)$;
finally, $Q$ maps $\mathfrak{S}(1,\ell)$ into $\mathfrak{S}(L,\ell)$, and $Q^{T}(-\,\cdot\,)$ maps $\mathfrak{S}(L,\ell)$ into $\mathfrak{S}(1,\ell)$.

\subsection*{4. Two lemmas on words ($k\ge1$)}

\begin{named}{Lemma 2}[covariant transport]
Let $k\ge1$ and $A,\tilde{A}\in\tilde{G}_k$, and let $y,\tilde{y}\in Y_{k}$ be either both real or both equal to $\infty$.
Let $\mathrm{Vf}$ be a bijection of $\mathfrak{S}^{\mathrm{ch}}(\eta_k,\ell)$ onto itself and $\mathrm{Vu}$ a bijection of $\mathfrak{S}(1,\ell)$ onto itself.
Assume that both maps are additive and commute with the multiplication by real numbers; they may be linear or antilinear.
Suppose that
\begin{itemize}
\item[(Hy1)] $F_{\tilde{y}}(\tilde{A})\cdot\mathrm{Vf}=\mathrm{Vf}\cdot F_{y}(A)$ and $F_{\infty}(\tilde{A})\cdot\mathrm{Vf}=\mathrm{Vf}\cdot F_{\infty}(A)$;
\item[(Hy2)] $Q_k(\tilde{A})\cdot J^{*}\cdot\mathrm{Vf}=\mathrm{Vu}\cdot Q_k(A)\cdot J^{*}$ and $J\cdot Q_k^{T}(-\tilde{A})\cdot\mathrm{Vu}=\mathrm{Vf}\cdot J\cdot Q_k^{T}(-A)$;
\item[(Hy3)] $P(\tilde{A})\cdot\mathrm{Vu}=\mathrm{Vu}\cdot P(A)$;
\item[(Hy4)] if $y$ is real, $\left(i\gamma_5\cdot\tilde{y}\right)\cdot\mathrm{Vu}=\mathrm{Vu}\cdot\left(i\gamma_5\cdot y\right)$.
\end{itemize}
Then the following relations hold.
\begin{itemize}
\item[(Co1)] $S_{\tilde{y}}(\tilde{A})\cdot\mathrm{Vf}=\mathrm{Vf}\cdot S_{y}(A)$ and $S_{\infty}(\tilde{A})\cdot\mathrm{Vf}=\mathrm{Vf}\cdot S_{\infty}(A)$.
\item[(Co2)] $H_{k}(\tilde{A})\cdot\mathrm{Vu}=\mathrm{Vf}\cdot H_{k}(A)$.
\item[(Co3)] $\mathfrak{H}_{k}(\tilde{A})\cdot\mathrm{Vf}=\mathrm{Vu}\cdot\mathfrak{H}_{k}(A)$.
\item[(Co4)] $T_k(\tilde{A})\cdot\mathrm{Vf}=\mathrm{Vu}\cdot T_k(A)$.
\item[(Co5)] $O(\tilde{A})\cdot\mathrm{Vu}=\mathrm{Vu}\cdot O(A)$ for $O=M_k$, $D_k$, $\dot{M}_k$.
\item[(Co6)] If $y$ is real, $B_{k,\tilde{y}}(\tilde{A})\cdot\mathrm{Vu}=\mathrm{Vu}\cdot B_{k,y}(A)$.
\item[(Co7)] If $y$ is real,
\[
\left(D_k(\tilde{A})+b_{L}\cdot P(\tilde{A})+i\gamma_5\cdot\tilde{y}\right)\cdot\mathrm{Vu}=\mathrm{Vu}\cdot\left(D_k(A)+b_{L}\cdot P(A)+i\gamma_5\cdot y\right)
\]
and $\Gamma_{k,\tilde{y}}(\tilde{A})\cdot\mathrm{Vu}=\mathrm{Vu}\cdot\Gamma_{k,y}(A)$.
\end{itemize}
\end{named}

\begin{proof}
All the inverses that occur exist by (Iv).
By hypothesis every real number commutes with $\mathrm{Vf}$ and with $\mathrm{Vu}$; in particular so do $b_{k}$, $b_{k}^{-1}$, $b_{k}^{2}$ and $b_{L}$.
The operators in question are the words (3.1).

(Co1) follows from (Hy1) and (Ru5).

(Co2).
By (Hy2) and (Co1),
\begin{align*}
H_{k}(\tilde{A})\cdot\mathrm{Vu}
&=b_{k}\cdot S_{\infty}(\tilde{A})\cdot\left[J\cdot Q_k^{T}(-\tilde{A})\cdot\mathrm{Vu}\right]
=b_{k}\cdot\left[S_{\infty}(\tilde{A})\cdot\mathrm{Vf}\right]\cdot J\cdot Q_k^{T}(-A)\\
&=b_{k}\cdot\mathrm{Vf}\cdot S_{\infty}(A)\cdot J\cdot Q_k^{T}(-A)
=\mathrm{Vf}\cdot H_{k}(A).
\end{align*}

(Co3).
By (Co1) and (Hy2),
\begin{align*}
\mathfrak{H}_{k}(\tilde{A})\cdot\mathrm{Vf}
&=b_{k}\cdot Q_k(\tilde{A})\cdot J^{*}\cdot\left[S_{\infty}(\tilde{A})\cdot\mathrm{Vf}\right]
=b_{k}\cdot\left[Q_k(\tilde{A})\cdot J^{*}\cdot\mathrm{Vf}\right]\cdot S_{\infty}(A)\\
&=b_{k}\cdot\mathrm{Vu}\cdot Q_k(A)\cdot J^{*}\cdot S_{\infty}(A)
=\mathrm{Vu}\cdot\mathfrak{H}_{k}(A).
\end{align*}

(Co4).
By (Hy1) and (Hy2),
\begin{align*}
T_k(\tilde{A})\cdot\mathrm{Vf}
&=b_{k}^{-1}\cdot Q_k(\tilde{A})\cdot J^{*}\cdot\left[F_{\infty}(\tilde{A})\cdot\mathrm{Vf}\right]
=b_{k}^{-1}\cdot\left[Q_k(\tilde{A})\cdot J^{*}\cdot\mathrm{Vf}\right]\cdot F_{\infty}(A)\\
&=b_{k}^{-1}\cdot\mathrm{Vu}\cdot Q_k(A)\cdot J^{*}\cdot F_{\infty}(A)
=\mathrm{Vu}\cdot T_k(A).
\end{align*}

(Co5).
Let $j\in\{1,2\}$.
Using (Hy2), then (Co1) $j$ times, then (Hy2) once more, one obtains
\begin{align*}
b_{k}^{2}\cdot Q_k(\tilde{A})\cdot J^{*}\cdot S_{\infty}(\tilde{A})^{j}\cdot\left[J\cdot Q_k^{T}(-\tilde{A})\cdot\mathrm{Vu}\right]
&=b_{k}^{2}\cdot Q_k(\tilde{A})\cdot J^{*}\cdot\left[S_{\infty}(\tilde{A})^{j}\cdot\mathrm{Vf}\right]\cdot J\cdot Q_k^{T}(-A)\\
&=b_{k}^{2}\cdot\left[Q_k(\tilde{A})\cdot J^{*}\cdot\mathrm{Vf}\right]\cdot S_{\infty}(A)^{j}\cdot J\cdot Q_k^{T}(-A)\\
&=\mathrm{Vu}\cdot b_{k}^{2}\cdot Q_k(A)\cdot J^{*}\cdot S_{\infty}(A)^{j}\cdot J\cdot Q_k^{T}(-A).
\end{align*}
For $j=1$ this is the relation for $M_k$, for $j=2$ the relation for $\dot{M}_k$.
It follows that
\[
D_k(\tilde{A})\cdot\mathrm{Vu}=b_{k}\cdot\mathrm{Vu}-M_k(\tilde{A})\cdot\mathrm{Vu}=\mathrm{Vu}\cdot\left(b_{k}-M_k(A)\right)=\mathrm{Vu}\cdot D_k(A).
\]

(Co6).
By (Hy3), (Hy4) and the additivity of $\mathrm{Vu}$,
\[
\left(b_{k}+b_{L}\cdot P(\tilde{A})+i\gamma_5\cdot\tilde{y}\right)\cdot\mathrm{Vu}=\mathrm{Vu}\cdot\left(b_{k}+b_{L}\cdot P(A)+i\gamma_5\cdot y\right),
\]
and (Ru5) gives (Co6).

(Co7).
The first relation is the sum of (Co5) for $D_k$, of $b_{L}$ times (Hy3), and of (Hy4).
The second relation follows from the first by (Ru5).
\end{proof}

\begin{named}{Lemma 3}[contravariant transport]
Let $k\ge1$, and let ${}^{\mathrm{dg}}$ be one of the two operations ${}^{T}$, ${}^{*}$ of Section~2.1.
Let $\mathrm{Nx}$ be an invertible constant $4\times4$ matrix, let $A,\tilde{A}\in\tilde{G}_k$, and let $y\in Y_{k}$.
Suppose that
\begin{itemize}
\item[(Hz1)] $F_{y}(A)^{\mathrm{dg}}=\mathrm{Nx}^{-1}\cdot F_{-y}(\tilde{A})\cdot\mathrm{Nx}$ and $F_{\infty}(A)^{\mathrm{dg}}=\mathrm{Nx}^{-1}\cdot F_{\infty}(\tilde{A})\cdot\mathrm{Nx}$;
\item[(Hz2)] $\left(Q_k(A)\cdot J^{*}\right)^{\mathrm{dg}}=J\cdot Q_k^{T}(-\tilde{A})$ and $\left(J\cdot Q_k^{T}(-A)\right)^{\mathrm{dg}}=Q_k(\tilde{A})\cdot J^{*}$;
\item[(Hz3)] $P(A)^{\mathrm{dg}}=P(\tilde{A})$;
\item[(Hz4)] if $y$ is real, $\left(i\gamma_5\cdot y\right)^{\mathrm{dg}}=\mathrm{Nx}^{-1}\cdot\left(-i\gamma_5\cdot y\right)\cdot\mathrm{Nx}$.
\end{itemize}
Then the following relations hold.
\begin{itemize}
\item[(Cz1)] $S_{y}(A)^{\mathrm{dg}}=\mathrm{Nx}^{-1}\cdot S_{-y}(\tilde{A})\cdot\mathrm{Nx}$ and $S_{\infty}(A)^{\mathrm{dg}}=\mathrm{Nx}^{-1}\cdot S_{\infty}(\tilde{A})\cdot\mathrm{Nx}$.
\item[(Cz2)] $H_{k}(A)^{\mathrm{dg}}=\mathrm{Nx}^{-1}\cdot\mathfrak{H}_{k}(\tilde{A})\cdot\mathrm{Nx}$.
\item[(Cz3)] $\mathfrak{H}_{k}(A)^{\mathrm{dg}}=\mathrm{Nx}^{-1}\cdot H_{k}(\tilde{A})\cdot\mathrm{Nx}$.
\item[(Cz4)] $O(A)^{\mathrm{dg}}=\mathrm{Nx}^{-1}\cdot O(\tilde{A})\cdot\mathrm{Nx}$ for $O=M_k$, $D_k$, $\dot{M}_k$.
\item[(Cz5)] If $y$ is real, $B_{k,y}(A)^{\mathrm{dg}}=\mathrm{Nx}^{-1}\cdot B_{k,-y}(\tilde{A})\cdot\mathrm{Nx}$.
\item[(Cz6)] If $y$ is real,
\[
\left(D_k(A)+b_{L}\cdot P(A)+i\gamma_5\cdot y\right)^{\mathrm{dg}}=\mathrm{Nx}^{-1}\cdot\left(D_k(\tilde{A})+b_{L}\cdot P(\tilde{A})-i\gamma_5\cdot y\right)\cdot\mathrm{Nx}
\]
and $\Gamma_{k,y}(A)^{\mathrm{dg}}=\mathrm{Nx}^{-1}\cdot\Gamma_{k,-y}(\tilde{A})\cdot\mathrm{Nx}$.
\end{itemize}
\end{named}

\begin{proof}
All the inverses that occur exist by (Iv).
By (Ru1) the operation ${}^{\mathrm{dg}}$ is additive, reverses the order of products and commutes with taking inverses, and $\left(\mathrm{cf}\cdot O\right)^{\mathrm{dg}}=\mathrm{cf}\cdot O^{\mathrm{dg}}$ for every \emph{real} number $\mathrm{cf}$; this is applied to $b_{k}$, $b_{k}^{2}$ and $b_{L}$.
By (Sca) the operators $Q_k(\tilde{A})\cdot J^{*}$, $J\cdot Q_k^{T}(-\tilde{A})$ and $P(\tilde{A})$ commute with $\mathrm{Nx}$ and with $\mathrm{Nx}^{-1}$.

(Cz1).
By (Ru1) and the first relation of (Hz1),
\[
S_{y}(A)^{\mathrm{dg}}=\left(F_{y}(A)^{\mathrm{dg}}\right)^{-1}=\left(\mathrm{Nx}^{-1}\cdot F_{-y}(\tilde{A})\cdot\mathrm{Nx}\right)^{-1}=\mathrm{Nx}^{-1}\cdot S_{-y}(\tilde{A})\cdot\mathrm{Nx};
\]
in the same way the second relation of (Hz1) gives
\[
S_{\infty}(A)^{\mathrm{dg}}=\left(\mathrm{Nx}^{-1}\cdot F_{\infty}(\tilde{A})\cdot\mathrm{Nx}\right)^{-1}=\mathrm{Nx}^{-1}\cdot S_{\infty}(\tilde{A})\cdot\mathrm{Nx}.
\]

(Cz2).
By (Hz2), (Cz1) and (Sca),
\begin{align*}
H_{k}(A)^{\mathrm{dg}}
&=b_{k}\cdot\left(J\cdot Q_k^{T}(-A)\right)^{\mathrm{dg}}\cdot S_{\infty}(A)^{\mathrm{dg}}
=b_{k}\cdot Q_k(\tilde{A})\cdot J^{*}\cdot\mathrm{Nx}^{-1}\cdot S_{\infty}(\tilde{A})\cdot\mathrm{Nx}\\
&=\mathrm{Nx}^{-1}\cdot b_{k}\cdot Q_k(\tilde{A})\cdot J^{*}\cdot S_{\infty}(\tilde{A})\cdot\mathrm{Nx}
=\mathrm{Nx}^{-1}\cdot\mathfrak{H}_{k}(\tilde{A})\cdot\mathrm{Nx}.
\end{align*}

(Cz3).
By (Cz1), (Hz2) and (Sca),
\begin{align*}
\mathfrak{H}_{k}(A)^{\mathrm{dg}}
&=b_{k}\cdot S_{\infty}(A)^{\mathrm{dg}}\cdot\left(Q_k(A)\cdot J^{*}\right)^{\mathrm{dg}}
=b_{k}\cdot\mathrm{Nx}^{-1}\cdot S_{\infty}(\tilde{A})\cdot\mathrm{Nx}\cdot J\cdot Q_k^{T}(-\tilde{A})\\
&=\mathrm{Nx}^{-1}\cdot b_{k}\cdot S_{\infty}(\tilde{A})\cdot J\cdot Q_k^{T}(-\tilde{A})\cdot\mathrm{Nx}
=\mathrm{Nx}^{-1}\cdot H_{k}(\tilde{A})\cdot\mathrm{Nx}.
\end{align*}

(Cz4).
Let $j\in\{1,2\}$.
By (Ru1), (Hz2), (Cz1) and (Sca),
\begin{align*}
\left(b_{k}^{2}\cdot Q_k(A)\cdot J^{*}\cdot S_{\infty}(A)^{j}\cdot J\cdot Q_k^{T}(-A)\right)^{\mathrm{dg}}
&=b_{k}^{2}\cdot\left(J\cdot Q_k^{T}(-A)\right)^{\mathrm{dg}}\cdot\left(S_{\infty}(A)^{\mathrm{dg}}\right)^{j}\cdot\left(Q_k(A)\cdot J^{*}\right)^{\mathrm{dg}}\\
&=b_{k}^{2}\cdot Q_k(\tilde{A})\cdot J^{*}\cdot\left(\mathrm{Nx}^{-1}\cdot S_{\infty}(\tilde{A})\cdot\mathrm{Nx}\right)^{j}\cdot J\cdot Q_k^{T}(-\tilde{A})\\
&=b_{k}^{2}\cdot Q_k(\tilde{A})\cdot J^{*}\cdot\mathrm{Nx}^{-1}\cdot S_{\infty}(\tilde{A})^{j}\cdot\mathrm{Nx}\cdot J\cdot Q_k^{T}(-\tilde{A})\\
&=\mathrm{Nx}^{-1}\cdot b_{k}^{2}\cdot Q_k(\tilde{A})\cdot J^{*}\cdot S_{\infty}(\tilde{A})^{j}\cdot J\cdot Q_k^{T}(-\tilde{A})\cdot\mathrm{Nx}.
\end{align*}
The case $j=1$ is the relation for $M_k$, the case $j=2$ the relation for $\dot{M}_k$.
Moreover
\[
D_k(A)^{\mathrm{dg}}=b_{k}-M_k(A)^{\mathrm{dg}}=\mathrm{Nx}^{-1}\cdot\left(b_{k}-M_k(\tilde{A})\right)\cdot\mathrm{Nx}.
\]

(Cz5).
By (Hz3), (Hz4) and (Sca),
\begin{align*}
\left(b_{k}+b_{L}\cdot P(A)+i\gamma_5\cdot y\right)^{\mathrm{dg}}
&=b_{k}+b_{L}\cdot P(\tilde{A})+\mathrm{Nx}^{-1}\cdot\left(-i\gamma_5\cdot y\right)\cdot\mathrm{Nx}\\
&=\mathrm{Nx}^{-1}\cdot\left(b_{k}+b_{L}\cdot P(\tilde{A})-i\gamma_5\cdot y\right)\cdot\mathrm{Nx};
\end{align*}
taking the inverse of both sides gives (Cz5).

(Cz6).
The first relation is the sum of (Cz4) for $D_k$, of $b_{L}$ times (Hz3) (together with (Sca)), and of (Hz4).
Taking the inverse of both sides gives the second relation.
\end{proof}

\subsection*{5. Proof of (a)--(f) ($k\ge1$)}

Proposition~2.2 is applied on the block torus $T_k=(0,\ell,-1)$ of Section~G1.2, with the numbers $e_k$ and $\mu$ of the data and with complex connections on $T_k$.
In Proposition~2.2\,(a), (b) the pair $(x,x')$ is an arbitrary ordered pair of distinct points of $\mathbb{L}_k\times O_{\mathrm{cub}}$; the pairs of $\mathrm{Pzt}_k$ are pairs of this kind (the data of Proposition~2.2; Section~G1.2).

\subsubsection*{5.1. Proof of (a)}

\begin{proof}[Proof of Proposition~10\,(a)]
\emph{The domain.}
By Lemma~1\,(ii), $A-\mathrm{Acl}$ lies in $\tilde{G}_k$.

\emph{The elementary relations.}
The function $\mathrm{uf}$ is a function on $T_k$, that is, an $\ell\mathbb{Z}^3$-periodic function on $\mathbb{L}_k$.
Hence $(\mathrm{Acl},\mathrm{uf})$ are data of Proposition~2.2\,(a) on $T_k$, and $\mathrm{uf}^{0}$, $\mathrm{uf}^{1}$ are the restrictions of $\mathrm{uf}$ to $\mathbb{Z}^3$ and to $L\mathbb{Z}^3$.
For every $y\in Y_{k}$, Proposition~2.2\,(a) yields
\begin{equation}\tag{5.1}
\begin{aligned}
&D_{A-\mathrm{Acl}}=\mathrm{uf}\cdot D_A\cdot\mathrm{uf}^{-1},\qquad
Q_k(A-\mathrm{Acl})\cdot J^{*}=\mathrm{uf}^{0}\cdot Q_k(A)\cdot J^{*}\cdot\mathrm{uf}^{-1},\\
&J\cdot Q_k^{T}\left(-(A-\mathrm{Acl})\right)=\mathrm{uf}\cdot J\cdot Q_k^{T}(-A)\cdot\left(\mathrm{uf}^{0}\right)^{-1},\qquad
Q(A-\mathrm{Acl})=\mathrm{uf}^{1}\cdot Q(A)\cdot\left(\mathrm{uf}^{0}\right)^{-1},\\
&Q^{T}\left(-(A-\mathrm{Acl})\right)=\mathrm{uf}^{0}\cdot Q^{T}(-A)\cdot\left(\mathrm{uf}^{1}\right)^{-1},\\
&O(A-\mathrm{Acl})=\mathrm{uf}\cdot O(A)\cdot\mathrm{uf}^{-1}\quad\text{for }O=J\cdot P_k\cdot J^{*},\ J\cdot P_{k+1}\cdot J^{*},\ \operatorname{Pen}_{y},\ F_{y},\\
&P(A-\mathrm{Acl})=\mathrm{uf}^{0}\cdot P(A)\cdot\left(\mathrm{uf}^{0}\right)^{-1},\qquad
\left(\delta_{\alpha,A-\mathrm{Acl}}(\mathrm{uf}\cdot f)\right)(x,x')=\mathrm{uf}(c)\cdot\left(\delta_{\alpha,A}f\right)(x,x').
\end{aligned}
\end{equation}

\emph{The composed operators.}
Lemma~2 is applied with $\tilde{A}:=A-\mathrm{Acl}$, $\tilde{y}:=y$, $\mathrm{Vf}:=\mathrm{uf}$ and $\mathrm{Vu}:=\mathrm{uf}^{0}$; since $|\mathrm{uf}|=1$, these are $\mathbb{C}$-linear bijections by (Ru3).
The relation of (5.1) for $F_{y}$, taken at $y$ and at $\infty$ and multiplied from the right by $\mathrm{uf}$, is (Hy1).
The second and the third relation of (5.1), multiplied from the right by $\mathrm{uf}$ and by $\mathrm{uf}^{0}$ respectively, constitute (Hy2).
The relation for $P$, multiplied from the right by $\mathrm{uf}^{0}$, is (Hy3).
(Hy4) holds with $\tilde{y}=y$, because the function $\mathrm{uf}^{0}$ commutes with the constant matrix $\gamma_5$ (Ru3).
Multiplying (Co1)--(Co7) from the right by $\mathrm{uf}^{-1}$, respectively by $\left(\mathrm{uf}^{0}\right)^{-1}$, one obtains
\begin{align*}
&S_{y}(A-\mathrm{Acl})=\mathrm{uf}\cdot S_{y}(A)\cdot\mathrm{uf}^{-1}\ \text{(Co1)},\qquad
H_{k}(A-\mathrm{Acl})=\mathrm{uf}\cdot H_{k}(A)\cdot\left(\mathrm{uf}^{0}\right)^{-1}\ \text{(Co2)},\\
&\mathfrak{H}_{k}(A-\mathrm{Acl})=\mathrm{uf}^{0}\cdot\mathfrak{H}_{k}(A)\cdot\mathrm{uf}^{-1}\ \text{(Co3)},\qquad
T_k(A-\mathrm{Acl})=\mathrm{uf}^{0}\cdot T_k(A)\cdot\mathrm{uf}^{-1}\ \text{(Co4)},\\
&O(A-\mathrm{Acl})=\mathrm{uf}^{0}\cdot O(A)\cdot\left(\mathrm{uf}^{0}\right)^{-1}\\
&\qquad\text{for }O=D_k,\ M_k,\ \dot{M}_k\ \text{(Co5)},\ B_{k,y}\ \text{(Co6)},\ \Gamma_{k,y}\ \text{(Co7)}\quad(y\text{ real}).
\end{align*}

\emph{The determinant.}
The first relation of (Co7) at $y=0$ reads
\[
D_k(A-\mathrm{Acl})+b_{L}\cdot P(A-\mathrm{Acl})=\mathrm{uf}^{0}\cdot\left(D_k(A)+b_{L}\cdot P(A)\right)\cdot\left(\mathrm{uf}^{0}\right)^{-1},
\]
and (Ru4)\,(i) gives $\delta Z_k(A-\mathrm{Acl})=\delta Z_k(A)$.

\emph{The barred operators.}
The triple $\left(-A,-\mathrm{Acl},\mathrm{uf}^{-1}\right)$ satisfies the hypotheses of (a).
Indeed, $-A$ lies in $\tilde{G}_k$ by (Dm); $-\mathrm{Acl}$ is real and $d(-\mathrm{Acl})=0$; $\left|\mathrm{uf}^{-1}\right|=1$, and taking inverses in the relation for $\mathrm{uf}$ gives
\[
\mathrm{uf}^{-1}(c')=\mathrm{uf}^{-1}(c)\cdot\exp\left(i e_k\cdot\eta_k\cdot(-\mathrm{Acl})(c,c')\right)
\]
on every bond; finally, the restriction of $\mathrm{uf}^{-1}$ to $T^0_k$ is $\left(\mathrm{uf}^{0}\right)^{-1}$.
For this triple the relations just proved read
\[
H_{k}\left((-A)-(-\mathrm{Acl})\right)=\mathrm{uf}^{-1}\cdot H_{k}(-A)\cdot\mathrm{uf}^{0},\qquad
T_k\left((-A)-(-\mathrm{Acl})\right)=\left(\mathrm{uf}^{0}\right)^{-1}\cdot T_k(-A)\cdot\mathrm{uf}.
\]
Since $-(A-\mathrm{Acl})=(-A)-(-\mathrm{Acl})$, and since $\mathsf{C}$ commutes with $\mathrm{uf}$, $\mathrm{uf}^{0}$ and their inverses (Ru3), it follows that
\begin{align*}
\bar{H}_{k}(A-\mathrm{Acl})
&=\mathsf{C}^{-1}\cdot H_{k}\left(-(A-\mathrm{Acl})\right)\cdot\mathsf{C}
=\mathsf{C}^{-1}\cdot\mathrm{uf}^{-1}\cdot H_{k}(-A)\cdot\mathrm{uf}^{0}\cdot\mathsf{C}\\
&=\mathrm{uf}^{-1}\cdot\left(\mathsf{C}^{-1}\cdot H_{k}(-A)\cdot\mathsf{C}\right)\cdot\mathrm{uf}^{0}
=\mathrm{uf}^{-1}\cdot\bar{H}_{k}(A)\cdot\mathrm{uf}^{0},\\
\bar{T}_k(A-\mathrm{Acl})
&=\mathsf{C}^{-1}\cdot T_k\left(-(A-\mathrm{Acl})\right)\cdot\mathsf{C}
=\mathsf{C}^{-1}\cdot\left(\mathrm{uf}^{0}\right)^{-1}\cdot T_k(-A)\cdot\mathrm{uf}\cdot\mathsf{C}\\
&=\left(\mathrm{uf}^{0}\right)^{-1}\cdot\left(\mathsf{C}^{-1}\cdot T_k(-A)\cdot\mathsf{C}\right)\cdot\mathrm{uf}
=\left(\mathrm{uf}^{0}\right)^{-1}\cdot\bar{T}_k(A)\cdot\mathrm{uf}.
\end{align*}

\emph{The examples.}
The verification that follows is valid for every $k\ge0$.

(1) Let $\boldsymbol{\lambda}$ be a real function of $c$ on $T_k$, and put $\mathrm{Acl}:=d\boldsymbol{\lambda}$ and $\mathrm{uf}:=\exp\left(i e_k\cdot\boldsymbol{\lambda}\right)$.
Then $\mathrm{Acl}$ is real.
On the oriented bond $(c,c')$ one has $(d\boldsymbol{\lambda})(c,c')=\left(\boldsymbol{\lambda}(c')-\boldsymbol{\lambda}(c)\right)/\eta_k$ (Section~I.1.1), so that
\[
\mathrm{uf}(c')=\mathrm{uf}(c)\cdot\exp\left(i e_k\cdot\left(\boldsymbol{\lambda}(c')-\boldsymbol{\lambda}(c)\right)\right)=\mathrm{uf}(c)\cdot\exp\left(i e_k\cdot\eta_k\cdot\mathrm{Acl}(c,c')\right).
\]
Further, $(dd\boldsymbol{\lambda})_{\mu\nu}=\partial_\mu\partial_\nu\boldsymbol{\lambda}-\partial_\nu\partial_\mu\boldsymbol{\lambda}=0$, $|\mathrm{uf}|=1$, and $\mathrm{uf}$ is periodic together with $\boldsymbol{\lambda}$.

(2) Let $\mathrm{Acl}_\mu$, $\mu=1,2,3$, be real numbers with $e_k\cdot\ell\cdot\mathrm{Acl}_\mu\in2\pi\mathbb{Z}$, and put
\[
\mathrm{Acl}(c,c+\eta_k e_{\mu}):=\mathrm{Acl}_\mu,\qquad
\mathrm{Acl}(c+\eta_k e_{\mu},c):=-\mathrm{Acl}_\mu,\qquad
\mathrm{uf}(c):=\exp\left(i e_k\cdot\mathrm{Acl}\cdot c\right),
\]
where $\mathrm{Acl}\cdot c:=\sum_\mu\mathrm{Acl}_\mu c_\mu$.
Then $(d\mathrm{Acl})_{\mu\nu}=\partial_\mu\mathrm{Acl}_\nu-\partial_\nu\mathrm{Acl}_\mu=0$, and
\[
\mathrm{uf}(c\pm\eta_k e_{\mu})=\mathrm{uf}(c)\cdot\exp\left(\pm i e_k\cdot\eta_k\cdot\mathrm{Acl}_\mu\right)=\mathrm{uf}(c)\cdot\exp\left(i e_k\cdot\eta_k\cdot\mathrm{Acl}(c,c\pm\eta_k e_{\mu})\right).
\]
Moreover $\mathrm{uf}(c+\ell e_{\mu})=\mathrm{uf}(c)\cdot\exp\left(i e_k\cdot\ell\cdot\mathrm{Acl}_\mu\right)=\mathrm{uf}(c)$, so that $\mathrm{uf}$ is a function on $T_k$; and $|\mathrm{uf}|=1$.
For the data $(k,\ell,\mu,e_k)=(k_{*},\ell_*,\mu_*,e_{*})$ the relation for $D_k$ in this example, namely
\[
D_k(A-\mathrm{Acl})=\mathrm{uf}^{0}\cdot D_k(A)\cdot\left(\mathrm{uf}^{0}\right)^{-1}\qquad\text{with }\mathrm{uf}^{0}(z)=\exp\left(i e_k\cdot\mathrm{Acl}\cdot z\right),
\]
is the form (G5.T-10) of Annex~A, Section~G5.7.
\end{proof}

\subsubsection*{5.2. Proof of (b)}

\begin{proof}[Proof of Proposition~10\,(b)]
Let $\tilde{g}=(S,a)\in\mathbb{G}_k$ and $r=r(S)$.
By (Dm), $\tilde{g}A$ lies in $\tilde{G}_k$.
For $y\in Y_{k}$ let $y':=y$ if $S\in\mathrm{Sp}_3$ or if $y$ is neutral, that is, $y\in\{0,\infty\}$, and let $y':=-y$ otherwise.
For neutral $y$ one has $-y=y$, since $-0=0$ and $-\infty:=\infty$; hence in every case $y'=y$ for $S\in\mathrm{Sp}_3$ and $y'=-y$ for $S\notin\mathrm{Sp}_3$.

\emph{The elementary relations.}
Proposition~2.2\,(b), applied to the morphism $c\mapsto rc+a$ of $T_k$ onto itself and to $\mathrm{mS}=\tilde{g}$ (Section~2.3), gives
\begin{equation}\tag{5.2}
\begin{aligned}
&D_{\tilde{g}A}\cdot\tilde{g}=\tilde{g}\cdot D_A,\qquad
Q_k(\tilde{g}A)\cdot J^{*}\cdot\tilde{g}=\tilde{g}\cdot Q_k(A)\cdot J^{*},\qquad
J\cdot Q_k^{T}(-\tilde{g}A)\cdot\tilde{g}=\tilde{g}\cdot J\cdot Q_k^{T}(-A),\\
&Q(\tilde{g}A)\cdot\tilde{g}=\tilde{g}\cdot Q(A),\qquad
Q^{T}(-\tilde{g}A)\cdot\tilde{g}=\tilde{g}\cdot Q^{T}(-A),\\
&O(\tilde{g}A)\cdot\tilde{g}=\tilde{g}\cdot O(A)\quad\text{for }O=J\cdot P_k\cdot J^{*},\ P,\ J\cdot P_{k+1}\cdot J^{*};\qquad
F_{y'}(\tilde{g}A)\cdot\tilde{g}=\tilde{g}\cdot F_{y}(A);\\
&\left(\delta_{\alpha,\tilde{g}A}\,\tilde{g}f\right)(\tilde{g}x,\tilde{g}x')=S\cdot\left(\delta_{\alpha,A}f\right)(x,x').
\end{aligned}
\end{equation}
These are the relations of (b) for $O=D$, $Q_k\cdot J^{*}$, $J\cdot Q_k^{T}(-\,\cdot\,)$, $Q$, $Q^{T}(-\,\cdot\,)$, $J\cdot P_k\cdot J^{*}$, $P$, $J\cdot P_{k+1}\cdot J^{*}$, the relation for $O_y=F_y$, and the last relation of (b); in the notation of Section~I.4.3 for functions of pairs the latter reads $\delta_{\alpha,\tilde{g}A}\cdot\tilde{g}=\tilde{g}\cdot\delta_{\alpha,A}$.

\emph{The operator $\operatorname{Pen}_{y}$.}
Since $\operatorname{Pen}_{y}(A)=F_{y}(A)-D_A-\mu$,
\[
\operatorname{Pen}_{y'}(\tilde{g}A)\cdot\tilde{g}=F_{y'}(\tilde{g}A)\cdot\tilde{g}-D_{\tilde{g}A}\cdot\tilde{g}-\mu\cdot\tilde{g}=\tilde{g}\cdot\left(F_{y}(A)-D_A-\mu\right)=\tilde{g}\cdot\operatorname{Pen}_{y}(A).
\]

\emph{The composed operators.}
Lemma~2 is applied with $\tilde{A}:=\tilde{g}A$, $\tilde{y}:=y'$, $\mathrm{Vf}:=\tilde{g}$ on $\mathfrak{S}^{\mathrm{ch}}(\eta_k,\ell)$ and $\mathrm{Vu}:=\tilde{g}$ on $\mathfrak{S}(1,\ell)$; these are $\mathbb{C}$-linear bijections (Section~2.3).
(Hy1) is the relation of (5.2) for $F_{y}$, at $y$ and at $\infty$ (note that $\infty'=\infty$).
(Hy2) and (Hy3) are contained in (5.2).
For (Hy4), let $y$ be real and $v\in\mathfrak{S}(1,\ell)$; by (Sn2),
\[
\left(\tilde{g}\left(i\gamma_5\cdot y\cdot v\right)\right)(z)=S\cdot i\gamma_5\cdot y\cdot v\left(\tilde{g}^{-1}z\right)=iy\cdot\left(S\cdot\gamma_5\cdot S^{-1}\right)\cdot S\cdot v\left(\tilde{g}^{-1}z\right)=\left(i\gamma_5\cdot y'\cdot\tilde{g}v\right)(z).
\]
The conclusions are
\begin{align*}
&S_{y'}(\tilde{g}A)\cdot\tilde{g}=\tilde{g}\cdot S_{y}(A)\ \text{(Co1)};\quad\text{in particular, at }y=\infty\text{ and at }y=0,\\
&\qquad S_k(\tilde{g}A)\cdot\tilde{g}=\tilde{g}\cdot S_k(A),\qquad S^{0}_{k+1}(\tilde{g}A)\cdot\tilde{g}=\tilde{g}\cdot S^{0}_{k+1}(A);\\
&H_{k}(\tilde{g}A)\cdot\tilde{g}=\tilde{g}\cdot H_{k}(A)\ \text{(Co2)},\qquad
\mathfrak{H}_{k}(\tilde{g}A)\cdot\tilde{g}=\tilde{g}\cdot\mathfrak{H}_{k}(A)\ \text{(Co3)},\\
&T_k(\tilde{g}A)\cdot\tilde{g}=\tilde{g}\cdot T_k(A)\ \text{(Co4)},\\
&O(\tilde{g}A)\cdot\tilde{g}=\tilde{g}\cdot O(A)\quad\text{for }O=D_k,\ M_k,\ \dot{M}_k\ \text{(Co5)};\\
&B_{k,y'}(\tilde{g}A)\cdot\tilde{g}=\tilde{g}\cdot B_{k,y}(A)\ \text{(Co6)},\qquad
\Gamma_{k,y'}(\tilde{g}A)\cdot\tilde{g}=\tilde{g}\cdot\Gamma_{k,y}(A)\ \text{(Co7)}\qquad(y\text{ real});\\
&\qquad\text{at }y=0:\quad B_{k}(\tilde{g}A)\cdot\tilde{g}=\tilde{g}\cdot B_{k}(A),\qquad \Gamma_k(\tilde{g}A)\cdot\tilde{g}=\tilde{g}\cdot\Gamma_k(A).
\end{align*}

\emph{The determinant.}
By the first relation of (Co7) at $y=0$ (note that $0'=0$ for every $S\in\mathrm{Sp}$),
\[
D_k(\tilde{g}A)+b_{L}\cdot P(\tilde{g}A)=\tilde{g}\cdot\left(D_k(A)+b_{L}\cdot P(A)\right)\cdot\tilde{g}^{-1}\qquad\text{on }\mathfrak{S}(1,\ell),
\]
and (Ru4)\,(i) gives $\delta Z_k(\tilde{g}A)=\delta Z_k(A)$.

\emph{The barred operators.}
By (Dm), $-A$ lies in $\tilde{G}_k$, and $-\tilde{g}A=\tilde{g}(-A)$ (Section~2.3).
From (Sn1) and from the relations for $H_{k}$ and $T_k$ just proved, taken at the connection $-A$, one obtains
\begin{align*}
\bar{H}_{k}(\tilde{g}A)\cdot\mathrm{glb}
&=\mathsf{C}^{-1}\cdot H_{k}\left(\tilde{g}(-A)\right)\cdot\mathsf{C}\cdot\mathrm{sgn}(S)\cdot\mathsf{C}^{-1}\cdot\tilde{g}\cdot\mathsf{C}
=\mathrm{sgn}(S)\cdot\mathsf{C}^{-1}\cdot H_{k}\left(\tilde{g}(-A)\right)\cdot\tilde{g}\cdot\mathsf{C}\\
&=\mathrm{sgn}(S)\cdot\mathsf{C}^{-1}\cdot\tilde{g}\cdot H_{k}(-A)\cdot\mathsf{C}\\
&=\left(\mathrm{sgn}(S)\cdot\mathsf{C}^{-1}\cdot\tilde{g}\cdot\mathsf{C}\right)\cdot\left(\mathsf{C}^{-1}\cdot H_{k}(-A)\cdot\mathsf{C}\right)
=\mathrm{glb}\cdot\bar{H}_{k}(A),\\
\bar{T}_k(\tilde{g}A)\cdot\mathrm{glb}
&=\mathsf{C}^{-1}\cdot T_k\left(\tilde{g}(-A)\right)\cdot\mathsf{C}\cdot\mathrm{sgn}(S)\cdot\mathsf{C}^{-1}\cdot\tilde{g}\cdot\mathsf{C}
=\mathrm{sgn}(S)\cdot\mathsf{C}^{-1}\cdot T_k\left(\tilde{g}(-A)\right)\cdot\tilde{g}\cdot\mathsf{C}\\
&=\mathrm{sgn}(S)\cdot\mathsf{C}^{-1}\cdot\tilde{g}\cdot T_k(-A)\cdot\mathsf{C}\\
&=\left(\mathrm{sgn}(S)\cdot\mathsf{C}^{-1}\cdot\tilde{g}\cdot\mathsf{C}\right)\cdot\left(\mathsf{C}^{-1}\cdot T_k(-A)\cdot\mathsf{C}\right)
=\mathrm{glb}\cdot\bar{T}_k(A).
\end{align*}
In the first of these two relations the map $\mathrm{glb}$ on the left-hand side acts on $\mathfrak{S}(1,\ell)$ and the one on the right-hand side on $\mathfrak{S}^{\mathrm{ch}}(\eta_k,\ell)$; in the second relation it is the other way round.
The sign $\mathrm{sgn}(S)$ is the same on both spaces.
\end{proof}

\subsubsection*{5.3. Proof of (c)}

\begin{proof}[Proof of Proposition~10\,(c)]
By (Dm), $-A$ lies in $\tilde{G}_k$.
Proposition~2.2\,(c) is valid for every complex connection, and it gives
\begin{equation}\tag{5.3}
\begin{aligned}
&D_A^{T}=\mathsf{C}^{-1}\cdot D_{-A}\cdot\mathsf{C},\qquad
\left(Q_k(A)\cdot J^{*}\right)^{T}=J\cdot Q_k^{T}(A),\qquad
Q(A)^{T}=Q^{T}(A),\\
&\left(J\cdot P_k(A)\cdot J^{*}\right)^{T}=J\cdot P_k(-A)\cdot J^{*},\qquad
P(A)^{T}=P(-A),\\
&\left(J\cdot P_{k+1}(A)\cdot J^{*}\right)^{T}=J\cdot P_{k+1}(-A)\cdot J^{*},\\
&F_{y}(A)^{T}=\mathsf{C}^{-1}\cdot F_{-y}(-A)\cdot\mathsf{C}\qquad(y\in Y_{k}).
\end{aligned}
\end{equation}

\emph{The operators $D$, $Q_k\cdot J^{*}$, $Q$, $F_{y}$.}
The relations of (c) for $O=D$, for $Q_k(A)\cdot J^{*}$, for $Q(A)$ and for $O_y=F_y$ are contained in (5.3).

\emph{The operators $J\cdot P_k\cdot J^{*}$, $J\cdot P_{k+1}\cdot J^{*}$, $P$.}
These operators have scalar coefficients (Section~2.2).
By (Sca),
\begin{align*}
J\cdot P_k(-A)\cdot J^{*}&=\mathsf{C}^{-1}\cdot J\cdot P_k(-A)\cdot J^{*}\cdot\mathsf{C},\\
J\cdot P_{k+1}(-A)\cdot J^{*}&=\mathsf{C}^{-1}\cdot J\cdot P_{k+1}(-A)\cdot J^{*}\cdot\mathsf{C},\\
P(-A)&=\mathsf{C}^{-1}\cdot P(-A)\cdot\mathsf{C}.
\end{align*}
Together with (5.3) this gives $O(A)^{T}=\mathsf{C}^{-1}\cdot O(-A)\cdot\mathsf{C}$ for $O=J\cdot P_k\cdot J^{*}$, $J\cdot P_{k+1}\cdot J^{*}$, $P$.

\emph{The operator $\operatorname{Pen}_{y}$.}
One has
\[
\operatorname{Pen}_{y}(A)^{T}=F_{y}(A)^{T}-D_A^{T}-\mu=\mathsf{C}^{-1}\cdot\left(F_{-y}(-A)-D_{-A}-\mu\right)\cdot\mathsf{C}=\mathsf{C}^{-1}\cdot\operatorname{Pen}_{-y}(-A)\cdot\mathsf{C}.
\]

\emph{The composed operators.}
Lemma~3 is applied with ${}^{\mathrm{dg}}:={}^{T}$, $\mathrm{Nx}:=\mathsf{C}$ and $\tilde{A}:=-A$.
(Hz1) is the last relation of (5.3), at $y$ and at $\infty$ (recall that $-\infty=\infty$).
For (Hz2), the second relation of (5.3) reads $\left(Q_k(A)\cdot J^{*}\right)^{T}=J\cdot Q_k^{T}\left(-(-A)\right)=J\cdot Q_k^{T}(-\tilde{A})$;
the same relation at the connection $-A$ reads $\left(Q_k(-A)\cdot J^{*}\right)^{T}=J\cdot Q_k^{T}(-A)$, and its transpose is, by (Ru1),
\[
\left(J\cdot Q_k^{T}(-A)\right)^{T}=Q_k(-A)\cdot J^{*}=Q_k(\tilde{A})\cdot J^{*}.
\]
(Hz3) is contained in (5.3).
For (Hz4), (Ru2) and (I.c3) give
\[
\left(i\gamma_5\cdot y\right)^{T}=iy\cdot\gamma_5^{T}=iy\cdot\left(-\mathsf{C}^{-1}\cdot\gamma_5\cdot\mathsf{C}\right)=\mathsf{C}^{-1}\cdot\left(-i\gamma_5\cdot y\right)\cdot\mathsf{C}.
\]
The conclusions are
\begin{align*}
&S_{y}(A)^{T}=\mathsf{C}^{-1}\cdot S_{-y}(-A)\cdot\mathsf{C}\ \text{(Cz1)};\quad\text{at }y=\infty\text{ and at }y=0,\\
&\qquad S_k(A)^{T}=\mathsf{C}^{-1}\cdot S_k(-A)\cdot\mathsf{C},\qquad S^{0}_{k+1}(A)^{T}=\mathsf{C}^{-1}\cdot S^{0}_{k+1}(-A)\cdot\mathsf{C};\\
&H_{k}(A)^{T}=\mathsf{C}^{-1}\cdot\mathfrak{H}_{k}(-A)\cdot\mathsf{C}\ \text{(Cz2)};\qquad
\mathfrak{H}_{k}(A)^{T}=\mathsf{C}^{-1}\cdot H_{k}(-A)\cdot\mathsf{C}\ \text{(Cz3)};\\
&O(A)^{T}=\mathsf{C}^{-1}\cdot O(-A)\cdot\mathsf{C}\quad\text{for }O=D_k,\ M_k,\ \dot{M}_k\ \text{(Cz4)};\\
&B_{k,y}(A)^{T}=\mathsf{C}^{-1}\cdot B_{k,-y}(-A)\cdot\mathsf{C}\ \text{(Cz5)},\qquad
\Gamma_{k,y}(A)^{T}=\mathsf{C}^{-1}\cdot\Gamma_{k,-y}(-A)\cdot\mathsf{C}\ \text{(Cz6)}\qquad(y\text{ real});\\
&\qquad\text{at }y=0:\quad B_{k}(A)^{T}=\mathsf{C}^{-1}\cdot B_{k}(-A)\cdot\mathsf{C},\qquad \Gamma_k(A)^{T}=\mathsf{C}^{-1}\cdot\Gamma_k(-A)\cdot\mathsf{C}.
\end{align*}

\emph{The operator $\bar{H}_{k}$.}
By its definition and by (Cz3), $\bar{H}_{k}(A)=\mathsf{C}^{-1}\cdot H_{k}(-A)\cdot\mathsf{C}=\mathfrak{H}_{k}(A)^{T}$.
By (3.1), (Ru1) and (5.3),
\[
\mathfrak{H}_{k}(A)^{T}=\left(b_{k}\cdot Q_k(A)\cdot J^{*}\cdot S_{\infty}(A)\right)^{T}=b_{k}\cdot S_{\infty}(A)^{T}\cdot\left(Q_k(A)\cdot J^{*}\right)^{T}=b_{k}\cdot S_k(A)^{T}\cdot J\cdot Q_k^{T}(A).
\]

\emph{The operator $\bar{T}_k$.}
By its definition, by (3.1), by (Sca) for $Q_k(-A)\cdot J^{*}$, and by (5.3) at $y=\infty$,
\begin{align*}
\bar{T}_k(A)
&=\mathsf{C}^{-1}\cdot T_k(-A)\cdot\mathsf{C}
=b_{k}^{-1}\cdot\mathsf{C}^{-1}\cdot Q_k(-A)\cdot J^{*}\cdot F_{\infty}(-A)\cdot\mathsf{C}\\
&=b_{k}^{-1}\cdot Q_k(-A)\cdot J^{*}\cdot\left(\mathsf{C}^{-1}\cdot F_{\infty}(-A)\cdot\mathsf{C}\right)
=b_{k}^{-1}\cdot Q_k(-A)\cdot J^{*}\cdot F_{\infty}(A)^{T}\\
&=b_{k}^{-1}\cdot Q_k(-A)\cdot J^{*}\cdot\left(D_A+\mu+b_{k}\cdot J\cdot P_k(A)\cdot J^{*}\right)^{T}.
\end{align*}
\end{proof}

\subsubsection*{5.4. Proof of (d)}

\begin{proof}[Proof of Proposition~10\,(d)]
The operators to which (d) refers are those of the two lists of (b), namely
$D$, $Q_k\cdot J^{*}$, $J\cdot Q_k^{T}(-\,\cdot\,)$, $Q$, $Q^{T}(-\,\cdot\,)$, $J\cdot P_k\cdot J^{*}$, $P$, $J\cdot P_{k+1}\cdot J^{*}$, $S_k$, $S^{0}_{k+1}$, $H_{k}$, $\mathfrak{H}_{k}$, $T_k$, $D_k$, $M_k$, $\dot{M}_k$, $B_{k}$, $\Gamma_k$ and $\operatorname{Pen}_{y}$, $F_{y}$, $S_{y}$ ($y\in Y_{k}$), $B_{k,y}$, $\Gamma_{k,y}$ ($y$ real);
further $\delta_{\alpha,\,\cdot\,}$, $\bar{H}_{k}$ and $\bar{T}_k$.
For $O=J\cdot Q_k^{T}(-\,\cdot\,)$, that is, $O(A)=J\cdot Q_k^{T}(-A)$, the relation reads $\overline{J\cdot Q_k^{T}(-\overline{A})}=\varkappa^{-1}\cdot J\cdot Q_k^{T}(A)\cdot\varkappa$;
for $O=Q^{T}(-\,\cdot\,)$ it reads $\overline{Q^{T}(-\overline{A})}=\varkappa^{-1}\cdot Q^{T}(A)\cdot\varkappa$.
The number $\delta Z_k$, which (b) also names, is the subject of (e); regarded as a multiple of the identity, it satisfies the relation of (d) in the form $\overline{\delta Z_k(\overline{A})}=\delta Z_k(-A)$, because a number commutes with $\varkappa$, and this relation is proved in Section~5.5.
By (Dm), $\overline{A}$ and $-A$ lie in $\tilde{G}_k$.

\emph{The elementary relations.}
Proposition~2.2\,(d) is valid for every complex connection, and it gives
\begin{equation}\tag{5.4}
\begin{aligned}
&\overline{O(\overline{A})}=\varkappa^{-1}\cdot O(-A)\cdot\varkappa\qquad\text{for}\\
&O=D,\ Q_k\cdot J^{*},\ J\cdot Q_k^{T}(-\,\cdot\,),\ Q,\ Q^{T}(-\,\cdot\,),\ J\cdot P_k\cdot J^{*},\ P,\ J\cdot P_{k+1}\cdot J^{*},\ F_{y}\\
&\qquad(y\in Y_{k},\ \text{the same }y\text{ on both sides}).
\end{aligned}
\end{equation}

\emph{The operator $\operatorname{Pen}_{y}$.}
Since $\mu$ is real,
\[
\overline{\operatorname{Pen}_{y}(\overline{A})}=\overline{F_{y}(\overline{A})}-\overline{D_{\overline{A}}}-\mu=\varkappa^{-1}\cdot\left(F_{y}(-A)-D_{-A}-\mu\right)\cdot\varkappa=\varkappa^{-1}\cdot\operatorname{Pen}_{y}(-A)\cdot\varkappa.
\]

\emph{The operator $\delta_{\alpha}$} (for every $k\ge0$).
Let $f\in\mathfrak{S}^{\mathrm{ch}}(\eta_k,\ell)$ and $(x,x')\in\mathrm{Pzt}_k$.
By (2.1),
\begin{align*}
\left(\overline{\delta_{\alpha,\overline{A}}}\,f\right)(x,x')
&=\overline{d_X(x,x')^{-\alpha}\cdot\left(U_{\overline{A}}(x,x')\cdot\overline{f(x')}-\overline{f(x)}\right)}\\
&=d_X(x,x')^{-\alpha}\cdot\left(U_{-A}(x,x')\cdot f(x')-f(x)\right)
=\left(\delta_{\alpha,-A}f\right)(x,x')
\end{align*}
(for a pair inside a cluster the three phases that occur are equal to $1$).
The operator $\delta_{\alpha,-A}$ has scalar coefficients, so that by (Sca) it is equal to $\varkappa^{-1}\cdot\delta_{\alpha,-A}\cdot\varkappa$.
Hence $\overline{\delta_{\alpha,\overline{A}}}=\varkappa^{-1}\cdot\delta_{\alpha,-A}\cdot\varkappa$.

\emph{The composed operators.}
Lemma~2 is applied with the pair $(\overline{A},-A)$ in the place of $(A,\tilde{A})$, with $\tilde{y}:=y$, and with $\mathrm{Vf}:=\mathrm{Kf}$, $\mathrm{Vu}:=\mathrm{Kf}$; by Section~2.3 these are additive bijections that commute with the multiplication by real numbers.
By (Kf1), each relation of (5.4) is equivalent to $O(-A)\cdot\mathrm{Kf}=\mathrm{Kf}\cdot O(\overline{A})$.
Accordingly, (Hy1) is (5.4) for $F_{y}$, at $y$ and at $\infty$.
(Hy2) is (5.4) for $Q_k\cdot J^{*}$, which reads $Q_k(-A)\cdot J^{*}\cdot\mathrm{Kf}=\mathrm{Kf}\cdot Q_k(\overline{A})\cdot J^{*}$, together with (5.4) for $J\cdot Q_k^{T}(-\,\cdot\,)$, which reads $J\cdot Q_k^{T}\left(-(-A)\right)\cdot\mathrm{Kf}=\mathrm{Kf}\cdot J\cdot Q_k^{T}(-\overline{A})$.
(Hy3) is (5.4) for $P$, and (Hy4) is (Kf2), with $\tilde{y}=y$.
The conclusions (Co1)--(Co7) state that $O(-A)\cdot\mathrm{Kf}=\mathrm{Kf}\cdot O(\overline{A})$ for $O=S_{y}$ ($y\in Y_{k}$; hence for $S_k$ and $S^{0}_{k+1}$), $H_{k}$, $\mathfrak{H}_{k}$, $T_k$, $M_k$, $D_k$, $\dot{M}_k$, $B_{k,y}$, $\Gamma_{k,y}$ ($y$ real; hence for $B_{k}$ and $\Gamma_k$), with the same $y$ on both sides.
Read backwards, (Kf1) turns this into
\[
\overline{O(\overline{A})}=\varkappa^{-1}\cdot O(-A)\cdot\varkappa\qquad\text{for these operators.}
\]

\emph{The barred operators.}
The matrix $\mathsf{C}$ is real (I.c3); therefore (Ru1) and (Ru2) give $\overline{\mathsf{C}^{-1}\cdot O\cdot\mathsf{C}}=\mathsf{C}^{-1}\cdot\overline{O}\cdot\mathsf{C}$.
The relations for $H_{k}$ and $T_k$ are now used at the connection $-A$, which lies in $\tilde{G}_k$; note that $\overline{-A}=-\overline{A}$ and $-(-A)=A$.
Together with (Kf4) they give
\begin{align*}
\overline{\bar{H}_{k}(\overline{A})}
&=\overline{\mathsf{C}^{-1}\cdot H_{k}(-\overline{A})\cdot\mathsf{C}}
=\mathsf{C}^{-1}\cdot\overline{H_{k}\left(\overline{-A}\right)}\cdot\mathsf{C}
=\mathsf{C}^{-1}\cdot\varkappa^{-1}\cdot H_{k}(A)\cdot\varkappa\cdot\mathsf{C}\\
&=\varkappa^{-1}\cdot\mathsf{C}^{-1}\cdot H_{k}(A)\cdot\mathsf{C}\cdot\varkappa
=\varkappa^{-1}\cdot\bar{H}_{k}(-A)\cdot\varkappa,\\
\overline{\bar{T}_k(\overline{A})}
&=\overline{\mathsf{C}^{-1}\cdot T_k(-\overline{A})\cdot\mathsf{C}}
=\mathsf{C}^{-1}\cdot\overline{T_k\left(\overline{-A}\right)}\cdot\mathsf{C}
=\mathsf{C}^{-1}\cdot\varkappa^{-1}\cdot T_k(A)\cdot\varkappa\cdot\mathsf{C}\\
&=\varkappa^{-1}\cdot\mathsf{C}^{-1}\cdot T_k(A)\cdot\mathsf{C}\cdot\varkappa
=\varkappa^{-1}\cdot\bar{T}_k(-A)\cdot\varkappa;
\end{align*}
in each of the two chains the last step is the definition $\bar{H}_{k}(-A)=\mathsf{C}^{-1}\cdot H_{k}(A)\cdot\mathsf{C}$, respectively $\bar{T}_k(-A)=\mathsf{C}^{-1}\cdot T_k(A)\cdot\mathsf{C}$.
\end{proof}

\subsubsection*{5.5. Proof of (e)}

\begin{proof}[Proof of Proposition~10\,(e)]
Take $y=0$ in the first relation of (Cz6), in the form in which it was obtained in Section~5.3:
\[
\left(D_k(A)+b_{L}\cdot P(A)\right)^{T}=\mathsf{C}^{-1}\cdot\left(D_k(-A)+b_{L}\cdot P(-A)\right)\cdot\mathsf{C}.
\]
By (Ru4)\,(iii) and (i) it follows that
\[
\delta Z_k(A)=\det\left(\left(D_k(A)+b_{L}\cdot P(A)\right)^{T}\right)=\det\left(D_k(-A)+b_{L}\cdot P(-A)\right)=\delta Z_k(-A).
\]
Next take $y=0$ in the first relation of (Co7), in the form in which it was obtained in Section~5.4:
\[
\left(D_k(-A)+b_{L}\cdot P(-A)\right)\cdot\mathrm{Kf}=\mathrm{Kf}\cdot\left(D_k\left(\overline{A}\right)+b_{L}\cdot P\left(\overline{A}\right)\right),
\]
that is,
\[
D_k(-A)+b_{L}\cdot P(-A)=\mathrm{Kf}\cdot\left(D_k\left(\overline{A}\right)+b_{L}\cdot P\left(\overline{A}\right)\right)\cdot\mathrm{Kf}^{-1}.
\]
Now (Kf3) gives $\delta Z_k(-A)=\overline{\delta Z_k\left(\overline{A}\right)}$.
Combined with the identity obtained first, this yields $\overline{\delta Z_k\left(\overline{A}\right)}=\delta Z_k(A)$.
\end{proof}

\subsubsection*{5.6. Proof of (f)}

\begin{proof}[Proof of Proposition~10\,(f)]
Let $A\in\tilde{G}_k$ be real.
By Proposition~2.2\,(e) one has
\begin{equation}\tag{5.5}
\begin{aligned}
&D_A^{*}=\gamma_5\cdot D_A\cdot\gamma_5,\qquad
\left(Q_k(A)\cdot J^{*}\right)^{*}=J\cdot Q_k^{T}(-A),\qquad
P(A)^{*}=P(A),\\
&F_{y}(A)^{*}=\gamma_5\cdot F_{-y}(A)\cdot\gamma_5\qquad(y\in Y_k);
\end{aligned}
\end{equation}
the third relation holds because $P(A)$ is an orthogonal projection.
For $O=D$ and for $O_{y}=F_{y}$ the relations of (f) are among those of (5.5).

Lemma~3 is applied with ${}^{\mathrm{dg}}:={}^{*}$, with $\mathrm{Nx}:=\gamma_5$, for which $\gamma_5^{-1}=\gamma_5$ by (I.c1), and with $\tilde{A}:=A$.
Its hypotheses are satisfied.
(Hz1) is the last relation of (5.5), taken at $y$ and at $\infty$.
(Hz2) consists of the second relation of (5.5) and of its adjoint $\left(J\cdot Q_k^{T}(-A)\right)^{*}=Q_k(A)\cdot J^{*}$, which follows by (Ru1).
(Hz3) is contained in (5.5).
(Hz4) follows from (Ru2) and (I.c1), since $\gamma_5$ is Hermitian and $\gamma_5^{2}=1$:
\[
\left(i\gamma_5\cdot y\right)^{*}=-iy\cdot\gamma_5=\gamma_5\cdot\left(-i\gamma_5\cdot y\right)\cdot\gamma_5.
\]
The conclusions of Lemma~3 are therefore the following.
By (Cz1),
\[
S_{y}(A)^{*}=\gamma_5\cdot S_{-y}(A)\cdot\gamma_5,
\]
and in particular, at $y=\infty$ and at $y=0$,
\[
S_k(A)^{*}=\gamma_5\cdot S_k(A)\cdot\gamma_5,\qquad
S^{0}_{k+1}(A)^{*}=\gamma_5\cdot S^{0}_{k+1}(A)\cdot\gamma_5.
\]
By (Cz2) and (Cz4),
\[
H_{k}(A)^{*}=\gamma_5\cdot\mathfrak{H}_{k}(A)\cdot\gamma_5,\qquad
O(A)^{*}=\gamma_5\cdot O(A)\cdot\gamma_5\quad\text{for }O=D_k,\ M_k,\ \dot{M}_k.
\]
By (Cz5) and (Cz6), for real $y$,
\[
B_{k,y}(A)^{*}=\gamma_5\cdot B_{k,-y}(A)\cdot\gamma_5,\qquad
\Gamma_{k,y}(A)^{*}=\gamma_5\cdot\Gamma_{k,-y}(A)\cdot\gamma_5,
\]
and at $y=0$
\[
B_k(A)^{*}=\gamma_5\cdot B_k(A)\cdot\gamma_5,\qquad
\Gamma_k(A)^{*}=\gamma_5\cdot\Gamma_k(A)\cdot\gamma_5.
\]
In addition, since $\mu$ is real and $\operatorname{Pen}_{y}(A)=F_{y}(A)-D_A-\mu$, subtracting the first relation of (5.5) from the last one gives
\[
\operatorname{Pen}_{y}(A)^{*}=\gamma_5\cdot\operatorname{Pen}_{-y}(A)\cdot\gamma_5;
\]
and the first relation of (Cz6) at $y=0$ gives
\[
\left(D_k(A)+b_{L}\cdot P(A)\right)^{*}=\gamma_5\cdot\left(D_k(A)+b_{L}\cdot P(A)\right)\cdot\gamma_5.
\]
\end{proof}

\subsection*{6. Proof of (g) (level 0)}

\subsubsection*{6.1. What is different, and what (g) says}

Let $k=0$.
Then $\eta_0=1$ and $\mathbb{L}_0=\mathbb{Z}^3$, and the fine lattice and the unit lattice coincide.
The spinor spaces are $\mathfrak{S}^{\mathrm{ch}}(1,\ell)$, with $192$ components per photon site and the pairing $\langle\,\cdot\,,\cdot\,\rangle^{\mathrm{ch}}$, which carries the weight $\frac{1}{48}$, and $\mathfrak{S}(L,\ell)$, with $4$ components and the pairing $\langle\,\cdot\,,\cdot\,\rangle_L$.
The parameter set is $Y_0=\mathbb{R}$, and its only neutral value is $y=0$.
Section~I.3.9, (I.L0.7) and (I.L0.8) give
\begin{equation}\tag{6.1}
\begin{aligned}
&\operatorname{Pen}_{y}(A)=b_{L}\cdot J\cdot P(A)\cdot J^{*}+i\gamma_5\cdot y,\qquad
F_{y}(A)=D_A+\mu+\operatorname{Pen}_{y}(A),\\
&S_{0,y}(A)=F_{y}(A)^{-1},\qquad
\Gamma_0(A)=S_{0,0}(A)=F_0(A)^{-1},\qquad
D_0(A)=D_A+\mu,\\
&\delta Z_0(A)=\det F_0(A)\quad(\text{of size }192\,n_0),\\
&H_0(A)=\bar{H}_0(A)=T_0(A)=\bar{T}_0(A)=1\quad\text{on }\mathfrak{S}^{\mathrm{ch}}(1,\ell).
\end{aligned}
\end{equation}
The objects $b_0$, $S_k(A)=S_{\infty}(A)$ at $k=0$, $M_0(A)$, $B_{0,y}(A)$, $\Gamma_{0,y}(A)$ and $\mathfrak{H}_0(A)$ are not defined (Section~I.3.9 and (I.L0.9); the operator $\mathfrak{H}_k$ is defined in Section~I.3.5 for $k\ge1$ only).
Neither $\dot{M}_0(A)=1$ nor $S^{0}_{1}(A)$, the second name of $\Gamma_0(A)$ (Section~I.3.9, (I.L0.7)), belongs to the list of (g); for these two operators the relations of (a)--(f) are the relations below for the identity $H_0(A)=1$ and for $\Gamma_0(A)$, respectively.
Sections~2 and 3.1--3.3 are valid for every $k\ge0$.

According to the reading rule stated in Proposition~10\,(g), the assertion that (a)--(f) hold at $k=0$ for the list given there consists of the relations (g-a)--(g-f) below.
These relations arise from (a)--(f) through the substitutions stated in (g) ((I.3.4), Section~I.3.9), in which the arrow means ``is replaced by'':
\begin{gather*}
Q(A)\rightsquigarrow Q(A)\cdot J^{*};\qquad
Q^{T}(-A)\rightsquigarrow J\cdot Q^{T}(-A);\\
P(A),\ J\cdot P_{k+1}(A)\cdot J^{*}\rightsquigarrow J\cdot P(A)\cdot J^{*};\qquad
D_k(A)\rightsquigarrow D_0(A);\qquad
S_{y}\rightsquigarrow S_{0,y};\\
S^{0}_{k+1},\ \Gamma_k\rightsquigarrow\Gamma_0;\qquad
H_k,\ \bar{H}_k,\ T_k,\ \bar{T}_k\rightsquigarrow1;\qquad
\mathrm{uf}^{0}\rightsquigarrow\mathrm{uf}.
\end{gather*}
The last substitution reflects that $T^{0}_{0}=T_0$ and that the operators ``of the unit lattice'' act on $\mathfrak{S}^{\mathrm{ch}}(1,\ell)$; the function $\mathrm{uf}^{1}$ stays on $\mathfrak{S}(L,\ell)$.
(The operators $Q_k(A)\cdot J^{*}$, $J\cdot Q_k^{T}(-A)$ and $J\cdot P_k(A)\cdot J^{*}$ become $Q_0=P_0=1$, which does not belong to the list.)
In (g-c) the relation $\left(Q(A)\cdot J^{*}\right)^{T}=J\cdot Q^{T}(A)$ takes the place of $Q(A)^{T}=Q^{T}(A)$.
In (g-f) the operator $\operatorname{Pen}_{y}$ is included, since it belongs to the list.
The data are $A\in\tilde{G}_0$ and a real number $y$.
\begin{itemize}
\item[(g-a)]
Let $\mathrm{Acl}$ be a real connection on $T_0$ with $d\,\mathrm{Acl}=0$, and let $\mathrm{uf}$ be a function of $c$ on $T_0$ with $|\mathrm{uf}|=1$ and
\[
\mathrm{uf}(c')=\mathrm{uf}(c)\cdot\exp\left\{ie_k\cdot\mathrm{Acl}(c,c')\right\}
\]
on every bond; let $\mathrm{uf}^{1}$ be its restriction to $L\mathbb{Z}^3/\ell\mathbb{Z}^3$.
Then $A-\mathrm{Acl}$ lies in $\tilde{G}_0$, and
\begin{align*}
&D_{A-\mathrm{Acl}}=\mathrm{uf}\cdot D_A\cdot\mathrm{uf}^{-1},\qquad
Q(A-\mathrm{Acl})\cdot J^{*}=\mathrm{uf}^{1}\cdot Q(A)\cdot J^{*}\cdot\mathrm{uf}^{-1},\\
&J\cdot Q^{T}\left(-(A-\mathrm{Acl})\right)=\mathrm{uf}\cdot J\cdot Q^{T}(-A)\cdot\left(\mathrm{uf}^{1}\right)^{-1},\\
&O(A-\mathrm{Acl})=\mathrm{uf}\cdot O(A)\cdot\mathrm{uf}^{-1}\quad\text{for }O=J\cdot P\cdot J^{*},\ D_0,\ \operatorname{Pen}_{y},\ F_{y},\ S_{0,y},\ \Gamma_0,\ H_0,\ T_0,\\
&O(A-\mathrm{Acl})=\mathrm{uf}^{-1}\cdot O(A)\cdot\mathrm{uf}\quad\text{for }O=\bar{H}_0,\ \bar{T}_0,\\
&\left(\delta_{\alpha,A-\mathrm{Acl}}\left(\mathrm{uf}\cdot f\right)\right)(x,x')=\mathrm{uf}(c)\cdot\left(\delta_{\alpha,A}f\right)(x,x'),\qquad
\delta Z_0(A-\mathrm{Acl})=\delta Z_0(A).
\end{align*}
\item[(g-b)]
For $\tilde{g}=(S,a)\in\mathbb{G}_0$ one has
\begin{align*}
&O(\tilde{g}A)\cdot\tilde{g}=\tilde{g}\cdot O(A)\quad\text{for }O=D,\ Q\cdot J^{*},\ J\cdot Q^{T}(-\,\cdot\,),\ J\cdot P\cdot J^{*},\ D_0,\ \Gamma_0,\ H_0,\ T_0;\\
&O_{y'}(\tilde{g}A)\cdot\tilde{g}=\tilde{g}\cdot O_{y}(A)\quad\text{for }O_{y}=\operatorname{Pen}_{y},\ F_{y},\ S_{0,y},\\
&\qquad\text{where }y'=y\text{ if }S\in\mathrm{Sp}_3\text{ or }y=0,\text{ and }y'=-y\text{ otherwise};\\
&\delta Z_0(\tilde{g}A)=\delta Z_0(A);\qquad
\bar{H}_0(\tilde{g}A)\cdot\mathrm{glb}=\mathrm{glb}\cdot\bar{H}_0(A),\qquad
\bar{T}_0(\tilde{g}A)\cdot\mathrm{glb}=\mathrm{glb}\cdot\bar{T}_0(A);\\
&\left(\delta_{\alpha,\tilde{g}A}\,\tilde{g}f\right)(\tilde{g}x,\tilde{g}x')=S\cdot\left(\delta_{\alpha,A}f\right)(x,x').
\end{align*}
\item[(g-c)]
One has
\begin{align*}
&O(A)^{T}=\mathsf{C}^{-1}\cdot O(-A)\cdot\mathsf{C}\quad\text{for }O=D,\ J\cdot P\cdot J^{*},\ D_0,\ \Gamma_0,\ H_0,\ T_0;\\
&O_{y}(A)^{T}=\mathsf{C}^{-1}\cdot O_{-y}(-A)\cdot\mathsf{C}\quad\text{for }O_{y}=\operatorname{Pen}_{y},\ F_{y},\ S_{0,y};\\
&\left(Q(A)\cdot J^{*}\right)^{T}=J\cdot Q^{T}(A);\qquad
\bar{H}_0(A)=\mathsf{C}^{-1}\cdot H_0(-A)\cdot\mathsf{C},\qquad
\bar{T}_0(A)=\mathsf{C}^{-1}\cdot T_0(-A)\cdot\mathsf{C}.
\end{align*}
\item[(g-d)]
One has
\[
\overline{O\left(\overline{A}\right)}=\varkappa^{-1}\cdot O(-A)\cdot\varkappa
\]
for $O=D$, $Q\cdot J^{*}$, $J\cdot Q^{T}(-\,\cdot\,)$, $J\cdot P\cdot J^{*}$, $D_0$, $\Gamma_0$, $H_0$, $T_0$, $\bar{H}_0$, $\bar{T}_0$, $\delta_{\alpha,\cdot}$, and for $O=\operatorname{Pen}_{y}$, $F_{y}$, $S_{0,y}$ with the same $y$ on both sides.
\item[(g-e)]
$\delta Z_0(-A)=\delta Z_0(A)$ and $\overline{\delta Z_0\left(\overline{A}\right)}=\delta Z_0(A)$.
\item[(g-f)]
For real $A\in\tilde{G}_0$,
\begin{gather*}
O(A)^{*}=\gamma_5\cdot O(A)\cdot\gamma_5\quad\text{for }O=D,\ D_0,\ \Gamma_0,\ H_0;\\
O_{y}(A)^{*}=\gamma_5\cdot O_{-y}(A)\cdot\gamma_5\quad\text{for }O_{y}=\operatorname{Pen}_{y},\ F_{y},\ S_{0,y}.
\end{gather*}
\end{itemize}
Four of the relations of (c) and (f) have on one side an operator for which a substitution at $k=0$ is given above, and on the other side $\mathfrak{H}_k$, $b_k$ or $S_k$, objects that are not defined at $k=0$ (I.L0.9):
in (c) the relation $H_k^{T}=\mathsf{C}^{-1}\cdot\mathfrak{H}_k(-\,\cdot\,)\cdot\mathsf{C}$, the formula for $\bar{H}_k$ and the formula for $\bar{T}_k$;
in (f) the relation $H_k^{*}=\gamma_5\cdot\mathfrak{H}_k\cdot\gamma_5$.
At $k=0$ these four relations have no counterpart; what remains of them are the relations for $H_0=T_0=1$ in (g-c) and (g-f).
The entries $O=S_k$ of the lists in (c) and (f), and the value $y=\infty$ of the family $S_{y}$, are not among these four.
They fall away because $S_k$ does not belong to the list of (g) and is not defined at $k=0$ (I.L0.9), in the same way as $M_k$, $B_k$, $B_{k,y}$ and $\Gamma_{k,y}$.

The relations displayed in Proposition~10\,(g) are contained in this list:
$S_{0,y}(A-\mathrm{Acl})=\mathrm{uf}\cdot S_{0,y}(A)\cdot\mathrm{uf}^{-1}$, $\Gamma_0(A-\mathrm{Acl})=\mathrm{uf}\cdot\Gamma_0(A)\cdot\mathrm{uf}^{-1}$ and $\delta Z_0(A-\mathrm{Acl})=\delta Z_0(A)$ in (g-a);
$\delta Z_0(\tilde{g}A)=\delta Z_0(A)$ in (g-b);
and $\delta Z_0(-A)=\delta Z_0(A)=\overline{\delta Z_0\left(\overline{A}\right)}$ in (g-e).

\subsubsection*{6.2. Proofs}

\begin{proof}[Proof of Proposition~10\,(g)]
Proposition~2.2 is applied on the block torus $T_0=(0,\ell,-1)$ at level $0$, with the pairs in its clauses (a) and (b) as in Section~5.
Every connection that occurs below lies in $\tilde{G}_0$ by (Dm), and by (Iv)\,(i) the operator $F_{y}(\,\cdot\,)$ is bijective at each of these connections for every real $y$.
The relations that concern $H_0$, $T_0$, $\bar{H}_0$, $\bar{T}_0$ state that the identity of $\mathfrak{S}^{\mathrm{ch}}(1,\ell)$ commutes with $\mathrm{uf}$, $\tilde{g}$, $\mathrm{glb}$, $\mathsf{C}$, $\varkappa$, $\gamma_5$ and coincides with its transpose, its conjugate and its adjoint.
Indeed $1\cdot\mathrm{Vm}=\mathrm{Vm}\cdot1$ for every map $\mathrm{Vm}$; further $1^{T}=1$, $\overline{1}=\mathrm{Cx}\cdot\mathrm{Cx}=1$ and $1^{*}=1$ by Section~2.1, and $\gamma_5\cdot1\cdot\gamma_5=\gamma_5^{2}=1$ by (I.c1).
The relation $\mathsf{C}^{-1}\cdot1\cdot\mathsf{C}=1$ expresses the agreement of (6.1) with the definitions $\bar{H}_k(A)=\mathsf{C}^{-1}\cdot H_k(-A)\cdot\mathsf{C}$ and $\bar{T}_k(A)=\mathsf{C}^{-1}\cdot T_k(-A)\cdot\mathsf{C}$ of Section~I.3.5.
This settles every relation of (g-a)--(g-f) that contains $H_0$, $T_0$, $\bar{H}_0$ or $\bar{T}_0$.

\emph{(g-a).}
By Lemma~1\,(ii) with $k=0$, the connection $A-\mathrm{Acl}$ lies in $\tilde{G}_0$.
Proposition~2.2\,(a) at $k=0$ gives
\begin{gather*}
D_{A-\mathrm{Acl}}=\mathrm{uf}\cdot D_A\cdot\mathrm{uf}^{-1},\qquad
Q(A-\mathrm{Acl})\cdot J^{*}=\mathrm{uf}^{1}\cdot Q(A)\cdot J^{*}\cdot\mathrm{uf}^{-1},\\
J\cdot Q^{T}\left(-(A-\mathrm{Acl})\right)=\mathrm{uf}\cdot J\cdot Q^{T}(-A)\cdot\left(\mathrm{uf}^{1}\right)^{-1},\\
J\cdot P(A-\mathrm{Acl})\cdot J^{*}=\mathrm{uf}\cdot J\cdot P(A)\cdot J^{*}\cdot\mathrm{uf}^{-1},\qquad
O(A-\mathrm{Acl})=\mathrm{uf}\cdot O(A)\cdot\mathrm{uf}^{-1}\quad\text{for }O=\operatorname{Pen}_{y},\ F_{y},
\end{gather*}
together with the relation for $\delta_{\alpha}$.
For $D_0$ one has $\mathrm{uf}\cdot\left(D_A+\mu\right)\cdot\mathrm{uf}^{-1}=\mathrm{uf}\cdot D_A\cdot\mathrm{uf}^{-1}+\mu$.
For $S_{0,y}$,
\[
S_{0,y}(A-\mathrm{Acl})=\left(\mathrm{uf}\cdot F_{y}(A)\cdot\mathrm{uf}^{-1}\right)^{-1}=\mathrm{uf}\cdot S_{0,y}(A)\cdot\mathrm{uf}^{-1},
\]
and $y=0$ gives the relation for $\Gamma_0$.
For the determinant, (Ru4)\,(i) on $\mathfrak{S}^{\mathrm{ch}}(1,\ell)$ gives
\[
\delta Z_0(A-\mathrm{Acl})=\det\left(\mathrm{uf}\cdot F_0(A)\cdot\mathrm{uf}^{-1}\right)=\det F_0(A).
\]
The verification of the two examples in Section~5.1 is valid for every $k\ge0$.

\emph{(g-b).}
Proposition~2.2\,(b) at $k=0$, applied to the morphism $c\mapsto rc+a$ and to $\mathrm{mS}=\tilde{g}$, gives
\begin{gather*}
D_{\tilde{g}A}\cdot\tilde{g}=\tilde{g}\cdot D_A,\qquad
Q(\tilde{g}A)\cdot J^{*}\cdot\tilde{g}=\tilde{g}\cdot Q(A)\cdot J^{*},\\
J\cdot Q^{T}(-\tilde{g}A)\cdot\tilde{g}=\tilde{g}\cdot J\cdot Q^{T}(-A),\qquad
F_{y'}(\tilde{g}A)\cdot\tilde{g}=\tilde{g}\cdot F_{y}(A),
\end{gather*}
together with the relation for $\delta_{\alpha}$; in the second and in the third relation one of the two factors $\tilde{g}$ acts on $\mathfrak{S}^{\mathrm{ch}}(1,\ell)$ and the other on $\mathfrak{S}(L,\ell)$.
For $J\cdot P\cdot J^{*}$, Section~2.2 gives
\begin{align*}
J\cdot P(\tilde{g}A)\cdot J^{*}\cdot\tilde{g}
&=\left(J\cdot Q^{T}(-\tilde{g}A)\right)\cdot\left(Q(\tilde{g}A)\cdot J^{*}\right)\cdot\tilde{g}
=\left(J\cdot Q^{T}(-\tilde{g}A)\right)\cdot\tilde{g}\cdot\left(Q(A)\cdot J^{*}\right)\\
&=\tilde{g}\cdot\left(J\cdot Q^{T}(-A)\right)\cdot\left(Q(A)\cdot J^{*}\right)
=\tilde{g}\cdot J\cdot P(A)\cdot J^{*}.
\end{align*}
For $D_0$ one has $\left(D_{\tilde{g}A}+\mu\right)\cdot\tilde{g}=\tilde{g}\cdot\left(D_A+\mu\right)$.
For $\operatorname{Pen}_{y}$,
\begin{align*}
\operatorname{Pen}_{y'}(\tilde{g}A)\cdot\tilde{g}
&=\left(F_{y'}(\tilde{g}A)-D_{\tilde{g}A}-\mu\right)\cdot\tilde{g}\\
&=\tilde{g}\cdot\left(F_{y}(A)-D_A-\mu\right)=\tilde{g}\cdot\operatorname{Pen}_{y}(A).
\end{align*}
For $S_{0,y}$, (Ru5) with $\tilde{O}=F_{y'}(\tilde{g}A)$, $O=F_{y}(A)$ and $\mathrm{Vm}=\tilde{g}$ gives $S_{0,y'}(\tilde{g}A)\cdot\tilde{g}=\tilde{g}\cdot S_{0,y}(A)$.
At $y=0$, where $y'=0$ for every $S\in\mathrm{Sp}$, this is $\Gamma_0(\tilde{g}A)\cdot\tilde{g}=\tilde{g}\cdot\Gamma_0(A)$.
For the determinant, $F_0(\tilde{g}A)=\tilde{g}\cdot F_0(A)\cdot\tilde{g}^{-1}$, hence $\delta Z_0(\tilde{g}A)=\delta Z_0(A)$ by (Ru4)\,(i).

\emph{(g-c).}
Proposition~2.2\,(c) at $k=0$ gives
\begin{gather*}
D_A^{T}=\mathsf{C}^{-1}\cdot D_{-A}\cdot\mathsf{C},\qquad
\left(Q(A)\cdot J^{*}\right)^{T}=J\cdot Q^{T}(A),\\
\left(J\cdot P(A)\cdot J^{*}\right)^{T}=J\cdot P(-A)\cdot J^{*},\qquad
F_{y}(A)^{T}=\mathsf{C}^{-1}\cdot F_{-y}(-A)\cdot\mathsf{C}.
\end{gather*}
The operator $J\cdot P(-A)\cdot J^{*}$ has scalar coefficients (Section~2.2); by (Sca) it equals $\mathsf{C}^{-1}\cdot J\cdot P(-A)\cdot J^{*}\cdot\mathsf{C}$.
Further,
\begin{gather*}
D_0(A)^{T}=D_A^{T}+\mu=\mathsf{C}^{-1}\cdot\left(D_{-A}+\mu\right)\cdot\mathsf{C},\\
\operatorname{Pen}_{y}(A)^{T}=F_{y}(A)^{T}-D_A^{T}-\mu=\mathsf{C}^{-1}\cdot\operatorname{Pen}_{-y}(-A)\cdot\mathsf{C},
\end{gather*}
and, by (Ru1),
\[
S_{0,y}(A)^{T}=\left(F_{y}(A)^{T}\right)^{-1}=\left(\mathsf{C}^{-1}\cdot F_{-y}(-A)\cdot\mathsf{C}\right)^{-1}=\mathsf{C}^{-1}\cdot S_{0,-y}(-A)\cdot\mathsf{C};
\]
$y=0$ gives the relation for $\Gamma_0$.

\emph{(g-d).}
For $O=D$, Proposition~2.2\,(d) at $k=0$ gives $\overline{D_{\overline{A}}}=\varkappa^{-1}\cdot D_{-A}\cdot\varkappa$; this relation is also (I.O5).
The three averaging operators of level $0$ are treated directly from their displays in Section~2.2, with the help of (2.1).
Let $f\in\mathfrak{S}^{\mathrm{ch}}(1,\ell)$, $v'\in\mathfrak{S}(L,\ell)$, $z'\in L\mathbb{Z}^3$ and $z\in B(z')$.
Then
\begin{align*}
\left(\overline{Q\left(\overline{A}\right)\cdot J^{*}}\,f\right)(z')
&=\overline{L^{-3}\sum_{z\in B(z')}U_{\overline{A}}(z',z)\cdot\frac{1}{48}\sum_{g\in O_{\mathrm{cub}}}\overline{f(z,g)}}\\
&=L^{-3}\sum_{z\in B(z')}U_{-A}(z',z)\cdot\frac{1}{48}\sum_{g\in O_{\mathrm{cub}}}f(z,g)
=\left(Q(-A)\cdot J^{*}f\right)(z'),\\
\left(\overline{J\cdot Q^{T}\left(-\overline{A}\right)}\,v'\right)(z,g)
&=\overline{U_{\overline{A}}(z',z)^{-1}\cdot\overline{v'(z')}}
=U_{-A}(z',z)^{-1}\cdot v'(z')
=\left(J\cdot Q^{T}\left(-(-A)\right)v'\right)(z,g).
\end{align*}
The operators $Q(-A)\cdot J^{*}$ and $J\cdot Q^{T}(A)$ have scalar coefficients; by (Sca) they equal $\varkappa^{-1}\cdot Q(-A)\cdot J^{*}\cdot\varkappa$ and $\varkappa^{-1}\cdot J\cdot Q^{T}(A)\cdot\varkappa$, respectively.
This proves (g-d) for $Q\cdot J^{*}$ and for $J\cdot Q^{T}(-\,\cdot\,)$.
By (Ru1), and by (Sca) in the last step,
\begin{align*}
\overline{J\cdot P\left(\overline{A}\right)\cdot J^{*}}
&=\overline{J\cdot Q^{T}\left(-\overline{A}\right)}\cdot\overline{Q\left(\overline{A}\right)\cdot J^{*}}
=J\cdot Q^{T}(A)\cdot Q(-A)\cdot J^{*}\\
&=J\cdot P(-A)\cdot J^{*}
=\varkappa^{-1}\cdot J\cdot P(-A)\cdot J^{*}\cdot\varkappa.
\end{align*}
For $D_0$: since $\mu$ is real, $\overline{D_0\left(\overline{A}\right)}=\overline{D_{\overline{A}}}+\mu=\varkappa^{-1}\cdot\left(D_{-A}+\mu\right)\cdot\varkappa$.
For $\operatorname{Pen}_{y}$: as $y$ is real, (Ru2) and (I.c4) give
\[
\overline{i\gamma_5\cdot y}=-iy\cdot\overline{\gamma_5}=-iy\cdot\left(-\varkappa^{-1}\cdot\gamma_5\cdot\varkappa\right)=\varkappa^{-1}\cdot\left(i\gamma_5\cdot y\right)\cdot\varkappa,
\]
and since $b_{L}$ is real, (6.1) and (Ru1) yield
\begin{align*}
\overline{\operatorname{Pen}_{y}\left(\overline{A}\right)}
&=b_{L}\cdot\overline{J\cdot P\left(\overline{A}\right)\cdot J^{*}}+\overline{i\gamma_5\cdot y}
=\varkappa^{-1}\cdot\left(b_{L}\cdot J\cdot P(-A)\cdot J^{*}+i\gamma_5\cdot y\right)\cdot\varkappa\\
&=\varkappa^{-1}\cdot\operatorname{Pen}_{y}(-A)\cdot\varkappa.
\end{align*}
For $F_{y}$,
\begin{align*}
\overline{F_{y}\left(\overline{A}\right)}
&=\overline{D_{\overline{A}}}+\mu+\overline{\operatorname{Pen}_{y}\left(\overline{A}\right)}
=\varkappa^{-1}\cdot\left(D_{-A}+\mu+\operatorname{Pen}_{y}(-A)\right)\cdot\varkappa\\
&=\varkappa^{-1}\cdot F_{y}(-A)\cdot\varkappa,
\end{align*}
with the same $y$ on both sides.
For $\delta_{\alpha,\cdot}$, the display in Section~5.4 is valid for every $k\ge0$: at $k=0$ one has $d_X\ge1$ for all pairs, and inside a cluster $d_X=1$ and no phase occurs (I.L0.14).
For $S_{0,y}$, (Ru1) gives
\[
\overline{S_{0,y}\left(\overline{A}\right)}=\left(\overline{F_{y}\left(\overline{A}\right)}\right)^{-1}=\left(\varkappa^{-1}\cdot F_{y}(-A)\cdot\varkappa\right)^{-1}=\varkappa^{-1}\cdot S_{0,y}(-A)\cdot\varkappa,
\]
and $y=0$ gives the relation for $\Gamma_0$.

\emph{(g-e).}
By (g-c) at $y=0$, $F_0(A)^{T}=\mathsf{C}^{-1}\cdot F_0(-A)\cdot\mathsf{C}$; hence, by (Ru4)\,(iii) and (i),
\[
\delta Z_0(A)=\det F_0(A)^{T}=\det F_0(-A)=\delta Z_0(-A).
\]
By (g-d) at $y=0$, $\overline{F_0\left(\overline{A}\right)}=\varkappa^{-1}\cdot F_0(-A)\cdot\varkappa$; hence, by (Ru4)\,(ii) and (i),
\[
\overline{\delta Z_0\left(\overline{A}\right)}=\det\left(\overline{F_0\left(\overline{A}\right)}\right)=\det F_0(-A)=\delta Z_0(-A)=\delta Z_0(A).
\]

\emph{(g-f).}
Let $A$ be real.
Proposition~2.2\,(e) gives $D_A^{*}=\gamma_5\cdot D_A\cdot\gamma_5$ and $F_{y}(A)^{*}=\gamma_5\cdot F_{-y}(A)\cdot\gamma_5$.
As $\mu$ is real and $\gamma_5^{2}=1$,
\begin{gather*}
D_0(A)^{*}=D_A^{*}+\mu=\gamma_5\cdot\left(D_A+\mu\right)\cdot\gamma_5,\\
\operatorname{Pen}_{y}(A)^{*}=F_{y}(A)^{*}-D_A^{*}-\mu=\gamma_5\cdot\operatorname{Pen}_{-y}(A)\cdot\gamma_5.
\end{gather*}
By (Ru1),
\[
S_{0,y}(A)^{*}=\left(F_{y}(A)^{*}\right)^{-1}=\left(\gamma_5\cdot F_{-y}(A)\cdot\gamma_5\right)^{-1}=\gamma_5\cdot S_{0,-y}(A)\cdot\gamma_5,
\]
and $y=0$ gives $\Gamma_0(A)^{*}=\gamma_5\cdot\Gamma_0(A)\cdot\gamma_5$.
\end{proof}

\subsection*{7. Remarks}

\begin{named}{Remark 1}[$d\,\mathrm{Acl}=0$ cannot be dropped in (a)]
Let $k=0$, $A=0$ and $\mathrm{uf}=1$, and let $\mathrm{Acl}$ take the value $2\pi/e_k$ on the bond $\left((0,1,0),(1,1,0)\right)$, the value $-2\pi/e_k$ on the reversed bond, and the value $0$ on every other bond.
Then $\mathrm{uf}(c')=\mathrm{uf}(c)\cdot\exp\left\{ie_k\cdot\mathrm{Acl}(c,c')\right\}$ on every bond, since $\exp(\pm2\pi i)=1$, whereas $d\,\mathrm{Acl}\ne0$: its modulus is $2\pi/e_k$ on each of the four plaquettes that contain the bond.
As $L\ge5$, the site $(1,1,0)$ lies in $B(0)$.
Among the six paths $\boldsymbol{\Gamma}^{\pi}\left((0,0,0),(1,1,0)\right)$, the bond belongs to exactly three, namely to those that bring the coordinate $2$ to its final value before the coordinate $1$.
Hence
\[
e_k\cdot(\tau\mathrm{Acl})\left((0,0,0),(1,1,0)\right)=\frac{1}{6}\cdot3\cdot2\pi=\pi,\qquad
U_{-\mathrm{Acl}}\left((0,0,0),(1,1,0)\right)=-1.
\]
By the display of Section~2.2, the coefficient of $f\left((1,1,0),g\right)$ in $\left(Q(-\mathrm{Acl})\cdot J^{*}f\right)(0)$ is $-L^{-3}/48$ for each of the $48$ labels $g$, while in $\left(Q(0)\cdot J^{*}f\right)(0)$ it is $+L^{-3}/48$.
Consequently $Q(A-\mathrm{Acl})\cdot J^{*}\ne\mathrm{uf}^{1}\cdot Q(A)\cdot J^{*}\cdot\mathrm{uf}^{-1}$.
\end{named}

\begin{named}{Remark 2}[the sign of $y$ cannot be dropped in (c), nor in (b) for $S$ outside $\mathrm{Sp}_3$]
Let $k\ge1$, let $y$ be real with $y\ne0$, and let $A=0$; the connection $0$ lies in $\tilde{G}_k$.
By Proposition~4\,(a),
\[
B_{k,y}(0)=\left(b_k+i\gamma_5\cdot y\right)^{-1}\cdot\left(1-P(0)\right)+\left(b_k+b_{L}+i\gamma_5\cdot y\right)^{-1}\cdot P(0),
\]
and for every real $\mathrm{cf}>0$ one has
\[
\left(\mathrm{cf}+i\gamma_5\cdot y\right)^{-1}=\frac{\mathrm{cf}-i\gamma_5\cdot y}{\mathrm{cf}^{2}+y^{2}},
\]
because $\gamma_5^{2}=1$.
By (5.3), $P(0)^{T}=P(0)$, and by (I.c3), $\mathsf{C}^{-1}\cdot\gamma_5\cdot\mathsf{C}=-\gamma_5^{T}$.
Therefore
\begin{multline*}
B_{k,y}(0)^{T}-\mathsf{C}^{-1}\cdot B_{k,y}(0)\cdot\mathsf{C}\\
=-2iy\cdot\gamma_5^{T}\cdot\left[\left(b_k^{2}+y^{2}\right)^{-1}\cdot\left(1-P(0)\right)+\left(\left(b_k+b_{L}\right)^{2}+y^{2}\right)^{-1}\cdot P(0)\right]\ne0.
\end{multline*}
The bracket is invertible: $P(0)^{2}=P(0)$, because $Q(0)\cdot Q^{T}(0)=1$ by the displays of Section~2.2, so that $1-P(0)$ and $P(0)$ are complementary idempotents, and their coefficients are positive.
Now let $\tilde{g}=(S,0)$ with $S\notin\mathrm{Sp}_3$.
Then $\tilde{g}0=0$, and (b) gives $\tilde{g}\cdot B_{k,y}(0)\cdot\tilde{g}^{-1}=B_{k,-y}(0)$, whereas
\[
B_{k,-y}(0)-B_{k,y}(0)=2iy\cdot\gamma_5\cdot\left[\text{the same bracket}\right]\ne0.
\]
\end{named}

\begin{named}{Remark 3}[(f) is also a consequence of (c) and (d)]
Let $A\in\tilde{G}_k$ be real, and let $O_{y}$ be one of the families of (c).
By (Ru1), by (c), by the reality of $\mathsf{C}$, and by (d) at the real connection $-A$ (note that $\overline{-A}=-A$),
\begin{align*}
O_{y}(A)^{*}
&=\overline{O_{y}(A)^{T}}
=\overline{\mathsf{C}^{-1}\cdot O_{-y}(-A)\cdot\mathsf{C}}
=\mathsf{C}^{-1}\cdot\overline{O_{-y}\left(\overline{-A}\right)}\cdot\mathsf{C}\\
&=\mathsf{C}^{-1}\cdot\varkappa^{-1}\cdot O_{-y}(A)\cdot\varkappa\cdot\mathsf{C}
=\left(\varkappa\cdot\mathsf{C}\right)^{-1}\cdot O_{-y}(A)\cdot\left(\varkappa\cdot\mathsf{C}\right).
\end{align*}
By Section~I.2.1, $\varkappa=\gamma_5\cdot\mathsf{C}$, and $\mathsf{C}^{2}=-1$ by (I.c3); hence $\varkappa\cdot\mathsf{C}=-\gamma_5$, and $\left(\varkappa\cdot\mathsf{C}\right)^{-1}\cdot O\cdot\left(\varkappa\cdot\mathsf{C}\right)=\gamma_5\cdot O\cdot\gamma_5$.
This is in accordance with Section~5.6.
\end{named}

\begin{named}{Remark 4}
By Section~5.3, $\bar{H}_k(A)=\mathfrak{H}_k(A)^{T}$.
The relation $Q(A)^{T}=Q^{T}(A)$ ($k\ge1$) is part of Proposition~2.2\,(c).
At $k=0$ its place is taken by $\left(Q(A)\cdot J^{*}\right)^{T}=J\cdot Q^{T}(A)$, which belongs to (g-c).
\end{named}

\begin{named}{Remark 5}[a variant of the definition of $\tilde{G}_k$]
Let the relation $A_0=A'+d\boldsymbol{\lambda}$ in the definition of $\tilde{G}_k$ (Section~I.5.3) be required on the bonds that meet $\square^{\sim 5}$ (Section~I.1.1) instead of the bonds contained in $\square^{\sim 5}$, with $A'$ a function on these bonds and $\boldsymbol{\lambda}$ a function on their end points.
For the set defined in this way Lemma~1 holds as well, with $\mathrm{lamc}$ a function on the end points of the bonds that meet $\square^{\sim 5}$ and $\mathrm{Acl}=d\,\mathrm{lamc}$ on these bonds.
In the proof of (i) the set $\mathrm{Wb}$ is replaced by
\[
\left\{c\in\mathbb{L}_k:|c-z_{\square}|<\frac{11M}{2}+\eta_k\right\},
\]
which is again denoted by $\mathrm{Wb}$.
Since $11ML^k$ is odd, this set is the product of the three intervals
\[
\left\{c_\mu\in\eta_k\mathbb{Z}:|c_\mu-z_{\square,\mu}|\le\frac{11M+\eta_k}{2}\right\},\qquad\mu=1,2,3,
\]
each of length $11M+\eta_k$, and
\[
\text{sup-diameter}(\mathrm{Wb})+\eta_k=11M+2\eta_k\le11M+2<13LM\le\ell.
\]
Hence $\mathrm{Wb}$ is a lift of its image in $T_k$, and the bonds of the image are the images of the bonds of $\mathrm{Wb}$ (Section~I.1.1).
The image contains $\square^{\sim 5}$ together with the six nearest neighbours of each site of $\square^{\sim 5}$; therefore it contains both end points of every bond that meets $\square^{\sim 5}$.
The gauge fixing lemma \cite[Lemma~16.26]{QEDI} is applied with $s_W:=11M+\eta_k$.
The other steps of the proof of (i), the proof of (ii), and the reasons given with (Dm) remain unchanged.
\end{named}

\subsection*{8. Conditions on the constants}

\emph{Conditions on $L$ and $M$.}
By the row (L5) of the table of Proposition~14.2\,(b) (line (C3) of Section~I.8.2), $L$ is odd.
By the row (M1) (line (C5)), $M$ is a power of $L$ with $M\ge L^{2}$; hence $M$ is a multiple of $L$, and $M$ is odd because $L$ is odd.
It follows that $11ML^k$ is odd, as required in Lemma~1 and in Remark~5 (Sections~I.0.5, I.1.1), and that $M\mathbb{Z}^3\subset L\mathbb{Z}^3$, as required in Section~2.3 (Section~I.0.5).
In Remark~1 the inequality $L\ge5$ of the row (L5) ensures that the site $(1,1,0)$ lies in $B(0)$.

\emph{The data and the ranges of the standing letters.}
The side satisfies $\ell\in\mathrm{Cl}$ (Section~I.1.1), which gives $13LM\le\ell$; hence $11M<13LM\le\ell$, as required in Lemma~1, and $11M+2<13LM\le\ell$, as required in Remark~5.
The condition $\ell\in\mathrm{Cl}$ also gives $\ell\ge3L$ and that $\ell$ is a multiple of $L$, so that $T_k$ is the block torus $(0,\ell,-1)$ (Section~G1.2).
The number $\alpha$ is real (Section~I.0.5; line (C1) of Section~I.8.2), so that $d_X(x,x')^{-\alpha}$ is real in the display of Section~5.4 for $\delta_{\alpha}$ and, through that display, in (g-d).
The numbers $\mu$, $e_k$, $\eta_k=L^{-k}$, $b_k$, $b_{L}$ are real, and $b_k>0$ ((SF) of Section~I.1.8; Sections~I.3.4, I.0.5).
The parameter $y$ is real or $\infty$ ($Y_k$: Sections~I.3.5, I.3.9).

\emph{Conditions of the propositions applied.}
Proposition~2.8\,(a), and Proposition~4 through Proposition~2.8\,(a), hold under the conditions stated with them at the head of the Propositions~2.m and with Proposition~2.8, namely, in the notation of the table of Proposition~14.2\,(b),
(L5) $L$ odd, $L\ge5$ (line (C3) of Section~I.8.2);
(M1) $M$ a power of $L$, $M\ge L^{2}$, (M2) $M>4L$, and $M\ge108\pi\cdot C'_{V}\cdot a'$, which is contained in (M5) (line (C5));
(Z2.s1) $6L\cdot e_0^{\varepsilon}\le1$ and (Z2.s2) $4C'_{V}\cdot a'\cdot e_0^{\varepsilon}\le1$ (line (C8));
and the condition $\alpha<1$ of Proposition~2.8, which holds by the range of the standing letter $\alpha$ (Section~I.0.5; line (C1) of Section~I.8.2).
By the table of Chapter~2a, Section~5, Proposition~2.2 as a whole requires that $L$ be odd, which is contained in (L5) (line (C3)), and that $b>0$ (line (C0)).
Its clauses (c), (d), (e) require in addition the closure relation $\sum_i w_i\hat n_i=0$ of the standard weights ((I.1.1); line (C0); (I.BARE)), by which the diagonal term of $D_A$ is the scalar $\frac{\rho}{2}\cdot\frac{\mathrm{sw}}{\eta_k}$.
Through Proposition~2.2\,(c) the closure relation enters the proofs of (c), (e), (g-c), (g-e); through Proposition~2.2\,(d) those of (d), (e), (g-d), (g-e); through Proposition~2.2\,(e) those of (f), (g-f).
The invariance property (G2) of the domain in \cite[Section~16.2.7]{QEDI} and the gauge fixing lemma \cite[Lemma~16.26]{QEDI} require no condition.

\section*{Chapter 11. Exact operators on the tori of the sides $\ell$ and $\mathrm{Jr}\,\ell$ (one formula; coverings of tori): proof of Proposition~11}

\subsection*{Proof}

\subsubsection*{0. Conventions of this chapter; blocks, paths, phases on the covering lattice}

\paragraph{0.0. Letters and labels.}
In this chapter the following notation is used.
For an operator $O$ of the lists of Proposition~11, $\mathrm{Op}_{\mathbb{L}}[O]$ is the operator $\mathrm{Op}_{\mathbb{L}}$ of Section~I.1.9 that belongs to $O$.
For a bond function $A$ on a torus, $\tilde A$ is its periodic extension, that is, the function written $\tilde B$ in Section~I.1.9 (see 6.0).
The covering maps $\mathbb{L}_k\to T_i$ (Section~6) are denoted $\mathrm{pr}_i$, and $\boldsymbol{\lambda}$ denotes a gauge function.
Unit sites are denoted $z$, $z_2$, $\tilde z$, and sites of $L\mathbb{Z}^3$ are denoted $z'$.
The centres of the blocks that contain a photon site $c$, which are the points $y_j$ of Section~I.3.2, are denoted $z_j(c)$; in 0.2 the letters $z$, $z'$ also stand for centres of blocks $\mathrm{Bl}_j$ of arbitrary side.
The letter $y$ always denotes the resolvent parameter.
The labels (f1)--(f4) refer to the four clauses of the Definition in Section~I.1.9.

\paragraph{0.1. Index sets, spaces, covering data.}
The index sets are those of (I.i1), (I.i4) and (I.i5) in Section~I.1.9, namely
\begin{gather*}
\mathbb{L}^{\mathrm{f}}_k = \mathbb{L}_k\times O_{\mathrm{cub}}\ \text{with}\ \operatorname{supp}x = \{c\};\qquad
\mathbb{Z}^3\ \text{with}\ \operatorname{supp}z = \{z\};\\
L\mathbb{Z}^3\ \text{with}\ \operatorname{supp}z' = \{z'\};\qquad
\mathrm{Pzt}_k\ \text{with}\ \operatorname{supp}(x,x') = \{c,c'\},
\end{gather*}
where $x = (c,g)$ and $x' = (c',g')$.
Both $\mathbb{Z}^3$ and $L\mathbb{Z}^3$ are subsets of $\mathbb{L}_k = \eta_k\mathbb{Z}^3$.
A vector $v\in\eta_k\mathbb{Z}^3$ acts by $(c,g) + v = (c + v,g)$ and $(x,x') + v = (x + v,x' + v)$; it maps $\mathbb{Z}^3$, resp.\ $L\mathbb{Z}^3$, into itself if $v\in\mathbb{Z}^3$, resp.\ $v\in L\mathbb{Z}^3$.
For an index set $\mathrm{Ind}$ let $\mathrm{Fu}(\mathrm{Ind})$ be the space of all functions $\mathrm{Ind}\to\mathbb{C}^4$, with no periodicity and no condition of growth, and put
\[
\left(T_v f\right)(t) := f(t - v),\qquad
\mathrm{Fu}_\ell(\mathrm{Ind}) := \left\{f\in\mathrm{Fu}(\mathrm{Ind}) : f(t + \ell e_\mu) = -f(t),\ \mu = 1,2,3\right\}
\]
for a side $\ell$.
By Sections~I.1.5 and I.1.7,
\[
\mathrm{Fu}_\ell\left(\mathbb{L}^{\mathrm{f}}_k\right) = \mathfrak{S}^{\mathrm{ch}}(\eta_k,\ell),\qquad
\mathrm{Fu}_\ell\left(\mathbb{Z}^3\right) = \mathfrak{S}(1,\ell),\qquad
\mathrm{Fu}_\ell\left(L\mathbb{Z}^3\right) = \mathfrak{S}(L,\ell),
\]
and $\mathrm{Fu}_\ell(\mathrm{Pzt}_k)$ is the space of the functions of pairs on the torus of side $\ell$.
A constant $4\times 4$ matrix $\mathrm{Mx}$ acts on every $\mathrm{Fu}(\mathrm{Ind})$ by $(\mathrm{Mx}\,f)(t) := \mathrm{Mx}\cdot f(t)$, a number acting as the corresponding multiple of the unit matrix.
These operators commute with every $T_v$ and leave every $\mathrm{Fu}_\ell(\mathrm{Ind})$ invariant.

The letter $\tilde B$ denotes a complex function on the oriented bonds of $\mathbb{L}_k$ with $\tilde B(c',c) = -\tilde B(c,c')$; no periodicity is required, and $(T_v\tilde B)(c,c') := \tilde B(c - v,c' - v)$.
The letter $\tilde Y$ denotes a union of $M$-cubes of $\mathbb{L}_k$; the empty union and $\mathbb{L}_k$ itself are admitted.
A bond function $B$ on $T_k = \mathbb{L}_k/\ell\mathbb{Z}^3$ is identified with the $\ell\mathbb{Z}^3$-periodic function $\tilde B$ on the bonds of $\mathbb{L}_k$, and a union $Y$ of $M$-cubes of $T_k$ with its preimage $\tilde Y$ in $\mathbb{L}_k$ (Section~I.6.3).
These are the $\ell\mathbb{Z}^3$-periodic extensions that enter clause (f3) of Section~I.1.9.

For a finite set $Z$ of sites of $\mathbb{L}_k$ put, in accordance with (I.1.3),
\[
\int_Z\mathrm{d}x\,f(x) := \eta_k^3\,\frac{1}{48}\sum_{c\in Z}\sum_{g\in O_{\mathrm{cub}}}f(c,g).
\]
For a site $\tilde c$ and finite sets $Z_1$, $Z_2$ of sites let
\[
\operatorname{dist}(\tilde c,Z_1) := \min\left\{|\tilde c - c| : c\in Z_1\right\},\qquad
\operatorname{dist}(Z_1,Z_2) := \min\left\{|c - c'| : c\in Z_1,\ c'\in Z_2\right\},
\]
where $|\cdot|$ is the sup norm.

\paragraph{0.2. Blocks.}
The integer $L$ is odd by (L5), and $M = L^{\mathrm{iM}}$ with an integer $\mathrm{iM}\ge 2$ by (M1).
For an integer $j\ge 0$ and a point $z\in L^j\eta_k\mathbb{Z}^3$ let
\[
\mathrm{Bl}_j(z) := \left\{c\in\mathbb{L}_k : |c - z| < \frac{L^j\eta_k}{2}\right\};
\]
this is the definition of Section~I.1.1, where it is made for $j\le k + 1$.
Thus $\Delta_z = \mathrm{Bl}_k(z)$ for $z\in\mathbb{Z}^3 = L^k\eta_k\mathbb{Z}^3$, and the $M$-cube with centre $z\in M\mathbb{Z}^3 = L^{k+\mathrm{iM}}\eta_k\mathbb{Z}^3$ is $\mathrm{Bl}_{k+\mathrm{iM}}(z)$.
\begin{itemize}
\item[(Bq1)]
One has
\[
\mathrm{Bl}_j(z) = \left\{z + \eta_k n : n\in\mathbb{Z}^3,\ |n|\le\frac{L^j - 1}{2}\right\}.
\]
Indeed, $c - z = \eta_k n$ with $n\in\mathbb{Z}^3$, and, $L^j$ being odd, the inequality $|n| < L^j/2$ is equivalent to $|n|\le (L^j - 1)/2$.
Consequently $\mathrm{Bl}_j(z)$ is a product of three intervals of $\eta_k\mathbb{Z}$ and has $L^{3j}$ points; each of its points has distance at most $(L^j - 1)\eta_k/2$ from $z$, and its sup-diameter equals $(L^j - 1)\eta_k$.
\item[(Bq2)]
For fixed $j$ the blocks $\mathrm{Bl}_j(z)$, $z\in L^j\eta_k\mathbb{Z}^3$, form a partition of $\mathbb{L}_k$.
Indeed, since $L^j$ is odd, every integer $n$ has exactly one representation $n = L^j n' + n''$ with integers $n'$, $n''$ and $|n''|\le (L^j - 1)/2$; one applies this to each of the three coordinates of $c/\eta_k$.
\item[(Bq3)]
(Nesting.)
If $\mathrm{Bl}_j(z)$ and $\mathrm{Bl}_{j+1}(z')$ have a common point $c$, then $\mathrm{Bl}_j(z)\subset\mathrm{Bl}_{j+1}(z')$ and $|z - z'|\le (L^{j+1} - L^j)\eta_k/2$.
To see this, note that $z'$ lies in $L^{j+1}\eta_k\mathbb{Z}^3$, which is a subset of $L^j\eta_k\mathbb{Z}^3$, so that $z - z' = L^j\eta_k n'$ with $n'\in\mathbb{Z}^3$.
By (Bq1),
\[
|z - z'|\le|z - c| + |c - z'|\le\left((L^j - 1) + (L^{j+1} - 1)\right)\frac{\eta_k}{2} < L^j\eta_k\,\frac{L + 1}{2}.
\]
Hence $|n'| < (L + 1)/2$, that is, $|n'|\le (L - 1)/2$, and $|z - z'| = L^j\eta_k|n'|\le (L^{j+1} - L^j)\eta_k/2$.
For $\tilde c\in\mathrm{Bl}_j(z)$ it follows that
\[
|\tilde c - z'|\le\frac{(L^j - 1)\eta_k}{2} + L^j\eta_k\,\frac{L - 1}{2} = \frac{(L^{j+1} - 1)\eta_k}{2},
\]
so that $\tilde c\in\mathrm{Bl}_{j+1}(z')$.
Induction on $j' - j$ gives: two blocks $\mathrm{Bl}_j(z)$, $\mathrm{Bl}_{j'}(z')$ with $j\le j'$ are either disjoint, or the first is contained in the second.
\item[(Bq4)]
(Consequences.)
By (Bq2), the point $[c]$, that is, the point $z$ of $\mathbb{Z}^3$ with $c\in\Delta_z$, is well defined.
One has $\#|\Delta_z| = \eta_k^{-3}$, hence $\int_{\Delta_z}\mathrm{d}x = \eta_k^3\cdot\frac{1}{48}\cdot 48\,\eta_k^{-3} = 1$; the points of $\Delta_z$ have distance at most $(1 - \eta_k)/2$ from $z$, the sup-diameter of $\Delta_z$ is $1 - \eta_k$, and $\Delta_z\cap\mathbb{Z}^3 = \{z\}$.
The set $B(z') = \{z\in\mathbb{Z}^3 : |z - z'|\le (L - 1)/2\}$ has $L^3$ points, $L$ being odd, and every $z\in\mathbb{Z}^3$ lies in exactly one $B(z')$, $z'\in L\mathbb{Z}^3$.
The block $\mathrm{Bl}_{k+1}(z')$ is the union of the $\Delta_z$, $z\in B(z')$: by (Bq2) and (Bq3) it is the union of those $\Delta_z$ that meet it, and these are the $\Delta_z$ with $z\in\mathrm{Bl}_{k+1}(z')\cap\mathbb{Z}^3 = B(z')$.
The sup-diameter of $\mathrm{Bl}_{k+1}(z')$ is $L - \eta_k$.
If a point $z$ of $\mathbb{Z}^3$ lies in $\tilde Y$, then $\Delta_z\subset\tilde Y$, because by (Bq3) the block $\Delta_z$ lies in the $M$-cube that contains $z$.
(That every unit block $\mathrm{Bl}_k(z)$ and every $L$-block $\mathrm{Bl}_{k+1}(z')$ lies in a single $M$-cube is part of (g0) in \cite[Section~17.4.1]{QEDI}; here it follows from (Bq2) and (Bq3), because $k + 1\le k + \mathrm{iM}$.)
\end{itemize}

\paragraph{0.3. Paths and phases.}
\begin{itemize}
\item[(Bq5)]
For $c,c'\in\mathbb{L}_k$ put
\[
\mathrm{bx}(c,c') := \left\{\tilde c\in\mathbb{L}_k : \min(c_\mu,c'_\mu)\le\tilde c_\mu\le\max(c_\mu,c'_\mu),\ \mu = 1,2,3\right\}.
\]
All sites of the six paths $\boldsymbol{\Gamma}^{\pi}(c,c')$ of Section~I.1.1 lie in $\mathrm{bx}(c,c')$, since along such a path each coordinate either equals $c_\mu$, or equals $c'_\mu$, or moves from $c_\mu$ to $c'_\mu$.
For $\tilde c\in\mathrm{bx}(c,c')$ one has $|\tilde c - c|\le|c' - c|$ and $|\tilde c - c'|\le|c' - c|$.
A product of three intervals of $\eta_k\mathbb{Z}$ that contains two points $c$, $c'$ contains $\mathrm{bx}(c,c')$.
\end{itemize}
The displays of Section~I.1.1, (I.1.6) and (I.3.3) are read on $\mathbb{L}_k$, with $\tilde B$ as the bond function.
For a path $\boldsymbol{\Gamma}$, $\tilde B(\boldsymbol{\Gamma})$ is the sum of $\tilde B(b)$ over the bonds $b$ of $\boldsymbol{\Gamma}$, each oriented along the path, and
\begin{gather*}
(\tau\tilde B)(c,c') := \frac{1}{6}\sum_{\pi}\tilde B\left(\boldsymbol{\Gamma}^{\pi}(c,c')\right),\qquad
(\tau\tilde B)(c,c) := 0;\\
u_{\tilde B}(c,c') := \exp\left(i\,e_k\eta_k\cdot\tilde B(c,c')\right)\quad\text{(bonds)};\\
U_{\tilde B}(c,c') := \exp\left(i\,e_k\eta_k\cdot(\tau\tilde B)(c,c')\right),\qquad
U_{\tilde B}(c,c) = 1;\qquad
U_{\tilde B}(x,x') := U_{\tilde B}(c,c').
\end{gather*}
For nearest neighbours the six paths reduce to the bond itself, so that $U_{\tilde B}(c,c') = u_{\tilde B}(c,c')$.
\begin{itemize}
\item[(Bq6)]
(Towers; $k\ge 1$.)
For $c\in\mathbb{L}_k$ and $0\le j\le k$ let $z_j(c)$ be the centre of the block $\mathrm{Bl}_j$ that contains $c$; thus $z_0(c) = c$ and $z_k(c) = [c]$.
By (Bq3) the blocks $\mathrm{Bl}_j(z_j(c))$ increase with $j$.
Therefore all $z_j(c)$ lie in $\Delta_{[c]}$, and by (Bq1) and (Bq5) all sites of the paths $\boldsymbol{\Gamma}^{\pi}(z_{j+1}(c),z_j(c))$, $j < k$, lie in $\Delta_{[c]}$.
(By (Bq3), $|z_{j+1}(c) - z_j(c)|\le (L^{j+1} - L^j)\eta_k/2 < 1/2$; these are the points and the paths of Section~I.3.2.)
For $c\in\Delta_z$ put
\[
\varphi_{\tilde B}(z,c) := \exp\left(i\,e_k\eta_k\cdot\sum_{j<k}(\tau\tilde B)\left(z_{j+1}(c),z_j(c)\right)\right).
\]
\item[(Bq7)]
For $z'\in L\mathbb{Z}^3$ put
\[
\mathrm{Bc}(z') := \left\{\tilde c\in\mathbb{L}_k : |\tilde c - z'|\le\frac{L - 1}{2}\right\}.
\]
This set is a product of three intervals of $\eta_k\mathbb{Z}$; it is contained in $\mathrm{Bl}_{k+1}(z')$, its sup-diameter is $L - 1$, and $\mathrm{Bc}(z')\cap\mathbb{Z}^3 = B(z')$, $\mathrm{Bc}(z')\cap L\mathbb{Z}^3 = \{z'\}$.
For $z\in B(z')$ all sites of the paths $\boldsymbol{\Gamma}^{\pi}(z',z)$ lie in $\mathrm{bx}(z',z)$, which is a subset of $\mathrm{Bc}(z')$.
\item[(Bq8)]
(What the phases read.)
All phases are exponentials and are therefore different from $0$.
The phase $u_{\tilde B}(c,c')$ is a function of the number $\tilde B(c,c')$.
The phase $U_{\tilde B}(c,c')$ is a function of the numbers $\tilde B(b)$, $b$ a bond with both ends in $\mathrm{bx}(c,c')$.
The phase $\varphi_{\tilde B}(z,c)$ is a function of the numbers $\tilde B(b)$, $b$ a bond with both ends in $\Delta_z$.
For $z\in B(z')$, the phase $U_{\tilde B}(z',z)$ is a function of the numbers $\tilde B(b)$, $b$ a bond with both ends in $\mathrm{bx}(z',z)$, a subset of $\mathrm{Bc}(z')$.
All this follows from (Bq5), (Bq6) and (Bq7).
\item[(Bq9)]
(Translations.)
Let $v\in LM\mathbb{Z}^3$; then $v\in L^j\eta_k\mathbb{Z}^3$ for $0\le j\le k + \mathrm{iM}$.
Consequently
\begin{gather*}
\mathrm{Bl}_j(z) + v = \mathrm{Bl}_j(z + v);\\
[c + v] = [c] + v,\qquad
z_j(c + v) = z_j(c) + v,\qquad
\Delta_z + v = \Delta_{z+v},\qquad
B(z') + v = B(z' + v),
\end{gather*}
and $\tilde Y + v$ is a union of $M$-cubes.
For every $v\in\eta_k\mathbb{Z}^3$ one has $\boldsymbol{\Gamma}^{\pi}(c + v,c' + v) = \boldsymbol{\Gamma}^{\pi}(c,c') + v$; further $(x + v)^{(i)} = x^{(i)} + v$ for $i = 0,\dots,3$ and $c_3(x + v) = c_3(x) + v$, by (I.1.2) and (I.s1); $M_i(x + v) = M_i(x)$, since by (I.2.1) the matrix $M_i(x)$ depends on the label only; and $d_X(x + v,x' + v) = d_X(x,x')$ and $\mathrm{Pzt}_k + v = \mathrm{Pzt}_k$, by (I.1.5) and Section~I.1.7.
Hence
\begin{gather*}
u_{T_v\tilde B}(c,c') = u_{\tilde B}(c - v,c' - v),\qquad
U_{T_v\tilde B}(c,c') = U_{\tilde B}(c - v,c' - v)\qquad\left(v\in\eta_k\mathbb{Z}^3\right);\\
\varphi_{T_v\tilde B}(z,c) = \varphi_{\tilde B}(z - v,c - v)\qquad\left(v\in LM\mathbb{Z}^3\right).
\end{gather*}
\end{itemize}

\subsubsection*{1. $k\ge 1$: the operators on the covering lattice}

For every $\tilde B$ and every $\tilde Y$ the displays below define linear operators between the spaces $\mathrm{Fu}(\cdot)$.
They are the displays (I.2.3), those of Sections~I.3.3 and I.3.5, (I.1.7) and (I.5.4), and, in (cv8), the formula for $B_{k,y}(A)$ named in clause (a) of Proposition~11, all read on the covering lattice.
Here $\mathrm{Op}_{\mathbb{L}}[O]$ is the operator that belongs to the operator $O$ of the list of (a); note that $\mathrm{Op}_{\mathbb{L}}[J Q_k^T](\tilde B)$ belongs to $J Q_k^T(-A)$, and $\mathrm{Op}_{\mathbb{L}}[Q^T](\tilde B)$ to $Q^T(-A)$.
Throughout,
\begin{gather*}
x = (c,g);\qquad
x^{(i)} = (c,g s_i)\ \ (i = 0,1,2);\qquad
x^{(3)} = (c_3,g s_3)\ \ \text{with}\ \ c_3 = c_3(x) = c + \eta_k\,g e_1;\\
M_i(x) = \frac{w_i}{2}\left(\rho + \gamma\cdot(g\hat n_i)\right),\quad\text{in particular}\quad M_3(x) = 3\left(\rho - \gamma\cdot(g e_1)\right)
\end{gather*}
(see (I.1.2) and (I.2.1)), and $f\in\mathrm{Fu}(\mathbb{L}^{\mathrm{f}}_k)$, $\mathrm{fu}\in\mathrm{Fu}(\mathbb{Z}^3)$, $\mathrm{fL}\in\mathrm{Fu}(L\mathbb{Z}^3)$, $F\in\mathrm{Fu}(\mathrm{Pzt}_k)$.
\begin{equation}
\begin{split}
\left(\mathrm{Op}_{\mathbb{L}}[D](\tilde B)\,f\right)(x) := -\eta_k^{-1}&\left\{\sum_{i=0}^{2}M_i(x)\cdot\left[f(x^{(i)}) - f(x)\right]\right.\\
&\left.{} + M_3(x)\cdot\left[u_{\tilde B}(c,c_3)\cdot f(x^{(3)}) - f(x)\right]\right\}\qquad\left(x\in\mathbb{L}^{\mathrm{f}}_k\right);
\end{split}
\tag{cv1}
\end{equation}
\begin{equation}
\begin{split}
\left(\mathrm{Op}_{\mathbb{L}}[Q_k J^*](\tilde B)\,f\right)(z) &:= \eta_k^3\sum_{c\in\Delta_z}\varphi_{\tilde B}(z,c)\cdot\frac{1}{48}\sum_{g\in O_{\mathrm{cub}}}f(c,g)\\
&\phantom{:}= \int_{\Delta_z}\mathrm{d}x\,\varphi_{\tilde B}(z,c)\cdot f(x)\qquad\left(z\in\mathbb{Z}^3\right);
\end{split}
\tag{cv2}
\end{equation}
\begin{equation}
\left(\mathrm{Op}_{\mathbb{L}}[J Q_k^T](\tilde B)\,\mathrm{fu}\right)(c,g) := \varphi_{\tilde B}([c],c)^{-1}\cdot\mathrm{fu}([c]),
\tag{cv3}
\end{equation}
which is the same value at each of the $48$ sites over $c$;
\begin{equation}
\left(\mathrm{Op}_{\mathbb{L}}[Q](\tilde B)\,\mathrm{fu}\right)(z') := L^{-3}\sum_{z\in B(z')}U_{\tilde B}(z',z)\cdot\mathrm{fu}(z)\qquad\left(z'\in L\mathbb{Z}^3\right);
\tag{cv4}
\end{equation}
\begin{equation}
\left(\mathrm{Op}_{\mathbb{L}}[Q^T](\tilde B)\,\mathrm{fL}\right)(z) := U_{\tilde B}(z',z)^{-1}\cdot\mathrm{fL}(z')\qquad\left(z\in\mathbb{Z}^3\right),
\tag{cv5}
\end{equation}
where $z'$ is the point of $L\mathbb{Z}^3$ with $z\in B(z')$;
\begin{equation}
\begin{split}
&\mathrm{Op}_{\mathbb{L}}[J P_k J^*] := \mathrm{Op}_{\mathbb{L}}[J Q_k^T]\cdot\mathrm{Op}_{\mathbb{L}}[Q_k J^*],\qquad
\mathrm{Op}_{\mathbb{L}}[P] := \mathrm{Op}_{\mathbb{L}}[Q^T]\cdot\mathrm{Op}_{\mathbb{L}}[Q],\\
&\mathrm{Op}_{\mathbb{L}}[J P_{k+1} J^*] := \mathrm{Op}_{\mathbb{L}}[J Q_k^T]\cdot\mathrm{Op}_{\mathbb{L}}[P]\cdot\mathrm{Op}_{\mathbb{L}}[Q_k J^*],
\end{split}
\tag{cv6}
\end{equation}
all factors being taken at the same $\tilde B$;
\begin{equation}
\begin{split}
&\mathrm{Op}_{\mathbb{L}}[\operatorname{Pen}_y] := \boldsymbol{\alpha}_y\cdot\mathrm{Op}_{\mathbb{L}}[J P_k J^*] + \boldsymbol{\beta}_y\cdot\mathrm{Op}_{\mathbb{L}}[J P_{k+1} J^*],\\
&\mathrm{Op}_{\mathbb{L}}[F_y] := \mathrm{Op}_{\mathbb{L}}[D] + \mu + \mathrm{Op}_{\mathbb{L}}[\operatorname{Pen}_y]\qquad\left(y\in Y_k = \mathbb{R}\cup\{\infty\}\right);
\end{split}
\tag{cv7}
\end{equation}
\begin{equation}
\begin{split}
\mathrm{Op}_{\mathbb{L}}[B_{k,y}] &:= \left(b_k + i\gamma_5 y\right)^{-1}\cdot\left(1 - \mathrm{Op}_{\mathbb{L}}[P]\right)\\
&\phantom{:=}{} + \left(b_k + b_L + i\gamma_5 y\right)^{-1}\cdot\mathrm{Op}_{\mathbb{L}}[P]\qquad(y\in\mathbb{R});
\end{split}
\tag{cv8}
\end{equation}
\begin{equation}
\mathrm{Op}_{\mathbb{L}}[T_k] := b_k^{-1}\,\mathrm{Op}_{\mathbb{L}}[Q_k J^*]\cdot\left(\mathrm{Op}_{\mathbb{L}}[D] + \mu\right) + \mathrm{Op}_{\mathbb{L}}[Q_k J^*];
\tag{cv9}
\end{equation}
\begin{equation}
\begin{split}
&\left(\mathrm{Op}_{\mathbb{L}}[T^{(1)}](\tilde B,\tilde Y)\,f\right)(z)\\
&\quad := \mathbf{1}[z\in\tilde Y]\,\eta_k^3\sum_{c\in\Delta_z}\varphi_{\tilde B}(z,c)\cdot\frac{1}{48}\sum_{g\in O_{\mathrm{cub}}}\left\{\left(1 + \frac{\mu}{b_k}\right)f(x)\right.\\
&\qquad\left.{} - (b_k\eta_k)^{-1}\left(\sum_{i=0}^{2}M_i(x)\cdot\left[f(x^{(i)}) - f(x)\right]\right.\right.\\
&\qquad\left.\left.{} + \mathbf{1}[c_3\in\tilde Y]\,M_3(x)\cdot\left[u_{\tilde B}(c,c_3)\cdot f(x^{(3)}) - f(x)\right]\right)\right\};
\end{split}
\tag{cv10}
\end{equation}
\begin{equation}
\begin{split}
&\left(\mathrm{Op}_{\mathbb{L}}[T^{(2)}](\tilde B,\tilde Y)\,F\right)(z)\\
&\quad := -\mathbf{1}[z\in\tilde Y]\,b_k^{-1}\eta_k^{\alpha-1}\,\eta_k^3\sum_{c\in\Delta_z}\varphi_{\tilde B}(z,c)\\
&\qquad\qquad{}\cdot\frac{1}{48}\sum_{g\in O_{\mathrm{cub}}}\mathbf{1}[c_3\notin\tilde Y]\,M_3(x)\cdot F(x,x^{(3)});
\end{split}
\tag{cv11}
\end{equation}
\begin{equation}
\begin{split}
&\mathrm{Op}_{\mathbb{L}}[\bar T_k](\tilde B) := \mathsf{C}^{-1}\cdot\mathrm{Op}_{\mathbb{L}}[T_k](-\tilde B)\cdot\mathsf{C},\\
&\mathrm{Op}_{\mathbb{L}}[\bar T^{(i)}](\tilde B,\tilde Y) := \mathsf{C}^{-1}\cdot\mathrm{Op}_{\mathbb{L}}[T^{(i)}](-\tilde B,\tilde Y)\cdot\mathsf{C}\qquad(i = 1,2);
\end{split}
\tag{cv12}
\end{equation}
\begin{equation}
\left(\mathrm{Op}_{\mathbb{L}}[\delta_\alpha](\tilde B)\,f\right)(x,x') := d_X(x,x')^{-\alpha}\left(U_{\tilde B}(c,c')\cdot f(x') - f(x)\right)
\tag{cv13}
\end{equation}
for $(x,x')\in\mathrm{Pzt}_k$, where $d_X(x,x') = \max\left(|c - c'|,\eta_k\right)$.

Some remarks on these displays.
In (cv7) and (cv8) the $4\times 4$ matrices are those of Section~I.3.5:
\begin{gather*}
\boldsymbol{\alpha}_y = b_k\,i\gamma_5 y\left(b_k + i\gamma_5 y\right)^{-1},\qquad
\boldsymbol{\beta}_y = b_k^2 b_L\left(b_k + b_L + i\gamma_5 y\right)^{-1}\left(b_k + i\gamma_5 y\right)^{-1}\quad(y\in\mathbb{R}),\\
\boldsymbol{\alpha}_\infty = b_k,\qquad\boldsymbol{\beta}_\infty = 0;
\end{gather*}
the inverses exist by (al4).
In (cv11), $(x,x^{(3)})$ is a pair of $\mathrm{Pzt}_k$, because $x^{(3)}\neq x$ and $d_X(x,x^{(3)}) = \eta_k\le M$.
In (cv10) the three differences inside the cluster ($i = 0,1,2$) carry no indicator.
For $z\in\tilde Y$ the sites $x^{(i)}$, $i\le 2$, lie over $c\in\Delta_z$, and $\Delta_z\subset\tilde Y$ by (Bq4); hence, for $z\in\tilde Y$, the sum in (cv10) is the sum over $\{i : x^{(i)}\ \text{over}\ \tilde Y\}$ of (I.5.4).
In (cv13), for a pair inside a cluster ($c = c'$, $g\neq g'$; there are $48\cdot 47$ such ordered pairs over each photon site) the display reduces to $\eta_k^{-\alpha}\left(f(c,g') - f(c,g)\right)$, which involves neither a phase nor a bond.

The following algebraic relations hold on the covering lattice, all operators being taken at one and the same $\tilde B$.
\begin{itemize}
\item[(al1)]
$\mathrm{Op}_{\mathbb{L}}[Q_k J^*]\cdot\mathrm{Op}_{\mathbb{L}}[J Q_k^T] = 1$ on $\mathrm{Fu}(\mathbb{Z}^3)$.
Indeed, applied to $\mathrm{fu}$, the left side takes at $z$ the value
\[
\int_{\Delta_z}\mathrm{d}x\,\varphi_{\tilde B}(z,c)\cdot\varphi_{\tilde B}(z,c)^{-1}\cdot\mathrm{fu}(z) = \mathrm{fu}(z)\int_{\Delta_z}\mathrm{d}x = \mathrm{fu}(z),
\]
by (Bq4).
It follows that
\[
\mathrm{Op}_{\mathbb{L}}[T_k] = b_k^{-1}\,\mathrm{Op}_{\mathbb{L}}[Q_k J^*]\cdot\left(\mathrm{Op}_{\mathbb{L}}[D] + \mu + b_k\,\mathrm{Op}_{\mathbb{L}}[J P_k J^*]\right),
\]
which is the first form of $T_k$ in Section~I.3.5.
\item[(al2)]
$\mathrm{Op}_{\mathbb{L}}[Q]\cdot\mathrm{Op}_{\mathbb{L}}[Q^T] = 1$ on $\mathrm{Fu}(L\mathbb{Z}^3)$, since
\[
L^{-3}\sum_{z\in B(z')}U_{\tilde B}(z',z)\cdot U_{\tilde B}(z',z)^{-1}\cdot\mathrm{fL}(z') = \mathrm{fL}(z'),
\]
the set $B(z')$ having $L^3$ points.
\item[(al3)]
By (al2), $\mathrm{Op}_{\mathbb{L}}[P]^2 = \mathrm{Op}_{\mathbb{L}}[P]$.
The coefficients of $\mathrm{Op}_{\mathbb{L}}[P]$ are numbers: for $z\in B(z')$,
\[
\left(\mathrm{Op}_{\mathbb{L}}[P]\,\mathrm{fu}\right)(z) = L^{-3}\sum_{z_2\in B(z')}U_{\tilde B}(z',z)^{-1}\cdot U_{\tilde B}(z',z_2)\cdot\mathrm{fu}(z_2).
\]
Consequently $\mathrm{Op}_{\mathbb{L}}[P]$ commutes with every constant $4\times 4$ matrix.
\item[(al4)]
Let $y$ be real, and let $\tilde b$ be one of the two numbers $b_k$, $b_k + b_L$.
Both numbers are positive: one has $b > 0$ and $L > 1$ (Section~I.0.5: $b = 1$, $L\ge 5$), so that $b_L = b/L > 0$ (Section~I.0.5) and $b_k\ge b_- = b(1 - L^{-1}) > 0$ (Section~I.3.4).
Since $\gamma_5^2 = 1$ by (I.c1),
\[
\left(\tilde b + i\gamma_5 y\right)\left(\tilde b - i\gamma_5 y\right) = \tilde b^2 + y^2 > 0.
\]
Hence $\tilde b + i\gamma_5 y$ is invertible.
\item[(al5)]
(The formula of (a) for $B_{k,y}$, on the covering lattice.)
For real $y$,
\begin{align*}
\left(b_k + b_L\,\mathrm{Op}_{\mathbb{L}}[P] + i\gamma_5 y\right)\cdot\mathrm{Op}_{\mathbb{L}}[B_{k,y}]
&= \mathrm{Op}_{\mathbb{L}}[B_{k,y}]\cdot\left(b_k + b_L\,\mathrm{Op}_{\mathbb{L}}[P] + i\gamma_5 y\right)\\
&= 1\quad\text{on}\ \mathrm{Fu}(\mathbb{Z}^3).
\end{align*}
To see this, write $\mathrm{Pl} := \mathrm{Op}_{\mathbb{L}}[P]$; then
\[
b_k + b_L\,\mathrm{Pl} + i\gamma_5 y = \left(b_k + i\gamma_5 y\right)\left(1 - \mathrm{Pl}\right) + \left(b_k + b_L + i\gamma_5 y\right)\mathrm{Pl}.
\]
By (al3), $\mathrm{Pl}$ and $1 - \mathrm{Pl}$ are idempotent, $\mathrm{Pl}\cdot(1 - \mathrm{Pl}) = 0$, and both commute with the functions of $\gamma_5$.
Hence each of the two products equals $(1 - \mathrm{Pl}) + \mathrm{Pl} = 1$.
\end{itemize}

\subsubsection*{2. $k\ge 1$: (f1), (f2), (f3)}

\paragraph{2.1. (f1).}
Each of the operators (cv1)--(cv5), (cv10), (cv11), (cv13) is given by a finite sum.
The row of (cv1) has the columns $x$, $x^{(0)}$, $x^{(1)}$, $x^{(2)}$, $x^{(3)}$; the row of (cv2) has the $48\,\eta_k^{-3}$ columns over $\Delta_z$; the rows of (cv3) and (cv5) have one column each; the row of (cv4) has $L^3$ columns; the rows of (cv10) and (cv11) have at most $2\cdot 48\,\eta_k^{-3}$ columns; and the row of (cv13) has the two columns $x$, $x'$.
Sums and products of operators whose rows have finite support, as well as their products with constant matrices, again have rows of finite support.

\paragraph{2.2. (f2).}
Let $v\in LM\mathbb{Z}^3$.
The relation
\[
T_v\cdot\mathrm{Op}_{\mathbb{L}}(\tilde B,\tilde Y) = \mathrm{Op}_{\mathbb{L}}(T_v\tilde B,\tilde Y + v)\cdot T_v
\]
is first verified for the operators that are defined by a display; each line below uses (Bq9).
In the integrals one substitutes $x = \tilde x - v$, that is, $c = \tilde c - v$ with the label unchanged; note that $\Delta_{z-v} = \Delta_z - v$.

For (cv1):
\begin{align*}
\left(T_v\,\mathrm{Op}_{\mathbb{L}}[D](\tilde B)\,f\right)(x)
&= \left(\mathrm{Op}_{\mathbb{L}}[D](\tilde B)\,f\right)(x - v)\\
&= -\eta_k^{-1}\left\{\sum_{i\le 2}M_i(x - v)\cdot\left[f\left((x - v)^{(i)}\right) - f(x - v)\right]\right.\\
&\qquad\left.{} + M_3(x - v)\cdot\left[u_{\tilde B}(c - v,c_3 - v)\cdot f\left((x - v)^{(3)}\right) - f(x - v)\right]\right\}\\
&= -\eta_k^{-1}\left\{\sum_{i\le 2}M_i(x)\cdot\left[(T_v f)(x^{(i)}) - (T_v f)(x)\right]\right.\\
&\qquad\left.{} + M_3(x)\cdot\left[u_{T_v\tilde B}(c,c_3)\cdot(T_v f)(x^{(3)}) - (T_v f)(x)\right]\right\}\\
&= \left(\mathrm{Op}_{\mathbb{L}}[D](T_v\tilde B)\,T_v f\right)(x).
\end{align*}
This computation is valid for every $v\in\eta_k\mathbb{Z}^3$; it is the translation covariance stated in (I.O7).

For (cv2):
\begin{align*}
\left(T_v\,\mathrm{Op}_{\mathbb{L}}[Q_k J^*](\tilde B)\,f\right)(z)
&= \int_{\Delta_{z-v}}\mathrm{d}x\,\varphi_{\tilde B}(z - v,c)\cdot f(x)\\
&= \int_{\Delta_z}\mathrm{d}\tilde x\,\varphi_{\tilde B}(z - v,\tilde c - v)\cdot f(\tilde x - v)\\
&= \int_{\Delta_z}\mathrm{d}\tilde x\,\varphi_{T_v\tilde B}(z,\tilde c)\cdot(T_v f)(\tilde x)\\
&= \left(\mathrm{Op}_{\mathbb{L}}[Q_k J^*](T_v\tilde B)\,T_v f\right)(z).
\end{align*}

For (cv3):
\begin{align*}
\left(T_v\,\mathrm{Op}_{\mathbb{L}}[J Q_k^T](\tilde B)\,\mathrm{fu}\right)(c,g)
&= \varphi_{\tilde B}([c - v],c - v)^{-1}\cdot\mathrm{fu}([c - v])\\
&= \varphi_{\tilde B}([c] - v,c - v)^{-1}\cdot\mathrm{fu}([c] - v)\\
&= \varphi_{T_v\tilde B}([c],c)^{-1}\cdot(T_v\mathrm{fu})([c])\\
&= \left(\mathrm{Op}_{\mathbb{L}}[J Q_k^T](T_v\tilde B)\,T_v\mathrm{fu}\right)(c,g).
\end{align*}

For (cv4):
\begin{align*}
\left(T_v\,\mathrm{Op}_{\mathbb{L}}[Q](\tilde B)\,\mathrm{fu}\right)(z')
&= L^{-3}\sum_{z\in B(z' - v)}U_{\tilde B}(z' - v,z)\cdot\mathrm{fu}(z)\\
&= L^{-3}\sum_{\tilde z\in B(z')}U_{\tilde B}(z' - v,\tilde z - v)\cdot\mathrm{fu}(\tilde z - v)\\
&= L^{-3}\sum_{\tilde z\in B(z')}U_{T_v\tilde B}(z',\tilde z)\cdot(T_v\mathrm{fu})(\tilde z)\\
&= \left(\mathrm{Op}_{\mathbb{L}}[Q](T_v\tilde B)\,T_v\mathrm{fu}\right)(z').
\end{align*}

For (cv5): if $z\in B(z')$, then $z - v\in B(z' - v)$; hence
\begin{align*}
\left(T_v\,\mathrm{Op}_{\mathbb{L}}[Q^T](\tilde B)\,\mathrm{fL}\right)(z)
&= U_{\tilde B}(z' - v,z - v)^{-1}\cdot\mathrm{fL}(z' - v)\\
&= U_{T_v\tilde B}(z',z)^{-1}\cdot(T_v\mathrm{fL})(z')\\
&= \left(\mathrm{Op}_{\mathbb{L}}[Q^T](T_v\tilde B)\,T_v\mathrm{fL}\right)(z).
\end{align*}

For (cv10):
\begin{align*}
&\left(T_v\,\mathrm{Op}_{\mathbb{L}}[T^{(1)}](\tilde B,\tilde Y)\,f\right)(z)\\
&\quad= \mathbf{1}[z - v\in\tilde Y]\int_{\Delta_{z-v}}\mathrm{d}x\,\varphi_{\tilde B}(z - v,c)\cdot\left\{\left(1 + \frac{\mu}{b_k}\right)f(x)\right.\\
&\qquad\left.{} - (b_k\eta_k)^{-1}\left(\sum_{i\le 2}M_i(x)\cdot\left[f(x^{(i)}) - f(x)\right]\right.\right.\\
&\qquad\left.\left.{} + \mathbf{1}[c_3(x)\in\tilde Y]\,M_3(x)\cdot\left[u_{\tilde B}\left(c,c_3(x)\right)\cdot f(x^{(3)}) - f(x)\right]\right)\right\}\\
&\quad= \mathbf{1}[z\in\tilde Y + v]\int_{\Delta_z}\mathrm{d}\tilde x\,\varphi_{T_v\tilde B}(z,\tilde c)\cdot\left\{\left(1 + \frac{\mu}{b_k}\right)(T_v f)(\tilde x)\right.\\
&\qquad\left.{} - (b_k\eta_k)^{-1}\left(\sum_{i\le 2}M_i(\tilde x)\cdot\left[(T_v f)(\tilde x^{(i)}) - (T_v f)(\tilde x)\right]\right.\right.\\
&\qquad\left.\left.{} + \mathbf{1}[c_3(\tilde x)\in\tilde Y + v]\,M_3(\tilde x)\cdot\left[u_{T_v\tilde B}\left(\tilde c,c_3(\tilde x)\right)\cdot(T_v f)(\tilde x^{(3)})\right.\right.\right.\\
&\qquad\qquad\left.\left.\left.{} - (T_v f)(\tilde x)\right]\right)\right\}\\
&\quad= \left(\mathrm{Op}_{\mathbb{L}}[T^{(1)}](T_v\tilde B,\tilde Y + v)\,T_v f\right)(z),
\end{align*}
because $(\tilde x - v)^{(i)} = \tilde x^{(i)} - v$, $c_3(\tilde x - v) = c_3(\tilde x) - v$, and $c_3(\tilde x) - v\in\tilde Y$ if and only if $c_3(\tilde x)\in\tilde Y + v$.

For (cv11):
\begin{align*}
&\left(T_v\,\mathrm{Op}_{\mathbb{L}}[T^{(2)}](\tilde B,\tilde Y)\,F\right)(z)\\
&\quad= -\mathbf{1}[z - v\in\tilde Y]\,b_k^{-1}\eta_k^{\alpha-1}\int_{\Delta_{z-v}}\mathrm{d}x\,\mathbf{1}[c_3(x)\notin\tilde Y]\\
&\qquad\qquad{}\cdot\varphi_{\tilde B}(z - v,c)\cdot M_3(x)\cdot F(x,x^{(3)})\\
&\quad= -\mathbf{1}[z\in\tilde Y + v]\,b_k^{-1}\eta_k^{\alpha-1}\int_{\Delta_z}\mathrm{d}\tilde x\,\mathbf{1}[c_3(\tilde x)\notin\tilde Y + v]\\
&\qquad\qquad{}\cdot\varphi_{T_v\tilde B}(z,\tilde c)\cdot M_3(\tilde x)\cdot(T_v F)(\tilde x,\tilde x^{(3)})\\
&\quad= \left(\mathrm{Op}_{\mathbb{L}}[T^{(2)}](T_v\tilde B,\tilde Y + v)\,T_v F\right)(z),
\end{align*}
because $F\left(\tilde x - v,(\tilde x - v)^{(3)}\right) = F\left(\tilde x - v,\tilde x^{(3)} - v\right) = (T_v F)(\tilde x,\tilde x^{(3)})$.

For (cv13):
\begin{align*}
\left(T_v\,\mathrm{Op}_{\mathbb{L}}[\delta_\alpha](\tilde B)\,f\right)(x,x')
&= d_X(x - v,x' - v)^{-\alpha}\left(U_{\tilde B}(c - v,c' - v)\cdot f(x' - v) - f(x - v)\right)\\
&= d_X(x,x')^{-\alpha}\left(U_{T_v\tilde B}(c,c')\cdot(T_v f)(x') - (T_v f)(x)\right)\\
&= \left(\mathrm{Op}_{\mathbb{L}}[\delta_\alpha](T_v\tilde B)\,T_v f\right)(x,x').
\end{align*}

The relation is preserved under sums; under products, since
\begin{align*}
T_v\cdot O_1(\tilde B,\tilde Y)\cdot O_2(\tilde B,\tilde Y)
&= O_1(T_v\tilde B,\tilde Y + v)\cdot T_v\cdot O_2(\tilde B,\tilde Y)\\
&= O_1(T_v\tilde B,\tilde Y + v)\cdot O_2(T_v\tilde B,\tilde Y + v)\cdot T_v;
\end{align*}
under products with numbers and with constant matrices, which commute with $T_v$; and under the substitution $\tilde B\mapsto -\tilde B$, since $T_v(-\tilde B) = -T_v\tilde B$.
Therefore it holds for (cv6)--(cv9) and (cv12) as well.

\paragraph{2.3. (f3).}
Let $\ell\in\mathrm{Cl}$, let $B$ be a complex bond function on $T_k = \mathbb{L}_k/\ell\mathbb{Z}^3$ and $Y$ a union of $M$-cubes of $T_k$, and let $\tilde B$, $\tilde Y$ be their $\ell\mathbb{Z}^3$-periodic extensions.

\emph{(1) $\mathrm{Op}_{\mathbb{L}}(\tilde B,\tilde Y)$ maps $\mathrm{Fu}_\ell$ into $\mathrm{Fu}_\ell$.}
Since $\ell$ is a multiple of $LM$, the vector $\ell e_\mu$ lies in $LM\mathbb{Z}^3$; moreover $T_{\ell e_\mu}\tilde B = \tilde B$ and $\tilde Y + \ell e_\mu = \tilde Y$.
By (f2), $T_{\ell e_\mu}$ therefore commutes with $\mathrm{Op}_{\mathbb{L}}(\tilde B,\tilde Y)$.
An element $f$ of $\mathrm{Fu}(\cdot)$ belongs to $\mathrm{Fu}_\ell(\cdot)$ if and only if $T_{\ell e_\mu}f = -f$ for $\mu = 1,2,3$.
Hence, for $f\in\mathrm{Fu}_\ell$,
\[
T_{\ell e_\mu}\,\mathrm{Op}_{\mathbb{L}}(\tilde B,\tilde Y)\,f = \mathrm{Op}_{\mathbb{L}}(\tilde B,\tilde Y)\,T_{\ell e_\mu}f = -\mathrm{Op}_{\mathbb{L}}(\tilde B,\tilde Y)\,f.
\]

\emph{(2) The restriction is the operator of Part~I.}
In Part~I the spinors and the functions of pairs on the torus are the elements of the spaces $\mathrm{Fu}_\ell$ (see 0.1).
The displays (I.2.3), (I.3.1), those of Sections~I.3.3 and I.3.5, (I.1.7) and (I.5.4) express the value of $\mathrm{Op}(B,Y)f$ at an index $t$ of the covering index set as a finite sum over indices of the covering index sets, namely over the neighbours $x^{(i)}$ of $x$, over the $48$ sites of a cluster, over the sites above the block $\Delta_z$ of $\mathbb{L}_k$, over the points of $B(z')$, or over the two points of a pair of $\mathrm{Pzt}_k$.
None of these sums is a sum over a fundamental domain.
The connection enters through $u_B(c,c_3)$, $\varphi_B(z,c)$, $U_B(z',z)$ and $U_B(c,c')$, which (I.1.6) and Section~I.3.2 form with the paths between the points of $\mathbb{L}_k$ just listed (see (Bq6) and (Bq7); for a pair of $\mathrm{Pzt}_k$ the paths are those of the same two points, by Section~I.1.8).
The value of $B$ on the image of a bond $b$ of $\mathbb{L}_k$ is $\tilde B(b)$.
Finally, the image of $z$, resp.\ of $c_3$, in $T_k$ lies in $Y$ if and only if $z$, resp.\ $c_3$, lies in $\tilde Y$.
Consequently, for $f\in\mathrm{Fu}_\ell$ the displays of Part~I are the displays (cv1)--(cv5), (cv10), (cv11), (cv13); in other words, the operators
\[
D_B,\quad Q_k(B)J^*,\quad J Q_k^T(-B),\quad Q(B),\quad Q^T(-B),\quad T^{(i)}_Y(B),\quad \delta_{\alpha,B}
\]
are, in this order, the restrictions of
\begin{gather*}
\mathrm{Op}_{\mathbb{L}}[D](\tilde B),\quad
\mathrm{Op}_{\mathbb{L}}[Q_k J^*](\tilde B),\quad
\mathrm{Op}_{\mathbb{L}}[J Q_k^T](\tilde B),\quad
\mathrm{Op}_{\mathbb{L}}[Q](\tilde B),\\
\mathrm{Op}_{\mathbb{L}}[Q^T](\tilde B),\quad
\mathrm{Op}_{\mathbb{L}}[T^{(i)}](\tilde B,\tilde Y),\quad
\mathrm{Op}_{\mathbb{L}}[\delta_\alpha](\tilde B).
\end{gather*}
For $Q_k(B)J^*$ and $J Q_k^T(-B)$ one uses the explicit forms given after (I.3.4); for $T^{(i)}_Y(B)$ one uses (I.5.4) at the unit sites of $Y$, the rows at the other unit sites being zero by Section~I.5.9.

\emph{(3) Composed operators.}
By (1), a product of restrictions is the restriction of the product.
Hence the operators
\begin{gather*}
J P_k(B)J^* = \left(J Q_k^T(-B)\right)\cdot\left(Q_k(B)J^*\right),\qquad
P(B) = Q^T(-B)\cdot Q(B),\\
J P_{k+1}(B)J^* = \left(J Q_k^T(-B)\right)\cdot P(B)\cdot\left(Q_k(B)J^*\right),
\end{gather*}
together with $\operatorname{Pen}_y(B)$, $F_y(B)$, $T_k(B)$, $\bar T_k(B) = \mathsf{C}^{-1}\cdot T_k(-B)\cdot\mathsf{C}$ and $\bar T^{(i)}_Y(B) = \mathsf{C}^{-1}\cdot T^{(i)}_Y(-B)\cdot\mathsf{C}$ (Sections~I.3.3, I.3.5 and (I.5.4)), are the restrictions of the operators (cv6), (cv7), (cv9), (cv12); here one uses that the periodic extension of $-B$ is $-\tilde B$.

\emph{(4) $B_{k,y}$: the formula named in (a) is the inverse of Section~I.3.5, for every complex bond function.}
For real $y$, Section~I.3.5 defines $B_{k,y}(B)$ as the inverse of $b_k + b_L P(B) + i\gamma_5 y$ on $\mathfrak{S}(1,\ell)$.
By (3), this operator is the restriction of $b_k + b_L\,\mathrm{Op}_{\mathbb{L}}[P](\tilde B) + i\gamma_5 y$ to $\mathrm{Fu}_\ell(\mathbb{Z}^3)$.
By (1), $\mathrm{Op}_{\mathbb{L}}[B_{k,y}](\tilde B)$ maps $\mathrm{Fu}_\ell(\mathbb{Z}^3)$ into itself, and by (al5) its restriction is a two-sided inverse of that operator.
Therefore, for every complex bond function $B$, the operator $b_k + b_L P(B) + i\gamma_5 y$ is bijective on $\mathfrak{S}(1,\ell)$, and $B_{k,y}(B)$ is the restriction of $\mathrm{Op}_{\mathbb{L}}[B_{k,y}](\tilde B)$.
By (cv8) and (3), this restriction is
\[
\left(b_k + i\gamma_5 y\right)^{-1}\cdot\left(1 - P(B)\right) + \left(b_k + b_L + i\gamma_5 y\right)^{-1}\cdot P(B),
\]
which is the formula named in (a).

\subsubsection*{3. $k\ge 1$: (f4) and the ranges}

\begin{named}{Definition}
Let $\mathrm{Ind}$ and $\mathrm{Ind}'$ be index sets.
For every complex bond function $\tilde{B}$ on $\mathbb{L}_k$ and every union $\tilde{Y}$ of $M$-cubes of $\mathbb{L}_k$ (0.1) let
\[
\mathrm{Op}_{\mathbb{L}}\left(\tilde{B},\tilde{Y}\right)\colon\mathrm{Fu}\left(\mathrm{Ind}'\right)\to\mathrm{Fu}(\mathrm{Ind})
\]
be a linear operator, and let $t\mapsto\mathrm{Rd}(t)$ assign to every $t\in\mathrm{Ind}$ a finite set $\mathrm{Rd}(t)$ of sites of $\mathbb{L}_k$.
The family $\mathrm{Op}_{\mathbb{L}}$ has \emph{property (R) with $\mathrm{Rd}$} if the following two conditions hold for every $t\in\mathrm{Ind}$.
\begin{itemize}
\item[(Rg1)] For all $\tilde{B}$, $\tilde{Y}$ the value $\left(\mathrm{Op}_{\mathbb{L}}\left(\tilde{B},\tilde{Y}\right)f\right)(t)$ depends on $f$ only through the values $f(t')$ at the indices $t'\in\mathrm{Ind}'$ with $\operatorname{supp}t'\subset\mathrm{Rd}(t)$.
\item[(Rg2)] For all $f$ one has
\[
\left(\mathrm{Op}_{\mathbb{L}}\left(\tilde{B},\tilde{Y}\right)f\right)(t)=\left(\mathrm{Op}_{\mathbb{L}}\left(\tilde{B}',\tilde{Y}'\right)f\right)(t)
\]
whenever $\tilde{B}$ and $\tilde{B}'$ coincide on the bonds with both ends in $\mathrm{Rd}(t)$ and $\tilde{Y}\cap\mathrm{Rd}(t)=\tilde{Y}'\cap\mathrm{Rd}(t)$.
\end{itemize}
\end{named}

\begin{named}{Lemma 3.1}
\begin{itemize}
\item[(i)] Let a family have property (R) with $\mathrm{Rd}$, and let $\mathrm{rg}$ be a number such that $\operatorname{dist}\left(\tilde{c},\operatorname{supp}t\right)\le\mathrm{rg}$ for all $t$ and all $\tilde{c}\in\mathrm{Rd}(t)$.
Then the family has the properties (f1) and (f4) of Section~I.1.9 with the range $\mathrm{rg}$, and with every larger range.
\item[(ii)] If two families $\mathrm{Fu}\left(\mathrm{Ind}'\right)\to\mathrm{Fu}(\mathrm{Ind})$ have property (R) with $\mathrm{Rd}_1$ and with $\mathrm{Rd}_2$, respectively, then their sum has property (R) with $t\mapsto\mathrm{Rd}_1(t)\cup\mathrm{Rd}_2(t)$.
\item[(iii)] Let $\mathrm{Op}_{\mathbb{L},1}\colon\mathrm{Fu}\left(\mathrm{Ind}''\right)\to\mathrm{Fu}(\mathrm{Ind})$ have property (R) with $\mathrm{Rd}_1$, and let $\mathrm{Op}_{\mathbb{L},2}\colon\mathrm{Fu}\left(\mathrm{Ind}'\right)\to\mathrm{Fu}\left(\mathrm{Ind}''\right)$ have property (R) with $\mathrm{Rd}_2$.
Then the product $\mathrm{Op}_{\mathbb{L},1}\left(\tilde{B},\tilde{Y}\right)\mathrm{Op}_{\mathbb{L},2}\left(\tilde{B},\tilde{Y}\right)$ has property (R) with
\[
\mathrm{Rd}(t):=\mathrm{Rd}_1(t)\cup\bigcup_{\substack{t''\in\mathrm{Ind}''\\ \operatorname{supp}t''\subset\mathrm{Rd}_1(t)}}\mathrm{Rd}_2\left(t''\right).
\]
\item[(iv)] If $\mathrm{Op}_{\mathbb{L}}$ has property (R) with $\mathrm{Rd}$, then so has $\mathrm{Mx}\,\mathrm{Op}_{\mathbb{L}}\,\mathrm{Mx}'$ for all constant matrices $\mathrm{Mx}$, $\mathrm{Mx}'$, and so has the family $\left(\tilde{B},\tilde{Y}\right)\mapsto\mathrm{Op}_{\mathbb{L}}\left(-\tilde{B},\tilde{Y}\right)$.
A constant matrix has property (R) with $\mathrm{Rd}(t)=\operatorname{supp}t$.
\end{itemize}
\end{named}

\begin{proof}
(i) Each of the four index sets $\mathbb{L}^{\mathrm{f}}_k$, $\mathbb{Z}^3$, $L\mathbb{Z}^3$, $\mathrm{Pzt}_k$ contains only finitely many indices whose support lies in a given finite set of sites; over a photon site there are $48$ fermion sites.
Hence, by (Rg1), the row at $t$ has finite support, which is (f1).
Suppose that the entry of the row at the column $t'$ is different from zero for some $\left(\tilde{B},\tilde{Y}\right)$.
Then $\operatorname{supp}t'\subset\mathrm{Rd}(t)$ by (Rg1); so every point of $\operatorname{supp}t'$ has distance at most $\mathrm{rg}$ from $\operatorname{supp}t$, and $\operatorname{dist}\left(\operatorname{supp}t,\operatorname{supp}t'\right)\le\mathrm{rg}$.
A bond with both ends in $\mathrm{Rd}(t)$ has both ends within the distance $\mathrm{rg}$ of $\operatorname{supp}t$, and every site of $\mathrm{Rd}(t)$ lies within the distance $\mathrm{rg}$ of $\operatorname{supp}t$.
Let now $\left(\tilde{B},\tilde{Y}\right)$ and $\left(\tilde{B}',\tilde{Y}'\right)$ be two pairs such that $\tilde{B}$ and $\tilde{B}'$ coincide on the bonds within the distance $\mathrm{rg}$ of $\operatorname{supp}t$, and such that $\tilde{Y}$ and $\tilde{Y}'$ contain the same sites within this distance.
By (Rg2) the row at $t$ takes the same value at the two pairs.
This means that the row at $t$ reads $\tilde{B}$ on the bonds, and $\tilde{Y}$ through the sites, within the distance $\mathrm{rg}$ of $\operatorname{supp}t$, in the sense of \cite[Definition~5.5(1)]{QEDI}; this is (f4).
As the hypothesis of (i) holds with every number larger than $\mathrm{rg}$ in place of $\mathrm{rg}$, so does the conclusion.

(ii) The value of the sum at $t$ is the sum of the two values at $t$.

(iv) One has
\[
\left(\mathrm{Mx}\,\mathrm{Op}_{\mathbb{L}}\left(\tilde{B},\tilde{Y}\right)\mathrm{Mx}'f\right)(t)=\mathrm{Mx}\cdot\left(\mathrm{Op}_{\mathbb{L}}\left(\tilde{B},\tilde{Y}\right)\left(\mathrm{Mx}'f\right)\right)(t),\qquad\left(\mathrm{Mx}'f\right)(t')=\mathrm{Mx}'\cdot f(t').
\]
Moreover, $\tilde{B}$ and $\tilde{B}'$ coincide on a bond if and only if $-\tilde{B}$ and $-\tilde{B}'$ coincide on it.

(iii) By (Rg1) and (Rg2) for $\mathrm{Op}_{\mathbb{L},1}$, the value $\left(\mathrm{Op}_{\mathbb{L},1}\mathrm{Op}_{\mathbb{L},2}f\right)(t)$ is a function of the values $\left(\mathrm{Op}_{\mathbb{L},2}f\right)(t'')$ with $\operatorname{supp}t''\subset\mathrm{Rd}_1(t)$, of the values of $\tilde{B}$ on the bonds with both ends in $\mathrm{Rd}_1(t)$, and of $\tilde{Y}\cap\mathrm{Rd}_1(t)$.
By (Rg1) and (Rg2) for $\mathrm{Op}_{\mathbb{L},2}$, each value $\left(\mathrm{Op}_{\mathbb{L},2}f\right)(t'')$ is in turn a function of the values $f(t')$ with $\operatorname{supp}t'\subset\mathrm{Rd}_2(t'')$, of the values of $\tilde{B}$ on the bonds with both ends in $\mathrm{Rd}_2(t'')$, and of $\tilde{Y}\cap\mathrm{Rd}_2(t'')$.
\end{proof}

\paragraph{3.2. The sets $\mathrm{Rd}$.}
Put
\[
\Delta_z^{+}:=\Delta_z\cup\left\{c_3(x)\colon x\text{ over }\Delta_z\right\}.
\]
As $|c_3(x)-c|=\eta_k$, every point of $\Delta_z^{+}$ has distance at most $(1-\eta_k)/2+\eta_k=(1+\eta_k)/2$ from $z$.
In the list below, $z'$ denotes the point of $L\mathbb{Z}^3$ with $[c]\in B(z')$ for the rows at $x=(c,g)$, and the point of $L\mathbb{Z}^3$ with $z\in B(z')$ for the rows at $z\in\mathbb{Z}^3$; in the first case $c\in\mathrm{Bl}_{k+1}(z')$, by (Bq4).
For every operator the list gives the row index $t$, the set $\mathrm{Rd}(t)$, the reason for property (R), a bound for $\operatorname{dist}\left(\tilde{c},\operatorname{supp}t\right)$, $\tilde{c}\in\mathrm{Rd}(t)$, and the range named in Proposition~11\,(a).
\begin{description}
\item[{$\mathrm{Op}_{\mathbb{L}}[D]$}] Row index $x$; $\mathrm{Rd}(x)=\{c,c_3\}$; reason: (cv1), whose columns are $x$, $x^{(0)}$, $x^{(1)}$, $x^{(2)}$ over $c$ and $x^{(3)}$ over $c_3$, and whose only bond is $(c,c_3)$; bound $\eta_k$; range $\eta_k$.
\item[{$\mathrm{Op}_{\mathbb{L}}\left[Q_kJ^{*}\right]$}] Row index $z$; $\mathrm{Rd}(z)=\Delta_z$; reason: (cv2), (Bq8); bound $(1-\eta_k)/2$; range $1/2$.
\item[{$\mathrm{Op}_{\mathbb{L}}\left[JQ_k^{T}\right]$}] Row index $x$; $\mathrm{Rd}(x)=\Delta_{[c]}$; reason: (cv3), whose only column is $[c]$, and (Bq8); bound $1-\eta_k$; range $1$.
\item[{$\mathrm{Op}_{\mathbb{L}}\left[JP_kJ^{*}\right]$}] Row index $x$; $\mathrm{Rd}(x)=\Delta_{[c]}$; reason: Lemma~3.1\,(iii), since $\Delta_{[c]}\cap\mathbb{Z}^3=\{[c]\}$; bound $1-\eta_k$; range $1$.
\item[{$\mathrm{Op}_{\mathbb{L}}[T_k]$, $\mathrm{Op}_{\mathbb{L}}\left[\bar{T}_k\right]$}] Row index $z$; $\mathrm{Rd}(z)=\Delta_z^{+}$; reason: Lemma~3.1\,(ii)--(iv), which give $\Delta_z$ together with the union of the sets $\{c,c_3\}$, $x$ over $\Delta_z$; bound $(1+\eta_k)/2$; range $1$.
\item[{$\mathrm{Op}_{\mathbb{L}}\left[T^{(1)}\right]$, $\mathrm{Op}_{\mathbb{L}}\left[\bar{T}^{(1)}\right]$}] Row index $z$; $\mathrm{Rd}(z)=\Delta_z^{+}$; reason: (cv10), whose columns are the sites over $\Delta_z$ and the sites $x^{(3)}$, whose bonds are the bonds with both ends in $\Delta_z$ (by (Bq8)) and the bonds $(c,c_3)$, and which reads $\tilde{Y}$ at $z$ and at the sites $c_3$; bound $(1+\eta_k)/2$; range $1$.
\item[{$\mathrm{Op}_{\mathbb{L}}\left[T^{(2)}\right]$, $\mathrm{Op}_{\mathbb{L}}\left[\bar{T}^{(2)}\right]$}] Row index $z$; $\mathrm{Rd}(z)=\Delta_z^{+}$; reason: (cv11), whose columns are the pairs $\left(x,x^{(3)}\right)$, with the support $\{c,c_3\}$, whose bonds are the bonds with both ends in $\Delta_z$, and which reads $\tilde{Y}$ at $z$ and at the sites $c_3$; bound $(1+\eta_k)/2$; range $1$.
\item[{$\mathrm{Op}_{\mathbb{L}}[Q]$}] Row index $z'$; $\mathrm{Rd}(z')=\mathrm{Bc}(z')$; reason: (cv4), whose columns lie in $B(z')$, and (Bq7), (Bq8); bound $(L-1)/2$; range $L/2$.
\item[{$\mathrm{Op}_{\mathbb{L}}\left[Q^{T}\right]$}] Row index $z$; $\mathrm{Rd}(z)=\mathrm{Bc}(z')$; reason: (cv5), whose only column is $z'$, and (Bq8); bound $L-1$; range $L$.
\item[{$\mathrm{Op}_{\mathbb{L}}[P]$, $\mathrm{Op}_{\mathbb{L}}\left[B_{k,y}\right]$}] Row index $z$; $\mathrm{Rd}(z)=\mathrm{Bc}(z')$; reason: Lemma~3.1\,(ii)--(iv), since $\mathrm{Bc}(z')\cap L\mathbb{Z}^3=\{z'\}$; bound $L-1$; range $L$.
\item[{$\mathrm{Op}_{\mathbb{L}}\left[JP_{k+1}J^{*}\right]$, $\mathrm{Op}_{\mathbb{L}}\left[\operatorname{Pen}_{y}\right]$}] Row index $x$; $\mathrm{Rd}(x)=\mathrm{Bl}_{k+1}(z')$; reason: Lemma~3.1\,(ii)--(iv) give the set
\[
\Delta_{[c]}\cup\mathrm{Bc}(z')\cup\bigcup_{z_2\in\mathrm{Bc}(z')\cap\mathbb{Z}^3}\Delta_{z_2},\qquad\mathrm{Bc}(z')\cap\mathbb{Z}^3=B(z'),
\]
which is $\mathrm{Bl}_{k+1}(z')$ by (Bq4), (Bq7); bound $L-\eta_k$; range $L$.
\item[{$\mathrm{Op}_{\mathbb{L}}[F_y]$}] Row index $x$; $\mathrm{Rd}(x)=\mathrm{Bl}_{k+1}(z')\cup\{c_3\}$; reason: Lemma~3.1\,(ii); bound $\max\left(L-\eta_k,\eta_k\right)=L-\eta_k$; range $L$.
\item[{$\mathrm{Op}_{\mathbb{L}}[\delta_{\alpha}]$}] Row index $(x,x')$; $\mathrm{Rd}(x,x')=\mathrm{bx}(c,c')$; reason: (cv13), whose columns are $x$ and $x'$, and (Bq8); for $c=c'$ one has $\mathrm{bx}(c,c')=\{c\}$, and no bond occurs; bound: the number $|c-c'|\le M$, by (Bq5); range $M$.
\end{description}
The bounds come from (Bq4) (the distance of the points of $\Delta_z$ from $z$; the sup-diameters of $\Delta_{[c]}$ and of $\mathrm{Bl}_{k+1}(z')$), from (Bq7) (the distance of the points of $\mathrm{Bc}(z')$ from $z'$; the sup-diameter of $\mathrm{Bc}(z')$) and from (Bq5).
Moreover $(1+\eta_k)/2\le1$, because $\eta_k\le1$.
By Lemma~3.1\,(i), every operator of the list has property (f4) with the range named in (a).
Together with Sections 1 and 2 this proves (a) for $k\ge1$: for the class $\mathrm{Cl}$, every operator of the list of (a) is given by one formula on $\mathbb{L}^{\mathrm{f}}_k$, of the range named there.

\begin{named}{Remark 3.3}
For the operators of the list in 3.2, with $\mathrm{rg}$ the range named in (a), every point of $\operatorname{supp}t$ also has distance at most $\mathrm{rg}$ from $\operatorname{supp}t'$ whenever the entry at $(t,t')$ is different from zero.
Indeed, $\operatorname{supp}t$ is a single point, except for $\mathrm{Op}_{\mathbb{L}}[\delta_{\alpha}]$; for this operator $t=(x,x')$, the column $t'$ is $x$ or $x'$, and $|c-c'|\le M$.
\end{named}

\begin{named}{Remark 3.4}[Smaller sets for two rows]
The operator $\mathrm{Op}_{\mathbb{L}}\left[Q^{T}\right]$ has property (R) also with $\mathrm{Rd}(z)=\mathrm{bx}(z',z)$, by (cv5) and (Bq8); the points of this set have distance at most $|z-z'|\le(L-1)/2$ from $z$, by (Bq5).
The operator $\mathrm{Op}_{\mathbb{L}}\left[JQ_k^{T}\right]$ has property (R) also with
\[
\mathrm{Rd}(x)=\bigcup_{j<k}\mathrm{bx}\left(z_{j+1}(c),z_j(c)\right),
\]
by (cv3) and (Bq5); this set contains $c=z_0(c)$ and $[c]=z_k(c)$.
Its points have distance at most $(1-\eta_k)/2$ from $c$.
Indeed, $c$ lies in $\mathrm{Bl}_j\left(z_j(c)\right)$, so that $|z_j(c)-c|\le(L^j-1)\cdot\eta_k/2\le(1-\eta_k)/2$ for $j\le k$, by (Bq1); and every coordinate of a point of $\mathrm{bx}\left(z_{j+1}(c),z_j(c)\right)$ lies between the corresponding coordinates of $z_{j+1}(c)$ and of $z_j(c)$.
\end{named}

\subsubsection*{4. Clause (a) at $k=0$}

Let $k=0$.
Then $\eta_0=1$, $\mathbb{L}_0=\mathbb{Z}^3$ and $\mathbb{L}^{\mathrm{f}}_0=\mathbb{Z}^3\times O_{\mathrm{cub}}$; the fine lattice and the unit lattice coincide, and $\Delta_z=\{z\}$.
There are no towers and no phases $\varphi$.
Further, $x^{(3)}=\left(c+g\,e_1,g\,s_3\right)$, $d_X(x,x')=\max\left(|c-c'|,1\right)$, and
\[
\int_Z\mathrm{d}x\,f(x)=\frac{1}{48}\sum_{c\in Z}\sum_{g\in O_{\mathrm{cub}}}f(c,g).
\]
The facts of 0.1 and (Bq1)--(Bq5), (Bq7)--(Bq9) hold with $k=0$, and $\mathrm{Bc}(z')=B(z')=\mathrm{Bl}_1(z')$, a set of sup-diameter $L-1$.
The index sets are $\mathbb{L}^{\mathrm{f}}_0$, $L\mathbb{Z}^3$ and $\mathrm{Pzt}_0$.
On the covering lattice the operators of Section~I.3.9 and of (I.3.4) are the following, for $f\in\mathrm{Fu}\left(\mathbb{L}^{\mathrm{f}}_0\right)$ and $\mathrm{fL}\in\mathrm{Fu}\left(L\mathbb{Z}^3\right)$.
The operator $\mathrm{Op}_{\mathbb{L}}[D](\tilde{B})$ is given by the display (cv1) at $\eta_0=1$, and $\mathrm{Op}_{\mathbb{L}}[D_0]$ is formed from it:
\begin{equation}\tag{cv0.1}
\begin{aligned}
\left(\mathrm{Op}_{\mathbb{L}}[D](\tilde{B})f\right)(x)&=-\sum_{i=0}^{2}M_{i}(x)\cdot\left[f\left(x^{(i)}\right)-f(x)\right]-M_{3}(x)\cdot\left[u_{\tilde{B}}(c,c_3)\cdot f\left(x^{(3)}\right)-f(x)\right],\\
\mathrm{Op}_{\mathbb{L}}[D_0]&:=\mathrm{Op}_{\mathbb{L}}[D]+\mu.
\end{aligned}
\end{equation}
The averaging operators are
\begin{equation}\tag{cv0.2}
\left(\mathrm{Op}_{\mathbb{L}}\left[QJ^{*}\right](\tilde{B})f\right)(z'):=L^{-3}\sum_{c\in B(z')}U_{\tilde{B}}(z',c)\cdot\frac{1}{48}\sum_{g\in O_{\mathrm{cub}}}f(c,g)\qquad\left(z'\in L\mathbb{Z}^3\right)
\end{equation}
and
\begin{equation}\tag{cv0.3}
\left(\mathrm{Op}_{\mathbb{L}}\left[JQ^{T}\right](\tilde{B})\,\mathrm{fL}\right)(c,g):=U_{\tilde{B}}(z',c)^{-1}\cdot\mathrm{fL}(z'),
\end{equation}
where $z'$ is the point of $L\mathbb{Z}^3$ with $c\in B(z')$; the value is the same at the $48$ sites over $c$, and this operator is the one that belongs to $JQ^{T}(-A)$.
The composed operators are
\begin{equation}\tag{cv0.4}
\begin{aligned}
\mathrm{Op}_{\mathbb{L}}\left[JPJ^{*}\right]&:=\mathrm{Op}_{\mathbb{L}}\left[JQ^{T}\right]\mathrm{Op}_{\mathbb{L}}\left[QJ^{*}\right],\qquad\mathrm{Op}_{\mathbb{L}}\left[\operatorname{Pen}_{y}\right]:=b_{L}\cdot\mathrm{Op}_{\mathbb{L}}\left[JPJ^{*}\right]+i\gamma_5\cdot y,\\
\mathrm{Op}_{\mathbb{L}}[F_y]&:=\mathrm{Op}_{\mathbb{L}}[D]+\mu+\mathrm{Op}_{\mathbb{L}}\left[\operatorname{Pen}_{y}\right]\qquad\left(y\in Y_0=\mathbb{R}\right).
\end{aligned}
\end{equation}
The operator $\mathrm{Op}_{\mathbb{L}}[\delta_{\alpha}](\tilde{B})$ is given by the display (cv13) at $\eta_0=1$:
\begin{equation}\tag{cv0.5}
\left(\mathrm{Op}_{\mathbb{L}}[\delta_{\alpha}](\tilde{B})f\right)(x,x')=\max\left(|c-c'|,1\right)^{-\alpha}\cdot\left(U_{\tilde{B}}(c,c')\cdot f(x')-f(x)\right);
\end{equation}
for a pair inside a cluster this is $f(c,g')-f(c,g)$.
Finally,
\begin{equation}\tag{cv0.6}
\mathrm{Op}_{\mathbb{L}}[H_0]=\mathrm{Op}_{\mathbb{L}}\left[\bar{H}_0\right]=\mathrm{Op}_{\mathbb{L}}[T_0]=\mathrm{Op}_{\mathbb{L}}\left[\bar{T}_0\right]=\mathrm{Op}_{\mathbb{L}}\left[\dot{M}_0\right]:=1\quad\text{on }\mathrm{Fu}\left(\mathbb{L}^{\mathrm{f}}_0\right).
\end{equation}

\emph{Property (f1).}
The displays (cv0.1), (cv0.2), (cv0.3) and (cv0.5) are finite sums: a row has $5$ columns, $48L^3$ columns, one column and two columns, respectively.

\emph{Property (f2).}
Let $v\in(L\cdot M)\mathbb{Z}^3$.
For (cv0.1) and (cv0.5) the computation is that of the lines (cv1) and (cv13) of 2.2 at $\eta_0=1$.
For (cv0.2),
\begin{align*}
\left(T_v\,\mathrm{Op}_{\mathbb{L}}\left[QJ^{*}\right](\tilde{B})f\right)(z')
&=L^{-3}\sum_{c\in B(z'-v)}U_{\tilde{B}}(z'-v,c)\cdot\frac{1}{48}\sum_{g}f(c,g)\\
&=L^{-3}\sum_{\tilde{c}\in B(z')}U_{\tilde{B}}(z'-v,\tilde{c}-v)\cdot\frac{1}{48}\sum_{g}f(\tilde{c}-v,g)\\
&=L^{-3}\sum_{\tilde{c}\in B(z')}U_{T_v\tilde{B}}(z',\tilde{c})\cdot\frac{1}{48}\sum_{g}\left(T_vf\right)(\tilde{c},g)
=\left(\mathrm{Op}_{\mathbb{L}}\left[QJ^{*}\right]\left(T_v\tilde{B}\right)T_vf\right)(z').
\end{align*}
For (cv0.3),
\[
\left(T_v\,\mathrm{Op}_{\mathbb{L}}\left[JQ^{T}\right](\tilde{B})\,\mathrm{fL}\right)(c,g)=U_{\tilde{B}}(z'-v,c-v)^{-1}\cdot\mathrm{fL}(z'-v)=U_{T_v\tilde{B}}(z',c)^{-1}\cdot\left(T_v\mathrm{fL}\right)(z'),
\]
because $c-v$ lies in $B(z'-v)$.
By the last paragraph of 2.2, the relation extends to sums, to products and to products with constant matrices.

\emph{Property (f3).}
The argument of 2.3\,(1) uses only (f2) and the fact that $\ell\cdot e_{\mu}\in(L\cdot M)\mathbb{Z}^3$; it holds at $k=0$.
Let $\ell\in\mathrm{Cl}$, let $B$ be a complex bond function on the torus $T_0$ of side $\ell$, and let $\tilde{B}$ be its periodic extension.
For $f\in\mathfrak{S}^{\mathrm{ch}}(1,\ell)=\mathrm{Fu}_{\ell}\left(\mathbb{L}^{\mathrm{f}}_0\right)$ and $\mathrm{fL}\in\mathfrak{S}(L,\ell)=\mathrm{Fu}_{\ell}\left(L\mathbb{Z}^3\right)$, the displays of (I.2.3) and of Section~I.3.3 with (I.3.1) and (I.1.7) for $D_Bf$, $Q(B)J^{*}f$, $JQ^{T}(-B)\,\mathrm{fL}$ and $\delta_{\alpha,B}f$ are (cv0.1)--(cv0.3) and (cv0.5) with $\tilde{B}$.
Indeed, they are sums over the neighbours of $x$, over $B(z')\times O_{\mathrm{cub}}$ and over the two points of a pair, and the paths are those between the points $z'$, $c$, respectively $c$, $c'$, of $\mathbb{Z}^3$.
The operators
\begin{gather*}
D_0(B)=D_B+\mu,\qquad JP(B)J^{*}=\left(JQ^{T}(-B)\right)\left(Q(B)J^{*}\right),\qquad\operatorname{Pen}_{y}(B)=b_{L}\cdot JP(B)J^{*}+i\gamma_5\cdot y,\\
F_y(B)=D_B+\mu+\operatorname{Pen}_{y}(B),\qquad H_0=\bar{H}_0=T_0=\bar{T}_0=\dot{M}_0=1
\end{gather*}
(Section~I.3.9) are the restrictions of the operators (cv0.1), (cv0.4), (cv0.6), since a product of restrictions is the restriction of the product.

\emph{Property (f4).}
Lemma~3.1 applies with the following sets, where $z'$ denotes the point of $L\mathbb{Z}^3$ with $c\in B(z')$; the entries are arranged as in 3.2.
\begin{description}
\item[{$\mathrm{Op}_{\mathbb{L}}[D]$, $\mathrm{Op}_{\mathbb{L}}[D_0]$}] Row index $x$; $\mathrm{Rd}(x)=\{c,c_3\}$; reason: (cv0.1); bound $1$; range $1$.
\item[{$\mathrm{Op}_{\mathbb{L}}\left[QJ^{*}\right]$}] Row index $z'$; $\mathrm{Rd}(z')=B(z')$; reason: (cv0.2), (Bq7), (Bq8); bound $(L-1)/2$; range $L/2$.
\item[{$\mathrm{Op}_{\mathbb{L}}\left[JQ^{T}\right]$}] Row index $x$; $\mathrm{Rd}(x)=B(z')$; reason: (cv0.3), whose only column is $z'$, and (Bq8); bound $L-1$; range $L$.
\item[{$\mathrm{Op}_{\mathbb{L}}\left[JPJ^{*}\right]$, $\mathrm{Op}_{\mathbb{L}}\left[\operatorname{Pen}_{y}\right]$}] Row index $x$; $\mathrm{Rd}(x)=B(z')$; reason: Lemma~3.1\,(ii)--(iv), since $B(z')\cap L\mathbb{Z}^3=\{z'\}$; bound $L-1$; range $L$.
\item[{$\mathrm{Op}_{\mathbb{L}}[F_y]$}] Row index $x$; $\mathrm{Rd}(x)=B(z')\cup\{c_3\}$; reason: Lemma~3.1\,(ii); bound $\max(L-1,1)=L-1$; range $L$.
\item[{$\mathrm{Op}_{\mathbb{L}}[\delta_{\alpha}]$}] Row index $(x,x')$; $\mathrm{Rd}(x,x')=\mathrm{bx}(c,c')$; reason: (cv0.5), (Bq8); bound: the number $|c-c'|\le M$; range $M$.
\item[{$\mathrm{Op}_{\mathbb{L}}[H_0]$, $\mathrm{Op}_{\mathbb{L}}\left[\bar{H}_0\right]$, $\mathrm{Op}_{\mathbb{L}}[T_0]$, $\mathrm{Op}_{\mathbb{L}}\left[\bar{T}_0\right]$, $\mathrm{Op}_{\mathbb{L}}\left[\dot{M}_0\right]$}] Row index $x$; $\mathrm{Rd}(x)=\{c\}$; reason: the identity; bound $0$; range $0$.
\end{description}
The operator $\mathrm{Op}_{\mathbb{L}}\left[JQ^{T}\right]$ has property (R) also with $\mathrm{Rd}(x)=\mathrm{bx}(z',c)$, by (cv0.3) and (Bq8); the points of this set have distance at most $|c-z'|\le(L-1)/2$ from $c$, by (Bq5).
This proves (a) for $k=0$.

\subsubsection*{6. Clause (b)}

The operators named in the last sentence of (a), namely $S_y(A)$, $\Gamma_{k,y}(A)$, $H_k$, $\bar{H}_k$, $\mathfrak{H}_k$, $D_k$, $M_k$, $\dot{M}_k$ for $k\ge1$ and $S_{0,y}(A)$, $\Gamma_0(A)$ for $k=0$, are those operators of the lists of (b) that do not belong to the lists of (a);
for them, clause (b) is obtained by Steps (ii) and (iii) (6.3, 6.4).
The data are those of the statement: $\ell_1$ and $\ell_2=\mathrm{Jr}\cdot\ell_1$ lie in $\mathrm{Cl}$, $\mathrm{Jr}$ is odd, $\mathrm{pr}_i\colon\mathbb{L}_k\to T_i$ are the covering maps, and $\mathrm{pj}\colon T_2\to T_1$ is the covering map with $\mathrm{pj}\circ\mathrm{pr}_2=\mathrm{pr}_1$.

\paragraph{6.0. The map $\mathrm{pj}^{*}$.}
For every index set $\mathrm{Ind}$ the space $\mathrm{Fu}_{\ell_1}(\mathrm{Ind})$ is a subspace of $\mathrm{Fu}_{\ell_2}(\mathrm{Ind})$: for $f\in\mathrm{Fu}_{\ell_1}(\mathrm{Ind})$,
\[
f\left(t+\ell_2\cdot e_{\mu}\right)=f\left(t+\mathrm{Jr}\cdot\ell_1\cdot e_{\mu}\right)=(-1)^{\mathrm{Jr}}f(t)=-f(t),
\]
since $\mathrm{Jr}$ is odd.
For $\mathrm{Ind}=\mathbb{L}^{\mathrm{f}}_k$, $\mathbb{Z}^3$, $L\mathbb{Z}^3$, $\mathrm{Pzt}_k$ the map $\mathrm{pj}^{*}$ is this inclusion; for pairs both points are shifted (0.1), and the label is a spectator (Section~I.6.3, (I.J3)).
An $\ell_1\mathbb{Z}^3$-periodic bond function, scalar function or union of $M$-cubes of $\mathbb{L}_k$ is $\ell_2\mathbb{Z}^3$-periodic.
Consequently, if $A$ is a bond function on $T_1$ with the $\ell_1\mathbb{Z}^3$-periodic extension $\tilde{A}$, then $\mathrm{pj}^{*}A$ is the bond function on $T_2$ whose $\ell_2\mathbb{Z}^3$-periodic extension is $\tilde{A}$.
If $\mathbf{Y}$ is a union of $M$-cubes of $T_1$ with the preimage $\tilde{Y}=\mathrm{pr}_1^{-1}\mathbf{Y}$ in $\mathbb{L}_k$, then the preimage of $\mathrm{pj}^{-1}\mathbf{Y}$ under $\mathrm{pr}_2$ is
\[
\mathrm{pr}_2^{-1}\,\mathrm{pj}^{-1}\mathbf{Y}=\mathrm{pr}_1^{-1}\mathbf{Y}=\tilde{Y},
\]
and $\mathrm{pj}^{-1}\mathbf{Y}$ is a union of $M$-cubes of $T_2$.
The map $\mathrm{pj}^{*}$ commutes with numbers and with constant $4\times4$ matrices, and $\mathrm{pj}^{*}(-A)=-\mathrm{pj}^{*}A$.

\paragraph{6.1. Domains.}
(1) \emph{If $A\in\tilde{G}_k(T_1)$, then $\mathrm{pj}^{*}A\in\tilde{G}_k(T_2)$.}
This is \cite[Lemma~5.20]{QEDI}, proved as \cite[Lemma~17.7]{QEDI}; it holds for every $k\ge0$, every $e_k>0$ and all $\ell_1,\ell_2\in\mathrm{Cl}$ with $\ell_2=\mathrm{Jr}\cdot\ell_1$, $\mathrm{Jr}$ odd, with $\mathrm{pj}^{*}$ the identification of 6.0, and $\tilde{G}_k$ is among the sets to which it applies.
For $\tilde{G}_k$ the proof is as follows.
Let $A=A_0+A_1$ be as in the definition of Section~I.5.3 on $T_1$.
Then $\mathrm{pj}^{*}A=\mathrm{pj}^{*}A_0+\mathrm{pj}^{*}A_1$, the bond function $\mathrm{pj}^{*}A_0$ is real, and $\mathrm{pj}^{*}A_1$ has the same set of values as $A_1$, so that $\left|\mathrm{pj}^{*}A_1\right|_{\infty}=|A_1|_{\infty}<e_k^{-1+\varepsilon}$.
Let $\square_2$ be an $M$-cube of $T_2$, let $\tilde{z}\in M\mathbb{Z}^3$ be a lift of its centre, put
\[
\tilde{C}:=\left\{c\in\mathbb{L}_k\colon|c-\tilde{z}|<11M/2\right\},
\]
and let $\square_1$ be the $M$-cube of $T_1$ with the centre $\mathrm{pr}_1(\tilde{z})$.
The set $\square_i^{\sim5}$ is the cube of side $11M$ with the centre of $\square_i$ (Section~I.1.1; \cite[Definition~5.7]{QEDI}), that is, $\square_i^{\sim5}=\mathrm{pr}_i(\tilde{C})$.
The sup-diameter of $\tilde{C}$ is $11M-\eta_k$ (the computation of (Bq1), with the odd integer $11M/\eta_k$ in place of $L^j$), and $11M<13\cdot L\cdot M\le\ell_1\le\ell_2$.
Therefore $\tilde{C}$ is a lift (Section~I.1.1) of $\square_1^{\sim5}$ and of $\square_2^{\sim5}$: the maps $\mathrm{pr}_1$ and $\mathrm{pr}_2$ are injective on $\tilde{C}$, and the bonds of $\mathrm{pr}_i(\tilde{C})$ are the images of the bonds of $\tilde{C}$ (Section~I.1.1 with $\mathrm{sd}=\eta_k$).
By the definition of $\tilde{G}_k(T_1)$ there are a real function $\boldsymbol{\lambda}$ on the sites and a real function $A'$ on the bonds of $\square_1^{\sim5}$ with $A_0=A'+d\boldsymbol{\lambda}$ and $|A'|_{\infty}<e_k^{-1+\varepsilon}$ there.
For the sites $c$ and the bonds $b$ of $\tilde{C}$ define
\[
\boldsymbol{\lambda}_2(\mathrm{pr}_2c):=\boldsymbol{\lambda}(\mathrm{pr}_1c),\qquad A'_2(\mathrm{pr}_2b):=A'(\mathrm{pr}_1b).
\]
For a bond $b=(c,c')$ of $\tilde{C}$ one has
\[
\left(\mathrm{pj}^{*}A_0\right)(\mathrm{pr}_2b)=A_0(\mathrm{pr}_1b)=A'(\mathrm{pr}_1b)+\eta_k^{-1}\cdot\left(\boldsymbol{\lambda}(\mathrm{pr}_1c')-\boldsymbol{\lambda}(\mathrm{pr}_1c)\right)=A'_2(\mathrm{pr}_2b)+\left(d\boldsymbol{\lambda}_2\right)(\mathrm{pr}_2b),
\]
and $|A'_2|_{\infty}=|A'|_{\infty}$.
Thus $\mathrm{pj}^{*}A_0=A'_2+d\boldsymbol{\lambda}_2$ on $\square_2^{\sim5}$, with $|A'_2|_{\infty}<e_k^{-1+\varepsilon}$, for every $M$-cube $\square_2$ of $T_2$.

(2) \emph{If $A\in\tilde{G}_k(T_i)$, then $-A\in\tilde{G}_k(T_i)$.}
This is the invariance under $A\mapsto-A$ in \cite[Lemma~5.19]{QEDI}, proved as \cite[Lemma~17.6]{QEDI}.
Indeed, $-A=(-A_0)+(-A_1)$, with $-A'$ and $-\boldsymbol{\lambda}$ on every $\square^{\sim5}$, and the conditions of Section~I.5.3 involve only reality and moduli.

\paragraph{6.2. Step (i): the operators of (a).}
Let $O$ be an operator of the list of (a) (for $k\ge1$: Sections 1--3; for $k=0$: Section 4), let $A$ be a complex bond function on $T_1$ and $\mathbf{Y}$ a union of $M$-cubes of $T_1$, and let $\tilde{A}$, $\tilde{Y}$ be as in 6.0.
By (f3) for the side $\ell_1$, the operator $O^{(1)}(A)$, respectively $O^{(1)}_{\mathbf{Y}}(A)$, is the restriction of $\mathrm{Op}_{\mathbb{L}}[O]\left(\tilde{A},\tilde{Y}\right)$ to $\mathrm{Fu}_{\ell_1}$.
By (f3) for the side $\ell_2$ and by 6.0, the operator $O^{(2)}\left(\mathrm{pj}^{*}A\right)$, respectively $O^{(2)}_{\mathrm{pj}^{-1}\mathbf{Y}}\left(\mathrm{pj}^{*}A\right)$, is the restriction of the \emph{same} operator $\mathrm{Op}_{\mathbb{L}}[O]\left(\tilde{A},\tilde{Y}\right)$ to $\mathrm{Fu}_{\ell_2}$.
Hence, for $f\in\mathrm{Fu}_{\ell_1}$,
\[
O^{(2)}\left(\mathrm{pj}^{*}A\right)\mathrm{pj}^{*}f=\mathrm{Op}_{\mathbb{L}}[O]\left(\tilde{A},\tilde{Y}\right)f=O^{(1)}(A)\,f=\mathrm{pj}^{*}O^{(1)}(A)\,f,
\]
and in the same way with $O^{(2)}_{\mathrm{pj}^{-1}\mathbf{Y}}$ and $O^{(1)}_{\mathbf{Y}}$ in place of $O^{(2)}$ and $O^{(1)}$.
This holds for every complex $A$, in particular for $A\in\tilde{G}_k(T_1)$.

\paragraph{6.3. Step (ii): sums and products.}
If $O_a^{(2)}\mathrm{pj}^{*}=\mathrm{pj}^{*}O_a^{(1)}$ and $O_b^{(2)}\mathrm{pj}^{*}=\mathrm{pj}^{*}O_b^{(1)}$, then
\[
\left(O_a^{(2)}+O_b^{(2)}\right)\mathrm{pj}^{*}=\mathrm{pj}^{*}\left(O_a^{(1)}+O_b^{(1)}\right),\qquad O_a^{(2)}O_b^{(2)}\,\mathrm{pj}^{*}=O_a^{(2)}\,\mathrm{pj}^{*}O_b^{(1)}=\mathrm{pj}^{*}O_a^{(1)}O_b^{(1)},
\]
whenever the sums and products are defined.
The relation holds for numbers and for constant matrices (6.0).

\paragraph{6.4. Step (iii): inverses.}
Let $O^{(2)}\mathrm{pj}^{*}=\mathrm{pj}^{*}O^{(1)}$, and let $O^{(1)}$ and $O^{(2)}$ be bijective.
For $f$ in the target space of $O^{(1)}$,
\[
O^{(2)}\,\mathrm{pj}^{*}\left(O^{(1)}\right)^{-1}f=\mathrm{pj}^{*}O^{(1)}\left(O^{(1)}\right)^{-1}f=\mathrm{pj}^{*}f,
\]
and applying $\left(O^{(2)}\right)^{-1}$ gives $\mathrm{pj}^{*}\left(O^{(1)}\right)^{-1}f=\left(O^{(2)}\right)^{-1}\mathrm{pj}^{*}f$.

\paragraph{6.5. The list for $k\ge1$.}
Let $A\in\tilde{G}_k(T_1)$.
By 6.1, the bond functions $\mathrm{pj}^{*}A$, $-A$ and $\mathrm{pj}^{*}(-A)=-\mathrm{pj}^{*}A$ lie in $\tilde{G}_k$ of their tori.
The data $(k,\ell_i,\mu,e_k)$, $i=1,2$, are data (SF).
\begin{itemize}
\item[(1)] For $D_A$, $Q_k(A)J^{*}$, $JQ_k^{T}(-A)$, $JP_k(A)J^{*}$, $JP_{k+1}(A)J^{*}$, $Q(A)$, $Q^{T}(-A)$, $P(A)$, $\operatorname{Pen}_{y}(A)$ and $F_y(A)$ ($y\in Y_k$), $B_{k,y}(A)$ ($y$ real; $B_k=B_{k,0}$), $T_k(A)$, $\bar{T}_k(A)$, $\delta_{\alpha,A}$, and for $T^{(i)}_{\mathbf{Y}}(A)$, $\bar{T}^{(i)}_{\mathbf{Y}}(A)$, with $\mathrm{pj}^{-1}\mathbf{Y}$ on $T_2$, the relation is Step (i).
\item[(2)] $S_y(A)=F_y(A)^{-1}$, $y\in Y_k$.
By Proposition~2.8\,(a) on the two tori, $F_y^{(1)}(A)$ is bijective on $\mathfrak{S}^{\mathrm{ch}}(\eta_k,\ell_1)$, and $F_y^{(2)}\left(\mathrm{pj}^{*}A\right)$ is bijective on $\mathfrak{S}^{\mathrm{ch}}(\eta_k,\ell_2)$; the relation follows from (1) and Step (iii).
In particular it holds for $S_k=S_{\infty}$, for $S_{k,y}=S_y$ ($y$ real) and for $S^{0}_{k+1}=S_0$.
\item[(3)] For the operators (Section~I.3.5)
\begin{gather*}
H_k(A)=b_k\cdot S_k(A)\,JQ_k^{T}(-A),\qquad\mathfrak{H}_k(A)=b_k\cdot Q_k(A)J^{*}S_k(A),\\
D_k(A)=b_k-b_k^2\cdot Q_k(A)J^{*}S_k(A)\,JQ_k^{T}(-A),\\
M_k(A)=b_k-D_k(A),\qquad\dot{M}_k(A)=b_k^2\cdot Q_k(A)J^{*}S_k(A)^2\,JQ_k^{T}(-A)
\end{gather*}
the relation follows from (1), (2) and Step (ii).
\item[(4)] $\bar{H}_k(A)=\mathsf{C}^{-1}H_k(-A)\,\mathsf{C}$.
By (3) at $-A$, $H_k^{(2)}\left(-\mathrm{pj}^{*}A\right)\mathrm{pj}^{*}=\mathrm{pj}^{*}H_k^{(1)}(-A)$, and $\mathsf{C}$ commutes with $\mathrm{pj}^{*}$.
Hence
\begin{align*}
\bar{H}_k^{(2)}\left(\mathrm{pj}^{*}A\right)\mathrm{pj}^{*}
&=\mathsf{C}^{-1}H_k^{(2)}\left(-\mathrm{pj}^{*}A\right)\mathsf{C}\,\mathrm{pj}^{*}
=\mathsf{C}^{-1}H_k^{(2)}\left(-\mathrm{pj}^{*}A\right)\mathrm{pj}^{*}\,\mathsf{C}\\
&=\mathrm{pj}^{*}\,\mathsf{C}^{-1}H_k^{(1)}(-A)\,\mathsf{C}
=\mathrm{pj}^{*}\bar{H}_k^{(1)}(A).
\end{align*}
\item[(5)] $\Gamma_{k,y}(A)=\left(D_k(A)+b_{L}\cdot P(A)+i\gamma_5\cdot y\right)^{-1}$, $y$ real.
The operator in the brackets satisfies the relation by (1), (3) and Step (ii).
By Proposition~4\,(b) on the two tori, it is bijective on $\mathfrak{S}(1,\ell_1)$ at $A$ and on $\mathfrak{S}(1,\ell_2)$ at $\mathrm{pj}^{*}A$; now Step (iii) applies.
Moreover $\Gamma_k=\Gamma_{k,0}$.
\end{itemize}

\paragraph{6.6. The list for $k=0$.}
Let $A\in\tilde{G}_0(T_1)$; by 6.1\,(1), $\mathrm{pj}^{*}A\in\tilde{G}_0(T_2)$.
The data $(0,\ell_i,\mu,e_k)$ are data (SF).
\begin{itemize}
\item[(1)] For $D_A$, $Q(A)J^{*}$, $JQ^{T}(-A)$, $JP(A)J^{*}$, $H_0=\bar{H}_0=T_0=\bar{T}_0=\dot{M}_0=1$, $D_0(A)=D_A+\mu$, $\operatorname{Pen}_{y}(A)$ and $F_y(A)$ ($y$ real), and $\delta_{\alpha,A}$, the relation is Step (i), together with Section 4.
\item[(2)] $S_{0,y}(A)=F_y(A)^{-1}$, $y$ real.
By Proposition~2.8\,(a) at $k=0$ on the two tori, $F_y^{(1)}(A)$ and $F_y^{(2)}\left(\mathrm{pj}^{*}A\right)$ are bijective on $\mathfrak{S}^{\mathrm{ch}}(1,\ell_1)$ and on $\mathfrak{S}^{\mathrm{ch}}(1,\ell_2)$, respectively; the relation follows from (1) and Step (iii).
Here $\Gamma_0(A)=S^{0}_1(A)=S_{0,0}(A)$.
\end{itemize}
This proves (b).
(For the numbers $\delta Z_k(A)$, which are determinants on the spaces $\mathfrak{S}(1,\ell_i)$, respectively $\mathfrak{S}^{\mathrm{ch}}(1,\ell_i)$, whose dimensions differ by the factor $\mathrm{Jr}^3$, no relation is asserted.)

\subsubsection*{7. Uniformity}

The operators $\mathrm{Op}_{\mathbb{L}}[O]\left(\tilde{B},\tilde{Y}\right)$ and the sets $\mathrm{Rd}(t)$ of Sections 1--4 are formed with $k$, $\eta_k$, $\mu$, $e_k$, $b_k$, $b_{L}$, $\alpha$, $\rho$, $w$, $L$, $M$ and the parameter $y$; the side $\ell$ does not enter them.
The ranges in (a) are the numbers $\eta_k$, $1/2$, $1$, $L/2$, $L$, $M$ for $k\ge1$ and $1$, $L/2$, $L$, $M$, $0$ for $k=0$.
They do not depend on $\ell$, $\mu$, $e_k$, $y$, $\tilde{B}$, $\tilde{Y}$, and only the first of them depends on $k$.
The properties (f1), (f2), (f4) hold for every complex $\tilde{B}$ and every $\tilde{Y}$, and (f3) holds for every $\ell\in\mathrm{Cl}$, every complex $B$ and every $Y$.
Clause (b) holds for all $\ell_1$ and $\ell_2=\mathrm{Jr}\cdot\ell_1$ in $\mathrm{Cl}$, every $k\ge0$, $\mu\in(0,1]$, $e_k\in(0,e_0]$, every $A\in\tilde{G}_k(T_1)$, every $y$ and every $\mathbf{Y}$; it contains no constant.
The cutoff $N$ and the charge $e$ do not occur: the data are scale-free (Section~I.1.8).

\subsubsection*{Conditions on the constants}

The table lists the conditions on the constants, the places in the proof where each of them enters, and the corresponding lines of Section~I.8.2 or of Proposition~14.2.

\noindent
\begin{tabular}{|p{2.5cm}|p{5.4cm}|p{2.1cm}|}
\hline
condition & where it is used & line of Section~I.8.2 / Proposition~14.2\\
\hline
$L$ odd & (Bq1)--(Bq3): the blocks are centred, partition $\mathbb{L}_k$ and are nested; 6.1\,(1): $11M/\eta_k$ is odd (with (M1)) & (C3); (L5)\\
\hline
$M$ a power of $L$, $M\ge L^2$ & 0.2: the $M$-cubes are blocks $\mathrm{Bl}_{k+\mathrm{iM}}$; $\Delta_z$ lies in one $M$-cube; $(L\cdot M)\mathbb{Z}^3$ lies in $(L^j\cdot\eta_k)\mathbb{Z}^3$ for $j\le k+\mathrm{iM}$; Section 1, remark on (cv11): $\eta_k\le1\le M$ & (C5); (M1)\\
\hline
$L\ge2$ & the equalities $\max\left(L-\eta_k,\eta_k\right)=L-\eta_k$ (list in 3.2) and $\max(L-1,1)=L-1$ (list in Section 4), where the ranges need only that both maxima are at most $L$; $\eta_k=L^{-k}\le1$ (Section 1, remark on (cv11); 3.2, after the list) & (L5)\\
\hline
$(L-1)/2\le M$ & $U_{\tilde{B}}(z',z)$ is a phase of (I.1.6), which is defined there for distances at most $M$ & follows from (M1)\\
\hline
$b>0$ and $L>1$ (they give $b_k\ge b_{-}>0$ and $b_{L}>0$: Sections~I.0.5, I.3.4) & (al4), and through it the inverses in (cv7), (cv8), (al5), 2.3\,(4); the factors $b_k^{-1}$ in (cv9)--(cv11), (al1) & (C0): $b=1$; (C3), (L5): $L\ge5$; (C4): $b_{L}$, $b_{-}$\\
\hline
the conditions of Propositions~2.8 and~4: (L5), (M1), (M2), (M5) ($M\ge108\pi\cdot C'_{V}\cdot a'$), (Z2.s1), (Z2.s2); $\alpha<1$ (through Proposition~2.7\,(d), (e)) & only through Proposition~2.8\,(a) (in 6.5\,(2) and 6.6\,(2)) and Proposition~4\,(b) (in 6.5\,(5)) & (C1), (C3), (C5), (C8)\\
\hline
\end{tabular}

\medskip
The following are properties of the data, not conditions on the constants: $\ell\in\mathrm{Cl}$ implies that $\ell$ is a multiple of $L\cdot M$ (2.3\,(1)) and that $11M<13\cdot L\cdot M\le\ell$ (6.1); $\mathrm{Jr}$ is odd (6.0).
Proposition~11 holds for all values of the constants that satisfy the conditions of the table.

\section*{Chapter 12. Integral representation of the logarithm of the one-step determinant, real fields: proof of Proposition 12}

\subsection*{0. Conventions of this chapter}

\emph{The first display.}
By \emph{the first display} we mean the display of Proposition~12 for $k\ge 1$.

\emph{Sites and cluster maps.}
As in Proposition~5, unit sites are denoted by $z$, with primes or indices, since the letter $y$ is reserved for the resolvent parameter;
as in Proposition~5\,(a), $y'$ denotes a site of $L\mathbb{Z}^3$.
The cluster maps $J$, $J^*$ are always written out (Section~I.3.3).

\emph{Further notation.}
$\mathrm{Mx}$ and $\mathrm{Mx}'$ denote $4\times 4$ matrices, and $P:=P(A)$, as in Proposition~4.
$R$ is the side of a block torus (Section~G1.2), and $\mathrm{cf}$ stands for a positive number.
The letter $K$ denotes the operator of Proposition~12, that is, $K=D_k(A)+b_L P(A)$ for $k\ge 1$ and $K=F_0(A)$ for $k=0$.
The letters $t$, $t'$ denote indices of the source or of the target space of an operator (the arguments of its kernel), and $v$, $v'$ denote spinors on the unit lattice.

\subsection*{Proof of Proposition 12}

\begin{proof}
Fix data $(k,\ell,\mu,e_k)$ as in (SF) (Sections~I.1.8, G1.1) and a real connection $A\in\tilde{G}_k$ (Section~I.5.3).
The side $\ell\in\mathrm{Cl}$ (Section~I.1.1) is an integer multiple of $L$, and $\ell\ge 13\,L\cdot M\ge 3L$.
Hence $T_k$ is the block torus $(0,\ell,-1)$ of Section~G1.2, and Proposition~2.2 applies on it with $R=\ell$.
For $k\ge 1$ the operators $S_k(A)=F_\infty(A)^{-1}$ and $S_{k,y}(A)=F_y(A)^{-1}$, $y\in\mathbb{R}$, exist by Proposition~2.8\,(a), and the inverse
\[
B_{k,y}(A)=\left(b_k+b_L P(A)+i\gamma_5 y\right)^{-1}
\]
exists by (6.2);
the proof of (6.2), given in 6.1, rests only on Proposition~2.2\,(f0), on $\gamma_5^2=1$ and on $b_k>0$, $b_L>0$.
Consequently $D_k(A)$, $K$ and the integrand of the first display are defined (Section~I.3.5).

\subsubsection*{Step 0. The space; bases, traces, kernels of products; the operator $\gamma_5$; inequalities for $4\times 4$ matrices}

By \emph{the space} we mean
\begin{align*}
k\ge 1:&\quad \mathfrak{S}(1,\ell),\ \text{of dimension}\ \mathrm{dsp}:=4n_k,\ \text{with}\ (v',v):=\sum_z\overline{v'(z)}\cdot v(z),\\
k=0:&\quad \mathfrak{S}^{\mathrm{ch}}(1,\ell),\ \text{of dimension}\ \mathrm{dsp}:=192\,n_0,\ \text{with}\ (f',f)^{\mathrm{ch}}=\frac{1}{48}\sum_c\sum_{g\in O_{\mathrm{cub}}}\overline{f'(c,g)}\cdot f(c,g).
\end{align*}
Here $z$, resp.\ $c$, ranges over a fundamental domain of $\ell\mathbb{Z}^3$ in $\mathbb{Z}^3$;
this domain has $n_k=\ell^3$ sites, resp.\ $n_0=\ell^3$ sites with $48$ labels over each
(Sections~I.1.5, I.1.6, and (I.1.3) at $\eta_0=1$).
The form $(v',v)=(v',v)_1=\langle\overline{v'},v\rangle_1$ is the sesquilinear pairing $(\cdot,\cdot)_a$ of Section~I.1.6 at $a=1$.
The two forms $(\cdot,\cdot)$ and $(\cdot,\cdot)^{\mathrm{ch}}$ are scalar products, that is, positive definite sesquilinear forms.
Adjoints, denoted by a star, are taken with respect to these forms.
For an operator on $\mathfrak{S}^{\mathrm{ch}}(\eta_k,\ell)$, or between $\mathfrak{S}^{\mathrm{ch}}(\eta_k,\ell)$ and $\mathfrak{S}(1,\ell)$ ($k\ge 1$), the adjoint refers to the form
\[
(f',f)^{\mathrm{ch}}=\eta_k^3\,\frac{1}{48}\sum_c\sum_{g\in O_{\mathrm{cub}}}\overline{f'(c,g)}\cdot f(c,g)
\]
on $\mathfrak{S}^{\mathrm{ch}}(\eta_k,\ell)$, where $c$ ranges over a fundamental domain of $\ell\mathbb{Z}^3$ in $\mathbb{L}_k$ (Section~I.1.6, (I.1.3)), and to the form $(\cdot,\cdot)$ on $\mathfrak{S}(1,\ell)$.
We write $\operatorname{Tr}$ and $\det$ for the trace and the determinant of linear maps from the space to itself;
both are independent of the basis used to compute them.

\emph{Levels $k\ge 1$.}
For a site $z$ of the fundamental domain and $\beta\in\{1,\dots,4\}$ let $\mathrm{bs}_{z,\beta}\in\mathfrak{S}(1,\ell)$ be the element whose component $\beta'$ at the site $z'$ of the fundamental domain equals $\mathbf{1}[z'=z,\ \beta'=\beta]$.
These $4n_k$ functions form an orthonormal basis.
Since $(Ov)(t)=\sum_{z'}O(t,z')\,v(z')$ (Section~I.1.6), one has
\[
O\,\mathrm{bs}_{z',\beta'}=\sum_{z,\beta}O(z,z')_{\beta\beta'}\,\mathrm{bs}_{z,\beta}.
\]
Thus the matrix of $O$ in this basis has the entries $O(z,z')_{\beta\beta'}$, and
\begin{equation}
\operatorname{Tr}O=\sum_z\operatorname{tr}O(z,z),\qquad (O_1O_2)(z,z'')=\sum_{z'}O_1(z,z')\,O_2(z',z'').
\tag{0.1}
\end{equation}

\emph{Level $0$.}
For a fermion site $x=(c,g)$ with $c$ in the fundamental domain and $g\in O_{\mathrm{cub}}$, and for $\beta\in\{1,\dots,4\}$, let $\mathrm{bs}_{x,\beta}\in\mathfrak{S}^{\mathrm{ch}}(1,\ell)$ be the element whose component $\beta'$ at the site $x'=(c',g')$, $c'$ in the fundamental domain, equals $\mathbf{1}[x'=x,\ \beta'=\beta]$.
There are $48\cdot 4\cdot n_0=192\,n_0$ such functions, and
\[
\left(\mathrm{bs}_{x,\beta},\mathrm{bs}_{x',\beta'}\right)^{\mathrm{ch}}=\frac{1}{48}\,\mathbf{1}[x=x',\ \beta=\beta'];
\]
they form an orthogonal basis of vectors of length $48^{-\frac{1}{2}}$.
By (I.1.4) and (I.1.3) at $\eta_0=1$,
\[
(Of)(x)=\frac{1}{48}\sum_{c',g'}O\left(x,(c',g')\right)f(c',g'),
\qquad\text{hence}\qquad
O\,\mathrm{bs}_{x',\beta'}=\sum_{x,\beta}\frac{1}{48}\,O(x,x')_{\beta\beta'}\,\mathrm{bs}_{x,\beta}.
\]
Thus the matrix of $O$ has the entries $\frac{1}{48}O(x,x')_{\beta\beta'}$, and
\begin{equation}
\begin{gathered}
\operatorname{Tr}O=\frac{1}{48}\sum_c\sum_{g\in O_{\mathrm{cub}}}\operatorname{tr}O\left((c,g),(c,g)\right)=\int\mathrm{d}x\,\operatorname{tr}O(x,x),\\
(O_1O_2)(x,x'')=\int\mathrm{d}x'\,O_1(x,x')\,O_2(x',x''),\qquad \int\mathrm{d}x=\frac{1}{48}\cdot 48\cdot n_0=n_0;
\end{gathered}
\tag{0.2}
\end{equation}
the first identity is (I.ke2).

Every kernel changes its sign when one of its arguments is shifted by $\ell e_\mu$ ((I.1.4), Section~G1.2);
hence the summands in (0.1) and (0.2) are $\ell\mathbb{Z}^3$-periodic, and these sums do not depend on the fundamental domain chosen.

\emph{The operator $\gamma_5$.}
On each of the spaces $\mathfrak{S}(1,\ell)$, $\mathfrak{S}^{\mathrm{ch}}(\eta_k,\ell)$ the operator $\gamma_5$ is defined by $(\gamma_5 f)(t):=\gamma_5 f(t)$;
the kernel of $\gamma_5 O$ is $\gamma_5 O(t,t')$.
The matrix $\gamma_5$ is Hermitian and satisfies $\gamma_5^2=1$ (I.c1), and the four components at one site enter the forms with the same weight.
Hence
\begin{equation}
\gamma_5^*=\gamma_5,\qquad \gamma_5^2=1\qquad\text{as operators; in particular } \gamma_5 \text{ is unitary}.
\tag{0.3}
\end{equation}
The matrix
\[
\gamma_5=\left(\begin{array}{cc}0&\mathbf{1}_2\\ \mathbf{1}_2&0\end{array}\right)
\]
(Section~I.2.1) is the permutation matrix of $(1\,3)(2\,4)$, which is a product of two transpositions;
thus $\det\gamma_5=1$ and $\operatorname{tr}\gamma_5=0$.
With respect to the bases above, the operator $\gamma_5$ is block diagonal:
it consists of $n_k$ diagonal blocks equal to $\gamma_5$ if $k\ge 1$, and of $48n_0$ such blocks, one for each fermion site, if $k=0$.
It follows that
\begin{equation}
\det\left(\gamma_5\ \text{on}\ \mathfrak{S}(1,\ell)\right)=1^{n_k}=1\quad(k\ge 1),\qquad
\det\left(\gamma_5\ \text{on}\ \mathfrak{S}^{\mathrm{ch}}(1,\ell)\right)=1^{48n_0}=1\quad(k=0).
\tag{0.4}
\end{equation}

\emph{Matrices.}
For a $4\times 4$ matrix, $|\cdot|_2$ denotes its operator norm with respect to the Euclidean norm of $\mathbb{C}^4$ (Section~I.1.5).
Each entry of a $4\times 4$ matrix is a component of the image of a unit vector, so its modulus does not exceed the norm $|\cdot|_2$ of the matrix;
moreover the matrix $\gamma_5$ is unitary.
Consequently, for $4\times 4$ matrices $\mathrm{Mx}$, $\mathrm{Mx}'$,
\begin{equation}
|\operatorname{tr}\mathrm{Mx}|\le 4\,|\mathrm{Mx}|_2,\qquad
|\operatorname{tr}(\mathrm{Mx}\,\mathrm{Mx}')|=\left|\sum_{\beta,\beta'}\mathrm{Mx}_{\beta\beta'}\,\mathrm{Mx}'_{\beta'\beta}\right|\le\left(\sum_{\beta,\beta'}|\mathrm{Mx}_{\beta\beta'}|\right)|\mathrm{Mx}'|_2,\qquad
|\gamma_5|_2=1.
\tag{0.5}
\end{equation}

\subsubsection*{Step 1. Two integrals}

Recall that $\arctan$ is a bijection of $\mathbb{R}$ onto $\left(-\frac{\pi}{2},\frac{\pi}{2}\right)$.
For a real number $\mathrm{ev}\ne 0$ and for $y\ge 0$ define
\begin{equation}
\begin{gathered}
\operatorname{Log}(\mathrm{ev}+iy):=\frac{1}{2}\log\left(\mathrm{ev}^2+y^2\right)+i\left[\arctan\frac{y}{\mathrm{ev}}+\pi\,\mathbf{1}[\mathrm{ev}<0]\right],\\
\operatorname{Log}\mathrm{ev}:=\operatorname{Log}(\mathrm{ev}+i0)=\log|\mathrm{ev}|+i\pi\,\mathbf{1}[\mathrm{ev}<0].
\end{gathered}
\tag{1.1}
\end{equation}
This function has the following three properties.
\begin{itemize}
\item[(a)] $\exp(\operatorname{Log}\mathrm{ev})=\mathrm{ev}$.
Indeed, for $\mathrm{ev}>0$ one has $\exp(\log\mathrm{ev})=\mathrm{ev}$, and for $\mathrm{ev}<0$ one has $\exp\left(\log|\mathrm{ev}|+i\pi\right)=-|\mathrm{ev}|=\mathrm{ev}$.
\item[(b)] The map $y\mapsto\operatorname{Log}(\mathrm{ev}+iy)$ is continuously differentiable on $[0,\infty)$, and
\[
\frac{\mathrm{d}}{\mathrm{d}y}\operatorname{Log}(\mathrm{ev}+iy)=\frac{i}{\mathrm{ev}+iy}.
\]
Indeed,
\[
\frac{\mathrm{d}}{\mathrm{d}y}\,\frac{1}{2}\log\left(\mathrm{ev}^2+y^2\right)=\frac{y}{\mathrm{ev}^2+y^2},\qquad
\frac{\mathrm{d}}{\mathrm{d}y}\arctan\frac{y}{\mathrm{ev}}=\frac{1}{\mathrm{ev}}\left(1+\frac{y^2}{\mathrm{ev}^2}\right)^{-1}=\frac{\mathrm{ev}}{\mathrm{ev}^2+y^2},
\]
and
\[
\frac{y+i\,\mathrm{ev}}{\mathrm{ev}^2+y^2}=\frac{i\,(\mathrm{ev}-iy)}{(\mathrm{ev}+iy)(\mathrm{ev}-iy)}=\frac{i}{\mathrm{ev}+iy}.
\]
\item[(c)] $\operatorname{Log}(\mathrm{ev}+iy)-\log y-\frac{i\pi}{2}\to 0$ as $y\to\infty$.
Indeed, $\frac{1}{2}\log\left(\mathrm{ev}^2+y^2\right)-\log y=\frac{1}{2}\log\left(1+\frac{\mathrm{ev}^2}{y^2}\right)\to 0$;
for $\mathrm{ev}>0$ one has $\arctan\frac{y}{\mathrm{ev}}\to\frac{\pi}{2}$, and for $\mathrm{ev}<0$ one has $\arctan\frac{y}{\mathrm{ev}}+\pi\to-\frac{\pi}{2}+\pi=\frac{\pi}{2}$.
\end{itemize}
From (b) and the fundamental theorem of calculus one obtains, for every $\mathrm{yc}\ge 0$,
\begin{equation}
-i\int_0^{\mathrm{yc}}\mathrm{d}y\,\frac{1}{\mathrm{ev}+iy}=\operatorname{Log}\mathrm{ev}-\operatorname{Log}(\mathrm{ev}+i\,\mathrm{yc}).
\tag{i}
\end{equation}
Furthermore
\begin{equation}
\int_1^{\infty}\frac{\mathrm{d}y}{y}\,\frac{\mathrm{ev}}{\mathrm{ev}+iy}-i\int_0^1\mathrm{d}y\,\frac{1}{\mathrm{ev}+iy}=\operatorname{Log}\mathrm{ev}-\frac{i\pi}{2},
\tag{ii}
\end{equation}
where the first integral converges absolutely, since
\[
\left|\frac{\mathrm{ev}}{y\,(\mathrm{ev}+iy)}\right|\le\frac{|\mathrm{ev}|}{y^2}\qquad\text{for } y\ge 1.
\]
To prove (ii), note that
\[
\frac{\mathrm{ev}}{y\,(\mathrm{ev}+iy)}=\frac{(\mathrm{ev}+iy)-iy}{y\,(\mathrm{ev}+iy)}=\frac{1}{y}-\frac{i}{\mathrm{ev}+iy}.
\]
By (b), for $\mathrm{yc}\ge 1$,
\[
\int_1^{\mathrm{yc}}\frac{\mathrm{d}y}{y}\,\frac{\mathrm{ev}}{\mathrm{ev}+iy}=\log\mathrm{yc}-\operatorname{Log}(\mathrm{ev}+i\,\mathrm{yc})+\operatorname{Log}(\mathrm{ev}+i),
\]
and by (c) the right side tends to $-\frac{i\pi}{2}+\operatorname{Log}(\mathrm{ev}+i)$ as $\mathrm{yc}\to\infty$.
By (i) at $\mathrm{yc}=1$,
\[
-i\int_0^1\mathrm{d}y\,\frac{1}{\mathrm{ev}+iy}=\operatorname{Log}\mathrm{ev}-\operatorname{Log}(\mathrm{ev}+i).
\]
Adding the last two results gives (ii).

\subsubsection*{Step 2. $K^*=\gamma_5K\gamma_5$}

\emph{Level $0$.}
Here $K=F_0(A)=D_A+\mu+b_L JP(A)J^*$ on $\mathfrak{S}^{\mathrm{ch}}(1,\ell)$ (Section~I.3.9).
Proposition~2.2\,(e), taken at $k=0$ and $y=0$, states that $F_0(A)^*=\gamma_5F_0(A)\gamma_5$.
(This can be checked term by term:
$D_A^*=\gamma_5D_A\gamma_5$ by Proposition~2.2\,(e);
$\mu$ and $b_L$ are real;
$JP(A)J^*$ is self-adjoint and commutes with $\gamma_5$, being an orthogonal projection (Proposition~2.2\,(e)) with a kernel that is a multiple of $\mathbf{1}_4$ (Proposition~2.2\,(f0));
and $\gamma_5^2=1$.)

\emph{Levels $k\ge 1$.}
Here $K=D_k(A)+b_LP(A)$ on $\mathfrak{S}(1,\ell)$, with
\[
D_k(A)=b_k-b_k^2\,Q_k(A)J^*S_k(A)JQ_k^T(-A),\qquad S_k(A)=F_\infty(A)^{-1},\qquad F_\infty(A)=D_A+\mu+b_kJP_k(A)J^*
\]
(Section~I.3.5).
\begin{itemize}
\item[(1)] Proposition~2.2\,(e) gives $D_A^*=\gamma_5D_A\gamma_5$ and shows that $JP_k(A)J^*$ is an orthogonal projection;
since its kernel is a multiple of $\mathbf{1}_4$ (Proposition~2.2\,(f0)), this projection commutes with $\gamma_5$.
The numbers $\mu$ and $b_k$ are real.
Together with (0.3) this gives
\[
F_\infty(A)^*=\gamma_5D_A\gamma_5+\mu+b_kJP_k(A)J^*=\gamma_5F_\infty(A)\gamma_5,
\]
and, taking inverses, $S_k(A)^*=\left(F_\infty(A)^*\right)^{-1}=\gamma_5S_k(A)\gamma_5$.
\item[(2)] By Proposition~2.2\,(e), $\left(Q_k(A)J^*\right)^*=JQ_k^T(-A)$, and therefore also $\left(JQ_k^T(-A)\right)^*=Q_k(A)J^*$.
This can also be seen directly, in the forms of Step~0 and with the two displays of (3) below:
\[
\left(v,Q_k(A)J^*f\right)=\sum_z\overline{v(z)}\cdot\int_{\Delta_z}\mathrm{d}x\,\varphi_A(z,c)f(x)=\sum_z\int_{\Delta_z}\mathrm{d}x\,\overline{\varphi_A(z,c)^{-1}v(z)}\cdot f(x)=\left(JQ_k^T(-A)v,f\right)^{\mathrm{ch}},
\]
because $\overline{\varphi_A(z,c)}=\varphi_A(z,c)^{-1}$ for real $A$ ((I.3.3)) and the sets $\Delta_z\times O_{\mathrm{cub}}$ partition the fermion sites (Section~I.1.4).
\item[(3)] By the displays of Section~I.3.3,
\[
\left(Q_k(A)J^*f\right)(z)=\int_{\Delta_z}\mathrm{d}x\,\varphi_A(z,c)f(x)=\eta_k^3\,\frac{1}{48}\sum_{c\in\Delta_z}\sum_{g\in O_{\mathrm{cub}}}\varphi_A(z,c)f(c,g)
\]
and $\left(JQ_k^T(-A)v\right)(x)=\varphi_A(z,c)^{-1}v(z)$ for $x=(c,g)$ over $\Delta_z$, where the $\varphi_A(z,c)$ are numbers ((I.3.3)).
Hence
\[
\gamma_5\,Q_k(A)J^*=Q_k(A)J^*\,\gamma_5,\qquad \gamma_5\,JQ_k^T(-A)=JQ_k^T(-A)\,\gamma_5,
\]
where $\gamma_5$ acts on $\mathfrak{S}(1,\ell)$ on one side and on $\mathfrak{S}^{\mathrm{ch}}(\eta_k,\ell)$ on the other.
\item[(4)] From (1)--(3),
\begin{align*}
D_k(A)^*&=b_k-b_k^2\left(JQ_k^T(-A)\right)^*S_k(A)^*\left(Q_k(A)J^*\right)^*
=b_k-b_k^2\,Q_k(A)J^*\gamma_5S_k(A)\gamma_5JQ_k^T(-A)\\
&=\gamma_5\left[b_k-b_k^2\,Q_k(A)J^*S_k(A)JQ_k^T(-A)\right]\gamma_5=\gamma_5D_k(A)\gamma_5.
\end{align*}
\item[(5)] Being an orthogonal projection (Proposition~2.2\,(e)) with a kernel that is a multiple of $\mathbf{1}_4$ (Proposition~2.2\,(f0)), $P(A)$ satisfies $P(A)^*=P(A)=\gamma_5P(A)\gamma_5$;
moreover $b_L$ is real.
\end{itemize}
From (4) and (5) one obtains $K^*=\gamma_5K\gamma_5$.

\subsubsection*{Step 3. $\mathrm{Kh}:=K\gamma_5$ is Hermitian and invertible; $\det\mathrm{Kh}=\det K$}

Put $\mathrm{Kh}:=K\gamma_5$.
By (0.3) and Step~2,
\[
\mathrm{Kh}^*=\gamma_5^*K^*=\gamma_5\left(\gamma_5K\gamma_5\right)=K\gamma_5=\mathrm{Kh}.
\]
Thus $\mathrm{Kh}$ is Hermitian with respect to the form of Step~0, that is, with respect to $(\cdot,\cdot)$ on $\mathfrak{S}(1,\ell)$ if $k\ge 1$ and with respect to $(\cdot,\cdot)^{\mathrm{ch}}$ on $\mathfrak{S}^{\mathrm{ch}}(1,\ell)$ if $k=0$.

The operator $K$ is invertible:
for $k\ge 1$ by Proposition~4\,(d), according to which $\Gamma_k(A)=\left(D_k(A)+b_LP(A)\right)^{-1}$ exists,
and for $k=0$ by Proposition~2.8\,(d), according to which $F_0(A)$ is invertible.
By (0.3) the operator $\gamma_5$ is invertible.
Hence $\mathrm{Kh}$ is invertible.

Moreover $\det\mathrm{Kh}=\det K\cdot\det\gamma_5=\det K$ by (0.4),
and $\det K=\delta Z_k(A)$ for every $k\ge 0$ by the definitions in Sections~I.3.5 and I.3.9.

\subsubsection*{Step 4. Spectral representation}

An orthonormal basis of the space is given by the vectors $\mathrm{bs}_{z,\beta}$ if $k\ge 1$, and by the vectors $\sqrt{48}\,\mathrm{bs}_{x,\beta}$ if $k=0$ (recall that $\mathrm{bs}_{x,\beta}$ has length $48^{-\frac{1}{2}}$).
In this basis the matrix of $\mathrm{Kh}$ has the entries $\mathrm{Kh}(z,z')_{\beta\beta'}$, resp.\ $\frac{1}{48}\mathrm{Kh}(x,x')_{\beta\beta'}$
(Step~0; multiplying all basis vectors by one common factor leaves the matrix of a map unchanged),
and this matrix is Hermitian, since $\mathrm{Kh}^*=\mathrm{Kh}$ (Step~3).
The spectral theorem for Hermitian matrices provides an orthonormal basis of the space consisting of eigenvectors of $\mathrm{Kh}$;
the corresponding eigenvalues $\mathrm{ev}_1,\dots,\mathrm{ev}_{\mathrm{dsp}}$, listed with multiplicity, are real,
where $\mathrm{dsp}=4n_k$ for $k\ge 1$ and $\mathrm{dsp}=192\,n_0$ for $k=0$.
Since $\mathrm{Kh}$ is invertible (Step~3), $\mathrm{ev}_j\ne 0$ for every $j$.
In this eigenbasis, for each real $y$, the operator $\mathrm{Kh}+iy$ is the diagonal matrix with the entries $\mathrm{ev}_j+iy$, none of which vanishes.
Since $\gamma_5^2=1$,
\begin{equation}
\begin{gathered}
K+i\gamma_5y=\left(K\gamma_5+iy\right)\gamma_5=\left(\mathrm{Kh}+iy\right)\gamma_5;
\qquad\text{so } K+i\gamma_5y \text{ is invertible for every real } y,\\
\left(K+i\gamma_5y\right)^{-1}=\gamma_5\left(\mathrm{Kh}+iy\right)^{-1},\qquad
\left(\mathrm{Kh}+iy\right)^{-1}=\gamma_5\left(K+i\gamma_5y\right)^{-1}.
\end{gathered}
\tag{4.1}
\end{equation}
Consequently
\begin{align}
\operatorname{Tr}\left[\gamma_5\left(K+i\gamma_5y\right)^{-1}\right]&=\operatorname{Tr}\left[\left(\mathrm{Kh}+iy\right)^{-1}\right]=\sum_{j=1}^{\mathrm{dsp}}\frac{1}{\mathrm{ev}_j+iy},
\tag{4.2}\\
\operatorname{Tr}\left[K\left(K+i\gamma_5y\right)^{-1}\right]&=\operatorname{Tr}\left[K\gamma_5\gamma_5\left(K+i\gamma_5y\right)^{-1}\right]=\operatorname{Tr}\left[\mathrm{Kh}\left(\mathrm{Kh}+iy\right)^{-1}\right]=\sum_{j=1}^{\mathrm{dsp}}\frac{\mathrm{ev}_j}{\mathrm{ev}_j+iy},
\tag{4.3}\\
\delta Z_k(A)&=\det K=\det\mathrm{Kh}=\prod_{j=1}^{\mathrm{dsp}}\mathrm{ev}_j=\exp\left(\sum_{j=1}^{\mathrm{dsp}}\operatorname{Log}\mathrm{ev}_j\right),
\tag{4.4}
\end{align}
the last line by Step~3 and Step~1\,(a).
Both right sides in (4.2) and (4.3) depend continuously on $y\in\mathbb{R}$.

\subsubsection*{Step 5. Level 0}

\paragraph{5.1. The display.}
By Section~I.3.9,
\[
F_y(A)=D_A+\mu+b_LJP(A)J^*+i\gamma_5y=K+i\gamma_5y .
\]
Hence, by (4.1), $S_{0,y}(A)=F_y(A)^{-1}=\left(K+i\gamma_5y\right)^{-1}$ for every real $y$, and (4.2), (4.3) become
\begin{equation}
\operatorname{Tr}\left[\gamma_5S_{0,y}(A)\right]=\sum_{j=1}^{192n_0}\frac{1}{\mathrm{ev}_j+iy},\qquad
\operatorname{Tr}\left[F_0(A)S_{0,y}(A)\right]=\sum_{j=1}^{192n_0}\frac{\mathrm{ev}_j}{\mathrm{ev}_j+iy}.
\tag{5.1}
\end{equation}
Now take $\mathrm{ev}=\mathrm{ev}_j$ in (ii) of Step~1 and sum over $j=1,\dots,192\,n_0$.
Only finitely many terms occur, and for each $j$ the integrands are continuous and satisfy
\[
\left|\frac{\mathrm{ev}_j}{y\,(\mathrm{ev}_j+iy)}\right|\le\frac{|\mathrm{ev}_j|}{y^2}\quad\text{on } [1,\infty),\qquad
\left|\frac{1}{\mathrm{ev}_j+iy}\right|\le\frac{1}{|\mathrm{ev}_j|}\quad\text{on } [0,1].
\]
Hence both integrals of the traces converge absolutely, and
\begin{equation}
\begin{split}
\int_1^{\infty}\frac{\mathrm{d}y}{y}\,\operatorname{Tr}\left[F_0(A)S_{0,y}(A)\right]-i\int_0^1\mathrm{d}y\,\operatorname{Tr}\left[\gamma_5S_{0,y}(A)\right]
&=\sum_{j=1}^{192n_0}\operatorname{Log}\mathrm{ev}_j-192\,n_0\cdot\frac{i\pi}{2}\\
&=\sum_j\operatorname{Log}\mathrm{ev}_j-96\pi i\,n_0.
\end{split}
\tag{5.2}
\end{equation}
As $n_0$ is an integer, $\exp(-96\pi i\,n_0)=1$;
taking exponentials in (5.2) and using (4.4), one obtains the display of Proposition~12 for $k=0$.

\paragraph{5.2. Bounds.}
Let $\mathrm{rw}_0:=59+b\cdot\mathrm{e}$.
Since $A\in\tilde{G}_0$ is real, $\mathfrak{i}(A)=e_k\cdot|\operatorname{Im}A|_\infty=0$ ((I.1.6)), and (I.L0.12) reads
\[
\sup_x\sup_\beta\sum_{\beta'}\int\mathrm{d}x'\,\left|F_0(A)(x,x')_{\beta\beta'}\right|\le 58+1+b\cdot\mathrm{e}=\mathrm{rw}_0 .
\]
For real $A$ this bound can also be proved directly, as follows.
Write $x=(c,g)$ (Section~I.1.3), and choose $y'\in L\mathbb{Z}^3$ such that $c$ lies in the $L$-block $B(y')$ (Section~I.1.1; at level $0$ the unit sites and the photon sites coincide).
By (I.2.3), (I.m1), (I.3.1) and the displays of Section~I.3.3,
\[
\begin{split}
\left(F_0(A)f\right)(x)={}&\left(\frac{\rho}{2}\cdot\mathrm{sw}+\mu\right)f(x)-\sum_{i=0}^{3}M_i(x)\,U_A\left(x,x^{(i)}\right)f\left(x^{(i)}\right)\\
&+b_L\sum_{c'\in B(y')}L^{-3}\,U_A(y',c)^{-1}\,U_A(y',c')\,\frac{1}{48}\sum_{g'\in O_{\mathrm{cub}}}f(c',g').
\end{split}
\]
According to (I.1.4), $\frac{1}{48}F_0(A)(x,x')_{\beta\beta'}$ is the coefficient with which $f(x')_{\beta'}$ enters the component $\beta$;
consequently $\sum_{\beta'}\int\mathrm{d}x'\,\left|F_0(A)(x,x')_{\beta\beta'}\right|$ does not exceed the sum of the moduli of the coefficients of the displayed expression in the row $\beta$
(a site $x^{(i)}$ or $(c',g')$ outside the fundamental domain is brought back by a translation in $\ell\mathbb{Z}^3$, at the cost of a sign).
Since $A$ is real, every phase $U_A$ has modulus one ((I.1.6)).
By (I.m7) the $l^1$ norm of every row is at most $3(1+\rho)$ for $M_0(x)$ and $M_3(x)$, and at most $3\sqrt{2}\,(\rho+\sqrt{2})$ for $M_1(x)$ and $M_2(x)$.
The contribution of the last term is $b_L\cdot L^3\cdot L^{-3}\cdot\frac{1}{48}\cdot 48=b_L$.
The total is
\[
\frac{\rho}{2}\cdot\mathrm{sw}+\mu+6(1+\rho)+6\sqrt{2}\,(\rho+\sqrt{2})+b_L=18+\rho\cdot\mathrm{sw}+\mu+b_L\le 18+29+1+b<\mathrm{rw}_0,
\]
because $\sqrt{2}<\frac{17}{12}$ (as $289>288$), so that $\mathrm{sw}=12+12\sqrt{2}<29$, and because $\rho<1$, $\mu\le 1$, $b_L\le b$.

Proposition~2.8\,(d) gives $\left|S_{0,y}(A)(x',x)\right|_2\le 32\,a'\,(1+y)^{-1}$ for all fermion sites $x$, $x'$ and every $y\ge 0$.
With (0.2) and (0.5) one obtains, for all $y\ge 0$,
\begin{align}
\left|\operatorname{Tr}\left[F_0(A)S_{0,y}(A)\right]\right|
&=\left|\int\mathrm{d}x\int\mathrm{d}x'\,\operatorname{tr}\left[F_0(A)(x,x')\,S_{0,y}(A)(x',x)\right]\right|\notag\\
&\le\int\mathrm{d}x\,\left\{\sum_\beta\sum_{\beta'}\int\mathrm{d}x'\,\left|F_0(A)(x,x')_{\beta\beta'}\right|\right\}\sup_{x'}\left|S_{0,y}(A)(x',x)\right|_2\notag\\
&\le n_0\cdot 4\,\mathrm{rw}_0\cdot 32\,a'\,(1+y)^{-1}=128\,\mathrm{rw}_0\cdot a'\cdot n_0\,(1+y)^{-1},
\tag{5.3}\\
\left|\operatorname{Tr}\left[\gamma_5S_{0,y}(A)\right]\right|
&=\left|\int\mathrm{d}x\,\operatorname{tr}\left[\gamma_5S_{0,y}(A)(x,x)\right]\right|
\le\int\mathrm{d}x\,4\,|\gamma_5|_2\,\left|S_{0,y}(A)(x,x)\right|_2\le 128\,a'\cdot n_0\,(1+y)^{-1}.
\tag{5.4}
\end{align}
In particular, (5.3) proves the assertion $\left|\operatorname{Tr}\left[F_0(A)S_{0,y}(A)\right]\right|=O\left((1+y)^{-1}\right)$ of Proposition~12.
Note that in $\int\mathrm{d}x=n_0$ the $48$ sites of each cluster compensate the weight $\frac{1}{48}$ of the measure;
the kernels are taken as in (I.1.4), and the factor $48$ coming from the inverse weight is part of the constant $a'$, which contains it through $C_{00}$.
Since
\[
\int_1^{\infty}\frac{\mathrm{d}y}{y\,(1+y)}=\left[\log\frac{y}{1+y}\right]_1^{\infty}=\log 2
\qquad\text{and}\qquad
\int_0^1\frac{\mathrm{d}y}{1+y}=\log 2,
\]
it follows that
\begin{equation}
\int_1^{\infty}\frac{\mathrm{d}y}{y}\,\left|\operatorname{Tr}\left[F_0(A)S_{0,y}(A)\right]\right|+\int_0^1\mathrm{d}y\,\left|\operatorname{Tr}\left[\gamma_5S_{0,y}(A)\right]\right|\le 128\,(\mathrm{rw}_0+1)\cdot a'\cdot n_0\log 2.
\tag{5.5}
\end{equation}

\subsubsection*{Step 6. Levels $k\ge 1$}

By Section~I.3.5,
\[
\Gamma_{k,y}(A)=\left(D_k(A)+b_LP(A)+i\gamma_5y\right)^{-1}=\left(K+i\gamma_5y\right)^{-1}.
\]
Write $P:=P(A)$, and let
\[
\mathrm{Ou}_y:=Q_k(A)J^*S_{k,y}(A)JQ_k^T(-A),
\]
an operator on $\mathfrak{S}(1,\ell)$.
Throughout this step $y$ denotes a real number with $y\ge 0$.

\paragraph{6.1. $B_{k,y}(A)$ and the trace of $\gamma_5B_{k,y}(A)$.}
For a number $\mathrm{cf}>0$ one has $(\mathrm{cf}+i\gamma_5y)(\mathrm{cf}-i\gamma_5y)=\mathrm{cf}^2+y^2\gamma_5^2=\mathrm{cf}^2+y^2$.
Hence $\mathrm{cf}+i\gamma_5y$ is invertible, and
\begin{equation}
\left(\mathrm{cf}+i\gamma_5y\right)^{-1}=\frac{\mathrm{cf}-i\gamma_5y}{\mathrm{cf}^2+y^2},\qquad
\gamma_5\left(\mathrm{cf}+i\gamma_5y\right)^{-1}=\frac{\mathrm{cf}\cdot\gamma_5-iy}{\mathrm{cf}^2+y^2},\qquad
\operatorname{tr}\left[\gamma_5\left(\mathrm{cf}+i\gamma_5y\right)^{-1}\right]=-\frac{4iy}{\mathrm{cf}^2+y^2}
\tag{6.1}
\end{equation}
($\operatorname{tr}\gamma_5=0$, $\operatorname{tr}\mathbf{1}_4=4$).
Since $P$ is idempotent with a kernel that is a multiple of $\mathbf{1}_4$ (Proposition~2.2\,(f0)), it commutes with $\gamma_5$, and
\[
\left(b_k+b_LP+i\gamma_5y\right)(1-P)=\left(b_k+i\gamma_5y\right)(1-P),\qquad
\left(b_k+b_LP+i\gamma_5y\right)P=\left(b_k+b_L+i\gamma_5y\right)P.
\]
Therefore
\[
\left(b_k+b_LP+i\gamma_5y\right)\left[\left(b_k+i\gamma_5y\right)^{-1}(1-P)+\left(b_k+b_L+i\gamma_5y\right)^{-1}P\right]=(1-P)+P=1.
\]
In finite dimension a right inverse is a two-sided inverse;
hence
\begin{equation}
B_{k,y}(A)=\left(b_k+i\gamma_5y\right)^{-1}(1-P)+\left(b_k+b_L+i\gamma_5y\right)^{-1}P
\qquad\left(b_k>0,\ b_L>0:\ \text{Section~I.3.4}\right).
\tag{6.2}
\end{equation}
The kernel of $P$ has the form $P(z,z')=\mathrm{pk}(z,z')\,\mathbf{1}_4$, where the $\mathrm{pk}(z,z')$ are numbers.
By (0.1) and by Proposition~2.2\,(f0) with $R=\ell$,
\[
4\sum_z\mathrm{pk}(z,z)=\operatorname{Tr}P=4\left(\frac{\ell}{L}\right)^3=4n_k\cdot L^{-3}.
\]
Hence, for every constant $4\times 4$ matrix $\mathrm{Mx}$,
\[
\operatorname{Tr}\left[\mathrm{Mx}\,P\right]=\sum_z\mathrm{pk}(z,z)\operatorname{tr}\mathrm{Mx}=n_k\cdot L^{-3}\operatorname{tr}\mathrm{Mx},\qquad
\operatorname{Tr}\left[\mathrm{Mx}\,(1-P)\right]=n_k\left(1-L^{-3}\right)\operatorname{tr}\mathrm{Mx}.
\]
Combined with (6.2) and (6.1) this yields
\begin{equation}
\operatorname{Tr}\left[\gamma_5B_{k,y}(A)\right]=-4iy\cdot n_k\left[\frac{1-L^{-3}}{b_k^2+y^2}+\frac{L^{-3}}{(b_k+b_L)^2+y^2}\right].
\tag{6.3}
\end{equation}
For $\mathrm{cf}>0$ and $\mathrm{yc}>0$,
\[
\int_0^{\mathrm{yc}}\mathrm{d}y\,\frac{y}{\mathrm{cf}^2+y^2}=\frac{1}{2}\log\left(1+\frac{\mathrm{yc}^2}{\mathrm{cf}^2}\right)=\log\mathrm{yc}-\log\mathrm{cf}+\frac{1}{2}\log\left(1+\frac{\mathrm{cf}^2}{\mathrm{yc}^2}\right).
\]
Using $(-i)(-4i)=-4$ and the definition of $c^{\mathrm{d}}_k$ (Section~I.3.4), one finds
\begin{equation}
\begin{split}
-i\int_0^{\mathrm{yc}}\mathrm{d}y\,\operatorname{Tr}\left[\gamma_5B_{k,y}(A)\right]
&=-4n_k\log\mathrm{yc}+4n_k\left[\left(1-L^{-3}\right)\log b_k+L^{-3}\log(b_k+b_L)\right]-\mathrm{rem}_1(\mathrm{yc})\\
&=-4n_k\log\mathrm{yc}+c^{\mathrm{d}}_k\cdot n_k-\mathrm{rem}_1(\mathrm{yc}),
\end{split}
\tag{6.4}
\end{equation}
where
\[
\mathrm{rem}_1(\mathrm{yc}):=2n_k\left(1-L^{-3}\right)\log\left(1+\frac{b_k^2}{\mathrm{yc}^2}\right)+2n_k\cdot L^{-3}\log\left(1+\frac{(b_k+b_L)^2}{\mathrm{yc}^2}\right)\to 0\qquad(\mathrm{yc}\to\infty).
\]

\paragraph{6.2. The trace of the first display.}
Let $z_1$, $z_2$ be unit sites.
Proposition~5\,(c), applied with $O=S_{k,y}(A)$, together with Proposition~2.8\,(c) and $w_k(y)=1$ (Section~I.3.5), gives
\[
\left|\mathrm{Ou}_y(z_1,z_2)\right|_2\le\mathrm{e}\,\left\|\mathbf{1}_{\Delta_{z_1}}S_{k,y}(A)\mathbf{1}_{\Delta_{z_2}}\right\|_{\mathrm{su}}\le 32\,\mathrm{e}\cdot a'\exp\left(-\gamma^Fd(z_1,z_2)\right)\le 32\,\mathrm{e}\cdot a',
\]
where $\|\cdot\|_{\mathrm{su}}$ and $\mathbf{1}_{\Delta_z}$ are as in Section~G1.3.
Proposition~5\,(b) bounds every row sum and every column sum of the matrix $\left(\left|B_{k,y}(A)(z,z')\right|_2\right)_{z,z'}$ by $(1+2\mathrm{e})\left(b_k^2+y^2\right)^{-\frac{1}{2}}$.
By (0.1) and (0.5),
\begin{equation}
\begin{split}
\left|\operatorname{Tr}\left[\gamma_5B_{k,y}(A)\,\mathrm{Ou}_y\,B_{k,y}(A)\right]\right|
&\le\sum_z4\left|\left(B_{k,y}(A)\,\mathrm{Ou}_y\,B_{k,y}(A)\right)(z,z)\right|_2\\
&\le 4\sum_z\sum_{z_1,z_2}\left|B_{k,y}(A)(z,z_1)\right|_2\left|\mathrm{Ou}_y(z_1,z_2)\right|_2\left|B_{k,y}(A)(z_2,z)\right|_2\\
&\le 4\cdot 32\,\mathrm{e}\cdot a'\sum_z\left(\sum_{z_1}\left|B_{k,y}(A)(z,z_1)\right|_2\right)\left(\sum_{z_2}\left|B_{k,y}(A)(z_2,z)\right|_2\right)\\
&\le 128\,\mathrm{e}\cdot(1+2\mathrm{e})^2\cdot a'\cdot n_k\left(b_k^2+y^2\right)^{-1}.
\end{split}
\tag{6.5}
\end{equation}
This proves the assertion of Proposition~12 that the trace in the first display is $O\left(\left(b_k^2+y^2\right)^{-1}\right)$ as $y\to\infty$.
By Proposition~4\,(b),
\begin{equation}
\begin{gathered}
\Gamma_{k,y}(A)=B_{k,y}(A)+b_k^2\,B_{k,y}(A)\,\mathrm{Ou}_y\,B_{k,y}(A),\qquad\text{so that}\\
\operatorname{Tr}\left[\gamma_5\Gamma_{k,y}(A)\right]=\operatorname{Tr}\left[\gamma_5B_{k,y}(A)\right]+b_k^2\operatorname{Tr}\left[\gamma_5B_{k,y}(A)\,\mathrm{Ou}_y\,B_{k,y}(A)\right].
\end{gathered}
\tag{6.6}
\end{equation}
The first two traces in (6.6) depend continuously on $y\in[0,\infty)$, by (4.2) and (6.3);
consequently so does the third.
Since
\[
\int_0^{\infty}\frac{\mathrm{d}y}{b_k^2+y^2}=\frac{\pi}{2b_k},
\]
the bound (6.5) shows that the integral in the first display converges absolutely, with
\begin{equation}
b_k^2\int_0^{\infty}\mathrm{d}y\,\left|\operatorname{Tr}\left[\gamma_5B_{k,y}(A)\,\mathrm{Ou}_y\,B_{k,y}(A)\right]\right|
\le 64\pi\,\mathrm{e}\cdot(1+2\mathrm{e})^2\cdot a'\cdot b_k\cdot n_k
\le 64\pi\,\mathrm{e}\cdot(1+2\mathrm{e})^2\cdot a'\cdot b\cdot n_k
\tag{6.7}
\end{equation}
($b_k\le b_+=b$: Section~I.3.4).

\paragraph{6.3. The display.}
Fix $\mathrm{yc}>0$.
Taking $\mathrm{ev}=\mathrm{ev}_j$ in (i) of Step~1, summing over $j=1,\dots,4n_k$ and inserting (4.2), (6.6) and (6.4), one gets
\begin{align*}
\sum_j\operatorname{Log}\mathrm{ev}_j
&=-i\int_0^{\mathrm{yc}}\mathrm{d}y\,\operatorname{Tr}\left[\gamma_5\Gamma_{k,y}(A)\right]+\sum_j\operatorname{Log}(\mathrm{ev}_j+i\,\mathrm{yc})\\
&=-4n_k\log\mathrm{yc}+c^{\mathrm{d}}_k\cdot n_k-\mathrm{rem}_1(\mathrm{yc})-i\,b_k^2\int_0^{\mathrm{yc}}\mathrm{d}y\,\operatorname{Tr}\left[\gamma_5B_{k,y}(A)\,\mathrm{Ou}_y\,B_{k,y}(A)\right]\\
&\qquad+4n_k\left(\log\mathrm{yc}+\frac{i\pi}{2}\right)+\mathrm{rem}_2(\mathrm{yc}),
\end{align*}
where
\[
\mathrm{rem}_2(\mathrm{yc}):=\sum_{j=1}^{4n_k}\left[\operatorname{Log}(\mathrm{ev}_j+i\,\mathrm{yc})-\log\mathrm{yc}-\frac{i\pi}{2}\right]\to 0\qquad(\mathrm{yc}\to\infty;\ \text{Step~1\,(c)}).
\]
The contributions $\pm 4n_k\log\mathrm{yc}$ cancel each other, and $4n_k\cdot\frac{i\pi}{2}=2\pi i\,n_k$.
Letting $\mathrm{yc}\to\infty$, and recalling from 6.2 that the integral converges, one arrives at
\begin{equation}
\sum_{j=1}^{4n_k}\operatorname{Log}\mathrm{ev}_j=c^{\mathrm{d}}_k\cdot n_k-i\,b_k^2\int_0^{\infty}\mathrm{d}y\,\operatorname{Tr}\left[\gamma_5B_{k,y}(A)\,Q_k(A)J^*S_{k,y}(A)JQ_k^T(-A)\,B_{k,y}(A)\right]+2\pi i\,n_k.
\tag{6.8}
\end{equation}
As $n_k$ is an integer, $\exp(2\pi i\,n_k)=1$;
taking exponentials in (6.8) and using (4.4), one obtains the display of Proposition~12 for $k\ge 1$.

\subsubsection*{Step 7. Uniformity}

All identities above are valid for every quadruple (SF) and for every real $A\in\tilde{G}_k$.
Consider the constants
\[
128\,\mathrm{rw}_0\cdot a',\qquad 128\,a',\qquad 128\,(\mathrm{rw}_0+1)\cdot a'\log 2
\]
of (5.3)--(5.5) and the constants
\[
128\,\mathrm{e}\cdot(1+2\mathrm{e})^2\cdot a',\qquad 64\pi\,\mathrm{e}\cdot(1+2\mathrm{e})^2\cdot a'\cdot b
\]
of (6.5), (6.7).
Apart from numerals, $\mathrm{e}$, $\pi$ and $\log 2$, the only constants they contain are $b$ and $a'$ (lines (C0), (C4r) of Section~I.8.2).
They depend neither on the data (SF), that is, on the level $k$, the side $\ell$ of the torus $T_k$, the mass $\mu\in(0,1]$ and the charge $e_k\in(0,e_0]$, nor on the connection $A$, nor on $y\ge 0$.
The dependence on the side is through the factor $n_k=\ell^3$ alone.
The eigenvalues $\mathrm{ev}_j$ depend on all the data, but the bounds (5.3)--(5.5), (6.5), (6.7) do not involve them:
the eigenvalues enter only the identities (4.2)--(4.4), (5.1), (5.2), (6.8), the remainder $\mathrm{rem}_2(\mathrm{yc})$, which tends to $0$ as $\mathrm{yc}\to\infty$, and the two majorants of 5.1, which serve to establish the absolute convergence of the integrals in (5.2).
\end{proof}

\subsection*{Conditions on the constants}

The proof requires no condition on the constants beyond those of Section~I.8.2 and Proposition~14.2\,(b).
It uses $b_k>0$, $0<b_L\le b$ and $b_k\le b_+=b$ (Sections~I.0.5, I.3.4; lines (C0), (C3), (C4) of Section~I.8.2: $b=1$, $L\ge 5$), $\rho<1$ (line (C0)) and the bound $\mu\le 1$ contained in the data (SF);
the inequalities $\rho<1$ and $\mu\le 1$ are needed only in 5.2, for the direct proof of (I.L0.12) at real $A$.
Besides $L\ge 5$, the steps of the proof use the following properties of $L$ and $M$, which follow from the conditions (L5) (on $L$, line (C3)) and (M1) (on $M$, line (C5)) of Proposition~14.2\,(b):
\begin{itemize}
\item[$\bullet$] $L$ is odd: in Step~2\,(2), (3), where the sets $\Delta_z\times O_{\mathrm{cub}}$ partition the fermion sites, and in the direct proof of (I.L0.12) in 5.2, where $c$ lies in exactly one $L$-block $B(y')$ and $B(y')$ has $L^3$ sites;
\item[$\bullet$] $M$ is an integer $\ge 1$: in the first paragraph of the proof, where $\ell\in\mathrm{Cl}$ is an integer multiple of $L$ and $13\,L\cdot M\ge 3L$;
\item[$\bullet$] $\frac{L-1}{2}\le M$: in the direct proof of (I.L0.12) in 5.2, where $|c-y'|\le\frac{L-1}{2}\le M$ for $c\in B(y')$, so that the phases $U_A(y',c)$ are defined ((I.1.6)).
\end{itemize}
The standard weights $w=(6,6\sqrt{2},6\sqrt{2},6)$, $\mathrm{sw}=12+12\sqrt{2}$, are fixed numbers, defined in the line (C0) of Section~I.8.2.
They enter the direct proof of (I.L0.12) in 5.2 through (I.m1), by which $\sum_iM_i(x)$ is the scalar $\frac{\rho}{2}\cdot\mathrm{sw}$, and through (I.m7);
they enter Step~2 through Proposition~2.2\,(e), namely through $D_A^*=\gamma_5D_A\gamma_5$.

The following conditions enter through the propositions on which the proof rests.
\begin{description}
\item[Through Proposition~2.8\,(a), (c), (d)]
The conditions (L5) (on $L$, line (C3)), (M1), (M2) and the part $M\ge 108\pi\cdot C'_V\cdot a'$ of (M5) (on $M$, line (C5)), (Z2.s1) and (Z2.s2) (on $e_0$, line (C8)) of Proposition~14.2\,(b);
the condition $\alpha<1$, which is part of the range of $\alpha$ (line (C1) of Section~I.8.2; the entry $\alpha$ of Proposition~14.2\,(a)).
These are the conditions stated with Proposition~2.8;
the condition $\alpha<1$ enters through Proposition~2.7\,(d), (e), and moreover $a'$ contains $(1-\alpha)^{-1}$, through $C_R$.
\item[Through Proposition~4\,(b), (d)]
The conditions of Proposition~2.8.
\item[Through Proposition~5\,(b), (c)]
The conditions (Z2.s1) (on $e_0$, line (C8)), (L5) (on $L$, line (C3)) and (M1) (on $M$, line (C5)) of Proposition~14.2\,(b).
\item[Through (I.L0.12)]
No condition, since the bound is used at $\mathfrak{i}(A)=0$, where $3(L-1)\,\mathfrak{i}(A)\le 1$ holds without (Z2.s1).
\item[Through Proposition~2.2\,(e), (f0)]
The part ``$L$ odd'' of the condition (L5) of Proposition~14.2\,(b) (on $L$, line (C3)), for the blocks $\mathrm{Bl}_j(z)$ of Section~G1.2: those of one side partition $\mathrm{Lab}$, they are nested, and each of them has $L^{3j}$ sites;
the condition $b>0$ (line (C0) of Section~I.8.2);
in the relations $D_A^*=\gamma_5D_A\gamma_5$ and $F_0(A)^*=\gamma_5F_0(A)\gamma_5$ of (e), the closure relation (I.1.1) of the standard weights (line (C0) of Section~I.8.2; (I.BARE))
(see the table in Chapter~2a, Section~5).
\end{description}

\section*{Chapter 13. Polymer expansion of the one-step determinant: proof of Proposition 13}

\subsection*{Proof}

\subsubsection*{0. Conventions of this chapter; standing facts}

\paragraph{0.1. Data and notation.}
Throughout the chapter the quadruple $(k,\ell,\mu,e_k)$ consists of data (SF), and $v_k:=-\log e_k$.
The conditions on $e_0$ listed in Proposition~14.2\,(b) are required for every $v\ge-\log e_0$;
in particular they hold at $v=v_k$, where $\exp(-\varepsilon v_k)=e_k^{\varepsilon}$.
The operators $S_{k,y}(X,A)$ are those of the family that belongs to the data of Proposition~13.
The letter $y$ stands for the resolvent parameter;
unit sites are denoted $z,z'$, and $y'$ denotes a site of the $L$-lattice.
Put $P:=P(A)$, $\mathrm{nX}:=d_M(X)$ and $\mathrm{e}:=\exp(1)$.

Let $\mathrm{Vc}$ be the complex vector space of the bond functions on $T_k$.
Its coordinates are the numbers $A(b)$, where $b$ ranges over the positively oriented bonds of $T_k$;
accordingly $\dim\mathrm{Vc}=3\left(\ell/\eta_k\right)^3$.
A bond function $A$ is real if all the numbers $A(b)$ are real, and $\overline{A}$ is the bond function with the coordinates $\overline{A(b)}$.
A function on an open subset of $\mathrm{Vc}$ is called analytic if it is holomorphic in these coordinates
(as in \cite[Section~5.3.1]{QEDI}; see Section~I.5.3).
Every space of spinors that occurs has finite dimension.
A function with values in operators is called analytic (resp.\ Borel) if the entries of its matrix with respect to fixed bases are analytic (resp.\ Borel).
With respect to other bases the entries are linear combinations of these, with coefficients that involve neither $A$ nor $y$;
hence the notion does not depend on the choice of the bases.
The set $D_k$ is finite, because $T_k$ contains $(\ell/M)^3$ $M$-cubes.

In this chapter the letters $m,m'$ always denote complex $4\times4$ matrices; no mass coupling occurs.
The symbols $\mathrm{Acl}$, $\mathrm{uf}$, $\mathrm{uf}^{0}$ (a closed real connection, its phase function, and the restriction of the latter to the unit lattice) have the meaning fixed in Section~G1.10.

\paragraph{0.2. The domain.}
The facts (G1), (G2), (G3) of \cite[Section~16.2.7]{QEDI}, combined with the conditions (Z2.s1) and (Z2.s4) of Proposition~14.2\,(b), yield the following.
\begin{itemize}
\item[(Do1)]
For $A\in\tilde{G}_k$ one has $\mathfrak{i}(A)<e_k^{\varepsilon}\le\frac{1}{6L}$.
Consequently $3\,\mathfrak{i}(A)\le1$, $3L\cdot\mathfrak{i}(A)\le\frac{1}{2}$ and $\exp\left(\mathfrak{i}(A)\right)\le\mathrm{e}$.
\item[(Do2)]
The set $\tilde{G}_k$ is open in $\mathrm{Vc}$ and star-shaped with respect to $0$, hence connected, and it contains $G_k$.
It is mapped into itself by each of the maps $A\mapsto A-d\boldsymbol{\lambda}$ (with $\boldsymbol{\lambda}$ a real function of $c$ on $T_k$), $A\mapsto-A$, $A\mapsto\overline{A}$ and $A\mapsto\tilde{g}A$ (with $\tilde{g}\in\mathbb{G}_k$).
\item[(Do3)]
One has $\mathfrak{X}_k\subset\tilde{G}_k$ and $\tilde{R}_k+\mathfrak{X}_k\subset\tilde{G}_k$.
\end{itemize}
Moreover, \cite[Lemma~16.2(a)]{QEDI} gives for every bond function $A$:
\begin{itemize}
\item[(Do4)]
The phases satisfy
\begin{gather*}
\left|u_A(c,c')^{\pm1}\right|\le\exp\left(\eta_k\cdot\mathfrak{i}(A)\right)\quad\text{for bonds }(c,c'),\qquad
\left|\varphi_A(z,c)^{\pm1}\right|\le\exp\left(\frac{3}{2}\,\mathfrak{i}(A)\right)\quad\text{for }c\in\Delta_z,\\
\left|U_A(c,c')^{\pm1}\right|\le\exp\left(3\,|c-c'|\cdot\mathfrak{i}(A)\right).
\end{gather*}
\end{itemize}

\paragraph{0.3. Geometry.}
By \cite[Section~17.4.1, (g0)]{QEDI} the $M$-cubes partition the set of the sites of $T_k$ and also the set of the sites of $T^0_k$,
and every unit block $\Delta_z$ and every $L$-block $B(y')$ is contained in a single $M$-cube.
Let $X$ be a union of $M$-cubes.
\begin{itemize}
\item[(Ge1)]
For a unit site $z$ the three conditions ``$z\in X$'', ``$\Delta_z\subset X$'' and ``$\Delta_z$ meets $X$'' are equivalent.
Indeed, $z$ belongs to $\Delta_z$, and $\Delta_z$ is contained in one $M$-cube, which is therefore the $M$-cube of $z$.
\item[(Ge2)]
If a site $z$ of $B(y')$ lies in $X$, then $B(y')\subset X$,
and the fine block $\mathrm{Bl}_{k+1}(y')$, that is, the union of the blocks $\Delta_{z'}$ with $z'\in B(y')$, is contained in $X$ as well.
Indeed, $B(y')$ is contained in one $M$-cube, namely the $M$-cube of $z$, and the second assertion follows from (Ge1).
For $k=0$ one has $\Delta_z=\{z\}$ and $\mathrm{Bl}_1(y')=B(y')$.
\item[(Ge3)]
The set $X$ contains exactly $M^3\cdot|X|_M$ unit sites:
as $M$ is odd, the $M$-cube $\{c:|c-z_0|<M/2\}$, $z_0\in M\mathbb{Z}^3$, contains precisely the $M^3$ unit sites $z$ with $|z_\mu-z_{0,\mu}|\le(M-1)/2$.
Moreover $\int_X\mathrm{d}x=M^3\cdot|X|_M$ (Section~I.1.6).
\item[(Ge4)]
A rectilinear path $\boldsymbol{\Gamma}^{\pi}(c,c')$ is contained in the lattice box spanned by $c$ and $c'$.
In the formula $(\tau_kA)(z,c)=\sum_{j<k}(\tau A)(y_{j+1},y_j)$ of (I.3.3) the points $y_j$, $y_{j+1}$ belong to $\Delta_z$, which is a lattice box.
In $(\tau A)(y',z)$ with $z\in B(y')$ both end points belong to $\mathrm{Bl}_{k+1}(y')$, which is a lattice box.
Hence, for $c\in\Delta_z$, the number $\varphi_A(z,c)^{\pm1}$ is the exponential of a linear form in the numbers $A(b)$, where $b$ ranges over the bonds with both end points in $\Delta_z$;
for $z\in B(y')$, the number $U_A(y',z)^{\pm1}$ is the exponential of a linear form in the numbers $A(b)$, where $b$ ranges over the bonds with both end points in $\mathrm{Bl}_{k+1}(y')$;
and $u_A(c,c_3)$ is the exponential of a multiple of $A(c,c_3)$.
In particular all three are entire functions on $\mathrm{Vc}$.
\end{itemize}

\paragraph{0.4. Kernels and traces.}
Let $\mathrm{FD}$ be a fundamental domain of $\ell\mathbb{Z}^3$ in $\mathbb{Z}^3$, and let $e_\beta$, $\beta=1,\dots,4$, be the unit vectors of $\mathbb{C}^4$.
Conditions such as ``$z\in X$'', ``$z=z'$'' or ``$z$ and $z'$ lie in one $L$-block'' are understood modulo $\ell\mathbb{Z}^3$ (Section~G1.2).

\emph{The case $k\ge1$, unit lattice.}
An operator $O$ on $\mathfrak{S}(1,\ell)$ has a kernel $O(z,z')$, defined for $z,z'\in\mathbb{Z}^3$ and with values in the $4\times4$ matrices, such that
\[
(Ov)(z)=\sum_{z'\in\mathrm{FD}}O(z,z')\,v(z'),\qquad O(z+\ell e_\mu,z')=O(z,z'+\ell e_\mu)=-O(z,z')
\]
(Section~G1.2, applied to the block torus $T_k=(0,\ell,-1)$).
For $z'\in\mathrm{FD}$ let $\delta_{z',\beta}$ be the element of $\mathfrak{S}(1,\ell)$ with $\delta_{z',\beta}(z)=\mathbf{1}[z=z']\,e_\beta$ for $z\in\mathrm{FD}$.
These $4n_k$ elements form a basis, and $(O\delta_{z',\beta})(z)=O(z,z')\,e_\beta$.
Therefore $\sum_{z\in\mathrm{FD}}\operatorname{tr}O(z,z)$ is the trace of the linear map $O$,
the kernel of a product is $(O_1O_2)(z,z'')=\sum_{z'\in\mathrm{FD}}O_1(z,z')\,O_2(z',z'')$, and
\begin{equation}
\operatorname{Tr}\left[O_1O_2\right]=\sum_{z,z'\in\mathrm{FD}}\operatorname{tr}\left[O_1(z,z')\,O_2(z',z)\right]=\operatorname{Tr}\left[O_2O_1\right],\qquad
\operatorname{Tr}\left[\mathrm{Um}\,O\,\mathrm{Um}^{-1}\right]=\operatorname{Tr}O
\tag{F0.1}
\end{equation}
for every invertible linear map $\mathrm{Um}$ of $\mathfrak{S}(1,\ell)$.

\emph{The case $k\ge1$, fine lattice.}
Let $\mathrm{FDf}$ be the union of the sets $\Delta_{z'}\times O_{\mathrm{cub}}$, $z'\in\mathrm{FD}$.
The blocks $\Delta_z$, $z\in\mathbb{Z}^3$, form a partition of $\mathbb{L}_k$ (Sections~I.1.1 and~I.1.4), and $\Delta_{z+\ell e_\mu}=\Delta_z+\ell e_\mu$.
So $\mathrm{FDf}$ is a fundamental domain of $\ell\mathbb{Z}^3$ in $\mathbb{L}^{\mathrm{f}}_k$, with $48\,\eta_k^{-3}\cdot n_k=48\left(\ell/\eta_k\right)^3$ sites.
An operator $O$ on $\mathfrak{S}^{\mathrm{ch}}(\eta_k,\ell)$ has the kernel (I.1.4),
\[
(Of)(x)=\int\mathrm{d}x'\,O(x,x')\,f(x')=\frac{\eta_k^3}{48}\sum_{x'\in\mathrm{FDf}}O(x,x')\,f(x'),
\]
subject to the same rule of signs.
For $x'\in\mathrm{FDf}$ let $\delta_{x',\beta}$ be the element of $\mathfrak{S}^{\mathrm{ch}}(\eta_k,\ell)$ with $\delta_{x',\beta}(x)=\mathbf{1}[x=x']\,e_\beta$ for $x\in\mathrm{FDf}$.
These $192\left(\ell/\eta_k\right)^3$ elements form a basis, and
\begin{equation}
(O\delta_{x',\beta})(x)=\frac{\eta_k^3}{48}\,O(x,x')\,e_\beta\qquad(x,x'\in\mathrm{FDf}).
\tag{F0.1a}
\end{equation}

\emph{The case $k=0$.}
An operator $O$ on $\mathfrak{S}^{\mathrm{ch}}(1,\ell)$ has the kernel (I.1.4),
\[
(Of)(x)=\int\mathrm{d}x'\,O(x,x')\,f(x')=\frac{1}{48}\sum_{c'\in\mathrm{FD}}\sum_{g'\in O_{\mathrm{cub}}}O\left(x,(c',g')\right)f(c',g'),
\]
subject to the same rule of signs.
For $x'\in\mathrm{FD}\times O_{\mathrm{cub}}$ let $\delta_{x',\beta}$ be the element of $\mathfrak{S}^{\mathrm{ch}}(1,\ell)$ with $\delta_{x',\beta}(x)=\mathbf{1}[x=x']\,e_\beta$ for $x\in\mathrm{FD}\times O_{\mathrm{cub}}$.
These $192\,n_0$ elements form a basis, and $(O\delta_{x',\beta})(x)=\frac{1}{48}\,O(x,x')\,e_\beta$.
Hence the trace of the linear map $O$ equals
$\frac{1}{48}\sum_{x\in\mathrm{FD}\times O_{\mathrm{cub}}}\operatorname{tr}O(x,x)=\int\mathrm{d}x\,\operatorname{tr}O(x,x)$, which is (I.ke2);
the kernel of a product is $(O_1O_2)(x,x'')=\int\mathrm{d}x'\,O_1(x,x')\,O_2(x',x'')$; and
\begin{equation}
\begin{aligned}
\operatorname{Tr}\left[O_1O_2\right]&=\int\mathrm{d}x\int\mathrm{d}x'\,\operatorname{tr}\left[O_1(x,x')\,O_2(x',x)\right]\\
&=\left(\frac{1}{48}\right)^2\sum_{c,c'\in\mathrm{FD}}\ \sum_{g,g'\in O_{\mathrm{cub}}}\operatorname{tr}\left[O_1\left((c,g),(c',g')\right)O_2\left((c',g'),(c,g)\right)\right],\\
\operatorname{Tr}\left[\mathrm{Um}\,O\,\mathrm{Um}^{-1}\right]&=\operatorname{Tr}O
\end{aligned}
\tag{F0.2}
\end{equation}
for every invertible linear map $\mathrm{Um}$ of $\mathfrak{S}^{\mathrm{ch}}(1,\ell)$.

\emph{Representatives.}
When a formula for a kernel is written down for particular representatives (for instance for $z,z'$ in one block $B(y')$ of $\mathbb{Z}^3$, or for $x'=x^{(3)}$ in $\mathbb{L}^{\mathrm{f}}_0$),
its values at the other representatives are given by the rule of signs.
The products $O_1(z,z')\,O_2(z',z)$ and $O_1(x,x')\,O_2(x',x)$ do not change when one of the two arguments is translated by $\ell e_\mu$, since two signs change;
so in (F0.1) and (F0.2) any representative of each argument may be used.
The moduli of the kernels are functions on the torus.

\emph{Entries.}
The identities $(O\delta_{z',\beta})(z)=O(z,z')\,e_\beta$, (F0.1a) and $(O\delta_{x',\beta})(x)=\frac{1}{48}\,O(x,x')\,e_\beta$ show the following.
In each of the three cases the entries of the kernel at points of the fundamental domain are the entries of the matrix of $O$ in the basis described above,
multiplied by $1$, by $48\,\eta_k^{-3}$ and by $48$, respectively; at other representatives a sign is added.
Consequently, if a function with values in operators is analytic in $A$ (resp.\ Borel in $y$) in the sense of Subsection~0.1, then so are the entries of its kernel.

\paragraph{0.5. Three inequalities.}
For complex $4\times4$ matrices $m,m'$ one has
\begin{equation}
\begin{aligned}
&\left|\operatorname{tr}(mm')\right|=\left|\sum_{\beta,\beta'}m_{\beta\beta'}\,m'_{\beta'\beta}\right|\le\left(\max_{\beta,\beta'}\left|m'_{\beta\beta'}\right|\right)\sum_{\beta,\beta'}\left|m_{\beta\beta'}\right|,\\
&\sum_{\beta,\beta'}\left|m_{\beta\beta'}\right|\le4\,|m|_c,\qquad\sum_{\beta,\beta'}\left|m_{\beta\beta'}\right|\le8\,|m|_2.
\end{aligned}
\tag{F0.3}
\end{equation}
Indeed, $|m|_c$ is the largest of the four row sums of $m$ (Section~I.1.5).
By the Cauchy--Schwarz inequality for four entries, the sum of the moduli of the entries of a row is at most twice its Euclidean norm;
and the row with the index $\beta$ is the complex conjugate of the vector $m^{*}e_\beta$, whose Euclidean norm is at most $|m^{*}|_2=|m|_2$.
Moreover, every entry of $m'$ has modulus at most $|m'|_c$ (Section~I.1.5).

\paragraph{0.6. Integrals that depend on a parameter.}
\begin{named}{Lemma A}
Let $\mathrm{dv}\ge1$ be an integer, $U$ an open subset of $\mathbb{C}^{\mathrm{dv}}$, $\mathrm{Iv}$ an interval of $\mathbb{R}$,
and $\tilde{f}\colon\mathrm{Iv}\times U\to\mathbb{C}$ a function with the following properties.
\begin{itemize}
\item[(i)] For every $y\in\mathrm{Iv}$ the function $\tilde{f}(y,\cdot\,)$ is holomorphic on $U$.
\item[(ii)] For every $A\in U$ the function $\tilde{f}(\cdot\,,A)$ is a Borel function.
\item[(iii)] There is a Borel function $\mathrm{maj}\ge0$ on $\mathrm{Iv}$ with $\int_{\mathrm{Iv}}\mathrm{d}y\,\mathrm{maj}(y)<\infty$ and $|\tilde{f}(y,A)|\le\mathrm{maj}(y)$ for all $(y,A)\in\mathrm{Iv}\times U$.
\end{itemize}
Then the integral $\mathrm{Fi}(A):=\int_{\mathrm{Iv}}\mathrm{d}y\,\tilde{f}(y,A)$ converges absolutely, $|\mathrm{Fi}(A)|\le\int_{\mathrm{Iv}}\mathrm{d}y\,\mathrm{maj}(y)$, and $\mathrm{Fi}$ is holomorphic on $U$.
\end{named}

\begin{proof}
The absolute convergence and the bound are consequences of (ii) and (iii).
Fix $A\in U$, and choose $\mathrm{rd}>0$ such that the closed polydisc $\{A':|A'_j-A_j|\le\mathrm{rd}\text{ for all }j\}$ is contained in $U$.

\emph{Continuity at $A$.}
Let $A^{(i)}\in U$ with $A^{(i)}\to A$.
By (i), $\tilde{f}(y,A^{(i)})\to\tilde{f}(y,A)$ for every $y$, and all these moduli are bounded by $\mathrm{maj}(y)$.
The dominated convergence theorem \cite[Theorem~1.34]{Rudin1987} gives $\mathrm{Fi}(A^{(i)})\to\mathrm{Fi}(A)$.

\emph{Partial derivatives at $A$.}
Fix $j$, let $e_j$ be the $j$-th unit vector of $\mathbb{C}^{\mathrm{dv}}$, and put $\mathrm{fq}_y(t):=\tilde{f}(y,A+te_j)$.
This function is holomorphic in a neighbourhood of $\{|t|\le\mathrm{rd}\}$, and $|\mathrm{fq}_y(t)|\le\mathrm{maj}(y)$.
Let $\mathrm{hq}$ be a complex number with $0<|\mathrm{hq}|\le\mathrm{rd}/2$.
Cauchy's formula \cite[Theorem~10.15]{Rudin1987}, applied at the points $\mathrm{hq}$ and $0$, together with the identity
$\frac{1}{t-\mathrm{hq}}-\frac{1}{t}=\frac{\mathrm{hq}}{t\,(t-\mathrm{hq})}$, yields
\[
\frac{\mathrm{fq}_y(\mathrm{hq})-\mathrm{fq}_y(0)}{\mathrm{hq}}=\frac{1}{2\pi i}\oint_{|t|=\mathrm{rd}}\mathrm{d}t\,\frac{\mathrm{fq}_y(t)}{t\,(t-\mathrm{hq})}.
\]
On the circle $|t|=\mathrm{rd}$ one has $|t-\mathrm{hq}|\ge\mathrm{rd}/2$;
so the modulus of the right side is at most
$\frac{2\pi\,\mathrm{rd}}{2\pi}\cdot\frac{2\,\mathrm{maj}(y)}{\mathrm{rd}^2}=\frac{2\,\mathrm{maj}(y)}{\mathrm{rd}}$.
As $\mathrm{hq}\to0$, the difference quotient tends to $\mathrm{fq}_y'(0)$.
Let $(\mathrm{hq}_i)$ be any sequence with $\mathrm{hq}_i\to0$ and $0<|\mathrm{hq}_i|\le\mathrm{rd}/2$.
The functions $y\mapsto\left(\mathrm{fq}_y(\mathrm{hq}_i)-\mathrm{fq}_y(0)\right)/\mathrm{hq}_i$ are Borel by (ii), they converge at every $y$,
and their moduli are bounded by $2\,\mathrm{maj}(y)/\mathrm{rd}$.
By dominated convergence, $\left(\mathrm{Fi}(A+\mathrm{hq}_ie_j)-\mathrm{Fi}(A)\right)/\mathrm{hq}_i$ converges to $\int_{\mathrm{Iv}}\mathrm{d}y\,\mathrm{fq}_y'(0)$.
Thus $\mathrm{Fi}$ possesses a complex derivative in each variable at every point of $U$.

Finally, a continuous function on an open subset of $\mathbb{C}^{\mathrm{dv}}$ that is holomorphic in each variable separately is holomorphic;
this is Osgood's lemma \cite[Chapter~I, Section~A]{GunningRossi1965}.
\end{proof}

\begin{named}{Remark B}
Let $\mathrm{fh}$ be holomorphic on an open set $U\subset\mathbb{C}^{\mathrm{dv}}$ that is mapped onto itself by $A\mapsto\overline{A}$.
Then $A\mapsto\overline{\mathrm{fh}(\overline{A})}$ is holomorphic on $U$.
Indeed, the quotient
\[
\frac{\overline{\mathrm{fh}\left(\overline{A}+\overline{\mathrm{hq}}\,e_j\right)}-\overline{\mathrm{fh}\left(\overline{A}\right)}}{\mathrm{hq}}
=\overline{\left(\frac{\mathrm{fh}\left(\overline{A}+\overline{\mathrm{hq}}\,e_j\right)-\mathrm{fh}\left(\overline{A}\right)}{\overline{\mathrm{hq}}}\right)}
\]
converges as $\mathrm{hq}\to0$, the function is continuous, and Osgood's lemma applies.
\end{named}

\subsubsection*{1. The integrand at the levels $k\ge1$}

Let $k\ge1$, $X\in D_k$, $A\in\tilde{G}_k$ and $y\ge0$.
Define the trace on $\mathfrak{S}(1,\ell)$
\[
\mathrm{Ig}_k(y;X,A):=\operatorname{Tr}\left[\gamma_5\,B_{k,y}(A)\,Q_k(A)J^{*}\,S_{k,y}(X,A)\,JQ_k^{T}(-A)\,B_{k,y}(A)\right],
\]
so that, by Section~I.5.7,
\[
\acute{E}_k(X,A)=-i\,b_k^2\int_0^{\infty}\mathrm{d}y\,\mathrm{Ig}_k(y;X,A).
\]

\paragraph{1.1. The operator $B_{k,y}(A)$.}
By Section~I.3.4 one has $b_k\ge b_{-}=b\left(1-1/L\right)>0$ and $b_L=b/L>0$.
The matrix $\gamma_5$ is Hermitian and $\gamma_5^2=1$ (I.c1);
hence $(b_k+i\gamma_5y)(b_k-i\gamma_5y)=b_k^2+y^2>0$, and the same holds with $b_k+b_L$ in the place of $b_k$.
Put
\[
\mathrm{bq}_1:=\left(b_k+i\gamma_5y\right)^{-1}=\frac{b_k-i\gamma_5y}{b_k^2+y^2},\qquad
\mathrm{bq}_2:=\left(b_k+b_L+i\gamma_5y\right)^{-1}=\frac{b_k+b_L-i\gamma_5y}{(b_k+b_L)^2+y^2}.
\]
These $4\times4$ matrices are polynomials in $\gamma_5$.
They commute with $\gamma_5$ and with each other, and they are normal,
with the eigenvalues $(b_k\pm iy)^{-1}$, resp.\ $(b_k+b_L\pm iy)^{-1}$, on the eigenspaces $\gamma_5=\pm1$.
It follows that
\[
|\mathrm{bq}_1|_2=\left(b_k^2+y^2\right)^{-\frac{1}{2}},\qquad
|\mathrm{bq}_2|_2=\left((b_k+b_L)^2+y^2\right)^{-\frac{1}{2}}\le\left(b_k^2+y^2\right)^{-\frac{1}{2}},\qquad
|\gamma_5|_2=1.
\]
By Proposition~2.2\,(f0) on $T_k$, $P^2=P$, and the kernel $P(A)(z,z')$ is a multiple of $1_4$;
so $P$ commutes with every constant $4\times4$ matrix acting at each site.
Consequently $P\,\mathrm{bq}_1(1-P)=\mathrm{bq}_1P(1-P)=0$ and $P\,\mathrm{bq}_2P=\mathrm{bq}_2P$, whence
\[
\left(b_k+b_LP+i\gamma_5y\right)\left[\mathrm{bq}_1(1-P)+\mathrm{bq}_2P\right]
=\left(b_k+i\gamma_5y\right)\mathrm{bq}_1(1-P)+\left(b_k+b_L+i\gamma_5y\right)\mathrm{bq}_2P=(1-P)+P=1.
\]
Since $\mathfrak{S}(1,\ell)$ has finite dimension, the operator $b_k+b_LP+i\gamma_5y$ is invertible,
and the definition of $B_{k,y}(A)$ in Section~I.3.5 gives
\begin{equation}
B_{k,y}(A)=\mathrm{bq}_1(1-P)+\mathrm{bq}_2P.
\tag{F1.1}
\end{equation}
Together with $(1-P)^2=1-P$ and $(1-P)P=P(1-P)=0$ this yields
\begin{equation}
\begin{aligned}
\mathrm{Gq}_y(A)&:=B_{k,y}(A)\,\gamma_5\,B_{k,y}(A)=\gamma_5\left[\mathrm{bq}_1^2(1-P)+\mathrm{bq}_2^2P\right],\\
\mathrm{Gq}_y(A)(z',z)&=\gamma_5\,\mathrm{bq}_1^2\,\mathbf{1}[z'=z]+\gamma_5\left(\mathrm{bq}_2^2-\mathrm{bq}_1^2\right)P(A)(z',z).
\end{aligned}
\tag{F1.2}
\end{equation}
By Proposition~5\,(a),
\[
P(A)(z',z)=\begin{cases}
L^{-3}\,U_A(y',z')^{-1}\,U_A(y',z)\,1_4&\text{if $z'$ and $z$ lie in one block $B(y')$,}\\
0&\text{otherwise,}
\end{cases}
\qquad\left|P(A)(z',z)\right|_2\le\mathrm{e}\cdot L^{-3}.
\]
A block $B(y')$ has $L^3$ sites.
Therefore $\left|\mathrm{Gq}_y(A)(z',z)\right|_2\le\left(b_k^2+y^2\right)^{-1}\left(\mathbf{1}[z'=z]+2\left|P(A)(z',z)\right|_2\right)$,
and for every unit site $z$
\begin{equation}
\sum_{z'\in\mathrm{FD}}\left|\mathrm{Gq}_y(A)(z',z)\right|_2\le\left(b_k^2+y^2\right)^{-1}\left(1+2L^3\cdot\mathrm{e}\cdot L^{-3}\right)=(1+2\mathrm{e})\left(b_k^2+y^2\right)^{-1}.
\tag{F1.3}
\end{equation}

\paragraph{1.2. The operator $\mathrm{Nq}_y(X,A):=Q_k(A)J^{*}\,S_{k,y}(X,A)\,JQ_k^{T}(-A)$ on $\mathfrak{S}(1,\ell)$.}
By Section~I.3.3,
\[
\left(Q_k(A)J^{*}f\right)(z)=\int_{\Delta_z}\mathrm{d}x\,\varphi_A(z,c)\,f(x)=\frac{\eta_k^3}{48}\sum_{c\in\Delta_z}\sum_{g\in O_{\mathrm{cub}}}\varphi_A(z,c)\,f(c,g),
\]
and $\left(JQ_k^{T}(-A)v\right)(x')=\varphi_A(z',c')^{-1}\,v(z')$ for $x'=(c',g')$ over $\Delta_{z'}$.
The integration over $x'$ extends over $\mathrm{FDf}$, the union of the sets $\Delta_{z'}\times O_{\mathrm{cub}}$ with $z'\in\mathrm{FD}$ (Subsection~0.4).
Writing $S_{k,y}(X,A)$ through its kernel (I.1.4), one obtains
\[
\left(\mathrm{Nq}_y(X,A)v\right)(z)=\int_{\Delta_z}\mathrm{d}x\,\varphi_A(z,c)\sum_{z'\in\mathrm{FD}}\int_{\Delta_{z'}}\mathrm{d}x'\,S_{k,y}(X,A)(x,x')\,\varphi_A(z',c')^{-1}\,v(z'),
\]
that is,
\begin{equation}
\begin{aligned}
\mathrm{Nq}_y(X,A)(z,z')&=\int_{\Delta_z}\mathrm{d}x\int_{\Delta_{z'}}\mathrm{d}x'\,\varphi_A(z,c)\,\varphi_A(z',c')^{-1}\,S_{k,y}(X,A)(x,x')\\
&=\left(\frac{\eta_k^3}{48}\right)^2\sum_{c\in\Delta_z}\ \sum_{g\in O_{\mathrm{cub}}}\ \sum_{c'\in\Delta_{z'}}\ \sum_{g'\in O_{\mathrm{cub}}}\varphi_A(z,c)\,\varphi_A(z',c')^{-1}\,S_{k,y}(X,A)\left((c,g),(c',g')\right).
\end{aligned}
\tag{F1.4}
\end{equation}

\emph{Support.}
By the property W12 of Section~G1.11 the kernel $S_{k,y}(X,A)(x,x')$ vanishes unless $c$ and $c'$ lie in $X$.
By (Ge1) it follows that $\mathrm{Nq}_y(X,A)(z,z')=0$ unless $z$ and $z'$ both lie in $X$.

\emph{Entries.}
Fix $z$, $z'$ and $\beta$, and let $f:=JQ_k^{T}(-A)\,\delta_{z',\beta}$.
Over $\Delta_{z'}$ this function equals $\varphi_A(z',c')^{-1}e_\beta$;
it is continued antiperiodically, and it vanishes off the orbit of $\Delta_{z'}\times O_{\mathrm{cub}}$.
Hence $f=\mathbf{1}_{\Delta_{z'}}f$, and $|f|_{\infty}\le\exp\left(\frac{3}{2}\,\mathfrak{i}(A)\right)$ by (Do4); this is the bound (I.Q3).
The column with the index $\beta$ of $\mathrm{Nq}_y(X,A)(z,z')$ is
\[
\mathrm{Nq}_y(X,A)(z,z')\,e_\beta=\left(Q_k(A)J^{*}S_{k,y}(X,A)f\right)(z)=\int_{\Delta_z}\mathrm{d}x\,\varphi_A(z,c)\left(\mathbf{1}_{\Delta_z}S_{k,y}(X,A)\mathbf{1}_{\Delta_{z'}}f\right)(x).
\]
Here $\int_{\Delta_z}\mathrm{d}x=1$, because each of the $48\,\eta_k^{-3}$ sites over $\Delta_z$ carries the weight $\eta_k^3/48$;
moreover $|\varphi_A(z,c)|\le\exp\left(\frac{3}{2}\,\mathfrak{i}(A)\right)$;
and the property W7 of Section~G1.11 for $\mathrm{par}=(\alpha_0,\gamma_{\mathrm{w}},\kappa_{\mathrm{w}};C_w)$ gives
\[
\left|\mathbf{1}_{\Delta_z}S_{k,y}(X,A)\mathbf{1}_{\Delta_{z'}}f\right|_{\infty}\le C_w\exp\left(-\gamma_{\mathrm{w}}\cdot d(z,z')-\kappa_{\mathrm{w}}\cdot\mathrm{nX}\right)|f|_{\infty}.
\]
Since $3\,\mathfrak{i}(A)\le1$ by (Do1) and $\gamma_{\mathrm{w}}\cdot d(z,z')\ge0$, every entry of $\mathrm{Nq}_y(X,A)(z,z')$ has modulus at most
\begin{equation}
\exp\left(3\,\mathfrak{i}(A)\right)C_w\exp\left(-\kappa_{\mathrm{w}}\cdot\mathrm{nX}\right)\le\mathrm{e}\cdot C_w\exp\left(-\kappa_{\mathrm{w}}\cdot\mathrm{nX}\right).
\tag{F1.5}
\end{equation}

\paragraph{1.3. The trace.}
The cyclicity (F0.1), the formula (F1.2) and the support of $\mathrm{Nq}_y(X,A)$ give
\begin{equation}
\mathrm{Ig}_k(y;X,A)=\operatorname{Tr}\left[\mathrm{Gq}_y(A)\,\mathrm{Nq}_y(X,A)\right]=\sum_{z,z'\in X\cap T^0_k}\operatorname{tr}\left[\mathrm{Gq}_y(A)(z',z)\,\mathrm{Nq}_y(X,A)(z,z')\right].
\tag{F1.6}
\end{equation}
The first and the third inequality of (F0.3), applied with $m=\mathrm{Gq}_y(A)(z',z)$ and $m'=\mathrm{Nq}_y(X,A)(z,z')$, together with (F1.5), (F1.3) and (Ge3), yield
\begin{equation}
\begin{aligned}
\left|\mathrm{Ig}_k(y;X,A)\right|&\le\sum_{z\in X\cap T^0_k}\ \sum_{z'\in\mathrm{FD}}8\left|\mathrm{Gq}_y(A)(z',z)\right|_2\cdot\mathrm{e}\cdot C_w\exp\left(-\kappa_{\mathrm{w}}\cdot\mathrm{nX}\right)\\
&\le8\,(1+2\mathrm{e})\,\mathrm{e}\cdot C_w\cdot M^3\cdot|X|_M\exp\left(-\kappa_{\mathrm{w}}\cdot\mathrm{nX}\right)\left(b_k^2+y^2\right)^{-1}.
\end{aligned}
\tag{F1.7}
\end{equation}
The right side does not depend on $A$.

\subsubsection*{2. The integrands at level $0$}

Let $k=0$; thus $\eta_0=1$, and $e_k$ is the number denoted $q$ in the flow.
Let $X\in D_0$, $A\in\tilde{G}_0$ and $y\ge0$.
All operators act on $\mathfrak{S}^{\mathrm{ch}}(1,\ell)$, a space of dimension $192\,n_0$.
Put
\[
\mathrm{Ig}'_0(y;X,A):=\operatorname{Tr}\left[F_0(A)\,S_{0,y}(X,A)\right],\qquad
\mathrm{Ig}''_0(y;X,A):=\operatorname{Tr}\left[\gamma_5\,S_{0,y}(X,A)\right],
\]
so that, by Section~I.5.7 and with the traces of (I.ke2),
\[
\acute{E}_0(X,A)=\int_1^{\infty}\frac{\mathrm{d}y}{y}\,\mathrm{Ig}'_0(y;X,A)-i\int_0^1\mathrm{d}y\,\mathrm{Ig}''_0(y;X,A).
\]

\paragraph{2.1. The operator $F_0(A)=D_A+\mu+b_LJP(A)J^{*}$.}
Let $f\in\mathfrak{S}^{\mathrm{ch}}(1,\ell)$ and $x=(c,g)$, and let $y'$ be the centre of the $L$-block that contains $c$.
By (I.2.3) at $\eta_0=1$ and by (I.m1), which gives $\sum_iM_i(x)=\frac{\rho}{2}\,\mathrm{sw}$,
and by (I.3.1) and the displays of Section~I.3.3, according to which
\[
\left(JP(A)J^{*}f\right)(c,g)=U_A(y',c)^{-1}L^{-3}\sum_{c'\in B(y')}U_A(y',c')\left(J^{*}f\right)(c'),
\]
one has
\begin{equation}
\begin{aligned}
\left(F_0(A)f\right)(x)={}&\left(\frac{\rho}{2}\,\mathrm{sw}+\mu\right)f(x)-\sum_{i=0}^{2}M_i(x)\,f(x^{(i)})-M_3(x)\,u_A(c,c_3)\,f(x^{(3)})\\
&+b_L\cdot L^{-3}\sum_{c'\in B(y')}U_A(y',c)^{-1}\,U_A(y',c')\,\frac{1}{48}\sum_{g'\in O_{\mathrm{cub}}}f(c',g').
\end{aligned}
\tag{F2.1}
\end{equation}
By (I.ke1), the evaluation at a point $x''$ has the kernel $48\cdot\mathbf{1}[x'=x'']$.
So the kernel (I.1.4) of $F_0(A)$ is, for $x'=(c',g')$,
\begin{equation}
\begin{aligned}
F_0(A)(x,x')={}&48\left\{\left(\frac{\rho}{2}\,\mathrm{sw}+\mu\right)\mathbf{1}[x'=x]-\sum_{i=0}^{2}M_i(x)\,\mathbf{1}[x'=x^{(i)}]-M_3(x)\,u_A(c,c_3)\,\mathbf{1}[x'=x^{(3)}]\right\}\\
&+b_L\cdot L^{-3}\,U_A(y',c)^{-1}\,U_A(y',c')\,\mathbf{1}[c'\in B(y')]\,1_4.
\end{aligned}
\tag{F2.2}
\end{equation}

\emph{Row sums.}
One has $\int\mathrm{d}x'\,48\cdot\mathbf{1}[x'=x'']=1$, and $\int\mathrm{d}x'\,\mathbf{1}[c'\in B(y')]\,(\cdot)=\sum_{c'\in B(y')}\frac{1}{48}\sum_{g'}(\cdot)$.
Hence, for every $x$ and every $\beta$,
\begin{align*}
\sum_{\beta'}\int\mathrm{d}x'\,\left|F_0(A)(x,x')_{\beta\beta'}\right|
\le{}&\frac{\rho}{2}\,\mathrm{sw}+\mu+|M_0(x)|_c+|M_1(x)|_c+|M_2(x)|_c+|M_3(x)|_c\cdot\left|u_A(c,c_3)\right|\\
&+b_L\cdot L^{-3}\sum_{c'\in B(y')}\left|U_A(y',c)^{-1}\right|\cdot\left|U_A(y',c')\right|.
\end{align*}
Since $\rho<1$, $\sqrt{2}<1.415$ (because $1.415^2=2.002225>2$), and by (I.m7),
\begin{gather*}
\frac{\rho}{2}\,\mathrm{sw}=\rho\left(6+6\sqrt{2}\right)<14.49,\qquad|M_0(x)|_c=3\,(1+\rho)<6,\\
|M_1(x)|_c+|M_2(x)|_c\le6\sqrt{2}\left(\rho+\sqrt{2}\right)<20.49,\qquad|M_3(x)|_c<6;
\end{gather*}
the first three of these quantities add up to less than $41$.
Further, $\mu\le1$, and $\left|u_A(c,c_3)\right|\le\exp\left(\mathfrak{i}(A)\right)$ by (Do4).
For $c,c'\in B(y')$ one has $|y'-c|,|y'-c'|\le(L-1)/2$;
hence $\left|U_A(y',c)^{-1}\right|\cdot\left|U_A(y',c')\right|\le\exp\left(3(L-1)\cdot\mathfrak{i}(A)\right)\le\mathrm{e}$ by (Do4) and (Do1).
Finally $b_L\le b$.
Since $1\le\exp\left(\mathfrak{i}(A)\right)\le\mathrm{e}$, one arrives at
\begin{equation}
\sup_x\sup_\beta\sum_{\beta'}\int\mathrm{d}x'\,\left|F_0(A)(x,x')_{\beta\beta'}\right|
\le41+1+6\exp\left(\mathfrak{i}(A)\right)+b\cdot\mathrm{e}
\le58\exp\left(\mathfrak{i}(A)\right)+1+b\cdot\mathrm{e}
\le\mathrm{rsF}:=58\,\mathrm{e}+1+b\cdot\mathrm{e}.
\tag{F2.3}
\end{equation}
This is the row-sum bound (I.L0.12) on $\tilde{G}_0$.

\paragraph{2.2. The traces.}
By W12 the kernel $S_{0,y}(X,A)(x',x)$ vanishes unless $c'$ and $c$ lie in $X$.
Therefore (F0.2) gives
\begin{equation}
\begin{aligned}
\mathrm{Ig}'_0(y;X,A)&=\int_X\mathrm{d}x\int_X\mathrm{d}x'\,\operatorname{tr}\left[F_0(A)(x,x')\,S_{0,y}(X,A)(x',x)\right],\\
\mathrm{Ig}''_0(y;X,A)&=\int_X\mathrm{d}x\,\operatorname{tr}\left[\gamma_5\,S_{0,y}(X,A)(x,x)\right],
\end{aligned}
\tag{F2.4}
\end{equation}
where $\int_X\mathrm{d}x=\frac{1}{48}\sum_{c\in X}\sum_{g\in O_{\mathrm{cub}}}$.
Let $y'(c)$ denote the centre of the $L$-block of $c$.
The sites $x^{(0)},x^{(1)},x^{(2)}$ lie over $c$, and $B(y'(c))\subset X$ by (Ge2).
Inserting (F2.2) into (F2.4), one obtains
\begin{equation}
\begin{aligned}
\mathrm{Ig}'_0(y;X,A)={}&\int_X\mathrm{d}x\left\{\left(\frac{\rho}{2}\,\mathrm{sw}+\mu\right)\operatorname{tr}S_{0,y}(X,A)(x,x)-\sum_{i=0}^{2}\operatorname{tr}\left[M_i(x)\,S_{0,y}(X,A)(x^{(i)},x)\right]\right.\\
&\qquad\left.-\mathbf{1}[c_3\in X]\,u_A(c,c_3)\operatorname{tr}\left[M_3(x)\,S_{0,y}(X,A)(x^{(3)},x)\right]\right\}\\
&+b_L\cdot L^{-3}\int_X\mathrm{d}x\sum_{c'\in B(y'(c))}\frac{1}{48}\sum_{g'\in O_{\mathrm{cub}}}U_A\left(y'(c),c\right)^{-1}\,U_A\left(y'(c),c'\right)\operatorname{tr}S_{0,y}(X,A)\left((c',g'),x\right).
\end{aligned}
\tag{F2.5}
\end{equation}

\emph{Bounds.}
By the property W7 for $\mathrm{par}$ in its form for $k=0$, every entry of $S_{0,y}(X,A)(x',x)$ has modulus at most
\[
\left|S_{0,y}(X,A)(x',x)\right|_c\le C_w\,(1+y)^{-1}\exp\left(-\gamma_{\mathrm{w}}\cdot d(x',x)-\kappa_{\mathrm{w}}\cdot\mathrm{nX}\right)\le C_w\,(1+y)^{-1}\exp\left(-\kappa_{\mathrm{w}}\cdot\mathrm{nX}\right).
\]
The first inequality of (F0.3), used under the integrals of (F2.4), then (F2.3) summed over the four values of $\beta$, and finally (Ge3), give
\begin{equation}
\begin{aligned}
\left|\mathrm{Ig}'_0(y;X,A)\right|&\le\int_X\mathrm{d}x\left[\sum_\beta\sum_{\beta'}\int\mathrm{d}x'\,\left|F_0(A)(x,x')_{\beta\beta'}\right|\right]C_w\,(1+y)^{-1}\exp\left(-\kappa_{\mathrm{w}}\cdot\mathrm{nX}\right)\\
&\le4\,\mathrm{rsF}\cdot C_w\cdot M^3\cdot|X|_M\,(1+y)^{-1}\exp\left(-\kappa_{\mathrm{w}}\cdot\mathrm{nX}\right),\\
\left|\mathrm{Ig}''_0(y;X,A)\right|&\le4\,C_w\cdot M^3\cdot|X|_M\,(1+y)^{-1}\exp\left(-\kappa_{\mathrm{w}}\cdot\mathrm{nX}\right).
\end{aligned}
\tag{F2.6}
\end{equation}
For the second bound one uses that $\gamma_5$ is monomial (I.c6), so that the moduli of its sixteen entries add up to $4$.
The right sides do not depend on $A$.
No factor $48$ enters (F2.6):
in (F2.3) the weight $1/48$ of the integration over $x'$ cancels the factor $48$ in the kernels of the point evaluations;
the weight $1/48$ of the integration over $x$ is absorbed by the sum over the $48$ labels, since $\int_X\mathrm{d}x=M^3\cdot|X|_M$ by (Ge3);
and the factor $48$ of a kernel bound at level $0$ is contained in the constant $C_w$ of W7 (see (I.L0.17) and (I.L0.26)).

\subsubsection*{3. Proof of (d1) and of the first assertion of (d0)}

\begin{proof}[Proof of Proposition~13\,(d1) and of the first assertion of Proposition~13\,(d0)]
\textbf{3.1.~The integrands as functions of $y$ and of $A$.}
Let $X\in D_k$ be fixed.

\emph{Borel dependence on $y$ at fixed $A$.}
By hypothesis the family is Borel in $y$: the map $y\mapsto S_{k,y}(X,A)$ is Borel on $[0,\infty)$.
By subsection~0.4 the entries of its kernel are therefore Borel functions of $y$.
The entries of $\mathrm{bq}_1$ and $\mathrm{bq}_2$ depend continuously on $y$, and none of the remaining factors in (F1.2), (F1.4), (F1.6) and (F2.4) depends on $y$.
Finite sums and finite products of Borel functions are Borel.
Hence the functions
\[
y\mapsto\mathrm{Ig}_k(y;X,A)\quad(k\ge1),\qquad y\mapsto\mathrm{Ig}'_0(y;X,A),\qquad y\mapsto\mathrm{Ig}''_0(y;X,A)
\]
are Borel.

\emph{Analytic dependence on $A$ at fixed $y$.}
By property W2 of Section~G1.11 the operator $S_{k,y}(X,\cdot)$ is analytic on $\tilde{G}_k$, and by subsection~0.4 so are the entries of its kernel.
By (Ge4) the numbers $\varphi_A(z,c)^{\pm1}$, $U_A(y',z)^{\pm1}$ and $u_A(c,c_{3})$ are entire functions on $\mathrm{Vc}$, whereas $\mathrm{bq}_1$, $\mathrm{bq}_2$, $\gamma_5$ and the matrices $M_{i}(x)$ do not depend on $A$.
For $k\ge1$ the formula (F1.2), read with Proposition~5\,(a), and the formulas (F1.4) and (F1.6) exhibit the integrand as a finite sum of products of functions of these kinds;
for $k=0$ the same follows from (F2.5) and from the formula for $\mathrm{Ig}''_0$ in (F2.4).
Consequently the integrands are analytic on $\tilde{G}_k$, which is open in $\mathrm{Vc}$ by (Do2).

\textbf{3.2.~The functions $\acute{E}_k$, $k\ge1$.}
Lemma~A is applied with $U=\tilde{G}_k$, $\mathrm{Iv}=[0,\infty)$, $\tilde{f}=\mathrm{Ig}_k(\,\cdot\,;X,\,\cdot\,)$, and with $\mathrm{maj}$ equal to the right side of (F1.7).
Hypotheses (i) and (ii) of the lemma hold by subsection~3.1, and (iii) holds because $b_k>0$ (Section~I.3.4) and
\[
\int_0^{\infty}\mathrm{d}y\,\left(b_k^2+y^2\right)^{-1}=\frac{\pi}{2b_k}.
\]
Hence the integral converges absolutely, which is the first assertion of (d0) for $k\ge1$.
Moreover $\acute{E}_k(X,\cdot)$ is analytic on $\tilde{G}_k$, and, since $\acute{E}_k(X,A)=-i\cdot b_k^2\int_0^{\infty}\mathrm{d}y\,\mathrm{Ig}_k(y;X,A)$,
\begin{align*}
\left|\acute{E}_k(X,A)\right|
&\le b_k^2\cdot8(1+2\mathrm{e})\cdot\mathrm{e}\cdot C_{w}\cdot M^3\cdot|X|_M\cdot\exp\left(-\kappa_{\mathrm{w}}\cdot d_M(X)\right)\cdot\frac{\pi}{2b_k}\\
&=8\cdot\frac{\pi}{2}\cdot b_k\cdot\mathrm{e}\cdot(1+2\mathrm{e})\cdot C_{w}\cdot M^3\cdot|X|_M\cdot\exp\left(-\kappa_{\mathrm{w}}\cdot d_M(X)\right).
\end{align*}
Now $b_k\le b$ by Section~I.3.4.
By \cite[Theorem~13.2(t1)]{QEDI} and the inequality $d_M(X)+1\le\exp\left(d_M(X)\right)$,
\[
|X|_M\le c_{X}\cdot\left(d_M(X)+1\right)\le c_{X}\cdot\exp\left(d_M(X)\right).
\]
Since $\kappa_{\mathrm{w}}=\kappa+\kappa_{0}+3$, one has $\kappa_{\mathrm{w}}-1=\kappa+\kappa_{0}+2\ge\kappa+1$.
With the definition of $C_{d}$ it follows that
\begin{equation}
\begin{split}
\left|\acute{E}_k(X,A)\right|
&\le8\cdot c_{X}\cdot C_{w}\cdot\frac{\pi}{2}\cdot b\cdot\mathrm{e}\cdot(1+2\mathrm{e})\cdot M^3\cdot\exp\left(-(\kappa_{\mathrm{w}}-1)\cdot d_M(X)\right)\\
&\le\frac{1}{2}\,C_{d}\cdot M^3\cdot\exp\left(-(\kappa+1)\cdot d_M(X)\right)\qquad(k\ge1).
\end{split}
\tag{F3.1}
\end{equation}

\textbf{3.3.~The function $\acute{E}_0$.}
Here Lemma~A is applied twice, with $U=\tilde{G}_0$.
The first application uses $\mathrm{Iv}=[1,\infty)$, $\tilde{f}(y,A)=y^{-1}\,\mathrm{Ig}'_0(y;X,A)$ and
\[
\mathrm{maj}(y)=4\,\mathrm{rsF}\cdot C_{w}\cdot M^3\cdot|X|_M\cdot\exp\left(-\kappa_{\mathrm{w}}\cdot d_M(X)\right)\cdot\frac{1}{y\,(1+y)};
\]
the second uses $\mathrm{Iv}=[0,1]$, $\tilde{f}(y,A)=\mathrm{Ig}''_0(y;X,A)$ and
\[
\mathrm{maj}(y)=4\,C_{w}\cdot M^3\cdot|X|_M\cdot\exp\left(-\kappa_{\mathrm{w}}\cdot d_M(X)\right)\cdot\frac{1}{1+y}.
\]
By (F2.6) these are majorants, and hypotheses (i), (ii) of Lemma~A hold by subsection~3.1.
One has
\[
\int_1^{\infty}\mathrm{d}y\,\frac{1}{y\,(1+y)}=\log2=\int_0^1\mathrm{d}y\,\frac{1}{1+y},\qquad\log2\le1.
\]
Both integrals therefore converge absolutely, which is the first assertion of (d0) for $k=0$, and $\acute{E}_0(X,\cdot)$ is analytic on $\tilde{G}_0$.
The modulus of $\acute{E}_0(X,A)$ is at most the sum of the integrals of the two majorants.
With $\mathrm{rsF}+1=2+58\mathrm{e}+b\cdot\mathrm{e}$ and with the bound $|X|_M\le c_{X}\cdot\exp\left(d_M(X)\right)$ of subsection~3.2 one obtains
\begin{equation}
\begin{split}
\left|\acute{E}_0(X,A)\right|
&\le4(\mathrm{rsF}+1)\cdot C_{w}\cdot M^3\cdot|X|_M\cdot\exp\left(-\kappa_{\mathrm{w}}\cdot d_M(X)\right)\\
&\le4\cdot c_{X}\cdot C_{w}\cdot(2+58\mathrm{e}+b\cdot\mathrm{e})\cdot M^3\cdot\exp\left(-(\kappa_{\mathrm{w}}-1)\cdot d_M(X)\right)\\
&\le\frac{1}{4}\,C_{d}\cdot M^3\cdot\exp\left(-(\kappa+1)\cdot d_M(X)\right).
\end{split}
\tag{F3.2}
\end{equation}

\textbf{3.4.~The functions $E^{d}_k$, $k\ge0$.}
By (Do2) the maps $A\mapsto-A$ and $A\mapsto\overline{A}$ send $\tilde{G}_k$ into itself; as they are involutions, they map $\tilde{G}_k$ onto itself.
The function $A\mapsto\acute{E}_k(X,-A)$ is analytic, being the composition of an analytic function with a linear map.
The functions $A\mapsto\overline{\acute{E}_k(X,\overline{A})}$ and $A\mapsto\overline{\acute{E}_k(X,-\overline{A})}$ are analytic by Remark~B, applied to $\acute{E}_k(X,\cdot)$ and to $A\mapsto\acute{E}_k(X,-A)$.
On $\tilde{G}_k$ each of these four functions is bounded by the right side of (F3.1), respectively of (F3.2).
By Section~I.5.7 the function $E^{d}_k(X,\cdot)$ is their mean.
It is therefore analytic on $\tilde{G}_k$ and obeys the same bound.
As $\frac{1}{2}C_{d}\le C_{d}$ and $\frac{1}{4}C_{d}\le C_{d}$, this proves (d1).
\end{proof}

\subsubsection*{4. Proof of (d2)}

\begin{proof}[Proof of Proposition~13\,(d2)]
Let $A,\tilde{A}\in\tilde{G}_k$ coincide on the bonds that meet $X$.
It will be shown that for every $y\ge0$ the integrands at $A$ and at $\tilde{A}$ are equal; the equality $\acute{E}_k(X,A)=\acute{E}_k(X,\tilde{A})$ follows.

\emph{The case $k\ge1$.}
In (F1.6) the sites $z$ and $z'$ range over the unit sites of $X$.
\begin{itemize}
\item[(i)] By (F1.4), the matrix $\mathrm{Nq}_y(X,A)(z,z')$ contains $A$ through $S_{k,y}(X,A)$, which equals $S_{k,y}(X,\tilde{A})$ by W12, and through the phases $\varphi_A(z,c)$, $c\in\Delta_z$, and $\varphi_A(z',c')^{-1}$, $c'\in\Delta_{z'}$.
By (Ge4) these phases depend only on the values of $A$ on the bonds with both end points in $\Delta_z$, respectively in $\Delta_{z'}$; by (Ge1) these blocks are contained in $X$.
\item[(ii)] By (F1.2), the matrix $\mathrm{Gq}_y(A)(z',z)$ contains $A$ only through $P(A)(z',z)$.
Whatever $A$ is, this matrix vanishes unless $z'$ and $z$ lie in one block $B(y')$.
In that case it equals $L^{-3}\cdot U_A(y',z')^{-1}\cdot U_A(y',z)\cdot1_4$ by Proposition~5\,(a), and by (Ge4) this depends only on the values of $A$ on the bonds with both end points in $\mathrm{Bl}_{k+1}(y')$; moreover $\mathrm{Bl}_{k+1}(y')\subset X$ by (Ge2).
\end{itemize}
A bond with both end points in $X$ meets $X$.
Hence $\mathrm{Ig}_k(y;X,A)=\mathrm{Ig}_k(y;X,\tilde{A})$.

\emph{The case $k=0$.}
By W12, $S_{0,y}(X,A)=S_{0,y}(X,\tilde{A})$, and so $\mathrm{Ig}''_0(y;X,A)=\mathrm{Ig}''_0(y;X,\tilde{A})$ by (F2.4).
In (F2.5) the site $x$ ranges over the sites over $X$, and apart from $S_{0,y}(X,A)$ the connection occurs only
\begin{itemize}
\item[(i)] in $u_A(c,c_{3})$, which depends only on the value of $A$ on the bond $(c,c_{3})$, a bond with the end point $c$ in $X$;
\item[(ii)] in $U_A(y'(c),c)^{-1}\cdot U_A(y'(c),c')$, $c'\in B(y'(c))$, which by (Ge4) depends only on the values of $A$ on the bonds with both end points in $B(y'(c))$, a subset of $X$ by (Ge2).
\end{itemize}
The hops of the types $0$, $1$, $2$ carry no phase.
Hence $\mathrm{Ig}'_0(y;X,A)=\mathrm{Ig}'_0(y;X,\tilde{A})$.

\emph{The functions $E^{d}_k$.}
By (Do2) the connections $-A$ and $-\tilde{A}$, $\overline{A}$ and $\overline{\tilde{A}}$, $-\overline{A}$ and $-\overline{\tilde{A}}$ lie in $\tilde{G}_k$, and the two connections of each of these pairs coincide on the bonds that meet $X$.
Therefore each of the four terms of $E^{d}_k(X,A)$ equals the corresponding term of $E^{d}_k(X,\tilde{A})$.
\end{proof}

\subsubsection*{5. Proof of (d3)}

\begin{proof}[Proof of Proposition~13\,(d3)]
\textbf{5.1.~Gauge invariance of $\acute{E}_k$.}
Let $\boldsymbol{\lambda}$ be a real function of $c$ on $T_k$, put $\mathrm{uf}:=\exp\left(i e_k\cdot\boldsymbol{\lambda}\right)$, and let $\mathrm{uf}^{0}$ be the restriction of $\mathrm{uf}$ to $T^0_k$.
Then $d\boldsymbol{\lambda}$ is a real connection with $d(d\boldsymbol{\lambda})=0$, and $|\mathrm{uf}|=1$.
Since $\eta_k\cdot(d\boldsymbol{\lambda})(c,c')=\boldsymbol{\lambda}(c')-\boldsymbol{\lambda}(c)$, one has
\[
\mathrm{uf}(c')=\mathrm{uf}(c)\cdot\exp\left\{i e_k\cdot\eta_k\cdot(d\boldsymbol{\lambda})(c,c')\right\}
\]
on every bond $(c,c')$.
By (Do2), $A-d\boldsymbol{\lambda}\in\tilde{G}_k$.
Proposition~2.2\,(a) therefore applies on $T_k$ with $\mathrm{Acl}=d\boldsymbol{\lambda}$.

\emph{The case $k\ge1$.}
Proposition~2.2\,(a) gives
\begin{gather*}
Q_k(A-d\boldsymbol{\lambda})J^{*}=\mathrm{uf}^{0}\cdot Q_k(A)J^{*}\cdot\mathrm{uf}^{-1},\qquad JQ_k^{T}\left(-(A-d\boldsymbol{\lambda})\right)=\mathrm{uf}\cdot JQ_k^{T}(-A)\cdot\left(\mathrm{uf}^{0}\right)^{-1},\\
P(A-d\boldsymbol{\lambda})=\mathrm{uf}^{0}\cdot P(A)\cdot\left(\mathrm{uf}^{0}\right)^{-1}.
\end{gather*}
Multiplication by the scalar function $\mathrm{uf}^{0}$ commutes with $\gamma_5$ and with numbers.
Hence
\[
b_k+b_{L}\cdot P(A-d\boldsymbol{\lambda})+i\gamma_5y=\mathrm{uf}^{0}\cdot\left(b_k+b_{L}\cdot P(A)+i\gamma_5y\right)\cdot\left(\mathrm{uf}^{0}\right)^{-1},
\]
and, by the definition of $B_{k,y}$ in Section~I.3.5, $B_{k,y}(A-d\boldsymbol{\lambda})=\mathrm{uf}^{0}\cdot B_{k,y}(A)\cdot\left(\mathrm{uf}^{0}\right)^{-1}$.
Property W13 gives $S_{k,y}(X,A-d\boldsymbol{\lambda})=\mathrm{uf}\cdot S_{k,y}(X,A)\cdot\mathrm{uf}^{-1}$.
In the product the inner factors cancel, and
\begin{equation}
\begin{split}
&\gamma_5\cdot B_{k,y}(A-d\boldsymbol{\lambda})\cdot Q_k(A-d\boldsymbol{\lambda})J^{*}\cdot S_{k,y}(X,A-d\boldsymbol{\lambda})\cdot JQ_k^{T}\left(-(A-d\boldsymbol{\lambda})\right)\cdot B_{k,y}(A-d\boldsymbol{\lambda})\\
&\qquad=\mathrm{uf}^{0}\cdot\left[\gamma_5\cdot B_{k,y}(A)\cdot Q_k(A)J^{*}\cdot S_{k,y}(X,A)\cdot JQ_k^{T}(-A)\cdot B_{k,y}(A)\right]\cdot\left(\mathrm{uf}^{0}\right)^{-1}.
\end{split}
\tag{F5.1}
\end{equation}
By (F0.1) it follows that $\mathrm{Ig}_k(y;X,A-d\boldsymbol{\lambda})=\mathrm{Ig}_k(y;X,A)$ for every $y\ge0$.

\emph{The case $k=0$.}
Proposition~2.2\,(a) at $k=0$, applied to $O=F_{y}$ at $y=0$, gives $F_{0}(A-d\boldsymbol{\lambda})=\mathrm{uf}\cdot F_{0}(A)\cdot\mathrm{uf}^{-1}$, and W13 gives $S_{0,y}(X,A-d\boldsymbol{\lambda})=\mathrm{uf}\cdot S_{0,y}(X,A)\cdot\mathrm{uf}^{-1}$.
As $\mathrm{uf}$ is a function of $c$, the $48$ sites of a cluster carry the same phase, and $\mathrm{uf}$ commutes with $\gamma_5$.
Consequently
\begin{equation}
\begin{aligned}
F_{0}(A-d\boldsymbol{\lambda})\cdot S_{0,y}(X,A-d\boldsymbol{\lambda})&=\mathrm{uf}\cdot F_{0}(A)\cdot S_{0,y}(X,A)\cdot\mathrm{uf}^{-1},\\
\gamma_5\cdot S_{0,y}(X,A-d\boldsymbol{\lambda})&=\mathrm{uf}\cdot\gamma_5\cdot S_{0,y}(X,A)\cdot\mathrm{uf}^{-1},
\end{aligned}
\tag{F5.2}
\end{equation}
and by (F0.2) each of the two traces takes the same value at $A-d\boldsymbol{\lambda}$ as at $A$.

Thus $\acute{E}_k(X,A-d\boldsymbol{\lambda})=\acute{E}_k(X,A)$ for every $k\ge0$, every $A\in\tilde{G}_k$ and every real function $\boldsymbol{\lambda}$ of $c$.

\textbf{5.2.~Gauge invariance, evenness and reality of $E^{d}_k$.}
As $\boldsymbol{\lambda}$ is real,
\[
-(A-d\boldsymbol{\lambda})=(-A)-d(-\boldsymbol{\lambda}),\qquad\overline{A-d\boldsymbol{\lambda}}=\overline{A}-d\boldsymbol{\lambda},\qquad-\overline{(A-d\boldsymbol{\lambda})}=(-\overline{A})-d(-\boldsymbol{\lambda}).
\]
Subsection~5.1, applied at $A$ with $\boldsymbol{\lambda}$ and at the points $-A$, $\overline{A}$, $-\overline{A}$ of $\tilde{G}_k$ with the real functions $-\boldsymbol{\lambda}$, $\boldsymbol{\lambda}$, $-\boldsymbol{\lambda}$, shows that each of the four terms of $E^{d}_k(X,A-d\boldsymbol{\lambda})$ equals the corresponding term of $E^{d}_k(X,A)$.

By Section~I.5.7,
\[
E^{d}_k(X,A)=\frac{1}{4}\left[\mathrm{Eq}_1+\mathrm{Eq}_2+\mathrm{Eq}_3+\mathrm{Eq}_4\right](A),
\]
where
\[
\mathrm{Eq}_1(A)=\acute{E}_k(X,A),\qquad\mathrm{Eq}_2(A)=\acute{E}_k(X,-A),\qquad\mathrm{Eq}_3(A)=\overline{\acute{E}_k(X,\overline{A})},\qquad\mathrm{Eq}_4(A)=\overline{\acute{E}_k(X,-\overline{A})}.
\]
Then
\[
\mathrm{Eq}_1(-A)=\mathrm{Eq}_2(A),\qquad\mathrm{Eq}_2(-A)=\mathrm{Eq}_1(A),\qquad\mathrm{Eq}_3(-A)=\mathrm{Eq}_4(A),\qquad\mathrm{Eq}_4(-A)=\mathrm{Eq}_3(A),
\]
so that $E^{d}_k(X,-A)=E^{d}_k(X,A)$; and
\[
\overline{\mathrm{Eq}_1(\overline{A})}=\mathrm{Eq}_3(A),\qquad\overline{\mathrm{Eq}_2(\overline{A})}=\mathrm{Eq}_4(A),\qquad\overline{\mathrm{Eq}_3(\overline{A})}=\mathrm{Eq}_1(A),\qquad\overline{\mathrm{Eq}_4(\overline{A})}=\mathrm{Eq}_2(A),
\]
so that $\overline{E^{d}_k(X,\overline{A})}=E^{d}_k(X,A)$.

\textbf{5.3.~Lattice symmetries.}
Let $\tilde{g}=(S,a)\in\mathbb{G}_k^{L}$ and $r:=r(S)$.
The polymer $\tilde{g}X=rX+a$ (a polymer by Proposition~16.1\,(a1)) and the connection $(\tilde{g}A)(c,c')=A\left(r(c-a),r(c'-a)\right)$ depend on $S$ only through $r$ (Section~I.4.3).
By Section~I.4.1, $r(\mathrm{Sp}_3)=\{\pm1\}^3=r(\mathrm{Sp})$.
Hence there is $S'\in\mathrm{Sp}_3$ with $r(S')=r$, and the element $\tilde{g}^{(3)}:=(S',a)$ of $\mathbb{G}^{L}_{k,3}$ (Section~I.4.2) satisfies $\tilde{g}^{(3)}X=\tilde{g}X$ and $\tilde{g}^{(3)}A=\tilde{g}A$.
It therefore suffices to consider $\tilde{g}\in\mathbb{G}^{L}_{k,3}$, that is, $S\gamma_5=\gamma_5S$.
By (Do2), $\tilde{g}A\in\tilde{G}_k$.
On $\mathfrak{S}(1,\ell)$ and on $\mathfrak{S}^{\mathrm{ch}}(\eta_k,\ell)$ the element $\tilde{g}$ acts by the invertible linear maps
\[
(\tilde{g}v)(z)=S\cdot v\left(r(z-a)\right),\qquad(\tilde{g}f)(c,g)=S\cdot f\left(r(c-a),rg\right),
\]
where $a$ is a representative in $LM\mathbb{Z}^3$; antiperiodic functions are mapped to antiperiodic functions, because $r\,\ell e_{\mu}=\pm\ell e_{\mu}$.
These maps commute with $\gamma_5$.

\emph{The case $k\ge1$.}
Proposition~10\,(b), for $S\in\mathrm{Sp}_3$ and with the same $y$ on both sides, gives
\[
Q_k(\tilde{g}A)J^{*}\cdot\tilde{g}=\tilde{g}\cdot Q_k(A)J^{*},\qquad JQ_k^{T}(-\tilde{g}A)\cdot\tilde{g}=\tilde{g}\cdot JQ_k^{T}(-A),\qquad B_{k,y}(\tilde{g}A)\cdot\tilde{g}=\tilde{g}\cdot B_{k,y}(A),
\]
and W14 for $\mathbb{G}^{L}_{k,3}$ gives $S_{k,y}(\tilde{g}X,\tilde{g}A)\cdot\tilde{g}=\tilde{g}\cdot S_{k,y}(X,A)$.
Consequently
\begin{equation}
\begin{split}
&\gamma_5\cdot B_{k,y}(\tilde{g}A)\cdot Q_k(\tilde{g}A)J^{*}\cdot S_{k,y}(\tilde{g}X,\tilde{g}A)\cdot JQ_k^{T}(-\tilde{g}A)\cdot B_{k,y}(\tilde{g}A)\\
&\qquad=\tilde{g}\cdot\left[\gamma_5\cdot B_{k,y}(A)\cdot Q_k(A)J^{*}\cdot S_{k,y}(X,A)\cdot JQ_k^{T}(-A)\cdot B_{k,y}(A)\right]\cdot\tilde{g}^{-1},
\end{split}
\tag{F5.3}
\end{equation}
and $\mathrm{Ig}_k(y;\tilde{g}X,\tilde{g}A)=\mathrm{Ig}_k(y;X,A)$ by (F0.1).

\emph{The case $k=0$.}
Proposition~10\,(g) together with Proposition~10\,(b), applied to $F_{y}$ at the neutral value $y=0$, gives $F_{0}(\tilde{g}A)\cdot\tilde{g}=\tilde{g}\cdot F_{0}(A)$, and W14 gives $S_{0,y}(\tilde{g}X,\tilde{g}A)\cdot\tilde{g}=\tilde{g}\cdot S_{0,y}(X,A)$.
Therefore
\[
F_{0}(\tilde{g}A)\cdot S_{0,y}(\tilde{g}X,\tilde{g}A)=\tilde{g}\cdot F_{0}(A)\cdot S_{0,y}(X,A)\cdot\tilde{g}^{-1},\qquad\gamma_5\cdot S_{0,y}(\tilde{g}X,\tilde{g}A)=\tilde{g}\cdot\gamma_5\cdot S_{0,y}(X,A)\cdot\tilde{g}^{-1},
\]
and the traces coincide by (F0.2).
(The element $\tilde{g}$ permutes the $48$ labels of a cluster by $g\mapsto rg$; the trace (I.ke2) is the trace of the linear map and is not affected by this permutation.)

Altogether,
\begin{equation}
\acute{E}_k(\tilde{g}X,\tilde{g}A)=\acute{E}_k(X,A)\qquad\text{for all }\tilde{g}\in\mathbb{G}_k^{L},\ A\in\tilde{G}_k,\ k\ge0.
\tag{F5.4}
\end{equation}
Finally, $\tilde{g}(-A)=-\tilde{g}A$ and $\tilde{g}\,\overline{A}=\overline{\tilde{g}A}$.
So the four terms of $E^{d}_k(\tilde{g}X,\tilde{g}A)$ are
\[
\acute{E}_k(\tilde{g}X,\tilde{g}A),\qquad\acute{E}_k\left(\tilde{g}X,\tilde{g}(-A)\right),\qquad\overline{\acute{E}_k\left(\tilde{g}X,\tilde{g}\,\overline{A}\right)},\qquad\overline{\acute{E}_k\left(\tilde{g}X,\tilde{g}(-\overline{A})\right)},
\]
and by (F5.4), applied at the points $A$, $-A$, $\overline{A}$, $-\overline{A}$ of $\tilde{G}_k$, they equal the four terms of $E^{d}_k(X,A)$.
This proves (d3).
\end{proof}

\subsubsection*{6. Proof of (d0)}

\begin{proof}[Proof of Proposition~13\,(d0)]
\emph{The first assertion of (d0).}
For every $X\in D_k$ and every $A\in\tilde{G}_k$ the integrals in the formulas for $\acute{E}_k(X,A)$ (Section~I.5.7) converge absolutely.
For $k\ge1$ this was shown in subsection~3.2, for $k=0$ in subsection~3.3, by means of the majorants (F1.7) and (F2.6), which do not depend on $A$.

Put
\[
\mathrm{LZ}(A):=\sum_{X\in D_k}\acute{E}_k(X,A).
\]
The sum is finite, since $D_k$ is a finite set (subsection~0.1), and $\mathrm{LZ}$ is analytic on $\tilde{G}_k$ by (d1).

\textbf{6.1.~Real connections.}
Let $A\in\tilde{G}_k$ be real.
By W1, $\sum_XS_{k,y}(X,A)=S_{k,y}(A)$ for every $y\ge0$.
The trace is linear, and each of the finitely many integrals converges absolutely (subsections~3.2 and~3.3).
Hence
\begin{align*}
k\ge1:&\quad\mathrm{LZ}(A)=-i\cdot b_k^2\int_0^{\infty}\mathrm{d}y\,\operatorname{Tr}\left[\gamma_5\cdot B_{k,y}(A)\cdot Q_k(A)J^{*}\cdot S_{k,y}(A)\cdot JQ_k^{T}(-A)\cdot B_{k,y}(A)\right],\\
k=0:&\quad\mathrm{LZ}(A)=\int_1^{\infty}\frac{\mathrm{d}y}{y}\,\operatorname{Tr}\left[F_{0}(A)\cdot S_{0,y}(A)\right]-i\int_0^1\mathrm{d}y\,\operatorname{Tr}\left[\gamma_5\cdot S_{0,y}(A)\right],
\end{align*}
and Proposition~12 yields
\[
\delta Z_k(A)=\exp\left\{c^{\mathrm{d}}_{k}\cdot n_k+\mathrm{LZ}(A)\right\}\qquad\left(c^{\mathrm{d}}_{0}=0\right).
\]

\textbf{6.2.~Analyticity of $\delta Z_k$ on $\tilde{G}_k$.}
\emph{The case $k\ge1$.}
By Section~I.3.5, $\delta Z_k(A)=\det\left(D_k(A)+b_{L}\cdot P(A)\right)$ with $D_k(A)=b_k-b_k^2\cdot Q_k(A)J^{*}\cdot S_k(A)\cdot JQ_k^{T}(-A)$.
By Proposition~2.8\,(a), (b) at $y=\infty$, the operator $S_k(A)=S_{\infty}(A)$ exists and is analytic on $\tilde{G}_k$.
In the bases of subsection~0.4 the two formulas of Section~I.3.3 recalled in subsection~1.2 give
\begin{align*}
\left(Q_k(A)J^{*}\delta_{x',\beta}\right)(z)&=\frac{\eta_k^3}{48}\cdot\varphi_A(z,c')\cdot\mathbf{1}[c'\in\Delta_z]\,e_{\beta}&&\left(z\in\mathrm{FD},\ x'=(c',g')\in\mathrm{FDf}\right),\\
\left(JQ_k^{T}(-A)\,\delta_{z',\beta}\right)(x)&=\varphi_A(z',c)^{-1}\cdot\mathbf{1}[c\in\Delta_{z'}]\,e_{\beta}&&\left(x=(c,g)\in\mathrm{FDf},\ z'\in\mathrm{FD}\right),
\end{align*}
and the entries of $P(A)$ are those given by Proposition~5\,(a).
By (Ge4) all these entries are entire functions of $A$.
Hence the entries of $D_k(A)+b_{L}\cdot P(A)$ are analytic on $\tilde{G}_k$, and so is its determinant, which is a polynomial in the entries.

\emph{The case $k=0$.}
Here $\delta Z_0(A)=\det F_{0}(A)$ (Section~I.3.9) is the determinant of a linear map of $\mathfrak{S}^{\mathrm{ch}}(1,\ell)$, a space of dimension $192\,n_0$.
By (I.J1) and (I.J4), $J$ is injective and $JJ^{*}$ is a projection; the range of $J$ has dimension $4n_0$, the kernel of $J^{*}$ has dimension $188\,n_0$, and the determinant is taken on the whole space, without a separation of the two parts.
The entries $\frac{1}{48}F_{0}(A)(x,x')$ of its matrix are entire functions of $A$ by (F2.2) and (Ge4).

\textbf{6.3.~Continuation.}
The functions $\delta Z_k$ and $\exp\left\{c^{\mathrm{d}}_{k}\cdot n_k+\mathrm{LZ}\right\}$ are analytic on the connected open set $\tilde{G}_k$ (Do2), and by subsection~6.1 they agree at its real points.
These real points form an open subset of the real subspace of $\mathrm{Vc}$, and this subset contains $0$.
At $0$ all partial derivatives of the two functions in the real directions agree, and for a holomorphic function these are its complex partial derivatives.
So the Taylor series of the two functions at $0$ coincide, the functions agree in a neighbourhood of $0$, and by the identity theorem \cite[Chapter~I, Section~A]{GunningRossi1965} they agree on $\tilde{G}_k$:
\begin{equation}
\delta Z_k(A)=\exp\left\{c^{\mathrm{d}}_{k}\cdot n_k+\sum_{X\in D_k}\acute{E}_k(X,A)\right\}\qquad\text{for all }A\in\tilde{G}_k.
\tag{F6.1}
\end{equation}
In particular, $\delta Z_k(A)\ne0$ for every $A\in\tilde{G}_k$.

\textbf{6.4.~The quotient.}
By Proposition~10\,(e) for $k\ge1$ and by Proposition~10\,(g) for $k=0$,
\[
\delta Z_k(-A)=\delta Z_k(A),\qquad\overline{\delta Z_k(\overline{A})}=\delta Z_k(A).
\]
The number $c^{\mathrm{d}}_{k}\cdot n_k$ is real: for $k\ge1$ because $b_k>0$ and $b_k+b_{L}>0$ (Section~I.3.4), and $c^{\mathrm{d}}_{0}=0$.
Together with (F6.1) this gives
\[
\exp\mathrm{LZ}(-A)=\exp\mathrm{LZ}(A),\qquad\exp\left(\overline{\mathrm{LZ}(\overline{A})}\right)=\overline{\exp\mathrm{LZ}(\overline{A})}=\exp\mathrm{LZ}(A).
\]
Thus $\mathrm{LZ}(-A)-\mathrm{LZ}(A)$ and $\overline{\mathrm{LZ}(\overline{A})}-\mathrm{LZ}(A)$ are continuous functions on the connected set $\tilde{G}_k$ with values in the discrete set $2\pi i\mathbb{Z}$.
They are constant, equal to their values at $A=0$:
\[
\mathrm{LZ}(-A)=\mathrm{LZ}(A),\qquad\overline{\mathrm{LZ}(\overline{A})}=\mathrm{LZ}(A)+\mathrm{cz},\qquad\mathrm{cz}:=\overline{\mathrm{LZ}(0)}-\mathrm{LZ}(0),
\]
and consequently $\overline{\mathrm{LZ}(-\overline{A})}=\mathrm{LZ}(-A)+\mathrm{cz}=\mathrm{LZ}(A)+\mathrm{cz}$.
Summing the definition of $E^{d}_k$ over $X$ one finds
\[
\sum_XE^{d}_k(X,A)=\frac{1}{4}\left[\mathrm{LZ}(A)+\mathrm{LZ}(-A)+\overline{\mathrm{LZ}(\overline{A})}+\overline{\mathrm{LZ}(-\overline{A})}\right]=\mathrm{LZ}(A)+\frac{\mathrm{cz}}{2}.
\]
Therefore
\[
\sum_XE_k^{\mathrm{det}}(X,A)=\sum_X\left[E^{d}_k(X,A)-E^{d}_k(X,0)\right]=\mathrm{LZ}(A)-\mathrm{LZ}(0),
\]
and (F6.1) gives
\begin{equation}
\frac{\delta Z_k(A)}{\delta Z_k(0)}=\exp\left\{\mathrm{LZ}(A)-\mathrm{LZ}(0)\right\}=\exp\left\{\sum_{X\in D_k}E_k^{\mathrm{det}}(X,A)\right\}\qquad\text{for all }A\in\tilde{G}_k.
\tag{F6.2}
\end{equation}
This proves (d0).
\end{proof}

\subsubsection*{7. Proof of (d4)}

\begin{proof}[Proof of Proposition~13\,(d4)]
The argument is the same at all levels $k\ge0$: apart from (d1), (d2) and (d3) it concerns bond functions only, and the level enters only through the inequality $\eta_k\le1$.

Let $\ell\ge M\cdot r_k$, $X\in D_k$ and $(\mathscr{A},\mathscr{C})\in\tilde{R}_k\times\mathfrak{X}_k$.
By (Do3) the connections $\mathscr{C}$ and $\mathscr{A}+\mathscr{C}$ lie in $\tilde{G}_k$; hence
\[
{\delta E^{\mathrm{det}}}_k(X;\mathscr{A},\mathscr{C})=E^{d}_k(X,\mathscr{A}+\mathscr{C})-E^{d}_k(X,\mathscr{C})
\]
is defined.
The conditions (A10) and (Z2.s4) of Proposition~14.2\,(b) are imposed for all $v\ge-\log e_0$, so they hold at $v=v_k$.
By the first inequality of (A10), $v_k\ge9$.

\textbf{Case 1: $d_M(X)\le v_k/4$.}

\emph{Step~1 (a box that carries $X$).}
By the definition of $r_k$ (Section~G1.1), since $r\ge2$ and $v_k\ge1$, one has $r_k\ge v_k^{r}\ge v_k^{2}$, and therefore $\ell\ge M\cdot v_k^2$.
For $v\ge9$ the numbers $v^2-\frac{v}{2}-4$ and $v^2-\frac{v}{4}-4$ are positive: both are increasing functions of $v$ on $[9,\infty)$, and at $v=9$ they take the values $72.5$ and $74.75$.
Hence
\begin{equation}
\begin{aligned}
2M\cdot\left(d_M(X)+2\right)&\le M\cdot\left(\frac{v_k}{2}+4\right)<M\cdot v_k^2\le\ell,\\
M\cdot\left(d_M(X)+4\right)&\le M\cdot\left(\frac{v_k}{4}+4\right)<M\cdot v_k^2\le\ell.
\end{aligned}
\tag{F7.1}
\end{equation}
Let $\mathrm{pr}\colon\mathbb{R}^3\to\mathbb{R}^3/\ell\mathbb{Z}^3$ be the projection.
The first inequality of (F7.1) and \cite[Theorem~13.2(t10)]{QEDI} give $X=\mathrm{pr}(\tilde{X})$, where $\tilde{X}$ is a finite union of $M$-cubes of $\mathbb{R}^3$ (centred at points of $M\mathbb{Z}^3$) that is connected through faces, $\mathrm{pr}$ is injective on $\tilde{X}$, and $d_M(\tilde{X})=d_M(X)$.
By \cite[Theorem~13.2(t7)]{QEDI}, in its form for finite unions of $M$-cubes of $\mathbb{R}^3$ connected through faces, with $d_M$ formed with nets in $\mathbb{R}^3$, any two points of $\tilde{X}$ have sup distance at most $M\cdot\left(d_M(\tilde{X})+2\right)=M\cdot\left(d_M(X)+2\right)$.
The bound extends to all points of the closed cubes of $\tilde{X}$, because the sup distance is continuous and each closed cube is the closure of the open cube of side $M$ with the same centre, which is contained in $\tilde{X}$.
For $\mu=1,2,3$ let $[\mathrm{am}_\mu,\mathrm{ap}_\mu]$ be the smallest closed interval that contains the $\mu$-th coordinates of all points of the closed cubes of $\tilde{X}$; then $\mathrm{ap}_\mu-\mathrm{am}_\mu\le M\cdot\left(d_M(X)+2\right)$.
Put
\[
\mathrm{Wb}:=\left\{c\in\mathbb{L}_k:\ \mathrm{am}_\mu-\eta_k\le c_\mu\le\mathrm{ap}_\mu+\eta_k,\ \mu=1,2,3\right\},\qquad s_W:=M\cdot\left(d_M(X)+3\right).
\]
The set $\mathrm{Wb}$ is a product of three intervals of $\eta_k\mathbb{Z}$.
The length of each of them, that is, its largest element minus its smallest one, is at most $M\cdot\left(d_M(X)+2\right)+2\eta_k\le s_W$, because $\eta_k\le1$ and $M\ge2$.
So the sup-diameter of $\mathrm{Wb}$, increased by $\eta_k$, is at most $s_W+1$, and by (F7.1)
\[
s_W+1<M\cdot\left(d_M(X)+4\right)<\ell.
\]
Hence $\mathrm{Wb}$ is a lift of $\mathrm{pr}(\mathrm{Wb})$ in the sense of Section~I.1.1: $\mathrm{pr}$ is injective on $\mathrm{Wb}$, and the bonds of $\mathrm{pr}(\mathrm{Wb})$ are the images of the bonds of $\mathrm{Wb}$.
Moreover $\mathrm{Wb}$ contains the sites of $\mathbb{L}_k$ that lie in $\tilde{X}$, together with all their nearest neighbours.
Consequently the following hold.
\begin{itemize}
\item[(b1)] Every bond $(c,c')$ of $T_k$ that meets $X$ is the image of exactly one bond $(\tilde{c},\tilde{c}')$ of $\mathrm{Wb}$.
Indeed, one end point, say $c$, lies in $X$; let $\tilde{c}$ be the site of $\tilde{X}$ over $c$; then $c'=c\pm\eta_ke_{\mu}$ modulo $\ell\mathbb{Z}^3$, and $\tilde{c}':=\tilde{c}\pm\eta_ke_{\mu}$ lies in $\mathrm{Wb}$.
There is no second such bond, since $\mathrm{pr}$ is injective on $\mathrm{Wb}$.
\item[(b2)] $\mathrm{pr}$ maps every plaquette of $\mathrm{Wb}$, that is, every plaquette with its four corners in $\mathrm{Wb}$, to a plaquette of $T_k$.
\end{itemize}

\emph{Step~2 (curvature of the real part).}
Let $\mathscr{A}=A_0+A_1$ be a splitting as in the definition of $\tilde{R}_k$ (Section~I.5.3; \cite[Definition~5.16(b)]{QEDI}).
Thus $|A_1|_{\infty}<e_k^{-\frac{3}{4}+\varepsilon}$, the bond function $A_0$ is real, and for every $M$-cube $\square$ there are a real function $\lambda$ on the sites of $\square^{\dagger}$ and a real function $A'$ on the bonds of $\square^{\dagger}$ such that $A_0=A'+d\lambda$ on $\square^{\dagger}$ and $|\partial A'|_{\infty}<e_k^{-\frac{3}{4}+2\varepsilon}$ on $\square^{\dagger}$.
Let $p=(c;\mu\nu)$ be a plaquette of $T_k$, and let $\square$ be the $M$-cube that contains $c$ \cite[Section~17.4.1(g0)]{QEDI}.
The four corners of $p$ have sup distance at most $\eta_k\le1<M$ from $c$; hence they lie in $\square^{\sim1}$, a subset of $\square^{\dagger}=\square^{\sim L}$.
The set $\square^{\dagger}$ does not wrap, since $(2L+1)M+\eta_k<13LM\le\ell$.
On $p$ one has $d(d\lambda)=0$, and therefore
\[
(dA_0)_{\mu\nu}(c)=(dA')_{\mu\nu}(c)=(\partial_\mu A'_\nu)(c)-(\partial_\nu A'_\mu)(c),
\]
where both difference quotients are formed inside $\square^{\dagger}$.
It follows that
\begin{equation}
|dA_0|_{\infty}<2e_k^{-\frac{3}{4}+2\varepsilon}.
\tag{F7.2}
\end{equation}

\emph{Step~3 (axial gauge).}
The function $\tilde{A}_0(\tilde{c},\tilde{c}'):=A_0(\mathrm{pr}\,\tilde{c},\mathrm{pr}\,\tilde{c}')$ is a real antisymmetric function on the bonds of $\mathrm{Wb}$, and $(d\tilde{A}_0)(p)=(dA_0)(\mathrm{pr}\,p)$ for the plaquettes $p$ of $\mathrm{Wb}$, by (b2).
Fix a site $o$ of $\mathrm{Wb}$, and for $\tilde{c}\in\mathrm{Wb}$ let $\boldsymbol{\Gamma}(o,\tilde{c})$ be the rectilinear path from $o$ to $\tilde{c}$ that changes the coordinates in the order $1,2,3$.
By \cite[Lemma~16.26]{QEDI}, applied to the box $\mathrm{Wb}$, the number $s_W$ and the function $\tilde{A}_0$ (the coordinates numbered $0,1,2$ there are those numbered $1,2,3$ here), the function $\tilde{\boldsymbol{\lambda}}(\tilde{c}):=\eta_k\cdot\tilde{A}_0\left(\boldsymbol{\Gamma}(o,\tilde{c})\right)$ is real, and
\[
\left|\left(\tilde{A}_0-d\tilde{\boldsymbol{\lambda}}\right)(\tilde{c},\tilde{c}')\right|\le2s_W\cdot\sup_p\left|(d\tilde{A}_0)(p)\right|<4s_W\cdot e_k^{-\frac{3}{4}+2\varepsilon}\qquad\text{for every bond }(\tilde{c},\tilde{c}')\text{ of }\mathrm{Wb};
\]
the second inequality follows from (F7.2).

Let $\boldsymbol{\lambda}$ be the real function on the sites of $T_k$ with $\boldsymbol{\lambda}(\mathrm{pr}\,\tilde{c}):=\tilde{\boldsymbol{\lambda}}(\tilde{c})$ for $\tilde{c}\in\mathrm{Wb}$ and $\boldsymbol{\lambda}:=0$ outside $\mathrm{pr}(\mathrm{Wb})$; it is well defined because $\mathrm{pr}$ is injective on $\mathrm{Wb}$.
For a bond $b=(c,c')$ that meets $X$, with its lift $(\tilde{c},\tilde{c}')$ from (b1), one has
\[
A_0(b)=\tilde{A}_0(\tilde{c},\tilde{c}'),\qquad(d\boldsymbol{\lambda})(b)=\frac{\boldsymbol{\lambda}(c')-\boldsymbol{\lambda}(c)}{\eta_k}=\frac{\tilde{\boldsymbol{\lambda}}(\tilde{c}')-\tilde{\boldsymbol{\lambda}}(\tilde{c})}{\eta_k}=(d\tilde{\boldsymbol{\lambda}})(\tilde{c},\tilde{c}').
\]
Since $s_W\le M\cdot\left(\frac{v_k}{4}+3\right)$ and $e_k^{\varepsilon}=\exp(-\varepsilon v_k)$, the second inequality of (A10) gives
\[
8s_W\cdot e_k^{\varepsilon}\le8M\cdot\left(\frac{v_k}{4}+3\right)\cdot\exp(-\varepsilon v_k)\le1.
\]
Hence
\begin{equation}
\left|(A_0-d\boldsymbol{\lambda})(b)\right|<4s_W\cdot e_k^{\varepsilon}\cdot e_k^{-\frac{3}{4}+\varepsilon}\le\frac{1}{2}\,e_k^{-\frac{3}{4}+\varepsilon}\qquad\text{for every bond }b\text{ that meets }X.
\tag{F7.3}
\end{equation}

\emph{Step~4 (a connection in $G_k$).}
Let $\mathrm{Bf}$ be the bond function with $\mathrm{Bf}:=A_0-d\boldsymbol{\lambda}+A_1$ on the bonds that meet $X$ and $\mathrm{Bf}:=0$ on all other bonds.
By (F7.3), $|\mathrm{Bf}|_{\infty}<\frac{3}{2}\,e_k^{-\frac{3}{4}+\varepsilon}$.
Let $t$ be a complex number with $|t|\le\frac{1}{3}\,e_k^{-\frac{1}{4}}$.
Since $|\mathscr{C}|_{\infty}<e_k^{-1+2\varepsilon}$ (definition of $\mathfrak{X}_k$, Section~I.5.3) and $e_k^{\varepsilon}\le\frac{1}{4}$ (first inequality of (Z2.s4)),
\[
\left|t\,\mathrm{Bf}+\mathscr{C}\right|_{\infty}<\frac{1}{2}\,e_k^{-1+\varepsilon}+e_k^{\varepsilon}\cdot e_k^{-1+\varepsilon}\le\frac{3}{4}\,e_k^{-1+\varepsilon}.
\]
So $t\,\mathrm{Bf}+\mathscr{C}$ lies in $G_k$ (Section~I.5.3), and $G_k\subset\tilde{G}_k$ by (Do2).
By (d1) the function
\[
\mathrm{fd}(t):=E^{d}_k\left(X,t\,\mathrm{Bf}+\mathscr{C}\right)
\]
is analytic on the disk $\left\{|t|<\frac{1}{3}\,e_k^{-\frac{1}{4}}\right\}$, as the composition of an analytic function with an affine map, and it is bounded there by $C_{d}\cdot M^3\cdot\exp\left(-(\kappa+1)\cdot d_M(X)\right)$.

\emph{Step~5 (identification).}
The connection $\mathscr{A}+\mathscr{C}-d\boldsymbol{\lambda}$ lies in $\tilde{G}_k$ (Do2) and coincides with $\mathrm{Bf}+\mathscr{C}$ on the bonds that meet $X$.
The point $t=1$ belongs to the disk of Step~4, because $\frac{1}{3}\,e_k^{-\frac{1}{4}}\ge\frac{1}{3}\exp\left(\frac{9}{4}\right)>1$.
By (d3), and then by (d2),
\[
E^{d}_k(X,\mathscr{A}+\mathscr{C})=E^{d}_k(X,\mathscr{A}+\mathscr{C}-d\boldsymbol{\lambda})=E^{d}_k(X,\mathrm{Bf}+\mathscr{C})=\mathrm{fd}(1),\qquad E^{d}_k(X,\mathscr{C})=\mathrm{fd}(0).
\]

\emph{Step~6 (Cauchy's formula).}
Put
\[
\mathrm{rad}:=\frac{1}{4}\,e_k^{-\frac{1}{4}}=\frac{1}{4}\exp\left(\frac{v_k}{4}\right)\ge\frac{1}{4}\exp\left(\frac{9}{4}\right)>2;
\]
the last inequality holds because the first five terms of the exponential series at $\frac{9}{4}$ add up to $1+2.25+2.53125+1.8984375+1.0678\ldots>8$.
The points $0$ and $1$ lie inside the circle $|t|=\mathrm{rad}$, and this circle lies in the disk of Step~4.
From $\frac{1}{t(t-1)}=\frac{1}{t-1}-\frac{1}{t}$ and Cauchy's formula,
\[
\mathrm{fd}(1)-\mathrm{fd}(0)=\frac{1}{2\pi i}\oint_{|t|=\mathrm{rad}}\mathrm{d}t\,\frac{\mathrm{fd}(t)}{t(t-1)},
\]
and therefore
\[
\left|\mathrm{fd}(1)-\mathrm{fd}(0)\right|\le\frac{1}{2\pi}\cdot2\pi\,\mathrm{rad}\cdot\frac{\sup|\mathrm{fd}|}{\mathrm{rad}\cdot(\mathrm{rad}-1)}=\frac{\sup|\mathrm{fd}|}{\mathrm{rad}-1}\le\frac{2\sup|\mathrm{fd}|}{\mathrm{rad}}=8e_k^{\frac{1}{4}}\cdot\sup|\mathrm{fd}|.
\]
Consequently
\[
\left|{\delta E^{\mathrm{det}}}_k(X;\mathscr{A},\mathscr{C})\right|\le8C_{d}\cdot M^3\cdot e_k^{\frac{1}{4}}\cdot\exp\left(-(\kappa+1)\cdot d_M(X)\right)\le8C_{d}\cdot M^3\cdot e_k^{\frac{1}{4}}\cdot\exp\left(-\kappa\cdot d_M(X)\right).
\]

\textbf{Case 2: $d_M(X)>v_k/4$.}
Then $\exp\left(-d_M(X)\right)<\exp\left(-\frac{v_k}{4}\right)=e_k^{\frac{1}{4}}$.
By (d1), applied to each of the two terms,
\[
\left|{\delta E^{\mathrm{det}}}_k(X;\mathscr{A},\mathscr{C})\right|\le2C_{d}\cdot M^3\cdot\exp\left(-(\kappa+1)\cdot d_M(X)\right)<2C_{d}\cdot M^3\cdot e_k^{\frac{1}{4}}\cdot\exp\left(-\kappa\cdot d_M(X)\right).
\]

In both cases the bound of (d4) holds.
For the last assertion of (d4), note that $0\in\mathfrak{X}_k$ and that $E_k^{\mathrm{det}}(X,\mathscr{A})=E^{d}_k(X,\mathscr{A})-E^{d}_k(X,0)={\delta E^{\mathrm{det}}}_k(X;\mathscr{A},0)$.
\end{proof}

\subsubsection*{8. Proof of (d5)}

\begin{proof}[Proof of Proposition~13\,(d5)]
Let the data be (SF+).
The level $k+1$ lives on the torus $T_{k+1}=\mathbb{L}_{k+1}/(\ell/L)\mathbb{Z}^3$, whose side $\ell/L$ belongs to $\mathrm{Cl}$, and its charge is $e_{k+1}=L^{\frac{1}{2}}\cdot e_k\le e_0$.
The argument is the same at every level $k\ge0$ (for $k=0$ it concerns the first step),
because it rests only on (d1)--(d3), on the bound (F9.1), which is derived from (d4), and on facts about sets of photon sites and about bond functions.

\paragraph{8.1. The two families.}
Since $\eta_{k+1}=\eta_k/L$, the map $c\mapsto c/L$ is a bijection of $T_k$ onto $T_{k+1}$.
For a set $Y$ of sites of $T_{k+1}$ one writes $LY:=\{c\in T_k:\ c/L\in Y\}$.
The map $b\mapsto b/L$ carries the bonds of $T_k$ onto those of $T_{k+1}$, and $A_L(b)=L^{-\frac{1}{2}}\cdot A(b/L)$ (Section~I.3.7).
The values of ${\delta E^{\mathrm{det}}}_k(X;\cdot,\cdot)$ and of $E_k^{\mathrm{det}}(X,\cdot\,)$ are numbers, that is, elements of degree zero of $\mathrm{Gr}[\psi,\chi]$,
and the homomorphism $S_L$ leaves them fixed.
Therefore Sections~I.3.7 and~I.5.6 (the operation $L$ on a family with two slots) and Section~I.5.11 (the operation $L$ on a family without slot) give, for every $Y\in D_{k+1}$,
\begin{equation}
\tag{F8.1}
\begin{aligned}
\left(L\,{\delta E^{\mathrm{det}}}_k\right)(Y;\mathscr{A},\mathscr{C})&=\sum_{X\in D_k:\ \bar{X}^{L}=LY}{\delta E^{\mathrm{det}}}_k(X;\mathscr{A}_L,\mathscr{C}_L),\\
\left(L\,E_k^{\mathrm{det}}\right)(Y;\mathscr{A})&=\sum_{X\in D_k:\ \bar{X}^{L}=LY}E_k^{\mathrm{det}}(X,\mathscr{A}_L)=\left(L\,{\delta E^{\mathrm{det}}}_k\right)(Y;\mathscr{A},0),
\end{aligned}
\end{equation}
where both sums are finite.
Both families consist of even elements without legs.
The number of the barred legs equals that of the unbarred legs, namely zero,
and the only kernel is the one with $n=m=0$, which is the number itself.

\emph{Domains.}
Let $(\mathscr{A},\mathscr{C})\in\tilde{R}_{k+1}\times\mathfrak{X}_{k+1}$.
As $\ell$ and $\ell/L$ lie in $\mathrm{Cl}$, $e_{k+1}=L^{\frac{1}{2}}\cdot e_k$ and $e_k\le e_0$, \cite[Lemma~17.5(a), (c)]{QEDI} applies:
$\mathscr{A}_L$ belongs to $L^{-\frac{3}{4}}\cdot\tilde{R}_k$, and $\mathscr{C}_L$ belongs to $\mathfrak{X}_k$.
Moreover $L^{-\frac{3}{4}}\cdot\tilde{R}_k\subset\tilde{R}_k$, since the thresholds in the definition of $\mathrm{cf}\cdot\tilde{R}_k$ increase with $\mathrm{cf}$.
By (Do3) the connections $\mathscr{A}_L+\mathscr{C}_L$ and $\mathscr{C}_L$ lie in $\tilde{G}_k$, so that every term of (F8.1) is defined.
At the level $k+1$ the sets $\tilde{R}_{k+1}$ and $\mathfrak{X}_{k+1}$ are open \cite[Lemma~17.3]{QEDI},
and they are invariant under $\mathbb{G}_{k+1}$ and under the maps $A\mapsto-A$ and $A\mapsto\overline{A}$ \cite[Lemma~17.6]{QEDI}.

\paragraph{8.2. The clauses (k1)--(k5), (k7) of Section~I.5.4 at the level $k+1$, for $L\,{\delta E^{\mathrm{det}}}_k$.}
\begin{itemize}
\item[(k1)]
The maps $(\mathscr{A},\mathscr{C})\mapsto\mathscr{A}_L+\mathscr{C}_L$ and $(\mathscr{A},\mathscr{C})\mapsto\mathscr{C}_L$ are linear,
and by subsection~8.1 they send $\tilde{R}_{k+1}\times\mathfrak{X}_{k+1}$ into $\tilde{G}_k$.
By (d1) the function $E^{d}_k(X,\cdot\,)$ is analytic on $\tilde{G}_k$.
Consequently every term of (F8.1), and hence the finite sum, is an analytic function on $\tilde{R}_{k+1}\times\mathfrak{X}_{k+1}$.
\item[(k2)]
The families have no legs, so the condition on the supports is void.
Let $(\mathscr{A},\mathscr{C})$ and $\left(\tilde{\mathscr{A}},\tilde{\mathscr{C}}\right)$ be pairs in $\tilde{R}_{k+1}\times\mathfrak{X}_{k+1}$ that coincide on the bonds of $T_{k+1}$ that meet $Y$,
and let $X\in D_k$ satisfy $\bar{X}^{L}=LY$.
Then $X\subset\bar{X}^{L}=LY$.
Indeed, by fact (g0) of \cite[Section~17.4.1]{QEDI} every $M$-cube of $X$ lies in an $LM$-cube,
and this $LM$-cube belongs to $\bar{X}^{L}$, which by definition \cite[Section~13.1.7]{QEDI} is the union of all $LM$-cubes in which at least one $M$-cube of $X$ lies.
So a bond $b$ of $T_k$ that meets $X$ has an end $c$ in $LY$;
then $c/L$ is an end of $b/L$ and lies in $Y$, that is, $b/L$ meets $Y$, and
\[
\mathscr{A}_L(b)=L^{-\frac{1}{2}}\cdot\mathscr{A}(b/L)=\tilde{\mathscr{A}}_L(b),\qquad\mathscr{C}_L(b)=\tilde{\mathscr{C}}_L(b).
\]
Now (d2) gives
\[
E^{d}_k\left(X,\mathscr{A}_L+\mathscr{C}_L\right)=E^{d}_k\left(X,\tilde{\mathscr{A}}_L+\tilde{\mathscr{C}}_L\right),\qquad E^{d}_k\left(X,\mathscr{C}_L\right)=E^{d}_k\left(X,\tilde{\mathscr{C}}_L\right),
\]
and every term of (F8.1) takes the same value at the two pairs.
\item[(k3)]
The automorphism $a_{\tilde{g}'}$ of Section~I.4.5 fixes numbers.
Let $\tilde{g}'\in\mathbb{G}_{k+1}$.
Then $\tilde{g}'Y\in D_{k+1}$:
by Proposition~16.1\,(a1) at the level $k+1$, the set $\tilde{g}'Y$ is a polymer of $T_{k+1}$,
the side $\ell/L$ lying in $\mathrm{Cl}$ by (SF+) and $e_{k+1}$ being positive.
By Proposition~16.1\,(d1) there is exactly one $\tilde{g}=(S,a)\in\mathbb{G}_k^{L}$ with $\tilde{g}'=(S,a/L)$.
Proposition~16.1\,(d2) gives $L\left(\tilde{g}'Y\right)=\tilde{g}(LY)$, and Proposition~16.1\,(d3) gives
\[
\left(\tilde{g}'\mathscr{A}\right)_L=\tilde{g}(\mathscr{A}_L),\qquad\left(\tilde{g}'\mathscr{C}\right)_L=\tilde{g}(\mathscr{C}_L).
\]
By Proposition~16.1\,(a1) the map $X\mapsto\tilde{g}X$ is a bijection of $D_k$ with $\overline{\tilde{g}X}^{L}=\tilde{g}\left(\bar{X}^{L}\right)$;
hence it maps the set $\left\{X:\ \bar{X}^{L}=LY\right\}$ onto the set $\left\{X':\ \bar{X}'^{L}=\tilde{g}(LY)=L\left(\tilde{g}'Y\right)\right\}$.
The map $\tilde{g}$ is linear on bond functions.
From (F8.1) and these relations, and then from (d3), one obtains
\begin{align*}
\left(L\,{\delta E^{\mathrm{det}}}_k\right)\left(\tilde{g}'Y;\tilde{g}'\mathscr{A},\tilde{g}'\mathscr{C}\right)
&=\sum_{X:\ \bar{X}^{L}=LY}\left[E^{d}_k\left(\tilde{g}X,\tilde{g}(\mathscr{A}_L+\mathscr{C}_L)\right)-E^{d}_k\left(\tilde{g}X,\tilde{g}(\mathscr{C}_L)\right)\right]\\
&=\sum_{X:\ \bar{X}^{L}=LY}{\delta E^{\mathrm{det}}}_k(X;\mathscr{A}_L,\mathscr{C}_L)=\left(L\,{\delta E^{\mathrm{det}}}_k\right)(Y;\mathscr{A},\mathscr{C}).
\end{align*}
\item[(k4)]
The automorphism $G_{\boldsymbol{\lambda}+\boldsymbol{\lambda}'}$ fixes numbers.
Let $(\mathscr{A},\mathscr{C})$ and $\left(\mathscr{A}^{1},\mathscr{C}^{1}\right)$ lie in $\tilde{R}_{k+1}\times\mathfrak{X}_{k+1}$,
let $\boldsymbol{\lambda}$, $\boldsymbol{\lambda}'$ be real functions on the sites of $T_{k+1}$,
and assume that $\mathscr{A}^{1}=\mathscr{A}-d\boldsymbol{\lambda}$ and $\mathscr{C}^{1}=\mathscr{C}-d\boldsymbol{\lambda}'$ on the bonds at distance at most $1$ from $Y$.
Put $\boldsymbol{\lambda}_L(c):=L^{\frac{1}{2}}\cdot\boldsymbol{\lambda}(c/L)$ and $\boldsymbol{\lambda}'_L(c):=L^{\frac{1}{2}}\cdot\boldsymbol{\lambda}'(c/L)$.
These are real functions of $c$ on $T_k$, and by Proposition~16.1\,(d4)
\[
(d\boldsymbol{\lambda})_L=d(\boldsymbol{\lambda}_L),\qquad(d\boldsymbol{\lambda}')_L=d(\boldsymbol{\lambda}'_L).
\]
Let $\bar{X}^{L}=LY$, and let $b$ be a bond of $T_k$ that meets $X$.
As in (k2), $b/L$ has an end in $Y$;
the other end of $b/L$ has the sup distance $\eta_{k+1}\le1$ from it, so that $b/L$ lies at distance at most $1$ from $Y$.
It follows that
\[
\mathscr{A}^{1}_L(b)=\mathscr{A}_L(b)-(d\boldsymbol{\lambda}_L)(b),\qquad\mathscr{C}^{1}_L(b)=\mathscr{C}_L(b)-(d\boldsymbol{\lambda}'_L)(b).
\]
The connections $\mathscr{A}^{1}_L+\mathscr{C}^{1}_L$ and $\mathscr{A}_L+\mathscr{C}_L-d(\boldsymbol{\lambda}_L+\boldsymbol{\lambda}'_L)$ lie in $\tilde{G}_k$ (subsection~8.1 and (Do2)),
and they coincide on the bonds that meet $X$;
the same holds for $\mathscr{C}^{1}_L$ and $\mathscr{C}_L-d\boldsymbol{\lambda}'_L$.
By (d2) and then (d3),
\begin{gather*}
E^{d}_k\left(X,\mathscr{A}^{1}_L+\mathscr{C}^{1}_L\right)=E^{d}_k\left(X,\mathscr{A}_L+\mathscr{C}_L-d(\boldsymbol{\lambda}_L+\boldsymbol{\lambda}'_L)\right)=E^{d}_k\left(X,\mathscr{A}_L+\mathscr{C}_L\right),\\
E^{d}_k\left(X,\mathscr{C}^{1}_L\right)=E^{d}_k\left(X,\mathscr{C}_L-d\boldsymbol{\lambda}'_L\right)=E^{d}_k\left(X,\mathscr{C}_L\right).
\end{gather*}
Thus every term of (F8.1) has the same value at $\left(\mathscr{A}^{1},\mathscr{C}^{1}\right)$ as at $(\mathscr{A},\mathscr{C})$.
\item[(k5)]
The automorphism $\mathrm{Cg}$ fixes numbers, one has $(-\mathscr{A})_L=-\mathscr{A}_L$, and $E^{d}_k(X,-A)=E^{d}_k(X,A)$ by (d3).
Hence
\[
\left(L\,{\delta E^{\mathrm{det}}}_k\right)(Y;-\mathscr{A},-\mathscr{C})=\left(L\,{\delta E^{\mathrm{det}}}_k\right)(Y;\mathscr{A},\mathscr{C}).
\]
\item[(k7)]
The automorphism $\mathbf{K}$ acts on numbers by complex conjugation, and $\left(-\overline{\mathscr{A}}\right)_L=-\overline{\mathscr{A}_L}$.
By (d3),
\[
\overline{E^{d}_k\left(X,-\overline{A}\right)}=\overline{E^{d}_k\left(X,\overline{A}\right)}=E^{d}_k(X,A).
\]
Therefore
\[
\overline{\left(L\,{\delta E^{\mathrm{det}}}_k\right)\left(Y;-\overline{\mathscr{A}},-\overline{\mathscr{C}}\right)}=\left(L\,{\delta E^{\mathrm{det}}}_k\right)(Y;\mathscr{A},\mathscr{C}).
\]
\end{itemize}
The clause (k6) follows from the bound (F9.1) below.
So $L\,{\delta E^{\mathrm{det}}}_k$ belongs to $K^{x,C}_{k+1}$, and, by (k7), to $K^{x}_{k+1}$.

\paragraph{8.3. $L\,E_k^{\mathrm{det}}$.}
At the level $k+1$ the clauses (k1)$^{\mathrm{o}}$--(k5)$^{\mathrm{o}}$, (k7)$^{\mathrm{o}}$ of Section~I.5.11 are the clauses (k1)--(k5), (k7) for a family on $\tilde{R}_{k+1}$ alone
(Section~I.5.4: the definition of $K_k^{C}$, $K_k$ and $\|\cdot\|_k$).
By the last equality of (F8.1) they are the statements proved in subsection~8.2, taken at $\mathscr{C}=\tilde{\mathscr{C}}=\mathscr{C}^{1}=0$ and $\boldsymbol{\lambda}'=0$;
here $0\in\mathfrak{X}_{k+1}$, $\tilde{g}'0=0$ and $-0=\overline{0}=0$.
The clause (k6)$^{\mathrm{o}}$ follows from (F9.1) at $\mathscr{C}=0$.
Hence $L\,E_k^{\mathrm{det}}$ belongs to $K_{k+1}$.
\end{proof}

\subsubsection*{9. Proof of (d6)}

\begin{proof}[Proof of Proposition~13\,(d6)]
Let the data be (SF+).
For a number the norm $\|\cdot\|_{h_{k+1}}$ is its modulus, since by (I.5.1) it is the modulus of the coefficient of the empty monomial.
Let $Y\in D_{k+1}$ and $(\mathscr{A},\mathscr{C})\in\tilde{R}_{k+1}\times\mathfrak{X}_{k+1}$.
By subsection~8.1 the pair $(\mathscr{A}_L,\mathscr{C}_L)$ lies in $\tilde{R}_k\times\mathfrak{X}_k$, and $\ell\ge M\cdot r_k$;
so (d4) applies to every term of (F8.1).
The sides $\ell$ and $\ell/L$ belong to $\mathrm{Cl}$, and the rate $\kappa$ satisfies $\kappa\ge\kappa_0$ (line (C2)).
Hence (d4) and the assertion (t8) of \cite[Theorem~13.2]{QEDI}, applied with the rate $\kappa$, give
\begin{align}
\exp\left(\kappa\cdot d_M(Y)\right)\cdot\left|\left(L\,{\delta E^{\mathrm{det}}}_k\right)(Y;\mathscr{A},\mathscr{C})\right|
&\le\exp\left(\kappa\cdot d_M(Y)\right)\cdot8\,C_d\cdot M^3\cdot e_k^{\frac{1}{4}}\sum_{X:\ \bar{X}^{L}=LY}\exp\left(-\kappa\cdot d_M(X)\right)\notag\\
&\le8\,C_d\cdot M^3\cdot e_k^{\frac{1}{4}}\cdot c_{\mathrm{rb}}\cdot L^3\cdot\exp\left(-\left(L\cdot(\kappa-\kappa_0-1)-\kappa\right)\cdot d_M(Y)\right)\notag\\
&\le8\,c_{\mathrm{rb}}\cdot L^3\cdot C_d\cdot M^3\cdot e_k^{\frac{1}{4}},\tag{F9.1}
\end{align}
where the last step uses $L\cdot(\kappa-\kappa_0-1)\ge\kappa$ and $d_M(Y)\ge0$.
At $v=v_k$ the condition (Z2.s5) reads $8\,c_{\mathrm{rb}}\cdot L^3\cdot C_d\cdot M^3\cdot e_k^{2\varepsilon}\le1$.
So the right side of (F9.1) does not exceed $e_k^{\frac{1}{4}-2\varepsilon}$.
Moreover $e_k\le L^{\frac{1}{2}}\cdot e_k=e_{k+1}$, and $\frac{1}{4}-2\varepsilon\ge0$ because $\varepsilon\le\bar{\varepsilon}=\frac{1}{200}$.
Forming the supremum over $(\mathscr{A},\mathscr{C})$ and $Y$, as in the definition of the norm in Section~I.5.4, one obtains
\[
\left\|L\,{\delta E^{\mathrm{det}}}_k\right\|^{x}_{k+1}\le e_k^{\frac{1}{4}-2\varepsilon}\le e_{k+1}^{\frac{1}{4}-2\varepsilon}.
\]
By (F8.1), $\left(L\,E_k^{\mathrm{det}}\right)(Y;\mathscr{A})=\left(L\,{\delta E^{\mathrm{det}}}_k\right)(Y;\mathscr{A},0)$, and $0\in\mathfrak{X}_{k+1}$.
Thus the supremum that defines $\left\|L\,E_k^{\mathrm{det}}\right\|_{k+1}$ ranges over a subset of the pairs $(\mathscr{A},\mathscr{C})$, and
\[
\left\|L\,E_k^{\mathrm{det}}\right\|_{k+1}\le\left\|L\,{\delta E^{\mathrm{det}}}_k\right\|^{x}_{k+1}.
\]
\end{proof}

\subsubsection*{10. Uniformity}

The bounds of (d1), (d4) and (d6) contain only the numbers $C_d$, $M$, $\kappa$, $\varepsilon$ and the datum $e_k$ (respectively $e_{k+1}$),
and $C_d$ is a function of $(c_X,C_w,b)$.
In the proofs of (d0)--(d6) the data and the variables enter as follows.
\begin{itemize}
\item The level $k$ enters through $b_{-}\le b_k\le b$, through $\eta_k\le1$, and through the two forms of the integrand (Sections~1 and~2).
\item The side $\ell$ enters through $\ell\in\mathrm{Cl}$
(fact (g0); the cubes $\square^{\dagger}$ do not wrap; the assertions (t1), (t7), (t8), (t10) of \cite[Theorem~13.2]{QEDI}; \cite[Lemma~17.5]{QEDI}),
through $\ell\ge M\cdot r_k$ in (F7.1), and through $\ell/L\in\mathrm{Cl}$ in Sections~8 and~9.
No bound contains $\ell$ or $n_k$.
\item The mass enters through $\mu\le1$ in (F2.3), and through Propositions~2.8, 10 and~12, which hold for all $\mu\in(0,1]$.
\item The charge enters through the conditions of Proposition~14.2\,(b) at $v=v_k$, which hold for every $e_k\in(0,e_0]$.
\item The family (p5) enters through the properties W1, W2, W7, W12, W13, W14 of Section~G1.11 for $\mathrm{par}$ and through the Borel measurability in $y$.
Of $\mathrm{par}$ only $C_w$ and $\kappa_{\mathrm{w}}$ occur in the bounds; $\gamma_{\mathrm{w}}$ enters through $\gamma_{\mathrm{w}}\ge0$ alone.
\item The majorants (F1.7) and (F2.6) do not contain $A$; the variable $y$ is integrated; and $X$ enters through $d_M(X)$.
\end{itemize}
Neither $N$ nor the torus of the flow is a datum of (SF) or (SF+); hence no bound contains them.
A twist is a real constant connection in the first slot (Section~I.6.6).
The sets $\tilde{G}_k$ and $\tilde{R}_k$ are invariant under the addition of real constant connections,
by the property (G2) of \cite[Section~16.2.7]{QEDI} and by \cite[Lemma~17.6]{QEDI},
and all clauses hold on the whole of $\tilde{G}_k$, respectively of $\tilde{R}_k\times\mathfrak{X}_k$.
Therefore all bounds hold uniformly in the twist.

\subsection*{Conditions on the constants}

Each condition on the constants that enters the proofs of this chapter is listed below,
together with the places where it is used and with its place in Section~I.8.2 or in Proposition~14.2;
all of them are among the conditions imposed there.
\begin{itemize}
\item $\rho<1$; $w=(6,6\sqrt{2},6\sqrt{2},6)$.
\emph{Used in:} the numerals $14.49$, $6$, $20.49$, $6$ in (F2.3);
moreover, the weights $w$ satisfy the closure relation (I.1.1), which in the form (I.m1) enters (F2.1) and hence (F2.2) and (F2.5):
by it the diagonal term of $F_0(A)$ is the scalar $\frac{\rho}{2}\,\mathrm{sw}+\mu$, and the numeral $14.49$ of (F2.3) bounds its part $\frac{\rho}{2}\,\mathrm{sw}$.
\emph{Place:} line (C0).
\item $b_k\le b$; $b_L=b/L\le b$.
\emph{Used in:} (F3.1); (F2.3).
\emph{Place:} Section~I.3.4; line (C4); line (C0), where $b:=1$, so that $b>0$.
\item $b_k\ge b_{-}=b\cdot\left(1-\frac{1}{L}\right)>0$; $b_L>0$.
\emph{Used in:} subsection~1.1 ($b_k^2+y^2>0$ at $y=0$; the norm of $\mathrm{bq}_2$ is at most that of $\mathrm{bq}_1$);
subsection~3.2 (the value $\pi/(2b_k)$ of the integral);
subsection~6.4 ($c^{\mathrm{d}}_{k}$ is real).
\emph{Place:} Section~I.3.4; line (C4), (L5); line (C0), where $b:=1$, so that $b>0$.
\item $L$ odd; $M$ a power of $L$, hence odd; $M\ge L^2$, hence $M\ge2$.
\emph{Used in:} (g0), (Ge1)--(Ge3); $s_W$ in Section~7, Step~1.
\emph{Place:} (L5), (M1).
\item $L\ge1$ (from $L\ge5$): $\eta_k\le1$, $\eta_{k+1}\le1$; $L^{-\frac{3}{4}}\le1$; $e_k\le L^{\frac{1}{2}}\cdot e_k$.
\emph{Used in:} Section~7, Steps~1 and~2; subsection~8.2, clause (k4); subsection~8.1; Section~9, last step.
\emph{Place:} (L5).
\item (Z2.s1): $6L\cdot\exp(-\varepsilon v)\le1$.
\emph{Used in:} (Do1), that is $3\,\mathfrak{i}(A)\le1$ in (F1.5), and $3(L-1)\cdot\mathfrak{i}(A)\le1$ and $\exp\left(\mathfrak{i}(A)\right)\le\mathrm{e}$ in (F2.3);
it is also a condition of Proposition~5\,(a).
\emph{Place:} Proposition~14.2\,(b), line (C8).
\item $\gamma_{\mathrm{w}}\ge0$.
\emph{Used in:} (F1.5) and (F2.6), where the factor $\exp\left(-\gamma_{\mathrm{w}}\cdot d\right)$ is dropped.
\emph{Place:} $\gamma_{\mathrm{w}}=\min\left(\gamma^{F},\gamma^{P}\right)>0$, line (C4r) of Section~I.8.2 and of Proposition~14.2\,(a).
\item $\kappa_{\mathrm{w}}-1\ge\kappa+1$.
\emph{Used in:} (F3.1), (F3.2).
\emph{Place:} $\kappa_{\mathrm{w}}:=\kappa+\kappa_0+3$, line (C4); Proposition~14.2\,(f).
\item (Z2.s4), all three inequalities.
\emph{Used in:} the property (G3) of \cite[Section~16.2.7]{QEDI}, that is (Do3): Sections~7, 8 and~9.
\emph{Place:} Proposition~14.2\,(b), line (C8).
\item (Z2.s4), first inequality: $\exp(-\varepsilon v)\le\frac{1}{4}$.
\emph{Used in:} Section~7, Step~4.
\emph{Place:} Proposition~14.2\,(b), line (C8).
\item (A10), first inequality: $v\ge9$.
\emph{Used in:} (F7.1); Section~7, Steps~5 and~6 ($\frac{1}{3}\exp\left(\frac{9}{4}\right)>1$; $\mathrm{rad}>2$).
\emph{Place:} Proposition~14.2\,(b), line (C8).
\item (A10), second inequality: $8M\cdot\left(\frac{v}{4}+4\right)\cdot\exp(-\varepsilon v)\le1$.
\emph{Used in:} (F7.3).
\emph{Place:} Proposition~14.2\,(b), line (C8).
\item $r\ge2$.
\emph{Used in:} $r_k\ge v_k^2$ in Section~7, Step~1.
\emph{Place:} line (C7).
\item $\kappa\ge\kappa_0$.
\emph{Used in:} the assertion (t8) in (F9.1).
\emph{Place:} (C2); Proposition~14.2\,(f).
\item $L\cdot(\kappa-\kappa_0-1)\ge\kappa$.
\emph{Used in:} (F9.1).
\emph{Place:} a consequence of (L3) and (C2): Proposition~14.2\,(b), row (L3), and (f).
\item (Z2.s5): $8\,c_{\mathrm{rb}}\cdot L^3\cdot C_d\cdot M^3\cdot\exp(-2\varepsilon v)\le1$, with the constant $C_d$ of Proposition~13.
\emph{Used in:} Section~9.
\emph{Place:} Proposition~14.2\,(b), line (C8).
\item $\frac{1}{4}-2\varepsilon\ge0$.
\emph{Used in:} Section~9, last step.
\emph{Place:} (S1) with $\bar{\varepsilon}=\frac{1}{200}$; lines (C0), (C7).
\item The conditions under which the statements applied in the proofs hold.
The statements applied are Propositions~2.2\,(f0), (a); 2.8\,(a), (b); 5\,(a); 10\,(b), (e), (g); 12; and 16.1\,(a1), (d1)--(d4),
besides the property (G3) of \cite[Section~16.2.7]{QEDI} and \cite[Lemma~17.5(a), (c)]{QEDI}.
The conditions of the Propositions just named are the following:
the six rows (L5), (M1), (M2), the part $M\ge108\pi\cdot C'_{V}\cdot a'$ of (M5), (Z2.s1) and (Z2.s2) of Proposition~14.2\,(b);
$\alpha<1$, the range of the line (C1), which enters Proposition~2.8 through Proposition~2.7\,(d), (e);
$\rho<1$ (line (C0)), which enters Proposition~12 through (I.L0.12);
$b>0$, that is $b:=1$ (line (C0)), a condition of Proposition~2.2 as a whole and hence of Propositions~10 and~12;
and the closure relation $\sum_{i}w_i\,\hat n_i=0$ (I.1.1) of the standard weights (line (C0)).
The closure relation enters Proposition~10\,(e), (g-e), which are applied in subsection~6.4, through Proposition~2.2\,(c), (d);
it enters Proposition~12, which is applied in subsection~6.1, through Proposition~2.2\,(e) and through (I.m1).
Of these conditions, Proposition~16.1 requires only the part ``$L$ odd'' of (L5) and the row (M1),
and Proposition~2.2\,(f0), (a) requires only ``$L$ odd'' and $b>0$.
The condition of the property (G3) of \cite[Section~16.2.7]{QEDI} is $L\ge5$, so that $\square^{\sim 5}$ lies in $\square^{\dagger}$;
those of \cite[Lemma~17.5(a), (c)]{QEDI} are $\varepsilon\le\frac{1}{4}$, $\alpha\ge0$ and (g0).
\emph{Used in:} the places where these statements are applied.
\emph{Place:} Proposition~14.2\,(a), (b).
\end{itemize}
The data $\ell\in\mathrm{Cl}$ (of which the proofs use that $\ell$ is a multiple of $L\cdot M$ and that $\ell\ge13L\cdot M$), $\ell\ge M\cdot r_k$ (in (d4)),
$\ell/L\in\mathrm{Cl}$ and $e_{k+1}=L^{\frac{1}{2}}\cdot e_k\le e_0$ (in (d5) and (d6)), and $\mu\le1$ are data of (SF) and (SF+);
they are not conditions on the constants.

\section*{Chapter 19. The block covariance at real fields: proofs of Proposition~19}

The statement proved in this chapter is Proposition~19.

\subsection*{Proof of Proposition~19}

Throughout the chapter the letters (a)--(f) denote the clauses of Proposition~19.
A number such as 1.3 or (19.2) without a prefix refers to a paragraph or to a display of this chapter.
The labels (w1) and (w2) denote the two facts on linear maps with which the proof of \cite[Theorem~13.18]{QEDI} begins:
(w1) is the bound for an operator whose real part is bounded below by a positive number,
and (w2) is the formula for the inverse of an invertible operator perturbed by $b\cdot Q'Q$, where $b>0$ and $Q$, $Q'$ are linear maps in opposite directions.

\subsubsection*{0. Conventions of this chapter}

\paragraph{0.1. Points, periods.}
All functions below are defined on the covering lattices $\mathbb{L}_k=\eta_k\mathbb{Z}^3$ and $\mathbb{L}^{\mathrm{f}}_k=\mathbb{L}_k\times O_{\mathrm{cub}}$, or on $\mathbb{Z}^3$ (Sections~I.1.1 and I.1.3), where $\eta_k=L^{-k}$.
Since $L$ is an integer (Section~I.0.5) and $k\ge0$, the number $L^k$ is an integer, so that $\mathbb{Z}^3=\eta_k(L^k\mathbb{Z}^3)$ is contained in $\mathbb{L}_k$.
The side $\ell$ is a positive integer: the condition $\ell\in\mathrm{Cl}$ means that $\ell/(LM)$ is an odd integer $\ge13$ (Section~I.1.1), where $L$ and $M$, a power of $L$, are positive integers (Section~I.0.5).
Hence $\ell\mathbb{Z}^3$ is a subgroup of $\mathbb{Z}^3$ and of $\mathbb{L}_k$.
For $z\in\mathbb{Z}^3$ let $\boldsymbol{\chi}(z)=(-1)^{z_1+z_2+z_3}$ (Section~I.6.3);
then $\boldsymbol{\chi}(z+z')=\boldsymbol{\chi}(z)\cdot\boldsymbol{\chi}(z')$ and $\boldsymbol{\chi}(z)^2=1$.
Translations act on the first coordinate of a fermion site only: $(c,g)+\ell z:=(c+\ell z,g)$.
A function is called periodic if it is invariant under the translations by the vectors $\ell z$, $z\in\mathbb{Z}^3$.
The sup distance on the torus is
\[
d(y,y')=\min\left\{|y-y'+\ell z|:z\in\mathbb{Z}^3\right\}
\]
(Section~I.1.1).

An element of $\mathfrak{S}^{\mathrm{ch}}(\eta_k,\ell)$ is a function $f\colon\mathbb{L}^{\mathrm{f}}_k\to\mathbb{C}^4$ such that $f(x+\ell z)=\boldsymbol{\chi}(z)\cdot f(x)$ for all $z\in\mathbb{Z}^3$,
and an element of $\mathfrak{S}(1,\ell)$ is a function $v\colon\mathbb{Z}^3\to\mathbb{C}^4$ such that $v(y+\ell z)=\boldsymbol{\chi}(z)\cdot v(y)$ for all $z\in\mathbb{Z}^3$.
Indeed, Section~I.1.5 requires $f(x+\ell e_\mu)=-f(x)$ and $v(y+\ell e_\mu)=-v(y)$ for the three generators $e_\mu$;
with $x-\ell e_\mu$ and $y-\ell e_\mu$ in place of $x$ and $y$ these relations become $f(x-\ell e_\mu)=-f(x)$ and $v(y-\ell e_\mu)=-v(y)$,
and the four relations, applied $|z_1|+|z_2|+|z_3|$ times, give the two conditions above.

The bond function $A$ is a real function on the bonds $(c,c')$ of $\mathbb{L}_k$ with $A(c',c)=-A(c,c')$ and $A(c+\ell z,c'+\ell z)=A(c,c')$ (Section~I.1.1).

\paragraph{0.2. The integral.}
For a periodic complex function $\mathrm{pf}$ on $\mathbb{L}^{\mathrm{f}}_k$ one has, by (I.1.3),
\[
\int\mathrm{d}x\,\mathrm{pf}(x)=\frac{\eta_k^3}{48}\sum_c\sum_{g\in O_{\mathrm{cub}}}\mathrm{pf}(c,g),
\]
where $c$ ranges over a set of representatives of $\mathbb{L}_k/\ell\mathbb{Z}^3$.
This is the sum of $\mathrm{pf}$ over the finite set $T^{\mathrm{f}}_k=\mathbb{L}^{\mathrm{f}}_k/\ell\mathbb{Z}^3$, each point having the weight $\eta_k^3/48$;
in particular it does not depend on the choice of the representatives.

\paragraph{0.3. Scalar products, norms, real part.}
The scalar products of Section~I.1.6 and of the data of Proposition~19 are
\begin{align*}
(f',f)^{\mathrm{ch}}&=\int\mathrm{d}x\,\overline{f'(x)}\cdot f(x)=\frac{\eta_k^3}{48}\sum_c\sum_{g\in O_{\mathrm{cub}}}\sum_{\beta=1}^{4}\overline{f'_\beta(c,g)}\cdot f_\beta(c,g)&&\text{on }\mathfrak{S}^{\mathrm{ch}}(\eta_k,\ell),\\
(v',v)&:=\sum_y\overline{v'(y)}\cdot v(y)=\sum_y\sum_{\beta=1}^{4}\overline{v'_\beta(y)}\cdot v_\beta(y)&&\text{on }\mathfrak{S}(1,\ell),
\end{align*}
where $y$ ranges over a set of representatives of $\mathbb{Z}^3/\ell\mathbb{Z}^3$.
The summands are periodic, since $\boldsymbol{\chi}(z)^2=1$.
Let $\|f\|:=\left((f,f)^{\mathrm{ch}}\right)^{\frac{1}{2}}$ and $\|v\|:=(v,v)^{\frac{1}{2}}$; thus $\|f\|^2=\int\mathrm{d}x\,|f(x)|_2^2$.
The two spaces have the dimensions $192(\ell/\eta_k)^3$ and $4\ell^3$; in particular they are finite dimensional and different from $\{0\}$.
The adjoint with respect to these scalar products is denoted by ${}^*$, and the operator norm with respect to these norms by $|\cdot|_{\mathrm{op}}$.
For a linear operator $O$ of one of the two spaces into itself and a real number $\mathrm{cf}$, the inequality $\operatorname{Re}O\ge\mathrm{cf}$ means that $\operatorname{Re}(f,Of)\ge\mathrm{cf}\cdot(f,f)$ for every $f$ in the space.

A periodic complex function $\psi$ on $\mathbb{L}_k$ acts on spinors by $(\psi f)(x):=\psi(c)\cdot f(x)$ for $x=(c,g)$ ((I.POS)),
and a function of the label $g$ with values in the $4\times4$ matrices acts by multiplication at each site.
Both kinds of operators map $\mathfrak{S}^{\mathrm{ch}}(\eta_k,\ell)$ into itself, because the multipliers do not change under $x\mapsto x+\ell z$.

\paragraph{0.4. Notation.}
Let $\mathrm{Rop}:=D_A+\mu$ for every $k\ge0$ and, for $k\ge1$, $\mathrm{Qa}:=Q_k(A)\cdot J^*$.
The other auxiliary letters of the chapter are introduced where they first occur.
The twist is denoted by $\theta$ (Section~9); the number introduced in 7.2 is a different object.
As in Part~I, $\boldsymbol{\chi}(\cdot)$ is the sign of Section~I.6.3, here the sign of the translation by $\ell z$;
$\mathrm{cf}$ denotes a real constant, $O$ an operator, and $c(\gamma)$ the lattice sum (G2.r6) of Section~G2.0.5, $\gamma$ being a dummy variable;
$m$ denotes a $4\times4$ block (Section~I.1.5).

\subsubsection*{1. Facts of Part I, derived from the definitions}

In 1.1--1.6 the level $k\ge0$ is arbitrary; in 1.7 and 1.8 it satisfies $k\ge1$.

\paragraph{1.1. Neighbours: (I.s1), (I.s2), (I.s3).}
Let $x=(c,g)$.
By (I.1.2),
\begin{gather*}
x^{(i)}=(c,gs_i)\quad(i=0,1,2),\qquad x^{(3)}=(c+\eta_k\cdot ge_1,\,gs_3),\\
c_3=c+\eta_k\cdot ge_1,\qquad\hat d_i(x)=-g\hat n_i.
\end{gather*}
According to Section~I.1.2, $s_i=1-2\hat n_i\hat n_i^{T}$ with $|\hat n_i|_2=1$.
Hence $s_i\hat n_i=\hat n_i-2\hat n_i=-\hat n_i$ and
\[
s_i^2=1-4\hat n_i\hat n_i^{T}+4\hat n_i\left(\hat n_i^{T}\hat n_i\right)\hat n_i^{T}=1;
\]
since $\hat n_3=-e_1$, in particular $s_3e_1=-e_1$.
Explicitly,
\[
s_0=1-2e_3e_3^{T}=\operatorname{diag}(1,1,-1),\qquad s_3=1-2e_1e_1^{T}=\operatorname{diag}(-1,1,1),
\]
while $s_1=1-(e_2-e_3)(e_2-e_3)^{T}$ fixes $e_1$ and interchanges $e_2$ and $e_3$,
and $s_2=1-(e_1-e_2)(e_1-e_2)^{T}$ fixes $e_3$ and interchanges $e_1$ and $e_2$.
Thus every $s_i$ is a signed permutation matrix.
The signed permutation matrices form the group $O_{\mathrm{cub}}$, and $\#|O_{\mathrm{cub}}|=3!\cdot2^3=48$;
hence $gs_i\in O_{\mathrm{cub}}$ for every $g\in O_{\mathrm{cub}}$.
\begin{itemize}
\item[(I.s1)] For $i\le2$ one has $\left(x^{(i)}\right)^{(i)}=(c,gs_is_i)=x$, and
\[
\left(x^{(3)}\right)^{(3)}=\left(c+\eta_k\cdot ge_1+\eta_k\cdot(gs_3)e_1,\,gs_3s_3\right)=\left(c+\eta_k\cdot ge_1-\eta_k\cdot ge_1,\,g\right)=x.
\]
So $x\mapsto x^{(i)}$ is an involution of $\mathbb{L}^{\mathrm{f}}_k$, and $(x+\ell z)^{(i)}=x^{(i)}+\ell z$.
In particular $\pi\left(\left(x^{(3)}\right)^{(3)}\right)=c$, that is, the photon site $c_3(x^{(3)})$ is $c$.
\item[(I.s2)] $\hat d_i(x^{(i)})=-(gs_i)\hat n_i=g\hat n_i=-\hat d_i(x)$.
\item[(I.s3)] With the standard weights (see (I.1.1)),
\begin{align*}
\sum_{i=0}^{3}w_i\hat d_i(x)&=-g\left[6e_3+6\sqrt2\cdot\frac{e_2-e_3}{\sqrt2}+6\sqrt2\cdot\frac{e_1-e_2}{\sqrt2}-6e_1\right]\\
&=-g\left[6e_3+6e_2-6e_3+6e_1-6e_2-6e_1\right]=0.
\end{align*}
\end{itemize}
\emph{Positions.}
For $i=0,1,2$ the site $x^{(i)}$ lies over $c$, and $g\mapsto gs_i$ permutes the $48$ sites of the cluster $\{c\}\times O_{\mathrm{cub}}$.
The site $x^{(3)}$ lies over $c_3$, and $|c_3-c|=\eta_k\cdot|ge_1|=\eta_k$, because $ge_1$ is one of the six vectors $\pm e_\mu$.

By (I.s1) the map $x\mapsto x^{(i)}$ induces a bijection of the finite set $T^{\mathrm{f}}_k$, all of whose points have the same weight (0.2).
Consequently, for $i=0,1,2,3$ and every periodic function $\mathrm{pf}$,
\begin{equation}
\int\mathrm{d}x\,\mathrm{pf}(x^{(i)})=\int\mathrm{d}x\,\mathrm{pf}(x).
\tag{19.1}
\end{equation}

\paragraph{1.2. Spin matrices: a part of (I.c1); (I.m6).}
By Section~I.2.1, $\gamma_\mu=\sigma_3\otimes\sigma_\mu$.
The Pauli matrices are Hermitian, $\sigma_3^2=1_2$, and $\sigma_\mu\sigma_\nu+\sigma_\nu\sigma_\mu=2\delta_{\mu\nu}1_2$.
Therefore $\gamma_\mu^*=\gamma_\mu$ and
\[
\gamma_\mu\gamma_\nu+\gamma_\nu\gamma_\mu=\sigma_3^2\otimes\left(\sigma_\mu\sigma_\nu+\sigma_\nu\sigma_\mu\right)=2\delta_{\mu\nu}1_4.
\]
For $x=(c,g)$ the vector $\hat d_i(x)=-g\hat n_i$ is real, and $|\hat d_i(x)|_2=|\hat n_i|_2=1$, since $g$ is orthogonal.
Hence the matrix $\gamma\cdot\hat d_i(x)=\sum_\mu\hat d_i(x)_\mu\gamma_\mu$ is Hermitian, and
\[
\left(\gamma\cdot\hat d_i(x)\right)^2=\frac12\sum_{\mu,\nu}\hat d_i(x)_\mu\cdot\hat d_i(x)_\nu\cdot\left(\gamma_\mu\gamma_\nu+\gamma_\nu\gamma_\mu\right)=|\hat d_i(x)|_2^2\,1_4=1_4,
\]
so that its eigenvalues lie in $\{1,-1\}$.

By (I.2.1) the matrix $M_i(x)=\frac{w_i}{2}\left(\rho-\gamma\cdot\hat d_i(x)\right)$ depends only on the label $g$;
it is Hermitian, and its eigenvalues lie in $\left\{\frac{w_i}{2}(\rho-1),\frac{w_i}{2}(\rho+1)\right\}$.
The norm $|\cdot|_2$ of a Hermitian matrix is the largest modulus of its eigenvalues.
As $\rho\ge0$, one has $|\rho-1|\le1+\rho$, and therefore $|M_i(x)|_2\le\frac{w_i}{2}(1+\rho)$.
(Equality holds: $\operatorname{tr}\gamma_\mu=\operatorname{tr}\sigma_3\cdot\operatorname{tr}\sigma_\mu=0$, so that both eigenvalues of $\gamma\cdot\hat d_i(x)$ occur.)
For $i=3$ one has $w_3=6$ and $\hat d_3(x)=ge_1$, and since $\rho\le1$,
\[
M_3(x)=3\left(\rho-\gamma\cdot(ge_1)\right),\qquad|M_3(x)|_2\le3(1+\rho)\le6.
\]

\paragraph{1.3. The hops: the first sentence of (I.O2).}
By (I.1.6) and (I.2.3), $U_A(x,x^{(i)})=1$ for $i=0,1,2$, and
\[
U_A(x,x^{(3)})=u_A(c,c_3)=\exp\left(i\,e_k\cdot\eta_k\cdot A(c,c_3)\right).
\]
Since $A$ is real, $|U_A(x,x^{(i)})|=1$.
By (I.s1) the pair $\left(x^{(i)},(x^{(i)})^{(i)}\right)$ is $(x^{(i)},x)$, and for $i=3$ one has $U_A(x^{(3)},x)=u_A(c_3,c)=\exp\left(-i\,e_k\cdot\eta_k\cdot A(c,c_3)\right)$.
Hence, for $i=0,\dots,3$,
\[
U_A(x^{(i)},x)=U_A(x,x^{(i)})^{-1}=\overline{U_A(x,x^{(i)})}.
\]
As $A$ is periodic and $(x+\ell z)^{(i)}=x^{(i)}+\ell z$, also $U_A(x+\ell z,x^{(i)}+\ell z)=U_A(x,x^{(i)})$.

Let $(R_i(A)f)(x):=U_A(x,x^{(i)})\cdot f(x^{(i)})$, as in (I.2.2).
The following three properties hold.
\begin{itemize}
\item[(i)] $(R_i(A)f)(x+\ell z)=U_A(x,x^{(i)})\cdot f(x^{(i)}+\ell z)=\boldsymbol{\chi}(z)\cdot(R_i(A)f)(x)$;
thus $R_i(A)$ maps $\mathfrak{S}^{\mathrm{ch}}(\eta_k,\ell)$ into itself.
\item[(ii)] $(R_i(A)^2f)(x)=U_A(x,x^{(i)})\cdot U_A(x^{(i)},x)\cdot f(x)=f(x)$; thus $R_i(A)^2=1$.
\item[(iii)] $R_i(A)^*=R_i(A)$.
To see this, let $\mathrm{pf}(x):=\overline{f'(x^{(i)})}\cdot U_A(x^{(i)},x)\,f(x)$.
This function is periodic, $\mathrm{pf}(x^{(i)})=\overline{f'(x)}\cdot U_A(x,x^{(i)})\,f(x^{(i)})$, and by (19.1)
\begin{align*}
(f',R_i(A)f)^{\mathrm{ch}}&=\int\mathrm{d}x\,\mathrm{pf}(x^{(i)})=\int\mathrm{d}x\,\mathrm{pf}(x)\\
&=\int\mathrm{d}x\,\overline{U_A(x,x^{(i)})\,f'(x^{(i)})}\cdot f(x)=(R_i(A)f',f)^{\mathrm{ch}}.
\end{align*}
\end{itemize}
So $R_i(A)$ is a Hermitian involution, and therefore unitary.
Since $(1-R_i(A))^*\cdot(1-R_i(A))=(1-R_i(A))^2=2(1-R_i(A))$, it follows that
\begin{equation}
\begin{split}
(f,(1-R_i(A))f)^{\mathrm{ch}}&=\frac12\left\|(1-R_i(A))f\right\|^2\\
&=\frac12\int\mathrm{d}x\,\left|U_A(x,x^{(i)})\cdot f(x^{(i)})-f(x)\right|_2^2\ge0.
\end{split}
\tag{19.2}
\end{equation}

\paragraph{1.4. The decomposition (I.2.4).}
By (I.2.3) and (I.2.1), $D_Af=\eta_k^{-1}\sum_iM_i\cdot\left(f-R_i(A)f\right)$;
by 0.3 and 1.3\,(i) the operator $D_A$ maps $\mathfrak{S}^{\mathrm{ch}}(\eta_k,\ell)$ into itself.
At a site $x$,
\begin{align*}
(D_Af)(x)&=\eta_k^{-1}\sum_{i=0}^{3}\frac{w_i}{2}\left(\rho-\gamma\cdot\hat d_i(x)\right)\left[f(x)-(R_i(A)f)(x)\right]\\
&=\eta_k^{-1}\Biggl\{\frac{\rho}{2}\sum_iw_i\left[f(x)-(R_i(A)f)(x)\right]-\frac12\,\gamma\cdot\left(\sum_iw_i\hat d_i(x)\right)f(x)\\
&\qquad\qquad+\frac12\sum_iw_i\left(\gamma\cdot\hat d_i(x)\right)\cdot(R_i(A)f)(x)\Biggr\}.
\end{align*}
The middle term vanishes by (I.s3).
Let
\[
\nabla_A^*\nabla_A:=\sum_iw_i\left(1-R_i(A)\right),\qquad(J_Af)(x):=\sum_iw_i\left(\gamma\cdot\hat d_i(x)\right)\cdot U_A(x,x^{(i)})\cdot f(x^{(i)}).
\]
Then
\begin{equation}
D_A=\eta_k^{-1}\left[\frac{\rho}{2}\,\nabla_A^*\nabla_A+\frac12\,J_A\right].
\tag{19.3}
\end{equation}
The assertions 1.3\,(i) and (ii), the computation of this paragraph and 1.6 do not use that $A$ is real:
they hold for every complex periodic bond function, every level $k\ge0$, every positive integer $\ell$ and every $e_k>0$.

\paragraph{1.5. (I.O2): $\nabla_A^*\nabla_A\ge0$, and $J_A$ is anti-Hermitian.}
The operator $\nabla_A^*\nabla_A$ is Hermitian by 1.3\,(iii), the $w_i$ being real, and by (19.2), since $w_i>0$,
\[
(f,\nabla_A^*\nabla_Af)^{\mathrm{ch}}=\sum_{i=0}^{3}\frac{w_i}{2}\int\mathrm{d}x\,\left|U_A(x,x^{(i)})\cdot f(x^{(i)})-f(x)\right|_2^2\ge0.
\]
Let $\mathrm{Gd}_i$ be the multiplication by the matrix $\gamma\cdot\hat d_i(x)$, which depends only on the label.
By 1.2, $\mathrm{Gd}_i^*=\mathrm{Gd}_i$, and $J_A=\sum_iw_i\,\mathrm{Gd}_i\cdot R_i(A)$.
From (I.s2) one obtains
\begin{align*}
(R_i(A)\cdot\mathrm{Gd}_if)(x)&=U_A(x,x^{(i)})\cdot\left(\gamma\cdot\hat d_i(x^{(i)})\right)\cdot f(x^{(i)})\\
&=-\left(\gamma\cdot\hat d_i(x)\right)\cdot U_A(x,x^{(i)})\cdot f(x^{(i)})=-(\mathrm{Gd}_i\cdot R_i(A)f)(x).
\end{align*}
Hence $(\mathrm{Gd}_i\cdot R_i(A))^*=R_i(A)\cdot\mathrm{Gd}_i=-\mathrm{Gd}_i\cdot R_i(A)$, that is, $J_A^*=-J_A$.
Consequently
\[
(f,J_Af)^{\mathrm{ch}}=(J_A^*f,f)^{\mathrm{ch}}=-(J_Af,f)^{\mathrm{ch}}=-\overline{(f,J_Af)^{\mathrm{ch}}},
\]
that is, $\operatorname{Re}(f,J_Af)^{\mathrm{ch}}=0$.

\paragraph{1.6. Conjugation by a weight: (I.O6).}
Let $\psi$ be a real periodic function on $\mathbb{L}_k$.
By 0.3, $\exp(\psi)$ and $\exp(-\psi)$ are multiplication operators on $\mathfrak{S}^{\mathrm{ch}}(\eta_k,\ell)$;
they are Hermitian, the functions being real, they are inverse to each other, and they commute with the matrices $M_i(x)$.
By (I.2.3),
\begin{align*}
&\left(\exp(\psi)\cdot D_A\cdot\exp(-\psi)f\right)(x)\\
&\qquad=-\eta_k^{-1}\sum_{i=0}^{3}M_i(x)\cdot\left[U_A(x,x^{(i)})\cdot\exp\left(\psi(c)-\psi(\pi(x^{(i)}))\right)\cdot f(x^{(i)})-f(x)\right].
\end{align*}
For $i=0,1,2$ one has $\pi(x^{(i)})=c$ (1.1), so that the three hops inside the cluster are unchanged; for $i=3$, $\pi(x^{(3)})=c_3$.
Therefore
\begin{equation}
\begin{split}
&\left(\left(\exp(\psi)\cdot D_A\cdot\exp(-\psi)-D_A\right)f\right)(x)\\
&\qquad=-\eta_k^{-1}\cdot M_3(x)\cdot u_A(c,c_3)\cdot\left(\exp\left(\psi(c)-\psi(c_3)\right)-1\right)\cdot f(x^{(3)}).
\end{split}
\tag{19.4}
\end{equation}

\paragraph{1.7. Unit blocks ($k\ge1$): Section~I.1.1 (``Blocks'', ``Faces''), Section~I.1.4.}
For $y\in\mathbb{Z}^3$ let $\Delta_y=\{c\in\mathbb{L}_k:|c-y|<1/2\}$.
A coordinate of a point of $\mathbb{L}_k$ has the form $j/L^k$ with an integer $j$, and its distance from an integer $y_\mu$ is never $1/2$:
the equality $|j/L^k-y_\mu|=1/2$ would give to the even number $2j$ the value $L^k(2y_\mu\pm1)$, which is odd because $L$ is odd.
A real number whose distance from $\mathbb{Z}$ differs from $1/2$ has exactly one integer at distance $<1/2$.
Hence every $c\in\mathbb{L}_k$ lies in exactly one block; its centre is denoted by $[c]$.
By 0.1, $y$ belongs to $\mathbb{L}_k$, so that the points of $\Delta_y$ are the points $y+\eta_k\cdot j$ with $j\in\mathbb{Z}^3$ and $|j_\mu|<L^k/2$, that is, $|j_\mu|\le(L^k-1)/2$, since $L^k$ is odd.
Thus $\#|\Delta_y|=(L^k)^3=\eta_k^{-3}$, and $\#|\Delta_y\times O_{\mathrm{cub}}|=48\eta_k^{-3}$.
For $z\in\mathbb{Z}^3$ one has $\Delta_{y+z}=\Delta_y+z$, and hence $[c+\ell z]=[c]+\ell z$.
If $y$ ranges over a set of representatives of $\mathbb{Z}^3/\ell\mathbb{Z}^3$, the points of the blocks $\Delta_y$ form a set of representatives of $\mathbb{L}_k/\ell\mathbb{Z}^3$:
for $c\in\mathbb{L}_k$ there is exactly one $z\in\mathbb{Z}^3$ such that $[c]-\ell z=[c-\ell z]$ is one of the chosen points $y$.

\paragraph{1.8. Phases ($k\ge1$): Section~I.3.2.}
\emph{Blocks of the sides $L^j\eta_k$.}
Let $0\le j\le k$.
Then $L^j\eta_k=L^{j-k}$, and $L^{j-k}\mathbb{Z}^3=\eta_k(L^j\mathbb{Z}^3)$ is contained in $\mathbb{L}_k$, since $L^j$ is an integer.
For $z'\in L^{j-k}\mathbb{Z}^3$ let, as in Section~I.1.1 (``Blocks''),
\[
\mathrm{Bl}_j(z')=\left\{c\in\mathbb{L}_k:|c-z'|<L^{j-k}/2\right\};
\]
thus $\mathrm{Bl}_k(y)=\Delta_y$ and $\mathrm{Bl}_0(z')=\{z'\}$.
Every $c\in\mathbb{L}_k$ lies in exactly one block $\mathrm{Bl}_j(z')$ with $z'\in L^{j-k}\mathbb{Z}^3$ (Section~I.1.1, ``Faces'', and Section~I.1.4).
Indeed, a coordinate of $c$ is $\mathrm{ja}/L^k$ and the same coordinate of $z'$ is $\mathrm{jb}\cdot L^{j-k}$, with integers $\mathrm{ja}$ and $\mathrm{jb}$;
the equality $|\mathrm{ja}/L^k-\mathrm{jb}\cdot L^{j-k}|=L^{j-k}/2$ would give to the even number $2\,\mathrm{ja}$ the value $L^j(2\,\mathrm{jb}\pm1)$, which is odd because $L$ is odd;
and a real number whose distance from $L^{j-k}\mathbb{Z}$ differs from $L^{j-k}/2$ has exactly one point of $L^{j-k}\mathbb{Z}$ at distance $<L^{j-k}/2$.
For $z\in\mathbb{Z}^3$ the vector $\ell z$ lies in $L^{j-k}\mathbb{Z}^3$, because $L^{k-j}\ell z$ is an integer vector ($L$ is an integer and $j\le k$);
therefore $\mathrm{Bl}_j(z')+\ell z=\mathrm{Bl}_j(z'+\ell z)$ is a block of the same side.

\emph{Centres.}
For $y\in\mathbb{Z}^3$, $c\in\Delta_y$ and $0\le j\le k$ let $y_j$ be the centre of the block $\mathrm{Bl}_j$ that contains $c$.
Thus $y_0=c$ and $y_k=y$, and the centres belonging to the point $c+\ell z$ of $\Delta_{y+\ell z}$ are the points $y_j+\ell z$.
For $j<k$ the vector $y_{j+1}-y_j$ lies in $L^{j-k}\mathbb{Z}^3$, which contains $L^{j+1-k}\mathbb{Z}^3$, and
\[
|y_{j+1}-y_j|\le|y_{j+1}-c|+|c-y_j|<\frac{L^{j+1-k}}{2}+\frac{L^{j-k}}{2}=L^{j-k}\cdot\frac{L+1}{2}.
\]
Since $(L+1)/2$ is an integer, it follows that
\[
|y_{j+1}-y_j|\le L^{j-k}\left(\frac{L+1}{2}-1\right)=L^{j-k}\cdot\frac{L-1}{2}<\frac{L^{j+1-k}}{2}\le\frac12.
\]
So $(y_{j+1},y_j)$ is the lift with $|y_{j+1}-y_j|<1/2$ on which Section~I.3.2 forms the paths.

\emph{Phases.}
By (I.3.3) and Section~I.1.1 (``Paths, axial functional''),
\begin{gather*}
\varphi_A(y,c)=\exp\left(i\,e_k\cdot\eta_k\cdot(\tau_kA)(y,c)\right),\qquad(\tau_kA)(y,c)=\sum_{j<k}(\tau A)(y_{j+1},y_j),\\
(\tau A)(c',c'')=\frac16\sum_\pi A\left(\boldsymbol{\Gamma}^{\pi}(c',c'')\right),\qquad(\tau A)(c',c')=0,
\end{gather*}
where $A(\boldsymbol{\Gamma})$ is the sum of $A(b)$ over the bonds $b$ of the path $\boldsymbol{\Gamma}$, oriented along the path.
So $(\tau_kA)(y,c)$ is a finite sum of terms $\frac16A(b)$ with bonds $b$ of $\mathbb{L}_k$.
Since $A$ is real, $(\tau_kA)(y,c)$ is real, and
\[
|\varphi_A(y,c)|=1,\qquad\varphi_A(y,c)^{-1}=\overline{\varphi_A(y,c)}.
\]
The rectilinear paths from $y_{j+1}+\ell z$ to $y_j+\ell z$ are the paths $\boldsymbol{\Gamma}^{\pi}(y_{j+1},y_j)+\ell z$, and $A$ is periodic; therefore
\[
\varphi_A(y+\ell z,c+\ell z)=\varphi_A(y,c)\qquad\text{for all }z\in\mathbb{Z}^3.
\]

\subsubsection*{2. Operators with positive real part}

\begin{named}{(w1)}
Let $O$ be a linear operator of a finite dimensional complex Hilbert space into itself,
and let $\mathrm{cf}>0$ be a number such that $\operatorname{Re}(f,Of)\ge\mathrm{cf}\cdot(f,f)$ for all $f$.
Then $O$ is invertible, $|O^{-1}|_{\mathrm{op}}\le1/\mathrm{cf}$, and $\operatorname{Re}O^{-1}\ge0$.
\end{named}

\begin{proof}
For every $f$,
\[
\|Of\|\cdot\|f\|\ge|(f,Of)|\ge\operatorname{Re}(f,Of)\ge\mathrm{cf}\cdot\|f\|^2,
\]
and hence $\|Of\|\ge\mathrm{cf}\cdot\|f\|$.
Therefore $O$ is injective and, the dimension being finite, bijective;
inserting $f=O^{-1}f'$ one obtains $\|O^{-1}f'\|\le\|f'\|/\mathrm{cf}$.
For $f'=Of$,
\[
\operatorname{Re}(f',O^{-1}f')=\operatorname{Re}(Of,f)=\operatorname{Re}\overline{(f,Of)}=\operatorname{Re}(f,Of)\ge0.
\]
\end{proof}

\subsubsection*{3. Proof of (a) ($k\ge0$)}

\begin{proof}[Proof of Proposition~19\,(a)]
Let $f\in\mathfrak{S}^{\mathrm{ch}}(\eta_k,\ell)$.
By (19.3), 1.5 and (19.2),
\begin{align*}
\operatorname{Re}(f,(D_A+\mu)f)^{\mathrm{ch}}&=\mu\cdot(f,f)^{\mathrm{ch}}+\frac{\rho}{2\eta_k}\cdot(f,\nabla_A^*\nabla_Af)^{\mathrm{ch}}+\frac{1}{2\eta_k}\cdot\operatorname{Re}(f,J_Af)^{\mathrm{ch}}\\
&=\mu\cdot(f,f)^{\mathrm{ch}}+\frac{\rho}{2\eta_k}\cdot(f,\nabla_A^*\nabla_Af)^{\mathrm{ch}},\\
\frac{\rho}{2\eta_k}\cdot(f,\nabla_A^*\nabla_Af)^{\mathrm{ch}}&=\frac{\rho}{4\eta_k}\sum_{i=0}^{3}w_i\cdot\frac{\eta_k^3}{48}\sum_c\sum_{g\in O_{\mathrm{cub}}}\left|U_A(x,x^{(i)})\cdot f(x^{(i)})-f(x)\right|_2^2\ge0,\tag{19.5}
\end{align*}
where $x=(c,g)$, because $\mu$ and $(f,\nabla_A^*\nabla_Af)^{\mathrm{ch}}$ are real, $\rho\ge0$, $\eta_k>0$ and $w_i>0$.
In (19.5) the summands with $i=0,1,2$ are $|f(c,gs_i)-f(c,g)|_2^2$; they belong to pairs inside the cluster over $c$ and contain no phase.
The summand with $i=3$ is $|u_A(c,c_3)\cdot f(c_3,gs_3)-f(c,g)|_2^2$; it joins the clusters over $c$ and over $c_3=c+\eta_k\cdot ge_1$.
This proves the identity and the inequality of (a).
Now (w1), applied to $O=\mathrm{Rop}=D_A+\mu$ on $\mathfrak{S}^{\mathrm{ch}}(\eta_k,\ell)$ with $\mathrm{cf}=\mu>0$, shows that $D_A+\mu$ is invertible,
that $|(D_A+\mu)^{-1}|_{\mathrm{op}}\le1/\mu$ and that $\operatorname{Re}(D_A+\mu)^{-1}\ge0$.
The determinant of an invertible linear map of a finite dimensional space into itself does not vanish; hence $\det(D_A+\mu)\ne0$.

\emph{Level $0$.}
At $k=0$ one has $\eta_0=1$, $\mathbb{L}_0=\mathbb{Z}^3$ and $\mathbb{L}^{\mathrm{f}}_0=\mathbb{Z}^3\times O_{\mathrm{cub}}$;
the number $e_k$ is the number $q$ (Section~I.2.5 and (I.L0.2));
the space is $\mathfrak{S}^{\mathrm{ch}}(1,\ell)$, of dimension $192\ell^3=192n_0$; and
\begin{gather*}
\int\mathrm{d}x\,\mathrm{pf}(x)=\frac{1}{48}\sum_c\sum_{g\in O_{\mathrm{cub}}}\mathrm{pf}(c,g),\qquad x^{(3)}=(c+ge_1,\,gs_3),\\
u_A(c,c_3)=\exp\left(i\,q\cdot A(c,c_3)\right),\qquad(D_Af)(x)=-\sum_{i=0}^{3}M_i(x)\cdot\left[U_A(x,x^{(i)})\cdot f(x^{(i)})-f(x)\right].
\end{gather*}
The paragraphs 0.1--0.3 and 1.1--1.5, Section~2 and the computation above contain the level only through the number $\eta_k>0$ and the lattice $\eta_k\mathbb{Z}^3$,
and of these they need only that $\ell\mathbb{Z}^3$ is contained in $\eta_k\mathbb{Z}^3$ and that $c+\eta_k\cdot ge_1$ lies in $\eta_k\mathbb{Z}^3$.
No block, no phase $\varphi_A$ and no number $b_k$ occurs in them.
They therefore hold with $\eta_0=1$ and $q$, and (19.5) reads
\begin{align*}
\operatorname{Re}(f,(D_A+\mu)f)^{\mathrm{ch}}&=\mu\cdot(f,f)^{\mathrm{ch}}+\frac{\rho}{2}\cdot(f,\nabla_A^*\nabla_Af)^{\mathrm{ch}}\\
&=\mu\cdot(f,f)^{\mathrm{ch}}+\frac{\rho}{4}\sum_{i=0}^{3}w_i\cdot\frac{1}{48}\sum_c\sum_{g\in O_{\mathrm{cub}}}\left|U_A(x,x^{(i)})\cdot f(x^{(i)})-f(x)\right|_2^2\\
&\ge\mu\cdot(f,f)^{\mathrm{ch}}.
\end{align*}
Hence $D_0(A)=D_A+\mu$ (Section~I.3.9) is invertible on $\mathfrak{S}^{\mathrm{ch}}(1,\ell)$, $|D_0(A)^{-1}|_{\mathrm{op}}\le1/\mu$, $\operatorname{Re}D_0(A)^{-1}\ge0$ and $\det D_0(A)\ne0$.
The clauses (b)--(f) have no meaning at $k=0$, for two reasons.
Each of the clauses (c)--(f) contains the number $b_k$, and (c) contains $S_k(A)$ as well;
neither $b_0$ nor $S_k(A)$ at $k=0$ is defined ((I.L0.9)).
The clause (b) concerns only the products $Q_k(A)\cdot J^*$, $J\cdot Q_k^{T}(-A)$ and $J\cdot P_k(A)\cdot J^*$, and these are formed for $k\ge1$ only,
while at $k=0$ one has $Q_0=P_0:=1$ on $\mathfrak{S}^{\mathrm{ch}}(1,\ell)$ ((I.3.4), Section~I.3.9 and (I.L0.7)).
\end{proof}

\subsubsection*{4. Proof of (b) ($k\ge1$)}

\begin{proof}[Proof of Proposition~19\,(b)]
By (I.3.1) and Section~I.3.3, for $f\in\mathfrak{S}^{\mathrm{ch}}(\eta_k,\ell)$, $v\in\mathfrak{S}(1,\ell)$, $y\in\mathbb{Z}^3$ and $x=(c,g)\in\mathbb{L}^{\mathrm{f}}_k$,
\begin{equation}
\begin{split}
(\mathrm{Qa}f)(y)&=(Q_k(A)\cdot J^*f)(y)=\eta_k^3\sum_{c\in\Delta_y}\varphi_A(y,c)\cdot\frac{1}{48}\sum_{g\in O_{\mathrm{cub}}}f(c,g),\\
(J\cdot Q_k^{T}(-A)v)(c,g)&=\varphi_A([c],c)^{-1}\cdot v([c]).
\end{split}
\tag{19.6}
\end{equation}

\emph{The two maps act between the two spaces.}
By 1.7 ($\Delta_{y+\ell z}=\Delta_y+\ell z$ and $[c+\ell z]=[c]+\ell z$) and by 1.8,
\begin{align*}
(\mathrm{Qa}f)(y+\ell z)&=\eta_k^3\sum_{c\in\Delta_y}\varphi_A(y+\ell z,c+\ell z)\cdot\frac{1}{48}\sum_{g\in O_{\mathrm{cub}}}f(c+\ell z,g)\\
&=\boldsymbol{\chi}(z)\cdot(\mathrm{Qa}f)(y),\\
(J\cdot Q_k^{T}(-A)v)(c+\ell z,g)&=\varphi_A([c]+\ell z,c+\ell z)^{-1}\cdot v([c]+\ell z)\\
&=\boldsymbol{\chi}(z)\cdot(J\cdot Q_k^{T}(-A)v)(c,g).
\end{align*}

\emph{The adjoint.}
Let $y$ range over a set of representatives of $\mathbb{Z}^3/\ell\mathbb{Z}^3$.
By (19.6), and since $\overline{\varphi_A(y,c)}=\varphi_A(y,c)^{-1}$ (1.8),
\begin{align*}
(v,\mathrm{Qa}f)&=\sum_y\overline{v(y)}\cdot\left[\eta_k^3\sum_{c\in\Delta_y}\varphi_A(y,c)\cdot\frac{1}{48}\sum_{g\in O_{\mathrm{cub}}}f(c,g)\right]\\
&=\frac{\eta_k^3}{48}\sum_y\sum_{c\in\Delta_y}\sum_{g\in O_{\mathrm{cub}}}\overline{\varphi_A([c],c)^{-1}\cdot v([c])}\cdot f(c,g).
\end{align*}
By 1.7 the points $c$ of the blocks $\Delta_y$ form a set of representatives of $\mathbb{L}_k/\ell\mathbb{Z}^3$.
Hence the right-hand side equals $\int\mathrm{d}x\,\overline{(J\cdot Q_k^{T}(-A)v)(x)}\cdot f(x)=(J\cdot Q_k^{T}(-A)v,f)^{\mathrm{ch}}$.
As this holds for all $v$ and $f$,
\[
\left(Q_k(A)\cdot J^*\right)^*=\mathrm{Qa}^*=J\cdot Q_k^{T}(-A).
\]

\emph{The product.}
By (19.6), since $[c]=y$ for $c\in\Delta_y$, and by 1.7,
\begin{align*}
(\mathrm{Qa}\cdot J\cdot Q_k^{T}(-A)v)(y)&=\eta_k^3\sum_{c\in\Delta_y}\varphi_A(y,c)\cdot\frac{1}{48}\sum_{g\in O_{\mathrm{cub}}}\varphi_A(y,c)^{-1}\cdot v(y)\\
&=\eta_k^3\cdot\eta_k^{-3}\cdot\frac{1}{48}\cdot48\cdot v(y)=v(y),
\end{align*}
since $\#|\Delta_y\times O_{\mathrm{cub}}|=48\eta_k^{-3}$ and every site has the weight $\eta_k^3/48$.
Thus
\[
\left(Q_k(A)\cdot J^*\right)\cdot\left(J\cdot Q_k^{T}(-A)\right)=\mathrm{Qa}\cdot\mathrm{Qa}^*=1\qquad\text{on }\mathfrak{S}(1,\ell).
\]

\emph{Norms.}
For every $v$ one has $\|\mathrm{Qa}^*v\|^2=(v,\mathrm{Qa}\cdot\mathrm{Qa}^*v)=\|v\|^2$, and $\mathfrak{S}(1,\ell)\ne\{0\}$; hence $|\mathrm{Qa}^*|_{\mathrm{op}}=1$.
With suprema over $v\ne0$ and $f\ne0$,
\[
|\mathrm{Qa}|_{\mathrm{op}}=\sup\frac{|(v,\mathrm{Qa}f)|}{\|v\|\cdot\|f\|}=\sup\frac{|(\mathrm{Qa}^*v,f)^{\mathrm{ch}}|}{\|v\|\cdot\|f\|}=|\mathrm{Qa}^*|_{\mathrm{op}}=1.
\]

\emph{The projection.}
By Section~I.3.3, $P_k(A)=Q_k^{T}(-A)\cdot Q_k(A)$, and therefore
\[
J\cdot P_k(A)\cdot J^*=\left(J\cdot Q_k^{T}(-A)\right)\cdot\left(Q_k(A)\cdot J^*\right)=\mathrm{Qa}^*\cdot\mathrm{Qa}.
\]
This operator is Hermitian, and $(\mathrm{Qa}^*\cdot\mathrm{Qa})^2=\mathrm{Qa}^*\cdot(\mathrm{Qa}\cdot\mathrm{Qa}^*)\cdot\mathrm{Qa}=\mathrm{Qa}^*\cdot\mathrm{Qa}$.
A Hermitian idempotent is an orthogonal projection.
\end{proof}

\subsubsection*{5. Proof of (c) ($k\ge1$)}

\begin{proof}[Proof of Proposition~19\,(c)]
The number $b_k$ of Section~I.3.4,
\[
b_k=\frac{b\cdot\left(1-L^{-1}\right)}{1-L^{-k}},
\]
is positive, since $b=1>0$, $L>1$ and $k\ge1$.
Let $f\in\mathfrak{S}^{\mathrm{ch}}(\eta_k,\ell)$.
Section~4 gives $J\cdot P_k(A)\cdot J^{*}=\mathrm{Qa}^{*}\cdot\mathrm{Qa}$, and together with (a) this yields
\begin{align*}
\operatorname{Re}\left(f,\left(D_A+\mu+b_k\cdot J\cdot P_k(A)\cdot J^{*}\right)f\right)^{\mathrm{ch}}
&=\operatorname{Re}\left(f,\left(D_A+\mu\right)f\right)^{\mathrm{ch}}+b_k\cdot\|\mathrm{Qa}\,f\|^2\\
&\ge\mu\cdot(f,f)^{\mathrm{ch}}.
\end{align*}
Applying (w1) with $\mathrm{cf}=\mu$, one finds that $D_A+\mu+b_k\cdot J\cdot P_k(A)\cdot J^{*}$ is invertible.
Consequently the two operators of Section~I.3.5,
\begin{gather*}
S_k(A)=\left(D_A+\mu+b_k\cdot J\cdot P_k(A)\cdot J^{*}\right)^{-1},\\
D_k(A)=b_k-b_k^2\cdot Q_k(A)\cdot J^{*}\cdot S_k(A)\cdot J\cdot Q_k^{T}(-A),
\end{gather*}
are defined.
\end{proof}

\subsubsection*{6. Proof of (d) ($k\ge1$)}

\begin{proof}[Proof of Proposition~19\,(d)]
The argument is (w2), the inversion formula from the proof of \cite[Theorem~13.18]{QEDI}, carried out for the operator $\mathrm{Rop}$, with $\mathrm{Qa}$ and $\mathrm{Qa}^{*}$ in the roles of $Q$ and $Q'$ and with the number $b_k$ in the role of $b$.
By (a), $\mathrm{Rop}=D_A+\mu$ is invertible.
On $\mathfrak{S}(1,\ell)$ define
\[
\mathrm{Cq}:=\mathrm{Qa}\cdot\mathrm{Rop}^{-1}\cdot\mathrm{Qa}^{*}=Q_k(A)\cdot J^{*}\cdot\left(D_A+\mu\right)^{-1}\cdot J\cdot Q_k^{T}(-A),\qquad
\mathrm{Nq}:=b_k^{-1}+\mathrm{Cq}.
\]

\emph{Real part and norm of $\mathrm{Cq}$.}
For $v\in\mathfrak{S}(1,\ell)$, clauses (b) and (a) give
\begin{gather*}
\operatorname{Re}\left(v,\mathrm{Cq}\,v\right)=\operatorname{Re}\left(\mathrm{Qa}^{*}v,\mathrm{Rop}^{-1}\cdot\mathrm{Qa}^{*}v\right)^{\mathrm{ch}}\ge0,\\
|\mathrm{Cq}|_{\mathrm{op}}\le|\mathrm{Qa}|_{\mathrm{op}}\cdot\left|\mathrm{Rop}^{-1}\right|_{\mathrm{op}}\cdot\left|\mathrm{Qa}^{*}\right|_{\mathrm{op}}\le\frac{1}{\mu}.
\end{gather*}
It follows that $\operatorname{Re}\mathrm{Nq}\ge b_k^{-1}>0$, and (w1) with $\mathrm{cf}=b_k^{-1}$ shows that $\mathrm{Nq}$ is invertible with $\left|\mathrm{Nq}^{-1}\right|_{\mathrm{op}}\le b_k$.

\emph{The inverse $S_k(A)$.}
By Section~4 one has $D_A+\mu+b_k\cdot J\cdot P_k(A)\cdot J^{*}=\mathrm{Rop}+b_k\cdot\mathrm{Qa}^{*}\cdot\mathrm{Qa}$.
Multiplying out and inserting $\mathrm{Qa}\cdot\mathrm{Rop}^{-1}\cdot\mathrm{Qa}^{*}=\mathrm{Cq}$, one obtains
\begin{align*}
&\left(\mathrm{Rop}+b_k\cdot\mathrm{Qa}^{*}\cdot\mathrm{Qa}\right)\cdot\left(\mathrm{Rop}^{-1}-\mathrm{Rop}^{-1}\cdot\mathrm{Qa}^{*}\cdot\mathrm{Nq}^{-1}\cdot\mathrm{Qa}\cdot\mathrm{Rop}^{-1}\right)\\
&\quad=1-\mathrm{Qa}^{*}\cdot\mathrm{Nq}^{-1}\cdot\mathrm{Qa}\cdot\mathrm{Rop}^{-1}+b_k\cdot\mathrm{Qa}^{*}\cdot\mathrm{Qa}\cdot\mathrm{Rop}^{-1}\\
&\qquad-b_k\cdot\mathrm{Qa}^{*}\cdot\mathrm{Cq}\cdot\mathrm{Nq}^{-1}\cdot\mathrm{Qa}\cdot\mathrm{Rop}^{-1}\\
&\quad=1+\mathrm{Qa}^{*}\cdot\left[b_k\cdot\mathrm{Nq}-1-b_k\cdot\mathrm{Cq}\right]\cdot\mathrm{Nq}^{-1}\cdot\mathrm{Qa}\cdot\mathrm{Rop}^{-1}=1,
\end{align*}
where the bracket vanishes because $b_k\cdot\mathrm{Nq}=1+b_k\cdot\mathrm{Cq}$.
On a space of finite dimension a right inverse is the inverse; therefore
\begin{equation}
S_k(A)=\mathrm{Rop}^{-1}-\mathrm{Rop}^{-1}\cdot\mathrm{Qa}^{*}\cdot\mathrm{Nq}^{-1}\cdot\mathrm{Qa}\cdot\mathrm{Rop}^{-1}.\tag{19.7}
\end{equation}

\emph{The operator $D_k(A)$.}
The operator $\mathrm{Nq}-\mathrm{Cq}$ is the number $b_k^{-1}$.
Multiplying (19.7) by $\mathrm{Qa}$ on the left and by $\mathrm{Qa}^{*}$ on the right, one gets
\[
\mathrm{Qa}\cdot S_k(A)\cdot\mathrm{Qa}^{*}=\mathrm{Cq}-\mathrm{Cq}\cdot\mathrm{Nq}^{-1}\cdot\mathrm{Cq}=\mathrm{Cq}\cdot\mathrm{Nq}^{-1}\cdot\left(\mathrm{Nq}-\mathrm{Cq}\right)=b_k^{-1}\cdot\mathrm{Cq}\cdot\mathrm{Nq}^{-1}.
\]
As $\mathrm{Qa}^{*}=J\cdot Q_k^{T}(-A)$ by (b), it follows that
\begin{align}
D_k(A)&=b_k-b_k^2\cdot\mathrm{Qa}\cdot S_k(A)\cdot\mathrm{Qa}^{*}=b_k-b_k\cdot\mathrm{Cq}\cdot\mathrm{Nq}^{-1}\notag\\
&=b_k\cdot\left(\mathrm{Nq}-\mathrm{Cq}\right)\cdot\mathrm{Nq}^{-1}=\mathrm{Nq}^{-1}.\tag{19.8}
\end{align}
Thus $D_k(A)$ is invertible, and in particular $\det D_k(A)\neq0$.
Moreover
\begin{gather*}
D_k(A)^{-1}=\mathrm{Nq}=b_k^{-1}+Q_k(A)\cdot J^{*}\cdot\left(D_A+\mu\right)^{-1}\cdot J\cdot Q_k^{T}(-A),\\
\left|D_k(A)^{-1}\right|_{\mathrm{op}}\le b_k^{-1}+|\mathrm{Cq}|_{\mathrm{op}}\le b_k^{-1}+\mu^{-1},\qquad
\left|D_k(A)\right|_{\mathrm{op}}=\left|\mathrm{Nq}^{-1}\right|_{\mathrm{op}}\le b_k,
\end{gather*}
and, for every $v\in\mathfrak{S}(1,\ell)$,
\[
\operatorname{Re}\left(v,D_k(A)^{-1}v\right)=b_k^{-1}\cdot(v,v)+\operatorname{Re}\left(v,\mathrm{Cq}\,v\right)\ge b_k^{-1}\cdot(v,v),
\]
which is the inequality $\operatorname{Re}D_k(A)^{-1}\ge b_k^{-1}$.
\end{proof}

\subsubsection*{7. Proof of (e) ($k\ge1$, $\mu\le12\mathrm{e}$)}

\begin{proof}[Proof of Proposition~19\,(e)]
\textbf{7.1. Blocks.}
For a point $y\in\mathbb{Z}^3$ and a vector $\mathrm{va}\in\mathbb{C}^4$ let $\mathrm{dlf}_{y,\mathrm{va}}$ be the element of $\mathfrak{S}(1,\ell)$ given by
\[
\mathrm{dlf}_{y,\mathrm{va}}(y+\ell z):=\boldsymbol{\chi}(z)\cdot\mathrm{va}\quad\left(z\in\mathbb{Z}^3\right),\qquad
\mathrm{dlf}_{y,\mathrm{va}}(y''):=0\quad\text{for }y''\notin y+\ell\mathbb{Z}^3 .
\]
Let $O$ be a linear operator of $\mathfrak{S}(1,\ell)$ into itself, and let $O(y,y')$ be its blocks in the sense of Section~I.1.6:
\[
(O\,v)(y)=\sum_{y''}O(y,y'')\cdot v(y''),\qquad O(y,y''+\ell e_\mu)=-O(y,y''),
\]
with $y''$ ranging over a set of representatives of $\mathbb{Z}^3/\ell\mathbb{Z}^3$.
Choosing the set of representatives so that it contains $y'$, one reads off $\left(O\,\mathrm{dlf}_{y',\mathrm{vb}}\right)(y)=O(y,y')\cdot\mathrm{vb}$.
Choosing the set of representatives in the scalar product $(v',v)$ so that it contains $y$, one reads off $\left(\mathrm{dlf}_{y,\mathrm{va}},v\right)=\overline{\mathrm{va}}\cdot v(y)$.
Hence, for all $y,y'\in\mathbb{Z}^3$ and all $\mathrm{va},\mathrm{vb}\in\mathbb{C}^4$,
\begin{equation}
\begin{aligned}
&\overline{\mathrm{va}}\cdot O(y,y')\cdot\mathrm{vb}=\left(\mathrm{dlf}_{y,\mathrm{va}},O\,\mathrm{dlf}_{y',\mathrm{vb}}\right),\qquad\left\|\mathrm{dlf}_{y,\mathrm{va}}\right\|=|\mathrm{va}|_2,\\
&|O(y,y')|_2=\sup\left\{\left|\left(\mathrm{dlf}_{y,\mathrm{va}},O\,\mathrm{dlf}_{y',\mathrm{vb}}\right)\right|:\ |\mathrm{va}|_2=|\mathrm{vb}|_2=1\right\}.
\end{aligned}
\tag{19.9}
\end{equation}
For $z\in\mathbb{Z}^3$ one has $\mathrm{dlf}_{y+\ell z,\mathrm{va}}=\boldsymbol{\chi}(z)\cdot\mathrm{dlf}_{y,\mathrm{va}}$.
Indeed, both functions vanish outside $y+\ell\mathbb{Z}^3$, and at the point $y+\ell z+\ell z''$, $z''\in\mathbb{Z}^3$, both take the value $\boldsymbol{\chi}(z'')\cdot\mathrm{va}$, because $\boldsymbol{\chi}(z)\cdot\boldsymbol{\chi}(z+z'')=\boldsymbol{\chi}(z'')$.
Since $\boldsymbol{\chi}(z)$ is real, (19.9) now gives
\[
O(y+\ell z,y'+\ell z')=\boldsymbol{\chi}(z)\cdot\boldsymbol{\chi}(z')\cdot O(y,y')\qquad\left(z,z'\in\mathbb{Z}^3\right).
\]
A change of the points $y$, $y'$ by vectors of $\ell\mathbb{Z}^3$ therefore multiplies $O(y,y')$ by $1$ or by $-1$.
In particular $|O(y,y')|_2$, $|O(y,y')|_c$ and the moduli of the entries of $O(y,y')$ are functions of the unit sites of $T^0_k$.
For $O=1$ one finds $\left(\mathrm{dlf}_{y,\mathrm{va}},\mathrm{dlf}_{y',\mathrm{vb}}\right)=0$ if $y-y'\notin\ell\mathbb{Z}^3$, and $\left(\mathrm{dlf}_{y,\mathrm{va}},\mathrm{dlf}_{y',\mathrm{vb}}\right)=\boldsymbol{\chi}(z)\cdot\overline{\mathrm{va}}\cdot\mathrm{vb}$ if $y=y'+\ell z$.
In both cases the modulus does not exceed $\mathbf{1}[y=y']\cdot|\mathrm{va}|_2\cdot|\mathrm{vb}|_2$, where the equality $y=y'$ is understood in $T^0_k$.

\textbf{7.2. The weight.}
Put $\vartheta:=\frac{\mu}{12\mathrm{e}}$.
The hypothesis $\mu\le12\mathrm{e}$ gives $0<\vartheta\le1$.
Fix $y'\in\mathbb{Z}^3$, and define for $c\in\mathbb{L}_k$
\[
\psi(c):=\vartheta\cdot\min\left\{|c-c'|:\ c'\in\Delta_{y'}+\ell\mathbb{Z}^3\right\};
\]
here the minimum is the sup distance on the torus from $c$ to the block of $y'$, and it is attained because the numbers $|c-c'|$ belong to $\eta_k\cdot\{0,1,2,\dots\}$.
The function $\psi$ is real and periodic, and $\psi=0$ on $\Delta_{y'}+\ell\mathbb{Z}^3$.
Let $c,\tilde c\in\mathbb{L}_k$, and let $c'$ be a point at which the minimum for $\tilde c$ is attained.
Then
\[
\psi(c)\le\vartheta\cdot|c-c'|\le\vartheta\cdot|c-\tilde c|+\vartheta\cdot|\tilde c-c'|=\vartheta\cdot|c-\tilde c|+\psi(\tilde c).
\]
This inequality holds for every pair of points, hence also with $c$ and $\tilde c$ interchanged, and so
\begin{gather*}
|\psi(c)-\psi(\tilde c)|\le\vartheta\cdot|c-\tilde c|;\\
\text{in particular}\quad|\psi(c)-\psi(c_{3})|\le\vartheta\cdot\eta_k\quad\text{for every }x=(c,g),
\end{gather*}
because $|c_{3}-c|=\eta_k$ (paragraph~1.1).
Now let $c\in\Delta_y+\ell\mathbb{Z}^3$ and $c'\in\Delta_{y'}+\ell\mathbb{Z}^3$.
There exist $z,z'\in\mathbb{Z}^3$ with $|c-y-\ell z|<\frac12$ and $|c'-y'-\ell z'|<\frac12$, and therefore
\[
|c-c'|>\left|y-y'+\ell\cdot(z-z')\right|-1\ge d(y,y')-1 .
\]
Consequently
\[
\psi(c)\ge\vartheta\cdot\left(d(y,y')-1\right)\qquad\text{for all }c\in\Delta_y+\ell\mathbb{Z}^3 .
\]

\textbf{7.3. The conjugated operator.}
Define
\[
\mathrm{Hps}:=\exp(\psi)\cdot D_A\cdot\exp(-\psi)-D_A,\qquad
\mathrm{Rps}:=\exp(\psi)\cdot\mathrm{Rop}\cdot\exp(-\psi)=\mathrm{Rop}+\mathrm{Hps};
\]
the last equality holds because the number $\mu$ commutes with $\exp(\psi)$.
For a real number $\mathrm{xr}$ one has $|\exp(\mathrm{xr})-1|\le\exp(|\mathrm{xr}|)-1$; for $\mathrm{xr}<0$ this is the inequality $1-\exp(-|\mathrm{xr}|)\le\exp(|\mathrm{xr}|)-1$, which holds because $\exp(|\mathrm{xr}|)+\exp(-|\mathrm{xr}|)\ge2$.
The right side increases with $|\mathrm{xr}|$.
Let $f\in\mathfrak{S}^{\mathrm{ch}}(\eta_k,\ell)$.
From (19.4), $|u_A(c,c_{3})|=1$, paragraphs~1.2 and~7.2, and (19.1) for $i=3$ and $\mathrm{pf}(x)=|f(x)|_2^2$, one obtains
\begin{align*}
\|\mathrm{Hps}\,f\|^2&=\int\mathrm{d}x\,\left|(\mathrm{Hps}\,f)(x)\right|_2^2\\
&\le\eta_k^{-2}\cdot\left(3(1+\rho)\right)^2\cdot\left(\exp(\vartheta\cdot\eta_k)-1\right)^2\cdot\int\mathrm{d}x\,\left|f\left(x^{(3)}\right)\right|_2^2\\
&=\left[3(1+\rho)\cdot\eta_k^{-1}\cdot\left(\exp(\vartheta\cdot\eta_k)-1\right)\right]^2\cdot\|f\|^2 .
\end{align*}
For $\mathrm{xr}\ge0$,
\[
\exp(\mathrm{xr})-1=\sum_{j\ge1}\frac{\mathrm{xr}^j}{j!}=\mathrm{xr}\cdot\sum_{j\ge0}\frac{\mathrm{xr}^j}{(j+1)!}\le\mathrm{xr}\cdot\exp(\mathrm{xr}).
\]
Moreover $\vartheta\cdot\eta_k\le1$, since $\vartheta\le1$ and $\eta_k=L^{-k}\le1$ ($L\ge1$, $k\ge0$).
Therefore
\begin{align}
|\mathrm{Hps}|_{\mathrm{op}}&\le3(1+\rho)\cdot\eta_k^{-1}\cdot\left(\exp(\vartheta\cdot\eta_k)-1\right)\le6\eta_k^{-1}\cdot\left(\exp(\vartheta\cdot\eta_k)-1\right)\notag\\
&\le6\vartheta\cdot\exp(\vartheta\cdot\eta_k)\le6\mathrm{e}\cdot\vartheta=\frac{\mu}{2}.\tag{19.10}
\end{align}
By (a) and (19.10),
\[
\operatorname{Re}\left(f,\mathrm{Rps}\,f\right)^{\mathrm{ch}}\ge\operatorname{Re}\left(f,\mathrm{Rop}\,f\right)^{\mathrm{ch}}-\left|\left(f,\mathrm{Hps}\,f\right)^{\mathrm{ch}}\right|\ge\mu\cdot\|f\|^2-\frac{\mu}{2}\cdot\|f\|^2=\frac{\mu}{2}\cdot\|f\|^2 .
\]
By (w1) with $\mathrm{cf}=\frac{\mu}{2}$, the operator $\mathrm{Rps}$ is invertible and $\left|\mathrm{Rps}^{-1}\right|_{\mathrm{op}}\le\frac{2}{\mu}$.
Moreover $\mathrm{Rps}^{-1}=\exp(\psi)\cdot\mathrm{Rop}^{-1}\cdot\exp(-\psi)$; equivalently, $\mathrm{Rop}^{-1}=\exp(-\psi)\cdot\mathrm{Rps}^{-1}\cdot\exp(\psi)$.

\textbf{7.4. The estimate.}
Let $y\in\mathbb{Z}^3$, let $\mathrm{va},\mathrm{vb}\in\mathbb{C}^4$ with $|\mathrm{va}|_2=|\mathrm{vb}|_2=1$, and put
\[
\mathrm{fa}:=\mathrm{Qa}^{*}\,\mathrm{dlf}_{y,\mathrm{va}},\qquad\mathrm{fb}:=\mathrm{Qa}^{*}\,\mathrm{dlf}_{y',\mathrm{vb}},
\]
with the point $y'$ fixed in paragraph~7.2.
By (19.6),
\[
\mathrm{fa}(c,g)=\varphi_A([c],c)^{-1}\cdot\mathrm{dlf}_{y,\mathrm{va}}([c]),\qquad
\mathrm{fb}(c,g)=\varphi_A([c],c)^{-1}\cdot\mathrm{dlf}_{y',\mathrm{vb}}([c]).
\]
So $\mathrm{fa}(c,g)=0$ unless $[c]\in y+\ell\mathbb{Z}^3$, that is, unless $c\in\Delta_y+\ell\mathbb{Z}^3$; and $\mathrm{fb}(c,g)=0$ unless $[c]\in y'+\ell\mathbb{Z}^3$, that is, unless $c\in\Delta_{y'}+\ell\mathbb{Z}^3$.
The operator $\mathrm{Qa}^{*}$ is isometric by (b); with (19.9) this gives $\|\mathrm{fa}\|=\|\mathrm{fb}\|=1$.
Paragraph~7.2 yields
\begin{gather*}
\exp(\psi)\,\mathrm{fb}=\mathrm{fb},\qquad\text{so that}\qquad\left\|\exp(\psi)\,\mathrm{fb}\right\|=1;\\
\left\|\exp(-\psi)\,\mathrm{fa}\right\|\le\exp\left(-\vartheta\cdot\left(d(y,y')-1\right)\right)\cdot\|\mathrm{fa}\|=\exp\left(-\vartheta\cdot\left(d(y,y')-1\right)\right).
\end{gather*}
The multiplication operator $\exp(-\psi)$ is Hermitian.
Hence, by paragraph~7.3,
\begin{align*}
\left|\left(\mathrm{dlf}_{y,\mathrm{va}},\mathrm{Cq}\,\mathrm{dlf}_{y',\mathrm{vb}}\right)\right|
&=\left|\left(\mathrm{fa},\mathrm{Rop}^{-1}\,\mathrm{fb}\right)^{\mathrm{ch}}\right|
=\left|\left(\exp(-\psi)\,\mathrm{fa},\mathrm{Rps}^{-1}\cdot\exp(\psi)\,\mathrm{fb}\right)^{\mathrm{ch}}\right|\\
&\le\frac{2}{\mu}\cdot\exp\left(-\vartheta\cdot\left(d(y,y')-1\right)\right)
=\frac{2}{\mu}\cdot\exp(\vartheta)\cdot\exp\left(-\vartheta\cdot d(y,y')\right)\\
&\le\frac{2\mathrm{e}}{\mu}\cdot\exp\left(-\vartheta\cdot d(y,y')\right),
\end{align*}
where the last step uses $\vartheta\le1$.
By (d), $D_k(A)^{-1}=b_k^{-1}+\mathrm{Cq}$; together with paragraph~7.1 this gives
\[
\left|\left(\mathrm{dlf}_{y,\mathrm{va}},D_k(A)^{-1}\,\mathrm{dlf}_{y',\mathrm{vb}}\right)\right|\le b_k^{-1}\cdot\mathbf{1}[y=y']+\frac{2\mathrm{e}}{\mu}\cdot\exp\left(-\frac{\mu}{12\mathrm{e}}\cdot d(y,y')\right).
\]
Taking the supremum over $\mathrm{va}$ and $\mathrm{vb}$ and using (19.9), one arrives at
\begin{equation}
\left|D_k(A)^{-1}(y,y')\right|_2\le b_k^{-1}\cdot\mathbf{1}[y=y']+\frac{2\mathrm{e}}{\mu}\cdot\exp\left(-\frac{\mu}{12\mathrm{e}}\cdot d(y,y')\right).\tag{19.11}
\end{equation}
As $y'$ was arbitrary, (19.11) holds for all $y,y'\in\mathbb{Z}^3$.

\textbf{7.5. The norm $|\cdot|_c$; entries.}
Let $m$ be a $4\times4$ matrix, and let $\mathrm{ub}_1,\dots,\mathrm{ub}_4$ be the standard basis of $\mathbb{C}^4$.
For an entry,
\[
\left|m_{\beta\beta'}\right|=\left|\overline{\mathrm{ub}_\beta}\cdot m\cdot\mathrm{ub}_{\beta'}\right|\le|m|_2 .
\]
For a row, the Cauchy--Schwarz inequality gives
\[
\sum_{\beta'=1}^{4}\left|m_{\beta\beta'}\right|\le4^{\frac12}\cdot\left(\sum_{\beta'}\left|m_{\beta\beta'}\right|^2\right)^{\frac12}=2\left|m^{*}\cdot\mathrm{ub}_\beta\right|_2\le2\left|m^{*}\right|_2=2|m|_2 ,
\]
and for a column,
\[
\sum_{\beta=1}^{4}\left|m_{\beta\beta'}\right|\le2\left(\sum_{\beta}\left|m_{\beta\beta'}\right|^2\right)^{\frac12}=2\left|m\cdot\mathrm{ub}_{\beta'}\right|_2\le2|m|_2 .
\]
Hence $|m|_c$, the largest row sum (Section~I.1.5), is at most $2|m|_2$.
Taking $m=D_k(A)^{-1}(y,y')$, one obtains from (19.11)
\[
\left|D_k(A)^{-1}(y,y')\right|_c\le2b_k^{-1}\cdot\mathbf{1}[y=y']+\frac{4\mathrm{e}}{\mu}\cdot\exp\left(-\frac{\mu}{12\mathrm{e}}\cdot d(y,y')\right),
\]
and every entry of the block $D_k(A)^{-1}(y,y')$ has modulus at most the right side of (19.11).
\end{proof}

\subsubsection*{8. Proof of (f) ($k\ge1$, $\mu\le12\mathrm{e}$)}

The sums over rows and columns are bounded by means of the lattice sum (G2.r6) of Section~G2.0.5.

\begin{named}{(G2.r6)}[{\cite[Section~13.1.6]{QEDI}}]
For $\gamma>0$ let $c(\gamma):=\sum_{z\in\mathbb{Z}^3}\exp(-\gamma\cdot|z|)$, with $|z|$ the sup norm.
\begin{itemize}
\item[(i)] For every unit site $y$,
\[
\sum_{y'\in T^0_k}\exp\left(-\gamma\cdot d(y,y')\right)\le c(\gamma).
\]
\item[(ii)] If $\gamma\le6$, then $c(\gamma)\le\left(\frac{12}{\gamma}\right)^3$.
\end{itemize}
\end{named}

\begin{proof}[Verification]
(i) For every class $y'$ of $\mathbb{Z}^3/\ell\mathbb{Z}^3$ choose a point $z(y')\in\mathbb{Z}^3$ that represents $y'-y$ and satisfies $|z(y')|=d(y,y')$.
Distinct classes give distinct points $z(y')$.
The sum in (i) is therefore a sum of some of the terms of the series $c(\gamma)$, all of which are positive.

(ii) The inequality $|z|\ge\frac13|z|_1$ gives
\[
c(\gamma)\le\left(\sum_{j\in\mathbb{Z}}\exp\left(-\frac{\gamma\cdot|j|}{3}\right)\right)^3 .
\]
Summing the geometric series,
\[
\sum_{j\in\mathbb{Z}}\exp\left(-\frac{\gamma\cdot|j|}{3}\right)=1+2\sum_{j\ge1}\exp\left(-\frac{\gamma\cdot j}{3}\right)=1+2\left(\exp\left(\frac{\gamma}{3}\right)-1\right)^{-1}\le1+\frac{6}{\gamma},
\]
since $\exp\left(\frac{\gamma}{3}\right)-1\ge\frac{\gamma}{3}$.
For $\gamma\le6$ one has $1\le\frac{6}{\gamma}$, hence $1+\frac{6}{\gamma}\le\frac{12}{\gamma}$.
\end{proof}

\begin{proof}[Proof of Proposition~19\,(f)]
\emph{Rows and columns.}
The entries of the matrix of $D_k(A)^{-1}$ are the numbers $D_k(A)^{-1}(y,y')_{\beta\beta'}$, where $y$ and $y'$ range over a set of representatives of $\mathbb{Z}^3/\ell\mathbb{Z}^3$ and $\beta,\beta'\in\{1,\dots,4\}$.
Let $(y,\beta)$ be an index.
Paragraph~7.5 (rows), (19.11) and (G2.r6) with $\gamma=\vartheta=\frac{\mu}{12\mathrm{e}}\le1\le6$ give
\begin{align*}
\sum_{y'}\sum_{\beta'=1}^{4}\left|D_k(A)^{-1}(y,y')_{\beta\beta'}\right|
&\le\sum_{y'}2\left|D_k(A)^{-1}(y,y')\right|_2
\le2b_k^{-1}+\frac{4\mathrm{e}}{\mu}\cdot\sum_{y'}\exp\left(-\vartheta\cdot d(y,y')\right)\\
&\le2b_k^{-1}+\frac{4\mathrm{e}}{\mu}\cdot c(\vartheta)
\le2b_k^{-1}+\frac{4\mathrm{e}}{\mu}\cdot\left(\frac{12}{\vartheta}\right)^3,
\end{align*}
and $\frac{12}{\vartheta}=\frac{144\mathrm{e}}{\mu}$.
Let $(y',\beta')$ be an index.
Paragraph~7.5 (columns), (19.11), the symmetry $d(y,y')=d(y',y)$ and (G2.r6) give
\begin{align*}
\sum_{y}\sum_{\beta=1}^{4}\left|D_k(A)^{-1}(y,y')_{\beta\beta'}\right|
&\le\sum_{y}2\left|D_k(A)^{-1}(y,y')\right|_2
\le2b_k^{-1}+\frac{4\mathrm{e}}{\mu}\cdot\sum_{y}\exp\left(-\vartheta\cdot d(y',y)\right)\\
&\le2b_k^{-1}+\frac{4\mathrm{e}}{\mu}\cdot c(\vartheta)
\le2b_k^{-1}+\frac{4\mathrm{e}}{\mu}\cdot\left(\frac{12}{\vartheta}\right)^3 .
\end{align*}
Hence every row and every column of the matrix of $D_k(A)^{-1}$ has $\ell^1$ norm at most
\[
2b_k^{-1}+4\mathrm{e}\cdot\mu^{-1}\cdot\left(\frac{144\mathrm{e}}{\mu}\right)^3 .
\]
\end{proof}

\subsubsection*{9. Uniformity; the twist}

\begin{proof}[Proof of the last sentence of Proposition~19]
The right sides of the bounds in (a), (d), (e) and (f) are, in this order,
\begin{gather*}
\frac{1}{\mu};\qquad b_k^{-1}+\mu^{-1},\quad b_k;\\
b_k^{-1}\cdot\mathbf{1}[y=y']+\frac{2\mathrm{e}}{\mu}\cdot\exp\left(-\frac{\mu}{12\mathrm{e}}\cdot d(y,y')\right)\ \text{and twice this expression};\\
2b_k^{-1}+4\mathrm{e}\cdot\mu^{-1}\cdot\left(\frac{144\mathrm{e}}{\mu}\right)^3;
\end{gather*}
the norms in (b) are equal to $1$.
These expressions are built from $\mu$, $b_k$, $d(y,y')$ and the numbers $2$, $4$, $12$, $144$, $\mathrm{e}$.
Neither $A$ nor $e_k$ nor $\rho$ occurs in them, and the side $\ell$ occurs only inside the distance $d(y,y')$ on the torus; no index $N$ of a spacing is among the data.
The level enters only through $b_k$.
For all $k\ge1$ one has $0<1-L^{-1}\le1-L^{-k}<1$, and Section~I.3.4 gives
\[
b_{-}=b\cdot\left(1-L^{-1}\right)\le b_k\le b=b_{+};
\]
hence $b_k^{-1}\le b_{-}^{-1}$ and $b_k\le b_{+}$ for all $k\ge1$.

In the proofs of (a)--(f) the bond function $A$ and the charge $e_k$ occur only in the numbers $u_A(c,c_{3})$ and $\varphi_A(y,c)$, and the arguments rest on three properties of these numbers alone: their modulus is $1$, $u_A(c_{3},c)=u_A(c,c_{3})^{-1}$, and they are periodic (paragraphs~1.3 and~1.8).
These properties hold for every real bond function $A$ on $T_k$ and every $e_k>0$.
Hence (a)--(f) hold, with the right sides listed above, for all real bond functions $A$ on $T_k$ and all $e_k>0$; in particular the bounds are uniform in $A$.

\emph{The twist.}
Let $\theta\in\mathbb{R}^3$.
The twist connection $c^{(k)}_{\theta}$ of Section~I.6.6 is the constant connection with the components $\frac{\theta_\mu}{e_k\cdot\ell_k}$, where $\ell_k>0$ is the side of the torus at the level $k$ (Section~I.6.1).
Thus it takes the value $\frac{\theta_\mu}{e_k\cdot\ell_k}$ on the bonds $(c,c+\eta_k e_\mu)$ and the value $-\frac{\theta_\mu}{e_k\cdot\ell_k}$ on the bonds $(c+\eta_k e_\mu,c)$.
It is a real periodic bond function.
Consequently, for every real bond function $A$ on $T_k$, the sum $A+c^{(k)}_{\theta}$ is again a real bond function on $T_k$, and (a)--(f) hold for it with the same right sides, for all $\theta\in\mathbb{R}^3$.
\end{proof}

\subsection*{Conditions on the constants}

\begin{description}
\item[$\rho\ge0$]
It enters in paragraph~1.2 ($|\rho-1|\le1+\rho$) and in Section~3 ($\frac{\rho}{2\eta_k}\cdot\left(f,\nabla_A^{*}\nabla_A f\right)^{\mathrm{ch}}\ge0$).
It is stated in Section~I.8.2, line (C0); in Proposition~14.2\,(a), row $\rho$ ($\rho\in(0,1)$); in Annex~B, Section~G2.9, first entry.
\item[$\rho\le1$]
It enters in paragraph~1.2 ($3(1+\rho)\le6$), hence in (19.10).
It is stated in the same places as $\rho\ge0$.
\item[The standard weights, with $w_i>0$, $\sum_iw_i\hat n_i=0$ (I.1.1) and $w_3=6$]
They enter in paragraph~1.5; in paragraph~1.1, (I.s3), hence in (19.3); in paragraph~1.2.
They are stated in Section~I.8.2, line (C0); in Proposition~14.2\,(a), row $w$; in Annex~B, Section~G2.9, second entry.
\item[$b>0$ ($b=1$) and $L>1$, so that $b_k>0$ and $b_{-}\le b_k\le b_{+}$]
They enter in paragraph~1.8 ($L^{j+1-k}\le1$ for $j<k$); in Sections~5 and~6; in Section~9; in paragraph~7.3 ($\eta_k=L^{-k}\le1$).
They are stated in Section~I.8.2, line (C0), row $b$, and lines (C3), (C4); in Section~I.3.4; in Proposition~14.2\,(a), rows $b$, $L$, $b_{-}$, $b_{+}$;
and in Annex~B, Section~G2.9, where the third entry gives $b=1$ and the entry on the first condition of (C3) gives that $L$ is odd and $L\ge5$, hence $L>1$.
\item[$L$ an odd integer]
It enters in paragraph~0.1 ($\mathbb{Z}^3\subset\mathbb{L}_k$); in paragraph~1.7 (no site of $\mathbb{L}_k$ lies on a face of a unit block; a unit block has $L^k$ points in each direction); in paragraph~1.8 ($L^{j-k}\mathbb{Z}^3\subset\mathbb{L}_k$; no site of $\mathbb{L}_k$ lies on a face of a block of the side $L^j\eta_k$, $0\le j\le k$; $L^{k-j}\ell z$ is an integer vector; $\frac{L+1}{2}$ is an integer).
It is stated in Section~I.0.5; in Section~I.8.2, line (C3); in Proposition~14.2\,(a), row $L$; and in Annex~B, Section~G2.9, in the entry on the first condition of (C3).
\item[$\ell\in\mathrm{Cl}$]
Of this condition the proofs need that $\ell$ is a positive integer: $\ell$ is an odd multiple of $LM$, where $M$, a power of $L$, is a positive integer.
It enters in paragraph~0.1, through the period lattice $\ell\mathbb{Z}^3$.
In Proposition~19 the side $\ell$ ranges over the class $\mathrm{Cl}$, which is described in Annex~B, Section~G2.9, in the entry on the class of sides;
the conditions on $M$ are stated in Section~I.0.5, in Section~I.8.2, line (C5), and in Proposition~14.2\,(a), row $M$.
\end{description}

\section*{Chapter 30. Free determinants and twists: proofs of Propositions 30.1--30.3}

\subsection*{0. Conventions of this chapter}

\paragraph{0.1.}
The conventions of Part~I hold throughout.
For $x\in\mathbb{R}^3$, $|x|=|x|_\infty$ denotes the maximum norm and $|x|_2$ the Euclidean norm;
$|v|_2$ is the Euclidean norm of a vector $v$ of $\mathbb{C}^4$ or of $\mathbb{C}^2$,
and $|m|_2$ is the operator norm of a matrix $m$ with respect to $|\cdot|_2$.
One sets $\mathrm{e}:=\exp(1)$;
thus $\mathrm{e}^{\pm i(\cdot)}$ and $\mathrm{e}^{\pm 2\pi i(\cdot)}$ are values of the exponential function.
Directions are denoted by $\mu,\nu\in\{1,2,3\}$.
For functions of a momentum $p$ one writes $\partial_\mu:=\frac{\partial}{\partial p_\mu}$,
and for a multi-index $\mathrm{ak}=(\mathrm{ak}_1,\mathrm{ak}_2,\mathrm{ak}_3)$
\[
\partial^{\mathrm{ak}}:=\partial_1^{\mathrm{ak}_1}\partial_2^{\mathrm{ak}_2}\partial_3^{\mathrm{ak}_3},\qquad
|\mathrm{ak}|:=\mathrm{ak}_1+\mathrm{ak}_2+\mathrm{ak}_3,\qquad
\mathrm{ak}!:=\mathrm{ak}_1!\,\mathrm{ak}_2!\,\mathrm{ak}_3!\,.
\]
If a function is holomorphic in $p$, then $\partial_\mu$ is the complex derivative;
restricted to a translate $ip''+\mathbb{R}^3$ of the real space, it is the derivative with respect to the real part.

\paragraph{0.2. Linear algebra.}
All Hilbert spaces are complex and of finite dimension.
For a linear map $O$ between two of them, $\|O\|$ is the operator norm and $O^*$ the adjoint.
For a linear map $O$ of a space into itself one puts $\operatorname{Re}O:=\frac12\left(O+O^*\right)$,
and for real $t$ the inequality $\operatorname{Re}O\ge t$ means that $\operatorname{Re}(f,Of)\ge t\,(f,f)$ for all $f$;
$\det O$ and $\operatorname{Tr}O$ are the determinant and the trace of $O$ as a linear map, so that neither a basis nor a scalar product enters.
\begin{itemize}
\item[(H1)]
If $\operatorname{Re}O\ge t>0$, then $O$ is bijective, $\|O^{-1}\|\le\frac{1}{t}$ and $\operatorname{Re}O^{-1}\ge0$.
Indeed, $t\,\|f\|^2\le\operatorname{Re}(f,Of)\le\|f\|\cdot\|Of\|$ for every $f$.
Hence $O$ is injective, therefore bijective, and $\|O^{-1}f'\|\le\frac{\|f'\|}{t}$.
For $f'=Of$ one has $\operatorname{Re}(f',O^{-1}f')=\operatorname{Re}(Of,f)=\operatorname{Re}(f,Of)\ge0$.
\item[(H2)]
For a linear map $O$ of a space $\mathrm{Vs}$ into itself,
\[
|\operatorname{Tr}O|\le\dim(\mathrm{Vs})\cdot\|O\|,
\]
because $\operatorname{Tr}O=\sum_i(\mathrm{oa}_i,O\,\mathrm{oa}_i)$ for every orthonormal basis $(\mathrm{oa}_i)$ of $\mathrm{Vs}$.
\item[(H3)]
(Jacobi's formula, \cite[Section~0.8]{HornJohnson2013}.)
Let $t\mapsto O(t)$ be a differentiable, resp.\ holomorphic, family of invertible maps, and $O':=\frac{\mathrm{d}}{\mathrm{d}t}O$.
Then
\[
\frac{\mathrm{d}}{\mathrm{d}t}\det O=\det O\cdot\operatorname{Tr}\left(O^{-1}O'\right),\qquad
\frac{\mathrm{d}}{\mathrm{d}t}O^{-1}=-O^{-1}O'O^{-1}.
\]
For the first formula, observe that $\det(1+O_1)-1-\operatorname{Tr}O_1$ is a polynomial in the entries of a matrix of $O_1$ which has no terms of degree $\le1$, and that
\[
\det O(t')=\det O(t)\cdot\det\left(1+O(t)^{-1}\left(O(t')-O(t)\right)\right);
\]
subtract $\det O(t)$, divide by $t'-t$ and let $t'\to t$.
By Cramer's rule $O^{-1}$ is differentiable, resp.\ holomorphic,
and the second formula follows on differentiating $O\,O^{-1}=1$.
\item[(H4)]
(Sylvester's determinant identity, \cite[Section~0.8.5]{HornJohnson2013}.)
For linear maps $O_1:\mathrm{Vs}_1\to\mathrm{Vs}_2$ and $O_2:\mathrm{Vs}_2\to\mathrm{Vs}_1$,
\[
\det\left(1+O_1O_2\right)=\det\left(1+O_2O_1\right),
\]
the determinant on the left being taken on $\mathrm{Vs}_2$ and the one on the right on $\mathrm{Vs}_1$.
Both sides are equal to the determinant of
\[
\left(\begin{array}{cc}1&-O_2\\ O_1&1\end{array}\right)
\]
on $\mathrm{Vs}_1\oplus\mathrm{Vs}_2$, as one sees by eliminating the one or the other off-diagonal block.
\item[(H5)]
(Singular values.)
Let $\mathrm{sgv}_1(O)\ge\mathrm{sgv}_2(O)\ge\cdots\ge0$ be the eigenvalues of $\left(O^*O\right)^{\frac12}$, counted with multiplicity, and for $\mathrm{sr}\ge1$ let
\[
\|O\|_{(\mathrm{sr})}:=\left(\sum_j\mathrm{sgv}_j(O)^{\mathrm{sr}}\right)^{\frac{1}{\mathrm{sr}}}.
\]
\begin{itemize}
\item[(i)]
$\mathrm{sgv}_j(O_1O_2)\le\|O_1\|\cdot\mathrm{sgv}_j(O_2)$ and $\mathrm{sgv}_j(O_1O_2)\le\mathrm{sgv}_j(O_1)\cdot\|O_2\|$.
To see this, recall that $\mathrm{sgv}_j(O)$ is the minimum, over the subspaces $\mathrm{Vs}$ of codimension $j-1$, of the maximum of $\frac{\|Of\|}{\|f\|}$ over $f\in\mathrm{Vs}\setminus\{0\}$
(the Courant--Fischer theorem for $O^*O$, \cite[Corollary~III.1.2]{Bhatia1997});
moreover $\|O_1O_2f\|\le\|O_1\|\cdot\|O_2f\|$, and $\mathrm{sgv}_j(O)=\mathrm{sgv}_j(O^*)$.
\item[(ii)]
$|\operatorname{Tr}O|\le\|O\|_{(1)}$;
$\|O_1O_2\|_{(1)}\le\|O_1\|_{(1)}\cdot\|O_2\|$ and $\|O_1O_2\|_{(1)}\le\|O_1\|\cdot\|O_2\|_{(1)}$;
$\|O\|_{(1)}\le\operatorname{rank}(O)\cdot\|O\|$.
For the proof, the singular value decomposition \cite[Section~2.6]{HornJohnson2013} gives
$O=\sum_j\mathrm{sgv}_j(O)\,\mathrm{oa}_j\,(\mathrm{ob}_j,\cdot\,)$ with orthonormal systems $(\mathrm{oa}_j)$ and $(\mathrm{ob}_j)$,
so that $\operatorname{Tr}O=\sum_j\mathrm{sgv}_j(O)\,(\mathrm{ob}_j,\mathrm{oa}_j)$;
the bounds for the products follow from (i);
and $\operatorname{rank}(O)$ is the number of the non-zero singular values, each of which is at most $\|O\|$.
\item[(iii)]
(H\"older's inequality.)
For every integer $\mathrm{sr}\ge1$,
\[
\|O_1O_2\cdots O_{\mathrm{sr}}\|_{(1)}\le\prod_{i=1}^{\mathrm{sr}}\|O_i\|_{(\mathrm{sr})}.
\]
This results from $\mathrm{sr}-1$ applications of \cite[Theorem~2.8]{Simon2005} (see also \cite[Section~IV.2]{Bhatia1997}), in the form
\[
\|O_1\cdots O_j\|_{(\mathrm{sr}/j)}\le\|O_1\cdots O_{j-1}\|_{(\mathrm{sr}/(j-1))}\cdot\|O_j\|_{(\mathrm{sr})}\qquad(2\le j\le\mathrm{sr}).
\]
\item[(iv)]
The singular values of a direct sum are those of the summands.
Singular values, norms, traces and determinants are invariant under $O\mapsto\mathrm{Vun}^{-1}O\,\mathrm{Vun}$ for unitary $\mathrm{Vun}$.
\end{itemize}
\end{itemize}

\paragraph{0.3. Fourier series.}
For a continuous $2\pi\mathbb{Z}^3$-periodic function $\mathrm{fper}$ on $\mathbb{R}^3$ let, as in Section~I.2.6,
\[
\mathrm{fperh}(z):=\int_{[-\pi,\pi]^3}\frac{\mathrm{d}^3p}{(2\pi)^3}\,\mathrm{fper}(p)\,\mathrm{e}^{-ip\cdot z}\qquad(z\in\mathbb{Z}^3).
\]
\begin{itemize}
\item[(F1)]
If $\sum_z|\mathrm{fperh}(z)|<\infty$, then $\mathrm{fper}(p)=\sum_z\mathrm{fperh}(z)\,\mathrm{e}^{ip\cdot z}$ for every $p$.
The series converges uniformly to a continuous periodic function with the same Fourier coefficients,
so that its difference with $\mathrm{fper}$ is a continuous periodic function all of whose Fourier coefficients vanish.
Such a function $\mathrm{fz}$ is zero.
By the Stone--Weierstrass theorem \cite[Theorem~7.33]{Rudin1976}, applied to the algebra of the trigonometric polynomials on the compact space $\mathbb{R}^3/2\pi\mathbb{Z}^3$
(this algebra is closed under conjugation, separates the points and contains $1$; in one variable see \cite[Theorem~4.25]{Rudin1987}),
there are trigonometric polynomials converging to $\mathrm{fz}$ uniformly.
Hence $\int_{[-\pi,\pi]^3}\mathrm{d}^3p\,|\mathrm{fz}(p)|^2$ is the limit of the integrals of $\mathrm{fz}$ against the complex conjugates of these polynomials, and all of these integrals vanish.
\item[(F2)]
For an integer $t\ge1$,
\[
\#|\{w\in\mathbb{Z}^3:|w|=t\}|=(2t+1)^3-(2t-1)^3=24t^2+2\le26t^2.
\]
\item[(F3)]
For $x>0$,
\[
\sum_{t\ge1}\frac{\exp(-xt)}{t}=\log\left(\frac{1}{1-\exp(-x)}\right).
\]
\item[(F4)]
By (F2),
\[
S_4:=\sum_{w\in\mathbb{Z}^3,\ w\neq0}|w|^{-4}=\sum_{t\ge1}\left(24t^2+2\right)t^{-4}=24\cdot\frac{\pi^2}{6}+2\cdot\frac{\pi^4}{90}=4\pi^2+\frac{\pi^4}{45}.
\]
\item[(F5)]
(Integration by parts.)
Let $\mathrm{fper}$ be $2\pi\mathbb{Z}^3$-periodic and of class $C^{\mathrm{kd}}$ on $\mathbb{R}^3$, $\mathrm{kd}\ge1$ an integer, and let $z\in\mathbb{Z}^3$ with $z_\nu\neq0$.
Then
\[
\mathrm{fperh}(z)=(iz_\nu)^{-\mathrm{kd}}\int_{[-\pi,\pi]^3}\frac{\mathrm{d}^3p}{(2\pi)^3}\left(\partial_\nu^{\mathrm{kd}}\mathrm{fper}\right)(p)\,\mathrm{e}^{-ip\cdot z}.
\]
This follows from $\mathrm{kd}$ integrations by parts in $p_\nu$ over $[-\pi,\pi]$;
the boundary terms cancel, since $\partial_\nu^j\mathrm{fper}$ and $\mathrm{e}^{-ip\cdot z}$ are $2\pi$-periodic in $p_\nu$.
\item[(F6)]
(Shift of one contour.)
Let $\nu\in\{1,2,3\}$ and $\mathrm{wd}>0$, and let $\mathrm{fhol}$ be a continuous function on the set
\[
\left\{p\in\mathbb{C}^3:\ p_\mu\in\mathbb{R}\ \text{for}\ \mu\neq\nu,\ |\operatorname{Im}p_\nu|\le\mathrm{wd}\right\}
\]
with $\mathrm{fhol}(p+2\pi e_\mu)=\mathrm{fhol}(p)$ for $\mu=1,2,3$,
such that for all fixed real $p_\mu$, $\mu\neq\nu$, the function which maps $\mathrm{zs}$ to the value of $\mathrm{fhol}$ at the point with the coordinates $p_\mu$ ($\mu\neq\nu$) and $p_\nu=\mathrm{zs}$ is holomorphic in a neighbourhood of the closed strip $|\operatorname{Im}\mathrm{zs}|\le\mathrm{wd}$.
Let $z\in\mathbb{Z}^3$ with $z_\nu\neq0$, and let $\mathrm{sgz}$ be the sign of $z_\nu$.
Then
\[
\int_{[-\pi,\pi]^3}\mathrm{d}^3p\,\mathrm{fhol}(p)\,\mathrm{e}^{-ip\cdot z}
=\exp\left(-\mathrm{wd}\,|z_\nu|\right)\int_{[-\pi,\pi]^3}\mathrm{d}^3p\,\mathrm{fhol}\left(p-i\,\mathrm{sgz}\cdot\mathrm{wd}\,e_\nu\right)\mathrm{e}^{-ip\cdot z}.
\]
For the proof fix $p_\mu$, $\mu\neq\nu$, and apply Cauchy's theorem \cite[Chapter~10]{Rudin1987} to the function $\mathrm{zs}\mapsto\mathrm{fhol}(\ldots)\,\mathrm{e}^{-iz_\nu\mathrm{zs}}$ on the rectangle with the corners $-\pi$, $\pi$, $\pi-i\,\mathrm{sgz}\cdot\mathrm{wd}$, $-\pi-i\,\mathrm{sgz}\cdot\mathrm{wd}$.
The integrals over the two vertical sides cancel, because the integrand has the period $2\pi$ in $\mathrm{zs}$.
On the lower side if $\mathrm{sgz}=1$, on the upper side if $\mathrm{sgz}=-1$, one has, with real $\mathrm{zs}$,
\[
\mathrm{e}^{-iz_\nu(\mathrm{zs}-i\,\mathrm{sgz}\cdot\mathrm{wd})}=\mathrm{e}^{-iz_\nu\mathrm{zs}}\exp\left(-\mathrm{wd}\,|z_\nu|\right).
\]
It remains to integrate over the two other coordinates.
\end{itemize}

\paragraph{0.4. Notation.}
In this chapter $\mathrm{me}>0$ denotes a mesh, and
$\mathrm{me}N:=\{\mathrm{me}\,n:n=1,2,\ldots\}$, $\mathrm{et}N:=\{\mathrm{et}\,n:n=1,2,\ldots\}$.
As in Section~G1.10 and Proposition~1.2, $\mathrm{ak}$ is a multi-index,
$\mathrm{tw}(p):=\sum_\mu\left(1-\cos p_\mu\right)$ and $\mathrm{rp}(p):=\left(|p|_2^2+\mathbf{m}^2\right)^{\frac12}$;
a complex momentum is written $p=p'+ip''$ with $p',p''\in\mathbb{R}^3$, as in Proposition~1.2\,(S7);
$D^{\mathrm{hh}}(p)$ and $\mathrm{Dsc}_{\mathbf{m}}(p)$ are the heavy block and the Schur symbol of Proposition~1.2.
In the proof of Proposition~30.1\,(b) in Section~2, however, $p'$ and $p''$ denote elements of $P_{n,0}$.
For multi-indices, $\binom{\mathrm{ak}}{\mathrm{ak}'}:=\prod_\mu\binom{\mathrm{ak}_\mu}{\mathrm{ak}'_\mu}$.
The operator $\mathrm{Gh}$ is $J_A$ of (I.2.4) at $A=0$;
for a bond function $\mathrm{cst}$, $D_{\mathrm{cst}}$ and $u_{\mathrm{cst}}$ are $D_A$ and $u_A$ of (I.2.3) and (I.1.6) at $A=\mathrm{cst}$.
The letters $\mathrm{Flg}$, $\mathrm{Flgh}$, $\mathrm{jf}$, $p_{\mathrm{jf}}$, $\mathrm{Uth}$, $\mathrm{Uth}^{\mathrm{u}}$, $\mathrm{Dfl}_\theta$, $\mathrm{Rtw}_\theta$, $D^{(p';\mathrm{ma})}$, $[c]$ and the constants of Propositions~30.1 and~30.3 have the meaning given to them there.
Further, $\hat n$ and $S_{\mathrm{nv}}$ are as in Section~I.4.1;
$\boldsymbol{\Gamma}$, $\boldsymbol{\Gamma}^{\pi}$, $\tau$, $\tau_k$ and $\varphi_A$ as in Sections~I.1.1 and~I.3.2;
$\xi$ as in Section~I.3.1;
and in $f_\beta$ the index $\beta$ is the spinor index of Section~I.3.6.
In~4.7 the letters $\alpha$ (a real number) and $\beta$ (a vector of $\mathbb{R}^3$) are those of Proposition~30.3\,(g2).
The Pauli matrices of Section~I.2.1 are $\sigma_1,\sigma_2,\sigma_3$, and $\sigma\cdot v:=\sum_\mu v_\mu\sigma_\mu$ for $v\in\mathbb{C}^3$, formed like $\gamma\cdot v$.
An upper index $+$ or $-$ on an operator that commutes with $\Sigma_3$ denotes its part in the eigenspace $\Sigma_3=1$, resp.\ $\Sigma_3=-1$;
the index $+$ is that of Proposition~30.3\,(a).
The upper indices $\mathrm{l}$ and $\mathrm{h}$ on an element of ${V_{\mathrm{v}}}$ denote its light and its heavy part;
$(0)$ and $(\mathrm{r})$ on an element of $\mathrm{Hb}(p)$ denote its mode $\mathrm{jf}=0$ and the sum of its other modes;
$(\mathrm{kd})$ on $\mathrm{Dpg}$, $\mathrm{Ppg}$, $\mathrm{Tpg}$ denotes the derivative of order $\mathrm{kd}$.
The letter $i$ has two meanings, which its position separates.
Where it is a lower index, the upper index in $x^{(i)}$, or the variable of a range of such an index (as in $i\le2$, $i=0,1,2$, $\sum_i$ and $\prod_{i=1}^{\mathrm{sr}}$), it is the type of a bond or an index of a sum or of a product.
Wherever it is a factor, on the line or in an exponent, it is $\sqrt{-1}$.
This is the case in front of $\gamma_\mu$, $\sigma_\mu$ and the matrices formed with them, such as $\gamma\cdot v$, $\sigma\cdot v$ and $\operatorname{diag}(\sigma\cdot v,-\sigma\cdot v)$;
in the argument of the exponential function in each of its forms, $\mathrm{e}^{\pm i(\cdot)}$, $\mathrm{e}^{\pm2\pi i(\cdot)}$ and $\exp(\pm i\,\cdot)$;
in complex momenta and in shifted contours, such as $p'+ip''$, $ip''+\mathbb{R}^3$, $p-i\,\mathrm{sgz}\cdot\mathrm{wd}\,e_\nu$ and the corners $\pm\pi-i\,\mathrm{sgz}\cdot\mathrm{wd}$;
and in factors such as $(iz_\nu)^{-\mathrm{kd}}$, $-i(ge_1)_\mu$, $\left(i\,\mathrm{et}\,(ge_1)_\nu\right)^{\mathrm{kd}}$, $\left(-i(c_\nu-c'_\nu)\right)^{\mathrm{kd}}$ and $(-i)^{\mathrm{kd}}$.
The letter $a$ denotes a vector of $\{\pm e_1,\pm e_2,\pm e_3\}$,
and a prime on $c$, $g$, $f$, $u$, $v$, $\mathrm{ak}$, $\mathrm{jf}$, $p$ denotes a further element of the same kind.
In this chapter (8.1) and (8.2) denote the displays so labelled in Proposition~30.1\,(c) and~(d), and (8.5) the display so labelled in Proposition~30.3\,(e);
(g1), (g2), (g3) are the three parts of clause~(g) of Proposition~30.3.

\paragraph{0.5. Classical theorems: references.}
Jacobi's formula for the derivative of a determinant, (H3): \cite[Section~0.8]{HornJohnson2013}.
Sylvester's determinant identity, (H4): \cite[Section~0.8.5]{HornJohnson2013}.
The Courant--Fischer theorem, (H5)\,(i): \cite[Corollary~III.1.2]{Bhatia1997}.
The singular value decomposition, (H5)\,(ii): \cite[Section~2.6]{HornJohnson2013}.
H\"older's inequality for Schatten norms, (H5)\,(iii): \cite[Theorem~2.8]{Simon2005} and \cite[Section~IV.2]{Bhatia1997}.
Simultaneous unitary diagonalisation of commuting normal matrices, Lemma~G\,(G5): \cite[Section~2.5]{HornJohnson2013}.
Cauchy's theorem in one variable, (F6): \cite[Chapter~10]{Rudin1987}.
Cauchy's integral formula and Cauchy's estimates on polydiscs, step~(c3) of Section~2: \cite[Section~2.2]{Hormander1990}.
Density of the trigonometric polynomials in the continuous periodic functions, (F1): \cite[Theorem~4.25]{Rudin1987} in one variable, and the Stone--Weierstrass theorem \cite[Theorem~7.33]{Rudin1976} in three variables.

\subsection*{1. Lemma A: the operator at zero connection on quasi-periodic functions}

\paragraph{Data and definitions.}
Let $\mathrm{me}>0$, $R\in\mathrm{me}N$ and $\theta\in\mathbb{R}^3$.
The space $\mathrm{HH}(\mathrm{me},R;\theta)$, with $R$ finite here,
the neighbours $x^{(i)}$ of $x=(c,g)$
and the operators $R_i$, $D$, $\nabla^{*}\nabla$ at the mesh $\mathrm{me}$
are those of Section~G1.4.
The scalar product $(f',f)$ is the product $(f',f)^{\mathrm{ch}}$ of Section~G1.4,
\[
(f',f)=\mathrm{me}^3\cdot\frac{1}{48}\sum_c\sum_{g\in O_{\mathrm{cub}}}\overline{f'(c,g)}\cdot f(c,g),
\]
where $c$ ranges over a period cell of $R\mathbb{Z}^3$ in $\mathrm{me}\,\mathbb{Z}^3$.
Since the multipliers $-\mathrm{e}^{i\theta_\mu}$ have modulus one, the summand is $R\mathbb{Z}^3$-periodic, and the choice of the cell does not matter.
The dimension of $\mathrm{HH}(\mathrm{me},R;\theta)$ is $192\left(\frac{R}{\mathrm{me}}\right)^3$, and $\|f\|:=(f,f)^{\frac12}$.
The space $\mathrm{HH2}(\mathrm{me},R;\theta)$ is defined in the same way, with values in $\mathbb{C}^2$ in the place of $\mathbb{C}^4$, and its scalar product is given by the same formula.

The vectors $\hat d_i(x)=-g\hat n_i$ and the matrices $M_i(x)=\frac{w_i}{2}\left(\rho-\gamma\cdot\hat d_i(x)\right)$ are those of (I.1.2) and (I.2.1);
neither contains the mesh.
On functions on $\mathrm{me}\,\mathbb{Z}^3\times O_{\mathrm{cub}}$ let
\begin{gather*}
(\mathrm{Gd}_if)(x):=\gamma\cdot\hat d_i(x)\,f(x),\qquad
\mathrm{Gh}:=\sum_{i=0}^{3}w_i\,\mathrm{Gd}_i\,R_i,\\
(J^{*}f)(c):=\frac{1}{48}\sum_{g\in O_{\mathrm{cub}}}f(c,g),\qquad
(J\xi)(c,g):=\xi(c),\qquad
\mathrm{Plt}:=JJ^{*},\qquad
\mathrm{Phv}:=1-\mathrm{Plt}.
\end{gather*}
The operator $D$ is (I.2.3) at the mesh $\mathrm{me}$ and at $A=0$.
The symbol $D^{(p;\mathrm{ma})}$ is (I.2.6) at the momentum $p\in\mathbb{C}^3$ and at $\boldsymbol{\mu}=\mathbf{m}$;
it acts on ${V_{\mathrm{v}}}=\{u:O_{\mathrm{cub}}\to\mathbb{C}^4\}$, which carries the scalar product
\[
(u',u):=\frac{1}{48}\sum_{g\in O_{\mathrm{cub}}}\overline{u'(g)}\cdot u(g),
\]
and $J$, $J^{*}$, $\mathrm{Plt}$, $\mathrm{Phv}$ act on ${V_{\mathrm{v}}}$ by the same formulas, without the variable $c$.

\begin{named}{Lemma A}
\begin{itemize}
\item[(Lz1)]
The map $x\mapsto x^{(i)}$ is an involution of $\mathrm{me}\,\mathbb{Z}^3\times O_{\mathrm{cub}}$ which commutes with the translations $(c,g)\mapsto(c+v,g)$, $v\in\mathrm{me}\,\mathbb{Z}^3$;
moreover $\hat d_i(x^{(i)})=-\hat d_i(x)$ and $\sum_{i=0}^{3}w_i\,\hat d_i(x)=0$.
The operators $R_i$, $\mathrm{Gd}_i$, $D$, $\nabla^{*}\nabla$, $\mathrm{Gh}$, $\mathrm{Plt}$, $\mathrm{Phv}$ map $\mathrm{HH}(\mathrm{me},R;\theta)$ into itself.
\item[(Lz2)]
On $\mathrm{HH}(\mathrm{me},R;\theta)$ the operator $R_i$ is unitary, $R_i^2=1$ and $R_i^*=R_i$;
further $\mathrm{Gd}_i^*=\mathrm{Gd}_i$, $R_i\,\mathrm{Gd}_i=-\mathrm{Gd}_i\,R_i$ and $\sum_iw_i\,\mathrm{Gd}_i=0$.
\item[(Lz3)]
One has
\[
D=\mathrm{me}^{-1}\left[\frac{\rho}{2}\nabla^{*}\nabla+\frac12\mathrm{Gh}\right];
\]
$\left(\nabla^{*}\nabla\right)^*=\nabla^{*}\nabla$ and $(f,\nabla^{*}\nabla f)=\sum_{i=0}^{3}\frac{w_i}{2}\|f-R_if\|^2\ge0$;
$\mathrm{Gh}^*=-\mathrm{Gh}$.
Consequently $\operatorname{Re}D=\frac{\rho}{2\,\mathrm{me}}\nabla^{*}\nabla\ge0$ and
\[
D^*=\mathrm{me}^{-1}\left[\frac{\rho}{2}\nabla^{*}\nabla-\frac12\mathrm{Gh}\right]=\gamma_5\,D\,\gamma_5.
\]
\item[(Lz4)]
$\mathrm{Plt}$ is the orthogonal projection onto the functions that do not depend on $g$, and it commutes with constant spin matrices.
One has $(f,\nabla^{*}\nabla f)\ge\lambda_{\mathrm{int}}\|\mathrm{Phv}\,f\|^2$, and therefore
\[
\operatorname{Re}(f,Df)\ge\frac{\hat g}{\mathrm{me}}\,\|\mathrm{Phv}\,f\|^2,\qquad\hat g=\frac{\rho\,\lambda_{\mathrm{int}}}{2}.
\]
\item[(Lz5)]
(Blocks.)
Along $\mathbb{C}^4=\mathbb{C}^2\oplus\mathbb{C}^2$ one has $\mathrm{HH}(\mathrm{me},R;\theta)=\mathrm{HH2}(\mathrm{me},R;\theta)\oplus\mathrm{HH2}(\mathrm{me},R;\theta)$;
the first summand is the eigenspace $\Sigma_3=1$, the second the eigenspace $\Sigma_3=-1$.
An operator $O$ which commutes with $\Sigma_3$ has the form $\operatorname{diag}(O^+,O^-)$ with operators $O^\pm$ on $\mathrm{HH2}(\mathrm{me},R;\theta)$.
The operator $D$ commutes with $\Sigma_3$, and
\begin{gather*}
D^+=\mathrm{me}^{-1}\left[\frac{\rho}{2}\nabla^{*}\nabla+\frac12\mathrm{Gh2}\right],\qquad
D^-=\mathrm{me}^{-1}\left[\frac{\rho}{2}\nabla^{*}\nabla-\frac12\mathrm{Gh2}\right]=\left(D^+\right)^*,\\
(\mathrm{Gh2}\,f)(x):=\sum_{i=0}^{3}w_i\,\sigma\cdot\hat d_i(x)\,f(x^{(i)}),\qquad
\mathrm{Gh2}^*=-\mathrm{Gh2}.
\end{gather*}
\item[(Lz6)]
Let $t>0$, and let $\mathrm{Pn}=\operatorname{diag}(\mathrm{Pn}^+,\mathrm{Pn}^+)$ with a Hermitian operator $\mathrm{Pn}^+\ge0$ on $\mathrm{HH2}(\mathrm{me},R;\theta)$.
Then $O:=D+t+\mathrm{Pn}$ has the following properties:
$\operatorname{Re}O=\frac{\rho}{2\,\mathrm{me}}\nabla^{*}\nabla+t+\mathrm{Pn}\ge t$;
$O$ is bijective and $\|O^{-1}\|\le\frac{1}{t}$;
$O=\operatorname{diag}\left(O^+,\left(O^+\right)^*\right)$;
$\det O=|\det O^+|^2>0$.
\item[(Lz7)]
(Plane waves.)
For $p\in\mathbb{C}^3$ and $u\in{V_{\mathrm{v}}}$ let $f(c,g):=\mathrm{e}^{ip\cdot c}\,u(g)$ on $\mathrm{me}\,\mathbb{Z}^3\times O_{\mathrm{cub}}$.
Then
\[
(Df)(c,g)=\mathrm{e}^{ip\cdot c}\,\mathrm{me}^{-1}\left(D^{(\mathrm{me}\,p;0)}u\right)(g).
\]
For real $p$ the map $\mathrm{Em}_p:{V_{\mathrm{v}}}\to\mathrm{HH}\left(1,1;p-(\pi,\pi,\pi)\right)$, $(\mathrm{Em}_pu)(c,g):=\mathrm{e}^{ip\cdot c}\,u(g)$, is unitary,
and $D^{(p;\mathrm{ma})}=\mathrm{Em}_p^{-1}(D+\mathbf{m})\,\mathrm{Em}_p$, where $D$ is taken at $\mathrm{me}=1$.
Consequently, for real $p$ and real $\mathbf{m}\ge0$,
\[
\operatorname{Re}\left(u,D^{(p;\mathrm{ma})}u\right)\ge\mathbf{m}\,\|u\|^2+\hat g\,\|\mathrm{Phv}\,u\|^2,\qquad
\left(D^{(p;\mathrm{ma})}\right)^*=\gamma_5\,D^{(p;\mathrm{ma})}\,\gamma_5;
\]
$D^{(p;\mathrm{ma})}$ commutes with $\Sigma_3$ and is equal to $\operatorname{diag}\left(\left(D^{(p;\mathrm{ma})}\right)^+,\left(\left(D^{(p;\mathrm{ma})}\right)^+\right)^*\right)$;
and $\det D^{(p;\mathrm{ma})}=\left|\det\left(D^{(p;\mathrm{ma})}\right)^+\right|^2$, a number which is $>0$ if $\mathbf{m}>0$.
\item[(Lz8)]
(A constant connection is a twist; (I.6.1) at the mesh $\mathrm{me}$.)
Let $\mathrm{ec}>0$, let $\mathrm{cst}$ be the constant real bond function on $\mathrm{me}\,\mathbb{Z}^3$ with $\mathrm{cst}(c,c+\mathrm{me}\,e_\mu)=\frac{\theta_\mu}{\mathrm{ec}\,R}$,
and let $D_{\mathrm{cst}}$ be the operator (I.2.3) with $(\eta_k,e_k)=(\mathrm{me},\mathrm{ec})$ at $A=\mathrm{cst}$.
The map $(\mathrm{Utw}f)(c,g):=\mathrm{e}^{i\frac{\theta\cdot c}{R}}f(c,g)$ maps $\mathrm{HH}(\mathrm{me},R;0)$ unitarily onto $\mathrm{HH}(\mathrm{me},R;\theta)$,
and $D_{\mathrm{cst}}=\mathrm{Utw}^{-1}D\,\mathrm{Utw}$ on $\mathrm{HH}(\mathrm{me},R;0)$.
In particular, $D_{\mathrm{cst}}$ does not depend on $\mathrm{ec}$.
\item[(Lz9)]
(Dilation.)
The map $(\mathrm{Sca}f)(c',g):=f(\mathrm{me}\,c',g)$, $c'\in\mathbb{Z}^3$, is a linear bijection of $\mathrm{HH}(\mathrm{me},R;\theta)$ onto $\mathrm{HH}(1,R/\mathrm{me};\theta)$,
and $\mathrm{Sca}\,D=\mathrm{me}^{-1}D'\,\mathrm{Sca}$, where $D'$ is the operator $D$ at the mesh $1$.
\end{itemize}
\end{named}

\begin{proof}
\begin{itemize}
\item[(Lz1)]
Since $s_i^2=1$, one has $\left(x^{(i)}\right)^{(i)}=x$ for $i\le2$.
Since $s_3e_1=-e_1$,
\[
\left(x^{(3)}\right)^{(3)}=\left(c+\mathrm{me}\,ge_1+\mathrm{me}\,gs_3e_1,\ gs_3s_3\right)=(c,g).
\]
The translations act on the first entry only, and the first entry of $x^{(i)}$ minus that of $x$ depends on $g$ alone.
Next, $\hat d_i(x^{(i)})=-gs_i\hat n_i=g\hat n_i$, because $s_i\hat n_i=-\hat n_i$ (Section~I.1.2);
and $\sum_iw_i\,\hat d_i(x)=-g\sum_iw_i\hat n_i=0$ by (I.1.1).
The coefficients $\gamma\cdot\hat d_i(x)$ and $M_i(x)$ depend on $g$ only.
Hence all the operators listed commute with the translations by $Re_\mu$, and they preserve the condition $f(c+Re_\mu,g)=-\mathrm{e}^{i\theta_\mu}f(c,g)$.
\item[(Lz2)]
For $f',f\in\mathrm{HH}(\mathrm{me},R;\theta)$ the function $\mathrm{Fs}(x):=\overline{f'(x)}\cdot f(x)$ is $R\mathbb{Z}^3$-periodic,
and by (Lz1) the map $x\mapsto x^{(i)}$ induces a bijection of $\left(\mathrm{me}\,\mathbb{Z}^3/R\mathbb{Z}^3\right)\times O_{\mathrm{cub}}$ onto itself.
Therefore, with $x$ ranging over this set,
\[
(R_if',R_if)=\mathrm{me}^3\cdot\frac{1}{48}\sum_x\mathrm{Fs}\left(x^{(i)}\right)=\mathrm{me}^3\cdot\frac{1}{48}\sum_x\mathrm{Fs}(x)=(f',f).
\]
By (Lz1), $R_i^2=1$, and a unitary involution is Hermitian.
For real $v$ the matrix $\gamma\cdot v$ is Hermitian, by (I.c1); this gives $\mathrm{Gd}_i^*=\mathrm{Gd}_i$.
Further
\[
(R_i\,\mathrm{Gd}_if)(x)=\gamma\cdot\hat d_i(x^{(i)})\,f(x^{(i)})=-\gamma\cdot\hat d_i(x)\,f(x^{(i)})=-(\mathrm{Gd}_i\,R_if)(x),
\]
and $\sum_iw_i\,\mathrm{Gd}_i$ is the multiplication by $\gamma\cdot\left(\sum_iw_i\,\hat d_i(x)\right)=0$.
\item[(Lz3)]
By (I.2.3) at $A=0$ and the definition of $M_i(x)$,
\[
\mathrm{me}\,D=\sum_i\frac{w_i}{2}\left(\rho-\mathrm{Gd}_i\right)\left(1-R_i\right)=\frac{\rho}{2}\nabla^{*}\nabla-\frac12\sum_iw_i\,\mathrm{Gd}_i+\frac12\sum_iw_i\,\mathrm{Gd}_i\,R_i,
\]
and this equals $\frac{\rho}{2}\nabla^{*}\nabla+\frac12\mathrm{Gh}$ by (Lz2).
The relation $\left(\nabla^{*}\nabla\right)^*=\nabla^{*}\nabla$ follows from (Lz2).
The number $(f,R_if)$ is real, and $\|f-R_if\|^2=2\|f\|^2-2(f,R_if)=2\left(f,(1-R_i)f\right)$.
Moreover $\mathrm{Gh}^*=\sum_iw_i\,R_i\,\mathrm{Gd}_i=-\sum_iw_i\,\mathrm{Gd}_i\,R_i=-\mathrm{Gh}$.
Finally, $\gamma_5$ commutes with $R_i$ and anticommutes with $\gamma_\mu$, and $\gamma_5^2=1$, by (I.c1);
so $\gamma_5\nabla^{*}\nabla\gamma_5=\nabla^{*}\nabla$ and $\gamma_5\,\mathrm{Gh}\,\gamma_5=-\mathrm{Gh}$.
\item[(Lz4)]
For a $\mathbb{C}^4$-valued function $\xi$ of $c$ with the same multipliers,
\[
(f,J\xi)=\mathrm{me}^3\cdot\frac{1}{48}\sum_{c,g}\overline{f(c,g)}\cdot\xi(c)=\mathrm{me}^3\sum_c\overline{(J^{*}f)(c)}\cdot\xi(c).
\]
Thus $J^{*}$ is the adjoint of $J$ for the scalar product $\mathrm{me}^3\sum_c$ on these functions;
since $J^{*}J=1$, the operator $\mathrm{Plt}=JJ^{*}$ is an orthogonal projection, with the range described.
Let $\ell^2(O_{\mathrm{cub}})$ be the space of the complex functions on $O_{\mathrm{cub}}$ with the scalar product $\frac{1}{48}\sum_g$, and on it $(R_i\mathrm{us})(g):=\mathrm{us}(gs_i)$.
The computation made in (Lz3) shows that $\sum_{i\le2}w_i(1-R_i)$ is Hermitian on $\ell^2(O_{\mathrm{cub}})$, with the quadratic form $\sum_{i\le2}\frac{w_i}{2}\|\mathrm{us}-R_i\mathrm{us}\|^2$.
This form vanishes only if $\mathrm{us}(gs_i)=\mathrm{us}(g)$ for all $g$ and $i=0,1,2$;
since $s_0,s_1,s_2$ generate $O_{\mathrm{cub}}$ (Section~I.1.2), this happens only for constant $\mathrm{us}$.
Hence the kernel consists of the constants, and
\[
\sum_{i\le2}\frac{w_i}{2}\|\mathrm{us}-R_i\mathrm{us}\|^2\ge\lambda_{\mathrm{int}}\left\|\mathrm{us}-\frac{1}{48}\sum_g\mathrm{us}(g)\right\|^2,
\]
where $\lambda_{\mathrm{int}}$ is the smallest non-zero eigenvalue, see (I.O8).
Apply this to $\mathrm{us}(g):=f_\beta(c,g)$ for every $c$ and every spinor index $\beta$, multiply by $\mathrm{me}^3$ and sum;
the result is $\sum_{i\le2}\frac{w_i}{2}\|f-R_if\|^2\ge\lambda_{\mathrm{int}}\|\mathrm{Phv}\,f\|^2$.
The term $i=3$ in (Lz3) is $\ge0$.
The last assertion follows from $\operatorname{Re}(f,Df)=\frac{\rho}{2\,\mathrm{me}}(f,\nabla^{*}\nabla f)$.
\item[(Lz5)]
By Section~I.2.1, $\gamma\cdot v=\operatorname{diag}(\sigma\cdot v,-\sigma\cdot v)$ and $\Sigma_3=\operatorname{diag}(1,1,-1,-1)$;
and $R_i$ does not act on the spinor index.
Hence $\nabla^{*}\nabla=\operatorname{diag}(\nabla^{*}\nabla,\nabla^{*}\nabla)$ and $\mathrm{Gh}=\operatorname{diag}(\mathrm{Gh2},-\mathrm{Gh2})$.
The two summands are orthogonal, and on each of them the scalar product of $\mathrm{HH}(\mathrm{me},R;\theta)$ restricts to that of $\mathrm{HH2}(\mathrm{me},R;\theta)$;
so the adjoint of $\operatorname{diag}(O^+,O^-)$ is $\operatorname{diag}\left((O^+)^*,(O^-)^*\right)$.
Together with $\mathrm{Gh}^*=-\mathrm{Gh}$ and $\left(\nabla^{*}\nabla\right)^*=\nabla^{*}\nabla$ from (Lz3), this gives $\mathrm{Gh2}^*=-\mathrm{Gh2}$, and $\nabla^{*}\nabla$ is Hermitian on $\mathrm{HH2}(\mathrm{me},R;\theta)$.
Hence $\left(D^+\right)^*=\mathrm{me}^{-1}\left[\frac{\rho}{2}\nabla^{*}\nabla-\frac12\mathrm{Gh2}\right]=D^-$.
\item[(Lz6)]
The formula for $\operatorname{Re}O$ follows from (Lz3), and (H1) gives the bijectivity and the bound on $O^{-1}$.
By (Lz5), $O^-=D^-+t+\mathrm{Pn}^+=\left(D^++t+\mathrm{Pn}^+\right)^*$.
Hence $\det O=\det O^+\cdot\det\left(O^+\right)^*=\det O^+\cdot\overline{\det O^+}$, and $\det O\neq0$.
\item[(Lz7)]
One has $f(x^{(i)})=\mathrm{e}^{ip\cdot c}\,u(gs_i)$ for $i\le2$ and $f(x^{(3)})=\mathrm{e}^{ip\cdot c}\,\mathrm{e}^{i\,\mathrm{me}\,p\cdot(ge_1)}\,u(gs_3)$;
the comparison with (I.2.6) at $\boldsymbol{\mu}=0$ and at the momentum $\mathrm{me}\,p$ gives the first assertion.
Now let $p$ be real and $\mathrm{me}=R=1$.
Since $-\mathrm{e}^{i(p_\mu-\pi)}=\mathrm{e}^{ip_\mu}$, every $f\in\mathrm{HH}\left(1,1;p-(\pi,\pi,\pi)\right)$ satisfies $f(c,g)=\mathrm{e}^{ip\cdot c}f(0,g)$ for $c\in\mathbb{Z}^3$,
and the period cell is $\{0\}$;
so $\mathrm{Em}_p$ is a unitary bijection.
The remaining assertions are (Lz3), (Lz4), (Lz5) on $\mathrm{HH}\left(1,1;p-(\pi,\pi,\pi)\right)$,
where the formula for the determinant follows from the block form, because $\det\left(O^+\right)^*=\overline{\det O^+}$,
together with (Lz6) for $\mathrm{Pn}=0$ and $t=\mathbf{m}$, which gives the strict inequality when $\mathbf{m}>0$;
here $\mathrm{Em}_p$ commutes with constant spin matrices and with $\mathrm{Plt}$.
\item[(Lz8)]
For $f\in\mathrm{HH}(\mathrm{me},R;0)$ one has $(\mathrm{Utw}f)(c+Re_\mu,g)=\mathrm{e}^{i\theta_\mu}\,\mathrm{e}^{i\frac{\theta\cdot c}{R}}\left(-f(c,g)\right)$.
Since $\left|\mathrm{e}^{i\frac{\theta\cdot c}{R}}\right|=1$ and the two scalar products are given by one formula, $\mathrm{Utw}$ is unitary onto $\mathrm{HH}(\mathrm{me},R;\theta)$;
its inverse is the multiplication by $\mathrm{e}^{-i\frac{\theta\cdot c}{R}}$.
For $a\in\{\pm e_\mu\}$ one has $\mathrm{cst}(c,c+\mathrm{me}\,a)=\frac{\theta\cdot a}{\mathrm{ec}\,R}$, bond functions being odd under reversal;
hence, by (I.1.6) with $(\eta_k,e_k)=(\mathrm{me},\mathrm{ec})$,
\[
u_{\mathrm{cst}}(c,c_3)=\exp\left\{i\,\mathrm{ec}\cdot\mathrm{me}\cdot\mathrm{cst}(c,c_3)\right\}=\mathrm{e}^{i\,\mathrm{me}\frac{\theta\cdot(ge_1)}{R}}.
\]
On the other hand $(\mathrm{Utw}f)(x^{(i)})=\mathrm{e}^{i\frac{\theta\cdot c}{R}}f(x^{(i)})$ for $i\le2$ and $(\mathrm{Utw}f)(x^{(3)})=\mathrm{e}^{i\frac{\theta\cdot c}{R}}\,\mathrm{e}^{i\,\mathrm{me}\frac{\theta\cdot(ge_1)}{R}}f(x^{(3)})$.
Therefore
\begin{align*}
(D\,\mathrm{Utw}f)(x)
&=-\mathrm{me}^{-1}\,\mathrm{e}^{i\frac{\theta\cdot c}{R}}\left\{\sum_{i\le2}M_i(x)\left[f(x^{(i)})-f(x)\right]+M_3(x)\left[u_{\mathrm{cst}}(c,c_3)\,f(x^{(3)})-f(x)\right]\right\}\\
&=(\mathrm{Utw}\,D_{\mathrm{cst}}f)(x).
\end{align*}
\item[(Lz9)]
One has $(\mathrm{Sca}f)\left(c'+\frac{R}{\mathrm{me}}e_\mu,g\right)=f(\mathrm{me}\,c'+Re_\mu,g)$
and $(\mathrm{Sca}f)(c'+ge_1,gs_3)=f(\mathrm{me}\,c'+\mathrm{me}\,ge_1,gs_3)$.
Thus $\mathrm{Sca}\,R_i=R_i'\,\mathrm{Sca}$, where $R_i'$ are the hops at the mesh $1$;
and $M_i$ contains no mesh.
\end{itemize}
\end{proof}

\subsection*{2. Proposition 30.1}

\subsubsection*{Proof of Proposition 30.1}

\begin{proof}
For $p\in\mathbb{C}^3$ the symbol $D^{(p)}$ of (I.2.6) at $\boldsymbol{\mu}=\mathbf{m}$ is denoted by $D^{(p;\mathrm{ma})}$;
in the indices of $D^{(p;\mathrm{ma})}$ and of $\mathfrak{d}_{\mathrm{ma}}$ the letters $\mathrm{ma}$ stand for the mass $\mathbf{m}$.
One has $\mathfrak{S}^{\mathrm{ch}}(1;n)=\mathrm{HH}(1,n;0)$ (Section~1), and, by Section~I.2.6,
\[
P_{n,\theta}=\left\{p:\ p_\mu=\frac{2\pi j_\mu+\pi+\theta_\mu}{n},\ j\in\{0,\dots,n-1\}^3\right\}.
\]
For every $n\ge1$, all functions are taken on the covering lattice $\mathbb{Z}^3\times O_{\mathrm{cub}}$.

\textbf{(a)}
For every $g\in O_{\mathrm{cub}}$ the vector $ge_1$ is one of the vectors $\pm e_\mu$.
By (I.2.6), therefore, the matrix of $D^{(p;\mathrm{ma})}$ with respect to the basis of ${V_{\mathrm{v}}}$ formed by the delta functions has entries that are polynomials in $\mathrm{e}^{ip_\mu}$ and $\mathrm{e}^{-ip_\mu}$.
Consequently $\mathfrak{d}_{\mathrm{ma}}$ is an entire function of $p\in\mathbb{C}^3$, and it is $2\pi\mathbb{Z}^3$-periodic.
Since $\mathbf{m}>0$, (Lz7) shows that $\mathfrak{d}_{\mathrm{ma}}$ is real and positive on $\mathbb{R}^3$;
being entire and real valued on $\mathbb{R}^3$, it is real analytic there.
As $\log\mathfrak{d}_{\mathrm{ma}}$ is continuous on $[-\pi,\pi]^3$, the integral that defines $p^{0}_{\mathrm{ch}}(\mathbf{m})$ exists and is a real number.

\textbf{(b)}
For $p'\in P_{n,0}$ and $u\in{V_{\mathrm{v}}}$ put
\[
f_{p',u}(c,g):=\mathrm{e}^{ip'\cdot c}\,u(g).
\]
Because $\mathrm{e}^{inp'_\mu}=\mathrm{e}^{i(2\pi j_\mu+\pi)}=-1$, the function $f_{p',u}$ belongs to $\mathfrak{S}^{\mathrm{ch}}(1;n)$.
The linear map
\[
\bigoplus_{p'\in P_{n,0}}{V_{\mathrm{v}}}\longrightarrow\mathfrak{S}^{\mathrm{ch}}(1;n),\qquad\left(u_{p'}\right)_{p'}\longmapsto\sum_{p'}f_{p',u_{p'}},
\]
is injective.
To see this, assume that the sum vanishes, multiply it by $\mathrm{e}^{-ip''\cdot c}$ with $p''\in P_{n,0}$, and sum over $c\in\{0,\dots,n-1\}^3$.
Since
\[
\sum_{c}\mathrm{e}^{i(p'-p'')\cdot c}=\prod_{\mu}\sum_{c_\mu=0}^{n-1}\mathrm{e}^{\frac{2\pi i(j_\mu-j''_\mu)c_\mu}{n}}=n^3\,\mathbf{1}[p'=p''],
\]
one obtains $u_{p''}=0$.
The two spaces both have dimension $192\,n^3$; hence the map is bijective.

Let $a\in\{\pm e_\mu\}$.
A cochain is odd under reversal of the bond, so $H_\theta(c,c+a)=\frac{\theta\cdot a}{qn}$.
The phase in (I.2.3) is therefore
\[
u_{H_\theta}(c,c_3)=\exp\left(iq\,H_\theta(c,c+ge_1)\right)=\mathrm{e}^{i\frac{\theta\cdot(ge_1)}{n}},
\]
and
\begin{multline*}
\left(\left(D_{H_\theta}+\mathbf{m}\right)f_{p',u}\right)(c,g)
=\mathrm{e}^{ip'\cdot c}\Bigl\{\mathbf{m}\,u(g)-\sum_{i\le2}M_{i}(g)\bigl(u(gs_i)-u(g)\bigr)\\
-M_{3}(g)\Bigl(\mathrm{e}^{i\frac{\theta\cdot(ge_1)}{n}}\,\mathrm{e}^{ip'\cdot(ge_1)}\,u(gs_3)-u(g)\Bigr)\Bigr\}
=f_{p',\,D^{(p'+\theta/n;\mathrm{ma})}u}(c,g).
\end{multline*}
In the decomposition above, $D_{H_\theta}+\mathbf{m}$ is thus the direct sum of the maps $D^{(p'+\theta/n;\mathrm{ma})}$ of ${V_{\mathrm{v}}}$, $p'\in P_{n,0}$.
The determinant of a direct sum is the product of the determinants of the summands, and $p'\mapsto p'+\theta/n$ is a bijection of $P_{n,0}$ onto $P_{n,\theta}$; this is the product formula.
The positivity follows from~(a).

\textbf{(d)}
By hypothesis, $\sum_{z}|\mathrm{fperh}(z)|<\infty$, and $\mathrm{fper}(p)=\sum_{z\in\mathbb{Z}^3}\mathrm{fperh}(z)\,\mathrm{e}^{ip\cdot z}$ for every real $p$.
(The assertion in parentheses in~(d), that for a continuous $\mathrm{fper}$ the second property is a consequence of the first, is~(F1).)
For $z\in\mathbb{Z}^3$ one has
\[
\sum_{p\in P_{n,\theta}}\mathrm{e}^{ip\cdot z}=\prod_{\mu}\left[\mathrm{e}^{i\frac{z_\mu(\pi+\theta_\mu)}{n}}\sum_{j=0}^{n-1}\mathrm{e}^{\frac{2\pi ijz_\mu}{n}}\right].
\]
The inner sum equals $n$ when $n$ divides $z_\mu$ and vanishes otherwise.
The product therefore vanishes unless $z=nw$ with $w\in\mathbb{Z}^3$, and in that case it equals
\[
n^3\,\mathrm{e}^{i\pi(w_1+w_2+w_3)}\,\mathrm{e}^{i\theta\cdot w}=n^3\,(-1)^{w_1+w_2+w_3}\,\mathrm{e}^{i\theta\cdot w}.
\]
Insert the Fourier series of $\mathrm{fper}$ into the finite sum over $P_{n,\theta}$ and interchange the two summations, which absolute convergence permits; the result is the first identity in~(8.2).
The term with $w=0$ equals $n^3\,\mathrm{fperh}(0)$, and the other terms have the moduli $n^3\,|\mathrm{fperh}(nw)|$; this gives the inequality in~(8.2).

\textbf{(c)}
\emph{(c0) Constants.}
The numbers $\hat{g}$, $\theta_{0}$, $K_{\mathrm{Str}}$, $\kappa_{(0)}$ and $A'_{\mathrm{ak}}$ are those of Proposition~1.2, and they are functions of $\rho$;
by the definition of $\kappa_{(0)}$ one has $0<\kappa_{(0)}\le\theta_{0}/2$.
The absolute number $\mathrm{cF13}$ and the numbers $C_{\mathrm{lH}}$, $C_{\mathrm{lS}}$, $C_{6}$ are those defined in Proposition~30.1; they are functions of $\rho$.

\emph{(c1) Two factors.}
Let $D^{\mathrm{hh}}(p)$ and $\mathrm{Dsc}_{\mathbf{m}}(p)$ be the heavy block and the Schur symbol of $D^{(p;\mathrm{ma})}$, as in Proposition~1.2;
the first is a linear map of $\ker J^{*}$, a space of dimension $192-4=188$, and the second is a $4\times4$ matrix.
For real $p$, Proposition~1.2\,(S9) gives $\det D^{\mathrm{hh}}(p)>0$ and $\det\mathrm{Dsc}_{\mathbf{m}}(p)>0$, because $\mathbf{m}>0$.
In particular $\mathrm{Dsc}_{\mathbf{m}}(p)$ is invertible, and Proposition~1.2\,(S8) yields, for real $p$ and with real logarithms,
\begin{gather*}
\log\mathfrak{d}_{\mathrm{ma}}(p)=\mathrm{lgH}(p)+\mathrm{lgS}(p),\\
\mathrm{lgH}(p):=\log\det D^{\mathrm{hh}}(p),\qquad\mathrm{lgS}(p):=\log\det\mathrm{Dsc}_{\mathbf{m}}(p).
\end{gather*}
Since $D^{(p;\mathrm{ma})}$ is an entire, $2\pi\mathbb{Z}^3$-periodic function of $p$, so is $D^{\mathrm{hh}}(p)$.
By Proposition~1.2\,(S3), $\mathrm{Dsc}_{\mathbf{m}}(p)$ is holomorphic and $2\pi\mathbb{Z}^3$-periodic on the strip $\{|\operatorname{Im}p_\mu|\le\theta_{0}\}$, and $|\mathrm{Dsc}_{\mathbf{m}}(p)|_2\le K_{\mathrm{Str}}$ there.
It is holomorphic even in a neighbourhood of this closed strip:
by Proposition~1.2\,(S1), $D^{\mathrm{hh}}(p)$ is invertible on the closed strip, and the set of the points at which an entire matrix function is invertible is open.
Hence $\mathrm{lgH}$ and $\mathrm{lgS}$ are $2\pi\mathbb{Z}^3$-periodic and of class $C^\infty$ on $\mathbb{R}^3$, and by~(H3), for real $p$,
\begin{gather*}
\partial_\mu\mathrm{lgH}(p)=\mathrm{dlH}_\mu(p):=\mathrm{Tr}\left(D^{\mathrm{hh}}(p)^{-1}\,\partial_\mu D^{\mathrm{hh}}(p)\right),\\
\partial_\mu\mathrm{lgS}(p)=\mathrm{dlS}_\mu(p):=\mathrm{tr}\left(\mathrm{Dsc}_{\mathbf{m}}(p)^{-1}\,\partial_\mu\mathrm{Dsc}_{\mathbf{m}}(p)\right).
\end{gather*}
The right-hand sides are defined, holomorphic and $2\pi\mathbb{Z}^3$-periodic at all points where $D^{\mathrm{hh}}(p)$, respectively $\mathrm{Dsc}_{\mathbf{m}}(p)$, is holomorphic and invertible.
Let $\mathrm{lgHh}(z)$ and $\mathrm{lgSh}(z)$ denote the Fourier coefficients of $\mathrm{lgH}$ and $\mathrm{lgS}$ in the sense of~0.3.

\emph{(c2) The heavy factor: $|\mathrm{lgHh}(z)|\le C_{\mathrm{lH}}\,|z|^{-1}\exp\left(-\theta_{0}|z|\right)$ for $z\ne0$.}
By (I.2.6),
\[
\left(\partial_\mu D^{(p;\mathrm{ma})}u\right)(g)=-i\,(ge_1)_\mu\,\mathrm{e}^{ip\cdot(ge_1)}\,M_{3}(g)\,u(gs_3).
\]
Here $u\mapsto u(\cdot\,s_3)$ is unitary on ${V_{\mathrm{v}}}$, since $g\mapsto gs_3$ is a bijection of $O_{\mathrm{cub}}$;
moreover $|(ge_1)_\mu|\le1$, $|M_{3}(g)|_2=3(1+\rho)$ by (I.m6), and $|\mathrm{e}^{ip\cdot(ge_1)}|\le\exp(\theta_{0})$ whenever $|\operatorname{Im}p_\nu|\le\theta_{0}$, because $ge_1$ is one of the vectors $\pm e_\nu$.
On the strip one therefore has $\|\partial_\mu D^{(p;\mathrm{ma})}\|\le3(1+\rho)\exp(\theta_{0})=3(1+\rho)+\hat{g}/2$, by the definition of $\theta_{0}$,
and the same bound holds for the compression $\partial_\mu D^{\mathrm{hh}}(p)=\mathrm{Phv}\left(\partial_\mu D^{(p;\mathrm{ma})}\right)\mathrm{Phv}$.
Together with $\|D^{\mathrm{hh}}(p)^{-1}\|\le2/\hat{g}$ (Proposition~1.2\,(S1)) and (H2) on $\ker J^{*}$, this gives
\[
|\mathrm{dlH}_\mu(p)|\le188\cdot\frac{2}{\hat{g}}\left(3(1+\rho)+\frac{\hat{g}}{2}\right)=C_{\mathrm{lH}}\qquad\text{for }|\operatorname{Im}p_\nu|\le\theta_{0}\ (\nu=1,2,3).
\]
Let $z\ne0$, and choose $\mu$ with $|z_\mu|=|z|$.
Apply (F5) with $\mathrm{kd}=1$, and then (F6) with $\mathrm{fhol}=\mathrm{dlH}_\mu$, $\nu=\mu$ and $\mathrm{wd}=\theta_{0}$;
this is legitimate, since $\mathrm{dlH}_\mu$ is holomorphic in a neighbourhood of the closed strip, as in~(c1).
The result is
\[
\mathrm{lgHh}(z)=(iz_\mu)^{-1}\exp\left(-\theta_{0}|z_\mu|\right)\int_{[-\pi,\pi]^3}\frac{\mathrm{d}^3p}{(2\pi)^3}\,\mathrm{dlH}_\mu\left(p-i\,\mathrm{sgz}\,\theta_{0}e_\mu\right)\mathrm{e}^{-ip\cdot z},
\]
and the modulus of the integral is at most $C_{\mathrm{lH}}$.

\emph{(c3) The light factor: $|\mathrm{lgSh}(z)|\le\mathrm{cF13}\cdot C_{\mathrm{lS}}\,|z|^{-3}\exp\left(-\kappa_{(0)}\cdot\mathbf{m}\cdot|z|\right)$ for $z\ne0$.}
Let $p=p'+ip''$ with $p'\in[-\pi,\pi]^3$ and $|p''_\nu|\le\kappa_{(0)}\cdot\mathbf{m}$ ($\nu=1,2,3$), and let $\mathrm{rp}(p')=\left(|p'|_2^2+\mathbf{m}^2\right)^{\frac12}$.
By Proposition~1.2\,(S7), $\mathrm{Dsc}_{\mathbf{m}}(p)$ is invertible, and
\[
\left|\partial^{\mathrm{ak}'}\mathrm{Dsc}_{\mathbf{m}}(p)^{-1}\right|_2\le A'_{\mathrm{ak}'}\,\mathrm{rp}(p')^{-1-|\mathrm{ak}'|}.
\]
As $\kappa_{(0)}\cdot\mathbf{m}\le\theta_{0}/2$, the closed polydisc of radius $\theta_{0}/2$ about $p$ is contained in the strip of width $\theta_{0}$.
Cauchy's estimates on this polydisc \cite[Section~2.2]{Hormander1990}, applied to the functions $(v',\mathrm{Dsc}_{\mathbf{m}}(\cdot)\,v)$ with $v,v'\in\mathbb{C}^4$, $|v|_2=|v'|_2=1$, together with Proposition~1.2\,(S3), give for every multi-index $\mathrm{aj}$ with $1\le|\mathrm{aj}|\le4$
\[
\left|\partial^{\mathrm{aj}}\mathrm{Dsc}_{\mathbf{m}}(p)\right|_2\le\mathrm{aj}!\,K_{\mathrm{Str}}\left(\frac{2}{\theta_{0}}\right)^{|\mathrm{aj}|}\le24\,K_{\mathrm{Str}}\left(\frac{2}{\theta_{0}}\right)^{|\mathrm{aj}|}.
\]
The Leibniz rule and $|\mathrm{tr}\,m|\le4\,|m|_2$ then give, for $|\mathrm{ak}|\le3$,
\[
\left|\partial^{\mathrm{ak}}\mathrm{dlS}_\mu(p)\right|\le4\sum_{\mathrm{ak}'\le\mathrm{ak}}\binom{\mathrm{ak}}{\mathrm{ak}'}A'_{\mathrm{ak}'}\,\mathrm{rp}(p')^{-1-|\mathrm{ak}'|}\cdot24\,K_{\mathrm{Str}}\left(\frac{2}{\theta_{0}}\right)^{|\mathrm{ak}|-|\mathrm{ak}'|+1}.
\]
From $\mathrm{rp}(p')^2\le3\pi^2+1$ one gets $\mathrm{rp}(p')^{-1-|\mathrm{ak}'|}\le\mathrm{rp}(p')^{-1-|\mathrm{ak}|}\,(3\pi^2+1)^{\frac{|\mathrm{ak}|-|\mathrm{ak}'|}{2}}$, and $\mathrm{rp}(p')\ge|p'|_2$.
It follows that
\begin{align}
\left|\partial^{\mathrm{ak}}\mathrm{dlS}_\mu(p'+ip'')\right|&\le C_{\mathrm{lS}}\,\mathrm{rp}(p')^{-1-|\mathrm{ak}|}\le C_{\mathrm{lS}}\,|p'|_2^{-1-|\mathrm{ak}|}\notag\\
&\qquad\left(|\mathrm{ak}|\le3,\ p'\in[-\pi,\pi]^3\setminus\{0\},\ |p''_\nu|\le\kappa_{(0)}\cdot\mathbf{m}\right).\tag{2.1}
\end{align}
Let $z\ne0$, and choose $\mu$ with $|z_\mu|=|z|$.
The function $\mathrm{dlS}_\mu$ is holomorphic in a neighbourhood of the closed strip $\{|\operatorname{Im}p_\nu|\le\kappa_{(0)}\cdot\mathbf{m}\}$,
because $\mathrm{Dsc}_{\mathbf{m}}$ is holomorphic near this strip and, by Proposition~1.2\,(S7) and periodicity, invertible on it.
By (F5) with $\mathrm{kd}=1$ and (F6) with $\mathrm{fhol}=\mathrm{dlS}_\mu$, $\nu=\mu$ and $\mathrm{wd}=\kappa_{(0)}\cdot\mathbf{m}$,
\[
\mathrm{lgSh}(z)=(iz_\mu)^{-1}\exp\left(-\kappa_{(0)}\cdot\mathbf{m}\cdot|z_\mu|\right)\int_{[-\pi,\pi]^3}\frac{\mathrm{d}^3p'}{(2\pi)^3}\,\mathrm{dlS}_\mu\left(p'-i\,\mathrm{sgz}\,\kappa_{(0)}\cdot\mathbf{m}\,e_\mu\right)\mathrm{e}^{-ip'\cdot z}.
\]
Recall \cite[Lemma~12.4]{QEDI}:
let $\mathrm{kk}<3$ be real, let $N>3-\mathrm{kk}$ be an integer, let $A\ge0$, and let $a$ be a $2\pi\mathbb{Z}^3$-periodic matrix valued function of class $C^N$ on $\mathbb{R}^3\setminus2\pi\mathbb{Z}^3$ such that $|\partial^{\beta}a(p)|\le A\,|p|_2^{-\mathrm{kk}-|\beta|}$ for $p\in[-\pi,\pi]^3\setminus\{0\}$ and $|\beta|\le N$;
then $a$ is integrable over $[-\pi,\pi]^3$, and
\[
\left|(2\pi)^{-3}\int_{[-\pi,\pi]^3}\mathrm{d}p\,\mathrm{e}^{ip\cdot x}\,a(p)\right|\le c(\mathrm{kk},N)\,A\,(1+|x|)^{-(3-\mathrm{kk})}\qquad(x\in\mathbb{Z}^3).
\]
In this lemma the modulus of the momentum is the Euclidean norm (Section~I.1.1): its proof integrates $|p|^{-\mathrm{kk}}$ over balls and uses that $\chi_0(|p|/t)$ is a smooth function of $p$.
The function $p'\mapsto\mathrm{dlS}_\mu\left(p'-i\,\mathrm{sgz}\,\kappa_{(0)}\cdot\mathbf{m}\,e_\mu\right)$ of the real variable $p'$ is $2\pi\mathbb{Z}^3$-periodic and of class $C^\infty$ on $\mathbb{R}^3$,
and by (2.1) it satisfies the hypothesis of the lemma with $(\mathrm{kk},N)=(1,3)$ and the constant $C_{\mathrm{lS}}$.
At $x=-z$ the lemma bounds the modulus of the integral above by $\mathrm{cF13}\cdot C_{\mathrm{lS}}\,(1+|z|)^{-2}$.
Hence
\[
|\mathrm{lgSh}(z)|\le\mathrm{cF13}\cdot C_{\mathrm{lS}}\,|z|^{-1}(1+|z|)^{-2}\exp\left(-\kappa_{(0)}\cdot\mathbf{m}\cdot|z|\right).
\]

\emph{(c4) Summation.}
The function $\mathrm{fper}:=\log\mathfrak{d}_{\mathrm{ma}}$ is continuous and periodic, $\mathrm{fperh}=\mathrm{lgHh}+\mathrm{lgSh}$, and $\sum_z|\mathrm{fperh}(z)|<\infty$ by (c2), (c3) and (F2).
By (F1), clause~(d) applies to $\mathrm{fper}$.
Moreover $\mathrm{fperh}(0)=p^{0}_{\mathrm{ch}}(\mathbf{m})$, and $\sum_{p\in P_{n,\theta}}\mathrm{fper}(p)=\log\det\left(D_{H_\theta}+\mathbf{m}\right)$ by~(b).
So (8.2) gives
\[
\left|\log\det\left(D_{H_\theta}+\mathbf{m}\right)-p^{0}_{\mathrm{ch}}(\mathbf{m})\cdot n^3\right|\le n^3\sum_{w\ne0}|\mathrm{lgSh}(nw)|+n^3\sum_{w\ne0}|\mathrm{lgHh}(nw)|.
\]
By (c3), (F2) and (F3), with $t=|w|$,
\begin{align*}
n^3\sum_{w\ne0}|\mathrm{lgSh}(nw)|
&\le n^3\sum_{t\ge1}26\,t^2\,\mathrm{cF13}\cdot C_{\mathrm{lS}}\,(nt)^{-3}\exp\left(-\kappa_{(0)}\cdot\mathbf{m}\cdot nt\right)\\
&=26\,\mathrm{cF13}\cdot C_{\mathrm{lS}}\sum_{t\ge1}\frac{\exp\left(-\kappa_{(0)}\cdot\mathbf{m}\cdot nt\right)}{t}
=26\,\mathrm{cF13}\cdot C_{\mathrm{lS}}\,\log\frac{1}{1-\exp\left(-\kappa_{(0)}\cdot\mathbf{m}\cdot n\right)}.
\end{align*}
By (c2) and (F2),
\[
n^3\sum_{w\ne0}|\mathrm{lgHh}(nw)|\le\sum_{t\ge1}26\,t^2\,n^3\,C_{\mathrm{lH}}\,(nt)^{-1}\exp\left(-\theta_{0}nt\right)=26\,C_{\mathrm{lH}}\sum_{t\ge1}t^{-1}(nt)^2\exp\left(-\theta_{0}nt\right).
\]
Because $\mathbf{m}\le1$, one has $\kappa_{(0)}\cdot\mathbf{m}\le\theta_{0}/2$, so that $\exp(-\theta_{0}nt)\le\exp\left(-\frac{\theta_{0}nt}{2}\right)\exp\left(-\kappa_{(0)}\cdot\mathbf{m}\cdot nt\right)$;
furthermore $\mathrm{xr}^2\exp\left(-\frac{\theta_{0}\,\mathrm{xr}}{2}\right)\le\left(\frac{4}{\theta_{0}}\right)^2\exp(-2)$ for real $\mathrm{xr}\ge0$, the maximum being attained at $\mathrm{xr}=4/\theta_{0}$.
With (F3) this yields
\[
n^3\sum_{w\ne0}|\mathrm{lgHh}(nw)|\le26\,C_{\mathrm{lH}}\,\frac{16}{(\mathrm{e}\,\theta_{0})^2}\,\log\frac{1}{1-\exp\left(-\kappa_{(0)}\cdot\mathbf{m}\cdot n\right)}.
\]
The sum of the two bounds is $C_{6}\log\frac{1}{1-\exp(-\kappa_{(0)}\cdot\mathbf{m}\cdot n)}=\varepsilon_5(\mathbf{m}\cdot n)$ (Section~I.8.2).

\textbf{Constants, uniformity.}
One has $\kappa_{(0)}>0$, because $\theta_{0}>0$, $c_{\mathrm{el}}>0$ and $K_{\mathrm{Str}}<\infty$ (Proposition~1.2).
By the formulas of Proposition~30.1, $C_{6}$ is formed from $\rho$, $\hat{g}$, $\theta_{0}$, $K_{\mathrm{Str}}$, the numbers $A'_{\mathrm{ak}}$ with $|\mathrm{ak}|\le3$, and the absolute number $\mathrm{cF13}$;
it is finite, since $\hat{g}>0$ and $\theta_{0}>0$, and it depends on $\rho$ only.
None of the constants depends on $n$, $\theta$, $\mathbf{m}$ or $q$:
the charge $q$ drops out in~(b);
$\theta$ enters only through the factors of modulus one in~(8.2);
and the bounds (c2), (c3) hold for all $\mathbf{m}\in(0,1]$, because Proposition~1.2\,(S1), (S3), (S7) are uniform in the mass.
The right-hand side of (8.1) depends on $n$ and $\mathbf{m}$ only through the product $\mathbf{m}\cdot n$.
\end{proof}

\subsection*{3. Proposition 30.2}

\subsubsection*{Proof of Proposition 30.2}

\begin{proof}
\textbf{(a)}
Both determinants are positive by Proposition~30.1\,(b).
Applying Proposition~30.1\,(c) at $\theta$ and at $\theta=0$, where $H_0=0$ and $D_{H_0}=D_0$, one obtains
\begin{multline*}
\left|\log\det\left(D_{H_\theta}+\mathbf{m}\right)-\log\det\left(D_0+\mathbf{m}\right)\right|\\
\le\left|\log\det\left(D_{H_\theta}+\mathbf{m}\right)-p^{0}_{\mathrm{ch}}(\mathbf{m})\cdot n^3\right|+\left|p^{0}_{\mathrm{ch}}(\mathbf{m})\cdot n^3-\log\det\left(D_0+\mathbf{m}\right)\right|
\le2\,\varepsilon_5(\mathbf{m}\cdot n).
\end{multline*}

\textbf{(b)}
The number $\mathbf{m}:=\bar{m}\cdot\eta$ lies in $(0,1]$, since $0<\bar{m}\le1$ and $0<\eta=L^{-N}\le1$, and $\mathbf{m}\cdot n=\bar{m}\cdot\eta\cdot R/\eta=\bar{m}\cdot R$.
The charge $q$ with which $H_\theta$ is formed does not occur on the right-hand side of Proposition~30.1\,(b); hence the determinant does not depend on it.
By Proposition~30.1\,(a), (b), the determinant $\det\left(D_{H_\theta}+\mathbf{m}\right)=\prod_{p'\in P_{n,0}}\mathfrak{d}_{\mathrm{ma}}(p'+\theta/n)$ is a continuous function of $\theta$,
and by Proposition~30.1\,(b), (c) one has, for every $\theta$,
\[
\exp\left(p^{0}_{\mathrm{ch}}(\mathbf{m})\cdot n^3-\varepsilon_5(\bar{m}\cdot R)\right)\le\det\left(D_{H_\theta}+\mathbf{m}\right)\le\exp\left(p^{0}_{\mathrm{ch}}(\mathbf{m})\cdot n^3+\varepsilon_5(\bar{m}\cdot R)\right).
\]
The measure $\mathrm{d}\theta/(2\pi)^3$ is a probability measure on $[-\pi,\pi]^3$, and averaging a continuous function whose values lie in an interval gives a number in that interval.
Therefore $\mathrm{zr}^{\mathrm{free}}_n$ lies in the same interval;
in particular $\mathrm{zr}^{\mathrm{free}}_n>0$, and taking logarithms gives the assertion.
\end{proof}

\begin{named}{Remark}
The bounds of Proposition~30.2 depend on $(N,R,\bar{m})$ only through $\bar{m}\cdot R$, and they do not contain the charge.
\end{named}

\subsection*{4. Proposition 30.3}

\subsubsection*{Proof of Proposition 30.3}

\paragraph{4.0 Setting}
The side $\ell=(\ell\,\mathfrak{N})\,\mathrm{et}$ is a positive integer multiple of $\mathrm{et}$.
Hence $\mathscr{H}_{\theta}=\mathrm{HH}(\mathrm{et},\ell;\theta)$, with its scalar product, is the space of Section~1 at $(\mathrm{me},R)=(\mathrm{et},\ell)$, and $D$ is the operator $D$ of Section~1 at the mesh $\mathrm{et}$;
by (Lz1), $D$ maps $\mathscr{H}_{\theta}$ into itself.
The norms $|\cdot|$ of the statement are written $\|\cdot\|$.

(B1) (Blocks.)
The sites of $\mathrm{et}\,\mathbb{Z}^3$ are the points $c=j/\mathfrak{N}$ with $j\in\mathbb{Z}^3$.
For $y\in\mathbb{Z}^3$ one has $|c_\mu-y_\mu|<\frac12$ if and only if $|j_\mu-\mathfrak{N}y_\mu|<\mathfrak{N}/2$, and, these being integers and $\mathfrak{N}$ being odd, this holds if and only if $|j_\mu-\mathfrak{N}y_\mu|\le(\mathfrak{N}-1)/2$.
Every integer has exactly one representation $\mathfrak{N}y_\mu+j'_\mu$ with $y_\mu\in\mathbb{Z}$ and $|j'_\mu|\le(\mathfrak{N}-1)/2$.
Consequently every $c$ lies in exactly one block $\Delta_y$, and $\Delta_y$ has $\mathfrak{N}^3=\mathrm{et}^{-3}$ sites;
$\Delta_0=\{j/\mathfrak{N}:j\in\mathrm{Mf}\}=-\Delta_0$;
$\Delta_{y+v}=\Delta_y+v$ and $[c+v]=[c]+v$ for $v\in\mathbb{Z}^3$;
$r\Delta_y=\Delta_{ry}$ and $[rc]=r[c]$ for $r\in O_{\mathrm{cub}}$, since the sup norm is invariant under signed permutations.
The block $\Delta_0$ is a period cell of $\mathbb{Z}^3$ in $\mathrm{et}\,\mathbb{Z}^3$, and if $\mathrm{Yc}$ is a period cell of $\ell\mathbb{Z}^3$ in $\mathbb{Z}^3$, then the union of the blocks $\Delta_y$, $y\in\mathrm{Yc}$, is a period cell of $\ell\mathbb{Z}^3$ in $\mathrm{et}\,\mathbb{Z}^3$.

(B2)
For every function $f:\mathrm{et}\,\mathbb{Z}^3\times O_{\mathrm{cub}}\to\mathbb{C}^4$ let
\[
(\mathrm{Pbl}\,f)(c,g):=\mathrm{et}^3\,\frac{1}{48}\sum_{c'\in\Delta_{[c]}}\sum_{g'\in O_{\mathrm{cub}}}f(c',g'),
\]
the mean of $f$ over the $\mathfrak{N}^3\cdot48$ sites of $\Delta_{[c]}\times O_{\mathrm{cub}}$.
The operator $\mathrm{Pbl}$ does not act on the spinor index, and by (B1) it commutes with the translations by $\mathbb{Z}^3$.

\paragraph{4.1 Clause (a)}
\begin{proof}[Proof of Proposition~30.3\,(a)]
\emph{The operators $Q$ and $Q^{*}$.}
By (B1),
\[
(Qf)(y+\ell e_\mu)=\mathrm{et}^3\,\frac{1}{48}\sum_{c\in\Delta_y}\sum_{g}f(c+\ell e_\mu,g)=-\mathrm{e}^{i\theta_\mu}(Qf)(y)
\]
and
\[
(Q^{*}v)(c+\ell e_\mu,g)=v([c]+\ell e_\mu)=-\mathrm{e}^{i\theta_\mu}(Q^{*}v)(c,g).
\]
So $Q$ maps $\mathscr{H}_{\theta}$ into $\mathscr{H}^{\mathrm{u}}_{\theta}$, and $Q^{*}$ maps $\mathscr{H}^{\mathrm{u}}_{\theta}$ into $\mathscr{H}_{\theta}$.
With the period cells of (B1),
\[
(Q^{*}v,f)=\mathrm{et}^3\,\frac{1}{48}\sum_{y\in\mathrm{Yc}}\sum_{c\in\Delta_y}\sum_{g\in O_{\mathrm{cub}}}\overline{v(y)}\cdot f(c,g)=\sum_{y\in\mathrm{Yc}}\overline{v(y)}\cdot(Qf)(y)=(v,Qf),
\]
so that $Q^{*}$ is the adjoint of $Q$.
Moreover
\[
(QQ^{*}v)(y)=\mathrm{et}^3\,\frac{1}{48}\sum_{c\in\Delta_y}\sum_{g}v(y)=\mathrm{et}^3\cdot\mathfrak{N}^3\cdot\frac{48}{48}\,v(y)=v(y).
\]
Hence $(Q^{*}Q)^2=Q^{*}(QQ^{*})Q=Q^{*}Q=(Q^{*}Q)^{*}$; thus $Q^{*}Q$ is an orthogonal projection, and its range is $\operatorname{ran}Q^{*}$.
Every element of $\operatorname{ran}Q^{*}$ is constant on each set $\Delta_y\times O_{\mathrm{cub}}$;
conversely, an $f\in\mathscr{H}_{\theta}$ with this property equals $Q^{*}v$, where $v(y)$ is the value of $f$ on $\Delta_y\times O_{\mathrm{cub}}$.
By the definitions, $Q^{*}Q=\mathrm{Pbl}$ on $\mathscr{H}_{\theta}$.
Finally $\|Q^{*}v\|^2=(v,QQ^{*}v)=\|v\|^2$, whence
\begin{equation}
\|Q\|=\|Q^{*}\|=1.\tag{4.1}
\end{equation}

\emph{Real parts, blocks, determinants.}
In the decomposition (Lz5) one has $\mathrm{Pbl}=\operatorname{diag}(\mathrm{Pbl},\mathrm{Pbl})$, and $\mathrm{Pbl}$ is Hermitian and $\ge0$.
Therefore (Lz6) applies on $\mathscr{H}_{\theta}$ with $t=\mu$ and $\mathrm{Pn}=0$, and also with $t=\mu$ and $\mathrm{Pn}=\mathsf{b}\,Q^{*}Q$.
It gives
\[
\operatorname{Re}(D+\mu)=\frac{\rho}{2\,\mathrm{et}}\,\nabla^{*}\nabla+\mu\ge\mu,\qquad\operatorname{Re}T_{\theta}=\frac{\rho}{2\,\mathrm{et}}\,\nabla^{*}\nabla+\mu+\mathsf{b}\,Q^{*}Q\ge\mu;
\]
$D+\mu$ and $T_{\theta}$ are bijective, and their inverses have norm at most $1/\mu$;
in the two eigenspaces of $\Sigma_3$, both identified with $\mathrm{HH2}(\mathrm{et},\ell;\theta)$,
\[
D+\mu=(D+\mu)^{+}\oplus\left((D+\mu)^{+}\right)^{*},\qquad T_{\theta}=T^{+}\oplus\left(T^{+}\right)^{*};
\]
and $\det(D+\mu)=\left|\det(D+\mu)^{+}\right|^2>0$, $\det T_{\theta}=\left|\det T^{+}\right|^2>0$.
For later use, (Lz4) gives
\begin{equation}
\operatorname{Re}(f,T_{\theta}f)\ge\frac{\hat{g}}{\mathrm{et}}\,\|\mathrm{Phv}\,f\|^2+\mu\,\|f\|^2+\mathsf{b}\,\|\mathrm{Pbl}\,f\|^2\qquad(f\in\mathscr{H}_{\theta}).\tag{4.2}
\end{equation}
\end{proof}

\paragraph{4.2 Clause (b)}
\begin{proof}[Proof of Proposition~30.3\,(b)]
By 4.1 and (H1), $G:=(D+\mu)^{-1}$ exists, $\|G\|\le1/\mu$ and $\operatorname{Re}G\ge0$.
The operator $\mathrm{Sfr}:=QGQ^{*}$ on $\mathscr{H}^{\mathrm{u}}_{\theta}$ satisfies $\operatorname{Re}\mathrm{Sfr}=Q(\operatorname{Re}G)Q^{*}\ge0$, and $\|\mathrm{Sfr}\|\le1/\mu$ by (4.1).
Hence $\operatorname{Re}(1/\mathsf{b}+\mathrm{Sfr})\ge1/\mathsf{b}>0$, and by (H1) the operator $\mathrm{Zfr}:=(1/\mathsf{b}+\mathrm{Sfr})^{-1}$ exists.
Multiplying out,
\begin{align*}
\left(G^{-1}+\mathsf{b}\,Q^{*}Q\right)\left(G-GQ^{*}\,\mathrm{Zfr}\,QG\right)
&=1-Q^{*}\,\mathrm{Zfr}\,QG+\mathsf{b}\,Q^{*}QG-\mathsf{b}\,Q^{*}\,\mathrm{Sfr}\,\mathrm{Zfr}\,QG\\
&=1+Q^{*}\left[\mathsf{b}-(1+\mathsf{b}\,\mathrm{Sfr})\,\mathrm{Zfr}\right]QG=1,
\end{align*}
because $(1+\mathsf{b}\,\mathrm{Sfr})\,\mathrm{Zfr}=\mathsf{b}\,(1/\mathsf{b}+\mathrm{Sfr})\,\mathrm{Zfr}=\mathsf{b}$.
As $G^{-1}+\mathsf{b}\,Q^{*}Q=T_{\theta}$, it follows that $T_{\theta}^{-1}=G-GQ^{*}\,\mathrm{Zfr}\,QG$.
The operators $\mathrm{Sfr}$ and $\mathrm{Zfr}$ commute, and therefore
\begin{gather*}
QT_{\theta}^{-1}Q^{*}=\mathrm{Sfr}-\mathrm{Sfr}\,\mathrm{Zfr}\,\mathrm{Sfr}=\mathrm{Sfr}\,\mathrm{Zfr}\left(\mathrm{Zfr}^{-1}-\mathrm{Sfr}\right)=\mathsf{b}^{-1}\,\mathrm{Sfr}\,\mathrm{Zfr},\\
\mathrm{Dfl}_{\theta}=\mathsf{b}-\mathsf{b}\,\mathrm{Sfr}\,\mathrm{Zfr}=\mathsf{b}\left(\mathrm{Zfr}^{-1}-\mathrm{Sfr}\right)\mathrm{Zfr}=\mathrm{Zfr}.
\end{gather*}
Consequently $\mathrm{Dfl}_{\theta}$ is invertible, and
\[
\mathrm{Dfl}_{\theta}^{-1}=\frac{1}{\mathsf{b}}+\mathrm{Sfr}=\frac{1}{\mathsf{b}}+Q(D+\mu)^{-1}Q^{*},\qquad\operatorname{Re}\mathrm{Dfl}_{\theta}^{-1}=\frac{1}{\mathsf{b}}+\operatorname{Re}\mathrm{Sfr}\ge\frac{1}{\mathsf{b}},\qquad\left\|\mathrm{Dfl}_{\theta}^{-1}\right\|\le\frac{1}{\mathsf{b}}+\frac{1}{\mu}.
\]
In addition, (H1) applied to $O=\mathrm{Dfl}_{\theta}^{-1}$ and $t=1/\mathsf{b}$ gives $\|\mathrm{Dfl}_{\theta}\|\le\mathsf{b}$.
\end{proof}

\paragraph{4.3 Clause (c)}
\begin{proof}[Proof of Proposition~30.3\,(c)]
The space $\mathscr{H}^{\mathrm{u}}_{\theta}$ has dimension $4\ell^3$.
By (H4) with $O_1=-\mathsf{b}\,Q$ and $O_2=T_{\theta}^{-1}Q^{*}$,
\begin{align*}
\det\mathrm{Dfl}_{\theta}&=\mathsf{b}^{4\ell^3}\det\left(1-\mathsf{b}\,QT_{\theta}^{-1}Q^{*}\right)=\mathsf{b}^{4\ell^3}\det\left(1-\mathsf{b}\,T_{\theta}^{-1}Q^{*}Q\right)\\
&=\mathsf{b}^{4\ell^3}\det\left(T_{\theta}^{-1}\left(T_{\theta}-\mathsf{b}\,Q^{*}Q\right)\right)=\mathsf{b}^{4\ell^3}\,\frac{\det(D+\mu)}{\det T_{\theta}}.
\end{align*}
By 4.1, both determinants on the right-hand side are positive.
Since this is valid for every $\theta$, and in particular for $\theta=0$, it follows that $\mathrm{Rtw}_{\theta}>0$.
\end{proof}

\paragraph{4.4 Clause (d)}
\begin{proof}[Proof of Proposition~30.3\,(d)]
Let $n:=\ell\,\mathfrak{N}$ and $\mathbf{m}:=\mathrm{et}\cdot\mu\in(0,1/3]$; then $\mathbf{m}\cdot n=\mu\cdot\ell$.
By (Lz9) with $\mathrm{me}=\mathrm{et}$ and $R=\ell$, the map $\mathrm{Sca}$ is a linear bijection of $\mathscr{H}_{\theta}$ onto $\mathrm{HH}(1,n;\theta)$, and
\[
\mathrm{Sca}\,(D+\mu)\,\mathrm{Sca}^{-1}=\mathrm{et}^{-1}\left(D'+\mathrm{et}\cdot\mu\right),
\]
where $D'$ is the operator $D$ at the mesh $1$.
By (Lz8) with $(\mathrm{me},R,\mathrm{ec})=(1,n,q)$, for which $\mathrm{cst}=H_{\theta}$, the operator $D'$ on $\mathrm{HH}(1,n;\theta)$ equals $\mathrm{Utw}\,D_{H_\theta}\,\mathrm{Utw}^{-1}$, with $D_{H_\theta}$ acting on $\mathrm{HH}(1,n;0)=\mathfrak{S}^{\mathrm{ch}}(1;n)$.
Similar operators have the same determinant, and the spaces involved have dimension $192\,n^3$.
Therefore
\[
\det\left(D+\mu\text{ on }\mathscr{H}_{\theta}\right)=\mathrm{et}^{-192n^3}\det\left(D_{H_\theta}+\mathrm{et}\cdot\mu\right),
\]
and at $\theta=0$
\[
\det\left(D+\mu\text{ on }\mathscr{H}_{\theta=0}\right)=\mathrm{et}^{-192n^3}\det\left(D_0+\mathrm{et}\cdot\mu\right).
\]
Insert these two identities into (c) at $\theta$ and at $\theta=0$, and divide.
The factors $\mathsf{b}^{4\ell^3}$ and $\mathrm{et}^{-192n^3}$ cancel, and
\[
\mathrm{Rtw}_{\theta}=\frac{\det\left(D_{H_\theta}+\mathrm{et}\cdot\mu\right)}{\det\left(D_0+\mathrm{et}\cdot\mu\right)}\cdot\frac{\det T_{\theta=0}}{\det T_{\theta}}.
\]
\end{proof}

\paragraph{4.5 Clause (f)}
\begin{proof}[Proof of Proposition~30.3\,(f)]
\emph{The maps $\mathrm{Uth}$, $\mathrm{Uth}^{\mathrm{u}}$ and the operator $D_{c_\theta}$.}
One has $\mathfrak{S}^{\mathrm{ch}}(\mathrm{et},\ell)=\mathrm{HH}(\mathrm{et},\ell;0)$.
Apply (Lz8) with $(\mathrm{me},R,\mathrm{ec})=(\mathrm{et},\ell,\mathrm{ec})$, for which $\mathrm{cst}=c_{\theta}$ and $\mathrm{Utw}=\mathrm{Uth}$:
the map $\mathrm{Uth}$ is a unitary map of $\mathfrak{S}^{\mathrm{ch}}(\mathrm{et},\ell)$ onto $\mathscr{H}_{\theta}$, and $D_{c_\theta}=\mathrm{Uth}^{-1}D\,\mathrm{Uth}$.
For $v\in\mathfrak{S}(1,\ell)$ one has
\[
(\mathrm{Uth}^{\mathrm{u}}v)(y+\ell e_\mu)=\mathrm{e}^{i\theta_\mu}\,\mathrm{e}^{i\frac{\theta\cdot y}{\ell}}\left(-v(y)\right);
\]
the factor $\mathrm{e}^{i\frac{\theta\cdot y}{\ell}}$ has modulus one, and the scalar products of $\mathfrak{S}(1,\ell)$ (Section~I.1.6: the sum over $y$) and of $\mathscr{H}^{\mathrm{u}}_{\theta}$ are given by the same formula.
Hence $\mathrm{Uth}^{\mathrm{u}}$ is a unitary map of $\mathfrak{S}(1,\ell)$ onto $\mathscr{H}^{\mathrm{u}}_{\theta}$.

\emph{The phases of the averages.}
Let $\mathfrak{N}=L^k$ with $k\ge1$, an odd integer $\ge3$ because $L$ is odd and $\ge5$; let $\ell\in\mathrm{Cl}$, $\mathrm{ec}=e_k$ and $\mathsf{b}=b_{k}$.
Then $\mathrm{et}=\eta_k$, and the blocks $\Delta_y$ are those of Section~I.1.1.
For $v\in\mathbb{R}^3$ write $c_{\theta}\cdot v:=\sum_\mu\frac{\theta_\mu}{\mathrm{ec}\cdot\ell}\,v_\mu$.
Let $c$, $c'$ be points of the covering lattice $\mathrm{et}\,\mathbb{Z}^3$, and let $\boldsymbol{\Gamma}$ be a rectilinear path $\boldsymbol{\Gamma}^{\pi}(c,c')$ (Section~I.1.1).
Along $\boldsymbol{\Gamma}$ the coordinate $\mu$ changes monotonically from $c_\mu$ to $c'_\mu$, through $|c'_\mu-c_\mu|/\mathrm{et}$ bonds, and each of these bonds contributes $\operatorname{sign}(c'_\mu-c_\mu)\,\theta_\mu/(\mathrm{ec}\cdot\ell)$ to $c_{\theta}(\boldsymbol{\Gamma})$.
Hence $c_{\theta}(\boldsymbol{\Gamma})=c_{\theta}\cdot(c'-c)/\mathrm{et}$ for each of the six orders $\pi$, and $(\tau c_{\theta})(c,c')=c_{\theta}\cdot(c'-c)/\mathrm{et}$.
For $c\in\Delta_y$ and the centres $c=y_0,y_1,\dots,y_k=y$ of Section~I.3.2, which are points of $\mathrm{et}\,\mathbb{Z}^3$, this gives
\begin{gather*}
\mathrm{et}\,(\tau_kc_{\theta})(y,c)=\mathrm{et}\sum_{j<k}(\tau c_{\theta})(y_{j+1},y_j)=\sum_{j<k}c_{\theta}\cdot(y_j-y_{j+1})=c_{\theta}\cdot(c-y),\\
\varphi_{c_\theta}(y,c)=\exp\left(i\,\mathrm{ec}\cdot\mathrm{et}\,(\tau_kc_{\theta})(y,c)\right)=\mathrm{e}^{i\frac{\theta\cdot(c-y)}{\ell}}.
\end{gather*}

\emph{The operators.}
By Section~I.3.3, for $f\in\mathfrak{S}^{\mathrm{ch}}(\mathrm{et},\ell)$ and $v\in\mathfrak{S}(1,\ell)$,
\begin{gather*}
\left(Q_k(c_{\theta})J^{*}f\right)(y)=\mathrm{et}^3\,\frac{1}{48}\sum_{c\in\Delta_y}\sum_{g\in O_{\mathrm{cub}}}\mathrm{e}^{i\frac{\theta\cdot(c-y)}{\ell}}f(c,g)=\mathrm{e}^{-i\frac{\theta\cdot y}{\ell}}\left(Q\,\mathrm{Uth}\,f\right)(y),\\
\left(JQ_k^{T}(-c_{\theta})v\right)(c,g)=\mathrm{e}^{-i\frac{\theta\cdot(c-[c])}{\ell}}v([c])=\mathrm{e}^{-i\frac{\theta\cdot c}{\ell}}\left(\mathrm{Uth}^{\mathrm{u}}v\right)([c])=\left(\mathrm{Uth}^{-1}Q^{*}\,\mathrm{Uth}^{\mathrm{u}}v\right)(c,g).
\end{gather*}
These are the first two identities of~(f).
They give $JP_k(c_{\theta})J^{*}=JQ_k^{T}(-c_{\theta})\,Q_k(c_{\theta})J^{*}=\mathrm{Uth}^{-1}Q^{*}Q\,\mathrm{Uth}$.
With $D_{c_\theta}=\mathrm{Uth}^{-1}D\,\mathrm{Uth}$ and Section~I.3.5 it follows that
\[
D_{c_\theta}+\mu+b_{k}\,JP_k(c_{\theta})J^{*}=\mathrm{Uth}^{-1}T_{\theta}\,\mathrm{Uth},\qquad\text{so that}\qquad S_k(c_{\theta})=\mathrm{Uth}^{-1}T_{\theta}^{-1}\,\mathrm{Uth}
\]
(the inverse exists by~(a)), and
\begin{align*}
D_k(c_{\theta})&=b_{k}-b_{k}^2\,Q_k(c_{\theta})J^{*}S_k(c_{\theta})\,JQ_k^{T}(-c_{\theta})\\
&=\left(\mathrm{Uth}^{\mathrm{u}}\right)^{-1}\left[\mathsf{b}-\mathsf{b}^2\,QT_{\theta}^{-1}Q^{*}\right]\mathrm{Uth}^{\mathrm{u}}=\left(\mathrm{Uth}^{\mathrm{u}}\right)^{-1}\mathrm{Dfl}_{\theta}\,\mathrm{Uth}^{\mathrm{u}}.
\end{align*}
Thus $\det D_k(c_{\theta})=\det\mathrm{Dfl}_{\theta}$.
At $\theta=0$ one has $c_{\theta}=0$, $\mathrm{Uth}=1$, $\mathrm{Uth}^{\mathrm{u}}=1$ and $D_k(0)=\mathrm{Dfl}_{\theta=0}$, and $\det\mathrm{Dfl}_{\theta=0}>0$ by~(c).
Hence $\det D_k(c_{\theta})/\det D_k(0)=\mathrm{Rtw}_{\theta}$.

\emph{In the flow.}
Let the data of (FL) be given (Sections~I.1.8 and~G4.0.1):
$\bar{m}\in(0,1]$; $e>0$ with $e\,h^{\frac12}\le e_1$; $R\in\mathrm{Tor}(e)$; an integer $N\ge N_{\min}(R,e)$.
These four data constitute the condition (S-4) of Section~I.1.8; the primitives that belong to (FL) do not enter.
The assertions (Sc1), (Sc3), (Sc4) are those of Section~G4.0.2, which hold under (S-4).
In \cite[Proposition~19.1(c), (e), (f)]{QEDI} these assertions are stated under the hypotheses $v_K\ge1$, $\left(1+\frac{\mathrm{lg}}{v_K}\right)^r\le L$ and $N\ge N_{\min}(R,e)$,
and, for~(f), under the condition $r_*>31L+16$ of \cite[Section~19.1.2]{QEDI}, which is the condition~(2.1) of Section~G4.0.3.
The hypothesis $N\ge N_{\min}(R,e)$ is one of the data of (FL).
The other three hold under (S-4) by (Sc6) of Section~G4.0.2:
$v_K\ge8$, $\left(1+\frac{\mathrm{lg}}{v_K}\right)^r\le\left(1+\frac{3\,\mathrm{lg}}{v_K}\right)^r<L$, and the condition~(2.1) of Section~G4.0.3.
Of these, $v_K\ge8$ is contained in Proposition~14.2\,(f), and the condition~(2.1) of Section~G4.0.3 in Proposition~14.2\,(c), both for $e\,h^{\frac12}\le e_1$;
the inequality $\left(1+\frac{3\,\mathrm{lg}}{v_K}\right)^r<L$ is the third inequality of the row~(e1) of Proposition~14.2\,(b) at $v=v_K$.
That this row holds at $v=v_K$ is seen as follows.
By the row $e_1$ of Proposition~14.2\,(a), $e_1$ is one half of the supremum of the numbers $e'\in\left(0,\min(e_0,1/3)\right]$ such that (e1)--(e8) hold for all $v\ge-\log e'$.
The data contain an $e>0$ with $e\,h^{\frac12}\le e_1$; hence $e_1>0$, the supremum equals $2e_1>e_1$, and there is such a number $e'$ with $e'>e_1$.
Consequently $v_K=-\log\left(e\,h^{\frac12}\right)\ge-\log e_1>-\log e'$.
(This is the first assertion of Proposition~14.2\,(c), taken at $v=v_K$: the datum $e>0$ with $e\,h^{\frac12}\le e_1$ gives $e_1>0$, so that the set in the definition of $e_1$ is not empty.)
The three assertions yield:
$k_{*}\ge1$, by (Sc1);
$\ell_*=\ell_{k_*}(R)$ lies in $\mathrm{Cl}$, by (Sc4) at $k=k_{*}$;
$0<\mu_*\le\mu_K\le1$, by (Sc3).
Further, $b_{-}\le b_{k_*}\le b_{+}$ (Section~I.3.4; $b_{k_*}$ is defined since $k_{*}\ge1$), and $e_{*}=e\,a_{*}^{\frac12}>0$.
Therefore $(\ell,\mathfrak{N},\mu,\mathsf{b},\mathrm{ec})=(\ell_*,L^{k_*},\mu_*,b_{k_*},e_{*})$ are data of clause~(f) with $k=k_{*}$;
the constant connection with the components $\theta_\mu/(e_{*}\cdot\ell_*)$ is $c^{(k_*)}_{\theta}$ (Section~I.6.6), and $D_k(\cdot)$ is $D_*(\cdot)=D_{k_*}(\cdot)$ (Section~I.6.2, the index $*$).
For these data, Section~I.6.1 ($a_{*}=L^{k_*-N}$, $\mu_*=\bar{m}\cdot a_{*}$, $\ell_*=R/a_{*}$, $\eta=L^{-N}$) gives
\[
\mathrm{et}\cdot\mu=L^{-k_*}\cdot\bar{m}\cdot L^{k_*-N}=\bar{m}\cdot\eta,\qquad\ell\,\mathfrak{N}=R\,L^{N-k_*}\,L^{k_*}=R/\eta=n;
\]
so the first ratio in~(d) is the ratio of Propositions~30.1 and~30.2 at $(n,\mathbf{m})=(R/\eta,\bar{m}\cdot\eta)$.
\end{proof}

\begin{named}{Remark}[Level 0]
Clause~(f) is a statement about the levels $k\ge1$.
At $k=0$ the objects $b_0$, $S_k(A)=S_\infty(A)$ and $M_0(A)$ are not defined, and $D_0(A)=D_A+\mu$ (Section~I.3.9);
no block average occurs, and the determinant of $D_{H_\theta}+\mathbf{m}$ is described by Proposition~30.1.
\end{named}

\paragraph{4.6 Bloch spaces, plane waves, the periodic gauge; clause (g1)}
The formulas by which clause (g) defines $p_{\mathrm{jf}}$, $u_1$, $\mathrm{uh}$ and $T(p)$ make sense for every $p\in\mathbb{C}^3$, and from now on they are taken for complex arguments;
$u_1$, $\mathrm{uh}$ and the matrix entries of $T(p)$ are then entire functions.
The index set $\{j\in\mathbb{Z}:|j|\le(\mathfrak{N}-1)/2\}$ is invariant under $j\mapsto-j$, so that $u_1(-z)=u_1(z)$.
Together with (B1) this gives, for every $\mathrm{pq}\in\mathbb{C}^3$,
\begin{equation}
\mathrm{et}^3\sum_{c\in\Delta_0}\mathrm{e}^{i\mathrm{pq}\cdot c}
=\mathrm{et}^3\sum_{c\in\Delta_0}\mathrm{e}^{-i\mathrm{pq}\cdot c}
=\prod_{\mu}\left(\mathfrak{N}^{-1}\sum_{|j|\le\frac{\mathfrak{N}-1}{2}}\mathrm{e}^{-\frac{i\mathrm{pq}_\mu j}{\mathfrak{N}}}\right)
=\mathrm{uh}(\mathrm{pq}).
\tag{4.3}
\end{equation}
For $\mathrm{jf},\mathrm{jf}'\in\mathrm{Mf}$ one has $|\mathrm{jf}_\mu-\mathrm{jf}'_\mu|\le\mathfrak{N}-1<\mathfrak{N}$, and therefore
\begin{equation}
\mathrm{et}^3\sum_{c\in\Delta_0}\mathrm{e}^{2\pi i(\mathrm{jf}-\mathrm{jf}')\cdot c}
=\prod_{\mu}\left(\mathfrak{N}^{-1}\sum_{|j|\le\frac{\mathfrak{N}-1}{2}}\mathrm{e}^{\frac{2\pi i(\mathrm{jf}_\mu-\mathrm{jf}'_\mu)j}{\mathfrak{N}}}\right)
=\mathbf{1}[\mathrm{jf}=\mathrm{jf}'].
\tag{4.4}
\end{equation}
The set $\mathrm{Mf}$ has $\#|\mathrm{Mf}|=\mathfrak{N}^3$ elements, and $\Delta_0$ is a period cell of $\mathbb{Z}^3$ in $\mathrm{et}\,\mathbb{Z}^3$ (B1).
Hence the $\mathfrak{N}^3$ functions $c\mapsto\mathrm{e}^{2\pi i\mathrm{jf}\cdot c}$, $\mathrm{jf}\in\mathrm{Mf}$, form an orthonormal basis of the space of $\mathbb{Z}^3$-periodic complex functions on $\mathrm{et}\,\mathbb{Z}^3$, equipped with the scalar product $\mathrm{et}^3\sum_{c\in\Delta_0}$.
Every such function $\mathrm{hs}$ is therefore expanded as
\begin{equation}
\mathrm{hs}(c)=\sum_{\mathrm{jf}\in\mathrm{Mf}}\left(\mathrm{et}^3\sum_{c'\in\Delta_0}\mathrm{e}^{-2\pi i\mathrm{jf}\cdot c'}\,\mathrm{hs}(c')\right)\mathrm{e}^{2\pi i\mathrm{jf}\cdot c}.
\tag{4.5}
\end{equation}
Let $p\in\mathbb{C}^3$.
By (B1) the function $c\mapsto\mathrm{e}^{-ip\cdot(c-[c])}$ is $\mathbb{Z}^3$-periodic, and on $\Delta_0$ it coincides with $\mathrm{e}^{-ip\cdot c}$.
Expanding it by (4.5) and computing the coefficients with (4.3), one obtains
\begin{equation}
\mathrm{e}^{-ip\cdot(c-[c])}=\sum_{\mathrm{jf}\in\mathrm{Mf}}\mathrm{uh}(p_{\mathrm{jf}})\,\mathrm{e}^{2\pi i\mathrm{jf}\cdot c}
\qquad(c\in\mathrm{et}\,\mathbb{Z}^3,\ p\in\mathbb{C}^3).
\tag{4.6}
\end{equation}

\emph{Definitions.}
For $p\in\mathbb{R}^3$ let
\[
\mathrm{Hb}(p):=\left\{f:\mathrm{et}\,\mathbb{Z}^3\times O_{\mathrm{cub}}\to\mathbb{C}^4\ :\ f(c+e_\mu,g)=\mathrm{e}^{ip_\mu}f(c,g)\ \ (\mu=1,2,3)\right\}=\mathrm{HH}\left(\mathrm{et},1;p-(\pi,\pi,\pi)\right),
\]
with the scalar product
\[
(f',f)=\mathrm{et}^3\cdot\frac{1}{48}\sum_{c\in\Delta_0}\sum_{g\in O_{\mathrm{cub}}}\overline{f'(c,g)}\cdot f(c,g).
\]
Its dimension is $192\,\mathfrak{N}^3$, and it depends on $p$ only modulo $2\pi\mathbb{Z}^3$.
For the side $1$ and the twist $\theta':=p-(\pi,\pi,\pi)$, the space $\mathrm{Hb}(p)$ is the space $\mathscr{H}_{\theta'}$ of Proposition~30.3; so Sections~4.0 and 4.1 apply to it.
Let $\mathrm{Tb}(p)$ be the restriction of $D+\mu+\mathsf{b}\,\mathrm{Pbl}$ to $\mathrm{Hb}(p)$.
It is the operator $T_{\theta'}$ at $\ell=1$, $\mathrm{Pbl}$ is an orthogonal projection of $\mathrm{Hb}(p)$, and
\begin{equation}
\operatorname{Re}\left(f,\mathrm{Tb}(p)f\right)\ge\frac{\hat{g}}{\mathrm{et}}\,\|\mathrm{Phv}\,f\|^2+\mu\,\|f\|^2+\mathsf{b}\,\|\mathrm{Pbl}\,f\|^2,
\qquad\det\mathrm{Tb}(p)>0
\qquad(p\ \text{real}).
\tag{4.7}
\end{equation}
Put $\mathrm{Hper}:=\mathrm{Hb}(0)$.
The space $\mathbb{C}^{\mathrm{Mf}}\otimes{V_{\mathrm{v}}}=\bigoplus_{\mathrm{jf}\in\mathrm{Mf}}{V_{\mathrm{v}}}$ carries the scalar product $\sum_{\mathrm{jf}}(u'_{\mathrm{jf}},u_{\mathrm{jf}})$.
For $p\in\mathbb{R}^3$ define
\begin{gather*}
\mathrm{Fw}_p:\mathbb{C}^{\mathrm{Mf}}\otimes{V_{\mathrm{v}}}\to\mathrm{Hb}(p),\qquad
\left(\mathrm{Fw}_p(u_{\mathrm{jf}})\right)(c,g):=\sum_{\mathrm{jf}\in\mathrm{Mf}}\mathrm{e}^{ip_{\mathrm{jf}}\cdot c}\,u_{\mathrm{jf}}(g),\\
U_p:\mathrm{Hper}\to\mathrm{Hb}(p),\qquad(U_p\mathrm{hp})(c,g):=\mathrm{e}^{ip\cdot c}\,\mathrm{hp}(c,g);
\end{gather*}
then $\mathrm{Fw}_p=U_p\mathrm{Fw}_0$.
For $p\in\mathbb{C}^3$ define operators on $\mathrm{Hper}$ by
\begin{align*}
(\mathrm{Dpg}(p)\mathrm{hp})(x)&:=-\mathrm{et}^{-1}\left\{\sum_{i\le2}M_i(x)\left[\mathrm{hp}(x^{(i)})-\mathrm{hp}(x)\right]+M_3(x)\left[\mathrm{e}^{i\,\mathrm{et}\cdot p\cdot(ge_1)}\,\mathrm{hp}(x^{(3)})-\mathrm{hp}(x)\right]\right\},\\
(\mathrm{Ppg}(p)\mathrm{hp})(c,g)&:=\mathrm{et}^3\cdot\frac{1}{48}\sum_{c'\in\Delta_{[c]}}\sum_{g'\in O_{\mathrm{cub}}}\mathrm{e}^{-ip\cdot(c-c')}\,\mathrm{hp}(c',g'),\\
\mathrm{Tpg}(p)&:=\mathrm{Dpg}(p)+\mu+\mathsf{b}\,\mathrm{Ppg}(p),
\end{align*}
where $x=(c,g)$, $x^{(i)}=(c,gs_i)$ for $i\le2$, and $x^{(3)}=(c+\mathrm{et}\cdot ge_1,gs_3)$.
Both $\mathrm{Dpg}(p)$ and $\mathrm{Ppg}(p)$ map $\mathrm{Hper}$ into itself, because $c-c'$ does not change when $c$ and $c'$ are translated by the same vector, and by (B1).
The formula $(\mathrm{Mper}_\nu\mathrm{hp})(c,g):=\mathrm{e}^{2\pi ic_\nu}\,\mathrm{hp}(c,g)$ defines a unitary map of $\mathrm{Hper}$, the function $\mathrm{e}^{2\pi ic_\nu}$ being $\mathbb{Z}^3$-periodic.
Finally let $(\mathrm{sh}_\mu f)(c,g):=f(c+e_\mu,g)$.

\begin{named}{Lemma G}
\begin{itemize}
\item[(G1)] For real $p$ the maps $\mathrm{Fw}_p$ and $U_p$ are unitary.
\item[(G2)] For every $p\in\mathbb{C}^3$ one has $T(p)=\mathrm{Fw}_0^{-1}\,\mathrm{Tpg}(p)\,\mathrm{Fw}_0$; more precisely,
\[
\mathrm{Fw}_0^{-1}\,\mathrm{Dpg}(p)\,\mathrm{Fw}_0=\bigoplus_{\mathrm{jf}}\mathrm{et}^{-1}\cdot D^{(\mathrm{et}\cdot p_{\mathrm{jf}};\,0)},
\qquad
\mathrm{Fw}_0^{-1}\,\mathrm{Ppg}(p)\,\mathrm{Fw}_0=\left(\mathrm{uh}(p_{\mathrm{jf}})\,\mathrm{uh}(p_{\mathrm{jf}'})\right)_{\mathrm{jf},\mathrm{jf}'}\otimes\mathrm{Plt}.
\]
\item[(G3)] For real $p$,
\[
\mathrm{Dpg}(p)=U_p^{-1}DU_p,\qquad\mathrm{Ppg}(p)=U_p^{-1}\,\mathrm{Pbl}\,U_p,\qquad\mathrm{Tpg}(p)=U_p^{-1}\,\mathrm{Tb}(p)\,U_p;
\]
consequently $T(p)=\mathrm{Fw}_p^{-1}\,\mathrm{Tb}(p)\,\mathrm{Fw}_p$, $\mathrm{Fw}_p^{-1}D\,\mathrm{Fw}_p=\bigoplus_{\mathrm{jf}}\mathrm{et}^{-1}\cdot D^{(\mathrm{et}\cdot p_{\mathrm{jf}};\,0)}$, $\mathrm{Fw}_p^{-1}\,\mathrm{Plt}\,\mathrm{Fw}_p=\bigoplus_{\mathrm{jf}}\mathrm{Plt}$, and $\mathrm{Ppg}(p)$ is an orthogonal projection.
\item[(G4)] For every $p\in\mathbb{C}^3$ and every $\nu$,
\[
\mathrm{Dpg}(p+2\pi e_\nu)=\mathrm{Mper}_\nu^{-1}\,\mathrm{Dpg}(p)\,\mathrm{Mper}_\nu,\qquad
\mathrm{Ppg}(p+2\pi e_\nu)=\mathrm{Mper}_\nu^{-1}\,\mathrm{Ppg}(p)\,\mathrm{Mper}_\nu.
\]
\item[(G5)] $\mathscr{H}_\theta$ is the orthogonal direct sum of its subspaces $\mathrm{Hb}(p)$, $p\in P_{\ell,\theta}$.
On $\mathrm{Hb}(p)$ the norm of $\mathscr{H}_\theta$ is $\ell^{\frac32}$ times the norm of $\mathrm{Hb}(p)$.
The operator $T_\theta$ maps every $\mathrm{Hb}(p)$ into itself and acts on it as $\mathrm{Tb}(p)$.
\end{itemize}
\end{named}

\begin{proof}
\emph{(G1)}
Since $\mathrm{e}^{ip_{\mathrm{jf}}\cdot e_\mu}=\mathrm{e}^{ip_\mu}$, the map $\mathrm{Fw}_p$ takes its values in $\mathrm{Hb}(p)$.
For real $p$ the factors $\mathrm{e}^{ip\cdot c}$ cancel in the scalar product, and (4.4) gives
\[
\left\|\mathrm{Fw}_p(u_{\mathrm{jf}})\right\|^2=\sum_{\mathrm{jf},\mathrm{jf}'}\left(\mathrm{et}^3\sum_{c\in\Delta_0}\mathrm{e}^{2\pi i(\mathrm{jf}'-\mathrm{jf})\cdot c}\right)(u_{\mathrm{jf}},u_{\mathrm{jf}'})=\sum_{\mathrm{jf}}\|u_{\mathrm{jf}}\|^2.
\]
Both spaces have the dimension $192\,\mathfrak{N}^3$; hence $\mathrm{Fw}_p$ is unitary.
The assertion on $U_p$ follows from $|\mathrm{e}^{ip\cdot c}|=1$.

\emph{(G2)}
Let $\mathrm{hp}(c,g)=\mathrm{e}^{2\pi i\mathrm{jf}'\cdot c}\,u(g)$ with $u\in{V_{\mathrm{v}}}$.
Then $\mathrm{hp}(x^{(i)})=\mathrm{e}^{2\pi i\mathrm{jf}'\cdot c}\,u(gs_i)$ for $i\le2$, and
\[
\mathrm{e}^{i\,\mathrm{et}\cdot p\cdot(ge_1)}\,\mathrm{hp}(x^{(3)})=\mathrm{e}^{2\pi i\mathrm{jf}'\cdot c}\,\mathrm{e}^{i\,\mathrm{et}\cdot p_{\mathrm{jf}'}\cdot(ge_1)}\,u(gs_3);
\]
so by (I.2.6)
\[
(\mathrm{Dpg}(p)\mathrm{hp})(c,g)=\mathrm{e}^{2\pi i\mathrm{jf}'\cdot c}\,\mathrm{et}^{-1}\left(D^{(\mathrm{et}\cdot p_{\mathrm{jf}'};\,0)}u\right)(g).
\]
Write $c'=[c]+c''$ with $c''\in\Delta_0$, and note that $\mathrm{e}^{2\pi i\mathrm{jf}'\cdot[c]}=1$.
By (4.3) and (4.6),
\begin{align*}
(\mathrm{Ppg}(p)\mathrm{hp})(c,g)
&=\mathrm{e}^{-ip\cdot(c-[c])}\left(\mathrm{et}^3\sum_{c''\in\Delta_0}\mathrm{e}^{ip_{\mathrm{jf}'}\cdot c''}\right)\frac{1}{48}\sum_{g'}u(g')
=\mathrm{e}^{-ip\cdot(c-[c])}\,\mathrm{uh}(p_{\mathrm{jf}'})\,J^*u\\
&=\sum_{\mathrm{jf}\in\mathrm{Mf}}\mathrm{e}^{2\pi i\mathrm{jf}\cdot c}\,\mathrm{uh}(p_{\mathrm{jf}})\,\mathrm{uh}(p_{\mathrm{jf}'})\,(\mathrm{Plt}\,u)(g).
\end{align*}
Finally $\mathrm{et}^{-1}\cdot D^{(\mathrm{pq};\,\mathrm{et}\cdot\mu)}=\mathrm{et}^{-1}\cdot D^{(\mathrm{pq};\,0)}+\mu$; with the definition of $T(p)$ this gives the first assertion.

\emph{(G3)}
One has $(U_p\mathrm{hp})(x^{(i)})=\mathrm{e}^{ip\cdot c}\,\mathrm{hp}(x^{(i)})$ for $i\le2$ and $(U_p\mathrm{hp})(x^{(3)})=\mathrm{e}^{ip\cdot c}\,\mathrm{e}^{i\,\mathrm{et}\cdot p\cdot(ge_1)}\,\mathrm{hp}(x^{(3)})$; hence $DU_p=U_p\,\mathrm{Dpg}(p)$.
Further
\[
(\mathrm{Pbl}\,U_p\mathrm{hp})(c,g)=\mathrm{et}^3\cdot\frac{1}{48}\sum_{c'\in\Delta_{[c]}}\sum_{g'}\mathrm{e}^{ip\cdot c'}\,\mathrm{hp}(c',g')=\mathrm{e}^{ip\cdot c}\,(\mathrm{Ppg}(p)\mathrm{hp})(c,g).
\]
The remaining assertions follow from (G2), (G1), the relation $\mathrm{Fw}_p=U_p\mathrm{Fw}_0$, and the fact that $\mathrm{Plt}$ commutes with the multiplication by functions of $c$.

\emph{(G4)}
The operator $\mathrm{Mper}_\nu^{-1}\,\mathrm{Dpg}(p)\,\mathrm{Mper}_\nu$ differs from $\mathrm{Dpg}(p)$ only in the hop of type $3$, where the additional factor
\[
\mathrm{e}^{-2\pi ic_\nu}\,\mathrm{e}^{2\pi i(c+\mathrm{et}\cdot ge_1)_\nu}=\mathrm{e}^{i\,\mathrm{et}\cdot2\pi e_\nu\cdot(ge_1)}
\]
appears.
The kernel of $\mathrm{Mper}_\nu^{-1}\,\mathrm{Ppg}(p)\,\mathrm{Mper}_\nu$ is
\[
\mathrm{e}^{-2\pi ic_\nu}\,\mathrm{e}^{-ip\cdot(c-c')}\,\mathrm{e}^{2\pi ic'_\nu}=\mathrm{e}^{-i(p+2\pi e_\nu)\cdot(c-c')}.
\]

\emph{(G5)}
By (Lz1) and (B2), $D$ and $\mathrm{Pbl}$ commute with every $\mathrm{sh}_\mu$.
On $\mathscr{H}_\theta$ the maps $\mathrm{sh}_\mu$ are commuting unitary maps (translations of the period cell), and $\mathrm{sh}_\mu^{\ell}=-\mathrm{e}^{i\theta_\mu}$.
Commuting unitary maps admit an orthogonal decomposition of the space into joint eigenspaces, since commuting normal matrices are diagonal in one orthonormal basis \cite[Section~2.5]{HornJohnson2013}.
The eigenvalues of $\mathrm{sh}_\mu$ are numbers $\mathrm{e}^{ip_\mu}$ with $\mathrm{e}^{i\ell p_\mu}=-\mathrm{e}^{i\theta_\mu}$, that is, $p_\mu\in(\theta_\mu+\pi+2\pi\mathbb{Z})/\ell$;
modulo $2\pi$ these are the $\ell$ values that occur in $P_{\ell,\theta}$.
The joint eigenspace for the eigenvalues $(\mathrm{e}^{ip_\mu})_\mu$ is
\[
\left\{f\in\mathscr{H}_\theta:\ f(c+e_\mu,g)=\mathrm{e}^{ip_\mu}f(c,g)\right\}=\mathrm{Hb}(p),
\]
since every $f\in\mathrm{Hb}(p)$ satisfies $f(c+\ell e_\mu,g)=\mathrm{e}^{i\ell p_\mu}f(c,g)=-\mathrm{e}^{i\theta_\mu}f(c,g)$.
For $f\in\mathrm{Hb}(p)$ the function $|f|_2^2$ is $\mathbb{Z}^3$-periodic, and a period cell of $\ell\mathbb{Z}^3$ consists of $\ell^3$ translates of $\Delta_0$; this gives the factor $\ell^{\frac32}$.
Finally, $T_\theta=D+\mu+\mathsf{b}\,\mathrm{Pbl}$ commutes with the $\mathrm{sh}_\mu$.
\end{proof}

\emph{Clause (g1): the operator $T_\theta$.}
By (G5), (G3) and (G1), the map
\[
\bigoplus_{p\in P_{\ell,\theta}}\left(\mathbb{C}^{\mathrm{Mf}}\otimes{V_{\mathrm{v}}}\right)\to\mathscr{H}_\theta,\qquad(\mathrm{ua}_p)_p\mapsto\ell^{-\frac32}\sum_p\mathrm{Fw}_p\,\mathrm{ua}_p,
\]
is unitary, and it transforms $T_\theta$ into the direct sum of the operators $T(p)$.

\emph{Clause (g1): the operator $\mathrm{Dfl}_\theta^{-1}$.}
For $p\in P_{\ell,\theta}$ and $v\in\mathbb{C}^4$ the function $y\mapsto\mathrm{e}^{ip\cdot y}v$ lies in $\mathscr{H}^{\mathrm{u}}_\theta$, because $\mathrm{e}^{i\ell p_\mu}=-\mathrm{e}^{i\theta_\mu}$.
For $p,\tilde p\in P_{\ell,\theta}$ one has
\[
\sum_{y\in\{0,\dots,\ell-1\}^3}\mathrm{e}^{i(p-\tilde p)\cdot y}=\ell^3\,\mathbf{1}[p=\tilde p].
\]
Hence the map
\[
\bigoplus_{p\in P_{\ell,\theta}}\mathbb{C}^4\to\mathscr{H}^{\mathrm{u}}_\theta,\qquad(v_p)_p\mapsto\ell^{-\frac32}\sum_p\mathrm{e}^{ip\cdot y}\,v_p,
\]
is isometric, and therefore unitary, both dimensions being $4\ell^3$.
Write $\mathrm{e}^{ip\cdot(\cdot)}v$ for the function $y\mapsto\mathrm{e}^{ip\cdot y}v$.
By (4.6) multiplied by $\mathrm{e}^{ip\cdot c}$, by (Lz7) together with $\mathrm{et}^{-1}\cdot D^{(\mathrm{pq};\,\mathrm{et}\cdot\mu)}=\mathrm{et}^{-1}\cdot D^{(\mathrm{pq};\,0)}+\mu$ (an invertible map of ${V_{\mathrm{v}}}$ for real $\mathrm{pq}$, by (Lz7) and (H1)), and by (4.3),
\begin{gather*}
\left(Q^*\,\mathrm{e}^{ip\cdot(\cdot)}v\right)(c,g)=\mathrm{e}^{ip\cdot[c]}\,v=\sum_{\mathrm{jf}}\mathrm{uh}(p_{\mathrm{jf}})\,\mathrm{e}^{ip_{\mathrm{jf}}\cdot c}\,(Jv)(g),\\
(D+\mu)^{-1}\left(\mathrm{e}^{ip_{\mathrm{jf}}\cdot c}\,u(g)\right)=\mathrm{e}^{ip_{\mathrm{jf}}\cdot c}\left(\left(\mathrm{et}^{-1}\cdot D^{(\mathrm{et}\cdot p_{\mathrm{jf}};\,\mathrm{et}\cdot\mu)}\right)^{-1}u\right)(g),\\
\left(Q\left(\mathrm{e}^{ip_{\mathrm{jf}}\cdot(\cdot)}u\right)\right)(y)=\mathrm{et}^3\cdot\frac{1}{48}\sum_{c\in\Delta_y}\sum_{g}\mathrm{e}^{ip_{\mathrm{jf}}\cdot c}\,u(g)=\mathrm{e}^{ip_{\mathrm{jf}}\cdot y}\,\mathrm{uh}(p_{\mathrm{jf}})\,J^*u=\mathrm{e}^{ip\cdot y}\,\mathrm{uh}(p_{\mathrm{jf}})\,J^*u.
\end{gather*}
Together with clause (b) this yields
\[
\mathrm{Dfl}_\theta^{-1}\left(\mathrm{e}^{ip\cdot(\cdot)}v\right)=\mathrm{e}^{ip\cdot(\cdot)}\left[\frac{1}{\mathsf{b}}+\sum_{\mathrm{jf}}\mathrm{uh}(p_{\mathrm{jf}})^2\,J^*\left(\mathrm{et}^{-1}\cdot D^{(\mathrm{et}\cdot p_{\mathrm{jf}};\,\mathrm{et}\cdot\mu)}\right)^{-1}J\right]v=\mathrm{e}^{ip\cdot(\cdot)}\,\mathrm{Ch}(p)\,v.
\]
Thus the second unitary map transforms $\mathrm{Dfl}_\theta^{-1}$ into the direct sum of the matrices $\mathrm{Ch}(p)$, $p\in P_{\ell,\theta}$.

\paragraph{4.7 Clause (g2)}
Let $p$ be real.
On $\mathbb{C}^{\mathrm{Mf}}\otimes{V_{\mathrm{v}}}$ put $\mathrm{Abd}:=\bigoplus_{\mathrm{jf}}\mathrm{et}^{-1}\cdot D^{(\mathrm{et}\cdot p_{\mathrm{jf}};\,\mathrm{et}\cdot\mu)}$, and define
\begin{gather*}
\mathrm{Ucol}:\mathbb{C}^4\to\mathbb{C}^{\mathrm{Mf}}\otimes{V_{\mathrm{v}}},\qquad\mathrm{Ucol}\,v:=\left(\mathrm{uh}(p_{\mathrm{jf}})\,Jv\right)_{\mathrm{jf}},\\
\mathrm{Urow}:\mathbb{C}^{\mathrm{Mf}}\otimes{V_{\mathrm{v}}}\to\mathbb{C}^4,\qquad\mathrm{Urow}(u_{\mathrm{jf}}):=\sum_{\mathrm{jf}}\mathrm{uh}(p_{\mathrm{jf}})\,J^*u_{\mathrm{jf}}.
\end{gather*}
Since $\mathrm{Plt}=JJ^*$, one has $\left(\mathrm{uh}(p_{\mathrm{jf}})\,\mathrm{uh}(p_{\mathrm{jf}'})\right)_{\mathrm{jf},\mathrm{jf}'}\otimes\mathrm{Plt}=\mathrm{Ucol}\,\mathrm{Urow}$, and so $T(p)=\mathrm{Abd}+\mathsf{b}\,\mathrm{Ucol}\,\mathrm{Urow}$.

\emph{The determinant.}
By (Lz7), $\operatorname{Re}\left(\mathrm{et}^{-1}\cdot D^{(\mathrm{pq};\,\mathrm{et}\cdot\mu)}\right)\ge\mu$ for real $\mathrm{pq}$.
Hence $\mathrm{Abd}$ is invertible, and
\[
\det\mathrm{Abd}=\prod_{\mathrm{jf}}\det\left(\mathrm{et}^{-1}\cdot D^{(\mathrm{et}\cdot p_{\mathrm{jf}};\,\mathrm{et}\cdot\mu)}\right)=\prod_{\mathrm{jf}}\mathrm{et}^{-192}\cdot\mathfrak{d}_{\mathrm{et}\cdot\mu}(\mathrm{et}\cdot p_{\mathrm{jf}}),
\]
each factor being positive by Proposition~30.1\,(a), as $0<\mathrm{et}\cdot\mu\le1$.
By (H4),
\begin{gather*}
\det T(p)=\det\mathrm{Abd}\cdot\det\left(1+\mathsf{b}\,\mathrm{Abd}^{-1}\,\mathrm{Ucol}\,\mathrm{Urow}\right)=\det\mathrm{Abd}\cdot\det\left(\mathbf{1}_4+\mathsf{b}\,\mathrm{Urow}\,\mathrm{Abd}^{-1}\,\mathrm{Ucol}\right),\\
\mathbf{1}_4+\mathsf{b}\,\mathrm{Urow}\,\mathrm{Abd}^{-1}\,\mathrm{Ucol}=\mathbf{1}_4+\mathsf{b}\sum_{\mathrm{jf}}\mathrm{uh}(p_{\mathrm{jf}})^2\,J^*\left(\mathrm{et}^{-1}\cdot D^{(\mathrm{et}\cdot p_{\mathrm{jf}};\,\mathrm{et}\cdot\mu)}\right)^{-1}J=\mathsf{b}\,\mathrm{Ch}(p).
\end{gather*}
Since $\det\left(\mathsf{b}\,\mathrm{Ch}(p)\right)=\mathsf{b}^4\det\mathrm{Ch}(p)$, this is the formula for $\det T(p)$ in (g2).
Its positivity follows from $\det T(p)=\det\mathrm{Tb}(p)>0$, by (G3) and (4.7).

\emph{The real part.}
Let
\[
\mathrm{Yq}_{\mathrm{jf}}:=J^*\left(\mathrm{et}^{-1}\cdot D^{(\mathrm{et}\cdot p_{\mathrm{jf}};\,\mathrm{et}\cdot\mu)}\right)^{-1}J,
\]
a $4\times4$ matrix.
By (H1), $\operatorname{Re}\left(\mathrm{et}^{-1}\cdot D^{(\mathrm{et}\cdot p_{\mathrm{jf}};\,\mathrm{et}\cdot\mu)}\right)^{-1}\ge0$; as $J^*$ is the adjoint of $J$, it follows that $\operatorname{Re}\mathrm{Yq}_{\mathrm{jf}}\ge0$.
The numbers $\mathrm{uh}(p_{\mathrm{jf}})$ are real, because
\[
u_1(z)=\mathfrak{N}^{-1}\left(1+2\sum_{j=1}^{(\mathfrak{N}-1)/2}\cos\frac{zj}{\mathfrak{N}}\right).
\]
Consequently
\[
\operatorname{Re}\mathrm{Ch}(p)=\frac{1}{\mathsf{b}}+\sum_{\mathrm{jf}}\mathrm{uh}(p_{\mathrm{jf}})^2\operatorname{Re}\mathrm{Yq}_{\mathrm{jf}}\ge\frac{1}{\mathsf{b}}.
\]

\emph{The form of $\mathrm{Ch}(p)$.}
Let $\mathrm{pq}$ be real and $\mathbf{m}>0$ real, and put $O:=D^{(\mathrm{pq};\,\mathbf{m})}$ and $\mathrm{Yq}:=J^*O^{-1}J$.
Three facts are needed.
\begin{itemize}
\item[(i)] $\mathrm{Yq}$ commutes with $\Sigma_3$: so does $O$, by (Lz7), and $J$, $J^*$ commute with constant spin matrices.
Hence $\mathrm{Yq}=\operatorname{diag}(\mathrm{Yq}^+,\mathrm{Yq}^-)$ with $2\times2$ matrices $\mathrm{Yq}^+$ and $\mathrm{Yq}^-$.
\item[(ii)] $\mathrm{Yq}^*=\gamma_5\,\mathrm{Yq}\,\gamma_5$.
Indeed, $O^*=\gamma_5O\gamma_5$ by (Lz7), so that $(O^{-1})^*=\gamma_5O^{-1}\gamma_5$.
The matrix
$\gamma_5=\left(\begin{array}{cc}0&\mathbf{1}_2\\ \mathbf{1}_2&0\end{array}\right)$
exchanges the two blocks: $\gamma_5\operatorname{diag}(\mathrm{Yq}^+,\mathrm{Yq}^-)\gamma_5=\operatorname{diag}(\mathrm{Yq}^-,\mathrm{Yq}^+)$.
Therefore $\mathrm{Yq}^-=(\mathrm{Yq}^+)^*$.
\item[(iii)] On ${V_{\mathrm{v}}}$ let $(\mathrm{Xa}\,u)(g):=\mathrm{Cc}\,\overline{u(-g)}$, an antilinear map; here $\mathrm{Cc}=\mathbf{1}_2\otimes\sigma_2$, and $-g:=(-1)g$, where $-1$ belongs to $O_{\mathrm{cub}}$ and commutes with all its elements.
On $\mathbb{C}^4$ let $\mathrm{cq}\,v:=\mathrm{Cc}\,\overline{v}$.
As $\gamma_\mu$ is Hermitian, $\overline{\gamma_\mu}=\gamma_\mu^T$, and (I.c5) gives $\mathrm{Cc}\,\overline{\gamma_\mu}\,\mathrm{Cc}^{-1}=-\gamma_\mu$.
One has $M_i(g)=\frac{w_i}{2}\left(\rho+\gamma\cdot(g\hat n_i)\right)$ with the real vector $g\hat n_i$; therefore
\[
\mathrm{Cc}\,\overline{M_i(-g)}=\frac{w_i}{2}\left(\rho-\mathrm{Cc}\,\overline{\gamma}\cdot(g\hat n_i)\,\mathrm{Cc}^{-1}\right)\mathrm{Cc}=\frac{w_i}{2}\left(\rho+\gamma\cdot(g\hat n_i)\right)\mathrm{Cc}=M_i(g)\,\mathrm{Cc}.
\]
Moreover $\overline{\mathrm{e}^{i\mathrm{pq}\cdot((-g)e_1)}}=\mathrm{e}^{i\mathrm{pq}\cdot(ge_1)}$ ($\mathrm{pq}$ being real), $(-g)s_i=-(gs_i)$, and $\mathbf{m}$ is real.
Applying $\mathrm{Cc}\,\overline{\,\cdot\,}$ to (I.2.6) at the label $-g$, one obtains
\begin{multline*}
(\mathrm{Xa}\,Ou)(g)=\mathbf{m}\,(\mathrm{Xa}\,u)(g)-\sum_{i\le2}M_i(g)\left((\mathrm{Xa}\,u)(gs_i)-(\mathrm{Xa}\,u)(g)\right)\\
-M_3(g)\left(\mathrm{e}^{i\mathrm{pq}\cdot(ge_1)}\,(\mathrm{Xa}\,u)(gs_3)-(\mathrm{Xa}\,u)(g)\right)=(O\,\mathrm{Xa}\,u)(g).
\end{multline*}
Thus $\mathrm{Xa}\,O=O\,\mathrm{Xa}$, and $\mathrm{Xa}\,O^{-1}=O^{-1}\,\mathrm{Xa}$.
Further $\mathrm{Xa}\,Jv=J\,\mathrm{cq}\,v$, and $J^*\mathrm{Xa}\,u=\frac{1}{48}\sum_g\mathrm{Cc}\,\overline{u(-g)}=\mathrm{cq}\,J^*u$.
Hence
\[
\mathrm{Yq}\,\mathrm{cq}=J^*O^{-1}\,\mathrm{Xa}\,J=J^*\mathrm{Xa}\,O^{-1}J=\mathrm{cq}\,\mathrm{Yq}.
\]
In the blocks this reads $\mathrm{Yq}^+\sigma_2\overline{v}=\sigma_2\overline{\mathrm{Yq}^+v}$ for all $v\in\mathbb{C}^2$, that is, $\mathrm{Yq}^+=\sigma_2\,\overline{\mathrm{Yq}^+}\,\sigma_2$.
Write $\mathrm{Yq}^+=\mathrm{yc}_0+\sum_\mu\mathrm{yc}_\mu\sigma_\mu$ with complex numbers $\mathrm{yc}_0$, $\mathrm{yc}_\mu$; this is possible because $\mathbf{1}_2$ and the $\sigma_\mu$ form a basis of the $2\times2$ matrices.
As $\sigma_2\overline{\sigma_\mu}\sigma_2=-\sigma_\mu$ for $\mu=1,2,3$, one finds $\sigma_2\,\overline{\mathrm{Yq}^+}\,\sigma_2=\overline{\mathrm{yc}_0}-\sum_\mu\overline{\mathrm{yc}_\mu}\,\sigma_\mu$.
So $\mathrm{yc}_0$ is real, and every $\mathrm{yc}_\mu$ is purely imaginary.
\end{itemize}
By (i)--(iii), $\mathrm{Yq}^+=\alpha'-i\sigma\cdot\beta'$ and $\mathrm{Yq}^-=(\mathrm{Yq}^+)^*=\alpha'+i\sigma\cdot\beta'$ with $\alpha'$ real and $\beta'\in\mathbb{R}^3$; hence
\[
\mathrm{Yq}=\alpha'-i\operatorname{diag}\left(\sigma\cdot\beta',-\sigma\cdot\beta'\right)=\alpha'-i\gamma\cdot\beta'.
\]
Apply this with $\mathrm{pq}=\mathrm{et}\cdot p_{\mathrm{jf}}$ and $\mathbf{m}=\mathrm{et}\cdot\mu$.
Then $\mathrm{et}^{-1}\,\mathrm{Yq}_{\mathrm{jf}}=\alpha'-i\gamma\cdot\beta'$, so that $\mathrm{Yq}_{\mathrm{jf}}=\alpha_{\mathrm{jf}}-i\gamma\cdot\beta_{\mathrm{jf}}$ with $\alpha_{\mathrm{jf}}$ real and $\beta_{\mathrm{jf}}\in\mathbb{R}^3$.
Since $i\gamma\cdot\beta_{\mathrm{jf}}$ is anti-Hermitian, $\alpha_{\mathrm{jf}}=\operatorname{Re}\mathrm{Yq}_{\mathrm{jf}}\ge0$.
Therefore
\[
\mathrm{Ch}(p)=\alpha-i\gamma\cdot\beta,\qquad\alpha:=\frac{1}{\mathsf{b}}+\sum_{\mathrm{jf}}\mathrm{uh}(p_{\mathrm{jf}})^2\,\alpha_{\mathrm{jf}}\ge\frac{1}{\mathsf{b}},\qquad\beta:=\sum_{\mathrm{jf}}\mathrm{uh}(p_{\mathrm{jf}})^2\,\beta_{\mathrm{jf}}\in\mathbb{R}^3.
\]
Moreover $\mathrm{Ch}(p)=\operatorname{diag}\left(\alpha-i\sigma\cdot\beta,\ \alpha+i\sigma\cdot\beta\right)$, and $\det\left(\alpha\mp i\sigma\cdot\beta\right)=\alpha^2+|\beta|_2^2$, because $\sigma\cdot\beta$ has the trace $0$ and the square $|\beta|_2^2$, hence the eigenvalues $\pm|\beta|_2$.
So $\det\mathrm{Ch}(p)=\left(\alpha^2+|\beta|_2^2\right)^2$.

\paragraph{4.8 Clause (g3)}
\emph{Analyticity, periodicity, independence.}
The function $p\mapsto\det T(p)$ is entire, since the entries of $T(p)$ are entire (beginning of 4.6), and it is positive on $\mathbb{R}^3$ by (g2).
Hence $\mathrm{Flg}=\log\det T$ is real analytic on $\mathbb{R}^3$.
By (G2) and (G4), $T(p+2\pi e_\nu)$ is similar to $T(p)$; so $\mathrm{Flg}$ is $2\pi\mathbb{Z}^3$-periodic.
The definition of $T(p)$ involves $\mathfrak{N}$, $\mu$, $\mathsf{b}$, $\rho$ and $p$; it involves neither $\ell$ nor $\theta$.

\emph{Invariance under $O_{\mathrm{cub}}$.}
For a unit vector $\hat n\in\mathbb{R}^3$ let $S_{\hat n}:=(\gamma\cdot\hat n)\,\gamma_5$; this is the matrix $S_{\mathrm{nv}}$ of Section~I.4.1.
The matrices $\gamma\cdot\hat n$ and $\gamma_5$ are Hermitian with square $1$, so $S_{\hat n}$ is unitary.
From $\gamma_5(\gamma\cdot v)\gamma_5=-\gamma\cdot v$ and $(\gamma\cdot\hat n)(\gamma\cdot v)=2(v\cdot\hat n)-(\gamma\cdot v)(\gamma\cdot\hat n)$ one obtains
\[
S_{\hat n}\,(\gamma\cdot v)\,S_{\hat n}^{-1}=-(\gamma\cdot\hat n)(\gamma\cdot v)(\gamma\cdot\hat n)=\gamma\cdot\left(v-2(v\cdot\hat n)\,\hat n\right);
\]
for $\hat n=\hat n_i$ this is $\gamma\cdot(s_iv)$.
Let $r\in O_{\mathrm{cub}}$.
Since $s_0$, $s_1$, $s_2$ generate $O_{\mathrm{cub}}$ (Section~I.1.2), one may choose a word $r=s_{i_1}\cdots s_{i_{\mathrm{iw}}}$ and put $\mathrm{Sw}_r:=S_{\hat n_{i_1}}\cdots S_{\hat n_{i_{\mathrm{iw}}}}$.
This is a unitary matrix with $\mathrm{Sw}_r\,(\gamma\cdot v)\,\mathrm{Sw}_r^{-1}=\gamma\cdot(rv)$ for $v\in\mathbb{R}^3$.
Define
\[
(\mathrm{Vw}_rf)(c,g):=\mathrm{Sw}_r\,f(r^{-1}c,r^{-1}g).
\]
\begin{itemize}
\item[(1)] $\mathrm{Vw}_r$ maps $\mathrm{Hb}(p)$ unitarily onto $\mathrm{Hb}(rp)$.
Indeed, for $f\in\mathrm{Hb}(p)$ and $z\in\mathbb{Z}^3$ one has $f(c+z,g)=\mathrm{e}^{ip\cdot z}f(c,g)$; hence
\[
(\mathrm{Vw}_rf)(c+e_\mu,g)=\mathrm{Sw}_r\,f(r^{-1}c+r^{-1}e_\mu,r^{-1}g)=\mathrm{e}^{ip\cdot(r^{-1}e_\mu)}\,(\mathrm{Vw}_rf)(c,g),
\]
and $p\cdot(r^{-1}e_\mu)=(rp)_\mu$, $r$ being orthogonal.
The map $(c,g)\mapsto(r^{-1}c,r^{-1}g)$ is a bijection of $\Delta_0\times O_{\mathrm{cub}}$ by (B1), and $\mathrm{Sw}_r$ is unitary.
The inverse is $f\mapsto\mathrm{Sw}_r^{-1}f(rc,rg)$.
\item[(2)] $D\,\mathrm{Vw}_r=\mathrm{Vw}_rD$.
With $r^{-1}x:=(r^{-1}c,r^{-1}g)$ one has $r^{-1}(x^{(i)})=(r^{-1}x)^{(i)}$: for $i\le2$ both sides equal $(r^{-1}c,r^{-1}gs_i)$, and for $i=3$ both equal $(r^{-1}c+\mathrm{et}\cdot(r^{-1}g)e_1,\ r^{-1}gs_3)$.
Further
\[
M_i(x)\,\mathrm{Sw}_r=\frac{w_i}{2}\,\mathrm{Sw}_r\left(\rho+\gamma\cdot(r^{-1}g\hat n_i)\right)=\mathrm{Sw}_r\,M_i(r^{-1}x).
\]
Therefore
\[
(D\,\mathrm{Vw}_rf)(x)=-\mathrm{et}^{-1}\sum_iM_i(x)\,\mathrm{Sw}_r\left[f\left((r^{-1}x)^{(i)}\right)-f(r^{-1}x)\right]=\mathrm{Sw}_r\,(Df)(r^{-1}x).
\]
\item[(3)] $\mathrm{Pbl}\,\mathrm{Vw}_r=\mathrm{Vw}_r\,\mathrm{Pbl}$, because by (B1) the map $(c',g')\mapsto(r^{-1}c',r^{-1}g')$ takes $\Delta_{[c]}\times O_{\mathrm{cub}}$ onto $\Delta_{[r^{-1}c]}\times O_{\mathrm{cub}}$.
\end{itemize}
Consequently $\mathrm{Tb}(rp)=\mathrm{Vw}_r\,\mathrm{Tb}(p)\,\mathrm{Vw}_r^{-1}$, and by (G3)
\[
\det T(rp)=\det\mathrm{Tb}(rp)=\det\mathrm{Tb}(p)=\det T(p).
\]

\emph{The summation formula.}
The function $\mathrm{Flg}$ is periodic and of class $C^4$.
For $z\neq0$ choose $\nu$ with $|z_\nu|=|z|$; by (F5) with $\mathrm{kd}=4$,
\[
|\mathrm{Flgh}(z)|\le|z|^{-4}\sup_{\mathbb{R}^3}\left|\partial_\nu^4\mathrm{Flg}\right|,
\]
and the supremum is finite, $\partial_\nu^4\mathrm{Flg}$ being continuous and periodic.
By (F4), $\sum_z|\mathrm{Flgh}(z)|<\infty$, and by (F1) Proposition~30.1\,(d) applies to $\mathrm{fper}=\mathrm{Flg}$, with $\ell$ in place of $n$.
By (g1), $\log\det T_\theta=\sum_{p\in P_{\ell,\theta}}\mathrm{Flg}(p)$.
So (8.2), at $\theta$ and at $\theta=0$, gives
\begin{equation}
\begin{aligned}
\log\det T_\theta&=\ell^3\sum_{w\in\mathbb{Z}^3}\mathrm{Flgh}(\ell w)\,(-1)^{w_1+w_2+w_3}\,\mathrm{e}^{i\theta\cdot w},\\
\log\frac{\det T_\theta}{\det T_{\theta=0}}&=\ell^3\sum_{w\neq0}\mathrm{Flgh}(\ell w)\,(-1)^{w_1+w_2+w_3}\left(\mathrm{e}^{i\theta\cdot w}-1\right).
\end{aligned}
\tag{4.8}
\end{equation}

\paragraph{4.9 Clause (e)}
Let $0<b_-\le\mathsf{b}\le b_+$.

\emph{Constants.}
The numbers $K_{\mathrm{iv}}$, $C_{\mathrm{iv}}$, $\gamma_{\Sigma}$, $C_{\mathrm{sv}}$, $C_{\mathrm{uv}}$, $S_{\mathrm{jf}}$, $C_{\mathrm{Flg}}$, $C_{\Sigma}$ are those of Proposition~30.3, formed in this order:
\begin{gather*}
K_{\mathrm{iv}}:=b_-^{-\frac12}+\frac{8}{\pi^3}\,(3\hat{g})^{-\frac12}+\left(1+\frac{8}{\pi^3}\right)\frac{2+b_+^{\frac12}}{\pi c_{\mathrm{N3}}},\qquad
C_{\mathrm{iv}}:=\max\left\{1,\left(\frac{\pi^3K_{\mathrm{iv}}}{8}\right)^2\right\},\\
\gamma_{\Sigma}:=\min\left\{1,\frac{1}{2\mathrm{e}\cdot(6+b_+)\cdot C_{\mathrm{iv}}}\right\},\qquad
C_{\mathrm{sv}}:=2\left(1+(1+b_+)\cdot C_{\mathrm{iv}}\right),\qquad
C_{\mathrm{uv}}:=\left(1+3\pi^2c_{\mathrm{N3}}^2\right)^{\frac12},\\
S_{\mathrm{jf}}:=1+\frac{4\pi^2+\pi^4/45}{(\pi c_{\mathrm{N3}})^4},\qquad
C_{\mathrm{Flg}}:=26\left[9^4\cdot192\cdot C_{\mathrm{sv}}^4\cdot C_{\mathrm{uv}}^3\cdot S_{\mathrm{jf}}+300\cdot\mathrm{e}\cdot b_+\cdot(2C_{\mathrm{iv}})^4\cdot\max(9,\mathrm{e}\cdot b_+)^3\right],\\
C_{\Sigma}:=2\left(4\pi^2+\frac{\pi^4}{45}\right)C_{\mathrm{Flg}},
\end{gather*}
where $\hat{g}=\rho\lambda_{\mathrm{int}}/2$ and $c_{\mathrm{N3}}=\min\left\{\rho/\pi,\ \rho^3\lambda_{\mathrm{int}}/(98\sqrt3\,\pi^3)\right\}$ (Section~I.8.2, line (C0r)).
By (F4) the number $4\pi^2+\pi^4/45$ is $S_4$; thus $S_{\mathrm{jf}}=1+S_4/(\pi c_{\mathrm{N3}})^4$ and $C_{\Sigma}=2S_4\cdot C_{\mathrm{Flg}}$.
These constants are functions of $(b_-,b_+,\rho,\lambda_{\mathrm{int}})$.

\begin{named}{Lemma E1}[the symbol has no zero but $\mathrm{pq}=0$]
For $\mathrm{pq}\in[-\pi,\pi]^3$ and $u\in{V_{\mathrm{v}}}$,
\[
\left\|D^{(\mathrm{pq};\,0)}u\right\|\ge c_{\mathrm{N3}}\,|\mathrm{pq}|_2\,\|u\|.
\]
Consequently, for $p\in[-\pi,\pi]^3$, $\mathrm{jf}\in\mathrm{Mf}$, $X_{\mathrm{jf}}:=\mathrm{et}^{-1}\cdot D^{(\mathrm{et}\cdot p_{\mathrm{jf}};\,0)}$ and $u\in{V_{\mathrm{v}}}$:
\begin{itemize}
\item[(i)] $\mathrm{et}\cdot p_{\mathrm{jf}}\in[-\pi,\pi]^3$;
\item[(ii)] $|p_{\mathrm{jf}}|_2\ge|p_{\mathrm{jf}}|\ge\pi|\mathrm{jf}|$ for $\mathrm{jf}\neq0$;
\item[(iii)] $\|X_{\mathrm{jf}}u\|\ge c_{\mathrm{N3}}\,|p_{\mathrm{jf}}|_2\,\|u\|$;
\item[(iv)] $\|(X_{\mathrm{jf}}+1)u\|^2\ge\left(1+c_{\mathrm{N3}}^2|p_{\mathrm{jf}}|_2^2\right)\|u\|^2$.
\end{itemize}
\end{named}

\begin{proof}
Let $\|u\|=1$, and put $\mathrm{dsy}:=\|D^{(\mathrm{pq};\,0)}u\|$, $u^{\mathrm{l}}:=\mathrm{Plt}\,u=JJ^*u$ and $u^{\mathrm{h}}:=\mathrm{Phv}\,u$.

(1) By (Lz7), $\hat{g}\,\|u^{\mathrm{h}}\|^2\le\operatorname{Re}\left(u,D^{(\mathrm{pq};\,0)}u\right)\le\mathrm{dsy}$.

(2) For $v\in\mathbb{C}^4$, (I.2.6) gives
\[
\left(D^{(\mathrm{pq};\,0)}Jv\right)(g)=-M_3(g)\left(\mathrm{e}^{i\mathrm{pq}\cdot(ge_1)}-1\right)v=3\left(\rho-\gamma\cdot(ge_1)\right)\left(1-\mathrm{e}^{i\mathrm{pq}\cdot(ge_1)}\right)v.
\]
Each of the six vectors $a\in\{\pm e_\mu\}$ equals $ge_1$ for exactly $8$ of the $48$ labels (I.s4).
Hence
\begin{align*}
J^*D^{(\mathrm{pq};\,0)}J
&=\frac{8\cdot3}{48}\sum_\mu\left[(\rho-\gamma_\mu)\left(1-\mathrm{e}^{i\mathrm{pq}_\mu}\right)+(\rho+\gamma_\mu)\left(1-\mathrm{e}^{-i\mathrm{pq}_\mu}\right)\right]\\
&=\sum_\mu\left[\rho\,(1-\cos\mathrm{pq}_\mu)+i\gamma_\mu\sin\mathrm{pq}_\mu\right]=:\mathrm{Wi}(\mathrm{pq}).
\end{align*}
Let $\mathrm{tw}:=\sum_\mu(1-\cos\mathrm{pq}_\mu)$, and let $\sin\mathrm{pq}$ denote the vector $(\sin\mathrm{pq}_\mu)_\mu$.
Since $(\gamma\cdot v)^2=|v|_2^2$ for real $v$ (I.c1),
\[
\mathrm{Wi}(\mathrm{pq})^*\,\mathrm{Wi}(\mathrm{pq})=\left(\rho\,\mathrm{tw}-i\gamma\cdot\sin\mathrm{pq}\right)\left(\rho\,\mathrm{tw}+i\gamma\cdot\sin\mathrm{pq}\right)=\rho^2\,\mathrm{tw}^2+|\sin\mathrm{pq}|_2^2,
\]
a scalar.
All the numbers $1-\cos\mathrm{pq}_\mu$ are non-negative, so $\mathrm{tw}^2\ge\sum_\mu(1-\cos\mathrm{pq}_\mu)^2$, and
\begin{align*}
\mathrm{Wi}(\mathrm{pq})^*\,\mathrm{Wi}(\mathrm{pq})
&\ge\sum_\mu\left[\sin^2\mathrm{pq}_\mu+\rho^2(1-\cos\mathrm{pq}_\mu)^2\right]
=\sum_\mu4\sin^2\frac{\mathrm{pq}_\mu}{2}\left[\cos^2\frac{\mathrm{pq}_\mu}{2}+\rho^2\sin^2\frac{\mathrm{pq}_\mu}{2}\right]\\
&\ge4\rho^2\sum_\mu\sin^2\frac{\mathrm{pq}_\mu}{2}\ge\left(\frac{2\rho}{\pi}\right)^2|\mathrm{pq}|_2^2;
\end{align*}
here $\rho\le1$ and $|\sin(t/2)|\ge|t|/\pi$ for $|t|\le\pi$ have been used.
As $J$ is isometric, it follows that
\[
\left\|\mathrm{Plt}\,D^{(\mathrm{pq};\,0)}u^{\mathrm{l}}\right\|=\left|\mathrm{Wi}(\mathrm{pq})\,J^*u\right|_2\ge\frac{2\rho}{\pi}\,|\mathrm{pq}|_2\,\|u^{\mathrm{l}}\|.
\]

(3) One has
\[
\left(\left(D^{(\mathrm{pq};\,0)}-D^{(0;\,0)}\right)u\right)(g)=-M_3(g)\left(\mathrm{e}^{i\mathrm{pq}\cdot(ge_1)}-1\right)u(gs_3).
\]
The map $u\mapsto u(\cdot\,s_3)$ is unitary, $|M_3(g)|_2=3(1+\rho)\le6$ (I.m6), and $|\mathrm{e}^{i\mathrm{pq}_\mu}-1|\le|\mathrm{pq}_\mu|$.
So $\|D^{(\mathrm{pq};\,0)}-D^{(0;\,0)}\|\le6|\mathrm{pq}|\le6|\mathrm{pq}|_2$.
Moreover $J^*D^{(0;\,0)}=0$: substituting $g\mapsto gs_i$ in the terms that contain $u(gs_i)$, one finds by (I.m1)
\[
J^*D^{(0;\,0)}u=-\frac{1}{48}\sum_g\sum_{i=0}^{3}M_i(g)\left(u(gs_i)-u(g)\right)=-\frac{1}{48}\sum_g\sum_{i=0}^{3}\left(M_i(gs_i)-M_i(g)\right)u(g)=0.
\]
Therefore
\[
\left\|\mathrm{Plt}\,D^{(\mathrm{pq};\,0)}u^{\mathrm{h}}\right\|=\left\|\mathrm{Plt}\left(D^{(\mathrm{pq};\,0)}-D^{(0;\,0)}\right)u^{\mathrm{h}}\right\|\le6|\mathrm{pq}|_2\,\|u^{\mathrm{h}}\|.
\]

(4) By (2) and (3), using $\|u^{\mathrm{l}}\|\ge1-\|u^{\mathrm{h}}\|$ and $2\rho/\pi\le1$,
\[
\mathrm{dsy}\ge\left\|\mathrm{Plt}\,D^{(\mathrm{pq};\,0)}u\right\|\ge|\mathrm{pq}|_2\left[\frac{2\rho}{\pi}\left(1-\|u^{\mathrm{h}}\|\right)-6\|u^{\mathrm{h}}\|\right]\ge|\mathrm{pq}|_2\left[\frac{2\rho}{\pi}-7\|u^{\mathrm{h}}\|\right].
\]
If $\|u^{\mathrm{h}}\|\le\rho/(7\pi)$, then $\mathrm{dsy}\ge(\rho/\pi)\,|\mathrm{pq}|_2$.
If $\|u^{\mathrm{h}}\|>\rho/(7\pi)$, then (1) gives, since $|\mathrm{pq}|_2\le\sqrt3\,\pi$,
\[
\mathrm{dsy}\ge\frac{\hat{g}\rho^2}{49\pi^2}\ge\frac{\hat{g}\rho^2}{49\pi^2}\cdot\frac{|\mathrm{pq}|_2}{\sqrt3\,\pi}=\frac{\rho^3\lambda_{\mathrm{int}}}{98\sqrt3\,\pi^3}\,|\mathrm{pq}|_2.
\]
In both cases $\mathrm{dsy}\ge c_{\mathrm{N3}}\,|\mathrm{pq}|_2$.

(i) $|\mathrm{et}\cdot p_{\mathrm{jf},\mu}|\le\left(\pi+2\pi(\mathfrak{N}-1)/2\right)/\mathfrak{N}=\pi$.
(ii) $|p_{\mathrm{jf}}|\ge2\pi|\mathrm{jf}|-\pi\ge\pi|\mathrm{jf}|$.
(iii) This is the first assertion at $\mathrm{pq}=\mathrm{et}\cdot p_{\mathrm{jf}}$.
(iv) $\|(X_{\mathrm{jf}}+1)u\|^2=\|X_{\mathrm{jf}}u\|^2+2\operatorname{Re}(u,X_{\mathrm{jf}}u)+\|u\|^2$, and $\operatorname{Re}(u,X_{\mathrm{jf}}u)\ge0$ by (Lz7).
\end{proof}

\begin{named}{Lemma E2}[the inverse at real momentum]
For all $p\in\mathbb{R}^3$, all odd $\mathfrak{N}\ge3$, all $\mu\in(0,1]$ and all $\mathsf{b}\in[b_-,b_+]$,
\[
\|\mathrm{Tb}(p)f\|\ge C_{\mathrm{iv}}^{-1}\,\|f\|\qquad(f\in\mathrm{Hb}(p)).
\]
Hence $\|T(p)^{-1}\|=\|\mathrm{Tpg}(p)^{-1}\|=\|\mathrm{Tb}(p)^{-1}\|\le C_{\mathrm{iv}}$.
\end{named}

\begin{proof}
Since $\mathrm{Hb}(p)$ and $\mathrm{Tb}(p)$ depend on $p$ only modulo $2\pi\mathbb{Z}^3$, one may assume $p\in[-\pi,\pi]^3$.
Let $\|f\|=1$ and $\mathrm{ef}:=\|\mathrm{Tb}(p)f\|$.
If $\mathrm{ef}>1$ there is nothing to prove, because $C_{\mathrm{iv}}\ge1$.
Let $\mathrm{ef}\le1$.
Then $\operatorname{Re}(f,\mathrm{Tb}(p)f)\le\mathrm{ef}$, and (4.7) together with $\mathrm{et}\le1/3$ yields
\[
\mu\le\mathrm{ef},\qquad\mathsf{b}\,\|\mathrm{Pbl}\,f\|^2\le\mathrm{ef},\qquad\|\mathrm{Phv}\,f\|^2\le\frac{\mathrm{ef}\cdot\mathrm{et}}{\hat{g}}\le\frac{\mathrm{ef}}{3\hat{g}}.
\]
So $\|\mathrm{Pbl}\,f\|\le(\mathrm{ef}/b_-)^{\frac12}$ and $\mathsf{b}\,\|\mathrm{Pbl}\,f\|=\mathsf{b}^{\frac12}\left(\mathsf{b}\,\|\mathrm{Pbl}\,f\|^2\right)^{\frac12}\le(b_+\cdot\mathrm{ef})^{\frac12}$, whence
\[
\|Df\|=\left\|\mathrm{Tb}(p)f-\mu f-\mathsf{b}\,\mathrm{Pbl}\,f\right\|\le2\,\mathrm{ef}+(b_+\cdot\mathrm{ef})^{\frac12}\le\left(2+b_+^{\frac12}\right)\mathrm{ef}^{\frac12}.
\]
Write $f=\mathrm{Fw}_p(u_{\mathrm{jf}})$ (G1), let $f^{(0)}(c,g):=\mathrm{e}^{ip\cdot c}\,u_{\mathrm{jf}=0}(g)$ be the term with $\mathrm{jf}=0$, and $f^{(\mathrm{r})}:=f-f^{(0)}$.
By (G3), (G1) and Lemma E1\,(iii), (ii),
\[
\|Df\|^2=\sum_{\mathrm{jf}}\|X_{\mathrm{jf}}u_{\mathrm{jf}}\|^2\ge c_{\mathrm{N3}}^2\pi^2\sum_{\mathrm{jf}\neq0}\|u_{\mathrm{jf}}\|^2=c_{\mathrm{N3}}^2\pi^2\,\|f^{(\mathrm{r})}\|^2;
\]
thus $\|f^{(\mathrm{r})}\|\le\|Df\|/(\pi c_{\mathrm{N3}})$.

By (4.3),
\[
\left(\mathrm{Pbl}\,f^{(0)}\right)(c,g)=\mathrm{et}^3\cdot\frac{1}{48}\sum_{c'\in\Delta_{[c]}}\sum_{g'}\mathrm{e}^{ip\cdot c'}\,u_{\mathrm{jf}=0}(g')=\left(\mathrm{et}^3\sum_{c'\in\Delta_{[c]}}\mathrm{e}^{ip\cdot c'}\right)J^*u_{\mathrm{jf}=0}=\mathrm{e}^{ip\cdot[c]}\,\mathrm{uh}(p)\,J^*u_{\mathrm{jf}=0}.
\]
Summing the geometric progression one finds $u_1(0)=1$ and
\[
u_1(t)=\sin\frac{t}{2}\cdot\left(\mathfrak{N}\sin\frac{t}{2\mathfrak{N}}\right)^{-1}\qquad(0<|t|\le\pi).
\]
For $0<|t|\le\pi$ one has $\sin(|t|/2)\ge|t|/\pi$ and $0<\mathfrak{N}\sin(|t|/(2\mathfrak{N}))\le|t|/2$; hence $u_1(t)\ge2/\pi$ on $[-\pi,\pi]$, and $\mathrm{uh}(p)\ge8/\pi^3$.
Since $(\mathrm{Plt}\,f^{(0)})(c,g)=\mathrm{e}^{ip\cdot c}\,J^*u_{\mathrm{jf}=0}$, it follows that
\[
\left\|\mathrm{Pbl}\,f^{(0)}\right\|=\mathrm{uh}(p)\,\left|J^*u_{\mathrm{jf}=0}\right|_2\ge\frac{8}{\pi^3}\,\left\|\mathrm{Plt}\,f^{(0)}\right\|.
\]
By (G3), $\mathrm{Phv}$ acts on each mode separately; so $\|\mathrm{Phv}\,f^{(0)}\|\le\|\mathrm{Phv}\,f\|$, and
\[
\left\|\mathrm{Plt}\,f^{(0)}\right\|\ge\left\|f^{(0)}\right\|-\left\|\mathrm{Phv}\,f^{(0)}\right\|\ge1-\|f^{(\mathrm{r})}\|-\|\mathrm{Phv}\,f\|.
\]
Using $\|\mathrm{Pbl}\,f\|\ge\|\mathrm{Pbl}\,f^{(0)}\|-\|f^{(\mathrm{r})}\|$ (recall $\|\mathrm{Pbl}\|=1$), one obtains
\begin{align*}
\left(\frac{\mathrm{ef}}{b_-}\right)^{\frac12}\ge\|\mathrm{Pbl}\,f\|
&\ge\frac{8}{\pi^3}\left(1-\|\mathrm{Phv}\,f\|\right)-\left(1+\frac{8}{\pi^3}\right)\|f^{(\mathrm{r})}\|\\
&\ge\frac{8}{\pi^3}-\frac{8}{\pi^3}\left(\frac{\mathrm{ef}}{3\hat{g}}\right)^{\frac12}-\left(1+\frac{8}{\pi^3}\right)\frac{\left(2+b_+^{\frac12}\right)\mathrm{ef}^{\frac12}}{\pi c_{\mathrm{N3}}}.
\end{align*}
This says that $8/\pi^3\le K_{\mathrm{iv}}\cdot\mathrm{ef}^{\frac12}$, that is, $\mathrm{ef}\ge\left(8/(\pi^3K_{\mathrm{iv}})\right)^2\ge1/C_{\mathrm{iv}}$.
The equalities of the norms follow from (G2), (G3) and (G1).
\end{proof}

\begin{named}{Lemma E3}[derivatives; a strip]
Fix $\nu\in\{1,2,3\}$.
For an integer $\mathrm{kd}\ge1$ let
\begin{gather*}
\mathrm{Dpg}^{(\mathrm{kd})}(p):=\partial_\nu^{\mathrm{kd}}\,\mathrm{Dpg}(p),\qquad\mathrm{Ppg}^{(\mathrm{kd})}(p):=\partial_\nu^{\mathrm{kd}}\,\mathrm{Ppg}(p),\\
\mathrm{Tpg}^{(\mathrm{kd})}(p):=\mathrm{Dpg}^{(\mathrm{kd})}(p)+\mathsf{b}\,\mathrm{Ppg}^{(\mathrm{kd})}(p)=\partial_\nu^{\mathrm{kd}}\,\mathrm{Tpg}(p).
\end{gather*}
Let $p=p'+i\,\mathrm{ti}\,e_\nu$ with $p'\in\mathbb{R}^3$ and $\mathrm{ti}$ real.
\begin{itemize}
\item[(i)] If $|\mathrm{ti}|\le1$, then
\begin{gather*}
\left\|\mathrm{Dpg}^{(\mathrm{kd})}(p)\right\|\le9\,\mathrm{et}^{\mathrm{kd}-1},\qquad\left\|\mathrm{Ppg}^{(\mathrm{kd})}(p)\right\|\le\mathrm{e},\\
\operatorname{rank}\mathrm{Ppg}^{(\mathrm{kd})}(p)\le4(\mathrm{kd}+1),\qquad\left\|\mathrm{Ppg}^{(\mathrm{kd})}(p)\right\|_{(1)}\le4(\mathrm{kd}+1)\cdot\mathrm{e}.
\end{gather*}
\item[(ii)] If $|\mathrm{ti}|\le\gamma_{\Sigma}$, then $\mathrm{Tpg}(p)$ is invertible; with $\mathrm{Gpg}(p):=\mathrm{Tpg}(p)^{-1}$ and $\mathrm{Epg}:=\mathrm{Tpg}(p)-\mathrm{Tpg}(p')$ one has
\begin{gather*}
\left\|\mathrm{Epg}\,\mathrm{Gpg}(p')\right\|\le\frac12,\qquad\mathrm{Gpg}(p)=\mathrm{Gpg}(p')\left(1+\mathrm{Epg}\,\mathrm{Gpg}(p')\right)^{-1},\\
\|\mathrm{Gpg}(p)\|\le2C_{\mathrm{iv}},\qquad\mathrm{sgv}_j\left(\mathrm{Gpg}(p)\right)\le2\,\mathrm{sgv}_j\left(\mathrm{Gpg}(p')\right)\ \text{for all}\ j.
\end{gather*}
\end{itemize}
\end{named}

\begin{proof}
Two remarks are used.

(E3.1) The map $\mathrm{hp}\mapsto\mathrm{hp}(x^{(3)})$ is the operator $R_3$ on $\mathrm{Hper}=\mathrm{HH}\left(\mathrm{et},1;-(\pi,\pi,\pi)\right)$, which is unitary by (Lz2).
The multiplication by a matrix-valued function of $g$ has as its norm the maximum over $g$ of the matrix norm $|\cdot|_2$.

(E3.2) Let
\[
(\mathrm{Ok}\,\mathrm{hp})(c,g):=\mathrm{et}^3\cdot\frac{1}{48}\sum_{c'\in\Delta_{[c]}}\sum_{g'}\mathrm{ks}(c,c')\,\mathrm{hp}(c',g'),
\]
where $\mathrm{ks}$ is a complex function of the pairs $(c,c')$ with $c'\in\Delta_{[c]}$ such that $\mathrm{ks}(c+v,c'+v)=\mathrm{ks}(c,c')$ for $v\in\mathbb{Z}^3$.
Then $\mathrm{Ok}$ maps $\mathrm{Hper}$ into itself (B1), and $\|\mathrm{Ok}\|\le\sup|\mathrm{ks}|$, where the supremum extends over these pairs.
Indeed, the Cauchy--Schwarz inequality for the measure $\mathrm{et}^3\cdot\frac{1}{48}\sum$ on $\Delta_{[c]}\times O_{\mathrm{cub}}$, which has mass one, gives
\[
\left|(\mathrm{Ok}\,\mathrm{hp})(c,g)\right|_2\le\sup|\mathrm{ks}|\cdot\left(\mathrm{et}^3\cdot\frac{1}{48}\sum_{c',g'}|\mathrm{hp}(c',g')|_2^2\right)^{\frac12}=\sup|\mathrm{ks}|\cdot\|\mathrm{hp}\|,
\]
because $|\mathrm{hp}|_2$ is $\mathbb{Z}^3$-periodic.

(i) Differentiating the definitions of 4.6, one finds
\begin{gather*}
\left(\mathrm{Dpg}^{(\mathrm{kd})}(p)\mathrm{hp}\right)(c,g)=-\mathrm{et}^{-1}\left(i\,\mathrm{et}\cdot(ge_1)_\nu\right)^{\mathrm{kd}}\mathrm{e}^{i\,\mathrm{et}\cdot p\cdot(ge_1)}\,M_3(g)\,\mathrm{hp}(x^{(3)}),\\
\left(\mathrm{Ppg}^{(\mathrm{kd})}(p)\mathrm{hp}\right)(c,g)=\mathrm{et}^3\cdot\frac{1}{48}\sum_{c'\in\Delta_{[c]}}\sum_{g'}\left(-i(c_\nu-c'_\nu)\right)^{\mathrm{kd}}\mathrm{e}^{-ip\cdot(c-c')}\,\mathrm{hp}(c',g').
\end{gather*}
Here $|(ge_1)_\nu|\le1$, $|M_3(g)|_2\le6$, and
\[
\left|\mathrm{e}^{i\,\mathrm{et}\cdot p\cdot(ge_1)}\right|=\exp\left(-\mathrm{et}\cdot\mathrm{ti}\,(ge_1)_\nu\right)\le\exp(\mathrm{et}\,|\mathrm{ti}|)\le\exp(1/3)\le3/2,
\]
since $(3/2)^3>\mathrm{e}$.
With (E3.1) this gives $\|\mathrm{Dpg}^{(\mathrm{kd})}(p)\|\le\mathrm{et}^{\mathrm{kd}-1}\cdot6\cdot\frac32$.
For $c$, $c'$ in one block one has $|c_\nu-c'_\nu|\le(\mathfrak{N}-1)/\mathfrak{N}<1$ and $|\mathrm{e}^{-ip\cdot(c-c')}|=\exp\left(\mathrm{ti}\,(c_\nu-c'_\nu)\right)\le\exp(|\mathrm{ti}|)$.
With (E3.2), $\|\mathrm{Ppg}^{(\mathrm{kd})}(p)\|\le\exp(|\mathrm{ti}|)\le\mathrm{e}$.
For $c'\in\Delta_{[c]}$ one has $[c']=[c]$, so that $c-c'=(c-[c])-(c'-[c'])$.
By the binomial formula the kernel equals $\sum_{\mathrm{ib}=0}^{\mathrm{kd}}\mathrm{ks1}_{\mathrm{ib}}(c)\,\mathrm{ks2}_{\mathrm{ib}}(c')$ with the $\mathbb{Z}^3$-periodic functions
\[
\mathrm{ks1}_{\mathrm{ib}}(c):=(-i)^{\mathrm{kd}}\binom{\mathrm{kd}}{\mathrm{ib}}\left(c_\nu-[c]_\nu\right)^{\mathrm{ib}}\mathrm{e}^{-ip\cdot(c-[c])},\qquad
\mathrm{ks2}_{\mathrm{ib}}(c'):=\left(-(c'_\nu-[c']_\nu)\right)^{\mathrm{kd}-\mathrm{ib}}\mathrm{e}^{ip\cdot(c'-[c'])}.
\]
The operator with the kernel $\mathrm{ks1}_{\mathrm{ib}}(c)\,\mathrm{ks2}_{\mathrm{ib}}(c')$ maps $\mathrm{hp}\in\mathrm{Hper}$ to the function $(c,g)\mapsto\mathrm{ks1}_{\mathrm{ib}}(c)\,v$, where
\[
v:=\mathrm{et}^3\cdot\frac{1}{48}\sum_{c'\in\Delta_0}\sum_{g'}\mathrm{ks2}_{\mathrm{ib}}(c')\,\mathrm{hp}(c',g')\in\mathbb{C}^4;
\]
here the sum over $\Delta_{[c]}$ has been replaced by the sum over $\Delta_0$, the summand being $\mathbb{Z}^3$-periodic.
The rank of this operator is at most $4$.
So $\operatorname{rank}\mathrm{Ppg}^{(\mathrm{kd})}(p)\le4(\mathrm{kd}+1)$, and (H5)\,(ii) gives the bound on the trace norm.

(ii) One has
\[
\left(\left(\mathrm{Dpg}(p)-\mathrm{Dpg}(p')\right)\mathrm{hp}\right)(c,g)=-\mathrm{et}^{-1}\,\mathrm{e}^{i\,\mathrm{et}\cdot p'\cdot(ge_1)}\left(\exp\left(-\mathrm{et}\cdot\mathrm{ti}\,(ge_1)_\nu\right)-1\right)M_3(g)\,\mathrm{hp}(x^{(3)}),
\]
and
\[
\left|\exp\left(-\mathrm{et}\cdot\mathrm{ti}\,(ge_1)_\nu\right)-1\right|\le\exp(\mathrm{et}\,|\mathrm{ti}|)-1\le\mathrm{et}\,|\mathrm{ti}|\exp(|\mathrm{ti}|),
\]
because $(ge_1)_\nu$ is $0$, $1$ or $-1$ and $\exp(\mathrm{xr})-1\le\mathrm{xr}\exp(\mathrm{xr})$ for $\mathrm{xr}\ge0$.
With (E3.1) the norm of this difference is at most $6|\mathrm{ti}|\exp(|\mathrm{ti}|)$.
The operator $\mathrm{Ppg}(p)-\mathrm{Ppg}(p')$ has the kernel $\mathrm{e}^{-ip'\cdot(c-c')}\left(\exp\left(\mathrm{ti}\,(c_\nu-c'_\nu)\right)-1\right)$, whose modulus is at most $\exp(|\mathrm{ti}|)-1\le|\mathrm{ti}|\exp(|\mathrm{ti}|)$.
With (E3.2),
\[
\|\mathrm{Epg}\|\le(6+\mathsf{b})\,|\mathrm{ti}|\exp(|\mathrm{ti}|)\le(6+b_+)\cdot\gamma_{\Sigma}\cdot\mathrm{e},
\]
and by Lemma E2
\[
\left\|\mathrm{Epg}\,\mathrm{Gpg}(p')\right\|\le(6+b_+)\cdot\mathrm{e}\cdot\gamma_{\Sigma}\cdot C_{\mathrm{iv}}\le\frac12.
\]
Now $\mathrm{Tpg}(p)=\left(1+\mathrm{Epg}\,\mathrm{Gpg}(p')\right)\mathrm{Tpg}(p')$, and the Neumann series gives $\left\|\left(1+\mathrm{Epg}\,\mathrm{Gpg}(p')\right)^{-1}\right\|\le2$; the remaining assertions follow with (H5)\,(i).
\end{proof}

\begin{named}{Lemma E4}[singular values]
For $p$ as in Lemma E3\,(ii) and $\mathrm{sr}\in\{1,2,3,4\}$,
\[
\mathrm{et}^{4-\mathrm{sr}}\,\|\mathrm{Gpg}(p)\|_{(\mathrm{sr})}^{\mathrm{sr}}\le192\cdot C_{\mathrm{sv}}^{\mathrm{sr}}\cdot C_{\mathrm{uv}}^{4-\mathrm{sr}}\cdot S_{\mathrm{jf}}.
\]
\end{named}

\begin{proof}
By (G4) the operators at $p+2\pi e_\mu$ and at $p$ are unitarily equivalent; one may therefore assume $p'\in[-\pi,\pi]^3$.
By (G3) and (Lz3), $\operatorname{Re}\mathrm{Dpg}(p')\ge0$; so $\mathrm{Gon}:=\left(\mathrm{Dpg}(p')+1\right)^{-1}$ exists (H1).
The operator
\[
\mathrm{Won}:=\mathrm{Tpg}(p')-\mathrm{Dpg}(p')-1=(\mu-1)+\mathsf{b}\,\mathrm{Ppg}(p')
\]
satisfies $\|\mathrm{Won}\|\le1+b_+$, since $0<\mu\le1$ and $\mathrm{Ppg}(p')$ is a projection (G3).
From $\mathrm{Gon}^{-1}=\mathrm{Tpg}(p')-\mathrm{Won}$ one gets $\mathrm{Gon}^{-1}\,\mathrm{Gpg}(p')=1-\mathrm{Won}\,\mathrm{Gpg}(p')$, that is, $\mathrm{Gpg}(p')=\mathrm{Gon}\left(1-\mathrm{Won}\,\mathrm{Gpg}(p')\right)$.
By (H5)\,(i), Lemma E2 and Lemma E3\,(ii),
\[
\mathrm{sgv}_j\left(\mathrm{Gpg}(p)\right)\le2\left(1+(1+b_+)\cdot C_{\mathrm{iv}}\right)\mathrm{sgv}_j(\mathrm{Gon})=C_{\mathrm{sv}}\,\mathrm{sgv}_j(\mathrm{Gon}).
\]
By (G2), $\mathrm{Gon}$ is unitarily equivalent to $\bigoplus_{\mathrm{jf}}(X_{\mathrm{jf}}+1)^{-1}$, where $X_{\mathrm{jf}}=\mathrm{et}^{-1}\cdot D^{(\mathrm{et}\cdot p'_{\mathrm{jf}};\,0)}$ is now formed with $p'$.
By Lemma E1\,(iv), each of the $192$ singular values of $(X_{\mathrm{jf}}+1)^{-1}$ is at most $\left(1+c_{\mathrm{N3}}^2|p'_{\mathrm{jf}}|_2^2\right)^{-\frac12}$.
With (H5)\,(iv),
\[
\|\mathrm{Gpg}(p)\|_{(\mathrm{sr})}^{\mathrm{sr}}\le C_{\mathrm{sv}}^{\mathrm{sr}}\cdot192\sum_{\mathrm{jf}\in\mathrm{Mf}}\left(1+c_{\mathrm{N3}}^2|p'_{\mathrm{jf}}|_2^2\right)^{-\frac{\mathrm{sr}}{2}}.
\]
By Lemma E1\,(i), $|\mathrm{et}\cdot p'_{\mathrm{jf}}|_2\le\sqrt3\,\pi$; therefore $\mathrm{et}^2\left(1+c_{\mathrm{N3}}^2|p'_{\mathrm{jf}}|_2^2\right)\le1+3\pi^2c_{\mathrm{N3}}^2=C_{\mathrm{uv}}^2$, that is,
\[
\mathrm{et}^{4-\mathrm{sr}}\le C_{\mathrm{uv}}^{4-\mathrm{sr}}\left(1+c_{\mathrm{N3}}^2|p'_{\mathrm{jf}}|_2^2\right)^{-\frac{4-\mathrm{sr}}{2}}.
\]
Hence
\[
\mathrm{et}^{4-\mathrm{sr}}\,\|\mathrm{Gpg}(p)\|_{(\mathrm{sr})}^{\mathrm{sr}}\le192\cdot C_{\mathrm{sv}}^{\mathrm{sr}}\cdot C_{\mathrm{uv}}^{4-\mathrm{sr}}\sum_{\mathrm{jf}}\left(1+c_{\mathrm{N3}}^2|p'_{\mathrm{jf}}|_2^2\right)^{-2},
\]
and by Lemma E1\,(ii) and (F4)
\[
\sum_{\mathrm{jf}}\left(1+c_{\mathrm{N3}}^2|p'_{\mathrm{jf}}|_2^2\right)^{-2}\le1+\sum_{\mathrm{jf}\neq0}\left(c_{\mathrm{N3}}\pi|\mathrm{jf}|\right)^{-4}\le S_{\mathrm{jf}}.
\]
\end{proof}

\begin{named}{Lemma E5}[the fourth derivative]
Fix $\nu$.
Wherever $\mathrm{Tpg}(p)$ is invertible, put $\mathrm{Gq}:=\mathrm{Gpg}(p)$, $\mathrm{tq}_{\mathrm{kd}}:=\mathrm{Tpg}^{(\mathrm{kd})}(p)$, and
\begin{multline*}
\mathrm{Fl4}(p):=\operatorname{Tr}(\mathrm{Gq}\,\mathrm{tq}_4)-4\operatorname{Tr}(\mathrm{Gq}\,\mathrm{tq}_1\,\mathrm{Gq}\,\mathrm{tq}_3)-3\operatorname{Tr}(\mathrm{Gq}\,\mathrm{tq}_2\,\mathrm{Gq}\,\mathrm{tq}_2)\\
+12\operatorname{Tr}(\mathrm{Gq}\,\mathrm{tq}_1\,\mathrm{Gq}\,\mathrm{tq}_1\,\mathrm{Gq}\,\mathrm{tq}_2)-6\operatorname{Tr}\left((\mathrm{Gq}\,\mathrm{tq}_1)^4\right).
\end{multline*}
\begin{itemize}
\item[(i)] For real $p$ one has $\partial_\nu^4\mathrm{Flg}(p)=\mathrm{Fl4}(p)$.
\item[(ii)] For $p=p'+i\,\mathrm{ti}\,e_\nu$ with $p'\in\mathbb{R}^3$ and $|\mathrm{ti}|\le\gamma_{\Sigma}$ one has $|\mathrm{Fl4}(p)|\le C_{\mathrm{Flg}}$.
\end{itemize}
\end{named}

\begin{proof}
(i) By (G2), $\mathrm{Flg}=\log\det\mathrm{Tpg}(p)$, with $\det\mathrm{Tpg}(p)>0$.
Write $\mathrm{pd}$ for $\partial_\nu$.
By (H3), $\mathrm{pd}\,\mathrm{Flg}=\operatorname{Tr}(\mathrm{Gq}\,\mathrm{tq}_1)$ and $\mathrm{pd}\,\mathrm{Gq}=-\mathrm{Gq}\,\mathrm{tq}_1\,\mathrm{Gq}$, so that
\[
\mathrm{pd}\left(\mathrm{Gq}\,\mathrm{tq}_{\mathrm{kd}}\right)=-\mathrm{Gq}\,\mathrm{tq}_1\,\mathrm{Gq}\,\mathrm{tq}_{\mathrm{kd}}+\mathrm{Gq}\,\mathrm{tq}_{\mathrm{kd}+1}.
\]
Using the cyclicity of the trace one obtains
\begin{align*}
\mathrm{pd}^2\,\mathrm{Flg}&=\operatorname{Tr}(\mathrm{Gq}\,\mathrm{tq}_2)-\operatorname{Tr}\left((\mathrm{Gq}\,\mathrm{tq}_1)^2\right),\\
\mathrm{pd}^3\,\mathrm{Flg}&=\operatorname{Tr}(\mathrm{Gq}\,\mathrm{tq}_3)-3\operatorname{Tr}(\mathrm{Gq}\,\mathrm{tq}_1\,\mathrm{Gq}\,\mathrm{tq}_2)+2\operatorname{Tr}\left((\mathrm{Gq}\,\mathrm{tq}_1)^3\right),\\
\mathrm{pd}\operatorname{Tr}(\mathrm{Gq}\,\mathrm{tq}_3)&=\operatorname{Tr}(\mathrm{Gq}\,\mathrm{tq}_4)-\operatorname{Tr}(\mathrm{Gq}\,\mathrm{tq}_1\,\mathrm{Gq}\,\mathrm{tq}_3),\\
\mathrm{pd}\operatorname{Tr}(\mathrm{Gq}\,\mathrm{tq}_1\,\mathrm{Gq}\,\mathrm{tq}_2)&=-2\operatorname{Tr}\left((\mathrm{Gq}\,\mathrm{tq}_1)^2\,\mathrm{Gq}\,\mathrm{tq}_2\right)+\operatorname{Tr}\left((\mathrm{Gq}\,\mathrm{tq}_2)^2\right)+\operatorname{Tr}(\mathrm{Gq}\,\mathrm{tq}_1\,\mathrm{Gq}\,\mathrm{tq}_3),\\
\mathrm{pd}\operatorname{Tr}\left((\mathrm{Gq}\,\mathrm{tq}_1)^3\right)&=3\operatorname{Tr}\left((\mathrm{Gq}\,\mathrm{tq}_1)^2\,\mathrm{Gq}\,\mathrm{tq}_2\right)-3\operatorname{Tr}\left((\mathrm{Gq}\,\mathrm{tq}_1)^4\right).
\end{align*}
Thus $\mathrm{pd}^4\,\mathrm{Flg}$ is the first of the last three right-hand sides, minus three times the second, plus twice the third; this is $\mathrm{Fl4}$.
The moduli of the five coefficients of $\mathrm{Fl4}$ add up to $1+4+3+12+6=26$.

(ii) Each of the five terms has the form $\operatorname{Tr}\prod_{i=1}^{\mathrm{sr}}\left(\mathrm{Gq}\,\mathrm{tq}_{\mathrm{kd}_i}\right)$ with $1\le\mathrm{sr}\le4$, $\mathrm{kd}_i\ge1$ and $\mathrm{kd}_1+\dots+\mathrm{kd}_{\mathrm{sr}}=4$.
Insert $\mathrm{tq}_{\mathrm{kd}}=\mathrm{Dpg}^{(\mathrm{kd})}+\mathsf{b}\,\mathrm{Ppg}^{(\mathrm{kd})}$ and expand into $2^{\mathrm{sr}}$ terms.
For the term without a factor $\mathrm{Ppg}$, (H5)\,(ii), (iii), (i), Lemma E3\,(i) (note that $\gamma_{\Sigma}\le1$) and Lemma E4 give
\begin{align*}
\left|\operatorname{Tr}\prod_i\left(\mathrm{Gq}\,\mathrm{Dpg}^{(\mathrm{kd}_i)}\right)\right|
&\le\prod_i\left\|\mathrm{Gq}\,\mathrm{Dpg}^{(\mathrm{kd}_i)}\right\|_{(\mathrm{sr})}
\le\|\mathrm{Gq}\|_{(\mathrm{sr})}^{\mathrm{sr}}\prod_i9\,\mathrm{et}^{\mathrm{kd}_i-1}
=9^{\mathrm{sr}}\,\mathrm{et}^{4-\mathrm{sr}}\,\|\mathrm{Gq}\|_{(\mathrm{sr})}^{\mathrm{sr}}\\
&\le9^{\mathrm{sr}}\cdot192\cdot C_{\mathrm{sv}}^{\mathrm{sr}}\cdot C_{\mathrm{uv}}^{4-\mathrm{sr}}\cdot S_{\mathrm{jf}}
\le9^4\cdot192\cdot C_{\mathrm{sv}}^4\cdot C_{\mathrm{uv}}^3\cdot S_{\mathrm{jf}},
\end{align*}
where the last step uses $C_{\mathrm{sv}}\ge1$ and $C_{\mathrm{uv}}\ge1$.
Each of the other $2^{\mathrm{sr}}-1\le15$ terms contains a factor $\mathsf{b}\,\mathrm{Ppg}^{(\mathrm{kd})}$ with $\mathrm{kd}\le4$.
Estimate the trace norm of one such factor by $b_+\cdot20\,\mathrm{e}$ (Lemma E3\,(i)), and the operator norms of all the other factors:
the $\mathrm{sr}\le4$ factors $\mathrm{Gq}$ have norm at most $2C_{\mathrm{iv}}$ (Lemma E3\,(ii); $2C_{\mathrm{iv}}\ge1$), and the $\mathrm{sr}-1\le3$ remaining factors $\mathrm{Dpg}^{(\mathrm{kd})}$, $\mathsf{b}\,\mathrm{Ppg}^{(\mathrm{kd})}$ have norm at most $\max(9,\mathrm{e}\cdot b_+)$ (as $\mathrm{et}\le1$).
By (H5)\,(ii) the modulus of such a term is at most $20\cdot\mathrm{e}\cdot b_+\cdot(2C_{\mathrm{iv}})^4\cdot\max(9,\mathrm{e}\cdot b_+)^3$.
Consequently each of the five traces has modulus at most
\[
9^4\cdot192\cdot C_{\mathrm{sv}}^4\cdot C_{\mathrm{uv}}^3\cdot S_{\mathrm{jf}}+15\cdot20\cdot\mathrm{e}\cdot b_+\cdot(2C_{\mathrm{iv}})^4\cdot\max(9,\mathrm{e}\cdot b_+)^3;
\]
multiplying by the moduli of the coefficients and adding, one arrives at the bound $C_{\mathrm{Flg}}$.
\end{proof}

\begin{named}{Lemma E6}[Fourier coefficients]
For $z\in\mathbb{Z}^3$, $z\neq0$,
\[
|\mathrm{Flgh}(z)|\le C_{\mathrm{Flg}}\,|z|^{-4}\exp(-\gamma_{\Sigma}|z|).
\]
\end{named}

\begin{proof}
Choose $\nu$ with $|z_\nu|=|z|$.
By (F5) with $\mathrm{kd}=4$ and Lemma E5\,(i),
\[
\mathrm{Flgh}(z)=(iz_\nu)^{-4}\int_{[-\pi,\pi]^3}\frac{\mathrm{d}^3p}{(2\pi)^3}\,\mathrm{Fl4}(p)\,\mathrm{e}^{-ip\cdot z}.
\]
The entries of $\mathrm{Tpg}(p)$ and of the $\mathrm{Tpg}^{(\mathrm{kd})}(p)$ in a fixed basis of $\mathrm{Hper}$ are entire functions of $p$.
By Lemma E3\,(ii), $\mathrm{Tpg}(p)$ is invertible on the set $\left\{p:\ p_\mu\ \text{real for}\ \mu\neq\nu,\ |\operatorname{Im}p_\nu|\le\gamma_{\Sigma}\right\}$, and the set on which it is invertible is open.
Hence $\mathrm{Fl4}$ is continuous on this set, and holomorphic in $p_\nu$ in a neighbourhood of the closed strip $|\operatorname{Im}p_\nu|\le\gamma_{\Sigma}$.
Differentiating (G4) one obtains $\mathrm{Tpg}^{(\mathrm{kd})}(p+2\pi e_\mu)=\mathrm{Mper}_\mu^{-1}\,\mathrm{Tpg}^{(\mathrm{kd})}(p)\,\mathrm{Mper}_\mu$, and also $\mathrm{Gpg}(p+2\pi e_\mu)=\mathrm{Mper}_\mu^{-1}\,\mathrm{Gpg}(p)\,\mathrm{Mper}_\mu$; so $\mathrm{Fl4}(p+2\pi e_\mu)=\mathrm{Fl4}(p)$.
Now (F6) with $\mathrm{fhol}=\mathrm{Fl4}$ and $\mathrm{wd}=\gamma_{\Sigma}$, together with Lemma E5\,(ii), gives
\[
\left|\int_{[-\pi,\pi]^3}\frac{\mathrm{d}^3p}{(2\pi)^3}\,\mathrm{Fl4}(p)\,\mathrm{e}^{-ip\cdot z}\right|\le\exp(-\gamma_{\Sigma}|z_\nu|)\,C_{\mathrm{Flg}}.
\]
\end{proof}

\emph{The first inequality of (8.5).}
By (4.8), $|\mathrm{e}^{i\theta\cdot w}-1|\le2$, Lemma E6 at $z=\ell w$ (where $|\ell w|=\ell|w|\ge\ell$) and (F4),
\begin{align*}
\left|\log\frac{\det T_\theta}{\det T_{\theta=0}}\right|
&\le2\ell^3\sum_{w\neq0}C_{\mathrm{Flg}}\,(\ell|w|)^{-4}\exp(-\gamma_{\Sigma}\ell|w|)\\
&\le2C_{\mathrm{Flg}}\,\ell^{-1}\exp(-\gamma_{\Sigma}\ell)\sum_{w\neq0}|w|^{-4}
=C_{\Sigma}\,\ell^{-1}\exp(-\gamma_{\Sigma}\ell).
\end{align*}

\emph{The second inequality of (8.5).}
All determinants in (d) are positive: $\det T_\theta$ and $\det T_{\theta=0}$ by clause (a), and $\det(D_{H_\theta}+\mathrm{et}\cdot\mu)$, $\det(D_0+\mathrm{et}\cdot\mu)$ by Proposition~30.1\,(b) (and also by clause (a) together with the two determinant identities of 4.4, which express $\det(D_{H_\theta}+\mathrm{et}\cdot\mu)$ and $\det(D_0+\mathrm{et}\cdot\mu)$ as positive multiples of $\det(D+\mu)$ on $\mathscr{H}_\theta$ and on $\mathscr{H}_{\theta=0}$).
By (d),
\[
\log\mathrm{Rtw}_\theta=\log\frac{\det(D_{H_\theta}+\mathrm{et}\cdot\mu)}{\det(D_0+\mathrm{et}\cdot\mu)}-\log\frac{\det T_\theta}{\det T_{\theta=0}}.
\]
By Proposition~30.2\,(a) with $n=\ell\mathfrak{N}$ and $\mathbf{m}=\mathrm{et}\cdot\mu\in(0,1]$, the first term has modulus at most $2\varepsilon_5(\mathrm{et}\cdot\mu\cdot\ell\mathfrak{N})=2\varepsilon_5(\mu\ell)$; the second term is bounded by the first inequality.

\emph{Constants, uniformity.}
By their formulas, $0<\gamma_{\Sigma}\le1$ and $C_{\Sigma}<\infty$, because $b_->0$, $\hat{g}>0$ and $c_{\mathrm{N3}}>0$.
These formulas contain $b_-$, $b_+$, $\rho$ and the number $\lambda_{\mathrm{int}}$, and nothing else.
The side $\ell$ enters only in (4.8), and $\theta$ only through the factors $\mathrm{e}^{i\theta\cdot w}-1$.
The mesh, that is $\mathfrak{N}$, enters only through the following facts: $\mathfrak{N}$ is odd; $\mathrm{et}\le1/3$; Lemma E1\,(i), (ii); the bound $u_1(t)\ge2/\pi$ on $[-\pi,\pi]$ (Lemma E2); the bound $|c_\nu-c'_\nu|\le(\mathfrak{N}-1)/\mathfrak{N}<1$ for $c$, $c'$ in one block (Lemma E3); and the convergent sum $S_{\mathrm{jf}}$.
The mass $\mu$ enters only through $0<\mu\le1$; in Lemma E2 only $\mu\ge0$ is needed, the mode at $p=0$ being bounded from below through $\mathsf{b}\,\mathrm{Pbl}$.
The number $\mathsf{b}$ enters only through $b_-\le\mathsf{b}\le b_+$, and $\mathrm{ec}$ occurs only in clause (f), where it drops out.
This completes the proof of Proposition~30.3.

\subsection*{5. Levels; the order of the proofs}

\paragraph{5.1 Levels}
Propositions~30.1 and 30.2 are statements at level $0$, in lattice units ($\eta_0=1$).
No block average occurs in them; the determinant is that of a linear map of a space of dimension $192n^3$ ($48$ sites with $4$ components over each of the $n^3$ photon sites), and the weight $1/48$ enters only the scalar products, on which no determinant depends.
Proposition~30.3 concerns a fine lattice of mesh $\mathrm{et}=1/\mathfrak{N}\le1/3$ together with its unit lattice; in clause (f) these are the levels $k\ge1$.
The fine space has dimension $192(\ell\mathfrak{N})^3$, the unit space dimension $4\ell^3$; the average $Q$ extends over the $\mathfrak{N}^3$ photon sites of a block and over the $48$ sites of each cluster, with the weight $\mathrm{et}^3/48$.
At level $0$, by contrast, the fine lattice and the unit lattice coincide ($\eta_0=1$ and $Q_0=P_0=1$), the objects $b_0$ and $S_k(A)=S_\infty(A)$ at $k=0$ are not defined, and $D_0(A)=D_A+\mu$ contains no block average (Section~I.3.9; see also the Remark on level~$0$ at the end of 4.5).
Since Proposition~30.3 requires an odd $\mathfrak{N}\ge3$, that is, a mesh $\mathrm{et}\le1/3$, the operators of level~$0$ of Section~I.3.9 are not instances of the definitions of Proposition~30.3.

\paragraph{5.2 The order of the proofs}
Propositions~30.1 and~30.2 are proved in Sections~2 and~3 without recourse to Proposition~30.3, the clauses of Proposition~30.1 in the order (a), (b), (d), (c).
The arrows of the following scheme give the order of dependence among the proofs of the assertions of Proposition~30.3:
the proof of an assertion in the scheme uses, among the assertions of Proposition~30.3, only those from which a succession of arrows leads to it, and not necessarily the one immediately before it.
\begin{gather*}
\text{(a)}\longrightarrow\text{(b)}\longrightarrow\text{(c)}\longrightarrow\text{(d)};\\
\text{(a), (b)}\longrightarrow\text{(g1)}\longrightarrow\text{(g2)}\longrightarrow\text{(g3)}\longrightarrow\text{(e), first inequality};\\
\left.\begin{array}{l}\text{(d)}\\ \text{(e), first inequality}\\ \text{Proposition~30.2\,(a)}\end{array}\right\}\longrightarrow\text{(e), second inequality}.
\end{gather*}
Besides, the proof of (g2) uses Proposition~30.1\,(a), the proof of (g3) uses Proposition~30.1\,(d), and in the proof of the second inequality of (e) the positivity of the four determinants comes from (a) and Proposition~30.1\,(b).
Clause (f) rests on (a) and on (c), the latter through $\det\mathrm{Dfl}_{\theta=0}>0$;
its last part, ``In the flow'', rests in addition on (Sc1), (Sc3), (Sc4) and (Sc6) of Section~G4.0.2.
The proofs of all the other clauses of Propositions~30.1, 30.2 and~30.3, among them those of Proposition~30.1\,(c) and Proposition~30.3\,(e), are independent of (f) (Section~IV.1\,(B0), (B5$'$)).

\subsection*{6. Conditions on the constants}
The proofs of Sections~1--4 require only the following conditions on the constants.
For each condition the places where it is used are indicated, together with the line of constants of Section~I.8.2, or of Proposition~14.2, in which it stands.
\begin{itemize}
\item
\emph{Condition:} $0<\rho<1$, used in the forms $\rho\ge0$ (in $\operatorname{Re}D=\frac{\rho}{2\,\mathrm{me}}\,\nabla^{*}\nabla\ge0$), $\rho>0$ (in $\hat{g}>0$ and $c_{\mathrm{N3}}>0$) and $\rho\le1$ (in $|M_3|_2=3(1+\rho)\le6$, in $\cos^2+\rho^2\sin^2\ge\rho^2$, and in $2\rho/\pi\le1$).
\emph{Used in:} (Lz3), (Lz4), (Lz6), (Lz7), (c2), Lemmas E1 and E3; in the form $\rho>0$ also in Lemmas E2 and E4 and in the paragraphs ``Constants, uniformity'' at the end of Section~2 and of 4.9.
\emph{Line:} (C0).
\item
\emph{Condition:} the standard weights, that is, $w_i>0$, $\sum_iw_i\hat n_i=0$ ((I.1.1)) and $w_3=6$.
\emph{Used in:} $w_i>0$ in (Lz3), (Lz4); the other two in (Lz1)--(Lz3) and in Lemma E1, steps (2) and (3); $w_3=6$, which gives $|M_3|_2=\frac{w_3}{2}(1+\rho)=3(1+\rho)$ by (I.m6), also in (c2) and in Lemma E3\,(i), (ii).
\emph{Line:} (C0).
\item
\emph{Condition:} $\lambda_{\mathrm{int}}>0$ and $\hat{g}=\rho\lambda_{\mathrm{int}}/2$.
\emph{Used in:} (Lz4); the identity $\hat{g}=\rho\lambda_{\mathrm{int}}/2$ also in Lemma E1, step (4); $\hat{g}>0$ also in (c2) and in Lemma E2.
\emph{Line:} (C0r).
\item
\emph{Condition:} $\theta_0>0$ with $3(1+\rho)\left(\exp(\theta_0)-1\right)=\hat{g}/2$.
\emph{Used in:} the identity in (c2); $\theta_0>0$ in (c2) (in (F6) with $\mathrm{wd}=\theta_0$), (c3) and (c4).
\emph{Line:} (C0r).
\item
\emph{Condition:} $K_{\mathrm{Str}}<\infty$.
\emph{Used in:} (c3).
\emph{Line:} (C0r); Proposition~1.2\,(S3).
\item
\emph{Condition:} $0<\kappa_{(0)}\le\theta_0/2$.
\emph{Used in:} (c3), (c4).
\emph{Line:} (C0r) (the definition of $\kappa_{(0)}$).
\item
\emph{Condition:} $A'_{\mathrm{ak}}<\infty$ for $|\mathrm{ak}|\le3$.
\emph{Used in:} (c3).
\emph{Line:} (C0r), row $A_{\mathrm{ak}},A'_{\mathrm{ak}}$; Proposition~1.2\,(S7).
\item
\emph{Condition:} $C_6$ is given by the formula in Proposition~30.1 (a definition).
\emph{Used in:} (c0), (c4).
\emph{Line:} (C0r), row $C_6$.
\item
\emph{Condition:} $c_{\mathrm{N3}}=\min\left\{\rho/\pi,\ \rho^3\lambda_{\mathrm{int}}/(98\sqrt3\,\pi^3)\right\}$ (a definition).
\emph{Used in:} the definition in Lemma E1, step (4); $c_{\mathrm{N3}}>0$ in Lemmas E2 and E4.
\emph{Line:} (C0r), row $c_{\mathrm{N3}}$.
\item
\emph{Condition:} $0<b_-\le b_+$.
\emph{Used in:} 4.9.
\emph{Line:} (C4).
\item
\emph{Condition:} $\gamma_{\Sigma}$ and $C_{\Sigma}$ are given by the formulas in Proposition~30.3 (definitions).
\emph{Used in:} 4.9.
\emph{Line:} (C6), rows $\gamma_{\Sigma}$, $C_{\Sigma}$.
\item
\emph{Condition:} $L$ odd and $L\ge5$, used in the form ``$L^k$ is an odd integer $\ge3$ for $k\ge1$'' and in the form $L\ge1$.
\emph{Used in:} 4.5; the form $L\ge1$, which gives $\eta=L^{-N}\le1$, in the proof of Proposition~30.2\,(b) in Section~3.
\emph{Line:} Section~I.0.5; (C3).
\item
\emph{Condition:} $e\,h^{\frac12}\le e_1$ (a datum of (FL)), where $e_1$ is the number of the line (C9), with (e1): $v\ge8$, $v^r>31L+16$ and $\left(1+3\,\mathrm{lg}/v\right)^r<L$ for all $v\ge-\log e_1$; further $L\ge33$ (hence $L\ge13$), and $M$ is a power of $L$.
These conditions are used through (Sc1), (Sc3), (Sc4) and (Sc6) of Section~G4.0.2, as follows.
The condition (e1) yields that part of (Sc6) which 4.5 uses, in three steps.
(i) Since $v_K=-\log\left(e\,h^{\frac12}\right)\ge-\log e_1$, (e1) at $v=v_K$ gives $v_K\ge8$ and $\left(1+3\,\mathrm{lg}/v_K\right)^r<L$; as $\mathrm{lg}>0$ and $r>0$, it follows that $v_K\ge1$ and $\left(1+\mathrm{lg}/v_K\right)^r\le L$, which are the hypotheses of \cite[Proposition~19.1]{QEDI}.
(ii) The clauses (c) and (e) of \cite[Proposition~19.1]{QEDI}, that is, (Sc1) and (Sc3), give $v_*\ge v_K\ge-\log e_1$.
(iii) Hence the second inequality of (e1) holds at $v=v_*$; together with $r_*\ge v_*^r$, which follows from the definition of $r_k$ in Section~I.6.1, it gives $r_*>31L+16$.
This is the condition (G4.2.1) of Section~G4.0.3, the hypothesis of \cite[Proposition~19.1(f)]{QEDI}, that is, of (Sc4).
The remaining assertions of (Sc6), namely $e\le1/3$ and the conditions of Section~G4.0.3 other than (G4.2.1), are not used in this chapter.
The proof of \cite[Proposition~19.1(f)]{QEDI} uses $L\ge33$, $L\ge13$ and that $M$ is a power of $L$.
\emph{Used in:} 4.5, paragraph ``In the flow'', which proves the last part of (f).
\emph{Line:} (C9), Proposition~14.2\,(b), row (e1); (C3) (Proposition~14.2\,(f): $L\ge2400^8$); (C5), Proposition~14.2\,(b), row (M1).
\end{itemize}
Propositions~30.1 and 30.2, the clauses (a)--(e) and (g) of Proposition~30.3, and the part of clause (f) that precedes ``In the flow'' hold without any condition on $\kappa$, $M$, $\varepsilon$, $e_0$, $e_1$;
conditions on these numbers enter only the last part of (f), through the last entry above.
The constants $C_6$, $\kappa_{(0)}$, $c_{\mathrm{N3}}$ are formed with $\rho$ (and the weights) only, and $\gamma_{\Sigma}$, $C_{\Sigma}$ with $(b_-,b_+,\rho)$ and the weights only;
so $C_6$, $\kappa_{(0)}$, $\gamma_{\Sigma}$, $C_{\Sigma}$ can be fixed before $e_0$ and $e_1$.
The number $e_1$ depends on $\gamma_{\Sigma}$ and $C_{\Sigma}$ through (e6), and among the assertions of Propositions~30.1--30.3 it occurs only in the hypothesis of the last part of (f), ``In the flow''.
Consequently the constants can be chosen in the order
\[
\text{(C0r), (C4), (C6)}\ \longrightarrow\ e_0\ \text{(C8)}\ \longrightarrow\ e_1\ \text{(C9)}\ \longrightarrow\ \text{(FL)}\ \longrightarrow\ \text{the part ``In the flow'' of (f)}.
\]

\section*{Chapter 14a. Assembly of the bounds on the fermion operators: proof of Proposition~14.1}

\subsection*{1. The statement}

The statement proved in this chapter is Proposition~14.1 (Z2$^{\mathrm{ch}}$).
It is asserted for all data (SF), with the constants of its paragraph ``Constants'', and it has five parts:
the properties FE1$^{\mathrm{ch}}$--FE6$^{\mathrm{ch}}$ and FE8$^{\mathrm{ch}}$ of the fermion operators (part~(i));
the properties W1, W2, W5--W17 (Section~G1.11) of the families of Propositions~7 and~8 for the parameters $\left(\alpha_0,\gamma^F,\kappa_{\mathrm{w}}+\kappa_0;C^F\right)$, $\mathrm{par}^0$ and $\mathrm{par}$ (part~(ii));
the clauses (d0)--(d6) of Proposition~13, which form FE7$^{\mathrm{ch}}$, for the families (p5) with the properties required there (part~(iii));
the independence of the constants of $k$, $\ell$, $\mu$, $e_k$, $A$, $y$, $s$, $X$, and the validity of all assertions at every sum of a connection of $\tilde{R}_k$, a connection of $\mathfrak{X}_k$ and a real constant connection (part~(iv));
and the range of the mass: the assertions of (i)--(iii) hold for every $\mu\in(0,1]$, with constants that do not depend on $\mu$ (part~(v)).

\subsection*{2. Proof}

\subsubsection*{2.0. Conventions of this chapter}

Throughout, $(k,\ell,\mu,e_k)$ are data (SF) (Section~I.1.8, Section~G1.1):
an integer $k\ge0$, a side $\ell\in\mathrm{Cl}$, a mass $\mu\in(0,1]$ and a number $e_k\in(0,e_0]$, with no relation assumed among the four.
The constants are the entries of Proposition~14.2\,(a), which is the list of Section~I.8.2;
they are chosen in the order in which they stand there, subject to the conditions of Proposition~14.2\,(b), and $e_0$ is the single ceiling of the line (C8).
The symbols (C0)--(C8), (C0r), (C4r) denote the lines of Proposition~14.2\,(a),
and (L1)--(L5), (M1)--(M6), (S1), (A1)--(A10), (Z2.s1)--(Z2.s5), (Z3.s1) denote the conditions listed in Proposition~14.2\,(b).
Proposition~2.8 and the Propositions~3--13, whose clauses are applied below at data (SF), hold for all data (SF) with this ceiling;
among them, Propositions~4, 5 and~7 concern $k\ge1$, Proposition~8 concerns $k=0$, Proposition~11 concerns two tori, and Proposition~13 has a family (p5) as an additional datum.

Fermion sites are written $x=(c,g)$ and $x'=(c',g')$, where $c$ is a site of $T_k$ and $g$ is one of the $48$ elements of $O_{\mathrm{cub}}$;
the letters $z,z'$ denote unit sites.
The resolvent parameter is $y\in Y_k$, with $Y_k=\mathbb{R}\cup\{\infty\}$ for $k\ge1$ and $Y_0=\mathbb{R}$.
The letter $\mathrm{cf}$ stands for a number $\ge0$, $\mathrm{vv}$ for a vector of $\mathbb{C}^4$ and $\mathrm{tc}$ for a complex number;
$\mathrm{cst}$ is a real constant bond function, as in 2.6.2.

\emph{Norms} (Section~I.1.5, Section~G1.3).
On $\mathbb{C}^4$ two norms are used,
\[
|v|_c=\max_{\beta\le4}|v_\beta|,\qquad |v|_2=\left(\sum_{\beta\le4}|v_\beta|^2\right)^{\frac{1}{2}},
\]
and for a $4\times4$ matrix $m$ the numbers $|m|_c$ and $|m|_2$ are the operator norms of $m$ with respect to them.
For $f\in\mathfrak{S}^{\mathrm{ch}}(\eta_k,\ell)$, the number $|f|_{\infty}$ is the supremum of $|f(c,g)_\beta|$ over the sites $c$, the $48$ labels $g$ and $\beta\in\{1,\dots,4\}$,
and $|f|_{\mathrm{su}}:=\sup_{(c,g)}|f(c,g)|_2$.
For $v\in\mathfrak{S}(1,\ell)$ the suprema extend over the unit sites, and for a function $F$ of pairs over the pairs.

\emph{Localisation} (Section~G1.3, Section~I.1.4).
The operator $\mathbf{1}_{\Delta_z}$ is the multiplication by the indicator of (the orbit of) $\Delta_z\times O_{\mathrm{cub}}$, a set of $48\eta_k^{-3}$ fermion sites;
on functions of pairs it restricts to the pairs whose first point lies over $\Delta_z$.
For $k=0$ one has $\Delta_z=\{z\}$, and $\Delta_z\times O_{\mathrm{cub}}$ is the cluster of $z$, which consists of $48$ sites.

\emph{Short pairs} (Section~I.1.7) are the pairs with $d_X(x,x')<1$.
For $k\ge1$ the $48\cdot47$ ordered pairs inside each cluster are short, with $d_X=\eta_k$ and $U_A(x,x')=1$.
For $k=0$ no pair is short.

\emph{Kernels} ((I.1.3), (I.1.4)).
For an operator $O$ whose source consists of spinors on fermion sites,
\[
(Of)(t)=\int\mathrm{d}x'\,O(t,x')\cdot f(x')=\frac{\eta_k^3}{48}\sum_{c'}\sum_{g'\in O_{\mathrm{cub}}}O\left(t,(c',g')\right)\cdot f(c',g'),
\]
where $c'$ ranges over a fundamental domain of $\ell\mathbb{Z}^3$.

\subsubsection*{2.1. The constants}

The following constants of Proposition~14.2\,(a) occur in Lemmas~1--3 and in 2.3--2.5;
each of them is defined by its entry there.
\begin{description}
\item[(C0), (C0r):]
$c_X$, $\kappa_0$, $K_0$, $b$ (one has $b=1$, written as a letter; the arguments below require only $b>0$), $w$ (with $w_3=6$);
$\lambda_{\mathrm{int}}$ (positive by (I.O8));
$\hat{g}$, $c_{\mathrm{hl}}$, $c_{*}$, $K_{\rho}$, $c_b$ (formed with the two numbers $A_1$, $A_2$ of its definition), $C_{\mathrm{Po}}$.
\item[(C4):]
$b_-$, $\alpha^P$, $\alpha^F$, $\alpha_0$, $\gamma^P$, $\kappa_{\mathrm{w}}$, $\kappa'$, $C'_V$, $C_T$.
\item[(C4r):]
$K_1$, $C_I$, $c_{\mathrm{wt}}$, $a_0$, $C_5$ (the constant of Proposition~1.6\,(c)), $\gamma_{(5)}$, $c_{\mathrm{re}}$, $a_z$, $k_1$, $C_{00}$, $C_z$, $\gamma_z$, $\gamma^F$,
$a'$ (formed with the function $c(\gamma):=\sum_{n\in\mathbb{Z}^3}\exp\left(-\gamma\cdot|n|\right)$, $\gamma>0$, defined in Proposition~14.2\,(a)),
$\gamma_{\mathrm{w}}$, $M^F$.
\item[(C6):]
$C_{\mathrm{By}}$, $C_{\Gamma}$, $C_1$, $C^F$, $C^P$, $C_w^{(0)}$, $C_w^{(1)}$, $C_w$, $C_d$.
\end{description}
The parameter sets
\[
\mathrm{par}^0=\left(\alpha_0,\gamma_{\mathrm{w}},\kappa_{\mathrm{w}}+\kappa_0;C_w^{(1)}\right),\qquad \mathrm{par}=\left(\alpha_0,\gamma_{\mathrm{w}},\kappa_{\mathrm{w}};C_w\right)
\]
are those defined in Proposition~14.2.

\begin{named}{Lemma 1}
\begin{itemize}
\item[(K1)]
$0<\gamma^F\le\frac{1}{6L}$;
$1\le192L\le C_{00}<\infty$;
$1\le a'<\infty$;
$\alpha_0=\alpha^F=\frac{1}{54}$ and $27\alpha_0=\frac{1}{2}$;
the numbers $C_{\mathrm{By}}$, $C_{\Gamma}$, $C_1$, $C^F$ are finite and positive.
\item[(K2)]
$64a'\le608a'\le3072a'\le C^F$.
\item[(K3)]
$2C_1\le C^F$,
$2C_{\mathrm{By}}\le C^F$,
$2C_{\Gamma}\le C^F$,
$2\mathrm{e}\cdot b^2\cdot(32a')^2\le C^F$,
$64\mathrm{e}\cdot b^2\cdot a'\le2C_1$,
$256\mathrm{e}\cdot(1+2\mathrm{e})^2\cdot a'\le2C_{\Gamma}$,
$\frac{2(1+2\mathrm{e})}{b_-}\le C_{\Gamma}$.
\item[(K4)]
$64\mathrm{e}\cdot b\cdot a'\le256\mathrm{e}\cdot b\cdot a'\le608\mathrm{e}\cdot b\cdot a'\le C^F$;
$4\le2\mathrm{e}\le608\mathrm{e}\le C^F$.
\item[(K5)]
$0<\gamma_{\mathrm{w}}\le\gamma^F$;
$0<\kappa_{\mathrm{w}}\le\kappa_{\mathrm{w}}+\kappa_0=\kappa'$;
$C^F\le C_w^{(0)}\le C_w^{(1)}\le C_w<\infty$;
$C_d<\infty$;
$M^{8\alpha_0}\le4\exp(2\gamma)\cdot M^{27\alpha_0}$ for every $\gamma\ge0$;
$4\exp(2\gamma_{\mathrm{w}})\cdot M^{27\alpha_0}\le C_w^{(1)}$.
\end{itemize}
\end{named}

\begin{proof}
(K1)
By (C0) one has $0<\rho<1$; moreover $\lambda_{\mathrm{int}}>0$, $w_3=6$, $b>0$, and $L\ge5$ by (L5).
Consequently the numbers $\hat{g}$, $c_{\mathrm{hl}}$, $c_{*}$, $K_{\rho}$, $A_1$, $A_2$, taken in this order, are finite and positive,
and therefore so are $c_b$, $C_{\mathrm{Po}}$, $b_-$, $K_1$, $C_I$, $c_{\mathrm{wt}}$.
It follows that
\begin{gather*}
0<a_0\le\frac{1}{3},\qquad 0<\gamma_{(5)}\le\frac{1}{3L},\qquad 0<c_{\mathrm{re}}\le b\left(1-\frac{1}{L}\right),\\
0<a_z\le\frac{1}{3L},\qquad 0<\gamma_z\le\frac{1}{3L},\qquad 0<\gamma^F\le\frac{1}{6L}.
\end{gather*}
As $c_{\mathrm{re}}>0$, the number $C_{00}$ is finite.
From $c_{\mathrm{re}}\le b(L-1)/L$ one obtains
\[
\frac{2L}{c_{\mathrm{re}}}\ge\frac{2L^2}{b(L-1)}\ge\frac{2L}{b},
\]
and $1+2k_1=121+2b\ge2b$; hence
\[
C_{00}\ge48\cdot\frac{2L}{b}\cdot2b=192L\ge1.
\]
The constant $C_5$ is a finite number (Proposition~1.6), so that $1\le C_{00}\le C_z<\infty$.
For an integer $j\ge1$ the number of the points $n\in\mathbb{Z}^3$ with $|n|=j$ equals $(2j+1)^3-(2j-1)^3=24j^2+2$.
Therefore, for $\gamma>0$,
\[
c(\gamma)=1+\sum_{j\ge1}\left(24j^2+2\right)\cdot\exp(-\gamma j),
\]
a convergent series whose value is at least $1$; hence $1\le a'<\infty$.
Further, $\min\left(\frac{1}{54},\frac{1}{4}\right)=\frac{1}{54}$.
The four remaining constants $C_{\mathrm{By}}$, $C_{\Gamma}$, $C_1$, $C^F$ are obtained from finitely many finite positive numbers by sums, products and maxima ($M$ and $\kappa'$ being finite).

(K2)
Since $27\alpha_0=\frac{1}{2}$, $M\ge L^2$ by (M1), and $L\ge5$, one has $M^{27\alpha_0}\ge L\ge5$;
moreover $\mathrm{e}\ge2$ and $\max(1,b)\ge1$.
The number $C_1$ is one of the entries of the maximum that defines $C^F$, and $304\mathrm{e}\cdot\max(1,b)\cdot a'\cdot M^{27\alpha_0}$ is one of the entries of the maximum that defines $C_1$.
Hence
\[
C^F\ge2C_1\ge2\cdot304\cdot2\cdot5\cdot a'=6080a'\ge3072a'.
\]

(K3)
The numbers $C_1$, $C_{\mathrm{By}}$, $C_{\Gamma}$ and $\mathrm{e}\cdot b^2\cdot(32a')^2$ are entries of the maximum that defines $C^F$.
One has $\kappa'=\kappa+2\kappa_0+3>0$, because $\kappa\ge7\kappa_0+7$ by (C2) and $\kappa_0=54\log72>0$;
thus $\exp(27\kappa')\ge1$, and
\[
32\mathrm{e}\cdot b^2\cdot a'\le32\mathrm{e}\cdot\max\left(1,b^2\right)\cdot a'\cdot\exp\left(27\kappa'\right)\le C_1.
\]
Finally, $\frac{2(1+2\mathrm{e})}{b_-}$ and $128\mathrm{e}\cdot(1+2\mathrm{e})^2\cdot a'$ are the two summands of $C_{\Gamma}$, and both are positive.

(K4)
As $b\le\max(1,b)$, one has $608\mathrm{e}\cdot b\cdot a'\le2\cdot304\mathrm{e}\cdot\max(1,b)\cdot a'\le C^F$.
Since $a'\ge1$ and $\max(1,b)\ge1$, it follows that $C^F\ge608\mathrm{e}$; and $2\mathrm{e}\ge4$.

(K5)
The number $\gamma_P$ lies in $(0,1]$ \cite[Section~15.7]{QEDI};
hence $\gamma^P>0$, and $\gamma_{\mathrm{w}}$ is the smaller of two positive numbers, one of which is $\gamma^F$.
By (C2) one has $\kappa\ge7\kappa_0+7$, and $\kappa_0=54\log72>0$;
so $\kappa_{\mathrm{w}}=\kappa+\kappa_0+3>0$ and $\kappa_{\mathrm{w}}\le\kappa_{\mathrm{w}}+\kappa_0$.
The number $C_w^{(0)}$ is a maximum one of whose entries is $C^F$, and $C_w^{(1)}$ is a maximum among whose entries are $C_w^{(0)}$ and $4\exp(2\gamma_{\mathrm{w}})\cdot M^{27\alpha_0}$.
Since $\exp(\kappa_0)=72^{54}\ge36$, one has $K_0\ge1$ and therefore $64K_0\ge1$, so that $C_w^{(1)}\le64K_0\cdot C_w^{(1)}=C_w$.
The constant $C_P$ defined in \cite[Section~15.7]{QEDI} is finite;
consequently $C^P$, $C_w^{(0)}$, $C_w^{(1)}$, $C_w$ are finite, and $C_d$ is a product of finite numbers.
Finally, $M\ge1$ by (M1), and $0<8\alpha_0\le27\alpha_0$; this gives $M^{8\alpha_0}\le M^{27\alpha_0}$, while $4\exp(2\gamma)\ge1$.
\end{proof}

\paragraph{The constants of Proposition~14.1.}
The constants $C^F$, $\alpha^F$, $M^F$, $C_T$ that occur in Proposition~14.1 are the entries of Proposition~14.2\,(a) that bear these names.
The constants $C_{\mathrm{By}}$, $C_{\Gamma}$ of Proposition~5, $C_1$ of Proposition~7, $C_T$ of Proposition~9 and $C_d$ of Proposition~13,
and the constants $a'$, $\gamma^F$ defined at the head of the Propositions~2.$m$, are given by the same formulas as the corresponding entries of Proposition~14.2\,(a).
The number $L^F=5$ is the lower bound for $L$ in (L5), and the inequality $0<\gamma^F\le\frac{1}{6L}$ is (K1).

\subsubsection*{2.2. The two norms}

\begin{named}{Lemma 2}
\begin{itemize}
\item[(a)]
$|v|_c\le|v|_2\le2|v|_c$ for every $v\in\mathbb{C}^4$.
\item[(b)]
$|m|_c\le2|m|_2$ for every $4\times4$ matrix $m$, and $\left|\mathrm{tc}\cdot\mathbf{1}_4\right|_c=\left|\mathrm{tc}\cdot\mathbf{1}_4\right|_2=|\mathrm{tc}|$ for every complex number $\mathrm{tc}$.
\item[(c)]
Let $O$ be a linear operator between two of the spaces $\mathfrak{S}^{\mathrm{ch}}(\eta_k,\ell)$, $\mathfrak{S}(1,\ell)$, the functions of short pairs and the functions of pairs.
If $|Of|_{\mathrm{su}}\le\mathrm{cf}\cdot|f|_{\mathrm{su}}$ for all $f$, then $|Of|_{\infty}\le2\mathrm{cf}\cdot|f|_{\infty}$ for all $f$.
\end{itemize}
\end{named}

\begin{proof}
(a)
One has $\max_{\beta}|v_\beta|^2\le\sum_{\beta\le4}|v_\beta|^2\le4\max_{\beta}|v_\beta|^2$.

(b)
For $v\in\mathbb{C}^4$ one has $|mv|_c\le|mv|_2\le|m|_2\cdot|v|_2\le2|m|_2\cdot|v|_c$ by (a);
moreover $|\mathrm{tc}\cdot v|_c=|\mathrm{tc}|\cdot|v|_c$ and $|\mathrm{tc}\cdot v|_2=|\mathrm{tc}|\cdot|v|_2$.

(c)
Applying (a) at every site, or at every pair, one finds $|F|_{\infty}\le|F|_{\mathrm{su}}\le2|F|_{\infty}$ for every function $F$ of each of the four kinds.
Hence $|Of|_{\infty}\le|Of|_{\mathrm{su}}\le\mathrm{cf}\cdot|f|_{\mathrm{su}}\le2\mathrm{cf}\cdot|f|_{\infty}$.
\end{proof}

\subsubsection*{2.3. Proof of (i)}

\begin{proof}[Proof of Proposition~14.1\,(i)]
\textbf{FE1$^{\mathrm{ch}}$, FE2$^{\mathrm{ch}}$ ($k\ge1$).}
Let $y\in Y_k$ and $A\in\tilde{G}_k$.
By Section~I.3.5,
\[
F_y(A)=D_A+\mu+\boldsymbol{\alpha}_y\cdot J\,P_k(A)\,J^{*}+\boldsymbol{\beta}_y\cdot J\,P_{k+1}(A)\,J^{*},\qquad \boldsymbol{\alpha}_{\infty}=b_k,\qquad \boldsymbol{\beta}_{\infty}=0.
\]
Hence $F_{\infty}(A)=D_A+\mu+b_k\cdot J\,P_k(A)\,J^{*}$, $S_k(A)=S_{\infty}(A)$, $S_{k,y}(A)=S_y(A)$ for real $y$, and $S^{0}_{k+1}(A)=S_0(A)=S_{k,0}(A)$.
By Proposition~2.8\,(a), $F_y(A)$ is bijective on $\mathfrak{S}^{\mathrm{ch}}(\eta_k,\ell)$.
By Proposition~2.8\,(b), the map $A\mapsto S_y(A)$ is analytic on $\tilde{G}_k$, an open set of complex bond functions \cite[Section~16.2.7, (G2)]{QEDI}.
By Proposition~2.8\,(c), in which $w_k(y)=1$ for $k\ge1$, one has for all unit sites $z,z'$ and all $f\in\mathfrak{S}^{\mathrm{ch}}(\eta_k,\ell)$
\begin{gather*}
\left|\mathbf{1}_{\Delta_z}\cdot S_y(A)\cdot\mathbf{1}_{\Delta_{z'}}\,f\right|_{\mathrm{su}}\le32a'\cdot\exp\left(-\gamma^F\cdot d(z,z')\right)\cdot|f|_{\mathrm{su}},\\
\left|\mathbf{1}_{\Delta_z}\cdot\delta_{\alpha,A}S_y(A)\cdot\mathbf{1}_{\Delta_{z'}}\,f\right|_{\mathrm{su}}\le304a'\cdot\exp\left(-\gamma^F\cdot d(z,z')\right)\cdot|f|_{\mathrm{su}},
\end{gather*}
where in the second bound the supremum extends over the short pairs, the pairs inside a cluster included.
Lemma~2\,(c), applied to the operator $\mathbf{1}_{\Delta_z}\cdot S_y(A)\cdot\mathbf{1}_{\Delta_{z'}}$ on $\mathfrak{S}^{\mathrm{ch}}(\eta_k,\ell)$
and to the operator $\mathbf{1}_{\Delta_z}\cdot\delta_{\alpha,A}S_y(A)\cdot\mathbf{1}_{\Delta_{z'}}$ from $\mathfrak{S}^{\mathrm{ch}}(\eta_k,\ell)$ to the functions of short pairs, together with (K2), yields
\begin{align*}
\left|\mathbf{1}_{\Delta_z}\cdot S_y(A)\cdot\mathbf{1}_{\Delta_{z'}}\,f\right|_{\infty}&\le64a'\cdot\exp\left(-\gamma^F\cdot d(z,z')\right)\cdot|f|_{\infty}\le C^F\cdot\exp\left(-\gamma^F\cdot d(z,z')\right)\cdot|f|_{\infty},\\
\left|\mathbf{1}_{\Delta_z}\cdot\delta_{\alpha,A}S_y(A)\cdot\mathbf{1}_{\Delta_{z'}}\,f\right|_{\infty}&\le608a'\cdot\exp\left(-\gamma^F\cdot d(z,z')\right)\cdot|f|_{\infty}\le C^F\cdot\exp\left(-\gamma^F\cdot d(z,z')\right)\cdot|f|_{\infty}.
\end{align*}
For $y=\infty$ this is FE1$^{\mathrm{ch}}$; for real $y$ it is FE2$^{\mathrm{ch}}$, and $y=0$ is the case of $S^{0}_{k+1}(A)$.
The numbers $64a'$, $608a'$, $C^F$, $\gamma^F$ do not depend on $y$.

\textbf{FE3$^{\mathrm{ch}}$, $k\ge1$.}
Propositions~4 and~5 hold for all data (SF) with $k\ge1$, all $A\in\tilde{G}_k$ and all real $y$.
Hence Proposition~4\,(a)--(d) and Proposition~5\,(e) hold.
The constant $C^F$ in Proposition~5\,(e) is the constant $C^F$ of Proposition~14.1.
According to Proposition~5\,(e), all bounds of Proposition~5\,(a)--(d) hold with $|\cdot|_c$ in place of $|\cdot|_2$ and with the constants doubled
(this is Lemma~2\,(b), applied block by block, and term by term in the row sums and in the column sums),
and the doubled constants are at most $C^F$.
The comparisons, bound by bound, are the following;
they use $b_k\ge b_-$ (Section~I.3.4), $y^2\ge0$ and $L^{-3}\le1$.
\begin{itemize}
\item[(a)]
$P(A)(z,z')$:
$2\mathrm{e}\cdot L^{-3}\le2\mathrm{e}\le C^F$ by (K4).
The block $P(A)(z,z')$ is the number displayed in Proposition~5\,(a) times $\mathbf{1}_4$, because the phases $U_A$ are numbers (Section~I.3.3);
hence its norm $|\cdot|_c$ is the modulus of this number (Lemma~2\,(b)), and the bound holds already with $\mathrm{e}\cdot L^{-3}$.
\item[(b)]
$B_{k,y}(A)(z,z')$, first bound:
\[
2\left(b_k^2+y^2\right)^{-\frac{1}{2}}\cdot\left(\mathbf{1}[z=z']+2\mathrm{e}\cdot L^{-3}\right)\le\frac{2(1+2\mathrm{e})}{b_-}\le C_{\Gamma}\le\frac{C^F}{2}\qquad\text{by (K3)}.
\]
\item[(b)]
$B_{k,y}(A)(z,z')$, second bound:
$2C_{\mathrm{By}}\cdot(1+|y|)^{-1}\le C^F\cdot(1+|y|)^{-1}$ by (K3).
\item[(b)]
Row sums and column sums:
\[
2(1+2\mathrm{e})\cdot\left(b_k^2+y^2\right)^{-\frac{1}{2}}\le\frac{2(1+2\mathrm{e})}{b_-}\le C_{\Gamma}\le\frac{C^F}{2}\qquad\text{by (K3)}.
\]
\item[(c)]
$\left|\left(Q_k(A)\,J^{*}\cdot O\cdot J\,Q_k^{T}(-A)\right)(z,z')\right|_c\le2\mathrm{e}\cdot\left\|\mathbf{1}_{\Delta_z}\cdot O\cdot\mathbf{1}_{\Delta_{z'}}\right\|_{\mathrm{su}}$, the norm on the right-hand side being that of Proposition~5\,(c):
$2\mathrm{e}\le C^F$ by (K4).
\item[(d)]
$M_k(A)(z,z')$:
$64\mathrm{e}\cdot b^2\cdot a'\le2C_1\le C^F$ by (K3).
\item[(d)]
$\dot{M}_k(A)(z,z')$:
$2\mathrm{e}\cdot b^2\cdot(32a')^2\le C^F$ by (K3).
\item[(d)]
$\left(\Gamma_{k,y}(A)-B_{k,y}(A)\right)(z,z')$:
\[
256\mathrm{e}\cdot(1+2\mathrm{e})^2\cdot a'\cdot b_k^2\cdot\left(b_k^2+y^2\right)^{-1}\le256\mathrm{e}\cdot(1+2\mathrm{e})^2\cdot a'\le2C_{\Gamma}\le C^F\qquad\text{by (K3)}.
\]
\item[(d)]
$\Gamma_k(A)(z,z')$:
$2C_{\Gamma}\le C^F$ by (K3).
\end{itemize}
These are all nine bounds of Proposition~5\,(a)--(d).
In the last bound of Proposition~5\,(e) one has $D_k(A)=b_k-M_k(A)$ (Section~I.3.5) and $\left|b_k\cdot\mathbf{1}_4\right|_c=b_k\le b_+$ (Lemma~2\,(b), Section~I.3.4).

\textbf{FE3$^{\mathrm{ch}}$, $k=0$.}
Here $Y_0=\mathbb{R}$, and $F_y(A)=D_A+\mu+b_L\cdot J\,P(A)\,J^{*}+i\gamma_5 y$ acts on $\mathfrak{S}^{\mathrm{ch}}(1,\ell)$, with $192$ components over every photon site (Section~I.3.9).
By Proposition~2.8\,(a), $F_y(A)$ is bijective for every real $y$.
Proposition~2.8\,(d) concerns the kernel (I.1.4), that is,
\[
\left(S_{0,y}(A)f\right)(x)=\frac{1}{48}\sum_{c'}\sum_{g'\in O_{\mathrm{cub}}}S_{0,y}(A)\left(x,(c',g')\right)\cdot f(c',g'),
\]
and states that at all pairs of fermion sites
\[
\left|S_{0,y}(A)(x,x')\right|_2\le32a'\cdot(1+|y|)^{-1}\cdot\exp\left(-\gamma^F\cdot d(x,x')\right).
\]
By Lemma~2\,(b) and (K2),
\[
\left|S_{0,y}(A)(x,x')\right|_c\le64a'\cdot(1+|y|)^{-1}\cdot\exp\left(-\gamma^F\cdot d(x,x')\right)\le C^F\cdot(1+|y|)^{-1}\cdot\exp\left(-\gamma^F\cdot d(x,x')\right).
\]
For $y=0$ one has $\Gamma_0(A)=S_{0,0}(A)$ (Section~I.3.9), and therefore $\left|\Gamma_0(A)(x,x')\right|_c\le C^F\cdot\exp\left(-\gamma^F\cdot d(x,x')\right)$.
Finally, $\delta Z_0(A)=\det F_0(A)$, a determinant of size $192n_0$, does not vanish, because $F_0(A)$ is bijective.

\emph{A second derivation of the bound at $k=0$, from Proposition~2.8\,(c).}
By (I.ke4), a pointwise bound on a kernel of level $0$ carries the factor $48$ of the inverse weight.
In the derivation above this factor is contained in the constant $32a'$ of Proposition~2.8\,(d):
it is contained in $C_{00}$, hence in $a'\ge C_{00}$ and in $C^F$;
the step from Proposition~2.8\,(d) to FE3$^{\mathrm{ch}}$ multiplies $32a'$ by the factor $2$ of Lemma~2\,(b).
In the following derivation, from the operator bound of Proposition~2.8\,(c), the factor $48$ appears explicitly.
For $k=0$ one has $w_0(y)=(1+|y|)^{-1}$, and $\mathbf{1}_{\Delta_z}$ is the indicator of the cluster of $z$.
By Proposition~2.8\,(c) and Lemma~2\,(c),
\[
\left|\mathbf{1}_{\Delta_z}\cdot S_{0,y}(A)\cdot\mathbf{1}_{\Delta_{z'}}\,f\right|_{\infty}\le64a'\cdot(1+|y|)^{-1}\cdot\exp\left(-\gamma^F\cdot d(z,z')\right)\cdot|f|_{\infty}.
\]
Let $x=(z,g)$, $x'=(z',g')$ and $\mathrm{vv}\in\mathbb{C}^4$, and let $f$ be the element of $\mathfrak{S}^{\mathrm{ch}}(1,\ell)$ with $f(x')=\mathrm{vv}$ that vanishes at the other $48n_0-1$ sites of a fundamental domain.
Then $|f|_{\infty}=|\mathrm{vv}|_c$ and $f=\mathbf{1}_{\Delta_{z'}}f$, and by (I.1.4), in which a single term, of weight $\frac{1}{48}$, survives,
$\left(S_{0,y}(A)f\right)(x)=\frac{1}{48}\,S_{0,y}(A)(x,x')\,\mathrm{vv}$.
Therefore
\[
\left|S_{0,y}(A)(x,x')\,\mathrm{vv}\right|_c=48\left|\left(\mathbf{1}_{\Delta_z}\cdot S_{0,y}(A)\cdot\mathbf{1}_{\Delta_{z'}}\,f\right)(x)\right|_c\le48\cdot64a'\cdot(1+|y|)^{-1}\cdot\exp\left(-\gamma^F\cdot d(x,x')\right)\cdot|\mathrm{vv}|_c,
\]
which is (I.ke4); and $48\cdot64a'=3072a'\le C^F$ by (K2).

\textbf{FE4$^{\mathrm{ch}}$.}
Proposition~6 holds for all data (SF) and all $A\in\tilde{G}_k$.
For $k\ge1$, FE4$^{\mathrm{ch}}$ is Proposition~6\,(f).
The comparisons behind it are $64\mathrm{e}\cdot b\cdot a'\le C^F$ and $608\mathrm{e}\cdot b\cdot a'\le C^F$ (the two bounds of Proposition~6\,(c) and the two bounds of Proposition~6\,(a), doubled)
and $256\mathrm{e}\cdot b\cdot a'\le C^F$ (the bound of Proposition~6\,(b) for long pairs, doubled);
all of them are contained in (K4).
For $k=0$, Proposition~6\,(e) applies:
$H_0(A)=\bar{H}_0(A)=1$ on $\mathfrak{S}^{\mathrm{ch}}(1,\ell)$;
there is no short pair;
a pair inside a cluster has $\left(\delta_{\alpha,A}f\right)(x,x')=f(x')-f(x)$;
and $\left|H_0^{\#}(A)f\right|_{\infty}\le(1+\mathrm{e})\cdot|f|_{\infty}\le4|f|_{\infty}$, together with the barred form.
Moreover $4\le C^F$ by (K4).
The localised clauses of FE4$^{\mathrm{ch}}$ concern $k\ge1$ only.

\textbf{FE5$^{\mathrm{ch}}$.}
Proposition~9 holds for all data (SF).
Its clause (a) concerns all $k\ge0$; for $k=0$ it reads $T_0=\bar{T}_0=H_0=\bar{H}_0=1$.
Its clauses (b), (d), (e) concern $k\ge1$; the operators $T^{(1)}$, $T^{(2)}$ are not defined for $k=0$ (I.L0.9).
The constant in (d) is $C_T$.

\textbf{FE6$^{\mathrm{ch}}$.}
Proposition~10\,(a)--(f) holds for all data (SF) with $k\ge1$, and Proposition~10\,(g) holds for $k=0$, for the list of operators given there.
No constant occurs.

\textbf{FE8$^{\mathrm{ch}}$.}
This is Proposition~11\,(b).
Its data are a level $k\ge0$, two sides $\ell_1$ and $\ell_2=\mathrm{Jr}\cdot\ell_1$ in $\mathrm{Cl}$ with $\mathrm{Jr}$ odd, and numbers $\mu\in(0,1]$, $e_k\in(0,e_0]$, the same on both tori;
thus $(k,\ell_1,\mu,e_k)$ and $(k,\ell_2,\mu,e_k)$ are data (SF).
The lists of operators for $k\ge1$ and for $k=0$ are given in Proposition~11\,(b) separately.
No constant occurs.
\end{proof}

\subsubsection*{2.4. Proof of (ii)}

\begin{named}{Lemma 3}[The parameters of a property]
Let families as in Section~G1.11 be given at data (SF), and let
\[
0\le\tilde{\gamma}\le\gamma,\qquad 0\le\tilde{\mathrm{kq}}\le\mathrm{kq},\qquad 0<C\le\tilde{C}.
\]
\begin{itemize}
\item[(a)]
Each of the properties W1, W2, W5--W16 for $(\alpha_0,\gamma,\mathrm{kq};C)$ implies the same property for $(\alpha_0,\tilde{\gamma},\tilde{\mathrm{kq}};\tilde{C})$.
\item[(b)]
($k=0$)
If W17 holds for $(\alpha_0,\gamma,\mathrm{kq};C)$, and if
\[
\left|H_0(s,A)f\right|_{\infty}\le4\exp(2\tilde{\gamma})\cdot M^{27\alpha_0}\cdot|f|_{\infty}\qquad\text{for all }s\in\mathrm{Pd}_0,\ A\in\tilde{G}_0,\ f\in\mathfrak{S}^{\mathrm{ch}}(1,\ell),
\]
then W17 holds for $(\alpha_0,\tilde{\gamma},\tilde{\mathrm{kq}};\tilde{C})$.
\end{itemize}
\end{named}

\begin{proof}
(a)
The polydisk $\mathrm{Pd}_k=\left\{|s_{\mathrm{pc}}|\le M^{\alpha_0}\right\}$ is the same for both parameter sets.
The properties W1, W2, W8--W16 contain none of $\gamma$, $\mathrm{kq}$, $C$.
The property W5 contains $C$ only, as the factor on the right-hand sides of its two inequalities.
The right-hand sides of W6 and W7 are
\[
C\cdot\exp\left(-\gamma\cdot d(t,t')\right),\qquad C\cdot\exp\left(-\gamma\cdot d(z,z')-\mathrm{kq}\cdot d_M(X)\right)
\]
and, for $k=0$,
\[
C\cdot(1+y)^{-1}\cdot\exp\left(-\gamma\cdot d(x,x')-\mathrm{kq}\cdot d_M(X)\right);
\]
in W7 for $k\ge1$ the second of these expressions is multiplied by $|f|_{\infty}$, a non-negative factor that does not depend on the parameter set.
Distances are non-negative, and $d_M(X)\ge0$, since $M\cdot d_M(X)$ is a length (Section~I.1.1).
Hence each right-hand side is at most the same expression formed with $(\tilde{\gamma},\tilde{\mathrm{kq}},\tilde{C})$.

(b)
The property W17 consists of three groups:
W1, W2, W8--W11, W13, W14, W16 for $H_0(s,A)$;
W5 for $H_0(s,A)$ with the constant $4\exp(2\gamma)\cdot M^{27\alpha_0}$, which for $k=0$, where no pair is short, is the single bound on $\left|H_0(s,A)f\right|_{\infty}$;
and W1, W2, W6--W9, W12--W16 for $\Gamma_0(s,A)$ and $S_{0,y}(X,A)$.
Part (a) applies to the first and to the third group;
the second group, formed with $\tilde{\gamma}$, is the hypothesis.
\end{proof}

\begin{proof}[Proof of Proposition~14.1\,(ii)]
\textbf{Step 1} (the families are of the kind required in Section~G1.11).
Let $k\ge1$, and let $s\in\mathrm{Pd}_k$, $A\in\tilde{G}_k$, $X\in D_k$ and $y\in Y_k$ (in the family (p5), $y\in[0,\infty)$).
The families $S_y(s,A)$, $S^{\mathrm{sy}}_y(s,A)$, $S_y(X,A)$, $S^{\mathrm{sy}}_y(X,A)$ of Section~G1.9 are sums over walks of products of operators on $\mathfrak{S}^{\mathrm{ch}}(\eta_k,\ell)$
(the sums over walks that define $S_{\infty}(s,A)$, $S^{\mathrm{sy}}_0(s,A)$, $S^{\mathrm{sy}}_{\infty}(X,A)$ and $S_y(X,A)$, which enter the families of Proposition~7, converge by Proposition~3\,(a)).
The operator $B_k(A)=B_{k,0}(A)$ exists (Proposition~4\,(a)).
By the table of Section~I.3.5, $J\,Q_k^{T}(-A)$ maps $\mathfrak{S}(1,\ell)$ to $\mathfrak{S}^{\mathrm{ch}}(\eta_k,\ell)$, $Q_k(A)\,J^{*}$ maps $\mathfrak{S}^{\mathrm{ch}}(\eta_k,\ell)$ to $\mathfrak{S}(1,\ell)$, and $B_k(A)$ acts on $\mathfrak{S}(1,\ell)$.
Hence, for the families of Proposition~7,
\begin{gather*}
H_k(s,A)=b_k\cdot S_{\infty}(s,A)\cdot J\,Q_k^{T}(-A)\colon\mathfrak{S}(1,\ell)\to\mathfrak{S}^{\mathrm{ch}}(\eta_k,\ell)\qquad(4\to192\text{ components});\\
\Gamma_k(s,A)=B_k(A)+b_k^2\cdot B_k(A)\cdot Q_k(A)\,J^{*}\cdot S^{\mathrm{sy}}_0(s,A)\cdot J\,Q_k^{T}(-A)\cdot B_k(A)\quad\text{acts on }\mathfrak{S}(1,\ell);\\
M_k(X,A)=b_k^2\cdot Q_k(A)\,J^{*}\cdot S^{\mathrm{sy}}_{\infty}(X,A)\cdot J\,Q_k^{T}(-A)\quad\text{acts on }\mathfrak{S}(1,\ell);\\
S_{k,y}(X,A)=S_y(X,A)\quad\text{acts on }\mathfrak{S}^{\mathrm{ch}}(\eta_k,\ell).
\end{gather*}
The connection $-A$ lies in $\tilde{G}_k$ \cite[Section~16.2.7, (G2)]{QEDI}; so $\bar{H}_k(s,A)=\mathsf{C}^{-1}\cdot H_k(s,-A)\cdot\mathsf{C}$ is defined.
The map $y\mapsto S_{k,y}(X,A)$ is continuous on $[0,\infty)$ (Proposition~7), hence Borel.

Let now $k=0$, and consider the families of Proposition~8.
The operators $\Gamma_0(s,A)=S^{\mathrm{sy}}_0(s,A)$ and $S_{0,y}(X,A)=S_y(X,A)$ are sums over walks of products of operators on $\mathfrak{S}^{\mathrm{ch}}(1,\ell)$
(the series converge by Proposition~3\,(a));
they act on $\mathfrak{S}^{\mathrm{ch}}(1,\ell)$, with $192$ components over every photon site.
The operator $H_0(s,A)$ is the multiplication by
\[
\prod\left\{s_{\mathrm{pc}}:\ \mathrm{pc}\in\pi_0,\ \mathrm{pc}\text{ meets }K_c\right\},
\]
the same factor on the $48$ sites over $c$, and $\bar{H}_0(s,A)=\mathsf{C}^{-1}\cdot H_0(s,-A)\cdot\mathsf{C}=H_0(s,A)$;
there is no $M_0(X,A)$.
The map $y\mapsto S_{0,y}(X,A)$ is continuous, hence Borel, on $[0,\infty)$ (Proposition~8\,(a)).

These are the fermionic parts (p3), (p4), (p5) of a system of primitives of level $k$ in the sense of Section~G1.11 (Section~I.5.8).

\textbf{Step 2} (the properties for $2C_1$).
Recall that $\kappa'=\kappa_{\mathrm{w}}+\kappa_0$ (Section~G1.1).

Let $k\ge1$.
By Proposition~7 the families have W1, W2, W5--W16 for $(\alpha_0,\gamma^F,\kappa';2C_1)$;
for $k\ge1$ the property W17 asserts nothing.

Let $k=0$.
By Proposition~8\,(a), the families $\Gamma_0(s,A)$ and $S_{0,y}(X,A)$ have W1, W2, W6--W9, W12--W16, each as far as it concerns the family, for $(\alpha_0,\gamma^F,\kappa';2C_1)$,
with W6 and W7 as bounds on the kernels (I.1.4) at pairs of fermion sites and with the factor $(1+y)^{-1}$ in W7.
By Proposition~8\,(b), $H_0(s,A)$ has W1, W2, W8--W11, W13, W14, W16, and
\begin{equation*}
\left|H_0(s,A)f\right|_{\infty}\le M^{8\alpha_0}\cdot|f|_{\infty}\le4\exp(2\gamma)\cdot M^{27\alpha_0}\cdot|f|_{\infty}\qquad\text{for every }\gamma\ge0.
\tag{1}
\end{equation*}
For $\gamma=\gamma^F$ this is W17 for $(\alpha_0,\gamma^F,\kappa';2C_1)$.

\emph{What W1, W2, W5--W16 say at level $0$.}
In Section~G1.11 the family $H_k(s,A)$ occurs in W1, W2, W5, W8--W11, W13, W14, W16;
the family $\Gamma_k(s,A)$ in W1, W2, W6, W8, W9, W13--W16;
the family $S_{k,y}(X,A)$ in W1, W2, W7, W12--W14;
the family $M_k(X,A)$ does not exist for $k=0$;
and W6, W7, W13, W15 are taken in the forms that Section~G1.11 gives for $k=0$.
All these entries occur in the lists of W17, with one difference:
in W17 the constant of W5 is $4\exp(2\gamma)\cdot M^{27\alpha_0}$, whereas in W5 it is the constant $C$ of the parameter set.
The property W5 holds with the constant of the parameter set as well.
Indeed, its second part is void, since no pair is short, and, as $M\ge1$ and $a'\ge1$,
\begin{equation*}
M^{8\alpha_0}\le M^{27\alpha_0}\le304\mathrm{e}\cdot\max(1,b)\cdot a'\cdot M^{27\alpha_0}\le C_1\le2C_1.
\tag{2}
\end{equation*}
Thus the families of level $0$ have W1, W2, W5--W16, each as far as it concerns a family of level $0$, and W17, for $(\alpha_0,\gamma^F,\kappa';2C_1)$.

\textbf{Step 3} (the constant $C^F$).
By (K3), $2C_1\le C^F$.
Lemma~3, applied with $(\tilde{\gamma},\tilde{\mathrm{kq}},\tilde{C})=(\gamma^F,\kappa',C^F)$, the hypothesis of Lemma~3\,(b) being (1) at $\gamma=\gamma^F$, shows that the families have W1, W2, W5--W17 for $(\alpha_0,\gamma^F,\kappa_{\mathrm{w}}+\kappa_0;C^F)$.

\textbf{Step 4} ($\mathrm{par}^0$ and $\mathrm{par}$).
By (K5), $\gamma_{\mathrm{w}}\le\gamma^F$, $\kappa_{\mathrm{w}}\le\kappa_{\mathrm{w}}+\kappa_0$ and $C^F\le C_w^{(1)}\le C_w$.
Lemma~3, applied to the result of Step~3 with $(\tilde{\gamma},\tilde{\mathrm{kq}},\tilde{C})=\left(\gamma_{\mathrm{w}},\kappa_{\mathrm{w}}+\kappa_0,C_w^{(1)}\right)$ and with $(\tilde{\gamma},\tilde{\mathrm{kq}},\tilde{C})=\left(\gamma_{\mathrm{w}},\kappa_{\mathrm{w}},C_w\right)$,
the hypothesis of Lemma~3\,(b) being in both cases (1) at $\gamma=\gamma_{\mathrm{w}}$, shows that the families have W1, W2, W5--W17 for $\mathrm{par}^0$ and for $\mathrm{par}$.
The constant $4\exp(2\gamma)\cdot M^{27\alpha_0}$ of W17 decreases when $\gamma$ decreases from $\gamma^F$ to $\gamma_{\mathrm{w}}$;
for this reason (1) is used at $\gamma=\gamma_{\mathrm{w}}$.
By (1) and (K5),
\[
M^{8\alpha_0}\le4\exp(2\gamma_{\mathrm{w}})\cdot M^{27\alpha_0}\le C_w^{(1)}\le C_w,
\]
which is the chain of inequalities in Proposition~14.1\,(ii).

\textbf{Step 5} (W14; translations).
By Proposition~7 ($k\ge1$) and Proposition~8\,(a), (b) ($k=0$), the property W14 holds for all $\tilde{g}\in\mathbb{G}_k$, and for the family (p5) for all $\tilde{g}\in\mathbb{G}_{k,3}$.
Let $\ell_1\in\mathrm{Cl}$ divide $\ell$, and let $a\in\ell_1\mathbb{Z}^3$.
Since $\ell_1/(LM)$ is an integer, $a$ lies in $LM\mathbb{Z}^3$, which is contained in $M\mathbb{Z}^3$;
so $\tilde{g}:=\left(1,\,a\bmod\ell\mathbb{Z}^3\right)$ is an element of $\mathbb{G}_k$, and even of $\mathbb{G}_k^{L}$ (Section~I.4.2).
The matrix $S=1$ commutes with $\gamma_5$; hence $\tilde{g}\in\mathbb{G}_{k,3}$ (Section~I.4.1).
Moreover $r(1)=1$, and Section~I.4.3 gives
\[
\tilde{g}(c,g)=(c+a,g),\quad (\tilde{g}f)(x)=f(x-a),\quad (\tilde{g}A)(c,c')=A(c-a,c'-a),\quad (\tilde{g}s)_{\mathrm{pc}}=s_{\mathrm{pc}-a},\quad \tilde{g}X=X+a.
\]
Therefore W14 holds for the translations by $\ell_1\mathbb{Z}^3$, for all families, (p5) included.
\end{proof}

\subsubsection*{2.5. Proof of (iii)}

By definition, FE7$^{\mathrm{ch}}$ consists of the clauses (d0)--(d6) of Proposition~13.
Recall the data of that proposition.
They are data (SF), together with operators $S_{k,y}(X,A)$ on $\mathfrak{S}^{\mathrm{ch}}(\eta_k,\ell)$, indexed by $X\in D_k$, $A\in\tilde{G}_k$ and $y\in[0,\infty)$, which are Borel in $y$ and have the properties W1, W2, W7, W12, W13 and W14 of Section~G1.11 for $\mathrm{par}$; here W7 is taken at $k=0$ in its form for $k=0$, and W14 is required for $\tilde{g}\in\mathbb{G}^{L}_{k,3}$.
In (d4) the side is subject to $\ell\ge M\cdot r_k$, and in (d5), (d6) the data are (SF+).
The constant of Proposition~13 is $C_{d}$, formed with the constant $C_{w}$ of $\mathrm{par}$; by (K5) it is one finite number, the same for all levels $k\ge0$.
At $k=0$ the traces in $\acute{E}_0$ are traces on $\mathfrak{S}^{\mathrm{ch}}(1,\ell)$ (Section~I.5.7, (I.ke2)), and $\delta Z_0(A)=\det F_0(A)$ is a determinant of size $192\,n_0$ (Proposition~13\,(d0)).
For every family with these properties Proposition~13 asserts the clauses (d0)--(d6).

\emph{The family of (ii).}
The operators $S_{k,y}(X,A)=S_{y}(X,A)$ act on $\mathfrak{S}^{\mathrm{ch}}(\eta_k,\ell)$ and are Borel in $y$ (2.4, Step~1).
They have W1, W2, W7, W12 and W13 for $\mathrm{par}$ (Step~4; at $k=0$ with W7 in its form for $k=0$), and W14 for all $\tilde{g}\in\mathbb{G}_{k,3}$ (Step~5).
Since $\mathbb{G}^{L}_{k,3}\subset\mathbb{G}_{k,3}$ (Section~I.4.2), this family satisfies the hypotheses of Proposition~13.
Hence (d0)--(d6) hold for it, with the same constant $C_{d}$.
The Remark ``The projected families'' that follows Proposition~14.1 is not part of (iii), and no step of the argument above rests on it; that the projected families (p5), together with their data, satisfy the hypotheses of Proposition~13 follows from Propositions~26.6\,(b) and 26.3\,(b), as is shown in that Remark and in the Remark ``How (a)--(d) are applied'' that follows the proof of Proposition~26.6 in Chapter~26b.
\qed

\subsubsection*{2.6. Proof of (iv)}

\paragraph{2.6.1. No constant depends on $k$, $\ell$, $\mu$, $e_k$, $A$, $y$, $s$, $X$.}
The constants occurring in (i)--(iii) are
$C^{F}$, $\gamma^{F}$, $\alpha^{F}=\alpha_{0}$, $L^{F}$, $M^{F}$, $C_{T}$, $C_{1}$, $C_{\mathrm{By}}$, $C_{\Gamma}$, $a'$;
$\kappa_{0}$, $\kappa_{\mathrm{w}}$, $\kappa'$, $\gamma_{\mathrm{w}}$, $C_{w}^{(0)}$, $C_{w}^{(1)}$, $C_{w}$ (the entries of $\mathrm{par}^{0}$ and $\mathrm{par}$);
$C_{d}$, $c_{X}$;
the numbers $b$, $b_{L}$, $b_{-}$, $b_{+}$, $L$, $M$, $\alpha$, $\kappa$, $\varepsilon$, $r$, which occur in the clauses applied;
and the ceiling $e_0$ of (SF).
The following list contains them, together with all constants that occur in their definitions.
Each entry gives the line of Proposition~14.2\,(a), the constants with their definitions, and the letters that the entry contains; the entries are arranged so that no entry contains a letter of a later entry.

\begin{description}
\item[1, (C0).] $\rho$, a number in $(0,1)$; $w=(6,6\sqrt{2},6\sqrt{2},6)$, $w_3=6$, $\mathrm{sw}=12+12\sqrt{2}$; $\bar{\varepsilon}=\frac{1}{200}$; $c_{X}=54$; $\kappa_{0}=54\log 72$; $b=1$.
Letters contained: none.
\item[2, (C0).] $K_{0}=\frac{\exp(\kappa_{0})}{36}$; $c_{\mathrm{rb}}=c_{X}\cdot K_{0}$; $c_{\mathrm{fl}}=61c_{X}\cdot K_{0}$.
Letters contained: $\kappa_{0}$, $c_{X}$.
\item[3, (C0r).] $\lambda_{\mathrm{int}}$, the least non-zero eigenvalue of $\sum_{i\le2}w_i(1-R_i)$ on $\ell^2(O_{\mathrm{cub}})$ ((I.O8)).
Letters contained: $w$.
\item[4, (C0r).] $\hat{g}$, $c_{\mathrm{hl}}$, $c_{*}$, $K_{\rho}$, $A_1$, $A_2$, $c_{b}$, $C_{\mathrm{Po}}$ (formulas in Proposition~14.2\,(a)).
Letters contained: $\rho$, $w_3$, $\lambda_{\mathrm{int}}$.
\item[5, (C1).] $\alpha$, a number in $(1-\bar{\varepsilon},1)$.
Letters contained: $\bar{\varepsilon}$.
\item[6, (C2).] $\kappa$, a number $\ge7\kappa_{0}+7$.
Letters contained: $\kappa_{0}$.
\item[7, (C3).] $L$, an odd integer with (L1)--(L5) ($L^{F}=5$ is the threshold in (L5), a numeral).
Letters contained: $\bar{\varepsilon}$, $c_{\mathrm{fl}}$, $c_{\mathrm{rb}}$, $\kappa$, $\kappa_{0}$.
\item[8, (C4).] $b_{L}=\frac{b}{L}$, $b_{-}=b\left(1-\frac{1}{L}\right)$, $b_{+}=b$; $N_c=(2L+1)^3$; $c_{C}=3L^{\frac{9}{2}}$.
Letters contained: $b$, $L$.
\item[9, (C4).] $\alpha^{P}=\frac{1}{4}$, $\alpha^{F}=\frac{1}{54}$, $\alpha_{0}=\min(\alpha^{F},\alpha^{P})=\frac{1}{54}$.
Letters contained: none.
\item[10, (C4).] $\gamma_{P}$, $C_{P}$, the constants defined in \cite[Section~15.7]{QEDI} before \cite[Proposition~15.24]{QEDI}; $\gamma^{P}=\frac{\gamma_{P}}{2}$.
Letters contained: in $\gamma_{P}$, the letter $L$; in $C_{P}$, the letters $L$, $\alpha$ (final paragraph of \cite[Section~15.7]{QEDI}).
\item[11, (C4).] $\kappa_{\mathrm{w}}=\kappa+\kappa_{0}+3$; $\kappa'=\kappa_{\mathrm{w}}+\kappa_{0}$; $\hat{\kappa}=N_c\cdot\kappa+\log(6c_{X})$.
Letters contained: $\kappa$, $\kappa_{0}$, $N_c$, $c_{X}$.
\item[12, (C4).] $C'_{V}=2\mathrm{e}\cdot(42+6b)$; $C_{T}=2\mathrm{e}^{\frac{3}{2}}\left[b_{-}^{-1}\left(54L\cdot\mathrm{e}^{6}+12\mathrm{e}+1\right)+1\right]$.
Letters contained: $b$, $b_{-}$, $L$.
\item[13, (C4r).] $K_{1}$, $C_{I}$, $c_{\mathrm{wt}}$, $a_{0}$, $\gamma_{(5)}$ (formulas in Proposition~14.2\,(a)).
Letters contained: $L$, $b$, $b_{-}$, $\rho$, $C_{\mathrm{Po}}$, $c_{b}$.
\item[14, (C4r).] $C_{5}=\mathrm{e}^{2}\cdot L\cdot C_{3}$, the constant of Proposition~1.6\,(c).
Letters contained: $b$, $\alpha$, $L$, $\rho$ (Proposition~1.6).
\item[15, (C4r).] $c_{\mathrm{re}}$, $a_{z}$, $k_{1}$, $C_{00}$ (formulas in Proposition~14.2\,(a)).
Letters contained: $b$, $L$, $\rho$, $C_{\mathrm{Po}}$.
\item[16, (C4r).] $C_{z}=\max(C_{5},C_{00})$, $\gamma_{z}=\min(\gamma_{(5)},a_{z})$, $\gamma^{F}=\frac{\gamma_{z}}{2}$, $a'=C_{z}\cdot c(\gamma^{F})$, $\gamma_{\mathrm{w}}=\min(\gamma^{F},\gamma^{P})$.
Letters contained: $C_{5}$, $C_{00}$, $\gamma_{(5)}$, $a_{z}$, $\gamma^{P}$.
\item[17, (C4r).] $M^{F}=\max\left\{4L+1,\ (108\pi\cdot C'_{V}\cdot a')^2,\ 108\pi\cdot C'_{V}\cdot a'\cdot\exp(27\kappa')\right\}$;
$M^{P}$, the least power $M_0\ge L^2$ of $L$ such that $2\alpha_{0}\cdot402\cdot\log x\le\gamma_{P}\cdot x$ for all $x\ge M_0$ (\cite[Section~15.10.3]{QEDI}; $x$ is a bound variable).
Letters contained: $L$, $C'_{V}$, $a'$, $\kappa'$, $\alpha_{0}$, $\gamma_{P}$.
\item[18, (C5).] $M$, a power of $L$ with (M1)--(M6).
Letters contained: $L$, $\gamma_{\mathrm{w}}$, $\alpha_{0}$, $\hat{\kappa}$, $M^{F}$, $M^{P}$.
\item[19, (C6).] $C_{\mathrm{By}}$, $C_{\Gamma}$, $C_{1}$, $C^{F}$ (formulas in Proposition~14.2\,(a)).
Letters contained: $b$, $b_{-}$, $a'$, $M$, $\alpha_{0}$, $\kappa'$.
\item[20, (C6).] $C^{P}=116\,C_{P}\cdot M^{402\alpha_{0}}$, $C_{w}^{(0)}$, $C_{w}^{(1)}$, $C_{w}$ (formulas in Proposition~14.2\,(a)).
Letters contained: $C_{P}$, $M$, $\alpha_{0}$, $C^{F}$, $\gamma_{\mathrm{w}}$, $K_{0}$.
\item[21, (C6).] $c_{\gamma}=\sup_{t}\sum_{t'}\nu(t')\cdot\exp\left(-\gamma_{\mathrm{w}}\cdot d(t,t')\right)$ (Section~I.5.8).
Letters contained: $\gamma_{\mathrm{w}}$, $L$, $M$ ($M$ through the class $\mathrm{Cl}$; $k$, $\ell$, $t$, $t'$ are bound variables, see below).
\item[22, (C6).] $C_{B}=C_{P}\cdot\left[(s_d+4)\cdot c_{\mathrm{bd}}(\gamma_{P})+\left(\frac{2}{\gamma_{P}}+\frac{62}{\gamma_{P}^{2}}\right)\cdot c_{\mathrm{bd}}\left(\frac{\gamma_{P}}{2}\right)\right]$, $s_d=(2L+1)M$, $c_{\mathrm{bd}}(\gamma)=3\exp(3\gamma)\cdot c(\gamma)$ (\cite[Section~15.9]{QEDI}, \cite[Section~15.1.1]{QEDI}).
Letters contained: $C_{P}$, $\gamma_{P}$, $L$, $M$.
\item[23, (C6).] $C_{d}=16c_{X}\cdot C_{w}\cdot\max\left\{\frac{\pi}{2}\cdot b\cdot\mathrm{e}\cdot(1+2\mathrm{e}),\ 2+58\mathrm{e}+b\cdot\mathrm{e}\right\}$.
Letters contained: $c_{X}$, $C_{w}$, $b$.
\item[24, (C6).] $c_{S}=10^{5}c_{X}^{5}\cdot K_{0}^{5}\cdot\exp(6\kappa)$.
Letters contained: $c_{X}$, $K_{0}$, $\kappa$.
\item[25, (C7).] $r$, an integer $\ge2$; $p_0$, an integer $>3r+2$; $p$, an integer $>2p_0$.
Letters contained: none.
\item[26, (C7).] $\varepsilon$, a number in $(1-\alpha,\bar{\varepsilon}]$.
Letters contained: $\alpha$, $\bar{\varepsilon}$.
\item[27, (C8).] $e_0$, a number $>0$ with (A1)--(A10), (Z2.s1)--(Z2.s5), (Z3.s1).
Letters contained: $\varepsilon$, $p_0$, $p$, $c_{\mathrm{fl}}$, $M$, $L$, $C_{w}$, $c_{\gamma}$, $c_{C}$, $c_{X}$, $C_{T}$, $c_{S}$, $K_{0}$, $C'_{V}$, $a'$, $c_{\mathrm{rb}}$, $C_{d}$, $C_{B}$ ($v$ is a bound variable, see below).
\end{description}

Each entry contains, besides letters defined earlier in the same entry, only letters of entries above it.
The entries are of four kinds.

\emph{Formulas} (the numbers of entry~1, and the entries 2--4, 8, 9, 11--13, 15--17, 19, 20, 22--24).
Each constant is a function of the letters named.

\emph{Constants of a statement} (entries 10 and 14).
Proposition~1.6 asserts that $C_{5}$ depends on $(b,\alpha,L,\rho)$ only, and not on $k$, $R$, $\theta$, $\mu$, $y$.
There $R$ is the period of the torus on which Proposition~1.6 is stated; hence $C_{5}$ contains neither the side $\ell$ nor $M$, as its position before the line (C5) requires.
(The constants $C_{3}$, $C_{R}$, $C_{G}$, $K_{\mathrm{Str}}$, $c_{\mathrm{el}}$, $t_{0}$, $\theta_{0}$, $\kappa_{(0)}$, $A_{\mathrm{ak}}$, $A'_{\mathrm{ak}}$ of Proposition~14.2\,(a) enter the list only through $C_{5}$.)
For $\gamma_{P}$ and $C_{P}$, the dependence stated in entry~10 is asserted in the final paragraph of \cite[Section~15.7]{QEDI}.
(The constant $c_{P}$ of Proposition~14.2\,(a), formed with the numeral $C_{Q}=1$, enters the list only through $\gamma_{P}$ and $C_{P}$; by the same paragraph of \cite{QEDI} it depends on $L$ only.)

\emph{Numbers that are chosen} ($\rho$, $\alpha$, $\kappa$, $L$, $M$, $r$, $p_0$, $p$, $\varepsilon$, $e_0$).
Each of them is one number, chosen once, in the range or under the conditions indicated.
Apart from the number itself, the conditions contain the following letters:
(L1): $\bar{\varepsilon}$; (L2): $c_{\mathrm{fl}}$, $\bar{\varepsilon}$; (L3): $\kappa$, $\kappa_{0}$; (L4): $c_{\mathrm{rb}}$, $\kappa$, $\kappa_{0}$; (L5): none;
(M1), (M2): $L$; (M3): $\gamma_{\mathrm{w}}$; (M4): $\alpha_{0}$, $\hat{\kappa}$; (M5): $M^{F}$; (M6): $M^{P}$;
(S1), the range of $\varepsilon$: $\alpha$, $\bar{\varepsilon}$.
The conditions on $e_0$ are treated below.

\emph{A supremum} (entry 21).
By Section~I.5.8 the supremum defining $c_{\gamma}$ is formed over all levels $k\ge0$, all sides $\ell\in\mathrm{Cl}$ and all indices $t$, and over the three index sets over which $t'$ ranges in the sum:
(i) the pairs $(y,\beta)$ with $y\in T^0_k$ and $\beta\in\{1,\dots,4\}$, with $\nu=1$ (for $k\ge1$), and at $k=0$ the pairs $(x,\beta)$ with $x\in T^{\mathrm{f}}_0$, with $\nu=\frac{1}{48}$;
(ii) the positively oriented unit bonds, with $\nu=1$;
(iii) the basis $\{\Upsilon\}$, with $\nu=1$.
Thus $k$, $\ell$, $t$, $t'$ are bound variables in $c_{\gamma}$, and the number is formed once, before data (SF) are given.
(The class $\mathrm{Cl}$ over which $\ell$ ranges is formed with $L$ and $M$, and the basis $\{\Upsilon\}$ with $L$; these are letters of entries above entry~21.)

\emph{The ceiling $e_0$} (entry 27).
Besides numerals and the variable $v$, the inequalities of Proposition~14.2\,(b) imposed on $e_0$ contain the following letters:
(A1): $p_0$, $\varepsilon$, $c_{\mathrm{fl}}$, $M$;
(A2): $C_{w}$, $c_{\gamma}$, $c_{C}$, $\varepsilon$;
(A3): $c_{X}$, $c_{\gamma}$, $C_{w}$, $C_{T}$, $M$, $\varepsilon$;
(A4): $C_{w}$, $C_{T}$, $\varepsilon$;
(A5): $M$, $\varepsilon$;
(A6): $C_{w}$, $c_{\gamma}$, $\varepsilon$;
(A7): $c_{S}$, $L$, $c_{X}$, $K_{0}$, $\varepsilon$;
(A8): $M$, $\varepsilon$;
(A9): $L$, $M$, $c_{X}$, $\varepsilon$;
(A10): $M$, $\varepsilon$, $C_{B}$, $p$;
(Z2.s1): $L$, $\varepsilon$;
(Z2.s2): $C'_{V}$, $a'$, $\varepsilon$;
(Z2.s3): $M$, $\varepsilon$;
(Z2.s4): $\varepsilon$;
(Z2.s5): $c_{\mathrm{rb}}$, $L$, $C_{d}$, $M$, $\varepsilon$;
(Z3.s1): none.
Together these are the eighteen letters of entry~27:
$\varepsilon$ (entry 26); $p_0$, $p$ (25); $c_{\mathrm{fl}}$, $c_{\mathrm{rb}}$, $K_{0}$ (2); $c_{X}$ (1); $L$ (7); $c_{C}$ (8); $C_{T}$, $C'_{V}$ (12); $a'$ (16); $M$ (18); $C_{w}$ (20); $c_{\gamma}$ (21); $C_{B}$ (22); $C_{d}$ (23); $c_{S}$ (24).
The variable $v$ is bound: each inequality is imposed for all $v\ge-\log e_0$.
Consequently no condition on $e_0$ contains one of the letters $k$, $\ell$, $\mu$, $e_k$, $A$, $y$, $s$, $X$, and $e_0$ is chosen before the data.
For data (SF) one has $v_k=-\log e_k\ge-\log e_0$; so each of the sixteen conditions holds at $v=v_k$, for every $e_k\in(0,e_0]$.

None of the letters $k$, $\ell$, $\mu$, $e_k$, $A$, $y$, $s$, $X$ occurs free in the list.
By induction along the list, no constant depends on them.
(The exponent $s$ of the line (C7) of Proposition~14.2\,(a) is distinct from the family $s=(s_{\mathrm{pc}})$ of (iv), and it occurs in no entry of the list.)
Independence of $k$, $\ell$, $\mu$, $e_k$, $A$, $y$, $s$, $X$ is what the word ``constant'' means in (I.U4)$^{\mathrm{ch}}$ (Sections~I.8.1 and G1.1).

The clauses of Propositions~2.8 and 3--13 that are applied in 2.3--2.5 hold for all data (SF), all $A\in\tilde{G}_k$ and all $y$, $s$, $X$, $z$, $z'$, $x$, $x'$, $f$, with the numbers above as constants: in each of these clauses the constants come first, and the quantifiers follow them.
To these clauses Lemmas~1--3 add only the factor $2$ of Lemma~2, the factor $48$ of (I.1.4) in the second derivation of FE3$^{\mathrm{ch}}$ at $k=0$ (from Proposition~2.8\,(c)), and comparisons between constants.
It follows that the bounds are uniform in the level $k\ge0$, in the side $\ell\in\mathrm{Cl}$ (hence in the torus), in the mass $\mu\in(0,1]$, in the charge $e_k\in(0,e_0]$, and in $A$, $y$, $s$, $X$.
No cutoff $N$ occurs in (SF).

The following numbers occur in these clauses as functions of the data, and they are not constants:
$\eta_k=L^{-k}$, $b_{k}\in[b_{-},b_{+}]$, $\boldsymbol{\alpha}_{y}$, $\boldsymbol{\beta}_{y}$, $\left(b_{k}^{2}+y^{2}\right)^{-\frac{1}{2}}$, $c^{\mathrm{d}}_{k}$, $n_k=\ell^{3}$, $w_{k}(y)$, $v_k$, $r_k$, the powers $e_k^{\frac{1}{4}}$ and $e_{k+1}^{\frac{1}{4}-2\varepsilon}$ in Proposition~13\,(d4), (d6), and the distances $d(\cdot,\cdot)$ and $d_M(X)$.

\paragraph{2.6.2. Real constants; the inclusion.}
A real constant bond function on $T_k$ is a bond function $\mathrm{cst}$ with
\[
\mathrm{cst}(c,c+\eta_k\cdot e_{\mu})=\mathrm{cst}_{\mu}=-\mathrm{cst}(c+\eta_k\cdot e_{\mu},c)
\]
for three real numbers $\mathrm{cst}_{1}$, $\mathrm{cst}_{2}$, $\mathrm{cst}_{3}$.
The domains $R_k$, $\tilde{R}_k$, $G_k$, $\tilde{G}_k$, $\mathfrak{X}_k$ are those of Section~I.5.3, that is, of \cite[Definition~5.16]{QEDI}.

\begin{named}{Lemma 4}
Let $\mathrm{cst}$ be a real constant bond function on $T_k$.
\begin{itemize}
\item[(a)] $\tilde{G}_k+\mathrm{cst}=\tilde{G}_k$.
\item[(b)] $\tilde{R}_k+\mathrm{cst}=\tilde{R}_k$.
\item[(c)] $\tilde{R}_k+\mathfrak{X}_k\subset\tilde{G}_k$.
\end{itemize}
\end{named}

Assertion (a) is contained in the display (G2) of \cite[Section~16.2.7]{QEDI}, assertion (b) in \cite[Lemma~17.6]{QEDI} (stated also as \cite[Lemma~5.19]{QEDI}), and assertion (c) in the display (G3) of \cite[Section~16.2.7]{QEDI}.
The three are statements on bond functions.
The displays (G2) and (G3) hold under the data of \cite[Section~16.2.1]{QEDI}, namely $k\ge0$, $\ell\in\mathrm{Cl}$, $e_k\in(0,e_0]$, and (G3) moreover under (Z2.s4); \cite[Lemma~17.6]{QEDI}, apart from its assertion on $\mathrm{Dom}_k$, holds for every $k\ge0$, $\ell\in\mathrm{Cl}$ and $e_k\in(0,e_0]$ (\cite[Section~17.3]{QEDI}).
These are the data (SF), and (Z2.s4) is a condition of the line (C8).
In the notation of this paper the arguments are as follows.

\begin{proof}
Let $\square$ be an $M$-cube and $i\in\{5,L\}$.
Let $z_{\square}$ be a lift of the centre of $\square$; then $\square^{\sim i}$ is the image of the set of the sites $c\in\mathbb{L}_k$ with $|c-z_{\square}|<\frac{(2i+1)M}{2}$.
The sup-diameter of this set is $(2i+1)M-\eta_k$, and
\[
(2i+1)M\le(2L+1)M<13LM\le\ell,
\]
because $5\le L$ ((L5)) and $\ell\in\mathrm{Cl}$.
So $\square^{\sim i}$ has a lift in $\mathbb{L}_k$ on which the covering map is injective and whose bonds are mapped onto the bonds of $\square^{\sim i}$ (Section~I.1.1).
Put $\boldsymbol{\lambda}_1(c):=\sum_{\mu}\mathrm{cst}_{\mu}\cdot\tilde{c}_{\mu}$, where $\tilde{c}$ is the point of the lift over $c$.
Then $\boldsymbol{\lambda}_1$ is a real function on the sites of $\square^{\sim i}$, and
\[
(d\boldsymbol{\lambda}_1)(c,c+\eta_k\cdot e_{\mu})=\eta_k^{-1}\left(\boldsymbol{\lambda}_1(c+\eta_k\cdot e_{\mu})-\boldsymbol{\lambda}_1(c)\right)=\mathrm{cst}_{\mu};
\]
that is, $d\boldsymbol{\lambda}_1=\mathrm{cst}$ on the bonds of $\square^{\sim i}$.

(a) Let $A=A_0+A_1\in\tilde{G}_k$, where $A_1\in G_k$, $A_0$ is real, and for every $\square$ one has $A_0=A'+d\boldsymbol{\lambda}$ on $\square^{\sim5}$ with real $\boldsymbol{\lambda}$, $A'$ and $|A'|_{\infty}<e_k^{-1+\varepsilon}$ on $\square^{\sim5}$.
Then $A_0+\mathrm{cst}$ is real, $A_0+\mathrm{cst}=A'+d(\boldsymbol{\lambda}+\boldsymbol{\lambda}_1)$ on $\square^{\sim5}$ with the same $A'$, and $A+\mathrm{cst}=(A_0+\mathrm{cst})+A_1$.
Hence $A+\mathrm{cst}\in\tilde{G}_k$.
The same argument with $-\mathrm{cst}$ in place of $\mathrm{cst}$ gives the equality.

(b) Let $\mathscr{A}=A_0+A_1\in\tilde{R}_k$, where $A_1\in R_k$, $A_0$ is real, and for every $\square$ one has $A_0=A'+d\boldsymbol{\lambda}$ on $\square^{\dagger}=\square^{\sim L}$ with real $\boldsymbol{\lambda}$ and $A'$, the three norms of $A'$ on $\square^{\dagger}$ being smaller than the thresholds of $R_k$.
Then $A_0+\mathrm{cst}=A'+d(\boldsymbol{\lambda}+\boldsymbol{\lambda}_1)$ on $\square^{\dagger}$ with the same $A'$; so $\mathscr{A}+\mathrm{cst}=(A_0+\mathrm{cst})+A_1$ lies in $\tilde{R}_k$.
The same argument with $-\mathrm{cst}$ gives the equality.

(c) Let $\mathscr{A}=A_0+A_1\in\tilde{R}_k$ be as in (b), and $\mathscr{C}\in\mathfrak{X}_k$.
Then $\mathscr{A}+\mathscr{C}=A_0+(A_1+\mathscr{C})$, and
\[
|A_1+\mathscr{C}|_{\infty}<e_k^{-\frac{3}{4}+\varepsilon}+e_k^{-1+2\varepsilon}=e_k^{-1+\varepsilon}\cdot\left(e_k^{\frac{1}{4}}+e_k^{\varepsilon}\right)\le e_k^{-1+\varepsilon}
\]
by the second inequality of (Z2.s4); so $A_1+\mathscr{C}\in G_k$.
The cube $\square^{\sim5}$ is contained in $\square^{\sim L}$ ($L\ge5$), and there $A_0=A'+d\boldsymbol{\lambda}$ with $|A'|_{\infty}<e_k^{-\frac{3}{4}+\varepsilon}\le e_k^{-1+\varepsilon}$, because $e_k\le1$ (first inequality of (Z2.s4), together with $\varepsilon>0$).
Hence $\mathscr{A}+\mathscr{C}\in\tilde{G}_k$.
\end{proof}

\emph{Consequence.}
Let $\mathscr{A}\in\tilde{R}_k$, $\mathscr{C}\in\mathfrak{X}_k$, let $\mathrm{cst}$ be a real constant bond function, and put $\mathscr{T}:=\mathscr{A}+\mathscr{C}+\mathrm{cst}$.
By Lemma~4\,(b), $\mathscr{A}+\mathrm{cst}\in\tilde{R}_k$; hence $\mathscr{T}=(\mathscr{A}+\mathrm{cst})+\mathscr{C}$ lies in $\tilde{R}_k+\mathfrak{X}_k$, and by Lemma~4\,(c) in $\tilde{G}_k$.
Therefore:
\begin{itemize}
\item[(1)] every clause of (i)--(iii) that is stated for $A\in\tilde{G}_k$ holds at $A=\mathscr{T}$;
\item[(2)] Proposition~9\,(d), which is stated for $A\in\tilde{R}_k+\mathfrak{X}_k$, holds at $A=\mathscr{T}$;
\item[(3)] Proposition~13\,(d4), which is stated for pairs in $\tilde{R}_k\times\mathfrak{X}_k$, holds at the pair $(\mathscr{A}+\mathrm{cst},\mathscr{C})$, where
\[
{\delta E^{\mathrm{det}}}_k(X;\mathscr{A}+\mathrm{cst},\mathscr{C})=E^{d}_k(X,\mathscr{T})-E^{d}_k(X,\mathscr{C})
\]
(Section~I.5.7);
\item[(4)] in Proposition~13\,(d5), (d6) (data (SF+)) the first slot of the activities of level $k+1$ ranges over $\tilde{R}_{k+1}$, formed with the side $\ell/L\in\mathrm{Cl}$ and the charge $e_{k+1}\le e_0$; the data $(k+1,\ell/L,\cdot\,,e_{k+1})$ satisfy the hypotheses of Lemma~4, so $\tilde{R}_{k+1}+\mathrm{cst}=\tilde{R}_{k+1}$.
\end{itemize}
In all four cases the constant is counted with $\mathscr{A}$.
\qed

\subsubsection*{2.7. Proof of (v)}

For a real number $\mu$ put
\[
F_{y}(A;\mu):=D_A+\mu+\operatorname{Pen}_{y}(A)
\]
(Sections~I.3.5 and I.3.9), and let $S_{y}(A;\mu)$ denote its inverse wherever it exists.
The assertions of (v) on the range of the mass are proved in (v.1), those at $\mu=0$ in (v.3), and the values $\mathfrak{m}_0=1$, $m_0=0$, $h'=h$ in (v.2); the Remark ``Where the mass enters'' that follows Proposition~14.1 is treated in (v.4).

\medskip
\noindent\textbf{(v.1) The range.}
Every clause of (i)--(iii) holds for every $\mu\in(0,1]$; the constants $C^{F}$, $\gamma^{F}$, $C_{T}$, $C_{d}$, $\mathrm{par}^{0}$, $\mathrm{par}$ and the ceiling $e_0$ are the same for all these $\mu$; and no relation between $\mu$ and $(k,\ell,e_k)$ is required.

\begin{proof}
The mass $\mu\in(0,1]$ is the third entry of the data (SF), on which 2.3--2.5 rest.
The second and the third assertion are 2.6.1 and the definition of (SF).
\end{proof}

\noindent\textbf{(v.2) $\mathfrak{m}_0=1$, $m_0=0$, $h'=h$.}

\begin{proof}
In \cite{QEDI} the number $\mathfrak{m}_0\in(0,1]$ is the upper end of the range $(0,\mathfrak{m}_0]$ of the masses of the scale-free setting of \cite[Section~5.4]{QEDI}, that is, of the masses for which the statements of \cite[Section~5.4.4]{QEDI} hold and families with the properties of \cite[Theorem~5.72]{QEDI} exist; $m_0$ is the least integer $\ge0$ with $L^{-m_0}\le\mathfrak{m}_0$, and $h':=h\cdot L^{-m_0}$ (\cite[Definition~5.98]{QEDI}).
By (v.1) the value $\mathfrak{m}_0=1$ is admissible for FE1$^{\mathrm{ch}}$--FE8$^{\mathrm{ch}}$ and for the families of (ii), and $1$ is the largest element of $(0,1]$.
Since $L^{0}=1\le1$, one obtains $m_0=0$, and $h'=h\cdot L^{0}=h$.
Thus each of the three letters is another name of a number that has a name already, namely $1$, $0$ and $h$; accordingly the formulas of \cite{QEDI} that contain them are read with $1$, $0$, $h$ in their places (Section~I.6.1, and Section~I.8.4 with $h$ for $h'$ and $0$ for $m_0$).
\end{proof}

\noindent\textbf{(v.3)}
The assertions of (v) at $\mu=0$ are the content of the following lemma.

\begin{named}{Lemma 5}[The end point $\mu=0$]
Let $k\ge0$, $\ell\in\mathrm{Cl}$, $e_k\in(0,e_0]$, $y\in Y_{k}$ and $A\in\tilde{G}_k$.
Then $F_{y}(A;0)$ is bijective on $\mathfrak{S}^{\mathrm{ch}}(\eta_k,\ell)$, $S_{y}(A;\mu)\to S_{y}(A;0)$ as $\mu\to0$, $\mu>0$, and the bounds of FE1$^{\mathrm{ch}}$, FE2$^{\mathrm{ch}}$ ($k\ge1$) and of FE3$^{\mathrm{ch}}$ at $k=0$ hold for $S_{y}(A;0)$ with the same $C^{F}$, $\gamma^{F}$.
\end{named}

\begin{proof}
Let $\mu\in(0,1]$.
Then $(k,\ell,\mu,e_k)$ are data (SF), and the domain $\tilde{G}_k$ does not depend on $\mu$.
By (G1.N1) of Section~G1.3 and Proposition~2.8\,(c),
\[
\|S_{y}(A;\mu)\|_{\mathrm{su}}\le N_{\gamma^{F}}\left(S_{y}(A;\mu)\right)\le32a'\cdot w_{k}(y)\le32a'.
\]
Hence, for every $f\in\mathfrak{S}^{\mathrm{ch}}(\eta_k,\ell)$,
\[
|f|_{\mathrm{su}}=\left|S_{y}(A;\mu)F_{y}(A;\mu)f\right|_{\mathrm{su}}\le32a'\cdot\left|F_{y}(A;0)f+\mu f\right|_{\mathrm{su}}\le32a'\cdot\left(\left|F_{y}(A;0)f\right|_{\mathrm{su}}+\mu\cdot|f|_{\mathrm{su}}\right).
\]
The constant $a'$ does not contain $\mu$ (2.6.1).
Letting $\mu\to0$ one obtains $|f|_{\mathrm{su}}\le32a'\cdot\left|F_{y}(A;0)f\right|_{\mathrm{su}}$.
So $F_{y}(A;0)$ is injective on $\mathfrak{S}^{\mathrm{ch}}(\eta_k,\ell)$, a space of the finite dimension $192\,(\ell/\eta_k)^3$; hence it is bijective, and $\|S_{y}(A;0)\|_{\mathrm{su}}\le32a'$.
From
\[
S_{y}(A;\mu)-S_{y}(A;0)=S_{y}(A;\mu)\left(F_{y}(A;0)-F_{y}(A;\mu)\right)S_{y}(A;0)=-\mu\,S_{y}(A;\mu)\,S_{y}(A;0)
\]
it follows that $\|S_{y}(A;\mu)-S_{y}(A;0)\|_{\mathrm{su}}\le\mu\cdot(32a')^2\to0$.

Fix $z$, $z'$, $f$ if $k\ge1$, and $x$, $x'$ if $k=0$.
The left sides of the bounds are
\[
\left|\mathbf{1}_{\Delta_z}S\,\mathbf{1}_{\Delta_{z'}}f\right|_{\infty},\qquad\left|\mathbf{1}_{\Delta_z}\delta_{\alpha,A}S\,\mathbf{1}_{\Delta_{z'}}f\right|_{\infty},\qquad|S(x,x')|_{c}
\]
at $S=S_{y}(A;\mu)$: norms of images of $S$ under linear maps between spaces of finite dimension, and these maps do not contain $\mu$.
So the left sides are continuous functions of $S$.
The right sides do not contain $\mu$.
The bounds hold for $\mu\in(0,1]$ (2.3); therefore they hold in the limit $\mu\to0$.
\end{proof}

\noindent\textbf{(v.4) The Remark ``Where the mass enters'': the mass in the operators and in the bounds, and its reality.}

\begin{proof}
That the bound $\mu\le1$ enters $C_{T}$ through the summand $1$ of its inner bracket, and $C_{d}$ through the summand $1$ of (I.L0.12), is shown in (3) and (4) below.
That it enters $K_{1}$ and $k_{1}$ is seen at one step of the proof of Proposition~1.6 and at one step of the proof of Proposition~2.3\,(b); these steps are recalled in (6).
That differences and commutators of $F_{y}$ contain no mass is (2).
That the relations with conjugation and with adjoints rest on the reality of the mass is (5) for the operators $F_{y}(A)$ on $T_k$, and (7) for the other operators.
That $\mathbf{m}>0$ is a hypothesis of Proposition~1.2\,(S7) and of Proposition~1.5\,(d) is read off from the data of these two clauses.
\begin{itemize}
\item[(1)] The operators $F_{y}(A;\mu)=F_{y}(A;0)+\mu$ and $D_0(A)=D_A+\mu$ are affine in $\mu$ with the coefficient $1$.
The operators
\[
T_k(A)=b_{k}^{-1}Q_k(A)J^{*}D_A+\left(1+\frac{\mu}{b_{k}}\right)Q_k(A)J^{*}
\]
(Section~I.3.5 with (I.Q1)) and $T^{(1)}(A)$ ((I.5.4)) are affine in $\mu$ with the coefficient $b_{k}^{-1}Q_k(A)J^{*}$ (for $T^{(1)}(A)$: its rows at the unit sites of $\mathbf{Y}$).
The operators $D_A$, $J$, $J^{*}$, $Q_k(A)J^{*}$, $JQ_k^{T}(-A)$, $Q(A)$, $Q^{T}(-A)$, $P(A)$, $JP_k(A)J^{*}$, $JP_{k+1}(A)J^{*}$, $JP(A)J^{*}$, $\operatorname{Pen}_{y}(A)$, $B_{k,y}(A)$, $\delta_{\alpha,A}$, $T^{(2)}(A)$ and the multiplications by functions of $c$ contain no mass ((I.2.3), Sections~I.3.1, I.3.3, I.3.5, I.3.9, (I.1.7), (I.5.4)).
\item[(2)] The multiple $\mu\cdot1$ cancels in differences, and it commutes with the multiplication by a function $\psi$ of $c$.
Hence, on every block torus, the operators
\begin{gather*}
F_{y}(B;\mu)-F_{y}(0;\mu)=(D_B-D_0)+\left(\operatorname{Pen}_{y}(B)-\operatorname{Pen}_{y}(0)\right),\qquad\left[F_{y}(A;\mu),\psi\right]=[D_A,\psi]+\left[\operatorname{Pen}_{y}(A),\psi\right],\\
\mathrm{Cm}_{\square,y}(A)=-\left[F_{y}(A;\mu),h_{\square}\right]\quad\text{(Section~G1.7)}
\end{gather*}
contain no mass.
\item[(3)] In Proposition~9\,(c) the mass occurs in the term $b_{-}^{-1}\cdot\mu$, which bounds $\frac{\mu}{b_{k}}$ (as $b_{k}\ge b_{-}$).
For $\mathfrak{i}\le1$, $\mathrm{qd}\le1$ and $\mu\le1$ one has
\begin{align*}
&\exp\left(\frac{3\mathfrak{i}}{2}\right)\left[b_{-}^{-1}\left(54L\cdot\mathrm{qd}\cdot\exp(6\,\mathrm{qd})+12+6\exp(\mathfrak{i})+\mu\right)+1\right]\\
&\qquad\le\mathrm{e}^{\frac{3}{2}}\left[b_{-}^{-1}\left(54L\cdot\mathrm{e}^{6}+12+6\mathrm{e}+1\right)+1\right]\le\mathrm{e}^{\frac{3}{2}}\left[b_{-}^{-1}\left(54L\cdot\mathrm{e}^{6}+12\mathrm{e}+1\right)+1\right]=\frac{C_{T}}{2},
\end{align*}
because $12+6\mathrm{e}\le12\mathrm{e}$ (as $\mathrm{e}\ge2$).
So the summand $1$ in the inner bracket of $C_{T}$ is the bound $\mu\le1$.
\item[(4)] ($k=0$, $\eta_0=1$, $A\in\tilde{G}_0$.)
For an operator $O$ on $\mathfrak{S}^{\mathrm{ch}}(1,\ell)$ the row sum
\[
\sup_{x}\sup_{\beta}\sum_{\beta'}\int\mathrm{d}x'\,\left|O(x,x')_{\beta\beta'}\right|
\]
is its norm for $|\cdot|_{\infty}$ ((I.ke3)).
The operator $F_0(A)=D_A+b_{L}JP(A)J^{*}+\mu$ has three parts.

\emph{The part $D_A$.}
By (I.2.3) and (I.m1),
\[
(D_Af)(x)=\frac{\rho}{2}\cdot\mathrm{sw}\cdot f(x)-\sum_{i=0}^{3}M_{i}(x)\cdot U_A(x,x^{(i)})\cdot f(x^{(i)}),
\]
with $U_A(x,x^{(i)})=1$ for $i\le2$ and $|U_A(x,x^{(3)})|=|u_A(c,c_{3})|\le\exp(\mathfrak{i}(A))$ (\cite[Lemma~16.2(a)]{QEDI}).
By (I.m7) and $\rho<1$,
\[
\frac{\rho}{2}\cdot\mathrm{sw}\le6+6\sqrt{2},\qquad|M_{0}(x)|_{c}\le6,\qquad|M_{1}(x)|_{c}+|M_{2}(x)|_{c}\le6\sqrt{2}\,(1+\sqrt{2})=12+6\sqrt{2},\qquad|M_{3}(x)|_{c}\le6.
\]
Hence
\[
|D_Af|_{\infty}\le\left(24+12\sqrt{2}+6\exp(\mathfrak{i}(A))\right)\cdot|f|_{\infty}\le58\exp(\mathfrak{i}(A))\cdot|f|_{\infty},
\]
because $24+12\sqrt{2}\le41$ and $\exp(\mathfrak{i}(A))\ge1$.

\emph{The part $b_{L}JP(A)J^{*}$.}
For $c\in B(y')$, $y'\in L\mathbb{Z}^3$ (Sections~I.3.1 and I.3.3),
\[
\left(JP(A)J^{*}f\right)(c,g)=U_A(y',c)^{-1}\cdot L^{-3}\sum_{c'\in B(y')}U_A(y',c')\cdot\frac{1}{48}\sum_{g'\in O_{\mathrm{cub}}}f(c',g').
\]
Since $|y'-c|,|y'-c'|\le\frac{L-1}{2}$, one has $\left|U_A(y',c)^{-1}U_A(y',c')\right|\le\exp(3L\cdot\mathfrak{i}(A))$ (\cite[Lemma~16.2(a)]{QEDI}), and $3L\cdot\mathfrak{i}(A)<3L\cdot e_k^{\varepsilon}\le\frac{1}{2}$ (the display (G1) of \cite[Section~16.2.7]{QEDI} and (Z2.s1)).
The sum over $c'$ has $L^3$ terms, and the mean over $g'$ has $48$ terms of the weight $\frac{1}{48}$.
So $\left|b_{L}JP(A)J^{*}f\right|_{\infty}\le b_{L}\cdot\mathrm{e}\cdot|f|_{\infty}\le b\cdot\mathrm{e}\cdot|f|_{\infty}$.

\emph{The part $\mu\cdot1$.}
Its kernel (I.1.4) is $48\mu\cdot\mathbf{1}[x=x']\cdot1_4$ ((I.ke1)), and its row sum is
\[
\int\mathrm{d}x'\,48\mu\cdot\mathbf{1}[x=x']=48\mu\cdot\frac{1}{48}=\mu.
\]

Altogether the row sum of $F_0(A)$ is at most $58\exp(\mathfrak{i}(A))+\mu+b\cdot\mathrm{e}$, where the first and the third summand bound operators that contain no mass.
At $\mu\le1$ this is the bound (I.L0.12).
Its summand $1$ is the first $1$ in
\[
2+58\mathrm{e}+b\cdot\mathrm{e}=(58\mathrm{e}+1+b\cdot\mathrm{e})+1,
\]
the second entry of the maximum in $C_{d}$.
\item[(5)] (The mass is real.)
Let $\mathrm{tc}$ be a complex number, and put $F_{y}(A;\mu+\mathrm{tc}):=F_{y}(A;\mu)+\mathrm{tc}$.
The operator $\mathrm{tc}\cdot1$ commutes with $\mathsf{C}$, $\varkappa$ and $\gamma_5$; moreover $(\mathrm{tc}\cdot1)^{T}=\mathrm{tc}\cdot1$ for $\langle\cdot,\cdot\rangle^{\mathrm{ch}}$, $\overline{\mathrm{tc}\cdot1}=\overline{\mathrm{tc}}\cdot1$ (here $\overline{O}$ is the operator $f\mapsto\overline{O\overline{f}}$, (I.O5)), and $(\mathrm{tc}\cdot1)^{*}=\overline{\mathrm{tc}}\cdot1$ for $(\cdot,\cdot)^{\mathrm{ch}}$.
The connections $-A$ and $\overline{A}$ lie in $\tilde{G}_k$ (the display (G2) of \cite[Section~16.2.7]{QEDI}).
Therefore Proposition~10\,(c), (d), (f) (for $k\ge1$) and Proposition~10\,(g) (for $k=0$) give, for $y\in Y_{k}$ (with $-\infty:=\infty$, Section~G1.9) and $A\in\tilde{G}_k$, with $A$ real in the third relation,
\begin{gather*}
F_{y}(A;\mu+\mathrm{tc})^{T}=\mathsf{C}^{-1}F_{-y}(-A;\mu+\mathrm{tc})\,\mathsf{C},\qquad\overline{F_{y}(\overline{A};\mu+\mathrm{tc})}=\varkappa^{-1}F_{y}(-A;\mu+\overline{\mathrm{tc}})\,\varkappa,\\
F_{y}(A;\mu+\mathrm{tc})^{*}=\gamma_5\,F_{-y}(A;\mu+\overline{\mathrm{tc}})\,\gamma_5.
\end{gather*}
The relation of transposition holds with every complex number in the place of the mass.
The relations of conjugation and of the adjoint hold with the same number on both sides if and only if $\overline{\mathrm{tc}}=\mathrm{tc}$.
\item[(6)] ($K_{1}$ and $k_{1}$.)
The steps at which the bound $\mu\le1$ enters $K_{1}$ and $k_{1}$ are the following.

\emph{The constant $K_{1}$.}
The step belongs to the proof of Lemma~6.2 in the proof of Proposition~1.6; this lemma states $\|u\|_{\mathrm{L2}}\le L\cdot C_{I}\cdot\|F_{y}(0)u\|_{\mathrm{L2}}$ for $u\in\mathrm{HH}(\eta_k,R;\theta)$ and $\mu\in[0,1]$.
There $F_{y}(0)=D+\mu+\operatorname{Pen}$, with the operator $D$ of Section~G1.4 and
\[
\operatorname{Pen}:=\boldsymbol{\alpha}_{y}\,JP_k(0)J^{*}+\boldsymbol{\beta}_{y}\,JP_{k+1}(0)J^{*},
\]
and $\mathrm{Xr}:=\operatorname{Re}\,(u,F_{y}(0)u)^{\mathrm{ch}}$.
In that proof $\mathrm{Xr}$ is a sum of three non-negative terms, one of which is $\mu\|u\|_{\mathrm{L2}}^{2}$, and the steps (ii), (iii) read
\[
\|\operatorname{Pen}u\|_{\mathrm{L2}}^{2}\le b\cdot\mathrm{Xr},\qquad\|\mu u\|_{\mathrm{L2}}^{2}=\mu\left(\mu\|u\|_{\mathrm{L2}}^{2}\right)\le\mu\cdot\mathrm{Xr}\le\mathrm{Xr}\quad\text{because }\mu\le1.
\]
Together with $Du=F_{y}(0)u-\mu u-\operatorname{Pen}u$ they give
\[
\|Du\|_{\mathrm{L2}}\le\|F_{y}(0)u\|_{\mathrm{L2}}+\left(1+b^{\frac{1}{2}}\right)\mathrm{Xr}^{\frac{1}{2}},\qquad1+b^{\frac{1}{2}}=K_{1}\cdot L^{-\frac{1}{2}}.
\]
So the $1$ in $1+b^{\frac{1}{2}}$, that is, the summand $L^{\frac{1}{2}}$ of $K_{1}=L^{\frac{1}{2}}+(L\cdot b)^{\frac{1}{2}}$, is the bound $\mu\le1$ of step (iii).
By the formulas of Proposition~14.2\,(a), $K_{1}$ occurs in $C_{I}$; $C_{I}$ in $a_{0}$ and in $C_{3}$; $a_{0}$ in $\gamma_{(5)}$; $C_{3}$ in $C_{5}$; $\gamma_{(5)}$ in $\gamma_{z}$, hence in $\gamma^{F}$; $C_{5}$ in $C_{z}$; and $C_{z}$ and $\gamma^{F}$ in $a'$.

\emph{The constant $k_{1}$.}
The step is Step~3 of the proof of Proposition~2.3\,(b) ($k=0$, a block torus, $\mu\in[0,1]$, $y$ real).
It uses the function $\mathrm{wf}$, the operators $\mathrm{Fw}$, $\mathrm{Rw}$ and the bound $\|\mathrm{Rw}\|_{\mathrm{L2}}\le\frac{c_{\mathrm{re}}}{2L}$ of Step~2 of that proof;
Step~2 does not use the bound $\mu\le1$, and since $\mu$ commutes with $\exp(\mathrm{wf})$, neither $\mathrm{Rw}$ nor its bound contains the mass.
There $\mathrm{Fw}:=\exp(\mathrm{wf})\,F_{y}(0)\exp(-\mathrm{wf})$ with a real function $\mathrm{wf}$ of $c$ such that $|\mathrm{wf}(c)-\mathrm{wf}(\tilde{c})|\le a_{z}\cdot|c-\tilde{c}|_{1}$, and $\mathrm{Rw}:=\mathrm{Fw}-F_{y}(0)$; thus
\[
\mathrm{Fw}-i\gamma_5y=D_0+\mu+b_{L}JP(0)J^{*}+\mathrm{Rw},
\]
where $D_0$ is $D_A$ at $A=0$.
One has $\|D_0\|_{\mathrm{L2}}<58$, $\|JP(0)J^{*}\|_{\mathrm{L2}}\le1$, $\|\mathrm{Rw}\|_{\mathrm{L2}}\le\frac{c_{\mathrm{re}}}{2L}\le\frac{b}{2L}$, and $\mu\le1$; hence
\[
\|\mathrm{Fw}-i\gamma_5y\|_{\mathrm{L2}}\le58+1+\frac{b}{L}+\frac{b}{2L}\le59+\frac{3}{10}b\le60+b=k_{1}\qquad(L\ge5,\ b>0).
\]
The summand $1$ is the bound $\mu\le1$.
\item[(7)] (Conjugation and adjoints, for the other operators.)

\emph{Where the mass enters.}
By (1), (2) and the definitions, the operators and numbers that occur in the relations with conjugation and with ${}^{*}$ at the places listed in the Remark are of four kinds.
\begin{itemize}
\item[(mk1)] Operators that contain no mass: those named as such in (1) and (2), and their products; $\mathrm{rs}$, $\mathrm{ext}$ and the linear map $\mathrm{Ecb}$ (Sections~G1.5 and G1.6); the constant matrices.
\item[(mk2)] $F_{y}$ on $\mathrm{HH}(\eta_k,R;\theta)$ (Proposition~1.6) and on a block torus (Proposition~2.2; Section~G1.2), among them $T_k$ and the mirror tori (Section~G1.5); $D_0(A)$; $T_k(A)$, $T^{(1)}(A)$.
They are affine in $\mu$, and the coefficient of $\mu$ contains no mass: this is (1), together with the fact that on $\mathrm{HH}(\eta_k,R;\theta)$ and on every block torus $F_{y}$ is given by the same display $D+\mu+\operatorname{Pen}_{y}$.
\item[(mk3)] Operators formed from those of (mk1), (mk2) by finitely many sums, products and inverses:
$S_{y}=F_{y}^{-1}$, also under the names $S_k(A)$, $S_{k,y}(A)$, $S^{0}_{k+1}(A)$, $S_{0,y}(A)$, $\Gamma_0(A)$;
$H_{k}(A)$, $\mathfrak{H}_{k}(A)$, $D_k(A)$, $M_k(A)$, $\dot{M}_k(A)$, $\Gamma_{k,y}(A)$, $\Gamma_k(A)$, $\bar{H}_{k}(A)$, $\bar{T}_k(A)$ (Sections~I.3.5 and I.3.9);
$\bar{T}^{(1)}(A)$ ((I.5.4));
the operator $K$ of Proposition~12;
$\mathrm{Sloc}_{\square,y}(A)=\mathrm{rs}\,F_{y}(A^{W})^{-1}\,\mathrm{ext}$ and its image under $\mathrm{Ecb}$ (Sections~G1.5 and G1.6);
$S_{y,\mathrm{wk}}(A)$, $S'_{y,\mathrm{wk}}(A)$, $S^{\mathrm{sy}}_{y,\mathrm{wk}}(A)$ (Sections~G1.8 and G1.9: products of functions $h_{\square}$ and of operators $\mathrm{Sloc}_{\square,y}(A)$, $\mathrm{Cm}_{\square,y}(A)$, and half the sum of two such products).
The numbers $\delta Z_k(A)$, $\delta Z_0(A)$ are determinants of operators of this kind.
\item[(mk4)] The sums over walks $S^{\#}_{y}(s,A)$, $S^{\#}_{y}(X,A)$ (Section~G1.9), and the families of Propositions~7 and 8 for which W16 is stated; these families are formed from the sums over walks and from $b_{k}$, $B_{k}(A)$, $Q_k(A)J^{*}$, $JQ_k^{T}(-A)$ by sums and products ($H_0(s,A)$ contains no mass).
\end{itemize}

\emph{Rational functions of the mass.}
Let $\mathrm{tc}$ be a complex number.
Form the operators (mk2) with $\mu+\mathrm{tc}$ in the place of $\mu$, and those of (mk3), (mk4) with these.
The spaces are of finite dimension, with the exception of $\mathrm{HH}(\eta_k,\infty;\theta)$.
In bases that do not depend on $\mathrm{tc}$, the entries of the operators (mk2) are polynomials in $\mathrm{tc}$ of degree at most $1$.
The determinant of $F_{y}+\mathrm{tc}$ is a polynomial in $\mathrm{tc}$ with the leading coefficient $1$; for $A\in\tilde{G}_k$ and real $y$, the determinant of $D_k(A)+b_{L}P(A)+i\gamma_5y$ is a rational function of $\mathrm{tc}$ which does not vanish at $\mathrm{tc}=0$ (Proposition~4\,(b)).
By Cramer's rule the entries of the operators (mk3) and the numbers $\delta Z_k(A)$, $\delta Z_0(A)$ are rational functions of $\mathrm{tc}$, and the inverses exist outside a finite set of values of $\mathrm{tc}$.

\emph{The relations.}
Write a relation between objects of (mk2), (mk3) in the form $\overline{O_1}=O_2$, respectively $O_1^{*}=O_2$ (the fields, $y$, $\square$ and $\mathrm{wk}$ being fixed, and $\varkappa$, $\gamma_5$ being counted with $O_2$), and write $O_1(\mathrm{tc})$, $O_2(\mathrm{tc})$ for $O_1$, $O_2$ formed with $\mu+\mathrm{tc}$.
Conjugation and ${}^{*}$ are additive and turn a factor $\mathrm{tc}$ into $\overline{\mathrm{tc}}$.
Hence the entries of $\overline{O_1(\overline{\mathrm{tc}})}$ in bases of real functions, respectively of $O_1(\overline{\mathrm{tc}})^{*}$ in orthonormal bases, are rational functions of $\mathrm{tc}$, and so are those of $O_2(\mathrm{tc})$.
Two rational functions that coincide at infinitely many points coincide.
Therefore the relation holds at all masses $\mu+\mathrm{tc}\in(0,1]$ at which its two sides are defined if and only if
\begin{equation}\tag{mk5}
\begin{split}
&\overline{O_1(\mathrm{tc})}=O_2(\overline{\mathrm{tc}}),\quad\text{respectively}\quad O_1(\mathrm{tc})^{*}=O_2(\overline{\mathrm{tc}}),\\
&\text{for all complex numbers }\mathrm{tc}\text{ at which the two sides are defined.}
\end{split}
\end{equation}
On $\mathrm{HH}(\eta_k,\infty;\theta)$ only $F_{y}(0)$ and its inverse occur.
The operator $\left(F_{y}(0)+\mathrm{tc}\right)^{*}-\gamma_5\left(F_{-y}(0)+\overline{\mathrm{tc}}\right)\gamma_5$ does not depend on $\mathrm{tc}$, because $(\mathrm{tc}\cdot1)^{*}=\overline{\mathrm{tc}}\cdot1$ and $\gamma_5^{2}=1$; moreover $\left(O^{-1}\right)^{*}=\left(O^{*}\right)^{-1}$ for a bounded bijection $O$.
So the equivalence holds there as well.

For (mk4) one has $(\overline{s})_{\mathrm{wk}}=\overline{s_{\mathrm{wk}}}$.
So (mk5) for the terms $S^{\#}_{y,\mathrm{wk}}$ states that $\overline{(\overline{s})_{\mathrm{wk}}\cdot S^{\#}_{y,\mathrm{wk}}(\overline{A})}$, formed with $\mu+\mathrm{tc}$, equals $s_{\mathrm{wk}}\cdot\varkappa^{-1}S^{\#}_{y,\mathrm{wk}}(-A)\,\varkappa$, formed with $\mu+\overline{\mathrm{tc}}$.
Conjugation and $O\mapsto\varkappa^{-1}O\varkappa$ are continuous bijections.
Consequently, if one of the series $S^{\#}_{y}(\overline{s},\overline{A})$ (with $\mu+\mathrm{tc}$) and $S^{\#}_{y}(s,-A)$ (with $\mu+\overline{\mathrm{tc}}$) converges, then so does the other, and (mk5) holds for the sums; the same holds for $S^{\#}_{y}(X,\cdot)$.
By $\overline{O_1O_2}=\overline{O_1}\;\overline{O_2}$ and by Proposition~10\,(d) for $B_{k}(A)$, $Q_k(A)J^{*}$, $JQ_k^{T}(-A)$, the relation (mk5) extends to the families of Propositions~7 and 8.

So the relations at these places, which are stated for real masses (among them all masses in $(0,1]$), are the relations (mk5) at real $\mathrm{tc}$.
In (mk5) a complex number $\mu+\mathrm{tc}$ in the place of the mass on the left side goes with $\mu+\overline{\mathrm{tc}}$ on the right side, as in (5); the form in which these relations are stated, with one number on both sides, is the case $\overline{\mathrm{tc}}=\mathrm{tc}$.
\qedhere
\end{itemize}
\end{proof}

\subsection*{3. Conditions on the constants}

\noindent\textbf{(A) Conditions used in the arguments of this chapter.}

\medskip
\noindent\begin{tabular}{p{0.32\textwidth}p{0.18\textwidth}p{0.38\textwidth}}
\hline
Condition & Line of Section~I.8.2 or Proposition~14.2 & Where it is used\\
\hline
$0<\rho<1$ & (C0) & Lemma~1 (K1): $\hat{g}$, $c_{b}$, $c_{\mathrm{re}}>0$; (v.4)\,(4): $\rho<1$ in the numbers of (I.m7)\\
$b>0$ ($b=1$) & (C0) & Lemma~1; (v.4)\,(6)\\
$\kappa_{0}=54\log72>0$, $K_{0}=\frac{\exp(\kappa_{0})}{36}\ge1$ & (C0) & Lemma~1 (K3), (K5)\\
$\lambda_{\mathrm{int}}>0$ & (C0r); (I.O8) & Lemma~1 (K1)\\
$\kappa\ge7\kappa_{0}+7$, which enters only through $\kappa>0$ & (C2) & Lemma~1 (K3), (K5)\\
$L\ge5$ & (C3), (L5) & Lemma~1 (K1): $b_{-}>0$; (K2): $M^{\frac{1}{2}}\ge L\ge5$; Lemma~4: $\square^{\sim5}\subset\square^{\sim L}$, $(2L+1)M<13LM$; (v.4)\,(6): $\frac{b}{L}+\frac{b}{2L}\le\frac{3}{10}b$\\
$\alpha_{0}=\min\left(\frac{1}{54},\frac{1}{4}\right)>0$, $27\alpha_{0}=\frac{1}{2}$ & (C4) & Lemma~1 (K1), (K2), (K5)\\
$0<\gamma_{P}\le1$, which gives $\gamma^{P}>0$ and, with $\gamma^{F}>0$, $\gamma_{\mathrm{w}}>0$; $C_{P}<\infty$, as $C_{P}$ is a maximum of finite numbers, among them $C_{\mathrm{Hk}}(\alpha)$, which is finite for $\alpha<1$ (line (C1)) & (C4), (C4r); \cite[Section~15.7]{QEDI}, final paragraph & Lemma~1 (K5), and through (K5) 2.4, Step~4 ($\gamma_{\mathrm{w}}\ge0$ in Lemma~3) and 2.5 ($C_{d}$ finite)\\
$M\ge L^{2}$, which enters through $M^{\frac{1}{2}}\ge L$ in (K2) and through $M\ge1$ at the other places & (C5), (M1) & Lemma~1 (K2), (K5); 2.4, Step~2\\
$\exp(-\varepsilon v)\le\frac{1}{4}$ and $\exp\left(-\frac{v}{4}\right)+\exp(-\varepsilon v)\le1$ for all $v\ge-\log e_0$; the first enters only through $e_k\le1$, which it gives together with $\varepsilon>0$ & (C8), (Z2.s4), first and second inequality; (C7), (S1): $\varepsilon>1-\alpha>0$ & Lemma~4\,(c)\\
$6L\cdot\exp(-\varepsilon v)\le1$ for all $v\ge-\log e_0$ & (C8), (Z2.s1) & (v.4)\,(4) only\\
\hline
\end{tabular}

\medskip
\noindent
Lemma~4 and 2.4, Step~5 use moreover the datum $\ell\in\mathrm{Cl}$ (Section~I.1.1): $\frac{\ell}{LM}$ is an integer and $\ell\ge13LM$.

\medskip

\noindent\textbf{(B) Conditions entering through the propositions applied.}
These are the conditions under which the propositions are stated; for a proposition that is stated without conditions of its own, they are the conditions of the propositions on which its proof rests.
For each proposition the lines of Proposition~14.2\,(a) to which the conditions belong are given at the end.
\begin{description}
\item[Proposition~1.6.] (L5) ($L$ odd, $L\ge5$); $\alpha<1$ (in $C_{R}$); $0<\rho<1$ (standing data).
Lines (C3), (C1), (C0).
\item[Proposition~2.8.] (L5); (M1); (M2); $M\ge108\pi\cdot C'_{V}\cdot a'$ (part of (M5)); (Z2.s1); (Z2.s2); $\alpha<1$ (range of the line (C1); through Proposition~2.7\,(d), (e)).
The first six are the conditions on which the Propositions~2.m rest.
Lines (C3), (C5), (C8), (C1).
\item[Propositions~4, 10, 11, 12.] Those of Proposition~2.8 (for Proposition~12 also (Z2.s1), (L5), (M1), through Proposition~5).
\item[Proposition~5.] (Z2.s1); (L5), (M1); the clauses (d) and (e) require in addition the conditions of Proposition~2.8 listed above, which enter through Proposition~2.8\,(a), (c) and Proposition~4\,(b).
Lines (C8), (C3), (C5), (C1).
\item[Proposition~6.] (Z2.s1), (Z2.s3); (L5), (M1); those of Proposition~2.8; $\alpha>0$ (range of the line (C1)), $\rho>0$ (range of the line (C0)).
Lines (C8), (C3), (C5), (C1), (C0).
\item[Propositions~3, 7, 8.] Those of Proposition~3, namely those of the Propositions~2.m, (M1), (M5) in full and (M2); also $\alpha<1$ (as for Proposition~2.8).
Lines (C3), (C5), (C8), (C1).
\item[Proposition~9.] In (a): those of Proposition~2.8; (Z2.s4); (L5); (M1); (S1), which, like (Z2.s4), enters (d) through the display (G3) of \cite[Section~16.2.7]{QEDI};
the range $\rho\in(0,1)$ of the line (C0), which is used in the form $0\le\rho\le1$: it enters (c), and through (c) also (d), the barred operators included, by way of $|M_3(x)|_2=|M_3(x)|_c=3(1+\rho)\le6$ ((I.m6), (I.m7));
the range $\alpha\in(1-\bar{\varepsilon},1)$ of the line (C1), in two forms: as $\alpha\ge0$ it enters (d) through $\eta_k^{\alpha}\le1$ in the bound of $T^{(2)}$, and as $0<\alpha<1$ it is among the data of \cite[Section~16.2.1]{QEDI} under which the display (G3) stands.
The display (G3) is stated for $0<\varepsilon\le\bar{\varepsilon}$ (\cite[Section~16.2.1]{QEDI}), and (S1) gives this range together with $\alpha<1$.
Lines (C8), (C3), (C5), (C1), (C7), (C0).
\item[Proposition~13.] $\rho<1$; (L5), (M1); (Z2.s1);
the conditions of Proposition~2.8, that is, in addition (M2), the part $M\ge108\pi\cdot C'_{V}\cdot a'$ of (M5), (Z2.s2) and $\alpha<1$ (range of the line (C1)), which enter the property W1 and (d0) through Proposition~2.8\,(a), (b), enter (d0) also through Propositions~12 and 10\,(e), (g), and enter (d3) through Proposition~10\,(b), (g);
in (d4)--(d6): (Z2.s4);
in (d4), and through (d4) also in (d5) and (d6): $r\ge2$ and the first two inequalities of (A10);
in (d5) and (d6): $\varepsilon\le\frac{1}{4}$ (from (S1)) and $\alpha\ge0$ (range of the line (C1)), which enter through \cite[Lemma~17.5(a), (c)]{QEDI}, and $L(\kappa-\kappa_{0}-1)\ge\kappa$ (from (L3)) and $\kappa\ge\kappa_{0}$;
in (d6): (Z2.s5) with the constant $C_{d}$ of Proposition~13, and $\varepsilon\le\frac{1}{8}$ (from (S1));
$\kappa_{\mathrm{w}}-1\ge\kappa+1$.
Lines (C0), (C3), (C5), (C8), (C1), (C7), (C2), (C4).
\end{description}
All conditions in (A) and (B) are among those of Proposition~14.2\,(a), (b).

\section*{Chapter 14b. The ordered list of the constants and conditions; admissible constants exist: proofs of Propositions 14.2, 14.3}

\subsection*{0. Conventions of this chapter}

The entries of Proposition~14.2\,(a) are called rows; they are numbered $1,\dots,136$ in the order in which they stand there.
``Row~33'' refers to this numbering, and $L\,[33]$ indicates that the constant $L$ stands in row~33.
Thus row~1 is $\rho$, row~21 is $A_{\mathrm{ak}},A'_{\mathrm{ak}}$, row~33 is $L$, row~78 is $M$, row~116 is $\sigma$, row~117 is $e_0$, row~118 is $e_1$, row~119 is $\bar{m}$, and row~136 is $e_2$; the numbers of all rows are listed in the table of Step~3.
A row that carries several names defines them in the order in which they are written.
If a condition consists of several inequalities, ``(A8), third'' denotes the third of them in the order in which they are written.
The letter $v$ denotes a real variable, $v>0$.
For $x>0$ put $\log^{+}x:=\max\{0,\log x\}$; further $\log^{+}0:=0$ and $0^{0}:=1$.

In this chapter $\tilde{t}$ denotes a lift of an index $t$, as in Proposition~27.3, and $x$ also denotes a positive real variable, as in $\varepsilon_5(x)$.
References of the form ``Step~$n$'' and ``5.1'', \dots, ``7.4'' are to the steps of Section~1 and to their paragraphs.

\subsection*{1. Proof of Proposition 14.2}

\subsubsection*{Step 0. What is to be proved}

Proposition~14.2\,(a)--(f) consists of definitions and of the following assertions; after each of them, the place of its proof is given in parentheses.
The clauses (a)--(d) consist of definitions and of the assertions (i)--(iv); the clauses (e) and (f) are assertions throughout, and they are treated in (v) and (vi).
\begin{itemize}
\item[(i)] The entries of the rows $\rho,\dots,\sigma$ of (a), that is, of rows~1--116, other than $C_E$, $\mathrm{Csy}_{\mathrm{ak}}$, $\mathrm{Cpv}_0$ (rows~27, 28), are defined, and the following assertions made in these entries hold:
$\alpha_{0}=\frac{1}{54}$, $C_{5}\ge1$, $\gamma^{F}\le\frac{1}{6L}$, $a'\ge1$, $\gamma_{\mathrm{w}}>0$, $c_{\kappa}\le K_{0}$, and, in the entry of $c_{\rho}$, that the inequality defining $c_{\rho}$ does not contain the point $v_0$ (Step~2).
Moreover $M$ drops out of $c_{\gamma}$, and $M$ and $k$ drop out of $c_{\kappa}$, while $L$ remains through the class $\mathrm{Cl}$ (Remarks~3.1, 3.2).
For the other rows the preamble of (a) refers to other statements:
the sets in which $L$, $M$, $(r,p_0,p,\varepsilon,s)$ and $e_0$ are chosen are not empty, and $e_1$, $e_G$, $e_2$ are positive, by Proposition~14.3;
the entries $C_E$, $\mathrm{Csy}_{\mathrm{ak}}$, $\mathrm{Cpv}_0$ and the entries of the line (C10) from $D_{\mathrm{be}}$ on (rows~27, 28 and the rows from~123 on, whose references are statements that stand after the Propositions~14.m in Part~II) are the definitions of their letters.
The preamble of (a) contains the following further assertions:
the sentences on the rows ``the constant(s) of \dots'' (Step~3, rows of the kind X; Lemma~1, row~26), on the letters named after ``Depends on'' (Remark~3.4), on a letter of a later row that occurs in the property defining a constant (Step~3, paragraph ``The two letters (x3)''), and on $K_{1}$, $k_{1}$, $M^{P}$ (Step~3, after the table);
and the sentence on $c_{\mathrm{cl}}$ in the last line of (a) (Step~3, after the table).
\item[(ii)] The two consequences of (L3) stated in its row (Step~8, (f2)).
The first sentence of the paragraph that follows the conditions of (b), together with the sentences of that paragraph on the labels of Section~I.8.4 that have no entry of their own (Step~9).
The second sentence of that paragraph is the rule (G1.T-12) for the proofs of all Propositions; that the proof of each Proposition obeys it is shown in the section ``Conditions on the constants'' of its chapter (end of Step~9; Section~IV.1, paragraph~E).
This rule is no part of (o1)--(o6), and no step of the proofs of Propositions~14.2 and 14.3 rests on it.
\item[(iii)] Clause (c): its first sentence, namely that (e1)--(e8) hold for all $v\ge-\log e_1$ if the set in the definition of $e_1$ is not empty (5.1); and the fact that the conditions hold for all the data specified in (c) (5.2, 5.3).
\item[(iv)] Clause (d): the quantities with the index $*$ do not depend on $N\ge N_0(e)$ (7.1);
(Z8.w), together with $e\,h^{\frac{1}{2}}\le e_1$ and (e8), gives (W3);
(Z8.w) is contained in $e\le e_G$;
the bound on $\mathrm{pp}$;
the two consequences of (Z10.s4) (Step~6).
The sentence of (d) concerning (g2) is Proposition~35.3\,(e) in the form
\[
d_{\mathrm{sep}}\ge\frac{\ell_*}{8L}-1=\frac{M\cdot r_{*}}{8L}-1
\]
(in the setting (ST*)$^{\mathrm{ch}}$, for every admitted placement).
This inequality is an assertion of Proposition~35.3\,(e), not of Proposition~14.2.
It is proved with the sixth entry, Propositions~35.1; 35.2; 35.3, of the order of the clauses given in the paragraph ``Other Propositions'' below; that entry takes of the Propositions~14.m only parts of (o1), and the inequality belongs to none of (o1)--(o6).
No step of the proofs of Propositions~14.2 and 14.3 uses it: paragraph~7.4 uses only that (g2) and (g5) are written with $\frac{M\cdot r_{*}}{8L}-1$ and with $\varpi_{\mathrm{sep}}^{+}$, without $d_{\mathrm{sep}}$.
The sentence of (d) on $\varpi_{\mathrm{sep}}^{+}$ names the same lower bound:
$\varpi_{\mathrm{sep}}^{+}$ arises from $\varpi_{\mathrm{sep}}$ when $d_{\mathrm{sep}}$ is replaced by it, because
\[
\left(\frac{M\cdot r_{*}}{8L}-1\right)-4M-6=\frac{M\cdot r_{*}}{8L}-4M-7.
\]
The sentence of (d) on $r_{*}$, $v_{*}$, $p_{*}$, $p_{0*}$ names Proposition~35.1\,(c) for comparison only;
that these numbers do not depend on $N\ge N_0(e)$ is proved in 7.1.
\item[(v)] Clause (e) (Steps~3, 4, 7).
\item[(vi)] Clause (f) (Step~8).
\end{itemize}
The indications after ``Reference:'' and ``Used in:'',
the four sentences at the head of (d) that name the statements in which the conditions of (d) enter and the places where these conditions are named (Propositions~34.1, 34.3, 34.4, 35.5--35.10, 42.2 with the Remark that follows it, 43.1, 43.2, Theorem~44 and Section~G6.Z4),
the words of the preamble of (a) that describe the lines (C0r), (C4r), (C6$'$),
and the clauses that compare an entry with the corresponding quantity of \cite{QEDI}
refer to other statements and to other places of the list.
The lines (C0r), (C4r), (C6$'$) consist of the entries listed under these marks.
Where an inequality that a Proposition derives from a condition of (b) is stated after ``Used in:'' in the entry of that condition, it is proved in Step~8 ((f1)--(f9), (f12)).

\subsubsection*{Order of proof of the two statements}

The parts of the proofs of Propositions~14.2 and 14.3 are arranged in the following order, in which each part uses, of the two proofs, only parts that precede it.
\begin{itemize}
\item[(o1)] Step~1; Lemma~1 (rows~1--116 without 27, 28); Step~4; the lists of letters in 7.1--7.3; Step~3; (f1)--(f8), (f10), (f11); Step~9.
These parts use the definitions, the conditions of (b) that restrict the constant under consideration or an earlier one, and the results cited in them.
Of 5.1, (f7) uses only the definition of the set $\mathrm{Sadm}_1$.
\item[(o2)] Proposition~14.3\,(a)--(e).
They rest on Lemma~1, (El7) and Step~4.
\item[(o3)] Paragraph 5.1, which proves (5.1) in the case that $\mathrm{Sadm}_1$ is not empty; this is the sentence of Proposition~14.2\,(c) on (e1)--(e8).
\item[(o4)] Proposition~14.3\,(f).
It rests on Lemma~1, (El6), (El7), Step~4 and (o3).
\item[(o5)] Paragraphs 5.2, 5.3; Step~6; the last sentence of 7.1, and 7.4; (f9), (f12).
Of the parts of the two proofs, they rest on (o1) and (o3) only.
Each of their sentences that contains $e_1$ is stated under the datum $e>0$ with $e\,h^{\frac{1}{2}}\le e_1$, or is void without such an $e$; under this datum $\mathrm{Sadm}_1$ is not empty, and (o3) applies.
\item[(o6)] Proposition~14.3\,(g) and the Consequence in Section~2.
They rest on 6.1, 6.4 and (o4);
on (o1), namely on Lemma~1: in (g) on row~30 together with rows~19, 20 and Proposition~30.1\,(c), for $\varepsilon_5\left(\frac{1}{L}\right)<\infty$, and in the Consequence on the rows~$\le116$;
on Propositions~35.10, 35.3\,(c1), 36\,(g), 37.2\,(a), (d), (e); and on \cite[Theorem~21.1(a)]{QEDI}.
The Consequence rests also on (o2), that is, on Proposition~14.3\,(a)--(e), and it refers, in addition, to the statements named as references of the rows of the kind X.
\end{itemize}

\paragraph{Other Propositions.}
Of the other Propositions, the proof of Proposition~14.2 uses only Propositions~1.2\,(S3), (S6), (S7), 1.3\,(b), 1.5, 1.6, 30.1\,(c), 30.3\,(e) and 46, which stand before the Propositions~14.m (Proposition~46 in Part~I, the others in Part~II); Propositions~1.3\,(b) and 46 enter through row~26 only.
Of the later Propositions that are named as references of rows, and of the displays of their chapters named in Step~3, the proof uses the definitions only, that is, which letters the property or the formula contains (Step~3).
These are Propositions~37.1\,(a), (f), (g) and 36\,(c) (rows~27, 28; Proposition~37.1\,(a) enters only through the definition of the matrix $(\nabla^{*}\nabla)(p)$, which the display [1.10] of Chapter~37 contains);
Propositions~35.3\,(c1), 36\,(e)--(g), 37.2\,(c)--(e) and 35.10 (rows~123--125, 127);
Proposition~35.3\,(e) (the sentences of (d) concerning (g2) and $\varpi_{\mathrm{sep}}^{+}$);
and the paragraphs of Propositions~43.1 and 43.2 that state the constant, the data, and the constants in their order (rows~122, 128, 130--136).
Proposition~35.1\,(c), which (d) names for comparison in the sentence on $r_{*}$, $v_{*}$, $p_{*}$, $p_{0*}$, is not used: that these numbers do not depend on $N\ge N_0(e)$ is proved in 7.1.
Nor does the proof use the four sentences at the head of (d) that name the statements in which the conditions of (d) enter.
Of the statements and places named in them, Propositions~34.1, 34.3, 34.4, 35.5--35.9, 42.2 with the Remark that follows it, Theorem~44 and Section~G6.Z4 are not used in the proof of Proposition~14.2;
of Propositions~35.10, 43.1 and 43.2, which these sentences name as well, the proof uses only what is listed above for the rows~122, 127, 128, 130--136.
The constants of rows~1--118 other than those of rows~27, 28 are defined without the constants of rows~26--28 ($C_E$ enters the numbers $\mathrm{Csy}_{\mathrm{ak}}$, and these enter $\mathrm{Cpv}_0$), and no condition of (b) contains the constants of rows~26--28 (tables of Steps~3, 4; end of Step~3).
The proof of Proposition~14.3 uses later Propositions only in (g), namely Propositions~35.10, 35.3\,(c1), 36\,(g) and 37.2\,(a), (d), (e).
Together with these Propositions, with all Propositions on which their proofs rest, directly or through other Propositions, and with Proposition~14.4, the clauses are proved in the following order.
Every entry uses earlier entries only, with two exceptions, which are stated in the entries concerned: the second entry uses entries of Lemma~1, which belong to (o1), and the eighth entry uses Proposition~14.4\,(a), which belongs to the last entry but one; these parts are proved earlier, at the places stated in the second and in the eighth entry.
\begin{itemize}
\item Proposition~46; Propositions~1.1--1.7.
Proposition~46 uses no other Proposition, and each of Propositions~1.1--1.7 uses, of the other Propositions, only earlier ones among them.
Of the Propositions~14.m they use rows of (a) as definitions, with the ranges stated there, and the condition (L5).
\item Propositions~30.1; 30.2; 30.3\,(a)--(e), (g), and (f) without its closing part ``In the flow''.
They use Proposition~1.2 and one another, and they contain no letter of the lines (C7)--(C9).
The signs and ranges of constants on which their proofs rest are entries of Lemma~1:
those of the rows~1, 2, 10, 12, 15, 19--21;
that of row~33, which the proof of Proposition~30.2\,(b) uses in the form $L\ge1$, so that $\eta=L^{-N}\le1$, and which Proposition~30.3 uses as well;
and, in Proposition~30.3 only, those of the rows~29, 35, 36.
Lemma~1 is proved, and is used, row by row: the entry of a row rests on entries of earlier rows and on the results cited in it.
The entries of the rows just listed, and those of the earlier rows on which they rest, use no Proposition other than 1.2\,(S3), (S6), (S7); they are proved before Proposition~30.1.
In particular, row~25, whose entry uses Proposition~30.1\,(c), row~30, whose entry rests on row~25, and rows~93, 94, whose entries use the list of constants of Proposition~30.3, are not among these rows.
\item (o1); (o2); (o3); (o4); (o5).
\item The closing part ``In the flow'' of Proposition~30.3\,(f).
It is stated under the datum $e\,h^{\frac{1}{2}}\le e_1$ and uses $v_K\ge8$, $\left(1+\frac{3\,\mathrm{lg}}{v_K}\right)^{r}<L$, (2.1) and $L\ge33$, that is, (o1), (o3), (o5): the inequality (e1) at $v=v_K$ is a case of (5.1), which (o3) derives from row~118 and the datum, and which Proposition~14.3\,(f) states as well.
\item Proposition~37.1.
Its proof uses no other Proposition.
\item Propositions~35.1; 35.2; 35.3.
Of the Propositions~14.m they use three things.
First, rows of Proposition~14.2\,(a) as definitions, with the ranges stated there: $b=1$, $\alpha<1$, $r\ge2$, and the definition of $e_1$ in row~118.
Second, the conditions (L1), (L5), (M1), (S1) and the rows (e1), (e4) of Proposition~14.2\,(b).
Third, the line $L\ge2400^8$ of Proposition~14.2\,(f), used as $L\ge13$ for the class $\mathrm{Cl}$.
All of these belong to (o1).
That $e_1>0$, and that (e1)--(e8) hold for all $v\ge-\log e_1$, follows from the definition of $e_1$ and the datum $e>0$ with $e\,h^{\frac{1}{2}}\le e_1$.
\item Propositions~36; 37.2.
They rest on Propositions~46, 37.1 and 35.3\,(b1).
\item Propositions~2.1--2.8, 3--8, 15, 16.1, 9--13, in this order, which is that of Part~II.
Of the Propositions~14.m their proofs use the following.
First, rows of Proposition~14.2\,(a), (b), as definitions and as conditions on the choice of the constants, and lines of Proposition~14.2\,(f); these are proved in (o1) and (o5).
For instance, the proofs of Propositions~5 and 8 take the definitions of $C_{1}$ and $C^{F}$ from the ordered list of Proposition~14.2.
Second, through Section~I.8.1, the proofs of Propositions~3 and 8 use Proposition~14.4\,(a) for the constants in their bounds.
This part of Proposition~14.4 belongs to the last entry but one; it rests on the first two entries and on (o1)--(o5) only, and it is proved before the present entry (Section~IV.1, paragraph ``Sentences of Part~I behind which a statement stands'', (6)).
Apart from these, each of them uses only the Propositions~1.m and earlier Propositions of this entry (Proposition~13 uses 16.1\,(a1), (d1)--(d4); Propositions~15 and 16.1 use no other Proposition), and none of their proofs uses one of the Propositions~30.m, 35.m, 36, 37.m.
The closing sentence of Proposition~2.8 names Proposition~14.1\,(i) only as the place where the clauses (a)--(c) of Proposition~2.8 are applied;
the inequalities $64a'\le608a'\le C^{F}$, which the proof of that sentence uses, follow from the entries $C_{1}$, $C^{F}$ of Proposition~14.2\,(a) and from (M1), (L5) (Chapter~14a, Lemma~1\,(K2)), and they rest on no Proposition.
\item Proposition~14.1.
It stands before Proposition~14.2 and uses Propositions~1.6, 2.8, 3\,(a) and 4--13; of Proposition~14.2 it uses the lists (a), (b) as definitions.
The Remark ``The projected families'' that follows Proposition~14.1 belongs neither to that Proposition nor to this entry.
It names the statements in which FE7$^{\mathrm{ch}}$ is applied to the projected families (Section~IV.1, paragraph~D).
Its assertions follow from clause (iii) of Proposition~14.1 as soon as the Propositions named in it, Propositions~26.1, 26.3\,(b), 26.5 and 26.6\,(b), (e), are proved;
these stand after Proposition~14.1, and of Proposition~14.1 they use only clause (ii) and facts on the constants.
No entry of this order uses the Remark: of Proposition~14.1, Proposition~35.10 takes only the facts on the constants of (G5.T-1).
\item Proposition~35.10.
It rests on Propositions~35.1\,(a), (b), (c), (f), 35.3\,(c1), 14.3\,(f), 14.2\,(a), (f) and 14.1, that is, on (o1), (o4), (o5) and the entries above.
\item (o6).
\item Proposition~14.4\,(b), (c), (d), and (a) for all letters other than $C_E$, $\mathrm{Csy}_{\mathrm{ak}}$, $\mathrm{Cpv}_0$.
They rest on Propositions~14.2, 14.3\,(c), (f), 1.2, 1.5, 30.1 and 30.3, that is, on (o1)--(o5).
\item Proposition~14.4\,(a) for the three letters $C_E$, $\mathrm{Csy}_{\mathrm{ak}}$, $\mathrm{Cpv}_0$.
It comes after Propositions~37.1\,(f), (g) and 36\,(c).
\end{itemize}

\subsubsection*{Step 1. Elementary inequalities}

\begin{itemize}
\item[(El1)] $2<\mathrm{e}<3$ and $3<\pi<4$.
Indeed, $\mathrm{e}=\sum_{j\ge0}\frac{1}{j!}>1+1+\frac{1}{2}$; and $\frac{1}{j!}\le2^{-(j-1)}$ for $j\ge1$, with strict inequality for $j\ge3$, gives $\mathrm{e}<1+1+\frac{1}{2}+\frac{1}{4}+\dots=3$.
The regular hexagon inscribed in the unit circle has the perimeter $6<2\pi$, and the square circumscribed about the unit circle has the perimeter $8>2\pi$.
\item[(El2)] $\log x\le x^{\frac{1}{2}}$ for $x\ge1$.
The function $x\mapsto x^{\frac{1}{2}}-\log x$ has the derivative $\frac{x^{\frac{1}{2}}-2}{2x}$; so on $[1,\infty)$ its minimum is its value at $x=4$, which is $2-2\log2>0$, because $2<\mathrm{e}$.
\item[(El3)] For real numbers $\mathrm{rt}>0$, $\mathrm{dg}>0$, $v>0$,
\[
v^{\mathrm{dg}}\cdot\exp(-\mathrm{rt}\cdot v)\le\left(\frac{\mathrm{dg}}{\mathrm{rt}\cdot\mathrm{e}}\right)^{\mathrm{dg}},
\qquad\text{hence}\qquad
v^{\mathrm{dg}}\cdot\exp(-\mathrm{rt}\cdot v)\le\left(\frac{2\,\mathrm{dg}}{\mathrm{rt}\cdot\mathrm{e}}\right)^{\mathrm{dg}}\exp\left(-\frac{\mathrm{rt}\cdot v}{2}\right).
\]
The function $v\mapsto\mathrm{dg}\cdot\log v-\mathrm{rt}\cdot v$ has the derivative $\frac{\mathrm{dg}}{v}-\mathrm{rt}$; its maximum, attained at $v=\frac{\mathrm{dg}}{\mathrm{rt}}$, is $\mathrm{dg}\cdot\log\frac{\mathrm{dg}}{\mathrm{rt}}-\mathrm{dg}$.
The second inequality is the first one with $\frac{\mathrm{rt}}{2}$ in the place of $\mathrm{rt}$, multiplied by $\exp\left(-\frac{\mathrm{rt}\cdot v}{2}\right)$.
\item[(El4)] $\log2\ge0.6927$.
Indeed,
\[
\log2=-\log\left(1-\frac{1}{2}\right)=\sum_{j\ge1}\frac{1}{j\cdot2^{j}}\ge\frac{1}{2}+\frac{1}{8}+\frac{1}{24}+\frac{1}{64}+\frac{1}{160}+\frac{1}{384}+\frac{1}{896}+\frac{1}{2048},
\]
and the right side is at least $0.5+0.125+0.041666+0.015625+0.00625+0.002604+0.001116+0.000488=0.692749$.
\item[(El5)] For $\gamma>0$ one has $1\le c(\gamma)\le\left(1+\frac{6}{\gamma}\right)^{3}$.
On the torus $\mathbb{Z}^3/\ell\mathbb{Z}^3$, $\ell\ge1$ an integer, every site $y$ satisfies
\[
\sum_{y'}\exp\left(-\gamma\,d(y,y')\right)\le c(\gamma).
\]
The term $z=0$ of $c(\gamma)$ equals $1$.
The inequality $|z|\ge\frac{|z|_1}{3}$ gives $c(\gamma)\le\left(\sum_{j\in\mathbb{Z}}\exp\left(-\frac{\gamma|j|}{3}\right)\right)^{3}=\left(\coth\frac{\gamma}{6}\right)^{3}$, and $\coth x=1+\frac{2}{\exp(2x)-1}\le1+\frac{1}{x}$ for $x>0$.
Every class $y'$ has a lift $n$ with $|n-y|=d(y,y')$, and different classes have different lifts.
\item[(El6)] $(1+x)^{\mathrm{dg}}\le\exp(\mathrm{dg}\cdot x)$ for $x\ge0$ and $\mathrm{dg}\ge0$, because $\log(1+x)\le x$.
\item[(El7)] Let $\mathrm{nt}\ge1$ be an integer and $\mathrm{tg}>0$, and for $i=1,\dots,\mathrm{nt}$ let $\mathrm{Cf}_i\in[0,\infty)$, $\mathrm{dg}_i\in[0,\infty)$, $\mathrm{rt}_i\in(0,\infty)$.
Put
\[
\mathrm{Vbd}\left((\mathrm{Cf}_i;\mathrm{dg}_i;\mathrm{rt}_i)_{i\le\mathrm{nt}};\mathrm{tg}\right):=\max_{i\le\mathrm{nt}}\frac{2}{\mathrm{rt}_i}\log^{+}\left(\frac{\mathrm{nt}\cdot\mathrm{Cf}_i}{\mathrm{tg}}\left(\frac{2\,\mathrm{dg}_i}{\mathrm{rt}_i\cdot\mathrm{e}}\right)^{\mathrm{dg}_i}\right)\qquad(<\infty).
\]
Then
\[
\sum_{i\le\mathrm{nt}}\mathrm{Cf}_i\cdot v^{\mathrm{dg}_i}\cdot\exp(-\mathrm{rt}_i\cdot v)\le\mathrm{tg}
\]
for \emph{all} $v>0$ with $v\ge\mathrm{Vbd}(\ldots;\mathrm{tg})$.
Indeed, by (El3) if $\mathrm{dg}_i>0$, and by $\exp(-\mathrm{rt}_i\cdot v)\le\exp\left(-\frac{\mathrm{rt}_i\cdot v}{2}\right)$ if $\mathrm{dg}_i=0$,
\[
\mathrm{Cf}_i\cdot v^{\mathrm{dg}_i}\cdot\exp(-\mathrm{rt}_i\cdot v)\le\mathrm{Cf}_i\left(\frac{2\,\mathrm{dg}_i}{\mathrm{rt}_i\cdot\mathrm{e}}\right)^{\mathrm{dg}_i}\exp\left(-\frac{\mathrm{rt}_i\cdot v}{2}\right),
\]
and the right side is at most $\frac{\mathrm{tg}}{\mathrm{nt}}$ as soon as
\[
\frac{\mathrm{rt}_i\cdot v}{2}\ge\log^{+}\left(\frac{\mathrm{nt}\cdot\mathrm{Cf}_i}{\mathrm{tg}}\left(\frac{2\,\mathrm{dg}_i}{\mathrm{rt}_i\cdot\mathrm{e}}\right)^{\mathrm{dg}_i}\right).
\]
\end{itemize}

\subsubsection*{Step 2. The entries of the rows 1--116 are defined; their ranges}

\begin{named}{Lemma 1}
Let one of the rows $1,\dots,116$ other than $27$, $28$ be given, and let the constants of this row and of all earlier rows be chosen as (a) prescribes: $L\,[33]$ with (L1)--(L5), $M\,[78]$ with (M1)--(M6), $\varepsilon\,[113]$ with (S1), $s\,[114]$ with (S2).
Then the entry of the row is defined and lies in the set named for it in the following list.
(A number ``in $(0,\infty)$'' is positive and finite.)
\end{named}

\paragraph{Rows 1--9.}
$\rho\in(0,1)$; $w_i>0$ and $\mathrm{sw}>0$; $\bar{\varepsilon}=\frac{1}{200}$; $c_{X}=54$.
Further $\kappa_{0}=54\log72\ge54=c_{X}$, because $72>3>\mathrm{e}$ (El1).
$K_{0}=\frac{\exp(54\log72)}{36}=\frac{72^{54}}{36}\ge2$.
$c_{\mathrm{rb}}=54\,K_{0}\ge1$.
$c_{\mathrm{fl}}=61\cdot54\cdot\frac{72^{54}}{36}\ge72^{54}\ge1$.
$b=1$.

\paragraph{Row 10.}
Let $\mathrm{Lcl}:=\sum_{i\le2}w_i(1-R_i)$ on $\ell^2(O_{\mathrm{cub}})$, where $(R_iu)(g)=u(g\,s_i)$ (I.O8).
Since $s_i^2=1$ (Section~I.1.2), the map $g\mapsto g\,s_i$ is an involution of the $48$ elements of $O_{\mathrm{cub}}$; so $R_i$ is a symmetric permutation matrix with $R_i^2=1$, its eigenvalues are $\pm1$, and $1-R_i\ge0$.
Hence $\mathrm{Lcl}$ is real symmetric, and $0\le\mathrm{Lcl}\le2(w_0+w_1+w_2)$.
For the indicator $u$ of the unit element and the counting scalar product one has $(u,u)=1$ and $(u,R_iu)=u(s_i)=0$, since $s_i\ne1$; so $(u,\mathrm{Lcl}\,u)=w_0+w_1+w_2>0$, and $\mathrm{Lcl}\ne0$.
A symmetric matrix that is $\ge0$ and $\ne0$ has a positive eigenvalue.
So the least non-zero eigenvalue exists, and
\[
\lambda_{\mathrm{int}}\in\left(0,2(w_0+w_1+w_2)\right].
\]

\paragraph{Row 11.}
$36\sqrt{2}<198$, so the square root is real; $\left(6+9\sqrt{2}\right)^2=198+108\sqrt{2}>198-36\sqrt{2}$; hence $\lambda_1\in(0,\infty)$.

\paragraph{Rows 12--18.}
$\hat{g}=\frac{\rho\cdot\lambda_{\mathrm{int}}}{2}\in(0,\infty)$.
$c_{\mathrm{hl}}=\sqrt{3}(1+\rho)\in\left(\sqrt{3},2\sqrt{3}\right)$.
$c_{*}=\frac{\min(6,\lambda_{\mathrm{int}})}{18}\in\left(0,\frac{1}{3}\right]$.
The function $x\mapsto3(1+\rho)\left(\exp(x)-1\right)$ is a strictly increasing bijection of $[0,\infty)$ onto itself, and $\frac{\hat{g}}{2}>0$; so $\theta_{0}$ exists, is unique, and
\[
\theta_{0}=\log\left(1+\frac{\hat{g}}{6(1+\rho)}\right)\in(0,\infty).
\]
$K_{\rho}=\rho+\frac{2c_{\mathrm{hl}}^2}{\hat{g}}\in(0,\infty)$.
$t_{0}=\frac{1}{4K_{\rho}^2}\in(0,\infty)$.
$c_{\mathrm{el}}\in(0,0.22]$.

\paragraph{Rows 19, 20.}
$K_{\mathrm{Str}}<\infty$ by Proposition~1.2\,(S3).
By the last sentence of (S3), at $p=0$, where $\sin0=0$ and $\mathrm{tw}(0)=0$, one has $\mathrm{Dsc}_{\mathbf{m}}(0)=\mathbf{m}\cdot1$; at $\mathbf{m}=1$ its norm is $1$.
So $K_{\mathrm{Str}}\in[1,\infty)$, and
\[
\kappa_{(0)}\in\left(0,\frac{\theta_{0}}{2}\right].
\]

\paragraph{Row 21.}
For every multi-index $\mathrm{ak}$, Proposition~1.2\,(S6), (S7) assert that there are numbers $A_{\mathrm{ak}}$, $A'_{\mathrm{ak}}$ with the bounds stated there.
In both bounds the left side is a norm, hence $\ge0$, and the factor beside the constant, $\mathrm{rp}(p)^{-1-|\mathrm{ak}|}$ resp.\ $\mathrm{rp}(p')^{-1-|\mathrm{ak}|}$, is positive and finite ($p\ne0$ in (S6), $\mathbf{m}>0$ in (S7)).
So $A_{\mathrm{ak}},A'_{\mathrm{ak}}\in[0,\infty)$.

\paragraph{Rows 22, 23.}
$A_1\in(0,\infty)$, because $\rho$, $c_{*}$, $\hat{g}$, $K_{\rho}$ are positive and finite; $A_2=\frac{1+c_{\mathrm{hl}}\cdot A_1}{\hat{g}}\in(0,\infty)$; so $c_{b}\in(0,\infty)$.
$C_{\mathrm{Po}}=\frac{1}{8c_{*}}+\frac{1}{\lambda_{\mathrm{int}}}\in(0,\infty)$.

\paragraph{Rows 24--26.}
Proposition~1.5 asserts a number $C_{G}$ with (a)--(f); in the first bound of its clause (b) the left side is a norm and the factor $(1+|n|)^{-2}$ is positive; so $C_{G}\in[0,\infty)$.
Further $\hat{g},\theta_{0}\in(0,\infty)$ (rows~12, 15), $K_{\mathrm{Str}}\in[1,\infty)$ (row~19), $A'_{\mathrm{ak}}\in[0,\infty)$ (row~21), and $\mathrm{cF13}=c(1,3)$ is a finite number (\cite[Lemma~12.4]{QEDI} with $(\mathrm{kk},N)=(1,3)$, where $\mathrm{kk}$ is the number written $\mathsf{k}$ there; by the list of constants of Proposition~30.1 it is an absolute number).
So $C_{\mathrm{lH}}$, $C_{\mathrm{lS}}$ and
\[
C_{6}=26\left[\mathrm{cF13}\cdot C_{\mathrm{lS}}+\frac{16\,C_{\mathrm{lH}}}{(\mathrm{e}\cdot\theta_{0})^2}\right]
\]
are defined and finite (list of constants of Proposition~30.1).
In the display (8.1) of Proposition~30.1\,(c) the left side is a modulus, and
\[
\log\frac{1}{1-\exp\left(-\kappa_{(0)}\cdot\mathbf{m}\cdot n\right)}>0,
\]
because $\kappa_{(0)}>0$ (row~20); so $C_{6}\in[0,\infty)$.

$c_{\rho}\in(0,\infty)$.
Let $X$ be the set of the fermion sites of \cite[Definition~2.4]{QEDII}, $D^{\mathrm{ch}}$ the chamber operator of \cite[Definition~2.10]{QEDII} with the standard weights, at $A=0$ and with the lattice mass $0$, and $\nabla^{*}\nabla$ the weighted Laplacian of \cite[Definition~3.17]{QEDII} at $A=0$.
Consider the set of the real numbers $\mathrm{cf}$ for which, at the given $\rho$, the inequality
\[
(D^{\mathrm{ch}})^{*}D^{\mathrm{ch}}\ge\mathrm{cf}\cdot\nabla^{*}\nabla
\]
holds on $\ell^2(X;\mathbb{C}^4)$ and on the spinors of every pair $(\Gamma,\chi)$, $\Gamma$ a subgroup of $\mathbb{Z}^3$ of rank three and $\chi$ a homomorphism of $\Gamma$ into the unit circle (the spaces of \cite[Proposition~3.18]{QEDII}).
This set is the set of the numbers with which the inequality of row~26 holds:
by Sections~I.1.1 and I.1.3 and by Proposition~46\,(a), $\mathrm{geo}$ is a bijection of $\mathbb{L}^{\mathrm{f}}_0=\mathbb{Z}^3\times O_{\mathrm{cub}}$ onto $X$, the point $v_0$ of \cite[Definition~2.4]{QEDII} being the point $v_0$ of Section~I.1.3, and $\mathrm{geo}(c+a,g)=\mathrm{geo}(c,g)+a$ for $a\in\mathbb{Z}^3$, so that $\psi\mapsto\psi\circ\mathrm{geo}$ maps the spinors of a pair $(\Gamma,\chi)$ on $X$ onto those on $\mathbb{L}^{\mathrm{f}}_0$;
transported to $\mathbb{L}^{\mathrm{f}}_0$ by $\mathrm{geo}$, the operator $D^{\mathrm{ch}}$ becomes the operator (I.2.3) at the level $0$ and at $A=0$, which contains no mass, that is, the operator $D_0$ of row~26, and $\nabla^{*}\nabla$ becomes $\sum_iw_i(1-R_i)$ (part (a) of the following Lemma and its proof);
and the replacement of a scalar product by a positive multiple of it leaves the adjoint of an operator unchanged and multiplies both quadratic forms of the inequality by the same factor.
By part (a) of the following Lemma, $D^{\mathrm{ch}}$ and $\nabla^{*}\nabla$ do not depend on the point $v_0$ of \cite[Definition~2.4]{QEDII}; so the inequality does not contain $v_0$, and the set is the same for every $v_0$, as row~26 asserts.
By part (b) of the following Lemma, the set contains the positive number $c_{b}$ of row~22.
It is bounded above: let $u\in\ell^2(X;\mathbb{C}^4)$ be the spinor whose value is a unit vector of $\mathbb{C}^4$ at one site and $0$ elsewhere, and let $(\cdot,\cdot)$ be the sesquilinear scalar product $\sum_x\overline{u'(x)}\cdot u(x)$; then $(u,\nabla^{*}\nabla u)=\mathrm{sw}>0$, because $\nabla^{*}\nabla=\sum_iw_i(1-R_i)$ with $(R_iu)(x)=u(x^{(i)})$ and $x^{(i)}\ne x$; hence $\mathrm{cf}\le\frac{\|D^{\mathrm{ch}}u\|^2}{\mathrm{sw}}$ for every element $\mathrm{cf}$ of the set.
It is closed: for a fixed spinor $u$, the inequality $\|D^{\mathrm{ch}}u\|^2\ge\mathrm{cf}\cdot(u,\nabla^{*}\nabla u)$, in the scalar product of the space of $u$, is preserved under limits of $\mathrm{cf}$.
So the set has a largest element; this is $c_{\rho}$ (row~26), and it is positive and finite.

\begin{named}{Lemma}
Let $\rho\in(0,1)$, and let $X$, $D^{\mathrm{ch}}$ and $\nabla^{*}\nabla$ be as above.
\begin{itemize}
\item[(a)] Transported to $\mathbb{Z}^3\times O_{\mathrm{cub}}$ by the bijection $\mathrm{geo}$ of Proposition~46\,(a), $D^{\mathrm{ch}}$ becomes the operator $D$ of Section~G1.4 at $\mathrm{me}=1$, and $\nabla^{*}\nabla$ becomes $\sum_iw_i(1-R_i)$ with $(R_iu)(x)=u(x^{(i)})$; neither of them contains $v_0$.
\item[(b)] $(D^{\mathrm{ch}})^{*}D^{\mathrm{ch}}\ge c_{b}\cdot\nabla^{*}\nabla$ on $\ell^2(X;\mathbb{C}^4)$ and on the spinors of every pair $(\Gamma,\chi)$, $\Gamma$ a subgroup of $\mathbb{Z}^3$ of rank three and $\chi$ a homomorphism of $\Gamma$ into the unit circle.
\end{itemize}
\end{named}
\begin{proof}
(a) By Proposition~46\,(a), $\mathrm{geo}$ is a bijection of $\mathbb{Z}^3\times O_{\mathrm{cub}}$ onto $X$ that maps the neighbour $x^{(i)}$ of a site $x$ (Section~I.1.3) to the neighbour of $\mathrm{geo}(x)$ across the wall of the type $i$.
By Proposition~46\,(b) at $A=0$ and with the lattice mass $0$, the transported $D^{\mathrm{ch}}$ is the operator (I.2.3) at the level $0$ and at $A=0$, that is, $D$ of Section~G1.4 at $\mathrm{me}=1$.
By \cite[Definition~3.17]{QEDII}, $\langle\psi,\nabla^{*}\nabla\psi\rangle=\frac{1}{2}\sum_x\sum_{i=0}^{3}w_i\left|\psi(x)-\psi(x^{(i)})\right|^2$; as $x\mapsto x^{(i)}$ is an involution, $\nabla^{*}\nabla=\sum_iw_i(1-R_i)$, and $\mathrm{geo}$ carries this operator to $\sum_iw_i(1-R_i)$ on $\mathbb{Z}^3\times O_{\mathrm{cub}}$.
The two operators so obtained are formed with $\rho$, the weights, the hopping matrices, which contain the label only, and the neighbour maps (I.O0); $v_0$ does not occur in them.

(b) By (a) the assertion is an inequality between the quadratic forms of $D$ and of $\sum_iw_i(1-R_i)$ for functions on $\mathbb{Z}^3\times O_{\mathrm{cub}}$; replacing the counting scalar product by $(\cdot,\cdot)^{\mathrm{ch}}$ multiplies both forms by $\frac{1}{48}$.
On $\ell^2(X;\mathbb{C}^4)$ the inequality is Proposition~1.3\,(b) at $\mathrm{me}=1$ and $R=\infty$.
Let $(\Gamma,\chi)$ be a pair as in (b), and $N:=[\mathbb{Z}^3:\Gamma]$; then $N\mathbb{Z}^3\subset\Gamma$.
A spinor $\psi$ of $(\Gamma,\chi)$, transported to a function $u$ on $\mathbb{Z}^3\times O_{\mathrm{cub}}$, satisfies $u(c+N\,e_{\mu},g)=\chi(N\,e_{\mu})\cdot u(c,g)$ with $|\chi(N\,e_{\mu})|=1$; so $u\in\mathrm{HH}(1,N;\theta)$ for the $\theta\in\mathbb{R}^3$ with $-\exp(i\theta_{\mu})=\chi(N\,e_{\mu})$.
By Proposition~1.3\,(b) at $\mathrm{me}=1$,
\[
\|Du\|^2\ge c_{b}\cdot(u,\nabla^{*}\nabla u)^{\mathrm{ch}},
\]
both sides being sums over a period cell of $N\mathbb{Z}^3$.
The operator $D$ commutes with the translations by $\mathbb{Z}^3$ (I.O0), so $Du$ is again transported from a spinor of $(\Gamma,\chi)$; the functions $|Du|_2^2$ and $w_i\left|u(x^{(i)})-u(x)\right|_2^2$ are $\Gamma$-periodic, and their sums over a period cell of $N\mathbb{Z}^3$ are $[\Gamma:N\mathbb{Z}^3]$ times their sums over a period cell of $\Gamma$.
Hence $\|D^{\mathrm{ch}}\psi\|^2\ge c_{b}\cdot\langle\psi,\nabla^{*}\nabla\psi\rangle$ for every spinor $\psi$ of $(\Gamma,\chi)$.
\end{proof}

\paragraph{Rows 27, 28.}
The constants $C_E$, $\mathrm{Csy}_{\mathrm{ak}}$, $\mathrm{Cpv}_0$ of these rows, which Lemma~1 excludes, exist and are finite by Propositions~37.1\,(f), (g) and 36\,(c), which are named as their references and stand later.
The constants of rows~1--118 other than those of rows~27, 28 are defined without these three letters, and no condition of (b) contains them (tables of Steps~3, 4; end of Step~3).

\paragraph{Rows 29, 30.}
$c_{\mathrm{N3}}\in(0,\infty)$.
For $x>0$ one has $0<\exp\left(-\kappa_{(0)}\cdot x\right)<1$, so that
\[
\log\frac{1}{1-\exp\left(-\kappa_{(0)}\cdot x\right)}\in(0,\infty),
\]
and this number decreases as $x$ increases.
So $\varepsilon_5$ is a non-increasing function $(0,\infty)\to[0,\infty)$.

\paragraph{Rows 31--33.}
$\alpha\in\left(\frac{199}{200},1\right)$.
$\kappa\ge7\kappa_{0}+7>0$.
$L$ is an odd integer $\ge5$ (L5).

\paragraph{Rows 34--48.}
$b_{L}=\frac{b}{L}\in(0,b)$.
$b_{-}=b\left(1-\frac{1}{L}\right)\in\left[\frac{4b}{5},b\right)$.
$b_{+}=b$.
$N_{*}=8L^3\ge1000$.
Since $L^3-1=(L-1)(L^2+L+1)$ and $L^2+L+1>(L-1)^2$, one has $a_{L}\in(1,\infty)$.
$G_{L},K_{L}\in(0,\infty)$.
$N_c$, $c_{l}$, $d_U$ are positive integers.
$c_{C}=3L^{\frac{9}{2}}$.
$C_{Q}=1$.
$\alpha^{P}=\frac{1}{4}$, $\alpha^{F}=\frac{1}{54}$, $\alpha_{0}=\min\left(\frac{1}{54},\frac{1}{4}\right)=\frac{1}{54}$.

\paragraph{Rows 49--52.}
Consider the numbers $g_1$, $c_u$, $B_D$, $B_X$, $B_Y$, $B_N$, $B_G$, $B_{\Pi}$, $c_N$, $c_G$, $c_{\Pi}$, $g_2$, $\gamma_H=\frac{g_2}{2}$ and $2c_{\mathrm{ch}}\cdot\exp\left(\frac{1}{8}\right)$ of \cite[Section~15.4.1]{QEDI} (the last one is denoted $c_{\rho}$ there; it is not the constant $c_{\rho}$ of row~26), the numbers $\mathrm{Ah}_1$, $\mathrm{Ah}_2$, $\mathrm{Ah}_3$, $S_{*}(\alpha)$ of \cite[Section~15.4.2]{QEDI} (the first three are written $\mathsf{A}_1$, $\mathsf{A}_2$, $\mathsf{A}_3$ there), and
\[
C_{\mathrm{De}}=\left(3+\frac{2}{\gamma_H}\right)^{3}\mathrm{Ah}_1,\qquad C_{\mathrm{Hk}}(\alpha)=2.5\left(3+\frac{2}{\gamma_H}\right)^{3}S_{*}(\alpha)
\]
of \cite[Theorem~15.14]{QEDI}.
They are formed by sums, products, quotients, powers and minima from positive numerals, from $\pi$, and from the following positive numbers:
$a_{\mathrm{sn}}=\frac{8}{5\pi}\sin\frac{5\pi}{8}$ (positive because $0<\frac{5\pi}{8}<\pi$, and at most $1$ because $\sin x\le x$ for $x\ge0$);
$c_{\mathrm{ch}}=\cosh\frac{1}{8}\ge1$;
$\exp\left(\frac{1}{8}\right)$;
$Z_2=1+\frac{\pi^2}{3}$;
$Z_2^3-1$ (positive because $Z_2>1$);
$1-\alpha$ (positive by row~31; it stands in the term $\frac{4}{1-\alpha}$ of $S_{*}(\alpha)$).
So these numbers lie in $(0,\infty)$.
The remaining difference among the numbers of \cite[Section~15.4.1]{QEDI}, $\mathrm{ar}_2=\mathrm{ar}_1-g_1=\frac{5\pi}{4}-\frac{1}{4}$, and the number $\mathrm{ar}_1$ describe domains; they enter none of the numbers above.
With $c_1=14L^{\frac{3}{2}}$ (\cite[Proposition~15.21(iv)]{QEDI}) and
\[
C_{\mathrm{Dt}}=c_1^2\cdot C_{\mathrm{De}}\cdot\exp(2\gamma_H\cdot L),\qquad
S_P=\exp\left(\frac{\gamma_H\cdot L}{2}\right)\left(L+\frac{2}{\gamma_H}\right)\cdot2L^3\exp\left(\frac{\gamma_H\cdot L}{2}\right)\cdot c\left(\frac{\gamma_H}{4}\right)
\]
(\cite[Section~15.7 and Section~15.1.1]{QEDI}; $c\left(\frac{\gamma_H}{4}\right)<\infty$ by (El5)) the same holds for
\begin{gather*}
c_{P}=\max\left\{9L^7\cdot C_{Q}^2,\ 15L^3\cdot\mathrm{Ah}_1\right\},\qquad
\gamma_{P}=\min\left\{\frac{\gamma_H}{2},\ \frac{1}{2c_{P}\cdot C_{\mathrm{Dt}}\cdot S_P}\right\},\\
C_{P}=\max\left\{C_{\mathrm{Hk}}(\alpha),\ C_{\mathrm{Dt}},\ (2c_{P})^{\frac{1}{2}},\ 4\mathrm{e}^2\right\}.
\end{gather*}
Moreover $\gamma_{P}\le1$ (\cite[Section~15.7]{QEDI}):
$g_2\le\frac{g_1\cdot c_N}{4B_N}\le\frac{1}{16}$, because $g_1=\frac{1}{4}$, $c_N=a_{\mathrm{sn}}^8\le1$ and $B_N=c_{\mathrm{ch}}^8+B_D\cdot B_X\ge1$;
so $\gamma_{P}\le\frac{\gamma_H}{2}=\frac{g_2}{4}\le\frac{1}{64}$.
Therefore
\[
c_{P}\in(0,\infty),\qquad\gamma_{P}\in(0,1],\qquad\gamma^{P}=\frac{\gamma_{P}}{2}\in\left(0,\frac{1}{2}\right],\qquad C_{P}\in\left[4\mathrm{e}^2,\infty\right).
\]

\paragraph{Rows 53--58.}
$\kappa_{\mathrm{w}},\kappa',\hat{\kappa}\in(0,\infty)$.
$C'_{V}=2\mathrm{e}(42+6b)=96\,\mathrm{e}$.
Since $b_{-}>0$, the number $C_{T}$ is finite, its bracket is $\ge1$, and $2\mathrm{e}^{\frac{3}{2}}\ge1$; so $C_{T}\in[1,\infty)$.
$\mu_-=L^{-3}\in(0,1)$.

\paragraph{Rows 59--62.}
$K_{1}\in(0,\infty)$.
Since $b_{-},c_{b},C_{\mathrm{Po}},K_{1}\in(0,\infty)$, also $C_{I}\in(0,\infty)$.
$c_{\mathrm{wt}}\in(0,\infty)$.
$a_{0}=\min\left\{\frac{1}{3},\frac{1}{2C_{I}\cdot c_{\mathrm{wt}}}\right\}\in\left(0,\frac{1}{3}\right]$.

\paragraph{Rows 63--66.}
$C_{G}\in[0,\infty)$ (row~24) and $1-\alpha>0$ (row~31); so $C_{R}=\frac{18\,000\,000\,C_{G}}{1-\alpha}\in[0,\infty)$.
The three summands of $C_{I}$ are $\ge0$, and $0<b_{-}<b=1$ (rows~9, 35); so $C_{I}\ge\frac{1}{b_{-}}>1$.
In the definition of $C_{3}$ (row~64) the first argument of the maximum is finite and $\ge0$ ($C_{R}\ge0$; $L,b,C_{I}>0$), the second one is $14L^3\cdot C_{I}$, and the factor $1+(1+\mathrm{e})L^2$ is $\ge1$; so $C_{3}\in\left[14L^3\cdot C_{I},\infty\right)$, in particular $C_{3}\ge14$.
As $L\ge1$ and $\mathrm{e}>1$, it follows that $C_{5}=\mathrm{e}^2\cdot L\cdot C_{3}\ge14$.
So $C_{5}\in[1,\infty)$, in accordance with the list of constants of Proposition~1.6.
$\gamma_{(5)}=\frac{a_{0}}{L}\in\left(0,\frac{1}{3L}\right]$.

\begin{named}{Remark}[A lower bound for the numbers $C_{3}$ of a general class of data]
In the notation of Section~G1.4 consider the following data:
a mesh $\mathrm{me}>0$;
a side $\mathrm{ell}=L^{j}\cdot\mathrm{me}$, $j\ge1$ an integer, and the blocks of side $\mathrm{ell}$, which partition $\mathrm{me}\,\mathbb{Z}^3$;
a period $R\in\mathrm{ell}\,\mathbb{N}$, or $R=\infty$;
a character $\theta\in\mathbb{R}^3$;
a number $\mu\ge0$ with $\mu\cdot\mathrm{ell}\le L$;
an operator $\mathrm{Pen}=\sum_E\mathrm{Pen}_E$, the sum over the blocks $E$ of side $\mathrm{ell}$, which, if $R<\infty$, commutes with the translations by $R\,e_{1}$, $R\,e_{2}$, $R\,e_{3}$, and where $\mathrm{Pen}_E$ acts on $L^2(E\times O_{\mathrm{cub}})$ and has the properties (P1)--(P3) stated below, with the penalty constants $c_1=b\left(1-\frac{1}{L}\right)$, $C_1=L\cdot b$, $C_2=2L^4\cdot b$ (in this remark and in its proof $c_1$, $C_1$, $C_2$ have only this meaning);
the operator $T:=D+\mu+\mathrm{Pen}$ on $\mathrm{HH}(\mathrm{me},R;\theta)$;
blocks $E$ and $E_1$ of side $\mathrm{ell}$;
a bounded function $f\in\mathfrak{S}^{\mathrm{ch}}(\mathrm{me},R;\theta)$ supported over the orbit of the block $E_1$, and, when $T$ is bijective, $u:=T^{-1}f$;
a point $x$ over $E$, and a point $\tilde{x}$ with $0<d_X(x,\tilde{x})\le\mathrm{ell}$.
Let $d(x,E_1)$ be the distance, measured with the $\ell^1$ norm in units of $\mathrm{ell}$ and modulo $R\mathbb{Z}^3$, of the photon site of $x$ from the orbit of $E_1$; it vanishes for $x$ over $E_1$.
\emph{Every} number $C_{3}$ such that, for all these data with $T$ bijective,
\[
|u(x)|_2+\mathrm{ell}^{\alpha}\cdot d_X(x,\tilde{x})^{-\alpha}\cdot|u(x)-u(\tilde{x})|_2\le C_{3}\cdot\mathrm{ell}\cdot\exp\left(-a_{0}\cdot d(x,E_1)\right)\cdot|f|_{\mathrm{su}},
\]
satisfies
\[
C_{3}\ge\frac{1}{b_{-}},
\]
and hence $\mathrm{e}^2\cdot L\cdot C_{3}>1$.
Here $C_{3}$ denotes any number with this property; for the number $C_{3}$ of row~64 the bound $C_{3}\ge14L^3\cdot C_{I}$, and with it $C_{5}\ge14$, was obtained above from its definition.
\end{named}
For the proof take a mesh $\mathrm{me}>0$, $\mathrm{ell}:=L\cdot\mathrm{me}$ (blocks of level one), $R:=\mathrm{ell}$, the periodic functions ($\theta=(\pi,\pi,\pi)$), $\mu:=0$ and
\[
\mathrm{Pen}:=\sum_E b_{-}\cdot\mathrm{ell}^{-1}\cdot P_E .
\]
This is the standard penalty (Std1) of \cite[Definition~12.9]{QEDI} with the number $b\left(1-\frac{1}{L}\right)=b_{-}$, written on $L^2(E\times O_{\mathrm{cub}})$.
On $L^2(E\times O_{\mathrm{cub}})$ it has the properties (P1)--(P3), which are the conditions (Pen1)--(Pen3) of \cite[Definition~12.8]{QEDI}, with $c_1=b\left(1-\frac{1}{L}\right)$, $C_1=L\cdot b$, $C_2=2L^4\cdot b$, namely
\begin{itemize}
\item[(P1)] $\mathrm{Re}\,(u,\mathrm{Pen}_Eu)\ge c_1\cdot\mathrm{ell}^{-1}\cdot\|P_Eu\|^2$;
\item[(P2)] $\|\mathrm{Pen}_Eu\|^2\le C_1\cdot\mathrm{ell}^{-1}\cdot\mathrm{Re}\,(u,\mathrm{Pen}_Eu)$, and the same for $\mathrm{Pen}_E^{*}$;
\item[(P3)] the kernel of $\mathrm{Pen}_E$ has modulus at most $C_2\cdot\mathrm{ell}^{-4}$.
\end{itemize}
Indeed, $P_E$ is the orthogonal projection of $L^2(E\times O_{\mathrm{cub}})$ onto the constants, with the kernel $\mathrm{ell}^{-3}$ (Section~G1.4); hence
\[
\mathrm{Re}\,(u,\mathrm{Pen}_Eu)=\frac{b_{-}}{\mathrm{ell}}\,\|P_Eu\|^2,\qquad
\|\mathrm{Pen}_Eu\|^2=\frac{b_{-}}{\mathrm{ell}}\,\mathrm{Re}\,(u,\mathrm{Pen}_Eu)\quad\text{with }b_{-}\le L\cdot b,
\]
$\mathrm{Pen}_E^{*}=\mathrm{Pen}_E$, and the kernel of $\mathrm{Pen}_E$ is $b_{-}\cdot\mathrm{ell}^{-4}\le2L^4\cdot b\cdot\mathrm{ell}^{-4}$.
For $\nu=1,2,3$ the translation by $R\,e_{\nu}=\mathrm{ell}\,e_{\nu}$ permutes the blocks of side $\mathrm{ell}$, and conjugation by it carries $P_E$ into $P_{E+R\,e_{\nu}}$; hence $\mathrm{Pen}$ commutes with it.
Moreover $\mu\cdot\mathrm{ell}=0\le L$.
So these are data of the kind described above, and $T$ is bijective by the following Lemma.

\begin{named}{Lemma}
Let $\mathrm{me}>0$, $\mathrm{ell}:=L\cdot\mathrm{me}$, $R:=\mathrm{ell}$ and $\theta:=(\pi,\pi,\pi)$, and let
\[
T:=D+\frac{b_{-}}{\mathrm{ell}}\sum_EP_E
\]
on $\mathrm{HH}(\mathrm{me},R;\theta)$, the sum over the blocks $E$ of side $\mathrm{ell}$ (Section~G1.4).
Then $T$ is bijective.
\end{named}
\begin{proof}
As $R<\infty$, the space $\mathrm{HH}(\mathrm{me},R;\theta)$ has finite dimension; so it suffices to show that $Tu=0$ implies $u=0$.
By (I.O2), which holds at the mesh $\mathrm{me}$ and for the period $R$ by (I.O0), the real part of $D$ is $\frac{\rho}{2\,\mathrm{me}}\nabla^{*}\nabla$, and
\[
(u,\nabla^{*}\nabla u)^{\mathrm{ch}}=\sum_{\{x,x^{(i)}\}}\frac{\mathrm{me}^3}{48}\,w_i\left|u\left(x^{(i)}\right)-u(x)\right|_2^2,
\]
the sum over the unoriented bonds counted modulo $R\mathbb{Z}^3$.
Each $P_E$ is an orthogonal projection, so that $(u,P_Eu)=\|P_Eu\|_{L^2,E}^2$.
Hence
\[
\mathrm{Re}\,(u,Tu)^{\mathrm{ch}}=\frac{\rho}{2\,\mathrm{me}}\,(u,\nabla^{*}\nabla u)^{\mathrm{ch}}+\frac{b_{-}}{\mathrm{ell}}\sum_E\|P_Eu\|_{L^2,E}^2,
\]
the last sum over the blocks modulo $R\mathbb{Z}^3$, and both terms on the right are $\ge0$.
If $Tu=0$, both terms vanish.
As $\rho>0$ and $w_i>0$, then $u(x^{(i)})=u(x)$ for every bond.
The bonds of the types $0$, $1$, $2$ join $(c,g)$ to $(c,g\,s_i)$, and $s_0$, $s_1$, $s_2$ generate $O_{\mathrm{cub}}$ (Section~I.1.2); the bond of the type $3$ joins $(c,g)$ to $(c+\mathrm{me}\cdot g\,e_1,g\,s_3)$, and as $g$ ranges over $O_{\mathrm{cub}}$, the vector $g\,e_1$ takes each of the six values $\pm e_{\mu}$.
So any two sites are joined by a path of bonds, and $u$ is constant.
Then $P_Eu=u$ for every block $E$, and $P_Eu=0$ gives $u=0$.
\end{proof}

Modulo $R\mathbb{Z}^3$ there is one block $E_1$, and its orbit is $\mathrm{me}\,\mathbb{Z}^3$.
Let $f$ be the constant function whose value is a unit vector of $\mathbb{C}^4$; it is bounded and periodic, so that $f\in\mathfrak{S}^{\mathrm{ch}}(\mathrm{me},R;\theta)$ for $\theta=(\pi,\pi,\pi)$, because $-\exp(i\pi)=1$.
By Section~G1.4,
\[
(Df)(x)=-\mathrm{me}^{-1}\sum_{i=0}^{3}M_{i}(x)\left[f\left(x^{(i)}\right)-f(x)\right]=0,
\]
and $P_Ef=f$ on $E\times O_{\mathrm{cub}}$.
Hence $Tf=\frac{b_{-}}{\mathrm{ell}}\,f$, $u=T^{-1}f=\frac{\mathrm{ell}}{b_{-}}\,f$ and $u(x)-u(\tilde{x})=0$, and at a point $x$ over $E=E_1$, where $d(x,E_1)=0$, the inequality reads
\[
\frac{\mathrm{ell}}{b_{-}}\le C_{3}\cdot\mathrm{ell}.
\]
(Since $u$ is a multiple of $f$, the choice of the norm on $\mathbb{C}^4$ does not matter.)
With $0<b_{-}<b=1$ (rows~9, 35), $L\ge1$ and $\mathrm{e}>1$ this gives
\[
\mathrm{e}^2\cdot L\cdot C_{3}\ge\frac{\mathrm{e}^2\cdot L}{b_{-}}>1
\]
for every such number $C_{3}$; a constant formed as $\mathrm{e}^2\cdot L\cdot C_{3}$ with such a number $C_{3}$ is therefore larger than $1$.

\paragraph{Rows 67--70.}
$c_{\mathrm{re}}=\min\left\{b_{-},\frac{\rho}{2C_{\mathrm{Po}}\cdot L}\right\}\in(0,b_{-}]$.
$a_{z}\in\left(0,\frac{1}{3L}\right]$.
$k_{1}=60+b=61$.
$C_{00}=48\left(\frac{2L}{c_{\mathrm{re}}}\right)(1+2k_{1})\in(0,\infty)$.

\paragraph{Rows 71--75.}
$C_z=\max(C_{5},C_{00})\in[1,\infty)$; also $C_{00}\ge1$ by itself, since $c_{\mathrm{re}}\le b_{-}<1$ and $1+2k_{1}=123$ give $C_{00}>48\cdot2L\cdot123$.
$\gamma_{z}=\min\left(\gamma_{(5)},a_{z}\right)\in\left(0,\frac{1}{3L}\right]$.
$\gamma^{F}=\frac{\gamma_{z}}{2}\in\left(0,\frac{1}{6L}\right]$.
By (El5), $c(\gamma^{F})\in[1,\infty)$; so $a'=C_z\cdot c(\gamma^{F})\in[1,\infty)$.
$\gamma_{\mathrm{w}}=\min\left(\gamma^{F},\gamma^{P}\right)\in\left(0,\frac{1}{6L}\right]$.

\paragraph{Row 76.}
$M^{F}\in[4L+1,\infty)$.

\paragraph{Row 77.}
$2\alpha_{0}\cdot402=804\,\alpha_{0}$.
The function $x\mapsto\frac{\log x}{x}$ has the derivative $\frac{1-\log x}{x^2}$, so it decreases on $[\mathrm{e},\infty)$; and $L^2\ge25>\mathrm{e}$.
Hence, for a number $M_0\ge L^2$, the property
\[
804\,\alpha_{0}\log x\le\gamma_{P}\cdot x\qquad\text{for all }x\ge M_0
\]
holds if and only if $804\,\alpha_{0}\log M_0\le\gamma_{P}\cdot M_0$.
By (El2), $\frac{\log M_0}{M_0}\le M_0^{-\frac{1}{2}}$; so the property holds if $M_0\ge\left(\frac{804\,\alpha_{0}}{\gamma_{P}}\right)^2$, which is a finite number, as $\gamma_{P}>0$.
The powers of $L$ are unbounded.
So the set of the powers $M_0\ge L^2$ of $L$ with the property is not empty, and $M^{P}$ is its least element.

\paragraph{Row 78.}
$M$ is a power of $L$ with $M\ge L^2$ (M1); $M$ is odd, and $M\ge25$.

\paragraph{Rows 79--86.}
$C_{\mathrm{By}},C_{\Gamma},C_{1},C^{F},C^{P},C_{w}^{(0)},C_{w}^{(1)},C_{w}\in(0,\infty)$: they are sums, products, powers and maxima of numbers of $(0,\infty)$, namely of $b_{-}$, $a'$, $M$, $C_{P}$, $K_{0}$, $\exp(27\kappa')$ and $\exp(2\gamma_{\mathrm{w}})$.

\paragraph{Row 87.}
The indices $t$ are those of Sections~I.5.8 and G4.0.4: a bond of $T_k$, a bond of $T^0_k$, a basis element $\Upsilon$, a unit site, a fermion site $x=(c,g)$.
The support $\operatorname{supp}t$ is the set of the end points of the bond or bonds, $\{z\}$, $\{c\}$, respectively; $d(t,t')$ is the sup distance of the supports; every support has sup-diameter at most $L$ (a bond: at most $1$; $\operatorname{supp}\Upsilon$ lies in one $L$-block or is one bond, by \cite[Proposition~15.21(ii)]{QEDI}).
To an index $t'$ of one of the three index sets of Section~I.5.8 assign a unit site $\mathrm{bpt}(t')$ as follows.
\begin{itemize}
\item[(i)] For $k\ge1$ and $t'=(y,\beta)$ put $\mathrm{bpt}(t'):=y$.
For $k=0$ and $t'=(x,\beta)$, $x=(c,g)$, put $\mathrm{bpt}(t'):=c$; as $\eta_0=1$, $c$ is a unit site.
The sum of the weights $\nu(t')$ over the $t'$ with $\mathrm{bpt}(t')=y$ is $4\cdot1=4$ for $k\ge1$, and $48\cdot4\cdot\frac{1}{48}=4$ for $k=0$ ($48$ labels $g$, four values of $\beta$, weight $\frac{1}{48}$).
\item[(ii)] For $t'=(y,y+e_{\mu})$ put $\mathrm{bpt}(t'):=y$; there are three such bonds, each of weight $1$.
\item[(iii)] For $\Upsilon=e^{y}_{i}$, $y\in L\mathbb{Z}^3$, put $\mathrm{bpt}(\Upsilon):=y$; there are $\dim\ker\tau_0=2L^3-3L^2+1$ of them (\cite[Proposition~15.21(i)]{QEDI}).
For $\Upsilon=\delta_b$ let $\mathrm{bpt}(\Upsilon)$ be the initial point of $b$; there are at most three such elements with a given $\mathrm{bpt}$.
So at most $2L^3-3L^2+4\le2L^3$ elements have the same $\mathrm{bpt}$.
\end{itemize}
In all cases the sum of the weights over the $t'$ with the same $\mathrm{bpt}$ is at most $2L^3$, and every point of $\operatorname{supp}t'$ has sup distance at most $L$ from $\mathrm{bpt}(t')$ (in (iii), at most $\frac{L-1}{2}$ within a block).
Let $t$ be an index, $c_t$ a point of $\operatorname{supp}t$, and $y_t$ the unit site with $c_t\in\Delta_{y_t}$; then $d(y_t,c_t)\le\frac{1}{2}$.
If $c\in\operatorname{supp}t$ and $c'\in\operatorname{supp}t'$ realise $d(t,t')$, then
\[
d\left(y_t,\mathrm{bpt}(t')\right)\le d(y_t,c_t)+d(c_t,c)+d(c,c')+d\left(c',\mathrm{bpt}(t')\right)\le\frac{1}{2}+L+d(t,t')+L .
\]
With (El5) it follows that
\begin{align*}
\sum_{t'}\nu(t')\exp\left(-\gamma_{\mathrm{w}}\cdot d(t,t')\right)
&\le\exp\left(\gamma_{\mathrm{w}}\left(2L+\frac{1}{2}\right)\right)\cdot2L^3\sum_{y\in T^0_k}\exp\left(-\gamma_{\mathrm{w}}\cdot d(y_t,y)\right)\\
&\le2L^3\exp\left((2L+1)\gamma_{\mathrm{w}}\right)\cdot c(\gamma_{\mathrm{w}})
\end{align*}
for all $k\ge0$, $\ell\in\mathrm{Cl}$, all indices $t$ and all three index sets.
For $k=1$, $t$ a unit site $y$ and the index set (i), the four terms with $t'=(y,\beta)$ are equal to $1$.
So
\[
c_{\gamma}\in\left[1,\ 2L^3\exp\left((2L+1)\gamma_{\mathrm{w}}\right)\cdot c(\gamma_{\mathrm{w}})\right].
\]

\paragraph{Rows 88--90.}
One has $C_W\in(0,\infty)$.
The constant of row 89 is the number defined at the beginning of \cite[Section~15.9]{QEDI}, formed with $s_d=(2L+1)\cdot M$:
\[
C_B=C_P\cdot\left[\left((2L+1)\cdot M+4\right)\cdot c_{\mathrm{bd}}(\gamma_P)+\left(\frac{2}{\gamma_P}+\frac{62}{\gamma_P^2}\right)\cdot c_{\mathrm{bd}}\left(\frac{\gamma_P}{2}\right)\right].
\]
Since $\gamma_P>0$, the estimate (El5) shows that $C_B\in(0,\infty)$.
Finally $C_d\in(0,\infty)$.

\paragraph{Row 91.}
The number $d_M(X)$ is a length divided by $M$ (\cite[Section~13.1.7]{QEDI}); hence $d_M(X)\ge0$.
By row 32 one has $\kappa\ge\kappa_0$ [5], so that $\exp(-\kappa\cdot d_M(X))\le\exp(-\kappa_0\cdot d_M(X))$.
The first of the two bounds on sums over polymers in \cite[Theorem~13.2]{QEDI} therefore gives, for every $k\ge0$, every $\ell\in\mathrm{Cl}$ and every $M$-cube $\square$,
\[
\sum_{X\in D_k,\ X\supset\square}\exp\left(-\kappa\cdot d_M(X)\right)\le K_0 .
\]
The term of the polymer $X=\square$ equals $1$, because $d_M(\square)=0$: a net that consists of a single point of the closed cube has length $0$.
Consequently $c_\kappa\in[1,K_0]$.

\paragraph{Rows 92--95.}
By (El5), $C_0\in(0,\infty)$.
Rows 35 and 36 give $0<b_-\le b_+$, and rows 12 and 29 give $\hat{g},c_{\mathrm{N3}}\in(0,\infty)$.
The numbers $K_{\mathrm{iv}}$, $C_{\mathrm{iv}}=\max\{1,(\pi^3\cdot K_{\mathrm{iv}}/8)^2\}\ge1$, $C_{\mathrm{sv}}$, $C_{\mathrm{uv}}$, $S_{\mathrm{jf}}$, $C_{\mathrm{Flg}}$ of the paragraph ``Constants'' of Proposition~30.3 are therefore obtained from numbers of $(0,\infty)$ by sums, products, powers and maxima, and
\[
\gamma_{\Sigma}=\min\left\{1,\frac{1}{2\mathrm{e}\cdot(6+b_+)\cdot C_{\mathrm{iv}}}\right\}\in(0,1],
\qquad
C_{\Sigma}=2\left(4\pi^2+\frac{\pi^4}{45}\right)\cdot C_{\mathrm{Flg}}\in(0,\infty).
\]
These are the constants with which Proposition~30.3\,(e) holds for the interval $[b_-,b_+]$.
Finally $c_S\in(0,\infty)$.

\paragraph{Rows 96--109.}
The numbers $\gamma^{F}$, $\mu_-$ and $b_-$ lie in $(0,\infty)$.
Hence
\[
\mathrm{ssum}\left(\frac{\gamma^{F}}{2}\right)=\left(\frac{24}{\gamma^{F}}\right)^3,
\qquad
\theta_{\mathrm{CT}}=\frac{\mu_-}{12\mathrm{e}},
\qquad
\mathrm{ssum}\left(\frac{\theta_{\mathrm{CT}}}{2}\right)=\left(\frac{24}{\theta_{\mathrm{CT}}}\right)^3
\]
lie in $(0,\infty)$.
Every entry of the line (C6$'$) arises from numbers of $(0,\infty)$ by sums, products, powers and minima, and so lies in $(0,\infty)$.

\paragraph{Rows 110--116.}
One has $r\ge2$; next $p_0>3r+2\ge8$, so that $p_0\ge9$; and $p>2p_0$, so that $p\ge19$.
Condition (S1) gives $0<1-\alpha<\varepsilon\le1/200$; the range of row 114 gives $0<s<1$, and (S2) gives $\sigma>0$.
Moreover $\mathfrak{e}=\sigma+2a_L/N_{*}>0$, and by row 115 $\mathfrak{e}\le s/4\le1/4$.
Altogether $0<\sigma<\mathfrak{e}\le1/4$.
\qed

Lemma~1 establishes the first assertion of (i) in Step~0, because the seven relations listed there are contained in the entries of Lemma~1 for the rows 26, 48, 65, 73, 74, 75 and 91.
The paragraph ``Ranges'' of (f) is obtained from the list of Lemma~1 in Step~8, (f11).

\subsubsection*{Step 3. (e), first sentence}

\emph{Kinds of rows.}
Every row of (a) is of one of the following five kinds.
\begin{itemize}
\item[N:] the definition in the entry consists of numerals.
\item[C:] the constant is chosen in a set; the set is described by a range given in the entry, or by the conditions of (b) that the entry names (in row 122: by the condition (Pb) of (d)).
\item[F:] the definition in the entry is a formula, or it names a formula, namely a display of \cite[Section~15]{QEDI} or the paragraph ``Constants'' of Proposition~30.1 or of Proposition~30.3.
\item[D:] the definition in the entry describes the number: an eigenvalue, a supremum, a least element or a largest element.
\item[X:] the number is one whose existence is asserted in the form ``there is a number with the property $\mathrm{Prp}$'', by the statement named as reference in the entry (in row 26: by Lemma~1, row 26).
\end{itemize}
In a row of kind X the property $\mathrm{Prp}$ contains:
\begin{itemize}
\item[(x1)] letters of (a) that are free, $\mathrm{Prp}$ being asserted for given values of them;
\item[(x2)] quantities over which $\mathrm{Prp}$ is universally quantified: levels, sides, meshes, running masses, twists, momenta, resolvent parameters, functions, fields;
\item[(x3)] in the rows 24 and 28, one letter of (a) that belongs to a \emph{later} row; $\mathrm{Prp}$ is quantified over the whole range of this letter as well (see the paragraph after the table).
\end{itemize}
By the preamble of (a), the constant of a row of kind X is a number with $\mathrm{Prp}$, fixed as soon as the letters (x1) are fixed.
In the rows 21, 24, 27, 28, 124, 125 it is the number given at the place named in the entry, that is, the number which the chapter containing that place defines and for which it proves $\mathrm{Prp}$; thus these rows are of the kinds F, D as well (see the list after the table).
In row 26 it is the largest number with $\mathrm{Prp}$, which exists by Lemma~1 (row 26); no other entry and no condition contains $c_{\rho}$.

\emph{Fifth column.}
In a row of kind N, C, F or D, the fifth column of the following table lists every letter of (a) that occurs in the definition in the entry, or in the conditions, displays and definitions that the entry names.
In a row of kind X it lists every letter (x1) of the property $\mathrm{Prp}$.
Each letter is followed by the number of its row in brackets.
The functions $c(\cdot)$, $c_{\mathrm{bd}}(\cdot)$, $\mathrm{ssum}(\cdot)$ contain no letter; $\mathrm{lg}$ contains $L$; $\mathrm{lh}(\cdot)$ contains $L$, $M$, $r$.
The numbers $\mathrm{e}$, $\pi$, $\sqrt{2}$, $\sqrt{3}$ count as numerals.

For each row of (a) the following table gives its line, its constant, its kind, the letters of (a) in its definition with their rows, and the place where these letters are read off.

\medskip
{\small
\noindent\begin{tabular}{@{}r@{\hspace{4pt}}l@{\hspace{4pt}}p{0.13\textwidth}@{\hspace{4pt}}c@{\hspace{4pt}}p{0.29\textwidth}@{\hspace{4pt}}p{0.29\textwidth}@{}}
\hline
row & line & constant & kind & letters of (a) in the definition [row] & read in \\
\hline
1 & (C0) & $\rho$ & C & -- & range \\
2 & (C0) & $w$ & N & -- & -- \\
3 & (C0) & $\bar{\varepsilon}$ & N & -- & -- \\
4 & (C0) & $c_{X}$ & N & -- & -- \\
5 & (C0) & $\kappa_0$ & N & -- & -- \\
6 & (C0) & $K_{0}$ & F & $\kappa_0$ [5] & -- \\
7 & (C0) & $c_{\mathrm{rb}}$ & F & $c_{X}$ [4], $K_{0}$ [6] & -- \\
8 & (C0) & $c_{\mathrm{fl}}$ & F & $c_{X}$ [4], $K_{0}$ [6] & -- \\
9 & (C0) & $b$ & N & -- & -- \\
10 & (C0r) & $\lambda_{\mathrm{int}}$ & D & $w$ [2] & (I.O8): the matrix contains $w_0$, $w_1$, $w_2$ and the maps $g\mapsto g\,s_i$ of $O_{\mathrm{cub}}$ \\
11 & (C0r) & $\lambda_1$ & N & -- & numerals \\
12 & (C0r) & $\hat{g}$ & F & $\rho$ [1], $\lambda_{\mathrm{int}}$ [10] & -- \\
13 & (C0r) & $c_{\mathrm{hl}}$ & F & $\rho$ [1] & -- \\
14 & (C0r) & $c_{*}$ & F & $w$ [2], $\lambda_{\mathrm{int}}$ [10] & $w_3$ is an entry of $w$ \\
15 & (C0r) & $\theta_{0}$ & F & $\rho$ [1], $\hat{g}$ [12] & Step~2: $\theta_0=\log\left(1+\hat{g}/(6(1+\rho))\right)$ \\
16 & (C0r) & $K_{\rho}$ & F & $\rho$ [1], $\hat{g}$ [12], $c_{\mathrm{hl}}$ [13] & -- \\
17 & (C0r) & $t_{0}$ & F & $K_{\rho}$ [16] & -- \\
18 & (C0r) & $c_{\mathrm{el}}$ & F & $\rho$ [1], $c_{*}$ [14], $t_{0}$ [17] & -- \\
19 & (C0r) & $K_{\mathrm{Str}}$ & D & $\rho$ [1], $w$ [2], $\theta_{0}$ [15] & Proposition~1.2\,(S3); $\mathrm{Dsc}_{\mathbf{m}}(p)$ is formed with the symbol (I.2.6), which contains $\rho$, $w$; $p$ and $\mathbf{m}$ are bound variables \\
20 & (C0r) & $\kappa_{(0)}$ & F & $\theta_{0}$ [15], $c_{\mathrm{el}}$ [18], $K_{\mathrm{Str}}$ [19] & -- \\
\hline
\end{tabular}
}

\medskip
{\small
\noindent\begin{tabular}{@{}r@{\hspace{4pt}}l@{\hspace{4pt}}p{0.13\textwidth}@{\hspace{4pt}}c@{\hspace{4pt}}p{0.29\textwidth}@{\hspace{4pt}}p{0.29\textwidth}@{}}
\hline
row & line & constant & kind & letters of (a) in the definition [row] & read in \\
\hline
21 & (C0r) & $A_{\mathrm{ak}}$, $A'_{\mathrm{ak}}$ & X & $\rho$ [1], $w$ [2], $\kappa_{(0)}$ [20] & Proposition~1.2\,(S6): ``there is $A_{\mathrm{ak}}=A_{\mathrm{ak}}(\rho)$''; (S7): ``$A'_{\mathrm{ak}}=A'_{\mathrm{ak}}(\rho)$''; the symbol contains $\rho$, $w$; (S7) contains $\kappa_{(0)}$ [20]; $p$ and $\mathbf{m}$ are quantified (``all bounds are uniform in $\mathbf{m}\in[0,1]$'') \\
22 & (C0r) & $c_{b}$ & F & $\rho$ [1], $\hat{g}$ [12], $c_{\mathrm{hl}}$ [13], $c_{*}$ [14], $K_{\rho}$ [16] & -- \\
23 & (C0r) & $C_{\mathrm{Po}}$ & F & $\lambda_{\mathrm{int}}$ [10], $c_{*}$ [14] & -- \\
24 & (C0r) & $C_{G}$ & X & $\rho$ [1], $w$ [2], $\kappa_{(0)}$ [20] & Proposition~1.5: ``there is a constant $C_G=C_G(\rho)$'' with (a)--(f); the operator contains $\rho$, $w$; (d) contains $\kappa_{(0)}$ [20]; (f) contains $\alpha$ [31]: see after the table \\
25 & (C0r) & $C_{6}$ & F & $\rho$ [1], $\hat{g}$ [12], $\theta_{0}$ [15], $K_{\mathrm{Str}}$ [19], $A'_{\mathrm{ak}}$ [21] & Proposition~30.1, paragraph ``Constants'': $C_{\mathrm{lH}}$ contains $\rho$, $\hat{g}$; $C_{\mathrm{lS}}$ contains $K_{\mathrm{Str}}$, $\theta_0$ and the $A'_{\mathrm{ak}}$ with $|\mathrm{ak}|\le3$; $\mathrm{cF13}$ is an absolute number, fixed once, which contains no letter of (a) (a number of kind X without letters (x1), whose property is asserted in \cite[Lemma~12.4]{QEDI}); in the definition in the entry, that is, in the part of the entry before ``depends on'', $c_{\mathrm{el}}$ occurs only in the last clause, which says that it enters $C_6$ only through the $A'_{\mathrm{ak}}$ \\
26 & (C0r) & $c_{\rho}$ & X & $\rho$ [1], $w$ [2] & the property: for the standard weights and at the given $\rho\in(0,1)$ of row 1, the inequality $(D^{\mathrm{ch}})^{*}D^{\mathrm{ch}}\ge c_{\rho}\cdot\nabla^{*}\nabla$ of Lemma~1, row~26 ($D^{\mathrm{ch}}$ and $\nabla^{*}\nabla$ at $A=0$, with the lattice mass $0$) holds on $\ell^2(X;\mathbb{C}^4)$ and on the spinors of every pair $(\Gamma,\chi)$, with $X$ and the pairs $(\Gamma,\chi)$ as there \\
27 & (C0r) & $C_E$ & X & $\rho$ [1], $w$ [2] & Proposition~37.1\,(f): ``there is $C_E=C_E(\rho)<\infty$'' \\
28 & (C0r) & $\mathrm{Csy}_{\mathrm{ak}}$, $\mathrm{Cpv}_0$ & X & $\rho$ [1], $w$ [2] & Proposition~37.1\,(g): $\mathrm{Csy}_{\mathrm{ak}}(\rho)$; Proposition~36\,(c): $\mathrm{Cpv}_0=\mathrm{Cpv}_0(\rho)$; (c) contains $\bar{m}$ [119]: see after the table \\
29 & (C0r) & $c_{\mathrm{N3}}$ & F & $\rho$ [1], $\lambda_{\mathrm{int}}$ [10] & -- \\
30 & (C0r) & $\varepsilon_5$ & F & $\kappa_{(0)}$ [20], $C_{6}$ [25] & -- \\
\hline
\end{tabular}
}

\medskip
{\small
\noindent\begin{tabular}{@{}r@{\hspace{4pt}}l@{\hspace{4pt}}p{0.13\textwidth}@{\hspace{4pt}}c@{\hspace{4pt}}p{0.29\textwidth}@{\hspace{4pt}}p{0.29\textwidth}@{}}
\hline
row & line & constant & kind & letters of (a) in the definition [row] & read in \\
\hline
31 & (C1) & $\alpha$ & C & $\bar{\varepsilon}$ [3] & range \\
32 & (C2) & $\kappa$ & C & $\kappa_0$ [5] & range \\
33 & (C3) & $L$ & C & $\bar{\varepsilon}$ [3], $\kappa_0$ [5], $c_{\mathrm{rb}}$ [7], $c_{\mathrm{fl}}$ [8], $\kappa$ [32] & (L1)--(L5): Step~4 \\
34 & (C4) & $b_{L}$ & F & $b$ [9], $L$ [33] & -- \\
35 & (C4) & $b_{-}$ & F & $b$ [9], $L$ [33] & -- \\
36 & (C4) & $b_{+}$ & F & $b$ [9] & -- \\
37 & (C4) & $N_{*}$ & F & $L$ [33] & -- \\
38 & (C4) & $a_{L}$ & F & $L$ [33] & -- \\
39 & (C4) & $G_{L}$ & F & $L$ [33] & -- \\
40 & (C4) & $K_{L}$ & F & $L$ [33], $G_{L}$ [39] & -- \\
41 & (C4) & $N_c$ & F & $L$ [33] & -- \\
42 & (C4) & $c_{l}$ & F & $L$ [33] & -- \\
43 & (C4) & $d_U$ & F & $L$ [33] & -- \\
44 & (C4) & $c_{C}$ & F & $L$ [33] & -- \\
45 & (C4) & $C_{Q}$ & N & -- & -- \\
46 & (C4) & $\alpha^{P}$ & N & -- & -- \\
47 & (C4) & $\alpha^{F}$ & N & -- & -- \\
48 & (C4) & $\alpha_{0}$ & F & $\alpha^{P}$ [46], $\alpha^{F}$ [47] & -- \\
49 & (C4) & $c_{P}$ & F & $L$ [33], $C_{Q}$ [45] & \cite[Section~15.7]{QEDI}: $\max\{9L^7\cdot C_Q^2,\ 15L^3\cdot\mathrm{Ah}_1\}$; $\mathrm{Ah}_1$ is a numeral of \cite[Section~15.4.2]{QEDI}, where it is denoted $\mathsf{A}_1$ \\
50 & (C4) & $\gamma_{P}$ & F & $L$ [33], $c_{P}$ [49] & \cite[Section~15.7]{QEDI}: $\min\{\gamma_H/2,\ 1/(2c_P\cdot C_{\mathrm{Dt}}\cdot S_P)\}$; $C_{\mathrm{Dt}}$, $S_P$ are formulas in $L$ and in the numerals $\gamma_H$, $C_{\mathrm{De}}$ of \cite[Section~15.4.1]{QEDI} and \cite[Theorem~15.14]{QEDI} \\
51 & (C4) & $\gamma^{P}$ & F & $\gamma_{P}$ [50] & -- \\
52 & (C4) & $C_{P}$ & F & $\alpha$ [31], $L$ [33], $c_{P}$ [49] & \cite[Section~15.7]{QEDI}: $\max\{C_{\mathrm{Hk}}(\alpha),\allowbreak\ C_{\mathrm{Dt}},\allowbreak\ (2c_P)^{\frac{1}{2}},\allowbreak\ 4\mathrm{e}^2\}$ \\
53 & (C4) & $\kappa_{\mathrm{w}}$ & F & $\kappa_0$ [5], $\kappa$ [32] & -- \\
54 & (C4) & $\kappa'$ & F & $\kappa_0$ [5], $\kappa_{\mathrm{w}}$ [53] & -- \\
55 & (C4) & $\hat{\kappa}$ & F & $c_{X}$ [4], $\kappa$ [32], $N_c$ [41] & -- \\
56 & (C4) & $C'_{V}$ & F & $b$ [9] & -- \\
57 & (C4) & $C_{T}$ & F & $L$ [33], $b_{-}$ [35] & -- \\
58 & (C4) & $\mu_-$ & F & $L$ [33] & -- \\
\hline
\end{tabular}
}

\medskip
{\small
\noindent\begin{tabular}{@{}r@{\hspace{4pt}}l@{\hspace{4pt}}p{0.13\textwidth}@{\hspace{4pt}}c@{\hspace{4pt}}p{0.29\textwidth}@{\hspace{4pt}}p{0.29\textwidth}@{}}
\hline
row & line & constant & kind & letters of (a) in the definition [row] & read in \\
\hline
59 & (C4r) & $K_{1}$ & F & $b$ [9], $L$ [33] & -- \\
60 & (C4r) & $C_{I}$ & F & $c_{b}$ [22], $C_{\mathrm{Po}}$ [23], $b_{-}$ [35], $K_{1}$ [59] & -- \\
61 & (C4r) & $c_{\mathrm{wt}}$ & F & $\rho$ [1], $b$ [9], $L$ [33] & -- \\
62 & (C4r) & $a_{0}$ & F & $C_{I}$ [60], $c_{\mathrm{wt}}$ [61] & -- \\
63 & (C4r) & $C_{R}$ & F & $C_{G}$ [24], $\alpha$ [31] & -- \\
64 & (C4r) & $C_{3}$ & F & $b$ [9], $L$ [33], $C_{I}$ [60], $C_{R}$ [63] & -- \\
65 & (C4r) & $C_{5}$ & F & $L$ [33], $C_{3}$ [64] & -- \\
66 & (C4r) & $\gamma_{(5)}$ & F & $L$ [33], $a_{0}$ [62] & -- \\
67 & (C4r) & $c_{\mathrm{re}}$ & F & $\rho$ [1], $b$ [9], $C_{\mathrm{Po}}$ [23], $L$ [33] & -- \\
68 & (C4r) & $a_{z}$ & F & $b$ [9], $L$ [33], $c_{\mathrm{re}}$ [67] & -- \\
69 & (C4r) & $k_{1}$ & F & $b$ [9] & -- \\
70 & (C4r) & $C_{00}$ & F & $L$ [33], $c_{\mathrm{re}}$ [67], $k_{1}$ [69] & -- \\
71 & (C4r) & $C_z$ & F & $C_{5}$ [65], $C_{00}$ [70] & -- \\
72 & (C4r) & $\gamma_{z}$ & F & $\gamma_{(5)}$ [66], $a_{z}$ [68] & -- \\
73 & (C4r) & $\gamma^{F}$ & F & $\gamma_{z}$ [72] & $L$ occurs only in the bound in brackets \\
74 & (C4r) & $a'$ & F & $C_z$ [71], $\gamma^{F}$ [73] & -- \\
75 & (C4r) & $\gamma_{\mathrm{w}}$ & F & $\gamma^{P}$ [51], $\gamma^{F}$ [73] & -- \\
76 & (C4r) & $M^{F}$ & F & $L$ [33], $\kappa'$ [54], $C'_{V}$ [56], $a'$ [74] & -- \\
77 & (C4r) & $M^{P}$ & D & $L$ [33], $\alpha_{0}$ [48], $\gamma_{P}$ [50] & in the entry, $x$ and $M_0$ are bound variables \\
78 & (C5) & $M$ & C & $L$ [33], $\alpha_{0}$ [48], $\hat{\kappa}$ [55], $\gamma_{\mathrm{w}}$ [75], $M^{F}$ [76], $M^{P}$ [77] & (M1)--(M6): Step~4 \\
\hline
\end{tabular}
}

\medskip
{\small
\noindent\begin{tabular}{@{}r@{\hspace{4pt}}l@{\hspace{4pt}}p{0.13\textwidth}@{\hspace{4pt}}c@{\hspace{4pt}}p{0.29\textwidth}@{\hspace{4pt}}p{0.29\textwidth}@{}}
\hline
row & line & constant & kind & letters of (a) in the definition [row] & read in \\
\hline
79 & (C6) & $C_{\mathrm{By}}$ & F & $b_{-}$ [35] & -- \\
80 & (C6) & $C_{\Gamma}$ & F & $b_{-}$ [35], $a'$ [74] & -- \\
81 & (C6) & $C_{1}$ & F & $b$ [9], $\alpha_{0}$ [48], $\kappa'$ [54], $a'$ [74], $M$ [78], $C_{\Gamma}$ [80] & -- \\
82 & (C6) & $C^{F}$ & F & $b$ [9], $a'$ [74], $C_{\mathrm{By}}$ [79], $C_{\Gamma}$ [80], $C_{1}$ [81] & -- \\
83 & (C6) & $C^{P}$ & F & $\alpha_{0}$ [48], $C_{P}$ [52], $M$ [78] & -- \\
84 & (C6) & $C_{w}^{(0)}$ & F & $C^{F}$ [82], $C^{P}$ [83] & -- \\
85 & (C6) & $C_{w}^{(1)}$ & F & $\alpha_{0}$ [48], $\gamma_{\mathrm{w}}$ [75], $M$ [78], $C_{w}^{(0)}$ [84] & -- \\
86 & (C6) & $C_{w}$ & F & $K_{0}$ [6], $C_{w}^{(1)}$ [85] & -- \\
87 & (C6) & $c_{\gamma}$ & D & $L$ [33], $\gamma_{\mathrm{w}}$ [75], $M$ [78] & Section~I.5.8; $L$ through the basis $\{\Upsilon\}$ and the fine lattices, $M$ through the class $\mathrm{Cl}$ (Remark~3.1) \\
88 & (C6) & $C_{W}$ & F & $C_{w}$ [86], $c_{\gamma}$ [87] & -- \\
89 & (C6) & $C_{B}$ & F & $L$ [33], $\gamma_{P}$ [50], $C_{P}$ [52], $M$ [78] & \cite[Section~15.9]{QEDI} \\
90 & (C6) & $C_{d}$ & F & $c_{X}$ [4], $b$ [9], $C_{w}$ [86] & -- \\
91 & (C6) & $c_{\kappa}$ & D & $\kappa$ [32], $L$ [33], $M$ [78] & $L$, $M$ through the class $\mathrm{Cl}$ and the $M$-cubes (Remark~3.2); $K_0$ occurs only in the bound $c_{\kappa}\le K_0$ stated in the entry \\
92 & (C6) & $C_{0}$ & F & $b_{+}$ [36], $\gamma_{P}$ [50], $C_{P}$ [52], $\gamma^{F}$ [73], $C^{F}$ [82] & -- \\
93 & (C6) & $\gamma_{\Sigma}$ & F & $\hat{g}$ [12], $c_{\mathrm{N3}}$ [29], $b_{-}$ [35], $b_{+}$ [36] & Proposition~30.3, paragraph ``Constants'': $C_{\mathrm{iv}}$ is formed with $K_{\mathrm{iv}}$, which contains $b_-$, $b_+$, $\hat{g}$, $c_{\mathrm{N3}}$ \\
94 & (C6) & $C_{\Sigma}$ & F & $\hat{g}$ [12], $c_{\mathrm{N3}}$ [29], $b_{-}$ [35], $b_{+}$ [36] & Proposition~30.3, paragraph ``Constants'': $C_{\mathrm{Flg}}$ is formed with $b_+$, $c_{\mathrm{N3}}$ and $C_{\mathrm{iv}}$ \\
95 & (C6) & $c_{S}$ & F & $c_{X}$ [4], $K_{0}$ [6], $\kappa$ [32] & -- \\
\hline
\end{tabular}
}

\medskip
{\small
\noindent\begin{tabular}{@{}r@{\hspace{4pt}}l@{\hspace{4pt}}p{0.13\textwidth}@{\hspace{4pt}}c@{\hspace{4pt}}p{0.29\textwidth}@{\hspace{4pt}}p{0.29\textwidth}@{}}
\hline
row & line & constant & kind & letters of (a) in the definition [row] & read in \\
\hline
96 & (C6$'$) & $\mathbf{S}$ & F & $\gamma^{F}$ [73], $C^{F}$ [82] & -- \\
97 & (C6$'$) & $\mathrm{hL}$ & F & $\mathbf{S}$ [96] & -- \\
98 & (C6$'$) & $C_K$ & F & $\mathbf{S}$ [96] & -- \\
99 & (C6$'$) & $C_K'$ & F & $\mathbf{S}$ [96] & -- \\
100 & (C6$'$) & $C_H$ & F & $M$ [78], $C^{F}$ [82] & -- \\
101 & (C6$'$) & $\mathrm{CVj}$ & F & $K_{0}$ [6], $\gamma^{F}$ [73], $C_H$ [100] & -- \\
102 & (C6$'$) & $\theta_{\mathrm{CT}}$ & F & $\mu_-$ [58] & -- \\
103 & (C6$'$) & $\mathrm{vt}$ & F & $\gamma^{F}$ [73], $M$ [78], $\theta_{\mathrm{CT}}$ [102] & -- \\
104 & (C6$'$) & $\mathrm{gGa}$ & F & $b_{-}$ [35], $\mu_-$ [58] & -- \\
105 & (C6$'$) & $h_{G}$ & F & $\mathrm{gGa}$ [104] & -- \\
106 & (C6$'$) & $\mathrm{gGa}_{1}$ & F & $b_{-}$ [35], $\mu_-$ [58] & -- \\
107 & (C6$'$) & $\mathrm{alG}$ & F & $b_{-}$ [35], $\mu_-$ [58], $\theta_{\mathrm{CT}}$ [102] & -- \\
108 & (C6$'$) & $c_{\tau}$ & F & $\mu_-$ [58], $\mathbf{S}$ [96] & -- \\
109 & (C6$'$) & $\mathrm{CWj}$ & F & $K_{0}$ [6], $M$ [78], $C_{d}$ [90], $\mathrm{CVj}$ [101], $h_{G}$ [105], $\mathrm{alG}$ [107] & -- \\
110 & (C7) & $r$ & C & -- & range \\
111 & (C7) & $p_0$ & C & $r$ [110] & range \\
112 & (C7) & $p$ & C & $p_0$ [111] & range \\
113 & (C7) & $\varepsilon$ & C & $\bar{\varepsilon}$ [3], $\alpha$ [31] & (S1) \\
114 & (C7) & $s$ & C & $N_{*}$ [37], $a_{L}$ [38], $\varepsilon$ [113] & (S2): Step~4 \\
115 & (C7) & $\mathfrak{e}$ & F & $\varepsilon$ [113], $s$ [114] & -- \\
116 & (C7) & $\sigma$ & F & $N_{*}$ [37], $a_{L}$ [38], $\mathfrak{e}$ [115] & -- \\
117 & (C8) & $e_0$ & C & $c_{X}$ [4], $K_{0}$ [6], $c_{\mathrm{rb}}$ [7], $c_{\mathrm{fl}}$ [8], $L$ [33], $c_{C}$ [44], $C'_{V}$ [56], $C_{T}$ [57], $a'$ [74], $M$ [78], $C_{w}$ [86], $c_{\gamma}$ [87], $C_{B}$ [89], $C_{d}$ [90], $c_{S}$ [95], $p_0$ [111], $p$ [112], $\varepsilon$ [113] & (A1)--(A10), (Z2.s1)--(Z2.s5), (Z3.s1): Step~4 \\
118 & (C9) & $e_1$ & D & $c_{\mathrm{fl}}$ [8], $L$ [33], $b_{-}$ [35], $N_{*}$ [37], $M$ [78], $c_{\kappa}$ [91], $C_{0}$ [92], $\gamma_{\Sigma}$ [93], $C_{\Sigma}$ [94], $r$ [110], $p_0$ [111], $p$ [112], $\varepsilon$ [113], $e_0$ [117] & (e1)--(e8): Step~4 \\
\hline
\end{tabular}
}

\medskip
{\small
\noindent\begin{tabular}{@{}r@{\hspace{4pt}}l@{\hspace{4pt}}p{0.13\textwidth}@{\hspace{4pt}}c@{\hspace{4pt}}p{0.29\textwidth}@{\hspace{4pt}}p{0.29\textwidth}@{}}
\hline
row & line & constant & kind & letters of (a) in the definition [row] & read in \\
\hline
119 & (C10) & $\bar{m}$ & C & -- & range \\
120 & (C10) & $h$ & F & $L$ [33], $\bar{m}$ [119] & -- \\
121 & (C10) & $e_G$ & D & $e_1$ [118], $h$ [120] & in the entry, $e$ is a bound variable \\
122 & (C10) & $\beta$, $x_0$ & C & $\bar{m}$ [119] & Section~I.7.6; (Pb) of (d): $b(f)$ and the inequality of \cite[Section~2.8]{QEDI} contain $\bar{m}$ (see 7.3) \\
123 & (C10) & $D_{\mathrm{be}}$, $D_\beta$, $C_{\beta}$, $\mathbf{D}$, $\mathrm{Cb}$ & D & $\alpha$ [31], $L$ [33], $h$ [120], $\beta$ [122], $x_0$ [122] & Section~G5.1.3; Proposition~35.3\,(c1): ``$C_\beta$ and $\mathbf{D}$ depend on $(\beta,h,L,\alpha)$ only'' \\
124 & (C10) & $\mathrm{Cah}_i$, $\mathrm{Cpv}_1$, $\mathrm{Cpv}_2$, $\mathrm{epv}$, $\mathrm{Cpv}_\beta$ & X & $\rho$ [1], $w$ [2], $\bar{m}$ [119], $\beta$ [122], $x_0$ [122] & Proposition~36\,(e): $\mathrm{Cah}_i(\beta)$; (f): $\mathrm{Cpv}_1(\rho,\beta)$, $\mathrm{Cpv}_2(\rho,\bar{m},\beta)$, $\mathrm{epv}$; (g): $\mathrm{Cpv}_\beta$ \\
125 & (C10) & $C_{\mathrm{ct}}$, $C_{1}^{\mathrm{N2}}$, $C_{2}^{\mathrm{N2}}$ & X & $\rho$ [1], $w$ [2], $\bar{m}$ [119], $\beta$ [122], $x_0$ [122], $D_{\mathrm{be}}$ [123] & Proposition~37.2 (setting $R\ge D_{\mathrm{be}}+1$), (c): $C_{\mathrm{ct}}(\rho,\beta)$; (d): $C_1^{\mathrm{N2}}(\rho,\bar{m},\beta)$; (e): $C_2^{\mathrm{N2}}(\bar{m},\beta)$, with $\exp(-\bar{m}\cdot(R-D_{\mathrm{be}}))$ \\
126 & (C10) & $C_{\mathrm{fr}}$ & F & $K_{0}$ [6], $L$ [33], $M$ [78], $C_{d}$ [90], $h$ [120], $D_\beta$ [123], $C_{\beta}$ [123] & -- \\
127 & (C10) & $e_{\sharp}$ & D & $K_{0}$ [6], $L$ [33], $N_{*}$ [37], $\mu_-$ [58], $M$ [78], $C_{d}$ [90], $C_{0}$ [92], $\mathrm{hL}$ [97], $C_K$ [98], $C_K'$ [99], $\mathrm{vt}$ [103], $h_{G}$ [105], $\mathrm{gGa}_{1}$ [106], $\mathrm{alG}$ [107], $c_{\tau}$ [108], $\mathrm{CWj}$ [109], $r$ [110], $p_0$ [111], $p$ [112], $\varepsilon$ [113], $s$ [114], $e_1$ [118], $h$ [120], $D_\beta$ [123], $C_{\beta}$ [123], $\mathbf{D}$ [123], $\mathrm{Cb}$ [123], $C_{\mathrm{fr}}$ [126] & (g1)--(g9): Step~7 \\
\hline
\end{tabular}
}

\medskip
{\small
\noindent\begin{tabular}{@{}r@{\hspace{4pt}}l@{\hspace{4pt}}p{0.13\textwidth}@{\hspace{4pt}}c@{\hspace{4pt}}p{0.29\textwidth}@{\hspace{4pt}}p{0.29\textwidth}@{}}
\hline
row & line & constant & kind & letters of (a) in the definition [row] & read in \\
\hline
128 & (C10) & $C^{G}_{\beta}$ & F & $\bar{m}$ [119], $h$ [120], $\beta$ [122], $x_0$ [122], $D_{\mathrm{be}}$ [123], $\mathrm{Cb}$ [123], $\mathrm{Cpv}_\beta$ [124], $C_{1}^{\mathrm{N2}}$ [125], $C_{2}^{\mathrm{N2}}$ [125] & $b(\cdot)$ is the functional of Section~I.7.6 (it contains $\bar{m}$), not the constant $b$ \\
129 & (C10) & $C_{\mathrm{mat}}$ & F & $\varepsilon_5$ [30], $L$ [33] & -- \\
130 & (C10) & $s_{B}$ & F & $\beta$ [122], $x_0$ [122] & $g_{\mathrm{ix}}$, $\beta_{\mathrm{ix}}$: Proposition~43.2, paragraph ``Data'' \\
131 & (C10) & $C_{3\mathrm{p}}$ & F & $C^{G}_{\beta}$ [128] & -- \\
132 & (C10) & $C_{\Gamma}^{(16)}$ & F & $L$ [33], $N_{*}$ [37], $a_{L}$ [38], $K_{L}$ [40], $C_{\mathrm{mat}}$ [129] & -- \\
133 & (C10) & $C_{U}$ & F & $N_{*}$ [37], $s_{B}$ [130], $C_{3\mathrm{p}}$ [131], $C_{\Gamma}^{(16)}$ [132] & -- \\
134 & (C10) & $C_{\bar{\kappa}}$ & F & $C_{U}$ [133] & -- \\
135 & (C10) & $C'_{\bar{\kappa}}$ & F & $C_{\bar{\kappa}}$ [134] & -- \\
136 & (C10) & $e_2$ & D & $\sigma$ [116], $\bar{m}$ [119], $e_G$ [121], $\beta$ [122], $x_0$ [122], $e_{\sharp}$ [127], $C'_{\bar{\kappa}}$ [135] & (Z10.s1), (Z10.s2), (Z10.s4); $b(f)$ contains $\bar{m}$ and $\beta$; the second formula of the display (Z10.s4) repeats the definition of row 135, and $C_{\bar{\kappa}}$ [134] occurs only there \\
\hline
\end{tabular}
}

\medskip
\paragraph{The two letters (x3).}

\emph{Row 24: $\alpha$ [31] in Proposition~1.5\,(f).}
Apart from the quantified masses $\mathbf{m}$, $\mu$, the mesh $\mathrm{me}$ and the sites, the clauses (a)--(e) of Proposition~1.5 contain $\rho$ and $w$ (through the operator and its symbol) and $\kappa_{(0)}$ [20] (in (d)), and no other letter of (a).
Let $C_G$ be a number with (a)--(e).
It is shown that (f) then holds with this $C_G$ for every $\alpha\in[0,1]$.
Let $x=(c,g)\ne\tilde{x}=(\tilde{c},\tilde{g})$ and $x'=(c',g')$ be points of $\mathrm{me}\,\mathbb{Z}^3\times O_{\mathrm{cub}}$, and put
\[
\mathrm{hd}:=d_X(x,\tilde{x})=\max(|c-\tilde{c}|,\mathrm{me}),
\qquad
\mathrm{rmn}:=\mathrm{me}+\min(|c-c'|,|\tilde{c}-c'|).
\]
Both sides of (f) are symmetric in $(x,\tilde{x})$, so it suffices to treat the case in which the minimum is attained at $c$, that is, $\mathrm{rmn}=\mathrm{me}+|c-c'|$.

\emph{Case $\mathrm{hd}\ge\mathrm{rmn}/2$.}
Here $\mathrm{me}+|c-c'|=\mathrm{rmn}\le2\,\mathrm{hd}$, and $\mathrm{me}+|\tilde{c}-c'|\le\mathrm{rmn}+|\tilde{c}-c|\le3\,\mathrm{hd}$.
Both numbers thus lie in $(0,3\,\mathrm{hd}]$, and every $t\in(0,3\,\mathrm{hd}]$ satisfies $t^{-2}=t^{\alpha}\cdot t^{-2-\alpha}\le(3\,\mathrm{hd})^{\alpha}\cdot t^{-2-\alpha}$.
The bound $|G(x,x')|_2\le2C_G\cdot(\mathrm{me}+|c-c'|)^{-2}$ of (e), used at $x$ and at $\tilde{x}$, therefore yields
\[
|G(x,x')-G(\tilde{x},x')|_2\le2C_G\cdot3^{\alpha}\cdot\mathrm{hd}^{\alpha}\cdot\left[(\mathrm{me}+|c-c'|)^{-2-\alpha}+(\mathrm{me}+|\tilde{c}-c'|)^{-2-\alpha}\right],
\]
and $2\cdot3^{\alpha}\le6\le6240$.

\emph{Case $\mathrm{hd}<\mathrm{rmn}/2$.}
The statement (I.s6) concerns $\mathbb{Z}^3\times O_{\mathrm{cub}}$ and is unchanged under a dilation of $c$.
By it, $x$ and $\tilde{x}$ are joined by a path of at most $39\,\mathrm{hd}/\mathrm{me}$ bonds whose photon sites $c''$ lie in the lattice box spanned by $c$ and $\tilde{c}$.
Hence $|c''-c|\le|\tilde{c}-c|\le\mathrm{hd}$, and $\mathrm{me}+|c''-c'|\ge\mathrm{rmn}-\mathrm{hd}>\mathrm{rmn}/2$.
For a site $x''$ over $c''$ and $i=0,\dots,3$, (e) contains the bound
\[
|G({x''}^{(i)},x')-G(x'',x')|_2\le20C_G\cdot\mathrm{me}\cdot(\mathrm{me}+|c''-c'|)^{-3}.
\]
Along each bond of the path the function $G(\cdot,x')$ thus changes by at most $20C_G\cdot\mathrm{me}\cdot(\mathrm{rmn}/2)^{-3}=160C_G\cdot\mathrm{me}\cdot\mathrm{rmn}^{-3}$, and
\begin{align*}
|G(x,x')-G(\tilde{x},x')|_2
&\le39\cdot\frac{\mathrm{hd}}{\mathrm{me}}\cdot160C_G\cdot\mathrm{me}\cdot\mathrm{rmn}^{-3}
=6240C_G\cdot\mathrm{hd}^{\alpha}\cdot\mathrm{rmn}^{-2-\alpha}\cdot\left(\frac{\mathrm{hd}}{\mathrm{rmn}}\right)^{1-\alpha}\\
&\le6240C_G\cdot\mathrm{hd}^{\alpha}\cdot(\mathrm{me}+|c-c'|)^{-2-\alpha},
\end{align*}
since $\mathrm{hd}/\mathrm{rmn}<1$ and $1-\alpha\ge0$.

Consequently, for a number $C_G$ the property ``(a)--(f), with (f) for every $\alpha\in[0,1]$'' coincides with the property ``(a)--(e)'', which does not contain $\alpha$.
Proposition~1.5 asserts it in this form: there is a constant $C_G=C_G(\rho)$ such that the clauses hold for all $\mathbf{m}\in[0,1]$ and, in (f), for every $\alpha\in[0,1]$
(cf.\ the paragraph preceding Proposition~1.1: the constants of Propositions~1.1--1.5 depend only on $\rho$ and the weights).

\emph{Row 28: $\bar{m}$ [119] in Proposition~36\,(c).}
The three bounds of Proposition~36\,(c) contain $\rho$ and $w$ (through the symbol), $\bar{m}$ (explicitly, and through $\mathbf{m}=\bar{m}\cdot\eta$), and the quantified variables $\mathrm{kph}$, $\eta$, $R$; they contain no probe.
The number $\bar{m}$ is a datum of the setting (FF), with the range $(0,1]$.
Proposition~36 asserts the bounds with a number $\mathrm{Cpv}_0=\mathrm{Cpv}_0(\rho)$ for all $\bar{m}\in(0,1]$
(cf.\ Section~I.8.2, line (C0r), where $\mathrm{Cpv}_0$ is among the constants that depend on $\rho$; (G1.T-11) in Annex~A, Section~G1.Z1, where $\mathrm{Csy}_{\mathrm{ak}}$, $\mathrm{Cpv}_0$, $C_E$ are functions of $\rho$; and Section~G5.9.1).
For given $\rho$, one number thus serves all $\bar{m}\in(0,1]$, and in the property that defines $\mathrm{Cpv}_0$ the letter $\bar{m}$ is quantified over its whole range.

In no other row of the kinds D, X does the property that defines the constant contain a letter of a later row (last column of the table; Remark~3.3).
This proves the assertion of the preamble of (a) on letters of later rows.

\paragraph{The rows of kind X other than row 26 as rows of the kinds F, D\@.}
In the rows 21, 24, 27, 28, 124, 125 the constant is the number given by the following display of the chapter named in the entry, the display being taken as the definition of that number.
After each display stand all letters of (a), with their rows, that occur in it and in the definitions of the auxiliary numbers it names.
\begin{itemize}
\item[Row 21.] Chapter~1a, Sections~2.6, 2.7 ($A_0=1/c_{\mathrm{el}}$, $A'_0=2/c_{\mathrm{el}}$, and a recursion on $|\mathrm{ak}|$): $\theta_0$ [15], $c_{\mathrm{el}}$ [18], $K_{\mathrm{Str}}$ [19].
\item[Row 24.] Chapter~1b, Section~1.4, the display of $C_G$, together with the families of Lemma~5.1 of that chapter: $\hat{g}$ [12], $c_{\mathrm{hl}}$ [13], $\theta_0$ [15], $\kappa_{(0)}$ [20], $A_{\mathrm{ak}}$, $A'_{\mathrm{ak}}$ [21]; the numbers $c(1,3)$, $c(0,4)$ in it are absolute numbers, like $\mathrm{cF13}$ (row 25).
\item[Row 27.] Chapter~37, [1.10]: $\rho$ [1], $w$ [2], $\lambda_1$ [11] ($w$ enters through the matrix $(\nabla^{*}\nabla)(p)$ of Proposition~37.1\,(a); the momentum $p$ is bound by a minimum).
\item[Row 28.] $\mathrm{Csy}_{\mathrm{ak}}$: Chapter~37, [1.12] ($\mathrm{Csy}_0=\sqrt{2}+7C_E$, and a recursion on $|\mathrm{ak}|$): $C_E$ [27].
$\mathrm{Cpv}_0$: Chapter~36, Section~4.7: the numbers $\mathrm{Csy}_{\mathrm{ak}}$, $|\mathrm{ak}|\le4$, of the same row.
\item[Row 124.] Chapter~36, whose table of the constants collects the values.
$\mathrm{Cah}_i$, (6.4): $\beta$, $x_0$ [122], $D_{\mathrm{be}}$ [123].
$\mathrm{Cpv}_1$, (7.6): $\mathrm{Cpv}_0$ [28], $\mathrm{Cah}_2$.
$\mathrm{Cpv}_2$, (7.7) with (7.2), (6.7): $\mathrm{Cpv}_0$ [28], $\bar{m}$ [119], $\beta$, $x_0$ [122], $D_{\mathrm{be}}$ [123], $\mathrm{Cah}_3$.
$\mathrm{epv}$, (7.8): the letters of $\mathrm{Cpv}_2$, and $\rho$ [1], $w$ [2] through the symbol.
$\mathrm{Cpv}_\beta$ is the supremum of Proposition~36\,(g): $\rho$ [1], $w$ [2], $\bar{m}$ [119], $\beta$, $x_0$ [122] ($\eta$ and $R$ are bound).
\item[Row 125.] Chapter~37, [2.1], with the numbers of its Section~2.0: $C_E$ [27], $\bar{m}$ [119], $\beta$, $x_0$ [122], $D_{\mathrm{be}}$ [123]; in addition $C_1^{\mathrm{N2}}$ contains $C_{\mathrm{ct}}$, an earlier name of the same row.
\end{itemize}
Each of these letters is named after ``depends on'' in the entry of its row, or is an earlier name of the row, or is one of $\rho$, $w$, or is one of the following three:
$\lambda_1$ (row 27), whose definition consists of numerals;
$C_E$ (row 28), which is determined by $\rho$, $w$ (row 27);
$D_{\mathrm{be}}$ (rows 124, 125), which is formed from $\beta$, $x_0$ (row 123).
With these numbers, every row other than row 26 is of one of the kinds N, C, F, D.

In each of the 136 rows, every letter of the fifth column carries a number smaller than the number of the row, and so does every letter of another row in the list just given for the rows 21, 24, 27, 28, 124, 125.
In a row with several names, a later name uses earlier names of the same row only (row 123: $D_\beta=D_{\mathrm{be}}+4$, $\mathrm{Cb}=1+h^{\frac{1}{2}}\cdot C_\beta$; rows 28, 124, 125: the list above).
The parameter sets $\mathrm{par}^{0}$, $\mathrm{par}$ defined after the list (a) are formed from $\alpha_0$ [48], $\gamma_{\mathrm{w}}$ [75], $\kappa_{\mathrm{w}}$ [53], $\kappa_0$ [5], $C_w^{(1)}$ [85], $C_w$ [86].
This proves the first half of the first sentence of (e); the second half is proved in Remark~3.4.
Consider the directed graph with an edge from each letter of the fifth column, and from each letter of another row in the list above, to the constant of the row.
Every edge increases the number of the row, and $\rho$ has the number 1; hence the graph contains no directed cycle.

The constants $K_1$ [59], $k_1$ [69] and $M^P$ [77], which stand under the mark (C4r), contain neither $\rho$ nor the weights.
Indeed, going back from them through the fifth column one meets $b$ [9], $L$ [33], $\alpha_0$ [48], $\gamma_P$ [50], then $C_Q$ [45], $\alpha^P$ [46], $\alpha^F$ [47], $c_P$ [49], and, from the row of $L$, $\bar{\varepsilon}$ [3], $c_X$ [4], $\kappa_0$ [5], $K_0$ [6], $c_{\mathrm{rb}}$ [7], $c_{\mathrm{fl}}$ [8], $\kappa$ [32]; neither $\rho$ [1] nor $w$ [2] is among them.
Every other row under the marks (C0r), (C4r) meets $\rho$ or $w$ in this way, with the exception of $\lambda_1$ [11], whose definition consists of numerals and which Section~I.8.2 (line (C0r)) lists among the constants that depend on the weights.
No row under the marks (C0), (C4) other than the rows 1, 2 meets $\rho$ or $w$.
The letters $c_\rho$ [26], $C_E$ [27], $\mathrm{Csy}_{\mathrm{ak}}$, $\mathrm{Cpv}_0$ [28] occur neither in the fifth column of any row nor in the third column of any row of the table of Step~4.
After ``depends on'', the entries of (a) name the last three only in the rows 124, 125, and the list for the rows of kind X names them only in the rows 28, 124, 125.
Hence, among the rows up to 118, these three occur only in their own rows 27, 28.
The absolute number $c_{\mathrm{cl}}$ at the end of (a) is formed from numerals ($\sqrt{2}$ and $\pi$ count as numerals: paragraph ``Fifth column''); so it contains no letter of (a).
It occurs in no entry of (a).
No condition contains it: the inequality of a condition of (b) contains only the constant on which the condition is imposed, the variable $v$, numerals and the letters of the third column of the table of Step~4, and 7.1--7.3 list everything that the conditions of (c), (d) contain.
This proves the assertion on $c_{\mathrm{cl}}$.

\begin{named}{Remark 3.1}[Row 87: $M$ drops out]
Let $k$ be a level.
Consider an index $\tilde{t}$ of the covering lattice, that is, a bond of $\mathbb{L}_k$, a bond of $\mathbb{Z}^3$, a point of $\mathbb{Z}^3$ (a unit site), a point of $\mathbb{L}^{\mathrm{f}}_k$ (a fermion site), or an element of the $L\mathbb{Z}^3$-periodic family $\{e^{y}_{i},\delta_b\}$ on $\mathbb{Z}^3$ of \cite[Proposition~15.21(i), (ii)]{QEDI}, with its support as in Section~G4.0.4, and let $j\in\{\mathrm{(i)},\mathrm{(ii)},\mathrm{(iii)}\}$.
Let $\mathrm{tsum}_j(\tilde{t})$ denote the sum of
\[
\nu(\tilde{t}')\cdot\exp\left(-\gamma_{\mathrm{w}}\cdot\operatorname{dist}(\operatorname{supp}\tilde{t},\operatorname{supp}\tilde{t}')\right)
\]
over all indices $\tilde{t}'$ of the $j$-th index set on the covering lattice.
This sum is finite by the estimate of Step~2, row 87, which uses the torus only through (El5).
On the torus of side $\ell\in\mathrm{Cl}$ every index $t'$ has a lift $\tilde{t}'$ with $\operatorname{dist}(\operatorname{supp}\tilde{t},\operatorname{supp}\tilde{t}')=d(t,t')$, and distinct indices have distinct lifts; so the sum on the torus is at most $\mathrm{tsum}_j(\tilde{t})$.
Conversely, let $\mathrm{Fs}$ be a finite set of indices $\tilde{t}'$, and let $\ell\in\mathrm{Cl}$ exceed twice the sup-diameter of the union of $\operatorname{supp}\tilde{t}$ with the supports of the elements of $\mathrm{Fs}$.
Then any two points $c$, $c'$ of this union satisfy $|c-c'|<\ell/2\le|c-c'+\ell n|$ for every $n\in\mathbb{Z}^3$, $n\ne0$.
Hence the elements of $\mathrm{Fs}$ have distinct images on the torus, and for them $d(t,t')=\operatorname{dist}(\operatorname{supp}\tilde{t},\operatorname{supp}\tilde{t}')$.
Since $\mathrm{Cl}$ is unbounded, the supremum over $\ell\in\mathrm{Cl}$ of the sum on the torus equals $\mathrm{tsum}_j(\tilde{t})$, and
\[
c_{\gamma}=\sup_{k,\tilde{t},j}\mathrm{tsum}_j(\tilde{t}).
\]
The right-hand side contains $\gamma_{\mathrm{w}}$ and $L$ (through $\mathbb{L}_k=L^{-k}\mathbb{Z}^3$ and through the family $\{e^{y}_{i},\delta_b\}$), but not $M$.
\end{named}

\begin{named}{Remark 3.2}[Row 91: $M$ drops out, $L$ stays]
The dilation by $1/M$ maps $\mathbb{R}^3/\ell\mathbb{Z}^3$ onto $\mathbb{R}^3/(\ell/M)\mathbb{Z}^3$, the closed $M$-cubes onto the closed unit cubes centred at the points of $\mathbb{Z}^3$, polymers onto unions of unit cubes connected through faces, and a net of length $\mathrm{len}$ onto a net of length $\mathrm{len}/M$.
Therefore $d_M(X)$ (\cite[Section~13.1.7]{QEDI}) is the infimum of the lengths of the nets that meet every closed unit cube of the image of $X$.
As a set of unions of $M$-cubes, $D_k$ does not depend on $k$, and the translations by $M\mathbb{Z}^3$ permute the cubes $\square$.
The sum in the entry of row 91 therefore depends on $(\kappa,\ell/M)$ only, and
\[
c_{\kappa}=\sup\left\{\text{the sum at }(\kappa,\ell/M):\ \ell/M\in L\cdot\{13,15,17,\dots\}\right\},
\]
because $\ell/(LM)$ is an odd integer $\ge13$ (Section~I.1.1).
This expression contains $\kappa$ and $L$.
\end{named}

\begin{named}{Remark 3.3}[Letters with another meaning]
\emph{In the entries of (a).}
\begin{itemize}
\item Row 19: $p$ is the momentum, a bound variable of the supremum, and not $p$ [112].
\item Row 22: $A_1$, $A_2$ are abbreviations defined in the entry, not numbers $A_{\mathrm{ak}}$ [21].
\item Row 25: $c(1,3)$ is the number $c(\mathrm{kk},N)$ of \cite[Lemma~12.4]{QEDI} at $(\mathrm{kk},N)=(1,3)$, not a value of the function $c(\cdot)$.
\item Row 77: $x$ and $M_0$ are bound variables.
\item Row 128: $b(\cdot)$ is the functional of Section~I.7.6, not $b$ [9].
\item The condition (S2), which the entry of row 114 names: there $\sigma$ abbreviates the expression $s\cdot\left(\frac{1}{4}-9\varepsilon\right)-12\varepsilon-2a_L/N_{*}$, whose letters are $s$, $\varepsilon$, $a_L$, $N_{*}$; row 116 then gives this expression the name $\sigma$.
\end{itemize}
\emph{In the statements and sections named in the entries of the rows of the kinds D and X, in the paragraphs ``Constants'' of Propositions~30.1 and~30.3, and in the Remark that follows the paragraph ``Rows 63--66'' in Step~2 (a lower bound for the numbers $C_3$ of a general class of data), together with its proof.}
The following letters denote quantities (x2) or other objects there, and not the constants of (a) with the same names.
\begin{itemize}
\item $p$ in Propositions~1.2, 1.5, 30.1, 36, 37.1: a momentum, not $p$ [112].
\item $e_1$ in ``$g\,e_1$'': the unit vector of Section~I.1.1, not $e_1$ [118].
\item $w$ in the paragraph ``Constants'' of Proposition~30.3: a vector of $\mathbb{Z}^3$, the variable of a summation, not $w$ [2].
\item $r$ in Section~G5.1.3: an element of $\{\pm1\}^3$ or of $O_{\mathrm{cub}}$.
\item $\beta\in\{1,\dots,4\}$ in Section~I.5.8\,(i) and in Step~2, row 87: the spinor index of Section~I.1.9, bound by the sum over $t'$; it is not $\beta$ [122].
\item $c_1$, $C_1$, $C_2$ in the properties (P1)--(P3) of the penalty used in the Remark that follows the paragraph ``Rows 63--66'' and in its proof (these are the inequalities (Pen1)--(Pen3) of \cite[Definition~12.8]{QEDI}): the numbers $b\cdot(1-1/L)$, $L\cdot b$, $2L^4\cdot b$; this $C_1$ is not $C_1$ [81].
\end{itemize}
In the display (8.1) of Proposition~30.1, the part ``$=\varepsilon_5(\mathbf{m}\cdot n)$'' is the definition of row 30.
In Proposition~30.3\,(e) and in its paragraph ``Constants'', $b_-$ and $b_+$ are the end points of an interval; they are taken at the values of the rows 35, 36.
\end{named}

\begin{named}{Remark 3.4}[The letters named after ``depends on'']
For every row other than the 11 rows 11, 19, 21, 24, 25, 87, 91, 93, 94, 124, 125, the letters that its entry in (a) names after ``depends on'' are exactly the letters of its fifth column.
In these 11 rows the two sets of letters differ as follows.

\emph{Letters named after ``depends on'' that are not in the fifth column:}
$w$ [2] (row 11);
$\hat{g}$ [12] (row 19);
$c_{\mathrm{el}}$ [18], $K_{\mathrm{Str}}$ [19], $\theta_0$ [15] (row 21);
$c_{\mathrm{el}}$, $K_{\mathrm{Str}}$, $\theta_0$, $\hat{g}$, $c_{\mathrm{hl}}$ [13], $A_{\mathrm{ak}}$, $A'_{\mathrm{ak}}$ [21] (row 24);
$c_{\mathrm{el}}$, $\kappa_{(0)}$ [20] (row 25);
$\rho$ [1], $w$ (rows 93, 94);
$\mathrm{Csy}_{\mathrm{ak}}$, $\mathrm{Cpv}_0$ [28] (row 124);
$C_E$ [27] (row 125).
Each of them carries a number smaller than the number of the row whose entry names it.
Every other letter named after ``depends on'' stands in the fifth column of its row, so that its number is smaller than the number of the row (Step~3, the paragraph on the 136 rows).
Hence every letter named after ``depends on'' in an entry stands before the constant of that entry; this proves the second half of the first sentence of (e).

\emph{Letters of the fifth column that are not named after ``depends on'':}
\begin{itemize}
\item $\rho$ (row 21) and $w$ (rows 21, 24, 124, 125): by Section~I.8.1, every constant below the line (C0) may depend on $\rho$ and on the weights.
\item $\kappa_{(0)}$ [20] (row 21): $\kappa_{(0)}=\min\left\{\frac{\theta_0}{2},\frac{c_{\mathrm{el}}\cdot\theta_0}{12K_{\mathrm{Str}}}\right\}$ is a function of $\theta_0$, $c_{\mathrm{el}}$, $K_{\mathrm{Str}}$, which the entry names.
\item $\hat{g}$, $c_{\mathrm{N3}}$ [29] (rows 93, 94): $\hat{g}=\rho\cdot\lambda_{\mathrm{int}}/2$ and $c_{\mathrm{N3}}=\min\left(\frac{\rho}{\pi},\frac{\rho^3\cdot\lambda_{\mathrm{int}}}{98\sqrt{3}\cdot\pi^3}\right)$ are functions of $\rho$ and of $\lambda_{\mathrm{int}}$ [10], which is formed with the weights; the entries name $\rho$, $w$.
\item $D_{\mathrm{be}}$ (row 125): it is formed from the probes, that is, from $\beta$, $x_0$, which the entry names.
\item $M$ (rows 87, 91): it drops out (Remarks~3.1, 3.2).
\end{itemize}
Consequently, in every row the constant (in a row of kind C: the set in which it is chosen) is determined by the letters named after ``depends on'' together with $\rho$ and $w$.
This proves the assertion of the preamble of (a) on the letters named after ``depends on''.
\end{named}

\subsubsection*{Step 4. (e), second sentence}

Consider a condition of Proposition~14.2\,(b).
Besides the constant that the condition restricts, the variable $v$ and numerals (among them $\mathrm{e}$ and $\pi$), the inequality which it imposes contains only letters of (a), namely those named after ``Letters contained'' in the entry of the condition, and none for (L5) and (Z3.s1).
Here $\mathrm{lh}(v)=L\cdot M\cdot v^r$ contributes the letters $L$, $M$, $r$, and $\mathrm{lg}$ contributes the letter $L$.
In (S2) the letter $\sigma$ is a name for the left side of the inequality (Remark 3.3).
These letters are listed, with their rows, in the third column of the following table.

\medskip
\noindent
\begin{tabular}{l l p{8cm}}
condition & restricted constant [row] & letters of (a) in the inequality [row]\\
\hline
(L1) & $L$\,[33] & $\bar{\varepsilon}$\,[3]\\
(L2) & $L$\,[33] & $\bar{\varepsilon}$\,[3], $c_{\mathrm{fl}}$\,[8]\\
(L3) & $L$\,[33] & $\kappa_{0}$\,[5], $\kappa$\,[32]\\
(L4) & $L$\,[33] & $\kappa_{0}$\,[5], $c_{\mathrm{rb}}$\,[7], $\kappa$\,[32]\\
(L5) & $L$\,[33] & none\\
(M1) & $M$\,[78] & $L$\,[33]\\
(M2) & $M$\,[78] & $L$\,[33]\\
(M3) & $M$\,[78] & $\gamma_{\mathrm{w}}$\,[75]\\
(M4) & $M$\,[78] & $\alpha_{0}$\,[48], $\hat{\kappa}$\,[55]\\
(M5) & $M$\,[78] & $M^{F}$\,[76]\\
(M6) & $M$\,[78] & $M^{P}$\,[77]\\
(S1) & $\varepsilon$\,[113] & $\bar{\varepsilon}$\,[3], $\alpha$\,[31]\\
(S2) & $s$\,[114] & $N_{*}$\,[37], $a_{L}$\,[38], $\varepsilon$\,[113]\\
(A1) & $e_0$\,[117] & $c_{\mathrm{fl}}$\,[8], $M$\,[78], $p_0$\,[111], $\varepsilon$\,[113]\\
(A2) & $e_0$\,[117] & $c_{C}$\,[44], $C_{w}$\,[86], $c_{\gamma}$\,[87], $\varepsilon$\,[113]\\
(A3) & $e_0$\,[117] & $c_{X}$\,[4], $C_{T}$\,[57], $M$\,[78], $C_{w}$\,[86], $c_{\gamma}$\,[87], $\varepsilon$\,[113]\\
(A4) & $e_0$\,[117] & $C_{T}$\,[57], $C_{w}$\,[86], $\varepsilon$\,[113]\\
(A5) & $e_0$\,[117] & $M$\,[78], $\varepsilon$\,[113]\\
(A6) & $e_0$\,[117] & $C_{w}$\,[86], $c_{\gamma}$\,[87], $\varepsilon$\,[113]\\
(A7) & $e_0$\,[117] & $c_{X}$\,[4], $K_{0}$\,[6], $L$\,[33], $c_{S}$\,[95], $\varepsilon$\,[113]\\
(A8) & $e_0$\,[117] & $M$\,[78], $\varepsilon$\,[113]\\
(A9) & $e_0$\,[117] & $c_{X}$\,[4], $L$\,[33], $M$\,[78], $\varepsilon$\,[113]\\
(A10) & $e_0$\,[117] & $M$\,[78], $C_{B}$\,[89], $p$\,[112], $\varepsilon$\,[113]\\
(Z2.s1) & $e_0$\,[117] & $L$\,[33], $\varepsilon$\,[113]\\
(Z2.s2) & $e_0$\,[117] & $C'_{V}$\,[56], $a'$\,[74], $\varepsilon$\,[113]\\
(Z2.s3) & $e_0$\,[117] & $M$\,[78], $\varepsilon$\,[113]\\
(Z2.s4) & $e_0$\,[117] & $\varepsilon$\,[113]\\
(Z2.s5) & $e_0$\,[117] & $c_{\mathrm{rb}}$\,[7], $L$\,[33], $M$\,[78], $C_{d}$\,[90], $\varepsilon$\,[113]\\
(Z3.s1) & $e_0$\,[117] & none\\
(e1) & $e_1$\,[118] & $L$\,[33], $r$\,[110], $p_0$\,[111]\\
(e2) & $e_1$\,[118] & $c_{\mathrm{fl}}$\,[8], $L$\,[33], $\varepsilon$\,[113]\\
(e3) & $e_1$\,[118] & $c_{\mathrm{fl}}$\,[8], $L$\,[33], $b_{-}$\,[35], $M$\,[78], $c_{\kappa}$\,[91], $C_{0}$\,[92], $r$\,[110], $\varepsilon$\,[113]\\
(e4) & $e_1$\,[118] & $L$\,[33], $M$\,[78], $C_{0}$\,[92], $r$\,[110], $p$\,[112], $\varepsilon$\,[113]\\
(e5) & $e_1$\,[118] & $L$\,[33], $M$\,[78], $C_{0}$\,[92], $r$\,[110], $p_0$\,[111], $p$\,[112]\\
(e6) & $e_1$\,[118] & $M$\,[78], $\gamma_{\Sigma}$\,[93], $C_{\Sigma}$\,[94], $r$\,[110]\\
(e7) & $e_1$\,[118] & $L$\,[33], $M$\,[78], $r$\,[110]\\
(e8) & $e_1$\,[118] & $N_{*}$\,[37], $M$\,[78], $r$\,[110]\\
\end{tabular}
\medskip

In each row, every number in the third column is smaller than the number of the restricted constant.
This proves the second sentence of (e).

\subsubsection*{Step 5. (c): the conditions hold whenever $e\cdot h^{\frac{1}{2}}\le e_1$}

\paragraph{5.1 (The number $e_1$).}
Let $\mathrm{Sadm}_1$ be the set of those $e'\in\left(0,\min\left(e_0,\frac{1}{3}\right)\right]$ for which (e1)--(e8) hold for all $v\ge-\log e'$.
If $e'\in\mathrm{Sadm}_1$ and $0<e''<e'$, then $-\log e''>-\log e'$, and therefore $e''\in\mathrm{Sadm}_1$.
According to row 118, $e_1=\frac{1}{2}\sup\mathrm{Sadm}_1$, and $e_1=0$ if $\mathrm{Sadm}_1$ is empty.
In the latter case no $e>0$ satisfies $e\cdot h^{\frac{1}{2}}\le e_1$, and (c) is void.
(The set $\mathrm{Sadm}_1$ is not empty, by Proposition~14.3\,(f).)
Otherwise $0<e_1<\sup\mathrm{Sadm}_1$, so that some $e'\in\mathrm{Sadm}_1$ satisfies $e'>e_1$; consequently $e_1\in\mathrm{Sadm}_1$, that is,
\begin{equation}\tag{5.1}
\text{(e1)--(e8) hold for all } v\ge-\log e_1 .
\end{equation}

\paragraph{5.2 (Data).}
Let $\bar{m}\in(0,1]$, let $e>0$ with $e_K=e\cdot h^{\frac{1}{2}}\le e_1$, let $T\ge0$ be an integer with $R=h\cdot L^T\in\mathrm{Tor}(e)$, and let $N\ge N_{\min}(R,e)$ be an integer.
The quantities $a_{k}$, $e_k$, $v_k$, $\mu_k$, $K$, $r_k$, $p_k$, $p_{0,k}$, $\ell_k$, $k_{*}=k_{*}(R)$ and the index $*$ are those of Sections~I.6.1 and I.6.2.
The number $q=e\cdot\eta^{\frac{1}{2}}$ is the number $e_k$ at $k=0$; it is not the ceiling $e_0$.
One has $v_K=-\log e_K\ge-\log e_1$.
By (5.1), the condition (e1) holds at $v=v_K$:
\[
v_K\ge8,\qquad\left(1+\frac{\mathrm{lg}}{v_K}\right)^r\le\left(1+\frac{3\,\mathrm{lg}}{v_K}\right)^r<L .
\]
Together with $R\in\mathrm{Tor}(e)$ they give the hypotheses of \cite[Proposition~19.1]{QEDI}, taken with $h'=h$ and $m_0=0$ (Section~I.6.1); the proof of that proposition uses that $M$ is a power of $L$, which is (M1).
Its clauses (c) and (e), which are (G4.Sc1) and (G4.Sc3), yield
\begin{equation}\tag{5.2}
\begin{aligned}
&1\le k_{*}\le K,\qquad v_{*}\ge v_K,\qquad e_k=e_{*}\cdot L^{-\frac{k_{*}-k}{2}}\le e_{*}\ \ \text{for } k\le k_{*},\qquad v_{*}^r\le r_{*}<L\cdot v_{*}^r,\\
&\ell_*=M\cdot r_{*},\ \text{ so that }\ M\cdot v_{*}^r\le\ell_*<L\cdot M\cdot v_{*}^r=\mathrm{lh}(v_{*}),\\
&\mu_*^{-1}<L\cdot\ell_*<L\cdot\mathrm{lh}(v_{*}),\qquad e^2\cdot R=e_{*}^2\cdot\ell_* .
\end{aligned}
\end{equation}
Here $e_k=e\cdot a_{k}^{\frac{1}{2}}$ and $a_{k}=L^{k-N}$ (Section~I.6.1).
Moreover $e_{*}=\exp(-v_{*})$, $p_{*}=v_{*}^p$ and $p_{0*}=v_{*}^{p_0}$.
Since $v_{*}\ge v_K\ge-\log e_1$, (5.1) gives (e1)--(e8) at $v=v_{*}$.

\paragraph{5.3 (The conditions).}
The conditions of (c) follow from (5.2) and from (e1)--(e7) at $v=v_{*}$.
\begin{itemize}
\item[(2.1)]
By (5.2) and the second inequality of (e1),
$r_{*}\ge v_{*}^r>31L+16$.
\item[(fl)]
By (5.2), $e_k^{\varepsilon}=e_{*}^{\varepsilon}\cdot L^{-\frac{(k_{*}-k)\varepsilon}{2}}$ for $k\le k_{*}$.
Summing the geometric series, whose ratio satisfies $L^{-\frac{\varepsilon}{2}}<1$, and using (e2), one obtains
\[
\sum_{k\le k_{*}}\mathrm{sm}_k=\left(c_{\mathrm{fl}}+2\right)\cdot e_{*}^{\varepsilon}\cdot\sum_{i=0}^{k_{*}}L^{-\frac{i\varepsilon}{2}}\le\left(c_{\mathrm{fl}}+2\right)\cdot\exp(-\varepsilon v_{*})\cdot\left(1-L^{-\frac{\varepsilon}{2}}\right)^{-1}\le\frac{1}{16}.
\]
The inequality $v_{*}\ge8$ is the first inequality of (e1).
\item[(T1)]
The inequalities $v_{*}\ge8$ and $p_0\cdot v_{*}^{p_0-1}\ge28$ are the first and the fourth inequality of (e1).
By (5.2) at $k=0$, $q\le e_{*}=\exp(-v_{*})<1$, and $\frac{3}{4}-8\varepsilon>0$ because $\varepsilon\le\frac{1}{200}$; hence $w^{\mathrm{m}}_0=q^{\frac{3}{4}-8\varepsilon}<1$.
\item[(T2)]
By (5.2) and by the first, the second and the third inequality of (e3), respectively,
\begin{align*}
8C_0^2\cdot\ell_*^6\cdot\left(b_{-}^{-1}+\mu_*^{-1}\right)&\le8C_0^2\cdot\mathrm{lh}(v_{*})^6\cdot\left(b_{-}^{-1}+L\cdot\mathrm{lh}(v_{*})\right)\le\exp\left(\frac{\varepsilon v_{*}}{2}\right)=e_{*}^{-\frac{\varepsilon}{2}},\\
4\left(2C_0+c_{\kappa}+5c_{\mathrm{fl}}+6\right)\cdot\ell_*^3&\le4\left(2C_0+c_{\kappa}+5c_{\mathrm{fl}}+6\right)\cdot\mathrm{lh}(v_{*})^3\le\exp(\varepsilon v_{*})=e_{*}^{-\varepsilon},\\
e_{*}^{\frac{1}{4}-9\varepsilon}&=\exp\left(-\left(\frac{1}{4}-9\varepsilon\right)v_{*}\right)\le\frac{1}{8}.
\end{align*}
\item[(T3)]
By (5.2) and (e4),
$4C_0\cdot\ell_*^2\cdot p_{*}\le4C_0\cdot\mathrm{lh}(v_{*})^2\cdot v_{*}^p<\exp(\varepsilon v_{*})=e_{*}^{-\varepsilon}$.
\item[(T4)]
By (5.2) and (e5),
\begin{align*}
&12\ell_*^3\cdot\left(\exp\left(-\frac{p_{0*}}{2}\right)+\exp\left(-\frac{p_{*}^2}{2C_0^2}\right)\right)\\
&\qquad\le12\,\mathrm{lh}(v_{*})^3\cdot\left(\exp\left(-\frac{v_{*}^{p_0}}{2}\right)+\exp\left(-\frac{v_{*}^{2p}}{2C_0^2}\right)\right)\\
&\qquad\le\exp(-v_{*})=e_{*}.
\end{align*}
\item[(T5)]
One has $\gamma_{\Sigma}>0$ (Lemma 1, row 93) and $\ell_*\ge M\cdot v_{*}^r$; hence, by (e6),
$C_{\Sigma}\cdot\exp(-\gamma_{\Sigma}\cdot\ell_*)\le C_{\Sigma}\cdot\exp\left(-\gamma_{\Sigma}\cdot M\cdot v_{*}^r\right)\le\exp(-v_{*})=e_{*}$.
\item[(T6)]
By (5.2) and (e7),
$L^2\cdot e^2\cdot R=L^2\cdot e_{*}^2\cdot\ell_*\le L^2\cdot\exp(-2v_{*})\cdot\mathrm{lh}(v_{*})\le1$.
\end{itemize}
The constants occurring in (e1)--(e7) contain neither $(e,T,N)$ nor $\bar{m}$.
Therefore the conditions hold for every $\bar{m}\in(0,1]$, every $e$ with $e\cdot h^{\frac{1}{2}}\le e_1$, every $R\in\mathrm{Tor}(e)$ and every $N\ge N_{\min}(R,e)$.
This proves (c).
Clause (c) corresponds to \cite[Proposition~19.2]{QEDI}; the latter is formulated with $\gamma_1$ in place of $\gamma_{\Sigma}$, with $L^{m_0+1}\cdot\mathrm{lh}(v)$ in place of $L\cdot\mathrm{lh}(v)$, and with the constants and the number $e_1$ of \cite[Section~19.1.2]{QEDI}.

\subsubsection*{Step 6. (d): (W3), (Z8.w), $e_G$, $\mathrm{pp}$, (Z10.s4)}

\paragraph{6.1 ($h$, $e_G$).}
Since $\bar{m}\in(0,1]$, one has $\log_L(1/\bar{m})\ge0$.
Hence $\lfloor\log_L(1/\bar{m})\rfloor$ is an integer $\ge0$, $h$ is one of the numbers $1,L,L^2,\dots$, and
\[
h\ge1,\qquad\log h\ge0 .
\]
For $e>0$ the two conditions in row 121 take the form
\[
e\le e_1\cdot h^{-\frac{1}{2}}\qquad\text{and}\qquad e\le\left(h^{\frac{1}{2}}\cdot\log h\right)^{-1},
\]
the second one being void if $h=1$.
Consequently, if $e_1>0$ (Proposition~14.3\,(f)), the set in row 121 is the interval $(0,e_G]$, where
\begin{equation}\tag{6.1}
e_G=\min\left\{e_1\cdot h^{-\frac{1}{2}},\ \left(h^{\frac{1}{2}}\cdot\log h\right)^{-1}\right\}\qquad(=e_1\ \text{if } h=1).
\end{equation}
It follows that $0<e\le e_G$ holds if and only if both $e\cdot h^{\frac{1}{2}}\le e_1$ and (Z8.w) hold; in particular (Z8.w) is contained in $e\le e_G$.

\paragraph{6.2 ((W3)).}
Let $e>0$ satisfy $e_K=e\cdot h^{\frac{1}{2}}\le e_1$ and (Z8.w), that is, $e_K\cdot\log h\le1$.
By (5.1), (e1) and (e8) hold at $v=v_K=-\log e_K$.
From $v_K\ge8\ge1$ one gets $v_K^r\ge1$.
The number $r_K=\min\{L^i:\ i\ge0,\ L^i\ge v_K^r\}$ satisfies $r_K/L<v_K^r$: if $r_K=L^i$ with $i\ge1$, then $L^{i-1}<v_K^r$ by minimality, and if $r_K=1$, then $1/L<1\le v_K^r$.
By Section~I.6.2, $R^{*}=R^{*}(e)=h\cdot M\cdot r_K/L$.
Therefore
\[
e^2\cdot R^{*}=e_K^2\cdot M\cdot\frac{r_K}{L}<e_K^2\cdot M\cdot v_K^r .
\]
Since $e=e_K\cdot h^{-\frac{1}{2}}$, one has $\log(1/e)=v_K+\frac{1}{2}\log h>0$.
With $e_K=\exp(-v_K)$, (Z8.w) and (e8) it follows that
\begin{align*}
6N_{*}\cdot e^2\cdot R^{*}\cdot\log\frac{1}{e}&\le6N_{*}\cdot M\cdot v_K^r\cdot e_K\cdot\left(e_K\cdot v_K+\frac{1}{2}\,e_K\cdot\log h\right)\\
&\le6N_{*}\cdot M\cdot v_K^r\cdot\exp(-v_K)\cdot\left(v_K\cdot\exp(-v_K)+\frac{1}{2}\right)\le\pi^2 .
\end{align*}
This is (W3).
By 6.1 it holds for every $e\in(0,e_G]$.

\paragraph{6.3 ($R^{*}(e)$, $a_{*}$, $\mathrm{pp}$).}
Let $e>0$ with $e\cdot h^{\frac{1}{2}}\le e_1$.
By (M1), $M$ is a power of $L$ with $M\ge L^2$, and $r_K$ is a power $L^i$ with $i\ge0$.
Hence $T^{*}(e)=\log_L(M\cdot r_K)-1$ is an integer $\ge1$, and $R^{*}(e)=h\cdot L^{T^{*}(e)}=R_{\mathrm{top}}(e)/L\le R_{\mathrm{top}}(e)$; thus $R^{*}(e)\in\mathrm{Tor}(e)$ (Section~I.6.2).
By (5.1), (e1) holds at $v=v_K$:
\[
v_K\ge8,\qquad\left(1+\frac{\mathrm{lg}}{v_K}\right)^r\le\left(1+\frac{3\,\mathrm{lg}}{v_K}\right)^r<L .
\]
These are the hypotheses of \cite[Proposition~19.1(g)]{QEDI}, which is (G4.Sc5); its proof uses $M\ge L^2$, that is, (M1).
On $R=R^{*}(e)$ and for $N\ge N_0(e)$ it gives
\[
a_{*}\in\left\{\frac{h}{L},\ \frac{h}{L^2}\right\}.
\]
The number $\mathrm{pp}:=\left(D_\beta/a_{*}+2M+4\right)^3$ (Section~G5.1.4) is formed on $R=R^{*}(e)$, and $D_\beta>0$.
Hence $D_\beta/a_{*}\le D_\beta\cdot L^2/h$, and
\[
\mathrm{pp}\le\left(\frac{D_\beta\cdot L^2}{h}+2M+4\right)^3 .
\]

\paragraph{6.4 ((Z10.s4)).}
One has $C_{\bar{\kappa}}\ge0$: by rows 130--134, $s_{B}\ge0$, $C_{3\mathrm{p}}\ge0$ and $C_{\Gamma}^{(16)}>0$, so that $C_{U}\ge0$.
Hence $C'_{\bar{\kappa}}=8C_{\bar{\kappa}}+3C_{\bar{\kappa}}^2\ge C_{\bar{\kappa}}$, and the inequality $\frac{5}{2}C'_{\bar{\kappa}}\cdot e_2^{\sigma}\le b(f)/4$ gives $\frac{5}{2}C_{\bar{\kappa}}\cdot e_2^{\sigma}\le b(f)/4$.
Multiplying the inequality $\frac{5}{2}C'_{\bar{\kappa}}\cdot e_2^{\sigma}\le b(f)/4$ by $\frac{2}{5}$ gives $C'_{\bar{\kappa}}\cdot e_2^{\sigma}\le b(f)/10$.

\subsubsection*{Step 7. (e), third sentence}

\paragraph{7.1 (Quantities formed with the data).}
By Sections~I.6.1 and I.6.2 the quantities $\eta=L^{-N}$, $q$, $a_{k}$, $e_k$, $v_k$, $\mu_k$, $K$, $r_k$, $p_k$, $p_{0,k}$, $R=h\cdot L^T$, $\ell_k$, $\mathfrak{n}_k$, $k_{*}(R)$, the corresponding quantities with the index $*$, $R^{*}(e)$, $r_K(e)$, the width $\ell_c=h\cdot M\cdot r_K(e)/(2L^2)$ (Section~I.6.4) and $w^{\mathrm{m}}_0=q^{\frac{3}{4}-8\varepsilon}$ (Section~I.5.4) are formed from the data $(e,T,N)$ and from the constants $L$\,[33], $M$\,[78], $r$\,[110], $p_0$\,[111], $p$\,[112], $\varepsilon$\,[113], $\bar{m}$\,[119], $h$\,[120].
In (d) the index $*$ refers to the torus $R=R^{*}(e)$ (Section~G5.1.1).
For $e\cdot h^{\frac{1}{2}}\le e_1$ this torus lies in $\mathrm{Tor}(e)$ (6.3).
On it the numbers $K-k_{*}$, $a_{*}$, $e_{*}$, $v_{*}=-\log e_{*}$, $\ell_*$, $r_{*}=\ell_*/M$, $p_{*}=v_{*}^p$, $p_{0*}=v_{*}^{p_0}$ do not depend on $N\ge N_0(e)$, by \cite[Proposition~19.1(c)]{QEDI}, which is (G4.Sc1); the hypotheses of that proposition were verified in 5.2.

\paragraph{7.2 ((c)).}
Besides the quantities of 7.1, the conditions and the error terms of (c) contain the following letters:
\begin{description}
\item[(2.1):] $L$;
\item[(fl):] $c_{\mathrm{fl}}$, $\varepsilon$;
\item[(T1):] $p_0$;
\item[(T2):] $C_0$, $b_{-}$, $c_{\kappa}$, $c_{\mathrm{fl}}$, $\varepsilon$;
\item[(T3):] $C_0$, $\varepsilon$;
\item[(T4):] $C_0$;
\item[(T5):] $C_{\Sigma}$, $\gamma_{\Sigma}$;
\item[(T6):] $L$;
\item[$\varepsilon_1$:] $\varepsilon$;
\item[$\varepsilon_2$:] $C_{\Sigma}$, $\gamma_{\Sigma}$;
\item[$\varepsilon_3$, $\varepsilon_4$:] $C_0$;
\item[$\varepsilon_5$:] the function of row 30.
\end{description}
All of them are constants of (a).

\paragraph{7.3 ((d)).}
The condition (W3) contains $N_{*}$, and (Z8.w) contains $h$.
The following quantities of (d) depend on $e$, with the exception of $\mathrm{mxL}$, which is a function of $\mathbf{D}$ alone.
Besides the quantities of 7.1, they contain the following letters:
\begin{description}
\item[$\tau_s$:] $c_{\tau}$;
\item[$\bar{\kappa}_1$:] $\varepsilon$;
\item[$\bar{r}$:] $h_{G}$, $C_K$, $K_{0}$, $\varepsilon$;
\item[$\mathrm{pp}$ (Section~G5.1.4):] $D_\beta$, $M$;
\item[$B_1$:] $\mathrm{gGa}_{1}$, $C_K'$, $\mathrm{pp}$, $\varepsilon$;
\item[$B_{\mathrm{wr}}$:] $C_{d}$, $M$, $K_{0}$, $\mathrm{pp}$;
\item[$B_0$:] $C_{\mathrm{fr}}$, $\varepsilon$, $B_{\mathrm{wr}}$, $B_1$;
\item[$\varpi_{\mathrm{sep}}^{+}$:] $N_{*}$, $\mathrm{CWj}$, $\mathrm{pp}$, $\mathrm{vt}$, $M$, $L$;
\item[$\varpi_0$:] $C_0$;
\item[$t_2$:] $N_{*}$, $C_{\beta}$, $\mathbf{D}$;
\item[$t_x$, $\mathbf{r}$:] $\varepsilon$, $\mathrm{Cb}$;
\item[$\varpi_L$:] $\varpi_0$, $\mathbf{D}$;
\item[$\mathrm{mxL}$:] $\mathbf{D}$;
\item[$\lambda_L$:] $\varpi_L$, $\mathrm{mxL}$, $N_{*}$, $\mathbf{r}$, $t_2$.
\end{description}
Besides these quantities, the conditions contain the following:
\begin{description}
\item[(g1):] $\mathrm{hL}$, $h_{G}$, $\varepsilon$, $M$, $\mathrm{alG}$;
\item[(g2):] $M$, $L$;
\item[(g3):] $C_0$, $\mathrm{gGa}_{1}$, $\varepsilon$, $C_K$, $K_{0}$;
\item[(g4):] $\varepsilon$, $\mu_-$, $\mathrm{gGa}_{1}$, $C_K'$;
\item[(g5):] $N_{*}$;
\item[(g6):] $N_{*}$, $s$;
\item[(g7):] $N_{*}$;
\item[(g8):] $s$;
\item[(g9):] $D_\beta$;
\item[(Pb):] $\bar{m}$; $\beta$, through $f=dd^{*}\beta$, $a=d^{*}\beta$, $\|a\|_{\infty}$ and the number $b(f)$, which is formed from $\bar{m}$ and $\beta$ (Section~I.7.6); $x_0$; and the radius $r_\beta$, also through $d_0=|x_0|_2-2r_\beta$;
\item[(Z10.s1):] $e_2$, $e_G$;
\item[(Z10.s2):] $e_2$, $e_{\sharp}$, and $\beta$, $x_0$ (through the three probes);
\item[(Z10.s4):] $C'_{\bar{\kappa}}$, $e_2$, $\sigma$, and the number $b(f)$, in the inequality;
further $C_{\bar{\kappa}}$\,[134], which occurs in the display only in its second formula $C'_{\bar{\kappa}}=8C_{\bar{\kappa}}+3C_{\bar{\kappa}}^2$, and this formula repeats the definition of row 135.
\end{description}
All of these are constants of (a), with one exception: the radius $r_\beta$, a number $>0$ such that $\operatorname{supp}\beta$ lies in a ball of radius $r_\beta$.
It is chosen together with $(\beta,x_0)$ (data of Proposition~43.2); it is not a row of (a), it occurs in no cell of (a), and among the conditions of (b), (c), (d) only (Pb) contains it.
The letters named here for (g1)--(g9) and for the quantities occurring in them are the following 23:
$K_{0}$, $L$, $N_{*}$, $\mu_-$, $M$, $C_{d}$, $C_0$, $\mathrm{hL}$, $C_K$, $C_K'$, $\mathrm{vt}$, $h_{G}$, $\mathrm{gGa}_{1}$, $\mathrm{alG}$, $c_{\tau}$, $\mathrm{CWj}$, $\varepsilon$, $s$, $D_\beta$, $C_{\beta}$, $\mathbf{D}$, $\mathrm{Cb}$, $C_{\mathrm{fr}}$.
Together with $r$, $p_0$, $p$, which enter through the quantities $r_{*}$, $\ell_*$, $\ell_c$, $p_{0*}$, $p_{*}$ of 7.1, and with $e_1$ and $h$, which enter through the range $\left(0,e_1\cdot h^{-\frac{1}{2}}\right]$ of $e'$ in the definition of $e_{\sharp}$ (and $h$ also through $a_{*}$, $e_{*}$, $\ell_c$), they are the 28 letters of row 127 in the table of Step 3.
The constant $\bar{m}$\,[119] enters (g1)--(g9) through $h$ only: $\mu_*$ occurs in none of them.

\paragraph{7.4 (No placement).}
Among the quantities of (d), only $d_{\mathrm{sep}}$ and, through it, $\varpi_{\mathrm{sep}}$ depend on the placement (Sections~G5.1.4 and G5.1.8).
Neither of them occurs in a condition: (g2) and (g5) are formulated with $M\cdot r_{*}/(8L)-1$ and with $\varpi_{\mathrm{sep}}^{+}$.
The numbers $\mathrm{pp}$, $t_x$, $t_2$ contain $D_\beta$, $\mathrm{Cb}$, $C_{\beta}$, $\mathbf{D}$, which are formed without a placement (Section~G5.1.3: suprema over $\mathrm{Ind}_\beta$, and the cube $Q_\beta$).
The condition (Pb) contains $\bar{m}$, $\beta$, $x_0$, $r_\beta$ only.
Hence no condition of (c), (d) contains a placement.
Together with 7.1--7.3 this shows that the conditions of (c), (d) contain only constants of (a) and the data $(e,T,N)$, and that (Pb) contains, besides these, the radius $r_\beta$, which is chosen together with $(\beta,x_0)$ and is not a row of (a).
This proves the third sentence of (e).

\subsubsection*{Step 8. (f)}

\paragraph{(f1) From (L1).}
One has $\frac{1}{8}-\frac{3\bar{\varepsilon}}{2}=\frac{47}{400}$.
Since $L\ge1$, (L1) says $L\ge2400^{\frac{400}{47}}$, and $\frac{400}{47}>8$, $2400>1$ give $L\ge2400^8$.
This has the following consequences:
\begin{gather*}
L\ge2400^8\ge4^8=65536>8>5,\qquad L^{\frac{1}{4}}\ge2400^2\ge2,\qquad L^{\frac{3}{4}}\ge2400^6\ge110,\\
L^{-\frac{5}{8}}\le2400^{-5}\le\frac{1}{4},\qquad L^{-\frac{3}{4}}\le2400^{-6}\le\frac{1}{2},\qquad L^{-3+\frac{1}{8}}=L^{-\frac{23}{8}}\le L^{-1}\le\frac{1}{2}.
\end{gather*}
The inequality $L^3\cdot\exp(-L)\le L^{-\frac{1}{4}}$ is equivalent to $\frac{13}{4}\log L\le L$; by (El2), $\frac{13}{4}\log L\le\frac{13}{4}L^{\frac{1}{2}}\le L$, because $L^{\frac{1}{2}}\ge2400^4\ge\frac{13}{4}$.
Let $0<\varepsilon\le\bar{\varepsilon}$.
Then $\frac{1}{4}-\frac{3\varepsilon}{2}\ge\frac{1}{4}-\frac{3}{400}=\frac{97}{400}\ge\frac{47}{400}$, so that $L^{\frac{1}{4}-\frac{3\varepsilon}{2}}\ge L^{\frac{47}{400}}\ge2400\ge13.6$.
Moreover $-\frac{7}{8}+\frac{\varepsilon}{2}\le-\frac{7}{8}+\frac{1}{400}\le-\frac{1}{2}$; therefore $4L^{-\frac{7}{8}+\frac{\varepsilon}{2}}\le4L^{-\frac{1}{2}}\le4\cdot2400^{-4}\le\frac{1}{4}$, further $L^{-\frac{7}{4}}\le\frac{1}{4}$, and the sum $L^{-\frac{7}{4}}+4L^{-\frac{7}{8}+\frac{\varepsilon}{2}}$ is at most $\frac{1}{2}<1$.

\paragraph{(f2) From (L3), (C2).}
One has $\kappa_{0}>0$ and $L\ge1$.
For $j=1,2$, (L3) gives
\[
L\cdot(\kappa-\kappa_{0}-j)=L\cdot(\kappa-6\kappa_{0}-4)+L\cdot(5\kappa_{0}+4-j)\ge\kappa+1+L\cdot(5\kappa_{0}+2)\ge\kappa .
\]
These are the inequalities $L\cdot(\kappa-\kappa_{0}-1)\ge\kappa$ and $L\cdot(\kappa-\kappa_{0}-2)\ge\kappa$, the two consequences stated in the row (L3).
Further $\kappa\ge7\kappa_{0}+7\ge\kappa_{0}$, and $\kappa_{\mathrm{w}}-1=\kappa+\kappa_{0}+2\ge\kappa+1$.

\paragraph{(f3) Rates.}
One has $\gamma^{F}\le1/(6L)$ (Lemma 1, row 73), that is, $\gamma^{F}\cdot L\le\frac{1}{6}$; and $\frac{1}{6}<0.6927\le\log2$ (El4).
The rate $\gamma_{\mathrm{w}}=\min(\gamma^{F},\gamma^{P})$ satisfies $\gamma_{\mathrm{w}}\le\gamma^{F}$ and $\gamma_{\mathrm{w}}\le\gamma^{P}$.
Next, $\gamma^{P}=\gamma_{P}/2\le\gamma_{P}\le1$, by Lemma 1, rows 50, 51, by the definitions of $\gamma_{P}$ and of $\gamma^{P}$ in \cite[Section~15.7]{QEDI} and \cite[Section~15.10.3]{QEDI}, and by the paragraph on the constants $C_{P}$, $\gamma_{P}$, $c_P$ at the end of \cite[Section~15.7]{QEDI}.
By (M3), $\gamma_{\mathrm{w}}\cdot M\ge16\ge1$, and so $\gamma^{F}\cdot M\ge\gamma_{\mathrm{w}}\cdot M\ge16$.
Finally $27\alpha_{0}=\frac{27}{54}=\frac{1}{2}$.

\paragraph{(f4) $a'$, $C_{5}$, $C'_{V}$, $M$.}
The inequalities $a'\ge1$ and $C_{5}\ge1$ are entries of Lemma 1, rows 74 and 65.
Since $b>0$, $C'_{V}=2\mathrm{e}\cdot(42+6b)\ge84\mathrm{e}\ge24\mathrm{e}$.
Put $\mathrm{xv}:=108\pi\cdot C'_{V}\cdot a'$.
By (El1), $\mathrm{xv}\ge108\cdot3\cdot24\cdot2\cdot1=15552\ge1600$.
By (M5) and $\mathrm{xv}\ge1$, $M\ge M^{F}\ge\mathrm{xv}^2\ge\mathrm{xv}$.
So $M\ge108\pi\cdot C'_{V}\cdot a'\ge1600$, and in particular $M\ge9$.
One has $\mathrm{qq}(\mathrm{wp})=54\pi\cdot C'_{V}\cdot a'\cdot\mathrm{wp}/M=\mathrm{xv}\cdot\mathrm{wp}/(2M)$ (Section~G1.10; in this form it enters Proposition~2.7\,(e)).
Further $M^{27\alpha_{0}}=M^{\frac{1}{2}}$, and $\mathrm{qq}\left(M^{\frac{1}{2}}\right)=\frac{1}{2}\,\mathrm{xv}\cdot M^{-\frac{1}{2}}\le\frac{1}{2}$, because $M\ge\mathrm{xv}^2$ by (M5).
Finally $\mathrm{qq}(\exp(27\kappa'))=\mathrm{xv}\cdot\exp(27\kappa')/(2M)\le\frac{1}{2}$, because $M\ge\mathrm{xv}\cdot\exp(27\kappa')$ by (M5).

\paragraph{(f5) Constants of the line (C6).}
One has $C^{F}=2\max\{C_{1},\dots\}\ge2C_{1}$, and $C^{F}\ge2\cdot304\mathrm{e}\cdot\max(1,b)\cdot a'\ge608\cdot2\ge4$.
Further $C_{w}^{(0)}=\max(C^{F},C^{P})\ge C^{F}$, $C_{w}^{(1)}=\max\left(C_{w}^{(0)},\dots\right)\ge C_{w}^{(0)}$, and $C_{w}=64K_{0}\cdot C_{w}^{(1)}\ge C_{w}^{(1)}$, because $K_{0}\ge1$ (Lemma 1, row 6) and $C_{w}^{(1)}>0$.
From $M\ge1$ and $402\alpha_{0}>0$ one gets $M^{402\alpha_{0}}\ge1$; hence $C^{P}=116C_{P}\cdot M^{402\alpha_{0}}\ge C_{P}$ (as $C_{P}>0$), and $C^{P}\le C_{w}^{(0)}\le C_{w}$.
By the definition of $C_{w}^{(1)}$, $4\exp(2\gamma_{\mathrm{w}})\cdot M^{27\alpha_{0}}\le C_{w}^{(1)}$.
Since $C_{0}\ge C^{F}\ge4$, one has $C_{0}\ge\max(1,C^{F})$.
The inequalities $c_{\kappa}\le K_{0}$, $C_{T}\ge1$, $c_{\gamma}\ge1$, $\kappa_{0}\ge c_{X}$ are entries of Lemma 1, rows 91, 57, 87, 5.

\paragraph{(f6) (C1), (C7).}
One has $0<\frac{199}{200}<\alpha<1$.
The inequality $1-\alpha<\varepsilon$ of (S1) means $1-\varepsilon<\alpha$, and $\varepsilon\le\bar{\varepsilon}=\frac{1}{200}$.
Further $p_0\ge9\ge2$, and $p>2p_0$ gives $2p>4p_0>p_0$ (Lemma 1, rows 111, 112).

\paragraph{(f7) $e_0$, $e_1$.}
The first inequality of (A8) at $v=-\log e_0$ says $e_0^{\varepsilon}\le\frac{1}{4}$, that is, $e_0\le4^{-\frac{1}{\varepsilon}}$; since $\frac{1}{\varepsilon}\ge200$, it follows that $e_0\le4^{-200}$.
For $e_k\in(0,e_0]$ one therefore has, by (El4),
\[
v=-\log e_k\ge200\log4=400\log2\ge400\cdot0.6927=277.08\ge277 .
\]
Hence $v\ge9$, $v\ge1$, and $p_{0,k}=v^{p_0}\ge1$.
By row 118, either $e_1=\frac{1}{2}\sup\mathrm{Sadm}_1$, where $\mathrm{Sadm}_1\subseteq\left(0,\min\left(e_0,\frac{1}{3}\right)\right]$ by its definition in 5.1, or $e_1=0$ if $\mathrm{Sadm}_1$ is empty.
In both cases $e_1\le\frac{1}{2}\min\left(e_0,\frac{1}{3}\right)$ (recall $e_0>0$), so that $e_1\le e_0$ and $e_1\le\frac{1}{3}$.

\paragraph{(f8) Conditions that are implied by others.}
\emph{(L1) follows from (L2).}
One has $\frac{3}{8}-6\bar{\varepsilon}=\frac{69}{200}$; so (L2) says $L^{\frac{69}{200}}\ge16c_{\mathrm{fl}}$, and $c_{\mathrm{fl}}\ge72^{54}$ (Lemma 1, row 8).
Moreover $\frac{47}{400}\cdot\frac{200}{69}=\frac{47}{138}\ge\frac{1}{3}$, as $141\ge138$.
Hence
\[
L^{\frac{47}{400}}=\left(L^{\frac{69}{200}}\right)^{\frac{47}{138}}\ge\left(72^{54}\right)^{\frac{47}{138}}\ge\left(72^{54}\right)^{\frac{1}{3}}=72^{18}\ge2400 .
\]
\emph{Further implications.}
The bound $L\ge5$ follows from (L1), by (f1).
The bound $\gamma_{\mathrm{w}}\cdot M\ge1$ follows from (M3), by (f3).
The condition (M2) follows from (M5): $M\ge M^{F}\ge4L+1>4L$.

The condition (Z2.s3), the third inequality of (A8) and the fifth inequality of (A1) are one and the same inequality $3M\cdot\exp(-\varepsilon v)\le1$.
The first inequality of (Z2.s4) and the first inequality of (A8) are both $\exp(-\varepsilon v)\le\frac{1}{4}$; the second inequality of (Z2.s4) and the fourth inequality of (A8) are both $\exp\left(-\frac{v}{4}\right)+\exp(-\varepsilon v)\le1$.

\emph{The second inequality of (A8) follows from its first and third inequality.}
By the first inequality of (A8), $\frac{1}{\varepsilon}\ge200$ and (El4), one has $v\ge277>0$ (see (f7)).
Moreover $\varepsilon\le\frac{1}{200}\le\frac{1}{4}$.
Hence
\begin{gather*}
3M\cdot\exp(-2\varepsilon v)=\left(3M\cdot\exp(-\varepsilon v)\right)\cdot\exp(-\varepsilon v)\le\frac{1}{4},\\
3M\cdot\exp\left(-\left(\frac{1}{4}+\varepsilon\right)v\right)=\left(3M\cdot\exp(-\varepsilon v)\right)\cdot\exp\left(-\frac{v}{4}\right)\le\exp\left(-\frac{v}{4}\right)\le\exp(-\varepsilon v)\le\frac{1}{4},
\end{gather*}
and the sum of the two left sides is at most $\frac{1}{2}\le1$.

\emph{(Z3.s1) and the first inequality of (A10) follow from the first inequality of (A8)}, by (f7).

\emph{The fourth inequality of (A9) follows from the third inequality of (A8) and (M4).}
The third inequality of (A8) says $\exp(\varepsilon v)\ge3M$; hence $\varepsilon v\ge\log(3M)\ge\log M\ge0$ (as $M\ge1$), and $v\ge200\log M$, because $0<\varepsilon\le\frac{1}{200}$.
(M4) says $M^{\alpha_{0}}\ge1+\exp(\hat{\kappa})>\exp(\hat{\kappa})$; with $\alpha_{0}=\frac{1}{54}$ (row 48) this gives $\log M>54\hat{\kappa}$, and $\hat{\kappa}=N_c\cdot\kappa+\log(6c_{X})\ge(2L+1)^3\cdot\kappa$ (rows 41, 55; $6c_{X}=324\ge1$).
Moreover $\kappa\ge7$ (row 32).
Therefore $v\ge10800\,(2L+1)^3\cdot\kappa\ge\max(1,10L)$, and $v^2\ge v\ge10L$.

\paragraph{(f9) The assertions of (f) that begin with $K_{0}\ge1$.}
They are proved together with three relations which Annex~B states beside them: the closure relation, $c_{\tau}\le\mu_-/48$, and $C_{0}\ge4b_{+}+16C^{F}\cdot c(\gamma^{F})$.
\begin{itemize}
\item $K_{0}\ge1$, $c_{\mathrm{fl}}\ge1$, and $M$ is odd: Lemma 1, rows 6, 8, 78.
\item The closure relation (I.1.1) (row 2; for $\hat n_i$ see Section~I.1.2):
\[
\sum_{i=0}^{3}w_i\cdot\hat n_i=6e_3+6\sqrt2\cdot\frac{e_2-e_3}{\sqrt2}+6\sqrt2\cdot\frac{e_1-e_2}{\sqrt2}-6e_1=0 .
\]
\item $12(1+\rho)\le24$, because $\rho<1$.
\item $L-1<M$, because $M\ge L^2$ (M1).
\item $b_{k}=b\cdot\left(1-L^{-1}\right)/\left(1-L^{-k}\right)\ge b_{-}$ for $k\ge1$, because $0<1-L^{-k}\le1$ (Section~I.3.4).
\item $C^{F}\ge1$, $C_{w}\ge1$, $0<\gamma^{F}\le1$: (f5) and Lemma 1, row 73.
\item $\kappa-1\ge\kappa_{0}$ and $\kappa-\kappa_{0}\ge2$, because $\kappa\ge7\kappa_{0}+7$.
\item $e_0\le1$: (f7).
\item $6\varepsilon\le\frac{6}{200}<1$.
\item $0<\sigma<\mathfrak{e}\le\frac{1}{4}$: Lemma 1, rows 115, 116.
\item $\mathrm{vt}\le\min\left(\gamma^{F}/8,1/M\right)$, $c_{\tau}\le\mu_-/48$ and $C_{0}\ge4b_{+}+16C^{F}\cdot c(\gamma^{F})$ hold by the definitions (rows 103, 108, 92).
\item $h\ge1$: 6.1.
\item Let $e\cdot h^{\frac{1}{2}}\le e_1$.
Then $e\le e_1\le\frac{1}{3}$, since $h\ge1$ and by (f7).
Further $v_K\ge8$ and $\left(1+3\,\mathrm{lg}/v_K\right)^r<L$, by (5.1): these are the first and the third inequality of (e1) at $v=v_K\ge-\log e_1$.
Finally, by (5.1) and the third inequality of (e3),
\[
3\left(e\cdot h^{\frac{1}{2}}\right)^{\frac{1}{4}-9\varepsilon}=3\exp\left(-\left(\frac{1}{4}-9\varepsilon\right)v_K\right)\le\frac{3}{8}\le1 .
\]
\item Let $e\cdot h^{\frac{1}{2}}\le e_1$, $R=R^{*}(e)$ and $N\ge N_0(e)$.
Then $\mu_-<\mu_*$, and hence $c_{\tau}\le\mu_-/48<\mu_*/48\le\mu_*/24$ (Section~G5.9.3).
Indeed, by row 120, $h=L^j$ with the integer $j=\lfloor\log_L(1/\bar{m})\rfloor>\log_L(1/\bar{m})-1$, so that $\bar{m}\cdot h>1/L$.
Moreover $a_{*}\ge h/L^2$ by 6.3.
Hence $\mu_*=\bar{m}\cdot a_{*}\ge\bar{m}\cdot h/L^2>L^{-3}=\mu_-$ (Section~I.6.1; row 58).
\end{itemize}

\paragraph{(f10) The inequalities between the constants of the line (C6') stated in Annex~B, Section~G5.9.3.}
Put $\mathrm{ssF}:=\mathrm{ssum}(\gamma^{F}/2)=\left(24/\gamma^{F}\right)^3$.
Since $\gamma^{F}\le1$, one has $\mathrm{ssF}\ge1$; further $C^{F}\ge4$, $M\ge25$ and $K_{0}\ge1$ ((f5); Lemma 1, rows 78, 6).
Hence $\mathbf{S}=C^{F}\cdot\mathrm{ssF}\ge4$, $C_H=203M\cdot(C^{F})^2\ge2$ and $h_{G}=1+\mathrm{gGa}\ge1$; all constants of the line are positive (Lemma 1).
\begin{itemize}
\item $C_K=226\mathbf{S}^2\le1328\mathbf{S}^2=C_K'$.
\item $C_K=226(C^{F})^2\cdot\mathrm{ssF}^2\le203M\cdot(C^{F})^2\cdot\mathrm{ssF}^2=C_H\cdot\mathrm{ssF}^2$, because $203M\ge226$.
\item $203^2=41209$; using $M\ge1$ and $C^{F}\ge1$,
\begin{gather*}
C_H^2=41209M^2\cdot(C^{F})^4\ge41209M\cdot(C^{F})^3\ge31095M\cdot(C^{F})^3,\\
C_H^2\ge41209(C^{F})^3\ge32803(C^{F})^3\ge21645(C^{F})^3 .
\end{gather*}
\item $149(C^{F})^2\le160(C^{F})^2\le203M\cdot(C^{F})^2=C_H$, because $M\ge1$.
\item $13.5\mathbf{S}\le14\mathbf{S}^2$, because $\mathbf{S}\ge1$; hence $13.5\mathbf{S}+1314\mathbf{S}^2\le1328\mathbf{S}^2=C_K'$.
\item $c_{\tau}\le1/(292\mathbf{S})$; hence $146c_{\tau}\cdot\mathbf{S}\le\frac{1}{2}$.
\item By $C_K\le C_H\cdot\mathrm{ssF}^2$, $1\le C_H\cdot\mathrm{ssF}^2$, $\mathrm{ssF}\le\mathrm{ssF}^2$ and $K_{0}\ge1$,
\begin{align*}
8C_K+3K_{0}+1200C_H\cdot\mathrm{ssF}\cdot K_{0}&\le(8+3+1200)\,C_H\cdot\mathrm{ssF}^2\cdot K_{0},\\
4C_K+82K_{0}+1200C_H\cdot\mathrm{ssF}^2\cdot K_{0}&\le(4+82+1200)\,C_H\cdot\mathrm{ssF}^2\cdot K_{0};
\end{align*}
both right sides are at most $1286C_H\cdot\mathrm{ssF}^2\cdot K_{0}\le10^3C_H^2\cdot\mathrm{ssF}^4\cdot K_{0}=\mathrm{CVj}$, because $C_H\cdot\mathrm{ssF}^2\ge2$.
\item $16h_{G}^2+256\le272h_{G}^2\le10^3h_{G}^2$, because $h_{G}\ge1$; hence $\left(16h_{G}^2+256\right)C_H^2\cdot\mathrm{ssF}^4\cdot K_{0}\le\mathrm{CVj}\cdot h_{G}^2$.
\item $2\le3$ and $2\mathrm{e}^4\le4\mathrm{e}^6$; hence, term by term,
$2\mathrm{CVj}\cdot h_{G}^2+2\mathrm{alG}\cdot\mathrm{CVj}^2\cdot h_{G}^4+2\mathrm{e}^4\cdot C_{d}\cdot M^3\cdot K_{0}\le\mathrm{CWj}$.
\item $\mathrm{gGa}^2=b_{-}^{-1}+\mu_-^{-1}\le b_{-}^{-1}+4/\mu_-=\mathrm{gGa}_{1}$, and $\mathrm{gGa}_{1}\ge b_{-}^{-1}\ge1$, because $0<b_{-}<b=1$ (Lemma 1, rows 9, 35).
\item $\mathrm{hL}=6\mathbf{S}=6C^{F}\cdot\mathrm{ssF}\ge6C^{F}$.
\item $\mathrm{vt}\le\theta_{\mathrm{CT}}/2$ and $\mathrm{vt}\le\gamma^{F}/8$ (row 103); hence $\theta_{\mathrm{CT}}-\mathrm{vt}\ge\theta_{\mathrm{CT}}/2$ and $\gamma^{F}-\mathrm{vt}\ge7\gamma^{F}/8$.
\item $c_{\tau}\le\mu_-/48$ (row 108).
\end{itemize}
The remaining inequality $\mu_-/48<\mu_*/24$ of Section~G5.9.3 contains the quantity $\mu_*$, which is formed with the data; it is proved at the end of (f9).

\paragraph{(f11) The ranges.}
Let every constant be chosen as (a) and (b) prescribe.
By Lemma 1 every entry of the rows 1--116 other than 27, 28 is defined and finite.
The sets named in Lemma 1 are contained in $(0,\infty)$, except those of $A_{\mathrm{ak}}$, $A'_{\mathrm{ak}}$\,[21], $C_{G}$\,[24], $C_{6}$\,[25], $C_{R}$\,[63] and of the values of $\varepsilon_5$\,[30], which are contained in $[0,\infty)$; the constant $C_{\Sigma}$\,[94] lies in $(0,\infty)$, hence in $[0,\infty)$.
The constants $C_E$\,[27], $\mathrm{Csy}_{\mathrm{ak}}$ and $\mathrm{Cpv}_0$\,[28] occur in no row up to $e_1$\,[118] other than their own rows 27, 28, and in no condition of (b), by the end of Step 3; this proves the assertion in parentheses in the part of (f) on ranges.
The special ranges are entries of Lemma 1:
\begin{itemize}
\item $\lambda_{\mathrm{int}}$\,[10] $\le2(w_0+w_1+w_2)$;
\item $c_{*}$\,[14] $\le\frac{1}{3}$;
\item $c_{\mathrm{el}}$\,[18] $\le0.22$;
\item $K_{\mathrm{Str}}$\,[19] $\ge1$;
\item $\kappa_{(0)}$\,[20] $\le\theta_{0}/2$;
\item $a_{L}$\,[38] $>1$;
\item $a_{0}$\,[62] $\le\frac{1}{3}$;
\item $\gamma_{(5)}$\,[66], $a_{z}$\,[68], $\gamma_{z}$\,[72] $\le1/(3L)$;
\item $c_{\mathrm{re}}$\,[67] $\le b_{-}$\,[35] $<b=1$;
\item $C_{00}$\,[70] $\ge1$ (stated in the entry for the rows 71--75);
\item $\gamma_{P}$\,[50] $\le1$;
\item $C_{P}$\,[52] $\ge4\mathrm{e}^2$;
\item $\gamma_{\Sigma}$\,[93] $\le1$;
\item $c_{\gamma}$\,[87] $\le2L^3\cdot\exp\left((2L+1)\gamma_{\mathrm{w}}\right)\cdot c(\gamma_{\mathrm{w}})$;
\item $c_{\kappa}$\,[91] $\ge1$;
\item $p_0$\,[111] $\ge9$ and $p$\,[112] $\ge19$.
\end{itemize}

\paragraph{(f12) Further inequalities stated after ``Used in'' in the entries of (b).}
\begin{itemize}
\item Row (L1): $110>2\mathrm{e}^4$.
Indeed, $\frac{1}{j!}\le\frac{1}{720}\cdot7^{-(j-6)}$ for $j\ge6$; hence
\[
\mathrm{e}\le1+1+\frac{1}{2}+\frac{1}{6}+\frac{1}{24}+\frac{1}{120}+\frac{1}{720}\cdot\frac{7}{6}=\frac{163}{60}+\frac{7}{4320}<2.7184 ,
\]
and $2.7184^2<7.39$, $7.39^2<54.62$, so that $2\mathrm{e}^4<109.24$.
Further $L\ge13$ and, in the row (L5), $L\ge3$ and $L>1$: (f1).
\item Row (L5), together with (S1): $1-\varepsilon\ge\frac{1}{2}$ and $L\ge4$ give $L^{1-\varepsilon}\ge L^{\frac{1}{2}}\ge2$.
\item Row (M1): $(L-1)/2\le L^2\le M$; $\eta_k=L^{-k}\le1\le L^2\le M$ for $k\ge0$ (as $L\ge1$); $M\cdot L^k$ ($k\ge0$) is a product of odd integers ((L5); Lemma 1, row 78), hence odd; $T^{*}(e)$ is an integer $\ge1$ for $e>0$ with $e\cdot h^{\frac{1}{2}}\le e_1$, by 6.3.
\item Row (M3): $\gamma_{\mathrm{w}}\cdot M\ge16$ and $M>0$ give $\gamma_{\mathrm{w}}>0$ (this is also Lemma 1, row 75).
\item Row (S1): $0<\varepsilon\le\frac{1}{200}\le\frac{1}{8}$ (Lemma 1, row 113).
\item Row (S2): $\mathfrak{e}>0$ (Lemma 1, row 115).
\item Row (A8): $e_k\le e_0\le4^{-200}<1$, by (f7).
\item Row (e3): the third inequality at $v=v_K$ is proved in (f9).
\end{itemize}

\subsubsection*{Step 9. The paragraph after the table (b)}

For $e_k\in(0,e_0]$ put $v=v_k=-\log e_k$.
Then $e_k^{\mathrm{cf}}=\exp(-\mathrm{cf}\cdot v)$ for every real number $\mathrm{cf}$, and $p_{0,k}=v^{p_0}$, $p_k=v^p$ (Section~I.6.1).
With this substitution the conditions of the stages 1--10 of \cite[Definition~5.97]{QEDI} and the conditions of the first kind of \cite[Section~5.6.2]{QEDI}, namely the conditions (A0)--(A10), (Z2.m), (Z2.s1)--(Z2.s5), (Z3.s1) and (Z8.s1) on the constants of these stages, correspond to the clauses (a), (b) of Proposition~14.2 as listed below.
The names are those of Section~I.8.4.
(A0)--(A10) are the conditions that carry these tags in \cite[Section~5.6.2]{QEDI};
(Z2.m) is the condition $M=L^m\ge M^F$ stated there, with $M^F$ formed with the constants $C_V$, $a$ of \cite{QEDI};
(Z2.s1)--(Z2.s5) are the five conditions on $e_0$ that follow it there;
(Z3.s1) is the condition $e_0\le\exp(-1)$;
(Z8.s1) is the condition that $e_1$ be the number defined by the inequalities (e1)--(e8) of \cite[Section~19.1.2]{QEDI};
(Z9.s1)--(Z9.s3) are the three conditions on $e_*$ stated after it.
In each entry of the list the conditions of \cite{QEDI} come first; after the word \emph{Here} follow the rows and the labels of Proposition~14.2\,(a), (b) that correspond to them.

\begin{description}
\item[Stage 1.]
$\bar{\varepsilon}=\frac{1}{200}$, $b=1$;
$c_{X}$, $\kappa_{0}$, $K_{0}$, $c_{\mathrm{rb}}$, $c_{\mathrm{fl}}$ with the values fixed in \cite[Section~13.1.6]{QEDI}.
\emph{Here:} the rows 3--9.
\item[Stages 2 and 3.]
$\alpha\in(1-\bar{\varepsilon},1)$;
$\kappa\ge7\kappa_{0}+7$.
\emph{Here:} the rows 31, 32.
\item[Stage 4.]
$L$ odd;
$L^{\frac18-\frac{3\bar{\varepsilon}}{2}}\ge2400$;
$c_{\mathrm{fl}}\cdot L^{-\frac38+6\bar{\varepsilon}}\le\frac{1}{16}$;
$L\cdot(\kappa-6\kappa_{0}-4)-1\ge\kappa$;
$c_{\mathrm{rb}}\cdot L^3\cdot\exp\left(-\left(L\cdot(\kappa-\kappa_{0}-1)-\kappa\right)\right)\le\frac12$;
$L\ge L^F$ and $L\ge L^P$, where $L^F=L^P=5$.
\emph{Here:} in this order, (L5); (L1); (L2); (L3); (L4); (L5).
\item[Stage 5.]
Definitions of constants, among them the numbers $\mathfrak{m}_0$, $m_0$, $h'$.
\emph{Here:} the rows 34--58, 73, 75;
$\mathfrak{m}_0=1$, $m_0=0$ and $h'=h$ (Section~I.6.1).
\item[Stage 6.]
$M=L^m$ and $M\ge L^2$;
$M>4L$;
$\gamma_{\mathrm{w}}\cdot M\ge16$;
$M^{\alpha_{0}}\ge1+\exp(\hat{\kappa})$;
$M\ge M^{F}$, with $M^F$ taken at $\kappa_{\mathrm{w}}+\kappa_{0}$;
$M\ge M^{P}$;
$M\ge L^{m_0+2}$.
\emph{Here:} in this order, (M1); (M2); (M3); (M4); (M5); (M6); and (M1) once more, because $m_0=0$.
\item[Stage 7.]
Definitions of constants.
\emph{Here:} the rows 79--95.
\item[Stage 8.]
$r\ge2$, $p_0>3r+2$, $p>2p_0$;
$\varepsilon\in(1-\alpha,\bar{\varepsilon}]$;
$s\in(0,1)$ with $\sigma>0$.
\emph{Here:} the rows 110--112;
(S1);
row 114 ($s\in(0,1)$) together with (S2).
\item[Stage 9.]
$e_0>0$, together with all conditions of \cite[Section~5.6.2]{QEDI} in which $e_0$ occurs: (A1)--(A10), (Z2.s1)--(Z2.s5) with the constant $C_{w}$ of $\mathrm{par}$, and (Z3.s1).
\emph{Here:} row 117 ($e_0>0$), and the entries of these labels below;
the constant $C_{d}$ of (Z2.s5) is that of row 90, which is formed with $C_{w}$ [86], the constant of $\mathrm{par}$.
\item[(A0)]
$L\cdot(\kappa-6\kappa_{0}-4)-1\ge\kappa$;
$\kappa_{\mathrm{w}}=\kappa+\kappa_{0}+3$;
$M\ge L^2$, $M>4L$, $\gamma_{\mathrm{w}}\cdot M\ge1$;
$M^{\alpha_{0}}\ge1+\exp(\hat{\kappa})$.
\emph{Here:} (L3); row 53; (M1), (M2), and (M3) together with (f3); (M4).
\item[(A1)]
$p_{0,k}\le e_k^{-\frac{\varepsilon}{4}}$;
$12\exp\left(-\frac{p_{0,k}^2}{2}\right)\le e_k^7$;
$e_k^{\frac{5\varepsilon}{2}}\le\frac12$;
$c_{\mathrm{fl}}\cdot e_k^{\varepsilon}\le1$;
$3M\cdot e_k^{\varepsilon}\le1$.
\emph{Here:} the five inequalities of (A1), in this order (note $p_{0,k}^2=v^{2p_0}$).
\item[(A2)]
$2\,C_{w}^2\cdot c_{\gamma}^2\cdot c_{C}\cdot e_k^{\frac{\varepsilon}{4}}\le1$.
\emph{Here:} (A2).
\item[(A3)]
$16\,c_{X}\cdot c_{\gamma}\cdot C_{w}\cdot C_{T}^2\cdot M^3\cdot e_k^{\frac{\varepsilon}{2}}\le1$;
$96\,M^3\cdot e_k^{\frac14}\le1$.
\emph{Here:} (A3), in this order.
\item[(A4)]
$32\,C_{w}\cdot C_{T}\cdot e_k^{\frac{3\varepsilon}{2}}\le1$.
\emph{Here:} (A4).
\item[(A5)]
$88\,M^3\cdot e_k^{\frac14-6\varepsilon}\le1$.
\emph{Here:} (A5).
\item[(A7)]
$c_{S}\cdot e_k^{\frac14-5\varepsilon}\le1$;
$c_{S}\cdot L^3\cdot e_k^{\varepsilon}\le1$;
$36\,c_{X}^2\cdot K_{0}^2\cdot e_k^{\frac14-5\varepsilon}\le1$.
\emph{Here:} (A7), in this order.
\item[(A6)]
$16\,C_{w}^2\cdot c_{\gamma}\cdot e_k^{2\varepsilon}\le3$.
\emph{Here:} (A6), with $112$ in place of $16$ (Section~I.8.4, (I.L0.19)).
\item[(A8)]
$e_0^{\varepsilon}\le\frac14$;
$3M\cdot\left(e_0^{\frac14+\varepsilon}+e_0^{2\varepsilon}\right)\le1$;
$3M\cdot e_0^{\varepsilon}\le1$;
$e_0^{-\frac34+\varepsilon}+e_0^{-1+2\varepsilon}\le e_0^{-1+\varepsilon}$.
\emph{Here:} (A8), in this order;
the fourth inequality after multiplication by $e_0^{1-\varepsilon}$.
\item[(A9)]
$160\,L^{1+\frac{3\varepsilon}{2}}\cdot M\cdot e_0^{\varepsilon}\le1$;
$90\,c_{X}\cdot L\cdot M\cdot e_0^{\varepsilon}\le1$;
$9M\cdot L^{\varepsilon}\cdot e_0^{2\varepsilon}\le1$;
$\left(-\log e_0\right)^2\ge10L$.
\emph{Here:} (A9), in this order.
\item[(A10)]
$e_0\le e^{F}_0$;
for every $e_k\in(0,e_0]$, with $v_k=-\log e_k$ and $p_k=v_k^p$,
$v_k\ge9$,
$8M\cdot\left(\frac{v_k}{4}+4\right)\cdot e_k^{\varepsilon}\le1$,
$8\,C_{B}\cdot p_k\le e_k^{-\frac34+3\varepsilon}$,
$8\,e_k^{\frac34-\frac{7\varepsilon}{2}}\le1$;
$8\,c_{\mathrm{rb}}\cdot L^3\cdot C_{d}\cdot M^3\cdot e_0^{2\varepsilon}\le1$.
\emph{Here:} the four inequalities imposed for every $e_k$ are the four inequalities of (A10), in this order (with $v=-\log e_k$ and $p_k=v^p$);
the last inequality is (Z2.s5), with the constant $C_{d}$ of row 90;
the inequality $e_0\le e^{F}_0$ has no row, and in its place the conditions (Z2.s1)--(Z2.s5) are imposed on $e_0$ itself.
\item[(Z2.m)]
$M=L^m\ge\max\left\{4L+1,\ \left(108\pi\cdot C_V\cdot a\right)^2,\ 108\pi\cdot C_V\cdot a\cdot\exp(27\kappa')\right\}$.
\emph{Here:} (M1), and (M5) with the row 76, which has $C'_{V}$, $a'$ in place of $C_V$, $a$.
\item[(Z2.s1), (Z2.s3), (Z2.s4)]
(Z2.s1): $6L\cdot e_0^{\varepsilon}\le1$;
(Z2.s3): $3M\cdot e_0^{\varepsilon}\le1$;
(Z2.s4): $e_0^{\varepsilon}\le\frac14$, $e_0^{\frac14}+e_0^{\varepsilon}\le1$, $4e_0^{\frac14+2\varepsilon}+2e_0^{3\varepsilon}\le1$.
\emph{Here:} the conditions with the same labels, in this order.
\item[(Z2.s2)]
$4\,C_V\cdot a\cdot e_0^{\varepsilon}\le1$.
\emph{Here:} (Z2.s2), with $C'_{V}$, $a'$ in place of $C_V$, $a$.
\item[(Z2.s5)]
$8\,c_{\mathrm{rb}}\cdot L^3\cdot C_{d}\cdot M^3\cdot e_0^{2\varepsilon}\le1$.
\emph{Here:} (Z2.s5), with the constant $C_{d}$ of row 90.
\item[(Z3.s1)]
$e_0\le\exp(-1)$.
\emph{Here:} (Z3.s1), that is, $v\ge1$.
\item[Stage 10 and (Z8.s1).]
$e_1$ is one half of the supremum of the numbers $e'\in\left(0,\min\left(e_0,\frac13\,L^{-\frac{m_0}{2}}\right)\right]$ such that (e1)--(e8) of \cite[Section~19.1.2]{QEDI} hold for all $v\ge-\log e'$.
There (e1), (e2), (e4), (e5), (e7), (e8) and the second and third inequalities of (e3) are stated as in Proposition~14.2\,(b), while the first inequality of (e3) and the inequality (e6) read
\begin{gather*}
8\,C_{0}^2\cdot\mathrm{lh}(v)^6\cdot\left(b_{-}^{-1}+L^{m_0+1}\cdot\mathrm{lh}(v)\right)\le\exp\left(\frac{\varepsilon v}{2}\right),\\
C_{\Sigma}\cdot\exp\left(-\gamma_1\cdot M\cdot v^r\right)\le\exp(-v).
\end{gather*}
\emph{Here:} row 118 and (e1)--(e8);
since $m_0=0$, the first inequality of (e3) contains $L\cdot\mathrm{lh}(v)$, and the upper end of the range of $e'$ is $\min\left(e_0,\frac13\right)$;
in (e6) the rate $\gamma_{\Sigma}$ stands in place of $\gamma_1$.
\end{description}

Every condition of \cite{QEDI} in this list has its counterpart after the word \emph{Here}, and every condition of Proposition~14.2\,(b) occurs among these counterparts.
In \cite[Section~5.6.2]{QEDI} the conditions (A1)--(A7) are imposed for every $e_k\in(0,e_0]$, whereas (A8), (A9), (Z2.s1)--(Z2.s5) and (Z3.s1) are imposed at $e_0$ only;
there (A10) imposes two inequalities at $e_0$ only, namely $e_0\le e^{F}_0$ and $8\,c_{\mathrm{rb}}\cdot L^3\cdot C_{d}\cdot M^3\cdot e_0^{2\varepsilon}\le1$, the latter being (Z2.s5), and its four other inequalities for every $e_k\in(0,e_0]$, and it is a condition on the constants alone.
In Proposition~14.2\,(b), (A10) consists of exactly these four inequalities, written with $v$ and imposed for every $v\ge-\log e_0$, and it is likewise a condition on the constants alone;
for these four inequalities the two forms coincide, as for (A1)--(A7), since $v=-\log e_k$ ranges over $[-\log e_0,\infty)$ when $e_k$ ranges over $(0,e_0]$;
the inequality (Z2.s5) is one of the conditions of the second group below;
the number $e^{F}_0$ does not occur, and all conditions are imposed for every $v\ge-\log e_0$.
For (A1)--(A7) the two forms coincide, since $v=-\log e_k$ ranges over $[-\log e_0,\infty)$ when $e_k$ ranges over $(0,e_0]$.
For the second group, (A8), (A9), (Z2.s1)--(Z2.s5), (Z3.s1), the two forms are equivalent.
Indeed, each inequality of this group is of one of two kinds.
Either it states that a sum of terms $c\cdot\exp(-\mathrm{cf}\cdot v)$, with constants $c\ge0$ and $\mathrm{cf}>0$, is at most a given number; then its left side does not increase with $v$.
Or it is a lower bound on $v$: this is the case for (Z3.s1), and for the fourth inequality of (A9), where $-\log e_0>0$ by the first inequality of (A8), which in both forms says $e_0^{\varepsilon}\le\frac14<1$, and where $v\mapsto v^2$ is increasing on $(0,\infty)$.
In both cases the inequality holds for all $v\ge-\log e_0$ as soon as it holds at $v=-\log e_0$.
The differences between the conditions of \cite{QEDI} and their counterparts are the changes listed in the Remark after Proposition~14.2.

\emph{Labels of Section~I.8.4 that have no row.}
The entries (A0), (Z2.m) and ``Stage 10 and (Z8.s1)'' of the list show the following.
(A0) consists of (L3), the definition of $\kappa_{\mathrm{w}}$ (row 53), (M1), (M2), the inequality $\gamma_{\mathrm{w}}\cdot M\ge1$, which follows from (M3) by (f3), and (M4).
The inequality $M\ge M^{F}$ of (Z2.m) is (M5), and its clause $M=L^m$ is contained in (M1).
(Z8.s1) is the definition of $e_1$ (row 118).
Further, \cite[Section~5.6.2]{QEDI} states the three conditions
\begin{gather*}
\text{(Z9.s1)}\quad\left(4+\frac{M^3}{2}\right)\cdot e_{*}^{\varepsilon}\le1,\\
\text{(Z9.s2)}\quad4\,C_{0}\cdot\ell_*^3\cdot g_1\le e_{*}^{-\frac14+\varepsilon},\\
\text{(Z9.s3)}\quad48\,e_{*}^{2\varepsilon}\le\mu_-;
\end{gather*}
with $\mathrm{gGa}_{1}$ in place of $g_1$ (Section~I.8.2, line (C6')) they are the second inequality of (g1), the first inequality of (g3) and the first inequality of (g4) of Proposition~14.2\,(d).
The label (G4.2.1) of Section~I.8.4 denotes the condition $r_{*}>31L+16$, which is the condition of Proposition~14.2\,(c) that carries the tag (2.1).
The condition (Z10.s3), the third of the four conditions on $e_2$ in \cite[Section~5.6.2]{QEDI}, namely the one that contains $C_Y$, is not imposed (Section~I.8.4), and it has no counterpart in Proposition~14.2.
Every other label named in Section~I.8.4 is a row of Proposition~14.2\,(b) or a label of Proposition~14.2\,(c), (d).

This proves the first sentence of the paragraph after the table (b) and its sentences on the labels of Section~I.8.4, and it completes the proof of Proposition~14.2.
The second sentence of that paragraph is the rule (G1.T-12) for the proofs of all Propositions: a further condition on the constants that the proof of a Proposition requires is of the type ``$e_0$ (respectively $e_1$, $e_{\sharp}$, $e_2$) small depending on earlier constants'', and it belongs to the list of Proposition~14.2.
That the proof of a Proposition satisfies this rule is shown in the section ``Conditions on the constants'' of its chapter, which names, for every inequality between the constants used in that proof, the row of Proposition~14.2\,(a), (b), (d) from which it comes, or the line of Proposition~14.2\,(c), (f) under which it stands (Section~IV.1, paragraph~E; Annex~B).
No step of the proofs of Propositions~14.2 and 14.3 rests on this rule.
\qed

\subsection*{2. Proof of Proposition 14.3}

The elementary inequalities (El1)--(El7), Lemma 1 and the tables of Steps 3 and 4 are those of Section 1.
A choice of the lines (C0)--(C$j$) is called \emph{admissible} if every constant of these lines is chosen as Proposition~14.2\,(a) prescribes, subject to the conditions of Proposition~14.2\,(b) that restrict it.
Lemma 1 applies to a given row as soon as this row and all earlier rows are chosen in this way.

\subsubsection*{(a) The lines (C0)--(C2)}

The interval $(0,1)$ contains $\frac12$.
The interval $(1-\bar{\varepsilon},1)=\left(\frac{199}{200},1\right)$ contains $\frac{399}{400}$.
The interval $[7\kappa_{0}+7,\infty)$ contains $7\kappa_{0}+7$, which is a finite number, since $\kappa_{0}=54\log72$.
The other entries of the lines (C0) and (C0r) are defined by Lemma 1 (the rows 2--30 other than 27, 28).

\subsubsection*{(b) The line (C3)}

Let $\kappa\ge7\kappa_{0}+7$.
Then
\[
\kappa-6\kappa_{0}-4\ge\kappa_{0}+3>0,\qquad\mathrm{dkap}:=\kappa-\kappa_{0}-1\ge6\kappa_{0}+6>0,
\]
and
\[
\frac18-\frac{3\bar{\varepsilon}}{2}=\frac{47}{400}=0.1175>0,\qquad\frac38-6\bar{\varepsilon}=\frac{69}{200}=0.345>0.
\]
Put
\begin{multline*}
L_{\min}:=\max\left\{5,\ 2400^{\left(\frac18-\frac{3\bar{\varepsilon}}{2}\right)^{-1}},\ \left(16\,c_{\mathrm{fl}}\right)^{\left(\frac38-6\bar{\varepsilon}\right)^{-1}},\ \frac{\kappa+1}{\kappa-6\kappa_{0}-4},\right.\\
\left.\frac{2}{\mathrm{dkap}}\cdot\log^{+}\left(2\,c_{\mathrm{rb}}\cdot\exp(\kappa)\cdot\left(\frac{6}{\mathrm{dkap}\cdot\mathrm{e}}\right)^3\right)\right\}.
\end{multline*}
The number $L_{\min}$ is finite, and it is a function of $(\bar{\varepsilon},c_{\mathrm{fl}},c_{\mathrm{rb}},\kappa,\kappa_{0})$.
Let $L$ be an odd integer with $L\ge L_{\min}$.
\begin{itemize}
\item[(L5)] $L$ is odd and $L\ge5$.
\item[(L1)] The map $x\mapsto x^{\frac{47}{400}}$ is increasing, and $L\ge2400^{\frac{400}{47}}$ gives $L^{\frac{47}{400}}\ge2400$.
\item[(L2)] $L\ge\left(16\,c_{\mathrm{fl}}\right)^{\frac{200}{69}}$ gives $L^{\frac{69}{200}}\ge16\,c_{\mathrm{fl}}$, that is, $c_{\mathrm{fl}}\cdot L^{-\frac38+6\bar{\varepsilon}}\le\frac{1}{16}$.
\item[(L3)] $L\cdot(\kappa-6\kappa_{0}-4)\ge\kappa+1$, because the factor $\kappa-6\kappa_{0}-4$ is positive and $L\ge\frac{\kappa+1}{\kappa-6\kappa_{0}-4}$.
\item[(L4)] The left side of (L4) equals $c_{\mathrm{rb}}\cdot\exp(\kappa)\cdot L^3\cdot\exp(-\mathrm{dkap}\cdot L)$.
By (El7) with $\mathrm{nt}=1$, $(\mathrm{Cf}_1;\mathrm{dg}_1;\mathrm{rt}_1)=\left(c_{\mathrm{rb}}\cdot\exp(\kappa);3;\mathrm{dkap}\right)$, $\mathrm{tg}=\frac12$ and $v=L$, it is at most $\frac12$ for all
\[
L\ge\frac{2}{\mathrm{dkap}}\cdot\log^{+}\left(2\,c_{\mathrm{rb}}\cdot\exp(\kappa)\cdot\left(\frac{6}{\mathrm{dkap}\cdot\mathrm{e}}\right)^3\right).
\]
\end{itemize}
Odd integers $\ge L_{\min}$ exist.
Hence the set of the admissible $L$ is not empty, and it contains every odd $L\ge L_{\min}$.

\subsubsection*{(c) The line (C5)}

Let the lines (C0)--(C4r) be chosen admissibly.
Lemma 1 gives the following.
First, $\gamma_{\mathrm{w}}\in\left(0,\frac{1}{6L}\right]$ (row 75), so that $\frac{16}{\gamma_{\mathrm{w}}}<\infty$.
Next, $\hat{\kappa}\in(0,\infty)$ (row 55) and $\alpha_{0}=\frac{1}{54}$ (row 48), so that $\left(1+\exp(\hat{\kappa})\right)^{\frac{1}{\alpha_{0}}}=\left(1+\exp(\hat{\kappa})\right)^{54}<\infty$.
Further, $M^{F}<\infty$ (row 76); this rests on $a'<\infty$ (row 74), that is, on $C_{5}<\infty$ (row 65), $c_{\mathrm{re}}>0$ (row 67), $\gamma^{F}>0$ (row 73) and (El5).
Finally, $M^{P}$ is defined and finite (row 77); this rests on $\gamma_{P}>0$, which holds by the formula for $\gamma_{P}$ in \cite[Section~15.7]{QEDI}.
Hence the number
\[
M_{\min}=\max\left\{L^2,\ 4L+1,\ \frac{16}{\gamma_{\mathrm{w}}},\ \left(1+\exp(\hat{\kappa})\right)^{54},\ M^{F},\ M^{P}\right\}
\]
of Proposition~14.3\,(c) is finite.
Let $M$ be a power of $L$ with $M\ge M_{\min}$.
\begin{itemize}
\item[(M1)] $M$ is a power of $L$, and $M\ge L^2$.
\item[(M2)] $M\ge4L+1>4L$.
\item[(M3)] $M\ge\frac{16}{\gamma_{\mathrm{w}}}$ and $\gamma_{\mathrm{w}}>0$ give $\gamma_{\mathrm{w}}\cdot M\ge16$.
\item[(M4)] The map $x\mapsto x^{\frac{1}{54}}$ is increasing, and $M\ge\left(1+\exp(\hat{\kappa})\right)^{54}$ gives $M^{\alpha_{0}}\ge1+\exp(\hat{\kappa})$.
\item[(M5)] $M\ge M^{F}$.
\item[(M6)] $M\ge M^{P}$.
\end{itemize}
Apart from the clause ``$M$ is a power of $L$'', each of (M1)--(M6) has the form $M\ge$ (a number formed with constants of the rows $<78$), or, in the case of (M2), $M>$ (such a number).
So all conditions on $M$ are lower bounds.
Since $L\ge5$, the powers of $L$ are unbounded, and there are powers of $L$ that are $\ge M_{\min}$.

\subsubsection*{(d) The line (C7)}

Let $\alpha\in(1-\bar{\varepsilon},1)$ and an admissible $L$ be given.
The triple $(r,p_0,p)=(2,9,19)$ satisfies $r\ge2$, $p_0>3r+2=8$ and $p>2p_0=18$.
Since $1-\alpha<\bar{\varepsilon}$, the interval $(1-\alpha,\bar{\varepsilon}]$ is not empty; it contains $\bar{\varepsilon}$.
Let $\varepsilon\in(0,\bar{\varepsilon}]$ and $s=\frac12$.
Then
\begin{gather*}
\frac12\cdot\left(\frac14-9\varepsilon\right)-12\varepsilon=\frac18-\frac{33}{2}\,\varepsilon\ge\frac18-\frac{33}{400}=\frac{17}{400}=0.0425,\\
\frac{2a_{L}}{N_{*}}=\frac{2\left(L^3-1\right)}{(L-1)^3\cdot8L^3}<\frac{1}{4\,(L-1)^3}\le\frac{1}{256}<\frac{17}{400}\qquad(L\ge5).
\end{gather*}
Hence $\sigma>0$ at $s=\frac12$: the condition (S2) holds, and $\frac12$ lies in $(0,1)$.

\subsubsection*{(e) The line (C8)}

Let the lines (C0)--(C7) be chosen admissibly.
By Lemma 1 the constants
\[
c_{X},\ K_{0},\ c_{\mathrm{rb}},\ c_{\mathrm{fl}},\ L,\ c_{C},\ C'_{V},\ C_{T},\ a',\ M,\ C_{w},\ c_{\gamma},\ C_{B},\ C_{d},\ c_{S},
\]
which are the letters of row 117 other than $p_0$, $p$, $\varepsilon$, lie in $(0,\infty)$;
$p_0$ and $p$ are positive integers;
and $0<\varepsilon\le\frac{1}{200}$.
Hence the exponents
\begin{gather*}
\frac{\varepsilon}{4},\ \frac{\varepsilon}{2},\ \varepsilon,\ \frac{3\varepsilon}{2},\ 2\varepsilon,\ \frac{5\varepsilon}{2},\ 3\varepsilon,\ \frac14,\ \frac14+\varepsilon,\ \frac14+2\varepsilon;\\
\frac14-5\varepsilon\ge\frac{9}{40};\qquad\frac14-6\varepsilon\ge\frac{11}{50};\qquad\frac34-\frac{7\varepsilon}{2}\ge\frac{293}{400};\qquad\frac34-3\varepsilon\ge\frac{147}{200}
\end{gather*}
are positive.
In the following table a term $(\mathrm{Cf};\mathrm{dg};\mathrm{rt})$ stands for $\mathrm{Cf}\cdot v^{\mathrm{dg}}\cdot\exp(-\mathrm{rt}\cdot v)$.
After division by its right side, where this right side is an exponential, each inequality named in the first column states that the sum of the terms in the second column is at most the number $\mathrm{tg}$ in the third column.
\[
\begin{array}{lll}
\text{inequality} & \text{terms }(\mathrm{Cf};\mathrm{dg};\mathrm{rt}) & \mathrm{tg}\\
\hline
\text{(A1), first} & \left(1;\,p_0;\,\frac{\varepsilon}{4}\right) & 1\\
\text{(A1), third} & \left(1;\,0;\,\frac{5\varepsilon}{2}\right) & \frac12\\
\text{(A1), fourth} & \left(c_{\mathrm{fl}};\,0;\,\varepsilon\right) & 1\\
\begin{array}[t]{@{}l@{}}\text{(A1), fifth $=$ (A8), third}\\ \text{$=$ (Z2.s3)}\end{array} & \left(3M;\,0;\,\varepsilon\right) & 1\\
\text{(A2)} & \left(2\,C_{w}^2\cdot c_{\gamma}^2\cdot c_{C};\,0;\,\frac{\varepsilon}{4}\right) & 1\\
\text{(A3), first} & \left(16\,c_{X}\cdot c_{\gamma}\cdot C_{w}\cdot C_{T}^2\cdot M^3;\,0;\,\frac{\varepsilon}{2}\right) & 1\\
\text{(A3), second} & \left(96\,M^3;\,0;\,\frac14\right) & 1\\
\text{(A4)} & \left(32\,C_{w}\cdot C_{T};\,0;\,\frac{3\varepsilon}{2}\right) & 1\\
\text{(A5)} & \left(88\,M^3;\,0;\,\frac14-6\varepsilon\right) & 1\\
\text{(A6)} & \left(112\,C_{w}^2\cdot c_{\gamma};\,0;\,2\varepsilon\right) & 3\\
\text{(A7), first} & \left(c_{S};\,0;\,\frac14-5\varepsilon\right) & 1\\
\text{(A7), second} & \left(c_{S}\cdot L^3;\,0;\,\varepsilon\right) & 1\\
\text{(A7), third} & \left(36\,c_{X}^2\cdot K_{0}^2;\,0;\,\frac14-5\varepsilon\right) & 1\\
\begin{array}[t]{@{}l@{}}\text{(A8), first}\\ \text{$=$ (Z2.s4), first}\end{array} & \left(1;\,0;\,\varepsilon\right) & \frac14\\
\text{(A8), second} & \left(3M;\,0;\,\frac14+\varepsilon\right),\ \left(3M;\,0;\,2\varepsilon\right) & 1\\
\begin{array}[t]{@{}l@{}}\text{(A8), fourth}\\ \text{$=$ (Z2.s4), second}\end{array} & \left(1;\,0;\,\frac14\right),\ \left(1;\,0;\,\varepsilon\right) & 1\\
\text{(A9), first} & \left(160\,L^{1+\frac{3\varepsilon}{2}}\cdot M;\,0;\,\varepsilon\right) & 1\\
\text{(A9), second} & \left(90\,c_{X}\cdot L\cdot M;\,0;\,\varepsilon\right) & 1\\
\text{(A9), third} & \left(9M\cdot L^{\varepsilon};\,0;\,2\varepsilon\right) & 1\\
\text{(A10), second} & \left(2M;\,1;\,\varepsilon\right),\ \left(32M;\,0;\,\varepsilon\right) & 1\\
\text{(A10), third} & \left(8\,C_{B};\,p;\,\frac34-3\varepsilon\right) & 1\\
\text{(A10), fourth} & \left(8;\,0;\,\frac34-\frac{7\varepsilon}{2}\right) & 1\\
\text{(Z2.s1)} & \left(6L;\,0;\,\varepsilon\right) & 1\\
\text{(Z2.s2)} & \left(4\,C'_{V}\cdot a';\,0;\,\varepsilon\right) & 1\\
\text{(Z2.s4), third} & \left(4;\,0;\,\frac14+2\varepsilon\right),\ \left(2;\,0;\,3\varepsilon\right) & 1\\
\text{(Z2.s5)} & \left(8\,c_{\mathrm{rb}}\cdot L^3\cdot C_{d}\cdot M^3;\,0;\,2\varepsilon\right) & 1
\end{array}
\]
In the row of the second inequality of (A10) the identity $8M\cdot\left(\frac{v}{4}+4\right)=2M\cdot v+32M$ is used.
By (El7) each of these inequalities holds for all $v$ with $v\ge\mathrm{Vbd}(\ldots;\mathrm{tg})$, a finite number attached to its row.
The remaining inequalities of the line (C8) are the following.
\begin{itemize}
\item[] (A1), second: it states $\frac{v^{2p_0}}{2}-7v\ge\log12$.
For $v\ge15$ one has $v^{2p_0}\ge v^2$ (as $v\ge1$ and $2p_0\ge2$), $\frac{v^2}{2}-7v=v\cdot\left(\frac{v}{2}-7\right)\ge\frac{15}{2}$, and $\log12\le4\le\frac{15}{2}$, because $12\le2^4\le\mathrm{e}^4$ by (El1).
So it holds for all $v\ge15$.
\item[] (A9), fourth: $v^2\ge10L$ for all $v\ge(10L)^{\frac12}$.
\item[] (A10), first: $v\ge9$.
\item[] (Z3.s1): $v\ge1$.
\end{itemize}
Let $\mathrm{Vbd}_0$ be the largest of the 26 numbers $\mathrm{Vbd}$ of the table and of the numbers $15$ and $(10L)^{\frac12}$.
Then $\mathrm{Vbd}_0<\infty$, and, since $15>9>1$, all the inequalities (A1)--(A10), (Z2.s1)--(Z2.s5), (Z3.s1) hold for all $v\ge\mathrm{Vbd}_0$.
Let $\mathrm{Sadm}_0$ be the set of the numbers $e_0>0$ such that all these inequalities hold for all $v\ge-\log e_0$.
\begin{itemize}
\item[(1)] The number $\exp(-\mathrm{Vbd}_0)$ lies in $\mathrm{Sadm}_0$.
\item[(2)] If $e_0\in\mathrm{Sadm}_0$ and $0<e_0'<e_0$, then $-\log e_0'>-\log e_0$, so that $e_0'\in\mathrm{Sadm}_0$.
\item[(3)] $\mathrm{Sadm}_0\subset\left(0,4^{-\frac{1}{\varepsilon}}\right]$, because the first inequality of (A8) at $v=-\log e_0$ states $e_0^{\varepsilon}\le\frac14$.
\end{itemize}
Put $e_{\max}:=\sup\mathrm{Sadm}_0$.
By (1)--(3),
\[
\exp(-\mathrm{Vbd}_0)\le e_{\max}\le4^{-\frac{1}{\varepsilon}}<1\qquad\text{and}\qquad(0,e_{\max})\subset\mathrm{Sadm}_0\subset(0,e_{\max}].
\]
So $\mathrm{Sadm}_0$ is an interval with the left end $0$ and with the positive length $e_{\max}$.

\subsubsection*{(f) The line (C9)}

Let the lines (C0)--(C8) be chosen admissibly.
Lemma 1 gives:
$c_{\mathrm{fl}}$, $L$, $b_{-}$, $N_{*}$, $M$, $c_{\kappa}$, $C_{0}$ lie in $(0,\infty)$;
$\gamma_{\Sigma}\in(0,1]$ and $C_{\Sigma}\in[0,\infty)$ (rows 93, 94, that is, Proposition~30.3\,(e) with the values of its constants);
$r\ge2$, $p_0\ge9$, $p\ge19$;
$0<\varepsilon\le\frac{1}{200}$, so that $\frac14-9\varepsilon\ge\frac{41}{200}>0$;
$L>1$, so that $\mathrm{lg}>0$ and $1-L^{-\frac{\varepsilon}{2}}>0$.
Recall that $\mathrm{lh}(v)=L\cdot M\cdot v^r$.
With terms as in (e), the inequalities (e2)--(e5), (e7), (e8) take the following form.
\[
\begin{array}{lll}
\text{inequality} & \text{terms }(\mathrm{Cf};\mathrm{dg};\mathrm{rt}) & \mathrm{tg}\\
\hline
\text{(e2)} & \left(c_{\mathrm{fl}}+2;\,0;\,\varepsilon\right) & \frac{1}{16}\left(1-L^{-\frac{\varepsilon}{2}}\right)\\
\text{(e3), first} & \begin{array}{@{}l@{}}\left(8\,C_{0}^2\cdot b_{-}^{-1}\cdot(L\cdot M)^6;\,6r;\,\frac{\varepsilon}{2}\right),\\ \left(8\,C_{0}^2\cdot L\cdot(L\cdot M)^7;\,7r;\,\frac{\varepsilon}{2}\right)\end{array} & 1\\
\text{(e3), second} & \left(4\left(2C_{0}+c_{\kappa}+5c_{\mathrm{fl}}+6\right)\cdot(L\cdot M)^3;\,3r;\,\varepsilon\right) & 1\\
\text{(e3), third} & \left(1;\,0;\,\frac14-9\varepsilon\right) & \frac18\\
\text{(e4)} & \left(4\,C_{0}\cdot(L\cdot M)^2;\,2r+p;\,\varepsilon\right) & \frac12\\
\begin{array}[t]{@{}l@{}}\text{(e5), for}\\ v\ge\max\{4,4C_{0}^2\}\end{array} & \left(24\,(L\cdot M)^3;\,3r;\,1\right) & 1\\
\text{(e7)} & \left(L^3\cdot M;\,r;\,2\right) & 1\\
\text{(e8)} & \left(6N_{*}\cdot M;\,r+1;\,2\right),\ \left(3N_{*}\cdot M;\,r;\,1\right) & \pi^2
\end{array}
\]
To the rows (e4), (e5), (e8) the following remarks apply.
\begin{itemize}
\item[] (e4): the sum of the terms is at most $\frac12<1$; so the strict inequality holds.
\item[] (e5): after multiplication by $\exp(v)$ it states
\[
12\,\mathrm{lh}(v)^3\cdot\left(\exp\left(v-\frac{v^{p_0}}{2}\right)+\exp\left(v-\frac{v^{2p}}{2C_{0}^2}\right)\right)\le1.
\]
Let $v\ge\max\{4,4C_{0}^2\}$.
From $v^{p_0}\ge v^2\ge4v$ (as $v\ge1$, $p_0\ge2$ and $v\ge4$) one obtains $v-\frac{v^{p_0}}{2}\le-v$, and from $v^{2p}\ge v^2\ge4C_{0}^2\cdot v$ (as $2p\ge2$ and $v\ge4C_{0}^2$) one obtains $v-\frac{v^{2p}}{2C_{0}^2}\le-v$.
So the left side is at most $24\,\mathrm{lh}(v)^3\cdot\exp(-v)=24\,(L\cdot M)^3\cdot v^{3r}\cdot\exp(-v)$.
\item[] (e8): $6N_{*}\cdot M\cdot v^r\cdot\exp(-v)\cdot\left(v\cdot\exp(-v)+\frac12\right)=6N_{*}\cdot M\cdot v^{r+1}\cdot\exp(-2v)+3N_{*}\cdot M\cdot v^r\cdot\exp(-v)$.
\end{itemize}
By (El7) each inequality of the table holds for all $v$ that are not smaller than its number $\mathrm{Vbd}$, for (e5) under the additional requirement $v\ge\max\{4,4C_{0}^2\}$.
The remaining inequalities are the following.
\begin{itemize}
\item[] (e1), first: $v\ge8$.
\item[] (e1), second: for $v\ge31L+17$ one has $v^r\ge v>31L+16$ (as $v\ge1$ and $r\ge1$).
\item[] (e1), third: for $v\ge2r$, by (El6), $\left(1+\frac{3\,\mathrm{lg}}{v}\right)^r\le\exp\left(\frac{3r\cdot\mathrm{lg}}{v}\right)\le\exp\left(\frac{3\,\mathrm{lg}}{2}\right)=L^{\frac34}<L$.
\item[] (e1), fourth: for $v\ge8$ one has $p_0\cdot v^{p_0-1}\ge9v\ge72\ge28$ (as $p_0\ge9$ and $v\ge1$).
\item[] (e6): after multiplication by $\exp(v)$ it states $C_{\Sigma}\cdot\exp\left(v-\gamma_{\Sigma}\cdot M\cdot v^r\right)\le1$.
Let $v\ge\max\left\{1,\frac{2}{\gamma_{\Sigma}\cdot M},\log^{+}C_{\Sigma}\right\}$.
Then $v^r\ge v^2$ (as $r\ge2$) and $\gamma_{\Sigma}\cdot M\cdot v^2\ge2v$, so that $v-\gamma_{\Sigma}\cdot M\cdot v^r\le-v$, and $C_{\Sigma}\cdot\exp(-v)\le1$.
\end{itemize}
Let $\mathrm{Vbd}_1$ be the largest of the 8 numbers $\mathrm{Vbd}$ of the table and of the numbers $8$, $31L+17$, $2r$, $4C_{0}^2$, $\frac{2}{\gamma_{\Sigma}\cdot M}$, $\log^{+}C_{\Sigma}$.
Then $\mathrm{Vbd}_1<\infty$, and (e1)--(e8) hold for all $v\ge\mathrm{Vbd}_1$.
Consequently the number
\[
e':=\min\left\{e_0,\ \frac13,\ \exp(-\mathrm{Vbd}_1)\right\}>0
\]
lies in the set $\mathrm{Sadm}_1$ defined in 5.1.
Hence $\sup\mathrm{Sadm}_1\ge e'>0$, that is, the supremum in the definition of $e_1$ is positive, and
\[
e_1=\frac12\sup\mathrm{Sadm}_1>0,\qquad e_1\le\frac12\min\left(e_0,\frac13\right),
\]
the second inequality because $\mathrm{Sadm}_1\subset\left(0,\min\left(e_0,\frac13\right)\right]$;
moreover (e1)--(e8) hold for all $v\ge-\log e_1$, by (5.1).

\emph{Dependence.}
The set $\mathrm{Sadm}_1$ is described by $e_0$, the numeral $\frac13$ and the inequalities (e1)--(e8).
By the table of Step 4 these inequalities contain, besides $v$ and numerals, exactly the letters $L$, $M$, $b_{-}$, $r$, $p_0$, $p$, $\varepsilon$, $c_{\mathrm{fl}}$, $c_{\kappa}$, $C_{0}$, $\gamma_{\Sigma}$, $C_{\Sigma}$, $N_{*}$.
So $e_1$ is a function of $e_0$ and of these constants.
The letter $\bar{m}$ has the row $119>118$.
The letters $N$, $T$, $\theta$ are not letters of Proposition~14.2\,(a); they occur in none of the rows 1--118 and in none of (e1)--(e8) (in the rows of the kind X they are bound variables of the property asserted by the statement named as the reference of the row; see Step 3).
So $e_1$ depends neither on $\bar{m}$ nor on $N$, $T$, $\theta$.

\subsubsection*{(g) The line (C10)}

\emph{The number $e_G$.}
Let $\bar{m}\in(0,1]$.
By 6.1, with the display (6.1) and $h\ge1$, the set
\[
\left\{e>0:\ e\cdot h^{\frac12}\le e_1,\ e\cdot h^{\frac12}\cdot\log h\le1\right\}
\]
is the interval $(0,e_G]$ with
\[
e_G=\min\left\{e_1\cdot h^{-\frac12},\ h^{-\frac12}\cdot(\log h)^{-1}\right\},
\]
the second entry being absent if $h=1$.
Since $e_1>0$ by (f) and $h<\infty$, one has $e_G>0$.
The left sides $e\cdot h^{\frac12}$ and $e\cdot h^{\frac12}\cdot\log h$ do not decrease with $e$, as $\log h\ge0$; so both conditions remain true for every smaller $e>0$.

\emph{The numbers $e_{\sharp}$.}
By Proposition~35.10, $e_{\sharp}(\beta')>0$ for every probe $\beta'$.
The three two-forms $\beta_0=\beta$ and $\beta_{\pm}=\beta\pm\tau_{x_0}\beta$ are probes (Section~I.7.6).
Indeed, they are real, smooth and compactly supported;
$\operatorname{supp}\beta$ lies in a ball of radius $r_\beta$, and $\operatorname{supp}\tau_{x_0}\beta$ in the translate of this ball by $x_0$;
a point of the first ball and a point of the second have a Euclidean distance at least $|x_0|_2-2r_\beta=d_0$, and $d_0>0$ by (Pb);
so the two supports are disjoint, $\beta_{\pm}=\beta$ on $\operatorname{supp}\beta$, and $\beta_{\pm}\ne0$ because $\beta\ne0$.

\emph{The number $e_2$.}
Let $e_G>0$, $e_{\sharp}(\beta_{\mathrm{ix}})>0$ for $\mathrm{ix}\in\{0,+,-\}$, $0\le C'_{\bar{\kappa}}<\infty$, $b(f)>0$ and $\sigma>0$.
The map $x\mapsto x^{\sigma}$ is continuous and strictly increasing on $(0,\infty)$, with the limit $0$ at $0$.
So, for $x>0$, the condition (Z10.s4), that is,
\[
\frac52\,C'_{\bar{\kappa}}\cdot x^{\sigma}\le\frac{b(f)}{4},
\]
holds if and only if
\[
x\le x_4:=\left(\frac{b(f)}{10\,C'_{\bar{\kappa}}}\right)^{\frac{1}{\sigma}}\in(0,\infty);
\]
if $C'_{\bar{\kappa}}=0$, it holds for all $x$, and then $x_4:=+\infty$.
Hence the set of the numbers $x>0$ that satisfy (Z10.s1), (Z10.s2) and (Z10.s4) is the interval
\[
(0,e_2],\qquad e_2=\min\left\{e_G,\ e_{\sharp}(\beta_0),\ e_{\sharp}(\beta_+),\ e_{\sharp}(\beta_-),\ x_4\right\}>0.
\]
So the three conditions hold for every $x\in(0,e_2]$, each of them remains true when a number satisfying it is replaced by a smaller positive number, and $e_2$ is the largest number that satisfies them.

\emph{The given numbers.}
$\sigma>0$ by (S2).
$b(f)>0$ is part of (Pb).
$C'_{\bar{\kappa}}<\infty$ holds for the following reasons.
For each of the three probes $\beta'$, the bounds $\mathrm{Cb}<\infty$ (Proposition~35.3\,(c1)), $\mathrm{Cpv}_{\beta}<\infty$ (Proposition~36\,(g)), $C_{1}^{\mathrm{N2}}<\infty$ and $C_{2}^{\mathrm{N2}}<\infty$ (Proposition~37.2\,(d), (e)), the fact that $b(dd^{*}\beta')$ is a real number (Proposition~37.2\,(a)), and $D_{\mathrm{be}}<\infty$ (as $\operatorname{supp}\beta'$ is compact) give $C^{G}_{\beta'}<\infty$ and $C_{3\mathrm{p}}<\infty$.
Further, $s_{B}<\infty$, because the $\beta_{\mathrm{ix}}$ are smooth with compact support.
Next, $\varepsilon_5\left(\frac1L\right)<\infty$ (Lemma 1, row 30: $C_{6}<\infty$, and $\kappa_{(0)}>0$ for the constant $\kappa_{(0)}$ of row 20, by the rows 19, 20; Proposition~30.1\,(c)) gives $C_{\mathrm{mat}}<\infty$ and $C_{\Gamma}^{(16)}<\infty$.
Then $C_{U}$, $C_{\bar{\kappa}}$, $C'_{\bar{\kappa}}$ are finite by their formulas (rows 133--135), and they are $\ge0$ (Step~6.4).

\emph{Pairs $(\beta,x_0)$ with (Pb).}
Let $\bar{m}\in(0,1]$.
Probes $\beta$ with $b(f)>0$, where $f=dd^{*}\beta$, exist by \cite[Theorem~21.1(a)]{QEDI}.
Fix such a $\beta$, and a number $r_\beta>0$ such that $\operatorname{supp}\beta$ lies in a ball of radius $r_\beta$; such a number exists because $\operatorname{supp}\beta$ is compact.
One has $\|a\|_{\infty}<\infty$, because $a=d^{*}\beta$ is smooth with compact support.
For $|x_0|_2>2r_\beta$ one has $d_0=|x_0|_2-2r_\beta>0$, and $d_0\to\infty$ as $|x_0|_2\to\infty$.
The inequality of \cite[Section~2.8]{QEDI} that enters (Pb) reads
\[
\frac{6}{\bar{m}^2}\cdot\|a\|_{\infty}^4\cdot\left(\frac{4\pi\cdot r_\beta^3}{3}\right)^2\cdot\left(\bar{m}^2+\left(\bar{m}+\frac{1}{d_0}\right)^2\right)\cdot\frac{\exp\left(-2\bar{m}\cdot d_0\right)}{4\pi^2\cdot d_0^2}\le\frac{b(f)}{12}.
\]
In its left side the factor $\frac{6}{\bar{m}^2}\cdot\|a\|_{\infty}^4\cdot\left(\frac{4\pi\cdot r_\beta^3}{3}\right)^2$ does not contain $x_0$, and for $d_0\ge1$ the remaining factor is at most
\[
\frac{\bar{m}^2+(\bar{m}+1)^2}{4\pi^2\cdot d_0^2};
\]
so the left side tends to $0$ as $|x_0|_2\to\infty$.
The right side $\frac{b(f)}{12}$ is positive.
So the inequality holds for all $x_0$ with $|x_0|_2$ larger than some number, and (Pb) holds for these pairs.
\qed

\subsubsection*{Consequence}

Consider the rows 1, \ldots, 136 in their order.
A row of the kind N, F or D is defined as soon as the earlier rows are chosen:
for the rows $\le116$ by Lemma 1;
for the rows 118, 121, 136 by (f) and (g);
for the rows 123, 127 by Propositions~35.3\,(c1) and 35.10;
for the rows 120, 126, 128--135 by their formulas.
A row of the kind X is defined by the assertion of the statement named as its reference.
In a row of the kind C the set of the admissible choices is not empty for every admissible choice of the earlier rows:
by (a) for the rows 1, 31, 32;
by (b) for row 33;
by (c) for row 78;
by (d) for the rows 110--114;
by (e) for row 117;
row 119 is the interval $(0,1]$;
for row 122 see the last paragraph of (g).
So admissible constants exist.

\subsection*{3. Conditions on the constants}

The proofs of Propositions~14.2 and 14.3 rest on the following conditions on the constants, and on no other condition on a constant.
Each of them is a condition of Proposition~14.2\,(a), (b) or (d).
Each condition is followed by the line of Section~I.8.2 or the label of Proposition~14.2 under which it stands, and by the places where it enters.

\begin{description}
\item[$0<\rho<1$.]
Line (C0), row 1.
Enters: $\rho>0$ in Lemma 1, rows 12, 15--18, 20, 22, 29, 61, 67;
$\rho<1$ in Lemma 1, row 13 (upper end of the range of $c_{\mathrm{hl}}$), and in (f9) ($12(1+\rho)\le24$).
This range of $\rho$ is also the range for which the Propositions of this paper that the two proofs cite are asserted.
\item[The weights $w$.]
Line (C0), row 2.
Enters: Lemma 1, rows 2, 10, 14, 26 ($\mathrm{sw}>0$);
(f9) (the closure relation).
\item[The values of $\bar{\varepsilon}$, $c_{X}$, $\kappa_{0}$, $K_{0}$, $c_{\mathrm{rb}}$, $c_{\mathrm{fl}}$.]
Line (C0), rows 3--8.
Enters: Lemma 1, rows 3--8, 31, 91, 113;
(f1), (f5), (f6), (f7), (f8), (f9), (f10);
the proofs of Proposition~14.3\,(a), (b), (d).
\item[$b=1$.]
Line (C0), row 9.
Enters: $b>0$ in Lemma 1, rows 34--36, 56, 59, 61, 64, 67--70, and in (f4), (f5);
$b=1$, through $b_{-}<1$, in Lemma 1, rows 65 ($C_{5}\ge1$), 69, 71, and in (f10) ($\mathrm{gGa}_{1}\ge1$).
\item[$1-\bar{\varepsilon}<\alpha<1$.]
Line (C1), row 31.
Enters: Lemma 1, rows 49--52, 63 ($\alpha<1$);
Step 3, row 24 ($0\le\alpha\le1$);
(f6);
the proof of Proposition~14.3\,(d) ($1-\alpha<\bar{\varepsilon}$).
\item[$\kappa\ge7\kappa_{0}+7$.]
Line (C2), row 32.
Enters: Lemma 1, rows 53--55 ($\kappa>0$) and 91 ($\kappa\ge\kappa_{0}$);
(f2), (f8), (f9);
the proof of Proposition~14.3\,(b).
\item[(L1)]
Line (C3).
Enters: (f1).
\item[(L2)]
Line (C3).
Enters: (f8).
\item[(L3)]
Line (C3).
Enters: (f2).
\item[(L5)]
Line (C3).
Enters: Lemma 1, rows 33--48, 58, 65, 77, 78, 87;
Step 5.3, (fl) ($L^{-\frac{\varepsilon}{2}}<1$);
6.1, 6.2 ($L>1$);
(f1), (f2), (f8), (f9) ($L\ge1$);
(f12) ($L$ odd, $L\ge4$);
the proofs of Proposition~14.3\,(c), (d), (f) ($L>1$).
\item[(M1)]
Line (C5).
Enters: Lemma 1, row 78;
Step 5.2 ($M$ a power of $L$, for \cite[Proposition~19.1]{QEDI});
Steps 6.3, 7.1 ($M\ge L^2$, for $R^{*}(e)\in\mathrm{Tor}(e)$ and \cite[Proposition~19.1(g)]{QEDI});
(f5), (f8) ($M\ge1$);
(f9) ($L-1<M$; $\mu_-<\mu_*$, through 6.3);
(f10) ($M\ge25$, hence $203M\ge226$ and $C_H\ge2$);
(f12) ($M\ge L^2$; $M$ odd).
\item[(M3)]
Line (C5).
Enters: (f3);
(f12) ($\gamma_{\mathrm{w}}>0$).
\item[(M4)]
Line (C5).
Enters: (f8) (the fourth inequality of (A9)).
\item[(M5)]
Line (C5).
Enters: (f4), (f8).
\item[{$r\ge2$, $p_0>3r+2$, $p>2p_0$.}]
Line (C7), rows 110--112.
Enters: Lemma 1, rows 110--112;
(f6);
the proofs of Proposition~14.3\,(e) ($2p_0\ge2$) and (f) ($r\ge2$, $p_0\ge9$, $2p\ge2$).
\item[{(S1) $1-\alpha<\varepsilon\le\bar{\varepsilon}$.}]
Line (C7).
Enters: Lemma 1, rows 113, 115;
Step 5.3, (fl) ($\varepsilon>0$) and (T1) ($\frac34-8\varepsilon>0$);
(f1), (f6), (f7), (f8), (f9), (f12) ($0<\varepsilon\le\frac{1}{200}$);
the proofs of Proposition~14.3\,(e), (f) ($0<\varepsilon\le\frac{1}{200}$).
\item[{(S2) $\sigma>0$; $0<s<1$.}]
Line (C7).
Enters: Lemma 1, rows 114--116;
(f9);
the proof of Proposition~14.3\,(g).
\item[$e_0>0$.]
Line (C8), row 117.
Enters: (f7), and Step 9 (there $-\log e_0$ is defined);
the proof of Proposition~14.3\,(f) ($e'>0$).
\item[(A8), first and third inequalities.]
Line (C8).
Enters: (f7), (f8);
Step 9 ($-\log e_0>0$);
the proof of Proposition~14.3\,(e), point (3).
\item[{(e1)--(e7), for all $v\ge-\log e_1$.}]
Line (C9).
Enters: Step 5.
\item[{(e1), (e8), and the third inequality of (e3).}]
Line (C9).
Enters: (e1) and (e8) in Step 6.2;
the first and third inequalities of (e1) in Steps 6.3, 7.1;
(e1) and the third inequality of (e3) in (f9).
\item[$0<\bar{m}\le1$.]
Line (C10), row 119.
Enters: Step 5.2 (a datum of \cite[Proposition~19.1]{QEDI});
6.1 ($\log_L(1/\bar{m})\ge0$, hence $h\ge1$ and $\log h\ge0$), and through it 6.2 ($\log(1/e)>0$);
(f9) ($h\ge1$; $\bar{m}\cdot h>1/L$);
the proof of Proposition~14.3\,(g) (through 6.1, and $\bar{m}>0$ in the inequality of \cite[Section~2.8]{QEDI}).
\item[{$e\cdot h^{\frac12}\le e_1$; (Z8.w); $e\le e_G$.}]
Line (C10), Proposition~14.2\,(d).
Enters: Steps 5, 6, 7.1;
(f9).
\item[(Z10.s4)]
Line (C10), Proposition~14.2\,(d).
Enters: Step 6.4;
the proof of Proposition~14.3\,(g).
\item[{(Pb): $b(f)>0$; $d_0>0$; the inequality of \cite[Section~2.8]{QEDI}.}]
Line (C10), row 122;
Proposition~14.2\,(d).
Enters: the proof of Proposition~14.3\,(g) ($b(f)>0$ for $e_2>0$; $d_0>0$ for $\beta_{\pm}\ne0$; the last paragraph of that proof shows that the whole of (Pb) can be satisfied);
the Consequence.
\end{description}

The other conditions of Proposition~14.2\,(b), namely (L4), (M2), (M6), (A1)--(A7), the second and fourth inequalities of (A8), (A9), (A10), (Z2.s1)--(Z2.s5) and (Z3.s1), can be satisfied by Proposition~14.3\,(b), (c), (e);
in (f8) and in Step 9 they are compared with one another and with \cite{QEDI}.
Among them, the fifth inequality of (A1) and (Z2.s3) are the third inequality of (A8), and the first inequality of (Z2.s4) is the first inequality of (A8) (see (f8));
these two inequalities of (A8) enter as stated under ``(A8), first and third inequalities''.

\section*{Chapter 14c. What the constants do not depend on: proofs of Proposition~14.4}

This chapter contains the proof of Proposition~14.4.

\subsection*{0. Conventions of this chapter}

Proposition~14.4 is an assertion about the list of Proposition~14.2, that is, about the definitions of the constants and about the conditions on them, in the form in which Proposition~14.2 states them.
The words ``letter'', ``chosen'', ``defined'', ``parent'', ``ancestor'', ``depends on'', ``contains'', ``enters'' and ``fermionic constant'' carry the meaning fixed in the paragraph of Proposition~14.4 that follows its clause~(d) and explains the terms used in (a)--(d);
in this chapter that paragraph is called \emph{Words}.

\subsubsection*{0.1. The list; letters}

By \emph{the list} is meant the ordered list of the constants in Proposition~14.2\,(a).
It consists of 136 entries, which name 149 letters.
The entries are numbered $1,\dots,136$ in the order in which they stand there, from the entry of $\rho$ to that of $e_2$;
the entry with the number $n$ is called row~$n$, and the entry that names a letter is called the row of this letter.
The words ``before'' and ``after'' refer to the order in which the letters stand in the list.
The lines (C0), (C0r), (C1), (C2), (C3), (C4), (C4r), (C5), (C6), (C6$'$), (C7), (C8), (C9) contain the 120 letters from $\rho$ to $e_1$, which occupy 118 rows;
these letters are called the constants of the lines (C0)--(C9).
The line (C10) contains the remaining 29 letters, from $\bar{m}$ to $e_2$.

The value of a letter is a number, with the following exceptions:
the value of $w$ is a vector of $\mathbb{R}^4$;
the values of $\varepsilon_5$ and $\mathrm{epv}$ are functions;
the values of $A_{\mathrm{ak}}$, $A'_{\mathrm{ak}}$, $\mathrm{Csy}_{\mathrm{ak}}$ and $\mathrm{Cah}_i$ are families of numbers;
the value of $\beta$ is a two-form;
the value of $x_0$ is a vector of $\mathbb{R}^3$;
finally, each of the 16 letters from $D_{\mathrm{be}}$ to $C^{G}_{\beta}$, which stand in the rows 123--128 of Proposition~14.2\,(a), has one value for each of the three probes
\[
\beta,\qquad\beta+\tau_{x_0}\beta,\qquad\beta-\tau_{x_0}\beta
\]
of (Z10.s2).

By the preamble of Proposition~14.2\,(a), resp.\ of Proposition~14.2\,(b), is meant the text of that clause that precedes its list;
the preamble of the Propositions~1.m is the text that stands under their common title, before Proposition~1.1.
The abbreviations $\mathrm{lg}$, $c$, $c_{\mathrm{bd}}$, $\mathrm{ssum}$, $\mathrm{lh}$ introduced in the preamble of Proposition~14.2\,(a) are not letters;
wherever one of them occurs, its definition is to be read in its place (\emph{Words}).
Accordingly, $\mathrm{lg}$ contains the letter $L$, and $\mathrm{lh}(v)$ contains the letters $L$, $M$, $r$, whereas $c$, $c_{\mathrm{bd}}$ and $\mathrm{ssum}$ contain no letter, since their variable of summation is bound.

\subsubsection*{0.2. Chosen letters; $\mathrm{Val}(X)$; admissible choices}

The chosen letters of the lines (C0)--(C9) are the eleven letters
\[
\rho,\quad\alpha,\quad\kappa,\quad L,\quad M,\quad r,\quad p_0,\quad p,\quad\varepsilon,\quad s,\quad e_0,
\]
and the chosen letters of the line (C10) are $\bar{m}$, $\beta$, $x_0$ (\emph{Words}).
For a chosen letter $X$ other than $\beta$ and $x_0$ let $\mathrm{Val}(X)$ denote the set of the values that satisfy the range given in the row of $X$ together with the conditions of Proposition~14.2\,(b) on $X$, that is, the conditions that are introduced there by the words ``On~$X$''.
The two letters $\beta$, $x_0$ of row~122 are chosen together, as one pair;
their condition is (Pb) of Proposition~14.2\,(d), which their row names.
Accordingly $\mathrm{Val}(\beta)$ and $\mathrm{Val}(x_0)$ denote one and the same set, namely the set of the pairs $(\beta,x_0)$, formed by a probe $\beta$ (Section~I.7.6) and a vector $x_0\in\mathbb{R}^3$, for which there is a number $r_\beta>0$ such that $\operatorname{supp}\beta$ lies in a ball of radius $r_\beta$ and the three requirements of (Pb) hold:
$b(f)>0$, $d_0=|x_0|_2-2r_\beta>0$, and the last inequality of (Pb), where $f=dd^{*}\beta$ and $a=d^{*}\beta$.
Thus $r_\beta$ is a bound variable of this definition.
The set is formed with the values of the letters of the rows before row~122, and among these letters only $\bar{m}$ occurs in it;
in its formation neither of the two letters $\beta$, $x_0$ counts among the letters before the other.
In other words, for every chosen letter $X$ the set $\mathrm{Val}(X)$ is the set of the admissible values in the sense of \emph{Words};
for $X=\beta$ and for $X=x_0$ it is the set of pairs just described.
A chosen letter $X$ other than $\beta$, $x_0$ is fixed once, at its place in the list, as an element of $\mathrm{Val}(X)$, and the pair of the values of $\beta$ and $x_0$ is fixed once, at row~122, as an element of $\mathrm{Val}(\beta)=\mathrm{Val}(x_0)$ (Sections~I.0.5 and~I.8.1).
All other letters are defined:
the row of such a letter prescribes a value as soon as the letters before it have values (preamble of Proposition~14.2\,(a); for the letters of the lines (C0)--(C9) see Section~1.3, for those of the line (C10) the proof of Lemma~0).

An admissible choice is an assignment of values to the chosen letters, made in the order of the list, with the following two properties.
The value of every chosen letter $X$ other than $\beta$, $x_0$ lies in $\mathrm{Val}(X)$, this set being formed with the values of the letters before $X$.
The pair of the values of $\beta$ and $x_0$ is an element of $\mathrm{Val}(\beta)=\mathrm{Val}(x_0)$, this set being formed with the values of the letters of the rows before row~122.
In Lemma~0, for $X=\beta$ and for $X=x_0$, the set $\mathrm{Val}(X)$ is this set of pairs.

\subsubsection*{0.3. The definition of a letter; kinds; parents}

If $X$ is a defined letter, the definition of $X$ is the definition that its row in Proposition~14.2\,(a) gives;
if $X$ is a chosen letter, the definition of $X$ is the description of $\mathrm{Val}(X)$.
Each letter belongs to one of five kinds.
\begin{description}
\item[N] The definition is a numeral expression.
\item[E] The definition is a formula in letters that stand before $X$.
For four letters, marked E$^{*}$, this formula is the one given in \cite[Section~15]{QEDI}; the four formulas are written out in Section~1.2.
\item[S] The letter $X$ is different from $c_{\rho}$, and it is a supremum, a least element or an eigenvalue of an object formed with letters before $X$; the variables over which the extremum is taken are bound.
\item[O] The row of $X$ reads ``the constant of'', or ``the constants of'', a statement; this statement is called the reference statement of $X$.
The letter $c_{\rho}$ of row~26 belongs to this kind as well: its row states an inequality for all the data that it names, and it defines $c_{\rho}$ as the largest number with which this inequality holds.
The letter $c_{\rho}$ has no reference statement, and for $c_{\rho}$ the word ``clause'' means this inequality.
\item[C] The letter $X$ is chosen.
\end{description}

Let $\mathrm{pa}(X)$ be the set of the parents of $X$, that is, of the letters named after ``depends on'' in its row (\emph{Words}), with the single exception of row~21 described below.
For a letter $X$ of the lines (C0)--(C9) this set is given in the column $\mathrm{pa}$ of Table~P (Section~1).
For the two letters $A_{\mathrm{ak}}$, $A'_{\mathrm{ak}}$ of row~21 the set $\mathrm{pa}$ consists of the letters named after ``Depend on'' in that row together with the letter $\rho$, which their reference statement names (Section~1.3, kind~O); in Table~P this letter carries the mark $+$.
Since $\rho$ is a parent of $c_{\mathrm{el}}$ (row~18) and $c_{\mathrm{el}}$ is a parent of $A_{\mathrm{ak}}$ and $A'_{\mathrm{ak}}$, the letter $\rho$ is an ancestor of these two letters in the sense of \emph{Words} as well.
Hence the sets of ancestors formed with $\mathrm{pa}$ coincide with those of \emph{Words}, and every chain in the sense of \emph{Words} is a chain in the sense of Section~0.6.
For a letter $X$ of the line (C10), $\mathrm{pa}(X)$ is the set of the letters named after ``Depends on'', resp.\ ``Depend on'', in its row of Proposition~14.2\,(a); a row that names several letters names one such set for all of them.
Concerning the parents of the letters of the line (C10), one has:
$e_1$ belongs to $\mathrm{pa}(e_G)$ and to $\mathrm{pa}(e_\sharp)$, and $e_G$, $e_\sharp$ belong to $\mathrm{pa}(e_2)$ (Section~5.4, last paragraph);
for further parents of these letters see the proof of Lemma~0;
and every parent stands before its letter.
The last fact is read off Proposition~14.2\,(a): in each of the rows 119--136, every letter named after ``Depends on'', resp.\ ``Depend on'', is a letter of an earlier row.

For the 120 letters of the lines (C0)--(C9) the sets $\mathrm{pa}(X)$ have the following two properties, which are proved in Section~1.3.
\begin{description}
\item[(Par1)] Every element of $\mathrm{pa}(X)$ stands before $X$, and the definition of $X$ gives the value of $X$ (for a chosen $X$: the set $\mathrm{Val}(X)$) as a function of the values of the letters of $\mathrm{pa}(X)$.

For the kinds N, E, S and C this holds because every letter other than $X$ that occurs in the definition of $X$ belongs to $\mathrm{pa}(X)$; there are four exceptions.
For $c_{\kappa}$ and $c_{\gamma}$ the letter $M$ occurs in the definition without belonging to $\mathrm{pa}(X)$;
Lemma~1 (Section~2) shows that the value of $c_{\kappa}$ is a function of $(\kappa,L)$, and for $c_{\gamma}$ see Section~1.3, kind S.
For $\gamma_{\Sigma}$ and $C_{\Sigma}$ the letters $\hat{g}$ and $c_{\mathrm{N3}}$ occur in the formula without belonging to $\mathrm{pa}(X)$;
they are functions of $(\rho,w)$ (Section~1.3, remark~(7)).

Let $X$ be of kind O and different from $c_{\rho}$.
Its reference statement asserts that there is a number which has the properties stated there for \emph{all} data of that statement and which depends only on letters of $\mathrm{pa}(X)$;
among these data there may be letters of the list that stand after $X$ ($\alpha$ in Proposition~1.5\,(f), $\bar{m}$ in Proposition~36\,(c); the preamble of Proposition~14.2\,(a) says the same).
The letter $X$ is such a number, fixed by a rule that assigns a number to every value of the letters of $\mathrm{pa}(X)$.
The rule is that of the preamble of Proposition~14.2\,(a) and of the row of $X$:
$X$ is the number given at the place that the row names.
Thus $A_{\mathrm{ak}}$ and $A'_{\mathrm{ak}}$ are given by the recursions of Chapter~1a, Sections~2.6 and~2.7;
$C_{G}$ is given by the display of Chapter~1b, Section~1.4;
$C_E$ and $\mathrm{Csy}_{\mathrm{ak}}$ are given by [1.10] and [1.12] of Chapter~37;
and
\[
\mathrm{Cpv}_0=10^8\cdot\max\left\{1,\ \max\left\{\mathrm{Csy}_{\mathrm{ak}}:\ |\mathrm{ak}|\le4\right\}\right\}^2
\]
(Chapter~36, Section~4.7; the table of the constants at the end of Chapter~36).
Here $A_{\mathrm{ak}}$ and $A'_{\mathrm{ak}}$ are formed with $c_{\mathrm{el}}$, $K_{\mathrm{Str}}$, $\theta_{0}$;
$C_{G}$ with $\hat{g}$, $c_{\mathrm{hl}}$, $\theta_{0}$, $\kappa_{(0)}$, $A_{\mathrm{ak}}$, $A'_{\mathrm{ak}}$, and with the absolute numbers $c(1,3)$, $c(0,4)$ (see (Par2));
$C_E$ with $\rho$, the weights and the numeral $\lambda_1$;
$\mathrm{Csy}_{\mathrm{ak}}$ with $C_E$, and $\mathrm{Cpv}_0$ with $\mathrm{Csy}_{\mathrm{ak}}$, so that both are formed with $\rho$ and the weights.
Each of these numbers is therefore a function of the letters of $\mathrm{pa}(X)$.

For $X=c_{\rho}$ the set of the numbers with which the inequality of row~26 holds is formed with $\rho$ and the weights alone (Section~1.3, kind~O), and this set has a largest element (Chapter~14b, Lemma~1, the paragraph on the rows 24--26).
The letter $c_{\rho}$ is this largest element; hence it is a function of $(\rho,w)$, the letters of $\mathrm{pa}(c_{\rho})$.
\item[(Par2)] Every variable in the definition of $X$ that is not a letter of the list is bound:
by a sum, a supremum, a maximum, a minimum, a quantifier ``for all'', the words ``the least'', or as the argument of a function, resp.\ the index of a family, that is the value of $X$.
The remaining symbols are numerals;
signs of analysis;
objects of Part~I that contain no parameter, among them $\mathbb{Z}^3$, $O_{\mathrm{cub}}$, $s_0,\dots,s_3$, the unit vectors $e_\mu$ and the normals $\hat{n}_i$ of Section~I.1.2, with which (I.1.2) defines $\hat{d}_i(x)=-g\,\hat{n}_i$ for a site $x=(c,g)$, the spin table of Section~I.2.1, and the maps $J$, $J^{*}$, $\mathrm{Phv}$ of Section~I.3.1;
the absolute numbers $c(1,3)=\mathrm{cF13}$ (rows 24, 25) and $c(0,4)$ (row 24) of \cite[Lemma~12.4]{QEDI}, which are fixed once (preamble of Proposition~14.2\,(a); for row 24 the display of Chapter~1b, Section~1.4; for row 25 the paragraph \emph{Constants} of Proposition~30.1);
the absolute numbers $\gamma_H$, $\mathrm{Ah}_1$, $C_{\mathrm{De}}$ and the function $C_{\mathrm{Hk}}$ of \cite[Section~15]{QEDI}, which enter the formulas of the rows 49, 50, 52 and are fixed once, $C_{\mathrm{Hk}}$ being evaluated at the letter $\alpha$ (Section~1.2);
and the basis $\{e^{0}_{i}\}$ of \cite[Section~5.2.5]{QEDI}, which is fixed once for each $L$.

A composite notation that occurs in a definition is read with its definition.
Such notations are the tori and the bonds, the index sets (i)--(iii) of Section~I.5.8, $\nu$, the supports, $d(t,t')$, $\mathrm{Cl}$, $D_k$, the $M$-cubes, $d_M$, the abbreviations of Section~1.2, and the operators of the inequality of row~26 together with the spaces on which this inequality is read.
Read in this way, a composite notation contains only letters that Sections~1.2 and~1.3 name for the row in question, variables that are bound in the definition, and objects of the preceding enumeration.
In particular, none of these objects, as it occurs in a definition, contains one of the eight quantities of Proposition~14.4\,(a) as a free variable.
\end{description}

\subsubsection*{0.4. Ancestors}

Let $\mathrm{anc}(X)$ denote the set of the ancestors of $X$ (\emph{Words}), that is, the smallest set of letters that contains $\mathrm{pa}(X)$ and contains $\mathrm{pa}(Y)$ for each of its elements $Y$.
Every ancestor of $X$ stands before $X$: for the lines (C0)--(C9) this follows from (Par1), and for the line (C10) from Section~0.3.

\subsubsection*{0.5. ``Depends on''}

The two cases of the definition in \emph{Words} are referred to as case~(i), in which $Q$ is a letter of the list, and case~(ii), in which $Q$ is not a letter of the list.
The sentence of \emph{Words} that lists the letters which occur in a definition without being named after ``depends on'' is a consequence of (Par1) and of Section~1.3 (for $\kappa_{(0)}$ in row~21 see Chapter~14b, Remark~3.4).
The assertion of \emph{Words} that the value of $X$ (for a chosen letter $X$: the set of its admissible values) is a function of the values of the chosen letters among the ancestors of $X$ is Lemma~0 of Section~2;
this lemma is proved for all 149 letters.
For a chosen letter the ancestors describe the dependence through the set of the admissible values;
it is in this sense that $M$ depends on $\rho$.

\subsubsection*{0.6. Conditions; ``occurs in''; chains}

A condition on a chosen letter $X$ is either the range given in its row or a condition of Proposition~14.2\,(b) that is introduced there by ``On~$X$'';
for $\beta$ and $x_0$ it is the condition (Pb) of Proposition~14.2\,(d).
The conditions (e1)--(e8) of Proposition~14.2\,(b), which are introduced there by ``On~$e_1$'', are the inequalities in the definition of the defined letter $e_1$.
The conditions of the line (C10) are the inequalities of Proposition~14.2\,(d).
With the exception of (Pb) and (W3), they occur in the definitions of $e_G$, $e_\sharp$, $e_2$;
the condition (W3) follows from (Z8.w), from $e\,h^{\frac{1}{2}}\le e_1$ and from (e8) (Proposition~14.2\,(d), first paragraph).
The displays (2.1), (fl), (T1)--(T6) of Proposition~14.2\,(c) are conditions on $(e,T,N)$:
they hold for every $\bar{m}\in(0,1]$, every $e>0$ with $e\,h^{\frac{1}{2}}\le e_1$, every $R=h\,L^{T}$ in $\mathrm{Tor}(e)$ and every $N\ge N_{\min}(R,e)$, as stated in the sentence of Proposition~14.2\,(c) according to which they are not conditions on constants.
In clause~(c) of Proposition~14.4 they do not count as conditions, with the exception of (T2)--(T5):
these hold by (e3)--(e6) (Section~I.8.2, the paragraph on the chain that carries $\rho$), and \emph{Words} counts them among the conditions on $e_1$.

A letter $Y$ \emph{occurs in} a condition if the condition contains $Y$ in the sense of \emph{Words}.
The quantities depending on $e$ that are replaced by their definitions in Proposition~14.2\,(c), (d) are those of the second paragraph of Proposition~14.2\,(d), of Section~G5.1.4 and of Sections~I.6.1, I.6.2.

A chain from $\rho$ to $M$ is a sequence of letters
\[
\rho=Y_0,\ Y_1,\ \dots,\ Y_m=M\qquad\text{with}\qquad Y_{i-1}\in\mathrm{pa}(Y_i)\quad(i=1,\dots,m).
\]

\subsubsection*{0.7. Notation}

In this chapter $X$ and $Y$ denote letters of the list and $Q$ denotes a quantity;
in Lemma~1, however, $X$ and $Y$ denote polymers.
The letter $\kappa_{(0)}$ of row~20,
\[
\kappa_{(0)}=\min\left\{\frac{\theta_{0}}{2},\ \frac{c_{\mathrm{el}}\cdot\theta_{0}}{12K_{\mathrm{Str}}}\right\},
\]
is a constant other than the letter $\kappa_{0}=54\log72$ of row~5.

\subsection*{1. The table of parents}

\subsubsection*{1.1. Table P: the 120 letters of the lines (C0)--(C9)}

The columns of Table~P contain, in this order:
the number of the row (Section~0.1);
the letter;
its line;
its kind (Section~0.3);
the set $\mathrm{pa}$ of Section~0.3, that is, the letters named after ``Depends on'' in the row, and in row~21 in addition the letter $\rho$, marked $+$ (Sections~0.3 and~1.3);
and a cross if the letter belongs to the set $\mathrm{NR}$ of Section~4.
Row~$n$ of Table~P concerns the letters of row~$n$ of Proposition~14.2\,(a), and their definitions are those that this row gives.

\medskip
\noindent
\begin{tabular}{rlllp{5.4cm}c}
no. & letter & line & kind & $\mathrm{pa}$ & $\mathrm{NR}$\\
\hline
1 & $\rho$ & (C0) & C & -- & \\
2 & $w$ & (C0) & N & -- & $\times$\\
3 & $\bar{\varepsilon}$ & (C0) & N & -- & $\times$\\
4 & $c_{X}$ & (C0) & N & -- & $\times$\\
5 & $\kappa_{0}$ & (C0) & N & -- & $\times$\\
6 & $K_{0}$ & (C0) & E & $\kappa_{0}$ & $\times$\\
7 & $c_{\mathrm{rb}}$ & (C0) & E & $c_{X}$, $K_{0}$ & $\times$\\
8 & $c_{\mathrm{fl}}$ & (C0) & E & $c_{X}$, $K_{0}$ & $\times$\\
9 & $b$ & (C0) & N & -- & $\times$\\
10 & $\lambda_{\mathrm{int}}$ & (C0r) & S & $w$ & $\times$\\
11 & $\lambda_1$ & (C0r) & N & $w$ & $\times$\\
12 & $\hat{g}$ & (C0r) & E & $\rho$, $\lambda_{\mathrm{int}}$ & \\
13 & $c_{\mathrm{hl}}$ & (C0r) & E & $\rho$ & \\
14 & $c_{*}$ & (C0r) & E & $w$, $\lambda_{\mathrm{int}}$ & $\times$\\
15 & $\theta_{0}$ & (C0r) & E & $\rho$, $\hat{g}$ & \\
16 & $K_{\rho}$ & (C0r) & E & $\rho$, $c_{\mathrm{hl}}$, $\hat{g}$ & \\
17 & $t_{0}$ & (C0r) & E & $K_{\rho}$ & \\
18 & $c_{\mathrm{el}}$ & (C0r) & E & $\rho$, $c_{*}$, $t_{0}$ & \\
19 & $K_{\mathrm{Str}}$ & (C0r) & S & $\rho$, $w$, $\theta_{0}$, $\hat{g}$ & \\
20 & $\kappa_{(0)}$ & (C0r) & E & $\theta_{0}$, $c_{\mathrm{el}}$, $K_{\mathrm{Str}}$ & \\
21 & $A_{\mathrm{ak}}$, $A'_{\mathrm{ak}}$ & (C0r) & O & $c_{\mathrm{el}}$, $K_{\mathrm{Str}}$, $\theta_{0}$, $+\rho$ & \\
22 & $c_{b}$ & (C0r) & E & $\rho$, $c_{*}$, $K_{\rho}$, $\hat{g}$, $c_{\mathrm{hl}}$ & \\
23 & $C_{\mathrm{Po}}$ & (C0r) & E & $c_{*}$, $\lambda_{\mathrm{int}}$ & $\times$\\
24 & $C_{G}$ & (C0r) & O & $\rho$, $c_{\mathrm{el}}$, $K_{\mathrm{Str}}$, $\theta_{0}$, $\hat{g}$, $c_{\mathrm{hl}}$, $\kappa_{(0)}$, $A_{\mathrm{ak}}$, $A'_{\mathrm{ak}}$ & \\
25 & $C_{6}$ & (C0r) & E & $\rho$, $c_{\mathrm{el}}$, $K_{\mathrm{Str}}$, $\theta_{0}$, $\hat{g}$, $\kappa_{(0)}$, $A'_{\mathrm{ak}}$ & \\
26 & $c_{\rho}$ & (C0r) & O & $\rho$, $w$ & \\
27 & $C_{E}$ & (C0r) & O & $\rho$, $w$ & \\
28 & $\mathrm{Csy}_{\mathrm{ak}}$, $\mathrm{Cpv}_{0}$ & (C0r) & O & $\rho$, $w$ & \\
29 & $c_{\mathrm{N3}}$ & (C0r) & E & $\rho$, $\lambda_{\mathrm{int}}$ & \\
30 & $\varepsilon_5$ & (C0r) & E & $C_{6}$, $\kappa_{(0)}$ & \\
\end{tabular}

\medskip
\noindent
\begin{tabular}{rlllp{5.4cm}c}
no. & letter & line & kind & $\mathrm{pa}$ & $\mathrm{NR}$\\
\hline
31 & $\alpha$ & (C1) & C & $\bar{\varepsilon}$ & $\times$\\
32 & $\kappa$ & (C2) & C & $\kappa_{0}$ & $\times$\\
33 & $L$ & (C3) & C & $\bar{\varepsilon}$, $c_{\mathrm{fl}}$, $c_{\mathrm{rb}}$, $\kappa$, $\kappa_{0}$ & $\times$\\
34 & $b_{L}$ & (C4) & E & $b$, $L$ & $\times$\\
35 & $b_{-}$ & (C4) & E & $b$, $L$ & $\times$\\
36 & $b_{+}$ & (C4) & E & $b$ & $\times$\\
37 & $N_{*}$ & (C4) & E & $L$ & $\times$\\
38 & $a_{L}$ & (C4) & E & $L$ & $\times$\\
39 & $G_{L}$ & (C4) & E & $L$ & $\times$\\
40 & $K_{L}$ & (C4) & E & $L$, $G_{L}$ & $\times$\\
41 & $N_c$ & (C4) & E & $L$ & $\times$\\
42 & $c_{l}$ & (C4) & E & $L$ & $\times$\\
43 & $d_U$ & (C4) & E & $L$ & $\times$\\
44 & $c_{C}$ & (C4) & E & $L$ & $\times$\\
45 & $C_{Q}$ & (C4) & N & -- & $\times$\\
46 & $\alpha^{P}$ & (C4) & N & -- & $\times$\\
47 & $\alpha^{F}$ & (C4) & N & -- & $\times$\\
48 & $\alpha_{0}$ & (C4) & E & $\alpha^{F}$, $\alpha^{P}$ & $\times$\\
49 & $c_{P}$ & (C4) & E$^{*}$ & $L$, $C_{Q}$ & $\times$\\
50 & $\gamma_{P}$ & (C4) & E$^{*}$ & $L$, $c_{P}$ & $\times$\\
51 & $\gamma^{P}$ & (C4) & E & $\gamma_{P}$ & $\times$\\
52 & $C_{P}$ & (C4) & E$^{*}$ & $L$, $\alpha$, $c_{P}$ & $\times$\\
53 & $\kappa_{\mathrm{w}}$ & (C4) & E & $\kappa$, $\kappa_{0}$ & $\times$\\
54 & $\kappa'$ & (C4) & E & $\kappa_{\mathrm{w}}$, $\kappa_{0}$ & $\times$\\
55 & $\hat{\kappa}$ & (C4) & E & $N_c$, $\kappa$, $c_{X}$ & $\times$\\
56 & $C'_{V}$ & (C4) & E & $b$ & $\times$\\
57 & $C_{T}$ & (C4) & E & $b_{-}$, $L$ & $\times$\\
58 & $\mu_-$ & (C4) & E & $L$ & $\times$\\
\end{tabular}

\medskip
\noindent
\begin{tabular}{rlllp{5.4cm}c}
no. & letter & line & kind & $\mathrm{pa}$ & $\mathrm{NR}$\\
\hline
59 & $K_{1}$ & (C4r) & E & $L$, $b$ & $\times$\\
60 & $C_{I}$ & (C4r) & E & $b_{-}$, $C_{\mathrm{Po}}$, $K_{1}$, $c_{b}$ & \\
61 & $c_{\mathrm{wt}}$ & (C4r) & E & $\rho$, $L$, $b$ & \\
62 & $a_{0}$ & (C4r) & E & $C_{I}$, $c_{\mathrm{wt}}$ & \\
63 & $C_{R}$ & (C4r) & E & $C_{G}$, $\alpha$ & \\
64 & $C_{3}$ & (C4r) & E & $L$, $b$, $C_{I}$, $C_{R}$ & \\
65 & $C_{5}$ & (C4r) & E & $L$, $C_{3}$ & \\
66 & $\gamma_{(5)}$ & (C4r) & E & $a_{0}$, $L$ & \\
67 & $c_{\mathrm{re}}$ & (C4r) & E & $b$, $L$, $\rho$, $C_{\mathrm{Po}}$ & \\
68 & $a_{z}$ & (C4r) & E & $L$, $c_{\mathrm{re}}$, $b$ & \\
69 & $k_{1}$ & (C4r) & E & $b$ & $\times$\\
70 & $C_{00}$ & (C4r) & E & $L$, $c_{\mathrm{re}}$, $k_{1}$ & \\
71 & $C_{z}$ & (C4r) & E & $C_{5}$, $C_{00}$ & \\
72 & $\gamma_{z}$ & (C4r) & E & $\gamma_{(5)}$, $a_{z}$ & \\
73 & $\gamma^{F}$ & (C4r) & E & $\gamma_{z}$ & \\
74 & $a'$ & (C4r) & E & $C_{z}$, $\gamma^{F}$ & \\
75 & $\gamma_{\mathrm{w}}$ & (C4r) & E & $\gamma^{F}$, $\gamma^{P}$ & \\
76 & $M^{F}$ & (C4r) & E & $L$, $C'_{V}$, $a'$, $\kappa'$ & \\
77 & $M^{P}$ & (C4r) & S & $L$, $\alpha_{0}$, $\gamma_{P}$ & $\times$\\
78 & $M$ & (C5) & C & $L$, $\gamma_{\mathrm{w}}$, $\alpha_{0}$, $\hat{\kappa}$, $M^{F}$, $M^{P}$ & \\
\end{tabular}

\medskip
\noindent
\begin{tabular}{rlllp{5.4cm}c}
no. & letter & line & kind & $\mathrm{pa}$ & $\mathrm{NR}$\\
\hline
79 & $C_{\mathrm{By}}$ & (C6) & E & $b_{-}$ & $\times$\\
80 & $C_{\Gamma}$ & (C6) & E & $b_{-}$, $a'$ & \\
81 & $C_{1}$ & (C6) & E & $b$, $a'$, $M$, $\alpha_{0}$, $C_{\Gamma}$, $\kappa'$ & \\
82 & $C^{F}$ & (C6) & E & $C_{1}$, $C_{\mathrm{By}}$, $b$, $a'$, $C_{\Gamma}$ & \\
83 & $C^{P}$ & (C6) & E & $C_{P}$, $M$, $\alpha_{0}$ & \\
84 & $C_{w}^{(0)}$ & (C6) & E & $C^{F}$, $C^{P}$ & \\
85 & $C_{w}^{(1)}$ & (C6) & E & $C_{w}^{(0)}$, $\gamma_{\mathrm{w}}$, $M$, $\alpha_{0}$ & \\
86 & $C_{w}$ & (C6) & E & $K_{0}$, $C_{w}^{(1)}$ & \\
87 & $c_{\gamma}$ & (C6) & S & $\gamma_{\mathrm{w}}$, $L$ & \\
88 & $C_{W}$ & (C6) & E & $C_{w}$, $c_{\gamma}$ & \\
89 & $C_{B}$ & (C6) & E$^{*}$ & $C_{P}$, $\gamma_{P}$, $L$, $M$ & \\
90 & $C_{d}$ & (C6) & E & $c_{X}$, $C_{w}$, $b$ & \\
91 & $c_{\kappa}$ & (C6) & S & $\kappa$, $L$ & $\times$\\
92 & $C_{0}$ & (C6) & E & $C_{P}$, $\gamma_{P}$, $C^{F}$, $b_{+}$, $\gamma^{F}$ & \\
93 & $\gamma_{\Sigma}$ & (C6) & E & $b_{-}$, $b_{+}$, $\rho$, $w$ & \\
94 & $C_{\Sigma}$ & (C6) & E & $b_{-}$, $b_{+}$, $\rho$, $w$ & \\
95 & $c_{S}$ & (C6) & E & $c_{X}$, $K_{0}$, $\kappa$ & $\times$\\
\end{tabular}

\medskip
\noindent
\begin{tabular}{rlllp{5.4cm}c}
no. & letter & line & kind & $\mathrm{pa}$ & $\mathrm{NR}$\\
\hline
96 & $\mathbf{S}$ & (C6$'$) & E & $C^{F}$, $\gamma^{F}$ & \\
97 & $\mathrm{hL}$ & (C6$'$) & E & $\mathbf{S}$ & \\
98 & $C_K$ & (C6$'$) & E & $\mathbf{S}$ & \\
99 & $C_K'$ & (C6$'$) & E & $\mathbf{S}$ & \\
100 & $C_H$ & (C6$'$) & E & $M$, $C^{F}$ & \\
101 & $\mathrm{CVj}$ & (C6$'$) & E & $C_H$, $\gamma^{F}$, $K_{0}$ & \\
102 & $\theta_{\mathrm{CT}}$ & (C6$'$) & E & $\mu_-$ & $\times$\\
103 & $\mathrm{vt}$ & (C6$'$) & E & $\theta_{\mathrm{CT}}$, $\gamma^{F}$, $M$ & \\
104 & $\mathrm{gGa}$ & (C6$'$) & E & $b_{-}$, $\mu_-$ & $\times$\\
105 & $h_{G}$ & (C6$'$) & E & $\mathrm{gGa}$ & $\times$\\
106 & $\mathrm{gGa}_{1}$ & (C6$'$) & E & $b_{-}$, $\mu_-$ & $\times$\\
107 & $\mathrm{alG}$ & (C6$'$) & E & $b_{-}$, $\mu_-$, $\theta_{\mathrm{CT}}$ & $\times$\\
108 & $c_{\tau}$ & (C6$'$) & E & $\mu_-$, $\mathbf{S}$ & \\
109 & $\mathrm{CWj}$ & (C6$'$) & E & $\mathrm{CVj}$, $h_{G}$, $\mathrm{alG}$, $C_{d}$, $M$, $K_{0}$ & \\
110 & $r$ & (C7) & C & -- & $\times$\\
111 & $p_0$ & (C7) & C & $r$ & $\times$\\
112 & $p$ & (C7) & C & $p_0$ & $\times$\\
113 & $\varepsilon$ & (C7) & C & $\alpha$, $\bar{\varepsilon}$ & $\times$\\
114 & $s$ & (C7) & C & $\varepsilon$, $a_{L}$, $N_{*}$ & $\times$\\
115 & $\mathfrak{e}$ & (C7) & E & $s$, $\varepsilon$ & $\times$\\
116 & $\sigma$ & (C7) & E & $\mathfrak{e}$, $a_{L}$, $N_{*}$ & $\times$\\
117 & $e_0$ & (C8) & C & $\varepsilon$, $p_0$, $p$, $c_{\mathrm{fl}}$, $M$, $L$, $C_{w}$, $c_{\gamma}$, $c_{C}$, $c_{X}$, $C_{T}$, $c_{S}$, $K_{0}$, $C'_{V}$, $a'$, $c_{\mathrm{rb}}$, $C_{d}$, $C_{B}$ & \\
118 & $e_1$ & (C9) & S & $e_0$, $L$, $M$, $b_{-}$, $r$, $p_0$, $p$, $\varepsilon$, $c_{\mathrm{fl}}$, $c_{\kappa}$, $C_{0}$, $\gamma_{\Sigma}$, $C_{\Sigma}$, $N_{*}$ & \\
\end{tabular}

\medskip

\subsubsection*{1.2. The four formulas of kind E$^{*}$}

By Section~I.9.2 the statements of \cite[Section~15]{QEDI} hold for the objects of this paper as they are stated there;
hence the formulas of that section which are written out below hold here, and they define the letters $c_{P}$, $\gamma_{P}$, $C_{P}$, $C_{B}$.
The objects of \cite[Section~15]{QEDI} that depend on the fermion, that is, on $\rho$, are
the group $\mathbb{G}_k$, which enters through $r(S)$ and $a$;
the numbers $C_{w}$, $\gamma_{\mathrm{w}}$, $c_{\gamma}$ in the hypothesis (A2) of \cite[Proposition~15.30]{QEDI}, that is, condition (A2) of \cite[Section~5.6.2]{QEDI};
the number $M$, together with the numbers formed with it, among them $s_d=(2L+1)\cdot M$ and the constant $C_{B}$ written out below;
and the number $e_0$, which enters through the condition $e_0\le\exp(-1)$ imposed at the head of \cite[Section~15]{QEDI}, that is, (Z3.s1) of Section~I.8.4, and through the range $0<e_k\le e_0$ of the charge.
The group $\mathbb{G}_k$ and the numbers $C_{w}$, $\gamma_{\mathrm{w}}$, $c_{\gamma}$, $M$, $e_0$ are those of Sections~I.4.2 and~I.8.2.
The following facts are stated in the notation of \cite[Section~15]{QEDI}.
The numbers $\gamma_H$, $\mathrm{Ah}_1$ and $C_{\mathrm{De}}$ are absolute numbers: they are defined in \cite[Sections~15.4.1, 15.4.2 and Theorem~15.14]{QEDI}, and the list of the constants in \cite[Section~15.11.1]{QEDI} gives them as absolute constants, depending on nothing.
The number $C_{\mathrm{Hk}}(\alpha)$ is a function of $\alpha$ alone (\cite[Section~15.4.2, Theorem~15.14 and Section~15.11.1]{QEDI}).
Further,
\[
c_1:=14L^{\frac{3}{2}}
\]
by \cite[Proposition~15.21\,(iv)]{QEDI}; it is a constant of that section, not a letter of the list.
(In \cite{QEDI} the numbers $\mathrm{Ah}_1$ and $\mathrm{sgm}(\gamma)$ are written $\mathsf{A}_1$ and $\boldsymbol{\varsigma}(\gamma)$.)
The definitions that lead to the four formulas are the following.
For $\gamma>0$,
\[
\mathrm{sgm}(\gamma):=\sum_{z\in\mathbb{Z}^3}\exp(-\gamma|z|)\quad(=c(\gamma)),\qquad
c_{\mathrm{bd}}(\gamma):=3\exp(3\gamma)\cdot\mathrm{sgm}(\gamma),\qquad
c_{\mathrm{bas}}(\gamma):=2L^3\exp(2\gamma L)\cdot\mathrm{sgm}(\gamma)
\]
(\cite[Section~15.1.1]{QEDI});
\[
c_{P}:=\max\left\{9L^7\cdot C_{Q}^2,\ 15L^3\cdot\mathrm{Ah}_1\right\},\qquad
C_{\mathrm{Dt}}:=c_1^2\cdot C_{\mathrm{De}}\cdot\exp(2\gamma_H L)
\]
(\cite[Section~15.7, the display preceding Proposition~15.22]{QEDI});
\begin{gather*}
S_P:=\exp\left(\frac{\gamma_H L}{2}\right)\cdot\left(L+\frac{2}{\gamma_H}\right)\cdot c_{\mathrm{bas}}\left(\frac{\gamma_H}{4}\right),\qquad
\gamma_{P}:=\min\left\{\frac{\gamma_H}{2},\ \frac{1}{2c_{P}\cdot C_{\mathrm{Dt}}\cdot S_P}\right\},\\
C_{P}:=\max\left\{C_{\mathrm{Hk}}(\alpha),\ C_{\mathrm{Dt}},\ (2c_{P})^{\frac{1}{2}},\ 4\exp(2)\right\}
\end{gather*}
(\cite[Section~15.7, the display preceding Proposition~15.24]{QEDI});
\[
C_{B}:=C_{P}\cdot\left[(s_d+4)\cdot c_{\mathrm{bd}}(\gamma_{P})+\left(\frac{2}{\gamma_{P}}+\frac{62}{\gamma_{P}^2}\right)\cdot c_{\mathrm{bd}}\left(\frac{\gamma_{P}}{2}\right)\right],\qquad
s_d:=(2L+1)M
\]
(\cite[Section~15.9]{QEDI}).
Consequently $c_{P}$ is a formula in $(L,C_{Q})$, $\gamma_{P}$ a formula in $(L,c_{P})$, $C_{P}$ a formula in $(L,\alpha,c_{P})$, and $C_{B}$ a formula in $(C_{P},\gamma_{P},L,M)$.
The list of \cite[Section~15.11.1]{QEDI} records the same dependences;
it states moreover that none of these constants depends on $k$, $\ell$, $N$, the twist, the harmonic part of a block field, a mass, $e_k$, $\varepsilon$ or $\kappa$.

In \cite[Section~15.4.1]{QEDI} the symbol $c_{\rho}$ denotes the number $2\cosh\left(\frac{1}{8}\right)\cdot\exp\left(\frac{1}{8}\right)$;
here this number is written $c_{\mathrm{rC}}$.
It is not the letter $c_{\rho}$ of the line (C0r), and it does not involve the Wilson parameter.

\subsubsection*{1.3. Verification of (Par1) and (Par2), kind by kind}

\emph{Every element of $\mathrm{pa}(X)$ stands before $X$.}
This is read off Table~P: each letter in the column $\mathrm{pa}$ of a row is the letter of a row with a smaller number.
The letters named after ``Depends on'' in the rows of Proposition~14.2\,(a) have the same property, which is read off in the same way.

\paragraph{Kinds N and E (9 and 83 letters).}
The definition given in the row is an arithmetical expression.
It is composed of numerals, $\pi$, $\mathrm{e}$, square roots, $\exp$, $\log$, $\min$, $\max$, the abbreviations $c$, $c_{\mathrm{bd}}$, $\mathrm{ssum}$ of Section~0.1, and letters;
apart from the exception described in remark~(7), these letters belong to the column $\mathrm{pa}$ of the same row of Table~P.
This is read off the rows one by one, with the following seven remarks.
\begin{itemize}
\item[(1)] Rows 2, 10, 14: $w_i$ and $w_3$ are components of the letter $w$, and $\mathrm{sw}=w_0+w_1+w_2+w_3$.
\item[(2)] Row 11: $\lambda_1$ is a number; the row names $w$ after ``Depends on'' because $\lambda_1$ is the gap of a matrix formed with $w$ (Proposition~37.1\,(b)).
\item[(3)] Row 15: the equation $3(1+\rho)\cdot\left(\exp(\theta_{0})-1\right)=\frac{\hat{g}}{2}$ has exactly one solution, namely
\[
\theta_{0}=\log\left(1+\frac{\hat{g}}{6(1+\rho)}\right).
\]
\item[(4)] Row 22: $A_1$ and $A_2$ are abbreviations defined in the same row.
\item[(5)] Row 30: the value of $\varepsilon_5$ is the function
\[
x\mapsto C_{6}\cdot\log\left(\frac{1}{1-\exp(-\kappa_{(0)}x)}\right)
\]
on $(0,\infty)$; here $x$ is its argument.
\item[(6)] The words that follow $b=1$ in row~9, the equality ``$=\frac{1}{54}$'' in row~48, and the parentheses ``(and $C_5\ge1$)'' in row~65, ``(and $\gamma^{F}\le\frac{1}{6L}$)'' in row~73, ``(and $a'\ge1$)'' in row~74 and ``(and $\gamma_{\mathrm{w}}>0$)'' in row~75 are remarks; they form no part of a definition.
\item[(7)] Rows 25, 93, 94 ($C_{6}$, $\gamma_{\Sigma}$, $C_{\Sigma}$): the row gives the formula.
In it $\mathrm{cF13}$, $C_{\mathrm{lS}}$, $C_{\mathrm{lH}}$ (row 25), resp.\ $C_{\mathrm{iv}}$ (row 93), resp.\ $C_{\mathrm{Flg}}$ (row 94) are abbreviations.
They are defined by arithmetical expressions in the paragraph \emph{Constants} of Proposition~30.1, resp.\ of Proposition~30.3, where $K_{\mathrm{iv}}$, $C_{\mathrm{sv}}$, $C_{\mathrm{uv}}$, $S_{\mathrm{jf}}$ enter as well;
the expression for $C_{\mathrm{lS}}$ contains a maximum over the multi-indices $\mathrm{ak}$ with $|\mathrm{ak}|\le3$ and a finite sum over $\mathrm{ak}'\le\mathrm{ak}$, whose variables are bound;
$\mathrm{cF13}$ is an absolute number.
When these definitions are inserted, the letters occurring in $C_{6}$ are $\rho$, $\hat{g}$, $\theta_{0}$, $K_{\mathrm{Str}}$ and $A'_{\mathrm{ak}}$, all of which belong to $\mathrm{pa}(C_{6})$.
The letters occurring in $\gamma_{\Sigma}$ and in $C_{\Sigma}$ are $b_{-}$, $b_{+}$, $\hat{g}$, $c_{\mathrm{N3}}$.
The last two do not belong to $\mathrm{pa}$;
however, $\hat{g}$ and $c_{\mathrm{N3}}$ are the formulas of the rows 12 and 29 in $\rho$ and $\lambda_{\mathrm{int}}$, and $\lambda_{\mathrm{int}}$ is determined by $w$ (row 10).
Hence $\gamma_{\Sigma}$ and $C_{\Sigma}$ are functions of $(b_{-},b_{+},\rho,w)$, which are the letters of $\mathrm{pa}$.
\end{itemize}
All denominators, and all arguments of $c$, $c_{\mathrm{bd}}$, $\mathrm{ssum}$ and $\log$, are positive.
Indeed:
\begin{itemize}
\item $\rho$, $\lambda_{\mathrm{int}}$, $c_{*}$, $\hat{g}$, $K_{\rho}$, $t_{0}$, $c_{\mathrm{el}}$, $\theta_{0}$, $c_{\mathrm{N3}}$, $C_{\mathrm{Po}}$ are positive;
\item $K_{\mathrm{Str}}\ge\left|\mathrm{Dsc}_{1}(0)\right|_2=1$ by Proposition~1.2\,(S3) at $p=0$, $\mathbf{m}=1$; hence $\kappa_{(0)}>0$, and $1-\exp(-\kappa_{(0)}x)\in(0,1)$ for $x>0$;
\item $1-\alpha>0$ (row 31);
\item $L\ge5$ and $L-1\ge4$ by (L5), and $N_{*}$, $b_{-}$, $b_{+}$, $\mu_-$, $\theta_{\mathrm{CT}}$ are positive;
\item $M\ge L^2>0$ by (M1);
\item $C_{\mathrm{iv}}\ge1$;
\item $c_{b}$, $C_{I}$, $c_{\mathrm{wt}}$, $a_{0}$, $c_{\mathrm{re}}$, $a_{z}$, $\gamma_{z}$, $\gamma^{F}$, $\gamma_{\mathrm{w}}$ are positive (proof of (6.1)--(6.3));
\item $\gamma_{P}\in(0,1]$ (\cite[Section~15.7]{QEDI}, the paragraph on the constants $C_{P}$, $\gamma_{P}$, $c_{P}$);
\item $C^{F}>4$ (Section~5.4, complement~(2)), hence $\mathbf{S}>0$.
\end{itemize}
Consequently each of the 92 expressions has a value.
(The positivity of these denominators and arguments, and with it the existence of the values, is contained in the paragraph \emph{Ranges} of Proposition~14.2\,(f).)

\paragraph{Kind E$^{*}$ (4 letters).}
These letters are treated in Section~1.2.

\paragraph{Kind S (6 letters).}
\begin{description}
\item[$\lambda_{\mathrm{int}}$ (row 10).]
On $\ell^2(O_{\mathrm{cub}})$ let $(R_iu)(g)=u(g\,s_i)$ (I.O8).
The matrix $\sum_{i\le2}w_i\,(1-R_i)$ is a real $48\times48$ matrix formed with $w_0$, $w_1$, $w_2$ and the group $O_{\mathrm{cub}}$.
Since $g\mapsto g\,s_i$ is an involution of $O_{\mathrm{cub}}$, each $R_i$ is a symmetric involution, so that $1-R_i\ge0$.
Hence the matrix is symmetric, non-negative and different from zero, and its least non-zero eigenvalue is a positive number.
The only letter that occurs is $w$.
The index $i$ is bound by the sum, and the variables $u$ and $g$ are bound in the definition of $R_i$;
no variable occurs free.
\item[$K_{\mathrm{Str}}$ (row 19).]
One has
\[
K_{\mathrm{Str}}=\sup\left\{\left|\mathrm{Dsc}_{\mathbf{m}}(p)\right|_2:\ |\operatorname{Im}p_\mu|\le\theta_{0},\ 0\le\mathbf{m}\le1\right\},
\]
which is finite by Proposition~1.2\,(S3).
The matrix $\mathrm{Dsc}_{\mathbf{m}}(p)$ is formed from the matrix $D^{(p)}$ of (I.2.6) at $\boldsymbol{\mu}=\mathbf{m}$;
this matrix contains $\rho$ and $w$ through
\[
M_{i}(g)=\frac{w_i}{2}\left(\rho-\gamma\cdot\hat{d}_i\right),
\]
the spin table and $s_0,\dots,s_3$.
The letters that occur are $\rho$, $w$, $\theta_{0}$;
the bound variables are $p$, a momentum and not the letter $p$ of row 112, and $\mathbf{m}$.
\item[$M^{P}$ (row 77).]
The letters that occur are $L$, $\alpha_{0}$, $\gamma_{P}$, and the bound variables are $M_0$ and $x$.
The set of the powers $M_0\ge L^2$ of $L$ that have the property of the definition is not empty, because $\frac{\log x}{x}\to0$ as $x\to\infty$;
hence it has a least element.
\item[$c_{\gamma}$ (row 87).]
This letter is the subject of Lemma~2 (Section~2).
The letters that occur are
$\gamma_{\mathrm{w}}$;
$L$, through the $L$-blocks and through the basis $\{\Upsilon\}$, which is formed with vectors $e^{0}_{i}$ that are fixed once for each $L$ and are the same for all $k$ and $\ell$ (\cite[Section~5.2.5 and Proposition~5.35]{QEDI});
and $L$ and $M$, through the class $\mathrm{Cl}$.
The bound variables are $k$, $\ell$, $t$, $t'$.
The letter $M$ occurs in the definition of $\mathrm{Cl}$ and does not belong to $\mathrm{pa}(c_{\gamma})$;
it drops out (Proposition~14.2\,(a), row 87), so that the value is a function of $(\gamma_{\mathrm{w}},L)$.
This is one of the exceptions named in (Par1).
\item[$c_{\kappa}$ (row 91).]
This letter is the subject of Lemma~1 (Section~2).
The definition contains $\kappa$, and it contains $M$ and $L$ through the $M$-cubes, through $d_M$ and through $\mathrm{Cl}$.
By Lemma~1 the value is a function of $(\kappa,L)$; this is one of the exceptions named in (Par1).
The bound variables are $k$, $\ell$, $\square$ and the polymer $X\in D_k$ of the sum in row~91;
this $X$ is bound by that sum and is not a letter of the list.
The remark that follows the formula in row~91, namely that $c_{\kappa}\le K_0$, that the letters $M$ and $k$ drop out and that $L$ remains through the class $\mathrm{Cl}$, is proved in Lemma~1.
\item[$e_1$ (row 118).]
The letters that occur are $e_0$ and the letters that occur in (e1)--(e8) (Table~Q, Section~5), and the bound variables are $e'$ and $v$.
The set of these $e'$ is not empty by Proposition~14.3\,(f);
if it were empty, the supremum would be read as $0$, as Section~G5.1.9 does for $e_\sharp$.
In either case the number is determined by $e_0$ and by the letters of (e1)--(e8).
\end{description}

\paragraph{Kind O (7 letters, 7 clauses).}
For each letter $X$ of this kind the following are given:
the clause of its reference statement, for $c_{\rho}$ the inequality of row~26;
the data of this clause;
and the part of the clause by which the reference statement asserts that $X$ depends only on letters of $\mathrm{pa}(X)$, for $c_{\rho}$ the reason why the inequality is formed with $\rho$ and the weights alone.
(The weights are numerals; hence a function of $\rho$ and the weights is a function of $\rho$.)
By this part of the clause a single number $X$ serves all data of the clause.
In five of the clauses, namely Proposition~1.5, the inequality that defines $c_{\rho}$ (row 26), Proposition~37.1\,(f), (g) and Proposition~36\,(c), the data are quantified by ``for all'', resp.\ ``every'', after the constant.
In Proposition~1.2\,(S6), (S7) the datum $\mathbf{m}$ stands before the constant;
there the last sentence of Proposition~1.2 states that the bounds are uniform in $\mathbf{m}$.
\begin{description}
\item[$A_{\mathrm{ak}}$, $A'_{\mathrm{ak}}$ (row 21).]
The reference statement is Proposition~1.2\,(S6), (S7).
Its data are $\mathbf{m}\in[0,1]$ (the data of Proposition~1.2), with $\mathbf{m}>0$ in (S7), and the momentum $p\in[-\pi,\pi]^3\setminus\{0\}$, resp.\ $p=p'+i\,p''$ in the strip of (S7);
the multi-index $\mathrm{ak}$ is the index of the family.
The clauses assert $A_{\mathrm{ak}}=A_{\mathrm{ak}}(\rho)$ and $A'_{\mathrm{ak}}=A'_{\mathrm{ak}}(\rho)$;
the last sentence of Proposition~1.2 states that all bounds are uniform in $\mathbf{m}\in[0,1]$;
and the second sentence of the preamble of the Propositions~1.m applies.
The set $\mathrm{pa}$ of row~21 contains $\rho$, the letter that the reference statement names (Section~0.3).
\item[$C_{G}$ (row 24).]
The reference statement is Proposition~1.5.
Its data are
$\mathbf{m}\in[0,1]$;
$x,x'\in\mathbb{Z}^3\times O_{\mathrm{cub}}$;
$n\in\mathbb{Z}^3$ and $a\in\{\pm e_\mu\}$;
a finitely supported $v$;
in (e) and (f) every mesh $\mathrm{me}>0$, every $\mu$ with $0\le\mu\le\frac{1}{\mathrm{me}}$ and all points of $\mathrm{me}\,\mathbb{Z}^3\times O_{\mathrm{cub}}$;
in (f), moreover, the H\"older order $\alpha$.
The letter $\alpha$ (row 31) stands \emph{after} $C_{G}$.
The sentence of Proposition~1.5 that introduces $C_{G}$ provides a constant $C_{G}=C_{G}(\rho)$ that serves all $\mathbf{m}\in[0,1]$ and, in (f), every $\alpha\in[0,1]$;
the second sentence of the preamble of the Propositions~1.m applies as well.
Hence a single number $C_{G}(\rho)$ serves all $\mathbf{m}\in[0,1]$ and all $\alpha\in(1-\bar{\varepsilon},1)$, the range of the letter $\alpha$.

That one number serves every $\alpha$ can also be seen directly:
with the constant $C_{G}$ of (e), the bound (f) holds for every $\alpha\in[0,1]$, as a consequence of (e) and (I.s6).
To see this, put
\[
\mathrm{rmn}:=\mathrm{me}+\min\left(|c-c'|,\ |\tilde{c}-c'|\right).
\]
Since $\mathrm{hd}=d_X(x,\tilde{x})=\max\left(|c-\tilde{c}|,\ \mathrm{me}\right)$, one has $|c-\tilde{c}|\le\mathrm{hd}$ and $\mathrm{me}\le\mathrm{hd}$.

\emph{The case $\mathrm{hd}\ge\frac{\mathrm{rmn}}{2}$.}
Each of the numbers $\mathrm{me}+|c-c'|$ and $\mathrm{me}+|\tilde{c}-c'|$ is at most $\mathrm{rmn}+\mathrm{hd}\le3\,\mathrm{hd}$.
Therefore (e) yields
\[
|G(x,x')|_2\le2C_{G}\cdot\left(\mathrm{me}+|c-c'|\right)^{-2}\le2C_{G}\cdot(3\,\mathrm{hd})^{\alpha}\cdot\left(\mathrm{me}+|c-c'|\right)^{-2-\alpha},
\]
and the same bound holds at $\tilde{x}$, with $\tilde{c}$ in place of $c$.
The norm $|G(x,x')-G(\tilde{x},x')|_2$ is at most the sum of the two bounds, and $2\cdot3^{\alpha}\le6\le6240$.

\emph{The case $\mathrm{hd}<\frac{\mathrm{rmn}}{2}$.}
By (I.s6), read at the mesh $\mathrm{me}$ as in the remark following Proposition~1.5, the sites $x$ and $\tilde{x}$ are joined by a path of at most $39\,\frac{\mathrm{hd}}{\mathrm{me}}$ bonds whose photon sites $y$ lie in the box spanned by $c$ and $\tilde{c}$.
It follows that $|y-c|\le\mathrm{hd}$, $|y-\tilde{c}|\le\mathrm{hd}$ and
\[
\mathrm{me}+|y-c'|\ge\mathrm{rmn}-\mathrm{hd}>\frac{\mathrm{rmn}}{2}.
\]
By (e), along each bond of the path the function $G(\,\cdot\,,x')$ changes by at most
\[
20\,C_{G}\cdot\mathrm{me}\cdot\left(\frac{\mathrm{rmn}}{2}\right)^{-3}=160\,C_{G}\cdot\mathrm{me}\cdot\mathrm{rmn}^{-3},
\]
and summation over the bonds of the path gives
\begin{align*}
|G(x,x')-G(\tilde{x},x')|_2&\le39\cdot160\,C_{G}\cdot\mathrm{hd}\cdot\mathrm{rmn}^{-3}\\
&=6240\,C_{G}\cdot\mathrm{hd}^{\alpha}\cdot\mathrm{rmn}^{-2-\alpha}\cdot\left(\frac{\mathrm{hd}}{\mathrm{rmn}}\right)^{1-\alpha}\\
&\le6240\,C_{G}\cdot\mathrm{hd}^{\alpha}\cdot\mathrm{rmn}^{-2-\alpha}.
\end{align*}
Finally, $\mathrm{rmn}^{-2-\alpha}$ is the larger of the two terms in the bracket of (f).
\item[$c_{\rho}$ (row 26).]
The clause is the inequality
\[
\left(D^{\mathrm{ch}}\right)^{*}D^{\mathrm{ch}}\ge c_{\rho}\,\nabla^{*}\nabla
\]
on $\ell^2(\mathbb{X};\mathbb{C}^4)$ and on the spinors of every pair $(\Gamma,\chi)$, where $\Gamma$ is a subgroup of $\mathbb{Z}^3$ of rank three and $\chi$ is a unitary character;
here $\mathbb{X}$ is the set of the fermion sites, $D^{\mathrm{ch}}$ is the operator $D_A$ of \cite[Definition~2.10]{QEDII} at $A=0$ and at lattice mass $m=0$, and $\nabla^{*}\nabla$ is the weighted Laplacian of \cite[Definition~3.17]{QEDII} at $A=0$.
By row 26 of Proposition~14.2\,(a), $c_{\rho}$ is the largest number with which this inequality holds on all these spaces.
The data, quantified after the constant, are the space $\ell^2(\mathbb{X};\mathbb{C}^4)$ and the spaces of spinors of all pairs $(\Gamma,\chi)$.
Read as operators on functions on the mirror group $W=\mathbb{Z}^3\rtimes O_{\mathrm{cub}}$ of \cite[Lemma~3.2]{QEDII}, the operators $D^{\mathrm{ch}}$ and $\nabla^{*}\nabla$, and with them the spaces on which the inequality is read, are the same for every choice of the point $v_0$ of the fundamental chamber (\cite[Remark~3.5]{QEDII});
they are formed with $\rho$ and the weights.
Hence the inequality does not contain $v_0$, which is not a letter of the list, and the set of the numbers with which it holds is formed with $\rho$ and the weights alone.
This set has a largest element (Chapter~14b, Lemma~1, the paragraph on the rows 24--26), and $c_{\rho}$, this element, depends on $\rho$ only.
\item[$C_{E}$ (row 27).]
The reference statement is the first sentence of Proposition~37.1\,(f).
Its data, quantified after the constant, are all $\mathbf{m}>0$ and all $p\in[-\pi,\pi]^3$, and the clause asserts $C_{E}=C_{E}(\rho)$.
\item[$\mathrm{Csy}_{\mathrm{ak}}$ (row 28).]
The reference statement is the first sentence of Proposition~37.1\,(g).
Its data, quantified after the constant, are all real $p$ and all $\mathbf{m}$ with $0<\mathbf{m}\le1$;
the multi-index $\mathrm{ak}$ is the index of the family.
The clause provides the numbers $\mathrm{Csy}_{\mathrm{ak}}(\rho)$.
\item[$\mathrm{Cpv}_{0}$ (row 28).]
The reference statement is the first sentence of Proposition~36\,(c), in the setting (FF) of Section~G5.6.0 with $\mathbf{m}=\bar{m}\eta$.
Its data, quantified after the constant, are all $\bar{m}\in(0,1]$, all real $\mathrm{kph}$, all $\eta\in(0,1]$ and all $R$ as in (FF).
The letter $\bar{m}$ occurs in $\mathrm{Vpl}^{\eta,R}$ and in $\mathrm{Vpl}^{\eta}$ through $\mathbf{m}=\bar{m}\eta$, and it occurs in the bounds through the explicit factors $\bar{m}^{-3}R^{-4}$, $\frac{1}{\bar{m}}$, $\frac{1}{\bar{m}}$;
the probe $\beta$, which also belongs to (FF), does not occur in $\mathrm{Vpl}^{\eta,R}$ or in $\mathrm{Vpl}^{\eta}$.
The letter $\bar{m}$ is the first letter of the line (C10), so it stands \emph{after} $\mathrm{Cpv}_{0}$;
it is one of the eight quantities of clause~(a) of Proposition~14.4.
The clause asserts $\mathrm{Cpv}_{0}=\mathrm{Cpv}_{0}(\rho)$.
\end{description}

\emph{The letters $C_E$, $\mathrm{Csy}_{\mathrm{ak}}$, $\mathrm{Cpv}_0$.}
No letter of the lines (C0)--(C9) has one of these three letters as a parent (Table~P), and none of the three occurs in a condition of Proposition~14.2\,(b) (Q.2).
If one of them were an ancestor of a letter $X$ of these lines, it would be a parent of $X$ or of an ancestor of $X$; but every ancestor of $X$ stands before $X$ and therefore belongs to the lines (C0)--(C9) as well.
Hence none of the three letters is an ancestor of another letter of the lines (C0)--(C9).
Consequently the values of the remaining $117$ letters of these lines, and the conditions of Proposition~14.2\,(b), do not involve the rules by which $C_E$, $\mathrm{Csy}_{\mathrm{ak}}$ and $\mathrm{Cpv}_0$ are fixed in Proposition~37.1\,(f), (g) and in Proposition~36\,(c).

\paragraph{Kind C (11 letters).}
For a chosen letter $X$, the set $\mathrm{Val}(X)$ is described by the range printed in the row of $X$ together with the rows of Proposition~14.2\,(b) that name $X$.
The letters occurring in these conditions are listed in Table~Q (Section~5).
Their union over all conditions on $X$ is the entry ``pa'' of the row of $X$:
\begin{center}
\begin{tabular}{ll}
$\rho$: none; & $r$: none;\\
$\alpha$: $\bar{\varepsilon}$; & $p_0$: $r$;\\
$\kappa$: $\kappa_{0}$; & $p$: $p_0$;\\
$L$: the letters of (L1)--(L5); & $\varepsilon$: $\alpha$, $\bar{\varepsilon}$ (the range of $\varepsilon$ is (S1));\\
$M$: the letters of (M1)--(M6); & $s$: the letters of (S2);
\end{tabular}
\end{center}
and, for $e_0$, the letters of (A1)--(A10), (Z2.s1)--(Z2.s5) and (Z3.s1).

Row 114, the row of $s$, refers to (S2).
In (S2) the symbol $\sigma$ serves only as a name for the left side
\[
s\cdot\left(\frac{1}{4}-9\varepsilon\right)-12\varepsilon-\frac{2a_{L}}{N_{*}},
\]
and the letters $\mathfrak{e}$, $\sigma$ of the rows 115 and 116 are defined after $s$ by the same expression.
Thus the definition of $s$ contains $\varepsilon$, $a_{L}$, $N_{*}$ and no letter that stands after $s$.

The conditions on $e_0$ contain no variable other than $v$.
By the preamble of Proposition~14.2\,(b) they are imposed for all $v\ge-\log e_0$, so that $v$ is a bound variable.

This establishes (Par1) and (Par2) for all $120$ letters of the lines (C0)--(C9).

\subsection*{2. Three lemmas}

\begin{named}{Lemma 0}
Let $X$ be a letter of the list, and let two admissible choices give the same values to those chosen letters which are ancestors of $X$.
If $X$ is defined, then $X$ takes the same value under both choices.
If $X$ is chosen, then $\mathrm{Val}(X)$ is the same set under both choices.
\end{named}

\begin{proof}
The proof is by induction on the number of the row of $X$ in Proposition~14.2\,(a).

Suppose first that $X$ is a letter of the lines (C0)--(C9).
By (Par1) and (Par2) the value of $X$, respectively the set $\mathrm{Val}(X)$, is a function of the values of the letters of $\mathrm{pa}(X)$ and of nothing else.
Let $Y\in\mathrm{pa}(X)$; then $Y$ stands before $X$.
If $Y$ is chosen, it is a chosen ancestor of $X$, and by hypothesis the two choices give it the same value.
If $Y$ is defined, then $\mathrm{anc}(Y)\subseteq\mathrm{anc}(X)$, so the two choices agree on the chosen ancestors of $Y$, and by the induction hypothesis $Y$ takes the same value under both.

\emph{The line (C10).}
Let $X$ be one of the $29$ letters of this line.
The letters of $\mathrm{pa}(X)$ stand in earlier rows (Section~0.3), and by the argument just given they take the same values under both choices.
It therefore suffices to show that the value of $X$, respectively the set $\mathrm{Val}(X)$, is a function of the values of the letters of $\mathrm{pa}(X)$ and of chosen ancestors of $X$, and of nothing else.

\emph{Ten letters: the rows 119, 120, 122, 123 and 130.}
For these letters every letter that occurs in the definition, apart from the letters of the row itself, belongs to $\mathrm{pa}(X)$, and every other variable is bound.
In detail:
\begin{description}
\item[$\bar{m}$.] The range $(0,1]$ contains no letter.
\item[$h$.] It is given by a formula in $L$ and $\bar{m}$.
\item[$\beta$, $x_0$.] These two letters are chosen together.
The set $\mathrm{Val}$ consists of the pairs $(\beta,x_0)$, formed by a probe $\beta$ (Section~I.7.6) and a vector $x_0\in\mathbb{R}^3$, such that $b(f)>0$ and there exists a number $r_\beta>0$ with the properties stated in (Pb).
Hence $r_\beta$ is bound.
The number
\[
b(f)=-6\,\mathrm{Tr}\left[\left(S_{\bar{m}}\,i\gamma\cdot a\right)^4\right]
\]
of Section~I.7.6 and the inequality of \cite[Section~2.8]{QEDI}, in the form in which (Pb) states it,
\[
\frac{6}{\bar{m}^2}\cdot\|a\|_{\infty}^4\cdot\left(\frac{4\pi\cdot r_\beta^3}{3}\right)^2\cdot\left(\bar{m}^2+\left(\bar{m}+\frac{1}{d_0}\right)^2\right)\cdot\frac{\exp\left(-2\bar{m}\cdot d_0\right)}{4\pi^2\cdot d_0^2}\le\frac{b(f)}{12},
\]
where $a=d^{*}\beta$, $f=dd^{*}\beta$, $\|a\|_{\infty}=\sup_y|a(y)|_2$ and $d_0=|x_0|_2-2r_\beta$, contain, besides $\beta$ and $x_0$, the letter $\bar{m}$ and no other letter.
\item[$D_{\mathrm{be}}$, $D_\beta$, $C_{\beta}$, $\mathbf{D}$, $\mathrm{Cb}$.] They are formed for each of the three probes $\beta$, $\beta+\tau_{x_0}\beta$, $\beta-\tau_{x_0}\beta$ of (Z10.s2), hence with $\beta$ and $x_0$.
The letters $C_{\beta}$ and $\mathbf{D}$ are formed in addition with $h$, $L$ and $\alpha$ (Section~G5.1.3: $\mathrm{asc}$ is $h/L$ or $h/L^2$, $\eta'=L^{-N'}$, and $\alpha$ is the order of the third norm), while the triples $(N',\mathrm{asc},r)$ of $\mathrm{Ind}_\beta$ are bound by the suprema; there $r$ is an element of $\{\pm1\}^3$, not the letter $r$ of row 110.
Finally, $D_\beta=D_{\mathrm{be}}+4$ and $\mathrm{Cb}=1+h^{\frac{1}{2}}\cdot C_{\beta}$ contain earlier names of the same row, which are functions of the letters of $\mathrm{pa}(X)$.
\item[$s_{B}$.] It is the maximum, over the three probes $\beta'$, of $\|dd^{*}\beta'\|_2+\|\beta'\|_{H^4}$; so it is formed with $\beta$ and $x_0$.
\end{description}

\emph{The remaining 19 letters: the rows 121, 124--129 and 131--136.}
By the preamble of Proposition~14.2\,(a), in its sentence on the letters named after ``Depends on'', the value of $X$ is determined by the values of the letters of $\mathrm{pa}(X)$, of $\rho$ and of $w$.
The letter $w$ has a single value (row 2).
The letter $\rho$ is chosen and is an ancestor of $X$, so that the two choices agree on it.
That $\rho$ is an ancestor of each of these letters is seen as follows.
\begin{itemize}
\item $\rho$ is an ancestor of $M$ (Table~P: $\rho\to c_{\mathrm{re}}\to a_{z}\to\gamma_{z}\to\gamma^{F}\to\gamma_{\mathrm{w}}\to M$); $M$ is a parent of $C_{\mathrm{fr}}$ and of $e_0$, $e_0$ is a parent of $e_1$, $e_1$ is a parent of $e_G$ and of $e_{\sharp}$, and these two are parents of $e_2$.
\item $\rho$ is a parent of the eight letters of the rows 124 and 125.
\item $\mathrm{Cpv}_\beta$ (row 124) is a parent of $C^{G}_{\beta}$, which is a parent of $C_{3\mathrm{p}}$; further, $C_{3\mathrm{p}}$ is a parent of $C_{U}$, $C_{U}$ of $C_{\bar{\kappa}}$, and $C_{\bar{\kappa}}$ of $C'_{\bar{\kappa}}$.
\item $\rho$ is a parent of $C_{6}$, $C_{6}$ of $\varepsilon_5$, $\varepsilon_5$ of $C_{\mathrm{mat}}$, and $C_{\mathrm{mat}}$ of $C_{\Gamma}^{(16)}$.
\end{itemize}
The rows 124 and 125 read ``the constants \dots\ of'' a statement, as the rows of kind O do.
For the eight letters of these two rows the word ``determined'' has the meaning fixed in (Par1): each of them is fixed by a rule that assigns a value to every value of the letters of $\mathrm{pa}(X)$.
Such rules are given by the table of the constants at the end of Chapter~36 and by [2.1] in Chapter~37.
\end{proof}

\begin{named}{Lemma 1}[$c_{\kappa}$]
Let $n\ge3$ be an integer.
The unit cubes of the torus $\mathbb{R}^3/n\mathbb{Z}^3$ are the sets $z+\left[-\frac{1}{2},\frac{1}{2}\right]^3$ with $z\in\mathbb{Z}^3/n\mathbb{Z}^3$.
Let $\mathrm{Plm}(n)$ be the set of the non-empty unions of unit cubes that are connected through common faces.
For $Y\in\mathrm{Plm}(n)$ let $d_u(Y)$ be the infimum of the lengths of the nets in $\mathbb{R}^3/n\mathbb{Z}^3$ that meet every cube of $Y$; here a net is a connected finite union of closed straight segments, points being admitted as segments, and its length is the sum of the Euclidean lengths of its segments.
Put
\[
\mathrm{Ksum}(\kappa,n):=\sum_{Y}\exp\left(-\kappa\cdot d_u(Y)\right),
\]
where the sum extends over the sets $Y\in\mathrm{Plm}(n)$ that contain the cube with the centre $0$.
Then, for every $k\ge0$, every $\ell\in\mathrm{Cl}$ and every $M$-cube $\square$ of $T_k$,
\[
\sum_{X\in D_k:\,X\supseteq\square}\exp\left(-\kappa\cdot d_M(X)\right)=\mathrm{Ksum}\left(\kappa,\frac{\ell}{M}\right).
\]
Consequently
\[
c_{\kappa}=\sup\left\{\mathrm{Ksum}(\kappa,L\cdot j):\ j\text{ an odd integer},\ j\ge13\right\},
\]
so that $c_{\kappa}$ is a function of $(\kappa,L)$ in which $M$ does not enter.
Moreover, $c_{\kappa}\le K_{0}$.
\end{named}

\begin{proof}
The objects of Section~I.1.1 are those of \cite[Section~5.2.3]{QEDI}.
The $M$-cubes of $T_k$ are centred at the points of $M\mathbb{Z}^3/\ell\mathbb{Z}^3$; a polymer is a non-empty union of $M$-cubes that is connected through faces; and $M\cdot d_M(X)$ is the infimum of the lengths of the nets in $\mathbb{R}^3/\ell\mathbb{Z}^3$ that meet every closed $M$-cube of $X$.
Let $\ell\in\mathrm{Cl}$, and put $n:=\ell/M=L\cdot j$, where $j=\ell/(L\cdot M)$ is an odd integer $\ge13$ (Section~I.1.1).
An $M$-cube is determined by its centre, and a polymer by the set of the centres of its $M$-cubes.
The dilation $y\mapsto y/M$ maps $\mathbb{R}^3/\ell\mathbb{Z}^3$ onto $\mathbb{R}^3/n\mathbb{Z}^3$ and the closed $M$-cube with the centre $z$ onto the unit cube with the centre $z/M$.
It carries common faces to common faces and nets to nets, and it divides all lengths by $M$.
It therefore induces a bijection $X\mapsto Y(X)$ of $D_k$ onto $\mathrm{Plm}(n)$ with $d_M(X)=d_u\left(Y(X)\right)$.
The level $k$ enters neither $\mathrm{Plm}(n)$ nor $d_u$.
The translation by minus the centre of $\square$, which is a vector of $M\mathbb{Z}^3$, maps the polymers containing $\square$ onto the polymers containing the $M$-cube with the centre $0$, and it preserves $d_M$.
Together these facts give the first identity.

As $\ell$ ranges over $\mathrm{Cl}=\{L\cdot M\cdot j:\ j\text{ odd},\ j\ge13\}$, the quotient $\ell/M$ ranges over $\{L\cdot j:\ j\text{ odd},\ j\ge13\}$, whatever $M$ is; this gives the formula for $c_{\kappa}$.

Finally, $\kappa\ge7\kappa_{0}+7\ge\kappa_{0}$ (line (C2)) and $d_M\ge0$.
Hence each of the sums is at most
\[
\sum_{X\in D_k:\,X\supseteq\square}\exp\left(-\kappa_{0}\cdot d_M(X)\right)\le K_{0},
\]
by \cite[Theorem~13.2(t2)]{QEDI}, which holds for all $k\ge0$ and all $\ell\in\mathrm{Cl}$.
\end{proof}

\begin{named}{Lemma 2}[$c_{\gamma}$]
For every $k\ge0$, every $\ell\in\mathrm{Cl}$, every index $t$ and each of the three index sets (i), (ii), (iii) of Section~I.5.8,
\[
\sum_{t'}\nu(t')\cdot\exp\left(-\gamma_{\mathrm{w}}\cdot d(t,t')\right)\le2L^3\cdot\exp\left(\gamma_{\mathrm{w}}\cdot(L-1)\right)\cdot c(\gamma_{\mathrm{w}}).
\]
Hence $c_{\gamma}$, the supremum of the left side over $k$, $\ell$, $t$ and the three index sets, is a finite number.
At $k=0$ the sum over the index set (i) equals
\[
\sum_{c\in T^0_0}4\exp\left(-\gamma_{\mathrm{w}}\cdot d(t,c)\right),
\]
which is the expression that this sum has for $k\ge1$.
\end{named}

\begin{proof}
By Section~I.5.8 and by the paragraph \emph{Indices} of Section~G4.0.4, which agrees with the paragraph \emph{Indices} of \cite[Section~5.2.3]{QEDI}, an index $t$ is one of the following: a fermion site or a unit site, each with a spinor component $\beta$ that enters neither supports nor distances; a fine bond; a unit bond; or a basis element $\Upsilon$.
Supports and distances are those of \cite[Section~5.2.3]{QEDI}, read with (I.POS).
The support of a site is the photon site under it, the support of a bond is the set of its two end points, and the support of $\Upsilon$ is the set of the end points of the bonds at which $\Upsilon$ does not vanish; $d(t,t')$ is the sup distance of the supports on the torus.

\emph{Level $0$, index set (i).}
The indices are the pairs $(x,\beta)$ with $x=(c,g)\in T^{\mathrm{f}}_0=T^0_0\times O_{\mathrm{cub}}$ and $\beta\in\{1,\dots,4\}$.
They carry the weight $\nu=\frac{1}{48}$, and $d\left(t,(x,\beta)\right)=d(t,c)$ by (I.POS).
Over one photon site $c$ there lie $48$ sites $x$, each with $4$ components, and their weights add up to $48\cdot4\cdot\frac{1}{48}=4$.
This proves the last assertion of the lemma.

\emph{A unit site attached to every index.}
Define $y_t$ as follows.
If $t$ is a site over $c$ or at $c$, or a bond with the first end point $c$ (for a positively oriented unit bond, its initial point), put $y_t:=[c]$.
If $t=e^{y}_{i}$ with $y\in L\mathbb{Z}^3/\ell\mathbb{Z}^3$, put $y_t:=y$.
If $t=\delta_b$, let $y_t$ be the initial point of $b$.
A point of the support of $t$ then lies at a distance from $y_t$ that is
\begin{itemize}
\item $<\frac{1}{2}$ for sites, since $|c-[c]|<\frac{1}{2}$;
\item $<\frac{3}{2}$ for bonds and for $\delta_b$, since the end points of a bond are at the distance $\eta_k\le1$, respectively $1$;
\item $\le\frac{L-1}{2}$ for $t=e^{y}_{i}$, since $e^{y}_{i}$ is a function on the interior bonds of the $L$-block $B(y)=\left\{y'':\ |y''-y|<\frac{L}{2}\right\}$ (\cite[Section~5.2.5]{QEDI}, \cite[Proposition~15.21(ii)]{QEDI}) and $L$ is odd.
\end{itemize}
As $L\ge5$ by (L5), each of these numbers is at most $\frac{L-1}{2}$.
Let $a\in\operatorname{supp}t$ and $a'\in\operatorname{supp}t'$ be points with $d(a,a')=d(t,t')$.
The triangle inequality gives
\[
d(y_t,y_{t'})\le d(y_t,a)+d(a,a')+d(a',y_{t'})\le d(t,t')+L-1,
\]
and therefore
\[
\exp\left(-\gamma_{\mathrm{w}}\cdot d(t,t')\right)\le\exp\left(\gamma_{\mathrm{w}}\cdot(L-1)\right)\cdot\exp\left(-\gamma_{\mathrm{w}}\cdot d(y_t,y_{t'})\right).
\]

\emph{Multiplicities.}
Let $y'$ be a unit site.
The sum of $\nu(t')$ over those indices $t'$ of one index set for which $y_{t'}=y'$ is as follows.
\begin{itemize}
\item[(i)] It equals $4$: for $k\ge1$ there are four components with $\nu=1$, and for $k=0$ the sum was computed above.
\item[(ii)] It equals $3$, because three positively oriented unit bonds start at $y'$.
\item[(iii)] It is at most $(2L^3-3L^2+1)+3$.
Indeed, if $y'$ is the centre of an $L$-block, the vectors $e^{y'}_{i}$ contribute their number $\dim\ker\tau_0=2L^3-3L^2+1$ (\cite[Proposition~15.21(i)]{QEDI}); and there are at most three elements $\delta_b$, because three positively oriented bonds start at $y'$.
\end{itemize}
Each of the three numbers is at most $2L^3$, since $3L^2\ge4$.

\emph{Summation.}
It follows that the left side of the asserted inequality is at most
\[
2L^3\cdot\exp\left(\gamma_{\mathrm{w}}\cdot(L-1)\right)\cdot\sum_{y'\in T^0_k}\exp\left(-\gamma_{\mathrm{w}}\cdot d(y_t,y')\right).
\]
Fix a representative of $y_t$ in $\mathbb{Z}^3$.
Every $y'\in T^0_k=\mathbb{Z}^3/\ell\mathbb{Z}^3$ has a representative $n'\in\mathbb{Z}^3$ with $|n'-y_t|=d(y_t,y')$, and distinct points $y'$ have distinct representatives.
Hence the last sum is at most
\[
\sum_{n\in\mathbb{Z}^3}\exp\left(-\gamma_{\mathrm{w}}\cdot|n|\right)=c(\gamma_{\mathrm{w}});
\]
this is the sentence on $c(\gamma)$ in Section~G1.3, applied to the block torus $T_k=(0,\ell,-1)$ of Section~G1.2.
Finally, $c(\gamma_{\mathrm{w}})$ is finite because $\gamma_{\mathrm{w}}>0$; this is (6.3), which follows from the definitions of the rows up to 75 alone.
\end{proof}

\subsection*{3. Proof of (a)}

\paragraph{3.1.}
Let $X$ be one of the $120$ letters of the lines (C0)--(C9).
The letter $\bar{m}$ is the first letter of the line (C10); it stands after $e_1$, and so after $X$.
By (Par1) every ancestor of $X$ stands before $X$.
Hence $\bar{m}$ is not an ancestor of $X$, that is, $X$ does not depend on $\bar{m}$ (Section~0.5, case~(i)).
(The letter $\bar{m}$ does occur in one of the clauses that define a letter, namely in Proposition~36\,(c), the clause of $\mathrm{Cpv}_0$; there it stands under ``for all'', see Section~1.3.)

The quantities $N$, the level $k$, the side $\ell$, the twist $\theta$, the running mass, the pair $(m_{\mathrm{in}},\varepsilon^{\mathrm{v}}_{\mathrm{in}})$ and $\mathscr{C}$ are not letters of the list.
Apply (Par2) to $X$ and to each ancestor of $X$, which is again one of the $120$ letters: no variable that is not a letter occurs free in these definitions.
So $X$ depends on none of these seven quantities (Section~0.5, case~(ii)).

By Lemma~0 the value of $X$, respectively the set $\mathrm{Val}(X)$, is a function of the values of those of the eleven letters
\[
\rho,\ \alpha,\ \kappa,\ L,\ M,\ r,\ p_0,\ p,\ \varepsilon,\ s,\ e_0
\]
which are ancestors of $X$.
These letters are fixed once, before $\bar{m}$, $e$, $T$, $N$ and before the data of any statement (Section~I.8.1; Proposition~14.2\,(a), line (C11)).
This proves (a).

The twist is $\theta$, and the independence of the twist includes that of the twist connection $c^{(k)}_{\theta}$.
The symbol $\mathscr{C}$ stands for an arbitrary element of $\mathfrak{X}_k$, along the flow for $a^{[k]}$; so the independence of $\mathscr{C}$ includes that of the external connection $a$ (Section~I.6.6).
None of these symbols occurs in a definition of the rows 1--118.

\paragraph{3.2. Where the quantities of (a) occur, and how they are bound.}
The table below records places at which one of the eight quantities of (a) occurs, or at which a datum occurs in whose place the quantity, or an expression containing it, is inserted when the statement is applied.
The places are taken from the definitions of the $120$ letters and from the clauses of the reference statements named in Section~1.3.
The entries are the bound variables and the data listed in Section~1.3 for the kinds S, O and C.
In the rows of the kinds N, E and E$^{*}$ none of the eight quantities occurs, except through the argument $x$ of $\varepsilon_5$ (row 30), in whose place $\mathbf{m}\cdot n$, $\bar{m}\cdot R$ and $\mu\cdot\ell$ are inserted ((8.1) in Proposition~30.1; Proposition~30.2; (8.5) in Proposition~30.3).
In the rows of kind C the only such variable is the variable $v$ of the conditions on $e_0$.
\begin{center}
\begin{tabular}{p{0.25\textwidth}p{0.68\textwidth}}
quantity & letter: how the quantity is bound\\
\hline
level $k$ (mesh $\eta_k=L^{-k}$) & $e_0$, $e_1$: the variable $v$, in whose place $v_k=-\log e_k$ is inserted, ranges over all $v\ge-\log e_0$, respectively over all $v\ge-\log e'$, where $e'$ is the variable of the supremum that defines $e_1$. $c_{\gamma}$, $c_{\kappa}$: supremum over all $k\ge0$. $C_{G}$: all meshes $\mathrm{me}>0$ (in Proposition~1.6, $\mathrm{me}=\eta_k$).\\
side $\ell$ & $c_{\gamma}$, $c_{\kappa}$: supremum over all $\ell\in\mathrm{Cl}$. $c_{\rho}$: the space $\ell^2$ of the infinite lattice and the spinors of every torus $(\Gamma,\chi)$. $\mathrm{Cpv}_0$: all $R$ as in the setting (FF) of Section~G5.6.0.\\
twist $\theta$ & $c_{\rho}$: every unitary character.\\
running mass & $K_{\mathrm{Str}}$: supremum over $0\le\mathbf{m}\le1$. $A_{\mathrm{ak}}$: all $\mathbf{m}\in[0,1]$. $A'_{\mathrm{ak}}$: all $\mathbf{m}\in(0,1]$. $C_{G}$: all $\mathbf{m}\in[0,1]$; all $\mu\in[0,1/\mathrm{me}]$. $C_E$: all $\mathbf{m}>0$. $\mathrm{Csy}_{\mathrm{ak}}$: all $\mathbf{m}\in(0,1]$. $\mathrm{Cpv}_0$: $\mathbf{m}=\bar{m}\cdot\eta$, see the rows of $N$ and $\bar{m}$.\\
$N$ ($\eta=L^{-N}$, $a_{k}=L^{k-N}$) & through $e_k=e\cdot a_{k}^{\frac{1}{2}}$, that is, through the variable $v$: see the row of the level $k$; through the side $\ell_k=R/a_{k}$ and the number $n=R/\eta$ of sites: see the row of the side; through the running mass $\mu_k=\bar{m}\cdot a_{k}$, at $k=0$ through $\mathbf{m}=\bar{m}\cdot\eta$: see the row of the running mass (Section~I.6.1; the setting (FF) of Section~G5.6.0). $\mathrm{Cpv}_0$: all $\eta\in(0,1]$.\\
$\bar{m}$ & through the running mass $\bar{m}\cdot a_{k}$: see the row of the running mass. $\mathrm{Cpv}_0$: all $\bar{m}\in(0,1]$; the factors $\bar{m}^{-3}$ and $1/\bar{m}$ occur in the bounds.\\
the pair $(m_{\mathrm{in}},\varepsilon^{\mathrm{v}}_{\mathrm{in}})$; $\mathscr{C}$ & nowhere.\\
(not among the eight) & the charge $e$, and $e_k$: $e_0$, $e_1$, through $v$ as above. The letter $\alpha$, which stands after $C_{G}$: $C_{G}$, every $\alpha\in[0,1]$.
\end{tabular}
\end{center}

\emph{Level $0$.}
Some letters serve level $0$ and the levels $k\ge1$ at once.
For each of them the entry that comes from level $0$ is the following.
\begin{itemize}
\item $C_z=\max\left(C_{5},C_{00}\right)$ and $\gamma_{z}=\min\left(\gamma_{(5)},a_{z}\right)$: the entries $C_{5}$ and $\gamma_{(5)}$ serve $k\ge1$ (Proposition~1.6\,(c), Proposition~2.3\,(a)), the entries $C_{00}$ and $a_{z}$ serve $k=0$ (Proposition~2.3\,(b); the factor $48$ of the kernels (I.1.4) is contained in $C_{00}$); through them $\gamma^{F}$ and $a'$ serve all $k\ge0$ (Proposition~2.3\,(c)).
\item $C_{1}$: it is the constant of Proposition~7 for $k\ge1$ and of Proposition~8 for $k=0$; in Proposition~8 the properties W6 and W7 are bounds on the kernels (I.1.4), with the constant $2C_{1}$.
\item $C^{F}$: it is the constant of FE1$^{\mathrm{ch}}$ and FE2$^{\mathrm{ch}}$ for $k\ge1$ and of FE3$^{\mathrm{ch}}$ at $k=0$ (Proposition~14.1\,(i)).
\item $C_{w}^{(1)}=\max\left(C_{w}^{(0)},4\exp\left(2\gamma_{\mathrm{w}}\right)\cdot M^{27\alpha_{0}}\right)$: the second entry is the constant of W3 and W5 for the diagonal families $H_{0}(s')$, $H_{0}(s'',A)$ of level $0$ at the rate $\gamma_{\mathrm{w}}$ (W17; Proposition~8\,(b), Proposition~14.1\,(ii)).
Here W1--W17 are the properties of the primitives displayed in Section~G1.11 and Section~G4.0.4; in \cite[Section~5.4.2]{QEDI} they are the properties numbered (1)--(17).
\item $C_{d}=16c_{X}\cdot C_{w}\cdot\max\left\{\frac{\pi}{2}\cdot b\cdot\mathrm{e}\cdot(1+2\mathrm{e}),\ 2+58\mathrm{e}+b\cdot\mathrm{e}\right\}$: the second entry is the row sum (I.L0.12) of level $0$ at $\exp\left(\mathfrak{i}(A)\right)\le\mathrm{e}$, which is $58\mathrm{e}+1+b\cdot\mathrm{e}$, increased by $1$ (see the remark after Proposition~13).
So $C_{d}$ is one constant for all levels $j\ge0$, which is the form (G5.T-1) of Annex~A, Section~G5.7.
\item $C_{W}=28C_{w}^2\cdot c_{\gamma}$ at all levels; here $28=4\cdot7$, and $7\ge\sqrt{48}=\mathsf{w}_W$ at level $0$ (I.L0.19).
\item $c_{\gamma}$: the weight $\frac{1}{48}$ of level $0$ cancels against the $48$ sites of a cluster (Lemma~2).
\item $e_0$: the condition (A3) consists of two inequalities, the first for $k\ge1$ and the second for $k=0$ (Proposition~14.2\,(b), row (A3)); the condition (A6) contains the numeral $112=16\cdot7$ (Section~I.8.4, (I.L0.19)).
\end{itemize}
That no letter takes one value at $k=0$ and another at $k\ge1$ follows from (Par2): the level occurs in no definition of the rows 1--118, except as a bound variable at the places listed in the row of the level $k$ in the table above.

\paragraph{3.3. Objects of the lines (C0)--(C9) that are not letters of the list.}
The symbols $\mathrm{par}^{0}$ and $\mathrm{par}$ are defined in Proposition~14.2\,(a), in the paragraph that follows the line (C11).
Their entries are letters of the rows 1--118, namely $\alpha_{0}$ (row 48), $\gamma_{\mathrm{w}}$ (row 75), $\kappa_{\mathrm{w}}$ (row 53), $C_{w}^{(1)}$ (row 85) and $C_{w}$ (row 86), with one exception: the third entry of $\mathrm{par}^{0}$ is the sum $\kappa_{\mathrm{w}}+\kappa_{0}$ of the letters of the rows 53 and 5, and this sum equals the letter $\kappa'$ of row 54.
The argument of 3.1 therefore applies to every entry.
The exponent $\mathfrak{e}_2=0$ and the thresholds $L^F=L^P=5$ of Section~I.8.2 are numerals, and $\mathrm{sw}=w_0+w_1+w_2+w_3$.
The number
\[
c_{\mathrm{cl}}=\frac{96\left(1+\sqrt{2}\right)\cdot\left(128+64\left(\frac{16}{\pi}+\frac{64}{\pi^2}\right)\right)}{4\pi^2}
\]
(Section~I.8.2, second row of the line (C0); Proposition~14.2\,(a), the paragraph that follows the line (C11)) is given by a numeral expression: it contains neither a letter nor a variable.
The functions $\mathfrak{d}_{\mathrm{ma}}(p)$ and $p^{0}_{\mathrm{ch}}(\mathbf{m})$ of Section~I.8.2, line (C0r), are formed from the matrix (I.2.6), which contains $\rho$, $w$ and objects without parameters.
In the index of $\mathfrak{d}_{\mathrm{ma}}$ the letters $\mathrm{ma}$ stand for the mass $\mathbf{m}$, so that the arguments of the two functions are $(\mathbf{m},p)$, respectively $\mathbf{m}$; the variable $p$ is bound in $p^{0}_{\mathrm{ch}}$.
As for $\varepsilon_5$, the constant is the function itself, not its value at a mass.
So the argument of 3.1 applies to these functions as well.

\paragraph{3.4. Remark.}
The same argument shows that no letter of the lines (C0)--(C9) depends on $e$, on $T$, on $e_k$, on a connection $A$, on $y$, on parameters $s'$, $s''$, on a polymer, on a system of primitives or on a placement.
Indeed, none of these is a letter standing before $e_1$, and by (Par2) none of them occurs free in a definition.

\subsection*{4. Proof of (b)}

\paragraph{4.1. The set $\mathrm{NR}$.}
Let $\mathrm{NR}$ be the set of the $58$ letters that carry a cross in Table~P.
They are the letter $w$ and the $57$ letters named in clause (b), namely
\begin{gather*}
\bar{\varepsilon},\ c_{X},\ \kappa_{0},\ K_{0},\ c_{\mathrm{rb}},\ c_{\mathrm{fl}},\ b\ \text{(7)};\qquad\lambda_{\mathrm{int}},\ \lambda_1,\ c_{*},\ C_{\mathrm{Po}}\ \text{(4)};\qquad\alpha,\ \kappa,\ L\ \text{(3)};\\
b_{L},\ b_{-},\ b_{+},\ N_{*},\ a_{L},\ G_{L},\ K_{L},\ N_c,\ c_{l},\ d_U,\ c_{C},\ C_{Q},\ \alpha^{P},\ \alpha^{F},\ \alpha_{0},\ c_{P},\ \gamma_{P},\ \gamma^{P},\ C_{P},\\
\kappa_{\mathrm{w}},\ \kappa',\ \hat{\kappa},\ C'_{V},\ C_{T},\ \mu_-,\ K_{1},\ k_{1},\ M^{P}\ \text{(28)};\\
C_{\mathrm{By}},\ c_{\kappa},\ c_{S},\ \theta_{\mathrm{CT}},\ \mathrm{gGa},\ h_{G},\ \mathrm{gGa}_{1},\ \mathrm{alG}\ \text{(8)};\qquad r,\ p_0,\ p,\ \varepsilon,\ s,\ \mathfrak{e},\ \sigma\ \text{(7)}.
\end{gather*}

\emph{The set $\mathrm{NR}$ is closed under $\mathrm{pa}$}: in Table~P every letter in the column ``pa'' of a row with a cross is the letter of a row with a cross.
The $58$ rows, with their parents, are the following.
\begin{center}
\begin{tabular}{p{0.5\textwidth}p{0.4\textwidth}}
letters & parents\\
\hline
rows 2--5, 9, 45--47, 110 ($w$, $\bar{\varepsilon}$, $c_{X}$, $\kappa_{0}$, $b$, $C_{Q}$, $\alpha^{P}$, $\alpha^{F}$, $r$) & none\\
$K_{0}$ & $\kappa_{0}$\\
$c_{\mathrm{rb}}$, $c_{\mathrm{fl}}$ & $c_{X}$, $K_{0}$\\
$\lambda_{\mathrm{int}}$, $\lambda_1$ & $w$\\
$c_{*}$ & $w$, $\lambda_{\mathrm{int}}$\\
$C_{\mathrm{Po}}$ & $c_{*}$, $\lambda_{\mathrm{int}}$\\
$\alpha$ & $\bar{\varepsilon}$\\
$\kappa$ & $\kappa_{0}$\\
$L$ & $\bar{\varepsilon}$, $c_{\mathrm{fl}}$, $c_{\mathrm{rb}}$, $\kappa$, $\kappa_{0}$\\
$b_{L}$, $b_{-}$ & $b$, $L$\\
$b_{+}$ & $b$\\
$N_{*}$, $a_{L}$, $G_{L}$, $N_c$, $c_{l}$, $d_U$, $c_{C}$, $\mu_-$ & $L$\\
$K_{L}$ & $L$, $G_{L}$\\
$\alpha_{0}$ & $\alpha^{F}$, $\alpha^{P}$\\
$c_{P}$ & $L$, $C_{Q}$\\
$\gamma_{P}$ & $L$, $c_{P}$\\
$\gamma^{P}$ & $\gamma_{P}$\\
$C_{P}$ & $L$, $\alpha$, $c_{P}$\\
$\kappa_{\mathrm{w}}$ & $\kappa$, $\kappa_{0}$\\
$\kappa'$ & $\kappa_{\mathrm{w}}$, $\kappa_{0}$\\
$\hat{\kappa}$ & $N_c$, $\kappa$, $c_{X}$\\
$C'_{V}$, $k_{1}$ & $b$\\
$C_{T}$ & $b_{-}$, $L$\\
$K_{1}$ & $L$, $b$\\
$M^{P}$ & $L$, $\alpha_{0}$, $\gamma_{P}$\\
$C_{\mathrm{By}}$ & $b_{-}$\\
$c_{\kappa}$ & $\kappa$, $L$\\
$c_{S}$ & $c_{X}$, $K_{0}$, $\kappa$\\
$\theta_{\mathrm{CT}}$ & $\mu_-$\\
$\mathrm{gGa}$, $\mathrm{gGa}_{1}$ & $b_{-}$, $\mu_-$\\
$h_{G}$ & $\mathrm{gGa}$\\
$\mathrm{alG}$ & $b_{-}$, $\mu_-$, $\theta_{\mathrm{CT}}$\\
$p_0$ & $r$\\
$p$ & $p_0$\\
$\varepsilon$ & $\alpha$, $\bar{\varepsilon}$\\
$s$ & $\varepsilon$, $a_{L}$, $N_{*}$\\
$\mathfrak{e}$ & $s$, $\varepsilon$\\
$\sigma$ & $\mathfrak{e}$, $a_{L}$, $N_{*}$
\end{tabular}
\end{center}
Hence $\mathrm{anc}(X)\subseteq\mathrm{NR}$ for every $X\in\mathrm{NR}$.
(The ancestors of $\lambda_{\mathrm{int}}$, $\lambda_1$, $c_{*}$, $C_{\mathrm{Po}}$ lie among $w$, $\lambda_{\mathrm{int}}$, $c_{*}$; these four letters are functions of the weights.)
The letter $\rho$ is not an element of $\mathrm{NR}$.
Therefore $\rho$ is an ancestor of no $X\in\mathrm{NR}$, that is, no letter of $\mathrm{NR}$ depends on $\rho$ (Section~0.5, case~(i)).
This proves (b) for the $57$ letters.

\emph{How the parents of the $58$ letters are read off.}
Forty-four of them are of the kinds N and E ($9+35$); their parents are read off the arithmetical expression that their rows in Proposition~14.2\,(a) give as the definition (Section~1.3; for $\lambda_1$ see remark~(2) there).
Six are chosen letters whose rows state the whole range: $\alpha$, $\kappa$, $r$, $p_0$, $p$, $\varepsilon$.
For the remaining eight letters the definitions that their rows give (for $L$ and $s$, the ranges) do not by themselves suffice to read off the parents:
\begin{itemize}
\item $L$ and $s$, whose rows refer to (L1)--(L5), respectively to (S2) (Section~1.3, kind C; Q.2);
\item $\lambda_{\mathrm{int}}$ and $M^{P}$ (Section~1.3, kind S);
\item $c_{P}$, $\gamma_{P}$, $C_{P}$ (Section~1.2);
\item $c_{\kappa}$, whose definition contains $M$, an ancestor of which is $\rho$; by Lemma~1 the value of $c_{\kappa}$ is a function of $(\kappa,L)$.
\end{itemize}

\paragraph{4.2. What it means.}
The chosen letters in $\mathrm{NR}$ are $\alpha$, $\kappa$, $L$, $r$, $p_0$, $p$, $\varepsilon$, $s$.
By Lemma~0 their admissible sets, and the values of the $50$ defined letters of $\mathrm{NR}$, are functions of these eight letters.
To state this explicitly, let $\mathrm{Val}_0$ be the set of the tuples $(\alpha,\kappa,L,r,p_0,p,\varepsilon,s)$ with
\begin{gather*}
1-\bar{\varepsilon}<\alpha<1;\qquad\kappa\ge7\kappa_{0}+7;\qquad L\text{ an odd integer with (L1)--(L5)};\\
r,\ p_0,\ p\text{ integers},\qquad r\ge2,\qquad p_0>3r+2,\qquad p>2p_0;\\
1-\alpha<\varepsilon\le\bar{\varepsilon};\qquad0<s<1,\qquad s\cdot\left(\frac{1}{4}-9\varepsilon\right)-12\varepsilon-\frac{2a_{L}}{N_{*}}>0,\quad\text{where }a_{L}=\frac{L^3-1}{(L-1)^3},\ N_{*}=8L^3.
\end{gather*}
Besides the eight letters, these conditions contain only the numbers
\[
\bar{\varepsilon}=\frac{1}{200},\qquad\kappa_{0}=54\log72,\qquad c_{\mathrm{rb}}=\frac{54\exp(\kappa_{0})}{36},\qquad c_{\mathrm{fl}}=\frac{61\cdot54\exp(\kappa_{0})}{36}.
\]
For every admissible choice, whatever the value of $\rho\in(0,1)$, the eight values form an element of $\mathrm{Val}_0$.
Every letter of $\mathrm{NR}$ is then either a number (for $w$, a vector of numbers), namely one of
\[
w,\ \bar{\varepsilon},\ c_{X},\ \kappa_{0},\ K_{0},\ c_{\mathrm{rb}},\ c_{\mathrm{fl}},\ b,\ \lambda_{\mathrm{int}},\ \lambda_1,\ c_{*},\ C_{\mathrm{Po}},\ C_{Q},\ \alpha^{P},\ \alpha^{F},\ \alpha_{0},\ C'_{V},\ k_{1},\ b_{+},
\]
or a function on $\mathrm{Val}_0$.
The conditions that define $\mathrm{Val}_0$ do not contain $\rho$; whether the eight values satisfy them is therefore independent of the value of $\rho$.
In particular, $(\alpha,\kappa,L)$ may be kept fixed while $\rho$ varies in $(0,1)$.

\paragraph{4.3. The exponents of Section~I.8.5.}
The table of Section~I.8.5 has ten rows.
The powers printed in its rows 1--9 are the following.
\begin{itemize}
\item[] Row 1 (thresholds): $e_k^{-\frac{3}{4}+\varepsilon}$, $e_k^{-\frac{3}{4}+2\varepsilon}$, $e_k^{-\frac{3}{4}+3\varepsilon}$; $e_k^{-1+\varepsilon}$; $e_k^{-1+2\varepsilon}$, $e_k^{-1+3\varepsilon}$.
\item[] Row 2 (weights, windows): $e_k^{-\frac{1}{4}}$, $e_k^{-\frac{1}{4}+\varepsilon}$; $e_k^{\frac{3}{4}-8\varepsilon}$, $e_k^{\frac{1}{4}-7\varepsilon}$.
\item[] Row 3 (polydisk, contours): $M^{\alpha_{0}}$; $e_k^{-\frac{3}{4}+\frac{9\varepsilon}{2}}$, $e_k^{-\frac{3}{4}+7\varepsilon}$, $e_k^{-\frac{1}{4}+3\varepsilon}$.
\item[] Row 4 (one step): $e_k^{\frac{1}{2}}$; $L^{-\frac{1}{4}+2\varepsilon}$; $e_{k+1}^{\frac{1}{4}-2\varepsilon}$; $e_k^{\frac{1}{4}-6\varepsilon}$; $e_k^{7}$.
\item[] Row 5 (trajectory): $e_k^{\varepsilon}\cdot w^{\mathrm{E}}_k$; $w^{\mathrm{m}}_k$; $e_k^{\varepsilon}\cdot w^{\mathrm{E}}_k$ (for $w^{\mathrm{E}}_k$ and $w^{\mathrm{m}}_k$ see row 2).
\item[] Row 6 (terminal): $e_{*}^{\frac{1}{4}-9\varepsilon}$ (three times); $e_{*}^{3\varepsilon}$.
\item[] Row 7 (final): $\mathfrak{e}=s\cdot\left(\frac{1}{4}-9\varepsilon\right)-12\varepsilon$; $\sigma=\mathfrak{e}-\frac{2a_{L}}{N_{*}}$.
\item[] Row 8: $\mathfrak{e}_2=0$.
\item[] Row 9 (rate in the spacing): $\eta^{\frac{1}{4}}$; and $\eta^{\frac{1}{2}}$, the power printed in \cite[Theorem~21.1(c2)]{QEDI}, \cite[Theorem~21.3(2)(c)]{QEDI} and \cite[Theorem~3.1, Part~II(ii)]{QEDI} at the places where $\eta^{\frac{1}{4}}$ stands here.
\end{itemize}
Every exponent is a rational number, or of the form $q_0+q_1\cdot\varepsilon$ with rational numbers $q_0$, $q_1$, or one of $\alpha_{0}$, $\mathfrak{e}$, $\sigma$, $\mathfrak{e}_2=0$.
So it is a formula in letters of $\mathrm{NR}$ (namely $\varepsilon$, $s$, $a_{L}$, $N_{*}$, $\alpha_{0}$), and $\rho$ is an ancestor of none of these.
The numerical factors printed in these rows are numerals and $c_{\mathrm{fl}}$, which is a letter of $\mathrm{NR}$.

Row 10 (rates) contains $\gamma_{\mathrm{w}}$ per unit of length; $\kappa_{\mathrm{w}}+\kappa_{0}$, $\kappa_{\mathrm{w}}$ and $\kappa$ per unit of $d_M$; and $\ell_H=M$.
The three rates per unit of $d_M$ are formulas in letters of $\mathrm{NR}$.
The letters $\gamma_{\mathrm{w}}$ and $M$ do not belong to $\mathrm{NR}$: $\rho$ is an ancestor of both ($\rho\to c_{\mathrm{re}}\to a_{z}\to\gamma_{z}\to\gamma^{F}\to\gamma_{\mathrm{w}}\to M$ in Table~P), and by Section~6 neither of them can be kept fixed as $\rho\to0$.
This is the content of the bracket at the end of the first sentence of (b): $\gamma_{\mathrm{w}}$ and $\ell_H=M$ carry $\rho$.

\paragraph{4.4. (L1)--(L5), (S1), (S2).}
By Table~Q (Section~5) these seven conditions contain, besides the letters $L$, $\varepsilon$, $s$ that they restrict, the following letters, in the order (L1), \dots, (L5), (S1), (S2):
\[
\bar{\varepsilon};\qquad c_{\mathrm{fl}},\ \bar{\varepsilon};\qquad\kappa,\ \kappa_{0};\qquad c_{\mathrm{rb}},\ \kappa,\ \kappa_{0};\qquad\text{none};\qquad\alpha,\ \bar{\varepsilon};\qquad\varepsilon,\ a_{L},\ N_{*}.
\]
None of them is $\rho$.
None of the nine letters $\bar{\varepsilon}$, $c_{\mathrm{fl}}$, $c_{\mathrm{rb}}$, $\kappa$, $\kappa_{0}$, $\alpha$, $\varepsilon$, $a_{L}$, $N_{*}$ is a fermionic constant in the sense of the paragraph \emph{Words} of Proposition~14.4.
Indeed, their lines are, in this order, (C0), (C0), (C0), (C2), (C0), (C1), (C7), (C4), (C4).
The references given in their rows of Proposition~14.2\,(a) are, in the same order: \cite[Definition~5.97, stage~1]{QEDI}, equal to the line (C0) of Section~I.8.2; \cite[Section~13.1.6]{QEDI} together with the Propositions 16.m and 18.m; \cite[Section~13.1.6]{QEDI}; Section~I.8.2; \cite[Section~13.1.6]{QEDI}; Section~I.8.2; Section~I.8.2; Section~I.7.6; Section~I.6.4.
Their ancestors lie among $\bar{\varepsilon}$, $c_{X}$, $\kappa_{0}$, $K_{0}$, $c_{\mathrm{rb}}$, $c_{\mathrm{fl}}$, $\alpha$, $\kappa$, $L$ (4.1), whose lines are (C0)--(C3) and whose references are \cite[Definition~5.97, stage~1]{QEDI}, \cite[Section~13.1.6]{QEDI}, the Propositions 16.m and 18.m, and Section~I.8.2; neither $w$ nor $\rho$ is among these ancestors.
None of these lines is (C0r) or (C4r), and none of these references is one of the statements named in the definition of a fermionic constant.

The same test applies to the three letters that are restricted.
The letter $\varepsilon$ is one of the nine letters.
The letter $L$ (line (C3), reference Section~I.8.2) is one of the nine ancestors just listed, and its own ancestors lie among them.
The letter $s$ has the line (C7), the reference Section~I.8.2 and the parents $\varepsilon$, $a_{L}$, $N_{*}$ (row 114); so its ancestors are these three letters and their ancestors.
Hence $L$, $\varepsilon$ and $s$ are not fermionic constants either.
In (S2) the symbol $\sigma$ is the name of the left side (Section~1.3, kind C) and adds no letter.

Finally, $c_{\mathrm{fl}}=\frac{61\cdot54\cdot72^{54}}{36}$ is given by a numeral expression (rows 4--6 and 8).
Section~I.8.2 lists it among the absolute numbers, with the value it has in \cite[Section~13.1.6]{QEDI}; besides that section, the reference in its row of Proposition~14.2\,(a) names the Propositions 16.m and 18.m, in which it occurs.
It is not formed with $D_A$.
This proves the last sentence of (b).

\paragraph{4.5. Remark.}
The set $\mathrm{NR}$ consists of all letters of the lines (C0)--(C9), other than $\rho$, that do not have $\rho$ among their ancestors.
Indeed, in Table~P every row without a cross, other than row 1, has in its column ``pa'' the letter $\rho$ or the letter of an earlier row without a cross; by induction on the number of the row, $\rho$ is an ancestor of the letter of every such row.

\subsection*{5. Proof of (c)}

\subsubsection*{5.1. Table Q: all conditions, with the letters that occur in them}

\paragraph{Q.1. The ranges.}
These are the ranges given in the entries of Proposition~14.2\,(a): the ranges of the chosen letters, and the range of the variable $e'$ in the entry of $e_1$.
Next to each range stand the letters that occur in it.

\noindent
\begin{tabular}{p{0.4\textwidth}p{0.5\textwidth}}
range & letters that occur in it\\
\hline
$\rho\in(0,1)$ & none\\
$\alpha\in(1-\bar{\varepsilon},1)$ & $\bar{\varepsilon}$\\
$\kappa\ge7\kappa_{0}+7$ & $\kappa_{0}$\\
$L$ an odd integer, $M$ a power of $L$ & contained in (L5), resp.\ in (M1)\\
$r\ge2$ & none\\
$p_0>3r+2$ & $r$\\
$p>2p_0$ & $p_0$\\
$\varepsilon\in(1-\alpha,\bar{\varepsilon}]$ & $\alpha$, $\bar{\varepsilon}$\\
$s\in(0,1)$ & none\\
$e_0>0$ & none\\
$e'\in\left(0,\min\left(e_0,\frac{1}{3}\right)\right]$, in the entry of $e_1$ & $e_0$, and no other letter: the variable $e'$ is bound, and $\frac{1}{3}$ is a numeral. This range is part of the definition of $e_1$, that is, of (Z8.s1), and it restricts no chosen letter.\\
$\bar{m}\in(0,1]$ & none\\
$\beta$, $x_0$ & their row names the condition (Pb), which is listed in Q.3
\end{tabular}

\paragraph{Q.2. The 37 rows of Proposition~14.2\,(b).}
The third column names the letters that occur in the row, the letter restricted by the row being left out.

\noindent
\begin{tabular}{lll}
label & condition on & letters that occur in it\\
\hline
(L1) & $L$ & $\bar{\varepsilon}$\\
(L2) & $L$ & $c_{\mathrm{fl}}$, $\bar{\varepsilon}$\\
(L3) & $L$ & $\kappa$, $\kappa_{0}$\\
(L4) & $L$ & $c_{\mathrm{rb}}$, $\kappa$, $\kappa_{0}$\\
(L5) & $L$ & none\\
(M1) & $M$ & $L$\\
(M2) & $M$ & $L$\\
(M3) & $M$ & $\gamma_{\mathrm{w}}$\\
(M4) & $M$ & $\alpha_{0}$, $\hat{\kappa}$\\
(M5) & $M$ & $M^{F}$\\
(M6) & $M$ & $M^{P}$\\
(S1) & $\varepsilon$ & $\alpha$, $\bar{\varepsilon}$\\
(S2) & $s$ & $\varepsilon$, $a_{L}$, $N_{*}$\\
(A1) & $e_0$ & $p_0$, $\varepsilon$, $c_{\mathrm{fl}}$, $M$\\
(A2) & $e_0$ & $C_{w}$, $c_{\gamma}$, $c_{C}$, $\varepsilon$\\
(A3) & $e_0$ & $c_{X}$, $c_{\gamma}$, $C_{w}$, $C_{T}$, $M$, $\varepsilon$\\
(A4) & $e_0$ & $C_{w}$, $C_{T}$, $\varepsilon$\\
(A5) & $e_0$ & $M$, $\varepsilon$\\
(A6) & $e_0$ & $C_{w}$, $c_{\gamma}$, $\varepsilon$\\
(A7) & $e_0$ & $c_{S}$, $L$, $c_{X}$, $K_{0}$, $\varepsilon$\\
(A8) & $e_0$ & $M$, $\varepsilon$\\
(A9) & $e_0$ & $L$, $M$, $c_{X}$, $\varepsilon$\\
(A10) & $e_0$ & $M$, $\varepsilon$, $C_{B}$, $p$\\
(Z2.s1) & $e_0$ & $L$, $\varepsilon$\\
(Z2.s2) & $e_0$ & $C'_{V}$, $a'$, $\varepsilon$\\
(Z2.s3) & $e_0$ & $M$, $\varepsilon$\\
(Z2.s4) & $e_0$ & $\varepsilon$\\
(Z2.s5) & $e_0$ & $c_{\mathrm{rb}}$, $L$, $C_{d}$, $M$, $\varepsilon$\\
(Z3.s1) & $e_0$ & none\\
(e1) & $e_1$ & $r$, $L$, $p_0$\\
(e2) & $e_1$ & $c_{\mathrm{fl}}$, $\varepsilon$, $L$\\
(e3) & $e_1$ & $C_{0}$, $L$, $M$, $r$, $b_{-}$, $c_{\kappa}$, $c_{\mathrm{fl}}$, $\varepsilon$\\
(e4) & $e_1$ & $C_{0}$, $L$, $M$, $r$, $p$, $\varepsilon$\\
(e5) & $e_1$ & $L$, $M$, $r$, $p_0$, $p$, $C_{0}$\\
(e6) & $e_1$ & $C_{\Sigma}$, $\gamma_{\Sigma}$, $M$, $r$\\
(e7) & $e_1$ & $L$, $M$, $r$\\
(e8) & $e_1$ & $N_{*}$, $M$, $r$
\end{tabular}

In each row the third column lists the letters that occur in the condition stated in the same row of Proposition~14.2\,(b), once $\mathrm{lg}$ is replaced by its definition in (e1) and $\mathrm{lh}(v)$ by its definition in (e3), (e4), (e5), (e7) (Section~0.1);
these are the letters named after ``Letters contained'' in that row of Proposition~14.2\,(b), and the two rows (L5) and (Z3.s1) name no letter.
In (S2) the symbol $\sigma$ is only the name of the left side.
The variable $v$ is bound (preamble of Proposition~14.2\,(b)).

\paragraph{Q.3. The conditions of the line (C10).}
These are the conditions of Proposition~14.2\,(d), together with the conditions in the rows $e_G$; $\beta$, $x_0$; $e_{\sharp}$; $e_2$ of Proposition~14.2\,(a).
In them every quantity that depends on $e$ is replaced by its definition (Proposition~14.2\,(d), second paragraph, that is, Section~G5.1.8; Section~G5.1.4).
The scales $a_{k}$, $e_k$, $v_k$, $r_k$, $p_k$, $p_{0,k}$, $K$, $R^{*}(e)$, $\ell_k$, $k_{*}$, $\ell_*$ are those of Sections~I.6.1 and~I.6.2, formed with $R=R^{*}(e)$, and an index $*$ means the value at $k_{*}$;
$\ell_c$ is the number of Section~I.6.4, and $\mathrm{pp}$ the number of Section~G5.1.4.
Apart from the data $e$ and $N$, the quantities $a_{*}$, $e_{*}$, $v_{*}$, $r_{*}$, $\ell_*$, $r_K$, $R^{*}(e)$, $\ell_c$ contain no letters other than $L$, $M$, $r$, $h$,
and the quantities $p_*$, $p_{0*}$, $\mathrm{pp}$ contain no letters other than $L$, $M$, $r$, $h$ and, respectively, $p$, $p_0$, $D_\beta$.
Below, $\mathrm{sc}$ stands for ``letters among $L$, $M$, $r$, $h$''.
For each condition the inequality is printed with all these definitions inserted, and then the letters that occur in it are named.
\begin{description}
\item[$e_G$.]
\[
e\cdot h^{\frac{1}{2}}\le e_1,\qquad e\cdot h^{\frac{1}{2}}\cdot\log h\le1;
\]
the second inequality is (Z8.w).
Letters: $e_1$, $h$.
\item[(W3).]
\[
6N_{*}\cdot e^2\cdot R^{*}(e)\cdot\log\frac{1}{e}\le\pi^2.
\]
Letters: $N_{*}$, $\mathrm{sc}$.
\item[$e_{\sharp}$, range.]
\[
e'\le e_1\cdot h^{-\frac{1}{2}}.
\]
Letters: $e_1$, $h$.
\item[(g1).]
\begin{gather*}
5\,\mathrm{hL}\cdot h_{G}\le e_{*}^{-\frac{1}{4}+\varepsilon},\qquad\left(4+\frac{M^3}{2}\right)\cdot e_{*}^{\varepsilon}\le1,\\
\mathrm{alG}\cdot\left(4\,h_{G}^2\cdot C_K\cdot e_{*}^{2\varepsilon}+650\,K_{0}\cdot e_{*}^{\frac{1}{4}-7\varepsilon}\right)\le\frac{1}{2}.
\end{gather*}
Letters: $\mathrm{hL}$, $h_{G}$, $\mathrm{alG}$, $C_K$, $K_{0}$, $\varepsilon$, $\mathrm{sc}$.
\item[(g2).]
\[
\frac{M\cdot r_{*}}{8L}-1\ge6M+8.
\]
Letters: $\mathrm{sc}$.
\item[(g3).]
\begin{gather*}
4\,C_{0}\cdot\ell_*^3\cdot\mathrm{gGa}_{1}\le e_{*}^{-\frac{1}{4}+\varepsilon},\qquad4\,\ell_*^3\cdot\left(2C_K+K_{0}+4\right)\le e_{*}^{-2\varepsilon},\\
e_{*}^{\frac{1}{4}-9\varepsilon}\le\frac{1}{8}.
\end{gather*}
Letters: $C_{0}$, $\mathrm{gGa}_{1}$, $C_K$, $K_{0}$, $\varepsilon$, $\mathrm{sc}$.
\item[(g4).]
\begin{gather*}
48\,e_{*}^{2\varepsilon}\le\mu_-,\qquad2\,\mathrm{gGa}_{1}\cdot C_K'\cdot e_{*}^{2\varepsilon}\le\frac{1}{2},\qquad c_{\tau}\cdot\ell_*\ge1,\\
52\,B_1\cdot\exp\left(-c_{\tau}\cdot\ell_*\right)\le e_{*}^{\frac{1}{4}-9\varepsilon},\\
C_{\mathrm{fr}}\cdot e_{*}^{4\varepsilon}+16\exp(6)\cdot C_{d}\cdot M^3\cdot K_{0}\cdot\mathrm{pp}\cdot\exp\left(-r_{*}\right)+B_1\le\frac{1}{4},
\end{gather*}
where $B_1=16\,\mathrm{gGa}_{1}\cdot C_K'\cdot\mathrm{pp}\cdot e_{*}^{2\varepsilon}$.
Letters: $\mu_-$, $\mathrm{gGa}_{1}$, $C_K'$, $c_{\tau}$, $C_{\mathrm{fr}}$, $C_{d}$, $K_{0}$, $D_\beta$, $\varepsilon$, $\mathrm{sc}$.
\item[(g5).]
\begin{multline*}
2\cdot2^{N_{*}}\cdot\mathrm{CWj}\cdot\mathrm{pp}\cdot\exp\left(-\mathrm{vt}\cdot\left(\frac{M\cdot r_{*}}{8L}-4M-7\right)\right)\\
\le\left(6\,e_{*}^{\frac{1}{4}-9\varepsilon}\right)^{N_{*}}.
\end{multline*}
Letters: $N_{*}$, $\mathrm{CWj}$, $\mathrm{vt}$, $D_\beta$, $\varepsilon$, $\mathrm{sc}$.
\item[(g6).]
\[
N_{*}\cdot\left(6\,e_{*}^{\frac{1}{4}-9\varepsilon}\right)^{1-s}\le\frac{1}{2}.
\]
Letters: $N_{*}$, $s$, $\varepsilon$, $\mathrm{sc}$.
\item[(g7).]
\begin{gather*}
\varpi_0\le\frac{1}{4},\qquad\varpi_0=18\,\ell_*^3\cdot\exp\left(-\frac{p_{0*}}{2}\right)+6\,\ell_*^3\cdot\exp\left(-\frac{(p_*-1)^2}{2C_{0}^2}\right);\\
8N_{*}\cdot t_2\le t_x,\\
t_2=e\cdot\min\left\{N_{*}^{-\frac{1}{2}},\ \frac{1}{C_{\beta}},\ \frac{1}{2\sqrt{N_{*}\cdot\mathbf{D}}}\right\},\qquad t_x=\frac{e_{*}^{3\varepsilon}}{\mathrm{Cb}}.
\end{gather*}
Letters: $C_{0}$, $p_0$, $p$, $N_{*}$, $C_{\beta}$, $\mathbf{D}$, $\mathrm{Cb}$, $\varepsilon$, $\mathrm{sc}$.
\item[(g8).]
\begin{gather*}
\left(\varpi_L^{\frac{1}{3}}\cdot\mathrm{mxL}^{\frac{2}{3}}\right)^{\frac{1}{N_{*}}}\cdot5\left(\frac{2N_{*}\cdot\mathbf{r}}{t_2}\right)^4\le\left(6\,e_{*}^{\frac{1}{4}-9\varepsilon}\right)^{s},\\
\varpi_L=\varpi_0^{\frac{1}{2}}\cdot\exp\left(\frac{\mathbf{D}}{2}\right),\qquad\mathrm{mxL}=2+\exp\left(\frac{\mathbf{D}}{2}\right),\qquad\mathbf{r}=\frac{t_x}{4}.
\end{gather*}
Letters: $C_{0}$, $p_0$, $p$, $N_{*}$, $C_{\beta}$, $\mathbf{D}$, $\mathrm{Cb}$, $s$, $\varepsilon$, $\mathrm{sc}$.
\item[(g9).]
\[
\frac{D_\beta}{2}+\frac{1}{2}\le\frac{3}{8}\,\ell_c,\qquad D_\beta+3\le\ell_c.
\]
Letters: $D_\beta$, $\mathrm{sc}$.
\item[(Pb).]
The three requirements $b(f)>0$, $d_0>0$ and
\begin{multline*}
\frac{6}{\bar{m}^2}\cdot\|a\|_{\infty}^4\cdot\left(\frac{4\pi\cdot r_\beta^3}{3}\right)^2\cdot\left(\bar{m}^2+\left(\bar{m}+\frac{1}{d_0}\right)^2\right)\\
\cdot\frac{\exp\left(-2\bar{m}\cdot d_0\right)}{4\pi^2\cdot d_0^2}\le\frac{b(f)}{12},
\end{multline*}
the last being the inequality displayed in \cite[Section~2.8]{QEDI}; here $f=dd^{*}\beta$, $a=d^{*}\beta$, $\|a\|_{\infty}=\sup_y|a(y)|_2$ and $d_0=|x_0|_2-2r_\beta$.
Letters: $\beta$, $x_0$ (the letters restricted), and $\bar{m}$, which occurs in the displayed inequality and, by Section~I.7.6, in $b(f)$.
Besides these the radius $r_\beta$ occurs: it is a datum chosen together with $(\beta,x_0)$, and not a letter of the list.
\item[(Z10.s1).]
\[
e_2\le e_G.
\]
Letters: $e_G$.
\item[(Z10.s2).]
\[
e_2\le e_{\sharp}(\beta')\quad\text{for the three probes}\quad\beta'=\beta,\ \beta+\tau_{x_0}\beta,\ \beta-\tau_{x_0}\beta.
\]
Letters: $e_{\sharp}$, and further $\beta$ and $x_0$, through the three probes.
\item[(Z10.s4).]
\[
\frac{5}{2}\,C'_{\bar{\kappa}}\cdot e_2^{\sigma}\le\frac{b(f)}{4}.
\]
Letters: $C'_{\bar{\kappa}}$, $\sigma$; further $\beta$ and $\bar{m}$, through $b(f)$ (Section~I.7.6).
\end{description}

\paragraph{Q.4. The displays of Proposition~14.2\,(c).}
They are conditions on $(e,T,N)$ (Section~0.6).
Besides letters of $\mathrm{sc}$, the letter $\varepsilon$ and the data, they contain the following letters:
\begin{itemize}
\item[] (2.1): none;\quad (fl): $c_{\mathrm{fl}}$;\quad (T1): $p_0$;\quad (T2): $C_{0}$, $b_{-}$, $c_{\kappa}$, $c_{\mathrm{fl}}$, and $\bar{m}$ through $\mu_*=\bar{m}\cdot a_{*}$;
\item[] (T3): $C_{0}$, $p$;\quad (T4): $C_{0}$, $p_0$, $p$;\quad (T5): $C_{\Sigma}$, $\gamma_{\Sigma}$;\quad (T6): none;
\item[] error terms: $\varepsilon_2$: $C_{\Sigma}$, $\gamma_{\Sigma}$;\quad $\varepsilon_3$: $C_{0}$, $p_0$, $p$;\quad $\varepsilon_4$: $C_{0}$, $p$.
\end{itemize}

\emph{Labels of Section~I.8.4 that have no row in Q.1--Q.4} (Proposition~14.2, the paragraph that follows (b)).
The label (A0) of Section~I.8.4 stands for (L3), the definition of $\kappa_{\mathrm{w}}$, (M1), (M2), the inequality $\gamma_{\mathrm{w}}\cdot M\ge1$ (a consequence of (M3)) and (M4);
(Z2.m) is (M5);
(Z8.s1) is the definition of $e_1$, whose conditions are the range $e'\in\left(0,\min\left(e_0,\frac{1}{3}\right)\right]$ of Q.1, which contains the letter $e_0$ only, and (e1)--(e8);
(Z9.s1), (Z9.s2), (Z9.s3) are the second inequality of (g1) and the first inequalities of (g3) and of (g4).
Hence every condition that stands under one of these labels is listed in Q.1--Q.4 or follows from a condition listed there:
the inequality $\gamma_{\mathrm{w}}\cdot M\ge1$ follows from (M3), which reads $\gamma_{\mathrm{w}}\cdot M\ge16$, and (Z9.s1)--(Z9.s3) are parts of (g1), (g3), (g4).
The equation $\kappa_{\mathrm{w}}=\kappa+\kappa_{0}+3$ is the definition of a letter (row~53); it contains the letters $\kappa$ and $\kappa_{0}$, and it is no condition.

\subsubsection*{5.2. First sentence: how $\rho$ enters $M$}

By Q.2 the conditions on $M$ are (M1)--(M6), and
\[
\mathrm{pa}({M})=\left\{L,\ \gamma_{\mathrm{w}},\ \alpha_{0},\ \hat{\kappa},\ M^{F},\ M^{P}\right\}.
\]
Of these letters, $L$ occurs in (M1) and (M2);
$\gamma_{\mathrm{w}}$ occurs in (M3) and in no other condition on $M$;
$\alpha_{0}$ and $\hat{\kappa}$ occur in (M4);
$M^{F}$ occurs in (M5) and in no other;
$M^{P}$ occurs in (M6).
The letters $L$, $\alpha_{0}$, $\hat{\kappa}$, $M^{P}$ belong to $\mathrm{NR}$; hence $\rho$ is none of them and is an ancestor of none of them (Section~4.1).
Let $\rho=Y_0,\dots,Y_m=M$ be a chain (Section~0.6).
Then $Y_{m-1}$ is an element of $\mathrm{pa}({M})$ which either is $\rho$ or has $\rho$ among its ancestors.
Consequently $Y_{m-1}$ is $\gamma_{\mathrm{w}}$ or $M^{F}$.
This is the first assertion: among the conditions on $M$, the letter $Y_{m-1}$ occurs in (M3) alone, or in (M5) alone.

One has $\mathrm{pa}(\gamma_{\mathrm{w}})=\{\gamma^{F},\gamma^{P}\}$, and $\gamma^{P}$ belongs to $\mathrm{NR}$.
Hence, if $Y_{m-1}=\gamma_{\mathrm{w}}$, then $Y_{m-2}=\gamma^{F}$, since $Y_{m-2}$ is $\rho$ or has $\rho$ among its ancestors.
Likewise $\mathrm{pa}(M^{F})=\{L,\ C'_{V},\ a',\ \kappa'\}$, and $L$, $C'_{V}$, $\kappa'$ belong to $\mathrm{NR}$; hence $Y_{m-1}=M^{F}$ forces $Y_{m-2}=a'$.
It follows that every chain contains $\gamma^{F}$ or $a'$.
This is the second assertion.

Chains exist, for instance
\begin{gather*}
\rho\to c_{\mathrm{re}}\to a_{z}\to\gamma_{z}\to\gamma^{F}\to\gamma_{\mathrm{w}}\to M,\\
\rho\to c_{\mathrm{re}}\to C_{00}\to C_z\to a'\to M^{F}\to M,
\end{gather*}
where in each chain every letter is a parent of the next one (Table P, rows 67, 68, 70--76, 78).

The conditions (M3) and (M5) are lower bounds.
Indeed, since $\gamma_{\mathrm{w}}>0$ by (6.3), the condition (M3) reads
\[
M\ge\frac{16}{\gamma_{\mathrm{w}}}=\max\left(\frac{16}{\gamma^{F}},\frac{16}{\gamma^{P}}\right),
\]
and (M5) reads
\[
M\ge M^{F}=\max\left\{4L+1,\ \left(108\pi\cdot C'_{V}\cdot a'\right)^2,\ 108\pi\cdot C'_{V}\cdot a'\cdot\exp\left(27\kappa'\right)\right\}.
\]

\subsubsection*{5.3. Second sentence: the conditions on $M$}

\emph{No condition on $M$ contains a letter of the lines (C6)--(C10).}
By Q.2 the letters occurring in (M1)--(M6), other than the restricted letter $M$ itself (line (C5)), are $L$ (line (C3)), $\alpha_{0}$ and $\hat{\kappa}$ (line (C4)), and $\gamma_{\mathrm{w}}$, $M^{F}$, $M^{P}$ (line (C4r)).
By (Par1) the ancestors of these letters stand before them.
Therefore, $M$ itself excepted, no letter of the lines (C5)--(C10) and no letter having such a letter among its ancestors occurs in a condition on $M$.

\emph{No condition on $M$ bounds $M$ from above.}
Let $M$ satisfy (M1)--(M6), and let $M'=L^jM$ with an integer $j\ge0$.
Then $M'$ satisfies (M1)--(M6):
\begin{itemize}
\item[(M1)] $M'$ is a power of $L$, and $M'\ge M\ge L^2$;
\item[(M2)] $M'\ge M>4L$;
\item[(M3)] $\gamma_{\mathrm{w}}\cdot M'\ge\gamma_{\mathrm{w}}\cdot M\ge16$, because $\gamma_{\mathrm{w}}>0$;
\item[(M4)] $M'^{\alpha_{0}}\ge M^{\alpha_{0}}\ge1+\exp(\hat{\kappa})$, because $\alpha_{0}=\frac{1}{54}>0$ and $M'\ge M\ge1$;
\item[(M5), (M6)] $M'\ge M\ge M^{F}$ and $M'\ge M\ge M^{P}$.
\end{itemize}
So $\mathrm{Val}({M})$ contains, together with $M$, every larger power of $L$.

The letter $M$ occurs in further conditions, namely in (A1), (A3), (A5), (A8), (A9), (A10), (Z2.s3), (Z2.s5), in (e3)--(e8), and in conditions of Q.3, for instance in (g2).
These are conditions on $e_0$, conditions on $e_1$, and conditions of the line (C10) (column ``condition on'' of Q.2; Q.3).
They restrict letters that stand after $M$, for the given value of $M$.
The condition (Pb) restricts $\beta$ and $x_0$; it contains neither $M$ nor $\rho$ (Q.3), and it is no condition on $M$.

\subsubsection*{5.4. Third sentence: nine letters}

By the last columns of Q.1 and Q.2, the lists of letters that close the entries of Q.3, and the definitions of the letters in Proposition~14.2\,(a), the following table names, for each of the nine letters, all conditions of Q.1--Q.3 in which the letter occurs, and all definitions of letters of the list in which it occurs, the definitions of $e_0$, $e_1$ and $e_{\sharp}$ excepted.
In these three definitions the nine letters occur through the conditions named in the second column (Sections~0.3 and~0.6):
$C_{w}$, $C_{T}$, $c_{\gamma}$, $C_{B}$, $C_{d}$ occur in the definition of $e_0$;
$C_{0}$, $\gamma_{\Sigma}$, $C_{\Sigma}$, $c_{\kappa}$ in that of $e_1$;
and $C_{d}$, $C_{0}$ in that of $e_{\sharp}$.

\noindent
\begin{tabular}{p{0.15\textwidth}p{0.47\textwidth}p{0.28\textwidth}}
letter (line) & occurs in the conditions & occurs in the definition of\\
\hline
$C_{w}$ (C6) & on $e_0$: (A2); (A3), first inequality; (A4); (A6) & $C_{W}$, $C_{d}$; the tuple $\mathrm{par}$\\
$C_{T}$ (C4) & on $e_0$: (A3), first inequality; (A4) & none\\
$c_{\gamma}$ (C6) & on $e_0$: (A2); (A3), first inequality; (A6) & $C_{W}$\\
$C_{B}$ (C6) & on $e_0$: (A10), third inequality & none\\
$C_{d}$ (C6) & on $e_0$: (Z2.s5); of the line (C10): (g4), fifth inequality (through $B_{\mathrm{wr}}$) & $\mathrm{CWj}$ (line (C6$'$)), $C_{\mathrm{fr}}$ (line (C10))\\
$C_{0}$ (C6) & on $e_1$: (e3), first and second inequality; (e4); (e5); of the line (C10): (g3), first inequality; (g7), first inequality, and (g8), both through $\varpi_0$ & none\\
$\gamma_{\Sigma}$ (C6) & on $e_1$: (e6) & none\\
$C_{\Sigma}$ (C6) & on $e_1$: (e6) & none\\
$c_{\kappa}$ (C6) & on $e_1$: (e3), second inequality & none
\end{tabular}

Consequently $C_{w}$, $C_{T}$, $c_{\gamma}$, $C_{B}$ occur in conditions on $e_0$ only;
$C_{d}$ occurs in a condition on $e_0$ and in a condition of the line (C10), and in no other condition;
$C_{0}$ occurs in conditions on $e_1$ and in conditions of the line (C10), and in no others;
$\gamma_{\Sigma}$, $C_{\Sigma}$, $c_{\kappa}$ occur in conditions on $e_1$, and in no others.
In particular, none of the nine letters occurs in a condition on $\alpha$, $\kappa$, $L$, $M$, $r$, $p_0$, $p$, $\varepsilon$ or $s$, nor in the range of $\bar{m}$, nor in (Pb).
By the paragraph \emph{Words} of Proposition~14.4, the displays (T2)--(T5) of Proposition~14.2\,(c) are counted among the conditions on $e_1$.
By Q.4, of the nine letters these displays contain $C_{0}$ (in (T2), (T3), (T4)), $c_{\kappa}$ (in (T2)), $C_{\Sigma}$ and $\gamma_{\Sigma}$ (in (T5)), and no other;
and these are precisely the four letters for which the third sentence admits conditions on $e_1$.
This proves the third sentence of Proposition~14.4\,(c).

Two complements follow.
\begin{itemize}
\item[(1)] The displays of Proposition~14.2\,(c) that contain one of the nine letters are (T2)--(T5) and the error terms $\varepsilon_2$, $\varepsilon_3$, $\varepsilon_4$ (Q.4); the letters among the nine which they contain are $C_{0}$, $c_{\kappa}$, $C_{\Sigma}$, $\gamma_{\Sigma}$.
The displays (T2)--(T5) hold for all data of Proposition~14.2\,(c), as consequences of (e3)--(e6), the conditions on $e_1$ that contain the same letters: (T2) follows from (e3), (T3) from (e4), (T4) from (e5), and (T5) from (e6) (Proposition~14.2\,(c), the sentence which says that these displays are no conditions on constants; Section~I.8.2, the paragraph after the table).
The error terms are definitions of quantities that depend on $e$.
\item[(2)] Some of the nine letters occur in inequalities of Proposition~14.2\,(f):
in six inequalities of the line that begins with $2C_{1}\le C^{F}$;
in the line ``Ranges'' (positivity; $C_{\Sigma}\ge0$; $\gamma_{\Sigma}\le1$; the upper bound of $c_{\gamma}$; $c_{\kappa}\ge1$);
in $C_{w}\ge1$;
and, through the sentence of (f) on the line (C6$'$), in the inequality of Annex~B, Section~G5.9.3, that contains $C_{d}$.
These inequalities are no conditions and restrict no chosen letter, since (f) is asserted for every choice that (a) and (b) admit (preamble of Proposition~14.2\,(a)).
The six inequalities of the first of these lines follow from the definitions of Proposition~14.2\,(a), as follows.
\begin{itemize}
\item $4\le C^{F}\le C_{w}^{(0)}\le C_{w}^{(1)}\le C_{w}$.
One has $C^{F}\ge2C_{\mathrm{By}}\ge2\sqrt{2}\,(1+2\mathrm{e})>4$, by the definitions of $C^{F}$ and $C_{\mathrm{By}}$ in Proposition~14.2\,(a) (rows 82 and 79);
$C_{w}^{(0)}=\max\left(C^{F},C^{P}\right)\ge C^{F}$;
$C_{w}^{(1)}=\max\left(C_{w}^{(0)},\dots\right)\ge C_{w}^{(0)}$;
and $C_{w}=64K_{0}\cdot C_{w}^{(1)}\ge C_{w}^{(1)}$, because $K_{0}=\frac{72^{54}}{36}\ge1$ and $C_{w}^{(1)}>0$.
\item $C_{P}\le C^{P}\le C_{w}$.
One has $C^{P}=116\,C_{P}\cdot M^{402\alpha_{0}}\ge C_{P}$, because $C_{P}\ge4\exp(2)>0$ (Section~1.2), $M\ge L^2\ge1$ by (M1), and $\alpha_{0}>0$;
and $C^{P}\le C_{w}^{(0)}\le C_{w}$ by the preceding inequality chain.
\item $C_{0}\ge\max\left(1,C^{F}\right)$.
One has $C_{0}\ge C^{F}$, because $C^{F}$ is an entry of the maximum that defines $C_{0}$ in Proposition~14.2\,(a) (row 92), and $C^{F}>4$.
\item $c_{\kappa}\le K_{0}$.
This is Lemma~1.
\item $C_{T}\ge1$.
One has
\[
C_{T}=2\mathrm{e}^{\frac{3}{2}}\cdot\left[b_{-}^{-1}\cdot\left(54L\cdot\mathrm{e}^6+12\mathrm{e}+1\right)+1\right]\ge2\mathrm{e}^{\frac{3}{2}}\ge1,
\]
because $b_{-}=b\cdot\left(1-\frac{1}{L}\right)>0$.
\item $c_{\gamma}\ge1$.
Let $k\ge1$, let $t=(y,\beta)$ with a unit site $y$, and take the index set (i).
The four terms $t'=(y,\beta')$, $\beta'=1,\dots,4$, have $\nu(t')=1$ and $d(t,t')=0$.
Hence $c_{\gamma}\ge4$.
\end{itemize}
\end{itemize}

\begin{named}{Remark}
In the third sentence of Proposition~14.4\,(c), occurrence cannot be replaced by ancestry.
Indeed, $C_{w}$ is a parent of $C_{d}$, which occurs in (g4).
Moreover, every letter that occurs in a condition on $e_0$ is a parent of $e_0$ (Section~1.3, kind C), hence an ancestor of $e_1$, $e_G$, $e_{\sharp}$ and $e_2$: for $e_0$ is a parent of $e_1$ (row 118 of Table P), $e_1$ is a parent of $e_G$ and of $e_{\sharp}$, and these two are parents of $e_2$.
All these relations follow the order in which the letters are chosen, since every parent stands before its letter ((Par1); Section~0.3); in particular, the parent relation has no cycle.
\end{named}

\subsection*{6. Proof of (d)}

Proposition~14.4\,(d) is proved in the following form.

\begin{named}{Lemma 3}
Let $\rho\in(0,1)$, $\alpha\in(1-\bar{\varepsilon},1)$ and $\kappa\ge7\kappa_{0}+7$, and let $L$ be an odd integer that satisfies (L1)--(L5).
Let $\lambda_{\mathrm{int}}$, $c_{*}$, $C_{\mathrm{Po}}$, $k_{1}$, $b$ be the numbers of Section~4.2.
Then
\begin{align}
&0<c_{b}\le\frac{\left(\lambda_{\mathrm{int}}\cdot c_{*}\right)^2\cdot\rho^4}{3456},\tag{6.1}\\
&0<c_{\mathrm{re}}\le\frac{\rho}{2C_{\mathrm{Po}}\cdot L},\tag{6.2}\\
&0<\gamma_{\mathrm{w}}\le\gamma^{F}\le\frac{\rho}{8\mathrm{e}\cdot(6+3b)\cdot C_{\mathrm{Po}}\cdot L^2},\tag{6.3}\\
&a'\ge C_{00}\ge\frac{192\,(1+2k_{1})\cdot C_{\mathrm{Po}}\cdot L^2}{\rho},\tag{6.4}\\
&M_{\mathrm{min}}\ge\frac{16}{\gamma_{\mathrm{w}}}\ge\frac{128\mathrm{e}\cdot(6+3b)\cdot C_{\mathrm{Po}}\cdot L^2}{\rho},\tag{6.5}
\end{align}
where $M_{\mathrm{min}}$ is the number defined in Proposition~14.3\,(c); and
\begin{equation}
\begin{aligned}
&c_{\mathrm{re}}\le\frac{\rho}{10\,C_{\mathrm{Po}}},\qquad\gamma_{\mathrm{w}}\le\gamma^{F}\le\frac{\rho}{200\mathrm{e}\cdot(6+3b)\cdot C_{\mathrm{Po}}},\\
&a'\ge\frac{4800\,(1+2k_{1})\cdot C_{\mathrm{Po}}}{\rho},\qquad M_{\mathrm{min}}\ge\frac{3200\mathrm{e}\cdot(6+3b)\cdot C_{\mathrm{Po}}}{\rho}.
\end{aligned}
\tag{6.6}
\end{equation}
\end{named}

Apart from $\rho$, the right sides of (6.1) and (6.6) contain numbers only.
Hence, as $\rho\to0$, the quantities $c_{b}$, $c_{\mathrm{re}}$, $\gamma^{F}$ (and with it $\gamma_{\mathrm{w}}$) tend to $0$, and $a'$, $M_{\mathrm{min}}$ tend to infinity, uniformly in the triples $(\alpha,\kappa,L)$ that satisfy the hypotheses above.
This holds when $(\alpha,\kappa,L)$ is kept fixed, which Section~4.2 permits, and equally when the triple is chosen in dependence on $\rho$.

Every $M$ that satisfies (M1)--(M6) satisfies $M\ge M_{\mathrm{min}}$, that is, $M$ is at least each of the six entries of the maximum in Proposition~14.3\,(c):
\begin{itemize}
\item $M\ge L^2$ by (M1);
\item $M\ge4L+1$ by (M2), since $M$ and $4L$ are integers;
\item $M\ge\frac{16}{\gamma_{\mathrm{w}}}$ by (M3), since $\gamma_{\mathrm{w}}>0$;
\item $M\ge\left(1+\exp(\hat{\kappa})\right)^{54}$ by (M4), since $\alpha_{0}=\frac{1}{54}$;
\item $M\ge M^{F}$ by (M5);
\item $M\ge M^{P}$ by (M6).
\end{itemize}
Therefore, by (6.6), no value of $M$ is admissible for all $\rho\in(0,1)$.

\begin{proof}[Proof of Lemma 3]
One has $\lambda_{\mathrm{int}}>0$ (Section~1.3) and $w_3=6$.
Hence $c_{*}=\frac{\min(w_3,\lambda_{\mathrm{int}})}{18}>0$ and
\[
0<C_{\mathrm{Po}}=\frac{1}{8c_{*}}+\frac{1}{\lambda_{\mathrm{int}}}<\infty.
\]
Moreover $b=1$, $L\ge5$ and $0<\rho<1$.
The formulas used below are those of the rows 12--14, 16, 22, 23, 35, 51, 59--62 and 66--75 of Proposition~14.2\,(a).

\emph{Proof of (6.1).}
One has $\hat{g}=\frac{\rho\cdot\lambda_{\mathrm{int}}}{2}>0$ and $c_{\mathrm{hl}}^2=3(1+\rho)^2\ge3$.
Hence
\[
K_{\rho}=\rho+\frac{2c_{\mathrm{hl}}^2}{\hat{g}}\ge\frac{6}{\hat{g}}=\frac{12}{\rho\cdot\lambda_{\mathrm{int}}}.
\]
Since $1+\frac{6(1+\rho)}{\hat{g}}\ge1$, it follows that
\begin{align*}
A_1&=\frac{2\sqrt{2}\,K_{\rho}}{\rho\cdot c_{*}}\cdot\left(1+\frac{6(1+\rho)}{\hat{g}}\right)\ge\frac{2\sqrt{2}\cdot12}{\rho^2\cdot\lambda_{\mathrm{int}}\cdot c_{*}}=\frac{24\sqrt{2}}{\rho^2\cdot\lambda_{\mathrm{int}}\cdot c_{*}},\\
A_1^2&\ge\frac{1152}{\rho^4\cdot\lambda_{\mathrm{int}}^2\cdot c_{*}^2}.
\end{align*}
The numbers $A_1$ and $A_2$ are finite, and $\left(30+24\sqrt{2}\right)A_2^2\ge0$.
Therefore
\[
0<c_{b}=\left[3A_1^2+\left(30+24\sqrt{2}\right)A_2^2\right]^{-1}\le\frac{1}{3A_1^2}\le\frac{\rho^4\cdot\lambda_{\mathrm{int}}^2\cdot c_{*}^2}{3456}.
\]

\emph{Proof of (6.2).}
By definition
\[
c_{\mathrm{re}}=\min\left\{b\cdot\left(1-\frac{1}{L}\right),\ \frac{\rho}{2C_{\mathrm{Po}}\cdot L}\right\};
\]
both entries are positive, and the minimum does not exceed the second entry.

\emph{Proof of (6.3).}
By (6.2),
\begin{align*}
a_{z}=\min\left\{\frac{1}{3L},\ \frac{c_{\mathrm{re}}}{2\mathrm{e}\cdot L\cdot(6+3b)}\right\}&\le\frac{c_{\mathrm{re}}}{2\mathrm{e}\cdot L\cdot(6+3b)}\\
&\le\frac{\rho}{4\mathrm{e}\cdot(6+3b)\cdot C_{\mathrm{Po}}\cdot L^2}.
\end{align*}
Further $\gamma_{z}=\min\left(\gamma_{(5)},a_{z}\right)\le a_{z}$, so that
\[
\gamma^{F}=\frac{\gamma_{z}}{2}\le\frac{\rho}{8\mathrm{e}\cdot(6+3b)\cdot C_{\mathrm{Po}}\cdot L^2},
\]
and $\gamma_{\mathrm{w}}=\min\left(\gamma^{F},\gamma^{P}\right)\le\gamma^{F}$.
It remains to show positivity.
One has $a_{z}>0$ by (6.2).
By (6.1), $c_{b}>0$; moreover $b_{-}=b\cdot\left(1-\frac{1}{L}\right)>0$.
Hence the numbers
\begin{gather*}
C_{I}=\frac{1}{b_{-}}+\frac{2C_{\mathrm{Po}}\cdot K_{1}^2}{c_{b}}+\left(\frac{2C_{\mathrm{Po}}}{c_{b}}\right)^{\frac{1}{2}},\\
c_{\mathrm{wt}}=\mathrm{e}\cdot\left(3(1+\rho)+6L^4\cdot b\right)
\end{gather*}
are finite and positive, $a_{0}=\min\left\{\frac{1}{3},\frac{1}{2C_{I}\cdot c_{\mathrm{wt}}}\right\}>0$, and $\gamma_{(5)}=\frac{a_{0}}{L}>0$.
Finally $\gamma^{P}=\frac{\gamma_{P}}{2}>0$, because $\gamma_{P}$ lies in $(0,1]$ \cite[Section~15.7]{QEDI}.
So $\gamma_{z}$, $\gamma^{F}$ and $\gamma_{\mathrm{w}}$ are positive.

\emph{Proof of (6.4).}
By (6.2),
\begin{align*}
C_{00}=48\cdot\frac{2L}{c_{\mathrm{re}}}\cdot(1+2k_{1})&\ge48\cdot2L\cdot\frac{2C_{\mathrm{Po}}\cdot L}{\rho}\cdot(1+2k_{1})\\
&=\frac{192\,(1+2k_{1})\cdot C_{\mathrm{Po}}\cdot L^2}{\rho}.
\end{align*}
Whatever the value of $C_{5}$ is, $C_z=\max\left(C_{5},C_{00}\right)\ge C_{00}$.
Moreover
\[
c(\gamma^{F})=\sum_{n\in\mathbb{Z}^3}\exp\left(-\gamma^{F}\cdot|n|\right)\ge1,
\]
the term $n=0$ being $1$.
Hence $a'=C_z\cdot c(\gamma^{F})\ge C_{00}$.

\emph{Proof of (6.5).}
The number $\frac{16}{\gamma_{\mathrm{w}}}$ is an entry of the maximum that defines $M_{\mathrm{min}}$; hence $M_{\mathrm{min}}\ge\frac{16}{\gamma_{\mathrm{w}}}$.
By (6.3),
\[
\frac{16}{\gamma_{\mathrm{w}}}\ge\frac{16}{\gamma^{F}}\ge\frac{128\mathrm{e}\cdot(6+3b)\cdot C_{\mathrm{Po}}\cdot L^2}{\rho}.
\]

\emph{Proof of (6.6).}
By (L5), $L\ge5$.
Hence $2L\ge10$ in (6.2), and $L^2\ge25$ in (6.3), (6.4), (6.5); and $8\cdot25=200$, $192\cdot25=4800$, $128\cdot25=3200$.
\end{proof}

\begin{named}{Remark}
The letters of kind O enter Lemma~3 only through the inequality $C_z\ge C_{00}$.
Hence Lemma~3 holds for every rule in the sense of (Par1).
Moreover $M_{\mathrm{min}}\ge M^{F}\ge\left(108\pi\cdot C'_{V}\cdot a'\right)^2$, and the right side tends to infinity as $\rho\to0$, by (6.6).
\end{named}

\subsection*{7. Conditions on the constants}

The proofs of this chapter use the following conditions on the constants.
Each of them is a range or a row of Proposition~14.2\,(a), (b), and no other condition is used.
For each condition the table gives the consequences through which it enters the proofs, the line of Section~I.8.2 in which it occurs or its label in Proposition~14.2, and the places where it is used.

\noindent
\begin{tabular}{p{0.26\textwidth}p{0.25\textwidth}p{0.12\textwidth}p{0.21\textwidth}}
condition & enters the proofs through & line / label & used in\\
\hline
$0<\rho<1$ & $\rho>0$; $1+\rho\ge1$ & (C0) & 1.3 (kinds N, E); Lemma~3\\
$b=1$ & $b>0$ & (C0) & 1.3 (kinds N, E); Lemma~3; 5.4\,(2)\\
$w=\left(6,6\sqrt{2},6\sqrt{2},6\right)$ & $w_i>0$; $w_3=6$ & (C0) & 1.3 ($\lambda_{\mathrm{int}}>0$); Lemma~3\\
$\kappa_{0}=54\log72$, $K_{0}=\frac{\exp(\kappa_{0})}{36}$ & $64K_{0}\ge1$ & (C0) & 5.4\,(2)\\
$\alpha\in(1-\bar{\varepsilon},1)$, $\bar{\varepsilon}=\frac{1}{200}$ & $0\le\alpha\le1$; $\alpha<1$ & (C1), (C0) & 1.3, $C_{G}$; 1.3 (kinds N, E)\\
$\kappa\ge7\kappa_{0}+7$ & $\kappa\ge\kappa_{0}$ & (C2) & Lemma~1 ($c_{\kappa}\le K_{0}$)\\
$L$ odd, $L\ge5$ & $L$ odd, $\frac{L-1}{2}\ge\frac{3}{2}$ and $3L^2\ge4$, in Lemma~2; $L>1$, in $b\cdot\left(1-\frac{1}{L}\right)>0$; $L\ge5$, in (6.6) & (C3), (L5) & 1.3 (kinds N, E); Lemma~2; Lemma~3; 5.4\,(2)\\
$\alpha_{0}=\frac{1}{54}$ & $\alpha_{0}>0$; $\frac{1}{\alpha_{0}}=54$ & (C4) & 5.3; 5.4\,(2); Lemma~3 ($M\ge M_{\mathrm{min}}$)\\
$M$ a power of $L$, $M\ge L^2$ & $M\ge1$; $M$ an integer & (C5), (M1) & 1.3 (kinds N, E); 5.3; 5.4\,(2); Lemma~3 ($M\ge M_{\mathrm{min}}$)\\
$\gamma_{\mathrm{w}}=\min\left(\gamma^{F},\gamma^{P}\right)$ & $\gamma_{\mathrm{w}}>0$, which follows from the definitions of the rows up to 75, by (6.3) & (C4r) & Lemma~2; 5.2, 5.3\\
(M1)--(M6) & the description of $\mathrm{Val}({M})$ & (C5) & 5.2, 5.3; Lemma~3 ($M\ge M_{\mathrm{min}}$)
\end{tabular}

\section*{Chapter 18a. Scaling of thresholds and norms, global fields on a small polymer, stability of locality: proofs of Propositions~18.1--18.4}

\subsection*{0. Conventions of this chapter; facts used in several of the four proofs}

Throughout the chapter $L$ is an odd integer with $L\ge 5$, and $M$ is a power of $L$ with $M\ge L^2$; in particular $M$ is odd.
The spacing of level $k$ is $\eta_k=L^{-k}\le 1$.
Hence $\eta_{k+1}=\eta_k/L$, and $1=L^k\cdot\eta_k$ is a multiple of $\eta_k$.

\subsubsection*{0.0. Labels and letters}

The settings (G2.S-1) and (G2.S-2) are those of Section~G2.0.4, and (G2.r2), (G2.r4), (G2.r5), (G2.r8) are the definitions recalled in Section~G2.0.5.
The clauses (I.k1), (I.k2), (I.k4) are those of Section~I.5.4.
Further, (I.AP) is the antiperiodicity of Section~I.3.6, (I.POS) is the rule of Section~I.1.4 on where positions are read, and (I.T-G) is the definition of the automorphism $G_{\Lambda}$ in Section~I.4.5.
The labels (C0)--(C8) denote the lines of the list of constants in Section~I.8.2, and (L1), (L5), (M1), (S1), (A8), (A9) denote rows of Proposition~14.2\,(b).
The enlarged cubes of Section~I.1.1, the field domains of Section~I.5.3 and the measure $\hat{\mu}$ of Section~I.5.7 are defined as in \cite[Definition~5.7]{QEDI}, \cite[Definition~5.16]{QEDI} and \cite[Section~5.2.5]{QEDI}, respectively.
Without a prefix, the labels (f1), (f2), (f3) refer to the three properties of the cut-off functions of \cite[Section~13.6]{QEDI}; they are recalled in (qu1) below.
The steps $0$--$7$ of the proof of Proposition~18.2 bear the numbers of the corresponding steps of the proof of \cite[Theorem~13.11]{QEDI}.

In this chapter the following notation is used besides the global one.
The symbols $\mathscr{A}(z)$ and $\mathscr{C}(z)$ denote the first and the second slot of a pair of level $k$ that depends on the point $z$ of Step~7; in Step~7 the letter $z=(z_1,\dots,z_4)$ denotes a point of $\mathbb{C}^4$ and not a site of $T^0_k$.
$K'_{nm}$ is a kernel, the kernels $K_n$ being those of Proposition~15\,(a).
$\mathrm{kc}_j^{\mathrm{R}}$ and $\mathrm{kc}_j^{\mathrm{I}}$ are the real and the imaginary part of $\mathrm{kc}_j$.
For an argument $\mathrm{ag}$, that is, for one of the $\xi_i$ or one of the $\boldsymbol{\zeta}_j$, the site $c(\mathrm{ag})$ is the photon site of $\mathrm{ag}$ defined in Section~2.4.
The letter $\mathrm{sn}$ denotes a sign, as in Proposition~18.4\,(a) and in Proposition~15\,(e).
The letter $J''$ denotes a subset of $\mathrm{Mo}\times\{0,1\}$, as in Proposition~15\,(D2), (D3), (e); it appears in (b4) only.
Moreover, $r$ is a radius, as in (G2.r2);
$v$ is a vector of $\mathbb{L}_k$ or of $\mathbb{R}^3$, as in Sections~I.1.3 and~I.1.7;
$n$ is a vector of $\mathbb{Z}^3$, as in Proposition~16.1\,(a1), or the number of the legs $\psi$, as in Section~I.5.1;
$b$ is a bond, as in Section~I.3.7 and in (I.k2);
$f$ is a function of the sites, as in Section~I.1.1, or the argument of $\oint_{\mathrm{rad}}$, as in Section~I.5.2;
$d$ is a value of $d_X$, and $Y$ is a polymer of level $k+1$, as in Proposition~18.3\,(e);
$G$ is an element of an algebra, or a map, beside $F$, as in Proposition~15\,(b).
The gauge function of Step~0 is $\boldsymbol{\lambda}_L$, and the direction indices are $\mu$, $\nu$.

\subsubsection*{0.1. Division by $L$ on the photon lattices}

Let an integer $k\ge 0$ and a side $\ell\in\mathrm{Cl}$ be given.
Then $\ell/L$ is a multiple of $M$, and $T_{k+1}=\mathbb{L}_{k+1}/(\ell/L)\mathbb{Z}^3$.
\begin{itemize}
\item[(sc1)] One has $\mathbb{L}_{k+1}=\eta_{k+1}\mathbb{Z}^3=L^{-1}\mathbb{L}_k$.
Hence $c\mapsto c/L$ is a bijection of $\mathbb{L}_k$ onto $\mathbb{L}_{k+1}$, with $|c/L-c'/L|=|c-c'|/L$ both for the sup norm $|\cdot|$ and for the norm $|\cdot|_1$, and with $(c+\ell\cdot n)/L=c/L+(\ell/L)\cdot n$ for $n\in\mathbb{Z}^3$.
It follows that the map induces a bijection of $T_k$ onto $T_{k+1}$, and that
\[
d(c/L,c'/L)=\min_{n\in\mathbb{Z}^3}\left|c/L-c'/L+(\ell/L)\cdot n\right|=\frac{d(c,c')}{L}.
\]
For a set $Z'$ of sites of $T_{k+1}$ put $LZ':=\{c\in T_k:\ c/L\in Z'\}$.
\item[(sc2)] For a bond $b=(c,c')$ of $T_k$ put $b/L:=(c/L,c'/L)$.
The bond $(c,c+\eta_k\cdot e_{\mu})$ is sent to the bond $(c/L,c/L+\eta_{k+1}\cdot e_{\mu})$ of $T_{k+1}$, and its reverse $(c+\eta_k\cdot e_{\mu},c)$ to the reverse of the latter.
Thus $b\mapsto b/L$ is a bijection between the bonds of $T_k$ and those of $T_{k+1}$ which preserves the direction and the orientation.
A bond $b$ has an end point (both end points) in $LZ'$ if and only if $b/L$ has an end point (both end points) in $Z'$.
\item[(sc3)] Let the directions $\mu$, $\nu$ be fixed.
The \emph{cell} of the difference quotient $(\partial_{\nu}A_{\mu})(c)$, or briefly the cell at $c$, is the set of the four points $c$, $c+\eta_k\cdot e_{\mu}$, $c+\eta_k\cdot e_{\nu}$, $c+\eta_k\cdot e_{\mu}+\eta_k\cdot e_{\nu}$ of (G2.r4), its corners; for $\nu=\mu$ the corners are the three points $c$, $c+\eta_k\cdot e_{\mu}$, $c+2\eta_k\cdot e_{\mu}$.
By (G2.r4) the quotient lies inside a set of sites when all corners of its cell belong to the set; the cell is then said to lie inside the set.
Division by $L$ sends the corners to $c/L$, $c/L+\eta_{k+1}\cdot e_{\mu}$, $c/L+\eta_{k+1}\cdot e_{\nu}$, $c/L+\eta_{k+1}\cdot e_{\mu}+\eta_{k+1}\cdot e_{\nu}$.
Consequently, in the sense of (G2.r4), a bond or a difference quotient lies inside $LZ'$ if and only if its image lies inside $Z'$.
If $0<d(c,c')\le 1$, then $0<d(c/L,c'/L)\le 1/L\le 1$; so the image of a H\"older difference inside $LZ'$ is a H\"older difference inside $Z'$.
\item[(sc4)] The side $\ell$ is a multiple of $LM$.
For $z'\in LM\mathbb{Z}^3$ the $LM$-cube $\{c:\ |c-z'|<L\cdot M/2\}$ coincides with $L\square'$, where $\square'$ is the $M$-cube $\{c':\ |c'-z'/L|<M/2\}$ of $T_{k+1}$ and $z'/L\in M\mathbb{Z}^3$; every $M$-cube of $T_{k+1}$ is of this form.
Since $L$ is odd, every $z\in M\mathbb{Z}^3$ has the form $z=z'+M\cdot j$ with exactly one $z'\in LM\mathbb{Z}^3$ and $j\in\mathbb{Z}^3$, $|j|\le(L-1)/2$.
The $M$-cube $\square$ with the centre $z$ is contained in the $LM$-cube with the centre $z'$, for every $c\in\square$ satisfies
\[
|c-z'|\le|c-z|+|z-z'|<\frac{M}{2}+\frac{(L-1)\cdot M}{2}=\frac{L\cdot M}{2}.
\]
Therefore every $M$-cube of $T_k$ is contained in exactly one of the sets $L\square'$, and $|z-z'|\le(L-1)\cdot M/2$.
\item[(sc5)] Let $c,c'\in\mathbb{L}_k$, and let $\pi$ be an order of the coordinates.
The sites of $\boldsymbol{\Gamma}^{\pi}(c/L,c'/L)$, formed in $\mathbb{L}_{k+1}$, are the sites of $\boldsymbol{\Gamma}^{\pi}(c,c')$ divided by $L$: both paths change the coordinates one after the other in the order $\pi$, the first one in steps of $\eta_{k+1}=\eta_k/L$, the second one in steps of $\eta_k$.
\end{itemize}

\subsubsection*{0.2. Division by $L$ on the fermion sites and on the pairs (Section~I.3.7: $x/L:=(c/L,g)$)}

The data are those of 0.1.

\emph{The objects of level $k+1$.}
In the setting (G2.S-2), hence in Section~3.4, one has $\ell/L\in\mathrm{Cl}$, and $T_{k+1}$, $T^{\mathrm{f}}_{k+1}$, $\mathrm{Pzt}_{k+1}$, $\mathrm{Pz}_{k+1}$ and the generators of level $k+1$ are the objects of Sections~I.1.1, I.1.3, I.1.7 and~I.3.6, formed for the level $k+1$ and the side $\ell/L$.
If only $\ell\in\mathrm{Cl}$ is assumed, as in Section~4, then the first two sentences of Proposition~16.1\,(d) show that $\ell/L$ is an odd multiple of $M$ with $\ell/L\ge 13M>2M$, and the same objects are defined by the same formulas.
Of the side $\ell/L$ the facts (sf1)--(sf4) require only that these objects be defined.
\begin{itemize}
\item[(sf1)] The map $x=(c,g)\mapsto x/L=(c/L,g)$ is a bijection of $\mathbb{L}^{\mathrm{f}}_k=\mathbb{L}_k\times O_{\mathrm{cub}}$ onto $\mathbb{L}^{\mathrm{f}}_{k+1}$, and $(x+\ell\cdot n)/L=x/L+(\ell/L)\cdot n$ by (sc1); so it induces a bijection of $T^{\mathrm{f}}_k$ onto $T^{\mathrm{f}}_{k+1}$.
One has $\pi(x/L)=\pi(x)/L$.
The label $g$ does not change, and the $48$ sites of the cluster of $c$ are mapped onto the $48$ sites of the cluster of $c/L$.
\item[(sf2)] For $i=0,1,2,3$ one has $x^{(i)}/L=(x/L)^{(i)}$, where the neighbour on the right is formed at level $k+1$ according to (I.1.2).
Indeed, $(c,g\,s_i)/L=(c/L,g\,s_i)$ for $i\le 2$, and $(c+\eta_k\cdot g\,e_1,g\,s_3)/L=(c/L+\eta_{k+1}\cdot g\,e_1,g\,s_3)$.
\item[(sf3)] For $x\ne x'$, the formula (I.1.5), applied at the two levels, gives
\[
d_X(x/L,x'/L)=\max\left(|c/L-c'/L|,\eta_{k+1}\right)=\frac{\max\left(|c-c'|,\eta_k\right)}{L}=\frac{d_X(x,x')}{L}.
\]
\item[(sf4)] By (sf1) and (sf3), the map $(x,x')\mapsto(x/L,x'/L)$ is a bijection of $\mathrm{Pzt}_k$, the set of the ordered pairs of distinct points of $\mathbb{L}^{\mathrm{f}}_k$ with $d_X\le M$, onto the set of the ordered pairs of distinct points of $\mathbb{L}^{\mathrm{f}}_{k+1}$ with $d_X\le M/L$; the latter set is contained in $\mathrm{Pzt}_{k+1}$, because $M/L\le M$.
The translation of both points by $\ell\cdot n$ corresponds under this map to the translation of both points by $(\ell/L)\cdot n$.
Hence the map induces a bijection of $\mathrm{Pz}_k$ onto the set of those pairs of $\mathrm{Pz}_{k+1}$ that satisfy $d_X\le M/L$.
A pair inside a cluster ($c=c'$, $g\ne g'$, $d_X=\eta_k$) is sent to a pair inside a cluster ($d_X=\eta_{k+1}$).
Every pair of $\mathrm{Pz}_{k+1}$ inside a cluster belongs to the image, since $\eta_{k+1}\le 1\le M/L$.
\end{itemize}

\subsubsection*{0.3. Balls}

Let $\tilde{c}$ be a photon site of $T_k$, and let $r>0$ be a real number with $2r+1<\ell$.
As in (G2.r2), $K(\tilde{c},r)=\{c:\ d(c,\tilde{c})\le r\}$.
Its lift is $\tilde{c}+\{v\in\mathbb{L}_k:\ |v|\le r\}$, a set of sup-diameter at most $2r$.
Section~I.1.1, applied with $\mathrm{sd}=1$, shows that for $c,c'\in K(\tilde{c},r)$ with $d(c,c')\le 1$ the vector $(c-\tilde{c})-(c'-\tilde{c})$ is the representative of $c-c'$ of sup norm at most $1$.
Consequently the following holds.
\begin{itemize}
\item[(bx1)] The bonds, the cells of the difference quotients and the pairs of the H\"older differences inside $K(\tilde{c},r)$ are those of the box $\{v\in\mathbb{L}_k:\ |v|\le r\}$, and $(c+\eta_k\cdot e_{\nu})-\tilde{c}=(c-\tilde{c})+\eta_k\cdot e_{\nu}$ whenever $c$ and $c+\eta_k\cdot e_{\nu}$ lie in $K(\tilde{c},r)$.
\item[(bx2)] The box is convex in the following sense: a point of $\mathbb{L}_k$ each of whose coordinates lies between the corresponding coordinates of two points of the box belongs to the box.
\end{itemize}

\subsubsection*{0.4. Numbers}

By the line (C0) one has $\bar{\varepsilon}=1/200$, and by the line (C7), which is the row (S1) of Proposition~14.2\,(b), one has $1-\alpha<\varepsilon\le\bar{\varepsilon}$.
Since $\alpha<1$ by the line (C1), it follows that $0<1-\alpha<\varepsilon\le\bar{\varepsilon}=1/200$.
The condition (L1) reads $L^{\frac{1}{8}-\frac{3\bar{\varepsilon}}{2}}\ge 2400$, and $\frac{1}{8}-\frac{3\bar{\varepsilon}}{2}=0.1175$.
As $L\ge 1$, every power of $L$ whose real exponent is at most $-0.1175$ is at most $1/2400$.
Each of the numbers
\[
\frac{3}{4},\qquad\frac{7}{4},\qquad\frac{7}{8}-\frac{\varepsilon}{2},\qquad 1-\varepsilon,\qquad\frac{1}{4}-\frac{3\varepsilon}{2},\qquad\frac{1}{8}-\frac{3\varepsilon}{2}
\]
is at least $0.1175$.
It follows that
\begin{equation}\tag{nm1}
L^{-\frac{3}{4}}\le\frac{1}{2400}\le\frac{1}{2},
\end{equation}
\begin{equation}\tag{nm2}
1382\,L^{-\frac{15}{8}+\frac{3\varepsilon}{2}}=1382\,L^{-\left(\frac{1}{8}-\frac{3\varepsilon}{2}\right)}\cdot L^{-\frac{7}{4}}\le\frac{1382}{2400}\,L^{-\frac{7}{4}}<L^{-\frac{7}{4}},
\end{equation}
\begin{equation}\tag{nm3}
13.6\,L^{-2+\frac{3\varepsilon}{2}}=13.6\,L^{-\left(\frac{1}{4}-\frac{3\varepsilon}{2}\right)}\cdot L^{-\frac{7}{4}}\le\frac{13.6}{2400}\,L^{-\frac{7}{4}}\le L^{-\frac{7}{4}},
\end{equation}
\begin{equation}\tag{nm4}
L^{-\frac{7}{4}}+4L^{-\frac{7}{8}+\frac{\varepsilon}{2}}\le\frac{5}{2400}<1\qquad\text{and}\qquad L^{-\frac{7}{4}}+4L^{-1+\varepsilon}\le\frac{5}{2400}<1.
\end{equation}

Now let $e_k\in(0,e_0]$.
Proposition~14.2\,(b) imposes the conditions on $e_0$ for all $v_k=-\log e_k\ge-\log e_0$, and $\exp(-\varepsilon\cdot v_k)=e_k^{\varepsilon}$.
Therefore the first inequality of (A8) gives (nm5), and the first and the third inequality of (A9) give (nm6) and (nm7), respectively:
\begin{equation}\tag{nm5}
e_k^{\varepsilon}\le\frac{1}{4},\qquad\text{hence}\qquad e_k<1;
\end{equation}
\begin{equation}\tag{nm6}
160\,L^{1+\frac{3\varepsilon}{2}}\cdot M\cdot e_k^{\varepsilon}\le 1,\qquad\text{hence}\qquad 160\,L\cdot M\cdot e_k^{\varepsilon}\le 1;
\end{equation}
\begin{equation}\tag{nm7}
9\,M\cdot L^{\varepsilon}\cdot e_k^{2\varepsilon}\le 1.
\end{equation}
The thresholds of (G2.r4) satisfy
\[
\mathrm{t2}_k=e_k^{\varepsilon}\cdot\mathrm{t1}_k,\qquad\mathrm{t3}_k=e_k^{\varepsilon}\cdot\mathrm{t2}_k,\qquad\mathrm{td2}_k=e_k^{\varepsilon}\cdot\mathrm{td1}_k,\qquad e_k\cdot\mathrm{t1}_k=e_k^{\frac{1}{4}+\varepsilon},\qquad e_k\cdot\mathrm{td1}_k=e_k^{2\varepsilon}.
\]
If $e_{k+1}=L^{\frac{1}{2}}\cdot e_k$, then
\begin{equation}\tag{nm8}
\begin{aligned}
&\mathrm{t1}_{k+1}=L^{-\frac{3}{8}+\frac{\varepsilon}{2}}\cdot\mathrm{t1}_k,\qquad\mathrm{t2}_{k+1}=L^{-\frac{3}{8}+\varepsilon}\cdot\mathrm{t2}_k,\qquad\mathrm{t3}_{k+1}=L^{-\frac{3}{8}+\frac{3\varepsilon}{2}}\cdot\mathrm{t3}_k,\\
&\mathrm{td1}_{k+1}=L^{-\frac{1}{2}+\varepsilon}\cdot\mathrm{td1}_k,\qquad\mathrm{td2}_{k+1}=L^{-\frac{1}{2}+\frac{3\varepsilon}{2}}\cdot\mathrm{td2}_k;
\end{aligned}
\end{equation}
the exponents are the halves of $-\frac{3}{4}+\varepsilon$, $-\frac{3}{4}+2\varepsilon$, $-\frac{3}{4}+3\varepsilon$, $-1+2\varepsilon$, $-1+3\varepsilon$.

\subsection*{1. Proposition 18.1}

\begin{proof}[Proof of Proposition~18.1]
The proof involves photon objects only.
By a bond function on $Z'$ is meant a complex function $A$ on the bonds inside $Z'$, that is, on the bonds with both end points in $Z'$, such that $A(c',c)=-A(c,c')$.
By (sc2) the function $A_L(b)=L^{-\frac{1}{2}}\cdot A(b/L)$ is defined on the bonds $b$ inside $LZ'$, and $A_L(c',c)=-A_L(c,c')$.
The three norms on $Z'$ and on $LZ'$ are those of (G2.r4).
A formula is asserted \emph{on $LZ'$} at every place at which its left-hand side is formed inside $LZ'$; by (sc3) its right-hand side is then formed inside $Z'$.

\textbf{(a)} \emph{The three formulas.}
Since $\eta_{k+1}=\eta_k/L$, (sc2) gives
\[
(A_L)_{\mu}(c)=A_L(c,c+\eta_k\cdot e_{\mu})=L^{-\frac{1}{2}}\cdot A\left(c/L,c/L+\eta_{k+1}\cdot e_{\mu}\right)=L^{-\frac{1}{2}}\cdot A_{\mu}(c/L)
\]
and
\begin{align*}
\left(\partial_{\nu}(A_L)_{\mu}\right)(c)&=\eta_k^{-1}\cdot\left[(A_L)_{\mu}(c+\eta_k\cdot e_{\nu})-(A_L)_{\mu}(c)\right]\\
&=\eta_k^{-1}\cdot L^{-\frac{1}{2}}\cdot\left[A_{\mu}(c/L+\eta_{k+1}\cdot e_{\nu})-A_{\mu}(c/L)\right]\\
&=\eta_k^{-1}\cdot L^{-\frac{1}{2}}\cdot\eta_{k+1}\cdot(\partial_{\nu}A_{\mu})(c/L)=L^{-\frac{3}{2}}\cdot(\partial_{\nu}A_{\mu})(c/L).
\end{align*}
Let $0<d(c,c')\le 1$.
By (sc1), $d(c,c')^{-\alpha}=L^{-\alpha}\cdot d(c/L,c'/L)^{-\alpha}$, and therefore
\begin{multline*}
d(c,c')^{-\alpha}\cdot\left[\left(\partial_{\nu}(A_L)_{\mu}\right)(c')-\left(\partial_{\nu}(A_L)_{\mu}\right)(c)\right]\\
=L^{-\frac{3}{2}-\alpha}\cdot d(c/L,c'/L)^{-\alpha}\cdot\left[(\partial_{\nu}A_{\mu})(c'/L)-(\partial_{\nu}A_{\mu})(c/L)\right].
\end{multline*}
Here $0<d(c/L,c'/L)\le 1$ by (sc3), so that the right-hand side is $L^{-\frac{3}{2}-\alpha}$ times a H\"older quotient of $\partial A$ inside $Z'$.

\emph{The norms.}
In view of (sc2) and (sc3), the first two formulas give $|A_L|=L^{-\frac{1}{2}}\cdot|A|$ and $|\partial A_L|=L^{-\frac{3}{2}}\cdot|\partial A|$, the norms of $A_L$ being formed on $LZ'$ and those of $A$ on $Z'$.
The third formula gives $|\delta_{\alpha}\partial A_L|\le L^{-\frac{3}{2}-\alpha}\cdot|\delta_{\alpha}\partial A|$; the supremum on the right-hand side extends over a larger set of pairs.
All these suprema are maxima over finite sets, the maximum over the empty set being $0$.
Suppose now that $|A|<\mathrm{t1}_{k+1}$, $|\partial A|<\mathrm{t2}_{k+1}$ and $|\delta_{\alpha}\partial A|<\mathrm{t3}_{k+1}$ on $Z'$.
Then (nm8) yields
\begin{gather*}
|A_L|<L^{-\frac{1}{2}}\cdot L^{-\frac{3}{8}+\frac{\varepsilon}{2}}\cdot\mathrm{t1}_k=L^{-\frac{7}{8}+\frac{\varepsilon}{2}}\cdot\mathrm{t1}_k,\qquad|\partial A_L|<L^{-\frac{3}{2}}\cdot L^{-\frac{3}{8}+\varepsilon}\cdot\mathrm{t2}_k=L^{-\frac{15}{8}+\varepsilon}\cdot\mathrm{t2}_k,\\
|\delta_{\alpha}\partial A_L|<L^{-\frac{3}{2}-\alpha}\cdot L^{-\frac{3}{8}+\frac{3\varepsilon}{2}}\cdot\mathrm{t3}_k=L^{-\alpha}\cdot L^{-\frac{15}{8}+\frac{3\varepsilon}{2}}\cdot\mathrm{t3}_k.
\end{gather*}

\textbf{(b)}
Take $Z'=T_{k+1}$ and $A=\mathscr{C}$ in the first two formulas of (a).
Together with (nm8) they give
\[
|\mathscr{C}_L|=L^{-\frac{1}{2}}\cdot|\mathscr{C}|<L^{-\frac{1}{2}}\cdot\mathrm{td1}_{k+1}=L^{-1+\varepsilon}\cdot\mathrm{td1}_k,\qquad|\partial\mathscr{C}_L|=L^{-\frac{3}{2}}\cdot|\partial\mathscr{C}|<L^{-\frac{3}{2}}\cdot\mathrm{td2}_{k+1}=L^{-2+\frac{3\varepsilon}{2}}\cdot\mathrm{td2}_k.
\]
As $\varepsilon\le 1/200$, one has $-1+\varepsilon\le-\frac{3}{4}$ and $-2+\frac{3\varepsilon}{2}\le-\frac{3}{4}$; moreover $L\ge 1$.
Hence both norms of $L^{\frac{3}{4}}\cdot\mathscr{C}_L$ are smaller than the thresholds of $\mathfrak{X}_k$, that is, $\mathscr{C}_L\in L^{-\frac{3}{4}}\cdot\mathfrak{X}_k$.
Finally $L^{-\frac{3}{4}}\cdot\mathfrak{X}_k\subset\mathfrak{X}_k$, since $L^{-\frac{3}{4}}\le 1$.

\textbf{(c)} \emph{The gauge function.}
Let $\boldsymbol{\lambda}$ be a function of the sites of a subset of $T_{k+1}$, and let $\boldsymbol{\lambda}_L(c):=L^{\frac{1}{2}}\cdot\boldsymbol{\lambda}(c/L)$, as in Section~I.3.7.
Then
\begin{align*}
(d\boldsymbol{\lambda})_L(c,c+\eta_k\cdot e_{\mu})&=L^{-\frac{1}{2}}\cdot\eta_{k+1}^{-1}\cdot\left[\boldsymbol{\lambda}(c/L+\eta_{k+1}\cdot e_{\mu})-\boldsymbol{\lambda}(c/L)\right]\\
&=L^{\frac{1}{2}}\cdot\eta_k^{-1}\cdot\left[\boldsymbol{\lambda}(c/L+\eta_{k+1}\cdot e_{\mu})-\boldsymbol{\lambda}(c/L)\right]\\
&=\eta_k^{-1}\cdot\left[\boldsymbol{\lambda}_L(c+\eta_k\cdot e_{\mu})-\boldsymbol{\lambda}_L(c)\right]=\left(d(\boldsymbol{\lambda}_L)\right)(c,c+\eta_k\cdot e_{\mu}).
\end{align*}
The map $A\mapsto A_L$ is linear.
Therefore, by (sc2), the identity $A_0=A'+d\boldsymbol{\lambda}$ on the bonds inside $\square'^{\dagger}$ implies $(A_0)_L=(A')_L+d(\boldsymbol{\lambda}_L)$ on the bonds inside $L\square'^{\dagger}$.
The functions $(A_0)_L$, $(A')_L$ and $\boldsymbol{\lambda}_L$ are real.

\emph{The set.}
Let $z'/L\in M\mathbb{Z}^3$ be the centre of $\square'$.
By \cite[Definition~5.7]{QEDI}, $\square'^{\dagger}=\square'^{\sim L}$ is the union of the $M$-cubes with the centres $z'/L+M\cdot j$, $|j|\le L$; as no site lies on a face (Section~I.1.1), this union is the set $\left\{c':\ |c'-z'/L|<\left(L+\frac{1}{2}\right)\cdot M\right\}$.
It does not wrap, because $(2L+1)\cdot M<13L\cdot M\le\ell/L$.
By (sc1), $L\square'^{\dagger}=\left\{c:\ |c-z'|<\left(L+\frac{1}{2}\right)\cdot L\cdot M\right\}$, and $z'$ is the centre of the $LM$-cube $L\square'$ by (sc4).

\emph{The bounds.}
By \cite[Definition~5.16]{QEDI}, the three norms of $A'$ on $\square'^{\dagger}$, as well as the three norms of $A_1$ on $T_{k+1}$, are smaller than the thresholds of $R_{k+1}$.
Part (a), applied with $Z'=\square'^{\dagger}$ and with $Z'=T_{k+1}$, respectively, furnishes the bounds of (a) for $(A')_L$ on $L\square'^{\dagger}$ and for $(A_1)_L$ on $T_k$.

\emph{The cubes of level $k$.}
Let $\square$ be an $M$-cube of $T_k$ contained in $L\square'$, and let $z$ be its centre.
By (sc4), $|z-z'|\le(L-1)\cdot M/2$.
For $c\in\square^{\dagger}=\left\{c:\ |c-z|<\left(L+\frac{1}{2}\right)\cdot M\right\}$ it follows that
\[
|c-z'|<\left(L+\frac{1}{2}\right)\cdot M+\frac{(L-1)\cdot M}{2}=\frac{3L}{2}\cdot M\le\left(L+\frac{1}{2}\right)\cdot L\cdot M.
\]
Thus $\square^{\dagger}$ is contained in $L\square'^{\dagger}$.

\emph{Conclusion.}
One has $\mathscr{A}_L=(A_0)_L+(A_1)_L$.
The exponents $-\frac{7}{8}+\frac{\varepsilon}{2}$, $-\frac{15}{8}+\varepsilon$ and $-\alpha-\frac{15}{8}+\frac{3\varepsilon}{2}$ are at most $-\frac{3}{4}$, because $\varepsilon\le 1/200$ and $\alpha>0$.
Hence the three norms of $(A_1)_L$ on $T_k$ are smaller than $L^{-\frac{3}{4}}\cdot\mathrm{t1}_k$, $L^{-\frac{3}{4}}\cdot\mathrm{t2}_k$ and $L^{-\frac{3}{4}}\cdot\mathrm{t3}_k$, that is, $(A_1)_L\in L^{-\frac{3}{4}}\cdot R_k$; and $(A_0)_L$ is real.
By (sc4), every $M$-cube $\square$ of $T_k$ is contained in $L\square'$ for exactly one $M$-cube $\square'$ of $T_{k+1}$.
For this cube $\square'$ let $(A',\boldsymbol{\lambda})$ be a pair as in the data of Proposition~18.1\,(c).
Thus $A'$ is a real function on the bonds inside $\square'^{\dagger}$ and $\boldsymbol{\lambda}$ is a real function on the sites of $\square'^{\dagger}$;
one has $A_0=A'+d\boldsymbol{\lambda}$ on the bonds inside $\square'^{\dagger}$, so that $(A',\boldsymbol{\lambda})$ is a representative of $A_0$ on $\square'^{\dagger}$ in the sense of \cite[Section~17.3]{QEDI};
and the three norms of $A'$ on $\square'^{\dagger}$ are smaller than the thresholds $\mathrm{t1}_{k+1}$, $\mathrm{t2}_{k+1}$, $\mathrm{t3}_{k+1}$ of $R_{k+1}$.
Such a pair exists by \cite[Definition~5.16\,(b)]{QEDI} with $\mathrm{cf}=1$, because $\mathscr{A}=A_0+A_1$ is split as in that definition.
The paragraphs on the gauge function, on the set, on the bounds and on the cubes of level $k$ apply to $\square'$, to $\square$ and to this pair.
The restrictions of $(A')_L$ and of $\boldsymbol{\lambda}_L$ to $\square^{\dagger}$ are real and satisfy $(A_0)_L=(A')_L+d(\boldsymbol{\lambda}_L)$ on $\square^{\dagger}$.
The three norms of $(A')_L$ on the subset $\square^{\dagger}$ of $L\square'^{\dagger}$ do not exceed those on $L\square'^{\dagger}$, and are therefore smaller than $L^{-\frac{3}{4}}\cdot\mathrm{t1}_k$, $L^{-\frac{3}{4}}\cdot\mathrm{t2}_k$ and $L^{-\frac{3}{4}}\cdot\mathrm{t3}_k$.
By the definition of $\mathrm{cf}\cdot\tilde{R}_k$ in \cite[Definition~5.16\,(b)]{QEDI}, with $\mathrm{cf}=L^{-\frac{3}{4}}\in(0,1]$, this means that $\mathscr{A}_L\in L^{-\frac{3}{4}}\cdot\tilde{R}_k$.

The set $L^{-\frac{3}{4}}\cdot\tilde{R}_k$ is contained in $\tilde{R}_k$, because $L^{-\frac{3}{4}}\le 1$, so that the conditions with $\mathrm{cf}=L^{-\frac{3}{4}}$ imply those with $\mathrm{cf}=1$.
It is also contained in $\tilde{R}_k^{-}$, which is defined with $\mathrm{cf}=1$ for $A'$ and with the condition $A_1\in\left(1-e_k^{\varepsilon}\right)\cdot R_k$ (\cite[Definition~5.16\,(c)]{QEDI}); indeed $L^{-\frac{3}{4}}\le\frac{1}{2}\le 1-e_k^{\varepsilon}$ by (nm1) and (nm5).
\end{proof}

\begin{named}{Remark}[Level $0$]
The case $k=0$ is covered by the proof above.
There $\eta_0=1$ and $\eta_1=1/L$, and $e_k$ is the number $q$.
The H\"older differences of level $0$ are formed at the pairs with $d(c,c')=1$, and division by $L$ sends these pairs to pairs of level $1$ at the distance $1/L$.
\end{named}

\begin{named}{Remark}[Uniformity]
The factors in Proposition~18.1 are powers of $L$ with the exponents displayed there.
The data $k$, $\ell$ and $e_k$ enter only through $\eta_{k+1}=\eta_k/L$, $e_{k+1}=L^{\frac{1}{2}}\cdot e_k$, the conditions $\ell\in\mathrm{Cl}$ and $\ell/L\in\mathrm{Cl}$, and the inequality $e_k\le e_0$, which gives (nm5) in the last paragraph of the proof of (c).
No quantity depends on $N$, on a torus of the scheme or on a twist.
\end{named}

\subsection*{2. Proposition 18.2}

\emph{Proof of Proposition~18.2.}

\subsubsection*{2.0. Notation}

For a bond function $\mathrm{bq}$ and a set $Z$ of sites, $|\mathrm{bq}|_Z$, $|\partial\mathrm{bq}|_Z$ and $|\delta_{\alpha}\partial\mathrm{bq}|_Z$ denote the three norms formed inside $Z$, as in (G2.r4); without the index $Z$ the norms are formed on $T_k$.
The components are $\mathrm{bq}_{\mu}(c)=\mathrm{bq}(c,c+\eta_k\cdot e_{\mu})$; a bond function is determined by its components, that is, by its values on the positively oriented bonds, since $\mathrm{bq}(c',c)=-\mathrm{bq}(c,c')$.
A bond function $\mathrm{kc}$ is called constant if its components $\mathrm{kc}_{\mu}$ do not depend on $c$; then $|\mathrm{kc}|=\max_{\mu}|\mathrm{kc}_{\mu}|$, and $\mathrm{kc}\cdot v:=\sum_{\mu}\mathrm{kc}_{\mu}\cdot v_{\mu}$ for $v\in\mathbb{R}^3$.
For $c$ in a ball $K(\tilde{c},r)$, $c-\tilde{c}$ is the vector of (G2.r2).
Finally, the following abbreviations are used:
\[
\mathrm{bd1}:=L^{-\frac{7}{8}+\frac{\varepsilon}{2}}\cdot\mathrm{t1}_k,\qquad\mathrm{bd2}:=L^{-\frac{15}{8}+\varepsilon}\cdot\mathrm{t2}_k,\qquad\mathrm{bd3}:=L^{-\frac{15}{8}+\frac{3\varepsilon}{2}}\cdot\mathrm{t3}_k.
\]

\subsubsection*{2.1. Results of \cite[Section~13]{QEDI} on photon objects, in the notation of Part~I}

\begin{itemize}
\item[(qu1)] \emph{Cut-off.}
For a real variable $\mathrm{rv}$ let
\[
\mathrm{ctf}(\mathrm{rv}):=\begin{cases}
1, & |\mathrm{rv}|\le 1,\\
1-2\left(|\mathrm{rv}|-1\right)^2, & 1\le|\mathrm{rv}|\le\frac{3}{2},\\
2\left(2-|\mathrm{rv}|\right)^2, & \frac{3}{2}\le|\mathrm{rv}|\le 2,\\
0, & |\mathrm{rv}|\ge 2.
\end{cases}
\]
The function $\mathrm{ctf}$ is continuously differentiable, $0\le\mathrm{ctf}\le 1$, $|\mathrm{ctf}'|\le 2$, and $\mathrm{ctf}'$ is Lipschitz continuous with the constant $4$.
Let $\tilde{c}$ be a photon site and $r>0$ a real number with $4r+11<\ell$.
For a point $\mathrm{pt}$ of the continuum torus $\mathbb{R}^3/\ell\mathbb{Z}^3$ put
\[
\mathrm{cut}_r(\mathrm{pt}):=\begin{cases}
\displaystyle\prod_{\mu}\mathrm{ctf}\left(\frac{(\mathrm{pt}-\tilde{c})_{\mu}}{r}\right), & |\mathrm{pt}-\tilde{c}|\le 2r+5,\\
0 & \text{otherwise},
\end{cases}
\]
where $(\mathrm{pt}-\tilde{c})_{\mu}$ are the coordinates of the representative of $\mathrm{pt}-\tilde{c}$ of sup norm at most $2r+5$.
Let $\mathrm{fc}_{\mu}(c):=\mathrm{cut}_r(c+\eta_k\cdot e_{\mu}/2)$ be the value of $\mathrm{cut}_r$ at the barycentre of the bond $(c,c+\eta_k\cdot e_{\mu})$ (\cite[Section~13.1.7]{QEDI}; Section~I.1.1).
Then
\begin{equation}\tag{f1}
\begin{aligned}
&0\le\mathrm{fc}_{\mu}\le 1;\qquad\mathrm{fc}_{\mu}(c)=1\ \text{ if $c$ and $c+\eta_k\cdot e_{\mu}$ lie in $K(\tilde{c},r)$};\\
&\mathrm{fc}_{\mu}(c)=0\ \text{ if }d(c,\tilde{c})\ge 2r+1;
\end{aligned}
\end{equation}
\begin{equation}\tag{f2}
|\mathrm{fc}_{\mu}(c)-\mathrm{fc}_{\mu}(c')|\le\frac{2}{r}\cdot|c-c'|_1\le\frac{6}{r}\cdot|c-c'|,\qquad|\partial_{\nu}\mathrm{fc}_{\mu}|\le\frac{2}{r};
\end{equation}
\begin{equation}\tag{f3}
|\partial_{\nu}\mathrm{fc}_{\mu}(c)-\partial_{\nu}\mathrm{fc}_{\mu}(c')|\le\frac{4}{r^2}\cdot|c-c'|_1\le\frac{12}{r^2}\cdot|c-c'|.
\end{equation}
\item[(qu2)] \emph{\cite[Lemma~13.9]{QEDI}.}
Let $r\ge 50$ with $4r+11<\ell$, and let $\mathrm{bq}$ be a complex function on the bonds inside $K(\tilde{c},2r+5)$.
Its three norms, formed inside $K(\tilde{c},2r+5)$, are denoted by $\mathrm{bn}_0:=|\mathrm{bq}|$, $\mathrm{bn}_1:=|\partial\mathrm{bq}|$ and $\mathrm{bn}_{\alpha}:=|\delta_{\alpha}\partial\mathrm{bq}|$.
Define a bond function $\mathrm{cut}_r\mathrm{bq}$ on $T_k$ by
\[
(\mathrm{cut}_r\mathrm{bq})_{\mu}(c):=\begin{cases}
\mathrm{fc}_{\mu}(c)\cdot\mathrm{bq}_{\mu}(c), & c\in K(\tilde{c},2r+1),\\
0 & \text{for all other sites $c$ of $T_k$}.
\end{cases}
\]
Then $\mathrm{cut}_r\mathrm{bq}=\mathrm{bq}$ on the bonds inside $K(\tilde{c},r)$, and on $T_k$
\[
|\mathrm{cut}_r\mathrm{bq}|\le\mathrm{bn}_0,\qquad|\partial(\mathrm{cut}_r\mathrm{bq})|\le\mathrm{bn}_1+\frac{2}{r}\cdot\mathrm{bn}_0,\qquad|\delta_{\alpha}\partial(\mathrm{cut}_r\mathrm{bq})|\le\mathrm{bn}_{\alpha}+\frac{12}{r}\cdot\mathrm{bn}_1+\frac{12}{r^2}\cdot\mathrm{bn}_0.
\]
By its definition, the map $\mathrm{bq}\mapsto\mathrm{cut}_r\mathrm{bq}$ is linear.
\item[(qu3)] \emph{\cite[Lemma~13.10]{QEDI}.}
Let $\mathrm{fn}$ be a complex function on the sites of $T_k$ with $\sum_c\mathrm{fn}(c)=0$, and let $\mathrm{osc}\ge 0$ be a number such that $|\mathrm{fn}(c)-\mathrm{fn}(c')|\le\mathrm{osc}$ whenever $d(c,c')\le 1$.
Then $|\mathrm{fn}|_{\infty}\le\ell\cdot\mathrm{osc}$.
For every bond function $\mathrm{bq}$ on $T_k$ the function $\mathrm{fn}=\partial_{\nu}\mathrm{bq}_{\mu}$ satisfies the hypothesis on the sum, since the sum over $c$ telescopes in the direction $\nu$.
\item[(qu4)] \emph{\cite[Theorem~13.2\,(t7)]{QEDI}.}
Any two points of a polymer $X$ have the sup distance at most $M\cdot\left(d_M(X)+2\right)$.
A small polymer is contained in $K\left(\tilde{c},M\cdot(L+2)\right)$ for each of its sites $\tilde{c}$.
\end{itemize}

\emph{Setting.}
In \cite{QEDI} the facts (qu1)--(qu3) are stated in the setting (Sb) of \cite[Section~13.1.1]{QEDI}, which coincides with the setting (G2.S-2) of Proposition~18.2.
They concern the single torus $T_k$, and no charge enters them.
In (f2) and (f3) the sites $c$, $c'$ lie in $K(\tilde{c},2r+4)$, and $c-c'$ is the difference of the vectors $c-\tilde{c}$ and $c'-\tilde{c}$ of (G2.r2).
The proof of (qu2) rests on $\alpha\le 1$ and on the fact that inside $K(\tilde{c},2r+5)$ the bonds, the difference quotients and the pairs at distance at most $1$ are those of a box of $\mathbb{L}_k$; by Section~0.3, applied with the radius $2r+5$, this holds as soon as $2(2r+5)+1<\ell$, which is the hypothesis $4r+11<\ell$ (Appendix~A).
The fact (qu4) holds for every $\ell\in\mathrm{Cl}$.

\subsubsection*{2.2. Construction (Steps 0--6)}

In this chapter a number of displays carry the tags (18a.1), (18a.2), \dots, assigned in the order in which the displays occur, from Step~0 below on.
Each of these tags is the name of the display that carries it.

\paragraph{Step 0 (data).}
Choose a splitting $\mathscr{A}=A_0+A_1$ as in \cite[Definition~5.16]{QEDI}.
Let $\square'$ be the $M$-cube of $T_{k+1}$ that contains $\tilde{c}/L$.
By (sc4), $L\square'$ is the $LM$-cube of $\tilde{c}$; its centre is denoted by $z'$.
Choose a representative of $A_0$ on $\square'^{\dagger}$, that is, a real bond function $A'$ and a real function $\boldsymbol{\lambda}$ such that $A_0=A'+d\boldsymbol{\lambda}$ on $\square'^{\dagger}$ and such that the three norms of $A'$ on $\square'^{\dagger}$ are below the thresholds of $R_{k+1}$.
Put
\begin{gather*}
\mathrm{ba0}:=(A_0)_L\ \ (\text{real}),\qquad
\mathrm{ba1}:=(A_1)_L,\qquad
\mathrm{ba0}':=(A')_L\ \ (\text{real}),\\
\mathrm{Zd}:=L\square'^{\dagger}=\left\{c:|c-z'|<\left(L+\frac{1}{2}\right)LM\right\},\qquad
\mathrm{Kb}:=K(\tilde{c},2\,\mathrm{rad}+5).
\end{gather*}
Proposition~18.1\,(c), (a) and the inequality $L^{-\alpha}\le1$ give
\begin{equation}\tag{18a.1}
\mathscr{A}_L=\mathrm{ba0}+\mathrm{ba1}\ \text{ on }T_k,\qquad
\mathrm{ba0}=\mathrm{ba0}'+d(\boldsymbol{\lambda}_L)\ \text{ on the bonds inside }\mathrm{Zd},
\end{equation}
\begin{equation}\tag{18a.2}
\begin{aligned}
&|\mathrm{ba1}|<\mathrm{bd1},\qquad|\partial\,\mathrm{ba1}|<\mathrm{bd2},\qquad|\delta_\alpha\partial\,\mathrm{ba1}|<\mathrm{bd3}\qquad\text{on }T_k,\\
&|\partial\,\mathrm{ba0}'|_{\mathrm{Zd}}<\mathrm{bd2},\qquad|\delta_\alpha\partial\,\mathrm{ba0}'|_{\mathrm{Zd}}<\mathrm{bd3}.
\end{aligned}
\end{equation}

\emph{Geometry.}
Recall that $L\ge5$ and $M\ge5$ are integers, and that $\ell\ge13LM$ and $\ell/L\ge13LM$, since both $\ell$ and $\ell/L$ belong to $\mathrm{Cl}$.
\begin{itemize}
\item[(ge1)] Let $\square''$ be an $M$-cube of $T_{k+1}$.
The set $L\square''^{\dagger}$ is a cube of side $(2L+1)LM$; the sup-diameter of its lift is $(2L+1)LM-\eta_k$, and $(2L+1)LM+1<13L^2M\le\ell$.
Hence $L\square''^{\dagger}$ does not wrap, and by Section~I.1.1, applied with $\mathrm{sd}=1$, its bonds, its difference quotients and its pairs of sites at distance at most $1$ are those of a cube of $\mathbb{L}_k$.
\item[(ge2)] One has $|\tilde{c}-z'|<LM/2$.
Consequently, if $0<r\le L^2M$ and $c\in K(\tilde{c},r)$, then $|c-z'|<r+LM/2\le\left(L+\frac{1}{2}\right)LM$; in other words, $K(\tilde{c},r)$ is contained in $\mathrm{Zd}$.
In particular $K(\tilde{c},2\,\mathrm{rad}+6)$ is contained in $\mathrm{Zd}$, because $2\,\mathrm{rad}+6=4LM+6\le5LM\le L^2M$.
\item[(ge3)] The number $\mathrm{rad}=2LM\ge50$ is an integer, and $4\,\mathrm{rad}+11=8LM+11<13LM\le\ell$.
Therefore (qu2) applies with $r=\mathrm{rad}$, and Section~0.3 applies to $\mathrm{Kb}$ and to $K(\tilde{c},2\,\mathrm{rad}+6)$, because $2(2\,\mathrm{rad}+6)+1=8LM+13<13LM$.
\item[(ge4)] By (qu4), $X$ is contained in $K(\tilde{c},M(L+2))$; moreover $M(L+2)+2\le2LM=\mathrm{rad}$, since $2M+2\le LM$.
If a bond has an end point at distance at most $1$ from $X$, then this end point lies in $K(\tilde{c},M(L+2)+1)$, and the other end point lies in $K(\tilde{c},M(L+2)+1+\eta_k)$.
Hence both end points lie in $K(\tilde{c},\mathrm{rad})$.
\end{itemize}

\paragraph{Step 1 (Taylor polynomials).}
Let $\mathrm{bq}$ denote one of the two bond functions $\mathrm{ba0}'$, $\mathrm{ba1}$.
By (ge2), $\mathrm{bq}$ is defined on the bonds inside $K(\tilde{c},2\,\mathrm{rad}+6)$.
For $\nu,\mu\in\{1,2,3\}$ and for every $c$ with $c,c+\eta_ke_\mu\in\mathrm{Kb}$ define
\begin{equation}\tag{18a.3}
\begin{aligned}
\mathrm{Vs}[\mathrm{bq}]_{\nu\mu}&:=(\partial_\nu\mathrm{bq}_\mu)(\tilde{c}),\qquad
\mathrm{Lin}[\mathrm{bq}]_\mu(c):=\sum_\nu\mathrm{Vs}[\mathrm{bq}]_{\nu\mu}\,(c-\tilde{c})_\nu,\\
\mathrm{rem}[\mathrm{bq}]_\mu(c)&:=\mathrm{bq}_\mu(c)-\mathrm{bq}_\mu(\tilde{c})-\mathrm{Lin}[\mathrm{bq}]_\mu(c).
\end{aligned}
\end{equation}
Here $\mathrm{bq}(\tilde{c})$ stands for the constant bond function with the components $\mathrm{bq}_\mu(\tilde{c})$.
The objects $\mathrm{Lin}[\mathrm{bq}]$ and $\mathrm{rem}[\mathrm{bq}]$, and likewise $\mathrm{Lp}$ and $\mathrm{qg}$ below, are defined by the displayed formulas at the point $c$, which for a bond is the initial site of the positively oriented bond; no covariance property of these auxiliary fields is needed.
\begin{itemize}
\item[(1a)] For every difference quotient inside $\mathrm{Kb}$,
\[
(\partial_\nu\mathrm{rem}[\mathrm{bq}]_\mu)(c)=(\partial_\nu\mathrm{bq}_\mu)(c)-(\partial_\nu\mathrm{bq}_\mu)(\tilde{c}).
\]
Indeed, by (bx1) one has $\mathrm{Lin}[\mathrm{bq}]_\mu(c+\eta_ke_\nu)-\mathrm{Lin}[\mathrm{bq}]_\mu(c)=\eta_k\,\mathrm{Vs}[\mathrm{bq}]_{\nu\mu}$, while the constant term does not contribute to a difference quotient.
\item[(1b)] $|\delta_\alpha\partial\,\mathrm{rem}[\mathrm{bq}]|_{\mathrm{Kb}}\le\mathrm{bd3}$.
By (1a), every H\"older difference of $\partial\,\mathrm{rem}[\mathrm{bq}]$ inside $\mathrm{Kb}$ coincides with the H\"older difference of $\partial\,\mathrm{bq}$ at the same two cells.
These cells lie inside $\mathrm{Zd}$ by (ge2), and the bound follows from (18a.2).
\item[(1c)] $|\partial\,\mathrm{rem}[\mathrm{bq}]|_{\mathrm{Kb}}\le(2\,\mathrm{rad}+5)\,\mathrm{bd3}\le3\,\mathrm{rad}\cdot\mathrm{bd3}$.
Let the quotient at $c$ lie inside $\mathrm{Kb}$, and put $v:=c-\tilde{c}$, so that $|v|\le2\,\mathrm{rad}+5$.
Set $c^{(0)}:=\tilde{c}$, and let $c^{(j+1)}$ be obtained from $c^{(j)}$ by moving each coordinate towards the corresponding coordinate of $c$ by the smaller of the two numbers $1$ and the remaining distance.
As $1$ and all coordinates of $v$ are multiples of $\eta_k$, every $c^{(j)}$ is a site.
Since $\mathrm{rad}$ is an integer, after $\mathrm{ns}:=\lceil|v|\rceil\le2\,\mathrm{rad}+5$ steps one arrives at $c^{(\mathrm{ns})}=c$, and every step has sup length in $(0,1]$.
Each coordinate of $c^{(j)}$ lies between the corresponding coordinates of $\tilde{c}$ and of $c$.
So the corners of the cell at $c^{(j)}$ lie between the corresponding corners of the cells at $\tilde{c}$ and at $c$, and all these corners belong to $\mathrm{Kb}$ (the corners of the cell at $\tilde{c}$ have distance at most $2\eta_k\le2$ from $\tilde{c}$).
By (bx2) the cell at $c^{(j)}$ lies inside $\mathrm{Kb}$, and hence inside $\mathrm{Zd}$.
Now (18a.2) gives
\[
\left|(\partial_\nu\mathrm{bq}_\mu)(c^{(j+1)})-(\partial_\nu\mathrm{bq}_\mu)(c^{(j)})\right|\le\left|c^{(j+1)}-c^{(j)}\right|^\alpha\,\mathrm{bd3}\le\mathrm{bd3} .
\]
Summation over $j<\mathrm{ns}$, together with (1a), yields the first inequality; the second one follows from $2\,\mathrm{rad}+5\le3\,\mathrm{rad}$.
\item[(1d)] $|\mathrm{rem}[\mathrm{bq}]|_{\mathrm{Kb}}\le3(2\,\mathrm{rad}+5)^2\,\mathrm{bd3}\le18\,\mathrm{rad}^2\cdot\mathrm{bd3}$.
By (18a.3), $\mathrm{rem}[\mathrm{bq}]_\mu(\tilde{c})=0$.
Let $c$ and $c+\eta_ke_\mu$ lie in $\mathrm{Kb}$.
A rectilinear path from $\tilde{c}$ to $c$ consists of $|c-\tilde{c}|_1/\eta_k\le3(2\,\mathrm{rad}+5)/\eta_k$ bonds.
Along a bond $(c'',c''+\eta_ke_\nu)$ of the path, traversed in either direction, the value of $\mathrm{rem}[\mathrm{bq}]_\mu$ changes by $\pm\eta_k(\partial_\nu\mathrm{rem}[\mathrm{bq}]_\mu)(c'')$.
The corners $c''$ and $c''+\eta_ke_\nu$ lie between $\tilde{c}$ and $c$, the corners $c''+\eta_ke_\mu$ and $c''+\eta_ke_\nu+\eta_ke_\mu$ lie between $\tilde{c}+\eta_ke_\mu$ and $c+\eta_ke_\mu$, and the four points $\tilde{c}$, $c$, $\tilde{c}+\eta_ke_\mu$, $c+\eta_ke_\mu$ belong to $\mathrm{Kb}$.
Therefore the quotient lies inside $\mathrm{Kb}$ by (bx2), and by (1c) its modulus is at most $(2\,\mathrm{rad}+5)\,\mathrm{bd3}$.
It follows that $|\mathrm{rem}[\mathrm{bq}]_\mu(c)|\le3(2\,\mathrm{rad}+5)\cdot(2\,\mathrm{rad}+5)\,\mathrm{bd3}$.
Finally, $3(2\,\mathrm{rad}+5)^2=12\,\mathrm{rad}^2+60\,\mathrm{rad}+75\le18\,\mathrm{rad}^2$, since $6\,\mathrm{rad}^2\ge60\,\mathrm{rad}+75$ for $\mathrm{rad}\ge12$.
\end{itemize}
The right-hand sides of the three bounds of (qu2) are increasing in $\mathrm{bn}_0$, $\mathrm{bn}_1$, $\mathrm{bn}_\alpha$.
Apply (qu2) with $r=\mathrm{rad}$, and use (1d), (1c), (1b) in the form $\mathrm{bn}_0\le18\,\mathrm{rad}^2\cdot\mathrm{bd3}$, $\mathrm{bn}_1\le3\,\mathrm{rad}\cdot\mathrm{bd3}$, $\mathrm{bn}_\alpha\le\mathrm{bd3}$.
On $T_k$ one obtains
\begin{equation}\tag{18a.4}
\begin{aligned}
|\mathrm{cut}_{\mathrm{rad}}\,\mathrm{rem}[\mathrm{bq}]|&\le18\,\mathrm{rad}^2\cdot\mathrm{bd3},\\
|\partial(\mathrm{cut}_{\mathrm{rad}}\,\mathrm{rem}[\mathrm{bq}])|&\le(3+36)\,\mathrm{rad}\cdot\mathrm{bd3}=39\,\mathrm{rad}\cdot\mathrm{bd3},\\
|\delta_\alpha\partial(\mathrm{cut}_{\mathrm{rad}}\,\mathrm{rem}[\mathrm{bq}])|&\le(1+36+216)\,\mathrm{bd3}=253\,\mathrm{bd3}.
\end{aligned}
\end{equation}

\paragraph{Step 2 (the real gauge function).}
The matrix $\mathrm{Vs}[\mathrm{ba0}']$ is a real $3\times3$ matrix.
Put
\begin{equation}\tag{18a.5}
\begin{aligned}
&\mathrm{sym}_{\nu\mu}:=\frac{1}{2}\left(\mathrm{Vs}[\mathrm{ba0}']_{\nu\mu}+\mathrm{Vs}[\mathrm{ba0}']_{\mu\nu}\right),\qquad
\mathrm{fa}_{\nu\mu}:=\mathrm{Vs}[\mathrm{ba0}']_{\nu\mu}-\mathrm{Vs}[\mathrm{ba0}']_{\mu\nu},\\
&\qquad\text{so that}\quad\mathrm{Vs}[\mathrm{ba0}']=\mathrm{sym}+\frac{1}{2}\,\mathrm{fa},\\
&\mathrm{qg}(c):=\sum_\mu\left((\mathrm{ba0}')_\mu(\tilde{c})-\frac{\eta_k}{2}\,\mathrm{sym}_{\mu\mu}\right)(c-\tilde{c})_\mu+\frac{1}{2}\sum_{\nu,\mu}\mathrm{sym}_{\nu\mu}\,(c-\tilde{c})_\nu(c-\tilde{c})_\mu\qquad(c\in\mathrm{Kb});
\end{aligned}
\end{equation}
the function $\mathrm{qg}$ is real.
Since $\mathrm{sym}$ is symmetric, for $v\in\mathbb{R}^3$ the quadratic form $\frac{1}{2}\sum_{\nu,\mu'}\mathrm{sym}_{\nu\mu'}v_\nu v_{\mu'}$ changes under $v\mapsto v+\eta_ke_\mu$ by $\eta_k\sum_\nu\mathrm{sym}_{\nu\mu}v_\nu+\frac{\eta_k^2}{2}\,\mathrm{sym}_{\mu\mu}$.
Together with (bx1) this gives, for $c,c+\eta_ke_\mu\in\mathrm{Kb}$,
\begin{equation}\tag{18a.6}
\begin{aligned}
(d\,\mathrm{qg})_\mu(c)&=\eta_k^{-1}\left[\mathrm{qg}(c+\eta_ke_\mu)-\mathrm{qg}(c)\right]\\
&=(\mathrm{ba0}')_\mu(\tilde{c})-\frac{\eta_k}{2}\,\mathrm{sym}_{\mu\mu}+\sum_\nu\mathrm{sym}_{\nu\mu}(c-\tilde{c})_\nu+\frac{\eta_k}{2}\,\mathrm{sym}_{\mu\mu}\\
&=(\mathrm{ba0}')_\mu(\tilde{c})+\sum_\nu\mathrm{sym}_{\nu\mu}(c-\tilde{c})_\nu .
\end{aligned}
\end{equation}
By (18a.3), on the bonds inside $\mathrm{Kb}$,
\[
\mathrm{ba0}'=\mathrm{ba0}'(\tilde{c})+\mathrm{Lin}[\mathrm{ba0}']+\mathrm{rem}[\mathrm{ba0}'],\qquad
\mathrm{ba1}=\mathrm{ba1}(\tilde{c})+\mathrm{Lin}[\mathrm{ba1}]+\mathrm{rem}[\mathrm{ba1}],
\]
and
\[
\mathrm{Lin}[\mathrm{ba0}']_\mu(c)=\sum_\nu\mathrm{sym}_{\nu\mu}(c-\tilde{c})_\nu+\sum_\nu\frac{\mathrm{fa}_{\nu\mu}}{2}\,(c-\tilde{c})_\nu .
\]
As $\mathrm{Kb}$ is contained in $\mathrm{Zd}$, the identities (18a.1) and (18a.6) give, on the bonds inside $\mathrm{Kb}$,
\begin{equation}\tag{18a.7}
\begin{aligned}
\mathrm{ba0}+\mathrm{ba1}&=d(\boldsymbol{\lambda}_L+\mathrm{qg})+\mathrm{ba1}(\tilde{c})+\mathrm{Lp}+\mathrm{rem}[\mathrm{ba0}']+\mathrm{rem}[\mathrm{ba1}],\\
\mathrm{Lp}_\mu(c)&:=\sum_\nu\left(\frac{\mathrm{fa}_{\nu\mu}}{2}+\mathrm{Vs}[\mathrm{ba1}]_{\nu\mu}\right)(c-\tilde{c})_\nu .
\end{aligned}
\end{equation}

\paragraph{Step 3 (slopes).}
Every coefficient of $\mathrm{Lp}$ has modulus at most
\[
\mathrm{slp}:=2\min\left(\mathrm{bd2},\ell\,\mathrm{bd3}\right);
\]
this follows from (3a) and (3b).
\begin{itemize}
\item[(3a)] $|\mathrm{Vs}[\mathrm{ba1}]_{\nu\mu}|\le\min(\mathrm{bd2},\ell\,\mathrm{bd3})$.
The bound by $\mathrm{bd2}$ holds because $|\partial\,\mathrm{ba1}|<\mathrm{bd2}$ by (18a.2).
The function $\mathrm{fn}:=\partial_\nu(\mathrm{ba1})_\mu$ on $T_k$ has sum zero, and by (18a.2) one has $|\mathrm{fn}(c)-\mathrm{fn}(c')|\le d(c,c')^\alpha\,\mathrm{bd3}\le\mathrm{bd3}$ whenever $0<d(c,c')\le1$.
So (qu3) with $\mathrm{osc}=\mathrm{bd3}$ gives the bound by $\ell\,\mathrm{bd3}$.
\item[(3b)] $|\mathrm{fa}_{\nu\mu}|\le2\min(\mathrm{bd2},\ell\,\mathrm{bd3})$.
For $\nu=\mu$ one has $\mathrm{fa}_{\mu\mu}=0$.
Let $\nu\ne\mu$, and put
\[
\mathrm{cur}_{\nu\mu}(c):=(\partial_\nu(\mathrm{ba0})_\mu)(c)-(\partial_\mu(\mathrm{ba0})_\nu)(c)\qquad(c\in T_k).
\]
The two quotients have the same cell, and its corners $c$, $c+\eta_ke_\mu$, $c+\eta_ke_\nu$, $c+\eta_ke_\mu+\eta_ke_\nu$ have sup distance at most $\eta_k\le1$ from $c$.
Each of the two terms has sum zero over $T_k$ (qu3).
For every function $f$ of the sites of a set that contains this cell,
\begin{equation}\tag{md}
(\partial_\nu(df)_\mu)(c)=\eta_k^{-2}\left[f(c+\eta_ke_\mu+\eta_ke_\nu)-f(c+\eta_ke_\nu)-f(c+\eta_ke_\mu)+f(c)\right]=(\partial_\mu(df)_\nu)(c).
\end{equation}
Let $d(c,c')\le1$.
Let $L\square''$ be the $LM$-cube of $c$ and $z''$ its centre, let $(A'',\boldsymbol{\lambda}'')$ be a representative of $A_0$ on $\square''^{\dagger}$ as in \cite[Definition~5.16]{QEDI}, and put $\mathrm{ba0}'':=(A'')_L$.
Since $|c-z''|<LM/2$, the corners of the cells at $c$ and at $c'$ have sup distance less than $LM/2+1+1\le\left(L+\frac{1}{2}\right)LM$ from $z''$.
So both cells lie inside $L\square''^{\dagger}$, which does not wrap by (ge1).
On this set $\mathrm{ba0}=\mathrm{ba0}''+d(\boldsymbol{\lambda}''_L)$ by Proposition~18.1\,(c), and (md) with $f=\boldsymbol{\lambda}''_L$ shows that $\mathrm{cur}_{\nu\mu}=\partial_\nu(\mathrm{ba0}'')_\mu-\partial_\mu(\mathrm{ba0}'')_\nu$ at $c$ and at $c'$.
By Proposition~18.1\,(c), (a) and $L^{-\alpha}\le1$, one obtains $|\mathrm{cur}_{\nu\mu}(c)|<2\,\mathrm{bd2}$ and, for $c'\ne c$,
\[
\left|\mathrm{cur}_{\nu\mu}(c)-\mathrm{cur}_{\nu\mu}(c')\right|\le2\,d(c,c')^\alpha\,\mathrm{bd3}\le2\,\mathrm{bd3} .
\]
By (qu3) with $\mathrm{fn}=\mathrm{cur}_{\nu\mu}$ and $\mathrm{osc}=2\,\mathrm{bd3}$ it follows that $|\mathrm{cur}_{\nu\mu}|_\infty\le2\ell\,\mathrm{bd3}$.
The cell at $\tilde{c}$ lies inside $\mathrm{Zd}$ by (ge2), and there $\mathrm{ba0}=\mathrm{ba0}'+d(\boldsymbol{\lambda}_L)$ by (18a.1); so (md) with $f=\boldsymbol{\lambda}_L$ gives
\[
\mathrm{cur}_{\nu\mu}(\tilde{c})=(\partial_\nu(\mathrm{ba0}')_\mu)(\tilde{c})-(\partial_\mu(\mathrm{ba0}')_\nu)(\tilde{c})=\mathrm{Vs}[\mathrm{ba0}']_{\nu\mu}-\mathrm{Vs}[\mathrm{ba0}']_{\mu\nu}=\mathrm{fa}_{\nu\mu}.
\]
Hence $|\mathrm{fa}_{\nu\mu}|=|\mathrm{cur}_{\nu\mu}(\tilde{c})|\le2\min(\mathrm{bd2},\ell\,\mathrm{bd3})$.
\end{itemize}
Consequently $\left|\mathrm{fa}_{\nu\mu}/2+\mathrm{Vs}[\mathrm{ba1}]_{\nu\mu}\right|\le2\min(\mathrm{bd2},\ell\,\mathrm{bd3})=\mathrm{slp}$.

\paragraph{Step 4 (definition).}
Let $\mathrm{rad}':=\min\left(e_k^{-\varepsilon},\ell/5\right)$.
Then
\begin{equation}\tag{18a.8}
\begin{aligned}
&\mathrm{rad}'\ge\mathrm{rad}\ge50\qquad\left(2LM\,e_k^{\varepsilon}\le1\ \text{by (nm6)};\ \ \ell/5\ge13LM/5\ge2LM\right),\\
&4\,\mathrm{rad}'+11\le\frac{4\ell}{5}+11<\ell\qquad(\ell\ge13LM>55).
\end{aligned}
\end{equation}
Hence (qu2) applies with $r=\mathrm{rad}'$, Section~0.3 applies to $K(\tilde{c},2\,\mathrm{rad}'+5)$, and the formula for $\mathrm{Lp}$ in (18a.7) defines a bond function on the bonds inside $K(\tilde{c},2\,\mathrm{rad}'+5)$.
Put
\begin{equation}\tag{18a.9}
\mathrm{bh}:=\mathrm{cut}_{\mathrm{rad}}\left(\mathrm{rem}[\mathrm{ba0}']+\mathrm{rem}[\mathrm{ba1}]\right)+\mathrm{cut}_{\mathrm{rad}'}\,\mathrm{Lp},\qquad
\mathrm{kc}_1:=\mathrm{ba1}(\tilde{c}),
\end{equation}
where $\mathrm{kc}_1$ is the constant bond function with the components $(\mathrm{ba1})_\mu(\tilde{c})=((A_1)_L)_\mu(\tilde{c})$, and
\begin{equation}\tag{18a.10}
\mathrm{mg}:=\boldsymbol{\lambda}_L+\mathrm{qg}\ \text{ on the sites of }\mathrm{Kb},\qquad
\mathrm{mg}:=0\ \text{ on the other sites of }T_k.
\end{equation}
Thus $\mathrm{bh}$ is a complex bond function on $T_k$, $\mathrm{kc}_1$ is a constant complex bond function, and $\mathrm{mg}$ is a real function of the sites of $T_k$.

\paragraph{Step 5 (norms of $\mathrm{bh}$).}
By Step 3, on $K(\tilde{c},2\,\mathrm{rad}'+5)$ one has
\[
|\mathrm{Lp}|\le3(2\,\mathrm{rad}'+5)\,\mathrm{slp}\le6.3\,\mathrm{rad}'\cdot\mathrm{slp}\quad(15\le0.3\,\mathrm{rad}'),\qquad
|\partial\,\mathrm{Lp}|\le\mathrm{slp},\qquad
|\delta_\alpha\partial\,\mathrm{Lp}|=0;
\]
for the second bound observe that, by (bx1), $\partial_\nu\mathrm{Lp}_\mu$ equals the coefficient $\mathrm{fa}_{\nu\mu}/2+\mathrm{Vs}[\mathrm{ba1}]_{\nu\mu}$.
By (qu2) with $r=\mathrm{rad}'$, on $T_k$:
\begin{equation}\tag{18a.11}
\begin{aligned}
|\mathrm{cut}_{\mathrm{rad}'}\,\mathrm{Lp}|&\le6.3\,\mathrm{rad}'\cdot\mathrm{slp},\\
|\partial(\mathrm{cut}_{\mathrm{rad}'}\,\mathrm{Lp})|&\le(1+12.6)\,\mathrm{slp}=13.6\,\mathrm{slp},\\
|\delta_\alpha\partial(\mathrm{cut}_{\mathrm{rad}'}\,\mathrm{Lp})|&\le(12+75.6)\,\frac{\mathrm{slp}}{\mathrm{rad}'}=87.6\,\frac{\mathrm{slp}}{\mathrm{rad}'}.
\end{aligned}
\end{equation}

\emph{Case $\mathrm{rad}'=e_k^{-\varepsilon}$.}
In this case one uses $\mathrm{slp}\le2\,\mathrm{bd2}=2L^{-\frac{15}{8}+\varepsilon}\cdot\mathrm{t2}_k$ together with $e_k^{-\varepsilon}\cdot\mathrm{t2}_k=\mathrm{t1}_k$ and $e_k^{\varepsilon}\cdot\mathrm{t2}_k=\mathrm{t3}_k$; the three right-hand sides of (18a.11) are then at most
\[
12.6\,L^{-\frac{15}{8}+\varepsilon}\cdot\mathrm{t1}_k,\qquad
27.2\,L^{-\frac{15}{8}+\varepsilon}\cdot\mathrm{t2}_k,\qquad
175.2\,L^{-\frac{15}{8}+\varepsilon}\cdot\mathrm{t3}_k.
\]

\emph{Case $\mathrm{rad}'=\ell/5<e_k^{-\varepsilon}$.}
With $\mathrm{slp}\le2\ell\,\mathrm{bd3}$ and $\ell<5e_k^{-\varepsilon}$ the three right-hand sides of (18a.11) are at most
\begin{gather*}
6.3\cdot\frac{\ell}{5}\cdot2\ell\,\mathrm{bd3}=2.52\,\ell^2\,\mathrm{bd3}\le63\,e_k^{-2\varepsilon}\,\mathrm{bd3}=63\,L^{-\frac{15}{8}+\frac{3\varepsilon}{2}}\cdot\mathrm{t1}_k,\\
27.2\,\ell\,\mathrm{bd3}\le136\,e_k^{-\varepsilon}\,\mathrm{bd3}=136\,L^{-\frac{15}{8}+\frac{3\varepsilon}{2}}\cdot\mathrm{t2}_k,\\
87.6\cdot2\ell\,\mathrm{bd3}\cdot\frac{5}{\ell}=876\,\mathrm{bd3}=876\,L^{-\frac{15}{8}+\frac{3\varepsilon}{2}}\cdot\mathrm{t3}_k.
\end{gather*}

\emph{The first term of $\mathrm{bh}$.}
The map $\mathrm{cut}_{\mathrm{rad}}$ is linear.
By (18a.4) for $\mathrm{bq}=\mathrm{ba0}'$ and for $\mathrm{bq}=\mathrm{ba1}$, and by $\mathrm{rad}=2LM$, the three norms of $\mathrm{cut}_{\mathrm{rad}}(\mathrm{rem}[\mathrm{ba0}']+\mathrm{rem}[\mathrm{ba1}])$ are at most
\begin{gather*}
36\,\mathrm{rad}^2\cdot\mathrm{bd3}=144\left(LM\,e_k^{\varepsilon}\right)^2L^{-\frac{15}{8}+\frac{3\varepsilon}{2}}\cdot\mathrm{t1}_k,\qquad
78\,\mathrm{rad}\cdot\mathrm{bd3}=156\,LM\,e_k^{\varepsilon}\,L^{-\frac{15}{8}+\frac{3\varepsilon}{2}}\cdot\mathrm{t2}_k,\\
506\,\mathrm{bd3}=506\,L^{-\frac{15}{8}+\frac{3\varepsilon}{2}}\cdot\mathrm{t3}_k;
\end{gather*}
by (nm6), $144\left(LM\,e_k^{\varepsilon}\right)^2\le144/160^2\le1$ and $156\,LM\,e_k^{\varepsilon}\le156/160\le1$.

\emph{The sum.}
Since $L^{-\frac{15}{8}+\varepsilon}\le L^{-\frac{15}{8}+\frac{3\varepsilon}{2}}$, in both cases
\begin{equation}\tag{18a.12}
\begin{aligned}
|\mathrm{bh}|&\le64\,L^{-\frac{15}{8}+\frac{3\varepsilon}{2}}\cdot\mathrm{t1}_k,\\
|\partial\,\mathrm{bh}|&\le137\,L^{-\frac{15}{8}+\frac{3\varepsilon}{2}}\cdot\mathrm{t2}_k,\\
|\delta_\alpha\partial\,\mathrm{bh}|&\le1382\,L^{-\frac{15}{8}+\frac{3\varepsilon}{2}}\cdot\mathrm{t3}_k,
\end{aligned}
\end{equation}
where $64=63+1$, $137=136+1$ and $1382=876+506$.

\paragraph{Step 6 (second slot).}
In the second slot nothing has to be chosen.
Put
\begin{equation}\tag{18a.13}
\mathrm{kc}_2:=\text{the constant bond function with the components }(\mathscr{C}_L)_\mu(\tilde{c}),\qquad
\mathscr{C}^b:=\mathrm{cut}_{\mathrm{rad}}\left(\mathscr{C}_L-\mathrm{kc}_2\right).
\end{equation}
For $c,c+\eta_ke_\mu\in\mathrm{Kb}$ the difference $(\mathscr{C}_L)_\mu(c)-(\mathscr{C}_L)_\mu(\tilde{c})$ is the sum of the terms $\pm\eta_k(\partial_\nu(\mathscr{C}_L)_\mu)(c'')$ over the at most $3(2\,\mathrm{rad}+5)/\eta_k$ bonds $(c'',c''+\eta_ke_\nu)$ of a rectilinear path from $\tilde{c}$ to $c$.
Therefore $|\mathscr{C}_L-\mathrm{kc}_2|_{\mathrm{Kb}}\le3(2\,\mathrm{rad}+5)\,|\partial\mathscr{C}_L|\le6.3\,\mathrm{rad}\cdot|\partial\mathscr{C}_L|$, because $\mathrm{rad}\ge50$, and $|\partial(\mathscr{C}_L-\mathrm{kc}_2)|_{\mathrm{Kb}}\le|\partial\mathscr{C}_L|$.
The first two bounds of (qu2) with $r=\mathrm{rad}$ (the number $\mathrm{bn}_\alpha$ is finite and does not enter them), Proposition~18.1\,(b), $\mathrm{rad}=2LM$ and $\mathrm{td2}_k=e_k^{\varepsilon}\,\mathrm{td1}_k$ now give
\begin{equation}\tag{18a.14}
\begin{aligned}
|\mathscr{C}^b|&\le6.3\,\mathrm{rad}\cdot|\partial\mathscr{C}_L|<12.6\,LM\,L^{-2+\frac{3\varepsilon}{2}}\cdot\mathrm{td2}_k=12.6\,M\,L^{-1+\frac{3\varepsilon}{2}}\,e_k^{\varepsilon}\,\mathrm{td1}_k\\
&=\frac{12.6}{160}\,L^{-\frac{1}{4}}\left[160\,L^{1+\frac{3\varepsilon}{2}}M\,e_k^{\varepsilon}\right]L^{-\frac{7}{4}}\cdot\mathrm{td1}_k\le L^{-\frac{7}{4}}\cdot\mathrm{td1}_k\qquad\text{(nm6)},\\
|\partial\mathscr{C}^b|&\le(1+12.6)\,|\partial\mathscr{C}_L|<13.6\,L^{-2+\frac{3\varepsilon}{2}}\cdot\mathrm{td2}_k\le L^{-\frac{7}{4}}\cdot\mathrm{td2}_k\qquad\text{(nm3)}.
\end{aligned}
\end{equation}

\subsubsection*{2.3. Proof of (o), (i), (ii)}

\begin{proof}[Proof of Proposition~18.2\,(o)]
The right-hand sides of (18a.9) and (18a.10) are built from $\tilde{c}$, $\mathrm{ba0}'=(A')_L$, $\mathrm{ba1}=(A_1)_L$, $\boldsymbol{\lambda}_L$ and the numbers $\mathrm{rad}=2LM$, $\mathrm{rad}'=\min(e_k^{-\varepsilon},\ell/5)$, and from nothing else: the objects $\mathrm{Vs}$, $\mathrm{Lin}$, $\mathrm{rem}$, $\mathrm{sym}$, $\mathrm{fa}$, $\mathrm{qg}$, $\mathrm{Lp}$ are formed from $\mathrm{ba0}'$, $\mathrm{ba1}$ and $\tilde{c}$ by (18a.3), (18a.5), (18a.7).
The triple $(A_1,A',\boldsymbol{\lambda})$ consists of a splitting of $\mathscr{A}$ and a representative on $\square'^{\dagger}$, where $L\square'$ is the $LM$-cube of $\tilde{c}$.
(The representatives of $A_0$ on the other cubes enter only the estimate (3b), not the construction.)
The right-hand sides of (18a.13) are built from $\mathscr{C}$, $\tilde{c}$ and $\mathrm{rad}$ only.
No family of activities enters Steps 0--6, and of the polymer $X$ only the site $\tilde{c}$ enters.
The bond function $\mathrm{kc}_2$ is the constant field described in (o).
If $\mathscr{C}=0$, then $\mathscr{C}_L=0$, $\mathrm{kc}_2=0$ and $\mathscr{C}^b=\mathrm{cut}_{\mathrm{rad}}0=0$.
\end{proof}

\begin{proof}[Proof of Proposition~18.2\,(i)]
The ball $K(\tilde{c},\mathrm{rad})$ is contained both in $K(\tilde{c},\mathrm{rad}')$ and in $\mathrm{Kb}$.
On the bonds inside $K(\tilde{c},\mathrm{rad})$, (qu2) gives $\mathrm{cut}_{\mathrm{rad}}(\mathrm{rem}[\mathrm{ba0}']+\mathrm{rem}[\mathrm{ba1}])=\mathrm{rem}[\mathrm{ba0}']+\mathrm{rem}[\mathrm{ba1}]$ and $\mathrm{cut}_{\mathrm{rad}'}\,\mathrm{Lp}=\mathrm{Lp}$; moreover $d\,\mathrm{mg}=d(\boldsymbol{\lambda}_L+\mathrm{qg})$ there, because both end points of such a bond lie in $\mathrm{Kb}$.
So (18a.1) and (18a.7) yield, on these bonds,
\[
\mathscr{A}_L-d\,\mathrm{mg}=\mathrm{ba0}+\mathrm{ba1}-d(\boldsymbol{\lambda}_L+\mathrm{qg})=\mathrm{kc}_1+\mathrm{Lp}+\mathrm{rem}[\mathrm{ba0}']+\mathrm{rem}[\mathrm{ba1}]=\mathrm{bh}+\mathrm{kc}_1 .
\]
By (qu2) with $r=\mathrm{rad}$ and $\mathrm{bq}=\mathscr{C}_L-\mathrm{kc}_2$, one has $\mathscr{C}^b=\mathrm{cut}_{\mathrm{rad}}(\mathscr{C}_L-\mathrm{kc}_2)=\mathscr{C}_L-\mathrm{kc}_2$ on the bonds inside $K(\tilde{c},\mathrm{rad})$.
The last assertion of (i) is (ge4).
\end{proof}

\begin{proof}[Proof of Proposition~18.2\,(ii)]
The bounds on $\mathrm{bh}$ are those of (18a.12), since $64\le1382$ and $137\le1382$.
The bounds on $\mathscr{C}^b$ are those of (18a.14).
Let $t$ be a complex number with $|t|\le L^{\frac{7}{4}}$.
By (nm2), the three norms of $t\cdot\mathrm{bh}$ are at most $L^{\frac{7}{4}}\cdot1382\,L^{-\frac{15}{8}+\frac{3\varepsilon}{2}}$ times $\mathrm{t1}_k$, $\mathrm{t2}_k$, $\mathrm{t3}_k$, respectively, and this is less than $\mathrm{t1}_k$, $\mathrm{t2}_k$, $\mathrm{t3}_k$, respectively; moreover $|t\cdot\mathscr{C}^b|<\mathrm{td1}_k$ and $|\partial(t\cdot\mathscr{C}^b)|<\mathrm{td2}_k$.
Thus $(t\cdot\mathrm{bh},t\cdot\mathscr{C}^b)$ lies in $R_k\times\mathfrak{X}_k$.
Finally, $|\mathrm{kc}_1|\le|\mathrm{ba1}|<\mathrm{bd1}=L^{-\frac{7}{8}+\frac{\varepsilon}{2}}\cdot\mathrm{t1}_k$ by (18a.2), and $|\mathrm{kc}_2|\le|\mathscr{C}_L|<L^{-1+\varepsilon}\cdot\mathrm{td1}_k$ by Proposition~18.1\,(b).
\end{proof}

\subsubsection*{2.4. Step 7 (phases): proof of (iii)}

\begin{proof}[Proof of Proposition~18.2\,(iii)]
Let $E$ be a family of level $k$ in the sense of (G2.r8) with the properties (I.k1), (I.k2), (I.k4).
Denote by $E_{nm}(X;\cdot,\cdot;\xi;\boldsymbol{\zeta})$ the antisymmetric kernels of the elements $E(X;\cdot,\cdot)$ of $\mathrm{Gr}[\psi,\chi]$; by Proposition~15\,(a) they are determined by the elements.
The arguments are $\xi=(x,\beta,\omega)=(c,g,\beta,\omega)$ with $x\in T^{\mathrm{f}}_k$, and $\boldsymbol{\zeta}=(x,x',\beta,\omega)$ with $(x,x')\in\mathrm{Pz}_k$ (Section~I.3.6).
For an argument $\mathrm{ag}$, that is, for one of the $\xi_i$ or one of the $\boldsymbol{\zeta}_j$, put
\[
c(\mathrm{ag}):=\pi(x),\qquad\omega(\mathrm{ag}):=\omega,
\]
where $\pi(x)$ is the photon site under the site $x$ of $\xi$, respectively under the first site $x$ of $\boldsymbol{\zeta}$.

\emph{(KG) Kernels under $G_\Lambda$.}
Let $F$ be an element of $\mathrm{Gr}[\psi,\chi]$ of level $k$ with the antisymmetric kernels $F_{nm}$, and let $\Lambda$ be a real function of the sites of $T_k$.
By (I.T-G), $G_\Lambda$ multiplies the generator with the index $\mathrm{ag}$ by $\exp\left\{(-1)^{\omega(\mathrm{ag})}ie_k\Lambda(c(\mathrm{ag}))\right\}$; the 48 sites of a cluster carry the same phase.
The map $G_\Lambda$ is a homomorphism.
Write
\[
F=\sum_{n,m}\frac{1}{n!\,m!}\int\mathrm{d}\xi\,\mathrm{d}\boldsymbol{\zeta}\,F_{nm}(\xi;\boldsymbol{\zeta})\,\psi(\xi_1)\cdots\psi(\xi_n)\,\chi(\boldsymbol{\zeta}_1)\cdots\chi(\boldsymbol{\zeta}_m),
\]
where $\int\mathrm{d}\xi$ is the sum over $\xi$ with the weight $\eta_k^3/48$, extended over $c\in T_k$, the 48 labels $g\in O_{\mathrm{cub}}$, $\beta\in\{1,\dots,4\}$ and $\omega\in\{0,1\}$, and where $\int\mathrm{d}\boldsymbol{\zeta}$ is the sum over $\boldsymbol{\zeta}$ with the weight $(\eta_k^3/48)^2$, extended over the pairs of $\mathrm{Pz}_k$, $\beta$ and $\omega$.
Applying $G_\Lambda$ to $F$ gives the same expression with $F_{nm}(\xi;\boldsymbol{\zeta})$ replaced by
\[
F_{nm}(\xi;\boldsymbol{\zeta})\prod_{\mathrm{ag}}\exp\left\{(-1)^{\omega(\mathrm{ag})}ie_k\Lambda(c(\mathrm{ag}))\right\},
\]
the product extending over the $n+m$ arguments.
This product is symmetric under permutations of the $\xi_i$ and of the $\boldsymbol{\zeta}_j$, so that the new kernel is antisymmetric in the $\xi_i$ and in the $\boldsymbol{\zeta}_j$; and since $\Lambda$ is $\ell\mathbb{Z}^3$-periodic, the new kernel is antiperiodic in each argument (I.AP).
By Proposition~15\,(a),
\begin{equation}\tag{18a.15}
(G_\Lambda F)_{nm}(\xi;\boldsymbol{\zeta})=F_{nm}(\xi;\boldsymbol{\zeta})\prod_{\mathrm{ag}}\exp\left\{(-1)^{\omega(\mathrm{ag})}ie_k\Lambda(c(\mathrm{ag}))\right\}.
\end{equation}

\emph{(7a) Points of the domain.}
By Proposition~18.1\,(b), (c), the pair $(\mathscr{A}_L,\mathscr{C}_L)$ lies in $\tilde{R}_k\times\mathfrak{X}_k$.
The set $R_k$ is contained in $\tilde{R}_k$ (\cite[Definition~5.16]{QEDI} with $A_0=0$, $A'=0$, $\boldsymbol{\lambda}=0$).
By (ii) with $t=1$, the pair $(\mathrm{bh},\mathscr{C}^b)$ lies in $R_k\times\mathfrak{X}_k$.
A constant bond function has vanishing difference quotients.
Hence (ii), (nm2) and (nm4) give
\[
|\mathrm{bh}+\mathrm{kc}_1|<\left(L^{-\frac{7}{4}}+L^{-\frac{7}{8}+\frac{\varepsilon}{2}}\right)\mathrm{t1}_k<\mathrm{t1}_k,\qquad\partial(\mathrm{bh}+\mathrm{kc}_1)=\partial\,\mathrm{bh},
\]
so that $\mathrm{bh}+\mathrm{kc}_1$ lies in $R_k$, and
\[
|\mathscr{C}^b+\mathrm{kc}_2|<\left(L^{-\frac{7}{4}}+L^{-1+\varepsilon}\right)\mathrm{td1}_k<\mathrm{td1}_k,\qquad\partial(\mathscr{C}^b+\mathrm{kc}_2)=\partial\mathscr{C}^b,
\]
so that $\mathscr{C}^b+\mathrm{kc}_2$ lies in $\mathfrak{X}_k$.

\emph{(7b) The real gauge function.}
By (i), $\mathrm{bh}+\mathrm{kc}_1=\mathscr{A}_L-d\,\mathrm{mg}$ and $\mathscr{C}^b+\mathrm{kc}_2=\mathscr{C}_L-d0$ on all bonds that have an end point at distance at most $1$ from $X$, in particular on the bonds at distance at most $1$ from $X$ (Section~I.1.1).
The functions $\mathrm{mg}$ and $0$ are real functions of $c$, and both pairs lie in $\tilde{R}_k\times\mathfrak{X}_k$ by (7a).
Apply (I.k4) with $(\mathscr{A}_L,\mathscr{C}_L)$ as the pair without index, $(\mathrm{bh}+\mathrm{kc}_1,\mathscr{C}^b+\mathrm{kc}_2)$ as the pair with the index $1$, $\boldsymbol{\lambda}=\mathrm{mg}$ and $\boldsymbol{\lambda}'=0$; it gives
\[
G_{\mathrm{mg}}\left[E(X;\mathrm{bh}+\mathrm{kc}_1,\mathscr{C}^b+\mathrm{kc}_2)\right]=E(X;\mathscr{A}_L,\mathscr{C}_L).
\]
By (18a.15), for all $(n,m)$ and all arguments,
\begin{equation}\tag{18a.16}
E_{nm}(X;\mathscr{A}_L,\mathscr{C}_L;\xi;\boldsymbol{\zeta})=E_{nm}(X;\mathrm{bh}+\mathrm{kc}_1,\mathscr{C}^b+\mathrm{kc}_2;\xi;\boldsymbol{\zeta})\prod_{\mathrm{ag}}\exp\left\{(-1)^{\omega(\mathrm{ag})}ie_k\,\mathrm{mg}(c(\mathrm{ag}))\right\}.
\end{equation}
Since $\mathrm{mg}$ is real, every factor has modulus one.
For $n=m=0$ there is no factor.

\emph{(7c) The constants.}
Write $\mathrm{kc}_j=\mathrm{kc}_j^{\mathrm{R}}+i\,\mathrm{kc}_j^{\mathrm{I}}$ with real constant bond functions $\mathrm{kc}_j^{\mathrm{R}}$, $\mathrm{kc}_j^{\mathrm{I}}$ ($j=1,2$); then $|\mathrm{kc}_j^{\mathrm{R}}|\le|\mathrm{kc}_j|$ and $|\mathrm{kc}_j^{\mathrm{I}}|\le|\mathrm{kc}_j|$.
Throughout this step $z=(z_1,\dots,z_4)$ denotes a point of $\mathbb{C}^4$.
Put
\begin{gather*}
\mathrm{Dz}:=\left\{z\in\mathbb{C}^4:|\operatorname{Im}z_1|<2,\ |\operatorname{Im}z_2|<2,\ |z_3|<2,\ |z_4|<2\right\},\\
\mathscr{A}(z):=\mathrm{bh}+z_1\mathrm{kc}_1^{\mathrm{R}}+z_2\mathrm{kc}_1^{\mathrm{I}},\qquad
\mathscr{C}(z):=\mathscr{C}^b+z_3\mathrm{kc}_2^{\mathrm{R}}+z_4\mathrm{kc}_2^{\mathrm{I}},\\
\Lambda_z(c):=-\left(z_1\mathrm{kc}_1^{\mathrm{R}}+z_2\mathrm{kc}_1^{\mathrm{I}}+z_3\mathrm{kc}_2^{\mathrm{R}}+z_4\mathrm{kc}_2^{\mathrm{I}}\right)\cdot(c-\tilde{c})\ \text{ for }c\in K(\tilde{c},\mathrm{rad}),\qquad
\Lambda_z(c):=0\ \text{ for the other sites of }T_k.
\end{gather*}
Here $\mathscr{A}(z)$ and $\mathscr{C}(z)$ are bond functions on $T_k$; the data $\mathscr{A}$, $\mathscr{C}$ of level $k+1$ are written without an argument.
The set $\mathrm{Dz}$ is open and convex, hence connected; the map $z\mapsto(\mathscr{A}(z),\mathscr{C}(z))$ is affine, and $z\mapsto\Lambda_z(c)$ is linear.

\emph{The pair $(\mathscr{A}(z),\mathscr{C}(z))$ lies in $\tilde{R}_k\times\mathfrak{X}_k$ for $z\in\mathrm{Dz}$.}
One has
\[
\mathscr{A}(z)=\left[\mathrm{bh}+i\left(\operatorname{Im}z_1\,\mathrm{kc}_1^{\mathrm{R}}+\operatorname{Im}z_2\,\mathrm{kc}_1^{\mathrm{I}}\right)\right]+\left[\operatorname{Re}z_1\,\mathrm{kc}_1^{\mathrm{R}}+\operatorname{Re}z_2\,\mathrm{kc}_1^{\mathrm{I}}\right].
\]
By (ii), (nm2) and (nm4), the first bracket has sup norm less than $\left(L^{-\frac{7}{4}}+4L^{-\frac{7}{8}+\frac{\varepsilon}{2}}\right)\mathrm{t1}_k<\mathrm{t1}_k$, and its difference quotients are those of $\mathrm{bh}$; so it lies in $R_k$, and hence in $\tilde{R}_k$.
The second bracket is a real constant bond function, and $\tilde{R}_k$ is invariant under the addition of real constant bond functions (Proposition~16.1\,(c) with $\mathrm{cf}=1$).
Moreover $|\mathscr{C}(z)|<\left(L^{-\frac{7}{4}}+4L^{-1+\varepsilon}\right)\mathrm{td1}_k<\mathrm{td1}_k$ and $\partial\mathscr{C}(z)=\partial\mathscr{C}^b$, so that $\mathscr{C}(z)$ lies in $\mathfrak{X}_k$.

\emph{The function $\mathrm{Kz}$.}
Fix $(n,m)$ and arguments $(\xi;\boldsymbol{\zeta})$, and put, for $z\in\mathrm{Dz}$,
\[
\mathrm{Kz}(z):=E_{nm}(X;\mathscr{A}(z),\mathscr{C}(z);\xi;\boldsymbol{\zeta})\prod_{\mathrm{ag}}\exp\left\{(-1)^{\omega(\mathrm{ag})}ie_k\Lambda_z(c(\mathrm{ag}))\right\}.
\]
The set $\tilde{R}_k\times\mathfrak{X}_k$ is open in the space of the pairs of complex bond functions \cite[Lemma~17.3]{QEDI}, the kernel is analytic on it by (I.k1), and the exponentials are entire functions of $z$.
Hence $\mathrm{Kz}$ is analytic on $\mathrm{Dz}$.

\emph{Real $z$.}
Let $z\in\mathrm{Dz}\cap\mathbb{R}^4$.
Put
\[
\boldsymbol{\lambda}(c):=-\left(z_1\mathrm{kc}_1^{\mathrm{R}}+z_2\mathrm{kc}_1^{\mathrm{I}}\right)\cdot(c-\tilde{c}),\qquad
\boldsymbol{\lambda}'(c):=-\left(z_3\mathrm{kc}_2^{\mathrm{R}}+z_4\mathrm{kc}_2^{\mathrm{I}}\right)\cdot(c-\tilde{c})\qquad\text{for }c\in K(\tilde{c},\mathrm{rad}),
\]
and $\boldsymbol{\lambda}:=\boldsymbol{\lambda}':=0$ on the other sites.
These functions are real, $\boldsymbol{\lambda}+\boldsymbol{\lambda}'=\Lambda_z$, and by (bx1), on the bonds inside $K(\tilde{c},\mathrm{rad})$,
\[
(d\boldsymbol{\lambda})_\mu=-\left(z_1\mathrm{kc}_1^{\mathrm{R}}+z_2\mathrm{kc}_1^{\mathrm{I}}\right)_\mu,\qquad
(d\boldsymbol{\lambda}')_\mu=-\left(z_3\mathrm{kc}_2^{\mathrm{R}}+z_4\mathrm{kc}_2^{\mathrm{I}}\right)_\mu,
\]
that is, $\mathscr{A}(z)=\mathrm{bh}-d\boldsymbol{\lambda}$ and $\mathscr{C}(z)=\mathscr{C}^b-d\boldsymbol{\lambda}'$ on these bonds.
By (ge4), the bonds at distance at most $1$ from $X$ are among them.
Now (I.k4), applied with $(\mathrm{bh},\mathscr{C}^b)$ as the pair without index and $(\mathscr{A}(z),\mathscr{C}(z))$ as the pair with the index $1$, gives
\[
G_{\Lambda_z}\left[E(X;\mathscr{A}(z),\mathscr{C}(z))\right]=E(X;\mathrm{bh},\mathscr{C}^b).
\]
By (18a.15) this says $\mathrm{Kz}(z)=E_{nm}(X;\mathrm{bh},\mathscr{C}^b;\xi;\boldsymbol{\zeta})=\mathrm{Kz}(0)$.

\emph{Continuation.}
The function $\mathrm{Kz}-\mathrm{Kz}(0)$ is analytic on the connected open set $\mathrm{Dz}$, and it vanishes on $\mathrm{Dz}\cap\mathbb{R}^4=\mathbb{R}^2\times(-2,2)^2$, which is a non-empty open subset of $\mathbb{R}^4$.
Hence all its partial derivatives at the point $0$ vanish, and its power series at $0$ is zero.
By the identity theorem for analytic functions of several complex variables (\cite[Section~13]{QEDI}; \cite[Chapter~I, Section~A]{GunningRossi1965}), $\mathrm{Kz}-\mathrm{Kz}(0)$ vanishes on $\mathrm{Dz}$.

\emph{Evaluation.}
The point $\mathrm{zs}:=(1,i,1,i)$ lies in $\mathrm{Dz}$, since its coordinates satisfy $\operatorname{Im}z_1=0$, $\operatorname{Im}z_2=1$ and $|z_3|=|z_4|=1$.
One has $\mathscr{A}(\mathrm{zs})=\mathrm{bh}+\mathrm{kc}_1$, $\mathscr{C}(\mathrm{zs})=\mathscr{C}^b+\mathrm{kc}_2$, and $\Lambda_{\mathrm{zs}}(c)=-(\mathrm{kc}_1+\mathrm{kc}_2)\cdot(c-\tilde{c})$ on $K(\tilde{c},\mathrm{rad})$.
So the identity $\mathrm{Kz}(\mathrm{zs})=\mathrm{Kz}(0)$ reads
\begin{equation}\tag{18a.17}
E_{nm}(X;\mathrm{bh}+\mathrm{kc}_1,\mathscr{C}^b+\mathrm{kc}_2;\xi;\boldsymbol{\zeta})\prod_{\mathrm{ag}}\exp\left\{(-1)^{\omega(\mathrm{ag})}ie_k\Lambda_{\mathrm{zs}}(c(\mathrm{ag}))\right\}=E_{nm}(X;\mathrm{bh},\mathscr{C}^b;\xi;\boldsymbol{\zeta}).
\end{equation}

\emph{(7d) The moduli.}
\emph{Position of the arguments.}
By (I.k2), the kernels in (18a.16) and (18a.17) vanish unless every $\xi_i$ lies over $X$ and every pair of the $\boldsymbol{\zeta}_j$ is admissible for $X$ (Section~I.1.7).
Assume that this is the case.
For $\mathrm{ag}=\xi_i$ the site $c(\mathrm{ag})$ lies in $X$.
For $\mathrm{ag}=\boldsymbol{\zeta}_j=(x,x',\cdot)$ of kind (I), the site $c(\mathrm{ag})=c$ lies in $X$; this includes the case of a pair inside a cluster, $c=c'$.
For kind (II) one has $x'=x^{(3)}$ and $x=(c,g)$, and $c$ or $c_3=c+\eta_kg\,e_1$ lies in $X$; so $d(c(\mathrm{ag}),X)\le\eta_k\le1$.
In all cases (qu4) gives $d(c(\mathrm{ag}),\tilde{c})\le M(L+2)+1\le\mathrm{rad}$, and
\[
|c(\mathrm{ag})-\tilde{c}|_1\le3\left(M(L+2)+1\right)\le9LM\qquad(2M+1\le2LM).
\]

\emph{The factors.}
One has
\[
\left|\exp\left\{(-1)^{\omega(\mathrm{ag})}ie_k\Lambda_{\mathrm{zs}}(c(\mathrm{ag}))\right\}\right|=\exp\left\{-(-1)^{\omega(\mathrm{ag})}e_k\operatorname{Im}\Lambda_{\mathrm{zs}}(c(\mathrm{ag}))\right\},\qquad
\operatorname{Im}\Lambda_{\mathrm{zs}}(c)=-\mathrm{kc}_1^{\mathrm{I}}\cdot(c-\tilde{c})-\mathrm{kc}_2^{\mathrm{I}}\cdot(c-\tilde{c}).
\]
By (ii), and since $e_k\cdot\mathrm{t1}_k=e_k^{\frac{1}{4}+\varepsilon}$ and $e_k\cdot\mathrm{td1}_k=e_k^{2\varepsilon}$,
\begin{gather*}
e_k\left|\mathrm{kc}_1^{\mathrm{I}}\cdot(c(\mathrm{ag})-\tilde{c})\right|\le e_k|\mathrm{kc}_1|\cdot9LM\le9LM\,L^{-\frac{7}{8}+\frac{\varepsilon}{2}}e_k^{\frac{1}{4}+\varepsilon}\le9LM\,e_k^{\varepsilon}\le1\qquad\left(L^{-\frac{7}{8}+\frac{\varepsilon}{2}}\le1,\ e_k\le1,\ \text{(nm6)}\right),\\
e_k\left|\mathrm{kc}_2^{\mathrm{I}}\cdot(c(\mathrm{ag})-\tilde{c})\right|\le e_k|\mathrm{kc}_2|\cdot9LM\le9LM\,L^{-1+\varepsilon}e_k^{2\varepsilon}=9M\,L^{\varepsilon}e_k^{2\varepsilon}\le1\qquad\text{(nm7)}.
\end{gather*}
So the factor of each argument is the product of a factor coming from the first slot and a factor coming from the second slot, each of modulus in $[\mathrm{e}^{-1},\mathrm{e}]$, and its modulus lies in $[\exp(-2),\exp(2)]$.

\emph{Conclusion.}
For arguments of this kind, that is, with every $\xi_i$ over $X$ and every pair of the $\boldsymbol{\zeta}_j$ admissible for $X$, (18a.16) and (18a.17) give
\begin{align*}
\left|E_{nm}(X;\mathscr{A}_L,\mathscr{C}_L;\xi;\boldsymbol{\zeta})\right|&=\left|E_{nm}(X;\mathrm{bh}+\mathrm{kc}_1,\mathscr{C}^b+\mathrm{kc}_2;\xi;\boldsymbol{\zeta})\right|\\
&=\left|E_{nm}(X;\mathrm{bh},\mathscr{C}^b;\xi;\boldsymbol{\zeta})\right|\prod_{\mathrm{ag}}\left|\exp\left\{(-1)^{\omega(\mathrm{ag})}ie_k\Lambda_{\mathrm{zs}}(c(\mathrm{ag}))\right\}\right|^{-1}\\
&\le\exp\left(2(n+m)\right)\left|E_{nm}(X;\mathrm{bh},\mathscr{C}^b;\xi;\boldsymbol{\zeta})\right|;
\end{align*}
for all other arguments the left-hand side vanishes.
For $n=m=0$ the products in (18a.16) and (18a.17) are empty, and
\[
E_{00}(X;\mathscr{A}_L,\mathscr{C}_L)=E_{00}(X;\mathrm{bh}+\mathrm{kc}_1,\mathscr{C}^b+\mathrm{kc}_2)=E_{00}(X;\mathrm{bh},\mathscr{C}^b).
\]
Integration of the inequality over all arguments with $\int\mathrm{d}\xi\,\mathrm{d}\boldsymbol{\zeta}$, that is, summation with the weights $(\eta_k^3/48)^n\left((\eta_k^3/48)^2\right)^m$ (by (I.AP) the moduli of the kernels are functions on $T^{\mathrm{f}}_k$ and on $\mathrm{Pz}_k$), yields
\[
\left\|E_{nm}(X;\mathscr{A}_L,\mathscr{C}_L)\right\|\le\exp\left(2(n+m)\right)\left\|E_{nm}(X;\mathrm{bh},\mathscr{C}^b)\right\|. \qedhere
\]
\end{proof}

\begin{named}{Remark}[Step 7 at $\mathscr{C}=0$]
Let $\mathscr{C}=0$.
By (o), $\mathscr{C}_L=0$, $\mathscr{C}^b=0$ and $\mathrm{kc}_2=0$.
Then $\mathscr{C}(z)=0$ for all $z$, $\boldsymbol{\lambda}'=0$, and $\mathscr{A}(z)$, $\Lambda_z$, $\mathrm{Kz}$ do not depend on $(z_3,z_4)$.
Thus in (7a)--(7d) the family is evaluated only at pairs whose second entry is $0$; the property (I.k4) is applied only with $\boldsymbol{\lambda}'=0$ and with both second entries equal to $0$; the continuation takes place on $\left\{(z_1,z_2)\in\mathbb{C}^2:|\operatorname{Im}z_1|<2,\ |\operatorname{Im}z_2|<2\right\}$, with the evaluation at $(1,i)$; and every argument carries a single factor, of modulus in $[\mathrm{e}^{-1},\mathrm{e}]$.
\end{named}

\paragraph{Level $0$.}
The case $k=0$ is included.
There $\eta_0=1$ and $\mathbb{L}_0=\mathbb{Z}^3$, and the charge $e_k$ of the level $k=0$ is the number $q$.
The steps of the chain in (1c) have sup length $1$, and the H\"older differences are formed at the pairs with $d(c,c')=1$.
The arguments are $\xi=(c,g,\beta,\omega)$ with $c\in\mathbb{Z}^3/\ell\mathbb{Z}^3$, with the weight $1/48$, and $\boldsymbol{\zeta}$ with the weight $(1/48)^2$; a pair inside a cluster has $d_X=1$; and for a pair of kind (II) one has $c_3=c+g\,e_1$ and $d(c,X)\le1$.

\paragraph{Uniformity.}
The numbers $1382$ and $2(n+m)$ and the exponents of $L$ are absolute.
The construction and the bounds hold for every $k\ge0$, for all $\ell$ with $\ell,\ell/L\in\mathrm{Cl}$, for every $e_k>0$ with $L^{\frac{1}{2}}e_k\le e_0$, for every small polymer $X$, for every site $\tilde{c}$ of $X$ and for every $(\mathscr{A},\mathscr{C})\in\tilde{R}_{k+1}\times\mathfrak{X}_{k+1}$.
The side $\ell$ enters only through $\mathrm{rad}'$ and through (3a), (3b), and it drops out in (18a.12).
The labels $g$, the spinor indices $\beta$ and the bar indices $\omega$ enter (iii) only through $(-1)^{\omega}$.
Since $\tilde{R}_{k+1}$ is invariant under the addition of real constant bond functions (Proposition~16.1\,(c)), fields with a twist in the first slot are among the bond functions $\mathscr{A}$ admitted in Proposition~18.2.

\subsection*{3. Proposition~18.3}

\subsubsection*{3.0. The algebra; a basis of $\mathrm{Gr}_k[Z]$; the clause (I.k2) in terms of locality}

Let $k\ge0$ be an integer, $\ell\in\mathrm{Cl}$ and $e_k>0$.
Denote by $\mathrm{Alg}$ the Grassmann algebra over $\mathbb{C}$ generated by the elements $\psi(\xi)$ and $\chi(\boldsymbol{\zeta})$ of level $k$ (Section~I.3.6) and by the spectators.
There are finitely many generators.
Hence $\mathrm{Alg}$ is a complex vector space of finite dimension, any two norms on it are equivalent, and convergence in $\mathrm{Alg}$ means convergence of every coefficient.
A subalgebra with unit is a linear subspace that contains $1$ and is closed under products.
Fix a fundamental domain of $\ell\mathbb{Z}^3$ in $\mathbb{L}_k$ and total orders of the index sets ((I.AP) in Section~I.3.6, and Section~I.5.1).
The basis monomials are the products
\[
\psi(I_1)\cdot\chi(I_2)\cdot\mathrm{sp}(\mathrm{Isp}),
\]
in which $I_1$ is a set of indices $\xi$, $I_2$ is a set of indices $\boldsymbol{\zeta}$, $\mathrm{Isp}$ is a set of indices of spectators, and $\mathrm{sp}(\mathrm{Isp})$ denotes the ordered product of the spectator generators whose indices lie in $\mathrm{Isp}$.
For $F\in\mathrm{Alg}$ let $c_F(I_1,I_2,\mathrm{Isp})$ be the coefficients of $F$ in this basis.

Let $Z$ be a set of sites of the torus $T_k$.
The properties ``$\xi=(x,\beta,\omega)$ lies over $Z$'' (that is, $\pi(x)\in Z$; see (I.POS)) and ``$\boldsymbol{\zeta}=(x,x',\beta,\omega)$ is admissible for $Z$'' (Section~I.1.7, the six paths being those of the shortest lift) are properties of the classes modulo $\ell\mathbb{Z}^3$.
By (I.AP), a generator at a translated index equals $+1$ or $-1$ times the generator at the index in the fundamental domain.
A triple $(I_1,I_2,\mathrm{Isp})$ is called \emph{allowed for} $Z$ if every $\xi\in I_1$ lies over $Z$ and every $\boldsymbol{\zeta}\in I_2$ is admissible for $Z$.

\begin{named}{Lemma (SP)}
The algebra $\mathrm{Gr}_k[Z]$ is the linear span of the basis monomials $\psi(I_1)\cdot\chi(I_2)\cdot\mathrm{sp}(\mathrm{Isp})$ for which the triple $(I_1,I_2,\mathrm{Isp})$ is allowed for $Z$.
Equivalently, an element $F$ of $\mathrm{Alg}$ lies in $\mathrm{Gr}_k[Z]$ if and only if $c_F(I_1,I_2,\mathrm{Isp})=0$ for every triple $(I_1,I_2,\mathrm{Isp})$ that is not allowed for $Z$.
\end{named}

\begin{proof}
The product of two basis monomials is either $0$ or $\pm$ the basis monomial whose three index sets are the unions of the corresponding index sets of the factors, and the triple of these unions is allowed for $Z$ whenever both factors belong to allowed triples.
Hence the span of the basis monomials of allowed triples is a subalgebra with unit, the monomial of the empty triple being $1$.
This span contains the generators named in the definition of $\mathrm{Gr}_k[Z]$, and therefore contains $\mathrm{Gr}_k[Z]$.
Conversely, the basis monomial of an allowed triple is a product of these generators, and so lies in $\mathrm{Gr}_k[Z]$.
\end{proof}

Consequently $\mathrm{Gr}_k[Z]$ is the common kernel of the linear functionals $F\mapsto c_F(I_1,I_2,\mathrm{Isp})$, where $(I_1,I_2,\mathrm{Isp})$ ranges over the triples that are not allowed for $Z$.
In particular $\mathrm{Gr}_k[Z]$ is a closed linear subspace of $\mathrm{Alg}$.

\emph{A family $\mathscr{E}$ of level $k$ has the property (I.k2) if and only if, for every $X\in D_k$, the map $\mathscr{E}(X;\cdot,\cdot)$ is local in $X$ on $\tilde{R}_k\times\mathfrak{X}_k$, without spectators.}

\begin{proof}
Let $\mathscr{E}$ be a family of level $k$ in the sense of (G2.r8), without spectators, and let $\mathrm{Dc}=\tilde{R}_k\times\mathfrak{X}_k$.
For $\#|I_1|=n$, $\#|I_2|=m$ and increasing arguments, Proposition~15\,(a) gives
\[
\mathscr{E}_{nm}(X;\mathscr{A},\mathscr{C};I_1;I_2)\cdot\nu(I_1)\cdot\nu(I_2)=c_F(I_1,I_2),
\]
the coefficient of $F=\mathscr{E}(X;\mathscr{A},\mathscr{C})$, where
\[
\nu(I_1)=\left(\frac{\eta_k^3}{48}\right)^n>0,\qquad\nu(I_2)=\left(\frac{\eta_k^3}{48}\right)^{2m}>0.
\]
An antisymmetric kernel vanishes when two of its arguments coincide, and it is determined, up to signs, by its values at increasing arguments.
Therefore the requirement of (I.k2) that the kernels vanish unless all $\xi_i$ lie over $X$ and all pairs $\boldsymbol{\zeta}_j$ are admissible for $X$ states, by (SP), that $\mathscr{E}(X;\mathscr{A},\mathscr{C})\in\mathrm{Gr}_k[X]$; this is condition~(i).
The requirement of (I.k2) that the kernels depend on $\mathscr{A}(b)$, $\mathscr{C}(b)$ only for the bonds $b$ that meet $X$ states that all kernels, or equivalently the element itself, take equal values at any two points of $\tilde{R}_k\times\mathfrak{X}_k$ that coincide on the bonds meeting $X$; this is condition~(ii).
(If the latter requirement is read for one slot at a time, one obtains the same condition, since the domain is a product: together with $(\mathscr{A},\mathscr{C})$ and $(\tilde{\mathscr{A}},\tilde{\mathscr{C}})$ it contains the point $(\tilde{\mathscr{A}},\mathscr{C})$.)
\end{proof}

\subsubsection*{3.1. (a)}

\begin{proof}[Proof of Proposition~18.3\,(a)]
Let $Z\subset Z'$.
For a pair $(x,x')\in\mathrm{Pz}_k$ write $c=\pi(x)$, $c'=\pi(x')$ and $c_3=\pi(x^{(3)})$.

\emph{Pairs.}
Kind (I) for $Z$ requires that $c$, $c'$ and the sites of the six paths $\boldsymbol{\Gamma}^{\pi}(c,c')$ lie in $Z$ (for $c=c'$: that $c\in Z$).
These sites then lie in $Z'$, so that the pair is of kind (I) for $Z'$.
Kind (II) for $Z$ requires that $x'=x^{(3)}$ and that $c$ or $c_3$ lies in $Z$.
Then $c$ or $c_3$ lies in $Z'$, and the pair is of kind (II) for $Z'$.
(A pair $(x,x^{(3)})$ with both $c$ and $c_3$ in $Z$ satisfies both conditions for $Z$, and both conditions for $Z'$.)

\emph{Bonds.}
A bond with an end in $Z$ has an end in $Z'$.

\emph{Algebras.}
A site over $Z$ lies over $Z'$.
Hence every triple allowed for $Z$ is allowed for $Z'$, and $\mathrm{Gr}_k[Z]\subset\mathrm{Gr}_k[Z']$ by (SP).

\emph{Maps.}
Let $F$ be local in $Z$ on $\mathrm{Dc}$.
As to (i): $F(\mathscr{A},\mathscr{C})$ lies in $\mathrm{Gr}_k[Z]$, hence in $\mathrm{Gr}_k[Z']$.
As to (ii): if two points of $\mathrm{Dc}$ coincide, in both components, on the bonds that meet $Z'$, then they coincide on the bonds that meet $Z$, since these are among the former; so $F$ takes the same value at the two points.
Thus $F$ is local in $Z'$.
\end{proof}

\subsubsection*{3.2. (b)}

\begin{proof}[Proof of Proposition~18.3\,(b)]
Fix $Z$ and $\mathrm{Dc}$.
Sums, products, multiples and limits of maps from $\mathrm{Dc}$ to $\mathrm{Alg}$ are taken pointwise.

(b1)~\emph{Algebra; limits.}
Let $F$ and $G$ be local in $Z$, and let $\mathrm{cf}$ be a complex number.
As to (i): $\mathrm{Gr}_k[Z]$ is a subalgebra with unit, so that $F+G$, $\mathrm{cf}\cdot F$, $F\cdot G$ and the constant map $1$ take their values in it.
As to (ii): if two points of $\mathrm{Dc}$ coincide on the bonds that meet $Z$, then $F$ and $G$ take equal values at them, and hence so do $F+G$, $\mathrm{cf}\cdot F$ and $F\cdot G$.
Next let $F_j$ be local in $Z$ for every $j$, and suppose that $F_j(\mathscr{A},\mathscr{C})\to F(\mathscr{A},\mathscr{C})$ at every point $(\mathscr{A},\mathscr{C})$ of $\mathrm{Dc}$.
As to (i): $\mathrm{Gr}_k[Z]$ is closed.
As to (ii): if $(\mathscr{A},\mathscr{C})$ and $(\tilde{\mathscr{A}},\tilde{\mathscr{C}})$ in $\mathrm{Dc}$ coincide on the bonds that meet $Z$, the two sequences $F_j(\mathscr{A},\mathscr{C})$ and $F_j(\tilde{\mathscr{A}},\tilde{\mathscr{C}})$ are equal, and so are their limits.

(b2)~\emph{Series; the exponential.}
The partial sums of a series of maps local in $Z$ are local in $Z$ by (b1).
If the series converges in norm at every point of $\mathrm{Dc}$, its sum is the pointwise limit of the partial sums, and is therefore local in $Z$ by (b1).
The multiplication of $\mathrm{Alg}$ is a bilinear map on a space of finite dimension.
So there is a number $\mathrm{Cn}$ with $\|F\cdot G\|\le\mathrm{Cn}\cdot\|F\|\cdot\|G\|$, whence $\|F^j\|\le\mathrm{Cn}^{j-1}\cdot\|F\|^j$, and the series
\[
\exp F=\sum_{j\ge0}\frac{1}{j!}\cdot F^j
\]
converges in norm for every $F\in\mathrm{Alg}$.
If $F$ is local in $Z$, then so are the terms $\frac{1}{j!}\cdot F^j$, and hence so is $\exp F$.

(b3)~\emph{Parameters.}
Let $\mathrm{Prm}$ be a set, and let $F_{\mathrm{pm}}$ be local in $Z$ on $\mathrm{Dc}$ for every $\mathrm{pm}\in\mathrm{Prm}$.
Consider a derivative of $\mathrm{pm}\mapsto F_{\mathrm{pm}}(\mathscr{A},\mathscr{C})$, where $\mathrm{Prm}$ is an open subset of a real or complex vector space of finite dimension and the derivative is a limit of difference quotients, or an iterated such limit.
Consider also an integral of $\mathrm{pm}\mapsto F_{\mathrm{pm}}(\mathscr{A},\mathscr{C})$ with respect to a complex measure on $\mathrm{Prm}$.
Both are formed coefficient by coefficient in the basis of 3.0, and each of them exists if and only if it exists for every coefficient.
Suppose that the derivative, resp.\ the integral, exists at every point of $\mathrm{Dc}$, and denote it by $G(\mathscr{A},\mathscr{C})$.
As to (i): for a triple $(I_1,I_2,\mathrm{Isp})$ that is not allowed for $Z$, the function $\mathrm{pm}\mapsto c_{F_{\mathrm{pm}}(\mathscr{A},\mathscr{C})}(I_1,I_2,\mathrm{Isp})$ vanishes identically by (SP).
Hence its derivative, resp.\ its integral, which is $c_{G(\mathscr{A},\mathscr{C})}(I_1,I_2,\mathrm{Isp})$, equals $0$, and $G(\mathscr{A},\mathscr{C})\in\mathrm{Gr}_k[Z]$ by (SP).
As to (ii): if $(\mathscr{A},\mathscr{C})$ and $(\tilde{\mathscr{A}},\tilde{\mathscr{C}})$ in $\mathrm{Dc}$ coincide on the bonds that meet $Z$, then $F_{\mathrm{pm}}(\mathscr{A},\mathscr{C})=F_{\mathrm{pm}}(\tilde{\mathscr{A}},\tilde{\mathscr{C}})$ for every $\mathrm{pm}$.
The two functions of $\mathrm{pm}$ are thus one and the same function, and $G(\mathscr{A},\mathscr{C})=G(\tilde{\mathscr{A}},\tilde{\mathscr{C}})$.
The four operations named in (b) are of this kind.
\begin{itemize}
\item[$\bullet$] $\frac{\partial}{\partial s_P}$ is an iterated partial derivative in the parameters $s_{\mathrm{pc}}$, $\mathrm{pc}\in P$.
\item[$\bullet$] $\int_{[0,1]^P}\mathrm{d}s_P$ is the integral with respect to the Lebesgue measure on $[0,1]^P$.
\item[$\bullet$] Let $\mathrm{rad}\ne1$.
With $t=\mathrm{rad}\cdot\mathrm{e}^{i\,\mathrm{ang}}$ and $\mathrm{d}t=i\cdot t\,\mathrm{d}\,\mathrm{ang}$ one has (Section~I.5.2)
\begin{align*}
\oint_{\mathrm{rad}}f&=(2\pi i)^{-1}\oint_{|t|=\mathrm{rad}}\mathrm{d}t\,\frac{f(t)}{t\cdot(t-1)}\\
&=(2\pi)^{-1}\int_0^{2\pi}\mathrm{d}\,\mathrm{ang}\,f\left(\mathrm{rad}\cdot\mathrm{e}^{i\,\mathrm{ang}}\right)\cdot\left(\mathrm{rad}\cdot\mathrm{e}^{i\,\mathrm{ang}}-1\right)^{-1},
\end{align*}
an integral over $[0,2\pi]$ with a bounded density.
\item[$\bullet$] The integral $\int\mathrm{d}\hat{\mu}(\mathsf{Z})$ exists only if the measure $\hat{\mu}$ is defined, that is, only if $p_{0,k}>0$; in that case
\[
\mathrm{d}\hat{\mu}=\frac{\mathbf{1}_{\mathrm{Box}_k}}{\mu_I(\mathrm{Box}_k)}\,\mathrm{d}\mu_I
\]
is a probability measure on $\mathrm{Box}_k$ \cite[Section~5.2.5]{QEDI}.
Indeed, $p_{0,k}=v_k^{p_0}$ with $v_k=-\log e_k$ and the integer $p_0$ (Section~I.6.1), so that $p_{0,k}>0$ if $e_k<1$.
Further,
\[
\mathrm{Box}_k=\left\{\mathsf{Z}\in V_k:\ |\mathsf{Z}(\Upsilon)|\le p_{0,k}\ \text{for all}\ \Upsilon\right\},
\]
where $\{\Upsilon\}$ is a basis of $V_k$ with $d_V\cdot n_k>0$ elements and $\mathsf{Z}(\Upsilon)=\langle\mathsf{Z},\Upsilon\rangle$, and $\mu_I$ is the standard Gaussian measure on $V_k$, which has a positive density.
If $p_{0,k}>0$, then $\mathrm{Box}_k$ contains the open set $\left\{\mathsf{Z}:\ |\mathsf{Z}(\Upsilon)|<p_{0,k}\ \text{for all}\ \Upsilon\right\}$, which contains $0$; hence $\mu_I(\mathrm{Box}_k)>0$.
If $p_{0,k}\le0$, then $\mathrm{Box}_k$ is empty or equal to $\{0\}$, and $\mu_I(\mathrm{Box}_k)=0$.
\end{itemize}

(b4)~\emph{The integral $\int\mathrm{d}\mu_I(W)$.}
Let $W$ be one of the spectators, on an index set $\mathrm{Mo}\times\{0,1\}$ with weights $\nu$ (Proposition~15\,(D2)).
In the basis monomials the family $W$ is placed last, as in Proposition~15\,(D3): $\mathrm{sp}(\mathrm{Isp})=\mathrm{sp}'(\mathrm{Isp}')\cdot W(J'')$, where $\mathrm{sp}'$ refers to the other spectators.
By Proposition~15\,(D3), $\int\mathrm{D}W$ is linear and maps $\phi(I)\cdot G$ to $\phi(I)\cdot\int\mathrm{D}W\,G$ for $G\in\mathrm{Gr}[W]$ and a monomial $\phi(I)$ in the other generators.
Moreover $\exp\left(-\langle\bar{W},W\rangle_{\nu}\right)$ is even and commutes with $\phi(I)$.
With the number
\[
m(J''):=\int\mathrm{d}\mu_I(W)\,W(J'')
\]
of Proposition~15\,(e) it follows that
\[
\int\mathrm{d}\mu_I(W)\,\psi(I_1)\cdot\chi(I_2)\cdot\mathrm{sp}'(\mathrm{Isp}')\cdot W(J'')=m(J'')\cdot\psi(I_1)\cdot\chi(I_2)\cdot\mathrm{sp}'(\mathrm{Isp}').
\]
The triple $(I_1,I_2,\mathrm{Isp}')$ is allowed for $Z$ if $(I_1,I_2,\mathrm{Isp}'\cup J'')$ is.
Therefore, by (SP), the linear map $F\mapsto\int\mathrm{d}\mu_I(W)\,F$ maps $\mathrm{Gr}_k[Z]$ into $\mathrm{Gr}_k[Z]$; this gives (i).
A fixed map takes equal values at equal arguments; this gives (ii).
(For $k\ge1$ one has $\nu=1$ and $m(J'')\in\{0,1,-1\}$.
At level~$0$ one has $\mathrm{Mo}=T^{\mathrm{f}}_0\times\{1,\dots,4\}$ and $\nu=\frac{1}{48}$, and $m(J'')=\mathrm{sn}(J'')\cdot48^{\frac{\#|J''|}{2}}$ for paired $J''$, $m(J'')=0$ otherwise.
The values of $m(J'')$ play no role here.)

(b5)~\emph{Products.}
Let $F_i$ be local in $Z_i$ for $i=1,\dots,n$.
The set $Z_1\cup\dots\cup Z_n$ is a union of $M$-cubes that contains every $Z_i$.
By (a), every $F_i$ is local in this union, and by (b1) so is the product $F_1\cdots F_n$.
\end{proof}

\subsubsection*{3.3. (c), (d)}

\begin{proof}[Proof of Proposition~18.3\,(c), (d)]
Assume (G2.S-1) and the data of Propositions~16.3 and~17.2.
Clause (c) is the last sentence of Proposition~16.3\,(d), and clause (d) is the last sentence of Proposition~17.2\,(c).
\end{proof}

\subsubsection*{3.4. (e)}

\begin{proof}[Proof of Proposition~18.3\,(e)]
Assume (G2.S-2).
The words ``level $k+1$'' refer to the data $(k+1,\ell/L,e_{k+1})$.
Here $\ell/L\in\mathrm{Cl}$ and $e_{k+1}>0$, so that the definition in Proposition~18.3, the facts of 3.0 and the clauses (a), (b) hold at level $k+1$.
The algebra $\mathrm{Gr}_{k+1}[Y]$ is formed in $\mathrm{Gr}[\psi,\chi]$ of level $k+1$, without spectators.
By (sc1), the inclusion $X\subset LY$ means that $c/L\in Y$ for every site $c$ of $X$.
The homomorphism $S_L$ is that of (G2.r5); the compatibility of the rule (G2.r5) with (I.AP) follows from (sf1) and (sf4), as shown in the paragraph on the objects at the beginning of the proof in Section~4 below.

Let $\mathrm{Dc}$ and $\mathrm{Dc}'$ be as in the first sentence of (e).
Denote by (1) the assertion of the third and fourth sentences of (e), which concerns a single map $F$, and by (2) the assertion of (e) on the sums over the $X\in D_k$ with $\bar{X}^{L}=LY$:
\begin{itemize}
\item[(1)] if $X\in D_k$ and $Y\in D_{k+1}$ satisfy $X\subset LY$, and $F$ is local in $X$ on $\mathrm{Dc}$, without spectators, then the map $(\mathscr{A},\mathscr{C})\mapsto S_L\left[F(\mathscr{A}_L,\mathscr{C}_L)\right]$ is local in $Y$ on $\mathrm{Dc}'$ at level $k+1$;
\item[(2)] if for every $X\in D_k$ the map $E'(X;\cdot,\cdot)$ is local in $X$ on $\mathrm{Dc}$, without spectators, then for every $Y\in D_{k+1}$ the map
\[
(\mathscr{A},\mathscr{C})\mapsto\sum_{X:\ \bar{X}^{L}=LY}S_L\left[E'(X;\mathscr{A}_L,\mathscr{C}_L)\right]
\]
is local in $Y$ on $\mathrm{Dc}'$.
\end{itemize}

\emph{The second sentence of (e): the three pairs of domains.}
In each of the three cases let $(\mathscr{A},\mathscr{C})\in\mathrm{Dc}'$.
\begin{itemize}
\item[(dm1)] $\mathrm{Dc}=\tilde{R}_k\times\mathfrak{X}_k$ and $\mathrm{Dc}'=\tilde{R}_{k+1}\times\mathfrak{X}_{k+1}$.
By Proposition~18.1\,(c), $\mathscr{A}_L\in\tilde{R}_k$.
By Proposition~18.1\,(b), $\mathscr{C}_L$ lies in $L^{-\frac{3}{4}}\mathfrak{X}_k$, which is contained in $\mathfrak{X}_k$ because $L^{-\frac{3}{4}}\le1$.
\item[(dm2)] $\mathrm{Dc}=\tilde{R}_k^{-}\times\mathfrak{X}_k$ and $\mathrm{Dc}'=\tilde{R}_{k+1}\times\mathfrak{X}_{k+1}$.
By Proposition~18.1\,(c), $\mathscr{A}_L\in\tilde{R}_k^{-}$; for $\mathscr{C}_L$ see (dm1).
\item[(dm3)] $\mathrm{Dc}=\tilde{R}_k\times\{0\}$ and $\mathrm{Dc}'=\tilde{R}_{k+1}\times\{0\}$.
By Proposition~18.1\,(c), $\mathscr{A}_L\in\tilde{R}_k$; and $0_L=0$.
\end{itemize}

\emph{Proof of (1), of (2), and of the assertions of (e) on the legs at pairs of kind (I) and of kind (II).}

\emph{The map is defined.}
For $(\mathscr{A},\mathscr{C})\in\mathrm{Dc}'$ the pair $(\mathscr{A}_L,\mathscr{C}_L)$ lies in $\mathrm{Dc}$, which is the domain of $F$.

\emph{Legs $\psi$.}
One has $S_L\psi(x,\beta,\omega)=L^{-1}\cdot\psi(x/L,\beta,\omega)$.
If $x=(c,g)$ lies over $X$, then $\pi(x/L)=c/L$ lies in $Y$; that is, $x/L=(c/L,g)$ lies over $Y$.

\emph{Legs $\chi$ of kind (I).}
One has $S_L\chi(x,x',\beta,\omega)=L^{-1-\alpha}\cdot\chi(x/L,x'/L,\beta,\omega)$.
Let $(x,x')\in\mathrm{Pz}_k$, with $c=\pi(x)$ and $c'=\pi(x')$, be of kind (I) for $X$, and let $d_X(x,x')=d$.
By (sf4) and (sf3), $(x/L,x'/L)$ is a pair of $\mathrm{Pz}_{k+1}$ with $d_X=d/L$.
The sites $c/L$ and $c'/L$ lie in $Y$.
By (sc5), the sites of the six paths $\boldsymbol{\Gamma}^{\pi}(c/L,c'/L)$ are the sites of $\boldsymbol{\Gamma}^{\pi}(c,c')$ divided by $L$; the latter lie in $X$, and so the paths $\boldsymbol{\Gamma}^{\pi}(c/L,c'/L)$ lie in $Y$.
(The paths of a pair of the torus are those of its lift in $\mathrm{Pzt}_k$, and by (sf4) the lift of the image pair is the image of the lift.)
If $c=c'$, that is, for a pair inside a cluster ($g\ne g'$, $d=\eta_k$), then $c/L=c'/L$ lies in $Y$, the labels $g\ne g'$ are not changed, and $d/L=\eta_{k+1}$; so the image is again a pair inside a cluster.
Hence $(x/L,x'/L)$ is of kind (I) for $Y$.

\emph{Legs $\chi$ of kind (II).}
Let $x'=x^{(3)}$, and let $c$ or $c_3=\pi(x^{(3)})$ lie in $X$.
By (sf2), $x^{(3)}/L=(x/L)^{(3)}$.
So the image pair is $(x/L,(x/L)^{(3)})$, and $\pi(x/L)=c/L$ or $\pi\left((x/L)^{(3)}\right)=c_3/L$ lies in $Y$: the image pair is of kind (II) for $Y$.
The two preceding paragraphs prove what (e) asserts on the legs: on a leg at a pair of kind (I) for $X$, including a pair inside a cluster, and on a leg at a pair of kind (II) for $X$.

\emph{Condition (i).}
The homomorphism $S_L$ maps $1$ to $1$, and it maps every generator of $\mathrm{Gr}_k[X]$ (there are no spectators) to a multiple of a generator of $\mathrm{Gr}_{k+1}[Y]$.
Therefore $S_L$ maps $\mathrm{Gr}_k[X]$ into $\mathrm{Gr}_{k+1}[Y]$, and $S_L\left[F(\mathscr{A}_L,\mathscr{C}_L)\right]$ lies in $\mathrm{Gr}_{k+1}[Y]$.

\emph{Condition (ii).}
Let $(\mathscr{A},\mathscr{C})$ and $(\tilde{\mathscr{A}},\tilde{\mathscr{C}})$ in $\mathrm{Dc}'$ coincide, in both components, on the bonds of $T_{k+1}$ that meet $Y$.
Let $b$ be a bond of $T_k$ that meets $X$.
It has an end in $X$, and this end lies in $LY$; by (sc2), $b/L$ has an end in $Y$.
Consequently
\begin{align*}
\mathscr{A}_L(b)&=L^{-\frac{1}{2}}\cdot\mathscr{A}(b/L)=L^{-\frac{1}{2}}\cdot\tilde{\mathscr{A}}(b/L)=\tilde{\mathscr{A}}_L(b),\\
\mathscr{C}_L(b)&=L^{-\frac{1}{2}}\cdot\mathscr{C}(b/L)=L^{-\frac{1}{2}}\cdot\tilde{\mathscr{C}}(b/L)=\tilde{\mathscr{C}}_L(b).
\end{align*}
So $(\mathscr{A}_L,\mathscr{C}_L)$ and $(\tilde{\mathscr{A}}_L,\tilde{\mathscr{C}}_L)$ are two points of $\mathrm{Dc}$ that coincide, in both components, on the bonds that meet $X$.
Since $F$ is local in $X$ on $\mathrm{Dc}$, one has $F(\mathscr{A}_L,\mathscr{C}_L)=F(\tilde{\mathscr{A}}_L,\tilde{\mathscr{C}}_L)$, and the images under $S_L$ are equal.
This proves (1).

\emph{Sums.}
Let $Y\in D_{k+1}$.
If $\bar{X}^{L}=LY$, then $X\subset LY$, because $\bar{X}^{L}$ is the union of the $LM$-cubes that contain an $M$-cube of $X$ \cite[Section~13.1.7]{QEDI}, and every $M$-cube of $X$ lies in an $LM$-cube (sc4).
The set $D_k$ is finite.
By (1), every term of the finite sum in (2) is local in $Y$ on $\mathrm{Dc}'$, and by (b1) at level $k+1$ so is the sum; the empty sum is the constant map $0$.
This proves (2).

\emph{The final assertion of (e): if $E'$ has (I.k2) at level $k$, then $L E'$ has (I.k2) at level $k+1$.}
Let the family $E'$ of level $k$ have the property (I.k2).
By 3.0, $E'(X;\cdot,\cdot)$ is local in $X$ on $\tilde{R}_k\times\mathfrak{X}_k$, without spectators, for every $X\in D_k$.
By (2) together with (dm1), $(L E')(Y;\cdot,\cdot)$ is local in $Y$ on $\tilde{R}_{k+1}\times\mathfrak{X}_{k+1}$ for every $Y\in D_{k+1}$.
By 3.0 at level $k+1$, this is the property (I.k2) for $L E'$.
(With (dm2) in the place of (dm1) one obtains: if the maps $E'(X;\cdot,\cdot)$, $X\in D_k$, are given, and local in $X$, on $\tilde{R}_k^{-}\times\mathfrak{X}_k$ only, then the family of level $k+1$ with the elements $\sum_{X:\ \bar{X}^{L}=LY}S_L\left[E'(X;\mathscr{A}_L,\mathscr{C}_L)\right]$ has the property (I.k2).)
\end{proof}

\paragraph{Level~0.}
The case $k=0$ is included in all clauses.
In (e) at $k=0$ one has $\eta_0=1$, and every pair of $\mathrm{Pz}_0$ has $d_X\ge1$.
A pair inside a cluster has $d_X=1$, and it is mapped to a pair inside a cluster of level~$1$, with $d_X=\eta_1=1/L$.
A pair $(x,x^{(3)})$ has $d_X=1$ and $c_3=c+g e_1$, and it is mapped to $(x/L,(x/L)^{(3)})$, where $\pi\left((x/L)^{(3)}\right)=c/L+\eta_1\cdot g e_1$.
In (b4) at $k=0$ the index set of $W$ is $T^{\mathrm{f}}_0\times\{1,\dots,4\}\times\{0,1\}$, with the weight $\frac{1}{48}$.

\paragraph{Uniformity.}
Proposition~18.3 involves no constant, and its clauses (a) and (b) hold, with the proofs of 3.1 and 3.2, for all data $(k,\ell,e_k)$ for which $T_k$, $T^{\mathrm{f}}_k$ and $\mathrm{Pz}_k$ are defined, that is, for $2M<\ell$.
These two clauses depend on $e_k$ only through the integral $\int\mathrm{d}\hat{\mu}(\mathsf{Z})$ in (b3), which exists only if $p_{0,k}>0$.
The assertions (1) and (2) hold, with the proofs of 3.4, for any two sets $\mathrm{Dc}$, $\mathrm{Dc}'$ as in the first sentence of (e), that is, with $(\mathscr{A}_L,\mathscr{C}_L)\in\mathrm{Dc}$ for every $(\mathscr{A},\mathscr{C})\in\mathrm{Dc}'$; the facts (dm1)--(dm3) follow from Proposition~18.1\,(b), (c).

\subsection*{4. Proposition~18.4}

\begin{proof}[Proof of Proposition~18.4]
\textbf{The objects.}
By Proposition~16.1\,(d), under the sole hypothesis $\ell\in\mathrm{Cl}$, the number $\ell/L$ is an odd multiple of $M$ with $\ell/L\ge13M>2M$.
Moreover $T_{k+1}=\mathbb{L}_{k+1}/(\ell/L)\mathbb{Z}^3$, $T^{\mathrm{f}}_{k+1}$, $\mathrm{Pz}_{k+1}$ with its shortest lifts, the generators $\psi$, $\chi$ of level $k+1$ and $S_L$ are given by the formulas of Sections~I.1.1, I.1.3, I.1.7, I.3.6 and of (G2.r5).
The homomorphism $S_L$ is well defined.
Indeed, the rule (G2.r5), stated on the generators indexed by $\mathbb{L}^{\mathrm{f}}_k$ and $\mathrm{Pzt}_k$, is compatible with (I.AP), because by (sf1)
\begin{align*}
S_L\psi(x+\ell\cdot e_{\mu},\beta,\omega)&=L^{-1}\cdot\psi\left(x/L+(\ell/L)\cdot e_{\mu},\beta,\omega\right)=-L^{-1}\cdot\psi(x/L,\beta,\omega)=-S_L\psi(x,\beta,\omega),\\
S_L\chi(x+\ell\cdot e_{\mu},x'+\ell\cdot e_{\mu},\beta,\omega)&=L^{-1-\alpha}\cdot\chi\left(x/L+(\ell/L)\cdot e_{\mu},x'/L+(\ell/L)\cdot e_{\mu},\beta,\omega\right)\\
&=-L^{-1-\alpha}\cdot\chi(x/L,x'/L,\beta,\omega)=-S_L\chi(x,x',\beta,\omega);
\end{align*}
by (sf4) the index on the right-hand side of (G2.r5) is a pair of $\mathrm{Pzt}_{k+1}$.
The images of the generators are odd elements, which anticommute; so, by the universal property of a Grassmann algebra, there is exactly one homomorphism with (G2.r5).

\textbf{The weights.}
One has $h_{k,1}=e_k^{-\frac{1}{4}}$ and $h_{k,2}=e_k^{-\frac{1}{4}+\varepsilon}$ (Section~I.5.1), and $e_{k+1}=L^{\frac{1}{2}}\cdot e_k$.
Hence
\[
h_{k+1,1}=L^{-\frac{1}{8}}\cdot e_k^{-\frac{1}{4}}=L^{-\frac{1}{8}}\cdot h_{k,1},\qquad h_{k+1,2}=L^{\frac{1}{2}\left(-\frac{1}{4}+\varepsilon\right)}\cdot e_k^{-\frac{1}{4}+\varepsilon}=L^{-\frac{1}{8}+\frac{\varepsilon}{2}}\cdot h_{k,2}.
\]

\textbf{(a)}
The assertion of (a) that $x\mapsto x/L$ is a bijection of $T^{\mathrm{f}}_k$ onto $T^{\mathrm{f}}_{k+1}$ is (sf1).
The assertion of (a) that $(x,x')\mapsto(x/L,x'/L)$ is a bijection of $\mathrm{Pz}_k$ onto the set of the pairs of $\mathrm{Pz}_{k+1}$ with $d_X\le M/L$, every pair inside a cluster belonging to this set, is (sf4).

\emph{Indices.}
For indices $\xi=(x,\beta,\omega)$ and $\boldsymbol{\zeta}=(x,x',\beta,\omega)$ of level $k$ put
\[
\xi/L:=(x/L,\beta,\omega),\qquad\boldsymbol{\zeta}/L:=(x/L,x'/L,\beta,\omega);
\]
then $L(\xi/L)=\xi$ and $L(\boldsymbol{\zeta}/L)=\boldsymbol{\zeta}$ in the notation of (G2.r5).
By (sf1) and (sf4), $\xi\mapsto\xi/L$ is a bijection of the index set of $\psi$ of level $k$ onto that of level $k+1$, and $\boldsymbol{\zeta}\mapsto\boldsymbol{\zeta}/L$ is a bijection of the index set of $\chi$ of level $k$ onto the set of the indices $\boldsymbol{\zeta}'$ of level $k+1$ whose pair has $d_X\le M/L$.
The label $g$, the spinor index $\beta$ and the bar index $\omega$ are not changed.
For a set $I$ of indices of one of the two families, $I/L$ denotes the set of the $\xi/L$, resp.\ of the $\boldsymbol{\zeta}/L$, with $\xi$, resp.\ $\boldsymbol{\zeta}$, in $I$.

\emph{Monomials.}
A fundamental domain and total orders being fixed at each of the two levels (Section~I.5.1), let $\psi(I_1)\cdot\chi(I_2)$ be a basis monomial of level $k$ with $\#|I_1|=n$ and $\#|I_2|=m$.
Since $S_L$ is a homomorphism, (G2.r5) yields the formula (18a.18) for the image of a basis monomial under $S_L$:
\begin{equation}
\begin{split}
S_L\left[\psi(I_1)\cdot\chi(I_2)\right]&=L^{-n}\cdot L^{-(1+\alpha)\cdot m}\cdot\prod_{\xi\in I_1}\psi(\xi/L)\cdot\prod_{\boldsymbol{\zeta}\in I_2}\chi(\boldsymbol{\zeta}/L)\\
&=\mathrm{sn}(I_1,I_2)\cdot L^{-n-(1+\alpha)\cdot m}\cdot\psi(I_1/L)\cdot\chi(I_2/L),
\end{split}
\tag{18a.18}
\end{equation}
where in the middle expression the factors stand in the order of $I_1$ and of $I_2$, and where $\mathrm{sn}(I_1,I_2)\in\{1,-1\}$.
Indeed, the $n+m$ generators are distinct; putting them into the increasing order of level $k+1$ produces the sign of a permutation, and replacing an index by its translate in the fundamental domain of level $k+1$ produces a sign by (I.AP).
If the fundamental domain of level $k+1$ is the image of that of level $k$ under $c\mapsto c/L$, and the orders of level $k+1$ are chosen so that $\xi\mapsto\xi/L$ and $\boldsymbol{\zeta}\mapsto\boldsymbol{\zeta}/L$ are increasing, then $\mathrm{sn}=1$ for all monomials.
The map $(I_1,I_2)\mapsto(I_1/L,I_2/L)$ is injective.
Consequently $S_L$ maps the basis monomials of level $k$ injectively to multiples of basis monomials of level $k+1$, with the factor $\mathrm{sn}(I_1,I_2)\cdot L^{-n-(1+\alpha)\cdot m}$ on bidegree $(n,m)$.
The assertions (b) and (c) do not depend on $\mathrm{sn}$: in the proof of (b) only $|\mathrm{sn}(I_1,I_2)|=1$ enters, and the proof of (c) does not use (18a.18).

\textbf{(b)}
Let $F=\sum c_F(I_1,I_2)\cdot\psi(I_1)\cdot\chi(I_2)$.
By (18a.18),
\[
S_LF=\sum c_F(I_1,I_2)\cdot\mathrm{sn}(I_1,I_2)\cdot L^{-\#|I_1|-(1+\alpha)\cdot\#|I_2|}\cdot\psi(I_1/L)\cdot\chi(I_2/L),
\]
and by the injectivity this is the expansion of $S_LF$ in the basis of level $k+1$, all other coefficients being $0$.
By (I.5.1),
\begin{align*}
\left\|S_LF\right\|_{(\mathsf{w}_1,\mathsf{w}_2)}&=\sum\left|c_F(I_1,I_2)\right|\cdot L^{-\#|I_1|-(1+\alpha)\cdot\#|I_2|}\cdot\mathsf{w}_1^{\#|I_1|}\cdot\mathsf{w}_2^{\#|I_2|}\\
&=\sum\left|c_F(I_1,I_2)\right|\cdot\left(L^{-1}\cdot\mathsf{w}_1\right)^{\#|I_1|}\cdot\left(L^{-1-\alpha}\cdot\mathsf{w}_2\right)^{\#|I_2|}\\
&=\|F\|_{\left(L^{-1}\cdot\mathsf{w}_1,\,L^{-1-\alpha}\cdot\mathsf{w}_2\right)}.
\end{align*}
For $F^{[n,m]}$ only the sets with $\#|I_1|=n$ and $\#|I_2|=m$ occur.
Taking $\mathsf{w}=h_{k+1}$ and using the relations between the weights, one obtains, with the sums over these sets,
\begin{align*}
\left\|S_LF^{[n,m]}\right\|_{h_{k+1}}&=L^{-n-(1+\alpha)\cdot m}\cdot h_{k+1,1}^n\cdot h_{k+1,2}^m\cdot\sum\left|c_F(I_1,I_2)\right|\\
&=L^{-n-(1+\alpha)\cdot m}\cdot L^{-\frac{n}{8}}\cdot L^{\left(-\frac{1}{8}+\frac{\varepsilon}{2}\right)\cdot m}\cdot h_{k,1}^n\cdot h_{k,2}^m\cdot\sum\left|c_F(I_1,I_2)\right|\\
&=L^{-n-(1+\alpha)\cdot m}\cdot L^{-\frac{n+m}{8}+\frac{m\cdot\varepsilon}{2}}\cdot\left\|F^{[n,m]}\right\|_{h_k}.
\end{align*}

\textbf{(c)}
Let $F_{nm}$ be the antisymmetric kernels of $F$:
\begin{multline*}
F=\sum_{n,m}\frac{1}{n!\,m!}\sum_{\xi_1,\dots,\xi_n}\ \sum_{\boldsymbol{\zeta}_1,\dots,\boldsymbol{\zeta}_m}\left(\frac{\eta_k^3}{48}\right)^n\cdot\left(\left(\frac{\eta_k^3}{48}\right)^2\right)^m\\
\cdot F_{nm}(\xi;\boldsymbol{\zeta})\cdot\psi(\xi_1)\cdots\psi(\xi_n)\cdot\chi(\boldsymbol{\zeta}_1)\cdots\chi(\boldsymbol{\zeta}_m).
\end{multline*}
Here, by Section~I.3.6, the weight of an index $\xi$ is $\frac{\eta_k^3}{48}$ and that of an index $\boldsymbol{\zeta}$ is $\left(\frac{\eta_k^3}{48}\right)^2$.
Each sum over an index $\xi$ contains the sum over the $48$ labels, and each sum over an index $\boldsymbol{\zeta}$ contains the sums over the labels of both points.
The products of the kernel with the generators are $\ell\mathbb{Z}^3$-periodic in each argument, so that by (I.AP) the sums do not depend on the fundamental domain.
Apply $S_L$ and (G2.r5), and substitute $\xi_i=L\xi'_i$ and $\boldsymbol{\zeta}_j=L\boldsymbol{\zeta}'_j$.
By (a), the $\xi'_i$ then range over all indices of $\psi$ of level $k+1$, and the $\boldsymbol{\zeta}'_j$ over the indices of $\chi$ of level $k+1$ whose pair has $d_X\le M/L$.
With $\frac{\eta_k^3}{48}=L^3\cdot\frac{\eta_{k+1}^3}{48}$ one obtains
\begin{multline*}
S_LF=\sum_{n,m}\frac{1}{n!\,m!}\sum_{\xi'}\ \sum_{\boldsymbol{\zeta}':\ d_X\le M/L}\left(\frac{\eta_{k+1}^3}{48}\right)^n\cdot\left(\left(\frac{\eta_{k+1}^3}{48}\right)^2\right)^m\cdot L^{3n+6m}\cdot L^{-n-(1+\alpha)\cdot m}\\
\cdot F_{nm}(L\xi';L\boldsymbol{\zeta}')\cdot\psi(\xi'_1)\cdots\psi(\xi'_n)\cdot\chi(\boldsymbol{\zeta}'_1)\cdots\chi(\boldsymbol{\zeta}'_m),
\end{multline*}
where the second sum extends over the tuples $\boldsymbol{\zeta}'$ all of whose pairs have $d_X\le M/L$.
This is a representation of $S_LF$ by kernels with respect to the weights of level $k+1$, with the kernel
\[
K'_{nm}(\xi';\boldsymbol{\zeta}'):=\begin{cases}L^{3n+6m}\cdot L^{-n-(1+\alpha)\cdot m}\cdot F_{nm}(L\xi';L\boldsymbol{\zeta}')&\text{if all pairs of }\boldsymbol{\zeta}'\text{ have }d_X\le M/L,\\0&\text{otherwise.}\end{cases}
\]
The kernel $K'_{nm}$ is antisymmetric in the $\xi'_i$ and in the $\boldsymbol{\zeta}'_j$, because $F_{nm}$ is and the condition on the pairs is symmetric.
It is antiperiodic in each argument under the translations by $(\ell/L)\cdot e_{\mu}$, because $L\left(\xi'+(\ell/L)\cdot e_{\mu}\right)=L\xi'+\ell\cdot e_{\mu}$.
By Proposition~15\,(a), the antisymmetric kernels are determined by the element; hence $(S_LF)_{nm}=K'_{nm}$.
Finally,
\begin{align*}
\left\|(S_LF)_{nm}\right\|&=\sum_{\xi'}\sum_{\boldsymbol{\zeta}'}\left(\frac{\eta_{k+1}^3}{48}\right)^{n+2m}\cdot\left|K'_{nm}(\xi';\boldsymbol{\zeta}')\right|\\
&=L^{-n-(1+\alpha)\cdot m}\cdot\sum_{\xi}\sum_{\boldsymbol{\zeta}}\left(\frac{L^3\cdot\eta_{k+1}^3}{48}\right)^{n+2m}\cdot\left|F_{nm}(\xi;\boldsymbol{\zeta})\right|\\
&=L^{-n-(1+\alpha)\cdot m}\cdot\sum_{\xi}\sum_{\boldsymbol{\zeta}}\left(\frac{\eta_k^3}{48}\right)^{n+2m}\cdot\left|F_{nm}(\xi;\boldsymbol{\zeta})\right|\\
&=L^{-n-(1+\alpha)\cdot m}\cdot\left\|F_{nm}\right\|.
\end{align*}
\end{proof}

\paragraph{Level~0.}
The case $k=0$ is included.
Then $\eta_0=1$; the weights are $\frac{1}{48}$ and $\left(\frac{1}{48}\right)^2$ at level~$0$, and $\frac{\eta_1^3}{48}=\frac{L^{-3}}{48}$ and $\left(\frac{L^{-3}}{48}\right)^2$ at level~$1$; moreover $\frac{1}{48}=L^3\cdot\frac{\eta_1^3}{48}$.
The pairs of $\mathrm{Pz}_0$, all of which satisfy $1\le d_X\le M$, are mapped onto the pairs of $\mathrm{Pz}_1$ with $\eta_1=1/L\le d_X\le M/L$.
At $k=0$ the identity (18a.18) and the identities proved in (b) and (c) hold with these values: $\eta_0=1$, and in (c) the weights $\frac{1}{48}$ and $\left(\frac{1}{48}\right)^2$ of level~$0$.

\paragraph{Uniformity.}
All relations proved above are identities, and the only numbers in them are the displayed powers of $L$.
They hold for every integer $k\ge0$, every $\ell\in\mathrm{Cl}$ and every $e_k>0$.

\subsection*{5. Conditions on the constants}

Besides the conditions on the constants under which Propositions~15\,(a), (e), 16.1\,(c), (d), 16.3\,(d), 17.2\,(c), \cite[Theorem~13.2\,(t7)]{QEDI} and \cite[Lemmas~13.9, 13.10 and~17.3]{QEDI} hold, the proofs of Sections~0--4 require the following conditions, each stated in the form in which it enters.

\begin{description}
\item[$L$ is an odd integer and $L\ge 5$]
It belongs to the line (C3) of Section~I.8.2 and coincides with the row (L5) of Proposition~14.2\,(b).
The inequality $L\ge 1$ enters all comparisons of powers of $L$;
the oddness of $L$ enters (sc4) and the paragraph on faces of Section~I.1.1;
the inequality $L\ge 5$ enters (ge2)--(ge4), (3b), (18a.8) and, through $\mathrm{rad}\ge 50$, (1c), (1d) and Steps~5 and~6 of the proof of Proposition~18.2.
The condition is applied in Sections~0.1, 0.2 and~0.4, in the proof of Proposition~18.1\,(b), (c), and in Steps~0, 1, 3, 4, 5, 6 and~7 of the proof of Proposition~18.2.
The proof of Proposition~18.3\,(e) applies it in two ways:
through (dm1) and (sf4), which rest on $L\ge 1$,
and through (sc4), which rests on the oddness of $L$, in the paragraph on the sums in Section~3.4, where every $M$-cube of $X$ lies in an $LM$-cube.
The proof of Proposition~18.4\,(a) applies it through (sf4).

\item[$L^{\frac{1}{8}-\frac{3\bar{\varepsilon}}{2}}\ge 2400$, through (nm1)--(nm4)]
This is the first condition of line (C3) and the row (L1) of Proposition~14.2\,(b).
It is applied in the last sentence of the proof of Proposition~18.1\,(c), through (nm1);
in Step~6 of the proof of Proposition~18.2, through (nm3), in the bound on the difference quotients of $\mathscr{C}^b$ in (18a.14);
in the proof of Proposition~18.2\,(ii), through (nm2), the bounds on $\mathscr{C}^b$ being those of (18a.14);
and in Step~7 of the proof of Proposition~18.2, through (nm2) and (nm4).

\item[$M$ is a power of $L$ (hence $M$ is odd) and $M\ge L^2$]
It belongs to the line (C5) of Section~I.8.2 and coincides with the row (M1) of Proposition~14.2\,(b).
The oddness of $M$ enters the paragraph on faces of Section~I.1.1;
the inequality $M\ge 5$ enters (ge2)--(ge4), (3b) and the inequality $\mathrm{rad}\ge 50$, hence (1c), (1d), (18a.8) and Steps~5 and~6;
the inequality $L\le M$ enters (sf4).
In Step~7 of the proof of Proposition~18.2 the condition enters through (ge4), which is applied in (7c),
and in (7d) through the inequalities $M(L+2)+1\le\mathrm{rad}$, a consequence of (ge4), and $2M+1\le 2LM$;
of $M$ these require only $M\ge 1$, which follows from $M\ge 5$.
The condition is applied in Section~0.2, in the proof of Proposition~18.1\,(c), in Steps~0, 1, 3, 4, 5, 6 and~7 of the proof of Proposition~18.2, and, through (sf4), in the proofs of Propositions~18.3\,(e) and~18.4\,(a).

\item[$0<\alpha$ and $\alpha\le 1$]
The condition follows from the line (C1) of Section~I.8.2, $\alpha\in(1-\bar{\varepsilon},1)$, because $1-\bar{\varepsilon}=\frac{199}{200}>0$ by the line (C0).
The inequality $0<\alpha$ enters in the form $L^{-\alpha}\le 1$, and in the form $d(c,c')^{\alpha}\le 1$ for $0<d(c,c')\le 1$ in (1c), (3a), (3b);
the inequality $\alpha\le 1$ enters the proof of (qu2), that is, of \cite[Lemma~13.9]{QEDI}.
In Section~0.4 the line (C1) is applied in the form $\alpha<1$: together with $1-\alpha<\varepsilon$ it gives $0<\varepsilon$, and for this $\alpha\le 1$ would suffice.
In Section~1 the condition is applied only in the proof of Proposition~18.1\,(c), where $0<\alpha$ gives $-\alpha-\frac{15}{8}+\frac{3\varepsilon}{2}\le-\frac{3}{4}$;
in the proof of Proposition~18.1\,(a) the number $\alpha$ enters only through identities, which hold for every real $\alpha$.
In the proof of Proposition~18.2 the condition is applied in Steps~0, 1, 3, 5 and~6.

\item[$0<\varepsilon\le\bar{\varepsilon}=\frac{1}{200}$]
The inequalities $1-\alpha<\varepsilon\le\bar{\varepsilon}$ belong to the line (C7) of Section~I.8.2 and coincide with the row (S1) of Proposition~14.2\,(b);
the number $\bar{\varepsilon}=\frac{1}{200}$ is defined in the line (C0);
and $0<\varepsilon$ follows from $1-\alpha<\varepsilon$ together with the inequality $\alpha<1$ of the line (C1).
The condition is applied in Section~0.4, in the proof of Proposition~18.1\,(b), (c), and in the proof of Proposition~18.2.

\item[$e_0^{\varepsilon}\le\frac{1}{4}$, in the form (nm5)]
It belongs to the line (C8) of Section~I.8.2 and, for $v=-\log e_0$, coincides with the first inequality of the row (A8) of Proposition~14.2\,(b).
It is applied in the last sentence of the proof of Proposition~18.1\,(c).

\item[$160\,L^{1+\frac{3\varepsilon}{2}}\cdot M\cdot e_0^{\varepsilon}\le 1$, in the form (nm6)]
It belongs to the line (C8) and, for $v=-\log e_0$, coincides with the first inequality of the row (A9).
It is applied in Steps~4, 5, 6 and~7 of the proof of Proposition~18.2; in Step~7 it enters (7d) for the first slot.
It also gives $e_k\le 1$.

\item[$9M\cdot L^{\varepsilon}\cdot e_0^{2\varepsilon}\le 1$, in the form (nm7)]
It belongs to the line (C8) and, for $v=-\log e_0$, coincides with the third inequality of the row (A9).
It is applied in Step~7 of the proof of Proposition~18.2, in (7d) for the second slot.
\end{description}

The hypotheses on the data $\ell$ and $e_k$ enter as follows.
The side $\ell$ lies in $\mathrm{Cl}$; this enters through the two consequences that $\ell$ is a multiple of $L\cdot M$, so that $\ell/L$ is a multiple of $M$, and that $\ell\ge 13L\cdot M$.
In the setting (G2.S-2), moreover, $\ell/L$ lies in $\mathrm{Cl}$; this enters through $\ell/L\ge 13L\cdot M$, that is $\ell\ge 13L^2\cdot M$, and through the formation of $D_{k+1}$ and of the domains of level $k+1$ with the side $\ell/L$.
In the setting (G2.S-2) also $e_{k+1}=L^{\frac{1}{2}}e_k\le e_0$, which enters through $e_k\le e_0$.
Propositions~18.3\,(a), (b) and~18.4\,(b), (c) hold for all values of the constants for which the objects occurring in them are defined.

\subsection*{Appendix A. Proofs of (f1)--(f3), (qu2) and (qu3)}

The statements (f1)--(f3), (qu2) and (qu3) of Section~2.1 are, respectively, the properties (f1)--(f3) of the cut-off functions in \cite[Section~13.6]{QEDI}, \cite[Lemma~13.9]{QEDI} and \cite[Lemma~13.10]{QEDI}.
In the setting of this chapter, where a bond, a difference quotient or a H\"older difference lies inside a set of sites in the sense of (G2.r4) and the balls are described by Section~0.3, their proofs are as follows.

Throughout, $k\ge 0$, $\ell$ is a side, $\tilde{c}$ is a photon site of $T_k$, and $r>0$ is a real number with $4r+11<\ell$.
The arguments below hold for every such $r$, and so they cover (qu2), which is stated for $r\ge 50$.
Put
\[
\mathrm{K5}:=K(\tilde{c},2r+5).
\]
Since $2(2r+5)+1<\ell$, Section~0.3 applies to $\mathrm{K5}$.

\paragraph{The cube $\mathrm{Cu5}$.}
Let $\mathrm{Cu5}$ be the set of the points $\mathrm{pt}$ of $\mathbb{R}^3/\ell\mathbb{Z}^3$ whose sup distance from $\tilde{c}$ is at most $2r+5+\frac{1}{2}$.
It is the image of the closed cube of $\mathbb{R}^3$ of side $4r+11$ centred at a lift of $\tilde{c}$, and the projection is injective on this cube, because $4r+11<\ell$.
For $\mathrm{pt}\in\mathrm{Cu5}$ let $\mathrm{pt}-\tilde{c}$ denote the vector of sup norm at most $2r+5+\frac{1}{2}$ that is read off in this cube.
As $2r+5+\frac{1}{2}<\frac{\ell}{2}$, its sup norm $\left|\mathrm{pt}-\tilde{c}\right|$ equals the sup distance of $\mathrm{pt}$ from $\tilde{c}$.
For $\mathrm{pt},\mathrm{pt}'\in\mathrm{Cu5}$ put
\[
\mathrm{pt}-\mathrm{pt}':=\left(\mathrm{pt}-\tilde{c}\right)-\left(\mathrm{pt}'-\tilde{c}\right).
\]
For a site $c$ of $\mathrm{K5}$ the vector $c-\tilde{c}$ is the one of (G2.r2).
The cube is a product of three intervals; consequently a point each of whose coordinates lies between the corresponding coordinates of two points of $\mathrm{Cu5}$ belongs to $\mathrm{Cu5}$.
Let $\mathrm{pt}\in\mathrm{Cu5}$ with $\left|\mathrm{pt}-\tilde{c}\right|\ge 2r$.
Then one of the three numbers $\left(\mathrm{pt}-\tilde{c}\right)_{\mu}/r$ has modulus at least $2$, and therefore
\[
\prod_{\mu}\mathrm{ctf}\left(\frac{\left(\mathrm{pt}-\tilde{c}\right)_{\mu}}{r}\right)=0.
\]
If $\mathrm{pt}\notin\mathrm{Cu5}$, then $\mathrm{cut}_r(\mathrm{pt})=0$ by the definition (qu1).
Hence
\begin{equation}
\begin{aligned}
&\mathrm{cut}_r(\mathrm{pt})=\prod_{\mu}\mathrm{ctf}\left(\frac{\left(\mathrm{pt}-\tilde{c}\right)_{\mu}}{r}\right)\qquad\text{for all }\mathrm{pt}\in\mathrm{Cu5},\\
&\mathrm{cut}_r(\mathrm{pt})=0\qquad\text{if the sup distance of }\mathrm{pt}\text{ from }\tilde{c}\text{ is at least }2r;
\end{aligned}
\tag{A.0}
\end{equation}
for $2r+5<\left|\mathrm{pt}-\tilde{c}\right|\le 2r+5+\frac{1}{2}$ both sides of the first identity vanish.

\emph{The points that occur.}
Let $c$ be a site of $K(\tilde{c},2r+4)$, let $\mu$, $\nu$ be directions, the case $\nu=\mu$ included, and let $0\le\mathrm{sx}\le 1$.
The point $c+\frac{\eta_k}{2}e_{\mu}+\mathrm{sx}\cdot\eta_k e_{\nu}$ has sup distance at most $2r+4+\frac{3\eta_k}{2}\le 2r+5+\frac{1}{2}$ from $\tilde{c}$; thus it lies in $\mathrm{Cu5}$.
At level~$0$ and for $\nu=\mu$ one has $\frac{3\eta_0}{2}=\frac{3}{2}$, and this distance can exceed $2r+5$; the proofs of (f2) and (f3) below therefore work with $\mathrm{Cu5}$, on which (A.0) gives $\mathrm{cut}_r$ by the product formula.
For $k\ge 1$ one has $\frac{3\eta_k}{2}<1$.

\begin{proof}[Proof of (f1)]
Since $0\le\mathrm{ctf}\le 1$, one has $0\le\mathrm{fc}_{\mu}\le 1$.
Suppose that $c$ and $c+\eta_k e_{\mu}$ lie in $K(\tilde{c},r)$.
By (bx1), every coordinate of $c+\frac{\eta_k}{2}e_{\mu}-\tilde{c}$ lies between the corresponding coordinates of $c-\tilde{c}$ and of $c+\eta_k e_{\mu}-\tilde{c}$, and hence has modulus at most $r$.
As $\mathrm{ctf}=1$ on $[-1,1]$, every factor of the product equals $1$, and $\mathrm{fc}_{\mu}(c)=1$.
Suppose now that $d(c,\tilde{c})\ge 2r+1$.
Then the barycentre $c+\frac{\eta_k}{2}e_{\mu}$ has sup distance at least $2r+1-\frac{\eta_k}{2}\ge 2r$ from $\tilde{c}$, and $\mathrm{fc}_{\mu}(c)=0$ by (A.0).
\end{proof}

\begin{proof}[Proof of (f2)]
Let $\mathrm{pt},\mathrm{pt}'\in\mathrm{Cu5}$, and go from $\mathrm{pt}'$ to $\mathrm{pt}$ by changing the coordinates one after the other; the intermediate points lie in $\mathrm{Cu5}$.
For real $\mathrm{rv}$, $\mathrm{rv}'$ one has $\left|\mathrm{ctf}(\mathrm{rv})-\mathrm{ctf}(\mathrm{rv}')\right|\le 2\left|\mathrm{rv}-\mathrm{rv}'\right|$, and the other two factors of the product lie in $[0,1]$.
By (A.0) it follows that
\[
\left|\mathrm{cut}_r(\mathrm{pt})-\mathrm{cut}_r(\mathrm{pt}')\right|\le\frac{2}{r}\left|\mathrm{pt}-\mathrm{pt}'\right|_1.
\]
Let $c,c'\in K(\tilde{c},2r+4)$.
The barycentres $\mathrm{pt}=c+\frac{\eta_k}{2}e_{\mu}$ and $\mathrm{pt}'=c'+\frac{\eta_k}{2}e_{\mu}$ lie in $\mathrm{Cu5}$, and $\mathrm{pt}-\mathrm{pt}'=c-c'$.
Therefore
\[
\left|\mathrm{fc}_{\mu}(c)-\mathrm{fc}_{\mu}(c')\right|\le\frac{2}{r}\left|c-c'\right|_1,\qquad\left|c-c'\right|_1\le 3\left|c-c'\right|.
\]
The points $c+\frac{\eta_k}{2}e_{\mu}$ and $c+\frac{\eta_k}{2}e_{\mu}+\eta_k e_{\nu}$ lie in $\mathrm{Cu5}$ and differ by $\eta_k e_{\nu}$.
Hence
\[
\left|\partial_{\nu}\mathrm{fc}_{\mu}(c)\right|\le\eta_k^{-1}\cdot\frac{2}{r}\cdot\eta_k=\frac{2}{r}.
\]
\end{proof}

\begin{proof}[Proof of (f3)]
For $\mathrm{pt}\in\mathrm{Cu5}$ define
\[
\mathrm{dcut}_{\nu}(\mathrm{pt}):=r^{-1}\,\mathrm{ctf}'\left(\frac{\left(\mathrm{pt}-\tilde{c}\right)_{\nu}}{r}\right)\cdot\prod_{\mu'\ne\nu}\mathrm{ctf}\left(\frac{\left(\mathrm{pt}-\tilde{c}\right)_{\mu'}}{r}\right).
\]
As a function of the vector $\mathrm{pt}-\tilde{c}$, this is the partial derivative with respect to the coordinate $\nu$ of the continuously differentiable function $\prod_{\mu}\mathrm{ctf}\left((\,\cdot\,)_{\mu}/r\right)$ on $\mathbb{R}^3$.
In the coordinate $\nu$ it is Lipschitz continuous with the constant $r^{-1}\cdot\frac{4}{r}$, since $\mathrm{ctf}'$ has the Lipschitz constant $4$ and the other factors lie in $[0,1]$.
In each of the two other coordinates it is Lipschitz continuous with the constant $r^{-1}\cdot 2\cdot\frac{2}{r}$, since $\left|\mathrm{ctf}'\right|\le 2$, $\mathrm{ctf}$ has the Lipschitz constant $2$ and the third factor lies in $[0,1]$.
Changing the coordinates one after the other, as in the proof of (f2), one obtains
\[
\left|\mathrm{dcut}_{\nu}(\mathrm{pt})-\mathrm{dcut}_{\nu}(\mathrm{pt}')\right|\le\frac{4}{r^2}\left|\mathrm{pt}-\mathrm{pt}'\right|_1\qquad\text{for }\mathrm{pt},\mathrm{pt}'\in\mathrm{Cu5}.
\]
Let $c,c'\in K(\tilde{c},2r+4)$.
The segment
\[
\left\{c+\frac{\eta_k}{2}e_{\mu}+\mathrm{sx}\cdot\eta_k e_{\nu}:0\le\mathrm{sx}\le 1\right\}
\]
lies in $\mathrm{Cu5}$.
By (A.0) and the fundamental theorem of calculus on this segment,
\begin{align*}
\left(\partial_{\nu}\mathrm{fc}_{\mu}\right)(c)
&=\eta_k^{-1}\left[\mathrm{cut}_r\left(c+\frac{\eta_k}{2}e_{\mu}+\eta_k e_{\nu}\right)-\mathrm{cut}_r\left(c+\frac{\eta_k}{2}e_{\mu}\right)\right]\\
&=\int_0^1\mathrm{d}\,\mathrm{sx}\;\mathrm{dcut}_{\nu}\left(c+\frac{\eta_k}{2}e_{\mu}+\mathrm{sx}\cdot\eta_k e_{\nu}\right).
\end{align*}
For each value of $\mathrm{sx}$, the point that belongs to $c$ and the point that belongs to $c'$ differ by $c-c'$.
Consequently
\[
\left|\partial_{\nu}\mathrm{fc}_{\mu}(c)-\partial_{\nu}\mathrm{fc}_{\mu}(c')\right|\le\frac{4}{r^2}\left|c-c'\right|_1\le\frac{12}{r^2}\left|c-c'\right|.
\]
\end{proof}

\begin{proof}[Proof of (qu2)]
If $d(c,\tilde{c})<2r+4$, the bond $(c,c+\eta_k e_{\mu})$ lies inside $\mathrm{K5}$, because its second end has distance less than $2r+4+\eta_k\le 2r+5$ from $\tilde{c}$.
Since $\left|\mathrm{fc}_{\mu}\right|\le 1$, this gives $\left|\mathrm{cut}_r\,\mathrm{bq}\right|\le\mathrm{bn}_0$.
By (f1), $\mathrm{cut}_r\,\mathrm{bq}=\mathrm{bq}$ on the bonds inside $K(\tilde{c},r)$.
For $2r+1\le d(c,\tilde{c})<2r+4$ both $(\mathrm{cut}_r\,\mathrm{bq})_{\mu}(c)$ and $\mathrm{fc}_{\mu}(c)\cdot\mathrm{bq}_{\mu}(c)$ vanish, by (f1).
Therefore
\begin{gather*}
(\mathrm{cut}_r\,\mathrm{bq})_{\mu}(c)=\mathrm{fc}_{\mu}(c)\cdot\mathrm{bq}_{\mu}(c)\qquad\text{if }d(c,\tilde{c})<2r+4,\\
(\mathrm{cut}_r\,\mathrm{bq})_{\mu}(c)=0\qquad\text{if }d(c,\tilde{c})\ge 2r+1.
\end{gather*}
Put $c^{+}:=c+\eta_k e_{\nu}$.

\emph{Second bound.}
The quotient $\left(\partial_{\nu}(\mathrm{cut}_r\,\mathrm{bq})_{\mu}\right)(c)$ is formed with the values at $c$ and at $c^{+}$.
If $d(c,\tilde{c})\ge 2r+2$, both sites have distance at least $2r+1$ from $\tilde{c}$, and both values vanish.
If $d(c,\tilde{c})<2r+3$, the sites $c$, $c^{+}$ have distance less than $2r+3+\eta_k\le 2r+4$ from $\tilde{c}$, and the corners $c$, $c+\eta_k e_{\mu}$, $c+\eta_k e_{\nu}$, $c+\eta_k e_{\mu}+\eta_k e_{\nu}$ of the cell at $c$ have distance less than $2r+3+2\eta_k\le 2r+5$ from $\tilde{c}$.
So the cell lies inside $\mathrm{K5}$, and
\begin{equation}
\begin{split}
\left(\partial_{\nu}(\mathrm{cut}_r\,\mathrm{bq})_{\mu}\right)(c)
&=\left(\partial_{\nu}(\mathrm{fc}_{\mu}\cdot\mathrm{bq}_{\mu})\right)(c)\\
&=\mathrm{fc}_{\mu}(c^{+})\cdot\left(\partial_{\nu}\mathrm{bq}_{\mu}\right)(c)+\mathrm{bq}_{\mu}(c)\cdot\left(\partial_{\nu}\mathrm{fc}_{\mu}\right)(c),
\end{split}
\tag{A.1}
\end{equation}
whose modulus is at most $\mathrm{bn}_1+\frac{2}{r}\mathrm{bn}_0$ by (f1) and (f2).
Every site satisfies $d(c,\tilde{c})\ge 2r+2$ or $d(c,\tilde{c})<2r+3$; this proves the second bound.
In the bound $2r+3+2\eta_k\le 2r+5$ for the distance of the corners, the number $2r+5$ cannot in general be replaced by $2r+4$: at level~$0$ and for $\nu=\mu$ the corner $c+2e_{\mu}$ can have distance at least $2r+4$ from $\tilde{c}$, for instance distance $104$ when $r=50$ and $c-\tilde{c}=102\,e_{\mu}$.
For $\nu\ne\mu$ at every level, and for $k\ge 1$, where $2\eta_k<1$, all corners have distance less than $2r+4$.

\emph{Third bound.}
Let $0<d(c,c')\le 1$.
If $d(c,\tilde{c})\ge 2r+2$ and $d(c',\tilde{c})\ge 2r+2$, both quotients vanish.
Otherwise both sites have distance less than $2r+3$ from $\tilde{c}$, (A.1) holds at both of them, and $\left|c-c'\right|=d(c,c')$ by Section~0.3.
Omit the indices $\mu$, $\nu$ and put $c'^{+}:=c'+\eta_k e_{\nu}$.
Then
\begin{multline*}
\partial(\mathrm{fc}\cdot\mathrm{bq})(c)-\partial(\mathrm{fc}\cdot\mathrm{bq})(c')
=\mathrm{fc}(c^{+})\cdot\left[\partial\mathrm{bq}(c)-\partial\mathrm{bq}(c')\right]+\left[\mathrm{fc}(c^{+})-\mathrm{fc}(c'^{+})\right]\cdot\partial\mathrm{bq}(c')\\
+\left[\mathrm{bq}(c)-\mathrm{bq}(c')\right]\cdot\partial\mathrm{fc}(c)+\mathrm{bq}(c')\cdot\left[\partial\mathrm{fc}(c)-\partial\mathrm{fc}(c')\right].
\end{multline*}
Moreover
\[
\left|\mathrm{bq}_{\mu}(c)-\mathrm{bq}_{\mu}(c')\right|\le\mathrm{bn}_1\cdot\left|c-c'\right|_1\le 3\,\mathrm{bn}_1\cdot\left|c-c'\right|.
\]
Indeed, take a rectilinear path from $c'$ to $c$; it has $\left|c-c'\right|_1/\eta_k$ bonds.
Along a bond $(c'',c''+\eta_k e_{\mu'})$ of the path, $\mu'$ being its direction, $\mathrm{bq}_{\mu}$ changes by $\pm\eta_k\left(\partial_{\mu'}\mathrm{bq}_{\mu}\right)(c'')$, in whichever sense the bond is traversed.
The points $c''$ and $c''+\eta_k e_{\mu'}$ are sites of the path, so each of their coordinates lies between the corresponding coordinates of $c'$ and $c$, and both have distance less than $2r+3$ from $\tilde{c}$.
The corners of the cell of the quotient $\left(\partial_{\mu'}\mathrm{bq}_{\mu}\right)(c'')$ arise from these two sites by adding $0$ or $\eta_k e_{\mu}$; hence they have distance less than $2r+3+\eta_k\le 2r+4$ from $\tilde{c}$, also for $\mu'=\mu$, and the cell lies inside $\mathrm{K5}$.
So each bond of the path contributes at most $\eta_k\,\mathrm{bn}_1$.
The first bound of (f2) is applied at $c^{+}$ and $c'^{+}$, which have distance less than $2r+4$ from $\tilde{c}$ and satisfy $c^{+}-c'^{+}=c-c'$.
The second bound of (f2) and the bound (f3) are applied at $c$ and $c'$, which have distance less than $2r+3$ from $\tilde{c}$.
All these sites lie in $K(\tilde{c},2r+4)$.
By (f1)--(f3), the moduli of the four terms are at most
\[
\mathrm{bn}_{\alpha}\cdot\left|c-c'\right|^{\alpha},\qquad
\frac{6}{r}\mathrm{bn}_1\cdot\left|c-c'\right|,\qquad
3\,\mathrm{bn}_1\cdot\left|c-c'\right|\cdot\frac{2}{r},\qquad
\mathrm{bn}_0\cdot\frac{12}{r^2}\left|c-c'\right|,
\]
and $\left|c-c'\right|\le\left|c-c'\right|^{\alpha}$, because $\left|c-c'\right|\le 1$ and $\alpha\le 1$.
Division by $d(c,c')^{\alpha}=\left|c-c'\right|^{\alpha}$ gives the third bound.
\end{proof}

\begin{proof}[Proof of (qu3)]
Two sites of $T_k$ have distance at most $\ell/2$.
Moving all coordinates simultaneously by at most $1$ in each step, $1$ being a multiple of $\eta_k$, one joins them by a chain of at most $\lceil\ell/2\rceil\le\ell$ steps of sup length at most $1$.
Hence $\left|\mathrm{fn}(c)-\mathrm{fn}(c')\right|\le\ell\cdot\mathrm{osc}$ for all sites $c$, $c'$.
Since $\sum_{c'}\mathrm{fn}(c')=0$, the number $\mathrm{fn}(c)$ is the mean of $\mathrm{fn}(c)-\mathrm{fn}(c')$ over the sites $c'$, and therefore $\left|\mathrm{fn}(c)\right|\le\ell\cdot\mathrm{osc}$.
For $\mathrm{fn}=\partial_{\nu}\mathrm{bq}_{\mu}$ one has
\[
\sum_{c}\mathrm{fn}(c)=\eta_k^{-1}\sum_{c}\left[\mathrm{bq}_{\mu}(c+\eta_k e_{\nu})-\mathrm{bq}_{\mu}(c)\right]=0,
\]
because $c\mapsto c+\eta_k e_{\nu}$ is a bijection of $T_k$ onto itself.
\end{proof}

\section*{Chapter 17. Transports, chains, the homomorphism $\boldsymbol{\Theta}$; the adjustment $E\mapsto E^{\mathrm{nat}}$: proofs of Propositions 17.1 and 17.2}

\subsection*{0. Conventions of this chapter; preliminaries for both proofs}

Throughout this chapter the setting is (G2.S-1) of Section~G2.0.4, with the data $(k,\ell,e_k)$, and the constants are those of Section~I.8.
For a side $\ell$, the torus $T_k$ formed with the side $\ell$ is denoted by $T_k(\ell)$.
A homomorphism is a map that is linear and multiplicative and sends $1$ to $1$; of the map $\mathrm{ev}$ in Proposition~17.2\,(a), only the linearity and the multiplicativity are needed.
The labels (3.1)--(3.4) refer to the displays of Proposition~17.1, and (G2.0.1) is the display of (G2.r7); this display and the identity $\|F\|_{\mathsf{w}}=\sum_{n,m}\|F^{[n,m]}\|_{\mathsf{w}}$ of (G2.r7) are assertions of Proposition~15\,(n).

In this chapter $z$, $z'$ denote points of $M\mathbb{Z}^3$, namely centres of $M$-cubes (Sections~I.4.2 and~I.5.6), and $w$ denotes a point of $\mathbb{Z}^3$ (Sections~I.1.9 and~I.6.3).
A fermion site $\tilde{x}$ is written $\tilde{x}=(\tilde{c},g')$, with its photon site $\tilde{c}=\pi(\tilde{x})$ and its label $g'\in O_{\mathrm{cub}}$; the letter $\tilde{g}$ denotes an element $(S,a)$ of $\mathbb{G}_k$ (Section~I.4.2).
Independent generators of a Grassmann algebra are written $\theta_i$, and $\theta(I)$ denotes their product, in increasing order, over a subset $I$ of the index set (Proposition~15\,(d), Section~I.5.1).
The letters $v$, $v'$, $v_i$, $v(I)$ denote elements of an algebra or of the linear span of the generators (Proposition~15\,(i)); the letter $v$ also denotes a vector of $\mathbb{R}^3$ or of $M\mathbb{Z}^3$ (Sections~I.1.3 and~I.1.9).
The number $\mathrm{sd}$ is that of the assertion on lifts in Section~I.1.1:
let $\tilde{Z}$ be a lift of a set $Z$ of sites of $T_k$, let $\mathrm{sd}>0$ satisfy $\operatorname{diam}(\tilde{Z})+\mathrm{sd}<\ell$, and let $c,c'\in Z$ satisfy $d(c,c')\le\mathrm{sd}$;
then, with $\tilde{c}$, $\tilde{c}'$ the points of $\tilde{Z}$ over $c$, $c'$, the difference $\tilde{c}-\tilde{c}'$ is the only representative of $c-c'$ modulo $\ell\mathbb{Z}^3$ of sup norm at most $\mathrm{sd}$, and $|\tilde{c}-\tilde{c}'|=d(c,c')$.
The objects $\mathrm{pr}$, $\tilde{\square}$, $\mathrm{Sm}(\tilde{\square})$ are those of Proposition~16.3\,(f), and $\tilde{B}$, $\mathrm{Grc}_k$, $\mathrm{cov}_\ell$, $T_v$, $\mathrm{rg}$ those of Section~I.1.9.
The letter $\tilde{U}$ denotes a finite set of sites of $\mathbb{L}_k$; in (Lf3), in (Ge1), in Proposition~17.2\,(c), and from Step~4 of the proof of Proposition~17.2\,(f) on, it is the $\eta_k$-neighbourhood of $\tilde{X}$.
The letters $A_0$, $A_1$, $A'$, $\lambda$ refer to a splitting as in \cite[Definition~5.16]{QEDI}.

\subsubsection*{0.1. Covering map, periodic extension, signs}

Let $\mathrm{pr}\colon\mathbb{L}_k\to T_k$ be the covering map, and on $\mathbb{L}^{\mathrm{f}}_k$ let $\mathrm{pr}(c,g):=(\mathrm{pr}\,c,g)$; a pair $(x,x')$ of $\mathrm{Pzt}_k$ is sent to its class modulo $\ell\mathbb{Z}^3$ in $\mathrm{Pz}_k$.
A function of the sites or of the bonds of $T_k$ is evaluated at the sites and bonds of $\mathbb{L}_k$ through $\mathrm{pr}$, that is, through its $\ell\mathbb{Z}^3$-periodic extension.
Accordingly, let $B$ be a bond function on $T_k$, let $c$, $c'$ be points of $\mathbb{L}_k$ with $|c-c'|\le M$, and let $\pi$ be an order of the three coordinates.
Then $B\left(\boldsymbol{\Gamma}^{\pi}(c,c')\right)$ is the sum of the numbers $B(\mathrm{pr}\,b)$ over the bonds $b$ of the path $\boldsymbol{\Gamma}^{\pi}(c,c')$ of $\mathbb{L}_k$, every bond being oriented along the path, and
\begin{equation}
(\tau B)(c,c')=\frac{1}{6}\sum_{\pi}B\left(\boldsymbol{\Gamma}^{\pi}(c,c')\right),\qquad(\tau B)(c,c)=0.
\tag{0.a}
\end{equation}
As $M<\ell/2$, the pair $(c,c')$ is a lift of $(\mathrm{pr}\,c,\mathrm{pr}\,c')$ with $|c-c'|<\ell/2$; hence (0.a) is the axial functional of Section~I.1.1.
The functional $\tau B$ depends linearly on $B$, with real coefficients.

For every oriented bond $(c,c')$ of $\mathbb{L}_k$,
\[
(d\Lambda)(c,c')=\eta_k^{-1}\cdot\left(\Lambda(c')-\Lambda(c)\right).
\]
For $c'=c+\eta_k\cdot e_\mu$ this is the definition $(d\Lambda)_\mu=\partial_\mu\Lambda$ of Section~I.1.1; for $c'=c-\eta_k\cdot e_\mu$ it follows from the rule $A(c',c)=-A(c,c')$ for bond functions.

For $w\in\mathbb{Z}^3$ let $\boldsymbol{\chi}(w):=(-1)^{w_1+w_2+w_3}$, as in Section~I.6.3.

\subsubsection*{0.2. $M$-cubes}

For $z\in M\mathbb{Z}^3$, the $M$-cube with centre $z$ is the set $\{c\in\mathbb{L}_k:|c-z|<M/2\}$.
By the lines (C3) and (C5) of Section~I.8.2, $M$ is a power of $L$, and $L$ is odd.
Hence $M\mathbb{Z}^3\subset\mathbb{Z}^3\subset\mathbb{L}_k=L^{-k}\mathbb{Z}^3$, and $M/\eta_k=M\cdot L^k$ is an odd integer.
Consequently, for $c\in\mathbb{L}_k$ every difference $c_\mu-z_\mu$ is a multiple of $\eta_k$ different from $\pm M/2$ (Section~I.1.1, paragraph ``Faces''), and the following assertions hold.
\begin{itemize}
\item[(Cu1)] The cube is the product of the three sets $\left\{z_\mu+\eta_k\cdot j:\ j\in\mathbb{Z},\ |j|\le\frac{1}{2}\left(M\cdot L^k-1\right)\right\}$, $\mu=1,2,3$, and its sites satisfy $|c-z|\le\frac{1}{2}\left(M-\eta_k\right)$.
\item[(Cu2)] The cube has $\left(M\cdot L^k\right)^3$ photon sites, and over them lie $48\left(M\cdot L^k\right)^3$ fermion sites (Sections~I.1.3 and~I.1.4).
\item[(Cu3)] Two different $M$-cubes of $\mathbb{L}_k$ with centres $z$, $z'$ have a common face if and only if $z'-z=M\cdot e_\mu$ or $z'-z=-M\cdot e_\mu$ for some direction $\mu$.
This is the definition of a common face in Section~I.1.1 (paragraph ``Cubes, polymers, parameter cubes''), and it coincides with the geometric notion.
Indeed, the closed $M$-cubes are the sets $z+[-M/2,M/2]^3$, $z\in M\mathbb{Z}^3$ \cite[Section~13.3]{QEDI}.
For $z\ne z'$ the intersection of the two closed cubes is the product over $\mu$ of the sets $\left[z_\mu-\frac{M}{2},z_\mu+\frac{M}{2}\right]\cap\left[z'_\mu-\frac{M}{2},z'_\mu+\frac{M}{2}\right]$; such a set is an interval of length $M$ if $z'_\mu=z_\mu$, a single point if $z'_\mu-z_\mu$ equals $M$ or $-M$, and empty otherwise.
So the intersection is a face, that is, a square of side $M$, if and only if exactly one coordinate of $z'-z$ equals $M$ or $-M$ and the other two vanish.
\end{itemize}

\subsubsection*{0.3. The lift of a small polymer}

Let $X\in D_k$ be small, that is, $d_M(X)\le L$.
Since $\ell\in\mathrm{Cl}$, one has $\ell\ge13L\cdot M$; moreover $L\ge5$ and $M\ge1\ge\eta_k$.
Therefore
\begin{equation}
3M\cdot(L+2)<\ell,\qquad M\cdot(L+2)+M+2\eta_k<\ell;
\tag{0.b}
\end{equation}
indeed, $3M\cdot(L+2)=3L\cdot M+6M\le5L\cdot M$, and $M\cdot(L+2)+M+2\eta_k\le M\cdot(L+3)+2\le M\cdot(L+5)\le2L\cdot M$, while $5L\cdot M$ and $2L\cdot M$ are smaller than $13L\cdot M\le\ell$.
\begin{itemize}
\item[(Lf1)] (Existence; diameter.)
By \cite[Theorem~13.2(t7)]{QEDI}, any two sites $c$, $c'$ of $X$ satisfy $d(c,c')\le M\cdot\left(d_M(X)+2\right)\le M\cdot(L+2)$.
Fix a site $c^0$ of $X$ and a point $\tilde{c}^0$ of $\mathbb{L}_k$ over it.
Every site $c$ of $X$ has a point $\tilde{c}$ of $\mathbb{L}_k$ over it with $|\tilde{c}-\tilde{c}^0|=d(c,c^0)\le M\cdot(L+2)$, and only one, because $2M\cdot(L+2)<\ell$ (two such points differ by a vector of $\ell\mathbb{Z}^3$ of sup norm at most $2M\cdot(L+2)$).
The set $\tilde{X}$ of these points is mapped onto $X$, and its sup-diameter is at most $2M\cdot(L+2)$; since $2M\cdot(L+2)+\eta_k\le3M\cdot(L+2)<\ell$ by (0.b), $\tilde{X}$ is a lift of $X$ in the sense of Section~I.1.1.
Put $\mathrm{sd}:=M\cdot(L+2)$.
Then the sup-diameter of $\tilde{X}$ plus $\mathrm{sd}$ is at most $3M\cdot(L+2)<\ell$, and Section~I.1.1 gives $|\tilde{c}-\tilde{c}'|=d(c,c')\le M\cdot(L+2)$ for any two points $\tilde{c}$, $\tilde{c}'$ of $\tilde{X}$, lying over the sites $c$, $c'$ of $X$.
Hence the sup-diameter of $\tilde{X}$ is at most $M\cdot(L+2)$.
\item[(Lf2)] (Uniqueness.)
A polymer is connected through nearest neighbours, so by Section~I.1.1 every lift of $X$ has the form $\tilde{X}+\ell\cdot w$ with $w\in\mathbb{Z}^3$.
Consequently every lift of $X$ has sup-diameter at most $M\cdot(L+2)$, and (Lf3)--(Lf5) hold for every lift $\tilde{X}$.
\item[(Lf3)] Let $\tilde{U}:=\left\{c\in\mathbb{L}_k:\operatorname{dist}(c,\tilde{X})\le\eta_k\right\}$, the distance being the sup distance.
The sup-diameter of $\tilde{U}$ is at most $M\cdot(L+2)+2\eta_k<\ell$.
Hence $\mathrm{pr}$ is injective on $\tilde{U}$, in particular on $\tilde{X}$, and $(c,g)\mapsto(\mathrm{pr}\,c,g)$ is a bijection of $\tilde{X}\times O_{\mathrm{cub}}$ onto $X\times O_{\mathrm{cub}}$.
\item[(Lf4)] By (0.b), the sup-diameter of $\tilde{X}$ plus $M$ is smaller than $\ell$.
Section~I.1.1, applied with $\mathrm{sd}=\eta_k$, shows that the bonds with both ends in $X$ are the images of the bonds with both ends in $\tilde{X}$.
Section~I.1.1, applied with $\mathrm{sd}=M$, shows that for sites $c$, $c'$ of $X$ with $d(c,c')\le M$ the points $\tilde{c}$, $\tilde{c}'$ of $\tilde{X}$ over them satisfy $|\tilde{c}-\tilde{c}'|=d(c,c')$, and that the six paths of the pair $(c,c')$ are the images of the paths $\boldsymbol{\Gamma}^{\pi}(\tilde{c},\tilde{c}')$ of $\mathbb{L}_k$.
\item[(Lf5)] By the paragraph on lifts in Section~I.1.1, a lift of a union of $M$-cubes is a union of $M$-cubes of $\mathbb{L}_k$ with the same common faces.
Applied to $X$, this shows that $\tilde{X}$ is a union of $M$-cubes of $\mathbb{L}_k$, that the map $\tilde{\square}\mapsto\mathrm{pr}(\tilde{\square})$ is a bijection of the $M$-cubes of $\tilde{X}$ onto the $M$-cubes of $X$, and that two cubes of $\tilde{X}$ have a common face if and only if their images have one.
Hence $|\tilde{X}|_M=|X|_M$, the set $\tilde{X}$ is connected through faces, and the chains of $M$-cubes of $X$ in which consecutive cubes have a common face are the images of the chains of $M$-cubes of $\tilde{X}$ with the same property.
\end{itemize}
For points over $\tilde{X}$, the words ``the cube of $c$'', the cubes $\square_i$ and the words ``the centre $z_i$'' always refer to $M$-cubes of $\tilde{X}$ and to their centres, which are points of $M\mathbb{Z}^3$.
Thus (Lf1), (Lf2) and the second inequality of (0.b) yield the following:
a small polymer $X$ has a lift $\tilde{X}$, unique up to a translation by a vector of $\ell\mathbb{Z}^3$, whose sup-diameter is at most $M\cdot(L+2)$, and $M\cdot(L+2)+M+2\eta_k<\ell$.

\subsubsection*{0.4. Generators over a lift; the subalgebra $\mathrm{Grp}[X]$; norms}

By Section~I.3.6 and (I.AP), the generators are $\psi(x,\beta,\omega)$, $x\in\mathbb{L}^{\mathrm{f}}_k$, and $\chi(x,x',\beta,\omega)$, $(x,x')\in\mathrm{Pzt}_k$, and for all $w\in\mathbb{Z}^3$
\begin{equation}
\begin{aligned}
\psi(x+\ell\cdot w,\beta,\omega)&=\boldsymbol{\chi}(w)\cdot\psi(x,\beta,\omega),\\
\chi(x+\ell\cdot w,x'+\ell\cdot w,\beta,\omega)&=\boldsymbol{\chi}(w)\cdot\chi(x,x',\beta,\omega);
\end{aligned}
\tag{0.c}
\end{equation}
independent generators are obtained by restricting the photon site of the first point $x$ to a fundamental domain of $\ell\mathbb{Z}^3$ in $\mathbb{L}_k$.
\begin{itemize}
\item[(Ge1)] Every set of sites of $\mathbb{L}_k$ on which $\mathrm{pr}$ is injective is contained in a fundamental domain: it suffices to add one point over each site of $T_k$ that is not in its image.
For a fundamental domain containing $\tilde{U}$ (see (Lf3)), the elements $\psi(x,\beta,\omega)$ with $x=(c,g)$, $c\in\tilde{U}$, and the elements $\chi(x,x',\beta,\omega)$ with $(x,x')\in\mathrm{Pzt}_k$, $x=(c,g)$, $c\in\tilde{U}$, belong to the independent generators.
A change of the fundamental domain multiplies generators by signs.
Hence none of the following depends on the fundamental domain: the bidegree; the coefficients of an element, up to signs; the norms (I.5.1); the linear span $\mathrm{Lin}$ of all generators $\psi(\xi)$, $\chi(\boldsymbol{\zeta})$; and the statement that a generator occurs in an element $F$, which means that the generator is, up to its sign, a factor of a basis monomial whose coefficient in $F$ is not zero.
\item[(Ge2)] For $v,v'\in\mathrm{Lin}$ one has $v\cdot v'+v'\cdot v=0$, in particular $v^2=0$.
Indeed, the expression $v\cdot v'+v'\cdot v$ is bilinear in $(v,v')$; the space $\mathrm{Lin}$ is spanned by independent generators $\theta_i$, and $\theta_i\cdot\theta_j+\theta_j\cdot\theta_i=0$ for all $i$, $j$ (for $i=j$ this reads $\theta_i^2=0$); so the expression vanishes on $\mathrm{Lin}\times\mathrm{Lin}$, and for $v'=v$ it equals $2v^2$.
\item[(Ge3)] Let $\mathrm{Grp}[X]$ be the subalgebra of $\mathrm{Gr}[\psi,\chi]$ generated by $1$ and the $\psi(\xi)$, $\xi$ over $X$.
By (0.c) and (Lf3), it is generated by $1$ and the independent generators $\psi(x,\beta,\omega)$ with $x$ over $\tilde{X}$, $\beta\in\{1,\dots,4\}$, $\omega\in\{0,1\}$.
\item[(Ge4)] The map $\|\cdot\|_{\mathsf{w}}$ is a norm on the vector space $\mathrm{Gr}[\psi,\chi]$, namely the sum of the moduli of the coefficients, a coefficient of bidegree $(n,m)$ being weighted by $\mathsf{w}_1^n\cdot\mathsf{w}_2^m$; moreover $\|\psi(x,\beta,\omega)\|_{\mathsf{w}}=\mathsf{w}_1$ and $\|\chi(x,x',\beta,\omega)\|_{\mathsf{w}}=\mathsf{w}_2$ for all indices on the covering lattices (Proposition~15\,(a) and (0.c)).
Hence the norm of a linear combination of elements of $\mathrm{Gr}[\psi,\chi]$ is at most the sum, over its terms, of the modulus of the coefficient times the norm of the element.
\item[(Ge5)] Let $\mathrm{Pty}$ be the automorphism of $\mathrm{Gr}[\psi,\chi]$ that multiplies every generator by $-1$.
\end{itemize}

\subsubsection*{0.5. Homomorphisms of Grassmann algebras}

\begin{named}{Lemma A}
Let $\mathrm{ng}\ge1$ be an integer, let $\theta_1,\dots,\theta_{\mathrm{ng}}$ be independent generators of a Grassmann algebra over $\mathbb{C}$, and let $\mathrm{Gr}'$ be the subalgebra generated by $1$ and $\theta_1,\dots,\theta_{\mathrm{ng}}$.
Let $\mathrm{Alg}$ be an associative algebra over $\mathbb{C}$ with identity element, and let $v_1,\dots,v_{\mathrm{ng}}\in\mathrm{Alg}$ satisfy $v_i\cdot v_j+v_j\cdot v_i=0$ for all $i$, $j$, the case $i=j$ included.
Then there is exactly one homomorphism $\mathrm{Hom}\colon\mathrm{Gr}'\to\mathrm{Alg}$ with $\mathrm{Hom}(\theta_i)=v_i$ for all $i$.
\end{named}

\begin{proof}
For a subset $I$ of $\{1,\dots,\mathrm{ng}\}$ let $\theta(I)$ be the product of the $\theta_i$, $i\in I$, taken in increasing order of $i$, and let $v(I)$ be the product of the $v_i$, $i\in I$, in the same order; put $\theta(\emptyset):=1$ and $v(\emptyset):=1$.
The monomials $\theta(I)$ form a basis of $\mathrm{Gr}'$.
A homomorphism with $\mathrm{Hom}(\theta_i)=v_i$ necessarily satisfies $\mathrm{Hom}(\theta(I))=v(I)$; this proves uniqueness.
For existence, let $\mathrm{Hom}$ be the linear map defined by $\mathrm{Hom}(\theta(I)):=v(I)$; in particular $\mathrm{Hom}(1)=v(\emptyset)=1$.
By hypothesis, interchanging two adjacent factors in a product of elements $v_j$ changes the sign of the product, and $v_i^2=0$ for every $i$; the same holds for products of the $\theta_j$.
Let $I$, $I'$ be two subsets.
If they have an element $i$ in common, then $\theta(I)\cdot\theta(I')=0$; likewise $v(I)\cdot v(I')=0$, since interchanges of adjacent factors bring the two factors $v_i$ next to each other.
If they are disjoint, then $\theta(I)\cdot\theta(I')=\mathrm{sn}\cdot\theta(I\cup I')$, where $\mathrm{sn}$ is the sign of the permutation that puts the sequence $(I,I')$ in increasing order; the same interchanges of adjacent factors give $v(I)\cdot v(I')=\mathrm{sn}\cdot v(I\cup I')$.
Thus $\mathrm{Hom}$ is multiplicative on the basis, and therefore, by bilinearity of the products, on $\mathrm{Gr}'$.
\end{proof}

\begin{named}{(Un)}
Let two maps of $\mathrm{Gr}'$ into an algebra over $\mathbb{C}$ be additive and multiplicative, let both be linear or both be antilinear over $\mathbb{C}$, and let them coincide on $\theta_1,\dots,\theta_{\mathrm{ng}}$.
Then they coincide on every linear combination of products of at least one generator; if both send $1$ to $1$, they coincide on $\mathrm{Gr}'$.
\end{named}

\begin{proof}
By multiplicativity, each of the two maps sends a product of at least one generator to the product, in the same order, of the images of its factors; so the two maps coincide on such products.
By additivity together with the common linearity, respectively antilinearity, each of them sends a finite linear combination of such products to the combination of their images with the same, respectively the complex conjugate, coefficients; so they coincide on such combinations.
The elements $\theta(I)$ span $\mathrm{Gr}'$, $\theta(\emptyset)=1$, and for $I\ne\emptyset$ the element $\theta(I)$ is a product of at least one generator (proof of Lemma~A); hence every element of $\mathrm{Gr}'$ is a multiple of $1$ plus a linear combination of such products.
If both maps send $1$ to $1$, the common linearity, respectively antilinearity, shows that they coincide on the multiples of $1$, and additivity shows that they coincide on $\mathrm{Gr}'$.
\end{proof}

\subsubsection*{0.6. The mean over the fermion sites over $X$}

Let $f$ be a function on $\tilde{X}\times O_{\mathrm{cub}}$ with values in a vector space.
By (I.1.3), by the identity $\operatorname{Vol}(X)=M^3\cdot|X|_M$, by (Cu2) and (Lf5), and since $\eta_k=L^{-k}$,
\begin{equation}
\begin{split}
\operatorname{Vol}(X)^{-1}\int_X\mathrm{d}\tilde{x}\,f(\tilde{x})
&=\left(M^3\cdot|X|_M\right)^{-1}\cdot\frac{\eta_k^3}{48}\sum_{\tilde{c}\in\tilde{X}}\ \sum_{g'\in O_{\mathrm{cub}}}f(\tilde{c},g')\\
&=\mathrm{cnt}(X)^{-1}\sum_{\tilde{c}\in\tilde{X}}\ \sum_{g'\in O_{\mathrm{cub}}}f(\tilde{c},g'),\qquad\mathrm{cnt}(X):=48\left(M\cdot L^k\right)^3\cdot|X|_M.
\end{split}
\tag{0.d}
\end{equation}
The number $\mathrm{cnt}(X)$ is the number of the fermion sites over $X$.
Thus $\operatorname{Vol}(X)^{-1}\int_X\mathrm{d}\tilde{x}$ is the arithmetic mean over these sites: over the $\left(M\cdot L^k\right)^3\cdot|X|_M$ photon sites of $\tilde{X}$ and, over each of them, over the $48$ labels.
In particular $\operatorname{Vol}(X)^{-1}\int_X\mathrm{d}\tilde{x}\,1=1$.

\subsection*{1. Proposition 17.1}

\subsubsection*{Proof of Proposition 17.1}

The notation is that of Section~0.
Let $\tilde{X}$ be the given lift of $X$, and for points $p$, $p'$ of $\mathbb{L}^{\mathrm{f}}_k$ let $c:=\pi(p)$ and $c':=\pi(p')$.

\begin{proof}[Proof of Proposition~17.1\,(a)]
Let $\mathscr{A}\in\tilde{R}_k$.
By \cite[Definition~5.16(b)]{QEDI}, taken with the factor $1$ (Section~I.5.3), one has $\mathscr{A}=A_0+A_1$ with $A_0$ real and $A_1\in R_k$; in particular $|A_1|_\infty<e_k^{-\frac{3}{4}+\varepsilon}$.
Hence $\operatorname{Im}\mathscr{A}=\operatorname{Im}A_1$, and
\[
e_k\cdot|\operatorname{Im}\mathscr{A}|_\infty\le e_k\cdot|A_1|_\infty<e_k^{\frac{1}{4}+\varepsilon}.
\]
Now let $\mathscr{C}\in\mathfrak{X}_k$.
The definition of $\mathfrak{X}_k$ (Section~I.5.3) gives $|\mathscr{C}|_\infty<e_k^{-1+2\varepsilon}$, so that $e_k\cdot|\operatorname{Im}\mathscr{C}|_\infty\le e_k\cdot|\mathscr{C}|_\infty<e_k^{2\varepsilon}$.
Since $\operatorname{Im}\mathscr{T}=\operatorname{Im}\mathscr{A}+\operatorname{Im}\mathscr{C}$, it follows that
\[
e_k\cdot|\operatorname{Im}\mathscr{T}|_\infty\le e_k\cdot|\operatorname{Im}\mathscr{A}|_\infty+e_k\cdot|\operatorname{Im}\mathscr{C}|_\infty<e_k^{\frac{1}{4}+\varepsilon}+e_k^{2\varepsilon}.
\]
Condition (A8), which stands in \cite[Section~5.6.2]{QEDI} (Section~I.8.4), contains the inequality $e_0^{\varepsilon}\le\frac{1}{4}$.
Hence $e_0<1$, and $e_k^{\varepsilon}\le e_0^{\varepsilon}\le\frac{1}{4}$, because $0<e_k\le e_0$ and $\varepsilon>0$.
As $0<e_k<1$ and $\varepsilon\le\frac{1}{200}\le\frac{1}{4}$, one has $e_k^{\frac{1}{4}}\le e_k^{\varepsilon}\le\frac{1}{4}$.
Therefore
\begin{gather*}
e_k^{\frac{1}{4}+\varepsilon}=e_k^{\frac{1}{4}}\cdot e_k^{\varepsilon}\le\frac{1}{4}\,e_k^{\varepsilon},\qquad e_k^{2\varepsilon}=e_k^{\varepsilon}\cdot e_k^{\varepsilon}\le\frac{1}{4}\,e_k^{\varepsilon},\\
e_k^{\frac{1}{4}+\varepsilon}+e_k^{2\varepsilon}\le\frac{1}{2}\,e_k^{\varepsilon}.
\end{gather*}
\end{proof}

\begin{proof}[Proof of Proposition~17.1\,(b)]
Let $(p,p')\in\mathrm{Pzt}_k$.
By (I.1.5), $|c-c'|\le d_X(p,p')\le M$; hence $(\tau B)(c,c')$ is given by (0.a), and $u_B(p,p')=U_B(c,c')$ is the number defined in (I.1.6).
\begin{itemize}
\item[(u1)] The path $\boldsymbol{\Gamma}^{\pi}(c,c')$ moves the coordinate $\mu$ from $c_\mu$ to $c'_\mu$ in $|c_\mu-c'_\mu|/\eta_k$ steps of length $\eta_k$; so it consists of $|c-c'|_1/\eta_k$ bonds.
By (0.a),
\begin{align*}
\left|\operatorname{Im}(\tau B)(c,c')\right|&\le\frac{1}{6}\sum_{\pi}\frac{|c-c'|_1}{\eta_k}\cdot|\operatorname{Im}B|_\infty=\frac{|c-c'|_1\cdot|\operatorname{Im}B|_\infty}{\eta_k},\\
\left|u_B(p,p')^{\pm1}\right|&=\exp\left(\mp e_k\cdot\eta_k\cdot\operatorname{Im}(\tau B)(c,c')\right)\le\exp\left(e_k\cdot|\operatorname{Im}B|_\infty\cdot|c-c'|_1\right).
\end{align*}
For $p=p'$ both sides of the inequality in (u1) equal $1$.
\item[(u2)] By (0.a), $e_k\cdot\eta_k\cdot(\tau B)(c,c')$ is a linear form in the numbers $B(\mathrm{pr}\,b)$, where $b$ ranges over the bonds of the six paths $\boldsymbol{\Gamma}^{\pi}(c,c')$; the number $u_B(p,p')$ is the exponential of $i$ times this form, hence an entire function of these numbers.
If $c=c'$, then $(\tau B)(c,c)=0$, and so $u_B(p,p')=1$.
\item[(u3)] The first relation: for $(p,p')\in\mathrm{Pzt}_k$ it is the identity $U_{\tilde{g}B}(\tilde{g}x,\tilde{g}x')=U_B(x,x')$ of Proposition~16.1\,(b); for $p=p'$ both sides equal $1$.

The second relation: by Section~0.1, the value of $d\Lambda$ on an oriented bond is $\eta_k^{-1}$ times the value of $\Lambda$ at the end point minus its value at the initial point.
Along a path $\boldsymbol{\Gamma}$ from $c$ to $c'$ the end point of each bond is the initial point of the next one, so the sum over the bonds of $\boldsymbol{\Gamma}$ telescopes: $(d\Lambda)(\boldsymbol{\Gamma})=\eta_k^{-1}\cdot\left(\Lambda(c')-\Lambda(c)\right)$.
This applies to each of the six paths $\boldsymbol{\Gamma}^{\pi}(c,c')$, and therefore $(\tau\,d\Lambda)(c,c')=\eta_k^{-1}\cdot\left(\Lambda(c')-\Lambda(c)\right)$.
If $B^1=B-d\Lambda$ on the bonds of the six paths, then (u2) and the linearity of $\tau$ give
\[
e_k\cdot\eta_k\cdot(\tau B^1)(c,c')=e_k\cdot\eta_k\cdot(\tau B)(c,c')-e_k\cdot\left(\Lambda(c')-\Lambda(c)\right),
\]
and the exponential of $i$ times this identity is the assertion.

The third relation follows from $\tau(-B)=-\tau B$.

The fourth relation: since $\tau$ has real coefficients, $\left(\tau(-\overline{B})\right)(c,c')=-\overline{(\tau B)(c,c')}$, and therefore
\begin{align*}
\overline{\exp\left\{i\cdot e_k\cdot\eta_k\cdot\left(\tau(-\overline{B})\right)(c,c')\right\}}
&=\exp\left\{-i\cdot e_k\cdot\eta_k\cdot\overline{\left(-\overline{(\tau B)(c,c')}\right)}\right\}\\
&=\exp\left\{i\cdot e_k\cdot\eta_k\cdot(\tau B)(c,c')\right\}.
\end{align*}
Taking inverses gives $\overline{u_{-\overline{B}}^{-1}}=u_B^{-1}$.
\item[(u4)] Multiplying the two relations by $d_X(p,p')^{\alpha}>0$ and solving for $\psi(p)$ and $\bar{\psi}(p)$ gives (3.1).
That the two relations read $\chi=\delta_{\alpha,B}\psi$ and $\bar{\chi}=\delta_{\alpha,-B}\bar{\psi}$ is (I.1.7), combined with $u_{-B}=u_B^{-1}$.
\end{itemize}
\end{proof}

\begin{proof}[Proof of Proposition~17.1\,(c)]
By \cite[Theorem~13.2(t1)]{QEDI} and $d_M(X)\le L$,
\[
|X|_M\le c_X\cdot\left(d_M(X)+1\right)\le c_X\cdot(L+1).
\]

\emph{The set $\mathrm{Ch}(x,\tilde{x})$.}
Its definition involves $x=(c,g)$ and $\tilde{x}=(\tilde{c},g')$ only through the cube of $c$ and the cube of $\tilde{c}$.
Since $\tilde{X}$ is connected through faces (Lf5), there is a chain of $M$-cubes of $\tilde{X}$, consecutive ones having a common face, from the cube of $c$ to the cube of $\tilde{c}$; hence there is a least number $n$ for which such a chain with $n+1$ members exists, and $\mathrm{Ch}(x,\tilde{x})$ is not empty.
The members of a shortest chain are pairwise different: if $\square_i=\square_j$ with $i<j$, then $(\square_0,\dots,\square_i,\square_{j+1},\dots,\square_n)$ is a chain with the same end cubes and fewer members.
Therefore $n+1\le|\tilde{X}|_M=|X|_M$, and $\mathrm{Ch}(x,\tilde{x})$, a set of sequences without repetitions in a finite set, is finite.
All elements of $\mathrm{Ch}(x,\tilde{x})$ have the same $n$.
By (Lf5), $\mathrm{pr}$ induces a bijection of $\mathrm{Ch}(x,\tilde{x})$ onto the set of the shortest chains of $M$-cubes of $X$ between the images of the two cubes.

The centres $z_i$ belong to $M\mathbb{Z}^3$, hence to $\mathbb{L}_k$ (Section~0.2); so the $p_i$ are points of $\mathbb{L}^{\mathrm{f}}_k$ over $\tilde{X}$.
\begin{itemize}
\item[(c1)] Let $(\square_0,\dots,\square_n)\in\mathrm{Ch}(x,\tilde{x})$.

\emph{The case $i=0$.}
Here $p_0=(c,g)$ and $p_1=(z_0,g)$ carry the same label.
The site $c$ lies in $\square_0$, so (Cu1) gives $|c-z_0|\le\frac{1}{2}\left(M-\eta_k\right)<M/2$.
If $p_0\ne p_1$, then $c\ne z_0$, hence $|c-z_0|\ge\eta_k$, and by (I.1.5) $d_X(p_0,p_1)=\max\left(|c-z_0|,\eta_k\right)=|c-z_0|$.
In each coordinate, every site of a path $\boldsymbol{\Gamma}^{\pi}(c,z_0)$ has a value in $\eta_k\mathbb{Z}$ between the values of $c$ and of $z_0$; as $\square_0$ is a product of three intervals (Cu1), the six paths lie in $\square_0$.

\emph{The case $1\le i\le n$.}
Here $p_i=(z_{i-1},g)$ and $p_{i+1}=(z_i,g)$ carry the same label.
The cubes $\square_{i-1}$ and $\square_i$ are different and have a common face; by (Cu3), $z_i-z_{i-1}$ equals $M\cdot e_\mu$ or $-M\cdot e_\mu$ for some direction $\mu$.
Hence $d_X(p_i,p_{i+1})=\max(M,\eta_k)=M$.
Only the coordinate $\mu$ changes, so each of the six paths is the straight path with the sites
\[
z_{i-1}+\frac{j}{M\cdot L^k}\cdot\left(z_i-z_{i-1}\right),\qquad j=0,1,\dots,M\cdot L^k.
\]
The site with the number $j$ has distance $\eta_k\cdot j$ from $z_{i-1}$ and distance $M-\eta_k\cdot j$ from $z_i$.
As $M\cdot L^k$ is odd (Section~0.2), $\eta_k\cdot j\ne M/2$.
If $\eta_k\cdot j<M/2$, the site lies in $\square_{i-1}$; if $\eta_k\cdot j>M/2$, then $M-\eta_k\cdot j<M/2$, and the site lies in $\square_i$.

\emph{The case $i=n+1$.}
Here $p_{n+1}=(z_n,g)$ and $p_{n+2}=(\tilde{c},g')$, where $\tilde{c}$ lies in $\square_n$ and $g'$ is an arbitrary element of $O_{\mathrm{cub}}$.
If $\tilde{c}\ne z_n$, then $\eta_k\le|\tilde{c}-z_n|\le\frac{1}{2}\left(M-\eta_k\right)<M/2$ by (Cu1), so that $d_X(p_{n+1},p_{n+2})=\max\left(|\tilde{c}-z_n|,\eta_k\right)=|\tilde{c}-z_n|$; in each coordinate, every site of a path $\boldsymbol{\Gamma}^{\pi}(z_n,\tilde{c})$ has a value in $\eta_k\mathbb{Z}$ between the values of $z_n$ and of $\tilde{c}$, and $\square_n$ is a product of three intervals (Cu1), so the six paths lie in $\square_n$.
If $\tilde{c}=z_n$ and $g'\ne g$, the two points are different points of the cluster of $z_n$; then $d_X=\max(0,\eta_k)=\eta_k\le M$, $u_B=1$ by (u2), there is no path, and the condition of kind (I) (Section~I.1.7) is that $z_n$, read through $\mathrm{pr}$, lies in $X$.
If $\tilde{c}=z_n$ and $g'=g$, the two points coincide.

\emph{Conclusion.}
In every case with $p_i\ne p_{i+1}$ one has $d_X(p_i,p_{i+1})\le M$, so the pair lies in $\mathrm{Pzt}_k$, and its two photon sites and its six paths lie in $\tilde{X}$.
The photon sites and the six paths of its class in $\mathrm{Pz}_k$ are the images under $\mathrm{pr}$ of those of the two points of the pair in $\mathbb{L}^{\mathrm{f}}_k$ (Sections~I.1.7 and~I.1.8, and (Lf4)); they lie in $\mathrm{pr}(\tilde{X})=X$.
Hence the pair is admissible for $X$ of kind (I) (Section~I.1.7).
\item[(c2)] For $v\in\mathbb{R}^3$ one has $|v|_1\le3|v|$.
By (c1),
\begin{gather*}
|\pi(p_0)-\pi(p_1)|_1\le3|c-z_0|<\frac{3M}{2},\qquad|\pi(p_i)-\pi(p_{i+1})|_1=M\quad(1\le i\le n),\\
|\pi(p_{n+1})-\pi(p_{n+2})|_1\le3|z_n-\tilde{c}|<\frac{3M}{2},
\end{gather*}
a term with $\pi(p_i)=\pi(p_{i+1})$ being $0$.
So the sum $\sum_{i=0}^{n+1}|\pi(p_i)-\pi(p_{i+1})|_1$ is at most $\frac{3M}{2}+n\cdot M+\frac{3M}{2}=(n+3)\cdot M$.
From $n+1\le c_X\cdot(L+1)$,
\[
n+3\le c_X\cdot(L+1)+2=c_X\cdot L+c_X+2\le3c_X\cdot L,
\]
because $c_X+2=56\le540\le2c_X\cdot L$ ($c_X=54$, $L\ge5$).
Let $e_k\cdot|\operatorname{Im}B|_\infty\le\frac{1}{2}\,e_k^{\varepsilon}$.
By (u1), for $0\le i\le n+2$,
\begin{align*}
|U_i|,\ |U_i^-|&\le\prod_{j<i}\exp\left(e_k\cdot|\operatorname{Im}B|_\infty\cdot|\pi(p_j)-\pi(p_{j+1})|_1\right)\\
&\le\exp\left(\frac{1}{2}\,e_k^{\varepsilon}\cdot(n+3)\cdot M\right)\le\exp\left(\frac{3}{2}\,c_X\cdot L\cdot M\cdot e_k^{\varepsilon}\right).
\end{align*}
Condition (A9), which stands in \cite[Section~5.6.2]{QEDI} (Section~I.8.4), contains the inequality $90c_X\cdot L\cdot M\cdot e_0^{\varepsilon}\le1$.
Hence $3c_X\cdot L\cdot M\cdot e_k^{\varepsilon}\le3c_X\cdot L\cdot M\cdot e_0^{\varepsilon}\le\frac{1}{30}$, and the exponent is at most $\frac{1}{60}\le1$.
\item[(c3)] Let $p_i\ne p_{i+1}$, and write $d_X$ for the number $d_X(p_i,p_{i+1})>0$.
If $d_X\le1$, then $d_X^{\alpha}\le1$; if $d_X\ge1$, then $d_X^{\alpha}\le d_X$, because $0<\alpha<1$ (C1).
So $d_X^{\alpha}\le\max(1,d_X)$.
By (c1), $d_X\le M$, and $M\ge1$; hence $d_X(p_i,p_{i+1})^{\alpha}\le M$.
The sum $\sum_{i=0}^{n+1}d_X(p_i,p_{i+1})^{\alpha}$ has at most $n+2$ terms, the terms with $p_i=p_{i+1}$ being omitted.
So it is at most $(n+2)\cdot M\le(n+3)\cdot M\le3c_X\cdot L\cdot M$.
\end{itemize}
\end{proof}

\begin{proof}[Proof of Proposition~17.1\,(d)]
\emph{The display in components.}
For $\omega\in\{0,1\}$ put $u^{(0)}:=u_B$, $u^{(1)}:=u_B^{-1}$, and $U^{(\omega)}_i:=\prod_{j<i}u^{(\omega)}(p_j,p_{j+1})$; thus $U^{(0)}_i=U_i$ and $U^{(1)}_i=U_i^-$.
The mean $\operatorname{Mn}$ is $\left(\#|\mathrm{Ch}(x,\tilde{x})|\right)^{-1}$ times the sum over $\mathrm{Ch}(x,\tilde{x})$.
By (0.d), the definitions (3.2) and (3.3) take the following form, for $x=(c,g)$ over $\tilde{X}$, $\beta\in\{1,\dots,4\}$ and $\omega\in\{0,1\}$, with $\tilde{x}=(\tilde{c},g')$:
\begin{equation}
\begin{split}
\boldsymbol{\Theta}\,\psi(x,\beta,\omega):={}&\mathrm{cnt}(X)^{-1}\sum_{\tilde{c}\in\tilde{X}}\ \sum_{g'\in O_{\mathrm{cub}}}\left(\#|\mathrm{Ch}(x,\tilde{x})|\right)^{-1}\sum_{(\square_0,\dots,\square_n)\in\mathrm{Ch}(x,\tilde{x})}\\
&\Bigl\{U^{(\omega)}_{n+2}\cdot\psi(\tilde{x},\beta,\omega)-\sum_{0\le i\le n+1:\ p_i\ne p_{i+1}}U^{(\omega)}_i\cdot d_X(p_i,p_{i+1})^{\alpha}\cdot\chi(p_i,p_{i+1},\beta,\omega)\Bigr\}.
\end{split}
\tag{1.a}
\end{equation}
For $\omega=0$ this is the first formula of (3.3), for $\omega=1$ the second.
Each site $\tilde{x}$ carries the weight $\mathrm{cnt}(X)^{-1}=\operatorname{Vol}(X)^{-1}\cdot\frac{\eta_k^3}{48}$: the sum over $g'$ with the factor $\frac{1}{48}$ is part of $\int_X\mathrm{d}\tilde{x}$.
For $B=\mathscr{T}$, (1.a) is (I.5.3).
By (c1), every pair $(p_i,p_{i+1})$ with $p_i\ne p_{i+1}$ lies in $\mathrm{Pzt}_k$; hence $\chi(p_i,p_{i+1},\beta,\omega)$ is a generator (Section~0.4), and the right side of (1.a) belongs to $\mathrm{Lin}$.
It decomposes as
\begin{equation}
\boldsymbol{\Theta}\,\psi(x,\beta,\omega)=\mathrm{Lpsi}(x,\beta,\omega)+\mathrm{Lchi}(x,\beta,\omega)
\tag{1.b}
\end{equation}
into its part $\mathrm{Lpsi}$, a linear combination of the generators $\psi(\tilde{x},\beta,\omega)$, $\tilde{x}$ over $\tilde{X}$ (bidegree $(1,0)$), and its part $\mathrm{Lchi}$, a linear combination of the generators $\chi(p_i,p_{i+1},\beta,\omega)$ at the chain pairs (bidegree $(0,1)$).
Both parts carry the same indices $\beta$ and $\omega$ as $\psi(x,\beta,\omega)$ itself.

\emph{Existence and uniqueness.}
By (Ge3), $\mathrm{Grp}[X]$ is generated by $1$ and the independent generators $\psi(x,\beta,\omega)$, $x$ over $\tilde{X}$.
The right sides of (1.a) lie in $\mathrm{Lin}$, and by (Ge2) they anticommute.
Lemma~A, applied with $\mathrm{Alg}=\mathrm{Gr}[\psi,\chi]$, yields exactly one homomorphism $\boldsymbol{\Theta}=\boldsymbol{\Theta}_{X,B}$ of $\mathrm{Grp}[X]$ into $\mathrm{Gr}[\psi,\chi]$ with (1.a).

\emph{Independence of the lift.}
By (Lf2), any other lift is $\tilde{X}'=\tilde{X}+\ell\cdot w$ with $w\in\mathbb{Z}^3$.
Let $\boldsymbol{\Theta}'$ be the homomorphism formed with $\tilde{X}'$.
For $x$ over $\tilde{X}$ the point $x+\ell\cdot w$ lies over $\tilde{X}'$, and $\tilde{x}\mapsto\tilde{x}+\ell\cdot w$ is a bijection of $\tilde{X}\times O_{\mathrm{cub}}$ onto $\tilde{X}'\times O_{\mathrm{cub}}$.
The translation by $\ell\cdot w$ (recall that $\ell$ is a multiple of $M$) maps the $M$-cubes of $\tilde{X}$ onto those of $\tilde{X}'$, centres onto centres, and common faces onto common faces; hence it maps $\mathrm{Ch}(x,\tilde{x})$ onto the set $\mathrm{Ch}(x+\ell\cdot w,\tilde{x}+\ell\cdot w)$ formed in $\tilde{X}'$, and the chain points of the image of a chain are the points $p_i+\ell\cdot w$.
As $B$ is read through $\mathrm{pr}$, one has $u_B(p_i+\ell\cdot w,p_{i+1}+\ell\cdot w)=u_B(p_i,p_{i+1})$; by (I.1.5), $d_X(p_i+\ell\cdot w,p_{i+1}+\ell\cdot w)=d_X(p_i,p_{i+1})$; and neither $\mathrm{cnt}(X)$ nor $\operatorname{Vol}(X)$ involves the lift.
By (0.c),
\begin{gather*}
\psi(\tilde{x}+\ell\cdot w,\beta,\omega)=\boldsymbol{\chi}(w)\cdot\psi(\tilde{x},\beta,\omega),\\
\chi(p_i+\ell\cdot w,p_{i+1}+\ell\cdot w,\beta,\omega)=\boldsymbol{\chi}(w)\cdot\chi(p_i,p_{i+1},\beta,\omega).
\end{gather*}
Inserting these identities into (1.a) for $\tilde{X}'$ at the point $x+\ell\cdot w$, one obtains
\[
\boldsymbol{\Theta}'\psi(x+\ell\cdot w,\beta,\omega)=\boldsymbol{\chi}(w)\cdot\boldsymbol{\Theta}\,\psi(x,\beta,\omega)=\boldsymbol{\Theta}\left(\boldsymbol{\chi}(w)\cdot\psi(x,\beta,\omega)\right)=\boldsymbol{\Theta}\,\psi(x+\ell\cdot w,\beta,\omega).
\]
The algebra $\mathrm{Grp}[X]$ is generated by $1$ and the elements $\psi(x+\ell\cdot w,\beta,\omega)$, $x$ over $\tilde{X}$; by (Un), $\boldsymbol{\Theta}'=\boldsymbol{\Theta}$.

\emph{Bars.}
By (1.a), $\boldsymbol{\Theta}\,\psi(x,\beta,0)$ is a linear combination of generators with $\omega=0$, and $\boldsymbol{\Theta}\,\psi(x,\beta,1)$ a linear combination of generators with $\omega=1$.
\end{proof}

In the proof of (e), the notation $U^{(\omega)}_i[B]$ is used where the dependence on $B$ matters.
Each of the four relations (3.4) is an equality of two maps that are additive and multiplicative, send $1$ to $1$, and are both linear (in the fourth relation, both antilinear).
By (Un), it suffices to prove each relation on generators of the subalgebra on which its two sides are defined.

\begin{proof}[Proof of Proposition~17.1\,(e), first relation]
Fix a representative $a\in M\mathbb{Z}^3$ of the translation part of $\tilde{g}=(S,a)$ (Section~I.4.3), and put $r:=r(S)$, $S^{(0)}:=S$, $S^{(1)}:=(S^{-1})^T$.
Then $\tilde{g}$ acts on $\mathbb{L}_k$, on $\mathbb{L}^{\mathrm{f}}_k$ and on $\mathrm{Pzt}_k$, and in components (I.T-g) reads
\begin{equation}
\begin{aligned}
a_{\tilde{g}}\psi(\tilde{g}x,\beta,\omega)&=\sum_{\beta'}S^{(\omega)}_{\beta\beta'}\cdot\psi(x,\beta',\omega),\\
a_{\tilde{g}}\chi(\tilde{g}x,\tilde{g}x',\beta,\omega)&=\sum_{\beta'}S^{(\omega)}_{\beta\beta'}\cdot\chi(x,x',\beta',\omega).
\end{aligned}
\tag{1.c}
\end{equation}
By Proposition~16.1\,(a1), $\tilde{g}X$ is a small polymer and $\tilde{g}\tilde{X}$ is a lift of $\tilde{g}X$.
So $\boldsymbol{\Theta}_{\tilde{g}X,\tilde{g}B}$ is defined on $\mathrm{Grp}[\tilde{g}X]$, which is generated by $1$ and the $\psi(\tilde{g}x,\beta,\omega)$, $x$ over $\tilde{X}$, and by clause (d) it may be formed with the lift $\tilde{g}\tilde{X}$.
By (1.c), $a_{\tilde{g}}$ maps $\mathrm{Grp}[\tilde{g}X]$ into $\mathrm{Grp}[X]$; so both sides of the relation are defined on $\mathrm{Grp}[\tilde{g}X]$.
Consider (1.a) for $(\tilde{g}\tilde{X},\tilde{g}B)$ at the site $\tilde{g}x=(r\cdot c+a,\,r\cdot g)$.
\begin{itemize}
\item The map $\tilde{x}=(\tilde{c},g')\mapsto\tilde{g}\tilde{x}=(r\cdot\tilde{c}+a,\,r\cdot g')$ is a bijection of $\tilde{X}\times O_{\mathrm{cub}}$ onto $(\tilde{g}\tilde{X})\times O_{\mathrm{cub}}$ (Proposition~16.1\,(a2)); so the summation may be indexed by the points $\tilde{x}$ over $\tilde{X}$.
Moreover $\mathrm{cnt}(\tilde{g}X)=\mathrm{cnt}(X)$, because $|\tilde{g}X|_M=|X|_M$ (Proposition~16.1\,(a1)).
\item The map $(\square_i)\mapsto(\tilde{g}\square_i)$ is a bijection of $\mathrm{Ch}(x,\tilde{x})$ onto the set $\mathrm{Ch}(\tilde{g}x,\tilde{g}\tilde{x})$ formed in $\tilde{g}\tilde{X}$ (Proposition~16.1\,(a1): $M$-cubes go to $M$-cubes, cubes with a common face to cubes with a common face, shortest chains to shortest chains, and the same holds for $\tilde{g}^{-1}$).
\item The chain points of $(\tilde{g}X,\tilde{g}x,\tilde{g}\tilde{x},(\tilde{g}\square_i))$ are the points $\tilde{g}p_i$: indeed $\tilde{g}(z_i,g)=(\text{the centre of }\tilde{g}\square_i,\,r\cdot g)$, and $r\cdot g$ is the label of $\tilde{g}x$ (Proposition~16.1\,(a2)).
Moreover $\tilde{g}p_i\ne\tilde{g}p_{i+1}$ if and only if $p_i\ne p_{i+1}$.
\item By (u3), $u_{\tilde{g}B}(\tilde{g}p_i,\tilde{g}p_{i+1})=u_B(p_i,p_{i+1})$; so $U^{(\omega)}_i[\tilde{g}B]$, formed with the points $\tilde{g}p_j$, equals $U^{(\omega)}_i[B]$.
Also $d_X(\tilde{g}p_i,\tilde{g}p_{i+1})=d_X(p_i,p_{i+1})$ (Proposition~16.1\,(a2)).
\end{itemize}
Consequently
\begin{equation*}
\begin{split}
\boldsymbol{\Theta}_{\tilde{g}X,\tilde{g}B}\,\psi(\tilde{g}x,\beta,\omega)={}&\mathrm{cnt}(X)^{-1}\sum_{\tilde{x}\text{ over }\tilde{X}}\left(\#|\mathrm{Ch}(x,\tilde{x})|\right)^{-1}\sum_{(\square_0,\dots,\square_n)\in\mathrm{Ch}(x,\tilde{x})}\\
&\Bigl\{U^{(\omega)}_{n+2}[B]\cdot\psi(\tilde{g}\tilde{x},\beta,\omega)-\sum_{i:\ p_i\ne p_{i+1}}U^{(\omega)}_i[B]\cdot d_X(p_i,p_{i+1})^{\alpha}\cdot\chi(\tilde{g}p_i,\tilde{g}p_{i+1},\beta,\omega)\Bigr\}.
\end{split}
\end{equation*}
Apply $a_{\tilde{g}}$ and use (1.c); the coefficients do not depend on $\beta$, and so
\begin{align*}
a_{\tilde{g}}\boldsymbol{\Theta}_{\tilde{g}X,\tilde{g}B}\,\psi(\tilde{g}x,\beta,\omega)
&=\sum_{\beta'}S^{(\omega)}_{\beta\beta'}\cdot\boldsymbol{\Theta}_{X,B}\,\psi(x,\beta',\omega)
=\boldsymbol{\Theta}_{X,B}\Bigl(\sum_{\beta'}S^{(\omega)}_{\beta\beta'}\cdot\psi(x,\beta',\omega)\Bigr)\\
&=\boldsymbol{\Theta}_{X,B}\,a_{\tilde{g}}\psi(\tilde{g}x,\beta,\omega).
\end{align*}

\emph{The representative.}
If the representative $a$ is replaced by $a+\ell\cdot w$, $w\in\mathbb{Z}^3$, then in (I.T-g) the point $\tilde{g}^{-1}x$ is replaced by $\tilde{g}^{-1}x-\ell\cdot rw$.
By (0.c), the image of every generator is thereby multiplied by $\boldsymbol{\chi}(-rw)$, and $\boldsymbol{\chi}(-rw)=\boldsymbol{\chi}(w)$, since $r$ is diagonal with entries $\pm1$ (Section~I.4.1).
So the automorphism belonging to $a+\ell\cdot w$ is $a_{\tilde{g}}\mathrm{Pty}^{w_1+w_2+w_3}$.
By (1.a), $\boldsymbol{\Theta}\,\mathrm{Pty}=\mathrm{Pty}\,\boldsymbol{\Theta}$ on the generators, for every small polymer and every bond function.
Hence the relation holds for every representative.
\end{proof}

\paragraph{Proof of Proposition~17.1\,(e), second relation.}
By (I.T-G), $G_\Lambda$ maps $\mathrm{Grp}[X]$ onto itself; so both sides are defined on $\mathrm{Grp}[X]$.
Let $B^1=B-d\Lambda$ on the bonds with both ends in $X$.
By (c1), the six paths of every chain pair lie in $\tilde{X}$; their bonds have both ends in $\tilde{X}$, and the images of these bonds have both ends in $X$.
By (u3), with $(-1)^{\omega}=1$ for $\omega=0$ and $(-1)^{\omega}=-1$ for $\omega=1$,
\begin{multline*}
u^{(\omega)}[B^1](p_j,p_{j+1})\\
=u^{(\omega)}[B](p_j,p_{j+1})\cdot\exp\left\{-i\cdot(-1)^{\omega}\cdot e_k\cdot\left(\Lambda(\pi(p_{j+1}))-\Lambda(\pi(p_j))\right)\right\},
\end{multline*}
and this holds also when $p_j=p_{j+1}$.
In the product over $j<i$ the exponents telescope:

\[
U^{(\omega)}_i[B^1]=U^{(\omega)}_i[B]\cdot\exp\left\{-i\cdot(-1)^{\omega}\cdot e_k\cdot\left(\Lambda(\pi(p_i))-\Lambda(\pi(p_0))\right)\right\}.
\]
In components the automorphism $G_\Lambda$ of (I.T-G) in Section~I.4.5 acts on the generators that occur in (1.a) by
\begin{align*}
G_\Lambda\,\psi(\tilde x,\beta,\omega)&=\exp\left\{i\cdot(-1)^{\omega}\cdot e_k\cdot\Lambda(\pi(p_{n+2}))\right\}\cdot\psi(\tilde x,\beta,\omega)\qquad(\tilde x=p_{n+2}),\\
G_\Lambda\,\chi(p_i,p_{i+1},\beta,\omega)&=\exp\left\{i\cdot(-1)^{\omega}\cdot e_k\cdot\Lambda(\pi(p_i))\right\}\cdot\chi(p_i,p_{i+1},\beta,\omega),
\end{align*}
the phase of a generator $\chi$ being the one of its first point.
Consequently, when $G_\Lambda$ is applied to (1.a) written for $B^1$, every term turns into the corresponding term of (1.a) for $B$, multiplied by $\exp\left\{i\cdot(-1)^{\omega}\cdot e_k\cdot\Lambda(\pi(p_0))\right\}$, where $\pi(p_0)=c$.
Hence
\[
G_\Lambda\,\boldsymbol{\Theta}_{X,B^1}\,\psi(x,\beta,\omega)=\exp\left\{i\cdot(-1)^{\omega}\cdot e_k\cdot\Lambda(c)\right\}\cdot\boldsymbol{\Theta}_{X,B}\,\psi(x,\beta,\omega)=\boldsymbol{\Theta}_{X,B}\,G_\Lambda\,\psi(x,\beta,\omega).
\]

\emph{Third relation.}
In components, (I.T-C) reads
\begin{align*}
\mathrm{Cg}\,\psi(x,\beta,\omega)&=(-1)^{\omega}\sum_{\beta'}\mathsf{C}_{\beta\beta'}\cdot\psi(x,\beta',1-\omega),\\
\mathrm{Cg}\,\chi(x,x',\beta,\omega)&=(-1)^{\omega}\sum_{\beta'}\mathsf{C}_{\beta\beta'}\cdot\chi(x,x',\beta',1-\omega).
\end{align*}
Therefore $\mathrm{Cg}$ maps $\mathrm{Grp}[X]$ onto itself, and both sides of the relation are defined on $\mathrm{Grp}[X]$.
By (u3) one has $u_{-B}=u_B^{-1}$, whence $U^{(\omega)}_i[-B]=U^{(1-\omega)}_i[B]$.
The chains, the chain points, the distances $d_X$ and the number $\mathrm{cnt}(X)$ involve neither $B$ nor $\omega$.
Applying $\mathrm{Cg}$ to (1.a) written for $-B$, one obtains
\begin{align*}
\mathrm{Cg}\,\boldsymbol{\Theta}_{X,-B}\,\psi(x,\beta,\omega)&=(-1)^{\omega}\sum_{\beta'}\mathsf{C}_{\beta\beta'}\cdot\mathrm{cnt}(X)^{-1}\sum_{\tilde x\text{ over }\tilde X}\#|\mathrm{Ch}(x,\tilde x)|^{-1}\sum_{\mathrm{Ch}(x,\tilde x)}\Bigl\{U^{(1-\omega)}_{n+2}[B]\cdot\psi(\tilde x,\beta',1-\omega)\\
&\qquad-\sum_{i:\,p_i\ne p_{i+1}}U^{(1-\omega)}_i[B]\cdot d_X(p_i,p_{i+1})^{\alpha}\cdot\chi(p_i,p_{i+1},\beta',1-\omega)\Bigr\}\\
&=(-1)^{\omega}\sum_{\beta'}\mathsf{C}_{\beta\beta'}\cdot\boldsymbol{\Theta}_{X,B}\,\psi(x,\beta',1-\omega)=\boldsymbol{\Theta}_{X,B}\,\mathrm{Cg}\,\psi(x,\beta,\omega),
\end{align*}
where $\sum_{\mathrm{Ch}(x,\tilde x)}$ denotes the sum over the chains $(\square_0,\dots,\square_n)\in\mathrm{Ch}(x,\tilde x)$.

\emph{Fourth relation.}
By (I.T-K), the map $\mathbf{K}$ is additive and multiplicative; it sends the product of a complex number and an element to the product of the conjugate number and the image of that element; and
\[
\mathbf{K}\,\psi(x,\beta,\omega)=\sum_{\beta'}\varkappa_{\beta\beta'}\cdot\psi(x,\beta',\omega),\qquad\mathbf{K}\,\chi(x,x',\beta,\omega)=\sum_{\beta'}\varkappa_{\beta\beta'}\cdot\chi(x,x',\beta',\omega).
\]
So $\mathbf{K}$ maps $\mathrm{Grp}[X]$ onto itself; both sides of the relation are defined on $\mathrm{Grp}[X]$, and both are antilinear.
From (u3) one gets $\overline{U^{(\omega)}_i[-\overline{B}]}=U^{(\omega)}_i[B]$, while the numbers $\mathrm{cnt}(X)^{-1}$, $\#|\mathrm{Ch}(x,\tilde x)|^{-1}$ and $d_X(p_i,p_{i+1})^{\alpha}$ are real.
Applying $\mathbf{K}$ to (1.a) written for $-\overline{B}$ therefore gives
\[
\mathbf{K}\,\boldsymbol{\Theta}_{X,-\overline{B}}\,\psi(x,\beta,\omega)=\sum_{\beta'}\varkappa_{\beta\beta'}\cdot\boldsymbol{\Theta}_{X,B}\,\psi(x,\beta',\omega)=\boldsymbol{\Theta}_{X,B}\,\mathbf{K}\,\psi(x,\beta,\omega).
\]

\paragraph{(f)}
Written out in components,
\[
\int_X\mathrm{d}x\int_X\mathrm{d}x'\;\bar\psi(x)^T\cdot K(x,x')\cdot\psi(x')=\int_X\mathrm{d}x\int_X\mathrm{d}x'\,\sum_{\beta,\beta'}K(x,x')_{\beta\beta'}\cdot\psi(x,\beta,1)\cdot\psi(x',\beta',0),
\]
and this is an element of $\mathrm{Grp}[X]$: a finite sum over the sites over $\tilde X$, each of which carries the weight $\frac{\eta_k^3}{48}$.
Since $\boldsymbol{\Theta}$ is linear and multiplicative, the image of this element under $\boldsymbol{\Theta}$ is
\[
\int_X\mathrm{d}x\int_X\mathrm{d}x'\,\sum_{\beta,\beta'}K(x,x')_{\beta\beta'}\cdot\boldsymbol{\Theta}\,\psi(x,\beta,1)\cdot\boldsymbol{\Theta}\,\psi(x',\beta',0).
\]
By (3.3),
\begin{align*}
\boldsymbol{\Theta}\,\psi(x,\beta,1)&=\operatorname{Vol}(X)^{-1}\cdot\int_X\mathrm{d}\tilde x'\,\left(\mathrm{tm}^{-}(x,\tilde x')\cdot\bar\psi_\beta(\tilde x')-\mathrm{rm}^{-}(x,\tilde x')[\bar\chi]_\beta\right),\\
\boldsymbol{\Theta}\,\psi(x',\beta',0)&=\operatorname{Vol}(X)^{-1}\cdot\int_X\mathrm{d}\tilde x\,\left(\mathrm{tm}(x',\tilde x)\cdot\psi_{\beta'}(\tilde x)-\mathrm{rm}(x',\tilde x)[\chi]_{\beta'}\right).
\end{align*}
We insert these two expressions, keep the two factors in their order, and expand the product.
The quantities $\mathrm{tm}$ and $\mathrm{tm}^{-}$ are numbers.
The four products so obtained are the four terms of Proposition~17.1\,(f): the product $(\mathrm{tm}^{-}\bar\psi)\cdot(\mathrm{tm}\,\psi)$, of bidegree $(2,0)$; the two products $-(\mathrm{tm}^{-}\bar\psi)\cdot(\mathrm{rm}[\chi])$ and $-(\mathrm{rm}^{-}[\bar\chi])\cdot(\mathrm{tm}\,\psi)$, of bidegree $(1,1)$; and the product $+(\mathrm{rm}^{-}[\bar\chi])\cdot(\mathrm{rm}[\chi])$, of bidegree $(0,2)$.
In the first of them the integrations over $x$ and $x'$ act only on the numbers $\mathrm{tm}^{-}(x,\tilde x')\cdot K(x,x')\cdot\mathrm{tm}(x',\tilde x)$.

Let now $B=0$.
Then every transport equals $1$, so that $U_i=U_i^{-}=1$, and $\mathrm{tm}=\mathrm{tm}^{-}=1$, each being a mean of ones.
The first term becomes
\[
\operatorname{Vol}(X)^{-2}\int_X\mathrm{d}\tilde x'\int_X\mathrm{d}\tilde x\;\bar\psi(\tilde x')^T\cdot\left[\int_X\mathrm{d}x\int_X\mathrm{d}x'\,K(x,x')\right]\cdot\psi(\tilde x)=\bar\psi_{\mathrm{av}}^T\cdot\left[\int_X\mathrm{d}x\int_X\mathrm{d}x'\,K(x,x')\right]\cdot\psi_{\mathrm{av}}.
\]
By (0.d),
\[
\psi_{\mathrm{av}}=\mathrm{cnt}(X)^{-1}\sum_{c\in\tilde X}\sum_{g\in O_{\mathrm{cub}}}\psi(c,g),
\]
which is the mean over the photon sites of $\tilde X$ and over the $48$ labels.
\qed

\begin{named}{Remark}[The level $k=0$]
For $k=0$ one has $\eta_0=1$ and $\mathbb{L}_0=\mathbb{Z}^3$, and $e_k$ is a number in $(0,e_0]$, namely the number $q$ of the flow.
An $M$-cube then contains $M^3$ photon sites and $48M^3$ fermion sites; $\mathrm{cnt}(X)=48M^3\cdot|X|_M$; a site has the weight $\frac{1}{48}$ in $\int\mathrm{d}x$; and the right side of (1.a) begins with $\left(48M^3\cdot|X|_M\right)^{-1}\sum_{\tilde x\text{ over }\tilde X}$, that is, with the sum over the photon sites $\tilde c\in\tilde X$ and, for each of them, over the label of $\tilde x$ in $O_{\mathrm{cub}}$.
The distances $d_X(p_i,p_{i+1})$ are integers between $1$ and $M$: for $i=0$ one has $|c-z_0|\in\left\{1,\dots,\frac{M-1}{2}\right\}$; for $1\le i\le n$ the distance equals $M$; for $i=n+1$ either $|\tilde c-z_n|\in\left\{1,\dots,\frac{M-1}{2}\right\}$, or the pair lies inside the cluster of $z_n$ and $d_X=\eta_0=1$.
In (c3) the case $d_X<1$ therefore does not arise, and a last link inside a cluster enters (1.a) with $d_X^{\alpha}=1$ and $u_B=1$; at such a link the relation of (u4) is $\chi(p,p')=\psi(p')-\psi(p)$, as in (I.L0.14).
Since $\eta_0=1$, the transports are $u_B(p,p')=\exp\left\{i\cdot e_k\cdot(\tau B)(c,c')\right\}$.
\end{named}

\begin{named}{Remark}[Dependence of the constants]
The constants $\frac12$, $\frac14$, $\mathrm{e}$ and $3c_X\cdot L\cdot M$ in the clauses of Proposition~17.1 are built from $c_X$ (C0), $L$ (C3) and $M$ (C5) alone.
They are the same for every level $k\ge0$, every side $\ell\in\mathrm{Cl}$ and every $e_k\in(0,e_0]$, and they do not depend on $B$, $X$, $\tilde g$ or $\Lambda$.
The side enters the argument only through the facts that $\ell$ is a multiple of $M$ and that $\ell\ge13L\cdot M$, in Section~0.3; and $e_k$ enters only through $0<e_k\le e_0$, which gives $e_k^{\varepsilon}\le e_0^{\varepsilon}$ in (a) and (c2).
\end{named}


\subsection*{2. Proposition 17.2}

\subsubsection*{Proof of Proposition 17.2}

\begin{proof}
We keep the notation of Section~0 and of the proof of Proposition~17.1, in particular the displays (1.a) and (1.b).
Let $\tilde X$ be a lift of the small polymer $X$ as in (Lf1), and put $\boldsymbol{\Theta}:=\boldsymbol{\Theta}_{X,\mathscr{T}}$; by Proposition~17.1\,(d) this homomorphism does not depend on the choice of the lift.

\paragraph{(0) Preparation.}
\begin{itemize}
\item[(Pp1)] Proposition~17.1\,(a) gives $e_k\cdot|\operatorname{Im}\mathscr{T}|_{\infty}<\frac12e_k^{\varepsilon}$.
Hence the bound of Proposition~17.1\,(c2) is available for $B=\mathscr{T}$:
\[
\left|U^{(\omega)}_i\right|\le\mathrm{e}\qquad\text{for }0\le i\le n+2\text{ and }\omega\in\{0,1\}.
\]
\item[(Pp2)] The kernels $E_{nm}(X;\mathscr{A},\mathscr{C};\cdot)$ are the antisymmetric kernels of the element $E(X;\mathscr{A},\mathscr{C})$ of $\mathrm{Gr}[\psi,\chi]$, and they are determined by this element (Proposition~15\,(a)).
They are antiperiodic in each argument, by (I.AP) in Section~I.3.6.
Together with (0.c) this shows that the product of a kernel with the corresponding generators is $\ell\mathbb{Z}^3$-periodic in each argument, so that each argument may be summed over an arbitrary system of representatives.
By the clause (k2) of Section~I.5.4, only the $\xi$ over $X$ and the $\boldsymbol{\zeta}$ admissible for $X$ contribute.
We write $\int_X\mathrm{d}\xi$ for the sum over the indices $\xi=(x,\beta,\omega)$ with $x$ over $\tilde X$, each taken with the weight $\frac{\eta_k^3}{48}$.
\item[(Pp3)] By (Pp2) and (Ge3) the element
\[
E^{[2,0]}(X;\mathscr{A},\mathscr{C})=\frac12\int_X\mathrm{d}\xi_1\int_X\mathrm{d}\xi_2\;E_{20}(X;\mathscr{A},\mathscr{C};\xi_1,\xi_2)\cdot\psi(\xi_1)\cdot\psi(\xi_2)
\]
belongs to $\mathrm{Grp}[X]$.
Therefore $\boldsymbol{\Theta}\left(E^{[2,0]}(X;\mathscr{A},\mathscr{C})\right)$ is defined, and, $\boldsymbol{\Theta}$ being linear and multiplicative,
\begin{equation}\tag{2.a}
\begin{aligned}
E^{\mathrm{nat}}(X;\mathscr{A},\mathscr{C})&=F_1+F_2,\\
F_1&:=E(X;\mathscr{A},\mathscr{C})-E^{[2,0]}(X;\mathscr{A},\mathscr{C})=\sum_{(n,m)\ne(2,0)}E^{[n,m]}(X;\mathscr{A},\mathscr{C}),\\
F_2&:=\boldsymbol{\Theta}\left(E^{[2,0]}(X;\mathscr{A},\mathscr{C})\right)\\
&=\frac12\int_X\mathrm{d}\xi_1\int_X\mathrm{d}\xi_2\;E_{20}(X;\mathscr{A},\mathscr{C};\xi_1,\xi_2)\cdot\boldsymbol{\Theta}\,\psi(\xi_1)\cdot\boldsymbol{\Theta}\,\psi(\xi_2).
\end{aligned}
\end{equation}
\item[(Pp4)] By (1.b),
\begin{align*}
\boldsymbol{\Theta}\,\psi(\xi_1)\cdot\boldsymbol{\Theta}\,\psi(\xi_2)={}&\mathrm{Lpsi}(\xi_1)\cdot\mathrm{Lpsi}(\xi_2)\\
&+\left[\mathrm{Lpsi}(\xi_1)\cdot\mathrm{Lchi}(\xi_2)+\mathrm{Lchi}(\xi_1)\cdot\mathrm{Lpsi}(\xi_2)\right]\\
&+\mathrm{Lchi}(\xi_1)\cdot\mathrm{Lchi}(\xi_2),
\end{align*}
the three groups being of the bidegrees $(2,0)$, $(1,1)$ and $(0,2)$.
Accordingly $F_2=F_2^{[2,0]}+F_2^{[1,1]}+F_2^{[0,2]}$.
\end{itemize}

\paragraph{(a)}
Let $x=(c,g)$ and $\tilde x$ lie over $\tilde X$, let $(\square_0,\dots,\square_n)\in\mathrm{Ch}(x,\tilde x)$, $\beta\in\{1,\dots,4\}$ and $\omega\in\{0,1\}$.
We claim that for $0\le i\le n+2$
\begin{equation}\tag{2.b}
\mathrm{ev}\,\psi(p_0,\beta,\omega)=U^{(\omega)}_i\cdot\mathrm{ev}\,\psi(p_i,\beta,\omega)-\sum_{j<i:\,p_j\ne p_{j+1}}U^{(\omega)}_j\cdot d_X(p_j,p_{j+1})^{\alpha}\cdot\mathrm{ev}\,\chi(p_j,p_{j+1},\beta,\omega).
\end{equation}
For $i=0$ this holds because $U^{(\omega)}_0=1$, an empty product, and the sum is empty.
Assume (2.b) for some $i\le n+1$.
If $p_i=p_{i+1}$, then $u^{(\omega)}(p_i,p_{i+1})=1$ and $U^{(\omega)}_{i+1}=U^{(\omega)}_i$, and no term is added to the sum.
If $p_i\ne p_{i+1}$, the pair $(p_i,p_{i+1})$ is of kind (I) for $X$ by Proposition~17.1\,(c1); for $i=n+1$ it may lie inside the cluster of $z_n$.
The images under $\mathrm{ev}$ therefore satisfy the relations of Proposition~17.1\,(u4) at this pair, and (3.1) yields
\[
\mathrm{ev}\,\psi(p_i,\beta,\omega)=u^{(\omega)}(p_i,p_{i+1})\cdot\mathrm{ev}\,\psi(p_{i+1},\beta,\omega)-d_X(p_i,p_{i+1})^{\alpha}\cdot\mathrm{ev}\,\chi(p_i,p_{i+1},\beta,\omega).
\]
Substituting this into (2.b) and using $U^{(\omega)}_i\cdot u^{(\omega)}(p_i,p_{i+1})=U^{(\omega)}_{i+1}$, one obtains (2.b) for $i+1$.
(By (0.c), the relations of (u4) at a pair of $\mathrm{Pzt}_k$ and at its translate by $\ell w$ differ by the factor $\boldsymbol{\chi}(w)$; so it does not matter which representative of a pair of $\mathrm{Pz}_k$ the hypothesis refers to.)

Now take $i=n+2$ in (2.b), so that $p_0=x$ and $p_{n+2}=\tilde x$, form the mean over $\mathrm{Ch}(x,\tilde x)$, and then the mean (0.d) over $\tilde x$.
The left side depends neither on the chain nor on $\tilde x$, and the weights of each of the two means add up to $1$.
Since $\mathrm{ev}$ is linear, the right side becomes $\mathrm{ev}$ applied to the right side of (1.a).
Thus
\[
\mathrm{ev}\,\psi(x,\beta,\omega)=\mathrm{ev}\left(\boldsymbol{\Theta}\,\psi(x,\beta,\omega)\right)\qquad\text{for all }x\text{ over }\tilde X\text{ and all }\beta,\omega.
\]
The maps $\mathrm{ev}\circ\boldsymbol{\Theta}$ and the restriction of $\mathrm{ev}$ to $\mathrm{Grp}[X]$ are linear and multiplicative, and they coincide on the generators of $\mathrm{Grp}[X]$.
By (Un) they coincide on $E^{[2,0]}(X;\mathscr{A},\mathscr{C})$, which is a linear combination of products of two generators.
With (2.a) it follows that
\[
\mathrm{ev}\left(E^{\mathrm{nat}}(X;\mathscr{A},\mathscr{C})\right)=\mathrm{ev}(F_1)+\mathrm{ev}\left(E^{[2,0]}(X;\mathscr{A},\mathscr{C})\right)=\mathrm{ev}\left(E(X;\mathscr{A},\mathscr{C})\right).
\]

\emph{The insertion.}
Let $\mathrm{sbs}_{\mathscr{T}}$ be the homomorphism of $\mathrm{Gr}[\psi,\chi]$ into its subalgebra generated by the $\psi(\xi)$ that is given by
\begin{align*}
\psi(\xi)&\mapsto\psi(\xi),\\
\chi(x,x',\beta,0)&\mapsto d_X(x,x')^{-\alpha}\cdot\left(u_{\mathscr{T}}(x,x')\cdot\psi(x',\beta,0)-\psi(x,\beta,0)\right),\\
\chi(x,x',\beta,1)&\mapsto d_X(x,x')^{-\alpha}\cdot\left(u_{\mathscr{T}}(x,x')^{-1}\cdot\psi(x',\beta,1)-\psi(x,\beta,1)\right).
\end{align*}
It exists and is unique by Lemma~A, applied to the independent generators that belong to a fundamental domain: the prescribed images lie in $\mathrm{Lin}$ and hence anticommute, by (Ge2).
The prescription then holds at all pairs of $\mathrm{Pzt}_k$: under $(x,x')\mapsto(x+\ell w,x'+\ell w)$ both sides are multiplied by $\boldsymbol{\chi}(w)$, by (0.c), because $u_{\mathscr{T}}$ and $d_X$ do not change.
In the notation of Section~I.5.4, $E(X;\mathscr{A},\mathscr{C},\psi,\delta_{\alpha,\mathscr{T}}\psi)$ is the element $\mathrm{sbs}_{\mathscr{T}}\left(E(X;\mathscr{A},\mathscr{C})\right)$: each generator $\chi$ is replaced by the H\"older quotient (I.1.7) of the generators $\psi$, with $\delta_{\alpha,-\mathscr{T}}$ on barred fields, and $u_{-\mathscr{T}}=u_{\mathscr{T}}^{-1}$.
The images under $\mathrm{sbs}_{\mathscr{T}}$ satisfy the relations of Proposition~17.1\,(u4) with $B=\mathscr{T}$ at all pairs of $\mathrm{Pzt}_k$, the pairs inside a cluster included; at the latter the relation reads $\chi(x,x')\mapsto\eta_k^{-\alpha}\cdot\left(\psi(x')-\psi(x)\right)$.
So $\mathrm{ev}=\mathrm{sbs}_{\mathscr{T}}$ satisfies the hypothesis of (a).
The same is true of $\mathrm{ev}=\mathrm{ev}_0\circ\mathrm{sbs}_{\mathscr{T}}$ for every homomorphism $\mathrm{ev}_0$ of the subalgebra generated by the $\psi(\xi)$ into an associative algebra; in particular it is true of the insertion of the fields of Section~I.5.4, $\psi=H_k(\mathscr{T})\Psi$ and $\bar\psi=\bar H_k(\mathscr{T})\bar\Psi$, wherever these operators are defined (the images are linear combinations of the generators $\Psi(u)$; Lemma~A).

\paragraph{(b)}
The pair $(0,0)$ lies in $\tilde R_k\times\mathfrak{X}_k$.
Indeed, in the notation of \cite[Definition~5.16]{QEDI} one has $0=A_0+A_1$ with $A_0=0$ (take $A'=0$ and $\lambda=0$ on every $\square^{\dagger}$) and $A_1=0$, and all the thresholds in that definition are positive.
At $(\mathscr{A},\mathscr{C})=(0,0)$ the total connection is $\mathscr{T}=0$.
Since $E$ is balanced, (G2.0.1) gives
\[
E^{[2,0]}(X;0,0)=\int_X\mathrm{d}x\int_X\mathrm{d}x'\;\bar\psi(x)^T\cdot E_{20}(X;0,0;x,x')\cdot\psi(x'),
\]
where the integrals extend over $X$ by (k2) and all arguments lie over $\tilde X$, by (I.AP).
Proposition~17.1\,(f), applied with $K=E_{20}(X;0,0;\cdot,\cdot)$ and $B=0$, and (Pp4) give
\[
F_2^{[2,0]}=\bar\psi_{\mathrm{av}}^T\cdot\left[\int_X\mathrm{d}x\int_X\mathrm{d}x'\,E_{20}(X;0,0;x,x')\right]\cdot\psi_{\mathrm{av}},
\]
and this vanishes when the matrix in the brackets vanishes.
Moreover $F_1^{[2,0]}=0$ by (2.a).
Hence the part of bidegree $(2,0)$ of $E^{\mathrm{nat}}(X;0,0)$ is the element $0$, and so is its antisymmetric kernel (Proposition~15\,(a)): $E^{\mathrm{nat}}_{20}(X;0,0;\xi_1,\xi_2)=0$ for all $\xi_1$, $\xi_2$.
In matrix notation, $E^{\mathrm{nat}}_{20}(X;0,0;x,x')=0$ for all $x=(c,g)$ and $x'=(c',g')$ over $X$, that is, on $(X\times O_{\mathrm{cub}})^2$.

The assertion of (b) that $E^{\mathrm{nat}}_{nm}=E_{nm}$ for $(n,m)\notin\{(2,0),(1,1),(0,2)\}$ holds at every $(\mathscr{A},\mathscr{C})$, and its proof does not use the balance of $E$.
By (2.a) and (Pp4), $(E^{\mathrm{nat}})^{[n,m]}=E^{[n,m]}$ whenever $(n,m)\notin\{(2,0),(1,1),(0,2)\}$, and the kernels are determined by the elements (Proposition~15\,(a)).
The bidegree $(0,0)$ is among these.

\paragraph{(c)}
Basis monomials are formed with a fundamental domain of $\ell\mathbb{Z}^3$ that contains the set $\tilde U$ of (Lf3), as in (Ge1).

\emph{The generators in $F_2$.}
By (2.a) and (1.a), $F_2$ is a linear combination of products of two elements of the set
\[
\mathrm{Gen}:=\left\{\psi(\tilde x,\beta,\omega):\tilde x\text{ over }\tilde X\right\}\cup\left\{\chi(p_i,p_{i+1},\beta,\omega):p_i\ne p_{i+1}\right\},
\]
where in the second set $(p_i,p_{i+1})$ ranges over the chain pairs of Proposition~17.1\,(c) formed for some $x$, $\tilde x$ over $\tilde X$ and some chain in $\mathrm{Ch}(x,\tilde x)$.
The elements of $\mathrm{Gen}$ are independent generators (Ge1).
A product of two of them is $0$ or, up to the sign, a basis monomial.
Consequently every basis monomial with a non-zero coefficient in $F_2$ is a product of two generators from $\mathrm{Gen}$.

\emph{Which generators occur in $E^{\mathrm{nat}}$.}
The coefficient of a basis monomial in $E^{\mathrm{nat}}(X;\mathscr{A},\mathscr{C})$ is the sum of its coefficients in $F_1$ and in $F_2$; if it does not vanish, one of the two does not vanish.
Suppose first that the coefficient in $F_1$ is not zero.
By (2.a) the monomial has a bidegree different from $(2,0)$ and the same coefficient in $E(X;\mathscr{A},\mathscr{C})$, so all its generators occur in $E(X;\mathscr{A},\mathscr{C})$.
A coefficient of $E(X;\mathscr{A},\mathscr{C})$ is a kernel multiplied by weights (Proposition~15\,(a)); so, by (k2), these generators are generators $\psi(\xi)$ with $\xi$ over $X$ and $\chi(\boldsymbol{\zeta})$ with $\boldsymbol{\zeta}$ admissible for $X$.
Suppose next that the coefficient in $F_2$ is not zero.
Then the generators of the monomial belong to $\mathrm{Gen}$; the generators $\psi(\tilde x,\beta,\omega)$ among them have $\tilde x$ over $X$, and the chain pairs are admissible for $X$, of kind (I), by Proposition~17.1\,(c1).

\emph{Location of the new generators $\chi(\boldsymbol{\zeta})$.}
Let $\chi(\boldsymbol{\zeta})$, $\boldsymbol{\zeta}=(p,p',\beta,\omega)$, occur in $E^{\mathrm{nat}}(X;\mathscr{A},\mathscr{C})$ but not in $E(X;\mathscr{A},\mathscr{C})$.
Then $\chi(\boldsymbol{\zeta})$ is a factor of a basis monomial whose coefficient in $E^{\mathrm{nat}}(X;\mathscr{A},\mathscr{C})$ is not zero.
The coefficient of this monomial in $F_1$ is zero, since otherwise $\chi(\boldsymbol{\zeta})$ would occur in $E(X;\mathscr{A},\mathscr{C})$.
So its coefficient in $F_2$ is not zero, and $\chi(\boldsymbol{\zeta})$ belongs to $\mathrm{Gen}$: $(p,p')=(p_i,p_{i+1})$ with $p_i\ne p_{i+1}$ for a chain as in Proposition~17.1\,(c).
By Proposition~17.1\,(c1) the pair is of kind (I) for $X$, $d_X(p,p')\le M$, and for $0\le i\le n$ the labels of $p$ and $p'$ are equal.
Different labels can therefore occur only for $i=n+1$, that is, at the last pair of the chain.

\emph{Dependence of the coefficients on $(\mathscr{A},\mathscr{C})$.}
By (2.a) and (1.a), the coefficients of $E^{\mathrm{nat}}(X;\mathscr{A},\mathscr{C})$ are built from the following quantities: the coefficients of $E(X;\mathscr{A},\mathscr{C})$, which are its kernels multiplied by products of weights (Proposition~15\,(a)); the numbers $\mathrm{cnt}(X)$, $\#|\mathrm{Ch}(x,\tilde x)|$ and $d_X(p_i,p_{i+1})^{\alpha}$, which involve no connection; and the transports $u_{\mathscr{T}}(p_i,p_{i+1})^{\pm1}$ at the chain pairs.
By Proposition~17.1\,(u2), (c1) and by (Lf4), these transports depend on $\mathscr{T}$ only through its values on bonds both of whose end points lie in $X$.

\emph{Property (k2) of $E^{\mathrm{nat}}$.}
For small $X$ the first part of (k2) is the assertion of (c) on the generators that occur in $E^{\mathrm{nat}}(X;\mathscr{A},\mathscr{C})$, proved above.
Let $(\mathscr{A},\mathscr{C})$ and $(\tilde{\mathscr{A}},\tilde{\mathscr{C}})$ in $\tilde R_k\times\mathfrak{X}_k$ coincide, in both components, on the bonds that meet $X$.
Then the kernels of $E(X;\cdot)$ at the two points coincide, by (k2) for $E$; and $\mathscr{A}+\mathscr{C}=\tilde{\mathscr{A}}+\tilde{\mathscr{C}}$ on the bonds with both end points in $X$, because these bonds meet $X$ (Section~I.1.1, ``Bonds and sets'').
By the assertion of (c) on the dependence of the coefficients on $(\mathscr{A},\mathscr{C})$, proved above, $E^{\mathrm{nat}}(X;\mathscr{A},\mathscr{C})=E^{\mathrm{nat}}(X;\tilde{\mathscr{A}},\tilde{\mathscr{C}})$.
For large $X$ one has $E^{\mathrm{nat}}(X)=E(X)$.

\paragraph{(d)}
Let $\mathsf{w}=(\mathsf{w}_1,\mathsf{w}_2)$ be an arbitrary weight vector.

\emph{Step 1 (one generator).}
From (1.a), (Ge4), (Pp1) and Proposition~17.1\,(c3), the weights of the means adding up to $1$, one obtains
\begin{equation}\tag{2.c}
\begin{aligned}
\left\|\mathrm{Lpsi}(x,\beta,\omega)\right\|_{\mathsf{w}}&\le\mathrm{cnt}(X)^{-1}\sum_{\tilde x\text{ over }\tilde X}\#|\mathrm{Ch}(x,\tilde x)|^{-1}\sum_{\mathrm{Ch}(x,\tilde x)}\left|U^{(\omega)}_{n+2}\right|\cdot\mathsf{w}_1\le\mathrm{e}\cdot\mathsf{w}_1,\\
\left\|\mathrm{Lchi}(x,\beta,\omega)\right\|_{\mathsf{w}}&\le\mathrm{cnt}(X)^{-1}\sum_{\tilde x\text{ over }\tilde X}\#|\mathrm{Ch}(x,\tilde x)|^{-1}\sum_{\mathrm{Ch}(x,\tilde x)}\ \sum_{i:\,p_i\ne p_{i+1}}\left|U^{(\omega)}_i\right|\cdot d_X(p_i,p_{i+1})^{\alpha}\cdot\mathsf{w}_2\\
&\le\mathrm{e}\cdot(3c_X\cdot L\cdot M)\cdot\mathsf{w}_2.
\end{aligned}
\end{equation}
(The sums over the photon site $\tilde c$ of $\tilde x$ and over the label of $\tilde x$ are means; so the bounds (2.c) are independent of $\eta_k$ and of the number $48$ of the labels.)

\emph{Step 2 (the quadratic part).}
The family $E$ is balanced.
By (G2.0.1), the integrals extending over $X$ by (k2),
\[
E^{[2,0]}(X)=\int_X\mathrm{d}x\int_X\mathrm{d}x'\,\sum_{\beta,\beta'}E_{20}(X;x,x')_{\beta\beta'}\cdot\psi(x,\beta,1)\cdot\psi(x',\beta',0),
\]
and
\begin{equation}\tag{2.d}
\mathrm{sE}:=\sum_{\beta,\beta'}\int_X\mathrm{d}x\int_X\mathrm{d}x'\,\left|E_{20}(X;\mathscr{A},\mathscr{C};x,x')_{\beta\beta'}\right|=\frac12\left\|E_{20}(X)\right\|,\qquad\left\|E^{[2,0]}(X)\right\|_{\mathsf{w}}=\mathsf{w}_1^2\cdot\mathrm{sE}.
\end{equation}
Consequently
\[
F_2=\int_X\mathrm{d}x\int_X\mathrm{d}x'\,\sum_{\beta,\beta'}E_{20}(X;x,x')_{\beta\beta'}\cdot\left(\mathrm{Lpsi}(x,\beta,1)+\mathrm{Lchi}(x,\beta,1)\right)\cdot\left(\mathrm{Lpsi}(x',\beta',0)+\mathrm{Lchi}(x',\beta',0)\right).
\]
By (Ge4), Proposition~15\,(b) and (2.c),
\begin{equation}\tag{2.e}
\begin{aligned}
\left\|F_2^{[2,0]}\right\|_{\mathsf{w}}&\le\mathrm{e}^2\cdot\mathsf{w}_1^2\cdot\mathrm{sE},\\
\left\|F_2^{[1,1]}\right\|_{\mathsf{w}}&\le2\mathrm{e}^2\cdot(3c_X\cdot L\cdot M)\cdot\mathsf{w}_1\cdot\mathsf{w}_2\cdot\mathrm{sE},\\
\left\|F_2^{[0,2]}\right\|_{\mathsf{w}}&\le\mathrm{e}^2\cdot(3c_X\cdot L\cdot M)^2\cdot\mathsf{w}_2^2\cdot\mathrm{sE}.
\end{aligned}
\end{equation}

\emph{Step 3 (kernels).}
An element of bidegree $(n,m)$ with the antisymmetric kernel $K_{nm}$ has the norm $\frac{\mathsf{w}_1^n\cdot\mathsf{w}_2^m}{n!\,m!}\cdot\left\|K_{nm}\right\|$; this is the case of equality in Proposition~15\,(a).
By (2.a) and (Pp4),
\[
(E^{\mathrm{nat}})^{[2,0]}=F_2^{[2,0]},\qquad(E^{\mathrm{nat}})^{[1,1]}=E^{[1,1]}+F_2^{[1,1]},\qquad(E^{\mathrm{nat}})^{[0,2]}=E^{[0,2]}+F_2^{[0,2]}.
\]
Together with (2.e) and (2.d) this gives
\begin{align*}
\frac{\mathsf{w}_1^2}{2}\cdot\left\|E^{\mathrm{nat}}_{20}(X)\right\|&\le\mathrm{e}^2\cdot\mathsf{w}_1^2\cdot\mathrm{sE}=\frac{\mathsf{w}_1^2}{2}\cdot\mathrm{e}^2\cdot\left\|E_{20}(X)\right\|,\\
\mathsf{w}_1\cdot\mathsf{w}_2\cdot\left\|E^{\mathrm{nat}}_{11}(X)\right\|&\le\mathsf{w}_1\cdot\mathsf{w}_2\cdot\left\|E_{11}(X)\right\|+2\mathrm{e}^2\cdot(3c_X\cdot L\cdot M)\cdot\mathsf{w}_1\cdot\mathsf{w}_2\cdot\mathrm{sE}\\
&=\mathsf{w}_1\cdot\mathsf{w}_2\cdot\left[\left\|E_{11}(X)\right\|+\mathrm{e}^2\cdot(3c_X\cdot L\cdot M)\cdot\left\|E_{20}(X)\right\|\right],\\
\frac{\mathsf{w}_2^2}{2}\cdot\left\|E^{\mathrm{nat}}_{02}(X)\right\|&\le\frac{\mathsf{w}_2^2}{2}\cdot\left\|E_{02}(X)\right\|+\mathrm{e}^2\cdot(3c_X\cdot L\cdot M)^2\cdot\mathsf{w}_2^2\cdot\mathrm{sE}\\
&=\frac{\mathsf{w}_2^2}{2}\cdot\left[\left\|E_{02}(X)\right\|+\mathrm{e}^2\cdot(3c_X\cdot L\cdot M)^2\cdot\left\|E_{20}(X)\right\|\right].
\end{align*}
Division by the positive factors on the left gives the three bounds on kernels.

\emph{Step 4 (the norm).}
Let $\mathsf{w}=h_k$; then $h_{k,2}=e_k^{\varepsilon}\cdot h_{k,1}$ (Section~I.5.1).
The norm is additive over the bidegrees, by (G2.r7).
By (2.e) and (2.d), with $\mathrm{tq}:=3c_X\cdot L\cdot M\cdot e_k^{\varepsilon}$,
\begin{align*}
\left\|F_2\right\|_{h_k}&\le\mathrm{e}^2\cdot\mathrm{sE}\cdot\left(h_{k,1}+3c_X\cdot L\cdot M\cdot h_{k,2}\right)^2\\
&=\mathrm{e}^2\cdot(1+\mathrm{tq})^2\cdot h_{k,1}^2\cdot\mathrm{sE}=\mathrm{e}^2\cdot(1+\mathrm{tq})^2\cdot\left\|E^{[2,0]}(X)\right\|_{h_k},\\
\left\|E^{\mathrm{nat}}(X)\right\|_{h_k}&\le\left\|F_1\right\|_{h_k}+\left\|F_2\right\|_{h_k}\\
&\le\left\|E(X)\right\|_{h_k}-\left\|E^{[2,0]}(X)\right\|_{h_k}+\mathrm{e}^2\cdot(1+\mathrm{tq})^2\cdot\left\|E^{[2,0]}(X)\right\|_{h_k}\\
&\le\mathrm{e}^2\cdot(1+\mathrm{tq})^2\cdot\left\|E(X)\right\|_{h_k},
\end{align*}
the last inequality because $\mathrm{e}^2\cdot(1+\mathrm{tq})^2\ge1$ and $\left\|E^{[2,0]}(X)\right\|_{h_k}\le\left\|E(X)\right\|_{h_k}$.
By (A9), $90c_X\cdot L\cdot M\cdot e_0^{\varepsilon}\le1$; hence $\mathrm{tq}\le3c_X\cdot L\cdot M\cdot e_0^{\varepsilon}\le\frac{1}{30}$.
Finally, $\mathrm{e}<2.71829$ gives $\mathrm{e}^2<7.38911$, and $7.38911\cdot\frac{961}{900}<7.89$; so $\mathrm{e}^2\cdot\left(\frac{31}{30}\right)^2<7.9\le10$.

\paragraph{(e)}
For large $X$ all clauses hold, because $E^{\mathrm{nat}}(X)=E(X)$; and $\tilde gX$ is large if and only if $X$ is, since $d_M(\tilde gX)=d_M(X)$ by Proposition~16.1\,(a1).
Let $X$ be small.

\emph{Linearity.}
The map $E\mapsto E^{[2,0]}(X;\mathscr{A},\mathscr{C})$ is linear, and $\boldsymbol{\Theta}_{X,\mathscr{T}}$ is a linear map that does not involve $E$.

\emph{Evenness and balance.}
$F_1$ is a sum of parts $E^{[n,m]}$ of an even balanced element, and each basis monomial of $F_1$ is a basis monomial of $E(X)$.
By (G2.0.1) and Proposition~17.1\,(d), by which barred generators are mapped to combinations of barred ones and unbarred generators to combinations of unbarred ones, every term of $F_2$ is the product of a combination of barred generators and a combination of unbarred ones: it has degree two, with one barred and one unbarred leg.

\emph{(k1).}
As shown in the proof of (c), under the dependence of the coefficients on $(\mathscr{A},\mathscr{C})$, every coefficient of $E^{\mathrm{nat}}(X;\mathscr{A},\mathscr{C})$ is a finite sum of products of coefficients of $E(X;\mathscr{A},\mathscr{C})$, of numbers that involve no connection, and of transports $u_{\mathscr{A}+\mathscr{C}}(p_i,p_{i+1})^{\pm1}$.
The coefficients of $E(X;\mathscr{A},\mathscr{C})$ are analytic on $\tilde R_k\times\mathfrak{X}_k$ by (k1) for $E$ and Proposition~15\,(a); the transports are entire functions of the numbers $\mathscr{A}(b)+\mathscr{C}(b)$, by Proposition~17.1\,(u2).
The kernels of $E^{\mathrm{nat}}(X)$ are its coefficients divided by products of weights (Proposition~15\,(a)).

\emph{(k2).}
This is clause (c).

\emph{(k6).}
By (d),
\[
\exp\left(\kappa\cdot d_M(X)\right)\cdot\left\|E^{\mathrm{nat}}(X;\mathscr{A},\mathscr{C})\right\|_{h_k}\le10\exp\left(\kappa\cdot d_M(X)\right)\cdot\left\|E(X;\mathscr{A},\mathscr{C})\right\|_{h_k}\le10\left\|E\right\|^{x}_{k};
\]
for large $X$ the factor $10$ is replaced by $1$.
Hence $\left\|E^{\mathrm{nat}}\right\|^{x}_{k}\le10\left\|E\right\|^{x}_{k}<\infty$.

\emph{(k3).}
Let $\tilde g\in\mathbb{G}_k$.
The pair $(\tilde g\mathscr{A},\tilde g\mathscr{C})$ lies in $\tilde R_k\times\mathfrak{X}_k$ (Proposition~16.1\,(c)), the polymer $\tilde gX$ is small, and $\tilde g\mathscr{A}+\tilde g\mathscr{C}=\tilde g\mathscr{T}$.
Since $a_{\tilde g}$ preserves the bidegree (Proposition~16.1\,(b)), the property (k3) of $E$ holds for each part $E^{[n,m]}$ separately:
\[
a_{\tilde g}\left[E^{[n,m]}(\tilde gX;\tilde g\mathscr{A},\tilde g\mathscr{C})\right]=E^{[n,m]}(X;\mathscr{A},\mathscr{C}).
\]
The first relation of (3.4), with $B=\mathscr{T}$ and applied to the element $E^{[2,0]}(\tilde gX;\tilde g\mathscr{A},\tilde g\mathscr{C})$ of $\mathrm{Grp}[\tilde gX]$, now gives
\begin{align*}
a_{\tilde g}\left[E^{\mathrm{nat}}(\tilde gX;\tilde g\mathscr{A},\tilde g\mathscr{C})\right]&=a_{\tilde g}\left[\left(E-E^{[2,0]}\right)(\tilde gX;\tilde g\mathscr{A},\tilde g\mathscr{C})\right]+a_{\tilde g}\,\boldsymbol{\Theta}_{\tilde gX,\tilde g\mathscr{T}}\left(E^{[2,0]}(\tilde gX;\tilde g\mathscr{A},\tilde g\mathscr{C})\right)\\
&=\left(E-E^{[2,0]}\right)(X;\mathscr{A},\mathscr{C})+\boldsymbol{\Theta}_{X,\mathscr{T}}\,a_{\tilde g}\left(E^{[2,0]}(\tilde gX;\tilde g\mathscr{A},\tilde g\mathscr{C})\right)\\
&=\left(E-E^{[2,0]}\right)(X;\mathscr{A},\mathscr{C})+\boldsymbol{\Theta}_{X,\mathscr{T}}\left(E^{[2,0]}(X;\mathscr{A},\mathscr{C})\right)=E^{\mathrm{nat}}(X;\mathscr{A},\mathscr{C}).
\end{align*}

\emph{(k4).}
Let $(\mathscr{A},\mathscr{C})$ and $(\mathscr{A}^1,\mathscr{C}^1)$ lie in $\tilde R_k\times\mathfrak{X}_k$, let $\boldsymbol{\lambda}$, $\boldsymbol{\lambda}'$ be real functions of $c$, and let $\mathscr{A}^1=\mathscr{A}-d\boldsymbol{\lambda}$ and $\mathscr{C}^1=\mathscr{C}-d\boldsymbol{\lambda}'$ on the bonds at distance $\le1$ from $X$.
Put $\Lambda:=\boldsymbol{\lambda}+\boldsymbol{\lambda}'$ and $\mathscr{T}^1:=\mathscr{A}^1+\mathscr{C}^1$.
A bond with both end points in $X$ is at distance $\le1$ from $X$, because both its end points have distance $0$ from $X$ (Section~I.1.1, ``Bonds and sets''); hence $\mathscr{T}^1=\mathscr{T}-d\Lambda$ on the bonds with both end points in $X$.
As $G_\Lambda$ preserves the bidegree, $G_\Lambda\left[E^{[n,m]}(X;\mathscr{A}^1,\mathscr{C}^1)\right]=E^{[n,m]}(X;\mathscr{A},\mathscr{C})$.
The second relation of (3.4), with $B=\mathscr{T}$ and $B^1=\mathscr{T}^1$, yields
\begin{align*}
G_\Lambda\left[E^{\mathrm{nat}}(X;\mathscr{A}^1,\mathscr{C}^1)\right]&=\left(E-E^{[2,0]}\right)(X;\mathscr{A},\mathscr{C})+G_\Lambda\,\boldsymbol{\Theta}_{X,\mathscr{T}^1}\left(E^{[2,0]}(X;\mathscr{A}^1,\mathscr{C}^1)\right)\\
&=\left(E-E^{[2,0]}\right)(X;\mathscr{A},\mathscr{C})+\boldsymbol{\Theta}_{X,\mathscr{T}}\,G_\Lambda\left(E^{[2,0]}(X;\mathscr{A}^1,\mathscr{C}^1)\right)=E^{\mathrm{nat}}(X;\mathscr{A},\mathscr{C}).
\end{align*}

\emph{(k5).}
The pair $(-\mathscr{A},-\mathscr{C})$ lies in $\tilde R_k\times\mathfrak{X}_k$ (Proposition~16.1\,(c)), and its total connection is $-\mathscr{T}$.
The automorphism $\mathrm{Cg}$ preserves the bidegree.
With the third relation of (3.4),
\begin{align*}
\mathrm{Cg}\left[E^{\mathrm{nat}}(X;-\mathscr{A},-\mathscr{C})\right]&=\left(E-E^{[2,0]}\right)(X;\mathscr{A},\mathscr{C})+\mathrm{Cg}\,\boldsymbol{\Theta}_{X,-\mathscr{T}}\left(E^{[2,0]}(X;-\mathscr{A},-\mathscr{C})\right)\\
&=\left(E-E^{[2,0]}\right)(X;\mathscr{A},\mathscr{C})+\boldsymbol{\Theta}_{X,\mathscr{T}}\,\mathrm{Cg}\left(E^{[2,0]}(X;-\mathscr{A},-\mathscr{C})\right)=E^{\mathrm{nat}}(X;\mathscr{A},\mathscr{C}).
\end{align*}
Therefore $E^{\mathrm{nat}}$ lies in $K^{x,C}_{k}$ if $E$ does.

\emph{(k7).}
Let $E$ lie in $K^{x}_{k}$.
The pair $(-\overline{\mathscr{A}},-\overline{\mathscr{C}})$ lies in $\tilde R_k\times\mathfrak{X}_k$ (Proposition~16.1\,(c)), and its total connection is $-\overline{\mathscr{T}}$.
The map $\mathbf{K}$ preserves the bidegree.
With the fourth relation of (3.4),
\begin{align*}
\mathbf{K}\left[E^{\mathrm{nat}}(X;-\overline{\mathscr{A}},-\overline{\mathscr{C}})\right]&=\left(E-E^{[2,0]}\right)(X;\mathscr{A},\mathscr{C})+\boldsymbol{\Theta}_{X,\mathscr{T}}\,\mathbf{K}\left(E^{[2,0]}(X;-\overline{\mathscr{A}},-\overline{\mathscr{C}})\right)\\
&=E^{\mathrm{nat}}(X;\mathscr{A},\mathscr{C}).
\end{align*}

\emph{Restriction to $\mathscr{C}=0$.}
The bond function $0$ lies in $\mathfrak{X}_k$, since the two thresholds in the definition of $\mathfrak{X}_k$ are positive (\cite[Definition~5.16\,(f)]{QEDI}).
By Section~I.5.4, $E|_0(X;\mathscr{A})=E(X;\mathscr{A},0)$.
Its kernels are those of $E(X;\mathscr{A},0)$; so they vanish unless all $\xi$ lie over $X$ and all $\boldsymbol{\zeta}$ are admissible for $X$, and they depend on $\mathscr{A}(b)$ only for the bonds $b$ that meet $X$: $E|_0$ has the property (k2)$^{\mathrm{o}}$ of Section~I.5.11.
At $\mathscr{C}=0$ one has $\mathscr{T}=\mathscr{A}$.
Hence, by the definition of $E^{\mathrm{nat}}$ for a family without slot (Section~I.5.11; it is the adjustment of Proposition~48.3), for small $X$
\begin{align*}
(E|_0)^{\mathrm{nat}}(X;\mathscr{A})&=E(X;\mathscr{A},0)-E^{[2,0]}(X;\mathscr{A},0)+\boldsymbol{\Theta}_{X,\mathscr{A}}\left(E^{[2,0]}(X;\mathscr{A},0)\right)\\
&=E^{\mathrm{nat}}(X;\mathscr{A},0)=E^{\mathrm{nat}}|_0(X;\mathscr{A}),
\end{align*}
and for large $X$ both sides equal $E(X;\mathscr{A},0)$.

\paragraph{(f)}
For a finite set $\tilde U$ of sites of $\mathbb{L}_k$, the algebra $\mathrm{Grc}_k(\tilde U)$ and the map $\mathrm{cov}_\ell$ are those of Section~I.1.9.
For $v\in M\mathbb{Z}^3$ let $\mathrm{tl}_v\colon\mathrm{Grc}_k(\tilde U)\to\mathrm{Grc}_k(\tilde U+v)$ be the isomorphism with
\[
\psi(x,\beta,\omega)\mapsto\psi(x+v,\beta,\omega),\qquad\chi(x,x',\beta,\omega)\mapsto\chi(x+v,x'+v,\beta,\omega);
\]
it maps the generators $\Psi(u)$, $W(u)$ of $\mathrm{Grc}_k(\tilde U)$ to the generators whose index has its site translated by $v$, but these two families do not occur below.
In the terms of Section~I.1.9, $T_v$ acts as $\mathrm{tl}_v$ on elements, as $\mathrm{Hom}\mapsto\mathrm{tl}_v\circ\mathrm{Hom}\circ\mathrm{tl}_v^{-1}$ on homomorphisms, and as $(T_vf)(t)=f(t-v)$ on bond functions and, in all arguments, on kernels.
Steps 1--3 prove the assertions of (f) on $\boldsymbol{\Theta}_{X,B}$ and on the corresponding map; Steps 4--6 prove the assertion on $E^{\mathrm{nat}}(X;\mathscr{A},\mathscr{C})$ together with the formula for its map, for every polymer $X$ that has a lift $\tilde X$ with $\operatorname{sup-diameter}(\tilde X)+2\eta_k+M<\ell$; Step 7 proves the consequence.

\emph{Step 1 (the homomorphism on the covering lattice).}
Let $\tilde X$ be a non-empty finite union of $M$-cubes of $\mathbb{L}_k$ that is connected through faces, and let $\tilde B$ be a complex function on the bonds of $\mathbb{L}_k$ with both end points in $\tilde X$ such that $\tilde B(c',c)=-\tilde B(c,c')$.
Let $\mathrm{Grcp}_k(\tilde X)$ be the subalgebra with unit of $\mathrm{Grc}_k(\tilde X)$ generated by the $\psi(x,\beta,\omega)$ with $c\in\tilde X$.
For $x$, $\tilde x$ over $\tilde X$ we form, as in Proposition~17.1\,(c), the set $\mathrm{Ch}(x,\tilde x)$ of the shortest chains of $M$-cubes of $\tilde X$ in which consecutive cubes have a common face, the centres $z_i$, the chain points $p_0=x=(c,g)$, $p_{i+1}=(z_i,g)$, $p_{n+2}=\tilde x$, and the distances $d_X(p_i,p_{i+1})$.
The argument given for the set $\mathrm{Ch}(x,\tilde x)$ in the proof of Proposition~17.1\,(c) requires of $\tilde X$ only that it is a finite union of $M$-cubes connected through faces; hence $\mathrm{Ch}(x,\tilde x)$ is finite and not empty, and the members of a shortest chain are pairwise different.
The three cases in the proof of Proposition~17.1\,(c1) require only that the $\square_i$ are $M$-cubes of $\tilde X$ and that consecutive ones are different and have a common face.
So every pair $(p_i,p_{i+1})$ with $p_i\ne p_{i+1}$ lies in $\mathrm{Pzt}_k$, and its photon sites and its six paths lie in $\tilde X$.
Hence the transports
\[
u_{\tilde B}(p_i,p_{i+1}):=\exp\left\{i\cdot e_k\cdot\eta_k\cdot\frac16\sum_{\pi}\tilde B\left(\boldsymbol{\Gamma}^{\pi}(\pi(p_i),\pi(p_{i+1}))\right)\right\}
\]
are defined and read $\tilde B$ on bonds with both end points in $\tilde X$ only, and $\chi(p_i,p_{i+1},\beta,\omega)$ is a generator of $\mathrm{Grc}_k(\tilde X)$.
Let
\[
\Phi^{\boldsymbol{\Theta}}(\tilde X,\tilde B):=\boldsymbol{\Theta}_{\tilde X,\tilde B}\colon\mathrm{Grcp}_k(\tilde X)\to\mathrm{Grc}_k(\tilde X)
\]
be the homomorphism given by (1.a), read in $\mathrm{Grc}_k(\tilde X)$, with the transports $u_{\tilde B}$ and with $48(M\cdot L^k)^3\cdot|\tilde X|_M$ in place of $\mathrm{cnt}(X)$.
It exists and is unique by Lemma~A, since the right sides anticommute pairwise, each also with itself.
Indeed, they lie in the linear span of the independent generators of $\mathrm{Grc}_k(\tilde X)$; for $v$, $v'$ in this span the expression $v\cdot v'+v'\cdot v$ is bilinear in $(v,v')$, and it vanishes when $v$ and $v'$ are independent generators of $\mathrm{Grc}_k(\tilde X)$, equal or not; hence it vanishes on the whole span, as in (Ge2).
It is formed with the cubes of $\tilde X$, their centres and common faces (through the chains and the chain points), the labels (the label $g$ of $x$ is kept along the chain, and the label of $\tilde x$ is averaged over), the weights of (I.1.3) (by (0.d)), and the transports.
It involves $k$, $M$, $\alpha$ and $e_k$, and no side $\ell$.

\emph{Step 2 (translations).}
Let $v\in M\mathbb{Z}^3$.
The translation by $v$ maps the $M$-cubes of $\tilde X$ onto those of $\tilde X+v$, centres onto centres, and common faces onto common faces; it therefore maps $\mathrm{Ch}(x,\tilde x)$ onto $\mathrm{Ch}(x+v,\tilde x+v)$, and the chain points of the image chain are the points $p_i+v$, with the same labels.
One has $d_X(p_i+v,p_{i+1}+v)=d_X(p_i,p_{i+1})$.
Since $\boldsymbol{\Gamma}^{\pi}(c+v,c'+v)=\boldsymbol{\Gamma}^{\pi}(c,c')+v$ and $(T_v\tilde B)(b+v)=\tilde B(b)$, also $u_{T_v\tilde B}(p_i+v,p_{i+1}+v)=u_{\tilde B}(p_i,p_{i+1})$; and $|\tilde X+v|_M=|\tilde X|_M$.
Hence (1.a) for $(\tilde X+v,T_v\tilde B)$ at $x+v$ is the image under $\mathrm{tl}_v$ of (1.a) for $(\tilde X,\tilde B)$ at $x$:
\[
\boldsymbol{\Theta}_{\tilde X+v,T_v\tilde B}\,\mathrm{tl}_v\,\psi(x,\beta,\omega)=\mathrm{tl}_v\,\boldsymbol{\Theta}_{\tilde X,\tilde B}\,\psi(x,\beta,\omega).
\]
By (Un),
\[
\Phi^{\boldsymbol{\Theta}}(\tilde X+v,T_v\tilde B)=\mathrm{tl}_v\circ\Phi^{\boldsymbol{\Theta}}(\tilde X,\tilde B)\circ\mathrm{tl}_v^{-1}\qquad\text{for all }v\in M\mathbb{Z}^3,\text{ in particular for }v\in(L\cdot M)\mathbb{Z}^3.
\]

\emph{Step 3 (the torus; $\mathrm{rg}=0$).}
Let $\ell\in\mathrm{Cl}$, let $X$ be a small polymer of $T_k=\mathbb{L}_k/\ell\mathbb{Z}^3$ with a lift $\tilde X$, let $B$ be a complex bond function on $T_k$, and put $\tilde B(b):=B(\mathrm{pr}\,b)$ on the bonds $b$ with both end points in $\tilde X$.
These are the bonds within the distance $0$ of $\tilde X$ in the sense of (I.f4) in Section~I.1.9.
By (0.b), $\operatorname{sup-diameter}(\tilde X)+2\cdot0+M<\ell$; so $X$ is admitted in Section~I.1.9, and $\mathrm{cov}_\ell$ is defined on $\mathrm{Grc}_k(\tilde X)$.
The map $\mathrm{cov}_\ell$ sends the generator $\psi(x,\beta,\omega)$, resp.\ $\chi(x,x',\beta,\omega)$, of $\mathrm{Grc}_k(\tilde X)$ to the generator of the torus with the same index.
By (Lf5) the chains of $\tilde X$ are those used in Proposition~17.1\,(c), and $u_{\tilde B}=u_B$ at the chain pairs (Section~0.1).
Hence $\mathrm{cov}_\ell$ maps the right side of (1.a) in $\mathrm{Grc}_k(\tilde X)$ to the right side of (1.a) in $\mathrm{Gr}[\psi,\chi]$:
\[
\mathrm{cov}_\ell\,\boldsymbol{\Theta}_{\tilde X,\tilde B}\,\psi(x,\beta,\omega)=\boldsymbol{\Theta}_{X,B}\,\psi(x,\beta,\omega)=\boldsymbol{\Theta}_{X,B}\,\mathrm{cov}_\ell\,\psi(x,\beta,\omega).
\]
By (Un),
\begin{equation}\tag{2.f}
\mathrm{cov}_\ell\circ\Phi^{\boldsymbol{\Theta}}(\tilde X,\tilde B)=\boldsymbol{\Theta}_{X,B}\circ\mathrm{cov}_\ell\qquad\text{on }\mathrm{Grcp}_k(\tilde X).
\end{equation}
Steps 1--3 establish the definition of Section~I.1.9 for homomorphisms, with $\mathrm{rg}=0$ and the data $\tilde B$.
(The right side of (2.f) does not depend on the lift, by Proposition~17.1\,(d).)

\emph{Step 4 (the indices of $E$ over a lift).}
Let $\ell\in\mathrm{Cl}$, and let $X$ be a polymer of $T_k$, small or large, with a lift $\tilde X$ such that $\operatorname{sup-diameter}(\tilde X)+2\eta_k+M<\ell$ (for small $X$ this is (0.b)).
Put $\tilde U:=\left\{c:\operatorname{dist}(c,\tilde X)\le\eta_k\right\}$.
Then $\operatorname{sup-diameter}(\tilde U)+M<\ell$, the covering map $\mathrm{pr}$ is injective on $\tilde U$, and $\mathrm{cov}_\ell$ is defined on $\mathrm{Grc}_k(\tilde U)$.
Let $\mathrm{Pad}(\tilde X)$ be the set of the pairs $(x,x')\in\mathrm{Pzt}_k$ for which either $c$, $c'$ and the six paths $\boldsymbol{\Gamma}^{\pi}(c,c')$ lie in $\tilde X$ (for $c=c'$: $c\in\tilde X$), or $x'=x^{(3)}$ and $c$ or $c_3$ lies in $\tilde X$.
The photon sites of these pairs lie in $\tilde U$.
We claim that the map sending $(x,x')$ to its class is a bijection of $\mathrm{Pad}(\tilde X)$ onto the set of the pairs of $\mathrm{Pz}_k$ that are admissible for $X$.
\begin{itemize}
\item The class of a pair of $\mathrm{Pad}(\tilde X)$ is admissible for $X$, of the same kind, because $\mathrm{pr}$ maps $\tilde X$ onto $X$ and the six paths onto the six paths of the class.
\item Injectivity: two pairs of $\mathrm{Pad}(\tilde X)$ in one class differ by a translation $\ell w$; their first points lie over $\tilde U$, on which $\mathrm{pr}$ is injective; hence $w=0$.
\item Surjectivity onto the pairs of kind (I): let $(x,x')\in\mathrm{Pzt}_k$ represent a pair of kind (I) for $X$; after a translation by an element of $\ell\mathbb{Z}^3$ we may assume $c\in\tilde X$.
Let $c''$ be $c'$ or a site of one of the six paths $\boldsymbol{\Gamma}^{\pi}(c,c')$; then $\mathrm{pr}\,c''$ lies in $X$ and $|c''-c|\le M$.
Let $\tilde c''$ be the point of $\tilde X$ over $\mathrm{pr}\,c''$.
By Section~I.1.1 (with $\mathrm{sd}=M$; here $\operatorname{sup-diameter}(\tilde X)+M<\ell$ and $d(\mathrm{pr}\,c'',\mathrm{pr}\,c)\le M$), the vector $\tilde c''-c$ is the only representative of $\mathrm{pr}\,c''-\mathrm{pr}\,c$ modulo $\ell\mathbb{Z}^3$ of sup norm $\le M$.
Since $c''-c$ is such a representative, $c''=\tilde c''$, a point of $\tilde X$.
\item Surjectivity onto the pairs of kind (II): here $x'=x^{(3)}$, and $\mathrm{pr}\,c$ or $\mathrm{pr}\,c_3$ lies in $X$; after a translation this point lies in $\tilde X$.
(If both lie in $X$, then both lie in $\tilde X$, by Section~I.1.1 with $\mathrm{sd}=\eta_k$.)
\end{itemize}

\emph{Step 5 (the map for $E^{\mathrm{nat}}$; $\mathrm{rg}=\eta_k$).}
The data are $\tilde B$ as in Step~1 and a family $K=(K_{nm})$ of kernels $K_{nm}(\xi_1,\dots,\xi_n;\boldsymbol{\zeta}_1,\dots,\boldsymbol{\zeta}_m)$, antisymmetric in the $\xi$ and in the $\boldsymbol{\zeta}$, at the indices $\xi=(x,\beta,\omega)$ with $x$ over $\tilde X$ and $\boldsymbol{\zeta}=(x,x',\beta,\omega)$ with $(x,x')\in\mathrm{Pad}(\tilde X)$.
All these indices and bonds lie within the distance $\eta_k$ of $\tilde X$.
Put
\begin{align*}
F[K]&:=\sum_{n,m}\frac{1}{n!\,m!}\sum_{\xi_1,\dots,\xi_n}\ \sum_{\boldsymbol{\zeta}_1,\dots,\boldsymbol{\zeta}_m}\left(\frac{\eta_k^3}{48}\right)^{n+2m}\cdot K_{nm}(\xi;\boldsymbol{\zeta})\cdot\psi(\xi_1)\cdots\psi(\xi_n)\cdot\chi(\boldsymbol{\zeta}_1)\cdots\chi(\boldsymbol{\zeta}_m),\\
\Phi^{\mathrm{nat}}(\tilde X;K,\tilde B)&:=F[K]-F[K]^{[2,0]}+\boldsymbol{\Theta}_{\tilde X,\tilde B}\left(F[K]^{[2,0]}\right)\qquad\text{if }d_M(\tilde X)\le L,\\
\Phi^{\mathrm{nat}}(\tilde X;K,\tilde B)&:=F[K]\qquad\text{if }d_M(\tilde X)>L.
\end{align*}
Here $F[K]$ is an element of $\mathrm{Grc}_k(\tilde U)$, $d_M(\tilde X)$ is formed with nets in $\mathbb{R}^3$, $F[K]^{[2,0]}$ lies in $\mathrm{Grcp}_k(\tilde X)$, and $\mathrm{Grc}_k(\tilde X)$ is a subalgebra of $\mathrm{Grc}_k(\tilde U)$.
The map $\Phi^{\mathrm{nat}}$ involves no side $\ell$.
For $v\in M\mathbb{Z}^3$ one has $\mathrm{tl}_v\,F[K]=F[T_vK]$, the right side formed with $\tilde X+v$: the weights are constants, and $\mathrm{Pad}(\tilde X+v)=\mathrm{Pad}(\tilde X)+v$ because $(x+v)^{(3)}=x^{(3)}+v$ by (I.1.2).
Moreover $\mathrm{tl}_v$ preserves the bidegree, and $d_M(\tilde X+v)=d_M(\tilde X)$.
Together with Step~2 this gives
\[
\Phi^{\mathrm{nat}}(\tilde X+v;T_vK,T_v\tilde B)=\mathrm{tl}_v\,\Phi^{\mathrm{nat}}(\tilde X;K,\tilde B).
\]

\emph{Step 6 (the torus).}
Let $E$ be a family on $T_k$ with (k2), let $(\mathscr{A},\mathscr{C})\in\tilde R_k\times\mathfrak{X}_k$, let $K$ consist of the kernels $E_{nm}(X;\mathscr{A},\mathscr{C};\cdot)$ at the indices of Step~5, and put $\tilde B(b):=\mathscr{T}(\mathrm{pr}\,b)$.
\begin{itemize}
\item[(i)] $E(X;\mathscr{A},\mathscr{C})=\mathrm{cov}_\ell\left(F[K]\right)$.
Indeed, by (Pp2), whose argument uses only (I.AP), (0.c) and (k2), each argument $\xi$ may be summed over its representatives over $\tilde X$ and each argument $\boldsymbol{\zeta}$ over its representatives in $\mathrm{Pad}(\tilde X)$ (Step~4), with the weights $\frac{\eta_k^3}{48}$ and $\left(\frac{\eta_k^3}{48}\right)^2$ of Section~I.3.6; and $\mathrm{cov}_\ell$ maps each generator to the generator with the same index.
Moreover $F[K]$ is the element of $\mathrm{Grc}_k(\tilde U)$ whose kernels are those of $E(X;\mathscr{A},\mathscr{C})$ restricted to the indices over $\tilde U$, that is, the element $F$ of Section~I.1.9 and of the statement.
For at the indices over $\tilde U$ that are not indices of Step~5 these kernels vanish by (k2): if $c$ lies in $\tilde U$ but not in $\tilde X$, then $\mathrm{pr}\,c$ does not lie in $X$, since $\mathrm{pr}$ is injective on $\tilde U$ and maps $\tilde X$ onto $X$; and if $c$, $c'$ lie in $\tilde U$ and the class of $(x,x')$ is admissible for $X$, then the pair of $\mathrm{Pad}(\tilde X)$ in this class (Step~4) has its first point over $\tilde U$ and over $\mathrm{pr}\,c$, and so it is $(x,x')$.
\item[(ii)] $\mathrm{cov}_\ell$ maps generators to generators of the same family; therefore $\mathrm{cov}_\ell\left(F^{[n,m]}\right)=\left(\mathrm{cov}_\ell\,F\right)^{[n,m]}$.
\item[(iii)] $d_M(\tilde X)\le L$ if and only if $X$ is small.
Indeed, $\tilde X$ is a non-empty finite union of $M$-cubes of $\mathbb{L}_k$, connected through faces (Section~I.1.1); let $\tilde\square$ be one of its cubes.
If $d_M(\tilde X)\le L$, then $\tilde X$ lies in $\mathrm{Sm}(\tilde\square)$, and $\mathrm{pr}(\tilde X)=X$ is small by Proposition~16.3\,(f1).
If $X$ is small, it contains $\mathrm{pr}(\tilde\square)$; by Proposition~16.3\,(f1) there is a $\tilde X'\in\mathrm{Sm}(\tilde\square)$ with $\mathrm{pr}(\tilde X')=X$, and $\tilde X'$ is a lift of $X$; by Section~I.1.1, $\tilde X=\tilde X'+\ell w$ for some $w\in\mathbb{Z}^3$; since $d_M$ is invariant under translations, $d_M(\tilde X)=d_M(\tilde X')\le L$.
\end{itemize}
By (i)--(iii) and (2.f), for small $X$
\begin{align*}
\mathrm{cov}_\ell\left(\Phi^{\mathrm{nat}}(\tilde X;K,\tilde B)\right)&=E(X)-E^{[2,0]}(X)+\boldsymbol{\Theta}_{X,\mathscr{T}}\left(\mathrm{cov}_\ell\,F[K]^{[2,0]}\right)\\
&=E(X)-E^{[2,0]}(X)+\boldsymbol{\Theta}_{X,\mathscr{T}}\left(E^{[2,0]}(X)\right)=E^{\mathrm{nat}}(X;\mathscr{A},\mathscr{C}),
\end{align*}
and for large $X$
\[
\mathrm{cov}_\ell\left(\Phi^{\mathrm{nat}}(\tilde X;K,\tilde B)\right)=E(X;\mathscr{A},\mathscr{C})=E^{\mathrm{nat}}(X;\mathscr{A},\mathscr{C}).
\]
Steps 4--6 establish the definition of Section~I.1.9 for elements, with $\mathrm{rg}=\eta_k$ and the data named in Step~5.

\emph{Step 7 (the consequence).}
Let $\ell_1,\ell_2\in\mathrm{Cl}$ and $\tilde X\in\mathrm{Sm}(\tilde\square)$, and for $i=1,2$ let $B_i$ be a complex bond function on $T_k(\ell_i)$, such that $B_1(\mathrm{pr}_1b)=B_2(\mathrm{pr}_2b)=:\tilde B(b)$ for the bonds $b$ with both end points in $\tilde X$.
By Proposition~16.3\,(f1), $\mathrm{pr}_i(\tilde X)$ is a small polymer of $T_k(\ell_i)$, and $\tilde X$ is its lift.
By (2.f), for $x$ over $\tilde X$,
\[
\boldsymbol{\Theta}_{\mathrm{pr}_i\tilde X,B_i}\,\psi(x,\beta,\omega)=\mathrm{cov}_{\ell_i}\left(\boldsymbol{\Theta}_{\tilde X,\tilde B}\,\psi(x,\beta,\omega)\right),\qquad i=1,2.
\]
The element $\boldsymbol{\Theta}_{\tilde X,\tilde B}\,\psi(x,\beta,\omega)$, the right side of (1.a), is one linear combination of the generators $\psi(\tilde x,\beta,\omega)$, $\tilde x$ over $\tilde X$, and $\chi(p,p',\beta,\omega)$, $p$, $p'$ over $\tilde X$, of $\mathrm{Grc}_k(\tilde X)$, and its coefficients contain neither $\ell_1$ nor $\ell_2$.
The map $\mathrm{cov}_{\ell_i}$ replaces each generator by the generator of the torus of side $\ell_i$ with the same index.
For $\omega=0$ this is the assertion for $\psi$, for $\omega=1$ the assertion for $\bar\psi$.
\end{proof}

\begin{named}{Remark}[Proposition~17.2 at $k=0$]
At $k=0$ the integral $\int_X\mathrm{d}\xi$ is the sum over the $\xi=(x,\beta,\omega)$ with $x$ over $\tilde X$, each with the weight $\frac{1}{48}$; an index $\boldsymbol{\zeta}$ has the weight $\left(\frac{1}{48}\right)^2$; the factor in $F[K]$ is $\left(\frac{1}{48}\right)^{n+2m}$; $\mathrm{cnt}(X)=48M^3\cdot|X|_M$; and in (f) one has $\mathrm{rg}=\eta_0=1$, so that $\tilde U$ is the $1$-neighbourhood of $\tilde X$.
At a pair inside a cluster the hypothesis of (a) reads $\mathrm{ev}\,\chi(x,x')=\mathrm{ev}\,\psi(x')-\mathrm{ev}\,\psi(x)$, with $d_X=1$ and no phase.
At $k=0$ the weights $\frac{1}{48}$ and $\left(\frac{1}{48}\right)^2$ of the indices enter the bounds of (d) only inside the kernel norms.
Indeed, the numerical factors in (2.c) and (2.e) are products of the numbers $2$, $\mathrm{e}$ and $3c_X\cdot L\cdot M$; and the weights of the indices enter only through the means in (2.c), whose weights add up to $1$, and through $\mathrm{sE}=\frac12\left\|E_{20}(X)\right\|$ in (2.d), that is, through the kernel norms on the two sides of the bounds of (d).
In the flow, the family of level $0$ is $E=0$ (I.L0.21), and then $E^{\mathrm{nat}}=0$ by linearity.
\end{named}

\begin{named}{Remark}[The constants of Proposition~17.2]
In Proposition~17.2 the constants are $\mathrm{e}^2$, $3c_X\cdot L\cdot M$, $\frac{1}{30}$ and $10$; like those of Proposition~17.1 they are built from $c_X$, $L$ and $M$ only, and none of them varies with the level, the side, $e_k$, $(\mathscr{A},\mathscr{C})$, $X$ or $E$.
The data $(k,\ell,e_k)$ enter the argument as in Proposition~17.1, and in addition through the relation $h_{k,2}=e_k^{\varepsilon}\cdot h_{k,1}$.
In (f) the maps $\Phi^{\boldsymbol{\Theta}}$ and $\Phi^{\mathrm{nat}}$ are the same for all sides $\ell\in\mathrm{Cl}$.
\end{named}

\subsection*{3. Conditions on the constants}

The proofs of Propositions~17.1 and~17.2 rest on the following conditions on the constants.
For each of them the second column names the places in Section~I.8 and in Proposition~14.2 where it is imposed, and the third column the clauses in which it is used; where a consequence of the condition suffices, that consequence is stated in parentheses.

\medskip
\noindent
\begin{tabular}{p{0.33\textwidth}p{0.35\textwidth}p{0.22\textwidth}}
\hline
condition & imposed in & used in\\
\hline
$e_0^{\varepsilon}\le\frac14$ & (C8) of Section~I.8.2; (A8) of Section~I.8.4; Proposition~14.2\,(b), (A8), first inequality & Proposition~17.1\,(a)\\
$0<\varepsilon\le\bar\varepsilon=\frac{1}{200}$ (it suffices that $0<\varepsilon\le\frac14$) & (C7), (C0), (C1) of Section~I.8.2; Proposition~14.2\,(b), (S1). Here $0<\varepsilon$ follows from the inequality $1-\alpha<\varepsilon$ of (C7) and (S1) together with $\alpha<1$ of (C1) & Proposition~17.1\,(a)\\
$90c_X\cdot L\cdot M\cdot e_0^{\varepsilon}\le1$ (it suffices that $3c_X\cdot L\cdot M\cdot e_k^{\varepsilon}\le\frac{1}{30}$) & (C8) of Section~I.8.2; (A9) of Section~I.8.4; Proposition~14.2\,(b), (A9), second inequality & Proposition~17.1\,(c2); Proposition~17.2\,(d)\\
$0<\alpha<1$ & (C1) of Section~I.8.2 & Proposition~17.1\,(a), through $\varepsilon>0$, and (c3)\\
$c_X=54$ (it suffices that $c_X$ is the constant of \cite[Theorem~13.2\,(t1)]{QEDI} and that $c_X+2\le2c_X\cdot L$) & (C0) of Section~I.8.2; Proposition~14.2\,(a) & Proposition~17.1\,(c), (c2)\\
$L$ odd, $L\ge5$ & (C3) of Section~I.8.2; Proposition~14.2\,(b), (L5) & Sections~0.2, 0.3; Proposition~17.1\,(c2)\\
$M$ a power of $L$ (it suffices that $M\mathbb{Z}^3\subset\mathbb{L}_k$, that $M\cdot L^k$ is odd and that $M\ge1$) & (C5) of Section~I.8.2; Proposition~14.2\,(b), (M1) & Sections~0.2, 0.3; Proposition~17.1\,(c1), (c3)\\
\hline
\end{tabular}
\medskip

The side $\ell$ ranges over $\mathrm{Cl}$; it enters the proofs only through the facts that $\ell$ is a multiple of $M$ and that $\ell\ge13L\cdot M$ (Section~0.3, (0.b)).
By Section~I.8.2 the constant $e_0$ is chosen after $\varepsilon$, $c_X$, $L$ and $M$, and each condition of the table bounds a constant only in terms of constants chosen before it.

\section*{Chapter 18b. Contraction of the map $L\,R$; bilinears inside a cluster: proofs of Propositions 18.5, 18.6}

\subsection*{0. Conventions of this chapter}

\subsubsection*{0.1. Setting}

Throughout the chapter the data are those of the setting (G2.S-2) of Section~G2.0.4.
The level $k$ is formed with $(k,\ell,e_k)$, the level $k+1$ with $(k+1,\ell/L,e_{k+1})$, and $\eta_{k+1}=\eta_k/L$.
The constants are those of Section~I.8.

\paragraph{Weights of the indices.}
By Section~I.3.6 and (I.1.3), at the level $k$ the indices carry the weights
\[
\nu(\xi)=\frac{\eta_k^3}{48}\quad\text{for }\xi=(x,\beta,\omega)=(c,g,\beta,\omega),
\qquad
\nu(\boldsymbol{\zeta})=\left(\frac{\eta_k^3}{48}\right)^2\quad\text{for }\boldsymbol{\zeta}=(x,x',\beta,\omega),
\]
and one has
\[
\int\mathrm{d}x=\eta_k^3\cdot\frac{1}{48}\sum_{c}\sum_{g\in O_{\mathrm{cub}}},
\qquad
\int\mathrm{d}\xi=\sum_{\xi}\nu(\xi),
\qquad
\int\mathrm{d}\boldsymbol{\zeta}=\sum_{\boldsymbol{\zeta}}\nu(\boldsymbol{\zeta}).
\]
At the level $k+1$ the same formulas hold with $\eta_{k+1}$ in place of $\eta_k$.

\paragraph{Norms.}
By (I.5.1), Section~I.5.1, and (G2.r7), (G2.0.1), which are clauses of Proposition~15\,(n),
one has for $F\in\mathrm{Gr}[\psi,\chi]$ and $\mathsf{w}=(\mathsf{w}_1,\mathsf{w}_2)$
\begin{gather*}
\|F\|_{\mathsf{w}}=\sum_{n,m}\left\|F^{[n,m]}\right\|_{\mathsf{w}},
\qquad
\left\|F^{[n,m]}\right\|_{\mathsf{w}}=\frac{\mathsf{w}_1^n\cdot\mathsf{w}_2^m}{n!\,m!}\cdot\|F_{nm}\|,
\\
\|F_{nm}\|=\int\mathrm{d}\xi\,\mathrm{d}\boldsymbol{\zeta}\,\left|F_{nm}(\xi_1,\dots,\xi_n;\boldsymbol{\zeta}_1,\dots,\boldsymbol{\zeta}_m)\right|
=\sum\left(\frac{\eta_k^3}{48}\right)^{n+2m}\left|F_{nm}(\xi_1,\dots,\xi_n;\boldsymbol{\zeta}_1,\dots,\boldsymbol{\zeta}_m)\right|,
\end{gather*}
where the last sum extends over all arguments and $F_{nm}$ are the antisymmetric kernels of $F$ (Proposition~15\,(a)).
The weights of the norm at the level $k$ are
\[
h_k=\left(e_k^{-\frac{1}{4}},\,e_k^{-\frac{1}{4}+\varepsilon}\right).
\]

\paragraph{Points of the covering lattices.}
By (I.AP), the generators and the kernels are indexed by the points of $\mathbb{L}^{\mathrm{f}}_k$ and of $\mathrm{Pzt}_k$, with
\[
\psi(x+\ell\cdot e_{\mu},\,\cdot\,)=-\psi(x,\,\cdot\,),
\qquad
\chi(x+\ell\cdot e_{\mu},x'+\ell\cdot e_{\mu},\,\cdot\,)=-\chi(x,x',\,\cdot\,),
\]
and kernels are antiperiodic in each argument.

\paragraph{Labels.}
The lines (C0)--(C7) are those of the ordered list of Section~I.8.2.
The conditions (L1), (L3), (L5), (M1), (S1), (A8), (A9) are rows of Proposition~14.2\,(b).
The labels (u2), (u4), (3.2), (3.3) stand in Proposition~17.1, and so do (c1)--(c3) wherever they are cited together with Proposition~17.1.
The conditions (n1)--(n3) are those in the data of Proposition~18.6.
The clauses (k1)--(k7) are those of the definition of the spaces of activities in Section~I.5.4.
The marks (t0), (t5), (t8), (t10) denote assertions of \cite[Theorem~13.2]{QEDI}; they are recalled in 0.3.

\subsubsection*{0.2. Notation}

In this chapter, for integers $n,m\ge0$,
\[
\mathrm{pw}(n,m):=n+(1+\alpha)\cdot m+\frac{n+m}{8}-\frac{m\cdot\varepsilon}{2},
\]
and
\[
\mathrm{rad}_t:=L^{\frac{7}{4}},
\qquad
\mathrm{Dk}:=\left\{t\in\mathbb{C}:\ |t|\le\mathrm{rad}_t\right\}.
\]
For the family $E$ of Section~2 and a polymer $X$ put
\[
\mathrm{Nm}(X):=\sup_{(\mathscr{A}',\mathscr{C}')\in\tilde{R}_k\times\mathfrak{X}_k}\left\|E(X;\mathscr{A}',\mathscr{C}')\right\|_{h_k}.
\]
In Section~3, $\mathrm{hp}(x)$ and $\bar{\mathrm{hp}}(x)$ denote the parts of $\boldsymbol{\Theta}\psi(x)$ and of $\boldsymbol{\Theta}\bar{\psi}(x)$ that contain the H\"older legs,
and, for the function $\mathrm{Ng}$ of Proposition~18.6,
\[
\mathrm{Nr}(g):=\frac{1}{48}\sum_{g''\in O_{\mathrm{cub}}}\mathrm{Ng}(g'',g).
\]
The symbols $\tilde{X}$, $\tilde{U}$, $\tilde{B}$, $\mathrm{Grc}_k$, $\mathrm{cov}_{\ell}$, $T_v$, $\mathrm{rg}$ have the meaning of Section~I.1.9,
with the two conventions on $\mathrm{Grc}_k(\tilde{U})$ and on $\mathrm{cov}_{\ell}$ that are made before Lemma~3.
As in Section~I.1.1, $\tilde{c}$ and $\tilde{c}'$ denote lifts of sites $c$ and $c'$.

\subsubsection*{0.3. Polymers and field domains}

The following assertions of \cite[Theorem~13.2]{QEDI} concern sets of photon sites only.
They hold for $k\ge0$ and $\ell\in\mathrm{Cl}$; the assertions (t5) and (t8) require in addition $\ell/L\in\mathrm{Cl}$.
The marks (t0), (t5), (t8), (t10) are those of the corresponding assertions of \cite[Proposition~5.58]{QEDI}; the paragraphs of \cite[Theorem~13.2]{QEDI} carry the same marks.
\begin{itemize}
\item[(t0)] $d_M(\square)=0$ for every $M$-cube $\square$.
\item[(t5)] For every $X\in D_k$ there is exactly one $Y\in D_{k+1}$ with $\bar{X}^{L}=LY$, and this $Y$ satisfies $L\cdot d_M(Y)\le d_M(X)$.
\item[(t8)] For every $Y\in D_{k+1}$ and every number $\mathrm{kd}\ge\kappa_0$,
\[
\sum_{X\in D_k:\ \bar{X}^{L}=LY}\exp\left(-\mathrm{kd}\cdot d_M(X)\right)
\le c_{\mathrm{rb}}\cdot L^3\cdot\exp\left(-L\cdot\left(\mathrm{kd}-\kappa_0-1\right)\cdot d_M(Y)\right).
\]
In \cite[Theorem~13.2]{QEDI} the number $\mathrm{kd}$ is called $\kappa'$; here $\kappa'$ denotes a constant of the line (C4).
\end{itemize}
The assertion (t5) rests on an elementary property of cubes: two $M$-cubes with a common face lie in one $LM$-cube or in two $LM$-cubes with a common face.
Indeed, $L$ is odd, so that the $LM$-cube with the centre $L\cdot M\cdot j$, $j\in\mathbb{Z}^3$, is the union of the $L^3$ $M$-cubes with the centres $M\cdot(Lj+i)$, $i\in\mathbb{Z}^3$, $|i|\le\frac{L-1}{2}$.
Two $M$-cubes with a common face have centres $M\cdot(Lj+i)$ and $M\cdot(Lj+i)+M\cdot e_{\mu}$, or $M\cdot(Lj+i)$ and $M\cdot(Lj+i)-M\cdot e_{\mu}$.
If $|i_{\mu}\pm1|\le\frac{L-1}{2}$, the second centre belongs to the same $j$.
Otherwise $i_{\mu}=\pm\frac{L-1}{2}$, and
\[
M\cdot(Lj+i)\pm M\cdot e_{\mu}=M\cdot\left(L\cdot(j\pm e_{\mu})+i'\right),
\qquad
i'_{\mu}=\mp\frac{L-1}{2},\quad i'_{\nu}=i_{\nu}\ \ (\nu\ne\mu);
\]
then the two $LM$-cubes have the centres $L\cdot M\cdot j$ and $L\cdot M\cdot(j\pm e_{\mu})$, and they have a common face.
(On the torus the centres are taken modulo $\ell\mathbb{Z}^3$; $\ell$ is a multiple of $L\cdot M$.)
Consequently, for $X\in D_k$, a chain of $M$-cubes of $X$ in which consecutive cubes have a common face is mapped to a chain of $LM$-cubes of $\bar{X}^{L}$ in which consecutive cubes are equal or have a common face.
Since every $LM$-cube of $\bar{X}^{L}$ contains an $M$-cube of $X$, the set $\bar{X}^{L}$ is connected through faces as a union of $LM$-cubes.
\begin{itemize}
\item[(t10)] Let $\mathrm{pr}$ be the projection of $\mathbb{R}^3$ onto $\mathbb{R}^3/\ell\mathbb{Z}^3$.
For a finite union $\tilde{X}$ of $M$-cubes of $\mathbb{R}^3$ (centred at the points of $M\mathbb{Z}^3$) that is connected through faces, $d_M(\tilde{X})$ is defined as $d_M$ in \cite[Section~13.1.7]{QEDI}, with nets in $\mathbb{R}^3$.
If $X\in D_k$ satisfies $2M\cdot\left(d_M(X)+2\right)<\ell$, then $X=\mathrm{pr}(\tilde{X})$ for such a union $\tilde{X}$ on which $\mathrm{pr}$ is injective;
$\tilde{X}$ is unique up to translations by vectors of $\ell\mathbb{Z}^3$, in the sense that every union of $M$-cubes of $\mathbb{R}^3$ that is connected through faces, on which $\mathrm{pr}$ is injective and whose image is $X$, is a translate of $\tilde{X}$ by a vector of $\ell\mathbb{Z}^3$;
and $d_M(\tilde{X})=d_M(X)$.
Conversely, if $\tilde{X}$ is non-empty and $2M\cdot\left(d_M(\tilde{X})+2\right)<\ell$, then $\mathrm{pr}(\tilde{X})$ is a polymer and $d_M\left(\mathrm{pr}(\tilde{X})\right)=d_M(\tilde{X})$.
\end{itemize}
By the definition in \cite[Section~13.1.7]{QEDI}, $\bar{X}^{L}$ is the union of the $LM$-cubes that contain an $M$-cube of $X$.

The domains of Section~I.5.3 are those of \cite[Definition~5.16]{QEDI}.
By that definition and \cite[Lemma~17.3]{QEDI}, the sets $R_k$, $\mathfrak{X}_k$, $\tilde{R}_k$ are open subsets of the space of the complex bond functions on $T_k$.
Moreover $R_k\subset\tilde{R}_k$: in the definition of $\tilde{R}_k$ take $A_0=0$, with the representative $A'=0$, $\lambda=0$ on every $\square^{\dagger}$, and $A_1\in R_k$.
The bond function $0$ lies in $R_k$ and in $\mathfrak{X}_k$.
``Analytic on $\tilde{R}_k\times\mathfrak{X}_k$'' means holomorphic on this open set.
Since \cite[Lemma~17.3]{QEDI} holds at every level, for every side in $\mathrm{Cl}$ and for every charge in $(0,e_0]$, all of this holds also at the level $k+1$, for the data $(k+1,\ell/L,e_{k+1})$.

\subsection*{1. Four lemmas}

\subsubsection*{Lemma 1 (numbers)}

\begin{named}{Lemma 1}
Let the constants be those of Section~I.8.2.
\begin{itemize}
\item[(i)] One has $L\ge2400^8$, $L^{\frac{1}{4}}\ge2$, $L^{\frac{3}{4}}\ge110>2\mathrm{e}^4$, and $\mathrm{rad}_t=L^{\frac{7}{4}}\ge2$, so that
\[
\frac{1}{\mathrm{rad}_t-1}\le2L^{-\frac{7}{4}},
\qquad
\frac{1}{\mathrm{rad}_t\cdot\left(\mathrm{rad}_t-1\right)}\le2L^{-\frac{7}{2}};
\]
moreover $L^3\cdot\exp(-L)\le L^{-\frac{1}{4}}$.
\item[(ii)] $\kappa-1\ge\kappa_0$, and $L\cdot\left(\kappa-\kappa_0-1\right)\ge L\cdot\left(\kappa-\kappa_0-2\right)\ge\kappa$.
\item[(iii)] $\mathrm{pw}(0,0)=0$ and $\mathrm{pw}(2,0)=\frac{9}{4}$;
\[
\mathrm{pw}(1,1)=2+\alpha+\frac{1}{4}-\frac{\varepsilon}{2}>\frac{13}{4}-\frac{3\varepsilon}{2},
\qquad
\mathrm{pw}(0,2)=2+2\alpha+\frac{1}{4}-\varepsilon>\frac{13}{4}-\frac{3\varepsilon}{2};
\]
$\mathrm{pw}(n,m)\ge n+m$ for all integers $n,m\ge0$, in particular $\mathrm{pw}(n,m)\ge4$ if $n+m\ge4$.
\item[(iv)] $h_{k+1,1}=L^{-\frac{1}{8}}\cdot h_{k,1}$ and $h_{k+1,2}=L^{-\frac{1}{8}+\frac{\varepsilon}{2}}\cdot h_{k,2}$;
$L^{-1}\cdot h_{k+1,1}\le h_{k,1}$ and $L^{-1-\alpha}\cdot h_{k+1,2}\le h_{k,2}$;
$w^{\mathrm{E}}_{k+1}=L^{\frac{1}{8}-\frac{7\varepsilon}{2}}\cdot w^{\mathrm{E}}_k$.
\end{itemize}
\end{named}

\begin{proof}
(i) The condition (L1) of Proposition~14.2\,(b), in the line (C3), reads
\[
L^{\frac{1}{8}-\frac{3\bar{\varepsilon}}{2}}\ge2400,
\qquad
\bar{\varepsilon}=\frac{1}{200}.
\]
Since $L\ge1$ and $\frac{1}{8}-\frac{3}{400}\le\frac{1}{8}$, it follows that $L^{\frac{1}{8}}\ge2400$.
Hence $L\ge2400^8$, $L^{\frac{1}{4}}\ge2400^2$, $L^{\frac{3}{4}}\ge2400^6$ and $L^{\frac{7}{4}}\ge2$.
Euler's number satisfies $\mathrm{e}=\sum_{j\ge0}\frac{1}{j!}<2.72$, because
\[
\sum_{j\le5}\frac{1}{j!}=\frac{163}{60}<2.7167,
\qquad
\sum_{j\ge6}\frac{1}{j!}\le\frac{1}{720}\sum_{i\ge0}7^{-i}=\frac{7}{4320}<0.002.
\]
As $2.72^2=7.3984$ and $7.3984^2<54.74$, one obtains $2\mathrm{e}^4<109.48<110$.
Every number $\mathrm{rad}\ge2$ satisfies $\mathrm{rad}-1\ge\mathrm{rad}/2$;
for $\mathrm{rad}=\mathrm{rad}_t=L^{\frac{7}{4}}$ this gives the two bounds on $\frac{1}{\mathrm{rad}_t-1}$ and on $\frac{1}{\mathrm{rad}_t\cdot\left(\mathrm{rad}_t-1\right)}$.
Finally, $\exp(L)\ge\frac{L^4}{4!}=\frac{L^4}{24}$ and $L^{\frac{3}{4}}\ge24$;
so $\exp(L)\ge L^{\frac{13}{4}}$, that is $L^3\cdot\exp(-L)\le L^{-\frac{1}{4}}$.

(ii) By the line (C2) of Section~I.8.2 (Proposition~14.2\,(a), row $\kappa$), $\kappa\ge7\kappa_0+7$;
by the line (C0) of Section~I.8.2 (Proposition~14.2\,(a), row $\kappa_0$), $\kappa_0=54\log72>0$.
Hence
\[
\kappa-1-\kappa_0\ge\left(7\kappa_0+7\right)-1-\kappa_0=6\kappa_0+6>0,
\]
that is, $\kappa-1\ge\kappa_0$.
The condition (L3) of Proposition~14.2\,(b), in the line (C3), reads $L\cdot\left(\kappa-6\kappa_0-4\right)-1\ge\kappa$.
Since $\kappa_0>0$ also gives $5\kappa_0+2\ge0$, one has
\[
\kappa-\kappa_0-2=\left(\kappa-6\kappa_0-4\right)+\left(5\kappa_0+2\right)\ge\kappa-6\kappa_0-4,
\]
and it follows that $L\cdot\left(\kappa-\kappa_0-2\right)\ge L\cdot\left(\kappa-6\kappa_0-4\right)\ge\kappa+1$.
The remaining inequality $L\cdot\left(\kappa-\kappa_0-1\right)\ge L\cdot\left(\kappa-\kappa_0-2\right)$ of the chain in (ii) holds because $L>0$.

(iii) By (C1) and by the condition (S1) of Proposition~14.2\,(b), in the line (C7),
\[
1-\varepsilon<\alpha<1,
\qquad
0<\varepsilon\le\frac{1}{200}.
\]
The four values are obtained by inserting the arguments into the definition of $\mathrm{pw}$.
Further,
\[
\mathrm{pw}(1,1)-\left(\frac{13}{4}-\frac{3\varepsilon}{2}\right)=\alpha-1+\varepsilon>0,
\qquad
\mathrm{pw}(0,2)-\left(\frac{13}{4}-\frac{3\varepsilon}{2}\right)=2\alpha-1+\frac{\varepsilon}{2}>1-\frac{3\varepsilon}{2}>0,
\]
and
\[
\mathrm{pw}(n,m)-(n+m)=m\cdot\left(\alpha-\frac{\varepsilon}{2}\right)+\frac{n+m}{8}\ge0,
\]
because $\alpha>1-\varepsilon>\frac{\varepsilon}{2}$.

(iv) By Sections~I.5.1 and I.5.4, $h_k=\left(e_k^{-\frac{1}{4}},e_k^{-\frac{1}{4}+\varepsilon}\right)$ and $w^{\mathrm{E}}_k=e_k^{\frac{1}{4}-7\varepsilon}$;
together with $e_{k+1}=L^{\frac{1}{2}}\cdot e_k$ this gives the three equalities.
The two inequalities hold because the exponents $-1-\frac{1}{8}$ and $-1-\alpha-\frac{1}{8}+\frac{\varepsilon}{2}$ are negative.
\end{proof}

\subsubsection*{Lemma 2 (bilinear elements and the four maps)}

Let the data be those of the setting (G2.S-1).
Let $K$ be a function on $\mathbb{L}^{\mathrm{f}}_k\times\mathbb{L}^{\mathrm{f}}_k$ with values in the complex $4\times4$ matrices that is antiperiodic in each argument:
\[
K(x+\ell\cdot e_{\mu},x')=K(x,x'+\ell\cdot e_{\mu})=-K(x,x').
\]
Put
\begin{align*}
\mathrm{Qd}[K]&:=\int\mathrm{d}x\int\mathrm{d}x'\,\bar{\psi}(x)^{T}K(x,x')\psi(x')\\
&\phantom{:}=\eta_k^6\cdot\left(\frac{1}{48}\right)^2\sum_{c,c'}\ \sum_{g,g'\in O_{\mathrm{cub}}}\ \sum_{\beta,\beta'}\bar{\psi}_{\beta}(c,g)\,K\left((c,g),(c',g')\right)_{\beta\beta'}\,\psi_{\beta'}(c',g'),
\end{align*}
where $c$ and $c'$ range over a fundamental domain.
By (I.AP) the summand is $\ell\mathbb{Z}^3$-periodic in $c$ and in $c'$.

\begin{named}{Lemma 2}
\begin{itemize}
\item[(a)] $\mathrm{Qd}[K]$ is an even, balanced element of $\mathrm{Gr}[\psi,\chi]$ of bidegree $(2,0)$.
In the matrix notation of Section~I.5.1 its kernel is $K$.
$\mathrm{Qd}[K]=0$ only if $K=0$.
Moreover
\[
\left\|\mathrm{Qd}[K]\right\|_{\mathsf{w}}=\mathsf{w}_1^2\sum_{\beta,\beta'}\int\mathrm{d}x\int\mathrm{d}x'\,\left|K(x,x')_{\beta\beta'}\right|.
\]
\item[(b)] Let $\tilde{g}=(S,a)\in\mathbb{G}_k$, and let $\Lambda$ be a real function on the sites of $T_k$, read as an $\ell\mathbb{Z}^3$-periodic function of $c\in\mathbb{L}_k$.
Then, with $x=(c,g)$ and $x'=(c',g')$,
\begin{align*}
a_{\tilde{g}}\,\mathrm{Qd}[K]&=\mathrm{Qd}[K^{\tilde{g}}],
&K^{\tilde{g}}(x,x')&:=S^{-1}K(\tilde{g}x,\tilde{g}x')S,
\tag{18b.1}\\
G_{\Lambda}\,\mathrm{Qd}[K]&=\mathrm{Qd}[K^{\Lambda}],
&K^{\Lambda}(x,x')&:=\exp\left\{i e_k\cdot\left(\Lambda(c')-\Lambda(c)\right)\right\}\cdot K(x,x'),
\tag{18b.2}\\
\mathrm{Cg}\,\mathrm{Qd}[K]&=\mathrm{Qd}[K^{\mathrm{Cg}}],
&K^{\mathrm{Cg}}(x,x')&:=\mathsf{C}^{-1}K(x',x)^{T}\mathsf{C},
\tag{18b.3}\\
\mathbf{K}\,\mathrm{Qd}[K]&=\mathrm{Qd}[K^{\mathbf{K}}],
&K^{\mathbf{K}}(x,x')&:=\varkappa^{-1}\,\overline{K(x,x')}\,\varkappa.
\tag{18b.4}
\end{align*}
\end{itemize}
\end{named}

\begin{proof}
(a) Let $x$, $x'$ range over a fundamental domain and $\beta,\beta'\in\{1,\dots,4\}$.
Up to the sign that brings the two generators into increasing order, the products $\bar{\psi}_{\beta}(x)\psi_{\beta'}(x')$ are pairwise different basis monomials:
each of them contains one barred and one unbarred generator, and it is determined by $(x,\beta)$ and $(x',\beta')$.
In particular $\mathrm{Qd}[K]$ is even, balanced and of bidegree $(2,0)$.
The coefficient of $\bar{\psi}_{\beta}(x)\psi_{\beta'}(x')$ in $\mathrm{Qd}[K]$ is $\left(\frac{\eta_k^3}{48}\right)^2\cdot K(x,x')_{\beta\beta'}$.
By (I.5.1) this gives the formula for the norm, and $\mathrm{Qd}[K]=0$ only if all these numbers vanish.
Let $E_{20}$ be the function of two indices with
\[
E_{20}\left((x,\beta,1),(x',\beta',0)\right):=K(x,x')_{\beta\beta'}=:-E_{20}\left((x',\beta',0),(x,\beta,1)\right),
\]
and $E_{20}:=0$ at two indices with equal $\omega$.
It is antisymmetric, and
\begin{align*}
\frac{1}{2}\int\mathrm{d}\xi_1\,\mathrm{d}\xi_2\,E_{20}(\xi_1,\xi_2)\,\psi(\xi_1)\psi(\xi_2)
&=\frac{1}{2}\,\mathrm{Qd}[K]-\frac{1}{2}\int\mathrm{d}x\int\mathrm{d}x'\sum_{\beta,\beta'}K(x,x')_{\beta\beta'}\,\psi_{\beta'}(x')\bar{\psi}_{\beta}(x)\\
&=\mathrm{Qd}[K].
\end{align*}
By Proposition~15\,(a), $E_{20}$ is the antisymmetric kernel of $\mathrm{Qd}[K]$; in the matrix notation of Section~I.5.1 it is the matrix kernel $K$.

(b) \emph{Proof of} (18b.1).
By (I.T-g), $a_{\tilde{g}}\bar{\psi}(x)=(S^{-1})^{T}\bar{\psi}(\tilde{g}^{-1}x)$ and $a_{\tilde{g}}\psi(x')=S\psi(\tilde{g}^{-1}x')$.
Hence
\[
a_{\tilde{g}}\left[\bar{\psi}(x)^{T}K(x,x')\psi(x')\right]
=\bar{\psi}(\tilde{g}^{-1}x)^{T}S^{-1}K(x,x')S\,\psi(\tilde{g}^{-1}x')
=f(\tilde{g}^{-1}x,\tilde{g}^{-1}x'),
\]
where
\[
f(\tilde{x},\tilde{x}'):=\bar{\psi}(\tilde{x})^{T}S^{-1}K(\tilde{g}\tilde{x},\tilde{g}\tilde{x}')S\,\psi(\tilde{x}')
=\bar{\psi}(\tilde{x})^{T}K^{\tilde{g}}(\tilde{x},\tilde{x}')\psi(\tilde{x}').
\]
By Proposition~16.1\,(a2), $\tilde{g}(\tilde{x}+\ell\cdot n)=\tilde{g}\tilde{x}+\ell\cdot rn$ for $n\in\mathbb{Z}^3$, with $r=r(S)$ diagonal with the entries $\pm1$.
Hence $K^{\tilde{g}}$ is antiperiodic in each argument, and $f$ is $\ell\mathbb{Z}^3$-periodic in $\tilde{x}$ and in $\tilde{x}'$.
Proposition~16.1\,(a2), applied to the element $\tilde{g}^{-1}$ of $\mathbb{G}_k$, gives
\[
\int\mathrm{d}x\,f'(\tilde{g}^{-1}x)=\int\mathrm{d}x\,f'(x)
\]
for every $\ell\mathbb{Z}^3$-periodic function $f'$; for a function with values in $\mathrm{Gr}[\psi,\chi]$ this identity is applied to the coefficients.
Applied first in $x$ and then in $x'$, it yields
\[
\int\mathrm{d}x\int\mathrm{d}x'\,f(\tilde{g}^{-1}x,\tilde{g}^{-1}x')=\int\mathrm{d}\tilde{x}\int\mathrm{d}\tilde{x}'\,f(\tilde{x},\tilde{x}')=\mathrm{Qd}[K^{\tilde{g}}].
\]

\emph{Proof of} (18b.2).
By (I.T-G), $G_{\Lambda}$ multiplies $\bar{\psi}(x)$ by $\exp\left(-i e_k\cdot\Lambda(c)\right)$ and $\psi(x')$ by $\exp\left(i e_k\cdot\Lambda(c')\right)$.

\emph{Proof of} (18b.3).
By (I.T-C), $\mathrm{Cg}\,\bar{\psi}(x)=-\mathsf{C}\psi(x)$ and $\mathrm{Cg}\,\psi(x')=\mathsf{C}\bar{\psi}(x')$.
Since two generators anticommute and $\left(\mathsf{C}^{T}K\mathsf{C}\right)^{T}=\mathsf{C}^{T}K^{T}\mathsf{C}$,
\begin{align*}
\mathrm{Cg}\left[\bar{\psi}(x)^{T}K(x,x')\psi(x')\right]
&=-\psi(x)^{T}\mathsf{C}^{T}K(x,x')\mathsf{C}\,\bar{\psi}(x')
=-\sum_{\beta,\beta'}\psi_{\beta}(x)\left(\mathsf{C}^{T}K(x,x')\mathsf{C}\right)_{\beta\beta'}\bar{\psi}_{\beta'}(x')\\
&=\sum_{\beta,\beta'}\bar{\psi}_{\beta'}(x')\left(\mathsf{C}^{T}K(x,x')\mathsf{C}\right)_{\beta\beta'}\psi_{\beta}(x)
=\bar{\psi}(x')^{T}\mathsf{C}^{T}K(x,x')^{T}\mathsf{C}\,\psi(x).
\end{align*}
By (I.c3), $\mathsf{C}^{T}=\mathsf{C}^{-1}$.
It remains to exchange the names $x$ and $x'$.

\emph{Proof of} (18b.4).
By (I.T-K), $\mathbf{K}$ is an antilinear ring automorphism with $\mathbf{K}\bar{\psi}=\varkappa\bar{\psi}$ and $\mathbf{K}\psi=\varkappa\psi$.
Hence
\[
\mathbf{K}\left[\bar{\psi}(x)^{T}K(x,x')\psi(x')\right]=\bar{\psi}(x)^{T}\varkappa^{T}\,\overline{K(x,x')}\,\varkappa\,\psi(x'),
\]
and $\varkappa^{T}=\varkappa^{-1}$ by (I.c4); the weights of the two integrals are real.

Finally, $K^{\Lambda}$, $K^{\mathrm{Cg}}$ and $K^{\mathbf{K}}$ are antiperiodic in each argument, because $K$ is antiperiodic and $\Lambda$ is periodic.
\end{proof}

\subsubsection*{Lemma 3 (elements and kernels in a lift)}

\paragraph{Notation.}
Let $k\ge0$.
For a finite set $\tilde{U}$ of sites of $\mathbb{L}_k$, $\mathrm{Grc}_k(\tilde{U})$ denotes in this chapter the subalgebra of the algebra $\mathrm{Grc}_k(\tilde{U})$ of Section~I.1.9
that is generated by the $\psi(x,\beta,\omega)$ with $c\in\tilde{U}$ and the $\chi(x,x',\beta,\omega)$ with $(x,x')\in\mathrm{Pzt}_k$ and $c,c'\in\tilde{U}$.
These indices are said to lie \emph{over} $\tilde{U}$; they carry the weights $\eta_k^3/48$ and $\left(\eta_k^3/48\right)^2$, respectively.
For a system $K=(K_{nm})$ of functions of $n$ indices $\xi$ and $m$ indices $\boldsymbol{\zeta}$ over $\tilde{U}$, antisymmetric in the $\xi$ and in the $\boldsymbol{\zeta}$ separately, put
\[
\mathrm{El}(K):=\sum_{n,m}\frac{1}{n!\,m!}\sum_{\xi_1,\dots,\xi_n,\,\boldsymbol{\zeta}_1,\dots,\boldsymbol{\zeta}_m}\left(\frac{\eta_k^3}{48}\right)^{n+2m}K_{nm}(\xi;\boldsymbol{\zeta})\,\psi(\xi_1)\cdots\psi(\xi_n)\,\chi(\boldsymbol{\zeta}_1)\cdots\chi(\boldsymbol{\zeta}_m)\ \in\mathrm{Grc}_k(\tilde{U}),
\]
where the inner sum extends over the indices over $\tilde{U}$.
By Proposition~15\,(a), applied to the two finite index sets over $\tilde{U}$, the map $K\mapsto\mathrm{El}(K)$ is a linear bijection onto $\mathrm{Grc}_k(\tilde{U})$,
and $K$ is the system of the antisymmetric kernels of $\mathrm{El}(K)$.

For $v\in\eta_k\mathbb{Z}^3$ put $(T_vK)(\xi;\boldsymbol{\zeta}):=K(\xi-v;\boldsymbol{\zeta}-v)$, all sites being shifted,
and let $T_v\colon\mathrm{Grc}_k(\tilde{U})\to\mathrm{Grc}_k(\tilde{U}+v)$ be the isomorphism with $\psi(x,\,\cdot\,)\mapsto\psi(x+v,\,\cdot\,)$ and $\chi(x,x',\,\cdot\,)\mapsto\chi(x+v,x'+v,\,\cdot\,)$.
Then $T_v\,\mathrm{El}(K)=\mathrm{El}(T_vK)$, as the substitution $\xi\mapsto\xi-v$, $\boldsymbol{\zeta}\mapsto\boldsymbol{\zeta}-v$ in the sum shows.

For a side $\ell\in\mathrm{Cl}$, $\mathrm{cov}_{\ell}\colon\mathrm{Grc}_k(\tilde{U})\to\mathrm{Gr}[\psi,\chi]$, the algebra on the torus of side $\ell$, is the homomorphism that maps every generator to the generator of the torus with the same index in the sense of (I.AP).
The images anticommute, so that this homomorphism exists for every finite set $\tilde{U}$.
If $\text{sup-diameter}(\tilde{U})+M<\ell$, it is the restriction of the map $\mathrm{cov}_{\ell}$ of Section~I.1.9.

For a finite union $\tilde{X}$ of $M$-cubes of $\mathbb{L}_k$ put
\[
\tilde{X}^{+\eta_k}:=\left\{c\in\mathbb{L}_k:\ \mathrm{dist}(c,\tilde{X})\le\eta_k\right\}
\qquad(\text{sup distance}).
\]
A pair $(x,x')\in\mathrm{Pzt}_k$ is \emph{admissible} for $\tilde{X}$ if $c$, $c'$ and the six paths $\boldsymbol{\Gamma}^{\pi}(c,c')$ lie in $\tilde{X}$ (kind (I); for $c=c'$ the condition is $c\in\tilde{X}$),
or if $x'=x^{(3)}$ and $c$ or $c_{3}$ lies in $\tilde{X}$ (kind (II)).
These are the notions of Section~I.1.7, formed on the covering lattice.

\begin{named}{Lemma 3}
Let $\ell\in\mathrm{Cl}$.
\begin{itemize}
\item[(a)] Let $\tilde{U}$ be a finite set of sites of $\mathbb{L}_k$ with $\text{sup-diameter}(\tilde{U})+\eta_k<\ell$.
Then $\mathrm{cov}_{\ell}$ is injective on $\mathrm{Grc}_k(\tilde{U})$ and preserves the bidegree.
For $F\in\mathrm{Grc}_k(\tilde{U})$ the antisymmetric kernels of $\mathrm{cov}_{\ell}F$, read by (I.AP) as functions of arguments on the covering lattices,
coincide with the antisymmetric kernels of $F$ at the arguments over $\tilde{U}$,
and they vanish at every system of arguments that contains an index which is not congruent modulo $\ell\mathbb{Z}^3$ to an index over $\tilde{U}$.
\item[(b)] Let $X$ be a polymer of $T_k$ with a lift $\tilde{X}$ such that $\text{sup-diameter}(\tilde{X})+2\eta_k+M<\ell$, and let $\tilde{U}:=\tilde{X}^{+\eta_k}$.
The covering map induces a bijection of the fermion sites over $\tilde{X}$ onto the fermion sites over $X$,
and a bijection of the pairs of $\mathrm{Pzt}_k$ that are admissible for $\tilde{X}$ onto the pairs of $\mathrm{Pz}_k$ that are admissible for $X$, which keeps the kind;
the pairs admissible for $\tilde{X}$ have both points over $\tilde{U}$.
Consequently, if $G\in\mathrm{Gr}[\psi,\chi]$ contains only generators $\psi(\xi)$ with $\xi$ over $X$ and $\chi(\boldsymbol{\zeta})$ with $\boldsymbol{\zeta}$ admissible for $X$,
then $G=\mathrm{cov}_{\ell}\left(\mathrm{El}(K)\right)$, where $K$ is the system of the antisymmetric kernels of $G$ restricted to the arguments over $\tilde{U}$;
and $K_{nm}(\xi;\boldsymbol{\zeta})=0$ unless all $\xi_i$ lie over $\tilde{X}$ and all $\boldsymbol{\zeta}_j$ are admissible for $\tilde{X}$.
\item[(c)] In the situation of (b), $X$ is small if and only if $d_M(\tilde{X})\le L$, where $d_M(\tilde{X})$ is formed with nets in $\mathbb{R}^3$.
\end{itemize}
\end{named}

\begin{proof}
(a) Two different points of $\tilde{U}$ that are congruent modulo $\ell\mathbb{Z}^3$ would have a sup distance at least $\ell$, which the hypothesis excludes.
So the covering map is injective on $\tilde{U}$, and there is a fundamental domain $\mathrm{Fu}$ of $\ell\mathbb{Z}^3$ in $\mathbb{L}_k$ that contains $\tilde{U}$.
By (I.AP), the $\psi(x,\beta,\omega)$ with $c\in\mathrm{Fu}$ and the $\chi(x,x',\beta,\omega)$ with $(x,x')\in\mathrm{Pzt}_k$, $c\in\mathrm{Fu}$, are independent generators of $\mathrm{Gr}[\psi,\chi]$,
and $\mathrm{cov}_{\ell}$ maps the generators of $\mathrm{Grc}_k(\tilde{U})$ to the generators of this list with the same indices.
Choose the total orders of the two index sets of the torus so that they extend the orders of the indices over $\tilde{U}$.
Then $\mathrm{cov}_{\ell}$ maps the basis monomials of $\mathrm{Grc}_k(\tilde{U})$ to the basis monomials with the same sets of indices.
Hence $\mathrm{cov}_{\ell}$ is injective and preserves the bidegree, and
\[
c_{\mathrm{cov}_{\ell}F}(I_1,I_2)=
\begin{cases}
c_F(I_1,I_2)&\text{if }I_1\text{ and }I_2\text{ consist of indices over }\tilde{U},\\
0&\text{otherwise.}
\end{cases}
\]
By Proposition~15\,(a), in both algebras and with the same weights,
\[
K_{nm}(I_1;I_2)\cdot\nu(I_1)\cdot\nu(I_2)=c(I_1,I_2)
\]
for increasing arguments.
So the two kernels coincide at increasing arguments over $\tilde{U}$, and by antisymmetry at all arguments over $\tilde{U}$;
and the kernel of $\mathrm{cov}_{\ell}F$ vanishes if one of its arguments, taken with its first point over $\mathrm{Fu}$, is not an index over $\tilde{U}$.
An index with its first point over $\mathrm{Fu}$ that is congruent to an index over $\tilde{U}$ coincides with this index, because $\tilde{U}\subset\mathrm{Fu}$.
Since the kernels of the torus are antiperiodic in each argument, (a) follows.

(b) One has $\text{sup-diameter}(\tilde{U})\le\text{sup-diameter}(\tilde{X})+2\eta_k$, and $2\eta_k+\eta_k\le2\eta_k+M$ because $M\ge\eta_k$.
Hence $\text{sup-diameter}(\tilde{U})+\eta_k<\ell$; the covering map is injective on $\tilde{U}$, and it maps $\tilde{X}$ onto $X$.
This gives the bijection of the fermion sites.

\emph{Pairs of kind} (I).
Let $c$, $c'$ lie in $X$ with $d(c,c')\le M$, and let $\tilde{c}$, $\tilde{c}'$ be their lifts in $\tilde{X}$.
By Section~I.1.1 (lifts, with $\mathrm{sd}=M$; note that $\text{sup-diameter}(\tilde{X})+M<\ell$), $\tilde{c}'-\tilde{c}$ is the only representative of $c'-c$ of sup norm at most $M$.
So $\left((\tilde{c},g),(\tilde{c}',g')\right)$ is the shortest lift of $\left((c,g),(c',g')\right)$, and the paths $\boldsymbol{\Gamma}^{\pi}(c,c')$ of the torus are the images of the paths $\boldsymbol{\Gamma}^{\pi}(\tilde{c},\tilde{c}')$.
A site $\mathrm{cs}$ of $\boldsymbol{\Gamma}^{\pi}(\tilde{c},\tilde{c}')$ satisfies $|\mathrm{cs}-\tilde{c}|\le M$.
If its image lies in $X$, with the lift $\mathrm{cs}'\in\tilde{X}$, then $\mathrm{cs}-\mathrm{cs}'\in\ell\mathbb{Z}^3$ and
\[
|\mathrm{cs}-\mathrm{cs}'|\le M+\text{sup-diameter}(\tilde{X})<\ell,
\]
so that $\mathrm{cs}=\mathrm{cs}'$.
Hence the six paths lie in $X$ if and only if their lifts lie in $\tilde{X}$.

\emph{Pairs of kind} (II).
If $c$ lies in $X$, the pair $\left((\tilde{c},g),(\tilde{c},g)^{(3)}\right)$ has the photon sites $\tilde{c}\in\tilde{X}$ and $\tilde{c}+\eta_k\cdot ge_1\in\tilde{U}$.
If $c_{3}$ lies in $X$, with the lift $\mathrm{cs}\in\tilde{X}$, the pair $\left((\mathrm{cs}-\eta_k\cdot ge_1,g),(\mathrm{cs},gs_3)\right)$ has its photon sites in $\tilde{U}$.
A site of $\tilde{U}$ whose image lies in $X$ is the lift of this image, by the injectivity on $\tilde{U}$;
so $c$, resp.\ $c_{3}$, lies in $X$ if and only if the lifted site lies in $\tilde{X}$.

Two pairs over $\tilde{U}$ that are congruent modulo $\ell\mathbb{Z}^3$ coincide.
This gives the bijection of the pairs, and it keeps the kind.
Therefore every generator that occurs in $G$ is, up to a sign, the image under $\mathrm{cov}_{\ell}$ of a generator of $\mathrm{Grc}_k(\tilde{U})$ whose index lies over $\tilde{X}$, resp.\ is admissible for $\tilde{X}$.
Thus $G=\mathrm{cov}_{\ell}F$ with an element $F\in\mathrm{Grc}_k(\tilde{U})$ that contains only such generators.
By (a), the kernels of $F$ are those of $G$ at the arguments over $\tilde{U}$, so that $F=\mathrm{El}(K)$,
and $K_{nm}(\xi;\boldsymbol{\zeta})$ vanishes unless all $\xi_i$ lie over $\tilde{X}$ and all $\boldsymbol{\zeta}_j$ are admissible for $\tilde{X}$.

(c) By Section~I.1.1, $\tilde{X}$ is a union of $M$-cubes of $\mathbb{L}_k$ with the same common faces as $X$;
so it is finite, non-empty and connected through faces, and $\mathrm{pr}(\tilde{X})=X$.
The number $M/\eta_k=M\cdot L^k$ is an odd integer, since $L$ is odd and $M$ is a power of $L$ (Section~I.1.1, faces).
Hence the sites of the $M$-cube with the centre $z\in M\mathbb{Z}^3$ are the points
\[
z+\eta_k\cdot j,
\qquad
j\in\mathbb{Z}^3,\quad|j|\le\frac{1}{2}\left(\frac{M}{\eta_k}-1\right);
\]
there are $\left(M/\eta_k\right)^3$ of them, and the extreme ones have the distance $\frac{M}{2}-\frac{\eta_k}{2}$ from $z$ in the direction concerned.
So a closed $M$-cube $\left\{y\in\mathbb{R}^3:\ |y-z|\le\frac{M}{2}\right\}$ of $\mathbb{R}^3$ reaches $\frac{\eta_k}{2}$ beyond its extreme sites in each direction,
and the union of the closed cubes of $\tilde{X}$ has the sup-diameter $\text{sup-diameter}(\tilde{X})+\eta_k<\ell$.
Consequently $\mathrm{pr}$ is injective on this union, since two different points that are congruent modulo $\ell\mathbb{Z}^3$ have a sup distance at least $\ell$.
Moreover, as $\ell\in\mathrm{Cl}$,
\[
2M\cdot(L+2)=2L\cdot M+4M\le6L\cdot M<13L\cdot M\le\ell.
\]
If $X$ is small, then $2M\cdot\left(d_M(X)+2\right)<\ell$.
By the first part of (t10) and its clause of uniqueness, $\tilde{X}$ is the translate by a vector of $\ell\mathbb{Z}^3$ of the union furnished by (t10),
and $d_M(\tilde{X})=d_M(X)\le L$, because a translation maps nets to nets of the same length.
If $d_M(\tilde{X})\le L$, then $2M\cdot\left(d_M(\tilde{X})+2\right)<\ell$,
and the converse part of (t10) gives $d_M(X)=d_M\left(\mathrm{pr}(\tilde{X})\right)=d_M(\tilde{X})\le L$.
\end{proof}

\subsubsection*{Lemma 4 (spin matrices and labels)}

For $v\in\mathbb{C}^3$ put
\[
(\mathbf{1}_2\otimes\sigma)\cdot v:=\sum_{\mu}v_\mu\,(\mathbf{1}_2\otimes\sigma_\mu).
\]

\begin{named}{Lemma 4}
Let $S\in\mathrm{Sp}$, let $r=r(S)=\operatorname{diag}(r_1,r_2,r_3)$, and let $v\in\mathbb{C}^3$.
\begin{itemize}
\item[(sp1)] $S\,(\gamma\cdot v)\,S^{-1}=\gamma\cdot(rv)$.
\item[(sp2)] $S\,\Sigma_3\,S^{-1}=S^{-1}\,\Sigma_3\,S=\det(r)\,\Sigma_3$.
\item[(sp3)] $S\left((\mathbf{1}_2\otimes\sigma)\cdot v\right)S^{-1}=\det(r)\,(\mathbf{1}_2\otimes\sigma)\cdot(rv)$.
\item[(sp4)] $\mathsf{C}^{-1}(\gamma\cdot v)^{T}\mathsf{C}=-\gamma\cdot v$,
\quad $\mathsf{C}^{-1}\Sigma_3^{T}\mathsf{C}=\Sigma_3$,
\quad $\mathsf{C}^{-1}\left((\mathbf{1}_2\otimes\sigma)\cdot v\right)^{T}\mathsf{C}=-(\mathbf{1}_2\otimes\sigma)\cdot v$.
\item[(sp5)] If $v$ is real, then
$\varkappa^{-1}\,\overline{\gamma\cdot v}\,\varkappa=\gamma\cdot v$,
\quad $\varkappa^{-1}\,\overline{\Sigma_3}\,\varkappa=-\Sigma_3$,
\quad $\varkappa^{-1}\,\overline{(\mathbf{1}_2\otimes\sigma)\cdot v}\,\varkappa=-(\mathbf{1}_2\otimes\sigma)\cdot v$.
\item[(sp6)] $\operatorname{tr}\gamma_\mu=\operatorname{tr}\Sigma_3=\operatorname{tr}(\mathbf{1}_2\otimes\sigma_\mu)=0$.
\item[(sp7)] The moduli of the $16$ entries of $\Sigma_3$ add up to $4$.
If $v$ is real, the moduli of the $16$ entries of $\gamma\cdot v$ add up to
\[
4\,|v_3|+4\left(v_1^2+v_2^2\right)^{\frac{1}{2}}\le4\sqrt{2}\cdot|v|_2 .
\]
\item[(sp8)] Let $r\in\{\pm1\}^3$ and $i\in\{0,1,2\}$.
The maps $g\mapsto rg$ and $g\mapsto gs_i$ are bijections of $O_{\mathrm{cub}}$.
For every $g\in O_{\mathrm{cub}}$ one has $(gs_i)s_i=g$ and $gs_i\ne g$, and the three elements $gs_0$, $gs_1$, $gs_2$ are pairwise distinct;
moreover $\det(rg)=\det(r)\det(g)$, $\det(gs_i)=-\det(g)$, $(gs_i)\hat n_i=-g\hat n_i$, and $(gs_i)v=gv$ for every $v\in\mathbb{R}^3$ with $v\cdot\hat n_i=0$.
\end{itemize}
\end{named}

\begin{proof}
(sp1) is (I.4.1).
For $v=e_\mu$ it reads $S\gamma_\mu S^{-1}=r_\mu\gamma_\mu$.

(sp2).
By the relation (c2) of Section~I.2.1, $i\Sigma_3=\gamma_1\gamma_2\gamma_3$.
Hence, by (sp1), $S\,(i\Sigma_3)\,S^{-1}=r_1r_2r_3\,\gamma_1\gamma_2\gamma_3$, that is, $S\Sigma_3S^{-1}=\det(r)\,\Sigma_3$.
Since $\det(r)=\pm1$, it follows that $S^{-1}\Sigma_3S=\det(r)^{-1}\Sigma_3=\det(r)\,\Sigma_3$.

(sp3).
Let $(\mu,\nu,\mathfrak{l})$ be a cyclic permutation of $(1,2,3)$.
Since $\gamma_\mu=\sigma_3\otimes\sigma_\mu$ and $\sigma_3^2=\mathbf{1}_2$,
\[
\gamma_\nu\gamma_{\mathfrak{l}}=(\sigma_3\otimes\sigma_\nu)(\sigma_3\otimes\sigma_{\mathfrak{l}})=\mathbf{1}_2\otimes\sigma_\nu\sigma_{\mathfrak{l}}=i\,(\mathbf{1}_2\otimes\sigma_\mu).
\]
Hence $\mathbf{1}_2\otimes\sigma_\mu=-i\gamma_\nu\gamma_{\mathfrak{l}}$, and by (sp1)
\[
S\,(\mathbf{1}_2\otimes\sigma_\mu)\,S^{-1}=r_\nu r_{\mathfrak{l}}\,(\mathbf{1}_2\otimes\sigma_\mu)=\det(r)\,r_\mu\,(\mathbf{1}_2\otimes\sigma_\mu),
\]
because $r_\mu^2=1$.
Multiplying by $v_\mu$ and summing over $\mu$ gives (sp3).

(sp4).
By the relation (c3) of Section~I.2.1, $-\gamma_\mu^{T}=\mathsf{C}^{-1}\gamma_\mu\mathsf{C}$ and $\mathsf{C}^{T}=\mathsf{C}^{-1}$.
Transposing the first of these identities gives
\[
-\gamma_\mu=\mathsf{C}^{T}\gamma_\mu^{T}\left(\mathsf{C}^{-1}\right)^{T}=\mathsf{C}^{-1}\gamma_\mu^{T}\mathsf{C},
\]
and by linearity this is the first relation of (sp4).
The matrices $\gamma_1$, $\gamma_2$, $\gamma_3$ anticommute with each other (relation (c1) of Section~I.2.1), and reversing the order of three of them takes three transpositions; therefore
\[
\mathsf{C}^{-1}(\gamma_1\gamma_2\gamma_3)^{T}\mathsf{C}=(-\gamma_3)(-\gamma_2)(-\gamma_1)=-\gamma_3\gamma_2\gamma_1=\gamma_1\gamma_2\gamma_3,
\]
and moreover
\[
\mathsf{C}^{-1}(\gamma_\nu\gamma_{\mathfrak{l}})^{T}\mathsf{C}=(-\gamma_{\mathfrak{l}})(-\gamma_\nu)=\gamma_{\mathfrak{l}}\gamma_\nu=-\gamma_\nu\gamma_{\mathfrak{l}}.
\]
Since $i\Sigma_3=\gamma_1\gamma_2\gamma_3$ and $\mathbf{1}_2\otimes\sigma_\mu=-i\gamma_\nu\gamma_{\mathfrak{l}}$, these two identities give the second and the third relation of (sp4).

(sp5).
By the relation (c4) of Section~I.2.1, $\overline{\gamma_\mu}=\varkappa^{-1}\gamma_\mu\varkappa$ and $\varkappa^2=1$; hence $\varkappa^{-1}\,\overline{\gamma_\mu}\,\varkappa=\gamma_\mu$, which gives the first relation for real $v$.
Next,
\[
\overline{i\Sigma_3}=\overline{\gamma_1}\,\overline{\gamma_2}\,\overline{\gamma_3}=\varkappa^{-1}(i\Sigma_3)\varkappa ,
\]
while $\overline{i\Sigma_3}=-i\Sigma_3$ because $\Sigma_3$ is real; so $\varkappa^{-1}\Sigma_3\varkappa=-\Sigma_3$, which is the second relation.
Finally,
\[
\overline{\mathbf{1}_2\otimes\sigma_\mu}=\overline{-i\gamma_\nu\gamma_{\mathfrak{l}}}=i\,\varkappa^{-1}\gamma_\nu\gamma_{\mathfrak{l}}\varkappa ,
\]
so that $\varkappa^{-1}\,\overline{\mathbf{1}_2\otimes\sigma_\mu}\,\varkappa=i\gamma_\nu\gamma_{\mathfrak{l}}=-(\mathbf{1}_2\otimes\sigma_\mu)$; for real $v$ this gives the third relation.

(sp6).
One has $\gamma_\mu=\sigma_3\otimes\sigma_\mu$ and $\Sigma_3=\sigma_3\otimes\mathbf{1}_2$.
The trace of the tensor product of two $2\times2$ matrices is the product of their traces, and $\operatorname{tr}\sigma_\mu=0$ for $\mu=1,2,3$.

(sp7).
One has $\Sigma_3=\operatorname{diag}(1,1,-1,-1)$ and $\gamma\cdot v=\operatorname{diag}(\sigma\cdot v,-\sigma\cdot v)$ with
\[
\sigma\cdot v=\left(\begin{array}{cc}v_3&v_1-iv_2\\ v_1+iv_2&-v_3\end{array}\right).
\]
For real $v$, the Cauchy--Schwarz inequality gives $|v_3|+\left(v_1^2+v_2^2\right)^{\frac{1}{2}}\le\sqrt{2}\cdot|v|_2$.

(sp8).
The set $O_{\mathrm{cub}}$ is a group that contains $\{\pm1\}^3$ and $s_0$, $s_1$, $s_2$.
By Section~I.1.2, $s_i^2=1$, $s_i\ne1$, the three matrices $s_0$, $s_1$, $s_2$ are distinct, $\det s_i=-1$, $s_i\hat n_i=-\hat n_i$, and $s_iv=v-2\,(v\cdot\hat n_i)\,\hat n_i$.
These facts give all assertions of (sp8).
\end{proof}

\subsection*{2. Proposition 18.5}

\subsubsection*{Proof of Proposition 18.5}

The clauses are proved in the order (d), (b), (a), (c), (e), (f).
Throughout,
\[
E:=(R\mathscr{E})^{\mathrm{nat}}.
\]

\paragraph{Step 0 (the connections of level $k$).}
Let $(\mathscr{A},\mathscr{C})\in\tilde{R}_{k+1}\times\mathfrak{X}_{k+1}$.
By Proposition~18.1\,(c), $\mathscr{A}_L\in\tilde{R}_k$.
By Proposition~18.1\,(b), $\mathscr{C}_L\in L^{-\frac{3}{4}}\mathfrak{X}_k$, and $L^{-\frac{3}{4}}\mathfrak{X}_k\subset\mathfrak{X}_k$, since $L^{-\frac{3}{4}}\le1$ and $\mathfrak{X}_k$ is defined by two bounds on norms.
The map $(\mathscr{A},\mathscr{C})\mapsto(\mathscr{A}_L,\mathscr{C}_L)$, where $A_L(b)=L^{-\frac{1}{2}}\cdot A(b/L)$, is complex linear;
one has $(\mathscr{A}+\mathscr{C})_L=\mathscr{A}_L+\mathscr{C}_L$ and $0_L=0$.
Consequently, for every family $E'$ of level $k$ in the sense of (G2.r8), the elements $E'(X;\mathscr{A}_L,\mathscr{C}_L)$, $X\in D_k$, are defined, and
\[
(LE')(Y;\mathscr{A},\mathscr{C})=\sum_{X:\,\bar{X}^{L}=LY}S_L\left[E'(X;\mathscr{A}_L,\mathscr{C}_L)\right]
\]
is a finite sum of elements of $\mathrm{Gr}[\psi,\chi]$ of level $k+1$.

\paragraph{Step 1 (the family $E$; the first two assertions of (d)).}
\emph{Membership and linearity.}
By Proposition~16.3\,(d), $R$ is linear, and $R\mathscr{E}$ lies in $K^{x,C}_k$, and in $K^{x}_k$ if $\mathscr{E}$ does.
By Proposition~17.2\,(e), applied to the family $R\mathscr{E}$, the family $E=(R\mathscr{E})^{\mathrm{nat}}$ lies in $K^{x,C}_k$, and in $K^{x}_k$ if $R\mathscr{E}$ does;
moreover the map ${E'}\mapsto{E'}^{\mathrm{nat}}$ is linear.
Hence $E$ depends linearly on $\mathscr{E}$, $E$ consists of even, balanced elements, $E$ has the clauses (k1)--(k6) of Section~I.5.4, and $E$ has (k7) if $\mathscr{E}\in K^{x}_k$.

\emph{Bounds.}
Let $(\mathscr{A}',\mathscr{C}')\in\tilde{R}_k\times\mathfrak{X}_k$.

Let $X$ be large.
By the definitions of Section~I.5.6, $(R\mathscr{E})(X)=\mathscr{E}(X)$ and $E(X)=(R\mathscr{E})(X)$.
Hence, by (k6),
\[
\left\|E(X;\mathscr{A}',\mathscr{C}')\right\|_{h_k}=\left\|\mathscr{E}(X;\mathscr{A}',\mathscr{C}')\right\|_{h_k}\le\|\mathscr{E}\|^{x}_k\cdot\exp\left(-\kappa\cdot d_M(X)\right).
\]

Let $X$ be small.
The family $R\mathscr{E}$ is balanced and has (k2).
By Proposition~17.2\,(d),
\[
\left\|E(X;\mathscr{A}',\mathscr{C}')\right\|_{h_k}\le10\left\|(R\mathscr{E})(X;\mathscr{A}',\mathscr{C}')\right\|_{h_k},
\]
and by Proposition~16.3\,(d),
\[
\left\|(R\mathscr{E})(X;\mathscr{A}',\mathscr{C}')\right\|_{h_k}\le\left\|\mathscr{E}(X;\mathscr{A}',\mathscr{C}')\right\|_{h_k}+\left\|\mathscr{E}(X;0,0)\right\|_{h_k}.
\]
The pair $(0,0)$ lies in $\tilde{R}_k\times\mathfrak{X}_k$ (Section~0.3).
By (k6), each of the two terms on the right is at most $\|\mathscr{E}\|^{x}_k\cdot\exp\left(-\kappa\cdot d_M(X)\right)$.

Recall that $\mathrm{Nm}(X)$ is the supremum of $\left\|E(X;\mathscr{A}',\mathscr{C}')\right\|_{h_k}$ over $(\mathscr{A}',\mathscr{C}')\in\tilde{R}_k\times\mathfrak{X}_k$.
It follows that
\begin{equation}
\begin{aligned}
\mathrm{Nm}(X)&\le20\,\|\mathscr{E}\|^{x}_k\cdot\exp\left(-\kappa\cdot d_M(X)\right)&&(X\text{ small}),\\
\mathrm{Nm}(X)&\le\|\mathscr{E}\|^{x}_k\cdot\exp\left(-\kappa\cdot d_M(X)\right)&&(X\text{ large}).
\end{aligned}
\tag{18b.5}
\end{equation}

\emph{Normalisation.}
Let $X$ be small.
By Proposition~16.3\,(d),
\[
(R\mathscr{E})_{00}(X;0,0)=0,\qquad\int_X\mathrm{d}x\int_X\mathrm{d}x'\,(R\mathscr{E})_{20}(X;0,0;x,x')=0;
\]
here each of the two integrals is $\frac{\eta_k^3}{48}$ times the sum over the photon sites $c$ of $X$ and the $48$ labels $g$.
By Proposition~17.2\,(b), applied to the balanced family $R\mathscr{E}$, one has $E_{00}=(R\mathscr{E})_{00}$, so that $E_{00}(X;0,0)=0$;
and $E_{20}(X;0,0;\cdot,\cdot)=0$ as a kernel on $(X\times O_{\mathrm{cub}})^2$.

\paragraph{Step 2 (small polymers; the bound of (d) and its factors by bidegree).}
Fix a small polymer $X$ and $(\mathscr{A},\mathscr{C})\in\tilde{R}_{k+1}\times\mathfrak{X}_{k+1}$.
Since $X$ is not empty, a photon site $\tilde{c}$ of $X$ can be chosen;
let $\mathrm{bh}$, $\mathscr{C}^{b}$ be the bond functions on $T_k$ that Proposition~18.2 provides for $(X,\tilde{c},\mathscr{A},\mathscr{C})$.
By Step 1, $E$ is a family of level $k$ with (k1), (k2), (k4).
Hence Proposition~18.2 yields the following two facts.
\begin{itemize}
\item[(ii)] $(t\cdot\mathrm{bh},t\cdot\mathscr{C}^{b})\in R_k\times\mathfrak{X}_k$ for every complex number $t$ with $|t|\le L^{\frac{7}{4}}=\mathrm{rad}_t$.
\item[(iii)] For all $(n,m)$ and all arguments,
\[
\left|E_{nm}(X;\mathscr{A}_L,\mathscr{C}_L;\xi,\boldsymbol{\zeta})\right|\le\exp\left(2(n+m)\right)\cdot\left|E_{nm}(X;\mathrm{bh},\mathscr{C}^{b};\xi,\boldsymbol{\zeta})\right|,
\]
and $E_{00}(X;\mathscr{A}_L,\mathscr{C}_L)=E_{00}(X;\mathrm{bh},\mathscr{C}^{b})$.
\end{itemize}

\emph{The disk.}
Put
\[
\mathrm{Ot}:=\left\{t\in\mathbb{C}:\left(t\cdot\mathrm{bh},t\cdot\mathscr{C}^{b}\right)\in R_k\times\mathfrak{X}_k\right\}.
\]
The set $R_k\times\mathfrak{X}_k$ is open (Section~0.3; \cite[Definition~5.16 and Lemma~17.3]{QEDI}), and $t\mapsto(t\cdot\mathrm{bh},t\cdot\mathscr{C}^{b})$ is linear;
hence $\mathrm{Ot}$ is open.
By Proposition~18.2\,(ii), $\mathrm{Ot}$ contains the compact disk $\mathrm{Dk}=\{t\in\mathbb{C}:|t|\le\mathrm{rad}_t\}$;
hence it contains the disk $\{t\in\mathbb{C}:|t|<\mathrm{rad}_t+\delta\}$ for some $\delta>0$.
Since $R_k$ and $\mathfrak{X}_k$ are defined by bounds on norms, $\mathrm{Ot}=-\mathrm{Ot}$.
For $t\in\mathrm{Ot}$ the pair $(t\cdot\mathrm{bh},t\cdot\mathscr{C}^{b})$ lies in $\tilde{R}_k\times\mathfrak{X}_k$ (Section~0.3).
Therefore, by (k1), every function $t\mapsto E_{nm}(X;t\cdot\mathrm{bh},t\cdot\mathscr{C}^{b};\xi;\boldsymbol{\zeta})$ is holomorphic on $\mathrm{Ot}$;
and by (G2.r7), which is part of Proposition~15\,(n), one has for $t\in\mathrm{Ot}$ and all $n$, $m$
\begin{equation}
\begin{split}
\frac{h_{k,1}^{n}\,h_{k,2}^{m}}{n!\,m!}\cdot\left\|E_{nm}(X;t\cdot\mathrm{bh},t\cdot\mathscr{C}^{b})\right\|
&=\left\|E^{[n,m]}(X;t\cdot\mathrm{bh},t\cdot\mathscr{C}^{b})\right\|_{h_k}\\
&\le\left\|E(X;t\cdot\mathrm{bh},t\cdot\mathscr{C}^{b})\right\|_{h_k}\le\mathrm{Nm}(X).
\end{split}
\tag{18b.6}
\end{equation}

\emph{Splitting by bidegrees.}
By Proposition~18.4\,(a), $S_L$ maps a basis monomial of bidegree $(n,m)$ to $\mathrm{sn}\cdot L^{-n-(1+\alpha)m}$, with $\mathrm{sn}\in\{1,-1\}$, times a basis monomial of the same bidegree;
hence $(S_LF)^{[n,m]}=S_L\left(F^{[n,m]}\right)$.
Recall that
\[
\mathrm{pw}(n,m)=n+(1+\alpha)m+\frac{n+m}{8}-\frac{m\varepsilon}{2},
\]
so that $L^{-n-(1+\alpha)m}\cdot L^{-\frac{n+m}{8}+\frac{m\varepsilon}{2}}=L^{-\mathrm{pw}(n,m)}$.
By (G2.r7) and Proposition~18.4\,(b), every $F\in\mathrm{Gr}[\psi,\chi]$ of level $k$ satisfies
\begin{equation}
\left\|S_LF\right\|_{h_{k+1}}=\sum_{n,m}\left\|S_LF^{[n,m]}\right\|_{h_{k+1}}=\sum_{n,m}L^{-\mathrm{pw}(n,m)}\cdot\left\|F^{[n,m]}\right\|_{h_k}.
\tag{18b.7}
\end{equation}
This identity is applied to $F=E(X;\mathscr{A}_L,\mathscr{C}_L)$.
This element is even;
so only the bidegrees $(n,m)$ with $n+m$ even occur, namely $(0,0)$, $(2,0)$, $(1,1)$, $(0,2)$ and the $(n,m)$ with $n+m\ge4$.

\emph{Bidegree $(0,0)$.}
Put $\mathrm{fn}_{00}(t):=E_{00}(X;t\cdot\mathrm{bh},t\cdot\mathscr{C}^{b})$ for $t\in\mathrm{Ot}$.
This function is holomorphic, and $\mathrm{fn}_{00}(0)=E_{00}(X;0,0)=0$ by Step 1.
The clause (k5) reads $\mathrm{Cg}\left[E(X;-\mathscr{A}',-\mathscr{C}')\right]=E(X;\mathscr{A}',\mathscr{C}')$ for $(\mathscr{A}',\mathscr{C}')\in\tilde{R}_k\times\mathfrak{X}_k$.
Since $\mathrm{Cg}$ is an automorphism that preserves the bidegree (Section~I.4.5), it fixes the part of bidegree $(0,0)$, which is a multiple of $1$.
Hence $E_{00}(X;-\mathscr{A}',-\mathscr{C}')=E_{00}(X;\mathscr{A}',\mathscr{C}')$;
as $\mathrm{Ot}=-\mathrm{Ot}$, this gives $\mathrm{fn}_{00}(-t)=\mathrm{fn}_{00}(t)$, and therefore $\mathrm{fn}_{00}'(0)=0$.
By Proposition~15\,(j) with $\mathrm{rad}=\mathrm{rad}_t>1$,
\[
\oint^{(2)}_{\mathrm{rad}_t}\mathrm{fn}_{00}=\mathrm{fn}_{00}(1)-\mathrm{fn}_{00}(0)-\mathrm{fn}_{00}'(0)=\mathrm{fn}_{00}(1),
\qquad
\left|\mathrm{fn}_{00}(1)\right|\le\frac{\sup_{|t|=\mathrm{rad}_t}\left|\mathrm{fn}_{00}(t)\right|}{\mathrm{rad}_t\cdot\left(\mathrm{rad}_t-1\right)}.
\]
By Lemma~1\,(i) and by (18b.6) at $(n,m)=(0,0)$, the right side is at most $2L^{-\frac{7}{2}}\cdot\mathrm{Nm}(X)$.
By Proposition~18.2\,(iii), $E_{00}(X;\mathscr{A}_L,\mathscr{C}_L)=\mathrm{fn}_{00}(1)$; and $\mathrm{pw}(0,0)=0$.
Hence
\begin{equation}
\left\|S_L\left[E^{[0,0]}(X;\mathscr{A}_L,\mathscr{C}_L)\right]\right\|_{h_{k+1}}=\left|E_{00}(X;\mathscr{A}_L,\mathscr{C}_L)\right|\le2L^{-\frac{7}{2}}\cdot\mathrm{Nm}(X).
\tag{18b.8}
\end{equation}

\emph{Bidegree $(2,0)$.}
Put $\mathrm{fn}_{20}(t):=E_{20}(X;t\cdot\mathrm{bh},t\cdot\mathscr{C}^{b};\cdot,\cdot)$ for $t\in\mathrm{Ot}$.
Its values lie in the finite-dimensional space of the functions of two indices $(\xi_1,\xi_2)$;
on this space $\|\cdot\|$ denotes the norm that assigns to a function the sum over all $(\xi_1,\xi_2)$ of $\left(\frac{\eta_k^3}{48}\right)^2$ times the modulus of its value at $(\xi_1,\xi_2)$.
The components of $\mathrm{fn}_{20}$ are holomorphic.

One has $\mathrm{fn}_{20}(0)=E_{20}(X;0,0;\cdot,\cdot)=0$ at all pairs of indices $\xi_1=(x,\beta,\omega_1)$, $\xi_2=(x',\beta',\omega_2)$.
Indeed, the matrix kernel $E_{20}(X;0,0;x,x')$ consists of the components at the bar indices $(\omega_1,\omega_2)=(1,0)$ (Section~I.5.1, matrix notation);
by Step 1 it vanishes for $x$, $x'$ over $X$, and by (k2) the kernel vanishes at arguments that do not lie over $X$.
Since $E$ is balanced (Step 1), the display (G2.0.1), which is part of Proposition~15\,(n), gives
\[
E^{[2,0]}(X;0,0)=\int\mathrm{d}x\int\mathrm{d}x'\,\bar{\psi}(x)^{T}E_{20}(X;0,0;x,x')\,\psi(x')=0
\]
and
\[
\frac{1}{2}\left\|\mathrm{fn}_{20}(0)\right\|=\frac{1}{2}\left\|E_{20}(X;0,0)\right\|=\sum_{\beta,\beta'}\int\mathrm{d}x\int\mathrm{d}x'\,\left|E_{20}(X;0,0;x,x')_{\beta\beta'}\right|=0.
\]
Component by component this reads as follows.
By antisymmetry, the components at $(\omega_1,\omega_2)=(0,1)$ are the negatives of those at $(1,0)$ with the two indices exchanged.
The components at $(0,0)$ and at $(1,1)$ vanish for every $t\in\mathrm{Ot}$:
in the balanced element $E(X;t\cdot\mathrm{bh},t\cdot\mathscr{C}^{b})$ the coefficient of a product of two unbarred generators $\psi$, or of two barred ones, is $0$,
and by Proposition~15\,(a) the antisymmetric kernel at increasing arguments is this coefficient divided by the two weights $\frac{\eta_k^3}{48}$.

By Proposition~15\,(j),
\[
\oint_{\mathrm{rad}_t}\mathrm{fn}_{20}=\mathrm{fn}_{20}(1)-\mathrm{fn}_{20}(0)=\mathrm{fn}_{20}(1),
\]
and, by Lemma~1\,(i),
\[
\left\|\mathrm{fn}_{20}(1)\right\|\le\frac{\sup_{|t|=\mathrm{rad}_t}\left\|\mathrm{fn}_{20}(t)\right\|}{\mathrm{rad}_t-1}\le2L^{-\frac{7}{4}}\cdot\sup_{|t|=\mathrm{rad}_t}\left\|\mathrm{fn}_{20}(t)\right\|.
\]
Summing the inequality of Proposition~18.2\,(iii) at $(n,m)=(2,0)$ over $(\xi_1,\xi_2)$ with the weights $\left(\frac{\eta_k^3}{48}\right)^2$ gives
\[
\left\|E_{20}(X;\mathscr{A}_L,\mathscr{C}_L)\right\|\le\mathrm{e}^4\cdot\left\|\mathrm{fn}_{20}(1)\right\|.
\]
With $\mathrm{pw}(2,0)=\frac{9}{4}$, (G2.0.1) and (18b.6) at $(n,m)=(2,0)$ it follows that
\begin{equation}
\begin{split}
\left\|S_L\left[E^{[2,0]}(X;\mathscr{A}_L,\mathscr{C}_L)\right]\right\|_{h_{k+1}}
&=L^{-\frac{9}{4}}\cdot\frac{h_{k,1}^2}{2}\cdot\left\|E_{20}(X;\mathscr{A}_L,\mathscr{C}_L)\right\|\\
&\le L^{-\frac{9}{4}}\cdot\mathrm{e}^4\cdot2L^{-\frac{7}{4}}\cdot\sup_{|t|=\mathrm{rad}_t}\frac{h_{k,1}^2}{2}\cdot\left\|\mathrm{fn}_{20}(t)\right\|\\
&\le2\,\mathrm{e}^4\cdot L^{-4}\cdot\mathrm{Nm}(X).
\end{split}
\tag{18b.9}
\end{equation}

\emph{The other bidegrees.}
By Lemma~1\,(iii),
\begin{gather*}
L^{-\mathrm{pw}(1,1)}=L^{-2-\alpha-\frac{1}{4}+\frac{\varepsilon}{2}}\le L^{-\frac{13}{4}+\frac{3\varepsilon}{2}},
\qquad
L^{-\mathrm{pw}(0,2)}=L^{-2-2\alpha-\frac{1}{4}+\varepsilon}\le L^{-\frac{13}{4}+\frac{3\varepsilon}{2}},\\
L^{-\mathrm{pw}(n,m)}\le L^{-4}\le L^{-\frac{13}{4}+\frac{3\varepsilon}{2}}\qquad(n+m\ge4).
\end{gather*}
The pair $(\mathscr{A}_L,\mathscr{C}_L)$ lies in $\tilde{R}_k\times\mathfrak{X}_k$ (Step 0).
Hence, by (G2.r7),
\begin{equation}
\begin{split}
\sum_{(n,m)\ne(0,0),(2,0)}\left\|S_L\left[E^{[n,m]}(X;\mathscr{A}_L,\mathscr{C}_L)\right]\right\|_{h_{k+1}}
&\le L^{-\frac{13}{4}+\frac{3\varepsilon}{2}}\cdot\sum_{n,m}\left\|E^{[n,m]}(X;\mathscr{A}_L,\mathscr{C}_L)\right\|_{h_k}\\
&=L^{-\frac{13}{4}+\frac{3\varepsilon}{2}}\cdot\left\|E(X;\mathscr{A}_L,\mathscr{C}_L)\right\|_{h_k}\\
&\le L^{-\frac{13}{4}+\frac{3\varepsilon}{2}}\cdot\mathrm{Nm}(X).
\end{split}
\tag{18b.10}
\end{equation}

\emph{Sum.}
By Lemma~1\,(i),
\[
2L^{-\frac{7}{2}}=\left(2L^{-\frac{1}{4}}\right)\cdot L^{-\frac{13}{4}}\le L^{-\frac{13}{4}},
\qquad
2\,\mathrm{e}^4\cdot L^{-4}=\left(2\,\mathrm{e}^4\cdot L^{-\frac{3}{4}}\right)\cdot L^{-\frac{13}{4}}\le L^{-\frac{13}{4}},
\]
and $L^{-\frac{13}{4}}\le L^{-\frac{13}{4}+\frac{3\varepsilon}{2}}$.
By (18b.7), the norm $\left\|S_L\left[E(X;\mathscr{A}_L,\mathscr{C}_L)\right]\right\|_{h_{k+1}}$ is the sum of the left sides of (18b.8), (18b.9) and (18b.10).
In (18b.8) and (18b.9) a single term of the sum in (G2.r7) is bounded by its supremum over the circle $|t|=\mathrm{rad}_t$, which is at most $\mathrm{Nm}(X)$ by (18b.6);
in (18b.10) the remaining terms are evaluated at the single point $(\mathscr{A}_L,\mathscr{C}_L)$.
In each of the three bounds, therefore, the supremum over the connections stands outside the sum over $(n,m)$, as in (k6), and adding them one obtains
\begin{equation}
\left\|S_L\left[E(X;\mathscr{A}_L,\mathscr{C}_L)\right]\right\|_{h_{k+1}}\le3L^{-\frac{13}{4}+\frac{3\varepsilon}{2}}\cdot\mathrm{Nm}(X)=3L^{-\frac{13}{4}+\frac{3\varepsilon}{2}}\cdot\sup_{\tilde{R}_k\times\mathfrak{X}_k}\left\|E(X;\cdot,\cdot)\right\|_{h_k}.
\tag{18b.11}
\end{equation}

\emph{The last assertion of (d).}
The factor $L^{-\mathrm{pw}(n,m)}$ in (18b.7) depends on $(n,m)$ alone:
$S_L$ multiplies $\chi(x,x',\beta,\omega)$ by $L^{-1-\alpha}$ for all labels $g$, $g'$ and for every value of $d_X(x,x')$ in $[\eta_k,M]$,
and the norms (I.5.1) give every generator $\chi(\boldsymbol{\zeta})$ the weight $h_{k,2}$, respectively $h_{k+1,2}$.
A leg at a pair inside a cluster ($c=c'$, $g\ne g'$, $d_X=\eta_k$) is counted like every other H\"older leg.
The factors $2L^{-\frac{7}{2}}$ and $2\,\mathrm{e}^4\cdot L^{-\frac{7}{4}}$ in (18b.8) and (18b.9) arise from the dependence on the connections.
No inequality of Steps 1 and 2 involves the number $\eta_k^{\alpha}$.
Together with Step 1, this proves (d).

\paragraph{Step 3 (clause (b)).}
Let $Y\in D_{k+1}$ and $(\mathscr{A},\mathscr{C})\in\tilde{R}_{k+1}\times\mathfrak{X}_{k+1}$.
Since the norm of a sum is at most the sum of the norms,
\[
\left\|(LR\mathscr{E})(Y;\mathscr{A},\mathscr{C})\right\|_{h_{k+1}}\le\sum_{X:\,\bar{X}^{L}=LY}\left\|S_L\left[E(X;\mathscr{A}_L,\mathscr{C}_L)\right]\right\|_{h_{k+1}}.
\]

\emph{Small $X$.}
The assertion (t8) of \cite[Theorem~13.2]{QEDI}, recalled in Section~0.3, is applied with $\mathrm{kd}=\kappa$; this is admissible, since $\kappa\ge\kappa_0$ by Lemma~1\,(ii).
Together with (18b.11), (18b.5) and the inequality $L\cdot(\kappa-\kappa_0-1)\ge\kappa$ of Lemma~1\,(ii), it gives
\begin{equation}
\begin{split}
\sum_{X\text{ small},\ \bar{X}^{L}=LY}\left\|S_L\left[E(X;\mathscr{A}_L,\mathscr{C}_L)\right]\right\|_{h_{k+1}}
&\le3L^{-\frac{13}{4}+\frac{3\varepsilon}{2}}\cdot20\,\|\mathscr{E}\|^{x}_k\cdot\sum_{X:\,\bar{X}^{L}=LY}\exp\left(-\kappa\cdot d_M(X)\right)\\
&\le60\,L^{-\frac{13}{4}+\frac{3\varepsilon}{2}}\cdot\|\mathscr{E}\|^{x}_k\cdot c_{\mathrm{rb}}\cdot L^{3}\cdot\exp\left(-\kappa\cdot d_M(Y)\right).
\end{split}
\tag{18b.12}
\end{equation}

\emph{Large $X$.}
By Proposition~18.4\,(b),
\[
\left\|S_LF\right\|_{h_{k+1}}=\left\|F\right\|_{\left(L^{-1}h_{k+1,1},\,L^{-1-\alpha}h_{k+1,2}\right)},
\]
and by Lemma~1\,(iv) and Proposition~15\,(k) this is at most $\|F\|_{h_k}$.
For large $X$ one has $d_M(X)>L$, hence
\[
\exp\left(-\kappa\cdot d_M(X)\right)\le\exp(-L)\cdot\exp\left(-(\kappa-1)\cdot d_M(X)\right).
\]
By (18b.5), by (t8) with $\mathrm{kd}=\kappa-1\ge\kappa_0$, and by the inequality $L\cdot(\kappa-\kappa_0-2)\ge\kappa$ of Lemma~1\,(ii),
\begin{equation}
\begin{split}
\sum_{X\text{ large},\ \bar{X}^{L}=LY}\left\|S_L\left[E(X;\mathscr{A}_L,\mathscr{C}_L)\right]\right\|_{h_{k+1}}
&\le\|\mathscr{E}\|^{x}_k\cdot\exp(-L)\cdot\sum_{X:\,\bar{X}^{L}=LY}\exp\left(-(\kappa-1)\cdot d_M(X)\right)\\
&\le\exp(-L)\cdot c_{\mathrm{rb}}\cdot L^{3}\cdot\|\mathscr{E}\|^{x}_k\cdot\exp\left(-\kappa\cdot d_M(Y)\right).
\end{split}
\tag{18b.13}
\end{equation}

\emph{Sum.}
By Lemma~1\,(i), $L^{3}\cdot\exp(-L)\le L^{-\frac{1}{4}}\le L^{-\frac{1}{4}+\frac{3\varepsilon}{2}}$.
Hence
\[
\left\|(LR\mathscr{E})(Y;\mathscr{A},\mathscr{C})\right\|_{h_{k+1}}\le61\,c_{\mathrm{rb}}\cdot L^{-\frac{1}{4}+\frac{3\varepsilon}{2}}\cdot\|\mathscr{E}\|^{x}_k\cdot\exp\left(-\kappa\cdot d_M(Y)\right).
\]
By the condition (C0) of Section~I.8.2, $61\,c_{\mathrm{rb}}=61\,c_{X}\cdot K_{0}=c_{\mathrm{fl}}$;
moreover $L^{\frac{3\varepsilon}{2}}\le L^{2\varepsilon}$.
Taking the supremum over $(\mathscr{A},\mathscr{C})$ and over $Y$, as in the definition (k6) of the norm at level $k+1$, one obtains
\begin{equation}
\|LR\mathscr{E}\|^{x}_{k+1}\le c_{\mathrm{fl}}\cdot L^{-\frac{1}{4}+\frac{3\varepsilon}{2}}\cdot\|\mathscr{E}\|^{x}_k\le c_{\mathrm{fl}}\cdot L^{-\frac{1}{4}+2\varepsilon}\cdot\|\mathscr{E}\|^{x}_k.
\tag{18b.14}
\end{equation}

\emph{Relative to the windows.}
By Lemma~1\,(iv), $w^{\mathrm{E}}_{k+1}=L^{\frac{1}{8}-\frac{7\varepsilon}{2}}\cdot w^{\mathrm{E}}_k$.
Division of (18b.14) by $w^{\mathrm{E}}_{k+1}$ gives
\[
\frac{\|LR\mathscr{E}\|^{x}_{k+1}}{w^{\mathrm{E}}_{k+1}}\le c_{\mathrm{fl}}\cdot L^{-\frac{1}{4}+2\varepsilon-\frac{1}{8}+\frac{7\varepsilon}{2}}\cdot\frac{\|\mathscr{E}\|^{x}_k}{w^{\mathrm{E}}_k}=c_{\mathrm{fl}}\cdot L^{-\frac{3}{8}+\frac{11\varepsilon}{2}}\cdot\frac{\|\mathscr{E}\|^{x}_k}{w^{\mathrm{E}}_k}.
\]

\paragraph{Step 4 (clause (a)).}
\emph{Linearity.}
The map $\mathscr{E}\mapsto E$ is linear (Step 1), and $S_L$ is linear;
hence the map $\mathscr{E}\mapsto(LR\mathscr{E})(Y;\mathscr{A},\mathscr{C})=\sum_{X}S_L\left[E(X;\mathscr{A}_L,\mathscr{C}_L)\right]$ is linear.

\emph{Form of the elements.}
The homomorphism $S_L$ maps $\psi(x,\beta,\omega)$ to a multiple of $\psi(x/L,\beta,\omega)$ and $\chi(x,x',\beta,\omega)$ to a multiple of $\chi(x/L,x'/L,\beta,\omega)$;
thus it preserves the bidegree and the bar index of every leg.
Hence $(LR\mathscr{E})(Y;\mathscr{A},\mathscr{C})$ is an even, balanced element of $\mathrm{Gr}[\psi,\chi]$ of level $k+1$.
\begin{itemize}
\item[(k1)] The map that assigns to an element its antisymmetric kernels is linear (Proposition~15\,(a)).
By Proposition~18.4\,(c), $(LR\mathscr{E})_{nm}(Y;\mathscr{A},\mathscr{C};\xi;\boldsymbol{\zeta})$ equals
\[
\sum_{X:\,\bar{X}^{L}=LY}L^{3n+6m}\cdot L^{-n-(1+\alpha)m}\cdot E_{nm}(X;\mathscr{A}_L,\mathscr{C}_L;L\xi;L\boldsymbol{\zeta})
\]
if all pairs of $\boldsymbol{\zeta}$ have $d_X\le M/L$, and $0$ otherwise.
The sum is finite.
Each term is the composition of the linear map of Step 0, which maps the open set $\tilde{R}_{k+1}\times\mathfrak{X}_{k+1}$ into $\tilde{R}_k\times\mathfrak{X}_k$, with a function that is holomorphic on $\tilde{R}_k\times\mathfrak{X}_k$, by (k1) for $E$.
\item[(k2)] $E$ has (k2) at level $k$; by the last sentence of Proposition~18.3\,(e), $LE$ has (k2) at level $k+1$.
\item[(k3)] $E$ has (k3) for all $\tilde{g}\in\mathbb{G}_k$, hence for the elements of the subgroup $\mathbb{G}_k^{L}$;
by Proposition~16.3\,(e), $LE$ has (k3) at level $k+1$, for the group $\mathbb{G}_{k+1}$.
\item[(k4), (k5)] $E$ has them; by Proposition~16.3\,(e), so has $LE$ at level $k+1$.
\item[(k6)] This is (18b.14).
\end{itemize}
Hence $LR\mathscr{E}\in K^{x,C}_{k+1}(\ell/L)$.
If $\mathscr{E}\in K^{x}_k$, then $E$ has (k7) (Step 1), and by Proposition~16.3\,(e) so has $LE$ at level $k+1$;
that is, $LR\mathscr{E}\in K^{x}_{k+1}(\ell/L)$.

\paragraph{Step 5 (clause (c)).}
The space $K_k^{C}$ and the family $\mathscr{E}|_0$ are defined in Section~I.5.4;
the families without slot, their clauses (k1)$^{\mathrm{o}}$--(k7)$^{\mathrm{o}}$, and the maps $R$, ${E'}\mapsto{E'}^{\mathrm{nat}}$, $L$ and $LR$ on such families are defined in Section~I.5.11.
By the data of Proposition~48.5, the map $LR$ of Proposition~48.5 is the map $LR$ of Section~I.5.11 on $K_k^{C}$.

\emph{The right side.}
The family $\mathscr{E}|_0$ lies in $K_k^{C}$.
This is part of the assertion of Section~I.5.4 on $\mathscr{E}|_0$, which is Proposition~48.1\,(a):
$\mathscr{E}|_0$ lies in $K_k^{C}$, and in $K_k$ if $\mathscr{E}\in K^{x}_k$, and $\|\mathscr{E}|_0\|_k\le\|\mathscr{E}\|^{x}_k$.
The reasons for the membership in $K_k^{C}$ are as follows.
The bond function $0$ lies in $\mathfrak{X}_k$ (Section~0.3), and $\mathscr{A}'\mapsto(\mathscr{A}',0)$ is a linear map of the open set $\tilde{R}_k$ into $\tilde{R}_k\times\mathfrak{X}_k$;
the elements $\mathscr{E}(X;\mathscr{A}',0)$ are even and balanced;
the clauses (k1)$^{\mathrm{o}}$, (k2)$^{\mathrm{o}}$, (k3)$^{\mathrm{o}}$, (k5)$^{\mathrm{o}}$, (k6)$^{\mathrm{o}}$ of Section~I.5.11 for $\mathscr{E}|_0$ are the clauses (k1), (k2), (k3), (k5), (k6) for $\mathscr{E}$ at $\mathscr{C}=0$, because $\tilde{g}\,0=0$ and $-0=0$;
one has $\|\mathscr{E}|_0\|_k\le\|\mathscr{E}\|^{x}_k$, since a supremum over a subset is at most the supremum over the whole set;
and (k4)$^{\mathrm{o}}$ is (k4) with $\mathscr{C}^{1}=\mathscr{C}=0$ and $\boldsymbol{\lambda}'=0$.
Hence the right side of (c) is the map $LR$ of Proposition~48.5, evaluated at a point of its domain $K_k^{C}$.

\emph{The identity.}
Let $Y\in D_{k+1}$ and $\mathscr{A}\in\tilde{R}_{k+1}$.
Then $0\in\mathfrak{X}_{k+1}$, and $0_L=0$ (Step 0).
Hence
\[
(LR\mathscr{E})|_0(Y;\mathscr{A})=(LR\mathscr{E})(Y;\mathscr{A},0)=\sum_{X:\,\bar{X}^{L}=LY}S_L\left[(R\mathscr{E})^{\mathrm{nat}}(X;\mathscr{A}_L,0)\right].
\]
By the last assertion of Proposition~17.2\,(e), applied to the family $R\mathscr{E}$ of $K^{x,C}_k$,
\[
(R\mathscr{E})^{\mathrm{nat}}(X;\mathscr{A}_L,0)=\left((R\mathscr{E})^{\mathrm{nat}}\right)\Big|_0(X;\mathscr{A}_L)=\left((R\mathscr{E})|_0\right)^{\mathrm{nat}}(X;\mathscr{A}_L).
\]
By the last but one assertion of Proposition~16.3\,(d), $(R\mathscr{E})|_0=R(\mathscr{E}|_0)$.
(This family without slot has the locality clause (k2)$^{\mathrm{o}}$ of Section~I.5.11, because $R\mathscr{E}$ has (k2);
so its part of bidegree $(2,0)$ lies in the domain of $\boldsymbol{\Theta}_{X,\mathscr{A}_L}$, and its adjustment is defined.)
Hence
\[
(LR\mathscr{E})|_0(Y;\mathscr{A})=\sum_{X:\,\bar{X}^{L}=LY}S_L\left[\left(R(\mathscr{E}|_0)\right)^{\mathrm{nat}}(X;\mathscr{A}_L)\right]=\left(LR(\mathscr{E}|_0)\right)(Y;\mathscr{A}).
\]

\paragraph{Step 6 (clause (e)).}
\emph{The first sentence.}
Let $E'$ be a family of level $k$, and let $(\mathscr{A},\mathscr{C})\in\tilde{R}_{k+1}\times\mathfrak{X}_{k+1}$.
By the assertion (t5) of \cite[Theorem~13.2]{QEDI}, recalled in Section~0.3, every $X\in D_k$ satisfies $\bar{X}^{L}=LY$ for exactly one $Y\in D_{k+1}$.
Hence the sets $\{X:\bar{X}^{L}=LY\}$, $Y\in D_{k+1}$, are pairwise disjoint, and their union is $D_k$.
Since $S_L$ is linear and all sums are finite,
\begin{equation}
\begin{split}
\sum_{Y\in D_{k+1}}(LE')(Y;\mathscr{A},\mathscr{C})
&=\sum_{Y}\ \sum_{X:\,\bar{X}^{L}=LY}S_L\left[E'(X;\mathscr{A}_L,\mathscr{C}_L)\right]\\
&=S_L\left[\sum_{X\in D_k}E'(X;\mathscr{A}_L,\mathscr{C}_L)\right].
\end{split}
\tag{18b.15}
\end{equation}

\emph{The relations at level $k$.}
Let $\mathscr{T}:=\mathscr{A}+\mathscr{C}$, and let $\mathrm{ev}'$ be a homomorphism of $\mathrm{Gr}[\psi,\chi]$ of level $k+1$ into an associative algebra as in (e).
Let $\psi'(x)$, $\bar{\psi}'(x)$, $\chi'(x,x')$, $\bar{\chi}'(x,x')$ be the columns of the images under $\mathrm{ev}'$ of the generators of level $k+1$, indexed by $\mathbb{L}^{\mathrm{f}}_{k+1}$ and $\mathrm{Pzt}_{k+1}$ as in (I.AP) of Section~I.3.6.
For $x\in\mathbb{L}^{\mathrm{f}}_k$ and $(x,x')\in\mathrm{Pzt}_k$ put (Section~I.3.7)
\begin{gather*}
\psi'_L(x):=L^{-1}\psi'(x/L),\qquad\bar{\psi}'_L(x):=L^{-1}\bar{\psi}'(x/L),\\
\chi'_L(x,x'):=L^{-1-\alpha}\chi'(x/L,x'/L),\qquad\bar{\chi}'_L(x,x'):=L^{-1-\alpha}\bar{\chi}'(x/L,x'/L).
\end{gather*}
The hypothesis $\chi'=\delta_{\alpha,\mathscr{T}}\psi'$, $\bar{\chi}'=\delta_{\alpha,-\mathscr{T}}\bar{\psi}'$ of (e) reads, for $(p,p')\in\mathrm{Pzt}_{k+1}$, with $d_X$ and $U$ formed at level $k+1$ ((I.1.5), (I.1.6), (I.1.7)), and with $U_{-\mathscr{T}}=U_{\mathscr{T}}^{-1}$,
\begin{gather*}
\chi'(p,p')=d_X(p,p')^{-\alpha}\left(U_{\mathscr{T}}(p,p')\,\psi'(p')-\psi'(p)\right),\\
\bar{\chi}'(p,p')=d_X(p,p')^{-\alpha}\left(U_{\mathscr{T}}(p,p')^{-1}\,\bar{\psi}'(p')-\bar{\psi}'(p)\right).
\end{gather*}
Let $(x,x')\in\mathrm{Pzt}_k$, $x=(c,g)$, $x'=(c',g')$.
\begin{itemize}
\item[(1)] The points $x/L=(c/L,g)$ and $x'/L=(c'/L,g')$ of $\mathbb{L}^{\mathrm{f}}_{k+1}$ are distinct, and
\[
d_X(x/L,x'/L)=\max\left(\frac{|c-c'|}{L},\eta_{k+1}\right)=\frac{\max\left(|c-c'|,\eta_k\right)}{L}=\frac{d_X(x,x')}{L}\le\frac{M}{L};
\]
hence $(x/L,x'/L)\in\mathrm{Pzt}_{k+1}$, and $d_X(x,x')^{-\alpha}=L^{-\alpha}\cdot d_X(x/L,x'/L)^{-\alpha}$.
A pair inside a cluster ($c=c'$, $d_X=\eta_k$) goes to a pair inside a cluster ($d_X=\eta_{k+1}$).
\item[(2)] Regard $\mathscr{T}$ as an $(\ell/L)\mathbb{Z}^3$-periodic function on the bonds of $\mathbb{L}_{k+1}$;
then $\mathscr{T}_L$ is an $\ell\mathbb{Z}^3$-periodic function on the bonds of $\mathbb{L}_k$.
Let $c\ne c'$.
For each of the six orders $\pi$, the map $b\mapsto b/L$ maps the bonds of the path $\boldsymbol{\Gamma}^{\pi}(c,c')$ of $\mathbb{L}_k$ onto the bonds of the path $\boldsymbol{\Gamma}^{\pi}(c/L,c'/L)$ of $\mathbb{L}_{k+1}$, and $\mathscr{T}_L(b)=L^{-\frac{1}{2}}\cdot\mathscr{T}(b/L)$.
Hence $(\tau\mathscr{T}_L)(c,c')=L^{-\frac{1}{2}}\cdot(\tau\mathscr{T})(c/L,c'/L)$.
Since
\[
e_k\eta_k\cdot L^{-\frac{1}{2}}=\left(L^{-\frac{1}{2}}e_{k+1}\right)\cdot\left(L\eta_{k+1}\right)\cdot L^{-\frac{1}{2}}=e_{k+1}\eta_{k+1},
\]
it follows that
\[
\exp\left\{ie_k\eta_k\cdot(\tau\mathscr{T}_L)(c,c')\right\}=\exp\left\{ie_{k+1}\eta_{k+1}\cdot(\tau\mathscr{T})(c/L,c'/L)\right\},
\]
that is, $U_{\mathscr{T}_L}(x,x')$, formed at level $k$, equals $U_{\mathscr{T}}(x/L,x'/L)$, formed at level $k+1$.
If $c=c'$, both numbers are $1$.
\item[(3)] By (1), (2) and the hypothesis at $(p,p')=(x/L,x'/L)$,
\begin{align*}
&d_X(x,x')^{-\alpha}\left(U_{\mathscr{T}_L}(x,x')\,\psi'_L(x')-\psi'_L(x)\right)\\
&\qquad=L^{-\alpha}d_X(x/L,x'/L)^{-\alpha}\cdot L^{-1}\left(U_{\mathscr{T}}(x/L,x'/L)\,\psi'(x'/L)-\psi'(x/L)\right)\\
&\qquad=L^{-1-\alpha}\chi'(x/L,x'/L)=\chi'_L(x,x'),\\
&d_X(x,x')^{-\alpha}\left(U_{\mathscr{T}_L}(x,x')^{-1}\,\bar{\psi}'_L(x')-\bar{\psi}'_L(x)\right)\\
&\qquad=L^{-\alpha}d_X(x/L,x'/L)^{-\alpha}\cdot L^{-1}\left(U_{\mathscr{T}}(x/L,x'/L)^{-1}\,\bar{\psi}'(x'/L)-\bar{\psi}'(x/L)\right)\\
&\qquad=L^{-1-\alpha}\bar{\chi}'(x/L,x'/L)=\bar{\chi}'_L(x,x').
\end{align*}
\end{itemize}
These are the relations $\chi'_L=\delta_{\alpha,\mathscr{T}_L}\psi'_L$ and $\bar{\chi}'_L=\delta_{\alpha,-\mathscr{T}_L}\bar{\psi}'_L$ at all pairs of $\mathrm{Pz}_k$;
this is the second relation of (I.O9), for functions with values in an algebra.

\emph{The identity.}
The map $\mathrm{ev}:=\mathrm{ev}'\circ S_L$ is a homomorphism of $\mathrm{Gr}[\psi,\chi]$ of level $k$ into the algebra.
By (G2.r5) it maps $\psi(x)$, $\bar{\psi}(x)$, $\chi(x,x')$, $\bar{\chi}(x,x')$ to $\psi'_L(x)$, $\bar{\psi}'_L(x)$, $\chi'_L(x,x')$, $\bar{\chi}'_L(x,x')$.
By (3), these images satisfy the relations of Proposition~17.1\,(u4) with $B=\mathscr{T}_L$ at all pairs of $\mathrm{Pz}_k$;
in particular they satisfy them at all pairs of kind (I) for every small $X$ (Section~I.1.7), the pairs inside the clusters over $X$ included.
The pair $(\mathscr{A}_L,\mathscr{C}_L)$ lies in $\tilde{R}_k\times\mathfrak{X}_k$, and $\mathscr{A}_L+\mathscr{C}_L=\mathscr{T}_L$ (Step 0);
$R\mathscr{E}$ is a family of level $k$ with (k2) (Step 1).
By Proposition~17.2\,(a), for small $X$,
\[
\mathrm{ev}\left(E(X;\mathscr{A}_L,\mathscr{C}_L)\right)=\mathrm{ev}\left((R\mathscr{E})(X;\mathscr{A}_L,\mathscr{C}_L)\right).
\]
For large $X$, $E(X)=(R\mathscr{E})(X)$ by definition.
By (18b.15) with $E'=E$, it follows that
\[
\begin{split}
\mathrm{ev}'\left[\sum_{Y}(LR\mathscr{E})(Y;\mathscr{A},\mathscr{C})\right]
&=\mathrm{ev}\left[\sum_{X\in D_k}E(X;\mathscr{A}_L,\mathscr{C}_L)\right]
=\mathrm{ev}\left[\sum_{X\in D_k}(R\mathscr{E})(X;\mathscr{A}_L,\mathscr{C}_L)\right]\\
&=\mathrm{ev}'\,S_L\left[\sum_{X}(R\mathscr{E})(X;\mathscr{A}_L,\mathscr{C}_L)\right].
\end{split}
\]

\paragraph{Step 7 ((f))}
In this step the class of sides consists of all $\ell$ with $\ell\in\mathrm{Cl}$ and $\ell/L\in\mathrm{Cl}$, and the notation of Lemma~3 is used.

\textbf{(7.1) The map $S_L$.}
Let $\tilde U$ be a finite set of sites of $\mathbb{L}_k$, and let $\mathrm{Slc}\colon\mathrm{Grc}_k(\tilde U)\to\mathrm{Grc}_{k+1}(\tilde U/L)$ be the homomorphism with
\[
\mathrm{Slc}\,\psi(x,\beta,\omega):=L^{-1}\cdot\psi(x/L,\beta,\omega),\qquad
\mathrm{Slc}\,\chi(x,x',\beta,\omega):=L^{-1-\alpha}\cdot\chi(x/L,x'/L,\beta,\omega),\qquad
x/L=(c/L,g).
\]
The right-hand sides make sense: $\tilde U/L$ is a finite set of sites of $\mathbb{L}_{k+1}=L^{-1}\mathbb{L}_k$, and $(x/L,x'/L)$ belongs to $\mathrm{Pzt}_{k+1}$ by (1) of Step~6.
These formulas are those of (G2.r5).
They involve $L$ and $\alpha$, but neither $\ell$ nor $\tilde U$, and no data enter.
The generator with the index $t$ is sent to a multiple of the generator with the index $t/L$, and $\operatorname{supp}(t/L)=(\operatorname{supp}t)/L$; so the range is $0$.

\emph{Translations.}
For $v\in\eta_k\mathbb{Z}^3$ one has $(x+v)/L=x/L+v/L$.
Consequently $\mathrm{Slc}\,T_v=T_{v/L}\,\mathrm{Slc}$ on the generators, and hence on $\mathrm{Grc}_k(\tilde U)$.
If $v\in L\cdot M\mathbb{Z}^3$, then $v/L\in M\mathbb{Z}^3$.

\emph{Tori.}
Assume that $\ell,\ell/L\in\mathrm{Cl}$ and $\operatorname{sup-diameter}(\tilde U)+\eta_k<\ell$.
Then $\operatorname{sup-diameter}(\tilde U/L)+\eta_{k+1}<\ell/L$.
Lemma~3\,(a), applied at the level $k$ with the side $\ell$ and at the level $k+1$ with the side $\ell/L$, shows that $\mathrm{cov}_\ell$ is injective on $\mathrm{Grc}_k(\tilde U)$ and that $\mathrm{cov}_{\ell/L}$ is injective on $\mathrm{Grc}_{k+1}(\tilde U/L)$.
Moreover,
\begin{equation}
\mathrm{cov}_{\ell/L}\circ\mathrm{Slc}=S_L\circ\mathrm{cov}_\ell\qquad\text{on }\mathrm{Grc}_k(\tilde U),
\tag{18b.16}
\end{equation}
since both sides are homomorphisms which map $\psi(x,\beta,\omega)$ to $L^{-1}$ times the generator $\psi(x/L,\beta,\omega)$ of the torus of side $\ell/L$, and $\chi(x,x',\beta,\omega)$ to $L^{-1-\alpha}$ times the generator $\chi(x/L,x'/L,\beta,\omega)$.
(The formula (G2.r5) is a formula for the generators indexed by the points of the covering lattices, and it is compatible with (I.AP), because $(x+\ell\cdot e_\mu)/L=x/L+(\ell/L)\cdot e_\mu$.)
The margin $\operatorname{sup-diameter}(\tilde U)+L\cdot M<\ell$ of Section~I.1.9 implies the hypothesis made here; under that margin, $\mathrm{cov}_\ell$ and $\mathrm{cov}_{\ell/L}$ are the restrictions of the maps of Section~I.1.9.
This is the first sentence of (f).
Lemma~3\,(a) and the identity (18b.16) together give the assertion of (f) that, as soon as $\operatorname{sup-diameter}(\tilde U)+\eta_k<\ell$, the rule of Section~I.1.9 defines injective homomorphisms $\mathrm{cov}_\ell$, $\mathrm{cov}_{\ell/L}$ on the subalgebras generated by $\psi$, $\chi$, and that these satisfy $\mathrm{cov}_{\ell/L}\circ\mathrm{Slc}=S_L\circ\mathrm{cov}_\ell$; see (7.7).

\textbf{(7.2) The maps.}
Let $\tilde X$ be a finite, non-empty union of $M$-cubes of $\mathbb{L}_k$ that is connected through faces, and let $\tilde U:=\tilde X^{+\eta_k}$.
A triple of data is a triple $\mathrm{dat}=(K^0,K,\tilde B)$, where $K^0$ and $K$ are systems of antisymmetric kernels on the indices over $\tilde U$, and $\tilde B$ is a complex function on the bonds of $\mathbb{L}_k$ with both ends in $\tilde X$ that is odd under reversal.
Denote by $K^0_{20}(x,x')$ the $4\times4$ matrix with the entries $K^0_{20}\left((x,\beta,1),(x',\beta',0)\right)$, and put
\begin{gather*}
T_v\,\mathrm{dat}:=\left(T_vK^0,T_vK,T_v\tilde B\right),\qquad\left(T_v\tilde B\right)(b):=\tilde B(b-v),\\
\int_{\tilde X}\mathrm{d}x:=\eta_k^3\cdot\frac{1}{48}\sum_{c\in\tilde X}\sum_{g\in O_{\mathrm{cub}}},\qquad\operatorname{Vol}(\tilde X):=M^3\cdot|\tilde X|_M.
\end{gather*}
Let $\boldsymbol{\Theta}_{\tilde X,\tilde B}$ be the homomorphism which Proposition~17.2\,(f) assigns to $(\tilde X,\tilde B)$: the display (3.3) of Proposition~17.1, formed with the cubes of $\tilde X$, their centres and their common faces (through the shortest chains), the labels, the measure $\int_{\tilde X}\mathrm{d}\tilde x$ and the transports $u_{\tilde B}$.
It maps the subalgebra of $\mathrm{Grc}_k(\tilde X)$ generated by the $\psi$ into $\mathrm{Grc}_k(\tilde X)$.
(By the data of Proposition~17.1 and by its clauses (b), (c) with (c1), and (d), all points in the displays (3.2), (3.3) of Proposition~17.1 are points of $\mathbb{L}^{\mathrm{f}}_k$ over $\tilde X$.
The chain points of a shortest chain $(\square_0,\dots,\square_n)$ of cubes of $\tilde X$ with the centres $z_i$ are $p_0=x=(c,g)$, $p_{i+1}=(z_i,g)$ and $p_{n+2}=\tilde x$.
Each chain pair $(p_i,p_{i+1})$ with $p_i\ne p_{i+1}$ belongs to $\mathrm{Pzt}_k$ and is of kind (I): its two photon sites lie in one cube of $\tilde X$ or are the centres of two cubes of $\tilde X$ with a common face, and its six paths lie in $\tilde X$.
Hence $\chi(p_i,p_{i+1},\beta,\omega)$ is a generator of $\mathrm{Grc}_k(\tilde X)$, and by (u2) of Proposition~17.1 the transport $u_{\tilde B}(p_i,p_{i+1})$ is formed with the values of $\tilde B$ on bonds with both ends in $\tilde X$.
All this concerns cubes and points of the covering lattice, and $\ell$ does not occur in it.)
For $F\in\mathrm{Grc}_k(\tilde U)$ let $F^{[2,0]}_{\tilde X}$ be the sum of the terms of $F^{[2,0]}$ both of whose generators have their sites over $\tilde X$.
If $d_M(\tilde X)\le L$, define
\begin{align*}
\Phi^R(\tilde X;K^0,K)&:=\mathrm{El}(K)-K^0_{00}-\int_{\tilde X}\mathrm{d}x\;\bar\psi(x)^T\cdot\left[\operatorname{Vol}(\tilde X)^{-1}\int_{\tilde X}\mathrm{d}x'\int_{\tilde X}\mathrm{d}x''\,K^0_{20}(x',x'')\right]\cdot\psi(x),\\
\Phi^E(\tilde X;\mathrm{dat})&:=F-F^{[2,0]}_{\tilde X}+\boldsymbol{\Theta}_{\tilde X,\tilde B}\left(F^{[2,0]}_{\tilde X}\right)\qquad\text{with }F:=\Phi^R(\tilde X;K^0,K).
\end{align*}
If $d_M(\tilde X)>L$, define
\[
\Phi^R(\tilde X;K^0,K):=\mathrm{El}(K),\qquad\Phi^E(\tilde X;\mathrm{dat}):=\mathrm{El}(K).
\]
In either case put
\[
\Phi^S(\tilde X;\mathrm{dat}):=\mathrm{Slc}\,\Phi^E(\tilde X;\mathrm{dat})\in\mathrm{Grc}_{k+1}(\tilde U/L).
\]
The maps $\Phi^R$ and $\Phi^E$ take their values in $\mathrm{Grc}_k(\tilde U)$.
All three maps depend on $k$, $\eta_k$, $e_k$, $\alpha$, $L$, $M$, and none of them depends on $\ell$.

\textbf{(7.3) Translations.}
Let $v\in M\mathbb{Z}^3$; this lattice contains $L\cdot M\mathbb{Z}^3$.
The set $\tilde X+v$ is a union of $M$-cubes of the same kind, $(\tilde X+v)^{+\eta_k}=\tilde U+v$, $|\tilde X+v|_M=|\tilde X|_M$, and $d_M(\tilde X+v)=d_M(\tilde X)$, since a translation of $\mathbb{R}^3$ maps nets to nets of the same length.
In the notation of Lemma~3, $T_v\,\mathrm{El}(K)=\mathrm{El}(T_vK)$.
For every matrix $\mathrm{Mx}$,
\[
T_v\left[\int_{\tilde X}\mathrm{d}x\;\bar\psi(x)^T\cdot\mathrm{Mx}\cdot\psi(x)\right]=\int_{\tilde X+v}\mathrm{d}x\;\bar\psi(x)^T\cdot\mathrm{Mx}\cdot\psi(x),
\]
and
\[
\int_{\tilde X+v}\mathrm{d}x\int_{\tilde X+v}\mathrm{d}x'\,\left(T_vK^0_{20}\right)(x,x')=\int_{\tilde X}\mathrm{d}x\int_{\tilde X}\mathrm{d}x'\,K^0_{20}(x,x').
\]
The translation $T_v$ preserves the bidegree, and $T_v\left(F^{[2,0]}_{\tilde X}\right)=\left(T_vF\right)^{[2,0]}_{\tilde X+v}$.
Further, $T_v\circ\boldsymbol{\Theta}_{\tilde X,\tilde B}=\boldsymbol{\Theta}_{\tilde X+v,T_v\tilde B}\circ T_v$.
This is the commutation of the map $(\tilde X,\tilde B)\mapsto\boldsymbol{\Theta}_{\tilde X,\tilde B}$ with the translations by $M\mathbb{Z}^3$ stated in Proposition~17.2\,(f): the cubes of $\tilde X+v$, their centres, their common faces and their shortest chains are the translates of those of $\tilde X$, the labels, $d_X$ and $\operatorname{Vol}$ do not change, and $\left(\tau T_v\tilde B\right)(c+v,c'+v)=(\tau\tilde B)(c,c')$.
Hence, by (7.1),
\begin{gather*}
\Phi^R(\tilde X+v;T_vK^0,T_vK)=T_v\,\Phi^R(\tilde X;K^0,K),\qquad\Phi^E(\tilde X+v;T_v\,\mathrm{dat})=T_v\,\Phi^E(\tilde X;\mathrm{dat}),\\
\Phi^S(\tilde X+v;T_v\,\mathrm{dat})=T_{v/L}\,\Phi^S(\tilde X;\mathrm{dat}).
\end{gather*}

\textbf{(7.4) Tori.}
Let $\ell,\ell/L\in\mathrm{Cl}$, $\mathscr{E}\in K^{x,C}_k(\ell)$, $(\mathscr{A},\mathscr{C})\in\tilde R_{k+1}\times\mathfrak{X}_{k+1}$ and $\mathscr{T}=\mathscr{A}+\mathscr{C}$, and suppose that $X\in D_k$ has a lift $\tilde X$ with $\operatorname{sup-diameter}(\tilde X)+2\eta_k+M<\ell$.
Let $\mathrm{dat}=(K^0,K,\tilde B)$ consist of the antisymmetric kernels of $\mathscr{E}(X;0,0)$ and of $\mathscr{E}(X;\mathscr{A}_L,\mathscr{C}_L)$, restricted to the arguments over $\tilde U$, and of the function $\tilde B$ which assigns to every bond $b$ with both ends in $\tilde X$ the value of $\mathscr{T}_L$ at the image of $b$.
Then
\begin{equation}
\begin{aligned}
&(R\mathscr{E})(X;\mathscr{A}_L,\mathscr{C}_L)=\mathrm{cov}_\ell\left(\Phi^R(\tilde X;K^0,K)\right),\qquad E(X;\mathscr{A}_L,\mathscr{C}_L)=\mathrm{cov}_\ell\left(\Phi^E(\tilde X;\mathrm{dat})\right),\\
&S_L\left[E(X;\mathscr{A}_L,\mathscr{C}_L)\right]=\mathrm{cov}_{\ell/L}\left(\Phi^S(\tilde X;\mathrm{dat})\right).
\end{aligned}
\tag{18b.17}
\end{equation}

\begin{proof}[Proof of (18b.17)]
The family $\mathscr{E}$ has (k2).
By Lemma~3\,(b), $\mathscr{E}(X;\mathscr{A}_L,\mathscr{C}_L)=\mathrm{cov}_\ell\left(\mathrm{El}(K)\right)$ and $\mathscr{E}(X;0,0)=\mathrm{cov}_\ell\left(\mathrm{El}(K^0)\right)$, and the kernels $K_{nm}$, $K^0_{nm}$ vanish unless all $\xi_i$ lie over $\tilde X$ and all $\boldsymbol{\zeta}_j$ are admissible for $\tilde X$.
By Lemma~3\,(c), the polymer $X$ is small if and only if $d_M(\tilde X)\le L$.

Let $X$ be large.
By definition $(R\mathscr{E})(X)=\mathscr{E}(X)$ and $E(X)=(R\mathscr{E})(X)$, while $\Phi^R=\Phi^E=\mathrm{El}(K)$.

Let $X$ be small.
Then $\operatorname{Vol}(X)=M^3\cdot|X|_M=\operatorname{Vol}(\tilde X)$, and $\mathrm{al}_0(\mathscr{E},X)\cdot\operatorname{Vol}(X)=\mathscr{E}_{00}(X;0,0)=K^0_{00}$.
In
\[
\mathrm{al}_2(\mathscr{E},X)=\operatorname{Vol}(X)^{-1}\int\mathrm{d}x\int\mathrm{d}x'\,\mathscr{E}_{20}(X;0,0;x,x')
\]
all arguments range over one lift of $X$ ((I.AP), Section~I.5.6).
When they range over $\tilde X$, the double integral equals $\int_{\tilde X}\mathrm{d}x\int_{\tilde X}\mathrm{d}x'\,K^0_{20}(x,x')$, because the sites over $\tilde X$ represent the sites over $X$ and the weights agree (Lemma~3\,(a), (b)).
For every matrix $\mathrm{Mx}$ one has
\[
\int_X\mathrm{d}x\;\bar\psi(x)^T\cdot\mathrm{Mx}\cdot\psi(x)=\mathrm{cov}_\ell\left(\int_{\tilde X}\mathrm{d}x\;\bar\psi(x)^T\cdot\mathrm{Mx}\cdot\psi(x)\right);
\]
indeed, the summand $\bar\psi(x)^T\cdot\mathrm{Mx}\cdot\psi(x)$ is $\ell\mathbb{Z}^3$-periodic in $c$, the sites over $\tilde X$ represent the sites over $X$ with the same weights $\eta_k^3/48$, and $\mathrm{cov}_\ell$ maps $\psi(x,\beta,\omega)$, $x$ over $\tilde X$, to the generator of the torus with the same index.
In view of the definition of $R$ in Section~I.5.6, this is the first identity.
Put $F:=\Phi^R(\tilde X;K^0,K)$.
Every generator occurring in $F^{[2,0]}$ has its site over $\tilde X$, so that $F^{[2,0]}_{\tilde X}=F^{[2,0]}$.
As $\mathrm{cov}_\ell$ preserves the bidegree (Lemma~3\,(a)), it follows that $(R\mathscr{E})^{[2,0]}(X;\mathscr{A}_L,\mathscr{C}_L)=\mathrm{cov}_\ell\left(F^{[2,0]}\right)$.
By the first sentence of Proposition~17.2\,(f), $\boldsymbol{\Theta}_{X,B}$ is given by one formula of range $0$ with the data $B$ on the bonds with both ends in $\tilde X$; since $\operatorname{sup-diameter}(\tilde X)+M<\ell$, the form that Section~I.1.9 gives to this statement for homomorphisms reads
\[
\mathrm{cov}_\ell\circ\boldsymbol{\Theta}_{\tilde X,\tilde B}=\boldsymbol{\Theta}_{X,\mathscr{T}_L}\circ\mathrm{cov}_\ell.
\]
(By the data of Proposition~17.1, the display (3.3) for $\boldsymbol{\Theta}_{X,B}$ is written with points of $\mathbb{L}^{\mathrm{f}}_k$ over the lift $\tilde X$, with the chains of cubes of $X$, which correspond to the chains of cubes of $\tilde X$ with the same common faces, and with the transports along paths in $\tilde X$; the identity compares two copies of this display on the generators $\psi(x,\cdot\,)$, $x$ over $\tilde X$.)
The total connection of the pair $(\mathscr{A}_L,\mathscr{C}_L)$ is $\mathscr{T}_L$ (Step~0).
Therefore, with the arguments $(\mathscr{A}_L,\mathscr{C}_L)$ suppressed on the left,
\[
E(X;\mathscr{A}_L,\mathscr{C}_L)=(R\mathscr{E})(X)-(R\mathscr{E})^{[2,0]}(X)+\boldsymbol{\Theta}_{X,\mathscr{T}_L}\left((R\mathscr{E})^{[2,0]}(X)\right)=\mathrm{cov}_\ell\left(F-F^{[2,0]}+\boldsymbol{\Theta}_{\tilde X,\tilde B}\left(F^{[2,0]}\right)\right),
\]
which is the second identity.

The third identity is a consequence of the second one and of (18b.16), because
\[
\operatorname{sup-diameter}(\tilde U)+\eta_k\le\operatorname{sup-diameter}(\tilde X)+3\eta_k<\ell,
\]
as $3\eta_k\le2\eta_k+M$.
\end{proof}

\textbf{(7.5) Kernels.}
By (18b.17) and by Lemma~3\,(a) at the level $k$, the antisymmetric kernels of $E(X;\mathscr{A}_L,\mathscr{C}_L)=(R\mathscr{E})^{\mathrm{nat}}(X;\mathscr{A}_L,\mathscr{C}_L)$ coincide, at the arguments over $\tilde U$, with the antisymmetric kernels of $\Phi^E(\tilde X;\mathrm{dat})$, and they vanish at the arguments that are not congruent modulo $\ell\mathbb{Z}^3$ to arguments over $\tilde U$.
Likewise, by (18b.17) and by Lemma~3\,(a) at the level $k+1$, with the side $\ell/L$ and the set $\tilde U/L$, the antisymmetric kernels of $S_L\left[(R\mathscr{E})^{\mathrm{nat}}(X;\mathscr{A}_L,\mathscr{C}_L)\right]$ coincide, at the arguments over $\tilde U/L$, with those of $\Phi^S(\tilde X;\mathrm{dat})$, and they vanish at the arguments that are not congruent modulo $(\ell/L)\mathbb{Z}^3$ to arguments over $\tilde U/L$.
The kernels of $\Phi^S$ are computed from those of $\Phi^E=:\mathrm{El}(K^E)$.
Since $\frac{\eta_k^3}{48}=\frac{L^3\cdot\eta_{k+1}^3}{48}$,
\begin{multline*}
\mathrm{Slc}\,\mathrm{El}(K^E)=\sum_{n,m}\frac{1}{n!\,m!}\sum_{\xi,\boldsymbol{\zeta}\text{ over }\tilde U}\left(\frac{L^3\cdot\eta_{k+1}^3}{48}\right)^{n+2m}K^E_{nm}(\xi;\boldsymbol{\zeta})\\
{}\cdot L^{-n-(1+\alpha)m}\cdot\psi(\xi_1/L)\cdots\psi(\xi_n/L)\cdot\chi(\boldsymbol{\zeta}_1/L)\cdots\chi(\boldsymbol{\zeta}_m/L).
\end{multline*}
By (1) of Step~6, $(\xi,\boldsymbol{\zeta})\mapsto(\xi/L,\boldsymbol{\zeta}/L)$ is a bijection of the indices over $\tilde U$ onto the indices $\xi'$ over $\tilde U/L$ and the indices $\boldsymbol{\zeta}'$ over $\tilde U/L$ whose pairs have $d_X\le M/L$.
Hence, with $L\xi'$, $L\boldsymbol{\zeta}'$ as in (G2.r5),
\[
\left(\Phi^S\right)_{nm}(\xi';\boldsymbol{\zeta}')=\begin{cases}L^{3n+6m}\cdot L^{-n-(1+\alpha)m}\cdot K^E_{nm}(L\xi';L\boldsymbol{\zeta}')&\text{if all pairs of $\boldsymbol{\zeta}'$ have $d_X\le M/L$},\\0&\text{otherwise}.\end{cases}
\]
These functions are antisymmetric, and by Proposition~15\,(a) they are the kernels of $\Phi^S(\tilde X;\mathrm{dat})$.
(On the torus this is Proposition~18.4\,(c).)

\textbf{(7.6) Conclusion.}
By (7.2)--(7.5), the two families of kernels are objects attached to a polymer in the sense of Section~I.1.9, given by the maps $\Phi^E$ and $\Phi^S$, which are the same for all $\ell$ of the class.
All indices that occur lie over $\tilde U$, the $\eta_k$-neighbourhood of $\tilde X$ (respectively over $\tilde U/L$); thus $\mathrm{rg}=\eta_k$, and the polymers are those with $\operatorname{sup-diameter}(\tilde X)+2\,\mathrm{rg}+M<\ell$.
The data are the following: the kernels of $\mathscr{E}(X;0,0)$, of which $K^0_{00}$ and $K^0_{20}$ at the sites over $\tilde X$ enter, and the kernels of $\mathscr{E}(X;\mathscr{A}_L,\mathscr{C}_L)$, both with their arguments over $\tilde U$; the set $\tilde X$ itself, that is, its cubes with their centres and common faces, which enter through $|\tilde X|_M$, $d_M(\tilde X)$, $\int_{\tilde X}$ and the chains of $\boldsymbol{\Theta}_{\tilde X,\tilde B}$; and $\mathscr{T}_L$ on the bonds with both ends in $\tilde X$.

\emph{The lift.}
Every other lift of $X$ is of the form $\tilde X+\ell\cdot w$ with $w\in\mathbb{Z}^3$ (Section~I.1.1).
Since $\mathscr{E}(X;\cdot\,,\cdot\,)$ is even, $K_{nm}$ and $K^0_{nm}$ vanish for odd $n+m$; by the antiperiodicity in each of the $n+m$ arguments, the data belonging to $\tilde X+\ell\cdot w$ are therefore $T_{\ell w}\,\mathrm{dat}$ (recall that $\mathscr{T}_L$ is $\ell\mathbb{Z}^3$-periodic).
The vector $\ell w$ lies in $M\mathbb{Z}^3$.
By (7.3), $\Phi^E(\tilde X+\ell\cdot w;T_{\ell w}\,\mathrm{dat})=T_{\ell w}\,\Phi^E(\tilde X;\mathrm{dat})$; moreover, $\mathrm{cov}_\ell\circ T_{\ell w}$ multiplies every generator by $(-1)^{w_1+w_2+w_3}$ (Section~I.1.9), and therefore agrees with $\mathrm{cov}_\ell$ on the even element $\Phi^E(\tilde X;\mathrm{dat})$.
At the level $k+1$, (7.3) gives $\Phi^S(\tilde X+\ell\cdot w;T_{\ell w}\,\mathrm{dat})=T_{(\ell/L)w}\,\Phi^S(\tilde X;\mathrm{dat})$, and $\mathrm{cov}_{\ell/L}\circ T_{(\ell/L)w}$ multiplies every generator by $(-1)^{w_1+w_2+w_3}$, hence agrees with $\mathrm{cov}_{\ell/L}$ on the even element $\Phi^S(\tilde X;\mathrm{dat})$.
Consequently the right-hand sides of (18b.17) do not depend on the choice of the lift.

\textbf{(7.7) The assertions of (f) on the class of data, on the maps $\mathrm{cov}_\ell$, $\mathrm{cov}_{\ell/L}$ and on the two margins.}
Let $\tilde X$ and $\tilde U=\tilde X^{+\eta_k}$ be as in (7.2); then $\operatorname{sup-diameter}(\tilde U)\le\operatorname{sup-diameter}(\tilde X)+2\eta_k$.

\emph{Kernels; the data.}
The second assertion of (f), the one on kernels, is established by (7.5), (7.6).
The kernels $K_{20}$, $K^0_{20}$ of the data in (7.4) vanish unless both arguments lie over $\tilde X$; this follows from Lemma~3\,(b) and is the clause (k2) for $\mathscr{E}$.
For data of this kind one has $F^{[2,0]}_{\tilde X}=F^{[2,0]}$ in (7.2), as in the proof of (18b.17).

\emph{The margin $L\cdot M$.}
Section~I.1.9 defines $\mathrm{cov}_{\ell/L}$ on its algebra over $\tilde U/L$, with all generators, provided that $\operatorname{sup-diameter}(\tilde U/L)+M<\ell/L$, that is, $\operatorname{sup-diameter}(\tilde U)+L\cdot M<\ell$; this is the margin required for $S_L$ in Section~I.1.9 and in Proposition~25.2\,(o7).
The arguments of (7.1)--(7.6) do not need this margin.
Indeed, by Lemma~3\,(a) and (18b.16), the rule of Section~I.1.9 defines, on the subalgebras generated by $\psi$, $\chi$, injective homomorphisms $\mathrm{cov}_\ell$, $\mathrm{cov}_{\ell/L}$ with $\mathrm{cov}_{\ell/L}\circ\mathrm{Slc}=S_L\circ\mathrm{cov}_\ell$ as soon as $\operatorname{sup-diameter}(\tilde U)+\eta_k<\ell$, and this inequality holds under the margin $\operatorname{sup-diameter}(\tilde X)+2\eta_k+M<\ell$ of (f), as shown at the end of (7.4).
Thus the margin $L\cdot M$ is needed only in order that $\mathrm{cov}_{\ell/L}$ be the map of Section~I.1.9 on all of $\mathrm{Grc}_{k+1}(\tilde U/L)$, that is, for the pairs with $M/L<d_X\le M$; by (1) of Step~6 these pairs do not occur in the image of $\mathrm{Slc}$.
(For the pairs with $d_X\le M/L$, the assertion of Section~I.1.1 on the only representative, applied with $\mathrm{sd}=M/L$, requires $\operatorname{sup-diameter}(\tilde U/L)+M/L<\ell/L$, that is, $\operatorname{sup-diameter}(\tilde U)+M<\ell$, and this is provided by the margin of (f).)
Under the margin $L\cdot M$, the maps $\mathrm{cov}_\ell$ and $\mathrm{cov}_{\ell/L}$ are the restrictions of the maps of Section~I.1.9, by (7.1).

\emph{Small $X$.}
By (G2.r3), a small polymer $X$ has a lift $\tilde X$ with $\operatorname{sup-diameter}(\tilde X)\le M\cdot(L+2)$.
Since $\eta_k\le M\le L\cdot M$ and $\ell\ge13L\cdot M$,
\[
\operatorname{sup-diameter}(\tilde X)+2\eta_k+L\cdot M\le M\cdot(L+2)+2M+L\cdot M=2L\cdot M+4M\le6L\cdot M<\ell.
\]
The left-hand sides of the two margins, $\operatorname{sup-diameter}(\tilde X)+2\eta_k+M$ and $\operatorname{sup-diameter}(\tilde U)+L\cdot M$, are both at most $\operatorname{sup-diameter}(\tilde X)+2\eta_k+L\cdot M$.
Hence both margins hold.

\emph{The admissible sets of the Propositions~27.m.}
An admissible set $\tilde X$ in the sense of the Definitions of Proposition~27.1 satisfies $\operatorname{diam}(\tilde X)\le\ell_1-14\,d_U\cdot M$ with $d_U=2L+1$; here $\operatorname{diam}(\tilde X)$ is the largest side of $\operatorname{box}(\tilde X)$, the smallest box (a product of three intervals that is a union of $M$-cubes) containing $\tilde X$.
The sides in the Propositions~27.m are $\ell_1$ and $\ell_2=\mathrm{Jr}\cdot\ell_1\ge\ell_1$ (the standing data (TT)).
Two sites of a box differ in each coordinate by less than its side; hence $\operatorname{sup-diameter}(\tilde X)<\operatorname{diam}(\tilde X)$, and both for $\ell=\ell_1$ and for $\ell=\ell_2$
\[
\operatorname{sup-diameter}(\tilde X)+2\eta_k+L\cdot M<\ell_1-14\cdot(2L+1)\cdot M+2M+L\cdot M<\ell_1\le\ell.
\]
As for small $X$, both margins hold.

\paragraph{Step 8 (level 0; uniformity)}
\emph{Level $0$.}
At $k=0$ one has $\eta_0=1$, $\mathbb{L}_0=\mathbb{Z}^3$, $T^{\mathrm{f}}_0=T_0\times O_{\mathrm{cub}}$, $\nu(\xi)=\frac{1}{48}$, $\nu(\boldsymbol{\zeta})=\left(\frac{1}{48}\right)^2$ and $\int\mathrm{d}x=\frac{1}{48}\sum_c\sum_g$; the charge $e_k$ of the level $k=0$ is the number denoted $q$ in the flow (Section~I.2.5).
At the level $1$ one has $\eta_1=1/L$, $\nu(\xi)=\frac{L^{-3}}{48}$ and $\nu(\boldsymbol{\zeta})=\left(\frac{L^{-3}}{48}\right)^2$.
Every pair of $\mathrm{Pz}_0$ has $d_X\ge1$; a pair inside a cluster has $d_X=1$, and $S_L$ maps a leg at such a pair to a leg at a pair inside a cluster of the level $1$, with $d_X=\eta_1=1/L$.
In Steps~0--7 and in Lemmas~1 and~3 the level $k$ enters only through $\eta_{k+1}=\eta_k/L$ and $e_{k+1}=L^{\frac12}e_k$ and, in Lemma~3 and in Step~7, through $\eta_k\le1\le M$ and the oddness of $M/\eta_k=M\cdot L^k$ (Section~4, row~8).
These relations hold for every $k\ge0$, because $\eta_k=L^{-k}$, $L$ is odd and $M$ is a power of $L$; and every statement invoked in these steps and lemmas holds for all $k\ge0$.
Hence Steps~0--7 remain valid without any change at $k=0$.
(In the flow the activity of the level $0$ is $\mathscr{E}=0$ by (I.L0.21), so that at $k=0$ the proposition is applied to $\mathscr{E}=0$ only.)

\emph{Uniformity.}
The numbers $c_{\mathrm{rb}}$ and $c_{\mathrm{fl}}$ are constants of the line (C0) of Section~I.8.2, and the exponents contain $\alpha$ and $\varepsilon$ only.
The level $k$, the side $\ell$ and the charge $e_k$ enter Steps~0--6 through the statements invoked there, which hold for all data of the setting (S-2) of Section~G2.0.4, and through the ratios
\[
\frac{\eta_{k+1}}{\eta_k}=\frac{1}{L},\qquad\frac{h_{k+1,1}}{h_{k,1}}=L^{-\frac18},\qquad\frac{h_{k+1,2}}{h_{k,2}}=L^{-\frac18+\frac{\varepsilon}{2}},\qquad\frac{w^{\mathrm{E}}_{k+1}}{w^{\mathrm{E}}_k}=L^{\frac18-\frac{7\varepsilon}{2}}.
\]
Therefore (b) and (d) hold with the same numbers for all $k\ge0$, all $\ell$ with $\ell,\ell/L\in\mathrm{Cl}$ and all $e_k\in\left(0,L^{-\frac12}e_0\right]$.
The data of the setting (S-2) are $k$, $\ell$ and $e_k$ alone.
The maps of (f) do not depend on $\ell$.
\qed

\subsection*{3. Proposition 18.6}

\subsubsection*{Proof of Proposition 18.6}
The clauses (a)--(d), (f) and (g) concern the level $k$ alone.

\emph{The kernel.}
For an $M$-cube $\square$ of $T_k$ and for $x=(c,g)$, $x'=(c',g')$ in $\mathbb{L}^{\mathrm{f}}_k$ put
\[
K^{\mathrm{Ng}}_{\square}(x,x'):=\begin{cases}\eta_k^{-3}\cdot(-1)^{w_1+w_2+w_3}\cdot\mathrm{Ng}(g,g')&\text{if the image of $c$ in $T_k$ lies in $\square$}\\&\text{and $c'=c+\ell\cdot w$ with $w\in\mathbb{Z}^3$},\\0&\text{otherwise}.\end{cases}
\]
This function is antiperiodic in each argument, since a shift of $c$ or of $c'$ by $\ell\cdot e_\mu$ changes $w$ by $\mp e_\mu$.
For $x$, $x'$ over one lift of $\square$ it equals $\eta_k^{-3}\cdot\mathbf{1}[c=c'\in\square]\cdot\mathrm{Ng}(g,g')$, because two sites of one lift that are congruent modulo $\ell\mathbb{Z}^3$ coincide.
In $\mathrm{Qd}\left[K^{\mathrm{Ng}}_{\square}\right]$ (Lemma~2) the sites $c$ and $c'$ range over one fundamental domain, and there $K^{\mathrm{Ng}}_{\square}(x,x')\ne0$ only if $c'=c$.
Therefore, with $x=(c,g)$,
\begin{equation}
\begin{split}
\mathrm{Qd}\left[K^{\mathrm{Ng}}_{\square}\right]&=\eta_k^6\cdot\left(\frac{1}{48}\right)^2\sum_{c\in\square}\sum_{g,g'\in O_{\mathrm{cub}}}\bar\psi(c,g)^T\cdot\eta_k^{-3}\cdot\mathrm{Ng}(g,g')\cdot\psi(c,g')\\
&=\int_{\square}\mathrm{d}x\;\frac{1}{48}\sum_{g'}\bar\psi(c,g)^T\cdot\mathrm{Ng}(g,g')\cdot\psi(c,g')=E^{\mathrm{Ng}}(\square).
\end{split}
\tag{18b.18}
\end{equation}

\paragraph{(a)}
\emph{Kernel and norm.}
By (18b.18) and Lemma~2\,(a), $E^{\mathrm{Ng}}(\square)$ is an even, balanced element of bidegree $(2,0)$, its only kernel is $E^{\mathrm{Ng}}_{20}(\square;x,x')=K^{\mathrm{Ng}}_{\square}(x,x')$, and
\begin{align*}
\left\|E^{\mathrm{Ng}}(\square)\right\|_{\mathsf{w}}&=\mathsf{w}_1^2\sum_{\beta,\beta'}\int\mathrm{d}x\int\mathrm{d}x'\,\left|K^{\mathrm{Ng}}_{\square}(x,x')_{\beta\beta'}\right|\\
&=\mathsf{w}_1^2\cdot\eta_k^3\sum_{c\in\square}\left(\frac{1}{48}\right)^2\sum_{g,g'}\sum_{\beta,\beta'}\left|\mathrm{Ng}(g,g')_{\beta\beta'}\right|=\mathsf{w}_1^2\cdot M^3\cdot|\mathrm{Ng}|_1,
\end{align*}
since $\square$ contains $(M/\eta_k)^3$ photon sites ($M/\eta_k=M\cdot L^k$ is odd; see the proof of Lemma~3\,(c)).
For $\mathsf{w}=h_k$ one has $h_{k,1}^2=e_k^{-\frac12}$.
The elements $E^{\mathrm{Ng}}(X)$ do not depend on $(\mathscr{A},\mathscr{C})$.
Since $d_M(\square)=0$ by \cite[Theorem~13.2\,(t0)]{QEDI} and $E^{\mathrm{Ng}}(X)=0$ for $|X|_M\ge2$,
\[
\left\|E^{\mathrm{Ng}}\right\|^x_k=\sup_X\exp\left(\kappa\cdot d_M(X)\right)\cdot\left\|E^{\mathrm{Ng}}(X)\right\|_{h_k}=M^3\cdot e_k^{-\frac12}\cdot|\mathrm{Ng}|_1<\infty;
\]
this is (k6).
\begin{itemize}
\item[(k1)] The kernels are constant functions of $(\mathscr{A},\mathscr{C})$.
\item[(k2)] $K^{\mathrm{Ng}}_{\square}(x,x')$ vanishes unless $x$ and $x'$ lie over $\square$; there is no leg $\chi$, and no value of a bond function enters.
\item[(k3)] Let $\tilde g=(S,a)\in\mathbb{G}_k$ and $r=r(S)$.
The map $\tilde g$ sends $M$-cubes onto $M$-cubes, and $|\tilde gX|_M=|X|_M$ (Proposition~16.1\,(a1)); hence both sides of (k3) vanish if $|X|_M\ge2$.
For $X=\square$, the identities (18b.18) and (18b.1) give $a_{\tilde g}\left[E^{\mathrm{Ng}}(\tilde g\square)\right]=\mathrm{Qd}\left[\left(K^{\mathrm{Ng}}_{\tilde g\square}\right)^{\tilde g}\right]$ with
\[
\left(K^{\mathrm{Ng}}_{\tilde g\square}\right)^{\tilde g}(x,x')=S^{-1}\cdot K^{\mathrm{Ng}}_{\tilde g\square}(\tilde gx,\tilde gx')\cdot S,\qquad\tilde gx=(rc+a,rg),\qquad\tilde gx'=(rc'+a,rg').
\]
The image of $rc+a$ lies in $\tilde g\square$ if and only if the image of $c$ lies in $\square$; one has $c'=c+\ell\cdot w$ if and only if $rc'+a=(rc+a)+\ell\cdot rw$; and the sum of the entries of $rw$ has the parity of $w_1+w_2+w_3$.
Hence $\left(K^{\mathrm{Ng}}_{\tilde g\square}\right)^{\tilde g}(x,x')=\eta_k^{-3}\cdot(-1)^{w_1+w_2+w_3}\cdot S^{-1}\cdot\mathrm{Ng}(rg,rg')\cdot S$ in the first case of the definition of the kernel, and $\left(K^{\mathrm{Ng}}_{\tilde g\square}\right)^{\tilde g}(x,x')=0$ otherwise.
By (n1), $S^{-1}\cdot\mathrm{Ng}(rg,rg')\cdot S=\mathrm{Ng}(g,g')$.
Consequently $a_{\tilde g}\left[E^{\mathrm{Ng}}(\tilde g\square)\right]=\mathrm{Qd}\left[K^{\mathrm{Ng}}_{\square}\right]=E^{\mathrm{Ng}}(\square)$.
\item[(k4)] By (18b.2), $G_\Lambda E^{\mathrm{Ng}}(\square)=\mathrm{Qd}\left[\left(K^{\mathrm{Ng}}_{\square}\right)^{\Lambda}\right]$, where $\left(K^{\mathrm{Ng}}_{\square}\right)^{\Lambda}$ carries the factor $\exp\left\{i\,e_k\cdot\left(\Lambda(c')-\Lambda(c)\right)\right\}$.
As $\Lambda$ is a function on $T_k$ and $K^{\mathrm{Ng}}_{\square}(x,x')\ne0$ only if $c'$ and $c$ have the same image, this factor equals $1$.
Hence $G_{\boldsymbol{\lambda}+\boldsymbol{\lambda}'}\left[E^{\mathrm{Ng}}(X;\mathscr{A}^1,\mathscr{C}^1)\right]=E^{\mathrm{Ng}}(X;\mathscr{A},\mathscr{C})$ for all real $\boldsymbol{\lambda}$, $\boldsymbol{\lambda}'$.
\item[(k5)] By (18b.3), $\mathrm{Cg}\,E^{\mathrm{Ng}}(\square)=\mathrm{Qd}\left[\left(K^{\mathrm{Ng}}_{\square}\right)^{\mathrm{Cg}}\right]$, where
\[
\left(K^{\mathrm{Ng}}_{\square}\right)^{\mathrm{Cg}}(x,x')=\mathsf{C}^{-1}\cdot K^{\mathrm{Ng}}_{\square}(x',x)^T\cdot\mathsf{C}=\eta_k^{-3}\cdot(-1)^{w_1+w_2+w_3}\cdot\mathsf{C}^{-1}\cdot\mathrm{Ng}(g',g)^T\cdot\mathsf{C}
\]
in the first case (then $c=c'-\ell\cdot w$, the parity of $-w$ is that of $w$, and the images of $c$ and $c'$ coincide), and $\left(K^{\mathrm{Ng}}_{\square}\right)^{\mathrm{Cg}}(x,x')=0$ otherwise.
By (n2) this is $K^{\mathrm{Ng}}_{\square}(x,x')$.
\end{itemize}
Thus $E^{\mathrm{Ng}}$ lies in $K^{x,C}_k$.
\begin{itemize}
\item[(k7)] Assume that (n3) holds.
By (18b.4), $\mathbf{K}\,E^{\mathrm{Ng}}(\square)=\mathrm{Qd}\left[\left(K^{\mathrm{Ng}}_{\square}\right)^{\mathbf{K}}\right]$ with
\[
\left(K^{\mathrm{Ng}}_{\square}\right)^{\mathbf{K}}(x,x')=\eta_k^{-3}\cdot(-1)^{w_1+w_2+w_3}\cdot\varkappa^{-1}\cdot\overline{\mathrm{Ng}(g,g')}\cdot\varkappa
\]
in the first case, and $\left(K^{\mathrm{Ng}}_{\square}\right)^{\mathbf{K}}(x,x')=0$ otherwise; by (n3) this is $K^{\mathrm{Ng}}_{\square}(x,x')$.
Hence $E^{\mathrm{Ng}}$ lies in $K^x_k$.
\end{itemize}

\emph{The sum.}
The sets $\square\times O_{\mathrm{cub}}$ partition $T^{\mathrm{f}}_k$ (Section~I.1.4).
With $(\mathrm{Ocl}\,\psi)(c,g):=\frac{1}{48}\sum_{g'}\mathrm{Ng}(g,g')\cdot\psi(c,g')$ one obtains
\[
\sum_{\square}E^{\mathrm{Ng}}(\square)=\int\mathrm{d}x\;\bar\psi(x)^T\cdot(\mathrm{Ocl}\,\psi)(x)=\left\langle\bar\psi,\mathrm{Ocl}\,\psi\right\rangle^{\mathrm{ch}}.
\]

\paragraph{(b)}
The cube $\square$ is a small polymer, because $d_M(\square)=0\le L$.
Since $E^{\mathrm{Ng}}(\square)$ has no part of bidegree $(0,0)$, one has $\mathrm{al}_{0}(E^{\mathrm{Ng}},\square)=0$.
With both arguments over one lift of $\square$, one obtains
\begin{align*}
\mathrm{al}_{2}(E^{\mathrm{Ng}},\square)
&=\operatorname{Vol}(\square)^{-1}\int_{\square}\mathrm{d}x\int_{\square}\mathrm{d}x'\,E^{\mathrm{Ng}}_{20}(\square;x,x')\\
&=M^{-3}\cdot\eta_k^{6}\cdot\left(\frac{1}{48}\right)^{2}\sum_{c,c'\in\square}\ \sum_{g,g'}\eta_k^{-3}\cdot\mathbf{1}[c=c']\cdot\mathrm{Ng}(g,g')\\
&=M^{-3}\cdot\eta_k^{3}\cdot\left(\frac{M}{\eta_k}\right)^{3}\left(\frac{1}{48}\right)^{2}\sum_{g,g'}\mathrm{Ng}(g,g')=\mathrm{Nav}.
\end{align*}
Let $S\in\mathrm{Sp}$ and $r=r(S)\in\{\pm1\}^3$.
By (n1) and (sp8),
\[
S\cdot\mathrm{Nav}\cdot S^{-1}=\left(\frac{1}{48}\right)^{2}\sum_{g,g'}\mathrm{Ng}(rg,rg')=\mathrm{Nav}.
\]
Hence $\mathrm{Nav}$ commutes with $\gamma_1$, $\gamma_2$, $\gamma_3$ and $\gamma_5$, and (I.Sch) gives $\mathrm{Nav}=\mathrm{na}\cdot1$ with a complex number $\mathrm{na}$.
If (n3) holds, then
\[
\mathrm{na}\cdot1=\mathrm{Nav}=\varkappa^{-1}\cdot\overline{\mathrm{Nav}}\cdot\varkappa=\overline{\mathrm{na}}\cdot1,
\]
so that $\mathrm{na}$ is real.
For a small polymer $X$ with $|X|_M\ge2$ one has $E^{\mathrm{Ng}}(X)=0$, and therefore $\mathrm{al}_{0}(E^{\mathrm{Ng}},X)=0$ and $\mathrm{al}_{2}(E^{\mathrm{Ng}},X)=0$.
Among the polymers $X$ with $|X|_M=1$, the only one that contains $\square$ is $\square$ itself.
The definitions in the data of Proposition~16.3 thus give
\[
\varepsilon(E^{\mathrm{Ng}},\square)=-\mathrm{al}_{0}(E^{\mathrm{Ng}},\square)=0,\qquad m(E^{\mathrm{Ng}},\square)=-\mathrm{al}_{2}(E^{\mathrm{Ng}},\square)=-\mathrm{na}\cdot1,
\]
that is, $\varepsilon(E^{\mathrm{Ng}})=0$ and $m(E^{\mathrm{Ng}})=-\mathrm{na}$, in agreement with Proposition~16.3\,(a).

\paragraph{(c)}
The operator $\mathrm{Ocl}$ maps $\mathfrak{S}^{\mathrm{ch}}(\eta_k,\ell)$ into itself, since it acts only on the label and on the spinor index.
Let $\xi\in\mathfrak{S}(\eta_k,\ell)$.
By (I.3.1),
\begin{align*}
\left(\mathrm{Ocl}\,J\xi\right)(c,g)&=\frac{1}{48}\sum_{g'}\mathrm{Ng}(g,g')\cdot(J\xi)(c,g')\\
&=\frac{1}{48}\sum_{g'}\mathrm{Ng}(g,g')\cdot\xi(c)=\mathrm{Nl}(g)\cdot\xi(c),
\end{align*}
and by (b)
\[
\left(J^{*}\mathrm{Ocl}\,J\xi\right)(c)=\frac{1}{48}\sum_{g}\mathrm{Nl}(g)\cdot\xi(c)=\mathrm{Nav}\cdot\xi(c)=\mathrm{na}\cdot\xi(c).
\]
As $\mathrm{Phv}=1-JJ^{*}$ by (I.J4), it follows that
\[
\left(\mathrm{Phv}\,\mathrm{Ocl}\,J\xi\right)(c,g)=\mathrm{Nl}(g)\cdot\xi(c)-\left(J(\mathrm{na}\cdot\xi)\right)(c,g)=\left(\mathrm{Nl}(g)-\mathrm{na}\right)\cdot\xi(c).
\]
By (b) the extracted number is $m(E^{\mathrm{Ng}})=-\mathrm{na}$, that is, $-J^{*}\mathrm{Ocl}\,J$.
The function $\mathrm{Nl}(\cdot)-\mathrm{na}$ enters neither $\varepsilon(E^{\mathrm{Ng}})$ nor $m(E^{\mathrm{Ng}})$.

\paragraph{(d)}
By the definition of $R$ in Section~I.5.6 and by (b),
\begin{align*}
\left(R E^{\mathrm{Ng}}\right)(\square)
&=E^{\mathrm{Ng}}(\square)-\mathrm{al}_{0}(E^{\mathrm{Ng}},\square)\cdot\operatorname{Vol}(\square)\\
&\quad-\int_{\square}\mathrm{d}x\,\bar{\psi}(x)^T\,\mathrm{al}_{2}(E^{\mathrm{Ng}},\square)\,\psi(x)\\
&=E^{\mathrm{Ng}}(\square)-\mathrm{na}\int_{\square}\mathrm{d}x\,\bar{\psi}(x)^T\psi(x).
\end{align*}
The last integral can be written in the form of $E^{\mathrm{Ng}}(\square)$:
\[
\int_{\square}\mathrm{d}x\,\bar{\psi}(c,g)^T\psi(c,g)=\int_{\square}\mathrm{d}x\,\frac{1}{48}\sum_{g'}\bar{\psi}(c,g)^T\left(48\cdot\mathbf{1}[g=g']\right)\psi(c,g').
\]
In terms of kernels: by Proposition~16.2\,(iv) the kernel of the left side is $48\eta_k^{-3}\cdot\mathbf{1}[x=x'\text{ over }\square]$, and this equals $\eta_k^{-3}\cdot\mathbf{1}[c=c'\in\square]\cdot48\cdot\mathbf{1}[g=g']$.
Since $E^{\mathrm{Ng}}(\square)$ is linear in $\mathrm{Ng}$, one obtains $\left(R E^{\mathrm{Ng}}\right)(\square)=E^{\mathrm{Ng}'}(\square)$ with
\[
\mathrm{Ng}'(g,g')=\mathrm{Ng}(g,g')-48\cdot\mathbf{1}[g=g']\cdot\mathrm{na}.
\]
For $|X|_M\ge2$ one has $\left(R E^{\mathrm{Ng}}\right)(X)=0=E^{\mathrm{Ng}'}(X)$: for small $X$ all three terms in the definition of $R$ vanish by (b), and for large $X$ one has $\left(R E^{\mathrm{Ng}}\right)(X)=E^{\mathrm{Ng}}(X)=0$.
Therefore $R E^{\mathrm{Ng}}=E^{\mathrm{Ng}'}$.

Let
\[
\mathrm{Ng}_{\mathrm{na}}(g,g'):=48\cdot\mathbf{1}[g=g']\cdot\mathrm{na}\cdot1.
\]
This function has (n1), because $S\cdot\mathrm{na}\cdot S^{-1}=\mathrm{na}$ and $\mathbf{1}[rg=rg']=\mathbf{1}[g=g']$;
it has (n2), because $\mathsf{C}^{-1}\left(\mathrm{na}\cdot1\right)^T\mathsf{C}=\mathrm{na}\cdot1$ and $\mathbf{1}[g'=g]=\mathbf{1}[g=g']$;
and it has (n3) if $\mathrm{na}$ is real, which by (b) is the case when $\mathrm{Ng}$ has (n3).
The conditions (n1) and (n2) are complex linear, and (n3) is a real linear condition on the function.
Consequently $\mathrm{Ng}'=\mathrm{Ng}-\mathrm{Ng}_{\mathrm{na}}$ has (n1) and (n2), and it has (n3) if $\mathrm{Ng}$ has.
For the quantities formed with $\mathrm{Ng}'$ one finds
\[
\mathrm{Nav}'=\mathrm{Nav}-\left(\frac{1}{48}\right)^{2}\sum_{g}48\cdot\mathrm{na}\cdot1=\mathrm{na}\cdot1-\mathrm{na}\cdot1=0,
\]
and moreover $\mathrm{Nl}'(g)=\mathrm{Nl}(g)-\mathrm{na}$ and $\mathrm{Nr}'(g)=\mathrm{Nr}(g)-\mathrm{na}$;
here $\mathrm{Nr}(g):=\frac{1}{48}\sum_{g''\in O_{\mathrm{cub}}}\mathrm{Ng}(g'',g)$, and $\mathrm{Nr}'$ is formed in the same way with $\mathrm{Ng}'$.

\paragraph{(e)}
Let $E:=\left(R E^{\mathrm{Ng}}\right)^{\mathrm{nat}}$.
By (a), $E^{\mathrm{Ng}}$ belongs to $K^{x,C}_{k}$, with $\left\|E^{\mathrm{Ng}}\right\|^{x}_{k}=M^3\cdot e_k^{-\frac{1}{2}}\cdot|\mathrm{Ng}|_1$;
the cube $\square$ is small, and $d_M(\square)=0$.
Proposition~18.5\,(d), applied with $\mathscr{E}=E^{\mathrm{Ng}}$ and $X=\square$, yields
\begin{equation}\tag{18b.19}
\begin{aligned}
&E_{20}(\square;0,0;\cdot,\cdot)=0;\\
&\sup_{\tilde{R}_k\times\mathfrak{X}_k}\left\|E(\square;\cdot,\cdot)\right\|_{h_k}\le20\left\|E^{\mathrm{Ng}}\right\|^{x}_{k}=20M^3\cdot e_k^{-\frac{1}{2}}\cdot|\mathrm{Ng}|_1;\\
&\left\|S_L\left[E(\square;\mathscr{A}_L,\mathscr{C}_L)\right]\right\|_{h_{k+1}}\le3L^{-\frac{13}{4}+\frac{3\varepsilon}{2}}\cdot20M^3\cdot e_k^{-\frac{1}{2}}\cdot|\mathrm{Ng}|_1\\
&\qquad=60L^{-\frac{13}{4}+\frac{3\varepsilon}{2}}\cdot M^3\cdot e_k^{-\frac{1}{2}}\cdot|\mathrm{Ng}|_1\qquad\text{for }(\mathscr{A},\mathscr{C})\in\tilde{R}_{k+1}\times\mathfrak{X}_{k+1}.
\end{aligned}
\end{equation}

\emph{H\"older legs.}
By (a) and (d), $R E^{\mathrm{Ng}}=E^{\mathrm{Ng}'}$ is balanced and has (k2), and its only kernel is the one of bidegree $(2,0)$.
Proposition~17.2\,(b) gives $E_{nm}=\left(R E^{\mathrm{Ng}}\right)_{nm}=0$ for $(n,m)\notin\{(2,0),(1,1),(0,2)\}$.
Together with $E_{20}(\square;0,0;\cdot,\cdot)=0$ this leaves
\[
E(\square;0,0)=E^{[1,1]}(\square;0,0)+E^{[0,2]}(\square;0,0).
\]
Hence every basis monomial with a non-zero coefficient in $E(\square;0,0)$ contains at least one generator $\chi(\boldsymbol{\zeta})$.

\emph{The element written out.}
The element $\left(R E^{\mathrm{Ng}}\right)(\square)=E^{\mathrm{Ng}'}(\square)$ is of bidegree $(2,0)$ and does not depend on the connections.
The definition of the adjustment therefore gives
\[
E(\square;\mathscr{A}',\mathscr{C}')=\boldsymbol{\Theta}_{\square,\mathscr{T}'}\left(E^{\mathrm{Ng}'}(\square)\right),\qquad\mathscr{T}'=\mathscr{A}'+\mathscr{C}'.
\]
Let $z$ be the centre of $\square$; all points below lie over one lift of $\square$.
For $X=\square$ the only chain of Proposition~17.1\,(c) from the cube of $c$ to the cube of $\tilde{c}$ is $(\square)$;
for it $n=0$, $p_0=x=(c,g)$, $p_1=(z,g)$ and $p_2=\tilde{x}=(\tilde{c},\tilde{g})$.
At $B=0$ all transports are equal to $1$, so that $\mathrm{tm}=\mathrm{tm}^{-}=1$.
With $\operatorname{Vol}(\square)=M^3$ and $d_X\left((c,g),(z,g)\right)=|c-z|$ for $c\ne z$, the formulas (3.2), (3.3) of Proposition~17.1 give
\begin{equation}\tag{18b.20}
\begin{aligned}
&\boldsymbol{\Theta}_{\square,0}\,\psi(x)=\psi_{\mathrm{av}}+\mathrm{hp}(x),\qquad\boldsymbol{\Theta}_{\square,0}\,\bar{\psi}(x)=\bar{\psi}_{\mathrm{av}}+\bar{\mathrm{hp}}(x),\\
&\psi_{\mathrm{av}}=M^{-3}\int_{\square}\mathrm{d}\tilde{x}\,\psi(\tilde{x}),\qquad\bar{\psi}_{\mathrm{av}}=M^{-3}\int_{\square}\mathrm{d}\tilde{x}\,\bar{\psi}(\tilde{x}),\\
&\mathrm{hp}(c,g):=-\mathbf{1}[c\ne z]\cdot|c-z|^{\alpha}\cdot\chi\left((c,g),(z,g)\right)\\
&\qquad\qquad-M^{-3}\int_{\square}\mathrm{d}\tilde{x}\,\mathbf{1}\left[\tilde{x}\ne(z,g)\right]\cdot d_X\left((z,g),\tilde{x}\right)^{\alpha}\cdot\chi\left((z,g),\tilde{x}\right),
\end{aligned}
\end{equation}
and $\bar{\mathrm{hp}}(c,g)$ is given by the same expression with $\bar{\chi}$ in the place of $\chi$.
Here $\int_{\square}\mathrm{d}\tilde{x}=\eta_k^3\cdot\frac{1}{48}\sum_{\tilde{c}\in\square}\sum_{\tilde{g}\in O_{\mathrm{cub}}}$.
The $47$ terms with $\tilde{c}=z$ and $\tilde{g}\ne g$ are legs inside the cluster of $z$, with $d_X=\eta_k$, whereas the first leg joins two sites with the same label $g$.
As $\boldsymbol{\Theta}_{\square,0}$ is a homomorphism, (18b.18) for $\mathrm{Ng}'$ gives
\begin{equation}\tag{18b.21}
\begin{aligned}
E(\square;0,0)&=\int_{\square}\mathrm{d}x\,\frac{1}{48}\sum_{g'}\left(\bar{\psi}_{\mathrm{av}}+\bar{\mathrm{hp}}(c,g)\right)^T\mathrm{Ng}'(g,g')\left(\psi_{\mathrm{av}}+\mathrm{hp}(c,g')\right)\\
&=M^3\cdot\bar{\psi}_{\mathrm{av}}^T\,\mathrm{Nav}'\,\psi_{\mathrm{av}}+\int_{\square}\mathrm{d}x\,\bar{\mathrm{hp}}(c,g)^T\left(\mathrm{Nl}(g)-\mathrm{na}\right)\psi_{\mathrm{av}}\\
&\quad+\int_{\square}\mathrm{d}x\,\bar{\psi}_{\mathrm{av}}^T\left(\mathrm{Nr}(g)-\mathrm{na}\right)\mathrm{hp}(c,g)\\
&\quad+\int_{\square}\mathrm{d}x\,\frac{1}{48}\sum_{g'}\bar{\mathrm{hp}}(c,g)^T\,\mathrm{Ng}'(g,g')\,\mathrm{hp}(c,g').
\end{aligned}
\end{equation}
For the second term note that $\frac{1}{48}\sum_{g'}\mathrm{Ng}'(g,g')=\mathrm{Nl}(g)-\mathrm{na}$.
For the third term,
\[
\eta_k^3\sum_{c}\left(\frac{1}{48}\right)^{2}\sum_{g,g'}\bar{\psi}_{\mathrm{av}}^T\,\mathrm{Ng}'(g,g')\,\mathrm{hp}(c,g')=\eta_k^3\sum_{c}\frac{1}{48}\sum_{g'}\bar{\psi}_{\mathrm{av}}^T\,\mathrm{Nr}'(g')\,\mathrm{hp}(c,g'),
\]
after which $g'$ is renamed $g$.
At $B=0$ the first term of (18b.21) coincides with the first term of Proposition~17.1\,(f); it is zero, since $\mathrm{Nav}'=0$ by (d).
In the second term the coefficient $\mathrm{Nl}(g)-\mathrm{na}$ is the coupling found in (c).
By (n2),
\[
\mathrm{Nr}(g)=\frac{1}{48}\sum_{g''}\mathrm{Ng}(g'',g)=\frac{1}{48}\sum_{g''}\mathsf{C}^{-1}\,\mathrm{Ng}(g,g'')^T\,\mathsf{C}=\mathsf{C}^{-1}\,\mathrm{Nl}(g)^T\,\mathsf{C}.
\]

\emph{All connections.}
Let $B$ be a complex bond function on $T_k$; again all points lie over one lift of $\square$.
By Proposition~17.1\,(b), $u_B(p,p')=U_B\left(\pi(p),\pi(p')\right)$ for $(p,p')\in\mathrm{Pzt}_k$, and $u_B(p,p)=1=U_B\left(\pi(p),\pi(p)\right)$;
thus the transports depend on the photon sites only, by (I.1.6).
Again $(\square)$ is the only chain for $X=\square$.
Hence (3.2) gives, for $x=(c,g)$, $x'=(c',g')$, $\tilde{x}=(\tilde{c},\tilde{g})$ and $\tilde{x}'=(\tilde{c}',\tilde{g}')$ over $\square$,
\begin{align*}
\mathrm{tm}(x',\tilde{x})&=u_B\left((c',g'),(z,g')\right)\cdot u_B\left((z,g'),(\tilde{c},\tilde{g})\right)=U_B(c',z)\cdot U_B(z,\tilde{c}),\\
\mathrm{tm}^{-}(x,\tilde{x}')&=u_B\left((c,g),(z,g)\right)^{-1}\cdot u_B\left((z,g),(\tilde{c}',\tilde{g}')\right)^{-1}=U_B(c,z)^{-1}\cdot U_B(z,\tilde{c}')^{-1}.
\end{align*}
The kernel $K^{\mathrm{Ng}'}_{\square}(x,x')=\eta_k^{-3}\cdot\mathbf{1}[c=c'\in\square]\cdot\mathrm{Ng}'(g,g')$ is zero for $c'\ne c$.
For $X=\square$ and $K=K^{\mathrm{Ng}'}_{\square}$ the matrix that stands in the first term of Proposition~17.1\,(f) is therefore, for all $\tilde{x}'$ and $\tilde{x}$ over $\square$,
\begin{equation}\tag{18b.22}
\begin{aligned}
&\int\mathrm{d}x\int\mathrm{d}x'\,\mathrm{tm}^{-}(x,\tilde{x}')\,K^{\mathrm{Ng}'}_{\square}(x,x')\,\mathrm{tm}(x',\tilde{x})\\
&\quad=\eta_k^{6}\cdot\left(\frac{1}{48}\right)^{2}\sum_{c\in\square}\sum_{g,g'}U_B(c,z)^{-1}\cdot U_B(z,\tilde{c}')^{-1}\cdot\eta_k^{-3}\cdot\mathrm{Ng}'(g,g')\\
&\qquad\qquad\cdot U_B(c,z)\cdot U_B(z,\tilde{c})\\
&\quad=U_B(z,\tilde{c}')^{-1}\cdot U_B(z,\tilde{c})\cdot\eta_k^{3}\cdot\left(\frac{M}{\eta_k}\right)^{3}\left(\frac{1}{48}\right)^{2}\sum_{g,g'}\mathrm{Ng}'(g,g')\\
&\quad=U_B(z,\tilde{c}')^{-1}\cdot U_B(z,\tilde{c})\cdot M^3\cdot\mathrm{Nav}'=0
\end{aligned}
\end{equation}
by (d).
The other three terms of Proposition~17.1\,(f) contain one or two H\"older generators; their bidegrees are $(1,1)$, $(1,1)$ and $(0,2)$.
Take $B=\mathscr{T}'=\mathscr{A}'+\mathscr{C}'$.
Since $E(\square;\mathscr{A}',\mathscr{C}')=\boldsymbol{\Theta}_{\square,\mathscr{T}'}\left(E^{\mathrm{Ng}'}(\square)\right)$, and $E^{\mathrm{Ng}'}(\square)=\mathrm{Qd}\left[K^{\mathrm{Ng}'}_{\square}\right]$ by (18b.18), it follows that, for all $(\mathscr{A}',\mathscr{C}')\in\tilde{R}_k\times\mathfrak{X}_k$,
\begin{gather*}
E^{[2,0]}(\square;\mathscr{A}',\mathscr{C}')=0,\quad\text{that is,}\quad E_{20}(\square;\mathscr{A}',\mathscr{C}';\cdot,\cdot)=0,\\
E(\square;\mathscr{A}',\mathscr{C}')=E^{[1,1]}(\square;\mathscr{A}',\mathscr{C}')+E^{[0,2]}(\square;\mathscr{A}',\mathscr{C}').
\end{gather*}
Thus, at every connection, every basis monomial with a non-zero coefficient in $E(\square;\mathscr{A}',\mathscr{C}')$ contains at least one generator $\chi(\boldsymbol{\zeta})$.
Now let $(\mathscr{A},\mathscr{C})\in\tilde{R}_{k+1}\times\mathfrak{X}_{k+1}$; then $(\mathscr{A}_L,\mathscr{C}_L)\in\tilde{R}_k\times\mathfrak{X}_k$ (Section~2, Step~0).
By (18b.7), Lemma~1\,(iii), (G2.r7) and the bound on the supremum in (18b.19),
\begin{equation}\tag{18b.23}
\begin{aligned}
\left\|S_L\left[E(\square;\mathscr{A}_L,\mathscr{C}_L)\right]\right\|_{h_{k+1}}
&=L^{-\mathrm{pw}(1,1)}\left\|E^{[1,1]}(\square;\mathscr{A}_L,\mathscr{C}_L)\right\|_{h_k}\\
&\quad+L^{-\mathrm{pw}(0,2)}\left\|E^{[0,2]}(\square;\mathscr{A}_L,\mathscr{C}_L)\right\|_{h_k}\\
&\le L^{-\frac{13}{4}+\frac{3\varepsilon}{2}}\left\|E(\square;\mathscr{A}_L,\mathscr{C}_L)\right\|_{h_k}\\
&\le20L^{-\frac{13}{4}+\frac{3\varepsilon}{2}}\cdot M^3\cdot e_k^{-\frac{1}{2}}\cdot|\mathrm{Ng}|_1.
\end{aligned}
\end{equation}
As $20\le60$, (18b.23) gives the bound on $S_L\left[E(\square;\mathscr{A}_L,\mathscr{C}_L)\right]$ in (18b.19) once more.
At $(\mathscr{A},\mathscr{C})=(0,0)$ one has $(0,0)_L=(0,0)$, and (18b.22) at $B=0$ is the vanishing of the first term of (18b.21).

\emph{Reblocking.}
If $|X|_M\ge2$, then $\left(R E^{\mathrm{Ng}}\right)(X)=0$ by (d), and hence $E(X)=0$; for small $X$ the three terms in the definition of the adjustment vanish.
For $X=\square$ the set $\bar{X}^{L}$ is the $LM$-cube that contains $\square$.
This follows from the definition in \cite[Section~13.1.7]{QEDI}, recalled in Section~0.3, by which $\bar{X}^{L}$ is the union of the $LM$-cubes that contain an $M$-cube of $X$,
and from the nesting of the $M$-cubes in the $LM$-cubes (Section~0.3).
Consequently, in
\[
\left(L R E^{\mathrm{Ng}}\right)(Y;\mathscr{A},\mathscr{C})=\sum_{X:\ \bar{X}^{L}=LY}S_L\left[E(X;\mathscr{A}_L,\mathscr{C}_L)\right]
\]
a non-zero term can occur only if $LY$ is a single $LM$-cube, that is, if $Y$ is an $M$-cube of level $k+1$;
in this case the $X$ with $\bar{X}^{L}=LY$ and $|X|_M=1$ are the $L^3$ $M$-cubes contained in $LY$.
By (18b.19),
\begin{align*}
\left\|\left(L R E^{\mathrm{Ng}}\right)(Y;\mathscr{A},\mathscr{C})\right\|_{h_{k+1}}&\le L^3\cdot60L^{-\frac{13}{4}+\frac{3\varepsilon}{2}}\cdot M^3\cdot e_k^{-\frac{1}{2}}\cdot|\mathrm{Ng}|_1\\
&=60L^{-\frac{1}{4}+\frac{3\varepsilon}{2}}\cdot M^3\cdot e_k^{-\frac{1}{2}}\cdot|\mathrm{Ng}|_1.
\end{align*}
Moreover $d_M(Y)=0$ by (t0) at the level $k+1$; here (t0) is the assertion of \cite[Theorem~13.2]{QEDI} that $d_M(\square)=0$ for every $M$-cube $\square$ (Section~0.3).
Therefore
\[
\left\|L R E^{\mathrm{Ng}}\right\|^{x}_{k+1}\le60L^{-\frac{1}{4}+\frac{3\varepsilon}{2}}\cdot M^3\cdot e_k^{-\frac{1}{2}}\cdot|\mathrm{Ng}|_1=60L^{-\frac{1}{4}+\frac{3\varepsilon}{2}}\cdot\left\|E^{\mathrm{Ng}}\right\|^{x}_{k}.
\]

\paragraph{(f)}
\emph{Existence of the two decompositions.}
One has
\[
e_3=\hat n_0,\qquad e_2=\sqrt{2}\,\hat n_1+\hat n_0,\qquad e_1=\sqrt{2}\,\hat n_2+\sqrt{2}\,\hat n_1+\hat n_0;
\]
so $\hat n_0$, $\hat n_1$, $\hat n_2$ span $\mathbb{R}^3$ and form a basis.
For $\mathrm{cv}=(v_1,v_2,v_3)$,
\[
\mathrm{cv}=(v_1,v_2-v_3,0)+(0,v_3,v_3)+0,
\]
where the first vector is orthogonal to $\hat n_0=e_3$ and the second one to $\hat n_1=\frac{e_2-e_3}{\sqrt{2}}$.

Let $S\in\mathrm{Sp}$, $r=r(S)$ and $g,g'\in O_{\mathrm{cub}}$.
By (sp8),
\begin{gather*}
\mathbf{1}[rg=rg']=\mathbf{1}[g=g'],\qquad\mathbf{1}[rg'=(rg)s_i]=\mathbf{1}[g'=gs_i],\\
\mathbf{1}[g=g's_i]=\mathbf{1}[g'=gs_i].
\end{gather*}

\emph{The function $\mathrm{Ng}_1$.}
(n1): by (sp2),
\begin{align*}
S\,\mathrm{Ng}_1(g,g')\,S^{-1}&=48\cdot\mathbf{1}[g=g']\cdot i\cdot\mathrm{co}_1\cdot\det(g)\cdot\det(r)\cdot\Sigma_3\\
&=48\cdot\mathbf{1}[rg=rg']\cdot i\cdot\mathrm{co}_1\cdot\det(rg)\cdot\Sigma_3=\mathrm{Ng}_1(rg,rg').
\end{align*}
(n2): by (sp4), $\mathsf{C}^{-1}\,\mathrm{Ng}_1(g',g)^T\,\mathsf{C}=48\cdot\mathbf{1}[g'=g]\cdot i\cdot\mathrm{co}_1\cdot\det(g')\cdot\Sigma_3=\mathrm{Ng}_1(g,g')$.
(n3): by (sp5), since $\mathrm{co}_1$ and $\det(g)$ are real,
\[
\varkappa^{-1}\cdot\overline{i\cdot\mathrm{co}_1\cdot\det(g)\cdot\Sigma_3}\cdot\varkappa=-i\cdot\mathrm{co}_1\cdot\det(g)\cdot\left(-\Sigma_3\right)=i\cdot\mathrm{co}_1\cdot\det(g)\cdot\Sigma_3.
\]
Further $\mathrm{Nl}_1(g)=\frac{1}{48}\cdot48\cdot i\cdot\mathrm{co}_1\cdot\det(g)\cdot\Sigma_3=i\cdot\mathrm{co}_1\cdot\det(g)\cdot\Sigma_3$, and (sp7) gives
\[
|\mathrm{Ng}_1|_1=\left(\frac{1}{48}\right)^{2}\sum_{g}48\cdot|\mathrm{co}_1|\cdot4=4|\mathrm{co}_1|.
\]

\emph{The function $\mathrm{Ng}_u$.}
(n1): by (sp1),
\begin{align*}
S\,\mathrm{Ng}_u(g,g')\,S^{-1}&=48\sum_{i\le2}\mathrm{ru}_i\cdot\mathbf{1}[g'=gs_i]\cdot\left(\gamma\cdot(rg\hat n_i)\right)\\
&=48\sum_{i\le2}\mathrm{ru}_i\cdot\mathbf{1}[rg'=(rg)s_i]\cdot\left(\gamma\cdot\left((rg)\hat n_i\right)\right)=\mathrm{Ng}_u(rg,rg').
\end{align*}
(n2): by (sp4),
\[
\mathsf{C}^{-1}\,\mathrm{Ng}_u(g',g)^T\,\mathsf{C}=-48\sum_{i\le2}\mathrm{ru}_i\cdot\mathbf{1}[g=g's_i]\cdot\left(\gamma\cdot(g'\hat n_i)\right);
\]
on the support of the indicator one has $g'=gs_i$, and $g'\hat n_i=-g\hat n_i$ by (sp8), which is (I.s2): $\hat d_i(x^{(i)})=-\hat d_i(x)$.
So the right side equals $48\sum_{i\le2}\mathrm{ru}_i\cdot\mathbf{1}[g'=gs_i]\cdot\left(\gamma\cdot(g\hat n_i)\right)=\mathrm{Ng}_u(g,g')$.
(n3) follows from (sp5), as $\mathrm{ru}_i$ and $g\hat n_i$ are real.
For every $i$ exactly one $g'$ satisfies $g'=gs_i$; hence
\[
\mathrm{Nl}_u(g)=\frac{1}{48}\sum_{g'}\mathrm{Ng}_u(g,g')=\sum_{i\le2}\mathrm{ru}_i\cdot\left(\gamma\cdot(g\hat n_i)\right)=\gamma\cdot\left(g\sum_i\mathrm{ru}_i\hat n_i\right)=\gamma\cdot(g\,\mathrm{cu}).
\]
By (I.1.2), $x^{(i)}=(c,gs_i)$ and $\hat d_i(x)=-g\hat n_i$ for $i\le2$; therefore
\begin{align*}
E^{\mathrm{Ng}_u}(\square)&=\int_{\square}\mathrm{d}x\,\frac{1}{48}\sum_{g'}\bar{\psi}(c,g)^T\,\mathrm{Ng}_u(g,g')\,\psi(c,g')\\
&=\sum_{i\le2}\mathrm{ru}_i\int_{\square}\mathrm{d}x\,\bar{\psi}(x)^T\left(\gamma\cdot(g\hat n_i)\right)\psi(c,gs_i)\\
&=-\sum_{i\le2}\mathrm{ru}_i\int_{\square}\mathrm{d}x\,\bar{\psi}(x)^T\left(\gamma\cdot\hat d_i(x)\right)\psi\left(x^{(i)}\right).
\end{align*}
The elements $gs_0$, $gs_1$, $gs_2$ are pairwise different by (sp8); so for every pair $(g,g')$ at most one term of the sum over $i$ is different from zero, and (sp7), with $|g\hat n_i|_2=1$, gives
\begin{align*}
|\mathrm{Ng}_u|_1&=\left(\frac{1}{48}\right)^{2}\sum_{g}\sum_{i\le2}48|\mathrm{ru}_i|\sum_{\beta,\beta'}\left|\left(\gamma\cdot(g\hat n_i)\right)_{\beta\beta'}\right|\\
&\le\frac{1}{48}\sum_{g}\sum_{i\le2}|\mathrm{ru}_i|\cdot4\sqrt{2}=4\sqrt{2}\sum_i|\mathrm{ru}_i|.
\end{align*}

\emph{The function $\mathrm{Ng}_v$.}
In the definition of $\mathrm{Ng}_v$ the factor in front of $\det(g)$ is the complex number $i=\sqrt{-1}$, the same in all three terms of the sum;
the index of summation, which labels the elements $s_j$ and the vectors $\mathrm{cv}_j$, is denoted $j\in\{0,1,2\}$, so that
\[
\mathrm{Ng}_v(g,g')=48\sum_{j\le2}\mathbf{1}[g'=gs_j]\cdot i\cdot\det(g)\cdot\left(\left(1_2\otimes\sigma\right)\cdot(g\,\mathrm{cv}_j)\right).
\]
(n1): by (sp3),
\begin{align*}
S\,\mathrm{Ng}_v(g,g')\,S^{-1}&=48\sum_{j\le2}\mathbf{1}[g'=gs_j]\cdot i\cdot\det(g)\cdot\det(r)\cdot\left(\left(1_2\otimes\sigma\right)\cdot(rg\,\mathrm{cv}_j)\right)\\
&=48\sum_{j\le2}\mathbf{1}[rg'=(rg)s_j]\cdot i\cdot\det(rg)\cdot\left(\left(1_2\otimes\sigma\right)\cdot\left((rg)\mathrm{cv}_j\right)\right)\\
&=\mathrm{Ng}_v(rg,rg').
\end{align*}
(n2): by (sp4),
\[
\mathsf{C}^{-1}\,\mathrm{Ng}_v(g',g)^T\,\mathsf{C}=-48\sum_{j\le2}\mathbf{1}[g=g's_j]\cdot i\cdot\det(g')\cdot\left(\left(1_2\otimes\sigma\right)\cdot(g'\mathrm{cv}_j)\right);
\]
on the support of the indicator $g'=gs_j$, so that $\det(g')=-\det(g)$, and $g'\mathrm{cv}_j=g\,\mathrm{cv}_j$ because $\mathrm{cv}_j\cdot\hat n_j=0$ (sp8).
The right side is therefore $48\sum_{j\le2}\mathbf{1}[g'=gs_j]\cdot i\cdot\det(g)\cdot\left(\left(1_2\otimes\sigma\right)\cdot(g\,\mathrm{cv}_j)\right)=\mathrm{Ng}_v(g,g')$.
(n3): by (sp5),
\begin{align*}
\varkappa^{-1}\cdot\overline{i\cdot\det(g)\cdot\left(\left(1_2\otimes\sigma\right)\cdot(g\,\mathrm{cv}_j)\right)}\cdot\varkappa&=-i\cdot\det(g)\cdot\left(-\left(1_2\otimes\sigma\right)\cdot(g\,\mathrm{cv}_j)\right)\\
&=i\cdot\det(g)\cdot\left(\left(1_2\otimes\sigma\right)\cdot(g\,\mathrm{cv}_j)\right).
\end{align*}
Finally $\mathrm{Nl}_v(g)=\sum_{j\le2}i\cdot\det(g)\cdot\left(\left(1_2\otimes\sigma\right)\cdot(g\,\mathrm{cv}_j)\right)=i\cdot\det(g)\cdot\left(\left(1_2\otimes\sigma\right)\cdot(g\,\mathrm{cv})\right)$.

\emph{$\mathrm{Nav}=0$.}
Each of the three functions has (n1); so $\mathrm{Nav}=\mathrm{na}\cdot1$ by (b), and
\[
4\,\mathrm{na}=\operatorname{tr}\mathrm{Nav}=\left(\frac{1}{48}\right)^{2}\sum_{g,g'}\operatorname{tr}\mathrm{Ng}(g,g').
\]
The values of $\mathrm{Ng}_1$ are multiples of $\Sigma_3$, those of $\mathrm{Ng}_u$ are linear combinations of the $\gamma_\mu$, and those of $\mathrm{Ng}_v$ are linear combinations of the $1_2\otimes\sigma_\mu$; all these matrices are traceless by (sp6).
Hence $\mathrm{na}=0$.

\emph{Consequence.}
The assertions (a)--(e) hold for $\mathrm{Ng}_1$, $\mathrm{Ng}_u$ and $\mathrm{Ng}_v$, with (n3) and $\mathrm{na}=0$.
By (b), $\varepsilon(E^{\mathrm{Ng}})=0$ and $m(E^{\mathrm{Ng}})=0$; by (d), $\mathrm{Ng}'=\mathrm{Ng}$ and $R E^{\mathrm{Ng}}=E^{\mathrm{Ng}}$.

\emph{The function $\mathrm{Ng}_0$ and the sum.}
The function $\mathrm{Ng}_0(g,g')=48\cdot\mathbf{1}[g=g']\cdot\mathrm{co}_0$ is the function $\mathrm{Ng}_{\mathrm{na}}$ of (d) with the real number $\mathrm{na}=\mathrm{co}_0$.
It has (n1), (n2) and (n3); moreover $\mathrm{Nl}_0(g)=\mathrm{co}_0$, $\mathrm{Nav}_0=\mathrm{co}_0$, and $m(E^{\mathrm{Ng}_0})=-\mathrm{co}_0$ by (b).
The conditions (n1)--(n3) are real linear, and the map $\mathrm{Ng}\mapsto\mathrm{Nl}$ is linear.
Therefore $\mathrm{Ng}_0+\mathrm{Ng}_1+\mathrm{Ng}_u+\mathrm{Ng}_v$ has (n1)--(n3), and
\[
\mathrm{Nl}(g)=\mathrm{co}_0+i\cdot\det(g)\cdot\mathrm{co}_1\cdot\Sigma_3+\gamma\cdot(g\,\mathrm{cu})+i\cdot\det(g)\cdot\left(\left(1_2\otimes\sigma\right)\cdot(g\,\mathrm{cv})\right).
\]
By (c) this is the matrix in $\left(\mathrm{Ocl}\,J\xi\right)(c,g)=\mathrm{Nl}(g)\cdot\xi(c)$, that is, the display of Section~I.4.6.
For the sum one has $\mathrm{na}=\mathrm{co}_0$.

\emph{The couplings in (18b.21).}
Let $\mathrm{Ng}=\mathrm{Ng}_1+\mathrm{Ng}_u+\mathrm{Ng}_v$, so that $\mathrm{na}=0$.
The relation $\mathrm{Nr}(g)=\mathsf{C}^{-1}\,\mathrm{Nl}(g)^T\,\mathsf{C}$ and (sp4) give
\begin{align*}
\mathrm{Nl}(g)&=i\cdot\det(g)\cdot\mathrm{co}_1\cdot\Sigma_3+\gamma\cdot(g\,\mathrm{cu})+i\cdot\det(g)\cdot\left(\left(1_2\otimes\sigma\right)\cdot(g\,\mathrm{cv})\right),\\
\mathrm{Nr}(g)&=i\cdot\det(g)\cdot\mathrm{co}_1\cdot\Sigma_3-\gamma\cdot(g\,\mathrm{cu})-i\cdot\det(g)\cdot\left(\left(1_2\otimes\sigma\right)\cdot(g\,\mathrm{cv})\right);
\end{align*}
they are the coefficients of the two terms of bidegree $(1,1)$ in (18b.21).

\paragraph{(g)}
For $\mathrm{cu}\in\mathbb{R}^3$ let
\[
\mathrm{Ng}_s(g,g'):=48\cdot\mathbf{1}[g=g']\cdot\left(\gamma\cdot(g\,\mathrm{cu})\right).
\]
(n1): by (sp1), $S\,\mathrm{Ng}_s(g,g')\,S^{-1}=48\cdot\mathbf{1}[g=g']\cdot\left(\gamma\cdot(rg\,\mathrm{cu})\right)=\mathrm{Ng}_s(rg,rg')$.
(n3) follows from (sp5), since $g\,\mathrm{cu}$ is real.
The function $\mathrm{Nl}$ is $\mathrm{Nl}(g)=\frac{1}{48}\cdot48\cdot\left(\gamma\cdot(g\,\mathrm{cu})\right)=\gamma\cdot(g\,\mathrm{cu})$.
By (sp4),
\[
\mathsf{C}^{-1}\,\mathrm{Ng}_s(g',g)^T\,\mathsf{C}=-48\cdot\mathbf{1}[g=g']\cdot\left(\gamma\cdot(g\,\mathrm{cu})\right)=-\mathrm{Ng}_s(g,g').
\]
Consequently (n2) forces $\mathrm{Ng}_s=0$, that is, $\gamma\cdot(g\,\mathrm{cu})=0$ for every $g$.
Since $(\gamma\cdot v)^2=|v|_2^2$ for real $v$ by (I.c1), and $|g\,\mathrm{cu}|_2=|\mathrm{cu}|_2$, this means $\mathrm{cu}=0$.

The corresponding family is given by $E^{\mathrm{Ng}_s}(\square)=\mathrm{Qd}\left[K^{\mathrm{Ng}_s}_{\square}\right]$, which holds because (18b.18) is valid for every function $\mathrm{Ng}$, and by $E^{\mathrm{Ng}_s}(X)=0$ for $|X|_M\ge2$.
By (18b.3) and the computation of $\left(K^{\mathrm{Ng}}_{\square}\right)^{\mathrm{Cg}}$ made for (k5) in (a),
\[
\mathrm{Cg}\left[E^{\mathrm{Ng}_s}(\square;-\mathscr{A},-\mathscr{C})\right]=\mathrm{Qd}\left[\left(K^{\mathrm{Ng}_s}_{\square}\right)^{\mathrm{Cg}}\right]=\mathrm{Qd}\left[-K^{\mathrm{Ng}_s}_{\square}\right]=-E^{\mathrm{Ng}_s}(\square;\mathscr{A},\mathscr{C}).
\]
If (k5) held, it would give $E^{\mathrm{Ng}_s}(\square)=-E^{\mathrm{Ng}_s}(\square)$, that is, $\mathrm{Qd}\left[K^{\mathrm{Ng}_s}_{\square}\right]=0$;
then $K^{\mathrm{Ng}_s}_{\square}=0$ by Lemma~2\,(a), and hence $\mathrm{cu}=0$.
For $\mathrm{cu}\ne0$ the family therefore violates (k5).

\paragraph{Level 0; uniformity}
\emph{Level 0.}
At $k=0$ one has $\eta_0=1$,
\begin{gather*}
E^{\mathrm{Ng}}(\square)=\left(\frac{1}{48}\right)^{2}\sum_{c\in\square}\sum_{g,g'}\bar{\psi}(c,g)^T\,\mathrm{Ng}(g,g')\,\psi(c,g'),\\
E^{\mathrm{Ng}}_{20}(\square;x,x')=\mathbf{1}[c=c'\in\square]\cdot\mathrm{Ng}(g,g'),
\end{gather*}
and $\left\|E^{\mathrm{Ng}}\right\|^{x}_{0}=M^3\cdot q^{-\frac{1}{2}}\cdot|\mathrm{Ng}|_1$, the number $e_k$ being read as $q$;
the cube $\square$ contains $M^3$ photon sites and $48M^3$ fermion sites.
In (18b.20) the $47$ legs inside the cluster of $z$ have $d_X=1$, and $S_L$ maps them to legs with $d_X=\eta_1=\frac{1}{L}$.
In (18b.22) the factor $\eta_k^3\cdot\left(\frac{M}{\eta_k}\right)^3$ equals $1\cdot M^3$.
Nothing else changes at $k=0$: the argument depends on $k$ only through $\eta_k^3\cdot\left(\frac{M}{\eta_k}\right)^3=M^3$, through Proposition~18.5\,(d) and through (18b.7).

\emph{Uniformity.}
The identities in (a)--(d), (f) and (g) contain $k$, $\ell$ and $e_k$ only through the explicit factors $\eta_k^{-3}$, $M^3$ and $e_k^{-\frac{1}{2}}$.
The bounds in (e) come from Proposition~18.5\,(d), and these hold with the same numbers for all data (G2.S-2) (Section~2, Step~8).
The identity (18b.22) holds for every complex bond function $B$, and the factors $L^{-\mathrm{pw}(1,1)}$, $L^{-\mathrm{pw}(0,2)}$ in (18b.23) come from (18b.7) and contain only $\alpha$, $\varepsilon$ and $L$.
\qed

\subsection*{4. Conditions on the constants}

The constants are chosen in the order of Section~I.8.2.
The arguments of this chapter apply the following conditions on them directly.

\noindent
\begin{tabular}{|p{0.04\textwidth}|p{0.24\textwidth}|p{0.21\textwidth}|p{0.31\textwidth}|}
\hline
no. & condition & stated in & applied in\\
\hline
1 & $c_{\mathrm{rb}}=c_{X}\cdot K_{0}$, $c_{\mathrm{fl}}=61c_{X}\cdot K_{0}$ (in the form $61c_{\mathrm{rb}}=c_{\mathrm{fl}}$) & Section~I.8.2, line (C0); Proposition~14.2\,(a) & Section~2, Step~3\\
\hline
2 & $0<\alpha<1$ (part of $\alpha\in(1-\bar{\varepsilon},1)$) & Section~I.8.2, line (C1); Proposition~14.2\,(a) & $\alpha<1$: Lemma~1\,(iii) (together with row~3 it gives $\varepsilon>1-\alpha>0$); $0<\alpha$: Lemma~1\,(iv); through $\varepsilon>0$: Section~2, Step~2 (paragraph Sum: $L^{-\frac{13}{4}}\le L^{-\frac{13}{4}+\frac{3\varepsilon}{2}}$) and Step~3 (paragraph Sum: $L^{-\frac{1}{4}}\le L^{-\frac{1}{4}+\frac{3\varepsilon}{2}}$, $L^{\frac{3\varepsilon}{2}}\le L^{2\varepsilon}$)\\
\hline
3 & $1-\alpha<\varepsilon\le\bar{\varepsilon}=\frac{1}{200}$ & Section~I.8.2, line (C7); Proposition~14.2\,(b), row (S1) & Lemma~1\,(iii), (iv); through (iii): (18b.10), (18b.11), (18b.23)\\
\hline
4 & $\kappa\ge7\kappa_{0}+7$ (in the form $\kappa-1\ge\kappa_{0}$), together with $\kappa_{0}=54\log72>0$ & Section~I.8.2, line (C2); Proposition~14.2\,(a); for $\kappa_{0}=54\log72$: Section~I.8.2, line (C0); Proposition~14.2\,(a), row $\kappa_{0}$ & Lemma~1\,(ii), where $\kappa_{0}>0$ yields both $\kappa-1\ge\kappa_{0}$ and, from row~6, $L\cdot\left(\kappa-\kappa_{0}-2\right)\ge\kappa$; (t8) of Section~0.3 (from \cite[Theorem~13.2]{QEDI}) in Section~2, Step~3\\
\hline
5 & $L^{\frac{1}{8}-\frac{3\bar{\varepsilon}}{2}}\ge2400$ (in the forms $L^{\frac{1}{4}}\ge2$, $L^{\frac{3}{4}}\ge110$, $L^{\frac{7}{4}}\ge2$, $L^3\cdot\exp(-L)\le L^{-\frac{1}{4}}$) & Section~I.8.2, line (C3), first condition; Proposition~14.2\,(b), row (L1) & Lemma~1\,(i); (18b.8), (18b.9), (18b.11), Section~2, Step~3\\
\hline
6 & $L\cdot\left(\kappa-6\kappa_{0}-4\right)-1\ge\kappa$ (in the form $L\cdot\left(\kappa-\kappa_{0}-2\right)\ge\kappa$, which follows with $\kappa_{0}>0$ of row~4) & Section~I.8.2, line (C3), third condition; Proposition~14.2\,(b), row (L3) & Lemma~1\,(ii); (18b.12), (18b.13)\\
\hline
7 & $L$ odd (part of ``$L$ odd, $L\ge5$'') & Section~I.8.2, line (C3); Proposition~14.2\,(b), row (L5) & Section~0.3 (an $LM$-cube is the union of $L^3$ $M$-cubes, which is the reason for (t5) of Section~0.3, from \cite[Theorem~13.2]{QEDI}); Section~3, (e), paragraph Reblocking (the same nesting; the $L^3$ $M$-cubes of $LY$); together with row~8: $M/\eta_k=M\cdot L^k$ is odd\\
\hline
8 & $M$ a power of $L$, $M\ge L^2$ (in the form $M\ge1\ge\eta_k$; together with row~7, $M$ and $M/\eta_k=M\cdot L^k$ are odd) & Section~I.8.2, line (C5); Proposition~14.2\,(b), row (M1) & $M\ge\eta_k$: Lemma~3\,(b), Section~2, sub-steps (7.4), (7.7); $M/\eta_k$ odd: Lemma~3\,(c) (a closed $M$-cube reaches $\eta_k/2$ beyond its extreme sites), the number $\left(M/\eta_k\right)^3$ of the photon sites of an $M$-cube in Section~3, (a), (b), (18b.22)\\
\hline
\end{tabular}

The side $\ell$ and the charge $e_k$ are data of the statements:
$\ell$ and $\ell/L$ lie in $\mathrm{Cl}$, and of this the arguments need only that $\ell$ is a multiple of $LM$ and that $\ell\ge13LM$;
further $0<e_k$ and $e_{k+1}=L^{\frac{1}{2}}e_k\le e_0$.
The number $d_U=2L+1$ is defined in Section~I.1.1 (Section~I.8.2, line (C4); Proposition~14.2\,(a), row $d_U$);
it enters Section~2, sub-step (7.7), through
\[
14\,d_U\cdot M=(28L+14)\cdot M>2M+L\cdot M,
\]
together with the inequality $\ell_1\le\ell_2$ between the sides of the Propositions~27.m (the standing data (TT)).

Through the Propositions applied in the proofs, the following conditions enter as well (Annex~B, Section~G2.9):
\begin{itemize}
\item the numbers $c_{X}=54$, $\kappa_{0}=54\log72$, $K_{0}=\exp(\kappa_{0})/36$, $c_{\mathrm{rb}}=c_{X}\cdot K_{0}$, $c_{\mathrm{fl}}=61c_{X}\cdot K_{0}$ of Section~I.8.2, line (C0): Proposition~16.3\,(d) and the Propositions~17.m;
\item $\exp(-\varepsilon v)\le\frac{1}{4}$ for $v\ge-\log e_0$ (the first condition of (A8)): Propositions~17.1\,(a) and 18.1\,(c);
\item the first three inequalities of (A9): Propositions~17.1\,(c2), 17.2\,(d) and 18.2;
\item (L1): Propositions~18.1\,(c) and 18.2\,(ii), (iii);
\item $L$ odd, $L\ge5$ (L5): Propositions~15\,(n), 16.1--16.3, 17.1, 17.2, 18.1\,(b), 18.3\,(e) and 18.4\,(a);
\item (C1), (S1): Propositions~17.1\,(a), (c3), 18.1 and 18.2;
\item $M$ a power of $L$, $M$ odd, $M\ge L^2$ (C5): Propositions~15\,(n), 16.1--16.3, 17.1, 17.2, 18.1\,(c), 18.2, 18.3\,(e) and 18.4\,(a).
\end{itemize}
The first condition of (A8) and the first three inequalities of (A9) enter in this way only.
The other conditions of the list are, in addition, applied directly: the numbers of line (C0) in the rows~1 and 4 of the table, (C1) and (S1) in the rows~2 and 3, (L1) in row~5, the oddness of $L$ from (L5) in row~7, and (C5) in row~8.
Only through the Propositions enter the inequality $L\ge5$ of (L5), the inequality $M\ge L^2$ of (C5) beyond $M\ge1$, and the consequences of (L1) other than those listed in row~5.

The proof of Proposition~18.6 requires the conditions under which Proposition~18.5\,(d) holds, namely the rows 2, 3 and 5 and the conditions that enter through the Propositions applied in Section~2, Steps~0--2, and, in addition to these, the rows 7 and 8:
that $L$ is odd (the nesting of the $M$-cubes in the $LM$-cubes, paragraph Reblocking of (e)),
and that $M/\eta_k=M\cdot L^k$ is odd (the number $\left(M/\eta_k\right)^3$ in (a), (b) and (18b.22)).

All the conditions above are among those of Sections~I.8.2 and I.8.4 (Proposition~14.2\,(a), (b)), and so they hold for the constants chosen there.

\section*{Chapter 20. Terminal bounds; the determinants of one step; the free determinant telescopes: proofs of Propositions 20.1--20.3}

\subsection*{0. Conventions of this chapter}

In Section~2 the letter $P$ stands for $P(B)$.
The display (6.1) is part of the statement of Proposition~20.2.
In Section~2 the steps of the proofs are labelled (n1)--(n5), Lemma~A and (s0)--(s7); in Section~3 they are labelled (D1), (D2).

\subsection*{1. Proposition 20.1}

\subsubsection*{Proof of Proposition 20.1}

\paragraph{1.1. Numbers.}
Condition (A8) of \cite[Section~5.6.2]{QEDI} is imposed in Section~I.8.4.
In the row (A8) of Proposition~14.2\,(b) its first inequality takes the form $\exp(-\varepsilon v)\le\frac{1}{4}$, required for every $v\ge-\log e_0$; at $v=-\log e_0$ it states that $e_0^{\varepsilon}\le\frac{1}{4}$.
As $\varepsilon>0$ (Section~I.0.5), it follows that $e_0<1$.
Hence $0<e_k\le e_0<1$, and, $\varepsilon$ being positive once more,
\begin{equation}\tag{20.1}
e_k^{\varepsilon}\le1,\qquad e_k^{2\varepsilon}\le1.
\end{equation}
Sections~I.3.4 and I.0.5 give $b_k\le b_{+}=b=1$.
The constant $C^{F}$ is non-negative, because it bounds a modulus in Proposition~5\,(e), and $c(\gamma^{F})$ is a sum of positive numbers by (G2.r6).
The line (C6) of Section~I.8.2 therefore yields $C_{0}\ge4b_{+}+16\,C^{F}\cdot c(\gamma^{F})\ge4$.
By Section~I.5.1, $h_k=\left(e_k^{-\frac{1}{4}},e_k^{-\frac{1}{4}+\varepsilon}\right)$; consequently
\begin{equation}\tag{20.2}
\mathrm{aw}_1=h_{k,1},\qquad C_{0}\cdot\mathrm{aw}_2=e_k^{-\frac{1}{4}+\varepsilon}=h_{k,2}=e_k^{\varepsilon}\cdot h_{k,1}\le h_{k,1},\qquad\mathrm{aw}_2=C_{0}^{-1}\cdot e_k^{\varepsilon}\cdot\mathrm{aw}_1\le\mathrm{aw}_1.
\end{equation}

\paragraph{1.2. Index sets; elements of degree two.}
Throughout this paragraph $k\ge1$.
The index sets and their weights are those of Section~I.3.6, with $\beta\in\{1,\dots,4\}$ and $\omega\in\{0,1\}$:
\begin{description}
\item[$\{u\}$] $u=(y,\beta,\omega)$ with $y\in T^0_k$ (the $\ell^3$ unit sites), and $\nu=1$;
\item[$\{\xi\}$] $\xi=(x,\beta,\omega)$ with $x=(c,g)\in T^{\mathrm{f}}_k=T_k\times O_{\mathrm{cub}}$ (the $(\ell/\eta_k)^3$ photon sites $c$, and over each of them the $48$ labels $g$), and $\nu=\eta_k^3/48$;
\item[$\{\boldsymbol{\zeta}\}$] $\boldsymbol{\zeta}=(x,x',\beta,\omega)$ with $(x,x')\in\mathrm{Pz}_k$, and $\nu=(\eta_k^3/48)^2$.
\end{description}
By (I.AP), a sum over $y$, over $x$ or over $(x,x')$ is a sum over a fundamental domain of $\ell\mathbb{Z}^3$; every summand below is $\ell\mathbb{Z}^3$-periodic.
Let $\phi$ be either the family $\Psi$ (with points $t=y$) or the family $\psi$ (with points $t=x$), let $\mathrm{Kj}(t,\beta;t',\beta')$ be complex numbers, and let $\mathsf{w}>0$.
For the norm (I.5.1) with the single weight $\mathsf{w}$ one has
\begin{equation}\tag{20.3}
\left\|\sum_{t,\beta,t',\beta'}\mathrm{Kj}(t,\beta;t',\beta')\cdot\bar{\phi}_{\beta}(t)\,\phi_{\beta'}(t')\right\|_{\mathsf{w}}=\mathsf{w}^2\cdot\sum_{t,\beta,t',\beta'}\left|\mathrm{Kj}(t,\beta;t',\beta')\right|.
\end{equation}
Indeed, $\bar{\phi}_{\beta}(t)=\phi(t,\beta,1)$ and $\phi_{\beta'}(t')=\phi(t',\beta',0)$ are two distinct generators.
Their product is $+1$ or $-1$ times the basis monomial $\phi(I)$ of the set $I=\{(t,\beta,1),(t',\beta',0)\}$, and the assignment $(t,\beta,t',\beta')\mapsto I$ is injective.
The coefficient of $\phi(I)$ is therefore $\pm\mathrm{Kj}(t,\beta;t',\beta')$, and all other coefficients are zero.

\paragraph{1.3. Proof of (a).}
Let $A\in\tilde{G}_k$.
Data (G2.S-3) are data (SF) with $k\ge1$.
Proposition~2.8\,(a), applied at $y=\infty$, shows that the operator $F_{\infty}(A)=D_A+\mu+b_k\cdot J\cdot P_k(A)\cdot J^{*}$ is a bijection of $\mathfrak{S}^{\mathrm{ch}}(\eta_k,\ell)$.
Hence $S_k(A)=F_{\infty}(A)^{-1}$ exists, and the operators
\[
D_k(A)=b_k-b_k^2\cdot Q_k(A)\cdot J^{*}\cdot S_k(A)\cdot J\cdot Q_k^{T}(-A),\qquad M_k(A)=b_k-D_k(A)
\]
on $\mathfrak{S}(1,\ell)$ (Section~I.3.5) are defined.

\emph{Analyticity.}
Choose the delta functions of fundamental domains as bases of $\mathfrak{S}^{\mathrm{ch}}(\eta_k,\ell)$ and of $\mathfrak{S}(1,\ell)$.
Section~I.3.3 gives
\begin{gather*}
\left(Q_k(A)\cdot J^{*}f\right)(y)=\sum_{c\in\Delta_y}\sum_{g\in O_{\mathrm{cub}}}\frac{\eta_k^3}{48}\cdot\varphi_A(y,c)\cdot f(c,g),\\
\left(J\cdot Q_k^{T}(-A)\,v\right)(c,g)=\varphi_A(y,c)^{-1}\cdot v(y)\quad(c\in\Delta_y).
\end{gather*}
Every matrix entry of $Q_k(A)\cdot J^{*}$ is thus $0$ or $\pm\frac{\eta_k^3}{48}\cdot\varphi_A(y,c)$, and every matrix entry of $J\cdot Q_k^{T}(-A)$ is $0$ or $\pm\varphi_A(y,c)^{-1}$.
Here $\varphi_A(y,c)^{\pm1}=\exp\left(\pm i\,e_k\cdot\eta_k\cdot(\tau_kA)(y,c)\right)$, and $(\tau_kA)(y,c)$ is a linear form in the finitely many numbers $A(b)$ ((I.3.3), Section~I.1.1).
So these entries are entire functions of $A$.
By Proposition~2.8\,(b) at $y=\infty$, the entries of $S_k(A)$ are analytic on $\tilde{G}_k$.
This set is open in the space of the complex bond functions (Section~I.5.3): the space is finite-dimensional, $G_k$ is defined by a strict inequality for a norm on it, and $\tilde{G}_k$ is the union of the translates $A_0+G_k$ over the real $A_0$ that its definition admits.
The entries of $D_k(A)$ are a constant minus finite sums of products of three functions of these kinds; hence they are analytic on $\tilde{G}_k$.

\emph{Bound.}
For $k\ge1$, Section~I.3.6 gives
\[
\langle\bar{\Psi},D_k(A)\Psi\rangle=\sum_{y,y'}\bar{\Psi}(y)\cdot D_k(A)(y,y')\,\Psi(y')=\sum_{y,\beta,y',\beta'}D_k(A)(y,y')_{\beta\beta'}\cdot\bar{\Psi}_{\beta}(y)\,\Psi_{\beta'}(y').
\]
Therefore (20.3) yields
\[
\left|\langle\bar{\Psi},D_k(A)\Psi\rangle\right|_{\mathrm{aw}_1}=\mathrm{aw}_1^2\cdot\sum_{y}\sum_{y'}\sum_{\beta,\beta'}\left|D_k(A)(y,y')_{\beta\beta'}\right|.
\]
The operator $b_k$ times the identity of $\mathfrak{S}(1,\ell)$ has the kernel $b_k\cdot\mathbf{1}[y=y']\cdot1_4$, so that $D_k(A)(y,y')=b_k\cdot\mathbf{1}[y=y']\cdot1_4-M_k(A)(y,y')$.
For a $4\times4$ block $m$, the number $|m|_c$ is the largest of its four row sums of moduli (Section~I.1.5); hence $\sum_{\beta,\beta'}|m_{\beta\beta'}|\le4\,|m|_c$.
Proposition~5\,(e) gives $|M_k(A)(y,y')|_c\le C^{F}\cdot\exp\left(-\gamma^{F}\cdot d(y,y')\right)$.
On the torus $T^0_k$ one has $\sum_{y'}\exp\left(-\gamma^{F}\cdot d(y,y')\right)\le c(\gamma^{F})$ by (G2.r6), since every class $y'$ has a representative $y+i$, $i\in\mathbb{Z}^3$, with $|i|=d(y,y')$, and different classes have different $i$.
Consequently, for every unit site $y$,
\begin{equation}\tag{20.4}
\begin{split}
\sum_{y'}\sum_{\beta,\beta'}\left|D_k(A)(y,y')_{\beta\beta'}\right|&\le4b_k+4\,C^{F}\cdot\sum_{y'}\exp\left(-\gamma^{F}\cdot d(y,y')\right)\le4b_{+}+4\,C^{F}\cdot c(\gamma^{F})\\
&\le4b_{+}+16\,C^{F}\cdot c(\gamma^{F})\le C_{0}.
\end{split}
\end{equation}
Since $\#|T^0_k|=\ell^3$ and $\mathrm{aw}_1^2=e_k^{-\frac{1}{2}}$, it follows that $\left|\langle\bar{\Psi},D_k(A)\Psi\rangle\right|_{\mathrm{aw}_1}\le C_{0}\cdot e_k^{-\frac{1}{2}}\cdot\ell^3$.

\paragraph{1.4. The substitution and its row sums.}
Let $B\in\tilde{G}_k$.
By Proposition~16.1\,(c), also $-B\in\tilde{G}_k$.
Proposition~2.8\,(a) at $y=\infty$, applied at $B$ and at $-B$, shows that $S_k(B)$ and $S_k(-B)$ exist; therefore the operators
\[
H_k(B)=b_k\cdot S_k(B)\cdot J\cdot Q_k^{T}(-B),\qquad\bar{H}_k(B)=\mathsf{C}^{-1}\cdot H_k(-B)\cdot\mathsf{C}
\]
from $\mathfrak{S}(1,\ell)$ to $\mathfrak{S}^{\mathrm{ch}}(\eta_k,\ell)$ (Section~I.3.5) are defined.
A function on $\{u\}$ is a pair $(v^0,v^1)$ of elements of $\mathfrak{S}(1,\ell)$, its value at $u=(y,\beta,\omega)$ being $v^{\omega}(y)_{\beta}$; recall that a function of $y$ on a fundamental domain has exactly one antiperiodic extension to $\mathbb{Z}^3$, and that this extension has the same sup norm.
Define operators $H^{(1)}$, $H^{(2)}$ from the functions on $\{u\}$ to the functions on $\{\xi\}$, resp. on $\{\boldsymbol{\zeta}\}$, by
\begin{equation}\tag{20.5}
\begin{aligned}
\left(H^{(1)}(v^0,v^1)\right)(x,\beta,0)&:=\left(H_k(B)\,v^0\right)(x)_{\beta},\\
\left(H^{(1)}(v^0,v^1)\right)(x,\beta,1)&:=\left(\bar{H}_k(B)\,v^1\right)(x)_{\beta},\\
\left(H^{(2)}(v^0,v^1)\right)(x,x',\beta,0)&:=\left(\delta_{\alpha,B}H_k(B)\,v^0\right)(x,x')_{\beta},\\
\left(H^{(2)}(v^0,v^1)\right)(x,x',\beta,1)&:=\left(\delta_{\alpha,-B}\bar{H}_k(B)\,v^1\right)(x,x')_{\beta},
\end{aligned}
\end{equation}
with the H\"older quotients $\delta_{\alpha,\pm B}$ of (I.1.7).
Since $H_k(B)\,v^0$ and $\bar{H}_k(B)\,v^1$ lie in $\mathfrak{S}^{\mathrm{ch}}(\eta_k,\ell)$, and $\delta_{\alpha,\pm B}$ maps $\mathfrak{S}^{\mathrm{ch}}(\eta_k,\ell)$ into the functions of pairs of Section~I.1.7, the right sides change sign under $x\mapsto x+\ell e_{\mu}$, resp. under $(x,x')\mapsto(x+\ell e_{\mu},x'+\ell e_{\mu})$, as (I.AP) requires of $\psi(\xi)$ and $\chi(\boldsymbol{\zeta})$.
The homomorphism $F\mapsto F\left(H_k^{\#}(B)\Psi\right)$ of the statement is
\begin{gather*}
\psi(\xi)\mapsto\left(H^{(1)}\Psi\right)(\xi)=\sum_{u}H^{(1)}(\xi,u)\cdot\Psi(u),\\
\chi(\boldsymbol{\zeta})\mapsto\left(H^{(2)}\Psi\right)(\boldsymbol{\zeta})=\sum_{u}H^{(2)}(\boldsymbol{\zeta},u)\cdot\Psi(u)\qquad(\nu(u)=1);
\end{gather*}
in the notation of Proposition~15\,(l)\,(iv) this reads $F\left(H_k^{\#}(B)\Psi\right)=F\left(H^{(1)}\Psi,H^{(2)}\Psi\right)$.

\emph{Row sums.}
As a function on $\{u\}$, the pair $(v^0,v^1)$ has the sup norm $\max\left(|v^0|_{\infty},|v^1|_{\infty}\right)$, where $|v|_{\infty}=\sup_y|v(y)|_c=\sup_{y,\beta}|v(y)_{\beta}|$ (Section~I.1.5).
Proposition~6\,(f), first clause, together with Proposition~6\,(d), at the connection $B\in\tilde{G}_k$, gives for every $v\in\mathfrak{S}(1,\ell)$
\begin{align*}
\left|H_k(B)\,v\right|_{\infty}&\le C^{F}\cdot|v|_{\infty},&\sup_{(x,x')\in\mathrm{Pz}_k}\left|\left(\delta_{\alpha,B}H_k(B)\,v\right)(x,x')\right|_c&\le C^{F}\cdot|v|_{\infty},\\
\left|\bar{H}_k(B)\,v\right|_{\infty}&\le C^{F}\cdot|v|_{\infty},&\sup_{(x,x')\in\mathrm{Pz}_k}\left|\left(\delta_{\alpha,-B}\bar{H}_k(B)\,v\right)(x,x')\right|_c&\le C^{F}\cdot|v|_{\infty}.
\end{align*}
The two suprema in the right column extend over all pairs of $\mathrm{Pz}_k$: over the $48\cdot47$ ordered pairs inside the cluster of each photon site ($d_X=\eta_k$, no phase), over the other short pairs ($d_X<1$) and over the long pairs ($1\le d_X\le M$).
The sup norms in the left column are suprema over the $(\ell/\eta_k)^3\cdot48$ fermion sites.
Neither $H^{(1)}$ nor $H^{(2)}$ mixes $\omega=0$ with $\omega=1$.
Hence $\left|H^{(i)}(v^0,v^1)\right|_{\infty}\le C^{F}\cdot\max\left(|v^0|_{\infty},|v^1|_{\infty}\right)$ for $i=1,2$, and Proposition~15\,(c) gives
\begin{equation}\tag{20.6}
\left|H^{(1)}\right|_{1,\infty}\le C^{F}\le C_{0},\qquad\left|H^{(2)}\right|_{1,\infty}\le C^{F}\le C_{0}.
\end{equation}

\paragraph{1.5. Proof of (b).}
Let $A\in\tilde{G}_k$, and let $H^{(1)}$ be formed with $B=A$.
In $\mathrm{Gr}[\psi]$, Section~I.3.6 and (I.1.3) give
\[
\langle\bar{\psi},\psi\rangle=\int\mathrm{d}x\,\bar{\psi}(x)\cdot\psi(x)=\sum_{c\in T_k}\sum_{g\in O_{\mathrm{cub}}}\frac{\eta_k^3}{48}\cdot\sum_{\beta}\bar{\psi}_{\beta}(c,g)\,\psi_{\beta}(c,g).
\]
Formula (20.3), applied with $\mathrm{Kj}(x,\beta;x',\beta')=\frac{\eta_k^3}{48}\cdot\mathbf{1}[x=x',\beta=\beta']$, gives for every $\mathsf{w}_1>0$
\begin{equation}\tag{20.7}
\left\|\langle\bar{\psi},\psi\rangle\right\|_{\mathsf{w}_1}=\mathsf{w}_1^2\cdot4\cdot\sum_{c\in T_k}\sum_{g\in O_{\mathrm{cub}}}\frac{\eta_k^3}{48}=\mathsf{w}_1^2\cdot4\cdot\left(\frac{\ell}{\eta_k}\right)^3\cdot48\cdot\frac{\eta_k^3}{48}=4\,\ell^3\cdot\mathsf{w}_1^2.
\end{equation}
The element ``$\langle\bar{\psi},\psi\rangle$ at $\psi=H_k(A)\Psi$, $\bar{\psi}=\bar{H}_k(A)\bar{\Psi}$'' is the image of $\langle\bar{\psi},\psi\rangle$ under the homomorphism $\psi(\xi)\mapsto\left(H^{(1)}\Psi\right)(\xi)$ of $\mathrm{Gr}[\psi]$ into $\mathrm{Gr}[\Psi]$.
Proposition~15\,(d) applies with $p=q=1$, $\phi_1=\psi$, $\theta_1=\Psi$, $O_{11}=H^{(1)}$, $C_{11}=C_{0}$ (admissible by (20.6)), $\mathsf{w}'_1=\mathrm{aw}_2$ and $\mathsf{w}_1=C_{0}\cdot\mathrm{aw}_2=e_k^{-\frac{1}{4}+\varepsilon}$; its condition $C_{11}\cdot\mathsf{w}'_1\le\mathsf{w}_1$ holds with equality.
By (20.7) and (20.1) it follows that
\[
\left|\langle\bar{\psi},\psi\rangle\right|_{\mathrm{aw}_2}\le\left\|\langle\bar{\psi},\psi\rangle\right\|_{\mathsf{w}_1}=4\,\ell^3\cdot e_k^{-\frac{1}{2}+2\varepsilon}\le4\,\ell^3\cdot e_k^{-\frac{1}{2}}.
\]

\paragraph{1.6. Proof of (c).}
Let $(\mathscr{A},\mathscr{C})\in\tilde{R}_k\times\mathfrak{X}_k$.
The inclusion $\tilde{R}_k+\mathfrak{X}_k\subset\tilde{G}_k$ of \cite[Lemma~17.4(a)]{QEDI} holds for every level $k\ge0$, every side $\ell\in\mathrm{Cl}$ and every $e_k\in(0,e_0]$.
Hence $\mathscr{T}=\mathscr{A}+\mathscr{C}$ lies in $\tilde{G}_k$.
Let $H^{(1)}$, $H^{(2)}$ be formed with $B=\mathscr{T}$, and apply Proposition~15\,(l)\,(iv) with $E=F\in\mathrm{Gr}[\psi,\chi]$, $\mathrm{Cr}=C_{0}$ (admissible by (20.6)), $\mathrm{aw}=\mathrm{aw}_2$ and $(\mathsf{w}_1,\mathsf{w}_2)=h_k$.
Its hypothesis is satisfied, since (20.2) gives $\mathrm{Cr}\cdot\mathrm{aw}=C_{0}\cdot\mathrm{aw}_2=h_{k,2}=\min(h_{k,1},h_{k,2})$.
Therefore
\[
\left|F\left(H_k^{\#}(\mathscr{T})\Psi\right)\right|_{\mathrm{aw}_2}=\left|F\left(H^{(1)}\Psi,H^{(2)}\Psi\right)\right|_{\mathrm{aw}_2}\le\|F\|_{h_k}.
\]
For $\mathscr{E}\in K^{x,C}_{k}$ and $X\in D_k$, the activity $\mathscr{E}(X;\mathscr{A},\mathscr{C})$ is an element of $\mathrm{Gr}[\psi,\chi]$ (Section~I.5.4), and the choice $F:=\mathscr{E}(X;\mathscr{A},\mathscr{C})$ gives the last assertion.

\paragraph{1.7. Uniformity.}
The numbers $C^{F}$, $\gamma^{F}$, $C_{0}$, $b_{+}$ are constants of the lines (C4), (C4r), (C6) of Section~I.8.2, and Proposition~5\,(e) and Proposition~6\,(f) hold with these constants for all data (SF) with $k\ge1$ and all connections in $\tilde{G}_k$.
So the three bounds hold, with the numbers as stated, for every $k\ge1$, every $\ell\in\mathrm{Cl}$, every $\mu\in(0,1]$, every $e_k\in(0,e_0]$, every $A\in\tilde{G}_k$, every $(\mathscr{A},\mathscr{C})\in\tilde{R}_k\times\mathfrak{X}_k$ and every $F$.
A twist, being a real constant connection, is included in $A$, resp. in $\mathscr{A}$; the sets $\tilde{G}_k$ and $\tilde{R}_k$ are invariant under the addition of real constants (Proposition~16.1\,(c)).
The factor $4$ in (20.4) and in (20.7) is the number of components of $\Psi(y)$, resp. of $\psi(x)$; in (20.7) the sum over the $48$ labels is compensated by the weight $1/48$.
\qed

\subsection*{2. Proposition 20.2}

\subsubsection*{Proof of Proposition 20.2}

Recall the determinant formula for block matrices \cite[Section~0.8.5]{HornJohnson2013}:
\begin{equation}\tag{cl3}
\begin{aligned}
\det\left(\begin{array}{cc}A_1&A_2\\ A_3&A_4\end{array}\right)&=\det A_1\cdot\det\left(A_4-A_3\cdot A_1^{-1}\cdot A_2\right)&&\text{if $A_1$ is invertible},\\
\det\left(\begin{array}{cc}A_1&A_2\\ A_3&A_4\end{array}\right)&=\det A_4\cdot\det\left(A_1-A_2\cdot A_4^{-1}\cdot A_3\right)&&\text{if $A_4$ is invertible}.
\end{aligned}
\end{equation}

\paragraph{2.1. Spaces, maps, numbers.}
All operators are formed at the connection $B$, with $\mathrm{zc}$ in the place of the mass.
In 2.1--2.6 the argument $B$ of $S_k(B)$, $D_k(B)$, $\delta Z_k(B)$, $S^{0}_{k+1}(B)$, $D_0(B)$, $\delta Z_0(B)$ is omitted.
Accordingly, $D_k$ and $D_0$ denote the operators $D_k(B)$ and $D_0(B)$, while $D_B$ and $D_{B'}$ denote the operator (I.2.3) at the connections $B$ and $B'$.
The numbers $e_k$ and $e_{k+1}$ are those of the data, also when $k=0$:
for every $k\ge0$, the case $k=0$ included, $e_k>0$ is the number given with the data of Proposition~20.2, and $e_{k+1}=L^{\frac{1}{2}}\cdot e_k$,
in agreement with the convention of Section~I.2.5 and (I.L0.2) for formulas written for a general level $k$.
Put
\[
\mathrm{Vf}:=\mathfrak{S}^{\mathrm{ch}}(\eta_k,\ell),\qquad\mathrm{Vu}:=\mathfrak{S}(1,\ell),\qquad\mathrm{Vl}:=\mathfrak{S}(L,\ell);
\]
by Section~I.1.5 their dimensions are $192\,(\ell/\eta_k)^3$, $4\ell^3=4n_k$ and $4(\ell/L)^3=4n_{k+1}$, respectively.
For $k\ge1$ define
\begin{gather*}
\mathrm{Qa}:=Q_k(B)\cdot J^{*}\colon\mathrm{Vf}\to\mathrm{Vu},\qquad\mathrm{Qa}':=J\cdot Q_k^{T}(-B)\colon\mathrm{Vu}\to\mathrm{Vf},\\
P:=P(B)=Q^{T}(-B)\cdot Q(B)\ \text{on }\mathrm{Vu},\\
\mathrm{Qd}:=Q(B)\cdot\mathrm{Qa}\colon\mathrm{Vf}\to\mathrm{Vl},\qquad\mathrm{Qd}':=\mathrm{Qa}'\cdot Q^{T}(-B)\colon\mathrm{Vl}\to\mathrm{Vf}.
\end{gather*}
For $k=0$ (level $0$) one has $\eta_0=1$, $\mathrm{Vf}=\mathfrak{S}^{\mathrm{ch}}(1,\ell)$ and $J^{*}\colon\mathfrak{S}^{\mathrm{ch}}(1,\ell)\to\mathfrak{S}(1,\ell)$, and one defines
\[
\mathrm{Qd}:=Q(B)\cdot J^{*}\colon\mathrm{Vf}\to\mathrm{Vl},\qquad\mathrm{Qd}':=J\cdot Q^{T}(-B)\colon\mathrm{Vl}\to\mathrm{Vf};
\]
at $k=0$ the maps $\mathrm{Qa}$, $\mathrm{Qa}'$ and the number $b_0$ do not occur, and $P(B)$ occurs only inside $J\cdot P(B)\cdot J^{*}$.
For every $k\ge0$ let
\[
\mathrm{bq}:=b_{k+1}/L.
\]

(n1) \emph{Numbers.}
Since $b>0$ and $L>1$, one has $b_L=b/L>0$, and $b_k=b\cdot(1-L^{-1})/(1-L^{-k})>0$ for $k\ge1$ (Section~I.3.4).
For $k\ge1$,
\[
\frac{b\cdot b_k}{b_k+b_L}=\frac{b\cdot(1-L^{-1})}{(1-L^{-1})+L^{-1}\cdot(1-L^{-k})}=\frac{b\cdot(1-L^{-1})}{1-L^{-k-1}}=b_{k+1}.
\]
Consequently
\begin{equation}\tag{20.8}
\mathrm{bq}=\frac{b_L\cdot b_k}{b_k+b_L}\quad(k\ge1),\qquad\mathrm{bq}=b_1/L=b_L\quad(k=0).
\end{equation}

(n2) \emph{Cluster maps.}
For every $\xi$ one has $(J^{*}J\xi)(c)=\frac{1}{48}\sum_{g\in O_{\mathrm{cub}}}\xi(c)=\xi(c)$, that is,
\begin{equation}\tag{20.9}
J^{*}\cdot J=1
\end{equation}
(compare (I.J1)).

(n3) \emph{Averages.}
For a unit site $y$,
\[
\Delta_y=\left\{c\in\mathbb{L}_k:|c-y|<\tfrac{1}{2}\right\}=\left\{y+\eta_k\cdot i:i\in\mathbb{Z}^3,\ |i_{\mu}|\le(L^k-1)/2\right\},
\]
and $\#|\Delta_y|=L^{3k}=\eta_k^{-3}$, because $L^k$ is odd.
For $y'\in L\mathbb{Z}^3$, the $L$-block $B(y')=\{y'+i:i\in\mathbb{Z}^3,\ |i_{\mu}|\le(L-1)/2\}$ has $\#|B(y')|=L^3$, and every $y\in B(y')$ satisfies $|y-y'|\le(L-1)/2\le M$, because $M\ge L^2$ (Section~I.0.5).
The number $U_B(y',y)$ of (I.1.6) is therefore defined; being an exponential, it is different from zero.
The displays of Section~I.3.3 give, for $v\in\mathrm{Vu}$, $v'\in\mathrm{Vl}$ and every complex $B$,
\begin{gather*}
\left(Q_k(B)\cdot Q_k^{T}(-B)\,v\right)(y)=\eta_k^3\cdot\sum_{c\in\Delta_y}\varphi_B(y,c)\cdot\varphi_B(y,c)^{-1}\cdot v(y)=v(y),\\
\left(Q(B)\cdot Q^{T}(-B)\,v'\right)(y')=L^{-3}\cdot\sum_{y\in B(y')}U_B(y',y)\cdot U_B(y',y)^{-1}\cdot v'(y')=v'(y').
\end{gather*}
Together with (20.9) this gives
\begin{equation}\tag{20.10}
\begin{aligned}
&Q(B)\cdot Q^{T}(-B)=1\ \text{on }\mathrm{Vl};\qquad\mathrm{Qa}\cdot\mathrm{Qa}'=Q_k(B)\cdot J^{*}\cdot J\cdot Q_k^{T}(-B)=1\ \text{on }\mathrm{Vu}\ (k\ge1);\\
&\mathrm{Qd}\cdot\mathrm{Qd}'=1\ \text{on }\mathrm{Vl}\ (k\ge0).
\end{aligned}
\end{equation}
Indeed, for $k\ge1$ one has $\mathrm{Qd}\cdot\mathrm{Qd}'=Q(B)\cdot\mathrm{Qa}\cdot\mathrm{Qa}'\cdot Q^{T}(-B)=Q(B)\cdot Q^{T}(-B)$, and for $k=0$ one has $\mathrm{Qd}\cdot\mathrm{Qd}'=Q(B)\cdot J^{*}\cdot J\cdot Q^{T}(-B)=Q(B)\cdot Q^{T}(-B)$.
The identities (I.Q1) are among the relations (20.10).

(n4) \emph{Projections.}
By Section~I.3.3, $P_k(B)=Q_k^{T}(-B)\cdot Q_k(B)$ and $P_{k+1}(B)=Q_k^{T}(-B)\cdot P(B)\cdot Q_k(B)$.
Hence
\begin{equation}\tag{20.11}
\begin{aligned}
&k\ge1:&&J\cdot P_k(B)\cdot J^{*}=\mathrm{Qa}'\cdot\mathrm{Qa},\qquad J\cdot P_{k+1}(B)\cdot J^{*}=\mathrm{Qa}'\cdot P\cdot\mathrm{Qa}=\mathrm{Qd}'\cdot\mathrm{Qd};\\
&k=0:&&J\cdot P(B)\cdot J^{*}=J\cdot Q^{T}(-B)\cdot Q(B)\cdot J^{*}=\mathrm{Qd}'\cdot\mathrm{Qd}.
\end{aligned}
\end{equation}
From (20.10) one obtains $P^2=Q^{T}(-B)\cdot\left(Q(B)\cdot Q^{T}(-B)\right)\cdot Q(B)=P$, $P\cdot Q^{T}(-B)=Q^{T}(-B)$ and $Q(B)\cdot P=Q(B)$; moreover $Q(B)$ maps onto $\mathrm{Vl}$, and $Q^{T}(-B)$ is injective.
So $P$ is idempotent of rank $\dim\mathrm{Vl}=4n_{k+1}$.
In the same way, for every $k\ge0$, (20.10) gives $\left(\mathrm{Qd}'\cdot\mathrm{Qd}\right)^2=\mathrm{Qd}'\cdot\left(\mathrm{Qd}\cdot\mathrm{Qd}'\right)\cdot\mathrm{Qd}=\mathrm{Qd}'\cdot\mathrm{Qd}$, where $\mathrm{Qd}$ maps onto $\mathrm{Vl}$ and $\mathrm{Qd}'$ is injective; so $\mathrm{Qd}'\cdot\mathrm{Qd}$ is an idempotent on $\mathrm{Vf}$ of rank $\dim\mathrm{Vl}=4n_{k+1}$.

(n5) \emph{The two fine operators.}
Section~I.3.5 (with $\boldsymbol{\alpha}_{\infty}=b_k$, $\boldsymbol{\beta}_{\infty}=0$; $\boldsymbol{\alpha}_{0}=0$, $\boldsymbol{\beta}_{0}=b_{k+1}/L$), Section~I.3.9 (at $y=0$), (20.11) and (20.8) give
\[
F_{\infty}(B)=\mathrm{Rz}+b_k\cdot\mathrm{Qa}'\cdot\mathrm{Qa}\quad(k\ge1),\qquad F_{0}(B)=\mathrm{Rz}+\mathrm{bq}\cdot\mathrm{Qd}'\cdot\mathrm{Qd}\quad(k\ge0);
\]
at $k=0$ the second operator is $\mathrm{Rz}+b_L\cdot J\cdot P(B)\cdot J^{*}$.
For $k\ge1$ one has $S_k=F_{\infty}(B)^{-1}$, $D_k=b_k-b_k^2\cdot\mathrm{Qa}\cdot S_k\cdot\mathrm{Qa}'$ and $\delta Z_k=\det\left(D_k+b_L\cdot P\right)$; further $\delta Z_0=\det F_{0}(B)$, and $S^{0}_{k+1}=F_{0}(B)^{-1}$ wherever this inverse exists.

\paragraph{2.2. One elimination.}

\begin{named}{Lemma A}
Let $\mathrm{Va}$, $\mathrm{Va}'$ be finite-dimensional complex vector spaces, let $\mathrm{Rx}\colon\mathrm{Va}\to\mathrm{Va}$, $\mathrm{Qx}\colon\mathrm{Va}\to\mathrm{Va}'$ and $\mathrm{Qx}'\colon\mathrm{Va}'\to\mathrm{Va}$ be linear maps, let $\mathrm{Ax}$ be an invertible linear map of $\mathrm{Va}'$, and let $\mathrm{bn}$ be a complex number.
Assume that $\mathrm{Fx}:=\mathrm{Rx}+\mathrm{bn}\cdot\mathrm{Qx}'\cdot\mathrm{Qx}$ is invertible.
Then
\begin{equation}\tag{20.12}
\det\left(\mathrm{Ax}-\mathrm{bn}^2\cdot\mathrm{Qx}\cdot\mathrm{Fx}^{-1}\cdot\mathrm{Qx}'\right)\cdot\det\mathrm{Fx}=\det\mathrm{Ax}\cdot\det\left(\mathrm{Rx}+\mathrm{Qx}'\cdot\left(\mathrm{bn}-\mathrm{bn}^2\cdot\mathrm{Ax}^{-1}\right)\cdot\mathrm{Qx}\right).
\end{equation}
\end{named}

\begin{proof}
Fix bases of $\mathrm{Va}$ and $\mathrm{Va}'$, and consider the linear map of $\mathrm{Va}\oplus\mathrm{Va}'$ whose block matrix has the blocks $A_1=\mathrm{Fx}$, $A_2=-\mathrm{bn}\cdot\mathrm{Qx}'$, $A_3=-\mathrm{bn}\cdot\mathrm{Qx}$, $A_4=\mathrm{Ax}$.
Since $A_1$ is invertible, (cl3) gives for its determinant the value $\det\mathrm{Fx}\cdot\det\left(\mathrm{Ax}-\mathrm{bn}^2\cdot\mathrm{Qx}\cdot\mathrm{Fx}^{-1}\cdot\mathrm{Qx}'\right)$.
Since $A_4$ is invertible, (cl3) gives for the same determinant the value $\det\mathrm{Ax}\cdot\det\left(\mathrm{Fx}-\mathrm{bn}^2\cdot\mathrm{Qx}'\cdot\mathrm{Ax}^{-1}\cdot\mathrm{Qx}\right)$.
Finally, $\mathrm{Fx}-\mathrm{bn}^2\cdot\mathrm{Qx}'\cdot\mathrm{Ax}^{-1}\cdot\mathrm{Qx}=\mathrm{Rx}+\mathrm{Qx}'\cdot\left(\mathrm{bn}-\mathrm{bn}^2\cdot\mathrm{Ax}^{-1}\right)\cdot\mathrm{Qx}$.
\end{proof}

\paragraph{2.3. Proof of (i) ($k\ge1$).}
Apply Lemma~A with $\mathrm{Va}=\mathrm{Vf}$, $\mathrm{Va}'=\mathrm{Vu}$, $\mathrm{Rx}=\mathrm{Rz}$, $\mathrm{Qx}=\mathrm{Qa}$, $\mathrm{Qx}'=\mathrm{Qa}'$, $\mathrm{bn}=b_k$ and $\mathrm{Ax}=b_k\cdot1$ on $\mathrm{Vu}$.
Then $\mathrm{Fx}=F_{\infty}(B)$, which is invertible by hypothesis; moreover $\mathrm{Ax}-\mathrm{bn}^2\cdot\mathrm{Qx}\cdot\mathrm{Fx}^{-1}\cdot\mathrm{Qx}'=D_k$, $\mathrm{bn}-\mathrm{bn}^2\cdot\mathrm{Ax}^{-1}=0$ and $\det\mathrm{Ax}=b_k^{\dim\mathrm{Vu}}=b_k^{4n_k}$.
Thus (20.12) becomes $\det D_k\cdot\det F_{\infty}(B)=b_k^{4n_k}\cdot\det\mathrm{Rz}$.

\paragraph{2.4. Proof of (ii) ($k\ge1$).}
By 2.1\,(n4), $P$ is an idempotent of rank $4n_{k+1}$ on $\mathrm{Vu}$, a space of dimension $4n_k$.
So $\mathrm{Vu}$ is the direct sum of the kernel of $P$, of dimension $4(n_k-n_{k+1})$, and of the image of $P$, of dimension $4n_{k+1}$; on these two subspaces $b_k+b_L\cdot P$ acts as multiplication by $b_k$, resp. by $b_k+b_L$, two positive numbers.
Hence $b_k+b_L\cdot P$ is invertible, and
\begin{equation}\tag{20.13}
\begin{aligned}
&\det\left(b_k+b_L\cdot P\right)=b_k^{4(n_k-n_{k+1})}\cdot(b_k+b_L)^{4n_{k+1}},\\
&\mathrm{Bo}:=\left(b_k+b_L\cdot P\right)^{-1}=b_k^{-1}\cdot(1-P)+(b_k+b_L)^{-1}\cdot P,\\
&b_k-b_k^2\cdot\mathrm{Bo}=b_k-b_k\cdot(1-P)-\frac{b_k^2}{b_k+b_L}\cdot P=\frac{b_k\cdot b_L}{b_k+b_L}\cdot P=\mathrm{bq}\cdot P\qquad(\text{by (20.8)}).
\end{aligned}
\end{equation}
Now apply Lemma~A with $\mathrm{Va}=\mathrm{Vf}$, $\mathrm{Va}'=\mathrm{Vu}$, $\mathrm{Rx}=\mathrm{Rz}$, $\mathrm{Qx}=\mathrm{Qa}$, $\mathrm{Qx}'=\mathrm{Qa}'$, $\mathrm{bn}=b_k$ and $\mathrm{Ax}=b_k+b_L\cdot P$.
Here $\mathrm{Fx}=F_{\infty}(B)$ and $\mathrm{Ax}-\mathrm{bn}^2\cdot\mathrm{Qx}\cdot\mathrm{Fx}^{-1}\cdot\mathrm{Qx}'=D_k+b_L\cdot P$, while (20.13) and (20.11) give $\mathrm{Rx}+\mathrm{Qx}'\cdot\left(\mathrm{bn}-\mathrm{bn}^2\cdot\mathrm{Ax}^{-1}\right)\cdot\mathrm{Qx}=\mathrm{Rz}+\mathrm{bq}\cdot\mathrm{Qa}'\cdot P\cdot\mathrm{Qa}=F_{0}(B)$.
So (20.12) becomes
\[
\delta Z_k\cdot\det F_{\infty}(B)=b_k^{4(n_k-n_{k+1})}\cdot(b_k+b_L)^{4n_{k+1}}\cdot\det F_{0}(B),\qquad F_{0}(B)=\mathrm{Rz}+(b_{k+1}/L)\cdot J\cdot P_{k+1}(B)\cdot J^{*}.
\]

\paragraph{2.5. Proof of (iii) ($k\ge1$, $\delta Z_k\ne0$).}
\emph{Invertibility.}
The left side of (ii) is a product of two non-zero numbers, and the factor in front of $\det F_{0}(B)$ on the right side is positive.
Hence $\det F_{0}(B)\ne0$: the inverse $S^{0}_{k+1}=F_{0}(B)^{-1}$ exists, and $D^{0}_{k+1}(B)=\mathrm{bq}-\mathrm{bq}^2\cdot\mathrm{Qd}\cdot S^{0}_{k+1}\cdot\mathrm{Qd}'$ is defined on $\mathrm{Vl}$.

\emph{The two forms of $D^{0}_{k+1}(B)$.}
From (20.13) and (20.11),
\begin{gather*}
F_{\infty}(B)-b_k^2\cdot\mathrm{Qa}'\cdot\mathrm{Bo}\cdot\mathrm{Qa}=\mathrm{Rz}+\mathrm{Qa}'\cdot\left(b_k-b_k^2\cdot\mathrm{Bo}\right)\cdot\mathrm{Qa}=F_{0}(B),\\
D_k+b_L\cdot P=\mathrm{Bo}^{-1}-b_k^2\cdot\mathrm{Qa}\cdot S_k\cdot\mathrm{Qa}'.
\end{gather*}
Put $\mathrm{Go}:=\mathrm{Bo}+b_k^2\cdot\mathrm{Bo}\cdot\mathrm{Qa}\cdot S^{0}_{k+1}\cdot\mathrm{Qa}'\cdot\mathrm{Bo}$.
Multiplying out,
\[
\begin{split}
\left(D_k+b_L\cdot P\right)\cdot\mathrm{Go}=1&+b_k^2\cdot\mathrm{Qa}\cdot S^{0}_{k+1}\cdot\mathrm{Qa}'\cdot\mathrm{Bo}-b_k^2\cdot\mathrm{Qa}\cdot S_k\cdot\mathrm{Qa}'\cdot\mathrm{Bo}\\
&-b_k^2\cdot\mathrm{Qa}\cdot S_k\cdot\left(b_k^2\cdot\mathrm{Qa}'\cdot\mathrm{Bo}\cdot\mathrm{Qa}\right)\cdot S^{0}_{k+1}\cdot\mathrm{Qa}'\cdot\mathrm{Bo}.
\end{split}
\]
In the last term, $b_k^2\cdot\mathrm{Qa}'\cdot\mathrm{Bo}\cdot\mathrm{Qa}=F_{\infty}(B)-F_{0}(B)$, and $S_k\cdot\left(F_{\infty}(B)-F_{0}(B)\right)\cdot S^{0}_{k+1}=S^{0}_{k+1}-S_k$.
The last term is therefore $-b_k^2\cdot\mathrm{Qa}\cdot S^{0}_{k+1}\cdot\mathrm{Qa}'\cdot\mathrm{Bo}+b_k^2\cdot\mathrm{Qa}\cdot S_k\cdot\mathrm{Qa}'\cdot\mathrm{Bo}$, and $\left(D_k+b_L\cdot P\right)\cdot\mathrm{Go}=1$.
As $\mathrm{Vu}$ is finite-dimensional, $\mathrm{Go}=\left(D_k+b_L\cdot P\right)^{-1}$.
From 2.1\,(n4) and the formula for $\mathrm{Bo}$ one obtains $\mathrm{Bo}\cdot Q^{T}(-B)=(b_k+b_L)^{-1}\cdot Q^{T}(-B)$ and $Q(B)\cdot\mathrm{Bo}=(b_k+b_L)^{-1}\cdot Q(B)$.
Together with (20.10) and (20.8) this gives
\begin{equation}\tag{20.14}
\begin{aligned}
b_L-b_L^2\cdot Q(B)\cdot\left(D_k+b_L\cdot P\right)^{-1}\cdot Q^{T}(-B)&=b_L-b_L^2\cdot(b_k+b_L)^{-1}\cdot Q(B)\cdot Q^{T}(-B)\\
&\qquad-\left(\frac{b_L\cdot b_k}{b_k+b_L}\right)^2\cdot Q(B)\cdot\mathrm{Qa}\cdot S^{0}_{k+1}\cdot\mathrm{Qa}'\cdot Q^{T}(-B)\\
&=\frac{b_L\cdot b_k}{b_k+b_L}-\mathrm{bq}^2\cdot\mathrm{Qd}\cdot S^{0}_{k+1}\cdot\mathrm{Qd}'\\
&=\mathrm{bq}-\mathrm{bq}^2\cdot\mathrm{Qd}\cdot S^{0}_{k+1}\cdot\mathrm{Qd}'=D^{0}_{k+1}(B).
\end{aligned}
\end{equation}

\emph{The determinant.}
Apply Lemma~A with $\mathrm{Va}=\mathrm{Vf}$, $\mathrm{Va}'=\mathrm{Vl}$, $\mathrm{Rx}=\mathrm{Rz}$, $\mathrm{Qx}=\mathrm{Qd}$, $\mathrm{Qx}'=\mathrm{Qd}'$, $\mathrm{bn}=\mathrm{bq}$ and $\mathrm{Ax}=\mathrm{bq}\cdot1$ on $\mathrm{Vl}$.
Here $\mathrm{Fx}=F_{0}(B)$, $\mathrm{Ax}-\mathrm{bn}^2\cdot\mathrm{Qx}\cdot\mathrm{Fx}^{-1}\cdot\mathrm{Qx}'=D^{0}_{k+1}(B)$, $\mathrm{bn}-\mathrm{bn}^2\cdot\mathrm{Ax}^{-1}=0$ and $\det\mathrm{Ax}=\mathrm{bq}^{4n_{k+1}}$, so that
\begin{equation}\tag{20.15}
\det D^{0}_{k+1}(B)\cdot\det F_{0}(B)=\mathrm{bq}^{4n_{k+1}}\cdot\det\mathrm{Rz}.
\end{equation}
The numbers $\det F_{\infty}(B)$, $\delta Z_k$ and $\det F_{0}(B)$ are different from zero.
Dividing (i) by (ii) gives
\[
\frac{\det D_k}{\delta Z_k}=\left(\frac{b_k}{b_k+b_L}\right)^{4n_{k+1}}\cdot\frac{\det\mathrm{Rz}}{\det F_{0}(B)}.
\]
Insert this into (20.15) and use $\mathrm{bq}\cdot(b_k+b_L)/b_k=b_L$, which holds by (20.8):
\begin{gather*}
\det D^{0}_{k+1}(B)=\left(\frac{\mathrm{bq}\cdot(b_k+b_L)}{b_k}\right)^{4n_{k+1}}\cdot\frac{\det D_k}{\delta Z_k}=b_L^{4n_{k+1}}\cdot\frac{\det D_k}{\delta Z_k},\\
L^{4n_{k+1}}\cdot\det D^{0}_{k+1}(B)=b^{4n_{k+1}}\cdot\frac{\det D_k}{\delta Z_k}.
\end{gather*}

\paragraph{2.6. Proof of (iii0) ($k=0$, $\delta Z_0\ne0$).}
At level $0$ the space $\mathrm{Vf}=\mathfrak{S}^{\mathrm{ch}}(1,\ell)$ has dimension $192\,n_0$, and the operators $D_0=\mathrm{Rz}$ and $F_{0}(B)=\mathrm{Rz}+b_L\cdot\mathrm{Qd}'\cdot\mathrm{Qd}=\mathrm{Rz}+b_L\cdot J\cdot P(B)\cdot J^{*}$ act on $\mathrm{Vf}$.
The map $\mathrm{Qd}=Q(B)\cdot J^{*}$ contains the mean $\frac{1}{48}\sum_{g}$ over each cluster, and $\mathrm{Vl}$ has dimension $4n_1$.
Since $\delta Z_0=\det F_{0}(B)\ne0$, the operator $F_{0}(B)$ is invertible.
Apply Lemma~A with $\mathrm{Va}=\mathrm{Vf}$, $\mathrm{Va}'=\mathrm{Vl}$, $\mathrm{Rx}=\mathrm{Rz}$, $\mathrm{Qx}=\mathrm{Qd}$, $\mathrm{Qx}'=\mathrm{Qd}'$, $\mathrm{bn}=b_L$ and $\mathrm{Ax}=b_L\cdot1$ on $\mathrm{Vl}$.
Then $\mathrm{Fx}=F_{0}(B)$,
\[
\mathrm{Ax}-\mathrm{bn}^2\cdot\mathrm{Qx}\cdot\mathrm{Fx}^{-1}\cdot\mathrm{Qx}'=b_L-b_L^2\cdot Q(B)\cdot J^{*}\cdot F_{0}(B)^{-1}\cdot J\cdot Q^{T}(-B)=D^{0}_1(B),
\]
$\mathrm{bn}-\mathrm{bn}^2\cdot\mathrm{Ax}^{-1}=0$ and $\det\mathrm{Ax}=b_L^{4n_1}$.
Consequently
\[
\det D^{0}_1(B)\cdot\delta Z_0=b_L^{4n_1}\cdot\det\mathrm{Rz}=b_L^{4n_1}\cdot\det D_0,\qquad L^{4n_1}\cdot\det D^{0}_1(B)=b^{4n_1}\cdot\frac{\det D_0}{\delta Z_0}.
\]
By (20.8), $\mathrm{bq}=b_L$ at $k=0$; hence $D^{0}_1(B)=\mathrm{bq}-\mathrm{bq}^2\cdot\mathrm{Qd}\cdot S^{0}_1\cdot\mathrm{Qd}'$ with $S^{0}_1=F_{0}(B)^{-1}$.
Written with $\mathrm{Qd}$, $\mathrm{Qd}'$ and $\mathrm{bq}$, the formula for $D^{0}_{k+1}(B)$ is thus the same for all $k\ge0$.

\paragraph{2.7. Scaling: the first display of (iv) ($k\ge0$).}
(s0) \emph{The side.}
By the definition of $\mathrm{Cl}$ in Section~I.1.1, $\ell/(LM)$ is an odd integer $\ge13$.
Hence $\ell/L=M\cdot\ell/(LM)$ is an odd multiple of $M$, in particular an integer, and $\ell/L\ge13M$.
Of the side $\ell/L$, the displays of level $k+1$ in the statement need only that it is a positive integer.
Then $(\ell/L)\mathbb{Z}^3$ is contained in $\mathbb{Z}^3$, hence in $L^j\eta_{k+1}\mathbb{Z}^3$ for $0\le j\le k+1$, and the translations by $(\ell/L)\mathbb{Z}^3$ carry $\mathbb{L}_{k+1}$, the blocks of the sides $L^j\cdot\eta_{k+1}$ and their centres to objects of the same kind.
No $L$-block of the unit lattice of level $k+1$ enters.
The operator $P_{k+1}(B'):=Q_{k+1}^{T}(-B')\cdot Q_{k+1}(B')$ is the operator $P_k(\,\cdot\,)$ of Section~I.3.3 at level $k+1$; it is distinguished from $P_{k+1}(B)=Q_k^{T}(-B)\cdot P(B)\cdot Q_k(B)$ by its argument.

(s1) \emph{The maps.}
One has $\eta_{k+1}=\eta_k/L$.
The map $c\mapsto c/L$ is a bijection of $\mathbb{L}_k$ onto $\mathbb{L}_{k+1}$, of $L\mathbb{Z}^3$ onto $\mathbb{Z}^3$ and of $\ell\mathbb{Z}^3$ onto $(\ell/L)\mathbb{Z}^3$; moreover $x/L=(c/L,g)$ (Section~I.3.7).
Put
\begin{gather*}
\mathrm{Scf}\,f':=f'_L,\quad f'_L(x)=L^{-1}\cdot f'(x/L)\quad\left(f'\in\mathfrak{S}^{\mathrm{ch}}(\eta_{k+1},\ell/L)\right);\\
\mathrm{Scv}\,v:=v_L,\quad v_L(y')=L^{-1}\cdot v(y'/L)\quad\left(v\in\mathfrak{S}(1,\ell/L),\ y'\in L\mathbb{Z}^3\right).
\end{gather*}
One has $f'_L(c+\ell e_{\mu},g)=L^{-1}\cdot f'\left(c/L+(\ell/L)e_{\mu},g\right)=-f'_L(c,g)$ and $v_L(y'+\ell e_{\mu})=L^{-1}\cdot v\left(y'/L+(\ell/L)e_{\mu}\right)=-v_L(y')$, so that $f'_L\in\mathrm{Vf}$ and $v_L\in\mathrm{Vl}$.
The maps $f\mapsto\left((c',g)\mapsto L\cdot f(Lc',g)\right)$ and $v'\mapsto\left(y\mapsto L\cdot v'(Ly)\right)$ are inverse to $\mathrm{Scf}$ and $\mathrm{Scv}$.
Hence $\mathrm{Scf}\colon\mathfrak{S}^{\mathrm{ch}}(\eta_{k+1},\ell/L)\to\mathrm{Vf}$ and $\mathrm{Scv}\colon\mathfrak{S}(1,\ell/L)\to\mathrm{Vl}$ are linear bijections.

(s2) \emph{Bonds.}
The map $b\mapsto b/L$ carries the bonds of $\mathbb{L}_k$ onto the bonds of $\mathbb{L}_{k+1}$.
By Section~I.3.7, $B(b)=B'_L(b)=L^{-\frac{1}{2}}\cdot B'(b/L)$; and $B$ is $\ell\mathbb{Z}^3$-periodic because $B'$ is $(\ell/L)\mathbb{Z}^3$-periodic.
Moreover $e_{k+1}\cdot\eta_{k+1}=L^{\frac{1}{2}}\cdot e_k\cdot\eta_k/L=L^{-\frac{1}{2}}\cdot e_k\cdot\eta_k$.
Therefore
\begin{equation}\tag{20.16}
e_k\cdot\eta_k\cdot B(b)=e_{k+1}\cdot\eta_{k+1}\cdot B'(b/L)\qquad\text{for every bond $b$ of $\mathbb{L}_k$}.
\end{equation}
The charges enter the scaling at this point only, and only through the relation $e_{k+1}=L^{\frac{1}{2}}\cdot e_k$.

(s3) \emph{The operator.}
Let $x=(c,g)\in\mathbb{L}^{\mathrm{f}}_k$.
Formula (I.1.2) at the two levels gives
\begin{gather*}
x^{(i)}/L=(c/L,g\,s_i)=(x/L)^{(i)}\quad(i=0,1,2),\\
x^{(3)}/L=\left((c+\eta_k\cdot g\,e_1)/L,g\,s_3\right)=\left(c/L+\eta_{k+1}\cdot g\,e_1,g\,s_3\right)=(x/L)^{(3)}.
\end{gather*}
The matrix $M_{i}(x)=\frac{w_i}{2}\cdot\left(\rho+\gamma\cdot(g\,\hat n_i)\right)$ depends on the label only (I.2.1), so that $M_{i}(x/L)=M_{i}(x)$.
For $i\le2$ one has $U_B(x,x^{(i)})=1=U_{B'}\left(x/L,(x/L)^{(i)}\right)$, and by (20.16)
\[
U_B(x,x^{(3)})=\exp\left(i\,e_k\cdot\eta_k\cdot B(c,c_{3})\right)=\exp\left(i\,e_{k+1}\cdot\eta_{k+1}\cdot B'(c/L,c_{3}/L)\right)=U_{B'}\left(x/L,(x/L)^{(3)}\right).
\]
Substituting $f'_L(x^{(i)})=L^{-1}\cdot f'\left((x/L)^{(i)}\right)$, $f'_L(x)=L^{-1}\cdot f'(x/L)$ and $\eta_k^{-1}=L^{-1}\cdot\eta_{k+1}^{-1}$ into (I.2.3), one obtains
\begin{align*}
\left(D_Bf'_L\right)(x)&=-\eta_k^{-1}\cdot\sum_{i=0}^{3}M_{i}(x)\cdot\left[U_B(x,x^{(i)})\cdot f'_L(x^{(i)})-f'_L(x)\right]\\
&=L^{-1}\cdot L^{-1}\cdot\left\{-\eta_{k+1}^{-1}\cdot\sum_{i=0}^{3}M_{i}(x/L)\cdot\left[U_{B'}\left(x/L,(x/L)^{(i)}\right)\cdot f'\left((x/L)^{(i)}\right)-f'(x/L)\right]\right\}\\
&=L^{-1}\cdot L^{-1}\cdot\left(D_{B'}f'\right)(x/L)=L^{-1}\cdot\left(D_{B'}f'\right)_L(x).
\end{align*}
Further, $\mathrm{zc}\cdot f'_L=L^{-1}\cdot\left(L\cdot\mathrm{zc}\cdot f'\right)_L=L^{-1}\cdot\left(\mathrm{zc}'\cdot f'\right)_L$.
Therefore
\begin{equation}\tag{20.17}
\left(D_B+\mathrm{zc}\right)f'_L=L^{-1}\cdot\left(\left(D_{B'}+\mathrm{zc}'\right)f'\right)_L,\qquad\text{that is,}\qquad\mathrm{Rz}\cdot\mathrm{Scf}=L^{-1}\cdot\mathrm{Scf}\cdot\left(D_{B'}+\mathrm{zc}'\right)
\end{equation}
(compare the first relation of (I.O9)).

(s4) \emph{Paths.}
For $c,\tilde{c}\in\mathbb{L}_k$ and an order $\pi$ of the coordinates, the sites of $\boldsymbol{\Gamma}^{\pi}(c,\tilde{c})$, divided by $L$, are the sites of $\boldsymbol{\Gamma}^{\pi}(c/L,\tilde{c}/L)$ in $\mathbb{L}_{k+1}$, in the same order (Section~I.1.1).
The path sums carry no weights.
Hence $B\left(\boldsymbol{\Gamma}^{\pi}(c,\tilde{c})\right)=L^{-\frac{1}{2}}\cdot B'\left(\boldsymbol{\Gamma}^{\pi}(c/L,\tilde{c}/L)\right)$ and $(\tau B)(c,\tilde{c})=L^{-\frac{1}{2}}\cdot(\tau B')(c/L,\tilde{c}/L)$, both sides vanishing for $c=\tilde{c}$; with $L^{-\frac{1}{2}}\cdot e_k\cdot\eta_k=e_{k+1}\cdot\eta_{k+1}$ from (s2) this gives
\begin{equation}\tag{20.18}
e_k\cdot\eta_k\cdot(\tau B)(c,\tilde{c})=e_{k+1}\cdot\eta_{k+1}\cdot(\tau B')(c/L,\tilde{c}/L).
\end{equation}

(s5) \emph{Centres of blocks.}
For $c\in\mathbb{L}_k$ and $0\le j\le k+1$, let $\tilde{z}_j(c)$ be the centre of the block of side $L^j\cdot\eta_k$ that contains $c$, that is, the point of $L^j\eta_k\mathbb{Z}^3$ with $|c-\tilde{z}_j(c)|<L^j\cdot\eta_k/2$.
This point exists and is unique.
Indeed, every coordinate of $c$ is an integer multiple of $\eta_k$, whereas the numbers $L^j\cdot\eta_k\cdot\left(i+\frac{1}{2}\right)$, $i\in\mathbb{Z}$, are not, because $L^j$ is odd; so each coordinate of $c$ lies in exactly one of the open intervals $L^j\cdot\eta_k\cdot\left(i-\frac{1}{2},i+\frac{1}{2}\right)$.
One has $\tilde{z}_0(c)=c$ and $\tilde{z}_k(c)=[c]$, and $\tilde{z}_{k+1}(c)$ lies in $L\mathbb{Z}^3$.

The vector $\tilde{z}_{j+1}(c)-\tilde{z}_j(c)$ lies in $L^j\eta_k\mathbb{Z}^3$, and its sup norm is smaller than $(L^{j+1}+L^j)\cdot\eta_k/2=\frac{L+1}{2}\cdot L^j\cdot\eta_k$, where $\frac{L+1}{2}$ is an integer.
Hence
\begin{equation}\tag{20.19}
|\tilde{z}_{j+1}(c)-\tilde{z}_j(c)|\le\frac{L-1}{2}\cdot L^j\cdot\eta_k\qquad\left(<\tfrac{1}{2}\ \text{for }j<k;\ \le\tfrac{L-1}{2}\ \text{for }j=k\right).
\end{equation}
In particular $[c]\in B\left(\tilde{z}_{k+1}(c)\right)$.
Conversely, let $y'\in L\mathbb{Z}^3$, $y\in B(y')$ and $c\in\Delta_y$; then $|c-y'|<\frac{1}{2}+\frac{L-1}{2}=\frac{L}{2}$, so that $y'=\tilde{z}_{k+1}(c)$ and $y=\tilde{z}_k(c)$.
Consequently, for every $y'\in L\mathbb{Z}^3$,
\begin{equation}\tag{20.20}
\begin{aligned}
&(y,c)\mapsto c\ \text{ is a bijection of }\ \{(y,c):y\in B(y'),\ c\in\Delta_y\}\\
&\text{onto }\ \mathrm{Bl}_{k+1}(y')=\{c\in\mathbb{L}_k:|c-y'|<L/2\}.
\end{aligned}
\end{equation}
For $k=0$ one has $\Delta_y=\{y\}$ and $\mathrm{Bl}_1(y')=B(y')$.

For $c'\in\mathbb{L}_{k+1}$ and $0\le j\le k+1$, let $\tilde{z}'_j(c')$ be the centre of the block of side $L^j\cdot\eta_{k+1}$ that contains $c'$.
Since $\tilde{z}_j(c)/L$ lies in $L^j\eta_{k+1}\mathbb{Z}^3$ and $|c/L-\tilde{z}_j(c)/L|<L^j\cdot\eta_{k+1}/2$, uniqueness gives
\begin{equation}\tag{20.21}
\tilde{z}'_j(c/L)=\tilde{z}_j(c)/L\qquad(0\le j\le k+1).
\end{equation}

(s6) \emph{Phases.}
All paths are formed on the covering lattices, between the points named, with the periodic extensions of $B$ and $B'$.
By (20.19) these points have sup distance at most $(L-1)/2$ at level $k$.
This number is smaller than $\ell/2$, as $\ell\ge13LM$, and at most $M$, as $M\ge L^2$ (Section~I.0.5); so $U_B(y',y)$ is the number (I.1.6), which is defined for $|c-c'|\le M$.
After division by $L$ the points have sup distance $<\frac{1}{2}\le(\ell/L)/2$ at level $k+1$.
They are therefore the lifts prescribed in Sections~I.1.1 and I.3.2.
Let $y'\in L\mathbb{Z}^3$, $c\in\mathrm{Bl}_{k+1}(y')$ and $y=[c]$.
By (I.3.3) and (I.1.6):

\begin{align*}
\varphi^{+}_{B}(y',c)=U_B(y',y)\cdot\varphi_B(y,c)
&=\exp\left(i\,e_k\eta_k\cdot\left[(\tau B)(\tilde{z}_{k+1}(c),\tilde{z}_k(c))+\sum_{j<k}(\tau B)(\tilde{z}_{j+1}(c),\tilde{z}_j(c))\right]\right)\\
&=\exp\left(i\,e_k\eta_k\cdot\sum_{j=0}^{k}(\tau B)(\tilde{z}_{j+1}(c),\tilde{z}_j(c))\right).
\end{align*}
At level~$0$ the sum over $j<k$ is empty, $\varphi_B(c,c)=1$, and $\varphi^{+}_B(y',c)=U_B(y',c)$.
At the level $k+1$, for a unit site $y_1\in\mathbb{Z}^3$ and a site $c'\in\mathbb{L}_{k+1}$ with $|c'-y_1|<1/2$, the display (I.3.3) takes the form
\begin{equation*}
\varphi_{B'}(y_1,c')=\exp\left(i\,e_{k+1}\eta_{k+1}\cdot\sum_{j=0}^{k}(\tau B')(\tilde{z}'_{j+1}(c'),\tilde{z}'_j(c'))\right).
\end{equation*}
Taking $c'=c/L$ and $y_1=y'/L$ and comparing the exponents by means of (20.18) and (20.21), one obtains
\begin{equation*}
\varphi^{+}_B(y',c)=\varphi_{B'}(y'/L,c/L)\qquad\left(y'\in L\mathbb{Z}^3,\ c\in\mathrm{Bl}_{k+1}(y')\right).
\tag{20.22}
\end{equation*}

(s7) \emph{The averages.}
Let $f'\in\mathfrak{S}^{\mathrm{ch}}(\eta_{k+1},\ell/L)$, $v\in\mathfrak{S}(1,\ell/L)$ and $y'\in L\mathbb{Z}^3$.
By (I.3.1),
\begin{equation*}
\left(J^{*}f'_L\right)(c)=\frac{1}{48}\sum_{g\in O_{\mathrm{cub}}}L^{-1}\cdot f'(c/L,g)=L^{-1}\cdot\left(J^{*}f'\right)(c/L).
\tag{20.23}
\end{equation*}
Let first $k\ge1$.
The map $c\mapsto c/L$ carries $\mathrm{Bl}_{k+1}(y')$ onto the set $\{c'\in\mathbb{L}_{k+1}:|c'-y'/L|<1/2\}$, which is the unit block of $y'/L$ at the level $k+1$.
The displays of Section~I.3.3, (20.20), the identity $L^{-3}\cdot\eta_k^3=\eta_{k+1}^3$, (20.23) and (20.22) therefore give
\begin{align*}
\left(\mathrm{Qd}\,f'_L\right)(y')
&=L^{-3}\sum_{y\in B(y')}U_B(y',y)\cdot\eta_k^3\sum_{c\in\Delta_y}\varphi_B(y,c)\cdot\left(J^{*}f'_L\right)(c)\\
&=\eta_{k+1}^3\sum_{c\in\mathrm{Bl}_{k+1}(y')}\varphi^{+}_B(y',c)\cdot L^{-1}\cdot\left(J^{*}f'\right)(c/L)\\
&=L^{-1}\cdot\eta_{k+1}^3\sum_{c':\,|c'-y'/L|<1/2}\varphi_{B'}(y'/L,c')\cdot\left(J^{*}f'\right)(c')\\
&=L^{-1}\cdot\left(Q_{k+1}(B')\cdot J^{*}f'\right)(y'/L)=\left(Q_{k+1}(B')\cdot J^{*}f'\right)_L(y').
\end{align*}
At level~$0$ one has $\eta_1^3=L^{-3}$, $\mathrm{Bl}_1(y')=B(y')$ and $\varphi^{+}_B(y',c)=U_B(y',c)$, and the display of $Q(B)$ in Section~I.3.3, (20.23) and (20.22) give
\begin{align*}
\left(\mathrm{Qd}\,f'_L\right)(y')
&=\left(Q(B)\cdot J^{*}f'_L\right)(y')=L^{-3}\sum_{c\in B(y')}U_B(y',c)\cdot\left(J^{*}f'_L\right)(c)\\
&=\eta_1^3\sum_{c\in\mathrm{Bl}_1(y')}\varphi^{+}_B(y',c)\cdot L^{-1}\cdot\left(J^{*}f'\right)(c/L)\\
&=L^{-1}\cdot\eta_1^3\sum_{c':\,|c'-y'/L|<1/2}\varphi_{B'}(y'/L,c')\cdot\left(J^{*}f'\right)(c')\\
&=L^{-1}\cdot\left(Q_1(B')\cdot J^{*}f'\right)(y'/L)=\left(Q_1(B')\cdot J^{*}f'\right)_L(y').
\end{align*}
Let now $x=(c,g)$ with $c\in\mathrm{Bl}_{k+1}(y')$, put $y=[c]$, and let $k\ge1$.
Since $(J\xi)(c,g)=\xi(c)$ by (I.3.1), the displays of $Q_k^{T}(-B)$ and $Q^{T}(-B)$ in Section~I.3.3 and (20.22) give
\begin{align*}
\left(\mathrm{Qd}'\,v_L\right)(x)
&=\left(Q_k^{T}(-B)\cdot Q^{T}(-B)\,v_L\right)(c)=\varphi_B(y,c)^{-1}\cdot U_B(y',y)^{-1}\cdot v_L(y')\\
&=\varphi^{+}_B(y',c)^{-1}\cdot L^{-1}\cdot v(y'/L)=L^{-1}\cdot\varphi_{B'}(y'/L,c/L)^{-1}\cdot v(y'/L)\\
&=L^{-1}\cdot\left(J\cdot Q_{k+1}^{T}(-B')\,v\right)(x/L)=\left(J\cdot Q_{k+1}^{T}(-B')\,v\right)_L(x).
\end{align*}
At level~$0$, for $x=(c,g)$ with $c\in B(y')$, the relation $(J\xi)(c,g)=\xi(c)$, the display of $Q^{T}(-B)$ in Section~I.3.3 and (20.22) give
\begin{align*}
\left(\mathrm{Qd}'\,v_L\right)(x)
&=\left(Q^{T}(-B)\,v_L\right)(c)=U_B(y',c)^{-1}\cdot v_L(y')=\varphi^{+}_B(y',c)^{-1}\cdot L^{-1}\cdot v(y'/L)\\
&=L^{-1}\cdot\varphi_{B'}(y'/L,c/L)^{-1}\cdot v(y'/L)=L^{-1}\cdot\left(J\cdot Q_1^{T}(-B')\,v\right)(x/L)=\left(J\cdot Q_1^{T}(-B')\,v\right)_L(x).
\end{align*}
Consequently, for every $k\ge0$,
\begin{equation*}
\mathrm{Qd}\cdot\mathrm{Scf}=\mathrm{Scv}\cdot Q_{k+1}(B')\cdot J^{*},\qquad\mathrm{Qd}'\cdot\mathrm{Scv}=\mathrm{Scf}\cdot J\cdot Q_{k+1}^{T}(-B').
\tag{20.24}
\end{equation*}
The relations (20.17) and (20.24) are the three relations of the first display of (iv).

\paragraph{2.8. The second display of (iv) ($k\ge0$).}
At the level $k+1$, which is at least $1$, Section~I.3.5 with $\mathrm{zc}'$ in the place of the mass gives
\begin{gather*}
F_\infty(B')=D_{B'}+\mathrm{zc}'+b_{k+1}\cdot J\cdot P_{k+1}(B')\cdot J^{*},\qquad J\cdot P_{k+1}(B')\cdot J^{*}=\left(J\cdot Q_{k+1}^{T}(-B')\right)\cdot\left(Q_{k+1}(B')\cdot J^{*}\right),\\
S_{k+1}(B')=F_\infty(B')^{-1},\qquad D_{k+1}(B')=b_{k+1}-b_{k+1}^2\cdot Q_{k+1}(B')\cdot J^{*}\cdot S_{k+1}(B')\cdot J\cdot Q_{k+1}^{T}(-B').
\end{gather*}
By (20.24),
\begin{equation*}
\mathrm{Qd}'\cdot\mathrm{Qd}\cdot\mathrm{Scf}=\mathrm{Qd}'\cdot\mathrm{Scv}\cdot Q_{k+1}(B')\cdot J^{*}=\mathrm{Scf}\cdot J\cdot P_{k+1}(B')\cdot J^{*}.
\end{equation*}
Together with (20.17) and $\mathrm{bq}=b_{k+1}/L$ this yields
\begin{align*}
F_0(B)\cdot\mathrm{Scf}&=\left(\mathrm{Rz}+\mathrm{bq}\cdot\mathrm{Qd}'\cdot\mathrm{Qd}\right)\cdot\mathrm{Scf}\\
&=L^{-1}\cdot\mathrm{Scf}\cdot\left(D_{B'}+\mathrm{zc}'+b_{k+1}\cdot J\cdot P_{k+1}(B')\cdot J^{*}\right)=L^{-1}\cdot\mathrm{Scf}\cdot F_\infty(B').
\tag{20.25}
\end{align*}
Since $\mathrm{Scf}$ is bijective, $F_0(B)$ is invertible if and only if $F_\infty(B')$ is invertible, that is, if and only if $S_{k+1}(B')$ exists.
Here $F_0(B)=\mathrm{Rz}+(b_{k+1}/L)\cdot J\cdot P_{k+1}(B)\cdot J^{*}$ for $k\ge1$, and $F_0(B)=\mathrm{Rz}+b_L\cdot J\cdot P(B)\cdot J^{*}$ for $k=0$ (2.1\,(n5)).
When these operators are invertible, (20.25) gives $S^{0}_{k+1}(B)\cdot\mathrm{Scf}=L\cdot\mathrm{Scf}\cdot S_{k+1}(B')$, that is,
\begin{equation*}
S^{0}_{k+1}(B)\,f'_L=L\cdot\left(S_{k+1}(B')\,f'\right)_L.
\end{equation*}
From (20.24) and this relation,
\begin{align*}
D^{0}_{k+1}(B)\cdot\mathrm{Scv}
&=\mathrm{bq}\cdot\mathrm{Scv}-\mathrm{bq}^2\cdot\mathrm{Qd}\cdot S^{0}_{k+1}(B)\cdot\mathrm{Qd}'\cdot\mathrm{Scv}\\
&=\mathrm{bq}\cdot\mathrm{Scv}-\mathrm{bq}^2\cdot\mathrm{Qd}\cdot S^{0}_{k+1}(B)\cdot\mathrm{Scf}\cdot J\cdot Q_{k+1}^{T}(-B')\\
&=\mathrm{bq}\cdot\mathrm{Scv}-\mathrm{bq}^2\cdot L\cdot\mathrm{Qd}\cdot\mathrm{Scf}\cdot S_{k+1}(B')\cdot J\cdot Q_{k+1}^{T}(-B')\\
&=\mathrm{Scv}\cdot\left[\frac{b_{k+1}}{L}-\frac{b_{k+1}^2}{L}\cdot Q_{k+1}(B')\cdot J^{*}\cdot S_{k+1}(B')\cdot J\cdot Q_{k+1}^{T}(-B')\right]=L^{-1}\cdot\mathrm{Scv}\cdot D_{k+1}(B'),
\end{align*}
that is, $D^{0}_{k+1}(B)\,v_L=L^{-1}\cdot\left(D_{k+1}(B')\,v\right)_L$.
(At $k=0$, $D^{0}_1(B)$ is the operator of (iii0); see the end of 2.6.)
Hence $D^{0}_{k+1}(B)=\mathrm{Scv}\cdot\left(L^{-1}\cdot D_{k+1}(B')\right)\cdot\mathrm{Scv}^{-1}$.
The determinant of a linear map is invariant under conjugation by a linear bijection, and $\mathfrak{S}(1,\ell/L)$ has the dimension $4(\ell/L)^3=4n_{k+1}$.
Therefore
\begin{equation*}
\det D^{0}_{k+1}(B)=L^{-4n_{k+1}}\cdot\det D_{k+1}(B').
\tag{20.26}
\end{equation*}

\paragraph{2.9. Proof of (6.1).}
Suppose that $S_k(B)$ exists (when $k\ge1$) and that $\delta Z_k(B)\ne0$.
By (iii), resp.\ (iii0), the operator $F_0(B)$ is invertible and
\begin{equation*}
L^{4n_{k+1}}\cdot\det D^{0}_{k+1}(B)=b^{4n_{k+1}}\cdot\frac{\det D_k(B)}{\delta Z_k(B)},
\end{equation*}
where $D_0(B)=\mathrm{Rz}=D_B+\mathrm{zc}$ at $k=0$.
By 2.8 the operator $S_{k+1}(B')$ exists, $D_{k+1}(B')$ is defined, and (20.26) holds.
It follows that
\begin{equation*}
\det D_{k+1}(B')=L^{4n_{k+1}}\cdot\det D^{0}_{k+1}(B)=b^{4n_{k+1}}\cdot\frac{\det D_k(B)}{\delta Z_k(B)}.
\end{equation*}

\paragraph{2.10. The exponents.}
It remains to prove the last sentence of Proposition~20.2.
The number $4n_k=\dim\mathrm{Vu}$ enters through the determinant of the multiple $b_k$ of the identity on $\mathrm{Vu}$ in 2.3 and through $\det\left(b_k+b_L\cdot P\right)$ in 2.4.
The number $4n_{k+1}$ enters as the rank of $P$ in 2.4, as $\dim\mathrm{Vl}$ through the determinant of the multiple $\mathrm{bq}$, resp.\ $b_L$, of the identity on $\mathrm{Vl}$ in 2.5 and 2.6, and as $\dim\mathfrak{S}(1,\ell/L)$ in (20.26); at $k=0$ it is also the rank of $J\cdot P(B)\cdot J^{*}=\mathrm{Qd}'\cdot\mathrm{Qd}$ (2.1\,(n4)).
On $\mathrm{Vf}$ the relation (20.25) gives the existence of the inverse $S_{k+1}(B')$ and its form $S^{0}_{k+1}(B)\cdot\mathrm{Scf}=L\cdot\mathrm{Scf}\cdot S_{k+1}(B')$, and no multiple of the identity on $\mathrm{Vf}$ is taken out of a determinant.
Hence, at $k=0$ as at $k\ge1$, the exponents in (i)--(iii), (iii0), (20.26) and (6.1) are formed from $4n_k$ and $4n_{k+1}$, and the dimension $192(\ell/\eta_k)^3$ of $\mathrm{Vf}$ does not enter them.

\paragraph{2.11. Uniformity.}
The assertions of Proposition~20.2 are identities and assertions that inverses exist; none of them is an estimate.
Of the constants of Section~I.8 the identities contain only $b$ and $L$; $\rho$ and the weights $w_i$ enter through $D_B$ alone, and $M$ enters only through the range $|c-c'|\le M$ of the symbol $U_B$ in (I.1.6) and through the class $\mathrm{Cl}$ of the sides.
The identities hold for every $k\ge0$, every $\ell\in\mathrm{Cl}$, every $e_k>0$, every complex $\mathrm{zc}$ and all complex $B$, $B'$ that satisfy the suppositions of the respective clauses.
\qed

\subsection*{3. Proposition 20.3}

\subsubsection*{Proof}

\paragraph{3.1. The data of the levels.}
In \cite[Section~19.2]{QEDI} one takes $h'=h$ and $m_0=0$, with $1$ as the upper end of the mass range (Section~I.6.1); the number $e_1$ there is the number $e_1$ of Proposition~14.2\,(a), line (C9).
By (G2.S-4), $e_K=e\cdot h^{\frac{1}{2}}\le e_1$, $R\in\mathrm{Tor}(e)$ and $N\ge N_{\min}(R,e)$.
By Proposition~14.3\,(f) the rows (e1)--(e8) of Proposition~14.2\,(b) hold for all $v\ge-\log e_1$, and $v_K=-\log e_K\ge-\log e_1$; so the row (e1) holds at $v=v_K$.
With it, (G4.Sc6) of Section~G4.0.2, that is Proposition~14.2\,(c), (f), gives $v_K\ge8$, $(1+\mathrm{lg}/v_K)^r\le(1+3\,\mathrm{lg}/v_K)^r<L$ and, for the given $R$ and $N$, the condition (G4.2.1) of Section~G4.0.3, namely $r_*>31L+16$.
These three assertions rest on the row (e1), which is the condition (e1) of \cite[Section~19.1.2]{QEDI}.
They are the hypotheses of \cite[Proposition~19.1]{QEDI} and of its clause (f).
The clauses (e) and (f) of that proposition give, for $0\le j\le k_*(R)$,
\begin{equation*}
\ell_j\in\mathrm{Cl},\qquad e_j\le e_*\le e_K,\qquad 0<\mu_j\le\mu_*\le\mu_K\le1.
\end{equation*}
Further, $e_K\le e_1$; $e_1\le e_0$ by Proposition~14.3\,(f); and $e_0<1$ by the first inequality of the row (A8) of Proposition~14.2\,(b) at $v=-\log e_0$, which reads $e_0^{\varepsilon}\le1/4$, together with $\varepsilon>0$.
Hence $0<e_j\le e_0<1$.
(At $j=0$ the symbol $e_j$ denotes $q=e\cdot\eta^{\frac{1}{2}}$, and $\mu_j$ denotes $\bar{m}\cdot\eta$ (Sections~I.2.5 and~I.6.1); for $j=1,2$ the symbol $e_j$ denotes the running charge $e\cdot a_j^{\frac{1}{2}}$.
In this section the symbols $e_0$, $e_1$, where they are written as such, denote the ceilings of Section~I.8.)
By Section~I.6.1 one has $a_{j+1}=L\cdot a_j$; hence
\begin{equation*}
\begin{aligned}
&\ell_{j+1}=\frac{\ell_j}{L},\qquad n_{j+1}=\frac{n_j}{L^3},\qquad e_{j+1}=L^{\frac{1}{2}}\cdot e_j,\\
&\mu_{j+1}=L\cdot\mu_j,\qquad\ell_0=\frac{R}{\eta}=n,\qquad L^j\cdot\ell_j=n.
\end{aligned}
\tag{20.27}
\end{equation*}
For the $k\le k_*(R)$ of the data two consequences follow.
\begin{itemize}
\item[(D1)] For $0\le j\le k$ the quadruple $(j,\ell_j,\mu_j,e_j)$ is a datum (SF) of Section~I.1.8, and for $j\ge1$ it is a datum (G2.S-3).
\item[(D2)] For $0\le j<k$ the sides $\ell_j$ and $\ell_j/L=\ell_{j+1}$ lie in $\mathrm{Cl}$, and the operators of level $j+1$ that occur in Proposition~20.3 are formed with the side $\ell_j/L$, the charge $L^{\frac{1}{2}}\cdot e_j$ and the mass $L\cdot\mu_j$.
They are therefore the operators with the argument $B'$ of Proposition~20.2 for the data $(k,\ell,e_k,\mathrm{zc}):=(j,\ell_j,e_j,\mu_j)$, with $\mathrm{zc}'=L\cdot\mu_j=\mu_{j+1}$, and $n_{j+1}$ is the number $n_{k+1}=\ell^3/L^3$ of Proposition~20.2.
\end{itemize}

\paragraph{3.2. The connections.}
One has $T_j(R)=\mathbb{L}_j/\ell_j\mathbb{Z}^3$.
Multiplication by $L^j$ maps $\mathbb{L}_j=L^{-j}\mathbb{Z}^3$ onto $\mathbb{Z}^3$, bonds onto bonds, and $\ell_j\mathbb{Z}^3$ onto $n\mathbb{Z}^3$ by (20.27).
Hence $a^{[j]}(b)=L^{\frac{j}{2}}\cdot a\left(L^j b\right)$ (Section~I.6.6) is a complex bond function on $T_j(R)$, and $a^{[0]}=a$.
For a bond $b$ of $T_j(R)$ with $j<k$, Section~I.3.7 gives
\begin{equation*}
\left(a^{[j+1]}\right)_L(b)=L^{-\frac{1}{2}}\cdot a^{[j+1]}(b/L)=L^{-\frac{1}{2}}\cdot L^{\frac{j+1}{2}}\cdot a\left(L^{j+1}b/L\right)=L^{\frac{j}{2}}\cdot a\left(L^j b\right)=a^{[j]}(b).
\tag{20.28}
\end{equation*}
Let $0\le j\le k$.
Every $\mathscr{C}\in\mathfrak{X}_j$ satisfies
\begin{equation*}
|\mathscr{C}|_\infty<e_j^{-1+2\varepsilon}=e_j^{\varepsilon}\cdot e_j^{-1+\varepsilon}\le e_j^{-1+\varepsilon},
\end{equation*}
because $e_j<1$ and $\varepsilon>0$; so $\mathscr{C}$ lies in $G_j$.
Every $A_1\in G_j$ lies in $\tilde{G}_j$: in the definition of $\tilde{G}_j$ (Section~I.5.3) one takes $A_0=0$, with $A'=0$ and $\boldsymbol{\lambda}=0$ on every $\square^{\sim5}$.
Thus $\mathfrak{X}_j$ is contained in $\tilde{G}_j$.
By the data $a^{[j]}$ lies in $\mathfrak{X}_j$, and so does the connection $0$.
Therefore
\begin{equation*}
a^{[j]}\in\tilde{G}_j\quad\text{and}\quad0\in\tilde{G}_j\qquad(0\le j\le k).
\tag{20.29}
\end{equation*}
The relations (20.29), (20.28) and the fourth relation of (20.27) make up the first sentence of (a).

\paragraph{3.3. Proof of (a), second sentence.}
\emph{The operators $S_j$ and $D_j$.}
Let $1\le j\le k$.
By (D1), (20.29) and Proposition~2.8\,(a) at $y=\infty$, the operator $F_\infty(A)=D_A+\mu_j+b_j\cdot J\cdot P_j(A)\cdot J^{*}$ is bijective on $\mathfrak{S}^{\mathrm{ch}}(\eta_j,\ell_j)$ for $A=a^{[j]}$ and for $A=0$.
So $S_j\left(a^{[j]}\right)$ and $S_j(0)$ exist, and $D_j\left(a^{[j]}\right)$, $D_j(0)$ are defined by the display of Section~I.3.5.
At $j=0$ the operator $D_0(A)=D_A+\bar{m}\cdot\eta$ is defined for every $A$ (Section~I.3.9).

\emph{The numbers $\delta Z_j$.}
Let $0\le j<k$.
By (D1) the data of Proposition~13 are at hand at level $j$: the quadruple $(j,\ell_j,\mu_j,e_j)$ and the family (p5) of the data of Proposition~20.3, with W1, W2, W7$^{\mathrm{ch}}$, W12, W13, W14 for $\mathrm{par}$.
The set $D_j$ is finite.
By Proposition~13\,(d0), for every $A\in\tilde{G}_j$,
\begin{equation*}
\begin{aligned}
&\delta Z_j(A)=\exp\left\{c^{\mathrm{d}}_j\cdot n_j+\sum_{X\in D_j}\acute{E}_j(X,A)\right\}\ne0,\\
&\frac{\delta Z_j(A)}{\delta Z_j(0)}=\exp\left\{\sum_{X\in D_j}E_j^{\mathrm{det}}(X,A)\right\}.
\end{aligned}
\tag{20.30}
\end{equation*}
By (20.29) this holds at $A=a^{[j]}$ and at $A=0$.

\emph{The determinants $\det D_j(0)$.}
The connection $0$ is real, and $\mu_j>0$.
For $j=0$, Proposition~19\,(a) at $k=0$ (side $n$, charge $q$, $\mu=\bar{m}\cdot\eta$, $A=0$) gives $\det D_0(0)=\det\left(D_0+\bar{m}\cdot\eta\right)\ne0$.
For $1\le j\le k$, Proposition~19\,(d) (side $\ell_j$, charge $e_j$, $\mu=\mu_j$, $A=0$) shows that $D_j(0)$ is invertible; so $\det D_j(0)\ne0$.
The operators of level $j$ contain no mass other than the running mass $\mu_j=\bar{m}\cdot a_j$ (at $j=0$ the mass $\bar{m}\cdot\eta$), and no lattice mass occurs (I.K10).
The only positivity of a mass that enters is $\mu_j>0$: in Proposition~19\,(a), (d) above, for $\det D_j(0)\ne0$, and as the lower end of the range $(0,1]$ of the mass in the data (SF) of (D1).
This is the last sentence of Proposition~20.3.
Besides its positivity, the mass enters through the bound $\mu_j\le1$, which (SF) requires and through which (D1) supplies the data of Proposition~2.8\,(a), of Proposition~13\,(d0) and of the displays of Section~I.5.7, and through the relation $\mu_{j+1}=L\cdot\mu_j$ of (20.27), with which (D2) and 3.4 apply Proposition~20.2.

\paragraph{3.4. Proof of (b).}
Let $0\le j<k$.
Proposition~20.2 is applied with the data of (D2), once with $B'=a^{[j+1]}$ and $B=B'_L=a^{[j]}$ (see (20.28)), and once with $B'=0$ and $B=0$.
The suppositions of (6.1) in Proposition~20.2 hold by 3.3: $S_j(B)$ exists if $j\ge1$, and $\delta Z_j(B)\ne0$.
Hence
\begin{equation*}
\begin{aligned}
&\det D_{j+1}\left(a^{[j+1]}\right)=b^{4n_{j+1}}\cdot\frac{\det D_j\left(a^{[j]}\right)}{\delta Z_j\left(a^{[j]}\right)},\\
&\det D_{j+1}(0)=b^{4n_{j+1}}\cdot\frac{\det D_j(0)}{\delta Z_j(0)}.
\end{aligned}
\tag{20.31}
\end{equation*}
By 3.3, since $j+1\le k$, both $\det D_{j+1}(0)$ and $\det D_j(0)$ are different from zero.
Dividing the first relation by the second and inserting the second relation of (20.30) at $A=a^{[j]}$, one finds
\begin{align*}
\frac{\det D_{j+1}\left(a^{[j+1]}\right)}{\det D_{j+1}(0)}
&=\frac{\det D_j\left(a^{[j]}\right)}{\det D_j(0)}\cdot\frac{\delta Z_j(0)}{\delta Z_j\left(a^{[j]}\right)}\\
&=\frac{\det D_j\left(a^{[j]}\right)}{\det D_j(0)}\cdot\exp\left\{-\sum_{X\in D_j}E_j^{\mathrm{det}}\left(X,a^{[j]}\right)\right\}.
\end{align*}

\emph{The step $j=0$.}
Here (6.1) is the case $k=0$ of Proposition~20.2, that is, (iii0) together with (iv).
The operators $D_0\left(a^{[0]}\right)=D_a+\bar{m}\cdot\eta$ and $D_0(0)=D_0+\bar{m}\cdot\eta$ act on $\mathfrak{S}^{\mathrm{ch}}(1,n)$, which has the dimension $192n^3=192n_0$.
The number $\delta Z_0(A)=\det\left(D_A+\bar{m}\cdot\eta+b_L\cdot J\cdot P(A)\cdot J^{*}\right)$ is a determinant on the same space (Section~I.3.9).
The operators $D_1\left(a^{[1]}\right)$ and $D_1(0)$ act on $\mathfrak{S}(1,\ell_1)$, which has the dimension $4n_1$, $n_1=n^3/L^3$; the exponent in (20.31) is $4n_1$.
In (20.30) at $j=0$ one has $c^{\mathrm{d}}_0=0$ (Section~I.3.4), and $E_0^{\mathrm{det}}$ is formed with the function $\acute{E}_0$ of Section~I.5.7, whose traces are traces on $\mathfrak{S}^{\mathrm{ch}}(1,n)$.
By (20.30),
\begin{equation*}
\exp\left\{\sum_{X\in D_0}E_0^{\mathrm{det}}(X,a)\right\}=\frac{\delta Z_0(a)}{\delta Z_0(0)}.
\end{equation*}
This is a quotient of two determinants, each of them a determinant on the whole space with all $192$ components per photon site, the light and the heavy components together: the $4$ components of the image of $\mathrm{Plt}=J\cdot J^{*}$, the light part (I.J4), and the $188$ components of the image of $\mathrm{Phv}=1-\mathrm{Plt}$, the heavy part.

\paragraph{3.5. Proof of (c).}
\emph{The first identity.}
By induction on $i$, (b) gives, for $0\le i\le k$,
\begin{equation*}
\frac{\det D_0\left(a^{[0]}\right)}{\det D_0(0)}=\exp\left\{\sum_{j<i}\sum_{X\in D_j}E_j^{\mathrm{det}}\left(X,a^{[j]}\right)\right\}\cdot\frac{\det D_i\left(a^{[i]}\right)}{\det D_i(0)}.
\tag{20.32}
\end{equation*}
For $i=0$ the two sides are equal, the sum being empty.
If (20.32) holds for some $i<k$, one inserts
\begin{equation*}
\frac{\det D_i\left(a^{[i]}\right)}{\det D_i(0)}=\exp\left\{\sum_{X\in D_i}E_i^{\mathrm{det}}\left(X,a^{[i]}\right)\right\}\cdot\frac{\det D_{i+1}\left(a^{[i+1]}\right)}{\det D_{i+1}(0)},
\end{equation*}
which is (b) at $j=i$, and obtains (20.32) for $i+1$.
At $i=k$, with $D_0\left(a^{[0]}\right)=D_a+\bar{m}\cdot\eta$ and $D_0(0)=D_0+\bar{m}\cdot\eta$, the identity (20.32) is the first identity of (c) with $\sum_{j<k}\sum_{X\in D_j}E_j^{\mathrm{det}}\left(X,a^{[j]}\right)$ in the exponent.

\emph{The second identity.}
Let $0\le j<k$.
By (D2) the sides $\ell_j$ and $\ell_j/L$ lie in $\mathrm{Cl}$.
By \cite[Theorem~13.2(t5)]{QEDI}, for every $X\in D_j$ there is exactly one $Y\in D_{j+1}$ with $\bar{X}^{L}=LY$.
The clause (t5) needs that $L$ is odd, through the following fact: two $M$-cubes with a common face lie in one $LM$-cube or in two $LM$-cubes with a common face.
Indeed, on a lift the centres of the two $M$-cubes are $z$ and $z+M e_\mu$ with $z\in M\mathbb{Z}^3$.
Since $L$ is odd, every coordinate of a point of $M\mathbb{Z}^3$ lies in exactly one open interval $LM\cdot\left(i-\frac{1}{2},i+\frac{1}{2}\right)$, $i\in\mathbb{Z}$, and the $LM$-cube that contains the $M$-cube with the centre $z$ has the centre $LM\cdot(i_1,i_2,i_3)$.
The two centres agree in the coordinates $\nu\ne\mu$ and differ by $M\le LM$ in the coordinate $\mu$.
So the two triples of indices agree at $\nu\ne\mu$, and at $\mu$ they are equal or consecutive.
By the existence and the uniqueness of $Y$, the sets $\{X\in D_j:\bar{X}^{L}=LY\}$, $Y\in D_{j+1}$, are pairwise disjoint, and their union is $D_j$.

The recursion of the data is read at $\mathscr{C}=a^{[j+1]}$; there $\mathscr{C}_L=a^{[j]}$ by (20.28).
The numbers $E_j^{\mathrm{det}}(X,A)$, $X\in D_j$, $A\in\tilde{G}_j$, are determined by the data of Proposition~20.3: they are formed by the displays of Section~I.5.7 under (SF), which holds by (D1), with the family (p5) of level $j$, and they are the numbers that appear in Proposition~13\,(d0).
At $A=a^{[j]}$ these displays involve $\acute{E}_j(X,\cdot)$ at the connections $a^{[j]}$, $-a^{[j]}$, $\overline{a^{[j]}}$, $-\overline{a^{[j]}}$ and $0$.
The first four have the same two norms $|\cdot|_\infty$ and $|\partial\,\cdot|_\infty$ as $a^{[j]}$, because a modulus is invariant under $A\mapsto-A$ and under $A\mapsto\overline{A}$, and $\partial(-A)=-\partial A$, $\partial\overline{A}=\overline{\partial A}$; so they lie in $\mathfrak{X}_j$, hence in $\tilde{G}_j$ by 3.2.
The connection $0$ lies in $\tilde{G}_j$ by (20.29).
At $\mathscr{C}=a^{[j+1]}$ the recursion reads
\begin{equation*}
\mathfrak{P}_{j+1}\left(Y,a^{[j+1]}\right)=\sum_{X:\,\bar{X}^{L}=LY}\left(\mathfrak{P}_j\left(X,a^{[j]}\right)+E_j^{\mathrm{det}}\left(X,a^{[j]}\right)\right)\qquad(Y\in D_{j+1}).
\end{equation*}
Starting from $\mathfrak{P}_0=0$, it defines the numbers $\mathfrak{P}_j\left(X,a^{[j]}\right)$, $X\in D_j$, for $j=0,1,\dots,k$ in turn.
Summation over $Y\in D_{j+1}$ gives
\begin{equation*}
\sum_{Y\in D_{j+1}}\mathfrak{P}_{j+1}\left(Y,a^{[j+1]}\right)=\sum_{X\in D_j}\mathfrak{P}_j\left(X,a^{[j]}\right)+\sum_{X\in D_j}E_j^{\mathrm{det}}\left(X,a^{[j]}\right)\qquad(0\le j<k).
\end{equation*}
Since $\sum_{X\in D_0}\mathfrak{P}_0\left(X,a^{[0]}\right)=0$, the sum of these relations over $j=0,\dots,k-1$ is
\begin{equation*}
\sum_{X\in D_k}\mathfrak{P}_k\left(X,a^{[k]}\right)=\sum_{j<k}\sum_{X\in D_j}E_j^{\mathrm{det}}\left(X,a^{[j]}\right).
\end{equation*}
This is the second identity; inserted into (20.32) at $i=k$, it gives the first identity in the form in which it is stated.
For $k=0$ both identities hold with empty sums, since $\mathfrak{P}_0=0$ and $D_0\left(a^{[0]}\right)=D_a+\bar{m}\cdot\eta$.

\emph{The free ratio.}
For $a=(t/e)\,\hat{\mathfrak{a}}$ the left side of (c) is
\begin{equation*}
\frac{\det\left(D_{(t/e)\hat{\mathfrak{a}}}+\bar{m}\cdot\eta\right)}{\det\left(D_0+\bar{m}\cdot\eta\right)}
\end{equation*}
on the torus with $n=R/\eta$ sites per side; this is the display that defines $N_0(t)$ in Section~I.7.6, as stated in the sentence after (c).

\paragraph{3.6. Uniformity.}
Apart from the membership $a^{[j]}\in\tilde{G}_j$ in (a), the assertions (a)--(c) are identities, assertions that inverses exist and assertions that determinants do not vanish; none of them is an estimate with a constant on its right side.
That membership is a bound with the threshold $e_j^{-1+\varepsilon}$ on the cubes $\square^{\sim5}$: by Section~I.5.3 it means that $a^{[j]}=A_0+A_1$ for some $A_1$ with $|A_1|_\infty<e_j^{-1+\varepsilon}$ and some real $A_0$ that, on every cube $\square^{\sim5}$, differs by the coboundary of a real function from a real field $A'$ with $|A'|_\infty<e_j^{-1+\varepsilon}$ on that cube.
The constants of the lines (C0)--(C9) of Section~I.8.2 that enter (a)--(c) are the following, and no others.
The constants $b$ and $L$ enter through $b_j$, $b_L$, $\eta_j$, the blocks and the scalings.
The constant $\rho$ and the weights $w_i$ enter through $D_A$ alone.
The constant $M$ enters through $\tilde{G}_j$, the sets $D_j$ of polymers, the closures $\bar{X}^{L}$, the range $|c-c'|\le M$ of $U_A$ in (I.1.6) and the class $\mathrm{Cl}$.
The constant $\varepsilon$ enters through $\tilde{G}_j$ and, in the data, through $\mathfrak{X}_j$.
The constants $e_1$, $r$ and $\mathrm{par}$ stand only in the data: in (G2.S-4), in $k_*(R)$ and in the properties of the families (p5).
The constants $L$ and $M$ stand in the data as well: in $\mathrm{Tor}(e)$ of (G2.S-4) through $R=h\cdot L^{T}\le R_{\mathrm{top}}(e)=h\cdot M\cdot r_K(e)$, in $N_{\min}(R,e)$ of (G2.S-4) through the condition $R\cdot L^{N}\ge L\cdot M\cdot r_0$, and in $k_*(R)$ through $\mathfrak{n}_k=\log_L\left(\ell_k/(M\cdot r_k)\right)$ (Sections~I.6.1 and~I.6.2).
By (I.U4)$^{\mathrm{ch}}$, none of these constants depends on $N$, $R$, $j$, $k$, $\bar{m}$, $e$ or $a$.
In the assertions (a)--(c) and in their data, where $a$ is an arbitrary complex one-cochain with $a^{[j]}\in\mathfrak{X}_j$ for $j\le k$, the only constants of the line (C10) of Section~I.8.2 that enter are $\bar{m}$, through $\mu_j$, and $h$, through (G2.S-4).
The probe $\beta$, which is also a constant of the line (C10), enters only the sentence after (c), which concerns $N_0(t)$, and its proof at the end of 3.5; it enters there through the special choice $a=(t/e)\,\hat{\mathfrak{a}}$ of the datum $a$, since $\hat{\mathfrak{a}}=\eta^{\frac{1}{2}}a_{\eta}$ with $a_{\eta}=d_\eta^{*}u$, and $u_{\mu\nu}$ is formed with $J_{\eta}\beta_{\mu\nu}$ (Section~I.7.6).
The number $\bar{m}$ is a datum of (G2.S-4); it enters (a)--(c) through $\mu_j=\bar{m}\cdot a_j$, and at $j=0$ through $\bar{m}\cdot\eta$.
The number $h=L^{\left\lfloor\log_L\frac{1}{\bar{m}}\right\rfloor}$ of Section~I.0.5 is a function of $(L,\bar{m})$, and it stands only in the data (G2.S-4).
Neither $\bar{m}$ nor $h$ depends on $N$, $R$, $j$, $k$, $e$ or $a$.
The assertions (a)--(c) hold for all data (G2.S-4), every $k\le k_*(R)$, all families (p5) with the properties named and every complex $a$ with $a^{[j]}\in\mathfrak{X}_j$ for $j\le k$.
The conditions on the constants needed in this proof are the definition of $e_1$, the inequalities $e_1\le e_0<1$ and $\varepsilon>0$, in 3.1 and 3.2, and the oddness of $L$, in the argument on $LM$-cubes in 3.5.
\qed

\subsection*{4. Conditions on the constants}

The three proofs use the following conditions on the constants directly; the rows follow the order of the lines of Section~I.8.2.

\noindent
\begin{tabular}{p{0.29\textwidth}p{0.36\textwidth}p{0.25\textwidth}}
\hline
condition & where it stands & used in\\
\hline
$b=1$; it enters through its consequences $b>0$, $b_k\le b_{+}=b$ and $4b_{+}\ge1$ & Section~I.8.2, line (C0); Section~I.3.4; Proposition~14.2\,(a), rows $b$, $b_{+}$ & Proposition~20.1: 1.1, (20.4); Proposition~20.2: (n1)\\
\hline
$L$ an odd integer, $L>1$ & Section~I.8.2, line (C3); Proposition~14.2\,(b), (L5) & Proposition~20.2: (n1), (n3), (s5); Proposition~20.3: 3.5 ($LM$-cubes)\\
\hline
$M\ge L^2$; it enters through its consequence $(L-1)/2\le M$ & Section~I.0.5; Section~I.8.2, line (C5); Proposition~14.2\,(b), (M1) & Proposition~20.2: (n3), (s6)\\
\hline
$\gamma^{F}>0$, which enters through its consequence $c(\gamma^{F})<\infty$; $C^{F}\ge0$ & Section~I.8.2, lines (C4r), (C6); the paragraph ``Constants'' at the head of Proposition~14.1 ($0<\gamma^{F}\le1/(6L)$); Proposition~14.2\,(f) ($4\le C^{F}$); both are presupposed by the data of Proposition~20.1 & Proposition~20.1: 1.1, (20.4)\\
\hline
$C_{0}\ge C^{F}$; $C_{0}\ge4b_{+}+16\,C^{F}\cdot c(\gamma^{F})$ & Section~I.8.2, line (C6); Proposition~14.2\,(a), row $C_{0}$ ($C_{0}$ is the maximum of three numbers, these two among them) & Proposition~20.1: (20.2), (20.4), (20.6)\\
\hline
$\varepsilon>0$ & Section~I.0.5; Section~I.8.2, line (C7); Proposition~14.2\,(b), (S1) ($\varepsilon>1-\alpha>0$) & Proposition~20.1: 1.1, (20.1), (20.2); Proposition~20.3: 3.1, 3.2\\
\hline
$\exp(-\varepsilon v)\le1/4$ for all $v\ge-\log e_0$; together with $\varepsilon>0$ it enters through its consequence $e_0<1$ & Section~I.8.2, line (C8); (A8), first inequality (Proposition~14.2\,(b)) & Proposition~20.1: (20.1); Proposition~20.3: 3.1, 3.2\\
\hline
$e_1\le e_0$ & Section~I.8.2, line (C9); Proposition~14.3\,(f) & Proposition~20.3: 3.1\\
\hline
$\ell\in\mathrm{Cl}$; it enters through its consequence that $\ell/(LM)$ is an odd integer $\ge13$, which gives in (s0) the assertion of Proposition~20.2 that $\ell/L$ is an odd multiple of $M$ and $\ell/L\ge13M$; the displays of (iv) need only its consequences that $\ell/L$ is a positive integer and that $L-1<\ell$ & class of the sides (a datum) & Proposition~20.2: (s0), (s6)\\
\hline
\end{tabular}

Each row bounds a constant only by constants chosen before it in Section~I.8.2: $M$ is bounded from below by $L$; $C_{0}$ from below by $C^{F}$, $b_{+}$, $\gamma^{F}$; $\varepsilon$ from below by $1-\alpha$; $e_0$ from above through $\varepsilon$; $e_1$ from above by $e_0$.

Further conditions enter the three proofs through the statements of \cite{QEDI} cited in them and through (G4.Sc6); these need the following conditions on the constants:

\noindent
\begin{tabular}{p{0.3\textwidth}p{0.37\textwidth}p{0.23\textwidth}}
\hline
statement & conditions it needs & where they stand\\
\hline
\cite[Lemma~17.4(a)]{QEDI} & $\exp(-v/4)+\exp(-\varepsilon v)\le1$; $L\ge5$; $e_k<1$ & (A8), fourth inequality; (L5); (A8), first inequality (also a consequence of the fourth)\\
\hline
\cite[Proposition~19.1(e), (f)]{QEDI}; (G4.Sc6) (the assertions applied in 3.1) & (e1): $v\ge8$, $v^r>31L+16$, $(1+3\,\mathrm{lg}/v)^r<L$ for $v\ge-\log e_1$; $L\ge33$, $L$ odd; $M$ a power of $L$ & Proposition~14.2\,(b), (e1); (L1), (L5); (M1)\\
\hline
\cite[Theorem~13.2(t5)]{QEDI} & $L$ odd; $\ell$, $\ell/L\in\mathrm{Cl}$ & (L5); data\\
\hline
(G2.r6), that is \cite[Section~13.1.6]{QEDI} & none & --\\
\hline
\end{tabular}

Propositions~2.8, 5, 6, 13, 15, 16.1 and~19 hold under the conditions on the constants that are stated with them.
Propositions~20.1--20.3 hold under the conditions on the constants of Section~I.8 and of Proposition~14.2.

\section*{Chapter 48a. The statements without the second slot (spaces; low kernels and extraction; adjustment; global fields and locality): proofs of Propositions 48.1--48.4}

\subsection*{0. Conventions of this chapter; four lemmas}

The notation is that of Part~I and of Sections~G2.0.4--G2.0.6.
For every integer $k\ge0$, every side $\ell\in\mathrm{Cl}$ and every $e_k\in(0,e_0]$, the clauses (t2), (t7) and (t10) of \cite[Theorem~13.2]{QEDI}, as well as \cite[Lemma~17.3]{QEDI} and \cite[Lemma~5.19]{QEDI}, hold in the form in which they are stated there.
The domains of \cite[Definition~5.16]{QEDI} are formed with the same data, and a phrase such as ``the kernels depend on $\mathscr{A}(b)$ only for the bonds $b$ that meet $X$'' has the meaning fixed in \cite[Definition~5.5\,(1)]{QEDI}.

In this chapter $\tilde{Z}$ is a lift of a set $Z$ of sites, and $\tilde{c}$, $\tilde{c}'$ are lifts of sites $c$, $c'$, in the sense of Section~I.1.1.
The symbol $\tilde{X}'$ denotes the union of $M$-cubes of $\mathbb{R}^3$ furnished by \cite[Theorem~13.2\,(t10)]{QEDI}, except in the proof of (f1) in Section~2\,(f), where it is a second element of $\mathrm{Sm}(\tilde{\square})$.
The objects $\tilde{X}$, $\tilde{U}$, $\tilde{B}$, $\mathrm{Grc}_k(\tilde{U})$, $\mathrm{cov}_\ell$, $T_v$, $\Phi$ and $\mathrm{rg}$ are those of Section~I.1.9, and $\tilde{\square}$, $\mathrm{Sm}(\tilde{\square})$ and $\mathrm{pr}$ are those of Proposition~16.3\,(f).
As in Section~G2.0.6, $\mathrm{Mx}$ denotes a $4\times4$ matrix.
The letter $\mathrm{rad}$ stands for a radius: in Section~1\,(d) for the radius of a polydisk, and in Section~4 for the number $2L\cdot M$ of Proposition~18.2.
In Sections~0--3 the letter $z$ is the centre of an $M$-cube (Sections~I.4.2 and~I.5.6), while in Section~4 it is a point of $\mathbb{C}^2$.
For a fermion site $x$, $x^{(3)}$ is the neighbour of type~3 of $x$, and $c_3=\pi\left(x^{(3)}\right)$ is the photon site of this neighbour (I.1.2).
In Section~4\,(b), on the other hand, the lower index of $x_i$, $x_j$, $x'_j$ and $c_j$ numbers the arguments: $x_i$ is the site of $\xi_i$, $x_j$ and $x'_j$ are the sites of $\boldsymbol{\zeta}_j$, and $c_j=\pi(x_j)$;
there the neighbour of type~3 of $x_j$ is written $x_j^{(3)}$, and the photon site of this neighbour is written $\pi\left(x_j^{(3)}\right)$.
In Section~4\,(b), Step~2, a set $\mathrm{Dz}\subset\mathbb{C}^2$, for $z\in\mathrm{Dz}$ a bond function $A(z)$ and a function $\Lambda_z$ on the sites of $T_k$, and a function $\mathrm{Kz}$ on $\mathrm{Dz}$ are defined.
They correspond to objects of Step~7 of the proof of \cite[Theorem~13.11]{QEDI}.
In that step a set $\mathscr{D}\subset\mathbb{C}^4$, and for $z\in\mathscr{D}$ bond functions $A(z)$, $C(z)$ and a function $\Lambda_z$, are defined;
a function $\mathscr{K}$ on $\mathscr{D}$ is formed first for the families with the second slot, and then once more in the last part of the step, which proves the assertions of \cite[Theorem~13.11\,(iii)]{QEDI} on the families without the second slot.
Compared with the first part of the step, the objects of Section~4\,(b) differ as follows:
the coordinates $z_3$, $z_4$ and the bond function $C(z)$ do not occur, $\Lambda_z$ contains no terms with $z_3$, $z_4$, and $\mathrm{Kz}(z)$ is formed with $E_{nm}(X;A(z);\xi,\boldsymbol{\zeta})$ in the place of $E_{nm}\left(X;A(z),C(z);\xi,\boldsymbol{\zeta}\right)$.
The last part of the step uses the bond functions $A(z)$ of the first part, functions $\Lambda_z$ that contain no terms with $z_3$, $z_4$, and the function $\mathscr{K}(z)$ formed with $E_{nm}(X;A(z);\xi,\boldsymbol{\zeta})$;
in these points $A(z)$, $\Lambda_z$ and $\mathrm{Kz}$ agree with the objects of the last part.
The two systems of objects differ in the following points.
In the last part, $z$ ranges over $\mathscr{D}$, whose conditions $|\operatorname{Im}z_1|<2$, $|\operatorname{Im}z_2|<2$ on the first two coordinates are those that define $\mathrm{Dz}$;
the coordinates $z_3$, $z_4$ enter neither $A(z)$ nor $\Lambda_z$, the bond function $C(z)$ vanishes for every $z\in\mathscr{D}$, and the identity $\mathscr{K}(z)=\mathscr{K}(0)$ is evaluated at $z=(1,i,1,i)$.
In Section~4\,(b) the set $\mathrm{Dz}$ lies in $\mathbb{C}^2$, no bond function $C(z)$ occurs, and the identity $\mathrm{Kz}(z)=\mathrm{Kz}(0)$ is evaluated at $z=(1,i)$.
In the last part the factor of each argument is formed with the value of $\Lambda_z$ at the site of the argument, or at the first site of its pair if the argument is one of $\chi$;
in Section~4\,(b) the arguments carry fermion sites, and the factor of the $j$-th argument is formed with the value of $\Lambda_z$ at the photon site $c_j=\pi(x_j)$.
Finally, the bond functions $\mathrm{bh}$, $\mathrm{kc}_1$, the site $\tilde{c}$ and the number $\mathrm{rad}=2L\cdot M$ of Section~4 take the places of the bond functions $\hat{b}$, $c_1$, the site $x_0$ and the number $\rho=2L\cdot M$ of \cite[Theorem~13.11]{QEDI}.

The labels (cl1) and (cl2) refer to two of the classical facts recalled at the beginning of \cite[Section~13]{QEDI}:
(cl1) is Cauchy's integral formula in one variable \cite[Chapter~4, Section~2.2]{Ahlfors1979},
and (cl2) is the identity theorem for analytic functions of several complex variables \cite[Chapter~I, Section~A]{GunningRossi1965}.

\subsubsection*{0.1. Setting, integrals, lifts}

The settings (S-1) and (S-2) are those of Section~G2.0.4.
Let $\mathrm{pr}\colon\mathbb{L}_k\to T_k$ denote the covering map; on fermion sites it acts by $\mathrm{pr}(c,g):=(\mathrm{pr}\,c,g)$.
In the places where the side has to be named, namely in Sections~2\,(f) and~3\,(f), $T_k(\ell)$ is the torus $T_k$ formed with the side $\ell$, and $\mathrm{pr}_\ell\colon\mathbb{L}_k\to T_k(\ell)$ is its covering map;
for two sides $\ell_1$, $\ell_2$ one puts $\mathrm{pr}_i:=\mathrm{pr}_{\ell_i}$.
Throughout, $\mathsf{w}=(\mathsf{w}_1,\mathsf{w}_2)$ is an arbitrary pair of positive numbers.

According to (I.AP) in Section~I.3.6, the generators and the kernels are indexed by the points of the covering lattices, with
\[
\psi(x+\ell\cdot e_\mu,\,\cdot\,)=-\psi(x,\,\cdot\,),\qquad\chi(x+\ell\cdot e_\mu,x'+\ell\cdot e_\mu,\,\cdot\,)=-\chi(x,x',\,\cdot\,),
\]
and every kernel is antiperiodic in each of its arguments.
It follows that a product of a kernel with generators is $\ell\mathbb{Z}^3$-periodic in each argument, so that
\[
\int\mathrm{d}\xi=\sum_{\xi}\nu(\xi),\qquad\int\mathrm{d}\boldsymbol{\zeta}=\sum_{\boldsymbol{\zeta}}\nu(\boldsymbol{\zeta}),\qquad\text{with}\quad\nu(\xi)=\frac{\eta_k^3}{48},\quad\nu(\boldsymbol{\zeta})=\left(\frac{\eta_k^3}{48}\right)^2,
\]
may be formed as sums over arbitrary sets of representatives.

Let $\tilde{Z}$ be a finite set of sites of $\mathbb{L}_k$, and let $f$ be a function on $\tilde{Z}\times O_{\mathrm{cub}}$ with values in a vector space.
Put
\begin{equation}\tag{0.a}
\int_{\tilde{Z}}\mathrm{d}x\,f(x):=\frac{\eta_k^3}{48}\cdot\sum_{c\in\tilde{Z}}\ \sum_{g\in O_{\mathrm{cub}}}f(c,g).
\end{equation}
For a union $Z$ of $M$-cubes of $T_k$ and an $\ell\mathbb{Z}^3$-periodic function $f$, the integral $\int_Z\mathrm{d}x\,f(x)$ of Section~I.1.6 equals the expression (0.a) formed with an arbitrary set $\tilde{Z}$ of representatives of the sites of $Z$.
Each $M$-cube contains $(M/\eta_k)^3$ photon sites, and over each photon site lie $48$ fermion sites of the weight $\eta_k^3/48$.
Consequently
\[
\int_Z\mathrm{d}x=M^3\cdot|Z|_M=\operatorname{Vol}(Z).
\]

Small and large polymers are defined as in (G2.r3) of Section~G2.0.5: a polymer $X\in D_k$ is small if $d_M(X)\le L$, and large otherwise.
Every small polymer $X$ has a lift $\tilde{X}$, unique up to $\ell\mathbb{Z}^3$, whose sup-diameter is at most $M\cdot(L+2)$; this is Lemma~D\,(Dl2) below.
A fermion site $x$ is said to lie over $\tilde{X}$ if $\pi(x)\in\tilde{X}$.
When a kernel is integrated without generators, all its arguments are chosen over $\tilde{X}$ (Section~I.3.6, (I.AP)).

Unless the slot is mentioned, the word family means a family without slot in the sense of Proposition~48.1.
The zero bond function belongs to $R_k$, and $R_k\subset\tilde{R}_k$: for $A\in R_k$ one takes $A_1=A$, $A_0=0$, $A'=0$ and $\lambda=0$ in \cite[Definition~5.16\,(b)]{QEDI}.

\subsubsection*{0.2. Lemma A (elements and kernels)}

\begin{named}{Lemma A}
Let $F\in\mathrm{Gr}[\psi,\chi]$.
\begin{itemize}
\item[(Ak1)] There is exactly one system of kernels $F_{nm}$, antisymmetric in the arguments $\xi$ and, separately, in the arguments $\boldsymbol{\zeta}$, such that
\begin{multline*}
F=\sum_{n,m}\frac{1}{n!\,m!}\int\mathrm{d}\xi\,\mathrm{d}\boldsymbol{\zeta}\,F_{nm}(\xi_1,\dots,\xi_n;\boldsymbol{\zeta}_1,\dots,\boldsymbol{\zeta}_m)\\
\cdot\psi(\xi_1)\cdots\psi(\xi_n)\cdot\chi(\boldsymbol{\zeta}_1)\cdots\chi(\boldsymbol{\zeta}_m).
\end{multline*}
At increasing arguments these kernels satisfy $F_{nm}(I_1;I_2)\cdot\nu(I_1)\cdot\nu(I_2)=c_F(I_1,I_2)$; here $c_F(I_1,I_2)$ is the coefficient of $F$ at the basis monomial $\psi(I_1)\cdot\chi(I_2)$ (Section~I.5.1), and $\nu(I):=\prod_{t\in I}\nu(t)$.
The term of index $(n,m)$ is denoted by $F^{[n,m]}$, and
\[
\|F\|_{\mathsf{w}}=\sum_{n,m}\left\|F^{[n,m]}\right\|_{\mathsf{w}},\qquad
\left\|F^{[n,m]}\right\|_{\mathsf{w}}=\frac{\mathsf{w}_1^n\cdot\mathsf{w}_2^m}{n!\,m!}\cdot\|F_{nm}\|,\qquad
\|F_{nm}\|=\int\mathrm{d}\xi\,\mathrm{d}\boldsymbol{\zeta}\,|F_{nm}|.
\]
In particular $|F_{00}|\le\|F\|_{\mathsf{w}}$.
A map of an open subset of a finite-dimensional complex vector space into $\mathrm{Gr}[\psi,\chi]$ is analytic if and only if all its kernels are analytic functions.
\item[(Ak2)] Let $\mathrm{phf}$ be a complex function of the indices $\xi$ and $\boldsymbol{\zeta}$ that is invariant under the translations by $\ell\mathbb{Z}^3$, and let $\mathrm{Gph}$ be the homomorphism of $\mathrm{Gr}[\psi,\chi]$ with $\psi(\xi)\mapsto\mathrm{phf}(\xi)\cdot\psi(\xi)$ and $\chi(\boldsymbol{\zeta})\mapsto\mathrm{phf}(\boldsymbol{\zeta})\cdot\chi(\boldsymbol{\zeta})$.
Then
\[
(\mathrm{Gph}\,F)_{nm}(\xi;\boldsymbol{\zeta})=F_{nm}(\xi;\boldsymbol{\zeta})\cdot\prod_i\mathrm{phf}(\xi_i)\cdot\prod_j\mathrm{phf}(\boldsymbol{\zeta}_j).
\]
For a real function $\Lambda$ of $c$, the automorphism $G_{\Lambda}$ of Section~I.4.5 is the map $\mathrm{Gph}$ with
\begin{align*}
\mathrm{phf}(\xi)&=\exp\left((-1)^{\omega}\,i\,e_k\cdot\Lambda(c)\right)&&\text{for }\xi=(x,\beta,\omega),\\
\mathrm{phf}(\boldsymbol{\zeta})&=\exp\left((-1)^{\omega}\,i\,e_k\cdot\Lambda(c)\right)&&\text{for }\boldsymbol{\zeta}=(x,x',\beta,\omega),
\end{align*}
where $c=\pi(x)$ is the photon site of the site $x$, respectively of the first site $x$ of the pair.
\item[(Ak3)] Let $F$ be balanced.
Then $F_{20}(\xi_1,\xi_2)=0$ whenever $\omega_1=\omega_2$.
Let $F_{20}(x,x')$ be the $4\times4$ matrix with the entries $F_{20}(x,x')_{\beta\beta'}:=F_{20}\left((x,\beta,1),(x',\beta',0)\right)$ (Section~I.5.1).
Then
\begin{gather*}
F^{[2,0]}=\int\mathrm{d}x\int\mathrm{d}x'\,\bar{\psi}(x)^T\cdot F_{20}(x,x')\cdot\psi(x'),\\
\frac{\mathsf{w}_1^2}{2}\cdot\|F_{20}\|=\left\|F^{[2,0]}\right\|_{\mathsf{w}}=\mathsf{w}_1^2\cdot\sum_{\beta,\beta'}\int\mathrm{d}x\int\mathrm{d}x'\,\left|F_{20}(x,x')_{\beta\beta'}\right|;
\end{gather*}
the two relations of this display together are the display (G2.0.1) of Section~G2.0.5, with $F$ in the place of $E(X)$.
\item[(Ak4)] Let $K(x,x')$ be a function of $x,x'\in\mathbb{L}^{\mathrm{f}}_k$ with values in the $4\times4$ matrices that is antiperiodic in $x$ and in $x'$.
Then the element
\[
\mathrm{Bq}[K]:=\int\mathrm{d}x\int\mathrm{d}x'\,\bar{\psi}(x)^T\cdot K(x,x')\cdot\psi(x')
\]
is balanced and of bidegree $(2,0)$, and $\mathrm{Bq}[K]_{20}(x,x')=K(x,x')$.
\end{itemize}
\end{named}

\begin{proof}[Proof of Lemma A]
(Ak1)
Proposition~15\,(a) is applied with $p=2$.
The index set $\mathrm{Ind}_1$ consists of the $\xi=(x,\beta,\omega)$ with $x$ in a set of representatives of $T^{\mathrm{f}}_k$, with the weight $\nu_1=\eta_k^3/48$;
the index set $\mathrm{Ind}_2$ consists of the $\boldsymbol{\zeta}=(x,x',\beta,\omega)$ with $(x,x')$ in a set of representatives of $\mathrm{Pz}_k$, with the weight $\nu_2=\left(\eta_k^3/48\right)^2$;
the orders on both sets are arbitrary.
To obtain kernels, define $F_{nm}(I_1;I_2):=c_F(I_1,I_2)/\left(\nu(I_1)\cdot\nu(I_2)\right)$ at increasing arguments, extend this function antisymmetrically to all orderings of pairwise different arguments, by $0$ to the arguments two of which coincide, and antiperiodically to the covering lattices.
A pair $(I_1,I_2)$ with $\#|I_1|=n$ and $\#|I_2|=m$ then appears $n!\,m!$ times in the integral of (Ak1), each time with the term $c_F(I_1,I_2)\cdot\psi(I_1)\cdot\chi(I_2)$; this proves the existence.
The uniqueness, and the formula for the norm, which holds with equality since the kernels are antisymmetric, are given by Proposition~15\,(a).
The norm (I.5.1) is a sum over the basis monomials; hence it is additive over the bidegrees.
Finally, the numbers $c_F(I_1,I_2)$ are the coordinates of $F$ in a basis, and they equal the values of the kernels at increasing arguments multiplied by positive numbers.
This gives the last assertion of (Ak1).

(Ak2)
Apply $\mathrm{Gph}$ to the representation of (Ak1).
The function $\prod_i\mathrm{phf}(\xi_i)\cdot\prod_j\mathrm{phf}(\boldsymbol{\zeta}_j)$ is invariant under the permutations of the $\xi$ and under those of the $\boldsymbol{\zeta}$.
The new kernels are therefore antisymmetric, and, $\mathrm{phf}$ being invariant under $\ell\mathbb{Z}^3$, they are antiperiodic in each argument.
By the uniqueness in (Ak1) they are the kernels of $\mathrm{Gph}\,F$.
The assertion on $G_{\Lambda}$ is (I.T-G) of Section~I.4.5.

(Ak3)
Let $\xi_1<\xi_2$ with $\omega_1=\omega_2$.
The monomial $\psi(\xi_1)\cdot\psi(\xi_2)$ is not balanced, so that $c_F\left(\{\xi_1,\xi_2\},\emptyset\right)=0$ and $F_{20}(\xi_1,\xi_2)=0$.
By antisymmetry the same holds for $\xi_1>\xi_2$, and $F_{20}(\xi,\xi)=0$.
Hence
\[
F^{[2,0]}=\frac{1}{2}\int\mathrm{d}\xi_1\int\mathrm{d}\xi_2\,F_{20}(\xi_1,\xi_2)\cdot\psi(\xi_1)\cdot\psi(\xi_2)
\]
is the sum of the two parts in which $(\omega_1,\omega_2)=(1,0)$ and $(\omega_1,\omega_2)=(0,1)$.
In the second part the names $\xi_1$, $\xi_2$ are exchanged; since
\[
F_{20}(\xi_2,\xi_1)\cdot\psi(\xi_2)\cdot\psi(\xi_1)=F_{20}(\xi_1,\xi_2)\cdot\psi(\xi_1)\cdot\psi(\xi_2),
\]
both factors having changed sign, the two parts coincide, and
\[
F^{[2,0]}=\sum_{\beta,\beta'}\int\mathrm{d}x\int\mathrm{d}x'\,F_{20}\left((x,\beta,1),(x',\beta',0)\right)\cdot\bar{\psi}_{\beta}(x)\cdot\psi_{\beta'}(x').
\]
In the same way $\|F_{20}\|=\int\mathrm{d}\xi_1\int\mathrm{d}\xi_2\,\left|F_{20}(\xi_1,\xi_2)\right|$ is the sum of its parts with $(\omega_1,\omega_2)=(1,0)$ and $(\omega_1,\omega_2)=(0,1)$.
These two parts are equal, because $\left|F_{20}(\xi_2,\xi_1)\right|=\left|F_{20}(\xi_1,\xi_2)\right|$, and the first of them is $\sum_{\beta,\beta'}\int\mathrm{d}x\int\mathrm{d}x'\,\left|F_{20}(x,x')_{\beta\beta'}\right|$.
Thus
\[
\|F_{20}\|=2\sum_{\beta,\beta'}\int\mathrm{d}x\int\mathrm{d}x'\,\left|F_{20}(x,x')_{\beta\beta'}\right|,
\]
and together with (Ak1) for $(n,m)=(2,0)$ this is the second display of (Ak3).

(Ak4)
Define
\[
\tilde{K}\left((x,\beta,1),(x',\beta',0)\right):=K(x,x')_{\beta\beta'},\qquad\tilde{K}\left((x',\beta',0),(x,\beta,1)\right):=-K(x,x')_{\beta\beta'},
\]
and $\tilde{K}(\xi_1,\xi_2):=0$ if $\omega_1=\omega_2$.
The function $\tilde{K}$ is antisymmetric and antiperiodic in each argument, and the computation made in the proof of (Ak3) shows
\[
\frac{1}{2}\int\mathrm{d}\xi_1\int\mathrm{d}\xi_2\,\tilde{K}(\xi_1,\xi_2)\cdot\psi(\xi_1)\cdot\psi(\xi_2)=\mathrm{Bq}[K].
\]
In particular $\mathrm{Bq}[K]$ is a linear combination of the monomials $\bar{\psi}_{\beta}(x)\cdot\psi_{\beta'}(x')$, hence balanced and of bidegree $(2,0)$.
By the uniqueness in (Ak1), $\tilde{K}$ is the kernel of $\mathrm{Bq}[K]$.
\end{proof}

\subsubsection*{0.3. Lemma B (the four maps on bilinear elements)}

\begin{named}{Lemma B}
Let $K$ be as in (Ak4), let $\tilde{g}=(S,a)\in\mathbb{G}_k$ with $a\in M\mathbb{Z}^3$, and let $\Lambda$ be a real function of $c$, that is, a function on the sites of $T_k$, regarded as an $\ell\mathbb{Z}^3$-periodic function on $\mathbb{L}_k$.
Then:
\begin{itemize}
\item[(Bm1)] $a_{\tilde{g}}\,\mathrm{Bq}[K]=\mathrm{Bq}[K']$ with $K'(x,x')=S^{-1}\cdot K(\tilde{g}x,\tilde{g}x')\cdot S$;
\item[(Bm2)] $G_{\Lambda}\,\mathrm{Bq}[K]=\mathrm{Bq}[K']$ with $K'(x,x')=\exp\left(i\,e_k\cdot\left(\Lambda(\pi x')-\Lambda(\pi x)\right)\right)\cdot K(x,x')$;
\item[(Bm3)] $\mathrm{Cg}\,\mathrm{Bq}[K]=\mathrm{Bq}[K']$ with $K'(x,x')=\mathsf{C}^{-1}\cdot K(x',x)^T\cdot\mathsf{C}$;
\item[(Bm4)] $\mathbf{K}\,\mathrm{Bq}[K]=\mathrm{Bq}[K']$ with $K'(x,x')=\varkappa^{-1}\cdot\overline{K(x,x')}\cdot\varkappa$.
\end{itemize}
In each of the four cases $K'$ is antiperiodic in $x$ and in $x'$.
\end{named}

\begin{proof}[Proof of Lemma B]
The four maps are those of (I.T-g), (I.T-G), (I.T-C) and (I.T-K) in Section~I.4.5.
Recall from Proposition~16.1\,(a2) that $\tilde{g}(x+\ell\cdot n)=\tilde{g}x+\ell\cdot r n$ for $n\in\mathbb{Z}^3$, where $r=r(S)$, and that $\int\mathrm{d}x\,f(\tilde{g}x)=\int\mathrm{d}x\,f(x)$ for every $\ell\mathbb{Z}^3$-periodic function $f$.

(Bm1)
By (I.T-g),
\[
\left(a_{\tilde{g}}\,\bar{\psi}(x)\right)^T=\left((S^{-1})^T\cdot\bar{\psi}\left(\tilde{g}^{-1}x\right)\right)^T=\bar{\psi}\left(\tilde{g}^{-1}x\right)^T\cdot S^{-1},\qquad a_{\tilde{g}}\,\psi(x')=S\cdot\psi\left(\tilde{g}^{-1}x'\right).
\]
Therefore
\[
a_{\tilde{g}}\,\mathrm{Bq}[K]=\int\mathrm{d}x\int\mathrm{d}x'\,\bar{\psi}\left(\tilde{g}^{-1}x\right)^T\cdot S^{-1}\cdot K(x,x')\cdot S\cdot\psi\left(\tilde{g}^{-1}x'\right).
\]
The integrand is $\ell\mathbb{Z}^3$-periodic in $x$ and in $x'$, so that $x$, $x'$ may be replaced by $\tilde{g}x$, $\tilde{g}x'$.
The function $K'$ is antiperiodic, since $\tilde{g}(x+\ell\cdot e_\mu)=\tilde{g}x+\ell\cdot r e_\mu=\tilde{g}x\pm\ell\cdot e_\mu$.

(Bm2)
The map $G_{\Lambda}$ multiplies $\bar{\psi}(x)$ by $\exp\left(-i\,e_k\cdot\Lambda(\pi x)\right)$ and $\psi(x')$ by $\exp\left(i\,e_k\cdot\Lambda(\pi x')\right)$.
As $\Lambda\circ\pi$ is $\ell\mathbb{Z}^3$-periodic, $K'$ is the product of $K$ with a function that is periodic in $x$ and in $x'$; hence $K'$ is antiperiodic in $x$ and in $x'$.

(Bm3)
By (I.T-C), $\mathrm{Cg}\,\bar{\psi}_{\beta}(x)=-\sum_a\mathsf{C}_{\beta a}\cdot\psi_a(x)$ and $\mathrm{Cg}\,\psi_{\beta'}(x')=\sum_b\mathsf{C}_{\beta' b}\cdot\bar{\psi}_b(x')$, with spinor indices $a$, $b$.
Since $\psi_a(x)$ and $\bar{\psi}_b(x')$ anticommute, one obtains, all sums ranging over $\beta$, $\beta'$, $a$, $b$,
\begin{align*}
\mathrm{Cg}\left[\sum_{\beta,\beta'}\bar{\psi}_{\beta}(x)\cdot K(x,x')_{\beta\beta'}\cdot\psi_{\beta'}(x')\right]
&=-\sum\mathsf{C}_{\beta a}\cdot K(x,x')_{\beta\beta'}\cdot\mathsf{C}_{\beta' b}\cdot\psi_a(x)\cdot\bar{\psi}_b(x')\\
&=\sum\bar{\psi}_b(x')\cdot\left[\mathsf{C}_{\beta' b}\cdot K(x,x')_{\beta\beta'}\cdot\mathsf{C}_{\beta a}\right]\cdot\psi_a(x)\\
&=\bar{\psi}(x')^T\cdot\mathsf{C}^T\cdot K(x,x')^T\cdot\mathsf{C}\cdot\psi(x).
\end{align*}
Exchange the names $x$, $x'$ in the double integral and use $\mathsf{C}^T=\mathsf{C}^{-1}$ (I.c3).
The function $K'$ arises from $K$ by exchanging the two arguments, transposing, and conjugating with constant matrices.
Hence it is antiperiodic in $x$, as $K$ is in its second argument, and in $x'$, as $K$ is in its first argument.

(Bm4)
The map $\mathbf{K}$ is antilinear and multiplicative, and both the weights $\left(\eta_k^3/48\right)^2$ and the matrix $\varkappa$ are real.
By (I.T-K),
\[
\mathbf{K}\left[\bar{\psi}(x)^T\cdot K(x,x')\cdot\psi(x')\right]=\left(\varkappa\cdot\bar{\psi}(x)\right)^T\cdot\overline{K(x,x')}\cdot\varkappa\cdot\psi(x')=\bar{\psi}(x)^T\cdot\varkappa^T\cdot\overline{K(x,x')}\cdot\varkappa\cdot\psi(x'),
\]
and $\varkappa^T=\varkappa=\varkappa^{-1}$ by (I.c4).
Complex conjugation and multiplication by constant matrices preserve the relations $K(x+\ell\cdot e_\mu,x')=-K(x,x')=K(x,x'+\ell\cdot e_\mu)$.
\end{proof}

\subsubsection*{0.4. Lemma C (local quadratic terms; this is Proposition~16.2\,(iv))}

\begin{named}{Lemma C}
Let $\tilde{g}=(S,a)\in\mathbb{G}_k$ and let $\Lambda$ be a real function of $c$, as in Lemma~B.
For a $4\times4$ matrix $\mathrm{Mx}$ and a union $Z$ of $M$-cubes of $T_k$ put
\[
\mathrm{Ql}_Z(\mathrm{Mx}):=\int_Z\mathrm{d}x\,\bar{\psi}(x)^T\cdot\mathrm{Mx}\cdot\psi(x).
\]
\begin{itemize}
\item[(Cq1)] $\mathrm{Ql}_Z(\mathrm{Mx})$ is balanced and of bidegree $(2,0)$, it contains only generators $\psi(\xi)$ with $\xi$ over $Z$, and
\[
\left\|\mathrm{Ql}_Z(\mathrm{Mx})\right\|_{\mathsf{w}}=\mathsf{w}_1^2\cdot\operatorname{Vol}(Z)\cdot\sum_{\beta,\beta'}\left|\mathrm{Mx}_{\beta\beta'}\right|.
\]
If $Z$ has a lift $\tilde{Z}$, then $\mathrm{Ql}_Z(\mathrm{Mx})_{20}(x,x')=48\eta_k^{-3}\cdot\mathbf{1}[x=x']\cdot\mathrm{Mx}$ for $x$, $x'$ over $\tilde{Z}$, and
\[
\int_{\tilde{Z}}\mathrm{d}x\int_{\tilde{Z}}\mathrm{d}x'\,\mathrm{Ql}_Z(\mathrm{Mx})_{20}(x,x')=\operatorname{Vol}(Z)\cdot\mathrm{Mx}.
\]
\item[(Cq2)] One has
\begin{gather*}
a_{\tilde{g}}\,\mathrm{Ql}_{\tilde{g}Z}\left(S\cdot\mathrm{Mx}\cdot S^{-1}\right)=\mathrm{Ql}_Z(\mathrm{Mx}),\qquad G_{\Lambda}\,\mathrm{Ql}_Z(\mathrm{Mx})=\mathrm{Ql}_Z(\mathrm{Mx}),\\
\mathrm{Cg}\,\mathrm{Ql}_Z(\mathrm{Mx})=\mathrm{Ql}_Z\left(\mathsf{C}^{-1}\cdot\mathrm{Mx}^T\cdot\mathsf{C}\right),\qquad\mathbf{K}\,\mathrm{Ql}_Z(\mathrm{Mx})=\mathrm{Ql}_Z\left(\varkappa^{-1}\cdot\overline{\mathrm{Mx}}\cdot\varkappa\right).
\end{gather*}
\end{itemize}
\end{named}

\begin{proof}[Proof of Lemma C]
Recall the rule $\int_{\tilde{g}Z}\mathrm{d}x\,f(x)=\int_Z\mathrm{d}x\,f(\tilde{g}x)$ of Proposition~16.1\,(a2).

(Cq1)
With $x$ ranging over a set of representatives of the fermion sites over $Z$,
\[
\mathrm{Ql}_Z(\mathrm{Mx})=\sum_{x\text{ over }Z}\frac{\eta_k^3}{48}\cdot\sum_{\beta,\beta'}\mathrm{Mx}_{\beta\beta'}\cdot\bar{\psi}_{\beta}(x)\cdot\psi_{\beta'}(x);
\]
in particular only generators $\psi(\xi)$ with $\xi$ over $Z$ occur.
For different triples $(x,\beta,\beta')$ the products $\bar{\psi}_{\beta}(x)\cdot\psi_{\beta'}(x)$ are, up to signs, different basis monomials.
The norm is therefore
\[
\sum_{x}\frac{\eta_k^3}{48}\cdot\sum_{\beta,\beta'}\left|\mathrm{Mx}_{\beta\beta'}\right|\cdot\mathsf{w}_1^2=\mathsf{w}_1^2\cdot\operatorname{Vol}(Z)\cdot\sum_{\beta,\beta'}\left|\mathrm{Mx}_{\beta\beta'}\right|.
\]
Now let
\[
K(x,x'):=48\eta_k^{-3}\cdot\mathbf{1}[\mathrm{pr}\,\pi(x)\in Z]\cdot\mathrm{Mx}\cdot\sum_{n\in\mathbb{Z}^3}(-1)^{n_1+n_2+n_3}\cdot\mathbf{1}[x'=x+\ell\cdot n].
\]
This function is antiperiodic in $x$ and in $x'$.
For $x$ over $Z$, only the one representative $x'=x+\ell\cdot n$ of the class of $x$ contributes to $\int\mathrm{d}x'\,K(x,x')\cdot\psi(x')$, and its contribution is
\[
\frac{\eta_k^3}{48}\cdot48\eta_k^{-3}\cdot(-1)^{n_1+n_2+n_3}\cdot\mathrm{Mx}\cdot\psi(x+\ell\cdot n)=\mathrm{Mx}\cdot\psi(x);
\]
for the other $x$ this integral vanishes.
Consequently $\mathrm{Bq}[K]=\mathrm{Ql}_Z(\mathrm{Mx})$, and by (Ak4) the element is balanced, of bidegree $(2,0)$, with the kernel $K$.
If $x$ and $x'$ lie over a lift $\tilde{Z}$ and $x'=x+\ell\cdot n$, then $n=0$, because $\mathrm{pr}$ is injective on $\tilde{Z}$; this gives the kernel over $\tilde{Z}$.
Over $\tilde{Z}$ lie $48(M/\eta_k)^3\cdot|Z|_M$ fermion sites, so that the double integral equals
\[
\left(\frac{\eta_k^3}{48}\right)^2\cdot48\eta_k^{-3}\cdot48\left(\frac{M}{\eta_k}\right)^3\cdot|Z|_M\cdot\mathrm{Mx}=M^3\cdot|Z|_M\cdot\mathrm{Mx}=\operatorname{Vol}(Z)\cdot\mathrm{Mx}.
\]

(Cq2)
By (I.T-g) and the rule of Proposition~16.1\,(a2) recalled above,
\begin{align*}
a_{\tilde{g}}\,\mathrm{Ql}_{\tilde{g}Z}\left(S\cdot\mathrm{Mx}\cdot S^{-1}\right)
&=\int_{\tilde{g}Z}\mathrm{d}x\,\bar{\psi}\left(\tilde{g}^{-1}x\right)^T\cdot S^{-1}\cdot\left(S\cdot\mathrm{Mx}\cdot S^{-1}\right)\cdot S\cdot\psi\left(\tilde{g}^{-1}x\right)\\
&=\int_Z\mathrm{d}x\,\bar{\psi}(x)^T\cdot\mathrm{Mx}\cdot\psi(x).
\end{align*}
Under $G_{\Lambda}$ the phase $\exp\left(-i\,e_k\cdot\Lambda(c)\right)$ of $\bar{\psi}(x)$ and the phase $\exp\left(i\,e_k\cdot\Lambda(c)\right)$ of $\psi(x)$ cancel.
The last two relations are the computations of (Bm3) and (Bm4) at $x'=x$:
\[
\mathrm{Cg}\left[\bar{\psi}(x)^T\cdot\mathrm{Mx}\cdot\psi(x)\right]=\bar{\psi}(x)^T\cdot\mathsf{C}^T\cdot\mathrm{Mx}^T\cdot\mathsf{C}\cdot\psi(x),\qquad
\mathbf{K}\left[\bar{\psi}(x)^T\cdot\mathrm{Mx}\cdot\psi(x)\right]=\bar{\psi}(x)^T\cdot\varkappa^T\cdot\overline{\mathrm{Mx}}\cdot\varkappa\cdot\psi(x),
\]
together with $\mathsf{C}^T=\mathsf{C}^{-1}$ and $\varkappa^T=\varkappa^{-1}$.
\end{proof}

\subsubsection*{0.5. Lemma D (a polymer and its elements in a lift)}

\begin{named}{Lemma D}
Let $X$ be a polymer of $T_k$ that has a lift $\tilde{X}$ (Section~I.1.1) with
\[
\operatorname{sup-diameter}(\tilde{X})+2\eta_k+M<\ell,
\]
and put $\tilde{U}:=\left\{c\in\mathbb{L}_k:\ |c-c'|\le\eta_k\text{ for some }c'\in\tilde{X}\right\}$.
\begin{itemize}
\item[(Dl1)] $\tilde{X}$ is a finite union of $M$-cubes of $\mathbb{L}_k$, connected through faces.
The map $\mathrm{pr}$ sends the cubes of $\tilde{X}$ bijectively onto the cubes of $X$, and centres onto centres; two cubes of $\tilde{X}$ have a common face if and only if their images have one.
The number $\operatorname{Vol}(\tilde{X})$, defined as $M^3$ times the number of the cubes of $\tilde{X}$, equals $\operatorname{Vol}(X)$.
Moreover $\operatorname{sup-diameter}(\tilde{U})+M<\ell$, and $\mathrm{pr}$ is injective on $\tilde{U}$.
A site of $\tilde{U}$ whose image lies in $X$ belongs to $\tilde{X}$.
\item[(Dl2)] $X$ is small if and only if $d_M(\tilde{X})\le L$, where $d_M(\tilde{X})$ is formed with nets in $\mathbb{R}^3$, as in \cite[Theorem~13.2\,(t10)]{QEDI}.
Moreover, for every small polymer $X$ of $T_k$, without the hypothesis of the lemma, the following holds:
$X$ has a lift of sup-diameter at most $M\cdot(L+2)$, namely the set of the sites of the set $\tilde{X}'$ of (t10);
every lift of $X$ is a translate of this one by a vector of $\ell\mathbb{Z}^3$;
hence every lift of $X$ has sup-diameter at most $M\cdot(L+2)$ and satisfies the hypothesis of the lemma.
\item[(Dl3)] Every fermion site over $X$ has exactly one representative over $\tilde{X}$.
Every pair of $\mathrm{Pz}_k$ that is admissible for $X$ has exactly one representative $(x,x')\in\mathrm{Pzt}_k$ with $\pi(x),\pi(x')\in\tilde{U}$;
if the pair is of kind (I) for $X$, then $\pi(x)$, $\pi(x')$ and the six paths $\boldsymbol{\Gamma}^{\pi}(\pi x,\pi x')$ lie in $\tilde{X}$.
\item[(Dl4)] Let $F\in\mathrm{Gr}[\psi,\chi]$ be such that its kernels vanish unless all $\xi_i$ lie over $X$ and all $\boldsymbol{\zeta}_j$ are admissible for $X$.
Let $\tilde{F}$ be the element of $\mathrm{Grc}_k(\tilde{U})$ (Section~I.1.9) that has the kernels of $F$, restricted to the arguments over $\tilde{U}$:
\begin{multline*}
\tilde{F}:=\sum_{n,m}\frac{1}{n!\,m!}\sum\nu(\xi_1)\cdots\nu(\xi_n)\cdot\nu(\boldsymbol{\zeta}_1)\cdots\nu(\boldsymbol{\zeta}_m)\cdot F_{nm}(\xi;\boldsymbol{\zeta})\\
\cdot\psi(\xi_1)\cdots\psi(\xi_n)\cdot\chi(\boldsymbol{\zeta}_1)\cdots\chi(\boldsymbol{\zeta}_m),
\end{multline*}
where the inner sum extends over the $\xi_i$ of the form $(x,\beta,\omega)$ with $\pi(x)\in\tilde{U}$ and over the $\boldsymbol{\zeta}_j$ of the form $(x,x',\beta,\omega)$ with $(x,x')\in\mathrm{Pzt}_k$ and $\pi(x),\pi(x')\in\tilde{U}$.
Then $\mathrm{cov}_\ell(\tilde{F})=F$ and $\mathrm{cov}_\ell\left(\tilde{F}^{[n,m]}\right)=F^{[n,m]}$;
$\tilde{F}$ is the only element of $\mathrm{Grc}_k(\tilde{U})$ with $\mathrm{cov}_\ell(\tilde{F})=F$;
and $\tilde{F}$ contains generators $\psi(\xi)$ only with $\xi$ over $\tilde{X}$.
\end{itemize}
\end{named}

\begin{proof}[Proof of Lemma D]
Recall that $\ell$ is a multiple of $\eta_k$ and of $M$, that $\ell\ge13L\cdot M$, that $\eta_k\le1\le M$, and that $M/\eta_k$ is an odd integer.
The sets $\mathrm{Pzt}_k$ and $\mathrm{Pz}_k$ of pairs, the admissibility of a pair for $X$ and the kinds (I) and (II) are those of Section~I.1.7.

(Dl1)
A lift of a union of $M$-cubes is a union of $M$-cubes of $\mathbb{L}_k$ with the same common faces (Section~I.1.1), and $\mathrm{pr}$ is injective on $\tilde{X}$.
Since $X$ is connected through faces, so is $\tilde{X}$; the cubes of $\tilde{X}$ and those of $X$ correspond bijectively, so that their numbers agree and $\operatorname{Vol}(\tilde{X})=\operatorname{Vol}(X)$.
As $\ell$ is a multiple of $M$, the map $\mathrm{pr}$ sends $M\mathbb{Z}^3$ into $M\mathbb{Z}^3/\ell\mathbb{Z}^3$, that is, centres to centres.
The sup-diameter of $\tilde{U}$ is at most $\operatorname{sup-diameter}(\tilde{X})+2\eta_k$, whence $\operatorname{sup-diameter}(\tilde{U})+M<\ell$.
Two points of $\tilde{U}$ that are congruent modulo $\ell\mathbb{Z}^3$ differ by less than $\ell$ in the sup norm, and therefore coincide.
Finally, let $c\in\tilde{U}$ with $\mathrm{pr}\,c\in X$.
There is a $c'\in\tilde{X}$ with $\mathrm{pr}\,c'=\mathrm{pr}\,c$; as $c'$ lies in $\tilde{U}$ as well, $c=c'$.

(Dl2)
The sites of the cube of centre $z$ fill the range $|c-z|\le(M-\eta_k)/2$.
Hence the closed cubes of $\tilde{X}$ form a set of sup-diameter $\operatorname{sup-diameter}(\tilde{X})+\eta_k<\ell$, and $\mathrm{pr}$ is injective on their union.
Let $X$ be small.
Then $2M\cdot\left(d_M(X)+2\right)\le2M\cdot(L+2)<\ell$, and by (t10) one has $X=\mathrm{pr}(\tilde{X}')$ for a finite union $\tilde{X}'$ of $M$-cubes, connected through faces, on which $\mathrm{pr}$ is injective, with $d_M(\tilde{X}')=d_M(X)$; and $\tilde{X}'$ is unique up to $\ell\mathbb{Z}^3$ among the sets of this kind.
By (Dl1), $\tilde{X}$ is a set of this kind.
Hence $\tilde{X}=\tilde{X}'+\ell\cdot n$ for some $n\in\mathbb{Z}^3$, and $d_M(\tilde{X})=d_M(\tilde{X}')=d_M(X)\le L$, since a translation maps nets to nets of the same length.
Conversely, let $d_M(\tilde{X})\le L$.
Then $2M\cdot\left(d_M(\tilde{X})+2\right)<\ell$, and the converse part of (t10) yields $d_M(X)=d_M(\mathrm{pr}\,\tilde{X})=d_M(\tilde{X})\le L$.

\emph{The small polymers.}
Let $X$ be an arbitrary small polymer of $T_k$.
Then $2M\cdot\left(d_M(X)+2\right)\le2M\cdot(L+2)<\ell$.
By the first part of (t10), $X=\mathrm{pr}(\tilde{X}')$ for a finite union $\tilde{X}'$ of $M$-cubes of $\mathbb{R}^3$ (centred at points of $M\mathbb{Z}^3$), connected through faces, on which $\mathrm{pr}$ is injective, and $d_M(\tilde{X}')=d_M(X)\le L$.
By \cite[Theorem~13.2\,(t7)]{QEDI}, applied in $\mathbb{R}^3$, any two sites of $\tilde{X}'$ have sup-distance at most $M\cdot\left(d_M(\tilde{X}')+2\right)\le M\cdot(L+2)$, and $M\cdot(L+2)+\eta_k<\ell$.
Since $\ell$ is a multiple of $\eta_k$, the lattice $\ell\mathbb{Z}^3$ is contained in $\mathbb{L}_k$.
The sites of the cube of $\tilde{X}'$ with the centre $z\in M\mathbb{Z}^3$ are the $c\in\mathbb{L}_k$ with $|c-z|<M/2$;
$\mathrm{pr}$ maps them onto the sites of $T_k$ at distance less than $M/2$ from $\mathrm{pr}\,z$, that is, onto the sites of the $M$-cube of $T_k$ with the centre $\mathrm{pr}\,z$, and these $M$-cubes are the cubes of $X$.
Thus $\mathrm{pr}$ maps the sites of $\tilde{X}'$ onto the sites of $X$, and these sites form a lift of $X$ in the sense of Section~I.1.1, of sup-diameter at most $M\cdot(L+2)$.
The sites of an $M$-cube form a product of three intervals of $\eta_k\mathbb{Z}$, so that an $M$-cube is connected through nearest neighbours.
Two $M$-cubes with a common face have centres represented by $z$ and $z+M\cdot e_\mu$ with $z\in M\mathbb{Z}^3$;
the first cube contains the site $z+\frac{M-\eta_k}{2}\cdot e_\mu$, and the second one contains its nearest neighbour $z+\frac{M+\eta_k}{2}\cdot e_\mu$, because $M/\eta_k$ is an odd integer.
As $X$ is connected through faces, it is connected through nearest neighbours, and by Section~I.1.1 every lift of $X$ is a translate of the lift just constructed by a vector of $\ell\mathbb{Z}^3$.
A translation does not change the sup-diameter, and
\[
M\cdot(L+2)+2\eta_k+M\le M\cdot(L+2)+2+M<13L\cdot M\le\ell.
\]

(Dl3)
The first assertion holds because $\mathrm{pr}$ maps $\tilde{X}$ bijectively onto $X$ and leaves the label untouched.
Let a pair of $\mathrm{Pz}_k$ be of kind (I) for $X$, with the photon sites $c,c'\in X$, and let $\tilde{c}$, $\tilde{c}'$ be their lifts in $\tilde{X}$.
As $\operatorname{sup-diameter}(\tilde{X})+M<\ell$ and $d(c,c')\le M$, the vector $\tilde{c}-\tilde{c}'$ is the only representative of $c-c'$ modulo $\ell\mathbb{Z}^3$ of sup norm at most $M$, and the pairs at distance at most $M$ inside $X$, together with their six paths, are those of $\tilde{X}$.
This is Section~I.1.1 applied with $\mathrm{sd}=M$: a site of a path $\boldsymbol{\Gamma}^{\pi}(\tilde{c},\tilde{c}')$ has distance at most $M$ from $\tilde{c}$, its image lies in $X$ by the definition of kind (I), and $\operatorname{sup-diameter}(\tilde{X})+M<\ell$; so it is the lift in $\tilde{X}$ of its image.
The pair has a representative in $\mathrm{Pzt}_k$; translate it so that its first photon site becomes $\tilde{c}$.
Its second photon site $\mathrm{cs}$ then satisfies $\mathrm{pr}\,\mathrm{cs}=c'$ and $|\mathrm{cs}-\tilde{c}|\le M$, so that $\mathrm{cs}-\tilde{c}$ is a representative of $c'-c$ of sup norm at most $M$, and $\mathrm{cs}=\tilde{c}'$.
For $c=c'$, that is, for a pair inside a cluster over $X$, this means: $\tilde{c}'=\tilde{c}$, the six paths reduce to the single site $\tilde{c}$, and $\mathrm{cs}-\tilde{c}$ is a representative of $0$ of sup norm at most $M<\ell$, whence $\mathrm{cs}=\tilde{c}$.
Now let the pair be of kind (II): $x'=x^{(3)}$, and $\pi(x)$ or $\pi\left(x^{(3)}\right)$ lies in $X$.
Choose the representative for which this site lies in $\tilde{X}$; the other site has distance $\eta_k$ from it (I.1.2) and therefore lies in $\tilde{U}$.
In both cases the representative is unique, since two representatives differ by a translation by some $\ell\cdot n$ and $\mathrm{pr}$ is injective on $\tilde{U}$.

(Dl4)
By (Dl1), $\operatorname{sup-diameter}(\tilde{U})+M<\ell$; hence $\mathrm{cov}_\ell$ is defined on $\mathrm{Grc}_k(\tilde{U})$ and is injective (Section~I.1.9).
It sends every generator to the generator of the torus with the same index.
Therefore $\mathrm{cov}_\ell(\tilde{F})$ is the sum displayed in (Dl4), read in $\mathrm{Gr}[\psi,\chi]$, and $\mathrm{cov}_\ell$ preserves the bidegree.
By the hypothesis on $F$ and the last assertion of (Dl1), a term of this sum vanishes unless all $\xi_i$ lie over $\tilde{X}$ and the pairs of all $\boldsymbol{\zeta}_j$ represent pairs that are admissible for $X$.
By (Dl3) the remaining arguments contain exactly one representative of each class at which the kernels of $F$ can be different from zero.
The products of kernels with generators do not depend on the choice of the representatives (Section~0.1).
Hence the sum is the representation (Ak1) of $F$.
The uniqueness of $\tilde{F}$ follows from the injectivity of $\mathrm{cov}_\ell$, and the last assertion from the vanishing of the terms in which some $\xi_i$ does not lie over $\tilde{X}$.
\end{proof}

\subsection*{1. Proposition~48.1}

\begin{proof}[Proof of Proposition~48.1]
Of the conditions on the constants, only the following enter:
the first condition of (A8) in Section~I.8.4, $0<\varepsilon\le\frac{1}{200}$, $\eta_k\le1<M/2$ and $\kappa\ge0$, and, from the line (C0) of Section~I.8.2, $0<\rho<1$ and $w_3=6$.
Of the range of $\rho$ only $0\le\rho\le1$ is needed, namely for the bound $24$ of (I.m7) in the proof of (k6)$^{\mathrm{o}}$ for the family $E^{\mathrm{bd}}$ of (c);
the value $w_3=6$ enters through the formula $M_3(x)=3\left(\rho-\gamma\cdot(ge_1)\right)$ of (I.2.1), through the bound $24$ and through the coefficient $-3\eta_k^3/48$ in the proof of (c).
Throughout, (k1)--(k7) are the clauses of Section~I.5.4, and (k1)$^{\mathrm{o}}$--(k7)$^{\mathrm{o}}$ are the clauses of the Definition in Proposition~48.1.

\emph{The clauses are meaningful.}
By \cite[Lemma~5.19]{QEDI}, the domains $\tilde{R}_k$ and $\mathfrak{X}_k$ of Section~I.5.3 are invariant under $A\mapsto\tilde{g}A$, under $A\mapsto-A$ and under $A\mapsto\overline{A}$;
moreover, $\tilde{g}$ maps $D_k$ onto $D_k$ (Proposition~16.1\,(a1)).
Hence the left-hand sides of (k3)$^{\mathrm{o}}$, (k5)$^{\mathrm{o}}$, (k7)$^{\mathrm{o}}$ and of (k3), (k5), (k7) are defined.
As the elements are even, $a_{\tilde{g}}[\,\cdot\,]$ does not depend on the representative $a\in M\mathbb{Z}^3$ of the translation part of $\tilde{g}$ (Section~I.4.3).
The words ``the kernels depend on $\mathscr{A}(b)$ only for the bonds $b$ that meet $X$'' mean that the kernels take equal values at any two points of the domain which coincide on these bonds \cite[Definition~5.5\,(1)]{QEDI}.

\emph{Proof of (a).}
Let $\mathrm{td1}_k:=e_k^{-1+2\varepsilon}$ and $\mathrm{td2}_k:=e_k^{-1+3\varepsilon}$ be the two thresholds of $\mathfrak{X}_k$ (G2.r4).
Since $|0|_\infty=0<\mathrm{td1}_k$ and $|\partial0|_\infty=0<\mathrm{td2}_k$, the bond function $0$ belongs to $\mathfrak{X}_k$.
Consequently $(\mathscr{A},0)\in\tilde{R}_k\times\mathfrak{X}_k$ for every $\mathscr{A}\in\tilde{R}_k$, and $\mathscr{E}|_0$ is a family without slot of even, balanced elements.
An element determines its kernels (Ak1), and therefore
\[
(\mathscr{E}|_0)_{nm}(X;\mathscr{A};\,\cdot\,)=\mathscr{E}_{nm}(X;\mathscr{A},0;\,\cdot\,).
\]
\begin{itemize}
\item[(k1)$^{\mathrm{o}}$] The map $\mathscr{A}\mapsto(\mathscr{A},0)$ is an affine map of the open set $\tilde{R}_k$ into the open set $\tilde{R}_k\times\mathfrak{X}_k$ \cite[Lemma~17.3]{QEDI}, and an analytic function composed with it is analytic.
\item[(k2)$^{\mathrm{o}}$] The first clause is the first clause of (k2) at the point $(\mathscr{A},0)$.
If $\mathscr{A},\tilde{\mathscr{A}}\in\tilde{R}_k$ coincide on the bonds that meet $X$, then $(\mathscr{A},0)$ and $(\tilde{\mathscr{A}},0)$ coincide there in both components, and the kernels agree by (k2).
\item[(k3)$^{\mathrm{o}}$] As $\tilde{g}0=0$, this is (k3) at $(\mathscr{A},0)$.
\item[(k4)$^{\mathrm{o}}$] Let $\mathscr{A}$, $\mathscr{A}^1$ and $\boldsymbol{\lambda}$ be as in (k4)$^{\mathrm{o}}$.
The pairs $(\mathscr{A},0)$ and $(\mathscr{A}^1,0)$ lie in $\tilde{R}_k\times\mathfrak{X}_k$, and with $\boldsymbol{\lambda}':=0$ one has $\mathscr{A}^1=\mathscr{A}-d\boldsymbol{\lambda}$ and $0=0-d\boldsymbol{\lambda}'$ on the bonds at distance $\le1$ from $X$.
Clause (k4) yields $G_{\boldsymbol{\lambda}+0}\left[\mathscr{E}(X;\mathscr{A}^1,0)\right]=\mathscr{E}(X;\mathscr{A},0)$.
\item[(k5)$^{\mathrm{o}}$, (k7)$^{\mathrm{o}}$] Since $-0=0$ and $-\overline{0}=0$, these are (k5) and (k7) at $(\mathscr{A},0)$.
\item[(k6)$^{\mathrm{o}}$] The supremum over $\tilde{R}_k\times\{0\}$ does not exceed the supremum over $\tilde{R}_k\times\mathfrak{X}_k$;
hence $\|\mathscr{E}|_0\|_k\le\|\mathscr{E}\|^{x}_{k}<\infty$.
\end{itemize}
Finally, $\mathscr{E}\mapsto\mathscr{E}|_0$ is an evaluation and is therefore linear.

\emph{Proof of (b).}
The elements $E^{\mathrm{c}}(X;\mathscr{A},\mathscr{C})$ and $E(X;\mathscr{A})$ are equal, so that by (Ak1) they have the same kernels.
Hence the first clause of (k2) and the first clause of (k2)$^{\mathrm{o}}$ are one and the same condition.
Suppose that $E$ has (k2)$^{\mathrm{o}}$, and let $(\mathscr{A},\mathscr{C})$ and $(\tilde{\mathscr{A}},\tilde{\mathscr{C}})$ in $\tilde{R}_k\times\mathfrak{X}_k$ coincide, in both components, on the bonds that meet $X$.
Then $\mathscr{A}$ and $\tilde{\mathscr{A}}$ coincide on these bonds, and
\[
E^{\mathrm{c}}(X;\mathscr{A},\mathscr{C})=E(X;\mathscr{A})=E(X;\tilde{\mathscr{A}})=E^{\mathrm{c}}(X;\tilde{\mathscr{A}},\tilde{\mathscr{C}}).
\]
Suppose conversely that $E^{\mathrm{c}}$ has (k2), and let $\mathscr{A},\tilde{\mathscr{A}}\in\tilde{R}_k$ coincide on the bonds that meet $X$.
By (a) the pairs $(\mathscr{A},0)$ and $(\tilde{\mathscr{A}},0)$ lie in $\tilde{R}_k\times\mathfrak{X}_k$, and they coincide on these bonds in both components;
thus
\[
E(X;\mathscr{A})=E^{\mathrm{c}}(X;\mathscr{A},0)=E^{\mathrm{c}}(X;\tilde{\mathscr{A}},0)=E(X;\tilde{\mathscr{A}}).
\]
Lastly, $E^{\mathrm{c}}|_0(X;\mathscr{A})=E^{\mathrm{c}}(X;\mathscr{A},0)=E(X;\mathscr{A})$.

\emph{Proof of (c): a family in $K_k$ whose constant extension violates (k4).}
Let $\square$ be an $M$-cube.
For a fermion site $x=(c,g)$ write
\[
c_3=\pi\left(x^{(3)}\right),\qquad
u_{\mathscr{A}}(c,c_3)=\exp\left\{i\,e_k\cdot\eta_k\cdot\mathscr{A}(c,c_3)\right\},\qquad
M_3(x)=3\left(\rho-\gamma\cdot(ge_1)\right)
\]
(see (I.1.2), (I.1.6) and (I.2.1)), and define
\[
E^{\mathrm{bd}}(\square;\mathscr{A}):=\int_{\square}\mathrm{d}x\,\mathbf{1}[c_3\in\square]\,\bar{\psi}(x)^{T}M_3(x)\,u_{\mathscr{A}}(c,c_3)\,\psi\left(x^{(3)}\right),
\qquad
E^{\mathrm{bd}}(X;\mathscr{A}):=0\quad\text{if }|X|_M\ge2.
\]
The integrand is $\ell\mathbb{Z}^3$-periodic in $x$:
one has $(x+\ell e_\mu)^{(3)}=x^{(3)}+\ell e_\mu$ (I.s1), the signs of $\bar{\psi}$ and of $\psi$ cancel, $M_3(x)$ depends on the label only, and $\mathscr{A}$ is a function on the bonds of $T_k$.

\emph{The family $E^{\mathrm{bd}}$ lies in $K_k$.}
Each monomial is a product $\bar{\psi}_\beta(x)\psi_{\beta'}\left(x^{(3)}\right)$; so the elements are even and balanced.
\begin{itemize}
\item[(k1)$^{\mathrm{o}}$] With respect to the basis of the monomials the coefficients are entire functions of $\mathscr{A}$; by (Ak1) so are the kernels.
\item[(k2)$^{\mathrm{o}}$] The sites $x$ and $x^{(3)}$ lie over $\square$, no generator $\chi$ occurs, and $\mathscr{A}$ is read on the bonds $(c,c_3)$, both ends of which lie in $\square$.
\item[(k3)$^{\mathrm{o}}$] The map $c\mapsto\tilde{g}c=rc+a$ is an isometry for the sup norm which maps $M\mathbb{Z}^3$ onto itself;
it therefore permutes the $M$-cubes, and $|\tilde{g}X|_M=|X|_M$.
By (I.T-g),
\begin{multline*}
a_{\tilde{g}}\left[E^{\mathrm{bd}}(\tilde{g}\square;\tilde{g}\mathscr{A})\right]\\
=\int_{\tilde{g}\square}\mathrm{d}x\,\mathbf{1}\left[\pi\left(x^{(3)}\right)\in\tilde{g}\square\right]\bar{\psi}\left(\tilde{g}^{-1}x\right)^{T}S^{-1}M_3(x)\,S\,u_{\tilde{g}\mathscr{A}}\left(\pi x,\pi x^{(3)}\right)\psi\left(\tilde{g}^{-1}\left(x^{(3)}\right)\right).
\end{multline*}
Replace $x$ by $\tilde{g}x$.
The map $x\mapsto\tilde{g}x$ carries $\square\times O_{\mathrm{cub}}$ onto $\tilde{g}\square\times O_{\mathrm{cub}}$ and preserves $\int\mathrm{d}x$ (I.a1).
Further, $(\tilde{g}x)^{(3)}=\tilde{g}\left(x^{(3)}\right)$ (I.a2), $S^{-1}M_3(\tilde{g}x)\,S=M_3(x)$ (I.m4), and $u_{\tilde{g}\mathscr{A}}(\tilde{g}c,\tilde{g}c_3)=u_{\mathscr{A}}(c,c_3)$, because $(\tilde{g}\mathscr{A})(\tilde{g}c,\tilde{g}c_3)=\mathscr{A}(c,c_3)$ (Section~I.4.3);
finally, $\tilde{g}c_3\in\tilde{g}\square$ if and only if $c_3\in\square$.
The result is $E^{\mathrm{bd}}(\square;\mathscr{A})$.
\item[(k4)$^{\mathrm{o}}$] Let $\mathscr{A}^1=\mathscr{A}-d\boldsymbol{\lambda}$ on the bonds at distance $\le1$ from $\square$.
For $c,c_3\in\square$ one has
\[
u_{\mathscr{A}^1}(c,c_3)=u_{\mathscr{A}}(c,c_3)\cdot\exp\left\{-i\,e_k\cdot\eta_k\cdot(d\boldsymbol{\lambda})(c,c_3)\right\}=u_{\mathscr{A}}(c,c_3)\cdot\exp\left\{-i\,e_k\cdot\left(\boldsymbol{\lambda}(c_3)-\boldsymbol{\lambda}(c)\right)\right\}.
\]
The automorphism $G_{\boldsymbol{\lambda}}$ multiplies $\bar{\psi}(x)$ by $\exp\left\{-i\,e_k\cdot\boldsymbol{\lambda}(c)\right\}$ and $\psi\left(x^{(3)}\right)$ by $\exp\left\{i\,e_k\cdot\boldsymbol{\lambda}(c_3)\right\}$.
The three phases cancel, and $G_{\boldsymbol{\lambda}}\left[E^{\mathrm{bd}}(\square;\mathscr{A}^1)\right]=E^{\mathrm{bd}}(\square;\mathscr{A})$.
\item[(k5)$^{\mathrm{o}}$] The computation in the proof of (Bm3), which rests on Section~I.4.5 and on (I.c3) only, gives
\[
\mathrm{Cg}\left[\bar{\psi}(x)^{T}M_3(x)\,\psi\left(x^{(3)}\right)\right]=\bar{\psi}\left(x^{(3)}\right)^{T}\mathsf{C}^{-1}M_3(x)^{T}\mathsf{C}\,\psi(x)=\bar{\psi}\left(x^{(3)}\right)^{T}M_3\left(x^{(3)}\right)\psi(x),
\]
the last equality by (I.m2);
moreover $u_{-\mathscr{A}}(c,c_3)=\exp\left\{-i\,e_k\cdot\eta_k\cdot\mathscr{A}(c,c_3)\right\}=u_{\mathscr{A}}(c_3,c)$.
Hence
\[
\mathrm{Cg}\left[E^{\mathrm{bd}}(\square;-\mathscr{A})\right]=\int_{\square}\mathrm{d}x\,\mathbf{1}[c_3\in\square]\,\bar{\psi}\left(x^{(3)}\right)^{T}M_3\left(x^{(3)}\right)u_{\mathscr{A}}(c_3,c)\,\psi(x).
\]
The map $x\mapsto x^{(3)}$ is an involution of $T^{\mathrm{f}}_k$ (I.s1);
it preserves the weights and maps the set $\{x:\ c\in\square,\ c_3\in\square\}$ onto itself.
Replacing $x$ by $x^{(3)}$ one obtains $E^{\mathrm{bd}}(\square;\mathscr{A})$.
\item[(k7)$^{\mathrm{o}}$] The computation in the proof of (Bm4), which rests on Section~I.4.5 and on (I.c4) only, gives
\[
\mathbf{K}\left[\bar{\psi}(x)^{T}M_3(x)\,u_{-\overline{\mathscr{A}}}(c,c_3)\,\psi\left(x^{(3)}\right)\right]=\bar{\psi}(x)^{T}\varkappa^{-1}\,\overline{M_3(x)}\,\varkappa\;\overline{u_{-\overline{\mathscr{A}}}(c,c_3)}\;\psi\left(x^{(3)}\right).
\]
By (I.m3) and (I.c4), $\varkappa^{-1}\,\overline{M_3(x)}\,\varkappa=\varkappa^{-2}M_3(x)\,\varkappa^{2}=M_3(x)$;
and $\overline{\exp\left\{-i\,e_k\cdot\eta_k\cdot\overline{\mathscr{A}(c,c_3)}\right\}}=u_{\mathscr{A}}(c,c_3)$.
Hence $\mathbf{K}\left[E^{\mathrm{bd}}(\square;-\overline{\mathscr{A}})\right]=E^{\mathrm{bd}}(\square;\mathscr{A})$.
\item[(k6)$^{\mathrm{o}}$] Write $\mathscr{A}=A_0+A_1$ as in \cite[Definition~5.16]{QEDI}.
Then $\operatorname{Im}\mathscr{A}=\operatorname{Im}A_1$ \cite[Lemma~17.3]{QEDI}, and $|A_1|_\infty<\mathrm{t1}_k$, where $\mathrm{t1}_k:=e_k^{-\frac{3}{4}+\varepsilon}$ is the first threshold of $R_k$ (G2.r4).
Moreover $\eta_k\le1$, and $e_k\cdot\mathrm{t1}_k=e_k^{\frac{1}{4}+\varepsilon}\le1$, since $e_k\le e_0<1$ by (A8).
Therefore
\[
\left|u_{\mathscr{A}}(c,c_3)\right|\le\exp\left(e_k\cdot\eta_k\cdot|\operatorname{Im}\mathscr{A}|_\infty\right)\le\mathrm{e}.
\]
The moduli of the $16$ entries of $M_3(x)$ add up to at most $24$ (I.m7), and a monomial $\bar{\psi}_\beta(x)\psi_{\beta'}\left(x^{(3)}\right)$ has the norm $h_{k,1}^2$.
Consequently $\left\|E^{\mathrm{bd}}(\square;\mathscr{A})\right\|_{h_k}\le24\,\mathrm{e}\cdot M^3\cdot h_{k,1}^2$ for all $\mathscr{A}\in\tilde{R}_k$;
as $D_k$ is finite, $\|E^{\mathrm{bd}}\|_k<\infty$.
\end{itemize}

\emph{The constant extension $(E^{\mathrm{bd}})^{\mathrm{c}}$ does not have (k4).}
Let $z$ be the centre of $\square$ and $x^0:=(z,1)$, where $1$ is the identity of $O_{\mathrm{cub}}$.
Then $(x^0)^{(3)}=(z+\eta_ke_1,s_3)$, and $z+\eta_ke_1\in\square$, because $\eta_k\le1<M/2$.
Put
\[
\delta:=\min\left\{\frac{\eta_k\cdot\mathrm{td1}_k}{2},\ \frac{\eta_k^2\cdot\mathrm{td2}_k}{4},\ \frac{\pi}{e_k}\right\},
\qquad
\boldsymbol{\lambda}'(c):=\delta\cdot\mathbf{1}[c=z],
\qquad
\mathscr{C}^1:=-d\boldsymbol{\lambda}';
\]
then $\boldsymbol{\lambda}'$ is a real function on the sites of $T_k$ which is not constant.
One has $|\mathscr{C}^1|_\infty=\delta/\eta_k<\mathrm{td1}_k$ and $|\partial\mathscr{C}^1|_\infty\le2\delta/\eta_k^2<\mathrm{td2}_k$, so that $\mathscr{C}^1\in\mathfrak{X}_k$.
The pairs $(\mathscr{A},\mathscr{C}):=(0,0)$ and $(\mathscr{A}^1,\mathscr{C}^1):=(0,\mathscr{C}^1)$ lie in $\tilde{R}_k\times\mathfrak{X}_k$, and with $\boldsymbol{\lambda}:=0$ one has $\mathscr{A}^1=\mathscr{A}-d\boldsymbol{\lambda}$ and $\mathscr{C}^1=\mathscr{C}-d\boldsymbol{\lambda}'$ on all bonds.
If $(E^{\mathrm{bd}})^{\mathrm{c}}$ had (k4), it would follow that
\[
G_{\boldsymbol{\lambda}'}\left[E^{\mathrm{bd}}(\square;0)\right]=E^{\mathrm{bd}}(\square;0).
\]
Up to signs, the products $\bar{\psi}_\beta(x)\psi_{\beta'}\left(x^{(3)}\right)$ that belong to different triples $(x,\beta,\beta')$ are different basis monomials, since $x^{(3)}$ is determined by $x$.
In $E^{\mathrm{bd}}(\square;0)$ the monomial $\bar{\psi}_1(x^0)\psi_2\left((x^0)^{(3)}\right)$ has the coefficient
\[
\frac{\eta_k^3}{48}\cdot M_3(x^0)_{12}=-3\cdot\frac{\eta_k^3}{48}\cdot(\gamma_1)_{12}=-\frac{3\eta_k^3}{48}\ne0;
\]
indeed $M_3(x^0)=3(\rho-\gamma_1)$, and $\gamma_1=\sigma_3\otimes\sigma_1=\operatorname{diag}(\sigma_1,-\sigma_1)$ has the entry $1$ at the place $(1,2)$.
The automorphism $G_{\boldsymbol{\lambda}'}$ multiplies this monomial by
\[
\exp\left\{i\,e_k\cdot\left(\boldsymbol{\lambda}'(z+\eta_ke_1)-\boldsymbol{\lambda}'(z)\right)\right\}=\exp(-i\,e_k\cdot\delta),
\]
and $0<e_k\cdot\delta\le\pi$.
So the two sides differ.
More generally, by (Bm2), $G_{\boldsymbol{\lambda}'}$ multiplies the kernel $K(x,x')$ of a term $\bar{\psi}(x)^{T}K(x,x')\,\psi(x')$ by $\exp\left\{i\,e_k\cdot\left(\boldsymbol{\lambda}'(\pi x')-\boldsymbol{\lambda}'(\pi x)\right)\right\}$.

\emph{The set $\tilde{R}_k+\mathfrak{X}_k$ is not contained in $\tilde{R}_k$.}
One has
\[
\frac{\mathrm{td1}_k}{\mathrm{t1}_k}=e_k^{-1+2\varepsilon+\frac{3}{4}-\varepsilon}=e_k^{-\left(\frac{1}{4}-\varepsilon\right)}.
\]
As $\varepsilon\le\frac{1}{200}$, one has $\frac{1}{4}-\varepsilon\ge\varepsilon$;
as $e_k\le e_0<1$ and $e_0^{\varepsilon}\le\frac{1}{4}$ (A8), it follows that $e_k^{\frac{1}{4}-\varepsilon}\le e_k^{\varepsilon}\le\frac{1}{4}$.
Thus $\mathrm{td1}_k/2\ge2\,\mathrm{t1}_k>\mathrm{t1}_k$.
Let $\mathscr{C}$ be the constant bond function whose three components are equal to $i\,\mathrm{td1}_k/2$.
Then $|\mathscr{C}|_\infty=\mathrm{td1}_k/2<\mathrm{td1}_k$ and $\partial\mathscr{C}=0$, so that $\mathscr{C}\in\mathfrak{X}_k$;
and $0\in\tilde{R}_k$.
If $0+\mathscr{C}$ belonged to $\tilde{R}_k$, that is, $\mathscr{C}=A_0+A_1$ with $A_0$ real and $A_1\in R_k$, then
\[
\frac{\mathrm{td1}_k}{2}=|\operatorname{Im}\mathscr{C}|_\infty=|\operatorname{Im}A_1|_\infty\le|A_1|_\infty<\mathrm{t1}_k,
\]
a contradiction.
Hence the expression $E(X;\mathscr{A}+\mathscr{C})$ has no meaning at the point $(0,\mathscr{C})$ of $\tilde{R}_k\times\mathfrak{X}_k$.

\emph{Proof of (d).}
Let $\mathrm{Do}$ stand for $\tilde{R}_k$, respectively for $\tilde{R}_k\times\mathfrak{X}_k$.
The complex bond functions on $T_k$ form a finite-dimensional complex vector space, and $\mathrm{Do}$ is an open subset of this space, respectively of its square \cite[Lemma~17.3]{QEDI}.
The algebra $\mathrm{Gr}[\psi,\chi]$ is finite-dimensional, and $\|\cdot\|_{h_k}$ is a norm on it.

\emph{Step 1.}
Let $\mathrm{Ba}(\mathrm{Do})$ be the complex vector space of the bounded analytic maps $f:\mathrm{Do}\to\mathrm{Gr}[\psi,\chi]$, with the norm $|f|_{\mathrm{Do}}:=\sup_{\mathrm{pt}\in\mathrm{Do}}\|f(\mathrm{pt})\|_{h_k}$.
This space is complete.
Indeed, let $(f_j)$ be a Cauchy sequence, and put $\delta_j:=\sup_{i\ge j}|f_i-f_j|_{\mathrm{Do}}$, so that $\delta_j\to0$.
For every $\mathrm{pt}\in\mathrm{Do}$ the sequence $(f_j(\mathrm{pt}))$ is a Cauchy sequence in a finite-dimensional space;
let $f(\mathrm{pt})$ be its limit.
Then
\[
\|f(\mathrm{pt})-f_j(\mathrm{pt})\|_{h_k}=\lim_{i\to\infty}\|f_i(\mathrm{pt})-f_j(\mathrm{pt})\|_{h_k}\le\delta_j.
\]
Hence $f$ is bounded, $|f-f_j|_{\mathrm{Do}}\le\delta_j\to0$, and $f$ is continuous, being a uniform limit of continuous maps.
The map $f$ is analytic.
To see this, choose linear coordinates $\mathrm{za}=(\mathrm{za}_1,\dots,\mathrm{za}_{\mathrm{nd}})$ on the vector space, $\mathrm{nd}$ being its dimension, a point $\mathrm{pt}^0$ of $\mathrm{Do}$ with the coordinates $\mathrm{za}^0$, and $\mathrm{rad}>0$ such that the closed polydisk $\left\{|\mathrm{za}_i-\mathrm{za}^0_i|\le\mathrm{rad},\ i=1,\dots,\mathrm{nd}\right\}$ is contained in $\mathrm{Do}$;
maps on $\mathrm{Do}$ are written as functions of the coordinates.
Cauchy's integral formula (cl1), applied to the variables $\mathrm{za}_1,\dots,\mathrm{za}_{\mathrm{nd}}$ one after the other, gives for $\mathrm{za}$ in the open polydisk
\[
f_j(\mathrm{za})=(2\pi i)^{-\mathrm{nd}}\oint\mathrm{d}\mathrm{za}'_1\cdots\oint\mathrm{d}\mathrm{za}'_{\mathrm{nd}}\;f_j(\mathrm{za}')\prod_{i=1}^{\mathrm{nd}}\frac{1}{\mathrm{za}'_i-\mathrm{za}_i},
\]
where each $\mathrm{za}'_i$ ranges over the circle $|\mathrm{za}'_i-\mathrm{za}^0_i|=\mathrm{rad}$ and $\mathrm{d}\mathrm{za}'_i$ is the differential.
The convergence $f_j\to f$ is uniform on the set of integration; hence the same formula holds for $f$.
For $|\mathrm{za}_i-\mathrm{za}^0_i|\le\mathrm{rad}/2$ the series
\[
\frac{1}{\mathrm{za}'_i-\mathrm{za}_i}=\sum_{n\ge0}\left(\mathrm{za}_i-\mathrm{za}^0_i\right)^n\left(\mathrm{za}'_i-\mathrm{za}^0_i\right)^{-n-1}
\]
converge uniformly in $\mathrm{za}'$.
Integrating term by term, one finds that near $\mathrm{pt}^0$ the map $f$ is the sum of a convergent power series in $\mathrm{za}-\mathrm{za}^0$.

\emph{Step 2.}
The set $D_k$ is finite, since $T_k$ has $(\ell/M)^3$ $M$-cubes.
Let $\mathrm{Pb}(\mathrm{Do})$ be the space of the families $E=\{E(X;\,\cdot\,):\ X\in D_k\}$ with $E(X;\,\cdot\,)\in\mathrm{Ba}(\mathrm{Do})$, normed by
\[
\max_X\exp\left(\kappa\cdot d_M(X)\right)\cdot\left|E(X;\,\cdot\,)\right|_{\mathrm{Do}}.
\]
It is a finite product of Banach spaces, and its norm is equivalent to the maximum of the norms of the factors;
so it is a complex Banach space.
By (Ak1), the families with (k1)$^{\mathrm{o}}$ and (k6)$^{\mathrm{o}}$ (for $\mathrm{Do}=\tilde{R}_k$), respectively with (k1) and (k6) (for $\mathrm{Do}=\tilde{R}_k\times\mathfrak{X}_k$), are exactly the elements of $\mathrm{Pb}(\mathrm{Do})$, and the norm is $\|\cdot\|_k$, respectively $\|\cdot\|^{x}_{k}$.

\emph{Step 3.}
For $X\in D_k$ and $\mathrm{pt}\in\mathrm{Do}$ the evaluation $E\mapsto E(X;\mathrm{pt})$ is a complex linear map of $\mathrm{Pb}(\mathrm{Do})$ into $\mathrm{Gr}[\psi,\chi]$;
as $\kappa\ge0$, one has $\|E(X;\mathrm{pt})\|_{h_k}\le\exp\left(-\kappa\cdot d_M(X)\right)\cdot\|E\|\le\|E\|$, and the evaluation is continuous.
The maps $a_{\tilde{g}}$, $G_\Lambda$ and $\mathrm{Cg}$ are complex linear maps of the finite-dimensional space $\mathrm{Gr}[\psi,\chi]$ into itself, hence continuous.
In the following list, $\mathrm{pt}=\mathscr{A}$, $\tilde{g}\mathrm{pt}=\tilde{g}\mathscr{A}$ and $\Lambda=\boldsymbol{\lambda}$ if $\mathrm{Do}=\tilde{R}_k$, whereas $\mathrm{pt}=(\mathscr{A},\mathscr{C})$, $\tilde{g}\mathrm{pt}=(\tilde{g}\mathscr{A},\tilde{g}\mathscr{C})$ and $\Lambda=\boldsymbol{\lambda}+\boldsymbol{\lambda}'$ if $\mathrm{Do}=\tilde{R}_k\times\mathfrak{X}_k$.
The space $K_{k}^{C}$, respectively $K^{x,C}_{k}$, is the set of the $E\in\mathrm{Pb}(\mathrm{Do})$ with the following properties.
\begin{itemize}
\item[(i)] For all $X$ and $\mathrm{pt}$, the element $E(X;\mathrm{pt})$ lies in the linear span of the balanced basis monomials.
(A balanced monomial has as many barred as unbarred factors, hence an even number of factors; so (i) includes evenness.)
\item[(ii)] For all $X$ and $\mathrm{pt}$, the element $E(X;\mathrm{pt})$ lies in the linear span of the basis monomials which contain only generators $\psi(\xi)$ with $\xi$ over $X$ and $\chi(\boldsymbol{\zeta})$ with $\boldsymbol{\zeta}$ admissible for $X$.
(By (Ak1) this is the first clause of (k2)$^{\mathrm{o}}$, respectively of (k2).)
\item[(iii)] $E(X;\mathrm{pt})-E(X;\tilde{\mathrm{pt}})=0$ for all $X$ and all $\mathrm{pt},\tilde{\mathrm{pt}}\in\mathrm{Do}$ which coincide, in both components, on the bonds that meet $X$.
\item[(iv)] $a_{\tilde{g}}\left[E(\tilde{g}X;\tilde{g}\mathrm{pt})\right]-E(X;\mathrm{pt})=0$ for all $X$, $\mathrm{pt}$ and $\tilde{g}\in\mathbb{G}_k$.
(Here $a_{\tilde{g}}$ is formed with any representative $a\in M\mathbb{Z}^3$ of the translation part; on the even elements, hence under (i), it does not depend on the representative, Section~I.4.3.)
\item[(v)] $G_\Lambda\left[E(X;\mathrm{pt}^1)\right]-E(X;\mathrm{pt})=0$ for all $X$ and all $(\mathrm{pt},\mathrm{pt}^1,\boldsymbol{\lambda})$, respectively $(\mathrm{pt},\mathrm{pt}^1,\boldsymbol{\lambda},\boldsymbol{\lambda}')$, as in (k4)$^{\mathrm{o}}$, respectively (k4).
\item[(vi)] $\mathrm{Cg}\left[E(X;-\mathrm{pt})\right]-E(X;\mathrm{pt})=0$ for all $X$ and $\mathrm{pt}$.
\end{itemize}
Each of these conditions demands that a continuous complex linear map of $\mathrm{Pb}(\mathrm{Do})$ into $\mathrm{Gr}[\psi,\chi]$ take its values in a linear subspace of this finite-dimensional space, and such a subspace is closed.
Hence each condition defines a closed complex linear subspace of $\mathrm{Pb}(\mathrm{Do})$, and $K_{k}^{C}$, respectively $K^{x,C}_{k}$, is the intersection of these subspaces:
a closed complex linear subspace of a Banach space, and therefore a complex Banach space.

\emph{Step 4.}
The map $\mathbf{K}$ is additive, $\mathbf{K}(tF)=\overline{t}\,\mathbf{K}F$ for complex numbers $t$, and $\|\mathbf{K}F\|_{h_k}=\|F\|_{h_k}$ (Section~I.4.5).
Therefore, for all $X$ and $\mathrm{pt}$, the map $E\mapsto\mathbf{K}\left[E(X;-\overline{\mathrm{pt}})\right]-E(X;\mathrm{pt})$ is continuous and real linear.
The space $K_k$, respectively $K^{x}_{k}$, is the set of the elements of $K_{k}^{C}$, respectively $K^{x,C}_{k}$, at which all these maps vanish:
a closed real linear subspace, hence a real Banach space.
\end{proof}

\begin{named}{Remark}
Assertions (a), (b) and (d) hold for all values of the constants that enter the definitions of the polymers, of the pairs, of the group $\mathbb{G}_k$, of $h_k$, of the domains and of the norms $\|\cdot\|_k$ and $\|\cdot\|^{x}_{k}$, subject in (d) to $\kappa\ge0$;
assertion (c) holds under the remaining conditions listed at the beginning of the proof.
Each of the four assertions holds for every $k\ge0$, every $\ell\in\mathrm{Cl}$ and every $e_k\in(0,e_0]$.
\end{named}

\subsection*{2. Proposition~48.2}

In this section $K_X(x,x'):=E_{20}(X;0;x,x')$, in the matrix notation of Section~I.5.1, for $x,x'\in\mathbb{L}^{\mathrm{f}}_k$;
this kernel has two arguments and is not the ball $K_x$ of Section~I.1.4.
As in Lemma~C, $\mathrm{Ql}_X(\mathrm{Mx})=\int_X\mathrm{d}x\,\bar{\psi}(x)^{T}\mathrm{Mx}\,\psi(x)$.
Thus, for small $X$,
\begin{gather*}
(R E)(X;\mathscr{A})=E(X;\mathscr{A})-\mathrm{al}_{0}(E,X)\cdot\operatorname{Vol}(X)-\mathrm{Ql}_X\left(\mathrm{al}_{2}(E,X)\right),
\\
\operatorname{Vol}(X)\cdot\mathrm{al}_{2}(E,X)=\int_{\tilde{X}}\mathrm{d}x\int_{\tilde{X}}\mathrm{d}x'\,K_X(x,x'),
\end{gather*}
where the integrals are those of (0.a), each of which contains the factor $\frac{1}{48}$ and the sum over the label.
Of the conditions on the constants, only $\kappa\ge\kappa_0$ (the line (C2) of Section~I.8.2), $c_{\mathrm{fl}}=61\,c_X\cdot K_0$ (the line (C0)), $M\ge1$ and $1/L\le1$ are needed.

\begin{proof}[Proof of Proposition~48.2\,(o)]
\emph{Lift.}
The definition of $\mathrm{al}_{0}$ involves no integration over a lift.
Two lifts of $X$ differ by a vector $\ell n$ with $n\in\mathbb{Z}^3$ (Dl2).
By (I.AP),
\[
K_X(x+\ell n,x'+\ell n)=(-1)^{2(n_1+n_2+n_3)}K_X(x,x')=K_X(x,x');
\]
the map $x\mapsto x+\ell n$ is a bijection of $\tilde{X}\times O_{\mathrm{cub}}$ onto $(\tilde{X}+\ell n)\times O_{\mathrm{cub}}$, and all weights are equal to $\eta_k^3/48$.
Therefore
\[
\int_{\tilde{X}+\ell n}\mathrm{d}x\int_{\tilde{X}+\ell n}\mathrm{d}x'\,K_X(x,x')=\int_{\tilde{X}}\mathrm{d}x\int_{\tilde{X}}\mathrm{d}x'\,K_X(x,x').
\]

\emph{Preparation.}
The maps $a_{\tilde{g}}$, $\mathrm{Cg}$ and $\mathbf{K}$ preserve the bidegree (Section~I.4.5);
$a_{\tilde{g}}$ and $\mathrm{Cg}$ are algebra automorphisms which fix the numbers, and $\mathbf{K}$ maps a number to its complex conjugate.
Since $E$ is balanced, (Ak3) gives $E^{[2,0]}(X;0)=\mathrm{Bq}[K_X]$.

\emph{If (k3)$^{\mathrm{o}}$ holds at $0$.}
Since $\tilde{g}0=0$, one has $a_{\tilde{g}}\left[E(\tilde{g}X;0)\right]=E(X;0)$.
In bidegree $(0,0)$ this reads $E_{00}(\tilde{g}X;0)=E_{00}(X;0)$.
In bidegree $(2,0)$ it reads $a_{\tilde{g}}\mathrm{Bq}[K_{\tilde{g}X}]=\mathrm{Bq}[K_X]$, and (Bm1) together with (Ak4) gives
\begin{equation}
S^{-1}K_{\tilde{g}X}(\tilde{g}x,\tilde{g}x')\,S=K_X(x,x')\qquad\text{for all }x,x'\in\mathbb{L}^{\mathrm{f}}_k.
\tag{2.a}
\end{equation}
The polymer $\tilde{g}X$ is small, $|\tilde{g}X|_M=|X|_M$, and $\tilde{g}\tilde{X}$ is a lift of $\tilde{g}X$ (Proposition~16.1\,(a1)).
The map $x=(c,g)\mapsto\tilde{g}x=(rc+a,rg)$ is a bijection of $\tilde{X}\times O_{\mathrm{cub}}$ onto $\tilde{g}\tilde{X}\times O_{\mathrm{cub}}$:
by Proposition~16.1\,(a2), $\pi(\tilde{g}x)=\tilde{g}\pi(x)$, the cluster of $c$ is mapped onto the cluster of $\tilde{g}c$, and the $48$ labels are permuted by $g\mapsto rg$.
All weights are equal.
With (2.a) it follows that
\begin{align*}
\operatorname{Vol}(\tilde{g}X)\cdot\mathrm{al}_{2}(E,\tilde{g}X)
&=\int_{\tilde{g}\tilde{X}}\mathrm{d}x\int_{\tilde{g}\tilde{X}}\mathrm{d}x'\,K_{\tilde{g}X}(x,x')
=\int_{\tilde{X}}\mathrm{d}x\int_{\tilde{X}}\mathrm{d}x'\,K_{\tilde{g}X}(\tilde{g}x,\tilde{g}x')\\
&=S\left[\int_{\tilde{X}}\mathrm{d}x\int_{\tilde{X}}\mathrm{d}x'\,K_X(x,x')\right]S^{-1}
=\operatorname{Vol}(X)\cdot S\,\mathrm{al}_{2}(E,X)\,S^{-1}.
\end{align*}
It remains to divide by $\operatorname{Vol}(\tilde{g}X)=\operatorname{Vol}(X)$.
Moreover,
\[
\mathrm{al}_{0}(E,\tilde{g}X)=\operatorname{Vol}(\tilde{g}X)^{-1}\cdot E_{00}(\tilde{g}X;0)=\operatorname{Vol}(X)^{-1}\cdot E_{00}(X;0)=\mathrm{al}_{0}(E,X).
\]

\emph{If (k5)$^{\mathrm{o}}$ holds at $0$.}
Since $-0=0$, one has $\mathrm{Cg}\left[E(X;0)\right]=E(X;0)$, hence $\mathrm{Cg}\,\mathrm{Bq}[K_X]=\mathrm{Bq}[K_X]$, and (Bm3) together with (Ak4) gives $K_X(x,x')=\mathsf{C}^{-1}K_X(x',x)^{T}\mathsf{C}$.
Integrate in $x$ and in $x'$ over $\tilde{X}$ (the double integral does not change when the names $x$ and $x'$ are exchanged), and divide by $\operatorname{Vol}(X)$:
$\mathrm{al}_{2}(E,X)=\mathsf{C}^{-1}\mathrm{al}_{2}(E,X)^{T}\mathsf{C}$.

\emph{If (k7)$^{\mathrm{o}}$ holds at $0$.}
Since $-\overline{0}=0$, one has $\mathbf{K}\left[E(X;0)\right]=E(X;0)$.
In bidegree $(0,0)$ this gives $\overline{E_{00}(X;0)}=E_{00}(X;0)$;
as $\operatorname{Vol}(X)$ is real, $\mathrm{al}_{0}(E,X)$ is real.
In bidegree $(2,0)$, (Bm4) together with (Ak4) gives $K_X(x,x')=\varkappa^{-1}\,\overline{K_X(x,x')}\,\varkappa$.
The weights are real;
integration yields $\mathrm{al}_{2}(E,X)=\varkappa^{-1}\,\overline{\mathrm{al}_{2}(E,X)}\,\varkappa$.

(Clause (o) is also Proposition~16.2\,(o) and (i)--(iii), applied to the family $\mathscr{E}(X;0,0):=E(X;0)$, which is defined at $(0,0)$.)
\end{proof}

\begin{proof}[Proof of Proposition~48.2\,(a)]
Every $E\in K_{k}^{C}$ has (k3)$^{\mathrm{o}}$, the clause (k2)$^{\mathrm{o}}$ gives the condition on $E_{20}$ that (o) requires, and $0\in\tilde{R}_k$.
Since $D_k$ is finite, all sums are finite.

\emph{Independence of $\square$.}
Let $a\in M\mathbb{Z}^3$ and $\tilde{g}=(1,a)$.
The map $X\mapsto\tilde{g}X=X+a$ is a bijection of $\{X\ \text{small}:\ X\supseteq\square\}$ onto $\{X\ \text{small}:\ X\supseteq\square+a\}$, with the inverse $X\mapsto X-a$ (Proposition~16.1\,(a1)).
By (o) with $S=1$, $\mathrm{al}_{0}(E,X+a)=\mathrm{al}_{0}(E,X)$ and $\mathrm{al}_{2}(E,X+a)=\mathrm{al}_{2}(E,X)$.
Hence $\varepsilon(E,\square+a)=\varepsilon(E,\square)$ and $m(E,\square+a)=m(E,\square)$.
The $M$-cubes are centred at the points of $M\mathbb{Z}^3/\ell\mathbb{Z}^3$;
so any two of them differ by a translation of this kind.

\emph{The matrix $m(E,\square)$ is a multiple of $1$.}
Let $z\in M\mathbb{Z}^3$ be the centre of $\square$, let $S\in\mathrm{Sp}$ and $r=r(S)$, and let $\tilde{g}_z=\left(S,(1-r)z\right)$ be the element of Section~I.4.2.
As $(1-r)z\in2M\mathbb{Z}^3$, the element $\tilde{g}_z$ lies in $\mathbb{G}_k$.
Since $\tilde{g}_zc=z+r(c-z)$ and $|\tilde{g}_zc-z|=|c-z|$, the element $\tilde{g}_z$ maps $\square$ onto itself, and hence maps the set $\{X\ \text{small}:\ X\supseteq\square\}$ bijectively onto itself (Proposition~16.1\,(a1)).
By (o),
\[
m(E,\square)=-\sum_X\mathrm{al}_{2}(E,\tilde{g}_zX)=-\sum_XS\,\mathrm{al}_{2}(E,X)\,S^{-1}=S\,m(E,\square)\,S^{-1}.
\]
For $S=\gamma_1,\gamma_2,\gamma_3,\gamma_5$ this says that $m(E,\square)$ commutes with these four matrices.
By (I.Sch), $m(E,\square)=m({E})\cdot1$ with a complex number $m({E})$.
(On the sites, $\tilde{g}_z$ acts by $(c,g)\mapsto\left(z+r(c-z),rg\right)$;
the labels are moved, and $\mathrm{al}_{2}$ is a mean over both labels.)

\emph{Linearity.}
The map $E\mapsto E(X;0)\mapsto\left(E_{00}(X;0),K_X\right)$ is complex linear, the kernels being coordinates (Ak1);
hence so are $\mathrm{al}_{0}(\,\cdot\,,X)$, $\mathrm{al}_{2}(\,\cdot\,,X)$, $\varepsilon$ and $m$.

\emph{Reality.}
Let $E\in K_k$.
By (o), $\mathrm{al}_{0}(E,X)$ is real, so that $\varepsilon({E})$ is real;
and $m({E})\cdot1=\varkappa^{-1}\,\overline{m({E})\cdot1}\,\varkappa=\overline{m({E})}\cdot1$, so that $m({E})$ is real.

\emph{With slot.}
Let $\mathscr{E}\in K^{x,C}_{k}$.
Then $\mathscr{E}|_0\in K_{k}^{C}$ (Proposition~48.1\,(a)) and $\mathscr{E}|_0(X;0)=\mathscr{E}(X;0,0)$.
With the definitions with slot of Section~I.5.6, which are those of the data of Propositions~16.2 and~16.3, one obtains
\[
\mathrm{al}_{0}(\mathscr{E},X)=\operatorname{Vol}(X)^{-1}\cdot\mathscr{E}_{00}(X;0,0)=\mathrm{al}_{0}(\mathscr{E}|_0,X),\qquad\mathrm{al}_{2}(\mathscr{E},X)=\mathrm{al}_{2}(\mathscr{E}|_0,X).
\]
Hence $\varepsilon(\mathscr{E},\square)=\varepsilon(\mathscr{E}|_0,\square)=\varepsilon(\mathscr{E}|_0)$ and $m(\mathscr{E},\square)=m(\mathscr{E}|_0,\square)=m(\mathscr{E}|_0)\cdot1$ for every $\square$;
that is, $\varepsilon(\mathscr{E})=\varepsilon(\mathscr{E}|_0)$ and $m(\mathscr{E})=m(\mathscr{E}|_0)$.
\end{proof}

\begin{proof}[Proof of Proposition~48.2\,(b)]
Recall that $0\in\tilde{R}_k$.
By (Ak1), $\operatorname{Vol}(X)\cdot|\mathrm{al}_{0}(E,X)|=|E_{00}(X;0)|\le\|E(X;0)\|_{h_k}$.
By (Ak3) with $\mathsf{w}=h_k$, and since $h_{k,1}^{-2}=e_k^{\frac{1}{2}}$,
\begin{equation}
\begin{split}
\operatorname{Vol}(X)\sum_{\beta,\beta'}\left|\mathrm{al}_{2}(E,X)_{\beta\beta'}\right|
&\le\sum_{\beta,\beta'}\int_{\tilde{X}}\mathrm{d}x\int_{\tilde{X}}\mathrm{d}x'\,\left|K_X(x,x')_{\beta\beta'}\right|\\
&=h_{k,1}^{-2}\cdot\left\|E^{[2,0]}(X;0)\right\|_{h_k}=e_k^{\frac{1}{2}}\cdot\left\|E^{[2,0]}(X;0)\right\|_{h_k};
\end{split}
\tag{2.b}
\end{equation}
here $K_X$ vanishes unless $x$ and $x'$ lie over $X$, so that the integrals over $\tilde{X}$ are the integrals of (Ak3).
This proves the last two bounds of (b).
Further, $\operatorname{Vol}(X)\ge M^3$;
by (Ak1) and (k6)$^{\mathrm{o}}$, $\left\|E^{[2,0]}(X;0)\right\|_{h_k}\le\|E(X;0)\|_{h_k}\le\|E\|_k\cdot\exp\left(-\kappa\cdot d_M(X)\right)$;
and by $\kappa\ge\kappa_0$ and \cite[Theorem~13.2\,(t2)]{QEDI},
\[
\sum_{X\supseteq\square}\exp\left(-\kappa\cdot d_M(X)\right)\le\sum_{X\supseteq\square}\exp\left(-\kappa_0\cdot d_M(X)\right)\le K_0.
\]
Consequently
\begin{gather*}
|\varepsilon({E})|\le\sum_{X\ \text{small},\ X\supseteq\square}|\mathrm{al}_{0}(E,X)|\le M^{-3}\cdot\|E\|_k\cdot K_0,
\\
|m({E})|=\left|m(E,\square)_{11}\right|\le\sum_{X\ \text{small},\ X\supseteq\square}\left|\mathrm{al}_{2}(E,X)_{11}\right|\le M^{-3}\cdot e_k^{\frac{1}{2}}\cdot\|E\|_k\cdot K_0.
\end{gather*}
Finally, $K_0\cdot M^{-3}\le K_0\le61\,c_X\cdot K_0=c_{\mathrm{fl}}$.
\end{proof}

\begin{proof}[Proof of Proposition~48.2\,(c)]
By the definition of $R$,
\[
\sum_XE(X;\mathscr{A})-\sum_X(R E)(X;\mathscr{A})=\sum_{X\ \text{small}}\left(\mathrm{al}_{0}(E,X)\cdot\operatorname{Vol}(X)+\mathrm{Ql}_X\left(\mathrm{al}_{2}(E,X)\right)\right).
\]
The sets $\square\times O_{\mathrm{cub}}$ form a partition of $T^{\mathrm{f}}_k$ (Section~I.1.4);
hence $\operatorname{Vol}(X)=\sum_{\square\subseteq X}M^3$ and $\mathrm{Ql}_X(\mathrm{Mx})=\sum_{\square\subseteq X}\mathrm{Ql}_{\square}(\mathrm{Mx})$.
Exchanging the sums over $X$ and over $\square$ and using (a), one obtains
\begin{align*}
&\sum_{\square}\left(M^3\sum_{X\ \text{small},\ X\supseteq\square}\mathrm{al}_{0}(E,X)+\mathrm{Ql}_{\square}\left(\sum_{X\ \text{small},\ X\supseteq\square}\mathrm{al}_{2}(E,X)\right)\right)\\
&\qquad=\sum_{\square}\left(-\varepsilon({E})\cdot M^3-m({E})\cdot\mathrm{Ql}_{\square}(1)\right)
=-\varepsilon({E})\cdot\ell^3-m({E})\cdot\langle\bar{\psi},\psi\rangle,
\end{align*}
because $T_k$ has $(\ell/M)^3$ $M$-cubes and $\sum_{\square}\mathrm{Ql}_{\square}(1)=\int\mathrm{d}x\,\bar{\psi}(x)\cdot\psi(x)$.
Moreover $n_k=\ell^3$.
\end{proof}

\begin{proof}[Proof of Proposition~48.2\,(d)]
\emph{Linearity.}
By (a), $\mathrm{al}_{0}$ and $\mathrm{al}_{2}$ are linear in $E$, and $\mathrm{Ql}_X(\mathrm{Mx})$ is linear in $\mathrm{Mx}$.

\emph{Bound.}
Let $X$ be small.
The two subtracted terms have the bidegrees $(0,0)$ and $(2,0)$.
The norm of the first one is $|E_{00}(X;0)|=\left\|E^{[0,0]}(X;0)\right\|_{h_k}$;
by (Cq1) and (2.b), the norm of the second one is $h_{k,1}^2\cdot\operatorname{Vol}(X)\sum_{\beta,\beta'}\left|\mathrm{al}_{2}(E,X)_{\beta\beta'}\right|\le\left\|E^{[2,0]}(X;0)\right\|_{h_k}$.
By (Ak1) their sum is at most $\|E(X;0)\|_{h_k}$.
Hence
\[
\left\|(R E)(X;\mathscr{A})\right\|_{h_k}\le\|E(X;\mathscr{A})\|_{h_k}+\|E(X;0)\|_{h_k};
\]
for large $X$ there is nothing to prove.
Each of the two terms on the right is at most $\|E\|_k\cdot\exp\left(-\kappa\cdot d_M(X)\right)$.
Therefore $\|R E\|_k\le2\|E\|_k$, and $2\le c_{\mathrm{fl}}$.

\emph{Normalisation.}
Let $X$ be small.
Then $(R E)_{00}(X;0)=E_{00}(X;0)-\mathrm{al}_{0}(E,X)\cdot\operatorname{Vol}(X)=0$.
By (Cq1) and (Ak4), for $x$ and $x'$ over $\tilde{X}$,
\[
(R E)_{20}(X;0;x,x')=K_X(x,x')-48\,\eta_k^{-3}\cdot\mathbf{1}[x=x']\cdot\mathrm{al}_{2}(E,X),
\]
and the double integral over $\tilde{X}$ is $\operatorname{Vol}(X)\cdot\mathrm{al}_{2}(E,X)-\operatorname{Vol}(X)\cdot\mathrm{al}_{2}(E,X)=0$.

\emph{The map $R$ preserves (k2)$^{\mathrm{o}}$}, for every family with (k2)$^{\mathrm{o}}$:
the subtracted terms do not depend on $\mathscr{A}$, and by (Cq1) they contain only generators $\psi(\xi)$ with $\xi$ over $X$.

\emph{The map $R$ sends $K_{k}^{C}$ into itself.}
The subtracted terms are even and balanced (Cq1), and they are constant in $\mathscr{A}$;
this gives (k1)$^{\mathrm{o}}$.
Clause (k2)$^{\mathrm{o}}$ has been shown above, and (k6)$^{\mathrm{o}}$ follows from the bound.
A polymer $X$ is small if and only if $\tilde{g}X$ is small (Proposition~16.1\,(a1)), and for large $X$ the clauses for $R E$ are those for $E$.
Let $X$ be small.
\begin{itemize}
\item[(k3)$^{\mathrm{o}}$] By (k3)$^{\mathrm{o}}$ for $E$, by (o) and by (Cq2),
\begin{align*}
a_{\tilde{g}}\left[(R E)(\tilde{g}X;\tilde{g}\mathscr{A})\right]
&=a_{\tilde{g}}\left[E(\tilde{g}X;\tilde{g}\mathscr{A})\right]-\mathrm{al}_{0}(E,\tilde{g}X)\cdot\operatorname{Vol}(\tilde{g}X)-a_{\tilde{g}}\mathrm{Ql}_{\tilde{g}X}\left(S\,\mathrm{al}_{2}(E,X)\,S^{-1}\right)\\
&=E(X;\mathscr{A})-\mathrm{al}_{0}(E,X)\cdot\operatorname{Vol}(X)-\mathrm{Ql}_X\left(\mathrm{al}_{2}(E,X)\right).
\end{align*}
\item[(k4)$^{\mathrm{o}}$] The automorphism $G_{\boldsymbol{\lambda}}$ fixes the numbers and $\mathrm{Ql}_X(\mathrm{Mx})$ (Cq2);
hence
\[
G_{\boldsymbol{\lambda}}\left[(R E)(X;\mathscr{A}^1)\right]=G_{\boldsymbol{\lambda}}\left[E(X;\mathscr{A}^1)\right]-\mathrm{al}_{0}(E,X)\cdot\operatorname{Vol}(X)-\mathrm{Ql}_X\left(\mathrm{al}_{2}(E,X)\right)=(R E)(X;\mathscr{A}).
\]
\item[(k5)$^{\mathrm{o}}$] The automorphism $\mathrm{Cg}$ fixes the numbers, and by (Cq2) and (o)
\[
\mathrm{Cg}\,\mathrm{Ql}_X\left(\mathrm{al}_{2}(E,X)\right)=\mathrm{Ql}_X\left(\mathsf{C}^{-1}\mathrm{al}_{2}(E,X)^{T}\mathsf{C}\right)=\mathrm{Ql}_X\left(\mathrm{al}_{2}(E,X)\right);
\]
hence
\[
\mathrm{Cg}\left[(R E)(X;-\mathscr{A})\right]=\mathrm{Cg}\left[E(X;-\mathscr{A})\right]-\mathrm{al}_{0}(E,X)\cdot\operatorname{Vol}(X)-\mathrm{Ql}_X\left(\mathrm{al}_{2}(E,X)\right)=(R E)(X;\mathscr{A}).
\]
\end{itemize}

\emph{The map $R$ sends $K_k$ into itself.}
For (k7)$^{\mathrm{o}}$ observe that $\mathbf{K}\left[\mathrm{al}_{0}(E,X)\cdot\operatorname{Vol}(X)\right]=\overline{\mathrm{al}_{0}(E,X)}\cdot\operatorname{Vol}(X)=\mathrm{al}_{0}(E,X)\cdot\operatorname{Vol}(X)$, and that by (Cq2) and (o)
\[
\mathbf{K}\,\mathrm{Ql}_X\left(\mathrm{al}_{2}(E,X)\right)=\mathrm{Ql}_X\left(\varkappa^{-1}\,\overline{\mathrm{al}_{2}(E,X)}\,\varkappa\right)=\mathrm{Ql}_X\left(\mathrm{al}_{2}(E,X)\right);
\]
hence
\[
\mathbf{K}\left[(R E)(X;-\overline{\mathscr{A}})\right]=\mathbf{K}\left[E(X;-\overline{\mathscr{A}})\right]-\mathrm{al}_{0}(E,X)\cdot\operatorname{Vol}(X)-\mathrm{Ql}_X\left(\mathrm{al}_{2}(E,X)\right)=(R E)(X;\mathscr{A}).
\]
\end{proof}

\begin{proof}[Proof of Proposition~48.2\,(e)]
The setting is (G2.S-2).
For $\mathscr{A}\in\tilde{R}_{k+1}$ one has $\mathscr{A}_L\in\tilde{R}_k$ (Proposition~18.1\,(c)), and the sum over $X$ is finite;
so $L E'$ is a family without slot of level $k+1$.
The set $\tilde{R}_{k+1}$ is invariant under $\mathscr{A}\mapsto\tilde{g}'\mathscr{A}$, $\mathscr{A}\mapsto-\mathscr{A}$ and $\mathscr{A}\mapsto\overline{\mathscr{A}}$ (Proposition~16.1\,(c) at the level $k+1$), and $\tilde{g}'$ maps $D_{k+1}$ onto itself (Proposition~16.1\,(a1) at the level $k+1$).
\begin{itemize}
\item[(k3)$^{\mathrm{o}}$] Let $\tilde{g}'=(S,a')\in\mathbb{G}_{k+1}$ and $\tilde{g}:=(S,La')\in\mathbb{G}_k^{L}$ (Proposition~16.1\,(d1)).
For $X\in D_k$ one has $\overline{\tilde{g}X}^{L}=\tilde{g}\left(\bar{X}^{L}\right)$ (Proposition~16.1\,(a1)) and $\tilde{g}(LY)=L(\tilde{g}'Y)$ (Proposition~16.1\,(d2)).
Hence $X\mapsto\tilde{g}X$ is a bijection of $\{X:\ \bar{X}^{L}=LY\}$ onto $\{X:\ \bar{X}^{L}=L\tilde{g}'Y\}$.
With $(\tilde{g}'\mathscr{A})_L=\tilde{g}(\mathscr{A}_L)$ (Proposition~16.1\,(d3)), with $a_{\tilde{g}'}S_L=S_La_{\tilde{g}}$ (Proposition~16.1\,(d4)) and with (k3)$^{\mathrm{o}}$ for $\tilde{g}\in\mathbb{G}_k^{L}$, one obtains
\begin{align*}
a_{\tilde{g}'}\left[(L E')(\tilde{g}'Y;\tilde{g}'\mathscr{A})\right]
&=\sum_{X:\ \bar{X}^{L}=LY}a_{\tilde{g}'}S_L\left[E'\left(\tilde{g}X;\tilde{g}(\mathscr{A}_L)\right)\right]
=\sum_XS_La_{\tilde{g}}\left[E'\left(\tilde{g}X;\tilde{g}(\mathscr{A}_L)\right)\right]\\
&=\sum_XS_L\left[E'(X;\mathscr{A}_L)\right]=(L E')(Y;\mathscr{A}).
\end{align*}
(The elements of $E'$ need not be even.
In that case $a_{\tilde{g}}$ depends on the representative in $LM\mathbb{Z}^3$ of the translation part of $\tilde{g}$, and $a_{\tilde{g}'}$ on the representative in $M\mathbb{Z}^3$ of the translation part of $\tilde{g}'$ (Section~I.4.3).
In Proposition~16.1\,(d4) the two representatives are $a$ and $a/L$.
The map $a'\mapsto La'$ carries the representatives of the translation part of $\tilde{g}'$ onto those of the translation part of $\tilde{g}$, because $L\left(a'+(\ell/L)n\right)=La'+\ell n$.
Clause (k3)$^{\mathrm{o}}$ is asserted for every vector of the class (Section~I.4.3);
so the computation holds for every representative $a'$.)
\item[(k4)$^{\mathrm{o}}$] Let $\mathscr{A},\mathscr{A}^1\in\tilde{R}_{k+1}$, let $\boldsymbol{\lambda}$ be a real function on the sites of $T_{k+1}$, and let $\mathscr{A}^1=\mathscr{A}-d\boldsymbol{\lambda}$ on the bonds at distance $\le1$ from $Y$.
Let $\bar{X}^{L}=LY$;
then $X\subseteq LY$.
Let $b$ be a bond of $T_k$ at distance $\le1$ from $X$.
Both ends of $b$ have sup-distance $\le1$ from $LY$;
so both ends of $b/L$ have sup-distance $\le1/L\le1$ from $Y$, and $\mathscr{A}^1(b/L)=\mathscr{A}(b/L)-(d\boldsymbol{\lambda})(b/L)$.
Multiplication by $L^{-\frac{1}{2}}$ gives
\[
\mathscr{A}^1_L(b)=\mathscr{A}_L(b)-(d\boldsymbol{\lambda})_L(b)=\mathscr{A}_L(b)-\left(d(\boldsymbol{\lambda}_L)\right)(b)
\]
(Proposition~16.1\,(d4)).
The function $\boldsymbol{\lambda}_L$ is real, and $\mathscr{A}_L$, $\mathscr{A}^1_L$ lie in $\tilde{R}_k$.
By Proposition~16.1\,(d4) and by (k4)$^{\mathrm{o}}$ for $E'$,
\[
G_{\boldsymbol{\lambda}}S_L\left[E'(X;\mathscr{A}^1_L)\right]=S_LG_{\boldsymbol{\lambda}_L}\left[E'(X;\mathscr{A}^1_L)\right]=S_L\left[E'(X;\mathscr{A}_L)\right].
\]
It remains to sum over $X$.
\item[(k5)$^{\mathrm{o}}$, (k7)$^{\mathrm{o}}$] One has $(-\mathscr{A})_L=-\mathscr{A}_L$ and $(\overline{\mathscr{A}})_L=\overline{\mathscr{A}_L}$, as $L^{-\frac{1}{2}}$ is real;
moreover $\mathrm{Cg}\,S_L=S_L\mathrm{Cg}$ and $\mathbf{K}S_L=S_L\mathbf{K}$ (Proposition~16.1\,(d4)).
Therefore
\begin{align*}
\mathrm{Cg}\left[(L E')(Y;-\mathscr{A})\right]&=\sum_XS_L\mathrm{Cg}\left[E'(X;-\mathscr{A}_L)\right]=(L E')(Y;\mathscr{A}),\\
\mathbf{K}\left[(L E')(Y;-\overline{\mathscr{A}})\right]&=\sum_XS_L\mathbf{K}\left[E'\left(X;-\overline{\mathscr{A}_L}\right)\right]=(L E')(Y;\mathscr{A}).
\end{align*}
\end{itemize}
\end{proof}

In clause (f) the class of sides is $\mathrm{Cl}$, the level $k$ and the number $e_k$ are fixed, $\mathrm{pr}=\mathrm{pr}_\ell:\mathbb{L}_k\to T_k(\ell)$ is the covering map, and $T_k(\ell)$ is as in Proposition~16.3\,(f3).
The clause consists of (f1), which is the assertion of Proposition~16.3\,(f1) itself, and of (f2)$^{\mathrm{o}}$ and (f3)$^{\mathrm{o}}$, which are Proposition~16.3\,(f2) and (f3) with the replacements named in Proposition~48.2\,(f).
The symbols $\mathrm{Sm}(\tilde{\square})$, $K$, $F$, $\ell_i$ and $\mathrm{pr}_i$ have the meaning that they have in these clauses, and the two families of (f3)$^{\mathrm{o}}$ are written $E^{(i)}$.

\begin{proof}[Proof of Proposition~48.2\,(f1)]
The definition of $\mathrm{Sm}(\tilde{\square})$ involves $\mathbb{L}_k$, $M$, $L$ and $\tilde{\square}$, but not $\ell$.
A translation by a vector $v\in M\mathbb{Z}^3$ maps $M$-cubes of $\mathbb{L}_k$ to $M$-cubes, common faces to common faces, and nets in $\mathbb{R}^3$ to nets of the same length;
hence $\mathrm{Sm}(\tilde{\square}+v)=\mathrm{Sm}(\tilde{\square})+v$.

Let $\tilde{X}\in\mathrm{Sm}(\tilde{\square})$.
Then $2M\left(d_M(\tilde{X})+2\right)\le2M(L+2)<\ell$, and by the converse part of \cite[Theorem~13.2\,(t10)]{QEDI} the map $\mathrm{pr}$ is injective on $\tilde{X}$, $\mathrm{pr}(\tilde{X})$ is a polymer, and $d_M(\mathrm{pr}\tilde{X})=d_M(\tilde{X})\le L$.
Thus $\mathrm{pr}(\tilde{X})$ is a small polymer which contains $\mathrm{pr}(\tilde{\square})$.
By \cite[Theorem~13.2\,(t7)]{QEDI}, in its form for unions of $M$-cubes of $\mathbb{R}^3$, any two points of $\tilde{X}$ have sup-distance at most $M\left(d_M(\tilde{X})+2\right)\le M(L+2)$, and $M(L+2)+\eta_k<\ell$.
Hence the set of the sites of $\tilde{X}$, which $\mathrm{pr}$ maps onto the set of the sites of $\mathrm{pr}(\tilde{X})$, is a lift of $\mathrm{pr}(\tilde{X})$ in the sense of Section~I.1.1.
By (Dl2), every lift of the small polymer $\mathrm{pr}(\tilde{X})$ is a translate of this lift by a vector of $\ell\mathbb{Z}^3$;
so $\tilde{X}$ is the lift of $\mathrm{pr}(\tilde{X})$, up to $\ell\mathbb{Z}^3$.

\emph{Injectivity.}
Let $\tilde{X},\tilde{X}'\in\mathrm{Sm}(\tilde{\square})$ have the same image $X$.
Both are lifts of $X$;
so $\tilde{X}'=\tilde{X}+\ell n$ with some $n\in\mathbb{Z}^3$ (Dl2).
Then $\tilde{\square}$ and $\tilde{\square}-\ell n$ both lie in $\tilde{X}$ and have the same image.
As $\mathrm{pr}$ is injective on $\tilde{X}$, it follows that $n=0$.

\emph{Surjectivity.}
Let $X$ be a small polymer which contains $\mathrm{pr}(\tilde{\square})$.
By (Dl2), $X$ has a lift $\tilde{X}$, and $\tilde{X}$ satisfies the hypothesis of Lemma~D.
By (Dl1) and (Dl2), $\tilde{X}$ is a non-empty finite union of $M$-cubes, connected through faces, with $d_M(\tilde{X})\le L$.
Its cube over $\mathrm{pr}(\tilde{\square})$ is $\tilde{\square}+\ell n$ for some $n\in\mathbb{Z}^3$.
The set $\tilde{X}-\ell n$ belongs to $\mathrm{Sm}(\tilde{\square})$, and its image is $X$.
\end{proof}

\begin{proof}[Proof of (f2)$^{\mathrm{o}}$]
\emph{The maps.}
Let $\tilde{X}$ be a non-empty finite union of $M$-cubes of $\mathbb{L}_k$ which is connected through faces, and let $\tilde{U}$ be its $\eta_k$-neighbourhood.
Let $K=(K_{nm})$ be a system of antisymmetric kernels whose arguments lie over $\tilde{U}$, and let $F\in\mathrm{Grc}_k(\tilde{U})$.
The integrals are those of (0.a), $\operatorname{Vol}(\tilde{X})$ is $M^3$ times the number of the cubes of $\tilde{X}$, and $K_{20}(x,x')$ is the $4\times4$ matrix with the entries
\[
K_{20}(x,x')_{\beta\beta'}=K_{20}\left((x,\beta,1),(x',\beta',0)\right).
\]
Define
\begin{gather*}
\Phi_{\mathrm{al}_0}(\tilde{X};K):=\operatorname{Vol}(\tilde{X})^{-1}K_{00},
\qquad
\Phi_{\mathrm{al}_2}(\tilde{X};K):=\operatorname{Vol}(\tilde{X})^{-1}\int_{\tilde{X}}\mathrm{d}x\int_{\tilde{X}}\mathrm{d}x'\,K_{20}(x,x'),
\\
\Phi_R(\tilde{X};K,F):=
\begin{cases}
F-K_{00}-\displaystyle\int_{\tilde{X}}\mathrm{d}x\,\bar{\psi}(x)^T\cdot\Phi_{\mathrm{al}_2}(\tilde{X};K)\cdot\psi(x) & \text{if } d_M(\tilde{X})\le L,\\
F & \text{otherwise};
\end{cases}
\end{gather*}
thus $\Phi_R(\tilde{X};K,F)$ is an element of $\mathrm{Grc}_k(\tilde{U})$.
These maps are built from $\tilde{X}$, $\eta_k$, $M$, $L$ and the number $48$ alone; the side $\ell$ does not enter them.

\emph{Translations.}
Let $v\in M\mathbb{Z}^3$; this lattice contains $LM\mathbb{Z}^3$, and it is contained in $\mathbb{Z}^3$, since $M$ is an integer.
Put $(T_vK)(\xi;\boldsymbol{\zeta}):=K(\xi-v;\boldsymbol{\zeta}-v)$.
The algebra $\mathrm{Grc}_k(\tilde{U})$ of Section~I.1.9 has four families of independent generators, $\psi(\xi)$, $\chi(\boldsymbol{\zeta})$, $\Psi(u)$ and $W(u)$;
here $u=(y,\beta,\omega)$ with $y\in\tilde{U}\cap\mathbb{Z}^3$ if $k\ge1$, and $u=(x,\beta,\omega)$ with $x$ over $\tilde{U}$ if $k=0$.
Put $u+v:=(y+v,\beta,\omega)$, resp.\ $u+v:=(x+v,\beta,\omega)$, and let $T_v\colon\mathrm{Grc}_k(\tilde{U})\to\mathrm{Grc}_k(\tilde{U}+v)$ be the relabelling of all four families,
\[
T_v\psi(\xi):=\psi(\xi+v),\qquad T_v\chi(\boldsymbol{\zeta}):=\chi(\boldsymbol{\zeta}+v),\qquad T_v\Psi(u):=\Psi(u+v),\qquad T_vW(u):=W(u+v).
\]
Since $v\in\mathbb{Z}^3$, for $k\ge1$ the point $y+v$ lies in $(\tilde{U}+v)\cap\mathbb{Z}^3$;
so the generators of $\mathrm{Grc}_k(\tilde{U})$ are mapped bijectively onto those of $\mathrm{Grc}_k(\tilde{U}+v)$, and $T_v$ is an isomorphism.
On numbers and on constant matrices $T_v$ acts as the identity (Section~I.1.9).
One has $\operatorname{Vol}(\tilde{X}+v)=\operatorname{Vol}(\tilde{X})$ and $d_M(\tilde{X}+v)=d_M(\tilde{X})$;
the $\eta_k$-neighbourhood of $\tilde{X}+v$ is $\tilde{U}+v$;
and $x\mapsto x+v=(c+v,g)$ is a bijection of $\tilde{X}\times O_{\mathrm{cub}}$ onto $(\tilde{X}+v)\times O_{\mathrm{cub}}$ which preserves the weights.
Consequently
\begin{align*}
\Phi_{\mathrm{al}_0}(\tilde{X}+v;T_vK)&=\Phi_{\mathrm{al}_0}(\tilde{X};K),\\
\Phi_{\mathrm{al}_2}(\tilde{X}+v;T_vK)&=\operatorname{Vol}(\tilde{X})^{-1}\int_{\tilde{X}+v}\mathrm{d}x\int_{\tilde{X}+v}\mathrm{d}x'\,K_{20}(x-v,x'-v)=\Phi_{\mathrm{al}_2}(\tilde{X};K),\\
T_v\int_{\tilde{X}}\mathrm{d}x\,\bar{\psi}(x)^T\cdot\mathrm{Mx}\cdot\psi(x)&=\int_{\tilde{X}+v}\mathrm{d}x\,\bar{\psi}(x)^T\cdot\mathrm{Mx}\cdot\psi(x)
\end{align*}
for every constant $4\times4$ matrix $\mathrm{Mx}$, and therefore
\[
\Phi_R(\tilde{X}+v;T_vK,T_vF)=T_v\Phi_R(\tilde{X};K,F).
\]

\emph{Identification.}
Let $\ell\in\mathrm{Cl}$ and $E\in K_{k}^{C}(\ell)$; of the properties of $E$ only (k2)$^{\mathrm{o}}$ is needed in what follows.
Let $X$ be a polymer of $T_k(\ell)$ which has a lift $\tilde{X}$ with
\[
\operatorname{sup\text{-}diameter}(\tilde{X})+2\eta_k+M<\ell,
\]
which is the condition of Section~I.1.9 for $\mathrm{rg}=\eta_k$, and let $\mathscr{A}\in\tilde{R}_k$.
As data take the kernels $K$ of $E(X;0)$ and the element $F:=\tilde{F}$ of (Dl4) for $E(X;\mathscr{A})$, that is, the element of $\mathrm{Grc}_k(\tilde{U})$ with the kernels of $E(X;\mathscr{A})$;
in both cases the antiperiodic kernels are restricted to the arguments over $\tilde{U}$ (Section~I.1.9).
If $X$ is small, then $\operatorname{Vol}(X)=\operatorname{Vol}(\tilde{X})$ by (Dl1);
hence $\mathrm{al}_{0}(E,X)=\Phi_{\mathrm{al}_0}(\tilde{X};K)$, and $\mathrm{al}_{2}(E,X)=\Phi_{\mathrm{al}_2}(\tilde{X};K)$ by the definition of $\mathrm{al}_{2}$ and Proposition~48.2\,(o).
By (Dl4), $\mathrm{cov}_\ell F=E(X;\mathscr{A})$.
The homomorphism $\mathrm{cov}_\ell$ leaves the numbers fixed and maps $\psi(x,\beta,\omega)$, $x$ over $\tilde{X}$, to the generator of the torus with the same index;
since $\tilde{X}$ is a set of representatives of the sites of $X$, it follows that
\[
\mathrm{cov}_\ell\left(\int_{\tilde{X}}\mathrm{d}x\,\bar{\psi}(x)^T\cdot\mathrm{Mx}\cdot\psi(x)\right)=\mathrm{Ql}_X(\mathrm{Mx}).
\]
By (Dl2), $d_M(\tilde{X})\le L$ holds if and only if $X$ is small.
For small $X$ the three terms of $\Phi_R(\tilde{X};K,F)$ are therefore mapped by $\mathrm{cov}_\ell$ to $E(X;\mathscr{A})$, $\mathrm{al}_{0}(E,X)\cdot\operatorname{Vol}(X)$ and $\mathrm{Ql}_X\left(\mathrm{al}_{2}(E,X)\right)$;
for large $X$ one has $\Phi_R(\tilde{X};K,F)=F$.
In both cases
\[
\mathrm{cov}_\ell\,\Phi_R(\tilde{X};K,F)=(RE)(X;\mathscr{A}).
\]
\end{proof}

\begin{proof}[Proof of (f3)$^{\mathrm{o}}$]
Let $\tilde{X}\in\mathrm{Sm}(\tilde{\square})$ and $X_i:=\mathrm{pr}_i\tilde{X}$.
By (f1), $X_i$ is a small polymer of $T_k(\ell_i)$ with the lift $\tilde{X}$, and $\operatorname{Vol}(X_i)=\operatorname{Vol}(\tilde{X})$ by (Dl1).
Hence
\begin{align*}
\mathrm{al}_{0}(E^{(i)},X_i)&=\operatorname{Vol}(\tilde{X})^{-1}E^{(i)}_{00}(X_i;0),\\
\mathrm{al}_{2}(E^{(i)},X_i)&=\operatorname{Vol}(\tilde{X})^{-1}\int_{\tilde{X}}\mathrm{d}x\int_{\tilde{X}}\mathrm{d}x'\,E^{(i)}_{20}(X_i;0;x,x'),
\end{align*}
where the integrals are those of (0.a), which do not contain the side.
By the hypothesis of (f3)$^{\mathrm{o}}$ the right-hand sides are the same for $i=1$ and for $i=2$.
By (f1),
\begin{align*}
\varepsilon(E^{(i)},\mathrm{pr}_i\tilde{\square})&=-\sum_{\tilde{X}\in\mathrm{Sm}(\tilde{\square})}\mathrm{al}_{0}(E^{(i)},\mathrm{pr}_i\tilde{X}),\\
m(E^{(i)},\mathrm{pr}_i\tilde{\square})&=-\sum_{\tilde{X}\in\mathrm{Sm}(\tilde{\square})}\mathrm{al}_{2}(E^{(i)},\mathrm{pr}_i\tilde{X}),
\end{align*}
and for $i=1,2$ these are sums over one and the same index set with equal terms.
\end{proof}

\begin{named}{Remark}[The case $k=0$; dependence on the parameters]
For $k=0$ one has $\int\mathrm{d}x=\frac{1}{48}\sum_c\sum_g$, the kernel of $\mathrm{Ql}_X(\mathrm{Mx})$ is $48\cdot\mathbf{1}[x=x']\,\mathrm{Mx}$, and $\tilde{U}$ is the $1$-neighbourhood of $\tilde{X}$.
The constants $K_{0}$, $c_{\mathrm{fl}}$ and $2$ of Proposition~48.2 are absolute numbers, and $M^{-3}\le1$.
The bounds of Proposition~48.2 hold for all $k\ge0$, $\ell\in\mathrm{Cl}$ and $e_k\in(0,e_0]$, without any further condition;
the charge $e_k$ enters the bounds of (b) only through $h_{k,1}^{-2}=e_k^{\frac{1}{2}}$.
The parameters of Proposition~48.2 are $k$, $\ell$ and $e_k$, and in (e) moreover $\ell/L$ and $e_{k+1}$; no other parameter enters.
By (f), the numbers $\mathrm{al}_{0}$, $\mathrm{al}_{2}$, $\varepsilon$ and $m$ are given by formulas which do not contain the side.
\end{named}

\subsection*{3. Proposition~48.3}

\begin{proof}[Proof of Proposition~48.3]
\emph{Preparation.}
Let $X$ be small and let $\tilde{X}$ be its lift (Lemma~D\,(Dl2); \cite[Theorem~13.2\,(t7), (t10)]{QEDI}).
All points $x$, $\tilde{x}$, $p_i$ below are points of $\mathbb{L}^{\mathrm{f}}_k$ over $\tilde{X}$.
The connection $\mathscr{A}\in\tilde{R}_k$ is a complex bond function on $T_k$;
so Proposition~17.1 applies with $B=\mathscr{A}$, and $u_{\mathscr{A}}(p,p')=U_{\mathscr{A}}(p,p')$, $\mathrm{tm}$, $\mathrm{rm}$, $\mathrm{tm}^{-}$, $\mathrm{rm}^{-}$, $U_i$, $U^{-}_i$ and $\boldsymbol{\Theta}=\boldsymbol{\Theta}_{X,\mathscr{A}}$ are formed with $\mathscr{A}$.
Parts (b)--(f) of Proposition~17.1, in particular the relations (3.4), hold for every complex bond function $B$ on $T_k$, and they are applied here with $B=\mathscr{A}$.
In this proof (3.1)--(3.4) denote the displays, and (u2)--(u4), (c1)--(c3) the clauses, of Proposition~17.1.
By (k2)$^{\mathrm{o}}$ the kernel $E_{20}(X;\mathscr{A};\xi_1,\xi_2)$ vanishes unless $\xi_1$ and $\xi_2$ lie over $X$;
hence $E^{[2,0]}(X;\mathscr{A})$ lies in the subalgebra generated by the $\psi(\xi)$, $\xi$ over $X$, on which $\boldsymbol{\Theta}_{X,\mathscr{A}}$ is defined (Proposition~17.1\,(d)), and $E^{\mathrm{nat}}(X;\mathscr{A})$ is defined.
By the second bound of Proposition~17.1\,(a), the bound for $\mathscr{A}\in\tilde{R}_k$ alone,
\[
e_k\cdot|\operatorname{Im}\mathscr{A}|_\infty<e_k^{\frac{1}{4}+\varepsilon}\le\tfrac{1}{4}e_k^{\varepsilon}\le\tfrac{1}{2}e_k^{\varepsilon},
\]
and therefore (c2) yields
\begin{equation}
|U_i|\le\mathrm{e},\qquad|U^{-}_i|\le\mathrm{e}\qquad\text{for all chains and all }i.
\tag{3.a}
\end{equation}
Throughout (a)--(d) the connection $\mathscr{A}$ is fixed.

\paragraph{(a)}
Fix $x$ and $\tilde{x}$ over $\tilde{X}$, a chain $(\square_0,\dots,\square_n)\in\mathrm{Ch}(x,\tilde{x})$ with the chain points $p_0=x,\dots,p_{n+2}=\tilde{x}$, and a spinor index, which is not written.
For $0\le i\le n+2$ one has
\begin{equation}
\mathrm{ev}\,\psi(p_0)=U_i\cdot\mathrm{ev}\,\psi(p_i)-\sum_{j<i,\ p_j\ne p_{j+1}}U_j\cdot d_X(p_j,p_{j+1})^{\alpha}\cdot\mathrm{ev}\,\chi(p_j,p_{j+1}).
\tag{3.b}
\end{equation}
For $i=0$ this is true, because $U_0=1$ is an empty product and the sum is empty.
Assume (3.b) for some $i\le n+1$.
If $p_i=p_{i+1}$, then $u_{\mathscr{A}}(p_i,p_i)=1$ and $U_{i+1}=U_i$, and no term is added.
If $p_i\ne p_{i+1}$, the pair $(p_i,p_{i+1})$ is of kind (I) for $X$ by (c1);
this includes the last pair $((z_n,g),\tilde{x})$ in the case where $\pi\tilde{x}=z_n$ and the label of $\tilde{x}$ differs from $g$: this pair lies inside the cluster of $z_n$, and $d_X=\eta_k$, $u_{\mathscr{A}}=1$ for it.
Hence the images of the generators satisfy the relations of (u4) at this pair, and (3.1) gives
\[
\mathrm{ev}\,\psi(p_i)=u_{\mathscr{A}}(p_i,p_{i+1})\cdot\mathrm{ev}\,\psi(p_{i+1})-d_X(p_i,p_{i+1})^{\alpha}\cdot\mathrm{ev}\,\chi(p_i,p_{i+1}).
\]
Multiplying by $U_i$ and using $U_i\cdot u_{\mathscr{A}}(p_i,p_{i+1})=U_{i+1}$, one obtains (3.b) for $i+1$.
This proves (3.b).
Now take $i=n+2$ and form the mean $\operatorname{Mn}$ over the finite, non-empty set $\mathrm{Ch}(x,\tilde{x})$.
By (3.2) and the linearity of $\mathrm{ev}$,
\[
\mathrm{ev}\,\psi(x)=\mathrm{tm}(x,\tilde{x})\cdot\mathrm{ev}\,\psi(\tilde{x})-\mathrm{ev}\left(\mathrm{rm}(x,\tilde{x})[\chi]\right).
\]
The left-hand side does not depend on $\tilde{x}$.
The operation $\operatorname{Vol}(X)^{-1}\int_X\mathrm{d}\tilde{x}$ is a mean, since $\int_X\mathrm{d}\tilde{x}=\operatorname{Vol}(X)$ (this integral contains $\frac{1}{48}$ times the sum over the label of $\tilde{x}$);
applying it, one obtains from (3.3)
\[
\mathrm{ev}\,\psi(x)=\mathrm{ev}\left(\boldsymbol{\Theta}\psi(x)\right).
\]

\emph{Barred generators.}
For $0\le i\le n+2$ one has
\begin{equation}
\mathrm{ev}\,\bar{\psi}(p_0)=U^{-}_i\cdot\mathrm{ev}\,\bar{\psi}(p_i)-\sum_{j<i,\ p_j\ne p_{j+1}}U^{-}_j\cdot d_X(p_j,p_{j+1})^{\alpha}\cdot\mathrm{ev}\,\bar{\chi}(p_j,p_{j+1}).
\tag{3.b$'$}
\end{equation}
For $i=0$ both sides equal $\mathrm{ev}\,\bar{\psi}(p_0)$.
If $p_i=p_{i+1}$, then $U^{-}_{i+1}=U^{-}_i$.
If $p_i\ne p_{i+1}$, the second relation of (3.1) gives
\[
\mathrm{ev}\,\bar{\psi}(p_i)=u_{\mathscr{A}}(p_i,p_{i+1})^{-1}\cdot\mathrm{ev}\,\bar{\psi}(p_{i+1})-d_X(p_i,p_{i+1})^{\alpha}\cdot\mathrm{ev}\,\bar{\chi}(p_i,p_{i+1});
\]
multiplication by $U^{-}_i$ and the identity $U^{-}_i\cdot u_{\mathscr{A}}(p_i,p_{i+1})^{-1}=U^{-}_{i+1}$ give the step from $i$ to $i+1$.
At $i=n+2$ the mean $\operatorname{Mn}$ yields, by (3.2),
\[
\mathrm{ev}\,\bar{\psi}(x)=\mathrm{tm}^{-}(x,\tilde{x})\cdot\mathrm{ev}\,\bar{\psi}(\tilde{x})-\mathrm{ev}\left(\mathrm{rm}^{-}(x,\tilde{x})[\bar{\chi}]\right),
\]
and the mean $\operatorname{Vol}(X)^{-1}\int_X\mathrm{d}\tilde{x}$ yields $\mathrm{ev}\,\bar{\psi}(x)=\mathrm{ev}\left(\boldsymbol{\Theta}\bar{\psi}(x)\right)$ by (3.3).
Thus the homomorphisms $\mathrm{ev}\circ\boldsymbol{\Theta}$ and $\mathrm{ev}$ agree on the generators $\psi(\xi)$, $\xi$ over $X$, and hence on the subalgebra generated by them, which contains $E^{[2,0]}(X;\mathscr{A})$.
Therefore
\begin{align*}
\mathrm{ev}\left(E^{\mathrm{nat}}(X;\mathscr{A})\right)&=\mathrm{ev}\left(E(X;\mathscr{A})\right)-\mathrm{ev}\left(E^{[2,0]}(X;\mathscr{A})\right)+\mathrm{ev}\left(\boldsymbol{\Theta}\left(E^{[2,0]}(X;\mathscr{A})\right)\right)\\
&=\mathrm{ev}\left(E(X;\mathscr{A})\right).
\end{align*}

\emph{The insertion.}
Let now $\mathrm{ev}$ be the homomorphism of $\mathrm{Gr}[\psi,\chi]$ into $\mathrm{Gr}[\psi]$ with $\mathrm{ev}\,\psi(\xi):=\psi(\xi)$ and, for $(x,x')\in\mathrm{Pzt}_k$,
\begin{align*}
\mathrm{ev}\,\chi(x,x',\beta,0)&:=d_X(x,x')^{-\alpha}\cdot\left(U_{\mathscr{A}}(x,x')\cdot\psi_\beta(x')-\psi_\beta(x)\right)=(\delta_{\alpha,\mathscr{A}}\psi)_\beta(x,x'),\\
\mathrm{ev}\,\chi(x,x',\beta,1)&:=d_X(x,x')^{-\alpha}\cdot\left(U_{\mathscr{A}}(x,x')^{-1}\cdot\bar{\psi}_\beta(x')-\bar{\psi}_\beta(x)\right)=(\delta_{\alpha,-\mathscr{A}}\bar{\psi})_\beta(x,x')
\end{align*}
(see (I.1.7); $U_{-\mathscr{A}}=U_{\mathscr{A}}^{-1}$).
This homomorphism exists.
Indeed, the prescribed images are odd elements of $\mathrm{Gr}[\psi]$;
they change their sign under $(x,x')\mapsto(x+\ell e_\mu,x'+\ell e_\mu)$, as $\chi$ does by (I.AP);
and odd elements of a Grassmann algebra anticommute and have the square $0$.
So the assignment on the independent generators extends to exactly one homomorphism.
By the definition of $\mathrm{ev}$, the relations of (u4) with $B=\mathscr{A}$ hold at all pairs of $\mathrm{Pzt}_k$.
Hence $\mathrm{ev}\left(E^{\mathrm{nat}}(X;\mathscr{A})\right)=\mathrm{ev}\left(E(X;\mathscr{A})\right)$.
The same identity holds for $\mathrm{ev}'\circ\mathrm{ev}$, for every homomorphism $\mathrm{ev}'$ of $\mathrm{Gr}[\psi]$ into an associative algebra.

\paragraph{(b)}
Let $E$ be balanced.
By (Ak3) and (k2)$^{\mathrm{o}}$,
\[
E^{[2,0]}(X;\mathscr{A})=\int\mathrm{d}x\int\mathrm{d}x'\,\bar{\psi}(x)^T\cdot E_{20}(X;\mathscr{A};x,x')\cdot\psi(x'),
\]
the integrals extending over $X$.
Proposition~17.1\,(f), applied to the kernel $K(x,x')=E_{20}(X;\mathscr{A};x,x')$, writes $\boldsymbol{\Theta}\left(E^{[2,0]}(X;\mathscr{A})\right)$ as a sum of four terms.
Here $\mathrm{tm}$ and $\mathrm{tm}^{-}$ are numbers, whereas $\mathrm{rm}(x',\tilde{x})[\chi]$ and $\mathrm{rm}^{-}(x,\tilde{x}')[\bar{\chi}]$ are linear in the generators $\chi$ by (3.2).
The four terms therefore have the bidegrees $(2,0)$, $(1,1)$, $(1,1)$ and $(0,2)$.
It follows that $(E^{\mathrm{nat}})^{[n,m]}(X;\mathscr{A})=E^{[n,m]}(X;\mathscr{A})$ for $(n,m)\notin\{(2,0),(1,1),(0,2)\}$, that is, $E^{\mathrm{nat}}_{nm}=E_{nm}$ for these $(n,m)$ by (Ak1);
in particular $E^{\mathrm{nat}}_{00}=E_{00}$.
Moreover
\[
(E^{\mathrm{nat}})^{[2,0]}(X;\mathscr{A})=E^{[2,0]}(X;\mathscr{A})-E^{[2,0]}(X;\mathscr{A})+(\text{the first term})=\text{the first term}.
\]
For the connection $\mathscr{A}=0$, which lies in $\tilde{R}_k$, the first term is
\[
\bar{\psi}_{\mathrm{av}}^T\cdot\left[\int\mathrm{d}x\int\mathrm{d}x'\,E_{20}(X;0;x,x')\right]\cdot\psi_{\mathrm{av}}
\]
by the last sentence of Proposition~17.1\,(f).
If the matrix in the brackets vanishes, then $(E^{\mathrm{nat}})^{[2,0]}(X;0)=0$, and $E^{\mathrm{nat}}_{20}(X;0;\cdot,\cdot)=0$ by (Ak1).

\paragraph{(c)}
One has
\[
E^{\mathrm{nat}}(X;\mathscr{A})=\left(E(X;\mathscr{A})-E^{[2,0]}(X;\mathscr{A})\right)+\boldsymbol{\Theta}\left(E^{[2,0]}(X;\mathscr{A})\right).
\]
The monomials of the first bracket are monomials of $E(X;\mathscr{A})$;
by (k2)$^{\mathrm{o}}$ they contain only generators $\psi(\xi)$ with $\xi$ over $X$ and $\chi(\boldsymbol{\zeta})$ with $\boldsymbol{\zeta}$ admissible for $X$.
The element $\boldsymbol{\Theta}\left(E^{[2,0]}(X;\mathscr{A})\right)$ is a polynomial in the elements $\boldsymbol{\Theta}\psi(\xi)$, $\xi$ over $X$.
By (3.2) and (3.3), $\boldsymbol{\Theta}\psi(x)$, resp.\ $\boldsymbol{\Theta}\bar{\psi}(x)$, is a linear combination of the generators $\psi(\tilde{x})$, resp.\ $\bar{\psi}(\tilde{x})$, with $\tilde{x}$ over $X$, and of the generators $\chi(p_i,p_{i+1})$, resp.\ $\bar{\chi}(p_i,p_{i+1})$, at chain pairs with $p_i\ne p_{i+1}$.
By (c1) such a pair lies in $\mathrm{Pzt}_k$, so that $d_X(p_i,p_{i+1})\le M$, and it is admissible for $X$ of kind (I);
its two points carry the label $g$ of $x$, except for the last pair $((z_n,g),\tilde{x})$, whose labels are arbitrary.
Hence a generator $\chi(\boldsymbol{\zeta})$ which occurs in $E^{\mathrm{nat}}(X;\mathscr{A})$ but not in $E(X;\mathscr{A})$ sits at such a pair.
The coefficients of these linear combinations are made of the numbers $\operatorname{Vol}(X)^{-1}$ and $\eta_k^3/48$, the weights of the means $\operatorname{Mn}$, the numbers $d_X(p_i,p_{i+1})^{\alpha}$, and products of transports $u_{\mathscr{A}}(p_j,p_{j+1})^{\pm1}$.
By (u2), $u_{\mathscr{A}}(p_j,p_{j+1})$ depends on $\mathscr{A}$ only through its values on the bonds of the six paths $\boldsymbol{\Gamma}^{\pi}(\pi p_j,\pi p_{j+1})$;
these paths lie in $X$ (kind (I)), so that both ends of each of these bonds lie in $X$.
Consequently the coefficients of $E^{\mathrm{nat}}(X;\mathscr{A})$ depend on $\mathscr{A}$ only through the kernels $E_{nm}(X;\mathscr{A})$ and through the values of $\mathscr{A}$ on the bonds with both ends in $X$.

\emph{The property (k2)$^{\mathrm{o}}$ for $E^{\mathrm{nat}}$.}
Its first clause is the assertion on the generators just proved.
A bond with both ends in $X$ meets $X$.
If $\mathscr{A},\tilde{\mathscr{A}}\in\tilde{R}_k$ coincide on the bonds that meet $X$, then $E_{nm}(X;\mathscr{A})=E_{nm}(X;\tilde{\mathscr{A}})$ by (k2)$^{\mathrm{o}}$ for $E$, and all transports coincide;
hence $E^{\mathrm{nat}}(X;\mathscr{A})=E^{\mathrm{nat}}(X;\tilde{\mathscr{A}})$.
For large $X$ one has $E^{\mathrm{nat}}(X)=E(X)$.

\paragraph{(d)}
Let $E$ be balanced, and let $\mathsf{w}=(\mathsf{w}_1,\mathsf{w}_2)$ be arbitrary.
For $x$ over $\tilde{X}$ put
\begin{align*}
\mathrm{Tha}(x,\beta,0)&:=\operatorname{Vol}(X)^{-1}\int_X\mathrm{d}\tilde{x}\,\mathrm{tm}(x,\tilde{x})\cdot\psi_\beta(\tilde{x}),\\
\mathrm{Thb}(x,\beta,0)&:=-\operatorname{Vol}(X)^{-1}\int_X\mathrm{d}\tilde{x}\,\operatorname{Mn}\sum_{i=0}^{n+1}U_i\cdot d_X(p_i,p_{i+1})^{\alpha}\cdot\chi_\beta(p_i,p_{i+1}),\\
\mathrm{Tha}(x,\beta,1)&:=\operatorname{Vol}(X)^{-1}\int_X\mathrm{d}\tilde{x}\,\mathrm{tm}^{-}(x,\tilde{x})\cdot\bar{\psi}_\beta(\tilde{x}),\\
\mathrm{Thb}(x,\beta,1)&:=-\operatorname{Vol}(X)^{-1}\int_X\mathrm{d}\tilde{x}\,\operatorname{Mn}\sum_{i=0}^{n+1}U^{-}_i\cdot d_X(p_i,p_{i+1})^{\alpha}\cdot\bar{\chi}_\beta(p_i,p_{i+1});
\end{align*}
the terms with $p_i=p_{i+1}$ are omitted, and $\int_X\mathrm{d}\tilde{x}$ is the integral (0.a) over $\tilde{X}$, that is, $\eta_k^3/48$ times the sum over the photon sites $\tilde{c}\in\tilde{X}$ and over the $48$ labels.
By (3.2) and (3.3),
\[
\boldsymbol{\Theta}\psi(x,\beta,\omega)=\mathrm{Tha}(x,\beta,\omega)+\mathrm{Thb}(x,\beta,\omega),
\]
where $\mathrm{Tha}$ has the bidegree $(1,0)$ and $\mathrm{Thb}$ the bidegree $(0,1)$.
A generator $\psi$ has the norm $\mathsf{w}_1$, a generator $\chi$ the norm $\mathsf{w}_2$.
By (3.a), $|\mathrm{tm}(x,\tilde{x})|=|\operatorname{Mn}U_{n+2}|\le\mathrm{e}$ and $|\mathrm{tm}^{-}(x,\tilde{x})|=|\operatorname{Mn}U^{-}_{n+2}|\le\mathrm{e}$;
by (3.a) and (c3),
\[
\operatorname{Mn}\sum_i|U_i|\cdot d_X(p_i,p_{i+1})^{\alpha}\le\mathrm{e}\cdot3c_X\cdot L\cdot M,
\qquad
\operatorname{Mn}\sum_i|U^{-}_i|\cdot d_X(p_i,p_{i+1})^{\alpha}\le\mathrm{e}\cdot3c_X\cdot L\cdot M;
\]
and $\operatorname{Vol}(X)^{-1}\int_X\mathrm{d}\tilde{x}$ is a mean.
Hence
\begin{equation}
\|\mathrm{Tha}(x,\beta,\omega)\|_{\mathsf{w}}\le\mathrm{e}\cdot\mathsf{w}_1,
\qquad
\|\mathrm{Thb}(x,\beta,\omega)\|_{\mathsf{w}}\le\mathrm{e}\cdot(3c_X\cdot L\cdot M)\cdot\mathsf{w}_2.
\tag{3.c}
\end{equation}
The map $\boldsymbol{\Theta}$ is linear and multiplicative.
By (Ak3), with $E_{20}(x,x'):=E_{20}(X;\mathscr{A};x,x')$ and all integrals over $X$,
\begin{multline*}
\boldsymbol{\Theta}\left(E^{[2,0]}(X;\mathscr{A})\right)=\sum_{\beta,\beta'}\int\mathrm{d}x\int\mathrm{d}x'\,E_{20}(x,x')_{\beta\beta'}\cdot\left(\mathrm{Tha}(x,\beta,1)+\mathrm{Thb}(x,\beta,1)\right)\\
\cdot\left(\mathrm{Tha}(x',\beta',0)+\mathrm{Thb}(x',\beta',0)\right).
\end{multline*}
Its parts of the bidegrees $(2,0)$, $(1,1)$ and $(0,2)$ are, respectively, the terms with $\mathrm{Tha}\cdot\mathrm{Tha}$, with $\mathrm{Tha}\cdot\mathrm{Thb}+\mathrm{Thb}\cdot\mathrm{Tha}$, and with $\mathrm{Thb}\cdot\mathrm{Thb}$.
Put
\[
s_{20}:=\sum_{\beta,\beta'}\int\mathrm{d}x\int\mathrm{d}x'\,\left|E_{20}(x,x')_{\beta\beta'}\right|=\tfrac{1}{2}\|E_{20}(X)\|
\]
(Ak3).
By Proposition~15\,(b) and (3.c), the norms of the three parts are at most
\begin{equation}
\mathrm{e}^2\cdot\mathsf{w}_1^2\cdot s_{20},
\qquad
2\,\mathrm{e}^2\cdot(3c_X\cdot L\cdot M)\cdot\mathsf{w}_1\mathsf{w}_2\cdot s_{20},
\qquad
\mathrm{e}^2\cdot(3c_X\cdot L\cdot M)^2\cdot\mathsf{w}_2^2\cdot s_{20}.
\tag{3.d}
\end{equation}
By (Ak1) an element of bidegree $(n,m)$ with the antisymmetric kernel $K$ has the norm $\frac{\mathsf{w}_1^n\mathsf{w}_2^m}{n!\,m!}\|K\|$.
By (b),
\begin{gather*}
(E^{\mathrm{nat}})^{[2,0]}=\boldsymbol{\Theta}(E^{[2,0]})^{[2,0]},
\qquad
(E^{\mathrm{nat}})^{[1,1]}=E^{[1,1]}+\boldsymbol{\Theta}(E^{[2,0]})^{[1,1]},
\\
(E^{\mathrm{nat}})^{[0,2]}=E^{[0,2]}+\boldsymbol{\Theta}(E^{[2,0]})^{[0,2]}.
\end{gather*}
Therefore
\begin{align*}
\tfrac{1}{2}\mathsf{w}_1^2\cdot\|E^{\mathrm{nat}}_{20}(X)\|&\le\mathrm{e}^2\cdot\mathsf{w}_1^2\cdot\tfrac{1}{2}\|E_{20}(X)\|,\\
\mathsf{w}_1\mathsf{w}_2\cdot\|E^{\mathrm{nat}}_{11}(X)\|&\le\mathsf{w}_1\mathsf{w}_2\cdot\|E_{11}(X)\|+\mathrm{e}^2\cdot(3c_X\cdot L\cdot M)\cdot\mathsf{w}_1\mathsf{w}_2\cdot\|E_{20}(X)\|,\\
\tfrac{1}{2}\mathsf{w}_2^2\cdot\|E^{\mathrm{nat}}_{02}(X)\|&\le\tfrac{1}{2}\mathsf{w}_2^2\cdot\|E_{02}(X)\|+\mathrm{e}^2\cdot(3c_X\cdot L\cdot M)^2\cdot\mathsf{w}_2^2\cdot\tfrac{1}{2}\|E_{20}(X)\|.
\end{align*}
Division by $\frac{1}{2}\mathsf{w}_1^2$, by $\mathsf{w}_1\mathsf{w}_2$ and by $\frac{1}{2}\mathsf{w}_2^2$ gives the three bounds on the kernels in (d).

Now let $\mathsf{w}=h_k$, and put $\mathrm{ya}:=3c_X\cdot L\cdot M\cdot e_k^{\varepsilon}$.
One has $h_{k,2}=e_k^{\varepsilon}h_{k,1}$, and $h_{k,1}^2\cdot s_{20}=\|E^{[2,0]}(X;\mathscr{A})\|_{h_k}$ by (Ak3).
The sum of the three bounds (3.d) is therefore $\mathrm{e}^2\cdot h_{k,1}^2\cdot s_{20}\cdot\left(1+2\,\mathrm{ya}+\mathrm{ya}^2\right)$, so that
\[
\left\|\boldsymbol{\Theta}\left(E^{[2,0]}(X;\mathscr{A})\right)\right\|_{h_k}\le\mathrm{e}^2\cdot(1+\mathrm{ya})^2\cdot\|E^{[2,0]}(X;\mathscr{A})\|_{h_k}.
\]
The norm is additive over the bidegrees (Ak1):
$\|E-E^{[2,0]}\|_{h_k}=\|E\|_{h_k}-\|E^{[2,0]}\|_{h_k}$.
Since $\mathrm{e}^2\cdot(1+\mathrm{ya})^2\ge1$, it follows that
\begin{align*}
\|E^{\mathrm{nat}}(X;\mathscr{A})\|_{h_k}&\le\|E(X;\mathscr{A})\|_{h_k}-\|E^{[2,0]}(X;\mathscr{A})\|_{h_k}+\mathrm{e}^2\cdot(1+\mathrm{ya})^2\cdot\|E^{[2,0]}(X;\mathscr{A})\|_{h_k}\\
&\le\mathrm{e}^2\cdot(1+\mathrm{ya})^2\cdot\|E(X;\mathscr{A})\|_{h_k}.
\end{align*}
By the second condition of (A9), $90c_X\cdot L\cdot M\cdot e_k^{\varepsilon}\le90c_X\cdot L\cdot M\cdot e_0^{\varepsilon}\le1$, that is, $\mathrm{ya}\le\frac{1}{30}$.
Since $\mathrm{e}^2<7.3891$ and $\left(\frac{31}{30}\right)^2<1.0678$, one concludes that $\mathrm{e}^2\cdot(1+\mathrm{ya})^2<7.9\le10$.

\paragraph{(e)}
Let $E\in K_{k}^{C}$.
For large $X$ one has $E^{\mathrm{nat}}(X)=E(X)$, and $X$ is large if and only if $\tilde{g}X$ is large, for every $\tilde{g}\in\mathbb{G}_k$ (Proposition~16.1\,(a1));
so all the clauses hold for large $X$.
Let $X$ be small.

\emph{Even and balanced.}
The homomorphism $\boldsymbol{\Theta}$ maps unbarred generators to linear combinations of unbarred ones, and barred generators to linear combinations of barred ones (Proposition~17.1\,(d));
hence $\boldsymbol{\Theta}\left(E^{[2,0]}(X;\mathscr{A})\right)$ is a linear combination of products of one barred and one unbarred generator.
So $E^{\mathrm{nat}}(X;\mathscr{A})$, the sum of this element and of $E(X;\mathscr{A})-E^{[2,0]}(X;\mathscr{A})$, is even and balanced.

\emph{(k1)$^{\mathrm{o}}$.}
By the display for $\boldsymbol{\Theta}\left(E^{[2,0]}(X;\mathscr{A})\right)$ in (d), every coefficient of $E^{\mathrm{nat}}(X;\mathscr{A})$ in the basis of the monomials is a finite sum of products of a value of a kernel of $E(X;\mathscr{A})$, which is analytic on $\tilde{R}_k$, of fixed numbers, and of transports $u_{\mathscr{A}}(p,p')^{\pm1}$, which are entire functions of $\mathscr{A}$ (by (u2), and since $u_{\mathscr{A}}^{-1}=u_{-\mathscr{A}}$ by (u3)).
By (Ak1) the kernels of $E^{\mathrm{nat}}$ are analytic on $\tilde{R}_k$.

\emph{(k2)$^{\mathrm{o}}$.}
This is (c).

\emph{(k6)$^{\mathrm{o}}$.}
By (d), $\|E^{\mathrm{nat}}(X;\mathscr{A})\|_{h_k}\le10\,\|E(X;\mathscr{A})\|_{h_k}$;
hence $\|E^{\mathrm{nat}}\|_k\le10\,\|E\|_k$.

The four maps $a_{\tilde{g}}$, $G_{\boldsymbol{\lambda}}$, $\mathrm{Cg}$ and $\mathbf{K}$ preserve the bidegree (Section~I.4.5):
\[
(a_{\tilde{g}}F)^{[2,0]}=a_{\tilde{g}}(F^{[2,0]}),\quad
(G_{\boldsymbol{\lambda}}F)^{[2,0]}=G_{\boldsymbol{\lambda}}(F^{[2,0]}),\quad
(\mathrm{Cg}\,F)^{[2,0]}=\mathrm{Cg}(F^{[2,0]}),\quad
(\mathbf{K}F)^{[2,0]}=\mathbf{K}(F^{[2,0]})
\]
for $F\in\mathrm{Gr}[\psi,\chi]$.
The domain $\tilde{R}_k$ is invariant under $\mathscr{A}\mapsto\tilde{g}\mathscr{A}$, $\mathscr{A}\mapsto-\mathscr{A}$ and $\mathscr{A}\mapsto\overline{\mathscr{A}}$ (Proposition~16.1\,(c)).

\emph{(k3)$^{\mathrm{o}}$.}
The first relation of (3.4) with $B=\mathscr{A}$, namely $a_{\tilde{g}}\boldsymbol{\Theta}_{\tilde{g}X,\tilde{g}\mathscr{A}}=\boldsymbol{\Theta}_{X,\mathscr{A}}a_{\tilde{g}}$, is applied to $E^{[2,0]}(\tilde{g}X;\tilde{g}\mathscr{A})$, which lies in the subalgebra generated by the $\psi(\xi)$, $\xi$ over $\tilde{g}X$.
Together with (k3)$^{\mathrm{o}}$ for $E$ this gives
\begin{align*}
a_{\tilde{g}}\left[E^{\mathrm{nat}}(\tilde{g}X;\tilde{g}\mathscr{A})\right]
&=a_{\tilde{g}}\left[(E-E^{[2,0]})(\tilde{g}X;\tilde{g}\mathscr{A})\right]+\boldsymbol{\Theta}_{X,\mathscr{A}}\left(a_{\tilde{g}}\left[E^{[2,0]}(\tilde{g}X;\tilde{g}\mathscr{A})\right]\right)\\
&=(E-E^{[2,0]})(X;\mathscr{A})+\boldsymbol{\Theta}_{X,\mathscr{A}}\left(E^{[2,0]}(X;\mathscr{A})\right)=E^{\mathrm{nat}}(X;\mathscr{A}).
\end{align*}

\emph{(k4)$^{\mathrm{o}}$.}
Let $\mathscr{A},\mathscr{A}^1\in\tilde{R}_k$, and let $\boldsymbol{\lambda}$ be a real function of $c$ with $\mathscr{A}^1=\mathscr{A}-d\boldsymbol{\lambda}$ on the bonds at distance $\le1$ from $X$.
The bonds with both ends in $X$ are among these bonds.
The second relation of (3.4) with $B=\mathscr{A}$, $B^1=\mathscr{A}^1$ and $\Lambda=\boldsymbol{\lambda}$ reads $G_{\boldsymbol{\lambda}}\boldsymbol{\Theta}_{X,\mathscr{A}^1}=\boldsymbol{\Theta}_{X,\mathscr{A}}G_{\boldsymbol{\lambda}}$.
Together with (k4)$^{\mathrm{o}}$ for $E$ it gives
\begin{align*}
G_{\boldsymbol{\lambda}}\left[E^{\mathrm{nat}}(X;\mathscr{A}^1)\right]
&=G_{\boldsymbol{\lambda}}\left[(E-E^{[2,0]})(X;\mathscr{A}^1)\right]+\boldsymbol{\Theta}_{X,\mathscr{A}}\left(G_{\boldsymbol{\lambda}}\left[E^{[2,0]}(X;\mathscr{A}^1)\right]\right)\\
&=(E-E^{[2,0]})(X;\mathscr{A})+\boldsymbol{\Theta}_{X,\mathscr{A}}\left(E^{[2,0]}(X;\mathscr{A})\right)=E^{\mathrm{nat}}(X;\mathscr{A}).
\end{align*}

\emph{(k5)$^{\mathrm{o}}$.}
By the third relation of (3.4), $\mathrm{Cg}\,\boldsymbol{\Theta}_{X,-\mathscr{A}}=\boldsymbol{\Theta}_{X,\mathscr{A}}\mathrm{Cg}$, and by (k5)$^{\mathrm{o}}$ for $E$,
\[
\mathrm{Cg}\left[E^{\mathrm{nat}}(X;-\mathscr{A})\right]=\mathrm{Cg}\left[(E-E^{[2,0]})(X;-\mathscr{A})\right]+\boldsymbol{\Theta}_{X,\mathscr{A}}\left(\mathrm{Cg}\left[E^{[2,0]}(X;-\mathscr{A})\right]\right)=E^{\mathrm{nat}}(X;\mathscr{A}).
\]
Thus $E^{\mathrm{nat}}$ lies in $K_{k}^{C}$.

\emph{(k7)$^{\mathrm{o}}$.}
Let now $E\in K_k$.
By the fourth relation of (3.4), $\mathbf{K}\boldsymbol{\Theta}_{X,-\overline{\mathscr{A}}}=\boldsymbol{\Theta}_{X,\mathscr{A}}\mathbf{K}$, and by (k7)$^{\mathrm{o}}$ for $E$,
\[
\mathbf{K}\left[E^{\mathrm{nat}}(X;-\overline{\mathscr{A}})\right]=\mathbf{K}\left[(E-E^{[2,0]})(X;-\overline{\mathscr{A}})\right]+\boldsymbol{\Theta}_{X,\mathscr{A}}\left(\mathbf{K}\left[E^{[2,0]}(X;-\overline{\mathscr{A}})\right]\right)=E^{\mathrm{nat}}(X;\mathscr{A}).
\]
Thus $E^{\mathrm{nat}}$ lies in $K_k$.

\emph{Linearity.}
The map $F\mapsto F^{[2,0]}$ is linear, and $\boldsymbol{\Theta}_{X,\mathscr{A}}$ is linear and does not depend on $E$.

\emph{The family with slot.}
Let $\mathscr{E}\in K^{x,C}_{k}$.
Then $\mathscr{E}|_0$ lies in $K_{k}^{C}$ (Proposition~48.1\,(a)).
At $\mathscr{C}=0$ the total connection is $\mathscr{T}=\mathscr{A}+0=\mathscr{A}$.
By the definition of $\mathscr{E}^{\mathrm{nat}}$ (Section~I.5.6, which is the definition in the data of Proposition~17.2), one has for small $X$
\[
\mathscr{E}^{\mathrm{nat}}(X;\mathscr{A},0)=\mathscr{E}(X;\mathscr{A},0)-\mathscr{E}^{[2,0]}(X;\mathscr{A},0)+\boldsymbol{\Theta}_{X,\mathscr{A}}\left(\mathscr{E}^{[2,0]}(X;\mathscr{A},0)\right)=(\mathscr{E}|_0)^{\mathrm{nat}}(X;\mathscr{A}),
\]
while for large $X$ both sides equal $\mathscr{E}(X;\mathscr{A},0)$.
That is, $(\mathscr{E}^{\mathrm{nat}})|_0=(\mathscr{E}|_0)^{\mathrm{nat}}$.

\paragraph{(f)}
The class of sides is $\mathrm{Cl}$, and $k$ and $e_k$ are fixed.
The assertion to be proved is Proposition~17.2\,(f) with $\mathscr{A}$ in the places of $B$ and of $\mathscr{T}$, and with $E(X;\mathscr{A})$, $E^{\mathrm{nat}}(X;\mathscr{A})$ in the places of $E(X;\mathscr{A},\mathscr{C})$, $E^{\mathrm{nat}}(X;\mathscr{A},\mathscr{C})$.
Its three parts are denoted as follows.
\begin{itemize}
\item[(f-1)] The assertion on $\boldsymbol{\Theta}_{X,B}$, $X$ small:
in the sense of Section~I.1.9 it is given by one formula of range $\mathrm{rg}=0$, the data being the values of $B$ on the bonds with both ends in $\tilde{X}$;
the corresponding map assigns to $(\tilde{X},\tilde{B})$ the homomorphism (3.3) formed with the cubes of $\tilde{X}$, their centres and common faces, the labels, the measure (I.1.3) and the transports;
and this map commutes with the translations by the vectors of $M\mathbb{Z}^3$.
This assertion is proved for every complex bond function $B$ on $T_k$.
\item[(f-2)] The assertion on $E^{\mathrm{nat}}(X;\mathscr{A})$:
for every polymer $X$ of $T_k$, small or large, which has a lift $\tilde{X}$ with $\operatorname{sup\text{-}diameter}(\tilde{X})+2\eta_k+M<\ell$, it is given by one formula of range $\mathrm{rg}=\eta_k$, the data being the kernels of $E(X;\mathscr{A})$ and the values of $\mathscr{A}$ on the bonds with both ends in $\tilde{X}$;
and the map is
\[
F\mapsto F-F^{[2,0]}+\boldsymbol{\Theta}_{\tilde{X},\tilde{B}}\left(F^{[2,0]}\right)\quad\text{if }d_M(\tilde{X})\le L,\qquad F\mapsto F\quad\text{otherwise},
\]
where $F$ is the element with these kernels, $\tilde{B}$ is the lift of $\mathscr{A}$, and $\boldsymbol{\Theta}_{\tilde{X},\tilde{B}}$ is the homomorphism (3.3) for $(\tilde{X},\tilde{B})$.
\item[(f-3)] The consequence for two sides $\ell_1,\ell_2\in\mathrm{Cl}$:
let $\tilde{X}\in\mathrm{Sm}(\tilde{\square})$, where $\mathrm{Sm}(\tilde{\square})$ is as in Proposition~16.3\,(f1), which is the first clause of Proposition~48.2\,(f), and let $B_1$, $B_2$ be bond functions on $T_k(\ell_1)$, $T_k(\ell_2)$ whose lifts coincide on the bonds with both ends in $\tilde{X}$;
then, for every $x$ over $\tilde{X}$, the elements $\boldsymbol{\Theta}_{\mathrm{pr}_1\tilde{X},B_1}\psi(x,\cdot)$ and $\boldsymbol{\Theta}_{\mathrm{pr}_2\tilde{X},B_2}\psi(x,\cdot)$ are the same linear combination of the generators $\psi(\tilde{x},\cdot)$, $\chi(p,p',\cdot)$ with indices over $\tilde{X}$, and the same holds for $\bar{\psi}$;
here $\mathrm{pr}_i\colon\mathbb{L}_k\to T_k(\ell_i)$ is the covering map.
\end{itemize}

\emph{The map $\Phi_{\boldsymbol{\Theta}}$.}
Let $\tilde{X}$ be a non-empty finite union of $M$-cubes of $\mathbb{L}_k$ which is connected through faces, and let $\tilde{B}$ be a complex function on the bonds of $\mathbb{L}_k$ with both ends in $\tilde{X}$ such that $\tilde{B}(c',c)=-\tilde{B}(c,c')$.
For $x=(c,g)$ and $\tilde{x}$ over $\tilde{X}$, with the photon site $\tilde{c}=\pi\tilde{x}$, let $\widetilde{\mathrm{Ch}}(x,\tilde{x})$ be the set of the shortest chains $(\square_0,\dots,\square_n)$ of $M$-cubes of $\tilde{X}$ from the cube of $c$ to the cube of $\tilde{c}$ in which consecutive cubes have a common face;
this set is finite and not empty.
Let $z_i$ be the centre of $\square_i$, and put $p_0:=x$, $p_{i+1}:=(z_i,g)$ for $0\le i\le n$, and $p_{n+2}:=\tilde{x}$.
For $i=0$ and for $i=n+1$ the photon sites of $p_i$ and $p_{i+1}$ lie in one $M$-cube;
this cube is a product of three intervals of $\eta_k\mathbb{Z}$, so that it contains the six rectilinear paths between any two of its sites, and the two photon sites have a sup-distance $<M$.
For $1\le i\le n$ the photon sites of $p_i$ and $p_{i+1}$ are the centres of two cubes with a common face;
they differ in one coordinate, by $M$, and the six paths are the straight lattice segment between them.
This segment lies in the union of the two cubes because $M/\eta_k$ is odd (Section~I.1.1, Faces):
a site of the segment has the sup-distances $t$ and $M-t$ from the two centres, where $t\in\eta_k\mathbb{Z}$ and $0\le t\le M$;
hence $t\ne M/2$, and $t<M/2$ or $M-t<M/2$.
Consequently, if $p_i\ne p_{i+1}$, then the pair $(p_i,p_{i+1})$ lies in $\mathrm{Pzt}_k$, both its photon sites and its six paths lie in $\tilde{X}$, $\chi(p_i,p_{i+1},\cdot)$ is a generator of $\mathrm{Grc}_k(\tilde{X})$, and
\[
u_{\tilde{B}}(p_i,p_{i+1}):=\exp\left\{i\cdot e_k\cdot\eta_k\cdot(\tau\tilde{B})(\pi p_i,\pi p_{i+1})\right\}
\]
is defined, with $u_{\tilde{B}}(p_i,p_{i+1}):=1$ if $\pi p_i=\pi p_{i+1}$.
Let $\Phi_{\boldsymbol{\Theta}}(\tilde{X},\tilde{B})$ be the homomorphism of the subalgebra of $\mathrm{Grc}_k(\tilde{X})$ generated by the $\psi(x,\beta,\omega)$, $x$ over $\tilde{X}$, into $\mathrm{Grc}_k(\tilde{X})$ which is given on these generators by (3.2) and (3.3), with $\widetilde{\mathrm{Ch}}$, $u_{\tilde{B}}$, $\operatorname{Vol}(\tilde{X})$ and the integral $\int_{\tilde{X}}\mathrm{d}\tilde{x}$ of (0.a) in the places of $\mathrm{Ch}$, $u_B$, $\operatorname{Vol}(X)$ and $\int_X\mathrm{d}\tilde{x}$.
It exists, since the prescribed images are odd;
it is the homomorphism (3.3) for $(\tilde{X},\tilde{B})$, which Proposition~17.2\,(f) denotes by $\boldsymbol{\Theta}_{\tilde{X},\tilde{B}}$.
It is built from $\tilde{X}$, $\tilde{B}$, $\eta_k$, $e_k$, $\alpha$, $M$ and the number $48$;
the side $\ell$ does not enter it.

\emph{Translations.}
Let $v\in M\mathbb{Z}^3$, and let $(T_v\tilde{B})(b):=\tilde{B}(b-v)$ on the bonds with both ends in $\tilde{X}+v$.
The translation by $v$ maps the cubes of $\tilde{X}$, their centres and their common faces onto those of $\tilde{X}+v$;
it maps $\widetilde{\mathrm{Ch}}(x,\tilde{x})$ onto the set $\widetilde{\mathrm{Ch}}(x+v,\tilde{x}+v)$ formed for $\tilde{X}+v$, the chain points $p_i$ onto $p_i+v$ (the labels are not changed: $(c,g)+v=(c+v,g)$), and paths onto paths.
Hence $u_{T_v\tilde{B}}(p+v,p'+v)=u_{\tilde{B}}(p,p')$, $d_X(p+v,p'+v)=d_X(p,p')$ and $\operatorname{Vol}(\tilde{X}+v)=\operatorname{Vol}(\tilde{X})$, and $\tilde{x}\mapsto\tilde{x}+v$ preserves the weights in (0.a).
With the relabelling $T_v\colon\mathrm{Grc}_k(\tilde{X})\to\mathrm{Grc}_k(\tilde{X}+v)$ of the generators it follows that
\[
\Phi_{\boldsymbol{\Theta}}(\tilde{X}+v,T_v\tilde{B})\,T_v=T_v\,\Phi_{\boldsymbol{\Theta}}(\tilde{X},\tilde{B}).
\]

\emph{Identification; proof of (f-1).}
Let $\ell\in\mathrm{Cl}$, let $X$ be a small polymer of $T_k(\ell)$ with the lift $\tilde{X}$, let $B$ be a complex bond function on $T_k(\ell)$, and let $\tilde{B}:=B\circ\mathrm{pr}$ on the bonds with both ends in $\tilde{X}$.
By (Dl2), $\operatorname{sup\text{-}diameter}(\tilde{X})+M\le M\cdot(L+2)+M<\ell$, and $\tilde{X}$ satisfies the hypothesis of Lemma~D;
so $\mathrm{cov}_\ell$ is defined on $\mathrm{Grc}_k(\tilde{X})$ (Section~I.1.9, with $\mathrm{rg}=0$).
By (Dl1), $\mathrm{pr}$ maps the cubes of $\tilde{X}$ bijectively onto the cubes of $X$ and centres onto centres, and two cubes of $\tilde{X}$ have a common face if and only if their images have one;
therefore $\mathrm{pr}$ maps $\widetilde{\mathrm{Ch}}(x,\tilde{x})$ bijectively onto $\mathrm{Ch}(x,\tilde{x})$.
In Proposition~17.1 the chain points are chosen over $\tilde{X}$;
they are the points $p_i$ above.
The value $(\tau B)(\pi p_i,\pi p_{i+1})$ is formed on the lift, with the paths of $\mathbb{L}_k$ between the two points of $\tilde{X}$ (Section~I.1.1), and on the bonds of these paths $B$ is given by $\tilde{B}$;
hence $u_B(p_i,p_{i+1})=u_{\tilde{B}}(p_i,p_{i+1})$.
Moreover $\operatorname{Vol}(X)=\operatorname{Vol}(\tilde{X})$ and $\int_X\mathrm{d}\tilde{x}=\int_{\tilde{X}}\mathrm{d}\tilde{x}$.
The map $\mathrm{cov}_\ell$ sends each generator of $\mathrm{Grc}_k(\tilde{X})$ to the generator of the torus with the same index.
So, for $x$ over $\tilde{X}$,
\[
\boldsymbol{\Theta}_{X,B}\left(\mathrm{cov}_\ell\,\psi(x,\beta,\omega)\right)=\boldsymbol{\Theta}_{X,B}\,\psi(x,\beta,\omega)=\mathrm{cov}_\ell\left(\Phi_{\boldsymbol{\Theta}}(\tilde{X},\tilde{B})\,\psi(x,\beta,\omega)\right).
\]
Both $\mathrm{cov}_\ell\circ\Phi_{\boldsymbol{\Theta}}(\tilde{X},\tilde{B})$ and $\boldsymbol{\Theta}_{X,B}\circ\mathrm{cov}_\ell$ are homomorphisms;
hence they coincide on the subalgebra generated by the $\psi(x,\beta,\omega)$, $x$ over $\tilde{X}$.
This is the relation which Section~I.1.9 requires of a homomorphism.

\emph{Proof of (f-3).}
By the first clause of Proposition~48.2\,(f), $\tilde{X}$ is the lift of each of the small polymers $\mathrm{pr}_i\tilde{X}$.
The two functions $\tilde{B}_i=B_i\circ\mathrm{pr}_i$ coincide;
denote them by $\tilde{B}$.
By (f-1),
\[
\boldsymbol{\Theta}_{\mathrm{pr}_i\tilde{X},B_i}\,\psi(x,\cdot)=\mathrm{cov}_{\ell_i}\left(\Phi_{\boldsymbol{\Theta}}(\tilde{X},\tilde{B})\,\psi(x,\cdot)\right),
\]
and $\mathrm{cov}_{\ell_i}$ keeps the indices and the coefficients.

\emph{Proof of (f-2).}
Let $\tilde{X}$ be as above and $\tilde{U}$ its $\eta_k$-neighbourhood;
$\mathrm{Grc}_k(\tilde{X})$ is the subalgebra of $\mathrm{Grc}_k(\tilde{U})$ generated by the generators all of whose photon sites lie in $\tilde{X}$.
For $F\in\mathrm{Grc}_k(\tilde{U})$, written in the basis of the monomials, let $\mathrm{Fs}$ be the sum of those terms whose monomial is a product $\psi(\xi_1)\psi(\xi_2)$ of two generators $\psi$, without a further factor, with $\xi_1$ and $\xi_2$ over $\tilde{X}$.
Then $\mathrm{Fs}$ lies in the subalgebra of $\mathrm{Grc}_k(\tilde{X})$ generated by the $\psi(x,\beta,\omega)$, $x$ over $\tilde{X}$, on which $\Phi_{\boldsymbol{\Theta}}(\tilde{X},\tilde{B})$ is defined.
For $\tilde{B}$ as above and all $F\in\mathrm{Grc}_k(\tilde{U})$ put
\[
\Phi_{\mathrm{nat}}(\tilde{X};F,\tilde{B}):=
\begin{cases}
F-\mathrm{Fs}+\Phi_{\boldsymbol{\Theta}}(\tilde{X},\tilde{B})(\mathrm{Fs}) & \text{if } d_M(\tilde{X})\le L,\\
F & \text{otherwise}.
\end{cases}
\]
The side $\ell$ does not enter this map.
The map $T_v$ sends every monomial to a monomial in the generators with the translated indices, up to a sign;
in particular it sends $\psi(\xi_1)\psi(\xi_2)$ to $\psi(\xi_1+v)\psi(\xi_2+v)$, and $\xi_i+v$ lies over $\tilde{X}+v$ if and only if $\xi_i$ lies over $\tilde{X}$.
So the element $\mathrm{Fs}$ which belongs to $(\tilde{X}+v,T_vF)$ is $T_v\mathrm{Fs}$.
Since $d_M(\tilde{X}+v)=d_M(\tilde{X})$, the relation for $\Phi_{\boldsymbol{\Theta}}$ gives
\[
\Phi_{\mathrm{nat}}(\tilde{X}+v;T_vF,T_v\tilde{B})=T_v\Phi_{\mathrm{nat}}(\tilde{X};F,\tilde{B})\qquad\text{for }v\in M\mathbb{Z}^3.
\]
Now let $\ell\in\mathrm{Cl}$, let $E$ be a family on $T_k(\ell)$ with (k2)$^{\mathrm{o}}$, let $\mathscr{A}\in\tilde{R}_k$, and let $X$ be a polymer of $T_k(\ell)$, small or large, with a lift $\tilde{X}$ such that $\operatorname{sup\text{-}diameter}(\tilde{X})+2\eta_k+M<\ell$.
Let $F$ be the element $\tilde{F}$ of (Dl4) for $E(X;\mathscr{A})$, that is, the element of $\mathrm{Grc}_k(\tilde{U})$ with the kernels of $E(X;\mathscr{A})$, and let $\tilde{B}:=\mathscr{A}\circ\mathrm{pr}$ on the bonds with both ends in $\tilde{X}$.
By (Dl4), $F$ contains only generators $\psi$ and $\chi$, and $\psi(\xi)$ only with $\xi$ over $\tilde{X}$;
hence $\mathrm{Fs}=F^{[2,0]}$.
So $\Phi_{\mathrm{nat}}(\tilde{X};F,\tilde{B})$ equals $F-F^{[2,0]}+\Phi_{\boldsymbol{\Theta}}(\tilde{X},\tilde{B})(F^{[2,0]})$ if $d_M(\tilde{X})\le L$, and $F$ otherwise:
on the elements with these kernels $\Phi_{\mathrm{nat}}$ is the map described in Proposition~17.2\,(f).
By (Dl4), $\mathrm{cov}_\ell F=E(X;\mathscr{A})$ and $\mathrm{cov}_\ell F^{[2,0]}=E^{[2,0]}(X;\mathscr{A})$.
By (Dl2), $d_M(\tilde{X})\le L$ holds if and only if $X$ is small.
For small $X$, (f-1) gives
\[
\mathrm{cov}_\ell\,\Phi_{\mathrm{nat}}(\tilde{X};F,\tilde{B})=E(X;\mathscr{A})-E^{[2,0]}(X;\mathscr{A})+\boldsymbol{\Theta}_{X,\mathscr{A}}\left(\mathrm{cov}_\ell F^{[2,0]}\right)=E^{\mathrm{nat}}(X;\mathscr{A});
\]
for large $X$, $\mathrm{cov}_\ell F=E(X;\mathscr{A})=E^{\mathrm{nat}}(X;\mathscr{A})$.
\end{proof}

\begin{named}{Remark}[Proposition~48.3 at $k=0$, and its constants]
For $k=0$ one has $\eta_0=1$, and every pair has $d_X\ge1$;
a last link $((z_n,g),\tilde{x})$ with $\pi\tilde{x}=z_n$ and with a label of $\tilde{x}$ different from $g$ lies inside a cluster, has $d_X=1$ and carries no phase;
and Proposition~17.1\,(c1)--(c3) hold for all $k\ge0$.
The constants $\mathrm{e}^2$, $3c_X\cdot L\cdot M$ and $10$ do not depend on $k$, $\ell$, $e_k$, $X$ or $\mathscr{A}$, and the charge $e_k$ enters only through $\mathrm{ya}\le\frac{1}{30}$.
Proposition~48.3 involves no parameters other than $k$, $\ell$ and $e_k$.
\end{named}

\subsection*{4. Proposition 48.4}

\begin{proof}[Proof of Proposition~48.4]
\emph{Preparation.}
The data are those of the setting (G2.S-2).
By Proposition~48.1\,(a), applied at the level $k+1$, the bond function $0$ lies in $\mathfrak{X}_{k+1}$.
Hence $(\mathscr{A},0)\in\tilde{R}_{k+1}\times\mathfrak{X}_{k+1}$, and Proposition~18.2, applied to this pair, provides $\mathrm{bh}$, $\mathscr{C}^{b}$, $\mathrm{mg}$, $\mathrm{kc}_1$, $\mathrm{kc}_2$ with the properties (o), (i) and (ii) of that Proposition for the pair $(\mathscr{A},0)$.
By (o), $\mathscr{C}^{b}=0$ and $\mathrm{kc}_2=0$, while $\mathrm{bh}$, $\mathrm{kc}_1$ and $\mathrm{mg}$ are constructed from $\mathscr{A}$ and $\tilde{c}$ alone; no family of activities enters their construction.
Here $\mathrm{rad}=2L\cdot M$.
Since $\eta_k\le1$ and $\ell\ge13L\cdot M$,
\[
2\,\mathrm{rad}+\eta_k\le4L\cdot M+1<13L\cdot M\le\ell,
\]
so that $K(\tilde{c},\mathrm{rad})$ is a ball in the sense of (G2.r2).

\emph{(a)}
The first assertion of (a) is the first relation in Proposition~18.2\,(i).
By Proposition~18.2\,(ii), the pair $(t\,\mathrm{bh},t\,\mathscr{C}^{b})$ lies in $R_k\times\mathfrak{X}_k$ for every complex number $t$ with $|t|\le L^{\frac{7}{4}}$.
In particular $t\,\mathrm{bh}\in R_k$, and $R_k\subset\tilde{R}_k$ by Section~0.1 (take $A_0=0$, $A'=0$, $\lambda=0$ in \cite[Definition~5.16\,(b)]{QEDI}).
For $t=L^{\frac{7}{4}}$ the inclusion $t\,\mathrm{bh}\in R_k$ means, in terms of the thresholds $\mathrm{t1}_k$, $\mathrm{t2}_k$, $\mathrm{t3}_k$ of $R_k$ (G2.r4),
\begin{equation}
\tag{4.a}
\begin{aligned}
&\text{the three norms of }\mathrm{bh}\text{ are smaller than}\\
&\qquad L^{-\frac{7}{4}}\cdot\mathrm{t1}_k,\ L^{-\frac{7}{4}}\cdot\mathrm{t2}_k,\ L^{-\frac{7}{4}}\cdot\mathrm{t3}_k\text{, respectively.}
\end{aligned}
\end{equation}

\emph{Numerical facts used in (b).}
The number $L$ is an integer with $L\ge1$, and $0<\frac{1}{8}-\frac{3\bar{\varepsilon}}{2}\le\frac{1}{8}$.
Hence the condition (L1) of Proposition~14.2\,(b), that is $L^{\frac{1}{8}-\frac{3\bar{\varepsilon}}{2}}\ge2400$, gives
\[
L^{\frac{1}{8}}\ge L^{\frac{1}{8}-\frac{3\bar{\varepsilon}}{2}}\ge2400,\qquad\text{that is,}\qquad L\ge2400^{8}.
\]
The condition (S1) of the same Proposition, which contains $\varepsilon\le\bar{\varepsilon}=\frac{1}{200}$, gives $-\frac{7}{8}+\frac{\varepsilon}{2}\le-\frac{3}{4}$.
Therefore
\begin{equation}
\tag{4.a$'$}
L\ge3,\qquad L^{-\frac{7}{8}+\frac{\varepsilon}{2}}\le L^{-\frac{3}{4}}\le1,\qquad L^{-\frac{7}{4}}+4L^{-\frac{7}{8}+\frac{\varepsilon}{2}}\le5L^{-\frac{3}{4}}\le5\cdot2400^{-6}<1.
\end{equation}
The second condition of (A9) in Proposition~14.2\,(b) reads $90c_X\cdot L\cdot M\cdot e_0^{\varepsilon}\le1$.
Since $c_X=54$ and $L,M\ge1$, one has $90c_X\cdot L\cdot M>1$ and $9\le90c_X$.
As $\varepsilon>0$, the first of these two inequalities and (A9) give $e_0^{\varepsilon}<1$, that is, $e_0<1$.
The second one and (A9) give
\[
9L\cdot M\cdot e_0^{\varepsilon}\le1.
\]

\emph{(b)}
Let $E$ be a family of level $k$ with the properties (k1)$^{\mathrm{o}}$, (k2)$^{\mathrm{o}}$ and (k4)$^{\mathrm{o}}$ of the definition in Proposition~48.1.
Fix $(n,m)$ and arguments
\[
\xi_i=(x_i,\beta_i,\omega_i)\quad(1\le i\le n),\qquad\boldsymbol{\zeta}_j=(x_j,x'_j,\beta_j,\omega_j)\quad(n<j\le n+m);
\]
here the lower index of $x_i$, $x_j$, $x'_j$ numbers the arguments, and the neighbour of type~3 of $x_j$ is written $x_j^{(3)}$.
For $1\le j\le n+m$ put $c_j:=\pi(x_j)$; this is the photon site of the site of the $j$-th argument if $j\le n$, and the photon site of the first site of the $j$-th argument if $j>n$.
Functions on the sites of $T_k$, such as $\mathrm{mg}$ and the functions $\Lambda_z$ below, are read as $\ell\mathbb{Z}^3$-periodic functions on $\mathbb{L}_k$.
In $d(c_j,\tilde{c})$ and $c_j-\tilde{c}$ the site $c_j$ stands for its image in $T_k$.

\emph{Step~1 (the real gauge function).}
By Proposition~18.1\,(c), $\mathscr{A}_L\in\tilde{R}_k$.
A constant bond function has vanishing difference quotients.
Hence (4.a) and the bound $|\mathrm{kc}_1|\le L^{-\frac{7}{8}+\frac{\varepsilon}{2}}\cdot\mathrm{t1}_k$ of Proposition~18.2\,(ii) show that the three norms of $\mathrm{bh}+\mathrm{kc}_1$ are smaller than
\[
\left(L^{-\frac{7}{4}}+L^{-\frac{7}{8}+\frac{\varepsilon}{2}}\right)\cdot\mathrm{t1}_k<\mathrm{t1}_k,\qquad L^{-\frac{7}{4}}\cdot\mathrm{t2}_k,\qquad L^{-\frac{7}{4}}\cdot\mathrm{t3}_k,
\]
respectively; the inequality on the left holds by (4.a$'$).
So $\mathrm{bh}+\mathrm{kc}_1\in R_k\subset\tilde{R}_k$.
By Proposition~18.2\,(i), $\mathrm{bh}+\mathrm{kc}_1=\mathscr{A}_L-d\,\mathrm{mg}$ on the bonds with both ends in $K(\tilde{c},\mathrm{rad})$.
These bonds include all bonds that have an end at distance at most $1$ from $X$, in particular all bonds at distance at most $1$ from $X$.
The function $\mathrm{mg}$ is real.
Therefore (k4)$^{\mathrm{o}}$, applied with $\mathscr{A}_L$, $\mathrm{bh}+\mathrm{kc}_1$, $\mathrm{mg}$ in the places of $\mathscr{A}$, $\mathscr{A}^1$, $\boldsymbol{\lambda}$, gives
\[
G_{\mathrm{mg}}\left[E(X;\mathrm{bh}+\mathrm{kc}_1)\right]=E(X;\mathscr{A}_L).
\]
By (Ak2) of Lemma~A, which is (I.T-G) of Section~I.4.5 written for the kernels, this identity reads
\begin{equation}
\tag{4.b}
E_{nm}(X;\mathscr{A}_L;\xi,\boldsymbol{\zeta})=E_{nm}(X;\mathrm{bh}+\mathrm{kc}_1;\xi,\boldsymbol{\zeta})\cdot\prod_{j=1}^{n+m}\exp\left\{(-1)^{\omega_j}\,i\,e_k\cdot\mathrm{mg}(c_j)\right\},
\end{equation}
and each factor of the product has modulus one.
For $n=m=0$ the product is empty.

\emph{Step~2 (the constant field).}
Write $\mathrm{kc}_1=\mathrm{kc}_1^{\mathrm{R}}+i\,\mathrm{kc}_1^{\mathrm{I}}$, where $\mathrm{kc}_1^{\mathrm{R}}$ and $\mathrm{kc}_1^{\mathrm{I}}$ are the real constant bond functions whose three components are the real parts, respectively the imaginary parts, of the three components of $\mathrm{kc}_1$.
Then
\[
|\mathrm{kc}_1^{\mathrm{R}}|,\ |\mathrm{kc}_1^{\mathrm{I}}|\le|\mathrm{kc}_1|\le L^{-\frac{7}{8}+\frac{\varepsilon}{2}}\cdot\mathrm{t1}_k.
\]
For a constant bond function $\mathrm{kc}$ and a vector $v\in\mathbb{R}^3$ put $\mathrm{kc}\cdot v:=\sum_{\mu}\mathrm{kc}_{\mu}v_{\mu}$.
Let
\[
\mathrm{Dz}:=\left\{z\in\mathbb{C}^2:\ |\operatorname{Im}z_1|<2,\ |\operatorname{Im}z_2|<2\right\},
\]
and for $z\in\mathrm{Dz}$ define
\begin{gather*}
A(z):=\mathrm{bh}+z_1\,\mathrm{kc}_1^{\mathrm{R}}+z_2\,\mathrm{kc}_1^{\mathrm{I}},\\
\Lambda_z(c):=-\left(z_1\,\mathrm{kc}_1^{\mathrm{R}}+z_2\,\mathrm{kc}_1^{\mathrm{I}}\right)\cdot(c-\tilde{c})\quad\text{for }c\in K(\tilde{c},\mathrm{rad}),\\
\Lambda_z(c):=0\quad\text{for the other sites }c\text{ of }T_k,
\end{gather*}
where $c-\tilde{c}$ is the vector of (G2.r2).
\begin{itemize}
\item[(2a)]
$A(z)\in\tilde{R}_k$.
Indeed,
\[
A(z)=\left[\operatorname{Re}z_1\,\mathrm{kc}_1^{\mathrm{R}}+\operatorname{Re}z_2\,\mathrm{kc}_1^{\mathrm{I}}\right]+\left[\mathrm{bh}+i\left(\operatorname{Im}z_1\,\mathrm{kc}_1^{\mathrm{R}}+\operatorname{Im}z_2\,\mathrm{kc}_1^{\mathrm{I}}\right)\right].
\]
The first bracket is a real constant bond function.
One has $|\operatorname{Im}z_1|\cdot|\mathrm{kc}_1^{\mathrm{R}}|+|\operatorname{Im}z_2|\cdot|\mathrm{kc}_1^{\mathrm{I}}|<4L^{-\frac{7}{8}+\frac{\varepsilon}{2}}\cdot\mathrm{t1}_k$, and constant bond functions have vanishing difference quotients.
Together with (4.a) and (4.a$'$) this shows that the three norms of the second bracket are smaller than
\[
\left(L^{-\frac{7}{4}}+4L^{-\frac{7}{8}+\frac{\varepsilon}{2}}\right)\cdot\mathrm{t1}_k<\mathrm{t1}_k,\qquad L^{-\frac{7}{4}}\cdot\mathrm{t2}_k,\qquad L^{-\frac{7}{4}}\cdot\mathrm{t3}_k,
\]
respectively.
So the second bracket lies in $R_k$, hence in $\tilde{R}_k$.
By Proposition~16.1\,(c), $\tilde{R}_k$ is invariant under the addition of real constant bond functions; therefore $A(z)\in\tilde{R}_k$.
\item[(2b)]
$A(z)=\mathrm{bh}-d\Lambda_z$ on the bonds with both ends in $K(\tilde{c},\mathrm{rad})$.
Indeed, let $c$ and $c+\eta_ke_{\mu}$ lie in $K(\tilde{c},\mathrm{rad})$.
The vectors $(c+\eta_ke_{\mu})-\tilde{c}$ and $(c-\tilde{c})+\eta_ke_{\mu}$ of $\mathbb{L}_k$ represent the same difference modulo $\ell\mathbb{Z}^3$, and in the sup norm they differ by at most $2\,\mathrm{rad}+\eta_k<\ell$; hence they are equal.
Consequently
\[
(d\Lambda_z)_{\mu}(c)=\eta_k^{-1}\left(\Lambda_z(c+\eta_ke_{\mu})-\Lambda_z(c)\right)=-\left(z_1\,\mathrm{kc}^{\mathrm{R}}_{1,\mu}+z_2\,\mathrm{kc}^{\mathrm{I}}_{1,\mu}\right),
\]
and $\mathrm{bh}-d\Lambda_z$ coincides with $A(z)$ on the bond from $c$ to $c+\eta_ke_{\mu}$.
\item[(2c)]
Put
\[
\mathrm{Kz}(z):=E_{nm}(X;A(z);\xi,\boldsymbol{\zeta})\cdot\prod_{j=1}^{n+m}\exp\left\{(-1)^{\omega_j}\,i\,e_k\cdot\Lambda_z(c_j)\right\}.
\]
The set $\mathrm{Dz}$ is open and convex.
The map $z\mapsto A(z)$ is affine, and its values lie in the open set $\tilde{R}_k$ by (2a) and \cite[Lemma~17.3]{QEDI}.
The kernel $\mathscr{A}\mapsto E_{nm}(X;\mathscr{A};\xi,\boldsymbol{\zeta})$ is analytic on $\tilde{R}_k$ by (k1)$^{\mathrm{o}}$, and $\Lambda_z(c_j)$ is linear in $z$.
Hence $\mathrm{Kz}$ is analytic on $\mathrm{Dz}$.
\item[(2d)]
Let $z\in\mathbb{R}^2$.
Then $\Lambda_z$ is a real function on the sites of $T_k$, and $\mathrm{bh}=A(0)$ and $A(z)$ belong to $\tilde{R}_k$.
The bonds at distance at most $1$ from $X$ have both ends in $K(\tilde{c},\mathrm{rad})$ (Step~1); so (2b) gives $A(z)=\mathrm{bh}-d\Lambda_z$ on these bonds.
Applying (k4)$^{\mathrm{o}}$ with $\mathrm{bh}$, $A(z)$, $\Lambda_z$ in the places of $\mathscr{A}$, $\mathscr{A}^1$, $\boldsymbol{\lambda}$, one obtains
\[
G_{\Lambda_z}\left[E(X;A(z))\right]=E(X;\mathrm{bh}).
\]
By (Ak2) this identity reads $\mathrm{Kz}(z)=E_{nm}(X;\mathrm{bh};\xi,\boldsymbol{\zeta})=\mathrm{Kz}(0)$.
\item[(2e)]
The function $\mathrm{Kz}-\mathrm{Kz}(0)$ is analytic on the connected open set $\mathrm{Dz}$ and vanishes on $\mathbb{R}^2$.
At the point $0$ each of its partial derivatives can be computed along real directions, and therefore vanishes.
So the power series of $\mathrm{Kz}-\mathrm{Kz}(0)$ at $0$ is zero.
By the identity theorem for analytic functions of several variables, (cl2), one has $\mathrm{Kz}=\mathrm{Kz}(0)$ on $\mathrm{Dz}$.
\item[(2f)]
The point $z=(1,i)$ lies in $\mathrm{Dz}$.
At this point $A(z)=\mathrm{bh}+\mathrm{kc}_1^{\mathrm{R}}+i\,\mathrm{kc}_1^{\mathrm{I}}=\mathrm{bh}+\mathrm{kc}_1$, and $\Lambda_{(1,i)}(c)=-\mathrm{kc}_1\cdot(c-\tilde{c})$ for $c\in K(\tilde{c},\mathrm{rad})$.
Thus (2e) gives
\begin{equation}
\tag{4.c}
E_{nm}(X;\mathrm{bh}+\mathrm{kc}_1;\xi,\boldsymbol{\zeta})\cdot\prod_{j=1}^{n+m}\exp\left\{(-1)^{\omega_j}\,i\,e_k\cdot\Lambda_{(1,i)}(c_j)\right\}=E_{nm}(X;\mathrm{bh};\xi,\boldsymbol{\zeta}).
\end{equation}
\end{itemize}

\emph{Step~3 (the bound).}
For $n=m=0$ neither (4.b) nor (4.c) contains a factor, and
\[
E_{00}(X;\mathscr{A}_L)=E_{00}(X;\mathrm{bh}+\mathrm{kc}_1)=E_{00}(X;\mathrm{bh}).
\]
Let $n+m\ge1$.
If some $\xi_i$ does not lie over $X$, or some $\boldsymbol{\zeta}_j$ is not admissible for $X$, then $E_{nm}(X;\mathscr{A}_L;\xi,\boldsymbol{\zeta})$ and $E_{nm}(X;\mathrm{bh};\xi,\boldsymbol{\zeta})$ both vanish by (k2)$^{\mathrm{o}}$.
Otherwise every $c_j$ has sup distance at most $\eta_k\le1$ from $X$.
Indeed, if the $j$-th argument is a $\xi_j$ or a pair of kind~(I), then $c_j\in X$; if it is a pair of kind~(II), then $x'_j=x_j^{(3)}$, and $c_j$ or its nearest neighbour $\pi\left(x_j^{(3)}\right)$ lies in $X$ (Section~I.1.7).
By \cite[Theorem~13.2\,(t7)]{QEDI}, $X\subset K\left(\tilde{c},M\cdot(L+2)\right)$.
Since $M\ge1$, and $L\ge3$ by (4.a$'$), it follows that
\[
d(c_j,\tilde{c})\le M\cdot(L+2)+1\le L\cdot M+3M\le2L\cdot M=\mathrm{rad},
\]
so that $c_j\in K(\tilde{c},\mathrm{rad})$, and that
\[
|c_j-\tilde{c}|_1\le3\,|c_j-\tilde{c}|\le3\left(M\cdot(L+2)+1\right)\le6L\cdot M\le9L\cdot M.
\]
Hence
\[
\left|\operatorname{Im}\Lambda_{(1,i)}(c_j)\right|=\left|\mathrm{kc}_1^{\mathrm{I}}\cdot(c_j-\tilde{c})\right|\le|\mathrm{kc}_1^{\mathrm{I}}|\cdot|c_j-\tilde{c}|_1.
\]
Now $|\mathrm{kc}_1^{\mathrm{I}}|\le L^{-\frac{7}{8}+\frac{\varepsilon}{2}}\cdot\mathrm{t1}_k$, $e_k\cdot\mathrm{t1}_k=e_k^{\frac{1}{4}+\varepsilon}$, and $L^{-\frac{7}{8}+\frac{\varepsilon}{2}}\le1$ by (4.a$'$); moreover $e_k\le e_0<1$ and $9L\cdot M\cdot e_0^{\varepsilon}\le1$, the two consequences of the second condition of (A9) obtained after (4.a$'$).
Therefore
\[
e_k\cdot\left|\operatorname{Im}\Lambda_{(1,i)}(c_j)\right|\le9L\cdot M\cdot L^{-\frac{7}{8}+\frac{\varepsilon}{2}}\cdot e_k^{\frac{1}{4}+\varepsilon}\le9L\cdot M\cdot e_0^{\varepsilon}\le1.
\]
Consequently each of the $n+m$ factors of the product in (4.c) has modulus between $\mathrm{e}^{-1}$ and $\mathrm{e}$.
Together with (4.b) this gives
\begin{align*}
\left|E_{nm}(X;\mathscr{A}_L;\xi,\boldsymbol{\zeta})\right|&=\left|E_{nm}(X;\mathrm{bh}+\mathrm{kc}_1;\xi,\boldsymbol{\zeta})\right|\le\mathrm{e}^{n+m}\cdot\left|E_{nm}(X;\mathrm{bh};\xi,\boldsymbol{\zeta})\right|\\
&=\exp(n+m)\cdot\left|E_{nm}(X;\mathrm{bh};\xi,\boldsymbol{\zeta})\right|\le\exp\left(2(n+m)\right)\cdot\left|E_{nm}(X;\mathrm{bh};\xi,\boldsymbol{\zeta})\right|.
\end{align*}
By (Ak1), the norms $\left\|E_{nm}(X;\mathscr{A}_L)\right\|$ and $\left\|E_{nm}(X;\mathrm{bh})\right\|$ are the integrals $\int\mathrm{d}\xi\,\mathrm{d}\boldsymbol{\zeta}$ of the moduli of the two kernels.
Integrating the last chain of inequalities over the arguments therefore gives
\[
\left\|E_{nm}(X;\mathscr{A}_L)\right\|\le\exp(n+m)\cdot\left\|E_{nm}(X;\mathrm{bh})\right\|\le\exp\left(2(n+m)\right)\cdot\left\|E_{nm}(X;\mathrm{bh})\right\|.
\]

\emph{(c)}
\emph{The first assertion of (c): the property (k2)$^{\mathrm{o}}$ is the locality of $E(X;\cdot)$ in $X$.}
Let $E$ be a family of level $k$ and $X\in D_k$; put $\mathrm{Dc}:=\tilde{R}_k\times\{0\}$ and $F_X(\mathscr{A},0):=E(X;\mathscr{A})$.
In the absence of spectators, $\mathrm{Gr}_k[X]$ is the linear span of the basis monomials that contain only generators $\psi(\xi)$ with $\xi$ over $X$ and $\chi(\boldsymbol{\zeta})$ with $\boldsymbol{\zeta}$ admissible for $X$.
By (Ak1), the coefficient of a basis monomial is the value of the corresponding kernel at the (increasing) arguments of the monomial, multiplied by a positive number, and the kernels are antisymmetric.
Hence the first clause of (k2)$^{\mathrm{o}}$ states that $F_X(\mathscr{A},0)\in\mathrm{Gr}_k[X]$ for all $\mathscr{A}$, which is condition~(i) of the definition in Proposition~18.3.
Two points $(\mathscr{A},0)$, $(\tilde{\mathscr{A}},0)$ of $\mathrm{Dc}$ coincide in both components on the bonds that meet $X$ if and only if $\mathscr{A}$ and $\tilde{\mathscr{A}}$ coincide there.
Since an element and its kernels determine each other (Ak1), the second clause of (k2)$^{\mathrm{o}}$ is therefore condition~(ii).
This argument involves the data of a single level only, namely (Ak1) and the definition in Proposition~18.3, whose data are those of one level.
Under (G2.S-2) the triple $(k+1,\ell/L,e_{k+1})$ satisfies (G2.S-1).
Hence the first assertion of (c) holds in the same way for a family of level $k+1$, a polymer of $D_{k+1}$ and $\mathrm{Dc}=\tilde{R}_{k+1}\times\{0\}$.

\emph{Restriction to a subset of the domain.}
Let a map be local in $Z$ on $\mathrm{Dc}$, and let $\mathrm{Dc}_0\subset\mathrm{Dc}$.
Then the restriction of the map to $\mathrm{Dc}_0$ is local in $Z$ on $\mathrm{Dc}_0$, because both conditions of the definition are assertions about all points, respectively all pairs of points, of the domain.

\emph{The second assertion of (c): the locality of $\mathscr{A}\mapsto S_L\left[F(\mathscr{A}_L)\right]$ in $Y$.}
Let $Y\in D_{k+1}$, let $X\in D_k$ be contained in $LY$, and let $F$ be local in $X$ on $\tilde{R}_k\times\{0\}$, without spectators; in the statement, $F(\mathscr{A}_L)$ stands for $F(\mathscr{A}_L,0)$.
Put $F^{\mathrm{c}}(\mathscr{A},\mathscr{C}):=F(\mathscr{A},0)$ for $(\mathscr{A},\mathscr{C})\in\tilde{R}_k\times\mathfrak{X}_k$.
Then $F^{\mathrm{c}}$ is local in $X$ on $\tilde{R}_k\times\mathfrak{X}_k$.
Indeed, its values are values of $F$, hence lie in $\mathrm{Gr}_k[X]$.
If $(\mathscr{A},\mathscr{C})$ and $(\tilde{\mathscr{A}},\tilde{\mathscr{C}})$ coincide in both components on the bonds that meet $X$, then so do $(\mathscr{A},0)$ and $(\tilde{\mathscr{A}},0)$, and
\[
F^{\mathrm{c}}(\mathscr{A},\mathscr{C})=F(\mathscr{A},0)=F(\tilde{\mathscr{A}},0)=F^{\mathrm{c}}(\tilde{\mathscr{A}},\tilde{\mathscr{C}}).
\]
(For $F=E(X;\,\cdot\,)$ this is Proposition~48.1\,(b).)
By Proposition~18.3\,(e), applied to the pair of domains $\left(\tilde{R}_k\times\mathfrak{X}_k,\tilde{R}_{k+1}\times\mathfrak{X}_{k+1}\right)$, the map
\[
(\mathscr{A},\mathscr{C})\mapsto S_L\left[F^{\mathrm{c}}(\mathscr{A}_L,\mathscr{C}_L)\right]=S_L\left[F(\mathscr{A}_L,0)\right]
\]
is local in $Y$ on $\tilde{R}_{k+1}\times\mathfrak{X}_{k+1}$, at the level $k+1$.
Since $0\in\mathfrak{X}_{k+1}$ (Proposition~48.1\,(a) at the level $k+1$), the set $\tilde{R}_{k+1}\times\{0\}$ is contained in this domain.
By the restriction property just proved, the map $\mathscr{A}\mapsto S_L\left[F(\mathscr{A}_L,0)\right]$ is local in $Y$ on $\tilde{R}_{k+1}\times\{0\}$.

\emph{The third assertion of (c): the property (k2)$^{\mathrm{o}}$ of $L\,E'$ at the level $k+1$.}
Let $E'$ have (k2)$^{\mathrm{o}}$ at the level $k$, and let $Y\in D_{k+1}$.
Every $X$ with $\bar{X}^{L}=LY$ is contained in $LY$.
By the two preceding assertions of (c), for each such $X$ the map $\mathscr{A}\mapsto S_L\left[E'(X;\mathscr{A}_L)\right]$ is local in $Y$ on $\tilde{R}_{k+1}\times\{0\}$.
The maps that are local in $Y$ on a fixed domain form an algebra (Proposition~18.3\,(b) at the level $k+1$).
Hence the finite sum $(LE')(Y;\,\cdot\,)$ of these maps is local in $Y$ on $\tilde{R}_{k+1}\times\{0\}$.
By the first assertion of (c), at the level $k+1$, this is (k2)$^{\mathrm{o}}$ for $LE'$.
\end{proof}

\begin{named}{Remark}
The bound of (b) holds with the factor $\exp(n+m)$ at every level $k\ge0$, including $k=0$, where $\eta_0=1$ and the first site of a pair of kind~(II) has distance at most $1$ from $X$.
The factor is the same for all data $(k,\ell,e_k)$ with $\ell/L\in\mathrm{Cl}$ and $e_{k+1}\le e_0$, for all $X$, $\tilde{c}$ and $\mathscr{A}$, and for every family $E$; the charge enters only through $e_k\le e_0<1$ and $9L\cdot M\cdot e_0^{\varepsilon}\le1$.
\end{named}

\subsection*{5. Conditions on the constants}

The proofs of Sections~0--4 rest on the following conditions on the constants and on the side~$\ell$.
The lines (C0)--(C8) named below are those of Section~I.8.2, and the labels (L1), (L5), (M1), (S1), (A8), (A9) are those of Proposition~14.2\,(b).

\begin{itemize}
\item
The condition $0<\rho<1$ of the line (C0) is needed only in the form $0\le\rho\le1$, by which the moduli of the $16$ entries of $M_{3}(x)$ add up to at most $24$ (I.m7).
In this form it is applied in Proposition~48.1\,(c), in the proof of the property (k6)$^{\mathrm{o}}$ of the family $E^{\mathrm{bd}}$.
The parameter $\rho$ itself occurs in Sections~0--4 only in the example of Proposition~48.1\,(c): it stands in the matrix $M_{3}(x)$ of (I.2.1), which enters the definition of the family $E^{\mathrm{bd}}$.
\item
Of the weights $w=(6,6\sqrt{2},6\sqrt{2},6)$ of the line (C0) only $w_3=6$ matters.
It enters through $M_{3}(x)=3\left(\rho-\gamma\cdot(g\,e_1)\right)$ (I.2.1), through the number $24$ of (I.m7), and through the coefficient $-3\eta_k^3/48\ne0$, and it concerns Proposition~48.1\,(c) alone: the definition of the family $E^{\mathrm{bd}}$, its property (k6)$^{\mathrm{o}}$, and the fact that $(E^{\mathrm{bd}})^{\mathrm{c}}$ does not have (k4).
\item
The numbers $c_X=54$, $\kappa_0=54\log72$, $K_0=\exp(\kappa_0)/36$ and $c_{\mathrm{fl}}=61c_X\cdot K_0$ of the line (C0) satisfy $K_0\le c_{\mathrm{fl}}$, $2\le c_{\mathrm{fl}}$ (as $K_0=72^{54}/36\ge1$) and $9\le90c_X$; moreover $90c_X\cdot L\cdot M>1$.
These inequalities serve in Proposition~48.2\,(b), (d), and in the proof of Proposition~48.4\,(b), where $e_0<1$ and $9L\cdot M\cdot e_0^{\varepsilon}\le1$ are derived after (4.a$'$).
\item
The condition $\kappa\ge7\kappa_0+7$ of the line (C2) gives $\kappa\ge\kappa_0$ and $\kappa\ge0$.
The inequality $\kappa\ge\kappa_0$ is applied in Proposition~48.2\,(b), and $\kappa\ge0$ in Step~3 of the proof of Proposition~48.1\,(d).
\item
The first condition of the line (C3), $L^{\frac{1}{8}-\frac{3\bar{\varepsilon}}{2}}\ge2400$, labelled (L1), gives $L\ge2400^8$.
Together with $\varepsilon\le\frac{1}{200}$ and $5\cdot2400^{-6}<1$ it yields the three inequalities (4.a$'$), namely $L\ge3$, $L^{-\frac{7}{8}+\frac{\varepsilon}{2}}\le1$ and $L^{-\frac{7}{4}}+4L^{-\frac{7}{8}+\frac{\varepsilon}{2}}<1$; they are derived in Section~4 and applied in Steps~1--3 of the proof of Proposition~48.4\,(b).
For $L\ge3$ and $M\ge1$ one has moreover $M\cdot(L+2)+1\le2L\cdot M$ and $3\left(M\cdot(L+2)+1\right)\le9L\cdot M$, which enter Step~3.
\item
The conditions that $L$ be odd with $L\ge5$ (line (C3), label (L5)) and that $M$ be a power of $L$ with $M\ge L^2$ (line (C5), label (M1)) give $L\ge1$, hence $\eta_k=L^{-k}\le1$ and $1/L\le1$; further, $M/\eta_k$ is an odd integer, $M\ge1$ and $1<M/2$.
These facts are applied in Section~0.1, in Lemma~D, in Proposition~48.1\,(c), in Proposition~48.2\,(b), (e), (f), in Proposition~48.3\,(f) (for the map $\Phi_{\boldsymbol{\Theta}}$), and in Proposition~48.4: in the Preparation, in (4.a$'$) and the derivation following it, and in Steps~2 and~3 of (b).
\item
The condition $1-\alpha<\varepsilon\le\bar{\varepsilon}=\frac{1}{200}$ of the line (C7), labelled (S1), where $\alpha<1$ by the line (C1), gives $\varepsilon>0$, $\frac{1}{4}-\varepsilon\ge\varepsilon$ and $-\frac{7}{8}+\frac{\varepsilon}{2}\le-\frac{3}{4}$.
It is needed in Proposition~48.1\,(c), in Proposition~48.3\,(d) (through $e_k^{\varepsilon}\le e_0^{\varepsilon}$), and in Proposition~48.4\,(b): in (4.a$'$), in the derivation following it, and in Step~3.
\item
The first condition of (A8), $e_0^{\varepsilon}\le\frac{1}{4}$, in the line (C8), is applied together with its consequence $e_0<1$ in Proposition~48.1\,(c).
\item
The second condition of (A9), $90c_X\cdot L\cdot M\cdot e_0^{\varepsilon}\le1$, also in the line (C8), gives $3c_X\cdot L\cdot M\cdot e_k^{\varepsilon}\le\frac{1}{30}$, $e_0<1$ and $9L\cdot M\cdot e_0^{\varepsilon}\le1$.
These are applied in Proposition~48.3\,(d), together with $\mathrm{e}^2\left(\frac{31}{30}\right)^2<7.9\le10$, and in Proposition~48.4\,(b), in the derivation following (4.a$'$) and in Step~3.
\item
The side $\ell$ is a datum, not a constant; it is subject to $\ell\in\mathrm{Cl}$, the class of sides.
Thus $\ell$ is an integer, hence a multiple of $\eta_k$, it is a multiple of $M$, and $\ell\ge13L\cdot M$; together with $L,M\ge1\ge\eta_k$ the last inequality gives
\[
2M\cdot(L+2)<\ell,\qquad M\cdot(L+2)+\eta_k<\ell,\qquad M\cdot(L+2)+2+M<\ell,\qquad 4L\cdot M+1<\ell.
\]
These facts are applied in Lemma~D, in Step~2 of the proof of Proposition~48.1\,(d), in Proposition~48.2\,(a), (c), (f), in Proposition~48.3\,(f), and in Proposition~48.4.
\end{itemize}

In Proposition~14.2\,(b) the conditions on $e_0$ in this list are imposed for every $e_k\in(0,e_0]$.
Since $\varepsilon>0$ and $e_k\le e_0$, one has $e_k^{\varepsilon}\le e_0^{\varepsilon}$, so that each of them holds at $e_k$ once it holds at $e_0$.

Each Proposition of Chapters~15, 16, 17 and~18a that is applied in Sections~0--4 holds under the conditions on the constants that are listed in the section on the constants of its chapter and in Section~G2.9.
Among these are (S1) and (A8) for Proposition~17.1\,(a), (A9) for Proposition~17.1\,(c2), and (L1), (S1), (A8) and (A9) for Propositions~18.1\,(c) and~18.2\,(ii).

The constants above are constants in the sense of (I.U4)$^{\mathrm{ch}}$ of Section~I.8.1.
In each of the conditions on them a constant is bounded only in terms of absolute numbers and of constants chosen before it in Section~I.8.2, and none of them depends on $N$, $k$, $\ell$, the twist, a mass, the pair $(m_{\mathrm{in}},\varepsilon^{\mathrm{v}}_{\mathrm{in}})$, $\mathscr{C}$ or $\bar{m}$.

\section*{Chapter 48b. Contraction without slot; the invariant $E_k=\mathscr{E}_k$ at $\mathscr{C}=0$: proofs of Propositions 48.5 and 48.6}

\subsection*{0. Conventions of this chapter}

\emph{Setting.}
Throughout the chapter the data are those of the setting (S-2) of Section~G2.0.4.
The letter $j$ stands for $k$ or for $k+1$, and the level $k+1$ carries the side $\ell/L$ and the charge $e_{k+1}$.
The constants are those of Section~I.8.
On the level $j$ the measures are
\[
\int\mathrm{d}x=\eta_j^3\cdot\frac{1}{48}\sum_{c}\sum_{g\in O_{\mathrm{cub}}}\quad\text{(I.1.3)},\qquad
\int\mathrm{d}\xi=\sum_{\xi}\frac{\eta_j^3}{48},\qquad
\int\mathrm{d}\boldsymbol{\zeta}=\sum_{\boldsymbol{\zeta}}\left(\frac{\eta_j^3}{48}\right)^2\quad\text{(Section~I.3.6)}.
\]
A polymer is called small or large in the sense of (G2.r3).

\emph{Notation.}
For $F\in\mathrm{Gr}[\psi,\chi]$ the part of bidegree $(n,m)$ is written $F^{[n,m]}$; accordingly $E^{a,[n,m]}$ and $E^{r,[2,0]}$ are the parts of bidegree $(n,m)$ and $(2,0)$ of the families $E^{a}$ and $E^{r}$ introduced in Section~3.1.
One puts $\mathrm{rad}_t:=L^{\frac{7}{4}}$.
For an $M$-cube $\square$ the symbol $\bar{\square}^{L}$ denotes the set $\bar{X}^{L}$ with $X=\square$.
The components of vectors $\mathrm{cu},\mathrm{cv}\in\mathbb{R}^3$ are written $\mathrm{cu}^{\mu}$, $\mathrm{cv}^{\mu}$.
In (48.2) below, $\mathrm{kq}$ denotes a number with $\mathrm{kq}\ge\kappa_0$; it takes the place of the number $\kappa'$ of \cite[Theorem~13.2(t8)]{QEDI}.

\emph{The mass coupling of level $0$.}
At $k=0$ the symbol $m_k$ stands for $m_{\mathrm{in}}$, a datum of Proposition~48.6\,(b) (Section~G4.0.5).
In Section~5.3 the symbol $m_0$ is the integer $m_0$ of \cite[Definition~5.98]{QEDI}, which enters \cite[Section~19]{QEDI}; it is given the value $0$ (Section~I.6.1).

\subsection*{1. Facts used in both proofs}

\begin{named}{Fact 1}[Domains]
The zero bond function belongs to $\tilde{R}_j$ and to $\mathfrak{X}_j$.
Indeed, the thresholds in \cite[Definition~5.16]{QEDI} (Section~I.5.3) are positive numbers, and $0=A_0+A_1$ with $A_1=0$ and $A_0=0=A'+d\boldsymbol{\lambda}$, where $A'=0$ and $\boldsymbol{\lambda}=0$ on every $\square^{\dagger}$.
By \cite[Lemma~17.3]{QEDI} the set $\tilde{R}_j$ is open in the space of the complex bond functions on $T_j$.
If $\mathscr{A}\in\tilde{R}_{k+1}$, then the bond function $\mathscr{A}_L$, $\mathscr{A}_L(b)=L^{-\frac{1}{2}}\cdot\mathscr{A}(b/L)$, belongs to $\tilde{R}_k$ (Proposition~18.1\,(c)).
The map $\mathscr{A}\mapsto\mathscr{A}_L$ is complex linear, and $0_L=0$.
\end{named}

\begin{named}{Fact 2}[Restriction]
For a family $F$ with slot of level $j$ (G2.r8) let
\[
F|_0(X;\mathscr{A}):=F(X;\mathscr{A},0),\qquad X\in D_j,\quad\mathscr{A}\in\tilde{R}_j;
\]
the right side is defined because $0\in\mathfrak{X}_j$ (Fact~1).
If $F\in K^{x,C}_j$, this is the family of Proposition~48.1\,(a).
The elements $F|_0(X;\mathscr{A})$ and $F(X;\mathscr{A},0)$ of $\mathrm{Gr}[\psi,\chi]$ are one and the same.
Consequently they have the same antisymmetric kernels (Proposition~15\,(a)), the same part of each bidegree and the same norms, and
\[
(F+F')|_0=F|_0+F'|_0,\qquad(z\cdot F)|_0=z\cdot F|_0\quad\text{for every complex number }z.
\]
\end{named}

\begin{named}{Fact 3}[Bidegrees]
Let $F\in\mathrm{Gr}[\psi,\chi]$ be of level $j$, and let $\mathsf{w}=(\mathsf{w}_1,\mathsf{w}_2)$ be a pair of weights ((G2.r7), that is Proposition~15\,(n); Section~I.5.1).
Then
\begin{gather*}
\|F\|_{\mathsf{w}}=\sum_{n,m}\left\|F^{[n,m]}\right\|_{\mathsf{w}},\qquad
\left\|F^{[n,m]}\right\|_{\mathsf{w}}=\frac{\mathsf{w}_1^n\cdot\mathsf{w}_2^m}{n!\,m!}\cdot\|F_{nm}\|,\\
\|F_{nm}\|=\int\mathrm{d}\xi\,\mathrm{d}\boldsymbol{\zeta}\,\left|F_{nm}(\xi_1,\dots,\xi_n;\boldsymbol{\zeta}_1,\dots,\boldsymbol{\zeta}_m)\right|=\sum_{\xi,\boldsymbol{\zeta}}\left(\frac{\eta_j^3}{48}\right)^{n+2m}\cdot\left|F_{nm}(\xi;\boldsymbol{\zeta})\right|.
\end{gather*}
In particular $\left\|F^{[0,0]}\right\|_{\mathsf{w}}=|F_{00}|$ and $\left\|F^{[2,0]}\right\|_{h_k}=\frac{h_{k,1}^2}{2}\cdot\|F_{20}\|$.
\end{named}

\begin{named}{Fact 4}[Scaling factors]
For integers $n,m\ge0$ let
\[
\mathrm{sca}(n,m):=L^{-n-(1+\alpha)m}\cdot L^{-\frac{n+m}{8}+\frac{m\varepsilon}{2}}.
\]
By (G2.r5) and Proposition~18.4\,(a) the element $S_LF^{[n,m]}$ is of bidegree $(n,m)$.
Hence $(S_LF)^{[n,m]}=S_LF^{[n,m]}$, and Proposition~18.4\,(b) together with Fact~3 gives, for every $F\in\mathrm{Gr}[\psi,\chi]$ of level $k$,
\begin{equation}
\|S_LF\|_{h_{k+1}}=\sum_{n,m}\mathrm{sca}(n,m)\cdot\left\|F^{[n,m]}\right\|_{h_k}.\tag{48.1}
\end{equation}
Since $0<\varepsilon\le1/200$ and $1-\varepsilon<\alpha<1$ ((C1), (C7)), the following hold.
\begin{itemize}
\item[(sc0)] $\mathrm{sca}(0,0)=1$ and $\mathrm{sca}(2,0)=L^{-2-\frac{2}{8}}=L^{-\frac{9}{4}}$.
\item[(sc1)] $\log_L\mathrm{sca}(1,1)=-2-\alpha-\frac{1}{4}+\frac{\varepsilon}{2}<-2-(1-\varepsilon)-\frac{1}{4}+\frac{\varepsilon}{2}=-\frac{13}{4}+\frac{3\varepsilon}{2}$.
\item[(sc2)] $\log_L\mathrm{sca}(0,2)=-2-2\alpha-\frac{1}{4}+\varepsilon<-2-2(1-\varepsilon)-\frac{1}{4}+\varepsilon=-\frac{17}{4}+3\varepsilon\le-\frac{13}{4}+\frac{3\varepsilon}{2}$, the last inequality because $3\varepsilon/2\le1$.
\item[(sc3)] $\log_L\mathrm{sca}(n,m)=-\frac{9}{8}(n+m)-m\left(\alpha-\frac{\varepsilon}{2}\right)\le-\frac{9}{8}(n+m)$, since $\alpha>1-\varepsilon>\varepsilon/2$; if $n+m\ge4$, this is at most $-\frac{9}{2}\le-4\le-\frac{13}{4}+\frac{3\varepsilon}{2}$.
\item[(sc4)] By (sc3), $\mathrm{sca}(n,m)\le1$ for all $(n,m)$; therefore (48.1) gives $\|S_LF\|_{h_{k+1}}\le\|F\|_{h_k}$.
\end{itemize}
\end{named}

\begin{named}{Fact 5}[The sum over $X$]
The set $D_k$ is finite.
By \cite[Theorem~13.2(t5)]{QEDI}, to every $X\in D_k$ there corresponds exactly one $Y\in D_{k+1}$ with $\bar{X}^{L}=LY$.
The proof of \cite[Theorem~13.2(t5)]{QEDI} rests on the elementary fact that two adjacent $M$-cubes lie in one $LM$-cube or in two $LM$-cubes with a common face, and on its consequence that $\bar{X}^{L}$ is connected through faces.
For completeness, the argument is as follows.
Since $L$ is odd, the $LM$-cube with the centre $z'\in LM\mathbb{Z}^3$ is the union of the $M$-cubes with the centres $z'+Mi$, $i\in\mathbb{Z}^3$, $|i|\le(L-1)/2$.
If $z=z'+Mi$, then
\begin{align*}
z+Me_{\mu}&=z'+M(i+e_{\mu})&&\text{if }i_{\mu}<\frac{L-1}{2},\\
z+Me_{\mu}&=(z'+LMe_{\mu})+M\left(i-(L-1)e_{\mu}\right)&&\text{if }i_{\mu}=\frac{L-1}{2},
\end{align*}
and in the second case the component $\mu$ of $i-(L-1)e_{\mu}$ equals $-(L-1)/2$.
So the $LM$-cubes of two $M$-cubes with the centres $z$ and $z+Me_{\mu}$ either coincide or have the centres $z'$ and $z'+LMe_{\mu}$.
Hence a chain of $M$-cubes of $X$ in which consecutive cubes have a common face is carried to a chain of $LM$-cubes of $\bar{X}^{L}$ in which consecutive cubes coincide or have a common face.
Every $LM$-cube of $\bar{X}^{L}$ contains an $M$-cube of $X$ (definition of $\bar{X}^{L}$ in \cite[Section~13.1.7]{QEDI}), and any two $M$-cubes of $X$ are joined by such a chain.
Therefore $\bar{X}^{L}$ is a union of $LM$-cubes that is connected through faces.
By \cite[Theorem~13.2(t8)]{QEDI}, for every $Y\in D_{k+1}$ and every number $\mathrm{kq}\ge\kappa_0$,
\begin{equation}
\sum_{X\in D_k:\,\bar{X}^{L}=LY}\exp\left(-\mathrm{kq}\cdot d_M(X)\right)\le c_{\mathrm{rb}}\cdot L^3\cdot\exp\left(-L\cdot(\mathrm{kq}-\kappa_0-1)\cdot d_M(Y)\right).\tag{48.2}
\end{equation}
Both assertions concern sets of photon sites; they hold for $k\ge0$ and $\ell,\ell/L\in\mathrm{Cl}$.
\end{named}

\subsection*{3. Proof of Proposition 48.5}

The conditions on the constants that enter the proof are collected in Section~6.

\subsubsection*{3.1. The family $E^{a}$. Proof of (d), first and second assertions}

\begin{proof}[Proof of Proposition~48.5\,(d), first and second assertions]
Let $E\in K_k^{C}$, and put $E^{r}:=RE$ and $E^{a}:=(E^{r})^{\mathrm{nat}}$; the clauses (k1)$^{\mathrm{o}}$--(k7)$^{\mathrm{o}}$ are those of the definition in Proposition~48.1.
Proposition~48.2\,(d) shows that $R$ is linear and that $E^{r}\in K_k^{C}$, with $E^{r}\in K_k$ if $E\in K_k$.
Thus $E^{r}$ is a family of even, balanced elements with (k1)$^{\mathrm{o}}$--(k6)$^{\mathrm{o}}$.
By Proposition~48.3\,(e) the map $\mathrm{nat}$ is linear and $E^{a}\in K_k^{C}$, with $E^{a}\in K_k$ if $E^{r}\in K_k$.
It follows that
\begin{equation}
E\mapsto E^{a}\ \text{is linear};\qquad E^{a}\in K_k^{C};\qquad E^{a}\in K_k\ \text{if}\ E\in K_k.\tag{48.3}
\end{equation}

\emph{Bounds.}
Let $A\in\tilde{R}_k$.
For small $X$ one applies Proposition~48.3\,(d) to the balanced family $E^{r}$, which has (k2)$^{\mathrm{o}}$, then Proposition~48.2\,(d), and finally (k6)$^{\mathrm{o}}$ at the points $A$ and $0$ of $\tilde{R}_k$ (Fact~1):
\[
\left\|E^{a}(X;A)\right\|_{h_k}\le10\left\|E^{r}(X;A)\right\|_{h_k}\le10\left(\|E(X;A)\|_{h_k}+\|E(X;0)\|_{h_k}\right)\le20\,\|E\|_k\cdot\exp\left(-\kappa\cdot d_M(X)\right).
\]
For large $X$ the definitions of $R$ (Proposition~48.2) and of $\mathrm{nat}$ (Proposition~48.3) give $E^{r}(X;A)=E(X;A)$ and $E^{a}(X;A)=E^{r}(X;A)$, whence $\left\|E^{a}(X;A)\right\|_{h_k}=\|E(X;A)\|_{h_k}\le\|E\|_k\cdot\exp\left(-\kappa\cdot d_M(X)\right)$.
In terms of
\[
\mathrm{sup}^{a}(X):=\sup_{A\in\tilde{R}_k}\left\|E^{a}(X;A)\right\|_{h_k}
\]
these bounds read
\begin{equation}
\begin{aligned}
\mathrm{sup}^{a}(X)&\le20\,\|E\|_k\cdot\exp\left(-\kappa\cdot d_M(X)\right)&&(X\ \text{small}),\\
\mathrm{sup}^{a}(X)&\le\|E\|_k\cdot\exp\left(-\kappa\cdot d_M(X)\right)&&(X\ \text{large}).
\end{aligned}\tag{48.4}
\end{equation}

\emph{Vanishing.}
Let $X$ be small.
Proposition~48.2\,(d) gives
\[
E^{r}_{00}(X;0)=0,\qquad\int\mathrm{d}x\int\mathrm{d}x'\,E^{r}_{20}(X;0;x,x')=0,
\]
where each of the two integrals is $\eta_k^3/48$ times the sum over the photon sites of one lift of $X$ and over the $48$ labels.
Proposition~48.3\,(b), applied to the balanced family $E^{r}$, yields $E^{a}_{00}=E^{r}_{00}$, hence $E^{a}_{00}(X;0)=0$, and also $E^{a}_{20}(X;0;\cdot,\cdot)=0$.
(The kernel $E^{a}_{20}(\xi_1,\xi_2)$ is determined by the matrices $E^{a}_{20}(x,x')$ of Section~I.5.1: its entries with the bar indices $(1,0)$ are the entries of these matrices, those with the bar indices $(0,1)$ are obtained from them by antisymmetry, and those with $(0,0)$ and $(1,1)$ vanish because the family is balanced.
So the kernel vanishes as a function of $(\xi_1,\xi_2)$ if and only if the matrices vanish on $(X\times O_{\mathrm{cub}})^2$.)
\end{proof}

\subsubsection*{3.2. Proof of (d), third assertion}

\begin{proof}[Proof of Proposition~48.5\,(d), third assertion]
Let $X$ be small and $\mathscr{A}\in\tilde{R}_{k+1}$, and put $B:=\mathscr{A}_L$; by Fact~1, $B\in\tilde{R}_k$.
Since $X$ is not empty, one may choose a photon site $\tilde{c}$ of $X$; let $\mathrm{bh}$ be the bond function of Proposition~48.4 for $(X,\tilde{c},\mathscr{A})$.
\begin{itemize}
\item[(i)] \emph{The disk.}
The set
\[
\mathrm{Dsk}:=\left\{t\in\mathbb{C}:t\cdot\mathrm{bh}\in\tilde{R}_k\right\}
\]
is open, as $\tilde{R}_k$ is open (Fact~1) and $t\mapsto t\cdot\mathrm{bh}$ is continuous.
By Proposition~48.4\,(a), $\mathrm{Dsk}$ contains the compact disk $\{|t|\le\mathrm{rad}_t\}$, where $\mathrm{rad}_t=L^{\frac{7}{4}}$.
A compact subset of an open subset of $\mathbb{C}$ has a positive distance from the complement of the open set; therefore $\mathrm{Dsk}$ contains a disk $\{|t|<\mathrm{rad}'\}$ with $\mathrm{rad}'>\mathrm{rad}_t$.
For every $(n,m)$ and all arguments the function $t\mapsto E^{a}_{nm}(X;t\cdot\mathrm{bh};\xi;\boldsymbol{\zeta})$ is analytic on $\mathrm{Dsk}$, being the composition of the complex linear map $t\mapsto t\cdot\mathrm{bh}$ with a function that is analytic on $\tilde{R}_k$ by (k1)$^{\mathrm{o}}$ (see (48.3)).
Moreover $\mathrm{rad}_t\ge2$ ($L^{\frac{1}{4}}\ge2$, (C3)); hence $\mathrm{rad}_t>1$ and $\mathrm{rad}_t-1\ge\mathrm{rad}_t/2$.
\item[(ii)] \emph{Bidegree $(0,0)$.}
Let $\mathrm{fa}_0(t):=E^{a}_{00}(X;t\cdot\mathrm{bh})$.
By Section~3.1, $\mathrm{fa}_0(0)=E^{a}_{00}(X;0)=0$.
Clause (k5)$^{\mathrm{o}}$ for $E^{a}$ gives $\mathrm{Cg}\left[E^{a}(X;-t\cdot\mathrm{bh})\right]=E^{a}(X;t\cdot\mathrm{bh})$ for $|t|<\mathrm{rad}'$.
Here $\mathrm{Cg}$ is an automorphism of the algebra $\mathrm{Gr}[\psi,\chi]$ that preserves the bidegree (Section~I.4.5), and it leaves the numbers fixed.
Comparing the parts of bidegree $(0,0)$ one finds $\mathrm{fa}_0(-t)=\mathrm{fa}_0(t)$, so that $\mathrm{fa}_0'(0)=0$.
Proposition~15\,(j) yields
\[
\mathrm{fa}_0(1)=\mathrm{fa}_0(1)-\mathrm{fa}_0(0)-\mathrm{fa}_0'(0)=\oint^{(2)}_{\mathrm{rad}_t}\mathrm{fa}_0
\]
and
\[
|\mathrm{fa}_0(1)|\le\frac{\sup_{|t|=\mathrm{rad}_t}|\mathrm{fa}_0(t)|}{\mathrm{rad}_t\cdot(\mathrm{rad}_t-1)}\le2L^{-\frac{7}{2}}\cdot\sup_{|t|=\mathrm{rad}_t}\left|E^{a}_{00}(X;t\cdot\mathrm{bh})\right|\le2L^{-\frac{7}{2}}\cdot\mathrm{sup}^{a}(X),
\]
where the last step uses Fact~3 in the form $\left|E^{a}_{00}(X;A)\right|=\left\|E^{a,[0,0]}(X;A)\right\|_{h_k}\le\left\|E^{a}(X;A)\right\|_{h_k}$.
As $E^{a}$ has (k1)$^{\mathrm{o}}$, (k2)$^{\mathrm{o}}$ and (k4)$^{\mathrm{o}}$, Proposition~48.4\,(b) gives $E^{a}_{00}(X;B)=E^{a}_{00}(X;\mathrm{bh})=\mathrm{fa}_0(1)$.
Together with (sc0) this yields
\begin{equation}
\mathrm{sca}(0,0)\cdot\left\|E^{a,[0,0]}(X;B)\right\|_{h_k}=\left|E^{a}_{00}(X;B)\right|\le2L^{-\frac{7}{2}}\cdot\mathrm{sup}^{a}(X).\tag{48.5}
\end{equation}
\item[(iii)] \emph{Bidegree $(2,0)$.}
Let $\mathrm{fa}_2(t):=E^{a}_{20}(X;t\cdot\mathrm{bh};\cdot,\cdot)$.
Its values lie in the finite-dimensional space of the functions $K(\xi_1,\xi_2)$ of two indices, which carries the norm
\[
\|K\|=\int\mathrm{d}\xi_1\,\mathrm{d}\xi_2\,|K(\xi_1,\xi_2)|=\sum_{\xi_1,\xi_2}\left(\frac{\eta_k^3}{48}\right)^2\cdot|K(\xi_1,\xi_2)|.
\]
By (i) the function $\mathrm{fa}_2$ is analytic on $\{|t|<\mathrm{rad}'\}$, and $\mathrm{fa}_2(0)=E^{a}_{20}(X;0;\cdot,\cdot)=0$ by Section~3.1.
Proposition~15\,(j) yields $\mathrm{fa}_2(1)=\mathrm{fa}_2(1)-\mathrm{fa}_2(0)=\oint_{\mathrm{rad}_t}\mathrm{fa}_2$ and
\[
\|\mathrm{fa}_2(1)\|\le\frac{\sup_{|t|=\mathrm{rad}_t}\|\mathrm{fa}_2(t)\|}{\mathrm{rad}_t-1}\le2L^{-\frac{7}{4}}\cdot\sup_{|t|=\mathrm{rad}_t}\left\|E^{a}_{20}(X;t\cdot\mathrm{bh})\right\|.
\]
Proposition~48.4\,(b) at $(n,m)=(2,0)$ gives $\left\|E^{a}_{20}(X;B)\right\|\le\exp(4)\cdot\left\|E^{a}_{20}(X;\mathrm{bh})\right\|=\mathrm{e}^4\cdot\|\mathrm{fa}_2(1)\|$.
By Fact~3,
\[
\frac{h_{k,1}^2}{2}\cdot\left\|E^{a}_{20}(X;A)\right\|=\left\|E^{a,[2,0]}(X;A)\right\|_{h_k}\le\left\|E^{a}(X;A)\right\|_{h_k}\le\mathrm{sup}^{a}(X)\qquad\text{for }A\in\tilde{R}_k.
\]
With (sc0) it follows that
\begin{equation}
\begin{split}
\mathrm{sca}(2,0)\cdot\left\|E^{a,[2,0]}(X;B)\right\|_{h_k}&=L^{-\frac{9}{4}}\cdot\frac{h_{k,1}^2}{2}\cdot\left\|E^{a}_{20}(X;B)\right\|\\
&\le L^{-\frac{9}{4}}\cdot\mathrm{e}^4\cdot2L^{-\frac{7}{4}}\cdot\mathrm{sup}^{a}(X)=2\mathrm{e}^4L^{-4}\cdot\mathrm{sup}^{a}(X).
\end{split}\tag{48.6}
\end{equation}
\item[(iv)] \emph{The other bidegrees.}
The element $E^{a}(X;B)$ is even; so $E^{a,[n,m]}(X;B)=0$ whenever $n+m$ is odd.
The pairs $(n,m)$ with even $n+m$, other than $(0,0)$ and $(2,0)$, are $(1,1)$, $(0,2)$ and those with $n+m\ge4$.
For all of them $\mathrm{sca}(n,m)\le L^{-\frac{13}{4}+\frac{3\varepsilon}{2}}$ by (sc1), (sc2) and (sc3).
By Fact~3,
\begin{equation}
\sum_{(n,m)\ne(0,0),(2,0)}\mathrm{sca}(n,m)\cdot\left\|E^{a,[n,m]}(X;B)\right\|_{h_k}\le L^{-\frac{13}{4}+\frac{3\varepsilon}{2}}\cdot\left\|E^{a}(X;B)\right\|_{h_k}\le L^{-\frac{13}{4}+\frac{3\varepsilon}{2}}\cdot\mathrm{sup}^{a}(X).\tag{48.7}
\end{equation}
\item[(v)] \emph{Sum.}
One has $2L^{-\frac{7}{2}}\le L^{-\frac{13}{4}}$ because $L^{\frac{1}{4}}\ge2$, and $2\mathrm{e}^4L^{-4}\le L^{-\frac{13}{4}}$ because $2\mathrm{e}^4<110\le L^{\frac{3}{4}}$ ((C3)).
From (48.1), (48.5), (48.6) and (48.7) one obtains
\[
\left\|S_L\left[E^{a}(X;\mathscr{A}_L)\right]\right\|_{h_{k+1}}\le\left(L^{-\frac{13}{4}}+L^{-\frac{13}{4}}+L^{-\frac{13}{4}+\frac{3\varepsilon}{2}}\right)\cdot\mathrm{sup}^{a}(X)\le3L^{-\frac{13}{4}+\frac{3\varepsilon}{2}}\cdot\mathrm{sup}^{a}(X).
\]
Only the quantity $\mathrm{sup}^{a}(X)$ enters these bounds, and in it the supremum over the connections is taken of the whole sum over $(n,m)$ in Fact~3: the parts of bidegree $(0,0)$ and $(2,0)$ at $B$ are bounded by suprema, over the circle $\{t\cdot\mathrm{bh}:|t|=\mathrm{rad}_t\}$ in $\tilde{R}_k$, of a single term of that sum, and the remaining parts are taken at the single point $B$.
\end{itemize}
This proves (d).
\end{proof}

\subsubsection*{3.3. Proof of (b)}

\begin{proof}[Proof of Proposition~48.5\,(b)]
Let $Y\in D_{k+1}$ and $\mathscr{A}\in\tilde{R}_{k+1}$.
The definition of $LRE$, the triangle inequality and Fact~5 give
\[
\left\|(LRE)(Y;\mathscr{A})\right\|_{h_{k+1}}\le\sum_{\substack{X\text{ small}\\\bar{X}^{L}=LY}}\left\|S_L\left[E^{a}(X;\mathscr{A}_L)\right]\right\|_{h_{k+1}}+\sum_{\substack{X\text{ large}\\\bar{X}^{L}=LY}}\left\|S_L\left[E^{a}(X;\mathscr{A}_L)\right]\right\|_{h_{k+1}}.
\]

\emph{Small $X$.}
By (d) and (48.4), each term of the first sum is at most $3L^{-\frac{13}{4}+\frac{3\varepsilon}{2}}\cdot20\,\|E\|_k\cdot\exp\left(-\kappa\cdot d_M(X)\right)$.
Apply (48.2) with $\mathrm{kq}=\kappa$ (note $\kappa\ge\kappa_0$ by (C2)) and use $L\cdot(\kappa-\kappa_0-1)\ge\kappa$ ((C3), third condition):
\begin{multline*}
\text{first sum}\le60L^{-\frac{13}{4}+\frac{3\varepsilon}{2}}\cdot\|E\|_k\cdot c_{\mathrm{rb}}\cdot L^3\cdot\exp\left(-L\cdot(\kappa-\kappa_0-1)\cdot d_M(Y)\right)\\
\le60\,c_{\mathrm{rb}}\cdot L^{-\frac{1}{4}+\frac{3\varepsilon}{2}}\cdot\|E\|_k\cdot\exp\left(-\kappa\cdot d_M(Y)\right).
\end{multline*}

\emph{Large $X$.}
By (sc4), Fact~1 and (48.4),
\[
\left\|S_L\left[E^{a}(X;\mathscr{A}_L)\right]\right\|_{h_{k+1}}\le\left\|E^{a}(X;\mathscr{A}_L)\right\|_{h_k}\le\|E\|_k\cdot\exp\left(-\kappa\cdot d_M(X)\right).
\]
As $d_M(X)>L$, one has $\exp\left(-\kappa\cdot d_M(X)\right)\le\exp(-L)\cdot\exp\left(-(\kappa-1)\cdot d_M(X)\right)$.
Now apply (48.2) with $\mathrm{kq}=\kappa-1$ (note $\kappa-1\ge\kappa_0$ by (C2)), and use $L\cdot(\kappa-\kappa_0-2)\ge\kappa$ ((C3), third condition) and $L^3\cdot\exp(-L)\le L^{-\frac{1}{4}}$ ((C3), first condition):
\[
\text{second sum}\le\exp(-L)\cdot\|E\|_k\cdot c_{\mathrm{rb}}\cdot L^3\cdot\exp\left(-L\cdot(\kappa-\kappa_0-2)\cdot d_M(Y)\right)\le c_{\mathrm{rb}}\cdot L^{-\frac{1}{4}}\cdot\|E\|_k\cdot\exp\left(-\kappa\cdot d_M(Y)\right).
\]

\emph{Together.}
Adding the two bounds,
\begin{multline*}
\left\|(LRE)(Y;\mathscr{A})\right\|_{h_{k+1}}\le\left(60L^{-\frac{1}{4}+\frac{3\varepsilon}{2}}+L^{-\frac{1}{4}}\right)\cdot c_{\mathrm{rb}}\cdot\|E\|_k\cdot\exp\left(-\kappa\cdot d_M(Y)\right)\\
\le61\,c_{\mathrm{rb}}\cdot L^{-\frac{1}{4}+\frac{3\varepsilon}{2}}\cdot\|E\|_k\cdot\exp\left(-\kappa\cdot d_M(Y)\right).
\end{multline*}
Here $61\,c_{\mathrm{rb}}=c_{\mathrm{fl}}$ (C0), and $L^{\frac{3\varepsilon}{2}}\le L^{2\varepsilon}$.
Multiplication by $\exp\left(\kappa\cdot d_M(Y)\right)$ and the suprema over $\mathscr{A}\in\tilde{R}_{k+1}$ and $Y\in D_{k+1}$ give
\[
\|LRE\|_{k+1}\le c_{\mathrm{fl}}\cdot L^{-\frac{1}{4}+\frac{3\varepsilon}{2}}\cdot\|E\|_k\le c_{\mathrm{fl}}\cdot L^{-\frac{1}{4}+2\varepsilon}\cdot\|E\|_k.
\]
\end{proof}

\subsubsection*{3.4. Proof of (a)}

\begin{proof}[Proof of Proposition~48.5\,(a)]
\emph{The map is defined and linear.}
For $Y\in D_{k+1}$ and $\mathscr{A}\in\tilde{R}_{k+1}$ the sum that defines $(LRE)(Y;\mathscr{A})$ is finite (Fact~5), and $\mathscr{A}_L$ lies in $\tilde{R}_k$ (Fact~1), where $E^{a}(X;\cdot)$ is defined.
The map $E\mapsto E^{a}(X;\mathscr{A}_L)$ is linear by (48.3), and $S_L$ is linear; hence $E\mapsto LRE$ is linear.

\emph{Even, balanced.}
By (G2.r5), $S_L$ sends every generator to a number times a generator of the same family with the same bar index $\omega$.
Consequently it sends a monomial to a number times a monomial with the same numbers of barred and of unbarred legs.
Since $E^{a}(X;\mathscr{A}_L)$ is even and balanced (48.3), so is $(LRE)(Y;\mathscr{A})$.

\emph{(k1)$^{\mathrm{o}}$.}
By Proposition~18.4\,(c), and because the antisymmetric kernels of a sum are the sums of the antisymmetric kernels (Proposition~15\,(a)), one has for indices $\xi$, $\boldsymbol{\zeta}$ of level $k+1$
\begin{equation}
\begin{aligned}
(LRE)_{nm}(Y;\mathscr{A};\xi;\boldsymbol{\zeta})
&=\sum_{X:\,\bar{X}^{L}=LY}L^{3n+6m}\cdot L^{-n-(1+\alpha)m}\cdot E^{a}_{nm}(X;\mathscr{A}_L;L\xi;L\boldsymbol{\zeta})\\
&\qquad\text{if all pairs of }\boldsymbol{\zeta}\text{ have }d_X\le M/L,\\
(LRE)_{nm}(Y;\mathscr{A};\xi;\boldsymbol{\zeta})
&=0\qquad\text{otherwise}.
\end{aligned}\tag{48.8}
\end{equation}
Here the factor $L^{3n+6m}$ is the quotient of the weights $(\eta_k^3/48)^{n+2m}$ and $(\eta_{k+1}^3/48)^{n+2m}$ of the $n+m$ arguments, $L\xi$ and $L\boldsymbol{\zeta}$ are as in (G2.r5), and the labels are left unchanged.
The set $\tilde{R}_{k+1}$ is open, the map $\mathscr{A}\mapsto\mathscr{A}_L$ is complex linear and maps it into $\tilde{R}_k$ (Fact~1), and $E^{a}_{nm}(X;\cdot;L\xi;L\boldsymbol{\zeta})$ is analytic on $\tilde{R}_k$.
So each term of the finite sum is analytic in $\mathscr{A}\in\tilde{R}_{k+1}$.

\emph{(k2)$^{\mathrm{o}}$.}
The family $E^{a}$ has (k2)$^{\mathrm{o}}$ at level $k$ (48.3); by Proposition~48.4\,(c), $LE^{a}=LRE$ has (k2)$^{\mathrm{o}}$ at level $k+1$.

\emph{(k3)$^{\mathrm{o}}$, (k4)$^{\mathrm{o}}$, (k5)$^{\mathrm{o}}$.}
The family $E^{a}$ has (k3)$^{\mathrm{o}}$ for all elements of $\mathbb{G}_k$, in particular for those of the subgroup $\mathbb{G}_k^{L}$, and it has (k4)$^{\mathrm{o}}$ and (k5)$^{\mathrm{o}}$.
Proposition~48.2\,(e) with $E'=E^{a}$ shows that $LRE$ has (k3)$^{\mathrm{o}}$ for the group $\mathbb{G}_{k+1}$ (side $\ell/L$), and (k4)$^{\mathrm{o}}$ and (k5)$^{\mathrm{o}}$ at level $k+1$.

\emph{(k6)$^{\mathrm{o}}$.}
By (b), $\|LRE\|_{k+1}<\infty$.

Therefore $LRE\in K_{k+1}^{C}(\ell/L)$.

\emph{(k7)$^{\mathrm{o}}$.}
If $E\in K_k$, then $E^{a}\in K_k$ (48.3), that is, $E^{a}$ has (k7)$^{\mathrm{o}}$; by Proposition~48.2\,(e), $LRE$ has (k7)$^{\mathrm{o}}$.
So $LR$ maps $K_k(\ell)$ into $K_{k+1}(\ell/L)$.

By linearity, (b) applies to differences:
\[
\|LRE_1-LRE_2\|_{k+1}\le c_{\mathrm{fl}}\cdot L^{-\frac{1}{4}+2\varepsilon}\cdot\|E_1-E_2\|_k\qquad\text{for }E_1,E_2\in K_k^{C}.
\]
\end{proof}

\subsubsection*{3.5. Proof of (c)}

\begin{proof}[Proof of Proposition~48.5\,(c)]
Let $\mathscr{E}\in K^{x,C}_k$ and $E:=\mathscr{E}|_0$.
Then $E\in K_k^{C}$ (Proposition~48.1\,(a)), so that $LRE$ is the map of Proposition~48.5.
The family $LR\mathscr{E}$ is the one defined in the data of Proposition~18.5, formed with $R$ of Section~I.5.6 and with $\mathrm{nat}$ of Proposition~17.2.
\begin{itemize}
\item[(i)] \emph{Extraction numbers.}
Let $X$ be small.
By Fact~2 the elements $\mathscr{E}(X;0,0)$ and $E(X;0)$ coincide, together with all their kernels.
Hence, with all arguments over one lift $\tilde{X}$ of $X$,
\begin{align*}
\mathrm{al}_0(\mathscr{E},X)&=\operatorname{Vol}(X)^{-1}\cdot\mathscr{E}_{00}(X;0,0)=\operatorname{Vol}(X)^{-1}\cdot E_{00}(X;0)=\mathrm{al}_0(E,X),\\
\mathrm{al}_2(\mathscr{E},X)&=\operatorname{Vol}(X)^{-1}\cdot\left(\frac{\eta_k^3}{48}\right)^2\sum_{c,c'\in\tilde{X}}\;\sum_{g,g'\in O_{\mathrm{cub}}}\mathscr{E}_{20}\left(X;0,0;(c,g),(c',g')\right)\\
&=\operatorname{Vol}(X)^{-1}\int\mathrm{d}x\int\mathrm{d}x'\,E_{20}(X;0;x,x')=\mathrm{al}_2(E,X).
\end{align*}
\item[(ii)] \emph{The map $R$.}
Let $A\in\tilde{R}_k$.
For small $X$, (i) and Fact~2 give
\begin{align*}
(R\mathscr{E})(X;A,0)&=\mathscr{E}(X;A,0)-\mathrm{al}_0(\mathscr{E},X)\cdot\operatorname{Vol}(X)-\int_X\mathrm{d}x\,\bar{\psi}(x)^T\,\mathrm{al}_2(\mathscr{E},X)\,\psi(x)\\
&=E(X;A)-\mathrm{al}_0(E,X)\cdot\operatorname{Vol}(X)-\int_X\mathrm{d}x\,\bar{\psi}(x)^T\,\mathrm{al}_2(E,X)\,\psi(x)=(RE)(X;A).
\end{align*}
For large $X$ one has $(R\mathscr{E})(X;A,0)=\mathscr{E}(X;A,0)=E(X;A)=(RE)(X;A)$.
Therefore
\begin{equation}
(R\mathscr{E})|_0=R(\mathscr{E}|_0).\tag{48.9}
\end{equation}
\item[(iii)] \emph{The map $\mathrm{nat}$.}
Let $A\in\tilde{R}_k$.
At the pair $(A,0)$ the total connection is $\mathscr{T}=A+0=A$.
For small $X$ one uses the definition of $\mathrm{nat}$ in Proposition~17.2, then (48.9), and then the definition of $\mathrm{nat}$ in Proposition~48.3.
In both definitions $\boldsymbol{\Theta}$ is the homomorphism $\boldsymbol{\Theta}_{X,B}$ of Proposition~17.1\,(d), here at the bond function $B=A$; the element $(R\mathscr{E})^{[2,0]}(X;A,0)=(RE)^{[2,0]}(X;A)$ lies in its domain, the subalgebra generated by the $\psi(\xi)$ with $\xi$ over $X$, because $RE$ has (k2)$^{\mathrm{o}}$ (Proposition~48.2\,(d)).
This gives
\begin{align*}
(R\mathscr{E})^{\mathrm{nat}}(X;A,0)&=(R\mathscr{E})(X;A,0)-(R\mathscr{E})^{[2,0]}(X;A,0)+\boldsymbol{\Theta}_{X,A}\left((R\mathscr{E})^{[2,0]}(X;A,0)\right)\\
&=(RE)(X;A)-(RE)^{[2,0]}(X;A)+\boldsymbol{\Theta}_{X,A}\left((RE)^{[2,0]}(X;A)\right)=(RE)^{\mathrm{nat}}(X;A).
\end{align*}
For large $X$ both sides equal $(RE)(X;A)$.
Therefore
\begin{equation}
\left((R\mathscr{E})^{\mathrm{nat}}\right)\big|_0=\left(R(\mathscr{E}|_0)\right)^{\mathrm{nat}}.\tag{48.10}
\end{equation}
\item[(iv)] \emph{The map $L$.}
Let $Y\in D_{k+1}$ and $\mathscr{A}\in\tilde{R}_{k+1}$.
The pair $(\mathscr{A},0)$ lies in $\tilde{R}_{k+1}\times\mathfrak{X}_{k+1}$, $0_L=0$, and $\mathscr{A}_L\in\tilde{R}_k$ (Fact~1).
By (48.10),
\begin{align*}
(LR\mathscr{E})|_0(Y;\mathscr{A})&=(LR\mathscr{E})(Y;\mathscr{A},0)=\sum_{X:\,\bar{X}^{L}=LY}S_L\left[(R\mathscr{E})^{\mathrm{nat}}(X;\mathscr{A}_L,0)\right]\\
&=\sum_{X:\,\bar{X}^{L}=LY}S_L\left[(RE)^{\mathrm{nat}}(X;\mathscr{A}_L)\right]=(LRE)(Y;\mathscr{A}).
\end{align*}
\end{itemize}
\end{proof}

\subsubsection*{3.6. Proof of (e)}

\emph{The clause.}
Clause (e) of Proposition~48.5 is the assertion of Proposition~18.5\,(f) in which $E(X;0)$, $E(X;\mathscr{A}_L)$, $(RE)^{\mathrm{nat}}(X;\mathscr{A}_L)$, $\mathscr{A}_L$ and (k2)$^{\mathrm{o}}$ stand in the places of $\mathscr{E}(X;0,0)$, $\mathscr{E}(X;\mathscr{A}_L,\mathscr{C}_L)$, $(R\mathscr{E})^{\mathrm{nat}}(X;\mathscr{A}_L,\mathscr{C}_L)$, $\mathscr{T}_L$ and (k2) (Section~I.5.4), respectively, with $E\in K_k^{C}$ and $\mathscr{A}\in\tilde{R}_{k+1}$.
Its first sentence (on $S_L$) is denoted by (e-i);
its second assertion, namely the hypothesis on the lift $\tilde{X}$ and the assertion on the kernels, together with the specification of the data, by (e-ii).
The rest of the clause, that is, everything that follows the specification of the data, is denoted by (e-iii);
it begins with the sentence on the data on which the formula is defined, namely the data in which the kernels of bidegree $(2,0)$ vanish unless both of their arguments lie over $\tilde{X}$, and it extends to the end of the clause.
The margin of (e-ii) is $\operatorname{sup-diameter}(\tilde{X})+2\eta_k+M<\ell$.

\emph{Covering objects} (Section~I.1.9, Section~G3.5.0).
Let $\tilde{U}$ be a finite set of sites of $\mathbb{L}_j$.
Here $\mathrm{Grc}_j(\tilde{U})$ denotes the Grassmann algebra with the independent generators $\psi(x,\beta,\omega)$, $x=(c,g)$ with $c\in\tilde{U}$, and $\chi(x,x',\beta,\omega)$, $(x,x')\in\mathrm{Pzt}_j$ with $c$ and $c'$ in $\tilde{U}$; the generators $\Psi$, $W$ of Section~I.1.9 do not occur.
These indices $\xi=(x,\beta,\omega)$ and $\boldsymbol{\zeta}=(x,x',\beta,\omega)$ are called the indices over $\tilde{U}$; they carry the weights $\eta_j^3/48$ and $(\eta_j^3/48)^2$, and $G_{nm}$ denotes the antisymmetric kernels of $G\in\mathrm{Grc}_j(\tilde{U})$ for these weights (Proposition~15\,(a)).
The map $T_v:\mathrm{Grc}_j(\tilde{U})\to\mathrm{Grc}_j(\tilde{U}+v)$ relabels the generators, $\psi(x,\cdot)\mapsto\psi(x+v,\cdot)$ and $\chi(x,x',\cdot)\mapsto\chi(x+v,x'+v,\cdot)$; on kernels it acts by $(T_vK)(\xi;\boldsymbol{\zeta})=K(\xi-v;\boldsymbol{\zeta}-v)$, on bond functions by $(T_v\tilde{B})(b)=\tilde{B}(b-v)$, and the kernels of $T_vG$ are the $T_vG_{nm}$.
If $\operatorname{sup-diameter}(\tilde{U})+M<\ell$, the homomorphism $\mathrm{cov}_{\ell}:\mathrm{Grc}_k(\tilde{U})\to\mathrm{Gr}[\psi,\chi]$ (torus of side $\ell$) maps each generator to the generator of the torus with the same index in the sense of (I.AP).
Finally, $\tilde{U}^{+a}:=\left\{c\in\mathbb{L}_j:\mathrm{dist}(c,\tilde{U})\le a\right\}$.

\begin{named}{(K)}[Kernels of $\mathrm{cov}_{\ell}G$]
Let $\operatorname{sup-diameter}(\tilde{U})+M<\ell$ and $G\in\mathrm{Grc}_k(\tilde{U})$.
Read as functions on the covering index sets that are antiperiodic in each argument (I.AP), the antisymmetric kernels of $\mathrm{cov}_{\ell}G$ are
\[
\left(\mathrm{cov}_{\ell}G\right)_{nm}(\xi;\boldsymbol{\zeta})=
\begin{cases}
G_{nm}(\xi;\boldsymbol{\zeta}) & \text{if all arguments are indices over }\tilde{U},\\
0 & \text{if some argument is not congruent modulo }\ell\mathbb{Z}^3\text{ to an index over }\tilde{U}.
\end{cases}
\]
\end{named}

\begin{proof}
If two indices over $\tilde{U}$ are congruent modulo $\ell\mathbb{Z}^3$, their first points differ by a vector of $\ell\mathbb{Z}^3$ whose sup norm is at most $\operatorname{sup-diameter}(\tilde{U})<\ell$; hence the two indices coincide.
Consequently there is exactly one function $K_{nm}$ on the covering index sets that is antiperiodic in each argument and has the two properties displayed above; it is antisymmetric in the $\xi$ and in the $\boldsymbol{\zeta}$.
The product $K_{nm}(\xi;\boldsymbol{\zeta})\,\psi(\xi_1)\cdots\chi(\boldsymbol{\zeta}_m)$ is $\ell\mathbb{Z}^3$-periodic in each argument.
Therefore the integral over a fundamental domain equals the sum over the indices over $\tilde{U}$:
\begin{multline*}
\sum_{n,m}\frac{1}{n!\,m!}\int\mathrm{d}\xi\,\mathrm{d}\boldsymbol{\zeta}\,K_{nm}(\xi;\boldsymbol{\zeta})\,\psi(\xi_1)\cdots\chi(\boldsymbol{\zeta}_m)\\
=\sum_{n,m}\frac{1}{n!\,m!}\sum_{\xi,\boldsymbol{\zeta}\text{ over }\tilde{U}}\left(\frac{\eta_k^3}{48}\right)^{n+2m}\cdot G_{nm}(\xi;\boldsymbol{\zeta})\cdot\mathrm{cov}_{\ell}\left(\psi(\xi_1)\right)\cdots\mathrm{cov}_{\ell}\left(\chi(\boldsymbol{\zeta}_m)\right)=\mathrm{cov}_{\ell}G.
\end{multline*}
The antisymmetric kernels are determined by the element (Proposition~15\,(a)).
\end{proof}

\begin{proof}[Proof of (e-i)]
\begin{itemize}
\item[(1)] Let $(x,x')\in\mathrm{Pzt}_k$ with $c$ and $c'$ in $\tilde{U}$.
The points $x/L=(c/L,g)$ and $x'/L=(c'/L,g')$ of $\mathbb{L}^{\mathrm{f}}_{k+1}$ are distinct, and $d_X(x/L,x'/L)=d_X(x,x')/L\le M/L$ (Section~I.3.7).
So $(x/L,x'/L)\in\mathrm{Pzt}_{k+1}$, with its photon sites in $\tilde{U}/L$, and (G2.r5) defines a homomorphism, again written $S_L$:
\begin{gather*}
S_L:\mathrm{Grc}_k(\tilde{U})\to\mathrm{Grc}_{k+1}(\tilde{U}/L),\\
\psi(x,\beta,\omega)\mapsto L^{-1}\psi(x/L,\beta,\omega),\\
\chi(x,x',\beta,\omega)\mapsto L^{-1-\alpha}\chi(x/L,x'/L,\beta,\omega).
\end{gather*}
The rule contains $L$ and $\alpha$ and does not contain $\ell$.
The generator with the index $t$ is sent to a number times the generator with the index $t/L$: the range is $0$, and there are no data.
\item[(2)] For $v\in LM\mathbb{Z}^3$,
\[
T_{v/L}S_L\psi(x,\cdot)=L^{-1}\psi(x/L+v/L,\cdot)=L^{-1}\psi\left((x+v)/L,\cdot\right)=S_LT_v\psi(x,\cdot)
\]
and
\[
T_{v/L}S_L\chi(x,x',\cdot)=L^{-1-\alpha}\chi\left((x+v)/L,(x'+v)/L,\cdot\right)=S_LT_v\chi(x,x',\cdot).
\]
Thus $T_{v/L}\circ S_L=S_L\circ T_v$.
\item[(3)] Let $\operatorname{sup-diameter}(\tilde{U})+LM<\ell$ (the margin of Section~I.1.9 for $S_L$).
Then $\operatorname{sup-diameter}(\tilde{U})+M<\ell$ and $\operatorname{sup-diameter}(\tilde{U}/L)+M<\ell/L$, so that $\mathrm{cov}_{\ell}$ on $\mathrm{Grc}_k(\tilde{U})$ and $\mathrm{cov}_{\ell/L}$ on $\mathrm{Grc}_{k+1}(\tilde{U}/L)$ are defined.
On the torus the rule (G2.r5) holds at every point $x$ of $\mathbb{L}^{\mathrm{f}}_k$: for $w\in\mathbb{Z}^3$ one has $\psi(x+\ell w,\cdot)=(-1)^{w_1+w_2+w_3}\psi(x,\cdot)$ at level $k$ and $\psi\left(x/L+(\ell/L)w,\cdot\right)=(-1)^{w_1+w_2+w_3}\psi(x/L,\cdot)$ at level $k+1$ (I.AP), and for $\chi$ both points are shifted.
Hence
\begin{gather*}
S_L\,\mathrm{cov}_{\ell}\,\psi(x,\cdot)=L^{-1}\psi(x/L,\cdot)=\mathrm{cov}_{\ell/L}\,S_L\,\psi(x,\cdot),\\
S_L\,\mathrm{cov}_{\ell}\,\chi(x,x',\cdot)=L^{-1-\alpha}\chi(x/L,x'/L,\cdot)=\mathrm{cov}_{\ell/L}\,S_L\,\chi(x,x',\cdot).
\end{gather*}
As both sides are homomorphisms, $\mathrm{cov}_{\ell/L}\circ S_L=S_L\circ\mathrm{cov}_{\ell}$ on $\mathrm{Grc}_k(\tilde{U})$.
\end{itemize}
\end{proof}

\emph{The sets of data and the maps.}
The following are objects of the covering lattice; none of them contains a side $\ell$.
Let $\tilde{X}$ be a non-empty finite union of $M$-cubes of $\mathbb{L}_k$ that is connected through faces, and let $\tilde{U}:=\tilde{X}^{+\eta_k}$.
One has
\[
\int_{\tilde{X}}\mathrm{d}x=\frac{\eta_k^3}{48}\sum_{c\in\tilde{X}}\;\sum_{g\in O_{\mathrm{cub}}},\qquad\operatorname{Vol}(\tilde{X})=M^3\cdot|\tilde{X}|_M,
\]
and $d_M(\tilde{X})$ is formed with nets in $\mathbb{R}^3$ (Proposition~16.3\,(f1)).
The generators of $\mathrm{Grc}_k(\tilde{X})$ are among those of $\mathrm{Grc}_k(\tilde{U})$; hence $\mathrm{Grc}_k(\tilde{X})$ is a subalgebra of $\mathrm{Grc}_k(\tilde{U})$.

\emph{The sets of data.}
The set $\mathrm{Dta}(\tilde{X})$ consists of the triples $(\mathrm{Kz},F,\tilde{B})$ of the following objects.
First, $\mathrm{Kz}=(\mathrm{Kz}_{00},\mathrm{Kz}_{20})$, where $\mathrm{Kz}_{00}$ is a complex number and $\mathrm{Kz}_{20}(x,x')$ is a function of the fermion sites $x$, $x'$ over $\tilde{X}$ with values in the complex $4\times4$ matrices; one puts $T_v\mathrm{Kz}:=(\mathrm{Kz}_{00},T_v\mathrm{Kz}_{20})$.
Second, $F$ is an element of $\mathrm{Grc}_k(\tilde{U})$ such that $F_{20}(\xi_1,\xi_2)=0$ unless $\xi_1$ and $\xi_2$ lie over $\tilde{X}$.
Third, $\tilde{B}$ is a complex function on the bonds of $\mathbb{L}_k$ with both ends in $\tilde{X}$, odd under reversal.
\begin{itemize}
\item[(Dt1)] For $v\in M\mathbb{Z}^3$ the map $(\mathrm{Kz},F,\tilde{B})\mapsto(T_v\mathrm{Kz},T_vF,T_v\tilde{B})$ is a bijection of $\mathrm{Dta}(\tilde{X})$ onto $\mathrm{Dta}(\tilde{X}+v)$.
Indeed, $(\tilde{X}+v)^{+\eta_k}=\tilde{U}+v$; moreover $(T_vF)_{20}(\xi_1,\xi_2)=F_{20}(\xi_1-v,\xi_2-v)$, and $\xi-v$ lies over $\tilde{X}$ if and only if $\xi$ lies over $\tilde{X}+v$; the inverse map is formed with $T_{-v}$.
\item[(Dt2)] For $(\mathrm{Kz},F,\tilde{B})\in\mathrm{Dta}(\tilde{X})$ and every complex $4\times4$ matrix $\mathrm{Mx}$, the elements
\[
F^{[2,0]}\qquad\text{and}\qquad F^{[2,0]}-\int_{\tilde{X}}\mathrm{d}x\,\bar{\psi}(x)^T\,\mathrm{Mx}\,\psi(x)
\]
lie in the subalgebra of $\mathrm{Grc}_k(\tilde{X})$ generated by the $\psi(\xi)$.
Indeed,
\[
F^{[2,0]}=\frac{1}{2}\sum_{\xi_1,\xi_2\text{ over }\tilde{U}}\left(\frac{\eta_k^3}{48}\right)^2\cdot F_{20}(\xi_1,\xi_2)\,\psi(\xi_1)\,\psi(\xi_2),
\]
and the terms in which an index does not lie over $\tilde{X}$ vanish.
\end{itemize}

\emph{The homomorphism $\boldsymbol{\Theta}_{\tilde{X},\tilde{B}}$} (for $d_M(\tilde{X})\le L$ and $\tilde{B}$ as in $\mathrm{Dta}(\tilde{X})$).
It is the homomorphism given by the display (3.3) of Proposition~17.1\,(d), formed on the covering lattice with the following objects: the shortest chains $(\square_0,\dots,\square_n)$ of $M$-cubes of $\tilde{X}$ in which consecutive cubes have a common face; their centres $z_i$; the chain points $p_0=x=(c,g)$, $p_{i+1}=(z_i,g)$ ($0\le i\le n$), $p_{n+2}=\tilde{x}=(\tilde{c},\tilde{g})$; the transports $U_{\tilde{B}}(p_i,p_{i+1})$; the numbers $d_X(p_i,p_{i+1})^{\alpha}$; the mean $\operatorname{Mn}$ over the chains; and
\[
\operatorname{Vol}(\tilde{X})^{-1}\int_{\tilde{X}}\mathrm{d}\tilde{x}=\operatorname{Vol}(\tilde{X})^{-1}\cdot\frac{\eta_k^3}{48}\sum_{\tilde{c}\in\tilde{X}}\;\sum_{\tilde{g}\in O_{\mathrm{cub}}}.
\]
The set $\tilde{X}$ belongs to $\mathrm{Sm}(\tilde{\square})$ for each of its $M$-cubes $\tilde{\square}$.
By Proposition~16.3\,(f1) it is therefore, for every side in $\mathrm{Cl}$, the lift of a small polymer, and Proposition~17.1\,(c), (c1) hold for it.
Thus the set of the chains is finite and not empty, and every pair $(p_i,p_{i+1})$ with $p_i\ne p_{i+1}$ lies in $\mathrm{Pzt}_k$ and is of kind (I) for $\tilde{X}$, that is, both its photon sites and its six paths lie in $\tilde{X}$.
(One has $|c-z_0|<M/2$, $|z_{i-1}-z_i|=M$ and $|z_n-\tilde{c}|<M/2$; a rectilinear path between two sites of one $M$-cube stays in this cube; and the six paths between the centres of two $M$-cubes with a common face are the straight segment, which lies in the two cubes.)
Hence the $\chi(p_i,p_{i+1},\cdot)$ are generators of $\mathrm{Grc}_k(\tilde{X})$, and the transports read $\tilde{B}$ on bonds with both ends in $\tilde{X}$.
The right sides of (3.3) are odd elements of $\mathrm{Grc}_k(\tilde{X})$.
It follows that there is exactly one homomorphism $\boldsymbol{\Theta}_{\tilde{X},\tilde{B}}$ of the subalgebra of $\mathrm{Grc}_k(\tilde{X})$ generated by the $\psi(\xi)$ into $\mathrm{Grc}_k(\tilde{X})$ that satisfies (3.3); this subalgebra is its domain.

\emph{The maps.}
For $(\mathrm{Kz},F,\tilde{B})\in\mathrm{Dta}(\tilde{X})$ put
\begin{gather*}
\mathrm{Phr}(\tilde{X};\mathrm{Kz},F):=
\begin{cases}
F-\mathrm{Kz}_{00}-\displaystyle\int_{\tilde{X}}\mathrm{d}x\,\bar{\psi}(x)^T\left[\operatorname{Vol}(\tilde{X})^{-1}\int_{\tilde{X}}\mathrm{d}x'\int_{\tilde{X}}\mathrm{d}x''\,\mathrm{Kz}_{20}(x',x'')\right]\psi(x) & \text{if }d_M(\tilde{X})\le L,\\
F & \text{otherwise};
\end{cases}\\
\mathrm{Pha}(\tilde{X};\mathrm{Kz},F,\tilde{B}):=
\begin{cases}
F^{r}-(F^{r})^{[2,0]}+\boldsymbol{\Theta}_{\tilde{X},\tilde{B}}\left((F^{r})^{[2,0]}\right)\ \text{with}\ F^{r}:=\mathrm{Phr}(\tilde{X};\mathrm{Kz},F) & \text{if }d_M(\tilde{X})\le L,\\
F & \text{otherwise}.
\end{cases}
\end{gather*}
Let $d_M(\tilde{X})\le L$.
The number $\mathrm{Kz}_{00}$ has the bidegree $(0,0)$; hence
\[
(F^{r})^{[2,0]}=F^{[2,0]}-\int_{\tilde{X}}\mathrm{d}x\,\bar{\psi}(x)^T\left[\operatorname{Vol}(\tilde{X})^{-1}\int_{\tilde{X}}\mathrm{d}x'\int_{\tilde{X}}\mathrm{d}x''\,\mathrm{Kz}_{20}(x',x'')\right]\psi(x),
\]
which lies in the domain of $\boldsymbol{\Theta}_{\tilde{X},\tilde{B}}$ by (Dt2).
So $\mathrm{Phr}(\tilde{X};\cdot)$ and $\mathrm{Pha}(\tilde{X};\cdot)$ are maps of $\mathrm{Dta}(\tilde{X})$ into $\mathrm{Grc}_k(\tilde{U})$, defined on all of $\mathrm{Dta}(\tilde{X})$: the pairs $(\tilde{X},\text{data})$ of Section~I.1.9 are the pairs $\left(\tilde{X},(\mathrm{Kz},F,\tilde{B})\right)$ with $(\mathrm{Kz},F,\tilde{B})\in\mathrm{Dta}(\tilde{X})$.
The two maps contain $\eta_k$, $M$, $L$, $\alpha$, $e_k$ and the number $48$; they do not contain $\ell$.

\begin{proof}[Proof of (e-ii)]
Let $X$ and $\tilde{X}$ be as in (e-ii).
By Section~I.1.1 the lift $\tilde{X}$ is a union of $M$-cubes of $\mathbb{L}_k$ with the same common faces as $X$.
Thus $\tilde{X}$ is a non-empty finite union of $M$-cubes of $\mathbb{L}_k$ that is connected through faces, and the set $\mathrm{Dta}(\tilde{X})$ and the maps $\mathrm{Phr}(\tilde{X};\,\cdot\,)$, $\mathrm{Pha}(\tilde{X};\,\cdot\,)$ are defined for it as above.
Let $\tilde{U}:=\tilde{X}^{+\eta_k}$; then
\[
\operatorname{sup-diameter}(\tilde{U})+M=\operatorname{sup-diameter}(\tilde{X})+2\eta_k+M<\ell .
\]
Put $B:=\mathscr{A}_L$; by (Fact~1) it lies in $\tilde{R}_k$.

\emph{The data of the scheme.}
Let
\[
\mathrm{Kz}_{00}:=E_{00}(X;0),\qquad
\mathrm{Kz}_{20}(x,x'):=E_{20}(X;0;x,x')\quad\left(x,x'\text{ over }\tilde{X}\right),\qquad
\mathrm{Kz}:=(\mathrm{Kz}_{00},\mathrm{Kz}_{20}),
\]
the values of $\mathrm{Kz}_{20}$ being $4\times4$ matrices in the sense of Section~I.5.1.
Let $F$ be the element of $\mathrm{Grc}_k(\tilde{U})$ whose antisymmetric kernels are the functions $E_{nm}(X;B;\xi;\boldsymbol{\zeta})$ of the indices $\xi$, $\boldsymbol{\zeta}$ over $\tilde{U}$.
Let $\tilde{B}$ be the $\ell\mathbb{Z}^3$-periodic extension of $B$, taken on the bonds both of whose ends lie in $\tilde{X}$.
Then
\begin{equation}\tag{Dt3}
(\mathrm{Kz},F,\tilde{B})\in\mathrm{Dta}(\tilde{X}).
\end{equation}
Indeed, by (k2)$^{\mathrm{o}}$ the kernel $F_{20}(\xi_1,\xi_2)=E_{20}(X;B;\xi_1,\xi_2)$ vanishes unless $\xi_1$ and $\xi_2$ lie over $X$.
If the image of a site $c$ of $\tilde{U}$ lies in $X$, then $c$ is congruent modulo $\ell\mathbb{Z}^3$ to a site of $\tilde{X}$;
both sites belong to $\tilde{U}$, on which the covering map is injective, and therefore $c\in\tilde{X}$.
Among the properties of the family $E$, this argument rests on (k2)$^{\mathrm{o}}$ alone.

Moreover, $E(X;B)=\mathrm{cov}_\ell\left(F\right)$.
By (k2)$^{\mathrm{o}}$ the kernels of $E(X;B)$ vanish unless all $\xi$ lie over $X$ and all pairs are admissible for $X$.
The two photon sites of an admissible pair lie in the image of $\tilde{U}$:
for a pair of kind (I) they lie in $X$;
for a pair of kind (II) one has $x'=x^{(3)}$, one of the sites $c$, $c_3$ lies in $X$, and the other one is a nearest neighbour of it.
Moreover their distance is at most $M$.
By Section~I.1.1 (lifts, with $\mathrm{sd}=M$, applied to the set $\tilde{U}$) such a pair is the image of a pair of $\mathrm{Pzt}_k$ both of whose points lie over $\tilde{U}$.
Consequently the kernels of $E(X;B)$ vanish at every argument that is not congruent to an index over $\tilde{U}$, and (K) yields the claim.

\emph{The first map.}
By Proposition~48.2\,(f), that is, by clause (f2) of Proposition~16.3 with $E(X;0)$ and $E(X;\mathscr{A})$ in the places of $\mathscr{E}(X;0,0)$ and $\mathscr{E}(X;\mathscr{A},\mathscr{C})$, taken at the connection $B$, one has
\[
E^{r}(X;B)=\mathrm{cov}_\ell\left(\mathrm{Phr}(\tilde{X};\mathrm{Kz},F)\right),
\]
and for all data in $\mathrm{Dta}(\tilde{X})$
\[
\mathrm{Phr}(\tilde{X}+v;T_v\mathrm{Kz},T_vF)=T_v\,\mathrm{Phr}(\tilde{X};\mathrm{Kz},F)\qquad\left(v\in LM\mathbb{Z}^3\right).
\]
The following facts lie behind these two relations.
For every matrix $\mathrm{Mx}$ the map $\mathrm{cov}_\ell$ sends $\int_{\tilde{X}}\mathrm{d}x\,\bar{\psi}(x)^T\cdot\mathrm{Mx}\cdot\psi(x)$ to $\int_{X}\mathrm{d}x\,\bar{\psi}(x)^T\cdot\mathrm{Mx}\cdot\psi(x)$:
the covering map is a bijection of $\tilde{X}\times O_{\mathrm{cub}}$ onto $X\times O_{\mathrm{cub}}$ that preserves the weights $\eta_k^3/48$, and every term contains two generators at one site, so that the signs of (I.AP) cancel.
Further $\operatorname{Vol}(\tilde{X})=\operatorname{Vol}(X)$, and
\[
\mathrm{al}_0(E,X)\cdot\operatorname{Vol}(X)=\mathrm{Kz}_{00},\qquad
\mathrm{al}_2(E,X)=\operatorname{Vol}(\tilde{X})^{-1}\int_{\tilde{X}}\mathrm{d}x'\int_{\tilde{X}}\mathrm{d}x''\,\mathrm{Kz}_{20}(x',x''),
\]
all arguments of $E_{20}(X;0;\,\cdot\,,\,\cdot\,)$ lying over the lift (the data of Proposition~48.2; (I.AP) in Section~I.3.6).
Finally, $T_v$ sends $\int_{\tilde{X}}\mathrm{d}x\,\bar{\psi}^T\cdot\mathrm{Mx}\cdot\psi$ to $\int_{\tilde{X}+v}\mathrm{d}x\,\bar{\psi}^T\cdot\mathrm{Mx}\cdot\psi$, and
\begin{gather*}
\int_{\tilde{X}+v}\mathrm{d}x'\int_{\tilde{X}+v}\mathrm{d}x''\,\left(T_v\mathrm{Kz}_{20}\right)(x',x'')=\int_{\tilde{X}}\mathrm{d}x'\int_{\tilde{X}}\mathrm{d}x''\,\mathrm{Kz}_{20}(x',x''),\\
\operatorname{Vol}(\tilde{X}+v)=\operatorname{Vol}(\tilde{X}),\qquad d_M(\tilde{X}+v)=d_M(\tilde{X}).
\end{gather*}

\emph{Small and large polymers.}
The polymer $X$ is small if and only if $d_M(\tilde{X})\le L$.
This follows from Proposition~16.3\,(f1), which holds in the present setting by Proposition~48.2\,(f).
In detail, let $\tilde{\square}$ be an $M$-cube of $\tilde{X}$.
If $d_M(\tilde{X})\le L$, then $\tilde{X}\in\mathrm{Sm}(\tilde{\square})$, and its image $X$ is small.
If $X$ is small, then $X$ is the image of some $\tilde{X}'\in\mathrm{Sm}(\tilde{\square})$, and $\tilde{X}'$ is a lift of $X$.
Two lifts of a polymer differ by a vector $\ell w$ of $\ell\mathbb{Z}^3$ (Section~I.1.1).
The cube $\tilde{\square}$ lies in $\tilde{X}$ and in $\tilde{X}'=\tilde{X}+\ell w$, and the covering map is injective on $\tilde{X}$; hence $w=0$ and $\tilde{X}'=\tilde{X}$.

\emph{The second map.}
By Proposition~48.3\,(f), that is, by Proposition~17.2\,(f) with $B=\mathscr{A}$, applied here at the connection $B=\mathscr{A}_L$, the homomorphism $\boldsymbol{\Theta}_{X,B}$ of a small polymer $X$ is given by one formula of range $0$ with the datum $\tilde{B}$.
This means that
\[
\mathrm{cov}_\ell\circ\boldsymbol{\Theta}_{\tilde{X},\tilde{B}}=\boldsymbol{\Theta}_{X,B}\circ\mathrm{cov}_\ell
\]
on the subalgebra of $\mathrm{Grc}_k(\tilde{X})$ generated by the $\psi(\xi)$, and that for every $\tilde{B}$ as in $\mathrm{Dta}(\tilde{X})$
\[
T_v\circ\boldsymbol{\Theta}_{\tilde{X},\tilde{B}}=\boldsymbol{\Theta}_{\tilde{X}+v,T_v\tilde{B}}\circ T_v\qquad\left(v\in M\mathbb{Z}^3\right)
\]
on the same subalgebra.
The reasons are these.
In Proposition~17.1 all the points $x$, $\tilde{x}$, $p_i$ are points of $\mathbb{L}^{\mathrm{f}}_k$ over the lift $\tilde{X}$.
The chains of $M$-cubes of $X$ are the images of the chains of $M$-cubes of $\tilde{X}$, since a lift has the same common faces (Section~I.1.1).
The pairs $(p_i,p_{i+1})$ are of kind (I) (Proposition~17.1\,(c1); Section~I.5.6);
hence $u_B(p_i,p_{i+1})$ is formed with the six paths between the two points in $\mathbb{L}_k$, which lie in $\tilde{X}$, that is, with $\tilde{B}$.
Therefore the display (3.3) of Proposition~17.1\,(d) for $\boldsymbol{\Theta}_{X,B}\psi(x)$ and $\boldsymbol{\Theta}_{X,B}\bar{\psi}(x)$ is the image under $\mathrm{cov}_\ell$ of the same display for $\boldsymbol{\Theta}_{\tilde{X},\tilde{B}}$.
The translation $T_v$ maps the cubes, the centres, the chains and the chain points of $(\tilde{X},x,\tilde{x})$ to those of $(\tilde{X}+v,x+v,\tilde{x}+v)$;
it leaves the labels, $d_X$ and the weights unchanged;
and $U_{T_v\tilde{B}}(p+v,p'+v)=U_{\tilde{B}}(p,p')$.

\emph{The element $E^{a}(X;B)$.}
Let $G:=\mathrm{Phr}(\tilde{X};\mathrm{Kz},F)$, so that $E^{r}(X;B)=\mathrm{cov}_\ell\left(G\right)$.
Since $\mathrm{cov}_\ell$ maps every generator to a generator of the same family, it preserves the bidegree, and $E^{r,[2,0]}(X;B)=\mathrm{cov}_\ell\left(G^{[2,0]}\right)$.
Let $X$ be small.
Then $d_M(\tilde{X})\le L$ and $G=F^{r}$.
By (Dt3) and (Dt2) the element $G^{[2,0]}=\left(F^{r}\right)^{[2,0]}$ lies in the subalgebra of $\mathrm{Grc}_k(\tilde{X})$ generated by the $\psi(\xi)$, on which the first relation of the second map holds.
The definition of $\mathrm{nat}$ (Proposition~48.3) gives
\begin{align*}
E^{a}(X;B)&=E^{r}(X;B)-E^{r,[2,0]}(X;B)+\boldsymbol{\Theta}_{X,B}\left(E^{r,[2,0]}(X;B)\right)\\
&=\mathrm{cov}_\ell\left(G-G^{[2,0]}+\boldsymbol{\Theta}_{\tilde{X},\tilde{B}}\left(G^{[2,0]}\right)\right)
=\mathrm{cov}_\ell\left(\mathrm{Pha}(\tilde{X};\mathrm{Kz},F,\tilde{B})\right).
\end{align*}
Let $X$ be large.
Then $d_M(\tilde{X})>L$, and
\[
E^{a}(X;B)=E^{r}(X;B)=E(X;B)=\mathrm{cov}_\ell\left(F\right)=\mathrm{cov}_\ell\left(\mathrm{Pha}(\tilde{X};\mathrm{Kz},F,\tilde{B})\right).
\]

\emph{Composition.}
Let $(\mathrm{Kz},F,\tilde{B})$ be an arbitrary element of $\mathrm{Dta}(\tilde{X})$, and let $v\in LM\mathbb{Z}^3$.
By (Dt1) the translated triple lies in $\mathrm{Dta}(\tilde{X}+v)$.
The map $T_v$ is an isomorphism of $\mathrm{Grc}_k(\tilde{U})$ onto $\mathrm{Grc}_k(\tilde{U}+v)$ that preserves the bidegree, and $d_M(\tilde{X}+v)=d_M(\tilde{X})$.
If $d_M(\tilde{X})\le L$, then, with $F^{r}=\mathrm{Phr}(\tilde{X};\mathrm{Kz},F)$ and the two relations for $T_v$ established above,
\begin{gather*}
T_v\,\mathrm{Pha}(\tilde{X};\mathrm{Kz},F,\tilde{B})=T_vF^{r}-\left(T_vF^{r}\right)^{[2,0]}+\boldsymbol{\Theta}_{\tilde{X}+v,T_v\tilde{B}}\left(\left(T_vF^{r}\right)^{[2,0]}\right),\\
T_vF^{r}=\mathrm{Phr}(\tilde{X}+v;T_v\mathrm{Kz},T_vF),
\end{gather*}
and the right-hand side of the first relation is $\mathrm{Pha}(\tilde{X}+v;T_v\mathrm{Kz},T_vF,T_v\tilde{B})$.
If $d_M(\tilde{X})>L$, both sides are equal to $T_vF$.
Taken together:
\begin{equation}\tag{48.11}
\begin{aligned}
E^{a}(X;\mathscr{A}_L)&=\mathrm{cov}_\ell\left(\mathrm{Pha}(\tilde{X};\mathrm{Kz},F,\tilde{B})\right)&&\text{for the data of the scheme},\\
\mathrm{Pha}(\tilde{X}+v;T_v\mathrm{Kz},T_vF,T_v\tilde{B})&=T_v\,\mathrm{Pha}(\tilde{X};\mathrm{Kz},F,\tilde{B})&&\text{for all }(\mathrm{Kz},F,\tilde{B})\in\mathrm{Dta}(\tilde{X}),\ v\in LM\mathbb{Z}^3.
\end{aligned}
\end{equation}

\emph{Kernels of level $k$.}
From (48.11) and (K) one obtains
\[
E^{a}_{nm}(X;\mathscr{A}_L;\xi;\boldsymbol{\zeta})=\mathrm{Pha}(\tilde{X};\mathrm{Kz},F,\tilde{B})_{nm}(\xi;\boldsymbol{\zeta})\qquad\text{for }\xi,\boldsymbol{\zeta}\text{ over }\tilde{U},
\]
and $E^{a}_{nm}(X;\mathscr{A}_L;\xi;\boldsymbol{\zeta})=0$ if one of the arguments is not congruent modulo $\ell\mathbb{Z}^3$ to an index over $\tilde{U}$.

\emph{Kernels of level $k+1$.}
Proposition~18.4\,(c), applied to the element $E^{a}(X;\mathscr{A}_L)$ of the algebra of the torus, together with the preceding paragraph, gives for indices $\xi$, $\boldsymbol{\zeta}$ of level $k+1$
\begin{equation}\tag{48.12}
\begin{aligned}
&\left(S_L\left[E^{a}(X;\mathscr{A}_L)\right]\right)_{nm}(\xi;\boldsymbol{\zeta})=L^{3n+6m}\cdot L^{-n-(1+\alpha)m}\cdot\mathrm{Pha}(\tilde{X};\mathrm{Kz},F,\tilde{B})_{nm}(L\xi;L\boldsymbol{\zeta})\\
&\qquad\text{if all arguments are indices over }\tilde{U}/L\text{ and all pairs of }\boldsymbol{\zeta}\text{ have }d_X\le M/L,\\[1ex]
&\left(S_L\left[E^{a}(X;\mathscr{A}_L)\right]\right)_{nm}(\xi;\boldsymbol{\zeta})=0\\
&\qquad\text{if a pair of }\boldsymbol{\zeta}\text{ has }d_X>M/L,\text{ or if an argument is not congruent}\\
&\qquad\text{modulo }(\ell/L)\mathbb{Z}^3\text{ to an index over }\tilde{U}/L.
\end{aligned}
\end{equation}
Here one uses that an index $t$ of level $k+1$, in the case of a pair with $d_X\le M/L$, is congruent modulo $(\ell/L)\mathbb{Z}^3$ to an index over $\tilde{U}/L$ if and only if the index $Lt$ of level $k$ is congruent modulo $\ell\mathbb{Z}^3$ to an index over $\tilde{U}$.
The right-hand side of (48.12) does not contain $\ell$.
By (48.11), a translation of $(\tilde{X},\mathrm{Kz},F,\tilde{B})$ by $v\in LM\mathbb{Z}^3$ translates the arguments in (48.12) by $v/L$.
If moreover $\operatorname{sup-diameter}(\tilde{X})+2\eta_k+LM<\ell$, then (48.11) and (e-i)\,(3) give
\[
S_L\left[E^{a}(X;\mathscr{A}_L)\right]=\mathrm{cov}_{\ell/L}\left(S_L\,\mathrm{Pha}(\tilde{X};\mathrm{Kz},F,\tilde{B})\right).
\]

\emph{Change of the lift.}
Every other lift of $X$ has the form $\tilde{X}+\ell w$ with $w\in\mathbb{Z}^3$ (Section~I.1.1).
The data formed with it are $T_{\ell w}\mathrm{Kz}$, $T_{\ell w}F$, $T_{\ell w}\tilde{B}$.
Indeed, the elements $E(X;0)$ and $E(X;B)$ are even, so that the sign $(-1)^{(n+m)(w_1+w_2+w_3)}$ which (I.AP) gives to a kernel with $n+m$ arguments equals $1$;
and the extension of $B$ is $\ell\mathbb{Z}^3$-periodic.
By (48.11) the value of $\mathrm{Pha}$ is replaced by its image under $T_{\ell w}$, and $\mathrm{cov}_\ell\circ T_{\ell w}=\mathrm{cov}_\ell$ on even elements (Section~I.1.9).
\end{proof}

\begin{proof}[Proof of (e-iii)]
\begin{itemize}
\item[(1)]
Under the margin of (e-ii), the kernels of $E^{a}(X;\mathscr{A}_L)$ and of $S_L\left[E^{a}(X;\mathscr{A}_L)\right]$ are given by the formulas of the paragraphs on the kernels of level $k$ and of level $k+1$, and the second relation of (48.11) holds;
all of these were derived above from that margin.
\item[(2)]
The maps $\mathrm{Phr}(\tilde{X};\,\cdot\,)$ and $\mathrm{Pha}(\tilde{X};\,\cdot\,)$ are defined on $\mathrm{Dta}(\tilde{X})$, that is, on data whose kernel $F_{20}$ vanishes unless both of its arguments lie over $\tilde{X}$.
The data of the scheme belong to $\mathrm{Dta}(\tilde{X})$ by (Dt3), which is a consequence of (k2)$^{\mathrm{o}}$.
\item[(3)]
The inequality $\operatorname{sup-diameter}(\tilde{U})+LM<\ell$ is equivalent to $\operatorname{sup-diameter}(\tilde{U}/L)+M<\ell/L$.
It is the condition under which Section~I.1.9 defines $\mathrm{cov}_{\ell/L}$ on the whole algebra $\mathrm{Grc}_{k+1}(\tilde{U}/L)$, with all its generators.
Apart from this, all the conclusions of (e-i) and (e-ii) hold under the margin of (e-ii) alone:
the formulas of (1), and, by (4), the identity of (e-i)\,(3) and the identity of the paragraph on the kernels of level $k+1$,
\begin{gather*}
\mathrm{cov}_{\ell/L}\circ S_L=S_L\circ\mathrm{cov}_\ell\quad\text{on }\mathrm{Grc}_k(\tilde{U}),\\
S_L\left[E^{a}(X;\mathscr{A}_L)\right]=\mathrm{cov}_{\ell/L}\left(S_L\,\mathrm{Pha}(\tilde{X};\mathrm{Kz},F,\tilde{B})\right),
\end{gather*}
with the homomorphism $\mathrm{cov}_{\ell/L}$ of (4).
\item[(4)]
Let $\operatorname{sup-diameter}(\tilde{U})+\eta_k<\ell$.
Then $\operatorname{sup-diameter}(\tilde{U})<\ell$ and $\operatorname{sup-diameter}(\tilde{U}/L)<\ell/L$.
The generators of $\mathrm{Grc}_k(\tilde{U})$ and of $\mathrm{Grc}_{k+1}(\tilde{U}/L)$ are independent, and the generators of the two tori are odd elements.
Hence the rule of Section~I.1.9, which sends every generator to the generator of the torus with the same index in the sense of (I.AP), defines homomorphisms
\begin{gather*}
\mathrm{cov}_\ell\colon\mathrm{Grc}_k(\tilde{U})\to\mathrm{Gr}[\psi,\chi]\quad(\text{level }k,\ \text{side }\ell),\\
\mathrm{cov}_{\ell/L}\colon\mathrm{Grc}_{k+1}(\tilde{U}/L)\to\mathrm{Gr}[\psi,\chi]\quad(\text{level }k+1,\ \text{side }\ell/L).
\end{gather*}
Under $\operatorname{sup-diameter}(\tilde{U})+M<\ell$, resp. $\operatorname{sup-diameter}(\tilde{U})+LM<\ell$, they coincide with the maps $\mathrm{cov}_\ell$, $\mathrm{cov}_{\ell/L}$ considered so far.
Two sites of $\tilde{U}$ differ by a vector of sup norm at most $\operatorname{sup-diameter}(\tilde{U})<\ell$.
Therefore distinct sites of $\tilde{U}$ are not congruent modulo $\ell\mathbb{Z}^3$, and some fundamental domain of $\ell\mathbb{Z}^3$ in $\mathbb{L}_k$ contains $\tilde{U}$.
By (I.AP) the generators of the torus whose (first) point lies over this domain are independent, and $\mathrm{cov}_\ell$ maps distinct generators of $\mathrm{Grc}_k(\tilde{U})$ to distinct members of this family.
Hence $\mathrm{cov}_\ell$ maps the basis monomials of $\mathrm{Grc}_k(\tilde{U})$ to signs times distinct basis monomials, and $\mathrm{cov}_\ell$ is injective.
In the same way $\mathrm{cov}_{\ell/L}$ is injective (Section~G3.5.0, the paragraph on the map $\mathrm{cov}_\ell$).
The computation in (e-i)\,(3) requires of its margin only that the two maps be defined.
Consequently
\[
\mathrm{cov}_{\ell/L}\circ S_L=S_L\circ\mathrm{cov}_\ell\qquad\text{on }\mathrm{Grc}_k(\tilde{U}).
\]
Under the margin of (e-ii), with $\tilde{U}=\tilde{X}^{+\eta_k}$,
\[
\operatorname{sup-diameter}(\tilde{U})+\eta_k=\operatorname{sup-diameter}(\tilde{X})+3\eta_k<\operatorname{sup-diameter}(\tilde{X})+2\eta_k+M<\ell\qquad\left(\eta_k\le1<M\right);
\]
so, by (48.11), the identity $S_L\left[E^{a}(X;\mathscr{A}_L)\right]=\mathrm{cov}_{\ell/L}\left(S_L\,\mathrm{Pha}(\tilde{X};\mathrm{Kz},F,\tilde{B})\right)$ holds there, with this map $\mathrm{cov}_{\ell/L}$.
\item[(5)]
The margin with $LM$ implies the margin with $M$.
Let $X$ be small and let $\tilde{X}$ be its lift.
Then $\operatorname{sup-diameter}(\tilde{X})\le M\cdot(L+2)$ by (G2.r3), and, since $\ell\in\mathrm{Cl}$,
\[
M\cdot(L+2)+2\eta_k+LM\le(2L+2)\cdot M+2<13LM\le\ell .
\]
Let $\tilde{X}$ be admissible in the sense of Proposition~27.1, with the sides $\ell_1\le\ell_2$ and with $d_U=2L+1$ of the standing data (TT).
Then $\operatorname{sup-diameter}(\tilde{X})\le\operatorname{diam}(\tilde{X})\le\ell_1-14d_U\cdot M$, and
\[
\ell_1-14\cdot(2L+1)\cdot M+2\eta_k+LM<\ell_1\le\ell_2 .
\]
\end{itemize}
\end{proof}

\subsubsection*{3.7. Proof of (f)}

\emph{Data.}
The data are those of Proposition~18.6: the setting (S-2), and a function $\mathrm{Ng}$ on $O_{\mathrm{cub}}\times O_{\mathrm{cub}}$ with values in the complex $4\times4$ matrices.
The conditions (n1), (n2), (n3) and the objects $\mathrm{Nl}$, $\mathrm{Nav}$, $|\mathrm{Ng}|_1$, $\mathrm{Ocl}$ are those of Proposition~18.6.
For every such function, no symmetry being assumed (as in the data of Proposition~18.6), and for every $M$-cube $\square$ let
\begin{align*}
\mathrm{El}(\mathrm{Ng},\square)&:=\int_{\square}\mathrm{d}x\,\frac{1}{48}\sum_{g'}\bar{\psi}(c,g)^T\cdot\mathrm{Ng}(g,g')\cdot\psi(c,g')\\
&=\frac{\eta_k^3}{48}\sum_{c\in\square}\ \sum_{g\in O_{\mathrm{cub}}}\frac{1}{48}\sum_{g'\in O_{\mathrm{cub}}}\ \sum_{\beta,\beta'}\bar{\psi}_{\beta}(c,g)\cdot\mathrm{Ng}(g,g')_{\beta\beta'}\cdot\psi_{\beta'}(c,g'),
\end{align*}
where $x=(c,g)$.
This is an element of $\mathrm{Gr}[\psi,\chi]$ of bidegree $(2,0)$, and it is linear in $\mathrm{Ng}$.
Read as a family without slot, $E^{\mathrm{Ng}}$ is the family
\[
E^{\mathrm{Ng}}|_0(\square;\mathscr{A}):=\mathrm{El}(\mathrm{Ng},\square),\qquad
E^{\mathrm{Ng}}|_0(X;\mathscr{A}):=0\quad\text{if }|X|_M\ge2,\qquad
\text{for all }\mathscr{A}\in\tilde{R}_k .
\]
This notation agrees with (Fact~2): the values of the family $E^{\mathrm{Ng}}$ of Proposition~18.6 are these elements, at all $(\mathscr{A},\mathscr{C})$.

\paragraph{(f-0) Kernel and norm.}
For $x=(c,g)$ and $x'=(c',g')$ over one lift of $\square$ let
\[
K(x,x'):=\eta_k^{-3}\cdot\mathbf{1}[c=c'\in\square]\cdot\mathrm{Ng}(g,g');
\]
$K$ is continued antiperiodically in each argument ((I.AP), Section~I.3.6), and it vanishes at all other arguments.
With the integrals formed on a fundamental domain that contains the lift, one has for $c\in\square$
\[
\int\mathrm{d}x'\,K(x,x')\cdot\psi(x')=\frac{\eta_k^3}{48}\sum_{c'}\sum_{g'}\eta_k^{-3}\cdot\mathbf{1}[c=c']\cdot\mathrm{Ng}(g,g')\cdot\psi(c',g')=\frac{1}{48}\sum_{g'}\mathrm{Ng}(g,g')\cdot\psi(c,g').
\]
Hence $\mathrm{El}(\mathrm{Ng},\square)=\int\mathrm{d}x\int\mathrm{d}x'\,\bar{\psi}(x)^T\cdot K(x,x')\cdot\psi(x')$, that is, $K$ is the matrix kernel of $\mathrm{El}(\mathrm{Ng},\square)$ in the sense of Section~I.5.1.
The $M$-cube $\square$ with the centre $z$ contains $(M/\eta_k)^3$ photon sites.
Indeed, $M/\eta_k=ML^k$ is an odd integer, $M$ being a power of the odd number $L$;
and in each direction $\mu$ the sites $c$ with $|c_\mu-z_\mu|<M/2$ are the $ML^k$ sites $z_\mu+\eta_kj$ with $j\in\mathbb{Z}$, $|j|\le(ML^k-1)/2$.
(This is the identity $\int_{\square}\mathrm{d}x=M^3$ of Section~I.1.6.)
By (G2.0.1), that is, by Proposition~15\,(n),
\begin{align*}
\left\|\mathrm{El}(\mathrm{Ng},\square)\right\|_{\mathsf{w}}
&=\mathsf{w}_1^2\cdot\sum_{\beta,\beta'}\int\mathrm{d}x\int\mathrm{d}x'\,\left|K(x,x')_{\beta\beta'}\right|
=\mathsf{w}_1^2\cdot\sum_{c\in\square}\left(\frac{\eta_k^3}{48}\right)^2\cdot\eta_k^{-3}\cdot\sum_{g,g'}\ \sum_{\beta,\beta'}\left|\mathrm{Ng}(g,g')_{\beta\beta'}\right|\\
&=\mathsf{w}_1^2\cdot\left(\frac{M}{\eta_k}\right)^3\cdot\eta_k^3\cdot|\mathrm{Ng}|_1
=\mathsf{w}_1^2\cdot M^3\cdot|\mathrm{Ng}|_1 .
\end{align*}
Applied to a difference, this shows that $\mathrm{El}(\mathrm{Ng},\square)=\mathrm{El}(\widetilde{\mathrm{Ng}},\square)$ if and only if $\mathrm{Ng}=\widetilde{\mathrm{Ng}}$.

\paragraph{(f-1) The four maps of Section~I.4.5.}
Put
\begin{gather*}
\mathrm{Ng}^{[S]}(g,g'):=S^{-1}\cdot\mathrm{Ng}\left(r(S)g,r(S)g'\right)\cdot S\quad\left(S\in\mathrm{Sp}\right),\qquad
\mathrm{Ng}^{[C]}(g,g'):=\mathsf{C}^{-1}\cdot\mathrm{Ng}(g',g)^T\cdot\mathsf{C},\\
\mathrm{Ng}^{[K]}(g,g'):=\varkappa^{-1}\cdot\overline{\mathrm{Ng}(g,g')}\cdot\varkappa .
\end{gather*}
Then (n1) states that $\mathrm{Ng}^{[S]}=\mathrm{Ng}$ for all $S\in\mathrm{Sp}$, (n2) states that $\mathrm{Ng}^{[C]}=\mathrm{Ng}$, and (n3) states that $\mathrm{Ng}^{[K]}=\mathrm{Ng}$.
For $\tilde{g}=(S,a)\in\mathbb{G}_k$ and real functions $\Lambda$ of $c$ one has
\begin{gather*}
a_{\tilde{g}}\,\mathrm{El}(\mathrm{Ng},\tilde{g}\square)=\mathrm{El}(\mathrm{Ng}^{[S]},\square),\qquad
G_{\Lambda}\,\mathrm{El}(\mathrm{Ng},\square)=\mathrm{El}(\mathrm{Ng},\square),\\
\mathrm{Cg}\,\mathrm{El}(\mathrm{Ng},\square)=\mathrm{El}(\mathrm{Ng}^{[C]},\square),\qquad
\mathbf{K}\,\mathrm{El}(\mathrm{Ng},\square)=\mathrm{El}(\mathrm{Ng}^{[K]},\square).
\end{gather*}
\begin{proof}
In every term the two generators sit over one photon site $c$.
A change of the lift of $c$ by $\ell w$ multiplies both of them by $(-1)^{w_1+w_2+w_3}$ (I.AP); so the terms are functions of the site $c$ of the torus.

For $a_{\tilde{g}}$: with $r=r(S)$ one has $a_{\tilde{g}}\psi(x)=S\cdot\psi(\tilde{g}^{-1}x)$ and $a_{\tilde{g}}\bar{\psi}(x)=(S^{-1})^T\cdot\bar{\psi}(\tilde{g}^{-1}x)$.
Hence
\[
a_{\tilde{g}}\left[\bar{\psi}(c,g)^T\cdot\mathrm{Ng}(g,g')\cdot\psi(c,g')\right]=\bar{\psi}\left(\tilde{g}^{-1}(c,g)\right)^T\cdot S^{-1}\cdot\mathrm{Ng}(g,g')\cdot S\cdot\psi\left(\tilde{g}^{-1}(c,g')\right).
\]
The substitution $c=\tilde{g}c''$, $g=rg_1$, $g'=rg_1'$ is a bijection of $\square\times O_{\mathrm{cub}}\times O_{\mathrm{cub}}$ onto $\tilde{g}\square\times O_{\mathrm{cub}}\times O_{\mathrm{cub}}$;
indeed, by Proposition~16.1\,(a1), (a2), $\tilde{g}$ maps the $M$-cube $\square$ onto the $M$-cube $\tilde{g}\square$, and the cluster of $c''$ onto the cluster of $\tilde{g}c''$.
Moreover $\tilde{g}^{-1}(c,g)=(c'',g_1)$ and $\tilde{g}^{-1}(c,g')=(c'',g_1')$.
This gives
\[
a_{\tilde{g}}\,\mathrm{El}(\mathrm{Ng},\tilde{g}\square)=\frac{\eta_k^3}{48}\sum_{c''\in\square}\ \sum_{g_1}\frac{1}{48}\sum_{g_1'}\bar{\psi}(c'',g_1)^T\cdot S^{-1}\cdot\mathrm{Ng}(rg_1,rg_1')\cdot S\cdot\psi(c'',g_1')=\mathrm{El}(\mathrm{Ng}^{[S]},\square).
\]

For $G_{\Lambda}$: the generators $\bar{\psi}(c,g)$ and $\psi(c,g')$ receive the factors $\exp\left(-i\,e_k\Lambda(c)\right)$ and $\exp\left(i\,e_k\Lambda(c)\right)$, whose product is $1$.

For $\mathrm{Cg}$: for a matrix $\mathrm{Mx}$,
\begin{align*}
\mathrm{Cg}\left[\bar{\psi}(x)^T\cdot\mathrm{Mx}\cdot\psi(x')\right]
&=\left(-\mathsf{C}\cdot\psi(x)\right)^T\cdot\mathrm{Mx}\cdot\mathsf{C}\cdot\bar{\psi}(x')
=-\sum_{\beta,\beta'}\psi_{\beta}(x)\cdot\left(\mathsf{C}^T\cdot\mathrm{Mx}\cdot\mathsf{C}\right)_{\beta\beta'}\cdot\bar{\psi}_{\beta'}(x')\\
&=\sum_{\beta,\beta'}\bar{\psi}_{\beta'}(x')\cdot\left(\mathsf{C}^T\cdot\mathrm{Mx}\cdot\mathsf{C}\right)_{\beta\beta'}\cdot\psi_{\beta}(x)
=\bar{\psi}(x')^T\cdot\mathsf{C}^T\cdot\mathrm{Mx}^T\cdot\mathsf{C}\cdot\psi(x)\\
&=\bar{\psi}(x')^T\cdot\mathsf{C}^{-1}\cdot\mathrm{Mx}^T\cdot\mathsf{C}\cdot\psi(x),
\end{align*}
because $\mathsf{C}^T=\mathsf{C}^{-1}$ (I.c3).
With $x=(c,g)$, $x'=(c,g')$ and $\mathrm{Mx}=\mathrm{Ng}(g,g')$, and after the names $g$ and $g'$ are exchanged in the double sum over $O_{\mathrm{cub}}\times O_{\mathrm{cub}}$, this is $\mathrm{El}(\mathrm{Ng}^{[C]},\square)$.

For $\mathbf{K}$:
\[
\mathbf{K}\left[\bar{\psi}(x)^T\cdot\mathrm{Mx}\cdot\psi(x')\right]
=\left(\varkappa\cdot\bar{\psi}(x)\right)^T\cdot\overline{\mathrm{Mx}}\cdot\varkappa\cdot\psi(x')
=\bar{\psi}(x)^T\cdot\varkappa^T\cdot\overline{\mathrm{Mx}}\cdot\varkappa\cdot\psi(x')
=\bar{\psi}(x)^T\cdot\varkappa^{-1}\cdot\overline{\mathrm{Mx}}\cdot\varkappa\cdot\psi(x'),
\]
because $\varkappa^T=\varkappa=\varkappa^{-1}$ (I.c4); the weights are real.
\end{proof}

\paragraph{(f-2) Spin relations.}
Let $S\in\mathrm{Sp}$, $r=r(S)=\mathrm{diag}(r_1,r_2,r_3)$ and $v\in\mathbb{C}^3$, and let $(\mu,\nu,\mathfrak{l})$ be a cyclic permutation of $(1,2,3)$.
Then
\begin{gather*}
1_2\otimes\sigma_\mu=-i\,\gamma_\nu\cdot\gamma_{\mathfrak{l}}\quad\left(\gamma_\nu\cdot\gamma_{\mathfrak{l}}=1_2\otimes\sigma_\nu\sigma_{\mathfrak{l}}=i\,1_2\otimes\sigma_\mu\right),\qquad
i\,\Sigma_3=\gamma_1\cdot\gamma_2\cdot\gamma_3\quad\text{(I.c2)},\\
S\cdot\gamma_\mu\cdot S^{-1}=r_\mu\,\gamma_\mu\quad\text{(I.4.1)},\qquad
r_\nu\,r_{\mathfrak{l}}=\det(r)\cdot r_\mu .
\end{gather*}
Hence
\begin{equation}\tag{sp1}
S\cdot(\gamma\cdot v)\cdot S^{-1}=\gamma\cdot(rv),\qquad
S\cdot\Sigma_3\cdot S^{-1}=\det(r)\cdot\Sigma_3,\qquad
S\cdot\left((1_2\otimes\sigma)\cdot v\right)\cdot S^{-1}=\det(r)\cdot\left((1_2\otimes\sigma)\cdot(rv)\right).
\end{equation}
By (I.c3), $-\gamma_\mu^T=\mathsf{C}^{-1}\cdot\gamma_\mu\cdot\mathsf{C}$;
transposing this relation one obtains $-\gamma_\mu=\mathsf{C}^T\cdot\gamma_\mu^T\cdot(\mathsf{C}^{-1})^T=\mathsf{C}^{-1}\cdot\gamma_\mu^T\cdot\mathsf{C}$.
A transposition reverses the order of a product, and $\gamma_{\mathfrak{l}}\cdot\gamma_\nu=-\gamma_\nu\cdot\gamma_{\mathfrak{l}}$, $\gamma_3\cdot\gamma_2\cdot\gamma_1=-\gamma_1\cdot\gamma_2\cdot\gamma_3$.
Hence
\begin{equation}\tag{sp2}
\mathsf{C}^{-1}\cdot(\gamma\cdot v)^T\cdot\mathsf{C}=-\gamma\cdot v,\qquad
\mathsf{C}^{-1}\cdot\Sigma_3^T\cdot\mathsf{C}=\Sigma_3,\qquad
\mathsf{C}^{-1}\cdot\left((1_2\otimes\sigma)\cdot v\right)^T\cdot\mathsf{C}=-(1_2\otimes\sigma)\cdot v,
\end{equation}
the second relation because $(-1)^3\,\gamma_3\cdot\gamma_2\cdot\gamma_1=\gamma_1\cdot\gamma_2\cdot\gamma_3$, the third because $-i\,(-1)^2\,\gamma_{\mathfrak{l}}\cdot\gamma_\nu=i\,\gamma_\nu\cdot\gamma_{\mathfrak{l}}$.
By (I.c4), $\overline{\gamma_\mu}=\varkappa^{-1}\cdot\gamma_\mu\cdot\varkappa$, and $\varkappa^2=1$; so $\varkappa^{-1}\cdot\overline{\gamma_\mu}\cdot\varkappa=\gamma_\mu$.
Since $\overline{i}=-i$, one has for \emph{real} $v$
\begin{equation}\tag{sp3}
\varkappa^{-1}\cdot\overline{\gamma\cdot v}\cdot\varkappa=\gamma\cdot v,\qquad
\varkappa^{-1}\cdot\overline{\Sigma_3}\cdot\varkappa=-\Sigma_3,\qquad
\varkappa^{-1}\cdot\overline{(1_2\otimes\sigma)\cdot v}\cdot\varkappa=-(1_2\otimes\sigma)\cdot v .
\end{equation}

\paragraph{(f-3) Three sums over $O_{\mathrm{cub}}$.}
One has
\[
\sum_g\det(g)=0,\qquad\sum_g g=0,\qquad\mathrm{Tm}:=\sum_g\det(g)\cdot g=0 .
\]
For the first sum, $g\mapsto gs_0$ is a bijection of $O_{\mathrm{cub}}$ and $\det(gs_0)=-\det(g)$.
For the second, $g\mapsto-g$ is a bijection.
For the third, the bijection $g\mapsto gs$ with $s=s_0$ or $s=s_3$ gives $\mathrm{Tm}=-\mathrm{Tm}\cdot s$, that is, $\mathrm{Tm}\cdot(1+s)=0$;
and $1+s_0=\mathrm{diag}(2,2,0)$, $1+s_3=\mathrm{diag}(0,2,2)$.

\paragraph{The clause.}
Let (a)$^{\mathrm{o}}$--(g)$^{\mathrm{o}}$ be the statements obtained from the clauses (a)--(g) of Proposition~18.6 by the following substitutions:
$E^{\mathrm{Ng}}|_0$ for $E^{\mathrm{Ng}}$ (for every function $\mathrm{Ng}$ that occurs),
$\mathscr{A}$ for $(\mathscr{A},\mathscr{C})$,
$0$ for $(0,0)$,
$\tilde{R}_k$ and $\tilde{R}_{k+1}$ for $\tilde{R}_k\times\mathfrak{X}_k$ and $\tilde{R}_{k+1}\times\mathfrak{X}_{k+1}$,
$K_{k}^{C}$ and $K_k$ for $K^{x,C}_{k}$ and $K^{x}_{k}$,
$\|\cdot\|_k$ and $\|\cdot\|_{k+1}$ for $\|\cdot\|^{x}_{k}$ and $\|\cdot\|^{x}_{k+1}$,
(k5)$^{\mathrm{o}}$ for (k5),
and $E^{\mathrm{o}}:=\left(R\left(E^{\mathrm{Ng}}|_0\right)\right)^{\mathrm{nat}}$ for the family $E$ of Proposition~18.6\,(e);
the maps $R$, $\mathrm{nat}$ and $L\,R$ are those without slot (Propositions~48.2, 48.3 and 48.5).
Clause (f) of Proposition~48.5 consists of the seven statements (a)$^{\mathrm{o}}$--(g)$^{\mathrm{o}}$.
In (a)$^{\mathrm{o}}$--(e)$^{\mathrm{o}}$ the function $\mathrm{Ng}$ has (n1) and (n2).

\begin{proof}[{Proof of (a)$^{\mathrm{o}}$, that is, of Proposition~18.6\,(a) for $E^{\mathrm{Ng}}|_0$}]
The elements are even and balanced (one barred and one unbarred leg), and they do not contain $\mathscr{A}$; so (k1)$^{\mathrm{o}}$ holds.
The kernel is given by (f-0); it vanishes unless $x$ and $x'$ lie over $\square$, and no leg $\chi$ occurs; this is (k2)$^{\mathrm{o}}$.
(k3)$^{\mathrm{o}}$: one has $|\tilde{g}X|_M=|X|_M$ (Proposition~16.1\,(a1)); for $|X|_M\ge2$ both sides vanish;
for $X=\square$, (f-1) and (n1) give $a_{\tilde{g}}\,\mathrm{El}(\mathrm{Ng},\tilde{g}\square)=\mathrm{El}(\mathrm{Ng}^{[S]},\square)=\mathrm{El}(\mathrm{Ng},\square)$.
(k4)$^{\mathrm{o}}$ follows from (f-1).
(k5)$^{\mathrm{o}}$: by (f-1) and (n2), $\mathrm{Cg}\,\mathrm{El}(\mathrm{Ng},\square)=\mathrm{El}(\mathrm{Ng}^{[C]},\square)=\mathrm{El}(\mathrm{Ng},\square)$.
(k7)$^{\mathrm{o}}$ under (n3): $\mathbf{K}\,\mathrm{El}(\mathrm{Ng},\square)=\mathrm{El}(\mathrm{Ng}^{[K]},\square)=\mathrm{El}(\mathrm{Ng},\square)$.
(k6)$^{\mathrm{o}}$: one has $d_M(\square)=0$ (\cite[Theorem~13.2]{QEDI}, the paragraph on single cubes) and $h_{k,1}^2=e_k^{-\frac{1}{2}}$; by (f-0),
\[
\left\|E^{\mathrm{Ng}}|_0\right\|_k=\left\|\mathrm{El}(\mathrm{Ng},\square)\right\|_{h_k}=M^3\cdot e_k^{-\frac{1}{2}}\cdot|\mathrm{Ng}|_1 .
\]
The sets $\square\times O_{\mathrm{cub}}$ partition $T^{\mathrm{f}}_k$ (Section~I.1.4), and $\frac{1}{48}\sum_{g'}\mathrm{Ng}(g,g')\cdot\psi(c,g')=(\mathrm{Ocl}\,\psi)(c,g)$; therefore
\[
\sum_{\square}\mathrm{El}(\mathrm{Ng},\square)=\int\mathrm{d}x\,\bar{\psi}(x)\cdot(\mathrm{Ocl}\,\psi)(x)=\langle\bar{\psi},\mathrm{Ocl}\,\psi\rangle^{\mathrm{ch}} .
\]
\end{proof}

\begin{proof}[{Proof of (b)$^{\mathrm{o}}$, that is, of Proposition~18.6\,(b) for $E^{\mathrm{Ng}}|_0$}]
The element $\mathrm{El}(\mathrm{Ng},\square)$ has no part of bidegree $(0,0)$; hence $\mathrm{al}_0=0$.
By (f-0),
\begin{align*}
\mathrm{al}_2(E^{\mathrm{Ng}}|_0,\square)&=\operatorname{Vol}(\square)^{-1}\int\mathrm{d}x\int\mathrm{d}x'\,K(x,x')\\
&=M^{-3}\sum_{c\in\square}\left(\frac{\eta_k^3}{48}\right)^2\cdot\eta_k^{-3}\cdot\sum_{g,g'}\mathrm{Ng}(g,g')\\
&=M^{-3}\cdot\left(\frac{M}{\eta_k}\right)^3\cdot\eta_k^3\cdot\mathrm{Nav}=\mathrm{Nav}.
\end{align*}
By (n1), the map $g\mapsto r(S)g$ being a bijection of $O_{\mathrm{cub}}$,
\[
S\cdot\mathrm{Nav}\cdot S^{-1}=\left(\frac{1}{48}\right)^2\sum_{g,g'}\mathrm{Ng}\left(r(S)g,r(S)g'\right)=\mathrm{Nav}\qquad\text{for }S=\gamma_1,\gamma_2,\gamma_3,\gamma_5;
\]
by (I.Sch), $\mathrm{Nav}=\mathrm{na}\cdot1$.
Under (n3), $\overline{\mathrm{na}}\cdot1=\varkappa^{-1}\cdot\overline{\mathrm{Nav}}\cdot\varkappa=\mathrm{Nav}=\mathrm{na}\cdot1$.
A small $X\ne\square$ that contains $\square$ has $|X|_M\ge2$; so $E^{\mathrm{Ng}}|_0(X)=0$, $\mathrm{al}_0=0$ and $\mathrm{al}_2=0$ for it.
Hence $\varepsilon(E^{\mathrm{Ng}}|_0,\square)=-\mathrm{al}_0(E^{\mathrm{Ng}}|_0,\square)=0$ and $m(E^{\mathrm{Ng}}|_0,\square)=-\mathrm{Nav}=-\mathrm{na}\cdot1$;
these are $\varepsilon(E^{\mathrm{Ng}}|_0)$ and $m(E^{\mathrm{Ng}}|_0)\cdot1$ (Proposition~48.2\,(a)).
\end{proof}

\begin{proof}[{Proof of (c)$^{\mathrm{o}}$, that is, of the three relations of Proposition~18.6\,(c), for $\xi\in\mathfrak{S}(\eta_k,\ell)$}]
The clause involves no family.
One has
\[
(\mathrm{Ocl}\,J\xi)(c,g)=\frac{1}{48}\sum_{g'}\mathrm{Ng}(g,g')\cdot\xi(c)=\mathrm{Nl}(g)\cdot\xi(c),
\qquad
(J^{*}\mathrm{Ocl}\,J\xi)(c)=\frac{1}{48}\sum_g\mathrm{Nl}(g)\cdot\xi(c)=\mathrm{Nav}\cdot\xi(c)=\mathrm{na}\cdot\xi(c)
\]
by (b)$^{\mathrm{o}}$.
With $\mathrm{Phv}=1-JJ^{*}$ (I.J4) it follows that
\[
\mathrm{Phv}\,\mathrm{Ocl}\,J\xi=\mathrm{Ocl}\,J\xi-J\left(\mathrm{na}\cdot\xi\right)=\left(\mathrm{Nl}(\,\cdot\,)-\mathrm{na}\right)\xi .
\]
\end{proof}

\begin{proof}[Proof of (d)$^{\mathrm{o}}$]
Clause (d)$^{\mathrm{o}}$ is Proposition~18.6\,(d) for $E^{\mathrm{Ng}}|_0$, with
\[
\mathrm{Ng}'(g,g')=\mathrm{Ng}(g,g')-48\cdot\mathbf{1}[g=g']\cdot\mathrm{na}.
\]
One has
\[
\int_{\square}\mathrm{d}x\,\bar{\psi}(x)^T\cdot\mathrm{na}\cdot\psi(x)=\int_{\square}\mathrm{d}x\,\frac{1}{48}\sum_{g'}\bar{\psi}(c,g)^T\cdot\left(48\cdot\mathbf{1}[g=g']\cdot\mathrm{na}\right)\cdot\psi(c,g')=\mathrm{El}\left(48\cdot\mathbf{1}[g=g']\cdot\mathrm{na},\square\right).
\]
By the definition of $R$ and by (b)$^{\mathrm{o}}$,
\[
\left(R\left(E^{\mathrm{Ng}}|_0\right)\right)(\square)=\mathrm{El}(\mathrm{Ng},\square)-0-\mathrm{El}\left(48\cdot\mathbf{1}[g=g']\cdot\mathrm{na},\square\right)=\mathrm{El}(\mathrm{Ng}',\square).
\]
The function $48\cdot\mathbf{1}[g=g']\cdot\mathrm{na}\cdot1$ has (n1), since $\mathbf{1}[rg=rg']=\mathbf{1}[g=g']$ and $S\cdot1\cdot S^{-1}=1$;
it has (n2), since $\mathsf{C}^{-1}\cdot1^T\cdot\mathsf{C}=1$ and $\mathbf{1}[g'=g]=\mathbf{1}[g=g']$;
and it has (n3) if $\mathrm{na}$ is real.
The conditions (n1), (n2) are complex linear in $\mathrm{Ng}$, and (n3) is real linear.
Hence $\mathrm{Ng}'$ has (n1) and (n2), and it has (n3) if $\mathrm{Ng}$ has it (then $\mathrm{na}$ is real by (b)$^{\mathrm{o}}$).
Finally
\[
\mathrm{Nav}'=\mathrm{Nav}-\left(\frac{1}{48}\right)^2\sum_g48\cdot\mathrm{na}=\mathrm{na}-\mathrm{na}=0 .
\]
\end{proof}

\begin{proof}[{Proof of (e)$^{\mathrm{o}}$, that is, of Proposition~18.6\,(e) for $E^{\mathrm{o}}$ and $L\,R\left(E^{\mathrm{Ng}}|_0\right)$}]
The family $E^{\mathrm{Ng}}|_0$ lies in $K_{k}^{C}$ by (a)$^{\mathrm{o}}$, and $\square$ is small, since $d_M(\square)=0\le L$.
So clause (d) of Proposition~48.5 applies with $E=E^{\mathrm{Ng}}|_0$, $E^{a}=E^{\mathrm{o}}$ and $X=\square$:
\begin{gather*}
E^{\mathrm{o}}_{20}(\square;0;\,\cdot\,,\,\cdot\,)=0,\qquad
\mathrm{sup}^{a}(\square)\le20\cdot\left\|E^{\mathrm{Ng}}|_0\right\|_k=20\cdot M^3\cdot e_k^{-\frac{1}{2}}\cdot|\mathrm{Ng}|_1,\\
\left\|S_L\left[E^{\mathrm{o}}(\square;\mathscr{A}_L)\right]\right\|_{h_{k+1}}\le3\cdot L^{-\frac{13}{4}+\frac{3\varepsilon}{2}}\cdot\mathrm{sup}^{a}(\square)\le60\cdot L^{-\frac{13}{4}+\frac{3\varepsilon}{2}}\cdot M^3\cdot e_k^{-\frac{1}{2}}\cdot|\mathrm{Ng}|_1 .
\end{gather*}
By (d)$^{\mathrm{o}}$ the element $\left(R\left(E^{\mathrm{Ng}}|_0\right)\right)(\square)$ has the bidegree $(2,0)$;
by Proposition~48.3\,(b), $E^{\mathrm{o}}_{nm}(\square)=\left(R\left(E^{\mathrm{Ng}}|_0\right)\right)_{nm}(\square)=0$ for $(n,m)\notin\{(2,0),(1,1),(0,2)\}$.
Since moreover $E^{\mathrm{o}}_{20}(\square;0;\,\cdot\,,\,\cdot\,)=0$, the element $E^{\mathrm{o}}(\square;0)$ is the sum of its parts of the bidegrees $(1,1)$ and $(0,2)$, and every monomial of these parts contains a leg $\chi$.

Let $|X|_M\ge2$.
Then $E^{\mathrm{Ng}}|_0(X)=0$, and $\mathrm{al}_0=\mathrm{al}_2=0$ if $X$ is small;
so $\left(R\left(E^{\mathrm{Ng}}|_0\right)\right)(X)=0$ for small and for large $X$, and $E^{\mathrm{o}}(X)=0$, because $\boldsymbol{\Theta}(0)=0$.
Hence
\[
\left(L\,R\left(E^{\mathrm{Ng}}|_0\right)\right)(Y;\mathscr{A})=\sum_{\square:\ \bar{\square}^{L}=LY}S_L\left[E^{\mathrm{o}}(\square;\mathscr{A}_L)\right].
\]
The set $\bar{\square}^{L}$ is the one $LM$-cube that contains $\square$; so the sum is empty unless $LY$ is an $LM$-cube, that is, unless $Y$ is an $M$-cube of level $k+1$.
For such a $Y$, with $z'$ the centre of $LY$, the $LM$-cube $LY$ contains $L^3$ $M$-cubes, namely those with the centres $z'+Mi$, $|i|\le(L-1)/2$ (Fact~5), and $d_M(Y)=0$.
Therefore
\[
\exp\left(\kappa\cdot d_M(Y)\right)\cdot\left\|\left(L\,R\left(E^{\mathrm{Ng}}|_0\right)\right)(Y;\mathscr{A})\right\|_{h_{k+1}}
\le L^3\cdot60\cdot L^{-\frac{13}{4}+\frac{3\varepsilon}{2}}\cdot M^3\cdot e_k^{-\frac{1}{2}}\cdot|\mathrm{Ng}|_1
=60\cdot L^{-\frac{1}{4}+\frac{3\varepsilon}{2}}\cdot M^3\cdot e_k^{-\frac{1}{2}}\cdot|\mathrm{Ng}|_1 .
\]
The last equality of the clause is (a)$^{\mathrm{o}}$.
\end{proof}

Clause (f)$^{\mathrm{o}}$ is Proposition~18.6\,(f) for the families $E^{\mathrm{Ng}}|_0$ of the functions $\mathrm{Ng}_1$, $\mathrm{Ng}_u$, $\mathrm{Ng}_v$, $\mathrm{Ng}_0$ displayed there, with $\mathrm{co}_0$, $\mathrm{co}_1$ real and $\mathrm{cu},\mathrm{cv}\in\mathbb{R}^3$.
It includes two facts on the decompositions of $\mathrm{cu}$ and $\mathrm{cv}$:
$\hat{n}_0,\hat{n}_1,\hat{n}_2$ form a basis of $\mathbb{R}^3$, so that there are unique real numbers $\mathrm{ru}_0,\mathrm{ru}_1,\mathrm{ru}_2$ with $\mathrm{cu}=\sum_{i\le2}\mathrm{ru}_i\,\hat{n}_i$;
and the three planes orthogonal to $\hat{n}_0,\hat{n}_1,\hat{n}_2$ span $\mathbb{R}^3$, so that there are vectors $\mathrm{cv}_0,\mathrm{cv}_1,\mathrm{cv}_2\in\mathbb{R}^3$ with $\mathrm{cv}=\mathrm{cv}_0+\mathrm{cv}_1+\mathrm{cv}_2$ and $\mathrm{cv}_i\cdot\hat{n}_i=0$.
The assertions of (f)$^{\mathrm{o}}$ hold for these $\mathrm{ru}_i$ and for every such choice of the $\mathrm{cv}_i$;
hence the display of Section~I.4.6 is attained for all real $\mathrm{co}_0$, $\mathrm{co}_1$ and all $\mathrm{cu},\mathrm{cv}\in\mathbb{R}^3$.

\begin{proof}[Proof of (f)$^{\mathrm{o}}$]
\emph{The decompositions.}
Let $\mathrm{cu}=(\mathrm{cu}^1,\mathrm{cu}^2,\mathrm{cu}^3)$ and $\mathrm{cv}=(\mathrm{cv}^1,\mathrm{cv}^2,\mathrm{cv}^3)$ be the components.
Since $\hat{n}_0=e_3$, $\hat{n}_1=(e_2-e_3)/\sqrt{2}$ and $\hat{n}_2=(e_1-e_2)/\sqrt{2}$ (Section~I.1.2), one has for real numbers $\mathrm{ru}_i$
\[
\mathrm{ru}_0\,\hat{n}_0+\mathrm{ru}_1\,\hat{n}_1+\mathrm{ru}_2\,\hat{n}_2=\left(\frac{\mathrm{ru}_2}{\sqrt{2}},\ \frac{\mathrm{ru}_1-\mathrm{ru}_2}{\sqrt{2}},\ \mathrm{ru}_0-\frac{\mathrm{ru}_1}{\sqrt{2}}\right).
\]
This vector equals $\mathrm{cu}$ if and only if
\[
\mathrm{ru}_2=\sqrt{2}\cdot\mathrm{cu}^1,\qquad\mathrm{ru}_1=\sqrt{2}\cdot\left(\mathrm{cu}^1+\mathrm{cu}^2\right),\qquad\mathrm{ru}_0=\mathrm{cu}^1+\mathrm{cu}^2+\mathrm{cu}^3;
\]
the first component determines $\mathrm{ru}_2$, then the second one determines $\mathrm{ru}_1$, then the third one determines $\mathrm{ru}_0$.
So the $\mathrm{ru}_i$ exist and are unique, and $\hat{n}_0,\hat{n}_1,\hat{n}_2$ form a basis.
The vectors
\[
\mathrm{cv}_0:=\mathrm{cv}^1e_1+\mathrm{cv}^2e_2,\qquad\mathrm{cv}_1:=0,\qquad\mathrm{cv}_2:=\mathrm{cv}^3e_3
\]
have the sum $\mathrm{cv}$, and
\[
\mathrm{cv}_0\cdot\hat{n}_0=\mathrm{cv}_0\cdot e_3=0,\qquad\mathrm{cv}_1\cdot\hat{n}_1=0,\qquad\mathrm{cv}_2\cdot\hat{n}_2=\mathrm{cv}^3e_3\cdot\frac{e_1-e_2}{\sqrt{2}}=0 .
\]
So vectors $\mathrm{cv}_i$ with the required properties exist; already the planes orthogonal to $\hat{n}_0$ and to $\hat{n}_2$ span $\mathbb{R}^3$.
They are not unique: $(\mathrm{cv}^1,\mathrm{cv}^2-\mathrm{cv}^3,0)$, $\mathrm{cv}^3\cdot(0,1,1)$, $0$ is another choice.
The arguments below involve the $\mathrm{cv}_i$ only through $\mathrm{cv}_i\cdot\hat{n}_i=0$ and $\mathrm{cv}_0+\mathrm{cv}_1+\mathrm{cv}_2=\mathrm{cv}$.

\emph{Elementary relations.}
Throughout, $r=r(S)$ lies in $\{\pm1\}^3$, a subgroup of $O_{\mathrm{cub}}$.
One has $\mathbf{1}[rg'=rgs_i]=\mathbf{1}[g'=gs_i]$ and $\det(rg)=\det(r)\cdot\det(g)$.
Moreover $g'=gs_i$ if and only if $g=g's_i$, since $s_i^2=1$; in this case
\[
\det(g')=-\det(g),\qquad g'\hat{n}_i=gs_i\hat{n}_i=-g\hat{n}_i,\qquad g'\mathrm{cv}_i=gs_i\mathrm{cv}_i=g\,\mathrm{cv}_i,
\]
because $s_i=1-2\hat{n}_i\hat{n}_i^T$ and $\mathrm{cv}_i\cdot\hat{n}_i=0$.

\emph{The function $\mathrm{Ng}_1$.}
(n1): by (sp1),
\[
S\cdot\mathrm{Ng}_1(g,g')\cdot S^{-1}=48\cdot\mathbf{1}[g=g']\cdot i\,\mathrm{co}_1\cdot\det(g)\cdot\det(r)\cdot\Sigma_3=\mathrm{Ng}_1(rg,rg').
\]
(n2): by (sp2), $\mathsf{C}^{-1}\cdot\mathrm{Ng}_1(g',g)^T\cdot\mathsf{C}=48\cdot\mathbf{1}[g=g']\cdot i\,\mathrm{co}_1\cdot\det(g)\cdot\Sigma_3$.
(n3): by (sp3),
\[
\varkappa^{-1}\cdot\overline{\mathrm{Ng}_1(g,g')}\cdot\varkappa=48\cdot\mathbf{1}[g=g']\cdot(-i)\,\mathrm{co}_1\cdot\det(g)\cdot\left(-\Sigma_3\right)=\mathrm{Ng}_1(g,g').
\]
Further $\mathrm{Nl}_1(g)=\frac{1}{48}\,\mathrm{Ng}_1(g,g)=i\,\mathrm{co}_1\cdot\det(g)\cdot\Sigma_3$, and $\mathrm{Nav}=i\,\mathrm{co}_1\cdot\Sigma_3\cdot\frac{1}{48}\sum_g\det(g)=0$ by (f-3).
The matrix $\Sigma_3$ has four entries of modulus one; hence
\[
|\mathrm{Ng}_1|_1=\left(\frac{1}{48}\right)^2\cdot48\cdot48\cdot|\mathrm{co}_1|\cdot4=4\,|\mathrm{co}_1| .
\]

\emph{The function $\mathrm{Ng}_u$.}
(n1): by (sp1),
\[
S\cdot\mathrm{Ng}_u(g,g')\cdot S^{-1}=48\sum_i\mathrm{ru}_i\cdot\mathbf{1}[g'=gs_i]\cdot\left(\gamma\cdot(rg\hat{n}_i)\right)=\mathrm{Ng}_u(rg,rg').
\]
(n2): by (sp2),
\[
\mathsf{C}^{-1}\cdot\mathrm{Ng}_u(g',g)^T\cdot\mathsf{C}=48\sum_i\mathrm{ru}_i\cdot\mathbf{1}[g=g's_i]\cdot\left(-\gamma\cdot(g'\hat{n}_i)\right)=48\sum_i\mathrm{ru}_i\cdot\mathbf{1}[g'=gs_i]\cdot\left(\gamma\cdot(g\hat{n}_i)\right).
\]
(n3) follows from (sp3), the vectors $g\hat{n}_i$ and the numbers $\mathrm{ru}_i$ being real.
Further $\mathrm{Nl}_u(g)=\sum_i\mathrm{ru}_i\cdot\left(\gamma\cdot(g\hat{n}_i)\right)=\gamma\cdot(g\,\mathrm{cu})$, and $\mathrm{Nav}=\gamma\cdot\left(\frac{1}{48}\sum_g g\,\mathrm{cu}\right)=0$ by (f-3).
For $x=(c,g)$ one has $x^{(i)}=(c,gs_i)$ and $\hat{d}_i(x)=-g\hat{n}_i$ for $i\le2$ (I.1.2); so
\[
\frac{1}{48}\sum_{g'}\mathrm{Ng}_u(g,g')\cdot\psi(c,g')=-\sum_i\mathrm{ru}_i\cdot\left(\gamma\cdot\hat{d}_i(x)\right)\cdot\psi\left(x^{(i)}\right),
\]
which gives the display of $E^{\mathrm{Ng}_u}|_0(\square)$.
The vector $g\hat{n}_0$ is $\pm e_\mu$, and $\gamma_\mu$ has four entries of modulus one (I.c6);
the vectors $g\hat{n}_1$, $g\hat{n}_2$ are of the form $(\pm e_\mu\pm e_\nu)/\sqrt{2}$ with $\mu\ne\nu$, and the moduli of the entries of $(\pm\gamma_\mu\pm\gamma_\nu)/\sqrt{2}$ add up to at most $8/\sqrt{2}=4\sqrt{2}$.
For fixed $g$ the three labels $gs_i$ are distinct.
Hence
\[
|\mathrm{Ng}_u|_1\le\left(\frac{1}{48}\right)^2\cdot48\cdot48\cdot\sum_i|\mathrm{ru}_i|\cdot4\sqrt{2}=4\sqrt{2}\sum_i|\mathrm{ru}_i| .
\]

\emph{The function $\mathrm{Ng}_v$.}
Here the index of summation, which labels the elements $s_j$ and the vectors $\mathrm{cv}_j$, is denoted $j\in\{0,1,2\}$, as in the definition of $\mathrm{Ng}_v$ in Proposition~18.6\,(f);
the factor $i$ in front of $\det(g)$ and $\det(g')$ is the imaginary unit, and the elementary relations above hold with $j$ in the place of $i$.
(n1): by (sp1),
\[
S\cdot\mathrm{Ng}_v(g,g')\cdot S^{-1}=48\sum_j\mathbf{1}[g'=gs_j]\cdot i\det(g)\cdot\det(r)\cdot\left((1_2\otimes\sigma)\cdot(rg\,\mathrm{cv}_j)\right)=\mathrm{Ng}_v(rg,rg').
\]
(n2): by (sp2),
\begin{align*}
\mathsf{C}^{-1}\cdot\mathrm{Ng}_v(g',g)^T\cdot\mathsf{C}&=48\sum_j\mathbf{1}[g=g's_j]\cdot i\det(g')\cdot\left(-(1_2\otimes\sigma)\cdot(g'\mathrm{cv}_j)\right)\\
&=48\sum_j\mathbf{1}[g'=gs_j]\cdot i\det(g)\cdot\left((1_2\otimes\sigma)\cdot(g\,\mathrm{cv}_j)\right).
\end{align*}
(n3): by (sp3),
\[
\varkappa^{-1}\cdot\overline{\mathrm{Ng}_v(g,g')}\cdot\varkappa=48\sum_j\mathbf{1}[g'=gs_j]\cdot(-i)\det(g)\cdot\left(-(1_2\otimes\sigma)\cdot(g\,\mathrm{cv}_j)\right)=\mathrm{Ng}_v(g,g').
\]
Further
\[
\mathrm{Nl}_v(g)=\sum_j i\det(g)\cdot\left((1_2\otimes\sigma)\cdot(g\,\mathrm{cv}_j)\right)=i\det(g)\cdot\left((1_2\otimes\sigma)\cdot(g\,\mathrm{cv})\right),\qquad
\mathrm{Nav}=i\cdot\left((1_2\otimes\sigma)\cdot\left(\frac{1}{48}\,\mathrm{Tm}\,\mathrm{cv}\right)\right)=0
\]
by (f-3).

\emph{Conclusion.}
For the three functions $\mathrm{Ng}_1$, $\mathrm{Ng}_u$, $\mathrm{Ng}_v$ one has $\mathrm{na}=0$.
With $\mathrm{na}=0$ one has $\mathrm{Ng}'=\mathrm{Ng}$ in (d)$^{\mathrm{o}}$, and $\left(R\left(E^{\mathrm{Ng}}|_0\right)\right)(X)=0=E^{\mathrm{Ng}}|_0(X)$ for $|X|_M\ge2$ (proof of (e)$^{\mathrm{o}}$); so $R\left(E^{\mathrm{Ng}}|_0\right)=E^{\mathrm{Ng}}|_0$.
The function $\mathrm{Ng}_0$ has (n1)--(n3) by the proof of (d)$^{\mathrm{o}}$, $\mathrm{co}_0$ being real; moreover $\mathrm{Nl}_0(g)=\mathrm{co}_0$ and $\mathrm{Nav}=\mathrm{co}_0\cdot1$.
So $\mathrm{na}=\mathrm{co}_0$, and $m(E^{\mathrm{Ng}_0}|_0)=-\mathrm{co}_0$ by (b)$^{\mathrm{o}}$.
Each of the four functions $\mathrm{Ng}_0$, $\mathrm{Ng}_1$, $\mathrm{Ng}_u$, $\mathrm{Ng}_v$ has (n1)--(n3);
the conditions (n1), (n2) are complex linear in $\mathrm{Ng}$, and (n3) is real linear (proof of (d)$^{\mathrm{o}}$);
so the sum $\mathrm{Ng}_0+\mathrm{Ng}_1+\mathrm{Ng}_u+\mathrm{Ng}_v$ has (n1)--(n3).
The map $\mathrm{Ng}\mapsto\mathrm{Nl}$ is linear; this gives the last display.
The $\mathrm{ru}_i$ and the $\mathrm{cv}_i$ exist for every $\mathrm{cu}$ and every $\mathrm{cv}$ (the decompositions);
hence the display is attained for all real $\mathrm{co}_0$, $\mathrm{co}_1$ and all $\mathrm{cu},\mathrm{cv}\in\mathbb{R}^3$.
\end{proof}

\begin{proof}[Proof of (g)$^{\mathrm{o}}$]
Clause (g)$^{\mathrm{o}}$ is Proposition~18.6\,(g) for the function
\[
\mathrm{Ng}(g,g'):=48\cdot\mathbf{1}[g=g']\cdot\left(\gamma\cdot(g\,\mathrm{cu})\right)\qquad\left(\mathrm{cu}\in\mathbb{R}^3\right),
\]
whose family is $E^{\mathrm{Ng}}|_0$.
(n1): by (sp1), $S\cdot\mathrm{Ng}(g,g')\cdot S^{-1}=48\cdot\mathbf{1}[g=g']\cdot\left(\gamma\cdot(rg\,\mathrm{cu})\right)=\mathrm{Ng}(rg,rg')$.
(n3) follows from (sp3).
Further $\mathrm{Nl}(g)=\frac{1}{48}\,\mathrm{Ng}(g,g)=\gamma\cdot(g\,\mathrm{cu})$.
By (sp2), $\mathrm{Ng}^{[C]}(g,g')=48\cdot\mathbf{1}[g=g']\cdot\left(-\gamma\cdot(g\,\mathrm{cu})\right)=-\mathrm{Ng}(g,g')$.
By (I.c1), $\left(\gamma\cdot(g\,\mathrm{cu})\right)^2=|g\,\mathrm{cu}|_2^2=|\mathrm{cu}|_2^2$; so $\mathrm{Ng}\ne0$ and $\mathrm{Ng}^{[C]}\ne\mathrm{Ng}$ if $\mathrm{cu}\ne0$.
By (f-1) and (f-0),
\[
\mathrm{Cg}\left[E^{\mathrm{Ng}}|_0(\square;-\mathscr{A})\right]=\mathrm{El}(\mathrm{Ng}^{[C]},\square)\ne\mathrm{El}(\mathrm{Ng},\square)=E^{\mathrm{Ng}}|_0(\square;\mathscr{A}).
\]
\end{proof}

\paragraph{Relation to the form with slot.}
For all $(\mathscr{A},\mathscr{C})\in\tilde{R}_k\times\mathfrak{X}_k$ the value $E^{\mathrm{Ng}}(X;\mathscr{A},\mathscr{C})$ of the family of Proposition~18.6 is the element $E^{\mathrm{Ng}}|_0(X;\mathscr{A})$.
The computations (i)--(iv) of Section~3.5 involve the family $\mathscr{E}$ only through its values at the pairs $(A,0)$ and through the definitions of $R$, $\mathrm{nat}$ and $L$.
Hence they apply to the family $E^{\mathrm{Ng}}$ with slot, which has (k2) by (f-0).
Consequently the families $R\,E^{\mathrm{Ng}}$, $\left(R\,E^{\mathrm{Ng}}\right)^{\mathrm{nat}}$ and $L\,R\,E^{\mathrm{Ng}}$ of Proposition~18.6 restrict at $\mathscr{C}=0$ to $R\left(E^{\mathrm{Ng}}|_0\right)$, $E^{\mathrm{o}}$ and $L\,R\left(E^{\mathrm{Ng}}|_0\right)$.
So each of (a)$^{\mathrm{o}}$--(g)$^{\mathrm{o}}$ is the clause of Proposition~18.6 with the same letter, read at $\mathscr{C}=0$, with the same numbers.

\subsubsection*{3.7a. Proof of (g)}

\begin{proof}[Proof of Proposition~48.5\,(g)]
Let $E'$ be a family without slot of level $k$, and let $\mathscr{A}\in\tilde{R}_{k+1}$.
By (Fact~1) the bond function $\mathscr{A}_L$ lies in $\tilde{R}_k$.
Hence $E'(X;\mathscr{A}_L)$ is defined for every $X\in D_k$; here $L\,E'$ is the family of Proposition~48.2\,(e).
By (Fact~5) the set $D_k$ is finite, and every $X\in D_k$ has exactly one $Y\in D_{k+1}$ with $\bar{X}^{L}=LY$.
Since $S_L$ is linear, one obtains
\begin{align*}
\sum_{Y\in D_{k+1}}\left(L\,E'\right)(Y;\mathscr{A})
&=\sum_{Y\in D_{k+1}}\ \sum_{X:\ \bar{X}^{L}=LY}S_L\left[E'(X;\mathscr{A}_L)\right]
=\sum_{X\in D_k}S_L\left[E'(X;\mathscr{A}_L)\right]\\
&=S_L\left[\sum_{X\in D_k}E'(X;\mathscr{A}_L)\right].
\end{align*}
\end{proof}

\subsubsection*{3.8. Level 0. Uniformity}

\paragraph{Level 0.}
The setting (S-2) includes the level $k=0$.
There one has $\eta_0=1$, the charge $e_k$ is the number $q$, and $T^{\mathrm{f}}_0=T^0_0\times O_{\mathrm{cub}}$.
The weights are $\nu(\xi)=1/48$ and $\nu(\boldsymbol{\zeta})=(1/48)^2$, and
\[
\int\mathrm{d}x=\frac{1}{48}\sum_{c}\ \sum_{g\in O_{\mathrm{cub}}}.
\]
Every pair has $d_X\ge1$, and a pair inside a cluster has $d_X=1$ and carries no phase ((I.L0.14)).
The scaling $S_L$ maps $\psi(x,\,\cdot\,)$, $x=(c,g)\in T^{\mathrm{f}}_0$, to $L^{-1}\cdot\psi\left((c/L,g),\,\cdot\,\right)$ on $T^{\mathrm{f}}_1$, whose spacing is $\eta_1=1/L$.
It maps a leg at a pair inside a cluster ($d_X=1$) to a leg at a pair inside a cluster ($d_X=\eta_1$), with the factor $L^{-1-\alpha}$.
In (48.8) and (48.12) the weights enter through the quotient
\[
\left(\frac{1}{48}\right)^{n+2m}\cdot\left(\frac{L^{-3}}{48}\right)^{-(n+2m)}=L^{3n+6m},
\]
in which the number $48$ cancels.
The arguments of 3.1--3.7a involve the level only through $\eta_k$, $e_k$ and the statements on which they rest, all of which hold for every $k\ge0$.
They rest neither on the condition $k\ge1$, nor on short pairs (Section~I.1.7), of which there are none at $k=0$ since every pair has $d_X\ge1$, nor on the gain $\eta_k^{\alpha}$ of a leg inside a cluster, which is trivial at $k=0$ since $\eta_0=1$.
Hence these arguments hold without change at $k=0$.
In the recursion of Section~I.5.7 the family of level~$0$ is $E_0=0$ ((I.L0.21)), and $L\,R\,0=0$ by (a).

\paragraph{Uniformity.}
The numbers that occur in (b) and (d) are
\[
20=10\cdot2,\qquad 3,\qquad 60=3\cdot20,\qquad 61\,c_{\mathrm{rb}}=c_{\mathrm{fl}}=61\,c_X\cdot K_0,
\]
together with powers of $L$ whose exponents contain no quantity other than $\varepsilon$.
These numbers do not depend on $k$, $\ell$, $e_k$, $E$, $X$, $Y$, $\mathscr{A}$ or $\tilde{c}$.
Since neither Proposition~48.5 nor its proof involves a mass, primitives, the cutoff index $N$, the charge $e$, $\rho$ or the weights $w_i$, these numbers are independent of such quantities as well.
In (f) the bounds contain, in addition, the number
\[
M^3\cdot e_k^{-\frac{1}{2}}\cdot|\mathrm{Ng}|_1=\left\|E^{\mathrm{Ng}}|_0\right\|_k.
\]

\subsection*{5. Proof of Proposition 48.6}

The conditions on the constants that enter the proof are collected in Section~6.

\paragraph{The systems of primitives.}
The systems of primitives that belong to the data of (a) and (b) have the properties W1--W8 and W10--W17 for $\mathrm{par}$, that is, the properties 1--8 and 10--17 of \cite[Section~5.4.2]{QEDI}, read as in Section~I.5.8 and listed in Section~G4.0.4.
These systems of primitives enter the proof only through the definitions of Section~I.5.7 and of Definition~24.D, and through Propositions~13\,(d5), 24.1\,(a) and 24.1\,(b).
The property W9 is not needed.

\subsubsection*{5.1. Proof of (a), every level $k\ge0$}

\begin{proof}[Proof of Proposition~48.6\,(a)]
\emph{The data.}
The quadruple $(k,\ell,\mu,e_k)$ and the system satisfy the following:
$k\ge0$;
$\ell$ and $\ell/L$ lie in $\mathrm{Cl}$;
$\ell\ge M\cdot r_k$;
$\mu\in(0,1/L]$;
$e_k\in(0,e_0]$ and $e_{k+1}=L^{\frac{1}{2}}\cdot e_k\le e_0$;
the system has W1--W8 and W10--W17 for $\mathrm{par}$.
These are the data (ST1)$^{\mathrm{ch}}$ of Section~G3.4.0.
Proposition~24.1 is stated under them, and under them Definition~24.D defines the step map with slot and, in the reading (D7), the step map without slot.
They include the data of Proposition~13\,(d5): the data (SF+), since $(0,1/L]$ is contained in $(0,1]$, and the family (p5) of the system, which has W1, W2, W7, W12, W13 and W14 for $\mathrm{par}$.
The functions $E^{d}_k$, $E^{\mathrm{det}}_k$ and ${\delta E^{\mathrm{det}}}_k$ are formed with this family (Section~I.5.7).
All families of level $k+1$ live on the torus of side $\ell/L$.

\emph{(a.1) The pair $(m,E)$ lies in $\mathrm{Ball}_k(4)$.}
By Proposition~48.1\,(a), $E=\mathscr{E}|_0$ lies in $K^{C}_{k}$, and
\[
\|E\|_k\le\|\mathscr{E}\|^{x}_{k}\le4\,w^{\mathrm{E}}_k;
\]
moreover $|m|\le4\,w^{\mathrm{m}}_k$.
For $k\ge1$ this is the definition of $\mathrm{Ball}_k(4)$ (Section~I.5.4).
At $k=0$ one has
\[
\mathrm{Ball}^{x}_{0}(4)=\mathrm{Ball}_0(4)=\left\{(m,0):\ |m|\le4\,w^{\mathrm{m}}_0\right\};
\]
then $\mathscr{E}=0$, so that $E=0|_0=0$, and $(m,E)=(m,0)$ lies in $\mathrm{Ball}_0(4)$ (see 5.2).
By (D1) and (D7) of Definition~24.D the data of the step map without slot are the elements of $\mathrm{Ball}_k(4)$.
By Proposition~24.1\,(b) the map $E^{*}_k$ is defined on $\mathrm{Ball}_k(4)$ and takes its values in $K^{C}_{k+1}$.

\emph{(a.2) The extraction numbers.}
Proposition~48.2\,(a), in its assertion on the families $\mathscr{E}\in K^{x,C}_k$, gives $\varepsilon(\mathscr{E})=\varepsilon(\mathscr{E}|_0)$ and $m(\mathscr{E})=m(\mathscr{E}|_0)$.
Since $E=\mathscr{E}|_0$, this means $\varepsilon({E})=\varepsilon(\mathscr{E})$ and $m({E})=m(\mathscr{E})$.

\emph{(a.3) The determinant.}
Let $\mathscr{A}\in\tilde{R}_k$.
By (Fact~1), $0$ lies in $\mathfrak{X}_k$ and in $\tilde{R}_k$.
Since $\tilde{R}_k$ is contained in $\tilde{G}_k$ \cite[Lemma~17.4(a)]{QEDI}, both $\mathscr{A}=\mathscr{A}+0$ and $0$ lie in $\tilde{G}_k$, which is the domain of $E^{d}_k(X,\,\cdot\,)$ (data of Proposition~13).
The two definitions of Section~I.5.7 therefore give
\begin{equation}\tag{48.13}
{\delta E^{\mathrm{det}}}_k(X;\mathscr{A},0)=E^{d}_k(X,\mathscr{A}+0)-E^{d}_k(X,0)=E^{\mathrm{det}}_k(X,\mathscr{A}).
\end{equation}

\emph{(a.4) The three terms.}
Let $Y\in D_{k+1}$ and $\mathscr{A}\in\tilde{R}_{k+1}$.
Then $(\mathscr{A},0)$ lies in $\tilde{R}_{k+1}\times\mathfrak{X}_{k+1}$, $\mathscr{A}_L$ lies in $\tilde{R}_k$, and $0_L=0$ (Fact~1).

\noindent\emph{First term.}
By Proposition~48.5\,(c),
\[
\left(L\,R\,\mathscr{E}\right)|_0=L\,R\,\left(\mathscr{E}|_0\right)=L\,R\,E.
\]

\noindent\emph{Second term.}
The values ${\delta E^{\mathrm{det}}}_k(X;\,\cdot\,,\,\cdot\,)$ and $E^{\mathrm{det}}_k(X,\,\cdot\,)$ are numbers, that is, elements of bidegree $(0,0)$, on which $S_L$ acts as the identity.
The definition of $L$ (Section~I.5.6), together with (48.13) at the point $\mathscr{A}_L$ of $\tilde{R}_k$, gives
\begin{align*}
\left(L\,{\delta E^{\mathrm{det}}}_k\right)\big|_0(Y;\mathscr{A})
&=\left(L\,{\delta E^{\mathrm{det}}}_k\right)(Y;\mathscr{A},0)
=\sum_{X:\ \bar{X}^{L}=LY}{\delta E^{\mathrm{det}}}_k(X;\mathscr{A}_L,0)\\
&=\sum_{X:\ \bar{X}^{L}=LY}E^{\mathrm{det}}_k(X,\mathscr{A}_L)
=\left(L\,E^{\mathrm{det}}_k\right)(Y;\mathscr{A}).
\end{align*}

\noindent\emph{Third term.}
Proposition~24.1\,(b), which holds for every $(m,\mathscr{E})\in\mathrm{Ball}^{x}_{k}(4)$, gives
\[
E^{*,x}_k(m,\mathscr{E})|_0=E^{*}_k\left(m,\mathscr{E}|_0\right)=E^{*}_k(m,E).
\]

\emph{(a.5) The sum.}
The restriction to the value $0$ of the second slot is additive (Fact~2).
Consequently
\begin{align*}
\left(L\,R\,\mathscr{E}+L\,{\delta E^{\mathrm{det}}}_k+E^{*,x}_k(m,\mathscr{E})\right)\big|_0
&=\left(L\,R\,\mathscr{E}\right)|_0+\left(L\,{\delta E^{\mathrm{det}}}_k\right)\big|_0+E^{*,x}_k(m,\mathscr{E})|_0\\
&=L\,R\,E+L\,E^{\mathrm{det}}_k+E^{*}_k(m,E).
\end{align*}
Each of the three terms on the right lies in $K^{C}_{k+1}=K^{C}_{k+1}(\ell/L)$:
$L\,R\,E$ by Proposition~48.5\,(a);
$L\,E^{\mathrm{det}}_k$, which even lies in $K_{k+1}$, by Proposition~13\,(d5);
$E^{*}_k(m,E)$ by Proposition~24.1\,(b).
The family that is restricted on the left lies in $K^{x,C}_{k+1}$:
$L\,R\,\mathscr{E}$ by Proposition~18.5\,(a),
$L\,{\delta E^{\mathrm{det}}}_k$ by Proposition~13\,(d5),
and $E^{*,x}_k(m,\mathscr{E})$ by Proposition~24.1\,(a).
\end{proof}

\subsubsection*{5.2. (a) at level 0}

At $k=0$ one has
\[
\mathrm{Ball}^{x}_{0}(4)=\mathrm{Ball}_0(4)=\left\{(m,0):\ |m|\le4\,w^{\mathrm{m}}_0\right\}
\]
(Section~I.5.4), where $w^{\mathrm{m}}_0$ is formed with $q$.
Hence $\mathscr{E}=0$ and $E=0|_0=0$, and $(m,E)=(m,0)$ lies in $\mathrm{Ball}_0(4)$.

All kernels of the families $\mathscr{E}=0$ and $E=0$ vanish.
Therefore $\mathrm{al}_0=0$ and $\mathrm{al}_2=0$ for every small $X$, and
\begin{gather*}
\varepsilon(\mathscr{E})=\varepsilon({E})=0,\qquad m(\mathscr{E})=m({E})=0,\qquad R\,\mathscr{E}=0,\qquad R\,E=0,\\
\left(R\,\mathscr{E}\right)^{\mathrm{nat}}=0,\qquad \left(R\,E\right)^{\mathrm{nat}}=0,\qquad L\,R\,\mathscr{E}=0,\qquad L\,R\,E=0;
\end{gather*}
for the adjustment one uses that $\boldsymbol{\Theta}$ is a homomorphism, so that $\boldsymbol{\Theta}(0)=0$, and for the last two identities that $S_L\,0=0$.

For $\mathscr{A}\in\tilde{R}_0$ the definitions of Section~I.5.7 give
\[
{\delta E^{\mathrm{det}}}_0(X;\mathscr{A},0)=E^{d}_0(X,\mathscr{A}+0)-E^{d}_0(X,0)=E^{\mathrm{det}}_0(X,\mathscr{A});
\]
here $\mathscr{A}$ and $0$ lie in $\tilde{G}_0$ by \cite[Lemma~17.4(a)]{QEDI} and (Fact~1).
Let $Y\in D_1$ (side $\ell/L$) and $\mathscr{A}\in\tilde{R}_1$.
Then $\mathscr{A}_L$ lies in $\tilde{R}_0$ and $0_L=0$ (Fact~1), and
\begin{align*}
\left(L\,{\delta E^{\mathrm{det}}}_0\right)\big|_0(Y;\mathscr{A})
&=\sum_{X\in D_0:\ \bar{X}^{L}=LY}{\delta E^{\mathrm{det}}}_0(X;\mathscr{A}_L,0)\\
&=\sum_{X\in D_0:\ \bar{X}^{L}=LY}E^{\mathrm{det}}_0(X,\mathscr{A}_L)
=\left(L\,E^{\mathrm{det}}_0\right)(Y;\mathscr{A}).
\end{align*}
Proposition~24.1\,(b) at $k=0$, applied to $(m,0)\in\mathrm{Ball}^{x}_{0}(4)$, gives $E^{*,x}_0(m,0)|_0=E^{*}_0(m,0)$.
By (Fact~2) it follows that
\[
\left(L\,{\delta E^{\mathrm{det}}}_0+E^{*,x}_0(m,0)\right)\big|_0=L\,E^{\mathrm{det}}_0+E^{*}_0(m,0),
\]
and both terms on the right lie in $K^{C}_{1}(\ell/L)$, by Proposition~13\,(d5) and Proposition~24.1\,(b) at $k=0$.

\begin{named}{Remark}
At level $0$ the objects of this relation are formed as follows.
The function $E^{d}_0$ is formed with $\acute{E}_0$, whose traces extend over $\mathfrak{S}^{\mathrm{ch}}(1,\ell)$, with $192$ components per photon site (Section~I.5.7 and (I.L0.26)).
The activities $E^{*,x}_0(m,0)$ and $E^{*}_0(m,0)$ are formed by (D4) of Definition~24.D at $k=0$, with the generators $W$ on $T^{\mathrm{f}}_0$ of weight $1/48$.
\end{named}

\subsubsection*{5.3. The scales of the flow}

\begin{named}{Lemma S}
Assume (S-4).
For every integer $k$ with $0\le k<k_*(R)$ the following hold:
$(k,\ell_k,e_k)$ are data (S-2), with $\ell_k/L=\ell_{k+1}$ and $L^{\frac{1}{2}}\cdot e_k=e_{k+1}$;
$\ell_k\ge M\cdot r_k$;
$0<\mu_k\le1/L$ and $L\cdot\mu_k=\mu_{k+1}$.
\end{named}

\begin{proof}
By Section~I.6.1 one has $a_{k+1}=L\cdot a_k$, and hence
\[
\ell_{k+1}=\frac{R}{a_{k+1}}=\frac{\ell_k}{L},\qquad
e_{k+1}=e\cdot a_{k+1}^{\frac{1}{2}}=L^{\frac{1}{2}}\cdot e_k,\qquad
\mu_{k+1}=\bar{m}\cdot a_{k+1}=L\cdot\mu_k.
\]
In \cite[Section~19.2]{QEDI} one reads $h'=h$ and $m_0=0$, and the upper end of the mass range is taken equal to $1$ (Section~I.6.1); the number $e_1$ is that of the line (C9).
The datum is $e_K=e\cdot h^{\frac{1}{2}}\le e_1$, so that $v_K\ge-\log e_1$.
By the second assertion of Proposition~14.3\,(f), the inequalities (e1)--(e8) hold for all $v\ge-\log e_1$; in particular (e1) holds at $v=v_K$.
With this, (Sc6) of Section~G4.0.2, that is, Proposition~14.2\,(c), (f), gives
\[
v_K\ge8,\qquad
\left(1+\frac{\mathrm{lg}}{v_K}\right)^{r}\le\left(1+\frac{3\,\mathrm{lg}}{v_K}\right)^{r}<L,
\]
and the condition (2.1) of Section~G4.0.3, namely $r_*>31L+16$, for $R\in\mathrm{Tor}(e)$ and $N\ge N_{\min}(R,e)$.
So the hypotheses of \cite[Proposition~19.1]{QEDI} are satisfied, together with the additional hypothesis of its clause (f).
By \cite[Proposition~19.1(c)]{QEDI}, $1\le k_*(R)\le K$.
By \cite[Proposition~19.1(f)]{QEDI}, $\ell_i$ lies in $\mathrm{Cl}$ for all $i\le k_*(R)$.
By \cite[Proposition~19.1(e)]{QEDI}, for $i\le k_*(R)$,
\[
\ell_i\ge M\cdot r_i,\qquad
e_i\le e_*\le e_K,\qquad
0<\mu_i\le\mu_*\le\mu_K\le1.
\]
Moreover $e_K\le e_1\le e_0$, by (S-4) and (C9).
Now let $k<k_*(R)$.
Then $k$ and $k+1$ are both at most $k_*(R)$.
Consequently $\ell_k$ and $\ell_k/L=\ell_{k+1}$ lie in $\mathrm{Cl}$;
$0<e_k\le e_0$ and $e_{k+1}\le e_0$;
$\ell_k\ge M\cdot r_k$;
and $\mu_k=\mu_{k+1}/L\le1/L$.
\end{proof}

\begin{named}{Remark}
The inequalities $v_K\ge8$ and $\left(1+\frac{3\,\mathrm{lg}}{v_K}\right)^{r}<L$, and the condition (2.1) of Section~G4.0.3, follow from the row (e1) of Proposition~14.2\,(b), which gives, at $v=v_K$ and at $v=v_*\ge v_K$,
\[
v\ge8,\qquad v^{r}>31L+16,\qquad\left(1+\frac{3\,\mathrm{lg}}{v}\right)^{r}<L;
\]
for (2.1) one uses, in addition, $r_*\ge v_*^{r}$.
\end{named}

\subsubsection*{5.4. Proof of (b)}

\begin{proof}[Proof of Proposition~48.6\,(b)]
\emph{The hypothesis.}
At the level $k$ the condition $(m_k,\mathscr{E}_k)\in\mathrm{Ball}^{x}_{k}(4)$ is a condition on the entries $m_k$ and $\mathscr{E}_k$.
Claim$(k)$\,(i) below establishes the clause in the data of (b) according to which the entries of level $k$, among them $m_k$ and $\mathscr{E}_k$, are defined as soon as $(m_j,\mathscr{E}_j)\in\mathrm{Ball}^{x}_{j}(4)$ holds for all $j<k$:
Claim$(0)$ holds without any condition, and for $j<k_{\mathrm{up}}$ Claim$(j)$ together with the single condition $(m_j,\mathscr{E}_j)\in\mathrm{Ball}^{x}_{j}(4)$ implies Claim$(j+1)$.
Hence the conditions of (b) can be imposed successively for $k=0,1,\dots$, and the hypothesis that the trajectory exists up to level $k_{\mathrm{up}}$ is well posed.

\emph{Claim$(k)$}, for $0\le k\le k_{\mathrm{up}}$:
\begin{itemize}
\item[(i)] $\varepsilon^{\mathrm{v},t}_k$ and $m_k$ are defined numbers;
$E_k$ is a defined family without slot and $\mathscr{E}_k$ a defined family with slot, both of level $k$ on $T_k(R)$;
$\mathfrak{P}_k=\left\{\mathfrak{P}_k(X,\,\cdot\,):\ X\in D_k\right\}$ is a defined family of complex functions on $\mathfrak{X}_k$;
\item[(ii)] $\varepsilon^{\mathrm{v},t}_k$ and $m_k$ are real;
\item[(iii)] $\mathscr{E}_k$ lies in $K^{x}_{k}=K^{x}_{k}(\ell_k)$;
\item[(iv)] $E_k=\mathscr{E}_k|_0$, and $E_k$ lies in $K_k=K_k(\ell_k)$;
\item[(v)] $\mathfrak{P}_k(X,0)=0$ for all $X\in D_k$.
\end{itemize}

\emph{Claim$(0)$.}
The number $\varepsilon^{\mathrm{v},t}_0=\varepsilon^{\mathrm{v}}_{\mathrm{in}}$ is real, and so is the entry $m_k$ at $k=0$, which is $m_{\mathrm{in}}$.
The family $\mathscr{E}_0=0$ lies in the real vector space $K^{x}_{0}$, and $E_0=0=0|_0$ lies in the space $K_k$ at $k=0$.
Finally $\mathfrak{P}_0=0$.

\emph{Claim$(k)$ implies Claim$(k+1)$, and the assertions for $k<k_{\mathrm{up}}$.}
Let $0\le k<k_{\mathrm{up}}$, and let Claim$(k)$ hold.
By (i), $m_k$ and $\mathscr{E}_k$ are defined, and $(m_k,\mathscr{E}_k)$ lies in $\mathrm{Ball}^{x}_{k}(4)$ by the hypothesis; at $k=0$ this means $|m_{\mathrm{in}}|\le4\,w^{\mathrm{m}}_0$, with $\mathscr{E}_0=0$.
Since $k<k_{\mathrm{up}}\le k_*(R)$, Lemma~S applies at $k$, and a system of primitives of level $k$ on $T_k(R)$ with W1--W8 and W10--W17 for $\mathrm{par}$ is given.
So the data of (a) are present, with $\ell=\ell_k$, $\mu=\mu_k$, $(m,\mathscr{E})=(m_k,\mathscr{E}_k)$ and, by (iv), $E=\mathscr{E}_k|_0=E_k$.
The level $k+1$ has the side $\ell_{k+1}$, the charge $e_{k+1}$ and the mass $\mu_{k+1}$.
By (a),
\begin{equation}\tag{48.14}
\begin{aligned}
&(m_k,E_k)\in\mathrm{Ball}_k(4),\qquad \varepsilon(E_k)=\varepsilon(\mathscr{E}_k),\qquad m(E_k)=m(\mathscr{E}_k),\\
&\left(L\,R\,\mathscr{E}_k+L\,{\delta E^{\mathrm{det}}}_k+E^{*,x}_k(m_k,\mathscr{E}_k)\right)\big|_0=L\,R\,E_k+L\,E^{\mathrm{det}}_k+E^{*}_k(m_k,E_k).
\end{aligned}
\end{equation}

\emph{(i) at $k+1$.}
The right-hand sides of the five equations of the recursion (Section~I.5.7) are defined, for the following reasons.
\begin{itemize}
\item[--] $\varepsilon(E_k)$ and $m(E_k)$: $E_k$ lies in $K^{C}_{k}$ (Proposition~48.2).
\item[--] $\delta\varepsilon^{0}_{k}=-n_k^{-1}\log\mu_I(\mathrm{Box}_k)$ (\cite[Section~5.2.5]{QEDI}; a photon object, Section~I.5.7): one has $0<\mu_I(\mathrm{Box}_k)\le1$, as shown below.
\item[--] $L\,R\,E_k$: Proposition~48.5.
\item[--] $L\,E^{\mathrm{det}}_k$ and $L\,{\delta E^{\mathrm{det}}}_k$: Proposition~13\,(d5).
\item[--] $E^{*}_k(m_k,E_k)$: $(m_k,E_k)$ lies in $\mathrm{Ball}_k(4)$ (Proposition~24.1\,(b)).
\item[--] $L\,R\,\mathscr{E}_k$: Proposition~18.5.
\item[--] $E^{*,x}_k(m_k,\mathscr{E}_k)$: $(m_k,\mathscr{E}_k)$ lies in $\mathrm{Ball}^{x}_{k}(4)$ (Proposition~24.1\,(a)).
\end{itemize}
The inequalities $0<\mu_I(\mathrm{Box}_k)\le1$ hold for the following reason.
By the proof of Lemma~S and by (C9), $e_k\le e_K\le e_1\le1/3$; hence $v_k=-\log e_k>0$ and $p_{0,k}=v_k^{p_0}>0$ (Section~I.6.1).
The measure $\mu_I$ is the standard Gaussian measure on $V_k$, a probability measure with a positive density.
The set
\[
\mathrm{Box}_k=\left\{\zeta\in V_k:\ |\zeta(\Upsilon)|\le p_{0,k}\ \text{for all }\Upsilon\right\}
\]
contains the open set $\left\{\zeta\in V_k:\ |\zeta(\Upsilon)|<p_{0,k}\ \text{for all }\Upsilon\right\}$, and this open set contains $0$.
Let further $\mathscr{C}\in\mathfrak{X}_{k+1}$.
Then $\mathscr{C}_L$ lies in $\mathfrak{X}_k$ (Proposition~18.1\,(b)), the domain of $\mathfrak{P}_k(X,\,\cdot\,)$; and $\mathfrak{X}_k=0+\mathfrak{X}_k$ is contained in $\tilde{G}_k$ ((Fact~1); \cite[Lemma~17.4(a)]{QEDI}), the domain of $E^{\mathrm{det}}_k(X,\,\cdot\,)$.
The sum over $X$ is finite (Fact~5).
So $\mathfrak{P}_{k+1}(Y,\,\cdot\,)$ is a defined function on $\mathfrak{X}_{k+1}$.

\emph{(ii) at $k+1$.}
Since $E_k$ lies in $K_k$ by (iv), the numbers $\varepsilon(E_k)$ and $m(E_k)$ are real (Proposition~48.2\,(a)).
The number $\delta\varepsilon^{0}_{k}$ is real, because $0<\mu_I(\mathrm{Box}_k)\le1$ by (i) at $k+1$; this also follows from the last sentence of Proposition~24.1\,(a).
Together with (ii) at $k$ it follows that
\[
\varepsilon^{\mathrm{v},t}_{k+1}=L^3\left(\varepsilon^{\mathrm{v},t}_k+\varepsilon(E_k)+\delta\varepsilon^{0}_{k}\right)
\qquad\text{and}\qquad
m_{k+1}=L\left(m_k+m(E_k)\right)
\]
are real.

\emph{(iii) at $k+1$.}
The family $\mathscr{E}_k$ lies in $K^{x}_{k}$, and $m_k$ is real.
One has $L\,R\,\mathscr{E}_k\in K^{x}_{k+1}$ by Proposition~18.5\,(a), $L\,{\delta E^{\mathrm{det}}}_k\in K^{x}_{k+1}$ by Proposition~13\,(d5), and $E^{*,x}_k(m_k,\mathscr{E}_k)\in K^{x}_{k+1}$ by Proposition~24.1\,(a), applied to a real $m$ and to $\mathscr{E}\in K^{x}_{k}$.
Since $K^{x}_{k+1}$ is a real vector space, it contains
\[
\mathscr{E}_{k+1}=L\,R\,\mathscr{E}_k+L\,{\delta E^{\mathrm{det}}}_k+E^{*,x}_k(m_k,\mathscr{E}_k).
\]

\emph{(iv) at $k+1$.}
The fourth and the third equation of the recursion, together with (48.14), give
\[
\mathscr{E}_{k+1}|_0=L\,R\,E_k+L\,E^{\mathrm{det}}_k+E^{*}_k(m_k,E_k)=E_{k+1}.
\]
By (iii) at $k+1$ and by Proposition~48.1\,(a) at the level $k+1$, the family $\mathscr{E}_{k+1}|_0$ lies in $K_{k+1}$.

\emph{(v) at $k+1$.}
By (Fact~1), $0$ lies in $\mathfrak{X}_{k+1}$ and $0_L=0$.
By Section~I.5.7, $E^{\mathrm{det}}_k(X,0)=E^{d}_k(X,0)-E^{d}_k(X,0)=0$.
With (v) at $k$ one obtains
\[
\mathfrak{P}_{k+1}(Y,0)=\sum_{X:\ \bar{X}^{L}=LY}\left(\mathfrak{P}_k(X,0)+E^{\mathrm{det}}_k(X,0)\right)=0.
\]

This proves Claim$(k)$ for all $k\le k_{\mathrm{up}}$, and (48.14) for all $k<k_{\mathrm{up}}$.

\emph{The first step, $k=0$} (if $k_{\mathrm{up}}\ge1$).
Here $\mathscr{E}_0=0$, $E_0=0$ and $\mathfrak{P}_0=0$; moreover $\varepsilon(E_0)=m(E_0)=0$, $L\,R\,E_0=0$ and $L\,R\,\mathscr{E}_0=0$ (by 5.2).
The five equations of the recursion become
\begin{gather*}
\varepsilon^{\mathrm{v},t}_1=L^3\left(\varepsilon^{\mathrm{v}}_{\mathrm{in}}+\delta\varepsilon^{0}_{0}\right),\qquad
m_1=L\cdot m_{\mathrm{in}},\qquad
E_1=L\,E^{\mathrm{det}}_0+E^{*}_0(m_{\mathrm{in}},0),\\
\mathscr{E}_1=L\,{\delta E^{\mathrm{det}}}_0+E^{*,x}_0(m_{\mathrm{in}},0),\qquad
\mathfrak{P}_1(Y,\mathscr{C})=\sum_{X:\ \bar{X}^{L}=LY}E^{\mathrm{det}}_0(X,\mathscr{C}_L),
\end{gather*}
and $E_1=\mathscr{E}_1|_0$ is the relation obtained in 5.2.

\emph{The last sentence of (b).}
Let $k<k_{\mathrm{up}}$, and put $(\varepsilon^{\mathrm{v}},m,\mathscr{E}):=(\varepsilon^{\mathrm{v},t}_k,m_k,\mathscr{E}_k)$.
By (8.6)$^{\mathrm{ch}}$ of Definition~24.D and by (48.14),
\begin{align*}
\varepsilon^{\mathrm{v},+}&=L^3\left(\varepsilon^{\mathrm{v},t}_k+\varepsilon(\mathscr{E}_k)+\delta\varepsilon^{0}_{k}\right)=L^3\left(\varepsilon^{\mathrm{v},t}_k+\varepsilon(E_k)+\delta\varepsilon^{0}_{k}\right)=\varepsilon^{\mathrm{v},t}_{k+1},\\
m^{+}&=L\left(m_k+m(\mathscr{E}_k)\right)=L\left(m_k+m(E_k)\right)=m_{k+1},\\
\mathscr{E}^{+}&=\mathscr{E}_{k+1},\qquad \mathscr{E}^{+}|_0=E_{k+1}.
\end{align*}

Clauses (a) and (b) therefore hold for every $k$, every side, mass and charge admitted by the data, every system with the properties named and every $(m,\mathscr{E})\in\mathrm{Ball}^{x}_{k}(4)$; clause (b) holds, moreover, for every $\bar{m}$, $e$, $R$, $N$ admitted by (S-4), every real pair $(m_{\mathrm{in}},\varepsilon^{\mathrm{v}}_{\mathrm{in}})$ and every $k_{\mathrm{up}}$ as in the data.
\end{proof}

\subsection*{6. Conditions on the constants}

The following conditions on the constants are used in Sections~1, 3 and~5.
Each of them belongs to the list of Section~I.8.2 and of Proposition~14.2.

\medskip
\noindent
\begin{tabular}{p{0.2\textwidth}p{0.3\textwidth}p{0.22\textwidth}p{0.13\textwidth}}
condition & form in which it is used & where it stands & place of use\\
\hline
$c_{\mathrm{rb}}=c_X\cdot K_0$, $c_{\mathrm{fl}}=61\,c_X\cdot K_0$
& $61\,c_{\mathrm{rb}}=c_{\mathrm{fl}}$
& Section~I.8.2, (C0); Proposition~14.2\,(a)
& 3.3\\
\hline
$\varepsilon\in(1-\alpha,\bar{\varepsilon}]$, $\bar{\varepsilon}=1/200$, $\alpha\in(1-\bar{\varepsilon},1)$
& $1-\varepsilon<\alpha$; $0<\varepsilon\le1/200$ (so that $3\varepsilon/2\le1$ and $1-\varepsilon>\varepsilon/2$)
& Section~I.8.2, (C1), (C7); Proposition~14.2\,(b), (S1)
& (Fact~4), (sc1)--(sc3)\\
\hline
$\kappa\ge7\kappa_0+7$
& $\kappa\ge\kappa_0$; $\kappa-1\ge\kappa_0$
& Section~I.8.2, (C2); Proposition~14.2\,(a)
& 3.3: (48.2) at $\mathrm{kq}=\kappa$ and at $\mathrm{kq}=\kappa-1$\\
\hline
$L^{\frac{1}{8}-\frac{3\bar{\varepsilon}}{2}}\ge2400$
& $L^{\frac{1}{4}}\ge2$ (hence $\mathrm{rad}_t=L^{\frac{7}{4}}\ge2$); $L^{\frac{3}{4}}\ge110>2\,\mathrm{e}^4$; $L^3\cdot\exp(-L)\le L^{-\frac{1}{4}}$
& Section~I.8.2, (C3), first condition; Proposition~14.2\,(b), (L1); consequences in Proposition~14.2\,(f)
& 3.2\,(i), (v); 3.3\\
\hline
$L\left(\kappa-6\kappa_0-4\right)-1\ge\kappa$
& $L\left(\kappa-\kappa_0-2\right)\ge\kappa$, hence $L\left(\kappa-\kappa_0-1\right)\ge\kappa$ (note $\kappa-\kappa_0-2\ge\kappa-6\kappa_0-4$)
& Section~I.8.2, (C3), third condition; Proposition~14.2\,(b), (L3)
& 3.3\\
\hline
$L$ odd
& the $LM$-cube with the centre $z'$ is the union of the $L^3$ $M$-cubes with the centres $z'+M\cdot i$, $|i|\le(L-1)/2$
& Section~I.0.5; Section~I.8.2, (C3); Proposition~14.2\,(b), (L5)
& (Fact~5); 3.7\,(e)$^{\mathrm{o}}$\\
\hline
\end{tabular}
\medskip

\noindent
\begin{tabular}{p{0.2\textwidth}p{0.3\textwidth}p{0.22\textwidth}p{0.13\textwidth}}
condition & form in which it is used & where it stands & place of use\\
\hline
$M$ a power of $L$, $L$ odd
& $M/\eta_k=M\cdot L^k$ is an odd integer; so an $M$-cube has $(M/\eta_k)^3$ photon sites, and $\int_{\square}\mathrm{d}x=M^3$
& Sections~I.0.5, I.1.1 (faces), I.1.6; Section~I.8.2, (C5), (C3); Proposition~14.2\,(b), (M1), (L5)
& 3.7\,(f-0), (b)$^{\mathrm{o}}$\\
\hline
$M\ge L^2$
& $\eta_k\le1<M$
& Section~I.0.5; Section~I.8.2, (C5); Proposition~14.2\,(b), (M1)
& 3.6\,(e-iii)\,(4)\\
\hline
$e_1\le e_0$
& $e_i\le e_K\le e_1\le e_0$ for $i\le k_*(R)$
& Section~I.8.2, (C9); Proposition~14.2\,(f)
& 5.3\\
\hline
$e_1\le1/3$
& $e_k<1$, so that $v_k>0$, $p_{0,k}>0$ and $\mu_I(\mathrm{Box}_k)>0$
& Section~I.8.2, (C9); Proposition~14.2\,(f)
& 5.4\,(i), (ii)\\
\hline
(e1): $v\ge8$, $v^{r}>31L+16$, $\left(1+\frac{3\,\mathrm{lg}}{v}\right)^{r}<L$ for all $v\ge-\log e_1$
& through (Sc6) of Section~G4.0.2, that is, Proposition~14.2\,(c), (f): $v_K\ge8$, $\left(1+\frac{\mathrm{lg}}{v_K}\right)^{r}<L$, and (2.1) of Section~G4.0.3
& Proposition~14.2\,(b), (e1), line (C9); Proposition~14.3\,(f)
& 5.3\\
\hline
$L\ge33$ (hence $L\ge13$); $M$ a power of $L$
& through \cite[Proposition~19.1(f)]{QEDI} and the relations that precede this proposition in \cite[Section~19.2]{QEDI}
& Proposition~14.2\,(b), (L1), (M1)
& 5.3\\
\hline
\end{tabular}
\medskip

The statements on which these proofs rest are subject to further conditions:
Propositions~48.3, 48.4 and 18.1 to (A8), (A9) and (L1);
\cite[Lemma~17.4(a)]{QEDI} to (A8), $L\ge5$ and $\varepsilon\le1/200$;
Propositions~13\,(d5), 24.1\,(a) and 24.1\,(b) to the conditions named with Proposition~13 and in Section~G3.6.1.

\section*{Chapter 21a. Identities of one step, first step included: proofs of Propositions 21.1--21.4}

\subsection*{0. Setting and preliminaries}

\subsubsection*{0.0. The setting}

Throughout the chapter the data are those of the setting (S21) of Section~G3.1.0.

\subsubsection*{0.1. Conventions of this chapter}

\emph{Covering lattices.}
All functions considered below are defined on covering index sets:
on $\mathbb{L}_k$, $\mathbb{Z}^3$, $L\mathbb{Z}^3$, $\mathbb{L}^{\mathrm{f}}_k=\mathbb{L}_k\times O_{\mathrm{cub}}$ and $\mathrm{Pzt}_k$ at the level $k$,
and on $\mathbb{L}_{k+1}$, $\mathbb{Z}^3$, $\mathbb{L}^{\mathrm{f}}_{k+1}$ and $\mathrm{Pzt}_{k+1}$ at the level $k+1$.
The letter $B$ is also used for the $\ell\mathbb{Z}^3$-periodic function on the bonds of $\mathbb{L}_k$ that $B$ induces,
and $B'$ for the $(\ell/L)\mathbb{Z}^3$-periodic function on the bonds of $\mathbb{L}_{k+1}$ that $B'$ induces.
For $n\in\mathbb{Z}^3$ let
\[
\mathrm{sgl}(n):=(-1)^{|n|_1};
\]
here $|n|_1$ has the same parity as the sum of the three components of $n$.
A function $f$ satisfies $f(t+\ell e_\mu)=-f(t)$ for $\mu=1,2,3$ if and only if $f(t+\ell\cdot n)=\mathrm{sgl}(n)\cdot f(t)$ for every $n\in\mathbb{Z}^3$;
such a function is called antiperiodic (for a function of pairs, both points are translated).
A function with $f(t+\ell\cdot n)=f(t)$ for every $n\in\mathbb{Z}^3$ is called periodic.

\emph{Arithmetic of the side.}
That $\ell$ lies in $\mathrm{Cl}$ means that $\ell/(LM)$ is an odd integer $\ge13$ (Section~I.1.1),
and $M$ is a power of $L$ with $M\ge L^2$ (Section~I.0.5).
Consequently, for $0\le j\le k+1$, the number
\[
\frac{\ell}{L^j\cdot\eta_k}=\frac{\ell}{LM}\cdot M\cdot L^{1+k-j}
\]
is an integer; in other words, $\ell\mathbb{Z}^3$ is contained in $L^j\eta_k\mathbb{Z}^3$.
Further, $\ell/L=\left(\ell/(LM)\right)\cdot M$ is an integer and an odd multiple of $M$, hence a multiple of $L$, and $\ell/L\ge13M>2M$.
Finally, $L/2<M$.

\emph{Arguments of either sign.}
The operators $Q_k^{T}(B)$ and $Q^{T}(B)$ are those of Section~I.3.3.
In the same manner, $Q_{k+1}^{T}(B)$ denotes the expression by which (S21) defines $Q_{k+1}^{T}(-B)$, with $-B$ put in the place of $B$:
\[
Q_{k+1}^{T}(B)=Q_k^{T}(B)\cdot Q^{T}(B)\quad(k\ge1),\qquad Q_1^{T}(B)=Q^{T}(B)\quad(k=0).
\]

\emph{Sesquilinear pairings.}
The pairing $(\cdot,\cdot)^{\mathrm{ch}}$ on $\mathfrak{S}^{\mathrm{ch}}(\eta_k,\ell)$ and the pairings $(\cdot,\cdot)_a$ on $\mathfrak{S}(a,\ell)$, $a\in\{\eta_k,1,L\}$, are those of Section~I.1.6.

\emph{Notation.}
In this chapter $n$, $n'$, $\tilde n$ denote vectors of $\mathbb{Z}^3$, and $\mathrm{cf}$ denotes a complex number.

\subsubsection*{0.2. Blocks}

Recall that $L$ is odd.
Let $0\le j\le k+1$, and put $\mathrm{bs}_j:=L^j\eta_k=L^{j-k}$.
\begin{itemize}
\item[(a)] \emph{Each $c\in\mathbb{L}_k$ belongs to exactly one block $\mathrm{Bl}_j(z)$ with $z\in\mathrm{bs}_j\mathbb{Z}^3$.}
The centre $z$ of this block is denoted by $y_j(c)$.
\begin{proof}
Write $c=\eta_k\cdot n$ with $n\in\mathbb{Z}^3$; then $c/\mathrm{bs}_j=n/L^j$.
No component of $n/L^j$ is one half of an odd integer:
otherwise twice the corresponding component of $n$, which is even, would be the product of $L^j$ and an odd integer, which is odd.
Hence for each component of $c/\mathrm{bs}_j$ there is exactly one integer at a distance $<1/2$ from it.
If $n'\in\mathbb{Z}^3$ is the vector formed by these integers, then $z=\mathrm{bs}_j\cdot n'$ is the only point of $\mathrm{bs}_j\mathbb{Z}^3$ with $|c-z|<\mathrm{bs}_j/2$.
\end{proof}
One has $y_0(c)=c$ and $y_k(c)=[c]$, and the points $y_0,\dots,y_k$ of Section~I.3.2 are $y_0(c),\dots,y_k(c)$.
\emph{Each $y\in\mathbb{Z}^3$ belongs to exactly one block $B(y')$ with $y'\in L\mathbb{Z}^3$.}
Indeed, no component of $y/L$ is one half of an odd integer, because twice a component of $y$ is even, whereas $L$ times an odd integer is odd;
hence there is exactly one $n'\in\mathbb{Z}^3$ with $|y/L-n'|<1/2$, and $y'=L\cdot n'$.
\item[(b)] \emph{One has $\#|\mathrm{Bl}_j(z)|=L^{3j}$}:
the points of $\mathrm{Bl}_j(z)$ are the points $z+\eta_k\cdot n$ with $n\in\mathbb{Z}^3$ and $|n|<L^j/2$, that is, with $|n|\le(L^j-1)/2$,
and each component of $n$ takes $L^j$ values.
Consequently $\#|\Delta_y|=L^{3k}=\eta_k^{-3}$ and $\#|\Delta_y\times O_{\mathrm{cub}}|=48\cdot\eta_k^{-3}$, and
\[
\int_{\Delta_y}\mathrm{d}x=\eta_k^3\cdot\frac{1}{48}\cdot48\cdot\eta_k^{-3}=1.
\]
\emph{One has $\#|B(y')|=L^3$}:
the points of $B(y')$ are the points $y'+n$ with $n\in\mathbb{Z}^3$ and $|n|<L/2$, that is, with $|n|\le(L-1)/2$,
and each component of $n$ takes $L$ values.
\item[(c)] \emph{The block $\mathrm{Bl}_{k+1}(y')$ is the disjoint union of the blocks $\Delta_y$, $y\in B(y')$};
equivalently, $y_{k+1}(c)=y'$ if and only if $[c]\in B(y')$.
\begin{proof}
The vector $[c]-y'$ has integer components, and $|c-[c]|<1/2$.
If $|[c]-y'|<L/2$, then $|[c]-y'|\le(L-1)/2$, and therefore $|c-y'|<(L-1)/2+1/2=L/2$.
If $|[c]-y'|\ge L/2$, then $|[c]_\mu-y'_\mu|\ge(L+1)/2$ for some $\mu$, and therefore $|c_\mu-y'_\mu|>(L+1)/2-1/2=L/2$.
By (a) the blocks $\Delta_y$ are pairwise disjoint.
\end{proof}
(For $k=0$ one has $\Delta_y=\{y\}$ and $\mathrm{Bl}_1(y')=B(y')$.)
\item[(d)] \emph{Translations and scaling.}
If $v\in\mathrm{bs}_j\mathbb{Z}^3$, then $\mathrm{Bl}_j(z+v)=\mathrm{Bl}_j(z)+v$, and hence $y_j(c+v)=y_j(c)+v$;
if $v\in L\mathbb{Z}^3$, then $B(y'+v)=B(y')+v$.
By 0.1 these relations hold for every $v\in\ell\mathbb{Z}^3$ and every $j\le k+1$.
Division by $L$ maps $\mathbb{L}_k$ onto $\mathbb{L}_{k+1}$,
it maps $\mathrm{bs}_j\mathbb{Z}^3=L^j\eta_k\mathbb{Z}^3$ onto $L^j\eta_{k+1}\mathbb{Z}^3$,
and it maps $\{c:|c-z|<L^j\eta_k/2\}$ onto $\{\tilde c:|\tilde c-z/L|<L^j\eta_{k+1}/2\}$, which is the block $\mathrm{Bl}_j(z/L)$ of the level $k+1$.
Hence $y_j(c)/L=y_j(c/L)$ for $0\le j\le k+1$, the right side being formed at the level $k+1$.
In particular, $\mathrm{Bl}_{k+1}(y')/L$ is the unit block $\Delta_{y'/L}$ of the level $k+1$.
\end{itemize}

\subsubsection*{0.3. Paths and phases}

\begin{itemize}
\item[(a)] \emph{Paths} (Section~I.1.1).
Let $c,c'\in\mathbb{L}_k$, and let $\pi=(\pi_1,\pi_2,\pi_3)$ be an order of the three coordinates.
The path $\boldsymbol{\Gamma}^{\pi}(c,c')$ is the sequence of sites
\[
c_{(0)}=c,\ c_{(1)},\ \dots,\ c_{(\mathrm{len})}=c',\qquad \mathrm{len}:=|c-c'|_1/\eta_k,
\]
along which first the coordinate $\pi_1$ moves from $c_{\pi_1}$ to $c'_{\pi_1}$ in steps of length $\eta_k$ while the two other coordinates stay fixed,
then the coordinate $\pi_2$ moves in the same manner, and finally the coordinate $\pi_3$.
For a bond function $A$,
\[
A\left(\boldsymbol{\Gamma}^{\pi}(c,c')\right)=\sum_{j<\mathrm{len}}A\left(c_{(j)},c_{(j+1)}\right).
\]
\item[(b)] \emph{Linearity.}
Both $(\tau A)(c,c')$ and $(\tau_kA)(y,c)$ are linear forms in the numbers $A(b)$, with real coefficients.
Therefore
\[
\tau(-A)=-\tau A,\qquad \tau_k(-A)=-\tau_kA,\qquad \overline{\left(\tau\overline{A}\right)(c,c')}=(\tau A)(c,c'),\qquad \overline{\left(\tau_k\overline{A}\right)(y,c)}=(\tau_kA)(y,c).
\]
\item[(c)] \emph{Translations.}
For $v\in\mathbb{L}_k$ one has $\boldsymbol{\Gamma}^{\pi}(c+v,c'+v)=\boldsymbol{\Gamma}^{\pi}(c,c')+v$.
Hence, if $A$ is periodic and $v\in\ell\mathbb{Z}^3$, then $(\tau A)(c+v,c'+v)=(\tau A)(c,c')$ and, by 0.2\,(d), $(\tau_kA)(y+v,c+v)=(\tau_kA)(y,c)$.
It follows that $u_B$, $U_B$, $\varphi_B$ and $\varphi^{+}_B$ do not change when both of their arguments are translated by one and the same vector of $\ell\mathbb{Z}^3$.
Being exponentials, these functions have no zeros.
\item[(d)] \emph{Bonds.}
If $|c-c'|_2=\eta_k$, each of the six paths from $c$ to $c'$ is the bond $(c,c')$;
thus $(\tau A)(c,c')=A(c,c')$ and $U_B(c,c')=u_B(c,c')$.
From $B(c',c)=-B(c,c')$ one obtains $u_B(c',c)=u_B(c,c')^{-1}=u_{-B}(c,c')$.
Combined with $U_B(x,x')=U_B(\pi(x),\pi(x'))$ and $U_B(c,c)=1$ (see (I.1.6)) and with $\left(x^{(i)}\right)^{(i)}=x$ (see (I.s1)), this yields
\begin{equation}
U_B\left(x^{(i)},x\right)=U_B\left(x,x^{(i)}\right)^{-1}=U_{-B}\left(x,x^{(i)}\right)\qquad(i=0,1,2,3).
\tag{21.12}
\end{equation}
For $i\le2$ the three numbers are equal to $1$, since $\pi\left(x^{(i)}\right)=c$;
for $i=3$ they are $u_B(c_3,c)$, $u_B(c,c_3)^{-1}$ and $u_{-B}(c,c_3)$.
\item[(e)] \emph{Scaling.}
Division by $L$ maps the bonds of $\mathbb{L}_k$ onto the bonds of $\mathbb{L}_{k+1}$ (because $\eta_k/L=\eta_{k+1}$),
and it maps the path $\boldsymbol{\Gamma}^{\pi}(c,c')$, site by site, onto the path $\boldsymbol{\Gamma}^{\pi}(c/L,c'/L)$ of $\mathbb{L}_{k+1}$.
\end{itemize}

\subsubsection*{0.4. Fundamental domains, integrals}

Fix a set $\mathrm{Fdm}^L$ of representatives of $L\mathbb{Z}^3/\ell\mathbb{Z}^3$; it has $(\ell/L)^3=n_{k+1}$ points.
Let $\mathrm{Fdm}^1$ be the union of the blocks $B(y')$, $y'\in\mathrm{Fdm}^L$,
and let $\mathrm{Fdm}$ be the union of the blocks $\Delta_y$, $y\in\mathrm{Fdm}^1$ (for $k=0$, $\mathrm{Fdm}=\mathrm{Fdm}^1$).
\begin{itemize}
\item[(a)] \emph{The set $\mathrm{Fdm}^1$ is a set of representatives of $\mathbb{Z}^3/\ell\mathbb{Z}^3$, with $n_k=\ell^3$ points,
and $\mathrm{Fdm}$ is a set of representatives of $\mathbb{L}_k/\ell\mathbb{Z}^3$.}
\begin{proof}
Consider first $\mathrm{Fdm}^1$.
Let $y\in\mathbb{Z}^3$, and let $y'\in L\mathbb{Z}^3$ be the point with $y\in B(y')$ (0.2\,(a)).
Write $y'=y'_0+\ell\cdot n$ with $y'_0\in\mathrm{Fdm}^L$ and $n\in\mathbb{Z}^3$; then $y-\ell\cdot n\in B(y'_0)$ by 0.2\,(d).
Suppose that $y-\ell\cdot n\in B(y'_1)$ and $y-\ell\cdot\tilde n\in B(y'_2)$ with $y'_1,y'_2\in\mathrm{Fdm}^L$.
Then $y$ belongs to $B(y'_1+\ell\cdot n)$ and to $B(y'_2+\ell\cdot\tilde n)$, so that $y'_1+\ell\cdot n=y'_2+\ell\cdot\tilde n$ by 0.2\,(a);
hence $y'_1=y'_2$ and $n=\tilde n$.
Consider now $\mathrm{Fdm}$.
Let $c\in\mathbb{L}_k$, and write $[c]=z+\ell\cdot n$ with $z\in\mathrm{Fdm}^1$ and $n\in\mathbb{Z}^3$; then $c-\ell\cdot n\in\Delta_z$ by 0.2\,(d).
If $c-\ell\cdot n\in\Delta_z$ and $c-\ell\cdot\tilde n\in\Delta_{z'}$ with $z,z'\in\mathrm{Fdm}^1$,
then $[c]=z+\ell\cdot n=z'+\ell\cdot\tilde n$ by 0.2\,(a), (d);
since $\mathrm{Fdm}^1$ is a set of representatives, $z=z'$ and $n=\tilde n$.
\end{proof}
\item[(b)] Let $\mathrm{Fp}$ be a periodic function on $\mathbb{L}^{\mathrm{f}}_k$.
By (I.1.3),
\[
\int\mathrm{d}x\,\mathrm{Fp}(x)=\eta_k^3\cdot\frac{1}{48}\sum_{c\in\mathrm{Fdm}}\sum_{g\in O_{\mathrm{cub}}}\mathrm{Fp}(c,g),
\]
and the value is the same if $\mathrm{Fdm}\times O_{\mathrm{cub}}$ is replaced by any other set of representatives of $\mathbb{L}^{\mathrm{f}}_k/\ell\mathbb{Z}^3$,
because all sites carry the same weight $\eta_k^3/48$ and $\mathrm{Fp}$ is periodic.
The choice of $\mathrm{Fdm}$ and $\mathrm{Fdm}^1$ gives
\begin{equation}
\int\mathrm{d}x\,\mathrm{Fp}(x)=\sum_{y\in\mathrm{Fdm}^1}\int_{\Delta_y}\mathrm{d}x\,\mathrm{Fp}(x),\qquad
\sum_{y\in\mathrm{Fdm}^1}\mathrm{Gp}(y)=\sum_{y'\in\mathrm{Fdm}^L}\sum_{y\in B(y')}\mathrm{Gp}(y)
\tag{21.13}
\end{equation}
for every periodic function $\mathrm{Gp}$ on $\mathbb{Z}^3$.
\item[(c)] \emph{Re-indexing.}
For $i=0,1,2,3$ and every periodic function $\mathrm{Fp}$ on $\mathbb{L}^{\mathrm{f}}_k$,
\[
\int\mathrm{d}x\,\mathrm{Fp}\left(x^{(i)}\right)=\int\mathrm{d}x\,\mathrm{Fp}(x).
\]
\begin{proof}
By (I.s1) the map $x\mapsto x^{(i)}$ is a bijection of $\mathbb{L}^{\mathrm{f}}_k$ that commutes with the translations.
It therefore maps the set of representatives $\mathrm{Fdm}\times O_{\mathrm{cub}}$ of $\mathbb{L}^{\mathrm{f}}_k/\ell\mathbb{Z}^3$ onto a set of representatives.
All sites carry the same weight $\eta_k^3/48$, and the assertion follows from (b).
\end{proof}
\item[(d)] If $f'$ and $f$ are antiperiodic, then the products $f'(x)\cdot f(x)$ and $\overline{f'(x)}\cdot f(x)$ are periodic, because $\mathrm{sgl}(n)^2=1$.
Hence the pairings of Section~I.1.6 are sums of the kind considered in (b).
\end{itemize}

\subsubsection*{0.5. Transposes and adjoints}

The spaces that enter are $\mathfrak{S}^{\mathrm{ch}}(\eta_k,\ell)$, $\mathfrak{S}(\eta_k,\ell)$, $\mathfrak{S}(1,\ell)$ and $\mathfrak{S}(L,\ell)$,
with the pairings $\langle\cdot,\cdot\rangle^{\mathrm{ch}}$, $\langle\cdot,\cdot\rangle_{\eta_k}$, $\langle\cdot,\cdot\rangle_1$ and $\langle\cdot,\cdot\rangle_L$,
together with the corresponding spaces of the level $k+1$.
Each of these pairings is bilinear and symmetric, and it is non-degenerate,
since $\langle\overline{f},f\rangle$ is a sum of the numbers $|f(t)|_2^2$ with positive weights.
Hence for every linear map $O\colon V\to V'$ between two of these spaces there is exactly one linear map $O^T\colon V'\to V$ such that
\[
\left\langle f'',Of'\right\rangle_{V'}=\left\langle O^Tf'',f'\right\rangle_V\qquad\text{for all }f''\in V',\ f'\in V,
\]
and exactly one linear map $O^*\colon V'\to V$ such that $\left(f'',Of'\right)_{V'}=\left(O^*f'',f'\right)_V$ for all such $f''$, $f'$.
The definition and the symmetry of the pairings give the following rules.
\begin{itemize}
\item[(a)] For linear maps $O$, $O_1$, $O_2$ and complex numbers $\mathrm{cf}$,
\begin{gather*}
\left(O_1+O_2\right)^T=O_1^T+O_2^T,\qquad \left(\mathrm{cf}\cdot O\right)^T=\mathrm{cf}\cdot O^T,\qquad \left(O_1\cdot O_2\right)^T=O_2^T\cdot O_1^T,\qquad \left(O^T\right)^T=O,\qquad 1^T=1;\\
\left(O_1\cdot O_2\right)^*=O_2^*\cdot O_1^*,\qquad \left(O^*\right)^*=O,\qquad \left(\mathrm{cf}\cdot O\right)^*=\overline{\mathrm{cf}}\cdot O^*.
\end{gather*}
\item[(b)] The map $O$ is bijective if and only if $O^T$ is bijective, and then $\left(O^{-1}\right)^T=\left(O^T\right)^{-1}$.
Indeed, if $O\cdot O_1=1$ and $O_1\cdot O=1$, then $O_1^T\cdot O^T=1$ and $O^T\cdot O_1^T=1$ by (a);
the converse follows from $\left(O^T\right)^T=O$.
\item[(c)] Let $\mathrm{Mx}$ be a constant $4\times4$ matrix, acting by $(\mathrm{Mx}\,f)(t)=\mathrm{Mx}\cdot f(t)$.
Its transpose is the action of the transposed matrix, and its adjoint is the action of the adjoint matrix,
because $f''\cdot\left(\mathrm{Mx}\,f'\right)=\left(\mathrm{Mx}^Tf''\right)\cdot f'$ and $\mathrm{Mx}^T\,\overline{f''}=\overline{\mathrm{Mx}^*f''}$.
For a complex periodic function $\psi$ of position one has $\psi^T=\psi$ and $\psi^*=\overline{\psi}$.
\item[(d)] One has $J^T=J^*$ and $\left(J^*\right)^T=J$ (see (I.J1)).
Indeed, for $a\in\{\eta_k,1\}$, $f\in\mathfrak{S}^{\mathrm{ch}}(a,\ell)$ and $\xi\in\mathfrak{S}(a,\ell)$, the sums over $c$ extending over a set of representatives of $a\mathbb{Z}^3/\ell\mathbb{Z}^3$,
\[
\left\langle f,J\xi\right\rangle^{\mathrm{ch}}=a^3\sum_c\frac{1}{48}\sum_{g\in O_{\mathrm{cub}}}f(c,g)\cdot\xi(c)=a^3\sum_c\left(J^*f\right)(c)\cdot\xi(c)=\left\langle J^*f,\xi\right\rangle_a.
\]
With $\overline{f}$ in the place of $f$ this becomes $(f,J\xi)^{\mathrm{ch}}=a^3\sum_c\overline{\left(J^*f\right)(c)}\cdot\xi(c)=\left(J^*f,\xi\right)_a$;
so the map $J^*$ of (I.3.1) is the adjoint of $J$.
Moreover $J^*\cdot J=1$, because $\left(J^*J\xi\right)(c)=\frac{1}{48}\sum_{g\in O_{\mathrm{cub}}}\xi(c)=\xi(c)$.
\item[(e)] \emph{Scalar operators.}
A linear map $O$ between two of these spaces is called scalar if
\[
(Of)(t)=\sum_{t'}o(t,t')\cdot f(t')
\]
with complex numbers $o(t,t')$, only finitely many of which are different from zero for each $t$; here $t$ and $t'$ range over the covering index sets.
Sums and products of scalar operators are scalar.
A scalar operator commutes with every constant $4\times4$ matrix $\mathrm{Mx}$, since
\[
\left(O\cdot\mathrm{Mx}\,f\right)(t)=\sum_{t'}o(t,t')\cdot\mathrm{Mx}\cdot f(t')=\mathrm{Mx}\cdot(Of)(t).
\]
Its kernel in the sense of (I.1.4) is a multiple of $1_4$ at every pair of arguments:
one collects the terms at $t'+\ell\cdot n$, $n\in\mathbb{Z}^3$, with the signs $\mathrm{sgl}(n)$, and divides by the weight of $t'$.
If a scalar operator is bijective, its inverse commutes with every $\mathrm{Mx}$.
\end{itemize}

\subsubsection*{0.6. Fields}

Let $\mathrm{Gr}^{-}$ be the odd part of $\mathrm{Gr}$, and let $(\mathrm{ob}_j)_j$ be a basis of the complex vector space $\mathrm{Gr}^{-}$.
A map $\mathrm{ff}$ from a covering index set (spinor index included) to $\mathrm{Gr}^{-}$ can be written in exactly one way as $\mathrm{ff}(t)=\sum_jf_j(t)\cdot\mathrm{ob}_j$ with complex functions $f_j$,
and $\mathrm{ff}$ is antiperiodic if and only if every $f_j$ is antiperiodic.
Thus a field on one of the lattices is $\mathrm{ff}=\sum_jf_j\cdot\mathrm{ob}_j$ with $f_j$ in the corresponding space $V$ of spinors,
and a field on $\mathrm{Pz}_k$, resp. on $\mathrm{Pz}_{k+1}$, is $\mathrm{ff}=\sum_jf_j\cdot\mathrm{ob}_j$ with functions of pairs $f_j$ (Section~I.1.7) of the level $k$, resp. of the level $k+1$ and the side $\ell/L$.
The words ``by linearity'' mean that $O\,\mathrm{ff}:=\sum_j\left(Of_j\right)\cdot\mathrm{ob}_j$ and $(\mathrm{ff})_L:=\sum_j\left(f_j\right)_L\cdot\mathrm{ob}_j$;
these definitions do not depend on the choice of the basis.
The pairing of two fields, with the factors in the order in which they are written, is
\begin{equation}
\left\langle\mathrm{ff}'',\mathrm{ff}'\right\rangle=\sum_t\nu(t)\sum_\beta\mathrm{ff}''_\beta(t)\cdot\mathrm{ff}'_\beta(t)=\sum_{j,j'}\left\langle f''_j,f'_{j'}\right\rangle\cdot\mathrm{ob}_j\cdot\mathrm{ob}_{j'},
\tag{21.14}
\end{equation}
where $t$ ranges over a set of representatives of the covering index set of the fields modulo the translations by $\ell\mathbb{Z}^3$ if the fields are of the level $k$ (the fields on $L\mathbb{Z}^3$ included),
and modulo the translations by $(\ell/L)\mathbb{Z}^3$ if they are of the level $k+1$, whose side is $\ell/L$;
$\nu(t)$ is the weight that Section~I.1.6, read at the level and the side of the fields, attaches to $t$.
Thus the sum over $t$ is the one by which Section~I.1.6 defines the pairing of the corresponding spaces of spinors at that level.
\begin{itemize}
\item[(a)] An identity $O_1=O_2$ between linear maps, and an identity $(Of)_L=O'f_L$, holds for fields whenever it holds for spinors.
\item[(b)] One has $\left\langle\mathrm{ff}'',O\,\mathrm{ff}'\right\rangle=\left\langle O^T\mathrm{ff}'',\mathrm{ff}'\right\rangle$:
both sides are equal to $\sum_{j,j'}\left\langle f''_j,Of'_{j'}\right\rangle\cdot\mathrm{ob}_j\cdot\mathrm{ob}_{j'}$, and no factors are interchanged.
\item[(c)] The pairing of fields is additive in each argument, and complex numbers can be taken out of each argument.
\item[(d)] By (21.14), an identity $\left\langle f''_L,f_L\right\rangle=L\cdot\left\langle f'',f\right\rangle$ between pairings holds for fields whenever it holds for spinors.
\end{itemize}

\subsubsection*{0.7. The elementary operators: explicit forms; they preserve antiperiodicity}

Let $f\in\mathfrak{S}^{\mathrm{ch}}(\eta_k,\ell)$, $v\in\mathfrak{S}(1,\ell)$ and $v'\in\mathfrak{S}(L,\ell)$, and write $x=(c,g)$.
Section~I.3.3, (I.3.1) and (I.1.3) give the following formulas.
For $k\ge1$,
\begin{equation}
\begin{split}
\left(Q_k(B)\cdot J^*f\right)(y)&=\eta_k^3\sum_{c\in\Delta_y}\varphi_B(y,c)\cdot\frac{1}{48}\sum_{g\in O_{\mathrm{cub}}}f(c,g)\\
&=\int_{\Delta_y}\mathrm{d}x\,\varphi_B(y,c)\cdot f(x),
\end{split}
\tag{21.15}
\end{equation}
\begin{equation}
\left(J\cdot Q_k^{T}(-B)\,v\right)(c,g)=\varphi_B(y,c)^{-1}\cdot v(y)\qquad\text{for }c\in\Delta_y\text{ and each of the }48\text{ labels }g.
\tag{21.16}
\end{equation}
For $k\ge0$,
\begin{equation}
\begin{aligned}
&\left(Q(B)\,v\right)(y')=L^{-3}\sum_{y\in B(y')}U_B(y',y)\cdot v(y),\\
&\left(Q^{T}(-B)\,v'\right)(y)=U_B(y',y)^{-1}\cdot v'(y')\qquad(y\in B(y')).
\end{aligned}
\tag{21.17}
\end{equation}
For $k=0$,
\begin{equation}
\begin{split}
\left(Q(B)\cdot J^*f\right)(y')&=L^{-3}\sum_{c\in B(y')}U_B(y',c)\cdot\frac{1}{48}\sum_{g\in O_{\mathrm{cub}}}f(c,g)\\
&=L^{-3}\int_{B(y')}\mathrm{d}x\,U_B(y',c)\cdot f(x),
\end{split}
\tag{21.18}
\end{equation}
\begin{equation}
\left(J\cdot Q^{T}(-B)\,v'\right)(c,g)=U_B(y',c)^{-1}\cdot v'(y')\qquad\text{for }c\in B(y')\text{ and each of the }48\text{ labels }g.
\tag{21.19}
\end{equation}
For $k=0$ the operators (21.17) between $\mathfrak{S}(1,\ell)$ and $\mathfrak{S}(L,\ell)$ enter only through the products $Q(B)\cdot J^*$ and $J\cdot Q^{T}(-B)$;
the formulas (21.18) and (21.19) are (21.17) composed with the maps (I.3.1).
The number $U_B(y',y)$ is defined, because $|y'-y|<L/2<M$ (see 0.1).

\emph{Each of the operators $D_B$, $J$, $J^*$, $Q_k(B)$, $Q_k^{T}(-B)$, $Q(B)$, $Q^{T}(-B)$ and $\delta_{\alpha,B}$ maps antiperiodic functions to antiperiodic functions.}
\begin{description}
\item[$D_B$:] the matrix $M_i(x)$ depends on the label only, $\left(x+\ell\cdot n\right)^{(i)}=x^{(i)}+\ell\cdot n$ by (I.s1),
and $U_B\left(x+\ell\cdot n,x^{(i)}+\ell\cdot n\right)=U_B\left(x,x^{(i)}\right)$ by 0.3\,(c);
hence each term of (I.2.3) at the site $x+\ell\cdot n$ is $\mathrm{sgl}(n)$ times the corresponding term at $x$.
\item[$Q_k(B)$:] for $\xi\in\mathfrak{S}(\eta_k,\ell)$ one has $\Delta_{y+\ell\cdot n}=\Delta_y+\ell\cdot n$ and $\varphi_B(y+\ell\cdot n,c+\ell\cdot n)=\varphi_B(y,c)$ by 0.2\,(d) and 0.3\,(c);
hence the sum that gives $\left(Q_k(B)\,\xi\right)(y+\ell\cdot n)$ is the sum over $c\in\Delta_y$ of $\varphi_B(y,c)\cdot\xi(c+\ell\cdot n)=\mathrm{sgl}(n)\cdot\varphi_B(y,c)\cdot\xi(c)$.
\item[$Q_k^{T}(-B)$:] one has $[c+\ell\cdot n]=[c]+\ell\cdot n$, so that the value at $c+\ell\cdot n$ is
$\varphi_B([c],c)^{-1}\cdot v\left([c]+\ell\cdot n\right)=\mathrm{sgl}(n)\cdot\varphi_B([c],c)^{-1}\cdot v([c])$.
\item[$Q(B)$:] one has $B(y'+\ell\cdot n)=B(y')+\ell\cdot n$ and $U_B(y'+\ell\cdot n,y+\ell\cdot n)=U_B(y',y)$;
hence the sum at $y'+\ell\cdot n$ is the sum over $y\in B(y')$ of $U_B(y',y)\cdot v(y+\ell\cdot n)=\mathrm{sgl}(n)\cdot U_B(y',y)\cdot v(y)$.
\item[$Q^{T}(-B)$:] the site $y+\ell\cdot n$ lies in $B(y'+\ell\cdot n)$, so that the value at $y+\ell\cdot n$ is
$U_B(y',y)^{-1}\cdot v'(y'+\ell\cdot n)=\mathrm{sgl}(n)\cdot U_B(y',y)^{-1}\cdot v'(y')$.
\item[$J$, $J^*$:] the translations act on $c$ only.
\item[$\delta_{\alpha,B}$:] neither $d_X$ nor $U_B$ changes when both points are translated by $\ell\cdot n$.
\end{description}
The operators $J$, $J^*$, $Q_k(\pm B)$, $Q_k^{T}(\pm B)$, $Q(\pm B)$ and $Q^{T}(\pm B)$ are scalar in the sense of 0.5\,(e).

\emph{The level $k+1$.}
Of the pair $(k,\ell)$, the arguments of Sections 0.2--0.5 and of the preceding part of 0.7 use only that $\ell$ is a multiple of $L^j\eta_k$ for $0\le j\le k+1$.
The number $\ell/L$ is a multiple of $M$, hence of $L$ (see 0.1), hence of $L^j\eta_{k+1}=L^{j-k-1}$ for $0\le j\le k+2$.
Therefore all of these facts hold with $(k+1,\ell/L,e_{k+1},B')$ in the place of $(k,\ell,e_k,B)$;
in particular they hold for the operators of the level $k+1$ in the list of (S21).
(These operators involve only the blocks $\mathrm{Bl}_j$ of the level $k+1$ with $j\le k+1$, whose sides are $L^{j-k-1}\le1$;
for these blocks it suffices that $\ell/L$ is an integer.)

\subsection*{1. Proposition 21.1}

\subsubsection*{Proof}

The clauses are proved in the order (b), (a), (c), (d0), (d), (e), (f).

\begin{proof}[Proof of Proposition~21.1\,(b)]
\emph{The case $k\ge1$.}
By (21.17), (21.15) and the relation $\varphi^{+}_B(y',c)=U_B(y',y)\cdot\varphi_B(y,c)$ for $c\in\Delta_y$, $y\in B(y')$ (see (I.3.3)),
\begin{align*}
\left(Q(B)\cdot Q_k(B)\cdot J^*f\right)(y')
&=L^{-3}\sum_{y\in B(y')}U_B(y',y)\cdot\int_{\Delta_y}\mathrm{d}x\,\varphi_B(y,c)\cdot f(x)\\
&=L^{-3}\sum_{y\in B(y')}\int_{\Delta_y}\mathrm{d}x\,\varphi^{+}_B(y',c)\cdot f(x)\\
&=L^{-3}\int_{\mathrm{Bl}_{k+1}(y')}\mathrm{d}x\,\varphi^{+}_B(y',c)\cdot f(x),
\end{align*}
where the last equality follows from 0.2\,(c).
Written out, the right side is $L^{-3}\cdot\eta_k^3\sum_{c\in\mathrm{Bl}_{k+1}(y')}\varphi^{+}_B(y',c)\cdot\frac{1}{48}\sum_{g\in O_{\mathrm{cub}}}f(c,g)$.
By (21.16) and (21.17), for $c\in\Delta_y$ and $y\in B(y')$, that is (by 0.2\,(c)) for $c\in\mathrm{Bl}_{k+1}(y')$,
\begin{align*}
\left(J\cdot Q_k^{T}(-B)\cdot Q^{T}(-B)\,v'\right)(c,g)
&=\varphi_B(y,c)^{-1}\cdot\left(Q^{T}(-B)\,v'\right)(y)=\varphi_B(y,c)^{-1}\cdot U_B(y',y)^{-1}\cdot v'(y')\\
&=\varphi^{+}_B(y',c)^{-1}\cdot v'(y').
\end{align*}
Since $Q_{k+1}(B)=Q(B)\cdot Q_k(B)$ and $Q_{k+1}^{T}(-B)=Q_k^{T}(-B)\cdot Q^{T}(-B)$ by (S21), these are the first two formulas of (b).
The definitions $P_{k+1}(B)=Q_k^{T}(-B)\cdot P(B)\cdot Q_k(B)$ and $P(B)=Q^{T}(-B)\cdot Q(B)$ of Section~I.3.3, together with the same two relations of (S21), give
\[
J\cdot P_{k+1}(B)\cdot J^*=J\cdot Q_k^{T}(-B)\cdot Q^{T}(-B)\cdot Q(B)\cdot Q_k(B)\cdot J^*=\left(J\cdot Q_{k+1}^{T}(-B)\right)\cdot\left(Q_{k+1}(B)\cdot J^*\right).
\]

\emph{The case $k=0$.}
Here $\eta_0=1$ and $\Delta_y=\{y\}$;
the sum that defines $(\tau_0B)(y,y)$ is empty, so that $\varphi_B(y,y)=1$ and $\varphi^{+}_B(y',c)=U_B(y',c)$;
moreover $\mathrm{Bl}_1(y')=B(y')$, $Q_1(B)=Q(B)$ and $Q_1^{T}(-B)=Q^{T}(-B)$.
Hence the first two formulas of (b) are (21.18) and (21.19),
and $J\cdot P_1(B)\cdot J^*=J\cdot P(B)\cdot J^*=\left(J\cdot Q^{T}(-B)\right)\cdot\left(Q(B)\cdot J^*\right)$, because $P(B)=Q^{T}(-B)\cdot Q(B)$.
\end{proof}

\begin{proof}[Proof of Proposition~21.1\,(a)]
\emph{Products of an average with its transpose at $-B$.}
Let $k\ge1$.
By (21.16), (21.15) and 0.2\,(b),
\begin{align*}
\left(\left(Q_k(B)\cdot J^*\right)\cdot\left(J\cdot Q_k^{T}(-B)\right)v\right)(y)
&=\eta_k^3\sum_{c\in\Delta_y}\varphi_B(y,c)\cdot\frac{1}{48}\sum_{g\in O_{\mathrm{cub}}}\varphi_B(y,c)^{-1}\cdot v(y)\\
&=\eta_k^3\cdot L^{3k}\cdot v(y)=v(y);
\end{align*}
here the sum over the cluster equals $48\cdot\frac{1}{48}=1$, which is the relation $J^*\cdot J=1$.
By (21.17),
\[
\left(Q(B)\cdot Q^{T}(-B)\,v'\right)(y')=L^{-3}\sum_{y\in B(y')}U_B(y',y)\cdot U_B(y',y)^{-1}\cdot v'(y')=L^{-3}\cdot L^3\cdot v'(y')=v'(y').
\]
It follows that
\[
P(B)\cdot Q^{T}(-B)=Q^{T}(-B)\cdot\left[Q(B)\cdot Q^{T}(-B)\right]=Q^{T}(-B),\qquad Q(B)\cdot P(B)=\left[Q(B)\cdot Q^{T}(-B)\right]\cdot Q(B)=Q(B).
\]
Let $k=0$.
By (21.19) and (21.18),
\begin{align*}
\left(\left(Q(B)\cdot J^*\right)\cdot\left(J\cdot Q^{T}(-B)\right)v'\right)(y')
&=L^{-3}\sum_{c\in B(y')}U_B(y',c)\cdot\frac{1}{48}\sum_{g\in O_{\mathrm{cub}}}U_B(y',c)^{-1}\cdot v'(y')\\
&=L^{-3}\cdot L^3\cdot v'(y')=v'(y').
\end{align*}
Together with $J\cdot P(B)\cdot J^*=\left(J\cdot Q^{T}(-B)\right)\cdot\left(Q(B)\cdot J^*\right)$ this gives
\begin{align*}
\left(J\cdot P(B)\cdot J^*\right)\cdot\left(J\cdot Q^{T}(-B)\right)&=J\cdot Q^{T}(-B)\cdot\left[\left(Q(B)\cdot J^*\right)\cdot\left(J\cdot Q^{T}(-B)\right)\right]=J\cdot Q^{T}(-B),\\
\left(Q(B)\cdot J^*\right)\cdot\left(J\cdot P(B)\cdot J^*\right)&=\left[\left(Q(B)\cdot J^*\right)\cdot\left(J\cdot Q^{T}(-B)\right)\right]\cdot Q(B)\cdot J^*=Q(B)\cdot J^*.
\end{align*}
For every $k\ge0$ one concludes that
\begin{equation}
\left(Q_{k+1}(B)\cdot J^*\right)\cdot\left(J\cdot Q_{k+1}^{T}(-B)\right)=1\qquad\text{on }\mathfrak{S}(L,\ell).
\tag{21.20}
\end{equation}
Indeed, for $k\ge1$ the product equals $Q(B)\cdot\left[\left(Q_k(B)\cdot J^*\right)\cdot\left(J\cdot Q_k^{T}(-B)\right)\right]\cdot Q^{T}(-B)=Q(B)\cdot Q^{T}(-B)=1$,
and for $k=0$ it is the relation just proved.

\emph{Idempotents.}
If $\mathrm{Xm}\colon V\to U$ and $\mathrm{Ym}\colon U\to V$ are linear maps with $\mathrm{Xm}\cdot\mathrm{Ym}=1_U$, then
\[
\left(\mathrm{Ym}\cdot\mathrm{Xm}\right)^2=\mathrm{Ym}\cdot\left(\mathrm{Xm}\cdot\mathrm{Ym}\right)\cdot\mathrm{Xm}=\mathrm{Ym}\cdot\mathrm{Xm}.
\]
This applies, for $k\ge1$, to the pairs
\[
(\mathrm{Xm},\mathrm{Ym})=\left(Q_k(B)\cdot J^*,\ J\cdot Q_k^{T}(-B)\right),\qquad \left(Q(B),\ Q^{T}(-B)\right),\qquad \left(Q_{k+1}(B)\cdot J^*,\ J\cdot Q_{k+1}^{T}(-B)\right),
\]
for the last of which $\mathrm{Xm}\cdot\mathrm{Ym}=1$ is (21.20), and, for $k=0$, to the pair $\left(Q(B)\cdot J^*,\ J\cdot Q^{T}(-B)\right)$.
The corresponding products $\mathrm{Ym}\cdot\mathrm{Xm}$ are $J\cdot P_k(B)\cdot J^*$, $P(B)$ and $J\cdot P_{k+1}(B)\cdot J^*$ (the last one by (b)) for $k\ge1$, and $J\cdot P(B)\cdot J^*$ for $k=0$;
so these operators are idempotent.

\emph{Kernels; constant matrices.}
The four operators are products of the scalar operators $J$, $J^*$, $Q_k(B)$, $Q_k^{T}(-B)$, $Q(B)$ and $Q^{T}(-B)$ (see 0.7).
Hence they are scalar, and by 0.5\,(e) their kernels are multiples of $1_4$ and they commute with every constant $4\times4$ matrix.
Explicitly, with $x=(c,g)$, $x'=(c',g')$ and the measure (I.1.3), the formulas (21.15)--(21.19) and (b) give
\begin{equation}
\begin{aligned}
&\left(J\cdot P_k(B)\cdot J^*f\right)(x)=\int_{\Delta_y}\mathrm{d}x'\,\varphi_B(y,c)^{-1}\cdot\varphi_B(y,c')\cdot f(x')\qquad(c\in\Delta_y),\\
&\left(J\cdot P_{k+1}(B)\cdot J^*f\right)(x)=\int_{\mathrm{Bl}_{k+1}(y')}\mathrm{d}x'\,L^{-3}\cdot\varphi^{+}_B(y',c)^{-1}\cdot\varphi^{+}_B(y',c')\cdot f(x')\qquad(c\in\mathrm{Bl}_{k+1}(y')),\\
&\left(P(B)\,v\right)(y)=\sum_{\tilde y\in B(y')}L^{-3}\cdot U_B(y',y)^{-1}\cdot U_B(y',\tilde y)\cdot v(\tilde y)\qquad(y\in B(y'));
\end{aligned}
\tag{21.21}
\end{equation}
for $k=0$ the second of these formulas, with $\varphi^{+}_B=U_B$ and $\mathrm{Bl}_1(y')=B(y')$, describes $J\cdot P(B)\cdot J^*$.
In (I.1.4) the variable $x'$ ranges over a fundamental domain;
a site over $\Delta_y$ has the form $x'-\ell\cdot n$ with $x'$ in this domain and $n\in\mathbb{Z}^3$, and $f(x'-\ell\cdot n)=\mathrm{sgl}(n)\cdot f(x')$.
Consequently the kernel (I.1.4) of $J\cdot P_k(B)\cdot J^*$ at $(x,x')$, $c\in\Delta_y$, is
\[
\mathrm{sgl}(n)\cdot\varphi_B(y,c)^{-1}\cdot\varphi_B(y,c'-\ell\cdot n)\cdot1_4
\]
if $c'-\ell\cdot n\in\Delta_y$ for some $n\in\mathbb{Z}^3$, and $0$ otherwise;
there is at most one such $n$, since the sup-diameter of $\Delta_y$ is $<1\le\ell$.
This kernel does not depend on the labels $g$, $g'$ (see (I.J5)).
By the second formula of (21.21), the kernel of $J\cdot P_{k+1}(B)\cdot J^*$ (for $k=0$, of $J\cdot P(B)\cdot J^*$) at $(x,x')$, $c\in\mathrm{Bl}_{k+1}(y')$, is
\[
\mathrm{sgl}(n)\cdot L^{-3}\cdot\varphi^{+}_B(y',c)^{-1}\cdot\varphi^{+}_B(y',c'-\ell\cdot n)\cdot1_4
\]
if $c'-\ell\cdot n\in\mathrm{Bl}_{k+1}(y')$ for some $n\in\mathbb{Z}^3$, and $0$ otherwise;
again there is at most one such $n$, since the sup-diameter of $\mathrm{Bl}_{k+1}(y')$ is $<L\le\ell$.
By the third formula, the kernel of $P(B)$ at $(y,\tilde y)$, $y\in B(y')$, is
\[
\mathrm{sgl}(n)\cdot L^{-3}\cdot U_B(y',y)^{-1}\cdot U_B(y',\tilde y-\ell\cdot n)\cdot1_4
\]
if $\tilde y-\ell\cdot n\in B(y')$ for some $n\in\mathbb{Z}^3$, and $0$ otherwise;
there is at most one such $n$, since the sup-diameter of $B(y')$ is $L-1<\ell$.
Here $P(B)$ acts on $\mathfrak{S}(1,\ell)$, and its kernel refers to the sum over the unit sites of a fundamental domain, without weight (Section~I.1.6).

\emph{Products of the two projections, $k\ge1$.}
Since $\left(Q_k(B)\cdot J^*\right)\cdot\left(J\cdot Q_k^{T}(-B)\right)=1$,
\begin{align*}
\left(J\cdot P_k(B)\cdot J^*\right)\cdot\left(J\cdot P_{k+1}(B)\cdot J^*\right)
&=J\cdot Q_k^{T}(-B)\cdot\left[Q_k(B)\cdot J^*\cdot J\cdot Q_k^{T}(-B)\right]\cdot P(B)\cdot Q_k(B)\cdot J^*\\
&=J\cdot Q_k^{T}(-B)\cdot P(B)\cdot Q_k(B)\cdot J^*=J\cdot P_{k+1}(B)\cdot J^*,\\
\left(J\cdot P_{k+1}(B)\cdot J^*\right)\cdot\left(J\cdot P_k(B)\cdot J^*\right)
&=J\cdot Q_k^{T}(-B)\cdot P(B)\cdot\left[Q_k(B)\cdot J^*\cdot J\cdot Q_k^{T}(-B)\right]\cdot Q_k(B)\cdot J^*\\
&=J\cdot P_{k+1}(B)\cdot J^*.
\end{align*}

\emph{Traces.}
Restriction to a set of representatives is an isomorphism of $\mathfrak{S}(1,\ell)$ onto $\left(\mathbb{C}^4\right)^{\mathrm{Fdm}^1}$, of $\mathfrak{S}(L,\ell)$ onto $\left(\mathbb{C}^4\right)^{\mathrm{Fdm}^L}$,
and, for $k=0$, of $\mathfrak{S}^{\mathrm{ch}}(1,\ell)$ onto $\left(\mathbb{C}^4\right)^{\mathrm{Fdm}^1\times O_{\mathrm{cub}}}$;
it transforms the action of $\mathrm{Mx}$ into the direct sum of $n_k$, $n_{k+1}$, resp. $48\cdot n_0$ copies of the matrix $\mathrm{Mx}$.
Hence $\operatorname{Tr}\mathrm{Mx}$ is equal to $n_k\cdot\operatorname{tr}\mathrm{Mx}$ on $\mathfrak{S}(1,\ell)$, to $n_{k+1}\cdot\operatorname{tr}\mathrm{Mx}$ on $\mathfrak{S}(L,\ell)$, and to $48\cdot n_0\cdot\operatorname{tr}\mathrm{Mx}$ on $\mathfrak{S}^{\mathrm{ch}}(1,\ell)$.
Moreover $\operatorname{Tr}\left(\mathrm{Ym}\cdot\mathrm{Xm}\right)=\operatorname{Tr}\left(\mathrm{Xm}\cdot\mathrm{Ym}\right)$ for linear maps $\mathrm{Xm}\colon V\to U$ and $\mathrm{Ym}\colon U\to V$.
For $k\ge1$, since $Q(B)$ commutes with $\mathrm{Mx}$,
\begin{align*}
\operatorname{Tr}\left[\mathrm{Mx}\cdot P(B)\right]&=\operatorname{Tr}\left[\left(\mathrm{Mx}\cdot Q^{T}(-B)\right)\cdot Q(B)\right]=\operatorname{Tr}\left[Q(B)\cdot\mathrm{Mx}\cdot Q^{T}(-B)\right]\\
&=\operatorname{Tr}\left[\mathrm{Mx}\cdot Q(B)\cdot Q^{T}(-B)\right]=\operatorname{Tr}_{\mathfrak{S}(L,\ell)}\mathrm{Mx}=n_{k+1}\cdot\operatorname{tr}\mathrm{Mx},
\end{align*}
and $\operatorname{Tr}\left[\mathrm{Mx}\cdot\left(1-P(B)\right)\right]=n_k\cdot\operatorname{tr}\mathrm{Mx}-n_{k+1}\cdot\operatorname{tr}\mathrm{Mx}$.
For $k=0$, since $Q(B)\cdot J^*$ commutes with $\mathrm{Mx}$,
\begin{align*}
\operatorname{Tr}\left[\mathrm{Mx}\cdot J\cdot P(B)\cdot J^*\right]&=\operatorname{Tr}\left[\left(\mathrm{Mx}\cdot J\cdot Q^{T}(-B)\right)\cdot\left(Q(B)\cdot J^*\right)\right]=\operatorname{Tr}\left[\mathrm{Mx}\cdot\left(Q(B)\cdot J^*\right)\cdot\left(J\cdot Q^{T}(-B)\right)\right]\\
&=\operatorname{Tr}_{\mathfrak{S}(L,\ell)}\mathrm{Mx}=n_1\cdot\operatorname{tr}\mathrm{Mx},
\end{align*}
and $\operatorname{Tr}\left[\mathrm{Mx}\cdot\left(1-J\cdot P(B)\cdot J^*\right)\right]=48\cdot n_0\cdot\operatorname{tr}\mathrm{Mx}-n_1\cdot\operatorname{tr}\mathrm{Mx}$.
The same numbers are obtained from (I.ke2) and (21.21):
the kernels on the diagonal are $L^{-3}\cdot1_4$, and $\sum_{y\in\mathrm{Fdm}^1}L^{-3}\cdot\operatorname{tr}\mathrm{Mx}=n_k\cdot L^{-3}\cdot\operatorname{tr}\mathrm{Mx}$ for $k\ge1$,
while $\int\mathrm{d}x\,L^{-3}\cdot\operatorname{tr}\mathrm{Mx}=n_0\cdot L^{-3}\cdot\operatorname{tr}\mathrm{Mx}$ for $k=0$, where $\int\mathrm{d}x=\frac{1}{48}\cdot48\cdot n_0$.

\emph{The relation $T_k\cdot H_k=1$.}
Let $k\ge1$, and suppose that $S_k(B)=F_{\infty}(B)^{-1}$ exists.
By Section~I.3.5,
\[
T_k(B)\cdot H_k(B)=b_k^{-1}\cdot Q_k(B)\cdot J^*\cdot F_{\infty}(B)\cdot b_k\cdot S_k(B)\cdot J\cdot Q_k^{T}(-B)=\left(Q_k(B)\cdot J^*\right)\cdot\left(J\cdot Q_k^{T}(-B)\right)=1.
\]
For $k=0$ one has $T_0(B)=H_0(B)=1$ on $\mathfrak{S}^{\mathrm{ch}}(1,\ell)$ by Section~I.3.9.
\end{proof}

\begin{proof}[Proof of Proposition~21.1\,(c) for $D_B$, the averaging operators, the projections, $F_{\infty}(B)$ and $S_k(B)$]
\emph{The hopping matrices.}
By (I.c3), $\gamma_\mu^T=-\mathsf{C}^{-1}\cdot\gamma_\mu\cdot\mathsf{C}$.
Write $\gamma^T\cdot v:=\sum_\mu v_\mu\cdot\gamma_\mu^T$ for $v\in\mathbb{R}^3$.
By (I.2.1) and (I.s2),
\[
M_i\left(x^{(i)}\right)^T=\frac{w_i}{2}\left(\rho-\gamma^T\cdot\hat{d}_i\left(x^{(i)}\right)\right)=\frac{w_i}{2}\left(\rho+\gamma^T\cdot\hat{d}_i(x)\right)=\frac{w_i}{2}\left(\rho-\mathsf{C}^{-1}\cdot\gamma\cdot\hat{d}_i(x)\cdot\mathsf{C}\right),
\]
that is,
\begin{equation}
M_i\left(x^{(i)}\right)^T=\mathsf{C}^{-1}\cdot M_i(x)\cdot\mathsf{C}\qquad(i=0,\dots,3).
\tag{21.22}
\end{equation}
(This is (I.m2) at the site $x^{(i)}$: since $\mathsf{C}^{-1}=-\mathsf{C}$, one has $\mathsf{C}\cdot\mathrm{Mx}\cdot\mathsf{C}^{-1}=\mathsf{C}^{-1}\cdot\mathrm{Mx}\cdot\mathsf{C}$ for every matrix $\mathrm{Mx}$.)
By (I.m1), which rests on the closure relation (I.1.1), $\sum_iM_i(x)=\frac{\rho}{2}\cdot\mathrm{sw}$ is a multiple of $1_4$; hence
\begin{equation}
\sum_iM_i(x)^T=\frac{\rho}{2}\cdot\mathrm{sw}=\mathsf{C}^{-1}\cdot\left(\sum_iM_i(x)\right)\cdot\mathsf{C}.
\tag{21.23}
\end{equation}

\emph{The operator $D_B$.}
Let $f'',f\in\mathfrak{S}^{\mathrm{ch}}(\eta_k,\ell)$.
By (I.2.3),
\begin{align*}
\left\langle f'',D_Bf\right\rangle^{\mathrm{ch}}
&=-\eta_k^{-1}\sum_{i=0}^{3}\int\mathrm{d}x\,f''(x)\cdot M_i(x)\cdot U_B\left(x,x^{(i)}\right)\cdot f\left(x^{(i)}\right)\\
&\quad+\eta_k^{-1}\int\mathrm{d}x\,f''(x)\cdot\left(\sum_iM_i(x)\right)f(x).
\end{align*}
Each integrand is periodic (see 0.7).
In the $i$-th term of the first sum, replace $x$ by $x^{(i)}$ (see 0.4\,(c));
since $\left(x^{(i)}\right)^{(i)}=x$, this term becomes
\[
\int\mathrm{d}x\,f''\left(x^{(i)}\right)\cdot M_i\left(x^{(i)}\right)\cdot U_B\left(x^{(i)},x\right)\cdot f(x)=\int\mathrm{d}x\,\left[U_{-B}\left(x,x^{(i)}\right)\cdot M_i\left(x^{(i)}\right)^Tf''\left(x^{(i)}\right)\right]\cdot f(x),
\]
by (21.12) and the rule $f''\cdot\left(\mathrm{Mx}\,f\right)=\left(\mathrm{Mx}^Tf''\right)\cdot f$.
Therefore, by (21.22) and (21.23),
\begin{align*}
\left(D_B^Tf''\right)(x)
&=-\eta_k^{-1}\sum_iU_{-B}\left(x,x^{(i)}\right)\cdot M_i\left(x^{(i)}\right)^Tf''\left(x^{(i)}\right)+\eta_k^{-1}\left(\sum_iM_i(x)^T\right)f''(x)\\
&=-\eta_k^{-1}\cdot\mathsf{C}^{-1}\sum_iM_i(x)\left[U_{-B}\left(x,x^{(i)}\right)\cdot\mathsf{C}f''\left(x^{(i)}\right)-\mathsf{C}f''(x)\right]\\
&=\left(\mathsf{C}^{-1}\cdot D_{-B}\cdot\mathsf{C}f''\right)(x).
\end{align*}
The sum over $i$ comprises the three hops inside the cluster ($i=0,1,2$, where $U_B\left(x,x^{(i)}\right)=U_B\left(x^{(i)},x\right)=U_{-B}\left(x,x^{(i)}\right)=1$ by (21.12)) and the hop of type $3$;
the computation treats all four hops alike.

\emph{The averaging operators.}
By 0.3\,(b), $\varphi_{-B}=\varphi_B^{-1}$ and $U_{-B}=U_B^{-1}$; hence, by Section~I.3.3,
\[
\left(Q_k^{T}(B)\,v\right)(c)=\varphi_B(y,c)\cdot v(y),\qquad \left(Q^{T}(B)\,v'\right)(y)=U_B(y',y)\cdot v'(y').
\]
For $k\ge1$, by (21.15) and (21.13),
\[
\left\langle v,Q_k(B)\cdot J^*f\right\rangle_1=\sum_{y\in\mathrm{Fdm}^1}v(y)\cdot\int_{\Delta_y}\mathrm{d}x\,\varphi_B(y,c)\cdot f(x)=\int\mathrm{d}x\,\left[\varphi_B([c],c)\cdot v([c])\right]\cdot f(x)=\left\langle J\cdot Q_k^{T}(B)\,v,f\right\rangle^{\mathrm{ch}}.
\]
For $k\ge0$, by (21.17) and (21.13),
\[
\left\langle v',Q(B)\,v\right\rangle_L=L^3\sum_{y'\in\mathrm{Fdm}^L}v'(y')\cdot L^{-3}\sum_{y\in B(y')}U_B(y',y)\cdot v(y)=\sum_{y\in\mathrm{Fdm}^1}\left[U_B(y',y)\cdot v'(y')\right]\cdot v(y)=\left\langle Q^{T}(B)\,v',v\right\rangle_1,
\]
where in the third expression $y'$ is the point of $\mathrm{Fdm}^L$ with $y\in B(y')$.
Hence $\left(Q_k(B)\cdot J^*\right)^T=J\cdot Q_k^{T}(B)$ for $k\ge1$, and $Q(B)^T=Q^{T}(B)$;
by 0.5\,(a), (d), $\left(Q(B)\cdot J^*\right)^T=\left(J^*\right)^T\cdot Q(B)^T=J\cdot Q^{T}(B)$.
(For $k=0$ the last relation reads
$\left\langle v',Q(B)\cdot J^*f\right\rangle_L=\sum_{y'}v'(y')\cdot\sum_{c\in B(y')}U_B(y',c)\cdot\frac{1}{48}\sum_gf(c,g)=\int\mathrm{d}x\,\left[U_B(y',c)\cdot v'(y')\right]\cdot f(x)$.)
Replacing $B$ by $-B$ and transposing (see 0.5\,(a)), one obtains
\begin{equation}
\begin{aligned}
&\left(J\cdot Q_k^{T}(-B)\right)^T=Q_k(-B)\cdot J^*,\qquad Q^{T}(-B)^T=Q(-B),\qquad \left(J\cdot Q^{T}(-B)\right)^T=Q(-B)\cdot J^*,\\
&\left(Q_{k+1}(B)\cdot J^*\right)^T=J\cdot Q_{k+1}^{T}(B),\qquad \left(J\cdot Q_{k+1}^{T}(-B)\right)^T=Q_{k+1}(-B)\cdot J^*;
\end{aligned}
\tag{21.24}
\end{equation}
for $k\ge1$ the last two relations follow from $\left(Q(B)\cdot Q_k(B)\cdot J^*\right)^T=J\cdot Q_k^{T}(B)\cdot Q^{T}(B)$,
and for $k=0$ they are the relations for $Q(\cdot)\cdot J^*$.

\emph{The projections.}
Since $P_k(-B)=Q_k^{T}(-(-B))\cdot Q_k(-B)$,
\[
\left(J\cdot P_k(B)\cdot J^*\right)^T=\left(Q_k(B)\cdot J^*\right)^T\cdot\left(J\cdot Q_k^{T}(-B)\right)^T=J\cdot Q_k^{T}(B)\cdot Q_k(-B)\cdot J^*=J\cdot P_k(-B)\cdot J^*.
\]
By (b) and (21.24),
\begin{gather*}
\left(J\cdot P_{k+1}(B)\cdot J^*\right)^T=\left(Q_{k+1}(B)\cdot J^*\right)^T\cdot\left(J\cdot Q_{k+1}^{T}(-B)\right)^T=J\cdot Q_{k+1}^{T}(B)\cdot Q_{k+1}(-B)\cdot J^*=J\cdot P_{k+1}(-B)\cdot J^*,\\
P(B)^T=Q^{T}(B)\cdot Q(-B)=P(-B),\qquad \left(J\cdot P(B)\cdot J^*\right)^T=J\cdot Q^{T}(B)\cdot Q(-B)\cdot J^*=J\cdot P(-B)\cdot J^*.
\end{gather*}

\emph{The operators $F_{\infty}$ and $S_k$, $k\ge1$.}
The quantities $\mu$ and $b_k$ are numbers, and $J\cdot P_k(-B)\cdot J^*$ commutes with $\mathsf{C}$ by clause (a). Hence
\[
F_{\infty}(B)^T=D_B^T+\mu+b_k\cdot\left(J\cdot P_k(B)\cdot J^*\right)^T=\mathsf{C}^{-1}\cdot D_{-B}\cdot\mathsf{C}+\mu+b_k\cdot J\cdot P_k(-B)\cdot J^*=\mathsf{C}^{-1}\cdot F_{\infty}(-B)\cdot\mathsf{C}.
\]
By 0.5\,(b), $F_{\infty}(B)$ is invertible if and only if $F_{\infty}(B)^T$ is invertible, that is, if and only if $F_{\infty}(-B)$ is invertible; in this case
\[
S_k(B)^T=\left(F_{\infty}(B)^T\right)^{-1}=\mathsf{C}^{-1}\cdot S_k(-B)\cdot\mathsf{C}.
\]
\end{proof}

\emph{The operators $D_k$, $M_k$, $\dot{M}_k$ ($k\ge1$; $S_k(B)$ exists, hence so does $S_k(-B)$).}
By (21.24), and because the scalar operators $Q_k(-B)J^{*}$ and $JQ_k^{T}(B)$ commute with $\mathsf{C}$ and with $\mathsf{C}^{-1}$,
\begin{align*}
D_k(B)^{T}
&=b_{k}-b_{k}^{2}\cdot\left(JQ_k^{T}(-B)\right)^{T}\cdot S_k(B)^{T}\cdot\left(Q_k(B)J^{*}\right)^{T}\\
&=b_{k}-b_{k}^{2}\cdot Q_k(-B)J^{*}\cdot\mathsf{C}^{-1}\cdot S_k(-B)\cdot\mathsf{C}\cdot JQ_k^{T}(B)\\
&=\mathsf{C}^{-1}\cdot\left[b_{k}-b_{k}^{2}\cdot Q_k(-B)J^{*}\cdot S_k(-B)\cdot JQ_k^{T}\left(-(-B)\right)\right]\cdot\mathsf{C}
=\mathsf{C}^{-1}\cdot D_k(-B)\cdot\mathsf{C}.
\end{align*}
Since $M_k=b_{k}-D_k$, it follows that $M_k(B)^{T}=b_{k}-\mathsf{C}^{-1}\cdot D_k(-B)\cdot\mathsf{C}=\mathsf{C}^{-1}\cdot M_k(-B)\cdot\mathsf{C}$.
Moreover $\left(S_k(B)^{2}\right)^{T}=\left(S_k(B)^{T}\right)^{2}=\mathsf{C}^{-1}\cdot S_k(-B)^{2}\cdot\mathsf{C}$, and therefore
\begin{align*}
\dot{M}_k(B)^{T}
&=b_{k}^{2}\cdot\left(JQ_k^{T}(-B)\right)^{T}\cdot\left(S_k(B)^{2}\right)^{T}\cdot\left(Q_k(B)J^{*}\right)^{T}\\
&=b_{k}^{2}\cdot Q_k(-B)J^{*}\cdot\mathsf{C}^{-1}\cdot S_k(-B)^{2}\cdot\mathsf{C}\cdot JQ_k^{T}(B)
=\mathsf{C}^{-1}\cdot\dot{M}_k(-B)\cdot\mathsf{C}.
\end{align*}

\emph{The operator $\Gamma_k$, and the equivalence in (c) that concerns $D_k(\pm B)+b_{L}P(\pm B)$ ($k\ge1$; $S_k(B)$ exists, hence so does $S_k(-B)$).}
By the relation $P(B)^{T}=P(-B)$ obtained above, and since $P(-B)$ commutes with $\mathsf{C}$ by clause~(a),
\[
\left(D_k(B)+b_{L}P(B)\right)^{T}=\mathsf{C}^{-1}\cdot D_k(-B)\cdot\mathsf{C}+b_{L}P(-B)=\mathsf{C}^{-1}\cdot\left(D_k(-B)+b_{L}P(-B)\right)\cdot\mathsf{C}.
\]
Consequently $D_k(B)+b_{L}P(B)$ is invertible if and only if $D_k(-B)+b_{L}P(-B)$ is (Section~0.5\,(b)), and in this case $\Gamma_k(B)^{T}=\mathsf{C}^{-1}\cdot\Gamma_k(-B)\cdot\mathsf{C}$.

\emph{The case $k=0$.}
One has
\begin{align*}
D_0(B)^{T}&=\left(D_B+\mu\right)^{T}=\mathsf{C}^{-1}\cdot\left(D_{-B}+\mu\right)\cdot\mathsf{C}=\mathsf{C}^{-1}\cdot D_0(-B)\cdot\mathsf{C},\\
F_0(B)^{T}&=D_B^{T}+\mu+b_{L}\left(JP(B)J^{*}\right)^{T}=\mathsf{C}^{-1}\cdot\left(D_{-B}+\mu+b_{L}\,JP(-B)J^{*}\right)\cdot\mathsf{C}=\mathsf{C}^{-1}\cdot F_0(-B)\cdot\mathsf{C}.
\end{align*}
Hence $F_0(B)$ is invertible if and only if $F_0(-B)$ is, and then $\Gamma_0(B)^{T}=\mathsf{C}^{-1}\cdot\Gamma_0(-B)\cdot\mathsf{C}$.

\emph{The barred operators ($k\ge1$; $S_k(B)$ exists).}
One computes
\begin{align*}
\mathfrak{H}_{k}(B)^{T}
&=b_{k}\cdot S_k(B)^{T}\cdot\left(Q_k(B)J^{*}\right)^{T}
=b_{k}\cdot\mathsf{C}^{-1}\cdot S_k(-B)\cdot\mathsf{C}\cdot JQ_k^{T}(B)\\
&=\mathsf{C}^{-1}\cdot\left[b_{k}\cdot S_k(-B)\cdot JQ_k^{T}\left(-(-B)\right)\right]\cdot\mathsf{C}
=\mathsf{C}^{-1}\cdot H_{k}(-B)\cdot\mathsf{C}=\bar{H}_{k}(B),\\
\bar{T}_k(B)
&=\mathsf{C}^{-1}\cdot T_k(-B)\cdot\mathsf{C}
=\mathsf{C}^{-1}\cdot b_{k}^{-1}\cdot Q_k(-B)J^{*}\cdot F_{\infty}(-B)\cdot\mathsf{C}\\
&=b_{k}^{-1}\cdot Q_k(-B)J^{*}\cdot\mathsf{C}^{-1}\cdot F_{\infty}(-B)\cdot\mathsf{C}
=b_{k}^{-1}\cdot Q_k(-B)J^{*}\cdot F_{\infty}(B)^{T},\\
\bar{T}_k(B)\cdot\bar{H}_{k}(B)
&=\mathsf{C}^{-1}\cdot T_k(-B)\cdot H_{k}(-B)\cdot\mathsf{C}=1\qquad\text{(clause~(a) at $-B$)},\\
\mathfrak{H}_{k}(B)\cdot H_{k}(B)
&=b_{k}\cdot Q_k(B)J^{*}\cdot S_k(B)\cdot b_{k}\cdot S_k(B)\cdot JQ_k^{T}(-B)=\dot{M}_k(B).
\end{align*}
Let $\Psi$, $\bar{\Psi}$ be fields on $T^0_k$.
Section~0.6\,(b), applied to $O=\mathfrak{H}_{k}(B)\colon\mathfrak{S}^{\mathrm{ch}}(\eta_k,\ell)\to\mathfrak{S}(1,\ell)$, gives
\[
\left\langle\bar{H}_{k}(B)\bar{\Psi},H_{k}(B)\Psi\right\rangle^{\mathrm{ch}}
=\left\langle\mathfrak{H}_{k}(B)^{T}\bar{\Psi},H_{k}(B)\Psi\right\rangle^{\mathrm{ch}}
=\left\langle\bar{\Psi},\mathfrak{H}_{k}(B)\cdot H_{k}(B)\Psi\right\rangle_1
=\left\langle\bar{\Psi},\dot{M}_k(B)\Psi\right\rangle_1.
\]
At $k=0$ one has $\bar{H}_0=\bar{T}_0=\dot{M}_0=1$ on $\mathfrak{S}^{\mathrm{ch}}(1,\ell)$ by Section~I.3.9.

\paragraph*{(d0) (phases).}
\emph{Inverses and conjugates.}
By Section~0.3\,(b), $e_k\cdot\eta_k\cdot\left(\tau_k(-B)\right)(y,c)=-e_k\cdot\eta_k\cdot(\tau_kB)(y,c)$, so that $\varphi_{-B}=\varphi_{B}^{-1}$;
likewise $e_k\cdot\eta_k\cdot\left(\tau(-B)\right)(c,c')=-e_k\cdot\eta_k\cdot(\tau B)(c,c')$, so that $U_{-B}=U_B^{-1}$.
Since $e_k$ and $\eta_k$ are real,
\begin{align*}
\overline{\varphi_{\overline{B}}(y,c)}
&=\exp\left\{-\mathrm{i}\,e_k\cdot\eta_k\cdot\overline{(\tau_k\overline{B})(y,c)}\right\}
=\exp\left\{-\mathrm{i}\,e_k\cdot\eta_k\cdot(\tau_kB)(y,c)\right\}=\varphi_{-B}(y,c),\\
\overline{U_{\overline{B}}(c,c')}
&=\exp\left\{-\mathrm{i}\,e_k\cdot\eta_k\cdot\overline{(\tau\overline{B})(c,c')}\right\}
=\exp\left\{-\mathrm{i}\,e_k\cdot\eta_k\cdot(\tau B)(c,c')\right\}=U_{-B}(c,c').
\end{align*}

\emph{Gauge transformations.}
Let $\boldsymbol{\lambda}$ be a complex function on $\mathbb{L}_k$ such that $d\boldsymbol{\lambda}$ is $\ell\mathbb{Z}^3$-periodic; then $B-d\boldsymbol{\lambda}$ is a bond function on $T_k$.
By Section~I.1.1, $(d\boldsymbol{\lambda})(c,c+\eta_ke_{\mu})=\eta_k^{-1}\left(\boldsymbol{\lambda}(c+\eta_ke_{\mu})-\boldsymbol{\lambda}(c)\right)$, and $d\boldsymbol{\lambda}$ is odd under reversal of the bond.
Hence
\[
\eta_k\cdot(d\boldsymbol{\lambda})(c_{(1)},c_{(2)})=\boldsymbol{\lambda}(c_{(2)})-\boldsymbol{\lambda}(c_{(1)})
\]
for every bond $(c_{(1)},c_{(2)})$.
Along a path $\boldsymbol{\Gamma}=\left(c_{(0)},\dots,c_{(\mathrm{len})}\right)$ the sum telescopes: $\eta_k\cdot(d\boldsymbol{\lambda})(\boldsymbol{\Gamma})=\boldsymbol{\lambda}(c_{(\mathrm{len})})-\boldsymbol{\lambda}(c_{(0)})$.
Taking the mean over the six paths from $c$ to $c'$ one obtains $\eta_k\cdot(\tau\,d\boldsymbol{\lambda})(c,c')=\boldsymbol{\lambda}(c')-\boldsymbol{\lambda}(c)$.
Therefore
\[
U_{B-d\boldsymbol{\lambda}}(c,c')
=\exp\left\{\mathrm{i}\,e_k\cdot\eta_k\cdot(\tau B)(c,c')-\mathrm{i}\,e_k\cdot\left(\boldsymbol{\lambda}(c')-\boldsymbol{\lambda}(c)\right)\right\}
=\mathrm{e}^{\mathrm{i}e_k\left(\boldsymbol{\lambda}(c)-\boldsymbol{\lambda}(c')\right)}\cdot U_B(c,c').
\]
With $y_j=y_j(c)$ one has
\[
\eta_k\cdot(\tau_k\,d\boldsymbol{\lambda})(y,c)
=\sum_{j<k}\eta_k\cdot(\tau\,d\boldsymbol{\lambda})(y_{j+1},y_j)
=\sum_{j<k}\left(\boldsymbol{\lambda}(y_j)-\boldsymbol{\lambda}(y_{j+1})\right)
=\boldsymbol{\lambda}(y_0)-\boldsymbol{\lambda}(y_k)
=\boldsymbol{\lambda}(c)-\boldsymbol{\lambda}(y),
\]
and consequently
\[
\varphi_{B-d\boldsymbol{\lambda}}(y,c)
=\exp\left\{-\mathrm{i}\,e_k\cdot\left(\boldsymbol{\lambda}(c)-\boldsymbol{\lambda}(y)\right)\right\}\cdot\varphi_B(y,c)
=\mathrm{e}^{\mathrm{i}e_k\left(\boldsymbol{\lambda}(y)-\boldsymbol{\lambda}(c)\right)}\cdot\varphi_B(y,c).
\]

\emph{Lattice symmetries.}
Let $\tilde{g}=(S,a)\in\mathbb{G}_k$, let $r=r(S)=\mathrm{diag}(r_1,r_2,r_3)$ with $r_\mu=\pm1$, and let $a\in M\mathbb{Z}^3$ (Sections~I.4.1--I.4.3).
Thus $\tilde{g}c=rc+a$ and $(\tilde{g}B)(\tilde{g}c_{(1)},\tilde{g}c_{(2)})=B(c_{(1)},c_{(2)})$.
Since $(\tilde{g}c)_\mu=r_\mu c_\mu+a_\mu$, each coordinate is transformed on its own: $c_\mu$ goes to $c_\mu+a_\mu$ or to $-c_\mu+a_\mu$.
Both maps are strictly monotone, preserve distances, and map $\eta_k\mathbb{Z}$ onto itself, because $a_\mu$ is an integer.
Hence $\tilde{g}$ maps bonds of $\mathbb{L}_k$ to bonds.
Now let $c_{(0)},\dots,c_{(\mathrm{len})}$ be the path $\boldsymbol{\Gamma}^{\pi}(c,c')$ of Section~0.3\,(a).
In the sequence $\tilde{g}c_{(0)},\dots,\tilde{g}c_{(\mathrm{len})}$ the coordinate $\pi_1$ first moves from $(\tilde{g}c)_{\pi_1}$ to $(\tilde{g}c')_{\pi_1}$ in steps of length $\eta_k$, then the coordinate $\pi_2$ does so, then $\pi_3$.
So this sequence is $\boldsymbol{\Gamma}^{\pi}(\tilde{g}c,\tilde{g}c')$, with the same order $\pi$; this is the assertion on paths in Proposition~16.1\,(a1).
It follows that
\[
(\tilde{g}B)\left(\boldsymbol{\Gamma}^{\pi}(\tilde{g}c,\tilde{g}c')\right)
=\sum_j(\tilde{g}B)(\tilde{g}c_{(j)},\tilde{g}c_{(j+1)})
=\sum_jB(c_{(j)},c_{(j+1)})
=B\left(\boldsymbol{\Gamma}^{\pi}(c,c')\right),
\]
and the mean over $\pi$ gives $(\tau\,\tilde{g}B)(\tilde{g}c,\tilde{g}c')=(\tau B)(c,c')$, the identity of Proposition~16.1\,(b).
Thus $U_{\tilde{g}B}(\tilde{g}c,\tilde{g}c')=U_B(c,c')$; here $|\tilde{g}c-\tilde{g}c'|=|c-c'|\le M$.

The map $\tilde{g}$ is an isometry for the sup norm, and for $j\le k$ it maps $\mathrm{bs}_j\mathbb{Z}^3=L^{j-k}\mathbb{Z}^3$ onto itself, since $a$ lies in $\mathbb{Z}^3$, which is contained in $\mathrm{bs}_j\mathbb{Z}^3$.
It therefore maps $\mathrm{Bl}_{j}(z)$ onto $\mathrm{Bl}_{j}(\tilde{g}z)$, and $y_j(\tilde{g}c)=\tilde{g}\,y_j(c)$ by Section~0.2\,(a); in particular $\tilde{g}c$ lies in $\Delta_{\tilde{g}y}$.
Hence
\[
(\tau_k\,\tilde{g}B)(\tilde{g}y,\tilde{g}c)
=\sum_{j<k}(\tau\,\tilde{g}B)(\tilde{g}y_{j+1},\tilde{g}y_j)
=\sum_{j<k}(\tau B)(y_{j+1},y_j)
=(\tau_kB)(y,c),
\]
and $\varphi_{\tilde{g}B}(\tilde{g}y,\tilde{g}c)=\varphi_B(y,c)$.
(If $a$ is replaced by another representative $a+\ell n$ of its class in $M\mathbb{Z}^3/\ell\mathbb{Z}^3$, then $\tilde{g}B$ does not change, while $\tilde{g}y$ and $\tilde{g}c$ are translated by $\ell n$; see Section~0.3\,(c).)

\paragraph*{(d) (gauge).}
Let $\boldsymbol{\lambda}$ be as in clause~(d), that is, a complex function on $\mathbb{L}_k$ such that $d\boldsymbol{\lambda}$ and $\mathrm{e}^{\mathrm{i}e_k\boldsymbol{\lambda}}$ are $\ell\mathbb{Z}^3$-periodic.
The function $\mathrm{e}^{\mathrm{i}e_k\boldsymbol{\lambda}}$ is periodic and has no zeros.
Hence the multiplication by it is a bijection of each of the spaces $\mathfrak{S}^{\mathrm{ch}}(\eta_k,\ell)$, $\mathfrak{S}(\eta_k,\ell)$, $\mathfrak{S}(1,\ell)$, $\mathfrak{S}(L,\ell)$, with inverse $\mathrm{e}^{-\mathrm{i}e_k\boldsymbol{\lambda}}$.
This multiplication commutes with $J$ and $J^{*}$, because it is a function of $c$ (all $48$ sites of a cluster have the same factor), and it commutes with constant matrices.
From (d0) one obtains
\begin{equation}
\begin{aligned}
&U_{B-d\boldsymbol{\lambda}}(c,c')\cdot\mathrm{e}^{\mathrm{i}e_k\boldsymbol{\lambda}(c')}=\mathrm{e}^{\mathrm{i}e_k\boldsymbol{\lambda}(c)}\cdot U_B(c,c'),\qquad
\varphi_{B-d\boldsymbol{\lambda}}(y,c)\cdot\mathrm{e}^{\mathrm{i}e_k\boldsymbol{\lambda}(c)}=\mathrm{e}^{\mathrm{i}e_k\boldsymbol{\lambda}(y)}\cdot\varphi_B(y,c),\\
&\varphi_{B-d\boldsymbol{\lambda}}(y,c)^{-1}\cdot\mathrm{e}^{\mathrm{i}e_k\boldsymbol{\lambda}(y)}=\mathrm{e}^{\mathrm{i}e_k\boldsymbol{\lambda}(c)}\cdot\varphi_B(y,c)^{-1}.
\end{aligned}
\tag{21.25}
\end{equation}

\emph{The operator $D_B$.}
Write $\pi(x^{(i)})$ for the photon site under $x^{(i)}$, so that $\pi(x^{(i)})=c$ for $i\le2$ and $\pi(x^{(3)})=c_{3}$.
By (I.2.3) and (I.1.6),
\begin{align*}
\left(D_{B-d\boldsymbol{\lambda}}\cdot\mathrm{e}^{\mathrm{i}e_k\boldsymbol{\lambda}}f\right)(x)
&=-\eta_k^{-1}\sum_{i}M_{i}(x)\left[U_{B-d\boldsymbol{\lambda}}\left(c,\pi(x^{(i)})\right)\cdot\mathrm{e}^{\mathrm{i}e_k\boldsymbol{\lambda}(\pi(x^{(i)}))}f(x^{(i)})-\mathrm{e}^{\mathrm{i}e_k\boldsymbol{\lambda}(c)}f(x)\right]\\
&=\mathrm{e}^{\mathrm{i}e_k\boldsymbol{\lambda}(c)}\left(D_Bf\right)(x).
\end{align*}
For $i\le2$ the first relation of (21.25) is used at $c'=c$, where $U_{B-d\boldsymbol{\lambda}}(c,c)=U_B(c,c)=1$ by (I.1.6); for $i=3$ it is the relation for the bond $(c,c_{3})$.

\emph{The averaging operators.}
By (21.15)--(21.17) and (21.25), with $x=(c,g)$, $c\in\Delta_y$, in the second identity and with $y\in B(y')$ in the fourth,
\begin{align*}
\left(Q_k(B-d\boldsymbol{\lambda})J^{*}\cdot\mathrm{e}^{\mathrm{i}e_k\boldsymbol{\lambda}}f\right)(y)
&=\int_{\Delta_y}\mathrm{d}x\,\varphi_{B-d\boldsymbol{\lambda}}(y,c)\cdot\mathrm{e}^{\mathrm{i}e_k\boldsymbol{\lambda}(c)}f(x)
=\mathrm{e}^{\mathrm{i}e_k\boldsymbol{\lambda}(y)}\left(Q_k(B)J^{*}f\right)(y),\\
\left(JQ_k^{T}\left(-(B-d\boldsymbol{\lambda})\right)\cdot\mathrm{e}^{\mathrm{i}e_k\boldsymbol{\lambda}}v\right)(x)
&=\varphi_{B-d\boldsymbol{\lambda}}(y,c)^{-1}\cdot\mathrm{e}^{\mathrm{i}e_k\boldsymbol{\lambda}(y)}v(y)
=\mathrm{e}^{\mathrm{i}e_k\boldsymbol{\lambda}(c)}\left(JQ_k^{T}(-B)v\right)(x),\\
\left(Q(B-d\boldsymbol{\lambda})\cdot\mathrm{e}^{\mathrm{i}e_k\boldsymbol{\lambda}}v\right)(y')
&=L^{-3}\sum_{y\in B(y')}U_{B-d\boldsymbol{\lambda}}(y',y)\cdot\mathrm{e}^{\mathrm{i}e_k\boldsymbol{\lambda}(y)}v(y)
=\mathrm{e}^{\mathrm{i}e_k\boldsymbol{\lambda}(y')}\left(Q(B)v\right)(y'),\\
\left(Q^{T}\left(-(B-d\boldsymbol{\lambda})\right)\cdot\mathrm{e}^{\mathrm{i}e_k\boldsymbol{\lambda}}v'\right)(y)
&=U_{B-d\boldsymbol{\lambda}}(y',y)^{-1}\cdot\mathrm{e}^{\mathrm{i}e_k\boldsymbol{\lambda}(y')}v'(y')
=\mathrm{e}^{\mathrm{i}e_k\boldsymbol{\lambda}(y)}\cdot U_B(y',y)^{-1}v'(y')\\
&=\mathrm{e}^{\mathrm{i}e_k\boldsymbol{\lambda}(y)}\left(Q^{T}(-B)v'\right)(y).
\end{align*}
At $k=0$, since $J$ and $J^{*}$ commute with $\mathrm{e}^{\mathrm{i}e_k\boldsymbol{\lambda}}$,
\begin{gather*}
Q(B-d\boldsymbol{\lambda})J^{*}\cdot\mathrm{e}^{\mathrm{i}e_k\boldsymbol{\lambda}}=Q(B-d\boldsymbol{\lambda})\cdot\mathrm{e}^{\mathrm{i}e_k\boldsymbol{\lambda}}\cdot J^{*}=\mathrm{e}^{\mathrm{i}e_k\boldsymbol{\lambda}}\cdot Q(B)J^{*},\\
JQ^{T}\left(-(B-d\boldsymbol{\lambda})\right)\cdot\mathrm{e}^{\mathrm{i}e_k\boldsymbol{\lambda}}=J\cdot\mathrm{e}^{\mathrm{i}e_k\boldsymbol{\lambda}}\cdot Q^{T}(-B)=\mathrm{e}^{\mathrm{i}e_k\boldsymbol{\lambda}}\cdot JQ^{T}(-B).
\end{gather*}

\emph{The H\"older quotient.}
Let $(x,x')\in\mathrm{Pzt}_k$, $x=(c,g)$, $x'=(c',g')$.
By (I.1.7) and (21.25),
\begin{align*}
\left(\delta_{\alpha,B-d\boldsymbol{\lambda}}\left(\mathrm{e}^{\mathrm{i}e_k\boldsymbol{\lambda}}f\right)\right)(x,x')
&=d_X(x,x')^{-\alpha}\left(U_{B-d\boldsymbol{\lambda}}(c,c')\cdot\mathrm{e}^{\mathrm{i}e_k\boldsymbol{\lambda}(c')}f(x')-\mathrm{e}^{\mathrm{i}e_k\boldsymbol{\lambda}(c)}f(x)\right)\\
&=\mathrm{e}^{\mathrm{i}e_k\boldsymbol{\lambda}(c)}\left(\delta_{\alpha,B}f\right)(x,x').
\end{align*}
For a pair inside a cluster ($c=c'$, $d_X=\eta_k$, $U_{B-d\boldsymbol{\lambda}}(c,c)=U_B(c,c)=1$) both sides are equal to $\eta_k^{-\alpha}\cdot\mathrm{e}^{\mathrm{i}e_k\boldsymbol{\lambda}(c)}\left(f(x')-f(x)\right)$.

\emph{Composed operators.}
Call a family $O(\cdot)$ of linear maps between the spaces named above \emph{intertwined} if $O(B-d\boldsymbol{\lambda})\cdot\mathrm{e}^{\mathrm{i}e_k\boldsymbol{\lambda}}=\mathrm{e}^{\mathrm{i}e_k\boldsymbol{\lambda}}\cdot O(B)$, that is, if $O(B-d\boldsymbol{\lambda})=\mathrm{e}^{\mathrm{i}e_k\boldsymbol{\lambda}}\cdot O(B)\cdot\mathrm{e}^{-\mathrm{i}e_k\boldsymbol{\lambda}}$.
Constant matrices are intertwined.
If $O_1$ and $O_2$ are intertwined, then so are $O_1+O_2$ and $O_1O_2$.
If $O$ is intertwined, then $O(B-d\boldsymbol{\lambda})$ is bijective if and only if $O(B)$ is, and $O(\cdot)^{-1}$ is intertwined.
The operators of Section~I.3.5 without bar are obtained by these operations from the families $B\mapsto D_B$, $Q_k(\cdot)J^{*}$, $JQ_k^{T}(-\,\cdot\,)$, $Q(\cdot)$, $Q^{T}(-\,\cdot\,)$, from numbers, and from the matrices $\boldsymbol{\alpha}_{y}$, $\boldsymbol{\beta}_{y}$, $\mathrm{i}\gamma_5y$, in the following order:
$JP_kJ^{*}$, $P$, $JP_{k+1}J^{*}$;
$\operatorname{Pen}_{y}$, $F_{y}$;
$S_{y}=F_{y}^{-1}$ (hence $S_k$, $S_{k,y}$, $S^{0}_{k+1}$);
$H_{k}$, $\mathfrak{H}_{k}$, $T_k$;
$D_k$, $M_k$, $\dot{M}_k$;
$D_k+b_{L}P+\mathrm{i}\gamma_5y$ and its inverse ($\Gamma_k$, $\Gamma_{k,y}$);
$b_{k}+b_{L}P+\mathrm{i}\gamma_5y$ and its inverse ($B_{k,y}$, $B_{k}$).
Those of Section~I.3.9 are obtained by the same operations:
$1$ ($=H_{0}$, $T_0$, $\dot{M}_0$, $Q_0$, $P_0$), $D_0$ (with $D_0(B)=D_B+\mu$), $JPJ^{*}$, $\operatorname{Pen}_{y}$, $F_{y}$, $S_{y}=F_{y}^{-1}$ (hence $\Gamma_0=S^{0}_{1}$ and $S_{0,y}$).
Consequently each of these operators exists at $B-d\boldsymbol{\lambda}$ if and only if it exists at $B$, and it is intertwined.
For $H_{k}^{\#}=\left(H_{k},\delta_{\alpha,\cdot}H_{k}\right)$ the second entry satisfies, by the relation for the H\"older quotient,
\[
\left(\delta_{\alpha,B-d\boldsymbol{\lambda}}H_{k}(B-d\boldsymbol{\lambda})\cdot\mathrm{e}^{\mathrm{i}e_k\boldsymbol{\lambda}}v\right)(x,x')
=\mathrm{e}^{\mathrm{i}e_k\boldsymbol{\lambda}(c)}\left(\delta_{\alpha,B}H_{k}(B)v\right)(x,x').
\]

\emph{Barred operators.}
The function $-\boldsymbol{\lambda}$ has the properties required of $\boldsymbol{\lambda}$: both $d(-\boldsymbol{\lambda})$ and $\mathrm{e}^{\mathrm{i}e_k(-\boldsymbol{\lambda})}=1/\mathrm{e}^{\mathrm{i}e_k\boldsymbol{\lambda}}$ are periodic.
Moreover $-(B-d\boldsymbol{\lambda})=(-B)-d(-\boldsymbol{\lambda})$.
The preceding paragraph, applied with $-B$ and $-\boldsymbol{\lambda}$ in place of $B$ and $\boldsymbol{\lambda}$, gives $H_{k}\left((-B)-d(-\boldsymbol{\lambda})\right)\cdot\mathrm{e}^{-\mathrm{i}e_k\boldsymbol{\lambda}}=\mathrm{e}^{-\mathrm{i}e_k\boldsymbol{\lambda}}\cdot H_{k}(-B)$.
Since $\mathsf{C}$ commutes with the multiplications, it follows that
\[
\bar{H}_{k}(B-d\boldsymbol{\lambda})\cdot\mathrm{e}^{-\mathrm{i}e_k\boldsymbol{\lambda}}
=\mathsf{C}^{-1}\cdot H_{k}\left((-B)-d(-\boldsymbol{\lambda})\right)\cdot\mathrm{e}^{-\mathrm{i}e_k\boldsymbol{\lambda}}\cdot\mathsf{C}
=\mathrm{e}^{-\mathrm{i}e_k\boldsymbol{\lambda}}\cdot\mathsf{C}^{-1}\cdot H_{k}(-B)\cdot\mathsf{C}
=\mathrm{e}^{-\mathrm{i}e_k\boldsymbol{\lambda}}\cdot\bar{H}_{k}(B).
\]
The same paragraph gives $T_k\left((-B)-d(-\boldsymbol{\lambda})\right)\cdot\mathrm{e}^{-\mathrm{i}e_k\boldsymbol{\lambda}}=\mathrm{e}^{-\mathrm{i}e_k\boldsymbol{\lambda}}\cdot T_k(-B)$, and this yields
\[
\bar{T}_k(B-d\boldsymbol{\lambda})\cdot\mathrm{e}^{-\mathrm{i}e_k\boldsymbol{\lambda}}
=\mathsf{C}^{-1}\cdot T_k\left((-B)-d(-\boldsymbol{\lambda})\right)\cdot\mathrm{e}^{-\mathrm{i}e_k\boldsymbol{\lambda}}\cdot\mathsf{C}
=\mathrm{e}^{-\mathrm{i}e_k\boldsymbol{\lambda}}\cdot\mathsf{C}^{-1}\cdot T_k(-B)\cdot\mathsf{C}
=\mathrm{e}^{-\mathrm{i}e_k\boldsymbol{\lambda}}\cdot\bar{T}_k(B).
\]
(For $k=0$: $\bar{H}_0=\bar{T}_0=1$.)

\emph{Determinants.}
By the paragraph on composed operators, $O(B-d\boldsymbol{\lambda})=\mathrm{e}^{\mathrm{i}e_k\boldsymbol{\lambda}}\cdot O(B)\cdot\mathrm{e}^{-\mathrm{i}e_k\boldsymbol{\lambda}}$ for $O=D_k$ and $O=D_k+b_{L}P$ when $k\ge1$, and for $O=D_0$ and $O=F_0$ when $k=0$.
These are similarities of endomorphisms of $\mathfrak{S}(1,\ell)$, respectively of $\mathfrak{S}^{\mathrm{ch}}(1,\ell)$.
Therefore $\det D_k(B-d\boldsymbol{\lambda})=\det D_k(B)$ and $\delta Z_k(B-d\boldsymbol{\lambda})=\delta Z_k(B)$.

\paragraph*{(e) (real $B$).}
If $B$ is real, the exponents in $\varphi_B$, $U_B$, $u_B$ are purely imaginary; hence $\overline{\varphi_B}=\varphi_B^{-1}$ and $\overline{U_B}=U_B^{-1}$.

\emph{The averaging operators.}
For $k\ge1$, by (21.15), (21.13) and (21.16),
\begin{align*}
\left(v,Q_k(B)J^{*}f\right)_1
&=\sum_{y\in\mathrm{Fdm}^1}\overline{v(y)}\cdot\int_{\Delta_y}\mathrm{d}x\,\varphi_B(y,c)f(x)
=\int\mathrm{d}x\,\overline{\varphi_B([c],c)^{-1}v([c])}\cdot f(x)\\
&=\left(JQ_k^{T}(-B)v,f\right)^{\mathrm{ch}}.
\end{align*}
For $k\ge0$, by (21.17) and (21.13),
\begin{align*}
\left(v',Q(B)v\right)_L
&=L^3\sum_{y'\in\mathrm{Fdm}^L}\overline{v'(y')}\cdot L^{-3}\sum_{y\in B(y')}U_B(y',y)v(y)
=\sum_{y\in\mathrm{Fdm}^1}\overline{U_B(y',y)^{-1}v'(y')}\cdot v(y)\\
&=\left(Q^{T}(-B)v',v\right)_1,
\end{align*}
where in the second sum $y'$ denotes the point of $L\mathbb{Z}^3$ with $y\in B(y')$.
Thus $\left(Q_k(B)J^{*}\right)^{*}=JQ_k^{T}(-B)$ for $k\ge1$ and $Q(B)^{*}=Q^{T}(-B)$; moreover $\left(Q(B)J^{*}\right)^{*}=(J^{*})^{*}\cdot Q(B)^{*}=JQ^{T}(-B)$ by Section~0.5\,(a), (d).
Consequently $\left(Q_{k+1}(B)J^{*}\right)^{*}=JQ_{k+1}^{T}(-B)$ for every $k\ge0$; for $k\ge1$ this is $\left(Q(B)\cdot Q_k(B)J^{*}\right)^{*}=JQ_k^{T}(-B)\cdot Q^{T}(-B)$.

\emph{Projections.}
If $\mathrm{Ym}=\mathrm{Xm}^{*}$, then $(\mathrm{Ym}\cdot\mathrm{Xm})^{*}=\mathrm{Xm}^{*}\cdot\mathrm{Ym}^{*}=\mathrm{Ym}\cdot\mathrm{Xm}$.
The operators $JP_k(B)J^{*}$, $JP_{k+1}(B)J^{*}$, $P(B)$, $JP(B)J^{*}$ are of the form $\mathrm{Ym}\cdot\mathrm{Xm}$ with $\mathrm{Ym}=\mathrm{Xm}^{*}$, and they are idempotent by clause~(a); hence they are orthogonal projections.

\emph{The operator $D_B$.}
By (I.2.1), (I.2.2) and (I.2.3),
\begin{align*}
(D_Bf)(x)
&=-\eta_k^{-1}\sum_{i}\frac{w_i}{2}\left(\rho-\gamma\cdot\hat{d}_i(x)\right)\left[\left(R_i(B)f\right)(x)-f(x)\right]\\
&=\eta_k^{-1}\biggl\{\frac{\rho}{2}\sum_{i}w_i\left(f-R_i(B)f\right)(x)+\frac{1}{2}\sum_{i}w_i\,\gamma\cdot\hat{d}_i(x)\left(R_i(B)f\right)(x)\\
&\qquad\qquad-\frac{1}{2}\,\gamma\cdot\left(\sum_{i}w_i\hat{d}_i(x)\right)f(x)\biggr\}.
\end{align*}
The last term vanishes by (I.s3), the closure relation.
This is (I.2.4):
\[
D_B=\eta_k^{-1}\left[\frac{\rho}{2}\nabla_B^{*}\nabla_B+\frac{1}{2}J_B\right].
\]
One has $R_i(B)^{2}=1$: by (I.s1) and (21.12), $\left(R_i(B)^{2}f\right)(x)=U_B(x,x^{(i)})\cdot U_B(x^{(i)},x)f(x)=f(x)$.
One has $R_i(B)^{*}=R_i(B)$: by Section~0.4\,(c) and $\overline{U_B(x^{(i)},x)}=U_B(x^{(i)},x)^{-1}=U_B(x,x^{(i)})$,
\[
\left(f'',R_i(B)f\right)^{\mathrm{ch}}
=\int\mathrm{d}x\,\overline{f''(x)}\cdot U_B(x,x^{(i)})f(x^{(i)})
=\int\mathrm{d}x\,\overline{f''(x^{(i)})}\cdot U_B(x^{(i)},x)f(x)
=\int\mathrm{d}x\,\overline{U_B(x,x^{(i)})f''(x^{(i)})}\cdot f(x).
\]
Hence $\left(1-R_i(B)\right)^{*}=1-R_i(B)$ and $\left(1-R_i(B)\right)^{2}=2\left(1-R_i(B)\right)$, so that
\[
\left(f,\left(1-R_i(B)\right)f\right)^{\mathrm{ch}}=\frac{1}{2}\left(\left(1-R_i(B)\right)f,\left(1-R_i(B)\right)f\right)^{\mathrm{ch}}\ge0.
\]
Since $w_i>0$, it follows that $\nabla_B^{*}\nabla_B=\left(\nabla_B^{*}\nabla_B\right)^{*}\ge0$.

Let $\mathrm{Gd}_i$ denote the multiplication by the matrix function $\gamma\cdot\hat{d}_i(x)$.
It is Hermitian, because $\hat{d}_i(x)$ is real and the $\gamma_\mu$ are Hermitian (I.c1).
By (I.s2),
\[
\left(R_i(B)\cdot\mathrm{Gd}_if\right)(x)=U_B(x,x^{(i)})\,\gamma\cdot\hat{d}_i(x^{(i)})f(x^{(i)})=-\gamma\cdot\hat{d}_i(x)\left(R_i(B)f\right)(x).
\]
Therefore
\[
J_B^{*}=\sum_{i}w_i\left(\mathrm{Gd}_i\cdot R_i(B)\right)^{*}=\sum_{i}w_i\,R_i(B)\cdot\mathrm{Gd}_i=-\sum_{i}w_i\,\mathrm{Gd}_i\cdot R_i(B)=-J_B.
\]
It follows that $D_B^{*}=\eta_k^{-1}\left[\frac{\rho}{2}\nabla_B^{*}\nabla_B-\frac{1}{2}J_B\right]$ and, because $\rho>0$,
\[
\operatorname{Re}D_B=\frac{1}{2}\left(D_B+D_B^{*}\right)=\frac{\rho}{2\eta_k}\nabla_B^{*}\nabla_B\ge0.
\]
The matrix $\gamma_5$ commutes with $R_i(B)$ and anticommutes with $\gamma\cdot\hat{d}_i(x)$ (I.c1), and $\gamma_5^2=1$.
Hence $\gamma_5\cdot\nabla_B^{*}\nabla_B\cdot\gamma_5=\nabla_B^{*}\nabla_B$, $\gamma_5\cdot J_B\cdot\gamma_5=-J_B$, and $\gamma_5\cdot D_B\cdot\gamma_5=D_B^{*}$.

\paragraph*{(f) (scaling).}
Let $B=B'_L$, that is, $B(b)=L^{-\frac{1}{2}}B'(b/L)$ (Section~I.3.7).

(f1) One has $e_{k+1}\cdot\eta_{k+1}=L^{\frac{1}{2}}\cdot e_k\cdot\eta_k/L=L^{-\frac{1}{2}}\cdot e_k\cdot\eta_k$.
Hence $e_k\cdot\eta_k\cdot B(b)=e_k\cdot\eta_k\cdot L^{-\frac{1}{2}}\cdot B'(b/L)=e_{k+1}\cdot\eta_{k+1}\cdot B'(b/L)$ for every bond $b$.

(f2) \emph{Phases.}
By (f1) and Section~0.3\,(e), $e_k\cdot\eta_k\cdot B\left(\boldsymbol{\Gamma}^{\pi}(c,c')\right)=e_{k+1}\cdot\eta_{k+1}\cdot B'\left(\boldsymbol{\Gamma}^{\pi}(c/L,c'/L)\right)$; hence $e_k\cdot\eta_k\cdot(\tau B)(c,c')=e_{k+1}\cdot\eta_{k+1}\cdot(\tau B')(c/L,c'/L)$.
Therefore, the right-hand sides being formed at the level $k+1$ (charge $e_{k+1}$, mesh $\eta_{k+1}$),
\begin{equation}
\begin{aligned}
&u_B(c,c')=u_{B'}(c/L,c'/L)\quad\text{(bonds)},\qquad U_B(c,c')=U_{B'}(c/L,c'/L)\quad(|c-c'|\le M),\\
&\varphi^{+}_B(y',c)=\varphi_{B'}(y'/L,c/L)\quad\left(c\in\mathrm{Bl}_{k+1}(y')\right).
\end{aligned}
\tag{21.26}
\end{equation}
For the third relation put $y_j=y_j(c)$ for $0\le j\le k+1$; then $y_k=[c]$ and $y_{k+1}=y'$ (Section~0.2\,(c)).
By (I.3.3),
\begin{align*}
\varphi^{+}_B(y',c)
&=U_B(y_{k+1},y_k)\cdot\varphi_B(y_k,c)
=\exp\left\{\mathrm{i}\,e_k\cdot\eta_k\sum_{j<k+1}(\tau B)(y_{j+1},y_j)\right\}\\
&=\exp\left\{\mathrm{i}\,e_{k+1}\cdot\eta_{k+1}\sum_{j<k+1}(\tau B')(y_{j+1}/L,y_j/L)\right\},
\end{align*}
and the points $y_j/L=y_j(c/L)$ are the centres of the blocks of the level $k+1$, of the sides $\eta_{k+1},L\eta_{k+1},\dots,1$, that contain $c/L$ (Section~0.2\,(d)).
The last expression is $\varphi_{B'}(y'/L,c/L)$, defined by (I.3.3) with $k+1$ in place of $k$.
(At $k=0$ the sum consists of the single term $j=0$: one has $\varphi^{+}_B(y',c)=U_B(y',c)$ and $\varphi_{B'}(y'/L,c/L)=\exp\left\{\mathrm{i}\,e_{k+1}\cdot\eta_1\cdot(\tau B')(y'/L,c/L)\right\}$.)

(f3) \emph{Sites.}
One has $x/L=(c/L,g)$.
By (I.1.2) at the two levels, $(x/L)^{(i)}=(c/L,gs_i)=x^{(i)}/L$ for $i\le2$, and
\[
(x/L)^{(3)}=\left(c/L+\eta_{k+1}ge_1,gs_3\right)=\left((c+\eta_kge_1)/L,gs_3\right)=x^{(3)}/L.
\]
Further, $M_{i}(x/L)=M_{i}(x)$, the label being the same, and
\[
d_X(x,x')=\max\left(|c-c'|,\eta_k\right)=L\max\left(|c/L-c'/L|,\eta_{k+1}\right)=L\,d_X(x/L,x'/L).
\]
Hence, if $(x,x')\in\mathrm{Pzt}_k$, then $x/L\ne x'/L$ and $d_X(x/L,x'/L)\le M/L\le M$, so that $(x/L,x'/L)$ lies in $\mathrm{Pzt}_{k+1}$.

(f4) \emph{Bijections.}
For $f\in\mathfrak{S}^{\mathrm{ch}}(\eta_{k+1},\ell/L)$ one has
\[
f_L(c+\ell e_{\mu},g)=L^{-1}f\left(c/L+(\ell/L)e_{\mu},g\right)=-f_L(c,g),
\]
so that $f_L\in\mathfrak{S}^{\mathrm{ch}}(\eta_k,\ell)$.
With $L(\tilde{c},g):=(L\tilde{c},g)$, the map that sends $\tilde{f}\in\mathfrak{S}^{\mathrm{ch}}(\eta_k,\ell)$ to the function $\tilde{x}\mapsto L\tilde{f}(L\tilde{x})$ in $\mathfrak{S}^{\mathrm{ch}}(\eta_{k+1},\ell/L)$ is the inverse of $f\mapsto f_L$.
For $v\in\mathfrak{S}(1,\ell/L)$ one has
\[
v_L(y'+\ell e_{\mu})=L^{-1}v\left(y'/L+(\ell/L)e_{\mu}\right)=-v_L(y'),
\]
so that $v_L\in\mathfrak{S}(L,\ell)$, and the map that sends $\tilde{v}\in\mathfrak{S}(L,\ell)$ to the function $y\mapsto L\tilde{v}(Ly)$ in $\mathfrak{S}(1,\ell/L)$ is the inverse of $v\mapsto v_L$.
Both maps $f\mapsto f_L$ and $v\mapsto v_L$ are denoted by $\mathrm{Rsc}$.
They commute with constant matrices.

(f5) \emph{Measures.}
The map $c\mapsto c/L$ sends a set of representatives of $\mathbb{L}_k/\ell\mathbb{Z}^3$ onto a set of representatives of $\mathbb{L}_{k+1}/(\ell/L)\mathbb{Z}^3$, and $\eta_k^3/48=L^3\eta_{k+1}^3/48$.
Hence, the integral on the left being taken at the level $k$ and the one on the right at the level $k+1$,
\[
\int\mathrm{d}x\,\mathrm{Fp}(x/L)=L^3\int\mathrm{d}\tilde{x}\,\mathrm{Fp}(\tilde{x})
\]
for every $(\ell/L)\mathbb{Z}^3$-periodic function $\mathrm{Fp}$ on $\mathbb{L}^{\mathrm{f}}_{k+1}$, and
\[
\int_{\mathrm{Bl}_{k+1}(y')}\mathrm{d}x\,\mathrm{Fp}(x/L)=L^3\int_{\Delta_{y'/L}}\mathrm{d}\tilde{x}\,\mathrm{Fp}(\tilde{x})
\]
for every function $\mathrm{Fp}$ (Section~0.2\,(d)).

\emph{The operator $D_B+\mu$.}
By (f3), (21.26), (I.2.3) and $\eta_k^{-1}=L^{-1}\cdot\eta_{k+1}^{-1}$,
\begin{align*}
(D_Bf_L)(x)
&=-\eta_k^{-1}\sum_{i}M_{i}(x)\left[U_B(x,x^{(i)})\cdot L^{-1}f(x^{(i)}/L)-L^{-1}f(x/L)\right]\\
&=L^{-1}\cdot L^{-1}\cdot\left(-\eta_{k+1}^{-1}\right)\sum_{i}M_{i}(x/L)\left[U_{B'}\left(x/L,(x/L)^{(i)}\right)f\left((x/L)^{(i)}\right)-f(x/L)\right]\\
&=L^{-1}\cdot L^{-1}\left(D_{B'}f\right)(x/L)=L^{-1}\left(\left(D_{B'}f\right)_L\right)(x),
\end{align*}
and $\mu f_L=L^{-1}(L\mu f)_L$.

\emph{The averaging operators.}
By clause~(b), (21.26) and (f5),
\begin{align*}
\left(Q_{k+1}(B)J^{*}f_L\right)(y')
&=L^{-3}\int_{\mathrm{Bl}_{k+1}(y')}\mathrm{d}x\,\varphi^{+}_B(y',c)\cdot L^{-1}f(x/L)
=L^{-1}\int_{\Delta_{y'/L}}\mathrm{d}\tilde{x}\,\varphi_{B'}(y'/L,\tilde{c})f(\tilde{x})\\
&=L^{-1}\left(Q_{k+1}(B')J^{*}f\right)(y'/L)=\left(\left(Q_{k+1}(B')J^{*}f\right)_L\right)(y'),
\end{align*}
where $\tilde{x}=(\tilde{c},g)$.
The second and the third expression agree by (21.15) at the level $k+1$, which gives
\[
\int_{\Delta_{y'/L}}\mathrm{d}\tilde{x}\,\varphi_{B'}(y'/L,\tilde{c})f(\tilde{x})
=\eta_{k+1}^3\sum_{\tilde{c}\in\Delta_{y'/L}}\varphi_{B'}(y'/L,\tilde{c})\,\frac{1}{48}\sum_{g\in O_{\mathrm{cub}}}f(\tilde{c},g)
=\left(Q_{k+1}(B')J^{*}f\right)(y'/L).
\]
At the level $0$ the computation, written out, reads
\begin{align*}
\left(Q(B)J^{*}f_L\right)(y')
&=L^{-3}\sum_{c\in B(y')}U_B(y',c)\,\frac{1}{48}\sum_{g\in O_{\mathrm{cub}}}L^{-1}f(c/L,g)\\
&=L^{-1}\cdot\eta_1^3\sum_{\tilde{c}\in\Delta_{y'/L}}\varphi_{B'}(y'/L,\tilde{c})\,\frac{1}{48}\sum_{g\in O_{\mathrm{cub}}}f(\tilde{c},g),
\end{align*}
because $\eta_1^3=L^{-3}$ and $B(y')/L=\Delta_{y'/L}$ is the unit block of the level $1$.
For $c\in\mathrm{Bl}_{k+1}(y')$, by clause~(b), (21.26) and (21.16) at the level $k+1$,
\begin{align*}
\left(JQ_{k+1}^{T}(-B)v_L\right)(x)
&=\varphi^{+}_B(y',c)^{-1}\cdot L^{-1}v(y'/L)
=L^{-1}\cdot\varphi_{B'}(y'/L,c/L)^{-1}v(y'/L)\\
&=L^{-1}\left(JQ_{k+1}^{T}(-B')v\right)(x/L)
=\left(\left(JQ_{k+1}^{T}(-B')v\right)_L\right)(x).
\end{align*}
With clause~(b) and $P_{k+1}(B')=Q_{k+1}^{T}(-B')\cdot Q_{k+1}(B')$ (Section~I.3.3 at the level $k+1$) this gives
\[
JP_{k+1}(B)J^{*}f_L=JQ_{k+1}^{T}(-B)\left(Q_{k+1}(B')J^{*}f\right)_L=\left(JQ_{k+1}^{T}(-B')\cdot Q_{k+1}(B')J^{*}f\right)_L=\left(JP_{k+1}(B')J^{*}f\right)_L.
\]

\emph{Pairings.}
By (f5),
\[
\left\langle f''_L,f_L\right\rangle^{\mathrm{ch}}
=\int\mathrm{d}x\,L^{-2}f''(x/L)\cdot f(x/L)
=L^{-2}\cdot L^{3}\int\mathrm{d}\tilde{x}\,f''(\tilde{x})\cdot f(\tilde{x})
=L\left\langle f'',f\right\rangle^{\mathrm{ch}}.
\]
For $v'',v\in\mathfrak{S}(1,\ell/L)$,
\[
\left\langle v''_L,v_L\right\rangle_L
=L^3\sum_{y'\in\mathrm{Fdm}^L}L^{-2}v''(y'/L)\cdot v(y'/L)
=L\sum_{y\in\mathrm{Fdm}^L/L}v''(y)\cdot v(y)
=L\left\langle v'',v\right\rangle_1,
\]
since $\mathrm{Fdm}^L/L$ is a set of representatives of $\mathbb{Z}^3/(\ell/L)\mathbb{Z}^3$.

\emph{The H\"older quotient.}
For $(x,x')\in\mathrm{Pzt}_k$, by (f3), (21.26) and (I.1.7),
\begin{align*}
\left(\delta_{\alpha,B}f_L\right)(x,x')
&=d_X(x,x')^{-\alpha}\left(U_B(x,x')\cdot L^{-1}f(x'/L)-L^{-1}f(x/L)\right)\\
&=L^{-\alpha}\cdot L^{-1}\cdot d_X(x/L,x'/L)^{-\alpha}\left(U_{B'}(x/L,x'/L)f(x'/L)-f(x/L)\right)\\
&=L^{-1-\alpha}\left(\delta_{\alpha,B'}f\right)(x/L,x'/L)=\left(\left(\delta_{\alpha,B'}f\right)_L\right)(x,x').
\end{align*}
For a pair inside a cluster this reads $\eta_k^{-\alpha}\cdot L^{-1}\left(f(x'/L)-f(x/L)\right)=L^{-1-\alpha}\cdot\eta_{k+1}^{-\alpha}\left(f(x'/L)-f(x/L)\right)$.
The quotient $\delta_{\alpha,B'}$ is defined on the whole of $\mathrm{Pzt}_{k+1}$, with unique shortest lifts on the torus of side $\ell/L$, because $2M<\ell/L$ (Section~0.1).

\emph{Fields.}
All these relations are linear in $f$, $v$, or bilinear in $(f'',f)$, $(v'',v)$ with the first factor on the left; by Section~0.6\,(a), (d) they hold for fields.
\qed

\subsection*{2. Proposition 21.2}

\begin{proof}[Proof of Proposition 21.2]
For fields the rules of Section~0.6\,(a)--(c) are used; the transposes are those of Proposition~21.1\,(c) in the form (21.24).

\textbf{(a)}
Let $k\ge1$ and assume that $S_k(B)$ exists.
By Section~0.6\,(c) the first term on the left-hand side is
\begin{multline*}
b_{k}\left\langle\bar{\Psi},\Psi\right\rangle_1-b_{k}\left\langle\bar{\Psi},Q_k(B)J^{*}\psi\right\rangle_1\\
-b_{k}\left\langle Q_k(-B)J^{*}\bar{\psi},\Psi\right\rangle_1+b_{k}\left\langle Q_k(-B)J^{*}\bar{\psi},Q_k(B)J^{*}\psi\right\rangle_1.
\end{multline*}
By (21.24), $Q_k(-B)J^{*}=\left(JQ_k^{T}(-B)\right)^{T}$.
Section~0.6\,(b), applied to $O=JQ_k^{T}(-B)\colon\mathfrak{S}(1,\ell)\to\mathfrak{S}^{\mathrm{ch}}(\eta_k,\ell)$, therefore gives for every field $\mathrm{Phf}$ on $T^0_k$
\begin{equation}
\left\langle Q_k(-B)J^{*}\bar{\psi},\mathrm{Phf}\right\rangle_1=\left\langle\bar{\psi},JQ_k^{T}(-B)\,\mathrm{Phf}\right\rangle^{\mathrm{ch}}.
\tag{21.27}
\end{equation}
With $\mathrm{Phf}=Q_k(B)J^{*}\psi$ the fourth term becomes
\[
b_{k}\left\langle\bar{\psi},JQ_k^{T}(-B)\cdot Q_k(B)J^{*}\psi\right\rangle^{\mathrm{ch}}=\left\langle\bar{\psi},b_{k}\,JP_k(B)J^{*}\psi\right\rangle^{\mathrm{ch}}.
\]
Added to $\left\langle\bar{\psi},(D_B+\mu)\psi\right\rangle^{\mathrm{ch}}$, it gives $\left\langle\bar{\psi},F_{\infty}(B)\psi\right\rangle^{\mathrm{ch}}$.
Since $\psi=H_{k}(B)\Psi=b_{k}\cdot S_k(B)\cdot JQ_k^{T}(-B)\Psi$, one has $F_{\infty}(B)\psi=b_{k}\,JQ_k^{T}(-B)\Psi$, and hence, by (21.27) with $\mathrm{Phf}=\Psi$,
\[
\left\langle\bar{\psi},F_{\infty}(B)\psi\right\rangle^{\mathrm{ch}}=b_{k}\left\langle\bar{\psi},JQ_k^{T}(-B)\Psi\right\rangle^{\mathrm{ch}}=b_{k}\left\langle Q_k(-B)J^{*}\bar{\psi},\Psi\right\rangle_1.
\]
This cancels the third term.
What remains is
\begin{align*}
&b_{k}\left\langle\bar{\Psi},\Psi\right\rangle_1-b_{k}\left\langle\bar{\Psi},Q_k(B)J^{*}\cdot b_{k}\cdot S_k(B)\cdot JQ_k^{T}(-B)\Psi\right\rangle_1\\
&\qquad=\left\langle\bar{\Psi},\left(b_{k}-b_{k}^{2}\cdot Q_k(B)J^{*}\cdot S_k(B)\cdot JQ_k^{T}(-B)\right)\Psi\right\rangle_1
=\left\langle\bar{\Psi},D_k(B)\Psi\right\rangle_1.
\end{align*}
No assumption on $\bar{\psi}$ has been used.

\textbf{(b)}
Let $k\ge0$, put $\mathrm{bLk}:=b_{k+1}/L$, and assume that $S^{0}_{k+1}(B)=\left(D_B+\mu+\mathrm{bLk}\cdot JP_{k+1}(B)J^{*}\right)^{-1}$ exists.
The computation uses (21.24) and Proposition~21.1\,(b), which hold for every $k\ge0$.
At $k=0$ its symbols are
\begin{gather*}
\mathrm{bLk}=b_1/L=b_{L},\quad Q_1(\pm B)J^{*}=Q(\pm B)J^{*},\quad JQ_1^{T}(-B)=JQ^{T}(-B),\\
JP_1(B)J^{*}=JP(B)J^{*},\quad S^{0}_{1}(B)=\Gamma_0(B);
\end{gather*}
moreover $\psi$ and $\bar{\psi}$ are then fields on $T^{\mathrm{f}}_0$, and $\langle\cdot,\cdot\rangle^{\mathrm{ch}}$ is the pairing with the weight $1/48$ on $\mathbb{Z}^3\times O_{\mathrm{cub}}$.
The first term on the left-hand side is
\begin{multline*}
\mathrm{bLk}\cdot\left\langle\bar{\Phi},\Phi\right\rangle_L-\mathrm{bLk}\cdot\left\langle\bar{\Phi},Q_{k+1}(B)J^{*}\psi\right\rangle_L\\
-\mathrm{bLk}\cdot\left\langle Q_{k+1}(-B)J^{*}\bar{\psi},\Phi\right\rangle_L+\mathrm{bLk}\cdot\left\langle Q_{k+1}(-B)J^{*}\bar{\psi},Q_{k+1}(B)J^{*}\psi\right\rangle_L.
\end{multline*}
By (21.24), $Q_{k+1}(-B)J^{*}=\left(JQ_{k+1}^{T}(-B)\right)^{T}$ for the pairings $\langle\cdot,\cdot\rangle_L$ and $\langle\cdot,\cdot\rangle^{\mathrm{ch}}$.
Hence, for every field $\mathrm{Phf}$ on $L\mathbb{Z}^3/\ell\mathbb{Z}^3$,
\[
\left\langle Q_{k+1}(-B)J^{*}\bar{\psi},\mathrm{Phf}\right\rangle_L=\left\langle\bar{\psi},JQ_{k+1}^{T}(-B)\,\mathrm{Phf}\right\rangle^{\mathrm{ch}}.
\]
By Proposition~21.1\,(b), $JQ_{k+1}^{T}(-B)\cdot Q_{k+1}(B)J^{*}=JP_{k+1}(B)J^{*}$.
Therefore the fourth term and $\left\langle\bar{\psi},(D_B+\mu)\psi\right\rangle^{\mathrm{ch}}$ add up to
\[
\left\langle\bar{\psi},\left(D_B+\mu+\mathrm{bLk}\cdot JP_{k+1}(B)J^{*}\right)\psi\right\rangle^{\mathrm{ch}}
=\mathrm{bLk}\cdot\left\langle\bar{\psi},JQ_{k+1}^{T}(-B)\Phi\right\rangle^{\mathrm{ch}}
=\mathrm{bLk}\cdot\left\langle Q_{k+1}(-B)J^{*}\bar{\psi},\Phi\right\rangle_L,
\]
by the definition of $\psi$ and the preceding relation.
This cancels the third term.
What remains is
\begin{align*}
&\mathrm{bLk}\cdot\left\langle\bar{\Phi},\Phi\right\rangle_L-\mathrm{bLk}\cdot\left\langle\bar{\Phi},Q_{k+1}(B)J^{*}\cdot\mathrm{bLk}\cdot S^{0}_{k+1}(B)\cdot JQ_{k+1}^{T}(-B)\Phi\right\rangle_L\\
&\qquad=\left\langle\bar{\Phi},\left(\mathrm{bLk}-\mathrm{bLk}^{2}\cdot Q_{k+1}(B)J^{*}\cdot S^{0}_{k+1}(B)\cdot JQ_{k+1}^{T}(-B)\right)\Phi\right\rangle_L.
\end{align*}
\end{proof}

\subsection*{3. Proposition 21.3}

\begin{proof}[Proof of Proposition~21.3]
Recall the Schur complement formulas of \cite[Lemma~14.2]{QEDI}, written here in the notation of this chapter.
Let $U$ and $V$ be finite-dimensional complex vector spaces, and let
\[
\mathrm{Mbl}=\left(\begin{array}{cc} m_{11} & m_{12}\\ m_{21} & m_{22}\end{array}\right)
\]
be a linear map of $U\oplus V$, with the blocks $m_{11}\colon U\to U$, $m_{12}\colon V\to U$, $m_{21}\colon U\to V$ and $m_{22}\colon V\to V$.
\begin{itemize}
\item[(LA2.i)] Suppose that $m_{22}$ is invertible, and set $\mathrm{scu}:=m_{11}-m_{12}\cdot m_{22}^{-1}\cdot m_{21}$.
Then $\det\mathrm{Mbl}=\det m_{22}\cdot\det\mathrm{scu}$;
the map $\mathrm{Mbl}$ is invertible if and only if $\mathrm{scu}$ is invertible, and in that case
\[
\mathrm{Mbl}^{-1}=\left(\begin{array}{cc}
\mathrm{scu}^{-1} & -\mathrm{scu}^{-1}\cdot m_{12}\cdot m_{22}^{-1}\\
-m_{22}^{-1}\cdot m_{21}\cdot\mathrm{scu}^{-1} & m_{22}^{-1}+m_{22}^{-1}\cdot m_{21}\cdot\mathrm{scu}^{-1}\cdot m_{12}\cdot m_{22}^{-1}
\end{array}\right).
\]
\item[(LA2.ii)] Suppose that $m_{11}$ is invertible, and set $\mathrm{scv}:=m_{22}-m_{21}\cdot m_{11}^{-1}\cdot m_{12}$.
Then $\det\mathrm{Mbl}=\det m_{11}\cdot\det\mathrm{scv}$;
the map $\mathrm{Mbl}$ is invertible if and only if $\mathrm{scv}$ is invertible, and in that case
\[
\mathrm{Mbl}^{-1}=\left(\begin{array}{cc}
m_{11}^{-1}+m_{11}^{-1}\cdot m_{12}\cdot\mathrm{scv}^{-1}\cdot m_{21}\cdot m_{11}^{-1} & -m_{11}^{-1}\cdot m_{12}\cdot\mathrm{scv}^{-1}\\
-\mathrm{scv}^{-1}\cdot m_{21}\cdot m_{11}^{-1} & \mathrm{scv}^{-1}
\end{array}\right).
\]
\end{itemize}

\emph{Abbreviations.}
All operators are taken at $B$.
Write
\[
\mathrm{Qk}:=Q_k(B)\cdot J^{*},\qquad
\mathrm{Qt}:=J\cdot Q_k^{T}(-B),\qquad
P:=P(B),\qquad
\mathrm{Yg}:=i\gamma_5 y,
\]
where $\mathrm{Yg}$ is a constant $4\times4$ matrix.
The operators $\mathrm{Qk}$ and $\mathrm{Qt}$ are scalar operators;
therefore Proposition~21.1\,(a), Section~I.3.3 and Section~0.7, together with 0.5\,(e), give
\begin{equation}\tag{21.28}
\begin{aligned}
&\mathrm{Qk}\cdot\mathrm{Qt}=1\ \text{on}\ \mathfrak{S}(1,\ell),\qquad
\mathrm{Qt}\cdot\mathrm{Qk}=J\cdot P_k\cdot J^{*},\qquad
\mathrm{Qt}\cdot P\cdot\mathrm{Qk}=J\cdot P_{k+1}\cdot J^{*},\qquad
P^2=P;\\
&\text{the operators }\mathrm{Qk},\ \mathrm{Qt}\text{ and }P\text{ commute with every constant matrix.}
\end{aligned}
\end{equation}

\emph{The matrices.}
By (I.c1) one has $\gamma_5^2=1$, so that $\mathrm{Yg}^2=-y^2$.
For a real number $\mathrm{cf}>0$ it follows that $(\mathrm{cf}+\mathrm{Yg})\cdot(\mathrm{cf}-\mathrm{Yg})=\mathrm{cf}^2+y^2>0$;
consequently $\mathrm{cf}+\mathrm{Yg}$ is invertible, and
\[
(\mathrm{cf}+\mathrm{Yg})^{-1}=\frac{\mathrm{cf}-\mathrm{Yg}}{\mathrm{cf}^2+y^2}.
\]
This is applied with $\mathrm{cf}=b_{k}$ and with $\mathrm{cf}=b_{k}+b_{L}$, which are positive by Section~I.3.4 ($b>0$, $L>1$).
Every matrix that occurs below is a polynomial in $\gamma_5$, and any two of them commute.

\emph{The inverse of $b_{k}+b_{L}\cdot P+\mathrm{Yg}$.}
One has
\[
b_{k}+b_{L}\cdot P+\mathrm{Yg}=(b_{k}+\mathrm{Yg})\cdot(1-P)+(b_{k}+b_{L}+\mathrm{Yg})\cdot P .
\]
From $P\cdot(1-P)=0=(1-P)\cdot P$, $(1-P)^2=1-P$, $P^2=P$ and (21.28) it follows that
\[
\left[(b_{k}+\mathrm{Yg})\cdot(1-P)+(b_{k}+b_{L}+\mathrm{Yg})\cdot P\right]\cdot\left[(b_{k}+\mathrm{Yg})^{-1}\cdot(1-P)+(b_{k}+b_{L}+\mathrm{Yg})^{-1}\cdot P\right]=(1-P)+P=1 .
\]
The factors $(b_{k}+\mathrm{Yg})^{\pm1}$, $(b_{k}+b_{L}+\mathrm{Yg})^{\pm1}$, $P$ and $1-P$ commute with each other:
polynomials in $\gamma_5$ commute, and $P$ commutes with constant matrices by (21.28).
Hence the two brackets commute, and their product in the opposite order is the identity as well.
It follows that $b_{k}+b_{L}\cdot P+i\gamma_5 y$ is invertible on $\mathfrak{S}(1,\ell)$, and that $B_{k,y}$ is given by the formula asserted in Proposition~21.3.

\emph{The operator $b_{k}-b_{k}^2\cdot B_{k,y}$.}
One has
\begin{align*}
b_{k}-b_{k}^2\cdot(b_{k}+\mathrm{Yg})^{-1}
&=b_{k}\cdot\left[(b_{k}+\mathrm{Yg})-b_{k}\right]\cdot(b_{k}+\mathrm{Yg})^{-1}
=b_{k}\cdot\mathrm{Yg}\cdot(b_{k}+\mathrm{Yg})^{-1}
=\boldsymbol{\alpha}_{y},\\
b_{k}-b_{k}^2\cdot(b_{k}+b_{L}+\mathrm{Yg})^{-1}
&=b_{k}\cdot(b_{L}+\mathrm{Yg})\cdot(b_{k}+b_{L}+\mathrm{Yg})^{-1}.
\end{align*}
Hence
\begin{align*}
b_{k}-b_{k}^2\cdot B_{k,y}
&=\boldsymbol{\alpha}_{y}\cdot(1-P)+b_{k}\cdot(b_{L}+\mathrm{Yg})\cdot(b_{k}+b_{L}+\mathrm{Yg})^{-1}\cdot P\\
&=\boldsymbol{\alpha}_{y}+\left[b_{k}\cdot(b_{L}+\mathrm{Yg})\cdot(b_{k}+b_{L}+\mathrm{Yg})^{-1}-b_{k}\cdot\mathrm{Yg}\cdot(b_{k}+\mathrm{Yg})^{-1}\right]\cdot P,
\end{align*}
and the matrix in brackets equals $\boldsymbol{\beta}_{y}$:
\begin{align*}
&b_{k}\cdot(b_{L}+\mathrm{Yg})\cdot(b_{k}+b_{L}+\mathrm{Yg})^{-1}-b_{k}\cdot\mathrm{Yg}\cdot(b_{k}+\mathrm{Yg})^{-1}\\
&\qquad=b_{k}\cdot\left[(b_{L}+\mathrm{Yg})\cdot(b_{k}+\mathrm{Yg})-\mathrm{Yg}\cdot(b_{k}+b_{L}+\mathrm{Yg})\right]\cdot(b_{k}+b_{L}+\mathrm{Yg})^{-1}\cdot(b_{k}+\mathrm{Yg})^{-1}\\
&\qquad=b_{k}\cdot\left[b_{L}\cdot b_{k}\right]\cdot(b_{k}+b_{L}+\mathrm{Yg})^{-1}\cdot(b_{k}+\mathrm{Yg})^{-1}
=\boldsymbol{\beta}_{y},
\end{align*}
because
$(b_{L}+\mathrm{Yg})\cdot(b_{k}+\mathrm{Yg})=b_{L}\cdot b_{k}+b_{L}\cdot\mathrm{Yg}+b_{k}\cdot\mathrm{Yg}+\mathrm{Yg}^2$
and
$\mathrm{Yg}\cdot(b_{k}+b_{L}+\mathrm{Yg})=b_{k}\cdot\mathrm{Yg}+b_{L}\cdot\mathrm{Yg}+\mathrm{Yg}^2$.

\emph{The identity (21.1).}
By Section~I.3.5, $F_{\infty}=D_B+\mu+b_{k}\cdot J\cdot P_k\cdot J^{*}=D_B+\mu+b_{k}\cdot\mathrm{Qt}\cdot\mathrm{Qk}$.
The identity just proved and (21.28) give
\begin{align*}
F_{\infty}-b_{k}^2\cdot\mathrm{Qt}\cdot B_{k,y}\cdot\mathrm{Qk}
&=D_B+\mu+\mathrm{Qt}\cdot\left(b_{k}-b_{k}^2\cdot B_{k,y}\right)\cdot\mathrm{Qk}
=D_B+\mu+\mathrm{Qt}\cdot\left(\boldsymbol{\alpha}_{y}+\boldsymbol{\beta}_{y}\cdot P\right)\cdot\mathrm{Qk}\\
&=D_B+\mu+\boldsymbol{\alpha}_{y}\cdot\mathrm{Qt}\cdot\mathrm{Qk}+\boldsymbol{\beta}_{y}\cdot\mathrm{Qt}\cdot P\cdot\mathrm{Qk}\\
&=D_B+\mu+\boldsymbol{\alpha}_{y}\cdot J\cdot P_k\cdot J^{*}+\boldsymbol{\beta}_{y}\cdot J\cdot P_{k+1}\cdot J^{*}
=F_{y}.
\end{align*}

\emph{The block map.}
Let $U:=\mathfrak{S}(1,\ell)$ and $V:=\mathfrak{S}^{\mathrm{ch}}(\eta_k,\ell)$, and define on $U\oplus V$
\[
\mathrm{Mbl}_y:=\left(\begin{array}{cc} B_{k,y}^{-1} & -b_{k}\cdot\mathrm{Qk}\\ -b_{k}\cdot\mathrm{Qt} & F_{\infty}\end{array}\right),
\qquad
B_{k,y}^{-1}=b_{k}+b_{L}\cdot P+i\gamma_5 y .
\]
The block $m_{22}=F_{\infty}$ is invertible by hypothesis, with the inverse $S_k$.
So (LA2.i) applies, with
\[
\mathrm{scu}=B_{k,y}^{-1}-(-b_{k}\cdot\mathrm{Qk})\cdot S_k\cdot(-b_{k}\cdot\mathrm{Qt})
=b_{k}+b_{L}\cdot P+i\gamma_5 y-b_{k}^2\cdot\mathrm{Qk}\cdot S_k\cdot\mathrm{Qt}
=D_k+b_{L}\cdot P+i\gamma_5 y,
\]
since $D_k=b_{k}-b_{k}^2\cdot\mathrm{Qk}\cdot S_k\cdot\mathrm{Qt}$ by Section~I.3.5.
The block $m_{11}=B_{k,y}^{-1}$ is invertible as well, and (LA2.ii) applies, with
\[
\mathrm{scv}=F_{\infty}-(-b_{k}\cdot\mathrm{Qt})\cdot B_{k,y}\cdot(-b_{k}\cdot\mathrm{Qk})
=F_{\infty}-b_{k}^2\cdot\mathrm{Qt}\cdot B_{k,y}\cdot\mathrm{Qk}
=F_{y}
\]
by (21.1).
Consequently (i) holds if and only if $\mathrm{scu}$ is invertible, that is, if and only if $\mathrm{Mbl}_y$ is invertible, that is, if and only if $\mathrm{scv}$ is invertible, that is, if and only if (ii) holds.
Moreover
\[
\det\mathrm{Mbl}_y=\det F_{\infty}\cdot\det\left(D_k+b_{L}\cdot P+i\gamma_5 y\right)=\det\left(b_{k}+b_{L}\cdot P+i\gamma_5 y\right)\cdot\det F_{y},
\]
which is (21.5); the four determinants are those of endomorphisms of $V$, $U$, $U$ and $V$, in this order.

\emph{The inverse in two forms.}
Suppose now that (i) and (ii) hold, so that $\mathrm{scu}^{-1}=\Gamma_{k,y}$ and $\mathrm{scv}^{-1}=S_{k,y}$.
Since $H_{k}=b_{k}\cdot S_k\cdot\mathrm{Qt}$ and $\mathfrak{H}_{k}=b_{k}\cdot\mathrm{Qk}\cdot S_k$ (Section~I.3.5), formula (LA2.i) yields
\begin{align*}
\mathrm{Mbl}_y^{-1}
&=\left(\begin{array}{cc}
\Gamma_{k,y} & \Gamma_{k,y}\cdot b_{k}\cdot\mathrm{Qk}\cdot S_k\\
S_k\cdot b_{k}\cdot\mathrm{Qt}\cdot\Gamma_{k,y} & S_k+S_k\cdot b_{k}\cdot\mathrm{Qt}\cdot\Gamma_{k,y}\cdot b_{k}\cdot\mathrm{Qk}\cdot S_k
\end{array}\right)\\
&=\left(\begin{array}{cc}
\Gamma_{k,y} & \Gamma_{k,y}\cdot\mathfrak{H}_{k}\\
H_{k}\cdot\Gamma_{k,y} & S_k+H_{k}\cdot\Gamma_{k,y}\cdot\mathfrak{H}_{k}
\end{array}\right),
\end{align*}
and formula (LA2.ii) yields
\[
\mathrm{Mbl}_y^{-1}=\left(\begin{array}{cc}
B_{k,y}+b_{k}^2\cdot B_{k,y}\cdot\mathrm{Qk}\cdot S_{k,y}\cdot\mathrm{Qt}\cdot B_{k,y} & b_{k}\cdot B_{k,y}\cdot\mathrm{Qk}\cdot S_{k,y}\\
b_{k}\cdot S_{k,y}\cdot\mathrm{Qt}\cdot B_{k,y} & S_{k,y}
\end{array}\right).
\]
Equating the entries $(1,1)$, $(2,2)$, $(2,1)$ and $(1,2)$ of the two expressions, one obtains (21.2), (21.3), and the first and the second relation of (21.4), respectively.

\emph{The value $y=0$.}
For $y=0$ one has $\mathrm{Yg}=0$; hence $\Gamma_{k,0}=(D_k+b_{L}\cdot P)^{-1}=\Gamma_k$, $\boldsymbol{\alpha}_{0}=0$, and
\[
\boldsymbol{\beta}_{0}=b_{k}^2\cdot b_{L}\cdot(b_{k}+b_{L})^{-1}\cdot b_{k}^{-1}=\frac{b_{k}\cdot b_{L}}{b_{k}+b_{L}}=\frac{b_{k+1}}{L},
\]
because $b_{k+1}=b\cdot b_{k}/(b_{k}+b_{L})$ and $b_{L}=b/L$ (Section~I.3.4).
Indeed,
\[
b_{k}+b_{L}=\frac{b\cdot\left[(1-L^{-1})+L^{-1}\cdot(1-L^{-k})\right]}{1-L^{-k}}=\frac{b\cdot(1-L^{-k-1})}{1-L^{-k}},
\qquad\text{so that}\qquad
\frac{b\cdot b_{k}}{b_{k}+b_{L}}=\frac{b\cdot(1-L^{-1})}{1-L^{-k-1}} .
\]
Therefore $F_{0}=D_B+\mu+(b_{k+1}/L)\cdot J\cdot P_{k+1}\cdot J^{*}$, and $S_{k,0}=F_{0}^{-1}=S^{0}_{k+1}$ by Section~I.3.5.

The computation is valid for every real $y$, whatever its sign.
Hence the assertions hold for every datum $(k,\ell,\mu,e_k,B,y)$ with $F_{\infty}(B)$ invertible.
Of the constants of Section~I.8 they involve $b$ and $L$ (through $b_{k}$, $b_{L}$, $\boldsymbol{\alpha}_{y}$, $\boldsymbol{\beta}_{y}$) and, through $D_B$ and $\mathrm{Cl}$, the parameter $\rho$, the weights and $M$;
none of these depends on the datum.
\end{proof}

\subsection*{4. Proposition 21.4}

\begin{proof}[Proof of Proposition~21.4]
Recall the determinant identity of Sylvester \cite[Lemma~14.3]{QEDI}:
for finite-dimensional complex vector spaces $U$, $V$ and linear maps $\mathrm{Xm}\colon V\to U$, $\mathrm{Ym}\colon U\to V$,
\begin{equation}\tag{LA3}
\det\left(1_U-\mathrm{Xm}\cdot\mathrm{Ym}\right)=\det\left(1_V-\mathrm{Ym}\cdot\mathrm{Xm}\right).
\end{equation}

\emph{Abbreviations.}
Let $\mathrm{bLk}:=b_{k+1}/L$.
By Section~I.3.4, $\mathrm{bLk}=b_{L}\cdot b_{k}/(b_{k}+b_{L})$ for $k\ge1$, and $b_{1}/L=b/L=b_{L}$.
For $k\ge0$ set
\[
\mathrm{Qkk}:=Q_{k+1}(B)\cdot J^{*},\qquad \mathrm{Qtt}:=J\cdot Q_{k+1}^{T}(-B);
\]
for $k\ge1$ set moreover $\mathrm{Qk}:=Q_k(B)\cdot J^{*}$ and $\mathrm{Qt}:=J\cdot Q_k^{T}(-B)$, so that $\mathrm{Qkk}=Q(B)\cdot\mathrm{Qk}$ and $\mathrm{Qtt}=\mathrm{Qt}\cdot Q^{T}(-B)$;
for $k=0$ one has $\mathrm{Qkk}=Q(B)\cdot J^{*}$ and $\mathrm{Qtt}=J\cdot Q^{T}(-B)$.
As in the proof of Proposition~21.1\,(f), $\mathrm{Rsc}$ denotes the map $f\mapsto f_L$, respectively $v\mapsto v_L$.

In Steps 0--6 and in Step 8, $B'$ is such that $\Gamma_k(B)$ exists and, if $k\ge1$, $S_k(B)$ exists;
this condition is called \emph{the supposition}.

\textbf{Step 0 (the supposition holds at $-B$).}
One has $(-B')_L=-B'_L=-B$.
By Proposition~21.1\,(c), for $k\ge1$ the operator $F_{\infty}(-B)$ is invertible because $F_{\infty}(B)$ is, and $D_k(-B)+b_{L}\cdot P(-B)$ is invertible because $D_k(B)+b_{L}\cdot P(B)$ is;
for $k=0$ the operator $F_{0}(-B)$ is invertible because $F_{0}(B)$ is.
Hence the pair $(-B',-B)$ satisfies the supposition, and every assertion proved below for $(B',B)$ holds also for $(-B',-B)$.

\textbf{Step 1 (existence; (21.6)).}
Let $k\ge1$.
Then $F_{\infty}(B)$ and $D_k(B)+b_{L}\cdot P(B)$ are invertible.
By Proposition~21.3 at $y=0$, in the direction (i)$\,\Rightarrow\,$(ii), the operator $F_{0}(B)=D_B+\mu+(b_{k+1}/L)\cdot J\cdot P_{k+1}(B)\cdot J^{*}$ is invertible, and its inverse is $S^{0}_{k+1}(B)$.
Let $k=0$.
By Section~I.3.9, $S^{0}_{1}(B)=\Gamma_0(B)=F_{0}(B)^{-1}$ with $F_{0}(B)=D_B+\mu+b_{L}\cdot J\cdot P(B)\cdot J^{*}$;
here $b_{L}=b_{1}/L$, and $J\cdot P(B)\cdot J^{*}=J\cdot P_{1}(B)\cdot J^{*}$ by (S21).
Thus, for every $k\ge0$,
\[
F_{0}(B)=D_B+\mu+\mathrm{bLk}\cdot J\cdot P_{k+1}(B)\cdot J^{*},\qquad S^{0}_{k+1}(B)=F_{0}(B)^{-1}.
\]
For $f\in\mathfrak{S}^{\mathrm{ch}}(\eta_{k+1},\ell/L)$, Proposition~21.1\,(f) gives
\begin{equation}\tag{21.29}
\begin{aligned}
F_{0}(B)\,f_L
&=L^{-1}\left(\left(D_{B'}+L\mu\right)f\right)_L+\frac{b_{k+1}}{L}\cdot\left(J\cdot P_{k+1}(B')\cdot J^{*}f\right)_L\\
&=L^{-1}\left(\left(D_{B'}+L\mu+b_{k+1}\cdot J\cdot P_{k+1}(B')\cdot J^{*}\right)f\right)_L
=L^{-1}\left(F_{\infty}(B')\,f\right)_L,
\end{aligned}
\end{equation}
that is, $F_{0}(B)\cdot\mathrm{Rsc}=L^{-1}\cdot\mathrm{Rsc}\cdot F_{\infty}(B')$.
The map $\mathrm{Rsc}$ is a bijection of $\mathfrak{S}^{\mathrm{ch}}(\eta_{k+1},\ell/L)$ onto $\mathfrak{S}^{\mathrm{ch}}(\eta_k,\ell)$.
Hence $F_{\infty}(B')=L\cdot\mathrm{Rsc}^{-1}\cdot F_{0}(B)\cdot\mathrm{Rsc}$ is invertible, $S_{k+1}(B')=L^{-1}\cdot\mathrm{Rsc}^{-1}\cdot S^{0}_{k+1}(B)\cdot\mathrm{Rsc}$, and $S^{0}_{k+1}(B)\cdot\mathrm{Rsc}=L\cdot\mathrm{Rsc}\cdot S_{k+1}(B')$;
this is (21.6).
(The identity (21.29) involves no inverse:
if $F_{0}(B)$ is read as $D_B+\mu+\mathrm{bLk}\cdot J\cdot P_{k+1}(B)\cdot J^{*}$, it holds for every complex bond function $B'$ on $T_{k+1}$.
Consequently $S_{k+1}(B')$ exists if and only if $F_{0}(B)$ is invertible.)
By Step 0, (21.29) holds for the pair $(-B',-B)$, and $S_{k+1}(-B')$ exists.
The displays of Section~I.3.5 at the level $k+1$,
\begin{gather*}
H_{k+1}(B')=b_{k+1}\cdot S_{k+1}(B')\cdot J\cdot Q_{k+1}^{T}(-B'),\\
D_{k+1}(B')=b_{k+1}-b_{k+1}^2\cdot Q_{k+1}(B')\cdot J^{*}\cdot S_{k+1}(B')\cdot J\cdot Q_{k+1}^{T}(-B'),\\
\bar{H}_{k+1}(B')=\mathsf{C}^{-1}\cdot H_{k+1}(-B')\cdot\mathsf{C},
\end{gather*}
contain no other inverse; hence these operators exist.
No condition on $B'$ beyond the supposition has been imposed.

\textbf{Step 2 ((21.7a), (21.8a)).}
Let $k\ge1$, and let $P=P(B)$.
By Proposition~21.3 at $y=0$, $B_{k}=b_{k}^{-1}\cdot(1-P)+(b_{k}+b_{L})^{-1}\cdot P$.
By Proposition~21.1\,(a), $Q(B)\cdot P=Q(B)$ and $P\cdot Q^{T}(-B)=Q^{T}(-B)$;
hence $Q(B)\cdot(1-P)=0$, $(1-P)\cdot Q^{T}(-B)=0$, and
\begin{equation}\tag{21.30}
Q(B)\cdot B_{k}=(b_{k}+b_{L})^{-1}\cdot Q(B),\qquad B_{k}\cdot Q^{T}(-B)=(b_{k}+b_{L})^{-1}\cdot Q^{T}(-B).
\end{equation}
By (21.2) at $y=0$, $\Gamma_k=B_{k}+b_{k}^2\cdot B_{k}\cdot\mathrm{Qk}\cdot S^{0}_{k+1}\cdot\mathrm{Qt}\cdot B_{k}$.
Together with (21.30) and $Q(B)\cdot Q^{T}(-B)=1$ this gives
\begin{align*}
Q(B)\cdot\Gamma_k\cdot Q^{T}(-B)
&=(b_{k}+b_{L})^{-1}+b_{k}^2\cdot(b_{k}+b_{L})^{-2}\cdot Q(B)\cdot\mathrm{Qk}\cdot S^{0}_{k+1}\cdot\mathrm{Qt}\cdot Q^{T}(-B)\\
&=(b_{k}+b_{L})^{-1}+b_{k}^2\cdot(b_{k}+b_{L})^{-2}\cdot\mathrm{Qkk}\cdot S^{0}_{k+1}\cdot\mathrm{Qtt}.
\end{align*}
It follows that
\begin{align*}
D^{0}_{k+1}(B)
&=b_{L}-b_{L}^2\cdot Q(B)\cdot\Gamma_k\cdot Q^{T}(-B)
=\left[b_{L}-\frac{b_{L}^2}{b_{k}+b_{L}}\right]-\left[\frac{b_{L}\cdot b_{k}}{b_{k}+b_{L}}\right]^2\cdot\mathrm{Qkk}\cdot S^{0}_{k+1}\cdot\mathrm{Qtt}\\
&=\mathrm{bLk}-\mathrm{bLk}^2\cdot\mathrm{Qkk}\cdot S^{0}_{k+1}\cdot\mathrm{Qtt},
\end{align*}
since $b_{L}-b_{L}^2/(b_{k}+b_{L})=b_{L}\cdot b_{k}/(b_{k}+b_{L})=\mathrm{bLk}$.
This is (21.7a).
By the first relation of (21.4) at $y=0$ and by (21.30),
\begin{align*}
b_{L}\cdot H_{k}\cdot\Gamma_k\cdot Q^{T}(-B)
&=b_{L}\cdot b_{k}\cdot S^{0}_{k+1}\cdot\mathrm{Qt}\cdot B_{k}\cdot Q^{T}(-B)
=b_{L}\cdot b_{k}\cdot(b_{k}+b_{L})^{-1}\cdot S^{0}_{k+1}\cdot\mathrm{Qt}\cdot Q^{T}(-B)\\
&=\mathrm{bLk}\cdot S^{0}_{k+1}\cdot\mathrm{Qtt},
\end{align*}
which is (21.8a).
Let $k=0$.
By (S21), $D^{0}_{1}(B)=b_{L}-b_{L}^2\cdot Q(B)\cdot J^{*}\cdot\Gamma_0(B)\cdot J\cdot Q^{T}(-B)$.
Since $b_{L}=b_{1}/L$, $Q(B)\cdot J^{*}=Q_{1}(B)\cdot J^{*}$, $\Gamma_0(B)=S^{0}_{1}(B)$ and $J\cdot Q^{T}(-B)=J\cdot Q_{1}^{T}(-B)$, this is, term by term, the right side of (21.7a).
By the first, the third and the fourth of these four equalities, $b_{L}\cdot\Gamma_0(B)\cdot J\cdot Q^{T}(-B)$ and $(b_{1}/L)\cdot S^{0}_{1}(B)\cdot J\cdot Q_{1}^{T}(-B)$ are one and the same expression.

\textbf{Step 3 ((21.7b), (21.8b)).}
Let $k\ge0$ and $v\in\mathfrak{S}(1,\ell/L)$.
By Proposition~21.1\,(f) and (21.6),
\begin{align*}
\mathrm{Qtt}\,v_L&=\left(J\cdot Q_{k+1}^{T}(-B')\,v\right)_L,\\
S^{0}_{k+1}(B)\cdot\mathrm{Qtt}\,v_L&=L\cdot\left(S_{k+1}(B')\cdot J\cdot Q_{k+1}^{T}(-B')\,v\right)_L,\\
\mathrm{Qkk}\cdot S^{0}_{k+1}(B)\cdot\mathrm{Qtt}\,v_L&=L\cdot\left(Q_{k+1}(B')\cdot J^{*}\cdot S_{k+1}(B')\cdot J\cdot Q_{k+1}^{T}(-B')\,v\right)_L.
\end{align*}
Inserting this into (21.7a), one obtains
\begin{align*}
D^{0}_{k+1}(B)\,v_L
&=\frac{b_{k+1}}{L}\cdot v_L-\left(\frac{b_{k+1}}{L}\right)^2\cdot L\cdot\left(Q_{k+1}(B')\cdot J^{*}\cdot S_{k+1}(B')\cdot J\cdot Q_{k+1}^{T}(-B')\,v\right)_L\\
&=L^{-1}\left(\left(b_{k+1}-b_{k+1}^2\cdot Q_{k+1}(B')\cdot J^{*}\cdot S_{k+1}(B')\cdot J\cdot Q_{k+1}^{T}(-B')\right)v\right)_L\\
&=L^{-1}\left(D_{k+1}(B')\,v\right)_L,
\end{align*}
which is (21.7b).
By (21.8a), the spinor $b_{L}\cdot H_{k}(B)\cdot\Gamma_k(B)\cdot Q^{T}(-B)\,v_L$ for $k\ge1$, respectively $b_{L}\cdot\Gamma_0(B)\cdot J\cdot Q^{T}(-B)\,v_L$ for $k=0$, equals
\[
\frac{b_{k+1}}{L}\cdot S^{0}_{k+1}(B)\cdot\mathrm{Qtt}\,v_L
=\frac{b_{k+1}}{L}\cdot L\cdot\left(S_{k+1}(B')\cdot J\cdot Q_{k+1}^{T}(-B')\,v\right)_L
=\left(H_{k+1}(B')\,v\right)_L ;
\]
this is (21.8b).
By 0.6\,(a), the identities (21.6), (21.7b) and (21.8b) hold also for fields $f$, $v$.
This proves (b).

\textbf{Step 4 ((21.10)).}
The rules for fields applied in this step are 0.6\,(a)--(d).

\emph{The case $k\ge1$.}
Set $\mathrm{Sr}:=b_{L}\cdot Q^{T}(-B)\,\Psi'_L$ and $\overline{\mathrm{Sr}}:=b_{L}\cdot Q^{T}(B)\,\bar{\Psi}'_L$; these are fields on $T^0_k$.
Then $\Psi^{c}=\Gamma_k\,\mathrm{Sr}$ and $\bar{\Psi}^{c}=\Gamma_k^{T}\,\overline{\mathrm{Sr}}$, where $\Gamma_k=\Gamma_k(B)$ and $\Gamma_k^{-1}=D_k(B)+b_{L}\cdot P(B)$.
By Proposition~21.1\,(c), $Q(B)^{T}=Q^{T}(B)$ and $Q^{T}(-B)^{T}=Q(-B)$ (see (21.24)).
Hence, by 0.6\,(b), for all fields $\Psi$, $\bar{\Psi}$ on $T^0_k$,
\begin{gather*}
\left\langle\bar{\Psi}'_L,Q(B)\,\Psi\right\rangle_L=\left\langle Q^{T}(B)\,\bar{\Psi}'_L,\Psi\right\rangle_1,
\qquad
\left\langle Q(-B)\,\bar{\Psi},\Psi'_L\right\rangle_L=\left\langle\bar{\Psi},Q^{T}(-B)\,\Psi'_L\right\rangle_1,\\
\left\langle Q(-B)\,\bar{\Psi},Q(B)\,\Psi\right\rangle_L=\left\langle\bar{\Psi},Q^{T}(-B)\cdot Q(B)\,\Psi\right\rangle_1=\left\langle\bar{\Psi},P(B)\,\Psi\right\rangle_1.
\end{gather*}
Expanding the second term on the left side of (21.10) by 0.6\,(c), one sees that the left side equals
\begin{equation}\tag{21.31}
\begin{aligned}
&\left\langle\bar{\Psi},D_k(B)\,\Psi\right\rangle_1+b_{L}\cdot\left\langle\bar{\Psi}'_L,\Psi'_L\right\rangle_L-\left\langle\overline{\mathrm{Sr}},\Psi\right\rangle_1-\left\langle\bar{\Psi},\mathrm{Sr}\right\rangle_1+b_{L}\cdot\left\langle\bar{\Psi},P(B)\,\Psi\right\rangle_1\\
&\qquad=\left\langle\bar{\Psi},\Gamma_k^{-1}\,\Psi\right\rangle_1-\left\langle\bar{\Psi},\mathrm{Sr}\right\rangle_1-\left\langle\overline{\mathrm{Sr}},\Psi\right\rangle_1+b_{L}\cdot\left\langle\bar{\Psi}'_L,\Psi'_L\right\rangle_L.
\end{aligned}
\end{equation}
Now insert $\Psi=\Gamma_k\,\mathrm{Sr}+W'$ and $\bar{\Psi}=\Gamma_k^{T}\,\overline{\mathrm{Sr}}+\bar{W}'$.
By 0.6\,(b), $\left\langle\Gamma_k^{T}\,\overline{\mathrm{Sr}},\mathrm{Phf}\right\rangle_1=\left\langle\overline{\mathrm{Sr}},\Gamma_k\,\mathrm{Phf}\right\rangle_1$ for every field $\mathrm{Phf}$.
Therefore
\begin{align*}
\left\langle\bar{\Psi},\Gamma_k^{-1}\,\Psi\right\rangle_1
&=\left\langle\overline{\mathrm{Sr}},\Gamma_k\cdot\Gamma_k^{-1}\left(\Gamma_k\,\mathrm{Sr}+W'\right)\right\rangle_1+\left\langle\bar{W}',\Gamma_k^{-1}\left(\Gamma_k\,\mathrm{Sr}+W'\right)\right\rangle_1\\
&=\left\langle\overline{\mathrm{Sr}},\Gamma_k\,\mathrm{Sr}\right\rangle_1+\left\langle\overline{\mathrm{Sr}},W'\right\rangle_1+\left\langle\bar{W}',\mathrm{Sr}\right\rangle_1+\left\langle\bar{W}',\Gamma_k^{-1}\,W'\right\rangle_1,\\
\left\langle\bar{\Psi},\mathrm{Sr}\right\rangle_1
&=\left\langle\overline{\mathrm{Sr}},\Gamma_k\,\mathrm{Sr}\right\rangle_1+\left\langle\bar{W}',\mathrm{Sr}\right\rangle_1,\\
\left\langle\overline{\mathrm{Sr}},\Psi\right\rangle_1
&=\left\langle\overline{\mathrm{Sr}},\Gamma_k\,\mathrm{Sr}\right\rangle_1+\left\langle\overline{\mathrm{Sr}},W'\right\rangle_1.
\end{align*}
Consequently (21.31) equals $\left\langle\bar{W}',\Gamma_k^{-1}\,W'\right\rangle_1-\left\langle\overline{\mathrm{Sr}},\Gamma_k\,\mathrm{Sr}\right\rangle_1+b_{L}\cdot\left\langle\bar{\Psi}'_L,\Psi'_L\right\rangle_L$.
Moreover
\[
\left\langle\overline{\mathrm{Sr}},\Gamma_k\,\mathrm{Sr}\right\rangle_1
=b_{L}^2\cdot\left\langle Q^{T}(B)\,\bar{\Psi}'_L,\Gamma_k\cdot Q^{T}(-B)\,\Psi'_L\right\rangle_1
=b_{L}^2\cdot\left\langle\bar{\Psi}'_L,Q(B)\cdot\Gamma_k\cdot Q^{T}(-B)\,\Psi'_L\right\rangle_L .
\]
So the left side of (21.10) is
\begin{align*}
&\left\langle\bar{W}',\left(D_k(B)+b_{L}\cdot P(B)\right)W'\right\rangle_1+\left\langle\bar{\Psi}'_L,\left(b_{L}-b_{L}^2\cdot Q(B)\cdot\Gamma_k(B)\cdot Q^{T}(-B)\right)\Psi'_L\right\rangle_L\\
&\qquad=\left\langle\bar{W}',\left(D_k(B)+b_{L}\cdot P(B)\right)W'\right\rangle_1+\left\langle\bar{\Psi}'_L,D^{0}_{k+1}(B)\,\Psi'_L\right\rangle_L.
\end{align*}
Finally, by (21.7b) and Proposition~21.1\,(f),
\[
\left\langle\bar{\Psi}'_L,D^{0}_{k+1}(B)\,\Psi'_L\right\rangle_L
=L^{-1}\cdot\left\langle\bar{\Psi}'_L,\left(D_{k+1}(B')\,\Psi'\right)_L\right\rangle_L
=L^{-1}\cdot L\cdot\left\langle\bar{\Psi}',D_{k+1}(B')\,\Psi'\right\rangle_1 .
\]
This is (21.10) for $k\ge1$.

\emph{The case $k=0$ (level 0).}
Here $\Psi$, $\bar{\Psi}$, $W'$, $\bar{W}'$ are fields on $T^{\mathrm{f}}_0$, with the index $(x,\beta)$, $x=(c,g)$, and $192$ components per unit site;
the pairing of the unit lattice is
\[
\left\langle\mathrm{ff}'',\mathrm{ff}'\right\rangle^{\mathrm{ch}}=\frac{1}{48}\sum_{c\in\mathrm{Fdm}}\sum_{g\in O_{\mathrm{cub}}}\mathrm{ff}''(c,g)\cdot\mathrm{ff}'(c,g).
\]
Set $\mathrm{Sr}:=b_{L}\cdot J\cdot Q^{T}(-B)\,\Psi'_L$ and $\overline{\mathrm{Sr}}:=b_{L}\cdot J\cdot Q^{T}(B)\,\bar{\Psi}'_L$;
these are fields on $T^{\mathrm{f}}_0$ that are constant on clusters.
Then $\Psi^{c}=\Gamma_0\,\mathrm{Sr}$ and $\bar{\Psi}^{c}=\Gamma_0^{T}\,\overline{\mathrm{Sr}}$, where $\Gamma_0=\Gamma_0(B)$ and $\Gamma_0^{-1}=F_{0}(B)=D_B+\mu+b_{L}\cdot J\cdot P(B)\cdot J^{*}$.
By Proposition~21.1\,(c), $\left(Q(B)\cdot J^{*}\right)^{T}=J\cdot Q^{T}(B)$ and $\left(J\cdot Q^{T}(-B)\right)^{T}=Q(-B)\cdot J^{*}$.
Hence
\begin{gather*}
\left\langle\bar{\Psi}'_L,Q(B)\cdot J^{*}\,\Psi\right\rangle_L=\left\langle J\cdot Q^{T}(B)\,\bar{\Psi}'_L,\Psi\right\rangle^{\mathrm{ch}},
\qquad
\left\langle Q(-B)\cdot J^{*}\,\bar{\Psi},\Psi'_L\right\rangle_L=\left\langle\bar{\Psi},J\cdot Q^{T}(-B)\,\Psi'_L\right\rangle^{\mathrm{ch}},\\
\left\langle Q(-B)\cdot J^{*}\,\bar{\Psi},Q(B)\cdot J^{*}\,\Psi\right\rangle_L
=\left\langle\bar{\Psi},J\cdot Q^{T}(-B)\cdot Q(B)\cdot J^{*}\,\Psi\right\rangle^{\mathrm{ch}}
=\left\langle\bar{\Psi},J\cdot P(B)\cdot J^{*}\,\Psi\right\rangle^{\mathrm{ch}},
\end{gather*}
and the left side of (21.10) equals
\begin{align*}
&\left\langle\bar{\Psi},\left(D_B+\mu\right)\Psi\right\rangle^{\mathrm{ch}}+b_{L}\cdot\left\langle\bar{\Psi}'_L,\Psi'_L\right\rangle_L-\left\langle\overline{\mathrm{Sr}},\Psi\right\rangle^{\mathrm{ch}}-\left\langle\bar{\Psi},\mathrm{Sr}\right\rangle^{\mathrm{ch}}+b_{L}\cdot\left\langle\bar{\Psi},J\cdot P(B)\cdot J^{*}\,\Psi\right\rangle^{\mathrm{ch}}\\
&\qquad=\left\langle\bar{\Psi},F_{0}(B)\,\Psi\right\rangle^{\mathrm{ch}}-\left\langle\bar{\Psi},\mathrm{Sr}\right\rangle^{\mathrm{ch}}-\left\langle\overline{\mathrm{Sr}},\Psi\right\rangle^{\mathrm{ch}}+b_{L}\cdot\left\langle\bar{\Psi}'_L,\Psi'_L\right\rangle_L.
\end{align*}
Insert $\Psi=\Gamma_0\,\mathrm{Sr}+W'$ and $\bar{\Psi}=\Gamma_0^{T}\,\overline{\mathrm{Sr}}+\bar{W}'$, and use $\left\langle\Gamma_0^{T}\,\overline{\mathrm{Sr}},\mathrm{Phf}\right\rangle^{\mathrm{ch}}=\left\langle\overline{\mathrm{Sr}},\Gamma_0\,\mathrm{Phf}\right\rangle^{\mathrm{ch}}$:
\begin{align*}
\left\langle\bar{\Psi},F_{0}(B)\,\Psi\right\rangle^{\mathrm{ch}}
&=\left\langle\overline{\mathrm{Sr}},\Gamma_0\,\mathrm{Sr}\right\rangle^{\mathrm{ch}}+\left\langle\overline{\mathrm{Sr}},W'\right\rangle^{\mathrm{ch}}+\left\langle\bar{W}',\mathrm{Sr}\right\rangle^{\mathrm{ch}}+\left\langle\bar{W}',F_{0}(B)\,W'\right\rangle^{\mathrm{ch}},\\
\left\langle\bar{\Psi},\mathrm{Sr}\right\rangle^{\mathrm{ch}}
&=\left\langle\overline{\mathrm{Sr}},\Gamma_0\,\mathrm{Sr}\right\rangle^{\mathrm{ch}}+\left\langle\bar{W}',\mathrm{Sr}\right\rangle^{\mathrm{ch}},\\
\left\langle\overline{\mathrm{Sr}},\Psi\right\rangle^{\mathrm{ch}}
&=\left\langle\overline{\mathrm{Sr}},\Gamma_0\,\mathrm{Sr}\right\rangle^{\mathrm{ch}}+\left\langle\overline{\mathrm{Sr}},W'\right\rangle^{\mathrm{ch}}.
\end{align*}
The sum is $\left\langle\bar{W}',F_{0}(B)\,W'\right\rangle^{\mathrm{ch}}-\left\langle\overline{\mathrm{Sr}},\Gamma_0\,\mathrm{Sr}\right\rangle^{\mathrm{ch}}+b_{L}\cdot\left\langle\bar{\Psi}'_L,\Psi'_L\right\rangle_L$, and
\[
\left\langle\overline{\mathrm{Sr}},\Gamma_0\,\mathrm{Sr}\right\rangle^{\mathrm{ch}}
=b_{L}^2\cdot\left\langle J\cdot Q^{T}(B)\,\bar{\Psi}'_L,\Gamma_0\cdot J\cdot Q^{T}(-B)\,\Psi'_L\right\rangle^{\mathrm{ch}}
=b_{L}^2\cdot\left\langle\bar{\Psi}'_L,Q(B)\cdot J^{*}\cdot\Gamma_0\cdot J\cdot Q^{T}(-B)\,\Psi'_L\right\rangle_L .
\]
Therefore the left side of (21.10) is
\begin{align*}
&\left\langle\bar{W}',\left(D_B+\mu+b_{L}\cdot J\cdot P(B)\cdot J^{*}\right)W'\right\rangle^{\mathrm{ch}}+\left\langle\bar{\Psi}'_L,\left(b_{L}-b_{L}^2\cdot Q(B)\cdot J^{*}\cdot\Gamma_0(B)\cdot J\cdot Q^{T}(-B)\right)\Psi'_L\right\rangle_L\\
&\qquad=\left\langle\bar{W}',\left(D_B+\mu+b_{L}\cdot J\cdot P(B)\cdot J^{*}\right)W'\right\rangle^{\mathrm{ch}}+\left\langle\bar{\Psi}'_L,D^{0}_{1}(B)\,\Psi'_L\right\rangle_L,
\end{align*}
and, by (21.7b) and Proposition~21.1\,(f),
\[
\left\langle\bar{\Psi}'_L,D^{0}_{1}(B)\,\Psi'_L\right\rangle_L
=L^{-1}\cdot\left\langle\bar{\Psi}'_L,\left(D_{1}(B')\,\Psi'\right)_L\right\rangle_L
=\left\langle\bar{\Psi}',D_{1}(B')\,\Psi'\right\rangle_1 .
\]
Here $\Psi'$ and $\bar{\Psi}'$ have $4$ components per site of $T^0_1$.

\textbf{Step 5 ((21.11); the last formula of (a)).}

\emph{Unbarred fields.}
Let $k\ge1$.
By (21.8b) for fields,
$\psi^{0}=H_{k}(B)\,\Psi^{c}=b_{L}\cdot H_{k}(B)\cdot\Gamma_k(B)\cdot Q^{T}(-B)\,\Psi'_L=\left(H_{k+1}(B')\,\Psi'\right)_L$,
and $T_k(B)\,\psi^{0}=T_k(B)\cdot H_{k}(B)\,\Psi^{c}=\Psi^{c}$ by Proposition~21.1\,(a).
Let $k=0$.
Then $H_{0}(B)=T_0(B)=1$, and $\psi^{0}=\Psi^{c}=b_{L}\cdot\Gamma_0(B)\cdot J\cdot Q^{T}(-B)\,\Psi'_L=\left(H_{1}(B')\,\Psi'\right)_L$ by (21.8b).

\emph{Barred fields.}
The matrix $\mathsf{C}$ commutes with $\mathrm{Rsc}$ and with the scalar operators $Q^{T}(B)$ and $J$.
By Proposition~21.1\,(c), $\Gamma_k(B)^{T}=\mathsf{C}^{-1}\cdot\Gamma_k(-B)\cdot\mathsf{C}$ for every $k\ge0$.
Hence
\begin{align*}
k\ge1\colon&\quad
\bar{\Psi}^{c}=b_{L}\cdot\mathsf{C}^{-1}\cdot\Gamma_k(-B)\cdot\mathsf{C}\cdot Q^{T}(B)\,\bar{\Psi}'_L
=\mathsf{C}^{-1}\left[b_{L}\cdot\Gamma_k(-B)\cdot Q^{T}(-(-B))\left(\mathsf{C}\,\bar{\Psi}'\right)_L\right],\\
k=0\colon&\quad
\bar{\Psi}^{c}=\mathsf{C}^{-1}\left[b_{L}\cdot\Gamma_0(-B)\cdot J\cdot Q^{T}(-(-B))\left(\mathsf{C}\,\bar{\Psi}'\right)_L\right].
\end{align*}
Here $\mathsf{C}\,\bar{\Psi}'$ is a field on $T^0_{k+1}$.
By Step 0, (21.8b) holds for the pair $(-B',-B)$.
With $\bar{H}_{k}(B)=\mathsf{C}^{-1}\cdot H_{k}(-B)\cdot\mathsf{C}$ for $k\ge1$ and $\bar{H}_{0}(B)=1$ it follows that
\begin{align*}
k\ge1\colon&\quad
\bar{\psi}^{0}=\bar{H}_{k}(B)\,\bar{\Psi}^{c}
=\mathsf{C}^{-1}\left[b_{L}\cdot H_{k}(-B)\cdot\Gamma_k(-B)\cdot Q^{T}(-(-B))\left(\mathsf{C}\,\bar{\Psi}'\right)_L\right]
=\mathsf{C}^{-1}\left(H_{k+1}(-B')\cdot\mathsf{C}\,\bar{\Psi}'\right)_L\\
&\qquad\quad=\left(\mathsf{C}^{-1}\cdot H_{k+1}(-B')\cdot\mathsf{C}\,\bar{\Psi}'\right)_L
=\left(\bar{H}_{k+1}(B')\,\bar{\Psi}'\right)_L,\\
k=0\colon&\quad
\bar{\psi}^{0}=\bar{\Psi}^{c}=\mathsf{C}^{-1}\left(H_{1}(-B')\cdot\mathsf{C}\,\bar{\Psi}'\right)_L=\left(\bar{H}_{1}(B')\,\bar{\Psi}'\right)_L.
\end{align*}
Moreover $\bar{T}_k(B)\,\bar{\psi}^{0}=\bar{T}_k(B)\cdot\bar{H}_{k}(B)\,\bar{\Psi}^{c}=\bar{\Psi}^{c}$ by Proposition~21.1\,(c) for $k\ge1$, and $\bar{T}_0(B)=1$.

\emph{H\"older quotients.}
Proposition~21.1\,(f), applied to the field $H_{k+1}(B')\,\Psi'$ on $T^{\mathrm{f}}_{k+1}$, and to the pair $(-B',-B)$ and the field $\bar{H}_{k+1}(B')\,\bar{\Psi}'$, gives
\[
\delta_{\alpha,B}\,\psi^{0}=\delta_{\alpha,B}\left(H_{k+1}(B')\,\Psi'\right)_L=\left(\delta_{\alpha,B'}\,H_{k+1}(B')\,\Psi'\right)_L,
\qquad
\delta_{\alpha,-B}\,\bar{\psi}^{0}=\left(\delta_{\alpha,-B'}\,\bar{H}_{k+1}(B')\,\bar{\Psi}'\right)_L
\]
at all pairs of $\mathrm{Pzt}_k$;
both sides are antiperiodic functions of pairs, that is, fields on $\mathrm{Pz}_k$.
At a pair inside a cluster, $x=(c,g)$, $x'=(c,g')$ with $g\ne g'$, the first relation reads
\[
\eta_k^{-\alpha}\cdot\left(\psi^{0}(c,g')-\psi^{0}(c,g)\right)
=L^{-1-\alpha}\cdot\eta_{k+1}^{-\alpha}\cdot\left(\left(H_{k+1}(B')\,\Psi'\right)(c/L,g')-\left(H_{k+1}(B')\,\Psi'\right)(c/L,g)\right).
\]

\emph{The formula for $k\ge1$.}
One has $\left(D_k(B)+b_{L}\cdot P(B)\right)\Psi^{c}=\mathrm{Sr}$, with $\mathrm{Sr}=b_{L}\cdot Q^{T}(-B)\,\Psi'_L$.
By Section~I.3.5, $D_k(B)=b_{k}-b_{k}^2\cdot\mathrm{Qk}\cdot S_k(B)\cdot\mathrm{Qt}=b_{k}-b_{k}\cdot\mathrm{Qk}\cdot H_{k}(B)$.
Hence $\left(b_{k}+b_{L}\cdot P\right)\Psi^{c}=\mathrm{Sr}+b_{k}\cdot\mathrm{Qk}\cdot H_{k}(B)\,\Psi^{c}=\mathrm{Sr}+b_{k}\cdot\mathrm{Qk}\,\psi^{0}$,
and with $B_{k}=\left(b_{k}+b_{L}\cdot P\right)^{-1}=b_{k}^{-1}\cdot(1-P)+(b_{k}+b_{L})^{-1}\cdot P$ and (21.30),
\begin{align*}
\Psi^{c}
&=B_{k}\,\mathrm{Sr}+b_{k}\cdot B_{k}\cdot\mathrm{Qk}\,\psi^{0}
=\frac{b_{L}}{b_{k}+b_{L}}\cdot Q^{T}(-B)\,\Psi'_L+(1-P)\cdot\mathrm{Qk}\,\psi^{0}+\frac{b_{k}}{b_{k}+b_{L}}\cdot P\cdot\mathrm{Qk}\,\psi^{0}\\
&=\mathrm{Qk}\,\psi^{0}+\frac{b_{L}}{b_{k}+b_{L}}\cdot Q^{T}(-B)\,\Psi'_L-\frac{b_{L}}{b_{k}+b_{L}}\cdot P\cdot\mathrm{Qk}\,\psi^{0}.
\end{align*}
Since $P\cdot\mathrm{Qk}=Q^{T}(-B)\cdot Q(B)\cdot Q_k(B)\cdot J^{*}=Q^{T}(-B)\cdot Q_{k+1}(B)\cdot J^{*}$, it follows that
\[
\Psi^{c}=Q_k(B)\cdot J^{*}\,\psi^{0}+\frac{b_{L}}{b_{k}+b_{L}}\cdot Q^{T}(-B)\left(\Psi'_L-Q_{k+1}(B)\cdot J^{*}\,\psi^{0}\right).
\]
Together with Step 8 this proves (a);
Step 8 establishes the assertions of (a) for $k=0$, namely $H_{0}=\bar{H}_{0}=T_0=\bar{T}_0=1$ and, at $B=0$ with $F_{0}(0)$ invertible, that $\Psi^{c}$ is constant on clusters only if $\Psi'=0$, and $\bar{\Psi}^{c}$ only if $\bar{\Psi}'=0$.

\textbf{Step 6 ((c)).}
By (21.7b), $D^{0}_{k+1}(B)\cdot\mathrm{Rsc}=L^{-1}\cdot\mathrm{Rsc}\cdot D_{k+1}(B')$, where $\mathrm{Rsc}\colon\mathfrak{S}(1,\ell/L)\to\mathfrak{S}(L,\ell)$ is a bijection;
both spaces have the dimension $4\cdot(\ell/L)^3=4n_{k+1}$.
Hence $D^{0}_{k+1}(B)=L^{-1}\cdot\mathrm{Rsc}\cdot D_{k+1}(B')\cdot\mathrm{Rsc}^{-1}$, and
\begin{equation}\tag{21.32}
\det D^{0}_{k+1}(B)=L^{-4n_{k+1}}\cdot\det D_{k+1}(B').
\end{equation}
Let $k\ge1$.
Apply (LA3) with $U=\mathfrak{S}(L,\ell)$, $V=\mathfrak{S}(1,\ell)$, $\mathrm{Xm}=b_{L}\cdot Q(B)\cdot\Gamma_k(B)$ and $\mathrm{Ym}=Q^{T}(-B)$.
Then $1_U-\mathrm{Xm}\cdot\mathrm{Ym}=b_{L}^{-1}\cdot D^{0}_{k+1}(B)$ and
\[
1_V-\mathrm{Ym}\cdot\mathrm{Xm}=1-b_{L}\cdot P(B)\cdot\Gamma_k(B)=\left(\Gamma_k(B)^{-1}-b_{L}\cdot P(B)\right)\cdot\Gamma_k(B)=D_k(B)\cdot\Gamma_k(B).
\]
Hence $b_{L}^{-4n_{k+1}}\cdot\det D^{0}_{k+1}(B)=\det D_k(B)\cdot\det\Gamma_k(B)=\det D_k(B)/\delta Z_k(B)$;
here $D_k(B)$ and $\Gamma_k(B)^{-1}=D_k(B)+b_{L}\cdot P(B)$ act on $\mathfrak{S}(1,\ell)$, which has the dimension $4n_k$.
Let $k=0$.
Apply (LA3) with $U=\mathfrak{S}(L,\ell)$, $V=\mathfrak{S}^{\mathrm{ch}}(1,\ell)$ (of dimension $192\,n_0$), $\mathrm{Xm}=b_{L}\cdot Q(B)\cdot J^{*}\cdot\Gamma_0(B)$ and $\mathrm{Ym}=J\cdot Q^{T}(-B)$.
Then $1_U-\mathrm{Xm}\cdot\mathrm{Ym}=b_{L}^{-1}\cdot D^{0}_{1}(B)$ and
\[
1_V-\mathrm{Ym}\cdot\mathrm{Xm}=1-b_{L}\cdot J\cdot P(B)\cdot J^{*}\cdot\Gamma_0(B)=\left(F_{0}(B)-b_{L}\cdot J\cdot P(B)\cdot J^{*}\right)\cdot\Gamma_0(B)=\left(D_B+\mu\right)\cdot\Gamma_0(B)=D_0(B)\cdot\Gamma_0(B).
\]
Hence $b_{L}^{-4n_1}\cdot\det D^{0}_{1}(B)=\det D_0(B)/\det F_{0}(B)=\det D_0(B)/\delta Z_0(B)$.
In both cases the exponent is the dimension of $U=\mathfrak{S}(L,\ell)$.
Combined with (21.32), this gives
\[
\det D_{k+1}(B')=L^{4n_{k+1}}\cdot b_{L}^{4n_{k+1}}\cdot\frac{\det D_k(B)}{\delta Z_k(B)}=b^{4n_{k+1}}\cdot\frac{\det D_k(B)}{\delta Z_k(B)},
\]
and $b^{4n_{k+1}}=c'_{k}$ by Section~I.3.4.

\textbf{Step 7 ((d)).}
In this step no supposition on $B$ is made.
The map $B'\mapsto B'_L$ is a linear bijection from the complex bond functions on $T_{k+1}=\mathbb{L}_{k+1}/(\ell/L)\mathbb{Z}^3$ onto those on $T_k$ (0.3\,(e));
its inverse is $B\mapsto B'$, $B'(b')=L^{\frac{1}{2}}\cdot B(Lb')$.
The complex bond functions on $T_k$ form the space $\mathbb{C}^{\mathrm{Nbd}}$, $\mathrm{Nbd}:=3\cdot(\ell/\eta_k)^3$, with the coordinates $B(b)$, where $b$ ranges over the positively oriented bonds of $T_k$.

\emph{Entire entries.}
In each of the spaces $\mathfrak{S}^{\mathrm{ch}}(\eta_k,\ell)$, $\mathfrak{S}(1,\ell)$, $\mathfrak{S}(L,\ell)$, $\mathfrak{S}^{\mathrm{ch}}(\eta_{k+1},\ell/L)$ and $\mathfrak{S}(1,\ell/L)$, choose the basis of the functions that vanish on a set of representatives (0.4) except at one point, where one component equals $1$;
these bases do not depend on $B$.
By (I.2.3) and (I.O1), the entries of the matrix of $D_B$ are finite sums of numbers that do not depend on $B$ and of such numbers multiplied by $u_B(c,c_{3})=\exp\left\{i\cdot e_k\cdot\eta_k\cdot B(c,c_{3})\right\}$.
By (21.15)--(21.19), the entries of $Q_k(\pm B)\cdot J^{*}$, $J\cdot Q_k^{T}(\pm B)$, $Q(\pm B)$ and $Q^{T}(\pm B)$ (for $k=0$: of $Q(\pm B)\cdot J^{*}$ and $J\cdot Q^{T}(\pm B)$) are finite sums of numbers multiplied by $\varphi_B(y,c)^{\pm1}$ or by $U_B(y',y)^{\pm1}$;
for a fixed pair $(x,x')$, the two coefficients of the map $f\mapsto\left(\delta_{\alpha,\pm B}f\right)(x,x')$ are $d_X(x,x')^{-\alpha}\cdot U_B(x,x')^{\pm1}$ and $-d_X(x,x')^{-\alpha}$.
The exponents are linear forms in the coordinates $B(b)$ (0.3\,(b)).
Hence all these entries are entire functions on $\mathbb{C}^{\mathrm{Nbd}}$.
Since sums and products of entire functions are entire, the entries of sums and products of these operators and of constant matrices are entire;
in particular this holds for $F_{\infty}(\pm B)$ ($k\ge1$), $F_{0}(\pm B)$, $P(\pm B)$, $T_k(B)$, $\bar{T}_k(B)$ and, $B'$ being a linear function of $B$, for $D_{B'}$, $Q_{k+1}(\pm B')\cdot J^{*}$, $J\cdot Q_{k+1}^{T}(\pm B')$, $F_{\infty}(\pm B')$ and $\delta_{\alpha,\pm B'}$.

\emph{Inverses.}
If the entries of a square matrix $O(B)$ are analytic on an open set $\mathrm{Opn}$, then $\det O(B)$ is analytic on $\mathrm{Opn}$, the set $\left\{B\in\mathrm{Opn}:\det O(B)\ne0\right\}$ is open, and on this set the entries of $O(B)^{-1}$ are analytic, since by Cramer's rule they are cofactors divided by $\det O(B)$.

\textbf{(d1).}
Let $k\ge1$.
The set $\mathrm{Opn}_1:=\left\{B:\det F_{\infty}(B)\ne0\right\}$ is open, and $S_k(B)$, hence also $D_k(B)$ and $\delta Z_k(B)=\det\left(D_k(B)+b_{L}\cdot P(B)\right)$, are analytic on $\mathrm{Opn}_1$.
The set $\mathrm{Iv}_k(\ell,\mu,e_k)=\left\{B\in\mathrm{Opn}_1:\delta Z_k(B)\ne0\right\}$ is open, and $\Gamma_k(B)$ is analytic on it.
Let $k=0$.
The set $\mathrm{Iv}_k(\ell,\mu,e_k)=\left\{B:\det F_{0}(B)\ne0\right\}$ is open, and $\Gamma_0(B)$ is analytic on it.
By Step 0, $\mathrm{Iv}_k$ is invariant under $B\mapsto-B$.

On $\mathrm{Iv}_k$ the operators $S_k(\pm B)$ and $\Gamma_k(\pm B)$ are analytic, because $B\mapsto-B$ is linear and maps $\mathrm{Iv}_k$ onto itself;
hence so are $H_{k}(B)$, $\bar{H}_{k}(B)=\mathsf{C}^{-1}\cdot H_{k}(-B)\cdot\mathsf{C}$, $D_k(B)$, $\det D_k(B)$, $\delta Z_k(B)$, $\Gamma_k(B)^{T}=\mathsf{C}^{-1}\cdot\Gamma_k(-B)\cdot\mathsf{C}$ and $D^{0}_{k+1}(B)$.
By Step 1, $F_{0}(\pm B)$ and $F_{\infty}(\pm B')$ are invertible for $B\in\mathrm{Iv}_k$, and their entries are entire;
consequently $S^{0}_{k+1}(B)$ and $S_{k+1}(\pm B')$, and with them $H_{k+1}(B')$, $\bar{H}_{k+1}(B')$, $D_{k+1}(B')$ and $\det D_{k+1}(B')$, are analytic functions of $B$ on $\mathrm{Iv}_k$.
These, together with the operators with entire entries listed above, are all the operators and determinants that occur in (a)--(c).
By Steps 1--6, (a)--(c) hold at every point of $\mathrm{Iv}_k$.
In (a) the fields $\Psi'$, $\bar{\Psi}'$, $W'$, $\bar{W}'$ do not depend on $B$, and both sides of each identity are polynomials in their components, whose coefficients are finite sums of products of entries of the operators named, of the weights of Section~I.1.6, of the numbers $b_{k}$, $b_{L}$ and of powers of $L$;
sums and products of analytic functions are analytic.
Hence (a)--(c) are identities between analytic functions on $\mathrm{Iv}_k$.

\textbf{(d2).}
Let $\mu\in(0,1]$ and $e_k\in(0,e_0]$; then $(k,\ell,\mu,e_k)$ are data (SF) of Section~I.1.8.
Let $B\in\tilde{G}_k$.
By Proposition~2.8\,(a), $F_{y}(B)$ is bijective on $\mathfrak{S}^{\mathrm{ch}}(\eta_k,\ell)$ for every $y\in Y_{k}$.
Let $k\ge1$.
At $y=\infty$, where $\boldsymbol{\alpha}_{\infty}=b_{k}$ and $\boldsymbol{\beta}_{\infty}=0$, the operator $F_{\infty}(B)=D_B+\mu+b_{k}\cdot J\cdot P_k(B)\cdot J^{*}$ is bijective, so that $S_k(B)$ exists.
At $y=0$ the operator $F_{0}(B)$ is bijective;
this is (ii) of Proposition~21.3 at $y=0$, so that (i) holds: $D_k(B)+b_{L}\cdot P(B)$ is invertible, and $\Gamma_k(B)$ exists.
Let $k=0$.
Then $Y_{0}=\mathbb{R}$, and at $y=0$ the operator $F_{0}(B)=D_B+\mu+b_{L}\cdot J\cdot P(B)\cdot J^{*}$ is bijective, so that $\Gamma_0(B)$ exists.
In either case $B$ lies in $\mathrm{Iv}_k(\ell,\mu,e_k)$.

\textbf{(d3).}
Let $\mu>0$, and let $B$ be real.
By Proposition~21.5\,(a), $F_{\infty}(B)$ is invertible if $k\ge1$;
by Proposition~21.5\,(c), $\Gamma_k(B)$ exists for every $k\ge0$.
Hence $B$ lies in $\mathrm{Iv}_k(\ell,\mu,e_k)$.

\emph{The three sources of the inverses.}
Level $k$ at $\pm B$:
if $B$ lies in $\tilde{G}_k$, then so does $-B$, by \cite[Lemma~17.6]{QEDI} and \cite[Lemma~5.19]{QEDI}, whose assertion on $\tilde{G}_k$ holds for every $k\ge0$, every $\ell\in\mathrm{Cl}$ and every $e_k\in(0,e_0]$;
if $B$ is real, then so is $-B$;
hence (d2), respectively (d3), applies to $B$ and to $-B$.
At $k=0$: $\Gamma_0(B)=S^{0}_{1}(B)=F_{0}(B)^{-1}$ by Section~I.3.9.
Level $k+1$ at $\pm B'$: (21.29) for $(B',B)$ and for $(-B',-B)$, together with the bijectivity of $\mathrm{Rsc}$ (Step 1).

\textbf{Step 8 (the assertions of (a) for $k=0$).}
By Section~I.3.9, $H_{0}=\bar{H}_{0}=T_0=\bar{T}_0=1$; hence $\psi^{0}=\Psi^{c}$ and $\bar{\psi}^{0}=\bar{\Psi}^{c}$.
It remains to show that at $B=0$, provided $F_{0}(0)$ is invertible, the field $\Psi^{c}$ is constant on clusters only if $\Psi'=0$, and $\bar{\Psi}^{c}$ only if $\bar{\Psi}'=0$.
Let $k=0$ and $B=0$ (so that $B'=0$), and let $F_{0}(0)$ be invertible.
\begin{itemize}
\item[(i)] Let $\xi\in\mathfrak{S}(1,\ell)$ and $f=J\xi$.
Then $f(x^{(i)})=\xi(c)=f(x)$ for $i=0,1,2$, and $f(x^{(3)})=\xi(c+g e_{1})$.
Let $\mathrm{Dz}$ denote the operator (I.2.3) at the connection $0$, without mass.
Then (I.2.3) and (I.2.1) at $\eta_0=1$ and $B=0$ give
\[
\left(\mathrm{Dz}\,J\xi\right)(c,g)=-M_{3}(x)\cdot\left[\xi(c+g e_{1})-\xi(c)\right]=-3\cdot\left(\rho-\gamma\cdot(g e_{1})\right)\cdot\left[\xi(c+g e_{1})-\xi(c)\right].
\]
\item[(ii)] Suppose that $\mathrm{Dz}\,J\xi$ is constant on clusters, and set $\mathrm{wv}(c):=-\frac{1}{3}\cdot\left(\mathrm{Dz}\,J\xi\right)(c,g)$;
this is a function $\mathbb{Z}^3\to\mathbb{C}^4$ that does not depend on $g$ (the number $3$ is $w_3/2$, by (I.2.1)).
Every $a\in\left\{\pm e_{1},\pm e_{2},\pm e_{3}\right\}$ equals $g e_{1}$ for some $g\in O_{\mathrm{cub}}$.
By (i), therefore, $\left(\rho-\gamma\cdot a\right)\cdot\left(\xi(c+a)-\xi(c)\right)=\mathrm{wv}(c)$ for all $c$ and all six vectors $a$.
By (I.c1), $\left(\rho+\gamma\cdot a\right)\cdot\left(\rho-\gamma\cdot a\right)=\rho^2-1$, which is not zero.
Hence
\begin{equation}\tag{21.33}
\xi(c+a)-\xi(c)=\left(\rho^2-1\right)^{-1}\cdot\left(\rho+\gamma\cdot a\right)\mathrm{wv}(c).
\end{equation}
At the site $c+a$ and for the direction $-a$, (21.33) reads $\xi(c)-\xi(c+a)=\left(\rho^2-1\right)^{-1}\cdot\left(\rho-\gamma\cdot a\right)\mathrm{wv}(c+a)$.
Adding the two relations gives $0=\left(\rho+\gamma\cdot a\right)\mathrm{wv}(c)+\left(\rho-\gamma\cdot a\right)\mathrm{wv}(c+a)$, and multiplication by $\rho+\gamma\cdot a$ yields
\[
\mathrm{wv}(c+a)=\mathrm{Ka}_a\,\mathrm{wv}(c),
\qquad
\mathrm{Ka}_a:=\left(1-\rho^2\right)^{-1}\cdot\left(\rho+\gamma\cdot a\right)^2=\left(1-\rho^2\right)^{-1}\cdot\left(1+\rho^2+2\rho\,(\gamma\cdot a)\right).
\]
Going from $c$ to $c+e_{1}+e_{2}$ in the two possible orders, one obtains
$\mathrm{Ka}_{e_{2}}\cdot\mathrm{Ka}_{e_{1}}\,\mathrm{wv}(c)=\mathrm{wv}(c+e_{1}+e_{2})=\mathrm{Ka}_{e_{1}}\cdot\mathrm{Ka}_{e_{2}}\,\mathrm{wv}(c)$.
But
\[
\mathrm{Ka}_{e_{1}}\cdot\mathrm{Ka}_{e_{2}}-\mathrm{Ka}_{e_{2}}\cdot\mathrm{Ka}_{e_{1}}=\left(1-\rho^2\right)^{-2}\cdot4\rho^2\cdot\left(\gamma_1\gamma_2-\gamma_2\gamma_1\right)=8\rho^2\cdot\left(1-\rho^2\right)^{-2}\cdot\gamma_1\gamma_2
\]
is an invertible matrix, since $(\gamma_1\gamma_2)^2=-1$ and $\rho\ne0$.
Hence $\mathrm{wv}=0$.
By (21.33), $\xi(c+a)=\xi(c)$ for all $c$ and $a$, so that $\xi$ is constant on $\mathbb{Z}^3$;
since $\xi(c+\ell e_{\mu})=-\xi(c)$, it follows that $\xi=0$.
\item[(iii)] Let $v'\in\mathfrak{S}(L,\ell)$, and suppose that $f:=\Gamma_0(0)\cdot J\cdot Q^{T}(0)\,v'$ is constant on clusters, that is, $f=J\xi$ with $\xi=J^{*}f\in\mathfrak{S}(1,\ell)$.
Then, since $J^{*}J=1$,
\[
J\cdot Q^{T}(0)\,v'=F_{0}(0)\,J\xi=\mathrm{Dz}\,J\xi+J\left(\mu\,\xi+b_{L}\cdot P(0)\,\xi\right).
\]
Thus $\mathrm{Dz}\,J\xi$ is constant on clusters, and $\xi=0$ by (ii).
Therefore $J\cdot Q^{T}(0)\,v'=0$, and $v'=\left(Q(0)\cdot J^{*}\right)\cdot\left(J\cdot Q^{T}(0)\right)v'=0$ by Proposition~21.1\,(a).
\item[(iv)] Write $\Psi'=\sum_j v_j\,\mathrm{ob}_j$ (0.6).
The field $\Psi^{c}=\sum_j b_{L}\cdot\Gamma_0(0)\cdot J\cdot Q^{T}(0)\,(v_j)_L\,\mathrm{ob}_j$ is constant on clusters if and only if each of its coefficients is;
by (iii) with $v'=b_{L}\cdot(v_j)_L$, this gives $b_{L}\cdot(v_j)_L=0$, hence $v_j=0$ for all $j$, because $b_{L}\ne0$ and $\mathrm{Rsc}$ is injective.
By Step 5, $\bar{\Psi}^{c}=\mathsf{C}^{-1}\left[b_{L}\cdot\Gamma_0(0)\cdot J\cdot Q^{T}(0)\left(\mathsf{C}\,\bar{\Psi}'\right)_L\right]$.
As $\mathsf{C}$ is a constant invertible matrix, $\bar{\Psi}^{c}$ is constant on clusters if and only if $b_{L}\cdot\Gamma_0(0)\cdot J\cdot Q^{T}(0)\left(\mathsf{C}\,\bar{\Psi}'\right)_L$ is.
What has just been shown, applied to the field $\mathsf{C}\,\bar{\Psi}'$ on $T^0_1$ in place of $\Psi'$, gives $\mathsf{C}\,\bar{\Psi}'=0$, that is, $\bar{\Psi}'=0$.
\end{itemize}

Thus the assertions hold for every datum $(k,\ell,\mu,e_k,B')$.
The constants of Section~I.8 that occur in (a)--(c), (d1) and (d3) are among those named in (S21) ($M$ enters through $\mathrm{Cl}$ and through the pairs of $\mathrm{Pz}_k$, where $d_X\le M$);
none of them depends on the datum.
The clause (d2) contains the ceiling $e_0$, which enters through Proposition~2.8\,(a), and, through $\tilde{G}_k$ (Section~I.5.3), the exponent $\varepsilon$;
this is the exception named in (S21).
\end{proof}

\subsection*{5. Conditions on the constants}

The proofs of Propositions~21.1--21.4 use the following conditions on the constants of Section~I.8, and no others;
each of them stands in a line of Section~I.8.2 or of Proposition~14.2.

\begin{description}
\item[$L$ is an odd integer, $L>1$.]
This enters 0.2\,(a)--(c) (one centre for every site, $L^{3j}$ sites in a block, nesting), hence Proposition~21.1\,(a), (b), (f) and all that follows;
the inequality $L>1$ enters $b_{k}$ (Section~I.3.4).
The condition stands in (C3), and in (L5) of Proposition~14.2\,(b).

\item[$b>0$ ($b=1$), hence $b_{k}>0$, $b_{L}>0$, $b_{k}+b_{L}>0$.]
This enters Proposition~21.3, where $b_{k}+i\gamma_5 y$ and $b_{k}+b_{L}+i\gamma_5 y$ are invertible for every real $y$, $y=0$ included;
the number $b_{k}^{-1}$, which occurs in $T_k$ (Proposition~21.1\,(a), (c));
and Steps 2, 5 and 6 of the proof of Proposition~21.4, together with Step 8\,(iv) there ($b_{L}\ne0$).
The condition stands in (C0) ($b$) and in (C4) ($b_{L}$, $b_{-}$, $b_{+}$).

\item[The standard weights: the closure relation (I.1.1).]
This enters Proposition~21.1\,(c), namely (21.23), the diagonal term of $D_B^{T}$;
and Proposition~21.1\,(e), namely the step from (I.2.3) to (I.2.4), hence $D_B^{*}=\gamma_5\cdot D_B\cdot\gamma_5$ and the formula for $\operatorname{Re}D_B$.
The condition stands in (C0) ($w$).

\item[$w_i>0$.]
This enters Proposition~21.1\,(e), namely $\nabla_B^{*}\nabla_B\ge0$;
and Step 8\,(i), (ii) of the proof of Proposition~21.4, through $w_3\ne0$ alone (division by $w_3/2=3$).
The condition stands in (C0) ($w$).

\item[$\rho\ge0$.]
This enters Proposition~21.1\,(e), namely $\operatorname{Re}D_B\ge0$.
The condition stands in (C0) ($\rho\in(0,1)$).

\item[$\rho\ne0$, $\rho^2\ne1$.]
Among the arguments of this chapter, these enter only Step 8 of the proof of Proposition~21.4, where the assertions of (a) for $k=0$ are proved ($H_{0}=\bar{H}_{0}=T_0=\bar{T}_0=1$; at $B=0$, with $F_{0}(0)$ invertible, $\Psi^{c}$ is constant on clusters only if $\Psi'=0$, and $\bar{\Psi}^{c}$ only if $\bar{\Psi}'=0$).
In (d2), moreover, the range $0<\rho<1$ enters through Proposition~2.8\,(a); see the entry on the conditions under which Proposition~2.8\,(a) holds.
The conditions stand in (C0) ($\rho\in(0,1)$).

\item[$M$ is a power of $L$, $M\ge L^2$.]
This enters 0.1, where $M$ is an integer multiple of $L$, so that, for $\ell\in\mathrm{Cl}$, $\ell/L$ is an integer multiple of $L$ (Section~0.7, level $k+1$);
the inequality $L/2<M$ ($U_B(y',y)$ is defined for $y\in B(y')$);
and the inclusion $M\mathbb{Z}^3\subset\mathbb{Z}^3$ (Proposition~21.1\,(d0), lattice symmetries).
The inequality $2M<\ell/L$, by contrast, follows from $\ell\in\mathrm{Cl}$ alone.
The condition stands in (C5), and in (M1) of Proposition~14.2\,(b).

\item[$\alpha$ is a real number.]
This enters as the exponent of $d_X$ in Proposition~21.1\,(d), (f) and in (21.11) of Proposition~21.4;
there $\alpha$ may be any real number (in (d2) the range $0\le\alpha<1$ enters through Proposition~2.8\,(a); see the next entry).
The condition stands in (C1).

\item[The conditions under which Proposition~2.8\,(a) holds.]
They are of two kinds.
The first kind consists of (L5), (M1), (M2), $M\ge108\pi\cdot C'_{V}\cdot a'$ (contained in (M5)), (Z2.s1), (Z2.s2), and $0\le\alpha<1$.
The first six of these are the conditions listed for the Propositions~2.m before Proposition~2.1, and the constants in Proposition~2.8 are subject to them and to $\alpha<1$
(the conditions under which Proposition~2.7 is stated include $\alpha<1$, and $a'$ contains $1/(1-\alpha)$).
The lower bound $0\le\alpha$ is used in Proposition~1.5\,(f) and in Lemma~6.5 of Chapter~1b, and it enters $a'$ through $C_{5}$.
Among the proofs of Propositions~21.1--21.4, only that of Proposition~21.4\,(d2) applies Proposition~2.8\,(a), and the conditions of the first kind enter it only through Proposition~2.8\,(a).
They stand in (C1), (C3), (C5), (C8), and in the entry on Proposition~21.4\,(d2) in Annex~B, Section~G3.6.1.
The second kind consists of further conditions that are used in the proofs on which Proposition~2.8\,(a) rests (see the sections on the conditions on the constants in Chapters~1a, 1b, 2a and~2b);
through Proposition~2.8\,(a) they enter (d2) as well.
From the line (C0) these are the range $0<\rho<1$, the standard weights (their values, $w_i>0$, and the closure relation (I.1.1)), and $b>0$.
Here $\rho>0$ gives $c_{\mathrm{re}}>0$ (Proposition~1.7), so that $C_{00}=48\cdot(2L/c_{\mathrm{re}})\cdot(1+2k_{1})$ and $a'$ are finite;
since $c_{\mathrm{re}}\le\rho/(2C_{\mathrm{Po}}\cdot L)$, one has $C_{00}\ge192\,C_{\mathrm{Po}}\cdot L^{2}\cdot(1+2k_{1})/\rho$, so that $a'$ contains $1/\rho$ through $C_{00}$, and $a'$ tends to infinity as $\rho\to0$ (Proposition~14.4\,(d)).
The inequality $\rho<1$ gives $|M_{3}(x)|_2=3(1+\rho)<6$ and $(1+\rho)\cdot\mathrm{sw}<58$, which enter the numerals of $C'_{V}=2\mathrm{e}\cdot(42+6b)$ and of $k_{1}=60+b$ (Propositions~2.1\,(a), (b) and~2.3\,(b)).
From the line (C7) it is $\varepsilon>0$, which gives $e_k^{\varepsilon}\le e_0^{\varepsilon}$ for $e_k\le e_0$.
Hence the conditions under which Proposition~2.8\,(a) holds stand in the lines (C0), (C1), (C3), (C5), (C7), (C8).

\item[The conditions that enter Proposition~21.5\,(a), (c).]
These are: $L$ odd, $L>1$, $M$ a power of $L$ with $L/2<M$ (in the definitions), $b>0$, $\rho\ge0$, and, through Proposition~21.1\,(e), the closure relation (I.1.1) and $w_i>0$.
Among the proofs of Propositions~21.1--21.4, only that of Proposition~21.4\,(d3) applies Proposition~21.5\,(a), (c), and these conditions enter it only through Proposition~21.5\,(a), (c).
They stand in (C0), (C3), (C4), (C5), and in the entry on Propositions~21.1--21.3, 21.4\,(a)--(c), (d1), (d3), 21.5 in Annex~B, Section~G3.6.1.
\end{description}

\begin{named}{Remark}
The side $\ell\in\mathrm{Cl}$ is a datum.
With one exception, the proofs use of it only the following properties:
$\ell$ is an integer multiple of $L$ (hence $L\le\ell$, which enters the description of the kernels in Proposition~21.1\,(a));
$2M<\ell$ and $2M<\ell/L$ where pairs occur.
The exception is Proposition~21.4\,(d2), together with the passage on the three sources of the inverses in Step~7 of the proof of Proposition~21.4, wherever $\tilde{G}_k$ occurs.
There the full condition $\ell\in\mathrm{Cl}$ enters in three ways:
\begin{itemize}
\item[(i)] as part of the data (SF) of Section~I.1.8, under which Proposition~2.8\,(a) is stated;
\item[(ii)] as a hypothesis of \cite[Lemma~5.19]{QEDI}, which gives the invariance of $\tilde{G}_k$ under $B\mapsto-B$;
\item[(iii)] through the definition of $\tilde{G}_k$ by the $M$-cubes of $T_k$ (Section~I.5.3), which requires that $\ell$ be a multiple of $M$.
\end{itemize}
Beyond this, the proof of (d2) uses nothing of $\ell$.
Likewise, the lattice symmetries of Proposition~21.1\,(d0) are defined with translation parts in $M\mathbb{Z}^3/\ell\mathbb{Z}^3$, which presupposes that $M$ divides $\ell$;
the proof of (d0) uses of a translation part $a$ only that $a\in\mathbb{Z}^3$.
That $\ell/L$ is a multiple of $L$ serves, in Section~0.7, only to extend Sections~0.2--0.5 to the level $k+1$;
for the operators of level $k+1$ in the list of (S21) it suffices that $\ell/L$ is an integer.
\end{named}

The mass $\mu$ is an arbitrary complex number and $e_k$ an arbitrary positive number, except in Proposition~21.4\,(d2), where $\mu\in(0,1]$ and $e_k\in(0,e_0]$, and in Proposition~21.4\,(d3), where $\mu>0$.

\section*{Chapter 21b. Real parts, $\gamma_5$-symmetry and invertibility of the operators of a level at a real background: proofs of Proposition~21.5}

\subsection*{0. Conventions of this chapter}

The following notation is used in this chapter.
The letters $\mathrm{Vs}$, $\mathrm{Vs}'$, $\mathrm{Vs}''$ stand for one of the spaces of 3.1 together with its scalar product.
Orthogonal projections are denoted by $\mathrm{Pj}$, $\mathrm{Pj}_1$, $\mathrm{Pj}_2$ (see (21.5.4) and (21.5.22)).
The operators $\mathrm{Kf}$ and $\mathrm{Kh}$ are introduced in Section~7, and $\mathrm{ew}_j$ denote the eigenvalues of $\mathrm{Kh}$.
The real numbers $\mathrm{pa}_y$, $\mathrm{qa}_y$, $\mathrm{pb}_y$, $\mathrm{qb}_y$ are those of (21.5.22), and $\mathrm{sf}$ denotes a periodic function on $\mathbb{L}^{\mathrm{f}}_k$ (see 8a.1 and 8a.2).
The letters $\mathrm{cf}$, $\mathrm{cf}_1$, $\mathrm{cf}_2$ denote numbers;
$\mathrm{Mx}$ is a constant $4\times4$ matrix, as in Proposition~21.1\,(a);
$O$, $O_1$, $O_2$ are linear maps;
$f$, $f'$, $f_1$, $f_2$ are elements of the spaces of 3.1;
and $v$, $v'$ are spinors on the unit lattice or on the $L$-lattice.
The clauses of Proposition~21.5 are denoted by (a), (b), (c), (d).
A number such as 3.1 without further indication refers to a subsection of this chapter.

\subsection*{1. The statement}

This chapter contains the proof of the following statement.

\begin{named}{Proposition 21.5}[Real backgrounds]
In the setting (S21) let $B$ be real and $\mu>0$.
All operators are taken at $B$, and $\operatorname{Re}O:=(O+O^*)/2$, the adjoint being formed with respect to the sesquilinear pairings of Section~I.1.6.
\begin{itemize}
\item[(a)] $\operatorname{Re}(D_B+\mu)\ge\mu$, and for $k\ge1$ also $\operatorname{Re}F_\infty\ge\mu$.
Hence $D_B+\mu$ and $F_\infty$ are invertible.
\item[(b)] $D_k$ is invertible and $\operatorname{Re}D_k>0$, so that $\det D_k\ne0$; for $k\ge1$,
\[
D_k^{-1}=b_k^{-1}+Q_kJ^*\cdot(D_B+\mu)^{-1}\cdot JQ_k^T(-B).
\]
\item[(c)] For $k\ge1$: $\operatorname{Re}(D_k+b_LP)>0$ and $(D_k+b_LP)^*=\gamma_5\,(D_k+b_LP)\,\gamma_5$.
For $k=0$: $\operatorname{Re}F_0\ge\mu$ and $F_0^*=\gamma_5F_0\gamma_5$.
Hence $\Gamma_k$ exists and $\delta Z_k\ne0$.
\item[(d)] For every real $y$ the operator $D_k+b_LP+i\gamma_5y$ ($k\ge1$), resp.\ $F_y=F_0+i\gamma_5y$ ($k=0$), is invertible;
for $y\ne0$ its inverse has operator norm at most $1/|y|$ with respect to the $L^2$ norms.
For $k\ge1$, $S_{k,y}$ exists.
\end{itemize}
\end{named}

\subsection*{2. The setting, the definitions and the recalled clauses}

\subsubsection*{2.1. The setting}

The setting is (S21) of Section~G3.1.0 with its data of level~$k$:
an integer $k\ge0$, a side $\ell\in\mathrm{Cl}$, a complex number $\mu$, a number $e_k>0$ and a complex bond function $B$ on $T_k$.
In Proposition~21.5 the number $\mu$ is real and positive, and $B$ is real.

\subsubsection*{2.2. The operators, all at $B$}

As in Proposition~21.5, operators written without argument are taken at $B$:
$Q_k$, $P_k$, $P_{k+1}$, $Q$, $P$, $F_y$, $S_k$, $D_k$, \dots\ stand for $Q_k(B)$, $P_k(B)$, $P_{k+1}(B)$, $Q(B)$, $P(B)$, $F_y(B)$, $S_k(B)$, $D_k(B)$, \dots,
whereas $Q_k^T(-B)$ and $Q^T(-B)$ are always written with their argument.
Since the letter $y$ is reserved for the resolvent parameter, a unit site is denoted by $z$ in the displays of this chapter.

\emph{Numbers.}
One has $b=1$, kept as a letter, and $L$ is an odd integer with $L\ge5$ (Section~I.0.5).
Consequently $b_L=b/L>0$, and for $k\ge1$
\[
b_k=\frac{b\,(1-L^{-1})}{1-L^{-k}}>0
\]
(Section~I.3.4), because $0<L^{-k}\le L^{-1}<1$.

\emph{The operator $D_B$, at every level $k\ge0$.}
The fermion sites $x=(c,g)$ and the objects $x^{(i)}$, $c_3$, $\hat{d}_i(x)$ are defined in Section~I.1.3 and in (I.1.2);
the phases $u_B$ and $U_B$ in (I.1.6);
the matrices $M_i(x)$ in (I.2.1);
and the operators $R_i(B)$, $D_B$, $\nabla_B^{*}\nabla_B$, $J_B$ on $\mathfrak{S}^{\mathrm{ch}}(\eta_k,\ell)$ in (I.2.2), (I.2.3), (I.2.4).

\emph{Levels $k\ge1$.}
The cluster maps $J$, $J^*$ are defined in (I.3.1), the phases $\varphi_B$ in (I.3.3), and the operators $Q_k$, $Q_k^T(-B)$, $P_k$, $Q$, $Q^T(-B)$, $P$, $P_{k+1}$ in Section~I.3.3.
The matrices $\boldsymbol{\alpha}_y$, $\boldsymbol{\beta}_y$, the operators $F_y$ ($y$ real or $y=\infty$), $S_k$, $S_{k,y}$, $D_k$, $\Gamma_k$, $\Gamma_{k,y}$, the determinant $\delta Z_k$, and the spaces between which all these operators act, are those of Section~I.3.5 and its table.
The space $\mathfrak{S}^{\mathrm{ch}}(\eta_k,\ell)$ has dimension $192\,(\ell/\eta_k)^3$, and $\mathfrak{S}(1,\ell)$ has dimension $4n_k$ (Sections~I.1.5 and I.1.1).
Since $\boldsymbol{\alpha}_\infty=b_k$ and $\boldsymbol{\beta}_\infty=0$ (Section~I.3.5),
\[
F_\infty=D_B+\mu+b_k\cdot JP_kJ^*.
\]
For real $y$ the inverses that enter the $4\times4$ matrices $\boldsymbol{\alpha}_y$, $\boldsymbol{\beta}_y$ exist: by $\gamma_5^2=1$,
\[
(\mathrm{cf}+i\gamma_5y)(\mathrm{cf}-i\gamma_5y)=\mathrm{cf}^2+y^2>0
\qquad\text{for }\mathrm{cf}=b_k\text{ and for }\mathrm{cf}=b_k+b_L.
\]
For $f\in\mathfrak{S}^{\mathrm{ch}}(\eta_k,\ell)$ and $v\in\mathfrak{S}(1,\ell)$ the operators between the two spaces act by (Section~I.3.3, (I.1.3), (I.3.1))
\begin{align*}
(Q_kJ^*f)(z)&=\int_{\Delta_z}\mathrm{d}x\,\varphi_B(z,c)\,f(x)
=\eta_k^3\sum_{c\in\Delta_z}\varphi_B(z,c)\,\frac{1}{48}\sum_{g\in O_{\mathrm{cub}}}f(c,g),\\
\left(JQ_k^T(-B)v\right)(c,g)&=\varphi_B(z,c)^{-1}\,v(z)
\qquad\text{for }c\in\Delta_z\text{ and all }48\text{ labels }g\in O_{\mathrm{cub}}.
\end{align*}

\emph{Level $0$.}
The operators $D_0$, $F_y$ ($y$ real), $S_{0,y}$, $\Gamma_0$ and the determinant $\delta Z_0$ are defined in Section~I.3.9.
There $\eta_0=1$, and all operators act on $\mathfrak{S}^{\mathrm{ch}}(1,\ell)$, a space of dimension $192\,n_0$.
The operator $D_0=D_B+\mu$ is $D_k$ at $k=0$, taken at the field $B$; it is not the operator (I.2.3) at zero field.
The symbols $b_0$, $F_\infty$, $S_k$, $Q_kJ^*$ do not appear at $k=0$ (Section~I.3.9: $Y_0=\mathbb{R}$, $Q_0=1$).

\subsubsection*{2.3. The recalled clauses}

For the real bond function $B$ of 2.1 the following clauses are recalled.
\begin{description}
\item[Proposition~21.1\,(a).]
For $k\ge1$, $\left(Q_kJ^*\right)\left(JQ_k^T(-B)\right)=1$ on $\mathfrak{S}(1,\ell)$.
The operators $JP_kJ^*$, $JP_{k+1}J^*$, $P$ ($k\ge1$) and $JPJ^*$ ($k=0$) are idempotent, their kernels are multiples of $1_4$, and they commute with every constant $4\times4$ matrix;
for $k\ge1$,
\[
\left(JP_kJ^*\right)\left(JP_{k+1}J^*\right)=\left(JP_{k+1}J^*\right)\left(JP_kJ^*\right)=JP_{k+1}J^*.
\]
\item[Proposition~21.1\,(e).]
One has $(Q_kJ^*)^*=JQ_k^T(-B)$ for $k\ge1$, $Q^*=Q^T(-B)$, and $(QJ^*)^*=JQ^T(-B)$ for $k=0$;
the operators $JP_kJ^*$, $JP_{k+1}J^*$ ($k\ge1$), $P$, and $JPJ^*$ ($k=0$) are orthogonal projections;
$D_B^*=\gamma_5D_B\gamma_5$;
and $\operatorname{Re}D_B=\frac{\rho}{2\eta_k}\cdot\nabla_B^{*}\nabla_B\ge0$.
The last two facts are part of (I.O2).
\item[Proposition~21.3, implication (i)$\Rightarrow$(ii).]
Let $k\ge1$, let $y$ be real, and let $F_\infty$ be invertible.
If $D_k+b_LP+i\gamma_5y$ is invertible on $\mathfrak{S}(1,\ell)$, then $F_y$ is invertible on $\mathfrak{S}^{\mathrm{ch}}(\eta_k,\ell)$.
\item[(I.c1).]
In particular, the matrices $\gamma_\mu$ ($\mu=1,2,3$) and $\gamma_5$ are Hermitian, $\gamma_\mu\gamma_\nu+\gamma_\nu\gamma_\mu=2\delta_{\mu\nu}$, $\gamma_5\gamma_\mu=-\gamma_\mu\gamma_5$, $\gamma_5^2=1$, and $(\gamma\cdot v)^2=|v|_2^2$ for real $v$.
\end{description}
In Section~8a the clauses of Proposition~21.1 recalled above are derived directly from the definitions of Part~I, with the sums over the clusters written out.

\subsection*{3. Preliminaries}

\subsubsection*{3.1. Scalar products and norms}

Let $c$ range over a fundamental domain of $\ell\mathbb{Z}^3$ in $\mathbb{L}_k$, and $z$ over a fundamental domain of $\ell\mathbb{Z}^3$ in $\mathbb{Z}^3$.
The scalar products are
\begin{align*}
(f',f)^{\mathrm{ch}}&=\int\mathrm{d}x\,\overline{f'(x)}\cdot f(x)
=\eta_k^3\,\frac{1}{48}\sum_c\sum_{g\in O_{\mathrm{cub}}}\sum_{\beta=1}^{4}\overline{f'_\beta(c,g)}\,f_\beta(c,g)
&&\text{on }\mathfrak{S}^{\mathrm{ch}}(\eta_k,\ell),\\
(v',v)_1&=\left\langle\overline{v'},v\right\rangle_1
=\sum_z\sum_{\beta=1}^{4}\overline{v'_\beta(z)}\,v_\beta(z)
&&\text{on }\mathfrak{S}(1,\ell);
\end{align*}
the first by (I.1.3) and Section~I.1.6, the factor $\eta_0^3$ being $1$ at $k=0$;
the second by Section~I.1.6 with $a=1$, every unit site carrying the weight $1$.
The summands are $\ell\mathbb{Z}^3$-periodic, as products of two antiperiodic factors, so that the choice of the fundamental domains is immaterial.
Both forms are linear in the second argument and conjugate linear in the first, and $(f',f)=\overline{(f,f')}$.
They are positive definite.
Indeed, $(f,f)^{\mathrm{ch}}=\eta_k^3\,\frac{1}{48}\sum_c\sum_g|f(c,g)|_2^2$ vanishes only if $f$ vanishes on the product of the fundamental domain with $O_{\mathrm{cub}}$, and then $f=0$ by antiperiodicity;
likewise $(v,v)_1=\sum_z|v(z)|_2^2$ vanishes only if $v$ vanishes on the fundamental domain, and then $v=0$ by antiperiodicity.
The $L^2$ norms, formed with $|\cdot|_2$ (Section~I.1.5), are
\[
\|f\|:=\left((f,f)^{\mathrm{ch}}\right)^{\frac{1}{2}},\qquad
\|v\|:=\left((v,v)_1\right)^{\frac{1}{2}},
\]
and $|O|_{\mathrm{op}}:=\sup\{\|Of\|:\|f\|=1\}$ is the associated operator norm.
The Cauchy--Schwarz inequality for a positive definite form \cite[Section~5.1]{HornJohnson2013} gives
\[
|(f',f)|\le\|f'\|\cdot\|f\|.
\]
In 3.2 and 3.3 the letters $\mathrm{Vs}$, $\mathrm{Vs}'$, $\mathrm{Vs}''$ denote either of the two spaces $\mathfrak{S}^{\mathrm{ch}}(\eta_k,\ell)$, $\mathfrak{S}(1,\ell)$, each with its form, which is written $(\cdot\,,\cdot)$.
Both spaces are finite dimensional.
(In 8a.7 the space $\mathfrak{S}(L,\ell)$ with its form $(\cdot\,,\cdot)_L$ is a third space of this kind.)

\subsubsection*{3.2. Adjoints and real parts}

Let $O:\mathrm{Vs}'\to\mathrm{Vs}''$ be linear.
Its adjoint $O^*:\mathrm{Vs}''\to\mathrm{Vs}'$ is the linear map with
\[
(f_2,Of_1)=(O^*f_2,f_1)\qquad\text{for all }f_1\in\mathrm{Vs}',\ f_2\in\mathrm{Vs}'';
\]
it exists and is unique, because the forms are positive definite and the dimensions are finite.
The definition yields the following rules:
$(O^*)^*=O$;
$(O_1O_2)^*=O_2^*O_1^*$;
$(\mathrm{cf}_1\cdot O_1+\mathrm{cf}_2\cdot O_2)^*=\overline{\mathrm{cf}_1}\cdot O_1^*+\overline{\mathrm{cf}_2}\cdot O_2^*$ for complex numbers $\mathrm{cf}_1$, $\mathrm{cf}_2$;
$1^*=1$;
and if $O$ is bijective, then so is $O^*$, and $(O^{-1})^*=(O^*)^{-1}$, as one sees by taking adjoints in $OO^{-1}=1$ and $O^{-1}O=1$.

Now let $O$ map $\mathrm{Vs}$ to itself.
Then $\operatorname{Re}O=(O+O^*)/2$ is self-adjoint.
For a real number $\mathrm{cf}$ the inequality $\operatorname{Re}O\ge\mathrm{cf}$ means that $(f,(\operatorname{Re}O)f)\ge\mathrm{cf}\cdot(f,f)$ for all $f\in\mathrm{Vs}$,
and $\operatorname{Re}O>0$ means that $(f,(\operatorname{Re}O)f)>0$ for all $f\ne0$.
Since $(f,O^*f)=(Of,f)=\overline{(f,Of)}$, one has
\begin{equation}
(f,(\operatorname{Re}O)f)=\operatorname{Re}(f,Of)\qquad\text{for all }f\in\mathrm{Vs}.
\tag{21.5.1}
\end{equation}
Moreover $\operatorname{Re}(O_1+O_2)=\operatorname{Re}O_1+\operatorname{Re}O_2$, and for a real number $\mathrm{cf}$ one has $\operatorname{Re}(\mathrm{cf}\cdot O)=\mathrm{cf}\cdot\operatorname{Re}O$ and $\operatorname{Re}(\mathrm{cf}\cdot1)=\mathrm{cf}$.

\subsubsection*{3.3. Three facts}

In (21.5.2)--(21.5.4) the map $O$ is a linear map of $\mathrm{Vs}$ to itself.

\begin{named}{(21.5.2)}
If $\operatorname{Re}O\ge\mathrm{cf}$ for a number $\mathrm{cf}>0$, then $\|Of\|\ge\mathrm{cf}\cdot\|f\|$ for all $f$;
moreover $O$ is bijective, $\det O\ne0$, and $|O^{-1}|_{\mathrm{op}}\le1/\mathrm{cf}$.
If $\operatorname{Re}O>0$, then $O$ is bijective and $\det O\ne0$.
\end{named}

\begin{proof}
By (21.5.1) and the Cauchy--Schwarz inequality,
\[
\mathrm{cf}\cdot\|f\|^2\le\operatorname{Re}(f,Of)\le|(f,Of)|\le\|f\|\cdot\|Of\|.
\]
For $f\ne0$ division by $\|f\|$ gives $\|Of\|\ge\mathrm{cf}\cdot\|f\|$, and for $f=0$ this inequality is trivial.
Hence $Of=0$ implies $f=0$.
An injective linear map of a finite dimensional space into itself is bijective, and its determinant does not vanish \cite[Section~0.5]{HornJohnson2013}.
Given $f'\in\mathrm{Vs}$, put $f:=O^{-1}f'$; then
\[
\|O^{-1}f'\|=\|f\|\le\frac{\|Of\|}{\mathrm{cf}}=\frac{\|f'\|}{\mathrm{cf}}.
\]
For the second sentence, let $Of=0$.
Then $(f,(\operatorname{Re}O)f)=\operatorname{Re}(f,Of)=0$, and $\operatorname{Re}O>0$ forces $f=0$.
So $O$ is injective, hence bijective, and $\det O\ne0$ \cite[Section~0.5]{HornJohnson2013}.
\end{proof}

\begin{named}{(21.5.3)}
If $O$ is bijective, $f'\in\mathrm{Vs}$ and $f:=O^{-1}f'$, then $\operatorname{Re}(f',O^{-1}f')=\operatorname{Re}(f,Of)$.
Consequently $\operatorname{Re}O\ge0$ implies $\operatorname{Re}O^{-1}\ge0$.
\end{named}

\begin{proof}
One has $(f',O^{-1}f')=(Of,f)=\overline{(f,Of)}$, and complex conjugation does not change the real part.
If $\operatorname{Re}O\ge0$, then by (21.5.1), applied to $O^{-1}$ and to $O$,
\[
\left(f',(\operatorname{Re}O^{-1})f'\right)=\operatorname{Re}(f',O^{-1}f')=\operatorname{Re}(f,Of)=\left(f,(\operatorname{Re}O)f\right)\ge0
\]
for every $f'\in\mathrm{Vs}$.
\end{proof}

\begin{named}{(21.5.4)}
If $\mathrm{Pj}$ is an orthogonal projection of $\mathrm{Vs}$, that is, $\mathrm{Pj}^*=\mathrm{Pj}=\mathrm{Pj}^2$, then $\operatorname{Re}\mathrm{Pj}=\mathrm{Pj}$ and $(f,\mathrm{Pj}\,f)=\|\mathrm{Pj}\,f\|^2\ge0$ for all $f$.
\end{named}

\begin{proof}
One has $\operatorname{Re}\mathrm{Pj}=\mathrm{Pj}$ because $\mathrm{Pj}^*=\mathrm{Pj}$, and
\[
(f,\mathrm{Pj}\,f)=(f,\mathrm{Pj}^*\cdot\mathrm{Pj}\,f)=(\mathrm{Pj}\,f,\mathrm{Pj}\,f).
\]
\end{proof}

\subsubsection*{3.4. The matrix $\gamma_5$ as an operator}

A constant $4\times4$ matrix $\mathrm{Mx}$ acts on $\mathfrak{S}^{\mathrm{ch}}(\eta_k,\ell)$ by $(\mathrm{Mx}\,f)(c,g):=\mathrm{Mx}\,f(c,g)$ and on $\mathfrak{S}(1,\ell)$ by $(\mathrm{Mx}\,v)(z):=\mathrm{Mx}\,v(z)$;
both actions preserve antiperiodicity.
Let $\mathrm{Mx}^*$ be the Hermitian conjugate matrix, $(\mathrm{Mx}^*)_{\beta'\beta}:=\overline{\mathrm{Mx}_{\beta\beta'}}$.
At every site $x=(c,g)$,
\[
\overline{f'(x)}\cdot\left(\mathrm{Mx}\,f(x)\right)
=\sum_{\beta,\beta'}\overline{f'_\beta(x)}\,\mathrm{Mx}_{\beta\beta'}\,f_{\beta'}(x)
=\sum_{\beta'}\overline{\left(\sum_{\beta}(\mathrm{Mx}^*)_{\beta'\beta}\,f'_\beta(x)\right)}\,f_{\beta'}(x)
=\overline{\mathrm{Mx}^*f'(x)}\cdot f(x).
\]
This identity holds for arbitrary vectors of $\mathbb{C}^4$;
for the vectors $v'(z)$, $v(z)$ at a unit site $z$ it reads
\[
\overline{v'(z)}\cdot\left(\mathrm{Mx}\,v(z)\right)=\overline{\mathrm{Mx}^*v'(z)}\cdot v(z).
\]
Inserting the two identities term by term into the two displays of 3.1, one finds that on both spaces the adjoint of the operator $\mathrm{Mx}$ is the operator $\mathrm{Mx}^*$.
By (I.c1) the matrix $\gamma_5$ is Hermitian, $\gamma_5^*=\gamma_5$, and $\gamma_5^2=1$.
Therefore, on both spaces,
\begin{equation}
\gamma_5^*=\gamma_5,\qquad\gamma_5^2=1,\qquad\|\gamma_5f\|^2=(f,\gamma_5^*\gamma_5f)=\|f\|^2.
\tag{21.5.5}
\end{equation}

\subsubsection*{3.5. The averaging operators and constant matrices ($k\ge1$)}

In the two displays for $Q_kJ^*f$ and $JQ_k^T(-B)v$ in 2.2 the numbers $\eta_k^3$, $1/48$, $\varphi_B(z,c)$, $\varphi_B(z,c)^{-1}$ are scalars.
Hence, for every constant $4\times4$ matrix $\mathrm{Mx}$, every $f\in\mathfrak{S}^{\mathrm{ch}}(\eta_k,\ell)$ and every $v\in\mathfrak{S}(1,\ell)$,
\begin{align*}
\left(Q_kJ^*(\mathrm{Mx}\,f)\right)(z)&=\eta_k^3\sum_{c\in\Delta_z}\varphi_B(z,c)\,\frac{1}{48}\sum_{g\in O_{\mathrm{cub}}}\mathrm{Mx}\,f(c,g)=\mathrm{Mx}\,(Q_kJ^*f)(z),\\
\left(JQ_k^T(-B)(\mathrm{Mx}\,v)\right)(c,g)&=\varphi_B(z,c)^{-1}\,\mathrm{Mx}\,v(z)=\mathrm{Mx}\left(JQ_k^T(-B)v\right)(c,g),
\end{align*}
that is,
\begin{equation}
Q_kJ^*\cdot\mathrm{Mx}=\mathrm{Mx}\cdot Q_kJ^*,\qquad
JQ_k^T(-B)\cdot\mathrm{Mx}=\mathrm{Mx}\cdot JQ_k^T(-B),
\tag{21.5.6}
\end{equation}
where in each of the two equations $\mathrm{Mx}$ acts on $\mathfrak{S}^{\mathrm{ch}}(\eta_k,\ell)$ on one side and on $\mathfrak{S}(1,\ell)$ on the other.

\subsection*{4. Proof of (a)}

\begin{proof}[Proof of Proposition~21.5\,(a)]
Let $k\ge0$.
The operator $D_B$ maps $\mathfrak{S}^{\mathrm{ch}}(\eta_k,\ell)$ into itself (Section~I.3.5 with its table, and Section~I.3.9);
at $k=0$ this space is $\mathfrak{S}^{\mathrm{ch}}(1,\ell)$ and $\eta_0=1$.
By Proposition~21.1\,(e), $\operatorname{Re}D_B=\frac{\rho}{2\eta_k}\cdot\nabla_B^{*}\nabla_B\ge0$.
As $\mu$ is real, $\operatorname{Re}(\mu\cdot1)=\mu$ by 3.2.
It follows that $\operatorname{Re}(D_B+\mu)=\frac{\rho}{2\eta_k}\cdot\nabla_B^{*}\nabla_B+\mu$;
by (21.5.1) this says that for all $f\in\mathfrak{S}^{\mathrm{ch}}(\eta_k,\ell)$
\begin{equation}
\operatorname{Re}\left(f,(D_B+\mu)f\right)^{\mathrm{ch}}
=\frac{\rho}{2\eta_k}\cdot\left(f,\nabla_B^{*}\nabla_Bf\right)^{\mathrm{ch}}+\mu\cdot(f,f)^{\mathrm{ch}}
\ge\mu\cdot(f,f)^{\mathrm{ch}}.
\tag{21.5.7}
\end{equation}
This is the inequality $\operatorname{Re}(D_B+\mu)\ge\mu$.

Now let $k\ge1$.
By Proposition~21.1\,(e), $JP_kJ^*$ is an orthogonal projection of $\mathfrak{S}^{\mathrm{ch}}(\eta_k,\ell)$, so that by (21.5.4)
\[
\operatorname{Re}\left(JP_kJ^*\right)=JP_kJ^*,\qquad(f,JP_kJ^*f)^{\mathrm{ch}}=\|JP_kJ^*f\|^2.
\]
The number $b_k$ is real and positive (see 2.2).
From $F_\infty=D_B+\mu+b_k\cdot JP_kJ^*$ and (21.5.7) one obtains, for all $f$,
\begin{equation}
\operatorname{Re}(f,F_\infty f)^{\mathrm{ch}}
=\operatorname{Re}\left(f,(D_B+\mu)f\right)^{\mathrm{ch}}+b_k\cdot\|JP_kJ^*f\|^2
\ge\mu\cdot(f,f)^{\mathrm{ch}},
\tag{21.5.8}
\end{equation}
that is, $\operatorname{Re}F_\infty\ge\mu$.

Since $\mu>0$, (21.5.2) with $\mathrm{cf}=\mu$ applies to $D_B+\mu$ (for every $k\ge0$) and to $F_\infty$ (for $k\ge1$), and shows that both operators are bijective.
In particular, for $k\ge1$ the operator $S_k=F_\infty^{-1}$ exists, and therefore $D_k$ is defined.
\end{proof}

\subsection*{5. Proof of (b)}

\begin{proof}[Proof of Proposition~21.5\,(b)]
\emph{The case $k=0$.}
Here $D_0=D_B+\mu$ on $\mathfrak{S}^{\mathrm{ch}}(1,\ell)$.
By (21.5.7), $\operatorname{Re}D_0\ge\mu>0$; in particular $\operatorname{Re}D_0>0$.
By (21.5.2), $D_0$ is bijective and $\det D_0\ne0$, a determinant of size $192\,n_0$.

\emph{The case $k\ge1$.}
The inverses $(D_B+\mu)^{-1}$ and $S_k$ exist by (a).

(i) \emph{The inverse.}
Since $P_k=Q_k^T(-B)\,Q_k$ by definition (Section~I.3.3),
\[
\left(JQ_k^T(-B)\right)\left(Q_kJ^*\right)=JP_kJ^*.
\]
Expanding the product below, using this relation in the fourth term, and then $F_\infty=D_B+\mu+b_k\cdot JP_kJ^*$ and $F_\infty S_k=1$, one obtains on $\mathfrak{S}(1,\ell)$
\begin{equation}
\begin{split}
&\left(b_k^{-1}+Q_kJ^*\cdot(D_B+\mu)^{-1}\cdot JQ_k^T(-B)\right)\left(b_k-b_k^2\,Q_kJ^*\cdot S_k\cdot JQ_k^T(-B)\right)\\
&\quad=1-b_k\,Q_kJ^*\cdot S_k\cdot JQ_k^T(-B)+b_k\,Q_kJ^*\cdot(D_B+\mu)^{-1}\cdot JQ_k^T(-B)\\
&\qquad\quad-b_k^2\,Q_kJ^*\cdot(D_B+\mu)^{-1}\cdot JP_kJ^*\cdot S_k\cdot JQ_k^T(-B)\\
&\quad=1+b_k\,Q_kJ^*\cdot(D_B+\mu)^{-1}\cdot\left[-(D_B+\mu)S_k+1-b_k\cdot JP_kJ^*\cdot S_k\right]\cdot JQ_k^T(-B)\\
&\quad=1+b_k\,Q_kJ^*\cdot(D_B+\mu)^{-1}\cdot\left[1-F_\infty S_k\right]\cdot JQ_k^T(-B)\\
&\quad=1.
\end{split}
\tag{21.5.9}
\end{equation}
The second factor on the left is $D_k$ (Section~I.3.5).
Hence $D_k$ is injective on $\mathfrak{S}(1,\ell)$, a space of the finite dimension $4n_k$;
consequently $D_k$ is bijective, $\det D_k\ne0$, and $D_k^{-1}$ is the first factor:
\[
D_k^{-1}=b_k^{-1}+Q_kJ^*\cdot(D_B+\mu)^{-1}\cdot JQ_k^T(-B).
\]

(ii) \emph{The real part of the inverse.}
Let $v\in\mathfrak{S}(1,\ell)$ and $f:=JQ_k^T(-B)v\in\mathfrak{S}^{\mathrm{ch}}(\eta_k,\ell)$.
By Proposition~21.1\,(e), $JQ_k^T(-B)=(Q_kJ^*)^*$.
Therefore
\[
\left(v,Q_kJ^*\cdot(D_B+\mu)^{-1}\cdot JQ_k^T(-B)v\right)_1
=\left((Q_kJ^*)^*v,(D_B+\mu)^{-1}f\right)^{\mathrm{ch}}
=\left(f,(D_B+\mu)^{-1}f\right)^{\mathrm{ch}}.
\]
By (21.5.7), $\operatorname{Re}(D_B+\mu)\ge\mu\ge0$;
so by (21.5.3) and (21.5.1) the real part of the last number is non-negative.
Together with (i), and since $b_k$ is real,
\begin{equation}
\begin{aligned}
&\operatorname{Re}(v,D_k^{-1}v)_1=b_k^{-1}\cdot(v,v)_1+\operatorname{Re}\left(f,(D_B+\mu)^{-1}f\right)^{\mathrm{ch}}\ge b_k^{-1}\cdot(v,v)_1,\\
&\text{that is,}\qquad\operatorname{Re}D_k^{-1}\ge b_k^{-1}.
\end{aligned}
\tag{21.5.10}
\end{equation}

(iii) \emph{The real part of $D_k$.}
Let $v\in\mathfrak{S}(1,\ell)$ and $v':=D_kv$.
By (21.5.3) with $O=D_k$, and by (21.5.10) at $v'$,
\begin{equation}
\operatorname{Re}(v,D_kv)_1=\operatorname{Re}(v',D_k^{-1}v')_1\ge b_k^{-1}\cdot\|D_kv\|^2.
\tag{21.5.11}
\end{equation}
If $v\ne0$, then $D_kv\ne0$ by (i), and $b_k^{-1}>0$; hence $\operatorname{Re}(v,D_kv)_1>0$.
By (21.5.1) this is the assertion $\operatorname{Re}D_k>0$.
\end{proof}

\subsection*{6. Proof of (c)}

\begin{proof}[Proof of Proposition~21.5\,(c)]
\emph{The case $k\ge1$.}
(i) \emph{Real part.}
By Proposition~21.1\,(e), $P$ is an orthogonal projection of $\mathfrak{S}(1,\ell)$, and $b_L$ is real and positive (see 2.2).
By (21.5.4), for all $v$,
\[
\operatorname{Re}\left(v,(D_k+b_LP)v\right)_1=\operatorname{Re}(v,D_kv)_1+b_L\cdot\|Pv\|^2\ge\operatorname{Re}(v,D_kv)_1,
\]
and by Section~5\,(iii) the right side is positive for $v\ne0$.
This is $\operatorname{Re}(D_k+b_LP)>0$.

(ii) \emph{Adjoint.}
Consider first the three terms of $F_\infty$.
By Proposition~21.1\,(e), $D_B^*=\gamma_5D_B\gamma_5$.
Since $\mu$ is real and $\gamma_5^2=1$, $(\mu\cdot1)^*=\mu\cdot1=\gamma_5(\mu\cdot1)\gamma_5$.
The operator $JP_kJ^*$ is self-adjoint by Proposition~21.1\,(e) and commutes with $\gamma_5$ by Proposition~21.1\,(a); as $b_k$ is real,
\[
\left(b_k\cdot JP_kJ^*\right)^*=b_k\cdot JP_kJ^*=b_k\cdot JP_kJ^*\cdot\gamma_5\gamma_5=\gamma_5\left(b_k\cdot JP_kJ^*\right)\gamma_5.
\]
Adding the three relations gives $F_\infty^*=\gamma_5F_\infty\gamma_5$.
Taking inverses (see 3.2, and note $\gamma_5^{-1}=\gamma_5$),
\[
S_k^*=\left(F_\infty^*\right)^{-1}=\left(\gamma_5F_\infty\gamma_5\right)^{-1}=\gamma_5S_k\gamma_5.
\]
By Proposition~21.1\,(e) and $(O^*)^*=O$,
\[
(Q_kJ^*)^*=JQ_k^T(-B),\qquad\left(JQ_k^T(-B)\right)^*=Q_kJ^*.
\]
With (21.5.6) for $\mathrm{Mx}=\gamma_5$ it follows that
\begin{align*}
\left(Q_kJ^*\cdot S_k\cdot JQ_k^T(-B)\right)^*
&=\left(JQ_k^T(-B)\right)^*S_k^*\left(Q_kJ^*\right)^*\\
&=Q_kJ^*\cdot\gamma_5S_k\gamma_5\cdot JQ_k^T(-B)\\
&=\gamma_5\cdot Q_kJ^*\cdot S_k\cdot JQ_k^T(-B)\cdot\gamma_5,
\end{align*}
where $\gamma_5$ acts on $\mathfrak{S}^{\mathrm{ch}}(\eta_k,\ell)$ in the middle member and on $\mathfrak{S}(1,\ell)$ in the last member.
The number $b_k$ is real, and $b_k\cdot1=\gamma_5(b_k\cdot1)\gamma_5$.
Hence
\[
D_k^*=b_k-b_k^2\,\gamma_5\cdot Q_kJ^*\cdot S_k\cdot JQ_k^T(-B)\cdot\gamma_5=\gamma_5D_k\gamma_5.
\]
The operator $P$ is self-adjoint by Proposition~21.1\,(e) and commutes with $\gamma_5$ by Proposition~21.1\,(a), so that $(b_LP)^*=b_LP=\gamma_5(b_LP)\gamma_5$.
Adding,
\begin{equation}
(D_k+b_LP)^*=\gamma_5\,(D_k+b_LP)\,\gamma_5\qquad\text{on }\mathfrak{S}(1,\ell).
\tag{21.5.12}
\end{equation}

\emph{The case $k=0$.}
Here $F_0=D_B+\mu+b_L\cdot JPJ^*$ on $\mathfrak{S}^{\mathrm{ch}}(1,\ell)$ (Section~I.3.9 at $y=0$).
By Proposition~21.1\,(e), $JPJ^*$ is an orthogonal projection of $\mathfrak{S}^{\mathrm{ch}}(1,\ell)$.
By (21.5.4), (21.5.7) and $b_L>0$, for all $f\in\mathfrak{S}^{\mathrm{ch}}(1,\ell)$,
\[
\operatorname{Re}(f,F_0f)^{\mathrm{ch}}=\operatorname{Re}\left(f,(D_B+\mu)f\right)^{\mathrm{ch}}+b_L\cdot\|JPJ^*f\|^2\ge\mu\cdot(f,f)^{\mathrm{ch}},
\]
that is, $\operatorname{Re}F_0\ge\mu$.
As for the three terms of $F_0$:
$D_B^*=\gamma_5D_B\gamma_5$ by Proposition~21.1\,(e);
$(\mu\cdot1)^*=\mu\cdot1=\gamma_5(\mu\cdot1)\gamma_5$;
and $JPJ^*$ is self-adjoint by Proposition~21.1\,(e) and commutes with $\gamma_5$ by Proposition~21.1\,(a), whence
\[
\left(b_L\cdot JPJ^*\right)^*=b_L\cdot JPJ^*=\gamma_5\left(b_L\cdot JPJ^*\right)\gamma_5.
\]
Adding,
\begin{equation}
F_0^*=\gamma_5F_0\gamma_5\qquad\text{on }\mathfrak{S}^{\mathrm{ch}}(1,\ell).
\tag{21.5.13}
\end{equation}

\emph{Conclusion.}
By (21.5.2), by its second sentence for $k\ge1$ and by its first sentence with $\mathrm{cf}=\mu$ for $k=0$, the operator $D_k+b_LP$, resp.\ $F_0$, is bijective and has a non-vanishing determinant.
Hence $\Gamma_k=(D_k+b_LP)^{-1}$ ($k\ge1$) and $\Gamma_0=F_0^{-1}$ exist, and
\[
\delta Z_k=\det(D_k+b_LP)\ne0\quad(\text{size }4n_k),\qquad
\delta Z_0=\det F_0\ne0\quad(\text{size }192\,n_0).
\]
\end{proof}

\subsection*{7. Proof of (d)}

\begin{proof}[Proof of Proposition~21.5\,(d)]
Put
\[
\mathrm{Kf}:=
\begin{cases}
D_k+b_LP\ \text{ on }\mathrm{Vs}:=\mathfrak{S}(1,\ell),\ \text{with the form }(\cdot\,,\cdot)_1, & k\ge1,\\
F_0\ \text{ on }\mathrm{Vs}:=\mathfrak{S}^{\mathrm{ch}}(1,\ell),\ \text{with the form }(\cdot\,,\cdot)^{\mathrm{ch}}, & k=0.
\end{cases}
\]
By (c), $\mathrm{Kf}$ is bijective, and by (21.5.12), (21.5.13), $\mathrm{Kf}^*=\gamma_5\cdot\mathrm{Kf}\cdot\gamma_5$.
Let $y$ be real.
The operator $\mathrm{Kf}+i\gamma_5y$ is the operator of (d);
for $k=0$ this follows from Section~I.3.9, according to which
\[
F_y=D_B+\mu+b_L\cdot JPJ^*+i\gamma_5y=F_0+i\gamma_5y.
\]
By 3.2 and (21.5.5), $y$ being real, $(\mathrm{Kf}+i\gamma_5y)^*=\mathrm{Kf}^*-iy\gamma_5$.
Moreover $\mathrm{Kf}^*\cdot\gamma_5=\gamma_5\cdot\mathrm{Kf}\cdot\gamma_5\gamma_5=\gamma_5\cdot\mathrm{Kf}$.
Hence
\begin{equation}
\begin{aligned}
(\mathrm{Kf}+i\gamma_5y)^*(\mathrm{Kf}+i\gamma_5y)
&=\mathrm{Kf}^*\cdot\mathrm{Kf}+iy\left(\mathrm{Kf}^*\cdot\gamma_5-\gamma_5\cdot\mathrm{Kf}\right)+y^2\gamma_5^2
=\mathrm{Kf}^*\cdot\mathrm{Kf}+y^2,\\
\text{so that}\quad\|(\mathrm{Kf}+i\gamma_5y)f\|^2&=\|\mathrm{Kf}\,f\|^2+y^2\cdot\|f\|^2\qquad\text{for all }f\in\mathrm{Vs}.
\end{aligned}
\tag{21.5.14}
\end{equation}

\emph{The case $y\ne0$.}
By (21.5.14), $\|(\mathrm{Kf}+i\gamma_5y)f\|\ge|y|\cdot\|f\|$.
Hence $\mathrm{Kf}+i\gamma_5y$ is injective, and therefore bijective, $\mathrm{Vs}$ being finite dimensional.
For $f'\in\mathrm{Vs}$ and $f:=(\mathrm{Kf}+i\gamma_5y)^{-1}f'$ one has $\|f\|\le\|f'\|/|y|$.
Thus $|(\mathrm{Kf}+i\gamma_5y)^{-1}|_{\mathrm{op}}\le1/|y|$.

\emph{The case $y=0$.}
Here $\mathrm{Kf}+i\gamma_5y=\mathrm{Kf}$, which is bijective by (c).

(In the notation of Sections~I.3.5 and I.3.9:
$\Gamma_{k,y}=(D_k+b_LP+i\gamma_5y)^{-1}$ for $k\ge1$ and $S_{0,y}=F_y^{-1}$ for $k=0$ exist for every real $y$, and their operator norms are at most $1/|y|$ for $y\ne0$.)

\emph{The operators $S_{k,y}$, $k\ge1$.}
By (a), $F_\infty$ is invertible, so that Proposition~21.3 applies at every real $y$.
Its condition (i) has just been established;
therefore (ii) holds: $F_y$ is invertible on $\mathfrak{S}^{\mathrm{ch}}(\eta_k,\ell)$, that is, $S_{k,y}=F_y^{-1}$ exists.
(A second proof, which rests on Proposition~21.1\,(a), (e) alone and yields a bound on $S_{k,y}$, is (21.5.22).)
\end{proof}

\emph{The Hermitian form of $\mathrm{Kf}$.}
Put $\mathrm{Kh}:=\mathrm{Kf}\cdot\gamma_5$.
By (21.5.5) and $\mathrm{Kf}^*=\gamma_5\cdot\mathrm{Kf}\cdot\gamma_5$,
\begin{equation}
\begin{aligned}
&\mathrm{Kh}^*=\gamma_5\cdot\mathrm{Kf}^*=\gamma_5\gamma_5\cdot\mathrm{Kf}\cdot\gamma_5=\mathrm{Kh},\\
&\mathrm{Kf}+i\gamma_5y=(\mathrm{Kh}+iy)\,\gamma_5,\qquad
(\mathrm{Kf}+i\gamma_5y)^{-1}=\gamma_5\,(\mathrm{Kh}+iy)^{-1}.
\end{aligned}
\tag{21.5.15}
\end{equation}
Thus $\mathrm{Kh}$ is Hermitian, and it is bijective as the product of two bijective maps.
(By the spectral theorem for Hermitian matrices \cite[Section~2.5]{HornJohnson2013}, $\mathrm{Kh}$ has an orthonormal basis of eigenvectors with real eigenvalues $\mathrm{ew}_j\ne0$;
the eigenvalues $\mathrm{ew}_j+iy$ of $\mathrm{Kh}+iy$ have modulus at least $|y|$, and $\gamma_5$ is unitary.
This yields (d) once more.)

\subsection*{8. Uniformity; explicit forms of the strict inequalities; a second proof of the last clause of (d)}

Proposition~21.5 and the facts established in Sections~4--7 hold for all data $k,\ell,e_k,B$ of the setting (S21) with real $B$, and for every $\mu>0$:
for every integer $k\ge0$, every $\ell\in\mathrm{Cl}$, every $e_k>0$, every real $B$ and every $\mu>0$.
No relation between these data is assumed, and no bound on $e_k$, on $\mu$ or on $B$ besides $e_k>0$ and $\mu>0$.
In the inequalities of Proposition~21.5\,(a), of (c) at $k=0$ and of (d) the right-hand sides are $\mu$, $\mu$ and $\frac{1}{|y|}$, respectively;
none of $k$, $\ell$, $e_k$, $B$, $\rho$, $w$, $L$, $b$ occurs in them.
A cutoff $N$ (Section~I.2.5) and a torus (Section~I.6.1) enter the assertions of Proposition~21.5, and these bounds, only through the data $(k,\ell,\mu,e_k,B)$.

For $k\ge1$ the strict inequalities $\mathrm{Re}\,D_k>0$ and $\mathrm{Re}\left(D_k+b_L\cdot P\right)>0$ hold with lower bounds in which only $\mu$ and $b_k$ occur.
These lower bounds and the remaining bounds among (21.5.16)--(21.5.22) are proved below, in addition to the assertions of Proposition~21.5.
The fact (21.5.22) gives a proof of the last clause of (d) that does not use the implication (i)$\Rightarrow$(ii) of Proposition~21.3.

\begin{named}{(21.5.16)}
Let $k\ge0$. Then
\[
\left|\left(D_B+\mu\right)^{-1}\right|_{\mathrm{op}}\le\frac{1}{\mu},\qquad
\mathrm{Re}\left(D_B+\mu\right)^{-1}\ge0,\qquad
\det\left(D_B+\mu\right)\ne0.
\]
Moreover $\left|S_k\right|_{\mathrm{op}}\le\frac{1}{\mu}$ if $k\ge1$, and $\left|\Gamma_0\right|_{\mathrm{op}}\le\frac{1}{\mu}$ if $k=0$.
\end{named}

\begin{proof}
The real part of $D_B+\mu$ is $\ge\mu$ by (21.5.7);
for $k\ge1$ the real part of $F_\infty$ is $\ge\mu$ by (21.5.8);
for $k=0$ the real part of $F_0$ is $\ge\mu$ by Section~6.
Hence (21.5.2) with $\mathrm{cf}=\mu$ applies to each of these three operators, and $S_k=F_\infty^{-1}$, $\Gamma_0=F_0^{-1}$.
Since $\mathrm{Re}\left(D_B+\mu\right)\ge\mu>0$, (21.5.3) gives $\mathrm{Re}\left(D_B+\mu\right)^{-1}\ge0$.
\end{proof}

\begin{named}{(21.5.17)}
Let $k\ge1$.
For $v\in\mathfrak{S}(1,\ell)$ one has $\left\|JQ_k^{T}(-B)\,v\right\|=\|v\|$, and for $f\in\mathfrak{S}^{\mathrm{ch}}(\eta_k,\ell)$
\[
\left\|Q_kJ^{*}f\right\|=\left\|JP_kJ^{*}f\right\|\le\|f\|.
\]
\end{named}

\begin{proof}
By Proposition~21.1\,(e) and the rule $\left(O^{*}\right)^{*}=O$ of subsection~3.2, $\left(JQ_k^{T}(-B)\right)^{*}=Q_kJ^{*}$;
by Proposition~21.1\,(a), $\left(Q_kJ^{*}\right)\left(JQ_k^{T}(-B)\right)=1$ on $\mathfrak{S}(1,\ell)$.
Therefore
\[
\left\|JQ_k^{T}(-B)\,v\right\|^2
=\left(v,\left(JQ_k^{T}(-B)\right)^{*}JQ_k^{T}(-B)\,v\right)_1
=\left(v,Q_kJ^{*}JQ_k^{T}(-B)\,v\right)_1
=(v,v)_1.
\]
Likewise, by Proposition~21.1\,(e) and the definition $P_k=Q_k^{T}(-B)Q_k$ (Section~I.3.3), $\left(Q_kJ^{*}\right)^{*}Q_kJ^{*}=JQ_k^{T}(-B)\,Q_kJ^{*}=JP_kJ^{*}$, and $JP_kJ^{*}$ is an orthogonal projection;
with (21.5.4) one obtains
\[
\left\|Q_kJ^{*}f\right\|^2
=\left(f,\left(Q_kJ^{*}\right)^{*}Q_kJ^{*}f\right)^{\mathrm{ch}}
=\left(f,JP_kJ^{*}f\right)^{\mathrm{ch}}
=\left\|JP_kJ^{*}f\right\|^2.
\]
Finally, if $\mathrm{Pj}$ is an orthogonal projection, then $\left(\mathrm{Pj}f,(1-\mathrm{Pj})f\right)=\left(f,\left(\mathrm{Pj}-\mathrm{Pj}^2\right)f\right)=0$;
hence $\|f\|^2=\|\mathrm{Pj}f\|^2+\|(1-\mathrm{Pj})f\|^2\ge\|\mathrm{Pj}f\|^2$.
\end{proof}

\begin{named}{(21.5.18)}
Let $k\ge1$. Then
\[
\left|D_k^{-1}\right|_{\mathrm{op}}\le b_k^{-1}+\mu^{-1},\qquad\left|D_k\right|_{\mathrm{op}}\le b_k.
\]
\end{named}

\begin{proof}
Let $v\in\mathfrak{S}(1,\ell)$.
One applies, in this order, the formula $D_k^{-1}=b_k^{-1}+Q_kJ^{*}\left(D_B+\mu\right)^{-1}JQ_k^{T}(-B)$ of Proposition~21.5\,(b), the inequality of (21.5.17), the bound on $\left(D_B+\mu\right)^{-1}$ in (21.5.16), and the identity of (21.5.17):
\begin{align*}
\left\|D_k^{-1}v\right\|
&\le b_k^{-1}\cdot\|v\|+\left\|Q_kJ^{*}\left(D_B+\mu\right)^{-1}JQ_k^{T}(-B)\,v\right\|\\
&\le b_k^{-1}\cdot\|v\|+\left\|\left(D_B+\mu\right)^{-1}JQ_k^{T}(-B)\,v\right\|\\
&\le b_k^{-1}\cdot\|v\|+\mu^{-1}\cdot\left\|JQ_k^{T}(-B)\,v\right\|
=\left(b_k^{-1}+\mu^{-1}\right)\cdot\|v\|.
\end{align*}
By (21.5.10), $\mathrm{Re}\,D_k^{-1}\ge b_k^{-1}$.
So (21.5.2) applies to $O=D_k^{-1}$ with $\mathrm{cf}=b_k^{-1}$, and it gives $\left|D_k\right|_{\mathrm{op}}=\left|\left(D_k^{-1}\right)^{-1}\right|_{\mathrm{op}}\le b_k$.
\end{proof}

\begin{named}{(21.5.19)}
Let $k\ge1$. Then
\[
\mathrm{Re}\,D_k\ge\frac{b_k\cdot\mu^2}{\left(\mu+b_k\right)^2}\ge\frac{b_-\cdot\mu^2}{\left(\mu+b_+\right)^2}.
\]
\end{named}

\begin{proof}
Let $v\in\mathfrak{S}(1,\ell)$.
By (21.5.18), $\|v\|=\left\|D_k^{-1}D_kv\right\|\le\left(b_k^{-1}+\mu^{-1}\right)\cdot\left\|D_kv\right\|$.
Inserting this into (21.5.11) one obtains
\[
\mathrm{Re}\,(v,D_kv)_1\ge b_k^{-1}\cdot\left(b_k^{-1}+\mu^{-1}\right)^{-2}\cdot\|v\|^2,
\]
and
\[
b_k^{-1}\cdot\left(b_k^{-1}+\mu^{-1}\right)^{-2}=b_k^{-1}\cdot\frac{\left(b_k\cdot\mu\right)^2}{\left(\mu+b_k\right)^2}=\frac{b_k\cdot\mu^2}{\left(\mu+b_k\right)^2};
\]
by (21.5.1) this is the first inequality.
For the second one recall that $b_-\le b_k\le b_+$ (Section~I.3.4, since $1-L^{-1}\le1-L^{-k}<1$):
the numerator is at least $b_-\cdot\mu^2$, and the denominator is at most $\left(\mu+b_+\right)^2$.
\end{proof}

\begin{named}{(21.5.20)}
Let $k\ge1$. Then
\[
\mathrm{Re}\left(D_k+b_L\cdot P\right)\ge\frac{b_k\cdot\mu^2}{\left(\mu+b_k\right)^2},\qquad
\left|\Gamma_k\right|_{\mathrm{op}}\le\frac{\left(\mu+b_k\right)^2}{b_k\cdot\mu^2}\le\frac{\left(\mu+b_+\right)^2}{b_-\cdot\mu^2}.
\]
\end{named}

\begin{proof}
Section~6\,(i) gives $\mathrm{Re}\left(v,\left(D_k+b_L\cdot P\right)v\right)_1\ge\mathrm{Re}\,(v,D_kv)_1$ for all $v\in\mathfrak{S}(1,\ell)$;
combined with (21.5.19), this is the first inequality.
The second follows from (21.5.2), applied to $D_k+b_L\cdot P$, whose inverse is $\Gamma_k$;
the third follows from $b_-\le b_k\le b_+$.
\end{proof}

\begin{named}{(21.5.21)}
Let $k\ge0$, let $y$ be real, and let $\mathrm{Kf}$ be the operator of Section~7 on the space $\mathrm{Vs}$ of Section~7. Then
\[
\left|\left(\mathrm{Kf}+i\gamma_5y\right)^{-1}\right|_{\mathrm{op}}
=\left(y^2+\left|\mathrm{Kf}^{-1}\right|_{\mathrm{op}}^{-2}\right)^{-\frac{1}{2}}
\le\min\left\{\frac{1}{|y|},\left|\mathrm{Kf}^{-1}\right|_{\mathrm{op}}\right\},
\]
with $\frac{1}{|y|}:=\infty$ at $y=0$.
\end{named}

\begin{proof}
Let $O$ be a bijective linear map of $\mathrm{Vs}$ onto itself.
Writing $f'=Of$, one has
\[
\left|O^{-1}\right|_{\mathrm{op}}
=\sup\left\{\frac{\left\|O^{-1}f'\right\|}{\|f'\|}:f'\ne0\right\}
=\sup\left\{\frac{\|f\|}{\|Of\|}:f\ne0\right\},
\]
and therefore $\left|O^{-1}\right|_{\mathrm{op}}^{-1}=\inf\left\{\|Of\|:\|f\|=1\right\}$.
The operators $\mathrm{Kf}$ and $\mathrm{Kf}+i\gamma_5y$ are bijective by Section~7.
For $\|f\|=1$, (21.5.14) reads $\left\|\left(\mathrm{Kf}+i\gamma_5y\right)f\right\|^2=\left\|\mathrm{Kf}\,f\right\|^2+y^2$.
Taking the infimum over these $f$ on both sides gives the identity;
the inequality holds because $y^2+\left|\mathrm{Kf}^{-1}\right|_{\mathrm{op}}^{-2}$ is at least $y^2$ and at least $\left|\mathrm{Kf}^{-1}\right|_{\mathrm{op}}^{-2}$.
\end{proof}

\begin{named}{(21.5.22)}
Let $k\ge1$ and let $y$ be real. Put
\[
\mathrm{Pj}_1:=JP_kJ^{*}-JP_{k+1}J^{*},\qquad\mathrm{Pj}_2:=JP_{k+1}J^{*},
\]
\[
\mathrm{pa}_y:=\frac{b_k\cdot y^2}{b_k^2+y^2},\qquad
\mathrm{pb}_y:=\frac{b_k\cdot\left(b_L\cdot\left(b_k+b_L\right)+y^2\right)}{\left(b_k+b_L\right)^2+y^2}.
\]
Then $\mathrm{pa}_y\ge0$, $\mathrm{pb}_y>0$, the operators $\mathrm{Pj}_1$ and $\mathrm{Pj}_2$ are orthogonal projections of $\mathfrak{S}^{\mathrm{ch}}(\eta_k,\ell)$, and every $f\in\mathfrak{S}^{\mathrm{ch}}(\eta_k,\ell)$ satisfies
\begin{multline*}
\mathrm{Re}\,(f,F_yf)^{\mathrm{ch}}
=\frac{\rho}{2\eta_k}\cdot\left(f,\nabla_B^{*}\nabla_Bf\right)^{\mathrm{ch}}+\mu\cdot(f,f)^{\mathrm{ch}}\\
+\mathrm{pa}_y\cdot\left\|\mathrm{Pj}_1f\right\|^2+\mathrm{pb}_y\cdot\left\|\mathrm{Pj}_2f\right\|^2
\ge\mu\cdot(f,f)^{\mathrm{ch}}.
\end{multline*}
Consequently $\mathrm{Re}\,F_y\ge\mu$, the operator $F_y$ is bijective, $S_{k,y}=F_y^{-1}$ exists, and $\left|S_{k,y}\right|_{\mathrm{op}}\le\frac{1}{\mu}$.
\end{named}

\begin{proof}
\emph{(1) The two matrices.}
Every matrix that occurs in this step is a polynomial in $\gamma_5$ or the inverse of such a polynomial;
any two of them commute.
From $\gamma_5^2=1$ one gets $\left(i\gamma_5y\right)^2=-y^2$ and, for every real number $\mathrm{cf}>0$ (see subsection~2.2),
\[
\left(\mathrm{cf}+i\gamma_5y\right)^{-1}=\frac{\mathrm{cf}-i\gamma_5y}{\mathrm{cf}^2+y^2}.
\]
With $\mathrm{cf}=b_k$, the definition of $\boldsymbol{\alpha}_y$ in Section~I.3.5 becomes
\[
\boldsymbol{\alpha}_y
=\frac{b_k\cdot i\gamma_5y\left(b_k-i\gamma_5y\right)}{b_k^2+y^2}
=\frac{b_k\cdot y^2+i\gamma_5\,b_k^2\cdot y}{b_k^2+y^2}
=\mathrm{pa}_y+i\gamma_5\,\mathrm{qa}_y,
\qquad
\mathrm{qa}_y:=\frac{b_k^2\cdot y}{b_k^2+y^2}.
\]
The definitions of $\boldsymbol{\alpha}_y$ and $\boldsymbol{\beta}_y$ in Section~I.3.5 give
\[
\boldsymbol{\alpha}_y+\boldsymbol{\beta}_y
=b_k\cdot\left[i\gamma_5y\left(b_k+b_L+i\gamma_5y\right)+b_k\cdot b_L\right]\left(b_k+b_L+i\gamma_5y\right)^{-1}\left(b_k+i\gamma_5y\right)^{-1},
\]
and the bracket factorises:
\[
b_k\cdot b_L+\left(b_k+b_L\right)\cdot i\gamma_5y+\left(i\gamma_5y\right)^2=\left(b_k+i\gamma_5y\right)\left(b_L+i\gamma_5y\right).
\]
Hence, by the formula for the inverse with $\mathrm{cf}=b_k+b_L$,
\begin{align*}
\boldsymbol{\alpha}_y+\boldsymbol{\beta}_y
&=b_k\cdot\left(b_L+i\gamma_5y\right)\left(b_k+b_L+i\gamma_5y\right)^{-1}
=\frac{b_k\cdot\left(b_L+i\gamma_5y\right)\left(b_k+b_L-i\gamma_5y\right)}{\left(b_k+b_L\right)^2+y^2}\\
&=\frac{b_k\cdot\left(b_L\cdot\left(b_k+b_L\right)+y^2+i\gamma_5\,b_k\cdot y\right)}{\left(b_k+b_L\right)^2+y^2}
=\mathrm{pb}_y+i\gamma_5\,\mathrm{qb}_y,
\qquad
\mathrm{qb}_y:=\frac{b_k^2\cdot y}{\left(b_k+b_L\right)^2+y^2};
\end{align*}
the third member rests on the expansion
\[
\left(b_L+i\gamma_5y\right)\left(b_k+b_L-i\gamma_5y\right)=b_L\cdot\left(b_k+b_L\right)+\left[\left(b_k+b_L\right)-b_L\right]\cdot i\gamma_5y-\left(i\gamma_5y\right)^2.
\]
The numbers $\mathrm{pa}_y$, $\mathrm{qa}_y$, $\mathrm{pb}_y$, $\mathrm{qb}_y$ are real;
moreover $\mathrm{pa}_y\ge0$ and $\mathrm{pb}_y>0$, since $b_k>0$ and $b_L>0$ (subsection~2.2).

\emph{(2) The two projections.}
By Proposition~21.1\,(e), $JP_kJ^{*}$ and $JP_{k+1}J^{*}$ are orthogonal projections of $\mathfrak{S}^{\mathrm{ch}}(\eta_k,\ell)$.
By Proposition~21.1\,(a), their products in either order equal $JP_{k+1}J^{*}$, and both operators commute with $\gamma_5$.
It follows that $\mathrm{Pj}_2$ is an orthogonal projection, that $\mathrm{Pj}_1^{*}=\mathrm{Pj}_1$, that
\[
\mathrm{Pj}_1^2=JP_kJ^{*}-JP_{k+1}J^{*}-JP_{k+1}J^{*}+JP_{k+1}J^{*}=\mathrm{Pj}_1,
\]
and that $\mathrm{Pj}_1$ and $\mathrm{Pj}_2$ commute with $\gamma_5$.
Since $JP_kJ^{*}=\mathrm{Pj}_1+\mathrm{Pj}_2$, step~(1) yields
\begin{align*}
\boldsymbol{\alpha}_y\,JP_kJ^{*}+\boldsymbol{\beta}_y\,JP_{k+1}J^{*}
&=\boldsymbol{\alpha}_y\,\mathrm{Pj}_1+\left(\boldsymbol{\alpha}_y+\boldsymbol{\beta}_y\right)\mathrm{Pj}_2\\
&=\left(\mathrm{pa}_y+i\gamma_5\,\mathrm{qa}_y\right)\mathrm{Pj}_1+\left(\mathrm{pb}_y+i\gamma_5\,\mathrm{qb}_y\right)\mathrm{Pj}_2.
\end{align*}

\emph{(3) Real parts.}
Let $\mathrm{cf}_1$, $\mathrm{cf}_2$ be real numbers, and let $\mathrm{Pj}$ be an orthogonal projection of $\mathfrak{S}^{\mathrm{ch}}(\eta_k,\ell)$ that commutes with $\gamma_5$.
By subsection~3.2 and (21.5.5),
\[
\left(\left(\mathrm{cf}_1+i\,\mathrm{cf}_2\cdot\gamma_5\right)\mathrm{Pj}\right)^{*}
=\mathrm{Pj}\left(\mathrm{cf}_1-i\,\mathrm{cf}_2\cdot\gamma_5\right)
=\left(\mathrm{cf}_1-i\,\mathrm{cf}_2\cdot\gamma_5\right)\mathrm{Pj},
\]
hence $\mathrm{Re}\left(\left(\mathrm{cf}_1+i\,\mathrm{cf}_2\cdot\gamma_5\right)\mathrm{Pj}\right)=\mathrm{cf}_1\cdot\mathrm{Pj}$.
By Section~I.3.5, $F_y=D_B+\mu+\boldsymbol{\alpha}_y\,JP_kJ^{*}+\boldsymbol{\beta}_y\,JP_{k+1}J^{*}$.
With step~(2) and Section~4 one obtains
\[
\mathrm{Re}\,F_y
=\mathrm{Re}\left(D_B+\mu\right)+\mathrm{pa}_y\cdot\mathrm{Pj}_1+\mathrm{pb}_y\cdot\mathrm{Pj}_2
=\frac{\rho}{2\eta_k}\cdot\nabla_B^{*}\nabla_B+\mu+\mathrm{pa}_y\cdot\mathrm{Pj}_1+\mathrm{pb}_y\cdot\mathrm{Pj}_2.
\]
In view of (21.5.1), (21.5.4) and (21.5.7), this is the display of (21.5.22), the inequality included.
The last sentence of (21.5.22) is (21.5.2) with $\mathrm{cf}=\mu$.

At $y=0$ one has $\mathrm{pa}_0=0$ and $\mathrm{pb}_0=\frac{b_k\cdot b_L}{b_k+b_L}=\frac{b_{k+1}}{L}$, by $b_{k+1}=\frac{b\cdot b_k}{b_k+b_L}$ (Section~I.3.4).
As $|y|\to\infty$, both $\mathrm{pa}_y$ and $\mathrm{pb}_y$ tend to $b_k$, which is the coefficient of $JP_kJ^{*}=\mathrm{Pj}_1+\mathrm{Pj}_2$ in $F_\infty$.
\end{proof}

The right-hand sides of (21.5.16)--(21.5.20) contain only $\mu$, $b_k$, $b_-$ and $b_+$, that is, only $\mu$, $b$, $L$ and, through $b_k$, the level;
the last members of (21.5.19) and (21.5.20) do not contain the level.
None of them contains $\ell$, $e_k$, $B$, $\rho$ or $w$.
The bound $\frac{1}{\mu}$ in (21.5.22) contains none of $y$, $k$, $\ell$, $e_k$, $B$, $\rho$, $w$, $L$, $b$.

\subsection*{8a. The clauses of Proposition~21.1\,(a), (e) read off the displays of Part~I ($B$ real)}

Throughout this section $B$ is real.
In this case the parts of Proposition~21.1\,(a), (e) recalled in subsection~2.3 can be read off directly from the displays of Part~I;
for completeness they are derived here in this way, in the setting in which they enter Sections~4--8, with the sums over the clusters written out.

\paragraph{8a.1. Sums over the torus.}
Let $\mathrm{sf}$ be an $\ell\mathbb{Z}^3$-periodic complex function on $\mathbb{L}^{\mathrm{f}}_k$.
Then
\[
\int\mathrm{d}x\,\mathrm{sf}(x)=\eta_k^3\cdot\frac{1}{48}\sum_c\sum_{g\in O_{\mathrm{cub}}}\mathrm{sf}(c,g),
\]
where $c$ ranges over a fundamental domain of $\ell\mathbb{Z}^3$ in $\mathbb{L}_k$, is a sum over $T^{\mathrm{f}}_k$.
By (I.s1) the map $x\mapsto x^{(i)}$ is a bijection of $T^{\mathrm{f}}_k$ for each $i\in\{0,1,2,3\}$.
Hence
\begin{equation}
\int\mathrm{d}x\,\mathrm{sf}\left(x^{(i)}\right)=\int\mathrm{d}x\,\mathrm{sf}(x).
\tag{21.5.23}
\end{equation}

\paragraph{8a.2. The hops.}
Let $x=(c,g)$.
By (I.2.2), $\left(R_i(B)f\right)(x)=U_B\left(x,x^{(i)}\right)f\left(x^{(i)}\right)$.
For $i=0,1,2$ the site $x^{(i)}=\left(c,gs_i\right)$ lies in the cluster of $c$, and $U_B\left(x,x^{(i)}\right)=1$.
For $i=3$ one has, by (I.1.2) and (I.1.6),
\begin{gather*}
x^{(3)}=\left(c+\eta_k\cdot ge_1,gs_3\right),\qquad
c_3=c+\eta_k\cdot ge_1,\\
U_B\left(x,x^{(3)}\right)=u_B(c,c_3)=\exp\left(ie_k\cdot\eta_k\cdot B(c,c_3)\right).
\end{gather*}
The function $R_i(B)f$ is antiperiodic, since $\left(x+\ell e_\mu\right)^{(i)}=x^{(i)}+\ell e_\mu$ (I.s1) and $B$ is periodic.
Furthermore $\left(x^{(i)}\right)^{(i)}=x$ by (I.s1), $B(c_3,c)=-B(c,c_3)$ by Section~I.1.1, and $e_k$, $\eta_k$, $B$ are real.
Consequently
\begin{equation}
U_B\left(x^{(i)},x\right)=U_B\left(x,x^{(i)}\right)^{-1}=\overline{U_B\left(x,x^{(i)}\right)}.
\tag{21.5.24}
\end{equation}
It follows that $\left(R_i(B)^2f\right)(x)=U_B\left(x,x^{(i)}\right)U_B\left(x^{(i)},x\right)f(x)=f(x)$.
Now let $f',f\in\mathfrak{S}^{\mathrm{ch}}(\eta_k,\ell)$, and put $\mathrm{sf}(x):=\overline{f'(x)}\cdot\left(U_B\left(x,x^{(i)}\right)f\left(x^{(i)}\right)\right)$.
This function is periodic, as the product of two antiperiodic factors and a periodic one, and (21.5.24) gives
\[
\mathrm{sf}\left(x^{(i)}\right)
=\overline{f'\left(x^{(i)}\right)}\cdot\left(U_B\left(x^{(i)},x\right)f(x)\right)
=\overline{U_B\left(x,x^{(i)}\right)f'\left(x^{(i)}\right)}\cdot f(x).
\]
By (21.5.23),
\[
\left(f',R_i(B)f\right)^{\mathrm{ch}}=\int\mathrm{d}x\,\mathrm{sf}(x)=\int\mathrm{d}x\,\mathrm{sf}\left(x^{(i)}\right)=\left(R_i(B)f',f\right)^{\mathrm{ch}}.
\]
Therefore
\begin{equation}
R_i(B)^{*}=R_i(B),\qquad R_i(B)^2=1,\qquad\left\|R_i(B)f\right\|^2=\left(f,R_i(B)^2f\right)^{\mathrm{ch}}=\|f\|^2.
\tag{21.5.25}
\end{equation}

\paragraph{8a.3. $\nabla_B^{*}\nabla_B\ge0$.}
By (I.2.4), $\nabla_B^{*}\nabla_B=\sum_iw_i\cdot\left(1-R_i(B)\right)$;
by (21.5.25) this operator is self-adjoint.
The number $\left(f,R_i(B)f\right)^{\mathrm{ch}}$ is real, and
\[
\left\|R_i(B)f-f\right\|^2=\left\|R_i(B)f\right\|^2+\|f\|^2-2\,\mathrm{Re}\left(f,R_i(B)f\right)^{\mathrm{ch}}=2\left(f,\left(1-R_i(B)\right)f\right)^{\mathrm{ch}}.
\]
As $w_i>0$, one obtains, with $x=(c,g)$,
\begin{equation}
\left(f,\nabla_B^{*}\nabla_Bf\right)^{\mathrm{ch}}
=\frac{1}{2}\sum_{i=0}^{3}w_i\cdot\eta_k^3\cdot\frac{1}{48}\sum_c\sum_{g\in O_{\mathrm{cub}}}\left|U_B\left(x,x^{(i)}\right)f\left(x^{(i)}\right)-f(x)\right|_2^2\ge0.
\tag{21.5.26}
\end{equation}
By (21.5.24) the terms of $x$ and of $x^{(i)}$ are equal;
thus (21.5.26) is the sum over the unoriented bonds $\left\{x,x^{(i)}\right\}$ that appears in (I.O2).

\paragraph{8a.4. $J_B$ is anti-self-adjoint.}
Let $\gamma\cdot\hat d_i$ also denote the operator $\left(\left(\gamma\cdot\hat d_i\right)f\right)(x):=\left(\gamma\cdot\hat d_i(x)\right)f(x)$, where $\hat d_i(x)=-g\hat n_i$ (I.1.2) is a real vector that depends on the label only.
The matrices $\gamma\cdot\hat d_i(x)$ are Hermitian (I.c1), and the computation of subsection~3.4 proceeds site by site;
therefore $\left(\gamma\cdot\hat d_i\right)^{*}=\gamma\cdot\hat d_i$.
By (I.s2), $\hat d_i\left(x^{(i)}\right)=-\hat d_i(x)$, so that
\[
\left(R_i(B)\left(\gamma\cdot\hat d_i\right)f\right)(x)
=U_B\left(x,x^{(i)}\right)\left(\gamma\cdot\hat d_i\left(x^{(i)}\right)\right)f\left(x^{(i)}\right)
=-\left(\gamma\cdot\hat d_i(x)\right)\left(R_i(B)f\right)(x).
\]
By (I.2.4), $J_B=\sum_iw_i\left(\gamma\cdot\hat d_i\right)R_i(B)$, and (21.5.25) gives
\begin{equation}
J_B^{*}=\sum_iw_i\,R_i(B)\left(\gamma\cdot\hat d_i\right)=-\sum_iw_i\left(\gamma\cdot\hat d_i\right)R_i(B)=-J_B.
\tag{21.5.27}
\end{equation}

\paragraph{8a.5. The operator $D_B$.}
By (I.2.3) and (I.2.1),
\begin{align*}
\left(D_Bf\right)(x)
&=\eta_k^{-1}\sum_i\frac{w_i}{2}\left(\rho-\gamma\cdot\hat d_i(x)\right)\left[f(x)-\left(R_i(B)f\right)(x)\right]\\
&=\eta_k^{-1}\left\{\frac{\rho}{2}\left(\nabla_B^{*}\nabla_Bf\right)(x)+\frac{1}{2}\left(J_Bf\right)(x)-\frac{1}{2}\left(\gamma\cdot\left(\sum_iw_i\,\hat d_i(x)\right)\right)f(x)\right\},
\end{align*}
and $\sum_iw_i\,\hat d_i(x)=0$ by (I.s3);
this is (I.2.4).
Since $\rho$ and $\eta_k$ are real, $\rho\ge0$ and $\eta_k>0$, the relations (21.5.25)--(21.5.27) yield
\begin{equation}
D_B^{*}=\eta_k^{-1}\left[\frac{\rho}{2}\cdot\nabla_B^{*}\nabla_B-\frac{1}{2}J_B\right],\qquad
\mathrm{Re}\,D_B=\frac{\rho}{2\eta_k}\cdot\nabla_B^{*}\nabla_B\ge0.
\tag{21.5.28}
\end{equation}
The matrix $\gamma_5$ commutes with $R_i(B)$, which acts by a scalar phase and a map of the sites, and $\gamma_5\left(\gamma\cdot\hat d_i(x)\right)=-\left(\gamma\cdot\hat d_i(x)\right)\gamma_5$ by (I.c1).
Together with $\gamma_5^2=1$ this gives $\gamma_5\nabla_B^{*}\nabla_B\gamma_5=\nabla_B^{*}\nabla_B$ and $\gamma_5J_B\gamma_5=-J_B$, hence
\begin{equation}
\gamma_5D_B\gamma_5=\eta_k^{-1}\left[\frac{\rho}{2}\cdot\nabla_B^{*}\nabla_B-\frac{1}{2}J_B\right]=D_B^{*}.
\tag{21.5.29}
\end{equation}

\paragraph{8a.6. $Q_kJ^{*}$ and $JQ_k^{T}(-B)$, $k\ge1$.}
The blocks $\Delta_z$, $z\in\mathbb{Z}^3$, form a partition of $\mathbb{L}_k$ (Sections~I.1.1, I.1.4; $L$ is odd);
moreover $\Delta_{z+\ell e_\mu}=\Delta_z+\ell e_\mu$ and $\#|\Delta_z|=\eta_k^{-3}=L^{3k}$.
Hence, if $z$ ranges over a fundamental domain of $\ell\mathbb{Z}^3$ in $\mathbb{Z}^3$, the union of the blocks $\Delta_z$ is a fundamental domain of $\ell\mathbb{Z}^3$ in $\mathbb{L}_k$.
For real $B$ the exponent of $\varphi_B(z,c)=\exp\left(ie_k\cdot\eta_k\cdot\left(\tau_kB\right)(z,c)\right)$ (I.3.3) is purely imaginary, so that $\overline{\varphi_B(z,c)}=\varphi_B(z,c)^{-1}$.
The two displays for $Q_kJ^{*}f$ and $JQ_k^{T}(-B)v$ in subsection~2.2, together with the displays of subsection~3.1, give
\begin{equation}
\begin{aligned}
\left(v,Q_kJ^{*}f\right)_1
&=\sum_z\overline{v(z)}\cdot\left[\eta_k^3\sum_{c\in\Delta_z}\varphi_B(z,c)\cdot\frac{1}{48}\sum_gf(c,g)\right]\\
&=\eta_k^3\cdot\frac{1}{48}\sum_z\sum_{c\in\Delta_z}\sum_g\overline{\varphi_B(z,c)^{-1}v(z)}\cdot f(c,g)
=\left(JQ_k^{T}(-B)v,f\right)^{\mathrm{ch}},\\
\left(Q_kJ^{*}JQ_k^{T}(-B)v\right)(z)
&=\eta_k^3\sum_{c\in\Delta_z}\varphi_B(z,c)\cdot\frac{1}{48}\sum_g\varphi_B(z,c)^{-1}v(z)\\
&=\eta_k^3\cdot\eta_k^{-3}\cdot\frac{48}{48}\cdot v(z)=v(z).
\end{aligned}
\tag{21.5.30}
\end{equation}
Therefore $\left(Q_kJ^{*}\right)^{*}=JQ_k^{T}(-B)$ and $\left(Q_kJ^{*}\right)\left(JQ_k^{T}(-B)\right)=1$.
The operator $JP_kJ^{*}=\left(Q_kJ^{*}\right)^{*}\left(Q_kJ^{*}\right)$ is self-adjoint,
\[
\left(JP_kJ^{*}\right)^2=\left(Q_kJ^{*}\right)^{*}\left[\left(Q_kJ^{*}\right)\left(Q_kJ^{*}\right)^{*}\right]\left(Q_kJ^{*}\right)=JP_kJ^{*},
\]
and $JP_kJ^{*}$ commutes with every constant matrix by (21.5.6).

\paragraph{8a.7. The one-step averages.}
The $L$-blocks $B(y')$, $y'\in L\mathbb{Z}^3$, form a partition of $\mathbb{Z}^3$;
one has $B\left(y'+\ell e_\mu\right)=B(y')+\ell e_\mu$, because $\ell$ is a multiple of $L$, and $\#|B(y')|=L^3$.
Hence, if $y'$ ranges over a fundamental domain of $\ell\mathbb{Z}^3$ in $L\mathbb{Z}^3$, the union of the $B(y')$ is a fundamental domain of $\ell\mathbb{Z}^3$ in $\mathbb{Z}^3$.
For real $B$ the exponent of $U_B(y',z)=\exp\left(ie_k\cdot\eta_k\cdot(\tau B)(y',z)\right)$ (I.1.6) is purely imaginary, and therefore $\overline{U_B(y',z)}=U_B(y',z)^{-1}$.
On $\mathfrak{S}(L,\ell)$ the form is $(v',v)_L=L^3\sum_{y'}\overline{v'(y')}\cdot v(y')$, with $y'$ ranging over a fundamental domain of $\ell\mathbb{Z}^3$ in $L\mathbb{Z}^3$ (Section~I.1.6 with $a=L$).

\emph{The case $k\ge1$.}
Let $v\in\mathfrak{S}(1,\ell)$ and $v'\in\mathfrak{S}(L,\ell)$.
By Section~I.3.3,
\begin{align*}
\left(v',Qv\right)_L
&=L^3\sum_{y'}\overline{v'(y')}\cdot\left[L^{-3}\sum_{z\in B(y')}U_B(y',z)\,v(z)\right]\\
&=\sum_{y'}\sum_{z\in B(y')}\overline{U_B(y',z)^{-1}v'(y')}\cdot v(z)
=\left(Q^{T}(-B)v',v\right)_1,\\
\left(QQ^{T}(-B)v'\right)(y')
&=L^{-3}\sum_{z\in B(y')}U_B(y',z)\,U_B(y',z)^{-1}v'(y')=v'(y').
\end{align*}
Thus $Q^{*}=Q^{T}(-B)$ and $QQ^{*}=1$, and $P=Q^{*}Q$ is self-adjoint and idempotent on $\mathfrak{S}(1,\ell)$.
The coefficients are scalars;
so $P$ commutes with every constant matrix.

\emph{The case $k=0$.}
Let $f\in\mathfrak{S}^{\mathrm{ch}}(1,\ell)$ and $v'\in\mathfrak{S}(L,\ell)$.
By Section~I.3.3 and (I.3.1),
\begin{align*}
\left(QJ^{*}f\right)(y')&=L^{-3}\sum_{c\in B(y')}U_B(y',c)\cdot\frac{1}{48}\sum_{g\in O_{\mathrm{cub}}}f(c,g),\\
\left(JQ^{T}(-B)v'\right)(c,g)&=U_B(y',c)^{-1}v'(y')
\end{align*}
for $c\in B(y')$ and all $48$ labels $g$.
Hence
\begin{align*}
\left(v',QJ^{*}f\right)_L
&=\frac{1}{48}\sum_{y'}\sum_{c\in B(y')}\sum_g\overline{U_B(y',c)^{-1}v'(y')}\cdot f(c,g)
=\left(JQ^{T}(-B)v',f\right)^{\mathrm{ch}},\\
\left(QJ^{*}JQ^{T}(-B)v'\right)(y')
&=L^{-3}\cdot L^3\cdot\frac{48}{48}\cdot v'(y')=v'(y').
\end{align*}
So $\left(QJ^{*}\right)^{*}=JQ^{T}(-B)$ and $\left(QJ^{*}\right)\left(QJ^{*}\right)^{*}=1$, and $JPJ^{*}=\left(QJ^{*}\right)^{*}\left(QJ^{*}\right)$ is self-adjoint and idempotent on $\mathfrak{S}^{\mathrm{ch}}(1,\ell)$.
Since the coefficients are scalars, $JPJ^{*}$ commutes with every constant matrix.

\paragraph{8a.8. $JP_{k+1}J^{*}$, $k\ge1$.}
By Section~I.3.3, $P_{k+1}=Q_k^{T}(-B)PQ_k$;
therefore, by subsection~8a.6,
\[
JP_{k+1}J^{*}=\left(JQ_k^{T}(-B)\right)P\left(Q_kJ^{*}\right)=\left(Q_kJ^{*}\right)^{*}P\left(Q_kJ^{*}\right).
\]
This operator is self-adjoint, because $P^{*}=P$ (subsection~8a.7).
Using $\left(Q_kJ^{*}\right)\left(JQ_k^{T}(-B)\right)=1$ (the second identity of (21.5.30)), $P^2=P$ (subsection~8a.7) and $JP_kJ^{*}=\left(JQ_k^{T}(-B)\right)\left(Q_kJ^{*}\right)$, one finds
\begin{equation}
\begin{aligned}
\left(JP_{k+1}J^{*}\right)^2
&=\left(JQ_k^{T}(-B)\right)P\left[\left(Q_kJ^{*}\right)\left(JQ_k^{T}(-B)\right)\right]P\left(Q_kJ^{*}\right)\\
&=\left(JQ_k^{T}(-B)\right)P^2\left(Q_kJ^{*}\right)=JP_{k+1}J^{*},\\
\left(JP_kJ^{*}\right)\left(JP_{k+1}J^{*}\right)
&=\left(JQ_k^{T}(-B)\right)\left[\left(Q_kJ^{*}\right)\left(JQ_k^{T}(-B)\right)\right]P\left(Q_kJ^{*}\right)=JP_{k+1}J^{*},\\
\left(JP_{k+1}J^{*}\right)\left(JP_kJ^{*}\right)
&=\left(JQ_k^{T}(-B)\right)P\left[\left(Q_kJ^{*}\right)\left(JQ_k^{T}(-B)\right)\right]\left(Q_kJ^{*}\right)=JP_{k+1}J^{*}.
\end{aligned}
\tag{21.5.31}
\end{equation}
Finally, $JP_{k+1}J^{*}$ commutes with every constant matrix, by (21.5.6) and because $P$ does (subsection~8a.7).

\subsection*{9. Conditions on the constants}
\mbox{}

\medskip
\noindent
\begin{tabular}{p{0.2\textwidth}p{0.44\textwidth}p{0.28\textwidth}}
\hline
condition & where it enters & line of the list of Section~I.8.2 and of Proposition~14.2\,(a)\\
\hline
$\rho\ge0$
& (21.5.7), through Proposition~21.1\,(e): $\mathrm{Re}\,D_B=\frac{\rho}{2\eta_k}\cdot\nabla_B^{*}\nabla_B\ge0$
& (C0), the entry $\rho\in(0,1)$\\
$b>0$
& subsection~2.2: $b_k>0$, $b_L>0$; these enter (21.5.8), (21.5.10), (21.5.11), Section~6 and (21.5.22)
& (C0), the entry $b=1$\\
$L>1$
& subsection~2.2: $0<L^{-k}\le L^{-1}<1$, hence $b_k>0$; Section~8: $b_-\le b_k\le b_+$
& (C3), the entry $L$ (an odd integer; $L\ge5$ by Section~I.0.5); (C4), the entries $b_L$, $b_-$, $b_+$\\
$L$ odd
& only through the definition of the blocks (Section~I.1.1), on which the partitions in subsections~8a.6 and~8a.7 rest, and through the clauses recalled in subsection~2.3
& (C3), the entry $L$\\
$M$ a power of $L$, $L/2<M$
& only through the definitions: $\ell\in\mathrm{Cl}$ is a multiple of $L$ (Section~I.1.1; subsection~8a.7), and $U_B(y',z)$ is defined for $z\in B(y')$, because the sup distance of $y'$ and $z$ is $<L/2<M$ (Section~I.1.1, (I.1.6))
& (C5), the entry $M$ ($M$ a power of $L$ with $M\ge L^2$, Section~I.0.5)\\
$w_i>0$ and the closure relation (I.1.1)
& only through Proposition~21.1\,(e) ($\nabla_B^{*}\nabla_B\ge0$; the form (I.2.4) of $D_B$), derived for real $B$ in subsections~8a.3 and~8a.5
& (C0), the entry $w$\\
\hline
\end{tabular}

\medskip
The closure relation (I.1.1) and the positivity $w_i>0$ enter Proposition~21.1\,(e) through (I.2.4) and (I.O2), and through it every clause of Proposition~21.5.
If (I.1.1) did not hold, $D_B$ would contain the additional term $-\frac{1}{2\eta_k}\left(\gamma\cdot\left(\sum_iw_i\,\hat d_i(x)\right)\right)$ (subsection~8a.5).
This term is self-adjoint, so that it would be a part of $\mathrm{Re}\,D_B$;
it has no sign, since for a real vector $v$ the eigenvalues of $\gamma\cdot v$ are $|v|_2$ and $-|v|_2$, by $\left(\gamma\cdot v\right)^2=|v|_2^2$ and $\gamma_5\left(\gamma\cdot v\right)\gamma_5=-\gamma\cdot v$ (I.c1);
and it anticommutes with $\gamma_5$, so that the identity $D_B^{*}=\gamma_5D_B\gamma_5$ would fail.
The statements of this chapter hold without any condition on $\alpha$, $\kappa$, $\varepsilon$, $e_0$, $e_1$ or on the constants that follow them in the list of Section~I.8.2, and they require of $M$ only the properties in the table.
They hold for every side $\ell\in\mathrm{Cl}$, all numbers $e_k>0$ and $\mu>0$, and every real $B$.

\section*{Chapter 22. The one-step formula for the density with external connection, all levels: proofs of Propositions 22.1--22.4}

\subsection*{0. Conventions of this chapter}

\subsubsection*{0.1. Setting and notation}

Throughout the chapter the setting (S22) is the one of Section~G3.2.0.
Fields on $T^0_{k+1}$, $T^0_k$, $T^{\mathrm{f}}_0$, $L\mathbb{Z}^3/\ell\mathbb{Z}^3$, $T^{\mathrm{f}}_k$, $T^{\mathrm{f}}_{k+1}$, $\mathrm{Pz}_k$ and $\mathrm{Pz}_{k+1}$, together with their pairings and their scalings, are understood as in Section~G3.1.0 (fields), at the sides named there.

The three cases listed in the paragraph \emph{Instances} of Proposition~22.1 are called (i), (ii), (iii), in the order in which they are listed there:
(i) is the family $(\Psi',\bar{\Psi}')$, for every $k\ge0$;
(ii) consists of the three families $(\Psi,\bar{\Psi})$, $(W,\bar{W})$, $(W',\bar{W}')$ for $k\ge1$;
(iii) consists of the same three families at $k=0$.

In this chapter $\mathrm{Jw}$, $\mathrm{Jw}'$, $\mathrm{Jw}_i$, $\mathrm{Jw}_i^{\mathrm{f}}$ denote sets of generators of a family,
$\mathrm{Ew}$, $\mathrm{Ew}'$ even elements of a Grassmann algebra,
$e_I$, $e'_I$ monomials, whose index $I$ is a set,
and $\mathrm{hm}$, $\mathrm{hm}_i$ homomorphisms.
The letters $\mathrm{kn}$, $\mathrm{kn}_i$, $\mathrm{ng}$, $\mathrm{nq}$, $\mathrm{Nz}_0$ stand for numbers,
the letters $\mathrm{Xw}$, $\mathrm{Eq}$, $\mathrm{Eq}_1$, $\mathrm{Rs}$, $\mathrm{Fp}$ for fields, exponents and integrands,
$\mathrm{Lm}$ for a linear map, $\mathrm{mq}$ for a measure and $\mathrm{Og}$ for an operator.
Without decoration, $J$ always denotes the cluster map and $w$ the quadruple of the standard weights.

\subsubsection*{0.2. Grassmann algebras}

Let $\mathrm{Gr}$ be a finite dimensional Grassmann algebra over $\mathbb{C}$ with the independent generators $\theta_1,\dots,\theta_{\mathrm{ng}}$.
For a subset $I$ of $\{1,\dots,\mathrm{ng}\}$ let $e_I$ denote the product of the $\theta_i$, $i\in I$, taken in increasing order, with $e_{\emptyset}=1$.
The $e_I$ form a basis of $\mathrm{Gr}$, and
\[
\theta_i\theta_j=-\theta_j\theta_i,\qquad \theta_i^2=0 .
\]
The algebra $\mathrm{Gr}$ is the direct sum of its even part, which is the span of the $e_I$ with $\#|I|$ even, and of its odd part.

\begin{named}{(Gr1)}
For subsets $I$, $I'$ of $\{1,\dots,\mathrm{ng}\}$ one has
\[
e_Ie_{I'}=(-1)^{\#|I|\,\#|I'|}\,e_{I'}e_I .
\]
Consequently even elements commute with every element, two odd elements anticommute, and $\theta'^2=0$ for every odd element $\theta'$, because $\theta'\theta'=-\theta'\theta'$.
If $I$ and $I'$ have a common index, then $e_Ie_{I'}=0$;
otherwise $e_Ie_{I'}=\mathrm{sn}(I,I')\,e_{I\cup I'}$, where $\mathrm{sn}(I,I')$ is the sign of the permutation that sorts the sequence $(I,I')$.
\end{named}

\begin{named}{(Gr2)}[homomorphisms]
Let $\theta'_1,\dots,\theta'_{\mathrm{ng}}$ be odd elements of a Grassmann algebra $\mathrm{Gr}'$.
There is exactly one homomorphism $\mathrm{hm}\colon\mathrm{Gr}\to\mathrm{Gr}'$ of algebras with unit such that $\mathrm{hm}(\theta_i)=\theta'_i$ for every $i$.
\end{named}

\begin{proof}
Let $e'_I$ be the product of the $\theta'_i$, $i\in I$, in increasing order;
set $\mathrm{hm}(e_I):=e'_I$ and extend linearly.
By (Gr1), applied in $\mathrm{Gr}'$, one has $\theta'_i\theta'_j=-\theta'_j\theta'_i$ and $\theta'^2_i=0$.
Therefore $e'_Ie'_{I'}=0$ if $I$ and $I'$ have a common index, and $e'_Ie'_{I'}=\mathrm{sn}(I,I')\,e'_{I\cup I'}$ otherwise.
This is the multiplication table of the $e_I$, and so $\mathrm{hm}$ is multiplicative.
It is unique, since the $\theta_i$ generate $\mathrm{Gr}$.
\end{proof}

The homomorphism $\mathrm{hm}$ maps even elements to even elements and odd elements to odd elements, because a product of odd elements has the parity of the number of its factors.
In these terms the symbol $G[\bar{W}\to\bar{h},\,W\to\mathrm{hf}]$ of Proposition~22.1 is $\mathrm{hm}(G)$, where $\mathrm{hm}$ is the homomorphism of $\mathrm{Gr}$ into itself that maps $\bar{W}_\beta(t)$ to $\bar{h}_\beta(t)$, maps $W_\beta(t)$ to $\mathrm{hf}_\beta(t)$, and satisfies $\mathrm{hm}(\theta_i)=\theta_i$ for the remaining generators $\theta_i$.
The symbol is thus defined for all maps $\bar{h}$, $\mathrm{hf}$ from the index set to the odd part of $\mathrm{Gr}$.

\begin{named}{(Gr3)}[exponential]
Every even element has the form $\mathrm{Ew}=\mathrm{Ew}_0\cdot1+\mathrm{Ew}_+$, where $\mathrm{Ew}_0$ is a number and $\mathrm{Ew}_+$ is a linear combination of the $e_I$ with $\#|I|\ge2$ even.
Then $\mathrm{Ew}_+^j=0$ for $2j>\mathrm{ng}$, and
\[
\exp(\mathrm{Ew}):=\sum_{j\ge0}\frac{\mathrm{Ew}^j}{j!}=\exp(\mathrm{Ew}_0)\sum_{2j\le\mathrm{ng}}\frac{\mathrm{Ew}_+^j}{j!} .
\]
The coefficients of $\exp(\mathrm{Ew})$ are $\exp(\mathrm{Ew}_0)$ times polynomials in the coefficients of $\mathrm{Ew}_+$;
hence $\exp$ is a continuous map of the even part into itself.
For even elements $\mathrm{Ew}$, $\mathrm{Ew}'$, which commute,
\[
\exp(\mathrm{Ew}+\mathrm{Ew}')=\exp(\mathrm{Ew})\cdot\exp(\mathrm{Ew}') .
\]
For $\mathrm{hm}$ as in (Gr2),
\[
\mathrm{hm}(\exp\mathrm{Ew})=\exp(\mathrm{hm}\,\mathrm{Ew}) ;
\]
indeed, $\mathrm{hm}$ is linear and multiplicative, $\mathrm{hm}(\mathrm{Ew}_+)$ is a combination of products of at least two odd elements and therefore has no term of degree $0$, and $\mathrm{hm}(\mathrm{Ew}_+)^j=\mathrm{hm}(\mathrm{Ew}_+^j)=0$ for $2j>\mathrm{ng}$.
\end{named}

\begin{named}{(Gr4)}[expansion relative to a family; the integral in an algebra with further generators]
Let $(W,\bar{W})$ be a family of generators of $\mathrm{Gr}$ indexed by a finite set $\mathrm{Mo}$ (see 0.3), let $\mathrm{nq}:=\#|\mathrm{Mo}|$, and let $\mathrm{Gr}^{\mathrm{o}}$ be the subalgebra generated by the generators of $\mathrm{Gr}$ that do not belong to the family.
For a set $\mathrm{Jw}$ of generators of the family let $W(\mathrm{Jw})$ be their product in increasing order.
Every $F\in\mathrm{Gr}$ has exactly one representation
\[
F=\sum_{\mathrm{Jw}}F_{\mathrm{Jw}}\,W(\mathrm{Jw}),\qquad F_{\mathrm{Jw}}\in\mathrm{Gr}^{\mathrm{o}},
\]
because the products $e_{I'}W(\mathrm{Jw})$, with $I'$ a set of indices of generators outside the family, are, up to signs, the basis elements $e_I$ of $\mathrm{Gr}$, each of them exactly once.
The number $\#|\mathrm{Jw}|$ is called the degree in the family.
The element $\mathrm{Top}=\prod_{\mathrm{mo}}W(\mathrm{mo})\bar{W}(\mathrm{mo})$ equals $\pm W(\mathrm{Jw})$ for the full set $\mathrm{Jw}$ of generators of the family.
The coefficient of $\mathrm{Top}$ is the element $[F]_{\mathrm{Top}}$ of $\mathrm{Gr}^{\mathrm{o}}$ for which
\[
F=[F]_{\mathrm{Top}}\cdot\mathrm{Top}+(\text{terms of degree}<2\,\mathrm{nq}\text{ in the family}).
\]
In $\mathrm{Gr}$ the definition (I.3.5), together with the sentence on algebras with further generators that follows it in Section~I.3.6, reads
\[
\int\mathrm{D}W\,F:=\left(\prod_{\mathrm{mo}}\nu(\mathrm{mo})^{-1}\right)[F]_{\mathrm{Top}},
\qquad
\int\mathrm{d}\mu_I(W)\,F:=\int\mathrm{D}W\,\exp\left(-\langle\bar{W},W\rangle_\nu\right)F ;
\]
here $F\mapsto\int\mathrm{D}W\,F$ is a linear map $\mathrm{Gr}\to\mathrm{Gr}^{\mathrm{o}}$.
\end{named}

\begin{named}{(Gr5)}
For $F^{\mathrm{o}}\in\mathrm{Gr}^{\mathrm{o}}$ and $F\in\mathrm{Gr}$,
\[
\int\mathrm{D}W\,(F^{\mathrm{o}}F)=F^{\mathrm{o}}\int\mathrm{D}W\,F,
\qquad
\int\mathrm{D}W\,(FF^{\mathrm{o}})=\left(\int\mathrm{D}W\,F\right)F^{\mathrm{o}} .
\]
\end{named}

\begin{proof}
One has $F^{\mathrm{o}}F=\sum_{\mathrm{Jw}}(F^{\mathrm{o}}F_{\mathrm{Jw}})\,W(\mathrm{Jw})$ with coefficients $F^{\mathrm{o}}F_{\mathrm{Jw}}$ in $\mathrm{Gr}^{\mathrm{o}}$, which gives the first relation.
For the second relation let first $F^{\mathrm{o}}$ be even or odd.
Then $W(\mathrm{Jw})F^{\mathrm{o}}=\pm F^{\mathrm{o}}W(\mathrm{Jw})$, with the sign $+$ for the full set $\mathrm{Jw}$, which has an even number of elements;
hence $FF^{\mathrm{o}}=[F]_{\mathrm{Top}}F^{\mathrm{o}}\cdot\mathrm{Top}+(\text{terms of degree}<2\,\mathrm{nq})$.
A general $F^{\mathrm{o}}$ is the sum of its even and its odd part, and both sides of the relation are linear in $F^{\mathrm{o}}$.
\end{proof}

\subsubsection*{0.3. Modes, matrices, transposes, fields}

Let $\mathrm{Mo}:=\{(t,\beta)\}$ be the index set of Proposition~22.1.
Its elements $\mathrm{mo}=(t,\beta)$ are called modes, and $\nu(\mathrm{mo}):=\nu(t)$.
A linear map $O$ of the functions on $\mathrm{Mo}$ has the matrix $O[\mathrm{mo},\mathrm{mo}']$ defined by
\[
(Of)(\mathrm{mo})=\sum_{\mathrm{mo}'}O[\mathrm{mo},\mathrm{mo}']\,f(\mathrm{mo}') ;
\]
thus $O[\mathrm{mo},\mathrm{mo}']=\nu(\mathrm{mo}')\,O(\mathrm{mo},\mathrm{mo}')$ for the kernel $O(\mathrm{mo},\mathrm{mo}')$ of Section~I.3.6, and $\det O$ is the determinant of this matrix.
The transpose of $O$ for the pairing $\langle f'',f'\rangle_\nu=\sum_{\mathrm{mo}}\nu(\mathrm{mo})f''(\mathrm{mo})f'(\mathrm{mo})$ has the matrix
\[
O^{T}[\mathrm{mo}',\mathrm{mo}]=\frac{\nu(\mathrm{mo})\,O[\mathrm{mo},\mathrm{mo}']}{\nu(\mathrm{mo}')} .
\]
One has $(O_1O_2)^{T}=O_2^{T}O_1^{T}$, and for invertible $O$ also $O^{T}(O^{-1})^{T}=(O^{-1}O)^{T}=1$.

For maps $f''$, $f'$ from $\mathrm{Mo}$ to the odd part of $\mathrm{Gr}$ put
\[
\langle f'',f'\rangle_\nu:=\sum_{\mathrm{mo}}\nu(\mathrm{mo})\,f''(\mathrm{mo})\,f'(\mathrm{mo}),
\qquad
(Of')(\mathrm{mo}):=\sum_{\mathrm{mo}'}O[\mathrm{mo},\mathrm{mo}']\,f'(\mathrm{mo}') ;
\]
in the pairing the factors stand in the order written, and the pairing is an even element.
Since no factor is exchanged,
\begin{equation}
\langle f'',Of'\rangle_\nu=\sum_{\mathrm{mo},\mathrm{mo}'}\nu(\mathrm{mo})\,O[\mathrm{mo},\mathrm{mo}']\,f''(\mathrm{mo})\,f'(\mathrm{mo}')=\langle O^{T}f'',f'\rangle_\nu .
\tag{P.0}
\end{equation}
For a homomorphism $\mathrm{hm}$ as in (Gr2),
\[
\mathrm{hm}\left(\langle f'',Of'\rangle_\nu\right)=\langle\mathrm{hm}\,f'',O\,\mathrm{hm}\,f'\rangle_\nu,
\qquad\text{where }(\mathrm{hm}\,f)(\mathrm{mo}):=\mathrm{hm}\left(f(\mathrm{mo})\right).
\]

\subsubsection*{0.4. The objects of Section~I.5.7, written with homomorphisms}

Let $B$ be a complex bond function on $T_k$ such that $S_k(B)$ and $S_k(-B)$ exist if $k\ge1$.
Let $\mathrm{ev}_B\colon\mathrm{Gr}[\psi,\chi]\to\mathrm{Gr}[\Psi]$ be the homomorphism with
\[
\psi\mapsto H_k(B)\Psi,\qquad
\bar{\psi}\mapsto\bar{H}_k(B)\bar{\Psi},\qquad
\chi\mapsto\delta_{\alpha,B}H_k(B)\Psi,\qquad
\bar{\chi}\mapsto\delta_{\alpha,-B}\bar{H}_k(B)\bar{\Psi} ;
\]
for $k=0$, since $H_0(B)=\bar{H}_0(B)=1$, these values are
\[
\psi\mapsto\Psi,\qquad
\bar{\psi}\mapsto\bar{\Psi},\qquad
\chi\mapsto\delta_{\alpha,B}\Psi,\qquad
\bar{\chi}\mapsto\delta_{\alpha,-B}\bar{\Psi} .
\]
The images are odd and have the antiperiodicity (I.AP) of the generators, so that (Gr2) defines $\mathrm{ev}_B$.
For $(\mathscr{A},\mathscr{C})\in\tilde{R}_k\times\mathfrak{X}_k$ and $\mathscr{T}=\mathscr{A}+\mathscr{C}$ the displays of Section~I.5.7 read
\begin{align*}
I^{x}_{k}[m,\mathscr{E}](\mathscr{A},\mathscr{C};\Psi)
&=\mathrm{ev}_{\mathscr{T}}\left(-m\langle\bar{\psi},\psi\rangle+\sum_X\mathscr{E}(X;\mathscr{A},\mathscr{C})\right),\\
\mathrm{Act}^{x}_{k}[m,\mathscr{E}](\mathscr{A},\mathscr{C};\Psi)
&=-\langle\bar{\Psi},D_k(\mathscr{T})\Psi\rangle+I^{x}_{k}[m,\mathscr{E}](\mathscr{A},\mathscr{C};\Psi),\\
\mathscr{F}^{x}_{k}[\varepsilon^{\mathrm{v}},m,\mathscr{E}](A,\mathscr{C};\Psi)
&=\det D_k(0)^{-1}\\
&\quad\cdot\exp\left\{\mathrm{Act}^{x}_{k}[m,\mathscr{E}](H_{k}A,\mathscr{C};\Psi)-\varepsilon^{\mathrm{v}}n_k+\sum_X\mathfrak{P}_k(X,\mathscr{C})\right\}
\end{align*}
for $A\in\mathrm{Dom}_k$.
Here, for every $k\ge0$,
\[
\langle\bar{\psi},\psi\rangle=\int\mathrm{d}x\,\bar{\psi}(x)\cdot\psi(x)=\eta_k^3\,\frac{1}{48}\sum_c\sum_{g\in O_{\mathrm{cub}}}\sum_\beta\bar{\psi}_\beta(c,g)\,\psi_\beta(c,g),
\]
where $c$ ranges over the photon sites of a fundamental domain and all $192$ components over each such site $c$ occur;
and, by Section~I.3.6, with $y$, respectively $x$, ranging over a fundamental domain,
\begin{gather*}
\langle\bar{\Psi},O\Psi\rangle=\langle\bar{\Psi},O\Psi\rangle_1=\sum_y\bar{\Psi}(y)\cdot(O\Psi)(y)\qquad(k\ge1),\\
\langle\bar{\Psi},O\Psi\rangle=\langle\bar{\Psi},O\Psi\rangle^{\mathrm{ch}}=\frac{1}{48}\sum_x\bar{\Psi}(x)\cdot(O\Psi)(x)\qquad(k=0).
\end{gather*}
For a pair $(\Phi,\bar{\Phi})$ of fields on the unit index set of level $k$, Section~G3.2.0 (paragraph \emph{Pairs of fields in the last slot}) sets
\[
I^{x}_{k}[m,\mathscr{E}](\mathscr{A},\mathscr{C};\Phi)=I^{x}_{k}[m,\mathscr{E}](\mathscr{A},\mathscr{C};\Psi)\,[\bar{\Psi}\to\bar{\Phi},\,\Psi\to\Phi],
\]
and defines $\mathrm{Act}^{x}_{k}[m,\mathscr{E}](\mathscr{A},\mathscr{C};\Phi)$ in the same way.
Since the operators are linear, $\mathrm{ev}_{\mathscr{T}}$ followed by $[\bar{\Psi}\to\bar{\Phi},\,\Psi\to\Phi]$ is the homomorphism of $\mathrm{Gr}[\psi,\chi]$ with
\[
\psi\mapsto H_k(\mathscr{T})\Phi,\qquad
\bar{\psi}\mapsto\bar{H}_k(\mathscr{T})\bar{\Phi},\qquad
\chi\mapsto\delta_{\alpha,\mathscr{T}}H_k(\mathscr{T})\Phi,\qquad
\bar{\chi}\mapsto\delta_{\alpha,-\mathscr{T}}\bar{H}_k(\mathscr{T})\bar{\Phi} .
\]

All these elements are even:
$\langle\bar{\psi},\psi\rangle$ and $\langle\bar{\Psi},O\Psi\rangle$ are sums of products of two odd elements, the $\mathscr{E}(X;\mathscr{A},\mathscr{C})$ are even (Section~I.5.4), and the homomorphisms of (Gr2) preserve the parity.
Consequently all exponentials below commute with each other, and (Gr3) applies to sums of exponents.

\subsection*{1. Proposition 22.1 (Berezin rules on an index set with weights)}

In the clauses (G0)$^{\mathrm{w}}$--(G3)$^{\mathrm{w}}$ the integrated family is $(W,\bar{W})$, and $\mathrm{Gr}^{\mathrm{o}}$, $\mathrm{nq}$ and $[\,\cdot\,]_{\mathrm{Top}}$ are formed for this family as in (Gr4).

\begin{named}{(L)}
For every linear map $U$ of the functions on $\mathrm{Mo}$,
\[
\prod_{\mathrm{mo}}(UW)(\mathrm{mo})\,\bar{W}(\mathrm{mo})=\det U\cdot\mathrm{Top} ;
\]
the factors $(UW)(\mathrm{mo})\,\bar{W}(\mathrm{mo})$ are even, so that their order does not matter.
\end{named}

\begin{proof}
Number the modes as $\mathrm{mo}_1,\dots,\mathrm{mo}_{\mathrm{nq}}$.
For odd elements $\theta'_1,\dots,\theta'_{\mathrm{nq}}$, $\theta''_1,\dots,\theta''_{\mathrm{nq}}$ of $\mathrm{Gr}$ one has
\[
\theta'_1\theta''_1\theta'_2\theta''_2\cdots\theta'_{\mathrm{nq}}\theta''_{\mathrm{nq}}
=(-1)^{\frac{\mathrm{nq}(\mathrm{nq}-1)}{2}}\,\theta'_1\theta'_2\cdots\theta'_{\mathrm{nq}}\,\theta''_1\theta''_2\cdots\theta''_{\mathrm{nq}} ;
\]
indeed, moving $\theta'_j$ to the left across $\theta''_1,\dots,\theta''_{j-1}$ produces $j-1$ changes of sign by (Gr1), and $\sum_{j=1}^{\mathrm{nq}}(j-1)=\frac{\mathrm{nq}(\mathrm{nq}-1)}{2}$.
The sign does not depend on the elements.
Apply this identity with $\theta''_j=\bar{W}(\mathrm{mo}_j)$, once with $\theta'_j=(UW)(\mathrm{mo}_j)$ and once with $\theta'_j=W(\mathrm{mo}_j)$.
It then remains to show that
\[
(UW)(\mathrm{mo}_1)\cdots(UW)(\mathrm{mo}_{\mathrm{nq}})=\det U\cdot W(\mathrm{mo}_1)\cdots W(\mathrm{mo}_{\mathrm{nq}}) .
\]
The left side equals
\[
\sum_{i_1,\dots,i_{\mathrm{nq}}}U[\mathrm{mo}_1,\mathrm{mo}_{i_1}]\cdots U[\mathrm{mo}_{\mathrm{nq}},\mathrm{mo}_{i_{\mathrm{nq}}}]\,W(\mathrm{mo}_{i_1})\cdots W(\mathrm{mo}_{i_{\mathrm{nq}}}) .
\]
By (Gr1) the product of the $W$ vanishes if two of the indices $i_j$ coincide;
otherwise $j\mapsto i_j$ is a permutation, and the product is its sign times $W(\mathrm{mo}_1)\cdots W(\mathrm{mo}_{\mathrm{nq}})$.
By the Leibniz formula for the determinant, the sum of the signed products of the entries is $\det U$.
\end{proof}

\begin{proof}[Proof of Proposition~22.1\,(G0)$^{\mathrm{w}}$]
Put $\mathrm{Xw}:=OW$.
The elements $\bar{W}(\mathrm{mo})\mathrm{Xw}(\mathrm{mo})$ are even, hence commute with each other, and their squares vanish:
\[
\bar{W}(\mathrm{mo})\mathrm{Xw}(\mathrm{mo})\bar{W}(\mathrm{mo})\mathrm{Xw}(\mathrm{mo})=-\bar{W}(\mathrm{mo})^2\,\mathrm{Xw}(\mathrm{mo})^2=0 .
\]
By (Gr3) it follows that
\begin{align*}
\exp\left(-\langle\bar{W},OW\rangle_\nu\right)
&=\prod_{\mathrm{mo}}\exp\left(-\nu(\mathrm{mo})\bar{W}(\mathrm{mo})\mathrm{Xw}(\mathrm{mo})\right)
=\prod_{\mathrm{mo}}\left(1-\nu(\mathrm{mo})\bar{W}(\mathrm{mo})\mathrm{Xw}(\mathrm{mo})\right)\\
&=\prod_{\mathrm{mo}}\left(1+\nu(\mathrm{mo})\mathrm{Xw}(\mathrm{mo})\bar{W}(\mathrm{mo})\right).
\end{align*}
Each of the $\mathrm{nq}$ factors is the sum of a term of degree $0$ and of a term of degree $2$ in the family.
The part of degree $2\,\mathrm{nq}$ of the product is therefore $\prod_{\mathrm{mo}}\nu(\mathrm{mo})\mathrm{Xw}(\mathrm{mo})\bar{W}(\mathrm{mo})$, and by (L) this equals $\left(\prod_{\mathrm{mo}}\nu(\mathrm{mo})\right)\det O\cdot\mathrm{Top}$.
Hence
\[
\left[\exp\left(-\langle\bar{W},OW\rangle_\nu\right)\right]_{\mathrm{Top}}=\left(\prod_{\mathrm{mo}}\nu(\mathrm{mo})\right)\det O,
\qquad
\int\mathrm{D}W\,\exp\left(-\langle\bar{W},OW\rangle_\nu\right)=\det O .
\]
No property of $O$ has entered.
For $O=1$ the value is $1$.

It remains to see that this single condition fixes the integral.
Let $\mathrm{kn}$ be a complex number and put $\int'F:=\mathrm{kn}\cdot[F]_{\mathrm{Top}}$.
The computation at $O=1$ gives $\int'\exp\left(-\langle\bar{W},W\rangle_\nu\right)=\mathrm{kn}\cdot\prod_{\mathrm{mo}}\nu(\mathrm{mo})$.
This number is $1$ if and only if $\mathrm{kn}=\prod_{\mathrm{mo}}\nu(\mathrm{mo})^{-1}$, that is, if and only if $\int'=\int\mathrm{D}W$.
\end{proof}

\begin{proof}[Proof of Proposition~22.1\,(G1)$^{\mathrm{w}}$]
Let $\mathrm{hm}:=[\bar{W}\to\bar{W}+\bar{g},\,W\to W+\mathrm{gf}]$;
this homomorphism fixes $\mathrm{Gr}^{\mathrm{o}}$.
Write $G=\sum_{\mathrm{Jw}}G_{\mathrm{Jw}}\,W(\mathrm{Jw})$ as in (Gr4).
The element $\mathrm{hm}\left(W(\mathrm{Jw})\right)$ is the product, in increasing order, of the elements $\mathrm{hm}(\theta_i)$, where $\theta_i$ ranges over the generators in $\mathrm{Jw}$;
each $\mathrm{hm}(\theta_i)$ is the sum of $\theta_i$ and of a value of $\mathrm{gf}$ or of $\bar{g}$, and such a value is an odd element of $\mathrm{Gr}^{\mathrm{o}}$.
Expand this product.
The term that contains no value of $\mathrm{gf}$ or $\bar{g}$ is $W(\mathrm{Jw})$.
In each of the other terms move the values of $\mathrm{gf}$, $\bar{g}$ to the left, which by (Gr1) costs a sign;
the term becomes an element of $\mathrm{Gr}^{\mathrm{o}}$ times some $W(\mathrm{Jw}')$ with $\#|\mathrm{Jw}'|<\#|\mathrm{Jw}|\le2\,\mathrm{nq}$.
Consequently
\[
\mathrm{hm}(G)=G+(\text{terms of degree}<2\,\mathrm{nq}\text{ in the family}),
\]
so that $[\mathrm{hm}(G)]_{\mathrm{Top}}=[G]_{\mathrm{Top}}$ and $\int\mathrm{D}W\,\mathrm{hm}(G)=\int\mathrm{D}W\,G$.
\end{proof}

\begin{proof}[Proof of Proposition~22.1\,(G2)$^{\mathrm{w}}$]
Put $\mathrm{gf}:=O^{-1}\mathrm{Sr}$ and $\bar{g}:=(O^{-1})^{T}\,\overline{\mathrm{Sr}}$;
these are fields.
Let
\[
\mathrm{Ew}:=-\langle\bar{W},OW\rangle_\nu+\langle\bar{W},\mathrm{Sr}\rangle_\nu+\langle\overline{\mathrm{Sr}},W\rangle_\nu ,
\]
which is an even element, and let $\mathrm{hm}$ be the homomorphism of the proof of (G1)$^{\mathrm{w}}$.
The fields $\mathrm{Sr}$, $\overline{\mathrm{Sr}}$ take their values in $\mathrm{Gr}^{\mathrm{o}}$ and are therefore fixed by $\mathrm{hm}$.
By 0.3,
\begin{multline*}
\mathrm{hm}(\mathrm{Ew})=-\langle\bar{W},OW\rangle_\nu-\langle\bar{W},O\,\mathrm{gf}\rangle_\nu-\langle\bar{g},OW\rangle_\nu-\langle\bar{g},O\,\mathrm{gf}\rangle_\nu\\
+\langle\bar{W},\mathrm{Sr}\rangle_\nu+\langle\bar{g},\mathrm{Sr}\rangle_\nu+\langle\overline{\mathrm{Sr}},W\rangle_\nu+\langle\overline{\mathrm{Sr}},\mathrm{gf}\rangle_\nu .
\end{multline*}
Since $O\,\mathrm{gf}=\mathrm{Sr}$, the second term cancels against the fifth, and the fourth against the sixth.
By (P.0), $\langle\bar{g},OW\rangle_\nu=\langle O^{T}\bar{g},W\rangle_\nu=\langle\overline{\mathrm{Sr}},W\rangle_\nu$, because $O^{T}(O^{-1})^{T}=1$;
so the third term cancels against the seventh.
Thus
\[
\mathrm{hm}(\mathrm{Ew})=-\langle\bar{W},OW\rangle_\nu+\langle\overline{\mathrm{Sr}},O^{-1}\mathrm{Sr}\rangle_\nu ,
\]
and the last term is an even element of $\mathrm{Gr}^{\mathrm{o}}$.
By (G1)$^{\mathrm{w}}$, (Gr3), (Gr5) and (G0)$^{\mathrm{w}}$,
\begin{align*}
\int\mathrm{D}W\,\exp(\mathrm{Ew})
&=\int\mathrm{D}W\,\exp\left(\mathrm{hm}(\mathrm{Ew})\right)
=\exp\langle\overline{\mathrm{Sr}},O^{-1}\mathrm{Sr}\rangle_\nu\cdot\int\mathrm{D}W\,\exp\left(-\langle\bar{W},OW\rangle_\nu\right)\\
&=\det O\cdot\exp\langle\overline{\mathrm{Sr}},O^{-1}\mathrm{Sr}\rangle_\nu .
\end{align*}
\end{proof}

\begin{proof}[Proof of Proposition~22.1\,(G3)$^{\mathrm{w}}$]
The homomorphism $\mathrm{hm}:=[\bar{W}\to u^{-1}\bar{W},\,W\to uW]$ fixes $\mathrm{Gr}^{\mathrm{o}}$ and multiplies each generator of the family by a number.
Hence $\mathrm{hm}\left(W(\mathrm{Jw})\right)$ is a multiple of $W(\mathrm{Jw})$, and
\[
\mathrm{hm}(\mathrm{Top})=\prod_{(t,\beta)}u(t)W_\beta(t)\,u(t)^{-1}\bar{W}_\beta(t)=\mathrm{Top} .
\]
Therefore $[\mathrm{hm}(G)]_{\mathrm{Top}}=[G]_{\mathrm{Top}}$, and the two integrals of (G3)$^{\mathrm{w}}$ coincide.
\end{proof}

\begin{named}{(L$'$)}
For every linear map $U$ of the functions on $\mathrm{Mo}$ and every $F\in\mathrm{Gr}$,
\[
\int\mathrm{D}W\,F[W\to UW]=\det U\cdot\int\mathrm{D}W\,F ;
\]
here $\bar{W}$ and all the other generators are kept fixed.
\end{named}

\begin{proof}
The homomorphism $\mathrm{hm}:=[W\to UW]$ fixes $\mathrm{Gr}^{\mathrm{o}}$ and maps each generator of the family to a linear combination of generators of the family;
hence it maps $W(\mathrm{Jw})$ to a combination of the $W(\mathrm{Jw}')$ with $\#|\mathrm{Jw}'|=\#|\mathrm{Jw}|$.
By (L), $\mathrm{hm}(\mathrm{Top})=\det U\cdot\mathrm{Top}$.
Therefore
\[
\mathrm{hm}(F)=\det U\cdot[F]_{\mathrm{Top}}\cdot\mathrm{Top}+(\text{terms of degree}<2\,\mathrm{nq}) .
\]
\end{proof}

\begin{proof}[Proof of Proposition~22.1\,(G4)$^{\mathrm{w}}$]
Let $\mathrm{Gr}^{(1)}$ be the subalgebra generated by all generators other than $W$, $\bar{W}$, and $\mathrm{Gr}^{\mathrm{oo}}$ the subalgebra generated by all generators other than $W$, $\bar{W}$, $W'$, $\bar{W}'$.
Put $\mathrm{Top}':=\prod_{\mathrm{mo}}W'(\mathrm{mo})\bar{W}'(\mathrm{mo})$.
For $F\in\mathrm{Gr}^{(1)}$ the coefficients of the expansion (Gr4) relative to the family $(W',\bar{W}')$ contain no $W$, $\bar{W}$, and so they lie in $\mathrm{Gr}^{\mathrm{oo}}$.

\emph{Renaming.}
The homomorphism $\mathrm{hm}_1:=[\bar{W}'\to\bar{W},\,W'\to W]$ fixes $\mathrm{Gr}^{\mathrm{oo}}$, maps a monomial of degree $j$ in the family $(W',\bar{W}')$ to a monomial of degree $j$ in the family $(W,\bar{W})$, and maps $\mathrm{Top}'$ to $\mathrm{Top}$.
Let $F\in\mathrm{Gr}^{(1)}$.
From $F=[F]_{\mathrm{Top}'}\cdot\mathrm{Top}'+(\text{terms of degree}<2\,\mathrm{nq}\text{ in }(W',\bar{W}'))$ one obtains
\[
\mathrm{hm}_1(F)=[F]_{\mathrm{Top}'}\cdot\mathrm{Top}+(\text{terms of degree}<2\,\mathrm{nq}\text{ in }(W,\bar{W})) .
\]
The two families have the same weights.
Hence
\begin{equation}
\int\mathrm{D}W'\,F=\int\mathrm{D}W\,\mathrm{hm}_1(F)\qquad(F\in\mathrm{Gr}^{(1)}).
\tag{P.1}
\end{equation}

\emph{Change of variables.}
Let $F:=\exp\left(-\langle\bar{W}',OW'\rangle_\nu\right)G$.
It lies in $\mathrm{Gr}^{(1)}$, and $\mathrm{hm}_1(F)=\exp\left(-\langle\bar{W},OW\rangle_\nu\right)\mathrm{hm}_1(G)$.
Let $\mathrm{hm}_2:=[W\to O^{-1}W]$.
Then $\mathrm{hm}_2\left(\langle\bar{W},OW\rangle_\nu\right)=\langle\bar{W},OO^{-1}W\rangle_\nu=\langle\bar{W},W\rangle_\nu$.
The homomorphisms $\mathrm{hm}_2\circ\mathrm{hm}_1$ and $[\bar{W}'\to\bar{W},\,W'\to O^{-1}W]$ agree on the generators of $\mathrm{Gr}^{(1)}$:
both fix the generators of $\mathrm{Gr}^{\mathrm{oo}}$, map $\bar{W}'$ to $\bar{W}$ and map $W'$ to $O^{-1}W$.
So they agree on $G$, and
\[
\mathrm{hm}_2\left(\mathrm{hm}_1(F)\right)=\exp\left(-\langle\bar{W},W\rangle_\nu\right)G[\bar{W}'\to\bar{W},\,W'\to O^{-1}W] .
\]
By (L$'$) with $U=O^{-1}$, and by (P.1),
\begin{align*}
\int\mathrm{d}\mu_I(W)\,G[\bar{W}'\to\bar{W},\,W'\to O^{-1}W]
&=\int\mathrm{D}W\,\mathrm{hm}_2\left(\mathrm{hm}_1(F)\right)
=\det(O^{-1})\cdot\int\mathrm{D}W\,\mathrm{hm}_1(F)\\
&=(\det O)^{-1}\int\mathrm{D}W'\,\exp\left(-\langle\bar{W}',OW'\rangle_\nu\right)G .
\end{align*}
\end{proof}

\begin{proof}[Proof of Proposition~22.1\,(G5)$^{\mathrm{w}}$]
\emph{Two families.}
Let $(W^1,\bar{W}^1)$ on $(\mathrm{Mo}_1,\nu_1)$ and $(W^2,\bar{W}^2)$ on $(\mathrm{Mo}_2,\nu_2)$ be two different families of generators of $\mathrm{Gr}$;
the index sets and the weights may differ.
Let $\mathrm{Gr}^{\mathrm{oo}}$ be the subalgebra generated by the remaining generators, and put
\[
\mathrm{kn}_i:=\prod_{\mathrm{mo}\in\mathrm{Mo}_i}\nu_i(\mathrm{mo})^{-1},
\qquad
\mathrm{Top}_i:=\prod_{\mathrm{mo}\in\mathrm{Mo}_i}W^i(\mathrm{mo})\bar{W}^i(\mathrm{mo}) .
\]
Let $\mathrm{Jw}_i^{\mathrm{f}}$ be the full set of the generators of the $i$-th family, and let $\mathrm{sn}_i\in\{1,-1\}$ be the sign with $W^i(\mathrm{Jw}_i^{\mathrm{f}})=\mathrm{sn}_i\,\mathrm{Top}_i$.
Every $F\in\mathrm{Gr}$ has exactly one representation
\[
F=\sum_{\mathrm{Jw}_1,\mathrm{Jw}_2}F_{\mathrm{Jw}_1\mathrm{Jw}_2}\,W^1(\mathrm{Jw}_1)\,W^2(\mathrm{Jw}_2),
\qquad F_{\mathrm{Jw}_1\mathrm{Jw}_2}\in\mathrm{Gr}^{\mathrm{oo}},
\]
where $\mathrm{Jw}_i$ ranges over the sets of generators of the $i$-th family;
indeed, the products of a monomial in the remaining generators, of a $W^1(\mathrm{Jw}_1)$ and of a $W^2(\mathrm{Jw}_2)$ are, up to signs, the basis elements $e_I$, each of them exactly once.

Relative to the second family, the coefficient of $W^2(\mathrm{Jw}_2)$ is $\sum_{\mathrm{Jw}_1}F_{\mathrm{Jw}_1\mathrm{Jw}_2}W^1(\mathrm{Jw}_1)$.
Hence
\[
\int\mathrm{D}W^2\,F=\mathrm{kn}_2\cdot\mathrm{sn}_2\sum_{\mathrm{Jw}_1}F_{\mathrm{Jw}_1\mathrm{Jw}_2^{\mathrm{f}}}W^1(\mathrm{Jw}_1),
\qquad
\int\mathrm{D}W^1\int\mathrm{D}W^2\,F=\mathrm{kn}_1\cdot\mathrm{kn}_2\cdot\mathrm{sn}_1\cdot\mathrm{sn}_2\,F_{\mathrm{Jw}_1^{\mathrm{f}}\mathrm{Jw}_2^{\mathrm{f}}} .
\]
Relative to the first family one uses
\[
W^1(\mathrm{Jw}_1)\,W^2(\mathrm{Jw}_2)=(-1)^{\#|\mathrm{Jw}_1|\,\#|\mathrm{Jw}_2|}\,W^2(\mathrm{Jw}_2)\,W^1(\mathrm{Jw}_1)
\]
and the fact that $\#|\mathrm{Jw}_1^{\mathrm{f}}|$ is even;
thus the coefficient of $W^1(\mathrm{Jw}_1^{\mathrm{f}})$ is $\sum_{\mathrm{Jw}_2}F_{\mathrm{Jw}_1^{\mathrm{f}}\mathrm{Jw}_2}W^2(\mathrm{Jw}_2)$.
Hence
\[
\int\mathrm{D}W^1\,F=\mathrm{kn}_1\cdot\mathrm{sn}_1\sum_{\mathrm{Jw}_2}F_{\mathrm{Jw}_1^{\mathrm{f}}\mathrm{Jw}_2}W^2(\mathrm{Jw}_2),
\qquad
\int\mathrm{D}W^2\int\mathrm{D}W^1\,F=\mathrm{kn}_1\cdot\mathrm{kn}_2\cdot\mathrm{sn}_1\cdot\mathrm{sn}_2\,F_{\mathrm{Jw}_1^{\mathrm{f}}\mathrm{Jw}_2^{\mathrm{f}}} .
\]
The two iterated integrals coincide for every $F\in\mathrm{Gr}$, in particular for the even ones.

\emph{A measure.}
Let $\mathrm{mq}$ be a finite positive measure on a measurable space, and let $\zeta\mapsto F(\zeta)=\sum_IF_I(\zeta)\,e_I$ be a function with values in $\mathrm{Gr}$ whose coefficients $F_I$ are integrable;
put $\int\mathrm{d}\,\mathrm{mq}\,F:=\sum_I\left(\int\mathrm{d}\,\mathrm{mq}\,F_I\right)e_I$.
Let $\mathrm{Lm}$ be a linear map of $\mathrm{Gr}$ into a finite dimensional vector space.
In a basis of this space the coefficients of $\mathrm{Lm}\left(F(\zeta)\right)=\sum_IF_I(\zeta)\,\mathrm{Lm}(e_I)$ are finite linear combinations of the $F_I$, hence integrable, and
\begin{equation}
\int\mathrm{d}\,\mathrm{mq}\,\mathrm{Lm}(F)=\sum_I\left(\int\mathrm{d}\,\mathrm{mq}\,F_I\right)\mathrm{Lm}(e_I)=\mathrm{Lm}\left(\int\mathrm{d}\,\mathrm{mq}\,F\right).
\tag{P.2}
\end{equation}
By (Gr4), $\int\mathrm{D}W$ is a map of this kind.
The relation (P.2) holds as well when $\mathrm{Lm}$ is the multiplication by a fixed element of $\mathrm{Gr}$ from the left or from the right.
\end{proof}

\begin{proof}[Proof of Proposition~22.1, \emph{Instances}]
Here $k\ge0$ is an arbitrary integer and $\ell$ an arbitrary integer $\ge1$, in the case~(i) below a multiple of $L$;
no further property of the side enters.

\emph{The index sets.}
\begin{itemize}
\item[(i)] The family $(\Psi',\bar{\Psi}')$, for every $k\ge0$.
The generators are $\Psi'_\beta(y')$, $\bar{\Psi}'_\beta(y')$, $y'\in\mathbb{Z}^3$, with the antiperiodicity (I.AP) for the side $\ell/L$.
A fundamental domain of $(\ell/L)\mathbb{Z}^3$ in $\mathbb{Z}^3$ has $n_{k+1}=(\ell/L)^3$ points;
so there are $4n_{k+1}$ modes $(y',\beta)$, and $\nu=1$.
\item[(ii)] The families $(\Psi,\bar{\Psi})$, $(W,\bar{W})$, $(W',\bar{W}')$, for $k\ge1$.
The generators sit at the points $y\in\mathbb{Z}^3$, with (I.AP) for the side $\ell$.
A fundamental domain has $n_k=\ell^3$ points;
so there are $4n_k$ modes $(y,\beta)$, $\nu=1$, and $\langle\,\cdot\,,\cdot\,\rangle_\nu=\langle\,\cdot\,,\cdot\,\rangle_1$.
\item[(iii)] The same three families at level $0$.
The generators sit at the points $x=(c,g)\in\mathbb{L}^{\mathrm{f}}_0=\mathbb{Z}^3\times O_{\mathrm{cub}}$, with (I.AP) for the side $\ell$ in the variable $c$.
A fundamental domain is the product of a fundamental domain of $\ell\mathbb{Z}^3$ in $\mathbb{Z}^3$ with $O_{\mathrm{cub}}$, and it has $48\,n_0$ points.
So there are $4\cdot48\,n_0=192\,n_0$ modes $(x,\beta)$, $\nu=1/48$, $\prod_{\mathrm{mo}}\nu(\mathrm{mo})^{-1}=48^{192n_0}$, and, by (I.1.3) with $\eta_0=1$,
\[
\langle f'',f'\rangle_\nu=\sum_{x,\beta}\frac{1}{48}\,f''_\beta(x)\,f'_\beta(x)=\int\mathrm{d}x\,f''(x)\cdot f'(x)=\langle f'',f'\rangle^{\mathrm{ch}} .
\]
\end{itemize}

\emph{Operators and fields.}
In the cases (i), (ii), (iii) the restriction to the fundamental domain is a linear bijection of $\mathfrak{S}(1,\ell/L)$, $\mathfrak{S}(1,\ell)$, $\mathfrak{S}^{\mathrm{ch}}(1,\ell)$, respectively, onto the functions on $\mathrm{Mo}$;
under it the pairing $\langle\,\cdot\,,\cdot\,\rangle_1$, $\langle\,\cdot\,,\cdot\,\rangle_1$, $\langle\,\cdot\,,\cdot\,\rangle^{\mathrm{ch}}$ of Section~I.1.6, respectively, becomes $\langle\,\cdot\,,\cdot\,\rangle_\nu$.
A linear map $O$ of one of these spaces is therefore a linear map of the functions on $\mathrm{Mo}$ with the same determinant, and its transpose for the pairing of Section~I.1.6 is $O^{T}$.
A field with the antiperiodicity (I.AP) is determined by its restriction to the fundamental domain, and $Of'$ and $\langle f'',f'\rangle$, formed on the covering index set, are those of 0.3.

\emph{The fundamental domain.}
Let a second fundamental domain be given.
Each of its points is $t+\ell v$ (in the case~(i): $t+(\ell/L)v$) for exactly one point $t$ of the first domain and one $v\in\mathbb{Z}^3$, and by (I.AP) the generators at the two points differ by the factor $(-1)^{v_1+v_2+v_3}$, the same for $W_\beta$ and $\bar{W}_\beta$.
Hence the subalgebras $\mathrm{Gr}^{\mathrm{o}}$ and the degree in the family are the same for both domains, $W_\beta(t+\ell v)\bar{W}_\beta(t+\ell v)=W_\beta(t)\bar{W}_\beta(t)$, the element $\mathrm{Top}$ is the same, and the weights are constant.
It follows that $[F]_{\mathrm{Top}}$ and $\int\mathrm{D}W\,F$ are the same for both domains.

\emph{The rule (G3)$^{\mathrm{w}}$.}
Let $u$ be a periodic function of the photon site without zeros, with the period lattice $\ell\mathbb{Z}^3$ in the cases (ii), (iii) and $(\ell/L)\mathbb{Z}^3$ in the case~(i).
In the case~(i) the relation $u\left(y'+(\ell/L)v\right)=u(y')$ for all $v\in\mathbb{Z}^3$ says that the function $Ly'\mapsto u(y')$ on the sites of $L\mathbb{Z}^3$ is $\ell\mathbb{Z}^3$-periodic.
Then $uW$ and $u^{-1}\bar{W}$ have the antiperiodicity (I.AP), and the homomorphism of (G3)$^{\mathrm{w}}$ does not depend on the fundamental domain.
In the case~(iii) one has $u(x)=u(c)$ for the $48$ sites $x$ over $c$, and $\mathrm{Top}$ is multiplied by the factor
\[
\prod_c\left(u(c)\,u(c)^{-1}\right)^{192}=1,
\]
since there are $192=4\cdot48$ modes over each photon site $c$.
\end{proof}

\subsection*{2. Proposition 22.2 (normalisation of the block transformation; the density of level 0)}

\begin{proof}[Proof of Proposition~22.2]
In this proof (G0)$^{\mathrm{w}}$, (G1)$^{\mathrm{w}}$, (G3)$^{\mathrm{w}}$ and (G5)$^{\mathrm{w}}$ denote the clauses of Proposition~22.1, the instances (i) and (iii) are those of its paragraph on instances, and (Gr3), (Gr5), (P.2) are the facts of these names in Section~0.2 and in the proof of Proposition~22.1.

\textbf{(a)}
The integration is over the family $(\Psi',\bar{\Psi}')$.
This is the instance (i) of Proposition~22.1 at the side $\ell/L$, which is a positive integer: there are $4n_{k+1}$ modes $(y',\beta)$, each of weight $\nu=1$.
Fields are understood in the sense of Section~G3.1.0, read at the side $\ell\in L\mathbb{N}$.
By Sections~I.3.7 and~I.5.7 one has, for $y\in L\mathbb{Z}^3$,
\[
\Psi'_L(y)=L^{-1}\cdot\Psi'\left(\frac{y}{L}\right),\qquad
\bar{\Psi}'_L(y)=L^{-1}\cdot\bar{\Psi}'\left(\frac{y}{L}\right),\qquad
\langle v'',v'\rangle_L=L^3\sum_{y\in L\mathbb{Z}^3/\ell\mathbb{Z}^3}v''(y)\cdot v'(y).
\]
Let
\[
\mathrm{Eq}:=-b_L\cdot\left\langle\exp\left(-i\cdot e_k\cdot\omega\right)\cdot\bar{\Psi}'_L-\overline{\mathrm{ff}},\ \exp\left(i\cdot e_k\cdot\omega\right)\cdot\Psi'_L-\mathrm{ff}\right\rangle_L .
\]
The summand at $y$ of this pairing is the product of two antiperiodic factors, the function $\omega$ being periodic; so the summands are $\ell\mathbb{Z}^3$-periodic, and $\mathrm{Eq}$ is an even element.

\emph{Phases.}
For $y'\in\mathbb{Z}^3$ let $u(y'):=\exp\left(-i\cdot e_k\cdot\omega(Ly')\right)$.
The function $u$ is $(\ell/L)\mathbb{Z}^3$-periodic, because $\omega$ is $\ell\mathbb{Z}^3$-periodic, and it has no zero.
For $y\in L\mathbb{Z}^3$ one finds
\[
(u\cdot\Psi')_L(y)=L^{-1}\cdot u\left(\frac{y}{L}\right)\cdot\Psi'\left(\frac{y}{L}\right)=\exp\left(-i\cdot e_k\cdot\omega(y)\right)\cdot\Psi'_L(y),
\qquad
(u^{-1}\cdot\bar{\Psi}')_L(y)=\exp\left(i\cdot e_k\cdot\omega(y)\right)\cdot\bar{\Psi}'_L(y).
\]
The fields $\mathrm{ff}$ and $\overline{\mathrm{ff}}$ contain no generator of the family $(\Psi',\bar{\Psi}')$.
Therefore
\[
\mathrm{Eq}\left[\bar{\Psi}'\to u^{-1}\cdot\bar{\Psi}',\ \Psi'\to u\cdot\Psi'\right]=-b_L\cdot\left\langle\bar{\Psi}'_L-\overline{\mathrm{ff}},\ \Psi'_L-\mathrm{ff}\right\rangle_L=:\mathrm{Eq}_1 ,
\]
and (G3)$^{\mathrm{w}}$, which applies to the $(\ell/L)\mathbb{Z}^3$-periodic function $u$ without zeros, gives together with (Gr3)
\[
\int\mathrm{D}\Psi'\,\exp(\mathrm{Eq})=\int\mathrm{D}\Psi'\,\exp(\mathrm{Eq}_1).
\]

\emph{Translation.}
For $y'\in\mathbb{Z}^3$ let $\mathrm{gf}(y'):=L\cdot\mathrm{ff}(Ly')$ and $\bar{g}(y'):=L\cdot\overline{\mathrm{ff}}(Ly')$.
These are fields on $T^0_{k+1}$ whose values lie in the subalgebra generated by the generators other than $\Psi'$, $\bar{\Psi}'$; their antiperiodicity follows from
\[
\mathrm{gf}\left(y'+\frac{\ell}{L}\cdot e_\mu\right)=L\cdot\mathrm{ff}(Ly'+\ell e_\mu)=-\mathrm{gf}(y'),
\]
and the same computation for $\bar{g}$.
Moreover $\mathrm{gf}_L=\mathrm{ff}$ and $\bar{g}_L=\overline{\mathrm{ff}}$.
Consequently
\[
\mathrm{Eq}_1\left[\bar{\Psi}'\to\bar{\Psi}'+\bar{g},\ \Psi'\to\Psi'+\mathrm{gf}\right]
=-b_L\cdot\left\langle\bar{\Psi}'_L+\bar{g}_L-\overline{\mathrm{ff}},\ \Psi'_L+\mathrm{gf}_L-\mathrm{ff}\right\rangle_L
=-b_L\cdot\left\langle\bar{\Psi}'_L,\Psi'_L\right\rangle_L ,
\]
and (G1)$^{\mathrm{w}}$, again with (Gr3), yields
\[
\int\mathrm{D}\Psi'\,\exp(\mathrm{Eq}_1)=\int\mathrm{D}\Psi'\,\exp\left(-b_L\cdot\left\langle\bar{\Psi}'_L,\Psi'_L\right\rangle_L\right).
\]

\emph{The Gaussian integral.}
Writing $y=Ly'$ one obtains
\[
\left\langle\bar{\Psi}'_L,\Psi'_L\right\rangle_L
=L^3\sum_{y\in L\mathbb{Z}^3/\ell\mathbb{Z}^3}L^{-2}\cdot\bar{\Psi}'\left(\frac{y}{L}\right)\cdot\Psi'\left(\frac{y}{L}\right)
=L\sum_{y'\in\mathbb{Z}^3/(\ell/L)\mathbb{Z}^3}\bar{\Psi}'(y')\cdot\Psi'(y')
=L\cdot\left\langle\bar{\Psi}',\Psi'\right\rangle_1 ,
\]
and $b_L\cdot L=b$.
Applied to the operator $O=b\cdot 1$ on the functions on the $4n_{k+1}$ modes, (G0)$^{\mathrm{w}}$ gives
\[
\int\mathrm{D}\Psi'\,\exp\left(-\left\langle\bar{\Psi}',b\cdot\Psi'\right\rangle_1\right)=\det(b\cdot 1)=b^{4n_{k+1}}=c_k^{-1}.
\]
The three steps together give $c_k\int\mathrm{D}\Psi'\,\exp(\mathrm{Eq})=c_k\cdot c_k^{-1}=1$, which is (a).

\textbf{(b)}
Let $\mathscr{C}:=\mathscr{C}'_L$, $\mathscr{B}_k:=H_kA^{\min}_k(A')$ and $\omega_k:=\omega_k(A')$.
All elements that occur below lie in $\mathrm{Gr}[\Psi',\Psi]$.

\emph{The integrand.}
Since $L>1$, one has $e_k=L^{-\frac{1}{2}}\cdot e_{k+1}<e_{k+1}\le e_0$.
The hypotheses of \cite[Proposition~15.30(ii)]{QEDI} (equivalently, of \cite[Proposition~5.38(ii)]{QEDI}) are satisfied at the data of (S22).
Indeed, the inequalities among them are conditions on the constants (the row of \cite[Proposition~15.30(ii)]{QEDI} in the list~II of Section~5 of this chapter), and those that concern $e_0$ are imposed for every $v\ge-\log e_0$, so that they hold at $v=v_k$;
the sides $\ell$ and $\ell/L$ lie in $\mathrm{Cl}$;
and $\mathrm{Dom}_{k+1}$ is formed with the side $\ell/L$ and the charge $e_{k+1}=L^{\frac{1}{2}}\cdot e_k$.
By \cite[Proposition~15.30(ii)]{QEDI} the field $A^{\min}_k(A')+\mathrm{Sq}_k\zeta$ lies in $\mathrm{Dom}_k$ for every $\zeta\in\mathrm{Box}_k$.
Hence
\[
F(\zeta):=F\left(A^{\min}_k(A')+\mathrm{Sq}_k\zeta,\mathscr{C};\Psi\right)
\]
is defined; it is an even element of $\mathrm{Gr}[\Psi]$ whose coefficients are integrable over $\mathrm{Box}_k$ with respect to $\mu_I$.
(If $F(\,\cdot\,,\mathscr{C};\Psi)$ is a bounded Borel map on $\mathrm{Dom}_k$, this is seen as follows: the map $\zeta\mapsto A^{\min}_k(A')+\mathrm{Sq}_k\zeta$ is affine, hence continuous, so that the coefficients of $F(\zeta)$ are bounded Borel functions on $\mathrm{Box}_k$, and $\mu_I$ is a probability measure.)
Put
\begin{align*}
k\ge 1:&\qquad\mathrm{ff}:=Q(\mathscr{B}_k+\mathscr{C})\Psi,\qquad\overline{\mathrm{ff}}:=Q(-\mathscr{B}_k-\mathscr{C})\bar{\Psi};\\
k=0:&\qquad\mathrm{ff}:=Q(\mathscr{B}_0+\mathscr{C})J^{*}\Psi,\qquad\overline{\mathrm{ff}}:=Q(-\mathscr{B}_0-\mathscr{C})J^{*}\bar{\Psi};
\end{align*}
\[
\mathrm{Eq}:=-b_L\cdot\left\langle\exp\left(-i\cdot e_k\cdot\omega_k\right)\cdot\bar{\Psi}'_L-\overline{\mathrm{ff}},\ \exp\left(i\cdot e_k\cdot\omega_k\right)\cdot\Psi'_L-\mathrm{ff}\right\rangle_L .
\]
For every complex bond function $B$ on $T_k$ the operator $Q(B)$ maps $\mathfrak{S}(1,\ell)$ into $\mathfrak{S}(L,\ell)$, and $J^{*}$ maps $\mathfrak{S}^{\mathrm{ch}}(1,\ell)$ into $\mathfrak{S}(1,\ell)$ (Sections~I.3.1 and~I.3.3).
So $\mathrm{ff}$ and $\overline{\mathrm{ff}}$ are fields on $L\mathbb{Z}^3/\ell\mathbb{Z}^3$ with four components, and their values lie in the subalgebra generated by $\Psi$, $\bar{\Psi}$.
At $k=0$ the field $\mathrm{ff}$ reads
\[
\mathrm{ff}(y')=L^{-3}\sum_{y\in B(y')}U_{\mathscr{B}_0+\mathscr{C}}(y',y)\cdot\frac{1}{48}\sum_{g\in O_{\mathrm{cub}}}\Psi(y,g).
\]
The restriction of $\omega_k$ to $L\mathbb{Z}^3$ is real and $\ell\mathbb{Z}^3$-periodic (Section~G3.2.0, the paragraph on $\omega_k$).
The element $\mathrm{Eq}$ does not depend on $\zeta$.
By Section~I.5.7,
\[
\left(\mathbf{T}^x_kF\right)(A',\mathscr{C}';\Psi')=c_k\int_{\mathrm{Box}_k}\mathrm{d}\mu_I(\zeta)\int\mathrm{D}\Psi\,\exp(\mathrm{Eq})\cdot F(\zeta).
\]
By (P.2) for the multiplication by $\exp(\mathrm{Eq})$, the coefficients of $\exp(\mathrm{Eq})\cdot F(\zeta)$ are finite linear combinations of those of $F(\zeta)$; hence they are integrable.

\emph{Exchange of the integrations.}
The restriction of $\mu_I$ to $\mathrm{Box}_k$ is a finite positive measure.
The side $\ell$ is a multiple of $L$, as $\ell/L\in\mathrm{Cl}$.
Apply (G5)$^{\mathrm{w}}$, first in the form (P.2) with $\mathrm{Lm}=\int\mathrm{D}\Psi'$ and then as the interchange of the two Berezin integrals;
then (Gr5), which applies because $F(\zeta)$ lies in the subalgebra generated by the generators other than $\Psi'$, $\bar{\Psi}'$;
and finally (a), with the restriction of $\omega_k$ to $L\mathbb{Z}^3$ in the place of $\omega$ and with the fields $\mathrm{ff}$, $\overline{\mathrm{ff}}$ above.
One obtains
\begin{align*}
\int\mathrm{D}\Psi'\,\left(\mathbf{T}^x_kF\right)(A',\mathscr{C}';\Psi')
&=c_k\int_{\mathrm{Box}_k}\mathrm{d}\mu_I(\zeta)\int\mathrm{D}\Psi'\int\mathrm{D}\Psi\,\exp(\mathrm{Eq})\cdot F(\zeta)\\
&=\int_{\mathrm{Box}_k}\mathrm{d}\mu_I(\zeta)\int\mathrm{D}\Psi\,\left[c_k\int\mathrm{D}\Psi'\,\exp(\mathrm{Eq})\right]\cdot F(\zeta)\\
&=\int_{\mathrm{Box}_k}\mathrm{d}\mu_I(\zeta)\int\mathrm{D}\Psi\,F(\zeta).
\end{align*}
At $k=0$ the integral $\int\mathrm{D}\Psi$ is that of the instance (iii) of Proposition~22.1: it concerns $192n_0$ pairs of generators on $T^{\mathrm{f}}_0$, with the weight $1/48$.

\textbf{(c)}
The datum $\mathrm{Bf}$ is an $\ell\mathbb{Z}^3$-periodic bond function on $\mathbb{Z}^3$; for $\ell\ge 3$ this is a bond function on $T_0$, and for $\ell=1,2$ it is the datum of the statement.
The operator $D_{\mathrm{Bf}}$ is given by the display (I.2.3) with $\eta_0=1$ and with $q$ in the place of $e_k$, on the space $\mathfrak{S}^{\mathrm{ch}}(1,\ell)$, which has dimension $192n_0$.
It maps this space into itself, because $\mathrm{Bf}$ is periodic and $x\mapsto x^{(i)}$ commutes with the translations (I.s1).

\emph{The determinant $\det(D_0+\mu)$ is not zero.}
Let $i\in\{0,1,2,3\}$.
By (I.s1) the map $x\mapsto x^{(i)}$ is a bijection of $\mathbb{L}^{\mathrm{f}}_0$ that commutes with the translations, and $\left(x^{(i)}\right)^{(i)}=x$.
Hence it maps a fundamental domain of $\ell\mathbb{Z}^3$ onto a fundamental domain, and every $\ell\mathbb{Z}^3$-periodic function $\mathrm{Fp}$ on $\mathbb{L}^{\mathrm{f}}_0$ with values in a vector space satisfies
\begin{equation}
\int\mathrm{d}x\,\mathrm{Fp}\left(x^{(i)}\right)=\int\mathrm{d}x\,\mathrm{Fp}(x).\tag{P.3}
\end{equation}
By (I.2.4) at $A=0$ and $\eta_0=1$,
\begin{gather*}
D_0=\frac{\rho}{2}\cdot\nabla_0^{*}\nabla_0+\frac{1}{2}\cdot\mathrm{Gh}_0,\qquad
\nabla_0^{*}\nabla_0=\sum_iw_i\cdot\left(1-R_i\right),\\
(R_if)(x)=f\left(x^{(i)}\right),\qquad
(\mathrm{Gh}_0f)(x)=\sum_iw_i\cdot\left(\gamma\cdot\hat{d}_i(x)\right)f\left(x^{(i)}\right).
\end{gather*}
Let $f',f\in\mathfrak{S}^{\mathrm{ch}}(1,\ell)$.
By (P.3),
\[
(f',R_if)^{\mathrm{ch}}=\int\mathrm{d}x\,\overline{f'(x)}\cdot f\left(x^{(i)}\right)=\int\mathrm{d}x\,\overline{f'\left(x^{(i)}\right)}\cdot f(x)=(R_if',f)^{\mathrm{ch}},\qquad R_i^2=1.
\]
So $R_i$ is Hermitian and unitary, $(f,R_if)^{\mathrm{ch}}$ is real and at most $(f,f)^{\mathrm{ch}}$, and, as $w_i>0$ (Section~I.0.5),
\[
(f,\nabla_0^{*}\nabla_0f)^{\mathrm{ch}}=\sum_iw_i\cdot\left[(f,f)^{\mathrm{ch}}-(f,R_if)^{\mathrm{ch}}\right]\ge 0.
\]
By (P.3), by $\hat{d}_i\left(x^{(i)}\right)=-\hat{d}_i(x)$ (I.s2), and because $\gamma\cdot v$ is Hermitian for real $v$ (I.c1),
\begin{align*}
(f',\mathrm{Gh}_0f)^{\mathrm{ch}}
&=\sum_iw_i\int\mathrm{d}x\,\overline{f'(x)}\cdot\left(\gamma\cdot\hat{d}_i(x)\right)f\left(x^{(i)}\right)
=-\sum_iw_i\int\mathrm{d}x\,\overline{f'\left(x^{(i)}\right)}\cdot\left(\gamma\cdot\hat{d}_i(x)\right)f(x)\\
&=-\sum_iw_i\int\mathrm{d}x\,\overline{\left(\gamma\cdot\hat{d}_i(x)\right)f'\left(x^{(i)}\right)}\cdot f(x)
=-(\mathrm{Gh}_0f',f)^{\mathrm{ch}}.
\end{align*}
It follows that $(f,\mathrm{Gh}_0f)^{\mathrm{ch}}$ is purely imaginary, and, $\rho$ being non-negative,
\[
\operatorname{Re}\left(f,(D_0+\mu)f\right)^{\mathrm{ch}}=\frac{\rho}{2}\cdot(f,\nabla_0^{*}\nabla_0f)^{\mathrm{ch}}+\mu\cdot(f,f)^{\mathrm{ch}}\ge\mu\cdot(f,f)^{\mathrm{ch}}.
\]
This inequality is (I.O2) at $A=0$, and it holds for every integer $\ell\ge 1$.
Since $\mu>0$, one has $(D_0+\mu)f\ne 0$ for $f\ne 0$, and so $\det(D_0+\mu)\ne 0$.

\emph{The integral.}
The mass term is $m$ times the identity on all $192$ components over each photon site:
\[
m\cdot\langle\bar{\Psi},\Psi\rangle^{\mathrm{ch}}=\frac{1}{48}\sum_x\sum_\beta\bar{\Psi}_\beta(x)\cdot m\cdot\Psi_\beta(x)=\langle\bar{\Psi},m\cdot 1\,\Psi\rangle^{\mathrm{ch}}.
\]
Hence the exponent equals $-\langle\bar{\Psi},(D_{\mathrm{Bf}}+\mu+m)\Psi\rangle^{\mathrm{ch}}-\varepsilon^{\mathrm{v}}\cdot n_0$.
By (Gr3), by the linearity of $\int\mathrm{D}\Psi$, and by (G0)$^{\mathrm{w}}$ at the instance (iii) (with $t=x\in T^{\mathrm{f}}_0$, $\nu=1/48$ and $\langle\cdot,\cdot\rangle_\nu=\langle\cdot,\cdot\rangle^{\mathrm{ch}}$), applied to $O=D_{\mathrm{Bf}}+\mu+m$,
\begin{multline*}
\int\mathrm{D}\Psi\,\det(D_0+\mu)^{-1}\cdot\exp\left\{-\left\langle\bar{\Psi},(D_{\mathrm{Bf}}+\mu)\Psi\right\rangle^{\mathrm{ch}}-m\cdot\langle\bar{\Psi},\Psi\rangle^{\mathrm{ch}}-\varepsilon^{\mathrm{v}}\cdot n_0\right\}\\
=\det(D_0+\mu)^{-1}\cdot\exp\left(-\varepsilon^{\mathrm{v}}\cdot n_0\right)\cdot\det(D_{\mathrm{Bf}}+\mu+m).
\end{multline*}
As (G0)$^{\mathrm{w}}$ holds for every operator $O$, nothing is required of $D_{\mathrm{Bf}}+\mu+m$: the numbers $m$, $\varepsilon^{\mathrm{v}}$ and the bond function $\mathrm{Bf}$ are arbitrary complex data.

\emph{The flow.}
At level $0$ the side is $\ell=n$, so that $n_0=n^3$; the mass is $\mu=\bar{m}\cdot\eta>0$, and the charge is $q$.
By Sections~I.3.9 and~I.5.7 one has $H_0A_0=A_0$ and $H_0(\mathscr{T})=\bar{H}_0(\mathscr{T})=1$, so that $\mathrm{ev}_{\mathscr{T}}$ maps $\psi$ to $\Psi$ and $\bar{\psi}$ to $\bar{\Psi}$;
moreover $D_0(\mathscr{T})=D_{\mathscr{T}}+\bar{m}\cdot\eta$, and $D_0(0)$ is the operator written $D_0+\mu$ above.
With $\mathscr{E}=0$ and $\mathscr{T}=A_0+a$, Section~0.4 gives
\begin{multline*}
\mathscr{F}^x_0[\varepsilon^{\mathrm{v}},m,0](A_0,a;\Psi)=\det\left(D_0+\bar{m}\cdot\eta\right)^{-1}\\
\cdot\exp\left\{-\left\langle\bar{\Psi},\left(D_{A_0+a}+\bar{m}\cdot\eta\right)\Psi\right\rangle^{\mathrm{ch}}-m\cdot\langle\bar{\Psi},\Psi\rangle^{\mathrm{ch}}-\varepsilon^{\mathrm{v}}\cdot n^3+\sum_X\mathfrak{P}_0(X,a)\right\},
\end{multline*}
which is (I.L0.22).
Further $y^x_0=\int\mathrm{D}\Psi\,\mathscr{F}^x_0$ (Section~I.5.7).
The number $\exp\left\{\sum_X\mathfrak{P}_0(X,a)\right\}$ is a factor, and the first part of (c), applied with $\mathrm{Bf}=A_0+a$, $m=m_{\mathrm{in}}$ and $\varepsilon^{\mathrm{v}}=\varepsilon^{\mathrm{v}}_{\mathrm{in}}$, gives
\[
y^x_0[\varepsilon^{\mathrm{v}}_{\mathrm{in}},m_{\mathrm{in}},0](A_0;a)=\exp\left\{\sum_X\mathfrak{P}_0(X,a)\right\}\cdot\exp\left(-\varepsilon^{\mathrm{v}}_{\mathrm{in}}\cdot n^3\right)\cdot\frac{\det\left(D_{A_0+a}+\bar{m}\cdot\eta+m_{\mathrm{in}}\right)}{\det\left(D_0+\bar{m}\cdot\eta\right)};
\]
in the flow $\mathfrak{P}_0=0$.
This identity holds whatever the sign and the size of $\bar{m}\cdot\eta+m_{\mathrm{in}}$.
\end{proof}

\subsection*{3. Proposition 22.3 (the one-step formula with slot)}

\subsubsection*{Proof}

\emph{Notation.}
All elements below lie in $\mathrm{Gr}:=\mathrm{Gr}[\Psi',\Psi,W]$, the Grassmann algebra generated by the three families $(\Psi',\bar{\Psi}')$, $(\Psi,\bar{\Psi})$ and $(W,\bar{W})$; these are the instances (i), and (ii) for $k\ge 1$, respectively (iii) for $k=0$, of Proposition~22.1.
The clauses (G1)$^{\mathrm{w}}$--(G5)$^{\mathrm{w}}$ are those of Proposition~22.1.
For $\zeta\in\mathrm{Box}_k$ let
\begin{gather*}
A_k:=A^{\min}_k(A')+\mathrm{Sq}_k\zeta,\qquad
\mathscr{Z}:=H_k\mathrm{Sq}_k\zeta=H_k\,C\,C_k^{\frac{1}{2}}\zeta,\\
\mathscr{T}^z:=\mathscr{T}+\mathscr{Z},\qquad
\mathscr{T}^g:=\mathscr{T}+\mathscr{Z}-d\omega_k .
\end{gather*}
Here $\mathscr{Z}$ is the field $\mathscr{Z}$ of the statement at $\mathsf{Z}=\zeta$; at $k=0$ one has $H_0=1$ and $\mathscr{Z}=C\,C_0^{\frac{1}{2}}\zeta$.
The field $\mathscr{B}_k=H_kA^{\min}_k(A')$ and the function $\omega_k=\omega_k(A')$ are those of Section~I.5.7.

\paragraph{(a) All objects are defined.}

\begin{proof}[Proof of Proposition~22.3\,(a)]
\textbf{(Ob1)} \emph{Data.}
Since $L>1$, one has $e_k=L^{-\frac{1}{2}}\cdot e_{k+1}<e_{k+1}\le e_0$ and $\mu\le 1/L<1$.
Hence $e_k\in(0,e_0]$ and $\mu\in(0,1]$.
The data $(k,\ell,\mu,e_k)$, together with the level $k+1$ of Section~G3.2.0, are data of the setting (S21), and $(k,\ell,e_k)$ are data (S-1) of Proposition~16.3.
They also meet the hypotheses of the results of~\cite{QEDI} that are applied below:
\cite[Theorem~14.15(b)]{QEDI} is stated for $k\ge 0$, $\ell\in L\mathbb{N}$ and every real bond function $A'$ on $T^0_{k+1}$ with arbitrary harmonic part;
\cite[Lemma~17.4]{QEDI} and \cite[Lemma~17.6]{QEDI}, the latter without its assertion on $\mathrm{Dom}_k$, are stated for $k\ge 0$, $\ell\in\mathrm{Cl}$ and $e_k\in(0,e_0]$, and \cite[Lemma~17.5]{QEDI} for $\ell,\ell/L\in\mathrm{Cl}$, $e_{k+1}=L^{\frac{1}{2}}\cdot e_k$ and $e_k\le e_0$ (the paragraph on hypotheses in \cite[Section~17.3]{QEDI});
\cite[Proposition~15.21(i), (ii)]{QEDI} is stated for every side $\ell$ that is a multiple of $L$;
\cite[Proposition~15.30(ii)]{QEDI} is stated for $\ell,\ell/L\in\mathrm{Cl}$, with $\mathrm{Dom}_{k+1}$ formed with the side $\ell/L$ and the charge $e_{k+1}$.
One has $e_0^{\varepsilon}\le\frac{1}{4}$, the first inequality of the condition (A8) of Section~I.8.4 (which is the condition (A8) of \cite[Section~5.6.2]{QEDI}), and $\varepsilon>0$.
Hence $e_0<1$, $v_k=-\log e_k>0$ and $p_{0,k}>0$.
By \cite[Proposition~15.21(ii)]{QEDI} the family $\{\Upsilon\}$ is a basis of $V_k$.
So $\zeta\mapsto\left(\zeta(\Upsilon)\right)_{\Upsilon}$ is a linear bijection of $V_k$ onto $\mathbb{R}^{\dim V_k}$, and $\mathrm{Box}_k$ is the inverse image of the cube $[-p_{0,k},p_{0,k}]^{\dim V_k}$; it is a compact set with interior points.
Therefore $\mu_I(\mathrm{Box}_k)>0$, the measure $\hat{\mu}$ and the number $\delta\varepsilon^0_k$ of Section~I.5.7 (\cite[Section~5.2.5]{QEDI}) are defined, and
\begin{equation}
\int_{\mathrm{Box}_k}\mathrm{d}\mu_I\,F=\mu_I(\mathrm{Box}_k)\cdot\int\mathrm{d}\hat{\mu}\,F=\exp\left(-\delta\varepsilon^0_k\cdot n_k\right)\cdot\int\mathrm{d}\hat{\mu}\,F.\tag{P.4}
\end{equation}

\textbf{(Ob2)} \emph{Domains.}
The field $\mathscr{A}'=H_{k+1}A'$ lies in $\frac{1}{4}\tilde{R}_{k+1}$: this is the definition of $\mathrm{Dom}_{k+1}$ (side $\ell/L$, charge $e_{k+1}$).
By \cite[Lemma~17.5(a)]{QEDI} with the factor $\frac{1}{4}$, the field $A=\mathscr{A}'_L$ lies in $\frac{1}{4}L^{-\frac{3}{4}}\cdot\tilde{R}_k$.
The sets $\mathrm{cf}\cdot\tilde{R}_k$, $0<\mathrm{cf}\le 1$, increase with $\mathrm{cf}$, since in their definition the thresholds are multiplied by $\mathrm{cf}$.
Hence $A$ lies in $L^{-\frac{3}{4}}\cdot\tilde{R}_k$, which is contained in $\tilde{R}_k^{-}$ by \cite[Lemma~17.5(b)]{QEDI}; and $\tilde{R}_k^{-}\subset\tilde{R}_k$ by the definition of $\tilde{R}_k^{-}$.
By \cite[Lemma~17.5(c)]{QEDI}, $\mathscr{C}=\mathscr{C}'_L$ lies in $\mathfrak{X}_k$.

Let $\zeta\in\mathrm{Box}_k$.
By \cite[Proposition~15.30(ii)]{QEDI} the field $A_k$ lies in $\mathrm{Dom}_k$, that is, $H_kA_k\in\frac{1}{4}\tilde{R}_k$.
The map $H_k$ is linear (\cite[Proposition~14.8(ii)]{QEDI}), and $H_kA^{\min}_k(A')=(H_{k+1}A')_L-d\omega_k$ by \cite[Theorem~14.15(b)]{QEDI}.
Consequently
\begin{equation}
H_kA_k=A-d\omega_k+\mathscr{Z},\qquad
H_kA_k+\mathscr{C}=\mathscr{T}^g,\qquad
\mathscr{B}_k+\mathscr{C}=H_kA^{\min}_k(A')+\mathscr{C}=\mathscr{T}-d\omega_k .\tag{P.5}
\end{equation}
By \cite[Theorem~14.15(b)]{QEDI}, $\omega_k$ is a real function on the sites of $T_k$.
By \cite[Lemma~17.6]{QEDI} the set $\frac{1}{4}\tilde{R}_k$ is invariant under $A\mapsto A-d\boldsymbol{\lambda}$ for every real function $\boldsymbol{\lambda}$ on the sites of $T_k$; here $\boldsymbol{\lambda}=-\omega_k$.
So $A+\mathscr{Z}=H_kA_k+d\omega_k$ lies in $\frac{1}{4}\tilde{R}_k$, which is a subset of $\tilde{R}_k$.
Thus the pairs $(A,\mathscr{C})$, $(A+\mathscr{Z},\mathscr{C})$ and $(H_kA_k,\mathscr{C})$ lie in $\tilde{R}_k\times\mathfrak{X}_k$.

By \cite[Lemma~17.4(a)]{QEDI} the set $\tilde{R}_k+\mathfrak{X}_k$ is contained in $\tilde{G}_k$; hence $\mathscr{T}=A+\mathscr{C}$, $\mathscr{T}^z=(A+\mathscr{Z})+\mathscr{C}$ and $\mathscr{T}^g=H_kA_k+\mathscr{C}$ lie in $\tilde{G}_k$.
By \cite[Lemma~17.6]{QEDI} their negatives lie in $\tilde{G}_k$ as well.
Finally $0\in\tilde{G}_k$, since all fields in the definition of $\tilde{G}_k$ (Section~I.5.3) may be taken to be $0$.

\textbf{(Ob3)} \emph{Operators.}
By (Ob1) and Proposition~21.4\,(d2), $\tilde{G}_k$ is contained in $\mathrm{Iv}_k(\ell,\mu,e_k)$: for every $B\in\tilde{G}_k$ the operator $\Gamma_k(B)$ exists, and, if $k\ge 1$, so does $S_k(B)$.
Let $B$ be one of $\mathscr{T}$, $\mathscr{T}^z$, $\mathscr{T}^g$, $0$; then $B$ and $-B$ lie in $\tilde{G}_k$.
For $k\ge 1$ the operators $S_k(B)$, $S_k(-B)$ and $\Gamma_k(B)$ exist; hence the operators
\[
H_k(B)=b_k\cdot S_k(B)JQ_k^T(-B),\qquad
D_k(B)=b_k-b_k^2\cdot Q_k(B)J^{*}S_k(B)JQ_k^T(-B),\qquad
\bar{H}_k(B)=\mathsf{C}^{-1}H_k(-B)\mathsf{C}
\]
exist.
For $k=0$ the operators $H_0(B)=\bar{H}_0(B)=1$ and $D_0(B)=D_B+\mu$ contain no inverse, and $\Gamma_0(B)=\left(D_B+\mu+b_L\cdot JP(B)J^{*}\right)^{-1}$ exists.
The operators $T_k(B)$, $\bar{T}_k(B)$ contain no inverse.

Apply Proposition~21.4\,(b) with $B'=\mathscr{T}'$, $B=\mathscr{T}'_L=\mathscr{T}$, and with $B'=0$, $B=0$; its supposition holds by what has just been shown.
It follows that $S_{k+1}(\pm\mathscr{T}')$ and $S_{k+1}(0)$ exist, and hence that $H_{k+1}$, $\bar{H}_{k+1}$ and $D_{k+1}$ exist at $\mathscr{T}'$ and at $0$.
So $\psi'$, $\bar{\psi}'$, $\chi'$, $\bar{\chi}'$, $\psi^0$, $\bar{\psi}^0$, $\chi^0$, $\bar{\chi}^0$, $\Psi^c$, $\bar{\Psi}^c$, $\Psi^f$, $\bar{\Psi}^f$ are defined.
They are fields with values in the subalgebra generated by $(\Psi',\bar{\Psi}')$, respectively, for $\Psi^f$ and $\bar{\Psi}^f$, in the subalgebra generated by $(\Psi',\bar{\Psi}')$ and $(W,\bar{W})$.
For $(x,x')\in\mathrm{Pzt}_k$ the pair $(x/L,x'/L)$ belongs to $\mathrm{Pzt}_{k+1}$ (Section~I.3.7); therefore $\chi^0=\chi'_L$ and $\bar{\chi}^0=\bar{\chi}'_L$ are defined on $\mathrm{Pz}_k$.

\textbf{(Ob4)} \emph{Determinants.}
The bond function $B=0$ is real, and $\mu>0$.
By Proposition~21.5\,(b), $\det D_k(0)\ne 0$, and by Proposition~21.5\,(c), $\delta Z_k(0)\ne 0$.
Proposition~21.4\,(c) at $B'=0$ gives
\begin{equation}
\det D_{k+1}(0)=b^{4n_{k+1}}\cdot\frac{\det D_k(0)}{\delta Z_k(0)}\ne 0.\tag{P.6}
\end{equation}

\textbf{(Ob5)} \emph{Dependence on the connection.}
In this step the entries of an operator are its matrix entries in the bases of the delta functions over fundamental domains.
The entries of $D_B$ are entire functions of the finitely many numbers $B(b)$ (I.O1).
The numbers $(\tau B)(c,c')$ and $(\tau_kB)(y,c)$ are linear forms in these numbers (Section~I.1.1 and (I.3.3)).
So $U_B(c,c')$, $\varphi_B(y,c)$ and their inverses are entire functions without zeros, and the entries of $Q_k(B)J^{*}$, $JQ_k^T(-B)$, $Q(B)$, $Q^T(-B)$, $P(B)$, $JP_k(B)J^{*}$ (for $k\ge 1$), of $Q(B)J^{*}$, $JQ^T(-B)$, $JP(B)J^{*}$ (for $k=0$), and the coefficients of $\delta_{\alpha,B}$, $\delta_{\alpha,-B}$ (I.1.7) are entire.
Let $k\ge 1$.
The entries of $F_{\infty}(B)=D_B+\mu+b_k\cdot JP_k(B)J^{*}$ and the determinant $\det F_{\infty}(B)$ are entire.
On the open set $\left\{B:\det F_{\infty}(B)\ne 0\right\}$ the entries of $S_k(B)$ are, by Cramer's rule, quotients of entire functions by $\det F_{\infty}(B)$; so they are analytic there, and so are the entries of $H_k(B)$ and $D_k(B)$.
The entries of $\bar{H}_k(B)$ are analytic on $\left\{B:\det F_{\infty}(-B)\ne 0\right\}$.
For $k=0$ the entries of $H_0=\bar{H}_0=1$ and of $D_0(B)$ are entire.
By Proposition~21.4\,(d2) and by \cite[Lemma~17.6]{QEDI}, which gives $\tilde{G}_k=-\tilde{G}_k$, both open sets contain $\tilde{G}_k$.
Consequently, for every $k\ge 0$,
\begin{equation}\tag{P.7}
\begin{split}
&\text{the entries of }D_k(B),\ H_k(B),\ \bar{H}_k(B)\text{ and the coefficients of }\delta_{\alpha,B}H_k(B),\ \delta_{\alpha,-B}\bar{H}_k(B)\\
&\text{are continuous functions of }B\text{ on }\tilde{G}_k.
\end{split}
\end{equation}

\textbf{(Ob6)} \emph{Continuity, parity.}
Let $(\tilde{m},\tilde{\mathscr{E}})$ be an arbitrary pair in $\mathbb{C}\times K^{x,C}_k$, and let $\zeta\mapsto\mathscr{A}(\zeta)$ be an affine map of $V_k$ into the bond functions on $T_k$ that maps $\mathrm{Box}_k$ into $\tilde{R}_k$.
Then $\mathscr{A}(\zeta)+\mathscr{C}$ lies in $\tilde{G}_k$ for $\zeta\in\mathrm{Box}_k$ (\cite[Lemma~17.4(a)]{QEDI}).
By the display of Section~I.5.4, the coefficients of
\[
-\tilde{m}\cdot\langle\bar{\psi},\psi\rangle+\sum_X\tilde{\mathscr{E}}\left(X;\mathscr{A}(\zeta),\mathscr{C}\right)
\]
in the basis of the ordered monomials of $\mathrm{Gr}[\psi,\chi]$ are finite linear combinations of the number $\tilde{m}$ and of values of the kernels $\tilde{\mathscr{E}}_{nm}\left(X;\mathscr{A}(\zeta),\mathscr{C};\,\cdot\,\right)$, with coefficients that do not depend on the connections (signs, factorials and the weights of Section~I.3.6).
By (I.k1) these kernels are analytic in the connections on $\tilde{R}_k\times\mathfrak{X}_k$; hence they are continuous in $\zeta$ on $\mathrm{Box}_k$.
The homomorphism $\mathrm{ev}_{\mathscr{A}(\zeta)+\mathscr{C}}$ of Section~0.4 maps every generator to a linear combination of the generators $\Psi$, $\bar{\Psi}$ whose coefficients are continuous in $\zeta$, by (P.7).
So the coefficients of $I^x_k[\tilde{m},\tilde{\mathscr{E}}](\mathscr{A}(\zeta),\mathscr{C};\Psi)$, which are polynomials in these two sets of functions, and those of $\left\langle\bar{\Psi},D_k\left(\mathscr{A}(\zeta)+\mathscr{C}\right)\Psi\right\rangle$, are continuous:
\begin{equation}\tag{P.8}
\zeta\mapsto\mathrm{Act}^x_k[\tilde{m},\tilde{\mathscr{E}}]\left(\mathscr{A}(\zeta),\mathscr{C};\Psi\right)\ \text{ is a continuous map of }\mathrm{Box}_k\text{ into the even part of }\mathrm{Gr}[\Psi].
\end{equation}

\emph{The integrand of $\mathbf{T}^x_k\mathscr{F}^x_k$.}
It is $\exp(\mathrm{Eq})\cdot\mathscr{F}^x_k[\varepsilon^{\mathrm{v}},m,\mathscr{E}](A_k,\mathscr{C};\Psi)$, where $\mathrm{Eq}$ is the exponent of the Gaussian factor (proof of Proposition~22.2\,(b)); $\mathrm{Eq}$ does not contain $\zeta$.
The map $\zeta\mapsto H_kA_k=A-d\omega_k+H_k\mathrm{Sq}_k\zeta$ is affine and maps $\mathrm{Box}_k$ into $\tilde{R}_k$ (Ob2).
By (P.8) with $(\tilde{m},\tilde{\mathscr{E}})=(m,\mathscr{E})$, by (Gr3) and by Section~0.4, the map $\zeta\mapsto\mathscr{F}^x_k[\varepsilon^{\mathrm{v}},m,\mathscr{E}](A_k,\mathscr{C};\Psi)$ is a continuous map of the compact set $\mathrm{Box}_k$ into the even part of $\mathrm{Gr}[\Psi]$.
Its coefficients are bounded and continuous, hence integrable with respect to $\mu_I$: the map $\mathscr{F}^x_k[\varepsilon^{\mathrm{v}},m,\mathscr{E}](\,\cdot\,,\mathscr{C};\Psi)$ belongs to the class of Section~G3.2.0.
The multiplication by the fixed element $\exp(\mathrm{Eq})$ and the map $\int\mathrm{D}\Psi$ are linear.
Hence
\[
\exp(\mathrm{Eq})\cdot\mathscr{F}^x_k[\varepsilon^{\mathrm{v}},m,\mathscr{E}](A_k,\mathscr{C};\Psi)
\qquad\text{and}\qquad
\int\mathrm{D}\Psi\,\exp(\mathrm{Eq})\cdot\mathscr{F}^x_k[\varepsilon^{\mathrm{v}},m,\mathscr{E}](A_k,\mathscr{C};\Psi)
\]
are continuous functions of $\zeta$ on $\mathrm{Box}_k$, with values in the even parts of $\mathrm{Gr}[\Psi',\Psi]$ and of $\mathrm{Gr}[\Psi']$ respectively, and $\mathbf{T}^x_k\mathscr{F}^x_k[\varepsilon^{\mathrm{v}},m,\mathscr{E}](A',\mathscr{C}';\Psi')$ is defined.

\emph{The element $V^{\mathrm{ex},x}_k$.}
For $k\ge 1$ the numbers $\varepsilon(\mathscr{E})$ and $m(\mathscr{E})$ are complex (Proposition~16.3\,(a)), and $\mathscr{E}^r=R\mathscr{E}$ lies in $K^{x,C}_k$ (Proposition~16.3\,(d)); for $k=0$ one has $\mathscr{E}^r=0$ and $m'=m$.
The four terms of $V^{\mathrm{ex},x}_k$ are the images of the even elements
\[
\mathrm{Act}^x_k[m',\mathscr{E}^r](A+\mathscr{Z},\mathscr{C};\Psi),\qquad
\mathrm{Act}^x_k[m',\mathscr{E}^r](A,\mathscr{C};\Psi),\qquad
I^x_k[m',\mathscr{E}^r](A,\mathscr{C};\Psi)
\]
under the homomorphisms $\left[\bar{\Psi}\to\bar{\Psi}^f,\ \Psi\to\Psi^f\right]$ and $\left[\bar{\Psi}\to\bar{\Psi}^c,\ \Psi\to\Psi^c\right]$, which do not depend on $\mathsf{Z}$ and preserve the parity (Gr2).
Only the first term contains $\mathsf{Z}$, and $\mathsf{Z}\mapsto A+\mathscr{Z}$ is affine and maps $\mathrm{Box}_k$ into $\tilde{R}_k$ (Ob2).
By (P.8) with $(\tilde{m},\tilde{\mathscr{E}})=(m',\mathscr{E}^r)$, the element $V^{\mathrm{ex},x}_k$ is an even element of $\mathrm{Gr}[\Psi',W]$ that is continuous in $\mathsf{Z}$ on $\mathrm{Box}_k$.
So $\exp\left(V^{\mathrm{ex},x}_k\right)$ is continuous in $\mathsf{Z}$ (Gr3), and $\Xi^x_k$ is defined; by (G5)$^{\mathrm{w}}$ the order of the two integrations in $\Xi^x_k$ is immaterial.
\end{proof}

\paragraph{(b) The fields.}

\begin{proof}[Proof of Proposition~22.3\,(b)]
Apply Proposition~21.4\,(a) with $B'=\mathscr{T}'$ and $B=\mathscr{T}$; its supposition holds by (Ob3).
Let $\Phi^c$, $\bar{\Phi}^c$ be the fields (21.9):
\begin{align*}
k\ge 1:&\qquad\Phi^c:=b_L\cdot\Gamma_k(\mathscr{T})Q^T(-\mathscr{T})\Psi'_L,\qquad\bar{\Phi}^c:=b_L\cdot\Gamma_k(\mathscr{T})^TQ^T(\mathscr{T})\bar{\Psi}'_L;\\
k=0:&\qquad\Phi^c:=b_L\cdot\Gamma_0(\mathscr{T})JQ^T(-\mathscr{T})\Psi'_L,\qquad\bar{\Phi}^c:=b_L\cdot\Gamma_0(\mathscr{T})^TJQ^T(\mathscr{T})\bar{\Psi}'_L.
\end{align*}
By (21.11),
\begin{gather*}
H_k(\mathscr{T})\Phi^c=\left(H_{k+1}(\mathscr{T}')\Psi'\right)_L=\psi'_L=\psi^0,\qquad
\Phi^c=T_k(\mathscr{T})H_k(\mathscr{T})\Phi^c=T_k(\mathscr{T})\psi^0=\Psi^c,\\
\bar{H}_k(\mathscr{T})\bar{\Phi}^c=\left(\bar{H}_{k+1}(\mathscr{T}')\bar{\Psi}'\right)_L=\bar{\psi}'_L=\bar{\psi}^0,\qquad
\bar{\Phi}^c=\bar{T}_k(\mathscr{T})\bar{\psi}^0=\bar{\Psi}^c.
\end{gather*}
Therefore
\begin{equation}\tag{P.9}
\Psi^c,\ \bar{\Psi}^c\ \text{ are the fields (21.9) at }B=\mathscr{T},\qquad
\psi^0=H_k(\mathscr{T})\Psi^c,\qquad
\bar{\psi}^0=\bar{H}_k(\mathscr{T})\bar{\Psi}^c .
\end{equation}
By the last line of (21.11), on $\mathrm{Pz}_k$,
\[
\delta_{\alpha,\mathscr{T}}\psi^0=\left(\delta_{\alpha,\mathscr{T}'}H_{k+1}(\mathscr{T}')\Psi'\right)_L=\chi'_L=\chi^0,\qquad
\delta_{\alpha,-\mathscr{T}}\bar{\psi}^0=\left(\delta_{\alpha,-\mathscr{T}'}\bar{H}_{k+1}(\mathscr{T}')\bar{\Psi}'\right)_L=\bar{\chi}'_L=\bar{\chi}^0 .
\]
With the definition of $H_k^{\#}$ (Section~I.3.5) this is the identity $\psi^{0,\#}=H_k^{\#}(\mathscr{T})\Psi^c$ on the unbarred and on the barred fields.
(At $k=0$ one has $H_0=\bar{H}_0=T_0=\bar{T}_0=1$, and $\Psi^c=\psi^0=\left(H_1(\mathscr{T}')\Psi'\right)_L$ and $\bar{\Psi}^c=\bar{\psi}^0=\left(\bar{H}_1(\mathscr{T}')\bar{\Psi}'\right)_L$ are fields on $T^{\mathrm{f}}_0$ with $192$ components over each site of $T^0_0$.)
\end{proof}

\emph{Consequence.}
By Section~0.4 and (P.9), the homomorphism $\mathrm{ev}_{\mathscr{T}}$ followed by $\left[\bar{\Psi}\to\bar{\Psi}^c,\ \Psi\to\Psi^c\right]$ is the homomorphism with $\psi\mapsto\psi^0$, $\bar{\psi}\mapsto\bar{\psi}^0$, $\chi\mapsto\chi^0$, $\bar{\chi}\mapsto\bar{\chi}^0$.
Hence
\begin{equation}\tag{P.10}
I^x_k[m',\mathscr{E}^r](A,\mathscr{C};\Psi^c)=-m'\cdot\langle\bar{\psi}^0,\psi^0\rangle+\sum_X\mathscr{E}^r\left(X;A,\mathscr{C},\psi^{0,\#}\right).
\end{equation}

\paragraph{(c) The formula.}
For $k\ge1$ put $Q^{b}(B):=Q(B)$, and for $k=0$ put $Q^{b}(B):=Q(B)J^{*}$;
these are operators from $\mathfrak{S}(1,\ell)$, resp.\ from $\mathfrak{S}^{\mathrm{ch}}(1,\ell)$, to $\mathfrak{S}(L,\ell)$.
By Section~I.5.7 and (P.5),
\begin{equation}\tag{P.11}
\begin{aligned}
&\left(\mathbf{T}^{x}_{k}\mathscr{F}^{x}_{k}[\varepsilon^{\mathrm{v}},m,\mathscr{E}]\right)(A',\mathscr{C}';\Psi')\\
&\qquad=c_k\int_{\mathrm{Box}_k}\mathrm{d}\mu_I(\zeta)\int\mathrm{D}\Psi\,\exp(\mathrm{Eq})\,\mathscr{F}^{x}_{k}[\varepsilon^{\mathrm{v}},m,\mathscr{E}](A_k,\mathscr{C};\Psi),\\
&\mathrm{Eq}=-b_L\cdot\left\langle\exp\left(-i\cdot e_k\cdot\omega_k\right)\bar{\Psi}'_L-Q^{b}(-\mathscr{T}+d\omega_k)\bar{\Psi},\right.\\
&\qquad\qquad\left.\exp\left(i\cdot e_k\cdot\omega_k\right)\Psi'_L-Q^{b}(\mathscr{T}-d\omega_k)\Psi\right\rangle_L.
\end{aligned}
\end{equation}
Throughout Steps~1--5 the variable $\zeta\in\mathrm{Box}_k$ is held fixed.

\emph{Step 1 (extraction).}
Let $k\ge1$.
For $(\mathscr{A},\mathscr{C})\in\tilde{R}_k\times\mathfrak{X}_k$, Proposition~16.3\,(c) gives in $\mathrm{Gr}[\psi,\chi]$
\[
\sum_X\mathscr{E}(X;\mathscr{A},\mathscr{C})=-\varepsilon(\mathscr{E})\cdot n_k-m(\mathscr{E})\cdot\langle\bar{\psi},\psi\rangle+\sum_X\mathscr{E}^{r}(X;\mathscr{A},\mathscr{C}).
\]
Adding $-m\cdot\langle\bar{\psi},\psi\rangle$ to both sides and applying $\mathrm{ev}_{\mathscr{A}+\mathscr{C}}$, one obtains
\[
I^{x}_{k}[m,\mathscr{E}](\mathscr{A},\mathscr{C};\Psi)=-\varepsilon(\mathscr{E})\cdot n_k+I^{x}_{k}[m',\mathscr{E}^{r}](\mathscr{A},\mathscr{C};\Psi),\qquad m'=m+m(\mathscr{E}).
\]
For $k=0$ the same identity holds with $\mathscr{E}=\mathscr{E}^{r}=0$, $m'=m$, and with $0$ in place of $\varepsilon(\mathscr{E})$;
accordingly, $\varepsilon(\mathscr{E}):=0$ for $k=0$.
Taking $(\mathscr{A},\mathscr{C})=(H_kA_k,\mathscr{C})$, one finds for every $k\ge0$
\begin{equation}\tag{P.12}
\begin{split}
\mathscr{F}^{x}_{k}[\varepsilon^{\mathrm{v}},m,\mathscr{E}](A_k,\mathscr{C};\Psi)
&=\det D_k(0)^{-1}\cdot\exp\left\{\mathrm{Act}^{x}_{k}[m',\mathscr{E}^{r}](H_kA_k,\mathscr{C};\Psi)\right.\\
&\qquad\left.-\left(\varepsilon^{\mathrm{v}}+\varepsilon(\mathscr{E})\right)\cdot n_k+\sum_X\mathfrak{P}_k(X,\mathscr{C})\right\}.
\end{split}
\end{equation}

\emph{Step 2 (the phases).}
Since $\omega_k$ is real and $\ell\mathbb{Z}^3$-periodic, the functions $\boldsymbol{\lambda}=\omega_k$ and $\boldsymbol{\lambda}=-\omega_k$ are of the kind admitted in Proposition~21.1\,(d).
Let $u$ be the function $\exp(i\cdot e_k\cdot\omega_k)$ on the index set of $(\Psi,\bar{\Psi})$, that is,
\[
u(y)=\exp\left(i\cdot e_k\cdot\omega_k(y)\right)\ \ (y\in\mathbb{Z}^3)\quad\text{for }k\ge1,\qquad
u(x)=\exp\left(i\cdot q\cdot\omega_0(c)\right)\quad\text{for }k=0,
\]
where for $k=0$ the same number stands at the $48$ sites $x$ over $c$.
For $G\in\mathrm{Gr}$ let $G^{u}:=G\left[\bar{\Psi}\to u^{-1}\bar{\Psi},\,\Psi\to u\Psi\right]$.
Then (G3)$^{\mathrm{w}}$ gives
\begin{equation}\tag{P.13}
\int\mathrm{D}\Psi\,G=\int\mathrm{D}\Psi\,G^{u}.
\end{equation}

(2a)
By Proposition~21.1\,(d), $D_k(\mathscr{T}^{z}-d\omega_k)\exp(i\cdot e_k\cdot\omega_k)=\exp(i\cdot e_k\cdot\omega_k)\,D_k(\mathscr{T}^{z})$;
hence $D_k(\mathscr{T}^{g})\,u\Psi=u\,D_k(\mathscr{T}^{z})\Psi$.
The factors $u(t)^{-1}$ and $u(t)$, whose product is $1$, stand at one and the same point $t$ of the index set;
therefore
\begin{align*}
k\ge1:\quad&\left(\langle\bar{\Psi},D_k(\mathscr{T}^{g})\Psi\rangle_1\right)^{u}
=\sum_y u(y)^{-1}\bar{\Psi}(y)\cdot u(y)\left(D_k(\mathscr{T}^{z})\Psi\right)(y)
=\langle\bar{\Psi},D_k(\mathscr{T}^{z})\Psi\rangle_1,\\
k=0:\quad&\left(\left\langle\bar{\Psi},\left(D_{\mathscr{T}^{g}}+\mu\right)\Psi\right\rangle^{\mathrm{ch}}\right)^{u}\\
&\quad=\frac{1}{48}\sum_x\exp\left(-i\cdot q\cdot\omega_0(c)\right)\bar{\Psi}(x)\cdot\exp\left(i\cdot q\cdot\omega_0(c)\right)\left(\left(D_{\mathscr{T}^{z}}+\mu\right)\Psi\right)(x)\\
&\quad=\left\langle\bar{\Psi},\left(D_{\mathscr{T}^{z}}+\mu\right)\Psi\right\rangle^{\mathrm{ch}}.
\end{align*}

(2b)
Denote by $G_{\omega_k}$ the automorphism (I.T-G) of $\mathrm{Gr}[\psi,\chi]$ (Section~I.4.5) with $\Lambda=\omega_k$.
Then
\begin{equation}\tag{P.14}
\left(\mathrm{ev}_{\mathscr{T}^{g}}F\right)^{u}=\mathrm{ev}_{\mathscr{T}^{z}}\left(G_{\omega_k}F\right)\qquad\text{for all }F\in\mathrm{Gr}[\psi,\chi].
\end{equation}
As both sides are homomorphisms of algebras in $F$, it is enough to check the identity on the generators.
Proposition~21.1\,(d), used for the operators without bar with $\boldsymbol{\lambda}=\omega_k$, for $\bar{H}_k$, and for the relation for $\delta_{\alpha,\cdot}$ with $(B,\boldsymbol{\lambda})=(\mathscr{T}^{z},\omega_k)$ and with $(B,\boldsymbol{\lambda})=(-\mathscr{T}^{z},-\omega_k)$, gives
\begin{align*}
\left(\mathrm{ev}_{\mathscr{T}^{g}}\psi(x)\right)^{u}
&=\left(H_k(\mathscr{T}^{z}-d\omega_k)\exp\left(i\cdot e_k\cdot\omega_k\right)\Psi\right)(x)
=\exp\left(i\cdot e_k\cdot\omega_k(c)\right)\left(H_k(\mathscr{T}^{z})\Psi\right)(x)\\
&=\mathrm{ev}_{\mathscr{T}^{z}}\left(\exp\left(i\cdot e_k\cdot\omega_k(c)\right)\psi(x)\right),\\
\left(\mathrm{ev}_{\mathscr{T}^{g}}\bar{\psi}(x)\right)^{u}
&=\left(\bar{H}_k(\mathscr{T}^{z}-d\omega_k)\exp\left(-i\cdot e_k\cdot\omega_k\right)\bar{\Psi}\right)(x)
=\exp\left(-i\cdot e_k\cdot\omega_k(c)\right)\left(\bar{H}_k(\mathscr{T}^{z})\bar{\Psi}\right)(x)\\
&=\mathrm{ev}_{\mathscr{T}^{z}}\left(\exp\left(-i\cdot e_k\cdot\omega_k(c)\right)\bar{\psi}(x)\right),\\
\left(\mathrm{ev}_{\mathscr{T}^{g}}\chi(x,x')\right)^{u}
&=\left(\delta_{\alpha,\mathscr{T}^{z}-d\omega_k}\left(\exp\left(i\cdot e_k\cdot\omega_k\right)H_k(\mathscr{T}^{z})\Psi\right)\right)(x,x')\\
&=\exp\left(i\cdot e_k\cdot\omega_k(c)\right)\left(\delta_{\alpha,\mathscr{T}^{z}}H_k(\mathscr{T}^{z})\Psi\right)(x,x')
=\mathrm{ev}_{\mathscr{T}^{z}}\left(\exp\left(i\cdot e_k\cdot\omega_k(c)\right)\chi(x,x')\right),\\
\left(\mathrm{ev}_{\mathscr{T}^{g}}\bar{\chi}(x,x')\right)^{u}
&=\left(\delta_{\alpha,-\mathscr{T}^{z}+d\omega_k}\left(\exp\left(-i\cdot e_k\cdot\omega_k\right)\bar{H}_k(\mathscr{T}^{z})\bar{\Psi}\right)\right)(x,x')\\
&=\exp\left(-i\cdot e_k\cdot\omega_k(c)\right)\left(\delta_{\alpha,-\mathscr{T}^{z}}\bar{H}_k(\mathscr{T}^{z})\bar{\Psi}\right)(x,x')\\
&=\mathrm{ev}_{\mathscr{T}^{z}}\left(\exp\left(-i\cdot e_k\cdot\omega_k(c)\right)\bar{\chi}(x,x')\right),
\end{align*}
where $c$ is the photon site under the first point $x$.
The right-hand sides are the images of the four generators under $\mathrm{ev}_{\mathscr{T}^{z}}\circ G_{\omega_k}$;
the pairs $(x,x')$ range over all of $\mathrm{Pzt}_k$.
(For a pair inside a cluster one has $c=c'$ and no phase $U$ occurs;
both points carry the number $\exp(i\cdot e_k\cdot\omega_k(c))$, and both sides of the identity for $\chi(x,x')$ equal $\eta_k^{-\alpha}\exp(i\cdot e_k\cdot\omega_k(c))\left(f(x')-f(x)\right)$, where $f=H_k(\mathscr{T}^{z})\Psi$.)

(2c)
One has
\[
G_{\omega_k}\langle\bar{\psi},\psi\rangle=\int\mathrm{d}x\,\exp\left(-i\cdot e_k\cdot\omega_k(c)\right)\bar{\psi}(x)\cdot\exp\left(i\cdot e_k\cdot\omega_k(c)\right)\psi(x)=\langle\bar{\psi},\psi\rangle .
\]
Let $k\ge1$.
By (Ob6) the family $\mathscr{E}^{r}$ lies in $K^{x,C}_{k}$, and by (Ob2) the pairs $(H_kA_k,\mathscr{C})=(A+\mathscr{Z}-d\omega_k,\mathscr{C})$ and $(A+\mathscr{Z},\mathscr{C})$ lie in $\tilde{R}_k\times\mathfrak{X}_k$.
The clause (k4) of Section~I.5.4, applied to $\mathscr{E}^{r}$ with $\mathscr{A}=A+\mathscr{Z}$, $\mathscr{A}^{1}=\mathscr{A}-d\omega_k$ on all bonds, $\mathscr{C}^{1}=\mathscr{C}$, $\boldsymbol{\lambda}=\omega_k$ and $\boldsymbol{\lambda}'=0$, gives
\[
G_{\omega_k}\left[\mathscr{E}^{r}(X;H_kA_k,\mathscr{C})\right]=\mathscr{E}^{r}(X;A+\mathscr{Z},\mathscr{C})\qquad\text{for every }X\in D_k .
\]
With (P.14), Section~0.4 and (P.5) it follows that
\begin{align*}
\left(I^{x}_{k}[m',\mathscr{E}^{r}](H_kA_k,\mathscr{C};\Psi)\right)^{u}
&=\mathrm{ev}_{\mathscr{T}^{z}}\left(-m'\cdot\langle\bar{\psi},\psi\rangle+\sum_X\mathscr{E}^{r}(X;A+\mathscr{Z},\mathscr{C})\right)\\
&=I^{x}_{k}[m',\mathscr{E}^{r}](A+\mathscr{Z},\mathscr{C};\Psi).
\end{align*}
At $k=0$ no activity is present, and the relation reduces to $\left(-m\cdot\langle\bar{\Psi},\Psi\rangle^{\mathrm{ch}}\right)^{u}=-m\cdot\langle\bar{\Psi},\Psi\rangle^{\mathrm{ch}}$.
Combined with (2a), this gives for every $k\ge0$
\begin{equation}\tag{P.15}
\left(\mathrm{Act}^{x}_{k}[m',\mathscr{E}^{r}](H_kA_k,\mathscr{C};\Psi)\right)^{u}=\mathrm{Act}^{x}_{k}[m',\mathscr{E}^{r}](A+\mathscr{Z},\mathscr{C};\Psi).
\end{equation}

(2d)
Proposition~21.1\,(d) with $(B,\boldsymbol{\lambda})=(\mathscr{T},\omega_k)$ and with $(B,\boldsymbol{\lambda})=(-\mathscr{T},-\omega_k)$ gives
\begin{gather*}
Q^{b}(\mathscr{T}-d\omega_k)\exp\left(i\cdot e_k\cdot\omega_k\right)=\exp\left(i\cdot e_k\cdot\omega_k\right)Q^{b}(\mathscr{T}),\\
Q^{b}(-\mathscr{T}+d\omega_k)\exp\left(-i\cdot e_k\cdot\omega_k\right)=\exp\left(-i\cdot e_k\cdot\omega_k\right)Q^{b}(-\mathscr{T}),
\end{gather*}
the function on the right-hand sides being evaluated at the sites of $L\mathbb{Z}^3$.
Hence
\begin{align*}
\mathrm{Eq}^{u}
&=-b_L\cdot L^3\sum_{y'\in L\mathbb{Z}^3/\ell\mathbb{Z}^3}\exp\left(-i\cdot e_k\cdot\omega_k(y')\right)\left(\bar{\Psi}'_L-Q^{b}(-\mathscr{T})\bar{\Psi}\right)(y')\\
&\qquad\qquad\cdot\exp\left(i\cdot e_k\cdot\omega_k(y')\right)\left(\Psi'_L-Q^{b}(\mathscr{T})\Psi\right)(y')\\
&=-b_L\cdot\left\langle\bar{\Psi}'_L-Q^{b}(-\mathscr{T})\bar{\Psi},\,\Psi'_L-Q^{b}(\mathscr{T})\Psi\right\rangle_L .
\end{align*}
By (P.12), (P.13), (P.15) and (Gr3),
\begin{equation}\tag{P.16}
\begin{split}
&\int\mathrm{D}\Psi\,\exp(\mathrm{Eq})\,\mathscr{F}^{x}_{k}[\varepsilon^{\mathrm{v}},m,\mathscr{E}](A_k,\mathscr{C};\Psi)\\
&\quad=\det D_k(0)^{-1}\cdot\exp\left\{-\left(\varepsilon^{\mathrm{v}}+\varepsilon(\mathscr{E})\right)\cdot n_k+\sum_X\mathfrak{P}_k(X,\mathscr{C})\right\}\\
&\qquad\cdot\int\mathrm{D}\Psi\,\exp\left\{-b_L\cdot\left\langle\bar{\Psi}'_L-Q^{b}(-\mathscr{T})\bar{\Psi},\,\Psi'_L-Q^{b}(\mathscr{T})\Psi\right\rangle_L+\mathrm{Act}^{x}_{k}[m',\mathscr{E}^{r}](A+\mathscr{Z},\mathscr{C};\Psi)\right\}.
\end{split}
\end{equation}

\emph{Step 3 (splitting).}
At $(A,\mathscr{C})$ the total connection is $\mathscr{T}$, so the definition of $\mathrm{Act}^{x}_{k}$ at $(A,\mathscr{C})$ gives
\begin{gather*}
\mathrm{Act}^{x}_{k}[m',\mathscr{E}^{r}](A+\mathscr{Z},\mathscr{C};\Psi)=-\langle\bar{\Psi},D_k(\mathscr{T})\Psi\rangle+\mathrm{Rs}(\Psi),\\
\mathrm{Rs}(\Psi):=I^{x}_{k}[m',\mathscr{E}^{r}](A,\mathscr{C};\Psi)+\mathrm{Act}^{x}_{k}[m',\mathscr{E}^{r}](A+\mathscr{Z},\mathscr{C};\Psi)-\mathrm{Act}^{x}_{k}[m',\mathscr{E}^{r}](A,\mathscr{C};\Psi).
\end{gather*}
For a pair of fields $(\Phi,\bar{\Phi})$ let $\mathrm{Rs}(\Phi):=\mathrm{Rs}(\Psi)\left[\bar{\Psi}\to\bar{\Phi},\,\Psi\to\Phi\right]$.

\emph{Step 4 (translation; completion of the square).}
The fields $\Psi^{c}$, $\bar{\Psi}^{c}$ take their values in the subalgebra generated by $\Psi'$, $\bar{\Psi}'$.
By (G1)$^{\mathrm{w}}$ one may replace the integrand of the integral over $(\Psi,\bar{\Psi})$ in (P.16) by its image under $\left[\bar{\Psi}\to\bar{\Psi}^{c}+\bar{\Psi},\,\Psi\to\Psi^{c}+\Psi\right]$ without changing the integral.
According to (P.9), the fields $\Psi^{c}$, $\bar{\Psi}^{c}$ coincide with the fields (21.9) taken at $B=\mathscr{T}$.
Proposition~21.4\,(a), display (21.10), taken with $B=\mathscr{T}$, $B'=\mathscr{T}'$ and with the fields $W':=\Psi$, $\bar{W}':=\bar{\Psi}$, states that this substitution maps
\begin{align*}
k\ge1:\quad&\langle\bar{\Psi},D_k(\mathscr{T})\Psi\rangle_1+b_L\cdot\left\langle\bar{\Psi}'_L-Q(-\mathscr{T})\bar{\Psi},\,\Psi'_L-Q(\mathscr{T})\Psi\right\rangle_L\\
&\qquad\longmapsto\ \langle\bar{\Psi}',D_{k+1}(\mathscr{T}')\Psi'\rangle_1+\left\langle\bar{\Psi},\left(D_k(\mathscr{T})+b_L\cdot P(\mathscr{T})\right)\Psi\right\rangle_1,\\
k=0:\quad&\left\langle\bar{\Psi},\left(D_{\mathscr{T}}+\mu\right)\Psi\right\rangle^{\mathrm{ch}}+b_L\cdot\left\langle\bar{\Psi}'_L-Q(-\mathscr{T})J^{*}\bar{\Psi},\,\Psi'_L-Q(\mathscr{T})J^{*}\Psi\right\rangle_L\\
&\qquad\longmapsto\ \langle\bar{\Psi}',D_1(\mathscr{T}')\Psi'\rangle_1+\left\langle\bar{\Psi},\left(D_{\mathscr{T}}+\mu+b_L\cdot JP(\mathscr{T})J^{*}\right)\Psi\right\rangle^{\mathrm{ch}}.
\end{align*}
Put $\mathrm{Og}:=\Gamma_k(\mathscr{T})^{-1}$;
thus $\mathrm{Og}=D_k(\mathscr{T})+b_L\cdot P(\mathscr{T})$ on $\mathfrak{S}(1,\ell)$ for $k\ge1$, and $\mathrm{Og}=F_0(\mathscr{T})=D_{\mathscr{T}}+\mu+b_L\cdot JP(\mathscr{T})J^{*}$ on $\mathfrak{S}^{\mathrm{ch}}(1,\ell)$ for $k=0$.
The element $\langle\bar{\Psi}',D_{k+1}(\mathscr{T}')\Psi'\rangle_1$ is even and contains neither $\Psi$ nor $\bar{\Psi}$.
By (Gr3) and (Gr5), the integral over $(\Psi,\bar{\Psi})$ in (P.16) therefore equals
\begin{equation}\tag{P.17}
\exp\left\{-\langle\bar{\Psi}',D_{k+1}(\mathscr{T}')\Psi'\rangle_1\right\}\int\mathrm{D}\Psi\,\exp\left(-\langle\bar{\Psi},\mathrm{Og}\,\Psi\rangle\right)\exp\left\{\mathrm{Rs}(\Psi^{c}+\Psi)\right\},
\end{equation}
where $\mathrm{Rs}(\Psi^{c}+\Psi)$ is formed with the pair $(\Psi^{c}+\Psi,\bar{\Psi}^{c}+\bar{\Psi})$.

\emph{Step 5 (change of variables).}
The rule (G4)$^{\mathrm{w}}$ is now applied with the family $(\Psi,\bar{\Psi})$ in place of $(W',\bar{W}')$, with $O=\mathrm{Og}$ and with $G=\exp\{\mathrm{Rs}(\Psi^{c}+\Psi)\}$.
Its hypotheses are satisfied:
$G$ lies in the subalgebra generated by $\Psi'$, $\bar{\Psi}'$, $\Psi$, $\bar{\Psi}$ and contains neither $W$ nor $\bar{W}$;
the families $(\Psi,\bar{\Psi})$ and $(W,\bar{W})$ have the same index set and the same weights ($\nu=1$ for $k\ge1$, $\nu=\frac{1}{48}$ on $T^{\mathrm{f}}_0$ for $k=0$);
and $\langle\bar{\Psi},\mathrm{Og}\,\Psi\rangle=\langle\bar{\Psi},\mathrm{Og}\,\Psi\rangle_\nu$.
By the definition of $\delta Z_k$ (Sections~I.3.5 and I.3.9), $\det\mathrm{Og}=\delta Z_k(\mathscr{T})$, a determinant of size $4n_k$ for $k\ge1$ and of size $192n_0$ for $k=0$;
moreover $\mathrm{Og}^{-1}=\Gamma_k(\mathscr{T})$.
The substitution $\left[\bar{\Psi}\to\bar{W},\,\Psi\to\Gamma_k(\mathscr{T})W\right]$ maps the pair $(\Psi^{c}+\Psi,\bar{\Psi}^{c}+\bar{\Psi})$ to $(\Psi^{c}+\Gamma_k(\mathscr{T})W,\bar{\Psi}^{c}+\bar{W})=(\Psi^{f},\bar{\Psi}^{f})$.
Consequently
\begin{equation}\tag{P.18}
\int\mathrm{D}\Psi\,\exp\left(-\langle\bar{\Psi},\mathrm{Og}\,\Psi\rangle\right)\exp\left\{\mathrm{Rs}(\Psi^{c}+\Psi)\right\}=\delta Z_k(\mathscr{T})\cdot\int\mathrm{d}\mu_I(W)\,\exp\left\{\mathrm{Rs}(\Psi^{f})\right\}.
\end{equation}

\emph{Step 6 (collection).}
By the definitions of $\mathrm{Rs}$ and of $V^{\mathrm{ex},x}_{k}$ (at $\mathsf{Z}=\zeta$),
\[
\mathrm{Rs}(\Psi^{f})=I^{x}_{k}[m',\mathscr{E}^{r}](A,\mathscr{C};\Psi^{c})+V^{\mathrm{ex},x}_{k},
\]
and the first summand is even and contains neither $W$, $\bar{W}$ nor $\zeta$.
Inserting (P.16), (P.17) and (P.18) into (P.11), moving the factors that are independent of $\zeta$ in front of the integral over $\zeta$ ((P.2) for the multiplication by a fixed element), and using (Gr5) and (P.4), one arrives, after renaming $\zeta$ as $\mathsf{Z}$, at
\begin{equation}\tag{P.19}
\begin{split}
&\left(\mathbf{T}^{x}_{k}\mathscr{F}^{x}_{k}[\varepsilon^{\mathrm{v}},m,\mathscr{E}]\right)(A',\mathscr{C}';\Psi')
=c_k\cdot\det D_k(0)^{-1}\cdot\delta Z_k(\mathscr{T})\\
&\qquad\cdot\exp\left\{-\langle\bar{\Psi}',D_{k+1}(\mathscr{T}')\Psi'\rangle_1+I^{x}_{k}[m',\mathscr{E}^{r}](A,\mathscr{C};\Psi^{c})\right.\\
&\qquad\qquad\left.-\left(\varepsilon^{\mathrm{v}}+\varepsilon(\mathscr{E})+\delta\varepsilon^{0}_{k}\right)\cdot n_k+\sum_X\mathfrak{P}_k(X,\mathscr{C})\right\}\cdot\Xi^{x}_{k}.
\end{split}
\end{equation}
Since $n_k=\ell^3=L^3\cdot n_{k+1}$, one has $\left(\varepsilon^{\mathrm{v}}+\varepsilon(\mathscr{E})+\delta\varepsilon^{0}_{k}\right)\cdot n_k=\varepsilon^{\mathrm{v},+}\cdot n_{k+1}$.
Proposition~21.1\,(f), read for fields, gives
\[
\langle\bar{\psi}^{0},\psi^{0}\rangle=\langle\bar{\psi}'_L,\psi'_L\rangle^{\mathrm{ch}}\ (\text{level }k)=L\cdot\langle\bar{\psi}',\psi'\rangle^{\mathrm{ch}}\ (\text{level }k+1)=L\cdot\langle\bar{\psi}',\psi'\rangle ;
\]
written out,
\[
\eta_k^3\cdot\frac{1}{48}\sum_{c,g}L^{-2}\,\bar{\psi}'(c/L,g)\cdot\psi'(c/L,g)=L\cdot\eta_{k+1}^3\cdot\frac{1}{48}\sum_{c',g}\bar{\psi}'(c',g)\cdot\psi'(c',g),
\]
because $\eta_k^3\cdot L^{-2}=L\cdot\eta_{k+1}^3$, the label $g$ being left untouched.
So (P.10) reads
\[
I^{x}_{k}[m',\mathscr{E}^{r}](A,\mathscr{C};\Psi^{c})=-m^{+}\cdot\langle\bar{\psi}',\psi'\rangle+\sum_X\mathscr{E}^{r}(X;A,\mathscr{C},\psi^{0,\#}),\qquad m^{+}=L\cdot m'.
\]

\emph{Step 7 (the constant).}
One has $c_k=b^{-4n_{k+1}}$, and by (P.6)
\[
c_k\cdot\det D_k(0)^{-1}=\det D_{k+1}(0)^{-1}\cdot\delta Z_k(0)^{-1}.
\]
Together with Step~6 this turns (P.19) into (2.4a$^{x}$)$^{\mathrm{ch}}$.

\emph{Level $0$, written out.}
Let $k=0$.
Then $e_k=q$, $\mathscr{E}=\mathscr{E}^{r}=0$, $m'=m$, $\varepsilon(\mathscr{E})=0$, $H_0=1$ (photon), $H_0(\cdot)=\bar{H}_0(\cdot)=T_0(\cdot)=\bar{T}_0(\cdot)=1$, and $D_0(B)=D_B+\mu$.
The families $(\Psi,\bar{\Psi})$ and $(W,\bar{W})$ consist of $192n_0$ pairs of generators on $T^{\mathrm{f}}_0$ with the weight $\frac{1}{48}$, and $(\Psi',\bar{\Psi}')$ consists of $4n_1$ pairs on $T^0_1$ with the weight $1$.
For a pair of fields $(\Phi,\bar{\Phi})$ on $T^{\mathrm{f}}_0$ and $(\mathscr{A},\mathscr{C})\in\tilde{R}_0\times\mathfrak{X}_0$ one has
\begin{equation}\tag{P.20}
\begin{aligned}
I^{x}_{0}[m,0](\mathscr{A},\mathscr{C};\Phi)&=-m\cdot\langle\bar{\Phi},\Phi\rangle^{\mathrm{ch}},\\
\mathrm{Act}^{x}_{0}[m,0](\mathscr{A},\mathscr{C};\Phi)&=-\left\langle\bar{\Phi},\left(D_{\mathscr{A}+\mathscr{C}}+\mu\right)\Phi\right\rangle^{\mathrm{ch}}-m\cdot\langle\bar{\Phi},\Phi\rangle^{\mathrm{ch}}.
\end{aligned}
\end{equation}
Put
\[
\mathrm{Nz}_0:=c_0\cdot\det D_0(0)^{-1}\cdot\exp\left\{-\varepsilon^{\mathrm{v}}\cdot n_0+\sum_X\mathfrak{P}_0(X,\mathscr{C})\right\},
\qquad\text{and recall}\qquad
\mathscr{Z}=C\,C_0^{\frac{1}{2}}\zeta .
\]
The chain of Steps~1--7 then reads
\begin{align*}
&\left(\mathbf{T}^{x}_{0}\mathscr{F}^{x}_{0}[\varepsilon^{\mathrm{v}},m,0]\right)(A',\mathscr{C}';\Psi')\\
&\ =\mathrm{Nz}_0\int_{\mathrm{Box}_0}\mathrm{d}\mu_I(\zeta)\int\mathrm{D}\Psi\,\exp\left\{-b_L\cdot\left\langle\exp\left(-i\cdot q\cdot\omega_0\right)\bar{\Psi}'_L-Q(-\mathscr{T}+d\omega_0)J^{*}\bar{\Psi},\right.\right.\\
&\qquad\left.\left.\exp\left(i\cdot q\cdot\omega_0\right)\Psi'_L-Q(\mathscr{T}-d\omega_0)J^{*}\Psi\right\rangle_L
-\left\langle\bar{\Psi},\left(D_{\mathscr{T}+\mathscr{Z}-d\omega_0}+\mu+m\right)\Psi\right\rangle^{\mathrm{ch}}\right\}
\qquad\bigl[\text{(P.11), (P.12), (P.20)}\bigr]\\
&\ =\mathrm{Nz}_0\int_{\mathrm{Box}_0}\mathrm{d}\mu_I(\zeta)\int\mathrm{D}\Psi\,\exp\left\{-b_L\cdot\left\langle\bar{\Psi}'_L-Q(-\mathscr{T})J^{*}\bar{\Psi},\,\Psi'_L-Q(\mathscr{T})J^{*}\Psi\right\rangle_L
-\left\langle\bar{\Psi},\left(D_{\mathscr{T}}+\mu\right)\Psi\right\rangle^{\mathrm{ch}}\right.\\
&\qquad\left.-\left\langle\bar{\Psi},\left(D_{\mathscr{T}+\mathscr{Z}}-D_{\mathscr{T}}\right)\Psi\right\rangle^{\mathrm{ch}}-m\cdot\langle\bar{\Psi},\Psi\rangle^{\mathrm{ch}}\right\}
\qquad\bigl[\text{Steps 2, 3}\bigr]\\
&\ =\mathrm{Nz}_0\cdot\exp\left\{-\langle\bar{\Psi}',D_1(\mathscr{T}')\Psi'\rangle_1\right\}\int_{\mathrm{Box}_0}\mathrm{d}\mu_I(\zeta)\int\mathrm{D}\Psi\,\exp\left\{-\langle\bar{\Psi},F_0(\mathscr{T})\Psi\rangle^{\mathrm{ch}}\right.\\
&\qquad\left.-\left\langle\bar{\Psi}^{c}+\bar{\Psi},\left(D_{\mathscr{T}+\mathscr{Z}}-D_{\mathscr{T}}\right)\left(\Psi^{c}+\Psi\right)\right\rangle^{\mathrm{ch}}-m\cdot\left\langle\bar{\Psi}^{c}+\bar{\Psi},\Psi^{c}+\Psi\right\rangle^{\mathrm{ch}}\right\}
\qquad\bigl[\text{Step 4}\bigr]\\
&\ =\mathrm{Nz}_0\cdot\exp\left\{-\langle\bar{\Psi}',D_1(\mathscr{T}')\Psi'\rangle_1\right\}\cdot\delta Z_0(\mathscr{T})\int_{\mathrm{Box}_0}\mathrm{d}\mu_I(\zeta)\int\mathrm{d}\mu_I(W)\\
&\qquad\cdot\exp\left\{-\left\langle\bar{\Psi}^{f},\left(D_{\mathscr{T}+\mathscr{Z}}-D_{\mathscr{T}}\right)\Psi^{f}\right\rangle^{\mathrm{ch}}-m\cdot\langle\bar{\Psi}^{f},\Psi^{f}\rangle^{\mathrm{ch}}\right\}
\qquad\bigl[\text{Step 5}\bigr],
\end{align*}
where $\mathrm{d}\mu_I(W)=\exp\left(-\langle\bar{W},W\rangle^{\mathrm{ch}}\right)\mathrm{D}W$, $\bar{\Psi}^{f}=\bar{\Psi}^{c}+\bar{W}$, and $\Psi^{f}=\Psi^{c}+\Gamma_0(\mathscr{T})W$ with
\[
\left(\Gamma_0(\mathscr{T})W\right)(x)=\int\mathrm{d}x'\,\Gamma_0(\mathscr{T})(x,x')\,W(x')=\frac{1}{48}\sum_{x'}\Gamma_0(\mathscr{T})(x,x')\,W(x')
\]
(the kernel being taken as in (I.1.4));
moreover $\delta Z_0(\mathscr{T})=\det F_0(\mathscr{T})=\det\left(D_{\mathscr{T}}+\mu+b_L\cdot JP(\mathscr{T})J^{*}\right)$ is a determinant of size $192n_0$.
Now write
\[
-m\cdot\langle\bar{\Psi}^{f},\Psi^{f}\rangle^{\mathrm{ch}}=-m\cdot\langle\bar{\Psi}^{c},\Psi^{c}\rangle^{\mathrm{ch}}-m\cdot\left[\langle\bar{\Psi}^{f},\Psi^{f}\rangle^{\mathrm{ch}}-\langle\bar{\Psi}^{c},\Psi^{c}\rangle^{\mathrm{ch}}\right],
\]
and use
\[
\langle\bar{\Psi}^{c},\Psi^{c}\rangle^{\mathrm{ch}}=\langle\bar{\psi}'_L,\psi'_L\rangle^{\mathrm{ch}}=L\cdot\langle\bar{\psi}',\psi'\rangle,\qquad
\int_{\mathrm{Box}_0}\mathrm{d}\mu_I=\exp\left(-\delta\varepsilon^{0}_{0}\cdot n_0\right)\int\mathrm{d}\hat{\mu},
\]
\[
n_0=L^3\cdot n_1,\qquad c_0\cdot\det D_0(0)^{-1}=\det D_1(0)^{-1}\cdot\delta Z_0(0)^{-1}.
\]
This yields the level-$0$ display of Proposition~22.3\,(c), with
\[
\Xi^{x}_{0}=\int\mathrm{d}\mu_I(W)\,\mathrm{d}\hat{\mu}(\mathsf{Z})\,\exp\left\{-\left\langle\bar{\Psi}^{f},\left(D_{\mathscr{T}+\mathscr{Z}}-D_{\mathscr{T}}\right)\Psi^{f}\right\rangle^{\mathrm{ch}}-m\cdot\left[\langle\bar{\Psi}^{f},\Psi^{f}\rangle^{\mathrm{ch}}-\langle\bar{\Psi}^{c},\Psi^{c}\rangle^{\mathrm{ch}}\right]\right\}.
\]
That the exponent here is $V^{\mathrm{ex},x}_{0}$ follows from (P.20):
\[
\mathrm{Act}^{x}_{0}[m,0](A+\mathscr{Z},\mathscr{C};\Psi^{f})-\mathrm{Act}^{x}_{0}[m,0](A,\mathscr{C};\Psi^{f})=-\left\langle\bar{\Psi}^{f},\left(D_{\mathscr{T}+\mathscr{Z}}-D_{\mathscr{T}}\right)\Psi^{f}\right\rangle^{\mathrm{ch}},
\]
since the terms with $\mu$ and with $m$ cancel, and
\[
I^{x}_{0}[m,0](A,\mathscr{C};\Psi^{f})-I^{x}_{0}[m,0](A,\mathscr{C};\Psi^{c})=-m\cdot\left[\langle\bar{\Psi}^{f},\Psi^{f}\rangle^{\mathrm{ch}}-\langle\bar{\Psi}^{c},\Psi^{c}\rangle^{\mathrm{ch}}\right].
\]

\paragraph{(d)}
\emph{First sentence.}
The right-hand side of (2.4a$^{x}$)$^{\mathrm{ch}}$ is built from the following objects:
the numbers $\det D_{k+1}(0)$, $\delta Z_k(0)$, $m^{+}=L\cdot\left(m+m(\mathscr{E})\right)$, $\varepsilon^{\mathrm{v},+}=L^3\cdot\left(\varepsilon^{\mathrm{v}}+\varepsilon(\mathscr{E})+\delta\varepsilon^{0}_{k}\right)$ and $n_{k+1}$, which contain no block field;
$\mathscr{T}'=\mathscr{A}'+\mathscr{C}'$ and, with it, $D_{k+1}(\mathscr{T}')$, $\psi'$, $\bar{\psi}'$, $\chi'$, $\bar{\chi}'$;
$A=\mathscr{A}'_L$, $\mathscr{C}=\mathscr{C}'_L$, $\mathscr{T}=\mathscr{T}'_L$ and, with them, $\psi^{0,\#}=(\psi',\chi')_L$ together with the barred fields, $\mathscr{E}^{r}(X;A,\mathscr{C},\psi^{0,\#})$, $\mathfrak{P}_k(X,\mathscr{C})$ and $\delta Z_k(\mathscr{T})$;
finally $\Xi^{x}_{k}$, which is formed from $A$, $\mathscr{C}$, $\mathscr{T}$, $\Psi^{c}=T_k(\mathscr{T})\psi^{0}$, $\bar{\Psi}^{c}=\bar{T}_k(\mathscr{T})\bar{\psi}^{0}$, $\Gamma_k(\mathscr{T})$, $m'$, $\mathscr{E}^{r}$, the measures $\mu_I(W)$, $\hat{\mu}$ and $\mathscr{Z}=H_k\,C\,C_k^{\frac{1}{2}}\mathsf{Z}$.
The linear maps $H_k$, $C$, $C_k^{\frac{1}{2}}$, the set $\mathrm{Box}_k$ and the measure $\hat{\mu}$ do not contain $A'$.
Hence $A'$ enters only through $\mathscr{A}'=H_{k+1}A'$, and none of $A^{\min}_k(A')$, $\mathscr{B}_k$, $\omega_k$ occurs.

\emph{Second sentence.}
Fix a real number $t$.
The function $\omega_k$ does not appear on the right-hand side.
On the left-hand side it enters only through the two phases in the display of $\mathbf{T}^{x}_{k}$ (Section~I.5.7);
the background $\mathscr{B}_k=H_kA^{\min}_k(A')$ is unaffected by the replacement of $\omega_k$ by $\omega_k+t$.
In the proofs of (a) and (c) the function $\omega_k$ enters only through the following two properties:
it is a real $\ell\mathbb{Z}^3$-periodic function on the sites of $T_k$, and $H_kA^{\min}_k(A')=(H_{k+1}A')_L-d\omega_k$.
The function $\omega_k+t$ has both properties, since $d(\omega_k+t)=d\omega_k$.
Consequently (2.4a$^{x}$)$^{\mathrm{ch}}$ remains valid when its left-hand side is formed with $\omega_k+t$, the right-hand side being the same;
the two left-hand sides are therefore equal.
\qed

\begin{named}{Remark}[range of the statement]
The assertions (a)--(d) hold for all data of (S22): every $k\ge0$, all sides $\ell$ with $\ell,\ell/L\in\mathrm{Cl}$, every $\mu\in(0,1/L]$, every $e_k$ with $L^{\frac{1}{2}}\cdot e_k\le e_0$, all complex numbers $\varepsilon^{\mathrm{v}}$, $m$, every $\mathscr{E}\in K^{x,C}_{k}$ if $k\ge1$ and $\mathscr{E}=0$ if $k=0$, every family $\mathfrak{P}_k$, every $A'\in\mathrm{Dom}_{k+1}$ and every $\mathscr{C}'\in\mathfrak{X}_{k+1}$.
The assertions (b), (c), (d) are identities between objects defined in Part~I.
The conditions on the constants under which (a)--(d) hold are those of Section~5.
They are conditions on the choice of the constants of Section~I.8 (the conditions on $e_0$ being imposed for all $v\ge-\log e_0$: Proposition~14.2\,(b)),
and none of them involves $k$, $\ell$, $\mu$, the cutoff $N$ (Section~I.2.5), the torus or the twist (Section~I.6.6);
a twist is contained in the harmonic part of $A'$, which is arbitrary.
\end{named}

\subsection*{4. Proposition 22.4 (the exponent of the first step)}

\begin{proof}[Proof of Proposition~22.4\,(a)]
Here $k=0$, $\mathscr{E}=\mathscr{E}^{r}=0$ and $m'=m$.
Apply (P.20) to the pairs $(A+\mathscr{Z},\mathscr{C})$ and $(A,\mathscr{C})$ of $\tilde{R}_0\times\mathfrak{X}_0$ (Proposition~22.3\,(a)), whose total connections are $\mathscr{T}+\mathscr{Z}$ and $\mathscr{T}$, and to the pairs of fields $(\Psi^{f},\bar{\Psi}^{f})$, $(\Psi^{c},\bar{\Psi}^{c})$ on $T^{\mathrm{f}}_0$.
This gives
\begin{align*}
&\mathrm{Act}^{x}_{0}[m,0](A+\mathscr{Z},\mathscr{C};\Psi^{f})-\mathrm{Act}^{x}_{0}[m,0](A,\mathscr{C};\Psi^{f})\\
&\quad=-\left\langle\bar{\Psi}^{f},\left(D_{\mathscr{T}+\mathscr{Z}}+\mu\right)\Psi^{f}\right\rangle^{\mathrm{ch}}-m\cdot\langle\bar{\Psi}^{f},\Psi^{f}\rangle^{\mathrm{ch}}
+\left\langle\bar{\Psi}^{f},\left(D_{\mathscr{T}}+\mu\right)\Psi^{f}\right\rangle^{\mathrm{ch}}+m\cdot\langle\bar{\Psi}^{f},\Psi^{f}\rangle^{\mathrm{ch}}\\
&\quad=-\left\langle\bar{\Psi}^{f},\left(D_{\mathscr{T}+\mathscr{Z}}-D_{\mathscr{T}}\right)\Psi^{f}\right\rangle^{\mathrm{ch}}=\mathrm{Ev}^{(1)},\\
&I^{x}_{0}[m,0](A,\mathscr{C};\Psi^{f})-I^{x}_{0}[m,0](A,\mathscr{C};\Psi^{c})
=-m\cdot\langle\bar{\Psi}^{f},\Psi^{f}\rangle^{\mathrm{ch}}+m\cdot\langle\bar{\Psi}^{c},\Psi^{c}\rangle^{\mathrm{ch}}=\mathrm{Ev}^{(5)}.
\end{align*}
By the definition of $V^{\mathrm{ex},x}_{0}$ in Proposition~22.3, the sum of these two differences is $V^{\mathrm{ex},x}_{0}$.
\end{proof}

\begin{proof}[Proof of Proposition~22.4\,(b)]
At level $0$ one has $\eta_0=1$, the charge is $q$, $\int\mathrm{d}x=\frac{1}{48}\sum_x$ with $x=(c,g)$ ranging over the fermion sites of a fundamental domain, $c_3=c+ge_1$ and $x^{(3)}=(c_3,gs_3)$.
Every integrand below is an $\ell\mathbb{Z}^3$-periodic function of $x$:
it contains one value of $\overline{\mathrm{ff}}$ and one value of $\mathrm{ff}$, both antiperiodic, while $\mathscr{T}$ and $\mathrm{Zf}$ are periodic and $M_3$ depends on the label only.
No two factors are interchanged anywhere below.

\emph{First display.}
By (I.2.3), for every bond function $B$,
\[
\left(D_B\mathrm{ff}\right)(x)=-\sum_{i\le2}M_i(x)\left[\mathrm{ff}(x^{(i)})-\mathrm{ff}(x)\right]-M_3(x)\left[u_B(c,c_3)\,\mathrm{ff}(x^{(3)})-\mathrm{ff}(x)\right].
\]
Neither the terms with $i\le2$ nor the term $M_3(x)\,\mathrm{ff}(x)$ contains $B$.
Hence
\[
\left(\left(D_{\mathscr{T}+\mathrm{Zf}}-D_{\mathscr{T}}\right)\mathrm{ff}\right)(x)=-M_3(x)\left(u_{\mathscr{T}+\mathrm{Zf}}(c,c_3)-u_{\mathscr{T}}(c,c_3)\right)\mathrm{ff}(x^{(3)}),
\]
which is (I.O6);
and by (I.1.6),
\[
u_{\mathscr{T}+\mathrm{Zf}}(c,c_3)=\exp\left\{i\cdot q\cdot\left(\mathscr{T}(c,c_3)+\mathrm{Zf}(c,c_3)\right)\right\}=\exp\left(i\cdot q\cdot\mathrm{Zf}(c,c_3)\right)u_{\mathscr{T}}(c,c_3).
\]
Pairing with $\overline{\mathrm{ff}}$ one obtains
\begin{gather*}
-\left\langle\overline{\mathrm{ff}},\left(D_{\mathscr{T}+\mathrm{Zf}}-D_{\mathscr{T}}\right)\mathrm{ff}\right\rangle^{\mathrm{ch}}=\int\mathrm{d}x\,\mathrm{Fp}(x),\\
\mathrm{Fp}(x):=\overline{\mathrm{ff}}(x)\cdot M_3(x)\left(\exp\left(i\cdot q\cdot\mathrm{Zf}(c,c_3)\right)-1\right)u_{\mathscr{T}}(c,c_3)\,\mathrm{ff}(x^{(3)}).
\end{gather*}
This is (22.1).

\emph{Second display.}
Since $x\ne x^{(3)}$ and $|c-c_3|=1\le M$ ($M$ is a positive integer), the pair $(x,x^{(3)})$ belongs to $\mathrm{Pzt}_0$;
moreover $d_X(x,x^{(3)})=\max(1,\eta_0)=1$, and $U_B(x,x^{(3)})=u_B(c,c_3)$, because for nearest neighbours the six paths reduce to the bond (Section~I.1.8).
By (I.1.7), with $u_{-\mathscr{T}}=u_{\mathscr{T}}^{-1}$,
\[
\mathrm{chf}(x,x^{(3)})=u_{\mathscr{T}}(c,c_3)\,\mathrm{ff}(x^{(3)})-\mathrm{ff}(x),\qquad
\mathrm{chfb}(x,x^{(3)})=u_{\mathscr{T}}(c,c_3)^{-1}\,\overline{\mathrm{ff}}(x^{(3)})-\overline{\mathrm{ff}}(x).
\]
Consequently $\mathrm{ff}(x)+\mathrm{chf}(x,x^{(3)})=u_{\mathscr{T}}(c,c_3)\,\mathrm{ff}(x^{(3)})$;
inserting this, one sees that the first term in the braces of the second expression of (b) equals $\mathrm{Fp}(x)$.
Further, $\overline{\mathrm{ff}}(x)+\mathrm{chfb}(x,x^{(3)})=u_{\mathscr{T}}(c,c_3)^{-1}\,\overline{\mathrm{ff}}(x^{(3)})$, so that the second term in these braces is
\[
u_{\mathscr{T}}(c,c_3)^{-1}\,\overline{\mathrm{ff}}(x^{(3)})\cdot M_3(x^{(3)})\left(\exp\left(-i\cdot q\cdot\mathrm{Zf}(c,c_3)\right)-1\right)\mathrm{ff}(x).
\]
Now consider $\mathrm{Fp}$ at the site $x^{(3)}$.
The photon site under $x^{(3)}$ is $c_3$, and by (I.s1) the neighbour of type $3$ of $x^{(3)}$ is $(x^{(3)})^{(3)}=x$, whose photon site is $c$.
Thus
\[
\mathrm{Fp}(x^{(3)})=\overline{\mathrm{ff}}(x^{(3)})\cdot M_3(x^{(3)})\left(\exp\left(i\cdot q\cdot\mathrm{Zf}(c_3,c)\right)-1\right)u_{\mathscr{T}}(c_3,c)\,\mathrm{ff}(x).
\]
Since bond functions are odd, $\mathrm{Zf}(c_3,c)=-\mathrm{Zf}(c,c_3)$ and $\mathscr{T}(c_3,c)=-\mathscr{T}(c,c_3)$, whence $u_{\mathscr{T}}(c_3,c)=u_{\mathscr{T}}(c,c_3)^{-1}$.
Therefore the second term in the braces is $\mathrm{Fp}(x^{(3)})$.
By (P.3) with $i=3$, $\int\mathrm{d}x\,\mathrm{Fp}(x^{(3)})=\int\mathrm{d}x\,\mathrm{Fp}(x)$.
It follows that
\[
\frac{1}{2}\int\mathrm{d}x\,\left\{\mathrm{Fp}(x)+\mathrm{Fp}(x^{(3)})\right\}=\int\mathrm{d}x\,\mathrm{Fp}(x),
\]
and each of the two halves equals one half of (22.1).

\emph{The vertex.}
Let $\mathrm{ev}$ be the homomorphism of $\mathrm{Gr}[\psi,\chi]$ with $\psi\to\mathrm{ff}$, $\bar{\psi}\to\overline{\mathrm{ff}}$, $\chi\to\mathrm{chf}$, $\bar{\chi}\to\mathrm{chfb}$ ((Gr2); the images are odd and have the antiperiodicity (I.AP)).
By its definition (23.3), stated with Proposition~23.2, $\mathrm{Vx}(\mathrm{Zf},\square)$ equals $\frac{1}{2}\int_{\square}\mathrm{d}x\,\{\,\cdots\}$, where the braces are those of the second expression of (b) with $\psi$, $\bar{\psi}$, $\chi$, $\bar{\chi}$ in place of $\mathrm{ff}$, $\overline{\mathrm{ff}}$, $\mathrm{chf}$, $\mathrm{chfb}$, and where $\int_{\square}\mathrm{d}x$ is $\frac{1}{48}$ times the sum over the $48M^3$ sites $x=(c,g)$ over $\square$.
Hence $\mathrm{ev}\left(\mathrm{Vx}(\mathrm{Zf},\square)\right)$ is the integral over $\square$ of the integrand of the second expression, that is, $\frac{1}{2}\int_{\square}\mathrm{d}x\,\{\mathrm{Fp}(x)+\mathrm{Fp}(x^{(3)})\}$.
An $M$-cube of $\mathbb{L}_0=\mathbb{Z}^3$ is the set of the sites in the cube of side $M$ with a centre $z\in M\mathbb{Z}^3$ (Section~I.1.1).
The integer $M$ is odd, being a power of the odd number $L$.
Hence $M/2\notin\mathbb{Z}$, no site lies on a face of an $M$-cube (Section~I.1.1, ``Faces''), and the $M$-cube centred at $z$ consists of the sites $c\in\mathbb{Z}^3$ with $|c-z|\le\frac{M-1}{2}$ (sup norm).
Each integer has a unique representation $Mj+r$ with $j\in\mathbb{Z}$ and $|r|\le\frac{M-1}{2}$.
Applying this to each coordinate of a site $c\in\mathbb{Z}^3$, one finds exactly one $z\in M\mathbb{Z}^3$ with $|c-z|\le\frac{M-1}{2}$;
so every site lies in exactly one $M$-cube, that is, the $M$-cubes partition $\mathbb{Z}^3$.
As $\ell\in\mathrm{Cl}$, the side $\ell$ is a multiple of $M$;
therefore $\ell\mathbb{Z}^3\subset M\mathbb{Z}^3$, and every translation by an element of $\ell\mathbb{Z}^3$ maps $M$-cubes onto $M$-cubes.
It follows that the $M$-cubes of $T_0$, that is, the images of these sets on the torus, form a partition of $T_0$, and that the sets $\square\times O_{\mathrm{cub}}$ form a partition of $T^{\mathrm{f}}_0$ (Section~I.1.4).
Since the integrand is $\ell\mathbb{Z}^3$-periodic, $\sum_{\square}\int_{\square}=\int$, and $\mathrm{ev}\left(\sum_{\square}\mathrm{Vx}(\mathrm{Zf},\square)\right)$ is the last expression of (b).
\end{proof}

\begin{named}{Remark}
In view of the first display, the last sentence of Proposition~22.4\,(b) is the identity
\[
\mathrm{ev}\left(\sum_{\square}\mathrm{Vx}(\mathrm{Zf},\square)\right)=-\left\langle\overline{\mathrm{ff}},\left(D_{\mathscr{T}+\mathrm{Zf}}-D_{\mathscr{T}}\right)\mathrm{ff}\right\rangle^{\mathrm{ch}}
\]
of Proposition~23.2\,(a), and the assertion there that each of the two halves has the same image is the relation $\int\mathrm{d}x\,\mathrm{Fp}(x^{(3)})=\int\mathrm{d}x\,\mathrm{Fp}(x)$.
\end{named}

\subsection*{5. Conditions on the constants}

The constants are those of Section~I.8, that is, of Proposition~14.2, chosen in the order given there;
the conditions below are conditions on this choice.
List~I gives the conditions on which the arguments of this chapter rest, together with the places where they enter.
List~II gives, for each result applied in this chapter, the conditions under which it holds;
the complete list of the conditions is Proposition~14.2.

\paragraph{I. Conditions entering the arguments of this chapter.}
\mbox{}

\medskip
\noindent
\begin{tabular}{p{0.24\textwidth}p{0.40\textwidth}p{0.22\textwidth}}
\hline
condition & where it enters & line\\
\hline
$b>0$ &
Proposition~22.2\,(a) ($\det(b\cdot1)=b^{4n_{k+1}}=c_k^{-1}$); Proposition~22.3: (Ob4), Step~7 &
(C0): $b=1$\\
$L$ an integer, $L>1$ &
Proposition~22.2\,(b) ($e_k<e_{k+1}\le e_0$); Proposition~22.3: (Ob1) ($e_k<e_{k+1}$, $\mu<1$); the scalings &
(C3); Proposition~14.2\,(b), (L5)\\
$\rho\ge0$; $w_i$ real and $w_i\ge0$ (implied by $w_i>0$); the closure relation (I.1.1) of the standard weights (through (I.2.4)) &
Proposition~22.2\,(c) ($\nabla_0^{*}\nabla_0\ge0$; $\operatorname{Re}(D_0+\mu)\ge\mu$) &
(C0): $\rho\in(0,1)$, $w=(6,6\sqrt{2},6\sqrt{2},6)$\\
$e_0<1$ &
Proposition~22.3: (Ob1) ($v_k>0$, $p_{0,k}>0$, $\mu_I(\mathrm{Box}_k)>0$) &
(C8): first inequality of (A8) ($e_0^{\varepsilon}\le\frac{1}{4}$), with $\varepsilon>0$ (C7)\\
$M$ odd, $M\ge1$ &
Proposition~22.4\,(b): $(x,x^{(3)})$ is a pair of $\mathrm{Pzt}_0$ (second display: $M\ge1$); the $M$-cubes partition $T_0$ and the sets $\square\times O_{\mathrm{cub}}$ partition $T^{\mathrm{f}}_0$ (last sentence: $M$ odd; Section~I.1.1, ``Faces'', and Section~I.1.4) &
(L5), (M1): $M$ is a power of the odd number $L$\\
\hline
\end{tabular}
\medskip

Proposition~22.1 and the first display of Proposition~22.4\,(b) require no condition on the constants, and Proposition~22.2\,(a) requires $b>0$ only.
The second display of Proposition~22.4\,(b) requires $M\ge1$, and the last sentence of Proposition~22.4\,(b) requires $M$ to be odd (last row of the table).
Proposition~22.4\,(a) requires no condition beyond those of Proposition~22.3 at $k=0$, whose objects it concerns.

\paragraph{II. Conditions entering through the results applied.}
These results are applied in Proposition~22.2\,(b) and in Proposition~22.3;
for a result of \cite{QEDI}, the conditions listed are those on which its proof in \cite{QEDI} rests.

\begin{description}
\item[{\cite[Lemma~17.4(a)]{QEDI}.}]
Conditions: $e_k^{\frac{1}{4}}+e_k^{\varepsilon}\le1$ for $e_k\le e_0$; $L\ge5$; $e_0\le1$.
Clauses (b) and (d) of \cite[Lemma~17.4]{QEDI} require in addition fact (g0) of \cite[Section~17.4.1]{QEDI} and $\varepsilon\le\frac{1}{200}$.
Lines: fourth inequality of (A8); (L5); first inequality of (A8).
\item[{\cite[Lemma~17.5]{QEDI}.}]
Conditions: $e_k^{\varepsilon}\le\frac{1}{4}$; $L^{-\frac{3}{4}}\le\frac{1}{2}$; $\varepsilon\le\frac{1}{200}$; $\alpha>0$; $L$ odd and $M$ a power of $L$ (fact (g0) of \cite[Section~17.4.1]{QEDI}: every $M$-cube lies in one $LM$-cube).
Lines: first inequality of (A8); Proposition~14.2\,(f), from (L1); (S1); (C1); (L5), (M1).
\item[{\cite[Lemma~17.6]{QEDI}, without its clause on $\mathrm{Dom}_k$.}]
Conditions: none ($\ell\in\mathrm{Cl}$ is a datum).
\item[{\cite[Proposition~15.30(ii)]{QEDI}.}]
Conditions:
$8\,C_B\cdot p_k\le e_k^{-\frac{3}{4}+3\varepsilon}$ and $8\cdot e_k^{\frac{3}{4}-\frac{7\varepsilon}{2}}\le1$;
$p_{0,k}\le e_k^{-\frac{\varepsilon}{4}}$;
$2\cdot C_{w}^2\cdot c_{\gamma}^2\cdot c_{C}\cdot e_k^{\frac{\varepsilon}{4}}\le1$;
$e_k<1$, $C_{w}\ge C_{P}$, $\gamma_{\mathrm{w}}\le\gamma_{P}$, $L^{\frac{3}{4}}\ge\frac{8}{3}$;
$\varepsilon\le\frac{1}{12}$;
$L$ odd, $L\ge3$ (\cite[Section~15]{QEDI}; the nesting of the $M$-cubes in the $LM$-cubes);
$0\le\alpha<1$ (a hypothesis of \cite[Theorem~15.14]{QEDI}, on which \cite[Proposition~15.30(iii)]{QEDI}, and hence (ii), rests; $C_{P}$ contains $1/(1-\alpha)$);
$\gamma_{\mathrm{w}}>0$ ($c_{\gamma}$ is finite).
\cite[Theorem~15.14]{QEDI} and \cite[Propositions~15.24, 15.29, 15.30]{QEDI} hold without the condition $M\ge M^{P}$ (Annex~B, Section~G3.6.1).
Lines:
third and fourth inequality of (A10);
first inequality of (A1);
(A2);
Proposition~14.2\,(f) (four entries: $e_0\le4^{-200}$; $C_{P}\le C^{P}\le C_{w}$; $\gamma_{\mathrm{w}}\le\gamma^{P}\le\gamma_{P}$; $L^{\frac{3}{4}}\ge110$);
(S1);
(L5);
(C1) and Proposition~14.2\,(f) ($0<\alpha<1$);
Proposition~14.2\,(f), ``Ranges''.
\item[{\cite[Proposition~15.21(i), (ii)]{QEDI} (existence of the basis $\{\Upsilon\}$).}]
Condition: $L$ odd.
Line: (L5).
\item[{\cite[Theorem~14.15(b)]{QEDI} and \cite[Proposition~14.8(ii)]{QEDI}.}]
Conditions: $L$ odd, $L\ge3$ (\cite[Section~14.1.1]{QEDI}).
Line: (L5).
\item[{Proposition~21.4\,(d2).}]
Conditions: those of Proposition~2.8\,(a), namely
$L\ge5$ odd; $M$ a power of $L$; $M>4L$; $M\ge108\pi\cdot C'_{V}\cdot a'$; $6L\cdot e_0^{\varepsilon}\le1$; $4\cdot C'_{V}\cdot a'\cdot e_0^{\varepsilon}\le1$; $0\le\alpha<1$ (the conditions stated with Proposition~2.7 include $\alpha<1$; $a'$ contains $1/(1-\alpha)$) (Annex~B, Section~G3.6.1).
Lines: (L5); (M1); (M2); (M5); (Z2.s1); (Z2.s2); (C1) and Proposition~14.2\,(f) ($0<\alpha<1$).
\item[{Proposition~21.1\,(d), (f), Proposition~21.4\,(a), (b), (c), Proposition~21.5\,(b), (c).}]
Conditions: $L$ odd, $L>1$; $M$ a power of $L$, $L/2<M$; $b>0$; $\rho\ge0$, $w_i>0$; (I.1.1)
(Annex~B, Section~G3.6.1; the conditions $\rho\ne0$, $\rho^2\ne1$, $w_3\ne0$ listed there serve the assertions of Proposition~21.4\,(a) for $k=0$ (Chapter~21a, Step~8 of the proof of Proposition~21.4), namely $H_{0}=\bar{H}_{0}=T_0=\bar{T}_0=1$ and, at $B=0$ with $F_{0}(0)$ invertible, that $\Psi^{c}$ is constant on clusters only if $\Psi'=0$, and $\bar{\Psi}^{c}$ only if $\bar{\Psi}'=0$).
Lines: (L5); (M1); (C0).
\item[{Proposition~16.3\,(a), (c), (d) (for $k\ge1$ only).}]
Conditions: $L$ odd, $L\ge5$; $M$ a power of $L$, $M$ odd, $M\ge L^2$; in (d): the numbers $c_{X}$, $\kappa_{0}$, $K_{0}$, $c_{\mathrm{rb}}$, $c_{\mathrm{fl}}$ of the line (C0)
(Annex~B, Section~G2.9; the condition $\kappa\ge7\kappa_{0}+7$ of the line (C2), which is listed there as well, concerns Proposition~16.3\,(b)).
Lines: (L5); (M1), (L5); (C0).
\end{description}

\paragraph{III. Data, not conditions on the constants.}
$\ell$ and $\ell/L$ lie in $\mathrm{Cl}$; $\mu\in(0,1/L]$; $e_{k+1}=L^{\frac{1}{2}}\cdot e_k\le e_0$.
In Proposition~22.4\,(b) the datum $\ell\in\mathrm{Cl}$ enters through the divisibility of $\ell$ by $M$.

\section*{Chapter 23. The two displays of the step (the substitution at the levels $k\ge 1$, the vertex of level $0$): proofs of Propositions 23.1 and 23.2}

\subsection*{0. Conventions of this chapter; elementary rules used in both proofs}

In this chapter $\mathrm{Vs}$ and $\mathrm{Vs}'$ are complex vector spaces, and $\theta$ is an additive map; in 1.2\,(f) it is one of the four maps of Section~I.4.5, as in Section~G3.4.2.
The letters $\mathrm{hf}$ and $\mathrm{hf}^{\omega}$ stand for auxiliary functions, $(\mathrm{bs}_j)$ for a basis, $\mathrm{id}$ for an idempotent and $\mathrm{id}_0$ for a number (both in 0.4), $\mathrm{od}(t)$ for odd elements of a Grassmann algebra, and $n$, $n'$, $n''$ for vectors of $\mathbb{Z}^3$.
Moreover $\mathrm{sgn}_S$ is a sign, $\mathrm{Vx}'$ and $\mathrm{Vx}''$ are the two halves of $\mathrm{Vx}$, and $\mathrm{dq}_x$, $\mathrm{dq}'_x$, $\mathrm{sz}$ are numbers defined in 2.2\,(d).

\subsubsection*{0.1. Sums}

Let $Z$ be a finite set of photon sites of $\mathbb{L}_k$, and let $\mathrm{hf}$ be a function on $Z\times O_{\mathrm{cub}}$ with values in a complex vector space.
By (I.1.3),
\begin{equation}\tag{i23.1}
\int_Z\mathrm{d}x\,\mathrm{hf}(x)=\frac{\eta_k^3}{48}\sum_{c\in Z}\sum_{g\in O_{\mathrm{cub}}}\mathrm{hf}(c,g).
\end{equation}

\emph{Levels $k\ge 1$.}
Let $y\in\mathbb{Z}^3$.
Since $L$ is odd (Section~I.0.5), so is the number $1/\eta_k=L^k$, and therefore
\[
\Delta_y=\left\{c\in\mathbb{L}_k:\ |c-y|<\frac{1}{2}\right\}=\left\{y+\eta_k n:\ n\in\mathbb{Z}^3,\ |n|\le\frac{L^k-1}{2}\right\}.
\]
Consequently $\#|\Delta_y|=L^{3k}=\eta_k^{-3}$, the set $\Delta_y\times O_{\mathrm{cub}}$ consists of $48\,\eta_k^{-3}$ fermion sites, $\int_{\Delta_y}\mathrm{d}x=1$, and
\begin{equation}\tag{i23.2}
|c-y|\le\frac{1-\eta_k}{2}\qquad\text{for } c\in\Delta_y .
\end{equation}

\emph{Level $0$.}
Here $\eta_0=1$, and $M$, being a power of $L$ (Section~I.0.5), is odd.
An $M$-cube $\square=\{c\in\mathbb{Z}^3:\ |c-z|<M/2\}$, $z\in M\mathbb{Z}^3$, contains $M^3$ photon sites, the set $\square\times O_{\mathrm{cub}}$ consists of $48\,M^3$ fermion sites, and
\[
\int_{\square}\mathrm{d}x\,\mathrm{hf}(x)=\frac{1}{48}\sum_{c\in\square}\sum_{g\in O_{\mathrm{cub}}}\mathrm{hf}(c,g).
\]

\subsubsection*{0.2. Rows; functions with values in a vector space}

Every operator $O$ of this chapter is given by rows: for every row index $t$ and every $\beta\in\{1,\dots,4\}$,
\[
(Of)(t)_{\beta}=\sum_{(t',\beta')}O(t,\beta;t',\beta')\cdot f(t')_{\beta'},
\]
where $t'$ ranges over finitely many column indices (fermion sites of $\mathbb{L}^{\mathrm{f}}_k$, pairs of $\mathrm{Pzt}_k$, or unit sites), and the complex numbers $O(t,\beta;t',\beta')$ are the \emph{entries} of the row at $t$.
Terms of a display with the same column index are collected, and the entries include the weights of the sums, such as $\eta_k^3/48$.
For the operators $H_{k}(B)$ and $\bar{H}_{k}(B)$ of 1.2\,(e) the column indices are the unit sites of a fundamental domain.
The row at $t$ \emph{contains} $f$ at $t'$ if an entry $O(t,\beta;t',\beta')$ occurs in the display.
It \emph{reads} a bond function $B$ on a set of bonds (a set $Y$ through a set of sites) in the sense of \cite[Definition~5.5\,(1)]{QEDI}, which is the sense of the word in (I.f4) of Section~I.1.9:
its entries take equal values at any two bond functions that coincide on these bonds (at any two sets that have the same intersection with this set of sites).

Let $\mathrm{Vs}$ be a complex vector space.
The spaces $\mathfrak{S}^{\mathrm{ch}}(\eta_k,\ell)$ and $\mathfrak{S}(1,\ell)$ and the space of functions of pairs (Sections~I.1.5 and~I.1.7) are formed also with $\mathrm{Vs}^4$ in place of $\mathbb{C}^4$, with the same antiperiodicity, and $O$ acts on them through the same rows.

\begin{named}{(i23.V1)}
If $(O_1f)(t)=(O_2f)(t)$ for all $f$ with values in $\mathbb{C}^4$, then $(O_1f)(t)=(O_2f)(t)$ for all $f$ with values in $\mathrm{Vs}^4$.
\end{named}

\begin{proof}
The function $f$ is determined by its values on a fundamental domain, which is a finite set; hence the components of its values lie in a finite-dimensional subspace of $\mathrm{Vs}$.
Choose a basis $(\mathrm{bs}_j)$ of this subspace and write $f=\sum_j f_j\cdot\mathrm{bs}_j$.
The functions $f_j$ take values in $\mathbb{C}^4$, and, coefficients with respect to a basis being unique, they inherit the antiperiodicity of $f$.
Finally $(O_if)(t)=\sum_j(O_if_j)(t)\cdot\mathrm{bs}_j$ for $i=1,2$, and the assertion follows.
\end{proof}

\begin{named}{(i23.V2)}
Let $\theta:\mathrm{Vs}\to\mathrm{Vs}'$ be additive, and put $\theta\circ f:=\left(t\mapsto\left(\theta\left(f(t)_{\beta}\right)\right)_{\beta}\right)$; this function has the antiperiodicity of $f$, since $\theta(-v)=-\theta(v)$.
If $\theta$ is linear, then
\[
\theta\left((Of)(t)_{\beta}\right)=\left(O(\theta\circ f)\right)(t)_{\beta}.
\]
If $\theta$ is antilinear, then
\[
\theta\left((Of)(t)_{\beta}\right)=\left(\overline{O}(\theta\circ f)\right)(t)_{\beta},
\]
where $\overline{O}$ denotes the operator whose entries are the complex conjugates of those of $O$; on functions with values in $\mathbb{C}^4$ one has $\overline{O}f=\overline{O\,\overline{f}}$, as in (I.O5).
\end{named}

\begin{proof}
In the linear case
\[
\theta\left(\sum O(t,\beta;t',\beta')\cdot f(t')_{\beta'}\right)=\sum O(t,\beta;t',\beta')\cdot\theta\left(f(t')_{\beta'}\right),
\]
and in the antilinear case the left side equals $\sum\overline{O(t,\beta;t',\beta')}\cdot\theta\left(f(t')_{\beta'}\right)$.
\end{proof}

A \emph{field} (Section~G3.1.0) is a function with values in the odd part of a finite-dimensional Grassmann algebra $\mathrm{Gr}$; this is the case $\mathrm{Vs}=\mathrm{Gr}$.
By (I.AP) of Section~I.3.6 the families of generators are functions of this kind.
Namely, $\psi=\left(x\mapsto(\psi(x,\beta,0))_{\beta}\right)$ and $\bar{\psi}=\left(x\mapsto(\psi(x,\beta,1))_{\beta}\right)$ have the antiperiodicity of $\mathfrak{S}^{\mathrm{ch}}(\eta_k,\ell)$, the families $\chi$ and $\bar{\chi}$ have that of the functions of pairs, and $\Psi$, $\bar{\Psi}$ ($k\ge 1$) have that of $\mathfrak{S}(1,\ell)$.

\subsubsection*{0.3. The product $v\cdot m\cdot v'$}

For columns $v$, $v'$ with four entries in an algebra over $\mathbb{C}$ and for a complex $4\times 4$ matrix $m$ let
\[
v\cdot m\cdot v':=\sum_{\beta,\beta'}v_{\beta}\cdot m_{\beta\beta'}\cdot v'_{\beta'},
\]
the factors being taken in this order.

\begin{named}{(i23.R1)}
For complex $4\times 4$ matrices $m_1$, $m_2$,
\[
(m_1\cdot v)\cdot m\cdot(m_2\cdot v')=v\cdot\left(m_1^{T}\cdot m\cdot m_2\right)\cdot v'.
\]
\end{named}

\begin{proof}
The left side is
\[
\sum\left((m_1)_{\beta\gamma}\cdot v_{\gamma}\right)\cdot m_{\beta\beta'}\cdot\left((m_2)_{\beta'\delta}\cdot v'_{\delta}\right)=\sum v_{\gamma}\cdot\left(m_1^{T}\cdot m\cdot m_2\right)_{\gamma\delta}\cdot v'_{\delta}.
\]
\end{proof}

\begin{named}{(i23.R2)}
If every entry of $v$ anticommutes with every entry of $v'$, then
\[
v\cdot m\cdot v'=-v'\cdot m^{T}\cdot v.
\]
\end{named}

\begin{proof}
For all $\beta$, $\beta'$ one has $v_{\beta}\cdot m_{\beta\beta'}\cdot v'_{\beta'}=-v'_{\beta'}\cdot(m^{T})_{\beta'\beta}\cdot v_{\beta}$; sum over $\beta$ and $\beta'$.
\end{proof}

\subsubsection*{0.4. Homomorphisms of Grassmann algebras}

According to Section~I.3.6, a homomorphism (an automorphism) of Grassmann algebras is a homomorphism of rings \emph{with unit}, that is, it maps $1$ to $1$; it is linear over $\mathbb{C}$, except for $\mathbf{K}$ and the compositions with $\mathbf{K}$, which are antilinear.
The condition that $1$ be mapped to $1$ excludes nothing but the zero map.
Indeed, the image of $1$ under an additive and multiplicative map into a Grassmann algebra $\mathrm{Gr}'$ is an element $\mathrm{id}$ with $\mathrm{id}^2=\mathrm{id}$, and its term of degree $0$ is a number $\mathrm{id}_0$ with $\mathrm{id}_0^2=\mathrm{id}_0$.
If $\mathrm{id}_0=0$, then $\mathrm{id}$ is nilpotent and equal to all its powers, so that $\mathrm{id}=0$; as $F=F\cdot 1$, the map sends $F$ to the image of $F$ times $\mathrm{id}$, and it is therefore the zero map.
If $\mathrm{id}_0=1$, then $(1-\mathrm{id})^2=1-\mathrm{id}$, and $1-\mathrm{id}$ has the term $0$ of degree $0$; hence $1-\mathrm{id}=0$.

Let $\mathrm{Gr}$ be the Grassmann algebra with the independent generators $\phi(t)$, $t$ in a finite totally ordered set, and let $\mathrm{Gr}'$ be a Grassmann algebra.

\begin{named}{(i23.H1)}
For odd elements $\mathrm{od}(t)$ of $\mathrm{Gr}'$ there is exactly one homomorphism of algebras with unit $\mathrm{Gr}\to\mathrm{Gr}'$ with $\phi(t)\mapsto\mathrm{od}(t)$.
\end{named}

\begin{proof}
The ordered monomials $\phi(I)$ form a basis of $\mathrm{Gr}$, and $\phi(\emptyset)=1$.
Let the linear map be defined by sending $\phi(I)$ to the product of the elements $\mathrm{od}(t)$, $t\in I$, taken in increasing order, and $1$ to $1$.
Odd elements of $\mathrm{Gr}'$ anticommute and have square $0$.
Therefore the product of the images of $\phi(I)$ and $\phi(I')$ vanishes if $I$ and $I'$ have a common element; otherwise it is brought into increasing order with the same sign as $\phi(I)\cdot\phi(I')$.
Hence the map is multiplicative.
It is unique, since the $\phi(t)$ and $1$ generate $\mathrm{Gr}$.
\end{proof}

\begin{named}{(i23.H2)}
Two ring homomorphisms with unit $\mathrm{Gr}\to\mathrm{Gr}'$ that are both linear, or both antilinear, over $\mathbb{C}$ and that coincide on the generators are equal.
\end{named}

\begin{proof}
They coincide on all monomials, the empty monomial $1$ included, and hence on all linear combinations of monomials.
\end{proof}

\begin{named}{(i23.H3)}
Let the generators be indexed by a covering index set, with $\phi(t+\ell e_{\mu})=-\phi(t)$ as in (I.AP), and let $\mathrm{od}(t)$ be a family of odd elements on the covering index set with $\mathrm{od}(t+\ell e_{\mu})=-\mathrm{od}(t)$.
Then the homomorphism of (i23.H1), formed with $t$ in a fundamental domain, maps $\phi(t)$ to $\mathrm{od}(t)$ for \emph{all} $t$, and it does not depend on the fundamental domain.
\end{named}

\begin{proof}
Write $t=t_0+\ell n$ with $t_0$ in the domain and $n\in\mathbb{Z}^3$.
By (I.AP), $\phi(t)=(-1)^{n_1+n_2+n_3}\phi(t_0)$, and the image of this element is $(-1)^{n_1+n_2+n_3}\mathrm{od}(t_0)=\mathrm{od}(t)$.
Consequently the homomorphism formed with a second domain maps the generators $\phi(t)$, $t$ in the first domain, to $\mathrm{od}(t)$, as does the homomorphism formed with the first domain.
By the uniqueness in (i23.H1) the two homomorphisms coincide.
\end{proof}

\subsubsection*{0.5. Facts on the spin matrices}

Recall from (I.c3) that $\mathsf{C}$ is real and $\mathsf{C}^{T}=\mathsf{C}^{-1}=-\mathsf{C}$, and from (I.c4) that $\varkappa$ is real, $\varkappa^{T}=\varkappa=\varkappa^{-1}$ and $\varkappa\cdot\mathsf{C}=-\mathsf{C}\cdot\varkappa$.
Recall also the relations (I.m2), $\mathsf{C}^{-1}\cdot M_{i}(x)^{T}\cdot\mathsf{C}=M_{i}(x^{(i)})$, and (I.m3), $\overline{M_{i}(x)}=\varkappa^{-1}\cdot M_{i}(x)\cdot\varkappa$, together with (I.m4) and (I.m7), and the identity $(S^{-1})^{T}=\pm\,\mathsf{C}^{-1}\cdot S\cdot\mathsf{C}$ for $S\in\mathrm{Sp}$ (Section~I.4.1).
Three consequences follow.

\begin{named}{(i23.C1)}
$\mathsf{C}^2=-1$.
\end{named}

\begin{proof}
By (I.c3), $\mathsf{C}^{-1}=-\mathsf{C}$.
\end{proof}

\begin{named}{(i23.C2)}
$\mathsf{C}^{-1}\cdot M_{i}(x)\cdot\mathsf{C}=M_{i}(x^{(i)})^{T}$.
\end{named}

\begin{proof}
Transposition of (I.m2), $\mathsf{C}^{-1}\cdot M_{i}(x)^{T}\cdot\mathsf{C}=M_{i}(x^{(i)})$, gives $\mathsf{C}^{T}\cdot M_{i}(x)\cdot(\mathsf{C}^{-1})^{T}=M_{i}(x^{(i)})^{T}$, and $\mathsf{C}^{T}=\mathsf{C}^{-1}$, $(\mathsf{C}^{-1})^{T}=\mathsf{C}$.
\end{proof}

\begin{named}{(i23.C3)}
$\varkappa^{T}\cdot\overline{M_{i}(x)}\cdot\varkappa=M_{i}(x)$.
\end{named}

\begin{proof}
By (I.m3) and (I.c4),
\[
\varkappa^{T}\cdot\overline{M_{i}(x)}\cdot\varkappa=\varkappa\cdot\left(\varkappa^{-1}\cdot M_{i}(x)\cdot\varkappa\right)\cdot\varkappa=M_{i}(x)\cdot\varkappa^2=M_{i}(x).
\]
\end{proof}

The letters $\mathsf{C}$, $\varkappa$ and $S$ also denote the maps $f\mapsto\left(t\mapsto\mathsf{C}\cdot f(t)\right)$, $f\mapsto\left(t\mapsto\varkappa\cdot f(t)\right)$ and $f\mapsto\left(t\mapsto S\cdot f(t)\right)$ on functions of sites, of unit sites and of pairs.
These maps commute with multiplication by scalar functions and with every map that only moves the index $t$.

\subsection*{1. Proposition 23.1 (levels $k\ge 1$: the rows of $T_k$, its two local pieces, the substitution)}

\subsubsection*{1.1. The statement}

The assertion to be proved is Proposition~23.1, with the data, the definition (23.1) of the four operators and the definition (8.3a)$^{\mathrm{ch}}$ of the substitution that are stated before it.
Its data are an integer $k\ge 1$, an element $\ell$ of $\mathrm{Cl}$, a complex number $\mu$, a number $e_k>0$, a complex bond function $B$ on $T_k$ and a union $Y$ of $M$-cubes of $T_k$; in the step map $Y$ is a polymer $\mathbf{Y}$.
At the unit sites of $Y$ the operators $T^{(1)}(B,Y)$ and $T^{(2)}(B,Y)$ are given by (23.1), and $\bar{T}^{(i)}(B,Y)=\mathsf{C}^{-1}\cdot T^{(i)}(-B,Y)\cdot\mathsf{C}$ for $i=1,2$; at the unit sites outside $Y$ all four operators have zero rows.
For a polymer $\mathbf{Y}$, the algebra $\mathrm{Gr}[\psi,\chi,\Psi\,|\,\mathbf{Y}]$ is the subalgebra of $\mathrm{Gr}[\psi,\chi,\Psi]$ generated by all $\psi(\xi)$, all $\chi(\boldsymbol{\zeta})$ and the $\Psi(u)$, $u=(y,\beta,\omega)$, with $y$ a unit site of $\mathbf{Y}$.
The substitution $\mathrm{sub}_{\mathbf{Y},B}\colon\mathrm{Gr}[\psi,\chi,\Psi\,|\,\mathbf{Y}]\to\mathrm{Gr}[\psi,\chi]$ is the homomorphism of algebras that fixes every $\psi(\xi)$ and every $\chi(\boldsymbol{\zeta})$ and maps $\Psi(y)$ and $\bar{\Psi}(y)$, for the unit sites $y$ of $\mathbf{Y}$, according to (8.3a)$^{\mathrm{ch}}$.
Proposition~23.1 asserts the following.
\begin{itemize}
\item[(a)] \emph{Rows of $T_k$; identity.}
For every $f\in\mathfrak{S}^{\mathrm{ch}}(\eta_k,\ell)$ and every unit site $y$ the value $\left(T_k(B)\,f\right)(y)$ is given by (23.2), and at the unit sites of $Y$
\[
T^{(1)}(B,Y)\,f+T^{(2)}(B,Y)\,\delta_{\alpha,B}f=T_k(B)\,f,\qquad\bar{T}^{(1)}(B,Y)\,f+\bar{T}^{(2)}(B,Y)\,\delta_{\alpha,-B}f=\bar{T}_k(B)\,f.
\]
The same holds for fields.
\item[(b)] \emph{What a row contains and what it reads.}
Let $y$ be a unit site of $Y$.
The row of $T^{(1)}(B,Y)$ at $y$ contains $f$ at the sites $x$ over $\Delta_y$ and at the sites $x^{(3)}$ with $c\in\Delta_y$ and $c_3\in Y$; the row of $T^{(2)}(B,Y)$ at $y$ contains $F$ at the pairs $\left(x,x^{(3)}\right)$ with $c\in\Delta_y$ and $c_3\notin Y$.
These sites lie over $Y$, and these pairs are of kind (II) for $Y$, with the first point over $Y$.
In the sense of Section~I.1.9 both rows read $B$ on the bonds contained in $\Delta_y$ and on the bonds $(c,c_3)$ with $c\in\Delta_y$, each of which has an end point in $\Delta_y$; they read $Y$ through the site $y$ and through the sites $c_3$ with $c\in\Delta_y$.
All these sites, pairs and bonds lie within sup-distance $\frac{1+\eta_k}{2}$ of $y$, and the entries of the rows are entire functions of the numbers $B(b)$.
The same holds for $\bar{T}^{(1)}(B,Y)$ and $\bar{T}^{(2)}(B,Y)$.
\item[(c)] \emph{Covariances of the operators} ($i=1,2$).
Let $\boldsymbol{\lambda}$ be as in Proposition~21.1\,(d), and let $\exp\left(i\,e_k\cdot\boldsymbol{\lambda}\right)$ act on functions of pairs by $\left(\exp\left(i\,e_k\cdot\boldsymbol{\lambda}\right)F\right)(x,x'):=\exp\left(i\,e_k\cdot\boldsymbol{\lambda}(c)\right)F(x,x')$.
Then
\begin{gather*}
T^{(i)}\left(B-d\boldsymbol{\lambda},Y\right)\exp\left(i\,e_k\cdot\boldsymbol{\lambda}\right)=\exp\left(i\,e_k\cdot\boldsymbol{\lambda}\right)T^{(i)}(B,Y),\\
\bar{T}^{(i)}\left(B-d\boldsymbol{\lambda},Y\right)\exp\left(-i\,e_k\cdot\boldsymbol{\lambda}\right)=\exp\left(-i\,e_k\cdot\boldsymbol{\lambda}\right)\bar{T}^{(i)}(B,Y).
\end{gather*}
Let $\tilde{g}=(S,a)\in\mathbb{G}_k$ act as in Section~I.4.3, and put $\left(\mathrm{glb}\,f\right)(x):=\left(S^{-1}\right)^{T}f\left(\tilde{g}^{-1}x\right)$, $\left(\mathrm{glb}\,F\right)(x,x'):=\left(S^{-1}\right)^{T}F\left(\tilde{g}^{-1}x,\tilde{g}^{-1}x'\right)$ and, on $\mathfrak{S}(1,\ell)$, $\left(\mathrm{glb}\,v\right)(y):=\left(S^{-1}\right)^{T}v\left(\tilde{g}^{-1}y\right)$.
Then
\[
T^{(i)}\left(\tilde{g}B,\tilde{g}Y\right)\tilde{g}=\tilde{g}\,T^{(i)}(B,Y),\qquad\bar{T}^{(i)}\left(\tilde{g}B,\tilde{g}Y\right)\mathrm{glb}=\mathrm{glb}\,\bar{T}^{(i)}(B,Y).
\]
Under charge conjugation,
\[
T^{(i)}(-B,Y)\cdot\mathsf{C}=\mathsf{C}\cdot\bar{T}^{(i)}(B,Y),\qquad\bar{T}^{(i)}(-B,Y)\cdot\mathsf{C}=\mathsf{C}\cdot T^{(i)}(B,Y),
\]
and, if $\mu$ is real,
\[
\overline{T^{(i)}\left(\overline{B},Y\right)}=\varkappa^{-1}\cdot T^{(i)}(-B,Y)\cdot\varkappa,\qquad\overline{\bar{T}^{(i)}\left(\overline{B},Y\right)}=\varkappa^{-1}\cdot\bar{T}^{(i)}(-B,Y)\cdot\varkappa,
\]
where $\overline{O}$ denotes the operator $f\mapsto\overline{O\,\overline{f}}$.
\item[(d)] \emph{The substitution: locality and parity.}
For $F\in\mathrm{Gr}[\psi,\chi,\Psi\,|\,\mathbf{Y}]$ the element $\mathrm{sub}_{\mathbf{Y},B}\,F$ contains, apart from the generators $\psi(\xi)$ and $\chi(\boldsymbol{\zeta})$ contained in $F$, only generators $\psi(\xi)$ with $\xi$ over $\mathbf{Y}$ and generators $\chi(\boldsymbol{\zeta})$ at pairs of kind (II) for $\mathbf{Y}$ whose first point lies over $\mathbf{Y}$.
Its coefficients are polynomials in the coefficients of $F$ and in entries of the rows of (b) at the unit sites of $\mathbf{Y}$, and these entries read $B$ on bonds that meet $\mathbf{Y}$.
Every generator is mapped by $\mathrm{sub}_{\mathbf{Y},B}$ to a linear combination of generators with the same bar index.
Hence every monomial of $\mathrm{Gr}[\psi,\chi,\Psi\,|\,\mathbf{Y}]$ is mapped to a linear combination of monomials with the same number of unbarred factors and the same number of barred factors, all factors $\psi(\xi)$, $\chi(\boldsymbol{\zeta})$ and $\Psi(u)$ being counted; in particular the difference of these two numbers and the parity of their sum are preserved, and even elements are mapped to even elements.
\item[(e)] \emph{The substitution: evaluation.}
Let $\mathrm{ev}$ be a homomorphism, linear over $\mathbb{C}$, of $\mathrm{Gr}[\psi,\chi]$ into a Grassmann algebra such that $\mathrm{ev}(\chi)=\delta_{\alpha,B}\,\mathrm{ev}(\psi)$ and $\mathrm{ev}(\bar{\chi})=\delta_{\alpha,-B}\,\mathrm{ev}(\bar{\psi})$ at the pairs of kind (II) for $\mathbf{Y}$.
Then $\mathrm{ev}\circ\mathrm{sub}_{\mathbf{Y},B}$ is the homomorphism, linear over $\mathbb{C}$, of $\mathrm{Gr}[\psi,\chi,\Psi\,|\,\mathbf{Y}]$ that coincides with $\mathrm{ev}$ on $\psi$ and $\chi$ and maps $\Psi(y)$ to $\left(T_k(B)\,\mathrm{ev}(\psi)\right)(y)$ and $\bar{\Psi}(y)$ to $\left(\bar{T}_k(B)\,\mathrm{ev}(\bar{\psi})\right)(y)$.
If, moreover, $S_k(B)$ exists and $\mathrm{ev}(\psi)=H_k(B)\,\Phi$, $\mathrm{ev}(\bar{\psi})=\bar{H}_k(B)\,\bar{\Phi}$ for fields $\Phi$, $\bar{\Phi}$ on $T^0_k$, then this homomorphism maps $\Psi(y)$ to $\Phi(y)$ and $\bar{\Psi}(y)$ to $\bar{\Phi}(y)$.
\item[(f)] \emph{The substitution: four covariances.}
For the automorphisms of Section~I.4.5, extended to $\Psi$ as described there,
\begin{align*}
a_{\tilde{g}}\circ\mathrm{sub}_{\tilde{g}\mathbf{Y},\tilde{g}B}&=\mathrm{sub}_{\mathbf{Y},B}\circ a_{\tilde{g}}\quad\text{on }\mathrm{Gr}[\psi,\chi,\Psi\,|\,\tilde{g}\mathbf{Y}],\ \text{for }\tilde{g}\in\mathbb{G}_k;\\
G_{\Lambda}\circ\mathrm{sub}_{\mathbf{Y},B^{1}}&=\mathrm{sub}_{\mathbf{Y},B}\circ G_{\Lambda}\quad\text{for every real function }\Lambda\text{ of }c\text{ and every }B^{1}\\
&\qquad\qquad\text{with }B^{1}=B-d\Lambda\text{ on the bonds that meet }\mathbf{Y};\\
\mathrm{Cg}\circ\mathrm{sub}_{\mathbf{Y},-B}&=\mathrm{sub}_{\mathbf{Y},B}\circ\mathrm{Cg};\\
\mathbf{K}\circ\mathrm{sub}_{\mathbf{Y},-\overline{B}}&=\mathrm{sub}_{\mathbf{Y},B}\circ\mathbf{K}\quad\text{if }\mu\text{ is real}.
\end{align*}
\item[(g)] \emph{The substitution: norm.}
Let the constants be those of Section~I.8, and let $\mu\in(0,1]$, $e_k\in(0,e_0]$ and $B\in\tilde{R}_k+\mathfrak{X}_k$.
For weights $\mathsf{w}_1,\mathsf{w}_2,\mathsf{w}_3>0$ with $C_T\cdot\left(\mathsf{w}_1+\mathsf{w}_2\right)\le\mathsf{w}_3$ and for every $F\in\mathrm{Gr}[\psi,\chi,\Psi\,|\,\mathbf{Y}]$,
\[
\left\|\mathrm{sub}_{\mathbf{Y},B}\,F\right\|_{(\mathsf{w}_1,\mathsf{w}_2)}\le\left\|F\right\|_{(\mathsf{w}_1,\mathsf{w}_2;\mathsf{w}_3)},
\]
where $\mathsf{w}_1$, $\mathsf{w}_2$, $\mathsf{w}_3$ are the weights of $\psi$, $\chi$, $\Psi$ in the norm (I.5.1).
\end{itemize}

\subsubsection*{1.2. Proof of Proposition 23.1}

\emph{Conventions of the proof.}
Throughout, $k\ge 1$.
The bond function $B$ and the set $Y$ are identified with their $\ell\mathbb{Z}^3$-periodic extensions to the bonds, respectively to the sites, of $\mathbb{L}_k$; the extension of $Y$ is a union of $M$-cubes of $\mathbb{L}_k$.
All displays are read on the covering lattices $\mathbb{L}_k$, $\mathbb{L}^{\mathrm{f}}_k$, $\mathrm{Pzt}_k$.
The letter $x=(c,g)$ denotes a fermion site, $x^{(i)}$ and $c_3=c+\eta_k ge_1$ are as in (I.1.2), and $y$ denotes a unit site.
Finally,
\[
b_k=\frac{b\cdot(1-L^{-1})}{1-L^{-k}}>0
\]
(Section~I.3.4).

\emph{Level $0$.}
Proposition~23.1 concerns $k\ge 1$.
At $k=0$ the number $b_0$ and the operators $T^{(1)}$, $T^{(2)}$ are not defined (I.L0.9), one has $T_0=\bar{T}_0=1$ and $\psi=\Psi$ (Section~I.3.9), and the step map contains no substitution (Definition~24.D\,(D4)).
At level $0$ the display of the step that contains rows of $D_A$ is the element $\mathrm{Vx}$ of Proposition~23.2.

\paragraph{Step 0. The displays (23.1) and (8.3a)$^{\mathrm{ch}}$ define what the statement says.}

\textbf{(0-1) Towers, paths, phases.}
Let $0\le j\le k$.
Write $c=\eta_k n$, and write a point $z$ of $L^j\eta_k\mathbb{Z}^3$ as $z=L^j\eta_k n'$, with $n,n'\in\mathbb{Z}^3$.
Since $L^j$ is odd, $c$ lies in $\mathrm{Bl}_{j}(z)$ if and only if $|n-L^jn'|\le(L^j-1)/2$, and for every $n$ there is exactly one such $n'$; thus the blocks of one side form a partition of $\mathbb{L}_k$.
Let now $j<k$, let $z'=L^{j+1}\eta_k n''$ with $n''\in\mathbb{Z}^3$, and let $z$ lie in $\mathrm{Bl}_{j+1}(z')$, that is, $|n'-Ln''|\le(L-1)/2$.
Then, for $c\in\mathrm{Bl}_{j}(z)$,
\[
|n-L^{j+1}n''|\le|n-L^jn'|+L^j\cdot|n'-Ln''|\le\frac{L^j-1}{2}+L^j\cdot\frac{L-1}{2}=\frac{L^{j+1}-1}{2},
\]
so that $\mathrm{Bl}_{j}(z)$ is contained in $\mathrm{Bl}_{j+1}(z')$.

Let $c\in\Delta_y=\mathrm{Bl}_{k}(y)$, and let $y_0=c,y_1,\dots,y_k$ be the centres of Section~I.3.2, so that $c\in\mathrm{Bl}_{j}(y_j)$; the sequence $(y_0,\dots,y_k)$ is called the \emph{tower} of $c$.
Then $y_k=y$.
The block $\mathrm{Bl}_{j}(y_j)$ is contained in the block of side $L^{j+1}\eta_k$ that contains $y_j$; as this block contains $c$, it is $\mathrm{Bl}_{j+1}(y_{j+1})$.
Hence
\[
y_j\in\mathrm{Bl}_{j}(y_j)\subset\mathrm{Bl}_{j+1}(y_{j+1})\subset\dots\subset\mathrm{Bl}_{k}(y)=\Delta_y ,
\]
and all the $y_j$ lie in $\Delta_y$.

Every site of the path $\boldsymbol{\Gamma}^{\pi}(y_{j+1},y_j)$ has, in each coordinate, a value between the corresponding values of $y_{j+1}$ and $y_j$, because the coordinates are brought to their final values one after the other (Section~I.1.1).
As $\Delta_y$ is a product of three intervals of $\eta_k\mathbb{Z}$, all bonds of these paths are contained in $\Delta_y$.
By (I.3.3), Section~I.1.1 and (I.1.6),
\begin{equation}\tag{i23.3}
\begin{aligned}
\varphi_B(y,c)&=\exp\left(ie_k\eta_k\sum_{j<k}\frac{1}{6}\sum_{\pi}\sum_{b\in\boldsymbol{\Gamma}^{\pi}(y_{j+1},y_j)}B(b)\right),\\
u_B(c,c_3)&=\exp\left(ie_k\eta_k\,B(c,c_3)\right),
\end{aligned}
\end{equation}
where the term of index $j$ vanishes if $y_{j+1}=y_j$ (Section~I.1.1: $(\tau B)(c,c):=0$).
Both quantities are exponentials of linear forms in the numbers $B(b)$; the first form contains only bonds $b$ that are contained in $\Delta_y$, the second only the bond $(c,c_3)$.
In particular $\varphi_{-B}=\varphi_B^{-1}$ and $u_{-B}=u_B^{-1}$.

\textbf{(0-2) $\Delta_y$ lies in $Y$ for every unit site $y$ of $Y$.}
Every unit block lies in one $M$-cube \cite[Section~17.4.1, (g0)]{QEDI}.
Indeed, let $z\in M\mathbb{Z}^3$ be the centre of the $M$-cube of $y$; since $M$ is odd and $y-z\in\mathbb{Z}^3$, one has $|y-z|\le(M-1)/2$, and by (i23.2)
\[
|c-z|\le\frac{M-1}{2}+\frac{1-\eta_k}{2}<\frac{M}{2}\qquad\text{for } c\in\Delta_y .
\]
The $M$-cube of $y$ is one of the cubes of $Y$.
Hence, for $y\in Y$ and $x$ over $\Delta_y$, the sites $x$, $x^{(0)}$, $x^{(1)}$, $x^{(2)}$, which lie over $c$ (I.s4), lie over $Y$, and the condition ``$x^{(i)}$ over $\mathbf{Y}$'' of (I.5.4) holds for $i\le 2$; for $i=3$ it is the condition $c_3\in Y$.
Since $U_B(x,x^{(i)})=1$ for $i\le 2$ and $U_B(x,x^{(3)})=u_B(c,c_3)$ (Section~I.1.8), for $y\in Y$ the display (23.1) is the display (I.5.4) with $\mathbf{Y}=Y$.

\textbf{(0-3) The displays as sums.}
By (i23.1), (I.1.2) and (I.2.1), for every unit site $y$ of $Y$,
\begin{equation}\tag{i23.4}
\begin{split}
\left(T^{(1)}(B,Y)f\right)(y)={}&\frac{\eta_k^3}{48}\sum_{c\in\Delta_y}\varphi_B(y,c)\sum_{g\in O_{\mathrm{cub}}}\Bigl\{\left(1+\frac{\mu}{b_k}\right)f(c,g)\\
&-(b_k\eta_k)^{-1}\Bigl[\sum_{i=0}^{2}\frac{w_i}{2}\left(\rho+\gamma\cdot(g\hat n_i)\right)\cdot\left(f(c,gs_i)-f(c,g)\right)\\
&+\mathbf{1}[c+\eta_k ge_1\in Y]\cdot 3\left(\rho-\gamma\cdot(ge_1)\right)\\
&\qquad\times\left(u_B(c,c+\eta_k ge_1)\cdot f(c+\eta_k ge_1,gs_3)-f(c,g)\right)\Bigr]\Bigr\},
\end{split}
\end{equation}
\begin{equation}\tag{i23.5}
\begin{split}
\left(T^{(2)}(B,Y)F\right)(y)={}&-b_k^{-1}\cdot\eta_k^{\alpha-1}\cdot\frac{\eta_k^3}{48}\sum_{c\in\Delta_y}\varphi_B(y,c)\sum_{g\in O_{\mathrm{cub}}}\mathbf{1}[c+\eta_k ge_1\notin Y]\\
&\times 3\left(\rho-\gamma\cdot(ge_1)\right)\cdot F\left((c,g),(c+\eta_k ge_1,gs_3)\right).
\end{split}
\end{equation}
Each of these sums has $48\,\eta_k^{-3}$ terms $(c,g)$.
The three differences $f(c,gs_i)-f(c,g)$, $i=0,1,2$, are differences inside the cluster of $c$; they carry no phase and are written with $f$.
The pair $(x,x^{(3)})=\left((c,g),(c+\eta_k ge_1,gs_3)\right)$ consists of two distinct points with
\[
d_X(x,x^{(3)})=\max\left(|c-c_3|,\eta_k\right)=\eta_k\le M
\]
(I.1.5); therefore it belongs to $\mathrm{Pzt}_k$, and $F(x,x^{(3)})$ is defined.

\textbf{(0-4) The values are antiperiodic.}
Let $v=\ell e_{\mu}$.
Since $\ell$ is an integer and $L^j\eta_k=L^{j-k}$ is the reciprocal of an integer for $0\le j\le k$, the vector $v$ lies in $\mathbb{Z}^3$ and in $L^j\eta_k\mathbb{Z}^3$.
The translation by $v$ maps $\mathbb{L}_k$ onto itself, the block $\mathrm{Bl}_{j}(z)$ onto $\mathrm{Bl}_{j}(z+v)$, the tower of $c$ onto the tower of $c+v$, and $\boldsymbol{\Gamma}^{\pi}(y_{j+1},y_j)$ onto $\boldsymbol{\Gamma}^{\pi}(y_{j+1}+v,y_j+v)$.
As $B$ is $\ell\mathbb{Z}^3$-periodic, (i23.3) yields
\[
\varphi_B(y+v,c+v)=\varphi_B(y,c),\qquad u_B(c+v,c_3+v)=u_B(c,c_3).
\]
Moreover $(x+v)^{(i)}=x^{(i)}+v$ (I.s1), the matrices in (i23.4) and (i23.5) depend on the label only, $Y+v=Y$, and
\[
f(x+v)=-f(x),\qquad F(x+v,x^{(3)}+v)=-F(x,x^{(3)}).
\]
Replacing $c$ by $c+v$ in (i23.4) and (i23.5) one obtains, for $y\in Y$,
\begin{gather*}
\left(T^{(1)}(B,Y)f\right)(y+v)=-\left(T^{(1)}(B,Y)f\right)(y),\\
\left(T^{(2)}(B,Y)F\right)(y+v)=-\left(T^{(2)}(B,Y)F\right)(y),
\end{gather*}
and $y+v$ lies in $Y$ if and only if $y$ does.
Therefore, with zero rows at the unit sites outside $Y$, the operator $T^{(1)}(B,Y)$ maps $\mathfrak{S}^{\mathrm{ch}}(\eta_k,\ell)$ to $\mathfrak{S}(1,\ell)$, and $T^{(2)}(B,Y)$ maps the functions of pairs to $\mathfrak{S}(1,\ell)$; both are linear.
(The three spaces are the spaces of antiperiodic functions of Section~I.1.9 on the index sets $\mathbb{L}^{\mathrm{f}}_k$, $\mathrm{Pzt}_k$, $\mathbb{Z}^3$.)
Since the rows outside $Y$ vanish, $|T^{(i)}(B,Y)|_{1,\infty}$ ((I.ke3), Section~I.5.2) is the supremum over the rows at the unit sites of $Y$.
By 0.2 the same holds for functions with values in $\mathrm{Vs}^4$.

\textbf{(0-5) The barred operators.}
By definition $\bar{T}^{(i)}(B,Y)=\mathsf{C}^{-1}\cdot T^{(i)}(-B,Y)\cdot\mathsf{C}$, with $\mathsf{C}$ acting on the values (see 0.5); hence these operators map the same spaces to $\mathfrak{S}(1,\ell)$ and have zero rows outside $Y$.
Insert (23.1) at $-B$, and use $\varphi_{-B}=\varphi_B^{-1}$, $u_{-B}=u_B^{-1}$ (0-1) together with $\mathsf{C}^{-1}\cdot M_{i}(x)\cdot\mathsf{C}=M_{i}(x^{(i)})^{T}$ (i23.C2).
For every unit site $y$ of $Y$ this gives
\begin{equation}\tag{i23.6}
\begin{aligned}
\left(\bar{T}^{(1)}(B,Y)f\right)(y)={}&\int_{\Delta_y}\mathrm{d}x\,\varphi_B(y,c)^{-1}\Bigl\{\left(1+\frac{\mu}{b_k}\right)f(x)\\
&-(b_k\eta_k)^{-1}\Bigl[\sum_{i\le 2}M_{i}(x^{(i)})^{T}\cdot\left(f(x^{(i)})-f(x)\right)\\
&+\mathbf{1}[c_3\in Y]\cdot M_{3}(x^{(3)})^{T}\cdot\left(u_B(c,c_3)^{-1}\cdot f(x^{(3)})-f(x)\right)\Bigr]\Bigr\},\\
\left(\bar{T}^{(2)}(B,Y)F\right)(y)={}&-b_k^{-1}\cdot\eta_k^{\alpha-1}\int_{\Delta_y}\mathrm{d}x\,\mathbf{1}[c_3\notin Y]\cdot\varphi_B(y,c)^{-1}\cdot M_{3}(x^{(3)})^{T}\cdot F(x,x^{(3)}),
\end{aligned}
\end{equation}
with
\[
M_{i}(x^{(i)})^{T}=\frac{w_i}{2}\left(\rho-\sum_{\nu}(g\hat n_i)_{\nu}\cdot\gamma_{\nu}^{T}\right),\qquad M_{3}(x^{(3)})^{T}=3\left(\rho+\sum_{\nu}(ge_1)_{\nu}\cdot\gamma_{\nu}^{T}\right).
\]
Indeed, (I.2.1) at the label $gs_i$ of $x^{(i)}$ gives
\[
M_{i}(x^{(i)})=\frac{w_i}{2}\left(\rho+\gamma\cdot(gs_i\hat n_i)\right)=\frac{w_i}{2}\left(\rho-\gamma\cdot(g\hat n_i)\right),
\]
because $s_i\hat n_i=-\hat n_i$ (Section~I.1.2); moreover $w_3/2=3$ and $\hat n_3=-e_1$.

\textbf{(0-6) The substitution exists and is unique.}
Let $\mathbf{Y}$ be a polymer.
The algebra $\mathrm{Gr}[\psi,\chi,\Psi\,|\,\mathbf{Y}]$ is the Grassmann algebra with the independent generators $\psi(\xi)$, $\chi(\boldsymbol{\zeta})$ and $\Psi(u)$, $u=(y,\beta,\omega)$, $y$ a unit site of $\mathbf{Y}$ (indices in a fundamental domain).
Consider the columns
\begin{equation}\tag{i23.7}
\begin{aligned}
\mathrm{sub}_{\mathbf{Y},B}\Psi(y)&:=\left(T^{(1)}(B,\mathbf{Y})\psi\right)(y)+\left(T^{(2)}(B,\mathbf{Y})\chi\right)(y),\\
\mathrm{sub}_{\mathbf{Y},B}\bar{\Psi}(y)&:=\left(\bar{T}^{(1)}(B,\mathbf{Y})\bar{\psi}\right)(y)+\left(\bar{T}^{(2)}(B,\mathbf{Y})\bar{\chi}\right)(y),
\end{aligned}
\end{equation}
in which the operators are applied to the functions $\psi$, $\chi$, $\bar{\psi}$, $\bar{\chi}$ of 0.2, with $\mathrm{Vs}=\mathrm{Gr}[\psi,\chi]$.
Their entries are linear combinations of generators of $\mathrm{Gr}[\psi,\chi]$, and hence odd.
By (0-4) and (0-5) they change sign under $y\mapsto y+\ell e_{\mu}$, as $\Psi(y)$ and $\bar{\Psi}(y)$ do.
The generators $\psi(\xi)$, $\chi(\boldsymbol{\zeta})$ are odd.
By (i23.H1) and (i23.H3) there is exactly one homomorphism
\[
\mathrm{sub}_{\mathbf{Y},B}:\mathrm{Gr}[\psi,\chi,\Psi\,|\,\mathbf{Y}]\to\mathrm{Gr}[\psi,\chi]
\]
that fixes $\psi(\xi)$ and $\chi(\boldsymbol{\zeta})$ and maps $\Psi(y)$, $\bar{\Psi}(y)$ to the columns (i23.7) for all unit sites $y$ of $\mathbf{Y}$.

The homomorphism $\mathrm{sub}_{\mathbf{Y},B}$ contains the four operators only through their rows at the unit sites of $\mathbf{Y}$.
The displays of Definition~24.D contain $T^{(i)}$ and $\bar{T}^{(i)}$ only through the homomorphisms $\mathrm{sub}_{\mathbf{Y},\mathscr{T}}$ of (D4), namely
\[
E^{(i),x}(\mathbf{Y})=\mathrm{sub}_{\mathbf{Y},\mathscr{T}}\left(E^{(i),*}(\mathbf{Y})\right),\qquad i\le 3,\quad\mathbf{Y}\in D_k .
\]
Consequently the step map depends on these operators only through these rows.

\textbf{(0-7) The relations for the phases.}
Proposition~21.1\,(d0) applies to the phases $\varphi_B(y,c)$, $c\in\Delta_y$, and to $u_B(c,c_3)=U_B(c,c_3)$ (nearest neighbours, Section~I.1.8).
For every complex function $\boldsymbol{\lambda}$ on $\mathbb{L}_k$ such that $d\boldsymbol{\lambda}$ is $\ell\mathbb{Z}^3$-periodic, and for every $\tilde{g}\in\mathbb{G}_k$, it gives
\begin{equation}\tag{ph-G}
\begin{aligned}
\varphi_{B-d\boldsymbol{\lambda}}(y,c)&=\exp\left(ie_k\left(\boldsymbol{\lambda}(y)-\boldsymbol{\lambda}(c)\right)\right)\cdot\varphi_B(y,c),\\
u_{B-d\boldsymbol{\lambda}}(c,c_3)&=\exp\left(ie_k\left(\boldsymbol{\lambda}(c)-\boldsymbol{\lambda}(c_3)\right)\right)\cdot u_B(c,c_3),
\end{aligned}
\end{equation}
\begin{equation}\tag{ph-g}
\varphi_{\tilde{g}B}(\tilde{g}y,\tilde{g}c)=\varphi_B(y,c),\qquad u_{\tilde{g}B}(\tilde{g}c,\tilde{g}c_3)=u_B(c,c_3),
\end{equation}
\begin{equation}\tag{ph-K}
\overline{\varphi_{\overline{B}}(y,c)}=\varphi_{-B}(y,c),\qquad\overline{u_{\overline{B}}(c,c_3)}=u_{-B}(c,c_3).
\end{equation}
These relations can also be read off from (i23.3).
For (ph-G): on a bond, $(d\boldsymbol{\lambda})(c',c'')=\eta_k^{-1}\left(\boldsymbol{\lambda}(c'')-\boldsymbol{\lambda}(c')\right)$, so that $\eta_k$ times the sum of $d\boldsymbol{\lambda}$ along a path equals the value of $\boldsymbol{\lambda}$ at the end point minus its value at the initial point.
Hence the exponent of $\varphi$ changes by
\[
-ie_k\sum_{j<k}\left(\boldsymbol{\lambda}(y_j)-\boldsymbol{\lambda}(y_{j+1})\right)=-ie_k\left(\boldsymbol{\lambda}(c)-\boldsymbol{\lambda}(y)\right),
\]
and the exponent of $u$ by $-ie_k\left(\boldsymbol{\lambda}(c_3)-\boldsymbol{\lambda}(c)\right)$.
For (ph-g): $\tilde{g}c=rc+a$, where $r$ is diagonal with entries $\pm 1$ and $a\in M\mathbb{Z}^3$.
This map is an isometry of the sup norm and maps $L^j\eta_k\mathbb{Z}^3$ onto itself.
It therefore maps $\mathrm{Bl}_{j}(z)$ onto $\mathrm{Bl}_{j}(\tilde{g}z)$, the tower of $c$ onto the tower of $\tilde{g}c$, and, $r$ being diagonal, $\boldsymbol{\Gamma}^{\pi}(c',c'')$ onto $\boldsymbol{\Gamma}^{\pi}(\tilde{g}c',\tilde{g}c'')$ with the same order $\pi$; moreover $(\tilde{g}B)(\tilde{g}b)=B(b)$.
For (ph-K): the exponents are $ie_k\eta_k$ times linear forms in $B$ with real coefficients, and $e_k\eta_k$ is real.

\paragraph{(a) Rows of $T_k$; the identity.}

\begin{proof}[Proof of Proposition~23.1\,(a)]
\textbf{(a-1) The display (23.2).}
By Section~I.3.3,
\[
\left(Q_k(B)\cdot J^{*}f\right)(y)=\int_{\Delta_y}\mathrm{d}x\,\varphi_B(y,c)\cdot f(x)=\frac{\eta_k^3}{48}\sum_{c\in\Delta_y}\sum_{g\in O_{\mathrm{cub}}}\varphi_B(y,c)\cdot f(c,g),
\]
and $\left(J\cdot Q_k^{T}(-B)v\right)(x)=\varphi_B(y,c)^{-1}\cdot v(y)$ for $x$ over $\Delta_y$.
Hence
\[
\left(Q_k(B)\cdot J^{*}\cdot J\cdot Q_k^{T}(-B)v\right)(y)=v(y)\int_{\Delta_y}\mathrm{d}x=v(y).
\]
As $P_k(B)=Q_k^{T}(-B)\cdot Q_k(B)$, it follows that $Q_k(B)\cdot J^{*}\cdot J\cdot P_k(B)\cdot J^{*}=Q_k(B)\cdot J^{*}$, so that the two expressions for $T_k(B)$ in Section~I.3.5 coincide:
\[
T_k(B)=b_k^{-1}\cdot Q_k(B)\cdot J^{*}\cdot\left(D_B+\mu\right)+Q_k(B)\cdot J^{*}.
\]
Therefore
\[
\left(T_k(B)f\right)(y)=\int_{\Delta_y}\mathrm{d}x\,\varphi_B(y,c)\cdot\left\{b_k^{-1}\cdot(D_Bf)(x)+\frac{\mu}{b_k}f(x)+f(x)\right\}.
\]
Inserting (I.2.3),
\[
(D_Bf)(x)=-\eta_k^{-1}\sum_{i=0}^{3}M_{i}(x)\cdot\left[U_B(x,x^{(i)})\cdot f(x^{(i)})-f(x)\right],
\]
one obtains (23.2), for every unit site $y$.

\textbf{(a-2) The unbarred identity.}
By Section~I.1.8, $U_B(x,x^{(i)})=1$ for $i=0,1,2$ and $U_B(x,x^{(3)})=u_B(c,c_3)$.
The sum over $i$ in (23.2) is therefore
\begin{multline*}
\sum_{i\le 2}M_{i}(x)\cdot\left(f(x^{(i)})-f(x)\right)\\
+\left(\mathbf{1}[c_3\in Y]+\mathbf{1}[c_3\notin Y]\right)\cdot M_{3}(x)\cdot\left(u_B(c,c_3)\cdot f(x^{(3)})-f(x)\right).
\end{multline*}
Let $y$ be a unit site of $Y$.
The term $(1+\mu/b_k)f(x)$, the terms with $i\le 2$ and the term with $\mathbf{1}[c_3\in Y]$ together make up $\left(T^{(1)}(B,Y)f\right)(y)$.
By (I.1.7) and $d_X(x,x^{(3)})=\eta_k$,
\[
u_B(c,c_3)\cdot f(x^{(3)})-f(x)=\eta_k^{\alpha}\cdot\left(\delta_{\alpha,B}f\right)(x,x^{(3)}).
\]
So the term with $\mathbf{1}[c_3\notin Y]$ contributes
\begin{multline*}
-(b_k\eta_k)^{-1}\cdot\eta_k^{\alpha}\int_{\Delta_y}\mathrm{d}x\,\mathbf{1}[c_3\notin Y]\cdot\varphi_B(y,c)\cdot M_{3}(x)\cdot\left(\delta_{\alpha,B}f\right)(x,x^{(3)})\\
=\left(T^{(2)}(B,Y)\,\delta_{\alpha,B}f\right)(y).
\end{multline*}
Hence $T^{(1)}(B,Y)f+T^{(2)}(B,Y)\,\delta_{\alpha,B}f=T_k(B)f$ at the unit sites of $Y$.

\textbf{(a-3) The barred identity.}
The display (I.1.7) multiplies the values of $f$ by scalars; hence $\delta_{\alpha,-B}(\mathsf{C}f)=\mathsf{C}\,\delta_{\alpha,-B}f$.
By the definition of $\bar{T}^{(i)}$, by (a-2) applied to the bond function $-B$ and to the function $\mathsf{C}f$, and by $\bar{T}_k(B)=\mathsf{C}^{-1}\cdot T_k(-B)\cdot\mathsf{C}$ (Section~I.3.5), one has at the unit sites of $Y$
\[
\begin{aligned}
\bar{T}^{(1)}(B,Y)f+\bar{T}^{(2)}(B,Y)\,\delta_{\alpha,-B}f&=\mathsf{C}^{-1}\left[T^{(1)}(-B,Y)(\mathsf{C}f)+T^{(2)}(-B,Y)\,\delta_{\alpha,-B}(\mathsf{C}f)\right]\\
&=\mathsf{C}^{-1}\cdot T_k(-B)\cdot\mathsf{C}f=\bar{T}_k(B)f.
\end{aligned}
\]

\textbf{(a-4) Fields.}
The display (23.2) and the two identities are identities between operators given by rows, namely $T_k(B)$, $\bar{T}_k(B)$, $T^{(i)}(B,Y)$, $\bar{T}^{(i)}(B,Y)$, $T^{(2)}(B,Y)\cdot\delta_{\alpha,B}$ and $\bar{T}^{(2)}(B,Y)\cdot\delta_{\alpha,-B}$.
By (i23.V1) they hold for functions with values in $\mathrm{Vs}^4$, in particular for fields.
\end{proof}

\paragraph{(b) What a row contains and reads}

\begin{proof}[Proof of Proposition~23.1\,(b)]
Fix a unit site $y$ of $Y$.

\textbf{(b-1) Sites and pairs.}
The display (i23.4) contains $f$ at the sites $(c,g)$ and $(c,g s_i)$, $i\le 2$, with $c\in\Delta_y$, which together are the sites over $\Delta_y$.
It contains $f$ moreover at the sites $x^{(3)}=(c+\eta_k\cdot g e_1,\,g s_3)$ with $c\in\Delta_y$ and $c_3\in Y$, and at no other site.
The sites of the first sort lie over $Y$ by (0-2); those of the second sort lie over $Y$ because of the indicator $\mathbf{1}[c_3\in Y]$.
The display (i23.5) contains $F$ exactly at the pairs $(x,x^{(3)})$ with $c\in\Delta_y$ and $c_3\notin Y$.
For such a pair the second point is $x^{(3)}$, and the photon site $c$ of the first point lies in $\Delta_y$, hence in $Y$ by (0-2).
So the pair is of kind (II) for $Y$ (Section~I.1.7), and its first point lies over $Y$.

\textbf{(b-2) Bonds.}
The entries of both rows are formed from the numbers $\varphi_B(y,c)$, $c\in\Delta_y$, and $u_B(c,c_3)$, $c\in\Delta_y$ (the latter occur in (i23.4) only, and there only for $c_3\in Y$), from the matrices $M_i(x)$, and from the numbers $\mu$, $b_k$, $\eta_k$, $\alpha$, $1/48$.
By (0-1), $\varphi_B(y,c)$ contains $B$ only on bonds contained in $\Delta_y$, and $u_B(c,c_3)$ contains $B$ only on the bond $(c,c_3)$.
Hence both rows read $B$ on the bonds contained in $\Delta_y$ and on the bonds $(c,c_3)$, $c\in\Delta_y$.
Each of these bonds has an end point in $\Delta_y$.

\textbf{(b-3) The set $Y$.}
The set $Y$ enters the rows at $y$ in two ways only:
through the alternative whether $y$ lies in $Y$, where the row is given by the display (23.1), or not, where the row is zero;
and through the indicators $\mathbf{1}[c_3\in Y]$, $\mathbf{1}[c_3\notin Y]$, $c\in\Delta_y$.
Consequently the rows read $Y$ through the site $y$ and the sites $c_3$, $c\in\Delta_y$.

\textbf{(b-4) Distances.}
Let $c\in\Delta_y$.
Then $|c-y|\le(1-\eta_k)/2$ by (i23.2), and
\[
|c_3-y|\le|c-y|+\eta_k\cdot|g e_1|\le\frac{1+\eta_k}{2}.
\]
It follows that the photon sites of all the sites named above, both points of every pair named above and both end points of every bond named above lie within the sup-distance $(1+\eta_k)/2$ of $y$.

\textbf{(b-5) Analyticity.}
Every entry is a finite sum of products, each of which consists of a number that does not depend on $B$, an entry of one of the matrices $M_i(x)$ or of $1_4$, a factor $\varphi_B(y,c)$ and, in some of the terms, a factor $u_B(c,c_3)$.
By (i23.3) these factors are exponentials of linear forms in the finitely many numbers $B(b)$.
Therefore the entries are entire functions of the numbers $B(b)$.

\textbf{(b-6) The barred operators.}
The rows (i23.6) contain $f$ and $F$ at the same sites and pairs as (i23.4) and (i23.5), with the same indicators.
Their entries are formed from $\varphi_B(y,c)^{-1}=\varphi_{-B}(y,c)$, from $u_B(c,c_3)^{-1}=u_{-B}(c,c_3)$, from the matrices $M_i(x^{(i)})^T$ and from the same numbers.
Accordingly the rows of $\bar{T}^{(1)}(B,Y)$ and of $\bar{T}^{(2)}(B,Y)$ at $y$ contain $f$, resp.\ $F$, at the sites, resp.\ pairs, of (b-1);
they read $B$ on the bonds of (b-2), since the inverse phases are formed with the same numbers $B(b)$;
they read $Y$ through the sites of (b-3);
the distances are those of (b-4);
and their entries are entire functions of the numbers $B(b)$, since by (i23.3) the inverse phases are again exponentials of linear forms.
\end{proof}

\paragraph{(c) Covariances of the operators}

\begin{proof}[Proof of Proposition~23.1\,(c)]
\textbf{(c-1) Gauge.}
Let $\boldsymbol{\lambda}$ be as in Proposition~21.1\,(d), that is, a complex function on $\mathbb{L}_k$ such that $d\boldsymbol{\lambda}$ and $\exp(i e_k\cdot\boldsymbol{\lambda})$ are $\ell\mathbb{Z}^3$-periodic.
Then $B-d\boldsymbol{\lambda}$ is a bond function on $T_k$, the function $\exp(i e_k\cdot\boldsymbol{\lambda})\cdot f$ lies in $\mathfrak{S}^{\mathrm{ch}}(\eta_k,\ell)$, and $\exp(i e_k\cdot\boldsymbol{\lambda})\cdot F$ is a function of pairs.
The relations (ph-G) of (0-7) hold.
Let $y$ be a unit site of $Y$.
The expressions
\[
\left(T^{(1)}(B-d\boldsymbol{\lambda},Y)\,\exp(i e_k\cdot\boldsymbol{\lambda})\,f\right)(y),
\qquad
\left(T^{(2)}(B-d\boldsymbol{\lambda},Y)\,\exp(i e_k\cdot\boldsymbol{\lambda})\,F\right)(y)
\]
consist of terms of four kinds, each multiplied by matrices and numbers that do not contain the connection.
By (ph-G) one has, for these four kinds:
\begin{itemize}
\item[(1)] $\varphi_{B-d\boldsymbol{\lambda}}(y,c)\cdot\exp\left(i e_k\cdot\boldsymbol{\lambda}(c)\right)\cdot f(x)=\exp\left(i e_k\cdot\boldsymbol{\lambda}(y)\right)\cdot\varphi_B(y,c)\cdot f(x)$;
\item[(2)] $\varphi_{B-d\boldsymbol{\lambda}}(y,c)\cdot\exp\left(i e_k\cdot\boldsymbol{\lambda}(c)\right)\cdot f(x^{(i)})=\exp\left(i e_k\cdot\boldsymbol{\lambda}(y)\right)\cdot\varphi_B(y,c)\cdot f(x^{(i)})$ for $i\le 2$, because $x^{(i)}$ lies over the same photon site $c$;
\item[(3)] for the terms with $f(x^{(3)})$,
\begin{align*}
&\varphi_{B-d\boldsymbol{\lambda}}(y,c)\cdot u_{B-d\boldsymbol{\lambda}}(c,c_3)\cdot\exp\left(i e_k\cdot\boldsymbol{\lambda}(c_3)\right)\cdot f(x^{(3)})\\
&\qquad=\exp\left(i e_k\cdot\left(\boldsymbol{\lambda}(y)-\boldsymbol{\lambda}(c)+\boldsymbol{\lambda}(c)-\boldsymbol{\lambda}(c_3)+\boldsymbol{\lambda}(c_3)\right)\right)\cdot\varphi_B(y,c)\cdot u_B(c,c_3)\cdot f(x^{(3)})\\
&\qquad=\exp\left(i e_k\cdot\boldsymbol{\lambda}(y)\right)\cdot\varphi_B(y,c)\cdot u_B(c,c_3)\cdot f(x^{(3)});
\end{align*}
\item[(4)] $\varphi_{B-d\boldsymbol{\lambda}}(y,c)\cdot\exp\left(i e_k\cdot\boldsymbol{\lambda}(c)\right)\cdot F(x,x^{(3)})=\exp\left(i e_k\cdot\boldsymbol{\lambda}(y)\right)\cdot\varphi_B(y,c)\cdot F(x,x^{(3)})$, the factor $\exp\left(i e_k\cdot\boldsymbol{\lambda}(c)\right)$ being the phase of the first point.
\end{itemize}
Thus every term of (23.1) is multiplied by $\exp\left(i e_k\cdot\boldsymbol{\lambda}(y)\right)$.
At the unit sites outside $Y$ both sides vanish.
Hence
\[
T^{(i)}(B-d\boldsymbol{\lambda},Y)\,\exp(i e_k\cdot\boldsymbol{\lambda})=\exp(i e_k\cdot\boldsymbol{\lambda})\,T^{(i)}(B,Y),\qquad i=1,2.
\]
For the barred operators, note that $-\boldsymbol{\lambda}$ has the two properties required of $\boldsymbol{\lambda}$, that $-(B-d\boldsymbol{\lambda})=(-B)-d(-\boldsymbol{\lambda})$, and that $\mathsf{C}$ commutes with the multiplication by scalar functions.
The relation just proved, applied to $(-B,-\boldsymbol{\lambda})$, gives
\begin{align*}
\bar{T}^{(i)}(B-d\boldsymbol{\lambda},Y)\,\exp(-i e_k\cdot\boldsymbol{\lambda})
&=\mathsf{C}^{-1}\,T^{(i)}\left((-B)-d(-\boldsymbol{\lambda}),Y\right)\exp\left(i e_k\cdot(-\boldsymbol{\lambda})\right)\mathsf{C}\\
&=\mathsf{C}^{-1}\exp\left(i e_k\cdot(-\boldsymbol{\lambda})\right)T^{(i)}(-B,Y)\,\mathsf{C}
=\exp(-i e_k\cdot\boldsymbol{\lambda})\,\bar{T}^{(i)}(B,Y).
\end{align*}

\textbf{(c-2) Lattice.}
Let $\tilde{g}=(S,a)\in\mathbb{G}_k$ (Section~I.4.2), and let $r=r(S)$, a diagonal matrix with the entries $\pm1$ (Section~I.4.1).
Let $a\in M\mathbb{Z}^3$ be a representative; one and the same representative is taken on both sides of each relation.
(Any other representative is $a+\ell\cdot n$ with $n\in\mathbb{Z}^3$.
It replaces $\tilde{g}^{-1}c=r(c-a)$ by $\tilde{g}^{-1}c-\ell\cdot rn$, and the entries of $rn$ are $\pm n_\mu$.
Hence, for antiperiodic $f$, $F$, $v$, it multiplies $\tilde{g}f$, $\tilde{g}F$, $\tilde{g}v$ and $\mathrm{glb}\,f$, $\mathrm{glb}\,F$, $\mathrm{glb}\,v$ by $(-1)^{n_1+n_2+n_3}$, and therefore both sides of each relation by the same sign (Section~I.4.3).
The bond function $\tilde{g}B$ and the set $\tilde{g}Y$ do not change.)
The following five facts hold.
\begin{itemize}
\item[(i23.L1)] The map $c\mapsto\tilde{g}c=rc+a$ is a bijection of $\mathbb{L}_k$ and of $\mathbb{Z}^3$, and $|\tilde{g}c-\tilde{g}c'|=|c-c'|$.
Hence $\tilde{g}\Delta_y=\Delta_{\tilde{g}y}$.
It maps the $M$-cube with the centre $z\in M\mathbb{Z}^3$ onto the $M$-cube with the centre $\tilde{g}z\in M\mathbb{Z}^3$, and $\tilde{g}(c+\ell\cdot n)=\tilde{g}c+\ell\cdot rn$ for $n\in\mathbb{Z}^3$.
Consequently $\tilde{g}Y$ is an $\ell\mathbb{Z}^3$-periodic union of $M$-cubes, $\tilde{g}B$ is an $\ell\mathbb{Z}^3$-periodic bond function, and $c$ lies in $Y$ if and only if $\tilde{g}c$ lies in $\tilde{g}Y$.
\item[(i23.L2)] The map $x\mapsto\tilde{g}x=(rc+a,\,rg)$ is a bijection of $\Delta_y\times O_{\mathrm{cub}}$ onto $\Delta_{\tilde{g}y}\times O_{\mathrm{cub}}$, since $g\mapsto rg$ is a bijection of $O_{\mathrm{cub}}$.
So, by (i23.1) (see (I.a1)), for every function $\mathrm{hf}$ on $\Delta_{\tilde{g}y}\times O_{\mathrm{cub}}$ with values in a complex vector space,
\[
\int_{\Delta_{\tilde{g}y}}\mathrm{d}\tilde{x}\,\mathrm{hf}(\tilde{x})=\int_{\Delta_y}\mathrm{d}x\,\mathrm{hf}(\tilde{g}x).
\]
Moreover $(\tilde{g}x)^{(i)}=\tilde{g}(x^{(i)})$ for $i=0,\ldots,3$, by (I.a2).
Hence $\pi\left((\tilde{g}x)^{(3)}\right)=\tilde{g}c_3$, and the pair $(x,x^{(3)})$ is mapped to $(\tilde{g}x,\tilde{g}x^{(3)})=\left(\tilde{g}x,(\tilde{g}x)^{(3)}\right)$.
\item[(i23.L3)] $M_i(\tilde{g}x)\cdot S=S\cdot M_i(x)$, by (I.m4).
\item[(i23.L4)] $\varphi_{\tilde{g}B}(\tilde{g}y,\tilde{g}c)=\varphi_B(y,c)$ and $u_{\tilde{g}B}(\tilde{g}c,\tilde{g}c_3)=u_B(c,c_3)$, by (ph-g) of (0-7).
\item[(i23.L5)] $(\tilde{g}f)(\tilde{g}x)=S\cdot f(x)$, $(\tilde{g}f)\left((\tilde{g}x)^{(i)}\right)=(\tilde{g}f)\left(\tilde{g}(x^{(i)})\right)=S\cdot f(x^{(i)})$, and $(\tilde{g}F)\left(\tilde{g}x,(\tilde{g}x)^{(3)}\right)=S\cdot F(x,x^{(3)})$, by Section~I.4.3 together with (i23.L2).
\end{itemize}
Let $y$ be a unit site of $Y$; then $\tilde{g}y$ is a unit site of $\tilde{g}Y$.
Write (23.1) for $(\tilde{g}B,\tilde{g}Y,\tilde{g}f)$ at the unit site $\tilde{g}y$ and substitute $\tilde{x}=\tilde{g}x$ as in (i23.L2).
By (i23.L1), (i23.L3), (i23.L4) and (i23.L5) one obtains
\begin{align*}
&\left(T^{(1)}(\tilde{g}B,\tilde{g}Y)\,\tilde{g}f\right)(\tilde{g}y)\\
&\quad=\int_{\Delta_y}\mathrm{d}x\,\varphi_{\tilde{g}B}(\tilde{g}y,\tilde{g}c)\cdot\left\{\left(1+\frac{\mu}{b_k}\right)\cdot(\tilde{g}f)(\tilde{g}x)\right.\\
&\qquad\left.{}-(b_k\cdot\eta_k)^{-1}\cdot\left[\sum_{i\le 2}M_i(\tilde{g}x)\cdot\left((\tilde{g}f)\left((\tilde{g}x)^{(i)}\right)-(\tilde{g}f)(\tilde{g}x)\right)\right.\right.\\
&\qquad\qquad\left.\left.{}+\mathbf{1}[\tilde{g}c_3\in\tilde{g}Y]\cdot M_3(\tilde{g}x)\cdot\left(u_{\tilde{g}B}(\tilde{g}c,\tilde{g}c_3)\cdot(\tilde{g}f)\left((\tilde{g}x)^{(3)}\right)-(\tilde{g}f)(\tilde{g}x)\right)\right]\right\}\\
&\quad=\int_{\Delta_y}\mathrm{d}x\,\varphi_B(y,c)\cdot\left\{\left(1+\frac{\mu}{b_k}\right)\cdot S\cdot f(x)-(b_k\cdot\eta_k)^{-1}\cdot\left[\sum_{i\le 2}S\cdot M_i(x)\cdot\left(f(x^{(i)})-f(x)\right)\right.\right.\\
&\qquad\qquad\left.\left.{}+\mathbf{1}[c_3\in Y]\cdot S\cdot M_3(x)\cdot\left(u_B(c,c_3)\cdot f(x^{(3)})-f(x)\right)\right]\right\}\\
&\quad=S\cdot\left(T^{(1)}(B,Y)\,f\right)(y)=\left(\tilde{g}\,T^{(1)}(B,Y)\,f\right)(\tilde{g}y).
\end{align*}
In the same way, the second display of (23.1) for $(\tilde{g}B,\tilde{g}Y,\tilde{g}F)$ at $\tilde{g}y$, the substitution $\tilde{x}=\tilde{g}x$ of (i23.L2) and then (i23.L1), (i23.L3), (i23.L4), (i23.L5) give
\begin{align*}
\left(T^{(2)}(\tilde{g}B,\tilde{g}Y)\,\tilde{g}F\right)(\tilde{g}y)
&=-b_k^{-1}\cdot\eta_k^{\alpha-1}\cdot\int_{\Delta_y}\mathrm{d}x\,\mathbf{1}[\tilde{g}c_3\notin\tilde{g}Y]\cdot\varphi_{\tilde{g}B}(\tilde{g}y,\tilde{g}c)\\
&\qquad{}\cdot M_3(\tilde{g}x)\cdot(\tilde{g}F)\left(\tilde{g}x,(\tilde{g}x)^{(3)}\right)\\
&=-b_k^{-1}\cdot\eta_k^{\alpha-1}\cdot\int_{\Delta_y}\mathrm{d}x\,\mathbf{1}[c_3\notin Y]\cdot\varphi_B(y,c)\cdot S\cdot M_3(x)\cdot F(x,x^{(3)})\\
&=S\cdot\left(T^{(2)}(B,Y)\,F\right)(y)=\left(\tilde{g}\,T^{(2)}(B,Y)\,F\right)(\tilde{g}y).
\end{align*}
Every unit site is $\tilde{g}y$ for exactly one unit site $y$, and for $y$ outside $Y$ the site $\tilde{g}y$ lies outside $\tilde{g}Y$, so that both sides vanish.
Hence $T^{(i)}(\tilde{g}B,\tilde{g}Y)\,\tilde{g}=\tilde{g}\,T^{(i)}(B,Y)$.

For the barred operators: by Section~I.4.1 there is a sign $\mathrm{sgn}_S\in\{1,-1\}$ with $(S^{-1})^T=\mathrm{sgn}_S\cdot\mathsf{C}^{-1}\cdot S\cdot\mathsf{C}$.
The matrix $\mathsf{C}$ acts on the values and $\tilde{g}^{-1}$ on the indices.
Therefore $\mathrm{glb}=\mathrm{sgn}_S\cdot\mathsf{C}^{-1}\,\tilde{g}\,\mathsf{C}$ on spinors on $T^{\mathrm{f}}_k$, on functions of pairs and on spinors on $T^0_k$, where $(\mathrm{glb}\,v)(y):=(S^{-1})^T\cdot v(\tilde{g}^{-1}y)$.
Moreover $-\tilde{g}B=\tilde{g}(-B)$.
The relation just shown, applied to $-B$, yields
\begin{align*}
\bar{T}^{(i)}(\tilde{g}B,\tilde{g}Y)\,\mathrm{glb}
&=\mathsf{C}^{-1}\,T^{(i)}\left(\tilde{g}(-B),\tilde{g}Y\right)\mathsf{C}\cdot\mathrm{sgn}_S\cdot\mathsf{C}^{-1}\,\tilde{g}\,\mathsf{C}
=\mathrm{sgn}_S\cdot\mathsf{C}^{-1}\,T^{(i)}\left(\tilde{g}(-B),\tilde{g}Y\right)\tilde{g}\,\mathsf{C}\\
&=\mathrm{sgn}_S\cdot\mathsf{C}^{-1}\,\tilde{g}\,T^{(i)}(-B,Y)\,\mathsf{C}
=\left(\mathrm{sgn}_S\cdot\mathsf{C}^{-1}\,\tilde{g}\,\mathsf{C}\right)\left(\mathsf{C}^{-1}\,T^{(i)}(-B,Y)\,\mathsf{C}\right)\\
&=\mathrm{glb}\,\bar{T}^{(i)}(B,Y).
\end{align*}

\textbf{(c-3) Charge conjugation.}
Multiplying the definition $\bar{T}^{(i)}(B,Y)=\mathsf{C}^{-1}\,T^{(i)}(-B,Y)\,\mathsf{C}$ from the left by $\mathsf{C}$ one obtains $T^{(i)}(-B,Y)\,\mathsf{C}=\mathsf{C}\,\bar{T}^{(i)}(B,Y)$.
The definition at $-B$ together with (i23.C1) gives
\[
\bar{T}^{(i)}(-B,Y)\,\mathsf{C}=\mathsf{C}^{-1}\,T^{(i)}(B,Y)\,\mathsf{C}^2=-\mathsf{C}^{-1}\,T^{(i)}(B,Y)=\mathsf{C}\,T^{(i)}(B,Y).
\]

\textbf{(c-4) Conjugation.}
Let $\mu$ be real.
The numbers $b_k$, $\eta_k$, $\alpha$, $\rho$, $w_i$ and the weight $\eta_k^3/48$ are real.
By (ph-K) of (0-7),
\[
\overline{\varphi_{\overline{B}}(y,c)}=\varphi_{-B}(y,c),\qquad\overline{u_{\overline{B}}(c,c_3)}=u_{-B}(c,c_3),
\]
and by (I.m3), $\overline{M_i(x)}=\varkappa^{-1}\cdot M_i(x)\cdot\varkappa$.
The entries of the operator $\overline{T^{(i)}(\overline{B},Y)}$ are the complex conjugates of those of $T^{(i)}(\overline{B},Y)$ (Section~0.2).
Hence, for a unit site $y$ of $Y$,
\begin{align*}
\left(\overline{T^{(1)}(\overline{B},Y)}\,f\right)(y)
&=\int_{\Delta_y}\mathrm{d}x\,\varphi_{-B}(y,c)\cdot\left\{\left(1+\frac{\mu}{b_k}\right)\cdot f(x)\right.\\
&\quad\left.{}-(b_k\cdot\eta_k)^{-1}\cdot\left[\sum_{i\le 2}\varkappa^{-1}\cdot M_i(x)\cdot\varkappa\cdot\left(f(x^{(i)})-f(x)\right)\right.\right.\\
&\qquad\left.\left.{}+\mathbf{1}[c_3\in Y]\cdot\varkappa^{-1}\cdot M_3(x)\cdot\varkappa\cdot\left(u_{-B}(c,c_3)\cdot f(x^{(3)})-f(x)\right)\right]\right\}\\
&=\varkappa^{-1}\left(T^{(1)}(-B,Y)\,\varkappa f\right)(y),\\
\left(\overline{T^{(2)}(\overline{B},Y)}\,F\right)(y)
&=-b_k^{-1}\cdot\eta_k^{\alpha-1}\cdot\int_{\Delta_y}\mathrm{d}x\,\mathbf{1}[c_3\notin Y]\cdot\varphi_{-B}(y,c)\cdot\varkappa^{-1}\cdot M_3(x)\cdot\varkappa\cdot F(x,x^{(3)})\\
&=\varkappa^{-1}\left(T^{(2)}(-B,Y)\,\varkappa F\right)(y).
\end{align*}
For the barred operators: $\mathsf{C}$ is real, so that $\overline{\mathsf{C}^{-1}\,O\,\mathsf{C}}=\mathsf{C}^{-1}\,\overline{O}\,\mathsf{C}$ for every operator $O$ given by rows.
By (I.c4), $\varkappa\,\mathsf{C}=-\mathsf{C}\,\varkappa$, and hence $\mathsf{C}^{-1}\varkappa^{-1}=(\varkappa\,\mathsf{C})^{-1}=-\varkappa^{-1}\mathsf{C}^{-1}$.
Apply the relation just shown with $-B$ in place of $B$; note that $-\overline{B}=\overline{-B}$.
This gives
\begin{align*}
\overline{\bar{T}^{(i)}(\overline{B},Y)}
&=\mathsf{C}^{-1}\left[\overline{T^{(i)}\left(\overline{-B},Y\right)}\right]\mathsf{C}
=\mathsf{C}^{-1}\varkappa^{-1}\,T^{(i)}(B,Y)\,\varkappa\,\mathsf{C}
=\left(-\varkappa^{-1}\mathsf{C}^{-1}\right)T^{(i)}(B,Y)\left(-\mathsf{C}\,\varkappa\right)\\
&=\varkappa^{-1}\left[\mathsf{C}^{-1}\,T^{(i)}(B,Y)\,\mathsf{C}\right]\varkappa
=\varkappa^{-1}\,\bar{T}^{(i)}(-B,Y)\,\varkappa.
\end{align*}
By (i23.V1) the relations (c-1), (c-2) and (c-3) hold also for functions with values in $\mathrm{Vs}^4$; (c-4) is a relation between the entries.
\end{proof}

\paragraph{(d) The substitution: locality, parity}

\begin{proof}[Proof of Proposition~23.1\,(d)]
Let
\[
F=\sum_{I_1,I_2,I_3}c_F(I_1,I_2,I_3)\cdot\psi(I_1)\cdot\chi(I_2)\cdot\Psi(I_3)\in\mathrm{Gr}[\psi,\chi,\Psi\,|\,\mathbf{Y}],
\]
where $I_3$ is a set of indices $u=(y,\beta,\omega)$ with $y$ a unit site of $\mathbf{Y}$.
The map $\mathrm{sub}_{\mathbf{Y},B}$ is a homomorphism which fixes $\psi$ and $\chi$; therefore
\begin{equation}
\mathrm{sub}_{\mathbf{Y},B}\left(F\right)=\sum_{I_1,I_2,I_3}c_F(I_1,I_2,I_3)\cdot\psi(I_1)\cdot\chi(I_2)\cdot\prod_{u\in I_3}\mathrm{sub}_{\mathbf{Y},B}\left(\Psi(u)\right),
\tag{i23.8}
\end{equation}
the factors of the product being taken in increasing order.
By (i23.7), $\mathrm{sub}_{\mathbf{Y},B}\left(\Psi(y,\beta,0)\right)$ is the linear combination of the generators $\psi(x',\beta',0)$, $\chi(x,x^{(3)},\beta',0)$ whose coefficients are the entries of the rows of $T^{(1)}(B,\mathbf{Y})$, $T^{(2)}(B,\mathbf{Y})$ at $y$; here a generator at a point of the covering lattice is $\pm$ a generator indexed by the fundamental domain, by (I.AP).
Likewise $\mathrm{sub}_{\mathbf{Y},B}\left(\Psi(y,\beta,1)\right)$ is the linear combination of the generators $\psi(x',\beta',1)$, $\chi(x,x^{(3)},\beta',1)$ whose coefficients are the entries of the rows of $\bar{T}^{(1)}(B,\mathbf{Y})$, $\bar{T}^{(2)}(B,\mathbf{Y})$ at $y$.

\emph{Locality.}
By (b), applied with $Y=\mathbf{Y}$, these generators are generators $\psi(\xi)$ with $\xi$ over $\mathbf{Y}$ and generators $\chi(\boldsymbol{\zeta})$ at pairs of kind (II) for $\mathbf{Y}$ whose first point lies over $\mathbf{Y}$.
Expand the products in (i23.8) and bring the generators into increasing order.
It follows that $\mathrm{sub}_{\mathbf{Y},B}\left(F\right)$ contains, in the sense of Section~I.5.1, generators that $F$ contains (only the terms of (i23.8) with $c_F(I_1,I_2,I_3)\ne0$ contribute) and generators of the two sorts just described, and no other generators.
Each of its coefficients is a finite sum of terms $\pm c_F(I_1,I_2,I_3)$ times a product of $\#|I_3|$ entries of rows at unit sites of $\mathbf{Y}$, hence a polynomial in the coefficients of $F$ and in these entries.
By (b) the rows at $y$ read $B$ on bonds with an end point in $\Delta_y$.
Since $\Delta_y$ lies in $\mathbf{Y}$ by (0-2), these bonds meet $\mathbf{Y}$.

\emph{Parity.}
The map $\mathrm{sub}_{\mathbf{Y},B}$ sends every generator to a linear combination of generators with the same bar index $\omega$.
Consequently it sends a monomial to a linear combination of monomials that have the same number of unbarred factors and the same number of barred factors as the monomial itself.
In particular the difference of these two numbers and the parity of their sum are preserved, and even elements are mapped to even elements.
\end{proof}

\paragraph{(e) The substitution: evaluation}

\begin{proof}[Proof of Proposition~23.1\,(e)]
Let $\mathrm{ev}:\mathrm{Gr}[\psi,\chi]\to\mathrm{Gr}'$ be a homomorphism, linear over $\mathbb{C}$, into a Grassmann algebra $\mathrm{Gr}'$, such that $\mathrm{ev}(\chi)=\delta_{\alpha,B}\,\mathrm{ev}(\psi)$ and $\mathrm{ev}(\bar{\chi})=\delta_{\alpha,-B}\,\mathrm{ev}(\bar{\psi})$ at the pairs of kind (II) for $\mathbf{Y}$.
Here, in the notation of (i23.V2),
\begin{align*}
\mathrm{ev}(\psi)&:=\mathrm{ev}\circ\psi=\left(x\mapsto\left(\mathrm{ev}\left(\psi(x,\beta,0)\right)\right)_\beta\right),\\
\mathrm{ev}(\bar{\psi})&:=\mathrm{ev}\circ\bar{\psi}=\left(x\mapsto\left(\mathrm{ev}\left(\psi(x,\beta,1)\right)\right)_\beta\right),\\
\mathrm{ev}(\chi)&:=\mathrm{ev}\circ\chi=\left((x,x')\mapsto\left(\mathrm{ev}\left(\chi(x,x',\beta,0)\right)\right)_\beta\right),\\
\mathrm{ev}(\bar{\chi})&:=\mathrm{ev}\circ\bar{\chi}=\left((x,x')\mapsto\left(\mathrm{ev}\left(\chi(x,x',\beta,1)\right)\right)_\beta\right);
\end{align*}
these functions take their values in $(\mathrm{Gr}')^4$ and have the antiperiodicity of $\psi$, resp.\ of $\chi$.
The composition $\mathrm{ev}\circ\mathrm{sub}_{\mathbf{Y},B}$ is a homomorphism of $\mathrm{Gr}[\psi,\chi,\Psi\,|\,\mathbf{Y}]$, linear over $\mathbb{C}$ because $\mathrm{ev}$ and $\mathrm{sub}_{\mathbf{Y},B}$ are ((0-6), (i23.H1)); on $\psi(\xi)$ and $\chi(\boldsymbol{\zeta})$ it coincides with $\mathrm{ev}$.
For a unit site $y$ of $\mathbf{Y}$, (i23.7) and the linear case of (i23.V2) give
\[
\mathrm{ev}\left(\mathrm{sub}_{\mathbf{Y},B}\Psi(y)\right)=\left(T^{(1)}(B,\mathbf{Y})\,\mathrm{ev}(\psi)\right)(y)+\left(T^{(2)}(B,\mathbf{Y})\,\mathrm{ev}(\chi)\right)(y).
\]
By (b) the row of $T^{(2)}(B,\mathbf{Y})$ at $y$ contains its argument only at pairs of kind (II) for $\mathbf{Y}$, and there $\mathrm{ev}(\chi)=\delta_{\alpha,B}\,\mathrm{ev}(\psi)$.
So the right-hand side equals
\[
\left(T^{(1)}(B,\mathbf{Y})\,\mathrm{ev}(\psi)+T^{(2)}(B,\mathbf{Y})\,\delta_{\alpha,B}\,\mathrm{ev}(\psi)\right)(y)=\left(T_k(B)\,\mathrm{ev}(\psi)\right)(y),
\]
by (a) for functions with values in $(\mathrm{Gr}')^4$ (see (a-4)).
For the barred generators, (i23.7) and the linear case of (i23.V2) give
\[
\mathrm{ev}\left(\mathrm{sub}_{\mathbf{Y},B}\bar{\Psi}(y)\right)=\left(\bar{T}^{(1)}(B,\mathbf{Y})\,\mathrm{ev}(\bar{\psi})\right)(y)+\left(\bar{T}^{(2)}(B,\mathbf{Y})\,\mathrm{ev}(\bar{\chi})\right)(y).
\]
By (b) the row of $\bar{T}^{(2)}(B,\mathbf{Y})$ at $y$ contains its argument only at pairs of kind (II) for $\mathbf{Y}$, where $\mathrm{ev}(\bar{\chi})=\delta_{\alpha,-B}\,\mathrm{ev}(\bar{\psi})$.
By the barred identity of (a) the sum is therefore
\[
\left(\bar{T}^{(1)}(B,\mathbf{Y})\,\mathrm{ev}(\bar{\psi})+\bar{T}^{(2)}(B,\mathbf{Y})\,\delta_{\alpha,-B}\,\mathrm{ev}(\bar{\psi})\right)(y)=\left(\bar{T}_k(B)\,\mathrm{ev}(\bar{\psi})\right)(y).
\]
By (i23.H2), applied to two linear maps, a homomorphism that is linear over $\mathbb{C}$ is determined by its values on the generators.
This proves the first assertion.

Assume now in addition that $S_k(B)$ exists and that $\mathrm{ev}(\psi)=H_k(B)\,\Phi$ and $\mathrm{ev}(\bar{\psi})=\bar{H}_k(B)\,\bar{\Phi}$ with fields $\Phi$, $\bar{\Phi}$ on $T^0_k$.
By Proposition~21.1\,(a), $T_k(B)\,H_k(B)=1$ on $\mathfrak{S}(1,\ell)$; indeed, by Section~I.3.5 and (a-1),
\[
T_k(B)\,H_k(B)=b_k^{-1}\cdot Q_k(B)\,J^*F_\infty(B)\cdot b_k\cdot S_k(B)\,J\,Q_k^T(-B)=Q_k(B)\,J^*J\,Q_k^T(-B)=1.
\]
By Proposition~21.1\,(c), $F_\infty(-B)$ is invertible because $F_\infty(B)$ is; so $\bar{H}_k(B)=\mathsf{C}^{-1}H_k(-B)\,\mathsf{C}$ exists, and $\bar{T}_k(B)\,\bar{H}_k(B)=1$.
By (i23.V1) both relations hold on fields.
Hence $\left(T_k(B)\,\mathrm{ev}(\psi)\right)(y)=\Phi(y)$ and $\left(\bar{T}_k(B)\,\mathrm{ev}(\bar{\psi})\right)(y)=\bar{\Phi}(y)$.
\end{proof}

\paragraph{(f) The substitution: four covariances}

\begin{proof}[Proof of Proposition~23.1\,(f)]
\emph{The four maps on the generators.}
They are those of Section~I.4.5, on $\Psi$ with the unit sites $y$ in place of $x$.
For an additive map $\theta$ and a family $\phi$ of generators write $\theta\circ\phi=\left(t\mapsto\left(\theta\left(\phi(t,\beta,\omega)\right)\right)_\beta\right)$, as in (i23.V2).
Then:
\begin{itemize}
\item[$a_{\tilde{g}}$:] $a_{\tilde{g}}\circ\psi=\tilde{g}\psi$, $a_{\tilde{g}}\circ\chi=\tilde{g}\chi$, $a_{\tilde{g}}\circ\Psi=\tilde{g}\Psi$, $a_{\tilde{g}}\circ\bar{\psi}=\mathrm{glb}\,\bar{\psi}$, $a_{\tilde{g}}\circ\bar{\chi}=\mathrm{glb}\,\bar{\chi}$, $a_{\tilde{g}}\circ\bar{\Psi}=\mathrm{glb}\,\bar{\Psi}$;
\item[$G_{\Lambda}$:] $G_{\Lambda}\circ\psi=\exp(i e_k\cdot\Lambda)\,\psi$, $G_{\Lambda}\circ\chi=\exp(i e_k\cdot\Lambda)\,\chi$ (with the phase of the first point), $G_{\Lambda}\circ\Psi=\exp(i e_k\cdot\Lambda)\,\Psi$, and the factor $\exp(-i e_k\cdot\Lambda)$ on $\bar{\psi}$, $\bar{\chi}$, $\bar{\Psi}$;
\item[$\mathrm{Cg}$:] $\mathrm{Cg}\circ\psi=\mathsf{C}\bar{\psi}$, $\mathrm{Cg}\circ\chi=\mathsf{C}\bar{\chi}$, $\mathrm{Cg}\circ\Psi=\mathsf{C}\bar{\Psi}$, $\mathrm{Cg}\circ\bar{\psi}=-\mathsf{C}\psi$, $\mathrm{Cg}\circ\bar{\chi}=-\mathsf{C}\chi$, $\mathrm{Cg}\circ\bar{\Psi}=-\mathsf{C}\Psi$;
\item[$\mathbf{K}$:] $\mathbf{K}\circ\phi=\varkappa\phi$ for each of the six families $\phi$; the map $\mathbf{K}$ is antilinear.
\end{itemize}
The maps $G_{\Lambda}$, $\mathrm{Cg}$ and $\mathbf{K}$ send every $\Psi(y,\beta,\omega)$ to a linear combination of the generators $\Psi(y,\beta',\omega')$ at the same unit site $y$, and the generators $\psi$, $\chi$ to linear combinations of generators $\psi$, $\chi$.
Hence they map $\mathrm{Gr}[\psi,\chi,\Psi\,|\,\mathbf{Y}]$ into itself and $\mathrm{Gr}[\psi,\chi]$ into itself, and both sides of the last three relations are defined on $\mathrm{Gr}[\psi,\chi,\Psi\,|\,\mathbf{Y}]$.
(For $a_{\tilde{g}}$ see (f-1).)

Let $\theta$ be one of the four maps.
In each of the four relations both sides are ring homomorphisms with unit into $\mathrm{Gr}[\psi,\chi]$, linear for $\theta=a_{\tilde{g}},G_{\Lambda},\mathrm{Cg}$ and antilinear for $\theta=\mathbf{K}$.
By (i23.H2) it suffices to compare them on the generators.
On $\psi(\xi)$ and $\chi(\boldsymbol{\zeta})$ both sides coincide with $\theta$, because the substitutions fix these generators and $\theta$ maps them to linear combinations of generators $\psi$, $\chi$.
It remains to consider $\Psi(y)$ and $\bar{\Psi}(y)$.

\textbf{(f-1) The map $a_{\tilde{g}}$.}
The maps $a_{\tilde{g}}$, $\tilde{g}$ and $\mathrm{glb}$ are formed with one and the same representative $a\in M\mathbb{Z}^3$.

\emph{The set $\tilde{g}\mathbf{Y}$ is a polymer.}
Two different $M$-cubes of $T_k$ have a common face if and only if their centres have representatives $z,z'\in M\mathbb{Z}^3$ with $z-z'=\pm M\cdot e_\mu$ for some direction $\mu$ (Section~I.1.1).
(For $z\ne z'$ in $M\mathbb{Z}^3$ the closed cubes $z+[-M/2,M/2]^3$ and $z'+[-M/2,M/2]^3$ of $\mathbb{R}^3$ have a square of side $M$ in common if $z-z'=\pm M\cdot e_\mu$;
they have only an edge or a corner in common if all coordinates of $z-z'$ have modulus at most $M$ and two or three of them are different from $0$;
and they have no common point otherwise.)
One has $\tilde{g}z-\tilde{g}z'=r(z-z')$, where $r$ is diagonal with the entries $\pm1$.
So $z-z'=\pm M\cdot e_\mu$ holds if and only if $\tilde{g}z-\tilde{g}z'=\pm M\cdot e_\mu$.
By (i23.L1), $\tilde{g}$ maps the $M$-cube with the centre $z$ onto the $M$-cube with the centre $\tilde{g}z$, and centres that are congruent modulo $\ell\mathbb{Z}^3$ to centres that are congruent modulo $\ell\mathbb{Z}^3$.
Therefore $\tilde{g}$ maps $M$-cubes of $T_k$ with a common face onto $M$-cubes of $T_k$ with a common face.
It follows that $\tilde{g}$ maps a non-empty union of $M$-cubes that is connected through faces, that is, a polymer (Section~I.1.1), onto a union of the same kind.

\emph{Both sides are defined.}
Let $y$ be a unit site of $\tilde{g}\mathbf{Y}$.
The map $a_{\tilde{g}}$ sends $\Psi(y)$ and $\bar{\Psi}(y)$ to linear combinations of the generators $\Psi(\tilde{g}^{-1}y,\,\cdot\,)$, and $\tilde{g}^{-1}y$ is a unit site of $\mathbf{Y}$; it sends the generators $\psi$, $\chi$ to linear combinations of generators $\psi$, $\chi$.
So $a_{\tilde{g}}$ maps $\mathrm{Gr}[\psi,\chi,\Psi\,|\,\tilde{g}\mathbf{Y}]$ into $\mathrm{Gr}[\psi,\chi,\Psi\,|\,\mathbf{Y}]$ and $\mathrm{Gr}[\psi,\chi]$ into itself, and both sides of the relation are defined on $\mathrm{Gr}[\psi,\chi,\Psi\,|\,\tilde{g}\mathbf{Y}]$.

\emph{The relation.}
Let $y$ be a unit site of $\tilde{g}\mathbf{Y}$.
By (i23.7), (i23.V2), and (c-2) for functions with values in $\mathrm{Gr}[\psi,\chi]^4$,
\begin{align*}
a_{\tilde{g}}\left(\mathrm{sub}_{\tilde{g}\mathbf{Y},\tilde{g}B}\Psi(y)\right)
&=\left(T^{(1)}(\tilde{g}B,\tilde{g}\mathbf{Y})\,\tilde{g}\psi\right)(y)+\left(T^{(2)}(\tilde{g}B,\tilde{g}\mathbf{Y})\,\tilde{g}\chi\right)(y)\\
&=\left(\tilde{g}\,T^{(1)}(B,\mathbf{Y})\,\psi\right)(y)+\left(\tilde{g}\,T^{(2)}(B,\mathbf{Y})\,\chi\right)(y)\\
&=S\cdot\left[\left(T^{(1)}(B,\mathbf{Y})\,\psi\right)(\tilde{g}^{-1}y)+\left(T^{(2)}(B,\mathbf{Y})\,\chi\right)(\tilde{g}^{-1}y)\right]\\
&=S\cdot\mathrm{sub}_{\mathbf{Y},B}\Psi(\tilde{g}^{-1}y)=\mathrm{sub}_{\mathbf{Y},B}\left(a_{\tilde{g}}\Psi(y)\right),\\
a_{\tilde{g}}\left(\mathrm{sub}_{\tilde{g}\mathbf{Y},\tilde{g}B}\bar{\Psi}(y)\right)
&=\left(\bar{T}^{(1)}(\tilde{g}B,\tilde{g}\mathbf{Y})\,\mathrm{glb}\,\bar{\psi}\right)(y)+\left(\bar{T}^{(2)}(\tilde{g}B,\tilde{g}\mathbf{Y})\,\mathrm{glb}\,\bar{\chi}\right)(y)\\
&=\left(\mathrm{glb}\,\bar{T}^{(1)}(B,\mathbf{Y})\,\bar{\psi}\right)(y)+\left(\mathrm{glb}\,\bar{T}^{(2)}(B,\mathbf{Y})\,\bar{\chi}\right)(y)\\
&=(S^{-1})^T\cdot\left[\left(\bar{T}^{(1)}(B,\mathbf{Y})\,\bar{\psi}\right)(\tilde{g}^{-1}y)+\left(\bar{T}^{(2)}(B,\mathbf{Y})\,\bar{\chi}\right)(\tilde{g}^{-1}y)\right]\\
&=(S^{-1})^T\cdot\mathrm{sub}_{\mathbf{Y},B}\bar{\Psi}(\tilde{g}^{-1}y)=\mathrm{sub}_{\mathbf{Y},B}\left(a_{\tilde{g}}\bar{\Psi}(y)\right).
\end{align*}

\emph{The representative.}
Let $a$ be replaced by $a+\ell\cdot n$, $n\in\mathbb{Z}^3$.
Then $\tilde{g}\mathbf{Y}$ and $\tilde{g}B$ do not change, while $\tilde{g}^{-1}x$ and $\tilde{g}^{-1}y$ are replaced by $\tilde{g}^{-1}x-\ell\cdot rn$ and $\tilde{g}^{-1}y-\ell\cdot rn$ (see (c-2)).
By (I.AP) the image under $a_{\tilde{g}}$ of every generator $\psi$, $\chi$, $\Psi$ is therefore multiplied by the one sign $(-1)^{n_1+n_2+n_3}$.
In other words, $a_{\tilde{g}}$ is replaced by its composition, from the left, with the automorphism of $\mathrm{Gr}[\psi,\chi,\Psi]$ that multiplies every generator by this sign; this automorphism exists by (i23.H1).
It maps $\mathrm{Gr}[\psi,\chi,\Psi\,|\,\mathbf{Y}]$ and $\mathrm{Gr}[\psi,\chi]$ into themselves, and it commutes with $\mathrm{sub}_{\mathbf{Y},B}$.
Indeed, on $\psi(\xi)$ and $\chi(\boldsymbol{\zeta})$ both compositions give the sign times the generator;
on $\Psi(u)$ both give the sign times $\mathrm{sub}_{\mathbf{Y},B}\left(\Psi(u)\right)$, because $\mathrm{sub}_{\mathbf{Y},B}\left(\Psi(u)\right)$ is a linear combination of generators (see (d));
and (i23.H2) applies.
Thus the change of the representative composes both sides of the relation, from the left, with one and the same automorphism, and the relation holds for every representative.

\textbf{(f-2) The map $G_{\Lambda}$.}
Let $\Lambda$ be a real function of $c$, that is, a function on the sites of $T_k$, and let $B^1=B-d\Lambda$ on the bonds that meet $\mathbf{Y}$.
The periodic extension of $\Lambda$ has the two properties required of $\boldsymbol{\lambda}$ in (c-1).
Let $y$ be a unit site of $\mathbf{Y}$.
By (b) the rows of $T^{(i)}(\,\cdot\,,\mathbf{Y})$ and $\bar{T}^{(i)}(\,\cdot\,,\mathbf{Y})$ at $y$ read the connection on bonds with an end point in $\Delta_y$, and these bonds meet $\mathbf{Y}$ by (0-2).
So these rows are the same for $B^1$ and for $B-d\Lambda$.
By (i23.7), (i23.V2) and (c-1),
\begin{align*}
G_{\Lambda}\left(\mathrm{sub}_{\mathbf{Y},B^1}\Psi(y)\right)
&=\left(T^{(1)}(B-d\Lambda,\mathbf{Y})\,\exp(i e_k\cdot\Lambda)\,\psi\right)(y)\\
&\qquad{}+\left(T^{(2)}(B-d\Lambda,\mathbf{Y})\,\exp(i e_k\cdot\Lambda)\,\chi\right)(y)\\
&=\exp\left(i e_k\cdot\Lambda(y)\right)\cdot\mathrm{sub}_{\mathbf{Y},B}\Psi(y)=\mathrm{sub}_{\mathbf{Y},B}\left(G_{\Lambda}\Psi(y)\right),\\
G_{\Lambda}\left(\mathrm{sub}_{\mathbf{Y},B^1}\bar{\Psi}(y)\right)
&=\left(\bar{T}^{(1)}(B-d\Lambda,\mathbf{Y})\,\exp(-i e_k\cdot\Lambda)\,\bar{\psi}\right)(y)\\
&\qquad{}+\left(\bar{T}^{(2)}(B-d\Lambda,\mathbf{Y})\,\exp(-i e_k\cdot\Lambda)\,\bar{\chi}\right)(y)\\
&=\exp\left(-i e_k\cdot\Lambda(y)\right)\cdot\mathrm{sub}_{\mathbf{Y},B}\bar{\Psi}(y)=\mathrm{sub}_{\mathbf{Y},B}\left(G_{\Lambda}\bar{\Psi}(y)\right).
\end{align*}

\textbf{(f-3) The map $\mathrm{Cg}$.}
For a unit site $y$ of $\mathbf{Y}$, by (i23.7), (i23.V2) and (c-3),
\begin{align*}
\mathrm{Cg}\left(\mathrm{sub}_{\mathbf{Y},-B}\Psi(y)\right)
&=\left(T^{(1)}(-B,\mathbf{Y})\,\mathsf{C}\bar{\psi}\right)(y)+\left(T^{(2)}(-B,\mathbf{Y})\,\mathsf{C}\bar{\chi}\right)(y)\\
&=\mathsf{C}\cdot\left[\left(\bar{T}^{(1)}(B,\mathbf{Y})\,\bar{\psi}\right)(y)+\left(\bar{T}^{(2)}(B,\mathbf{Y})\,\bar{\chi}\right)(y)\right]\\
&=\mathsf{C}\cdot\mathrm{sub}_{\mathbf{Y},B}\bar{\Psi}(y)=\mathrm{sub}_{\mathbf{Y},B}\left(\mathrm{Cg}\,\Psi(y)\right),\\
\mathrm{Cg}\left(\mathrm{sub}_{\mathbf{Y},-B}\bar{\Psi}(y)\right)
&=-\left(\bar{T}^{(1)}(-B,\mathbf{Y})\,\mathsf{C}\psi\right)(y)-\left(\bar{T}^{(2)}(-B,\mathbf{Y})\,\mathsf{C}\chi\right)(y)\\
&=-\mathsf{C}\cdot\left[\left(T^{(1)}(B,\mathbf{Y})\,\psi\right)(y)+\left(T^{(2)}(B,\mathbf{Y})\,\chi\right)(y)\right]\\
&=-\mathsf{C}\cdot\mathrm{sub}_{\mathbf{Y},B}\Psi(y)=\mathrm{sub}_{\mathbf{Y},B}\left(\mathrm{Cg}\,\bar{\Psi}(y)\right).
\end{align*}

\textbf{(f-4) The map $\mathbf{K}$.}
Let $\mu$ be real.
With $-B$ in place of $B$, (c-4) reads
\[
\overline{T^{(i)}(-\overline{B},Y)}=\varkappa^{-1}\,T^{(i)}(B,Y)\,\varkappa,\qquad\overline{\bar{T}^{(i)}(-\overline{B},Y)}=\varkappa^{-1}\,\bar{T}^{(i)}(B,Y)\,\varkappa.
\]
For a unit site $y$ of $\mathbf{Y}$, by (i23.7), by (i23.V2) for the antilinear map $\mathbf{K}$, and by $\varkappa^2=1$, $\varkappa^{-1}=\varkappa$,
\begin{align*}
\mathbf{K}\left(\mathrm{sub}_{\mathbf{Y},-\overline{B}}\Psi(y)\right)
&=\left(\left[\overline{T^{(1)}(-\overline{B},\mathbf{Y})}\right]\varkappa\psi\right)(y)+\left(\left[\overline{T^{(2)}(-\overline{B},\mathbf{Y})}\right]\varkappa\chi\right)(y)\\
&=\varkappa^{-1}\cdot\left[\left(T^{(1)}(B,\mathbf{Y})\,\varkappa\varkappa\psi\right)(y)+\left(T^{(2)}(B,\mathbf{Y})\,\varkappa\varkappa\chi\right)(y)\right]\\
&=\varkappa\cdot\mathrm{sub}_{\mathbf{Y},B}\Psi(y)=\mathrm{sub}_{\mathbf{Y},B}\left(\mathbf{K}\Psi(y)\right),\\
\mathbf{K}\left(\mathrm{sub}_{\mathbf{Y},-\overline{B}}\bar{\Psi}(y)\right)
&=\left(\left[\overline{\bar{T}^{(1)}(-\overline{B},\mathbf{Y})}\right]\varkappa\bar{\psi}\right)(y)+\left(\left[\overline{\bar{T}^{(2)}(-\overline{B},\mathbf{Y})}\right]\varkappa\bar{\chi}\right)(y)\\
&=\varkappa^{-1}\cdot\left[\left(\bar{T}^{(1)}(B,\mathbf{Y})\,\varkappa\varkappa\bar{\psi}\right)(y)+\left(\bar{T}^{(2)}(B,\mathbf{Y})\,\varkappa\varkappa\bar{\chi}\right)(y)\right]\\
&=\varkappa\cdot\mathrm{sub}_{\mathbf{Y},B}\bar{\Psi}(y)=\mathrm{sub}_{\mathbf{Y},B}\left(\mathbf{K}\bar{\Psi}(y)\right).
\end{align*}
\end{proof}

\paragraph{(g) The substitution: norm}

\begin{proof}[Proof of Proposition~23.1\,(g)]
Fix $\mu\in(0,1]$, $e_k\in(0,e_0]$, $B\in\tilde{R}_k+\mathfrak{X}_k$ and a polymer $\mathbf{Y}$; the constants are those of Section~I.8.
Clause (g) is the rule (N2) of Proposition~15\,(d), applied to the data described below.
The letters $p$, $q$, $\phi_i$, $\theta_j$, $\mathrm{Ind}_i$, $\nu$, $O_{ij}$, $C_{ij}$, $\mathsf{w}$, $\mathsf{w}'$ have the meaning that Proposition~15 gives them.

\emph{The families.}
Let $p=3$, and let
$\phi_1=\psi$ on $\mathrm{Ind}_1=\{\xi\}$, with $\nu=\eta_k^3/48$;
$\phi_2=\chi$ on $\mathrm{Ind}_2=\{\boldsymbol{\zeta}\}$, with $\nu=\left(\eta_k^3/48\right)^2$;
and let $\phi_3$ be the family of the generators $\Psi(u)$ on $\mathrm{Ind}_3=\{u=(y,\beta,\omega):\ y\text{ a unit site of }\mathbf{Y}\}$, with $\nu=1$
(Section~I.3.6; all indices are taken in a fundamental domain).
Then $\mathrm{Gr}[\phi_1,\phi_2,\phi_3]=\mathrm{Gr}[\psi,\chi,\Psi\,|\,\mathbf{Y}]$.
An element $F$ of this subalgebra has there the same coefficients as in $\mathrm{Gr}[\psi,\chi,\Psi]$;
consequently the norm (I.5.1) of $F$ for the weights $(\mathsf{w}_1,\mathsf{w}_2,\mathsf{w}_3)$ is $\|F\|_{(\mathsf{w}_1,\mathsf{w}_2;\mathsf{w}_3)}$.
Further, let $q=2$, $\theta_1=\psi$ and $\theta_2=\chi$, with the weights $\mathsf{w}'=(\mathsf{w}_1,\mathsf{w}_2)$.

\emph{The operators.}
Let $O_{11}=1$, $O_{22}=1$, $O_{12}=0$ and $O_{21}=0$.
For a function $\mathrm{hf}$ on $\{\xi\}$ and for $\omega\in\{0,1\}$ denote by $\mathrm{hf}^{\omega}$ the antiperiodic function $x\mapsto\left(\mathrm{hf}(x,\beta',\omega)\right)_{\beta'}$;
the function $O_{31}\mathrm{hf}$ on $\mathrm{Ind}_3$ takes the value $\left(T^{(1)}(B,\mathbf{Y})\,\mathrm{hf}^{0}\right)(y)_{\beta}$ at $(y,\beta,0)$ and the value $\left(\bar{T}^{(1)}(B,\mathbf{Y})\,\mathrm{hf}^{1}\right)(y)_{\beta}$ at $(y,\beta,1)$.
For a function $\mathrm{hf}$ on $\{\boldsymbol{\zeta}\}$ denote by $\mathrm{hf}^{\omega}$ the function of pairs $(x,x')\mapsto\left(\mathrm{hf}(x,x',\beta',\omega)\right)_{\beta'}$;
the function $O_{32}\mathrm{hf}$ on $\mathrm{Ind}_3$ takes the value $\left(T^{(2)}(B,\mathbf{Y})\,\mathrm{hf}^{0}\right)(y)_{\beta}$ at $(y,\beta,0)$ and the value $\left(\bar{T}^{(2)}(B,\mathbf{Y})\,\mathrm{hf}^{1}\right)(y)_{\beta}$ at $(y,\beta,1)$.
For these operators the homomorphism $\phi_i(t)\mapsto\sum_j\left(O_{ij}\theta_j\right)(t)$ of Proposition~15\,(d) fixes $\psi(\xi)$ and $\chi(\boldsymbol{\zeta})$ and sends $\Psi(y)$, $\bar{\Psi}(y)$ to the columns (i23.7);
it is therefore the homomorphism $\mathrm{sub}_{\mathbf{Y},B}$ of (0-6).

\emph{The row sums.}
One has $|O_{11}|_{1,\infty}=|O_{22}|_{1,\infty}=1$, since the kernel of the identity is $\mathbf{1}[t=s]/\nu(s)$.
Let $t=(y,\beta,\omega)\in\mathrm{Ind}_3$, and let $T$ denote $T^{(1)}(B,\mathbf{Y})$ if $\omega=0$ and $\bar{T}^{(1)}(B,\mathbf{Y})$ if $\omega=1$.
The row of $O_{31}$ at $t$ consists of those entries $T(y,\beta;t',\beta')$, in the sense of 0.2, of the row of $T$ at the unit site $y$ of $\mathbf{Y}$ whose first spinor index is the given $\beta$, the columns carrying the same bar index $\omega$.
In the same way, with $T$ now denoting $T^{(2)}(B,\mathbf{Y})$ if $\omega=0$ and $\bar{T}^{(2)}(B,\mathbf{Y})$ if $\omega=1$, the row of $O_{32}$ at $t$ consists of those entries $T(y,\beta;t',\beta')$ of the row of $T$ at $y$ whose first spinor index is $\beta$, with columns of the same bar index.
According to (b), every column index of a row at $y$ has sup-distance at most $(1+\eta_k)/2$ from $y$.
As $1+\eta_k<\ell$, no two of these indices are congruent modulo $\ell\mathbb{Z}^3$, and the replacement of the column indices by indices in a fundamental domain alters at most the signs of entries (I.AP).
The number $|O|_{1,\infty}$ of (D4) in Proposition~15 is the supremum over $t$ of $\sum_s|O(t,s)|\cdot\nu(s)$.
For $O=O_{31}$ and for $O=O_{32}$, at $t=(y,\beta,\omega)$ this sum is
\begin{equation*}
\sum_{(t',\beta')}\left|T(y,\beta;t',\beta')\right|\qquad(\beta\ \text{fixed}),
\end{equation*}
the entries being those of 0.2, that is, with the weights included;
the supremum is formed over all $(y,\beta,\omega)$.
Consequently, for $O_{31}$ the supremum over $(y,\beta)$ at a fixed $\omega$ is the number
\begin{equation*}
\sup_{y}\ \sup_{\beta}\ \sum_{\beta'}\int\mathrm{d}x'\,\left|T(y,x')_{\beta\beta'}\right|
\end{equation*}
of (I.ke3), formed with the kernels of $T$ with respect to $\int\mathrm{d}x'$;
for $O_{32}$ it is the number of Section~I.5.2, that is, the supremum over $(y,\beta)$ of the sum of the moduli of the coefficients of the values $F\left(x,x^{(3)}\right)_{\beta'}$ in $(TF)(y)_{\beta}$.
In Proposition~9 the operators $T^{(i)}_{\mathbf{Y}}(A)$, $\bar{T}^{(i)}(A)$ are those of (23.1).
Hence Proposition~9\,(d), in the form (G3.T-9) of Section~G3.0.4, yields
\begin{align*}
|O_{31}|_{1,\infty}&=\max\left(\left|T^{(1)}(B,\mathbf{Y})\right|_{1,\infty},\left|\bar{T}^{(1)}(B,\mathbf{Y})\right|_{1,\infty}\right)\le C_{T},\\
|O_{32}|_{1,\infty}&=\max\left(\left|T^{(2)}(B,\mathbf{Y})\right|_{1,\infty},\left|\bar{T}^{(2)}(B,\mathbf{Y})\right|_{1,\infty}\right)\le C_{T}.
\end{align*}

\emph{The condition of Proposition~15\,(d).}
Choose $C_{11}=C_{22}=1$, $C_{12}=C_{21}=0$ and $C_{31}=C_{32}=C_{T}$.
The condition $\sum_jC_{ij}\cdot\mathsf{w}'_j\le\mathsf{w}_i$ then reads
$\mathsf{w}_1\le\mathsf{w}_1$ for $i=1$,
$\mathsf{w}_2\le\mathsf{w}_2$ for $i=2$,
and $C_{T}\cdot\mathsf{w}_1+C_{T}\cdot\mathsf{w}_2\le\mathsf{w}_3$ for $i=3$;
the last inequality is the hypothesis of (g).
Proposition~15\,(d) gives
\begin{equation*}
\left\|\mathrm{sub}_{\mathbf{Y},B}\left(F\right)\right\|_{(\mathsf{w}_1,\mathsf{w}_2)}\le\|F\|_{(\mathsf{w}_1,\mathsf{w}_2;\mathsf{w}_3)}.\qedhere
\end{equation*}
\end{proof}

\subsection*{2. Proposition 23.2 (level $0$: the vertex of the first step)}

\subsubsection*{2.1. The statement}

The assertion to be proved is Proposition~23.2, for the data and with the definition (23.3) of $\mathrm{Vx}(\mathrm{Zf},\square)$ that precede it.
These data are the following:
$k=0$, so that $\eta_0=1$;
a side $\ell\in\mathrm{Cl}$;
a number $q>0$, the charge of level~$0$;
a complex bond function $\mathrm{Zf}$ on $T_0$;
and an $M$-cube $\square$ of $T_0$.
The element $\mathrm{Vx}(\mathrm{Zf},\square)$ of $\mathrm{Gr}[\psi,\chi]$ is defined by (23.3), and in this notation the display (I.5.5) takes the form (8.3$_0$)$^{\mathrm{ch}}$,
\begin{equation*}
G^{(1)}(s,\square)=\mathrm{Vx}\left(\mathscr{Z}(s'),\square\right).
\end{equation*}
For these data Proposition~23.2 asserts the following.
\begin{itemize}
\item[(a)] \emph{Identity.}
Let $\mathscr{T}$ be a complex bond function on $T_0$, let $\mathrm{ff}$ and $\overline{\mathrm{ff}}$ be fields on $T^{\mathrm{f}}_0$, and let $\mathrm{ev}$ be the homomorphism of $\mathrm{Gr}[\psi,\chi]$ with $\psi\mapsto\mathrm{ff}$, $\bar{\psi}\mapsto\overline{\mathrm{ff}}$, $\chi\mapsto\delta_{\alpha,\mathscr{T}}\,\mathrm{ff}$ and $\bar{\chi}\mapsto\delta_{\alpha,-\mathscr{T}}\,\overline{\mathrm{ff}}$.
Then the two halves of (23.3), each summed over the $M$-cubes, have the same image under $\mathrm{ev}$, and
\begin{equation*}
\mathrm{ev}\left(\sum_{\square}\mathrm{Vx}(\mathrm{Zf},\square)\right)=-\left\langle\overline{\mathrm{ff}},\left(D_{\mathscr{T}+\mathrm{Zf}}-D_{\mathscr{T}}\right)\mathrm{ff}\right\rangle^{\mathrm{ch}}.
\end{equation*}
\item[(b)] \emph{Four symmetries.}
For every $\tilde{g}\in\mathbb{G}_0$ and every real function $\Lambda$ of $c$,
\begin{gather*}
a_{\tilde{g}}\left[\mathrm{Vx}\left(\tilde{g}\,\mathrm{Zf},\tilde{g}\square\right)\right]=\mathrm{Vx}(\mathrm{Zf},\square),\qquad G_{\Lambda}\left[\mathrm{Vx}(\mathrm{Zf},\square)\right]=\mathrm{Vx}(\mathrm{Zf},\square),\\
\mathrm{Cg}\left[\mathrm{Vx}(-\mathrm{Zf},\square)\right]=\mathrm{Vx}(\mathrm{Zf},\square),\qquad\mathbf{K}\left[\mathrm{Vx}\left(-\overline{\mathrm{Zf}},\square\right)\right]=\mathrm{Vx}(\mathrm{Zf},\square).
\end{gather*}
The map $\mathrm{Cg}$ sends the first half of $\mathrm{Vx}(-\mathrm{Zf},\square)$ to the second half of $\mathrm{Vx}(\mathrm{Zf},\square)$, and the second half of $\mathrm{Vx}(-\mathrm{Zf},\square)$ to the first half of $\mathrm{Vx}(\mathrm{Zf},\square)$;
each of the other three maps preserves each half.
\item[(c)] \emph{Locality.}
Apart from $\psi$, $\bar{\psi}$ at the sites over $\square$ and $\chi$, $\bar{\chi}$ at the pairs $\left(x,x^{(3)}\right)$ with $c\in\square$, the element $\mathrm{Vx}(\mathrm{Zf},\square)$ contains no generators.
Each such pair is admissible for $\square$: it satisfies the condition of kind (II), and, if $c_3$ lies in $\square$, also that of kind (I).
The element reads $\mathrm{Zf}$ on the bonds $(c,c_3)$ with $c\in\square$, which are the oriented bonds whose initial point lies in $\square$;
these bonds together with their reverses $(c_3,c)$, on which $\mathrm{Zf}(c_3,c)=-\mathrm{Zf}(c,c_3)$, are the bonds that meet $\square$.
It contains neither a connection $A$ or $\mathscr{C}$ nor the mass;
each of its monomials has one barred and one unbarred generator;
and its coefficients are entire functions of the numbers $\mathrm{Zf}(b)$.
\item[(d)] \emph{Norm.}
For all weights $\mathsf{w}=(\mathsf{w}_1,\mathsf{w}_2)$,
\begin{multline*}
\left\|\mathrm{Vx}(\mathrm{Zf},\square)\right\|_{\mathsf{w}}\le12(1+\rho)\cdot M^3\cdot\mathsf{w}_1\cdot\left(\mathsf{w}_1+\mathsf{w}_2\right)\\
\cdot\sup\left\{\left|\exp\left(iq\,\mathrm{Zf}(b)\right)-1\right|:\ b\text{ an oriented bond that meets }\square\right\}.
\end{multline*}
If $q\le1$ and $q\cdot\left|\mathrm{Zf}(b)\right|\le\frac{1}{2}q^{1-\frac{\varepsilon}{2}}$ on these bonds, then
\begin{equation*}
\left\|\mathrm{Vx}(\mathrm{Zf},\square)\right\|_{h_0/2}\le6(1+\rho)\cdot M^3\cdot q^{\frac{1}{2}-\frac{\varepsilon}{2}}\le12\,M^3\cdot q^{\frac{1}{2}-\frac{\varepsilon}{2}},\qquad h_0/2=\left(\tfrac{1}{2}q^{-\frac{1}{4}},\ \tfrac{1}{2}q^{-\frac{1}{4}+\varepsilon}\right).
\end{equation*}
\item[(e)] \emph{In the step map.}
Consider the setting (ST1)$^{\mathrm{ch}}$ of Section~G3.4.0 at $k=0$.
For $s\in\mathrm{Pd}_0$ and $\mathsf{Z}\in\mathrm{Box}_0$ the field $\mathrm{Zf}=\mathscr{Z}(s')$ satisfies the condition of (d) (Proposition~24.3\,(q1));
the function $G^{(1)}(\,\cdot\,,\square)$ does not depend on $s''$, it is continuous in $(s',\mathsf{Z})$ on $\mathrm{Pd}_0\times\mathrm{Box}_0$, and it is analytic in $s'$ in the interior;
and
\begin{equation*}
\left\|G^{(1)}(s,\square)\right\|_{h_0/2}\le12\,M^3\cdot q^{\frac{1}{2}-\frac{\varepsilon}{2}}\le q^{\frac{1}{4}-5\varepsilon}.
\end{equation*}
\end{itemize}

\subsubsection*{2.2. Proof of Proposition 23.2}

\emph{Conventions of the proof.}
Throughout, $k=0$.
Thus $\eta_0=1$ and $\mathbb{L}_0=\mathbb{Z}^3$, the fine lattice coincides with the unit lattice, and $e_k$ is written $q$ ((I.L0.1), (I.L0.2)).
For a fermion site $x=(c,g)$ one has $x^{(3)}=(c+ge_1,gs_3)$ and $c_3=c+ge_1$ (I.1.2).
Bond functions on $T_0$ are regarded as $\ell\mathbb{Z}^3$-periodic bond functions on $\mathbb{Z}^3$, and $u_B(c,c_3)=\exp\left(iq\,B(c,c_3)\right)$ (I.1.6).

\paragraph{Step 0. The display (23.3)}

The two points of the pair $\left(x,x^{(3)}\right)$ are distinct, and $d_X\left(x,x^{(3)}\right)=\max\left(|c-c_3|,1\right)=1\le M$ by (I.1.5).
Hence the pair is an element of $\mathrm{Pzt}_0$ (Section~I.1.7), and $\chi\left(x,x^{(3)}\right)$, $\bar{\chi}\left(x,x^{(3)}\right)$ belong to the generators (Section~I.3.6).
Denote the two halves of (23.3) by
\begin{align*}
\mathrm{Vx}'(\mathrm{Zf},\square)&:=\frac{1}{2}\int_{\square}\mathrm{d}x\,\bar{\psi}(x)\cdot M_3(x)\left(\exp\left(iq\,\mathrm{Zf}(c,c_3)\right)-1\right)\cdot\left(\psi(x)+\chi\left(x,x^{(3)}\right)\right),\\
\mathrm{Vx}''(\mathrm{Zf},\square)&:=\frac{1}{2}\int_{\square}\mathrm{d}x\,\left(\bar{\psi}(x)+\bar{\chi}\left(x,x^{(3)}\right)\right)\cdot M_3\left(x^{(3)}\right)\left(\exp\left(-iq\,\mathrm{Zf}(c,c_3)\right)-1\right)\cdot\psi(x),
\end{align*}
so that $\mathrm{Vx}=\mathrm{Vx}'+\mathrm{Vx}''$.
Using 0.1 and (I.2.1), one writes them as sums over the $48\,M^3$ fermion sites over $\square$:
\begin{align}
\mathrm{Vx}'(\mathrm{Zf},\square)
&=\frac{1}{96}\sum_{c\in\square}\sum_{g\in O_{\mathrm{cub}}}\sum_{\beta,\beta'}3\left(\rho-\gamma\cdot(ge_1)\right)_{\beta\beta'}\left(\exp\left(iq\,\mathrm{Zf}(c,c+ge_1)\right)-1\right)\notag\\
&\qquad\times\bar{\psi}_{\beta}(c,g)\left[\psi_{\beta'}(c,g)+\chi_{\beta'}\left((c,g),(c+ge_1,gs_3)\right)\right],\notag\\
\mathrm{Vx}''(\mathrm{Zf},\square)
&=\frac{1}{96}\sum_{c\in\square}\sum_{g\in O_{\mathrm{cub}}}\sum_{\beta,\beta'}3\left(\rho+\gamma\cdot(ge_1)\right)_{\beta\beta'}\left(\exp\left(-iq\,\mathrm{Zf}(c,c+ge_1)\right)-1\right)\notag\\
&\qquad\times\left[\bar{\psi}_{\beta}(c,g)+\bar{\chi}_{\beta}\left((c,g),(c+ge_1,gs_3)\right)\right]\psi_{\beta'}(c,g).\tag{i23.9}
\end{align}
For each of the six vectors $\pm e_1$, $\pm e_2$, $\pm e_3$ there are $8$ among the $48$ labels $g$ for which $ge_1$ is that vector (I.s4).
The photon site $c$ ranges over a lift of $\square$, which is an $M$-cube of $\mathbb{Z}^3$.
A translation of the lift by $\ell e_{\mu}$ reverses the sign of both generators in every term (I.AP) and changes neither $\mathrm{Zf}$ nor the matrices;
so the sums do not depend on the choice of the lift.
No bond of type $0$, $1$ or $2$ occurs.

At $\mathrm{Zf}=\mathscr{Z}(s')$ the display (23.3) is the display (I.5.5).
Indeed, in (I.5.5) the generators $\bar{\psi}$, $\bar{\chi}$ are rows, so that the expression $\bar{\psi}(x)\,M_3(x)\,(\,\cdot\,)$ there is the expression $\bar{\psi}(x)\cdot M_3(x)\cdot(\,\cdot\,)$ here;
the phases, the matrices $M_3(x)$ and $M_3\left(x^{(3)}\right)$, the factor $\frac{1}{2}$ and $\int_{\square}\mathrm{d}x$ agree.
This is the relation (8.3$_0$)$^{\mathrm{ch}}$.

\paragraph{(a) Identity}

\begin{proof}[Proof of Proposition~23.2\,(a)]
The values of $\mathrm{ff}$ and $\overline{\mathrm{ff}}$ are odd, and these fields are antiperiodic in the same way as $\psi$ and $\bar{\psi}$;
likewise $\delta_{\alpha,\mathscr{T}}\mathrm{ff}$ and $\delta_{\alpha,-\mathscr{T}}\overline{\mathrm{ff}}$ are functions of pairs (Section~I.1.8) whose values are odd.
The existence and the uniqueness of $\mathrm{ev}$ follow from (i23.H1) and (i23.H3).
By (I.L0.14), which is (I.1.7) with $d_X\left(x,x^{(3)}\right)=1$ and $U_B\left(x,x^{(3)}\right)=u_B(c,c_3)$, and by $u_{-\mathscr{T}}=u_{\mathscr{T}}^{-1}$, one has
\begin{align*}
\mathrm{ff}(x)+\left(\delta_{\alpha,\mathscr{T}}\mathrm{ff}\right)\left(x,x^{(3)}\right)&=u_{\mathscr{T}}(c,c_3)\,\mathrm{ff}\left(x^{(3)}\right),\\
\overline{\mathrm{ff}}(x)+\left(\delta_{\alpha,-\mathscr{T}}\overline{\mathrm{ff}}\right)\left(x,x^{(3)}\right)&=u_{\mathscr{T}}(c,c_3)^{-1}\,\overline{\mathrm{ff}}\left(x^{(3)}\right).
\end{align*}
The $M$-cubes form a partition of the sites of $T_0$ \cite[Section~17.4.1, (g0)]{QEDI};
hence $\sum_{\square}\int_{\square}=\int$, the integral over a fundamental domain.
Since $\mathrm{ev}$ is linear, it follows that
\begin{align}
\mathrm{ev}\left(\sum_{\square}\mathrm{Vx}'(\mathrm{Zf},\square)\right)&=\frac{1}{2}\int\mathrm{d}x\,\overline{\mathrm{ff}}(x)\cdot M_3(x)\left(\exp\left(iq\,\mathrm{Zf}(c,c_3)\right)-1\right)u_{\mathscr{T}}(c,c_3)\cdot\mathrm{ff}\left(x^{(3)}\right),\tag{i23.10}\\
\mathrm{ev}\left(\sum_{\square}\mathrm{Vx}''(\mathrm{Zf},\square)\right)&=\frac{1}{2}\int\mathrm{d}x\,u_{\mathscr{T}}(c,c_3)^{-1}\,\overline{\mathrm{ff}}\left(x^{(3)}\right)\cdot M_3\left(x^{(3)}\right)\left(\exp\left(-iq\,\mathrm{Zf}(c,c_3)\right)-1\right)\cdot\mathrm{ff}(x).\tag{i23.11}
\end{align}
As functions of $x\in\mathbb{L}^{\mathrm{f}}_0$ the two integrands are $\ell\mathbb{Z}^3$-periodic:
each of them contains two antiperiodic factors, $\mathscr{T}$ and $\mathrm{Zf}$ are periodic, and $(x+v)^{(3)}=x^{(3)}+v$ for $v\in\ell\mathbb{Z}^3$ (I.s1).
The map $x\mapsto x^{(3)}$ is a bijection of $T^{\mathrm{f}}_0$ onto itself with $\left(x^{(3)}\right)^{(3)}=x$ (I.s1), and all sites carry the weight $1/48$.
Substituting $x^{(3)}$ for the integration variable $x$ in (i23.11) turns $x^{(3)}$, $c$, $c_3$ into $x$, $c_3$, $c$; this gives
\begin{equation*}
\mathrm{ev}\left(\sum_{\square}\mathrm{Vx}''(\mathrm{Zf},\square)\right)=\frac{1}{2}\int\mathrm{d}x\,u_{\mathscr{T}}(c_3,c)^{-1}\,\overline{\mathrm{ff}}(x)\cdot M_3(x)\left(\exp\left(-iq\,\mathrm{Zf}(c_3,c)\right)-1\right)\cdot\mathrm{ff}\left(x^{(3)}\right).
\end{equation*}
As $\mathscr{T}(c_3,c)=-\mathscr{T}(c,c_3)$ and $\mathrm{Zf}(c_3,c)=-\mathrm{Zf}(c,c_3)$ (Section~I.1.1), one has $u_{\mathscr{T}}(c_3,c)^{-1}=u_{\mathscr{T}}(c,c_3)$ and $\exp\left(-iq\,\mathrm{Zf}(c_3,c)\right)=\exp\left(iq\,\mathrm{Zf}(c,c_3)\right)$.
The right side is therefore that of (i23.10): the two halves, summed over the $M$-cubes, have the same image.
In this computation no two factors have been interchanged.

The relation (I.O6) at $\eta_0=1$, valid for fields by (i23.V1), gives
\begin{equation*}
\left(\left(D_{\mathscr{T}+\mathrm{Zf}}-D_{\mathscr{T}}\right)\mathrm{ff}\right)(x)=-M_3(x)\left(u_{\mathscr{T}+\mathrm{Zf}}-u_{\mathscr{T}}\right)(c,c_3)\,\mathrm{ff}\left(x^{(3)}\right),
\end{equation*}
and $u_{\mathscr{T}+\mathrm{Zf}}(c,c_3)=\exp\left(iq\,\mathrm{Zf}(c,c_3)\right)\cdot u_{\mathscr{T}}(c,c_3)$.
Consider the pairing $\langle f',f\rangle^{\mathrm{ch}}=\int\mathrm{d}x\,f'(x)\cdot f(x)$ of Section~I.1.6, formed for fields with the factors in the order in which they are written (Section~G3.1.0).
Then
\begin{equation*}
-\left\langle\overline{\mathrm{ff}},\left(D_{\mathscr{T}+\mathrm{Zf}}-D_{\mathscr{T}}\right)\mathrm{ff}\right\rangle^{\mathrm{ch}}=\int\mathrm{d}x\,\overline{\mathrm{ff}}(x)\cdot M_3(x)\left(\exp\left(iq\,\mathrm{Zf}(c,c_3)\right)-1\right)u_{\mathscr{T}}(c,c_3)\cdot\mathrm{ff}\left(x^{(3)}\right),
\end{equation*}
which is twice the right side of (i23.10), and hence the sum of the two images.
\end{proof}

\paragraph{(b) Four symmetries}

\begin{proof}[Proof of Proposition~23.2\,(b)]
\textbf{(b-1) The map $a_{\tilde{g}}$.}
Let $\tilde{g}=(S,a)\in\mathbb{G}_0$ (Section~I.4.2), let $r=r(S)$ (Section~I.4.1), and let $a\in M\mathbb{Z}^3$ be a representative.
(If $a$ is replaced by another representative $a+\ell n$, $n\in\mathbb{Z}^3$, then $\tilde{g}^{-1}x$ is replaced by $\tilde{g}^{-1}x-\ell rn$, and by (I.AP) the image of every generator is multiplied by $(-1)^{n_1+n_2+n_3}$.
Every term below contains two generators, so that its image does not change.)
Under $c\mapsto\tilde{g}c=rc+a$ the $M$-cube with centre $z$ is mapped onto the $M$-cube whose centre is $\tilde{g}z\in M\mathbb{Z}^3$, and $(\tilde{g}\,\mathrm{Zf})(c,c')=\mathrm{Zf}\left(\tilde{g}^{-1}c,\tilde{g}^{-1}c'\right)$ (Section~I.4.3).
Since $\tilde{g}^{-1}$ is an element of $\mathbb{G}_0$, (I.a2) gives $\tilde{g}^{-1}\left(x^{(3)}\right)=\left(\tilde{g}^{-1}x\right)^{(3)}$, and therefore $\tilde{g}^{-1}c_3=\pi\left(\left(\tilde{g}^{-1}x\right)^{(3)}\right)$.
By (I.T-g) of Section~I.4.5 and by (i23.R1) with $m_1=(S^{-1})^{T}$, $m_2=S$, one has for every $4\times4$ matrix $m$
\begin{align*}
a_{\tilde{g}}\left[\bar{\psi}(x)\cdot m\cdot\left(\psi(x)+\chi\left(x,x^{(3)}\right)\right)\right]
&=\bar{\psi}\left(\tilde{g}^{-1}x\right)\cdot\left(S^{-1}\cdot m\cdot S\right)\cdot\left(\psi\left(\tilde{g}^{-1}x\right)+\chi\left(\tilde{g}^{-1}x,\left(\tilde{g}^{-1}x\right)^{(3)}\right)\right),\\
a_{\tilde{g}}\left[\left(\bar{\psi}(x)+\bar{\chi}\left(x,x^{(3)}\right)\right)\cdot m\cdot\psi(x)\right]
&=\left(\bar{\psi}\left(\tilde{g}^{-1}x\right)+\bar{\chi}\left(\tilde{g}^{-1}x,\left(\tilde{g}^{-1}x\right)^{(3)}\right)\right)\cdot\left(S^{-1}\cdot m\cdot S\right)\cdot\psi\left(\tilde{g}^{-1}x\right).
\end{align*}
By (I.m4) at the sites $\tilde{g}^{-1}x$ and $\left(\tilde{g}^{-1}x\right)^{(3)}$,
\begin{equation*}
S^{-1}\cdot M_3(x)\cdot S=M_3\left(\tilde{g}^{-1}x\right),\qquad S^{-1}\cdot M_3\left(x^{(3)}\right)\cdot S=M_3\left(\left(\tilde{g}^{-1}x\right)^{(3)}\right).
\end{equation*}
So $a_{\tilde{g}}$ maps the integrand of $\mathrm{Vx}'(\tilde{g}\,\mathrm{Zf},\tilde{g}\square)$ at $x$ to the integrand of $\mathrm{Vx}'(\mathrm{Zf},\square)$ at $\tilde{g}^{-1}x$, and the integrand of $\mathrm{Vx}''(\tilde{g}\,\mathrm{Zf},\tilde{g}\square)$ at $x$ to the integrand of $\mathrm{Vx}''(\mathrm{Zf},\square)$ at $\tilde{g}^{-1}x$.
The map $x\mapsto\tilde{g}^{-1}x$ is a bijection of $\tilde{g}\square\times O_{\mathrm{cub}}$ onto $\square\times O_{\mathrm{cub}}$ (I.a1), and all weights equal $1/48$.
Consequently
\begin{equation*}
a_{\tilde{g}}\left[\mathrm{Vx}'(\tilde{g}\,\mathrm{Zf},\tilde{g}\square)\right]=\mathrm{Vx}'(\mathrm{Zf},\square),\qquad a_{\tilde{g}}\left[\mathrm{Vx}''(\tilde{g}\,\mathrm{Zf},\tilde{g}\square)\right]=\mathrm{Vx}''(\mathrm{Zf},\square).
\end{equation*}

\textbf{(b-2) The map $G_{\Lambda}$.}
Every term of (i23.9) is a number times the product of one barred generator, $\bar{\psi}_{\beta}(x)$ or $\bar{\chi}_{\beta}\left(x,x^{(3)}\right)$, and one unbarred generator, $\psi_{\beta'}(x)$ or $\chi_{\beta'}\left(x,x^{(3)}\right)$, both of which have the same site $x$ as their site, respectively as their first point.
By (I.T-G) the first factor is multiplied by $\exp\left(-iq\,\Lambda(c)\right)$ and the second by $\exp\left(iq\,\Lambda(c)\right)$.
Thus every term is invariant under $G_{\Lambda}$, and so are $\mathrm{Vx}'(\mathrm{Zf},\square)$ and $\mathrm{Vx}''(\mathrm{Zf},\square)$.

\textbf{(b-3) The map $\mathrm{Cg}$.}
Being odd, the entries of $\psi$, $\chi$, $\bar{\psi}$, $\bar{\chi}$ anticommute with each other.
Let $m$ be a $4\times4$ matrix, and put $v:=\bar{\psi}(x)+\bar{\chi}\left(x,x^{(3)}\right)$ and $v':=\psi(x)+\chi\left(x,x^{(3)}\right)$.
By (I.T-C), (i23.R1), (i23.R2) and $\mathsf{C}^{T}=\mathsf{C}^{-1}$,
\begin{align*}
\mathrm{Cg}\left[\bar{\psi}(x)\cdot m\cdot v'\right]
&=\left(-\mathsf{C}\psi(x)\right)\cdot m\cdot\left(\mathsf{C}v\right)
=-\psi(x)\cdot\left(\mathsf{C}^{T}\cdot m\cdot\mathsf{C}\right)\cdot v\\
&=v\cdot\left(\mathsf{C}^{T}\cdot m\cdot\mathsf{C}\right)^{T}\cdot\psi(x)
=v\cdot\left(\mathsf{C}^{-1}\cdot m^{T}\cdot\mathsf{C}\right)\cdot\psi(x),\\
\mathrm{Cg}\left[v\cdot m\cdot\psi(x)\right]
&=\left(-\mathsf{C}v'\right)\cdot m\cdot\left(\mathsf{C}\bar{\psi}(x)\right)
=-v'\cdot\left(\mathsf{C}^{T}\cdot m\cdot\mathsf{C}\right)\cdot\bar{\psi}(x)
=\bar{\psi}(x)\cdot\left(\mathsf{C}^{-1}\cdot m^{T}\cdot\mathsf{C}\right)\cdot v'.
\end{align*}
By (I.m2) at $x$, and at $x^{(3)}$ together with $\left(x^{(3)}\right)^{(3)}=x$,
\begin{equation*}
\mathsf{C}^{-1}\cdot M_3(x)^{T}\cdot\mathsf{C}=M_3\left(x^{(3)}\right),\qquad\mathsf{C}^{-1}\cdot M_3\left(x^{(3)}\right)^{T}\cdot\mathsf{C}=M_3(x).
\end{equation*}
Under $\mathrm{Cg}$ the integrand $\frac{1}{2}\bar{\psi}(x)\cdot M_3(x)\left(\exp\left(-iq\,\mathrm{Zf}(c,c_3)\right)-1\right)\cdot v'$ of $\mathrm{Vx}'(-\mathrm{Zf},\square)$ goes over into $\frac{1}{2}v\cdot M_3\left(x^{(3)}\right)\left(\exp\left(-iq\,\mathrm{Zf}(c,c_3)\right)-1\right)\cdot\psi(x)$, the integrand of $\mathrm{Vx}''(\mathrm{Zf},\square)$;
likewise the integrand $\frac{1}{2}v\cdot M_3\left(x^{(3)}\right)\left(\exp\left(iq\,\mathrm{Zf}(c,c_3)\right)-1\right)\cdot\psi(x)$ of $\mathrm{Vx}''(-\mathrm{Zf},\square)$ goes over into $\frac{1}{2}\bar{\psi}(x)\cdot M_3(x)\left(\exp\left(iq\,\mathrm{Zf}(c,c_3)\right)-1\right)\cdot v'$, the integrand of $\mathrm{Vx}'(\mathrm{Zf},\square)$.
It follows that $\mathrm{Cg}\left[\mathrm{Vx}'(-\mathrm{Zf},\square)\right]=\mathrm{Vx}''(\mathrm{Zf},\square)$ and $\mathrm{Cg}\left[\mathrm{Vx}''(-\mathrm{Zf},\square)\right]=\mathrm{Vx}'(\mathrm{Zf},\square)$, and hence $\mathrm{Cg}\left[\mathrm{Vx}(-\mathrm{Zf},\square)\right]=\mathrm{Vx}(\mathrm{Zf},\square)$.

\textbf{(b-4) The map $\mathbf{K}$.}
The antilinear map $\mathbf{K}$ conjugates every coefficient and multiplies each of the four fields $\psi$, $\bar{\psi}$, $\chi$, $\bar{\chi}$ by $\varkappa$ (I.T-K).
Since $q$ is real,
\begin{equation*}
\overline{\exp\left(\pm iq\left(-\overline{\mathrm{Zf}(c,c_3)}\right)\right)-1}=\exp\left(\pm iq\,\mathrm{Zf}(c,c_3)\right)-1.
\end{equation*}
By (i23.R1) and (i23.C3), for any two columns $v_1$, $v_2$,
\begin{equation*}
\left(\varkappa v_1\right)\cdot\overline{M_3(\,\cdot\,)}\cdot\left(\varkappa v_2\right)=v_1\cdot\left(\varkappa^{T}\cdot\overline{M_3(\,\cdot\,)}\cdot\varkappa\right)\cdot v_2=v_1\cdot M_3(\,\cdot\,)\cdot v_2
\end{equation*}
at $x$ and at $x^{(3)}$.
Therefore $\mathbf{K}$ maps the integrand of $\mathrm{Vx}'(-\overline{\mathrm{Zf}},\square)$ to that of $\mathrm{Vx}'(\mathrm{Zf},\square)$, and the integrand of $\mathrm{Vx}''(-\overline{\mathrm{Zf}},\square)$ to that of $\mathrm{Vx}''(\mathrm{Zf},\square)$;
the weight $1/96$ is real.
\end{proof}

\paragraph{(c) Locality}

\begin{proof}[Proof of Proposition~23.2\,(c)]
Everything is read off (i23.9).
The generators appearing there are $\bar{\psi}_{\beta}(x)$, $\psi_{\beta'}(x)$ with $x$ over $\square$, and $\chi_{\beta'}\left(x,x^{(3)}\right)$, $\bar{\chi}_{\beta}\left(x,x^{(3)}\right)$ with $c\in\square$.

More precisely, in the sense of Section~I.5.1 the element $\mathrm{Vx}(\mathrm{Zf},\square)$ contains no other generator.
By (d-1) below, the products $\bar{\psi}_{\beta}(x)\,\chi_{\beta}\left(x,x^{(3)}\right)$ and $\bar{\chi}_{\beta}\left(x,x^{(3)}\right)\psi_{\beta}(x)$ have, up to signs, the coefficients
\begin{equation*}
\frac{1}{96}M_3(x)_{\beta\beta}\left(\exp\left(iq\,\mathrm{Zf}(c,c_3)\right)-1\right)\qquad\text{and}\qquad\frac{1}{96}M_3\left(x^{(3)}\right)_{\beta\beta}\left(\exp\left(-iq\,\mathrm{Zf}(c,c_3)\right)-1\right).
\end{equation*}
The diagonal entries of $M_3(x)$ and of $M_3\left(x^{(3)}\right)$ are $3\rho$ or $3(\rho\pm1)$ (Section~I.2.1, (I.2.1)), and they are different from $0$ because $0<\rho<1$.
Consequently, if $\exp\left(iq\,\mathrm{Zf}(c,c_3)\right)\neq1$, the sixteen generators $\psi_{\beta}(x)$, $\bar{\psi}_{\beta}(x)$, $\chi_{\beta}\left(x,x^{(3)}\right)$, $\bar{\chi}_{\beta}\left(x,x^{(3)}\right)$ at a site $x$ over $\square$ are all contained in $\mathrm{Vx}(\mathrm{Zf},\square)$.
If $\exp\left(iq\,\mathrm{Zf}(c,c_3)\right)=1$, none of them is contained in it, because they occur only in the products of (d-1) at $x$, whose coefficients then vanish.
The function $\mathrm{Zf}\mapsto\mathrm{Vx}(\mathrm{Zf},\square)$ contains all of them.

The second point of the pair $\left(x,x^{(3)}\right)$ is $x^{(3)}$, and $c\in\square$; so the pair satisfies the condition of kind (II) for $\square$ (Section~I.1.7).
Should $c_3$ lie in $\square$ as well, the sites $c$ and $c_3$ are nearest neighbours, so that each of the six paths $\boldsymbol{\Gamma}^{\pi}(c,c_3)$ reduces to the bond $(c,c_3)$, whose two sites lie in $\square$;
then the pair fulfils the condition of kind (I) too.
In either case the pair is admissible for $\square$.

Of the bond function $\mathrm{Zf}$ only the numbers $\mathrm{Zf}(c,c_3)$, $c\in\square$, enter:
two bond functions that agree on these bonds give the same element $\mathrm{Vx}(\mathrm{Zf},\square)$, which means that $\mathrm{Vx}(\mathrm{Zf},\square)$ reads $\mathrm{Zf}$ on these bonds in the sense of \cite[Definition~5.5\,(1)]{QEDI}, as in 0.2.
For $c\in\square$, as $g$ ranges over $O_{\mathrm{cub}}$, the sites $c_3=c+ge_1$ exhaust the six nearest neighbours of $c$ (I.s4);
consequently the oriented bonds with initial point in $\square$ are exactly the bonds $(c,c_3)$, $c\in\square$.
Since a bond meets $\square$ when one of its end points lies in $\square$ (Section~I.1.1), and since $\mathrm{Zf}(c',c)=-\mathrm{Zf}(c,c')$, these bonds are, apart from the orientation, the bonds that meet $\square$.

Apart from the generators, the display (i23.9) is built from the numbers $1/96$, $\rho$, $q$, the matrices $\gamma_{\mu}$ and the bond function $\mathrm{Zf}$, and from nothing else;
in particular it contains neither a connection $A$ or $\mathscr{C}$ nor the mass.
Each term is a number times a barred generator times an unbarred generator.
Each coefficient is a finite sum of products of constants with functions $\exp\left(\pm iq\,\mathrm{Zf}(b)\right)-1$, and is therefore an entire function of the numbers $\mathrm{Zf}(b)$.
\end{proof}

\paragraph{(d) Norm}

\begin{proof}[Proof of Proposition~23.2\,(d)]
\textbf{(d-1) The coefficients.}
Put
\begin{equation*}
\mathrm{dq}_x:=\exp\left(iq\,\mathrm{Zf}(c,c_3)\right)-1,\qquad\mathrm{dq}'_x:=\exp\left(-iq\,\mathrm{Zf}(c,c_3)\right)-1=\exp\left(iq\,\mathrm{Zf}(c_3,c)\right)-1.
\end{equation*}
Collecting the terms of (i23.9), one finds that for every $x$ over $\square$ and all $\beta,\beta'\in\{1,\dots,4\}$ the element $\mathrm{Vx}(\mathrm{Zf},\square)$ contains the three products
\begin{align*}
\bar{\psi}_{\beta}(x)\,\psi_{\beta'}(x)\quad&\text{with the coefficient}\quad\frac{1}{96}\left[M_3(x)_{\beta\beta'}\cdot\mathrm{dq}_x+M_3\left(x^{(3)}\right)_{\beta\beta'}\cdot\mathrm{dq}'_x\right],\\
\bar{\psi}_{\beta}(x)\,\chi_{\beta'}\left(x,x^{(3)}\right)\quad&\text{with the coefficient}\quad\frac{1}{96}M_3(x)_{\beta\beta'}\cdot\mathrm{dq}_x,\\
\bar{\chi}_{\beta}\left(x,x^{(3)}\right)\psi_{\beta'}(x)\quad&\text{with the coefficient}\quad\frac{1}{96}M_3\left(x^{(3)}\right)_{\beta\beta'}\cdot\mathrm{dq}'_x.
\end{align*}
Since a lift of $\square$ has sup-diameter $M-1$ and $M-1+\eta_0=M<\ell$ (Section~I.1.1), the covering map is injective on it.
Hence the $48\,M^3$ sites $x$ are distinct sites of $T^{\mathrm{f}}_0$, and the pairs $\left(x,x^{(3)}\right)$ are distinct elements of $\mathrm{Pz}_0$ (Section~I.1.7), their first points being distinct.
So the $3\cdot16\cdot48\,M^3$ products above are products of two distinct generators, and no two of them consist of the same two generators.
Up to signs, which come from the order of the two factors and from (I.AP), the numbers displayed are the coefficients $c_F$ of (I.5.1), and all other coefficients vanish.
In the norm, each of $\psi$, $\bar{\psi}$ carries the weight $\mathsf{w}_1$ and each of $\chi$, $\bar{\chi}$ the weight $\mathsf{w}_2$.
Therefore
\begin{align*}
\left\|\mathrm{Vx}(\mathrm{Zf},\square)\right\|_{\mathsf{w}}
&=\frac{1}{96}\sum_{x\text{ over }\square}\sum_{\beta,\beta'}\Bigl\{\left|M_3(x)_{\beta\beta'}\cdot\mathrm{dq}_x+M_3\left(x^{(3)}\right)_{\beta\beta'}\cdot\mathrm{dq}'_x\right|\cdot\mathsf{w}_1^2\\
&\qquad+\left|M_3(x)_{\beta\beta'}\right|\cdot\left|\mathrm{dq}_x\right|\cdot\mathsf{w}_1\cdot\mathsf{w}_2+\left|M_3\left(x^{(3)}\right)_{\beta\beta'}\right|\cdot\left|\mathrm{dq}'_x\right|\cdot\mathsf{w}_1\cdot\mathsf{w}_2\Bigr\}\\
&\le\frac{1}{96}\sum_{x\text{ over }\square}\Bigl\{\left|\mathrm{dq}_x\right|\sum_{\beta,\beta'}\left|M_3(x)_{\beta\beta'}\right|+\left|\mathrm{dq}'_x\right|\sum_{\beta,\beta'}\left|M_3\left(x^{(3)}\right)_{\beta\beta'}\right|\Bigr\}\cdot\mathsf{w}_1\cdot\left(\mathsf{w}_1+\mathsf{w}_2\right).
\end{align*}

\textbf{(d-2) The count.}
Applied at $x$ and at $x^{(3)}$, (I.m7) bounds each of the two sums over $(\beta,\beta')$ by $12(1+\rho)$.
Indeed, $M_3=3(\rho\mp\gamma_{\mu})$ with the matrices $\gamma_{\mu}$ of Section~I.2.1.
For $\mu=1,2$ the matrix $\gamma_{\mu}$ has a vanishing diagonal and exactly one entry of modulus $1$ in each row;
so $M_3$ has four entries $3\rho$ and four entries of modulus $3$, which add up to $12(1+\rho)$.
For $\mu=3$ one has $\gamma_3=\mathrm{diag}(1,-1,-1,1)$, and the moduli of the entries are $3(1-\rho)$ and $3(1+\rho)$, each twice, which add up to $12$, since $0<\rho<1$ (Section~I.0.5).
The bonds $(c,c_3)$ and $(c_3,c)$ are oriented bonds that meet $\square$; therefore $\left|\mathrm{dq}_x\right|,\left|\mathrm{dq}'_x\right|\le\mathrm{sz}$, where
\begin{equation*}
\mathrm{sz}:=\sup\left\{\left|\exp\left(iq\,\mathrm{Zf}(b)\right)-1\right|:\ b\text{ an oriented bond that meets }\square\right\}.
\end{equation*}
Since there are $48\,M^3$ sites $x$ over $\square$, it follows that
\begin{align*}
\left\|\mathrm{Vx}(\mathrm{Zf},\square)\right\|_{\mathsf{w}}
&\le\frac{1}{96}\cdot48\,M^3\cdot2\cdot12(1+\rho)\cdot\mathrm{sz}\cdot\mathsf{w}_1\cdot\left(\mathsf{w}_1+\mathsf{w}_2\right)\\
&=12(1+\rho)\cdot M^3\cdot\mathsf{w}_1\cdot\left(\mathsf{w}_1+\mathsf{w}_2\right)\cdot\mathrm{sz}.
\end{align*}

\textbf{(d-3) The weights $h_0/2$.}
Let $q\le1$, and let $q\cdot|\mathrm{Zf}(b)|\le\frac{1}{2}q^{1-\frac{\varepsilon}{2}}$ on the bonds $b$ that meet $\square$; here $0<\varepsilon\le1/200$ (Section~I.8.2, line (C7)).
For a complex number $t$ with $|t|\le\frac{1}{2}$,
\begin{equation*}
\left|\exp(t)-1\right|\le\sum_{j\ge1}\frac{|t|^j}{j!}\le|t|\cdot\exp(|t|)\le\mathrm{e}^{\frac{1}{2}}\,|t|\le2\,|t|,
\end{equation*}
because $\mathrm{e}<4$.
This is applied at $t=iq\,\mathrm{Zf}(b)$, for which $|t|\le\frac{1}{2}q^{1-\frac{\varepsilon}{2}}\le\frac{1}{2}$, since $q\le1$ and $1-\frac{\varepsilon}{2}>0$.
Thus $\mathrm{sz}\le q^{1-\frac{\varepsilon}{2}}$.
For $\mathsf{w}=h_0/2$ one has
\begin{equation*}
\mathsf{w}_1\cdot\left(\mathsf{w}_1+\mathsf{w}_2\right)=\frac{1}{4}q^{-\frac{1}{2}}\left(1+q^{\varepsilon}\right)\le\frac{1}{2}q^{-\frac{1}{2}},
\end{equation*}
because $q^{\varepsilon}\le1$.
Together with (d-2) and $\rho<1$ this gives
\begin{equation*}
\left\|\mathrm{Vx}(\mathrm{Zf},\square)\right\|_{h_0/2}\le12(1+\rho)\cdot M^3\cdot\frac{1}{2}q^{-\frac{1}{2}}\cdot q^{1-\frac{\varepsilon}{2}}=6(1+\rho)\cdot M^3\cdot q^{\frac{1}{2}-\frac{\varepsilon}{2}}\le12\,M^3\cdot q^{\frac{1}{2}-\frac{\varepsilon}{2}}.\qedhere
\end{equation*}
\end{proof}

\paragraph{(e) In the step map}

\begin{proof}[Proof of Proposition~23.2\,(e)]
Consider the setting (ST1)$^{\mathrm{ch}}$ of Section~G3.4.0 at $k=0$.
In it $q>0$ and $L^{\frac{1}{2}}q\le e_0$, so that $q<e_0$;
a system of primitives of level $0$ with the properties W1--W8, W10--W17 for $\mathrm{par}$ is given (these are the properties 1--8 and 10--17 of \cite[Section~5.4.2]{QEDI});
moreover $\mathscr{Z}(s')=H_0(s')\cdot C\cdot C_{0}^{1/2}(s')\cdot\mathsf{Z}$ (Definition~24.D, (D2)) and $G^{(1)}(s,\square)=\mathrm{Vx}\left(\mathscr{Z}(s'),\square\right)$.
Since the conditions on $e_0$ are required for every $v\ge-\log e_0$ (Proposition~14.2\,(b)), they hold in particular at $v=-\log q$.

\textbf{(e-1) The hypothesis of (d).}
From the second part of (A3), $96\,M^3\cdot q^{\frac{1}{4}}\le1$, one gets $q\le1$.
By Proposition~24.3\,(q1), for $s\in\mathrm{Pd}_0$ and $\mathsf{Z}\in\mathrm{Box}_0$ the three norms of $\mathscr{Z}(s')$ are at most $\frac{1}{2}q^{-\frac{\varepsilon}{2}}$;
the first of them is the sup norm $\left|\mathscr{Z}(s')\right|_{\infty}$ over the bonds (Section~I.5.3; (G2.r4) of Section~G2.0.5).
Hence $q\cdot\left|\mathscr{Z}(s')(b)\right|\le\frac{1}{2}q^{1-\frac{\varepsilon}{2}}$ for every bond $b$, and $\mathrm{Zf}=\mathscr{Z}(s')$ satisfies the hypothesis of (d).

\textbf{(e-2) The bound.}
By (d-3), $\left\|G^{(1)}(s,\square)\right\|_{h_0/2}\le12\,M^3\cdot q^{\frac{1}{2}-\frac{\varepsilon}{2}}$.
By the second part of (A3),
\begin{equation*}
12\,M^3\cdot q^{\frac{1}{2}-\frac{\varepsilon}{2}}=\frac{1}{8}\left(96\,M^3\cdot q^{\frac{1}{4}}\right)q^{\frac{1}{4}-\frac{\varepsilon}{2}}\le\frac{1}{8}q^{\frac{1}{4}-\frac{\varepsilon}{2}}\le q^{\frac{1}{4}-5\varepsilon},
\end{equation*}
because $q\le1$ and $\frac{1}{4}-\frac{\varepsilon}{2}\ge\frac{1}{4}-5\varepsilon$.

\textbf{(e-3) Regularity.}
For a bond $b$ of $T_0$,
\begin{equation*}
\mathscr{Z}(s')(b)=\sum_{b'}\sum_{\Upsilon,\Upsilon'}H_0(s')(b,b')\cdot C(b',\Upsilon)\cdot C_{0}^{1/2}(s')(\Upsilon,\Upsilon')\cdot\mathsf{Z}(\Upsilon')\qquad\text{(finite sums)}.
\end{equation*}
The factor
\begin{equation*}
H_0(s')(b,b')=\mathbf{1}[b=b']\prod\left\{s'_{\mathrm{pc}}:\ \mathrm{pc}\in\pi'_0,\ b\text{ contained in }\mathrm{pc}\right\}
\end{equation*}
(Section~I.5.8) is a polynomial in $s'$.
The factor $C(b',\Upsilon)$ depends neither on $s'$ nor on $\mathsf{Z}$.
The entries $C_{0}^{1/2}(s')(\Upsilon,\Upsilon')$ are continuous on $\mathrm{Pd}_0$ and analytic in its interior:
this is the property W2 of the family $C_{0}^{1/2}(s')$ (the family (p2) of Section~I.5.8) of the given system, a part of the setting (ST1)$^{\mathrm{ch}}$ in the form (G3.T-W) of Section~G3.0.4;
it is property~2 of \cite[Section~5.4.2]{QEDI}, which property~17 there requires of this family at level $0$.
The dependence on $\mathsf{Z}$ is linear.
Therefore $(s',\mathsf{Z})\mapsto\mathscr{Z}(s')(b)$ is continuous on $\mathrm{Pd}_0\times\mathrm{Box}_0$ and, for every $\mathsf{Z}$, analytic in $s'$ in the interior of $\mathrm{Pd}_0$.
By (c), the coefficients of $\mathrm{Vx}(\mathrm{Zf},\square)$ are entire functions of the finitely many numbers $\mathrm{Zf}(b)$, $b$ meeting $\square$.
Hence the composition $G^{(1)}(s,\square)=\mathrm{Vx}\left(\mathscr{Z}(s'),\square\right)$, a map into the span of the finitely many monomials of (d-1), is continuous in $(s',\mathsf{Z})$ on $\mathrm{Pd}_0\times\mathrm{Box}_0$ and analytic in $s'$ in the interior.
It does not depend on $s''$.
\end{proof}

\subsection*{3. Uniformity}

\paragraph{Proposition 23.1\,(a)--(f)}
These clauses are identities and assertions about supports.
They hold for every integer $k\ge1$, every $\ell\in\mathrm{Cl}$, every complex $\mu$ (real $\mu$ in the two clauses on conjugation), every $e_k>0$, every complex $B$ and every union $Y$ of $M$-cubes (every polymer $\mathbf{Y}$).
Their proofs rest on two kinds of inequalities:
the inequalities between sup-distances in 0.1, (0-1), (0-2), (0-3), (b-4) and (f-1), which involve $L^j$, $M$ and $\eta_k$ alone,
and the inequality $b_k>0$ of the conventions of 1.2, of which only the consequence that $b_k$ is defined, real and different from $0$ matters; it rests on $b>0$ and $L>1$.
Of the constants of Section~I.8 the proofs involve $L$, $M$, $b$, $\alpha$, $\rho$ and $w$ alone;
they require of $L$ that it be odd and $L>1$, and of $M$ that it be odd (a power of $L$) and $M\ge1$, so that $\eta_k\le M$ in (0-3).
Of the weights $w$ they require the closure relation (I.1.1), and this at one place only, in the second part of (e), through Proposition~21.1\,(c):
there the relation $D_B^{T}=\mathsf{C}^{-1}\cdot D_{-B}\cdot\mathsf{C}$, which rests on (I.1.1) through (21.23), gives $F_\infty(B)^{T}=\mathsf{C}^{-1}\cdot F_\infty(-B)\cdot\mathsf{C}$;
hence $F_\infty(-B)$ is invertible together with $F_\infty(B)$, so that $S_k(-B)$ exists and $\bar{H}_k(B)=\mathsf{C}^{-1}H_k(-B)\,\mathsf{C}$ is defined.
Of $\rho$ and $\alpha$ nothing is required beyond their being real numbers.
The side $\ell$ enters only through periodicity and antiperiodicity;
examples are the antiperiodicity of the values in Step (0-4), Step (0-6) together with (i23.H3), the class of $\boldsymbol{\lambda}$ in (c-1), (i23.L1) and the representative $a$ in (c-2), and the centres modulo $\ell\mathbb{Z}^3$ and the representative $a$ in (f-1).
Each row of $T_k(B)$, $\bar{T}_k(B)$, $T^{(i)}(B,Y)$, $\bar{T}^{(i)}(B,Y)$ is a sum over the $48\,\eta_k^{-3}$ fermion sites over a single unit block, not over the whole torus.
The operators $H_{k}(B)$, $\bar{H}_{k}(B)$, whose column indices are the unit sites of a fundamental domain (0.2), enter only through the two identities $T_k(B)\cdot H_{k}(B)=1$ and $\bar{T}_k(B)\cdot\bar{H}_{k}(B)=1$, in the second part of (e).
The level $k$ enters through $\eta_k$, through $b_k$ and through the length $k$ of the towers.

\paragraph{Proposition 23.1\,(g)}
The constant of (g) is
\begin{equation*}
C_{T}=2\,\mathrm{e}^{\frac{3}{2}}\left[b_{-}^{-1}\left(54L\cdot\mathrm{e}^6+12\,\mathrm{e}+1\right)+1\right],
\end{equation*}
a function of $(b_{-},L)$ (Section~I.8.2, line (C4)).
Proposition~9\,(d) provides the bound on the row sums with this number for all $k\ge1$, $\ell\in\mathrm{Cl}$, $\mu\in(0,1]$, $e_k\in(0,e_0]$, $B\in\tilde{R}_k+\mathfrak{X}_k$ and all polymers $\mathbf{Y}$.
Hence the inequality of (g), whose only constant is $C_{T}$, holds for all these data, for every $F$ and for all weights with $C_{T}\cdot(\mathsf{w}_1+\mathsf{w}_2)\le\mathsf{w}_3$.
The inequality $1+\eta_k<\ell$ on the side $\ell$ is needed for (g): it makes the column indices of a row pairwise incongruent modulo $\ell\mathbb{Z}^3$.

\paragraph{Proposition 23.2\,(a)--(c)}
These are identities and assertions about supports, valid for every $\ell\in\mathrm{Cl}$, every $q>0$, all complex $\mathrm{Zf}$ and $\mathscr{T}$, and every $M$-cube $\square$.

\paragraph{Proposition 23.2\,(d)}
Besides the weights, the bounds depend only on $\rho$, $M$, $q$, $\varepsilon$ and on the supremum over the bonds that meet $\square$;
in particular they depend neither on $\ell$, nor on $\square$, nor on the torus, nor on $\mathrm{Zf}$ beyond this supremum.
The inequality $M-1+\eta_0=M<\ell$ on the side $\ell$ is needed for (d): in (d-1) it makes the covering map injective on a lift of $\square$.
In the number $12(1+\rho)\cdot M^3=\frac{1}{96}\cdot48\,M^3\cdot2\cdot12(1+\rho)$ the weight $1/48$ of a site compensates the $48$ sites of a cluster.

\paragraph{Proposition 23.2\,(e)}
The two bounds are valid for all data $(\ell,\mu,q)$ of (ST1)$^{\mathrm{ch}}$ at $k=0$, for every system with the properties required there, for every $s\in\mathrm{Pd}_0$ and $\mathsf{Z}\in\mathrm{Box}_0$, and for every $M$-cube $\square$.
They depend on the cutoff $N$ and on the charge $e$ only through the datum $q$ (along the flow $q=e\cdot\eta^{\frac{1}{2}}$), and they do not depend on the mass $\mu$.

\subsection*{4. Conditions on the constants}

\begin{description}
\item[Place in the proofs.] 0.1; 1.2, (0-1), (0-2); 2.2, Step~0, (a), (d-1) (the numbers $L^j$ and $M$ are odd; the blocks and the $M$-cubes partition the sites; a unit block lies in one $M$-cube; an $M$-cube of level $0$ has $M^3$ sites); 1.2, (0-3), and 2.2, Step~0: $M\ge1$ ((M1)).
\item[Condition.] $L$ odd; $M$ a power of $L$ (hence odd, and $M\ge1$).
\item[Where it stands.] Section~I.8.2, lines (C3), (C5); Proposition~14.2\,(b), (L5), (M1); Annex~B, Section~G3.6.1, row of Propositions~23.1\,(a)--(f) and 23.2\,(a)--(d).
\end{description}

\begin{description}
\item[Place in the proofs.] 1.2, all clauses ($b_k$ is defined and positive).
\item[Condition.] $b=1>0$, $L>1$.
\item[Where it stands.] Section~I.8.2, lines (C0), (C3), (C4); Section~I.3.4.
\end{description}

\begin{description}
\item[Place in the proofs.] 1.2, (e), second part, through Proposition~21.1\,(a), (c) (the existence of $S_k(-B)$ and of $\bar{H}_k(B)$).
\item[Condition.] The closure relation (I.1.1) of the standard weights, which enters through Proposition~21.1\,(c).
\item[Where it stands.] Section~I.8.2, line (C0); (I.1.1); Chapter~21a, Section~5; Annex~B, Section~G3.6.1, first entry.
\end{description}

\begin{description}
\item[Place in the proofs.] 1.2, (g), through Proposition~9\,(d).
\item[Condition.] Those of Proposition~9\,(d), as listed in Chapter~9, Sections~8.2 and 9:
(Z2.s4), which gives $\mathfrak{i}(B)\le1$ and $e_k\cdot\sup|dB|\le1$ for $B\in\tilde{R}_k+\mathfrak{X}_k$;
(L5);
(M1), which together with (L5) makes $M$ odd with $M>4L$;
$0\le\rho\le1$, so that $3(1+\rho)\le6$ (contained in $0<\rho<1$);
the standard weights, through the closure relation (I.1.1) and the value $w_3=6$;
$\alpha\in(0,1)$, by line (C1), used in two forms:
as $\alpha\ge0$, which gives $\eta_k^{\alpha}\le1$,
and as $0<\alpha<1$, both for the domain $\tilde{R}_k+\mathfrak{X}_k$, through its summand $\tilde{R}_k$ of \cite[Definition~5.16]{QEDI}, which depends on $\alpha$ through the norm $|\delta_{\alpha}\partial A|_{\infty}$ (the other summand, the set $\mathfrak{X}_k$ of the same definition, does not involve $\alpha$),
and for the display (G3) of \cite[Section~16.2.7]{QEDI}, which is stated under the data of \cite[Section~16.2.1]{QEDI}, $0<\alpha<1$ among them;
$b=1>0$ and $L>1$, so that $b_k$ is defined and positive and $b_{-}>0$, as for all clauses of 1.2;
$b_{-}\le b_k$;
$\varepsilon$ as fixed in line (C7), that is, under (S1) of Proposition~14.2\,(b): $1-\alpha<\varepsilon\le\bar{\varepsilon}$.
\item[Where it stands.] Proposition~14.2\,(b) and the sentence on the conditions after Proposition~9, for (Z2.s4), (L5), (M1), (S1); Section~I.8.2, lines (C0) ($\rho$, the weights, $b=1$), (C1) ($\alpha\in(1-\bar{\varepsilon},1)$), (C3) ($L>1$), (C4) ($b_{-}$), (C7) ($\varepsilon$), and Proposition~14.2\,(a); Section~I.3.4 ($b_k$, and $b_{-}\le b_k$); Chapter~9, Sections~8.2 and 9; Annex~B, Section~G3.6.1, row of Proposition~23.1\,(g).
\end{description}

\begin{description}
\item[Place in the proofs.] 2.2, (c) (the diagonal entries of $M_3$ do not vanish); (d-2), (d-3) ($12(1+\rho)$ as a bound for the sums of the moduli of the entries; $1+\rho\le2$).
\item[Condition.] $0<\rho<1$.
\item[Where it stands.] Section~I.8.2, line (C0); Annex~B, Section~G3.6.1.
\end{description}

\begin{description}
\item[Place in the proofs.] 2.2, (d-3), (e-2) ($q^{\varepsilon}\le1$, $q^{1-\frac{\varepsilon}{2}}\le1$, $\frac{1}{4}-\frac{\varepsilon}{2}\ge\frac{1}{4}-5\varepsilon$ for $q\le1$).
\item[Condition.] $0<\varepsilon\le2$; it holds because $\varepsilon\in(1-\alpha,\bar{\varepsilon}]$, $\bar{\varepsilon}=1/200$.
\item[Where it stands.] Section~I.8.2, line (C7); Proposition~14.2\,(b), (S1).
\end{description}

\begin{description}
\item[Place in the proofs.] 2.2, (e-1), (e-2).
\item[Condition.] (A3), second part: $96\,M^3\cdot\exp\left(-\frac{v}{4}\right)\le1$ for all $v\ge-\log e_0$.
\item[Where it stands.] Proposition~14.2\,(b), (A3); Section~I.8.4.
\end{description}

\begin{description}
\item[Place in the proofs.] 2.2, (e): the setting and the steps (e-1)--(e-3); through them, 2.2, (c) and (d), and the proof of Proposition~24.3\,(q1) in Chapter~24, Section~2.2.
\item[Condition.] In the proof of (e) itself:
the second part of (A3) at $v=-\log q$, that is, $96\,M^3\cdot q^{\frac{1}{4}}\le1$, in (e-1) and (e-2);
$\varepsilon\ge0$, contained in the range $0<\varepsilon\le2$, which holds by line (C7), in (e-2), where $\frac{1}{4}-\frac{\varepsilon}{2}\ge\frac{1}{4}-5\varepsilon$ is used;
$L\ge1$, by which the inequality $L^{\frac{1}{2}}\cdot q\le e_0$ of (ST1)$^{\mathrm{ch}}$ at $k=0$ gives $q\le e_0$, so that the conditions on $e_0$ hold at $v=-\log q$.
Through Proposition~23.2\,(c) and (d), which (e-2) and (e-3) apply, the conditions of these two clauses:
$L$ and $M$ odd;
$M\ge1$;
$0<\rho<1$;
$0<\varepsilon\le2$.
Through Proposition~24.3\,(q1), which (e-1) applies at $k=0$, where $e_k=q$.
The assertion of (q1) applied there, the bound $\frac{1}{2}e_k^{-\frac{\varepsilon}{2}}$ on the three norms of $\mathscr{Z}(s')$, rests on the following conditions:
the first inequality of (A1), that is, $p_{0,k}\le e_k^{-\frac{\varepsilon}{4}}$;
(A2);
$C_w\ge4\exp\left(2\gamma_{\mathrm{w}}\right)\cdot M^{27\alpha_0}$, which holds by the definition of $C_w^{(1)}$ and which the bound needs at level $0$, the only level of (e): by W17 the diagonal family $H_0(s')$ has the property W3 with the constant $4\exp\left(2\gamma_{\mathrm{w}}\right)\cdot M^{27\alpha_0}$, hence with the constant $C_w$;
$L$ odd, $L\ge3$ and $M$ a power of $L$, the setting of \cite[Proposition~5.35]{QEDI}, that is, of \cite[Proposition~15.21]{QEDI}, from which the fact (P2) of Chapter~24 is taken; (P2) enters the bound through the inequality $\sum_{\Upsilon}|C(b',\Upsilon)|\le c_C$.
The two other assertions of (q1), the inequalities between the radii $\mathrm{rad}_1$, $\mathrm{rad}_2$, $\mathrm{rad}_4$ and the assertion on the domains, are not needed in the proof of (e); they rest moreover on the following conditions:
the third inequality of (A1), that is, $e_k^{\frac{5\varepsilon}{2}}\le\frac{1}{2}$, the condition (A5) together with $M\ge1$, and $\varepsilon\le\frac{1}{200}$ in the form $\frac{1}{2}-4\varepsilon>0$, in the inequalities between the radii;
$\varepsilon\ge0$ and $e_k\le1$, the latter by the first inequality of (A8), in the inequalities between the radii and in the inequalities $\mathrm{t3}_k\le\mathrm{t2}_k\le\mathrm{t1}_k$ between the thresholds;
the first and the fourth inequality of (A8), together with $L\ge5$, on which the proof of \cite[Lemma~5.17(a)]{QEDI}, that is, of \cite[Lemma~17.4(a)]{QEDI}, rests; by this lemma the set $\tilde{R}_k+\mathfrak{X}_k$ is contained in $\tilde{G}_k$.
\item[Where it stands.] Proposition~14.2\,(b) and Section~I.8.4, for (A1), (A2), (A3), (A5), (A8);
Section~I.8.2, lines (C0), (C7), and Proposition~14.2\,(b), (S1), for $\varepsilon$;
Section~I.8.2, line (C0), for $\rho$;
Section~I.8.2, lines (C3), (C5), Proposition~14.2\,(b), (L5), (M1), and Proposition~14.2\,(f), for $L$ and $M$;
for the inequality on $C_w$: Section~I.8.2, line (C6); Proposition~14.2\,(a), rows [85], [86]; Proposition~14.2\,(f); (K6) of Chapter~24, Section~0.3;
Chapter~24, Section~11; Annex~B, Section~G3.6.1, row of Proposition~23.2\,(e).
\end{description}

\medskip
The side $\ell\in\mathrm{Cl}$ and the inequality $L^{\frac{1}{2}}q\le e_0$ of (ST1)$^{\mathrm{ch}}$ belong to the data.
Of $\ell$ the proofs require that it be an integer multiple of $M$, that $M-1+\eta_0=M<\ell$ and that $1+\eta_k<\ell$.
Each of the conditions on the constants above is one of those of Section~I.8, that is, of Proposition~14.2 (values and ranges in (a), conditions in (b), consequences in (f)).

\section*{Chapter 25. The formal locality lemma: proofs of Propositions 25.1--25.3}

\subsection*{0. Conventions of this chapter}

\emph{Setting.}
Throughout this chapter the data and the covering objects are those of Section~G3.5.0.
That section also recalls three clauses of Section~I.1.9, under the names (G3.r1)--(G3.r3).

\emph{References for the definitions.}
The notion of an operator given by one formula, the index sets and the properties (I.f1)--(I.f4) are those of Section~I.1.9.
Definition~24.D, with its clauses (D0)--(D8) and the displays (8.1)$^{\mathrm{ch}}$--(8.6)$^{\mathrm{ch}}$, is stated in Section~G3.4.1.
The displays (23.1) and (8.3a)$^{\mathrm{ch}}$ are the definitions that precede Proposition~23.1, and (23.3) is the definition that precedes Proposition~23.2.
By (23.2) is meant the display in Proposition~23.1\,(a); it is established below as (25.2).
The property W12 is the one stated in Section~G1.11.
The clauses (t7) and (t10) are those of \cite[Proposition~5.58]{QEDI} (see also \cite[Theorem~13.2]{QEDI}).
The facts on the basis $\{\Upsilon\}$ and on the map $C$ are those of \cite[Proposition~15.21]{QEDI}, summarised in \cite[Proposition~5.35]{QEDI}.

\emph{Supports.}
For a basis element $\Upsilon$, the set $\mathrm{supp}\,\Upsilon$ consists of the end points of the bonds that belong to the support of $\Upsilon$ (Section~G4.0.4, paragraph on indices; \cite[Section~5.2.3]{QEDI}).

\subsection*{1. Preliminaries}

\subsubsection*{1.1. Notation}

The level $k$, the mass $\mu$, the charge $e_k$ and all constants are those of Section~G3.5.0.
One has $M=L^{\mathrm{jm}}$ for an integer $\mathrm{jm}\ge1$, and both $L$ and $M$ are odd (Section~I.0.5).
The side $\ell$ belongs to $\mathrm{Cl}$; in particular, $\ell$ is a multiple of $L\cdot M$ and $\ell\ge13L\cdot M$.
Consequently
\[
\ell\mathbb{Z}^3\subset(L\cdot M)\mathbb{Z}^3\subset M\mathbb{Z}^3\subset L\mathbb{Z}^3\subset\mathbb{Z}^3\subset\eta_k\mathbb{Z}^3.
\]

Let $\tilde{U}$, $\tilde{U}'$ be finite non-empty sets of sites of $\mathbb{L}_k$ (in Sections~1.5, 3.9 and~4.3 also of $\mathbb{L}_{k+1}$), and let $p$ be a site.
Put
\begin{align*}
\mathrm{sdm}(\tilde{U})&:=\max\left\{|c-c'|:\ c,c'\in\tilde{U}\right\},\\
\mathrm{dist}(p,\tilde{U})&:=\min\left\{|p-c|:\ c\in\tilde{U}\right\},\\
\mathrm{dist}(\tilde{U},\tilde{U}')&:=\min\left\{|c-c'|:\ c\in\tilde{U},\ c'\in\tilde{U}'\right\};
\end{align*}
thus $\mathrm{sdm}(\tilde{U})$ is the sup-diameter of $\tilde{U}$.
For the neighbourhoods $\tilde{U}^{+a}$ of Section~G3.5.0 one has $\mathrm{sdm}\left(\tilde{U}^{+a}\right)\le\mathrm{sdm}(\tilde{U})+2a$.

The covering map $\mathbb{L}_k\to T_k$ is denoted by $\mathrm{pr}_\ell$, and so are the maps that it induces on the index sets of Section~I.1.9.
For $w\in\mathbb{Z}^3$ let $\boldsymbol{\chi}(w):=(-1)^{w_1+w_2+w_3}$ (Section~I.6.3).
In this chapter the letter $w$ stands for a vector of $\mathbb{Z}^3$ with the components $w_1$, $w_2$, $w_3$, as in Sections~I.1.9 and~I.6.3; in the paragraph on the map $\mathrm{cov}_\ell$ of Section~G3.5.0 this vector is called $n$.
This vector is not to be confused with the weights of Section~I.0.5, which are denoted $w_i$ in the matrices $M_i(x)$ of (Of1) and form the quadruple $w$ among the constants listed in Section~2.3.

A lattice box is a product of three intervals of $\eta_k\mathbb{Z}$.
For sites $c$, $c'$ of $\mathbb{L}_k$ let $\mathrm{hul}(c,c')$ be the set of the sites $p\in\mathbb{L}_k$ such that, for $\mu=1,2,3$, the number $p_\mu$ lies in the closed interval with the end points $c_\mu$ and $c'_\mu$.
It is a lattice box.

\subsubsection*{1.2. Blocks, paths, towers, translations}

\begin{named}{Fact 1}[Blocks]
For an integer $j\ge0$ and a point $z\in L^j\eta_k\mathbb{Z}^3$ put
\[
\mathrm{Bl}_j(z):=\left\{c\in\mathbb{L}_k:\ |c-z|<\frac{L^j\eta_k}{2}\right\};
\]
for $j\le k+1$ these are the blocks of Section~I.1.1.
Since $M=L^{k+\mathrm{jm}}\eta_k$, the $M$-cube with the centre $z\in M\mathbb{Z}^3$ is $\mathrm{Bl}_{k+\mathrm{jm}}(z)$, and the $LM$-cube with the centre $z\in(L\cdot M)\mathbb{Z}^3$ is $\mathrm{Bl}_{k+\mathrm{jm}+1}(z)$.
\begin{itemize}
\item[(i)] One has
\[
\mathrm{Bl}_j(z)=\left\{c\in\mathbb{L}_k:\ |c-z|\le\frac{(L^j-1)\cdot\eta_k}{2}\right\},
\]
and this is a lattice box with $L^j$ sites in each direction.
\item[(ii)] For fixed $j$ the sets $\mathrm{Bl}_j(z)$, $z\in L^j\eta_k\mathbb{Z}^3$, form a partition of $\mathbb{L}_k$.
\item[(iii)] If $j\le j'$, then every set $\mathrm{Bl}_j(z)$ is contained in exactly one set $\mathrm{Bl}_{j'}(z')$.
\end{itemize}
\end{named}

\begin{proof}
Write $c=\eta_k\cdot n$ and $z=\eta_k\cdot L^j\cdot a$ with $n,a\in\mathbb{Z}^3$.
The condition $|c-z|<\frac{L^j\eta_k}{2}$ means $\left|n_\mu-L^j\cdot a_\mu\right|<\frac{L^j}{2}$ for $\mu=1,2,3$.
As $n_\mu-L^j\cdot a_\mu$ is an integer and $L^j$ is odd, this is equivalent to $\left|n_\mu-L^j\cdot a_\mu\right|\le\frac{L^j-1}{2}$, which is (i).
Every integer $n_\mu$ has exactly one representation $n_\mu=L^j\cdot a_\mu+r_\mu$ with integers $a_\mu$, $r_\mu$ and $|r_\mu|\le\frac{L^j-1}{2}$, because there are $L^j$ such numbers $r_\mu$, one in each residue class modulo $L^j$; this is (ii).
For (iii) it suffices to treat the case $j'=j+1$.
Write $a_\mu=L\cdot a'_\mu+r'_\mu$ with $|r'_\mu|\le\frac{L-1}{2}$.
Every $c\in\mathrm{Bl}_j(z)$ satisfies
\[
\left|n_\mu-L^{j+1}\cdot a'_\mu\right|\le\left|n_\mu-L^j\cdot a_\mu\right|+L^j\cdot|r'_\mu|\le\frac{L^j-1}{2}+L^j\cdot\frac{L-1}{2}=\frac{L^{j+1}-1}{2}.
\]
Therefore $\mathrm{Bl}_j(z)$ is contained in $\mathrm{Bl}_{j+1}\left(\eta_k\cdot L^{j+1}\cdot a'\right)$, and by (ii) in no other block of this side.
\end{proof}

\emph{Consequences and notation.}
For $c\in\mathbb{L}_k$ let $[c]$ be the unit site with $c\in\Delta_{[c]}$ (Section~I.1.1), and let $[c]^L$ be the point $y'$ of $L\mathbb{Z}^3$ with $c\in\mathrm{Bl}_{k+1}(y')$.
\begin{itemize}
\item[(1.a)] The block
\[
\Delta_y=\mathrm{Bl}_k(y)=\left\{c:\ |c-y|\le\frac{1-\eta_k}{2}\right\}
\]
has $\eta_k^{-3}$ sites, and $\Delta_y\times O_{\mathrm{cub}}$ has $48\,\eta_k^{-3}$ sites; accordingly
\[
\int_{\Delta_y}\mathrm{d}x=\frac{\eta_k^3}{48}\cdot48\,\eta_k^{-3}=1.
\]
For $k=0$ one has $\Delta_y=\{y\}$ and $[c]=c$.
\item[(1.b)] One has $\mathrm{Bl}_{k+1}(y')=\left\{c:\ |c-y'|\le\frac{L-\eta_k}{2}\right\}$, and the set
\[
B(y')=\mathrm{Bl}_{k+1}(y')\cap\mathbb{Z}^3=\left\{y\in\mathbb{Z}^3:\ |y-y'|\le\frac{L-1}{2}\right\}
\]
has $L^3$ sites.
The set $\mathrm{Bl}_{k+1}(y')$ is the disjoint union of the blocks $\Delta_y$, $y\in B(y')$: a unit site $y$ belongs to $\Delta_y$, and by (iii) and (ii) the block $\Delta_y$ is contained in $\mathrm{Bl}_{k+1}(y')$ if and only if $y$ is.
\item[(1.c)] Each block $\Delta_y$ and each block $\mathrm{Bl}_{k+1}(y')$ lies in one $M$-cube, and each $M$-cube lies in one $LM$-cube.
This is (iii) with $k\le k+1\le k+\mathrm{jm}\le k+\mathrm{jm}+1$, and it is the statement (g0) of \cite[Section~17.4.1]{QEDI}.
Hence a union of $M$-cubes is the disjoint union of the blocks $\Delta_y$ over its unit sites $y$, and also the disjoint union of the blocks $\mathrm{Bl}_{k+1}(y')$ over its points $y'$ of $L\mathbb{Z}^3$.
\end{itemize}

\begin{named}{Fact 2}[Paths]
Along the path $\boldsymbol{\Gamma}^{\pi}(c,c')$ of Section~I.1.1 the coordinates change one after the other, each of them monotonically from its value at $c$ to its value at $c'$.
Hence every site of $\boldsymbol{\Gamma}^{\pi}(c,c')$ belongs to $\mathrm{hul}(c,c')$.
Each $p\in\mathrm{hul}(c,c')$ satisfies $|p-c|\le|c-c'|$ and $|p-c'|\le|c-c'|$, and every lattice box that contains $c$ and $c'$ contains $\mathrm{hul}(c,c')$.
Let $\mathrm{Pth}(c,c')$ denote the set of the bonds of the six paths $\boldsymbol{\Gamma}^{\pi}(c,c')$.
It is empty if $c=c'$, and it consists of the bond $(c,c')$ if $|c-c'|_2=\eta_k$.
The number
\[
(\tau\tilde{B})(c,c')=\frac{1}{6}\sum_{\pi}\tilde{B}\left(\boldsymbol{\Gamma}^{\pi}(c,c')\right)
\]
is a linear form in the numbers $\tilde{B}(b)$, $b\in\mathrm{Pth}(c,c')$.
\end{named}

\begin{named}{Fact 3}[Towers; $k\ge1$]
For $c\in\mathbb{L}_k$ and $0\le j\le k$ let $y_j(c)$ be the point of $L^j\eta_k\mathbb{Z}^3$ with $c\in\mathrm{Bl}_j\left(y_j(c)\right)$ (Fact~1\,(ii)).
Then $y_0(c)=c$ and $y_k(c)=[c]$; the points $y_j(c)$ form the tower of Section~I.3.2.
Let $\mathrm{Twr}(c)$ be the union of the sets $\mathrm{Pth}\left(y_{j+1}(c),y_j(c)\right)$, $0\le j<k$.
Every end point $p$ of a bond of $\mathrm{Twr}(c)$ satisfies
\[
p\in\Delta_{[c]},\qquad|p-[c]|\le\frac{1-\eta_k}{2},\qquad|p-c|\le\frac{1-\eta_k}{2}.
\]
\end{named}

\begin{proof}
The block $\mathrm{Bl}_j\left(y_j(c)\right)$ contains $c$ and its centre $y_j(c)$, and by Fact~1\,(iii) and~(ii) it is contained in $\mathrm{Bl}_k\left(y_k(c)\right)=\Delta_{[c]}$.
Thus $y_j(c)$ and $y_{j+1}(c)$ lie in the lattice box $\Delta_{[c]}$.
The point $p$ belongs to $\mathrm{hul}\left(y_{j+1}(c),y_j(c)\right)$ for some $j<k$, and therefore lies in $\Delta_{[c]}$ by Fact~2.
By (1.a), $|p-[c]|\le\frac{1-\eta_k}{2}$.
For every $i\le k$, Fact~1\,(i) gives
\[
|y_i(c)-c|\le\frac{(L^i-1)\cdot\eta_k}{2}\le\frac{(L^k-1)\cdot\eta_k}{2}=\frac{1-\eta_k}{2}.
\]
Since $p_\mu$ lies between $y_{j+1}(c)_\mu$ and $y_j(c)_\mu$, it follows that
\[
|p_\mu-c_\mu|\le\max\left\{\left|y_{j+1}(c)_\mu-c_\mu\right|,\left|y_j(c)_\mu-c_\mu\right|\right\}\le\frac{1-\eta_k}{2}.
\]
\end{proof}

For $c\in\Delta_y$ the definitions (I.3.3) take the form
\[
(\tau_k\tilde{B})(y,c)=\sum_{j<k}(\tau\tilde{B})\left(y_{j+1}(c),y_j(c)\right),\qquad
\varphi_{\tilde{B}}(y,c)=\exp\left(i\,e_k\cdot\eta_k\cdot(\tau_k\tilde{B})(y,c)\right).
\]
Thus $\varphi_{\tilde{B}}(y,c)$ is the exponential of a linear form in the numbers $\tilde{B}(b)$, $b\in\mathrm{Twr}(c)$.

\begin{named}{Fact 4}[Paths in an $L$-block]
Let $y'\in L\mathbb{Z}^3$ and $y\in B(y')$.
Then $|y-y'|\le\frac{L-1}{2}\le M$, so that
\[
U_{\tilde{B}}(y',y)=\exp\left(i\,e_k\cdot\eta_k\cdot(\tau\tilde{B})(y',y)\right)
\]
is defined by (I.1.6); it is the exponential of a linear form in the numbers $\tilde{B}(b)$, $b\in\mathrm{Pth}(y',y)$, and these $b$ are fine bonds.
Every end point $p$ of a bond of $\mathrm{Pth}(y',y)$ lies in $\mathrm{hul}(y',y)$, which is contained in the lattice box $\mathrm{Bl}_{k+1}(y')$, and satisfies $|p-y'|\le\frac{L-1}{2}$ and $|p-y|\le\frac{L-1}{2}$ (Fact~2).
\end{named}

\begin{named}{Fact 5}[Translations]
Translations act by $(T_vf)(t)=f(t-v)$; in particular $(T_v\tilde{B})(b+v)=\tilde{B}(b)$.
Let $v\in\eta_k\mathbb{Z}^3$.
\begin{itemize}
\item[(5.a)] One has
\[
(x+v)^{(i)}=x^{(i)}+v,\qquad c_3(x+v)=c_3(x)+v,\qquad M_i(x+v)=M_i(x),\qquad d_X(x+v,x'+v)=d_X(x,x'),
\]
and the map $(x,x')\mapsto(x+v,x'+v)$ maps $\mathrm{Pzt}_k$ onto itself.
This follows from (I.1.2), (I.2.1) and (I.1.5), since a translation does not change the label.
\item[(5.b)] One has $\boldsymbol{\Gamma}^{\pi}(c+v,c'+v)=\boldsymbol{\Gamma}^{\pi}(c,c')+v$.
Hence $\mathrm{Pth}(c+v,c'+v)=\mathrm{Pth}(c,c')+v$, and
\[
(\tau T_v\tilde{B})(c+v,c'+v)=(\tau\tilde{B})(c,c'),\qquad u_{T_v\tilde{B}}(c+v,c'+v)=u_{\tilde{B}}(c,c'),\qquad U_{T_v\tilde{B}}(x+v,x'+v)=U_{\tilde{B}}(x,x').
\]
\item[(5.c)] If $v\in L^j\eta_k\mathbb{Z}^3$, then $\mathrm{Bl}_j(z)+v=\mathrm{Bl}_j(z+v)$.
Let $v\in\mathbb{Z}^3$; this lattice is contained in $L^j\eta_k\mathbb{Z}^3=L^{j-k}\mathbb{Z}^3$ for every $j\le k$.
Then $y_j(c+v)=y_j(c)+v$ for $j\le k$, $[c+v]=[c]+v$, $\Delta_y+v=\Delta_{y+v}$, $\mathrm{Twr}(c+v)=\mathrm{Twr}(c)+v$ and $\varphi_{T_v\tilde{B}}(y+v,c+v)=\varphi_{\tilde{B}}(y,c)$.
If in addition $v\in L\mathbb{Z}^3$, then $[c+v]^L=[c]^L+v$, $B(y')+v=B(y'+v)$ and $\mathrm{Bl}_{k+1}(y')+v=\mathrm{Bl}_{k+1}(y'+v)$.
A translation by $v\in M\mathbb{Z}^3$ maps $M$-cubes to $M$-cubes, centres to centres and common faces to common faces, and $\tilde{Y}+v$ is then a union of $M$-cubes.
A translation by $v\in(L\cdot M)\mathbb{Z}^3$ maps $LM$-cubes to $LM$-cubes.
\item[(5.d)] For every finite set $Z$ of photon sites and every function $\mathrm{hf}$ on $\mathbb{L}^{\mathrm{f}}_k$,
\[
\int_{Z+v}\mathrm{d}x\,\mathrm{hf}(x)=\int_Z\mathrm{d}x\,\mathrm{hf}(x+v).
\]
\end{itemize}
\end{named}

\subsubsection*{1.3. Lifts}

Within Lemma~1 the letters $\tilde{X}$ and $\tilde{\mathbf{Y}}$ stand for unions of $M$-cubes of $\mathbb{L}_k$.
They are required to be finite, or to be connected through faces, only where this is said explicitly.

\begin{named}{Lemma 1}
\begin{itemize}
\item[(a)] Let $\tilde{U}$ be a finite non-empty subset of $\mathbb{R}^3$ (in the applications, a set of sites of $\mathbb{L}_k$ or of $\mathbb{L}_{k+1}$), and let $\ell'>0$ and $a\ge0$ be real numbers with $\mathrm{sdm}(\tilde{U})+a<\ell'$.
Then
\[
|c-c'-\ell'\cdot w|>a\qquad\text{for all }c,c'\in\tilde{U}\text{ and all }w\in\mathbb{Z}^3\setminus\{0\}.
\]
Consequently the map $c\mapsto c+\ell'\mathbb{Z}^3$ is injective on $\tilde{U}$; one has $\mathrm{dist}(\tilde{U},\tilde{U}+\ell'\cdot w)>a$ for $w\ne0$; and a point $p$ with $\mathrm{dist}(p,\tilde{U})\le a$ belongs to $\tilde{U}+\ell'\mathbb{Z}^3$ if and only if it belongs to $\tilde{U}$.
\item[(b)] \emph{Polymers.}
Let $\tilde{\mathbf{Y}}$ be a finite non-empty union of $M$-cubes of $\mathbb{L}_k$ with $\mathrm{sdm}(\tilde{\mathbf{Y}})+\eta_k<\ell$, and let $\mathbf{Y}:=\mathrm{pr}_\ell(\tilde{\mathbf{Y}})$.
The map $\mathrm{pr}_\ell$ maps the $M$-cubes of $\tilde{\mathbf{Y}}$ bijectively onto the $M$-cubes of $\mathbf{Y}$, and two $M$-cubes of $\tilde{\mathbf{Y}}$ have a common face if and only if their images have one.
Hence a union $\tilde{X}$ of $M$-cubes contained in $\tilde{\mathbf{Y}}$ is connected through faces if and only if $\mathrm{pr}_\ell(\tilde{X})$ is.
The map $\tilde{X}\mapsto\mathrm{pr}_\ell(\tilde{X})$ is a bijection of the set of the non-empty unions of $M$-cubes contained in $\tilde{\mathbf{Y}}$ that are connected through faces onto the set of the polymers contained in $\mathbf{Y}$; every such $\tilde{X}$ is a lift of $\mathrm{pr}_\ell(\tilde{X})$ (Section~I.1.1), and $\left|\mathrm{pr}_\ell(\tilde{X})\right|_M=|\tilde{X}|_M$.
If moreover $\tilde{\mathbf{Y}}$ is connected through faces, then $\mathbf{Y}$ is a polymer and $\tilde{\mathbf{Y}}$ is a lift of $\mathbf{Y}$.
Conversely (Section~I.1.1), every lift of a polymer $X$ is a finite union of $M$-cubes of $\mathbb{L}_k$ that is connected through faces, and two lifts of $X$ differ by a vector of $\ell\mathbb{Z}^3$.
\item[(c)] \emph{Small polymers.}
A polymer $X$ of $T_k$ that has a lift $\tilde{X}$ is small if and only if $d_M(\tilde{X})\le L$, where $d_M(\tilde{X})$ is formed with nets in $\mathbb{R}^3$, as in \cite[Proposition~5.58\,(t10)]{QEDI}.
Every small polymer has a lift, and every lift $\tilde{X}$ of a small polymer satisfies $\mathrm{sdm}(\tilde{X})\le M\cdot(L+2)$, hence $\mathrm{sdm}(\tilde{X})+2+M<\ell$.
Every finite non-empty union $\tilde{X}$ of $M$-cubes of $\mathbb{L}_k$ that is connected through faces and satisfies $d_M(\tilde{X})\le L$ is a lift of the small polymer $\mathrm{pr}_\ell(\tilde{X})$.
For an $M$-cube $\tilde{\square}$ of $\mathbb{L}_k$ let $\mathrm{Sm}(\tilde{\square})$ be the set of those unions $\tilde{X}$ of the kind just described that contain $\tilde{\square}$.
The set $\mathrm{Sm}(\tilde{\square})$ does not depend on $\ell$; one has $\mathrm{Sm}(\tilde{\square}+v)=\mathrm{Sm}(\tilde{\square})+v$ for $v\in M\mathbb{Z}^3$; and $\tilde{X}\mapsto\mathrm{pr}_\ell(\tilde{X})$ is a bijection of $\mathrm{Sm}(\tilde{\square})$ onto the set of the small polymers of $T_k$ that contain $\mathrm{pr}_\ell(\tilde{\square})$.
\item[(d)] \emph{Admissible pairs.}
The two conditions of Section~I.1.7 are read on the covering lattice as follows.
Let $\tilde{X}$ be a union of $M$-cubes of $\mathbb{L}_k$, and let $(x,x')\in\mathrm{Pzt}_k$ with $x=(c,g)$ and $x'=(c',g')$.
The pair satisfies (I) for $\tilde{X}$ if $c$, $c'$ and all sites of the six paths $\boldsymbol{\Gamma}^{\pi}(c,c')$ lie in $\tilde{X}$ (for $c=c'$: if $c$ lies in $\tilde{X}$); it satisfies (II) for $\tilde{X}$ if $x'=x^{(3)}$ and $c$ or $c_3$ lies in $\tilde{X}$; and it is admissible for $\tilde{X}$ if it satisfies (I) or (II).
Both photon sites of a pair that is admissible for $\tilde{X}$ lie in $\tilde{X}^{+\eta_k}$.
If the pair satisfies (I), the bonds of $\mathrm{Pth}(c,c')$ are contained in $\tilde{X}$; if it satisfies (II), then $\mathrm{Pth}(c,c')$ consists of the bond $(c,c_3)$, which meets $\tilde{X}$.
Now let $\tilde{X}$ be finite with $\mathrm{sdm}(\tilde{X})+2\eta_k<\ell$, and put $X:=\mathrm{pr}_\ell(\tilde{X})$.
Then $\mathrm{pr}_\ell$ maps the set of the pairs of $\mathrm{Pzt}_k$ that are admissible for $\tilde{X}$ bijectively onto the set of the pairs of $\mathrm{Pz}_k$ that are admissible for $X$, and a pair satisfies (I), resp.\ (II), for $\tilde{X}$ if and only if its image satisfies (I), resp.\ (II), for $X$.
\end{itemize}
\end{named}

\begin{proof}
(a) One has $|\ell'\cdot w|\ge\ell'$ and $|c-c'|\le\mathrm{sdm}(\tilde{U})$, whence $|c-c'-\ell'\cdot w|\ge\ell'-\mathrm{sdm}(\tilde{U})>a$.
An equality $c=c'+\ell'\cdot w$ with $w\ne0$ would give $0>a$; so the map is injective.
The number $\mathrm{dist}(\tilde{U},\tilde{U}+\ell'\cdot w)$ is the minimum of the numbers $|c-c'-\ell'\cdot w|$, $c,c'\in\tilde{U}$.
For the last assertion let $p=c'+\ell'\cdot w$ with $c'\in\tilde{U}$, and let $|p-c|\le a$ for some $c\in\tilde{U}$; then $|c-c'-\ell'\cdot w|\le a$, and therefore $w=0$.

(b) Since $\ell\mathbb{Z}^3\subset M\mathbb{Z}^3$, the map $\mathrm{pr}_\ell$ maps $M$-cubes onto $M$-cubes, and by (a) it is injective on $\tilde{\mathbf{Y}}$.
Let $z,z'\in M\mathbb{Z}^3$ be the centres of two $M$-cubes of $\tilde{\mathbf{Y}}$.
The two cubes have a common face if and only if $z-z'=\pm M\cdot e_\mu$ for some $\mu$, and their images have one if and only if $z-z'\in\pm M\cdot e_\mu+\ell\mathbb{Z}^3$ for some $\mu$.
Suppose that $z-z'=\pm M\cdot e_\mu+\ell\cdot w$ with $w\ne0$.
Then $|z-z'|\ge\ell-M$.
Choose a direction $\nu$ with $|z_\nu-z'_\nu|=|z-z'|$, and let $\boldsymbol{\sigma}$ be the sign of $z_\nu-z'_\nu$.
The sites
\[
z+\boldsymbol{\sigma}\cdot\frac{M-\eta_k}{2}\cdot e_\nu\qquad\text{and}\qquad z'-\boldsymbol{\sigma}\cdot\frac{M-\eta_k}{2}\cdot e_\nu
\]
lie in the two cubes, by Fact~1\,(i), because $\frac{M-\eta_k}{2}=\frac{L^{k+\mathrm{jm}}-1}{2}\cdot\eta_k$ is a multiple of $\eta_k$.
Their distance is at least $|z-z'|+M-\eta_k\ge\ell-\eta_k$, which contradicts $\mathrm{sdm}(\tilde{\mathbf{Y}})+\eta_k<\ell$.
Hence common faces correspond in both directions.
Consider the graph whose vertices are the $M$-cubes of $\tilde{\mathbf{Y}}$, two of them being joined by an edge when they have a common face, and the graph formed in the same way for $\mathbf{Y}$; then $\mathrm{pr}_\ell$ is an isomorphism of the first graph onto the second, and a union of cubes is connected through faces exactly when the part of the graph that belongs to it is connected.
A polymer contained in $\mathbf{Y}$ is a union of $M$-cubes of $\mathbf{Y}$; hence it is the image of the union $\tilde{X}$ of the corresponding cubes of $\tilde{\mathbf{Y}}$, and of no other union of cubes of $\tilde{\mathbf{Y}}$.
As $\mathrm{sdm}(\tilde{X})\le\mathrm{sdm}(\tilde{\mathbf{Y}})$, the set $\tilde{X}$ is a lift of its image.
This proves the assertions on the sets $\tilde{X}$, and, with $\tilde{X}=\tilde{\mathbf{Y}}$, those on $\mathbf{Y}$.

The converse is stated in Section~I.1.1; the reason is the following.
Let $\tilde{X}$ be a lift of the polymer $X$; by (a), $\mathrm{pr}_\ell$ is a bijection of $\tilde{X}$ onto $X$.
Let $c,c'\in X$ be at the distance $\eta_k$, let $r$ be the representative of $c-c'$ with $|r|=\eta_k$, and let $\tilde{c}$, $\tilde{c}'$ be the points of $\tilde{X}$ over $c$, $c'$.
Then $\tilde{c}-\tilde{c}'=r+\ell\cdot w$, and $w\ne0$ would give $|\tilde{c}-\tilde{c}'|\ge\ell-\eta_k>\mathrm{sdm}(\tilde{X})$.
So $\tilde{c}-\tilde{c}'=r$: nearest neighbours in $X$ are lifted to nearest neighbours with the same difference.
An $M$-cube $\square$ of $X$ is connected through nearest neighbours.
Let $\tilde{c}_0\in\tilde{X}$ lie over $c_0\in\square$, and let $\tilde{\square}$ be the $M$-cube of $\mathbb{L}_k$ that contains $\tilde{c}_0$; following paths of nearest neighbours in $\square$ one sees that, for every $c\in\square$, the point of $\tilde{X}$ over $c$ is the point of $\tilde{\square}$ over $c$.
Hence $\tilde{X}$ is a union of $M$-cubes; it is finite, as $X$ is, and $\mathrm{sdm}(\tilde{X})+\eta_k<\ell$.
The first part, applied with $\tilde{\mathbf{Y}}:=\tilde{X}$ (connectedness of $\tilde{\mathbf{Y}}$ was not assumed there), shows that $\tilde{X}$ is connected through faces, because $X=\mathrm{pr}_\ell(\tilde{X})$ is.
Moreover $X$ is connected through nearest neighbours, since two $M$-cubes with a common face contain nearest neighbours.
Let $\tilde{X}_1$, $\tilde{X}_2$ be two lifts of $X$, and for a site $c$ of $X$ let $\tilde{c}_1$, $\tilde{c}_2$ be the points of $\tilde{X}_1$, $\tilde{X}_2$ over $c$.
The vector $\tilde{c}_2-\tilde{c}_1\in\ell\mathbb{Z}^3$ does not change when $c$ is replaced by a nearest neighbour in $X$; so it is constant.

(c) Let $X$ be small.
Then
\[
2M\cdot\left(d_M(X)+2\right)\le2M\cdot(L+2)<13L\cdot M\le\ell.
\]
By \cite[Proposition~5.58\,(t10)]{QEDI} one has $X=\mathrm{pr}_\ell(\tilde{X}_0)$ for a finite union $\tilde{X}_0$ of $M$-cubes that is connected through faces and on which $\mathrm{pr}_\ell$ is injective, and $d_M(\tilde{X}_0)=d_M(X)\le L$.
By \cite[Proposition~5.58\,(t7)]{QEDI}, in its form for unions of cubes of $\mathbb{R}^3$,
\[
\mathrm{sdm}(\tilde{X}_0)\le M\cdot\left(d_M(\tilde{X}_0)+2\right)\le M\cdot(L+2),
\]
and $M\cdot(L+2)+2+M\le M\cdot(L+5)<13L\cdot M\le\ell$.
So $\tilde{X}_0$ is a lift in the sense of Section~I.1.1.
By (b), every lift $\tilde{X}$ of $X$ is of the form $\tilde{X}_0+\ell\cdot w$; a translation maps nets to nets of the same length; hence $d_M(\tilde{X})=d_M(\tilde{X}_0)\le L$ and $\mathrm{sdm}(\tilde{X})=\mathrm{sdm}(\tilde{X}_0)$.
Conversely, let $\tilde{X}$ be non-empty and connected through faces, with $d_M(\tilde{X})\le L$.
Then $2M\cdot\left(d_M(\tilde{X})+2\right)<\ell$, and by the converse part of \cite[Proposition~5.58\,(t10)]{QEDI} the map $\mathrm{pr}_\ell$ is injective on $\tilde{X}$, $\mathrm{pr}_\ell(\tilde{X})$ is a polymer, and $d_M\left(\mathrm{pr}_\ell(\tilde{X})\right)=d_M(\tilde{X})\le L$.
By \cite[Proposition~5.58\,(t7)]{QEDI}, $\mathrm{sdm}(\tilde{X})\le M\cdot(L+2)$; therefore $\tilde{X}$ is a lift of the small polymer $\mathrm{pr}_\ell(\tilde{X})$.
This proves the first sentence of (c) as well.
The definition of $\mathrm{Sm}(\tilde{\square})$ involves only cubes, faces and nets of $\mathbb{R}^3$, and the covariance is (5.c).
As to the bijection: by what has been shown, $\mathrm{pr}_\ell$ maps $\mathrm{Sm}(\tilde{\square})$ into the set of the small polymers that contain $\mathrm{pr}_\ell(\tilde{\square})$.
A small polymer $X$ that contains $\mathrm{pr}_\ell(\tilde{\square})$ has a lift, and exactly one of its translates by vectors of $\ell\mathbb{Z}^3$ contains $\tilde{\square}$: the lifts of $\mathrm{pr}_\ell(\tilde{\square})$ are the cubes $\tilde{\square}+\ell\cdot w$, and a lift of $X$ contains exactly one of them, by (a).
This is the last assertion of \cite[Proposition~5.58\,(t10)]{QEDI}.

(d) If (I) holds, then $c$ and $c'$ lie in $\tilde{X}$, and the bonds of the paths have both end points in $\tilde{X}$.
If (II) holds, then $|c-c_3|=\eta_k$, one of the two sites lies in $\tilde{X}$, and the six paths reduce to the bond $(c,c_3)$ (Fact~2).
A pair of $\mathrm{Pz}_k$ is a class $(x,x')+\ell\mathbb{Z}^3$ with $(x,x')\in\mathrm{Pzt}_k$ (Section~I.1.7).
It satisfies (I) for $X$ if and only if $c$, $c'$ and the sites of the paths lie in $\tilde{X}+\ell\mathbb{Z}^3$, and it satisfies (II) for $X$ if and only if $x'=x^{(3)}$ and $c$ or $c_3$ lies in $\tilde{X}+\ell\mathbb{Z}^3$.
Hence the image of a pair that satisfies (I), resp.\ (II), for $\tilde{X}$ satisfies (I), resp.\ (II), for $X$.
Injectivity: if $(x,x')$ and $(x+\ell\cdot w,x'+\ell\cdot w)$ are both admissible for $\tilde{X}$, then $c$ and $c+\ell\cdot w$ lie in $\tilde{X}^{+\eta_k}$, and $\mathrm{sdm}\left(\tilde{X}^{+\eta_k}\right)\le\mathrm{sdm}(\tilde{X})+2\eta_k<\ell$; so $w=0$ by (a).
Surjectivity, together with the kinds: let the class of $(x,x')$ satisfy (I) for $X$.
After a translation by a vector of $\ell\mathbb{Z}^3$ one may assume $c\in\tilde{X}$.
Consecutive sites of a path are at the distance $\eta_k$, and by (a), with $a=\eta_k$, a site of $\tilde{X}+\ell\mathbb{Z}^3$ at the distance $\eta_k$ from a site of $\tilde{X}$ lies in $\tilde{X}$.
Hence all sites of the six paths, and $c'$, lie in $\tilde{X}$.
Let the class satisfy (II) for $X$; after a translation, that one of the sites $c$, $c_3$ which lies in $\tilde{X}+\ell\mathbb{Z}^3$ lies in $\tilde{X}$.
If a class satisfies both conditions for $X$, then the representative with $c\in\tilde{X}$ has $c_3\in\tilde{X}$, by (a); it is the only representative that is admissible for $\tilde{X}$, and it satisfies both conditions for $\tilde{X}$.
\end{proof}

\subsubsection*{1.4. Operators with rows of finite support; local families; passage to a torus}

\emph{Functions.}
For an index set $\mathrm{Ind}$ of Section~I.1.9 let $\mathrm{Fn}(\mathrm{Ind})$ be the space of all functions on $\mathrm{Ind}$ with values in $\mathbb{C}^4$ (when the spinor index is carried) or in $\mathbb{C}$.
A space of bond functions is indexed by the positively oriented bonds $(c,c+a\cdot e_\mu)$, where $a$ is the mesh of the lattice under consideration, and one sets $f(c',c):=-f(c,c')$.
This convention changes neither the supports nor the action of the translations nor the weight $1$; a translation maps positively oriented bonds to positively oriented bonds, and so does the map $b\mapsto b/L$.
Consequently, for every index set $\mathrm{Ind}$, the values $f(t)$, $t\in\mathrm{Ind}$, of a function $f\in\mathrm{Fn}(\mathrm{Ind})$ are independent of one another.
The weights of the indices are
\[
\nu=\frac{\eta_k^3}{48}\ \text{on}\ \mathbb{L}^{\mathrm{f}}_k,\qquad\left(\frac{\eta_k^3}{48}\right)^2\ \text{on}\ \mathrm{Pzt}_k,\qquad\eta_k^3\ \text{on the sites of}\ \mathbb{L}_k,
\]
\[
1\ \text{on}\ \mathbb{Z}^3,\ \text{on bonds and on the basis elements},\qquad L^3\ \text{on}\ L\mathbb{Z}^3
\]
(Sections~I.1.6 and~I.3.6; for the unit bonds and the basis elements, the paragraph on $c_{\gamma}$ in Section~I.5.8; for the bonds of the fine lattices $\mathbb{L}_k$ and $\mathbb{L}_{k+1}$ the weight $1$ is a convention of this chapter).
On the index sets of level $k+1$, which are the sources of the maps of no.~14 in the table of Proposition~25.1, the weights are the same, with $\eta_{k+1}$ in place of $\eta_k$.
All these weights are invariant under translations and do not depend on $\ell$.

\emph{Rows.}
An operator $O:\mathrm{Fn}(\mathrm{Ind}')\to\mathrm{Fn}(\mathrm{Ind})$ has rows of finite support if
\[
(Of)(t)=\sum_{t'\in\mathrm{Ind}'}O[t,t']\cdot f(t')
\]
with coefficients $O[t,t']$, which are $4\times4$ matrices or numbers, such that for every $t$ only finitely many of them differ from $0$.
The row at $t$ is $O[t,\cdot\,]$, and
\[
O(t,t'):=\frac{O[t,t']}{\nu(t')}
\]
is the kernel; for $\mathrm{Ind}'=\mathbb{L}^{\mathrm{f}}_k$ this is the kernel (I.1.4), with the sum extended over all of $\mathbb{L}^{\mathrm{f}}_k$, and for $\mathrm{Ind}'=\mathbb{Z}^3$ it is the coefficient.

\begin{named}{Definition}[Local family]
A local family consists of the following data:
\begin{itemize}
\item[$\bullet$] operators $\mathrm{Op}_{\mathbb{L}}(\tilde{B},\tilde{Y}):\mathrm{Fn}(\mathrm{Ind}')\to\mathrm{Fn}(\mathrm{Ind})$ with rows of finite support, one for every bond function $\tilde{B}$ on $\mathbb{L}_k$ and every union $\tilde{Y}$ of $M$-cubes of $\mathbb{L}_k$;
\item[$\bullet$] a subgroup $\mathrm{Gv}$ of $\eta_k\mathbb{Z}^3$ that acts on $\mathrm{Ind}$ and on $\mathrm{Ind}'$, and that is contained in $M\mathbb{Z}^3$ if the operators depend on $\tilde{Y}$;
\item[$\bullet$] for every $t\in\mathrm{Ind}$, a finite subset $\mathrm{Col}(t)$ of $\mathrm{Ind}'$, a finite set $\mathrm{RdB}(t)$ of bonds of $\mathbb{L}_k$ and a finite set $\mathrm{RdY}(t)$ of sites of $\mathbb{L}_k$, none of which depends on $(\tilde{B},\tilde{Y})$.
\end{itemize}
These data are subject to the following conditions.
\begin{itemize}
\item[(l1)] $\mathrm{Op}_{\mathbb{L}}(\tilde{B},\tilde{Y})[t,t']=0$ for $t'\notin\mathrm{Col}(t)$.
\item[(l2)] $\mathrm{Op}_{\mathbb{L}}(\tilde{B},\tilde{Y})[t,\cdot\,]=\mathrm{Op}_{\mathbb{L}}(\tilde{B}_1,\tilde{Y}_1)[t,\cdot\,]$ whenever $\tilde{B}=\tilde{B}_1$ on $\mathrm{RdB}(t)$ and $\tilde{Y}$, $\tilde{Y}_1$ contain the same sites of $\mathrm{RdY}(t)$.
\item[(l3)] For fixed $\tilde{Y}$, every entry of $\mathrm{Op}_{\mathbb{L}}(\tilde{B},\tilde{Y})[t,t']$ is an entire function of the numbers $\tilde{B}(b)$, $b\in\mathrm{RdB}(t)$.
\item[(l4)] For $v\in\mathrm{Gv}$,
\[
\mathrm{Col}(t+v)=\mathrm{Col}(t)+v,\qquad\mathrm{RdB}(t+v)=\mathrm{RdB}(t)+v,\qquad\mathrm{RdY}(t+v)=\mathrm{RdY}(t)+v,
\]
and
\[
\mathrm{Op}_{\mathbb{L}}(T_v\tilde{B},\tilde{Y}+v)[t+v,t'+v]=\mathrm{Op}_{\mathbb{L}}(\tilde{B},\tilde{Y})[t,t'].
\]
\end{itemize}
The range number of the family is the supremum, over $t\in\mathrm{Ind}$, of the numbers $\mathrm{dist}(\mathrm{supp}\,t,\mathrm{supp}\,t')$ with $t'\in\mathrm{Col}(t)$ and of the numbers $\mathrm{dist}(p,\mathrm{supp}\,t)$ with $p$ an end point of a bond of $\mathrm{RdB}(t)$ or $p\in\mathrm{RdY}(t)$.
\end{named}

\emph{Relation with Section~I.1.9.}
Property (I.f1) is the finiteness of $\mathrm{Col}(t)$ together with (l1).
Property (I.f2), for $v\in\mathrm{Gv}$, is the last relation of (l4), because
\[
\left(T_v\,\mathrm{Op}_{\mathbb{L}}(\tilde{B},\tilde{Y})f\right)(t+v)=\sum_{t'}\mathrm{Op}_{\mathbb{L}}(\tilde{B},\tilde{Y})[t,t']\cdot f(t')
\]
and
\[
\left(\mathrm{Op}_{\mathbb{L}}(T_v\tilde{B},\tilde{Y}+v)\,T_vf\right)(t+v)=\sum_{t'}\mathrm{Op}_{\mathbb{L}}(T_v\tilde{B},\tilde{Y}+v)[t+v,t'+v]\cdot f(t').
\]
Property (I.f4) holds with every $\mathrm{rg}$ that is not smaller than the range number.
Indeed, a column $t'$ with $\mathrm{dist}(\mathrm{supp}\,t,\mathrm{supp}\,t')>\mathrm{rg}$ lies outside $\mathrm{Col}(t)$, and by (l2) the row at $t$ reads $\tilde{B}$ on $\mathrm{RdB}(t)$ and $\tilde{Y}$ through $\mathrm{RdY}(t)$, which are subsets of the bonds and of the sites within the distance $\mathrm{rg}$ of $\mathrm{supp}\,t$; the word ``reads'' is used in the sense of \cite[Definition~5.5\,(1)]{QEDI}.

\begin{named}{Lemma 2}[Algebra of local families]
\begin{itemize}
\item[(a)] A linear combination of local families $\mathrm{Fn}(\mathrm{Ind}')\to\mathrm{Fn}(\mathrm{Ind})$, with constant $4\times4$ matrices as coefficients on either side, is a local family; its sets $\mathrm{Col}$, $\mathrm{RdB}$, $\mathrm{RdY}$ are the unions, and its group $\mathrm{Gv}$ is the intersection, of those of the given families.
\item[(b)] If $\mathrm{Op}_{\mathbb{L}}$ is a local family, then so is $(\tilde{B},\tilde{Y})\mapsto\mathrm{Op}_{\mathbb{L}}(-\tilde{B},\tilde{Y})$, with the same $\mathrm{Col}$, $\mathrm{RdB}$, $\mathrm{RdY}$ and $\mathrm{Gv}$.
\item[(c)] Let $\mathrm{Op}_{\mathbb{L},1}:\mathrm{Fn}(\mathrm{Ind}_1)\to\mathrm{Fn}(\mathrm{Ind}_2)$ and $\mathrm{Op}_{\mathbb{L},2}:\mathrm{Fn}(\mathrm{Ind}_2)\to\mathrm{Fn}(\mathrm{Ind}_3)$ be local families.
Then the composition $\mathrm{Op}_{\mathbb{L},2}\circ\mathrm{Op}_{\mathbb{L},1}$, formed at the same $(\tilde{B},\tilde{Y})$, is a local family, with
\begin{gather*}
(\mathrm{Op}_{\mathbb{L},2}\circ\mathrm{Op}_{\mathbb{L},1})[t,t']=\sum_{t''\in\mathrm{Col}_2(t)}\mathrm{Op}_{\mathbb{L},2}[t,t'']\cdot\mathrm{Op}_{\mathbb{L},1}[t'',t'],\qquad
\mathrm{Col}(t)=\bigcup_{t''\in\mathrm{Col}_2(t)}\mathrm{Col}_1(t''),\\
\mathrm{RdB}(t)=\mathrm{RdB}_2(t)\cup\bigcup_{t''\in\mathrm{Col}_2(t)}\mathrm{RdB}_1(t''),\qquad
\mathrm{RdY}(t)=\mathrm{RdY}_2(t)\cup\bigcup_{t''\in\mathrm{Col}_2(t)}\mathrm{RdY}_1(t''),\qquad
\mathrm{Gv}=\mathrm{Gv}_1\cap\mathrm{Gv}_2.
\end{gather*}
\end{itemize}
\end{named}

\begin{proof}
(a), (b): each of the conditions (l1)--(l4) is inherited term by term; for (b) one uses $T_v(-\tilde{B})=-T_v\tilde{B}$.
(c) The sums are finite.
As to (l1): if $t'$ lies outside all the sets $\mathrm{Col}_1(t'')$, $t''\in\mathrm{Col}_2(t)$, then every term vanishes.
As to (l2): if $\tilde{B}=\tilde{B}_1$ on $\mathrm{RdB}(t)$ and $\tilde{Y}$, $\tilde{Y}_1$ contain the same sites of $\mathrm{RdY}(t)$, then the factors $\mathrm{Op}_{\mathbb{L},2}[t,t'']$ coincide, and so do the rows $\mathrm{Op}_{\mathbb{L},1}[t'',\cdot\,]$ for $t''\in\mathrm{Col}_2(t)$.
As to (l3): the entries are finite sums of products of entire functions.
As to (l4): one replaces $t''$ by $t''+v$ in the sum.
\end{proof}

\begin{named}{Lemma 3}[Passage to a torus]
Let $\mathrm{Op}_{\mathbb{L}}$ be a local family with $\ell\mathbb{Z}^3\subset\mathrm{Gv}$, let $\mathrm{rg}$ be at least its range number, and let $\tilde{B}$ and $\tilde{Y}$ be $\ell\mathbb{Z}^3$-periodic.
Let $\epsilon$ be one of the two characters $\epsilon(w)=1$ and $\epsilon(w)=\boldsymbol{\chi}(w)$ of $\mathbb{Z}^3$, and put
\[
\mathrm{Fn}_{\epsilon}(\mathrm{Ind}):=\left\{f\in\mathrm{Fn}(\mathrm{Ind}):\ f(t+\ell\cdot w)=\epsilon(w)\cdot f(t)\right\},
\]
the space of the periodic, resp.\ antiperiodic, functions of Section~I.1.9.
\begin{itemize}
\item[(a)] The operator $\mathrm{Op}_{\mathbb{L}}(\tilde{B},\tilde{Y})$ commutes with $T_{\ell\cdot w}$ and maps $\mathrm{Fn}_{\epsilon}(\mathrm{Ind}')$ into $\mathrm{Fn}_{\epsilon}(\mathrm{Ind})$.
\item[(b)] Let $\mathrm{FD}'$ be a fundamental domain of $\ell\mathbb{Z}^3$ in $\mathrm{Ind}'$; the action is free, because $\mathrm{supp}(t'+v)=\mathrm{supp}\,t'+v$ and the supports are finite and non-empty.
For $f\in\mathrm{Fn}_{\epsilon}(\mathrm{Ind}')$,
\[
\left(\mathrm{Op}_{\mathbb{L}}(\tilde{B},\tilde{Y})f\right)(t)=\sum_{t'\in\mathrm{FD}'}\mathrm{Op}[t,t']\cdot f(t'),\qquad
\mathrm{Op}[t,t']:=\sum_{w\in\mathbb{Z}^3}\epsilon(w)\cdot\mathrm{Op}_{\mathbb{L}}(\tilde{B},\tilde{Y})[t,t'+\ell\cdot w],
\]
where the last sum is finite.
The values $f(t')$, $t'\in\mathrm{FD}'$, are free: every function on $\mathrm{FD}'$ has exactly one extension in $\mathrm{Fn}_{\epsilon}(\mathrm{Ind}')$, and for bond functions $\mathrm{FD}'$ consists of positively oriented bonds, by the paragraph on functions above.
Therefore the numbers, resp.\ matrices, $\mathrm{Op}[t,t']$, $t'\in\mathrm{FD}'$, are determined by the operator: they are the coefficients of the restriction of $\mathrm{Op}_{\mathbb{L}}(\tilde{B},\tilde{Y})$ to $\mathrm{Fn}_{\epsilon}(\mathrm{Ind}')$, and $\frac{\mathrm{Op}[t,t']}{\nu(t')}$ is the kernel of this restriction.
Moreover $\mathrm{Op}[t,t'+\ell\cdot w]=\epsilon(w)\cdot\mathrm{Op}[t,t']$.
\item[(c)] If $\mathrm{supp}\,t$ and $\mathrm{supp}\,t'$ are contained in a set $\tilde{U}$ with $\mathrm{sdm}(\tilde{U})+\mathrm{rg}<\ell$, then $\mathrm{Op}[t,t']=\mathrm{Op}_{\mathbb{L}}(\tilde{B},\tilde{Y})[t,t']$.
\end{itemize}
\end{named}

\begin{proof}
(a) This is (l4) with $v=\ell\cdot w$, since $T_{\ell\cdot w}\tilde{B}=\tilde{B}$ and $\tilde{Y}+\ell\cdot w=\tilde{Y}$.
If $T_{-\ell\cdot w}f=\epsilon(w)\cdot f$, then
\[
T_{-\ell\cdot w}\,\mathrm{Op}_{\mathbb{L}}f=\mathrm{Op}_{\mathbb{L}}\,T_{-\ell\cdot w}f=\epsilon(w)\cdot\mathrm{Op}_{\mathbb{L}}f.
\]
(b) Every index of $\mathrm{Ind}'$ can be written as $t'+\ell\cdot w$ with exactly one $t'\in\mathrm{FD}'$ and $w\in\mathbb{Z}^3$, and $f(t'+\ell\cdot w)=\epsilon(w)\cdot f(t')$.
(c) For $w\ne0$, $p\in\mathrm{supp}\,t$ and $p'\in\mathrm{supp}\,t'$, Lemma~1\,(a) gives $|p-p'-\ell\cdot w|>\mathrm{rg}$.
Hence $\mathrm{dist}\left(\mathrm{supp}\,t,\mathrm{supp}(t'+\ell\cdot w)\right)>\mathrm{rg}$, and $t'+\ell\cdot w$ lies outside $\mathrm{Col}(t)$.
\end{proof}

\subsubsection*{1.5. The algebras $\mathrm{Grc}_k(\tilde U)$ and the maps $\mathrm{cov}_{\ell}$}

In this subsection the level is $j\in\{k,k+1\}$, and $\ell'$ is a side that is a multiple of $M$ with $\ell'>2M$.
Both $(j,\ell')=(k,\ell)$ with $\ell\in\mathrm{Cl}$ and $(j,\ell')=(k+1,\ell/L)$ with $\ell\in\mathrm{Cl}$ are of this kind; in the second case $\ell/L$ is an odd multiple of $M$ and $\ell/L\ge13M$.
For such $(j,\ell')$ the set $\mathrm{Pz}_j$ and the generators of the torus of side $\ell'$ are defined in Sections~I.1.7 and~I.3.6; for the level $k+1$ and the side $\ell/L$ see Sections~I.3.7 and~I.5.11.
The letter $\varphi$ denotes any one of $\psi$, $\chi$, $\Psi$, $W$, and $t$ denotes its index; $\varphi^{(\ell')}(t)$ is the generator of the torus of side $\ell'$ with the covering index $t$.
By (I.AP),
\[
\varphi^{(\ell')}(t+\ell'w)=\boldsymbol{\chi}(w)\cdot\varphi^{(\ell')}(t)\qquad(w\in\mathbb{Z}^3),
\]
and the $\varphi^{(\ell')}(t)$ with $t$ in a fundamental domain of $\ell'\mathbb{Z}^3$ are independent generators.

Let $\mathrm{Par}$ denote the automorphism of a Grassmann algebra that maps every generator to its negative.
An element $F$ is called even if $\mathrm{Par}\,F=F$, and odd if $\mathrm{Par}\,F=-F$.
For $v\in\mathbb{Z}^3$, and on subalgebras generated by $\psi$, $\chi$ also for $v\in\eta_j\mathbb{Z}^3$, let $T_v$ be the isomorphism $\mathrm{Grc}_j(\tilde U)\to\mathrm{Grc}_j(\tilde U+v)$ with $T_v\varphi(t)=\varphi(t+v)$.
Denote by $\mathrm{As}_j(\tilde U)$ the subalgebra of $\mathrm{Grc}_j(\tilde U)$ generated by those generators of $\mathrm{Grc}_j(\tilde U)$ that belong to $\psi$ or to $\chi$.
For $\tilde U\subset\tilde U_1$ the algebra $\mathrm{Grc}_j(\tilde U)$ is a subalgebra of $\mathrm{Grc}_j(\tilde U_1)$, and $\mathrm{As}_j(\tilde U)$ is a subalgebra of $\mathrm{As}_j(\tilde U_1)$.

\begin{named}{Lemma 4}
Let $\tilde U$ be a finite non-empty set of sites of $\mathbb{L}_j$ with $\mathrm{sdm}(\tilde U)<\ell'$.
\begin{itemize}
\item[(a)] The assignment $\varphi(t)\mapsto\varphi^{(\ell')}(t)$ defines a homomorphism $\mathrm{cov}_{\ell'}:\mathrm{Grc}_j(\tilde U)\to\mathrm{Gr}[\psi,\chi,\Psi,W]$ into the algebra of the torus of side $\ell'$.
Up to signs, the images of distinct generators of $\mathrm{Grc}_j(\tilde U)$ are distinct members of one system of independent generators.
Hence $\mathrm{cov}_{\ell'}$ is injective; up to signs it maps monomials in the generators to monomials, distinct ones to distinct ones; and it preserves the degree in each family, the bar index and the parity.
It is the injective homomorphism of the last sentence of the paragraph ``The map $\mathrm{cov}_{\ell}$'' in Section~G3.5.0.
If $\tilde U\subset\tilde U_1$ with $\mathrm{sdm}(\tilde U_1)<\ell'$, the map for $\tilde U$ is the restriction of the map for $\tilde U_1$.
If $\mathrm{sdm}(\tilde U)+M<\ell'$, the map is the map $\mathrm{cov}_{\ell'}$ of Section~G3.5.0, and every pair of $\mathrm{Pz}_j$ whose two photon sites lie in $\mathrm{pr}_{\ell'}(\tilde U)$ is the image of exactly one pair of $\mathrm{Pzt}_j$ whose two photon sites lie in $\tilde U$.
\item[(b)] For $w\in\mathbb{Z}^3$,
\[
\mathrm{cov}_{\ell'}\circ T_{\ell'w}=\mathrm{cov}_{\ell'}\circ\mathrm{Par}^{w_1+w_2+w_3}\qquad\text{on }\mathrm{Grc}_j(\tilde U).
\]
In particular $\mathrm{cov}_{\ell'}(T_{\ell'w}F)=\mathrm{cov}_{\ell'}\left(F\right)$ for every even $F$.
\item[(c)] (Periodic summands.) Let $X$ be a set of photon sites of the torus that has a lift $\tilde X$ (Section~I.1.1), and let $\mathrm{hf}$ be a function on $\mathbb{L}^{\mathrm{f}}_j$ with values in a vector space such that $\mathrm{hf}(x+\ell'w)=\mathrm{hf}(x)$ for all $x$ and all $w\in\mathbb{Z}^3$.
Then
\[
\int_X\mathrm{d}x\,\mathrm{hf}(x)=\int_{\tilde X}\mathrm{d}x\,\mathrm{hf}(x).
\]
If $\mathrm{hf}$ is a function on $\mathbb{Z}^3$ with $\mathrm{hf}(y+\ell'w)=\mathrm{hf}(y)$ for all $y$ and $w$, then the sum of $\mathrm{hf}$ over the unit sites of $X$ equals the sum of $\mathrm{hf}$ over $\tilde X\cap\mathbb{Z}^3$.
\item[(d)] (Change of the lift.) Let $w\in\mathbb{Z}^3$.
\begin{itemize}
\item[(d1)] If $F\in\mathrm{Grc}_j(\tilde U)$ is even, then $\mathrm{cov}_{\ell'}(T_{\ell'w}F)=\mathrm{cov}_{\ell'}\left(F\right)$.
\item[(d2)] Let $\mathrm{As}$ be a subalgebra of $\mathrm{Grc}_j(\tilde U)$ that is generated by generators, let $\mathrm{hm}:\mathrm{As}\to\mathrm{Grc}_j(\tilde U)$ be a linear map with $\mathrm{hm}\circ\mathrm{Par}=\mathrm{Par}\circ\mathrm{hm}$, and let $O$, $O'$ be linear maps on $\mathrm{cov}_{\ell'}(\mathrm{As})$ with
\[
\mathrm{cov}_{\ell'}\circ\mathrm{hm}=O\circ\mathrm{cov}_{\ell'}\ \text{ on }\mathrm{As},\qquad\mathrm{cov}_{\ell'}\circ\left(T_{\ell'w}\,\mathrm{hm}\,T_{\ell'w}^{-1}\right)=O'\circ\mathrm{cov}_{\ell'}\ \text{ on }T_{\ell'w}\mathrm{As}.
\]
Then $O'=O$.
\end{itemize}
\item[(e)] (Operations.) The map $\mathrm{cov}_{\ell'}$ commutes with finite sums, with products, with multiples, with $F\mapsto\exp F$, and with $F\mapsto F^{[n,n']}$, the part of bidegree $(n,n')$ in $(\psi,\chi)$.
Let $\mathrm{Is}$ be a set of indices of the generators $\psi$, $\chi$ over $\tilde U$, let $\mathrm{As}$ be the subalgebra generated by the $\varphi(t)$, $t\in\mathrm{Is}$, and let $v(t)$, $t\in\mathrm{Is}$, be odd elements of $\mathrm{Grc}_j(\tilde U)$.
Let $\mathrm{ad}_v:\mathrm{As}\to\mathrm{Grc}_j(\tilde U)$ be the homomorphism with $\mathrm{ad}_v\varphi(t)=\varphi(t)+v(t)$.
If $v^{(\ell')}$ is a family of odd elements of the algebra of the torus, indexed like $\psi^{\#}=(\psi,\chi)$ and subject to (I.AP), with $v^{(\ell')}(t)=\mathrm{cov}_{\ell'}(v(t))$ for $t\in\mathrm{Is}$, then
\[
\mathrm{cov}_{\ell'}\circ\mathrm{ad}_v=\mathrm{ad}_{v^{(\ell')}}\circ\mathrm{cov}_{\ell'}\qquad\text{on }\mathrm{As},
\]
where $\mathrm{ad}_{v^{(\ell')}}$ is the homomorphism of (D1) in Definition~24.D that adds $v^{(\ell')}$ to the generators.
\end{itemize}
\end{named}

\begin{proof}
(a) The algebra $\mathrm{Grc}_j(\tilde U)$ is the free Grassmann algebra on its generators, and the elements $\varphi^{(\ell')}(t)$ anticommute; therefore the homomorphism exists.
If two indices of one family are congruent modulo $\ell'\mathbb{Z}^3$, they are of the form $t$ and $t+\ell'w$, with the first photon sites $c$ and $c+\ell'w$; when both of these sites lie in $\tilde U$, Lemma~1\,(a) with $a=0$ gives $w=0$.
Distinct indices over $\tilde U$ therefore belong to distinct classes modulo $\ell'\mathbb{Z}^3$.
After the choice of a fundamental domain, (I.AP) shows that each $\varphi^{(\ell')}(t)$ is plus or minus the generator at the representative of $t$.
A homomorphism that sends the free generators injectively, up to signs, to members of a system of free generators sends the basis of the monomials injectively, up to signs, to members of the basis of the monomials; such a homomorphism is injective.
The degrees and the bar indices are those of the generators.
For the last sentence of (a), recall that a pair of $\mathrm{Pz}_j$ is the class of a pair $(x,x')$ of $\mathrm{Pzt}_j$, with $|c-c'|\le M$.
Assume that $\mathrm{pr}_{\ell'}(c)$ and $\mathrm{pr}_{\ell'}(c')$ lie in $\mathrm{pr}_{\ell'}(\tilde U)$.
After a translation one may take $c\in\tilde U$; write $c'=c''+\ell'w$ with $c''\in\tilde U$.
Then $|c-c''-\ell'w|\le M$, and Lemma~1\,(a) with $a=M$ yields $w=0$.
Uniqueness is the fact, shown above, that two indices over $\tilde U$ which are congruent modulo $\ell'\mathbb{Z}^3$ coincide.

(b) Both sides are homomorphisms, and on a generator
\[
\mathrm{cov}_{\ell'}\left(\varphi(t+\ell'w)\right)=\varphi^{(\ell')}(t+\ell'w)=\boldsymbol{\chi}(w)\cdot\varphi^{(\ell')}(t)=\mathrm{cov}_{\ell'}\left(\mathrm{Par}^{w_1+w_2+w_3}\varphi(t)\right).
\]

(c) The integral $\int_X\mathrm{d}x\,\mathrm{hf}(x)$ equals $\frac{\eta_j^3}{48}$ times the sum of $\mathrm{hf}$ over an arbitrary set of representatives of the fermion sites over $X$.
The fermion sites over $\tilde X$ form such a set, because the covering map is a bijection of $\tilde X$ onto $X$.
Since $\ell'\mathbb{Z}^3\subset\mathbb{Z}^3$, the unit sites of $\tilde X$ form a set of representatives of the unit sites of $X$.

(d) Clause (d1) is (b).
For (d2) let $F\in\mathrm{As}$ and $n:=w_1+w_2+w_3$; then
\begin{align*}
O'\,\mathrm{cov}_{\ell'}\left(\mathrm{Par}^nF\right)&=O'\,\mathrm{cov}_{\ell'}\left(T_{\ell'w}F\right)=\mathrm{cov}_{\ell'}\left(T_{\ell'w}\,\mathrm{hm}\,F\right)=\mathrm{cov}_{\ell'}\left(\mathrm{Par}^n\,\mathrm{hm}\,F\right)\\
&=\mathrm{cov}_{\ell'}\left(\mathrm{hm}\,\mathrm{Par}^nF\right)=O\,\mathrm{cov}_{\ell'}\left(\mathrm{Par}^nF\right),
\end{align*}
and $\mathrm{Par}^n$ maps $\mathrm{As}$ onto itself.

(e) The map $\mathrm{cov}_{\ell'}$ is a homomorphism of finite dimensional algebras that preserves the degrees.
The two sides of the last identity are homomorphisms on $\mathrm{As}$; on a generator $\varphi(t)$, $t\in\mathrm{Is}$, both take the value $\varphi^{(\ell')}(t)+\mathrm{cov}_{\ell'}\left(v(t)\right)$.
\end{proof}

\subsection*{2. Proposition~25.1}

\subsubsection*{2.1. The statement}

Proposition~25.1 is stated in Part~II, for the data of Section~G3.5.0, and it asserts the following.
For every side $\ell\in\mathrm{Cl}$, every complex bond function $B$ on $T_k$ and, where $Y$ occurs, every union $Y$ of $M$-cubes of $T_k$, each operator $\mathrm{Op}$ of its table is defined by the display named in the corresponding row.
Read on the covering lattice, the same display defines, for all $\tilde B$ and $\tilde Y$, a linear operator $\mathrm{Op}_{\mathbb{L}}(\tilde B,\tilde Y)$ between the spaces of all functions on the index sets named in the row.
These operators have the properties (I.f1)--(I.f4) with the range $\mathrm{rg}$ given in the table; that is, $\mathrm{Op}$ is given by one formula on $\mathbb{L}^{\mathrm{f}}_k$, of range $\mathrm{rg}$.
The property (I.f2) holds for all $v\in L\mathbb{Z}^3$ if $\mathrm{Op}$ does not depend on $Y$, and for all $v\in M\mathbb{Z}^3$ otherwise.
The formulas contain $k$, $\eta_k$, $\mu$, $e_k$, $b_k$, $b_L$, $\alpha$, $\rho$, the quadruple $w=(w_0,w_1,w_2,w_3)$ of the weights of Section~I.0.5, $L$, $M$, the constant matrices of Section~I.2.1, the basis $\{e^{0}_{i}\}$ of \cite[Section~5.2.5]{QEDI}, which is chosen once for each $L$ and is the same for all $k$ and $\ell$, and the resolvent parameter; they do not contain $\ell$.
The entries of every row are entire functions of those numbers $\tilde B(b)$ which the row reads according to the table.
The table has fourteen rows, numbered 1 to 14; each of them names one or several operators and gives the levels, the index sets between which the operators act, the defining display, the range $\mathrm{rg}$ and what the row of such an operator at the index $t$ reads.
Moreover, the proposition gives on the covering lattice the formulas (K1)--(K4), in which a kernel is taken by (I.1.4) where the operator acts on spinors on fermion sites, and with the sum over $y'$ where it acts on $\mathfrak{S}(1,\ell)$.
Here (K1) gives the kernels of $Q_k(\tilde B)\cdot J^{*}$ and of $J\cdot Q_k^{T}(-\tilde B)$; (K2) the kernel of $P(\tilde B)$ and, for $k=0$, that of $J\cdot P(\tilde B)\cdot J^{*}$; (K3) the expressions of $B_{k,y}(\tilde B)$ and of $B_{k,y}(\tilde B)\cdot\gamma_5\cdot B_{k,y}(\tilde B)$ through $P(\tilde B)$; and (K4), for $k=0$, the kernel of $F_0(\tilde B)$.
The last paragraph of the proposition deals with the operators $S_k(A)$, $S_{k,y}(A)$, $H_{k}(A)$, $\bar{H}_{k}(A)$, $\mathfrak{H}_{k}(A)$, $D_k(A)$, $M_k(A)$, $\dot{M}_k(A)$, $\Gamma_k(A)$, $\Gamma_{k,y}(A)$ and with the photon operators $H_{k}$, $C_{k}^{1/2}$, $\mathrm{Sq}_{k}$, $A^{\min}_{k}$.
It asserts that none of these operators is an entry of the table, and that none of them enters Definition~24.D, the definitions of Section~I.5.6 or the definition of $\acute{E}_k$.

\subsubsection*{2.2. Proof, first part: the operators on the covering lattice, row by row}

For every number of the table of Proposition~25.1 the display named there is read for an arbitrary $f\in\mathrm{Fn}(\cdot)$, an arbitrary bond function $\tilde B$ on $\mathbb{L}_k$ and an arbitrary union $\tilde Y$ of $M$-cubes of $\mathbb{L}_k$.
All sums that occur below are finite; hence each display defines an operator with rows of finite support between the spaces of all functions.
For each operator the following data are determined: the coefficients; the sets $\mathrm{Col}(t)$, $\mathrm{RdB}(t)$, $\mathrm{RdY}(t)$, each of which is empty when it is not named; the largest distance, that is, the range number of 1.4; and the group $\mathrm{Gv}$.
The conditions (l1) and (l2) can be read off the coefficients.
As to (l3), every coefficient is a finite sum of products of constant matrices and of numbers $u_{\tilde B}$, $U_{\tilde B}$, $\varphi_{\tilde B}$ and their inverses, and by Facts~2--4 these numbers are exponentials of linear forms in the $\tilde B(b)$, $b\in\mathrm{RdB}(t)$.
For the operators (Of1)--(Of5), (Of10), (Of11), (Of13), whose coefficients are given by explicit formulas, (l4) follows from Fact~5; for (Of12) and (Of14) it is verified in the corresponding paragraphs; for the remaining operators it follows from Lemma~2.

\paragraph{(Of1) [no.~1]}
On $\mathrm{Fn}(\mathbb{L}^{\mathrm{f}}_k)$ let $(M_if)(x):=M_i(x)\cdot f(x)$, where by (I.2.1)
\[
M_i(x)=\frac{w_i}{2}\cdot\left(\rho+\gamma\cdot(g\hat n_i)\right),
\]
a function of the label $g$ alone.
For a constant $4\times4$ matrix $\mathrm{Mx}$ and any index set that carries the spinor index let $(\mathrm{Mx}\,f)(t):=\mathrm{Mx}\cdot f(t)$.
The coefficients are $M_i[x,x']=\mathbf{1}[x'=x]\cdot M_i(x)$ and $\mathrm{Mx}[t,t']=\mathbf{1}[t'=t]\cdot\mathrm{Mx}$; further $\mathrm{Col}(t)=\{t\}$, and the distance is $0$.
For $M_i$ one has $\mathrm{Gv}=\eta_k\mathbb{Z}^3$ by (5.a); for $\mathrm{Mx}$, $\mathrm{Gv}$ is the group that acts on the index set.

\paragraph{(Of2) [no.~2]}
Let $(J\xi)(c,g):=\xi(c)$ and $(J^{*}f)(c):=\frac{1}{48}\sum_{g\in O_{\mathrm{cub}}}f(c,g)$.
Then
\[
J[x,c']=\mathbf{1}[c'=c]\cdot\mathbf{1}_4,\quad\mathrm{Col}(x)=\{c\};\qquad J^{*}[c,x']=\frac{1}{48}\cdot\mathbf{1}[c'=c]\cdot\mathbf{1}_4,\quad\mathrm{Col}(c)=\{c\}\times O_{\mathrm{cub}},
\]
so that the row of $J^{*}$ has $48$ columns.
The distance is $0$, and $\mathrm{Gv}=\eta_k\mathbb{Z}^3$.
On the functions of the sites of $\mathbb{L}_k$ one has $J^{*}J=1$, since $\frac{1}{48}\sum_g\xi(c)=\xi(c)$.

\paragraph{(Of3) [no.~3]}
Let $(R_i(\tilde B)f)(x):=U_{\tilde B}(x,x^{(i)})\cdot f(x^{(i)})$, where $U_{\tilde B}(x,x^{(i)})=1$ for $i=0,1,2$ and
\[
U_{\tilde B}(x,x^{(3)})=u_{\tilde B}(c,c_3)=\exp\left\{i\cdot e_k\cdot\eta_k\cdot\tilde B(c,c_3)\right\}
\]
((I.1.6), (I.2.2)).
Thus $R_i(\tilde B)[x,x']=\mathbf{1}[x'=x^{(i)}]\cdot U_{\tilde B}(x,x^{(i)})\cdot\mathbf{1}_4$.
By (I.2.3),
\[
(D_{\tilde B}f)(x):=-\eta_k^{-1}\sum_{i=0}^{3}M_i(x)\cdot\left[U_{\tilde B}(x,x^{(i)})\cdot f(x^{(i)})-f(x)\right].
\]
The five sites $x$, $x^{(0)}$, $x^{(1)}$, $x^{(2)}$, $x^{(3)}$ are distinct, because $gs_i\ne g$, the matrices $s_0$, $s_1$, $s_2$ are distinct, and $c_3\ne c$; moreover $\sum_iM_i(x)=\frac{\rho}{2}\cdot\mathrm{sw}\cdot\mathbf{1}_4$ by (I.m1).
Consequently
\begin{gather*}
D_{\tilde B}[x,x]=\eta_k^{-1}\cdot\frac{\rho}{2}\cdot\mathrm{sw}\cdot\mathbf{1}_4,\qquad D_{\tilde B}[x,x^{(i)}]=-\eta_k^{-1}\cdot M_i(x)\quad(i=0,1,2),\\
D_{\tilde B}[x,x^{(3)}]=-\eta_k^{-1}\cdot M_3(x)\cdot u_{\tilde B}(c,c_3),
\end{gather*}
and all other coefficients vanish.
For $R_i$ one has $\mathrm{Col}(x)=\{x^{(i)}\}$, and for $D_{\tilde B}$ one has $\mathrm{Col}(x)=\{x,x^{(0)},x^{(1)},x^{(2)},x^{(3)}\}$; further $\mathrm{RdB}(x)=\{(c,c_3)\}$ for $R_3$ and $D_{\tilde B}$, while $\mathrm{RdB}(x)$ is empty for $R_0$, $R_1$, $R_2$.
The sites $x$, $x^{(0)}$, $x^{(1)}$, $x^{(2)}$ lie over $c$, the site $x^{(3)}$ lies over $c_3=c+\eta_k\cdot ge_1$, and $|c_3-c|=\eta_k$.
The largest distance is $\eta_k$, and $\mathrm{Gv}=\eta_k\mathbb{Z}^3$ by (5.a), (5.b).

\paragraph{(Of4) [no.~4; $k\ge1$]}
By Section~I.3.3,
\[
(Q_k(\tilde B)\xi)(y):=\eta_k^3\sum_{c\in\Delta_y}\varphi_{\tilde B}(y,c)\cdot\xi(c),\qquad\left(Q_k^{T}(-\tilde B)v\right)(c):=\varphi_{\tilde B}([c],c)^{-1}\cdot v([c]).
\]
Together with (Of2) this gives
\begin{align*}
\left(Q_k(\tilde B)J^{*}f\right)(y)&=\frac{\eta_k^3}{48}\sum_{c\in\Delta_y}\sum_{g\in O_{\mathrm{cub}}}\varphi_{\tilde B}(y,c)\cdot f(c,g)=\int_{\Delta_y}\mathrm{d}x\,\varphi_{\tilde B}(y,c)\cdot f(x),\\
\left(JQ_k^{T}(-\tilde B)v\right)(x)&=\varphi_{\tilde B}([c],c)^{-1}\cdot v([c]),
\end{align*}
and therefore
\[
\left(Q_k(\tilde B)J^{*}\right)[y,x']=\frac{\eta_k^3}{48}\cdot\mathbf{1}[c'\in\Delta_y]\cdot\varphi_{\tilde B}(y,c')\cdot\mathbf{1}_4,\qquad\left(JQ_k^{T}(-\tilde B)\right)[x,y']=\mathbf{1}\left[y'=[c]\right]\cdot\varphi_{\tilde B}([c],c)^{-1}\cdot\mathbf{1}_4.
\]
For the row of $Q_k(\tilde B)J^{*}$ at $y$: $\mathrm{Col}(y)=\Delta_y\times O_{\mathrm{cub}}$, and $\mathrm{RdB}(y)$ is the union of the sets $\mathrm{Twr}(c)$, $c\in\Delta_y$.
For the row of $JQ_k^{T}(-\tilde B)$ at $x$: $\mathrm{Col}(x)=\{[c]\}$ and $\mathrm{RdB}(x)=\mathrm{Twr}(c)$.
By (1.a) and Fact~3 all distances are at most $\frac{1-\eta_k}{2}$, and all bonds named are contained in $\Delta_y$, respectively in $\Delta_{[c]}$.
By (5.c), $\mathrm{Gv}=\mathbb{Z}^3$.
Finally $Q_k(\tilde B)Q_k^{T}(-\tilde B)=1$, because
\[
\eta_k^3\sum_{c\in\Delta_y}\varphi_{\tilde B}(y,c)\cdot\varphi_{\tilde B}(y,c)^{-1}\cdot v(y)=\eta_k^3\cdot\eta_k^{-3}\cdot v(y).
\]

\paragraph{(Of5) [no.~5]}
By Section~I.3.3,
\[
(Q(\tilde B)v)(y'):=L^{-3}\sum_{y\in B(y')}U_{\tilde B}(y',y)\cdot v(y),\qquad\left(Q^{T}(-\tilde B)v'\right)(y):=U_{\tilde B}([y]^L,y)^{-1}\cdot v'([y]^L),
\]
so that
\[
Q(\tilde B)[y',y]=L^{-3}\cdot\mathbf{1}[y\in B(y')]\cdot U_{\tilde B}(y',y)\cdot\mathbf{1}_4,\qquad Q^{T}(-\tilde B)[y,y'']=\mathbf{1}\left[y''=[y]^L\right]\cdot U_{\tilde B}(y'',y)^{-1}\cdot\mathbf{1}_4.
\]
For the row of $Q(\tilde B)$ at $y'$: $\mathrm{Col}(y')=B(y')$, and $\mathrm{RdB}(y')$ is the union of the sets $\mathrm{Pth}(y',y)$, $y\in B(y')$.
For the row of $Q^{T}(-\tilde B)$ at $y$: $\mathrm{Col}(y)=\{[y]^L\}$ and $\mathrm{RdB}(y)=\mathrm{Pth}([y]^L,y)$.
By (1.b) and Fact~4 all distances are at most $\frac{L-1}{2}$, and $\mathrm{Gv}=L\mathbb{Z}^3$ by (5.b), (5.c).
Moreover $Q(\tilde B)Q^{T}(-\tilde B)=1$, since
\[
L^{-3}\sum_{y\in B(y')}U_{\tilde B}(y',y)\cdot U_{\tilde B}(y',y)^{-1}\cdot v'(y')=v'(y').
\]
\emph{Level~0.}
Here
\begin{align*}
\left(Q(\tilde B)J^{*}f\right)(y')&=L^{-3}\sum_{c\in B(y')}U_{\tilde B}(y',c)\cdot\frac{1}{48}\sum_{g\in O_{\mathrm{cub}}}f(c,g),\\
\left(JQ^{T}(-\tilde B)v'\right)(x)&=U_{\tilde B}([c]^L,c)^{-1}\cdot v'([c]^L).
\end{align*}
The row at $y'$ has $\mathrm{Col}(y')=B(y')\times O_{\mathrm{cub}}$, that is $48L^3$ columns, with the coefficients $\frac{L^{-3}}{48}\cdot U_{\tilde B}(y',c')\cdot\mathbf{1}_4$; the row at $x$ has $\mathrm{Col}(x)=\{[c]^L\}$.
The sets $\mathrm{RdB}$, the distances and $\mathrm{Gv}$ are as before.

\paragraph{(Of6) [no.~6]}
Let $P(\tilde B):=Q^{T}(-\tilde B)Q(\tilde B)$.
By Lemma~2\,(c) and (Of5), with $y':=[y_1]^L$,
\[
P(\tilde B)[y_1,y_2]=\mathbf{1}\left[[y_2]^L=y'\right]\cdot L^{-3}\cdot U_{\tilde B}(y',y_1)^{-1}\cdot U_{\tilde B}(y',y_2)\cdot\mathbf{1}_4,\qquad\mathrm{Col}(y_1)=B(y'),
\]
and $\mathrm{RdB}(y_1)$ is the union of the sets $\mathrm{Pth}(y',y)$, $y\in B(y')$.
For $y_2\in B(y')$ one has $|y_1-y_2|\le L-1$.
For an end $p$ of a bond of $\mathrm{RdB}(y_1)$, Fact~4 gives $|p-y_1|\le|p-y'|+|y'-y_1|\le L-1$.
The largest distance is $L-1$, and $\mathrm{Gv}=L\mathbb{Z}^3$.
Further
\[
P(\tilde B)^2=Q^{T}(-\tilde B)\left[Q(\tilde B)Q^{T}(-\tilde B)\right]Q(\tilde B)=P(\tilde B),
\]
and $P(\tilde B)$ commutes with every constant matrix, since its coefficients are multiples of $\mathbf{1}_4$.

\emph{The operator $B_{k,y}(\tilde B)$ for $k\ge1$ and real $y$.}
One has $b_k>0$ and $b_L>0$ (Section~I.3.4), and $\gamma_5^2=1$ by (I.c1).
Hence, for every $a>0$,
\[
(a+i\cdot\gamma_5\cdot y)(a-i\cdot\gamma_5\cdot y)=a^2+y^2>0,\qquad(a+i\cdot\gamma_5\cdot y)^{-1}=\left(a^2+y^2\right)^{-1}\cdot(a-i\cdot\gamma_5\cdot y).
\]
On $\mathrm{Fn}(\mathbb{Z}^3)$,
\[
b_k+b_L\cdot P(\tilde B)+i\cdot\gamma_5\cdot y=(b_k+i\cdot\gamma_5\cdot y)\left(1-P(\tilde B)\right)+(b_k+b_L+i\cdot\gamma_5\cdot y)\,P(\tilde B).
\]
Since $P(\tilde B)$ is idempotent and commutes with $\gamma_5$, the operator
\begin{equation}
(b_k+i\cdot\gamma_5\cdot y)^{-1}\left(1-P(\tilde B)\right)+(b_k+b_L+i\cdot\gamma_5\cdot y)^{-1}P(\tilde B)\tag{25.1}
\end{equation}
gives $\left(1-P(\tilde B)\right)+P(\tilde B)=1$ when it is multiplied from either side by $b_k+b_L\cdot P(\tilde B)+i\cdot\gamma_5\cdot y$.
Therefore $b_k+b_L\cdot P(\tilde B)+i\cdot\gamma_5\cdot y$ is a bijection of $\mathrm{Fn}(\mathbb{Z}^3)$, the display of Section~I.3.5 defines $B_{k,y}(\tilde B)$ on $\mathrm{Fn}(\mathbb{Z}^3)$, and $B_{k,y}(\tilde B)$ is the operator (25.1).
By Lemma~2\,(a) it is a local family with the sets $\mathrm{Col}$ and $\mathrm{RdB}$, the group $\mathrm{Gv}$ and the largest distance of $P(\tilde B)$.

\emph{Level~0.}
For $x_1=(c_1,g_1)$ and $y'=[c_1]^L$,
\[
\left(JP(\tilde B)J^{*}f\right)(x_1)=U_{\tilde B}(y',c_1)^{-1}\cdot L^{-3}\sum_{c_2\in B(y')}U_{\tilde B}(y',c_2)\cdot\frac{1}{48}\sum_{g_2\in O_{\mathrm{cub}}}f(c_2,g_2).
\]
Hence $\left(JP(\tilde B)J^{*}\right)[x_1,x_2]=\frac{1}{48}\cdot P(\tilde B)[c_1,c_2]$ for $x_2=(c_2,g_2)$, and $\mathrm{Col}(x_1)=B(y')\times O_{\mathrm{cub}}$; the set $\mathrm{RdB}(x_1)$, the largest distance $L-1$ and $\mathrm{Gv}=L\mathbb{Z}^3$ are those of $P(\tilde B)$.

\paragraph{(Of7) [no.~7; $k\ge1$]}
One has $JP_k(\tilde B)J^{*}=\left(JQ_k^{T}(-\tilde B)\right)\left(Q_k(\tilde B)J^{*}\right)$.
By Lemma~2\,(c) and (Of4), with $y:=[c]$,
\[
\left(JP_k(\tilde B)J^{*}\right)[x,x']=\frac{\eta_k^3}{48}\cdot\mathbf{1}[c'\in\Delta_y]\cdot\varphi_{\tilde B}(y,c)^{-1}\cdot\varphi_{\tilde B}(y,c')\cdot\mathbf{1}_4,\qquad\mathrm{Col}(x)=\Delta_y\times O_{\mathrm{cub}},
\]
and $\mathrm{RdB}(x)$ is the union of the sets $\mathrm{Twr}(c')$, $c'\in\Delta_y$.
The sites $c$, $c'$ and the ends of these bonds lie in $\Delta_y$ (Fact~3), which is a lattice box of sup-diameter $1-\eta_k$.
The largest distance is $1-\eta_k$, and $\mathrm{Gv}=\mathbb{Z}^3$.

Next, $JP_{k+1}(\tilde B)J^{*}=\left(JQ_k^{T}(-\tilde B)\right)P(\tilde B)\left(Q_k(\tilde B)J^{*}\right)$.
By Lemma~2\,(c), (Of4) and (Of6), with $y':=[c]^L$ and $\varphi^{+}_{\tilde B}(y',c)=U_{\tilde B}(y',[c])\cdot\varphi_{\tilde B}([c],c)$ (see (I.3.3)),
\begin{gather*}
\left(JP_{k+1}(\tilde B)J^{*}\right)[x,x']=\frac{\eta_k^3}{48}\cdot L^{-3}\cdot\mathbf{1}\left[c'\in\mathrm{Bl}_{k+1}(y')\right]\cdot\varphi^{+}_{\tilde B}(y',c)^{-1}\cdot\varphi^{+}_{\tilde B}(y',c')\cdot\mathbf{1}_4,\\
\mathrm{Col}(x)=\mathrm{Bl}_{k+1}(y')\times O_{\mathrm{cub}},
\end{gather*}
and $\mathrm{RdB}(x)$ is the union of the sets $\mathrm{Twr}(c')$, $c'\in\mathrm{Bl}_{k+1}(y')$, and of the sets $\mathrm{Pth}(y',y)$, $y\in B(y')$.
By (1.b) and Facts~3 and~4 all these sites and ends of bonds lie in $\mathrm{Bl}_{k+1}(y')$, a lattice box of sup-diameter $L-\eta_k$ that contains $c$.
The largest distance is $L-\eta_k$, and $\mathrm{Gv}=L\mathbb{Z}^3$.

For $y\in Y_k$ let
\[
\operatorname{Pen}_{y}(\tilde B):=\boldsymbol{\alpha}_{y}\cdot JP_k(\tilde B)J^{*}+\boldsymbol{\beta}_{y}\cdot JP_{k+1}(\tilde B)J^{*},\qquad F_{y}(\tilde B):=D_{\tilde B}+\mu+\operatorname{Pen}_{y}(\tilde B),
\]
where $\boldsymbol{\alpha}_{y}$, $\boldsymbol{\beta}_{y}$ are the constant matrices of Section~I.3.5, which exist by the formula for $(a+i\cdot\gamma_5\cdot y)^{-1}$ above, and $\boldsymbol{\alpha}_{\infty}=b_k$, $\boldsymbol{\beta}_{\infty}=0$.
By Lemma~2\,(a) these are local families with $\mathrm{Col}(x)\subset\left(\mathrm{Bl}_{k+1}(y')\times O_{\mathrm{cub}}\right)\cup\{x^{(3)}\}$, and $\mathrm{RdB}(x)$ is the set just described, to which for $F_{y}(\tilde B)$ the bond $(c,c_3)$ is added.
Since $|c_3-c|=\eta_k\le L-\eta_k$, the largest distance is $L-\eta_k$, and $\mathrm{Gv}=L\mathbb{Z}^3$.

\paragraph{(Of8) [no.~8; $k=0$]}
For real $y$ let $F_{y}(\tilde B):=D_{\tilde B}+\mu+b_L\cdot JP(\tilde B)J^{*}+i\cdot\gamma_5\cdot y$ (Section~I.3.9).
By Lemma~2\,(a), (Of3) and (Of6), $\mathrm{Col}(x)=\left(B([c]^L)\times O_{\mathrm{cub}}\right)\cup\{x^{(3)}\}$, and $\mathrm{RdB}(x)$ is the union of $\{(c,c_3)\}$ and of the sets $\mathrm{Pth}(y',c')$, $c'\in B(y')$, where $y'=[c]^L$.
Since $|c_3-c|=1\le L-1$, the largest distance is $L-1$, and $\mathrm{Gv}=L\mathbb{Z}^3$.

\paragraph{(Of9) [no.~9, first half; $k\ge1$]}
By Section~I.3.5, $T_k(\tilde B):=b_k^{-1}\cdot Q_k(\tilde B)J^{*}\left(D_{\tilde B}+\mu+b_k\cdot JP_k(\tilde B)J^{*}\right)$.
From $J^{*}J=1$ and $Q_k(\tilde B)Q_k^{T}(-\tilde B)=1$ one obtains
\[
Q_k(\tilde B)J^{*}JP_k(\tilde B)J^{*}=Q_k(\tilde B)Q_k^{T}(-\tilde B)Q_k(\tilde B)J^{*}=Q_k(\tilde B)J^{*}.
\]
Hence $T_k(\tilde B)=b_k^{-1}\cdot Q_k(\tilde B)J^{*}\left(D_{\tilde B}+\mu\right)+Q_k(\tilde B)J^{*}$, and (Of3), (Of4) give
\begin{equation}
\begin{split}
\left(T_k(\tilde B)f\right)(y)=\int_{\Delta_y}\mathrm{d}x\,\varphi_{\tilde B}(y,c)\cdot&\left\{\left(1+\frac{\mu}{b_k}\right)f(x)\right.\\
&\left.\quad-(b_k\cdot\eta_k)^{-1}\sum_{i=0}^{3}M_i(x)\cdot\left[U_{\tilde B}(x,x^{(i)})\cdot f(x^{(i)})-f(x)\right]\right\},
\end{split}
\tag{25.2}
\end{equation}
which is the display (23.2) on the covering lattice.
By Lemma~2\,(a), (c),
\[
\mathrm{Col}(y)=\left(\Delta_y\times O_{\mathrm{cub}}\right)\cup\left\{x^{(3)}:c\in\Delta_y\right\},
\]
and $\mathrm{RdB}(y)$ is the union of the sets $\mathrm{Twr}(c)$, $c\in\Delta_y$, and of the bonds $(c,c_3(x))$, $x$ over $\Delta_y$.
The bonds of the first kind are contained in $\Delta_y$, and $|c_3-y|\le|c-y|+\eta_k\le\frac{1+\eta_k}{2}$.
The largest distance is $\frac{1+\eta_k}{2}$, and $\mathrm{Gv}=\mathbb{Z}^3$.
The operator $\bar{T}_k(\tilde B):=\mathsf{C}^{-1}T_k(-\tilde B)\mathsf{C}$ has the same sets, by Lemma~2\,(a), (b).

\paragraph{(Of10) [nos.~9, second half, and~10; $k\ge1$]}
Read on the covering lattice, the display (23.1) becomes, for $y\in\tilde Y\cap\mathbb{Z}^3$,
\begin{align*}
\left(T^{(1)}(\tilde B,\tilde Y)f\right)(y)&:=\int_{\Delta_y}\mathrm{d}x\,\varphi_{\tilde B}(y,c)\cdot\left\{\left(1+\frac{\mu}{b_k}\right)f(x)\right.\\
&\qquad-(b_k\cdot\eta_k)^{-1}\left[\sum_{i\le2}M_i(x)\cdot\left(f(x^{(i)})-f(x)\right)\right.\\
&\qquad\qquad\left.\left.+\mathbf{1}[c_3\in\tilde Y]\cdot M_3(x)\cdot\left(u_{\tilde B}(c,c_3)\cdot f(x^{(3)})-f(x)\right)\right]\right\},\\
\left(T^{(2)}(\tilde B,\tilde Y)F\right)(y)&:=-b_k^{-1}\cdot\eta_k^{\alpha-1}\int_{\Delta_y}\mathrm{d}x\,\mathbf{1}[c_3\notin\tilde Y]\cdot\varphi_{\tilde B}(y,c)\cdot M_3(x)\cdot F(x,x^{(3)}),
\end{align*}
while both rows vanish for $y\notin\tilde Y$; moreover $\bar{T}^{(i)}(\tilde B,\tilde Y):=\mathsf{C}^{-1}T^{(i)}(-\tilde B,\tilde Y)\mathsf{C}$.
The pair $(x,x^{(3)})$ lies in $\mathrm{Pzt}_k$, since $x^{(3)}\ne x$ and $d_X(x,x^{(3)})=\eta_k\le M$.
To obtain the coefficients, put $x'=x^{(i)}$ in the terms that contain $f(x^{(i)})$, that is, $x=x'^{(i)}$ by (I.s1); for $i\le2$ the site $x'^{(i)}$ lies over $c'$, and for $i=3$ the site $x'^{(3)}$ lies over $c_3(x')$, with $c_3(x'^{(3)})=c'$.
Writing $\int\mathrm{d}x=\frac{\eta_k^3}{48}\sum$, one obtains
\begin{align*}
T^{(1)}(\tilde B,\tilde Y)[y,x']&=\mathbf{1}[y\in\tilde Y]\cdot\frac{\eta_k^3}{48}\cdot\left\{\mathbf{1}[c'\in\Delta_y]\cdot\varphi_{\tilde B}(y,c')\cdot\left[\left(1+\frac{\mu}{b_k}\right)\mathbf{1}_4\right.\right.\\
&\qquad\left.+(b_k\cdot\eta_k)^{-1}\left(\sum_{i\le2}\left(M_i(x')-M_i(x'^{(i)})\right)+\mathbf{1}[c_3(x')\in\tilde Y]\cdot M_3(x')\right)\right]\\
&\qquad-(b_k\cdot\eta_k)^{-1}\cdot\mathbf{1}[c_3(x')\in\Delta_y]\cdot\mathbf{1}[c'\in\tilde Y]\\
&\qquad\qquad\left.\cdot\,\varphi_{\tilde B}\left(y,c_3(x')\right)\cdot M_3(x'^{(3)})\cdot u_{\tilde B}\left(c_3(x'),c'\right)\right\},\\
T^{(2)}(\tilde B,\tilde Y)[y,(x,x')]&=-\mathbf{1}[y\in\tilde Y]\cdot\mathbf{1}[x'=x^{(3)}]\cdot\mathbf{1}[c\in\Delta_y]\cdot\mathbf{1}[c_3\notin\tilde Y]\\
&\qquad\cdot\,b_k^{-1}\cdot\eta_k^{\alpha-1}\cdot\frac{\eta_k^3}{48}\cdot\varphi_{\tilde B}(y,c)\cdot M_3(x).
\end{align*}
For $T^{(1)}$ the sets $\mathrm{Col}(y)$ and $\mathrm{RdB}(y)$ are those of $T_k(\tilde B)$.
For $T^{(2)}$ one has $\mathrm{Col}(y)=\left\{(x,x^{(3)}):c\in\Delta_y\right\}$, and $\mathrm{RdB}(y)$ is the union of the sets $\mathrm{Twr}(c)$, $c\in\Delta_y$, whose bonds are contained in $\Delta_y$.
For both operators $\mathrm{RdY}(y)=\{y\}\cup\left\{c_3(x):x\text{ over }\Delta_y\right\}$.
One has $\mathrm{dist}\left(\{y\},\{c,c_3\}\right)\le|y-c|\le\frac{1-\eta_k}{2}$, and the two points of a pair, the ends of the bonds and the sites of $\mathrm{RdY}(y)$ have distance at most $\frac{1+\eta_k}{2}$ from $y$.
The largest distance is $\frac{1+\eta_k}{2}$.
Further $\mathrm{Gv}=M\mathbb{Z}^3$: for $v\in M\mathbb{Z}^3$ one has $\mathbf{1}[y+v\in\tilde Y+v]=\mathbf{1}[y\in\tilde Y]$ and $\mathbf{1}[c_3(x+v)\in\tilde Y+v]=\mathbf{1}[c_3(x)\in\tilde Y]$, and (5.a)--(5.c) apply.
For $\bar{T}^{(i)}$ one uses Lemma~2\,(a), (b).

\paragraph{(Of11) [no.~11]}
For $(x,x')\in\mathrm{Pzt}_k$ let
\[
(\delta_{\alpha,\tilde B}f)(x,x'):=d_X(x,x')^{-\alpha}\cdot\left(U_{\tilde B}(x,x')\cdot f(x')-f(x)\right)
\]
as in (I.1.7); here $|c-c'|\le M$, so that $U_{\tilde B}(x,x')=U_{\tilde B}(c,c')$ is defined by (I.1.6), and $U_{\tilde B}=1$ if $c=c'$.
Thus
\begin{gather*}
\delta_{\alpha,\tilde B}[(x,x'),x'']=d_X(x,x')^{-\alpha}\cdot\left(\mathbf{1}[x''=x']\cdot U_{\tilde B}(x,x')-\mathbf{1}[x''=x]\right)\cdot\mathbf{1}_4,\\
\mathrm{Col}(x,x')=\{x,x'\},\qquad\mathrm{RdB}(x,x')=\mathrm{Pth}(c,c').
\end{gather*}
Since $\operatorname{supp}(x,x')=\{c,c'\}$, the columns have distance $0$.
An end $p$ of a bond of $\mathrm{Pth}(c,c')$ lies in $\mathrm{hul}(c,c')$, so that $|p-c|\le|c-c'|\le M$ by Fact~2.
The largest distance is at most $M$, and $\mathrm{Gv}=\eta_k\mathbb{Z}^3$ by (5.a), (5.b).
The definitions of $\mathrm{Pzt}_k$ and of $d_X$ do not contain $\ell$ ((I.1.5), Section~I.1.7).

\paragraph{(Of12) [no.~12]}
On $\mathbb{Z}^3$ consider the $L$-blocks $B(y)$, $y\in L\mathbb{Z}^3$ \cite[Section~5.2.5]{QEDI}.
A unit bond is called interior if both its ends lie in one block, and it is called a face bond of the face $\mathrm{fc}=\langle y,y+Le_{\mu}\rangle$ if its ends lie in $B(y)$ and in $B(y+Le_{\mu})$; the central bond of this face is
\[
b_{\mathrm{fc}}:=\left(y+\frac{L-1}{2}e_{\mu},\ y+\frac{L+1}{2}e_{\mu}\right).
\]
The covering family of the basis elements is
\begin{align*}
\mathrm{Upc}:={}&\left\{e^{y}_{i}:y\in L\mathbb{Z}^3,\ i\text{ in the index set of the basis }\{e^{0}_{i}\}\right\}\\
&\cup\left\{\delta_b:b\text{ a positively oriented face bond of }\mathbb{Z}^3\text{ that is not central}\right\},\qquad e^{y}_{i}(b):=e^{0}_{i}(b-y),
\end{align*}
where $\{e^{0}_{i}\}$ is the finite basis of \cite[Section~5.2.5]{QEDI}; it consists of bond functions on the interior bonds of $B(0)$ and is the same for all $k$ and $\ell$.
The elements of $\mathrm{Upc}$ are bond functions on $\mathbb{Z}^3$ with finite supports.
A vector $v\in L\mathbb{Z}^3$ acts by $e^{y}_{i}+v:=e^{y+v}_{i}$ and $\delta_b+v:=\delta_{b+v}$; as bond functions, $\Upsilon+v=T_v\Upsilon$.
The support of $e^{y}_{i}$ lies in $B(y)$, and that of $\delta_b$ consists of the two ends of $b$ \cite[Proposition~5.35 and Proposition~15.21(ii)]{QEDI}.
For $\tilde z\in\mathrm{Fn}(\mathrm{Upc})$ the sum $\mathsf{Z}:=\sum_{\Upsilon}\tilde z(\Upsilon)\cdot\Upsilon$ has only finitely many non-zero terms at every bond: an interior bond of $B(y)$ lies only in the supports of the $e^{y}_{i}$, a non-central face bond $b$ only in that of $\delta_b$, and a central bond in no support.
With all bonds oriented along $+e_{\mu}$, the display of $C$ \cite[Section~5.2.5 and Proposition~15.21(iv)]{QEDI} is
\begin{gather*}
(C\tilde z)(b'):=\mathsf{Z}(b')\qquad\text{for }b'\text{ not central},\\
(C\tilde z)(b_{\mathrm{fc}}):=-\sum\left\{\mathsf{Z}(b):b\text{ a non-central face bond of }\mathrm{fc}\right\}-\sum_b\frac{n_{\mathrm{fc}}(b)}{L}\cdot\mathsf{Z}(b),
\end{gather*}
where the second sum extends over the interior bonds $b$ of direction $\mu$ of $B(y)$ and of $B(y+Le_{\mu})$, and $n_{\mathrm{fc}}(b)$ is the number of the straight paths $\boldsymbol{\Gamma}_{p,p+Le_{\mu}}$, $p\in B(y)$, that contain $b$.
Let $C(b',\Upsilon)$ be the coefficient of $\tilde z(\Upsilon)$ in $(C\tilde z)(b')$; the row index $b'$ is a positively oriented unit bond (see 1.4), and at the reversed bond the row is the negative.
The coefficients are as follows:
\begin{itemize}
\item if $b'$ is an interior bond of $B(y)$, then $C(b',e^{y}_{i})=e^{y}_{i}(b')$;
\item if $b'$ is a positively oriented non-central face bond, then $C(b',\delta_{b'})=1$;
\item if $b'=b_{\mathrm{fc}}$, then $C(b_{\mathrm{fc}},\delta_b)=-1$ for the non-central face bonds $b$ of $\mathrm{fc}$, and
\[
C(b_{\mathrm{fc}},e^{y''}_{i})=-\sum_b\frac{n_{\mathrm{fc}}(b)}{L}\cdot e^{y''}_{i}(b)\qquad\text{for }y''\in\{y,y+Le_{\mu}\},
\]
the sum extending over the interior bonds $b$ of direction $\mu$ of $B(y'')$;
\end{itemize}
all other coefficients vanish.
Let $\mathrm{Col}(b')$ be the set of those $\Upsilon$ whose support lies in the union of the $L$-blocks of the two ends of $b'$; it is finite and contains all the $\Upsilon$ named above.
Neither a connection nor $\tilde Y$ occurs.
As to the distance, let $\Upsilon\in\mathrm{Col}(b')$; a point $p'$ of $\operatorname{supp}\Upsilon$ lies in the $L$-block of an end $p$ of $b'$, so that $|p-p'|\le L-1$ by (1.b), and $\mathrm{dist}(\operatorname{supp}b',\operatorname{supp}\Upsilon)\le L-1$.
(All points of $\operatorname{supp}b'$ and of $\operatorname{supp}\Upsilon$ lie in one $L$-block, or in two $L$-blocks with a common face.)
Finally $\mathrm{Gv}=L\mathbb{Z}^3$: a translation by $v\in L\mathbb{Z}^3$ maps blocks, faces, central bonds and the paths $\boldsymbol{\Gamma}_{p,p+Le_{\mu}}$ to those of the translated face, $n_{\mathrm{fc}+v}(b+v)=n_{\mathrm{fc}}(b)$, and $e^{y+v}_{i}(b+v)=e^{y}_{i}(b)$; hence $C(b'+v,\Upsilon+v)=C(b',\Upsilon)$.

\paragraph{(Of13) [no.~13]}
Let $(\mathbf{1}_{\tilde Y}f)(x):=\mathbf{1}[c\in\tilde Y]\cdot f(x)$ (Section~I.1.4).
Then $\mathrm{Col}(x)=\{x\}$ and $\mathrm{RdY}(x)=\{c\}$; the distance is $0$, and $\mathrm{Gv}=M\mathbb{Z}^3$.

\paragraph{(Of14) [no.~14]}
The map $c\mapsto c/L$ is a bijection $\mathbb{L}_k\to\mathbb{L}_{k+1}$, because $\eta_{k+1}=\eta_k/L$; it maps bonds to bonds, positively oriented bonds to positively oriented bonds, and $L\mathbb{Z}^3$ onto $\mathbb{Z}^3$.
The map $x\mapsto x/L=(c/L,g)$ is a bijection $\mathbb{L}^{\mathrm{f}}_k\to\mathbb{L}^{\mathrm{f}}_{k+1}$.
For $(x,x')\in\mathrm{Pzt}_k$ the pair $(x/L,x'/L)$ consists of distinct points, and at the level $k+1$
\[
d_X(x/L,x'/L)=\max\left(\frac{|c-c'|}{L},\frac{\eta_k}{L}\right)=\frac{d_X(x,x')}{L}\le\frac{M}{L}\le M;
\]
so it lies in $\mathrm{Pzt}_{k+1}$.
The displays of Section~I.3.7 are
\begin{gather*}
A'_L(b):=L^{-\frac{1}{2}}\cdot A'(b/L),\qquad f'_L(x):=L^{-1}\cdot f'(x/L),\\
F'_L(x,x'):=L^{-1-\alpha}\cdot F'(x/L,x'/L),\qquad v'_L(y'):=L^{-1}\cdot v'(y'/L)\quad(y'\in L\mathbb{Z}^3),
\end{gather*}
where $A'$ is a bond function on $\mathbb{L}_{k+1}$.
For a function $\Lambda'$ of the sites of $\mathbb{L}_{k+1}$, Section~I.3.7 has the display $\Lambda'_L(c):=L^{\frac{1}{2}}\cdot\Lambda'(c/L)$; this map is treated together with the other four.
The row at $t$ has the single column $t/L$, with the coefficient $L^{-\frac{1}{2}}$, $L^{-1}$, $L^{-1-\alpha}$, $L^{-1}$, respectively, and $L^{\frac{1}{2}}$ for $\Lambda'\mapsto\Lambda'_L$.
For $A'\mapsto A'_L$ the row index is a positively oriented bond $b$ (see 1.4), and the column $b/L$ is positively oriented.
The notions of 1.4 are used with the action $t'\mapsto t'+v/L$ of $v$ on the index sets of the level $k+1$ and with $L\operatorname{supp}t'$ in the place of $\operatorname{supp}t'$; then $\mathrm{Col}(t)=\{t/L\}$, $L\operatorname{supp}(t/L)=\operatorname{supp}t$, and the distance is $0$.
Nothing is read: $\mathrm{RdB}(t)$ and $\mathrm{RdY}(t)$ are empty.
One has $\mathrm{Gv}=\eta_k\mathbb{Z}^3$ for the first three maps and for $\Lambda'\mapsto\Lambda'_L$, and $\mathrm{Gv}=L\mathbb{Z}^3$ for the fourth:
\begin{align*}
(T_vA'_L)(b)&=L^{-\frac{1}{2}}\cdot A'\left(\frac{b-v}{L}\right)=L^{-\frac{1}{2}}\cdot(T_{v/L}A')(b/L)=\left((T_{v/L}A')_L\right)(b),\\
(T_vf'_L)(x)&=L^{-1}\cdot f'\left(\frac{x-v}{L}\right)=L^{-1}\cdot(T_{v/L}f')(x/L)=\left((T_{v/L}f')_L\right)(x),\\
(T_vF'_L)(x,x')&=L^{-1-\alpha}\cdot F'\left(\frac{x-v}{L},\frac{x'-v}{L}\right)=\left((T_{v/L}F')_L\right)(x,x'),\\
(T_vv'_L)(y')&=L^{-1}\cdot v'\left(\frac{y'-v}{L}\right)=\left((T_{v/L}v')_L\right)(y')\qquad(v\in L\mathbb{Z}^3,\text{ so that }v/L\in\mathbb{Z}^3),\\
(T_v\Lambda'_L)(c)&=L^{\frac{1}{2}}\cdot\Lambda'\left(\frac{c-v}{L}\right)=L^{\frac{1}{2}}\cdot(T_{v/L}\Lambda')(c/L)=\left((T_{v/L}\Lambda')_L\right)(c).
\end{align*}

\subsubsection*{2.3. Proof, second part: (f1), (f2), (f4), the range $\mathrm{rg}$ and what the row at the index $t$ reads, the entire functions, the letters}

By 2.2, each operator of the table is a local family in the sense of 1.4.
The results of 2.2 are collected in the following table.

\noindent
\begin{tabular}{|p{0.05\textwidth}|p{0.12\textwidth}|p{0.12\textwidth}|p{0.21\textwidth}|p{0.29\textwidth}|}
\hline
no. & largest distance found in 2.2 & $\mathrm{rg}$ in the table of Proposition~25.1 & $\mathrm{Gv}$ & the sets $\mathrm{RdB}(t)$, $\mathrm{RdY}(t)$ of 2.2\\
\hline
1, 2 & $0$ & $0$ & $\eta_k\mathbb{Z}^3$ (for a constant matrix: the group that acts on the index set, $\eta_k\mathbb{Z}^3$, $\mathbb{Z}^3$ or $L\mathbb{Z}^3$; it contains $L\mathbb{Z}^3$) & empty\\
\hline
3 & $\eta_k$ & $\eta_k$ & $\eta_k\mathbb{Z}^3$ & the bond $(c,c_3)$\\
\hline
4 & $\frac{1-\eta_k}{2}$ & $\frac{1-\eta_k}{2}$ & $\mathbb{Z}^3$ & row at $y$: the $\mathrm{Twr}(c)$, $c\in\Delta_y$; row at $x$: $\mathrm{Twr}(c)$\\
\hline
5 & $\frac{L-1}{2}$ & $\frac{L-1}{2}$ & $L\mathbb{Z}^3$ & row at $y'$: the $\mathrm{Pth}(y',y)$, $y\in B(y')$; row at $y$ (for $k=0$: at $x=(y,g)$): $\mathrm{Pth}([y]^L,y)$\\
\hline
6 & $L-1$ & $L-1$ & $L\mathbb{Z}^3$ & the $\mathrm{Pth}(y',y)$, $y\in B(y')$, where $y'$ is the centre of the $L$-block of the row\\
\hline
7 & $1-\eta_k$; $L-\eta_k$ & $1-\eta_k$; $L-\eta_k$ & $\mathbb{Z}^3$; $L\mathbb{Z}^3$ & the $\mathrm{Twr}(c')$, $c'\in\Delta_{[c]}$; the $\mathrm{Twr}(c')$, $c'\in\mathrm{Bl}_{k+1}(y')$, and the $\mathrm{Pth}(y',y)$, $y\in B(y')$, where $y'=[c]^L$; for $F_y$ also $(c,c_3)$\\
\hline
8 & $L-1$ & $L-1$ & $L\mathbb{Z}^3$ & $(c,c_3)$ and the set of no.~6\\
\hline
9 & $\frac{1+\eta_k}{2}$ & $\frac{1+\eta_k}{2}$ & $\mathbb{Z}^3$; $M\mathbb{Z}^3$ & the $\mathrm{Twr}(c)$ and the bonds $(c,c_3(x))$, $c\in\Delta_y$; $\mathrm{RdY}(y)=\{y\}\cup\{c_3(x):c\in\Delta_y\}$\\
\hline
10 & $\frac{1+\eta_k}{2}$ & $\frac{1+\eta_k}{2}$ & $M\mathbb{Z}^3$ & the $\mathrm{Twr}(c)$, $c\in\Delta_y$; $\mathrm{RdY}(y)$ as in no.~9\\
\hline
11 & $\le M$ & $M$ & $\eta_k\mathbb{Z}^3$ & $\mathrm{Pth}(c,c')$\\
\hline
12 & $\le L-1$ & $L-1$ & $L\mathbb{Z}^3$ & empty\\
\hline
13 & $0$ & $0$ & $M\mathbb{Z}^3$ & $\mathrm{RdY}(x)=\{c\}$\\
\hline
14 & $0$ & $0$ & $\eta_k\mathbb{Z}^3$; $L\mathbb{Z}^3$ & empty\\
\hline
\end{tabular}

\medskip
\emph{(f1) and (f4).}
By the paragraph of 1.4 on the relation with Section~I.1.9, (f1) holds, and (f4) holds with $\mathrm{rg}$ equal to the number given in the table of Proposition~25.1; in every row this number is at least the largest distance.

\emph{(f2).}
The relation (f2) holds for every $v\in\mathrm{Gv}$.
For the operators that do not depend on $\tilde{Y}$ the group $\mathrm{Gv}$ is $\eta_k\mathbb{Z}^3$, $\mathbb{Z}^3$ or $L\mathbb{Z}^3$, and each of these groups contains $L\mathbb{Z}^3$.
For the operators $T^{(1)}$, $\bar{T}^{(1)}$ of no.~9 and for the nos.~10 and 13 one has $\mathrm{Gv}=M\mathbb{Z}^3$.
Both $L\mathbb{Z}^3$ and $M\mathbb{Z}^3$ contain $LM\mathbb{Z}^3$; hence (f2) of Section~I.1.9 holds.

\emph{What the row at the index $t$ reads.}
By (l2), the row at $t$ reads $\tilde{B}$ on $\mathrm{RdB}(t)$ and $\tilde{Y}$ through $\mathrm{RdY}(t)$.
For each row of the table these sets are compared with the statement of Proposition~25.1 on what the row at the index $t$ reads.
For the nos.~1, 2, 12 and 14 nothing is read.
For no.~3 the row reads the single bond $(c,c_3)$.
For no.~4, $\mathrm{Twr}(c)$ is the set of the bonds of the paths $\boldsymbol{\Gamma}^{\pi}(y_{j+1},y_j)$, $j<k$, of the tower of $c$; by Fact~3 these bonds are contained in $\Delta_y$ and lie within the distance $(1-\eta_k)/2$ of the row index, $y$ or $c$ respectively.
For no.~5 the statement follows from Fact~4, and for no.~6 from (Of6).
For no.~7 it follows from (Of7), since the bonds of $\mathrm{Twr}(c')$ are contained in $\Delta_{[c']}$, which lies in $\Delta_{[c]}$, resp.\ in $\mathrm{Bl}_{k+1}(y')$.
For no.~8 the statement follows from (Of8), and for the nos.~9 and 10 from (Of9) and (Of10).
For no.~11 it follows from (Of11): $\mathrm{Pth}(c,c')$ is empty if $c=c'$, and the columns of the row are $x$ and $x'$.
For no.~13 it follows from (Of13).
The assertion that $M_i(x)$ depends on the label $g$ only is contained in (Of1), and the assertion that $C(b',\Upsilon)=0$ unless $\operatorname{supp}\Upsilon$ lies in the union of the $L$-blocks of the ends of $b'$ is contained in (Of12); the assertions that $T_v$ at level $k$ goes with $T_{v/L}$ at level $k+1$ and that $\operatorname{supp}(t/L)=\operatorname{supp}(t)/L$ are contained in (Of14).
The last two clauses of the entry of no.~14 are conventions; the second of them is the convention of Section~G3.5.0.

\emph{Entire functions.}
That the entries of every row are entire functions of the numbers $\tilde{B}(b)$ which the row reads is (l3).

\emph{The letters.}
The coefficients given in 2.2 and the sets $\Delta_y$, $B(y')$, $\mathrm{Bl}_j(z)$, $\mathrm{Twr}(c)$, $\mathrm{Pth}(c,c')$, $\mathrm{Pzt}_k$, $\mathrm{Upc}$ are formed with $k$, $\eta_k=L^{-k}$, $\mu$, $e_k$, $b_k$, $b_L$, $\alpha$, $\rho$, the weights $w$ (through $M_i$ and $\mathrm{sw}$), $L$, $M$ (through $\mathrm{Pzt}_k$ and the $M$-cubes of $\tilde{Y}$), the matrices of Section~I.2.1, the basis $\{e^{0}_{i}\}$ of \cite[Section~5.2.5]{QEDI} (through $C$ and $\mathrm{Upc}$; it is chosen once for each $L$ and is the same for all $k$ and $\ell$, see (Of12)) and the resolvent parameter $y$.
None of them contains the side $\ell$.

\subsubsection*{2.4. Proof, third part: (f3)}

Let $\ell\in\mathrm{Cl}$, let $B$ be a complex bond function on $T_k$ and $Y$ a union of $M$-cubes of $T_k$, and let $\tilde{B}$, $\tilde{Y}$ be their $\ell\mathbb{Z}^3$-periodic extensions.
In every row of the table the group $\mathrm{Gv}$ contains $\ell\mathbb{Z}^3$.
By Lemma~3\,(a), therefore, $\mathrm{Op}_{\mathbb{L}}(\tilde{B},\tilde{Y})$ maps periodic functions to periodic functions and antiperiodic functions to antiperiodic functions.
It remains to show that its restriction is the operator $\mathrm{Op}(B,Y)$ defined on the torus.

\emph{(i) The operators (Of1)--(Of5), (Of10), (Of11), (Of13), (Of14).}
By Sections~I.1.5, I.1.7 and I.1.9, the spaces $\mathfrak{S}(a,\ell)$, $\mathfrak{S}^{\mathrm{ch}}(a,\ell)$ and the space of the functions of pairs are, by definition, the spaces of the antiperiodic functions on $a\mathbb{Z}^3$, on $a\mathbb{Z}^3\times O_{\mathrm{cub}}$ and on $\mathrm{Pzt}_k$; a bond function on $T_k$ and a union of $M$-cubes of $T_k$ are identified with their periodic extensions.
The displays (I.2.1)--(I.2.3), (I.3.1), the displays of Section~I.3.3, (I.3.4), (I.1.7), the display of Section~I.1.4, the displays of Section~I.3.7 and (23.1) are written in terms of the values of $f$ at points of these covering index sets, of the sets $\Delta_y$ and $B(y')$, which Section~I.1.1 defines as subsets of $\mathbb{L}_k$ and of $\mathbb{Z}^3$, of the neighbours (I.1.2), which are defined on $\mathbb{L}^{\mathrm{f}}_k$, of the towers of Section~I.3.2 and of the paths of Section~I.1.1.
On the torus, Sections~I.1.1 and I.3.2 form a path on a lift with $|c-c'|<\ell/2$.
The pairs of end points that occur are $(y_{j+1}(c),y_j(c))$, $(y',y)$ with $y\in B(y')$, and $(c,c')$ with $|c-c'|\le M$; their distance is at most $M<\ell/2$.
Hence the path is the path $\boldsymbol{\Gamma}^{\pi}$ of $\mathbb{L}_k$ between the two points named, and $B$ is summed over its image, that is, $\tilde{B}$ is summed over the path itself.
The conditions $y\in Y$, $c\in Y$, $c_3\in Y$ mean membership in $\tilde{Y}$.
Consequently, for $f$ in the space of the torus, the display that defines $(\mathrm{Op}(B,Y)f)(t)$ and the display of 2.2 for $(\mathrm{Op}_{\mathbb{L}}(\tilde{B},\tilde{Y})f)(t)$ are the same expression.
For (Of14) one has in addition: if $f'$ is antiperiodic under $(\ell/L)\mathbb{Z}^3$, then $f'_L$ is antiperiodic under $\ell\mathbb{Z}^3$, because $(t+\ell w)/L=t/L+(\ell/L)w$; the same holds for $F'$ and $v'$, and, with ``periodic'' in the place of ``antiperiodic'', for $A'$ and $\Lambda'$.
The level $k+1$ has the side $\ell/L$ (Section~G3.5.0).

\emph{(ii) Sums and products.}
These are $P(B)$, $JP(B)J^{*}$, $JP_k(B)J^{*}$, $JP_{k+1}(B)J^{*}$, $\operatorname{Pen}_{y}(B)$, $F_{y}(B)$, $T_k(B)$, $\bar{T}_k(B)$ and $\bar{T}^{(i)}(B,Y)$.
Sections~I.3.3, I.3.5 and I.3.9, and, for $\bar{T}^{(i)}(B,Y)$, the display (23.1), which is (I.5.4), define them as sums and products of operators of (i) and of constant matrices.
Each factor maps the spaces of the torus into one another (Lemma~3\,(a)); hence the restriction of a sum or product formed on the covering lattice is the sum or product of the restrictions.

\emph{(iii) The operator $B_{k,y}(B)$.}
By (ii), the restriction of $b_k+b_L\cdot P(\tilde{B})+i\gamma_5y$ to $\mathfrak{S}(1,\ell)$ is $b_k+b_L\cdot P(B)+i\gamma_5y$, and the restriction of (25.1) is a two-sided inverse of it on $\mathfrak{S}(1,\ell)$.
Thus $B_{k,y}(B)$ exists for every complex $B$ and every real $y$, and it is the restriction of $B_{k,y}(\tilde{B})$.

\emph{(iv) The map $C$.}
The side $\ell$ is a multiple of $L$.
The $L$-blocks, the faces, and the interior, face and central bonds of $T^0_k$ are the images of those of $\mathbb{Z}^3$, and the basis $\{\Upsilon\}$ of $V_{k}$ is the image of $\mathrm{Upc}$: the element $e^{y}_{i}$ is the translate of $e^{0}_{i}$ by $y\in L\mathbb{Z}^3/\ell\mathbb{Z}^3$ (\cite[Section~5.2.5]{QEDI}, \cite[Proposition~15.21(ii)]{QEDI}), and the periodic extension of the basis element with the index $\mathrm{pr}_{\ell}(\Upsilon)$ is $\sum_{w}T_{\ell w}\Upsilon$.
A function on the basis of the torus is an $\ell\mathbb{Z}^3$-periodic function $\tilde{z}\in\mathrm{Fn}(\mathrm{Upc})$, and the periodic extension of the element $\sum\mathsf{Z}(\Upsilon)\,\Upsilon$ of $V_{k}$ with these coordinates is $\sum_{\Upsilon\in\mathrm{Upc}}\tilde{z}(\Upsilon)\,\Upsilon$.
The union of two $L$-blocks with a common face has the sup-diameter $2L-1$, and $2L-1+1<\ell$; hence $\mathrm{pr}_{\ell}$ maps this union, its bonds and the lines $\Gamma_{p,p+Le_{\mu}}$ contained in it bijectively onto their images (Lemma~1\,(a); Section~I.1.1 with $\mathrm{sd}=1$).
Therefore the two sums in the formula for $C$ of \cite[Proposition~15.21(iv)]{QEDI}, formed on the torus, are the two sums of (Of12), and away from the central bonds both operators are the identity.

\emph{Existence on the torus.}
No display of the table contains an inverse other than $B_{k,y}(B)$ and the constant matrices $\boldsymbol{\alpha}_{y}$, $\boldsymbol{\beta}_{y}$; so each $\mathrm{Op}$ is defined for every $\ell\in\mathrm{Cl}$, every complex $B$ and every union $Y$ of $M$-cubes.
This is the first assertion of Proposition~25.1.
Together with 2.3 it follows that $\mathrm{Op}$ is given by one formula on $\mathbb{L}^{\mathrm{f}}_k$, of range $\mathrm{rg}$.

\emph{Uniformity.}
All assertions of 2.2--2.4 are identities and inclusions between finite sets.
They hold for every integer $k\ge0$ (at the levels named in the table), every complex $\mu$, every $e_k>0$, every admitted value of the resolvent parameter, all $\tilde{B}$ and $\tilde{Y}$, and every $\ell\in\mathrm{Cl}$; no estimate enters.

\subsubsection*{2.5. Proof, fourth part: the kernels (K1)--(K4)}

By 1.4, the kernel is the coefficient divided by the weight of the column; this weight is $\eta_k^3/48$ on $\mathbb{L}^{\mathrm{f}}_k$ and $1$ on $\mathbb{Z}^3$.
\begin{itemize}
\item[(K1)]
By (Of4),
\begin{gather*}
\left(Q_k(\tilde{B})J^{*}\right)(y,x)=\frac{48}{\eta_k^3}\cdot\frac{\eta_k^3}{48}\cdot\mathbf{1}[c\in\Delta_y]\cdot\varphi_{\tilde{B}}(y,c)\cdot1_4,\\
\left(JQ_k^{T}(-\tilde{B})\right)(x,y)=\mathbf{1}\left[y=[c]\right]\cdot\varphi_{\tilde{B}}([c],c)^{-1}\cdot1_4,
\end{gather*}
and $y=[c]$ means $c\in\Delta_y$.
\item[(K2)]
By (Of6), $P(\tilde{B})(y_1,y_2)=P(\tilde{B})[y_1,y_2]$.
At $k=0$, where $\eta_0=1$,
\[
\left(JP(\tilde{B})J^{*}\right)(x_1,x_2)=48\cdot\left(JP(\tilde{B})J^{*}\right)[x_1,x_2]=48\cdot\frac{1}{48}\cdot P(\tilde{B})[c_1,c_2]=P(\tilde{B})(c_1,c_2).
\]
\item[(K3)]
The first formula is (25.1).
Put $P:=P(\tilde{B})$, $m_1:=(b_k+i\gamma_5y)^{-1}$ and $m_2:=(b_k+b_L+i\gamma_5y)^{-1}$.
The operators $m_1$, $m_2$, $\gamma_5$ and $P$ commute with one another, $P^2=P$ and $P\cdot(1-P)=0$.
Hence
\[
B_{k,y}(\tilde{B})\,\gamma_5\,B_{k,y}(\tilde{B})=\gamma_5\cdot\left[m_1\cdot(1-P)+m_2\cdot P\right]^2=\gamma_5\cdot\left[m_1^2\cdot(1-P)+m_2^2\cdot P\right].
\]
\item[(K4)]
Here $k=0$, $\eta_0=1$, and the weight is $1/48$.
By (Of3), (Of8) and (K2),
\begin{align*}
F_0(\tilde{B})(x,x')&=48\cdot\left(D_{\tilde{B}}+\mu\right)[x,x']+b_L\cdot\left(JP(\tilde{B})J^{*}\right)(x,x')\\
&=48\cdot\left\{\mathbf{1}[x'=x]\cdot\left(\frac{\rho}{2}\cdot\mathrm{sw}+\mu\right)\cdot1_4-\sum_{i=0}^{3}\mathbf{1}\left[x'=x^{(i)}\right]\cdot M_{i}(x)\cdot U_{\tilde{B}}\left(x,x^{(i)}\right)\right\}\\
&\qquad+b_L\cdot P(\tilde{B})(c,c').
\end{align*}
\end{itemize}

\subsubsection*{2.6. Level 0, collected}

At $k=0$ one has $\eta_0=1$, $\mathbb{L}_0=\mathbb{Z}^3$ and $\Delta_y=\{y\}$, and $\mathrm{Twr}(c)$ is empty; the nos.~4, 7, 9 and 10 do not occur.
The rows are the following.
\begin{itemize}
\item
The row of $D_{\tilde{B}}$ at $x$ has the five columns $x$, $x^{(0)}$, $x^{(1)}$, $x^{(2)}$ (over $c$) and $x^{(3)}$ (over $c+ge_1$), with the coefficients $\frac{\rho}{2}\cdot\mathrm{sw}\cdot1_4$, $-M_{i}(x)$ ($i\le2$) and $-M_{3}(x)\cdot\exp\left(i\,e_k\cdot\tilde{B}(c,c_3)\right)$.
\item
The row of $Q(\tilde{B})J^{*}$ at $y'$ has the $48L^3$ columns over $B(y')$, with the coefficients $\frac{L^{-3}}{48}\cdot U_{\tilde{B}}(y',c')\cdot1_4$.
\item
The row of $JQ^{T}(-\tilde{B})$ at $x$ has the single column $[c]^L$.
\item
The row of $JP(\tilde{B})J^{*}$ at $x$ has the $48L^3$ columns over $B([c]^L)$, with the coefficients $\frac{L^{-3}}{48}\cdot U_{\tilde{B}}(y',c)^{-1}\cdot U_{\tilde{B}}(y',c')\cdot1_4$, where $y'=[c]^L$.
\item
The row of $F_{y}(\tilde{B})$ at $x$ is given, with $y'=[c]^L$, by
\begin{align*}
F_{y}(\tilde{B})[x,x']&=D_{\tilde{B}}[x,x']+b_L\cdot\frac{L^{-3}}{48}\cdot\mathbf{1}[c'\in B(y')]\cdot U_{\tilde{B}}(y',c)^{-1}\cdot U_{\tilde{B}}(y',c')\cdot1_4\\
&\qquad+\mathbf{1}[x'=x]\cdot\left(\mu+i\gamma_5y\right).
\end{align*}
\item
The row of $\delta_{\alpha,\tilde{B}}$ at $(x,x')$ has the columns $x$ and $x'$, with $d_X(x,x')=\max(|c-c'|,1)$; for a pair inside a cluster ($c=c'$, $g\ne g'$) the row is $f(x')-f(x)$, without phase.
\end{itemize}
In the kernels (I.1.4) the weight $1/48$ of a column appears as the factor $48$ of (K4), and it cancels in (K2).

\emph{Operators of Section~I.3.9 that are not rows of the table.}
These are
\begin{gather*}
D_0(\tilde{B}):=D_{\tilde{B}}+\mu,
\qquad
\operatorname{Pen}_{y}(\tilde{B}):=b_L\cdot JP(\tilde{B})J^{*}+i\gamma_5y\quad(y\ \text{real}),\\
H_{0}=\bar{H}_0=T_0=\bar{T}_0:=1\ \text{on}\ \mathrm{Fn}(\mathbb{L}^{\mathrm{f}}_0).
\end{gather*}
They are linear combinations of the nos.~1, 3 and 6 ((Of1), (Of3), (Of6)).
By Lemma~2\,(a) they are local families, with the largest distances $1$, $L-1$, $0$, the groups $\mathrm{Gv}=\mathbb{Z}^3$, $L\mathbb{Z}^3$, $\mathbb{Z}^3$, and the sets $\mathrm{RdB}(x)=\{(c,c_3)\}$, the set of no.~6, the empty set, respectively; by 2.4\,(ii), (f3) holds for them.

\subsubsection*{2.7. Proof, fifth part: the last paragraph of Proposition~25.1}

None of the operators listed in the last paragraph of Proposition~25.1 is an entry of the table.
Nor does any of them enter Definition~24.D, the definitions of Section~I.5.6 or the definition of $\acute{E}_k$, for the following reasons.
\begin{itemize}
\item[(1)]
In Definition~24.D the operators $H_{k}$, $C_{k}^{1/2}$, $\Gamma_k$, $M_k$, $\bar{H}_{k}$ enter only through $H_{k}(s')$, $C_{k}^{1/2}(s')$, $H_{k}(s'',B)$, $\bar{H}_{k}(s'',B)$, $\Gamma_k(s'',B)$ and $M_k(X,\cdot)$, each of which has a parameter family or a polymer as first argument; these are members of the families (p1)--(p4) of the system of Section~I.5.8.
The clauses (D0)--(D8), in which $D_k$ without argument denotes the set of the polymers, involve none of the operators $S_k$, $S_{k,y}$, $\mathfrak{H}_{k}$, $D_k(\cdot)$, $\dot{M}_k$, $\Gamma_{k,y}$, $\mathrm{Sq}_{k}$, $A^{\min}_{k}$.
\item[(2)]
The definitions of Section~I.5.6 are formed with $\operatorname{Vol}(X)$, integrals over $X$, the transports $U_{\mathscr{T}}$, the chains and the scalings $\mathscr{A}_L$, $\mathscr{C}_L$, $\psi_L$, $\chi_L$, and involve no operator of the list.
\item[(3)]
The definitions of $\acute{E}_k$, $E^{d}_k$, $E_k^{\mathrm{det}}$, ${\delta E^{\mathrm{det}}}_k$ in Section~I.5.7 are formed with $\gamma_5$, $B_{k,y}(A)$, $Q_k(A)J^{*}$, $JQ_k^{T}(-A)$ and $F_0(A)$, which are the nos.~1, 4, 6 and 8 of the table, and with $S_{k,y}(X,A)$, which has a polymer as first argument and is a member of (p5).
\end{itemize}
This completes the proof of Proposition~25.1.
\qed

\subsection*{3. Proposition~25.2}

\subsubsection*{3.1. The statement}

Proposition~25.2 is stated in Part~II, for the data of Section~G3.5.0, and it asserts the following.
Each of the families $O_{\ell}(X;\mathrm{data})$ of its clauses (o1)--(o9) is given by one formula of the range $\mathrm{rg}$ named in the clause, in the sense of Section~I.1.9 together with (G3.r1)--(G3.r3).
That is, there is one map $\Phi$, which does not contain $\ell$ and is defined for all lifts $\tilde{X}$ admitted in the clause and for all data on the $\mathrm{rg}$-neighbourhood of $\tilde{X}$, such that
\[
\Phi\left(\tilde{X}+v,T_v\,\mathrm{data}\right)=T_v\,\Phi\left(\tilde{X},\mathrm{data}\right)
\]
for all $v\in LM\mathbb{Z}^3$, and for (o1)--(o6), (o8), (o9) also for all $v\in M\mathbb{Z}^3$, and such that $O_{\ell}(X;\mathrm{data})$ is the value of $\Phi$ at $\tilde{X}$ and at the data composed with the covering map, for every $\ell\in\mathrm{Cl}$ and every polymer $X$ of $T_k$ with a lift $\tilde{X}$ whose sup-diameter satisfies
\[
\mathrm{sdm}(\tilde{X})+2\,\mathrm{rg}+M<\ell.
\]
For an element of a Grassmann algebra this identity means $O_{\ell}(X;\mathrm{data})=\mathrm{cov}_{\ell}\left(\Phi(\tilde{X},\dots)\right)$, and for a homomorphism it means $\mathrm{cov}_{\ell}\circ\Phi(\tilde{X},\dots)=O_{\ell}(X;\mathrm{data})\circ\mathrm{cov}_{\ell}$.
The value does not depend on the choice of the lift.
For the operation $L$ of (o7), in which $D_{k+1}$ occurs, the side $\ell$ is any side with $\ell\in\mathrm{Cl}$ and $\ell/L\in\mathrm{Cl}$, and a polymer $Y$ of level $k+1$, a lift $\tilde{Y}'$ of $Y$ and the side $\ell/L$ take the places of $X$, $\tilde{X}$ and $\ell$.
Three of the objects are not attached to a polymer: the homomorphism $S_L$ of (o7) and the integration of (o9) are indexed by a finite set $\tilde{U}$ of sites of $\mathbb{L}_k$, which takes the place of $\tilde{X}$, under the condition on $(\tilde{U},\ell)$ stated in the clause, and (o10) contains neither a polymer nor a range; for (o10) the assertion is the statement made in that clause.
The set $\bar{X}^{L}$ of (o7) is a set of sites, namely the image of $\Phi(\tilde{X})$ under the covering map.

The ten clauses concern, in this order, the following objects:
(o1) the volume $\operatorname{Vol}(X)=M^3\cdot|\tilde{X}|_M$ and, for a constant $4\times4$ matrix $\mathrm{Mx}$, the form $\int_X\mathrm{d}x\,\bar{\psi}(x)\cdot\mathrm{Mx}\cdot\psi(x)$, with $\mathrm{rg}=0$;
(o2) the pairing $\langle\bar{\Psi},O\Psi\rangle$ of the unit lattice, for an operator $O$ of the torus whose kernel vanishes unless both indices lie in (over) $X$, with $\mathrm{rg}=0$;
(o3) the extraction, that is, the numbers $\mathrm{al}_{0}(E,X)$, $\mathrm{al}_{2}(E,X)$ of Section~I.5.6 for a family $E$ with (k2) and a small polymer $X$, and the element $(RE)(X)$, with $\mathrm{rg}=0$ and $d_M(\tilde{X})\le L$;
(o4) the adjustment, that is, the homomorphism $\boldsymbol{\Theta}_{X,\mathscr{T}}$ of the display (I.5.3), with $\mathrm{rg}=0$ and $d_M(\tilde{X})\le L$;
(o5) the substitution, that is, the homomorphism (G3.8.3a)$^{\mathrm{ch}}$ formed with the rows nos.~9 and 10 of Proposition~25.1, for $k\ge1$, with $\mathrm{rg}=\eta_k$;
(o6) the vertex $\mathrm{Vx}(\mathrm{Zf},\tilde{\square})$ of the display (23.3), for $k=0$, with $\mathrm{rg}=1$;
(o7) the reblocking, that is, the homomorphism $S_L$ and the set $\bar{X}^{L}$, both with $\mathrm{rg}=0$, and the operation $L$ of Section~I.5.6, which is attached to a polymer of level $k+1$, has the range $\eta_{k+1}$ at level $k+1$ and is considered for $\ell\in\mathrm{Cl}$ with $\ell/L\in\mathrm{Cl}$;
(o8) the traces of $\acute{E}_k$, with $\mathrm{rg}=0$, through which, in the setting in which Section~I.5.7 forms these numbers, $\acute{E}_k(X,B)$ is the value of a functional that does not contain $\ell$, from which $E^{d}_k$, $E_k^{\mathrm{det}}$ and ${\delta E^{\mathrm{det}}}_k$ are formed;
(o9) the integration over $W$, with $\mathrm{rg}=0$, which commutes with $\mathrm{cov}_{\ell}$ on $\mathrm{Grc}_k(\tilde{U})$ for every finite set $\tilde{U}$ of sites of $\mathbb{L}_k$ and every $\ell\in\mathrm{Cl}$ with $\mathrm{sdm}(\tilde{U})+M<\ell$;
(o10) the integration over $\mathsf{Z}$, which is not attached to a polymer: for $p_{0,k}>0$, a condition that follows from $e_k<1$, the measure $\hat{\mu}$ is, in the coordinates $\mathsf{Z}(\Upsilon)$, the product over the basis of copies of one probability measure on $\mathbb{R}$ that depends on $p_{0,k}$ only, and
\[
\delta\varepsilon^{0}_{k}=-d_V\cdot\log\mathbb{P}\left(|\mathrm{gn}|\le p_{0,k}\right),
\]
where $\mathrm{gn}$ is a standard normal random variable and $d_V=2-2L^{-3}$.

In its closing paragraph the proposition determines which objects and operations occur in the displays (G3.8.1)$^{\mathrm{ch}}$--(G3.8.6)$^{\mathrm{ch}}$ of Definition~24.D, in the operations $\varepsilon$, $m$, $R$, $\mathrm{nat}$, $L$ of Section~I.5.6 and in the displays of $\acute{E}_k$, $E^{d}_k$, $E_k^{\mathrm{det}}$, ${\delta E^{\mathrm{det}}}_k$ of Section~I.5.7.
Left aside are the primitives, the data $(m,\mathscr{E})$, $\varepsilon^{\mathrm{v}}$ and the family $\mathfrak{P}_k$, the connections with their pointwise sums and their multiples by numbers, the parameters, the variables of integration, the generators, and the operations of a Grassmann algebra over $\mathbb{C}$ (finite sums, products, multiples by numbers, $\exp$, the series of (G3.8.4), the part $F^{[2,0]}$ of bidegree $(2,0)$, the application of a row of an operator to a family of generators, and the homomorphism of Definition~24.D\,(D1) which adds a family of odd elements to the generators); $\mathrm{cov}_{\ell}$ commutes with each of these operations.
Apart from these, the displays and operations named contain the following objects and operations, and no others:
in (G3.8.1)$^{\mathrm{ch}}$, the operators $C$, $\delta_{\alpha,\pm B}$, the matrix $\mathsf{C}$ and the map $B\mapsto-B$;
in (G3.8.3)$^{\mathrm{ch}}$, the objects (o2) (in $G^{(1)}$, for $k\ge1$), (o1) (in $G^{(3)}$ and $G^{(5)}$), (o5) and (o6), the functionals $\oint_{\mathrm{rad}}$, and $m'=m+m(\mathscr{E})$, $\mathscr{E}^{r}=R\mathscr{E}$;
in (G3.8.2), the combinatorial data $P_0(X)$, $\mathrm{Yp}(X,P)$, $Y'(X,P)$, the sum over the pairs $(X,P)$ with $Y'(X,P)=\mathbf{Y}$, the operation $\int\mathrm{d}s_P\,\frac{\partial}{\partial s_P}$ and the distinction, made in Definition~24.D\,(D3), between single $M$-cubes and other polymers;
in (G3.8.4), the objects (o9) and (o10), the coefficients $\mathrm{urs}^{T}$, the index sets (the sets $\{X_i\}$ with the union $X$ and the families $(X_1,\dots,X_n)$ with the union $X$) and the numbers $1/n!$;
in (G3.8.6)$^{\mathrm{ch}}$, the object (o7) together with the row no.~14 of Proposition~25.1, the operations $R$ and $\mathrm{nat}$, the displays of $\acute{E}_k$, and the numbers $L^3$, $L$, $\varepsilon(\mathscr{E})$, $m(\mathscr{E})$ and $\delta\varepsilon^{0}_{k}$ of (o10);
in $R$, the distinction between small and large polymers, (o1) and (o3);
in $\varepsilon$ and $m$, the object (o3) and the sum over the small polymers that contain a given $M$-cube, which, for a lift $\tilde{\square}$ of the cube, are the images of the finite non-empty unions $\tilde{X}$ of $M$-cubes of $\mathbb{L}_k$ that are connected through faces, contain $\tilde{\square}$ and satisfy $d_M(\tilde{X})\le L$, a set of unions which does not contain $\ell$;
in $\mathrm{nat}$, the distinction between small and large polymers and (o4);
in $L$, the object (o7), that is, the set $\bar{X}^{L}$, the index set of the sum over the polymers $X$ with $\bar{X}^{L}=LY$ and the homomorphism $S_L$, together with the row no.~14 of Proposition~25.1, that is, $\mathscr{A}_L$ and $\mathscr{C}_L$;
in $\acute{E}_k$, the object (o8) and the integrations over $y$ with the factors $-i\,b_k^2$, $1/y$ and $-i$;
in $E^{d}_k$, $E_k^{\mathrm{det}}$, ${\delta E^{\mathrm{det}}}_k$, the maps $B\mapsto-B$, $B\mapsto\overline{B}$, complex conjugation, the mean of four numbers and differences.

\subsubsection*{3.2. Scheme of the proof}

For each of the families (o1)--(o9) four points are established.
The family (o9) is attached to a finite set $\tilde{U}$ of sites of $\mathbb{L}_k$ with $\mathrm{sdm}(\tilde{U})+M<\ell$, which takes the place of the $\mathrm{rg}$-neighbourhood of a lift; consequently it holds with $\tilde{U}=\tilde{X}^{+\mathrm{rg}}$ for every $\mathrm{rg}$ and every lift $\tilde{X}$ with the margin of (C).
The assertion (o10) does not concern an object attached to a polymer; it is proved, as formulated there, in 3.11.
\begin{itemize}
\item[(A)]
$\Phi$ is given by a finite expression on the covering lattice, for all admitted $\tilde{X}$ and all data named; the side $\ell$ does not occur in it.
\item[(B)]
$\Phi(\tilde{X}+v,T_v\,\mathrm{data})=T_v\,\Phi(\tilde{X},\mathrm{data})$ for the $v$ named.
Here $(T_v\,\mathrm{data})(t_1,t_2,\dots):=\mathrm{data}(t_1-v,t_2-v,\dots)$; the translation $T_v$ acts trivially on numbers and on the parameters $y$, $t$, on elements by the relabelling $T_v$ of 1.5, and on a linear map $\mathrm{hm}$ between subalgebras by $\mathrm{hm}\mapsto T_v\,\mathrm{hm}\,T_v^{-1}$.
\item[(C)]
The identity with the object on the torus holds for every $\ell\in\mathrm{Cl}$ and every lift $\tilde{X}$ of $X$ with $\mathrm{sdm}(\tilde{X})+2\,\mathrm{rg}+M<\ell$.
Under this condition $\mathrm{sdm}(\tilde{X}^{+\mathrm{rg}})+M<\ell$, so that $\mathrm{cov}_{\ell}$ is defined on $\mathrm{Grc}_k(\tilde{X}^{+\mathrm{rg}})$ (Section~G3.5.0, Lemma~4\,(a)).
(For the operation $L$ of (o7) the polymer is a polymer $Y$ of level $k+1$ and the side is $\ell/L$; the four points are established at that level, in 3.9.)
\item[(D)]
The value does not depend on the lift.
\end{itemize}

\emph{Point (D), for all families at once.}
Two lifts of $X$ are $\tilde{X}$ and $\tilde{X}+\ell w$ (Lemma~1\,(b)).
The data lifted to the neighbourhood of $\tilde{X}+\ell w$ are the images under $T_{\ell w}$ of the data lifted to the neighbourhood of $\tilde{X}$: for a bond function, $b\mapsto B(\mathrm{pr}_{\ell}b)$ is $\ell\mathbb{Z}^3$-periodic; for a kernel of the torus with two arguments, $O(t_1+\ell w,t_2+\ell w)=\boldsymbol{\chi}(w)^2\cdot O(t_1,t_2)=O(t_1,t_2)$ by (I.AP).
Since $\ell w$ lies in $LM\mathbb{Z}^3$, (B) gives
\[
\Phi(\tilde{X}+\ell w,\text{lifted data})=T_{\ell w}\,\Phi(\tilde{X},\text{lifted data}).
\]
If the value is a number, the two values are equal.
If it is an even element $F$, then $\mathrm{cov}_{\ell}\left(T_{\ell w}F\right)=\mathrm{cov}_{\ell}\left(F\right)$ (Lemma~4\,(d1)).
If it is a linear map $\mathrm{hm}$ that commutes with $\mathrm{Par}$, the two maps $O$, $O'$ of the torus determined by $\mathrm{cov}_{\ell}\circ\mathrm{hm}=O\circ\mathrm{cov}_{\ell}$ and $\mathrm{cov}_{\ell}\circ\left(T_{\ell w}\,\mathrm{hm}\,T_{\ell w}^{-1}\right)=O'\circ\mathrm{cov}_{\ell}$ coincide (Lemma~4\,(d2)).
In each case below it therefore suffices to note which of these three cases occurs.
A homomorphism that maps every generator to a linear combination of generators commutes with $\mathrm{Par}$.
(For (o3) and (o4) the object on the torus is \emph{defined} through a lift ((I.AP) and Section~I.5.6); there (D) shows that it is well defined.
For the other families the object on the torus is defined without a lift, and (D) also follows from (C).)

\emph{Uniformity.}
All assertions are identities.
They hold for every integer $k\ge0$, every complex $\mu$, every $e_k>0$ and every $\ell\in\mathrm{Cl}$, with three restrictions on the data: (o10) is stated for $p_{0,k}>0$ (3.11); the assertions of (o8) on $\acute{E}_k$, $E^{d}_k$, $E_k^{\mathrm{det}}$, ${\delta E^{\mathrm{det}}}_k$ are stated in the setting in which Section~I.5.7 forms these numbers (3.10, the paragraph on $\acute{E}_k$); the operation $L$ of (o7) is stated for $\ell$ and $\ell/L$ in $\mathrm{Cl}$ (3.9).
No estimate and no further constant enters.

\subsubsection*{3.3. (o1)}

\begin{itemize}
\item[(A)]
Put
\begin{gather*}
\Phi'(\tilde{X}):=M^3\cdot|\tilde{X}|_M,\\
\Phi(\tilde{X};\mathrm{Mx}):=\int_{\tilde{X}}\mathrm{d}x\,\bar{\psi}(x)\cdot\mathrm{Mx}\,\psi(x)=\frac{\eta_k^3}{48}\sum_{c\in\tilde{X}}\ \sum_{g\in O_{\mathrm{cub}}}\ \sum_{\beta,\beta'}\bar{\psi}_{\beta}(c,g)\cdot\mathrm{Mx}_{\beta\beta'}\cdot\psi_{\beta'}(c,g),
\end{gather*}
an element of $\mathrm{Grc}_k(\tilde{X})$.
\item[(B)]
Let $v\in M\mathbb{Z}^3$.
Then $|\tilde{X}+v|_M=|\tilde{X}|_M$ by (5.c), and by (5.d)
\[
\int_{\tilde{X}+v}\mathrm{d}x\,\bar{\psi}(x)\cdot\mathrm{Mx}\,\psi(x)=\int_{\tilde{X}}\mathrm{d}x\,\bar{\psi}(x+v)\cdot\mathrm{Mx}\,\psi(x+v)=T_v\,\Phi(\tilde{X};\mathrm{Mx}).
\]
\item[(C)]
One has $\operatorname{Vol}(X)=M^3\cdot|X|_M=M^3\cdot|\tilde{X}|_M$ (Lemma~1\,(b)).
On the torus the function $\mathrm{hf}(x):=\bar{\psi}^{(\ell)}(x)\cdot\mathrm{Mx}\,\psi^{(\ell)}(x)$ satisfies $\mathrm{hf}(x+\ell w)=\boldsymbol{\chi}(w)^2\cdot\mathrm{hf}(x)=\mathrm{hf}(x)$, and Lemma~4\,(c) gives
\[
\int_X\mathrm{d}x\,\mathrm{hf}(x)=\int_{\tilde{X}}\mathrm{d}x\,\mathrm{hf}(x)=\mathrm{cov}_{\ell}\left(\Phi(\tilde{X};\mathrm{Mx})\right).
\]
This is a sum over the $48\,\eta_k^{-3}\cdot M^3\cdot|\tilde{X}|_M$ fermion sites over the lift, not over the periodic extension of $X$.
The case $\mathrm{Mx}=1$, $\tilde{X}=\tilde{\square}$ is $\langle\bar{\psi},\psi\rangle_{\square}$.

\emph{Kernel} (Section~I.5.1, matrix notation).
With $K(x,x'):=48\,\eta_k^{-3}\cdot\mathbf{1}[x=x']\cdot\mathbf{1}[c\in\tilde{X}]\cdot\mathrm{Mx}$ one has
\[
\int\mathrm{d}x'\,K(x,x')\,\psi(x')=\frac{\eta_k^3}{48}\cdot48\,\eta_k^{-3}\cdot\mathbf{1}[c\in\tilde{X}]\cdot\mathrm{Mx}\,\psi(x),
\]
and hence
\[
\int\mathrm{d}x\int\mathrm{d}x'\,\bar{\psi}(x)^{T}K(x,x')\,\psi(x')=\int_{\tilde{X}}\mathrm{d}x\,\bar{\psi}(x)^{T}\mathrm{Mx}\,\psi(x).
\]
\item[(D)]
The values are a number and an even element.
\end{itemize}

\subsubsection*{3.4. (o2)}

\begin{itemize}
\item[(A)]
For a kernel $O(u,u')$ with values in the $4\times4$ matrices, given at the unit indices in $\tilde{X}$ (for $k=0$: over $\tilde{X}$), put
\[
\Phi(\tilde{X};O):=\sum_{y,y'\in\tilde{X}\cap\mathbb{Z}^3}\ \sum_{\beta,\beta'}\bar{\Psi}_{\beta}(y)\cdot O(y,y')_{\beta\beta'}\cdot\Psi_{\beta'}(y')\qquad(k\ge1),
\]
\begin{align*}
\Phi(\tilde{X};O)&:=\left(\frac{1}{48}\right)^2\sum_{x,x'\text{ over }\tilde{X}}\ \sum_{\beta,\beta'}\bar{\Psi}_{\beta}(x)\cdot O(x,x')_{\beta\beta'}\cdot\Psi_{\beta'}(x')\\
&\phantom{:}=\int_{\tilde{X}}\mathrm{d}x\int_{\tilde{X}}\mathrm{d}x'\,\bar{\Psi}(x)\cdot O(x,x')\,\Psi(x')\qquad(k=0).
\end{align*}
\item[(B)]
Let $v\in M\mathbb{Z}^3$; this group is contained in $\mathbb{Z}^3$.
Replace $(y,y')$ by $(y+v,y'+v)$; one has $(T_vO)(y+v,y'+v)=O(y,y')$.
\item[(C)]
Let $O$ be an operator on $\mathfrak{S}(1,\ell)$ (for $k=0$: on $\mathfrak{S}^{\mathrm{ch}}(1,\ell)$).
Its kernel satisfies $O(t,t'+\ell w)=\boldsymbol{\chi}(w)\cdot O(t,t')$ (Section~I.1.6), and also $O(t+\ell w,t')=\boldsymbol{\chi}(w)\cdot O(t,t')$, because $Of$ is antiperiodic for every $f$.
Hence $\bar{\Psi}^{(\ell)}(t)\cdot O(t,t')\,\Psi^{(\ell)}(t')$ is $\ell\mathbb{Z}^3$-periodic in $t$ and in $t'$ separately, and $\langle\bar{\Psi},O\Psi\rangle$ (Section~I.3.6) is its sum (for $k=0$: its integral $\int\mathrm{d}x\int\mathrm{d}x'$) over any two sets of representatives.
If the kernel vanishes unless both indices lie in $X$ (for $k=0$: over $X$), only the indices in (over) $X$ contribute, and their representatives are taken in (over) $\tilde{X}$ (Lemma~4\,(c)).
The result is $\mathrm{cov}_{\ell}\left(\Phi(\tilde{X};O)\right)$, where $O(u,u')$ are the values of the kernel at the indices in (over) $\tilde{X}$.
\item[(D)]
The value is an even element.
\end{itemize}

\subsubsection*{3.5. (o3)}

\begin{itemize}
\item[(A)]
For $\tilde{X}$ with $d_M(\tilde{X})\le L$, a number $\mathrm{Ec}_{00}$ and matrices $\mathrm{Ec}_{20}(x,x')$, $x,x'$ over $\tilde{X}$, put
\begin{gather*}
\mathrm{al}_{0}:=\operatorname{Vol}(\tilde{X})^{-1}\cdot\mathrm{Ec}_{00},
\qquad
\mathrm{al}_{2}:=\operatorname{Vol}(\tilde{X})^{-1}\cdot\left(\frac{\eta_k^3}{48}\right)^2\sum_{x,x'\text{ over }\tilde{X}}\mathrm{Ec}_{20}(x,x'),\\
\operatorname{Vol}(\tilde{X}):=M^3\cdot|\tilde{X}|_M.
\end{gather*}
Both sums include the sum over the $48$ labels.
\item[(B)]
For $v\in M\mathbb{Z}^3$ one has $d_M(\tilde{X}+v)=d_M(\tilde{X})$ and $\operatorname{Vol}(\tilde{X}+v)=\operatorname{Vol}(\tilde{X})$, and the sum over $(x+v,x'+v)$ of $(T_v\mathrm{Ec}_{20})(x+v,x'+v)=\mathrm{Ec}_{20}(x,x')$ is unchanged.
\item[(C)]
Let $E$ be a family with the property (k2) of Section~I.5.4.

\emph{The element over the lift} (for every polymer, small or large, that has a lift as follows).
Let $X$ be a polymer of $T_k$ with a lift $\tilde{X}$ such that $\mathrm{sdm}(\tilde{X})+2\eta_k+M<\ell$.
Then $\mathrm{sdm}(\tilde{X}^{+\eta_k})+M\le\mathrm{sdm}(\tilde{X})+2\eta_k+M<\ell$, so that $\mathrm{cov}_{\ell}$ is defined on $\mathrm{Grc}_k(\tilde{X}^{+\eta_k})$ (Section~G3.5.0, Lemma~4\,(a)).
By Section~I.5.4, for all $(\mathscr{A},\mathscr{C})$,
\begin{multline*}
E(X;\mathscr{A},\mathscr{C})=\sum_{n,n'}\frac{1}{n!\,n'!}\int\mathrm{d}\xi\,\mathrm{d}\boldsymbol{\zeta}\,E_{nn'}(X;\mathscr{A},\mathscr{C};\xi_1,\dots,\xi_n;\boldsymbol{\zeta}_1,\dots,\boldsymbol{\zeta}_{n'})\\
\times\psi^{(\ell)}(\xi_1)\cdots\psi^{(\ell)}(\xi_n)\,\chi^{(\ell)}(\boldsymbol{\zeta}_1)\cdots\chi^{(\ell)}(\boldsymbol{\zeta}_{n'}),
\end{multline*}
where the integrals extend over fundamental domains.
The kernels and the generators are antiperiodic in each argument ((I.AP)); hence the integrand is $\ell\mathbb{Z}^3$-periodic in each argument separately, and the representatives may be chosen freely.
By (k2) the integrand vanishes unless all $\xi_i$ lie over $X$ and all $\boldsymbol{\zeta}_j$ are admissible for $X$.
The sites over $X$ have their representatives over $\tilde{X}$; the pairs admissible for $X$ have their representatives among the pairs of $\mathrm{Pzt}_k$ that are admissible for $\tilde{X}$, one for each (Lemma~1\,(d), which requires $\mathrm{sdm}(\tilde{X})+2\eta_k<\ell$), and both photon sites of such a pair lie in $\tilde{X}^{+\eta_k}$.
Hence $E(X;\mathscr{A},\mathscr{C})=\mathrm{cov}_{\ell}(\mathrm{Ecv})$, where $\mathrm{Ecv}=\mathrm{Ecv}(\tilde{X};\mathscr{A},\mathscr{C})\in\mathrm{As}_k(\tilde{X}^{+\eta_k})$ is the same expression with the generators of the covering algebra, with the $\xi_i$ ranging over the indices over $\tilde{X}$ and the pairs over the pairs admissible for $\tilde{X}$, and with the values of the kernels at these arguments.
Since $\mathrm{cov}_{\ell}$ is injective and preserves the bidegree and the parity (Lemma~4\,(a)), one has $\mathrm{Ecv}=\mathrm{cov}_{\ell}^{-1}\left(E(X;\mathscr{A},\mathscr{C})\right)$; moreover $\mathrm{Ecv}$ is even if $E(X;\mathscr{A},\mathscr{C})$ is even, its part of bidegree $(0,0)$ is the number $E_{00}(X;\mathscr{A},\mathscr{C})$, and its part of bidegree $(2,0)$ is
\[
\frac{1}{2}\int\mathrm{d}\xi_1\,\mathrm{d}\xi_2\,E_{20}(\xi_1,\xi_2)\,\psi(\xi_1)\,\psi(\xi_2),
\]
with $\xi_1$, $\xi_2$ over $\tilde{X}$.
In this part the terms with the bar indices $(\omega_1,\omega_2)=(1,0)$ and $(0,1)$ are equal, by the antisymmetry of the kernel and of the product; their sum is
\[
\int_{\tilde{X}}\mathrm{d}x\int_{\tilde{X}}\mathrm{d}x'\,\bar{\psi}(x)^{T}E_{20}(X;\mathscr{A},\mathscr{C};x,x')\,\psi(x'),
\]
with the matrix $E_{20}(x,x')$ of Section~I.5.1, whose entries are $E_{20}\left((x,\beta,1),(x',\beta',0)\right)$.
For the families of Section~I.5.4, which have equally many barred and unbarred legs, this is the whole part of bidegree $(2,0)$.
The monomials $\bar{\psi}_{\beta}(x)\,\psi_{\beta'}(x')$, $x,x'$ over $\tilde{X}$, are linearly independent of one another and of the monomials with two unbarred or two barred generators; hence the element determines the matrices $E_{20}(X;\mathscr{A},\mathscr{C};x,x')$ at the arguments over $\tilde{X}$.

\emph{Another lift.}
Let $E(X;\mathscr{A},\mathscr{C})$ be even and $w\in\mathbb{Z}^3$.
Then $T_{\ell w}\mathrm{Ecv}(\tilde{X};\mathscr{A},\mathscr{C})$ lies in $\mathrm{As}_k\left((\tilde{X}+\ell w)^{+\eta_k}\right)$, and by Lemma~4\,(b)
\[
\mathrm{cov}_{\ell}\left(T_{\ell w}\mathrm{Ecv}(\tilde{X};\mathscr{A},\mathscr{C})\right)=\mathrm{cov}_{\ell}\left(\mathrm{Ecv}(\tilde{X};\mathscr{A},\mathscr{C})\right)=E(X;\mathscr{A},\mathscr{C}).
\]
As $\mathrm{cov}_{\ell}$ is injective on $\mathrm{Grc}_k\left((\tilde{X}+\ell w)^{+\eta_k}\right)$, it follows that $\mathrm{Ecv}(\tilde{X}+\ell w;\mathscr{A},\mathscr{C})=T_{\ell w}\mathrm{Ecv}(\tilde{X};\mathscr{A},\mathscr{C})$.

\emph{The numbers.}
Let $X$ be small.
By Lemma~1\,(c), $X$ has a lift $\tilde{X}$, with $d_M(\tilde{X})\le L$ and $\mathrm{sdm}(\tilde{X})+2+M<\ell$; since $\eta_k\le1$, the paragraph on the element over the lift applies to $X$ and $\tilde{X}$.
One has $\mathrm{al}_{0}(E,X)=\operatorname{Vol}(X)^{-1}\cdot E_{00}(X;0,0)$, and $\operatorname{Vol}(X)=\operatorname{Vol}(\tilde{X})$ by (o1).
In
\[
\mathrm{al}_{2}(E,X)=\operatorname{Vol}(X)^{-1}\int\mathrm{d}x\int\mathrm{d}x'\,E_{20}(X;0,0;x,x')
\]
the kernel is integrated without generators; by the last sentence of (I.AP) it is evaluated with both arguments over one lift of $X$, and by (k2) it vanishes off $X$.
Hence
\[
\mathrm{al}_{2}(E,X)=\operatorname{Vol}(\tilde{X})^{-1}\int_{\tilde{X}}\mathrm{d}x\int_{\tilde{X}}\mathrm{d}x'\,E_{20}(X;0,0;x,x').
\]
These are the numbers of (A) at $\mathrm{Ec}_{00}=E_{00}(X;0,0)$ and $\mathrm{Ec}_{20}=E_{20}(X;0,0;\cdot,\cdot)$ over $\tilde{X}$, that is, at the parts of bidegree $(0,0)$ and $(2,0)$ of $\mathrm{Ecv}(\tilde{X};0,0)=\mathrm{cov}_{\ell}^{-1}\left(E(X;0,0)\right)$.

\emph{The operation $R$.}
By Section~I.5.6, and by (o1) with $\mathrm{Mx}=\mathrm{al}_{2}$, one has for small $X$
\begin{align*}
\mathrm{cov}_{\ell}^{-1}\left((R E)(X;\mathscr{A},\mathscr{C})\right)&=\mathrm{Ecv}(\tilde{X};\mathscr{A},\mathscr{C})-\mathrm{al}_{0}\cdot\operatorname{Vol}(\tilde{X})\\
&\qquad-\int_{\tilde{X}}\mathrm{d}x\,\bar{\psi}(x)\cdot\mathrm{al}_{2}\,\psi(x).
\end{align*}
For a large $X$ with a lift as in the first paragraph, $(R E)(X;\mathscr{A},\mathscr{C})=E(X;\mathscr{A},\mathscr{C})=\mathrm{cov}_{\ell}\left(\mathrm{Ecv}(\tilde{X};\mathscr{A},\mathscr{C})\right)$.
\item[(D)]
The values are numbers.
(By this, $\mathrm{al}_{2}(E,X)$ is well defined.)
\end{itemize}

\subsubsection*{3.6. (o4)}

\begin{description}
\item[(A)]
Let $d_M(\tilde X)\le L$, and let $\tilde B$ be given on the bonds contained in $\tilde X$.
For fermion sites $x=(c,g)$ and $\tilde x=(\tilde c,\tilde g)$ over $\tilde X$, let $\mathrm{Ch}(x,\tilde x)$ denote the set of all shortest chains $(\tilde\square_0,\dots,\tilde\square_n)$ of $M$-cubes of $\tilde X$ which lead from the cube of $c$ to the cube of $\tilde c$ and in which any two consecutive cubes have a common face.
This set is finite, and it is not empty because $\tilde X$ is connected through faces.
The centre $z_i\in M\mathbb{Z}^3$ of $\tilde\square_i$ is a site of $\mathbb{L}_k$.
As in (I.SY6)$^{\mathrm{ch}}$, set
\[
p_0:=x,\qquad p_{i+1}:=(z_i,g)\quad(0\le i\le n),\qquad p_{n+2}:=\tilde x .
\]

\emph{The pairs.}
Let $0\le i\le n+1$ with $p_i\ne p_{i+1}$.
Then $(p_i,p_{i+1})$ belongs to $\mathrm{Pzt}_k$ and satisfies (I) for $\tilde X$, in the sense of Lemma~1\,(d); this is seen in three cases.

\emph{Case $i=0$.}
The sites $c$ and $z_0$ lie in the lattice box $\tilde\square_0$.
Hence $\mathrm{hul}(c,z_0)\subset\tilde\square_0$ by Fact~2, and $|c-z_0|\le(M-\eta_k)/2$ by Fact~1\,(i).

\emph{Case $1\le i\le n$.}
One has $z_i-z_{i-1}=\pm M e_{\mu}$, and each of the six paths is the straight segment from $z_{i-1}$ to $z_i$.
A site $z_{i-1}\pm a e_{\mu}$ of this segment, $0\le a\le M$, lies in $\tilde\square_{i-1}$ if $a<M/2$ and in $\tilde\square_i$ if $M-a<M/2$; the value $a=M/2$ does not occur, since $M/(2\eta_k)$ is not an integer.
Here $d_X(p_i,p_{i+1})=M$.

\emph{Case $i=n+1$.}
The sites $z_n$ and $\tilde c$ lie in the lattice box $\tilde\square_n$; hence $\mathrm{hul}(z_n,\tilde c)\subset\tilde\square_n$ and $|z_n-\tilde c|\le(M-\eta_k)/2$.
If $\tilde c=z_n$ and $\tilde g\ne g$, the pair lies inside the cluster over $z_n$, one has $d_X(p_{n+1},p_{n+2})=\eta_k$, and there is no path.

Consequently $d_X(p_i,p_{i+1})\le M$, and both photon sites of the pair lie in $\tilde X$.
The transport
\[
u(p_i,p_{i+1}):=U_{\tilde B}(p_i,p_{i+1})
\]
of (I.1.6) is the exponential of a linear form in the numbers $\tilde B(b)$, $b\in\mathrm{Pth}\left(\pi(p_i),\pi(p_{i+1})\right)$, and these bonds are contained in $\tilde X$.
If the two points lie over one photon site, then $u(p_i,p_{i+1})=1$.

\emph{The map.}
Put
\[
U_i:=\prod_{j<i}u(p_j,p_{j+1}),\qquad U^{-}_i:=\prod_{j<i}u(p_j,p_{j+1})^{-1},
\]
let $\operatorname{Mn}$ be the mean over $\mathrm{Ch}(x,\tilde x)$, and write
\[
\int_{\tilde X}\mathrm{d}\tilde x:=\frac{\eta_k^3}{48}\sum_{\tilde c\in\tilde X}\ \sum_{\tilde g\in O_{\mathrm{cub}}}.
\]
For $x$ over $\tilde X$ define, componentwise in the spinor index,
\begin{align*}
\boldsymbol{\Theta}_{\tilde X,\tilde B}\,\psi(x)&:=\operatorname{Vol}(\tilde X)^{-1}\int_{\tilde X}\mathrm{d}\tilde x\,\operatorname{Mn}\left\{U_{n+2}\cdot\psi(\tilde x)-\sum_{i=0}^{n+1}U_i\cdot d_X(p_i,p_{i+1})^{\alpha}\cdot\chi(p_i,p_{i+1})\right\},\\
\boldsymbol{\Theta}_{\tilde X,\tilde B}\,\bar{\psi}(x)&:=\operatorname{Vol}(\tilde X)^{-1}\int_{\tilde X}\mathrm{d}\tilde x\,\operatorname{Mn}\left\{U^{-}_{n+2}\cdot\bar{\psi}(\tilde x)-\sum_{i=0}^{n+1}U^{-}_i\cdot d_X(p_i,p_{i+1})^{\alpha}\cdot\bar{\chi}(p_i,p_{i+1})\right\},
\end{align*}
where the terms with $p_i=p_{i+1}$ are omitted.
These are the displays (I.5.3), written on the covering lattice.
Their right-hand sides are linear combinations of generators of $\mathrm{Grc}_k(\tilde X)$; hence they are odd, and any two of them anticommute.
The subalgebra $\mathrm{As}$ of $\mathrm{Grc}_k(\tilde X)$ generated by the $\psi(\xi)$, $\xi$ over $\tilde X$, is free on these generators.
Therefore there is exactly one homomorphism $\boldsymbol{\Theta}_{\tilde X,\tilde B}\colon\mathrm{As}\to\mathrm{Grc}_k(\tilde X)$ that takes the values displayed.

\item[(B)]
Let $v\in M\mathbb{Z}^3$.
By (5.c) the translation by $v$ carries cubes, common faces, centres and chains to cubes, common faces, centres and chains; thus $\mathrm{Ch}(x+v,\tilde x+v)=\mathrm{Ch}(x,\tilde x)+v$, and the chain points belonging to $(x+v,\tilde x+v)$ are the points $p_i+v$.
Neither $d_X$ nor the labels change, by (5.a), and $U_{T_v\tilde B}(p+v,p'+v)=U_{\tilde B}(p,p')$ by (5.b).
Moreover $\operatorname{Vol}(\tilde X+v)=\operatorname{Vol}(\tilde X)$, and (5.d) applies to the integral over $\tilde x$.
Hence
\[
\boldsymbol{\Theta}_{\tilde X+v,T_v\tilde B}\,\psi(x+v)=T_v\,\boldsymbol{\Theta}_{\tilde X,\tilde B}\,\psi(x),\qquad
\boldsymbol{\Theta}_{\tilde X+v,T_v\tilde B}\,\bar{\psi}(x+v)=T_v\,\boldsymbol{\Theta}_{\tilde X,\tilde B}\,\bar{\psi}(x),
\]
that is, $\boldsymbol{\Theta}_{\tilde X+v,T_v\tilde B}=T_v\,\boldsymbol{\Theta}_{\tilde X,\tilde B}\,T_v^{-1}$.

\item[(C)]
Let $X$ be small and $\tilde X$ a lift of $X$; by Lemma~1\,(c), $\mathrm{sdm}(\tilde X)+M<\ell$.
Let $\mathscr{T}$ be a bond function on $T_k$, and put $\tilde B(b):=\mathscr{T}(\mathrm{pr}_\ell b)$ for the bonds $b$ contained in $\tilde X$.
Section~I.5.6 defines $\boldsymbol{\Theta}_{X,\mathscr{T}}$ by (I.5.3), all points being chosen over one lift of $X$.
By Lemma~1\,(b) the $M$-cubes of $X$ and their common faces are the images of the $M$-cubes of $\tilde X$ and of their common faces; therefore the shortest chains of $X$ between two of its cubes are the images of the shortest chains of $\tilde X$.
For a pair $(p_i,p_{i+1})$ of a chain one has $|\pi(p_i)-\pi(p_{i+1})|\le M<\ell/2$.
So the transport $U_{\mathscr{T}}(p_i,p_{i+1})$ of the torus is formed with the paths of $\mathbb{L}_k$ between the two points (Section~I.1.1), and it equals $U_{\tilde B}(p_i,p_{i+1})$.
Further, $\operatorname{Vol}(X)=\operatorname{Vol}(\tilde X)$, and $\int_X\mathrm{d}\tilde x$ is the integral $\int_{\tilde X}\mathrm{d}\tilde x$ over the lift.
It follows that, for $x$ over $\tilde X$, the element $\boldsymbol{\Theta}_{X,\mathscr{T}}\,\psi^{(\ell)}(x)$ is the first expression of (A) with the generators $\psi^{(\ell)}(\tilde x)$ and $\chi^{(\ell)}(p_i,p_{i+1})$ in place of $\psi(\tilde x)$ and $\chi(p_i,p_{i+1})$, and that $\boldsymbol{\Theta}_{X,\mathscr{T}}\,\bar{\psi}^{(\ell)}(x)$ is the second expression of (A) with $\bar{\psi}^{(\ell)}(\tilde x)$ and $\bar{\chi}^{(\ell)}(p_i,p_{i+1})$; in other words,
\[
\boldsymbol{\Theta}_{X,\mathscr{T}}\,\mathrm{cov}_\ell\left(\psi(x)\right)=\mathrm{cov}_\ell\left(\boldsymbol{\Theta}_{\tilde X,\tilde B}\,\psi(x)\right),\qquad
\boldsymbol{\Theta}_{X,\mathscr{T}}\,\mathrm{cov}_\ell\left(\bar{\psi}(x)\right)=\mathrm{cov}_\ell\left(\boldsymbol{\Theta}_{\tilde X,\tilde B}\,\bar{\psi}(x)\right).
\]
The two sides of $\mathrm{cov}_\ell\circ\boldsymbol{\Theta}_{\tilde X,\tilde B}=\boldsymbol{\Theta}_{X,\mathscr{T}}\circ\mathrm{cov}_\ell$ are homomorphisms on $\mathrm{As}$ which agree on its generators; hence they coincide.
The image $\mathrm{cov}_\ell(\mathrm{As})$ is the subalgebra generated by the $\psi(\xi)$, $\xi$ over $X$.

\item[(D)]
$\boldsymbol{\Theta}_{\tilde X,\tilde B}$ is a homomorphism that maps generators to linear combinations of generators.
(This shows that $\boldsymbol{\Theta}_{X,\mathscr{T}}$ is well defined.)
\end{description}

\paragraph{The operation nat.}
Let $E$ be a family with the property (k2) of Section~I.5.4, let $X$ be small with the lift $\tilde X$, put $\mathscr{T}:=\mathscr{A}+\mathscr{C}$, let $\tilde B$ be the lift of $\mathscr{T}$ on the bonds contained in $\tilde X$, and let $\mathrm{Ecv}:=\mathrm{Ecv}(\tilde X;\mathscr{A},\mathscr{C})$ be the element over the lift (Section~3.5\,(C)).
By (k2) the part $\mathrm{Ecv}^{[2,0]}$ contains generators $\psi(\xi)$ only with $\xi$ over $\tilde X$; so it lies in $\mathrm{As}$.
The map $\mathrm{cov}_\ell$ commutes with $F\mapsto F^{[2,0]}$ (Lemma~4\,(e)), and $\mathrm{cov}_\ell$ on $\mathrm{Grc}_k(\tilde X)$ is the restriction of $\mathrm{cov}_\ell$ on $\mathrm{Grc}_k(\tilde X^{+\eta_k})$ (Lemma~4\,(a)).
With (C) and the display of Section~I.5.6 one obtains
\begin{align*}
&\mathrm{cov}_\ell\left(\mathrm{Ecv}-\mathrm{Ecv}^{[2,0]}+\boldsymbol{\Theta}_{\tilde X,\tilde B}\left(\mathrm{Ecv}^{[2,0]}\right)\right)\\
&\qquad=E(X;\mathscr{A},\mathscr{C})-E^{[2,0]}(X;\mathscr{A},\mathscr{C})+\boldsymbol{\Theta}_{X,\mathscr{T}}\left(E^{[2,0]}(X;\mathscr{A},\mathscr{C})\right)
=E^{\mathrm{nat}}(X;\mathscr{A},\mathscr{C}).
\end{align*}
The argument of $\mathrm{cov}_\ell$ on the left lies in $\mathrm{As}_k(\tilde X^{+\eta_k})$; it is formed from $\mathrm{Ecv}$ and from the values of $\tilde B$ on the bonds contained in $\tilde X$, and $\ell$ does not occur in it.
For a large $X$ with a lift as in Section~3.5\,(C) one has $E^{\mathrm{nat}}(X;\mathscr{A},\mathscr{C})=E(X;\mathscr{A},\mathscr{C})=\mathrm{cov}_\ell\left(\mathrm{Ecv}(\tilde X;\mathscr{A},\mathscr{C})\right)$.

\subsubsection*{3.7. (o5) ($k\ge1$)}

\begin{description}
\item[(A)]
Let $\tilde B$ be given on the bonds that meet $\tilde{\mathbf{Y}}$, and extend it in an arbitrary way to a bond function on $\mathbb{L}_k$.
For a unit site $y$ of $\tilde{\mathbf{Y}}$ one has $\Delta_y\subset\tilde{\mathbf{Y}}$ by (1.c).
According to (Of10), the rows of $T^{(1)}(\tilde B,\tilde{\mathbf{Y}})$, $T^{(2)}(\tilde B,\tilde{\mathbf{Y}})$, $\bar{T}^{(1)}(\tilde B,\tilde{\mathbf{Y}})$ and $\bar{T}^{(2)}(\tilde B,\tilde{\mathbf{Y}})$ at $y$ read $\tilde B$ on bonds that are contained in $\Delta_y$ or have the form $(c,c_3)$ with $c\in\Delta_y$; all these bonds meet $\tilde{\mathbf{Y}}$.
By (l2), therefore, the rows do not depend on the extension.
The coefficients in (Of10), taken with $\tilde Y=\tilde{\mathbf{Y}}$, show where the columns are.
The row of $T^{(1)}$ has its columns at sites $x'$ with $c'\in\Delta_y$, or with $c_3(x')\in\Delta_y$ and $c'\in\tilde{\mathbf{Y}}$; these are sites over $\tilde{\mathbf{Y}}$.
The row of $T^{(2)}$ has its columns at pairs $(x,x^{(3)})$ with $c\in\Delta_y$ and $c_3\notin\tilde{\mathbf{Y}}$; for such a pair $c_3\in\tilde{\mathbf{Y}}^{+\eta_k}$, and the pair satisfies (II) for $\tilde{\mathbf{Y}}$.
Let $\mathrm{As}$ be the subalgebra of $\mathrm{Grc}_k(\tilde{\mathbf{Y}}^{+\eta_k})$ generated by $\psi$, $\chi$ and the $\Psi(u)$, $u=(y,\beta,\omega)$ with $y$ a unit site of $\tilde{\mathbf{Y}}$, and let $\mathrm{As}'$ be the subalgebra generated by $\psi$ and $\chi$.
The map $\mathrm{sub}_{\tilde{\mathbf{Y}},\tilde B}\colon\mathrm{As}\to\mathrm{As}'$ is the homomorphism which fixes $\psi$ and $\chi$ and which maps, for the unit sites $y$ of $\tilde{\mathbf{Y}}$,
\begin{align*}
\Psi(y)&\mapsto\sum_{x'}T^{(1)}(\tilde B,\tilde{\mathbf{Y}})[y,x']\cdot\psi(x')+\sum_{(x,x^{(3)})}T^{(2)}(\tilde B,\tilde{\mathbf{Y}})[y,(x,x^{(3)})]\cdot\chi(x,x^{(3)}),\\
\bar{\Psi}(y)&\mapsto\sum_{x'}\bar{T}^{(1)}(\tilde B,\tilde{\mathbf{Y}})[y,x']\cdot\bar{\psi}(x')+\sum_{(x,x^{(3)})}\bar{T}^{(2)}(\tilde B,\tilde{\mathbf{Y}})[y,(x,x^{(3)})]\cdot\bar{\chi}(x,x^{(3)}).
\end{align*}
Here a coefficient is a $4\times4$ matrix which acts on the column of the four components, and $\bar{T}^{(j)}(\tilde B,\tilde{\mathbf{Y}})[y,\cdot\,]=\mathsf{C}^{-1}\,T^{(j)}(-\tilde B,\tilde{\mathbf{Y}})[y,\cdot\,]\,\mathsf{C}$ has the same columns as $T^{(j)}(\tilde B,\tilde{\mathbf{Y}})[y,\cdot\,]$.
The images are linear combinations of generators of $\mathrm{As}'$, and $\mathrm{As}$ is free on its generators; so the homomorphism exists and is unique.
A bond that meets $\tilde{\mathbf{Y}}$ has both ends in $\tilde{\mathbf{Y}}^{+\eta_k}$; thus the data lie within the distance $\mathrm{rg}=\eta_k$ of $\tilde{\mathbf{Y}}$.

\item[(B)]
For $v\in M\mathbb{Z}^3$ the covariance is (l4) for (Of10), applied with $\tilde Y+v=\tilde{\mathbf{Y}}+v$.

\item[(C)]
Let $\mathbf{Y}$ be a polymer with the lift $\tilde{\mathbf{Y}}$, where $\mathrm{sdm}(\tilde{\mathbf{Y}})+2\eta_k+M<\ell$; let $B$ be a bond function on $T_k$, and $\tilde B$ its lift on the bonds that meet $\tilde{\mathbf{Y}}$.
Denote by $\tilde B_{\mathrm{p}}$ the periodic extension of $B$, and put $\tilde Y_{\mathrm{p}}:=\tilde{\mathbf{Y}}+\ell\mathbb{Z}^3$.
By (f3) of Section~I.1.9, which holds by Section~2.4, the rows of $T^{(j)}(B,\mathbf{Y})$ on the torus at the covering index $y$ are the rows of $T^{(j)}(\tilde B_{\mathrm{p}},\tilde Y_{\mathrm{p}})$.
For $y\in\tilde{\mathbf{Y}}\cap\mathbb{Z}^3$ the sites of $\mathrm{RdY}(y)=\{y\}\cup\{c_3(x):c\in\Delta_y\}$ have distance at most $\eta_k$ from $\tilde{\mathbf{Y}}$; by Lemma~1\,(a) with $a=\eta_k$, such a site lies in $\tilde Y_{\mathrm{p}}$ if and only if it lies in $\tilde{\mathbf{Y}}$.
Moreover $\tilde B_{\mathrm{p}}=\tilde B$ on $\mathrm{RdB}(y)$.
Hence (l2) gives
\[
T^{(j)}(\tilde B_{\mathrm{p}},\tilde Y_{\mathrm{p}})[y,\cdot\,]=T^{(j)}(\tilde B,\tilde{\mathbf{Y}})[y,\cdot\,],
\]
and the same holds for $\bar{T}^{(j)}$.
The field $\left(T^{(1)}(B,\mathbf{Y})\psi\right)(y)$ of (8.3a)$^{\mathrm{ch}}$ is $\sum_{x'}T^{(1)}(\tilde B_{\mathrm{p}},\tilde Y_{\mathrm{p}})[y,x']\cdot\psi^{(\ell)}(x')$, operators acting on fields by linearity and the sum extending over the covering indices of the row.
Likewise $\left(T^{(2)}(B,\mathbf{Y})\chi\right)(y)$ is $\sum T^{(2)}(\tilde B_{\mathrm{p}},\tilde Y_{\mathrm{p}})[y,(x,x^{(3)})]\cdot\chi^{(\ell)}(x,x^{(3)})$, and the two barred fields are the same sums formed with $\bar{T}^{(j)}$, $\bar{\psi}^{(\ell)}$ and $\bar{\chi}^{(\ell)}$.
Therefore
\[
\mathrm{sub}_{\mathbf{Y},B}\,\mathrm{cov}_\ell\left(\Psi(y)\right)=\mathrm{cov}_\ell\left(\mathrm{sub}_{\tilde{\mathbf{Y}},\tilde B}\,\Psi(y)\right),\qquad
\mathrm{sub}_{\mathbf{Y},B}\,\mathrm{cov}_\ell\left(\bar{\Psi}(y)\right)=\mathrm{cov}_\ell\left(\mathrm{sub}_{\tilde{\mathbf{Y}},\tilde B}\,\bar{\Psi}(y)\right),
\]
and both maps fix $\psi$ and $\chi$.
Hence $\mathrm{cov}_\ell\circ\mathrm{sub}_{\tilde{\mathbf{Y}},\tilde B}=\mathrm{sub}_{\mathbf{Y},B}\circ\mathrm{cov}_\ell$ on $\mathrm{As}$.
The unit sites of $\mathbf{Y}$ are the images of the unit sites of $\tilde{\mathbf{Y}}$; so $\mathrm{cov}_\ell(\mathrm{As})$ contains all $\Psi(u)$ with $u$ in $\mathbf{Y}$.

\item[(D)]
$\mathrm{sub}_{\tilde{\mathbf{Y}},\tilde B}$ is a homomorphism that maps generators to linear combinations of generators.
\end{description}

\subsubsection*{3.8. (o6) ($k=0$; $e_k$ is the charge $q$ of (23.3))}

\begin{description}
\item[(A)]
For a bond function $\mathrm{Zf}$ on the bonds that meet $\tilde\square$ let
\begin{multline*}
\mathrm{Vx}(\mathrm{Zf},\tilde\square):=\frac{1}{2}\cdot\frac{1}{48}\sum_{c\in\tilde\square}\ \sum_{g\in O_{\mathrm{cub}}}\left\{\bar{\psi}(x)\cdot M_{3}(x)\cdot\left(\exp\left(\mathrm{i}q\,\mathrm{Zf}(c,c_{3})\right)-1\right)\cdot\left(\psi(x)+\chi(x,x^{(3)})\right)\right.\\
\left.+\left(\bar{\psi}(x)+\bar{\chi}(x,x^{(3)})\right)\cdot M_{3}(x^{(3)})\cdot\left(\exp\left(-\mathrm{i}q\,\mathrm{Zf}(c,c_{3})\right)-1\right)\cdot\psi(x)\right\},
\end{multline*}
where $x=(c,g)$, $c_3=c+g e_1$, $M_3(x)=3\left(\rho-\gamma\cdot(g e_1)\right)$ and $M_3(x^{(3)})=3\left(\rho+\gamma\cdot(g e_1)\right)$.
Since $c_3\in\tilde\square^{+1}$ and $(x,x^{(3)})\in\mathrm{Pzt}_0$, this is an element of $\mathrm{Grc}_0(\tilde\square^{+1})$.
It reads $\mathrm{Zf}$ on the bonds $(c,c_3)$, $c\in\tilde\square$, which meet $\tilde\square$ and have both ends in $\tilde\square^{+1}$; hence $\mathrm{rg}=1$.

\item[(B)]
For $v\in M\mathbb{Z}^3$ the covariance follows from (5.a), (5.d) and $(T_v\mathrm{Zf})(c+v,c_3+v)=\mathrm{Zf}(c,c_3)$.

\item[(C)]
Let $\square$ be an $M$-cube of $T_0$ with the lift $\tilde\square$, and $\mathrm{Zf}$ a bond function on $T_0$.
One has $\mathrm{sdm}(\tilde\square)+2+M=2M+1<\ell$.
On the torus every term of (23.3) is the product of one barred and one unbarred generator, whose indices are formed from $x$, with a function of $\mathrm{Zf}(c,c_3)$; under $x\mapsto x+\ell w$ such a term is multiplied by $\boldsymbol{\chi}(w)^2=1$.
By Lemma~4\,(c) the integral $\int_{\square}$ is the sum over the $x$ over $\tilde\square$, and
\[
\mathrm{Vx}(\mathrm{Zf},\square)=\mathrm{cov}_\ell\left(\mathrm{Vx}(\widetilde{\mathrm{Zf}},\tilde\square)\right),
\]
where $\widetilde{\mathrm{Zf}}$ is the lift of $\mathrm{Zf}$.

\item[(D)]
$\mathrm{Vx}(\mathrm{Zf},\tilde\square)$ is an even element.
\end{description}

\subsubsection*{3.9. (o7)}

The clause (o7) concerns three objects: the homomorphism $S_L$, the set $\bar{X}^{L}$ and the operation $L$.
As in Section~1.5, $\mathrm{As}_k(\tilde U)$ and $\mathrm{As}_{k+1}(\tilde U/L)$ denote the subalgebras generated by $\psi$ and $\chi$.
Section~I.3.7 and (o7) prescribe the images of $\psi$ and $\chi$ only; accordingly $S_L$ is regarded as a map on $\mathrm{As}_k(\tilde U)$.

\paragraph{The homomorphism $S_L$} ($\mathrm{rg}=0$; no data).

\begin{description}
\item[(A)]
The set $\tilde U/L$ is a finite set of sites of $\mathbb{L}_{k+1}$.
For $(x,x')\in\mathrm{Pzt}_k$ with $c,c'\in\tilde U$ the pair $(x/L,x'/L)$ belongs to $\mathrm{Pzt}_{k+1}$ by (Of14), and its photon sites lie in $\tilde U/L$.
Let $S_L\colon\mathrm{As}_k(\tilde U)\to\mathrm{As}_{k+1}(\tilde U/L)$ be the homomorphism with
\[
S_L\,\psi(x,\beta,\omega):=L^{-1}\,\psi(x/L,\beta,\omega),\qquad S_L\,\chi(x,x',\beta,\omega):=L^{-1-\alpha}\,\chi(x/L,x'/L,\beta,\omega).
\]
It maps generators to multiples of generators, and hence it preserves the parity.
If $\tilde U\subset\tilde U_1$, the map belonging to $\tilde U$ is the restriction of the map belonging to $\tilde U_1$.

\item[(B)]
For $v\in\eta_k\mathbb{Z}^3$, in particular for $v\in LM\mathbb{Z}^3$, one has $(x+v)/L=x/L+v/L$.
Hence $T_{v/L}S_L\,\phi(t)=S_L\,T_v\,\phi(t)$ for the generators $\phi(t)$, and consequently on all of $\mathrm{As}_k(\tilde U)$.

\item[(C)]
Let $\ell\in\mathrm{Cl}$ and $\mathrm{sdm}(\tilde U)+LM<\ell$.
Then $\mathrm{sdm}(\tilde U)+M<\ell$, and
\[
\mathrm{sdm}(\tilde U/L)+M=\frac{\mathrm{sdm}(\tilde U)+LM}{L}<\frac{\ell}{L}.
\]
So $\mathrm{cov}_\ell$ on $\mathrm{Grc}_k(\tilde U)$ and $\mathrm{cov}_{\ell/L}$ on $\mathrm{Grc}_{k+1}(\tilde U/L)$ are defined (Section~G3.5.0 and Section~1.5; $\ell/L$ is an odd multiple of $M$ with $\ell/L\ge13M$, but it need not belong to $\mathrm{Cl}$).
On the torus $S_L$ is the homomorphism with
\[
S_L\,\psi^{(\ell)}(x,\cdot\,)=L^{-1}\,\psi^{(\ell/L)}(x/L,\cdot\,),\qquad S_L\,\chi^{(\ell)}(x,x',\cdot\,)=L^{-1-\alpha}\,\chi^{(\ell/L)}(x/L,x'/L,\cdot\,)
\]
(Section~I.3.7); this rule is compatible with (I.AP), because $(x+\ell w)/L=x/L+(\ell/L)w$.
On the generators,
\begin{align*}
\mathrm{cov}_{\ell/L}\left(S_L\,\psi(x,\cdot\,)\right)&=L^{-1}\,\psi^{(\ell/L)}(x/L,\cdot\,)=S_L\left(\mathrm{cov}_\ell\,\psi(x,\cdot\,)\right),\\
\mathrm{cov}_{\ell/L}\left(S_L\,\chi(x,x',\cdot\,)\right)&=L^{-1-\alpha}\,\chi^{(\ell/L)}(x/L,x'/L,\cdot\,)=S_L\left(\mathrm{cov}_\ell\,\chi(x,x',\cdot\,)\right).
\end{align*}
The two sides of $\mathrm{cov}_{\ell/L}\circ S_L=S_L\circ\mathrm{cov}_\ell$ are homomorphisms on $\mathrm{As}_k(\tilde U)$ which agree on the generators; therefore they are equal.

\item[(C$'$)]
Let $\ell\in\mathrm{Cl}$, and assume only $\mathrm{sdm}(\tilde U)<\ell$.
Then $\mathrm{sdm}(\tilde U/L)<\ell/L$.
By Lemma~4\,(a) the rule $\phi(t)\mapsto\phi^{(\ell')}(t)$ defines injective homomorphisms $\mathrm{cov}_\ell$ on $\mathrm{Grc}_k(\tilde U)$ and $\mathrm{cov}_{\ell/L}$ on $\mathrm{Grc}_{k+1}(\tilde U/L)$, and for these maps the computation of (C) gives $\mathrm{cov}_{\ell/L}\circ S_L=S_L\circ\mathrm{cov}_\ell$ on $\mathrm{As}_k(\tilde U)$.

\item[(D)]
Section~I.3.7 defines $S_L$ on the torus without reference to a set $\tilde U$; (C) and (C$'$) hold for every $\tilde U$ that satisfies the condition stated there.
\end{description}

\paragraph{The set $\bar{X}^{L}$} ($\mathrm{rg}=0$; no data).
For a finite union $\tilde X$ of $M$-cubes let $\mathrm{clL}(\tilde X)$ be the union of the $LM$-cubes of $\mathbb{L}_k$ that contain an $M$-cube of $\tilde X$.
The side $\ell$ does not occur in $\mathrm{clL}(\tilde X)$, and $\mathrm{clL}(\tilde X+v)=\mathrm{clL}(\tilde X)+v$ for $v\in LM\mathbb{Z}^3$ by (5.c).
Since $\ell$ is a multiple of $LM$, the map $\mathrm{pr}_\ell$ carries $LM$-cubes onto $LM$-cubes.
Let $\tilde X$ be a lift of $X$.
Every $M$-cube lies in exactly one $LM$-cube, on $\mathbb{L}_k$ as well as on $T_k$ (1.c); and the $LM$-cube of $T_k$ that contains the $M$-cube $\mathrm{pr}_\ell(\tilde\square)$, $\tilde\square$ an $M$-cube of $\tilde X$, is the image of the $LM$-cube that contains $\tilde\square$.
Hence
\[
\bar{X}^{L}=\mathrm{pr}_\ell\left(\mathrm{clL}(\tilde X)\right)\qquad\text{for every lift }\tilde X\text{ of }X.
\]

\paragraph{The operation $L$}
(an object attached to a polymer $Y$ of level $k+1$; $\mathrm{rg}=\eta_{k+1}$ at level $k+1$; the data are the elements over the lifts of the polymers $X$ with $\bar{X}^{L}=LY$).

\emph{The index set.}
Let $\tilde Y'$ be a finite non-empty union of $M$-cubes of $\mathbb{L}_{k+1}$ that is connected through faces.
The map $c'\mapsto Lc'$ carries the $M$-cube $\{c'\in\mathbb{L}_{k+1}:|c'-z'|<M/2\}$, $z'\in M\mathbb{Z}^3$, onto the $LM$-cube $\{c\in\mathbb{L}_k:|c-Lz'|<LM/2\}$, and it carries cubes with a common face ($z'_1-z'_2=\pm M e_{\mu}$) to $LM$-cubes with a common face.
Hence $L\tilde Y':=\{Lc':c'\in\tilde Y'\}$ is a finite union of $LM$-cubes of $\mathbb{L}_k$ which is connected through faces.
It is also a finite union of $M$-cubes of $\mathbb{L}_k$ which is connected through faces.
Indeed, the $LM$-cube with the centre $z$ is the union of the $L^3$ $M$-cubes with the centres $z+Ma$, $a\in\mathbb{Z}^3$, $|a|\le(L-1)/2$ (Fact~1), and these are connected through faces; moreover the $M$-cubes with the centres $z+\frac{L-1}{2}M e_{\mu}$ and $z+\frac{L+1}{2}M e_{\mu}$, which lie in the $LM$-cubes with the centres $z$ and $z+LM e_{\mu}$ respectively, have a common face.
Let $\mathrm{Ix}(\tilde Y')$ be the set of the non-empty unions $\tilde X$ of $M$-cubes of $\mathbb{L}_k$ which are contained in $L\tilde Y'$, are connected through faces and satisfy $\mathrm{clL}(\tilde X)=L\tilde Y'$.
This set is finite, and $\ell$ does not occur in it.
For $v\in LM\mathbb{Z}^3$,
\[
L\left(\tilde Y'+v/L\right)=L\tilde Y'+v,\qquad\mathrm{Ix}\left(\tilde Y'+v/L\right)=\mathrm{Ix}(\tilde Y')+v,
\]
by (5.c) and $\mathrm{clL}(\tilde X+v)=\mathrm{clL}(\tilde X)+v$.

\emph{The sum.}
Let $\ell$ and $\ell/L$ belong to $\mathrm{Cl}$, and let $Y\in D_{k+1}$ (side $\ell/L$) have the lift $\tilde Y'$.
Lemma~1 holds at every level and for every side in $\mathrm{Cl}$, hence with $(k+1,\ell/L)$ in place of $(k,\ell)$: $\tilde Y'$ is a finite non-empty union of $M$-cubes of $\mathbb{L}_{k+1}$ connected through faces, and $\mathrm{sdm}(\tilde Y')+\eta_{k+1}<\ell/L$.
Therefore
\[
\mathrm{sdm}(L\tilde Y')+\eta_k=L\left(\mathrm{sdm}(\tilde Y')+\eta_{k+1}\right)<\ell,
\]
and $\mathrm{pr}_\ell(L\tilde Y')=LY$, since $\mathrm{pr}_\ell(Lc')=L\,\mathrm{pr}_{\ell/L}(c')$.
Every polymer $X$ with $\bar{X}^{L}=LY$ is a subset of $LY$.
By Lemma~1\,(b), applied with $\tilde{\mathbf{Y}}=L\tilde Y'$, the map $\tilde X\mapsto\mathrm{pr}_\ell(\tilde X)$ is a bijection of the set of the non-empty unions of $M$-cubes contained in $L\tilde Y'$ that are connected through faces onto the set of the polymers contained in $LY$, and $\tilde X$ is a lift of $\mathrm{pr}_\ell(\tilde X)$.
For such $\tilde X$ one has $\mathrm{clL}(\tilde X)\subset L\tilde Y'$, because $L\tilde Y'$ is a union of $LM$-cubes and an $M$-cube lies in only one $LM$-cube; moreover $\overline{\mathrm{pr}_\ell(\tilde X)}^{\,L}=\mathrm{pr}_\ell\left(\mathrm{clL}(\tilde X)\right)$, and $\mathrm{pr}_\ell$ is injective on $L\tilde Y'$.
So $\overline{\mathrm{pr}_\ell(\tilde X)}^{\,L}=LY$ if and only if $\mathrm{clL}(\tilde X)=L\tilde Y'$.
By the display of Sections~I.3.7 and I.5.6 it follows that
\begin{equation}\tag{25.3}
\begin{split}
(LE)(Y;\mathscr{A},\mathscr{C},\cdot\,)&=\sum_{X:\,\bar{X}^{L}=LY}S_L\left[E(X;\mathscr{A}_L,\mathscr{C}_L,\cdot\,)\right]\\
&=\sum_{\tilde X\in\mathrm{Ix}(\tilde Y')}S_L\left[E(\mathrm{pr}_\ell\tilde X;\mathscr{A}_L,\mathscr{C}_L,\cdot\,)\right],
\end{split}
\end{equation}
where $\mathscr{A}_L$ and $\mathscr{C}_L$ are given by no.~14 of Proposition~25.1.

\begin{description}
\item[(A)]
Let $\tilde Y'$ be as in ``The index set'', and let $\mathrm{Ffm}=(F_{\tilde X})_{\tilde X\in\mathrm{Ix}(\tilde Y')}$ be a family of elements $F_{\tilde X}$ of $\mathrm{As}_k(\tilde X^{+\eta_k})$.
Put
\[
\Phi_L(\tilde Y';\mathrm{Ffm}):=\sum_{\tilde X\in\mathrm{Ix}(\tilde Y')}S_L\,F_{\tilde X}
\qquad\text{and}\qquad
\tilde U:=(L\tilde Y')^{+\eta_k}.
\]
A site $c$ of $\mathbb{L}_k$ lies in $\tilde U$ if and only if $\mathrm{dist}(c/L,\tilde Y')\le\eta_k/L=\eta_{k+1}$.
Hence
\[
\tilde U/L=\tilde Y'^{+\eta_{k+1}}:=\left\{c'\in\mathbb{L}_{k+1}:\mathrm{dist}(c',\tilde Y')\le\eta_{k+1}\right\},
\]
the neighbourhood of Section~G3.5.0, formed in $\mathbb{L}_{k+1}$.
As $\tilde X^{+\eta_k}\subset\tilde U$, every $F_{\tilde X}$ lies in $\mathrm{As}_k(\tilde U)$, and $\Phi_L(\tilde Y';\mathrm{Ffm})$ is an element of $\mathrm{As}_{k+1}(\tilde Y'^{+\eta_{k+1}})$.
The construction does not involve $\ell$.
All arguments of the data lie over $\tilde U=L(\tilde Y'^{+\eta_{k+1}})$; thus the range is $\eta_{k+1}$ at level $k+1$.

\item[(B)]
For $v\in LM\mathbb{Z}^3$ let $T_v\mathrm{Ffm}$ denote the family with $(T_v\mathrm{Ffm})_{\tilde X+v}:=T_vF_{\tilde X}$, indexed by $\mathrm{Ix}(\tilde Y')+v=\mathrm{Ix}(\tilde Y'+v/L)$.
By (B) for $S_L$,
\[
\Phi_L(\tilde Y'+v/L;T_v\mathrm{Ffm})=\sum_{\tilde X\in\mathrm{Ix}(\tilde Y')}S_L\,T_vF_{\tilde X}=T_{v/L}\sum_{\tilde X\in\mathrm{Ix}(\tilde Y')}S_L\,F_{\tilde X}=T_{v/L}\,\Phi_L(\tilde Y';\mathrm{Ffm}).
\]
Here $v/L$ ranges over $M\mathbb{Z}^3$.

\item[(C)]
Let $\ell$ and $\ell/L$ belong to $\mathrm{Cl}$.
Let $Y\in D_{k+1}$ have a lift $\tilde Y'$ with $\mathrm{sdm}(\tilde Y')+2\eta_{k+1}+M<\ell/L$; this is the margin of the first paragraph of Proposition~25.2 at level $k+1$, with the side $\ell/L$ and $\mathrm{rg}=\eta_{k+1}$.
Let $E$ be a family of even elements of $\mathrm{Gr}[\psi,\chi]$ of level $k$ on the torus of side $\ell$ with (k2), and let $(\mathscr{A},\mathscr{C})$ be a pair of bond functions on $T_{k+1}=\mathbb{L}_{k+1}/(\ell/L)\mathbb{Z}^3$ for which the elements $E(X;\mathscr{A}_L,\mathscr{C}_L)$ are defined.
Then
\[
\mathrm{sdm}(\tilde U)+LM\le\mathrm{sdm}(L\tilde Y')+2\eta_k+LM=L\left(\mathrm{sdm}(\tilde Y')+2\eta_{k+1}+M\right)<\ell .
\]
Hence (C) for $S_L$ holds on $\mathrm{As}_k(\tilde U)$; and since $\ell/L\in\mathrm{Cl}$ and $\mathrm{sdm}(\tilde U/L)+M<\ell/L$, the map $\mathrm{cov}_{\ell/L}$ on $\mathrm{Grc}_{k+1}(\tilde U/L)$ is the map of Section~G3.5.0 at level $k+1$.
For $\tilde X\in\mathrm{Ix}(\tilde Y')$ one has
\[
\mathrm{sdm}(\tilde X)+2\eta_k+M\le\mathrm{sdm}(L\tilde Y')+2\eta_k+LM<\ell,
\]
and $\tilde X$ is a lift of $\mathrm{pr}_\ell(\tilde X)$; by Section~3.5\,(C), paragraph ``The element over the lift'',
\[
E(\mathrm{pr}_\ell\tilde X;\mathscr{A}_L,\mathscr{C}_L)=\mathrm{cov}_\ell\left(\mathrm{Ecv}(\tilde X;\mathscr{A}_L,\mathscr{C}_L)\right),
\]
with $\mathrm{Ecv}(\tilde X;\mathscr{A}_L,\mathscr{C}_L)$ in $\mathrm{As}_k(\tilde X^{+\eta_k})$, which is a subalgebra of $\mathrm{As}_k(\tilde U)$.
Combined with (25.3) this yields
\begin{align*}
(LE)(Y;\mathscr{A},\mathscr{C})&=\sum_{\tilde X\in\mathrm{Ix}(\tilde Y')}S_L\,\mathrm{cov}_\ell\left(\mathrm{Ecv}(\tilde X;\mathscr{A}_L,\mathscr{C}_L)\right)
=\sum_{\tilde X\in\mathrm{Ix}(\tilde Y')}\mathrm{cov}_{\ell/L}\left(S_L\,\mathrm{Ecv}(\tilde X;\mathscr{A}_L,\mathscr{C}_L)\right)\\
&=\mathrm{cov}_{\ell/L}\left(\Phi_L\left(\tilde Y';\left(\mathrm{Ecv}(\tilde X;\mathscr{A}_L,\mathscr{C}_L)\right)_{\tilde X\in\mathrm{Ix}(\tilde Y')}\right)\right).
\end{align*}

\item[(C$'$)]
Suppose that, instead of the margin of (C), only $\mathrm{sdm}(L\tilde Y')+2\eta_k+M<\ell$ is given.
Then $\mathrm{sdm}(\tilde Y')+\eta_{k+1}<\ell/L$; so, for $\tilde Y'$ as in ``The index set'', $\mathrm{pr}_{\ell/L}(\tilde Y')$ is a polymer of level $k+1$ with the lift $\tilde Y'$, by Lemma~1\,(b) at level $k+1$, and (25.3) holds for it.
Moreover $\mathrm{sdm}(\tilde X)+2\eta_k+M<\ell$ for $\tilde X\in\mathrm{Ix}(\tilde Y')$, and $\mathrm{sdm}(\tilde U)+M<\ell$.
The computation of (C), with (C$'$) for $S_L$ in place of (C), gives the same identity, in which $\mathrm{cov}_{\ell/L}$ is now the injective homomorphism of Lemma~4\,(a) on $\mathrm{Grc}_{k+1}(\tilde U/L)$.

\item[(D)]
Section~I.5.6 defines $(LE)(Y)$ without reference to a lift, and (C) holds for every lift with the margin stated in (C).
The independence of the lift can also be checked directly.
Two lifts of $Y$ are $\tilde Y'$ and $\tilde Y'+(\ell/L)w$ with $w\in\mathbb{Z}^3$ (Lemma~1\,(b) at level $k+1$).
With $v=\ell w\in LM\mathbb{Z}^3$ one has $\mathrm{Ix}\left(\tilde Y'+(\ell/L)w\right)=\mathrm{Ix}(\tilde Y')+\ell w$; the element over the lift $\tilde X+\ell w$ of $\mathrm{pr}_\ell(\tilde X)$ is $T_{\ell w}\mathrm{Ecv}(\tilde X;\mathscr{A}_L,\mathscr{C}_L)$ (Section~3.5\,(C), paragraph ``Another lift''); and by (B) the value of $\Phi_L$ at the second lift is $T_{(\ell/L)w}$ applied to its value at the first lift.
This value is even, because $S_L$ preserves the parity; so $\mathrm{cov}_{\ell/L}$ maps both values to the same element, by Lemma~4\,(d1) with $(j,\ell')=(k+1,\ell/L)$.
\end{description}

\paragraph{Level $0$.}
The three objects of (o7) involve only the generators $\psi$ and $\chi$, which are indexed by $\mathbb{L}^{\mathrm{f}}_k$ and $\mathrm{Pzt}_k$ at every level.
At $k=0$, where $\eta_0=1$ and $\eta_1=1/L$, all assertions of this subsection hold as stated.

\subsubsection*{3.10. (o8)}

\begin{description}
\item[(A)]
Let $\tilde B$ be given on the bonds contained in $\tilde X$, and extend it arbitrarily.

\emph{What the kernels read.}
Let $y_1$ and $y_2$ be unit sites of $\tilde X$ (for $k=0$: $c_1,c_2\in\tilde X$).
By (K2) of Proposition~25.1, $P(\tilde B)(y_1,y_2)=0$ unless $[y_1]^L=[y_2]^L=:y'$; in that case it reads $\tilde B$ on $\mathrm{Pth}(y',y_1)\cup\mathrm{Pth}(y',y_2)$, a set of bonds with both ends in $\mathrm{Bl}_{k+1}(y')$ (Fact~4).
By (1.c) the block $\mathrm{Bl}_{k+1}(y')$ lies in the $M$-cube of $y_1$, hence in $\tilde X$.
So $P(\tilde B)(y_1,y_2)$ reads $\tilde B$ on bonds contained in $\tilde X$; by (K3) the same holds for $B_{k,y}(\tilde B)(y_1,y_2)$ and for $\left(B_{k,y}(\tilde B)\,\gamma_5\,B_{k,y}(\tilde B)\right)(y_1,y_2)$.
For $c\in\Delta_{y_1}$ the phase $\varphi_{\tilde B}(y_1,c)$ reads $\tilde B$ on $\mathrm{Twr}(c)$, whose bonds are contained in $\Delta_{y_1}$, and $\Delta_{y_1}$ lies in $\tilde X$ (Fact~3, (1.c)).
By (K4), for $x,x'$ over $\tilde X$ the kernel $F_0(\tilde B)(x,x')$ contains $\tilde B$ through $P(\tilde B)(c,c')$ and, if $x'=x^{(3)}$, through $u_{\tilde B}(c,c_3)$, where $c$ and $c_3=c'$ lie in $\tilde X$.
Consequently the following numbers are functions of the data named, and $\ell$ does not occur in them.
For a kernel $O(x,x')$ ($4\times4$ matrices) at the fermion sites over $\tilde X$ and for real $y$, put, if $k\ge1$,
\begin{multline*}
\mathrm{Trk}_k(\tilde X;\tilde B,O,y):=\sum_{y_1,y_2\in\tilde X\cap\mathbb{Z}^3}\operatorname{tr}\left[\left(B_{k,y}(\tilde B)\,\gamma_5\,B_{k,y}(\tilde B)\right)(y_2,y_1)\right.\\
\left.\cdot\left(\frac{\eta_k^3}{48}\right)^{2}\sum_{x\text{ over }\Delta_{y_1}}\ \sum_{x'\text{ over }\Delta_{y_2}}\varphi_{\tilde B}(y_1,c)\cdot O(x,x')\cdot\varphi_{\tilde B}(y_2,c')^{-1}\right],
\end{multline*}
and, if $k=0$,
\begin{align*}
\mathrm{Trk}_0'(\tilde X;\tilde B,O)&:=\left(\frac{1}{48}\right)^{2}\sum_{x,x'\text{ over }\tilde X}\operatorname{tr}\left[F_0(\tilde B)(x,x')\cdot O(x',x)\right],\\
\mathrm{Trk}_0''(\tilde X;O)&:=\frac{1}{48}\sum_{x\text{ over }\tilde X}\operatorname{tr}\left[\gamma_5\cdot O(x,x)\right].
\end{align*}

\item[(B)]
Let $v\in M\mathbb{Z}^3$; note that $M\mathbb{Z}^3\subset L\mathbb{Z}^3$.
By (l4) for the nos.~4, 6 and 8 of the table of Proposition~25.1 the kernels satisfy $\mathrm{Op}_{\mathbb{L}}(T_v\tilde B)(t+v,t'+v)=\mathrm{Op}_{\mathbb{L}}(\tilde B)(t,t')$; moreover $(T_vO)(x+v,x'+v)=O(x,x')$ and $\Delta_y+v=\Delta_{y+v}$.
Hence the three numbers do not change.

\item[(C)]
Let $X$ have the lift $\tilde X$ with $\mathrm{sdm}(\tilde X)+M<\ell$, let $B$ be a bond function on $T_k$ with the periodic extension $\tilde B_{\mathrm{p}}$, and let $O$ be an operator on $\mathfrak{S}^{\mathrm{ch}}(\eta_k,\ell)$ whose kernel vanishes unless both arguments lie over $X$.
Choose fundamental domains of $\ell\mathbb{Z}^3$ in $\mathbb{L}^{\mathrm{f}}_k$ and in $\mathbb{Z}^3$ that contain the sites over $\tilde X$, respectively the unit sites of $\tilde X$; this is possible by Lemma~1\,(a).
A point of such a fundamental domain that lies in (over) $\tilde X+\ell\mathbb{Z}^3$ lies in (over) $\tilde X$, because it is the only representative of its class.
So $O(x,x')=0$ for $x,x'$ in the fundamental domain unless both lie over $\tilde X$.
Since $L-1<M$, one has $\mathrm{sdm}(\tilde X)+(L-1)<\ell$.
By Lemma~3\,(c), the kernels on the torus of $P(B)$, of $1$ and of $F_0(B)$, at indices in (over) $\tilde X$, are the kernels of the corresponding operators on the covering lattice at $\tilde B_{\mathrm{p}}$; hence so is the kernel of $B_{k,y}(B)\,\gamma_5\,B_{k,y}(B)$, by (K3), which holds on the torus by Section~2.4\,(ii), (iii).

\emph{The case $k\ge1$.}
Put $\mathrm{Oav}:=Q_k(B)J^{*}\,O\,JQ_k^{T}(-B)$, an operator on $\mathfrak{S}(1,\ell)$.
For $v\in\mathfrak{S}(1,\ell)$ and $x\in\mathbb{L}^{\mathrm{f}}_k$,
\begin{align*}
\left(O\,JQ_k^{T}(-B)\,v\right)(x)&=\int\mathrm{d}x'\,O(x,x')\cdot\varphi_{\tilde B_{\mathrm{p}}}([c'],c')^{-1}\cdot v([c'])\\
&=\sum_{y_2\in\tilde X\cap\mathbb{Z}^3}\left[\int_{\Delta_{y_2}}\mathrm{d}x'\,O(x,x')\cdot\varphi_{\tilde B_{\mathrm{p}}}(y_2,c')^{-1}\right]\cdot v(y_2),
\end{align*}
where $x'$ ranges over the fundamental domain in the first integral; the second equality holds because $\tilde X$ is the disjoint union of the $\Delta_{y_2}$ (1.c).
Applying $Q_k(B)J^{*}$ by (Of4), one finds
\[
\mathrm{Oav}(y_1,y_2)=\int_{\Delta_{y_1}}\mathrm{d}x\int_{\Delta_{y_2}}\mathrm{d}x'\,\varphi_{\tilde B_{\mathrm{p}}}(y_1,c)\cdot O(x,x')\cdot\varphi_{\tilde B_{\mathrm{p}}}(y_2,c')^{-1}
\]
for $y_2\in\tilde X\cap\mathbb{Z}^3$, and $\mathrm{Oav}(y_1,y_2)=0$ for the other $y_2$ of the fundamental domain.
Moreover, for $y_1$ in the fundamental domain, $\mathrm{Oav}(y_1,y_2)=0$ unless $y_1\in\tilde X$.
Indeed, otherwise $y_1$ lies outside $\tilde X+\ell\mathbb{Z}^3$; then the $M$-cube of $y_1$, which contains $\Delta_{y_1}$ by (1.c), is not one of the cubes of the union $\tilde X+\ell\mathbb{Z}^3$ of $M$-cubes, and no $x$ over $\Delta_{y_1}$ lies over $X$.
On the finite-dimensional space $\mathfrak{S}(1,\ell)$ the trace is cyclic, and $\operatorname{Tr}O'=\sum_{z}\operatorname{tr}O'(z,z)$, the sum extending over the fundamental domain.
Therefore
\begin{align*}
\operatorname{Tr}\left[\gamma_5\,B_{k,y}(B)\,\mathrm{Oav}\,B_{k,y}(B)\right]
&=\operatorname{Tr}\left[B_{k,y}(B)\,\gamma_5\,B_{k,y}(B)\,\mathrm{Oav}\right]\\
&=\sum_{y_1,y_2\in\tilde X\cap\mathbb{Z}^3}\operatorname{tr}\left[\left(B_{k,y}(B)\,\gamma_5\,B_{k,y}(B)\right)(y_2,y_1)\cdot\mathrm{Oav}(y_1,y_2)\right]\\
&=\mathrm{Trk}_k(\tilde X;\tilde B,O,y),
\end{align*}
where $\tilde B$ is the restriction of $\tilde B_{\mathrm{p}}$ to the bonds contained in $\tilde X$.

\emph{The case $k=0$.}
By (I.ke2) and the composition of kernels (I.1.4),
\[
\operatorname{Tr}\left[F_0(B)\,O\right]=\int\mathrm{d}x\,\operatorname{tr}\left(F_0(B)\,O\right)(x,x)=\int\mathrm{d}x\int\mathrm{d}x'\,\operatorname{tr}\left[F_0(B)(x,x')\cdot O(x',x)\right];
\]
only $x,x'$ over $\tilde X$ contribute, and for these $F_0(B)(x,x')=F_0(\tilde B_{\mathrm{p}})(x,x')$.
This is $\mathrm{Trk}_0'(\tilde X;\tilde B,O)$.
The kernel of $\gamma_5O$ is $\gamma_5\cdot O(x,x')$; hence $\operatorname{Tr}\left[\gamma_5O\right]=\int\mathrm{d}x\,\operatorname{tr}\left[\gamma_5\cdot O(x,x)\right]=\mathrm{Trk}_0''(\tilde X;O)$.

\emph{The numbers $\acute{E}_k$, $E^{d}_k$, $E_k^{\mathrm{det}}$, ${\delta E^{\mathrm{det}}}_k$.}
Let the data be (SF) (Section~I.1.8: $\ell\in\mathrm{Cl}$, $\mu\in(0,1]$, $e_k\in(0,e_0]$), let $B\in\tilde{G}_k$, and let $S_{k,y}(X,A)$, $X\in D_k$, $A\in\tilde{G}_k$, $y\ge0$, be a family (p5) of Section~I.5.8 as in the data of Proposition~13.
This is the setting in which Section~I.5.7 (paragraph ``Settings'') forms $\acute{E}_k$, $E^{d}_k$, $E_k^{\mathrm{det}}$ and ${\delta E^{\mathrm{det}}}_k$.
Together with $B$, also $-B$, $\overline{B}$ and $-\overline{B}$ lie in $\tilde{G}_k$, and $0\in\tilde{G}_k$.
Indeed, let $B=A_0+A_1$ be a splitting as in Section~I.5.3: $A_1\in G_k$, $A_0$ is real, and for every $M$-cube $\square$ one has $A_0=A'+d\boldsymbol{\lambda}$ on $\square^{\sim 5}$ with real $\boldsymbol{\lambda}$ and $A'$ and with $|A'|_\infty<e_k^{-1+\varepsilon}$ on $\square^{\sim 5}$.
Then $-B=(-A_0)+(-A_1)$ with $-A_0=(-A')+d(-\boldsymbol{\lambda})$, and $\overline{B}=A_0+\overline{A_1}$, with the same moduli.
Among the properties of the datum (p5) is the first half of W12 (Section~G1.11): the kernel of $S_{k,y}(X,B)$ vanishes unless both arguments lie over $X$.
Hence, for every $y\ge0$, the operator $O=S_{k,y}(X,B)$ satisfies the hypothesis made on $O$ at the beginning of (C), and the integrands of Section~I.5.7 are
\begin{gather*}
\operatorname{Tr}\left[\gamma_5\,B_{k,y}(B)\,Q_k(B)J^{*}\,S_{k,y}(X,B)\,JQ_k^{T}(-B)\,B_{k,y}(B)\right]=\mathrm{Trk}_k(\tilde X;\tilde B,O_y,y)\qquad(k\ge1),\\
\operatorname{Tr}\left[F_0(B)\,S_{0,y}(X,B)\right]=\mathrm{Trk}_0'(\tilde X;\tilde B,O_y),\qquad
\operatorname{Tr}\left[\gamma_5\,S_{0,y}(X,B)\right]=\mathrm{Trk}_0''(\tilde X;O_y)\qquad(k=0),
\end{gather*}
where $O_y$ is the kernel of $S_{k,y}(X,B)$ at the arguments over $\tilde X$.
Since the integrands agree for every $y$, it follows that
\begin{align*}
\acute{E}_k(X,B)&=-\mathrm{i}\,b_{k}^{2}\int_0^\infty\mathrm{d}y\,\mathrm{Trk}_k(\tilde X;\tilde B,O_y,y)\qquad(k\ge1),\\
\acute{E}_0(X,B)&=\int_1^\infty\frac{\mathrm{d}y}{y}\,\mathrm{Trk}_0'(\tilde X;\tilde B,O_y)-\mathrm{i}\int_0^1\mathrm{d}y\,\mathrm{Trk}_0''(\tilde X;O_y),
\end{align*}
in the sense in which the integrals of Section~I.5.7 are taken.
The right-hand sides are the values of one functional of $(\tilde X;\tilde B;(O_y)_{y\ge0})$ in which $\ell$ does not occur.
The lifts of $-B$, $\overline{B}$ and $-\overline{B}$ are $-\tilde B$, $\overline{\tilde B}$ and $-\overline{\tilde B}$.
The number $E^{d}_k(X,B)$ is the mean of $\acute{E}_k(X,B)$, $\acute{E}_k(X,-B)$, $\overline{\acute{E}_k(X,\overline{B})}$ and $\overline{\acute{E}_k(X,-\overline{B})}$, each formed with the kernels of $S_{k,y}(X,\cdot\,)$ at the connection named; and
\[
E_k^{\mathrm{det}}(X,B)=E^{d}_k(X,B)-E^{d}_k(X,0),\qquad
{\delta E^{\mathrm{det}}}_k(X;\mathscr{A},\mathscr{C})=E^{d}_k(X,\mathscr{A}+\mathscr{C})-E^{d}_k(X,\mathscr{C})
\]
(for $\mathscr{A}+\mathscr{C}$ and $\mathscr{C}$ in $\tilde{G}_k$) are differences of values of $E^{d}_k$ (Section~I.5.7).

\item[(D)]
The values are numbers.
\end{description}

\subsubsection*{3.11. (o9), (o10)}

\paragraph{(o9).}
Let $\tilde{U}$ be a finite set of sites of $\mathbb{L}_k$ with $\mathrm{sdm}(\tilde{U})+M<\ell$.
Put $\nu:=1$ for $k\ge1$ and $\nu:=1/48$ for $k=0$.
A \emph{mode} is a pair $(t,\beta)$ with $\beta\in\{1,\dots,4\}$, where $t$ is a unit site of $\tilde{U}$ if $k\ge1$, and $t=x$ is a fermion site over $\tilde{U}$ if $k=0$.
To a mode belong the two generators $W_{\beta}(t)=W(t,\beta,0)$ and $\bar{W}_{\beta}(t)=W(t,\beta,1)$.
Denote by $\mathrm{IW}_c$ the set of the modes of $\mathrm{Grc}_k(\tilde{U})$, and by $\mathrm{IW}_\ell$ the set of the modes of the torus, with $t$ in a fundamental domain.
With each mode associate
\[
\mathrm{qm}:=W_{\beta}(t)\cdot\bar{W}_{\beta}(t),
\]
an even element with $\mathrm{qm}^2=0$.
Then $-\nu\cdot\bar{W}_{\beta}(t)\cdot W_{\beta}(t)=\nu\cdot\mathrm{qm}$, and, with all products extending over the modes (Section~I.3.6),
\[
\exp\left(-\langle\bar{W},W\rangle_{\nu}\right)=\prod\left(1+\nu\cdot\mathrm{qm}\right),\qquad \mathrm{Top}=\prod\mathrm{qm}.
\]

\emph{The coefficient of $\mathrm{Top}$.}
Let $\mathrm{Gr}$ be a Grassmann algebra, $\mathrm{IW}$ a finite set of modes among its generators, and $\mathrm{Gr}'$ the subalgebra generated by the remaining generators.
For a subset $\mathrm{Jw}$ of the set of the generators of the modes of $\mathrm{IW}$ let $W(\mathrm{Jw})$ denote the product of its members in a fixed order.
Each $F\in\mathrm{Gr}$ can be written in exactly one way as
\[
F=\sum_{\mathrm{Jw}}F_{\mathrm{Jw}}\cdot W(\mathrm{Jw}),\qquad F_{\mathrm{Jw}}\in\mathrm{Gr}',
\]
where $\mathrm{Jw}$ ranges over these subsets.
When $\mathrm{Jw}$ is the full set, $W(\mathrm{Jw})=\boldsymbol{\sigma}\cdot\mathrm{Top}$ with a sign $\boldsymbol{\sigma}$, and the coefficient of $\mathrm{Top}$ in $F$ is $\boldsymbol{\sigma}\cdot F_{\mathrm{Jw}}$ for this $\mathrm{Jw}$.
By (I.3.5), $\int\mathrm{D}W\,F$ equals $\nu^{-\#|\mathrm{IW}|}$ times this coefficient, and
\[
\int\mathrm{d}\mu_I(W)\,F=\int\mathrm{D}W\,\prod\left(1+\nu\cdot\mathrm{qm}\right)\cdot F.
\]
\begin{itemize}
\item[(A)]
Let $\int_c$ be this integral for $\mathrm{Gr}=\mathrm{Grc}_k(\tilde{U})$ and $\mathrm{IW}=\mathrm{IW}_c$.
It is formed with $\nu$ and with the anticommutation relations; the side $\ell$ does not enter it.
Let $W(\mathrm{Jw})$ be a monomial.
Among the terms of $\prod\left(1+\nu\cdot\mathrm{qm}\right)\cdot W(\mathrm{Jw})$ there is a multiple of $\mathrm{Top}$ if and only if $\mathrm{Jw}$ is a union of modes, that is, the set of all generators of some set of modes.
In that case $W(\mathrm{Jw})=\pm\prod_{\text{modes in }\mathrm{Jw}}\mathrm{qm}$, the sign being the sign of the reordering; with $n$ the number of these modes, the term is $\pm\nu^{\#|\mathrm{IW}_c|-n}\cdot\mathrm{Top}$, and
\[
\int_c W(\mathrm{Jw})=\pm\nu^{-n}.
\]
These are the moments (I.B4).
\item[(B)]
Let $v\in\mathbb{Z}^3$.
The translation $T_v$ maps the modes of $\tilde{U}$ onto those of $\tilde{U}+v$, with unchanged weights.
Hence $T_v\int_c F=\int_c T_vF$.
\item[(C)]
By Lemma~4\,(a), $\mathrm{cov}_\ell$ induces an injection of $\mathrm{IW}_c$ into $\mathrm{IW}_\ell$; call its image $\mathrm{IW}_j$.
The element $\mathrm{cov}_\ell(\mathrm{qm})=W^{(\ell)}_{\beta}(t)\cdot\bar{W}^{(\ell)}_{\beta}(t)$ is the element $\mathrm{qm}$ of the mode of the torus in the class of $(t,\beta)$: by (I.AP), a change of the representative multiplies both factors by the same sign.
For $F\in\mathrm{Grc}_k(\tilde{U})$ set
\[
G:=\mathrm{cov}_\ell\left(F\right),\qquad G':=\prod_{\mathrm{IW}_j}\left(1+\nu\cdot\mathrm{qm}\right)\cdot G=\mathrm{cov}_\ell\left(\prod_{\mathrm{IW}_c}\left(1+\nu\cdot\mathrm{qm}\right)\cdot F\right).
\]
No generator of a mode outside $\mathrm{IW}_j$ occurs in $G'$.
Expanding $\prod_{\mathrm{IW}_\ell\setminus\mathrm{IW}_j}\left(1+\nu\cdot\mathrm{qm}\right)\cdot G'$, one finds a single term that contains all generators of the modes outside $\mathrm{IW}_j$, namely $\prod_{\mathrm{IW}_\ell\setminus\mathrm{IW}_j}\left(\nu\cdot\mathrm{qm}\right)\cdot G'$.
Therefore the coefficient of $\mathrm{Top}$, formed for $\mathrm{IW}_\ell$, in $\prod_{\mathrm{IW}_\ell}\left(1+\nu\cdot\mathrm{qm}\right)\cdot G$ is $\nu^{\#|\mathrm{IW}_\ell|-\#|\mathrm{IW}_j|}$ times the coefficient of
\[
\mathrm{Top}_j:=\prod_{\mathrm{IW}_j}\mathrm{qm}
\]
in $G'$, and $\int\mathrm{d}\mu_I(W)\,G$ is $\nu^{-\#|\mathrm{IW}_j|}$ times the coefficient of $\mathrm{Top}_j$ in $G'$.
Let $\mathrm{Gr}'$ be formed in $\mathrm{Grc}_k(\tilde{U})$ with $\mathrm{IW}=\mathrm{IW}_c$, and $\mathrm{Top}_c:=\prod_{\mathrm{IW}_c}\mathrm{qm}$.
The map $\mathrm{cov}_\ell$ sends $\mathrm{Gr}'$ into the subalgebra generated by the generators outside $\mathrm{IW}_j$; it sends each $W(\mathrm{Jw})$ to $\pm$ the product of the images of the members of $\mathrm{Jw}$, and different sets $\mathrm{Jw}$ to different sets; and $\mathrm{cov}_\ell(\mathrm{Top}_c)=\mathrm{Top}_j$.
Consequently, for every $H\in\mathrm{Grc}_k(\tilde{U})$, the coefficient of $\mathrm{Top}_j$ in $\mathrm{cov}_\ell(H)$ is the image under $\mathrm{cov}_\ell$ of the coefficient of $\mathrm{Top}_c$ in $H$.
Taking $H=\prod_{\mathrm{IW}_c}\left(1+\nu\cdot\mathrm{qm}\right)\cdot F$ and using $\#|\mathrm{IW}_j|=\#|\mathrm{IW}_c|$, one obtains
\[
\int\mathrm{d}\mu_I(W)\,\mathrm{cov}_\ell\left(F\right)=\mathrm{cov}_\ell\left(\int_c F\right).
\]
\item[(D)]
The map $\int_c$ commutes with $\mathrm{Par}$, because $\mathrm{Par}$ fixes every $\mathrm{qm}$ and the full set of generators of the modes of $\mathrm{IW}_c$ has an even number of members.
\end{itemize}

\paragraph{(o10).}
Let $p_{0,k}>0$.
(By Section~I.6.1, $p_{0,k}=v_k^{p_0}$ with $v_k=-\log e_k$ and an integer $p_0$; thus $p_{0,k}>0$ whenever $e_k<1$.)
The space $V_k$ carries the counting scalar product, $\mu_I$ is the standard Gaussian measure on $V_k$, $\mathsf{Z}(\Upsilon)=\langle\mathsf{Z},\Upsilon\rangle$, and the definitions of \cite[Section~5.2.5]{QEDI} read
\[
\mathrm{Box}_k=\left\{\mathsf{Z}:\ |\mathsf{Z}(\Upsilon)|\le p_{0,k}\text{ for all }\Upsilon\right\},\qquad
\mathrm{d}\hat{\mu}=\frac{\mathbf{1}_{\mathrm{Box}_k}}{\mu_I(\mathrm{Box}_k)}\,\mathrm{d}\mu_I,\qquad
\delta\varepsilon^{0}_{k}=-n_k^{-1}\cdot\log\mu_I(\mathrm{Box}_k).
\]
By \cite[Proposition~15.21(ii)]{QEDI}, $\{\Upsilon\}$ is an orthonormal basis of $V_k$ with $d_V\cdot n_k$ members, where $d_V=2-2L^{-3}$.
The map $\mathsf{Z}\mapsto\left(\mathsf{Z}(\Upsilon)\right)_{\Upsilon}$ is therefore an isometry of $V_k$ onto $\mathbb{R}^{d_V\cdot n_k}$; it carries $\mu_I$ to the product of $d_V\cdot n_k$ copies of the standard normal law, and $\mathrm{Box}_k$ to the product of the intervals $[-p_{0,k},p_{0,k}]$.
Hence, with $\mathrm{gn}$ a standard normal variable,
\[
\mu_I(\mathrm{Box}_k)=P\left(|\mathrm{gn}|\le p_{0,k}\right)^{d_V\cdot n_k}>0;
\]
the measure $\hat{\mu}$ is the product, over the basis, of copies of the standard normal law restricted to $[-p_{0,k},p_{0,k}]$ and normalised, and this law is a probability measure on $\mathbb{R}$ that depends on $p_{0,k}$ alone; and
\[
\delta\varepsilon^{0}_{k}=-n_k^{-1}\cdot d_V\cdot n_k\cdot\log P\left(|\mathrm{gn}|\le p_{0,k}\right)=-d_V\cdot\log P\left(|\mathrm{gn}|\le p_{0,k}\right).
\]

\subsubsection*{3.12. Completeness}

The clause ``Completeness'' of Proposition~25.2 is verified line by line on the displays to which it refers.
These are the displays (8.1)$^{\mathrm{ch}}$--(8.6)$^{\mathrm{ch}}$ in (D2)--(D6) of Definition~24.D;
the displays of Section~I.5.6 that define $\mathrm{al}_{0}$, $\mathrm{al}_{2}$ and the operations $\varepsilon$, $m$, $R$, $\mathrm{nat}$, $L$, together with (I.5.3);
and, of Section~I.5.7, the displays of $\acute{E}_k$, $E^{d}_k$, $E_k^{\mathrm{det}}$, ${\delta E^{\mathrm{det}}}_k$ in the paragraph ``Determinant'', together with the radii $\mathrm{rad}_1$, $\mathrm{rad}_2$, $\mathrm{rad}_4$ of the paragraph ``Step map''.
The other displays of Section~I.5.7 lie outside the range of the clause.
The following objects are set aside:
the primitives (the members of (p1)--(p5), and at $k=0$ also the diagonal families $H_0(s')$, $H_0(s'',B)$ of Section~I.5.8);
the data $(m,\mathscr{E})$, $\varepsilon^{\mathrm{v}}$ and the family $\mathfrak{P}_k$;
the connections $A$, $\mathscr{C}$, $\mathscr{T}$, $\mathscr{A}$ together with their pointwise sums and their multiples by numbers ($A+\mathscr{C}$, $\mathscr{T}+t\,\mathscr{Z}(s')$, $A+t\,\mathscr{Z}(s')$, $B+\mathrm{Zf}$, $\mathscr{A}+\mathscr{C}$);
the parameters $s$, $t$, $y$;
the variables $\mathsf{Z}$, $W$;
and the generators.
Also set aside are the following operations of a Grassmann algebra over $\mathbb{C}$, called the \emph{algebra operations} below: finite sums, products, multiples by numbers, $\exp$, the application of a row to a family of generators (a finite linear combination), and the homomorphism $\mathrm{ad}_v$ of Lemma~4\,(e), which is the homomorphism of (D1) of Definition~24.D.
The two remaining operations of this kind, the series of (8.4) and the part $F^{[2,0]}$, are named in the entries of the lines in which they occur.
Inspection of the displays of this range shows that every symbol in a line is either one of the objects set aside or one of the objects named in the entry of that line in the following list.
\begin{description}
\item[Display (8.1)$^{\mathrm{ch}}$]
\begin{itemize}
\item $\mathscr{Z}(s')=H_k(s')\cdot C\cdot C_k^{1/2}(s')\,\mathsf{Z}$: the map $C$ (no.~12 of the table of Proposition~25.1); algebra operations.
\item $H_k^{\#}(s'',B)\,\Phi$: $\delta_{\alpha,B}$, $\delta_{\alpha,-B}$ (no.~11); $\mathsf{C}$, $\mathsf{C}^{-1}$ (no.~1) and $B\mapsto-B$ in $\bar{H}_k(s'',B)$; algebra operations.
\item $J_k^{\#}$, $\tilde{\Gamma}_k$, $W_k^{\#}$: algebra operations (at $k=0$ the rows of $\Gamma_0(s'',B)$ are applied with $\int\mathrm{d}x'=\frac{1}{48}\sum$).
\end{itemize}
\item[Display (8.2)]
\begin{itemize}
\item $\mathrm{Loc}[G](\mathbf{Y})$: $P_0(X)$, $\mathrm{Yp}(X,P)$, $Y'(X,P)$; the sum over the pairs $(X,P)$ with $Y'(X,P)=\mathbf{Y}$; the values $1$ and $0$ of $s$; $\int_{[0,1]^P}\mathrm{d}s_P\,\partial/\partial s_P$; the distinction between a single $M$-cube and the other polymers (last sentence of (D3) of Definition~24.D: a function $G(\,\cdot\,,\square)$ is set equal to $0$ on the other polymers).
\end{itemize}
\item[Display (8.3)$^{\mathrm{ch}}$, $k\ge1$]
\begin{itemize}
\item $G^{(1)}$: $\oint_{\mathrm{rad}_1}$; (o2) with $O=M_k(X,\mathscr{T}+t\,\mathscr{Z}(s'))$.
\item $G^{(2)}$, $G^{(4)}$: $\oint_{\mathrm{rad}_2}$, $\oint_{\mathrm{rad}_4}$; $\mathscr{E}^{r}=R\,\mathscr{E}$ (the line $R$ below); algebra operations ($\mathrm{ad}_v$ with $v=t\,J_k^{\#}\Psi$ and with $v=t\,W_k^{\#}$).
\item $G^{(3)}$, $G^{(5)}$: $\oint_{\mathrm{rad}_2}$, $\oint_{\mathrm{rad}_4}$; $m'=m+m(\mathscr{E})$ (the line $\varepsilon$, $m$ below); (o1) with $\mathrm{Mx}=1$ and $\tilde{X}=\tilde{\square}$; algebra operations ($\mathrm{ad}_v$).
\item $E^{(i),*}$, $E^{(\le5),x}$: (8.2).
\item $E^{(i),x}(\mathbf{Y})$: (o5).
\item $G^{(\le3)}$: algebra operations ($\mathrm{ad}_v$ with $v=W_k^{\#}(s'',\mathscr{T})$).
\end{itemize}
\item[Display (8.3)$^{\mathrm{ch}}$, $k=0$]
\begin{itemize}
\item $G^{(1)}$: (o6) at $\mathrm{Zf}=\mathscr{Z}(s')$.
\item $G^{(\le3)}$, $G^{(5)}$: as for $k\ge1$, with $m$ in the place of $m'$.
\item $E^{(1),*}=E^{(1),x}=E^{(\le3),x}$, $E^{(\le5),x}$: (8.2).
\end{itemize}
\item[Display (8.4)]
\begin{itemize}
\item $\mathrm{Km}(X)$: the sets $\{X_i\}$ of distinct polymers with union $X$; algebra operations ($\exp(\,\cdot\,)-1$, products).
\item $\mathrm{Hm}(X)$: (o9), (o10).
\item $E^{\#,x}_k$: the families $(X_1,\dots,X_n)$ with union $X$; the numbers $1/n!$; $\mathrm{urs}^{T}$; the series over $n$; algebra operations.
\end{itemize}
\item[Display (8.6)$^{\mathrm{ch}}$]
\begin{itemize}
\item $E^{*,x}_k$, $L\,R\,\mathscr{E}$, $L\,{\delta E^{\mathrm{det}}}_k$, $\mathfrak{P}_{k+1}$: (o7) with no.~14 ($\mathscr{A}_L$, $\mathscr{C}_L$; $\psi_L$ and $\chi_L$ are given by $S_L$); the lines $R$, $\mathrm{nat}$, $\acute{E}_k$ below.
\item $\varepsilon^{\mathrm{v},+}$, $m^{+}$: the numbers $L^3$, $L$; $\varepsilon(\mathscr{E})$, $m(\mathscr{E})$ (the line $\varepsilon$, $m$ below); $\delta\varepsilon^{0}_{k}$ ((o10)); here $\varepsilon^{\mathrm{v}}$ and $m$ belong to the data $(m,\mathscr{E})$, $\varepsilon^{\mathrm{v}}$ set aside above.
\end{itemize}
\item[Operations of Section~I.5.6]
\begin{itemize}
\item $R$: the distinction small / large; (o3); (o1) ($\operatorname{Vol}(X)$, and the form with $\mathrm{Mx}=\mathrm{al}_{2}$).
\item $\varepsilon$, $m$: (o3); the sum over the small polymers that contain one $M$-cube $\square$.
\item $\mathrm{nat}$: the distinction small / large; $F\mapsto F^{[2,0]}$; (o4).
\item $L$: (o7).
\end{itemize}
\item[Displays of Section~I.5.7]
\begin{itemize}
\item $\acute{E}_k$: (o8), which contains nos.~1, 4, 6, 8; the integrations over $y$ with the factors $-i\,b_k^2$, $1/y$ and $-i$.
\item $E^{d}_k$, $E_k^{\mathrm{det}}$, ${\delta E^{\mathrm{det}}}_k$: $B\mapsto-B$, $B\mapsto\overline{B}$, complex conjugation, the mean of four numbers, differences.
\end{itemize}
\end{description}
The matrices $M_i(x)$ occur inside (o5), (o6) and $F_0$, and $\gamma_5$ occurs inside (o8).
The radii $\mathrm{rad}_1$, $\mathrm{rad}_2$, $\mathrm{rad}_4$ are formed with $e_k$ and $\varepsilon$ (Section~I.5.7), $\oint_{\mathrm{rad}}$ is the functional of Section~I.5.2, and $\mathrm{urs}^{T}(X_1,\dots,X_n)$ is a function of the graph of the incompatibility relations \cite[Section~13.3]{QEDI}.
None of these objects contains $\ell$.

\emph{The side $\ell$ in the entries of the list that are neither rows of the table of Proposition~25.1, nor objects (o1)--(o10), nor the index sets of (8.2) and (8.4).}
\begin{description}
\item[Algebra operations] $\mathrm{cov}_\ell$ commutes with these operations (Lemma~4\,(e)); a row applied to a family of generators is a linear combination of them, and $\mathrm{cov}_\ell$ maps it to the same combination of the generators of the torus.
\item[{$B\mapsto-B$, and the pointwise sums and multiples of the connections set aside}] the lift $B\mapsto\left(b\mapsto B(\mathrm{pr}_\ell b)\right)$ is linear, and so it commutes with them.
\item[{$F\mapsto F^{[2,0]}$}] Lemma~4\,(e).
\item[The series of (8.4)] $\mathrm{cov}_\ell$ is a linear map between finite-dimensional spaces, hence continuous; it therefore maps the sum of a convergent series to the sum of the series of the images.
\item[{$\int_{[0,1]^P}\mathrm{d}s_P\,\partial/\partial s_P$ and $\oint_{\mathrm{rad}}$}] these are a derivative and integrals in auxiliary parameters of functions with values in one finite-dimensional complex vector space ((D3) of Definition~24.D), and they are formed coefficient by coefficient with respect to a basis of that space.
Since $\mathrm{cov}_\ell$ is a linear map between finite-dimensional spaces, it commutes with them.
They contain $\ell$ at most through the set $P$, which belongs to the index sets of (8.2).
\item[The numbers] the values $1$ and $0$ of $s$ and the numbers $1/n!$, $L^3$, $L$ do not contain $\ell$.
\item[Small / large] a polymer with a lift $\tilde{X}$ is small if and only if $d_M(\tilde{X})\le L$ (Lemma~1\,(c)).
\item[Single $M$-cube / other polymer] a polymer with a lift $\tilde{X}$ is a single $M$-cube if and only if $|\tilde{X}|_M=1$, since $|\mathrm{pr}_\ell(\tilde{X})|_M=|\tilde{X}|_M$ (Lemma~1\,(b)).
\item[{$\varepsilon$, $m$}] by Lemma~1\,(c), the small polymers that contain $\mathrm{pr}_\ell(\tilde{\square})$ are the sets $\mathrm{pr}_\ell(\tilde{X})$, $\tilde{X}\in\mathrm{Sm}(\tilde{\square})$, and $\mathrm{Sm}(\tilde{\square})$ does not contain $\ell$; hence
\[
\varepsilon\left(E\right)=-\sum_{\tilde{X}\in\mathrm{Sm}(\tilde{\square})}\mathrm{al}_{0},\qquad m\left(E\right)\cdot\mathbf{1}_4=-\sum_{\tilde{X}\in\mathrm{Sm}(\tilde{\square})}\mathrm{al}_{2},
\]
with the numbers of (o3) formed at the kernels of $E(\mathrm{pr}_\ell\tilde{X};0,0)$ over $\tilde{X}$.
\item[{$\delta\varepsilon^{0}_{k}$}] (o10).
\item[The integrations over $y$] (o8).
\item[{The sums over $\{X:\ \bar{X}^{L}=LY\}$ in (8.6)$^{\mathrm{ch}}$ and in $L$}] for a polymer $Y$ of level $k+1$ with a lift $\tilde{Y}'$ they range over the sets $\mathrm{pr}_\ell(\tilde{X})$, $\tilde{X}\in\mathrm{Ix}(\tilde{Y}')$, and $\mathrm{Ix}(\tilde{Y}')$ does not contain $\ell$ ((25.3)); in $\mathfrak{P}_{k+1}$ the summands are numbers.
\end{description}
This completes the proof of Proposition~25.2.\qed

\subsection*{4. Proposition 25.3}

\subsubsection*{4.1. The statement}

Proposition~25.3 is stated in Part~II, for the data of Section~G3.5.0, two sides $\ell_1,\ell_2\in\mathrm{Cl}$ and the tori $T_i:=\mathbb{L}_k/\ell_i\mathbb{Z}^3$, $i=1,2$, with the same $k$, $\mu$ and $e_k$ on both tori.
It consists of three clauses, whose content is the following.

Clause~(a) concerns the rows.
Let $\mathrm{Op}$ be one of the operators of the table of Proposition~25.1, let $\mathrm{rg}$ be its range, and let $t$ be a row index.
For $i=1,2$ let $B_i$ be a complex bond function on $T_i$ and $Y_i$ a union of $M$-cubes of $T_i$, and let $\left(\tilde{B}_i,\tilde{Y}_i\right)$ be their periodic extensions.
If $\left(\tilde{B}_1,\tilde{Y}_1\right)$ and $\left(\tilde{B}_2,\tilde{Y}_2\right)$ coincide on the bonds, respectively at the sites, that lie within the distance $\mathrm{rg}$ of $\mathrm{supp}\,t$, then the rows of $\mathrm{Op}_{\mathbb{L}}\left(\tilde{B}_1,\tilde{Y}_1\right)$ and of $\mathrm{Op}_{\mathbb{L}}\left(\tilde{B}_2,\tilde{Y}_2\right)$ at $t$ coincide.
If, in addition, $t$ and $t'$ are contained in a set $\tilde{U}$ with $\mathrm{sdm}(\tilde{U})+\mathrm{rg}<\ell_i$ for $i=1,2$, then the kernel of $\mathrm{Op}(B_1,Y_1)$ on $T_1$ and the kernel of $\mathrm{Op}(B_2,Y_2)$ on $T_2$ at $(t,t')$ are both equal to the kernel $\mathrm{Op}_{\mathbb{L}}\left(\tilde{B}_1,\tilde{Y}_1\right)(t,t')$ on the covering lattice.

Clause~(b) concerns the operators of Definition~24.D near a polymer.
Its hypotheses are these: $\tilde{\mathbf{Y}}$ is a finite non-empty union of $M$-cubes of $\mathbb{L}_k$ that is connected through faces and satisfies
\[
\mathrm{sdm}(\tilde{\mathbf{Y}})+3M<\ell_i\qquad(i=1,2);
\]
$\mathbf{Y}_i$ is the image of $\tilde{\mathbf{Y}}$ in $T_i$;
and, for $i=1,2$, $B_i$ is a complex bond function on $T_i$, where $B_1$ and $B_2$ have equal lifts on the bonds $b$ that meet $\tilde{\mathbf{Y}}$, that is, $B_1(\mathrm{pr}_{\ell_1}b)=B_2(\mathrm{pr}_{\ell_2}b)$ for these bonds.
Under these hypotheses five assertions hold, all indices being taken in $\tilde{\mathbf{Y}}^{+\eta_k}$.

Assertion~(b1) describes the polymers and the admissible pairs.
The polymers contained in $\mathbf{Y}_i$ are the images of the non-empty unions $\tilde{X}$ of $M$-cubes contained in $\tilde{\mathbf{Y}}$ that are connected through faces.
The pairs of $\mathrm{Pz}_k(T_i)$ that are admissible for the image of such a set $\tilde{X}$ are the images of the pairs of $\mathrm{Pzt}_k$ that are admissible for $\tilde{X}$, and each of the latter has for $\tilde{X}$ the kind, (I) or (II), that its image has for the image of $\tilde{X}$; here the two kinds of Section~I.1.7 are read on the covering lattice as in Lemma~1\,(d).
A site at distance at most $\eta_k$ from $\tilde{\mathbf{Y}}$ lies in the periodic extension of $\mathbf{Y}_i$ if and only if it lies in $\tilde{\mathbf{Y}}$.

Assertion~(b2) states that at these pairs the rows of $\delta_{\alpha,B_1}$ and of $\delta_{\alpha,B_2}$ coincide, and so do the rows of $\delta_{\alpha,-B_1}$ and of $\delta_{\alpha,-B_2}$.

Assertion~(b3) concerns $k\ge1$.
For $j=1,2$ the rows of $T^{(j)}(B_1,\mathbf{Y}_1)$ and of $T^{(j)}(B_2,\mathbf{Y}_2)$ at the unit sites of $\tilde{\mathbf{Y}}$ coincide, and so do the rows of $\bar{T}^{(j)}(B_1,\mathbf{Y}_1)$ and of $\bar{T}^{(j)}(B_2,\mathbf{Y}_2)$ at these sites.
Consequently, with $\tilde{B}$ the common lift of $B_1$ and $B_2$ on the bonds that meet $\tilde{\mathbf{Y}}$,
\[
\mathrm{cov}_{\ell_1}^{-1}\circ\mathrm{sub}_{\mathbf{Y}_1,B_1}\circ\mathrm{cov}_{\ell_1}
=\mathrm{cov}_{\ell_2}^{-1}\circ\mathrm{sub}_{\mathbf{Y}_2,B_2}\circ\mathrm{cov}_{\ell_2}
=\mathrm{sub}_{\tilde{\mathbf{Y}},\tilde{B}}
\]
on the subalgebra on which $\mathrm{sub}_{\tilde{\mathbf{Y}},\tilde{B}}$ is defined in (o5).

Assertion~(b4) states that the entries $C(b',\Upsilon)$ with $b'$ and $\Upsilon$ contained in $\tilde{\mathbf{Y}}$ are the same on $T_1$ and on $T_2$.

Assertion~(b5) concerns the objects of Proposition~25.2.
Let $\tilde{X}$ be a set as in (b1).
The objects (o1)--(o4), (o6), the set $\bar{X}^{L}$ of (o7) and the objects (o8), (o9), formed on $T_1$ and on $T_2$ at data with equal lifts (data whose compositions with the covering maps coincide, in the sense of (G3.r2)), are, in the sense of Proposition~25.2, the images under $\mathrm{cov}_{\ell_1}$ and under $\mathrm{cov}_{\ell_2}$ of one and the same object on the covering lattice; where these objects are numbers, the two numbers are equal.
If $p_{0,k}>0$, then in the coordinates $\mathsf{Z}(\Upsilon)$ the measure $\hat{\mu}$ is, on each of the two tori, the product over the basis of that torus of copies of one and the same probability measure on $\mathbb{R}$, and $\delta\varepsilon^{0}_{k}$ is the same number on both tori (see (o10)).

For the homomorphism $S_L$ of (o7) and a finite set $\tilde{U}$ contained in $\tilde{\mathbf{Y}}^{+\eta_k}$, assertion~(b5) contains the identity
\[
\mathrm{cov}_{\ell_i/L}\circ S_L=S_L\circ\mathrm{cov}_{\ell_i}\qquad(i=1,2)
\]
on the subalgebra of $\mathrm{Grc}_k(\tilde{U})$ generated by $\psi$ and $\chi$.
In it, $\mathrm{cov}_{\ell_i/L}$ is the map of Section~G3.5.0 when $\mathrm{sdm}(\tilde{U})+L\cdot M<\ell_i$; under the condition $\mathrm{sdm}(\tilde{\mathbf{Y}})+3M<\ell_i$ alone, the identity holds with the injective homomorphisms $\mathrm{cov}_{\ell_i/L}$ of Section~G3.5.0.

For the operation $L$ of (o7), assume in addition that $\ell_1/L$ and $\ell_2/L$ lie in $\mathrm{Cl}$, and let $\tilde{Y}'$ be a finite non-empty union of $M$-cubes of $\mathbb{L}_{k+1}$ that is connected through faces and satisfies $L\tilde{Y}'\subset\tilde{\mathbf{Y}}$.
For $i=1,2$ let $Y'_i$ be the image of $\tilde{Y}'$ in $\mathbb{L}_{k+1}/(\ell_i/L)\mathbb{Z}^3$, let $\mathrm{Ef}^{(i)}$ be a family of even elements of $\mathrm{Gr}[\psi,\chi]$ of level $k$ on $T_i$ with the property (k2) of Section~I.5.4, and let $(\mathscr{A}_i,\mathscr{C}_i)$ be a pair of bond functions on $\mathbb{L}_{k+1}/(\ell_i/L)\mathbb{Z}^3$ for which the elements $\mathrm{Ef}^{(i)}\left(X;(\mathscr{A}_i)_L,(\mathscr{C}_i)_L\right)$ are defined.
For $\tilde{X}\in\mathrm{Ix}(\tilde{Y}')$ let $X_i$ be the image of $\tilde{X}$ in $T_i$.
If the elements
\[
\mathrm{cov}_{\ell_i}^{-1}\left(\mathrm{Ef}^{(i)}\left(X_i;(\mathscr{A}_i)_L,(\mathscr{C}_i)_L\right)\right),\qquad\tilde{X}\in\mathrm{Ix}(\tilde{Y}'),
\]
do not depend on $i$, then
\[
\left(L\,\mathrm{Ef}^{(i)}\right)\left(Y'_i;\mathscr{A}_i,\mathscr{C}_i\right)=\mathrm{cov}_{\ell_i/L}\left(\Phi_L\right)\qquad(i=1,2),
\]
with one index set $\mathrm{Ix}(\tilde{Y}')$ and one and the same element $\Phi_L$ for $i=1$ and for $i=2$.
Here $\mathrm{cov}_{\ell_i/L}$ is the map of Section~G3.5.0 at the level $k+1$ if $\mathrm{sdm}(\tilde{Y}')+2\eta_{k+1}+M<\ell_i/L$, and the injective homomorphism of Section~G3.5.0 otherwise.

The numbers of Section~I.5.7 are compared in the setting described at the end of (o8), given on each of the two tori with a family (p5), written $S^{(i)}_{k,y}$, on $T_i$.
Let $\tilde{X}$ be contained in $\tilde{\mathbf{Y}}$, let $X_i$ be the image of $\tilde{X}$ in $T_i$, and let every connection on $T_i$ that occurs below lie in $\tilde{G}_k(T_i)$.
For a pair $(B'_1,B'_2)$ of such connections two conditions are considered:
(i) $B'_1$ and $B'_2$ have equal lifts on the bonds contained in $\tilde{X}$;
(ii) for every $y\ge0$ the kernels $S^{(i)}_{k,y}(X_i,B'_i)(x,x')$, $x$ and $x'$ over $\tilde{X}$, do not depend on $i$.
By (o8), $\acute{E}_k(X_i,B'_i)$ is the value of one functional at the lift of $B'_i$ on the bonds contained in $\tilde{X}$ and at the lifts of these kernels; hence (i) and (ii) give $\acute{E}_k(X_1,B'_1)=\acute{E}_k(X_2,B'_2)$.
Condition~(i) holds at the five pairs $(B_1,B_2)$, $(-B_1,-B_2)$, $\left(\overline{B_1},\overline{B_2}\right)$, $\left(-\overline{B_1},-\overline{B_2}\right)$ and $(0,0)$, by the hypothesis of (b), since a bond contained in $\tilde{X}$ meets $\tilde{\mathbf{Y}}$.
It follows that $\acute{E}_k(X_i,B_i)$, $E^{d}_k(X_i,B_i)$ and $E_k^{\mathrm{det}}(X_i,B_i)$ do not depend on $i$ if (ii) holds at each pair of connections at which $\acute{E}_k$ is evaluated in them: at $(B_1,B_2)$ for $\acute{E}_k$, at the first four of the five pairs for $E^{d}_k$, and at all five for $E_k^{\mathrm{det}}$.
For a pair $(A_i,C_i)$ of connections on $T_i$, $i=1,2$, the number
\[
{\delta E^{\mathrm{det}}}_k(X_i;A_i,C_i)=E^{d}_k(X_i,A_i+C_i)-E^{d}_k(X_i,C_i)
\]
is evaluated at $\pm(A_i+C_i)$, $\pm\overline{A_i+C_i}$, $\pm C_i$ and $\pm\overline{C_i}$.
If each of the pairs $(A_1+C_1,A_2+C_2)$ and $(C_1,C_2)$ has equal lifts on the bonds contained in $\tilde{X}$, and (ii) holds at the pairs $(B'_1,B'_2)$, $(-B'_1,-B'_2)$, $\left(\overline{B'_1},\overline{B'_2}\right)$ and $\left(-\overline{B'_1},-\overline{B'_2}\right)$ for $B'_i=A_i+C_i$ and for $B'_i=C_i$, then (i) holds at these eight pairs as well; so $E^{d}_k(X_i,A_i+C_i)$ and $E^{d}_k(X_i,C_i)$ do not depend on $i$, and ${\delta E^{\mathrm{det}}}_k(X_1;A_1,C_1)={\delta E^{\mathrm{det}}}_k(X_2;A_2,C_2)$.

The last part of (b5) concerns the sums that define $\varepsilon$ and $m$ in Section~I.5.6; it holds for all $\ell_1,\ell_2\in\mathrm{Cl}$ and makes no reference to the set $\tilde{\mathbf{Y}}$.
Let $\tilde{\square}$ be an $M$-cube of $\mathbb{L}_k$, and let $\mathrm{Sm}(\tilde{\square})$ be the set of the finite non-empty unions $\tilde{X}$ of $M$-cubes of $\mathbb{L}_k$ that are connected through faces, contain $\tilde{\square}$ and satisfy $d_M(\tilde{X})\le L$.
For $i=1,2$ let $\mathrm{Ef}^{(i)}$ be a family on $T_i$ with (k2).
Assume that for every $\tilde{X}\in\mathrm{Sm}(\tilde{\square})$, with $X_i$ the image of $\tilde{X}$ in $T_i$, neither the number $\mathrm{Ef}^{(i)}_{00}(X_i;0,0)$ nor the kernel $\mathrm{Ef}^{(i)}_{20}(X_i;0,0;x,x')$, $x$ and $x'$ over $\tilde{X}$, depends on $i$.
Then $\mathrm{al}_{0}\left(\mathrm{Ef}^{(i)},X_i\right)$ and $\mathrm{al}_{2}\left(\mathrm{Ef}^{(i)},X_i\right)$, $\tilde{X}\in\mathrm{Sm}(\tilde{\square})$, do not depend on $i$, and neither do the sums
\[
-\sum_{X}\mathrm{al}_{0}\left(\mathrm{Ef}^{(i)},X\right)
\qquad\text{and}\qquad
-\sum_{X}\mathrm{al}_{2}\left(\mathrm{Ef}^{(i)},X\right),
\]
where $X$ ranges over the small polymers of $T_i$ that contain the image of $\tilde{\square}$.
These sums are the right-hand sides of the display of Section~I.5.6 that defines $\varepsilon(E)$ and $m(E)\cdot\mathbf{1}_4$, formed with the image $\square_i$ of $\tilde{\square}$ in $T_i$;
in the notation of Proposition~16.3 and of Section~I.5.11 they are the number $\varepsilon\left(\mathrm{Ef}^{(i)},\square_i\right)$ and the $4\times4$ matrix $m\left(\mathrm{Ef}^{(i)},\square_i\right)$.
For a family with (k2) alone this matrix need not be a multiple of $\mathbf{1}_4$, and both quantities may depend on the cube.
Where Proposition~16.3\,(a) applies to $\mathrm{Ef}^{(i)}$, or Proposition~48.2\,(a) for a family without slot, that is, for $\mathrm{Ef}^{(i)}$ in $K^{x,C}_{k}(\ell_i)$ or in $K^{C}_{k}(\ell_i)$, they are equal to $\varepsilon\left(\mathrm{Ef}^{(i)}\right)$ and $m\left(\mathrm{Ef}^{(i)}\right)\cdot\mathbf{1}_4$, and the conclusion reads
\[
\varepsilon\left(\mathrm{Ef}^{(1)}\right)=\varepsilon\left(\mathrm{Ef}^{(2)}\right),\qquad m\left(\mathrm{Ef}^{(1)}\right)=m\left(\mathrm{Ef}^{(2)}\right).
\]

Clause~(c) is the periodisation.
Let $\ell_2=\mathrm{Jr}\cdot\ell_1$ with $\mathrm{Jr}$ odd, and let $\mathrm{pj}\colon T_2\to T_1$ and $\mathrm{pj}^{*}$ be as in Section~I.6.3.
Then
\[
\mathrm{Op}\left(\mathrm{pj}^{*}B,\mathrm{pj}^{-1}Y\right)\circ\mathrm{pj}^{*}=\mathrm{pj}^{*}\circ\mathrm{Op}(B,Y)
\]
for every operator $\mathrm{Op}$ of the table of Proposition~25.1, every complex bond function $B$ on $T_1$ and every union $Y$ of $M$-cubes of $T_1$.
For the operator of no.~14 of that table, whose source consists of objects of level $k+1$, the map $\mathrm{pj}^{*}$ to the right of $\mathrm{Op}$ is the inclusion of the periodic objects (bond functions), respectively of the antiperiodic objects (spinors and functions of pairs), of level $k+1$ and side $\ell_1/L$ into those of side $\ell_2/L=\mathrm{Jr}\cdot\ell_1/L$; this is the definition of Section~I.6.3, read at the level $k+1$ with these two sides, which need not lie in $\mathrm{Cl}$.
Since $\mathrm{Jr}$ is odd, an object that is antiperiodic for the side $\ell_1/L$ is antiperiodic for the side $\ell_2/L$.

\subsubsection*{4.2. Proof of (a)}

\begin{proof}[Proof of Proposition~25.3\,(a)]
By 2.2 and 2.3, $\mathrm{Op}$ is a local family whose sets $\mathrm{RdB}(t)$ and $\mathrm{RdY}(t)$ consist of bonds, respectively sites, within the distance $\mathrm{rg}$ of $\mathrm{supp}\,t$.
Assume that $(\tilde{B}_1,\tilde{Y}_1)$ and $(\tilde{B}_2,\tilde{Y}_2)$ agree on all bonds and at all sites within this distance.
Then $\tilde{B}_1=\tilde{B}_2$ on $\mathrm{RdB}(t)$, and $\tilde{Y}_1$, $\tilde{Y}_2$ contain the same sites of $\mathrm{RdY}(t)$; so (l2) gives
\[
\mathrm{Op}_{\mathbb{L}}(\tilde{B}_1,\tilde{Y}_1)[t,\,\cdot\,]=\mathrm{Op}_{\mathbb{L}}(\tilde{B}_2,\tilde{Y}_2)[t,\,\cdot\,],
\]
which is the first assertion.

Now let $\mathrm{supp}\,t$ and $\mathrm{supp}\,t'$ be contained in a set $\tilde{U}$ with $\mathrm{sdm}(\tilde{U})+\mathrm{rg}<\ell_i$ for $i=1,2$.
By (f3), established in 2.4, $\mathrm{Op}(B_i,Y_i)$ is the restriction of $\mathrm{Op}_{\mathbb{L}}(\tilde{B}_i,\tilde{Y}_i)$ to the periodic, respectively antiperiodic, functions of side $\ell_i$.
Lemma~3\,(b), (c), applied with $\ell=\ell_i$, shows that the coefficient of this restriction at $(t,t')$ is $\mathrm{Op}_{\mathbb{L}}(\tilde{B}_i,\tilde{Y}_i)[t,t']$; the kernel is the coefficient divided by the weight $\nu(t')$, which is the same number on the covering lattice and on both tori (paragraph ``Functions'' of 1.4).
The two coefficients coincide by the first assertion.
For no.~14 the column index is of level $k+1$, the sides at that level are $\ell_i/L$, and ``$t'$ is contained in $\tilde{U}$'' means that $L\,\mathrm{supp}\,t'$ is contained in $\tilde{U}$ (last column of no.~14 in the table of Proposition~25.1; (Of14)); Lemma~3\,(c) applies with the shifts $t'+(\ell_i/L)w$, for which
\[
L\,\mathrm{supp}\left(t'+(\ell_i/L)w\right)=L\,\mathrm{supp}\,t'+\ell_i w.
\]
\end{proof}

\subsubsection*{4.3. Proof of (b)}

\begin{proof}[Proof of Proposition~25.3\,(b)]
\emph{Notation.}
Write $\mathrm{pr}_i:=\mathrm{pr}_{\ell_i}$.
Let $\tilde{B}$ be the common lift of $B_1$ and $B_2$ on the bonds that meet $\tilde{\mathbf{Y}}$, that is, $\tilde{B}(b)=B_1(\mathrm{pr}_1b)=B_2(\mathrm{pr}_2b)$.
Let $\tilde{B}_i$ be the $\ell_i\mathbb{Z}^3$-periodic extension of $B_i$; then $\tilde{B}_i=\tilde{B}$ on the bonds that meet $\tilde{\mathbf{Y}}$.
The set $\tilde{Y}_i:=\tilde{\mathbf{Y}}+\ell_i\mathbb{Z}^3$ is the periodic extension of $\mathbf{Y}_i$.
Since $\mathrm{sdm}(\tilde{\mathbf{Y}})+3M<\ell_i$ and $\eta_k\le1\le M$,
\[
\mathrm{sdm}\left(\tilde{\mathbf{Y}}^{+\eta_k}\right)+M\le\mathrm{sdm}(\tilde{\mathbf{Y}})+2\eta_k+M<\ell_i.
\]
Hence $\mathrm{pr}_i$ is injective on $\tilde{\mathbf{Y}}^{+\eta_k}$ and on the indices contained in it (Lemma~1\,(a)), $\mathrm{cov}_{\ell_i}$ is defined on $\mathrm{Grc}_k(\tilde{\mathbf{Y}}^{+\eta_k})$, and every union $\tilde{X}$ of $M$-cubes contained in $\tilde{\mathbf{Y}}$ satisfies, for every $\mathrm{rg}\le1$,
\[
\mathrm{sdm}(\tilde{X})+2\,\mathrm{rg}+M<\ell_i.
\]
\begin{itemize}
\item[(b1)]
The first clause is Lemma~1\,(b) with $\ell=\ell_i$.
The second clause is Lemma~1\,(d) with $\ell=\ell_i$ (note that $\mathrm{sdm}(\tilde{X})+2\eta_k<\ell_i$): the two conditions (I) and (II) are preserved in both directions, and the lift of an admissible pair that is admissible for $\tilde{X}$ is the unique lift with both photon sites in $\tilde{\mathbf{Y}}^{+\eta_k}$.
The third clause is Lemma~1\,(a) with $\tilde{U}=\tilde{\mathbf{Y}}$, $a=\eta_k$ and $\ell'=\ell_i$.
\item[(b2)]
Let $(x,x')\in\mathrm{Pzt}_k$ be admissible for a union $\tilde{X}$ of $M$-cubes contained in $\tilde{\mathbf{Y}}$.
By (f3), the row of $\delta_{\alpha,B_i}$ at $(x,x')$ is the row of $\delta_{\alpha,\tilde{B}_i}$.
According to (Of11), this row has the columns $x$ and $x'$, with the coefficient $-d_X(x,x')^{-\alpha}\,\mathbf{1}_4$ at $x$ and $d_X(x,x')^{-\alpha}\cdot U_{\tilde{B}_i}(x,x')\,\mathbf{1}_4$ at $x'$, and it reads $\tilde{B}_i$ on the set $\mathrm{RdB}(x,x')=\mathrm{Pth}(c,c')$ (last column of the table of Proposition~25.1).
By Lemma~1\,(d), either the bonds of $\mathrm{Pth}(c,c')$ are contained in $\tilde{X}$, or $\mathrm{Pth}(c,c')$ consists of the bond $(c,c_3)$, which meets $\tilde{X}$; in both cases these bonds meet $\tilde{\mathbf{Y}}$, where $\tilde{B}_1=\tilde{B}_2=\tilde{B}$.
Both rows therefore coincide with the row of $\delta_{\alpha,\tilde{B}}$.
The same holds with $-B_i$ and $-\tilde{B}$.
The kernels on the two tori at $((x,x'),x'')$, $x''$ over $\tilde{\mathbf{Y}}^{+\eta_k}$, are these coefficients divided by $\eta_k^3/48$ (Lemma~3\,(c) with $\tilde{U}=\tilde{\mathbf{Y}}^{+\eta_k}$ and $\mathrm{rg}=M$).
\item[(b3)]
Apply 3.7\,(C) with $\ell=\ell_1$ and with $\ell=\ell_2$: for a unit site $y$ of $\tilde{\mathbf{Y}}$ and $j=1,2$,
\[
T^{(j)}(\tilde{B}_i,\tilde{Y}_i)[y,\,\cdot\,]=T^{(j)}(\tilde{B},\tilde{\mathbf{Y}})[y,\,\cdot\,],\qquad
\bar{T}^{(j)}(\tilde{B}_i,\tilde{Y}_i)[y,\,\cdot\,]=\bar{T}^{(j)}(\tilde{B},\tilde{\mathbf{Y}})[y,\,\cdot\,],
\]
and the right-hand sides do not contain $i$.
The columns lie over $\tilde{\mathbf{Y}}$ for $j=1$; for $j=2$ they are pairs $(x,x^{(3)})$ with $c\in\tilde{\mathbf{Y}}$ and $c_3\in\tilde{\mathbf{Y}}^{+\eta_k}$.
By 3.7\,(C),
\[
\mathrm{cov}_{\ell_i}\circ\mathrm{sub}_{\tilde{\mathbf{Y}},\tilde{B}}=\mathrm{sub}_{\mathbf{Y}_i,B_i}\circ\mathrm{cov}_{\ell_i}
\]
on the subalgebra $\mathrm{As}$ of $\mathrm{Grc}_k(\tilde{\mathbf{Y}}^{+\eta_k})$ generated by $\psi$, $\chi$ and the $\Psi(u)$, $u\in\tilde{\mathbf{Y}}$.
Since $\mathrm{cov}_{\ell_i}$ is injective, $\mathrm{cov}_{\ell_i}^{-1}\circ\mathrm{sub}_{\mathbf{Y}_i,B_i}\circ\mathrm{cov}_{\ell_i}=\mathrm{sub}_{\tilde{\mathbf{Y}},\tilde{B}}$ on $\mathrm{As}$, for $i=1$ and for $i=2$.
\item[(b4)]
By (Of12), $C$ is a local family of range $L-1$ without connection.
For $b'$ and $\Upsilon$ contained in $\tilde{\mathbf{Y}}$ one has $\mathrm{sdm}(\tilde{\mathbf{Y}})+(L-1)<\ell_i$, since $L-1<M$.
By Lemma~3\,(c), with the character $\epsilon=1$, the entry of $C$ on $T_i$ at $(b',\Upsilon)$ is the coefficient $C(b',\Upsilon)$ of (Of12).
\item[(b5)]
Let $\tilde{X}$ be a non-empty union of $M$-cubes contained in $\tilde{\mathbf{Y}}$ and connected through faces, and let $X_i:=\mathrm{pr}_i(\tilde{X})$.
By (b1), $\tilde{X}$ is a lift of $X_i$, and $\mathrm{sdm}(\tilde{X})+2\,\mathrm{rg}+M<\ell_i$ for $\mathrm{rg}\le1$.
Data on $T_1$ and on $T_2$ are said to have equal lifts if, lifted to the $\mathrm{rg}$-neighbourhood of $\tilde{X}$ in the sense of (G3.r2) in Section~G3.5.0, they give one and the same datum on the covering lattice.
The data of the list below are the following, for $i=1,2$, and they are assumed to have equal lifts; (o1) and (o9) have no data.
In (o2), $O_i$ is an operator on $\mathfrak{S}(1,\ell_i)$ (for $k=0$: on $\mathfrak{S}^{\mathrm{ch}}(1,\ell_i)$) whose kernel vanishes unless both indices lie in $X_i$ (for $k=0$: over $X_i$); the kernels of $O_1$ and $O_2$ coincide at the unit indices in $\tilde{X}$ (for $k=0$: at the fermion sites over $\tilde{X}$), and $O$ denotes their common values there.
In (o3), where $d_M(\tilde{X})\le L$, the data on $T_i$ are a number $\mathrm{Ec}^{(i)}_{00}$ and a kernel $\mathrm{Ec}^{(i)}_{20}$ of $T_i$, antiperiodic in each argument, that vanishes unless both arguments lie over $X_i$, with
\[
\mathrm{Ec}^{(1)}_{00}=\mathrm{Ec}^{(2)}_{00}=:\mathrm{Ec}_{00},\qquad
\mathrm{Ec}^{(1)}_{20}(x,x')=\mathrm{Ec}^{(2)}_{20}(x,x')=:\mathrm{Ec}_{20}(x,x')\quad\text{for }x,x'\text{ over }\tilde{X}.
\]
The numbers $\mathrm{al}_{0}$ and $\mathrm{al}_{2}$ on $T_i$ formed with these data are given by the display of Section~I.5.6 that defines $\mathrm{al}_{0}(E,X)$ and $\mathrm{al}_{2}(E,X)$, with $X_i$ in the place of $X$ and with $\mathrm{Ec}^{(i)}_{00}$ and $\mathrm{Ec}^{(i)}_{20}$ in the places of $E_{00}(X;0,0)$ and $E_{20}(X;0,0;\cdot,\cdot)$, both arguments ranging over one lift of $X_i$ (last sentence of (I.AP)).
Let $E^{(i)}$ be a family on $T_i$ with (k2); by (k2) the kernel $E^{(i)}_{20}(X_i;0,0;\cdot,\cdot)$ vanishes unless both arguments lie over $X_i$, and the data $\mathrm{Ec}^{(i)}_{00}=E^{(i)}_{00}(X_i;0,0)$ and $\mathrm{Ec}^{(i)}_{20}=E^{(i)}_{20}(X_i;0,0;\cdot,\cdot)$ give the numbers $\mathrm{al}_{0}\left(E^{(i)},X_i\right)$ and $\mathrm{al}_{2}\left(E^{(i)},X_i\right)$.
In (o4), where $d_M(\tilde{X})\le L$, $\mathscr{T}_i$ is a bond function on $T_i$, and
\[
\mathscr{T}_1(\mathrm{pr}_1b)=\mathscr{T}_2(\mathrm{pr}_2b)=:\tilde{\mathscr{T}}(b)\qquad\text{for the bonds }b\text{ contained in }\tilde{X}.
\]
In (o6), $k=0$, $\tilde{X}=\tilde{\square}$ is an $M$-cube contained in $\tilde{\mathbf{Y}}$, and $\mathrm{Zf}_i$ is a bond function on $T_i$ with $\mathrm{Zf}_1(\mathrm{pr}_1b)=\mathrm{Zf}_2(\mathrm{pr}_2b)=:\widetilde{\mathrm{Zf}}(b)$ for the bonds $b$ that meet $\tilde{\square}$.
In (o8), $B'_i$ is a complex bond function on $T_i$, and
\[
B'_1(\mathrm{pr}_1b)=B'_2(\mathrm{pr}_2b)=:\tilde{B}'(b)\qquad\text{for the bonds }b\text{ contained in }\tilde{X};
\]
the bond functions $B_i$ of the hypothesis of (b) are an instance, with $\tilde{B}'=\tilde{B}$, because a bond contained in $\tilde{X}$ meets $\tilde{\mathbf{Y}}$.
Further, $O_i$ is an operator on $\mathfrak{S}^{\mathrm{ch}}(\eta_k,\ell_i)$ whose kernel vanishes unless both arguments lie over $X_i$; the kernels of $O_1$ and $O_2$ coincide at the fermion sites over $\tilde{X}$, and $O$ denotes their common values there; and $y$ is real.
By the points (C) of 3.3--3.6, 3.8, 3.10 and 3.11, the following holds for $i=1,2$.
\begin{itemize}
\item[(o1)]
$\operatorname{Vol}(X_i)=M^3\cdot|\tilde{X}|_M$, and
\[
\int_{X_i}\mathrm{d}x\,\bar{\psi}(x)\cdot\mathrm{Mx}\,\psi(x)=\mathrm{cov}_{\ell_i}\left(\int_{\tilde{X}}\mathrm{d}x\,\bar{\psi}(x)\cdot\mathrm{Mx}\,\psi(x)\right).
\]
\item[(o2)]
$\langle\bar{\Psi},O_i\Psi\rangle=\mathrm{cov}_{\ell_i}\left(\Phi(\tilde{X};O)\right)$.
\item[(o3)]
$X_1$ is small if and only if $d_M(\tilde{X})\le L$, and this holds if and only if $X_2$ is small (Lemma~1\,(c)); in that case $\mathrm{al}_{0}$ and $\mathrm{al}_{2}$ on $T_i$, formed with $\mathrm{Ec}^{(i)}_{00}$ and $\mathrm{Ec}^{(i)}_{20}$, are the numbers of 3.5\,(A) formed with $\mathrm{Ec}_{00}$ and $\mathrm{Ec}_{20}$, by the argument of the paragraph ``The numbers'' of 3.5\,(C) with $\ell=\ell_i$, in which the vanishing of $\mathrm{Ec}^{(i)}_{20}$ unless both arguments lie over $X_i$ takes the place of (k2).
\item[(o4)]
$\mathrm{cov}_{\ell_i}\circ\boldsymbol{\Theta}_{\tilde{X},\tilde{\mathscr{T}}}=\boldsymbol{\Theta}_{X_i,\mathscr{T}_i}\circ\mathrm{cov}_{\ell_i}$ on the subalgebra of $\mathrm{Grc}_k(\tilde{X})$ generated by the $\psi(\xi)$, $\xi$ over $\tilde{X}$.
\item[(o6)]
$\mathrm{Vx}(\mathrm{Zf}_i,\mathrm{pr}_i\tilde{\square})=\mathrm{cov}_{\ell_i}\left(\mathrm{Vx}(\widetilde{\mathrm{Zf}},\tilde{\square})\right)$.
\item[(o8)]
The three traces of (o8) on $T_i$, formed with $B'_i$, $O_i$ and $y$, are $\mathrm{Trk}_k(\tilde{X};\tilde{B}',O,y)$ for $k\ge1$, and $\mathrm{Trk}_0'(\tilde{X};\tilde{B}',O)$ and $\mathrm{Trk}_0''(\tilde{X};O)$ for $k=0$ (3.10\,(C) with $\ell=\ell_i$).
\item[(o9)]
$\int\mathrm{d}\mu_I(W)\,\mathrm{cov}_{\ell_i}\left(F\right)=\mathrm{cov}_{\ell_i}\left(\int_c F\right)$ for $F\in\mathrm{Grc}_k(\tilde{\mathbf{Y}}^{+\eta_k})$.
\end{itemize}
In each case the object on the covering lattice does not contain $i$.

\emph{(o3): the sums that define $\varepsilon$ and $m$ in Section~I.5.6.}
Let $\tilde{\square}$ be an $M$-cube of $\mathbb{L}_k$, and let $\mathrm{Ef}^{(i)}$ be a family on $T_i$ with (k2).
By Lemma~1\,(c), every $\tilde{X}\in\mathrm{Sm}(\tilde{\square})$ satisfies
\[
\mathrm{sdm}(\tilde{X})+3M\le M\cdot(L+2)+3M=M\cdot(L+5)<\ell_i;
\]
hence the assertion (o3) above holds with $\tilde{X}$ in the place of $\tilde{\mathbf{Y}}$.
Suppose that for every $\tilde{X}\in\mathrm{Sm}(\tilde{\square})$ the number $\mathrm{Ef}^{(i)}_{00}(\mathrm{pr}_i\tilde{X};0,0)$ and the kernel $\mathrm{Ef}^{(i)}_{20}(\mathrm{pr}_i\tilde{X};0,0;x,x')$, $x,x'$ over $\tilde{X}$, do not depend on $i$.
Then $\mathrm{al}_{0}(\mathrm{Ef}^{(i)},\mathrm{pr}_i\tilde{X})$ and $\mathrm{al}_{2}(\mathrm{Ef}^{(i)},\mathrm{pr}_i\tilde{X})$ do not depend on $i$.
The map $\tilde{X}\mapsto\mathrm{pr}_i(\tilde{X})$ is a bijection of $\mathrm{Sm}(\tilde{\square})$ onto the set of the small polymers of $T_i$ that contain $\mathrm{pr}_i(\tilde{\square})$ (Lemma~1\,(c)).
Hence the sums $-\sum\mathrm{al}_{0}(\mathrm{Ef}^{(i)},X)$ and $-\sum\mathrm{al}_{2}(\mathrm{Ef}^{(i)},X)$ over these polymers do not depend on $i$.
They are the right-hand sides of the display of Section~I.5.6 that defines $\varepsilon$ and $m$, formed with the cube $\mathrm{pr}_i(\tilde{\square})$, that is, the number $\varepsilon\left(\mathrm{Ef}^{(i)},\mathrm{pr}_i\tilde{\square}\right)$ and the matrix $m\left(\mathrm{Ef}^{(i)},\mathrm{pr}_i\tilde{\square}\right)$ in the notation of Proposition~16.3 and of Section~I.5.11.
Where Proposition~16.3\,(a) or Proposition~48.2\,(a) applies to $\mathrm{Ef}^{(i)}$, these are $\varepsilon\left(\mathrm{Ef}^{(i)}\right)$ and $m\left(\mathrm{Ef}^{(i)}\right)\cdot\mathbf{1}_4$, and one obtains $\varepsilon\left(\mathrm{Ef}^{(1)}\right)=\varepsilon\left(\mathrm{Ef}^{(2)}\right)$ and $m\left(\mathrm{Ef}^{(1)}\right)=m\left(\mathrm{Ef}^{(2)}\right)$.

\emph{(o8): the numbers of Section~I.5.7.}
Let the setting of the paragraph of 3.10 on the numbers $\acute{E}_k$, $E^{d}_k$, $E_k^{\mathrm{det}}$, ${\delta E^{\mathrm{det}}}_k$ be given on both tori, with families $S^{(i)}_{k,y}$, and let $(B'_1,B'_2)$, $B'_i\in\tilde{G}_k(T_i)$, be a pair that satisfies the conditions (i) and (ii) of the statement.
For every $y\ge0$ the kernel of $O_i:=S^{(i)}_{k,y}(X_i,B'_i)$ vanishes unless both arguments lie over $X_i$, by the first half of W12 (Section~G1.11);
by (ii) the kernels of $O_1$ and $O_2$ coincide at the fermion sites over $\tilde{X}$, and by (i) the functions $B'_1$ and $B'_2$ have equal lifts on the bonds contained in $\tilde{X}$.
So (o8) above shows that the integrands of $\acute{E}_k(X_1,B'_1)$ and of $\acute{E}_k(X_2,B'_2)$ coincide for every $y\ge0$, and therefore $\acute{E}_k(X_1,B'_1)=\acute{E}_k(X_2,B'_2)$.
The lift $B'\mapsto\left(b\mapsto B'(\mathrm{pr}_ib)\right)$ is linear and commutes with complex conjugation.
Hence condition (i) at $(B'_1,B'_2)$ implies condition (i) at $(-B'_1,-B'_2)$, $\left(\overline{B'_1},\overline{B'_2}\right)$ and $\left(-\overline{B'_1},-\overline{B'_2}\right)$, and condition (i) holds at $(0,0)$.
It holds at $(B_1,B_2)$ by the hypothesis of (b), and at $(A_1+C_1,A_2+C_2)$ and $(C_1,C_2)$ by the hypothesis of the assertion on ${\delta E^{\mathrm{det}}}_k$.
Consequently $\acute{E}_k(X_i,B_i)$, $E^{d}_k(X_i,B_i)$, $E_k^{\mathrm{det}}(X_i,B_i)$ and ${\delta E^{\mathrm{det}}}_k(X_i;A_i,C_i)$ do not depend on $i$ if (ii) holds at each pair at which $\acute{E}_k$ is evaluated in them:
at $(B_1,B_2)$ for $\acute{E}_k$;
at the four pairs $(B_1,B_2)$, $(-B_1,-B_2)$, $\left(\overline{B_1},\overline{B_2}\right)$, $\left(-\overline{B_1},-\overline{B_2}\right)$ for $E^{d}_k$;
at these four and at $(0,0)$ for $E_k^{\mathrm{det}}$;
and, for ${\delta E^{\mathrm{det}}}_k$, at the eight pairs $(B'_1,B'_2)$, $(-B'_1,-B'_2)$, $\left(\overline{B'_1},\overline{B'_2}\right)$, $\left(-\overline{B'_1},-\overline{B'_2}\right)$ with $B'_i=A_i+C_i$ and with $B'_i=C_i$.

\emph{(o7): $\bar{X}^{L}$.}
One has $\bar{X}_i^{L}=\mathrm{pr}_i\left(\mathrm{clL}(\tilde{X})\right)$ by 3.9; no margin beyond that of a lift is needed.

\emph{(o7): $S_L$.}
Let $\tilde{U}$ be a finite set contained in $\tilde{\mathbf{Y}}^{+\eta_k}$.
If $\mathrm{sdm}(\tilde{U})+L\cdot M<\ell_i$ for $i=1,2$, then
\[
\mathrm{cov}_{\ell_i/L}\circ S_L=S_L\circ\mathrm{cov}_{\ell_i}\qquad\text{on }\mathrm{As}_k(\tilde{U})
\]
(3.9, point (C) for $S_L$).
Without this condition one still has $\mathrm{sdm}(\tilde{U})\le\mathrm{sdm}(\tilde{\mathbf{Y}})+2\eta_k<\ell_i$, and the same identity holds with the injective homomorphisms $\mathrm{cov}_{\ell_i/L}$ of Lemma~4\,(a) on $\mathrm{Grc}_{k+1}(\tilde{U}/L)$ (3.9, point (C$'$) for $S_L$).
The map $S_L:\mathrm{As}_k(\tilde{U})\to\mathrm{As}_{k+1}(\tilde{U}/L)$ does not contain $i$.

\emph{(o7): the operation $L$.}
Assume in addition that $\ell_1/L$ and $\ell_2/L$ lie in $\mathrm{Cl}$ ($D_{k+1}$ is formed on the sides $\ell_i/L$).
Let $\tilde{Y}'$ be a finite non-empty union of $M$-cubes of $\mathbb{L}_{k+1}$, connected through faces, with $L\tilde{Y}'$ contained in $\tilde{\mathbf{Y}}$;
let $Y_i:=\mathrm{pr}_{\ell_i/L}(\tilde{Y}')$;
let $\mathrm{Ef}^{(i)}$ be a family of even elements of $\mathrm{Gr}[\psi,\chi]$ of level $k$ on $T_i$ with (k2);
and let $(\mathscr{A}_i,\mathscr{C}_i)$ be a pair of bond functions on $\mathbb{L}_{k+1}/(\ell_i/L)\mathbb{Z}^3$ for which the elements $\mathrm{Ef}^{(i)}(X;(\mathscr{A}_i)_L,(\mathscr{C}_i)_L)$ are defined.
Then
\[
\mathrm{sdm}(L\tilde{Y}')+2\eta_k+M\le\mathrm{sdm}(\tilde{\mathbf{Y}})+2\eta_k+M<\ell_i.
\]
Consequently $\mathrm{sdm}(\tilde{Y}')+\eta_{k+1}=\left(\mathrm{sdm}(L\tilde{Y}')+\eta_k\right)/L<\ell_i/L$, $Y_i$ is a polymer of level $k+1$ with the lift $\tilde{Y}'$ (Lemma~1\,(b) at the level $k+1$), and (25.3) with $\ell=\ell_i$ yields
\[
(L\,\mathrm{Ef}^{(i)})(Y_i;\mathscr{A}_i,\mathscr{C}_i)=\sum_{\tilde{X}\in\mathrm{Ix}(\tilde{Y}')}S_L\left[\mathrm{Ef}^{(i)}\left(\mathrm{pr}_i\tilde{X};(\mathscr{A}_i)_L,(\mathscr{C}_i)_L\right)\right],
\]
with one index set $\mathrm{Ix}(\tilde{Y}')$ for $i=1,2$.
The data are the elements
\[
\mathrm{Ecv}^{(i)}(\tilde{X}):=\mathrm{cov}_{\ell_i}^{-1}\left(\mathrm{Ef}^{(i)}\left(\mathrm{pr}_i\tilde{X};(\mathscr{A}_i)_L,(\mathscr{C}_i)_L\right)\right)\in\mathrm{As}_k(\tilde{X}^{+\eta_k}),\qquad\tilde{X}\in\mathrm{Ix}(\tilde{Y}')
\]
(3.5\,(C), paragraph ``The element over the lift''; note that $\mathrm{sdm}(\tilde{X})+2\eta_k+M<\ell_i$).
They have equal lifts if $\mathrm{Ecv}^{(1)}(\tilde{X})=\mathrm{Ecv}^{(2)}(\tilde{X})=:\mathrm{Ecv}(\tilde{X})$ for all $\tilde{X}\in\mathrm{Ix}(\tilde{Y}')$.
In that case 3.9, point (C$'$) for $L$, applied with $\ell=\ell_i$, gives
\[
(L\,\mathrm{Ef}^{(i)})(Y_i;\mathscr{A}_i,\mathscr{C}_i)=\mathrm{cov}_{\ell_i/L}\left(\Phi_L\left(\tilde{Y}';\left(\mathrm{Ecv}(\tilde{X})\right)_{\tilde{X}\in\mathrm{Ix}(\tilde{Y}')}\right)\right),\qquad i=1,2,
\]
with one and the same element $\Phi_L(\dots)$ of $\mathrm{As}_{k+1}\left({\tilde{Y}'}^{+\eta_{k+1}}\right)$, where $\mathrm{cov}_{\ell_i/L}$ is the injective homomorphism of Lemma~4\,(a) on $\mathrm{Grc}_{k+1}\left({\tilde{Y}'}^{+\eta_{k+1}}\right)$.
If moreover $\mathrm{sdm}(\tilde{Y}')+2\eta_{k+1}+M<\ell_i/L$, this is the map $\mathrm{cov}_{\ell_i/L}$ of Section~G3.5.0 at the level $k+1$ (3.9, point (C) for $L$).

\emph{(o10).}
By 3.11, in the coordinates $\mathsf{Z}(\Upsilon)$ the measure $\hat{\mu}$ is, on both tori, the product over the respective basis of copies of one and the same probability measure on $\mathbb{R}$, and $\delta\varepsilon^{0}_{k}=-d_V\cdot\log P\left(|\mathrm{gn}|\le p_{0,k}\right)$ is the same number, because $p_{0,k}$ is formed with $e_k$, which is the same on both tori.
\end{itemize}
\end{proof}

\subsubsection*{4.4. Proof of (c)}

\begin{proof}[Proof of Proposition~25.3\,(c)]
Let $\ell_2=\mathrm{Jr}\cdot\ell_1$ with $\mathrm{Jr}$ odd, let $B$ be a bond function on $T_1$ and $Y$ a union of $M$-cubes of $T_1$, and let $(\tilde{B},\tilde{Y})$ be their $\ell_1\mathbb{Z}^3$-periodic extensions.
Since $\ell_2\mathbb{Z}^3$ is contained in $\ell_1\mathbb{Z}^3$, the pair $(\tilde{B},\tilde{Y})$ is $\ell_2\mathbb{Z}^3$-periodic, and it is the $\ell_2\mathbb{Z}^3$-periodic extension of $(\mathrm{pj}^{*}B,\mathrm{pj}^{-1}Y)$ (Section~I.6.3).
A function $f$ with $f(t+\ell_1w)=\boldsymbol{\chi}(w)\,f(t)$ satisfies
\[
f(t+\ell_2w)=f(t+\ell_1\mathrm{Jr}\,w)=\boldsymbol{\chi}(\mathrm{Jr}\,w)\,f(t)=\boldsymbol{\chi}(w)\,f(t),
\]
because $\mathrm{Jr}$ is odd; and a periodic function of side $\ell_1$ is periodic of side $\ell_2$.
The map $\mathrm{pj}^{*}$ is this inclusion (Section~I.6.3).
By (f3) for $\ell_1$ and for $\ell_2$, both of which lie in $\mathrm{Cl}$,
\[
\mathrm{Op}(\mathrm{pj}^{*}B,\mathrm{pj}^{-1}Y)\,\mathrm{pj}^{*}f=\mathrm{Op}_{\mathbb{L}}(\tilde{B},\tilde{Y})\,f=\mathrm{pj}^{*}\left(\mathrm{Op}(B,Y)\,f\right).
\]
For no.~14 the sides at level $k+1$ are $\ell_1/L$ and $\ell_2/L=\mathrm{Jr}\cdot\ell_1/L$.
Let $\epsilon$ be one of the two characters of Lemma~3, and $f'$ a function of level $k+1$ with $f'\left(t'+(\ell_1/L)w\right)=\epsilon(w)\,f'(t')$.
Then
\[
f'\left(t'+(\ell_2/L)w\right)=f'\left(t'+(\ell_1/L)\mathrm{Jr}\,w\right)=\epsilon(\mathrm{Jr}\,w)\,f'(t')=\epsilon(w)\,f'(t'),
\]
and $\mathrm{pj}^{*}$ at level $k+1$ is this inclusion.
By 2.4\,(i) and (Of14), the operator on the side $\ell_i$ is the restriction of $f'\mapsto f'_L$ on the covering lattice; hence
\[
(\mathrm{pj}^{*}f')_L=f'_L=\mathrm{pj}^{*}(f'_L).
\]
\end{proof}

\subsection*{5. Conditions on the constants}

\noindent
\begin{tabular}{p{0.24\textwidth}p{0.42\textwidth}p{0.22\textwidth}}
\hline
Condition
& Where it is used
& Line of Section~I.8.2; label of Proposition~14.2\,(b)\\
\hline
$L$ odd, $L>1$
& $L$ odd: Fact~1 (the blocks partition and are nested; no site lies on a face), and through it all of Sections 1--4.
$L>1$: (Of7) ($\eta_k\le L-\eta_k$), (Of8) ($1\le L-1$), and $b_k>0$ ((Of6); Section~I.3.4)
& (C3); (L5)\\
$M$ a power of $L$, $M\ge L$
& Fact~1 ($M$-cubes are blocks $\mathrm{Bl}_{k+\mathrm{jm}}$), (1.c); $M$ odd in 3.6\,(A)
& (C5); (M1)\\
$L-1<M$
& Fact~4 ($(L-1)/2\le M$); 3.10\,(C); 4.3\,(b4)
& (C5); from (M1): $M\ge L^2$\\
$b>0$, hence (with $L>1$) $b_k>0$ and $b_L>0$
& (Of6), (K3): existence of $(b_k+i\,\gamma_5\,y)^{-1}$ and $(b_k+b_L+i\,\gamma_5\,y)^{-1}$; (Of7): $\boldsymbol{\alpha}_{y}$, $\boldsymbol{\beta}_{y}$; (Of9), (Of10): $b_k^{-1}$, $\mu/b_k$, $(b_k\cdot\eta_k)^{-1}$, and through them 3.7 and 4.3\,(b3)
& (C0): $b=1$; (C4): $b_L$; Section~I.3.4; (L5)\\
the standard weights, through $\sum_iM_i(x)=(\rho/2)\cdot\mathrm{sw}$ ((I.m1), which rests on (I.1.1))
& (Of3), (K4)
& (C0)\\
$p_{0,k}>0$ (a consequence of $e_k<1$)
& (o10) only (3.11)
& (C8); (A10), first inequality ($v\ge9$ on $(0,e_0]$)\\
$\ell\in\mathrm{Cl}$, used in the forms: $\ell$ a multiple of $L\cdot M$; $M<\ell/2$; $2L<\ell$; $2M+1<\ell$; $2M\cdot(L+2)<\ell$; $M\cdot(L+5)<\ell$; $\ell/L$ a multiple of $M$ with $2M<\ell/L$ (that is, $2L\cdot M<\ell$)
& 2.4; Lemma~1\,(c); 3.8; 4.3 ((b5), the sums that define $\varepsilon$, $m$); 1.5 and 3.9 (the last form; in 3.9 also $\ell$ a multiple of $L\cdot M$); 3.2\,(D) ($\ell$ a multiple of $L\cdot M$); 3.6\,(C) ($M<\ell/2$)
& class of sides (Section~I.1.1)\\
\hline
\end{tabular}

\medskip
The numbers $\alpha$, $\rho$, $\mu$, $e_k$ enter the formulas only as numbers.
As in the paragraph ``Uniformity'' of 3.2, Propositions~25.1--25.3 hold for every complex $\mu$, every $e_k>0$, every $\alpha>0$ and every real $\rho$, with the following exceptions; the third of them concerns the side, not the constants.
\begin{itemize}
\item[(i)]
(o10) (3.11; the last paragraph of 4.3\,(b5)) requires $p_{0,k}>0$.
This holds for $e_k<1$, hence for every $e_k\in(0,e_0]$ by the first inequality of (A10) of Proposition~14.2\,(b); it is the row $p_{0,k}>0$ of the table.
\item[(ii)]
The last two sentences of (o8) of Proposition~25.2, and the assertion of Proposition~25.3\,(b5) on $\acute{E}_k$, $E^{d}_k$, $E_k^{\mathrm{det}}$, ${\delta E^{\mathrm{det}}}_k$ (3.10\,(C), the paragraph on these numbers; 4.3\,(b5), the paragraph on the numbers of Section~I.5.7), are stated in the setting of Section~I.5.7: data (SF) (Section~I.1.8: $\ell\in\mathrm{Cl}$, $\mu\in(0,1]$, $e_k\in(0,e_0]$), a connection $B\in\tilde{G}_k$, and a family (p5) of Section~I.5.8.
\item[(iii)]
The operation $L$ of (o7) (3.9; 4.3\,(b5)) is stated for $\ell$ and $\ell/L$ in $\mathrm{Cl}$.
\end{itemize}

\section*{Chapter 24. One renormalisation step without large fields, with an external connection, at every level $k\ge0$: proofs of Propositions 24.1--24.3}

\subsection*{0. Conventions of this chapter}

\subsubsection*{0.1. Setting and definitions}

Throughout the chapter the setting is (ST1)$^{\mathrm{ch}}$ of Section~G3.4.0.
The step map of level $k$ is introduced in Definition~24.D of Section~G3.4.1.
That definition consists of the clauses (D0)--(D8) and the displays (8.1)$^{\mathrm{ch}}$--(8.6)$^{\mathrm{ch}}$, and it ends with the order in which the clauses of the Propositions 23.m and 24.m are proved.
Definition~24.I of Section~G3.4.2 introduces the four transformations on arguments, together with the notions of invariance relative to a polymer and relative to a set of arguments.

\subsubsection*{0.2. Arrangement}

The clauses of Propositions~24.1, 24.2 and~24.3 are established in the order fixed at the end of Definition~24.D.
The sections of this chapter, their content, and the step of that order to which each of them belongs are as follows.

\medskip
\noindent
\begin{tabular}{@{}p{0.09\linewidth}p{0.6\linewidth}p{0.24\linewidth}@{}}
Section & Content & Step of the order of Definition~24.D\\
1 & statement of Proposition~24.3 & \\
2 & Proposition~24.3\,(adm), (q1) & (2)\\
3 & Proposition~24.3\,(q5), (al), (ga), (be), (de), (ze), (et); among these, (ze), the fourth sentence of (al) and the last sentence of (ga4) only for $E^{(i),*}$, $E^{(i),x}$, $E^{(\le3),x}$, $E^{(\le5),x}$, $\mathrm{Km}$, $\mathrm{Hm}$ and for the terms of the series (8.4) & (4)\\
4 & statement of Proposition~24.1 & \\
5 & Proposition~24.1\,(a), first sentence together with its display; Proposition~24.1\,(b), first sentence (the same in the reading (o)) & (5)\\
6 & Proposition~24.3\,(ze), the fourth sentence of (al) and the last sentence of (ga4) for $E^{\#,x}_k(m,\mathscr{E})(\mathbf{Y})$; the last sentences of Proposition~24.3 & (6)\\
7 & statement of Proposition~24.2 & \\
8 & Proposition~24.2, with the rules (Pr1)--(Pr5) & (6)\\
9 & Proposition~24.1\,(f), (g); the remaining sentences of (a); (b) from its second sentence on; the first clause of (d) (in 9.6\,(0)); (c); the rest of (d); (e); (i) & (7)\\
10 & the case (o): the places at which the second slot is used & \\
11 & conditions on the constants & \\
\end{tabular}
\medskip

The displays (8.1)$^{\mathrm{ch}}$, (8.2), (8.3)$^{\mathrm{ch}}$, (8.3a)$^{\mathrm{ch}}$, (8.3$_0$)$^{\mathrm{ch}}$, (8.4) and (8.6)$^{\mathrm{ch}}$ are those of Definition~24.D and of Propositions~23.1 and~23.2.
The display (24.1) is the one in Proposition~24.3\,(ga1), and the displays of this chapter are numbered (24.2), (24.3), and so on.
A section number $n.m$ without prefix, preceded or not by the word Section, refers to a section of this chapter.

\subsubsection*{0.3. Constants and conditions in the form in which they are used}

Let $e_k$ be the number of (ST1)$^{\mathrm{ch}}$.
The conditions that Proposition~14.2\,(b) imposes on $e_0$ are required for all $v\ge-\log e_0$.
As $e_k<e_{k+1}\le e_0$, they hold at $v=-\log e_k$ and at $v=-\log e_{k+1}$.
Among them are the following, written with $e_k$ and labelled as in Proposition~14.2.
\begin{itemize}
\item[(A1)] $p_{0,k}\le e_k^{-\frac{\varepsilon}{4}}$;\quad $12\exp\left(-\frac{p_{0,k}^2}{2}\right)\le e_k^7$;\quad $e_k^{\frac{5\varepsilon}{2}}\le\frac12$;\quad $c_{\mathrm{fl}}\cdot e_k^{\varepsilon}\le1$;\quad $3M\cdot e_k^{\varepsilon}\le1$.
\item[(A2)] $2\,C_w^2\cdot c_{\gamma}^2\cdot c_C\cdot e_k^{\frac{\varepsilon}{4}}\le1$.
\item[(A3)] $16\,c_X\cdot c_{\gamma}\cdot C_w\cdot C_T^2\cdot M^3\cdot e_k^{\frac{\varepsilon}{2}}\le1$;\quad $96\,M^3\cdot e_k^{\frac14}\le1$.
\item[(A4)] $32\,C_w\cdot C_T\cdot e_k^{\frac{3\varepsilon}{2}}\le1$.
\item[(A5)] $88\,M^3\cdot e_k^{\frac14-6\varepsilon}\le1$.
\item[(A6)] $112\,C_w^2\cdot c_{\gamma}\cdot e_k^{2\varepsilon}\le3$.
\item[(A7)] $c_S\cdot e_k^{\frac14-5\varepsilon}\le1$;\quad $c_S\cdot L^3\cdot e_k^{\varepsilon}\le1$;\quad $\mathbf{H}_{0}=36\,c_X^2\cdot K_0^2\cdot e_k^{\frac14-5\varepsilon}\le1$;\quad here $c_S=10^5\,c_X^5\cdot K_0^5\cdot\exp(6\kappa)$.
\item[(A8)] $e_k^{\varepsilon}\le\frac14$;\quad $e_k^{\frac14}+e_k^{\varepsilon}\le1$.
\item[(A10)] the third and the fourth inequality: $8\,C_B\cdot p_k\le e_k^{-\frac34+3\varepsilon}$;\quad $8\,e_k^{\frac34-\frac{7\varepsilon}{2}}\le1$;\quad here $p_k=v_k^p$ and $v_k=-\log e_k$.
\end{itemize}
The constants chosen earlier satisfy the following conditions; after each of them the conditions from which it follows are named, where (L1)--(L5) and (M1)--(M6) are the conditions on $L$ and on $M$ in Proposition~14.2\,(b), and (C0)--(C4) are the lines of Section~I.8.2.
\begin{itemize}
\item[(K1)] $L$ is odd, $M$ is a power of $L$ (so that $M$ is odd), and $M>4L$, by (L5), (M1), (M2); moreover $L\ge4^8$ and $L^{-\frac34}\le\frac12$, by Proposition~14.2\,(f), as a consequence of (L1).
\item[(K2)] $c_X=54$, $\kappa_0=54\log72$ and $K_0=\exp(\kappa_0)/36$, so that $\kappa_0\ge c_X\ge1$ and $K_0\ge1$, by (C0); and $\kappa\ge7\kappa_0+7$, by (C2).
\item[(K3)] $L\cdot(\kappa-6\kappa_0-4)-1\ge\kappa$, by (L3).
\item[(K4)] $\kappa_{\mathrm{w}}=\kappa+\kappa_0+3$, by (C4).
\item[(K5)] $M^{\alpha_0}\ge1+\exp(\hat{\kappa})$, where $\hat{\kappa}=N_c\cdot\kappa+\log(6c_X)$ and $N_c=(2L+1)^3$, by (M4).
\item[(K6)] $C_w\ge4\exp(2\gamma_{\mathrm{w}})\cdot M^{27\alpha_0}$, $C_w\ge1$ and $c_{\gamma}\ge1$, by Proposition~14.2\,(a) (the definitions of $C_{w}^{(1)}$ and of $C_w$) and Proposition~14.2\,(f).
\item[(K7)] $0<1-\alpha<\varepsilon\le\frac{1}{200}$, by (C0) (where $\bar{\varepsilon}=\frac{1}{200}$), by (C1) and by the condition (S1) of Proposition~14.2\,(b).
\item[(K8)] $C_P\le C^P\le C_w$ and $\gamma_{\mathrm{w}}\le\gamma^P\le\gamma_P$, by Proposition~14.2\,(f); $L^{\frac34}\ge110\ge\frac83$, by Proposition~14.2\,(f), as a consequence of (L1); and $M\ge M^P$, by (M6).
\end{itemize}
By (A8) one has $e_k<1$; consequently $e_k^{a}\le e_k^{a'}$ whenever $a\ge a'$.
Moreover
\[
h_{k,2}=e_k^{\varepsilon}\cdot h_{k,1}\le h_{k,1},\qquad h_{k,1}^2=e_k^{-\frac12},\qquad\sqrt{48}\le7.
\]

In Sections~5 and~9 the following rates and the following number are used:
\begin{gather*}
\mathrm{rt}_1:=\kappa-\kappa_0-1,\qquad\mathrm{rt}_2:=\kappa-2\kappa_0-2,\qquad\mathrm{rt}_4:=\kappa-3\kappa_0-3,\qquad\mathrm{rt}_3:=\kappa-5\kappa_0-4,\\
\mathrm{rt}_5:=\mathrm{rt}_2-\kappa_0=\kappa-3\kappa_0-2,\qquad\mathrm{rt}_6:=\mathrm{rt}_4-2\kappa_0=\kappa-5\kappa_0-3,\qquad c_S':=14\,c_X^2\cdot K_0^2\cdot\exp(3\kappa).
\end{gather*}
By (K2) the six rates are positive, and $\kappa_0\le\mathrm{rt}_1\le\kappa$.

The weights $h_k/2$ and $h_k/4$ denote the pairs $(h_{k,1}/2,h_{k,2}/2)$ and $(h_{k,1}/4,h_{k,2}/4)$, which are the weights of $(\psi,\chi)$.
The weight $(h_k/2;\mathsf{w}_3)$ carries in addition the weight $\mathsf{w}_3$ of $\Psi$, and $(h_k/4;\mathsf{w}_W)$ carries in addition the weight $\mathsf{w}_W$ of $W$, where $\mathsf{w}_W=1$ for $k\ge1$ and $\mathsf{w}_W=\sqrt{48}$ for $k=0$.
The number $\mathsf{w}_W$ is the number $\max\nu^{-\frac12}$ of (I.5.2).

\subsubsection*{0.4. Words and notation}

\emph{Unit points.}
For $k\ge1$ a unit point $\mathrm{tu}$ is a unit site $y$ of $T^0_k$; its weight is $\nu=1$, and $\int\mathrm{d}\mathrm{tu}:=\sum_y$.
For $k=0$ a unit point $\mathrm{tu}$ is a fermion site $x$ of $T^{\mathrm{f}}_0$; its weight is $\nu=1/48$, and
\[
\int\mathrm{d}\mathrm{tu}:=\int\mathrm{d}x=\frac{1}{48}\sum_c\sum_g.
\]
A unit index is a triple $u=(\mathrm{tu},\beta,\omega)$.
In Proposition~24.3\,(ga2), (ga4) the unit point itself, an element of the unit index set of Section~G3.2.0, is called a unit index and is written $u$.
The letter $u$ has this meaning also in the entries $H_k(s'',B)(x,u)$, $\bar{H}_k(s'',B)(x,u)$ and $\Gamma_k(s'',\mathscr{T})(u,u')$ in (e2) of Section~3.9, which are $4\times4$ blocks.
Everywhere else in this chapter an index $u$ is the index $(\mathrm{tu},\beta,\omega)$ of Section~I.3.6.
For a union $Z$ of $M$-cubes, the phrase ``$\mathrm{tu}$ (or $u$) lies in $Z$'' means that $y\in Z$, respectively that $x$ lies over $Z$.
The spinors on the unit points form the space $\mathfrak{S}(1,\ell)$ for $k\ge1$ and the space $\mathfrak{S}^{\mathrm{ch}}(1,\ell)$ for $k=0$.
The kernel of an operator $O$ defined on this space is taken with respect to $\int\mathrm{d}\mathrm{tu}$, that is,
\[
(Ov)(\cdot)=\int\mathrm{d}\mathrm{tu}'\,O(\cdot,\mathrm{tu}')\,v(\mathrm{tu}')
\]
(Section~I.1.6, (I.1.4), (I.L0.4)).

\emph{Contains; lies in.}
That an element contains a generator is meant in the sense of Section~I.5.1, and that a path lies in $Z$ in the sense of Section~I.1.7.

\emph{Reads; depends through.}
These words have the meaning of \cite[Definition~5.5(1)]{QEDI}: when all other objects are held fixed, the value is the same at any two admitted arguments that coincide at the places named.
The admitted arguments are specified wherever the words are used.

\emph{Parity.}
Let $\mathrm{Pty}$ be the automorphism that multiplies every generator by $-1$.
An element $F$ is called even if $\mathrm{Pty}\,F=F$.
A homomorphism that maps every generator to a linear combination of generators commutes with $\mathrm{Pty}$, and therefore maps even elements to even elements.

\emph{The algebras $\mathrm{Gl}_k(Z)$; local maps.}
Let $Z$ be a union of $M$-cubes.
Denote by $\mathrm{Gl}_k(Z)$ the subalgebra of $\mathrm{Gr}[\psi,\chi,\Psi,W]$ generated by the $\psi(\xi)$ with $\xi$ over $Z$, the $\chi(\boldsymbol{\zeta})$ with $\boldsymbol{\zeta}$ admissible for $Z$ (Section~I.1.7), and the $\Psi(u)$, $W(u)$ with $u$ in $Z$; at $k=0$ no $\Psi$ occurs.
A map $F$ from $\tilde{R}_k^{-}\times\mathrm{Sd}_k\times\mathrm{Box}_{k}$ to $\mathrm{Gr}[\psi,\chi,\Psi,W]$, the set $\mathrm{Sd}_k$ being that of Section~0.5, is called \emph{local in $Z$} if the following two conditions hold.
\begin{itemize}
\item[(l1)] The values of $F$ lie in $\mathrm{Gl}_k(Z)$.
\item[(l2)] $F(A,\mathscr{C},\mathsf{Z})=F(\tilde{A},\tilde{\mathscr{C}},\tilde{\mathsf{Z}})$ whenever $A=\tilde{A}$ and $\mathscr{C}=\tilde{\mathscr{C}}$ on the bonds that meet $Z$, and $\mathsf{Z}(\Upsilon)=\tilde{\mathsf{Z}}(\Upsilon)$ for every $\Upsilon$ contained in $Z$.
\end{itemize}
Such maps have the following three properties.
\begin{itemize}
\item[(Lo1)] If $Z\subseteq Z'$, every map that is local in $Z$ is local in $Z'$.
Indeed, by Section~I.1.7 the condition of kind (I) requires that $c$, $c'$ and the sites of the six paths lie in $Z$, and the condition of kind (II) requires that $x'=x^{(3)}$ and that $c$ or $c_{3}$ lie in $Z$.
Since a site of $Z$ is a site of $Z'$, a pair admissible for $Z$ is admissible for $Z'$.
Moreover a bond that meets $Z$ meets $Z'$, and an $\Upsilon$ contained in $Z$ is contained in $Z'$.
\item[(Lo2)] The maps that are local in $Z$ form an algebra, and this algebra is closed under pointwise limits.
Indeed, $\mathrm{Gl}_k(Z)$ is a subalgebra and a linear subspace of a finite dimensional space, hence closed, and the condition (l2) is preserved under sums, under products and under pointwise limits.
It follows that the algebra is stable under $\exp$, under series that converge at every argument, and under derivatives and integrals with respect to auxiliary parameters ($t$, $s$), provided that the function is local in $Z$ for every value of the parameter and that the derivative or the integral exists; for such derivatives and integrals are pointwise limits of linear combinations.
\item[(Lo3)] If $F_i$ is local in $Z_i$ for $i=1,\dots,n$, then $\prod_iF_i$ is local in $Z_1\cup\dots\cup Z_n$.
This follows from (Lo1) and (Lo2).
\end{itemize}

\emph{Families of parameters; the terms of (8.2).}
For a set $\mathrm{Ps}$ of parameter cubes, let $\mathbf{1}[\mathrm{Ps}]$ be the family whose value is $1$ at the cubes of $\mathrm{Ps}$ and $0$ at all other cubes.
For a union $U$ of $M$-cubes, the family $\mathbf{1}_U$ of (D1) of Definition~24.D is $\mathbf{1}[\mathrm{Ps}]$, where $\mathrm{Ps}$ is the set of the parameter cubes contained in $U$.
Let $X$ be a polymer, and let $P$ be a set of parameter cubes, of the kinds that occur, lying outside $P_0(X)$.
Put
\begin{gather*}
\mathrm{Slc}(X,P):=\left\{s\in\mathrm{Pd}_k:\ s_{\mathrm{pc}}=0\ \text{for every}\ \mathrm{pc}\ \text{outside}\ P_0(X)\cup P\right\},\\
\mathrm{Tm}[G](X,P):=\int_{[0,1]^P}\mathrm{d}s_P\,\frac{\partial}{\partial s_P}\left[G(s,X)\right]_{s=1\ \text{on}\ P_0(X),\ s=0\ \text{outside}\ P_0(X)\cup P},
\end{gather*}
so that
\[
\mathrm{Loc}[G](\mathbf{Y})=\sum_{(X,P):\,Y'(X,P)=\mathbf{Y}}\mathrm{Tm}[G](X,P).
\]
For $G(\cdot,X)$ as in (D3) one has
\begin{equation}
\mathrm{Tm}[G](X,P)=\sum_{\mathrm{Psb}\subseteq P}(-1)^{\#|P|-\#|\mathrm{Psb}|}\,G\left(\mathbf{1}[P_0(X)\cup\mathrm{Psb}],X\right).\tag{V}
\end{equation}
Indeed, the function $s_P\mapsto G(s,X)$ is analytic in a neighbourhood of $[0,1]^P$, and for a function $g$ of one variable $\int_0^1\mathrm{d}s\,g'(s)=g(1)-g(0)$.
Applying this identity to the variables $s_{\mathrm{pc}}$, $\mathrm{pc}\in P$, one after the other, and expanding, one obtains (V), where $\mathrm{Psb}$ is the set of those variables that are set equal to $1$.
All families occurring in (V) belong to $\mathrm{Slc}(X,P)$.

\emph{Adding fields.}
For a family $v^{\#}$ as in (D1), let $\mathrm{sb}_v$ be the homomorphism that maps $\psi(\xi)$ to $\psi(\xi)+v(\xi)$ and $\chi(\boldsymbol{\zeta})$ to $\chi(\boldsymbol{\zeta})+v(\boldsymbol{\zeta})$ and fixes the other generators.
Then $F(\psi^{\#}+v^{\#})=\mathrm{sb}_v\,F$.

\subsubsection*{0.5. The two cases}

Sections~2--9 treat two cases at once.
The parts of the argument that concern the case (x) alone, and the places at which the two cases rest on different statements, are listed in Section~10.
The two cases are the following.
\begin{itemize}
\item[(x)] $\mathrm{Sd}_j:=\mathfrak{X}_j$ for $j=k,k+1$.
\item[(o)] $\mathrm{Sd}_j:=\{0\}$.
\end{itemize}
In the case (o), a family without slot $E=\{E(X;\mathscr{A})\}$ (Proposition~48.1) is written as the family $\mathscr{E}(X;\mathscr{A},0):=E(X;\mathscr{A})$ on $D_j\times\tilde{R}_j\times\{0\}$.
The symbols $K^{x,C}_j$, $K^{x}_j$, $\|\cdot\|^{x}_j$, $\mathrm{Ball}^{x}_{k}$ then denote $K^{C}_j$, $K_j$, $\|\cdot\|_j$, $\mathrm{Ball}_{k}$; the symbols $\varepsilon$, $m$, $R$, $\mathrm{nat}$, $L$ denote the operations of Propositions~48.2, 48.3 and~48.5; ${\delta E^{\mathrm{det}}}_k$ denotes $E_k^{\mathrm{det}}$; and $\boldsymbol{\lambda}'=0$ and $\mathscr{C}=\mathscr{C}^1=0$ in all gauge relations.
In this notation the clauses (k1)$^{\mathrm{o}}$--(k7)$^{\mathrm{o}}$ of Proposition~48.1 are the clauses (k1)--(k7) of Section~I.5.4 with $\mathrm{Sd}_j$ in place of $\mathfrak{X}_j$.
Likewise the displays (D1)--(D6) at $\mathscr{C}=0$ are the displays of the reading (o) of (D7), without the family $\mathfrak{P}_k$ and the entry $\mathfrak{P}_{k+1}$ of (8.6)$^{\mathrm{ch}}$, which are void in that reading.
The objects of the case (o) are $E^{(i)}$, $E^{\#}_k(m,E)$ and $E^{*}_k(m,E)$.

\emph{Meaning of analyticity.}
Let $\mathrm{Dm}$ be one of the open sets $\tilde{R}_k$, $\tilde{R}_k^{-}$, $\tilde{R}_{k+1}$ (these sets are open by DF0 of Section~0.6).
The phrase ``analytic (continuous) in $(A,\mathscr{C})$ on $\mathrm{Dm}\times\mathrm{Sd}_j$'' means, in the case (x), holomorphic (continuous) on the open set $\mathrm{Dm}\times\mathfrak{X}_j$, and, in the case (o), holomorphic (continuous) in $A$ on $\mathrm{Dm}$.
The phrase is read in the same way when further variables ($s$, $t$, $\mathsf{Z}$) are present.

\emph{Properties of the pair $(\mathrm{Sd}_k,\mathrm{Sd}_{k+1})$.}
In either case the pair has the following six properties.
\begin{itemize}
\item[(Sd1)] $0$ belongs to $\mathrm{Sd}_k$ and to $\mathrm{Sd}_{k+1}$.
\item[(Sd2)] $\tilde{R}_k+\mathrm{Sd}_k$ is contained in $\tilde{R}_k+\mathfrak{X}_k$; the latter set is contained in $\tilde{G}_k$ by DF1, and it is the set on which Proposition~9\,(d) holds.
The inclusion holds because $0\in\mathfrak{X}_k$.
\item[(Sd3)] For $j=k$ and for $j=k+1$ the set $\mathrm{Sd}_j$ is invariant under $\mathbb{G}_j$, under $\mathscr{C}\mapsto-\mathscr{C}$ and under $\mathscr{C}\mapsto\overline{\mathscr{C}}$.
In the case (x) this is DF3 at the levels $k$ and $k+1$; in the case (o) it holds because these maps are linear or antilinear.
\item[(Sd4)] If $\mathscr{C}'\in\mathrm{Sd}_{k+1}$, then $\mathscr{C}'_L\in\mathrm{Sd}_k$.
In the case (x) this is DF2\,(c); in the case (o) it holds because $0_L=0$.
\item[(Sd5)] The clauses (k1)--(k7) for the data and for the values, and, for $k\ge1$, the following two facts on the extraction.
\begin{itemize}
\item[(X1)] $m(\mathscr{E})$ is a complex number; it is complex linear in $\mathscr{E}$, and real for $\mathscr{E}\in K^{x}_k$; moreover
\[
|m(\mathscr{E})|\le c_{\mathrm{fl}}\cdot e_k^{\frac12}\cdot\|\mathscr{E}\|^{x}_k,\qquad|\varepsilon(\mathscr{E})|\le c_{\mathrm{fl}}\cdot\|\mathscr{E}\|^{x}_k.
\]
\item[(X2)] $R$ is linear and maps $K^{x,C}_k$ and $K^{x}_k$ into themselves; moreover, for $(\mathscr{A},\mathscr{C})\in\tilde{R}_k\times\mathrm{Sd}_k$,
\[
\|(R\mathscr{E})(X;\mathscr{A},\mathscr{C})\|_{h_k}\le\|\mathscr{E}(X;\mathscr{A},\mathscr{C})\|_{h_k}+\|\mathscr{E}(X;0,0)\|_{h_k}.
\]
\end{itemize}
In the case (x) these facts are Proposition~16.3\,(a), (b), (d); in the case (o) they are Proposition~48.2\,(a), (b), (d).
\item[(Sd6)] $\mathrm{Sd}_j$ is open in the case (x), by DF0, and equal to $\{0\}$ in the case (o).
Hence $\mathrm{Dm}\times\mathrm{Sd}_j$ is open in the case (x), and in both cases the words ``analytic in $(A,\mathscr{C})$'' have the meaning fixed above.
\end{itemize}

\subsubsection*{0.6. Statements that are quoted}

\paragraph{Domains.}
The facts DF0--DF3 below are \cite[Lemmas~17.3--17.6]{QEDI}.
They involve no fermion object.
They hold at every level $j\ge0$, for every side in $\mathrm{Cl}$ and for every $e_j\in(0,e_0]$, and therefore at the levels $k$ and $k+1$ of (ST1)$^{\mathrm{ch}}$.
\begin{itemize}
\item[DF0] (\cite[Lemma~17.3]{QEDI}.)
The sets $R_j$, $G_j$, $\mathfrak{X}_j$, $\mathrm{cf}\cdot\tilde{R}_j$ (for $0<\mathrm{cf}\le1$), $\tilde{R}_j^{-}$ and $\tilde{G}_j$ are open subsets of the space of the complex bond functions on $T_j$.
\item[DF1] (\cite[Lemma~17.4(a)]{QEDI}.)
The sets $\tilde{R}_j$ and $\tilde{R}_j+\mathfrak{X}_j$ are contained in $\tilde{G}_j$.
\item[DF2] (\cite[Lemma~17.5]{QEDI}.)
Let $\ell$ and $\ell/L$ lie in $\mathrm{Cl}$, and let $e_{k+1}=L^{\frac12}\cdot e_k$.
\begin{itemize}
\item[(a)] If $\mathscr{A}\in\tilde{R}_{k+1}$, then $\mathscr{A}_L\in L^{-\frac34}\cdot\tilde{R}_k$.
\item[(b)] The set $L^{-\frac34}\cdot\tilde{R}_k$ is contained in $\tilde{R}_k^{-}$.
\item[(c)] If $\mathscr{C}\in\mathfrak{X}_{k+1}$, then $\mathscr{C}_L\in\mathfrak{X}_k$.
\end{itemize}
In the proof one uses $(d\boldsymbol{\lambda})_L=d(\boldsymbol{\lambda}_L)$, where $\boldsymbol{\lambda}_L(c):=L^{\frac12}\cdot\boldsymbol{\lambda}(c/L)$.
\item[DF3] (\cite[Lemma~17.6]{QEDI}, without its clause on $\mathrm{Dom}_j$.)
The sets $R_j$, $\mathrm{cf}\cdot\tilde{R}_j$, $\tilde{R}_j^{-}$, $G_j$, $\tilde{G}_j$ and $\mathfrak{X}_j$ are invariant under $\mathbb{G}_j$, under $A\mapsto-A$ and under $A\mapsto\overline{A}$.
The sets $\mathrm{cf}\cdot\tilde{R}_j$, $\tilde{R}_j^{-}$ and $\tilde{G}_j$ are invariant under $A\mapsto A-d\boldsymbol{\lambda}$ for every real function $\boldsymbol{\lambda}$ of $c$.
\end{itemize}
By Section~I.5.3, that is, by \cite[Definition~5.16]{QEDI}, the set $\tilde{R}_k^{-}$ consists of the sums $A_0+A_1$ with $A_1\in(1-e_k^{\varepsilon})\cdot R_k$ and $A_0$ as in the definition of $\tilde{R}_k$; hence $\tilde{R}_k^{-}\subseteq\tilde{R}_k$.
The thresholds of $R_k$ are
\[
\mathrm{t1}_k:=e_k^{-\frac34+\varepsilon}\ \ge\ \mathrm{t2}_k:=e_k^{-\frac34+2\varepsilon}\ \ge\ \mathrm{t3}_k:=e_k^{-\frac34+3\varepsilon}.
\]

\paragraph{Geometry.}
The following facts are those of \cite[Section~17.4.1]{QEDI}, where they bear the labels (g0)--(g4).
They concern sets of lattice sites, and $\mathrm{pc}$ denotes a parameter cube.
\begin{itemize}
\item[(g0)] No site of $\mathbb{L}_k$ lies on a face of an $M$-cube, and the $M$-cubes $\{c:|c-z|<M/2\}$, $z\in M\mathbb{Z}^3$, form a partition of the set of the sites of $T_k$ and of that of $T^0_k$.
Every $\Delta_y$ and every $L$-block lies in one $M$-cube, and every $M$-cube lies in one $LM$-cube.
\item[(g1)] For a cube $\mathrm{pc}$ of $\pi'_k$ (of $\pi''_k$) with centre $z$, let $\mathrm{pc}^{+}$ be the union of the $M$-cubes whose centres are at sup-distance at most $M$ (at most $LM$) from $z$.
Every lattice site at sup-distance at most $2L$ from a lattice site of $\mathrm{pc}$ lies in $\mathrm{pc}^{+}$.
The $M$-cubes that meet $\mathrm{pc}$ are precisely the $M$-cubes of $\mathrm{pc}^{+}$.
A parameter cube of any kind meets at most $N_c$ $M$-cubes.
\item[(g2)] For a non-empty link-connected set $\mathrm{Lnk}$ of parameter cubes of one kind, let $\mathrm{pol}(\mathrm{Lnk})$ be the union of the $M$-cubes that meet some cube of $\mathrm{Lnk}$.
Then $\mathrm{pol}(\mathrm{Lnk})$ is a polymer, and it contains the union of $\mathrm{Lnk}$.
\item[(g3)] Distinct link-components of a union of parameter cubes of one kind have no lattice site in common.
A union $\mathbf{Y}$ of $M$-cubes contains at most $|\mathbf{Y}|_M$ parameter cubes of each kind.
\item[(g4)] For every site $z$, the cubes of $\pi_k$ that meet $K_{z}$ form a link-connected set.
\end{itemize}
The link-components of a set of parameter cubes of one kind are the classes of the equivalence relation under which two of its cubes are equivalent if they are joined by a finite sequence of its cubes in which any two consecutive cubes are linked.
They are non-empty and link-connected, so that (g2) applies to each of them, and an empty set of parameter cubes has no link-component.

\paragraph{Interpolation and resummation.}
Lemmas~6.1 and~6.2 below are \cite[Lemmas~17.12 and~17.13]{QEDI}, respectively.
They involve no fermion object, and their values may lie in any finite dimensional normed space.

\begin{named}{Lemma 6.1}[{\cite[Lemma~17.12]{QEDI}}]
Let $\mathrm{Pc}$ be a finite set, let $P_0\subseteq\mathrm{Pc}$ and $\mathrm{rad}>1$, and let $f$ be a continuous map from $\left\{s\in\mathbb{C}^{\mathrm{Pc}}:\ |s_{\mathrm{pc}}|\le\mathrm{rad}\ \text{for all}\ \mathrm{pc}\in\mathrm{Pc}\right\}$ to a finite dimensional normed space, analytic in the interior, with $\|f\|\le\mathrm{Bq}$.
For a subset $P$ of $\mathrm{Pc}\setminus P_0$ let
\[
f_P:=\int_{[0,1]^P}\mathrm{d}s_P\,\frac{\partial}{\partial s_P}f(s)\qquad\text{at}\ s=1\ \text{on}\ P_0,\ s=0\ \text{outside}\ P_0\cup P.
\]
Then
\[
f(1)=\sum_Pf_P,\qquad\|f_P\|\le\mathrm{Bq}\cdot(\mathrm{rad}-1)^{-\#|P|}.
\]
\end{named}

\begin{named}{Lemma 6.2}[{\cite[Lemma~17.13]{QEDI}}]
Let $\kappa_0\le\mathrm{kq}\le\kappa$ and $\mathrm{Bq}_0\ge0$, and let the parameter cubes be those of $\pi_k$, or those of $\pi_k$, $\pi'_k$ and $\pi''_k$.
For every $X\in D_k$ let $f^X$ be a map as in Lemma~6.1, with $\mathrm{Pc}$ the set of all parameter cubes, $\mathrm{rad}=M^{\alpha_0}$, $P_0=P_0(X)$ and $\mathrm{Bq}=\mathrm{Bq}_0\cdot\exp\left(-\mathrm{kq}\cdot d_M(X)\right)$, and such that $f^X_P=0$ unless $Y'(X,P)$ is a polymer.
Then
\[
F(\mathbf{Y}):=\sum_{(X,P):\,Y'(X,P)=\mathbf{Y}}f^X_P
\]
satisfies
\[
\sum_{\mathbf{Y}\in D_k}F(\mathbf{Y})=\sum_{X\in D_k}f^X(1),\qquad\|F(\mathbf{Y})\|\le4\,c_X\cdot K_0\cdot\mathrm{Bq}_0\cdot\exp\left(-(\mathrm{kq}-\kappa_0-1)\cdot d_M(\mathbf{Y})\right).
\]
\end{named}
The conditions on the constants that Lemma~6.2 requires are (K5) and (K1).

\paragraph{Polymers.}
The following assertions are contained in \cite[Proposition~5.58]{QEDI}, where they bear the same labels, and they are proved in \cite[Theorem~13.2]{QEDI}.
Here $\ell\in\mathrm{Cl}$, and for (t5) also $\ell/L\in\mathrm{Cl}$.
\begin{itemize}
\item[(t0)] $d_M(\square)=0$.
\item[(t1)] $|X|_M\le c_X\cdot\left(d_M(X)+1\right)$.
\item[(t2)] $\sum_{X\in D_k,\,X\supseteq\square}\exp\left(-\kappa_0\cdot d_M(X)\right)\le K_0$; hence $\sum_{X\in D_k,\,X\subseteq Z}\exp\left(-\kappa_0\cdot d_M(X)\right)\le K_0\cdot|Z|_M$ for every union $Z$ of $M$-cubes.
\item[(t4)] Let $X_1,\dots,X_n$ be polymers, and assume that the graph on $\{1,\dots,n\}$ in which $i$ and $j$ are joined whenever $X_i\cup X_j$ is connected (through faces) is connected.
Then $X_1\cup\dots\cup X_n$ is a polymer, and
\[
d_M(X_1\cup\dots\cup X_n)\le\sum_id_M(X_i)+3(n-1).
\]
\item[(t5)] For every $X\in D_k$ there is exactly one $Y\in D_{k+1}$ with $\bar{X}^{L}=LY$, and $L\cdot d_M(Y)\le d_M(X)$.
\item[(t6)] $d+1\le\exp(d)$ and $d+1\le2\exp\left(\frac{d}{2}\right)$ for $d\ge0$.
\end{itemize}
Two polymers are called incompatible if their union is connected (through faces), and compatible otherwise.
Put $\mathrm{urs}^{T}(X_1):=1$ and
\[
\mathrm{urs}^{T}(X_1,\dots,X_n):=\sum_G(-1)^{\#|G|},
\]
where $G$ ranges over the connected graphs on $\{1,\dots,n\}$ all of whose edges join incompatible polymers, and $\#|G|$ is the number of the edges of $G$; the sum is $0$ if there is no such graph.
\begin{itemize}
\item[(CE1)] In the ring of the formal power series in the commuting variables $z_X$, $X\in D_k$,
\[
\sum_{\{X_i\}\ \text{pairwise compatible}}\ \prod_iz_{X_i}=\exp\left\{\sum_{n\ge1}\frac{1}{n!}\sum_{(X_1,\dots,X_n)}\mathrm{urs}^{T}(X_1,\dots,X_n)\prod_iz_{X_i}\right\},
\]
where on the left the sum ranges over the finite sets of distinct polymers, the empty set included.
\item[(CE2)] $|\mathrm{urs}^{T}(X_1,\dots,X_n)|$ is at most the number of the trees on $\{1,\dots,n\}$ all of whose edges join incompatible polymers.
\end{itemize}

\paragraph{Photon objects.}
The following facts are taken from \cite[Propositions~15.21, 15.26 and~15.30(ii)]{QEDI}, with the definitions of Sections~I.5.7 and~I.5.3, which are those of \cite[Sections~5.2.5 and~5.3.1]{QEDI}.
\begin{itemize}
\item[(P1)] $\{\Upsilon\}$ is an orthonormal basis of $V_{k}$ for the counting scalar product, and $\mathsf{Z}(\Upsilon)=\langle\mathsf{Z},\Upsilon\rangle$.
The set $\operatorname{supp}\Upsilon$, which is the set of the end points of the bonds in the support of $\Upsilon$, lies in one $L$-block or consists of the two end points of one bond that joins two $L$-blocks.
\item[(P2)] The number $C(b',\Upsilon):=(C\Upsilon)(b')$ is real, and $C(b',\Upsilon)=0$ unless $\operatorname{supp}\Upsilon$ lies in the union of the $L$-blocks (one or two) of the ends of $b'$.
Moreover $|C\mathsf{Z}|_\infty\le c_{C}\cdot\sup_\Upsilon|\mathsf{Z}(\Upsilon)|$ for $\mathsf{Z}\in V_{k}$.
\item[(P3)] Every $\tilde{g}\in\mathbb{G}_k$ maps $V_{k}$ isometrically onto itself and permutes the set $\{\pm\Upsilon\}$; moreover $\tilde{g}C=C\tilde{g}$.
The set $\mathrm{Box}_{k}$ and the measures $\mu_I$, $\hat{\mu}$ are invariant under $\mathbb{G}_k$ and under $\mathsf{Z}\mapsto-\mathsf{Z}$.
\item[(P4)] Let $\mathrm{gn}$ be a standard normal random variable, and let $d_V=2-2L^{-3}$.
Then $\delta\varepsilon^{0}_{k}=-d_V\cdot\log P\left(|\mathrm{gn}|\le p_{0,k}\right)$, and, if $p_{0,k}\ge1$,
\[
0\le\delta\varepsilon^{0}_{k}\le12\exp\left(-\frac{p_{0,k}^2}{2}\right).
\]
\item[(P5)] (\cite[Proposition~15.30(ii)]{QEDI}.)
$\mathrm{Dom}_j$ is the set of the real bond functions $A$ on $T^0_j$, with arbitrary harmonic part, such that $H_jA\in\frac14\cdot\tilde{R}_j$.
Let $\ell$ and $\ell/L$ lie in $\mathrm{Cl}$, and let $\mathrm{Dom}_{k+1}$ be formed with the side $\ell/L$ and with $e_{k+1}=L^{\frac12}\cdot e_k$.
If $A'\in\mathrm{Dom}_{k+1}$ and $\zeta\in\mathrm{Box}_{k}$, then $A^{\min}_{k}(A')+\mathrm{Sq}_{k}\zeta$ lies in $\mathrm{Dom}_k$.
The hypotheses of \cite[Proposition~15.30]{QEDI} are satisfied at $e_k$: the two inequalities of (A10) written in Section~0.3; $e_k<1$, by (A8); $\varepsilon\le\frac{1}{200}$, by (K7); $L^{\frac34}\ge\frac83$, by (K8); $C_w\ge C_P$ and $\gamma_{\mathrm{w}}\le\gamma_P$, by (K8); and the conditions (A1) and (A2) of \cite[Section~5.6.2]{QEDI} at $e_k$, which are the conditions (A1), with all five of its inequalities, and (A2) of Section~0.3, with the same constant $c_{\mathrm{fl}}$.
Of (A1), the proof of \cite[Proposition~15.30(ii)]{QEDI} requires the first inequality alone.
That proof rests on \cite[Propositions~5.31, 5.32, 5.34, 5.35 and~5.36]{QEDI}, the last of which gives the reality of $A^{\min}_{k}$ and of $\mathrm{Sq}_{k}$, and on \cite[Theorem~14.15(b)]{QEDI}.
The five propositions hold by \cite[Theorem~15.37]{QEDI} for $\alpha\in\left(1-\frac{1}{200},1\right)$, $L$ odd, $L\ge5$, $\alpha_0\le\alpha^{P}$, $M$ a power of $L$ and $M\ge M^{P}$; these hypotheses hold by (K7), by (K1), by $\alpha_0=\min(\alpha^{F},\alpha^{P})$ (see (C4)) and by (K8).
\cite[Theorem~14.15]{QEDI} holds for $k\ge0$, for odd $L\ge3$ and for every side that is a multiple of $L$; its hypotheses hold by (K1) and because $\ell\in\mathrm{Cl}$ is a multiple of $L$, so that it imposes no further condition on the constants.
\end{itemize}

\paragraph{Primitives.}
The properties W1--W8, W10--W17 for $\mathrm{par}=(\alpha_0,\gamma_{\mathrm{w}},\kappa_{\mathrm{w}};C_w)$, which are the properties 1--8 and 10--17 of \cite[Section~5.4.2]{QEDI} read by the rules of Section~I.5.8, are stated in Sections~G1.11 and~G4.0.4; W3, W4 and the kernels of (p1), (p2) are stated in Section~G4.0.4.
At the places where a fermion index occurs they are used in the forms (G3.T-W) of Annex~A, Section~G3.0.4.
At $k=0$ the diagonal family $H_0(s')$ has the property W3, and $H_0(s'',A)$ has the property W5, with the constant $4\exp(2\gamma_{\mathrm{w}})\cdot M^{27\alpha_0}\le C_w$, by W17 and (K6).

\paragraph{The constant $c_{\gamma}$.}
By Section~I.5.8, every index $\mathrm{ti}$ (a fine or unit bond, a basis element, or a unit point) satisfies
\begin{gather*}
\sum_{\beta'}\int\mathrm{d}\mathrm{tu}'\,\exp\left(-\gamma_{\mathrm{w}}\cdot d(\mathrm{ti},\mathrm{tu}')\right)\le c_{\gamma},\\
\sum_{b'}\exp\left(-\gamma_{\mathrm{w}}\cdot d(\mathrm{ti},b')\right)\le c_{\gamma},\qquad\sum_{\Upsilon'}\exp\left(-\gamma_{\mathrm{w}}\cdot d(\mathrm{ti},\Upsilon')\right)\le c_{\gamma};
\end{gather*}
in the first inequality $\mathrm{tu}'$ ranges over the unit points, with the weight $1/48$ at $k=0$, and in the second $b'$ ranges over the positively oriented unit bonds.

\paragraph{Grassmann calculus.}
In this chapter a clause of Proposition~15 is denoted by its tag: (N0), (N1), (N2), (N3)$^{\mathrm{w}}$, (N4), (N4$'$), (N5) and (N6) denote the clauses (a), (b), (d), (f), (g), (h), (j) and (k), which bear these tags, and (n) denotes the clause (n).
The clause (c), on row sums, is denoted (rs); the clause (e), on moments, which bears the tag (m)$^{\mathrm{w}}$, is denoted (mo); and the clause (i), Wick's theorem, is denoted (Wi).
In (rs) and (N2) the index of a spinor function consists of a site, a pair or a unit point together with a spinor index $\beta$, and $|f|_\infty=\sup|f(\cdot)|_c$ (Section~I.1.5).
By (rs) the norm $|O|_{1,\infty}$ is submultiplicative.

\paragraph{Classical facts.}
Cauchy's formula and Cauchy's estimates in one and in several variables, the identity theorem, and the analyticity of uniform limits of analytic functions are taken from \cite[Chapter~10]{Rudin1987} and \cite[Section~2.2]{Hormander1990}, and Fubini's theorem from \cite[Chapter~8]{Rudin1987}.
M\"obius inversion on the subsets of a finite set \cite[Section~2.1]{Stanley2012} states: if $F(U)=\sum_{Z\subseteq U}f(Z)$ for all subsets $U$, then
\[
f(Z)=\sum_{U\subseteq Z}(-1)^{\#|Z\setminus U|}F(U).
\]
By Cayley's formula \cite[Chapter~2]{vanLintWilson2001}, if $\mathrm{dg}_1,\dots,\mathrm{dg}_n\ge1$ and $\sum_j\mathrm{dg}_j=2n-2$, the number of the trees on $\{1,\dots,n\}$ in which the vertex $j$ has the degree $\mathrm{dg}_j$ for every $j$ is
\[
\frac{(n-2)!}{\prod_j(\mathrm{dg}_j-1)!}.
\]
Wick's theorem for the integral (I.3.5) is the clause (Wi) of Proposition~15.

\subsection*{1. Statement of Proposition~24.3}

Proposition~24.3 of Part~II describes the structure of the step map of Definition~24.D, that is, what the objects attached to a polymer contain and what they read.
It is formulated in the setting (ST1)$^{\mathrm{ch}}$ of Section~G3.4.0, with the following data:
a pair $(m,\mathscr{E})\in\mathrm{Ball}^{x}_{k}(4)$;
$A\in\tilde{R}_k^{-}$ and $\mathscr{C}\in\mathfrak{X}_k$, with $\mathscr{T}=A+\mathscr{C}$;
an element $\mathsf{Z}\in\mathrm{Box}_{k}$;
and a polymer $\mathbf{Y}\in D_k$.

Its clauses are stated in the following order, and they have the following subjects.
Clause (adm) concerns the pairs that are admissible for a polymer in the sense of Section~I.1.7:
their admissibility for the polymers that contain the given one, their paths, and the sup-distance of their photon sites from the polymer.
Clause (q1) concerns the domains:
the inequalities between the radii $\mathrm{rad}_1$, $\mathrm{rad}_2$ and $\mathrm{rad}_4$,
a bound on the three norms of $\mathscr{Z}(s')$ for $s\in\mathrm{Pd}_k$,
and the sets that contain $A+t\cdot\mathrm{Zf}$ and $\mathscr{T}+t\cdot\mathrm{Zf}$ when $|t|\le\mathrm{rad}_1$ and $\mathrm{Zf}$ is a complex bond function whose three norms are at most $\frac{1}{2}e_k^{-\frac{\varepsilon}{2}}$.
Clause (q5) concerns regularity:
that of $\mathscr{Z}(s')$,
that of the entries of $H_{k}^{\#}(s'',B)$, of $\Gamma_k(s'',B)$ and, for $k\ge1$, of $M_k(X,B)$ at $B=\mathscr{T}+t\cdot\mathscr{Z}(s')$,
and the hypothesis of Definition~24.D\,(D3) for every function $G(\,\cdot\,,X)$ to which $\mathrm{Loc}$ is applied in Definition~24.D\,(D4).
Clause (al) concerns sub-polymers and the datum:
the polymers and the parameter cubes of the terms of (8.2) attributed to $\mathbf{Y}$,
the polymers of the sets and families of (8.4) whose union is $\mathbf{Y}$,
the polymers $X\in D_k$ with $\bar{X}^{L}=LY$ for $Y\in D_{k+1}$,
and the fact that the objects attached to $\mathbf{Y}$ depend on the datum $(m,\mathscr{E})\in\mathrm{Ball}^{x}_{k}(4)$ only through $m+m(\mathscr{E})$ and the $(R\mathscr{E})(X)$ with $X\subseteq\mathbf{Y}$,
together with a remark on the kernels of $\mathscr{E}$ through which $m(\mathscr{E})$ and $(R\mathscr{E})(X)-\mathscr{E}(X)$ depend on $\mathscr{E}$.
Clause (be) concerns the places where the exact operators and forms stand in the displays for an object attached to $\mathbf{Y}$:
the H\"older quotients $\delta_{\alpha,B}$ and $\delta_{\alpha,-B}$, the substitutions, the rows of $T^{(i)}(\mathscr{T},\cdot\,)$ and $\bar{T}^{(i)}(\mathscr{T},\cdot\,)$ for $k\ge1$, the forms $\langle\bar{\psi},\psi\rangle_{\square}$, the elements $\mathrm{Vx}(\mathscr{Z}(s'),\square)$ for $k=0$, and the map $C$;
it also says which of these read the connection or $\mathscr{Z}(s')$, and on which bonds.
Clause (ga) concerns the primitives at the families $s=(s',s'')\in\mathrm{Pd}_k$ that vanish at every parameter cube not contained in $\mathbf{Y}$, and it has four parts:
(ga1) concerns the entries of $H_{k}(s')$, $C$ and $C_{k}^{1/2}(s')$ that vanish, and the resulting display (24.1) for $\mathscr{Z}(s')(b)$, $b$ a fine bond that meets $\mathbf{Y}$;
(ga2) concerns the rows of $H_{k}(s'',B)$, $\bar{H}_{k}(s'',B)$ and $\Gamma_k(s'',B)$ and, for $k\ge1$, the kernel of $M_k(X,B)$;
(ga3) concerns the bonds through which these rows and $M_k(X,B)$ depend on $B$;
(ga4) lists the entries of the primitives that stand in the displays (8.1)$^{\mathrm{ch}}$--(8.4) for an object attached to $\mathbf{Y}$, and compares the objects attached to $\mathbf{Y}$ that are formed with two systems of primitives.
Clause (de) concerns the integrations in the displays for an object attached to $\mathbf{Y}$, in $t$, in $s_P$, in the generators $W(u)$ and in the coordinates $\mathsf{Z}(\Upsilon)$, with their ranges and their measures.
Clause (ze) concerns the generators that the objects attached to $\mathbf{Y}$ contain, the parity of these objects, the coordinates of $\mathsf{Z}$ and the bonds through which they depend on $\mathsf{Z}$ and on $(A,\mathscr{C})$, and the continuity and analyticity of their coefficients;
it also contains the fact that two polymers without a common $M$-cube have no fermion site, no unit site and no contained basis element in common.
Clause (et) concerns the terms $(X,P)$ of (8.2) that vanish because $Y'(X,P)$ is not connected;
the identity
\[
\sum_{\mathbf{Y}\in D_k}\mathrm{Loc}[G](\mathbf{Y})=\sum_{X\in D_k}G(1,X),
\]
which follows from this for every function $G(\,\cdot\,,X)$ to which $\mathrm{Loc}$ is applied in Definition~24.D\,(D4), $1$ being the family whose value is $1$ at all parameter cubes;
and the expansion of the exponential of $\sum_{\mathbf{Y}\in D_k,\ \mathbf{Y}\subseteq U}E^{(\le5),x}(\mathbf{Y})$, for a union $U$ of $M$-cubes, as the sum of the $\mathrm{Km}(Z)$ over the unions $Z$ of $M$-cubes contained in $U$.

The last assertions of Proposition~24.3 concern the two other parts of $\mathscr{E}^{+}(Y)$, which fall outside (ga).
The element $(R\mathscr{E})^{\mathrm{nat}}(X)$ contains the numbers $\mathrm{al}_{0}$, $\mathrm{al}_{2}$ and the homomorphism $\boldsymbol{\Theta}_{X,\mathscr{T}}$ of Section~I.5.6.
The element ${\delta E^{\mathrm{det}}}_k(X)$ contains the family (p5) of Section~I.5.8, through the kernel of $S_{k,y}(X,B)$ at indices over $X$ (W12), and it contains the exact operators in the displays of $\acute{E}_k$ (Section~I.5.7).
The exact operators of all three parts are described in Propositions~25.1--25.3.
Finally, all assertions of Proposition~24.3 hold in the reading~(o) as well.

\subsection*{2. Proof of Proposition~24.3: (adm), (q1)}

\subsubsection*{2.1. (adm)}

\begin{proof}[Proof of Proposition~24.3\,(adm)]
Let $Z$ and $Z'$ be unions of $M$-cubes with $Z\subseteq Z'$,
and let the pair $(x,x')\in\mathrm{Pz}_k$, $x=(c,g)$, $x'=(c',g')$, be admissible for $Z$ in the sense of Section~I.1.7.

\emph{Pairs of kind (I).}
Here $c$ and $c'$ lie in $Z$, and so do the six paths $\boldsymbol{\Gamma}^{\pi}(c,c')$.
As $Z\subseteq Z'$, the same holds with $Z'$ in the place of $Z$;
thus the pair fulfils the condition of kind (I) for $Z'$.
Every bond of one of these paths joins two consecutive sites of the path, and both of them lie in $Z$;
so the bond is contained in $Z$.
If $c=c'$, that is, for a pair inside a cluster, the condition reduces to $c\in Z$;
then $(\tau B)(c,c)=0$ and $U_B(x,x')=1$ (Section~I.1.1, (I.1.6)), and no path enters.

\emph{Pairs of kind (II).}
Here $x'=x^{(3)}$, so that $c'=c_{3}=c+\eta_k\,g e_1$ by (I.1.2), and $c$ or $c_{3}$ lies in $Z$, hence in $Z'$;
thus the pair fulfils the condition of kind (II) for $Z'$.
The sites $c$ and $c_{3}$ differ in a single coordinate, by $\eta_k$.
Consequently each of the six paths $\boldsymbol{\Gamma}^{\pi}(c,c_{3})$ is the photon bond $(c,c_{3})$.
This bond has an end in $Z$, that is, it meets $Z$.

\emph{Distances.}
For a pair of kind (I) both photon sites belong to $Z$.
For a pair of kind (II) one of them belongs to $Z$, and the other one is at sup-distance $\eta_k$ from it.

Taking for $Z$ a polymer $X$ and for $Z'$ a polymer that contains $X$, one obtains the three assertions of (adm).
\end{proof}

\subsubsection*{2.2. (q1)}

\begin{proof}[Proof of Proposition~24.3\,(q1)]
\emph{The radii.}
By (A1) one has
\[
\frac{\mathrm{rad}_2}{\mathrm{rad}_1}=2e_k^{\frac{5\varepsilon}{2}}\le1,
\qquad
\frac{\mathrm{rad}_4}{\mathrm{rad}_1}=2e_k^{\frac{1}{2}-\frac{3\varepsilon}{2}}\le2e_k^{\frac{5\varepsilon}{2}}\le1,
\]
the middle inequality holding because $\frac{1}{2}-4\varepsilon>0$ by (K7).
Since $M\ge1$, (A5) gives $e_k^{\frac{1}{4}-6\varepsilon}\le\frac{1}{88}$.
Hence
\[
\mathrm{rad}_4^{-1}=e_k^{\frac{1}{4}-3\varepsilon}\le e_k^{\frac{1}{4}-6\varepsilon}\le\frac{1}{88},
\qquad
\mathrm{rad}_2^{-1}=e_k^{\frac{3}{4}-7\varepsilon}=e_k^{\frac{1}{4}-3\varepsilon}\cdot e_k^{\frac{1}{2}-4\varepsilon}\le\mathrm{rad}_4^{-1},
\]
where the last inequality uses once more $\frac{1}{2}-4\varepsilon>0$ (K7).
Therefore
\[
2\le88\le\mathrm{rad}_4\le\mathrm{rad}_1,
\qquad
2\le88\le\mathrm{rad}_2\le\mathrm{rad}_1,
\]
and consequently, for $i=1,2,4$,
\begin{equation}\tag{24.2}
\frac{1}{\mathrm{rad}_i-1}\le\frac{2}{\mathrm{rad}_i}.
\end{equation}

\emph{The three norms of $\mathscr{Z}(s')$.}
Let $s\in\mathrm{Pd}_k$ and $\mathsf{Z}\in\mathrm{Box}_{k}$.
By the definition of $\mathrm{Box}_{k}$ (Section~I.5.7; \cite[Section~5.2.5]{QEDI}) and by (A1),
\[
|\mathsf{Z}(\Upsilon)|\le p_{0,k}\le e_k^{-\frac{\varepsilon}{4}}
\]
for every $\Upsilon$.
By W4 (Section~G4.0.4) and the bound for $c_{\gamma}$ in Section~0.6 with the sum over the basis elements $\Upsilon'$, it follows that, for every $\Upsilon$,
\[
\left|\left(C_{k}^{1/2}(s')\,\mathsf{Z}\right)(\Upsilon)\right|
\le\sum_{\Upsilon'}C_{w}\cdot\exp\left(-\gamma_{\mathrm{w}}\cdot d(\Upsilon,\Upsilon')\right)\cdot e_k^{-\frac{\varepsilon}{4}}
\le C_{w}\cdot c_{\gamma}\cdot e_k^{-\frac{\varepsilon}{4}};
\]
for complex $s'$ these coordinates are complex numbers.
The numbers $C(b',\Upsilon)$ are real.
Evaluating the bound of (P2) at the element of $V_{k}$ whose coordinates are the numbers $\mathrm{sign}\,C(b',\Upsilon)$, one obtains
\[
\sum_{\Upsilon}|C(b',\Upsilon)|\le c_{C}
\]
for every unit bond $b'$.
Hence $v:=C\cdot C_{k}^{1/2}(s')\,\mathsf{Z}$ satisfies
\[
|v(b')|\le c_{C}\cdot C_{w}\cdot c_{\gamma}\cdot e_k^{-\frac{\varepsilon}{4}}.
\]
For a positively oriented fine bond $b=(c,c+\eta_ke_{\mu})$ one has
\[
\mathscr{Z}(s')_{\mu}(c)=\sum_{b'}H_{k}(s')(b,b')\cdot v(b'),
\qquad
\left(\partial_{\nu}\mathscr{Z}(s')_{\mu}\right)(c)=\sum_{b'}\left(\partial_{\nu}H_{k}(s')\right)(b,b')\cdot v(b'),
\]
and, for two sites $c$, $\tilde{c}$ with $0<|c-\tilde{c}|\le1$ and $\tilde{b}:=(\tilde{c},\tilde{c}+\eta_ke_{\mu})$,
the H\"older quotient of $\partial_{\nu}\mathscr{Z}(s')_{\mu}$ at $c$, $\tilde{c}$ is
\[
\sum_{b'}\left(\delta_{\alpha}\partial_{\nu}H_{k}(s')\right)\left((b,\tilde{b});b'\right)\cdot v(b').
\]
By W3 and the bound for $c_{\gamma}$ in Section~0.6 with the sum over the positively oriented unit bonds $b'$, each of the three norms of $\mathscr{Z}(s')$ (Section~I.5.3; \cite[Section~5.3.1]{QEDI}) is therefore at most
\[
C_{w}\cdot c_{\gamma}\cdot c_{C}\cdot C_{w}\cdot c_{\gamma}\cdot e_k^{-\frac{\varepsilon}{4}}
=C_{w}^2\cdot c_{\gamma}^2\cdot c_{C}\cdot e_k^{-\frac{\varepsilon}{4}}
\le\frac{1}{2}e_k^{-\frac{\varepsilon}{4}}\cdot e_k^{-\frac{\varepsilon}{4}}
=\frac{1}{2}e_k^{-\frac{\varepsilon}{2}},
\]
the inequality being (A2).

\emph{Level $0$.}
Here $H_{0}(s')$ is the diagonal family of Section~I.5.8.
By W17 it has W3 with the constant $4\exp(2\gamma_{\mathrm{w}})\cdot M^{27\alpha_{0}}\le C_{w}$ (see (K6)), and the family (p2) has W4 (W17).
So the three sums over $b'$ are again at most $C_{w}\cdot c_{\gamma}\cdot\sup_{b'}|v(b')|$,
with $|v(b')|\le c_{C}\cdot C_{w}\cdot c_{\gamma}\cdot e_k^{-\frac{\varepsilon}{4}}$,
and the bound displayed above holds at $k=0$ as well.

\emph{The domains.}
Let $\mathrm{Zf}$ be a complex bond function whose three norms are at most $\frac{1}{2}e_k^{-\frac{\varepsilon}{2}}$, and let $|t|\le\mathrm{rad}_1$.
The three norms of $t\cdot\mathrm{Zf}$ are at most
\[
\frac{1}{2}e_k^{-\frac{3}{4}+\frac{9\varepsilon}{2}}\cdot\frac{1}{2}e_k^{-\frac{\varepsilon}{2}}
=\frac{1}{4}e_k^{-\frac{3}{4}+4\varepsilon}
<e_k^{\varepsilon}\cdot\mathrm{t3}_k,
\]
and $\mathrm{t3}_k\le\mathrm{t2}_k\le\mathrm{t1}_k$.
Let $A=A_{0}+A_{1}$ be a decomposition as in the definition of $\tilde{R}_k^{-}$ (Section~0.6):
for $i=1,2,3$ the $i$-th norm of $A_{1}$ is smaller than $\left(1-e_k^{\varepsilon}\right)$ times the $i$-th of the thresholds $\mathrm{t1}_k$, $\mathrm{t2}_k$, $\mathrm{t3}_k$.
The $i$-th norm of $t\cdot\mathrm{Zf}$ is smaller than $e_k^{\varepsilon}\cdot\mathrm{t3}_k$, hence smaller than $e_k^{\varepsilon}$ times the $i$-th threshold.
It follows that the $i$-th norm of $A_{1}+t\cdot\mathrm{Zf}$ is smaller than the $i$-th threshold.
Thus $A_{1}+t\cdot\mathrm{Zf}\in R_k$, and
\[
A+t\cdot\mathrm{Zf}=A_{0}+\left(A_{1}+t\cdot\mathrm{Zf}\right)\in\tilde{R}_k.
\]
Finally $\mathscr{T}+t\cdot\mathrm{Zf}=\left(A+t\cdot\mathrm{Zf}\right)+\mathscr{C}$ lies in $\tilde{R}_k+\mathrm{Sd}_k$,
which is contained in $\tilde{R}_k+\mathfrak{X}_k$ and in $\tilde{G}_k$ by (Sd2).
\end{proof}

\subsection*{3. Proof of Proposition~24.3: (q5), (al), (ga), (be), (de), (et), and (ze) for the finite sums and integrals}

\subsubsection*{3.1. (q5), first and second sentence}

\begin{proof}[Proof of Proposition~24.3\,(q5), first and second sentence]
The field $\mathscr{Z}(s')=H_{k}(s')\cdot C\cdot C_{k}^{1/2}(s')\,\mathsf{Z}$ depends linearly on $\mathsf{Z}$.
By W2 for (p1) and (p2) (at $k=0$: W17) the entries of $H_{k}(s')$ and of $C_{k}^{1/2}(s')$ are continuous on $\mathrm{Pd}_k$ and analytic in its interior;
$C$ is a fixed matrix.
For every fine bond $b$ the number $\mathscr{Z}(s')(b)$ is a finite sum of products of such entries and of coordinates $\mathsf{Z}(\Upsilon)$.
Hence $\mathscr{Z}(s')$ is continuous in $(s',\mathsf{Z})$ on $\mathrm{Pd}_k\times\mathrm{Box}_{k}$ and analytic in $s'$ in the interior.

Put
\[
\mathrm{Dq}:=\mathrm{Pd}_k\times\{|t|\le\mathrm{rad}_1\}\times\tilde{R}_k^{-}\times\mathrm{Sd}_k\times\mathrm{Box}_{k},
\qquad
\mathrm{Bm}(s',t,A,\mathscr{C},\mathsf{Z}):=A+\mathscr{C}+t\cdot\mathscr{Z}(s').
\]
The map $\mathrm{Bm}$ is continuous on $\mathrm{Dq}$;
it is a polynomial in $t$, $A$, $\mathscr{C}$ and in the entries just named, and by (q1) its values lie in $\tilde{G}_k$.
In the variables $(s,t,A,\mathscr{C})$ the interior of $\mathrm{Dq}$ is
\[
\left\{|s_{\mathrm{pc}}|<M^{\alpha_{0}}\right\}\times\{|t|<\mathrm{rad}_1\}\times\tilde{R}_k^{-}\times\mathrm{Sd}_k
\]
(DF0, (Sd6)).
By W2 (at $k=0$: W17) the entries of $H_{k}(s'',B)$ and of $\Gamma_k(s'',B)$ are continuous on $\mathrm{Pd}_k\times\tilde{G}_k$ and analytic in $(s'',B)$ in the interior,
and for $k\ge1$ the entries of $M_k(X,B)$ are analytic in $B$ on $\tilde{G}_k$.
(A system of level $0$ contains no family $M_0$.)
Their compositions with $\mathrm{Bm}$ are therefore continuous on $\mathrm{Dq}$ and analytic in $(s,t,A,\mathscr{C})$ in the interior.

The remaining entries of $H_{k}^{\#}(s'',B)$ are of two kinds.
First, $\bar{H}_{k}(s'',B)=\mathsf{C}^{-1}H_{k}(s'',-B)\,\mathsf{C}$, where $-B\in\tilde{G}_k$ by DF3.
Second, for a pair $(x,x')\in\mathrm{Pz}_k$ and a unit point $\mathrm{tu}'$,
\begin{equation}\tag{24.3}
\begin{aligned}
\left(\delta_{\alpha,B}H_{k}(s'',B)\right)\left((x,x'),\mathrm{tu}'\right)&=d_X(x,x')^{-\alpha}\cdot\left(U_B(x,x')\cdot H_{k}(s'',B)(x',\mathrm{tu}')-H_{k}(s'',B)(x,\mathrm{tu}')\right),\\
\left(\delta_{\alpha,-B}\bar{H}_{k}(s'',B)\right)\left((x,x'),\mathrm{tu}'\right)&=d_X(x,x')^{-\alpha}\cdot\left(U_B(x,x')^{-1}\cdot\bar{H}_{k}(s'',B)(x',\mathrm{tu}')-\bar{H}_{k}(s'',B)(x,\mathrm{tu}')\right),
\end{aligned}
\end{equation}
where $U_B(x,x')=\exp\left\{i\,e_k\eta_k\cdot(\tau B)(c,c')\right\}$ is an entire function of $B$ ((I.1.6), (I.1.7)).
So these entries have the same regularity.

\emph{Fixed $s\in\mathrm{Pd}_k$.}
The set of the $(t,A,\mathscr{C})$ (in the case (o): of the $(t,A)$, with $\mathscr{C}=0$) for which $A+\mathscr{C}+t\cdot\mathscr{Z}(s')\in\tilde{G}_k$ is open,
as the preimage of an open set (DF0) under a continuous map,
and by (q1) it contains $\{|t|\le\mathrm{rad}_1\}\times\tilde{R}_k^{-}\times\mathrm{Sd}_k$.
On this set the entries are compositions of a polynomial map with functions that are analytic in $B$ on $\tilde{G}_k$
(W2: for every fixed $s\in\mathrm{Pd}_k$, the boundary included).

\emph{Fixed $t$ with $|t|\le\mathrm{rad}_1$, the circle $|t|=\mathrm{rad}_1$ included, and fixed $\mathsf{Z}\in\mathrm{Box}_{k}$.}
The map
\[
(s',A,\mathscr{C})\longmapsto\mathrm{Bm}(s',t,A,\mathscr{C},\mathsf{Z})=A+\mathscr{C}+t\cdot\mathscr{Z}(s')
\]
is analytic on $\left\{|s'_{\mathrm{pc}}|<M^{\alpha_{0}}\right\}\times\tilde{R}_k^{-}\times\mathrm{Sd}_k$:
it is a polynomial in $A$, $\mathscr{C}$ and in the entries of $H_{k}(s')$ and $C_{k}^{1/2}(s')$, which are analytic in $s'$ there.
Its values lie in the open set $\tilde{G}_k$ ((q1), DF0).
The entries of $H_{k}(s'',B)$ and of $\Gamma_k(s'',B)$ are analytic in $(s'',B)$ on $\left\{|s''_{\mathrm{pc}}|<M^{\alpha_{0}}\right\}\times\tilde{G}_k$,
and for $k\ge1$ the entries of $M_k(X,B)$ are analytic in $B$ on $\tilde{G}_k$ (W2; at $k=0$: W17);
the remaining entries of $H_{k}^{\#}(s'',B)$ arise from these through $B\mapsto-B$ (DF3) and through (24.3).
Consequently every entry named in (q5), taken at $B=\mathrm{Bm}(s',t,A,\mathscr{C},\mathsf{Z})$,
is analytic in $(s,A,\mathscr{C})$ on $\left\{|s_{\mathrm{pc}}|<M^{\alpha_{0}}\right\}\times\tilde{R}_k^{-}\times\mathrm{Sd}_k$.
\end{proof}

\subsubsection*{3.2. (al)}

\begin{proof}[Proof of Proposition~24.3\,(al)]
\emph{First sentence.}
Let $Y'(X,P)=\mathbf{Y}$.
A lattice site of $X$, or of a cube of $P_0(X)\cup P$, lies in exactly one $M$-cube (g0), and this $M$-cube meets $X\cup\mathrm{Yp}(X,P)$;
so it is one of the $M$-cubes of $Y'(X,P)=\mathbf{Y}$.
Hence $X$ and all cubes of $P_0(X)\cup P$ are contained in $\mathbf{Y}$.

\emph{Second and third sentence.}
If a set $\{X_i\}$ or a family $(X_1,\dots,X_n)$ has the union $\mathbf{Y}$, each of its members is a subset of $\mathbf{Y}$.
Every $M$-cube of $X$ lies in an $LM$-cube (g0), and this $LM$-cube belongs to $\bar{X}^{L}$ by the definition of $\bar{X}^{L}$;
so $X$ is contained in $\bar{X}^{L}=LY$.

\emph{The datum.}
In (D4) the datum $(m,\mathscr{E})$ enters only through the number $m'=m+m(\mathscr{E})$
(where $m(\mathscr{E})$ is a complex number by (X1); at $k=0$ one has $m'=m$ and $\mathscr{E}^{r}=0$)
and through the elements $\mathscr{E}^{r}(X;\,\cdot\,)=(R\mathscr{E})(X;\,\cdot\,)$.
Indeed, $G^{(1)}$ does not contain the datum;
$G^{(2)}(s,X)$ and $G^{(4)}(s'',X)$ contain $\mathscr{E}^{r}(X;\,\cdot\,)$ with the same $X$;
$G^{(3)}$ and $G^{(5)}$ contain $m'$.
Let $(m_1,\mathscr{E}_1)$ and $(m_2,\mathscr{E}_2)$ be as in (al).
By the first sentence, $E^{(i),*}(\mathbf{Y})=\mathrm{Loc}[G^{(i)}](\mathbf{Y})$ is formed from the functions $G^{(i)}(\,\cdot\,,X)$ with $X\subseteq\mathbf{Y}$,
and these coincide for the two data.
Therefore, for every polymer $\mathbf{Y}'\subseteq\mathbf{Y}$, the two data give the same elements
\[
E^{(i),*}(\mathbf{Y}'),\qquad
E^{(i),x}(\mathbf{Y}')=\mathrm{sub}_{\mathbf{Y}',\mathscr{T}}\left(E^{(i),*}(\mathbf{Y}')\right),\qquad
E^{(\le3),x}(\mathbf{Y}')
\]
(at $k=0$: $E^{(1),x}(\mathbf{Y}')=E^{(1),*}(\mathbf{Y}')$).
Since $G^{(\le3)}(s'',X)$ is formed from $E^{(\le3),x}(X)$, the three functions of the second localisation coincide for $X\subseteq\mathbf{Y}$,
and so does $E^{(\le5),x}(\mathbf{Y}')$ for $\mathbf{Y}'\subseteq\mathbf{Y}$.
The elements $\mathrm{Km}(\mathbf{Y}')$ and $\mathrm{Hm}(\mathbf{Y}')$, $\mathbf{Y}'\subseteq\mathbf{Y}$, are formed by (8.4) from the $E^{(\le5),x}(X_i)$ with $X_i\subseteq\mathbf{Y}'$ (second sentence);
so they coincide as well.
The terms of the series (8.4) for $E^{\#,x}_k(m,\mathscr{E})(\mathbf{Y})$ are the numbers $\frac{1}{n!}\mathrm{urs}^{T}(X_1,\dots,X_n)$ times the products of the $\mathrm{Hm}(X_i)$, with $X_i\subseteq\mathbf{Y}$;
hence the two series agree term by term.
For $E^{\#,x}_k(m,\mathscr{E})(\mathbf{Y})$ itself the assertion follows from this once the series (8.4) are known to converge (Section~5.4); see Section~6.

\emph{The parenthetical sentence of (al).}
By Section~I.5.6, with $\mathbf{1}$ the unit $4\times4$ matrix and $\square$ an $M$-cube,
\[
m(\mathscr{E})\cdot\mathbf{1}=-\sum_{X\text{ small},\ X\supseteq\square}\mathrm{al}_{2}(\mathscr{E},X),
\]
and
\[
(R\mathscr{E})(X)-\mathscr{E}(X)=-\mathrm{al}_{0}(\mathscr{E},X)\cdot\operatorname{Vol}(X)-\int_X\mathrm{d}x\,\bar{\psi}(x)^{T}\,\mathrm{al}_{2}(\mathscr{E},X)\,\psi(x)
\]
for small $X$, while $(R\mathscr{E})(X)-\mathscr{E}(X)=0$ for large $X$;
here $\mathrm{al}_{0}(\mathscr{E},X)$ and $\mathrm{al}_{2}(\mathscr{E},X)$ are formed from $\mathscr{E}_{00}(X;0,0)$ and $\mathscr{E}_{20}(X;0,0;\,\cdot\,,\,\cdot\,)$.
(The pair $(0,0)$ lies in $\tilde{R}_k\times\mathrm{Sd}_k$ by (Sd1).)
\end{proof}

\subsubsection*{3.3. Rows of the primitives}

\begin{named}{Lemma 3.A}
Let $\mathrm{Ps}$ be a set of parameter cubes (of the kinds $\pi_k$, $\pi'_k$, $\pi''_k$),
let $\mathrm{Wz}$ be a union of $M$-cubes that contains every $M$-cube meeting a cube of $\mathrm{Ps}$,
let $\mathrm{Wz}_1$ be a component of $\mathrm{Wz}$,
and let $s=(s',s'')\in\mathrm{Pd}_k$ vanish outside $\mathrm{Ps}$.
Let $\mathrm{Lnks}$ be the set of the link-components of the three sets
\[
\{\mathrm{pc}\in\mathrm{Ps}:\ \mathrm{pc}\in\pi_k\},\qquad
\{\mathrm{pc}\in\mathrm{Ps}:\ \mathrm{pc}\in\pi'_k\},\qquad
\{\mathrm{pc}\in\mathrm{Ps}:\ \mathrm{pc}\in\pi''_k\}.
\]
For $\mathrm{Lnk}\in\mathrm{Lnks}$ the symbol $\mathrm{Lnk}$ also denotes the union of the cubes of $\mathrm{Lnk}$ (as in W8).
\begin{itemize}
\item[(r0)] For $\mathrm{Lnk}\in\mathrm{Lnks}$ the polymer $\mathrm{pol}(\mathrm{Lnk})$ is contained in $\mathrm{Wz}$.
Either it lies in $\mathrm{Wz}_1$, and then $\mathrm{Lnk}$ is called \emph{attached} and every cube of $\mathrm{Lnk}$ is contained in $\mathrm{Wz}_1$;
or $\mathrm{pol}(\mathrm{Lnk})$ and $\mathrm{Wz}_1$ have no $M$-cube in common.
If a photon site $\tilde{c}$ of $\mathrm{Wz}_1$ lies in a cube $\mathrm{pc}$ of $\mathrm{Lnk}$,
or lies in $\mathrm{pc}^{+}$ for a cube $\mathrm{pc}$ of $\mathrm{Lnk}$ of the kind $\pi'_k$ or $\pi''_k$,
then $\mathrm{Lnk}$ is attached.
\item[(r1)] Let $b$ be a fine bond that has an end in $\mathrm{Wz}_1$.
The row $H_{k}(s')(b,\,\cdot\,)$ vanishes if $b$ is not contained in $\mathrm{Wz}_1$,
and $H_{k}(s')(b,b')=0$ unless the unit bond $b'$ is contained in $\mathrm{Wz}_1$.
For a unit bond $b'$ contained in $\mathrm{Wz}_1$, $C(b',\Upsilon)=0$ unless $\Upsilon$ is contained in $\mathrm{Wz}_1$.
For an $\Upsilon$ with a point of its support in $\mathrm{Wz}_1$, the row $C_{k}^{1/2}(s')(\Upsilon,\,\cdot\,)$ vanishes if $\Upsilon$ is not contained in $\mathrm{Wz}_1$,
and $C_{k}^{1/2}(s')(\Upsilon,\Upsilon')=0$ unless $\Upsilon'$ is contained in $\mathrm{Wz}_1$.
Hence
\begin{equation}\tag{24.4}
\mathscr{Z}(s')(b)=\sum_{b',\,\Upsilon,\,\Upsilon'\subseteq\mathrm{Wz}_1}H_{k}(s')(b,b')\cdot C(b',\Upsilon)\cdot C_{k}^{1/2}(s')(\Upsilon,\Upsilon')\cdot\mathsf{Z}(\Upsilon'),
\end{equation}
where the summation extends over the unit bonds $b'$ and the basis elements $\Upsilon$, $\Upsilon'$ that are contained in $\mathrm{Wz}_1$.
The entries of $H_{k}(s')$ and of $C_{k}^{1/2}(s')$ in (24.4) do not change if $s'$ is replaced by a family $\tilde{s}'\in\mathrm{Pd}_k$
that vanishes outside $\mathrm{Ps}$ and agrees with $s'$ at the cubes contained in $\mathrm{Wz}_1$.
\item[(r2)] Let $B\in\tilde{G}_k$, and let $x=(c,g)$ be a fermion site of $T^{\mathrm{f}}_k$ such that $c$ has sup-distance at most $1$ from a photon site $\tilde{c}$ of $\mathrm{Wz}_1$.
The rows $H_{k}(s'',B)(x,\,\cdot\,)$ and $\bar{H}_{k}(s'',B)(x,\,\cdot\,)$ vanish unless all cubes of $\pi_k$ that meet $K_{c}$ belong to $\mathrm{Ps}$.
If they do, these cubes belong to one $\mathrm{Lnk}_3\in\mathrm{Lnks}$, the link-component $\mathrm{Lnk}_3$ is attached, $c$ lies in $\mathrm{Lnk}_3$ (so $c$ lies in $\mathrm{Wz}_1$),
and the columns of the two rows are contained in $\mathrm{Lnk}_3$.
The two rows do not change if $(s'',B)$ is replaced by $(\tilde{s}'',\tilde{B})$,
where $\tilde{s}''\in\mathrm{Pd}_k$ vanishes outside $\mathrm{Ps}$ and agrees with $s''$ at the cubes contained in $\mathrm{Wz}_1$,
and $\tilde{B}\in\tilde{G}_k$ agrees with $B$ on the bonds that meet $\mathrm{Wz}_1$.
\item[(r3)] Let $B\in\tilde{G}_k$, and let the unit point $\mathrm{tu}$ be contained in an attached $\mathrm{Lnk}_3\in\mathrm{Lnks}$ of the kind $\pi_k$.
The row $\Gamma_k(s'',B)(\mathrm{tu},\,\cdot\,)$ has its columns in $\mathrm{Lnk}_3$, and it does not change under the replacements of (r2).
\end{itemize}
\end{named}

\begin{proof}
(r0)
By (g2), $\mathrm{pol}(\mathrm{Lnk})$ is a polymer that contains the union of $\mathrm{Lnk}$.
Each of its $M$-cubes meets a cube of $\mathrm{Ps}$, and so it belongs to $\mathrm{Wz}$.
A connected union of $M$-cubes of $\mathrm{Wz}$ lies in one component of $\mathrm{Wz}$; this gives the alternative.
In the first case the lattice sites of the cubes of $\mathrm{Lnk}$ lie in $\mathrm{pol}(\mathrm{Lnk})$, hence in $\mathrm{Wz}_1$.
Now let $\tilde{c}$ be a photon site of $\mathrm{Wz}_1$; its $M$-cube $\tilde{\square}$ is an $M$-cube of $\mathrm{Wz}_1$ (g0).
If $\tilde{c}$ lies in a cube $\mathrm{pc}$ of $\mathrm{Lnk}$, then $\tilde{\square}$ meets $\mathrm{pc}$.
If $\tilde{c}$ lies in $\mathrm{pc}^{+}$, then $\tilde{\square}$ is an $M$-cube of $\mathrm{pc}^{+}$, and it meets $\mathrm{pc}$ by (g1).
In both cases $\tilde{\square}$ belongs to $\mathrm{pol}(\mathrm{Lnk})$ and to $\mathrm{Wz}_1$, so that $\mathrm{Lnk}$ is attached by the alternative.

(r1)
Apply W8 and W10 to (p1), with $\mathscr{Y}$ the set of the cubes of $\pi'_k$ in $\mathrm{Ps}$ (at $k=0$: W17).
By W10 the row at $b$ vanishes unless $b$ is contained in some $\mathrm{Lnk}_1\in\mathrm{Lnks}$ of the kind $\pi'_k$.
In that case the end of $b$ in $\mathrm{Wz}_1$ lies in a cube of $\mathrm{Lnk}_1$;
by (r0) the link-component $\mathrm{Lnk}_1$ is attached, and so $b$ is contained in $\mathrm{Wz}_1$.
By W8, $H_{k}(s')(b,b')=0$ unless $b'$ is contained in $\mathrm{Lnk}_1$, which lies in $\mathrm{Wz}_1$;
and the row at $b$ is the same for $\tilde{s}'$, since $\tilde{s}'$ vanishes outside $\mathscr{Y}$ and equals $s'$ on the cubes of $\mathrm{Lnk}_1$, which are contained in $\mathrm{Wz}_1$.
Next let $b'$ be contained in $\mathrm{Wz}_1$.
By (P2), $C(b',\Upsilon)=0$ unless $\operatorname{supp}\Upsilon$ lies in the union of the $L$-blocks of the ends of $b'$;
by (g0) each of these $L$-blocks lies in the $M$-cube of the corresponding end, which is an $M$-cube of $\mathrm{Wz}_1$.
Now apply W8 and W10 to (p2), with $\mathscr{Y}$ the set of the cubes of $\pi''_k$ in $\mathrm{Ps}$.
The row at $\Upsilon$ vanishes unless $\Upsilon$ is contained in some $\mathrm{Lnk}_2\in\mathrm{Lnks}$ of the kind $\pi''_k$.
The point of $\operatorname{supp}\Upsilon$ in $\mathrm{Wz}_1$ then lies in a cube of $\mathrm{Lnk}_2$;
so $\mathrm{Lnk}_2$ is attached, $\Upsilon$ is contained in $\mathrm{Wz}_1$, the columns $\Upsilon'$ are contained in $\mathrm{Lnk}_2$, and the row is the same for $\tilde{s}'$.
Finally, (24.4) is the definition
\[
\mathscr{Z}(s')(b)=\sum_{b',\,\Upsilon,\,\Upsilon'}H_{k}(s')(b,b')\cdot C(b',\Upsilon)\cdot C_{k}^{1/2}(s')(\Upsilon,\Upsilon')\cdot\mathsf{Z}(\Upsilon')
\]
with the vanishing terms left out.

(r2)
Apply W8 and W11 to (p3), with $\mathscr{Y}$ the set of the cubes of $\pi_k$ in $\mathrm{Ps}$ (at $k=0$: W17);
here $\operatorname{supp}x=\{c\}$ and $K_{x}=K_{c}$ (POS).
By W11 the row of $H_{k}(s'',B)$ at $x$ vanishes unless $\mathscr{Y}$ contains all cubes of $\pi_k$ that meet $K_{c}$.
By (g4) these cubes form a link-connected set; so they belong to one $\mathrm{Lnk}_3$.
The $M$-cube of $c$ meets $K_{c}$; so $c$ lies in $\mathrm{Lnk}_3$, that is, the index $x$ is contained in $\mathrm{Lnk}_3$.
Since $|\tilde{c}-c|\le1\le2L$, the site $\tilde{c}$ lies in $K_{c}$;
so the $M$-cube of $\tilde{c}$ meets $K_{c}$ and is a cube of $\mathrm{Lnk}_3$, and $\mathrm{Lnk}_3$ is attached by (r0).
By W8 the columns are contained in $\mathrm{Lnk}_3$, and the row is the same for $(\tilde{s}'',\tilde{B})$:
$\tilde{s}''$ equals $s''$ on the cubes of $\mathrm{Lnk}_3$, and a bond that meets $\mathrm{Lnk}_3$ meets $\mathrm{Wz}_1$.
As for $\bar{H}_{k}(s'',B)=\mathsf{C}^{-1}H_{k}(s'',-B)\,\mathsf{C}$:
the bond functions $-B$ and $-\tilde{B}$ lie in $\tilde{G}_k$ (DF3) and coincide where $B$ and $\tilde{B}$ do,
and the row of $\bar{H}_{k}(s'',B)$ at $x$ is $\mathsf{C}^{-1}H_{k}(s'',-B)(x,\,\cdot\,)\,\mathsf{C}$.

(r3)
This is W8 for $\Gamma_k$, with the same $\mathscr{Y}$ and with $Y_1=\mathrm{Lnk}_3$ (at $k=0$: W17).
\end{proof}

\subsubsection*{3.4. (ga): head, (ga1), (ga2), (ga3)}

\begin{proof}[Proof of Proposition~24.3\,(ga): head, (ga1), (ga2), (ga3)]
Let $\mathrm{Ps}(\mathbf{Y})$ be the set of the parameter cubes (of the three kinds) that are contained in $\mathbf{Y}$.
An $M$-cube that meets a cube $\mathrm{pc}$ of $\mathrm{Ps}(\mathbf{Y})$ has a lattice site in common with $\mathrm{pc}$;
this site lies in $\mathbf{Y}$, and so the $M$-cube is an $M$-cube of $\mathbf{Y}$ (g0).
Lemma~3.A therefore applies with $\mathrm{Ps}=\mathrm{Ps}(\mathbf{Y})$ and $\mathrm{Wz}=\mathrm{Wz}_1=\mathbf{Y}$ (the polymer $\mathbf{Y}$ is connected),
and to every $s\in\mathrm{Pd}_k$ that vanishes at the parameter cubes not contained in $\mathbf{Y}$.
All elements of $\mathrm{Lnks}$ are attached.
The cubes of $\pi_k$ in $\mathrm{Ps}(\mathbf{Y})$ are the $M$-cubes of $\mathbf{Y}$; they form one link-component, $\mathrm{Lnk}_3=\mathbf{Y}$.

\emph{Head.}
In a term $(X,P)$ of (8.2) attributed to $\mathbf{Y}$, the families $s$ that occur lie in $\mathrm{Slc}(X,P)$,
and the cubes of $P_0(X)\cup P$ are contained in $\mathbf{Y}$ (Section~3.2).

\emph{(ga1).}
The first three sentences of (ga1) are those of (r1).
The fourth sentence follows from (r1) for a fine bond or a basis element that has a point of its support in $\mathbf{Y}$.
If $b$ or $\Upsilon$ has no point of its support in $\mathbf{Y}$, it is contained in no element of $\mathrm{Lnks}$ (these lie in $\mathbf{Y}$), and the row vanishes by W10.
The display (24.1) is (24.4).

\emph{(ga2).}
If $c$ has sup-distance at most $1$ from $\mathbf{Y}$, then (r2) applies:
the rows at $x$ vanish unless all cubes of $\pi_k$ that meet $K_{c}$ are $M$-cubes of $\mathbf{Y}$,
and the $M$-cube of $c$ is one of these cubes, so that in particular $c$ lies in $\mathbf{Y}$.
If $c$ is farther from $\mathbf{Y}$, the $M$-cube of $c$ meets $K_{c}$ and is not a cube of $\mathbf{Y}$, and the rows vanish by W11.
The columns are contained in $\mathrm{Lnk}_3=\mathbf{Y}$.
At $k=0$, $H_{0}(s'',B)$ is the multiplication by the number
\[
\prod\left\{s''_{\mathrm{pc}}:\ \mathrm{pc}\in\pi_0,\ \mathrm{pc}\text{ meets }K_{c}\right\};
\]
by (I.ke1) its kernel in the sense of (I.1.4) is $48\cdot\mathbf{1}[x=x']$ times this number times $\mathbf{1}_4$,
and $\bar{H}_{0}(s'',B)=\mathsf{C}^{-1}H_{0}(s'',-B)\,\mathsf{C}$ is the same operator; so the column is $x$.
The assertion on $\Gamma_k$ is (r3) with $\mathrm{Lnk}_3=\mathbf{Y}$.
The assertion on $M_k(X,B)$ is W12.

\emph{(ga3).}
This is the last sentence of (r2) and of (r3), with $\tilde{s}''=s''$;
for the rows of (p3) at points $x$ with $c$ outside $\mathbf{Y}$ both sides vanish.
The assertion on $M_k$ is W12.
\end{proof}

\subsubsection*{3.5. The functions of the first localisation}

\emph{(Reg).}
A function $G(s,X)$ of (D4), with values in one of the finite dimensional spaces $\mathrm{Gr}[\psi,\chi,\Psi]$, $\mathrm{Gr}[\psi,\chi,W]$,
is said to have (Reg) if its coefficients are
\begin{itemize}
\item[(i)] continuous in $(s,A,\mathscr{C},\mathsf{Z})$ on $\mathrm{Pd}_k\times\tilde{R}_k^{-}\times\mathrm{Sd}_k\times\mathrm{Box}_{k}$;
\item[(ii)] analytic in $s$ in the interior of $\mathrm{Pd}_k$, for fixed $(A,\mathscr{C},\mathsf{Z})$;
\item[(iii)] analytic in $(A,\mathscr{C})$ on $\tilde{R}_k^{-}\times\mathrm{Sd}_k$, for fixed $s\in\mathrm{Pd}_k$ and fixed $\mathsf{Z}\in\mathrm{Box}_{k}$.
\end{itemize}
Properties (i) and (ii) together are the hypothesis of (D3).

\emph{(Int).}
Let $\mathrm{Omq}$ be a subset of a finite dimensional real or complex vector space,
let $\{|t-\mathrm{tce}|=\mathrm{rad}\}$ be a circle in the $t$-plane, with a complex number $\mathrm{tce}$ and $\mathrm{rad}>0$,
let $\mathrm{krq}$ be a continuous complex function on this circle,
and let $f(t,w)$ be continuous on $\{|t-\mathrm{tce}|=\mathrm{rad}\}\times\mathrm{Omq}$.
Then the function
\[
w\longmapsto\oint_{|t-\mathrm{tce}|=\mathrm{rad}}\mathrm{d}t\,f(t,w)\cdot\mathrm{krq}(t)
\]
is continuous on $\mathrm{Omq}$.
If $\mathrm{Omq}$ is an open subset of some $\mathbb{C}^n$ and $f(t,\,\cdot\,)$ is analytic for every $t$, then this function is analytic.
Indeed, $f\cdot\mathrm{krq}$ is uniformly continuous on the product of the circle with any compact subset of $\mathrm{Omq}$;
so the Riemann sums of the integral converge uniformly on the compact subsets of $\mathrm{Omq}$, and they are continuous, respectively analytic.
For $\mathrm{rad}>1$ the functional $\oint_{\mathrm{rad}}$ of Section~I.5.2 is an integral of this kind, with
\[
\mathrm{tce}=0,\qquad\mathrm{krq}(t)=\frac{1}{2\pi i\,t(t-1)};
\]
so are the Cauchy coefficients of Section~9.4\,(3), with $\mathrm{krq}(t)=\frac{(t-\mathrm{tce})^{-j-1}}{2\pi i}$.
The same conclusions hold for the function $w\mapsto\int\mathrm{d}\hat{\mu}(\mathsf{Z})\,f(\mathsf{Z},w)$, where $f$ is continuous on $\mathrm{Box}_{k}\times\mathrm{Omq}$.
Indeed, take partitions of the compact set $\mathrm{Box}_{k}$ into finitely many Borel sets $\mathrm{Bs}_j$ of diameter at most $\delta$, and points $\mathsf{Z}_j\in\mathrm{Bs}_j$;
then the sums
\[
\sum_j\hat{\mu}(\mathrm{Bs}_j)\cdot f(\mathsf{Z}_j,w)
\]
converge to the integral as $\delta\to0$, uniformly on the compact subsets of $\mathrm{Omq}$,
by uniform continuity and because $\hat{\mu}$ has mass one.

\begin{named}{Lemma 3.B}
Let $G$ be one of $G^{(1)}$, $G^{(2)}$, $G^{(3)}$ if $k\ge1$, and let $G=G^{(1)}$ if $k=0$.
Let $X\in D_k$, where $X=\square$ is an $M$-cube for $G^{(3)}$ and at $k=0$.
The kinds of the parameter cubes are $\pi_k$, $\pi'_k$, $\pi''_k$.
\begin{itemize}
\item[(reg)] $G(s,X)$ is defined for all $s\in\mathrm{Pd}_k$, $(A,\mathscr{C})\in\tilde{R}_k^{-}\times\mathrm{Sd}_k$ and $\mathsf{Z}\in\mathrm{Box}_{k}$;
it has (Reg), and it is even.
\end{itemize}
Let $P$ be a set of parameter cubes outside $P_0(X)$, put $Y':=Y'(X,P)$, and let $Y'_1$ be the component of $Y'$ that contains $X$.
\begin{itemize}
\item[(a)] For $s\in\mathrm{Slc}(X,P)$, $G(s,X)$ lies in the subalgebra generated by the $\psi(\xi)$ with $\xi$ over $X$,
the $\chi(\boldsymbol{\zeta})$ with $\boldsymbol{\zeta}$ admissible for $X$, and the $\Psi(u)$ with $u$ in $Y'_1$.
It takes the same value at two sets of arguments
\[
(A,\mathscr{C},\mathsf{Z},s),\ (\tilde{A},\tilde{\mathscr{C}},\tilde{\mathsf{Z}},\tilde{s})\ \in\ \tilde{R}_k^{-}\times\mathrm{Sd}_k\times\mathrm{Box}_{k}\times\mathrm{Slc}(X,P)
\]
that agree in the following entries:
$A$ and $\mathscr{C}$ on the bonds that meet $Y'_1$;
the $\mathsf{Z}(\Upsilon)$ with $\Upsilon$ contained in $Y'_1$;
the $s_{\mathrm{pc}}$ with $\mathrm{pc}$ contained in $Y'_1$.
\item[(b)] (The first assertion of (et).)
If $Y'$ is not connected, there is a cube $\mathrm{pc}\in P$ such that $G(s,X)=G(\tilde{s},X)$ for all $s,\tilde{s}\in\mathrm{Slc}(X,P)$ that differ at $\mathrm{pc}$ only;
and $\mathrm{Tm}[G](X,P)=0$.
\end{itemize}
\end{named}

\begin{proof}
Every $M$-cube of $X$ meets $X$; so $X$ lies in $Y'$, and since $X$ is connected, the component $Y'_1$ exists.
By its definition, $Y'$ contains every $M$-cube that meets a cube of $P_0(X)\cup P$.
Hence Lemma~3.A applies with $\mathrm{Ps}=P_0(X)\cup P$, $\mathrm{Wz}=Y'$, $\mathrm{Wz}_1=Y'_1$, and with $s\in\mathrm{Slc}(X,P)$.
\begin{itemize}
\item[(c0)] \emph{Connections.}
For $s\in\mathrm{Pd}_k$ and $|t|\le\mathrm{rad}_1$, (q1) gives $A+t\cdot\mathscr{Z}(s')\in\tilde{R}_k$ and $\mathscr{T}+t\cdot\mathscr{Z}(s')\in\tilde{G}_k$.
For $s\in\mathrm{Slc}(X,P)$ and a bond $b$ that meets $Y'_1$, the value $\mathscr{Z}(s')(b)$ is given by (24.4).
So the values of these two connections on the bonds that meet $Y'_1$ agree for two sets of arguments as in (a).
\item[(c1)] \emph{$G^{(1)}$, $k\ge1$.}
The integrand is
\[
\langle\bar{\Psi},M_k(X,B)\Psi\rangle_1=\sum_{y,y'}\sum_{\beta,\beta'}\bar{\Psi}_{\beta}(y)\cdot M_k(X,B)(y,y')_{\beta\beta'}\cdot\Psi_{\beta'}(y')
\]
at $B=\mathscr{T}+t\cdot\mathscr{Z}(s')$, $|t|=\mathrm{rad}_1$.
The circle $|t|=\mathrm{rad}_1$ is the boundary of the $t$-disk of $\mathrm{Dq}$.
By Section~3.1 the coefficients of the integrand are continuous on $\mathrm{Dq}$;
for every fixed $t$ on the circle and fixed $(A,\mathscr{C},\mathsf{Z})$ they are analytic in $s$ in the interior of $\mathrm{Pd}_k$ (Section~3.1, fixed $t$);
for every fixed $t$ on the circle, fixed $s\in\mathrm{Pd}_k$ and fixed $\mathsf{Z}$ they are analytic in $(A,\mathscr{C})$ on $\tilde{R}_k^{-}\times\mathrm{Sd}_k$ (Section~3.1, fixed $s$).
Now (Int) gives (Reg).
The integrand is bilinear, hence even.
By W12 the sum extends over $y,y'\in X$ only, and $M_k(X,B)=M_k(X,\tilde{B})$ whenever $B,\tilde{B}\in\tilde{G}_k$ agree on the bonds that meet $X$;
since $X$ lies in $Y'_1$, (c0) yields (a).
\item[(c2)] \emph{The fields $J_k^{\#}(s'',\mathscr{T},\mathscr{Z}(s'))\Psi$, $k\ge1$.}
By (8.1)$^{\mathrm{ch}}$ they are $H_{k}^{\#}(s'',B)\Psi$ at $B=\mathscr{T}+\mathscr{Z}(s')$ minus the same at $B=\mathscr{T}$;
both connections lie in $\tilde{G}_k$, by (q1) at $t=1\le\mathrm{rad}_1$ and at $t=0$.
Their entries are continuous on $\mathrm{Pd}_k\times\tilde{R}_k^{-}\times\mathrm{Sd}_k\times\mathrm{Box}_{k}$,
analytic in $s$ in the interior and, for fixed $s$, analytic in $(A,\mathscr{C})$ (Section~3.1).
Let $s\in\mathrm{Slc}(X,P)$, and let $x=(c,g)$ be a site over $X$ or one of the two points of a pair admissible for $X$.
By (adm), $c$ has sup-distance at most $\eta_k\le1$ from $X$, and $X$ lies in $Y'_1$.
By (r2) the rows of $H_{k}(s'',B)$ and of $\bar{H}_{k}(s'',B)$ at $x$ have their columns in $Y'_1$,
and, together with (c0), they agree for two sets of arguments as in (a).
By (24.3) the rows of $\delta_{\alpha,B}H_{k}(s'',B)$ and of $\delta_{\alpha,-B}\bar{H}_{k}(s'',B)$ at a pair $(x,x')$ admissible for $X$
are combinations of the rows at $x$ and at $x'$ with the phase $U_B(x,x')^{\pm1}$,
which reads $B$ on the bonds of the paths of the pair;
by (adm) these bonds meet $X$, and (c0) applies.
\item[(c3)] \emph{$G^{(2)}$, $k\ge1$.}
The family $\mathscr{E}^{r}=R\mathscr{E}$ lies in $K^{x,C}_{k}$ ((X2), (Sd5)).
For $|t|\le\mathrm{rad}_2\le\mathrm{rad}_1$ the pair $\left(A+t\cdot\mathscr{Z}(s'),\mathscr{C}\right)$ lies in $\tilde{R}_k\times\mathrm{Sd}_k$ (q1).
By (k1) and (N0) the coefficients of $\mathscr{E}^{r}(X;B_1,\mathscr{C})$ are analytic in $(B_1,\mathscr{C})$ on $\tilde{R}_k\times\mathrm{Sd}_k$.
The map $(t,s',A,\mathsf{Z})\mapsto A+t\cdot\mathscr{Z}(s')$ is continuous;
for fixed $(t,A,\mathsf{Z})$ it is analytic in $s'$ in the interior of $\mathrm{Pd}_k$ (Section~3.1), and for fixed $(t,s',\mathsf{Z})$ it is affine in $A$.
Consequently the coefficients of $\mathscr{E}^{r}\left(X;A+t\cdot\mathscr{Z}(s'),\mathscr{C}\right)$ are continuous in $(t,s,A,\mathscr{C},\mathsf{Z})$ on
\[
\{|t|=\mathrm{rad}_2\}\times\mathrm{Pd}_k\times\tilde{R}_k^{-}\times\mathrm{Sd}_k\times\mathrm{Box}_{k},
\]
analytic in $s$ in the interior of $\mathrm{Pd}_k$ for fixed $(t,A,\mathscr{C},\mathsf{Z})$, and analytic in $(A,\mathscr{C})$ for fixed $(t,s,\mathsf{Z})$.
The integrand is
\[
\mathrm{sb}_{tv}\left(\mathscr{E}^{r}\left(X;A+t\cdot\mathscr{Z}(s'),\mathscr{C}\right)\right),\qquad v=J_k^{\#}(s'',\mathscr{T},\mathscr{Z}(s'))\Psi;
\]
its coefficients are polynomials in those of $\mathscr{E}^{r}\left(X;A+t\cdot\mathscr{Z}(s'),\mathscr{C}\right)$, in $t$ and in the entries of (c2).
With (Int) this gives (Reg).
The element $\mathscr{E}^{r}(X;\,\cdot\,)$ is even, $\mathrm{sb}_{tv}$ maps generators to linear combinations of generators,
and $\oint_{\mathrm{rad}}$ commutes with the linear map $\mathrm{Pty}$; so $G^{(2)}(s,X)$ is even.
By (k2), $\mathscr{E}^{r}(X;B_1,\mathscr{C})$ contains only the $\psi(\xi)$ with $\xi$ over $X$ and the $\chi(\boldsymbol{\zeta})$ with $\boldsymbol{\zeta}$ admissible for $X$,
and it takes the same value at two pairs of $\tilde{R}_k\times\mathrm{Sd}_k$ that agree on the bonds meeting $X$.
The homomorphism $\mathrm{sb}_{tv}$ acts on it through the images of these generators.
Together with (c0) and (c2) this gives (a).
\item[(c4)] \emph{$G^{(3)}$, $k\ge1$, $X=\square$.}
The form
\[
\langle\bar{\psi},\psi\rangle_{\square}=\int_{\square}\mathrm{d}x\,\bar{\psi}(x)\cdot\psi(x)=\frac{\eta_k^3}{48}\sum_{c\in\square}\sum_{g}\sum_{\beta}\bar{\psi}_{\beta}(c,g)\,\psi_{\beta}(c,g)
\]
contains only the $\psi(\xi)$ with $\xi$ over $\square$, and no connection.
The integrand is $-m'\cdot\mathrm{sb}_{tv}\left(\langle\bar{\psi},\psi\rangle_{\square}\right)$, $v=J_k^{\#}(s'',\mathscr{T},\mathscr{Z}(s'))\Psi$,
with the complex number $m'=m+m(\mathscr{E})$ (X1).
Its coefficients are polynomials in $m'$, $t$ and the entries of (c2); with (Int) this gives (Reg).
It is the image of a bilinear element under a homomorphism that maps generators to linear combinations of generators; hence it is even.
The homomorphism $\mathrm{sb}_{tv}$ acts through the images $\psi(\xi)+t\cdot v(\xi)$ of the generators $\psi(\xi)$, $\xi$ over $\square$;
by (c2) these images lie in the subalgebra named in (a), and they agree for two sets of arguments as in (a).
\item[(c5)] \emph{$G^{(1)}$, $k=0$, $X=\square$.}
Here $G^{(1)}(s,\square)=\mathrm{Vx}(\mathscr{Z}(s'),\square)$.
By Proposition~23.2\,(e) it is continuous in $(s',\mathsf{Z})$ on $\mathrm{Pd}_0\times\mathrm{Box}_{0}$, analytic in $s'$ in the interior, and independent of $s''$;
by Proposition~23.2\,(c) it contains neither $A$ nor $\mathscr{C}$.
This gives (Reg).
By Proposition~23.2\,(c) every monomial has one barred and one unbarred generator, so that the element is even;
it contains no generator other than $\psi$, $\bar{\psi}$ at the sites over $\square$ and $\chi$, $\bar{\chi}$ at the pairs $(x,x^{(3)})$ with $c\in\square$, which are admissible for $\square$;
and it reads $\mathrm{Zf}=\mathscr{Z}(s')$ on the bonds that meet $\square$, which meet $Y'_1$; now (c0) applies.
\item[(b)] Let $Y'$ be not connected.
Then $Y'$ contains an $M$-cube $\square'$ outside $Y'_1$, and $\square'$ meets $X\cup\mathrm{Yp}(X,P)$.
It has no site in common with $Y'_1$ (g0), hence none with $X$ and none with a cube of $P_0(X)$, the cubes of $P_0(X)$ being contained in $X$.
So $\square'$ meets a cube $\mathrm{pc}$ of $P$, and $\mathrm{pc}$ is not contained in $Y'_1$.
By (a), $G(s,X)=G(\tilde{s},X)$ whenever $s,\tilde{s}\in\mathrm{Slc}(X,P)$ differ at $\mathrm{pc}$ only.
In (V) the subsets of $P$ form pairs $\{\mathrm{Psb},\mathrm{Psb}\cup\{\mathrm{pc}\}\}$ with $\mathrm{pc}\notin\mathrm{Psb}$;
the two families $\mathbf{1}[P_0(X)\cup\mathrm{Psb}]$ and $\mathbf{1}[P_0(X)\cup\mathrm{Psb}\cup\{\mathrm{pc}\}]$ differ at $\mathrm{pc}$ only,
and the two corresponding terms carry opposite signs.
Hence $\mathrm{Tm}[G](X,P)=0$.
\end{itemize}
\end{proof}

\subsubsection*{3.6. $E^{(i),*}(\mathbf{Y})$, $E^{(i),x}(\mathbf{Y})$, $E^{(\le3),x}(\mathbf{Y})$}

By Lemma~3.B\,(reg), every function $G^{(i)}(\,\cdot\,,X)$ satisfies the hypothesis of (D3) of Definition~24.D.
Consequently $E^{(i),*}(\mathbf{Y})=\mathrm{Loc}\left[G^{(i)}\right](\mathbf{Y})$ is defined for $i\le3$; at $k=0$ only $i=1$ occurs.
Let $(X,P)$ be a term attributed to $\mathbf{Y}$.
Then $Y'_1=Y'=\mathbf{Y}$, and $X$ lies in $\mathbf{Y}$ (Section~3.2).
By (V), $\mathrm{Tm}\left[G^{(i)}\right](X,P)$ is a finite linear combination, with coefficients $\pm1$, of the elements
\[
G^{(i)}\left(\mathbf{1}\left[P_0(X)\cup\mathrm{Psb}\right],X\right),\qquad\mathrm{Psb}\subseteq P.
\]
Each of these elements is a map that is local in $\mathbf{Y}$ and contains no $W$.
This follows from Lemma~3.B\,(a), applied with $\tilde{s}=s$, and from (Lo1): a site over $X$ lies over $\mathbf{Y}$, and by (adm) a pair admissible for $X$ is admissible for $\mathbf{Y}$.
Lemma~3.B\,(reg) shows moreover that each of them is even, continuous in $(A,\mathscr{C},\mathsf{Z})$ and analytic in $(A,\mathscr{C})$.
Hence
\begin{equation}\tag{24.5}
\parbox{0.8\linewidth}{$E^{(i),*}(\mathbf{Y})$ is local in $\mathbf{Y}$, contains no $W$, is even, is continuous in $(A,\mathscr{C},\mathsf{Z})$ on $\tilde{R}_k^{-}\times\mathrm{Sd}_k\times\mathrm{Box}_k$ and is analytic in $(A,\mathscr{C})$.}
\end{equation}
In particular $E^{(i),*}(\mathbf{Y})$ is an element of the subalgebra $\mathrm{Gr}[\psi,\chi,\Psi|\mathbf{Y}]$ of Proposition~23.1.
This is the first sentence of (ze).

\emph{The case $k\ge1$.}
Here $E^{(i),x}(\mathbf{Y})=\mathrm{sub}_{\mathbf{Y},\mathscr{T}}\left(E^{(i),*}(\mathbf{Y})\right)$ is defined, and by Proposition~23.1\,(d) it belongs to $\mathrm{Gr}[\psi,\chi]$.
Apart from the generators $\psi$, $\chi$ that $E^{(i),*}(\mathbf{Y})$ contains, it contains only generators $\psi(\xi)$ with $\xi$ over $\mathbf{Y}$ and $\chi(\boldsymbol{\zeta})$ at pairs of kind (II) for $\mathbf{Y}$.
Its coefficients are polynomials in the coefficients of $E^{(i),*}(\mathbf{Y})$ and in entries of the rows of $T^{(j)}(\mathscr{T},\mathbf{Y})$ and $\bar{T}^{(j)}(\mathscr{T},\mathbf{Y})$ at the unit sites of $\mathbf{Y}$; these entries are entire functions of $\mathscr{T}$, and they read $\mathscr{T}$ on bonds that meet $\mathbf{Y}$ (Proposition~23.1\,(b), (d)).
The homomorphism $\mathrm{sub}_{\mathbf{Y},\mathscr{T}}$ maps even elements to even elements.
\emph{The case $k=0$.}
Here $E^{(1),x}=E^{(1),*}$.
For every $k\ge0$ it follows that
\begin{equation}\tag{24.6}
\parbox{0.8\linewidth}{$E^{(i),x}(\mathbf{Y})$ ($i\le3$) and $E^{(\le3),x}(\mathbf{Y})$ are local in $\mathbf{Y}$, contain neither $\Psi$ nor $W$, are even, are continuous in $(A,\mathscr{C},\mathsf{Z})$ and are analytic in $(A,\mathscr{C})$.}
\end{equation}

\subsubsection*{3.7. The functions of the second localisation}

\begin{named}{Lemma 3.C}
Let $G$ be one of the functions $G^{(\le3)}$, $G^{(4)}$, $G^{(5)}$, where $G^{(4)}=0$ at $k=0$.
Let $X\in D_k$, with $X=\square$ an $M$-cube in the case of $G^{(5)}$; the kind of the parameter cubes is $\pi_k$.
Then the assertions (reg), (a) and (b) of Lemma~3.B hold with $s''$ in place of $s$ and with $W(u)$ in place of $\Psi(u)$.
\end{named}

\begin{proof}
In the second localisation $P_0(X)$ is the set of the $M$-cubes of $X$, and $P$ is a set of $M$-cubes not contained in $X$.
An $M$-cube meets another $M$-cube only if the two are equal, by (g0) of Section~0.6 (\cite[Section~17.4.1]{QEDI}); therefore $Y'(X,P)$ is the union of $X$ and of the cubes of $P$.
Lemma~3.A applies as in Section~3.5.

(c6) \emph{The fields $W_k^{\#}(s'',\mathscr{T})$.}
By Proposition~24.3\,(q1) at $t=0$, the connection $\mathscr{T}$ lies in $\tilde{G}_k$.
According to (8.1)$^{\mathrm{ch}}$, the unbarred components of $W_k^{\#}(s'',\mathscr{T})$ are $H_k(s'',\mathscr{T})\cdot\Gamma_k(s'',\mathscr{T})\,W$ and its image under $\delta_{\alpha,\mathscr{T}}$, and the barred components are $\bar{H}_k(s'',\mathscr{T})\,\bar{W}$ and its image under $\delta_{\alpha,-\mathscr{T}}$.
Their coefficients are finite sums of products of entries of $H_k^{\#}(s'',\mathscr{T})$ and of $\Gamma_k(s'',\mathscr{T})$, at $k=0$ with the weight $\frac{1}{48}$.
By Section~3.1, taken at $t=0$, these coefficients are continuous on $\mathrm{Pd}_k\times\tilde{R}_k^{-}\times\mathrm{Sd}_k$, analytic in $s''$ in the interior, and, for fixed $s''$, analytic in $(A,\mathscr{C})$.
Let $s''\in\mathrm{Slc}(X,P)$, and let $x$ be a site over $X$ or a point of a pair admissible for $X$.
By (adm) and (r2), the rows of $H_k(s'',\mathscr{T})$ and of $\bar{H}_k(s'',\mathscr{T})$ at $x$ either vanish or have their columns in an attached $\mathrm{Lnk}_3$.
By (r3), the rows of $\Gamma_k(s'',\mathscr{T})$ at these columns have their columns in $\mathrm{Lnk}_3$, which lies in $Y'_1$.
All these rows take the same values at two sets of arguments as in (a).
The phases $U_{\mathscr{T}}(x,x')^{\pm1}$ at a pair admissible for $X$ read $\mathscr{T}$ on bonds that meet $X$ (adm).

(c7) \emph{The three functions.}
Put $v:=W_k^{\#}(s'',\mathscr{T})$.
Then
\begin{align*}
G^{(\le3)}(s'',X)&=\mathrm{sb}_v\left(E^{(\le3),x}(X;A,\mathscr{C},\mathsf{Z})\right),\\
G^{(4)}(s'',X)&=\oint_{\mathrm{rad}_4}\mathrm{sb}_{tv}\left(\mathscr{E}^{r}(X;A,\mathscr{C})\right),\\
G^{(5)}(s'',\square)&=\oint_{\mathrm{rad}_4}(-m')\cdot\mathrm{sb}_{tv}\left(\langle\bar{\psi},\psi\rangle_{\square}\right),
\end{align*}
where $(A,\mathscr{C})$ lies in $\tilde{R}_k\times\mathrm{Sd}_k$ and $m'$ is a complex number ((X1); at $k=0$, $m'=m$).
By (24.6) with $X$ in place of $\mathbf{Y}$, by (k1), (k2), and by step (c4) of Section~3.5, the three elements to which $\mathrm{sb}$ is applied are even elements of $\mathrm{Gr}[\psi,\chi]$ that are local in $X$, continuous in $(A,\mathscr{C},\mathsf{Z})$ and analytic in $(A,\mathscr{C})$.
The coefficients of their images are polynomials in their coefficients, in $t$ and in the coefficients of the fields of (c6).
Combined with (Int), this gives (reg); combined with (c6), it gives (a).
The paragraph (b) of the proof in Section~3.5 uses, of the function $G$, only the assertion (a), the definition of $Y'(X,P)$ and (V).
It therefore proves (b) for the three functions of the second localisation.
\end{proof}

\subsubsection*{3.8. $E^{(\le5),x}(\mathbf{Y})$, $\mathrm{Km}(\mathbf{Y})$, $\mathrm{Hm}(\mathbf{Y})$}

By Lemma~3.C\,(reg), the element $E^{(\le5),x}(\mathbf{Y})=\mathrm{Loc}\left[G^{(\le3)}+G^{(4)}+G^{(5)}\right](\mathbf{Y})$ is defined.
Let $(X,P)$ be attributed to $\mathbf{Y}$; then $Y'_1=Y'=\mathbf{Y}$, and $X$ lies in $\mathbf{Y}$.
By (V), $\mathrm{Tm}[\,\cdot\,](X,P)$ is a linear combination, with coefficients $\pm1$, of the values of $G^{(\le3)}+G^{(4)}+G^{(5)}$ at the families $\mathbf{1}\left[P_0(X)\cup\mathrm{Psb}\right]$, $\mathrm{Psb}\subseteq P$.
By Lemma~3.C\,(a), applied with $\tilde{s}''=s''$, and by (Lo1), each of these values is a map that is local in $\mathbf{Y}$ and contains no $\Psi$; by Lemma~3.C\,(reg), it is even, continuous in $(A,\mathscr{C},\mathsf{Z})$ and analytic in $(A,\mathscr{C})$.
Therefore
\begin{equation}\tag{24.7}
\parbox{0.8\linewidth}{$E^{(\le5),x}(\mathbf{Y})$ is local in $\mathbf{Y}$, contains no $\Psi$, is even, is continuous in $(A,\mathscr{C},\mathsf{Z})$ and is analytic in $(A,\mathscr{C})$.}
\end{equation}

Let $F$ be an even element, and let $\mathrm{cf}$ be the coefficient of the monomial $1$ in $F$.
The element $F-\mathrm{cf}$ is a linear combination of monomials of degree at least $2$; it is therefore nilpotent, and it commutes with the number $\mathrm{cf}$.
Hence
\[
\exp F=\exp(\mathrm{cf})\cdot\sum_{j}\frac{(F-\mathrm{cf})^{j}}{j!},
\]
the sum being finite.
The coefficients of $\exp F$ are thus $\exp(\mathrm{cf})$ times polynomials in the coefficients of $F$, that is, entire functions of the coefficients of $F$.
As $\mathrm{Pty}$ is an automorphism, one has $\mathrm{Pty}\,\exp F=\exp\left(\mathrm{Pty}\,F\right)=\exp F$, so that $\exp F$ is even.
Finite sums and finite products of even elements are even, and their coefficients are polynomials in the coefficients of the summands and of the factors; continuity and analyticity are therefore preserved.
Even elements commute, so that the products in $\mathrm{Km}(\mathbf{Y})$ are defined; the sum is finite, $D_k$ being a finite set.
The polymers $X_i$ of a set whose union is $\mathbf{Y}$ lie in $\mathbf{Y}$.
From (24.7), these facts and (Lo1)--(Lo3) one obtains
\begin{equation}\tag{24.8}
\parbox{0.8\linewidth}{$\mathrm{Km}(\mathbf{Y})$ is local in $\mathbf{Y}$, contains no $\Psi$, is even, is continuous in $(A,\mathscr{C},\mathsf{Z})$ and is analytic in $(A,\mathscr{C})$.}
\end{equation}

The set $\mathrm{Box}_k$ is compact and $\hat{\mu}$ is a probability measure on it; hence
\[
\mathrm{Hm}(\mathbf{Y})=\int\mathrm{d}\mu_I(W)\,\mathrm{d}\hat{\mu}(\mathsf{Z})\,\mathrm{Km}(\mathbf{Y})
\]
is defined.
The functional $\int\mathrm{d}\mu_I(W)$ is the linear map that sends $c_J\cdot W(J)$, $c_J\in\mathrm{Gr}[\psi,\chi]$, to $c_J\cdot m(J)$; this is the identity for $\int\mathrm{d}\mu_I(W)\,F$ in (N3)$^{\mathrm{w}}$, that is, in Proposition~15\,(f).
By (mo), that is, Proposition~15\,(e), one has $m(J)=0$ unless $J$ is paired, and for paired $J$ the number $\#|J|$ is even.
This map therefore sends $\mathrm{Gl}_k(\mathbf{Y})\cap\mathrm{Gr}[\psi,\chi,W]$ into $\mathrm{Gl}_k(\mathbf{Y})\cap\mathrm{Gr}[\psi,\chi]$, and even elements to even elements.
The integral over $\mathsf{Z}$ is a limit of linear combinations of values of the integrand.
If $(A,\mathscr{C})$ and $(A^1,\mathscr{C}^1)$ coincide on the bonds that meet $\mathbf{Y}$, then by (24.8) the integrands coincide at every $\mathsf{Z}$.
Applying (Int) to the integral over $\mathrm{Box}_k$, one obtains
\begin{equation}\tag{24.9}
\parbox{0.8\linewidth}{$\mathrm{Hm}(\mathbf{Y})$ is an even element of $\mathrm{Gr}[\psi,\chi]$ that contains no generators other than $\psi(\xi)$ with $\xi$ over $\mathbf{Y}$ and $\chi(\boldsymbol{\zeta})$ with $\boldsymbol{\zeta}$ admissible for $\mathbf{Y}$; it takes the same value at two pairs of $\tilde{R}_k^{-}\times\mathrm{Sd}_k$ that coincide on the bonds that meet $\mathbf{Y}$; it is analytic in $(A,\mathscr{C})$.}
\end{equation}
The relations (24.5)--(24.9) contain the sentences of (ze) on the generators, on evenness, on the dependence on $\mathsf{Z}$ and on $(A,\mathscr{C})$, and on continuity and analyticity, for the objects $E^{(i),*}(\mathbf{Y})$, $E^{(i),x}(\mathbf{Y})$, $E^{(\le3),x}(\mathbf{Y})$, $E^{(\le5),x}(\mathbf{Y})$, $\mathrm{Km}(\mathbf{Y})$ and $\mathrm{Hm}(\mathbf{Y})$.

\emph{The last sentence of (ze).}
By (g0) the $M$-cubes partition the photon sites and the unit sites.
A fermion site over two polymers, a unit site that lies in both, or a point of the support (which is not empty) of a basis element contained in both, lies in an $M$-cube that belongs to both polymers.

\emph{The last sentence of (q5)} is contained in Lemma~3.B\,(reg) and Lemma~3.C\,(reg).
\qed

\subsubsection*{3.9. (ga4)}

\begin{proof}[Proof of Proposition~24.3\,(ga4)]
Let $(X,P)$ be a term of (8.2), of the first or of the second localisation, attributed to $\mathbf{Y}$; let $s\in\mathrm{Slc}(X,P)$, and let $B$ be one of the connections $\mathscr{T}+\mathscr{Z}(s')$ and $\mathscr{T}$.
The family $s$ vanishes at the parameter cubes that are not contained in $\mathbf{Y}$, since the cubes of $P_0(X)\cup P$ are contained in $\mathbf{Y}$ (Section~3.4, paragraph \emph{Head}).
After the omission of the terms that vanish by (ga1) and (ga2), the fields of (8.1)$^{\mathrm{ch}}$ are as follows; $x=(c,g)$ denotes a fermion site.
\begin{itemize}
\item For a bond $b$ that meets $\mathbf{Y}$, the value $\mathscr{Z}(s')(b)$ is given by the display (24.1).
\item Let $k\ge1$.
If $c$ lies in $\mathbf{Y}$, then
\begin{align*}
\left(H_k(s'',B)\,\Psi\right)(x)&=\sum_{y\in\mathbf{Y}}H_k(s'',B)(x,y)\cdot\Psi(y),\\
\left(\bar{H}_k(s'',B)\,\bar{\Psi}\right)(x)&=\sum_{y\in\mathbf{Y}}\bar{H}_k(s'',B)(x,y)\cdot\bar{\Psi}(y),\\
\left(H_k(s'',\mathscr{T})\cdot\Gamma_k(s'',\mathscr{T})\,W\right)(x)&=\sum_{y,y'\in\mathbf{Y}}H_k(s'',\mathscr{T})(x,y)\cdot\Gamma_k(s'',\mathscr{T})(y,y')\cdot W(y'),\\
\left(\bar{H}_k(s'',\mathscr{T})\,\bar{W}\right)(x)&=\sum_{y\in\mathbf{Y}}\bar{H}_k(s'',\mathscr{T})(x,y)\cdot\bar{W}(y);
\end{align*}
if $c$ does not lie in $\mathbf{Y}$, these four fields vanish at $x$.
\item Let $k=0$, and put
\[
\mathrm{hq}(c):=\prod_{\mathrm{pc}\in\pi_0:\ \mathrm{pc}\text{ meets }K_c}s''_{\mathrm{pc}};
\]
this number vanishes unless $c$ lies in $\mathbf{Y}$.
Then
\begin{align*}
\left(H_0(s'',\mathscr{T})\cdot\Gamma_0(s'',\mathscr{T})\,W\right)(x)&=\mathrm{hq}(c)\cdot\int_{\mathbf{Y}}\mathrm{d}x'\,\Gamma_0(s'',\mathscr{T})(x,x')\cdot W(x')\\
&=\mathrm{hq}(c)\cdot\frac{1}{48}\sum_{x'\text{ over }\mathbf{Y}}\Gamma_0(s'',\mathscr{T})(x,x')\cdot W(x'),\\
\left(\bar{H}_0(s'',\mathscr{T})\,\bar{W}\right)(x)&=\mathrm{hq}(c)\cdot\bar{W}(x),
\end{align*}
since the entry $H_0(s'',B)(x,x')$ of the kernel (I.1.4) equals $48\cdot\mathbf{1}[x=x']\cdot\mathrm{hq}(c)\cdot\mathbf{1}_4$.
\item The components with a H\"older quotient are, at a pair $(x,x')$,
\[
d_X(x,x')^{-\alpha}\cdot\left(U_B(x,x')^{\pm1}\cdot f(x')-f(x)\right),
\]
where $f$ is one of the fields just displayed; compare (24.3).
\end{itemize}
In (D4) these fields occur at the sites over $X$ and at the pairs admissible for $X$ (Section~3.5, steps (c3), (c4); Section~3.7, step (c7)), and $X$ is contained in $\mathbf{Y}$.
For $k\ge1$, $G^{(1)}(s,X)$ contains the entries $M_k\left(X,\mathscr{T}+t\cdot\mathscr{Z}(s')\right)(y,y')$ with $y,y'\in X$ and $|t|=\mathrm{rad}_1$ (Section~3.5, step (c1)).
The first slot of $\mathscr{E}^{r}(X;\,\cdot\,,\mathscr{C},\,\cdot\,)$ is $A+t\cdot\mathscr{Z}(s')$ with $|t|=\mathrm{rad}_2$ in $G^{(2)}$, and $A$ in $G^{(4)}$.
No primitive occurs in the substitution, in the forms, in $\mathrm{Loc}$ or in (8.4).
The objects attached to $\mathbf{Y}$ are formed from terms attributed to polymers contained in $\mathbf{Y}$ ((al)).
Hence the displays contain the primitives through the entries listed in (ga4), and through no other entries.

\emph{The connections.}
By (24.1), on a bond $b$ that meets $\mathbf{Y}$ the connections $\mathscr{T}+t\cdot\mathscr{Z}(s')$ and $A+t\cdot\mathscr{Z}(s')$ take the values
\[
A(b)+\mathscr{C}(b)+t\cdot\mathscr{Z}(s')(b)\qquad\text{and}\qquad A(b)+t\cdot\mathscr{Z}(s')(b),
\]
respectively, where $\mathscr{Z}(s')(b)$ is formed from the first two kinds of entries listed in (ga4), $H_k(s')(b,b')$ and $C_{k}^{1/2}(s')(\Upsilon,\Upsilon')$, from the matrix $C$ and from the numbers $\mathsf{Z}(\Upsilon')$ with $\Upsilon'$ contained in $\mathbf{Y}$.
By Proposition~24.3\,(q1) these connections lie in $\tilde{G}_k$ and in $\tilde{R}_k$, respectively.

\emph{Two systems.}
Let $\mathrm{Sy}^1$ and $\mathrm{Sy}^2$ be two systems of primitives on $T_k$ with the properties of (ST1)$^{\mathrm{ch}}$.
An upper index $j$ marks the objects formed with $\mathrm{Sy}^j$; in particular $C_{k}^{1/2,j}(s')$ denotes the family $C_{k}^{1/2}(s')$ of $\mathrm{Sy}^j$.
In (ga4) the two systems are assumed to coincide in the entries listed there, each entry being taken at the connections $\mathscr{T}+t\cdot\mathscr{Z}(s')$ and $\mathscr{T}$ formed with the respective system, for all $s$, $(A,\mathscr{C})$ and $\mathsf{Z}$ under consideration.
Written out, this hypothesis says: for all $s\in\mathrm{Pd}_k$ that vanish at the parameter cubes not contained in $\mathbf{Y}$, all $(A,\mathscr{C})\in\tilde{R}_k^{-}\times\mathrm{Sd}_k$, all $\mathsf{Z}\in\mathrm{Box}_k$ and all indices named in (ga4),
\begin{equation}\tag{e1}
H_k^1(s')(b,b')=H_k^2(s')(b,b'),\qquad C_{k}^{1/2,1}(s')(\Upsilon,\Upsilon')=C_{k}^{1/2,2}(s')(\Upsilon,\Upsilon');
\end{equation}
\begin{equation}\tag{e2}
\begin{aligned}
&H_k^1(s'',B^1)(x,u)=H_k^2(s'',B^2)(x,u),\qquad\bar{H}_k^1(s'',B^1)(x,u)=\bar{H}_k^2(s'',B^2)(x,u)\\
&\qquad\text{for }B^j=\mathscr{T}+\mathscr{Z}^j(s')\text{ and for }B^j=\mathscr{T};\\
&\Gamma_k^1(s'',\mathscr{T})(u,u')=\Gamma_k^2(s'',\mathscr{T})(u,u');
\end{aligned}
\end{equation}
\begin{equation}\tag{e3}
M_k^1\left(X,\mathscr{T}+t\cdot\mathscr{Z}^1(s')\right)(y,y')=M_k^2\left(X,\mathscr{T}+t\cdot\mathscr{Z}^2(s')\right)(y,y')\quad\text{for }|t|=\mathrm{rad}_1\qquad(k\ge1).
\end{equation}
In (e2) the letters $u$, $u'$ denote unit points, and the entries are $4\times4$ blocks.

\emph{A sufficient condition.}
According to (ga4), the hypothesis just written out is satisfied if the entries of the first line coincide and the other entries coincide as functions of $(s'',B)$ on $\mathrm{Pd}_k\times\tilde{G}_k$.
To prove this, assume that (e1) holds and that the entries of the family (p3), of $\Gamma_k$ and of $M_k$ at these indices coincide for the two systems as functions of $(s'',B)$ on $\mathrm{Pd}_k\times\tilde{G}_k$.
Then (e2) and (e3) hold.
Indeed, (e1) and (24.1) give $\mathscr{Z}^1(s')=\mathscr{Z}^2(s')$ on the bonds that meet $\mathbf{Y}$.
The connections $B^1$, $B^2$ (in (e3): $B^j=\mathscr{T}+t\cdot\mathscr{Z}^j(s')$ with $|t|=\mathrm{rad}_1$) lie in $\tilde{G}_k$ by Proposition~24.3\,(q1).
By (ga3) for $\mathrm{Sy}^1$, the entries of $\mathrm{Sy}^1$ at $B^1$ and at $B^2$ coincide, and by assumption the entries of $\mathrm{Sy}^1$ and of $\mathrm{Sy}^2$ at $B^2$ coincide.

\emph{The objects coincide.}
Assume (e1)--(e3).
By (e1) and (24.1), $\mathscr{Z}^1(s')=\mathscr{Z}^2(s')$ on the bonds that meet $\mathbf{Y}$.
Consequently the phases $U_B(x,x')$ at the pairs admissible for polymers $X\subseteq\mathbf{Y}$ coincide, because they read $B$ on bonds that meet $X$ (adm).
The elements $\mathscr{E}^{r}\left(X;A+t\cdot\mathscr{Z}^j(s'),\mathscr{C}\right)$ coincide by (k2), both first slots lying in $\tilde{R}_k$.
At $k=0$ the elements $\mathrm{Vx}\left(\mathscr{Z}^j(s'),\square\right)$, $\square\subseteq\mathbf{Y}$, coincide by Proposition~23.2\,(c).
Together with (e2), (e3) and the reduced displays above, this shows that the functions $G^{(i)}(s,X)$, $i\le3$, coincide for all $X\subseteq\mathbf{Y}$ and all $s\in\mathrm{Slc}(X,P)$ with $Y'(X,P)\subseteq\mathbf{Y}$.
Since the substitution contains no primitive, it follows that $E^{(i),*}(\mathbf{Y}')$, $E^{(i),x}(\mathbf{Y}')$ and $E^{(\le3),x}(\mathbf{Y}')$ coincide for every polymer $\mathbf{Y}'\subseteq\mathbf{Y}$.
Hence so do the functions of the second localisation, the elements $E^{(\le5),x}(\mathbf{Y}')$, $\mathrm{Km}(\mathbf{Y}')$ and $\mathrm{Hm}(\mathbf{Y}')$ for $\mathbf{Y}'\subseteq\mathbf{Y}$, and the terms
\[
\frac{1}{n!}\,\mathrm{urs}^{T}(X_1,\dots,X_n)\cdot\mathrm{Hm}(X_1)\cdots\mathrm{Hm}(X_n),\qquad X_i\subseteq\mathbf{Y},
\]
of the series (8.4) for $E^{\#,x}_k(m,\mathscr{E})(\mathbf{Y})$.
(For $E^{\#,x}_k(m,\mathscr{E})(\mathbf{Y})$ itself, the coincidence of the terms gives the assertion as soon as the series (8.4) are known to converge, Section~5.4; see Section~6.)
\end{proof}

\subsubsection*{3.10. (be)}

\begin{proof}[Proof of Proposition~24.3\,(be)]
Besides the primitives, the display (8.1)$^{\mathrm{ch}}$ contains $C$, $\delta_{\alpha,\pm B}$ and $\mathsf{C}$.
The operators $\delta_{\alpha,B}$ and $\delta_{\alpha,-B}$ occur in $H_k^{\#}(s'',B)$, hence in $J_k^{\#}\left(s'',\mathscr{T},\mathscr{Z}(s')\right)$, with $B=\mathscr{T}+\mathscr{Z}(s')$ and $B=\mathscr{T}$, and in $W_k^{\#}(s'',\mathscr{T})$, with $B=\mathscr{T}$.
These fields are added to the generators of $\mathscr{E}^{r}(X)$, $\langle\bar{\psi},\psi\rangle_{\square}$ and $E^{(\le3),x}(X)$, which contain $\chi(\boldsymbol{\zeta})$ only at pairs admissible for $X$ ((k2), (24.6)).
A homomorphism $\mathrm{sb}_v$ acts on an element only through the images of the generators that the element contains.
In a term attributed to a polymer $\mathbf{Y}'\subseteq\mathbf{Y}$, the polymer $X$ lies in $\mathbf{Y}'$ (Section~3.2).
This proves the first sentence of (be).

\emph{The case $k\ge1$.}
By (D4), $E^{(i),*}(\mathbf{Y})=\mathrm{Loc}\left[G^{(i)}\right](\mathbf{Y})$ contains no substitution, whereas
\[
E^{(i),x}(\mathbf{Y})=\mathrm{sub}_{\mathbf{Y},\mathscr{T}}\left(E^{(i),*}(\mathbf{Y})\right),\qquad E^{(\le3),x}(\mathbf{Y})=\sum_i\mathrm{sub}_{\mathbf{Y},\mathscr{T}}\left(E^{(i),*}(\mathbf{Y})\right).
\]
In $E^{(\le5),x}(\mathbf{Y})$, and through it in $\mathrm{Km}$, $\mathrm{Hm}$ and $E^{\#,x}_k$, the substitution occurs inside
\[
G^{(\le3)}(s'',X)=\mathrm{sb}_v\left(E^{(\le3),x}(X)\right),\qquad E^{(\le3),x}(X)=\sum_i\mathrm{sub}_{X,\mathscr{T}}\left(E^{(i),*}(X)\right),
\]
with $X\subseteq\mathbf{Y}$.
By (8.3a)$^{\mathrm{ch}}$, $\mathrm{sub}_{X,\mathscr{T}}$ contains the rows of $T^{(i)}(\mathscr{T},X)$ and of $\bar{T}^{(i)}(\mathscr{T},X)$ at the unit sites of $X$.
This proves the assertions of (be) on the substitution and on the rows of $T^{(i)}(\mathscr{T},\cdot\,)$ and $\bar{T}^{(i)}(\mathscr{T},\cdot\,)$.
\emph{The case $k=0$.}
No substitution occurs in the displays: $E^{(1),x}=E^{(1),*}$ and $G^{(\le3)}(s'',X)=\mathrm{sb}_v\left(E^{(1),*}(X)\right)$ (D4); the assertions of (be) on the substitution and on the rows of $T^{(i)}(\mathscr{T},\cdot\,)$ and $\bar{T}^{(i)}(\mathscr{T},\cdot\,)$ are void.

The functions $G^{(3)}(s,\square)$ and $G^{(5)}(s'',\square)$ contain $\langle\bar{\psi},\psi\rangle_{\square}$, and at $k=0$ one has $G^{(1)}(s,\square)=\mathrm{Vx}\left(\mathscr{Z}(s'),\square\right)$.
Being the polymers $X=\square$ of terms attributed to polymers contained in $\mathbf{Y}$, these cubes lie in $\mathbf{Y}$.
The matrix $C$ occurs in $\mathscr{Z}(s')$, that is, in (24.1).

\emph{What is read.}
The row of $\delta_{\alpha,\pm B}$ at a pair admissible for $X$ reads $B$ through $U_B(x,x')$, on the bonds of the paths of the pair; these are bonds contained in $X$, or the bond $(c,c_{3})$, which meets $X$ (adm).
For $k\ge1$ the rows of $T^{(i)}(\mathscr{T},X)$ and of $\bar{T}^{(i)}(\mathscr{T},X)$ at a unit site $y$ of $X$ read $\mathscr{T}$ on bonds that have an end in $\Delta_y$ (Proposition~23.1\,(b)), and $\Delta_y$ lies in $X$ (g0).
The element $\mathrm{Vx}(\mathrm{Zf},\square)$ reads $\mathrm{Zf}$ on the bonds that meet $\square$ (Proposition~23.2\,(c)).
Both $X$ and $\square$ lie in $\mathbf{Y}$.
The displays of $C$ (Section~I.5.7; \cite[Section~5.2.5]{QEDI}), of $M_i(x)$ (I.2.1), of $\mathsf{C}$ (Section~I.2.1) and of $\langle\bar{\psi},\psi\rangle_{\square}$ (Section~I.3.6) contain no connection.
\end{proof}

\subsubsection*{3.11. (de)}

\begin{proof}[Proof of Proposition~24.3\,(de)]
The integrations that occur in (D4), (8.2) and (8.4) are the following: the functionals $\oint_{\mathrm{rad}_i}$ over the circles $|t|=\mathrm{rad}_i$; the integrals $\int_{[0,1]^P}\mathrm{d}s_P$, where $P$ consists of parameter cubes contained in $\mathbf{Y}$ whenever $(X,P)$ is attributed to a polymer contained in $\mathbf{Y}$ (Section~3.2); and $\int\mathrm{d}\mu_I(W)\,\mathrm{d}\hat{\mu}(\mathsf{Z})$ in $\mathrm{Hm}$.

\emph{The integral over $W$.}
By (24.8), $\mathrm{Km}(\mathbf{Y})=\sum_Jc_J\cdot W(J)$ with $c_J\in\mathrm{Gr}[\psi,\chi]$, where $J$ ranges over sets of unit indices $u$ in $\mathbf{Y}$.
The identity for $\int\mathrm{d}\mu_I(W)\,F$ in (N3)$^{\mathrm{w}}$ (Proposition~15\,(f)) gives
\[
\int\mathrm{d}\mu_I(W)\,\mathrm{Km}(\mathbf{Y})=\sum_Jc_J\cdot m(J).
\]
By (mo), that is, (I.B4), one has $m(J)=0$ unless $J$ is paired; for paired $J$,
\[
m(J)=\pm\prod\nu^{-1},
\]
the product extending over the pairs $\left\{W_\beta(\mathrm{tu}),\bar{W}_\beta(\mathrm{tu})\right\}$ contained in $J$, with $\nu=1$ for $k\ge1$ and $\nu=\frac{1}{48}$ for $k=0$; at $k=0$ this reads $m(J)=\pm48^{\frac{\#|J|}{2}}$.
By (N3)$^{\mathrm{w}}$ the integral does not see the modes that do not occur.
Applied with $\mathrm{Mo}^0$ the set of the modes $(\mathrm{tu},\beta)$ with $\mathrm{tu}$ in $\mathbf{Y}$, this assertion shows that the integral is the integral (I.3.5) over the generators $W(u)$, $u$ in $\mathbf{Y}$.
The factorisation over the unit points is (N4), that is, Proposition~15\,(g).

\emph{The integral over $\mathsf{Z}$.}
The space $V_k$ is equipped with the counting scalar product, $\mu_I$ is the standard Gaussian measure on $V_k$, and $\mathsf{Z}(\Upsilon)=\langle\mathsf{Z},\Upsilon\rangle$; moreover
\[
\mathrm{Box}_k=\left\{\mathsf{Z}:\ |\mathsf{Z}(\Upsilon)|\le p_{0,k}\text{ for all }\Upsilon\right\},\qquad\mathrm{d}\hat{\mu}=\frac{\mathbf{1}_{\mathrm{Box}_k}\,\mathrm{d}\mu_I}{\mu_I(\mathrm{Box}_k)}
\]
(Section~I.5.7; \cite[Section~5.2.5]{QEDI}).
By (P1) of Section~0.6 (\cite[Proposition~15.21]{QEDI}), the family $\{\Upsilon\}$ is an orthonormal basis.
Hence $|\mathsf{Z}|^2=\sum_\Upsilon\mathsf{Z}(\Upsilon)^2$, and in the coordinates $\mathsf{Z}(\Upsilon)$ the measure $\mu_I$ is the product of standard normal laws, $\mathrm{Box}_k$ is the product of the intervals $[-p_{0,k},p_{0,k}]$, and $\hat{\mu}$ is the product of copies of the standard normal law restricted to $[-p_{0,k},p_{0,k}]$ and normalised.
By (24.8) the integrand depends on the coordinates $\mathsf{Z}(\Upsilon)$ with $\Upsilon$ contained in $\mathbf{Y}$ only; the other factors have mass one.
\end{proof}

\subsubsection*{3.12. (et)}

\begin{proof}[Proof of Proposition~24.3\,(et)]
In (D4), $\mathrm{Loc}$ is applied to the functions $G^{(i)}$, $i\le3$ (at $k=0$ only $i=1$), with the kinds $\pi_k$, $\pi'_k$, $\pi''_k$, and to $G^{(\le3)}+G^{(4)}+G^{(5)}$ (at $k=0$, $G^{(4)}=0$), with the kind $\pi_k$.
Let $G$ be one of them, let $X\in D_k$, and let $P$ be a set of parameter cubes, of the kinds that occur, outside $P_0(X)$.
The families $s\in\mathrm{Pd}_k$ that vanish outside $P_0(X)\cup P$ are precisely the elements of $\mathrm{Slc}(X,P)$.

\emph{The terms that vanish.}
Suppose that $Y'(X,P)$ is not connected.
For $G=G^{(i)}$, $i\le3$, the assertion is Lemma~3.B\,(b).
The cube $\mathrm{pc}$ produced in the paragraph (b) of the proof in Section~3.5 is a cube of $P$ that is not contained in $Y'_1$, and its choice depends on $(X,P)$ alone.
By Lemma~3.C\,(a), one and the same cube $\mathrm{pc}$ therefore serves for $G^{(\le3)}$, $G^{(4)}$ and $G^{(5)}$, hence for their sum, and $\mathrm{Tm}[G](X,P)=0$ by the pairing of the terms of (V) used in that paragraph.
(A function $G(\,\cdot\,,\square)$ is $0$ on the polymers that are not $M$-cubes (D3); for these polymers there is nothing to prove.)

\emph{The identity.}
By Lemma~3.B\,(reg) and Lemma~3.C\,(reg), $G(\,\cdot\,,X)$ is continuous on the compact set $\mathrm{Pd}_k$, hence bounded, and analytic in its interior; moreover $M^{\alpha_0}>1$ by (K5).
Lemma~6.1 of Section~0.6 (\cite[Lemma~17.12]{QEDI}), applied with $\mathrm{Pc}$ the set of all parameter cubes of the kinds that occur, $P_0=P_0(X)$, $\mathrm{rad}=M^{\alpha_0}$, $f=G(\,\cdot\,,X)$ and $f_P=\mathrm{Tm}[G](X,P)$, yields
\[
G(1,X)=\sum_P\mathrm{Tm}[G](X,P),
\]
the sum extending over the sets $P$ of parameter cubes outside $P_0(X)$.
The terms for which $Y'(X,P)$ is not connected vanish.
For each of the remaining terms, $Y'(X,P)$ is a polymer, and the term is attributed to this polymer and to no other.
As $D_k$ is finite, summation over $X$ gives
\[
\sum_{X\in D_k}G(1,X)=\sum_{\mathbf{Y}\in D_k}\ \sum_{(X,P):\ Y'(X,P)=\mathbf{Y}}\mathrm{Tm}[G](X,P)=\sum_{\mathbf{Y}\in D_k}\mathrm{Loc}[G](\mathbf{Y}).
\]
This is the first identity of Lemma~6.2 of Section~0.6 (\cite[Lemma~17.13]{QEDI}).

\emph{The expansion.}
By (24.7) the elements $E^{(\le5),x}(\mathbf{Y})$ are even, and even elements commute by (N1), that is, Proposition~15\,(b).
Let $U$ be a union of $M$-cubes.
Then
\[
\exp\left\{\sum_{\mathbf{Y}\in D_k,\ \mathbf{Y}\subseteq U}E^{(\le5),x}(\mathbf{Y})\right\}=\prod_{\mathbf{Y}\in D_k,\ \mathbf{Y}\subseteq U}\left(1+\left(\exp E^{(\le5),x}(\mathbf{Y})-1\right)\right)=\sum_Z\mathrm{Km}(Z),
\]
where $Z$ ranges over the unions of $M$-cubes contained in $U$, the empty union included.
For the last equality one expands the product: a set $\{X_i\}$ of distinct polymers contained in $U$ is counted in $\mathrm{Km}(Z)$, $Z$ being the union of the $X_i$ (the empty set is counted in $\mathrm{Km}(\emptyset)=1$), and the members of a set whose union is $Z\subseteq U$ are contained in $U$.
Finally, $\mathrm{Km}(Z)$ is the product of the $\mathrm{Km}(Z_a)$ over the components $Z_a$ of $Z$.
Indeed, each polymer of a set whose union is $Z$ lies in one component, so that the sets whose union is $Z$ are the unions of sets whose unions are the $Z_a$.
\end{proof}

\subsection*{4. The statement of Proposition~24.1}

Proposition~24.1, which is stated in Part~II, assumes the setting (ST1)$^{\mathrm{ch}}$ of Section~G3.4.0 together with the step map of Definition~24.D.
It consists of the clauses (a)--(g) and~(i), with the following content.

Clause~(a) concerns the data $(m,\mathscr{E})\in\mathrm{Ball}^{x}_{k}(4)$, with $\mathscr{E}=0$ when $k=0$, and the arguments $A\in\tilde{R}_k^{-}$ and $\mathscr{C}\in\mathfrak{X}_k$.
For these, every display of Definition~24.D is defined, in the order fixed at the end of that definition, the series (8.4) converge absolutely in $\|\cdot\|_{h_k/4}$, and for every $X\in D_k$ one has
\[
\left\|E^{\#,x}_k(m,\mathscr{E})(X;A,\mathscr{C})\right\|_{h_k/4}\le14\,c_X^2\cdot K_0^2\cdot\exp(3\kappa)\cdot\mathbf{H}_0\cdot\exp\left(-(\kappa-5\kappa_0-4)\cdot d_M(X)\right).
\]
The family $E^{*,x}_k(m,\mathscr{E})$ is an element of $K^{x,C}_{k+1}=K^{x,C}_{k+1}(\ell/L)$, and $\left\|E^{*,x}_k(m,\mathscr{E})\right\|^{x}_{k+1}\le e_k^{\frac{1}{4}-6\varepsilon}$.
Its dependence on the couplings is analytic in the following sense.
For $(m_a,\mathscr{E}_a)$ and $(m_b,\mathscr{E}_b)$ in $\mathbb{C}\times K^{x,C}_k$, with $\mathscr{E}_a=\mathscr{E}_b=0$ when $k=0$, the map $t\mapsto E^{*,x}_k\left(m_a+t\cdot m_b,\,\mathscr{E}_a+t\cdot\mathscr{E}_b\right)$ into $K^{x,C}_{k+1}$ is continuous on the closed set of the complex numbers $t$ for which the argument lies in $\mathrm{Ball}^{x}_{k}(4)$, and analytic in the interior of this set: near every interior point it is the sum of a power series in $t$ with coefficients in $K^{x,C}_{k+1}$ that converges in $\|\cdot\|^{x}_{k+1}$.
For real $m$ and $\mathscr{E}\in K^{x}_{k}$ the family $E^{*,x}_k(m,\mathscr{E})$ lies in $K^{x}_{k+1}$.
The last assertion of (a) is the bound $0\le\delta\varepsilon^{0}_{k}\le e_k^{7}$.

Clause~(b) concerns the reading~(o) of Definition~24.D\,(D7) and the restriction $\,\cdot\,|_0$ to $\mathscr{C}=0$ of Section~I.5.4.
In the reading~(o) the first assertion of (a) holds: for every $(m,E)\in\mathrm{Ball}_{k}(4)$, with $E=0$ when $k=0$, and every $A\in\tilde{R}_k^{-}$, the displays of Definition~24.D, read according to (D7), are defined in the order fixed there, and the series (8.4) converge absolutely in $\|\cdot\|_{h_k/4}$, with the bound displayed above.
The map $E^{*}_k$ of the reading~(o) sends $\mathrm{Ball}_{k}(4)$ into $K^{C}_{k+1}$, and the pairs $(m,E)\in\mathrm{Ball}_{k}(4)$ with real $m$ and $E\in K_k$ into $K_{k+1}$;
it satisfies $\left\|E^{*}_k(m,E)\right\|_{k+1}\le e_k^{\frac{1}{4}-6\varepsilon}$, and it is analytic in the sense of (a), with $\mathbb{C}\times K^{C}_{k}$, $\mathrm{Ball}_{k}(4)$, $K^{C}_{k+1}$ and $\|\cdot\|_{k+1}$ in the places of $\mathbb{C}\times K^{x,C}_{k}$, $\mathrm{Ball}^{x}_{k}(4)$, $K^{x,C}_{k+1}$ and $\|\cdot\|^{x}_{k+1}$.
For $(m,\mathscr{E})\in\mathrm{Ball}^{x}_{k}(4)$ the pair $(m,\mathscr{E}|_0)$ lies in $\mathrm{Ball}_{k}(4)$, and $E^{*,x}_k(m,\mathscr{E})|_0=E^{*}_k\left(m,\mathscr{E}|_0\right)$.

Clause~(c) is the identity of the step.
Its hypotheses are the following: $\varepsilon^{\mathrm{v}}$ and $m$ are real numbers with $|m|\le w^{\mathrm{m}}_k$;
$\mathscr{E}\in K^{x}_{k}$ satisfies $\|\mathscr{E}\|^{x}_{k}\le w^{\mathrm{E}}_k$, and $\mathscr{E}=0$ if $k=0$;
$\mathfrak{P}_k=\left\{\mathfrak{P}_k(X,\cdot\,):X\in D_k\right\}$ is an arbitrary family of complex functions on $\mathfrak{X}_k$;
$A'\in\mathrm{Dom}_{k+1}$ has an arbitrary harmonic part, and $\mathscr{C}'\in\mathfrak{X}_{k+1}$.
The triple $(\varepsilon^{\mathrm{v},+},m^{+},\mathscr{E}^{+})$ and the family $\mathfrak{P}_{k+1}$ are given by (8.6)$^{\mathrm{ch}}$, and $\mathscr{F}^{x}_{k}$ is formed with $\mathfrak{P}_k$, $\mathscr{F}^{x}_{k+1}$ with $\mathfrak{P}_{k+1}$ (Section~I.5.7).
The clause asserts the identity of the step (CV.4x)$^{\mathrm{ch}}$ and the identity (2.4x)$^{\mathrm{ch}}$ for $y^{x}$, which read
\begin{align*}
\mathbf{T}^{x}_{k}\mathscr{F}^{x}_{k}[\varepsilon^{\mathrm{v}},m,\mathscr{E}](A',\mathscr{C}';\,\cdot\,)
&=\mathscr{F}^{x}_{k+1}[\varepsilon^{\mathrm{v},+},m^{+},\mathscr{E}^{+}](A',\mathscr{C}';\,\cdot\,),\\
y^{x}_{k+1}[\varepsilon^{\mathrm{v},+},m^{+},\mathscr{E}^{+}](A';\mathscr{C}')
&=\int_{\mathrm{Box}_k}\mathrm{d}\mu_I(\zeta)\,y^{x}_{k}[\varepsilon^{\mathrm{v}},m,\mathscr{E}]\left(A^{\min}_{k}(A')+\mathrm{Sq}_{k}\,\zeta;\,\mathscr{C}'_L\right).
\end{align*}
At level $0$ the first identity takes the form
\[
\mathbf{T}^{x}_{0}\mathscr{F}^{x}_{0}[\varepsilon^{\mathrm{v}},m,0](A',\mathscr{C}';\,\cdot\,)
=\mathscr{F}^{x}_{1}\left[L^3\cdot\left(\varepsilon^{\mathrm{v}}+\delta\varepsilon^{0}_{0}\right),\,L\cdot m,\,L\,{\delta E^{\mathrm{det}}}_0+E^{*,x}_0(m,0)\right](A',\mathscr{C}';\,\cdot\,),
\]
where the left side is an integral over $192\,n_0$ pairs of generators $\Psi$ on $T^{\mathrm{f}}_0$ and the right side is an element of $\mathrm{Gr}[\Psi']$ with $4\,n_1$ pairs of generators.
At $\mathscr{C}'=0$, and for a family with $\mathfrak{P}_k(\cdot\,,0)=0$, the two identities are the identities for the objects $\mathscr{F}_k$, $y_k$, $\mathbf{T}_k$ without slot, with the data $(\varepsilon^{\mathrm{v}},m,\mathscr{E}|_0)$ and their image $(\varepsilon^{\mathrm{v},+},m^{+},\mathscr{E}^{+}|_0)$;
in this case $\mathfrak{P}_{k+1}(\cdot\,,0)=0$.

The full step is described by clause~(d).
In the situation of (c) the family $\mathscr{E}^{+}$ lies in $K^{x}_{k+1}$,
\[
\mathscr{E}^{+}|_0=L R\,(\mathscr{E}|_0)+L\,E_k^{\mathrm{det}}+E^{*}_k\left(m,\mathscr{E}|_0\right),
\]
and
\begin{gather*}
\left\|\mathscr{E}^{+}\right\|^{x}_{k+1}\le c_{\mathrm{fl}}\cdot L^{-\frac{1}{4}+2\varepsilon}\cdot\|\mathscr{E}\|^{x}_{k}+e_{k+1}^{\frac{1}{4}-2\varepsilon}+e_k^{\frac{1}{4}-6\varepsilon},\\
\left|m^{+}-L\cdot m\right|\le L\cdot c_{\mathrm{fl}}\cdot e_k^{\frac{1}{2}}\cdot\|\mathscr{E}\|^{x}_{k},\qquad
\left|\varepsilon^{\mathrm{v},+}-L^3\cdot\varepsilon^{\mathrm{v}}\right|\le L^3\cdot\left(c_{\mathrm{fl}}\cdot\|\mathscr{E}\|^{x}_{k}+e_k^{7}\right).
\end{gather*}
For $k=0$ the clause states $\left\|\mathscr{E}^{+}\right\|^{x}_{1}\le\left(L^{\frac{1}{2}}\cdot q\right)^{\frac{1}{4}-2\varepsilon}+q^{\frac{1}{4}-6\varepsilon}$, $m^{+}=L\cdot m$ and $0\le\varepsilon^{\mathrm{v},+}-L^3\cdot\varepsilon^{\mathrm{v}}\le L^3\cdot q^{7}$.

Clause~(e) bounds the dependence of $E^{*,x}_k$ on the couplings: for $(m_1,\mathscr{E}_1)$ and $(m_2,\mathscr{E}_2)$ in $\mathrm{Ball}^{x}_{k}(1)$,
\[
\left\|E^{*,x}_k(m_1,\mathscr{E}_1)-E^{*,x}_k(m_2,\mathscr{E}_2)\right\|^{x}_{k+1}
\le e_k^{\frac{1}{4}-6\varepsilon}\cdot\left(\frac{|m_1-m_2|}{w^{\mathrm{m}}_k}+\frac{\|\mathscr{E}_1-\mathscr{E}_2\|^{x}_{k}}{w^{\mathrm{E}}_k}\right),
\]
and the same bound holds for $E^{*}_k$ and for $(m_1,E_1)$, $(m_2,E_2)$ in $\mathrm{Ball}_{k}(1)$, with the norms $\|\cdot\|_k$ and $\|\cdot\|_{k+1}$.
For $k=0$ the second quotient is absent from these bounds.

Clause~(f) lists symmetries.
For $(m,\mathscr{E})\in\mathrm{Ball}^{x}_{k}(4)$ the family $E^{*,x}_k(m,\mathscr{E})$ has, at level $k+1$ and in the terms of Section~I.5.4, the clause (k3) for the group $\mathbb{G}_{k+1}$ of the side $\ell/L$, which acts on the fermion sites by (I.4.2), the clause (k4) in both connections, and the clause (k5).
Every monomial that occurs in it has as many barred legs as unbarred legs, and for real $m$ and $\mathscr{E}\in K^{x}_{k}$ it has the clause (k7) as well.

Locality and analyticity form clause~(g).
For $(m,\mathscr{E})\in\mathrm{Ball}^{x}_{k}(4)$ the family $E^{*,x}_k(m,\mathscr{E})$ has the clauses (k1) and (k2) at level $k+1$: its kernels are analytic on $\tilde{R}_{k+1}\times\mathfrak{X}_{k+1}$, they vanish unless all their sites lie over the polymer and all their pairs are of kind (I) or of kind (II) for it (Section~I.1.7), and they depend on the two connections through the bonds that meet the polymer.
For $Y\in D_{k+1}$, each H\"older leg of $E^{*,x}_k(m,\mathscr{E})(Y)$ sits at the image, under $(x,x')\mapsto(x/L,x'/L)$, either of a pair at which some $\mathscr{E}^{r}(X;B_1,\mathscr{C})$ with $X\subseteq LY$ and $(B_1,\mathscr{C})\in\tilde{R}_k\times\mathfrak{X}_k$ has a H\"older leg, or of a pair $(x,x^{(3)})$, $x=(c,g)$, with $c\in LY$;
the pairs of the second sort arise from (8.3a)$^{\mathrm{ch}}$ when $k\ge1$ and from (8.3$_0$)$^{\mathrm{ch}}$ when $k=0$.
In particular every H\"older leg of $E^{*,x}_0(m,0)$ is a pair $(x',{x'}^{(3)})$ of $\mathrm{Pz}_1$, of length $d_X=\eta_1$.
The structure of the objects attached to a polymer of level $k$ is described in Proposition~24.3, and the exact operators that occur in them in Propositions~25.1--25.3.

Clause~(i), finally, concerns uniformity.
Apart from numerals, from the norms of the data and from $d_M(X)$, and apart from $\varepsilon$, which enters through exponents only, the bounds displayed in (a)--(e) contain the datum $e_k$ (also by way of $e_{k+1}=L^{\frac{1}{2}}\cdot e_k$, $w^{\mathrm{m}}_k$, $w^{\mathrm{E}}_k$ and $\mathbf{H}_0$) and the constants $L$, $c_{\mathrm{fl}}$, $c_X$, $K_0$, $\kappa$, $\kappa_0$, and nothing else.
With these numbers they hold for all data $(k,\ell,\mu,e_k)$ of (ST1)$^{\mathrm{ch}}$, for every system of primitives with the properties required there, for every harmonic part of $A'$, for every $\mathscr{C}'$ and for all $(m,\mathscr{E})$ in the balls named.

\subsection*{5. Proof of Proposition~24.1\,(a), first sentence with its display, and of Proposition~24.1\,(b), first sentence}

\subsubsection*{5.1. Three operator bounds}

The operator $H_k^{\#}(s'',B)$ consists of four operators from the spinors on the unit points to the spinors on $T^{\mathrm{f}}_k$, respectively to the functions of pairs: $H_k(s'',B)$ and $\delta_{\alpha,B}H_k(s'',B)$ (unbarred fields), $\bar{H}_k(s'',B)$ and $\delta_{\alpha,-B}\bar{H}_k(s'',B)$ (barred fields).
The number $\left|H_k^{\#}(s'',B)\right|_{1,\infty}$ is the largest of their four norms $|\cdot|_{1,\infty}$, where
\[
|O|_{1,\infty}=\sup\ \sum_{\beta'}\int\mathrm{d}\mathrm{tu}'\,\left|O(\,\cdot\,,\mathrm{tu}')_{\beta\beta'}\right|,
\]
the supremum extending over the rows $(x,\beta)$, respectively $\left((x,x'),\beta\right)$ with $(x,x')\in\mathrm{Pz}_k$, which are inserted in the place of the dot and of $\beta$ ((I.ke3), Section~I.5.2; at $k=0$, $\int\mathrm{d}\mathrm{tu}'=\frac{1}{48}\sum_{x'}$).
The same convention is used for $J_k^{\#}$ and for $H_k^{\#}\cdot\tilde{\Gamma}_k$.
The properties W2, W5, W6, W12 and W17 of the system of primitives of (ST1)$^{\mathrm{ch}}$ are the properties 2, 5, 6, 12 and 17 of \cite[Section~5.4.2]{QEDI}, read as in Section~I.5.8, for the parameters $(\alpha_0,\gamma_{\mathrm{w}},\kappa_{\mathrm{w}};C_{w})$.

\begin{named}{Lemma 5.A}
Let $s\in\mathrm{Pd}_k$, $A\in\tilde{R}_k^{-}$, $\mathscr{C}\in\mathrm{Sd}_k$ and $\mathscr{T}=A+\mathscr{C}$.
\begin{itemize}
\item[(q2)] For $k\ge1$, $B\in\tilde{G}_k$, $X\in D_k$ and $\mathsf{w}_3>0$,
\[
\left\|\langle\bar{\Psi},M_k(X,B)\Psi\rangle_1\right\|_{\mathsf{w}_3}\le4c_{X}\cdot c_{\gamma}\cdot C_{w}\cdot M^3\cdot\mathsf{w}_3^2\cdot\exp\left(-(\kappa_{\mathrm{w}}-1)\cdot d_M(X)\right).
\]
\item[(q3)] For $B\in\tilde{G}_k$,
\[
\left|H_k^{\#}(s'',B)\right|_{1,\infty}\le4C_{w}.
\]
For a complex bond function $\mathrm{Zf}$ whose three norms (Section~I.5.3) are at most $\frac{1}{2}e_k^{-\frac{\varepsilon}{2}}$,
\[
\left|J_k^{\#}(s'',\mathscr{T},\mathrm{Zf})\right|_{1,\infty}\le16C_{w}\cdot e_k^{\frac{3}{4}-\frac{9\varepsilon}{2}}.
\]
\item[(q4)] One has
\[
\left|H_k^{\#}(s'',\mathscr{T})\cdot\tilde{\Gamma}_k(s'',\mathscr{T})\right|_{1,\infty}\le4C_{w}^2\cdot c_{\gamma},
\]
and hence, for every $k\ge0$,
\[
\mathsf{w}_W\cdot\left|H_k^{\#}(s'',\mathscr{T})\cdot\tilde{\Gamma}_k(s'',\mathscr{T})\right|_{1,\infty}\le28C_{w}^2\cdot c_{\gamma}=C_{W}.
\]
\end{itemize}
\end{named}

\begin{proof}
(q2)
The element in question is
\[
\langle\bar{\Psi},M_k(X,B)\Psi\rangle_1=\sum_{y,y'}\ \sum_{\beta,\beta'}M_k(X,B)(y,y')_{\beta\beta'}\cdot\bar{\Psi}_\beta(y)\cdot\Psi_{\beta'}(y').
\]
Up to signs, the products $\bar{\Psi}_\beta(y)\cdot\Psi_{\beta'}(y')$ are distinct basis monomials of degree two; hence the norm (I.5.1) equals $\mathsf{w}_3^2$ times the sum of the moduli of the entries.
By W12 the entries vanish unless $y$ and $y'$ lie in $X$.
The polymer $X$ contains $M^3\cdot|X|_M$ unit sites, so that there are $4M^3\cdot|X|_M$ rows $(y,\beta)$.
By W6,
\[
\left|M_k(X,B)(y,y')\right|_c\le C_{w}\cdot\exp\left(-\gamma_{\mathrm{w}}\cdot d(y,y')-\kappa_{\mathrm{w}}\cdot d_M(X)\right),
\]
and $|\cdot|_c$ is the largest row sum of the $4\times4$ block (Section~I.1.5).
Hence in the row $(y,\beta)$ the moduli add up to at most
\[
\sum_{y'}C_{w}\cdot\exp\left(-\gamma_{\mathrm{w}}\cdot d(y,y')\right)\cdot\exp\left(-\kappa_{\mathrm{w}}\cdot d_M(X)\right)\le c_{\gamma}\cdot C_{w}\cdot\exp\left(-\kappa_{\mathrm{w}}\cdot d_M(X)\right).
\]
Finally, $|X|_M\le c_{X}\cdot\exp\left(d_M(X)\right)$ by (t1) and (t6) of Section~0.6 (\cite[Theorem~13.2]{QEDI}).
The product of these bounds is the bound (q2).

(q3)
By (rs), that is, Proposition~15\,(c), it suffices to bound $|O\,v|_\infty$, for each of the four operators $O$, for a spinor $v$ on the unit points with $|v|_\infty\le1$.
Let $f:=H_k(s'',B)\,v$.

\emph{The case $k\ge1$.}
By W5, $|f|_\infty\le C_{w}$, and $\left|(\delta_{\alpha,B}f)(x,x')\right|_c\le C_{w}$ at every short pair, that is, for $d_X<1$; the pairs inside a cluster are among them, with $d_X=\eta_k$.
At a long pair, $1\le d_X\le M$, the bound (I.Lc2)$^{\mathrm{ch}}$ gives $\left|(\delta_{\alpha,B}f)(x,x')\right|_c\le4|f|_\infty\le4C_{w}$.

\emph{The case $k=0$.}
Here $f(x)=\mathrm{hq}(c)\cdot v(x)$.
By W17, which gives W5 with the constant $4\exp\left(2\gamma_{\mathrm{w}}\right)\cdot M^{27\alpha_0}$, and by (K6), one has $|f|_\infty\le C_{w}$.
Every pair of $\mathrm{Pz}_0$ is long: $d_X\ge\eta_0=1$, and for a pair inside a cluster $d_X=1$ and $(\delta_{\alpha,B}f)(x,x')=f(x')-f(x)$.
By (I.Lc2)$^{\mathrm{ch}}$, $\left|(\delta_{\alpha,B}f)(x,x')\right|_c\le4C_{w}$.

\emph{The barred fields.}
One has $\bar{H}_k(s'',B)=\mathsf{C}^{-1}\cdot H_k(s'',-B)\cdot\mathsf{C}$ and, since $\delta_{\alpha,-B}$ commutes with constant matrices,
\[
\delta_{\alpha,-B}\bar{H}_k(s'',B)=\mathsf{C}^{-1}\cdot\left[\delta_{\alpha,-B}H_k(s'',-B)\right]\cdot\mathsf{C}.
\]
The connection $-B$ lies in $\tilde{G}_k$ by DF3 (\cite[Lemma~17.6]{QEDI}).
The matrices $\mathsf{C}$ and $\mathsf{C}^{-1}=-\mathsf{C}$ are monomial with entries of modulus one ((I.c3), (I.c6)); so they preserve $|\cdot|_c$.
The bounds for $-B$ therefore give the bounds for the barred fields.

\emph{The second bound.}
By Proposition~24.3\,(q1), $\mathscr{T}+t\cdot\mathrm{Zf}$ lies in $\tilde{G}_k$ for $|t|\le\mathrm{rad}_1$, and the set of the $t$ with $\mathscr{T}+t\cdot\mathrm{Zf}\in\tilde{G}_k$ is open by DF0 (\cite[Lemma~17.3]{QEDI}).
By W2, for fixed $s''\in\mathrm{Pd}_k$, and by (24.3), the map $t\mapsto H_k^{\#}(s'',\mathscr{T}+t\cdot\mathrm{Zf})\,v$ is analytic in a neighbourhood of $\{|t|\le\mathrm{rad}_1\}$.
By (N5), that is, Proposition~15\,(j), and by (24.2), the field
\[
J_k^{\#}(s'',\mathscr{T},\mathrm{Zf})\,v=\oint_{\mathrm{rad}_1}H_k^{\#}(s'',\mathscr{T}+t\cdot\mathrm{Zf})\,v
\]
has sup norm at most
\[
\frac{4C_{w}}{\mathrm{rad}_1-1}\le\frac{8C_{w}}{\mathrm{rad}_1}=16C_{w}\cdot e_k^{\frac{3}{4}-\frac{9\varepsilon}{2}},
\]
the first bound being applied at $B=\mathscr{T}+t\cdot\mathrm{Zf}$, $|t|=\mathrm{rad}_1$.
By (rs) this gives the bound, on the unbarred and on the barred fields.

(q4)
Consider first the field $W$.
By W6 and the definition of $c_{\gamma}$ in Section~0.6, for every unit point $\mathrm{tu}$ and every $\beta$,
\begin{align*}
\sum_{\beta'}\int\mathrm{d}\mathrm{tu}'\,\left|\Gamma_k(s'',\mathscr{T})(\mathrm{tu},\mathrm{tu}')_{\beta\beta'}\right|
&\le\int\mathrm{d}\mathrm{tu}'\,\left|\Gamma_k(s'',\mathscr{T})(\mathrm{tu},\mathrm{tu}')\right|_c\\
&\le C_{w}\cdot\int\mathrm{d}\mathrm{tu}'\,\exp\left(-\gamma_{\mathrm{w}}\cdot d(\mathrm{tu},\mathrm{tu}')\right)\le C_{w}\cdot c_{\gamma};
\end{align*}
at $k=0$ the kernel is that of (I.1.4), $\int\mathrm{d}\mathrm{tu}'=\frac{1}{48}\sum_{x'}$, and the weight $\frac{1}{48}$ is the weight in the definition of $c_{\gamma}$.
Hence $\left|\Gamma_k(s'',\mathscr{T})\right|_{1,\infty}\le C_{w}\cdot c_{\gamma}$, and (q3) together with the submultiplicativity of $|\cdot|_{1,\infty}$ gives
\[
\left|H_k^{\#}(s'',\mathscr{T})\cdot\Gamma_k(s'',\mathscr{T})\right|_{1,\infty}\le4C_{w}^2\cdot c_{\gamma}.
\]
On $\bar{W}$ the second factor of $\tilde{\Gamma}_k$ is the identity, and $4C_{w}\le4C_{w}^2\cdot c_{\gamma}$ by (K6).
Finally, $\mathsf{w}_W\le7$.
\end{proof}

\subsubsection*{5.2. The localised functions}

\begin{named}{Proposition 5.B}
Let $(m,\mathscr{E})\in\mathrm{Ball}^{x}_{k}(4)$, $A\in\tilde{R}_k^{-}$, $\mathscr{C}\in\mathrm{Sd}_k$, $\mathsf{Z}\in\mathrm{Box}_{k}$ and $\mathbf{Y}\in D_k$.
\begin{itemize}
\item[(a)] For every $X\in D_k$ and every $(B_1,\mathscr{C})\in\tilde{R}_k\times\mathrm{Sd}_k$,
\[
\left\|\mathscr{E}^{r}(X;B_1,\mathscr{C})\right\|_{h_k}
\le2\,\|\mathscr{E}\|^{x}_{k}\cdot\exp\left(-\kappa\cdot d_M(X)\right)
\le8\,w^{\mathrm{E}}_k\cdot\exp\left(-\kappa\cdot d_M(X)\right)
\le12\,w^{\mathrm{E}}_k\cdot\exp\left(-\kappa\cdot d_M(X)\right);
\]
moreover $|m'|\le8\,w^{\mathrm{m}}_k$.
\item[(b)] One has
\[
\left\|E^{(\le3),x}(\mathbf{Y})\right\|_{h_k/2}\le8\,c_X\cdot K_0\cdot e_k^{\frac{1}{4}-5\varepsilon}\cdot\exp\left(-\mathrm{rt}_1\cdot d_M(\mathbf{Y})\right).
\]
\item[(c)] One has
\[
\left\|E^{(\le5),x}(\mathbf{Y})\right\|_{(h_k/4;\,\mathsf{w}_W)}\le\mathbf{H}_{0}\cdot\exp\left(-\mathrm{rt}_2\cdot d_M(\mathbf{Y})\right).
\]
\item[(d)] (The identity of Proposition~24.3\,(et).)
Let $1$ denote the family that takes the value $1$ at every parameter cube.
Then
\[
\sum_{\mathbf{Y}\in D_k}E^{(i),*}(\mathbf{Y})=\sum_{X\in D_k}G^{(i)}(1,X)\quad(i\le3),
\qquad
\sum_{\mathbf{Y}\in D_k}E^{(\le5),x}(\mathbf{Y})=\sum_{X\in D_k}\left(G^{(\le3)}+G^{(4)}+G^{(5)}\right)(1,X).
\]
\item[(e)] For every $(m,\mathscr{E})\in\mathbb{C}\times K^{x,C}_{k}$ (with $\mathscr{E}=0$ if $k=0$) the displays (8.1)$^{\mathrm{ch}}$--(8.3)$^{\mathrm{ch}}$ define $E^{(\le3),x}(\mathbf{Y})$ and $E^{(\le5),x}(\mathbf{Y})$, and (24.5)--(24.7) hold.
The element $E^{(\le5),x}(\mathbf{Y})$ is the sum of an element that does not contain $(m,\mathscr{E})$ and an element that is complex linear in $(m,\mathscr{E})$.
\end{itemize}
\end{named}

\begin{proof}
\emph{(a).}
If $k=0$, then $\mathscr{E}^{r}=0$ and $m'=m$, and $|m|\le4\,w^{\mathrm{m}}_0$.
Let $k\ge1$.
By (Sd1) the pair $(0,0)$ belongs to $\tilde{R}_k\times\mathrm{Sd}_k$.
Hence (X2) and (k6) of Section~I.5.4 give
\[
\left\|\mathscr{E}^{r}(X;B_1,\mathscr{C})\right\|_{h_k}
\le\left\|\mathscr{E}(X;B_1,\mathscr{C})\right\|_{h_k}+\left\|\mathscr{E}(X;0,0)\right\|_{h_k}
\le2\,\|\mathscr{E}\|^{x}_{k}\cdot\exp\left(-\kappa\cdot d_M(X)\right),
\]
while $\|\mathscr{E}\|^{x}_{k}\le4\,w^{\mathrm{E}}_k$, and $2\cdot4=8\le12$.
By (X1),
\[
|m'|\le4\,w^{\mathrm{m}}_k+c_{\mathrm{fl}}\cdot e_k^{\frac{1}{2}}\cdot4\,w^{\mathrm{E}}_k
=4\,w^{\mathrm{m}}_k\cdot\left(1+c_{\mathrm{fl}}\cdot e_k^{\frac{1}{2}+\frac{1}{4}-7\varepsilon-\frac{3}{4}+8\varepsilon}\right)
=4\,w^{\mathrm{m}}_k\cdot\left(1+c_{\mathrm{fl}}\cdot e_k^{\varepsilon}\right)
\le8\,w^{\mathrm{m}}_k,
\]
where the last step is (A1).

\emph{(f0) The form $\langle\bar{\psi},\psi\rangle_{\square}$, for every $k\ge0$.}
One has
\[
\langle\bar{\psi},\psi\rangle_{\square}=\frac{\eta_k^3}{48}\sum_{c\in\square}\sum_{g\in O_{\mathrm{cub}}}\sum_{\beta}\bar{\psi}_{\beta}(c,g)\cdot\psi_{\beta}(c,g).
\]
The right side consists of $4\cdot48\cdot\eta_k^{-3}\cdot M^3$ distinct basis monomials, and each of them carries the coefficient $\eta_k^3/48$.
Therefore
\[
\left\|\langle\bar{\psi},\psi\rangle_{\square}\right\|_{(\mathsf{w}_1,\mathsf{w}_2)}=4\,M^3\cdot\mathsf{w}_1^2,
\qquad\text{and in particular}\qquad
\left\|\langle\bar{\psi},\psi\rangle_{\square}\right\|_{h_k}=4\,M^3\cdot e_k^{-\frac{1}{2}}.
\]

\emph{(b).}
In this part the generators $\psi$, $\chi$ carry the weights $h_k/2$, the generators $\Psi$ carry the weight
\[
\mathsf{w}_3:=C_T\cdot h_{k,1},
\]
norms without index are the norms $\|\cdot\|_{(h_k/2;\,\mathsf{w}_3)}$, and $s\in\mathrm{Pd}_k$.

\emph{The function $G^{(1)}$ for $k\ge1$.}
For $|t|=\mathrm{rad}_1$ the connection $\mathscr{T}+t\cdot\mathscr{Z}(s')$ belongs to $\tilde{G}_k$, by (q1).
The norm of an integral is bounded by the integral of the norm, and on the circle $|t|=\mathrm{rad}_1$ one has $|t\cdot(t-1)|\ge\mathrm{rad}_1\cdot(\mathrm{rad}_1-1)$.
With (q2) of Lemma~5.A, with (24.2) and with $2/\mathrm{rad}_1=4\,e_k^{\frac{3}{4}-\frac{9\varepsilon}{2}}$ this yields
\begin{multline*}
\left\|G^{(1)}(s,X)\right\|
\le4\,c_X\cdot c_{\gamma}\cdot C_{w}\cdot M^3\cdot C_{T}^2\cdot e_k^{-\frac{1}{2}}\cdot4\,e_k^{\frac{3}{4}-\frac{9\varepsilon}{2}}\cdot\exp\left(-(\kappa_{\mathrm{w}}-1)\cdot d_M(X)\right)\\
=\left[16\,c_X\cdot c_{\gamma}\cdot C_{w}\cdot C_{T}^2\cdot M^3\cdot e_k^{\frac{\varepsilon}{2}}\right]\cdot e_k^{\frac{1}{4}-5\varepsilon}\cdot\exp\left(-(\kappa_{\mathrm{w}}-1)\cdot d_M(X)\right)
\le e_k^{\frac{1}{4}-5\varepsilon}\cdot\exp\left(-\kappa\cdot d_M(X)\right).
\end{multline*}
The last inequality follows from the first part of (A3) and from $\kappa_{\mathrm{w}}-1=\kappa+\kappa_{0}+2\ge\kappa$, see (K4).

\emph{The function $G^{(1)}$ for $k=0$.}
Proposition~23.2\,(e) gives
\[
\left\|G^{(1)}(s,\square)\right\|_{h_0/2}\le12\,M^3\cdot e_k^{\frac{1}{2}-\frac{\varepsilon}{2}}\le e_k^{\frac{1}{4}-5\varepsilon},
\]
and $d_M(\square)=0$ by (t0).
Here and below (t0)--(t6), (CE1) and (CE2) denote the assertions on polymers of \cite[Proposition~5.58]{QEDI}, recalled in Section~0.6.
At $k=0$ the element $G^{(1)}$ contains no generator $\Psi$.

\emph{The function $G^{(2)}$ for $k\ge1$.}
Let $|t|=\mathrm{rad}_2$.
By (q1) each of the three norms of $\mathscr{Z}(s')$ is at most $\frac{1}{2}\,e_k^{-\frac{\varepsilon}{2}}$.
So (q3) and (A4) give
\[
\mathrm{rad}_2\cdot\left|J_k^{\#}\left(s'',\mathscr{T},\mathscr{Z}(s')\right)\right|_{1,\infty}\cdot\mathsf{w}_3
\le e_k^{-\frac{3}{4}+7\varepsilon}\cdot16\,C_{w}\cdot e_k^{\frac{3}{4}-\frac{9\varepsilon}{2}}\cdot C_{T}\cdot h_{k,1}
=16\,C_{w}\cdot C_{T}\cdot e_k^{\frac{3\varepsilon}{2}}\cdot h_{k,2}
\le\frac{1}{2}\,h_{k,2}.
\]
Put $v:=J_k^{\#}(s'',\mathscr{T},\mathscr{Z}(s'))\Psi$.
The homomorphism $\mathrm{sb}_{tv}$ sends $\psi(\xi)$ to $\psi(\xi)+t\cdot v(\xi)$ and $\chi(\boldsymbol{\zeta})$ to $\chi(\boldsymbol{\zeta})+t\cdot v(\boldsymbol{\zeta})$.
Apply (N2), that is Proposition~15\,(d): the row sums are $1$ for the identity on $\psi$, resp. on $\chi$, and at most $|t|\cdot|J_k^{\#}|_{1,\infty}$ for the part acting on $\Psi$.
Since
\[
\frac{h_{k,1}}{2}+\frac{h_{k,2}}{2}\le h_{k,1},\qquad\frac{h_{k,2}}{2}+\frac{h_{k,2}}{2}=h_{k,2},
\]
the norm $\|\cdot\|_{(h_k/2;\,\mathsf{w}_3)}$ of the integrand is at most $\left\|\mathscr{E}^{r}\left(X;A+t\cdot\mathscr{Z}(s'),\mathscr{C}\right)\right\|_{h_k}$, and the pair $\left(A+t\cdot\mathscr{Z}(s'),\mathscr{C}\right)$ lies in $\tilde{R}_k\times\mathrm{Sd}_k$ by (q1).
Now (a) and (24.2) give
\[
\left\|G^{(2)}(s,X)\right\|\le2\,e_k^{\frac{3}{4}-7\varepsilon}\cdot12\,e_k^{\frac{1}{4}-7\varepsilon}\cdot\exp\left(-\kappa\cdot d_M(X)\right)=24\,e_k^{1-14\varepsilon}\cdot\exp\left(-\kappa\cdot d_M(X)\right).
\]

\emph{The function $G^{(3)}$ for $k\ge1$.}
The integrand equals $-m'\cdot\mathrm{sb}_{tv}\left(\langle\bar{\psi},\psi\rangle_{\square}\right)$ with the same $v$.
The application of (N2) just made shows that its norm $\|\cdot\|_{(h_k/2;\,\mathsf{w}_3)}$ does not exceed $|m'|\cdot\left\|\langle\bar{\psi},\psi\rangle_{\square}\right\|_{h_k}$.
By (f0), (a) and (24.2),
\[
\left\|G^{(3)}(s,\square)\right\|\le8\,e_k^{\frac{3}{4}-8\varepsilon}\cdot4\,M^3\cdot e_k^{-\frac{1}{2}}\cdot2\,e_k^{\frac{3}{4}-7\varepsilon}=64\,M^3\cdot e_k^{1-15\varepsilon}.
\]

\emph{The sum.}
Since $\frac{1}{2}-4\varepsilon\ge0$, (A5) gives
\[
24\,e_k^{\frac{3}{4}-9\varepsilon}+64\,M^3\cdot e_k^{\frac{3}{4}-10\varepsilon}\le88\,M^3\cdot e_k^{\frac{1}{4}-6\varepsilon}\le1.
\]
A function $G(\cdot,\square)$ vanishes on the polymers that are not single $M$-cubes, and $d_M(\square)=0$.
Consequently $f^X:=\sum_{i\le3}G^{(i)}(\cdot,X)$ satisfies
\[
\left\|f^X(s)\right\|_{(h_k/2;\,\mathsf{w}_3)}\le2\,e_k^{\frac{1}{4}-5\varepsilon}\cdot\exp\left(-\kappa\cdot d_M(X)\right)\qquad(s\in\mathrm{Pd}_k).
\]
By Lemma~3.B the function $f^X$ is continuous on $\mathrm{Pd}_k$ and analytic in its interior, and $f^X_P=\mathrm{Tm}[f^X](X,P)$ vanishes unless $Y'(X,P)$ is connected.
Lemma~6.2 of Section~0.6 (\cite[Lemma~17.13]{QEDI}) therefore applies with $\mathrm{kq}=\kappa$ and $\mathrm{Bq}_0=2\,e_k^{\frac{1}{4}-5\varepsilon}$, and it gives
\[
\left\|\sum_{i\le3}E^{(i),*}(\mathbf{Y})\right\|_{(h_k/2;\,C_T\cdot h_{k,1})}\le8\,c_X\cdot K_0\cdot e_k^{\frac{1}{4}-5\varepsilon}\cdot\exp\left(-\mathrm{rt}_1\cdot d_M(\mathbf{Y})\right).
\]

\emph{The substitution ($k\ge1$).}
By (24.5) the element $\sum_iE^{(i),*}(\mathbf{Y})$ lies in $\mathrm{Gr}[\psi,\chi,\Psi|\mathbf{Y}]$.
The connection $\mathscr{T}$ lies in $\tilde{R}_k+\mathfrak{X}_k$, by (Sd2), and $C_T\cdot\left(h_{k,1}/2+h_{k,2}/2\right)\le C_T\cdot h_{k,1}$.
Hence Proposition~23.1\,(g) yields
\[
\left\|E^{(\le3),x}(\mathbf{Y})\right\|_{h_k/2}\le\left\|\sum_iE^{(i),*}(\mathbf{Y})\right\|_{(h_k/2;\,C_T\cdot h_{k,1})}.
\]
At $k=0$ no substitution occurs: $E^{(\le3),x}=E^{(1),*}$.

\emph{(c).}
By (A6), $C_W=28\,C_{w}^2\cdot c_{\gamma}\le\frac{3}{4}\,e_k^{-2\varepsilon}$.
Therefore
\[
\mathrm{rad}_4\cdot C_W\le\frac{3}{4}\,e_k^{-\frac{1}{4}+\varepsilon}=\frac{3}{4}\,h_{k,2},
\qquad
C_W\le\frac{3}{4}\,e_k^{\frac{1}{4}-3\varepsilon}\cdot h_{k,2}\le\frac{1}{4}\,h_{k,2},
\]
the last inequality because $3\,e_k^{\frac{1}{4}-3\varepsilon}\le1$ by (A5).
In this part the generators $\psi$, $\chi$ carry the weights $h_k/4$ and the generators $W$ the weight $\mathsf{w}_W$; moreover $s''\in\mathrm{Pd}_k$ and $v:=W_k^{\#}(s'',\mathscr{T})$.
In (N2), applied to $\mathrm{sb}_{tv}$, the row sums are $1$ on $\psi$, resp. on $\chi$, and at most $|t|\cdot\left|H_{k}^{\#}(s'',\mathscr{T})\,\tilde{\Gamma}_k(s'',\mathscr{T})\right|_{1,\infty}$ on $W$;
by (q4) the latter, multiplied by the weight $\mathsf{w}_W$, is at most $|t|\cdot C_W$.
For $|t|=1$ one gets
\[
\frac{h_{k,i}}{4}+C_W\le\frac{h_{k,i}}{4}+\frac{h_{k,2}}{4}\le\frac{h_{k,i}}{2},
\]
and for $|t|=\mathrm{rad}_4$
\[
\frac{h_{k,i}}{4}+\mathrm{rad}_4\cdot C_W\le\frac{h_{k,i}}{4}+\frac{3}{4}\,h_{k,2}\le h_{k,i}\qquad(i=1,2).
\]
With (a), (b), (f0), (24.2) and $2/\mathrm{rad}_4=2\,e_k^{\frac{1}{4}-3\varepsilon}$ it follows that
\begin{align*}
\left\|G^{(\le3)}(s'',X)\right\|_{(h_k/4;\,\mathsf{w}_W)}
&\le\left\|E^{(\le3),x}(X)\right\|_{h_k/2}\le8\,c_X\cdot K_0\cdot e_k^{\frac{1}{4}-5\varepsilon}\cdot\exp\left(-\mathrm{rt}_1\cdot d_M(X)\right),\\
\left\|G^{(4)}(s'',X)\right\|_{(h_k/4;\,\mathsf{w}_W)}
&\le2\,e_k^{\frac{1}{4}-3\varepsilon}\cdot12\,e_k^{\frac{1}{4}-7\varepsilon}\cdot\exp\left(-\kappa\cdot d_M(X)\right)=24\,e_k^{\frac{1}{2}-10\varepsilon}\cdot\exp\left(-\kappa\cdot d_M(X)\right),\\
\left\|G^{(5)}(s'',\square)\right\|_{(h_k/4;\,\mathsf{w}_W)}
&\le8\,e_k^{\frac{3}{4}-8\varepsilon}\cdot4\,M^3\cdot e_k^{-\frac{1}{2}}\cdot2\,e_k^{\frac{1}{4}-3\varepsilon}=64\,M^3\cdot e_k^{\frac{1}{2}-11\varepsilon}.
\end{align*}
By (A5),
\[
24\,e_k^{\frac{1}{4}-5\varepsilon}+64\,M^3\cdot e_k^{\frac{1}{4}-6\varepsilon}\le88\,M^3\cdot e_k^{\frac{1}{4}-6\varepsilon}\le1.
\]
Hence the sum $f^X$ of the three functions is bounded on $\mathrm{Pd}_k$ by
\[
\left(8\,c_X\cdot K_0+1\right)\cdot e_k^{\frac{1}{4}-5\varepsilon}\cdot\exp\left(-\mathrm{rt}_1\cdot d_M(X)\right)\le9\,c_X\cdot K_0\cdot e_k^{\frac{1}{4}-5\varepsilon}\cdot\exp\left(-\mathrm{rt}_1\cdot d_M(X)\right).
\]
By Lemma~3.C it fulfils the hypotheses of Lemma~6.2 with $\mathrm{kq}=\mathrm{rt}_1$, which is admissible because $\kappa_{0}\le\mathrm{rt}_1\le\kappa$, and with $\mathrm{Bq}_0=9\,c_X\cdot K_0\cdot e_k^{\frac{1}{4}-5\varepsilon}$.
Since
\[
4\,c_X\cdot K_0\cdot9\,c_X\cdot K_0\cdot e_k^{\frac{1}{4}-5\varepsilon}=\mathbf{H}_{0},
\qquad
\mathrm{rt}_1-\kappa_{0}-1=\mathrm{rt}_2,
\]
the bound of Lemma~6.2 is (c).

\emph{Level $0$.}
At $k=0$ one has $G^{(4)}=0$; the generators $W$ are indexed by $T^{\mathrm{f}}_0$ with the weight $\nu=1/48$, and $\mathsf{w}_W=\sqrt{48}\le7$; the factor $7$ is contained in $C_W$, see (q4).
Every inequality of this part is written with $\mathsf{w}_W$ and $C_W$; so all of them hold at $k=0$.

\emph{(d).}
This is the first identity of Lemma~6.2, applied to $f^X=G^{(i)}(\cdot,X)$, respectively to the sum of the three functions of the second localisation (Lemmas~3.B and~3.C).

\emph{(e).}
In Sections~3.1--3.8 the datum $(m,\mathscr{E})$ enters only through two of its properties: $m'$ is a complex number, and $\mathscr{E}^{r}$ belongs to $K^{x,C}_{k}$.
By (X1) and (X2) both properties hold for every $(m,\mathscr{E})\in\mathbb{C}\times K^{x,C}_{k}$.
The function $G^{(1)}$ does not contain $(m,\mathscr{E})$.
The functions $G^{(2)}$ and $G^{(4)}$ are linear in $\mathscr{E}^{r}=R\mathscr{E}$, which is linear in $\mathscr{E}$ by (X2).
The functions $G^{(3)}$ and $G^{(5)}$ are linear in $m'=m+m(\mathscr{E})$, which is linear in $(m,\mathscr{E})$ by (X1).
Finally, $\mathrm{Loc}$, $\mathrm{sub}_{\mathbf{Y},\mathscr{T}}$, $\mathrm{sb}_v$ and $\oint_{\mathrm{rad}}$ are linear maps.
\end{proof}

\subsubsection*{5.3. Mayer expansion, integration, logarithm}

\begin{named}{Lemma 5.C}
Let $\mathsf{w}=(\mathsf{w}_1,\mathsf{w}_2)$ with $\mathsf{w}_1,\mathsf{w}_2>0$.
For every $X\in D_k$ let $V(X)$ be an even element of $\mathrm{Gr}[\psi,\chi,W]$ that depends continuously on $\mathsf{Z}\in\mathrm{Box}_{k}$ and satisfies
\[
\left\|V(X)\right\|_{(\mathsf{w};\,\mathsf{w}_W)}\le\mathbf{H}_{0}\cdot\exp\left(-\mathrm{rt}_2\cdot d_M(X)\right).
\]
Assume that $V(X)$ contains $W(u)$ only for $u$ in $X$ (for $k=0$: only for $x$ over $X$), and that it depends on $\mathsf{Z}$ only through the coordinates $\mathsf{Z}(\Upsilon)$ with $\Upsilon$ contained in $X$.
Let $\mathrm{Km}$, $\mathrm{Hm}$ and $E^{\#}$ be formed from $V$ by (8.4), with $V$ in the place of $E^{(\le5),x}$; thus, for a union $X$ of $M$-cubes,
\[
\mathrm{Km}(X)=\sum_{\{X_i\}:\,\bigcup_iX_i=X}\ \prod_i\left(\exp V(X_i)-1\right),
\qquad
\mathrm{Km}(\emptyset)=1,
\qquad
\mathrm{Hm}(X)=\int\mathrm{d}\mu_I(W)\,\mathrm{d}\hat{\mu}(\mathsf{Z})\,\mathrm{Km}(X),
\]
the sum extending over the sets of distinct polymers with the union $X$, and, for $X\in D_k$,
\[
E^{\#}(X)=\sum_{n\ge1}\frac{1}{n!}\sum_{(X_1,\dots,X_n):\,\bigcup_iX_i=X}\mathrm{urs}^{T}(X_1,\dots,X_n)\cdot\prod_{i=1}^{n}\mathrm{Hm}(X_i).
\]
Then the series for $E^{\#}(X)$ converges absolutely in $\|\cdot\|_{\mathsf{w}}$, and
\[
\int\mathrm{d}\mu_I(W)\,\mathrm{d}\hat{\mu}(\mathsf{Z})\,\exp\left\{\sum_{X\in D_k}V(X)\right\}=\exp\left\{\sum_{X\in D_k}E^{\#}(X)\right\},
\qquad
\left\|E^{\#}(X)\right\|_{\mathsf{w}}\le c_{S}'\cdot\mathbf{H}_{0}\cdot\exp\left(-\mathrm{rt}_3\cdot d_M(X)\right).
\]
For every integer $\mathrm{nq}\ge0$ the norms $\|\cdot\|_{\mathsf{w}}$ of the terms with $n>\mathrm{nq}$ of the series for $E^{\#}(X)$ add up to at most $2^{-\mathrm{nq}}\cdot c_{S}'\cdot\mathbf{H}_{0}\cdot\exp\left(-\mathrm{rt}_3\cdot d_M(X)\right)$.
For an integer $\mathrm{nq}\ge1$ let
\[
\mathrm{pl}_{\mathrm{nq}}(V):=\sum_{1\le j\le\mathrm{nq}}\frac{V^j}{j!},
\]
and let $\mathrm{Km}_{\mathrm{nq}}$, $\mathrm{Hm}_{\mathrm{nq}}$ and $E^{\#}_{\mathrm{nq}}(X)$ be formed in the same way, with $\mathrm{pl}_{\mathrm{nq}}(V(X_i))$ in the place of $\exp V(X_i)-1$ and with the series (8.4) stopped at $n=\mathrm{nq}$.
Then $E^{\#}_{\mathrm{nq}}(X)$ is a polynomial in integrals of products of the elements $V(X_i)$, $X_i\subseteq X$;
\[
\left\|E^{\#}_{\mathrm{nq}}(X)\right\|_{\mathsf{w}}\le c_{S}'\cdot\mathbf{H}_{0}\cdot\exp\left(-\mathrm{rt}_3\cdot d_M(X)\right);
\]
and there are numbers $\mathrm{dlq}_{\mathrm{nq}}$, depending only on $\ell$, $e_k$ and the constants, with $\mathrm{dlq}_{\mathrm{nq}}\to0$ as $\mathrm{nq}\to\infty$ and
\[
\left\|E^{\#}(X)-E^{\#}_{\mathrm{nq}}(X)\right\|_{\mathsf{w}}\le\mathrm{dlq}_{\mathrm{nq}}
\]
for all $X\in D_k$ and all families $V$ with the properties named.
\end{named}

\begin{named}{Remark}
For $k\ge1$ one has $\mathsf{w}_W=1$, and for these levels Lemma~5.C is \cite[Lemma~17.17]{QEDI}.
Lemma~5.C holds for every $k\ge0$:
the generators $W$ carry the weight $\mathsf{w}_W$ in the norms, the integral over $W$ is estimated by (N3)$^{\mathrm{w}}$, that is Proposition~15\,(f),
and the factorisation of this integral over the unit points is (N4), that is Proposition~15\,(g).
\end{named}

\begin{proof}
Put
\[
\mathbf{H}_{1}:=7\,c_X\cdot K_0\cdot\exp\left(3\,\mathrm{rt}_5\right)\cdot\mathbf{H}_{0},
\qquad
z_S:=28\,\mathbf{H}_{1}\cdot\exp(\kappa_{0})\cdot K_0\cdot\exp\left(3\,\mathrm{rt}_6\right).
\]
One has $\mathrm{rt}_5\le\kappa$ and $3\,\mathrm{rt}_5+3\,\mathrm{rt}_6+\kappa_{0}=6\kappa-23\kappa_{0}-15\le6\kappa$.
With (A7) it follows that
\begin{equation}
\begin{aligned}
&\mathbf{H}_{0}\le1,\qquad
4\,\mathbf{H}_{0}\cdot K_0\cdot c_X\cdot\exp\left(3\,\mathrm{rt}_5\right)=144\,c_X^3\cdot K_0^3\cdot\exp\left(3\,\mathrm{rt}_5\right)\cdot e_k^{\frac{1}{4}-5\varepsilon}\le c_{S}\cdot e_k^{\frac{1}{4}-5\varepsilon}\le1,\\
&2\,z_S=14112\,c_X^3\cdot K_0^4\cdot\exp\left(3\,\mathrm{rt}_5+3\,\mathrm{rt}_6+\kappa_{0}\right)\cdot e_k^{\frac{1}{4}-5\varepsilon}\le c_{S}\cdot e_k^{\frac{1}{4}-5\varepsilon}\le1.
\end{aligned}
\tag{24.10}
\end{equation}

\emph{Step 1 (Mayer expansion).}
For $V=E^{(\le5),x}$ the next paragraph is the expansion asserted in Proposition~24.3\,(et).

By (N1), that is Proposition~15\,(b), even elements commute.
Expanding the product one obtains
\[
\exp\left\{\sum_XV(X)\right\}=\prod_X\left(1+\left(\exp V(X)-1\right)\right)=\sum_Z\mathrm{Km}(Z),
\]
where $Z$ ranges over the unions of $M$-cubes, the empty union included;
a subset $\{X_i\}$ of $D_k$ is counted in $\mathrm{Km}(Z)$ for $Z$ the union of the $X_i$.
The element $\mathrm{Km}(Z)$ is the product of the $\mathrm{Km}(Z_a)$ over the components $Z_a$ of $Z$:
every polymer of a set with the union $Z$ lies in one component, and so the sets with the union $Z$ are exactly the unions of sets with the unions $Z_a$.

Let $\mathbf{Y}$ be a polymer and $\{X_1,\dots,X_n\}$ a set of distinct polymers with the union $\mathbf{Y}$.
The graph on $\{1,\dots,n\}$ in which $i$ and $j$ are joined whenever $X_i$ and $X_j$ are incompatible is connected.
Indeed, otherwise the index set splits into two non-empty parts with no edge between them.
The unions $U_1$ and $U_2$ of the polymers of the two parts cover the connected set $\mathbf{Y}$;
so there are $M$-cubes $\square_1\subseteq U_1$ and $\square_2\subseteq U_2$ that coincide or have a common face,
and then a polymer $X_i$ containing $\square_1$ and a polymer $X_j$ containing $\square_2$ are incompatible, although $i$ and $j$ lie in different parts.
By (t4),
\[
d_M(\mathbf{Y})\le\sum_id_M(X_i)+3\,(n-1).
\]
Since $\exp(a)-1\le2a$ for $a\in[0,1]$, (N1) and $\mathbf{H}_{0}\le1$ give $\left\|\exp V(X)-1\right\|\le\exp\left(\|V(X)\|\right)-1\le2\,\|V(X)\|$.
As $\mathrm{rt}_2=\mathrm{rt}_5+\kappa_{0}$, one finds
\begin{multline*}
\prod_i\left\|\exp V(X_i)-1\right\|_{(\mathsf{w};\,\mathsf{w}_W)}
\le\left(2\,\mathbf{H}_{0}\right)^n\cdot\prod_i\exp\left(-\mathrm{rt}_5\cdot d_M(X_i)\right)\cdot\exp\left(-\kappa_{0}\cdot d_M(X_i)\right)\\
\le\exp\left(-\mathrm{rt}_5\cdot d_M(\mathbf{Y})\right)\cdot\left(2\,\mathbf{H}_{0}\cdot\exp\left(3\,\mathrm{rt}_5\right)\right)^n\cdot\prod_i\exp\left(-\kappa_{0}\cdot d_M(X_i)\right).
\end{multline*}
To a set of $n$ distinct polymers there correspond $n!$ ordered $n$-tuples.
Together with (t2) this gives
\begin{multline*}
\left\|\mathrm{Km}(\mathbf{Y})\right\|_{(\mathsf{w};\,\mathsf{w}_W)}
\le\exp\left(-\mathrm{rt}_5\cdot d_M(\mathbf{Y})\right)\cdot\sum_{n\ge1}\frac{1}{n!}\left(2\,\mathbf{H}_{0}\cdot\exp\left(3\,\mathrm{rt}_5\right)\cdot\sum_{X\subseteq\mathbf{Y}}\exp\left(-\kappa_{0}\cdot d_M(X)\right)\right)^n\\
\le\exp\left(-\mathrm{rt}_5\cdot d_M(\mathbf{Y})\right)\cdot\mathrm{xq}\cdot\exp(\mathrm{xq}),
\qquad
\mathrm{xq}:=2\,\mathbf{H}_{0}\cdot K_0\cdot\exp\left(3\,\mathrm{rt}_5\right)\cdot|\mathbf{Y}|_M.
\end{multline*}
By (24.10) and (t1),
\[
\mathrm{xq}\le\frac{|\mathbf{Y}|_M}{2\,c_X}\le\frac{d_M(\mathbf{Y})+1}{2},
\]
and by (t1) and (t6),
\[
\mathrm{xq}\le4\,\mathbf{H}_{0}\cdot K_0\cdot c_X\cdot\exp\left(3\,\mathrm{rt}_5\right)\cdot\exp\left(\frac{d_M(\mathbf{Y})}{2}\right).
\]
Consequently
\[
\mathrm{xq}\cdot\exp(\mathrm{xq})\le4\exp\left(\frac{1}{2}\right)\cdot c_X\cdot K_0\cdot\exp\left(3\,\mathrm{rt}_5\right)\cdot\mathbf{H}_{0}\cdot\exp\left(d_M(\mathbf{Y})\right).
\]
Since $4\exp\left(\frac{1}{2}\right)\le7$ and $\mathrm{rt}_5-1=\mathrm{rt}_4$, it follows that
\begin{equation}
\left\|\mathrm{Km}(\mathbf{Y})\right\|_{(\mathsf{w};\,\mathsf{w}_W)}\le\mathbf{H}_{1}\cdot\exp\left(-\mathrm{rt}_4\cdot d_M(\mathbf{Y})\right)=:\mathrm{mj}(\mathbf{Y}).
\tag{24.11}
\end{equation}

\emph{Step 2 (integration).}
The element $\mathrm{Km}(\mathbf{Y})$ is continuous in $\mathsf{Z}$, and $\hat{\mu}$ is a probability measure on the compact set $\mathrm{Box}_{k}$.
By (N3)$^{\mathrm{w}}$,
\[
\left\|\int\mathrm{d}\mu_I(W)\,F\right\|_{\mathsf{w}}\le\|F\|_{(\mathsf{w};\,\mathsf{w}_W)},
\]
where $\mathsf{w}_W=\max\nu^{-\frac{1}{2}}$ equals $1$ for $k\ge1$ and $\sqrt{48}$ for $k=0$; at $k=0$ each integrated generator $W_{\beta}(x)$, $\bar{W}_{\beta}(x)$ costs the factor $\sqrt{48}$.
The norm of the integral over $\mathsf{Z}$ is bounded by the integral of the norms.
Hence $\left\|\mathrm{Hm}(\mathbf{Y})\right\|_{\mathsf{w}}\le\mathrm{mj}(\mathbf{Y})$.

\emph{The elements $\mathrm{Hm}(\mathbf{Y})$ and $\mathrm{Hm}_{\mathrm{nq}}(\mathbf{Y})$ are even elements of $\mathrm{Gr}[\psi,\chi]$.}
As $V(X)$ is even, so are $\exp V(X)-1$ (Section~3.8) and the polynomial $\mathrm{pl}_{\mathrm{nq}}(V(X))$, and hence also $\mathrm{Km}(\mathbf{Y})$ and $\mathrm{Km}_{\mathrm{nq}}(\mathbf{Y})$.
Let $F=\sum_Jc_J\,W(J)$, with $c_J\in\mathrm{Gr}[\psi,\chi]$, be even; then $\mathrm{Pty}\,c_J=(-1)^{\#|J|}\,c_J$.
By the identity in the first sentence of (N3)$^{\mathrm{w}}$, that is Proposition~15\,(f), one has $\int\mathrm{d}\mu_I(W)\,F=\sum_Jc_J\,m(J)$.
By (mo), that is Proposition~15\,(e), $m(J)=0$ unless $J$ is paired; for paired $J$ the number $\#|J|$ is even, and so $c_J$ is even.
Thus $\int\mathrm{d}\mu_I(W)$ maps even elements to even elements of $\mathrm{Gr}[\psi,\chi]$.
The integral over $\mathsf{Z}$ is a limit of linear combinations of values, and the even elements form a closed linear subspace.

Let $\mathbf{Y}_1$ and $\mathbf{Y}_2$ be compatible polymers; then they have no $M$-cube in common.
Hence $\mathrm{Km}(\mathbf{Y}_1)$ and $\mathrm{Km}(\mathbf{Y}_2)$ contain $W$ at disjoint sets of unit points (for $k=0$: of fermion sites of $T^{\mathrm{f}}_0$), and they depend on disjoint sets of coordinates $\mathsf{Z}(\Upsilon)$;
this follows from the last sentence of Proposition~24.3\,(ze) (Section~3.8), since the polymers $X_i$ that occur in $\mathrm{Km}(\mathbf{Y}_a)$ lie in $\mathbf{Y}_a$.
The measure $\hat{\mu}$ is a product measure over the coordinates (Section~3.11).
By (N4), that is Proposition~15\,(g), and by Fubini's theorem, the integral of a product of elements $\mathrm{Km}(\mathbf{Y}_i)$ over pairwise compatible polymers $\mathbf{Y}_i$ is the product of the $\mathrm{Hm}(\mathbf{Y}_i)$.
The components of a union $Z$ of $M$-cubes are pairwise compatible polymers;
conversely, a set of pairwise compatible polymers is the set of the components of its union.
Together with Step~1 this yields
\begin{equation}
\int\mathrm{d}\mu_I(W)\,\mathrm{d}\hat{\mu}(\mathsf{Z})\,\exp\left\{\sum_XV(X)\right\}=\sum_{\{\mathbf{Y}_i\}\ \text{pairwise compatible}}\ \prod_i\mathrm{Hm}(\mathbf{Y}_i).
\tag{24.12}
\end{equation}

\emph{Step 3 (the series).}
For a polymer $\mathbf{Y}$ put
\[
\mathrm{Sg}_n(\mathbf{Y}):=\sum_{(\mathbf{Y}_1,\dots,\mathbf{Y}_n):\,\bigcup_i\mathbf{Y}_i=\mathbf{Y}}\left|\mathrm{urs}^{T}(\mathbf{Y}_1,\dots,\mathbf{Y}_n)\right|\cdot\prod_i\mathrm{mj}(\mathbf{Y}_i).
\]
Then $\mathrm{Sg}_1(\mathbf{Y})=\mathrm{mj}(\mathbf{Y})$.
Let $n\ge2$.
By (CE2), $\mathrm{Sg}_n(\mathbf{Y})$ does not exceed the sum of $\prod_i\mathrm{mj}(\mathbf{Y}_i)$ over the trees $\mathrm{tre}$ on $\{1,\dots,n\}$ and over the $n$-tuples with the union $\mathbf{Y}$ in which $\mathbf{Y}_i$ and $\mathbf{Y}_j$ are incompatible for every edge $\{i,j\}$ of $\mathrm{tre}$.
For such an $n$-tuple, (t4) and $\mathrm{rt}_4=\mathrm{rt}_6+2\kappa_{0}$ give
\[
\prod_i\exp\left(-\mathrm{rt}_4\cdot d_M(\mathbf{Y}_i)\right)\le\exp\left(-\mathrm{rt}_6\cdot d_M(\mathbf{Y})\right)\cdot\exp\left(3\,\mathrm{rt}_6\cdot(n-1)\right)\cdot\prod_i\exp\left(-2\kappa_{0}\cdot d_M(\mathbf{Y}_i)\right).
\]
Now the condition on the union is dropped, and only $\mathbf{Y}_i\subseteq\mathbf{Y}$ is kept.
Root $\mathrm{tre}$ at the vertex $1$, and let $\mathrm{kd}_j$ be the number of the children of $j$.
A polymer that is incompatible with a polymer $\mathbf{Y}'$ contains an $M$-cube of $\mathbf{Y}'$ or an $M$-cube that has a face in common with an $M$-cube of $\mathbf{Y}'$, because the union is connected;
that is, it contains one of at most $7\,|\mathbf{Y}'|_M$ cubes.
By $a^{\mathrm{kd}}\le\mathrm{kd}!\cdot\exp(a)$, by $|Z|_M\le c_X\cdot\left(d_M(Z)+1\right)\le\kappa_{0}\cdot\left(d_M(Z)+1\right)$ (see (t1) and (K2)), and by (t2),
\begin{gather*}
\sum_{\mathbf{Y}_j\ \text{incompatible with}\ \mathbf{Y}'}\exp\left(-2\kappa_{0}\cdot d_M(\mathbf{Y}_j)\right)\cdot|\mathbf{Y}_j|_M^{\mathrm{kd}_j}\le7\,\mathrm{kd}_j!\cdot\exp(\kappa_{0})\cdot K_0\cdot|\mathbf{Y}'|_M,\\
\sum_{\mathbf{Y}_1\subseteq\mathbf{Y}}\exp\left(-2\kappa_{0}\cdot d_M(\mathbf{Y}_1)\right)\cdot|\mathbf{Y}_1|_M^{\mathrm{kd}_1}\le\mathrm{kd}_1!\cdot\exp(\kappa_{0})\cdot K_0\cdot|\mathbf{Y}|_M.
\end{gather*}
Sum over the $\mathbf{Y}_j$ from the leaves towards the root; when $\mathbf{Y}_j$ is summed, its $\mathrm{kd}_j$ children have contributed the factor $|\mathbf{Y}_j|_M^{\mathrm{kd}_j}$.
In this way the sum of $\prod_i\exp\left(-2\kappa_{0}\cdot d_M(\mathbf{Y}_i)\right)$ is bounded by
\[
\left(7\exp(\kappa_{0})\cdot K_0\right)^n\cdot|\mathbf{Y}|_M\cdot\prod_j\mathrm{kd}_j!.
\]
The vertices of $\mathrm{tre}$ have the degrees $\mathrm{kd}_1$ and $\mathrm{kd}_j+1$ ($j\ge2$).
By Cayley's formula the number of the trees with prescribed $(\mathrm{kd}_j)$, where $\sum_j\mathrm{kd}_j=n-1$ and $\mathrm{kd}_1\ge1$, equals
\[
\frac{(n-2)!}{(\mathrm{kd}_1-1)!\cdot\prod_{j\ge2}\mathrm{kd}_j!}.
\]
Therefore
\[
\sum_{\mathrm{tre}}\prod_j\mathrm{kd}_j!=(n-2)!\cdot\sum_{(\mathrm{kd}_j)}\mathrm{kd}_1\le(n-2)!\cdot(n-1)\cdot\binom{2n-2}{n-1}\le(n-1)!\cdot4^{n-1}.
\]
Collecting these bounds one obtains, also for $n=1$,
\begin{equation}
\frac{\mathrm{Sg}_n(\mathbf{Y})}{n!}\le7\,\mathbf{H}_{1}\cdot\exp(\kappa_{0})\cdot K_0\cdot|\mathbf{Y}|_M\cdot\exp\left(-\mathrm{rt}_6\cdot d_M(\mathbf{Y})\right)\cdot z_S^{n-1}.
\tag{24.13}
\end{equation}
By (24.10), $z_S\le\frac{1}{2}$.
With (t1), (t6) and $\mathrm{rt}_6-1=\mathrm{rt}_3$ it follows that
\begin{multline*}
\sum_{n\ge1}\frac{\mathrm{Sg}_n(\mathbf{Y})}{n!}\le14\,c_X\cdot\exp(\kappa_{0})\cdot K_0\cdot\mathbf{H}_{1}\cdot\exp\left(-\mathrm{rt}_3\cdot d_M(\mathbf{Y})\right)\\
=98\,c_X^2\cdot K_0^2\cdot\exp\left(\kappa_{0}+3\,\mathrm{rt}_5\right)\cdot\mathbf{H}_{0}\cdot\exp\left(-\mathrm{rt}_3\cdot d_M(\mathbf{Y})\right)
\le c_{S}'\cdot\mathbf{H}_{0}\cdot\exp\left(-\mathrm{rt}_3\cdot d_M(\mathbf{Y})\right),
\end{multline*}
because $\kappa_{0}+3\,\mathrm{rt}_5=3\kappa-8\kappa_{0}-6$, $c_{S}'=14\,c_X^2\cdot K_0^2\cdot\exp(3\kappa)$ and $98\exp\left(-8\kappa_{0}-6\right)\le14$.
For $\mathrm{nq}\ge0$ one has $\sum_{n>\mathrm{nq}}z_S^{n-1}\le2\cdot2^{-\mathrm{nq}}$, and hence
\[
\sum_{n>\mathrm{nq}}\frac{\mathrm{Sg}_n(\mathbf{Y})}{n!}\le2^{-\mathrm{nq}}\cdot c_{S}'\cdot\mathbf{H}_{0}\cdot\exp\left(-\mathrm{rt}_3\cdot d_M(\mathbf{Y})\right).
\]

\emph{Step 4 (the logarithm).}
The set $D_k$ is finite.
By (N1) the even part of $\mathrm{Gr}[\psi,\chi]$, normed by $\|\cdot\|_{\mathsf{w}}$, is a commutative Banach algebra with unit.
Let $\mathrm{Ph}(z)$ be the formal power series in the exponent of (CE1), regarded as the family of its terms
\[
\mathrm{trm}=\frac{1}{n!}\,\mathrm{urs}^{T}(\mathbf{Y}_1,\dots,\mathbf{Y}_n)\cdot\prod_iz_{\mathbf{Y}_i},
\]
one term for each $n$ and each $n$-tuple.
Put $z_{\mathbf{Y}}:=\mathrm{Hm}(\mathbf{Y})$.
By Steps~2 and~3 the norms of the terms $\mathrm{trm}(\mathrm{Hm})$ add up to at most
\[
\mathrm{Sgt}:=\sum_{\mathbf{Y}\in D_k}\sum_n\frac{\mathrm{Sg}_n(\mathbf{Y})}{n!}<\infty.
\]
Consider the family of the elements $\frac{1}{j!}\,\mathrm{trm}_1(\mathrm{Hm})\cdots\mathrm{trm}_j(\mathrm{Hm})$, indexed by $j\ge0$ and by the $j$-tuples $(\mathrm{trm}_1,\dots,\mathrm{trm}_j)$ of terms.
The sum of their norms is at most $\exp(\mathrm{Sgt})$; so the family is absolutely summable, and it may be summed in groups.
Grouped according to $j$, its sum is $\exp\left(\mathrm{Ph}(\mathrm{Hm})\right)$, where $\mathrm{Ph}(\mathrm{Hm}):=\sum_{\mathrm{trm}}\mathrm{trm}(\mathrm{Hm})$.
Grouped according to the monomials $\prod_{\mathbf{Y}}z_{\mathbf{Y}}^{i_{\mathbf{Y}}}$, the numerical coefficients within one group add up to the coefficient of that monomial in the formal series $\exp\mathrm{Ph}(z)$;
by (CE1) this coefficient is $1$ if the monomial is the product of the variables of a set of pairwise compatible polymers, and $0$ otherwise.
So $\exp\left(\mathrm{Ph}(\mathrm{Hm})\right)$ is the right side of (24.12).
In $\mathrm{Ph}(\mathrm{Hm})$ collect the terms with a given union:
$\mathrm{urs}^{T}(\mathbf{Y}_1,\dots,\mathbf{Y}_n)$ vanishes unless the incompatibility graph is connected, and then the union is a polymer by (t4).
Hence $\mathrm{Ph}(\mathrm{Hm})=\sum_{X\in D_k}E^{\#}(X)$.
The absolute convergence, the bound on $E^{\#}(X)$ and the bound on the terms with $n>\mathrm{nq}$ are contained in Step~3.

\emph{Step 5 (the approximants).}
One has $\left\|\mathrm{pl}_{\mathrm{nq}}(V(X))\right\|\le\exp\left(\|V(X)\|\right)-1$.
So Steps~1--3 hold for $\mathrm{Km}_{\mathrm{nq}}$, $\mathrm{Hm}_{\mathrm{nq}}$ and $E^{\#}_{\mathrm{nq}}$ with the same constants; in particular
\[
\left\|\mathrm{Hm}_{\mathrm{nq}}(\mathbf{Y})\right\|_{\mathsf{w}}\le\mathrm{mj}(\mathbf{Y}),
\qquad
\left\|E^{\#}_{\mathrm{nq}}(X)\right\|_{\mathsf{w}}\le c_{S}'\cdot\mathbf{H}_{0}\cdot\exp\left(-\mathrm{rt}_3\cdot d_M(X)\right).
\]
By definition, $\mathrm{Hm}_{\mathrm{nq}}(\mathbf{Y})$ is a finite sum of integrals of products of the $V(X_i)$, $X_i\subseteq\mathbf{Y}$, and $E^{\#}_{\mathrm{nq}}(X)$ is a polynomial in the $\mathrm{Hm}_{\mathrm{nq}}(X_i)$, $X_i\subseteq X$.
Since $\|V(X)\|\le1$,
\[
\left\|\exp V(X)-1-\mathrm{pl}_{\mathrm{nq}}(V(X))\right\|\le\sum_{j>\mathrm{nq}}\frac{1}{j!}\le\frac{\exp(1)}{(\mathrm{nq}+1)!}.
\]
The difference $\mathrm{Km}(\mathbf{Y})-\mathrm{Km}_{\mathrm{nq}}(\mathbf{Y})$ is a finite sum of differences of products, and
\[
a_1\cdots a_n-b_1\cdots b_n=\sum_ja_1\cdots a_{j-1}\cdot(a_j-b_j)\cdot b_{j+1}\cdots b_n,
\]
where all factors have norm at most $\exp(1)-1\le2$.
There are at most $2^{\#|D_k|}$ sets $\{X_i\}$, each of them with $n\le\#|D_k|$ members; hence
\[
\left\|\mathrm{Km}(\mathbf{Y})-\mathrm{Km}_{\mathrm{nq}}(\mathbf{Y})\right\|_{(\mathsf{w};\,\mathsf{w}_W)}\le2^{\#|D_k|}\cdot\#|D_k|\cdot2^{\#|D_k|}\cdot\frac{\exp(1)}{(\mathrm{nq}+1)!},
\]
and by Step~2 the same bound holds for $\left\|\mathrm{Hm}(\mathbf{Y})-\mathrm{Hm}_{\mathrm{nq}}(\mathbf{Y})\right\|_{\mathsf{w}}$.
The numbers $\mathrm{mj}(\mathbf{Y})$ are positive and independent of $V$.
Put
\[
\delta_{\mathrm{nq}}:=\frac{4^{\#|D_k|}\cdot\#|D_k|\cdot\exp(1)}{(\mathrm{nq}+1)!\cdot\min_{\mathbf{Y}\in D_k}\mathrm{mj}(\mathbf{Y})}.
\]
Then $\left\|\mathrm{Hm}(\mathbf{Y})-\mathrm{Hm}_{\mathrm{nq}}(\mathbf{Y})\right\|_{\mathsf{w}}\le\delta_{\mathrm{nq}}\cdot\mathrm{mj}(\mathbf{Y})$ for all $\mathbf{Y}\in D_k$ and all $V$, and $\delta_{\mathrm{nq}}\to0$.
The number $\delta_{\mathrm{nq}}$ depends only on $\ell$, $e_k$ and the constants.
Indeed, an $M$-cube is determined by its centre, a point of $M\mathbb{Z}^3/\ell\mathbb{Z}^3$.
The dilation by $1/M$ maps the closed $M$-cubes of $\mathbb{R}^3/\ell\mathbb{Z}^3$ onto the closed unit cubes of $\mathbb{R}^3/(\ell/M)\mathbb{Z}^3$ with centres in $\mathbb{Z}^3/(\ell/M)\mathbb{Z}^3$, common faces onto common faces and nets onto nets, and it divides lengths by $M$;
by Section~I.1.1, $M\cdot d_M(\mathbf{Y})$ is the infimum of the lengths of the nets that meet every closed $M$-cube of $\mathbf{Y}$.
So $\#|D_k|$ and the numbers $d_M(\mathbf{Y})$, $\mathbf{Y}\in D_k$, are determined by $\ell/M$;
and $\mathrm{mj}(\mathbf{Y})=\mathbf{H}_{1}\cdot\exp\left(-\mathrm{rt}_4\cdot d_M(\mathbf{Y})\right)$ contains, besides $d_M(\mathbf{Y})$, only $e_k$ (through $\mathbf{H}_{0}$) and the constants.
In a term of (8.4) with $n\le\mathrm{nq}$, replacing the $\mathrm{Hm}(X_i)$ by the $\mathrm{Hm}_{\mathrm{nq}}(X_i)$ changes the product by at most $n\cdot\delta_{\mathrm{nq}}\cdot\prod_i\mathrm{mj}(X_i)$.
By (24.13), and since $\sum_nn\cdot2^{-(n-1)}=4$, these changes add up to at most
\[
\delta_{\mathrm{nq}}\cdot\sum_nn\cdot\frac{\mathrm{Sg}_n(X)}{n!}\le28\,\delta_{\mathrm{nq}}\cdot\mathbf{H}_{1}\cdot\exp(\kappa_{0})\cdot K_0\cdot|X|_M.
\]
The terms with $n>\mathrm{nq}$ add up to at most $2^{-\mathrm{nq}}\cdot c_{S}'\cdot\mathbf{H}_{0}$.
As $|X|_M\le(\ell/M)^3$, the number
\[
\mathrm{dlq}_{\mathrm{nq}}:=28\,\delta_{\mathrm{nq}}\cdot\mathbf{H}_{1}\cdot\exp(\kappa_{0})\cdot K_0\cdot\left(\frac{\ell}{M}\right)^3+2^{-\mathrm{nq}}\cdot c_{S}'\cdot\mathbf{H}_{0}
\]
has the property asserted.
\end{proof}

\subsubsection*{5.4. The first sentence of (a); the first sentence of (b)}

\begin{proof}[Proof of the first sentence of Proposition~24.1\,(a), with its display, and of the first sentence of Proposition~24.1\,(b)]
Let $(m,\mathscr{E})\in\mathrm{Ball}^{x}_{k}(4)$, with $\mathscr{E}=0$ when $k=0$, and let $A\in\tilde{R}_k^{-}$ and $\mathscr{C}\in\mathrm{Sd}_k$.

\emph{All displays are defined, in the order of Definition~24.D.}
In (8.1)$^{\mathrm{ch}}$ the primitives are defined on $\mathrm{Pd}_k$, respectively on $\mathrm{Pd}_k\times\tilde{G}_k$, and by (q1) the connections $\mathscr{T}$ and $\mathscr{T}+\mathscr{Z}(s')$ at which they are evaluated lie in $\tilde{G}_k$.
The clauses (q1) and (q5) hold by Sections~2.2 and~3.1.
The functions $G^{(1)}$, $G^{(2)}$, $G^{(3)}$ are defined by Lemma~3.B\,(reg), and the first localisation is defined by Section~3.6.
The display (8.3a)$^{\mathrm{ch}}$ has a meaning by Section~3.6, since $E^{(i),*}(\mathbf{Y})$ lies in $\mathrm{Gr}[\psi,\chi,\Psi|\mathbf{Y}]$.
The functions $G^{(\le 3)}$, $G^{(4)}$, $G^{(5)}$ are defined by Lemma~3.C\,(reg); the second localisation and the elements $\mathrm{Km}$, $\mathrm{Hm}$ are defined by Section~3.8.
The convergence of the series in (8.4) is proved below.
Consider (8.6)$^{\mathrm{ch}}$, and let $(\mathscr{A},\mathscr{C})\in\tilde{R}_{k+1}\times\mathrm{Sd}_{k+1}$.
By DF2\,(a), (b), that is \cite[Lemma~17.5(a),(b)]{QEDI}, and by (Sd4), the pair $(\mathscr{A}_L,\mathscr{C}_L)$ lies in $\tilde{R}_k^{-}\times\mathrm{Sd}_k$; hence $E^{*,x}_k(m,\mathscr{E})(Y;\mathscr{A},\mathscr{C})$ is a finite sum of elements that are defined.
The family $LR\mathscr{E}$ is defined by Proposition~18.5\,(a) (in the case (o), by Proposition~48.5\,(a)); the numbers $m(\mathscr{E})$ and $\varepsilon(\mathscr{E})$ are defined by (X1), and the number $\delta\varepsilon^{0}_{k}$ by Section~I.5.7, that is by \cite[Section~5.2.5]{QEDI}.
The connections $\mathscr{A}_L+\mathscr{C}_L$ and $\mathscr{C}_L$ lie in $\tilde{G}_k$ by (Sd2), because $0\in\tilde{R}_k$; and $0\in\tilde{G}_k$ by DF1, that is \cite[Lemma~17.4(a)]{QEDI}.
Together with $B$, the connections $-B$, $\overline{B}$ and $-\overline{B}$ lie in $\tilde{G}_k$ by DF3, that is \cite[Lemma~17.6]{QEDI}; and by Proposition~13\,(d0) the integrals in the displays of $\acute{E}_k(X,\cdot\,)$ converge absolutely at every point of $\tilde{G}_k$.
Consequently, for every $X\in D_k$ and every $B\in\tilde{G}_k$ the number $\acute{E}_k(X,B)$ is defined, and so is $E^{d}_k(X,B)$, the mean of the four numbers
\[
\acute{E}_k(X,B),\qquad\acute{E}_k(X,-B),\qquad\overline{\acute{E}_k(X,\overline{B})},\qquad\overline{\acute{E}_k(X,-\overline{B})}
\]
(Section~I.5.7).
It follows that the numbers
\[
{\delta E^{\mathrm{det}}}_k(X;\mathscr{A}_L,\mathscr{C}_L)=E^{d}_k(X,\mathscr{A}_L+\mathscr{C}_L)-E^{d}_k(X,\mathscr{C}_L),\qquad E_k^{\mathrm{det}}(X,\mathscr{C}_L)=E^{d}_k(X,\mathscr{C}_L)-E^{d}_k(X,0)
\]
are defined, and that
\[
\left(L\,{\delta E^{\mathrm{det}}}_k\right)(Y;\mathscr{A},\mathscr{C})=\sum_{X:\,\bar{X}^{L}=LY}{\delta E^{\mathrm{det}}}_k(X;\mathscr{A}_L,\mathscr{C}_L)
\]
is a finite sum of such numbers.
By Proposition~13\,(d5) the family $L\,{\delta E^{\mathrm{det}}}_k$ lies in $K^{x}_{k+1}$ (in the case (o): the family $L\,E_k^{\mathrm{det}}$ lies in $K_{k+1}$).
In the case (x), the sum in
\[
\mathfrak{P}_{k+1}(Y,\mathscr{C})=\sum_{X:\,\bar{X}^{L}=LY}\left(\mathfrak{P}_k(X,\mathscr{C}_L)+E_k^{\mathrm{det}}(X,\mathscr{C}_L)\right)
\]
is finite, and by (Sd4) the connection $\mathscr{C}_L$ lies in $\mathrm{Sd}_k=\mathfrak{X}_k$, which is the domain of the functions $\mathfrak{P}_k(X,\cdot\,)$ of the given family (Definition~24.D\,(D6)).
In the case (o), the family $\mathfrak{P}_k$ and the entry $\mathfrak{P}_{k+1}$ of (8.6)$^{\mathrm{ch}}$ are void (Definition~24.D\,(D7)).

\emph{Convergence and bound.}
By Proposition~5.B\,(c) and by (24.7), the elements $V(X):=E^{(\le 5),x}(X)$ satisfy the hypotheses of Lemma~5.C with $\mathsf{w}=h_k/4$.
Indeed, they are even elements of $\mathrm{Gr}[\psi,\chi,W]$ and continuous in $\mathsf{Z}$; they contain $W(u)$ only for $u$ in $X$; they depend on $\mathsf{Z}$ only through the coordinates $\mathsf{Z}(\Upsilon)$ with $\Upsilon$ contained in $X$; and
\[
\left\|V(X)\right\|_{(h_k/4;\,\mathsf{w}_W)}\le\mathbf{H}_{0}\cdot\exp\left(-\mathrm{rt}_2\cdot d_M(X)\right).
\]
Lemma~5.C therefore shows that the series (8.4) converge absolutely in $\|\cdot\|_{h_k/4}$, and that
\begin{equation}
\begin{split}
\left\|E^{\#,x}_k(m,\mathscr{E})(X;A,\mathscr{C})\right\|_{h_k/4}&\le c_{S}'\cdot\mathbf{H}_{0}\cdot\exp\left(-\mathrm{rt}_3\cdot d_M(X)\right)\\
&=14\,c_{X}^2\cdot K_{0}^2\cdot\exp(3\kappa)\cdot\mathbf{H}_{0}\cdot\exp\left(-(\kappa-5\kappa_{0}-4)\cdot d_M(X)\right),
\end{split}
\tag{24.14}
\end{equation}
\begin{equation}
\int\mathrm{d}\mu_I(W)\,\mathrm{d}\hat{\mu}(\mathsf{Z})\,\exp\left\{\sum_{X}E^{(\le 5),x}(X)\right\}=\exp\left\{\sum_{X}E^{\#,x}_k(m,\mathscr{E})(X;A,\mathscr{C})\right\}\qquad\text{in }\mathrm{Gr}[\psi,\chi].
\tag{24.15}
\end{equation}
In the case (x) this is the first sentence of Proposition~24.1\,(a), with its display.
In the case (o) (Section~0.5: data $(m,E)\in\mathrm{Ball}_{k}(4)$, with $E=0$ if $k=0$; $A\in\tilde{R}_k^{-}$; the displays of the reading (o) of Definition~24.D\,(D7)) it is the first sentence of Proposition~24.1\,(b).
\end{proof}

\subsection*{6. Proof of Proposition 24.3: (ze), (al), (ga4) for $E^{\#,x}_k(m,\mathscr{E})(\mathbf{Y})$; the last sentences}

\begin{proof}
By Section~5.4 the series (8.4) for $E^{\#,x}_k(m,\mathscr{E})(\mathbf{Y})$ converges in norm at every $(A,\mathscr{C})\in\tilde{R}_k^{-}\times\mathrm{Sd}_k$.
Its terms are real multiples of products $\mathrm{Hm}(X_1)\cdots\mathrm{Hm}(X_n)$ in which every $X_i$ is contained in $\mathbf{Y}$, by (al).
By (24.9) and (Lo1)--(Lo3), each term is an even element of $\mathrm{Gr}[\psi,\chi]$ which contains $\psi(\xi)$ only with $\xi$ over $\mathbf{Y}$ and $\chi(\boldsymbol{\zeta})$ only with $\boldsymbol{\zeta}$ admissible for $\mathbf{Y}$, and which takes the same value at any two pairs of $\tilde{R}_k^{-}\times\mathrm{Sd}_k$ that coincide on the bonds meeting $\mathbf{Y}$.
By (Lo2) the sum of the series inherits these properties.
This completes the proof of (ze).

\emph{The fourth sentence of (al) and the last sentence of (ga4), for $E^{\#,x}_k(m,\mathscr{E})(\mathbf{Y})$.}
Let $(m_1,\mathscr{E}_1)$, $(m_2,\mathscr{E}_2)$ be two data as in (al); both lie in $\mathrm{Ball}^{x}_{k}(4)$.
Respectively, let $\mathrm{Sy}^1$, $\mathrm{Sy}^2$ be two systems as in (ga4); both have the properties of (ST1)$^{\mathrm{ch}}$.
Section~5.4, applied to each of the two data, respectively with each of the two systems, shows that both series (8.4) for the polymer $\mathbf{Y}$ converge at every $(A,\mathscr{C})\in\tilde{R}_k^{-}\times\mathrm{Sd}_k$.
By Section~3.2, respectively by Section~3.9, the two series coincide term by term.
Hence their sums coincide.
This completes the proof of (al) and of (ga4).

\emph{The two other parts of $\mathscr{E}^{+}(Y)$: $(R\mathscr{E})^{\mathrm{nat}}(X)$ and ${\delta E^{\mathrm{det}}}_k(X)$.}
By (8.6)$^{\mathrm{ch}}$,
\[
\mathscr{E}^{+}(Y)=\sum_{X:\,\bar{X}^{L}=LY}S_L\left[(R\mathscr{E})^{\mathrm{nat}}(X)+{\delta E^{\mathrm{det}}}_k(X)+E^{\#,x}_k(m,\mathscr{E})(X)\right]\qquad\text{at }(\mathscr{A}_L,\mathscr{C}_L).
\]
By Section~I.5.6, the element $(R\mathscr{E})^{\mathrm{nat}}(X)$ is formed from $\mathscr{E}(X)$, from the numbers $\mathrm{al}_{0}(\mathscr{E},X)$, $\mathrm{al}_{2}(\mathscr{E},X)$ and from the homomorphism $\boldsymbol{\Theta}_{X,\mathscr{T}}$.
By Section~I.5.7,
\[
{\delta E^{\mathrm{det}}}_k(X;\mathscr{A},\mathscr{C})=E^{d}_k(X,\mathscr{A}+\mathscr{C})-E^{d}_k(X,\mathscr{C}),
\]
and $E^{d}_k$ is the mean of four values of $\acute{E}_k$.
The displays of $\acute{E}_k$ contain $S_{k,y}(X,B)$ under traces, together with the exact operators $B_{k,y}(B)$, $Q_k(B)J^{*}$, $JQ_k^{T}(-B)$, $\gamma_5$ for $k\ge1$, and $F_0(B)$, $\gamma_5$ for $k=0$; by W12 (Section~G1.11) the kernel of $S_{k,y}(X,B)$ vanishes unless both indices lie over $X$.

\emph{The reading (o).}
The arguments of Sections~2, 3, 5 and~6 apply in both cases of Section~0.5.
This completes the proof of Proposition~24.3.
\end{proof}

\subsection*{7. Statement of Proposition 24.2}

Proposition~24.2 states invariances of $E^{\#,x}_k$ at level~$k$.
Its setting is (ST1)$^{\mathrm{ch}}$, with $(m,\mathscr{E})\in\mathrm{Ball}^{x}_{k}(4)$, $X\in D_k$ and $(A,\mathscr{C})\in\tilde{R}_k^{-}\times\mathfrak{X}_k$, and it asserts four identities in $\mathrm{Gr}[\psi,\chi]$ for the element $E^{\#,x}_k(m,\mathscr{E})(X;A,\mathscr{C})$.
In clause~(a) the transformation is $a_{\tilde{g}}$: for every $\tilde{g}\in\mathbb{G}_k^{L}$,
\[
a_{\tilde{g}}\left[E^{\#,x}_k(m,\mathscr{E})\left(\tilde{g}X;\tilde{g}A,\tilde{g}\mathscr{C}\right)\right]=E^{\#,x}_k(m,\mathscr{E})(X;A,\mathscr{C}).
\]
In clause~(b) it is $G_{\boldsymbol{\lambda}+\boldsymbol{\lambda}'}$, where $\boldsymbol{\lambda}$ and $\boldsymbol{\lambda}'$ are real functions of $c$:
whenever $(A^1,\mathscr{C}^1)$ lies in $\tilde{R}_k^{-}\times\mathfrak{X}_k$ and $A^1=A-d\boldsymbol{\lambda}$, $\mathscr{C}^1=\mathscr{C}-d\boldsymbol{\lambda}'$ on the bonds at distance at most~$1$ from $X$, one has
\[
G_{\boldsymbol{\lambda}+\boldsymbol{\lambda}'}\left[E^{\#,x}_k(m,\mathscr{E})(X;A^1,\mathscr{C}^1)\right]=E^{\#,x}_k(m,\mathscr{E})(X;A,\mathscr{C});
\]
this is the invariance under (T-G) relative to the polymer~$X$.
In clause~(c) the map $\mathrm{Cg}$ is applied to the value at the pair $(-A,-\mathscr{C})$:
\[
\mathrm{Cg}\left[E^{\#,x}_k(m,\mathscr{E})(X;-A,-\mathscr{C})\right]=E^{\#,x}_k(m,\mathscr{E})(X;A,\mathscr{C}).
\]
Clause~(d) assumes in addition that $m$ is real and that $\mathscr{E}\in K^{x}_{k}$;
under these hypotheses the map $\mathbf{K}$, applied to the value at the pair $\left(-\overline{A},-\overline{\mathscr{C}}\right)$, gives
\[
\mathbf{K}\left[E^{\#,x}_k(m,\mathscr{E})\left(X;-\overline{A},-\overline{\mathscr{C}}\right)\right]=E^{\#,x}_k(m,\mathscr{E})(X;A,\mathscr{C}).
\]
The same assertions hold in the reading~(o).

\subsection*{8. Proof of Proposition 24.2}

The notions $\mathrm{ar}^{\theta}$, ``invariant on $\mathrm{Ar}$'', ``covariant on $\mathrm{Ar}$ (at an index)'', ``relative to a polymer'', and the domains, are those of Definition~24.I (Section~0.1).
Throughout, $(m,\mathscr{E})\in\mathrm{Ball}^{x}_{k}(4)$.

\subsubsection*{8.1. The four maps $\theta$; the maps of the arguments}

\begin{description}
\item[(th0)] \emph{The maps of the algebra.}
Let $\theta$ be one of the maps $a_{\tilde{g}}$ ($\tilde{g}\in\mathbb{G}_k^{L}$), $G_{\Lambda}$ ($\Lambda=\boldsymbol{\lambda}+\boldsymbol{\lambda}'$ real), $\mathrm{Cg}$, $\mathbf{K}$, acting on $\mathrm{Gr}[\psi,\chi,\Psi,W]$ as in Section~I.4.5; on $\Psi$ and $W$ it acts, for $k\ge1$, with the unit sites $y$ in the place of $x$, and, for $k=0$, as on $\psi$.
The map $\mathrm{Cg}$ is extended by the identity on $W$ (Definition~24.I).
Then $\theta$ is a ring automorphism, linear over $\mathbb{C}$ for $a_{\tilde{g}}$, $G_{\Lambda}$, $\mathrm{Cg}$ and antilinear for $\mathbf{K}$.
It maps every generator, bijectively, to a generator of the same family times a number of modulus one, and it preserves all the norms (I.5.1).
By Section~I.4.5 and Proposition~16.1\,(b), applied to the families $\Psi$, $W$, the map $\theta$ has these properties and is compatible with (I.AP) on $\mathrm{Gr}[\psi,\chi,\Psi,W]$;
for $\theta=\mathrm{Cg}$ both facts hold for its extension by the identity on $W$, the map which is $\mathrm{Cg}$ on $\psi$, $\chi$, $\Psi$ and the identity on $W$.
Hence $\theta$ is continuous.
It commutes with real-linear limits, and therefore with integrals of continuous functions with respect to real measures.
\item[(th1)] \emph{The maps of the connections and of the variables.}
Put $B^{\theta}:=\tilde{g}B$, $B^{\theta}:=-B$, $B^{\theta}:=-\overline{B}$ for (T-g), (T-C), (T-K) respectively.
The sets $\tilde{R}_k^{-}$, $\tilde{R}_k$, $\tilde{G}_k$ are invariant under these three maps by DF3, and so is $\mathrm{Sd}_k$, by (Sd3).
The set $\mathrm{Box}_k$ is invariant under $\mathsf{Z}\mapsto\tilde{g}\mathsf{Z}$ and under $\mathsf{Z}\mapsto-\mathsf{Z}$ by (P3) of Section~0.6, that is by \cite[Proposition~15.26]{QEDI}; the polydisk $\mathrm{Pd}_k$ is invariant under $s\mapsto\tilde{g}s$ and under $s\mapsto\overline{s}$; the circles $|t|=\mathrm{rad}$ are invariant under $t\mapsto\overline{t}$.
Hence, for (T-g), (T-C), (T-K), every $\mathrm{ar}^{\theta}$ lies in the domain in which $\mathrm{ar}$ lies.
For (T-G), the condition that $(A^1,\mathscr{C}^1)$ lie in the domain is part of the notion (Definition~24.I).
\item[(th2)] \emph{Geometry for $\tilde{g}=(S,a)\in\mathbb{G}_k^{L}$.}
Let $r=r(S)=\mathrm{diag}(\pm1)$ and $a\in LM\mathbb{Z}^3$.
The map $c\mapsto\tilde{g}c=rc+a$ is an isometry of $\mathbb{L}_k$ for the sup norm, with $\tilde{g}(c+\ell n)=\tilde{g}c+\ell\cdot rn$; it induces bijections of $T_k$ and of $T^0_k$.
It maps $M\mathbb{Z}^3$ onto $M\mathbb{Z}^3$ and $LM\mathbb{Z}^3$ onto $LM\mathbb{Z}^3$, and
\[
\tilde{g}\{c:|c-z|<d\}=\{c:|c-\tilde{g}z|<d\},\qquad\tilde{g}\{c:|c-z|\le d\}=\{c:|c-\tilde{g}z|\le d\}.
\]
So it maps $M$-cubes, $LM$-cubes and the cubes of $\pi_k$, $\pi'_k$, $\pi''_k$ onto cubes of the same kind, and centres onto centres.
Two $M$-cubes have a common face if and only if their centres differ by $\pm Me_{\mu}$, and $r$ maps $\pm Me_{\mu}$ to $\pm Me_{\mu}$; so $\tilde{g}$ preserves the relation ``having a common face''.
Being a bijection of the lattice sites, $\tilde{g}$ preserves the relations ``having a common lattice site'', ``contained in'' and ``meets''.
Consequently polymers are mapped to polymers and components to components, and
\[
P_0(\tilde{g}X)=\tilde{g}P_0(X),\qquad\mathrm{Yp}(\tilde{g}X,\tilde{g}P)=\tilde{g}\,\mathrm{Yp}(X,P),\qquad Y'(\tilde{g}X,\tilde{g}P)=\tilde{g}\,Y'(X,P);
\]
incompatible polymers are mapped to incompatible polymers, so that $\mathrm{urs}^{T}(\tilde{g}X_1,\dots,\tilde{g}X_n)=\mathrm{urs}^{T}(X_1,\dots,X_n)$; and $\tilde{g}\,\mathbf{1}[\mathrm{Ps}]=\mathbf{1}[\tilde{g}\,\mathrm{Ps}]$, because $(\tilde{g}s)_{\mathrm{pc}}=s_{\tilde{g}^{-1}\mathrm{pc}}$.
On the fermion sites, $\tilde{g}x=(rc+a,rg)$; by (I.a1), (I.a2) this map preserves $\int\mathrm{d}x$ and $\mathrm{Pz}_k$, maps the sites over $Z$ to the sites over $\tilde{g}Z$, and maps the pairs admissible for $Z$ to the pairs admissible for $\tilde{g}Z$.
\end{description}

\subsubsection*{8.2. The integral over $W$}

\begin{named}{Lemma 8.A}
\begin{itemize}
\item[(a)] Let $\mathrm{Gr}[\psi,\chi,\phi]$ be the Grassmann algebra with further independent generators $\phi(i)$, $i$ in a finite set, and let $\mathrm{iw}_1$, $\mathrm{iw}_2$ be homomorphisms of it into $\mathrm{Gr}[\psi,\chi,W]$ which fix $\psi$, $\chi$ and map every $\phi(i)$ to a linear combination of the generators $W(u)$, barred and unbarred.
If
\[
\int\mathrm{d}\mu_I(W)\,\mathrm{iw}_1(\phi(i))\cdot\mathrm{iw}_1(\phi(i'))=\int\mathrm{d}\mu_I(W)\,\mathrm{iw}_2(\phi(i))\cdot\mathrm{iw}_2(\phi(i'))\qquad\text{for all }i,i',
\]
then $\int\mathrm{d}\mu_I(W)\,\mathrm{iw}_1\left(F\right)=\int\mathrm{d}\mu_I(W)\,\mathrm{iw}_2\left(F\right)$ for all $F\in\mathrm{Gr}[\psi,\chi,\phi]$.
\item[(b)] Let $\theta$ be a ring homomorphism of $\mathrm{Gr}[\psi,\chi,W]$ into itself, linear or antilinear over $\mathbb{C}$, which maps $\mathrm{Gr}[\psi,\chi]$ into itself and acts on the generators $W$ by
\[
W(\mathrm{tu})\mapsto(\mathrm{Uw}\,W)(\mathrm{tu}),\qquad\bar{W}(\mathrm{tu})\mapsto\left((\mathrm{Uw}^{-1})^{T}\bar{W}\right)(\mathrm{tu}),
\]
where $\mathrm{Uw}$ is an invertible linear map of the spinors on the unit points and the transpose is taken for $\langle\,,\,\rangle_{\nu}$ (with $\nu=1$ for $k\ge1$; for $k=0$ the pairing is $\langle\,,\,\rangle^{\mathrm{ch}}$, with $\nu=\frac{1}{48}$).
Then
\[
\int\mathrm{d}\mu_I(W)\,\theta F=\theta\left(\int\mathrm{d}\mu_I(W)\,F\right)\qquad\text{for all }F\in\mathrm{Gr}[\psi,\chi,W].
\]
\item[(c)] The hypothesis of (b) holds for $a_{\tilde{g}}$, with $(\mathrm{Uw}\,f)(\mathrm{tu})=S\,f(\tilde{g}^{-1}\mathrm{tu})$; for $G_{\Lambda}$, with $\mathrm{Uw}$ the multiplication by $\exp(ie_k\cdot\Lambda)$ (read at $y$, respectively at $c=\pi(x)$); for $\mathbf{K}$, with $\mathrm{Uw}=\varkappa$; and for $\mathrm{Cg}$, extended by the identity on $W$, with $\mathrm{Uw}=1$.
\end{itemize}
\end{named}

\begin{proof}
(a) The element $F$ is a linear combination of elements $F_0\cdot\phi(i_1)\cdots\phi(i_n)$, where $F_0$ is a monomial in $\psi$, $\chi$, and $\mathrm{iw}_a$ maps such an element to $F_0\cdot\mathrm{iw}_a(\phi(i_1))\cdots\mathrm{iw}_a(\phi(i_n))$.
By (Wi), that is Proposition~15\,(i), for elements $v_1,\dots,v_n$ of the linear span of the $W(u)$ the integral $\int\mathrm{d}\mu_I(W)\,F_0\cdot v_1\cdots v_n$ equals $F_0$ times the Pfaffian of the matrix $\left(\int\mathrm{d}\mu_I(W)\,v_j\cdot v_{j'}\right)$ if $n$ is even, and it vanishes if $n$ is odd.
For $v_j=\mathrm{iw}_a(\phi(i_j))$ these matrices are the same for $a=1$ and $a=2$, by hypothesis.

(b) Write $F=\sum_Jc_J\cdot W(J)$ with $c_J\in\mathrm{Gr}[\psi,\chi]$.
Then $\theta F=\sum_J\theta(c_J)\cdot\theta(W(J))$; here $\theta(c_J)$ lies in $\mathrm{Gr}[\psi,\chi]$, and $\theta(W(J))$ is the product of the elements $\theta(W(u))$, $u\in J$, that is, $s_{\mathrm{Uw}}(W(J))$ in the notation of (N4$'$), Proposition~15\,(h).
(The coefficient of $W(J)$ is $1$; so this holds for antilinear $\theta$ as well.)
By (N4$'$),
\[
\int\mathrm{d}\mu_I(W)\,\theta(W(J))=\int\mathrm{d}\mu_I(W)\,W(J)=m(J).
\]
By the identity in the first sentence of (N3)$^{w}$, that is Proposition~15\,(f), a factor in $\mathrm{Gr}[\psi,\chi]$ that stands to the left of a linear combination of the $W(J')$ may be taken out of the integral; in particular $\int\mathrm{d}\mu_I(W)\,F=\sum_Jc_J\cdot m(J)$.
The numbers $m(J)$ are real by (mo), that is Proposition~15\,(e).
Therefore
\[
\int\mathrm{d}\mu_I(W)\,\theta F=\sum_J\theta(c_J)\cdot m(J)=\theta\left(\sum_Jc_J\cdot m(J)\right)=\theta\left(\int\mathrm{d}\mu_I(W)\,F\right).
\]

(c) \emph{The map $a_{\tilde{g}}$.}
By Section~I.4.5, $a_{\tilde{g}}W(\mathrm{tu})=S\,W(\tilde{g}^{-1}\mathrm{tu})=(\mathrm{Uw}\,W)(\mathrm{tu})$.
For spinors $f'$, $f$ on the unit points,
\[
\langle f',\mathrm{Uw}\,f\rangle_{\nu}=\int\mathrm{d}\,\mathrm{tu}\;f'(\mathrm{tu})\cdot S\,f(\tilde{g}^{-1}\mathrm{tu})=\int\mathrm{d}\,\mathrm{tu}\;\left(S^{T}f'(\tilde{g}\,\mathrm{tu})\right)\cdot f(\mathrm{tu}),
\]
because $\tilde{g}$ preserves $\int\mathrm{d}\,\mathrm{tu}$ by (th2) (the products are $\ell\mathbb{Z}^3$-periodic).
Hence
\[
(\mathrm{Uw}^{T}f')(\mathrm{tu})=S^{T}f'(\tilde{g}\,\mathrm{tu}),\qquad(\mathrm{Uw}^{-1}f)(\mathrm{tu})=S^{-1}f(\tilde{g}\,\mathrm{tu}),\qquad\left((\mathrm{Uw}^{-1})^{T}f'\right)(\mathrm{tu})=(S^{-1})^{T}f'(\tilde{g}^{-1}\mathrm{tu}),
\]
and at $f'=\bar{W}$ the last expression is $a_{\tilde{g}}\bar{W}(\mathrm{tu})$.
\emph{The map $G_{\Lambda}$.}
A multiplication operator is its own transpose, and $(\mathrm{Uw}^{-1})^{T}=\exp(-ie_k\cdot\Lambda)$.
\emph{The map $\mathbf{K}$.}
One has $(\varkappa^{-1})^{T}=\varkappa$ by (I.c4).
\emph{The map $\mathrm{Cg}$.}
Here $\mathrm{Uw}=1$.
In all four cases $\theta$ maps $\mathrm{Gr}[\psi,\chi]$ into itself.
\end{proof}

\subsubsection*{8.3. The rules (Pr1)--(Pr5)}

Let $\theta$ be one of the four maps, $\mathrm{ar}\mapsto\mathrm{ar}^{\theta}$ the corresponding map of the arguments, and $\mathrm{Ar}$ a set of arguments which contains, together with $\mathrm{ar}$, every admitted $\mathrm{ar}^{\theta}$.

For $s=(s',s'')\in\mathrm{Pd}_k$ and $\mathsf{Z}'\in\mathrm{Box}_k$ put
\[
\mathscr{Z}(s')[\mathsf{Z}']:=H_k(s')\cdot C\cdot C_{k}^{1/2}(s')\cdot\mathsf{Z}'.
\]
This is the field of (8.1)$^{\mathrm{ch}}$ formed with $\mathsf{Z}'$ in the place of $\mathsf{Z}$, and it is linear in $\mathsf{Z}'$.
Without a bracket, $\mathscr{Z}(s')$ denotes $\mathscr{Z}(s')[\mathsf{Z}]$, where $\mathsf{Z}$ is the entry of $\mathrm{ar}$.
For $\mathrm{ar}=(X,P,s,t;A,\mathscr{C},\mathsf{Z})$ and a transform $\mathrm{ar}^{\theta}$, the entries of $\mathrm{ar}^{\theta}$ are written with the superscript $\theta$:
\[
\mathrm{ar}^{\theta}=\left(X^{\theta},P^{\theta},s^{\theta},t^{\theta};A^{\theta},\mathscr{C}^{\theta},\mathsf{Z}^{\theta}\right),\qquad s^{\theta}=\left((s')^{\theta},(s'')^{\theta}\right),\qquad\mathscr{T}^{\theta}:=A^{\theta}+\mathscr{C}^{\theta}.
\]
For (T-G), $\mathrm{ar}^{\theta}$ is any admitted transform, the same one throughout a formula.
For a part $o$ of the entries of $\mathrm{ar}$, $o^{\theta}$ is the corresponding part of the entries of $\mathrm{ar}^{\theta}$.
For a polymer, a cube, a set of cubes, or an index made of these (such as $\mathbf{Y}$, $\square$, $i$), the superscript $\theta$ denotes the image under $\tilde{g}$ (see (th2)) for (T-g), and the object itself for (T-G), (T-C) and (T-K).
Explicitly, for (T-G) the entries $X$, $P$, $s$, $t$, $\mathsf{Z}$ are unchanged, and
\[
A^{\theta}=A^{1},\qquad\mathscr{C}^{\theta}=\mathscr{C}^{1},\qquad\mathscr{T}^{\theta}=\mathscr{T}^{1}:=A^{1}+\mathscr{C}^{1}.
\]
For (T-g), (T-C) and (T-K), the connections $A^{\theta}$, $\mathscr{C}^{\theta}$, $\mathscr{T}^{\theta}$ are the connections $B^{\theta}$ of (th1) at $B=A$, $\mathscr{C}$, $\mathscr{T}$ respectively.

\begin{named}{(Pr1)}[Rules that hold point by point]
\begin{itemize}
\item[(i)] If $F_1$, $F_2$ are invariant on $\mathrm{Ar}$, then so are $F_1+F_2$, $F_1F_2$ and $\mathrm{cf}\cdot F_1$, for every complex number $\mathrm{cf}$ (for (T-K): for every real number $\mathrm{cf}$).
If $F_1$ is even, then $\exp F_1$ is invariant on $\mathrm{Ar}$.
If the $F_n$ are invariant on $\mathrm{Ar}$ and $\sum_nF_n$ converges at every point of $\mathrm{Ar}$, then the sum is invariant on $\mathrm{Ar}$.
\item[(ii)] (Fields.)
Let $\phi$ be one of the families of generators, or the pair $\psi^{\#}=(\psi,\chi)$, and let $\theta(\phi(i))=\sum_jO_{ij}\cdot\phi(j)$.
Let $v$ be a family of odd elements that stands in the place of $\phi$.
For a value $v(\mathrm{ar})$ let $\mathrm{sbs}_{v(\mathrm{ar})}$ be the homomorphism that maps $\phi(i)$ to $v(\mathrm{ar})(i)$, or the homomorphism that maps $\phi(i)$ to $\phi(i)+v(\mathrm{ar})(i)$, in either case fixing the generators of the other families.
Let $F$ be invariant on $\mathrm{Ar}$, and let $v$ be covariant on $\mathrm{Ar}$ at the indices $i$ of the generators $\phi(i)$ that the values $F(\mathrm{ar}^{\theta})$ contain, that is,
\[
\theta\left(v(\mathrm{ar}^{\theta})(i)\right)=\sum_jO_{ij}\cdot v(\mathrm{ar})(j).
\]
Then $\mathrm{ar}\mapsto\mathrm{sbs}_{v(\mathrm{ar})}\left(F(\mathrm{ar})\right)$ is invariant on $\mathrm{Ar}$.
Together with $v$, the family $t\cdot v$ is covariant.
\item[(iii)] (Connections.)
Let $F_1(X;A_1,\mathscr{C})$ be invariant on a set $\mathrm{Arq}$ of arguments $(X;A_1,\mathscr{C})$.
Assume that for every $\mathrm{ar}\in\mathrm{Ar}$ and every admitted $\mathrm{ar}^{\theta}$ the argument $\left(X;A+t\cdot\mathscr{Z}(s')[\mathsf{Z}],\mathscr{C}\right)$ lies in $\mathrm{Arq}$ and $\left(X^{\theta};A^{\theta}+t^{\theta}\cdot\mathscr{Z}((s')^{\theta})[\mathsf{Z}^{\theta}],\mathscr{C}^{\theta}\right)$ is an admitted transform of it.
Then $\mathrm{ar}\mapsto F_1\left(X;A+t\cdot\mathscr{Z}(s')[\mathsf{Z}],\mathscr{C}\right)$ is invariant on $\mathrm{Ar}$.
\end{itemize}
\end{named}

\begin{proof}
(i) The map $\theta$ is a continuous ring homomorphism which is linear over $\mathbb{R}$ (over $\mathbb{C}$ for the first three maps).
Hence
\[
\theta\left(F_1(\mathrm{ar}^{\theta})\cdot F_2(\mathrm{ar}^{\theta})\right)=\theta\left(F_1(\mathrm{ar}^{\theta})\right)\cdot\theta\left(F_2(\mathrm{ar}^{\theta})\right)=F_1(\mathrm{ar})\cdot F_2(\mathrm{ar}),
\]
and the same argument applies to sums, to the exponential series and to limits.

(ii) Let $\mathrm{Gq}$ be the subalgebra generated by the generators that $F(\mathrm{ar}^{\theta})$ contains.
The maps $\theta\circ\mathrm{sbs}_{v(\mathrm{ar}^{\theta})}$ and $\mathrm{sbs}_{v(\mathrm{ar})}\circ\theta$ are ring homomorphisms on $\mathrm{Gq}$, both linear or both antilinear.
On a generator $\phi(i)$ of $\mathrm{Gq}$,
\[
\theta\left(\mathrm{sbs}_{v(\mathrm{ar}^{\theta})}\phi(i)\right)=\theta\left(v(\mathrm{ar}^{\theta})(i)\right)\left[+\theta(\phi(i))\right]=\sum_jO_{ij}\cdot\left(v(\mathrm{ar})(j)\left[+\phi(j)\right]\right)=\mathrm{sbs}_{v(\mathrm{ar})}\left(\theta\phi(i)\right),
\]
where the terms in brackets occur for the homomorphism of the second kind.
A generator of another family is fixed by $\mathrm{sbs}_{v(\mathrm{ar}^{\theta})}$, and $\theta$ maps it to a multiple of a generator of its own family, which $\mathrm{sbs}_{v(\mathrm{ar})}$ fixes.
So the two homomorphisms coincide on $\mathrm{Gq}$, and
\[
\theta\left(\mathrm{sbs}_{v(\mathrm{ar}^{\theta})}F(\mathrm{ar}^{\theta})\right)=\mathrm{sbs}_{v(\mathrm{ar})}\left(\theta F(\mathrm{ar}^{\theta})\right)=\mathrm{sbs}_{v(\mathrm{ar})}F(\mathrm{ar}).
\]
As to the factor $t$: for linear $\theta$ one has $t^{\theta}=t$; for (T-K) one has $t^{\theta}=\overline{t}$ and $\theta(\overline{t}\cdot w)=t\cdot\theta(w)$.

(iii) This is the definition of the invariance of $F_1$, applied at the two arguments named.
\end{proof}

\begin{named}{(Pr2)}[Contours]
Let $\mathrm{rad}>1$, let $F$ be invariant on $\mathrm{Ar}$ and continuous in $t$ on $|t|=\mathrm{rad}$, and let $\mathrm{Ar}$ contain, together with $\mathrm{ar}$, all arguments that differ from $\mathrm{ar}$ in $t$ only and have $|t|=\mathrm{rad}$.
Then $\oint_{\mathrm{rad}}F$, as a function of the other arguments, is invariant.
\end{named}

\begin{proof}
With $t=\mathrm{rad}\cdot\exp(i\,\mathrm{ag})$ one has
\[
\oint_{\mathrm{rad}}f=(2\pi)^{-1}\int_0^{2\pi}\mathrm{d}\,\mathrm{ag}\;f\left(\mathrm{rad}\cdot\exp(i\,\mathrm{ag})\right)\cdot\left(\mathrm{rad}\cdot\exp(i\,\mathrm{ag})-1\right)^{-1}.
\]
For linear $\theta$, $t^{\theta}=t$, and $\theta$ commutes with the integral by (th0).
For (T-K), let $f^{\theta}$ denote $F$ at the transformed other arguments; then
\[
\begin{split}
\theta\left(f^{\theta}\left(\mathrm{rad}\cdot\exp(i\,\mathrm{ag})\right)\cdot\left(\mathrm{rad}\cdot\exp(i\,\mathrm{ag})-1\right)^{-1}\right)&=\left(\mathrm{rad}\cdot\exp(-i\,\mathrm{ag})-1\right)^{-1}\cdot\theta\left(f^{\theta}\left(\overline{\mathrm{rad}\cdot\exp(-i\,\mathrm{ag})}\right)\right)\\
&=\left(\mathrm{rad}\cdot\exp(-i\,\mathrm{ag})-1\right)^{-1}\cdot f\left(\mathrm{rad}\cdot\exp(-i\,\mathrm{ag})\right),
\end{split}
\]
and the substitution $\mathrm{ag}\mapsto-\mathrm{ag}$ gives the assertion.
\end{proof}

\begin{named}{(Pr3)}[The operation of (8.2)]
Let $G(\cdot\,,X)$ be as in (D3), and let $o$ denote the arguments other than $(X,P,s)$.
If $\theta\left(G(s^{\theta},X^{\theta};o^{\theta})\right)=G(s,X;o)$ for the families $s=\mathbf{1}[P_0(X)\cup\mathrm{Psb}]$, $\mathrm{Psb}$ a subset of $P$, then
\[
\theta\left(\mathrm{Tm}[G](X^{\theta},P^{\theta};o^{\theta})\right)=\mathrm{Tm}[G](X,P;o).
\]
(These families have $s=1$ on $P_0(X)$, $s=0$ outside $P_0(X)\cup P$, and $s_P\in[0,1]^P$.)
\end{named}

\begin{proof}
By (V), $\mathrm{Tm}[G](X,P;o)$ is a combination with real coefficients of the values $G(\mathbf{1}[P_0(X)\cup\mathrm{Psb}],X;o)$, $\mathrm{Psb}\subset P$.
For (T-g): $\left(\mathbf{1}[P_0(X)\cup\mathrm{Psb}]\right)^{\theta}=\mathbf{1}[P_0(\tilde{g}X)\cup\tilde{g}\,\mathrm{Psb}]$ by (th2), and $\mathrm{Psb}\mapsto\tilde{g}\,\mathrm{Psb}$ is a bijection of the subsets of $P$ onto the subsets of $\tilde{g}P$ that preserves $\#|\mathrm{Psb}|$.
For (T-G) and (T-C): $s^{\theta}=s$.
For (T-K): $s^{\theta}=\overline{s}=s$, because the families are real.
\end{proof}

\begin{named}{(Pr4)}[Sums over polymers and over sets of cubes]
Let $\mathrm{Ixs}(\mathbf{Y})$ be one of the finite index sets
\begin{gather*}
\left\{(X,P):Y'(X,P)=\mathbf{Y}\right\},\\
\left\{\{X_i\}:\text{a set of distinct polymers with the union }\mathbf{Y}\right\},\\
\left\{(X_1,\dots,X_n):X_1\cup\dots\cup X_n=\mathbf{Y}\right\},
\end{gather*}
and let $\mathrm{cf}(i)$ be real coefficients that depend on the incidence relations of the polymers only ($1$, or $\mathrm{urs}^{T}(X_1,\dots,X_n)/n!$).
If $\theta\left(F_{i^{\theta}}(o^{\theta})\right)=F_i(o)$ for all $i\in\mathrm{Ixs}(\mathbf{Y})$, then
\[
\theta\left(\sum_{i'\in\mathrm{Ixs}(\mathbf{Y}^{\theta})}\mathrm{cf}(i')\cdot F_{i'}(o^{\theta})\right)=\sum_{i\in\mathrm{Ixs}(\mathbf{Y})}\mathrm{cf}(i)\cdot F_i(o).
\]
\end{named}

\begin{proof}
For (T-G), (T-C), (T-K) one has $i^{\theta}=i$ and $\mathbf{Y}^{\theta}=\mathbf{Y}$.
For (T-g), by (th2) the map $i\mapsto\tilde{g}i$ is a bijection of $\mathrm{Ixs}(\mathbf{Y})$ onto $\mathrm{Ixs}(\tilde{g}\mathbf{Y})$ with $\mathrm{cf}(\tilde{g}i)=\mathrm{cf}(i)$.
For (T-G): a bond at distance at most $1$ from a polymer $\mathbf{Y}'$ contained in $\mathbf{Y}$ is at distance at most $1$ from $\mathbf{Y}$; so a pair of arguments that is admitted relative to $\mathbf{Y}$ is admitted relative to $\mathbf{Y}'$, and invariance relative to $\mathbf{Y}'$ implies invariance relative to $\mathbf{Y}$.
\end{proof}

\begin{named}{(Pr5)}[The fluctuation integral]
Let $F$ have its values in $\mathrm{Gr}[\psi,\chi,W]$, be continuous in $\mathsf{Z}$ on $\mathrm{Box}_k$, and satisfy $\theta\left(F(o^{\theta},\mathsf{Z}^{\theta})\right)=F(o,\mathsf{Z})$ for all $\mathsf{Z}\in\mathrm{Box}_k$.
Let $\mathsf{Z}\mapsto\mathsf{Z}^{\theta}$ be a bijection of $\mathrm{Box}_k$ that preserves $\hat{\mu}$, and let $\theta$ act on $(W,\bar{W})$ as in Lemma~8.A\,(b).
Then
\[
\theta\left(\int\mathrm{d}\mu_I(W)\,\mathrm{d}\hat{\mu}(\mathsf{Z})\,F(o^{\theta},\mathsf{Z})\right)=\int\mathrm{d}\mu_I(W)\,\mathrm{d}\hat{\mu}(\mathsf{Z})\,F(o,\mathsf{Z}).
\]
\end{named}

\begin{proof}
By the invariance of $\hat{\mu}$,
\[
\int\mathrm{d}\mu_I(W)\,\mathrm{d}\hat{\mu}(\mathsf{Z})\,F(o^{\theta},\mathsf{Z})=\int\mathrm{d}\mu_I(W)\,\mathrm{d}\hat{\mu}(\mathsf{Z})\,F(o^{\theta},\mathsf{Z}^{\theta}).
\]
The map $\theta$ commutes with the integral over $\mathsf{Z}$ by (th0), and with the integral over $W$ by Lemma~8.A\,(b); so $\theta$ of the right side is the integral of $\theta\left(F(o^{\theta},\mathsf{Z}^{\theta})\right)=F(o,\mathsf{Z})$.
\end{proof}

\subsubsection*{8.4. The chain}

\begin{named}{Lemma 8.B}
Let $\theta$ be one of the four transformations.
For (T-g), (T-C), (T-K), let $\mathrm{Ar}$ be the set of all arguments in the domains of Definition~24.I.
For (T-G), the assertions are understood relative to the polymer that belongs to the object (Definition~24.I), on the sets $\mathrm{Ar}(\mathbf{Y})$: for every polymer $\mathbf{Y}$, the argument $\mathrm{ar}$ ranges over $\mathrm{Ar}(\mathbf{Y})$, and $\mathrm{ar}^{\theta}$ ranges over the transforms of $\mathrm{ar}$ that are admitted relative to $\mathbf{Y}$.
For (T-G), moreover, the polymer in (h1) and (h3) is $X$, the polymer of $\mathscr{E}^{r}(X)$;
in (h2) the polymer of Definition~24.I is $X$, and (h2) is assumed for the transforms that are admitted relative to $\mathbf{Y}$, which are admitted relative to $X$ because $X$ lies in $\mathbf{Y}$;
in (h4) and (h6) the polymer is $\mathbf{Y}=Y'(X,P)$, the polymer of the function $G(s,X)$ into which the fields $v$ are inserted, and the covariance of these fields is understood on the set $\mathrm{Ar}(Y'(X,P))$ of Definition~24.I.
Assume that, for all polymers $X$, $\mathbf{Y}$ with $X$ contained in $\mathbf{Y}$, all cubes $\square$ contained in $\mathbf{Y}$, all $(X,P)$ with $Y'(X,P)=\mathbf{Y}$, all $s\in\mathrm{Slc}(X,P)$ and all $t$ with $|t|\le\mathrm{rad}_1$:
\begin{itemize}
\item[(h1)] $\left(X^{\theta};A^{\theta}+t^{\theta}\cdot\mathscr{Z}((s')^{\theta})[\mathsf{Z}^{\theta}],\mathscr{C}^{\theta}\right)$ is an admitted transform of $\left(X;A+t\cdot\mathscr{Z}(s')[\mathsf{Z}],\mathscr{C}\right)$, and both lie in $D_k\times\tilde{R}_k\times\mathrm{Sd}_k$;
\item[(h2)] for $k\ge1$,
\[
\theta\left\langle\bar{\Psi},M_k\left(X^{\theta},\mathscr{T}^{\theta}+t^{\theta}\cdot\mathscr{Z}((s')^{\theta})[\mathsf{Z}^{\theta}]\right)\Psi\right\rangle_1=\left\langle\bar{\Psi},M_k\left(X,\mathscr{T}+t\cdot\mathscr{Z}(s')[\mathsf{Z}]\right)\Psi\right\rangle_1,
\]
and for $k=0$,
\[
\theta\,\mathrm{Vx}\left(\mathscr{Z}((s')^{\theta})[\mathsf{Z}^{\theta}],\square^{\theta}\right)=\mathrm{Vx}\left(\mathscr{Z}(s')[\mathsf{Z}],\square\right);
\]
\item[(h3)] $\mathscr{E}^{r}(X;\cdot\,,\cdot\,)$ is invariant on $\tilde{R}_k\times\mathrm{Sd}_k$ (for (T-G): relative to $X$), and $\theta\left(m'\cdot\langle\bar{\psi},\psi\rangle_{\square^{\theta}}\right)=m'\cdot\langle\bar{\psi},\psi\rangle_{\square}$;
\item[(h4)] ($k\ge1$) the family $v=J_k^{\#}(s'',\mathscr{T},\mathscr{Z}(s'))\Psi$ is covariant at the sites over $X^{\theta}$ and at the pairs admissible for $X^{\theta}$, that is, $\theta\left(v(\mathrm{ar}^{\theta})(i)\right)=\sum_jO_{ij}\cdot v(\mathrm{ar})(j)$ for these indices $i$, in the notation of (Pr1)\,(ii);
\item[(h5)] ($k\ge1$) $\theta\circ\mathrm{sub}_{\mathbf{Y}^{\theta},\mathscr{T}^{\theta}}=\mathrm{sub}_{\mathbf{Y},\mathscr{T}}\circ\theta$ on $\mathrm{Gr}[\psi,\chi,\Psi|\mathbf{Y}^{\theta}]$.
\end{itemize}
Then $E^{(i),*}(\mathbf{Y})$, $E^{(i),x}(\mathbf{Y})$ ($i\le3$) and $E^{(\le 3),x}(\mathbf{Y})$ are invariant.
If moreover
\begin{itemize}
\item[(h6)] $W_k^{\#}(s'',\mathscr{T})$ is covariant, in the same sense, at the sites over $X^{\theta}$ and at the pairs admissible for $X^{\theta}$, and
\item[(h7)] $\theta$ acts on $(W,\bar{W})$ as in Lemma~8.A\,(b), and $\mathsf{Z}\mapsto\mathsf{Z}^{\theta}$ is a bijection of $\mathrm{Box}_k$ that preserves $\hat{\mu}$,
\end{itemize}
then $E^{(\le 5),x}(\mathbf{Y})$, $\mathrm{Km}(\mathbf{Y})$, $\mathrm{Hm}(\mathbf{Y})$ and $E^{\#,x}_k(m,\mathscr{E})(\mathbf{Y})$ are invariant.
\end{named}

\begin{proof}
(1) $G^{(1)}$.
By (h2) the integrand (for $k\ge1$), respectively the function itself (for $k=0$), is invariant; apply (Pr2) with $\mathrm{rad}=\mathrm{rad}_1$.

(2) $G^{(2)}$ ($k\ge1$).
By (h3), (h1) and (Pr1)\,(iii), the map $\mathrm{ar}\mapsto\mathscr{E}^{r}\left(X;A+t\cdot\mathscr{Z}(s'),\mathscr{C}\right)$ is invariant for $|t|\le\mathrm{rad}_2\le\mathrm{rad}_1$.
By (k2) of Section~I.5.4, its values contain $\psi(\xi)$ only with $\xi$ over $X^{\theta}$, respectively over $X$, and $\chi(\boldsymbol{\zeta})$ only with $\boldsymbol{\zeta}$ admissible for $X^{\theta}$, respectively for $X$.
By (h4) and (Pr1)\,(ii), applied to $v=t\cdot J_k^{\#}(s'',\mathscr{T},\mathscr{Z}(s'))\Psi$ and to the homomorphism that adds $v$, the integrand of $G^{(2)}$ is invariant; apply (Pr2).

(3) $G^{(3)}$ ($k\ge1$).
The element $m'\cdot\langle\bar{\psi},\psi\rangle_{\square}$ is invariant by (h3), and it contains $\psi(\xi)$ only with $\xi$ over $\square$.
By (h4) and (Pr1)\,(ii), with the same $v$, the integrand $-m'\cdot\mathrm{sb}_v\left(\langle\bar{\psi},\psi\rangle_{\square}\right)$ is invariant; apply (Pr2).

(4) By (Pr3), $\mathrm{Tm}[G^{(i)}](X,P)$ is invariant for every $(X,P)$ with $Y'(X,P)=\mathbf{Y}$ (the families of (V) lie in $\mathrm{Slc}(X,P)$); by (Pr4), $E^{(i),*}(\mathbf{Y})$ is invariant.

(5) By (24.5), $E^{(i),*}(\mathbf{Y}^{\theta};\cdot\,)$ lies in $\mathrm{Gr}[\psi,\chi,\Psi|\mathbf{Y}^{\theta}]$.
By (h5),
\[
\theta\left(E^{(i),x}(\mathbf{Y}^{\theta};o^{\theta})\right)=\theta\,\mathrm{sub}_{\mathbf{Y}^{\theta},\mathscr{T}^{\theta}}\left(E^{(i),*}(\mathbf{Y}^{\theta};o^{\theta})\right)=\mathrm{sub}_{\mathbf{Y},\mathscr{T}}\left(\theta\,E^{(i),*}(\mathbf{Y}^{\theta};o^{\theta})\right)=E^{(i),x}(\mathbf{Y};o).
\]
At $k=0$ one has $E^{(1),x}=E^{(1),*}$.
By (Pr1)\,(i), $E^{(\le 3),x}(\mathbf{Y})$ is invariant.

(6) Let $X$ lie in $\mathbf{Y}$.
By (5), $E^{(\le 3),x}(X)$ is invariant (for (T-G): relative to $X$, hence, by (Pr4), relative to $\mathbf{Y}$), and by (24.6) it contains $\psi$ only over $X$ and $\chi$ only at pairs admissible for $X$.
By (h6) and (Pr1)\,(ii), $G^{(\le 3)}(s'',X)=\mathrm{sb}_v\left(E^{(\le 3),x}(X)\right)$, with $v=W_k^{\#}(s'',\mathscr{T})$, is invariant for $s''\in\mathrm{Slc}(X,P)$, $Y'(X,P)=\mathbf{Y}$.
By (h3), (h6), (Pr1)\,(ii) with $t\cdot v$, and (Pr2) with $\mathrm{rad}=\mathrm{rad}_4$ (here $\mathrm{rad}_4\le\mathrm{rad}_1$ by (q1)), the functions $G^{(4)}$ and $G^{(5)}$ are invariant.

(7) By (Pr3) and (Pr4), $E^{(\le 5),x}(\mathbf{Y})$ is invariant; by (Pr1)\,(i) and (Pr4), so is $\mathrm{Km}(\mathbf{Y})$.
By (24.8), $\mathrm{Km}(\mathbf{Y})$ has its values in $\mathrm{Gr}[\psi,\chi,W]$ and is continuous in $\mathsf{Z}$; by (h7) and (Pr5), $\mathrm{Hm}(\mathbf{Y})$ is invariant.
The series (8.4) converges at every argument (Section~5.4), and its coefficients $\mathrm{urs}^{T}/n!$ are real; by (Pr1)\,(i) and (Pr4), $E^{\#,x}_k(m,\mathscr{E})(\mathbf{Y})$ is invariant.
\end{proof}

\emph{Level $0$.}
There is no field $\Psi$ and no substitution: (h4), (h5) and the steps (2), (3) are void, $G^{(1)}(s,\square)=\mathrm{Vx}(\mathscr{Z}(s'),\square)$, and $G^{(4)}=0$.
The generators $W$, $\bar{W}$ are indexed by $T^{\mathrm{f}}_0$ ($192$ components per unit site, weight $\frac{1}{48}$), and $\theta=a_{\tilde{g}}$, $G_{\Lambda}$, $\mathbf{K}$ acts on $W(x)$ as on $\psi(x)$; for $a_{\tilde{g}}$ this includes the permutation $g\mapsto rg$ of the labels.
The map $\mathrm{Cg}$ is the identity on $W$ at every level, $k=0$ included ((th0), Definition~24.I).
The operator $H_0(s'',B)$ is the multiplication by $\mathrm{hq}(c)$ (Section~3.9), the same number on the $48$ sites over $c$; the operator $\Gamma_0(s'',B)$ acts on $\mathfrak{S}^{\mathrm{ch}}(1,\ell)$, with the kernel (I.1.4).
In Sections~8.5--8.8 the properties W13--W16 of these two families are those given by W17.

\subsubsection*{8.5. Proof of (a): (T-g), $\tilde{g}=(S,a)\in\mathbb{G}_k^{L}$}

\begin{proof}[Proof of Proposition~24.2\,(a)]
Let $\tilde{g}=(S,a)\in\mathbb{G}_k^{L}$.
On barred fields $\tilde{g}$ acts through the map $\mathrm{glb}$,
\[
\left(\mathrm{glb}\,f\right)(x):=\left(S^{-1}\right)^{T}f\left(\tilde{g}^{-1}x\right),
\]
and by the same formula on functions of pairs and on functions of the unit points.
By Section~I.4.1 there is a sign $\mathrm{sn}\in\{1,-1\}$ with $\left(S^{-1}\right)^{T}=\mathrm{sn}\cdot\mathsf{C}^{-1}S\,\mathsf{C}$; consequently
\[
\mathrm{glb}=\mathrm{sn}\cdot\mathsf{C}^{-1}\,\tilde{g}\,\mathsf{C}.
\]
The hypotheses (h1)--(h7) of Lemma~8.B hold for $\theta=a_{\tilde{g}}$, as follows.
\begin{description}
\item[(h1)]
By W14 of Section~G4.0.4 for the families (p1), (p2) of Section~I.5.8 (for $k=0$: by W17 of Section~G4.0.4), and by (P3) of Section~0.6,
\begin{align*}
\mathscr{Z}(\tilde{g}s')[\tilde{g}\mathsf{Z}]
&=H_k(\tilde{g}s')\,C\,C_{k}^{1/2}(\tilde{g}s')\,\tilde{g}\mathsf{Z}
=H_k(\tilde{g}s')\,C\,\tilde{g}\,C_{k}^{1/2}(s')\,\mathsf{Z}\\
&=H_k(\tilde{g}s')\,\tilde{g}\,C\,C_{k}^{1/2}(s')\,\mathsf{Z}
=\tilde{g}\left(\mathscr{Z}(s')[\mathsf{Z}]\right);
\end{align*}
for complex coordinates the identity follows by linearity.
Hence
\[
\tilde{g}A+t\cdot\mathscr{Z}(\tilde{g}s')[\tilde{g}\mathsf{Z}]=\tilde{g}\left(A+t\cdot\mathscr{Z}(s')[\mathsf{Z}]\right),
\]
and this field lies in $\tilde{R}_k$ by (q1) of Proposition~24.3 and by \cite[Lemma~5.19]{QEDI}; moreover $\tilde{g}\mathscr{C}\in\mathrm{Sd}_k$ by (Sd3) of Section~0.5.
In the same way $\tilde{g}\mathscr{T}+t\cdot\mathscr{Z}(\tilde{g}s')[\tilde{g}\mathsf{Z}]=\tilde{g}\left(\mathscr{T}+t\cdot\mathscr{Z}(s')[\mathsf{Z}]\right)$, that is,
\[
\mathscr{T}^{\theta}+t\cdot\mathscr{Z}(s'^{\theta})[\mathsf{Z}^{\theta}]=\tilde{g}B,\qquad B:=\mathscr{T}+t\cdot\mathscr{Z}(s')[\mathsf{Z}].
\]
\item[(h2)]
Let $k\ge1$, and let $O$ be an operator on $\mathfrak{S}(1,\ell)$.
By Section~I.4.3,
\begin{align*}
a_{\tilde{g}}\left\langle\bar{\Psi},O\Psi\right\rangle_1
&=\sum_{y,y'}\left(\left(S^{-1}\right)^{T}\bar{\Psi}\left(\tilde{g}^{-1}y\right)\right)\cdot O(y,y')\,S\,\Psi\left(\tilde{g}^{-1}y'\right)\\
&=\sum_{y,y'}\bar{\Psi}(y)\cdot S^{-1}O\left(\tilde{g}y,\tilde{g}y'\right)S\,\Psi(y')
=\left\langle\bar{\Psi},\tilde{g}^{-1}O\,\tilde{g}\,\Psi\right\rangle_1;
\end{align*}
the substitution $y\mapsto\tilde{g}y$, $y'\mapsto\tilde{g}y'$ is legitimate because the summands are $\ell\mathbb{Z}^3$-periodic in $y$ and in $y'$.
By W14, $M_k(\tilde{g}X,\tilde{g}B)=\tilde{g}\,M_k(X,B)\,\tilde{g}^{-1}$, and (h2) follows.
For $k=0$, (h2) is Proposition~23.2\,(b) with $\mathrm{Zf}=\mathscr{Z}(s')[\mathsf{Z}]$, for which $\tilde{g}\,\mathrm{Zf}=\mathscr{Z}(\tilde{g}s')[\tilde{g}\mathsf{Z}]$.
\item[(h3)]
By (X2) of Section~0.5, $\mathscr{E}^{r}$ lies in $K^{x,C}_{k}$, and clause (k3) of Section~I.5.4 holds for the group $\mathbb{G}_k$, which contains $\mathbb{G}_k^{L}$.
For the form of an $M$-cube one has, by (I.a1),
\begin{align*}
a_{\tilde{g}}\left\langle\bar{\psi},\psi\right\rangle_{\tilde{g}\square}
&=\int_{\tilde{g}\square}\mathrm{d}x\,\left(\left(S^{-1}\right)^{T}\bar{\psi}\left(\tilde{g}^{-1}x\right)\right)\cdot\left(S\,\psi\left(\tilde{g}^{-1}x\right)\right)\\
&=\int_{\tilde{g}\square}\mathrm{d}x\,\bar{\psi}\left(\tilde{g}^{-1}x\right)\cdot\psi\left(\tilde{g}^{-1}x\right)
=\left\langle\bar{\psi},\psi\right\rangle_{\square},
\end{align*}
and since $m'$ is a number and $a_{\tilde{g}}$ is linear, $a_{\tilde{g}}\left(m'\cdot\left\langle\bar{\psi},\psi\right\rangle_{\tilde{g}\square}\right)=m'\cdot\left\langle\bar{\psi},\psi\right\rangle_{\square}$.
\item[(h4), (h6)]
By W14 (for $k=0$: by W17),
\[
H_k(\tilde{g}s'',\tilde{g}B)\,\tilde{g}=\tilde{g}\,H_k(s'',B),\qquad\Gamma_k(\tilde{g}s'',\tilde{g}B)\,\tilde{g}=\tilde{g}\,\Gamma_k(s'',B).
\]
The automorphism $a_{\tilde{g}}$ maps the family $\Psi$ to $\tilde{g}\Psi$, $\bar{\Psi}$ to $\mathrm{glb}\,\bar{\Psi}$, $W$ to $\tilde{g}W$ and $\bar{W}$ to $\mathrm{glb}\,\bar{W}$.
Therefore
\[
a_{\tilde{g}}\left[\left(H_k(\tilde{g}s'',\tilde{g}B)\Psi\right)(x)\right]=\left(H_k(\tilde{g}s'',\tilde{g}B)\,\tilde{g}\Psi\right)(x)=S\left(H_k(s'',B)\Psi\right)\left(\tilde{g}^{-1}x\right),
\]
which is the combination that $a_{\tilde{g}}$ assigns to $\psi(x)$.
By Section~I.4.4, $\delta_{\alpha,\tilde{g}B}\left(\tilde{g}f\right)=\tilde{g}\left(\delta_{\alpha,B}f\right)$; hence
\[
a_{\tilde{g}}\left[\left(\delta_{\alpha,\tilde{g}B}H_k(\tilde{g}s'',\tilde{g}B)\Psi\right)(x,x')\right]=S\left(\delta_{\alpha,B}H_k(s'',B)\Psi\right)\left(\tilde{g}^{-1}x,\tilde{g}^{-1}x'\right).
\]
For the barred fields, W14 and $-\tilde{g}B=\tilde{g}(-B)$ give
\begin{align*}
\bar{H}_k(\tilde{g}s'',\tilde{g}B)\,\mathrm{glb}
&=\mathsf{C}^{-1}H_k\left(\tilde{g}s'',\tilde{g}(-B)\right)\mathsf{C}\cdot\mathrm{sn}\cdot\mathsf{C}^{-1}\tilde{g}\,\mathsf{C}
=\mathrm{sn}\cdot\mathsf{C}^{-1}\tilde{g}\,H_k(s'',-B)\,\mathsf{C}\\
&=\left(\mathrm{sn}\cdot\mathsf{C}^{-1}\tilde{g}\,\mathsf{C}\right)\left(\mathsf{C}^{-1}H_k(s'',-B)\,\mathsf{C}\right)
=\mathrm{glb}\,\bar{H}_k(s'',B),
\end{align*}
and $\delta_{\alpha,-\tilde{g}B}\left(\mathrm{glb}\,f\right)=\mathrm{glb}\left(\delta_{\alpha,-B}f\right)$, because the H\"older quotient commutes with constant matrices.
With $B=\mathscr{T}+\mathscr{Z}(s')$ and with $B=\mathscr{T}$ these relations show that $J_k^{\#}\Psi$ is covariant.
Together with the relation for $\Gamma_k$ at $B=\mathscr{T}$ they show that $W_k^{\#}(s'',\mathscr{T})$ is covariant.
\item[(h5)]
This is the first relation of Proposition~23.1\,(f), taken at $B=\mathscr{T}$.
\item[(h7)]
This follows from Lemma~8.A\,(c) and (P3).
\end{description}
Lemma~8.B now yields
\[
a_{\tilde{g}}\left[E^{\#,x}_k(m,\mathscr{E})(\tilde{g}X;\tilde{g}A,\tilde{g}\mathscr{C})\right]=E^{\#,x}_k(m,\mathscr{E})(X;A,\mathscr{C}).
\]
\end{proof}

\subsubsection*{8.6. Proof of (b): (T-G)}

\begin{proof}[Proof of Proposition~24.2\,(b)]
Let $\mathbf{Y}$ be a polymer, let $(A,\mathscr{C})$ and $(A^1,\mathscr{C}^1)$ be elements of $\tilde{R}_k^{-}\times\mathrm{Sd}_k$, and let $\boldsymbol{\lambda}$, $\boldsymbol{\lambda}'$ be real functions of $c$ (in the case (o): $\boldsymbol{\lambda}'=0$ and $\mathscr{C}=\mathscr{C}^1=0$).
Put $\Lambda:=\boldsymbol{\lambda}+\boldsymbol{\lambda}'$, and assume that $A^1=A-d\boldsymbol{\lambda}$ and $\mathscr{C}^1=\mathscr{C}-d\boldsymbol{\lambda}'$ on the bonds at distance at most $1$ from $\mathbf{Y}$.
Put $\mathscr{T}^1:=A^1+\mathscr{C}^1$.
Let $(X,P)$ satisfy $Y'(X,P)=\mathbf{Y}$, let $s\in\mathrm{Slc}(X,P)$ and $|t|\le\mathrm{rad}_1$, and put
\[
B:=\mathscr{T}+t\cdot\mathscr{Z}(s'),\qquad B^1:=\mathscr{T}^1+t\cdot\mathscr{Z}(s').
\]
By Section~3.4, $s$ vanishes at every parameter cube that is not contained in $\mathbf{Y}$.
The connections $B$ and $B^1$ lie in $\tilde{G}_k$ by (q1), and $B-d\Lambda$ lies in $\tilde{G}_k$ by \cite[Lemma~5.19]{QEDI}.
A bond that meets $\mathbf{Y}$ has one end in $\mathbf{Y}$, and its other end lies at the sup-distance $\eta_k\le1$ from the first one.
Such a bond is therefore at distance at most $1$ from $\mathbf{Y}$, and $B^1=B-d\Lambda$ on every bond that meets $\mathbf{Y}$.
Let $\Lambda^{\mathrm{u}}$ be the restriction of $\Lambda$ to the unit points (for $k=0$: $\Lambda$ read at $c$).

\emph{Phases.}
For a path $\boldsymbol{\Gamma}$ from $c$ to $c'$ the sum that defines $(d\Lambda)(\boldsymbol{\Gamma})$ telescopes, and $\eta_k\cdot(d\Lambda)(\boldsymbol{\Gamma})=\Lambda(c')-\Lambda(c)$.
Consequently, if $B^1=B-d\Lambda$ on the bonds of the six paths of a pair $(x,x')$, then
\begin{equation}
U_{B^1}(x,x')=U_B(x,x')\cdot\exp\left(-i\cdot e_k\cdot\left(\Lambda(c')-\Lambda(c)\right)\right).
\tag{24.16}
\end{equation}
By (adm) of Proposition~24.3 this holds at every pair that is admissible for $X$: the bonds of its paths meet $X$, and for a pair inside a cluster both sides are equal to $1$.

The hypotheses (h1)--(h7) of Lemma~8.B hold for $\theta=G_{\Lambda}$, relative to $\mathbf{Y}$, as follows.
\begin{description}
\item[(h1)]
The field $\mathscr{Z}(s')$ does not contain the connections.
On the bonds at distance at most $1$ from $\mathbf{Y}$, hence on those at distance at most $1$ from $X$, one has $A^1+t\cdot\mathscr{Z}(s')=\left(A+t\cdot\mathscr{Z}(s')\right)-d\boldsymbol{\lambda}$, and both fields lie in $\tilde{R}_k$ by (q1).
\item[(h2)]
Let $k\ge1$.
By W12 and W13,
\[
M_k(X,B^1)=M_k(X,B-d\Lambda)=\exp\left(i\cdot e_k\cdot\Lambda^{\mathrm{u}}\right)M_k(X,B)\exp\left(-i\cdot e_k\cdot\Lambda^{\mathrm{u}}\right).
\]
For an operator $O$ on $\mathfrak{S}(1,\ell)$,
\begin{align*}
G_{\Lambda}\left\langle\bar{\Psi},O\Psi\right\rangle_1
&=\sum_{y,y'}\exp\left(-i\cdot e_k\cdot\Lambda(y)\right)\cdot\bar{\Psi}(y)\cdot O(y,y')\exp\left(i\cdot e_k\cdot\Lambda(y')\right)\Psi(y')\\
&=\left\langle\bar{\Psi},\exp\left(-i\cdot e_k\cdot\Lambda^{\mathrm{u}}\right)O\exp\left(i\cdot e_k\cdot\Lambda^{\mathrm{u}}\right)\Psi\right\rangle_1,
\end{align*}
and (h2) follows with $O=M_k(X,B^1)$.
For $k=0$, (h2) is the second relation of Proposition~23.2\,(b).
\item[(h3)]
Since $\mathscr{E}^{r}\in K^{x,C}_{k}$ by (X2), clause (k4) of Section~I.5.4 gives the invariance of $\mathscr{E}^{r}$ on $\tilde{R}_k\times\mathrm{Sd}_k$ relative to $X$.
In $\left\langle\bar{\psi},\psi\right\rangle_{\square}=\int_{\square}\mathrm{d}x\,\bar{\psi}(x)\cdot\psi(x)$ the two factors acquire the phases $\exp\left(-i\cdot e_k\cdot\Lambda(c)\right)$ and $\exp\left(i\cdot e_k\cdot\Lambda(c)\right)$ at the same photon site, so that the form is left unchanged.
\item[(h4), (h6)]
Let $x=(c,g)$ be a site over $X$ or a point of a pair that is admissible for $X$.
If $c\notin\mathbf{Y}$, the rows at $x$ of $H_k(s'',\cdot\,)$ and of $\bar{H}_k(s'',\cdot\,)$ vanish at $B$ and at $B^1$, by (ga2) of Proposition~24.3.
If $c\in\mathbf{Y}$, then (ga3) and W13 (for $k=0$: W17) give
\begin{align*}
H_k(s'',B^1)(x,\cdot\,)&=H_k(s'',B-d\Lambda)(x,\cdot\,)=\exp\left(i\cdot e_k\cdot\Lambda(c)\right)H_k(s'',B)(x,\cdot\,)\exp\left(-i\cdot e_k\cdot\Lambda^{\mathrm{u}}\right),\\
\bar{H}_k(s'',B^1)(x,\cdot\,)&=\exp\left(-i\cdot e_k\cdot\Lambda(c)\right)\bar{H}_k(s'',B)(x,\cdot\,)\exp\left(i\cdot e_k\cdot\Lambda^{\mathrm{u}}\right);
\end{align*}
the second relation follows from the first one, taken at $-B$ and $-\Lambda$, because $\bar{H}_k(s'',B)=\mathsf{C}^{-1}H_k(s'',-B)\,\mathsf{C}$ and $-(B-d\Lambda)=-B-d(-\Lambda)$.
For a unit point $\mathrm{tu}$ of $\mathbf{Y}$, (ga3) and W13 give
\[
\Gamma_k(s'',\mathscr{T}^1)(\mathrm{tu},\cdot\,)=\exp\left(i\cdot e_k\cdot\Lambda^{\mathrm{u}}(\mathrm{tu})\right)\Gamma_k(s'',\mathscr{T})(\mathrm{tu},\cdot\,)\exp\left(-i\cdot e_k\cdot\Lambda^{\mathrm{u}}\right).
\]
The automorphism $G_{\Lambda}$ maps $\Psi$ to $\exp\left(i\cdot e_k\cdot\Lambda^{\mathrm{u}}\right)\Psi$ and $\bar{\Psi}$ to $\exp\left(-i\cdot e_k\cdot\Lambda^{\mathrm{u}}\right)\bar{\Psi}$, and it acts on $W$, $\bar{W}$ in the same way.
In both cases, therefore,
\[
G_{\Lambda}\left[\left(H_k(s'',B^1)\Psi\right)(x)\right]=\exp\left(i\cdot e_k\cdot\Lambda(c)\right)\cdot\left(H_k(s'',B)\Psi\right)(x).
\]
Let $(x,x')$, $x'=(c',g')$, be a pair admissible for $X$.
By (24.3) and (24.16),
\begin{align*}
&G_{\Lambda}\left[\left(\delta_{\alpha,B^1}H_k(s'',B^1)\Psi\right)(x,x')\right]\\
&\quad=d_X(x,x')^{-\alpha}\cdot\Bigl(U_B(x,x')\cdot\exp\left(-i\cdot e_k\cdot\left(\Lambda(c')-\Lambda(c)\right)\right)\cdot\exp\left(i\cdot e_k\cdot\Lambda(c')\right)\cdot\left(H_k(s'',B)\Psi\right)(x')\\
&\qquad\qquad-\exp\left(i\cdot e_k\cdot\Lambda(c)\right)\cdot\left(H_k(s'',B)\Psi\right)(x)\Bigr)\\
&\quad=\exp\left(i\cdot e_k\cdot\Lambda(c)\right)\cdot\left(\delta_{\alpha,B}H_k(s'',B)\Psi\right)(x,x'),
\end{align*}
that is, the phase of the first point, which is the phase that $G_{\Lambda}$ assigns to $\chi(x,x')$.
For the barred fields one has $G_{\Lambda}\left[\left(\bar{H}_k(s'',B^1)\bar{\Psi}\right)(x)\right]=\exp\left(-i\cdot e_k\cdot\Lambda(c)\right)\cdot\left(\bar{H}_k(s'',B)\bar{\Psi}\right)(x)$.
Moreover $U_{B^1}(x,x')^{-1}=U_B(x,x')^{-1}\cdot\exp\left(i\cdot e_k\cdot\left(\Lambda(c')-\Lambda(c)\right)\right)$, and hence
\begin{align*}
&G_{\Lambda}\left[\left(\delta_{\alpha,-B^1}\bar{H}_k(s'',B^1)\bar{\Psi}\right)(x,x')\right]\\
&\quad=d_X(x,x')^{-\alpha}\cdot\Bigl(U_B(x,x')^{-1}\cdot\exp\left(i\cdot e_k\cdot\left(\Lambda(c')-\Lambda(c)\right)\right)\cdot\exp\left(-i\cdot e_k\cdot\Lambda(c')\right)\cdot\left(\bar{H}_k(s'',B)\bar{\Psi}\right)(x')\\
&\qquad\qquad-\exp\left(-i\cdot e_k\cdot\Lambda(c)\right)\cdot\left(\bar{H}_k(s'',B)\bar{\Psi}\right)(x)\Bigr)\\
&\quad=\exp\left(-i\cdot e_k\cdot\Lambda(c)\right)\cdot\left(\delta_{\alpha,-B}\bar{H}_k(s'',B)\bar{\Psi}\right)(x,x').
\end{align*}
These relations at $t=1$ and at $t=0$ show that $J_k^{\#}\Psi$ is covariant.
By (ga2) the columns of the rows of $H_k(s'',\mathscr{T})$ and of $\bar{H}_k(s'',\mathscr{T})$ lie in $\mathbf{Y}$; together with the relation for $\Gamma_k$ this gives
\[
G_{\Lambda}\left[\left(H_k(s'',\mathscr{T}^1)\,\Gamma_k(s'',\mathscr{T}^1)\,W\right)(x)\right]=\exp\left(i\cdot e_k\cdot\Lambda(c)\right)\cdot\left(H_k(s'',\mathscr{T})\,\Gamma_k(s'',\mathscr{T})\,W\right)(x),
\]
so that $W_k^{\#}(s'',\mathscr{T})$ is covariant.
\item[(h5)]
This is the second relation of Proposition~23.1\,(f), since $\mathscr{T}^1=\mathscr{T}-d\Lambda$ on the bonds that meet $\mathbf{Y}$.
\item[(h7)]
This follows from Lemma~8.A\,(c); here $\mathsf{Z}^{\theta}=\mathsf{Z}$.
\end{description}
Lemma~8.B, applied with $\mathbf{Y}=X$, yields
\[
G_{\boldsymbol{\lambda}+\boldsymbol{\lambda}'}\left[E^{\#,x}_k(m,\mathscr{E})(X;A^1,\mathscr{C}^1)\right]=E^{\#,x}_k(m,\mathscr{E})(X;A,\mathscr{C}).
\]
\end{proof}

\subsubsection*{8.7. Proof of (d): (T-K); $m$ real, $\mathscr{E}\in K^{x}_{k}$}

\begin{proof}[Proof of Proposition~24.2\,(d)]
Let $m$ be real and $\mathscr{E}\in K^{x}_{k}$.
Then $m'=m+m(\mathscr{E})$ is real by (X1), and $\mathscr{E}^{r}\in K^{x}_{k}$ by (X2).
For an operator $O$ let $\overline{O}$ denote the operator $f\mapsto\overline{O\overline{f}}$; as the measures are real, its kernel is the complex conjugate of the kernel of $O$.
The automorphism $\mathbf{K}$ maps a field $O\Psi$ to $\left(\overline{O}\right)(\varkappa\Psi)$.
The hypotheses (h1)--(h7) of Lemma~8.B hold for $\theta=\mathbf{K}$, as follows.
\begin{description}
\item[(h1)]
The families (p1), (p2) are real for real $s$ (W16; for $k=0$: W17).
For $O=H_k$ and for $O=C_{k}^{1/2}$ the maps $s'\mapsto\overline{O\left(\overline{s'}\right)}$ and $s'\mapsto O(s')$ are continuous on $\mathrm{Pd}_k$ and analytic in its interior (W2), and they coincide at the real points.
By the identity theorem (the Taylor coefficients at $0$ coincide, and the open polydisk is connected) and by continuity, they coincide on $\mathrm{Pd}_k$.
As $C$ is real (P2), it follows that $\overline{\mathscr{Z}\left(\overline{s'}\right)[-\mathsf{Z}]}=-\mathscr{Z}(s')[\mathsf{Z}]$, and therefore
\begin{gather*}
-\overline{A}+\overline{t}\cdot\mathscr{Z}\left(\overline{s'}\right)[-\mathsf{Z}]=-\overline{A+t\cdot\mathscr{Z}(s')[\mathsf{Z}]},\\
\mathscr{T}^{\theta}+t^{\theta}\cdot\mathscr{Z}(s'^{\theta})[\mathsf{Z}^{\theta}]=-\overline{B}\qquad\text{for }B:=\mathscr{T}+t\cdot\mathscr{Z}(s')[\mathsf{Z}].
\end{gather*}
\item[(h2)]
Let $k\ge1$, and let $O$ be an operator on $\mathfrak{S}(1,\ell)$.
Then
\[
\mathbf{K}\left\langle\bar{\Psi},O\Psi\right\rangle_1=\sum_{y,y'}\left(\varkappa\bar{\Psi}(y)\right)\cdot\overline{O(y,y')}\,\varkappa\Psi(y')=\left\langle\bar{\Psi},\varkappa^{T}\,\overline{O}\,\varkappa\Psi\right\rangle_1.
\]
W16 at the connection $-B$ reads $\overline{M_k\left(X,-\overline{B}\right)}=\varkappa^{-1}M_k(X,B)\,\varkappa$, and $\varkappa^{T}\varkappa^{-1}=1=\varkappa\varkappa$ by (I.c4); this gives (h2).
For $k=0$, (h2) is the fourth relation of Proposition~23.2\,(b) with $\mathrm{Zf}=\mathscr{Z}(s')[\mathsf{Z}]$, for which $-\overline{\mathrm{Zf}}=\mathscr{Z}\left(\overline{s'}\right)[-\mathsf{Z}]$.
\item[(h3)]
The invariance of $\mathscr{E}^{r}$ is clause (k7) of Section~I.5.4.
Since the weight $\eta_k^3/48$ is real,
\[
\mathbf{K}\left(m'\cdot\left\langle\bar{\psi},\psi\right\rangle_{\square}\right)=\overline{m'}\cdot\int_{\square}\mathrm{d}x\,\left(\varkappa\bar{\psi}(x)\right)\cdot\left(\varkappa\psi(x)\right)=m'\cdot\int_{\square}\mathrm{d}x\,\bar{\psi}(x)\cdot\varkappa^{T}\varkappa\,\psi(x)=m'\cdot\left\langle\bar{\psi},\psi\right\rangle_{\square}.
\]
\item[(h4), (h6)]
W16 at $-B$ gives
\[
\overline{H_k\left(\overline{s''},-\overline{B}\right)}=\varkappa^{-1}H_k(s'',B)\,\varkappa,\qquad\overline{\Gamma_k\left(\overline{s''},-\overline{B}\right)}=\varkappa^{-1}\Gamma_k(s'',B)\,\varkappa.
\]
Hence
\[
\mathbf{K}\left[\left(H_k\left(\overline{s''},-\overline{B}\right)\Psi\right)(x)\right]=\left(\varkappa^{-1}H_k(s'',B)\,\varkappa\varkappa\Psi\right)(x)=\varkappa\left(H_k(s'',B)\Psi\right)(x).
\]
For the barred fields, as $\mathsf{C}$ is real and $\varkappa\mathsf{C}=-\mathsf{C}\varkappa$,
\begin{align*}
\overline{\bar{H}_k\left(\overline{s''},-\overline{B}\right)}
&=\mathsf{C}^{-1}\,\overline{H_k\left(\overline{s''},\overline{B}\right)}\,\mathsf{C}
=\mathsf{C}^{-1}\varkappa^{-1}H_k(s'',-B)\,\varkappa\,\mathsf{C}\\
&=\varkappa^{-1}\mathsf{C}^{-1}H_k(s'',-B)\,\mathsf{C}\,\varkappa
=\varkappa^{-1}\bar{H}_k(s'',B)\,\varkappa.
\end{align*}
Moreover
\[
\overline{U_{-\overline{B}}(x,x')}=\overline{\exp\left(-i\cdot e_k\cdot\eta_k\cdot\overline{(\tau B)(c,c')}\right)}=U_B(x,x'),\qquad\overline{U_{\overline{B}}(x,x')}=U_{-B}(x,x').
\]
By (24.3), therefore, $\mathbf{K}$ maps the components with a H\"older quotient at $\left(\overline{s''},-\overline{B}\right)$ to $\varkappa$ times the corresponding components at $(s'',B)$.
Hence $J_k^{\#}\Psi$ and $W_k^{\#}(s'',\mathscr{T})$ are covariant.
\item[(h5)]
This is the fourth relation of Proposition~23.1\,(f); the mass $\mu$ is real.
\item[(h7)]
This follows from Lemma~8.A\,(c); the map $\mathsf{Z}\mapsto-\mathsf{Z}$ preserves $\mathrm{Box}_{k}$ and $\hat{\mu}$ by (P3).
\end{description}
Lemma~8.B now yields
\[
\mathbf{K}\left[E^{\#,x}_k(m,\mathscr{E})\left(X;-\overline{A},-\overline{\mathscr{C}}\right)\right]=E^{\#,x}_k(m,\mathscr{E})(X;A,\mathscr{C}).
\]
\end{proof}

\subsubsection*{8.8. Proof of (c): (T-C)}

\begin{proof}[Proof of Proposition~24.2\,(c)]
\emph{Step 1: the hypotheses (h1)--(h5) for $\theta=\mathrm{Cg}$.}
By (I.c3),
\[
\mathsf{C}^{T}=\mathsf{C}^{-1}=-\mathsf{C},\qquad\mathsf{C}^2=-1.
\]
From $\bar{H}_k(s'',B)=\mathsf{C}^{-1}H_k(s'',-B)\,\mathsf{C}$ one obtains
\begin{equation}
H_k(s'',-B)\,\mathsf{C}=\mathsf{C}\,\bar{H}_k(s'',B),\qquad\bar{H}_k(s'',-B)\,\mathsf{C}=\mathsf{C}^{-1}H_k(s'',B)\,\mathsf{C}^2=\mathsf{C}\,H_k(s'',B).
\tag{24.17}
\end{equation}
\begin{description}
\item[(h1)]
The field $\mathscr{Z}(s')$ is linear in $\mathsf{Z}$; so $-A+t\cdot\mathscr{Z}(s')[-\mathsf{Z}]=-\left(A+t\cdot\mathscr{Z}(s')[\mathsf{Z}]\right)$, and $\mathscr{T}^{\theta}+t\cdot\mathscr{Z}(s')[\mathsf{Z}^{\theta}]=-B$ for $B:=\mathscr{T}+t\cdot\mathscr{Z}(s')[\mathsf{Z}]$.
\item[(h2)]
Let $k\ge1$, and let $O$ be an operator on $\mathfrak{S}(1,\ell)$.
Since the generators anticommute,
\begin{align*}
\mathrm{Cg}\left\langle\bar{\Psi},O\Psi\right\rangle_1
&=\sum\left(-\mathsf{C}\Psi(y)\right)_a\,O(y,y')_{ab}\left(\mathsf{C}\bar{\Psi}(y')\right)_b
=\sum\left(\mathsf{C}\bar{\Psi}(y')\right)_b\,O(y,y')_{ab}\left(\mathsf{C}\Psi(y)\right)_a\\
&=\left\langle\mathsf{C}\bar{\Psi},O^{T}\mathsf{C}\Psi\right\rangle_1
=\left\langle\bar{\Psi},\mathsf{C}^{-1}O^{T}\mathsf{C}\Psi\right\rangle_1,
\end{align*}
the sums extending over $y$, $y'$ and over the spinor indices $a$, $b$.
Take $O=M_k(X,-B)$.
By W15, $M_k(X,-B)^{T}=\mathsf{C}^{-1}M_k(X,B)\,\mathsf{C}$, so that $\mathsf{C}^{-1}O^{T}\mathsf{C}=\mathsf{C}^{-2}M_k(X,B)\,\mathsf{C}^2=M_k(X,B)$.
For $k=0$, (h2) is the third relation of Proposition~23.2\,(b) with $-\mathrm{Zf}=\mathscr{Z}(s')[-\mathsf{Z}]$.
\item[(h3)]
The invariance of $\mathscr{E}^{r}$ is clause (k5) of Section~I.5.4 (by (X2), $\mathscr{E}^{r}\in K^{x,C}_{k}$).
Moreover
\begin{align*}
\mathrm{Cg}\left\langle\bar{\psi},\psi\right\rangle_{\square}
&=\int_{\square}\mathrm{d}x\,\left(-\mathsf{C}\psi(x)\right)\cdot\left(\mathsf{C}\bar{\psi}(x)\right)
=\int_{\square}\mathrm{d}x\,\left(\mathsf{C}\bar{\psi}(x)\right)\cdot\left(\mathsf{C}\psi(x)\right)\\
&=\int_{\square}\mathrm{d}x\,\bar{\psi}(x)\cdot\mathsf{C}^{T}\mathsf{C}\,\psi(x)
=\left\langle\bar{\psi},\psi\right\rangle_{\square}.
\end{align*}
\item[(h4)]
The automorphism $\mathrm{Cg}$ maps $\Psi$ to $\mathsf{C}\bar{\Psi}$ and $\bar{\Psi}$ to $-\mathsf{C}\Psi$.
By (24.17),
\begin{align*}
\mathrm{Cg}\left[\left(H_k(s'',-B)\Psi\right)(x)\right]&=\left(H_k(s'',-B)\,\mathsf{C}\bar{\Psi}\right)(x)=\mathsf{C}\left(\bar{H}_k(s'',B)\bar{\Psi}\right)(x),\\
\mathrm{Cg}\left[\left(\bar{H}_k(s'',-B)\bar{\Psi}\right)(x)\right]&=-\left(\bar{H}_k(s'',-B)\,\mathsf{C}\Psi\right)(x)=-\mathsf{C}\left(H_k(s'',B)\Psi\right)(x),
\end{align*}
and these are the combinations that $\mathrm{Cg}$ assigns to $\psi(x)$ and to $\bar{\psi}(x)$.
At the arguments $\mathrm{ar}^{\theta}$ the unbarred component with a H\"older quotient is $\delta_{\alpha,-B}H_k(s'',-B)\Psi$, and the barred one is $\delta_{\alpha,B}\bar{H}_k(s'',-B)\bar{\Psi}$.
The H\"older quotient commutes with $\mathsf{C}$; so $\mathrm{Cg}$ maps these two components to $\mathsf{C}\,\delta_{\alpha,-B}\bar{H}_k(s'',B)\bar{\Psi}$ and to $-\mathsf{C}\,\delta_{\alpha,B}H_k(s'',B)\Psi$, respectively.
\item[(h5)]
This is the third relation of Proposition~23.1\,(f).
\end{description}
The part of Lemma~8.B that requires (h1)--(h5) only gives
\begin{equation}
\mathrm{Cg}\left[E^{(\le3),x}(X;-A,-\mathscr{C},-\mathsf{Z})\right]=E^{(\le3),x}(X;A,\mathscr{C},\mathsf{Z})\qquad\text{for all }X\in D_k.
\tag{24.18}
\end{equation}

\emph{Step 2: the second localisation as one substitution.}
Let $\phi^{\#}=\left(\phi(\xi),\mathrm{phc}(\boldsymbol{\zeta})\right)$ be further independent generators, with the same indices as $\psi^{\#}=\left(\psi(\xi),\chi(\boldsymbol{\zeta})\right)$.
Extend $\mathrm{Cg}$ to $\mathrm{Gr}[\psi,\chi,\phi^{\#}]$ by
\[
\phi\mapsto\mathsf{C}\bar{\phi},\qquad\bar{\phi}\mapsto-\mathsf{C}\phi,\qquad\mathrm{phc}\mapsto\mathsf{C}\,\overline{\mathrm{phc}},\qquad\overline{\mathrm{phc}}\mapsto-\mathsf{C}\,\mathrm{phc}.
\]
Let $\mathrm{ad}:\mathrm{Gr}[\psi,\chi]\to\mathrm{Gr}[\psi,\chi,\phi^{\#}]$ be the homomorphism with
\[
\psi(\xi)\mapsto\psi(\xi)+\phi(\xi),\qquad\chi(\boldsymbol{\zeta})\mapsto\chi(\boldsymbol{\zeta})+\mathrm{phc}(\boldsymbol{\zeta}).
\]
The identity $\mathrm{Cg}\circ\mathrm{ad}=\mathrm{ad}\circ\mathrm{Cg}$ holds on the generators, and hence on $\mathrm{Gr}[\psi,\chi]$.
Put
\begin{align*}
\mathrm{Fv}(X;A,\mathscr{C},\mathsf{Z}):={}&\mathrm{ad}\left(E^{(\le3),x}(X;A,\mathscr{C},\mathsf{Z})\right)+\mathrm{ad}\left(\mathscr{E}^{r}(X;A,\mathscr{C})\right)-\mathscr{E}^{r}(X;A,\mathscr{C})\\
&-\mathbf{1}[X=\square]\cdot m'\cdot\left(\mathrm{ad}\left(\left\langle\bar{\psi},\psi\right\rangle_{\square}\right)-\left\langle\bar{\psi},\psi\right\rangle_{\square}\right).
\end{align*}
By (24.18) and (h3),
\begin{equation}
\mathrm{Cg}\left[\mathrm{Fv}(X;-A,-\mathscr{C},-\mathsf{Z})\right]=\mathrm{Fv}(X;A,\mathscr{C},\mathsf{Z}).
\tag{24.19}
\end{equation}
For $s''\in\mathrm{Pd}_k$ and $B\in\tilde{G}_k$ let $\mathrm{iw}_{s'',B}:\mathrm{Gr}[\psi,\chi,\phi^{\#}]\to\mathrm{Gr}[\psi,\chi,W]$ be the homomorphism that fixes $\psi$ and $\chi$ and maps $\phi^{\#}$ to $W_k^{\#}(s'',B)$.
Then $\mathrm{iw}_{s'',\mathscr{T}}\circ\mathrm{ad}=\mathrm{sb}_v$ with $v=W_k^{\#}(s'',\mathscr{T})$.
The integrands of $G^{(4)}$ and $G^{(5)}$ are polynomials in $t$, and for such a function $f$ one has $\oint_{\mathrm{rad}_4}f=f(1)-f(0)$ by (N5) (Proposition~15\,(j)).
Consequently
\[
\left(G^{(\le3)}+G^{(4)}+G^{(5)}\right)(s'',X)=\mathrm{iw}_{s'',\mathscr{T}}\left(\mathrm{Fv}(X;A,\mathscr{C},\mathsf{Z})\right).
\]
Let $U$ be a union of $M$-cubes and $X$ a polymer contained in $U$.
Setting all parameters other than those of the $M$-cubes of $U$ equal to $0$ turns the function $s''\mapsto\left(G^{(\le3)}+G^{(4)}+G^{(5)}\right)(s'',X)$ into a function of the parameters of the $M$-cubes of $U$.
Apply \cite[Lemma~17.12]{QEDI} to this function, with the set of the $M$-cubes of $U$ as the finite index set and the set of the $M$-cubes of $X$ as $P_0$.
Its value at $1$ is $\mathrm{iw}_{\mathbf{1}_U,\mathscr{T}}\left(\mathrm{Fv}(X)\right)$, and for a set $P$ of $M$-cubes of $U$ outside $X$ the term $f_P$ of that lemma is $\mathrm{Tm}\left[G^{(\le3)}+G^{(4)}+G^{(5)}\right](X,P)$.
By Section~3.7, $Y'(X,P)$ is the union of $X$ and of the cubes of $P$.
The terms for which it is not connected vanish (Lemma~3.C\,(b)), and every other term is attributed to a polymer $\mathbf{Y}$ contained in $U$.
Conversely, a term attributed to a polymer $\mathbf{Y}$ contained in $U$ has $X$ and the cubes of $P$ in $U$ (Section~3.2).
Summation over $X$ gives
\begin{equation}
\sum_{\mathbf{Y}\in D_k,\ \mathbf{Y}\subseteq U}E^{(\le5),x}(\mathbf{Y})=\sum_{X\in D_k,\ X\subseteq U}\mathrm{iw}_{\mathbf{1}_U,\mathscr{T}}\left(\mathrm{Fv}(X;A,\mathscr{C},\mathsf{Z})\right).
\tag{24.20}
\end{equation}

\emph{Step 3: the function}
\[
\Xi_U(A,\mathscr{C}):=\int\mathrm{d}\mu_I(W)\,\mathrm{d}\hat{\mu}(\mathsf{Z})\,\exp\left(\sum_{\mathbf{Y}\subseteq U}E^{(\le5),x}(\mathbf{Y})\right)
\]
\emph{satisfies} $\mathrm{Cg}\left[\Xi_U(-A,-\mathscr{C})\right]=\Xi_U(A,\mathscr{C})$.
The elements $\mathrm{Fv}(X)$ are even; put
\[
\mathrm{Fq}(A,\mathscr{C},\mathsf{Z}):=\exp\left(\sum_{X\subseteq U}\mathrm{Fv}(X;A,\mathscr{C},\mathsf{Z})\right).
\]
By (24.20) the integrand of $\Xi_U(A,\mathscr{C})$ is $\mathrm{iw}_{\mathbf{1}_U,\mathscr{T}}\left(\mathrm{Fq}(A,\mathscr{C},\mathsf{Z})\right)$, and by (24.19), $\mathrm{Cg}\left[\mathrm{Fq}(-A,-\mathscr{C},-\mathsf{Z})\right]=\mathrm{Fq}(A,\mathscr{C},\mathsf{Z})$.
At the parameters $\mathbf{1}_U$ write
\begin{gather*}
\phi_+:=H_k(\mathbf{1}_U,\mathscr{T})\,\Gamma_k(\mathbf{1}_U,\mathscr{T})\,W,\qquad\bar{\phi}_+:=\bar{H}_k(\mathbf{1}_U,\mathscr{T})\,\bar{W},\\
\phi_-:=H_k(\mathbf{1}_U,-\mathscr{T})\,\Gamma_k(\mathbf{1}_U,-\mathscr{T})\,W,\qquad\bar{\phi}_-:=\bar{H}_k(\mathbf{1}_U,-\mathscr{T})\,\bar{W},\\
\phi':=\mathsf{C}\bar{\phi}_-,\qquad\bar{\phi}':=\mathsf{C}^{-1}\phi_-.
\end{gather*}
The homomorphism $\mathrm{iw}_{\mathbf{1}_U,\mathscr{T}}$ maps $\left(\phi,\mathrm{phc},\bar{\phi},\overline{\mathrm{phc}}\right)$ to $\left(\phi_+,\delta_{\alpha,\mathscr{T}}\phi_+,\bar{\phi}_+,\delta_{\alpha,-\mathscr{T}}\bar{\phi}_+\right)$, and $\mathrm{iw}_{\mathbf{1}_U,-\mathscr{T}}$ maps it to $\left(\phi_-,\delta_{\alpha,-\mathscr{T}}\phi_-,\bar{\phi}_-,\delta_{\alpha,\mathscr{T}}\bar{\phi}_-\right)$.
Let $\mathrm{iw}'$ be the homomorphism that fixes $\psi$ and $\chi$ and maps $\left(\phi,\mathrm{phc},\bar{\phi},\overline{\mathrm{phc}}\right)$ to $\left(\phi',\delta_{\alpha,\mathscr{T}}\phi',\bar{\phi}',\delta_{\alpha,-\mathscr{T}}\bar{\phi}'\right)$.
\begin{itemize}
\item[(i)]
On $\mathrm{Gr}[\psi,\chi,\phi^{\#}]$ one has $\mathrm{Cg}\circ\mathrm{iw}_{\mathbf{1}_U,-\mathscr{T}}=\mathrm{iw}'\circ\mathrm{Cg}$, $\mathrm{Cg}$ being the identity on $W$.
On $\psi$ and $\chi$ both sides act as $\mathrm{Cg}$.
On the other generators,
\begin{align*}
\mathrm{Cg}\left(\mathrm{iw}_{\mathbf{1}_U,-\mathscr{T}}\,\phi\right)&=\phi_-=\mathsf{C}\bar{\phi}'=\mathrm{iw}'\left(\mathsf{C}\bar{\phi}\right)=\mathrm{iw}'\left(\mathrm{Cg}\,\phi\right),\\
\mathrm{Cg}\left(\mathrm{iw}_{\mathbf{1}_U,-\mathscr{T}}\,\bar{\phi}\right)&=\bar{\phi}_-=-\mathsf{C}\mathsf{C}\bar{\phi}_-=-\mathsf{C}\phi'=\mathrm{iw}'\left(\mathrm{Cg}\,\bar{\phi}\right),\\
\mathrm{Cg}\left(\mathrm{iw}_{\mathbf{1}_U,-\mathscr{T}}\,\mathrm{phc}\right)&=\delta_{\alpha,-\mathscr{T}}\phi_-=\mathsf{C}\,\delta_{\alpha,-\mathscr{T}}\bar{\phi}'=\mathrm{iw}'\left(\mathsf{C}\,\overline{\mathrm{phc}}\right)=\mathrm{iw}'\left(\mathrm{Cg}\,\mathrm{phc}\right),\\
\mathrm{Cg}\left(\mathrm{iw}_{\mathbf{1}_U,-\mathscr{T}}\,\overline{\mathrm{phc}}\right)&=\delta_{\alpha,\mathscr{T}}\bar{\phi}_-=-\mathsf{C}\,\delta_{\alpha,\mathscr{T}}\phi'=\mathrm{iw}'\left(\mathrm{Cg}\,\overline{\mathrm{phc}}\right).
\end{align*}
\item[(ii)]
\emph{Pair expectations.}
By (mo) and (n), that is, by Proposition~15\,(e) and (n),
\[
\int\mathrm{d}\mu_I\,W_a(\mathrm{tu})\,\bar{W}_{a'}(\mathrm{tu}')=\frac{\mathbf{1}[\mathrm{tu}=\mathrm{tu}',\,a=a']}{\nu(\mathrm{tu})}
\]
(for $k=0$ this is $48\cdot\mathbf{1}[x=x',\,a=a']$, the kernel (I.1.4) of the identity), and
\[
\int\mathrm{d}\mu_I\,W_a(\mathrm{tu})\,W_{a'}(\mathrm{tu}')=\int\mathrm{d}\mu_I\,\bar{W}_a(\mathrm{tu})\,\bar{W}_{a'}(\mathrm{tu}')=0.
\]
With the kernels of Section~0.4, for which $\int\mathrm{d}\,\mathrm{tu}=\sum_{\mathrm{tu}}\nu(\mathrm{tu})$, it follows that
\begin{align*}
\int\mathrm{d}\mu_I\,\phi_+(x)_{\beta}\,\bar{\phi}_+(x')_{\beta'}
&=\sum_{\mathrm{tu},\mathrm{tu}'}\nu(\mathrm{tu})\cdot\nu(\mathrm{tu}')\sum_{a,a'}\left(H_k\Gamma_k\right)(x,\mathrm{tu})_{\beta a}\,\bar{H}_k(x',\mathrm{tu}')_{\beta'a'}\cdot\frac{\mathbf{1}[\mathrm{tu}=\mathrm{tu}',\,a=a']}{\nu(\mathrm{tu})}\\
&=\int\mathrm{d}\,\mathrm{tu}\,\sum_{a}\left(H_k\Gamma_k\right)(x,\mathrm{tu})_{\beta a}\,\bar{H}_k(x',\mathrm{tu})_{\beta'a}
=\mathrm{Gz}(\mathscr{T})(x,x')_{\beta\beta'},
\end{align*}
the kernel (I.1.4) of the operator
\[
\mathrm{Gz}(\mathscr{T}):=H_k(\mathbf{1}_U,\mathscr{T})\,\Gamma_k(\mathbf{1}_U,\mathscr{T})\,\bar{H}_k(\mathbf{1}_U,\mathscr{T})^{T}\qquad\text{on }\mathfrak{S}^{\mathrm{ch}}(\eta_k,\ell)
\]
(the kernel of $\bar{H}_k^{T}$ is $\bar{H}_k^{T}(\mathrm{tu},x')=\bar{H}_k(x',\mathrm{tu})^{T}$).
The expectations of $\phi_+\phi_+$ and of $\bar{\phi}_+\bar{\phi}_+$ vanish.
Next, $\bar{H}_k(\mathbf{1}_U,-\mathscr{T})=\mathsf{C}^{-1}H_k(\mathbf{1}_U,\mathscr{T})\,\mathsf{C}$, $\Gamma_k(\mathbf{1}_U,-\mathscr{T})^{T}=\mathsf{C}^{-1}\Gamma_k(\mathbf{1}_U,\mathscr{T})\,\mathsf{C}$ (W15; for $k=0$: W17, the transposes being taken for $\left\langle\,\cdot\,,\cdot\,\right\rangle^{\mathrm{ch}}$), and
\[
H_k(\mathbf{1}_U,-\mathscr{T})^{T}=\left(\mathsf{C}\bar{H}_k(\mathbf{1}_U,\mathscr{T})\,\mathsf{C}^{-1}\right)^{T}=\mathsf{C}\,\bar{H}_k(\mathbf{1}_U,\mathscr{T})^{T}\mathsf{C}^{-1}.
\]
With all operators on the right taken at $(\mathbf{1}_U,\mathscr{T})$, this gives
\[
\mathrm{Gz}(-\mathscr{T})^{T}=\bar{H}_k(\mathbf{1}_U,-\mathscr{T})\,\Gamma_k(\mathbf{1}_U,-\mathscr{T})^{T}H_k(\mathbf{1}_U,-\mathscr{T})^{T}=\mathsf{C}^{-1}H_k\Gamma_k\,\mathsf{C}^2\bar{H}_k^{T}\mathsf{C}^{-1}=-\mathsf{C}^{-1}\mathrm{Gz}(\mathscr{T})\,\mathsf{C}^{-1}.
\]
The element $\phi'$ is linear in $\bar{W}$ and $\bar{\phi}'$ is linear in $W$; so the expectations of $\phi'\phi'$ and of $\bar{\phi}'\bar{\phi}'$ vanish.
Further, $\int\mathrm{d}\mu_I\,\bar{\phi}_-(x)_a\,\phi_-(x')_b=-\mathrm{Gz}(-\mathscr{T})(x',x)_{ba}=-\left(\mathrm{Gz}(-\mathscr{T})^{T}\right)(x,x')_{ab}$ and $\left(\mathsf{C}^{-1}\right)^{T}=\mathsf{C}$; hence
\[
\int\mathrm{d}\mu_I\,\phi'(x)\,\bar{\phi}'(x')^{T}=\mathsf{C}\left[-\mathrm{Gz}(-\mathscr{T})^{T}(x,x')\right]\mathsf{C}=\mathsf{C}\mathsf{C}^{-1}\mathrm{Gz}(\mathscr{T})(x,x')\,\mathsf{C}^{-1}\mathsf{C}=\mathrm{Gz}(\mathscr{T})(x,x').
\]
Under $\mathrm{iw}_{\mathbf{1}_U,\mathscr{T}}$ as well as under $\mathrm{iw}'$, the images of $\mathrm{phc}$ and $\overline{\mathrm{phc}}$ are obtained from the images of $\phi$ and $\bar{\phi}$ by the same linear maps $\delta_{\alpha,\mathscr{T}}$ and $\delta_{\alpha,-\mathscr{T}}$.
Therefore all pair expectations of the images of the generators $\phi^{\#}$ under $\mathrm{iw}_{\mathbf{1}_U,\mathscr{T}}$ and under $\mathrm{iw}'$ coincide, and Lemma~8.A\,(a) gives
\[
\int\mathrm{d}\mu_I(W)\,\mathrm{iw}'\left(F\right)=\int\mathrm{d}\mu_I(W)\,\mathrm{iw}_{\mathbf{1}_U,\mathscr{T}}\left(F\right)\qquad\text{for all }F\in\mathrm{Gr}[\psi,\chi,\phi^{\#}].
\]
\item[(iii)]
The measure $\hat{\mu}$ is invariant under $\mathsf{Z}\mapsto-\mathsf{Z}$ (P3).
This invariance, Lemma~8.A\,(b), (c) for $\mathrm{Cg}$ (with $\mathrm{Uw}=1$) and (th0) of Section~8.1, followed by (i), (24.19) and (ii), give
\begin{align*}
\mathrm{Cg}\left[\Xi_U(-A,-\mathscr{C})\right]
&=\int\mathrm{d}\mu_I(W)\,\mathrm{d}\hat{\mu}(\mathsf{Z})\,\mathrm{Cg}\left[\mathrm{iw}_{\mathbf{1}_U,-\mathscr{T}}\left(\mathrm{Fq}(-A,-\mathscr{C},-\mathsf{Z})\right)\right]\\
&=\int\mathrm{d}\mu_I(W)\,\mathrm{d}\hat{\mu}(\mathsf{Z})\,\mathrm{iw}'\left(\mathrm{Cg}\,\mathrm{Fq}(-A,-\mathscr{C},-\mathsf{Z})\right)\\
&=\int\mathrm{d}\mu_I(W)\,\mathrm{d}\hat{\mu}(\mathsf{Z})\,\mathrm{iw}'\left(\mathrm{Fq}(A,\mathscr{C},\mathsf{Z})\right)\\
&=\int\mathrm{d}\mu_I(W)\,\mathrm{d}\hat{\mu}(\mathsf{Z})\,\mathrm{iw}_{\mathbf{1}_U,\mathscr{T}}\left(\mathrm{Fq}(A,\mathscr{C},\mathsf{Z})\right)
=\Xi_U(A,\mathscr{C}).
\end{align*}
\end{itemize}

\emph{Step 4.}
By (et) of Proposition~24.3 (Section~3.12),
\[
\exp\left(\sum_{\mathbf{Y}\subseteq U}E^{(\le5),x}(\mathbf{Y})\right)=\sum_{Z\subseteq U}\mathrm{Km}(Z),
\]
where $Z$ ranges over the unions of $M$-cubes contained in $U$, the empty union included.
Hence $\Xi_U=\sum_{Z\subseteq U}\mathrm{Hm}(Z)$.
M\"obius inversion on the subsets of the set of the $M$-cubes gives
\[
\mathrm{Hm}(Z)=\sum_{U\subseteq Z}(-1)^{|Z\setminus U|_M}\,\Xi_U,
\]
and Step~3 yields $\mathrm{Cg}\left[\mathrm{Hm}(Z;-A,-\mathscr{C})\right]=\mathrm{Hm}(Z;A,\mathscr{C})$.
By the series (8.4) of Definition~24.D and by the rules (Pr1)\,(i) and (Pr4) of Section~8.3,
\[
\mathrm{Cg}\left[E^{\#,x}_k(m,\mathscr{E})(X;-A,-\mathscr{C})\right]=E^{\#,x}_k(m,\mathscr{E})(X;A,\mathscr{C}).
\]

\emph{The reading (o).}
The arguments of Sections~8.1--8.8 hold in both cases (x) and (o) of Section~0.5; in particular they give the reading (o) of Proposition~24.2.
This completes the proof of Proposition~24.2.
\end{proof}

\subsection*{9. Proof of Proposition 24.1: (f), (g), the other sentences of (a) and of (b), (c), (d), (e), (i)}

\subsubsection*{9.1. Scaling, reblocking, analyticity of $E^{\#,x}_k$}

\begin{description}
\item[(Sc1)]
For every $X\in D_k$ there is exactly one $Y\in D_{k+1}$ with $\bar{X}^{L}=LY$, by (t5) (\cite[Proposition~5.58]{QEDI}), and $X$ is contained in $LY$ (Section~3.2).
Consequently, for every family $E'=\left\{E'(X):X\in D_k\right\}$ of elements of $\mathrm{Gr}[\psi,\chi]$ of level $k$,
\[
\sum_{Y\in D_{k+1}}\ \sum_{X:\,\bar{X}^{L}=LY}S_L\left[E'(X)\right]=S_L\left[\sum_{X\in D_k}E'(X)\right].
\]
\item[(Sc2)]
Let $Y\in D_{k+1}$, and let $X$ be contained in $LY$.
The map $c\mapsto c/L$ is a bijection of $T_k$ onto $T_{k+1}=\mathbb{L}_{k+1}/(\ell/L)\mathbb{Z}^3$ which divides distances by $L$.
It maps the sites of $X$ into $Y$, bonds to bonds, a bond that meets $X$ to a bond that meets $Y$, and the path $\boldsymbol{\Gamma}^{\pi}(c,c')$ onto $\boldsymbol{\Gamma}^{\pi}(c/L,c'/L)$, since in both paths the coordinates are brought to their final values in the same order.
Moreover $x/L=(c/L,g)$, $(x/L)^{(3)}=x^{(3)}/L$, and $d_X(x/L,x'/L)$ at level $k+1$ equals $d_X(x,x')/L$ (Section~I.3.7).
It follows that a site over $X$ is mapped to a site over $Y$.
A pair of kind (I) for $X$ (the sites $c$, $c'$ and the six paths lie in $X$; for $c=c'$: $c\in X$) is mapped to a pair of $\mathrm{Pz}_{k+1}$ of kind (I) for $Y$, and a pair inside a cluster is mapped to a pair inside a cluster.
A pair $\left(x,x^{(3)}\right)$ of kind (II) for $X$ is mapped to $\left(x/L,(x/L)^{(3)}\right)$, where $c/L$ or $c_3/L$ lies in $Y$; this pair is of kind (II) for $Y$.
\item[(Sc3)]
By Proposition~18.4\,(b), every $F\in\mathrm{Gr}[\psi,\chi]$ of level $k$ satisfies
\[
\left\|S_LF\right\|_{h_{k+1}}=\left\|F\right\|_{\left(L^{-1}h_{k+1,1},\,L^{-1-\alpha}h_{k+1,2}\right)},\qquad h_{k+1,1}=L^{-\frac{1}{8}}h_{k,1},\quad h_{k+1,2}=L^{-\frac{1}{8}+\frac{\varepsilon}{2}}h_{k,2}.
\]
By (K1) and (K7) of Section~0.3, $L^{-1}\le\frac{1}{4}$, $L^{-1-\alpha}\le\frac{1}{4}$ and $L^{-\frac{1}{8}+\frac{\varepsilon}{2}}\le1$.
Hence, by (N6) (Proposition~15\,(k)),
\[
\left\|S_LF\right\|_{h_{k+1}}\le\left\|F\right\|_{h_k/4}.
\]
\item[(Sc4)]
By (K2), $\mathrm{rt}_3-\kappa_0=\kappa-6\kappa_0-4\ge\kappa_0+3>0$.
Let $Y\in D_{k+1}$.
Write $\mathrm{rt}_3=(\mathrm{rt}_3-\kappa_0)+\kappa_0$, and apply in turn (t5), (t2) (the set $LY$ consists of $L^3\cdot|Y|_M$ $M$-cubes), (t1), (t6) of \cite[Proposition~5.58]{QEDI} and (K3); this gives
\begin{align*}
\sum_{X:\,\bar{X}^{L}=LY}\exp\left(-\mathrm{rt}_3\cdot d_M(X)\right)
&\le\exp\left(-L\cdot(\mathrm{rt}_3-\kappa_0)\cdot d_M(Y)\right)\cdot\sum_{X\subseteq LY}\exp\left(-\kappa_0\cdot d_M(X)\right)\\
&\le\exp\left(-L\cdot(\mathrm{rt}_3-\kappa_0)\cdot d_M(Y)\right)\cdot K_0\cdot L^3\cdot|Y|_M\\
&\le c_X\cdot K_0\cdot L^3\cdot\exp\left(-\left(L\cdot(\mathrm{rt}_3-\kappa_0)-1\right)\cdot d_M(Y)\right)\\
&\le c_X\cdot K_0\cdot L^3\cdot\exp\left(-\kappa\cdot d_M(Y)\right).
\end{align*}
\item[(Sc5)]
By Proposition~16.1\,(a1), (d), the map $\tilde{g}=(S,a)\mapsto\tilde{g}':=(S,a/L)$ is an isomorphism of $\mathbb{G}_k^{L}$ onto $\mathbb{G}_{k+1}$, and
\begin{gather*}
L\left(\tilde{g}'Y\right)=\tilde{g}(LY),\qquad\left(\tilde{g}'A\right)_L=\tilde{g}\left(A_L\right),\qquad\overline{\tilde{g}X}^{L}=\tilde{g}\left(\bar{X}^{L}\right),\\
S_L\,a_{\tilde{g}}=a_{\tilde{g}'}\,S_L,\qquad S_L\,\mathrm{Cg}=\mathrm{Cg}\,S_L,\qquad S_L\,\mathbf{K}=\mathbf{K}\,S_L,\qquad G_{\Lambda}\,S_L=S_L\,G_{\Lambda_L},\qquad(d\Lambda)_L=d\left(\Lambda_L\right),
\end{gather*}
where $\Lambda$ is a real function on the sites of $T_{k+1}$ and $\Lambda_L(c)=L^{\frac{1}{2}}\cdot\Lambda(c/L)$.
\item[(An)]
For $(m,\mathscr{E})\in\mathrm{Ball}^{x}_{k}(4)$ and $X\in D_k$ the coefficients of $E^{\#,x}_k(m,\mathscr{E})(X;A,\mathscr{C})$ are analytic in $(A,\mathscr{C})$ on $\tilde{R}_k^{-}\times\mathrm{Sd}_k$, in the sense fixed by (Sd6).
\begin{proof}
Put $V(X):=E^{(\le5),x}(X)$.
By (24.7) the elements $V(X_i)$ are continuous in $(A,\mathscr{C},\mathsf{Z})$ and analytic in $(A,\mathscr{C})$.
The element $\mathrm{Hm}_{\mathrm{nq}}(X_i)$ of Lemma~5.C is the integral over $(W,\mathsf{Z})$ of a finite sum of products of them.
The integral over $W$ is a linear map, and the integral over $\mathsf{Z}$ preserves analyticity in $(A,\mathscr{C})$ by (Int) of Section~3.5.
Hence $E^{\#}_{\mathrm{nq}}(X)$, a polynomial in the elements $\mathrm{Hm}_{\mathrm{nq}}(X_i)$, is analytic.
By Lemma~5.C, $E^{\#}_{\mathrm{nq}}(X)\to E^{\#,x}_k(m,\mathscr{E})(X)$ uniformly on $\tilde{R}_k^{-}\times\mathrm{Sd}_k$, and a uniform limit of analytic functions is analytic.
\end{proof}
\end{description}

\subsubsection*{9.2. (f)}

\begin{proof}[Proof of Proposition~24.1\,(f)]
Let $(m,\mathscr{E})\in\mathrm{Ball}^{x}_{k}(4)$, $Y\in D_{k+1}$ and $(\mathscr{A},\mathscr{C})\in\tilde{R}_{k+1}\times\mathrm{Sd}_{k+1}$.
By DF2 (\cite[Lemma~17.5]{QEDI}) and (Sd4) the pair $(\mathscr{A}_L,\mathscr{C}_L)$ lies in $\tilde{R}_k^{-}\times\mathrm{Sd}_k$, and
\[
E^{*,x}_k(m,\mathscr{E})(Y;\mathscr{A},\mathscr{C})=\sum_{X:\,\bar{X}^{L}=LY}S_L\left[E^{\#,x}_k(m,\mathscr{E})(X;\mathscr{A}_L,\mathscr{C}_L)\right].
\]
In this section $E^{\#}$ stands for $E^{\#,x}_k(m,\mathscr{E})$ and $E^{*}$ for $E^{*,x}_k(m,\mathscr{E})$.
Unless the range is written out, a sum over $X$ extends over the polymers $X\in D_k$ with $\bar{X}^{L}=LY$.
The clauses (k1)--(k7) are those of Section~I.5.4, read at level $k+1$.

\emph{(k3).}
Let $\tilde{g}'\in\mathbb{G}_{k+1}$, and let $\tilde{g}=(S,a)\in\mathbb{G}_k^{L}$ be its preimage under the isomorphism $\tilde{g}\mapsto\tilde{g}'=(S,a/L)$ of (Sc5).
By DF3 (\cite[Lemma~17.6]{QEDI}) at level $k+1$ and by (Sd3), the pair $(\tilde{g}'\mathscr{A},\tilde{g}'\mathscr{C})$ lies in $\tilde{R}_{k+1}\times\mathrm{Sd}_{k+1}$.
By (Sc5) one has $(\tilde{g}'\mathscr{A})_L=\tilde{g}(\mathscr{A}_L)$ and $(\tilde{g}'\mathscr{C})_L=\tilde{g}(\mathscr{C}_L)$; moreover $\overline{\tilde{g}X}^{L}=\tilde{g}(\bar{X}^{L})$ and $\tilde{g}(LY)=L(\tilde{g}'Y)$, so that $X\mapsto\tilde{g}X$ is a bijection of $\{X:\bar{X}^{L}=LY\}$ onto $\{X:\bar{X}^{L}=L(\tilde{g}'Y)\}$.
With Proposition~24.2\,(a) and the relation $S_L\circ a_{\tilde{g}}=a_{\tilde{g}'}\circ S_L$ of (Sc5) it follows that
\begin{align*}
a_{\tilde{g}'}\left[E^{*}(\tilde{g}'Y;\tilde{g}'\mathscr{A},\tilde{g}'\mathscr{C})\right]
&=\sum_{X:\,\bar{X}^{L}=LY}a_{\tilde{g}'}\,S_L\left[E^{\#}\left(\tilde{g}X;\tilde{g}(\mathscr{A}_L),\tilde{g}(\mathscr{C}_L)\right)\right]\\
&=\sum_{X:\,\bar{X}^{L}=LY}S_L\,a_{\tilde{g}}\left[E^{\#}\left(\tilde{g}X;\tilde{g}(\mathscr{A}_L),\tilde{g}(\mathscr{C}_L)\right)\right]\\
&=\sum_{X:\,\bar{X}^{L}=LY}S_L\left[E^{\#}(X;\mathscr{A}_L,\mathscr{C}_L)\right]=E^{*}(Y;\mathscr{A},\mathscr{C}).
\end{align*}
The element $\tilde{g}'$ acts on the fermion sites of $T^{\mathrm{f}}_{k+1}$ by (I.4.2); with $r=r(S)$ and by (I.a4),
\[
\tilde{g}'(c',g)=\left(rc'+a/L,\ rg\right)=\tilde{g}(Lc',g)/L .
\]

\emph{(k4).}
Let $(\mathscr{A},\mathscr{C})$ and $(\mathscr{A}^1,\mathscr{C}^1)$ lie in $\tilde{R}_{k+1}\times\mathrm{Sd}_{k+1}$, let $\boldsymbol{\lambda}$, $\boldsymbol{\lambda}'$ be real functions on the sites of $T_{k+1}$ (in the case (o): $\boldsymbol{\lambda}'=0$), and assume that $\mathscr{A}^1=\mathscr{A}-d\boldsymbol{\lambda}$ and $\mathscr{C}^1=\mathscr{C}-d\boldsymbol{\lambda}'$ on the bonds at distance $\le1$ from $Y$.
By (Sc5),
\[
(\mathscr{A}^1)_L=\mathscr{A}_L-d(\boldsymbol{\lambda}_L),\qquad(\mathscr{C}^1)_L=\mathscr{C}_L-d(\boldsymbol{\lambda}'_L)
\]
on those bonds $b$ of $T_k$ for which $b/L$ is at distance $\le1$ from $Y$, that is, for which both end points of $b$ have sup-distance at most $L$ from $LY$.
For a polymer $X\subset LY$ these bonds include all bonds at distance $\le1$ from $X$; and the polymers $X$ with $\bar{X}^{L}=LY$ are contained in $LY$ (Sc1).
The four scaled fields $\mathscr{A}_L$, $(\mathscr{A}^1)_L$, resp. $\mathscr{C}_L$, $(\mathscr{C}^1)_L$, lie in $\tilde{R}_k^{-}$, resp. in $\mathrm{Sd}_k$.
Put $\Lambda:=\boldsymbol{\lambda}+\boldsymbol{\lambda}'$, so that $\Lambda_L=\boldsymbol{\lambda}_L+\boldsymbol{\lambda}'_L$.
Proposition~24.2\,(b), applied to each polymer $X$ of the sum, and (Sc5) give
\begin{align*}
G_{\Lambda}\left[E^{*}(Y;\mathscr{A}^1,\mathscr{C}^1)\right]
&=\sum_XG_{\Lambda}\,S_L\left[E^{\#}\left(X;(\mathscr{A}^1)_L,(\mathscr{C}^1)_L\right)\right]
=\sum_XS_L\,G_{\Lambda_L}\left[E^{\#}\left(X;(\mathscr{A}^1)_L,(\mathscr{C}^1)_L\right)\right]\\
&=\sum_XS_L\left[E^{\#}(X;\mathscr{A}_L,\mathscr{C}_L)\right]=E^{*}(Y;\mathscr{A},\mathscr{C});
\end{align*}
here $G_{\Lambda}$ is the automorphism of level $k+1$ and $G_{\Lambda_L}$ that of level $k$.

\emph{(k5) and (k7).}
One has $(-\mathscr{A})_L=-\mathscr{A}_L$ and $(\overline{\mathscr{A}})_L=\overline{\mathscr{A}_L}$, and the same holds for $\mathscr{C}$.
The pairs $(-\mathscr{A},-\mathscr{C})$ and $(-\overline{\mathscr{A}},-\overline{\mathscr{C}})$ lie in $\tilde{R}_{k+1}\times\mathrm{Sd}_{k+1}$ (DF3, (Sd3)), and $S_L\circ\mathrm{Cg}=\mathrm{Cg}\circ S_L$, $S_L\circ\mathbf{K}=\mathbf{K}\circ S_L$ (Sc5).
Proposition~24.2\,(c) therefore gives
\[
\mathrm{Cg}\left[E^{*}(Y;-\mathscr{A},-\mathscr{C})\right]=\sum_XS_L\,\mathrm{Cg}\left[E^{\#}(X;-\mathscr{A}_L,-\mathscr{C}_L)\right]=\sum_XS_L\left[E^{\#}(X;\mathscr{A}_L,\mathscr{C}_L)\right]=E^{*}(Y;\mathscr{A},\mathscr{C}),
\]
which is (k5).
For real $m$ and $\mathscr{E}\in K^{x}_{k}$, Proposition~24.2\,(d) gives in the same way
\[
\mathbf{K}\left[E^{*}(Y;-\overline{\mathscr{A}},-\overline{\mathscr{C}})\right]=\sum_XS_L\,\mathbf{K}\left[E^{\#}\left(X;-\overline{\mathscr{A}_L},-\overline{\mathscr{C}_L}\right)\right]=\sum_XS_L\left[E^{\#}(X;\mathscr{A}_L,\mathscr{C}_L)\right]=E^{*}(Y;\mathscr{A},\mathscr{C}),
\]
which is (k7).

\emph{Legs.}
Take in (k4) $\mathscr{A}^1=\mathscr{A}$, $\mathscr{C}^1=\mathscr{C}$, $\boldsymbol{\lambda}$ a real constant and $\boldsymbol{\lambda}'=0$; then $d\boldsymbol{\lambda}=0$, and (k4) says that $G_{\boldsymbol{\lambda}}\left[E^{*}(Y;\mathscr{A},\mathscr{C})\right]=E^{*}(Y;\mathscr{A},\mathscr{C})$.
The automorphism $G_{\boldsymbol{\lambda}}$ multiplies a monomial with $n_u$ unbarred and $n_b$ barred legs by $\exp\left(i\,e_{k+1}\cdot\boldsymbol{\lambda}\cdot(n_u-n_b)\right)$.
The basis monomials are linearly independent.
Hence the coefficient of a monomial in $E^{*}(Y;\mathscr{A},\mathscr{C})$ vanishes unless $\exp\left(i\,e_{k+1}\cdot\boldsymbol{\lambda}\cdot(n_u-n_b)\right)=1$ for all real $\boldsymbol{\lambda}$, that is, unless $n_u=n_b$.
In particular $E^{*}(Y;\mathscr{A},\mathscr{C})$ is even.
\end{proof}

\subsubsection*{9.3. (g)}

\begin{proof}[Proof of Proposition~24.1\,(g)]
The abbreviations of Section~9.2 are kept.

\emph{(k1).}
The map $(\mathscr{A},\mathscr{C})\mapsto(\mathscr{A}_L,\mathscr{C}_L)$ is linear, and it maps $\tilde{R}_{k+1}\times\mathrm{Sd}_{k+1}$ into $\tilde{R}_k^{-}\times\mathrm{Sd}_k$.
By (An), the coefficients of $E^{\#}(X;\mathscr{A}_L,\mathscr{C}_L)$ are therefore analytic on $\tilde{R}_{k+1}\times\mathrm{Sd}_{k+1}$; so are the coefficients of their images under $S_L$ and those of the finite sum $E^{*}(Y;\mathscr{A},\mathscr{C})$.
By (N0) (Proposition~15\,(a)) the antisymmetric kernels are the coefficients divided by the weights; so they are analytic as well.

\emph{(k2).}
By Proposition~24.3\,(ze) (Section~6), $E^{\#}(X;A,\mathscr{C})$ contains no generators other than the $\psi(\xi)$ with $\xi$ over $X$ and the $\chi(\boldsymbol{\zeta})$ with $\boldsymbol{\zeta}$ admissible for $X$.
The homomorphism $S_L$ sends $\psi(x,\beta,\omega)$ to $L^{-1}\psi(x/L,\beta,\omega)$ and $\chi(x,x',\beta,\omega)$ to $L^{-1-\alpha}\chi(x/L,x'/L,\beta,\omega)$.
For $X\subset LY$ the images stand, by (Sc2), at sites over $Y$ and at pairs admissible for $Y$.
Consequently the kernels of $E^{*}(Y)$ vanish unless all sites lie over $Y$ and all pairs are of the kinds (I), (II) for $Y$.
Again by Proposition~24.3\,(ze), $E^{\#}(X;\mathscr{A}_L,\mathscr{C}_L)$ depends on $\mathscr{A}_L$ and $\mathscr{C}_L$ only through their values at the bonds $b$ that meet $X$.
Here $\mathscr{A}_L(b)=L^{-\frac12}\mathscr{A}(b/L)$ and $\mathscr{C}_L(b)=L^{-\frac12}\mathscr{C}(b/L)$, and $b/L$ meets $Y$ (Sc2).
Hence $E^{*}(Y;\mathscr{A},\mathscr{C})$ depends on $\mathscr{A}$ and $\mathscr{C}$ only through the bonds that meet $Y$.

\emph{The H\"older legs.}
For a polymer $\mathbf{Y}$ of level $k$ let $\mathrm{Hlg}(\mathbf{Y})$ be the set of the pairs $(x,x')\in\mathrm{Pz}_k$, $x=(c,g)$, with one of the following two properties:
$\mathscr{E}^{r}(X;B_1,\mathscr{C})$ contains a generator $\chi(x,x',\beta,\omega)$ for some polymer $X\subset\mathbf{Y}$ and some $(B_1,\mathscr{C})\in\tilde{R}_k\times\mathrm{Sd}_k$;
or $x'=x^{(3)}$ and $c\in\mathbf{Y}$.
If $\mathbf{Y}'\subset\mathbf{Y}$, then $\mathrm{Hlg}(\mathbf{Y}')\subset\mathrm{Hlg}(\mathbf{Y})$.
The elements of $\mathrm{Gr}[\psi,\chi,\Psi,W]$ that contain $\chi$ only at pairs of $\mathrm{Hlg}(\mathbf{Y})$ form a subalgebra, and this subalgebra is a closed linear subspace.
Every object attached to $\mathbf{Y}$ belongs to it, for the following reasons.
\begin{itemize}
\item For $k\ge1$, $G^{(1)}$ contains only $\Psi$ and $\bar{\Psi}$ ((D4) of Definition~24.D): neither $\psi$ nor $\chi$.
\item $G^{(3)}$ and $G^{(5)}$ arise from elements without $\chi$ under homomorphisms that send $\psi(\xi)$ to $\psi(\xi)$ plus a linear combination of the $\Psi(u)$, resp. of the $W(u)$; so they contain no $\chi$.
\item $G^{(2)}(s,X)$ and $G^{(4)}(s'',X)$ are integrals of $\mathrm{sb}_{tv}\left(\mathscr{E}^{r}(X;\,\cdot\,)\right)$.
The homomorphism $\mathrm{sb}_{tv}$ sends $\chi(\boldsymbol{\zeta})$ to $\chi(\boldsymbol{\zeta})+tv(\boldsymbol{\zeta})$, and neither $v(\boldsymbol{\zeta})$ nor $v(\xi)$ contains $\chi$; so $\chi$ stands only at pairs at which $\mathscr{E}^{r}(X;\,\cdot\,)$ contains it.
\item At $k=0$, $G^{(1)}(s,\square)=\mathrm{Vx}(\mathscr{Z}(s'),\square)$ contains $\chi$ only at pairs $(x,x^{(3)})$ with $c\in\square$ (Proposition~23.2\,(c); (8.3$_0$)$^{\mathrm{ch}}$).
\item $\mathrm{sub}_{\mathbf{Y},\mathscr{T}}$ adds $\chi$ only at pairs of kind (II) for $\mathbf{Y}$ whose first point lies over $\mathbf{Y}$, that is, at pairs $(x,x^{(3)})$ with $c\in\mathbf{Y}$ (Proposition~23.1\,(d); (8.3a)$^{\mathrm{ch}}$).
\item $G^{(\le3)}(s'',X)=\mathrm{sb}_v\left(E^{(\le3),x}(X)\right)$ contains $\chi$ only at pairs at which $E^{(\le3),x}(X)$ contains it.
\item $\mathrm{Loc}$, sums, products, exponentials, the integrals and the series of (8.4) do not lead out of the subalgebra.
\end{itemize}
The element $E^{\#}(X)$ is attached to $X$; so it contains $\chi$ only at pairs of $\mathrm{Hlg}(X)$.
From the action of $S_L$ on the generators $\chi$ and from Proposition~24.3\,(al) (the polymers $X$ with $\bar{X}^{L}=LY$ are contained in $LY$) it follows that every H\"older leg of $E^{*}(Y)$ stands at a pair $(x/L,x'/L)$ with $(x,x')\in\mathrm{Hlg}(X)$ for some polymer $X\subset LY$.
At $k=0$ one has $\mathscr{E}^{r}=0$.
The pairs are then the pairs $(x,x^{(3)})$ with $c\in LY$, and their images $\left(x/L,(x/L)^{(3)}\right)$ are pairs $(x',x'^{(3)})$ of $\mathrm{Pz}_1$ with $d_X=\eta_1$ ((I.1.5), Section~I.3.7).
\end{proof}

\subsubsection*{9.4. The other sentences of (a)}

\begin{proof}[Proof of the remaining sentences of Proposition~24.1\,(a)]
The abbreviations of Section~9.2 are kept.

\emph{$E^{*,x}_k(m,\mathscr{E})$ lies in $K^{x,C}_{k+1}=K^{x,C}_{k+1}(\ell/L)$; the bound.}
Let $Y\in D_{k+1}$ and $(\mathscr{A},\mathscr{C})\in\tilde{R}_{k+1}\times\mathrm{Sd}_{k+1}$.
By (Sc3), (24.14) and (Sc4),
\[
\left\|E^{*}(Y;\mathscr{A},\mathscr{C})\right\|_{h_{k+1}}\le\sum_{X:\,\bar{X}^{L}=LY}\left\|E^{\#}(X;\mathscr{A}_L,\mathscr{C}_L)\right\|_{h_k/4}\le c_S'\cdot\mathbf{H}_{0}\cdot c_{X}\cdot K_{0}\cdot L^3\cdot\exp\left(-\kappa\cdot d_M(Y)\right),
\]
and, by (A7) of Section~0.3,
\[
c_{X}\cdot K_{0}\cdot L^3\cdot c_S'\cdot\mathbf{H}_{0}=504\,c_{X}^5\cdot K_{0}^5\cdot\exp(3\kappa)\cdot L^3\cdot e_k^{\frac14-5\varepsilon}\le c_{S}\cdot L^3\cdot e_k^{\varepsilon}\cdot e_k^{\frac14-6\varepsilon}\le e_k^{\frac14-6\varepsilon}.
\]
This is (k6), together with the bound $\left\|E^{*,x}_k(m,\mathscr{E})\right\|^{x}_{k+1}\le e_k^{\frac14-6\varepsilon}$.
The elements $E^{*}(Y;\mathscr{A},\mathscr{C})$ are even and have equally many barred and unbarred legs, and they satisfy (k3), (k4), (k5) (Section~9.2) as well as (k1) and (k2) (Section~9.3).

\emph{Real data.}
For real $m$ and $\mathscr{E}\in K^{x}_{k}$ the family $E^{*,x}_k(m,\mathscr{E})$ also satisfies (k7) (Section~9.2); so it lies in $K^{x}_{k+1}$.

\emph{Dependence on the datum.}
Let $(m_a,\mathscr{E}_a)$ and $(m_b,\mathscr{E}_b)$ lie in $\mathbb{C}\times K^{x,C}_{k}$ (at $k=0$: $\mathscr{E}_a=\mathscr{E}_b=0$), and put
\[
\mathrm{Tq}:=\left\{t\in\mathbb{C}:\ (m_a+tm_b,\mathscr{E}_a+t\mathscr{E}_b)\in\mathrm{Ball}^{x}_{k}(4)\right\}.
\]
Since $|m_a+tm_b|$ and $\|\mathscr{E}_a+t\mathscr{E}_b\|^{x}_{k}$ are continuous convex functions of $t$, the set $\mathrm{Tq}$ is closed and convex.
If $\mathrm{Tq}$ has at most one point, there is nothing to prove.
Otherwise $\mathrm{Tq}$ contains a segment, and it is infinite.
For $t\in\mathrm{Tq}$ write $f(t):=E^{*,x}_k(m_a+tm_b,\mathscr{E}_a+t\mathscr{E}_b)$.

(1)
Let $V_t(X):=E^{(\le5),x}(X)$ be the element formed with the datum $(m_a+tm_b,\mathscr{E}_a+t\mathscr{E}_b)$.
By Proposition~5.B\,(e), $V_t(X)=V^0(X)+tV^1(X)$, where the elements $V^0(X)$, $V^1(X)$ are continuous in $(A,\mathscr{C},\mathsf{Z})$ and analytic in $(A,\mathscr{C})$.
Hence, for every integer $\mathrm{nq}\ge1$, the element $E^{\#}_{\mathrm{nq}}(X)$ of Lemma~5.C, formed with $V_t$, is a polynomial in $t$ whose coefficients are analytic in $(A,\mathscr{C})$ (as in (An)).
Put
\[
f_{\mathrm{nq}}(t)(Y;\mathscr{A},\mathscr{C}):=\sum_{X:\,\bar{X}^{L}=LY}S_L\left[E^{\#}_{\mathrm{nq}}(X;\mathscr{A}_L,\mathscr{C}_L)\right];
\]
this is a polynomial in $t$, of some degree $\mathrm{dg}$.
For $t\in\mathrm{Tq}$ the elements $V_t(X)$ satisfy the hypotheses of Lemma~5.C.
For a family $\Phi=\{\Phi(Y;\mathscr{A},\mathscr{C})\}$ of elements of $\mathrm{Gr}[\psi,\chi]$ of level $k+1$ let
\[
\mathrm{Nq}(\Phi):=\sup_{Y\in D_{k+1}}\exp\left(\kappa\cdot d_M(Y)\right)\sup_{(\mathscr{A},\mathscr{C})\in\tilde{R}_{k+1}\times\mathrm{Sd}_{k+1}}\left\|\Phi(Y;\mathscr{A},\mathscr{C})\right\|_{h_{k+1}}.
\]
Lemma~5.C, (Sc3), (Sc4) and the computation above give $\mathrm{Nq}\left(f_{\mathrm{nq}}(t)\right)\le e_k^{\frac14-6\varepsilon}$ for $t\in\mathrm{Tq}$, and
\[
\mathrm{Nq}\left(f(t)-f_{\mathrm{nq}}(t)\right)\le\max_{Y\in D_{k+1}}\exp\left(\kappa\cdot d_M(Y)\right)\cdot\#|D_k|\cdot\mathrm{dlq}_{\mathrm{nq}}\longrightarrow0\qquad(\mathrm{nq}\to\infty),
\]
uniformly in $t\in\mathrm{Tq}$, because $D_k$ and $D_{k+1}$ are finite.

(2)
The coefficients of a polynomial of degree $\mathrm{dg}$ are linear combinations, with numerical coefficients, of its values at any $\mathrm{dg}+1$ distinct points; let these points be taken in $\mathrm{Tq}$.
Then the coefficients of $f_{\mathrm{nq}}$ have finite $\mathrm{Nq}$, and consequently $\mathrm{Nq}\left(f_{\mathrm{nq}}(t)-f_{\mathrm{nq}}(t')\right)\to0$ as $t'\to t$.
Together with (1) this shows that $t\mapsto f(t)$ is continuous on $\mathrm{Tq}$ for $\mathrm{Nq}=\|\cdot\|^{x}_{k+1}$.
Moreover $f(t)(Y;\mathscr{A},\mathscr{C})$ is jointly continuous in $(t,\mathscr{A},\mathscr{C})$ on $\mathrm{Tq}\times\tilde{R}_{k+1}\times\mathrm{Sd}_{k+1}$ and analytic in $t$ in the interior of $\mathrm{Tq}$, being the uniform limit of the $f_{\mathrm{nq}}(t)(Y;\mathscr{A},\mathscr{C})$, which are polynomials in $t$ with coefficients analytic in $(\mathscr{A},\mathscr{C})$.

(3)
\emph{The power series.}
It is now shown that, as asserted in (a), near every interior point of $\mathrm{Tq}$ the map $t\mapsto f(t)$ is the sum of a power series in $t$ with coefficients in $K^{x,C}_{k+1}$ that converges in $\|\cdot\|^{x}_{k+1}$.
Let the closed disk $\{|t-\mathrm{tce}|\le\mathrm{rd}\}$ lie in $\mathrm{Tq}$, and define, coefficient by coefficient,
\[
\mathrm{fc}_j(Y;\mathscr{A},\mathscr{C}):=\frac{1}{2\pi i}\oint_{|t-\mathrm{tce}|=\mathrm{rd}}\mathrm{d}t\,f(t)(Y;\mathscr{A},\mathscr{C})\,(t-\mathrm{tce})^{-j-1}\qquad(j\ge0).
\]
By (2) and by (Int) of Section~3.5, the coefficients of $\mathrm{fc}_j$ are analytic in $(\mathscr{A},\mathscr{C})$; this is (k1).
The clauses (k2)--(k5), the evenness and the balance of the legs are complex linear conditions on the values at fixed arguments, and they are stable under limits.
They hold for $f(t)$ at every point $t$ of the circle; hence they hold for $\mathrm{fc}_j$.
The norm of an integral is at most the integral of the norms; therefore $\|\mathrm{fc}_j\|^{x}_{k+1}\le e_k^{\frac14-6\varepsilon}\cdot\mathrm{rd}^{-j}$, and $\mathrm{fc}_j\in K^{x,C}_{k+1}$.
Cauchy's formula, applied to each coefficient, gives $f(t)(Y;\mathscr{A},\mathscr{C})=\sum_{j\ge0}\mathrm{fc}_j(Y;\mathscr{A},\mathscr{C})\,(t-\mathrm{tce})^j$ for $|t-\mathrm{tce}|<\mathrm{rd}$, and
\[
\left\|f(t)-\sum_{j\le j_0}\mathrm{fc}_j\,(t-\mathrm{tce})^j\right\|^{x}_{k+1}\le\sum_{j>j_0}\|\mathrm{fc}_j\|^{x}_{k+1}\,|t-\mathrm{tce}|^j\longrightarrow0\qquad(j_0\to\infty).
\]
Thus, near every interior point of $\mathrm{Tq}$, $f$ is the sum of a power series in $t$ with coefficients in $K^{x,C}_{k+1}$ that converges in $\|\cdot\|^{x}_{k+1}$; that is, $f$ is analytic as a map into $K^{x,C}_{k+1}$.

\emph{The number $\delta\varepsilon^{0}_{k}$.}
By (A8) and (K7), $e_k\le4^{-\frac{1}{\varepsilon}}\le4^{-200}$.
Hence $v_k\ge1$, and $p_{0,k}=v_k^{p_0}\ge1$, because $p_0>3r+2>0$ ((C7) in Section~I.8.2; cf. Proposition~14.2\,(f)).
By (P4) of Section~0.6 and by (A1),
\[
0\le\delta\varepsilon^{0}_{k}\le12\exp\left(-\frac{p_{0,k}^2}{2}\right)\le e_k^7.
\]
\end{proof}

\subsubsection*{9.5. (b)}

\begin{proof}[Proof of Proposition~24.1\,(b)]
\emph{First and second sentence.}
In the case (o) the arguments of Sections~2--9.4 give the following.
For $(m,E)\in\mathrm{Ball}_{k}(4)$ ($E=0$ if $k=0$) and $A\in\tilde{R}_k^{-}$ all displays of the reading (o) are defined, in the order fixed in Definition~24.D, and the series (8.4) converge absolutely in $\|\cdot\|_{h_k/4}$, with the bound (24.14); this is the first sentence (Section~5.4).
Moreover $E^{*}_k(m,E)\in K_{k+1}^{C}$ and $\|E^{*}_k(m,E)\|_{k+1}\le e_k^{\frac14-6\varepsilon}$; the map $t\mapsto E^{*}_k(m_a+tm_b,E_a+tE_b)$ has the continuity and the analyticity stated in (a); and for real $m$ and $E\in K_{k}$ the value lies in $K_{k+1}$.

\emph{Third sentence.}
Let $(m,\mathscr{E})\in\mathrm{Ball}^{x}_{k}(4)$ (the case (x)).
By Proposition~48.1\,(a) (cf. Section~I.5.4), $\mathscr{E}|_0$ lies in $K_{k}^{C}$ and $\|\mathscr{E}|_0\|_k\le\|\mathscr{E}\|^{x}_{k}$; hence $(m,\mathscr{E}|_0)\in\mathrm{Ball}_{k}(4)$.
Proposition~16.3\,(a), (d) gives
\[
m(\mathscr{E})=m(\mathscr{E}|_0),\qquad(R\mathscr{E})|_0=R(\mathscr{E}|_0).
\]
Let $\mathscr{C}=0$, an element of $\mathfrak{X}_k$; then $\mathscr{T}=A$.
According to the rule (SLOT) of Definition~24.D, $\mathscr{C}$ enters the displays (8.1)$^{\mathrm{ch}}$--(8.4) only through $\mathscr{T}=A+\mathscr{C}$ in the fermionic operators and through the second slot of $\mathscr{E}^{r}$.
At $\mathscr{C}=0$ one has
\[
\mathscr{E}^{r}(X;B_1,0)=\left(R(\mathscr{E}|_0)\right)(X;B_1),\qquad m'=m+m(\mathscr{E}|_0).
\]
So each of the displays (8.1)$^{\mathrm{ch}}$--(8.4), formed with the datum $(m,\mathscr{E})$ at $\mathscr{C}=0$, is the same display of the case (o), formed with the datum $(m,\mathscr{E}|_0)$; in particular
\[
E^{\#,x}_k(m,\mathscr{E})(X;A,0)=E^{\#}_k(m,\mathscr{E}|_0)(X;A)\qquad(A\in\tilde{R}_k^{-}).
\]
For $\mathscr{C}=0\in\mathfrak{X}_{k+1}$ one has $\mathscr{C}_L=0$.
Therefore
\[
E^{*,x}_k(m,\mathscr{E})(Y;\mathscr{A},0)=\sum_{X:\,\bar{X}^{L}=LY}S_L\left[E^{\#}_k(m,\mathscr{E}|_0)(X;\mathscr{A}_L)\right]=E^{*}_k(m,\mathscr{E}|_0)(Y;\mathscr{A}),
\]
which is the third sentence of (b).
\end{proof}

\subsubsection*{9.6. (c)}

\begin{proof}[Proof of Proposition~24.1\,(c)]
Here the case (x) is considered.
Let $\varepsilon^{\mathrm{v}}$, $m$, $\mathscr{E}$, $\mathfrak{P}_k$, $A'$, $\mathscr{C}'$ be as in (c).
Then $(m,\mathscr{E})$ lies in $\mathrm{Ball}^{x}_{k}(1)$, which is contained in $\mathrm{Ball}^{x}_{k}(4)$.
The symbols $m'$, $\mathscr{E}^{r}$, $m^{+}$, $\varepsilon^{\mathrm{v},+}$, $\mathscr{A}'$, $\mathscr{T}'$, $\psi'$, $\bar{\psi}'$, $\chi'$, $\bar{\chi}'$, $A$, $\mathscr{C}$, $\mathscr{T}$, $\psi^{0}$, $\bar{\psi}^{0}$, $\chi^{0}$, $\bar{\chi}^{0}$, $\psi^{0,\#}$, $\Psi^{c}$, $\bar{\Psi}^{c}$, $\mathscr{Z}$, $\Psi^{f}$, $\bar{\Psi}^{f}$, $V^{\mathrm{ex},x}_k$, $\Xi^{x}_{k}$ have the meaning fixed in Proposition~22.3.

\emph{(0) Preliminaries.}
By Proposition~22.3\,(a), $A\in\tilde{R}_k^{-}$ and $\mathscr{C}\in\mathfrak{X}_k$; the connections $\mathscr{T}$, $\mathscr{T}+\mathscr{Z}$ and their negatives lie in $\tilde{G}_k$; and all operators that occur in the displays of Proposition~22.3 exist.
So the step map is defined at the background $(A,\mathscr{C})$ (Section~5.4).
By the definition of $\mathrm{Dom}_{k+1}$, $\mathscr{A}'=H_{k+1}A'$ lies in $\frac14\tilde{R}_{k+1}$, a subset of $\tilde{R}_{k+1}$.
By Proposition~18.5\,(a), Proposition~13\,(d5) and Section~9.4 (real data), the family
\[
\mathscr{E}^{+}=LR\mathscr{E}+L\,{\delta E^{\mathrm{det}}}_k+E^{*,x}_k(m,\mathscr{E})
\]
lies in $K^{x}_{k+1}$; this is the first clause of (d).
The data $(k+1,\ell/L,L\mu,e_{k+1})$ are data (SF) of Section~I.1.8, since $\ell/L\in\mathrm{Cl}$, $L\mu\in(0,1]$ and $e_{k+1}\le e_0$.
Hence $\mathscr{F}^{x}_{k+1}[\varepsilon^{\mathrm{v},+},m^{+},\mathscr{E}^{+}]$, formed with $\mathfrak{P}_{k+1}$, is defined (Section~I.5.7).
As $0\in\tilde{R}_k$, DF1 (\cite[Lemma~17.4(a)]{QEDI}) shows that $\mathscr{C}\in\tilde{G}_k$; so $E_k^{\mathrm{det}}(X,\mathscr{C})$, ${\delta E^{\mathrm{det}}}_k(X;A,\mathscr{C})$ and $\mathfrak{P}_{k+1}$ are defined.
By Proposition~22.3\,(b), on the unbarred and on the barred fields,
\begin{equation}\tag{24.21}
\psi^{0,\#}=H_k^{\#}(\mathscr{T})\,\Psi^{c};\quad\text{in particular}\quad\chi^{0}=\delta_{\alpha,\mathscr{T}}\psi^{0},\quad\bar{\chi}^{0}=\delta_{\alpha,-\mathscr{T}}\bar{\psi}^{0}\quad\text{at all pairs of }\mathrm{Pz}_k.
\end{equation}

\emph{Full strength.}
The following identities hold by the properties W1 (at $k=0$: W17) of Sections~G1.11 and~G4.0.4.
The families (p1), (p2) of Section~G4.0.4 satisfy $H_k(1)=H_k$ and $C_{k}^{1/2}(1)=C_{k}^{1/2}$; hence $\mathscr{Z}(1)=H_k\,C\,C_{k}^{1/2}\mathsf{Z}=\mathscr{Z}$.
For $B\in\tilde{G}_k$ one has $H_k(1,B)=H_k(B)$, hence $\bar{H}_k(1,B)=\mathsf{C}^{-1}H_k(-B)\,\mathsf{C}=\bar{H}_k(B)$ and $H_k^{\#}(1,B)=H_k^{\#}(B)$; further $\Gamma_k(1,B)=\Gamma_k(B)$, and, for $k\ge1$, $\sum_XM_k(X,B)=M_k(B)=b_k-D_k(B)$.
With $\Psi^{f}=\Psi^{c}+\Gamma_k(\mathscr{T})W$ and $\bar{\Psi}^{f}=\bar{\Psi}^{c}+\bar{W}$, the identity (24.21) and the linearity of $H_k^{\#}(\mathscr{T})$ give
\begin{equation}\tag{24.22}
H_k^{\#}(\mathscr{T})\,\Psi^{f}=\psi^{0,\#}+W_k^{\#}(1,\mathscr{T}).
\end{equation}

\emph{Three homomorphisms.}
\begin{itemize}
\item $\mathrm{ev}:\mathrm{Gr}[\psi,\chi,W]\to\mathrm{Gr}[\Psi',W]$ is the homomorphism that sends $\psi^{\#}$ to $\psi^{0,\#}$ and fixes $W$.
\item $\mathrm{ev}_1:\mathrm{Gr}[\psi,\chi]\to\mathrm{Gr}[\Psi',W]$ is the homomorphism that sends $\psi^{\#}$ to $H_k^{\#}(\mathscr{T})\Psi^{f}$.
\item For $k\ge1$, $\mathrm{ev}_f:\mathrm{Gr}[\psi,\chi,\Psi]\to\mathrm{Gr}[\Psi',W]$ is the homomorphism that coincides with $\mathrm{ev}_1$ on $\psi$, $\chi$ and sends $\Psi$ to $\Psi^{f}$, $\bar{\Psi}$ to $\bar{\Psi}^{f}$.
For $k=0$ put $\mathrm{ev}_f:=\mathrm{ev}_1$; there $H_0^{\#}(\mathscr{T})\Psi^{f}=\left(\Psi^{f},\delta_{\alpha,\mathscr{T}}\Psi^{f};\ \bar{\Psi}^{f},\delta_{\alpha,-\mathscr{T}}\bar{\Psi}^{f}\right)$.
\end{itemize}
The generators are sent to odd elements with the antiperiodicity (I.AP); hence these homomorphisms exist.
Being linear maps between finite dimensional spaces, they commute with the integrals over $t$, $s$ and $\mathsf{Z}$.
The map $\mathrm{ev}$ commutes also with the integral over $W$: for $F=\sum_Jc_JW(J)$ with $c_J\in\mathrm{Gr}[\psi,\chi]$, and with the moments $m(J)$ of Proposition~15\,(e),
\[
\int\mathrm{d}\mu_I(W)\,\mathrm{ev}(F)=\sum_J\mathrm{ev}(c_J)\,m(J)=\mathrm{ev}\left(\int\mathrm{d}\mu_I(W)\,F\right).
\]
By (24.22),
\begin{equation}\tag{24.23}
\mathrm{ev}\circ\mathrm{sb}_v=\mathrm{ev}_1\quad\text{on }\mathrm{Gr}[\psi,\chi],\qquad v:=W_k^{\#}(1,\mathscr{T}).
\end{equation}

\emph{(1)}
The following identity holds:
\[
\mathrm{ev}_f\left(\sum_X\sum_{i\le3}G^{(i)}(1,X)\right)=\mathrm{Act}^{x}_{k}[m',\mathscr{E}^{r}](A+\mathscr{Z},\mathscr{C};\Psi^{f})-\mathrm{Act}^{x}_{k}[m',\mathscr{E}^{r}](A,\mathscr{C};\Psi^{f}).
\]

\emph{The case $k\ge1$.}
At $s=1$ the integrand of $G^{(i)}(1,X)$ is analytic in $t$ in a neighbourhood of the closed disk bounded by the circle of integration of $G^{(i)}$ in (D4) of Definition~24.D, that is, of $\{|t|\le\mathrm{rad}_1\}$ for $i=1$ and of $\{|t|\le\mathrm{rad}_2\}$ for $i=2,3$.
For the entries of the operators this is Proposition~24.3\,(q5), third assertion.
For $\mathscr{E}^{r}$ it follows from (k1) together with (q1) and DF0 (\cite[Lemma~17.3]{QEDI}): the compact disk $\{|t|\le\mathrm{rad}_2\}$, where $\mathrm{rad}_2\le\mathrm{rad}_1$, is mapped into the open set $\tilde{R}_k$.
By (N5) (Proposition~15\,(j)), applied with $\mathrm{rad}=\mathrm{rad}_1$ for $i=1$ and with $\mathrm{rad}=\mathrm{rad}_2$ for $i=2,3$, one has $\oint_{\mathrm{rad}}f=f(1)-f(0)$ for each of these integrands $f$.

For $i=1$,
\[
\sum_XG^{(1)}(1,X)=\left\langle\bar{\Psi},\left(M_k(\mathscr{T}+\mathscr{Z})-M_k(\mathscr{T})\right)\Psi\right\rangle_1=-\left\langle\bar{\Psi},\left(D_k(\mathscr{T}+\mathscr{Z})-D_k(\mathscr{T})\right)\Psi\right\rangle_1,
\]
and $\mathrm{ev}_f$ sends this element to $-\left\langle\bar{\Psi}^{f},\left(D_k(\mathscr{T}+\mathscr{Z})-D_k(\mathscr{T})\right)\Psi^{f}\right\rangle_1$.

For $i=2,3$, the homomorphism $\mathrm{ev}_f$ sends $\psi^{\#}+J_k^{\#}(1,\mathscr{T},\mathscr{Z})\Psi$ to
\[
H_k^{\#}(\mathscr{T})\Psi^{f}+\left(H_k^{\#}(\mathscr{T}+\mathscr{Z})-H_k^{\#}(\mathscr{T})\right)\Psi^{f}=H_k^{\#}(\mathscr{T}+\mathscr{Z})\Psi^{f}.
\]
Consequently
\[
\mathrm{ev}_f\sum_XG^{(2)}(1,X)=\sum_X\left[\mathscr{E}^{r}\left(X;A+\mathscr{Z},\mathscr{C},H_k^{\#}(\mathscr{T}+\mathscr{Z})\Psi^{f}\right)-\mathscr{E}^{r}\left(X;A,\mathscr{C},H_k^{\#}(\mathscr{T})\Psi^{f}\right)\right]
\]
and
\[
\mathrm{ev}_f\sum_{\square}G^{(3)}(1,\square)=-m'\,\langle\bar{\psi},\psi\rangle\big|_{\psi=H_k(\mathscr{T}+\mathscr{Z})\Psi^{f},\ \bar{\psi}=\bar{H}_k(\mathscr{T}+\mathscr{Z})\bar{\Psi}^{f}}+m'\,\langle\bar{\psi},\psi\rangle\big|_{\psi=H_k(\mathscr{T})\Psi^{f},\ \bar{\psi}=\bar{H}_k(\mathscr{T})\bar{\Psi}^{f}};
\]
here the form $\langle\bar{\psi},\psi\rangle$ of Section~I.3.6 is evaluated at the fields indicated, and $\sum_{\square}\langle\bar{\psi},\psi\rangle_{\square}=\langle\bar{\psi},\psi\rangle$, because the sets $\square\times O_{\mathrm{cub}}$ partition $T^{\mathrm{f}}_k$ (Section~I.1.4).
The total connection of the pair $(A+\mathscr{Z},\mathscr{C})$ is $\mathscr{T}+\mathscr{Z}$.
By the definitions of $\mathrm{Act}^{x}_{k}$ and $I^{x}_{k}$ (Sections~I.5.7 and G3.2.0) the sum of the three contributions is the difference asserted in (1).

\emph{The case $k=0$.}
Here $G^{(2)}=G^{(3)}=0$, $\mathscr{E}^{r}=0$, $m'=m$ and $G^{(1)}(1,\square)=\mathrm{Vx}(\mathscr{Z},\square)$.
Proposition~23.2\,(a), applied with $\mathscr{Z}$, $\Psi^{f}$, $\bar{\Psi}^{f}$ in the places of $\mathrm{Zf}$, $\mathrm{ff}$, $\overline{\mathrm{ff}}$, gives
\[
\mathrm{ev}_f\left(\sum_{\square}\mathrm{Vx}(\mathscr{Z},\square)\right)=-\left\langle\bar{\Psi}^{f},\left(D_{\mathscr{T}+\mathscr{Z}}-D_{\mathscr{T}}\right)\Psi^{f}\right\rangle^{\mathrm{ch}}.
\]
On the other side, $D_0(B)=D_B+\mu$ and $H_0=\bar{H}_0=1$; therefore
\[
\mathrm{Act}^{x}_{0}[m,0](A+\mathscr{Z},\mathscr{C};\Psi^{f})-\mathrm{Act}^{x}_{0}[m,0](A,\mathscr{C};\Psi^{f})=-\left\langle\bar{\Psi}^{f},\left(D_{\mathscr{T}+\mathscr{Z}}-D_{\mathscr{T}}\right)\Psi^{f}\right\rangle^{\mathrm{ch}},
\]
since the term $-m\langle\bar{\Psi}^{f},\Psi^{f}\rangle^{\mathrm{ch}}$ does not contain the connection.
All these pairings are formed on $T^{\mathrm{f}}_0$ with the weight $1/48$.

\emph{(2)}
For $k\ge1$ and every polymer $\mathbf{Y}$,
\[
\mathrm{ev}_1\circ\mathrm{sub}_{\mathbf{Y},\mathscr{T}}=\mathrm{ev}_f\quad\text{on }\mathrm{Gr}[\psi,\chi,\Psi|\mathbf{Y}].
\]
Indeed, $\mathrm{ev}_1(\chi)=\delta_{\alpha,\mathscr{T}}\mathrm{ev}_1(\psi)$, $\mathrm{ev}_1(\bar{\chi})=\delta_{\alpha,-\mathscr{T}}\mathrm{ev}_1(\bar{\psi})$, $\mathrm{ev}_1(\psi)=H_k(\mathscr{T})\Psi^{f}$ and $\mathrm{ev}_1(\bar{\psi})=\bar{H}_k(\mathscr{T})\bar{\Psi}^{f}$, and $S_k(\mathscr{T})$ exists.
By Proposition~23.1\,(e), $\mathrm{ev}_1\circ\mathrm{sub}_{\mathbf{Y},\mathscr{T}}$ coincides with $\mathrm{ev}_1$ on $\psi$, $\chi$ and sends $\Psi(y)$ to $\Psi^{f}(y)$ and $\bar{\Psi}(y)$ to $\bar{\Psi}^{f}(y)$ at the unit sites $y$ of $\mathbf{Y}$; so the two homomorphisms agree on the generators.

\emph{(3)}
The identity $\mathrm{ev}\left(\sum_{\mathbf{Y}}E^{(\le5),x}(\mathbf{Y})\right)=V^{\mathrm{ex},x}_k$ holds.
By Proposition~5.B\,(d) the left side equals $\mathrm{ev}\sum_X\left(G^{(\le3)}+G^{(4)}+G^{(5)}\right)(1,X)$.
By (24.23), and then by (2) and (24.5) (at $k=0$: $E^{(\le3),x}=E^{(1),*}$ and $\mathrm{ev}_1=\mathrm{ev}_f$),
\[
\mathrm{ev}\,G^{(\le3)}(1,X)=\mathrm{ev}\,\mathrm{sb}_v\left(E^{(\le3),x}(X)\right)=\mathrm{ev}_1\left(\sum_i\mathrm{sub}_{X,\mathscr{T}}\left(E^{(i),*}(X)\right)\right)=\mathrm{ev}_f\left(\sum_{i\le3}E^{(i),*}(X)\right).
\]
By Proposition~5.B\,(d) and (1), the sum of these elements over $X$ equals $\mathrm{Act}^{x}_{k}[m',\mathscr{E}^{r}](A+\mathscr{Z},\mathscr{C};\Psi^{f})-\mathrm{Act}^{x}_{k}[m',\mathscr{E}^{r}](A,\mathscr{C};\Psi^{f})$.
The integrands of $G^{(4)}$ and $G^{(5)}$ are polynomials in $t$; by (N5), (24.23) and (24.21),
\begin{align*}
\mathrm{ev}\sum_X\left(G^{(4)}+G^{(5)}\right)(1,X)
&=\sum_X\left[\mathscr{E}^{r}\left(X;A,\mathscr{C},H_k^{\#}(\mathscr{T})\Psi^{f}\right)-\mathscr{E}^{r}\left(X;A,\mathscr{C},H_k^{\#}(\mathscr{T})\Psi^{c}\right)\right]\\
&\quad-m'\,\langle\bar{\psi},\psi\rangle\big|_{\psi=H_k(\mathscr{T})\Psi^{f},\ \bar{\psi}=\bar{H}_k(\mathscr{T})\bar{\Psi}^{f}}+m'\,\langle\bar{\psi},\psi\rangle\big|_{\psi=H_k(\mathscr{T})\Psi^{c},\ \bar{\psi}=\bar{H}_k(\mathscr{T})\bar{\Psi}^{c}}\\
&=I^{x}_{k}[m',\mathscr{E}^{r}](A,\mathscr{C};\Psi^{f})-I^{x}_{k}[m',\mathscr{E}^{r}](A,\mathscr{C};\Psi^{c}).
\end{align*}
The sum of the two parts is $V^{\mathrm{ex},x}_k$.
(At $k=0$ the two parts are $\mathrm{Ev}^{(1)}$ and $\mathrm{Ev}^{(5)}=-m\left[\langle\bar{\Psi}^{f},\Psi^{f}\rangle^{\mathrm{ch}}-\langle\bar{\Psi}^{c},\Psi^{c}\rangle^{\mathrm{ch}}\right]$ of Proposition~22.4\,(a).)

\emph{(4)}
Apply $\mathrm{ev}$ to (24.15).
By (0) and (3),
\[
\Xi^{x}_{k}=\int\mathrm{d}\mu_I(W)\,\mathrm{d}\hat{\mu}(\mathsf{Z})\,\exp\left(V^{\mathrm{ex},x}_k\right)=\exp\left\{\sum_X\mathrm{ev}\left(E^{\#,x}_k(m,\mathscr{E})(X;A,\mathscr{C})\right)\right\}.
\]

\emph{(5) The other factors.}
By Proposition~13\,(d0) at $B=\mathscr{T}$ and by the definitions of Section~I.5.7,
\[
\frac{\delta Z_k(\mathscr{T})}{\delta Z_k(0)}=\exp\left\{\sum_X\left(E^{d}_k(X,\mathscr{T})-E^{d}_k(X,0)\right)\right\}=\exp\left\{\sum_X{\delta E^{\mathrm{det}}}_k(X;A,\mathscr{C})+\sum_XE_k^{\mathrm{det}}(X,\mathscr{C})\right\}.
\]
The family $R\mathscr{E}$ has (k2), by (X2).
By (24.21), $\mathrm{ev}(\chi)=\delta_{\alpha,\mathscr{T}}\mathrm{ev}(\psi)$ and $\mathrm{ev}(\bar{\chi})=\delta_{\alpha,-\mathscr{T}}\mathrm{ev}(\bar{\psi})$ at all pairs.
Proposition~17.2\,(a) for small $X$, and the relation $E^{\mathrm{nat}}(X)=E(X)$ for large $X$, give
\[
\mathrm{ev}\left(\mathscr{E}^{r}(X;A,\mathscr{C})\right)=\mathrm{ev}\left((R\mathscr{E})^{\mathrm{nat}}(X;A,\mathscr{C})\right).
\]
Let $\mathrm{ev}'$ be the homomorphism from $\mathrm{Gr}[\psi,\chi]$ of level $k+1$ to $\mathrm{Gr}[\Psi']$ with $\psi\mapsto\psi'$, $\chi\mapsto\chi'$, $\bar{\psi}\mapsto\bar{\psi}'$, $\bar{\chi}\mapsto\bar{\chi}'$.
Then $\mathrm{ev}=\mathrm{ev}'\circ S_L$ on $\mathrm{Gr}[\psi,\chi]$ of level $k$, since
\[
\mathrm{ev}'\left(S_L\psi(x)\right)=L^{-1}\psi'(x/L)=\psi'_L(x)=\psi^{0}(x),\qquad\mathrm{ev}'\left(S_L\chi(x,x')\right)=L^{-1-\alpha}\chi'(x/L,x'/L)=\chi^{0}(x,x'),
\]
and in the same way for the barred generators.
With $A=\mathscr{A}'_L$ and $\mathscr{C}=\mathscr{C}'_L$, (Sc1) and (8.6)$^{\mathrm{ch}}$ give
\begin{align*}
\sum_X\mathrm{ev}\left[(R\mathscr{E})^{\mathrm{nat}}(X;A,\mathscr{C})+{\delta E^{\mathrm{det}}}_k(X;A,\mathscr{C})+E^{\#,x}_k(m,\mathscr{E})(X;A,\mathscr{C})\right]
&=\mathrm{ev}'\left(\sum_{Y\in D_{k+1}}\mathscr{E}^{+}(Y;\mathscr{A}',\mathscr{C}')\right)\\
&=\sum_Y\mathscr{E}^{+}(Y;\mathscr{A}',\mathscr{C}',\psi',\chi'),\\
\sum_X\left(\mathfrak{P}_k(X,\mathscr{C})+E_k^{\mathrm{det}}(X,\mathscr{C})\right)&=\sum_Y\mathfrak{P}_{k+1}(Y,\mathscr{C}').
\end{align*}

\emph{(6) The identity (CV.4x)$^{\mathrm{ch}}$.}
Insert (4) and (5) into the identity (2.4a$^{x}$)$^{\mathrm{ch}}$ of Proposition~22.3\,(c); the exponents are even and commute.
In the exponent the three sums $\sum_X\mathscr{E}^{r}(X;A,\mathscr{C},\psi^{0,\#})=\sum_X\mathrm{ev}\left(\mathscr{E}^{r}(X;A,\mathscr{C})\right)$, $\sum_X{\delta E^{\mathrm{det}}}_k(X;A,\mathscr{C})$ (from the factor $\delta Z_k(\mathscr{T})/\delta Z_k(0)$) and $\sum_X\mathrm{ev}\left(E^{\#,x}_k(m,\mathscr{E})(X;A,\mathscr{C})\right)$ (from $\Xi^{x}_{k}$) add up, by (5), to $\sum_Y\mathscr{E}^{+}(Y;\mathscr{A}',\mathscr{C}',\psi',\chi')$.
Likewise $\sum_X\mathfrak{P}_k(X,\mathscr{C})+\sum_XE_k^{\mathrm{det}}(X,\mathscr{C})=\sum_Y\mathfrak{P}_{k+1}(Y,\mathscr{C}')$.
This gives
\begin{multline*}
\mathbf{T}^{x}_{k}\mathscr{F}^{x}_{k}[\varepsilon^{\mathrm{v}},m,\mathscr{E}](A',\mathscr{C}';\Psi')
=\det D_{k+1}(0)^{-1}\exp\Big\{-\left\langle\bar{\Psi}',D_{k+1}(\mathscr{T}')\Psi'\right\rangle_1-m^{+}\langle\bar{\psi}',\psi'\rangle\\
+\sum_Y\mathscr{E}^{+}(Y;\mathscr{A}',\mathscr{C}',\psi',\chi')-\varepsilon^{\mathrm{v},+}\cdot n_{k+1}+\sum_Y\mathfrak{P}_{k+1}(Y,\mathscr{C}')\Big\}
=\mathscr{F}^{x}_{k+1}[\varepsilon^{\mathrm{v},+},m^{+},\mathscr{E}^{+}](A',\mathscr{C}';\Psi'),
\end{multline*}
by the definitions of $I^{x}_{k+1}$, $\mathrm{Act}^{x}_{k+1}$ and $\mathscr{F}^{x}_{k+1}$ (Section~I.5.7), in which $\psi'=H_{k+1}(\mathscr{T}')\Psi'$, $\chi'$, $\bar{\psi}'$, $\bar{\chi}'$ are the fields to be inserted.

\emph{Level $0$.}
Here $\mathscr{E}=0$, $m'=m$, $\varepsilon(\mathscr{E})=m(\mathscr{E})=0$ and $LR\mathscr{E}=0$, so that
\[
\left(\varepsilon^{\mathrm{v},+},m^{+},\mathscr{E}^{+}\right)=\left(L^3\left(\varepsilon^{\mathrm{v}}+\delta\varepsilon^{0}_{0}\right),\ Lm,\ L\,{\delta E^{\mathrm{det}}}_0+E^{*,x}_0(m,0)\right).
\]
In $\mathbf{T}^{x}_{0}$ the integral $\int\mathrm{D}\Psi$ extends over the generators $\Psi_\beta(x)$, $\bar{\Psi}_\beta(x)$, $x\in T^{\mathrm{f}}_0=T^0_0\times O_{\mathrm{cub}}$, that is, over $4\cdot48\,n_0=192\,n_0$ pairs, with the weight $1/48$ (Section~I.3.6, (I.L0.3), (I.L0.5)).
The right side lies in $\mathrm{Gr}[\Psi']$, where $\Psi'$ is indexed by $T^0_1$, with $4\,n_1$ pairs of generators.

\emph{(7) The identity (2.4x)$^{\mathrm{ch}}$.}
Integrate (CV.4x)$^{\mathrm{ch}}$ with $\int\mathrm{D}\Psi'$.
On the right side this gives $y^{x}_{k+1}[\varepsilon^{\mathrm{v},+},m^{+},\mathscr{E}^{+}](A';\mathscr{C}')$.
To the left side Proposition~22.2\,(b) is applied, with $F:=\mathscr{F}^{x}_{k}[\varepsilon^{\mathrm{v}},m,\mathscr{E}](\,\cdot\,,\mathscr{C};\,\cdot\,)$ and $\mathscr{C}=\mathscr{C}'_L$.
The values of $F$ are even.
For $\zeta\in\mathrm{Box}_{k}$ the field $A_z:=A^{\min}_{k}(A')+\mathrm{Sq}_{k}\zeta$ lies in $\mathrm{Dom}_k$ by (P5) of Section~0.6 (\cite[Proposition~15.30(ii)]{QEDI}).
Indeed, $\ell$ and $\ell/L$ lie in $\mathrm{Cl}$, and $\mathrm{Dom}_{k+1}$ is formed with the side $\ell/L$ and with $e_{k+1}=L^{\frac12}e_k$ ((ST1)$^{\mathrm{ch}}$); the conditions on the constants named in (P5) are (A10) (third and fourth inequality), (A8), (A1) (first inequality), (A2), (K1), (K7), (K8) of Section~0.3, at $v=-\log e_k$.
Hence $H_kA_z\in\frac14\tilde{R}_k\subset\tilde{R}_k$.
Since $H_k=H_k(1)$ is a linear operator (W1; at $k=0$: $H_0=1$), the map $\zeta\mapsto H_kA_z$ is affine, hence continuous.
The coefficients of $\mathrm{Act}^{x}_{k}[m,\mathscr{E}](\mathscr{A},\mathscr{C};\Psi)$ are continuous functions of $\mathscr{A}$ on $\tilde{R}_k$:
\begin{itemize}
\item for $k\ge1$, $D_k(\mathscr{T})=b_k-\sum_XM_k(X,\mathscr{T})$ (W1), and its entries are analytic on $\tilde{G}_k$ (W2);
\item for $k=0$, the entries of $D_0(\mathscr{T})=D_{\mathscr{T}}+\mu$ are entire functions (I.O1);
\item $I^{x}_{k}[m,\mathscr{E}](\mathscr{A},\mathscr{C};\Psi)$ is the image of $-m\langle\bar{\psi},\psi\rangle+\sum_X\mathscr{E}(X;\mathscr{A},\mathscr{C})$ under $\psi^{\#}\mapsto H_k^{\#}(\mathscr{T})\Psi$; so its coefficients are polynomials in the coefficients of the elements $\mathscr{E}(X;\mathscr{A},\mathscr{C})$, which are analytic by (k1), and in the entries of $H_k^{\#}(\mathscr{T})=H_k^{\#}(1,\mathscr{T})$, which are analytic on $\tilde{G}_k$ (W1, W2, (24.3)).
\end{itemize}
The coefficients of the exponential of an even element are entire functions of its coefficients (Section~3.8).
So the coefficients of $\zeta\mapsto F(A_z,\mathscr{C};\Psi)$ are continuous on the compact set $\mathrm{Box}_{k}$, hence integrable, and $F$ belongs to the class of Section~G3.2.0.
Proposition~22.2\,(b) now gives
\begin{align*}
\int\mathrm{D}\Psi'\,(\mathbf{T}^{x}_{k}F)(A',\mathscr{C}';\Psi')
&=\int_{\mathrm{Box}_{k}}\mathrm{d}\mu_I(\zeta)\int\mathrm{D}\Psi\,F\left(A^{\min}_{k}(A')+\mathrm{Sq}_{k}\zeta,\mathscr{C}'_L;\Psi\right)\\
&=\int_{\mathrm{Box}_{k}}\mathrm{d}\mu_I(\zeta)\,y^{x}_{k}[\varepsilon^{\mathrm{v}},m,\mathscr{E}]\left(A^{\min}_{k}(A')+\mathrm{Sq}_{k}\zeta;\mathscr{C}'_L\right).
\end{align*}
Together with the right side this is (2.4x)$^{\mathrm{ch}}$.

\emph{(8) The case $\mathscr{C}'=0$.}
The objects $\mathscr{F}_j$, $y_j$ ($j=k,k+1$) and $\mathbf{T}_k$ are the objects with slot, taken at the value $0$ of the second slot (Section~I.5.7).
Let $\mathfrak{P}_k(\,\cdot\,,0)=0$ and $\mathscr{C}'=0$; then $\mathscr{C}=0$.
By Section~I.5.7, $\mathscr{F}^{x}_{k}[\varepsilon^{\mathrm{v}},m,\mathscr{E}](\,\cdot\,,0;\,\cdot\,)$ contains $\mathscr{E}$ only through $\mathscr{E}(\,\cdot\,;\,\cdot\,,0,\,\cdot\,)=\mathscr{E}|_0$, and $\mathscr{F}^{x}_{k+1}[\varepsilon^{\mathrm{v},+},m^{+},\mathscr{E}^{+}](\,\cdot\,,0;\,\cdot\,)$ contains $\mathscr{E}^{+}$ only through $\mathscr{E}^{+}|_0$.
By Proposition~16.3\,(a), $\varepsilon(\mathscr{E})=\varepsilon(\mathscr{E}|_0)$ and $m(\mathscr{E})=m(\mathscr{E}|_0)$; so $(\varepsilon^{\mathrm{v},+},m^{+})$ is formed from $(\varepsilon^{\mathrm{v}},m,\mathscr{E}|_0)$.
Furthermore $E_k^{\mathrm{det}}(X,0)=E^{d}_k(X,0)-E^{d}_k(X,0)=0$, and hence
\[
\mathfrak{P}_{k+1}(Y,0)=\sum_{X:\,\bar{X}^{L}=LY}\left(\mathfrak{P}_k(X,0)+E_k^{\mathrm{det}}(X,0)\right)=0.
\]
Consequently (CV.4x)$^{\mathrm{ch}}$ and (2.4x)$^{\mathrm{ch}}$ at $\mathscr{C}'=0$ are the two identities for $\mathscr{F}_k$, $y_k$, $\mathbf{T}_k$ with the data $(\varepsilon^{\mathrm{v}},m,\mathscr{E}|_0)$ and $(\varepsilon^{\mathrm{v},+},m^{+},\mathscr{E}^{+}|_0)$.

The argument uses no property of the harmonic part of $A'$; so the identities hold for every harmonic part of $A'$.
\end{proof}

\subsubsection*{9.7. (d)}

\begin{proof}[Proof of Proposition~24.1\,(d)]
Let the data be those of (c).
In particular $m$ is real, $\mathscr{E}$ lies in $K^{x}_{k}$, and $(m,\mathscr{E})$ lies in $\mathrm{Ball}^{x}_{k}(1)$.
By Proposition~18.5\,(a) the family $LR\mathscr{E}$ lies in $K^{x}_{k+1}$;
by Proposition~13\,(d5) the family $L\,{\delta E^{\mathrm{det}}}_k$ lies in $K^{x}_{k+1}$;
by 9.4, applied to the real datum $(m,\mathscr{E})\in\mathrm{Ball}^{x}_{k}(1)$, the family $E^{*,x}_k(m,\mathscr{E})$ lies in $K^{x}_{k+1}$.
By (8.6)$^{\mathrm{ch}}$ the family $\mathscr{E}^{+}$ is the sum of these three.
Hence $\mathscr{E}^{+}$ lies in $K^{x}_{k+1}$, and Proposition~18.5\,(b), Proposition~13\,(d6) and 9.4 give
\[
\left\|\mathscr{E}^{+}\right\|^{x}_{k+1}\le c_{\mathrm{fl}}\cdot L^{-\frac{1}{4}+2\varepsilon}\cdot\left\|\mathscr{E}\right\|^{x}_{k}+e_{k+1}^{\frac{1}{4}-2\varepsilon}+e_k^{\frac{1}{4}-6\varepsilon}.
\]

\emph{Restriction to $\mathscr{C}=0$.}
By Proposition~18.5\,(c) one has $\left(LR\mathscr{E}\right)|_0=LR\left(\mathscr{E}|_0\right)$.
For $\mathscr{C}=0$ also $\mathscr{C}_L=0$, and
\[
{\delta E^{\mathrm{det}}}_k(X;\mathscr{A}_L,0)=E^{d}_k(X,\mathscr{A}_L)-E^{d}_k(X,0)=E_k^{\mathrm{det}}(X,\mathscr{A}_L);
\]
consequently $\left(L\,{\delta E^{\mathrm{det}}}_k\right)|_0=L\,E_k^{\mathrm{det}}$.
By 9.5, $E^{*,x}_k(m,\mathscr{E})|_0=E^{*}_k\left(m,\mathscr{E}|_0\right)$.
The three relations together give $\mathscr{E}^{+}|_0=LR\left(\mathscr{E}|_0\right)+L\,E_k^{\mathrm{det}}+E^{*}_k\left(m,\mathscr{E}|_0\right)$.

\emph{The couplings.}
By (8.6)$^{\mathrm{ch}}$,
\[
m^{+}-L\cdot m=L\cdot m(\mathscr{E}),\qquad
\varepsilon^{\mathrm{v},+}-L^3\cdot\varepsilon^{\mathrm{v}}=L^3\left(\varepsilon(\mathscr{E})+\delta\varepsilon^{0}_{k}\right).
\]
From (X1) and from $0\le\delta\varepsilon^{0}_{k}\le e_k^{7}$ (9.4) it follows that
\[
\left|m^{+}-L\cdot m\right|\le L\cdot c_{\mathrm{fl}}\cdot e_k^{\frac{1}{2}}\cdot\left\|\mathscr{E}\right\|^{x}_{k},\qquad
\left|\varepsilon^{\mathrm{v},+}-L^3\cdot\varepsilon^{\mathrm{v}}\right|\le L^3\left(c_{\mathrm{fl}}\cdot\left\|\mathscr{E}\right\|^{x}_{k}+e_k^{7}\right).
\]

\emph{The case $k=0$.}
Here $\mathscr{E}=0$, so that $LR\mathscr{E}=0$, $m(\mathscr{E})=0$ and $\varepsilon(\mathscr{E})=0$;
moreover $e_k=q$ and $e_{k+1}=L^{\frac{1}{2}}\cdot q$.
The relations above become
\[
\left\|\mathscr{E}^{+}\right\|^{x}_{1}\le\left(L^{\frac{1}{2}}\cdot q\right)^{\frac{1}{4}-2\varepsilon}+q^{\frac{1}{4}-6\varepsilon},\qquad
m^{+}=L\cdot m,\qquad
\varepsilon^{\mathrm{v},+}-L^3\cdot\varepsilon^{\mathrm{v}}=L^3\cdot\delta\varepsilon^{0}_{0}\in\left[0,L^3\cdot q^{7}\right].
\]
\end{proof}

\subsubsection*{9.8. (e)}

\begin{proof}[Proof of Proposition~24.1\,(e)]
Let $(m_i,\mathscr{E}_i)\in\mathrm{Ball}^{x}_{k}(1)$ for $i=1,2$.
Write the difference $E^{*,x}_k(m_1,\mathscr{E}_1)-E^{*,x}_k(m_2,\mathscr{E}_2)$ as the sum of two brackets:
\[
\left[E^{*,x}_k(m_1,\mathscr{E}_1)-E^{*,x}_k(m_2,\mathscr{E}_1)\right]+\left[E^{*,x}_k(m_2,\mathscr{E}_1)-E^{*,x}_k(m_2,\mathscr{E}_2)\right].
\]

\emph{The first bracket.}
Let $m_1\neq m_2$, and put
\[
f(t):=E^{*,x}_k\left(m_2+t\cdot(m_1-m_2),\mathscr{E}_1\right),\qquad
\mathrm{rd}:=\frac{3\,w^{\mathrm{m}}_k}{|m_1-m_2|};
\]
since $|m_1-m_2|\le 2\,w^{\mathrm{m}}_k$, one has $\mathrm{rd}\ge\frac{3}{2}$.
For $|t|\le\mathrm{rd}$,
\[
\left|m_2+t\cdot(m_1-m_2)\right|\le w^{\mathrm{m}}_k+3\,w^{\mathrm{m}}_k=4\,w^{\mathrm{m}}_k;
\]
so the argument of $f(t)$ lies in $\mathrm{Ball}^{x}_{k}(4)$, $\left\|f(t)\right\|^{x}_{k+1}\le e_k^{\frac{1}{4}-6\varepsilon}$,
and every coefficient of $f(t)(Y;\mathscr{A},\mathscr{C})$ is continuous on $|t|\le\mathrm{rd}$ and analytic in $|t|<\mathrm{rd}$ (9.4).
Let $1<\mathrm{rd}'<\mathrm{rd}$.
For every $Y$, every $(\mathscr{A},\mathscr{C})$ and every coefficient, (N5) gives $f(1)-f(0)=\oint_{\mathrm{rd}'}f$;
and the norm $\|\cdot\|_{h_{k+1}}$ of an integral is at most the integral of the norms.
Therefore $\left\|f(1)-f(0)\right\|^{x}_{k+1}\le e_k^{\frac{1}{4}-6\varepsilon}/(\mathrm{rd}'-1)$.
Letting $\mathrm{rd}'$ tend to $\mathrm{rd}$ and using $\mathrm{rd}-1\ge\mathrm{rd}/3$, one obtains
\[
\left\|f(1)-f(0)\right\|^{x}_{k+1}\le\frac{3\,e_k^{\frac{1}{4}-6\varepsilon}}{\mathrm{rd}}=e_k^{\frac{1}{4}-6\varepsilon}\cdot\frac{|m_1-m_2|}{w^{\mathrm{m}}_k}.
\]
Since $f(1)$ and $f(0)$ are the two terms of the first bracket, this bounds the first bracket.

\emph{The second bracket} ($k\ge1$).
Let $\mathscr{E}_1\neq\mathscr{E}_2$, and put
\[
f(t):=E^{*,x}_k\left(m_2,\mathscr{E}_2+t\cdot(\mathscr{E}_1-\mathscr{E}_2)\right),\qquad
\mathrm{rd}:=\frac{3\,w^{\mathrm{E}}_k}{\left\|\mathscr{E}_1-\mathscr{E}_2\right\|^{x}_{k}}\ge\frac{3}{2}.
\]
For $|t|\le\mathrm{rd}$,
\[
\left\|\mathscr{E}_2+t\cdot(\mathscr{E}_1-\mathscr{E}_2)\right\|^{x}_{k}\le w^{\mathrm{E}}_k+3\,w^{\mathrm{E}}_k=4\,w^{\mathrm{E}}_k,
\]
and $|m_2|\le w^{\mathrm{m}}_k$;
so the argument lies in $\mathrm{Ball}^{x}_{k}(4)$, $\left\|f(t)\right\|^{x}_{k+1}\le e_k^{\frac{1}{4}-6\varepsilon}$,
and the coefficients of $f(t)$ are continuous on $|t|\le\mathrm{rd}$ and analytic in the open disk (9.4).
Applying (N5) on the circles $|t|=\mathrm{rd}'<\mathrm{rd}$ and letting $\mathrm{rd}'$ tend to $\mathrm{rd}$, one obtains
\[
\left\|f(1)-f(0)\right\|^{x}_{k+1}\le\frac{e_k^{\frac{1}{4}-6\varepsilon}}{\mathrm{rd}-1}\le\frac{3\,e_k^{\frac{1}{4}-6\varepsilon}}{\mathrm{rd}}=e_k^{\frac{1}{4}-6\varepsilon}\cdot\frac{\left\|\mathscr{E}_1-\mathscr{E}_2\right\|^{x}_{k}}{w^{\mathrm{E}}_k},
\]
and $f(1)$, $f(0)$ are the two terms of the second bracket.
At $k=0$ one has $\mathscr{E}_1=\mathscr{E}_2=0$, and the second bracket vanishes.

The sum of the two bounds is the inequality asserted in (e).
In the case (o) the same argument is the proof of the bound for $E^{*}_k$, for $(m_i,E_i)\in\mathrm{Ball}_k(1)$, with the norms $\|\cdot\|_k$ and $\|\cdot\|_{k+1}$.
\end{proof}

\subsubsection*{9.9. (i)}

\begin{proof}[Proof of Proposition~24.1\,(i)]
The bounds asserted in (a)--(e) are the following:
\begin{gather*}
14\,c_X^2\cdot K_0^2\cdot\exp(3\kappa)\cdot\mathbf{H}_{0}\cdot\exp\left(-\left(\kappa-5\kappa_0-4\right)d_M(X)\right),\qquad
\mathbf{H}_{0}=36\,c_X^2\cdot K_0^2\cdot e_k^{\frac{1}{4}-5\varepsilon};\\
e_k^{\frac{1}{4}-6\varepsilon};\qquad e_k^{7};\\
c_{\mathrm{fl}}\cdot L^{-\frac{1}{4}+2\varepsilon}\cdot\left\|\mathscr{E}\right\|^{x}_{k}+e_{k+1}^{\frac{1}{4}-2\varepsilon}+e_k^{\frac{1}{4}-6\varepsilon};\qquad
L\cdot c_{\mathrm{fl}}\cdot e_k^{\frac{1}{2}}\cdot\left\|\mathscr{E}\right\|^{x}_{k};\qquad
L^3\left(c_{\mathrm{fl}}\cdot\left\|\mathscr{E}\right\|^{x}_{k}+e_k^{7}\right);\\
e_k^{\frac{1}{4}-6\varepsilon}\left(\frac{|m_1-m_2|}{w^{\mathrm{m}}_k}+\frac{\left\|\mathscr{E}_1-\mathscr{E}_2\right\|^{x}_{k}}{w^{\mathrm{E}}_k}\right).
\end{gather*}
Apart from numerals, the norms of the data and $d_M(X)$, they contain $e_k$ (also through $e_{k+1}=L^{\frac{1}{2}}\cdot e_k$, $w^{\mathrm{m}}_k=e_k^{\frac{3}{4}-8\varepsilon}$, $w^{\mathrm{E}}_k=e_k^{\frac{1}{4}-7\varepsilon}$ and $\mathbf{H}_{0}$),
the constants $L$, $c_{\mathrm{fl}}$, $c_X$, $K_0$, $\kappa$, $\kappa_0$, and $\varepsilon$ in exponents.

These bounds were obtained in Sections~2--9 for arbitrary data $(k,\ell,\mu,e_k)$ of (ST1)$^{\mathrm{ch}}$,
an arbitrary system with W1--W8, W10--W17 for $\mathrm{par}$,
arbitrary $(A,\mathscr{C})\in\tilde{R}_k^{-}\times\mathrm{Sd}_k$, resp.\ $(\mathscr{A},\mathscr{C})\in\tilde{R}_{k+1}\times\mathrm{Sd}_{k+1}$,
an arbitrary $A'\in\mathrm{Dom}_{k+1}$ (no step of the arguments involves its harmonic part),
an arbitrary $\mathscr{C}'$, and arbitrary $(m,\mathscr{E})$ in the balls named.
The data enter the estimates through three channels:
through the inequalities of Section~0.3, which are imposed on the constants for every $e_k\le e_0$ and contain neither $k$, $\ell$, $\mu$ nor the system;
through the properties W of the system for $\mathrm{par}$, whose constants $(\alpha_0,\gamma_{\mathrm{w}},\kappa_{\mathrm{w}};C_w)$ are those of Section~I.8.2;
and through the results applied, which carry the same quantifiers.
A cutoff $N$, a charge $e$, a twist and a physical side of the torus do not occur in Sections~2--9: the data are $(k,\ell,\mu,e_k)$.
The level enters only through the distinction between $k=0$ and $k\ge1$;
at $k=0$ the weight $1/48$ of the unit points enters through $\mathsf{w}_W=\sqrt{48}\le7$ in $C_W$ and through $c_{\gamma}$.
\end{proof}

\subsection*{10. The case (o): the places at which the second slot is used}

In Sections~2--9 the second slot enters in five ways:
\begin{itemize}
\item[(i)] in the total connection $\mathscr{T}=A+\mathscr{C}$, at which the fermionic operators are taken;
\item[(ii)] as the second argument of $\mathscr{E}$ and of $\mathscr{E}^{r}$;
\item[(iii)] where two values of $\mathscr{C}$ are compared or $\mathscr{C}$ is transformed ((l2) of Section~0.4; Section~8; 9.2);
\item[(iv)] as a variable in assertions of continuity and analyticity;
\item[(v)] alone, as the connection in $E^{d}_k(X,\mathscr{C})$, inside
\[
{\delta E^{\mathrm{det}}}_k(X;\mathscr{A},\mathscr{C})=E^{d}_k(X,\mathscr{A}+\mathscr{C})-E^{d}_k(X,\mathscr{C})
\quad\text{and inside}\quad
E_k^{\mathrm{det}}(X,\mathscr{C})=E^{d}_k(X,\mathscr{C})-E^{d}_k(X,0):
\]
in 5.4 and in Section~6 in both cases, and in 9.6~(0), (5), (6), (8) and in 9.7 in the case (x);
and, in the case (x) only, as the argument of $\mathfrak{P}_k(X,\mathscr{C})$: in 5.4 and in 9.6~(0), (5), (6), (8).
\end{itemize}
What the occurrences (v) require is that $\mathscr{C}$, resp.\ $\mathscr{C}_L$, lie in $\tilde{G}_k$, which is (Sd2) together with $0\in\tilde{R}_k$,
and that $\mathscr{C}_L$ lie in $\mathrm{Sd}_k$, which is (Sd4).
In the case (o) the display is that of (D7) of Definition~24.D:
${\delta E^{\mathrm{det}}}_k$ is read $E_k^{\mathrm{det}}$, with $E_k^{\mathrm{det}}(X,\mathscr{A})=E^{d}_k(X,\mathscr{A})-E^{d}_k(X,0)$,
and the family $\mathfrak{P}_k$ and the entry $\mathfrak{P}_{k+1}$ of (8.6)$^{\mathrm{ch}}$ are void.
Accordingly, among the occurrences listed under (v), that of $E_k^{\mathrm{det}}(X,\mathscr{C}_L)$ inside $\mathfrak{P}_{k+1}(Y,\mathscr{C})$ in 5.4 belongs to the case (x) only.

The properties of $(\mathrm{Sd}_k,\mathrm{Sd}_{k+1})$ that enter the arguments, and the places where they enter, are the following.
\begin{description}
\item[{(Sd1)}] 3.2 and 5.2~(a): $\mathrm{al}_{0}$, $\mathrm{al}_{2}$ and the bound (X2) are formed at $(0,0)$;
in the case (x): 9.5 (third sentence), 9.6~(8), 9.7: the second slot is set equal to $0$ at the levels $k$ and $k+1$.
\item[{(Sd2)}] 2.2 (q1): $\mathscr{T}+t\cdot\mathrm{Zf}$ lies in $\tilde{G}_k$;
through (q1) in 3.1, 3.5, 3.7, 3.9, 5.1, 5.2, 5.4, 8.5--8.8;
5.2~(b): Proposition~23.1\,(g) at $B=\mathscr{T}$;
5.4: $\mathscr{A}_L+\mathscr{C}_L$ and $\mathscr{C}_L$ lie in $\tilde{G}_k$ ($0$ lies in $\tilde{R}_k$);
in the case (x): 9.6~(0): $\mathscr{C}$ lies in $\tilde{G}_k$.
\item[{(Sd3), level $k$}] 8.1 (th1), 8.5, 8.7, 8.8: $\tilde{g}\mathscr{C}$, $-\mathscr{C}$ and $-\overline{\mathscr{C}}$ lie in $\mathrm{Sd}_k$.
\item[{(Sd3), level $k+1$}] 9.2: (k3), (k5), (k7) are evaluated at $\tilde{g}'\mathscr{C}$, $-\mathscr{C}$ and $-\overline{\mathscr{C}}$.
\item[{(Sd4)}] 5.4, 9.2, 9.3, 9.4: $(\mathscr{A}_L,\mathscr{C}_L)$ lies in $\tilde{R}_k^{-}\times\mathrm{Sd}_k$.
\item[{(Sd5)}] (k1): 3.5 (c3), 3.7 (c7), 9.6~(1), (7);
(k2): 3.5 (c3), 3.7 (c7), 3.9, 3.10, 8.4, 9.3, 9.6~(5);
(k3): 8.5;
(k4): 8.6;
(k5): 8.8;
(k6): 5.2~(a);
(k7): 8.7;
(X1): 3.2, 3.5 (c4), 3.7 (c7), 5.2, 5.4, 8.7, 9.7;
(X2): 3.5, 5.2, 8.5--8.8, 9.6.
\item[{(Sd6)}] 3.1, 3.5--3.8, 9.1 (An), 9.3, 9.4: the meaning of the words ``analytic in $(A,\mathscr{C})$''.
\end{description}

All contours and all Cauchy formulas are in the variable $t$ ((N5), 5.1, 9.4, 9.8),
and all interpolations are in $s$ ((V), Lemmas~6.1 and 6.2 of Section~0.6, that is \cite[Lemmas~17.12, 17.13]{QEDI});
the identity theorem is applied in the variable $s'$ only (8.7 (h1)).

\emph{Statements applied that differ in the two cases.}
The facts (X1), (X2) are Proposition~16.3\,(a), (b), (d) in the case (x) and Proposition~48.2\,(a), (b), (d) in the case (o).
In 5.4, Proposition~18.5\,(a), resp.\ Proposition~48.5\,(a), is applied.
Proposition~13\,(d5) contains both $L\,{\delta E^{\mathrm{det}}}_k$ and $L\,E_k^{\mathrm{det}}$.
Proposition~13\,(d0) is a statement on all $B\in\tilde{G}_k$ and involves no slot;
in 5.4 it is applied at the connections of (i) and (v) and at $0$.
In the parts common to the two cases, that is, in Sections~2, 3, 5, 6, 8, 9.1--9.4 and 9.8, the other results applied are
Propositions~15, 16.1, 18.4, 23.1, 23.2,
the properties W of the system,
DF0--DF3 of Section~0.6 (\cite[Lemmas~17.3, 17.4(a), 17.5, 17.6]{QEDI}),
the facts (g0)--(g4) of Section~0.6 (\cite[Section~17.4.1]{QEDI}),
Lemmas~6.1 and 6.2 of Section~0.6,
the polymer facts of \cite[Theorem~13.2]{QEDI},
\cite[Propositions~15.21, 15.26]{QEDI},
and (I.Lc2)$^{\mathrm{ch}}$, a statement on all $B\in\tilde{G}_k$, which 5.1 (q3) applies at an arbitrary $B\in\tilde{G}_k$ and at $-B$;
none of them involves a slot.
The fact (P5) of Section~0.6, that is \cite[Proposition~15.30(ii)]{QEDI}, involves no slot either; it is applied only in 9.6~(7).
Where a connection containing $\mathscr{C}$ enters one of the results of the common parts, it is
$B=\mathscr{T}+t\cdot\mathrm{Zf}$ with $|t|\le\mathrm{rad}_1$ and $\mathrm{Zf}$ as in (q1) (in particular $\mathrm{Zf}=\mathscr{Z}(s')$; for $t=0$, $B=\mathscr{T}$),
or an image of such a $B$ under $\tilde{g}$, under $B\mapsto-B$, under $B\mapsto\overline{B}$ and under their compositions ((th1)),
or $B-d\Lambda$ (8.6);
in Proposition~23.2 the field is $\mathrm{Zf}=\mathscr{Z}(s')$, which does not contain $\mathscr{C}$.
In both cases $B$ and these images lie in $\tilde{R}_k+\mathrm{Sd}_k$, hence in $\tilde{R}_k+\mathfrak{X}_k$ and in $\tilde{G}_k$ ((q1), (Sd2), (Sd3), DF3),
and $B-d\Lambda$ lies in $\tilde{G}_k$ (DF3).

\emph{Parts that concern the case (x) only:} in 9.5, the proof of the third sentence of Proposition~24.1\,(b); 9.6; 9.7.
Besides results named above, these parts apply Proposition~48.1\,(a),
the identities $m(\mathscr{E})=m(\mathscr{E}|_0)$, $\varepsilon(\mathscr{E})=\varepsilon(\mathscr{E}|_0)$ and $(R\mathscr{E})|_0=R(\mathscr{E}|_0)$ of Proposition~16.3\,(a), (d),
and Propositions~17.2\,(a), 22.2\,(b), 22.3\,(a)--(c), 22.4\,(a), 18.5\,(b), (c) and 13\,(d6);
the statements of all these contain the second slot.

For the pair $(\{0\},\{0\})$ the properties (Sd1)--(Sd5) hold for the reasons given in Section~0.5, and (Sd6) is the reading (D7).
Hence, in the reading (o),
Sections~2, 3, 5, 6 prove Proposition~24.3,
Section~8 proves Proposition~24.2,
and Sections~5, 9.1--9.4, 9.8 prove the first two sentences of Proposition~24.1\,(b) and the second sentence of Proposition~24.1\,(e),
that is, the bound for $E^{*}_k$ on $\mathrm{Ball}_{k}(1)$ with the norms $\|\cdot\|_k$ and $\|\cdot\|_{k+1}$,
together with the parenthesis on $k=0$ of (e) as far as it concerns this bound.

\subsection*{11. Conditions on the constants}

The list below contains the conditions on the constants that are used in the arguments of this chapter
and in the results of \cite{QEDI} and of Section~I.5.3 that are applied directly
(DF0--DF3, (g0)--(g4), Lemmas~6.1 and 6.2, \cite[Theorem~13.2]{QEDI}, \cite[Propositions~15.21, 15.26, 15.30(ii)]{QEDI}, (I.Lc2)$^{\mathrm{ch}}$).
Every condition of the list is among the conditions of Sections~I.8.2 and I.8.4, that is, of Proposition~14.2.
For each condition the list gives the line of Section~I.8.2, or the label in Proposition~14.2, under which it stands, and the places where it is used.
\begin{description}
\item[{$L$ odd; $M$ a power of $L$}]
\emph{Stands in:} (C3) (L5); (C5) (M1).
\emph{Used:} (g0); \cite[Theorem~13.2]{QEDI}, \cite[Propositions~15.21, 15.26, 15.30(ii)]{QEDI}.
\item[{$L\ge4$ (from $L\ge4^{8}$)}]
\emph{Stands in:} (C3); Proposition~14.2\,(f), from (L1).
\emph{Used:} 9.1 (Sc3).
\item[{$L^{-\frac{3}{4}}\le\frac{1}{2}$, $L\ge5$}]
\emph{Stands in:} (C3); Proposition~14.2\,(f).
\emph{Used:} DF1, DF2; $L\ge5$ is also a hypothesis of \cite[Theorem~15.37]{QEDI}, the theorem in which \cite[Propositions~15.21, 15.26, 15.30(ii)]{QEDI} are collected.
\item[{$\kappa_0\ge c_X\ge1$, $K_0\ge1$}]
\emph{Stands in:} (C0), values; Proposition~14.2\,(f).
\emph{Used:} 5.3 Step~3; 5.2~(c) ($8\,c_X\cdot K_0+1\le9\,c_X\cdot K_0$); 5.3 (24.10) and 9.4 ($144\,c_X^3\cdot K_0^3$, $14112\,c_X^3\cdot K_0^4$, $504\,c_X^5\cdot K_0^5\le10^5\,c_X^5\cdot K_0^5$).
\item[{$\kappa\ge7\kappa_0+7$}]
\emph{Stands in:} (C2).
\emph{Used:} 0.3: the rates are positive, and $\kappa_0\le\mathrm{rt}_1$ (5.2~(c), 5.3); 9.1 (Sc4).
\item[{$L\left(\kappa-6\kappa_0-4\right)-1\ge\kappa$}]
\emph{Stands in:} (C3) (L3).
\emph{Used:} 9.1 (Sc4).
\item[{$\kappa_{\mathrm{w}}=\kappa+\kappa_0+3$}]
\emph{Stands in:} (C4).
\emph{Used:} 5.2~(b), $i=1$.
\item[{$M>4L$}]
\emph{Stands in:} (C5) (M2).
\emph{Used:} (g1), (g4).
\item[{$M^{\alpha_0}\ge1+\exp(\hat{\kappa})$}]
\emph{Stands in:} (C5) (M4).
\emph{Used:} Lemma~6.2; $M^{\alpha_0}>1$ in (D3), (V) and Lemma~6.1 (in 3.12 and 8.8).
\item[{$C_w\ge4\exp(2\gamma_{\mathrm{w}})\cdot M^{27\alpha_0}$}]
\emph{Stands in:} (C6), definitions of $C_{w}^{(1)}$ and $C_w$.
\emph{Used:} 2.2, 5.1 (q3), level $0$.
\item[{$C_w\ge1$, $c_{\gamma}\ge1$}]
\emph{Stands in:} (C6); Proposition~14.2\,(f).
\emph{Used:} 5.1 (q4).
\item[{$C_W=28\,C_w^2\cdot c_{\gamma}$; $\sqrt{48}\le7$}]
\emph{Stands in:} (C6); (I.L0.19).
\emph{Used:} 5.1 (q4).
\item[{$0<1-\alpha<\varepsilon\le\frac{1}{200}$}]
\emph{Stands in:} (C0), (C1), (C7) (S1).
\emph{Used:} comparisons of powers of $e_k$ (2.2, 5.2); 9.1 (Sc3);
9.4 ($p_{0,k}\ge1$, with the range $p_0>3r+2>0$ of the line (C7); Proposition~14.2\,(f));
DF1, DF2; (P5), that is \cite[Proposition~15.30(ii)]{QEDI}, in 9.6~(7).
\item[{(A1), first inequality, $p_{0,k}\le e_k^{-\frac{\varepsilon}{4}}$}]
\emph{Stands in:} (C8).
\emph{Used:} 2.2; 9.6~(7), through (P5).
\item[{(A1), second inequality, $12\exp\left(-\frac{p_{0,k}^2}{2}\right)\le e_k^{7}$}]
\emph{Stands in:} (C8).
\emph{Used:} 9.4.
\item[{(A1), third inequality, $e_k^{\frac{5\varepsilon}{2}}\le\frac{1}{2}$}]
\emph{Stands in:} (C8).
\emph{Used:} 2.2 ($\mathrm{rad}_2\le\mathrm{rad}_1$, $\mathrm{rad}_4\le\mathrm{rad}_1$).
\item[{(A1), fourth inequality, $c_{\mathrm{fl}}\cdot e_k^{\varepsilon}\le1$}]
\emph{Stands in:} (C8).
\emph{Used:} 5.2~(a).
\item[{(A1), fifth inequality, $3M\cdot e_k^{\varepsilon}\le1$}]
\emph{Stands in:} (C8).
\emph{Used:} (I.Lc2)$^{\mathrm{ch}}$ in 5.1 (q3).
\item[{(A2)}]
\emph{Stands in:} (C8).
\emph{Used:} 2.2; 9.6~(7), through (P5).
\item[{(A3), first part}]
\emph{Stands in:} (C8).
\emph{Used:} 5.2~(b), $i=1$, $k\ge1$.
\item[{(A3), second part}]
\emph{Stands in:} (C8).
\emph{Used:} through Proposition~23.2\,(e): 5.2~(b), $i=1$, $k=0$.
\item[{(A4)}]
\emph{Stands in:} (C8).
\emph{Used:} 5.2~(b), $i=2,3$.
\item[{(A5)}]
\emph{Stands in:} (C8).
\emph{Used:} 2.2 ($\mathrm{rad}_i\ge2$); 5.2~(b), (c).
\item[{(A6), with the number $112$}]
\emph{Stands in:} (C8).
\emph{Used:} 5.2~(c).
\item[{(A7), its three inequalities}]
\emph{Stands in:} (C8).
\emph{Used:} 5.3 (24.10); 9.4.
\item[{(A8), $e_k^{\varepsilon}\le\frac{1}{4}$ and $e_k^{\frac{1}{4}}+e_k^{\varepsilon}\le1$}]
\emph{Stands in:} (C8).
\emph{Used:} $e_k<1$ (throughout; (P5) in 9.6~(7)); DF1, DF2; 9.4.
\item[{(A10), third and fourth inequality}]
For $v\ge-\log e_0$:
\[
8\,C_B\cdot v^{p}\le\exp\left(\left(\frac{3}{4}-3\varepsilon\right)v\right),\qquad
8\exp\left(-\left(\frac{3}{4}-\frac{7\varepsilon}{2}\right)v\right)\le1.
\]
\emph{Stands in:} (C8).
\emph{Used:} (P5) in 9.6~(7).
\item[{$C_P\le C^{P}\le C_w$, $\gamma_{\mathrm{w}}\le\gamma^{P}\le\gamma_{P}$; $L^{\frac{3}{4}}\ge110\ge\frac{8}{3}$; $M\ge M^{P}$; $\alpha_0\le\alpha^{P}$}]
\emph{Stands in:} (C6), (C4), (C4r): Proposition~14.2\,(f); (C3): Proposition~14.2\,(f), from (L1); (C5) (M6); (C4), definition of $\alpha_0$.
\emph{Used:} (P5) in 9.6~(7): the first three through \cite[Proposition~15.30]{QEDI} and its proof; the last two as hypotheses of \cite[Theorem~15.37]{QEDI}.
\item[{$\ell$, $\ell/L$ in $\mathrm{Cl}$ (data)}]
\emph{Stands in:} (ST1)$^{\mathrm{ch}}$.
\emph{Used:} $M$ divides $\ell$ and $(2L+1)M<\ell$ ((g0)--(g4)); $\ell/(LM)$ is an integer ((t5), 9.1); $2M<\ell/L$ (pairs of level $k+1$); DF2, (P5).
\item[{$\ell\ge M\cdot r_k$, $\mu\le1/L$ (data)}]
\emph{Stands in:} (ST1)$^{\mathrm{ch}}$.
\emph{Used:} $\ell\ge M\cdot r_k$: only through Proposition~13\,(d5), (d6) (data (SF+)).
$\mu\le1/L$: directly in 9.6~(0) ($L\cdot\mu\le1$: the data $(k+1,\ell/L,L\cdot\mu,e_{k+1})$ are data (SF));
through the setting (S22) of Proposition~22.2\,(b) and Proposition~22.3 (in 9.6);
in the form $\mu\le1$ through the data of Proposition~23.1\,(g) (in 5.2~(b)) and the data (SF), (SF+) of Proposition~13\,(d0), (d5), (d6) (in 5.4, 9.6, 9.7).
\end{description}

\emph{The properties of the given system, and the places where they are used.}

\medskip
\noindent
\begin{tabular}{p{0.08\textwidth}p{0.34\textwidth}p{0.08\textwidth}p{0.34\textwidth}}
property & used in & property & used in\\
\hline
W1 & 9.6~(0), (7) & W10 & 3.3 (r1), 3.4\\
W2 & 3.1; 5.1 (q3); 8.7 (h1); 9.6~(7) & W11 & 3.3 (r2), 3.4\\
W3, W4 & 2.2 & W12 & 3.4; 3.5 (c1); 5.1 (q2); Section~6; 8.6 (h2)\\
W5 & 5.1 (q3) & W13 & 8.6\\
W6 & 5.1 (q2), (q4) & W14 & 8.5\\
W7 & only through Proposition~13, which requires of the family (p5) also W1, W2, W12--W14 & W15 & 8.8 (h2), Step~3~(ii)\\
W8 & 3.3 (r1)--(r3) & W16 & 8.7\\
 & & W17 & level $0$: 2.2, 3.1, 3.3, 3.4, 5.1, 8.4--8.8, 9.6
\end{tabular}

\section*{Chapter 26a. Projection of the localised operators to smaller tori: proofs of Propositions 26.1--26.4}

\subsection*{0. Conventions of this chapter}

\subsubsection*{0.1. Letters and numbers}

The notation is that of Part~I, in particular the bound $\mathrm{sd}$ of Section~I.1.1 and the index sets $\mathrm{Ind}$ of Section~I.1.9, together with that of Sections~G4.0.0 and G4.0.4.
As in Section~I.1.1 and in Proposition~27.1, a tilde marks objects of the covering lattice:
sites $\tilde{c}$, $\tilde{z}$, parameter cubes $\widetilde{\mathrm{pc}}$, $M$-cubes $\tilde{\square}$, and sets $\tilde{Z}$, $\tilde{X}$.
The letter $j$ stands for an integer, the letters $\mathrm{dk}$, $\mathrm{dk}'$ for vectors of $\mathbb{Z}^3$, and $\mathsf{Z}$, $\mathsf{Z}'$ for elements of $V_k$.
Further letters are introduced at the place where they are first needed:
$\mathrm{Fnd}_i$ and $\mathrm{Rw}$ in Section~0.4;
$\delta_{t',\mathrm{vl}}$, $\mathrm{vl}$ and $\mathrm{Kr}$ in Sections~0.5 and~3.2;
$\mathrm{dist}$, the sets $\tilde{S}$, $\tilde{S}'$, the half side $\mathrm{hs}$ of a cube and the relation $\mathrm{rel}$, which is one of $<$ and $\le$, in Section~0.6;
the constant $4\times4$ matrices $\mathrm{Mx}$ in Section~1;
$\mathrm{prc}_i$, $\mathrm{pjc}$, $\mathrm{Seg}$, $\mathrm{Sgs}$ and $\tilde{K}$ in Section~3.1;
$\mathrm{trm}_j$ in Section~3.4;
$\mathrm{mph}$, $\mathrm{Kn}$, $\mathrm{Yl}$, $\mathrm{Yl}_w$ and $\mathrm{pol}$ in Sections~3.7 and~4.
The map $\mathrm{mph}$ is the one that bears the same name in Section~G1.2, and $\mathrm{pol}$ is the map of (g2) in \cite[Section~17.4.1]{QEDI}.

In Sections~0, 1 and~3 the data are those of Proposition~26.1, and $i$ is $1$ or $2$.
The letter $a$ denotes a mesh in $\{\eta_k,1\}$ (in (1.6) and in Section~2, as in Section~I.1.5 and in (I.3.1)), except in a pair $(S,a)$ of W14, where it is the translation part (Section~I.4.2).
The letter $v$ denotes a translation vector in $t+v$, in $T_v$ and in the element $\tilde{g}_v$ of the groups of Section~I.4.2 (in Proposition~26.1, $v\in\ell_1\mathbb{Z}^3$);
where an operator or a pairing is applied to $v$, it denotes an element of $\mathfrak{S}(1,\ell_i)$ (Section~I.1.6).
The letters $t$, $t'$, $t''$ denote indices, $b$, $b'$ bonds (Section~G4.0.4), $w$, $w'$ vectors of $\mathbb{Z}^3$ (Proposition~26.1), and $Y$, $Z$, $Z'$ sets of photon sites (unions of cubes in Propositions~26.2 and~26.4);
in this chapter $t$ is never the contour variable, $b$ never the constant $b=1$, and $w$ never a flux.
For an index $t$, $\nu(t)$ is the weight of Section~I.3.6 (see Section~0.4); the letter $\nu$ does not denote a direction in this chapter.
The symbol $C_{0}^{1/2}(s')$ denotes the family $C_{k}^{1/2}(s')$ of (p2) in Section~G4.0.4, taken at $k=0$; it has nothing to do with the constant $C_{0}$ of the line (C6).

The displays (8.0) and (8.0$'$) are those of Proposition~26.3.
Proposition~26.3\,(a) and~(d) refer to the indices and the kernels of Section~G4.0.4; these indices belong to the index sets listed in Section~0.2, and these kernels are described in Section~0.5.
The labels (g0)--(g3) refer to the facts so labelled in \cite[Section~17.4.1]{QEDI}; they are not the conditions (g1)--(g9) of Section~I.8.4.

\subsubsection*{0.2. Index sets of the covering lattice}

The following eight sets are index sets of level $k$ in the sense of Section~I.1.9, and the indices of Section~G4.0.4 are among them:
each of them carries the translations $t\mapsto t+v$, $v\in LM\mathbb{Z}^3$, and a map $\operatorname{supp}$ into the finite non-empty subsets of $\mathbb{L}_k$ such that $\operatorname{supp}(t+v)=\operatorname{supp}t+v$.
The sets (x1)--(x7) are among the index sets listed in Section~I.1.9: the sets (x1)--(x4) consist of cells of $\mathbb{L}_k$ and of $\mathbb{Z}^3$ and belong to (I.i1), and (x5), (x6), (x7) are (I.i2), (I.i4), (I.i5), respectively.
The set (x8) is a further index set of level $k$ of the same sort.
\begin{itemize}
\item[(x1)] The sites $\tilde{c}$ of $\mathbb{L}_k$, with $\operatorname{supp}\tilde{c}=\{\tilde{c}\}$.
\item[(x2)] The fine bonds, that is, the positively oriented bonds $(\tilde{c},\tilde{c}+\eta_k\cdot e_{\mu})$ of $\mathbb{L}_k$; the support of such a bond is the set of its two end points.
\item[(x3)] The unit sites $\tilde{z}\in\mathbb{Z}^3$, with $\operatorname{supp}\tilde{z}=\{\tilde{z}\}$.
\item[(x4)] The unit bonds, that is, the positively oriented bonds $(\tilde{z},\tilde{z}+e_{\mu})$ of $\mathbb{Z}^3$, with the support $\{\tilde{z},\tilde{z}+e_{\mu}\}$, the set of the two end points.
\item[(x5)] The basis elements $\Upsilon$: the bond functions $e^{y}_{i}$, $y\in L\mathbb{Z}^3$, and $\delta_b$, $b$ a positively oriented face bond of $\mathbb{Z}^3$ that is not central \cite[Section~5.2.5]{QEDI};
here $\operatorname{supp}\Upsilon$ is the set of the end points of the bonds at which $\Upsilon$ does not vanish, and $(\Upsilon+v)(b):=\Upsilon(b-v)$.
\item[(x6)] The fermion sites $x=(\tilde{c},g)\in\mathbb{L}^{\mathrm{f}}_k$, with $\operatorname{supp}x=\{\tilde{c}\}$ and $x+v=(\tilde{c}+v,g)$ (Section~I.1.3, (I.POS)).
\item[(x7)] The pairs $(x,x')\in\mathrm{Pzt}_k$, $x=(\tilde{c},g)$, $x'=(\tilde{c}',g')$, with $\operatorname{supp}(x,x')=\{\tilde{c},\tilde{c}'\}$ and $(x,x')+v=(x+v,x'+v)$ (Section~I.1.7).
\item[(x8)] The parameter cubes $\widetilde{\mathrm{pc}}$ of $\mathbb{L}_k$ of one kind, named by their centres $\tilde{z}$, where $\tilde{z}\in M\mathbb{Z}^3$ for $\pi_k$ and $\pi'_k$ and $\tilde{z}\in LM\mathbb{Z}^3$ for $\pi''_k$;
$\operatorname{supp}\widetilde{\mathrm{pc}}$ is the set of the sites of $\widetilde{\mathrm{pc}}$.
\end{itemize}
Concerning (x5): the vectors $e^{0}_{i}$ form the orthonormal basis, fixed in \cite[Section~5.2.5]{QEDI}, of the kernel of $\tau$ in the space of the real functions on the interior bonds of the $L$-block $B(0)$ with the centre $0$.
In the symbols $e^{0}_{i}$ and $e^{y}_{i}$ the upper index is a point of $L\mathbb{Z}^3$, not an exponent, and the lower index numbers the vectors of this basis;
this lower index is not the index $i\in\{1,2\}$ of the torus $T_i$.
The function $e^{y}_{i}$ is the translate of $e^{0}_{i}$ by $y$, and the $L$-blocks, their faces and the central bonds are invariant under $L\mathbb{Z}^3$ (\cite[Proposition~5.35]{QEDI}, \cite[Proposition~15.21\,(i), (ii)]{QEDI});
hence $\Upsilon+v$ is a basis element for every $v\in L\mathbb{Z}^3$.
By \cite[Proposition~5.35]{QEDI}, $\operatorname{supp}\Upsilon$ is contained in one $L$-block or consists of the two end points of one bond.
The unit lattice, the $L$-blocks, the $M$-cubes, the $LM$-cubes and the parameter cubes are invariant under $LM\mathbb{Z}^3$.
Since $\ell_1$ and $\ell_2$ are multiples of $L\cdot M$, the groups $\ell_2\mathbb{Z}^3$ and $\ell_1\mathbb{Z}^3$, the first contained in the second, act on all these sets.
The action is free: for $v\ne0$ the finite set $\operatorname{supp}t+v$ is different from $\operatorname{supp}t$.

The supports of the indices (x1)--(x6), which include the indices of Section~G4.0.4, have sup-diameter at most $L$ (Section~G4.0.4).
More precisely, the sup-diameter is $\eta_k$ or $1$ for a bond, at most $L-1$ for a basis element $e^{y}_{i}$ (an $L$-block has the sup-diameter $L-1$), $1$ for a basis element $\delta_b$, and $0$ for a site.

For an index $t$ of the covering lattice, $\mathrm{pr}_i t$ denotes the index of $T_i$ that lies under $t$; thus $\operatorname{supp}(\mathrm{pr}_i t)=\mathrm{pr}_i(\operatorname{supp}t)$.
Let $Z$ be a set of sites of $T_i$.
The index $t$ is said to be contained in $Z$ (on $T_i$) if $\mathrm{pr}_i(\operatorname{supp}t)\subset Z$; for a unit site or a fermion site this means that $t$ lies in (over) $Z$.
If $t$ is one lift of an index of $T_i$, the lifts of this index are the indices $t+\ell_i\cdot\mathrm{dk}$, $\mathrm{dk}\in\mathbb{Z}^3$.

\subsubsection*{0.3. Objects on $T_i$}

In both cases of Proposition~26.1 the map $\boldsymbol{\chi}\colon\mathbb{Z}^3\to\{1,-1\}$ is a homomorphism.
An object on $T_i$ with the index set $\mathrm{Ind}$ is a function $f$ on $\mathrm{Ind}$ with values in $\mathbb{C}$ (bond functions, scalars, coordinates of elements of $V_k$, parameter families) or in $\mathbb{C}^4$ (spinors, functions of pairs) such that
\begin{equation}
f(t+\ell_i\cdot\mathrm{dk})=\boldsymbol{\chi}(\mathrm{dk})\cdot f(t)\qquad\left(t\in\mathrm{Ind},\ \mathrm{dk}\in\mathbb{Z}^3\right)
\tag{0.a}
\end{equation}
(Sections~I.1.5, I.1.7, I.1.9 and~I.6.3).
A bond function is extended to the negatively oriented bonds by $A(c',c)=-A(c,c')$.
An element $\mathsf{Z}$ of $V_k(T_i)$, or of its complexification when $s'$ is complex, is at the same time two things:
the $\ell_i\mathbb{Z}^3$-periodic bond function on $\mathbb{Z}^3$ with $\tau\mathsf{Z}=0$ in every $L$-block and $\mathsf{Z}=0$ on the central bonds,
and the $\ell_i\mathbb{Z}^3$-periodic function
\[
\Upsilon\mapsto\mathsf{Z}(\Upsilon)=\sum_b\mathsf{Z}(b)\cdot\Upsilon(b)
\]
on (x5), where $b$ ranges over the positively oriented bonds at which $\Upsilon$ does not vanish.
Because $L+1<\ell_i$, the map $\mathrm{pr}_i$ is injective on these bonds, so that $\mathsf{Z}(\Upsilon)$ is the coordinate $\langle\mathsf{Z},\mathrm{pr}_i\Upsilon\rangle$ of \cite[Section~5.2.5]{QEDI} on $T_i$.
As $\{\mathrm{pr}_i\Upsilon\}$ is a basis \cite[Proposition~5.35]{QEDI}, every $\ell_i\mathbb{Z}^3$-periodic complex function on (x5) is the family of the coordinates of exactly one element of the complexification of $V_k(T_i)$.
The spaces that occur, and their index sets, are:
bond functions on $T_i$: (x2);
bond functions on $T^0_i$: (x4);
$V_k(T_i)$: (x5);
scalars: (x1);
$\mathfrak{S}(1,\ell_i)$: (x3);
$\mathfrak{S}^{\mathrm{ch}}(\eta_k,\ell_i)$: (x6);
functions of pairs: (x7);
parameter families: (x8).

\emph{Level~$0$.}
Here $\eta_0=1$ and $\mathbb{L}_0=\mathbb{Z}^3$, the sets (x2) and (x4) coincide, and the index set of $\mathfrak{S}^{\mathrm{ch}}(1,\ell_i)$ is (x6) $=\mathbb{Z}^3\times O_{\mathrm{cub}}$.

\subsubsection*{0.4. Weights, fundamental sets, pairings}

Put
\[
\nu(t):=1\ \text{on (x3), (x4), (x5)},\qquad\nu(t):=\eta_k^3\ \text{on (x1), (x2)},\qquad\nu(t):=\frac{\eta_k^3}{48}\ \text{on (x6)}.
\]
The weight $\nu$ is invariant under translations and does not depend on the side.
A fundamental set $\mathrm{Fnd}_i$ of an index set $\mathrm{Ind}$ is a subset that contains exactly one point of every orbit of $\ell_i\mathbb{Z}^3$;
for (x6) it is chosen of the form $\mathrm{Fnd}^0_i\times O_{\mathrm{cub}}$ with a fundamental set $\mathrm{Fnd}^0_i$ of (x1).
Let $\mathrm{Rw}$ be a set of representatives of $\mathbb{Z}^3/\mathrm{Jr}\,\mathbb{Z}^3$; then $\#|\mathrm{Rw}|=\mathrm{Jr}^3$.
\begin{itemize}
\item[(fs)] If $\mathrm{Fnd}_1$ is a fundamental set for $\ell_1\mathbb{Z}^3$, then the sets $\mathrm{Fnd}_1+\ell_1\cdot w$, $w\in\mathrm{Rw}$, are pairwise disjoint, and their union is a fundamental set for $\ell_2\mathbb{Z}^3$.
Indeed, the action is free.
Hence every $t\in\mathrm{Ind}$ can be written as $t=t_1+\ell_1\cdot\mathrm{dk}$ with exactly one pair $(t_1,\mathrm{dk})\in\mathrm{Fnd}_1\times\mathbb{Z}^3$,
and $\mathrm{dk}=w+\mathrm{Jr}\cdot\mathrm{dk}'$ with exactly one pair $(w,\mathrm{dk}')\in\mathrm{Rw}\times\mathbb{Z}^3$;
therefore $t=(t_1+\ell_1\cdot w)+\ell_2\cdot\mathrm{dk}'$ with exactly one triple $(t_1,w,\mathrm{dk}')$.
\end{itemize}
Let $f'$ and $f$ be objects on $T_i$ with the same index set, one of (x1)--(x6) (the sets on which $\nu$ is defined), and with the same character.
Put
\begin{equation}
\langle f',f\rangle_{T_i}:=\sum_{t\in\mathrm{Fnd}_i}\nu(t)\cdot f'(t)\cdot f(t),
\tag{0.b}
\end{equation}
where $f'(t)\cdot f(t)$ is the product of two numbers or, for values in $\mathbb{C}^4$, the sum $\sum_{\beta}f'_{\beta}(t)\cdot f_{\beta}(t)$.
By (0.a) and $\boldsymbol{\chi}(\mathrm{dk})^2=1$ the summand is an $\ell_i\mathbb{Z}^3$-periodic function of $t$; therefore (0.b) does not depend on the choice of $\mathrm{Fnd}_i$.
The pairings named in Proposition~26.1 are of this form:
\begin{align*}
\langle f',f\rangle^{\mathrm{ch}}&=\eta_k^3\cdot\frac{1}{48}\sum_{\tilde{c}}\sum_{g\in O_{\mathrm{cub}}}\sum_{\beta}f'_{\beta}(\tilde{c},g)\cdot f_{\beta}(\tilde{c},g)&&\text{on }\mathfrak{S}^{\mathrm{ch}}(\eta_k,\ell_i),\\
\langle v',v\rangle_1&=\sum_{\tilde{z}}\sum_{\beta}v'_{\beta}(\tilde{z})\cdot v_{\beta}(\tilde{z})&&\text{on }\mathfrak{S}(1,\ell_i),\\
\langle A',A\rangle&=\sum_{\tilde{c}}\eta_k^3\sum_{\mu}A'_{\mu}(\tilde{c})\cdot A_{\mu}(\tilde{c})&&\text{for bond functions on }T_i,\\
\langle A',A\rangle&=\sum_{\tilde{z}}\sum_{\mu}A'_{\mu}(\tilde{z})\cdot A_{\mu}(\tilde{z})&&\text{for bond functions on }T^0_i,\\
\langle\boldsymbol{\lambda}',\boldsymbol{\lambda}\rangle&=\sum_{\tilde{c}}\eta_k^3\cdot\boldsymbol{\lambda}'(\tilde{c})\cdot\boldsymbol{\lambda}(\tilde{c})&&\text{for scalars},\\
\langle\mathsf{Z}',\mathsf{Z}\rangle&=\sum_b\mathsf{Z}'(b)\cdot\mathsf{Z}(b)=\sum_{\Upsilon}\mathsf{Z}'(\Upsilon)\cdot\mathsf{Z}(\Upsilon)&&\text{on }V_k(T_i),
\end{align*}
where $\tilde{c}$, $\tilde{z}$, $b$ and $\Upsilon$ range over fundamental sets.
The first pairing is formed with the measure (I.1.3), as in Section~I.1.6, and the second is the pairing $\langle\,,\rangle_1$ of Section~I.1.6;
the three pairings of bond functions and scalars are those of \cite[Section~5.2.5]{QEDI};
the last one is the counting scalar product of \cite[Section~5.2.5]{QEDI}, for which $\{\mathrm{pr}_i\Upsilon\}$ is an orthonormal basis \cite[Proposition~5.35]{QEDI}.
Each of these pairings is non-degenerate, being a sum of products of coordinates with positive weights.

\subsubsection*{0.5. Kernels}

The kernels of Section~G4.0.4 and of (I.1.4) are described as follows.
Let $O$ be a linear operator from the objects on $T_i$ with the index set $\mathrm{Ind}'$ (the columns), one of (x1)--(x6), to the objects on $T_i$ with the index set $\mathrm{Ind}$ (the rows), both with the character $\boldsymbol{\chi}$.
For $t'\in\mathrm{Ind}'$ and an element $\mathrm{vl}$ of the space of values ($\mathbb{C}$ or $\mathbb{C}^4$) let $\delta_{t',\mathrm{vl}}$ be the object on $T_i$ that vanishes off the orbit of $t'$ and satisfies
\[
\delta_{t',\mathrm{vl}}(t'+\ell_i\cdot\mathrm{dk}):=\boldsymbol{\chi}(\mathrm{dk})\cdot\nu(t')^{-1}\cdot\mathrm{vl}\qquad\left(\mathrm{dk}\in\mathbb{Z}^3\right).
\]
The kernel of $O$ is the function $O(t,t')$ on $\mathrm{Ind}\times\mathrm{Ind}'$, a number or a $4\times4$ matrix, given by $O(t,t')\cdot\mathrm{vl}:=\left(O\,\delta_{t',\mathrm{vl}}\right)(t)$.
It is the only function on $\mathrm{Ind}\times\mathrm{Ind}'$ with the following two properties:
\begin{itemize}
\item[(kn1)] $(Of)(t)=\sum_{t'\in\mathrm{Fnd}_i}\nu(t')\cdot O(t,t')\cdot f(t')$ for all $f$ and for one fundamental set $\mathrm{Fnd}_i$ (and then for every fundamental set);
\item[(kn2)] $O(t,t'+\ell_i\cdot\mathrm{dk})=\boldsymbol{\chi}(\mathrm{dk})\cdot O(t,t')$.
\end{itemize}
Moreover it has the property
\begin{itemize}
\item[(kn3)] $O(t+\ell_i\cdot\mathrm{dk},t')=\boldsymbol{\chi}(\mathrm{dk})\cdot O(t,t')$.
\end{itemize}
Indeed, (kn1) follows from $f=\sum_{t'\in\mathrm{Fnd}_i}\nu(t')\cdot\delta_{t',f(t')}$,
(kn2) follows from $\delta_{t'+\ell_i\cdot\mathrm{dk},\mathrm{vl}}=\boldsymbol{\chi}(\mathrm{dk})\cdot\delta_{t',\mathrm{vl}}$,
and (kn3) holds because $O\,\delta_{t',\mathrm{vl}}$ is an object on $T_i$.
For the uniqueness, let $\mathrm{Kr}$ be a function with (kn1) and (kn2), and write $t'=t'_0+\ell_i\cdot\mathrm{dk}$ with $t'_0\in\mathrm{Fnd}_i$.
Then $\delta_{t',\mathrm{vl}}(t'_0)=\boldsymbol{\chi}(\mathrm{dk})\cdot\nu(t')^{-1}\cdot\mathrm{vl}$, and (kn1) for $f=\delta_{t',\mathrm{vl}}$ reads, since $\nu(t'_0)=\nu(t')$,
\[
\left(O\,\delta_{t',\mathrm{vl}}\right)(t)=\boldsymbol{\chi}(\mathrm{dk})\cdot\mathrm{Kr}(t,t'_0)\cdot\mathrm{vl}=\mathrm{Kr}(t,t')\cdot\mathrm{vl}.
\]
With $\nu=1$ on unit bonds, basis elements and unit sites, and $\nu=\frac{\eta_k^3}{48}$ on fermion sites, the properties (kn1)--(kn3) are the description of the kernels given in Section~G4.0.4.
For a source $\mathfrak{S}^{\mathrm{ch}}(\eta_k,\ell_i)$, (kn1) is (I.1.4):
\[
(Of)(t)=\eta_k^3\cdot\frac{1}{48}\sum_{\tilde{c}'}\sum_{g'\in O_{\mathrm{cub}}}O\left(t,(\tilde{c}',g')\right)\cdot f(\tilde{c}',g').
\]
By (kn2) and (kn3), the modulus of $O(t,t')$, and whether $O(t,t')$ vanishes, depend on $\mathrm{pr}_i t$ and $\mathrm{pr}_i t'$ only.

\subsubsection*{0.6. Distances, cubes}

Let $d_i$ be the sup distance of the torus $T_i$ (Section~I.1.1, Proposition~26.1); for sites $\tilde{c}$, $\tilde{c}'$ of $\mathbb{L}_k$ it is given by
\begin{equation}
d_i(\mathrm{pr}_i\tilde{c},\mathrm{pr}_i\tilde{c}')=\min_{\mathrm{dk}\in\mathbb{Z}^3}\left|\tilde{c}-\tilde{c}'-\ell_i\cdot\mathrm{dk}\right|.
\tag{0.c}
\end{equation}
\begin{itemize}
\item[(ds1)] If $|\tilde{c}-\tilde{c}'|<\frac{\ell_i}{2}$, then $d_i(\mathrm{pr}_i\tilde{c},\mathrm{pr}_i\tilde{c}')=|\tilde{c}-\tilde{c}'|$, and $\tilde{c}-\tilde{c}'$ is the only element of its class modulo $\ell_i\mathbb{Z}^3$ with sup norm smaller than $\frac{\ell_i}{2}$.
Indeed, for $\mathrm{dk}\ne0$ one has $\left|\tilde{c}-\tilde{c}'-\ell_i\cdot\mathrm{dk}\right|\ge\ell_i-|\tilde{c}-\tilde{c}'|>\frac{\ell_i}{2}$.
\item[(ds2)] $d_1(\mathrm{pr}_1\tilde{c},\mathrm{pr}_1\tilde{c}')\le d_2(\mathrm{pr}_2\tilde{c},\mathrm{pr}_2\tilde{c}')$; that is, $d_1(\mathrm{pj}\,c,\mathrm{pj}\,c')\le d_2(c,c')$ for sites $c$, $c'$ of $T_2$.
This holds because $\ell_2\mathbb{Z}^3\subset\ell_1\mathbb{Z}^3$.
\item[(ds3)] For finite sets $\tilde{S}$, $\tilde{S}'$ of sites of $\mathbb{L}_k$ put
\[
\mathrm{dist}(\tilde{S},\tilde{S}'):=\min\left\{|\tilde{c}-\tilde{c}'|:\ \tilde{c}\in\tilde{S},\ \tilde{c}'\in\tilde{S}'\right\}.
\]
For indices $t$, $t'$ of the covering lattice, $d_i(t,t')$ is the sup distance on $T_i$ of the sets $\mathrm{pr}_i\operatorname{supp}t$ and $\mathrm{pr}_i\operatorname{supp}t'$ (Section~G4.0.4); by (0.c),
\[
d_i(t,t')=\min_{\mathrm{dk}\in\mathbb{Z}^3}\mathrm{dist}\left(\operatorname{supp}t,\operatorname{supp}t'+\ell_i\cdot\mathrm{dk}\right).
\]
\item[(ds4)] (Cubes, unit blocks.)
Let $\tilde{z}\in M\mathbb{Z}^3$ and $(\mathrm{hs},\mathrm{rel})=\left(\frac{M}{2},<\right)$ or $(M,\le)$;
or let $\tilde{z}\in LM\mathbb{Z}^3$ and $(\mathrm{hs},\mathrm{rel})=\left(\frac{L\cdot M}{2},<\right)$ or $(L\cdot M,\le)$;
or let $\tilde{z}\in\mathbb{Z}^3$ and $(\mathrm{hs},\mathrm{rel})=\left(\frac{1}{2},<\right)$.
Then, by (0.c),
\[
\mathrm{pr}_i\left\{\tilde{c}:\ |\tilde{c}-\tilde{z}|\mathrel{\mathrm{rel}}\mathrm{hs}\right\}=\left\{c\in T_i:\ d_i(c,\mathrm{pr}_i\tilde{z})\mathrel{\mathrm{rel}}\mathrm{hs}\right\}.
\]
Indeed, if $|\tilde{c}-\tilde{z}|\mathrel{\mathrm{rel}}\mathrm{hs}$, then $d_i(\mathrm{pr}_i\tilde{c},\mathrm{pr}_i\tilde{z})\le|\tilde{c}-\tilde{z}|$, and hence $d_i(\mathrm{pr}_i\tilde{c},\mathrm{pr}_i\tilde{z})\mathrel{\mathrm{rel}}\mathrm{hs}$.
Conversely, if $d_i(c,\mathrm{pr}_i\tilde{z})\mathrel{\mathrm{rel}}\mathrm{hs}$, then by (0.c) the site $c$ has a lift $\tilde{c}$ with $|\tilde{c}-\tilde{z}|=d_i(c,\mathrm{pr}_i\tilde{z})$.
Consequently $\mathrm{pr}_i$ maps the $M$-cubes, the cubes of $\pi'_k$, the $LM$-cubes, the cubes of $\pi''_k$ and the unit blocks $\Delta_{\tilde{z}}=\left\{\tilde{c}:|\tilde{c}-\tilde{z}|<\frac{1}{2}\right\}$ (Section~I.1.1) of $\mathbb{L}_k$ onto the sets of the same kind of $T_i$, centre to centre;
in particular $\mathrm{pr}_i\Delta_{\tilde{z}}=\Delta_{\mathrm{pr}_i\tilde{z}}$.
Since $\mathrm{pj}\circ\mathrm{pr}_2=\mathrm{pr}_1$, the map $\mathrm{pj}$ carries the cube of $T_2$ with the centre $z'$ onto the cube of the same kind of $T_1$ with the centre $\mathrm{pj}\,z'$.
Cubes are named by their centres.
\end{itemize}

\subsection*{1. Proposition 26.1}

\begin{proof}[Verification of the facts read off in Proposition~26.1, and of the identity stated with the definition of $\bar{H}^{\mathrm{pj}}_k$]
\textbf{(1.0) The data are (SF).}
The integer $k\ge0$, the side $\ell_i\in\mathrm{Cl}$, the mass $\mu\in(0,1]$ and the number $e_k\in(0,e_0]$ are exactly what the definition of (SF) in Section~I.1.8 requires, for each of the two quadruples $(k,\ell_i,\mu,e_k)$.

\textbf{(1.1) Every object on $T_1$ is an object on $T_2$; the image of $\mathrm{pj}^{*}$ is $\mathrm{Per}$.}
Let $f$ be an object on $T_1$, so that $f(t+\ell_1\cdot\mathrm{dk})=\boldsymbol{\chi}(\mathrm{dk})\cdot f(t)$ for all $\mathrm{dk}\in\mathbb{Z}^3$.
Then
\[
f(t+\ell_2\cdot\mathrm{dk})=f(t+\ell_1\cdot\mathrm{Jr}\cdot\mathrm{dk})=\boldsymbol{\chi}(\mathrm{Jr}\cdot\mathrm{dk})\cdot f(t)=\boldsymbol{\chi}(\mathrm{dk})^{\mathrm{Jr}}\cdot f(t)=\boldsymbol{\chi}(\mathrm{dk})\cdot f(t),
\]
because $\boldsymbol{\chi}$ is a homomorphism with values $\pm1$ and $\mathrm{Jr}$ is odd.
This is (0.a) for $i=2$.
The remaining conditions in the definitions of the spaces (a bond function is odd under the reversal of a bond; an element $\mathsf{Z}$ of $V_k$ satisfies $\tau\mathsf{Z}=0$ in every $L$-block and vanishes on the central bonds) are conditions on the function on the covering index set, and no side enters them.
Hence $f$ is an object on $T_2$ of the same type; this object is $\mathrm{pj}^{*}f$ (Section~I.6.3), and it satisfies the condition that defines $\mathrm{Per}$.
Conversely, let $f$ be an object on $T_2$ that lies in $\mathrm{Per}$.
The condition defining $\mathrm{Per}$ is (0.a) for $i=1$, and the remaining conditions are those on $T_1$; hence $f$ is an object on $T_1$, and $f=\mathrm{pj}^{*}f$.
On fermion sites one has $(\mathrm{pj}^{*}f)(\tilde{c},g)=f(\tilde{c},g)$ for each of the $48$ labels $g$ (Section~I.1.3),
and on pairs $(\mathrm{pj}^{*}F)(x,x')=F(x,x')$ for $(x,x')\in\mathrm{Pzt}_k$, the pairs inside a cluster included.
The set $\mathrm{Pzt}_k$ is defined on $\mathbb{L}^{\mathrm{f}}_k$ by the condition $d_X\le M$ (Section~I.1.7), in which no side occurs.

\textbf{(1.2) $\mathrm{pj}^{*}$ is linear and injective, and $|\mathrm{pj}^{*}f|_{\infty}=|f|_{\infty}$.}
The functions $\mathrm{pj}^{*}f$ and $f$ coincide on the covering index set.
The norm $|f|_{\infty}$ is the supremum over this index set of the modulus of $f(t)$ (for values in $\mathbb{C}^4$, of $|f(t)|_c$, Section~I.1.5), and no side enters it.

\textbf{(1.3) The pairings.}
Let $f'$ and $f$ be objects on $T_1$ with the same index set $\mathrm{Ind}$, one of (x1)--(x6), and with the same character, and let $\mathrm{Fnd}_1$ be a fundamental set for $\ell_1\mathbb{Z}^3$.
By (fs), the union of the sets $\mathrm{Fnd}_1+\ell_1\cdot w$, $w\in\mathrm{Rw}$, may be taken as $\mathrm{Fnd}_2$ in (0.b).
Since $\nu(t+\ell_1\cdot w)=\nu(t)$ and $f'(t+\ell_1\cdot w)\cdot f(t+\ell_1\cdot w)=\boldsymbol{\chi}(w)^2\cdot f'(t)\cdot f(t)=f'(t)\cdot f(t)$, one obtains
\begin{align*}
\langle\mathrm{pj}^{*}f',\mathrm{pj}^{*}f\rangle_{T_2}&=\sum_{w\in\mathrm{Rw}}\sum_{t\in\mathrm{Fnd}_1}\nu(t+\ell_1\cdot w)\cdot f'(t+\ell_1\cdot w)\cdot f(t+\ell_1\cdot w)\\
&=\sum_{w\in\mathrm{Rw}}\sum_{t\in\mathrm{Fnd}_1}\nu(t)\cdot f'(t)\cdot f(t)=\mathrm{Jr}^3\cdot\langle f',f\rangle_{T_1}.
\end{align*}
For $\langle\,,\rangle^{\mathrm{ch}}$ on $\mathfrak{S}^{\mathrm{ch}}(\eta_k,\cdot)$ this display reads, with $\mathrm{Fnd}_1=\mathrm{Fnd}^0_1\times O_{\mathrm{cub}}$ for a fundamental set $\mathrm{Fnd}^0_1$ of the sites of $\mathbb{L}_k$,
\begin{multline*}
\eta_k^3\cdot\frac{1}{48}\sum_{w\in\mathrm{Rw}}\sum_{\tilde{c}\in\mathrm{Fnd}^0_1}\sum_{g\in O_{\mathrm{cub}}}\sum_{\beta}f'_{\beta}(\tilde{c}+\ell_1\cdot w,g)\cdot f_{\beta}(\tilde{c}+\ell_1\cdot w,g)\\
=\mathrm{Jr}^3\cdot\eta_k^3\cdot\frac{1}{48}\sum_{\tilde{c}\in\mathrm{Fnd}^0_1}\sum_{g\in O_{\mathrm{cub}}}\sum_{\beta}f'_{\beta}(\tilde{c},g)\cdot f_{\beta}(\tilde{c},g):
\end{multline*}
the translation moves only $\tilde{c}$, the two signs $(-1)^{w_1+w_2+w_3}$ multiply to $1$, and the weight $\frac{\eta_k^3}{48}$ is the same on both tori.
For $\langle\,,\rangle_1$ on $\mathfrak{S}(1,\cdot)$ one has $\mathrm{Ind}=$ (x3) and $\nu=1$.
For bond functions on $T_i$, bond functions on $T^0_i$ and scalars one has $\mathrm{Ind}=$ (x2), (x4), (x1), respectively $\nu=\eta_k^3$, $1$, $\eta_k^3$, and $\boldsymbol{\chi}=1$.
For $V_k$ one has $\mathrm{Ind}=$ (x5), $\nu=1$ and $\boldsymbol{\chi}=1$.

\emph{Level~$0$.}
Here $\eta_0=1$, and the display for $\langle\,,\rangle^{\mathrm{ch}}$, with the weight $\frac{1}{48}$, is also the pairing of the unit lattice $T^{\mathrm{f}}_0$ (Section~I.1.6, (I.L0.4)).

\textbf{(1.4) Parameter families.}
By (ds4), the parameter cubes of $T_i$ are the images $\mathrm{pr}_i(\widetilde{\mathrm{pc}})$, and $\mathrm{pj}(\mathrm{pr}_2\widetilde{\mathrm{pc}})=\mathrm{pr}_1\widetilde{\mathrm{pc}}$.
A parameter family on $T_i$ is an $\ell_i\mathbb{Z}^3$-periodic function $\widetilde{\mathrm{pc}}\mapsto s_{\widetilde{\mathrm{pc}}}$ on (x8), and $s_{\mathrm{pr}_i\widetilde{\mathrm{pc}}}:=s_{\widetilde{\mathrm{pc}}}$.
For a family $s$ on $T_1$ and $\mathrm{pc}=\mathrm{pr}_2(\widetilde{\mathrm{pc}})$ this gives
\[
(\mathrm{pj}^{*}s)_{\mathrm{pc}}=s_{\widetilde{\mathrm{pc}}}=s_{\mathrm{pr}_1\widetilde{\mathrm{pc}}}=s_{\mathrm{pj}\,\mathrm{pc}}.
\]
Hence $\sup_{\mathrm{pc}}\left|(\mathrm{pj}^{*}s)_{\mathrm{pc}}\right|=\sup_{\mathrm{pc}}|s_{\mathrm{pc}}|$.
The polydisk $\mathrm{Pd}_k(T_i)$ is the product of the closed disks $|s_{\mathrm{pc}}|\le M^{\alpha_{0}}$ \cite[Definition~5.8]{QEDI}, and its interior is the product of the open disks.
Therefore $\mathrm{pj}^{*}$ maps $\mathrm{Pd}_k(T_1)$ into $\mathrm{Pd}_k(T_2)$, and the interior into the interior.

\textbf{(1.5) $\tilde{G}_k$.}
The lemma \cite[Lemma~17.7]{QEDI}, which reappears as \cite[Lemma~5.20]{QEDI}, holds for every $k\ge0$, every $e_k>0$ and all $\ell_1,\ell_2\in\mathrm{Cl}$ with $\ell_2=\mathrm{Jr}\cdot\ell_1$, $\mathrm{Jr}$ odd;
it asserts that $\mathrm{pj}^{*}$ maps $\tilde{G}_k$ on $T_1$ into $\tilde{G}_k$ on $T_2$.
The data of Proposition~26.1 are of this kind, and on bond functions $\mathrm{pj}^{*}$ is the identification used in that lemma.

\textbf{(1.6) $J$ and $J^{*}$.}
Let $a\in\{\eta_k,1\}$, $\xi\in\mathfrak{S}(a,\ell_1)$ and $f\in\mathfrak{S}^{\mathrm{ch}}(a,\ell_1)$.
By (I.3.1),
\begin{gather*}
(J\,\mathrm{pj}^{*}\xi)(\tilde{c},g)=\xi(\tilde{c})=(\mathrm{pj}^{*}J\xi)(\tilde{c},g)\quad\text{for all }48\text{ labels }g,\\
(J^{*}\,\mathrm{pj}^{*}f)(\tilde{c})=\frac{1}{48}\sum_{g\in O_{\mathrm{cub}}}f(\tilde{c},g)=(\mathrm{pj}^{*}J^{*}f)(\tilde{c}).
\end{gather*}
No side occurs in the two displays of (I.3.1); this is the commutation with the periodic extension recorded in (I.J3).

\textbf{(1.7) $A\mapsto-A$, complex conjugation, constant spin matrices.}
The values $(-A)(b)$, $\overline{f}(t)$ and $\mathrm{Mx}\cdot f(t)$, with $\mathrm{Mx}$ a constant $4\times4$ matrix, are formed from the value at the same index.
These operations preserve (0.a) for $i=1$ and for $i=2$.
Hence
\[
\mathrm{pj}^{*}(-A)=-\mathrm{pj}^{*}A,\qquad\mathrm{pj}^{*}\overline{f}=\overline{\mathrm{pj}^{*}f},\qquad\mathrm{pj}^{*}(\mathrm{Mx}\cdot f)=\mathrm{Mx}\cdot\mathrm{pj}^{*}f.
\]

\textbf{(1.8) $\bar{H}^{\mathrm{pj}}_k$.}
Let $A\in\tilde{G}_k(T_1)$.
Then $-A\in\tilde{G}_k(T_1)$: by \cite[Lemma~17.6]{QEDI} (see also \cite[Lemma~5.19]{QEDI}), the set $\tilde{G}_k$ is invariant under $A\mapsto-A$, in the setting $k\ge0$, $\ell\in\mathrm{Cl}$, $e_k\in(0,e_0]$.
So $H^{\mathrm{pj}}_k(s'',-A)$ is defined.
For $f$ in the source space on $T_1$, the relations (1.7) with $\mathrm{Mx}=\mathsf{C}^{-1}$ and with $\mathrm{Mx}=\mathsf{C}$, together with the definition of $H^{\mathrm{pj}}_k$, give
\begin{align*}
\mathrm{pj}^{*}\left(\mathsf{C}^{-1}H^{\mathrm{pj}}_k(s'',-A)\,\mathsf{C}f\right)&=\mathsf{C}^{-1}\,\mathrm{pj}^{*}\left(H^{\mathrm{pj}}_k(s'',-A)\,\mathsf{C}f\right)\\
&=\mathsf{C}^{-1}H^{(2)}_k\left(\mathrm{pj}^{*}s'',\mathrm{pj}^{*}(-A)\right)\mathrm{pj}^{*}(\mathsf{C}f)\\
&=\mathsf{C}^{-1}H^{(2)}_k\left(\mathrm{pj}^{*}s'',-\mathrm{pj}^{*}A\right)\mathsf{C}\,\mathrm{pj}^{*}f\\
&=\bar{H}^{(2)}_k\left(\mathrm{pj}^{*}s'',\mathrm{pj}^{*}A\right)\mathrm{pj}^{*}f,
\end{align*}
where the last equality is the definition of $\bar{H}^{(2)}_k$ in (p3) of Section~G4.0.4.

\textbf{(1.9) The definition.}
By (1.1) and (1.2), $\mathrm{pj}^{*}$ is injective with the image $\mathrm{Per}$.
By the first sentence of Proposition~26.3\,(a), the right-hand sides of the two defining relations, the one for $\mathrm{pj}^{*}\left(\mathrm{Op}^{\mathrm{pj}}(s,A)f\right)$ and the one for $\mathrm{pj}^{*}\left(\mathrm{Op}^{\mathrm{pj}}(X,A)f\right)$, lie in $\mathrm{Per}$.
That sentence concerns $\widetilde{\mathrm{Op}}$ alone: it follows from W14$^{\mathrm{ch}}$ (Section~3.2) and does not involve the projection.
Hence each of the two relations determines exactly one object $\mathrm{Op}^{\mathrm{pj}}(\,\cdot\,)f$ on $T_1$, and the map $f\mapsto\mathrm{Op}^{\mathrm{pj}}(\,\cdot\,)f$ is linear,
being the composition of $\mathrm{pj}^{*}$, of the linear operator $\widetilde{\mathrm{Op}}(\,\cdot\,)$ and of the inverse of $\mathrm{pj}^{*}$ on $\mathrm{Per}$.

\emph{Level~$0$.}
At level $0$ the system on $T_2$ consists of a family (p2), of $\Gamma_0(s'',A)$, of $S_{0,y}(X,A)$ and of the diagonal families (Section~G4.0.4).
The relation for $\mathrm{Op}^{\mathrm{pj}}(s,A)$ is applied to $H_{0}(s')$, $C_{0}^{1/2}(s')$, $H_{0}(s'',A)$ and $\Gamma_0(s'',A)$, and the relation for $\mathrm{Op}^{\mathrm{pj}}(X,A)$ to $S_{0,y}(X,A)$ only.
For $H_{0}(s'',A)$, $\Gamma_0$ and $S_{0,y}$ the objects $f$ and $\mathrm{Op}^{\mathrm{pj}}(\,\cdot\,)f$ are elements of $\mathfrak{S}^{\mathrm{ch}}(1,\ell_1)$, that is, functions on $\mathbb{Z}^3\times O_{\mathrm{cub}}$.
\end{proof}

\subsection*{2. Proposition 26.2}

\begin{proof}[Proof of Proposition~26.2]
In this proof $\mathrm{pr}\colon\mathbb{L}_k\to T_k$ is the covering map and $d$ is the sup distance of $T_k$.
Let $\mathrm{hs}$ be the half side of the cubes of the kind under consideration:
the cube of $\pi_k$ with the centre $z$ is $\left\{c:d(c,z)<\frac{M}{2}\right\}$,
the cube of $\pi'_k$ with the centre $z$ is $\{c:d(c,z)\le M\}$, $z\in M\mathbb{Z}^3/\ell\mathbb{Z}^3$,
and the cube of $\pi''_k$ with the centre $z$ is $\{c:d(c,z)\le L\cdot M\}$, $z\in LM\mathbb{Z}^3/\ell\mathbb{Z}^3$
(\cite[Definition~5.8]{QEDI}; (g0) of \cite[Section~17.4.1]{QEDI});
thus $\mathrm{hs}$ is $\frac{M}{2}$, $M$ or $L\cdot M$, respectively.
A lattice site is a site of $T_k$.
A link-component of $Y$ is the union of the cubes of a link-component of $\mathscr{Y}$ (Section~I.1.1).

\textbf{Step 0 (reduction to $\mathbb{L}_k$).}
Let $c_{\circ}$ be a point of $\operatorname{supp}t$, let $\tilde{c}_{\circ}$ be a lift of $c_{\circ}$, and put
\begin{gather*}
\tilde{Z}:=\left\{\tilde{c}\in\mathbb{L}_k:\ |\tilde{c}-\tilde{c}_{\circ}|\le2L\cdot M+L\right\},\\
Z:=\mathrm{pr}(\tilde{Z})=\left\{c\in T_k:\ d(c,c_{\circ})\le2L\cdot M+L\right\},\qquad\mathrm{sd}:=4L\cdot M+2L,
\end{gather*}
the second description of $Z$ being a consequence of (0.c) for the side $\ell$.
The sup-diameter of $\tilde{Z}$ is $\mathrm{sd}$, and
\[
\text{sup-diameter}(\tilde{Z})+\mathrm{sd}=8L\cdot M+4L<13L\cdot M\le\ell,
\]
because $4L<5L\cdot M$ and $\ell\in\mathrm{Cl}$.
Hence $\tilde{Z}$ is a lift of $Z$ in the sense of Section~I.1.1, and $\mathrm{pr}$ is injective on $\tilde{Z}$.
For $c,c'\in Z$ with their lifts $\tilde{c},\tilde{c}'\in\tilde{Z}$ one has $d(c,c')\le|\tilde{c}-\tilde{c}'|\le\mathrm{sd}$, and Section~I.1.1 yields
\begin{equation}
\begin{aligned}
&|\tilde{c}-\tilde{c}'|=d(c,c'),\ \text{and }\tilde{c}-\tilde{c}'\text{ is the only representative of }c-c'\\
&\text{modulo }\ell\mathbb{Z}^3\text{ of sup norm at most }\mathrm{sd}.
\end{aligned}
\tag{2.a}
\end{equation}
\begin{itemize}
\item[(rd1)] The support of $t$ has sup-diameter at most $L$ (Section~0.2); so $\operatorname{supp}t\subset Z$.
If $t$ is a bond with the end points $c$, $c'$, and $c'-c=a\cdot e_{\mu}$ modulo $\ell\mathbb{Z}^3$ with $a\in\{\eta_k,1\}$, then (2.a) gives $\tilde{c}'-\tilde{c}=a\cdot e_{\mu}$, because $|a\cdot e_{\mu}|\le1\le\mathrm{sd}$; thus the lifted end points are $\tilde{c}$ and $\tilde{c}+a\cdot e_{\mu}$.
If $t=\Upsilon$, any two points $c$, $c'$ of $\operatorname{supp}\Upsilon$ satisfy $d(c,c')\le L-1$ (Section~0.2); hence $|\tilde{c}-\tilde{c}'|\le L-1$.
\item[(rd2)] Let $\mathrm{pc}$ be a cube of the kind, with the centre $z$, that contains a point $c$ of $\operatorname{supp}t$.
Then $d(z,c_{\circ})\le d(z,c)+d(c,c_{\circ})\le\mathrm{hs}+L\le2L\cdot M+L$, so that $z\in Z$.
The lift $\tilde{z}\in\tilde{Z}$ of $z$ lies in $M\mathbb{Z}^3$ (in $LM\mathbb{Z}^3$ for $\pi''_k$), because $\ell$ is a multiple of $L\cdot M$.
For a point $c'$ of $Z$, (2.a) shows that $c'$ lies in $\mathrm{pc}$ if and only if $|\tilde{c}'-\tilde{z}|\le\mathrm{hs}$ (for $\pi_k$: $|\tilde{c}'-\tilde{z}|<\mathrm{hs}$).
\end{itemize}

\textbf{(a)}
Let $\mathrm{pc}_1$ and $\mathrm{pc}_2$ be cubes of $\mathscr{Y}$ that contain a point of $\operatorname{supp}t$, with the centres $z_1$, $z_2$ and the lifts $\tilde{z}_1$, $\tilde{z}_2$ of (rd2).

\emph{Kind $\pi_k$.}
Here $t$ is a unit site or a fermion site, and $\operatorname{supp}t$ consists of one photon site ((I.POS): $\operatorname{supp}x=\{c\}$).
The $M$-cubes partition the sites of $T_k$ ((g0) of \cite[Section~17.4.1]{QEDI}).
Hence $\mathrm{pc}_1=\mathrm{pc}_2$.

\emph{Kind $\pi'_k$.}
Here $t$ is a fine bond ($a=\eta_k$) or a unit bond ($a=1$), with the lifted end points $\tilde{c}$ and $\tilde{c}'=\tilde{c}+a\cdot e_{\mu}$ (rd1).
If $\mathrm{pc}_1$ and $\mathrm{pc}_2$ contain the same end point, they have a lattice site in common, and so they are equal or linked \cite[Definition~5.8]{QEDI}.
Otherwise, after the names $\mathrm{pc}_1$, $\mathrm{pc}_2$ have been exchanged if necessary, $\mathrm{pc}_1$ contains $c$ and not $c'$, and $\mathrm{pc}_2$ contains $c'$ and not $c$.
By (rd2), $|\tilde{c}-\tilde{z}_1|\le M$ and $|\tilde{c}+a\cdot e_{\mu}-\tilde{z}_1|>M$.
The coordinates of $\tilde{c}$ and of $\tilde{c}+a\cdot e_{\mu}$ other than the $\mu$-th one coincide.
Hence $|\tilde{c}_{\mu}+a-\tilde{z}_{1,\mu}|>M$, while $|\tilde{c}_{\mu}-\tilde{z}_{1,\mu}|\le M$.
As $\tilde{c}_{\mu}+a-\tilde{z}_{1,\mu}\ge-M+a>-M$, it follows that $\tilde{c}_{\mu}+a-\tilde{z}_{1,\mu}>M$, and therefore
\[
\tilde{c}_{\mu}\le\tilde{z}_{1,\mu}+M<\tilde{c}_{\mu}+a.
\]
The three numbers are multiples of $a$: for $a=\eta_k$, $\tilde{c}_{\mu}$ lies in $\eta_k\mathbb{Z}$ and the integer $\tilde{z}_{1,\mu}+M$ is a multiple of $\eta_k=L^{-k}$; for $a=1$ all three are integers.
So $\tilde{c}_{\mu}=\tilde{z}_{1,\mu}+M$.
By (rd2) for $\mathrm{pc}_2$, $|\tilde{c}+a\cdot e_{\mu}-\tilde{z}_2|\le M$ and $|\tilde{c}-\tilde{z}_2|>M$; hence $|\tilde{c}_{\mu}+a-\tilde{z}_{2,\mu}|\le M$ and $|\tilde{c}_{\mu}-\tilde{z}_{2,\mu}|>M$.
As $\tilde{c}_{\mu}-\tilde{z}_{2,\mu}\le M-a<M$, it follows that $\tilde{c}_{\mu}-\tilde{z}_{2,\mu}<-M$, and therefore
\[
\tilde{c}_{\mu}<\tilde{z}_{2,\mu}-M\le\tilde{c}_{\mu}+a.
\]
These three numbers are multiples of $a$ as well; so $\tilde{c}_{\mu}+a=\tilde{z}_{2,\mu}-M$.
Subtracting the two equalities gives $\tilde{z}_{2,\mu}-\tilde{z}_{1,\mu}=2M+a$.
The left-hand side is a multiple of $M$, whereas the right-hand side lies strictly between $2M$ and $3M$, because $0<a\le1<M$.
This is a contradiction.
Hence $\mathrm{pc}_1$ and $\mathrm{pc}_2$ have $c$ or $c'$ in common, and they are equal or linked.

\emph{Kind $\pi''_k$.}
Here $t=\Upsilon$.
Let $c$ and $c'$ be points of $\operatorname{supp}\Upsilon$ with $c\in\mathrm{pc}_1$ and $c'\in\mathrm{pc}_2$.
By (rd1) and (rd2), for every direction $\mu$,
\[
|\tilde{z}_{1,\mu}-\tilde{z}_{2,\mu}|\le|\tilde{z}_{1,\mu}-\tilde{c}_{\mu}|+|\tilde{c}_{\mu}-\tilde{c}'_{\mu}|+|\tilde{c}'_{\mu}-\tilde{z}_{2,\mu}|\le2L\cdot M+L-1<3L\cdot M.
\]
Consequently this multiple of $L\cdot M$ is $0$, $\pm L\cdot M$ or $\pm2L\cdot M$.
Let $\tilde{z}$ be the point with the coordinates
\[
\tilde{z}_{\mu}:=\begin{cases}\frac{1}{2}\left(\tilde{z}_{1,\mu}+\tilde{z}_{2,\mu}\right)&\text{if }\tilde{z}_{1,\mu}-\tilde{z}_{2,\mu}\in\{0,2L\cdot M,-2L\cdot M\},\\ \tilde{z}_{1,\mu}&\text{if }\tilde{z}_{1,\mu}-\tilde{z}_{2,\mu}\in\{L\cdot M,-L\cdot M\}.\end{cases}
\]
Then $\tilde{z}$ lies in $LM\mathbb{Z}^3$, a subset of $\mathbb{L}_k$, and $|\tilde{z}-\tilde{z}_1|\le L\cdot M$, $|\tilde{z}-\tilde{z}_2|\le L\cdot M$.
By (0.c) for the side $\ell$, $d(\mathrm{pr}\,\tilde{z},z_j)\le|\tilde{z}-\tilde{z}_j|\le L\cdot M$ for $j=1,2$; thus $\mathrm{pr}\,\tilde{z}$ is a lattice site of $\mathrm{pc}_1$ and of $\mathrm{pc}_2$.
Hence the two cubes are equal or linked.

\textbf{(b)}
As $\operatorname{supp}t\subset Y$, every point of $\operatorname{supp}t$ lies in a cube of $\mathscr{Y}$.
By (a), the cubes of $\mathscr{Y}$ that contain a point of $\operatorname{supp}t$ are pairwise equal or linked; therefore
\begin{itemize}
\item[(b1)] the cubes of $\mathscr{Y}$ that contain a point of $\operatorname{supp}t$ belong to one link-component $\mathscr{Y}_1$ of $\mathscr{Y}$.
\end{itemize}
These cubes cover $\operatorname{supp}t$.
So $t$ is contained in the union $Y_1$ of $\mathscr{Y}_1$, which is a link-component of $Y$.
Different link-components of $Y$ have no lattice site in common ((g3) of \cite[Section~17.4.1]{QEDI}), and $\operatorname{supp}t$ is not empty.
Hence $t$ is contained in no other link-component of $Y$.

\textbf{The final assertion of Proposition~26.2.}
If $\operatorname{supp}t\subset Y$, then $t$ is contained in a link-component of $Y$, by (b).
If $t$ is contained in a link-component $Y_1$ of $Y$, then $\operatorname{supp}t\subset Y_1\subset Y$.
Hence an index of the kind is contained in no link-component of $Y$ if and only if its support is not contained in $Y$.

The case $k=0$ is included in the argument above: there $\eta_0=1$, the fine bonds and the unit bonds coincide, so that $a=1$, and the indices of the kind $\pi_0$ are the fermion sites of $T^{\mathrm{f}}_0$ and the unit sites.
\end{proof}

\subsection*{3. Proposition~26.3}

This section is devoted to the proof of Proposition~26.3.
The constants are subject to the conditions listed in Section~5.

For $s\in\mathrm{Pd}_k(T_1)$ and $A\in\tilde{G}_k(T_1)$ put, throughout the section,
\[
s_2:=\mathrm{pj}^{*}s,\qquad A_2:=\mathrm{pj}^{*}A;
\]
by (1.4) and (1.5), $s_2$ lies in $\mathrm{Pd}_k(T_2)$ and $A_2$ lies in $\tilde{G}_k(T_2)$.
The letters $t$, $t'$ always denote indices of the covering lattice in the sense of Section~0.2.

\subsubsection*{3.0. The families, their spaces, indices and weights (Section~G4.0.4)}

For each family of a system of level $k$ on $T_i$ the following list records the source space and the target space, the character $\boldsymbol{\chi}$, the column indices $t'$ with their weights $\nu(t')$, the row indices $t$, and the kind.

\begin{description}
\item[(p1) $H_k(s')$ ($k\ge1$), $H_0(s')$ ($k=0$).]
Source: the bond functions on $T^0_i$; target: the bond functions on $T_i$; $\boldsymbol{\chi}=1$;
columns: the unit bonds, with $\nu(t')=1$; rows: the fine bonds; kind $\pi'_k$.
\item[(p2) $C_k^{1/2}(s')$ ($k\ge0$).]
Source and target: $V_k(T_i)$; $\boldsymbol{\chi}=1$;
columns: the basis elements, with $\nu(t')=1$; rows: the basis elements; kind $\pi''_k$.
\item[(p3) $H_k(s'',A)$, $k\ge1$.]
Source: $\mathfrak{S}(1,\ell_i)$; target: $\mathfrak{S}^{\mathrm{ch}}(\eta_k,\ell_i)$; $\boldsymbol{\chi}(w)=(-1)^{w_1+w_2+w_3}$;
columns: the unit sites, with $\nu(t')=1$; rows: the fermion sites; kind $\pi_k$.
\item[(p3) $H_0(s'',A)$, $k=0$.]
Source and target: $\mathfrak{S}^{\mathrm{ch}}(1,\ell_i)$; $\boldsymbol{\chi}(w)=(-1)^{w_1+w_2+w_3}$;
columns: the fermion sites of $\mathbb{L}^{\mathrm{f}}_0$, with $\nu(t')=1/48$; rows: the fermion sites; kind $\pi_0$.
\item[$\Gamma_k(s'',A)$, $k\ge1$.]
Source and target: $\mathfrak{S}(1,\ell_i)$; $\boldsymbol{\chi}(w)=(-1)^{w_1+w_2+w_3}$;
columns: the unit sites, with $\nu(t')=1$; rows: the unit sites; kind $\pi_k$.
\item[$\Gamma_0(s'',A)$, $k=0$.]
Source and target: $\mathfrak{S}^{\mathrm{ch}}(1,\ell_i)$; $\boldsymbol{\chi}(w)=(-1)^{w_1+w_2+w_3}$;
columns: the fermion sites of $\mathbb{L}^{\mathrm{f}}_0$, with $\nu(t')=1/48$; rows: the fermion sites; kind $\pi_0$.
\item[$M_k(X,A)$, $k\ge1$.]
Source and target: $\mathfrak{S}(1,\ell_i)$; $\boldsymbol{\chi}(w)=(-1)^{w_1+w_2+w_3}$;
columns: the unit sites, with $\nu(t')=1$; rows: the unit sites; no kind.
\item[$S_{k,y}(X,A)$, $k\ge0$.]
Source and target: $\mathfrak{S}^{\mathrm{ch}}(\eta_k,\ell_i)$; $\boldsymbol{\chi}(w)=(-1)^{w_1+w_2+w_3}$;
columns: the fermion sites, with $\nu(t')=\eta_k^3/48$; rows: the fermion sites; no kind.
\end{description}

In each case the source space and the target space have the same character.
For a family with parameters, every row index and every column index belongs to the kind of the family in the sense of Proposition~26.2.

\subsubsection*{3.1. Geometry of the covering map $\mathrm{pj}$}

\begin{description}
\item[(cv1)] \emph{Near points.}
Let $\tilde{c}$ be a site of $\mathbb{L}_k$ and $c'$ a site of $T_1$.
By (0.c) the site $c'$ has a lift $\tilde{c}'$ with $|\tilde{c}'-\tilde{c}|=d_1(\mathrm{pr}_1\tilde{c},c')$, and for this lift
\[
d_2(\mathrm{pr}_2\tilde{c},\mathrm{pr}_2\tilde{c}')=d_1(\mathrm{pr}_1\tilde{c},c').
\]
Indeed, (0.c) for $i=2$ gives $d_2(\mathrm{pr}_2\tilde{c},\mathrm{pr}_2\tilde{c}')\le|\tilde{c}'-\tilde{c}|$, and the reverse inequality is (ds2).
\item[(cv2)] \emph{Faces and links.}
Two $M$-cubes of $T_i$ have a common face if and only if their centres differ, modulo $\ell_i\mathbb{Z}^3$, by $Me_\mu$ or by $-Me_\mu$ for some direction $\mu$ (Section~I.1.1).
If this is the case for the centres $z'$, $z''$ of two $M$-cubes of $T_2$, then it is the case modulo $\ell_1\mathbb{Z}^3$ for $\mathrm{pj}\,z'$ and $\mathrm{pj}\,z''$, so that the two image cubes have a common face.
If two cubes of $\pi'_k$ (resp.\ of $\pi''_k$) of $T_2$ have the lattice site $c$ in common, then, by (ds4), $\mathrm{pj}\,c$ is a lattice site of both image cubes.
Hence $\mathrm{pj}$ maps two linked cubes of one kind to two cubes that are equal or linked, and it maps a link-connected set of cubes to a link-connected set of cubes.
\item[(cv3)] \emph{Preimages.}
Let $\mathscr{Y}$ be a set of parameter cubes of one kind of $T_1$ and $Y$ its union, let
\[
\mathscr{Y}':=\left\{\mathrm{pc}':\ \mathrm{pc}'\text{ is a parameter cube of this kind of }T_2\text{ and }\mathrm{pj}\,\mathrm{pc}'\in\mathscr{Y}\right\},
\]
and let $Y'$ be the union of $\mathscr{Y}'$.
Then $Y'=\mathrm{pj}^{-1}(Y)$.
Indeed, if $c$ lies in a cube $\mathrm{pc}'$ of $\mathscr{Y}'$, then $\mathrm{pj}\,c$ lies in $\mathrm{pj}\,\mathrm{pc}'$ by (ds4), and this is a cube of $\mathscr{Y}$.
Conversely, let $\mathrm{pj}\,c$ lie in the cube $\mathrm{pc}$ of $\mathscr{Y}$ with the centre $z$, and let $\tilde{c}$ be a lift of $c$.
By (cv1), $z$ has a lift $\tilde{z}$ with $|\tilde{z}-\tilde{c}|=d_1(\mathrm{pj}\,c,z)$;
it lies in $M\mathbb{Z}^3$, resp.\ in $LM\mathbb{Z}^3$ for the kind $\pi''_k$, because $\ell_1$ is a multiple of $LM$.
The cube $\mathrm{pc}'$ of $T_2$ with the centre $\mathrm{pr}_2\tilde{z}$ satisfies $\mathrm{pj}\,\mathrm{pc}'=\mathrm{pc}$ by (ds4), hence belongs to $\mathscr{Y}'$;
and it contains $c$, because $d_2(c,\mathrm{pr}_2\tilde{z})=d_1(\mathrm{pj}\,c,z)$.
\item[(cv4)] \emph{Containment; bonds.}
One has $\mathrm{pr}_1(\operatorname{supp}t)=\mathrm{pj}\left(\mathrm{pr}_2(\operatorname{supp}t)\right)$.
Consequently, if $t$ is contained in a set $Z'$ of sites of $T_2$, then $t$ is contained in $\mathrm{pj}\,Z'$ on $T_1$;
and for a set $Y$ of sites of $T_1$, the index $t$ is contained in $Y$ on $T_1$ if and only if it is contained in $\mathrm{pj}^{-1}(Y)$ on $T_2$.
A bond $b'$ of $T_2$ is of the form $\mathrm{pr}_2(\widetilde{\mathrm{bd}})$ with a bond $\widetilde{\mathrm{bd}}$ of $\mathbb{L}_k$;
then $\mathrm{pj}\,b':=\mathrm{pr}_1(\widetilde{\mathrm{bd}})$ is a bond of $T_1$, and $\left(\mathrm{pj}^{*}A\right)(b')=A(\mathrm{pj}\,b')$.
If $b'$ has an end point in $Z'$, then $\mathrm{pj}\,b'$ has an end point in $\mathrm{pj}\,Z'$;
thus $\mathrm{pj}$ maps the bonds that meet $Z'$ to bonds that meet $\mathrm{pj}\,Z'$.
\item[(cv5)] \emph{Polymers.}
Let $X'$ be a polymer of $T_2$.
By (ds4), $\mathrm{pj}\,X'$ is a non-empty union of $M$-cubes.
Any two $M$-cubes of $X'$ are joined by a chain of $M$-cubes of $X'$ in which consecutive cubes have a common face;
by (cv2) the images form such a chain in $\mathrm{pj}\,X'$.
Therefore $\mathrm{pj}\,X'\in D_k(T_1)$, and $D_k(T_2)$ is the disjoint union of the sets $\{X':\mathrm{pj}\,X'=X\}$, $X\in D_k(T_1)$.
The set $D_k(T_2)$ is finite.
For $v\in\ell_1\mathbb{Z}^3$ one has $\mathrm{pj}(X'+v)=\mathrm{pj}\,X'$.
\item[(cv6)] \emph{The number $d_M$.}
For every $X'\in D_k(T_2)$,
\[
d_M(\mathrm{pj}\,X')\le d_M(X').
\]
The number $d_M$ is defined in \cite[Section~13.1.7]{QEDI}, and the conventions of Section~I.1.1 agree with that definition:
a closed straight segment of $\mathbb{R}^3/\ell_i\mathbb{Z}^3$ is the image of a closed straight segment of $\mathbb{R}^3$ and has the Euclidean length of the latter, and the length of a net is the sum of the lengths of the segments in its writing as a finite union of segments.
By this definition, $M\cdot d_M(X)$ is the infimum of the lengths of the nets that meet every closed $M$-cube of $X$, the infimum extending over the nets and over their writings as finite unions of segments.
Let $\mathrm{prc}_i:\mathbb{R}^3\to\mathbb{R}^3/\ell_i\mathbb{Z}^3$ and $\mathrm{pjc}:\mathbb{R}^3/\ell_2\mathbb{Z}^3\to\mathbb{R}^3/\ell_1\mathbb{Z}^3$ be the covering maps, so that $\mathrm{pjc}\circ\mathrm{prc}_2=\mathrm{prc}_1$.
Let $\mathrm{Sgs}$ be a finite set of closed straight segments of $\mathbb{R}^3$ such that the union of the sets $\mathrm{prc}_2(\mathrm{Seg})$, $\mathrm{Seg}\in\mathrm{Sgs}$, is a net that meets every closed $M$-cube of $X'$.
The image of this first net under $\mathrm{pjc}$ is the union of the sets $\mathrm{prc}_1(\mathrm{Seg})$, $\mathrm{Seg}\in\mathrm{Sgs}$;
as the continuous image of a connected set it is connected, so that it is a net, written with segments of the same lengths; call it the second net.
The closed $M$-cube of $\mathbb{R}^3/\ell_i\mathbb{Z}^3$ with the centre $\mathrm{pr}_i\tilde{z}$, $\tilde{z}\in M\mathbb{Z}^3$, is
\[
\mathrm{prc}_i\left\{y\in\mathbb{R}^3:|y-\tilde{z}|\le\frac{M}{2}\right\},
\]
by (0.c) read for points of $\mathbb{R}^3$.
So $\mathrm{pjc}$ maps the closed $M$-cube with the centre $z'$ onto the closed $M$-cube with the centre $\mathrm{pj}\,z'$;
and by (ds4) the $M$-cubes of $\mathrm{pj}\,X'$ are the images of the $M$-cubes of $X'$.
Hence every closed $M$-cube of $\mathrm{pj}\,X'$ is the image under $\mathrm{pjc}$ of a closed $M$-cube of $X'$, and the image of a point of the first net that lies in the latter is a point of the second net that lies in the former.
It follows that $M\cdot d_M(\mathrm{pj}\,X')$ is at most the length of the first net;
the infimum of these lengths over all such nets and writings is $M\cdot d_M(X')$.
\item[(cv7)] \emph{The $M$-cubes that meet $K_c$.}
Let $\tilde{c}$ be a site of $\mathbb{L}_k$ and
\[
\tilde{K}:=\left\{\tilde{c}'\in\mathbb{L}_k:|\tilde{c}'-\tilde{c}|\le2L\right\}.
\]
By (0.c), $\mathrm{pr}_i(\tilde{K})=K_{\mathrm{pr}_i\tilde{c}}$.
An $M$-cube of $T_i$ meets $K_{\mathrm{pr}_i\tilde{c}}$ if and only if it is $\mathrm{pr}_i(\tilde{\square})$ for an $M$-cube $\tilde{\square}$ of $\mathbb{L}_k$ that meets $\tilde{K}$:
if it contains $\mathrm{pr}_i\tilde{c}'$ with $\tilde{c}'\in\tilde{K}$, then by (ds4) it is the image of the $M$-cube of $\tilde{c}'$.
The centres of two $M$-cubes of $\mathbb{L}_k$ that meet $\tilde{K}$ differ by less than $2\left(2L+\frac{M}{2}\right)=4L+M<\ell_i$, so that different ones have different images.
Hence $\mathrm{pr}_i$ maps the set of the $M$-cubes of $\mathbb{L}_k$ that meet $\tilde{K}$ bijectively onto the set of the $M$-cubes of $T_i$ that meet $K_{\mathrm{pr}_i\tilde{c}}$,
and $\mathrm{pj}$ maps the latter set for $i=2$ bijectively onto the latter set for $i=1$.
\item[(cv8)] \emph{Level $0$: the cubes of $\pi'_0$ that contain a bond.}
Let $k=0$, and let $\widetilde{\mathrm{bd}}=(\tilde{c},\tilde{c}+e_\mu)$ be a bond of $\mathbb{Z}^3$.
The cube of $\pi'_0(T_i)$ with the centre $z$ contains $\mathrm{pr}_i\widetilde{\mathrm{bd}}$ (that is, both end points) if and only if
\[
z=\mathrm{pr}_i\tilde{z}\quad\text{with}\quad \tilde{z}\in M\mathbb{Z}^3,\qquad |\tilde{z}-\tilde{c}|\le M,\qquad |\tilde{z}-\tilde{c}-e_\mu|\le M.
\]
Indeed, let $d_i(z,\mathrm{pr}_i\tilde{c})\le M$ and $d_i(z,\mathrm{pr}_i\tilde{c}+e_\mu)\le M$.
By (0.c) there are a lift $\tilde{z}$ of $z$ with $|\tilde{z}-\tilde{c}|\le M$ and a vector $\mathrm{dk}\in\mathbb{Z}^3$ with $|\tilde{z}+\ell_i\,\mathrm{dk}-\tilde{c}-e_\mu|\le M$.
Then $\ell_i|\mathrm{dk}|\le2M+1<\ell_i$, whence $\mathrm{dk}=0$.
The converse follows from (0.c).
Two such points $\tilde{z}$ differ by at most $2M<\ell_i$.
Hence $\mathrm{pr}_i$ maps the set of these $\tilde{z}$ bijectively onto the set of the centres of the cubes of $\pi'_0(T_i)$ that contain $\mathrm{pr}_i\widetilde{\mathrm{bd}}$,
and $\mathrm{pj}$ maps the cubes of $\pi'_0(T_2)$ that contain $\mathrm{pr}_2\widetilde{\mathrm{bd}}$ bijectively onto the cubes of $\pi'_0(T_1)$ that contain $\mathrm{pr}_1\widetilde{\mathrm{bd}}$.
\end{description}

\subsubsection*{3.2. Proof of (a)}

\begin{proof}[Proof of Proposition~26.3\,(a)]
In the argument below the only property of the system on $T_2$ that enters is W14$^{\mathrm{ch}}$, and it enters only at the elements $\tilde{g}_v=(1_4,v)$, $v\in\ell_1\mathbb{Z}^3$, defined in the next paragraph.
Hence the argument proves the assertion of Proposition~26.3 on clause~(a):
clause~(a), together with (8.0), holds for every system on $T_2$ with W14$^{\mathrm{ch}}$, and it requires W14$^{\mathrm{ch}}$ only at the elements $\tilde{g}_v$, $v\in\ell_1\mathbb{Z}^3$.

\emph{The translations as elements of the step group.}
Let $v\in\ell_1\mathbb{Z}^3$ and $\tilde{g}_v:=(1_4,v\bmod\ell_2\mathbb{Z}^3)$.
The matrix $1_4$ commutes with $\gamma_5$, hence lies in $\mathrm{Sp}_3$, and $r(1_4)=1$;
moreover $v\in LM\mathbb{Z}^3$, because $\ell_1$ is a multiple of $LM$.
Consequently $\tilde{g}_v$ belongs to $\mathbb{G}^{L}_{k,3}(T_2)$, which is a subgroup of $\mathbb{G}_k^{L}(T_2)$ (Section~I.4.2).
For $\tilde{g}_v$, that is for $S=1_4$, with the representative $v$ of the translation part taken as the vector $a$, the actions of $\tilde{g}=(S,a)$ defined in Section~I.4.3 (on fermion sites by (I.4.2)) give
\begin{gather*}
\tilde{g}_v\,c=c+v,\qquad \tilde{g}_v\,(c,g)=(c+v,g),\qquad (\tilde{g}_vA)(c,c')=A(c-v,c'-v),\\
(\tilde{g}_vf)(x)=1_4\,f(x-v)\ \text{on }\mathfrak{S}^{\mathrm{ch}},\qquad (\tilde{g}_vf)(y)=f(y-v)\ \text{on }\mathfrak{S},\\
(\tilde{g}_vs)_{\mathrm{pc}}=s_{\mathrm{pc}-v},\qquad \tilde{g}_v\,X'=X'+v;
\end{gather*}
in the second identity the label is multiplied by $r=1$.
On $V_k$ the element $\tilde{g}_v$ acts on the bond function $\mathsf{Z}$, and in coordinates
\[
(\tilde{g}_v\mathsf{Z})(\Upsilon)=\sum_b\mathsf{Z}(b-v)\cdot\Upsilon(b)=\sum_b\mathsf{Z}(b)\cdot\Upsilon(b+v)=\mathsf{Z}(\Upsilon-v),
\]
where $\Upsilon-v$ is again a basis element (Section~0.2).
Thus $\tilde{g}_v$ acts as $T_v$ on each of the spaces of Section~3.0, on connections and on parameter families.

\emph{Commutation.}
The parameter family $s_2$ and the connection $A_2$ are $\ell_1\mathbb{Z}^3$-periodic: $T_vs_2=s_2$ and $T_vA_2=A_2$.
For (p1)--(p3) and $\Gamma_k$ (for (p1), (p2) without the argument $A$), property W14$^{\mathrm{ch}}$ on $T_2$ at $\tilde{g}=\tilde{g}_v$ reads
\[
\mathrm{Op}^{(2)}(T_vs_2,T_vA_2)\,T_v=T_v\,\mathrm{Op}^{(2)}(s_2,A_2),
\]
and therefore
\[
\widetilde{\mathrm{Op}}(s,A)\,T_v=T_v\,\widetilde{\mathrm{Op}}(s,A).
\]
For $M_k$ ($k\ge1$) and $S_{k,y}$ it reads
\[
\mathrm{Op}^{(2)}(X'+v,T_vA_2)\,T_v=T_v\,\mathrm{Op}^{(2)}(X',A_2).
\]
By (cv5) the map $X'\mapsto X'+v$ is a bijection of $\{X'\in D_k(T_2):\mathrm{pj}\,X'=X\}$ onto itself, with the inverse $X'\mapsto X'-v$.
Summing over this set one obtains
\[
T_v\,\widetilde{\mathrm{Op}}(X,A)=\sum_{X':\ \mathrm{pj}\,X'=X}\mathrm{Op}^{(2)}(X'+v,A_2)\,T_v=\widetilde{\mathrm{Op}}(X,A)\,T_v.
\]

\emph{$\mathrm{Per}$ is mapped into $\mathrm{Per}$.}
An object $f$ on $T_2$ lies in $\mathrm{Per}$ if and only if $T_{-\ell_1w}f=\boldsymbol{\chi}(w)\cdot f$ for all $w\in\mathbb{Z}^3$, because $(T_{-\ell_1w}f)(t)=f(t+\ell_1w)$.
Let $f$ lie in $\mathrm{Per}$ of the source space.
Then
\[
T_{-\ell_1w}\,\widetilde{\mathrm{Op}}(\,\cdot\,)f=\widetilde{\mathrm{Op}}(\,\cdot\,)\,T_{-\ell_1w}f=\boldsymbol{\chi}(w)\,\widetilde{\mathrm{Op}}(\,\cdot\,)f,
\]
and the target space has the same character (Section~3.0).
So $\widetilde{\mathrm{Op}}(\,\cdot\,)f$ lies in $\mathrm{Per}$ of the target space.
Together with (1.9) this shows that the projection is well defined.

\emph{The product in (8.0) depends only on the class of $w$.}
By (kn2) on $T_2$, by $\ell_2=\mathrm{Jr}\,\ell_1$, and by $\boldsymbol{\chi}(\mathrm{Jr}\,e_\mu)=\boldsymbol{\chi}(e_\mu)^{\mathrm{Jr}}=\boldsymbol{\chi}(e_\mu)$,
\begin{align*}
\boldsymbol{\chi}(w+\mathrm{Jr}\,e_\mu)\cdot\widetilde{\mathrm{Op}}(\,\cdot\,)(t,t'+\ell_1w+\ell_2e_\mu)
&=\boldsymbol{\chi}(w)\cdot\boldsymbol{\chi}(e_\mu)\cdot\boldsymbol{\chi}(e_\mu)\cdot\widetilde{\mathrm{Op}}(\,\cdot\,)(t,t'+\ell_1w)\\
&=\boldsymbol{\chi}(w)\cdot\widetilde{\mathrm{Op}}(\,\cdot\,)(t,t'+\ell_1w).
\end{align*}
This holds for all $w\in\mathbb{Z}^3$ and $\mu=1,2,3$, and the vectors $\mathrm{Jr}\,e_\mu$ generate $\mathrm{Jr}\,\mathbb{Z}^3$.
For the spinor families each of the two factors changes its sign under $w\mapsto w+\mathrm{Jr}\,e_\mu$;
for (p1), (p2) one has $\boldsymbol{\chi}=1$, and each factor depends only on the class of $w$.

\emph{The kernel.}
Put
\[
\mathrm{Kr}(t,t'):=\sum_{w\in\mathrm{Rw}}\boldsymbol{\chi}(w)\cdot\widetilde{\mathrm{Op}}(\,\cdot\,)(t,t'+\ell_1w);
\]
by the preceding display this does not depend on the choice of $\mathrm{Rw}$.
Let $f$ be an element of the source space on $T_1$, and let $\mathrm{Fnd}_1$ be a fundamental set of the column indices for $\ell_1\mathbb{Z}^3$.
Applying (kn1) on $T_2$ with the fundamental set of (fs), and using $\nu(t'+\ell_1w)=\nu(t')$ and $f(t'+\ell_1w)=\boldsymbol{\chi}(w)\cdot f(t')$, one finds
\begin{align*}
\left(\mathrm{pj}^{*}\,\mathrm{Op}^{\mathrm{pj}}(\,\cdot\,)f\right)(t)
&=\left(\widetilde{\mathrm{Op}}(\,\cdot\,)\,\mathrm{pj}^{*}f\right)(t)\\
&=\sum_{w\in\mathrm{Rw}}\sum_{t'\in\mathrm{Fnd}_1}\nu(t')\cdot\widetilde{\mathrm{Op}}(\,\cdot\,)(t,t'+\ell_1w)\cdot\boldsymbol{\chi}(w)\cdot f(t')\\
&=\sum_{t'\in\mathrm{Fnd}_1}\nu(t')\cdot\mathrm{Kr}(t,t')\cdot f(t').
\end{align*}
The left side is $\left(\mathrm{Op}^{\mathrm{pj}}(\,\cdot\,)f\right)(t)$.
Hence $\mathrm{Kr}$ has the property (kn1) on $T_1$, with the same weight $\nu$ as on $T_2$;
for columns on fermion sites this weight is $\eta_k^3/48$, the sum over $t'$ being $\eta_k^3\cdot\frac{1}{48}\sum_{\tilde{c}'}\sum_{g'\in O_{\mathrm{cub}}}$.
The function $\mathrm{Kr}$ also has (kn2) on $T_1$.
Indeed, the set $\mathrm{Rw}+e_\mu$ is again a set of representatives, and therefore
\begin{align*}
\mathrm{Kr}(t,t'+\ell_1e_\mu)
&=\sum_{w\in\mathrm{Rw}}\boldsymbol{\chi}(w)\cdot\widetilde{\mathrm{Op}}(\,\cdot\,)\left(t,t'+\ell_1(w+e_\mu)\right)\\
&=\boldsymbol{\chi}(e_\mu)\sum_{w\in\mathrm{Rw}}\boldsymbol{\chi}(w+e_\mu)\cdot\widetilde{\mathrm{Op}}(\,\cdot\,)\left(t,t'+\ell_1(w+e_\mu)\right)
=\boldsymbol{\chi}(e_\mu)\cdot\mathrm{Kr}(t,t').
\end{align*}
This holds for all $t'$ and $\mu=1,2,3$.
Since $\boldsymbol{\chi}$ is a homomorphism with $\boldsymbol{\chi}(e_\mu)^2=1$, it follows that $\mathrm{Kr}(t,t'+\ell_1\mathrm{dk})=\boldsymbol{\chi}(\mathrm{dk})\cdot\mathrm{Kr}(t,t')$ for all $\mathrm{dk}\in\mathbb{Z}^3$.
By the uniqueness assertion of Section~0.5, $\mathrm{Kr}$ is the kernel of $\mathrm{Op}^{\mathrm{pj}}(\,\cdot\,)$, and this is (8.0).

\emph{Level $0$.}
For $\Gamma_0$, $S_{0,y}$ and $H_0(s'',A)$ the indices are $t=(\tilde{c},g)$ and $t'=(\tilde{c}',g')$ in $\mathbb{Z}^3\times O_{\mathrm{cub}}$, with $t'+\ell_1w=(\tilde{c}'+\ell_1w,g')$, and $\nu=1/48$ on both tori;
here (8.0) is an identity between kernels in the sense of (I.1.4).
\end{proof}

\subsubsection*{3.3. Proof of (e)}

\begin{proof}[Proof of Proposition~26.3\,(e)]
No property of a system enters; accordingly clause~(e) holds without any assumption on a system, as asserted in Proposition~26.3.
Let $O^{(2)}$ map the objects on $T_2$ with the index set $\mathrm{Ind}'$ to the objects on $T_2$ with the index set $\mathrm{Ind}$, and let $O^{(1)}$ do the same on $T_1$, both spaces carrying one of the pairings (0.b).
The transpose $O^T$ is defined by $\langle O^Tf_1,f\rangle=\langle f_1,Of\rangle$.
\begin{itemize}
\item[(i)] Suppose that the two spaces have different characters.
Then $O^{(2)}=0$ and $O^{(1)}=0$, so that both sides of the assertion vanish.
Indeed, for $v=\ell_2e_\mu=\ell_1\mathrm{Jr}\,e_\mu$ the operator $T_{-v}$ is the multiplication by $1$ on one of the two spaces and by $-1$ on the other;
hence the relation $O^{(2)}T_{-v}=T_{-v}O^{(2)}$ reads $O^{(2)}=-O^{(2)}$.
Then $\mathrm{pj}^{*}O^{(1)}=O^{(2)}\mathrm{pj}^{*}=0$, and $\mathrm{pj}^{*}$ is injective.
From now on both spaces have the character $\boldsymbol{\chi}$.
\item[(ii)] For objects $f'$, $f$ on $T_2$,
\[
\left\langle T_vf',T_vf\right\rangle_{T_2}=\sum_{t\in\mathrm{Fnd}_2}\nu(t)\cdot f'(t-v)\cdot f(t-v)=\sum_{t\in\mathrm{Fnd}_2-v}\nu(t)\cdot f'(t)\cdot f(t)=\left\langle f',f\right\rangle_{T_2},
\]
because $\mathrm{Fnd}_2-v$ is a fundamental set and (0.b) does not depend on the fundamental set.
So $T_v^T=T_{-v}$.
Transposing $O^{(2)}T_{-v}=T_{-v}O^{(2)}$ one obtains $T_v\left(O^{(2)}\right)^T=\left(O^{(2)}\right)^TT_v$ for all $v\in\ell_1\mathbb{Z}^3$.
\item[(iii)] Let $f_1$ be an object on $T_1$ with the index set $\mathrm{Ind}$.
Then $T_{-\ell_1w}\mathrm{pj}^{*}f_1=\boldsymbol{\chi}(w)\,\mathrm{pj}^{*}f_1$, and by (ii)
\[
T_{-\ell_1w}\left(O^{(2)}\right)^T\mathrm{pj}^{*}f_1=\boldsymbol{\chi}(w)\left(O^{(2)}\right)^T\mathrm{pj}^{*}f_1.
\]
Thus $\left(O^{(2)}\right)^T\mathrm{pj}^{*}f_1$ lies in $\mathrm{Per}$: it equals $\mathrm{pj}^{*}f_2$ for an object $f_2$ on $T_1$ with the index set $\mathrm{Ind}'$.
\item[(iv)] For every object $f$ on $T_1$ with the index set $\mathrm{Ind}'$, (1.3) gives
\begin{align*}
\mathrm{Jr}^3\cdot\left\langle f_2,f\right\rangle_{T_1}
&=\left\langle\mathrm{pj}^{*}f_2,\mathrm{pj}^{*}f\right\rangle_{T_2}
=\left\langle\left(O^{(2)}\right)^T\mathrm{pj}^{*}f_1,\mathrm{pj}^{*}f\right\rangle_{T_2}
=\left\langle\mathrm{pj}^{*}f_1,O^{(2)}\mathrm{pj}^{*}f\right\rangle_{T_2}\\
&=\left\langle\mathrm{pj}^{*}f_1,\mathrm{pj}^{*}O^{(1)}f\right\rangle_{T_2}
=\mathrm{Jr}^3\cdot\left\langle f_1,O^{(1)}f\right\rangle_{T_1}
=\mathrm{Jr}^3\cdot\left\langle\left(O^{(1)}\right)^Tf_1,f\right\rangle_{T_1}.
\end{align*}
The pairing is non-degenerate (Section~0.4).
Therefore $f_2=\left(O^{(1)}\right)^Tf_1$, that is, $\left(O^{(2)}\right)^T\mathrm{pj}^{*}f_1=\mathrm{pj}^{*}\left(O^{(1)}\right)^Tf_1$.
\end{itemize}
\end{proof}

\subsubsection*{3.4. Proof of (d)}

\begin{proof}[Proof of Proposition~26.3\,(d)]
Here, too, no property of a system enters; accordingly clause~(d) holds without any assumption on a system, as asserted in Proposition~26.3.
Let $\tilde{S}:=\operatorname{supp}t$ and $\tilde{S}':=\operatorname{supp}t'$;
both sets have sup-diameter at most $L$ (Section~0.2).
By (ds3) the summand $\exp\left(-\gamma_{\mathrm{w}}\cdot d_2(t,t'+\ell_1w)\right)$ depends only on the class of $w$ modulo $\mathrm{Jr}$.

\emph{Normalisation.}
If $t'$ is replaced by $t'+\ell_1\mathrm{dk}$, the right side of (8.0$'$) does not change, while on the left side the classes $w$ are permuted by $w\mapsto w+\mathrm{dk}$.
By (ds3) the vector $\mathrm{dk}$ can be chosen such that $\mathrm{dist}(\tilde{S},\tilde{S}'+\ell_1\mathrm{dk})=d_1(t,t')$.
It may therefore be assumed that
\[
\mathrm{dist}(\tilde{S},\tilde{S}')=d_1:=d_1(t,t')=\min_{\mathrm{dk}\in\mathbb{Z}^3}\mathrm{dist}(\tilde{S},\tilde{S}'+\ell_1\mathrm{dk}).
\]
Moreover $d_1\le\ell_1/2$, since for $\tilde{c}\in\tilde{S}$ and $\tilde{c}'\in\tilde{S}'$ there is a vector $\mathrm{dk}\in\mathbb{Z}^3$ with $|\tilde{c}-\tilde{c}'-\ell_1\mathrm{dk}|\le\ell_1/2$.

\emph{One vector for each class.}
By (ds3),
\[
d_2(t,t'+\ell_1w)=\min_{\mathrm{dk}\in\mathbb{Z}^3}\mathrm{dist}(\tilde{S},\tilde{S}'+\ell_1w+\ell_2\mathrm{dk})=\min\left\{\mathrm{dist}(\tilde{S},\tilde{S}'+\ell_1w'):\ w'\in w+\mathrm{Jr}\,\mathbb{Z}^3\right\}.
\]
For every class $w$ choose $w'\in w+\mathrm{Jr}\,\mathbb{Z}^3$ with $d_2(t,t'+\ell_1w)=\mathrm{dist}(\tilde{S},\tilde{S}'+\ell_1w')$.
Different classes have different vectors $w'$.

\emph{Lower bounds.}
The definition of $d_1$ gives $\mathrm{dist}(\tilde{S},\tilde{S}'+\ell_1w')\ge d_1$.
Let $\tilde{c}\in\tilde{S}$ and $\tilde{c}'\in\tilde{S}'$ with $|\tilde{c}-\tilde{c}'|=d_1$, and let $\tilde{\tilde{c}}\in\tilde{S}$ and $\tilde{\tilde{c}}'\in\tilde{S}'$ with $|\tilde{\tilde{c}}-\tilde{\tilde{c}}'-\ell_1w'|=\mathrm{dist}(\tilde{S},\tilde{S}'+\ell_1w')$.
Then
\begin{align*}
\ell_1|w'|&\le|\tilde{\tilde{c}}-\tilde{\tilde{c}}'-\ell_1w'|+|\tilde{\tilde{c}}-\tilde{c}|+|\tilde{c}-\tilde{c}'|+|\tilde{c}'-\tilde{\tilde{c}}'|\\
&\le\mathrm{dist}(\tilde{S},\tilde{S}'+\ell_1w')+L+d_1+L.
\end{align*}
Hence, for $|w'|=j\ge2$, using $\ell_1-d_1\ge d_1$ and $j-1\ge1$,
\[
\mathrm{dist}(\tilde{S},\tilde{S}'+\ell_1w')\ge\ell_1j-d_1-2L\ge d_1+\ell_1(j-1)-2L\ge d_1+(\ell_1-2L)(j-1).
\]

\emph{Counting.}
There are $27$ vectors $w'\in\mathbb{Z}^3$ with $|w'|\le1$, and there are $(2j+1)^3-(2j-1)^3=24j^2+2\le26j^2$ vectors with $|w'|=j\ge1$.
Further $\ell_1-2L\ge13LM-2L\ge6M$, so that $\gamma_{\mathrm{w}}\cdot(\ell_1-2L)\ge6\gamma_{\mathrm{w}}\cdot M\ge6$ by $\gamma_{\mathrm{w}}\cdot M\ge1$.
Since $\gamma_{\mathrm{w}}>0$, again by $\gamma_{\mathrm{w}}\cdot M\ge1$, the lower bounds for the distances yield upper bounds for the summands.
Consequently
\[
\sum_{w\in(\mathbb{Z}/\mathrm{Jr}\mathbb{Z})^3}\exp\left(-\gamma_{\mathrm{w}}\cdot d_2(t,t'+\ell_1w)\right)\le\exp\left(-\gamma_{\mathrm{w}}\cdot d_1\right)\cdot\left[27+\sum_{j\ge2}26j^2\cdot\exp\left(-6(j-1)\right)\right].
\]

\emph{The number.}
Put $\mathrm{trm}_j:=26j^2\cdot\exp\left(-6(j-1)\right)$.
For $j\ge2$,
\[
\frac{\mathrm{trm}_{j+1}}{\mathrm{trm}_j}=\left(\frac{j+1}{j}\right)^2\cdot\exp(-6)\le\frac{9}{4}\cdot\exp(-6)\le\frac{1}{100},
\]
and, as $\mathrm{e}>2.7$,
\[
\mathrm{trm}_2=104\cdot\exp(-6)\le\frac{104}{2.7^6}=\frac{104}{387.420489}<0.269.
\]
Hence
\[
\sum_{j\ge2}\mathrm{trm}_j\le\left(1-\frac{1}{100}\right)^{-1}\mathrm{trm}_2<0.272<0.3,
\]
and the bracket is at most $27.3\le64$.
This is (8.0$'$).
\end{proof}

\subsubsection*{3.5. A sum over polymers}

Let $t$ be a unit site or a fermion site of the covering lattice.
The set $\mathrm{pr}_2(\operatorname{supp}t)$ consists of a single photon site of $T_2$;
let $\square_t$ be the $M$-cube of $T_2$ that contains this site, which is unique by the fact (g0) of \cite[Section~17.4.1]{QEDI}.
For every $X\in D_k(T_1)$,
\begin{equation}
\sum_{\substack{X'\in D_k(T_2):\ \mathrm{pj}\,X'=X,\\ X'\supseteq\square_t}}\exp\left(-(\kappa_{\mathrm{w}}+\kappa_{0})\cdot d_M(X')\right)\le K_{0}\cdot\exp\left(-\kappa_{\mathrm{w}}\cdot d_M(X)\right).\tag{3.a}
\end{equation}
To see this, note first that $d_M(X')\ge d_M(X)$ for every term, by (cv6), and that $\kappa_{\mathrm{w}}\ge0$, since $\kappa_{\mathrm{w}}=\kappa+\kappa_{0}+3$ with $\kappa\ge7\kappa_{0}+7$ and $\kappa_{0}=54\log72>0$.
Hence $\exp\left(-\kappa_{\mathrm{w}}\cdot d_M(X')\right)\le\exp\left(-\kappa_{\mathrm{w}}\cdot d_M(X)\right)$, and each term
\[
\exp\left(-\kappa_{\mathrm{w}}\cdot d_M(X')\right)\cdot\exp\left(-\kappa_{0}\cdot d_M(X')\right)
\]
is at most $\exp\left(-\kappa_{\mathrm{w}}\cdot d_M(X)\right)\cdot\exp\left(-\kappa_{0}\cdot d_M(X')\right)$.
Assertion (t2) of \cite[Theorem~13.2]{QEDI}, applied on $T_2$ in the setting of that theorem ($k\ge0$, $\ell_2\in\mathrm{Cl}$), gives
\[
\sum_{X'\in D_k(T_2):\ X'\supseteq\square_t}\exp\left(-\kappa_{0}\cdot d_M(X')\right)\le K_{0},
\]
and (3.a) follows.

\subsubsection*{3.6. Proof of (c)}

\begin{proof}[Proof of Proposition~26.3\,(c)]
Here $k=0$ and $\eta_0=1$.
By Section~G4.0.4 the families (p1) and (p3) of the system on $T_2$ are the diagonal families of Section~I.5.8 on $T_2$.

\emph{The family $H_0(s')$.}
Let $A'$ be a bond function on $T^0_1=T_1$, let $s'\in\mathrm{Pd}_0(T_1)$, and let $\widetilde{\mathrm{bd}}$ be a bond of $\mathbb{Z}^3$.
The display of Section~I.5.8 on $T_2$, the relation $\left(\mathrm{pj}^{*}s'\right)_{\mathrm{pc}'}=s'_{\mathrm{pj}\,\mathrm{pc}'}$ of (1.4) and the bijection of (cv8) give
\begin{align*}
\left(H_0^{(2)}(\mathrm{pj}^{*}s')\,\mathrm{pj}^{*}A'\right)(\widetilde{\mathrm{bd}})
&=\prod\left\{\left(\mathrm{pj}^{*}s'\right)_{\mathrm{pc}'}:\ \mathrm{pc}'\in\pi'_0(T_2),\ \mathrm{pr}_2\widetilde{\mathrm{bd}}\text{ contained in }\mathrm{pc}'\right\}\cdot A'(\widetilde{\mathrm{bd}})\\
&=\prod\left\{s'_{\mathrm{pc}}:\ \mathrm{pc}\in\pi'_0(T_1),\ \mathrm{pr}_1\widetilde{\mathrm{bd}}\text{ contained in }\mathrm{pc}\right\}\cdot A'(\widetilde{\mathrm{bd}})\\
&=\left(\mathrm{pj}^{*}\left(H_0^{(1)}(s')A'\right)\right)(\widetilde{\mathrm{bd}}).
\end{align*}
By the definition of the projection and the injectivity of $\mathrm{pj}^{*}$, it follows that $H_0^{\mathrm{pj}}(s')=H_0^{(1)}(s')$, the diagonal family on $T_1$.

\emph{The family $H_0(s'',A)$.}
Let $f\in\mathfrak{S}^{\mathrm{ch}}(1,\ell_1)$, $s''\in\mathrm{Pd}_0(T_1)$, $A\in\tilde{G}_0(T_1)$ and $x=(\tilde{c},g)$.
Section~I.5.8 on $T_2$, (1.4) and the bijection of (cv7) give
\begin{align*}
\left(H_0^{(2)}(\mathrm{pj}^{*}s'',\mathrm{pj}^{*}A)\,\mathrm{pj}^{*}f\right)(\tilde{c},g)
&=\prod\left\{\left(\mathrm{pj}^{*}s''\right)_{\mathrm{pc}'}:\ \mathrm{pc}'\in\pi_0(T_2),\ \mathrm{pc}'\text{ meets }K_{\mathrm{pr}_2\tilde{c}}\right\}\cdot f(\tilde{c},g)\\
&=\prod\left\{s''_{\mathrm{pc}}:\ \mathrm{pc}\in\pi_0(T_1),\ \mathrm{pc}\text{ meets }K_{\mathrm{pr}_1\tilde{c}}\right\}\cdot f(\tilde{c},g)\\
&=\left(\mathrm{pj}^{*}\left(H_0^{(1)}(s'',A)f\right)\right)(\tilde{c},g),
\end{align*}
with the same factor at the $48$ sites $(\tilde{c},g)$, $g\in O_{\mathrm{cub}}$.
Hence $H_0^{\mathrm{pj}}(s'',A)=H_0^{(1)}(s'',A)$.
The factor is a number; therefore
\[
\bar{H}_0^{\mathrm{pj}}(s'',A)=\mathsf{C}^{-1}H_0^{(1)}(s'',-A)\,\mathsf{C}=H_0^{(1)}(s'',A).
\]

\emph{Their properties on $T_1$.}
By (1.0) the data $(0,\ell_1,\mu,e_k)$ are (SF).

To the family $H_0^{(1)}(s')$ applies \cite[Proposition~15.36]{QEDI}.
The setting of that proposition is the one of \cite[Section~15.10]{QEDI}:
$\ell\in\mathrm{Cl}$, $M$ a power of $L$ with $M\ge L^2$, $0<\alpha_{0}\le1/4$, and, from \cite[Section~15.10.3]{QEDI} on, $M\ge M^{P}$.
Here $\ell_1\in\mathrm{Cl}$ is a datum, and the three conditions on $M$ and $\alpha_{0}$ are among the conditions of Section~5.
By \cite[Proposition~15.36]{QEDI} the family has the properties W1, W2, W8, W9, W10, W14 for all $\tilde{g}\in\mathbb{G}_0$ (a group which contains $\mathbb{G}_0^{L}$), W16, and W3 with the constant $4\exp(2\gamma)\cdot M^{27\alpha_{0}}$ for every rate $\gamma>0$;
these are the properties with the same numbers in \cite[Section~5.4.2]{QEDI}.
The rate is taken to be $\gamma=\gamma_{\mathrm{w}}>0$.

The family $H_0^{(1)}(s'',A)$ has, by Proposition~8\,(b), the properties W1, W2, W8--W11, W13, W14 for all $\tilde{g}\in\mathbb{G}_0$, W16, and W5 with the constant $M^{8\alpha_{0}}\le4\exp(2\gamma)\cdot M^{27\alpha_{0}}$ for every $\gamma\ge0$;
again the rate is taken to be $\gamma=\gamma_{\mathrm{w}}$.
At level $0$ there is no short pair, so that the second part of W5 is void.

Finally, $4\exp(2\gamma_{\mathrm{w}})\cdot M^{27\alpha_{0}}\le C_{w}^{(1)}$ by the definition
\[
C_{w}^{(1)}=\max\left(C_{w}^{(0)},4\exp(2\gamma_{\mathrm{w}})\cdot M^{27\alpha_{0}}\right).
\]
For these two families the bound (8.0$'$) does not enter, and so no factor $64$ arises.
\end{proof}

\subsubsection*{3.7. Proof of (b)}

\begin{proof}[Proof of Proposition~26.3\,(b)]
\textbf{The projection is a system of primitives of level $k$ on $T_1$.}
By (1.9), the operators $\mathrm{Op}^{\mathrm{pj}}(s,A)$ and $\mathrm{Op}^{\mathrm{pj}}(X,A)$ are defined for all $s\in\mathrm{Pd}_k(T_1)$, $A\in\tilde{G}_k(T_1)$, $X\in D_k(T_1)$ and $y\in[0,\infty)$, and they are linear maps between the spaces on $T_1$ listed in Section~3.0.
The barred family is given by its definition, $\bar{H}^{\mathrm{pj}}_k(s'',A)=\mathsf{C}^{-1}H^{\mathrm{pj}}_k(s'',-A)\,\mathsf{C}$.

\emph{Measurability in $y$.}
Fix $X$, $A$ and two fermion sites $x$, $x'$.
By (8.0) and by the definition of $\widetilde{\mathrm{Op}}(X,A)$,
\[
S^{\mathrm{pj}}_{k,y}(X,A)(x,x')=\sum_{w\in\mathrm{Rw}}\boldsymbol{\chi}(w)\cdot\sum_{X'\in D_k(T_2):\ \mathrm{pj}(X')=X}S^{(2)}_{k,y}(X',A_2)(x,x'+\ell_1\cdot w).
\]
Both sums are finite, by (cv5), and the signs $\boldsymbol{\chi}(w)$ do not depend on $y$.
Each summand is a Borel function of $y\in[0,\infty)$, because measurability in $y$ belongs to the definition of a system on $T_2$ in Section~G4.0.4, for $k\ge1$ as well as for $k=0$.
The space $\mathfrak{S}^{\mathrm{ch}}(\eta_k,\ell_1)$ is finite-dimensional, so that by (kn1) an operator on it is determined by finitely many entries of its kernel.
Hence $y\mapsto S^{\mathrm{pj}}_{k,y}(X,A)$ is a Borel function on $[0,\infty)$.

\emph{Level $0$.}
By (c), the projections of the families (p1), (p3) are the diagonal families on $T_1$, and there is no family $M_0$.
Therefore $\mathrm{Op}^{\mathrm{pj}}$ is a system of level $0$ in the sense of Section~G4.0.4.

In the arguments below for W1--W16, which are properties from the list W1--W17 in Section~G4.0.4, the level $k\ge0$ is arbitrary.
At $k=0$ the family $M_k$ does not occur, and the families (p1), (p3) are those of (c); they are dealt with under W17 below.

\textbf{W1.}
By (1.4), $\mathrm{pj}^*$ maps the parameter family on $T_1$ with all entries equal to $1$ to the parameter family on $T_2$ with all entries equal to $1$; both are written $1$.

\emph{(p1), $k\ge1$.}
Let $A'$ be a bond function on $T^0_1$.
The definition of the projection, W1 on $T_2$ and \cite[Proposition~15.28]{QEDI} for $O=H_k$ give
\[
\mathrm{pj}^*H^{\mathrm{pj}}_k(1)A'=H^{(2)}_k(1)\,\mathrm{pj}^*A'=H^{(2)}_k\,\mathrm{pj}^*A'=\mathrm{pj}^*H^{(1)}_kA'.
\]
The hypotheses of \cite[Proposition~15.28]{QEDI} are satisfied here: $k\ge0$, $\ell_1,\ell_2\in\mathrm{Cl}$ and $\ell_2=\mathrm{Jr}\cdot\ell_1$ with $\mathrm{Jr}$ odd.
In that proposition, as in \cite[Proposition~15.5]{QEDI}, the operator $H_k$ acts on the real bond functions.
On the complex bond functions of Section~I.1.1 both sides of the identity $H^{(2)}_k\,\mathrm{pj}^*=\mathrm{pj}^*H^{(1)}_k$ are understood as the complex linear continuations, and the identity persists, because $\mathrm{pj}^*$ is complex linear and maps real bond functions to real bond functions.

\emph{(p2), $k\ge0$.}
For $\mathsf{Z}\in V_k(T_1)$, W1 on $T_2$ and \cite[Proposition~15.28]{QEDI} for $O=C_{k}^{1/2}$ give
\[
\mathrm{pj}^*C_{k}^{1/2,\mathrm{pj}}(1)\mathsf{Z}=C_{k}^{1/2,(2)}(1)\,\mathrm{pj}^*\mathsf{Z}=C_{k}^{1/2,(2)}\,\mathrm{pj}^*\mathsf{Z}=\mathrm{pj}^*C_{k}^{1/2,(1)}\mathsf{Z};
\]
on the complexification both sides are again the complex linear continuations.

\emph{(p3) for $k\ge1$; $\Gamma_k$ for $k\ge0$.}
Let $A\in\tilde{G}_k(T_1)$, and let $v$ lie in the source space on $T_1$.
From W1 on $T_2$ and Proposition~11\,(b), applied to $O=H_k(A)$ and to $O=\Gamma_k(A)$ if $k\ge1$, and to $O=\Gamma_0(A)$ if $k=0$, one obtains
\begin{align*}
\mathrm{pj}^*H^{\mathrm{pj}}_k(1,A)v&=H^{(2)}_k(1,A_2)\,\mathrm{pj}^*v=H^{(2)}_k(A_2)\,\mathrm{pj}^*v=\mathrm{pj}^*H^{(1)}_k(A)v,\\
\mathrm{pj}^*\Gamma^{\mathrm{pj}}_k(1,A)v&=\Gamma^{(2)}_k(1,A_2)\,\mathrm{pj}^*v=\Gamma^{(2)}_k(A_2)\,\mathrm{pj}^*v=\mathrm{pj}^*\Gamma^{(1)}_k(A)v.
\end{align*}

\emph{$M_k$ for $k\ge1$; $S_{k,y}$ for $k\ge0$ and $y\ge0$.}
Let $f$ lie in the source space on $T_1$.
By (cv5), W1 on $T_2$ and Proposition~11\,(b), applied to $O=M_k(A)$, resp. $O=S_{k,y}(A)$,
\begin{align*}
\mathrm{pj}^*\sum_{X\in D_k(T_1)}\mathrm{Op}^{\mathrm{pj}}(X,A)f&=\sum_{X\in D_k(T_1)}\ \sum_{X':\ \mathrm{pj}(X')=X}\mathrm{Op}^{(2)}(X',A_2)\,\mathrm{pj}^*f=\sum_{X'\in D_k(T_2)}\mathrm{Op}^{(2)}(X',A_2)\,\mathrm{pj}^*f\\
&=O^{(2)}(A_2)\,\mathrm{pj}^*f=\mathrm{pj}^*O^{(1)}(A)f.
\end{align*}
Since $\mathrm{pj}^*$ is injective, it follows that
\begin{gather*}
H^{\mathrm{pj}}_k(1)=H_k,\qquad C_{k}^{1/2,\mathrm{pj}}(1)=C_{k}^{1/2},\qquad H^{\mathrm{pj}}_k(1,A)=H_k(A),\qquad\Gamma^{\mathrm{pj}}_k(1,A)=\Gamma_k(A),\\
\sum_{X}M^{\mathrm{pj}}_k(X,A)=M_k(A),\qquad\sum_{X}S^{\mathrm{pj}}_{k,y}(X,A)=S_{k,y}(A),
\end{gather*}
where the right sides are the exact operators on $T_1$; the fermionic ones are those of Section~I.3.5 for $k\ge1$ and those of Section~I.3.9 for $k=0$.

\emph{Level $0$.}
On $\mathfrak{S}^{\mathrm{ch}}(1,\ell_i)$ one has $\Gamma_0(A)=F_0(A)^{-1}$ and $S_{0,y}(A)=F_y(A)^{-1}$ by Section~I.3.9, and both operators belong to the list of Proposition~11\,(b) for $k=0$.

\textbf{W2.}
The map $(s,A)\mapsto(s_2,A_2)$ is complex linear, and it maps $\mathrm{Pd}_k(T_1)\times\tilde{G}_k(T_1)$ into $\mathrm{Pd}_k(T_2)\times\tilde{G}_k(T_2)$.
By \cite[Lemma~17.3]{QEDI}, $\tilde{G}_k(T_i)$ is an open subset of the space of the complex bond functions on $T_i$.
Consequently the interior of $\mathrm{Pd}_k(T_i)\times\tilde{G}_k(T_i)$ is the product of the interior of $\mathrm{Pd}_k(T_i)$ with $\tilde{G}_k(T_i)$, and by (1.4) the interior for $i=1$ is mapped into the interior for $i=2$.
According to (8.0), every entry of the kernel of $\mathrm{Op}^{\mathrm{pj}}(s,A)$ is the sum over $w\in\mathrm{Rw}$ of the entries $\mathrm{Op}^{(2)}(s_2,A_2)(t,t'+\ell_1\cdot w)$, multiplied by the signs $\boldsymbol{\chi}(w)$, which are independent of $(s,A)$.
Likewise, every entry of the kernel of $\mathrm{Op}^{\mathrm{pj}}(X,A)$ is the sum, over $w\in\mathrm{Rw}$ and over the finitely many $X'$ with $\mathrm{pj}(X')=X$, of the entries $\mathrm{Op}^{(2)}(X',A_2)(t,t'+\ell_1\cdot w)$, multiplied by the signs $\boldsymbol{\chi}(w)$.
By W2 on $T_2$, for (p1)--(p3) and $\Gamma_k$ each summand is the composition of the linear map $(s,A)\mapsto(s_2,A_2)$ with a function that is continuous, resp. analytic in the interior; therefore the entries of $\mathrm{Op}^{\mathrm{pj}}(s,A)$ are continuous on $\mathrm{Pd}_k(T_1)\times\tilde{G}_k(T_1)$ and analytic in its interior.
Let $s\in\mathrm{Pd}_k(T_1)$ be fixed, boundary points being admitted.
Then $s_2$ is a fixed point of $\mathrm{Pd}_k(T_2)$, the map $A\mapsto A_2$ is linear from $\tilde{G}_k(T_1)$ into $\tilde{G}_k(T_2)$, and $A_2\mapsto\mathrm{Op}^{(2)}(s_2,A_2)$ is analytic on $\tilde{G}_k(T_2)$; so $A\mapsto\mathrm{Op}^{\mathrm{pj}}(s,A)$ is analytic on $\tilde{G}_k(T_1)$.
Finally, the entries of $M^{\mathrm{pj}}_k(X,\,\cdot\,)$ and of $S^{\mathrm{pj}}_{k,y}(X,\,\cdot\,)$ are finite sums of the functions $A\mapsto\mathrm{Op}^{(2)}(X',A_2)(t,t'+\ell_1\cdot w)$, and each of these is analytic on $\tilde{G}_k(T_1)$.

\textbf{W3} ((p1), $k\ge1$).
The bonds $b'$, $b$, $\tilde{b}$ and the kernels $\partial_{\mathfrak{l}}O$, $\delta_{\alpha}\partial_{\mathfrak{l}}O$ are those of property W3 in the list of Section~G4.0.4.
Let $b=(c,c+\eta_k\cdot e_{\mu})$ and $\tilde{b}=(\tilde{c},\tilde{c}+\eta_k\cdot e_{\mu})$ be fine bonds of $T_1$ as in W3; here $c$ and $\tilde{c}$ are sites of $T_1$.
By (0.c) the two bonds have lifts to $\mathbb{L}_k$ whose initial points are at the sup distance $d_1(c,\tilde{c})\le1$ from each other; in the rest of this paragraph $b$ and $\tilde{b}$ stand for these lifts.
Since $1<\ell_2/2$, (ds1) shows that $d_1(c,\tilde{c})$ is also the distance of the images of the two initial points in $T_2$.
The two difference operations act on the first argument of the kernel.
Hence (8.0), in which $\boldsymbol{\chi}=1$, yields for each of the three kernels $\mathrm{Kn}\in\{O,\ \partial_{\mathfrak{l}}O,\ \delta_{\alpha}\partial_{\mathfrak{l}}O\}$
\[
\mathrm{Kn}\left[H^{\mathrm{pj}}_k(s')\right](\,\cdot\,;b')=\sum_{w\in\mathrm{Rw}}\mathrm{Kn}\left[\widetilde{\mathrm{Op}}(s')\right](\,\cdot\,;b'+\ell_1\cdot w),\qquad\widetilde{\mathrm{Op}}(s')=H^{(2)}_k(s'_2),
\]
the first argument being $b$, resp. $(b,\tilde{b})$, on both sides.
Here and in the estimates below,
\[
C_{w}^{(1)}=\max\left(C_{w}^{(0)},\,4\cdot\exp(2\gamma_{\mathrm{w}})\cdot M^{27\alpha_{0}}\right)
\]
is the constant of the line (C6) of Section~I.8.2 and the fourth entry of $\mathrm{par}^0$; it is a single number, the same for every $w\in\mathrm{Rw}$, and so it can be taken out of the sums over $w\in\mathrm{Rw}$.
In $C_{w}$, $C_{w}^{(0)}$ and $C_{w}^{(1)}$ the subscript is part of the name, as is the upright subscript of $\gamma_{\mathrm{w}}$ and $\kappa_{\mathrm{w}}$; it does not denote a vector $w\in\mathbb{Z}^3$, unlike the index of $\mathrm{Yl}_w$ in W9 below.
Property W3 of the system on $T_2$, with the parameters $\mathrm{par}^0$, together with (8.0$'$), gives
\[
\left|\mathrm{Kn}\left[H^{\mathrm{pj}}_k(s')\right](\,\cdot\,;b')\right|\le C_{w}^{(1)}\cdot\sum_{w\in\mathrm{Rw}}\exp\left(-\gamma_{\mathrm{w}}\cdot d_2(b,b'+\ell_1\cdot w)\right)\le64\cdot C_{w}^{(1)}\cdot\exp\left(-\gamma_{\mathrm{w}}\cdot d_1(b,b')\right).
\]
Finally $64\cdot C_{w}^{(1)}\le64\cdot K_{0}\cdot C_{w}^{(1)}=C_{w}$, because $K_{0}\ge1$ and $C_{w}^{(1)}>0$.

\textbf{W4} ((p2), $k\ge0$).
By (8.0) with $\boldsymbol{\chi}=1$, W4 on $T_2$ and (8.0$'$),
\begin{align*}
\left|C_{k}^{1/2,\mathrm{pj}}(s')(\Upsilon,\Upsilon')\right|&\le\sum_{w\in\mathrm{Rw}}\left|\widetilde{\mathrm{Op}}(s')(\Upsilon,\Upsilon'+\ell_1\cdot w)\right|\le C_{w}^{(1)}\cdot\sum_{w\in\mathrm{Rw}}\exp\left(-\gamma_{\mathrm{w}}\cdot d_2(\Upsilon,\Upsilon'+\ell_1\cdot w)\right)\\
&\le64\cdot C_{w}^{(1)}\cdot\exp\left(-\gamma_{\mathrm{w}}\cdot d_1(\Upsilon,\Upsilon')\right)\le C_{w}\cdot\exp\left(-\gamma_{\mathrm{w}}\cdot d_1(\Upsilon,\Upsilon')\right).
\end{align*}

\textbf{W5} ((p3), $k\ge1$).
Let $v\in\mathfrak{S}(1,\ell_1)$.
By (1.2), the definition of the projection and W5 on $T_2$,
\[
\left|H^{\mathrm{pj}}_k(s'',A)v\right|_{\infty}=\left|\mathrm{pj}^*H^{\mathrm{pj}}_k(s'',A)v\right|_{\infty}=\left|H^{(2)}_k(s''_2,A_2)\,\mathrm{pj}^*v\right|_{\infty}\le C_{w}^{(1)}\cdot\left|\mathrm{pj}^*v\right|_{\infty}=C_{w}^{(1)}\cdot|v|_{\infty}.
\]
Put $f:=H^{\mathrm{pj}}_k(s'',A)v\in\mathfrak{S}^{\mathrm{ch}}(\eta_k,\ell_1)$, and let $(x,x')\in\mathrm{Pzt}_k$, with $x=(\tilde{c},g)$ and $x'=(\tilde{c}',g')$, be a short pair, that is, $d_X(x,x')<1$.
By (I.1.7) and (I.1.6),
\[
(\delta_{\alpha,A}f)(x,x')=d_X(x,x')^{-\alpha}\cdot\left(\exp\left(i\cdot e_k\cdot\eta_k\cdot(\tau A)(\tilde{c},\tilde{c}')\right)\cdot f(x')-f(x)\right),
\]
where $f(x)$, $f(x')$ are the values of $f$ at the two points of $\mathbb{L}^{\mathrm{f}}_k$, and $(\tau A)(\tilde{c},\tilde{c}')$ is the mean of the sums of $A$ along the six paths $\boldsymbol{\Gamma}^{\pi}(\tilde{c},\tilde{c}')$ of $\mathbb{L}_k$.
On $\mathbb{L}^{\mathrm{f}}_k$ the functions $f$ and $\mathrm{pj}^*f$ are the same, and on the bonds of $\mathbb{L}_k$ the functions $A$ and $A_2$ are the same.
Therefore
\[
(\delta_{\alpha,A}f)(x,x')=\left(\delta_{\alpha,A_2}\mathrm{pj}^*f\right)(x,x')=\left(\delta_{\alpha,A_2}H^{(2)}_k(s''_2,A_2)\,\mathrm{pj}^*v\right)(x,x').
\]
For a pair inside a cluster, $\tilde{c}=\tilde{c}'$ and $g\ne g'$, both sides equal $\eta_k^{-\alpha}\cdot\left(f(\tilde{c},g')-f(\tilde{c},g)\right)$.
The short pairs of $T_1$ and those of $T_2$ are the classes of the same pairs of $\mathrm{Pzt}_k$.
By W5 on $T_2$ the norm $|\cdot|_c$ of the right side is at most $C_{w}^{(1)}\cdot\left|\mathrm{pj}^*v\right|_{\infty}=C_{w}^{(1)}\cdot|v|_{\infty}$.
Here $C_{w}^{(1)}\le64\cdot K_{0}\cdot C_{w}^{(1)}=C_{w}$, as $K_{0}\ge1$ and $C_{w}^{(1)}>0$.
\emph{Level $0$.} The assertion is contained in (c).

\textbf{W6 for $\Gamma_k$} ($k\ge0$).
By Section~I.1.5, $|\cdot|_c$ is a norm on the $4\times4$ matrices.
From (8.0), W6 on $T_2$ and (8.0$'$),
\begin{align*}
\left|\Gamma^{\mathrm{pj}}_k(s'',A)(t,t')\right|_c&\le\sum_{w\in\mathrm{Rw}}\left|\Gamma^{(2)}_k(s''_2,A_2)(t,t'+\ell_1\cdot w)\right|_c\le C_{w}^{(1)}\cdot\sum_{w\in\mathrm{Rw}}\exp\left(-\gamma_{\mathrm{w}}\cdot d_2(t,t'+\ell_1\cdot w)\right)\\
&\le64\cdot C_{w}^{(1)}\cdot\exp\left(-\gamma_{\mathrm{w}}\cdot d_1(t,t')\right)\le C_{w}\cdot\exp\left(-\gamma_{\mathrm{w}}\cdot d_1(t,t')\right).
\end{align*}
For $k\ge1$ the indices $t$, $t'$ are unit sites.
\emph{Level $0$.} Here $t=(\tilde{c},g)$ and $t'=(\tilde{c}',g')$ are fermion sites, both kernels are kernels in the sense of (I.1.4) with the weight $\frac{1}{48}$, and $d_i(t,t')$ is the distance of the photon sites, by (I.POS).

\textbf{W6 for $M_k$} ($k\ge1$).
Let $z$, $z'$ be unit sites.
By W12 on $T_2$, $M^{(2)}_k(X',A_2)(z,z'')=0$ unless $\mathrm{pr}_2z\in X'$, that is, unless $X'$ contains the $M$-cube $\square_z$ of Section~3.5.
Combining (8.0), W6 on $T_2$ with the parameters $\mathrm{par}^0$, (3.a) and (8.0$'$), one obtains
\begin{align*}
\left|M^{\mathrm{pj}}_k(X,A)(z,z')\right|_c&\le C_{w}^{(1)}\cdot\sum_{w\in\mathrm{Rw}}\ \sum_{\substack{X':\ \mathrm{pj}(X')=X,\\ \square_z\subseteq X'}}\exp\left(-\gamma_{\mathrm{w}}\cdot d_2(z,z'+\ell_1\cdot w)-(\kappa_{\mathrm{w}}+\kappa_{0})\cdot d_M(X')\right)\\
&\le C_{w}^{(1)}\cdot64\cdot\exp\left(-\gamma_{\mathrm{w}}\cdot d_1(z,z')\right)\cdot K_{0}\cdot\exp\left(-\kappa_{\mathrm{w}}\cdot d_M(X)\right)\\
&=C_{w}\cdot\exp\left(-\gamma_{\mathrm{w}}\cdot d_1(z,z')-\kappa_{\mathrm{w}}\cdot d_M(X)\right).
\end{align*}

\textbf{W7, $k\ge1$.}
Let $z$, $z'$ be unit sites of the covering lattice, $f\in\mathfrak{S}^{\mathrm{ch}}(\eta_k,\ell_1)$ and $y\ge0$.
On $T_i$ the indicator $\mathbf{1}_{\Delta_{z'}}$ is, by (I.POS), the indicator of the $\ell_i\mathbb{Z}^3$-orbit of $\Delta_{z'}\times O_{\mathrm{cub}}$; it is denoted $\mathbf{1}^{(i)}_{\Delta_{z'}}$.
By (fs) for the unit sites, the $\ell_1\mathbb{Z}^3$-orbit of $\Delta_{z'}$ is the disjoint union, over $w\in\mathrm{Rw}$, of the $\ell_2\mathbb{Z}^3$-orbits of $\Delta_{z'+\ell_1\cdot w}$.
Therefore $\mathrm{pj}^*\left(\mathbf{1}^{(1)}_{\Delta_{z'}}f\right)=\sum_{w\in\mathrm{Rw}}\mathbf{1}^{(2)}_{\Delta_{z'+\ell_1\cdot w}}\mathrm{pj}^*f$, and
\[
\mathrm{pj}^*\left(S^{\mathrm{pj}}_{k,y}(X,A)\,\mathbf{1}^{(1)}_{\Delta_{z'}}f\right)=\sum_{X':\ \mathrm{pj}(X')=X}\ \sum_{w\in\mathrm{Rw}}S^{(2)}_{k,y}(X',A_2)\,\mathbf{1}^{(2)}_{\Delta_{z'+\ell_1\cdot w}}\mathrm{pj}^*f.
\]
For every $f_1\in\mathfrak{S}^{\mathrm{ch}}(\eta_k,\ell_1)$ the number $\left|\mathbf{1}^{(1)}_{\Delta_z}f_1\right|_{\infty}$ is the supremum of $|f_1(x)|_c$ over the $48\cdot\eta_k^{-3}$ sites $x$ over $\Delta_z$, and it equals $\left|\mathbf{1}^{(2)}_{\Delta_z}\mathrm{pj}^*f_1\right|_{\infty}$.
By (g0) of \cite[Section~17.4.1]{QEDI}, $\Delta_z$ lies in a single $M$-cube, the $M$-cube of $z$.
By W12 on $T_2$, the rows of $S^{(2)}_{k,y}(X',A_2)$ at the sites over $\Delta_z$ vanish unless $\square_z\subseteq X'$.
Hence W7 on $T_2$ with the parameters $\mathrm{par}^0$, (3.a), (8.0$'$) and (1.2) give
\begin{align*}
\left|\mathbf{1}^{(1)}_{\Delta_z}S^{\mathrm{pj}}_{k,y}(X,A)\,\mathbf{1}^{(1)}_{\Delta_{z'}}f\right|_{\infty}&\le C_{w}^{(1)}\cdot\sum_{\substack{X':\ \mathrm{pj}(X')=X,\\ \square_z\subseteq X'}}\ \sum_{w\in\mathrm{Rw}}\exp\left(-\gamma_{\mathrm{w}}\cdot d_2(z,z'+\ell_1\cdot w)-(\kappa_{\mathrm{w}}+\kappa_{0})\cdot d_M(X')\right)\cdot\left|\mathrm{pj}^*f\right|_{\infty}\\
&\le64\cdot K_{0}\cdot C_{w}^{(1)}\cdot\exp\left(-\gamma_{\mathrm{w}}\cdot d_1(z,z')-\kappa_{\mathrm{w}}\cdot d_M(X)\right)\cdot|f|_{\infty}.
\end{align*}

\textbf{W7, $k=0$.}
Let $x$, $x'$ be fermion sites of $\mathbb{L}^{\mathrm{f}}_0$ and $y\ge0$.
By W12 on $T_2$, $S^{(2)}_{0,y}(X',A_2)(x,x'')=0$ unless $\square_x\subseteq X'$.
By (8.0), the kernels being those of (I.1.4) with the weight $\frac{1}{48}$ on both tori, by W7 on $T_2$ in its form for $k=0$, by (3.a) and by (8.0$'$),
\begin{align*}
\left|S^{\mathrm{pj}}_{0,y}(X,A)(x,x')\right|_c&\le C_{w}^{(1)}\cdot\sum_{w\in\mathrm{Rw}}\ \sum_{\substack{X':\ \mathrm{pj}(X')=X,\\ \square_x\subseteq X'}}(1+y)^{-1}\cdot\exp\left(-\gamma_{\mathrm{w}}\cdot d_2(x,x'+\ell_1\cdot w)-(\kappa_{\mathrm{w}}+\kappa_{0})\cdot d_M(X')\right)\\
&\le C_{w}\cdot(1+y)^{-1}\cdot\exp\left(-\gamma_{\mathrm{w}}\cdot d_1(x,x')-\kappa_{\mathrm{w}}\cdot d_M(X)\right),
\end{align*}
the factor $(1+y)^{-1}$ being common to all terms.

\textbf{Setting for W8--W11.}
Let $\mathrm{Op}$ be one of the families to which the property refers, $\mathscr{Y}$ a set of parameter cubes of $T_1$ of the kind of $\mathrm{Op}$, $Y$ its union, and $s\in\mathrm{Pd}_k(T_1)$ with $s_{\mathrm{pc}}=0$ for every $\mathrm{pc}\notin\mathscr{Y}$.
Let $\mathscr{Y}'$ and $Y'=\mathrm{pj}^{-1}(Y)$ be as in (cv3).
As $(s_2)_{\mathrm{pc}'}=s_{\mathrm{pj}(\mathrm{pc}')}$ by (1.4), the family $s_2$ vanishes outside $\mathscr{Y}'$.
\emph{Level $0$.} W8 and W9 concern $C_{0}^{1/2}(s')$, of the kind $\pi''_0$, whose rows and columns are basis elements, and $\Gamma_0(s'',A)$, of the kind $\pi_0$, whose rows and columns are fermion sites $x=(\tilde{c},g)$ of $\mathbb{L}^{\mathrm{f}}_0$ with $\operatorname{supp}x=\{\tilde{c}\}$; W10 concerns $C_{0}^{1/2}(s')$; the diagonal families are treated in (c).
\begin{itemize}
\item[(lk1)] Let $\mathrm{Ysl}$ be a link-component of $\mathscr{Y}'$ and $\mathrm{Yl}$ its union.
By (cv2), $\{\mathrm{pj}(\mathrm{pc}'):\mathrm{pc}'\in\mathrm{Ysl}\}$ is a link-connected subset of $\mathscr{Y}$, and it is therefore contained in one link-component of $\mathscr{Y}$.
\item[(lk2)] Let $\mathscr{Y}_1$ be a link-component of $\mathscr{Y}$, $Y_1$ its union, and $t$ an index belonging to the kind that is contained in $Y_1$ on $T_1$.
By (cv4), $t$ is contained in $Y'$ on $T_2$.
By Proposition~26.2\,(b) on $T_2$, for the side $\ell_2\in\mathrm{Cl}$ and the set $\mathscr{Y}'$, the index $t$ is contained in exactly one link-component $\mathrm{Yl}$ of $Y'$.
Let $\mathrm{Ysl}$ be the link-component of $\mathscr{Y}'$ whose union is $\mathrm{Yl}$; it is unique, because different link-components have no lattice site in common, by (g3) of \cite[Section~17.4.1]{QEDI}, and their unions are not empty.
Then
\[
\mathrm{pj}(\mathrm{pc}')\in\mathscr{Y}_1\ \text{ for every }\mathrm{pc}'\in\mathrm{Ysl},\qquad\mathrm{pj}(\mathrm{Yl})\subseteq Y_1.
\]
To see this, note that $\operatorname{supp}t$ is not empty and that $t$ is contained in $\mathrm{Yl}$ on $T_2$; so there are a cube $\mathrm{pc}'_1\in\mathrm{Ysl}$ and a point $\tilde{c}\in\operatorname{supp}t$ with $\mathrm{pr}_2\tilde{c}\in\mathrm{pc}'_1$.
The cube $\mathrm{pj}(\mathrm{pc}'_1)$ belongs to $\mathscr{Y}$ and contains $\mathrm{pr}_1\tilde{c}$, a point of the support of $t$ on $T_1$.
Since $t$ is contained in $Y_1$ on $T_1$, some cube of $\mathscr{Y}_1$ also contains a point of this support.
The support of $t$ on $T_1$ is contained in $Y$; so (b1) in Section~2, applied on $T_1$, shows that the two cubes belong to the same link-component of $\mathscr{Y}$, whence $\mathrm{pj}(\mathrm{pc}'_1)\in\mathscr{Y}_1$.
By (lk1), all cubes of $\mathrm{Ysl}$ are mapped into the link-component of $\mathrm{pj}(\mathrm{pc}'_1)$, which is $\mathscr{Y}_1$.
By (ds4), $\mathrm{pj}(\mathrm{Yl})$ is the union of the images of these cubes; hence $\mathrm{pj}(\mathrm{Yl})\subseteq Y_1$.
\item[(lk3)] In the situation of (lk2), let $\tilde{s}\in\mathrm{Pd}_k(T_1)$ vanish outside $\mathscr{Y}$ and agree with $s$ on the cubes of $\mathscr{Y}_1$, and let $\tilde{A}\in\tilde{G}_k(T_1)$ agree with $A$ on the bonds that meet $Y_1$.
Then $\tilde{s}_2:=\mathrm{pj}^*\tilde{s}$ lies in $\mathrm{Pd}_k(T_2)$, vanishes outside $\mathscr{Y}'$, and agrees with $s_2$ on the cubes of $\mathrm{Ysl}$: for $\mathrm{pc}'\in\mathrm{Ysl}$ one has, by (lk2),
\[
(\tilde{s}_2)_{\mathrm{pc}'}=\tilde{s}_{\mathrm{pj}(\mathrm{pc}')}=s_{\mathrm{pj}(\mathrm{pc}')}=(s_2)_{\mathrm{pc}'}.
\]
Further, $\tilde{A}_2:=\mathrm{pj}^*\tilde{A}$ lies in $\tilde{G}_k(T_2)$ and agrees with $A_2$ on the bonds of $T_2$ that meet $\mathrm{Yl}$: such a bond $b'$ is mapped to a bond $\mathrm{pj}(b')$ that meets $\mathrm{pj}(\mathrm{Yl})\subseteq Y_1$, and by (cv4) one has $\tilde{A}_2(b')=\tilde{A}(\mathrm{pj}(b'))=A(\mathrm{pj}(b'))=A_2(b')$.
\end{itemize}

\textbf{W8} ((p1)--(p3), $\Gamma_k$).
Let the row index $t$ be contained in $Y_1$, and let $\tilde{s}$, $\tilde{A}$ be as in (lk3).
Property W8 on $T_2$, for the set $\mathscr{Y}'$, the family $s_2$, the link-component $\mathrm{Yl}$, the row index $t$ and the pair $(\tilde{s}_2,\tilde{A}_2)$, gives
\[
\widetilde{\mathrm{Op}}(s,A)(t,t'')=0\ \text{ unless $t''$ is contained in $\mathrm{Yl}$ on $T_2$},\qquad\widetilde{\mathrm{Op}}(s,A)(t,t'')=\widetilde{\mathrm{Op}}(\tilde{s},\tilde{A})(t,t'')\ \text{ for all }t''.
\]
By (8.0), for all $t'$,
\[
\mathrm{Op}^{\mathrm{pj}}(s,A)(t,t')=\sum_{w\in\mathrm{Rw}}\boldsymbol{\chi}(w)\cdot\widetilde{\mathrm{Op}}(s,A)(t,t'+\ell_1\cdot w)=\sum_{w\in\mathrm{Rw}}\boldsymbol{\chi}(w)\cdot\widetilde{\mathrm{Op}}(\tilde{s},\tilde{A})(t,t'+\ell_1\cdot w)=\mathrm{Op}^{\mathrm{pj}}(\tilde{s},\tilde{A})(t,t').
\]
If $\mathrm{Op}^{\mathrm{pj}}(s,A)(t,t')\ne0$, then $\widetilde{\mathrm{Op}}(s,A)(t,t'+\ell_1\cdot w)\ne0$ for some $w$; so $t'+\ell_1\cdot w$ is contained in $\mathrm{Yl}$ on $T_2$, and by (cv4) and (lk2) it is contained in $\mathrm{pj}(\mathrm{Yl})$, hence in $Y_1$, on $T_1$.
As an index of $T_1$, $t'+\ell_1\cdot w$ is the same as $t'$.

\textbf{W9} ((p1)--(p3), $\Gamma_k$).
Let the column index $t'$ be contained in $Y_1$, and let $\tilde{s}$, $\tilde{A}$ be as in (lk3).
For every $w\in\mathrm{Rw}$ the index $t'+\ell_1\cdot w$ is a lift of the same index of $T_1$ as $t'$; therefore it is contained in $Y_1$ on $T_1$.
Applied to $t'+\ell_1\cdot w$, (lk2) and (lk3) furnish a link-component $\mathrm{Yl}_w$ of $Y'$ that contains $t'+\ell_1\cdot w$ on $T_2$ and satisfies $\mathrm{pj}(\mathrm{Yl}_w)\subseteq Y_1$, such that $\tilde{s}_2$ vanishes outside $\mathscr{Y}'$ and agrees with $s_2$ on the cubes of the link-component of $\mathscr{Y}'$ with the union $\mathrm{Yl}_w$, and $\tilde{A}_2$ agrees with $A_2$ on the bonds that meet $\mathrm{Yl}_w$.
Property W9 on $T_2$, for the column index $t'+\ell_1\cdot w$, gives
\begin{gather*}
\widetilde{\mathrm{Op}}(s,A)(t,t'+\ell_1\cdot w)=0\ \text{ unless $t$ is contained in $\mathrm{Yl}_w$ on $T_2$},\\
\widetilde{\mathrm{Op}}(s,A)(t,t'+\ell_1\cdot w)=\widetilde{\mathrm{Op}}(\tilde{s},\tilde{A})(t,t'+\ell_1\cdot w)\ \text{ for all }t.
\end{gather*}
Multiplication by $\boldsymbol{\chi}(w)$ and summation over $w\in\mathrm{Rw}$ give, by (8.0), $\mathrm{Op}^{\mathrm{pj}}(s,A)(t,t')=\mathrm{Op}^{\mathrm{pj}}(\tilde{s},\tilde{A})(t,t')$ for all $t$.
If $\mathrm{Op}^{\mathrm{pj}}(s,A)(t,t')\ne0$, then $\widetilde{\mathrm{Op}}(s,A)(t,t'+\ell_1\cdot w)\ne0$ for some $w$; hence $t$ is contained in $\mathrm{Yl}_w$ on $T_2$, and by (cv4) it is contained in $\mathrm{pj}(\mathrm{Yl}_w)$, hence in $Y_1$, on $T_1$.

\textbf{W10} ((p1)--(p3)).
Let the row index $t$ be contained in no link-component of $Y$ on $T_1$.
By the final assertion of Proposition~26.2, applied on $T_1$, the support of $t$ on $T_1$ is not contained in $Y$.
By (cv4), the support of $t$ on $T_2$ is not contained in $Y'$, and the same assertion of Proposition~26.2, applied on $T_2$, shows that $t$ is contained in no link-component of $Y'$.
Property W10 on $T_2$, for the set $\mathscr{Y}'$ and the family $s_2$, gives $\widetilde{\mathrm{Op}}(s,A)(t,t'')=0$ for all $t''$, and then (8.0) gives $\mathrm{Op}^{\mathrm{pj}}(s,A)(t,t')=0$ for all $t'$.

\textbf{W11} ((p3)).
Let $x=(\tilde{c},g)$ be a fermion site, and let $\mathrm{pc}$ be a cube of $\pi_k(T_1)$ that meets $K_{\mathrm{pr}_1\tilde{c}}$ and does not belong to $\mathscr{Y}$.
By (cv7), $\mathrm{pc}=\mathrm{pj}(\mathrm{pc}')$ for a cube $\mathrm{pc}'$ of $\pi_k(T_2)$ that meets $K_{\mathrm{pr}_2\tilde{c}}$, and $\mathrm{pc}'\notin\mathscr{Y}'$.
By W11 on $T_2$, for the set $\mathscr{Y}'$ and the family $s''_2$, the row of $H^{(2)}_k(s''_2,A_2)$ at $x$ vanishes; by (8.0), so does the row of $H^{\mathrm{pj}}_k(s'',A)$ at $x$.
(By (I.POS), $K_x=K_c$ for each of the $48$ sites $x$ over a photon site $c$.)

\textbf{W12} ($M_k$, $k\ge1$; $S_{k,y}$, $k\ge0$).
Suppose that $\mathrm{Op}^{\mathrm{pj}}(X,A)(t,t')\ne0$.
By (8.0) and the definition of $\widetilde{\mathrm{Op}}(X,A)$ there are $w\in\mathrm{Rw}$ and $X'$ with $\mathrm{pj}(X')=X$ and $\mathrm{Op}^{(2)}(X',A_2)(t,t'+\ell_1\cdot w)\ne0$.
By W12 on $T_2$, $t$ and $t'+\ell_1\cdot w$ lie over $X'$ on $T_2$; by (cv4) they lie over $\mathrm{pj}(X')=X$ on $T_1$, so that $t$ and $t'$ lie over $X$.
Next let $A,\tilde{A}\in\tilde{G}_k(T_1)$ agree on the bonds that meet $X$.
If $\mathrm{pj}(X')=X$, then a bond $b'$ of $T_2$ that meets $X'$ is mapped to a bond $\mathrm{pj}(b')$ that meets $X$, by (cv4); hence $A_2(b')=A(\mathrm{pj}(b'))=\tilde{A}(\mathrm{pj}(b'))=\tilde{A}_2(b')$.
By W12 on $T_2$, $\mathrm{Op}^{(2)}(X',A_2)=\mathrm{Op}^{(2)}(X',\tilde{A}_2)$ for each such $X'$.
Summation gives $\widetilde{\mathrm{Op}}(X,A)=\widetilde{\mathrm{Op}}(X,\tilde{A})$, and consequently $\mathrm{Op}^{\mathrm{pj}}(X,A)=\mathrm{Op}^{\mathrm{pj}}(X,\tilde{A})$.

\textbf{W13} ((p3)--(p5)).
Let $\boldsymbol{\lambda}$ be a real function of $c$ on $T_1$, that is, a real $\ell_1\mathbb{Z}^3$-periodic function on $\mathbb{L}_k$, and let $\boldsymbol{\lambda}^{u}$ be its restriction to $\mathbb{Z}^3$.
By \cite[Lemma~17.6]{QEDI}, $\tilde{G}_k$ is invariant under $A\mapsto A-d\boldsymbol{\lambda}$ for real $\boldsymbol{\lambda}$; thus $A-d\boldsymbol{\lambda}\in\tilde{G}_k(T_1)$.
The function $\mathrm{pj}^*\boldsymbol{\lambda}$ is a real function of $c$ on $T_2$.
The coboundary $d\boldsymbol{\lambda}$, the restriction to $\mathbb{Z}^3$, and the products $\exp\left(\pm i\cdot e_k\cdot\boldsymbol{\lambda}(\tilde{c})\right)\cdot f(\tilde{c},g)$, with the same phase at the $48$ sites over $\tilde{c}$ by (I.POS), and $\exp\left(\pm i\cdot e_k\cdot\boldsymbol{\lambda}^{u}(\tilde{z})\right)\cdot v(\tilde{z})$ are defined on the covering lattice by formulas in which the side does not occur.
Hence
\begin{gather*}
\mathrm{pj}^*(A-d\boldsymbol{\lambda})=A_2-d\,\mathrm{pj}^*\boldsymbol{\lambda},\qquad(\mathrm{pj}^*\boldsymbol{\lambda})^{u}=\mathrm{pj}^*(\boldsymbol{\lambda}^{u}),\\
\mathrm{pj}^*\left(\exp\left(\pm i\cdot e_k\cdot\boldsymbol{\lambda}\right)\cdot f\right)=\exp\left(\pm i\cdot e_k\cdot\mathrm{pj}^*\boldsymbol{\lambda}\right)\cdot\mathrm{pj}^*f,\qquad\mathrm{pj}^*\left(\exp\left(\pm i\cdot e_k\cdot\boldsymbol{\lambda}^{u}\right)\cdot v\right)=\exp\left(\pm i\cdot e_k\cdot(\mathrm{pj}^*\boldsymbol{\lambda})^{u}\right)\cdot\mathrm{pj}^*v.
\end{gather*}
In each of the five chains of equalities below, the first equality is the definition of the projection combined with the first of these relations, and the second is W13 on $T_2$ at $(s''_2,A_2,\mathrm{pj}^*\boldsymbol{\lambda})$, resp. at $(X',A_2,\mathrm{pj}^*\boldsymbol{\lambda})$ for every $X'$ with $\mathrm{pj}(X')=X$.
The third equality is the third or the fourth relation, applied to the phase on the right; the fourth equality is the definition of the projection; the fifth is the third or the fourth relation, applied to the phase on the left.

\emph{(p3), $k\ge1$.} For $v\in\mathfrak{S}(1,\ell_1)$,
\begin{align*}
\mathrm{pj}^*\left(H^{\mathrm{pj}}_k(s'',A-d\boldsymbol{\lambda})v\right)&=H^{(2)}_k(s''_2,A_2-d\,\mathrm{pj}^*\boldsymbol{\lambda})\,\mathrm{pj}^*v\\
&=\exp\left(i\cdot e_k\cdot\mathrm{pj}^*\boldsymbol{\lambda}\right)\cdot H^{(2)}_k(s''_2,A_2)\cdot\exp\left(-i\cdot e_k\cdot(\mathrm{pj}^*\boldsymbol{\lambda})^{u}\right)\cdot\mathrm{pj}^*v\\
&=\exp\left(i\cdot e_k\cdot\mathrm{pj}^*\boldsymbol{\lambda}\right)\cdot H^{(2)}_k(s''_2,A_2)\,\mathrm{pj}^*\left(\exp\left(-i\cdot e_k\cdot\boldsymbol{\lambda}^{u}\right)\cdot v\right)\\
&=\exp\left(i\cdot e_k\cdot\mathrm{pj}^*\boldsymbol{\lambda}\right)\cdot\mathrm{pj}^*\left(H^{\mathrm{pj}}_k(s'',A)\exp\left(-i\cdot e_k\cdot\boldsymbol{\lambda}^{u}\right)\cdot v\right)\\
&=\mathrm{pj}^*\left(\exp\left(i\cdot e_k\cdot\boldsymbol{\lambda}\right)\cdot H^{\mathrm{pj}}_k(s'',A)\exp\left(-i\cdot e_k\cdot\boldsymbol{\lambda}^{u}\right)\cdot v\right).
\end{align*}

\emph{$\Gamma_k$, $k\ge1$.} For $v\in\mathfrak{S}(1,\ell_1)$,
\begin{align*}
\mathrm{pj}^*\left(\Gamma^{\mathrm{pj}}_k(s'',A-d\boldsymbol{\lambda})v\right)&=\Gamma^{(2)}_k(s''_2,A_2-d\,\mathrm{pj}^*\boldsymbol{\lambda})\,\mathrm{pj}^*v\\
&=\exp\left(i\cdot e_k\cdot(\mathrm{pj}^*\boldsymbol{\lambda})^{u}\right)\cdot\Gamma^{(2)}_k(s''_2,A_2)\cdot\exp\left(-i\cdot e_k\cdot(\mathrm{pj}^*\boldsymbol{\lambda})^{u}\right)\cdot\mathrm{pj}^*v\\
&=\exp\left(i\cdot e_k\cdot(\mathrm{pj}^*\boldsymbol{\lambda})^{u}\right)\cdot\Gamma^{(2)}_k(s''_2,A_2)\,\mathrm{pj}^*\left(\exp\left(-i\cdot e_k\cdot\boldsymbol{\lambda}^{u}\right)\cdot v\right)\\
&=\exp\left(i\cdot e_k\cdot(\mathrm{pj}^*\boldsymbol{\lambda})^{u}\right)\cdot\mathrm{pj}^*\left(\Gamma^{\mathrm{pj}}_k(s'',A)\exp\left(-i\cdot e_k\cdot\boldsymbol{\lambda}^{u}\right)\cdot v\right)\\
&=\mathrm{pj}^*\left(\exp\left(i\cdot e_k\cdot\boldsymbol{\lambda}^{u}\right)\cdot\Gamma^{\mathrm{pj}}_k(s'',A)\exp\left(-i\cdot e_k\cdot\boldsymbol{\lambda}^{u}\right)\cdot v\right).
\end{align*}

\emph{$M_k$, $k\ge1$.} For $v\in\mathfrak{S}(1,\ell_1)$, all sums extending over the $X'\in D_k(T_2)$ with $\mathrm{pj}(X')=X$,
\begin{align*}
\mathrm{pj}^*\left(M^{\mathrm{pj}}_k(X,A-d\boldsymbol{\lambda})v\right)&=\sum_{X'}M^{(2)}_k(X',A_2-d\,\mathrm{pj}^*\boldsymbol{\lambda})\,\mathrm{pj}^*v\\
&=\exp\left(i\cdot e_k\cdot(\mathrm{pj}^*\boldsymbol{\lambda})^{u}\right)\cdot\sum_{X'}M^{(2)}_k(X',A_2)\cdot\exp\left(-i\cdot e_k\cdot(\mathrm{pj}^*\boldsymbol{\lambda})^{u}\right)\cdot\mathrm{pj}^*v\\
&=\exp\left(i\cdot e_k\cdot(\mathrm{pj}^*\boldsymbol{\lambda})^{u}\right)\cdot\sum_{X'}M^{(2)}_k(X',A_2)\,\mathrm{pj}^*\left(\exp\left(-i\cdot e_k\cdot\boldsymbol{\lambda}^{u}\right)\cdot v\right)\\
&=\exp\left(i\cdot e_k\cdot(\mathrm{pj}^*\boldsymbol{\lambda})^{u}\right)\cdot\mathrm{pj}^*\left(M^{\mathrm{pj}}_k(X,A)\exp\left(-i\cdot e_k\cdot\boldsymbol{\lambda}^{u}\right)\cdot v\right)\\
&=\mathrm{pj}^*\left(\exp\left(i\cdot e_k\cdot\boldsymbol{\lambda}^{u}\right)\cdot M^{\mathrm{pj}}_k(X,A)\exp\left(-i\cdot e_k\cdot\boldsymbol{\lambda}^{u}\right)\cdot v\right).
\end{align*}

\emph{$S_{k,y}$, $k\ge0$.} For $f\in\mathfrak{S}^{\mathrm{ch}}(\eta_k,\ell_1)$, with the same range of the sums over $X'$,
\begin{align*}
\mathrm{pj}^*\left(S^{\mathrm{pj}}_{k,y}(X,A-d\boldsymbol{\lambda})f\right)&=\sum_{X'}S^{(2)}_{k,y}(X',A_2-d\,\mathrm{pj}^*\boldsymbol{\lambda})\,\mathrm{pj}^*f\\
&=\exp\left(i\cdot e_k\cdot\mathrm{pj}^*\boldsymbol{\lambda}\right)\cdot\sum_{X'}S^{(2)}_{k,y}(X',A_2)\cdot\exp\left(-i\cdot e_k\cdot\mathrm{pj}^*\boldsymbol{\lambda}\right)\cdot\mathrm{pj}^*f\\
&=\exp\left(i\cdot e_k\cdot\mathrm{pj}^*\boldsymbol{\lambda}\right)\cdot\sum_{X'}S^{(2)}_{k,y}(X',A_2)\,\mathrm{pj}^*\left(\exp\left(-i\cdot e_k\cdot\boldsymbol{\lambda}\right)\cdot f\right)\\
&=\exp\left(i\cdot e_k\cdot\mathrm{pj}^*\boldsymbol{\lambda}\right)\cdot\mathrm{pj}^*\left(S^{\mathrm{pj}}_{k,y}(X,A)\exp\left(-i\cdot e_k\cdot\boldsymbol{\lambda}\right)\cdot f\right)\\
&=\mathrm{pj}^*\left(\exp\left(i\cdot e_k\cdot\boldsymbol{\lambda}\right)\cdot S^{\mathrm{pj}}_{k,y}(X,A)\exp\left(-i\cdot e_k\cdot\boldsymbol{\lambda}\right)\cdot f\right).
\end{align*}

\emph{Level $0$.} Here $\boldsymbol{\lambda}^{u}=\boldsymbol{\lambda}$ acts on $\mathfrak{S}^{\mathrm{ch}}(1,\ell_i)$ by (I.POS), and for $f\in\mathfrak{S}^{\mathrm{ch}}(1,\ell_1)$
\begin{align*}
\mathrm{pj}^*\left(\Gamma^{\mathrm{pj}}_0(s'',A-d\boldsymbol{\lambda})f\right)&=\Gamma^{(2)}_0(s''_2,A_2-d\,\mathrm{pj}^*\boldsymbol{\lambda})\,\mathrm{pj}^*f\\
&=\exp\left(i\cdot e_k\cdot\mathrm{pj}^*\boldsymbol{\lambda}\right)\cdot\Gamma^{(2)}_0(s''_2,A_2)\cdot\exp\left(-i\cdot e_k\cdot\mathrm{pj}^*\boldsymbol{\lambda}\right)\cdot\mathrm{pj}^*f\\
&=\exp\left(i\cdot e_k\cdot\mathrm{pj}^*\boldsymbol{\lambda}\right)\cdot\Gamma^{(2)}_0(s''_2,A_2)\,\mathrm{pj}^*\left(\exp\left(-i\cdot e_k\cdot\boldsymbol{\lambda}\right)\cdot f\right)\\
&=\exp\left(i\cdot e_k\cdot\mathrm{pj}^*\boldsymbol{\lambda}\right)\cdot\mathrm{pj}^*\left(\Gamma^{\mathrm{pj}}_0(s'',A)\exp\left(-i\cdot e_k\cdot\boldsymbol{\lambda}\right)\cdot f\right)\\
&=\mathrm{pj}^*\left(\exp\left(i\cdot e_k\cdot\boldsymbol{\lambda}\right)\cdot\Gamma^{\mathrm{pj}}_0(s'',A)\exp\left(-i\cdot e_k\cdot\boldsymbol{\lambda}\right)\cdot f\right).
\end{align*}
Since $\mathrm{pj}^*$ is injective, the relations of W13 hold on $T_1$.

\textbf{W14.}
Let $\tilde{g}=(S,a)\in\mathbb{G}_k^{L}(T_1)$, and for (p5) let $\tilde{g}\in\mathbb{G}^{L}_{k,3}(T_1)$, that is, $S\in\mathrm{Sp}_3$.
Put $r:=r(S)\in\{\pm1\}^3$, choose a representative $a_1\in(L\cdot M)\mathbb{Z}^3$ of $a$, and define, in accordance with (I.4.2),
\[
\mathrm{mph}(\tilde{c}):=r\cdot\tilde{c}+a_1\ \text{ on }\mathbb{L}_k,\qquad\mathrm{mph}(\tilde{c},g):=(r\cdot\tilde{c}+a_1,\ r\cdot g)\ \text{ on }\mathbb{L}^{\mathrm{f}}_k.
\]
Since $r$ is diagonal with the entries $\pm1$ and $a_1\in(L\cdot M)\mathbb{Z}^3$, the map $\mathrm{mph}$ maps each of the following sets onto itself: $\mathbb{L}_k$, $\mathbb{Z}^3$, the set of the bonds, and the sets of the $L$-blocks, of the $M$-cubes, of the $LM$-cubes and of the parameter cubes of each kind.
On $\mathbb{L}^{\mathrm{f}}_k$ it is a bijection with $\pi(\mathrm{mph}\,x)=\mathrm{mph}(\pi(x))$, by (I.a1); here $r\cdot g\in O_{\mathrm{cub}}$, because $\{\pm1\}^3$ is a subgroup of $O_{\mathrm{cub}}$.
Moreover $\mathrm{mph}(t+\ell_i\cdot\mathrm{dk})=\mathrm{mph}(t)+\ell_i\cdot r\cdot\mathrm{dk}$.
By Section~I.4.3, the action of $\tilde{g}$ on the objects on $T_1$, formed with the representative $a_1$, is given on the covering index sets by
\begin{gather*}
(\tilde{g}A)(\tilde{c},\tilde{c}')=A(\mathrm{mph}^{-1}\tilde{c},\mathrm{mph}^{-1}\tilde{c}'),\qquad(\tilde{g}f)(x)=S\cdot f(\mathrm{mph}^{-1}x),\qquad(\tilde{g}v)(\tilde{z})=S\cdot v(\mathrm{mph}^{-1}\tilde{z}),\\
(\tilde{g}s)_{\widetilde{\mathrm{pc}}}=s_{\mathrm{mph}^{-1}\widetilde{\mathrm{pc}}},\qquad\tilde{g}\,\mathrm{pr}_1(\tilde{X})=\mathrm{pr}_1(\mathrm{mph}\,\tilde{X});
\end{gather*}
on $V_k$ it is the action on the bond function $\mathsf{Z}$, which maps $V_k(T_1)$ onto itself by \cite[Proposition~15.21(ii)]{QEDI}.
The element $\tilde{g}_2:=(S,\,a_1\bmod\ell_2\mathbb{Z}^3)$ belongs to $\mathbb{G}_k^{L}(T_2)$, and to $\mathbb{G}^{L}_{k,3}(T_2)$ if $S\in\mathrm{Sp}_3$; its action on the objects on $T_2$, formed with the same representative $a_1$, is given by the same five formulas with $\mathrm{pr}_2$ in place of $\mathrm{pr}_1$.
Consequently, for every object $f$ on $T_1$ (a function, a connection or a parameter family) and every $X'\in D_k(T_2)$,
\begin{equation}
\tilde{g}_2\,\mathrm{pj}^*f=\mathrm{pj}^*\,\tilde{g}f,\qquad\mathrm{pj}(\tilde{g}_2X')=\tilde{g}\,\mathrm{pj}(X').
\tag{3.b}
\end{equation}
Indeed, the two sides of the first relation are one and the same function on the covering index set; and if $X'=\mathrm{pr}_2(\tilde{X})$, then
\[
\mathrm{pj}(\tilde{g}_2X')=\mathrm{pj}\left(\mathrm{pr}_2(\mathrm{mph}\,\tilde{X})\right)=\mathrm{pr}_1(\mathrm{mph}\,\tilde{X})=\tilde{g}\,\mathrm{pr}_1(\tilde{X})=\tilde{g}\,\mathrm{pj}(X').
\]
If $a_1$ is replaced by another representative $a_1+\ell_1\cdot\mathrm{dk}$, then $\tilde{g}f$ is multiplied by $\boldsymbol{\chi}(r\cdot\mathrm{dk})=\boldsymbol{\chi}(\mathrm{dk})$, on the source space as well as on the target space of each family (Section~3.0); hence the relations W14 on $T_1$ do not depend on the representative.
Moreover $\tilde{g}s\in\mathrm{Pd}_k(T_1)$, its entries being those of $s$, and $\tilde{g}A\in\tilde{G}_k(T_1)$, because $\tilde{G}_k$ is invariant under $\mathbb{G}_k$ by \cite[Lemma~17.6]{QEDI}.

\emph{Families with parameters} ((p1)--(p3) and $\Gamma_k$; for (p1), (p2) the argument $A$ is absent).
For $f$ in the source space on $T_1$, the definition of the projection, (3.b) and W14 on $T_2$ at $\tilde{g}_2$ give
\begin{align*}
\mathrm{pj}^*\left(\mathrm{Op}^{\mathrm{pj}}(\tilde{g}s,\tilde{g}A)\,\tilde{g}f\right)&=\mathrm{Op}^{(2)}(\mathrm{pj}^*\tilde{g}s,\mathrm{pj}^*\tilde{g}A)\,\mathrm{pj}^*\tilde{g}f=\mathrm{Op}^{(2)}(\tilde{g}_2s_2,\tilde{g}_2A_2)\,\tilde{g}_2\,\mathrm{pj}^*f\\
&=\tilde{g}_2\,\mathrm{Op}^{(2)}(s_2,A_2)\,\mathrm{pj}^*f=\tilde{g}_2\,\mathrm{pj}^*\left(\mathrm{Op}^{\mathrm{pj}}(s,A)f\right)=\mathrm{pj}^*\left(\tilde{g}\,\mathrm{Op}^{\mathrm{pj}}(s,A)f\right).
\end{align*}

\emph{$M_k$ ($k\ge1$) and $S_{k,y}$.}
By (3.b), $X'\mapsto\tilde{g}_2X'$ maps $\{X':\mathrm{pj}(X')=X\}$ into $\{X'':\mathrm{pj}(X'')=\tilde{g}X\}$, and $X''\mapsto\tilde{g}_2^{-1}X''$ maps the second set into the first; so the first map is a bijection between the two sets.
Together with W14 on $T_2$ at $\tilde{g}_2$, where $\tilde{g}_2\in\mathbb{G}^{L}_{k,3}(T_2)$ in the case of $S_{k,y}$, this yields
\begin{align*}
\mathrm{pj}^*\left(\mathrm{Op}^{\mathrm{pj}}(\tilde{g}X,\tilde{g}A)\,\tilde{g}f\right)&=\sum_{X'':\ \mathrm{pj}(X'')=\tilde{g}X}\mathrm{Op}^{(2)}(X'',\tilde{g}_2A_2)\,\tilde{g}_2\,\mathrm{pj}^*f=\sum_{X':\ \mathrm{pj}(X')=X}\mathrm{Op}^{(2)}(\tilde{g}_2X',\tilde{g}_2A_2)\,\tilde{g}_2\,\mathrm{pj}^*f\\
&=\tilde{g}_2\sum_{X':\ \mathrm{pj}(X')=X}\mathrm{Op}^{(2)}(X',A_2)\,\mathrm{pj}^*f=\tilde{g}_2\,\mathrm{pj}^*\left(\mathrm{Op}^{\mathrm{pj}}(X,A)f\right)=\mathrm{pj}^*\left(\tilde{g}\,\mathrm{Op}^{\mathrm{pj}}(X,A)f\right).
\end{align*}
As $\mathrm{pj}^*$ is injective, $\mathrm{Op}^{\mathrm{pj}}(\tilde{g}s,\tilde{g}A)\,\tilde{g}=\tilde{g}\,\mathrm{Op}^{\mathrm{pj}}(s,A)$ and $\mathrm{Op}^{\mathrm{pj}}(\tilde{g}X,\tilde{g}A)\,\tilde{g}=\tilde{g}\,\mathrm{Op}^{\mathrm{pj}}(X,A)$.
\emph{Level $0$.} On $\mathfrak{S}^{\mathrm{ch}}(1,\ell_i)$ the action reads $(\tilde{g}f)(\tilde{c},g)=S\cdot f\left(r\cdot(\tilde{c}-a_1),\,r\cdot g\right)$: the label is multiplied by $r$.

\textbf{W15} ((p4)).
By \cite[Lemma~17.6]{QEDI}, $-A\in\tilde{G}_k(T_1)$, and $\mathrm{pj}^*(-A)=-A_2$ by (1.7).
By (a), the operator $O^{(2)}:=\Gamma^{(2)}_k(s''_2,A_2)=\widetilde{\mathrm{Op}}(s'',A)$ commutes with the $T_v$, $v\in\ell_1\mathbb{Z}^3$, and $O^{(2)}\,\mathrm{pj}^*=\mathrm{pj}^*O^{(1)}$ with $O^{(1)}:=\Gamma^{\mathrm{pj}}_k(s'',A)$.
Using (e), W15 on $T_2$, (1.7) for the constant matrices $\mathsf{C}$ and $\mathsf{C}^{-1}$, and the definition of the projection at $-A$, one finds
\begin{align*}
\mathrm{pj}^*\left(\Gamma^{\mathrm{pj}}_k(s'',A)\right)^T&=\left(\Gamma^{(2)}_k(s''_2,A_2)\right)^T\mathrm{pj}^*=\mathsf{C}^{-1}\Gamma^{(2)}_k(s''_2,-A_2)\,\mathsf{C}\,\mathrm{pj}^*=\mathsf{C}^{-1}\Gamma^{(2)}_k\left(s''_2,\mathrm{pj}^*(-A)\right)\mathrm{pj}^*\,\mathsf{C}\\
&=\mathsf{C}^{-1}\mathrm{pj}^*\,\Gamma^{\mathrm{pj}}_k(s'',-A)\,\mathsf{C}=\mathrm{pj}^*\,\mathsf{C}^{-1}\Gamma^{\mathrm{pj}}_k(s'',-A)\,\mathsf{C}.
\end{align*}
For $k\ge1$ put $O^{(2)}:=\widetilde{\mathrm{Op}}(X,A)$ and $O^{(1)}:=M^{\mathrm{pj}}_k(X,A)$; with W15 on $T_2$ for every $X'$ with $\mathrm{pj}(X')=X$,
\begin{align*}
\mathrm{pj}^*\left(M^{\mathrm{pj}}_k(X,A)\right)^T&=\sum_{X':\ \mathrm{pj}(X')=X}\left(M^{(2)}_k(X',A_2)\right)^T\mathrm{pj}^*=\sum_{X':\ \mathrm{pj}(X')=X}\mathsf{C}^{-1}M^{(2)}_k(X',-A_2)\,\mathsf{C}\,\mathrm{pj}^*\\
&=\mathsf{C}^{-1}\mathrm{pj}^*\,M^{\mathrm{pj}}_k(X,-A)\,\mathsf{C}=\mathrm{pj}^*\,\mathsf{C}^{-1}M^{\mathrm{pj}}_k(X,-A)\,\mathsf{C}.
\end{align*}
Since $\mathrm{pj}^*$ is injective, W15 holds on $T_1$.
For $k\ge1$ the transposes are taken with respect to $\langle\,\cdot\,,\,\cdot\,\rangle_1$ on $\mathfrak{S}(1,\ell_i)$.
\emph{Level $0$.} The transposes are taken with respect to $\langle\,\cdot\,,\,\cdot\,\rangle^{\mathrm{ch}}$ on $\mathfrak{S}^{\mathrm{ch}}(1,\ell_i)$, with the weight $\frac{1}{48}$; (e) and (1.3) apply to this pairing.

\textbf{W16.}
For an operator $O$, $\overline{O}$ denotes the operator $f\mapsto\overline{O\,\overline{f}}$, as in (I.O5).
One has $\overline{s}\in\mathrm{Pd}_k(T_1)$, and $\overline{A}$, $-A$ belong to $\tilde{G}_k(T_1)$ by \cite[Lemma~17.6]{QEDI}.
By (1.7), $\mathrm{pj}^*\overline{s}=\overline{s_2}$, $\mathrm{pj}^*\overline{A}=\overline{A_2}$, $\mathrm{pj}^*\overline{f}=\overline{\mathrm{pj}^*f}$ and $\mathrm{pj}^*\varkappa f=\varkappa\,\mathrm{pj}^*f$.
For $O=H_k$ ($k\ge1$) and for $O=\Gamma_k$ ($k\ge0$), W16 on $T_2$ gives
\[
\mathrm{pj}^*\left(\overline{O^{\mathrm{pj}}(\overline{s''},\overline{A})\,\overline{f}}\right)=\overline{O^{(2)}(\overline{s''_2},\overline{A_2})\,\overline{\mathrm{pj}^*f}}=\varkappa^{-1}O^{(2)}(s''_2,-A_2)\,\varkappa\,\mathrm{pj}^*f=\mathrm{pj}^*\left(\varkappa^{-1}O^{\mathrm{pj}}(s'',-A)\,\varkappa f\right).
\]
For $M_k$ ($k\ge1$), the sums extending over the $X'$ with $\mathrm{pj}(X')=X$,
\[
\mathrm{pj}^*\left(\overline{M^{\mathrm{pj}}_k(X,\overline{A})\,\overline{f}}\right)=\sum_{X'}\overline{M^{(2)}_k(X',\overline{A_2})\,\overline{\mathrm{pj}^*f}}=\sum_{X'}\varkappa^{-1}M^{(2)}_k(X',-A_2)\,\varkappa\,\mathrm{pj}^*f=\mathrm{pj}^*\left(\varkappa^{-1}M^{\mathrm{pj}}_k(X,-A)\,\varkappa f\right).
\]
Since $\mathrm{pj}^*$ is injective, W16 holds on $T_1$ for (p3) and (p4).
For (p1) ($k\ge1$) and (p2) ($k\ge0$): if $s'$ is real, then so is $s'_2$; the kernels $\widetilde{\mathrm{Op}}(s')(t,t'')$ are real by W16 on $T_2$, and by (8.0), in which $\boldsymbol{\chi}=1$, the kernels of $\mathrm{Op}^{\mathrm{pj}}(s')$ are sums of real numbers.

\textbf{W17} ($k=0$).
The diagonal families are, by (c), the projections of (p1) and (p3); they have W1, W2, W8--W11, W13, W14 and W16, each in so far as it concerns the family, and W3, W5 with the constant $4\cdot\exp(2\gamma_{\mathrm{w}})\cdot M^{27\alpha_{0}}$.
For the other families the arguments above, taken at $k=0$, give the following properties: for (p2), W1, W2, W4, W8, W9, W10, W14 and W16; for $\Gamma^{\mathrm{pj}}_0$, W1, W2, W6, W8, W9, W13, W14, W15 and W16; for $S^{\mathrm{pj}}_{0,y}$, W1, W2, W7 in its form for $k=0$, W12, W13, W14 (for $\mathbb{G}^{L}_{0,3}$) and the Borel measurability in $y$.
At $k=0$ these arguments rest, as far as the system on $T_2$ is concerned, on the property with the same number, on W14 (through (a)) and, for W7, on W12; for the families in question all of them are contained in W17 on $T_2$.
\end{proof}

\subsubsection*{3.8. Constants; uniformity}

In going from $\mathrm{par}^{0}$ to $\mathrm{par}$, the exponent $\alpha_{0}$ and the rate $\gamma_{\mathrm{w}}$ remain unchanged.
In (3.a) the coefficient of $d_M$ in the exponent decreases from $\kappa_{\mathrm{w}}+\kappa_{0}$ to $\kappa_{\mathrm{w}}$;
the difference $\kappa_{0}$ is absorbed by (t2) of \cite[Theorem~13.2]{QEDI}.
As to the constant, the bounds of Section~3.7 carry $C_{w}^{(1)}$ in W5,
$64C_{w}^{(1)}$ in W3, in W4 and in W6 for $\Gamma_k$,
and $64K_{0}\cdot C_{w}^{(1)}$ in W6 for $M_k$ and in W7;
each of these is at most $C_{w}=64K_{0}\cdot C_{w}^{(1)}$.
Here $64$ and $27.3$ are absolute numbers, and $K_{0}$ is the number $\frac{\exp\left(\kappa_{0}\right)}{36}$ of the line (C0).
The proof depends on the data $(k,\ell_1,\ell_2,\mathrm{Jr},\mu,e_k)$ only through the following facts:
$\ell_1$ and $\ell_2$ are multiples of $LM$ with $\ell_i\ge13L\cdot M$;
$\mathrm{Jr}$ is odd;
the system on $T_2$ satisfies the hypotheses imposed on it;
and the data are (SF), which is the setting of the statements applied in the proof.
No bound contains $k$, $\ell_1$, $\ell_2$, $\mathrm{Jr}$, $\mu$, $e_k$, $y$, $s$, $A$ or $X$ otherwise than as displayed.
At level~$0$ the factor $48$ of (I.ke4) is contained in $C_{w}^{(1)}$, through $C^{F}$.
Moreover (8.0) is an identity between kernels that carry the same weight $\frac{1}{48}$ on both tori,
so that no further factor $48$ arises.

\subsection*{4. Proposition~26.4}

\begin{proof}[Proof of Proposition~26.4]
Throughout, $d$ denotes the sup distance of $T_k$; it satisfies the triangle inequality.
For a union $X$ of $M$-cubes, $X^{\sim i}$ is the union of the $M$-cubes whose centres have sup-distance at most $i\cdot M$ from the centre of some $M$-cube of $X$ (Section~I.1.1).
Consequently ${X}^{\sim i}\subseteq{X'}^{\sim i'}$ whenever $X\subseteq X'$ and $i\le i'$.
Since $\ell$ is a multiple of $M$, the distance of two points of $M\mathbb{Z}^3/\ell\mathbb{Z}^3$ is a multiple of $M$.
The indices $t$, their supports and the rows of the families are those of Section~G4.0.4,
and $\square(t)$ denotes, as there, the union of the $M$-cubes that $t$ meets;
the parameter cubes and the polydisk $\mathrm{Pd}_k$ are those of \cite[Definition~5.8]{QEDI}.
A parameter cube is contained in a set of sites if all its lattice sites lie in this set,
and it meets an $M$-cube if the two have a lattice site in common \cite[Section~5.2.3]{QEDI}.

Let $\mathscr{Y}$ be the set of the parameter cubes of the kind of $\mathrm{Op}$ that are contained in $Z$, and let $Y$ be its union.
As $s$ vanishes off $Z$, one has $s_{\mathrm{pc}}=0$ for every parameter cube $\mathrm{pc}$ outside $\mathscr{Y}$.

\emph{Step~1: there is a cube $\mathrm{pc}$ of $\mathscr{Y}$ that contains $\operatorname{supp}t$.}
In each of the three cases below a cube $\mathrm{pc}$ of the kind of $\mathrm{Op}$ is exhibited
which contains $\operatorname{supp}t$ and all of whose lattice sites lie in $\square(t)^{\sim c_{l}}$.
Since $\square(t)^{\sim c_{l}}\subseteq Z_1\subseteq Z$, such a cube is contained in $Z$;
that is, $\mathrm{pc}$ belongs to $\mathscr{Y}$.

\emph{Kind $\pi_k$} ($\mathrm{Op}=H_{k}(s'',A)$ or $\mathrm{Op}=\Gamma_k(s'',A)$).
The row index $t$ is either a fermion site $x=(c,g)$, as for the rows of (p3) at every level $k\ge0$ and for the rows of $\Gamma_0$,
or a unit site $z$, as for the rows of $\Gamma_k$ with $k\ge1$.
In either case $\operatorname{supp}t$ consists of one photon site, by (I.POS): $\operatorname{supp}x=\{c\}$.
By the fact (g0) of \cite[Section~17.4.1]{QEDI} this site lies in exactly one $M$-cube,
and this $M$-cube is $\square(t)$; note that $\square(x)=\square(c)$.
Put $\mathrm{pc}:=\square(t)$.
It is a cube of $\pi_k$, and it is contained in $\square(t)^{\sim c_{l}}$.

\emph{Kind $\pi'_k$} ($\mathrm{Op}=H_{k}(s')$).
Here $t$ is a fine bond with the end points $c$ and $c'$, and $d(c,c')=\eta_k\le1$.
Let $z$ be the centre of the $M$-cube $\square(c)$ of $c$, and put
\[
\mathrm{pc}:=\left\{c'':d(c'',z)\le M\right\},
\]
a cube of $\pi'_k$.
One has $d(c,z)<\frac{M}{2}$ and $d(c',z)<\frac{M}{2}+1\le M$, the last inequality because $M\ge2$;
so $\operatorname{supp}t\subseteq\mathrm{pc}$.
By the fact (g1) of \cite[Section~17.4.1]{QEDI}, the lattice sites of $\mathrm{pc}$ lie in the union of the $M$-cubes whose centres have sup-distance at most $M$ from $z$,
and this union is $\square(c)^{\sim1}$.
Since $\square(c)$ is one of the $M$-cubes that $t$ meets, and $c_{l}=2L\ge1$, one obtains
\[
\square(c)^{\sim1}\subseteq\square(t)^{\sim1}\subseteq\square(t)^{\sim c_{l}}.
\]

\emph{Kind $\pi''_k$} ($\mathrm{Op}=C_{k}^{1/2}(s')$).
Here $t=\Upsilon$.
Let $c$ be a point of $\operatorname{supp}\Upsilon$,
and let $z\in LM\mathbb{Z}^3/\ell\mathbb{Z}^3$ be the centre of the $LM$-cube that contains $c$;
the $LM$-cubes partition the sites (Section~I.1.1, on faces, and Section~I.1.4).
Then $d(c,z)<\frac{L\cdot M}{2}$.
Put
\[
\mathrm{pc}:=\left\{c'':d(c'',z)\le L\cdot M\right\},
\]
a cube of $\pi''_k$.
The sup-diameter of $\operatorname{supp}\Upsilon$ is at most $L$ \cite[Proposition~5.35]{QEDI},
and $L\le\frac{L\cdot M}{2}$ because $M\ge2$.
Hence every $c'\in\operatorname{supp}\Upsilon$ satisfies
\[
d(c',z)\le d(c',c)+d(c,z)<L+\frac{L\cdot M}{2}\le L\cdot M,
\]
and therefore $\operatorname{supp}t\subseteq\mathrm{pc}$.
Now let $c''$ be a lattice site of $\mathrm{pc}$, and let $z''$ and $z_c$ be the centres of the $M$-cubes that contain $c''$ and $c$, respectively.
Then
\[
d(z'',z_c)\le d(z'',c'')+d(c'',z)+d(z,c)+d(c,z_c)<\frac{M}{2}+L\cdot M+\frac{L\cdot M}{2}+\frac{M}{2}=\left(\frac{3L}{2}+1\right)\cdot M\le(2L+1)\cdot M.
\]
The number $d(z'',z_c)$ is a multiple of $M$; it follows that $d(z'',z_c)\le2L\cdot M$.
Hence the $M$-cube of $c''$ is one of the $M$-cubes of $\square(c)^{\sim2L}$.
Finally, $\square(c)^{\sim2L}\subseteq\square(t)^{\sim c_{l}}$, because $\square(c)$ is one of the $M$-cubes that $t$ meets and $c_{l}=2L$.

\emph{Step~2.}
Let $\mathscr{Y}_1$ be the link-component of $\mathscr{Y}$ that contains the cube $\mathrm{pc}$ of Step~1, and let $Y_1$ be its union.
By Section~I.1.1, $Y_1$ is a link-component of $Y$.
Since $\operatorname{supp}t$ lies in $\mathrm{pc}$, the index $t$ is contained in $Y_1$.

\emph{Step~3: $Y_1\subseteq Z_1$.}
Let $\mathrm{pol}(\mathscr{Y}_1)$ be the union of the $M$-cubes that meet a cube of $\mathscr{Y}_1$.
By the fact (g2) of \cite[Section~17.4.1]{QEDI}, $\mathrm{pol}(\mathscr{Y}_1)$ is a polymer, and it contains $Y_1$.
Each $M$-cube of $\mathrm{pol}(\mathscr{Y}_1)$ has a lattice site in common with a cube of $\mathscr{Y}$,
and this site lies in $Z$, because the cubes of $\mathscr{Y}$ are contained in $Z$.
As $Z$ is a union of $M$-cubes and the $M$-cubes partition the sites (g0),
each $M$-cube of $\mathrm{pol}(\mathscr{Y}_1)$ is one of the $M$-cubes of $Z$.
The polymer $\mathrm{pol}(\mathscr{Y}_1)$ contains $\operatorname{supp}t$;
so it contains the $M$-cube of a point of $\operatorname{supp}t$, which is an $M$-cube of $\square(t)$, hence of $Z_1$.
Both $Z_1$ and $\mathrm{pol}(\mathscr{Y}_1)$ are connected through faces, and they have an $M$-cube in common;
therefore their union is a union of $M$-cubes of $Z$ that is connected through faces.
Since $Z_1$ is a maximal subset of $Z$ that is connected through faces, it follows that $\mathrm{pol}(\mathscr{Y}_1)\subseteq Z_1$,
and thus $Y_1\subseteq Z_1$.

\emph{Step~4: W8.}
Put $\tilde{s}:=s\,\mathbf{1}_{Z_1}$ and $\tilde{A}:=A_1$.
The entries of $\tilde{s}$ are entries of $s$ or $0$; so $\tilde{s}$ lies in $\mathrm{Pd}_k$.
The family $\tilde{s}$ vanishes outside $\mathscr{Y}$:
a parameter cube that does not belong to $\mathscr{Y}$ is not contained in $Z$, hence not contained in $Z_1$.
On the cubes of $\mathscr{Y}_1$ the family $\tilde{s}$ is equal to $s$,
since these cubes are contained in $Y_1$, hence in $Z_1$ by Step~3.
The connection $A_1$ lies in $\tilde{G}_k$, and it is equal to $A$ on the bonds that meet $Y_1$:
such a bond has an end point in $Y_1$, hence in $Z_1$, so that it meets $Z_1$ (Section~I.1.1).
The property W8$^{\mathrm{ch}}$ of the system (Section~G4.0.4), applied to $(\mathscr{Y},Y,s,Y_1)$, to the row index $t$ (Step~2) and to the pair $(\tilde{s},\tilde{A})$, gives
\begin{gather*}
\mathrm{Op}(s,A)(t,t')=0\quad\text{unless $t'$ is contained in $Y_1$,}\\
\mathrm{Op}(s,A)(t,\,\cdot\,)=\mathrm{Op}\left(s\,\mathbf{1}_{Z_1},A_1\right)(t,\,\cdot\,).
\end{gather*}
An index that is contained in $Y_1$ is contained in $Z_1$, by Step~3.
These are the two assertions of the display of Proposition~26.4.
For (p1) and (p2) the argument $A$ is absent.

\emph{Step~5.}
The sentence that follows the display of Proposition~26.4 is the property W12$^{\mathrm{ch}}$ of the system (Section~G4.0.4) for $M_k(X,A)$ ($k\ge1$) and for $S_{k,y}(X,A)$, with $\tilde{A}=A_1$.

\emph{Level~$0$.}
The families with parameters are $H_{0}(s')$, whose rows are the bonds of $T_0$ (in the case $\pi'_0$ one has $d(c,c')=\eta_0=1$),
$C_{0}^{1/2}(s')$,
and $H_{0}(s'',A)$ and $\Gamma_0(s'',A)$, whose rows are the fermion sites of $T^{\mathrm{f}}_0$ (case $\pi_0$).
There is no family $M_0$.
None of the steps changes.

\emph{Uniformity.}
The statement contains no constant.
The proof depends on $(k,\ell)$ only through the facts that $\eta_k\le1$, that $\ell$ is a multiple of $LM$, and that $\ell\in\mathrm{Cl}$, the last one in (g0)--(g2).
The numbers $\mu$ and $e_k$ enter only through the hypotheses on the system.
\end{proof}

\subsection*{5. Conditions on the constants}

The following conditions on the constants enter the proofs of this chapter (cf.\ Section~I.8.2 and Proposition~14.2\,(a), (b)).

\begin{description}
\item[$L$ odd ($L\ge5$).]
This condition is part of the line (C3) of Section~I.8.2; in Proposition~14.2\,(b) it is the condition (L5).
It enters through the fact (g0) of \cite[Section~17.4.1]{QEDI}: no photon site lies on a face, and the $M$-cubes and the $LM$-cubes partition the sites;
(g0) is applied in Sections~2, 3.1 (in (cv6): the $M$-cubes of $\mathrm{pj}\,X'$; in (cv7): the $M$-cube of $\tilde{c}'$), 3.5, 3.7 (W7) and~4.
Moreover, it belongs to the setting of \cite[Section~15]{QEDI} ($L$ odd, $L\ge3$),
and so enters through \cite[Proposition~5.35]{QEDI}, which is contained in \cite[Proposition~15.21]{QEDI} (Sections~0.2--0.4, 1, 2, 3.7 (W14) and~4),
through \cite[Propositions~5.31 and~5.39]{QEDI}, which are contained in \cite[Propositions~15.5 and~15.28]{QEDI} (Section~3.7, W1),
and through \cite[Proposition~15.36]{QEDI} (Section~3.6).
Finally, it belongs to the setting of \cite[Theorem~13.2]{QEDI} ($L$ odd), and so enters through (t2) of that theorem (Section~3.5).

\item[$M$ a power of $L$, $M\ge L^2$.]
Both conditions are part of the line (C5) of Section~I.8.2; together they form the condition (M1) of Proposition~14.2\,(b).
They enter through (g0) ($M$ is odd; the blocks and the cubes are nested)
and through the setting of \cite[Section~15.10]{QEDI}, for \cite[Proposition~15.36]{QEDI} (Section~3.6).
In the explicit inequalities of this chapter they enter only through the bound $M\ge2$.

\item[$M>4L$.]
The line (C5) contains this condition as well; in Proposition~14.2\,(b) it is the condition (M2).
It belongs to the setting of the facts (g0)--(g3) of \cite[Section~17.4.1]{QEDI}, which are applied in Sections~2, 3 and~4.

\item[$M\ge M^{P}$.]
This is a further condition of the line (C5), the condition (M6) of Proposition~14.2\,(b).
It belongs to the setting of \cite[Sections~15.10.3--15.10.5]{QEDI}, for \cite[Proposition~15.36]{QEDI} (Section~3.6).

\item[$\gamma_{\mathrm{w}}\cdot M\ge1$ (hence $\gamma_{\mathrm{w}}>0$).]
The line (C5) imposes it in the form $\gamma_{\mathrm{w}}\cdot M\ge16$, which is the condition (M3) of Proposition~14.2\,(b); it also stands in (A0) of Section~I.8.4.
It is applied in Section~3.4, where $\gamma_{\mathrm{w}}\cdot\left(\ell_1-2L\right)\ge6$ and the function $u\mapsto\exp\left(-\gamma_{\mathrm{w}}\cdot u\right)$ is decreasing,
and in Section~3.6, where the rate in \cite[Proposition~15.36]{QEDI} is $\gamma=\gamma_{\mathrm{w}}>0$.

\item[$\kappa_{\mathrm{w}}\ge0$.]
It follows from the lines (C4), (C2) and (C0) of Section~I.8.2:
(C4) defines $\kappa_{\mathrm{w}}:=\kappa+\kappa_{0}+3$ (a definition that also stands in (A0) of Section~I.8.4),
(C2) requires $\kappa\ge7\kappa_{0}+7$,
and (C0) gives $\kappa_{0}=54\log72>0$.
It is applied in (3.a), in the form $\exp\left(-\kappa_{\mathrm{w}}\cdot d_M(X')\right)\le\exp\left(-\kappa_{\mathrm{w}}\cdot d_M(X)\right)$.

\item[$K_{0}\ge1$.]
By the line (C0), $K_{0}=\frac{\exp\left(\kappa_{0}\right)}{36}=\frac{72^{54}}{36}$.
In Section~3.7 the condition enters through the inequality $64C_{w}^{(1)}\le C_{w}$ in W3, in W4 and in W6 for $\Gamma_k$, and through the inequality $C_{w}^{(1)}\le C_{w}$ in W5.

\item[$\kappa_{0}=54\log72$, $K_{0}=\frac{\exp\left(\kappa_{0}\right)}{36}$.]
Both are definitions of the line (C0); they enter through (t2) in (3.a).

\item[$C_{w}^{(1)}\ge4\exp\left(2\gamma_{\mathrm{w}}\right)\cdot M^{27\alpha_{0}}$ (hence $C_{w}^{(1)}>0$).]
This follows from the definition of $C_{w}^{(1)}$ in the line (C6).
The inequality is applied in Section~3.6, and the positivity $C_{w}^{(1)}>0$, together with $K_{0}\ge1$, in Section~3.7, W3--W6.

\item[$C_{w}=64K_{0}\cdot C_{w}^{(1)}$.]
This is the definition of $C_{w}$ in the line (C6); it is applied in Section~3.7, in W6 for $M_k$ and in W7.

\item[$0<\alpha_{0}\le\frac{1}{4}$.]
In the line (C4) one has $\alpha_{0}=\min\left(\alpha^{F},\alpha^{P}\right)$ with $\alpha^{F}=\frac{1}{54}$ and $\alpha^{P}=\frac{1}{4}$, which gives the condition.
It is part of the setting of \cite[Section~15.10]{QEDI}, on which Section~3.6 rests.

\item[Conditions of Propositions~11 and~8.]
The conditions of Proposition~11 are those of Proposition~2.8, and the conditions of Proposition~8 are those of Proposition~3; they are stated with these Propositions.
They enter through Proposition~11\,(b), in Section~3.7 (W1), and through Proposition~8\,(b), in Section~3.6.
\end{description}

The sides enter through the hypothesis $\ell\in\mathrm{Cl}$, which is part of the data:
$\ell$ is an odd multiple of $LM$ with $\ell\ge13L\cdot M$.
The inequalities between $\ell$, $L$ and $M$ that are used are consequences of $\ell\ge13L\cdot M$, $L\ge1$ and $M\ge2$; they are
\[
L+1<\ell_i\quad\text{(Section~0.3)},\qquad 8L\cdot M+4L<\ell\quad\text{(Section~2)},
\]
and
\[
2M+1<\ell_i,\qquad 4L+M<\ell_i,\qquad 1<\frac{\ell_i}{2},\qquad \ell_1-2L\ge6M\quad\text{(Section~3)}.
\]

\section*{Chapter 26b. The systems of primitives on the top torus; the data of the flow are data of one step: proofs of Propositions~26.5 and~26.6}

\subsection*{0. Conventions of this chapter}

The notation is that of Part~I and of the definitions in the Sections~G1 and~G4 of Part~II.

In this chapter $\gamma_{\mathrm{old}}$, $\mathrm{kq}_{\mathrm{old}}$, $C_{\mathrm{old}}$ and $\gamma_{\mathrm{new}}$, $\mathrm{kq}_{\mathrm{new}}$, $C_{\mathrm{new}}$ denote the entries of two sets of parameters that are compared with each other.
The labels (10.2) and (10.3) are labels of this chapter:
they mark the two displays of 1.4 that define the photon families (p1) and (p2), that is, the kernels of the weakened photon operators of \cite[Section~15.10.3]{QEDI}.
The operator $C_k^{1/2}$ and the family (10.3) at $k=0$ are written $C_0^{1/2}$ and $C_0^{1/2}(s')$;
the constant $C_0$ of the line (C6) of Section~I.8.2 does not occur in this chapter.
The letter $R'$ denotes an element of $\mathrm{Tor}(e)$ (in 2.4 and 2.8 it is $R$ or $R_{\mathrm{top}}$),
and $T$, $T'$ are the exponents in $R=hL^{T}$ and $R'=hL^{T'}$ (Section~I.6.2).
In 2.10 the letters $\mathrm{vh}_j$, $\mathrm{rh}_j$, $\mathrm{nch}_j(R)$ and $\mathrm{js}(R)$ denote the numbers $\hat{v}_j$, $\hat{r}_j$, $\hat{\mathfrak{n}}_j$ and $j_{*}(R,e)$ of \cite[Section~19.2]{QEDI};
thus $\mathrm{js}(R)$ is the integer $\mathrm{js}(R,e)$ of the assertion (Sc0) of Section~G4.0.2.
In 2.10, moreover, $N'$ is an integer $\ge0$ that stands in the place of $N$, $K':=N'+\log_L h$, and $t:=\log_L\left(\frac{1}{\bar{m}}\right)$.
The letter $t$ in the inequality of 1.4 that defines $M^{P}$ is a variable of that inequality; it is not the number $t$ of 2.10.
Finally, (2.1) is the condition (2.1) of Section~G4.0.3.

\subsection*{1. Proposition 26.5 (the systems on the top torus)}

\begin{proof}[Proof of Proposition~26.5]
Let $(k,\ell,\mu,e_k)$ be data (SF), and let $T_k=\mathbb{L}_k/\ell\mathbb{Z}^3$.

\paragraph{1.1. What has to be shown.}
For each family under consideration, $\mathrm{Pd}_k$ is the closed polydisk
\[
\mathrm{Pd}_k=\left\{s=(s_{\mathrm{pc}}):\ s_{\mathrm{pc}}\in\mathbb{C},\ |s_{\mathrm{pc}}|\le M^{\alpha_0}\right\}
\]
over the index set of that family,
namely $\pi'_k$ for (p1), $\pi''_k$ for (p2), and $\pi_k$ for (p3) and for $\Gamma_k$ (Section~I.1.1; \cite[Definition~5.8]{QEDI}).
As in Section~G1.11, a set of parameters is written $(\alpha_0,\gamma,\mathrm{kq};C)$:
$M^{\alpha_0}$ is the radius of $\mathrm{Pd}_k$, $\gamma$ a rate per unit length, $\mathrm{kq}$ a rate per unit of $d_M$, and $C$ a constant.
That a family has the property W$n$ for $(\alpha_0,\gamma,\mathrm{kq};C)$ is meant in the sense of Section~G4.0.4.
No polymer occurs in the families (p1) and (p2);
for them the entry $\mathrm{kq}$ is void, and a point is written in its place.
The two sets of parameters of the assertion are
\[
\mathrm{par}^{0}=\left(\alpha_0,\gamma_{\mathrm{w}},\kappa_{\mathrm{w}}+\kappa_0;C_{w}^{(1)}\right),\qquad
\mathrm{par}=\left(\alpha_0,\gamma_{\mathrm{w}},\kappa_{\mathrm{w}};C_{w}\right).
\]
Three assertions have to be proved.
\begin{itemize}
\item[(A)] The families named in Proposition~26.5 are families of linear operators with the parameters, the source spaces and the target spaces that Section~G4.0.4 prescribes for a system of level $k$.
\item[(B)] For $1\le n\le17$, every family that is named among the families concerned at the end of the property W$n$ of Section~G4.0.4 has the property W$n$ for $\mathrm{par}^{0}$.
\item[(C)] The assertion (B) holds with $\mathrm{par}$ in the place of $\mathrm{par}^{0}$.
\end{itemize}

\paragraph{1.2. The constants.}
The following facts are consequences of the definitions of the constants in Proposition~26.5 and in the lines (C0), (C4), (C4r), (C5), (C6) of Section~I.8.2.
The constants $\alpha^{P}$ of the line (C4) and $C^{P}$ of the line (C6) have the values $\alpha^{P}=\frac{1}{4}$ and $C^{P}=116\,C_{P}\cdot M^{402\alpha_0}$ given in Proposition~14.2\,(a);
by the line (C5) and the condition (M1) of Proposition~14.2\,(b), $M$ is a power of $L$ with $M\ge L^2$, hence $M\ge1$.
\begin{itemize}
\item[(cs1)] $\alpha_0=\min\left(\alpha^{F},\alpha^{P}\right)=\min\left(\frac{1}{54},\frac{1}{4}\right)=\frac{1}{54}$.
Hence $0<\alpha_0\le\alpha^{P}=\frac{1}{4}$.
\item[(cs2)] One has $\gamma^{F}>0$ (Proposition~14.1), and $\gamma^{P}=\frac{\gamma_{P}}{2}$ with $\gamma_{P}\in(0,1]$:
the number $\gamma^{P}$ is defined in \cite[Section~15.10.3]{QEDI}, the number $\gamma_{P}$ in \cite[Section~15.7]{QEDI},
and the range of $\gamma_{P}$ is established in the paragraph on the constants $C_{P}$, $\gamma_{P}$, $c_P$ at the end of \cite[Section~15.7]{QEDI}.
Hence $\gamma_{\mathrm{w}}=\min\left(\gamma^{F},\gamma^{P}\right)$ satisfies
\[
0<\gamma_{\mathrm{w}},\qquad\gamma_{\mathrm{w}}\le\gamma^{F},\qquad\gamma_{\mathrm{w}}\le\gamma^{P}.
\]
\item[(cs3)] $C_{w}^{(0)}=\max\left(C^{F},C^{P}\right)$ and $C_{w}^{(1)}=\max\left(C_{w}^{(0)},4\exp(2\gamma_{\mathrm{w}})\cdot M^{27\alpha_0}\right)$.
Hence
\[
C^{F}\le C_{w}^{(1)},\qquad C^{P}\le C_{w}^{(1)},\qquad 4\exp(2\gamma_{\mathrm{w}})\cdot M^{27\alpha_0}\le C_{w}^{(1)}.
\]
Moreover $C^{P}=116\,C_{P}\cdot M^{402\alpha_0}>0$, because $C_{P}\ge4\exp(2)$ by the definition of $C_{P}$ in \cite[Section~15.7]{QEDI} and $M\ge1$;
and
\[
C^{F}\ge2C_{\mathrm{By}}=2(1+2\mathrm{e})\sqrt{2}\cdot\max\left(1,b_{-}^{-1}\right)>0
\]
by the definition of $C^{F}$ in Proposition~14.1 and the definition of $C_{\mathrm{By}}$ in Proposition~14.2\,(a).
\item[(cs4)] $\kappa_{0}=54\log72>0$ and $K_{0}=\frac{\exp(\kappa_0)}{36}=\frac{72^{54}}{36}=2\cdot72^{53}\ge1$.
Hence $\kappa_{\mathrm{w}}\le\kappa_{\mathrm{w}}+\kappa_0$ and $C_{w}^{(1)}\le64K_{0}\cdot C_{w}^{(1)}=C_{w}$.
\item[(cs5)] $M\ge1$, $0<8\alpha_0\le27\alpha_0$ and $\exp(2\gamma_{\mathrm{w}})\ge1$.
Hence $M^{8\alpha_0}\le M^{27\alpha_0}\le4\exp(2\gamma_{\mathrm{w}})\cdot M^{27\alpha_0}$.
\end{itemize}

\paragraph{1.3. Change of the set of parameters.}

\begin{named}{Lemma A}[Change of the set of parameters]
Let $\gamma_{\mathrm{new}}\le\gamma_{\mathrm{old}}$, $\mathrm{kq}_{\mathrm{new}}\le\mathrm{kq}_{\mathrm{old}}$ and $0<C_{\mathrm{old}}\le C_{\mathrm{new}}$.
If a family has one of the properties W1--W16 for $\left(\alpha_0,\gamma_{\mathrm{old}},\mathrm{kq}_{\mathrm{old}};C_{\mathrm{old}}\right)$,
then it has this property for $\left(\alpha_0,\gamma_{\mathrm{new}},\mathrm{kq}_{\mathrm{new}};C_{\mathrm{new}}\right)$.
\end{named}

\begin{proof}
In the properties W1, W2 and W8--W16 of Section~G4.0.4 the set of parameters enters only through the radius $M^{\alpha_0}$ of $\mathrm{Pd}_k$, and this radius is the same for both sets.
Each of the properties W3--W7 is a family of inequalities whose left sides do not contain the set of parameters;
their right sides are
\begin{align*}
\text{W3, W4:}&\quad C\cdot\exp\left(-\gamma\cdot d(t,t')\right);\\
\text{W5:}&\quad C\cdot|v|_\infty;\\
\text{W6:}&\quad C\cdot\exp\left(-\gamma\cdot d(t,t')\right),\qquad C\cdot\exp\left(-\gamma\cdot d(z,z')-\mathrm{kq}\cdot d_M(X)\right);\\
\text{W7:}&\quad C\cdot\exp\left(-\gamma\cdot d(z,z')-\mathrm{kq}\cdot d_M(X)\right)\cdot|f|_\infty\quad(k\ge1),\\
&\quad C\cdot(1+y)^{-1}\cdot\exp\left(-\gamma\cdot d(x,x')-\mathrm{kq}\cdot d_M(X)\right)\quad(k=0).
\end{align*}
The function $d(\cdot,\cdot)\ge0$ is a distance, and $d_M(X)\ge0$, since $M\cdot d_M(X)$ is an infimum of lengths (Section~I.1.1; \cite[Section~5.2.3]{QEDI}).
Therefore, for all values $d\ge0$ and $d_M\ge0$,
\[
\gamma_{\mathrm{old}}\cdot d+\mathrm{kq}_{\mathrm{old}}\cdot d_M\ge\gamma_{\mathrm{new}}\cdot d+\mathrm{kq}_{\mathrm{new}}\cdot d_M ,
\]
and consequently
\[
C_{\mathrm{old}}\cdot\exp\left(-\gamma_{\mathrm{old}}\cdot d-\mathrm{kq}_{\mathrm{old}}\cdot d_M\right)
\le C_{\mathrm{old}}\cdot\exp\left(-\gamma_{\mathrm{new}}\cdot d-\mathrm{kq}_{\mathrm{new}}\cdot d_M\right)
\le C_{\mathrm{new}}\cdot\exp\left(-\gamma_{\mathrm{new}}\cdot d-\mathrm{kq}_{\mathrm{new}}\cdot d_M\right).
\]
So every right side formed with the old set of parameters is at most the corresponding right side formed with the new set.
\end{proof}

The property W17 does not fall under Lemma A:
its constant $4\exp(2\gamma)\cdot M^{27\alpha_0}$ increases with $\gamma$.
It is dealt with in 1.7.

\paragraph{1.4. The photon families.}
\emph{The two families.}
Let $H_k(b,b')$ and $C_k^{1/2}(\Upsilon,\Upsilon')$ be the kernels of the exact operators $H_k$ and $C_k^{1/2}$ of \cite[Section~5.2.5]{QEDI} (Section~I.5.7).
Let $\chi_{\mathrm{pc}}$, $\mathrm{pc}\in\pi'_k$, be the partition of unity on the fine bonds of \cite[Section~15.10.2]{QEDI};
thus $\chi_{\mathrm{pc}}(b)\ge0$ and $\sum_{\mathrm{pc}}\chi_{\mathrm{pc}}(b)=1$.
Let $\mathrm{Pt}'(\mathrm{pc},b')$, a finite set of cubes of $\pi'_k$, and $\mathrm{Pt}''(\Upsilon,\Upsilon')$, a finite set of cubes of $\pi''_k$, be the sets of \cite[Section~15.10.3]{QEDI}.
The families are defined by
\begin{equation}
H_k(s')(b,b'):=H_k(b,b')\cdot\sum_{\mathrm{pc}\in\pi'_k}\chi_{\mathrm{pc}}(b)\cdot\prod_{\mathrm{pc}'\in\mathrm{Pt}'(\mathrm{pc},b')}s'_{\mathrm{pc}'}\tag{10.2}
\end{equation}
for $k\ge1$, $s'=(s'_{\mathrm{pc}})_{\mathrm{pc}\in\pi'_k}$, a fine bond $b$ and a unit bond $b'$, and by
\begin{equation}
C_k^{1/2}(s')(\Upsilon,\Upsilon'):=C_k^{1/2}(\Upsilon,\Upsilon')\cdot\prod_{\mathrm{pc}\in\mathrm{Pt}''(\Upsilon,\Upsilon')}s'_{\mathrm{pc}}\tag{10.3}
\end{equation}
for $k\ge0$ and $s'=(s'_{\mathrm{pc}})_{\mathrm{pc}\in\pi''_k}$.
At level $0$ the family $H_0(s')$ is the diagonal family of Section~I.5.8, which is the family of \cite[Proposition~15.36]{QEDI}.
All entries are polynomials in $s'$ with real coefficients, and none of them contains $A$, $\mu$ or $e_k$.

\emph{Part (A) for these families.}
Let $s'\in\mathrm{Pd}_k$.
With the convention
\[
(OA)(b)=\sum_{b'}O(b,b')\cdot A(b'),\qquad (OA)(c',c):=-(OA)(c,c'),
\]
where $b$ is a positively oriented fine bond and $b'$ ranges over the positively oriented unit bonds,
the kernel (10.2) defines a linear operator from the bond functions on $T^0_k$ to the bond functions on $T_k$.
The matrix (10.3) in the basis $\{\Upsilon\}$ defines a linear operator on $V_{k}$ for real $s'$, and on the complexification of $V_{k}$ for complex $s'$;
on the complexification the element $\tilde{g}$ acts complex-linearly, and the scalar product of $V_{k}$ is extended complex-bilinearly.
Here $V_{k}$ is the space of the real bond functions $\mathsf{Z}$ on $T^0_k$ with $\tau\mathsf{Z}=0$ that vanish on the central bonds, with the counting scalar product (\cite[Section~5.2.5]{QEDI}),
and its complexification consists of the complex bond functions with these two properties.
The element $\tilde{g}$ acts on them by $(\tilde{g}A)(c,c')=A\left(\tilde{g}^{-1}c,\tilde{g}^{-1}c'\right)$ (Section~I.4.3);
it maps $V_{k}$ onto itself (\cite[Proposition~15.21(ii)]{QEDI}), and hence also the complexification.
Consequently, for every $s'\in\mathrm{Pd}_k$, real or not, the relation W14 for (p2),
\[
C_k^{1/2}(\tilde{g}s')\cdot\tilde{g}=\tilde{g}\cdot C_k^{1/2}(s'),
\]
is a relation between complex-linear operators;
and in the proof of \cite[Theorem~15.35]{QEDI}, where this relation is obtained from the matrix $\left\langle\tilde{g}^{-1}\Upsilon,O\,\tilde{g}^{-1}\Upsilon'\right\rangle$ of $\tilde{g}O\tilde{g}^{-1}$, the bracket is the bilinear one.
At $k=0$ the fine bonds are the unit bonds, and $H_0(s')$ is a diagonal operator on the bond functions on $T^0_0$.
These are the spaces and the parameters of (p1) and (p2) in Section~G4.0.4.

\emph{The hypotheses of \cite[Theorem~15.35]{QEDI} and \cite[Proposition~15.36]{QEDI}.}
These two statements assume:
$0\le\alpha<1$;
$L$ odd and $L\ge3$;
$0<\alpha_0\le\frac{1}{4}$;
$M$ a power of $L$ with $M\ge L^2$ and $M\ge M^{P}(\alpha,L,\alpha_0)$;
$\ell\in\mathrm{Cl}$;
$k\ge1$ for (10.2), $k\ge0$ for (10.3), $k=0$ for $H_0(s')$;
and $\mathrm{Pd}_k$ is the polydisk of radius $M^{\alpha_0}$ formed with the $\alpha_0$ of these data.
In the present setting $\alpha\in(1-\bar{\varepsilon},1)$ by the line (C1) of Section~I.8.2;
$L$ is odd and $L\ge5$ by the condition (L5) of Proposition~14.2\,(b);
$\alpha_0=\frac{1}{54}$ by (cs1);
$M$ is a power of $L$ with $M\ge L^2$ by (M1);
and $M\ge M^{P}$ by (M6), where $M^{P}$ (Proposition~14.2\,(a)) is the number of \cite[Section~15.10.3]{QEDI} at $\alpha_0=\frac{1}{54}$,
that is, the least power $M_0\ge L^2$ of $L$ such that
\[
2\alpha_0\cdot402\cdot\log t\le\gamma_{P}\cdot t\qquad\text{for all real }t\ge M_0 .
\]
Since $t\mapsto\frac{\log t}{t}$ is decreasing on $[\mathrm{e},\infty)$ and $L^2\ge\mathrm{e}$,
this condition is equivalent to the single inequality
\[
2\alpha_0\cdot402\cdot\log M_0\le\gamma_{P}\cdot M_0 ,
\]
as is noted in \cite[Section~15.10.3]{QEDI} after the definition of $M^{P}$.
Finally $\ell\in\mathrm{Cl}$ by (SF).
The numbers $\gamma^{P}=\frac{\gamma_{P}}{2}$ and $C^{P}=116\,C_{P}\cdot M^{402\alpha_0}$ of Proposition~14.2\,(a) are those of \cite[Section~15.10.3]{QEDI}.
Hence both statements apply.
In them the directions are $0,1,2$ (here $1,2,3$; Section~I.1.1), a fine site is written $x$ (here $c$), a parameter cube is denoted by its centre $c$ (here $\mathrm{pc}$), and an element of $\mathbb{G}_k$ is written $g$ (here $\tilde{g}$).

\emph{The properties supplied by \cite[Theorem~15.35]{QEDI} and \cite[Proposition~15.36]{QEDI}.}
\begin{itemize}
\item[(tp1)] (\cite[Theorem~15.35]{QEDI}.)
For $k\ge1$ the family (10.2) has the properties W1, W2, W3, W8, W9, W10 and W16, and the property W14 for all $\tilde{g}\in\mathbb{G}_k$, for $\left(\alpha_0,\gamma^{P},\cdot\,;C^{P}\right)$.
\item[(tp2)] (\cite[Theorem~15.35]{QEDI}.)
For $k\ge0$ the family (10.3) has the properties W1, W2, W4, W8, W9, W10, W14 (for $\tilde{g}\in\mathbb{G}_k^{L}$) and W16 for $\left(\alpha_0,\gamma^{P},\cdot\,;C^{P}\right)$.
\item[(tp3)] (\cite[Proposition~15.36]{QEDI}.)
The family $H_0(s')$ has the properties W1, W2, W8, W9, W10 and W16, the property W14 for all $\tilde{g}\in\mathbb{G}_0$,
and, for every rate $\gamma>0$, the property W3 with the constant $4\exp(2\gamma)\cdot M^{27\alpha_0}$.
\end{itemize}

\emph{Comparison of the two definitions.}
The properties of systems of primitives in \cite[Section~5.4.2]{QEDI} form a numbered list of eighteen entries;
in (tp1)--(tp3), W$n$ denotes the $n$-th entry of that list.
For the families (p1) and (p2), the properties W1, W2, W8, W10, W14 and W16 of Section~G4.0.4 are these entries of \cite[Section~5.4.2]{QEDI},
with $t$ in the place of $x$, $\mathrm{pc}$ in the place of $c$ and $\tilde{g}$ in the place of $g$;
since the two families do not depend on $A$, the clauses on $A$ and $\tilde{A}$ are void.
In W1 the exact operators $H_k$ and $C_k^{1/2}$ on the right-hand side are the photon operators of \cite[Section~5.2.5]{QEDI} (Section~I.5.7), with $H_0=1$, as in the entry W1 of \cite[Section~5.4.2]{QEDI}.
The groups $\mathbb{G}_k$, $\mathbb{G}_k^{L}$ and the actions of $\tilde{g}=(S,a)$ on photon sites, cells, bond functions ($V_{k}$ is a space of bond functions on $T^0_k$) and parameter families are those of \cite[Section~5.2.4]{QEDI}:
\[
\tilde{g}c=r(S)c+a,\qquad(\tilde{g}A)(c,c')=A\left(\tilde{g}^{-1}c,\tilde{g}^{-1}c'\right),\qquad(\tilde{g}s)_{\mathrm{pc}}=s_{\tilde{g}^{-1}\mathrm{pc}}
\]
(Sections~I.4.2 and~I.4.3).
In W9, Section~G4.0.4 requires, for $\mathscr{Y}$, $Y$, $s$, $Y_1$, $\tilde{s}$ as in W8 and a column index $t'$ contained in $Y_1$,
that $\mathrm{Op}(s)(t,t')=0$ unless $t$ is contained in $Y_1$, and that $\mathrm{Op}(s)(\cdot,t')=\mathrm{Op}(\tilde{s})(\cdot,t')$;
in \cite[Section~5.4.2]{QEDI} the last equation is expressed by saying that the column at $t'$ depends on $s$ only through its values on the cubes of $Y_1$, $s$ ranging over the families in $\mathrm{Pd}_k$ that vanish outside $\mathscr{Y}$.
The two requirements are the same.
The properties W3 and W4 of Section~G4.0.4 are the entries W3 and W4 of \cite[Section~5.4.2]{QEDI},
read with the definition of the kernels $O(b,b')$ and of the differences $\partial_{\mathfrak{l}}O$ and $\delta_\alpha\partial_{\mathfrak{l}}O$ given there,
with $c$, $\tilde{c}$ in the place of the initial points $x$, $\tilde{x}$ of $b$, $\tilde{b}$, and with $\mathfrak{l}=1,2,3$ in the place of $\mathfrak{l}=0,1,2$.
Since $\mathbb{G}_k^{L}$ is a subgroup of $\mathbb{G}_k$ (Section~I.4.2), (tp1) and (tp3) contain W14 in the form of Section~G4.0.4, that is, for $\tilde{g}\in\mathbb{G}_k^{L}$.

\paragraph{1.5. The fermionic families.}
Here $\kappa'=\kappa_{\mathrm{w}}+\kappa_0$ (Section~G1.1; line (C4) of Section~I.8.2).
The formulas that define the families (Section~G1.9, Propositions~7 and~8) do not contain $\kappa'$.
The number $\kappa'$ enters in three ways:
through the number $M^{F}$ (defined in the line (C5) of Section~I.8.2 and in Proposition~14.2\,(a)), which enters through the condition (M5);
through the constants $C_{1}$ and $C^{F}$ of the line (C6), which are formed with this value of $\kappa'$;
and as the rate per unit of $d_M$ in W6 and W7.

\emph{The properties supplied by Propositions~7, 8 and~14.1.}
\begin{itemize}
\item[(tf1)] (Proposition~14.1\,(ii), the assertions that the families of Proposition~7 are the fermionic part of a system of primitives and have the properties W1, W2, W5--W17 for $\left(\alpha_0,\gamma^{F},\kappa_{\mathrm{w}}+\kappa_0;C^{F}\right)$; Proposition~7.)
For $k\ge1$ the families (p3) $H_k(s'',A)$, $\bar{H}_k(s'',A)$, (p4) $\Gamma_k(s'',A)$, $M_k(X,A)$ and (p5) $S_{k,y}(X,A)$ of Proposition~7 (where $s$ stands for $s''$),
defined for $s''=(s''_{\mathrm{pc}})_{\mathrm{pc}\in\pi_k}$ in $\mathrm{Pd}_k$, $A\in\tilde{G}_k$, $X\in D_k$ and $y\ge0$,
are the fermionic part of a system of primitives of level $k$ on $T_k$ (Section~G1.11),
and they have the properties W1, W2, W5--W16 of Section~G1.11 for $\left(\alpha_0,\gamma^{F},\kappa_{\mathrm{w}}+\kappa_0;C^{F}\right)$.
The property W14 holds for all $\tilde{g}\in\mathbb{G}_k$, and for (p5) for all $\tilde{g}\in\mathbb{G}_{k,3}$.
The map $y\mapsto S_{k,y}(X,A)$ is continuous on $[0,\infty)$.
\item[(tf2)] (Proposition~14.1\,(ii), the assertions that the families of Proposition~8 are the fermionic part of a system of primitives and have the properties W1, W2, W5--W17 for $\left(\alpha_0,\gamma^{F},\kappa_{\mathrm{w}}+\kappa_0;C^{F}\right)$; Proposition~8\,(a).)
For $k=0$ the families $\Gamma_0(s'',A)$ and $S_{0,y}(X,A)$ of Proposition~8, which act on $\mathfrak{S}^{\mathrm{ch}}(1,\ell)$,
have the properties W1, W2, W6--W9, W12--W16 of Section~G1.11, each in so far as it concerns the family, for $\left(\alpha_0,\gamma^{F},\kappa_{\mathrm{w}}+\kappa_0;C^{F}\right)$;
here W6 and W7 are bounds of the kernels (I.1.4) at pairs of fermion sites, W7 with the factor $(1+y)^{-1}$.
The property W14 holds for all $\tilde{g}\in\mathbb{G}_0$, and for $S_{0,y}(X,A)$ for all $\tilde{g}\in\mathbb{G}_{0,3}$.
The map $y\mapsto S_{0,y}(X,A)$ is continuous on $[0,\infty)$.
\item[(tf3)] (Proposition~8\,(b).)
For $k=0$ the family $H_0(s'',A)$ of Proposition~8, which is the diagonal family of Section~I.5.8, together with $\bar{H}_0(s'',A)=H_0(s'',A)$,
has the properties W1, W2, W8--W11, W13 and W16, the property W14 for all $\tilde{g}\in\mathbb{G}_0$, and the property W5 with the constant $M^{8\alpha_0}$.
\end{itemize}

\emph{The properties of Section~G1.11 as properties of Section~G4.0.4.}
Restricted to the families (p3)--(p5), the properties W1, W2, W5--W16 of Section~G4.0.4 are the properties W1, W2, W5--W16 of Section~G1.11, with $s''$ in the place of $s$;
here (p3), (p4), (p5) and $\mathrm{Op}$ stand for the families that (tf1)--(tf3) denote by their letters.
The equations coincide, and so do
the index sets (unit sites $z$; fermion sites $x=(c,g)$ with $\mathrm{supp}\,x=\{c\}$ and $K_x=K_c$),
the kernels ((I.1.4) for a source space of spinors on fermion sites, and $(Ov)(t)=\sum_{z'}O(t,z')\,v(z')$ for the source space $\mathfrak{S}(1,\ell)$),
the norms ($|\cdot|_c$ on $\mathbb{C}^4$ and on $4\times4$ blocks, and $|\cdot|_\infty$, of Section~I.1.5)
and the groups in W14.
For the individual properties the identification rests on the following facts.
\begin{description}
\item[W1] The clause on $M_k$ concerns the levels $k\ge1$ only, since a system of level $0$ contains no family $M_0(X,A)$;
the exact operators on the right-hand sides are those of Section~I.3.5 for $k\ge1$ and those of Section~I.3.9 for $k=0$.
\item[W5] A bound by $C\cdot|v|_\infty$ holds for the supremum over the short pairs if and only if it holds at every short pair.
\item[W6] The indices are unit sites for $k\ge1$;
for $k=0$ they are fermion sites, and the kernels are those of (I.1.4).
\item[W7] The operator $\mathbf{1}_{\Delta_z}$ is the multiplication by the indicator of $\Delta_z\times O_{\mathrm{cub}}$, by the rule (I.POS) of Section~I.1.4.
The clause for $k=0$ is required for every $y\ge0$ and concerns the kernels (I.1.4):
the resolvent parameter $y$ of a system ranges over $[0,\infty)$, and at level $0$ the family $S_{0,y}(X,A)$ is indexed as (p5).
\item[W8] The set $\mathscr{Y}$ is a set of parameter cubes of the kind of $\mathrm{Op}$;
since the kind of (p3) and of $\Gamma_k$ is $\pi_k$ (Section~G4.0.4), $\mathscr{Y}$ is a subset of $\pi_k$.
\item[W9] The property W9 is the property W8 with rows and columns exchanged.
\item[W12] Both indices are required to lie over $X$.
A fermion site $x=(c,g)$ lies in $X$, or over $X$, if and only if $c$ lies in $X$ (I.POS);
for $M_k$ the indices are unit sites, and to lie over $X$ means to lie in $X$.
\item[W13] At $k=0$ one has $\boldsymbol{\lambda}^{\mathrm{u}}=\boldsymbol{\lambda}$, and the relation for $\Gamma_0$ is the relation of W13 with this $\boldsymbol{\lambda}^{\mathrm{u}}$.
The relation is required for all real functions $\boldsymbol{\lambda}$.
\item[W14] The actions of $\tilde{g}$ are those defined in Section~I.4.3, with (I.4.2);
the groups for which the relation is required are those named in Section~I.5.8.
\item[W15] The pairings are those of Section~I.1.6, $\langle\cdot,\cdot\rangle_1$ for $k\ge1$ and $\langle\cdot,\cdot\rangle^{\mathrm{ch}}$ for $k=0$;
with them the form $\langle\bar{\Psi},O\Psi\rangle$ of Section~I.3.6 is formed.
\item[W16] The family $M_k(X,A)$ has no parameter $s$, and for it the relation holds in the form without $s$.
\item[$\bar{H}_k$] The relations for the barred family are consequences of its definition $\bar{H}_k(s'',A)=\mathsf{C}^{-1}H_k(s'',-A)\mathsf{C}$ and impose no further condition.
\item[W17] The clauses of the property W17 of Section~G4.0.4 on $H_0(s'',A)$, $\Gamma_0$ and $S_{0,y}$ are the property W17 of Section~G1.11.
Each property is required of a family only in so far as it concerns that family, and for $k=0$ the properties W6 and W7 have their form for $k=0$ (for W6 see above).
\end{description}
In W10 the row index $t$ of (p3) is the fermion site $x$.
Since $\mathbb{G}_k^{L}$ and $\mathbb{G}^{L}_{k,3}$ are subgroups of $\mathbb{G}_k$ and $\mathbb{G}_{k,3}$ (Section~I.4.2),
(tf1)--(tf3) contain W14 in the form of Section~G4.0.4.

\emph{Part (A) for these families.}
The fermionic part of a system of level $k$ in Section~G1.11 consists of the families (p3), (p4), (p5) of Section~G4.0.4 if $k\ge1$,
and of $\Gamma_0$, $S_{0,y}$ and the diagonal family $H_0(s'',A)$ of Section~I.5.8, without a family $M_0$, if $k=0$;
the spaces and the parameters are the same.
A function that is continuous in $y$ is Borel measurable in $y$.

\paragraph{1.6. The levels $k\ge1$.}
Part (A) holds by 1.4 and 1.5:
the families (10.2), (10.3) and the families of (tf1) form a system of primitives of level $k$ on $T_k$.
For part (B), the property W17 concerns level $0$ and is void for $k\ge1$.
For $1\le n\le16$ the following list gives, for every family that Section~G4.0.4 names at the end of the property W$n$, the statements from which W$n$ is obtained;
a family that is not named there for W$n$ does not occur under W$n$.
\begin{description}
\item[W1, W2] for (p1) $H_k(s')$: (tp1); for (p2) $C_k^{1/2}(s')$: (tp2); for (p3) $H_k(s'',A)$, for $\Gamma_k(s'',A)$, for $M_k(X,A)$ and for (p5) $S_{k,y}(X,A)$: (tf1).
\item[W3] for (p1) $H_k(s')$: (tp1) and (26.5.1).
\item[W4] for (p2) $C_k^{1/2}(s')$: (tp2) and (26.5.2).
\item[W5] for (p3) $H_k(s'',A)$: (tf1) and (26.5.3).
\item[W6] for $\Gamma_k(s'',A)$ and for $M_k(X,A)$: (tf1) and (26.5.4).
\item[W7] for (p5) $S_{k,y}(X,A)$: (tf1) and (26.5.5).
\item[W8, W9] for (p1) $H_k(s')$: (tp1); for (p2) $C_k^{1/2}(s')$: (tp2); for (p3) $H_k(s'',A)$ and for $\Gamma_k(s'',A)$: (tf1).
\item[W10] for (p1) $H_k(s')$: (tp1); for (p2) $C_k^{1/2}(s')$: (tp2); for (p3) $H_k(s'',A)$: (tf1).
\item[W11] for (p3) $H_k(s'',A)$: (tf1).
\item[W12] for $M_k(X,A)$ and for (p5) $S_{k,y}(X,A)$: (tf1).
\item[W13] for (p3) $H_k(s'',A)$, for $\Gamma_k(s'',A)$, for $M_k(X,A)$ and for (p5) $S_{k,y}(X,A)$: (tf1).
\item[W14] for (p1) $H_k(s')$: (tp1); for (p2) $C_k^{1/2}(s')$: (tp2); for (p3) $H_k(s'',A)$, for $\Gamma_k(s'',A)$ and for $M_k(X,A)$: (tf1); for (p5) $S_{k,y}(X,A)$: (tf1), with $\tilde{g}\in\mathbb{G}^{L}_{k,3}$.
\item[W15] for $\Gamma_k(s'',A)$ and for $M_k(X,A)$: (tf1).
\item[W16] for (p1) $H_k(s')$: (tp1); for (p2) $C_k^{1/2}(s')$: (tp2); for (p3) $H_k(s'',A)$, for $\Gamma_k(s'',A)$ and for $M_k(X,A)$: (tf1).
\end{description}
The statements (tp1), (tp2) give the properties for $\left(\alpha_0,\gamma^{P},\cdot\,;C^{P}\right)$, and (tf1) gives them for $\left(\alpha_0,\gamma^{F},\kappa_{\mathrm{w}}+\kappa_0;C^{F}\right)$.
By (cs2), (cs3) and Lemma A, applied with $\left(\gamma_{\mathrm{old}},C_{\mathrm{old}}\right)=\left(\gamma^{P},C^{P}\right)$, resp.\ $\left(\gamma_{\mathrm{old}},\mathrm{kq}_{\mathrm{old}},C_{\mathrm{old}}\right)=\left(\gamma^{F},\kappa_{\mathrm{w}}+\kappa_0,C^{F}\right)$,
and with $\left(\gamma_{\mathrm{new}},\mathrm{kq}_{\mathrm{new}},C_{\mathrm{new}}\right)=\left(\gamma_{\mathrm{w}},\kappa_{\mathrm{w}}+\kappa_0,C_{w}^{(1)}\right)$,
these properties hold for $\mathrm{par}^{0}$.
Explicitly, for all $s',s''\in\mathrm{Pd}_k$, $A\in\tilde{G}_k$, $X\in D_k$ and $y\ge0$,
for $b$, $\tilde{b}$, $b'$, $\Upsilon$, $\Upsilon'$, $\mathfrak{l}$ as in the properties W3 and W4 of Section~G4.0.4,
for unit sites $z$, $z'$, and for $v\in\mathfrak{S}(1,\ell)$ and $f\in\mathfrak{S}^{\mathrm{ch}}(\eta_k,\ell)$, one has
\begin{equation}
\begin{split}
&\left|H_k(s')(b,b')\right|,\quad\left|\left(\partial_{\mathfrak{l}}H_k(s')\right)(b,b')\right|,\quad\left|\left(\delta_\alpha\partial_{\mathfrak{l}}H_k(s')\right)\left((b,\tilde{b});b'\right)\right|\\
&\qquad\le C^{P}\cdot\exp\left(-\gamma^{P}\cdot d(b,b')\right)\le C_{w}^{(1)}\cdot\exp\left(-\gamma_{\mathrm{w}}\cdot d(b,b')\right);
\end{split}\tag{26.5.1}
\end{equation}
\begin{equation}
\left|C_k^{1/2}(s')(\Upsilon,\Upsilon')\right|\le C^{P}\cdot\exp\left(-\gamma^{P}\cdot d(\Upsilon,\Upsilon')\right)\le C_{w}^{(1)}\cdot\exp\left(-\gamma_{\mathrm{w}}\cdot d(\Upsilon,\Upsilon')\right);\tag{26.5.2}
\end{equation}
\begin{equation}
\begin{split}
&\left|H_k(s'',A)\,v\right|_\infty\le C^{F}\cdot|v|_\infty\le C_{w}^{(1)}\cdot|v|_\infty,\\
&\left|\left(\delta_{\alpha,A}H_k(s'',A)\,v\right)(x,x')\right|_c\le C^{F}\cdot|v|_\infty\le C_{w}^{(1)}\cdot|v|_\infty
\quad\text{for all }(x,x')\in\mathrm{Pzt}_k\text{ with }d_X(x,x')<1;
\end{split}\tag{26.5.3}
\end{equation}
\begin{equation}
\begin{split}
\left|\Gamma_k(s'',A)(z,z')\right|_c&\le C^{F}\cdot\exp\left(-\gamma^{F}\cdot d(z,z')\right)\le C_{w}^{(1)}\cdot\exp\left(-\gamma_{\mathrm{w}}\cdot d(z,z')\right),\\
\left|M_k(X,A)(z,z')\right|_c&\le C^{F}\cdot\exp\left(-\gamma^{F}\cdot d(z,z')-(\kappa_{\mathrm{w}}+\kappa_0)\cdot d_M(X)\right)\\
&\le C_{w}^{(1)}\cdot\exp\left(-\gamma_{\mathrm{w}}\cdot d(z,z')-(\kappa_{\mathrm{w}}+\kappa_0)\cdot d_M(X)\right);
\end{split}\tag{26.5.4}
\end{equation}
\begin{equation}
\begin{split}
\left|\mathbf{1}_{\Delta_z}S_{k,y}(X,A)\,\mathbf{1}_{\Delta_{z'}}f\right|_\infty
&\le C^{F}\cdot\exp\left(-\gamma^{F}\cdot d(z,z')-(\kappa_{\mathrm{w}}+\kappa_0)\cdot d_M(X)\right)\cdot|f|_\infty\\
&\le C_{w}^{(1)}\cdot\exp\left(-\gamma_{\mathrm{w}}\cdot d(z,z')-(\kappa_{\mathrm{w}}+\kappa_0)\cdot d_M(X)\right)\cdot|f|_\infty .
\end{split}\tag{26.5.5}
\end{equation}
In (26.5.3) the pairs with $d_X(x,x')=\max\left(|c-c'|,\eta_k\right)<1$, the short pairs of Section~I.1.7, include, over every photon site $c$,
the $48\cdot47$ ordered pairs $x=(c,g)$, $x'=(c,g')$ with $g\ne g'$ in $O_{\mathrm{cub}}$;
for them $d_X=\eta_k<1$, because $k\ge1$.
At such a pair the quotient is
\[
\eta_k^{-\alpha}\left(\left(H_k(s'',A)\,v\right)(c,g')-\left(H_k(s'',A)\,v\right)(c,g)\right),
\]
without a phase ((I.1.7)).
In (26.5.4) the kernels are those of $(Ov)(z)=\sum_{z'}O(z,z')\,v(z')$, that is, $4\times4$ blocks at pairs of unit sites.
In (26.5.5) the operator $\mathbf{1}_{\Delta_z}$ is the multiplication by the indicator of $\Delta_z\times O_{\mathrm{cub}}$,
a set of $48\eta_k^{-3}$ fermion sites of measure
\[
\int_{\Delta_z}\mathrm{d}x=\eta_k^{3}\cdot\frac{1}{48}\cdot48\eta_k^{-3}=1
\]
(Section~I.1.4 and (I.1.3)).

\paragraph{1.7. Level $0$.}
At level $0$ one has $\eta_0=1$; the fine lattice and the unit lattice coincide, $T^{\mathrm{f}}_0=T^0_0\times O_{\mathrm{cub}}$,
and the spinors are the elements of $\mathfrak{S}^{\mathrm{ch}}(1,\ell)$, with $192$ components per photon site (Section~I.3.9 and (I.L0.1)).

Part (A): by 1.4 and 1.5, the family (10.3) at $k=0$, the families $\Gamma_0(s'',A)$ and $S_{0,y}(X,A)$ on $\mathfrak{S}^{\mathrm{ch}}(1,\ell)$, the second Borel measurable in $y$,
and the two diagonal families $H_0(s')$, $H_0(s'',A)$ of Section~I.5.8 form what Section~G4.0.4 calls a system of level $0$;
there is no family $M_0$.

Part (B): at level $0$ the properties W1--W16 are asserted in the form of the property W17.
The property W17 states that the diagonal families have W1, W2, W8--W11, W13, W14, W16, and W3, W5 with the constant $4\exp(2\gamma)\cdot M^{27\alpha_0}$, where $\gamma$ is the rate of the set of parameters;
and that (p2), $\Gamma_0$ and $S_{0,y}$ have W1, W2, W4, W6--W9, W10 (for (p2)) and W12--W16, with W7 in the form for $k=0$;
each property W$n$ is required of a family only in so far as Section~G4.0.4 names that family at the end of W$n$, (p1) standing for $H_0(s')$ and (p3) for $H_0(s'',A)$.
The following list gives the statements from which each property is obtained.
\begin{description}
\item[W1, W2] for $H_0(s')$: (tp3); for (p2) $C_0^{1/2}(s')$: (tp2); for $H_0(s'',A)$: (tf3); for $\Gamma_0(s'',A)$ and for $S_{0,y}(X,A)$: (tf2).
\item[W3] for $H_0(s')$: (tp3) at $\gamma=\gamma_{\mathrm{w}}$, and (26.5.6).
\item[W4] for (p2) $C_0^{1/2}(s')$: (tp2) and (26.5.7).
\item[W5] for $H_0(s'',A)$: (tf3) and (26.5.8).
\item[W6] for $\Gamma_0(s'',A)$: (tf2) and (26.5.9).
\item[W7] for $S_{0,y}(X,A)$: (tf2) and (26.5.10).
\item[W8, W9] for $H_0(s')$: (tp3); for (p2) $C_0^{1/2}(s')$: (tp2); for $H_0(s'',A)$: (tf3); for $\Gamma_0(s'',A)$: (tf2).
\item[W10] for $H_0(s')$: (tp3); for (p2) $C_0^{1/2}(s')$: (tp2); for $H_0(s'',A)$: (tf3).
\item[W11] for $H_0(s'',A)$: (tf3).
\item[W12] for $S_{0,y}(X,A)$: (tf2).
\item[W13] for $H_0(s'',A)$: (tf3); for $\Gamma_0(s'',A)$ and for $S_{0,y}(X,A)$: (tf2).
\item[W14] for $H_0(s')$: (tp3); for (p2) $C_0^{1/2}(s')$: (tp2); for $H_0(s'',A)$: (tf3); for $\Gamma_0(s'',A)$: (tf2); for $S_{0,y}(X,A)$: (tf2), with $\tilde{g}\in\mathbb{G}^{L}_{0,3}$.
\item[W15] for $\Gamma_0(s'',A)$: (tf2).
\item[W16] for $H_0(s')$: (tp3); for (p2) $C_0^{1/2}(s')$: (tp2); for $H_0(s'',A)$: (tf3); for $\Gamma_0(s'',A)$: (tf2).
\end{description}

\emph{The properties without a constant} (W1, W2, W8--W16) contain, of the set of parameters, only $\alpha_0$;
hence they hold for $\mathrm{par}^{0}$.
At level $0$ they are read as follows.
\begin{description}
\item[W1] $H_0(1)=1=H_0$ on the bond functions, $C_0^{1/2}(1)=C_0^{1/2}$, $H_0(1,A)=1=H_0(A)$ on $\mathfrak{S}^{\mathrm{ch}}(1,\ell)$, $\Gamma_0(1,A)=\Gamma_0(A)=F_{0}(A)^{-1}$, and
\[
\sum_XS_{0,y}(X,A)=S_{0,y}(A)=\left(F_{0}(A)+i\gamma_5\cdot y\right)^{-1},\qquad F_{0}(A)=D_A+\mu+b_{L}\cdot J\,P(A)\,J^{*}
\]
(Section~I.3.9).
\item[W13] $\boldsymbol{\lambda}^{\mathrm{u}}=\boldsymbol{\lambda}$, and $\exp\left(\pm i e_k\cdot\boldsymbol{\lambda}\right)$ multiplies $f(c,g)$ by $\exp\left(\pm i e_k\cdot\boldsymbol{\lambda}(c)\right)$, the same phase on the $48$ sites over $c$.
\item[W14] $\tilde{g}(c,g)=\left(r(S)c+a,\,r(S)g\right)$ ((I.4.2)).
\item[W15] the transpose is taken with respect to $\langle\cdot,\cdot\rangle^{\mathrm{ch}}$, where $\langle f',f\rangle^{\mathrm{ch}}=\frac{1}{48}\sum_c\sum_{g\in O_{\mathrm{cub}}}f'(c,g)\cdot f(c,g)$.
\end{description}

\emph{The bounds.}
For all $s',s''\in\mathrm{Pd}_0$, $A\in\tilde{G}_0$, $X\in D_0$ and $y\ge0$,
for bonds $b$, $\tilde{b}$, $b'$ of $T^0_0$ as in the property W3 of Section~G4.0.4 (here $\eta_0=1$ and $|c-\tilde{c}|=1$), for $\mathfrak{l}=1,2,3$,
for fermion sites $x=(c,g)$, $x'=(c',g')$ of $T^{\mathrm{f}}_0$ and for $f\in\mathfrak{S}^{\mathrm{ch}}(1,\ell)$, one has
\begin{equation}
\begin{split}
&\left|H_0(s')(b,b')\right|,\quad\left|\left(\partial_{\mathfrak{l}}H_0(s')\right)(b,b')\right|,\quad\left|\left(\delta_\alpha\partial_{\mathfrak{l}}H_0(s')\right)\left((b,\tilde{b});b'\right)\right|\\
&\qquad\le4\exp(2\gamma_{\mathrm{w}})\cdot M^{27\alpha_0}\cdot\exp\left(-\gamma_{\mathrm{w}}\cdot d(b,b')\right)\le C_{w}^{(1)}\cdot\exp\left(-\gamma_{\mathrm{w}}\cdot d(b,b')\right);
\end{split}\tag{26.5.6}
\end{equation}
\begin{equation}
\left|C_0^{1/2}(s')(\Upsilon,\Upsilon')\right|\le C^{P}\cdot\exp\left(-\gamma^{P}\cdot d(\Upsilon,\Upsilon')\right)\le C_{w}^{(1)}\cdot\exp\left(-\gamma_{\mathrm{w}}\cdot d(\Upsilon,\Upsilon')\right);\tag{26.5.7}
\end{equation}
\begin{equation}
\begin{split}
&\left(H_0(s'',A)f\right)(c,g)=\prod_{\substack{\mathrm{pc}\in\pi_0\\ \mathrm{pc}\text{ meets }K_c}}s''_{\mathrm{pc}}\cdot f(c,g)\qquad\text{(one factor for the $48$ labels $g$)},\\
&\left|H_0(s'',A)f\right|_\infty\le M^{8\alpha_0}\cdot|f|_\infty\le4\exp(2\gamma_{\mathrm{w}})\cdot M^{27\alpha_0}\cdot|f|_\infty\le C_{w}^{(1)}\cdot|f|_\infty;
\end{split}\tag{26.5.8}
\end{equation}
\begin{equation}
\left|\Gamma_0(s'',A)(x,x')\right|_c\le C^{F}\cdot\exp\left(-\gamma^{F}\cdot d(x,x')\right)\le C_{w}^{(1)}\cdot\exp\left(-\gamma_{\mathrm{w}}\cdot d(x,x')\right);\tag{26.5.9}
\end{equation}
\begin{equation}
\begin{split}
\left|S_{0,y}(X,A)(x,x')\right|_c&\le C^{F}\cdot(1+y)^{-1}\cdot\exp\left(-\gamma^{F}\cdot d(x,x')-(\kappa_{\mathrm{w}}+\kappa_0)\cdot d_M(X)\right)\\
&\le C_{w}^{(1)}\cdot(1+y)^{-1}\cdot\exp\left(-\gamma_{\mathrm{w}}\cdot d(x,x')-(\kappa_{\mathrm{w}}+\kappa_0)\cdot d_M(X)\right).
\end{split}\tag{26.5.10}
\end{equation}
The bound (26.5.6) is (tp3) at the rate $\gamma=\gamma_{\mathrm{w}}$, which is positive by (cs2); its second inequality is (cs3).
The bound (26.5.7) follows from (tp2), (cs2) and (cs3).
The bound (26.5.8) follows from (tf3), (cs5) and (cs3);
the second half of W5 is void, because $d_X(x,x')=\max\left(|c-c'|,\eta_0\right)\ge1$ for all pairs, so that at level $0$ there is no short pair.
Hence the two diagonal families have W3, resp.\ W5, with the constant $4\exp(2\gamma_{\mathrm{w}})\cdot M^{27\alpha_0}$ that the property W17 requires at the rate $\gamma_{\mathrm{w}}$ of $\mathrm{par}^{0}$,
and also with the constant $C_{w}^{(1)}$ of $\mathrm{par}^{0}$.
The bounds (26.5.9) and (26.5.10) follow from (tf2), (cs2), (cs3) and from $d,d_M\ge0$ (Lemma A).
In (26.5.9) and (26.5.10) one has $d(x,x')=d(c,c')$ (I.POS), and the kernels are those of (I.1.4) at $\eta_0=1$:
\[
(Of)(x)=\int\mathrm{d}x'\,O(x,x')\,f(x')=\frac{1}{48}\sum_{c'}\sum_{g'\in O_{\mathrm{cub}}}O\left(x,(c',g')\right)f(c',g'),
\]
with $c'$ ranging over a fundamental domain of $\ell\mathbb{Z}^3$ in $\mathbb{Z}^3$;
they are $4\times4$ blocks at pairs of fermion sites.
The factor $48$ of the inverse weight ((I.ke4)) is contained in $C^{F}$, through $C_{00}$, $a'$ and $C_{1}$ (lines (C4r) and (C6) of Section~I.8.2);
replacing $C^{F}$ by $C_{w}^{(1)}$ introduces no further factor.
This is W17 for $\mathrm{par}^{0}$, and with it W1--W16 at level $0$ in the form in which Section~G4.0.4 states them.

\paragraph{1.8. From $\mathrm{par}^{0}$ to $\mathrm{par}$.}
The sets $\mathrm{par}^{0}$ and $\mathrm{par}$ have the same $\alpha_0$ and the same rate $\gamma_{\mathrm{w}}$,
and $\kappa_{\mathrm{w}}\le\kappa_{\mathrm{w}}+\kappa_0$ and $C_{w}^{(1)}\le C_{w}$ by (cs4).
By Lemma A, applied with $\left(\gamma_{\mathrm{old}},\mathrm{kq}_{\mathrm{old}},C_{\mathrm{old}}\right)=\left(\gamma_{\mathrm{w}},\kappa_{\mathrm{w}}+\kappa_0,C_{w}^{(1)}\right)$
and $\left(\gamma_{\mathrm{new}},\mathrm{kq}_{\mathrm{new}},C_{\mathrm{new}}\right)=\left(\gamma_{\mathrm{w}},\kappa_{\mathrm{w}},C_{w}\right)$,
every property W1--W16 that holds for $\mathrm{par}^{0}$ holds for $\mathrm{par}$.
The constant $4\exp(2\gamma_{\mathrm{w}})\cdot M^{27\alpha_0}$ of the property W17 contains, of the set of parameters, only $\gamma_{\mathrm{w}}$ and $\alpha_0$;
hence it is the same for $\mathrm{par}^{0}$ and for $\mathrm{par}$, and (26.5.6), (26.5.8) are W3, W5 of the diagonal families for $\mathrm{par}$.
The other clauses of W17 are instances of W1, W2, W4, W6--W16.
This proves (C).

\paragraph{1.9. Uniformity.}
The numbers $\alpha_0$, $\gamma_{\mathrm{w}}$, $\kappa_{\mathrm{w}}$, $\kappa_0$, $C^{P}$, $C^{F}$, $C_{w}^{(1)}$ and $C_{w}$ are constants of the lines (C0), (C4), (C4r), (C6) of Section~I.8.2;
their definitions contain only constants that precede them in the list of Proposition~14.2\,(a), and none of $k$, $\ell$, $\mu$, $e_k$.
The statements (tp1)--(tp3) hold for every level $k$ named in them and every $\ell\in\mathrm{Cl}$, and (tf1)--(tf3) for all data (SF), with these numbers.
Hence $\mathrm{par}^{0}$ and $\mathrm{par}$ are the same sets of parameters for all $(k,\ell,\mu,e_k)$,
and (26.5.1)--(26.5.10) hold with constants that do not depend on $k$, $\ell$, $\mu$, $e_k$, $s'$, $s''$, $A$, $X$, $y$.
\end{proof}

\subsection*{2. Proposition 26.6 (the data of the flow are data of one step)}

\begin{proof}[Proof of Proposition~26.6]\leavevmode
\paragraph{2.1. The numbers.}
By Section~I.1.8, (S-4) is the condition (FL) with the primitives left out.
It concerns a quadruple $(\bar{m}, e, R, N)$ such that
\[
\bar{m} \in (0,1], \qquad e > 0, \qquad e_K = e h^{\frac{1}{2}} \le e_1, \qquad R \in \mathrm{Tor}(e),
\]
and $N$ is an integer with $N \ge 0$ and $N \ge N_{\min}(R,e)$.
The primitives that (FL) names (Section~G4.0.1) are objects formed from such a quadruple; their existence is proved in 2.7.
In 2.2--2.6 nothing beyond (S-4) is assumed of the quadruple.

The following definitions are those of Sections~I.0.5, I.6.1 and I.6.2, and, for $w^{\mathrm{m}}_k$ and $w^{\mathrm{E}}_k$, of Section~I.5.4.
The number $h = L^{\lfloor \log_L(1/\bar{m}) \rfloor}$ is a power of $L$ with a non-negative exponent; further $\eta = L^{-N}$ and $K = N + \log_L h$, and for every integer $k \ge 0$
\begin{gather*}
a_k = L^{k-N}, \qquad e_k = e\, a_k^{\frac{1}{2}}, \qquad v_k = -\log e_k, \qquad \mu_k = \bar{m}\, a_k,\\
r_k = \min\left\{ L^i : i \ge 0 \text{ an integer},\ L^i \ge v_k^r \right\}, \qquad p_k = v_k^{p}, \qquad p_{0,k} = v_k^{p_0},\\
w^{\mathrm{m}}_k = e_k^{\frac{3}{4} - 8\varepsilon}, \qquad w^{\mathrm{E}}_k = e_k^{\frac{1}{4} - 7\varepsilon}.
\end{gather*}
Moreover
\[
R_{\mathrm{top}} = R_{\mathrm{top}}(e) = h M r_K(e), \qquad \mathrm{Tor}(e) = \left\{ h L^T : T \ge 0 \text{ an integer},\ h L^T \le R_{\mathrm{top}} \right\},
\]
and for every $R' \in \mathrm{Tor}(e)$
\[
\ell_k(R') = \frac{R'}{a_k}, \qquad T_k(R') = \mathbb{L}_k / \ell_k(R') \mathbb{Z}^3, \qquad \mathfrak{n}_k = \log_L \frac{\ell_k(R')}{M r_k};
\]
$k_{*}(R')$ denotes the first $k$ with $\mathfrak{n}_k = 0$, and $k_{\mathrm{top}} = k_{*}(R_{\mathrm{top}})$.
The class of the sides is
\[
\mathrm{Cl} = \left\{ \ell : \ell/(LM) \text{ is an odd integer } \ge 13 \right\}
\]
(Section~I.1.1).
At $k = 0$ one has $a_0 = \eta$; the charge $e_k$ of level $0$ is $q = e \eta^{\frac{1}{2}}$, the mass $\mu_k$ of level $0$ is $\bar{m}\eta$, and $\ell_0(R') = R'/\eta$.
From these formulas one reads off, for all $k \ge 0$ and $R' \in \mathrm{Tor}(e)$,
\begin{equation}
\tag{26.6.1}
a_{k+1} = L a_k, \qquad e_{k+1} = L^{\frac{1}{2}} e_k, \qquad \mu_{k+1} = L \mu_k, \qquad \ell_{k+1}(R') = \frac{\ell_k(R')}{L},
\end{equation}
and
\begin{equation}
\tag{26.6.2}
\parbox{0.8\linewidth}{the numbers $a_k$, $e_k$, $v_k$, $\mu_k$, $r_k$, $p_k$, $p_{0,k}$, $w^{\mathrm{m}}_k$, $w^{\mathrm{E}}_k$, $K$, $R_{\mathrm{top}}$ and the set $\mathrm{Tor}(e)$ are formed from $(\bar{m}, e, N)$ and the constants; $R'$ does not enter them.}
\end{equation}

\paragraph{2.2. The results of \cite[Section~19.2]{QEDI}.}
The assertions (Sc1)--(Sc4) of Section~G4.0.2 are the clauses (c), (d), (e), (f) of \cite[Proposition~19.1]{QEDI}, read with $h' = h$, $m_0 = 0$ and with the upper end of the mass range equal to $1$ (Section~I.6.1).
The assertion (Sc6) of Section~G4.0.2, that is, Proposition~14.2\,(c), (f) together with the row (e1) at $v = v_K$, which holds by the second assertion of Proposition~14.3\,(f), takes the place of \cite[Proposition~19.2]{QEDI}.
Apart from the assertions of \cite[Proposition~19.2]{QEDI} on the condition (W3) of Proposition~14.2\,(d), which (Sc6) does not contain, it is that proposition, read with $h' = h$, $m_0 = 0$, the upper end $1$ of the mass range, and with the constants and the number $e_1$ of Section~I.8.2 as defined in Proposition~14.2\,(a).
Of the rows (e1)--(e8) of Proposition~14.2\,(b), the two parts of (Sc6) that enter below, in (1) and in (3), need the row (e1) only; the arguments in brackets derive them from this row.
In \cite[Proposition~19.1]{QEDI} and in (Sc6) there occur only the numbers $(\bar{m}, e, R, N)$, constants and the conditions of Section~G4.0.3; neither an operator nor a primitive occurs in them.
Hence (Sc1)--(Sc4) and (Sc6), in the form of Section~G4.0.2, hold for every quadruple with (S-4).
Their hypotheses are verified in the following order.
\begin{itemize}
\item[(1)] The data of (Sc6) are $e > 0$ and $e_K \le e_1$.
Among its first assertions are $v_K \ge 8$ and $(1 + 3\,\mathrm{lg}/v_K)^r < L$; since $\mathrm{lg}$, $v_K$ and $r$ are positive, it follows that
\[
\left( 1 + \frac{\mathrm{lg}}{v_K} \right)^r \le \left( 1 + \frac{3\,\mathrm{lg}}{v_K} \right)^r < L.
\]
[One has $v_K = -\log e_K \ge -\log e_1$, and the row (e1) of Proposition~14.2\,(b) holds for all $v \ge -\log e_1$ by Proposition~14.3\,(f).]
\item[(2)] The hypotheses of \cite[Proposition~19.1]{QEDI} are $v_K \ge 1$, $(1 + \mathrm{lg}/v_K)^r \le L$ and $R \in \mathrm{Tor}(e)$; they hold by (1) and (S-4).
The clauses (c) and (e) of that proposition are stated for $N \ge N_{\min}(R,e)$, its clause (d) for $R \le R'$ in $\mathrm{Tor}(e)$ and $N \ge N_{\min}(R,e)$.
\item[(3)] By (Sc6), the condition (2.1) of Section~G4.0.3, namely $r_{*} > 31L + 16$, holds for every $R \in \mathrm{Tor}(e)$ and every $N \ge N_{\min}(R,e)$.
[By clause (e) of the proposition in (2), $v_{*} \ge v_K \ge -\log e_1$; by the definition of $r_{*}$, $r_{*} \ge v_{*}^r$; the row (e1) is then applied at $v = v_{*}$.]
\item[(4)] Clause (f) of \cite[Proposition~19.1]{QEDI} has the further hypothesis (2.1), which holds by (3); this clause is applied for $N \ge N_{\min}(R,e)$.
\end{itemize}
From now on (Sc1)--(Sc4) are applied in the formulation of Section~G4.0.2, to a quadruple $(\bar{m}, e, R, N)$ with (S-4).

\paragraph{2.3. $R_{\mathrm{top}}$ satisfies (S-4) with the same $(\bar{m}, e, N)$.}
By (M1), $M$ is a power of $L$, and by its definition $r_K(e) = L^i$ with an integer $i \ge 0$.
Thus $M r_K(e)$ is a power of $L$ with a non-negative integral exponent; in the notation of Section~I.6.2 this exponent is $T^{*}(e) + 1$, so that
\[
R_{\mathrm{top}} = h L^{T^{*}(e)+1}.
\]
As $R_{\mathrm{top}} \le R_{\mathrm{top}}$, the number $R_{\mathrm{top}}$ belongs to $\mathrm{Tor}(e)$, and it is the largest element of $\mathrm{Tor}(e)$; in particular $R \le R_{\mathrm{top}}$.
Taking $R' = R_{\mathrm{top}}$ in (Sc2), one obtains
\begin{equation}
\tag{26.6.3}
N_{\min}(R_{\mathrm{top}},e) \le N_{\min}(R,e) \le N, \qquad k_{*}(R) \le k_{\mathrm{top}}.
\end{equation}
The conditions imposed on $\bar{m}$ and on $e$ do not involve the torus.
Therefore $(\bar{m}, e, R_{\mathrm{top}}, N)$ satisfies (S-4), and (Sc1)--(Sc4) hold on $R$ as well as on $R_{\mathrm{top}}$.

\paragraph{2.4. Ranges.}
Let $R'$ be $R$ or $R_{\mathrm{top}}$.
For every $k \le k_{*}(R')$ one has the following.
\begin{itemize}
\item[(ra1)] $\ell_k(R')$ lies in $\mathrm{Cl}$, by (Sc4).
\item[(ra2)] $\ell_k(R') \ge M r_k$, by (Sc3).
\item[(ra3)] $0 < e_k \le e_K \le e_1 \le e_0$. Here $e_k > 0$ because $e > 0$ and $a_k > 0$; the three inequalities follow from (Sc3), from (S-4) and from $e_1 \le e_0$ (C9), in this order.
\item[(ra4)] $0 < \mu_k \le \mu_K \le 1$, by (Sc3).
\end{itemize}
For every $k < k_{*}(R')$, so that $k + 1 \le k_{*}(R')$, one has moreover the following.
\begin{itemize}
\item[(ra5)] $\ell_k(R')/L = \ell_{k+1}(R')$ lies in $\mathrm{Cl}$, by (26.6.1) and (ra1) at $k+1$.
\item[(ra6)] $e_k > 0$ and $L^{\frac{1}{2}} e_k = e_{k+1} \le e_0$, by (26.6.1) and (ra3) at $k+1$.
\item[(ra7)] $0 < \mu_k = \mu_{k+1}/L \le 1/L$, by (26.6.1) and (ra4) at $k+1$.
\end{itemize}
From (ra1), (ra3) and (ra4) one obtains
\begin{equation}
\tag{26.6.4}
\text{for every } k \le k_{*}(R') \text{ the quadruple } (k, \ell_k(R'), \mu_k, e_k) \text{ is (SF).}
\end{equation}
The numbers $v_k$, $r_k$, $p_k$, $p_{0,k}$, $w^{\mathrm{m}}_k$, $w^{\mathrm{E}}_k$ are formed with $e_k$ by the formulas of Sections~I.6.1 and I.5.4 that occur in (OS); they are displayed in 2.1.
Together with (ra1), (ra5), (ra7), (ra6) and (ra2) this yields
\begin{equation}
\tag{26.6.5}
\parbox{0.8\linewidth}{for every $k < k_{*}(R')$ the quadruple $(k, \ell_k(R'), \mu_k, e_k)$ is (OS-) with $\ell_k(R') \ge M r_k$; the level $k+1$ has the side $\ell_k(R')/L = \ell_{k+1}(R')$, the mass $L\mu_k = \mu_{k+1}$ and the charge $L^{\frac{1}{2}} e_k = e_{k+1}$.}
\end{equation}

\paragraph{2.5. Clause (c).}
By (Sc1), $k_{*}(R) \ge 1$.
For $R' = R$ and $k = k_{*}(R)$, the statement (26.6.4) says that $(k_{*}, \ell_*, \mu_*, e_{*})$ is (SF), that is, $\ell_* \in \mathrm{Cl}$, $\mu_* \in (0,1]$ and $e_{*} \in (0, e_0]$.

\paragraph{2.6. Clause (d).}
Let $R < R_{\mathrm{top}}$.
Then $R = h L^T$ with an integer $T \ge 0$, and $R_{\mathrm{top}} = h L^{T^{*}(e)+1}$ by 2.3; hence $T \le T^{*}(e)$, and
\[
\mathrm{Jr} := \frac{R_{\mathrm{top}}}{R} = L^{T^{*}(e)+1-T}
\]
is a power of $L$ with an exponent $\ge 1$.
Since $L$ is odd (L5), $\mathrm{Jr}$ is an odd integer $\ge 1$.
For every $k \ge 0$,
\[
\ell_k(R_{\mathrm{top}}) = \frac{R_{\mathrm{top}}}{a_k} = \frac{\mathrm{Jr} \cdot R}{a_k} = \mathrm{Jr} \cdot \ell_k(R).
\]
Let $k \le k_{*}(R)$.
Then $k \le k_{\mathrm{top}}$ by (26.6.3), and (26.6.4), applied on $R$ and on $R_{\mathrm{top}}$, shows that the quadruples $(k, \ell_k(R), \mu_k, e_k)$ and $(k, \ell_k(R_{\mathrm{top}}), \mu_k, e_k)$ are (SF); in particular $\ell_k(R)$ and $\ell_k(R_{\mathrm{top}})$ lie in $\mathrm{Cl}$.
By (26.6.2), the mass $\mu_k$ and the charge $e_k$ are the same numbers in the two quadruples.
These are the data of Proposition~26.1: a level $k \ge 0$; sides $\ell_1 := \ell_k(R)$ and $\ell_2 := \ell_k(R_{\mathrm{top}})$ in $\mathrm{Cl}$ with $\ell_2 = \mathrm{Jr} \cdot \ell_1$, where $\mathrm{Jr}$ is an odd integer $\ge 1$; and $\mu := \mu_k \in (0,1]$ and $e_k \in (0, e_0]$, the same on both tori.

\paragraph{2.7. The primitives; clauses (e), (a) and (b).}
\textit{(i) The top torus.}
Let $k < k_{\mathrm{top}}$.
By (26.6.4) for $R' = R_{\mathrm{top}}$, the quadruple $(k, \ell_k(R_{\mathrm{top}}), \mu_k, e_k)$ is (SF).
By Proposition~26.5, the families named there form a system of primitives of level $k$ on $T_k(R_{\mathrm{top}})$ with W1--W17$^{\mathrm{ch}}$ for $\mathrm{par}^{0}$ and for $\mathrm{par}$.
Thus the systems that (FL) names on $T_k(R_{\mathrm{top}})$, $k < k_{\mathrm{top}}$, exist.
For the quadruple $(\bar{m}, e, R_{\mathrm{top}}, N)$ they are all the primitives that belong to (FL), because by Section~G4.0.1 projections occur only for tori smaller than $R_{\mathrm{top}}$.
Together with 2.3 it follows that $(\bar{m}, e, R_{\mathrm{top}}, N)$ is (FL).
This is (a).

\textit{(ii) A smaller torus.}
Let $R < R_{\mathrm{top}}$ and $k < k_{*}(R)$.
Then $k < k_{\mathrm{top}}$ by (26.6.3), so that the system $\mathrm{Op}^{(2)}$ of (i) on $T_2 := T_k(R_{\mathrm{top}})$ is available.
By 2.6, the data of Proposition~26.1 are given, with $T_1 := T_k(R)$ and $\mathrm{Jr} = R_{\mathrm{top}}/R$.
The system $\mathrm{Op}^{(2)}$ has W14$^{\mathrm{ch}}$; hence the definition of the projection in Proposition~26.1 applies to it.
As $\mathrm{Op}^{(2)}$ has W1--W17$^{\mathrm{ch}}$ for $\mathrm{par}^{0}$, Proposition~26.3\,(b) shows that its projection $\mathrm{Op}^{\mathrm{pj}}$ is a system of primitives of level $k$ on $T_k(R)$, with W1--W17$^{\mathrm{ch}}$ for $\mathrm{par}$.
So the primitives that (FL) names on $T_k(R)$, $k < k_{*}(R)$, exist.
Together with 2.3, (26.6.3), (26.6.4) for $R' = R_{\mathrm{top}}$ and (i), this is (e).

\textit{(iii) Clause (b).}
Let $k < k_{*}(R)$.
By (26.6.5) for $R' = R$, the quadruple $(k, \ell_k(R), \mu_k, e_k)$ is (OS-) with $\ell_k(R) \ge M r_k$, and the next level has the side, the mass and the charge asserted in (b).
The primitives of (FL) on $T_k(R)$ are the system of (i) if $R = R_{\mathrm{top}}$ (then $k < k_{\mathrm{top}}$), and the system of (ii) if $R < R_{\mathrm{top}}$.
In both cases they form a system of primitives of level $k$ on $T_k(R)$ with W1--W17$^{\mathrm{ch}}$ for $\mathrm{par}$.
The data are therefore (OS) in the sense of Sections~I.1.8 and G4.0.1, and (b) is proved.

\paragraph{2.8. Level 0.}
By (Sc1), $k_{*}(R') \ge 1$ for $R' = R$ and for $R' = R_{\mathrm{top}}$; hence $k = 0$ is one of the levels of (b) on both tori.
There the data are $(0, \ell_0(R'), \bar{m}\eta, q)$ with $\ell_0(R') = R'/\eta$, and (26.6.5) reads as follows: $R'/\eta$ and $R'/(L\eta)$ lie in $\mathrm{Cl}$; $0 < \bar{m}\eta \le 1/L$; $q > 0$ and $L^{\frac{1}{2}} q \le e_0$, with the ceiling $e_0$ of the line (C8); and $R'/\eta \ge M r_0$, where
\[
r_0 = \min\left\{ L^i : i \ge 0 \text{ an integer},\ L^i \ge (-\log q)^r \right\}.
\]
The level $1$ has the side $R'/(L\eta)$, the mass $L\bar{m}\eta$ and the charge $L^{\frac{1}{2}} q$.
The system of (b) is a system of level $0$ in the sense of Section~G4.0.4: the family (p2) $C_{0}^{1/2}(s')$; the families $\Gamma_0(s'',A)$ and $S_{0,y}(X,A)$ on $\mathfrak{S}^{\mathrm{ch}}(1, R'/\eta)$ (Section~I.1.5), the second Borel measurable in $y$; the diagonal families $H_{0}(s')$ and $H_{0}(s'',A)$; and no family $M_0$.
On $R_{\mathrm{top}}$ these are the families of 1.7.
On $R < R_{\mathrm{top}}$ they are their projections (Proposition~26.3\,(b); explicitly, Proposition~26.3\,(c)).
Indeed, by Proposition~26.3\,(b) the projection of the system of 1.7 on $T_0(R_{\mathrm{top}})$ is a system of level $0$ on $T_0(R)$, and in a system of level $0$ the families (p1), (p3) are the diagonal families of Section~I.5.8 (Section~G4.0.4); so the projections of $H_{0}(s')$ and $H_{0}(s'',A)$ on $T_0(R_{\mathrm{top}})$ are the diagonal families $H_{0}(s')$ and $H_{0}(s'',A)$ on $T_0(R)$, which is the first assertion of Proposition~26.3\,(c).
By Proposition~26.3\,(b), the bounds W6 for $\Gamma_0^{\mathrm{pj}}$ and W7 for $S^{\mathrm{pj}}_{0,y}$ are bounds, for $\mathrm{par}$, of the kernels (I.1.4) at pairs of fermion sites, the second with the factor $(1+y)^{-1}$.
In (d), the data of Proposition~26.1 at $k = 0$ are $\ell_1 = R/\eta$, $\ell_2 = R_{\mathrm{top}}/\eta = \mathrm{Jr} \cdot R/\eta$, $\mu = \bar{m}\eta$, and $q$ in the place of $e_k$; the source and target spaces of $\Gamma_0$, $S_{0,y}$ and $H_{0}(s'',A)$ are $\mathfrak{S}^{\mathrm{ch}}(1, \ell_i)$.

\paragraph{2.9. Uniformity.}
The set of parameters $\mathrm{par} = (\alpha_{0}, \gamma_{\mathrm{w}}, \kappa_{\mathrm{w}}; C_{w})$ consists of constants of the lines (C0)--(C6) of Section~I.8.2.
In Proposition~26.5 it does not depend on $(k, \ell, \mu, e_k)$ (see 1.9), and in Proposition~26.3\,(b) it does not depend on $(k, \ell_1, \ell_2, \mathrm{Jr}, \mu, e_k)$.
Consequently, in (b) the properties W1--W17$^{\mathrm{ch}}$ hold with one set of parameters for all $\bar{m} \in (0,1]$, all $e$ with $e h^{\frac{1}{2}} \le e_1$, all $R \in \mathrm{Tor}(e)$, all $N \ge N_{\min}(R,e)$ and all levels $k < k_{*}(R)$.
Neither the ceilings $e_0$, $e_1$ nor the bounds $1$, $1/L$ on the masses in (ra3), (ra4), (ra6), (ra7) depend on these data.
\end{proof}

\begin{named}{Remark}[How (a)--(d) are applied]
Let (FL) hold.
By (b), a statement whose data are (OS) applies to $(k, \ell_k(R), \mu_k, e_k)$, with the primitives of (FL), for every $k < k_{*}(R)$; by (a), by (b) on $R_{\mathrm{top}}$ and by (26.6.3), it applies also to $(k, \ell_k(R_{\mathrm{top}}), \mu_k, e_k)$, with the systems of Proposition~26.5, for these $k$.
Data (OS) are data (SF+) of Section~I.1.8, since $\mu_k \le 1/L \le 1$ and $e_k < L^{\frac{1}{2}} e_k \le e_0$; and the family (p5) of the system has, among the properties W1--W17$^{\mathrm{ch}}$ for $\mathrm{par}$, the properties W1, W2, W7, W12, W13, W14 for $\mathrm{par}$ (on $R < R_{\mathrm{top}}$ by Proposition~26.3\,(b)).
Hence a statement whose data are (SF) or (SF+), together with a family (p5) that has these properties for $\mathrm{par}$, applies at the same levels.
In particular, Proposition~13 applies at these levels, with the family (p5) of the system as its family (p5); for $R < R_{\mathrm{top}}$ this is the family (p5) of the projected system.
By (c), a statement whose data are (SF) with $k \ge 1$ applies at $k = k_{*}(R)$.
By (d), if $R < R_{\mathrm{top}}$, a statement whose data are those of Proposition~26.1 applies to the pair $T_k(R)$, $T_k(R_{\mathrm{top}})$ for every $k \le k_{*}(R)$.
\end{named}

\paragraph{2.10. Appendix: the arithmetic behind (Sc1)--(Sc4), in the notation of Part~I.}
What follows is the part of \cite[Section~19.2]{QEDI} that precedes \cite[Proposition~19.1]{QEDI}, namely the definition of the numbers of depth $j$ together with the expression of $a_k$, $e_k$, $v_k$, $r_k$, $\ell_k$, $\mathfrak{n}_k$ through them at $j = K - k$, followed by \cite[Proposition~19.1(a)--(f)]{QEDI} and the part of (Sc6) that enters 2.2; all of it is written with $h' = h$, $m_0 = 0$, the number $e_1$ of Proposition~14.2\,(a) and the data (S-4).
Throughout, $R = h L^T$ and $\mathrm{lg} = \frac{1}{2} \log L$.
The numbers $\mathrm{vh}_j$, $\mathrm{rh}_j$, $\mathrm{nch}_j(R)$ and $\mathrm{js}(R)$ defined below are the numbers $\hat{v}_j$, $\hat{r}_j$, $\hat{\mathfrak{n}}_j$ and $j_{*}$ of \cite[Section~19.2]{QEDI}.
\begin{itemize}
\item[(ar0)] \textit{Depth.}
For integers $j \ge 0$ let
\begin{gather*}
\mathrm{vh}_j := v_K + j \cdot \mathrm{lg}, \qquad \mathrm{rh}_j := \min\left\{ L^i : i \ge 0 \text{ an integer},\ L^i \ge \mathrm{vh}_j^r \right\},\\
\mathrm{nch}_j(R) := T + j - \log_L\left( M \cdot \mathrm{rh}_j \right).
\end{gather*}
These numbers are formed from $(e_K, L, M, r, T, j)$; $N$ does not enter them.
As $M$ and $\mathrm{rh}_j$ are powers of $L$, $\mathrm{nch}_j(R)$ is an integer.
Let $0 \le k \le K$ and $j = K - k$.
Then $a_k = L^{k-N} = h L^{-j}$, because $a_K = L^{K-N} = h$; further $e_k = e\, a_k^{\frac{1}{2}} = e_K L^{-\frac{j}{2}}$, so that $v_k = v_K + j \cdot \mathrm{lg} = \mathrm{vh}_j$ and $r_k = \mathrm{rh}_j$; finally $\ell_k(R) = R/a_k = L^{T+j}$ and
\[
\mathfrak{n}_k = \log_L \frac{L^{T+j}}{M \cdot \mathrm{rh}_j} = \mathrm{nch}_j(R).
\]
\item[(ar1)] One has $v_K \ge 8$ and
\[
\left( 1 + \frac{\mathrm{lg}}{v_K} \right)^r \le \left( 1 + \frac{3\,\mathrm{lg}}{v_K} \right)^r < L.
\]
[This is the row (e1), which holds for all $v \ge -\log e_1$ (Proposition~14.3\,(f)), taken at $v = v_K = -\log e_K \ge -\log e_1$.]
\item[(ar2)] For $j \ge 0$,
\[
\mathrm{rh}_{j+1} \in \left\{ \mathrm{rh}_j,\ L \cdot \mathrm{rh}_j \right\} \qquad \text{and} \qquad \mathrm{nch}_{j+1}(R) - \mathrm{nch}_j(R) \in \{0, 1\}.
\]
[Since $\mathrm{vh}_{j+1} > \mathrm{vh}_j > 0$, one has $\mathrm{rh}_{j+1} \ge \mathrm{rh}_j$.
By $\mathrm{vh}_j \ge v_K$ and (ar1),
\[
\mathrm{vh}_{j+1}^r = \mathrm{vh}_j^r \cdot \left( 1 + \frac{\mathrm{lg}}{\mathrm{vh}_j} \right)^r \le \mathrm{vh}_j^r \cdot \left( 1 + \frac{\mathrm{lg}}{v_K} \right)^r \le L \cdot \mathrm{vh}_j^r \le L \cdot \mathrm{rh}_j;
\]
as $L \cdot \mathrm{rh}_j$ is a power of $L$ with a non-negative exponent, $\mathrm{rh}_{j+1} \le L \cdot \mathrm{rh}_j$.
Both numbers are powers of $L$, and $\mathrm{nch}_{j+1}(R) - \mathrm{nch}_j(R) = 1 - \log_L( \mathrm{rh}_{j+1}/\mathrm{rh}_j )$.]
\item[(ar3)] One has $\mathrm{nch}_0(R) \le 0$, and $\mathrm{nch}_j(R) \to \infty$ as $j \to \infty$.
Hence
\[
\mathrm{js}(R) := \max\left\{ j \ge 0 : \mathrm{nch}_j(R) = 0 \right\}
\]
exists, $\mathrm{nch}_j(R) \le 0$ for $j \le \mathrm{js}(R)$, and $\mathrm{nch}_j(R) \ge 1$ for $j > \mathrm{js}(R)$.
[The inequality $R \le R_{\mathrm{top}} = h M r_K = h M \cdot \mathrm{rh}_0$ says $L^T \le M \cdot \mathrm{rh}_0$.
By (ar1), $\mathrm{vh}_j = v_K + j \cdot \mathrm{lg} \ge v_K \ge 8$, and $r > 0$; so $\mathrm{vh}_j^r \ge 1$.
Hence $\mathrm{rh}_j < L \cdot \mathrm{vh}_j^r$: if $\mathrm{rh}_j = 1$, then $\mathrm{rh}_j \le \mathrm{vh}_j^r$; if $\mathrm{rh}_j = L^i$ with $i \ge 1$, then $L^{i-1} < \mathrm{vh}_j^r$ by the minimality of $i$.
Consequently
\[
\mathrm{nch}_j(R) > T + j - 1 - \log_L M - r \log_L\left( v_K + j \cdot \mathrm{lg} \right),
\]
and the right-hand side tends to infinity.
By (ar2), the integers $\mathrm{nch}_j(R)$ form a non-decreasing sequence with steps $0$ or $1$; it starts at a value $\le 0$ and tends to infinity, so it takes the value $0$.]
\item[(ar4)] \textit{(Sc1).}
For an integer $N' \ge 0$ put $K' := N' + \log_L h$.
At the cutoff $N'$ the number $R L^{N'}$ is $\ell_0(R)$; so the inequality $R L^{N'} \ge L M r_0$ says $\mathfrak{n}_0 \ge 1$, that is, $\mathrm{nch}_{K'}(R) \ge 1$ by (ar0), that is, $K' > \mathrm{js}(R)$ by (ar3).
The number $\mathrm{js}(R)$ does not involve $N'$.
Therefore
\[
N_{\min}(R,e) = \max\left\{ 0,\ \mathrm{js}(R) + 1 - \log_L h \right\}.
\]
Let $N \ge N_{\min}(R,e)$.
Then $K > \mathrm{js}(R)$; $\mathfrak{n}_k = \mathrm{nch}_{K-k}(R) \ge 1$ for $0 \le k < K - \mathrm{js}(R)$, and $\mathfrak{n}_k = 0$ at $k = K - \mathrm{js}(R)$.
Hence $k_{*}(R) = K - \mathrm{js}(R)$, and $1 \le k_{*}(R) \le K$.
\item[(ar5)] \textit{(Sc2).}
Let $R \le R' = h L^{T'}$ in $\mathrm{Tor}(e)$.
Then $\mathrm{nch}_j(R') = \mathrm{nch}_j(R) + T' - T \ge \mathrm{nch}_j(R)$; so $\mathrm{nch}_j(R') \ge 1$ for $j > \mathrm{js}(R)$, that is, $\mathrm{js}(R') \le \mathrm{js}(R)$.
By (ar4), $N_{\min}(R',e) \le N_{\min}(R,e)$, and for $N \ge N_{\min}(R,e)$
\[
k_{*}(R') = K - \mathrm{js}(R') \ge K - \mathrm{js}(R) = k_{*}(R).
\]
\item[(ar6)] \textit{(Sc3).}
Let $k \le k_{*}(R)$ and $j = K - k \ge \mathrm{js}(R)$.
By (ar3), $\mathfrak{n}_k = \mathrm{nch}_j(R) \ge 0$, that is, $\ell_k(R) \ge M r_k$.
The number $\mathrm{vh}_j$ increases with $j$; thus $v_k = \mathrm{vh}_j \ge \mathrm{vh}_{\mathrm{js}(R)} = v_{*} \ge \mathrm{vh}_0 = v_K$, and so $e_k \le e_{*} \le e_K$.
By (ar2), $\mathrm{rh}_j$ is non-decreasing in $j$, whence $r_k \ge r_{*}$.
The mass $\mu_k = \bar{m} a_k > 0$ increases with $k$, and $k \le k_{*}(R) \le K$; so $\mu_k \le \mu_* \le \mu_K$.
With $t := \log_L(1/\bar{m}) \ge 0$ one has $\mu_K = \bar{m} h = L^{\lfloor t \rfloor - t}$, a number in $(1/L, 1]$.
\item[(ar7)] \textit{(2.1) and (Sc4).}
By the definition of $r_{*}$, $r_{*} \ge v_{*}^r$, and $v_{*} \ge v_K \ge -\log e_1$ by (ar6); the row (e1) at $v = v_{*}$ gives $r_{*} > 31L + 16$.
Since $31L + 16 > L$ and $r_{*}$ is a power of $L$, it follows that $r_{*} \ge L^2$.
For $k \le k_{*}(R)$ one has $\ell_k(R) = M r_k L^{\mathfrak{n}_k}$, so that
\[
\frac{\ell_k(R)}{LM} = r_k L^{\mathfrak{n}_k - 1},
\]
with $r_k \ge r_{*} \ge L^2$ and $\mathfrak{n}_k \ge 0$ by (ar6).
This is $L^i$ with an integer $i \ge 1$.
Moreover $L$ is odd by (L5), and $L \ge 13$.
[The row (L1) reads $L^{\frac{1}{8} - \frac{3\bar{\varepsilon}}{2}} \ge 2400$ with $\bar{\varepsilon} = 1/200$ (C0); the exponent is $\frac{47}{400}$, a number in $(0,1]$, and $L \ge 1$; hence $L \ge L^{\frac{47}{400}} \ge 2400$.]
So $\ell_k(R)/(LM)$ is an odd integer $\ge 13$, that is, $\ell_k(R)$ lies in $\mathrm{Cl}$.
\end{itemize}

\subsection*{3. Conditions on the constants}

Throughout this chapter the constants, their definitions and the conditions imposed on them are those of Section~I.8.2 and of Proposition~14.2.
For each condition or definition in the tables below, the column ``where'' gives the line of Section~I.8.2 and the row of Proposition~14.2 in which it is stated.

\subsubsection*{3.1. Used in the arguments of Sections 1 and 2 themselves}

\begin{tabular}{p{0.40\textwidth}p{0.24\textwidth}p{0.26\textwidth}}
\hline
condition or definition & where & used in\\
\hline
$\alpha_{0} = \min(\alpha^{F}, \alpha^{P})$, $\alpha^{F} = 1/54$, $\alpha^{P} = 1/4$ & (C4); Proposition~14.2\,(a) & (cs1)\\
$\gamma_{\mathrm{w}} = \min(\gamma^{F}, \gamma^{P})$; $\gamma^{F} > 0$; $\gamma^{P} = \gamma_{P}/2 > 0$ & (C4), (C4r); Proposition~14.2\,(a) & 1.4 (the definition $\gamma^{P} = \gamma_{P}/2$); (cs2); (26.5.1)--(26.5.10); the rate in (tp3)\\
$C_{w}^{(0)} = \max(C^{F}, C^{P})$, $C_{w}^{(1)} = \max\left( C_{w}^{(0)}, 4\exp(2\gamma_{\mathrm{w}}) M^{27\alpha_{0}} \right)$, $C_{w} = 64 K_{0} C_{w}^{(1)}$ & (C6); Proposition~14.2\,(a) & (cs3), (cs4)\\
$C^{P} = 116\, C_{P} M^{402\alpha_{0}}$, $C_{P} \ge 4\exp(2)$; $C^{F} \ge 2 C_{\mathrm{By}}$, $C_{\mathrm{By}} = (1 + 2\mathrm{e})\sqrt{2}\max(1, 1/b_{-})$;
the definition of $C^{P}$ enters in two ways: in 1.4 it shows that the constant $C^{P}$ of Proposition~14.2\,(a) is the constant $C^{P}$ of \cite[Section~15.10.3]{QEDI}, so that (tp1), (tp2), and with them the first inequalities of (26.5.1), (26.5.2), (26.5.7), hold with this $C^{P}$; in (cs3) it enters through $C^{P} > 0$;
the bounds $C_{P} \ge 4\exp(2)$, $C^{F} \ge 2 C_{\mathrm{By}}$ and the formula for $C_{\mathrm{By}}$ enter only through $C^{P} > 0$ and $C^{F} > 0$
& (C4), (C6); Proposition~14.2\,(a); \cite[Section~15.7]{QEDI}
& 1.4; (cs3); the hypothesis $0 < C_{\mathrm{old}}$ of Lemma~A in 1.6, 1.7;
through $C_{w}^{(1)} \ge C^{F} > 0$, the hypothesis $0 < C_{\mathrm{old}} = C_{w}^{(1)}$ of Lemma~A in 1.8 and the inequality $C_{w}^{(1)} \le 64 K_{0} \cdot C_{w}^{(1)}$ of (cs4)\\
$\kappa_{0} = 54 \log 72$, which enters only through $\kappa_{0} \ge 0$; $K_{0} = \exp(\kappa_{0})/36$, which enters only through $64 K_{0} \ge 1$ & (C0); Proposition~14.2\,(a) & (cs4), 1.8\\
$\kappa' = \kappa_{\mathrm{w}} + \kappa_{0}$, $\kappa_{\mathrm{w}} = \kappa + \kappa_{0} + 3$ & (C4); Proposition~14.2\,(a) & 1.5: the value at which $M^{F}$, $C_1$, $C^{F}$ are formed\\
$M \ge 1$ & (C5), (M1) & (cs3), (cs5)\\
$M$ a power of $L$ & (C5), (M1) & 2.3; 2.10 (ar0) ($\mathrm{nch}_j(R)$ is an integer), and through it (ar3), (ar7)\\
$L$ odd & (C3), (L5) & 2.6 ($\mathrm{Jr}$ odd)\\
$e_1 \le e_0$ & (C9); Proposition~14.2\,(a), row $e_1$ & (ra3), (ra6)\\
\hline
\end{tabular}

\subsubsection*{3.2. Hypotheses of the statements of \cite{QEDI} that are applied}

\begin{tabular}{p{0.40\textwidth}p{0.24\textwidth}p{0.26\textwidth}}
\hline
condition & where & statement\\
\hline
$0 \le \alpha < 1$ & (C1) & \cite[Theorem~15.35]{QEDI}\\
$L$ odd, $L \ge 3$ & (C3), (L5) & \cite[Theorem~15.35, Proposition~15.36]{QEDI}\\
$0 < \alpha_{0} \le 1/4$ & (C4) & \cite[Theorem~15.35, Proposition~15.36]{QEDI}\\
$M$ a power of $L$, $M \ge L^2$ & (C5), (M1) & the opening paragraph of \cite[Section~15.10]{QEDI}\\
$M \ge M^{P}$ & (C5), (M6) & \cite[Theorem~15.35]{QEDI}\\
$M$ a power of $L$ & (C5), (M1) & \cite[Proposition~19.1]{QEDI} (a standing assumption of \cite[Section~5.2.3]{QEDI}, invoked in the proof of \cite[Proposition~19.1]{QEDI}; here 2.10 (ar0), (ar7))\\
for all $v \ge -\log e_1$: $v \ge 8$, $v^r > 31L + 16$, $(1 + 3\,\mathrm{lg}/v)^r < L$ & (C9), (e1); the number $e_1$ of Proposition~14.2\,(a) satisfies it by Proposition~14.3\,(f) & (Sc6) of Section~G4.0.2, in the place of \cite[Proposition~19.2]{QEDI} (the part that enters: $v_K \ge 8$, $(1 + \mathrm{lg}/v_K)^r \le L$, (2.1)); through it \cite[Proposition~19.1]{QEDI}\\
$L$ odd, $L \ge 13$ (the proof of \cite[Proposition~19.1(f)]{QEDI} uses in addition $L \ge 33$, which 2.10 (ar7) does not need) & (C3), (L5), (L1) with $\bar{\varepsilon} = 1/200$ (C0) (it gives $L \ge 2400$: 2.10 (ar7)) & \cite[Proposition~19.1(f)]{QEDI}, that is (Sc4)\\
$r$ an integer $\ge 2$, which enters only through $r > 0$, that is, $r \ge 1$ & (C7) & \cite[Proposition~19.1]{QEDI}: the hypothesis $(1 + \mathrm{lg}/v_K)^r \le L$, from (Sc6) (2.2 (1), 2.10 (ar1)); clauses (a), (b) (2.10 (ar2), (ar3))\\
\hline
\end{tabular}

\subsubsection*{3.3. Conditions that enter through the Propositions applied (as stated with these Propositions)}

\begin{tabular}{p{0.40\textwidth}p{0.24\textwidth}p{0.26\textwidth}}
\hline
condition & where & statement\\
\hline
$M \ge M^{F}$ at $\kappa' = \kappa_{\mathrm{w}} + \kappa_{0}$, where $M^{F}$ is the largest of the numbers $4L + 1$, $(108\pi C'_{V} a')^2$ and $108\pi C'_{V} a' \exp(27\kappa')$ (this contains $M > 4L$) & (C5), (M5), (M2) & Proposition~14.1\,(ii) (Propositions~3, 7, 8)\\
$L$ odd and $L \ge 5$ (L5); $M$ a power of $L$ and $M \ge L^2$ (M1); (Z2.s1), (Z2.s2); $\alpha < 1$
& (C3), (L5); (C5), (M1); (C8) with the rows (Z2.s1), (Z2.s2) of Proposition~14.2\,(b); (C1)
& Proposition~14.1\,(ii) (Propositions~3, 4\,(a), 7, 8)\\
(Z2.s3), (Z2.s4)
& (C8); Proposition~14.2\,(b)
& they are among the conditions on the constants of Proposition~26.5, but enter neither Proposition~14.1\,(ii) nor any step of the arguments of Sections~1 and~2;
they are used in Proposition~6\,(b)--(f), Proposition~9\,(d), Proposition~13\,(d4)--(d6) and in the assertions (G3) of \cite[Section~16.2.7]{QEDI}\\
$L$ odd; $M$ a power of $L$, $M \ge L^2$, $M > 4L$, $\gamma_{\mathrm{w}} M \ge 1$, $M \ge M^{P}$; $0 < \alpha_{0} \le 1/4$, $\kappa_{\mathrm{w}} \ge 0$; $K_{0} \ge 1$; $C_{w}^{(1)} \ge 4\exp(2\gamma_{\mathrm{w}}) M^{27\alpha_{0}}$;
in addition, at W1 and through Proposition~11\,(b), the conditions of Proposition~11, namely (L5), (M1), (M2), the part $M \ge 108\pi \cdot C'_{V} \cdot a'$ of (M5), (Z2.s1), (Z2.s2) and $\alpha < 1$;
and, at W17 and through Proposition~26.3\,(c) and Proposition~8\,(b), the condition (M5) in full
& (C3), (L5); (C5), (M1)--(M3), (M5), (M6); (C4) with (C2), (C0); (C0); (C6); (C8) with the rows (Z2.s1), (Z2.s2) of Proposition~14.2\,(b); (C1)
& Proposition~26.3\,(b)\\
\hline
\end{tabular}

In \cite[Theorem~5.72]{QEDI} the charge is bounded by a number $e^{F}_{0}$ of that theorem; here (Z2.s1)--(Z2.s4) are conditions on the ceiling $e_0$ of the line (C8) itself, and $e^{F}_{0}$ occurs in none of the conditions of this section (Section~I.8.4).
No condition of this section contains $k$, $\ell$, $\mu$, $e_k$, $\bar{m}$, $e$, $R$ or $N$.

\section*{Chapter 27a. Two tori compared near a flat set (lifts, generators, primitives, exact operators): proofs of Propositions 27.1--27.4}

\subsection*{0. Conventions of this chapter; elementary facts on cubes, boxes and covering maps}

Throughout this chapter the standing data (TT), fixed at the head of the Propositions~27.$m$, are in force.

\paragraph{Notation.}
The index $i$ takes the values $1$ and $2$.
The letters $w,w'$ denote vectors of $\mathbb{Z}^3$.
The letters $v,v'$ denote vectors of the lattice that is named where they occur:
of $M\mathbb{Z}^3$ in (0.3b), of $LM\mathbb{Z}^3$ in the first paragraph and in the table of (0.5b), and of $\ell_i\mathbb{Z}^3$ or $\ell_1\mathbb{Z}^3$ elsewhere.
Photon sites of $\mathbb{L}_k$ are denoted by $u,u',c,c'$, and centres of cubes or blocks by $z,z'$.
Where $a$ or $a'$ occurs in an exponent ${}^{\sim a}$, it is an integer $\ge0$.
An index of a torus $T_i$ is written $\mathrm{tq}$, and a fermion site of $T_i$ is written $\mathrm{xq}$;
the plain letters $t$ and $x$ denote covering indices and sites of $\mathbb{L}^{\mathrm{f}}_k$,
except in the statement and in the proof of Proposition~27.3, where $t$ is an index of $T_1$ and $\tilde{t}$ is its lift.
Three letters are also used with a second meaning.
As the argument of $\Psi$ or of $W$, and in 4\,(f), the letters $u,u_1,\dots,u_p$ are indices $(\mathrm{mo},\omega)$ of these generators (Section~I.3.6; Proposition~15\,(D2)).
In 4\,(a-3) the $w_j$ are the weights of Section~I.0.5.
In 4.2 the letter $\tilde{U}$ denotes a finite set of sites of $\mathbb{L}_k$, as in Section~I.1.9,
and in 4\,(e) it denotes the set of the statement, a union of $M$-cubes of $\mathbb{L}_{k+1}$.
The index set introduced in 2\,(a) is denoted by $\mathrm{Ix}(\tilde{X})$.

In this chapter, numbers such as 0.1--0.5 and 4.0--4.2 without further reference denote the subsections of Section~0 and the paragraphs 4.0--4.2 of Section~4.
The label (2.1) refers to the condition $r_{*}>31L+16$, which is (G4.2.1) of Section~G4.0.3;
the numbered display of Section~2 is (2.0).
The steps of the proof of Proposition~27.2\,(e) are labelled (e1)--(e9);
the condition on $e_1$ that carries the same label is referred to as the row (e1) of Proposition~14.2\,(b).
The condition (k2) is (I.k2) of Section~I.5.4, and (AP) is (I.AP) of Section~I.3.6.
The definitions (D1)--(D4) are those of Proposition~27.2;
(8.0) is the display of Proposition~26.3\,(a);
(K1)--(K4) are the kernels displayed in Proposition~25.1;
(23.1) is the display of Proposition~23.1, and (23.3) is the display of Proposition~23.2.

Cubes, polymers and parameter cubes are those of Section~I.1.1,
which agree with the definitions of \cite[Section~5.2.3, Definitions~5.7 and~5.8]{QEDI};
the objects $V_k$, $\{\Upsilon\}$, $C$, $\mathrm{Box}_k$, $\mu_I$, $\hat{\mu}$ and $\delta\varepsilon^{0}_{k}$ are those of Section~I.5.7,
which agree with the definitions of \cite[Section~5.2.5]{QEDI}.
The indices, their supports and the map $\mathrm{pr}_i$ on them are those of the Definition that precedes Proposition~27.1.
That Definition rests on the index sets (I.i1)--(I.i5) of Section~I.1.9, on Section~G4.0.4 for the supports of bonds and of basis elements, and on Section~I.1.7 for pairs;
the sorts of indices, their supports and the translations acting on them are listed in (0.5b) below.
These indices comprise the indices of \cite[Section~5.2.3]{QEDI}, namely the sites, the bonds and the basis elements, with the same supports,
and in addition the fermion sites, the pairs of $\mathrm{Pzt}_k$, the blocks $\mathrm{Bl}_{j}(z)$, the $M$-cubes, the $LM$-cubes and the parameter cubes.

\subsubsection*{0.1. Arithmetic of $\ell_1$, $\ell_2$, $L$ and $M$}

\begin{itemize}
\item[(0.1a)] The sides $\ell_1$ and $\ell_2$ belong to $\mathrm{Cl}$, so that they are integers and multiples of $L\cdot M$.
Moreover $\ell_2=\mathrm{Jr}\cdot\ell_1$ with an odd integer $\mathrm{Jr}\ge1$ (Proposition~26.6\,(d)).
Consequently $\ell_2\mathbb{Z}^3\subset\ell_1\mathbb{Z}^3\subset LM\mathbb{Z}^3$.
\item[(0.1b)] $\ell_2\ge\ell_1>(31L+16)\cdot M$.
\item[(0.1c)] $L$ is odd, and $M$ is a power of $L$ (Section~I.0.5); further $\eta_k=L^{-k}\le1\le M$.
\end{itemize}
For a set $\tilde{S}$ of sites of $\mathbb{L}_k$ put
\[
\mathrm{sdiam}(\tilde{S}):=\sup\left\{|u-u'|:\ u,u'\in\tilde{S}\right\}.
\]

\subsubsection*{0.2. Blocks and cubes}

Let $a$ be one of the numbers $L^j\eta_k$ with $0\le j\le k+1$, $M$ and $L\cdot M$, and let $z\in a\mathbb{Z}^3$.
Put
\[
\mathrm{Kub}_a(z):=\left\{c\in\mathbb{L}_k:\ |c-z|<\frac{a}{2}\right\}.
\]
Then $\mathrm{Kub}_{L^j\eta_k}(z)=\mathrm{Bl}_{j}(z)$ and $\mathrm{Kub}_1(y)=\Delta_y$,
while $\mathrm{Kub}_M(z)$ is the $M$-cube and $\mathrm{Kub}_{LM}(z)$ is the $LM$-cube with the centre $z$ (Section~I.1.1).
By (0.1c) the number $a/\eta_k$ is an odd integer, and $c_\mu-z_\mu\in\eta_k\mathbb{Z}$ for every $c\in\mathbb{L}_k$.
Recall that, for an odd integer $\mathrm{od}\ge1$, the sets
\[
\left\{j_1\in\mathbb{Z}:\ |j_1-\mathrm{od}\cdot j_2|\le\frac{\mathrm{od}-1}{2}\right\},\qquad j_2\in\mathbb{Z},
\]
form a partition of $\mathbb{Z}$,
and that $|j_1-\mathrm{od}\cdot j_2|<\frac{\mathrm{od}}{2}$ holds if and only if $|j_1-\mathrm{od}\cdot j_2|\le\frac{\mathrm{od}-1}{2}$.
Applied with $\mathrm{od}=a/\eta_k$ in each coordinate, this gives the following facts.
\begin{itemize}
\item[(0.2a)] For $c\in\mathbb{L}_k$ one has $|c-z|<\frac{a}{2}$ if and only if $|c-z|\le\frac{a-\eta_k}{2}$.
Moreover $\mathrm{sdiam}(\mathrm{Kub}_a(z))=a-\eta_k$,
and each site of $\mathbb{L}_k$ lies in exactly one of the sets $\mathrm{Kub}_a(z)$, $z\in a\mathbb{Z}^3$.
\item[(0.2b)] (Nesting.) Let $a,a'$ be two of these numbers such that $a'/a$ is an odd integer,
let $z\in a\mathbb{Z}^3$ and $z'\in a'\mathbb{Z}^3$, and suppose that $z$ lies in $\mathrm{Kub}_{a'}(z')$.
Then $|z-z'|\le\frac{a'-a}{2}$, and $\mathrm{Kub}_a(z)\subset\mathrm{Kub}_{a'}(z')$.
Indeed, each difference $z_\mu-z'_\mu$ belongs to $a\mathbb{Z}$ and has modulus less than $\frac{a'}{2}$;
by the fact recalled above, applied with $\mathrm{od}=a'/a$, its modulus is at most $\frac{a}{2}\cdot\left(\frac{a'}{a}-1\right)=\frac{a'-a}{2}$.
For $c\in\mathrm{Kub}_a(z)$ it follows that
\[
|c-z'|\le\frac{a-\eta_k}{2}+\frac{a'-a}{2}=\frac{a'-\eta_k}{2}.
\]

\emph{Consequences.}
Every $\Delta_y$ and every $\mathrm{Bl}_{k+1}(y')$, $y'\in L\mathbb{Z}^3$, lies in one $M$-cube, namely in the $M$-cube of its centre.
Further, $B(y')=\mathrm{Bl}_{k+1}(y')\cap\mathbb{Z}^3$;
every $M$-cube lies in one $LM$-cube;
and $\mathrm{Kub}_{LM}(z')$ is the union of the $L^3$ $M$-cubes $\mathrm{Kub}_M(z'+Mw)$ with $|w|\le\frac{L-1}{2}$.
For the last assertion, observe first that $|Mw|\le\frac{(L-1)\cdot M}{2}<\frac{L\cdot M}{2}$,
so that these $L^3$ cubes lie in $\mathrm{Kub}_{LM}(z')$ by the first part.
Conversely, let $c$ belong to $\mathrm{Kub}_{LM}(z')$ and to $\mathrm{Kub}_M(z)$.
The site $z$ lies in some $\mathrm{Kub}_{LM}(z'')$, which then contains $\mathrm{Kub}_M(z)$ and hence $c$;
by (0.2a), $z''=z'$, and therefore $|z-z'|\le\frac{L\cdot M-M}{2}$.
These consequences are contained in the fact (g0) of \cite[Section~17.4.1]{QEDI}.
\item[(0.2c)] The sets $\mathrm{Kub}_a(z)$ and the closed cubes $\{c:|c-z|\le M\}$, $\{c:|c-z|\le L\cdot M\}$ are products of three intervals of $\eta_k\mathbb{Z}$.
Hence such a set contains, together with two sites $c,c'$, all sites of the lattice box spanned by $c$ and $c'$,
and in particular the sites of the six paths $\boldsymbol{\Gamma}^{\pi}(c,c')$:
each coordinate of a point of such a path lies between the corresponding coordinates of $c$ and $c'$ (Section~I.1.1).
\end{itemize}
\emph{Parameter cubes} (\cite[Definition~5.8]{QEDI}).
The cubes of $\pi_k$ are the $M$-cubes $\mathrm{Kub}_M(z)$, $z\in M\mathbb{Z}^3$;
the cubes of $\pi'_k$ are the sets $\{c:|c-z|\le M\}$, $z\in M\mathbb{Z}^3$;
the cubes of $\pi''_k$ are the sets $\{c:|c-z|\le L\cdot M\}$, $z\in LM\mathbb{Z}^3$.
Their sup-diameters are $M-\eta_k$, $2M$ and $2L\cdot M$, respectively.

\subsubsection*{0.3. Unions of $M$-cubes; enlargements}

A union $Z$ of $M$-cubes of $\mathbb{L}_k$ is identified with the set of its sites,
and $\mathrm{Ctr}(Z)\subset M\mathbb{Z}^3$ denotes the set of the centres of its cubes.
\begin{itemize}
\item[(0.3a)] By (0.2a), an $M$-cube that has a site in common with $Z$ is one of the cubes of $Z$.
\end{itemize}
For an integer $a\ge0$, the set $Z^{\sim a}$ is the union of the $M$-cubes $\mathrm{Kub}_M(z')$ with $z'\in M\mathbb{Z}^3$ and $|z'-z|\le a\cdot M$ for some $z\in\mathrm{Ctr}(Z)$ (Section~I.1.1);
on $T_i$ it is defined in the same way, with $d_i$ in the place of $|\cdot|$.
\begin{itemize}
\item[(0.3b)] If $Z\subset Z'$, then $Z^{\sim a}\subset Z'^{\sim a}$.
For $v\in M\mathbb{Z}^3$ one has $(Z+v)^{\sim a}=Z^{\sim a}+v$.
Moreover $(Z^{\sim a})^{\sim a'}=Z^{\sim(a+a')}$.
For the last identity, note that $|z''-z'|\le a'\cdot M$ and $|z'-z|\le a\cdot M$ imply $|z''-z|\le(a+a')\cdot M$.
Conversely, let $|z''-z|\le(a+a')\cdot M$, and for $\mu=1,2,3$ let $z'_\mu$ be the point between $z_\mu$ and $z''_\mu$ whose distance from $z_\mu$ is $\min\left(|z''_\mu-z_\mu|,a\cdot M\right)$.
Then $z'\in M\mathbb{Z}^3$, $|z'-z|\le a\cdot M$ and $|z''-z'|\le a'\cdot M$.
\item[(0.3c)] If $c'\in Z$ and $|c-c'|\le a\cdot M$, then $c\in Z^{\sim a}$.
Indeed, let $c'\in\mathrm{Kub}_M(z)$ and $c\in\mathrm{Kub}_M(z')$;
then $z\in\mathrm{Ctr}(Z)$ by (0.3a), $|z'-z|<\frac{M}{2}+a\cdot M+\frac{M}{2}$, and $z'-z\in M\mathbb{Z}^3$, so that $|z'-z|\le a\cdot M$.
In particular, since $\eta_k\le1\le M$, a bond of $\mathbb{L}_k$ or of $\mathbb{Z}^3$ that meets $Z$ is contained in $Z^{\sim1}$,
and if $x$ lies over $Z$, then $c_{3}(x)=c+\eta_k\,ge_1$ lies in $Z^{\sim1}$.
\end{itemize}

\subsubsection*{0.4. Boxes}

Let $Z$ be a non-empty finite union of $M$-cubes of $\mathbb{L}_k$, and for $\mu=1,2,3$ put
\[
\mathrm{lo}_\mu:=\min\left\{z_\mu:\ z\in\mathrm{Ctr}(Z)\right\},\qquad
\mathrm{hi}_\mu:=\max\left\{z_\mu:\ z\in\mathrm{Ctr}(Z)\right\},\qquad
\mathrm{sid}_\mu:=\mathrm{hi}_\mu-\mathrm{lo}_\mu+M.
\]
\begin{itemize}
\item[(0.4a)] One has
\[
\mathrm{box}(Z)=\left\{c\in\mathbb{L}_k:\ \mathrm{lo}_\mu-\frac{M}{2}<c_\mu<\mathrm{hi}_\mu+\frac{M}{2}\ \text{for}\ \mu=1,2,3\right\};
\]
its sides are the numbers $\mathrm{sid}_\mu$, and $\mathrm{diam}(Z)=\max_\mu\mathrm{sid}_\mu$.
Indeed, the set on the right is a product of three intervals,
and by (0.2a) it is the union of the $M$-cubes $\mathrm{Kub}_M(z)$ with $\mathrm{lo}_\mu\le z_\mu\le\mathrm{hi}_\mu$ for $\mu=1,2,3$;
in particular it contains $Z$.
On the other hand, a product of three intervals that contains $Z$ contains, in the coordinate $\mu$, the numbers $\mathrm{lo}_\mu-\frac{M-\eta_k}{2}$ and $\mathrm{hi}_\mu+\frac{M-\eta_k}{2}$,
which are coordinates of sites of $Z$, and hence all points of $\eta_k\mathbb{Z}$ between them.
\item[(0.4b)] One has
\[
\mathrm{box}(Z)^{\sim a}=\left\{c:\ \mathrm{lo}_\mu-\frac{M}{2}-a\cdot M<c_\mu<\mathrm{hi}_\mu+\frac{M}{2}+a\cdot M\right\},
\]
a box of $M$-cubes with the sides $\mathrm{sid}_\mu+2a\cdot M$.
Two sites $u,u'$ of a box with the sides $\mathrm{sid}_\mu$ satisfy $|u_\mu-u'_\mu|\le\mathrm{sid}_\mu-\eta_k$;
indeed, by (0.2a) the extreme values of the coordinate $\mu$ on the box are $\mathrm{lo}_\mu-\frac{M-\eta_k}{2}$ and $\mathrm{hi}_\mu+\frac{M-\eta_k}{2}$.
\item[(0.4c)] Let $Z\subset Z'$, both non-empty and finite.
Then $\mathrm{lo}_\mu(Z')\le\mathrm{lo}_\mu(Z)$ and $\mathrm{hi}_\mu(Z)\le\mathrm{hi}_\mu(Z')$;
hence $\mathrm{box}(Z)\subset\mathrm{box}(Z')$ and $\mathrm{diam}(Z)\le\mathrm{diam}(Z')$.
Moreover $Z^{\sim a}\subset\mathrm{box}(Z)^{\sim a}$ by (0.3b).
\item[(0.4d)] A box of $M$-cubes is connected through faces:
any two of its cubes are joined by a chain of cubes of the box in which consecutive centres differ by $+Me_\mu$ or by $-Me_\mu$;
one brings the three coordinates to their final values one after the other.
\end{itemize}

\subsubsection*{0.5. Covering maps, indices, lifts}

Let $\mathrm{pr}_i:\mathbb{L}_k\to T_i=\mathbb{L}_k/\ell_i\mathbb{Z}^3$ be the quotient map.
\begin{itemize}
\item[(0.5a)] The sup distance of the torus is given by
\[
d_i(\mathrm{pr}_i c,\mathrm{pr}_i c')=\min_{w\in\mathbb{Z}^3}|c-c'-\ell_i w|.
\]
\item[(0.5b)] \emph{Indices.}
The covering indices of Proposition~27.1 form sets $\mathrm{Ind}$, one for each sort, on which $LM\mathbb{Z}^3$ acts.
Every index $t$ has a support $\operatorname{supp}t$, a finite non-empty set of sites of $\mathbb{L}_k$, and $\operatorname{supp}(t+v)=\operatorname{supp}t+v$.
The sorts, the action and the supports are the following.

\noindent
\begin{tabular}{p{2.5cm}p{3.1cm}p{2.6cm}p{2.0cm}}
\hline
sort & index $t\mapsto t+v$ & $\operatorname{supp}t$ & $\mathrm{sdiam}(\operatorname{supp}t)$\\
\hline
site of $\mathbb{L}_k$; unit site & $c\mapsto c+v$; $y\mapsto y+v$ & $\{c\}$; $\{y\}$ & $0$\\
bond of $\mathbb{L}_k$; unit bond & $(c,c')\mapsto(c+v,c'+v)$ & $\{c,c'\}$ & $\eta_k$; $1$\\
basis element & $\Upsilon$ (a bond function on $\mathbb{Z}^3$) $\mapsto$ its translate & the end points of the bonds $b'$ with $\Upsilon(b')\ne0$ & at most $L-1$\\
fermion site & $x=(c,g)\mapsto(c+v,g)$ & $\{c\}$ & $0$\\
pair of $\mathrm{Pzt}_k$ & $(x,x')\mapsto(x+v,x'+v)$ & $\{c,c'\}$ & at most $M$\\
block $\mathrm{Bl}_{j}(z)$, $0\le j\le k+1$ & $(\text{sort},z)\mapsto(\text{sort},z+v)$ & $\mathrm{Kub}_{L^j\eta_k}(z)$ & $L^j\eta_k-\eta_k<L$\\
$M$-cube; $LM$-cube & $(\text{sort},z)\mapsto(\text{sort},z+v)$ & $\mathrm{Kub}_M(z)$; $\mathrm{Kub}_{LM}(z)$ & $M-\eta_k$; $L\cdot M-\eta_k$\\
parameter cube of $\pi_k$, $\pi'_k$, $\pi''_k$ & $(\text{sort},z)\mapsto(\text{sort},z+v)$ & its sites (see 0.2) & $M-\eta_k$; $2M$; $2L\cdot M$\\
\hline
\end{tabular}

For the basis elements recall that the family $\{\Upsilon\}$ on $\mathbb{Z}^3$ consists of the translates $e^{y}_{n}$ of the $e^{0}_{n}$ by the points $y\in L\mathbb{Z}^3$
and of the functions $\delta_b$, $b$ a positively oriented face bond that is not central;
here $\{e^{0}_{n}\}$ is the orthonormal basis of $\ker\tau_0$ fixed in \cite[Section~5.2.5]{QEDI}, whose elements are denoted $e^{0}_{i}$ there;
it has $\dim\ker\tau_0=2L^3-3L^2+1$ elements, and the same basis is taken on $T_1$ and on $T_2$ \cite[Proposition~15.21(i)]{QEDI}.
The family $\{\Upsilon\}$ is invariant under $L\mathbb{Z}^3$,
and $\operatorname{supp}\Upsilon$ lies in one $L$-block or is one bond
(\cite[Section~5.2.5]{QEDI}; \cite[Proposition~5.35]{QEDI}, \cite[Proposition~15.21(ii)]{QEDI}).
The grids $L^j\eta_k\mathbb{Z}^3$ ($j\le k+1$), $M\mathbb{Z}^3$ and $LM\mathbb{Z}^3$ are invariant under $LM\mathbb{Z}^3$,
because $L\cdot M/(L^j\eta_k)=L^{1+k-j}\cdot M$ is an integer.
Consequently every support has sup-diameter at most $2L\cdot M$.

The indices of $T_i$ of a given sort are the orbits of $\ell_i\mathbb{Z}^3$ in $\mathrm{Ind}$, and $\mathrm{pr}_i(t)$ is the orbit of $t$
(Proposition~26.1; Section~I.1.7 for pairs; Section~I.1.9; for blocks and cubes, the definitions in Proposition~27.1, by which the orbit of $(\text{sort},z)$ is $(\text{sort},\mathrm{pr}_i(z))$).
Hence
\[
\mathrm{pr}_i(t)=\mathrm{pr}_i(t')\ \text{ if and only if }\ t'=t+v\ \text{ for a vector }v\in\ell_i\mathbb{Z}^3,
\qquad
\operatorname{supp}\mathrm{pr}_i(t):=\mathrm{pr}_i(\operatorname{supp}t).
\]
A bond, a unit bond or a pair of $T_i$ is at the same time an ordered pair of sites of $T_i$, resp.\ of $T^{\mathrm{f}}_i$, and the orbit is determined by this ordered pair:
if $c_1=c+v$ and $c_1'=c'+v'$ with $v,v'\in\ell_i\mathbb{Z}^3$, $|c-c'|\le M$ and $|c_1-c_1'|\le M$,
then $|v-v'|\le2M<\ell_i$, hence $v=v'$.
A \emph{lift} of an index $\mathrm{tq}$ of $T_i$ is a covering index $t$ with $\mathrm{pr}_i(t)=\mathrm{tq}$.
An index is \emph{contained in} a set of sites if its support is a subset of this set.
\item[(0.5c)] The set of the sites of a block or of a cube of $T_i$ is the image of the set of the sites of any of its lifts:
\[
\left\{\mathrm{pr}_i c:\ d_i(\mathrm{pr}_i c,\mathrm{pr}_i z)<\frac{a}{2}\right\}=\mathrm{pr}_i\left(\mathrm{Kub}_a(z)\right),
\]
and likewise with $\le M$ and $\le L\cdot M$ in the place of $<\frac{a}{2}$ for $\pi'_k$ and $\pi''_k$.
Indeed, by (0.5a) the inequality $d_i(\mathrm{pr}_i c,\mathrm{pr}_i z)<\frac{a}{2}$ means that $c-\ell_i w\in\mathrm{Kub}_a(z)$ for some $w$.
\item[(0.5d)] The side $\ell_i$ is a multiple of $L\cdot M$.
Therefore $\mathrm{pr}_i$ maps the $M$-cubes ($LM$-cubes) of $\mathbb{L}_k$ onto the $M$-cubes ($LM$-cubes) of $T_i$ (0.5c),
and the $M$-cubes ($LM$-cubes) of $T_i$ form a partition of the set of the sites of $T_i$.
Indeed, let $a=M$ or $a=L\cdot M$, and let $\mathrm{pr}_i(c)$ lie in $\mathrm{pr}_i(\mathrm{Kub}_a(z))$ and in $\mathrm{pr}_i(\mathrm{Kub}_a(z'))$;
then $c$ lies in $\mathrm{Kub}_a(z+v)$ and in $\mathrm{Kub}_a(z'+v')$ for some $v,v'\in\ell_i\mathbb{Z}^3$,
so that $z+v=z'+v'$ by (0.2a), and the two cubes of $T_i$ coincide.
The $M$-cube of $T_i$ that contains $\mathrm{pr}_i(c)$ is the image of the $M$-cube of $c$.
\item[(0.5e)] Let $\tilde{S}$ be a set of sites of $\mathbb{L}_k$.
Every bond of $T_i$ that meets $\mathrm{pr}_i(\tilde{S})$ is of the form $\mathrm{pr}_i(b)$ with a bond $b$ of $\mathbb{L}_k$ that meets $\tilde{S}$.
Indeed, the bond is the image of a bond $b_0$ of $\mathbb{L}_k$ (Section~I.1.1), and one of its end points is $\mathrm{pr}_i(c)$ with $c\in\tilde{S}$;
the corresponding end point $c_0$ of $b_0$ satisfies $c-c_0\in\ell_i\mathbb{Z}^3$, and $b:=b_0+(c-c_0)$ has the required properties.
\item[(0.5f)] For a bond function $B$ on $T_1$ and a bond $b$ of $\mathbb{L}_k$,
\[
(\mathrm{pj}^{*}B)(\mathrm{pr}_2 b)=B(\mathrm{pj}\,\mathrm{pr}_2 b)=B(\mathrm{pr}_1 b),
\]
since $\mathrm{pj}^{*}$ is the inclusion of the $\ell_1\mathbb{Z}^3$-periodic objects and $\mathrm{pj}\circ\mathrm{pr}_2=\mathrm{pr}_1$ (Proposition~26.1).
\end{itemize}

\paragraph{Level 0.}
At $k=0$ one has $\eta_0=1$, $\mathbb{L}_0=\mathbb{Z}^3$, $\Delta_y=\{y\}$ and $\mathrm{Bl}_1(y')=B(y')$.
All statements of 0.2--0.5 hold as written, since $a/\eta_0=a$ is odd.

\subsection*{1. Proposition 27.1}

\subsubsection*{Proof of Proposition 27.1}

\begin{proof}
Let $\mathrm{lo}_\mu$, $\mathrm{hi}_\mu$ and $\mathrm{sid}_\mu$ be the numbers of 0.4 for $Z=\tilde{X}$.

\textbf{(a)} By (0.4b) with $a=6d_U$,
\[
N(\tilde{X})=\left\{c\in\mathbb{L}_k:\ \mathrm{lo}_\mu-\frac{M}{2}-6d_U\cdot M<c_\mu<\mathrm{hi}_\mu+\frac{M}{2}+6d_U\cdot M\ \ (\mu=1,2,3)\right\}.
\]
This is a box of $M$-cubes, and its sides satisfy
\[
\mathrm{sid}_\mu+12d_U\cdot M\le\mathrm{diam}(\tilde{X})+12d_U\cdot M\le\ell_1-14d_U\cdot M+12d_U\cdot M=\ell_1-2d_U\cdot M.
\]
The second sentence of (0.4b) therefore gives
\begin{equation}
|u-u'|\le\ell_1-2d_U\cdot M-\eta_k\qquad\text{for all sites }u,u'\text{ of }N(\tilde{X}).\tag{1.1}
\end{equation}
Let $\tilde{X}'$ be a non-empty union of $M$-cubes contained in $\tilde{X}$; it is finite.
By (0.4c), $\mathrm{diam}(\tilde{X}')\le\mathrm{diam}(\tilde{X})\le\ell_1-14d_U\cdot M$,
so that $\tilde{X}'$ is admissible, and $\mathrm{box}(\tilde{X}')\subset\mathrm{box}(\tilde{X})$.
By (0.3b) it follows that
\[
N(\tilde{X}')=\mathrm{box}(\tilde{X}')^{\sim6d_U}\subset\mathrm{box}(\tilde{X})^{\sim6d_U}=N(\tilde{X}).
\]

\textbf{(b)} Let $0\ne v\in\ell_1\mathbb{Z}^3$; then $|v|\ge\ell_1$.
If $u$ and $u+v$ were both sites of $N(\tilde{X})$, then (1.1) would give $|v|\le\ell_1-2d_U\cdot M-\eta_k<\ell_1$.
Hence $N(\tilde{X})$ and $N(\tilde{X})+v$ are disjoint.
Next let $u,u'$ be sites of $N(\tilde{X})$ with $\mathrm{pr}_i(u)=\mathrm{pr}_i(u')$.
The vector $v:=u'-u$ lies in $\ell_i\mathbb{Z}^3$, hence in $\ell_1\mathbb{Z}^3$ by (0.1a),
and $u'$ belongs both to $N(\tilde{X})$ and to $N(\tilde{X})+v$; therefore $v=0$.
Finally let $t,t'$ be indices contained in $N(\tilde{X})$ with $\mathrm{pr}_i(t)=\mathrm{pr}_i(t')$.
By (0.5b), $t'=t+v$ with $v\in\ell_i\mathbb{Z}^3$.
For a point $u$ of $\operatorname{supp}t$, which exists because the support is not empty, both $u$ and $u+v$ are sites of $N(\tilde{X})$;
so $v=0$, and $t'=t$.

\textbf{(c)} Let $u,u'$ be sites of $N(\tilde{X})$ and let $0\ne w\in\mathbb{Z}^3$.
By (0.1b) and (1.1),
\begin{equation}
|u-u'-\ell_i w|\ \ge\ \ell_i\cdot|w|-|u-u'|\ \ge\ \ell_1-\left(\ell_1-2d_U\cdot M-\eta_k\right)=2d_U\cdot M+\eta_k.\tag{1.2}
\end{equation}
By (0.5a), $d_i(\mathrm{pr}_i u,\mathrm{pr}_i u')$ is the least of the numbers $|u-u'-\ell_i w|$, $w\in\mathbb{Z}^3$.
If it is less than $2d_U\cdot M$, then by (1.2) it is not attained at any $w\ne0$, and so it equals $|u-u'|$.
The representatives of the class of $u-u'$ modulo $\ell_i\mathbb{Z}^3$ are the vectors $u-u'-\ell_i w$, $w\in\mathbb{Z}^3$,
and by (1.2) those with $w\ne0$ have sup norm greater than $2d_U\cdot M$.

\textbf{(d)} \emph{Existence.}
Let $\mathrm{tq}$ be an index of $T_i$ that is contained in $\mathrm{pr}_i(\tilde{S})$, let $t_0$ be a lift of $\mathrm{tq}$, and let $u_0'$ be a point of $\operatorname{supp}t_0$.
By (0.5b), $\mathrm{pr}_i(u_0')$ lies in $\operatorname{supp}\mathrm{tq}$, hence in $\mathrm{pr}_i(\tilde{S})$;
thus there is a site $u_0\in\tilde{S}$ with $u_0-u_0'\in\ell_i\mathbb{Z}^3$.
The index $t:=t_0+(u_0-u_0')$ is then a lift of $\mathrm{tq}$, and $u_0\in\operatorname{supp}t$.
Let $u'$ be an arbitrary point of $\operatorname{supp}t$.
Since $\mathrm{pr}_i(u')$ lies in $\mathrm{pr}_i(\tilde{S})$, there are a site $u\in\tilde{S}$ and a vector $w\in\mathbb{Z}^3$ with $u=u'+\ell_i w$.
The sites $u$ and $u_0$ belong to $N(\tilde{X})$, and by (0.5b)
\[
|u-u_0-\ell_i w|=|u'-u_0|\le\mathrm{sdiam}(\operatorname{supp}t)\le2L\cdot M<(4L+2)\cdot M=2d_U\cdot M.
\]
By (1.2) this excludes $w\ne0$.
Hence $u'=u$ lies in $\tilde{S}$, and $t$ is contained in $\tilde{S}$.

\emph{Uniqueness.}
Two lifts of $\mathrm{tq}$ that are contained in $\tilde{S}$ are contained in $N(\tilde{X})$ and have the same image; by (b) they are equal.

\emph{Pairs.}
A pair $\mathrm{tq}$ of $\mathrm{Pz}_k(T_i)=\mathrm{Pzt}_k/\ell_i\mathbb{Z}^3$ is a class of elements of $\mathrm{Pzt}_k$, and its lifts are the elements of this class;
they constitute what Section~I.1.7 calls its shortest lift.
By definition, $d_X(\mathrm{tq})$ and the six paths of $\mathrm{tq}$ are $d_X$ and the images under $\mathrm{pr}_i$ of the six paths of a lift
(Section~I.1.7; Section~I.1.1, paths on the torus).
Hence the lift $(x,x')$ of $\mathrm{tq}$ that is contained in $\tilde{S}$ is an element of $\mathrm{Pzt}_k$ with $d_X(x,x')=d_X(\mathrm{tq})$,
and $\mathrm{pr}_i$ maps $\boldsymbol{\Gamma}^{\pi}(c,c')$ onto the path of $\mathrm{tq}$ that belongs to the order $\pi$.
In terms of the two sites:
let $\mathrm{xq}\ne\mathrm{xq}'$ be sites of $T^{\mathrm{f}}_i$ over $\mathrm{pr}_i(\tilde{S})$ with $d_i(\pi(\mathrm{xq}),\pi(\mathrm{xq}'))\le M$,
and let $x=(c,g)$, $x'=(c',g')$ be their lifts over $\tilde{S}$;
then $|c-c'|=d_i(\pi(\mathrm{xq}),\pi(\mathrm{xq}'))\le M$ by (c), because $M<2d_U\cdot M$, and $x\ne x'$;
consequently $(x,x')$ belongs to $\mathrm{Pzt}_k$, and it is the lift contained in $\tilde{S}$ of the pair $(\mathrm{xq},\mathrm{xq}')$.

\emph{Enlargements.}
Since $\ell_i$ is a multiple of $M$, the map $\mathrm{pr}_i$ maps $M\mathbb{Z}^3$ onto the set of the centres of the $M$-cubes of $T_i$,
and $\mathrm{Kub}_M(z)$ onto the $M$-cube with the centre $\mathrm{pr}_i(z)$ (0.5c).
Let $\mathrm{Kub}_M(z')$ be a cube of $\tilde{X}^{\sim a}$, so that $|z'-z|\le a\cdot M$ for some $z\in\mathrm{Ctr}(\tilde{X})$.
Then $d_i(\mathrm{pr}_i z',\mathrm{pr}_i z)\le a\cdot M$ by (0.5a), and $\mathrm{pr}_i(\mathrm{Kub}_M(z'))$ is a cube of $(\mathrm{pr}_i\tilde{X})^{\sim a}$.
Conversely, let an $M$-cube of $T_i$ belong to $(\mathrm{pr}_i\tilde{X})^{\sim a}$:
its centre is $\mathrm{pr}_i(z'')$ with $z''\in M\mathbb{Z}^3$ and $d_i(\mathrm{pr}_i z'',\mathrm{pr}_i z)\le a\cdot M$ for some $z\in\mathrm{Ctr}(\tilde{X})$.
By (0.5a) there is a vector $w$ with $|z''-\ell_i w-z|\le a\cdot M$;
the point $z':=z''-\ell_i w$ lies in $M\mathbb{Z}^3$, $\mathrm{Kub}_M(z')$ is a cube of $\tilde{X}^{\sim a}$, and its image is the given cube.
This argument uses neither the admissibility of $\tilde{X}$ nor any bound on $a$.

\textbf{(e)} If $(t_1,\dots,t_p)\in\mathrm{Rel}$, then $(\mathrm{pr}_i t_1,\dots,\mathrm{pr}_i t_p)\in\mathrm{Rel}_i$ by the definition of $\mathrm{Rel}_i$.
Conversely, let $(\mathrm{pr}_i t_1,\dots,\mathrm{pr}_i t_p)\in\mathrm{Rel}_i$.
There is a tuple $(t'_1,\dots,t'_p)\in\mathrm{Rel}$ with $\mathrm{pr}_i(t'_j)=\mathrm{pr}_i(t_j)$, that is, $t'_j=t_j+v_j$ with $v_j\in\ell_i\mathbb{Z}^3$ (0.5b).
Since $\mathrm{Rel}$ is invariant under the common translation by $-v_1$, which is a vector of $\ell_1\mathbb{Z}^3$ by (0.1a), one may assume $v_1=0$.
Let $j\ge2$, write $v_j=\ell_i w_j$, and let $u$ be a point of $\operatorname{supp}t_1$ and $u'$ a point of $\operatorname{supp}t_j$.
The points $u$ and $u'+v_j$ belong to the union of the supports of a tuple of $\mathrm{Rel}$, so that $|u-u'-\ell_i w_j|<2d_U\cdot M$.
As $u$ and $u'$ are sites of $N(\tilde{X})$, (1.2) gives $w_j=0$.
Thus $t'_j=t_j$ for all $j$, and $(t_1,\dots,t_p)\in\mathrm{Rel}$.

\textbf{(f)} \emph{How the objects of $T_i$ are defined.}
Three remarks are read off the definitions.
\begin{itemize}
\item[(i)] The maps $x\mapsto x^{(j)}$ ($j=0,1,2,3$; see (I.1.2)), $x\mapsto\pi(x)$, and the map that assigns to an $M$-cube its centre are given on $T_i$ by the formulas that define them on the covering lattice,
and these formulas commute with the translations by $\ell_i\mathbb{Z}^3$ (Section~I.1.3, (I.s1)).
Hence $\mathrm{fm}(\mathrm{pr}_i t)=\mathrm{pr}_i(\mathrm{fm}(t))$ for each of these maps $\mathrm{fm}$.
\item[(ii)] One has $\operatorname{supp}\mathrm{pr}_i(t)=\mathrm{pr}_i(\operatorname{supp}t)$,
and the set of the sites of a block or of a cube of $T_i$ is the image of the set of the sites of a lift ((0.5b), (0.5c)).
Consequently, for covering indices $t_1,t_2$, the supports of $\mathrm{pr}_i(t_1)$ and $\mathrm{pr}_i(t_2)$ have a common site
if and only if $\operatorname{supp}t_1$ and $\operatorname{supp}(t_2+v)$ have a common site for some $v\in\ell_i\mathbb{Z}^3$;
indeed, $\mathrm{pr}_i(u_1)=\mathrm{pr}_i(u_2)$ means $u_1=u_2+v$ with such a $v$.
\item[(iii)] By their definitions, the following objects of $T_i$ are the images of the corresponding covering objects:
the bonds of $T_i$ and of $T^0_i$ (Section~I.1.1: a bond is an ordered pair of sites whose difference, a class modulo $\ell_i\mathbb{Z}^3$, contains a vector of Euclidean norm $\eta_k$, resp.\ $1$);
the pairs of $\mathrm{Pz}_k(T_i)$ (Section~I.1.7);
the six paths of a pair (Section~I.1.1: they are formed on a lift with $|c-c'|<\frac{\ell_i}{2}$, and the lifts of a pair in $\mathrm{Pzt}_k$ differ by translations by $\ell_i\mathbb{Z}^3$);
the basis elements of $T_i$ and the bonds of their supports
(\cite[Section~5.2.5]{QEDI}: for every $n$ and every $y\in L\mathbb{Z}^3/\ell_i\mathbb{Z}^3$ the basis element $e^{y}_{n}$ of $T_i$ is the translate of $e^{0}_{n}$ by $y$,
that is, the image under $\mathrm{pr}_i$ of the element $e^{\tilde{y}}_{n}$ of the family on $\mathbb{Z}^3$ for any $\tilde{y}\in L\mathbb{Z}^3$ over $y$;
the block $B(0)$ embeds in $T^0_i$ because $L-1<\ell_i$).
\end{itemize}
By the definition of $\mathrm{Rel}_i$, a tuple $(\mathrm{tq}_1,\dots,\mathrm{tq}_p)$ belongs to $\mathrm{Rel}_i$ if and only if the $\mathrm{tq}_j$ have lifts $t_j$ with $(t_1,\dots,t_p)\in\mathrm{Rel}$.
In the following list each relation $\mathrm{Rel}$ is the set of the tuples of covering indices with the property stated first;
the entry ``Diameter'' bounds the sup-diameter of the union of the supports of a tuple of $\mathrm{Rel}$,
and the entry ``On $T_i$'' says how the relation on $T_i$ is defined and why it is $\mathrm{Rel}_i$.
\begin{description}
\item[No.~1.] The pairs $(c,c')$ with $c,c'\in\mathbb{L}_k$ and $|c-c'|_2=\eta_k$; the pairs $(y,y')$ with $y,y'\in\mathbb{Z}^3$ and $|y-y'|_2=1$.
Diameter: $\eta_k$, resp.\ $1$.
On $T_i$: ``is a bond''; by (iii).
\item[No.~2.] For each $j=0,1,2,3$ one relation: the pairs $(x,x')$ with $x'=x^{(j)}$.
Diameter: $0$ for $j\le2$, and $\eta_k$ for $j=3$.
On $T_i$: $\mathrm{xq}'=\mathrm{xq}^{(j)}$; by (i), $\mathrm{Rel}_i=\left\{(\mathrm{pr}_i x,(\mathrm{pr}_i x)^{(j)})\right\}$, and $\mathrm{pr}_i$ is onto.
\item[No.~3.] The pairs $(x,x')$ with $\pi(x)=\pi(x')$.
Diameter: $0$.
On $T_i$: $\pi(\mathrm{xq})=\pi(\mathrm{xq}')$;
if both are equal to $\mathrm{pr}_i(c)$, then $\mathrm{xq}=\mathrm{pr}_i(c,g)$ and $\mathrm{xq}'=\mathrm{pr}_i(c,g')$; the converse holds by (i).
\item[No.~4.] The pairs $(x,x')\in\mathrm{Pzt}_k$, that is, $x\ne x'$ and $\max\left(|c-c'|,\eta_k\right)\le M$.
Diameter: at most $M$.
On $T_i$: ``$(\mathrm{xq},\mathrm{xq}')$ is a pair of $\mathrm{Pz}_k(T_i)$''; by (iii).
\item[No.~5.] For each of the six orders $\pi$ one relation: the triples $(b,x,x')$ with $(x,x')\in\mathrm{Pzt}_k$ and $b$ a bond of $\boldsymbol{\Gamma}^{\pi}(c,c')$.
Diameter: at most $M$, by (0.2c).
On $T_i$: ``$b$ lies on the path of the order $\pi$ of the pair''; by (iii).
\item[No.~6.] The pairs $(c,t)$ with $t$ one of $\Delta_y$, $\mathrm{Bl}_{j}(z)$, an $M$-cube, an $LM$-cube, a parameter cube, and $c\in\operatorname{supp}t$.
Diameter: $1-\eta_k$; less than $L$; less than $M$; less than $L\cdot M$; at most $2L\cdot M$, in the order of the five kinds of $t$.
On $T_i$: membership; by (ii) with $t_1=c$.
\item[No.~7.] The pairs $(y,y')$ with $y'\in L\mathbb{Z}^3$ and $y\in B(y')$.
Diameter: at most $\frac{L-1}{2}$.
On $T_i$: membership in the $L$-block, that is, $d_i(\mathrm{pr}_i y,\mathrm{pr}_i y')<\frac{L}{2}$; by (0.5a).
\item[No.~8.] The pairs $(z,t)$ with $t=\mathrm{Kub}_M(z)$.
Diameter: less than $M$.
On $T_i$: by (i).
\item[No.~9.] The pairs $(t,t')$ with $t=\mathrm{Kub}_M(z)$, $t'=\mathrm{Kub}_M(z')$, and $z'-z=+Me_\mu$ or $z'-z=-Me_\mu$ for some $\mu$.
Diameter: $2M-\eta_k$.
On $T_i$: common face; see below.
\item[No.~10.] The pairs $(t,t')$ of linked parameter cubes of one kind (\cite[Definition~5.8]{QEDI}):
for $\pi_k$, as in no.~9; for $\pi'_k$ and $\pi''_k$, $\operatorname{supp}t$ and $\operatorname{supp}t'$ have a common site.
Diameter: $2M-\eta_k$; at most $4M$; at most $4L\cdot M$.
On $T_i$: for $\pi_k$, as in no.~9; for $\pi'_k$ and $\pi''_k$, a common lattice site; by (ii).
\item[No.~11.] The pairs $(t,t')$ with $t$ a parameter cube and $t'$ an $M$-cube that have a common site.
Diameter: at most $2L\cdot M+2M-2\eta_k$.
On $T_i$: ``meets''; by (ii).
\item[No.~12.] The pairs $(b',\Upsilon)$ with $b'$ a unit bond such that $\Upsilon(b')\ne0$.
Diameter: at most $L-1$.
On $T_i$: by (iii).
\end{description}
The bounds of nos.~10 and~11 are obtained as follows.
For no.~10: if $u$ is a common site, every site of the union has sup-distance at most $2M$, resp.\ $2L\cdot M$, from $u$.
For no.~11: if $u$ is a common site, the sites of $t'$ have sup-distance at most $M-\eta_k$ from $u$, and $u$ has distance at most $L\cdot M$ from the centre of $t$.

\emph{No.~9.}
By Section~I.1.1, two $M$-cubes of $T_i$ have a common face if their centres differ by $+Me_\mu$ or by $-Me_\mu$ modulo $\ell_i\mathbb{Z}^3$,
that is, if they have lifts $\mathrm{Kub}_M(z)$, $\mathrm{Kub}_M(z')$ with $z'-z=+Me_\mu$ or $z'-z=-Me_\mu$; this relation is $\mathrm{Rel}_i$.
It is the relation ``adjacent'' of \cite[Section~13.3]{QEDI}:
on $T_i$ two $M$-cubes are adjacent if the corresponding closed cubes of the continuum torus $\mathbb{R}^3/\ell_i\mathbb{Z}^3$ have a common face,
the closed cubes of the torus being the images of the closed cubes $\left\{y\in\mathbb{R}^3:\ |y-z|\le\frac{M}{2}\right\}$, $z\in M\mathbb{Z}^3$.
The image of a common face is a common face, and $\mathrm{Kub}_M(z)$, $\mathrm{Kub}_M(z+Me_\mu)$ have different images because $M<\ell_i$.
Conversely, let two different $M$-cubes with the centres $\mathrm{pr}_i(z)$ and $\mathrm{pr}_i(z')$ be adjacent.
The midpoint of their common face has the lift $z+\frac{M}{2}\,\mathrm{ae}$ with a vector $\mathrm{ae}\in\{e_\mu,-e_\mu\}$,
and this point lies in a lift $\left\{y:\ |y-z'-v|\le\frac{M}{2}\right\}$, $v\in\ell_i\mathbb{Z}^3$, of the second closed cube.
A point $z'+v$ of $M\mathbb{Z}^3$ with $\left|z+\frac{M}{2}\,\mathrm{ae}-z'-v\right|\le\frac{M}{2}$ is either $z$ or $z+M\,\mathrm{ae}$:
the coordinates $\nu\ne\mu$ of $z'+v-z$ are multiples of $M$ of modulus at most $\frac{M}{2}$, hence $0$,
and the coordinate $\mu$ of $z'+v-z$ lies between $0$ and the coordinate $\mu$ of $M\,\mathrm{ae}$.
The first case is excluded, since the two cubes are different.
Hence the lifts $\mathrm{Kub}_M(z)$ and $\mathrm{Kub}_M(z'+v)$ are in $\mathrm{Rel}$.

\emph{Range.}
The largest of the bounds in the list is $4L\cdot M$, and $4L\cdot M<(4L+2)\cdot M=2d_U\cdot M$.
By (0.1b),
\[
(4L+2)\cdot M<\frac{(31L+16)\cdot M}{2}<\frac{\ell_i}{2}.
\]

\emph{Invariance.}
Each relation $\mathrm{Rel}$ of the list is defined in terms of differences of coordinates, the labels, the maps of (i),
the grids $\eta_k\mathbb{Z}^3$, $\mathbb{Z}^3$, $L^j\eta_k\mathbb{Z}^3$ ($j\le k+1$), $L\mathbb{Z}^3$, $M\mathbb{Z}^3$, $LM\mathbb{Z}^3$, and the family $\{\Upsilon\}$.
All of these are invariant under the translations by $LM\mathbb{Z}^3$, and these translations do not move the labels (0.5b).
Therefore each $\mathrm{Rel}$ is invariant under $LM\mathbb{Z}^3$, in particular under $\ell_1\mathbb{Z}^3$ (0.1a).
\end{proof}

\paragraph{Level 0.}
Nothing changes at $k=0$:
$\eta_0=1$; the two relations of no.~1 coincide;
in no.~4 every pair has $d_X\ge1$, and a pair inside a cluster ($c=c'$, $g\ne g'$) has the support $\{c\}$.

\paragraph{Uniformity.}
The clauses of Proposition~27.1 hold for all data (TT): for every $(\bar{m},e,R,N)$ with (FL) and $R<R_{\mathrm{top}}$, and at every level $k\le k_m$.
They hold uniformly in these data.
Indeed, Proposition~27.1 introduces no constant and proves no bound;
its clauses are inclusions, injectivity statements, identities and equivalences.
The constants that occur in them are $L$, $M$ and $d_U=2L+1$ of Section~I.8.2,
and these do not depend on the data (TT), by the order in which the constants are chosen (Section~I.8.1).

\subsection*{2. Proposition 27.2}

\subsubsection*{Proof of Proposition 27.2}

For pairs of $\mathrm{Pzt}_k$ and unions $\tilde X$ of $M$-cubes of $\mathbb{L}_k$, the expressions ``admissible for $\tilde X$'', ``of kind (I)'' and ``of kind (II)'' are understood with the wording of Section~I.1.7, in which $Z$ stands for an arbitrary union of $M$-cubes.

\begin{proof}[Proof of Proposition~27.2\,(a)]
\emph{Independence.}
By Proposition~27.1\,(b), no two distinct sites of $N(\tilde X)$ are congruent modulo $\ell_i\mathbb{Z}^3$.
Choose a set $\mathrm{Rep}_i$ of representatives of $\mathbb{L}_k$ modulo $\ell_i\mathbb{Z}^3$ that contains all sites of $N(\tilde X)$.
Since $\ell_i$ is an integer, $\mathbb{Z}^3$ is a union of classes modulo $\ell_i\mathbb{Z}^3$; consequently $\mathrm{Rep}_i\cap\mathbb{Z}^3$ is a set of representatives of $\mathbb{Z}^3$ modulo $\ell_i\mathbb{Z}^3$.
By (I.AP), the following elements are independent generators of the Grassmann algebra of $T_i$:
\begin{equation}\tag{2.0}
\begin{aligned}
&\psi(x,\beta,\omega)\ \text{with}\ c\in\mathrm{Rep}_i;\qquad
\chi(x,x',\beta,\omega)\ \text{with}\ (x,x')\in\mathrm{Pzt}_k\ \text{and}\ c\in\mathrm{Rep}_i;\\
&k\ge1\colon\ \Psi(y,\beta,\omega),\ W(y,\beta,\omega)\ \text{with}\ y\in\mathrm{Rep}_i\cap\mathbb{Z}^3;\qquad
k=0\colon\ W(x,\beta,\omega)\ \text{with}\ c\in\mathrm{Rep}_i;
\end{aligned}
\end{equation}
here $x=(c,g)$, $x'=(c',g')$, $g,g'\in O_{\mathrm{cub}}$, $\beta\in\{1,\dots,4\}$, $\omega\in\{0,1\}$, and for the generators $\chi$ the site $c=\pi(x)$ is the photon site of the \emph{first} point.
Let $\mathrm{Ix}(\tilde X)$ be the set of the indices of the generators named in (D2), namely:
the indices $(\psi;x,\beta,\omega)$ with $c\in N(\tilde X)$;
the indices $(\chi;x,x',\beta,\omega)$ with $(x,x')\in\mathrm{Pzt}_k$ and $c,c'\in N(\tilde X)$, among which are, over every photon site $c$ of $N(\tilde X)$, the $48\cdot47$ ordered pairs inside the cluster of $c$;
for $k\ge1$, the indices $(\Psi;y,\beta,\omega)$ and $(W;y,\beta,\omega)$ with $y\in N(\tilde X)\cap\mathbb{Z}^3$;
for $k=0$, the indices $(W;x,\beta,\omega)$ with $c\in N(\tilde X)$.
This set is the same for $i=1$ and for $i=2$.
As $N(\tilde X)\subset\mathrm{Rep}_i$, the generators with indices in $\mathrm{Ix}(\tilde X)$ are pairwise distinct members of the system (2.0), and therefore they are independent.
Thus $\mathrm{Gr}_i(\tilde X)$ is the Grassmann algebra over $\mathbb{C}$ with these generators, and the ordered monomials in them form a basis of it.
A bijection between the free generators of two Grassmann algebras extends to exactly one isomorphism of algebras.
For the bijection that matches the generators with the same index in $\mathrm{Ix}(\tilde X)$, this isomorphism is $\mathrm{idn}[\tilde X]$.

\emph{Monotonicity.}
Let $\tilde X$ be contained in an admissible set $\tilde X'$.
Proposition~27.1\,(a) gives $N(\tilde X)\subset N(\tilde X')$.
Hence $\mathrm{Ix}(\tilde X)\subset\mathrm{Ix}(\tilde X')$, the algebra $\mathrm{Gr}_i(\tilde X)$ lies in $\mathrm{Gr}_i(\tilde X')$, and $\mathrm{idn}[\tilde X']$ maps each generator with index in $\mathrm{Ix}(\tilde X)$ to the generator of $T_2$ with the same index.
Two homomorphisms that agree on the generators are equal; so $\mathrm{idn}[\tilde X]$ is the restriction of $\mathrm{idn}[\tilde X']$.

\emph{Bidegree and norms.}
Fix a total order on $\mathrm{Ix}(\tilde X)$ and, for $i=1,2$, a total order on the index set of the system (2.0) that extends it.
Every $F\in\mathrm{Gr}_1(\tilde X)$ can be written as
\[
F=\sum_I c_F(I)\cdot\phi(I),
\]
where $I$ ranges over the subsets of $\mathrm{Ix}(\tilde X)$ and $\phi(I)$ is the ordered product of the generators of $T_1$ with indices in $I$.
Then $\mathrm{idn}\,F=\sum_I c_F(I)\cdot\phi(I)$, with the generators of $T_2$ in the place of those of $T_1$.
So the coefficients of $F$ and of $\mathrm{idn}\,F$ with respect to the systems (2.0) are the same numbers, attached to the same subsets.
It follows that $\mathrm{idn}$ preserves the degree in each of the families $\psi$, $\chi$, $\Psi$, $W$ (in particular the bidegree in $(\psi,\chi)$) and the numbers of barred and of unbarred factors.
Let $I_1,\dots,I_4$ denote the parts of $I$ in the families $\psi$, $\chi$, $\Psi$, $W$.
By (I.5.1), every weight vector $\mathsf{w}=(\mathsf{w}_1,\dots,\mathsf{w}_4)$ satisfies
\[
\|\mathrm{idn}\,F\|_{\mathsf{w}}=\sum_I\left|c_F(I)\right|\cdot\mathsf{w}_1^{\#|I_1|}\cdot\mathsf{w}_2^{\#|I_2|}\cdot\mathsf{w}_3^{\#|I_3|}\cdot\mathsf{w}_4^{\#|I_4|}=\|F\|_{\mathsf{w}}.
\]
(The norm (I.5.1) depends neither on the orders nor on $\mathrm{Rep}_i$: by (I.AP), another choice multiplies every monomial by a sign.)

\emph{The three maps.}
Let $\Lambda_1$, $\Lambda_2$ have equal lifts on $N(\tilde X)$ in the sense of (D1).
Written in spinor columns, the maps $G_{\Lambda_i}$, $\mathrm{Cg}$, $\mathbf{K}$ of (I.T-G), (I.T-C), (I.T-K) in Section~I.4.5 act, for $i=1,2$, on the generators with indices in $\mathrm{Ix}(\tilde X)$ as follows.
The automorphism $G_{\Lambda_i}$ is given by
\begin{align*}
\psi(x)&\mapsto\exp\left(\mathrm{i}\cdot e_k\cdot\Lambda_i(\mathrm{pr}_i c)\right)\cdot\psi(x),&
\bar\psi(x)&\mapsto\exp\left(-\mathrm{i}\cdot e_k\cdot\Lambda_i(\mathrm{pr}_i c)\right)\cdot\bar\psi(x),\\
\chi(x,x')&\mapsto\exp\left(\mathrm{i}\cdot e_k\cdot\Lambda_i(\mathrm{pr}_i c)\right)\cdot\chi(x,x'),&
\bar\chi(x,x')&\mapsto\exp\left(-\mathrm{i}\cdot e_k\cdot\Lambda_i(\mathrm{pr}_i c)\right)\cdot\bar\chi(x,x'),
\end{align*}
and in the same way on $\Psi(y)$, $W(y)$ with $\Lambda_i(\mathrm{pr}_i y)$ if $k\ge1$, and on $W(x)$ with $\Lambda_i(\mathrm{pr}_i c)$ if $k=0$.
The automorphism $\mathrm{Cg}$ is given by
\begin{align*}
\psi(x)&\mapsto\mathsf{C}\cdot\bar\psi(x),&
\bar\psi(x)&\mapsto-\mathsf{C}\cdot\psi(x),\\
\chi(x,x')&\mapsto\mathsf{C}\cdot\bar\chi(x,x'),&
\bar\chi(x,x')&\mapsto-\mathsf{C}\cdot\chi(x,x'),
\end{align*}
and in the same way on $\Psi$, $W$.
The map $\mathbf{K}$ replaces every coefficient by its complex conjugate, and each of $\psi$, $\bar\psi$, $\chi$, $\bar\chi$ (and $\Psi$, $W$) at an index by $\varkappa$ times the same field at the same index.
In each of the three cases a generator with index in $\mathrm{Ix}(\tilde X)$ is sent to a linear combination of generators at the same site or at the same pair, whose indices therefore lie in $\mathrm{Ix}(\tilde X)$.
The coefficients are entries of $\mathsf{C}$, $-\mathsf{C}$, $\varkappa$, or the numbers $\exp\left(\pm\mathrm{i}\cdot e_k\cdot\Lambda_i(\mathrm{pr}_i c)\right)$.
By (TT) the number $e_k$ is the same on both tori, and $\Lambda_1(\mathrm{pr}_1 c)=\Lambda_2(\mathrm{pr}_2 c)$ for $c\in N(\tilde X)$.
Hence each of the three maps sends $\mathrm{Gr}_i(\tilde X)$ into itself, and on every generator of $\mathrm{Gr}_1(\tilde X)$ one has
\[
\mathrm{idn}\circ G_{\Lambda_1}=G_{\Lambda_2}\circ\mathrm{idn},\qquad
\mathrm{idn}\circ\mathrm{Cg}=\mathrm{Cg}\circ\mathrm{idn},\qquad
\mathrm{idn}\circ\mathbf{K}=\mathbf{K}\circ\mathrm{idn}.
\]
In the first two identities both sides are homomorphisms of algebras; in the third both sides are antilinear ring homomorphisms, $\mathrm{idn}$ being linear.
Consequently the three identities hold on all of $\mathrm{Gr}_1(\tilde X)$.
\end{proof}

\begin{proof}[Proof of Proposition~27.2\,(b)]
\emph{Sites.}
A fermion site $\mathrm{xq}$ of $T^{\mathrm{f}}_i$ over $X=\mathrm{pr}_i(\tilde X)$ is an index contained in $\mathrm{pr}_i(\tilde X)$.
By Proposition~27.1\,(d) with $\tilde S=\tilde X$, it has exactly one lift over $\tilde X$.

\emph{Pairs of kind (I).}
Let the pair $\mathrm{tq}$ of $\mathrm{Pz}_k(T_i)$ be of kind (I) for $X$: its two photon sites and its six paths lie in $X$.
Proposition~27.1\,(d) with $\tilde S=\tilde X$ provides exactly one lift $(x,x')$ of $\mathrm{tq}$ contained in $\tilde X$; it is an element of $\mathrm{Pzt}_k$.
Let $u$ be a site of one of the paths $\boldsymbol{\Gamma}^{\pi}(c,c')$.
Each coordinate of $u$ lies between the corresponding coordinates of $c$ and $c'$ (Section~I.1.1), and $\mathrm{box}(\tilde X)$ is a product of three intervals that contains $c$ and $c'$ (0.4a).
So $u$ lies in $\mathrm{box}(\tilde X)$, and hence in $N(\tilde X)$.
The site $\mathrm{pr}_i(u)$ lies on the corresponding path of $\mathrm{tq}$, which lies in $X=\mathrm{pr}_i(\tilde X)$; therefore $\mathrm{pr}_i(u)=\mathrm{pr}_i(u')$ for some site $u'$ of $\tilde X$.
Proposition~27.1\,(b) gives $u=u'$.
Thus the six paths lie in $\tilde X$, and $(x,x')$ is of kind (I) for $\tilde X$.
(If $c=c'$, there is no path, and the condition reads $c\in\tilde X$.)

\emph{Pairs of kind (II).}
Let $\mathrm{tq}=(\mathrm{xq},\mathrm{xq}^{(3)})$ with $\pi(\mathrm{xq})\in X$ or $\pi(\mathrm{xq}^{(3)})\in X$.
If $\pi(\mathrm{xq})\in X$, let $x$ be the lift of $\mathrm{xq}$ over $\tilde X$, and put $x':=x^{(3)}$.
If $\pi(\mathrm{xq}^{(3)})\in X$, let $x'$ be the lift of $\mathrm{xq}^{(3)}$ over $\tilde X$, and put $x:=(x')^{(3)}$; then $x'=x^{(3)}$ by (I.s1).
In both cases $\mathrm{pr}_i(x)=\mathrm{xq}$ and $\mathrm{pr}_i(x^{(3)})=\mathrm{xq}^{(3)}$, since $\mathrm{pr}_i(x^{(3)})=(\mathrm{pr}_i x)^{(3)}$ (relation no.~2 of the list of relations in the proof of Proposition~27.1\,(f)).
Moreover $d_X(x,x^{(3)})=\eta_k\le M$.
So $(x,x^{(3)})$ is a lift of $\mathrm{tq}$ in $\mathrm{Pzt}_k$.
Since $c$ or $c_3$ lies in $\tilde X$, this lift is of kind (II) for $\tilde X$, and by (0.3c) both of its photon sites lie in $\tilde X^{\sim1}$.

\emph{Uniqueness; shortest lift.}
By (0.4c), and because $1\le6\cdot d_U$, the set $\tilde X^{\sim1}$ is contained in $N(\tilde X)$.
By Proposition~27.1\,(b), a pair therefore has at most one lift with both photon sites in $\tilde X^{\sim1}$.
Every lift in $\mathrm{Pzt}_k$ is the shortest lift (Section~I.1.7).
The two conditions do not exclude each other (Section~I.1.7).
For a pair that satisfies both of them for $X$, the two constructions above yield lifts with both photon sites in $\tilde X^{\sim1}$, hence one and the same lift, and this lift satisfies both conditions for $\tilde X$.

\emph{$F(X;\mathscr{A},\mathscr{C})$ lies in $\mathrm{Gr}_i(\tilde X)$.}
Write $F(X;\mathscr{A},\mathscr{C})$ with its antisymmetric kernels $F_{nm}$ (Proposition~15\,(a)).
The integrals $\int\mathrm{d}\xi$ and $\int\mathrm{d}\boldsymbol{\zeta}$ are the sums over the indices of the system (2.0), with the weights $\frac{\eta_k^3}{48}$ and $\left(\frac{\eta_k^3}{48}\right)^2$ (Section~I.3.6).
By (I.k2), $F_{nm}$ vanishes unless all $\xi_j$ lie over $X$ and all pairs $\boldsymbol{\zeta}_j$ are admissible for $X$.
Let $\xi_j=(x,\beta,\omega)$, with $c\in\mathrm{Rep}_i$, lie over $X$.
The lift of $\mathrm{pr}_i(x)$ over $\tilde X$ has its photon site in $N(\tilde X)$, which is a subset of $\mathrm{Rep}_i$, and it is congruent to $x$.
Hence this lift is $x$ itself, and $c\in\tilde X$.
Let $\boldsymbol{\zeta}_j=(x,x',\beta,\omega)$, with $(x,x')\in\mathrm{Pzt}_k$ and $c\in\mathrm{Rep}_i$, be admissible for $X$.
The lift constructed above has its first photon site in $\tilde X^{\sim1}$, a subset of $\mathrm{Rep}_i$, and differs from $(x,x')$ by a vector of $\ell_i\mathbb{Z}^3$.
Hence this lift is $(x,x')$ itself, and $c$, $c'$ lie in $\tilde X^{\sim1}$.
Consequently only generators with indices in $\mathrm{Ix}(\tilde X)$ occur.
\end{proof}

\begin{proof}[Proof of Proposition~27.2\,(c)]
\emph{Only if.}
Suppose that $F^{(1)}\sim F^{(2)}$, let $\tilde X$ be connected and admissible, and let $(\mathscr{A},\mathscr{C})\in\mathrm{Dm}(T_1)$.
By \cite[Lemma~17.7]{QEDI}, the map $\mathrm{pj}^{*}$ carries each of the sets $\tilde R_k$, $\tilde R_k^{-}$, $\mathfrak{X}_k$ of $T_1$ into the set of the same name of $T_2$; hence $\mathrm{pj}^{*}(\mathscr{A},\mathscr{C})\in\mathrm{Dm}(T_2)$.
By (0.5f), $(\mathscr{A},\mathscr{C})$ and $\mathrm{pj}^{*}(\mathscr{A},\mathscr{C})$ have equal lifts on all bonds of $\mathbb{L}_k$.
Therefore (D4) gives
\[
\mathrm{idn}\,F^{(1)}(\mathrm{pr}_1\tilde X;\mathscr{A},\mathscr{C})=F^{(2)}(\mathrm{pr}_2\tilde X;\mathrm{pj}^{*}\mathscr{A},\mathrm{pj}^{*}\mathscr{C}).
\]

\emph{If, and the last assertion of (c).}
Assume that the first identity displayed in (c), the one that contains $\mathrm{pj}^{*}\mathscr{A}$ and $\mathrm{pj}^{*}\mathscr{C}$, holds for every connected admissible $\tilde X$ and all $(\mathscr{A},\mathscr{C})\in\mathrm{Dm}(T_1)$.
Let $(\mathscr{A}_i,\mathscr{C}_i)\in\mathrm{Dm}(T_i)$ have equal lifts on the bonds that meet $\tilde X$.
Every bond of $T_2$ that meets $\mathrm{pr}_2(\tilde X)$ is of the form $\mathrm{pr}_2(b)$ with a bond $b$ that meets $\tilde X$ (0.5e), and the hypothesis together with (0.5f) gives
\[
\mathscr{A}_2(\mathrm{pr}_2 b)=\mathscr{A}_1(\mathrm{pr}_1 b)=\left(\mathrm{pj}^{*}\mathscr{A}_1\right)(\mathrm{pr}_2 b),
\]
and the same for $\mathscr{C}$.
So $(\mathscr{A}_2,\mathscr{C}_2)$ and $\mathrm{pj}^{*}(\mathscr{A}_1,\mathscr{C}_1)$, which both lie in $\mathrm{Dm}(T_2)$, coincide on the bonds that meet $\mathrm{pr}_2(\tilde X)$.
By (I.k2) for $F^{(2)}$, and by the assumed identity,
\begin{align*}
F^{(2)}(\mathrm{pr}_2\tilde X;\mathscr{A}_2,\mathscr{C}_2)
&=F^{(2)}(\mathrm{pr}_2\tilde X;\mathrm{pj}^{*}\mathscr{A}_1,\mathrm{pj}^{*}\mathscr{C}_1)\\
&=\mathrm{idn}\,F^{(1)}(\mathrm{pr}_1\tilde X;\mathscr{A}_1,\mathscr{C}_1).
\end{align*}
This is the last assertion of (c).
A bond that meets $\tilde X$ is contained in $\tilde X^{\sim1}$ (0.3c), hence in $N(\tilde X)$.
So connections with equal lifts on $N(\tilde X)$, that is, on the bonds contained in $N(\tilde X)$ (D1), have equal lifts on the bonds that meet $\tilde X$, and $F^{(1)}\sim F^{(2)}$ follows.
\end{proof}

\begin{proof}[Proof of Proposition~27.2\,(d)]
The assertions (t7) and (t10) of \cite[Theorem~13.2]{QEDI} (listed under the same marks in \cite[Proposition~5.58]{QEDI}) are made for sides in $\mathrm{Cl}$, and $\ell_1$, $\ell_2$ lie in $\mathrm{Cl}$ by (0.1a).
A union of $M$-cubes of $\mathbb{L}_k$ is identified with the union of the $M$-cubes of $\mathbb{R}^3$ with the same centres, and $d_M(\tilde X)$ is formed with nets in $\mathbb{R}^3$.

\emph{Admissibility.}
Let $\tilde X$ be non-empty and connected, with $d_M(\tilde X)\le L$.
Then $\tilde X$ is finite.
Indeed, the definition of \cite[Section~13.1.7]{QEDI}, formed with nets in $\mathbb{R}^3$, applies to an arbitrary union of $M$-cubes: $M\cdot d_M(\tilde X)$ is the infimum of the lengths of the nets in $\mathbb{R}^3$ that meet every closed $M$-cube of $\tilde X$; for finite unions that are connected through faces it gives the number $d_M$ of (t10).
A net, being a finite union of closed segments, is bounded and meets only finitely many closed $M$-cubes; thus $d_M(\tilde X)<\infty$ forces $\tilde X$ to be finite.
Let $\mathrm{lo}_\mu$, $\mathrm{hi}_\mu$, $\mathrm{sid}_\mu$ be the numbers of subsection~0.4 for $Z=\tilde X$, and choose centres $z^{-}$, $z^{+}$ of cubes of $\tilde X$ with $z^{-}_\mu=\mathrm{lo}_\mu$ and $z^{+}_\mu=\mathrm{hi}_\mu$.
The points $z^{-}-\frac{M}{2}\cdot e_\mu$ and $z^{+}+\frac{M}{2}\cdot e_\mu$ belong to closed cubes of $\tilde X$, and their coordinates of index $\mu$ differ by $\mathrm{sid}_\mu$.
By (t7) in $\mathbb{R}^3$,
\[
\mathrm{sid}_\mu\le M\cdot\left(d_M(\tilde X)+2\right)\le(L+2)\cdot M.
\]
By (0.1b),
\[
\ell_1-14\cdot d_U\cdot M>(31L+16)\cdot M-(28L+14)\cdot M=(3L+2)\cdot M\ge(L+2)\cdot M.
\]
Hence $\mathrm{diam}(\tilde X)\le\ell_1-14\cdot d_U\cdot M$, that is, $\tilde X$ is admissible.

\emph{Small polymers are flat; the bijection.}
Let $X$ be a small polymer of $T_i$.
By (0.1b),
\[
2M\cdot\left(d_M(X)+2\right)\le(2L+4)\cdot M<\ell_i.
\]
So (t10) applies: $X=\mathrm{pr}_i(\tilde X)$ for a finite union $\tilde X$ of $M$-cubes that is connected through faces, and $d_M(\tilde X)=d_M(X)\le L$.
The union $\tilde X$ is non-empty, and it is admissible by the first part.
Therefore $X$ is flat.
The bijection of (d) is the last assertion of (t10).

\emph{The equivalences.}
Let $\tilde X$ be connected and admissible.
Then $\mathrm{pr}_i(\tilde X)$ is a polymer: it is non-empty, and cubes with a common face are mapped to cubes with a common face (relation no.~9 of the list of relations in the proof of Proposition~27.1\,(f)).
Moreover $\mathrm{pr}_i$ is injective on $\tilde X$ (Proposition~27.1\,(b)).
The same holds in $\mathbb{R}^3$.
Indeed, for $z\in\mathrm{Ctr}(\tilde X)$ the coordinate of index $\mu$ of a point of the closed cube $\left\{y\in\mathbb{R}^3:|y-z|\le\frac{M}{2}\right\}$ lies in the interval $\left[\mathrm{lo}_\mu-\frac{M}{2},\mathrm{hi}_\mu+\frac{M}{2}\right]$, formed with the numbers of subsection~0.4 for this $\tilde X$; its length is
\[
\mathrm{sid}_\mu\le\mathrm{diam}(\tilde X)\le\ell_1-14\cdot d_U\cdot M<\ell_i
\]
by (0.4a).
So two distinct points of the closed cubes of $\tilde X$ never differ by a vector of $\ell_i\mathbb{Z}^3$, and the projection of $\mathbb{R}^3$ onto $\mathbb{R}^3/\ell_i\mathbb{Z}^3$ is injective on the union of these closed cubes.
Suppose that $\mathrm{pr}_i(\tilde X)$ is small.
By (t10), applicable as in the preceding paragraph, $\mathrm{pr}_i(\tilde X)=\mathrm{pr}_i(\tilde X')$ for a union $\tilde X'$ that is connected through faces, with $\mathrm{pr}_i$ injective on $\tilde X'$ and $d_M(\tilde X')=d_M(\mathrm{pr}_i\tilde X)\le L$.
By the uniqueness clause of (t10), according to which a connected union of $M$-cubes on which $\mathrm{pr}_i$ is injective and which has a given image $X$ is unique up to translations by $\ell_i\mathbb{Z}^3$, the union $\tilde X$ is a translate of $\tilde X'$.
Since $d_M$ is invariant under translations, $d_M(\tilde X)\le L$.
Conversely, let $d_M(\tilde X)\le L$.
Then $2M\cdot\left(d_M(\tilde X)+2\right)<\ell_i$, the union $\tilde X$ is non-empty, and the converse part of (t10) gives $d_M(\mathrm{pr}_i\tilde X)=d_M(\tilde X)\le L$.
Consequently, for $i=1$ as well as for $i=2$, the polymer $\mathrm{pr}_i(\tilde X)$ is small if and only if $d_M(\tilde X)\le L$, and in this case $d_M(\mathrm{pr}_i\tilde X)=d_M(\tilde X)$.
Finally, $\mathrm{pr}_i$ is injective on the $M$-cubes of $\tilde X$ (Proposition~27.1\,(b)); hence $|\mathrm{pr}_i\tilde X|_M=|\tilde X|_M$ and
\[
\operatorname{Vol}(\mathrm{pr}_i\tilde X)=M^3\cdot|\tilde X|_M=\operatorname{Vol}(\tilde X).
\]
\end{proof}

\begin{proof}[Proof of Proposition~27.2\,(e)]
\begin{itemize}
\item[(e1)] \emph{Indices.}
The image of an index contained in $\tilde Y$ is contained in $\mathrm{pr}_i(\tilde Y)$.
By Proposition~27.1\,(d), applied with the admissible set $\tilde Y$ in the place of $\tilde X$ and with $\tilde S:=\tilde Y$, every index of $T_i$ contained in $\mathrm{pr}_i(\tilde Y)$ has exactly one lift contained in $\tilde Y$.
Hence, for each sort of indices, $\mathrm{pr}_i$ is a bijection of the indices contained in $\tilde Y$ onto the indices of $T_i$ contained in $\mathrm{pr}_i(\tilde Y)$.
This holds for the $M$-cubes, for the parameter cubes of each of the three kinds (a parameter cube is contained in a set if all its lattice sites lie in it: \cite[Section~5.2.3, Definition~5.8]{QEDI}), for the unit sites, for the fermion sites (since $\operatorname{supp}x=\{c\}$, ``contained in'' means ``over'') and for the basis elements.

\item[(e2)] \emph{Common faces.}
Two $M$-cubes contained in $\tilde Y$ have a common face if and only if their images have one (Proposition~27.1\,(e), and relation no.~9 of the list of relations in the proof of Proposition~27.1\,(f)).
Together with (e1) this shows that, for every union $\tilde S'$ of $M$-cubes contained in $\tilde Y$, the map $\mathrm{pr}_i$ is an isomorphism of the graph whose vertices are the $M$-cubes of $\tilde S'$, two of them being joined when they have a common face, onto the graph formed in the same way with the $M$-cubes of $\mathrm{pr}_i(\tilde S')$.

\item[(e3)] \emph{Polymers.}
By (e2), the image of a non-empty connected union of $M$-cubes contained in $\tilde Y$ is a polymer contained in $\mathrm{pr}_i(\tilde Y)$.
Let $X$ be a polymer of $T_i$ contained in $\mathrm{pr}_i(\tilde Y)$, and let $\tilde X$ be the union of the lifts in $\tilde Y$ of its $M$-cubes (e1).
Then $\mathrm{pr}_i(\tilde X)=X$, the union $\tilde X$ is non-empty, and it is connected by (e2).
If $\mathrm{pr}_i(\tilde X)=\mathrm{pr}_i(\tilde X')$ for unions $\tilde X$, $\tilde X'$ contained in $\tilde Y$, then $\tilde X=\tilde X'$, because $\mathrm{pr}_i$ is injective on the $M$-cubes of $\tilde Y$.

\item[(e4)] The set $\tilde Y^{\sim(L+1)}$ is contained in $\mathrm{box}(\tilde Y)^{\sim(L+1)}$ (0.4c).
This set is contained in $\mathrm{box}(\tilde Y)^{\sim6d_U}=N(\tilde Y)$, because $L+1\le6\cdot(2L+1)$; recall that, by the definition in subsection~0.3, $Z^{\sim a}\subset Z^{\sim a'}$ for $a\le a'$.
The map $\mathrm{pr}_i$ is injective on $N(\tilde Y)$ (Proposition~27.1\,(b)).

\item[(e5)] \emph{$P_0$; the assertions of (e) hold for every choice of a set of kinds.}
The objects $P_0$, $\tilde P$, $\mathrm{Yp}$, $Y'$ are formed with the parameter cubes of a set of kinds that is chosen among $\pi_k$, $\pi'_k$, $\pi''_k$ and then kept fixed (``the kinds that occur'' of Definition~24.D\,(D0); compare \cite[Section~17.4.3]{QEDI}).
The steps (e5), (e6), (e9) hold for every choice of this set, since they depend on the kinds only through (e1), which holds for each kind separately.
By Proposition~27.1\,(a), $\tilde X$ is admissible.
Applied with $\tilde X$ in the place of $\tilde Y$, (e1) shows that the parameter cubes of each kind contained in $\mathrm{pr}_i(\tilde X)$ are the images of those contained in $\tilde X$:
\[
P_0(\mathrm{pr}_i\tilde X)=\mathrm{pr}_iP_0(\tilde X).
\]

\item[(e6)] \emph{$Y'$.}
All sites of $\tilde X$ and of the cubes of $P_0(\tilde X)\cup\tilde P$ lie in $\tilde Y$; so $\tilde X\cup\mathrm{Yp}(\tilde X,\tilde P)$ is contained in $\tilde Y$.
An $M$-cube that meets this set has a site in $\tilde Y$, and by (0.3a) it is a cube of $\tilde Y$.
Hence $Y'(\tilde X,\tilde P)$ is contained in $\tilde Y$.
The cubes of $\mathrm{pr}_i\tilde P$ do not belong to $P_0(\mathrm{pr}_i\tilde X)=\mathrm{pr}_iP_0(\tilde X)$, because $\mathrm{pr}_i$ is injective on the parameter cubes contained in $\tilde Y$ (e1); so $Y'(\mathrm{pr}_i\tilde X,\mathrm{pr}_i\tilde P)$ is defined.
Since the image of a union is the union of the images,
\[
\mathrm{pr}_i\left(\tilde X\cup\mathrm{Yp}(\tilde X,\tilde P)\right)=\mathrm{pr}_i\tilde X\cup\mathrm{Yp}(\mathrm{pr}_i\tilde X,\mathrm{pr}_i\tilde P).
\]
If an $M$-cube $\tilde\square$ meets $\tilde X\cup\mathrm{Yp}(\tilde X,\tilde P)$, then $\mathrm{pr}_i(\tilde\square)$ meets the image of this set.
Conversely, let an $M$-cube $\square$ of $T_i$ contain a site $\mathrm{pr}_i(u)$ with $u\in\tilde X\cup\mathrm{Yp}(\tilde X,\tilde P)$.
The $M$-cube $\tilde\square$ of $\mathbb{L}_k$ that contains $u$ meets $\tilde X\cup\mathrm{Yp}(\tilde X,\tilde P)$, and $\mathrm{pr}_i(\tilde\square)$ is the $M$-cube of $T_i$ that contains $\mathrm{pr}_i(u)$, that is, $\square$ (0.5d).
Therefore
\[
\mathrm{pr}_iY'(\tilde X,\tilde P)=Y'(\mathrm{pr}_i\tilde X,\mathrm{pr}_i\tilde P).
\]

\item[(e7)] \emph{Closure.}
The $LM$-closure of $\tilde X$ is the union of the $LM$-cubes that contain an $M$-cube of $\tilde X$ (\cite[Section~13.1.7]{QEDI}).
Let $\mathrm{Kub}_{LM}(z')$ contain the cube $\mathrm{Kub}_M(z)$ of $\tilde X$.
By (0.2b), the $M$-cubes of $\mathrm{Kub}_{LM}(z')$ have the centres $z'+M\cdot w$ with $|w|\le\frac{L-1}{2}$, and $z$ is one of these centres.
So all of them have sup-distance at most $(L-1)\cdot M$ from $z$, and $\mathrm{Kub}_{LM}(z')$ lies in $\tilde X^{\sim(L-1)}$, hence in $\tilde Y^{\sim(L-1)}$ (0.3b).
The map $\mathrm{pr}_i$ sends $LM$-cubes onto $LM$-cubes of $T_i$, and the $LM$-cubes of $T_i$ partition the sites of $T_i$ (0.5d).
If $\mathrm{Kub}_{LM}(z')$ contains the cube $\tilde\square$ of $\tilde X$, its image is the $LM$-cube of $T_i$ that contains $\mathrm{pr}_i(\tilde\square)$.
Every $M$-cube of $\mathrm{pr}_i(\tilde X)$ is of this form $\mathrm{pr}_i(\tilde\square)$.
Hence the image of the $LM$-closure of $\tilde X$ is the union of the $LM$-cubes of $T_i$ that contain an $M$-cube of $\mathrm{pr}_i(\tilde X)$, that is, $\overline{\mathrm{pr}_i\tilde X}^{L}$.

\item[(e8)] \emph{Connectedness and the other relations.}
By (e2) with $\tilde S'=\tilde X_1\cup\tilde X_2$, the union $\tilde X_1\cup\tilde X_2$ is connected if and only if $\mathrm{pr}_i\tilde X_1\cup\mathrm{pr}_i\tilde X_2$ is connected.
Two polymers are incompatible if their union is connected (\cite[Section~13.3]{QEDI}); so incompatibility is preserved in both directions.
Let $\tilde S_1$, $\tilde S_2$ be sets of sites contained in $\tilde Y$ (unions of $M$-cubes, parameter cubes, supports of indices).
Since $\mathrm{pr}_i$ is injective on $\tilde Y$, the set $\tilde S_1$ is contained in $\tilde S_2$ if and only if $\mathrm{pr}_i\tilde S_1$ is contained in $\mathrm{pr}_i\tilde S_2$; moreover $\mathrm{pr}_i(\tilde S_1\cup\tilde S_2)=\mathrm{pr}_i\tilde S_1\cup\mathrm{pr}_i\tilde S_2$, and $\tilde S_1$, $\tilde S_2$ have a common site (they ``meet'') if and only if their images have one.
Common faces are dealt with in (e2).
Links between two parameter cubes of one kind contained in $\tilde Y$ are preserved in both directions by Proposition~27.1\,(e) and relation no.~10 of the list of relations in the proof of Proposition~27.1\,(f).

\item[(e9)] \emph{The last assertion of (e).}
Let $Y'(X,P)=\mathbf{Y}$, on $T_i$ or on $\mathbb{L}_k$.
Every site of $X\cup\mathrm{Yp}(X,P)$ lies in an $M$-cube ((0.2a), (0.5d)), and this cube meets $X\cup\mathrm{Yp}(X,P)$; hence it is one of the cubes of $Y'(X,P)$.
Consequently $X\cup\mathrm{Yp}(X,P)$ is contained in $\mathbf{Y}$: the set $X$ and all sites of all cubes of $P_0(X)\cup P$ lie in $\mathbf{Y}$.
\end{itemize}
\end{proof}

\begin{proof}[Proof of Proposition~27.2\,(f)]
The family $\mathscr{E}^{(i)}|_0=\mathscr{E}^{(i)}(\,\cdot\,;\,\cdot\,,0)$ is a family of even elements of $\mathrm{Gr}[\psi,\chi]$ on $\tilde R_k(T_i)$ (resp.\ on $\tilde R_k^{-}(T_i)$) with property (k2)$^{\mathrm{o}}$ of the Definition in Proposition~48.1, which is (I.k2) at $\mathscr{C}=0$.
So it is a local family without slot in the sense of (D3).
Let $\tilde X$ be connected and admissible, and let $\mathscr{A}_1$, $\mathscr{A}_2$, with $\mathscr{A}_i\in\mathrm{Dm}^{\mathrm{o}}(T_i)$, have equal lifts on $N(\tilde X)$.
Since $0\in\mathfrak{X}_k(T_i)$ (Section~I.5.3), the pair $(\mathscr{A}_i,0)$ lies in $\mathrm{Dm}(T_i)$, and the pairs $(\mathscr{A}_1,0)$, $(\mathscr{A}_2,0)$ have equal lifts on $N(\tilde X)$.
By (D4) for $\mathscr{E}^{(1)}\sim\mathscr{E}^{(2)}$,
\[
\mathrm{idn}\,\mathscr{E}^{(1)}(\mathrm{pr}_1\tilde X;\mathscr{A}_1,0)=\mathscr{E}^{(2)}(\mathrm{pr}_2\tilde X;\mathscr{A}_2,0).
\]
This is the display of (D4) for the families without slot.
\end{proof}

\begin{proof}[Proof of Proposition~27.2\,(g)]
\emph{$F^{\mathrm{c}}$ is a local family.}
By Proposition~26.6\,(d), the quadruples $(k,\ell_i,\mu_k,e_k)$ are (SF); so the data (S-1) of Proposition~48.1 are given on $T_i$.
The element $F^{\mathrm{c}}(X;\mathscr{A},\mathscr{C})=F(X;\mathscr{A})$ is an even element of $\mathrm{Gr}[\psi,\chi]$, and it does not depend on $\mathscr{C}$.
Its antisymmetric kernels are those of $F(X;\mathscr{A})$ (Proposition~15\,(a)), and the points $(\mathscr{A},0)$ with $\mathscr{A}\in\mathrm{Dm}^{\mathrm{o}}(T_i)$ lie in $\mathrm{Dm}^{\mathrm{o}}(T_i)\times\mathfrak{X}_k(T_i)$, because $0\in\mathfrak{X}_k$ (Section~I.5.3).
Hence the first clause of (I.k2) for $F^{\mathrm{c}}$, read on $\mathrm{Dm}^{\mathrm{o}}\times\mathfrak{X}_k$, and the first clause of (k2)$^{\mathrm{o}}$ for $F$, read on $\mathrm{Dm}^{\mathrm{o}}$, are one and the same condition: the kernels vanish unless all $\xi_j$ lie over $X$ and all pairs $\boldsymbol{\zeta}_j$ are admissible for $X$.
Consider now the second clauses.
If $(\mathscr{A},\mathscr{C})$ and $(\tilde{\mathscr{A}},\tilde{\mathscr{C}})$ in $\mathrm{Dm}^{\mathrm{o}}(T_i)\times\mathfrak{X}_k(T_i)$ coincide, in both components, on the bonds that meet $X$, then $\mathscr{A}$ and $\tilde{\mathscr{A}}$ coincide there, and (k2)$^{\mathrm{o}}$ gives
\[
F^{\mathrm{c}}(X;\mathscr{A},\mathscr{C})=F(X;\mathscr{A})=F(X;\tilde{\mathscr{A}})=F^{\mathrm{c}}(X;\tilde{\mathscr{A}},\tilde{\mathscr{C}}).
\]
Conversely, let $F^{\mathrm{c}}$ have (I.k2), and let $\mathscr{A},\tilde{\mathscr{A}}\in\mathrm{Dm}^{\mathrm{o}}(T_i)$ coincide on the bonds that meet $X$.
Then the pairs $(\mathscr{A},0)$ and $(\tilde{\mathscr{A}},0)$ coincide there in both components, and therefore
\[
F(X;\mathscr{A})=F^{\mathrm{c}}(X;\mathscr{A},0)=F^{\mathrm{c}}(X;\tilde{\mathscr{A}},0)=F(X;\tilde{\mathscr{A}}).
\]
Thus (I.k2) for $F^{\mathrm{c}}$ is (k2)$^{\mathrm{o}}$ for $F$, read on $\mathrm{Dm}^{\mathrm{o}}$; for $\mathrm{Dm}^{\mathrm{o}}=\tilde R_k$ this is Proposition~48.1\,(b), which is stated for families on $\tilde R_k$ only.
So $F^{\mathrm{c}}$ is a local family in the sense of (D3), with $\mathrm{Dm}=\mathrm{Dm}^{\mathrm{o}}\times\mathfrak{X}_k$.
As $0\in\mathfrak{X}_k(T_i)$ (Section~I.5.3), one has $F(X;\mathscr{A})=F^{\mathrm{c}}(X;\mathscr{A},0)$, and by (b) for $F^{\mathrm{c}}$ this element lies in $\mathrm{Gr}_i(\tilde X)$ whenever $X=\mathrm{pr}_i(\tilde X)$ is flat.
Hence the left side of the last display of (D4), the one for local families without slot, is defined.

\emph{The two relations.}
Let $\tilde X$ be connected and admissible.
Suppose that $F^{(1)}\sim F^{(2)}$, and let $(\mathscr{A}_i,\mathscr{C}_i)\in\mathrm{Dm}(T_i)$ have equal lifts on $N(\tilde X)$.
Then $\mathscr{A}_1$ and $\mathscr{A}_2$ have equal lifts on $N(\tilde X)$, and
\begin{align*}
\mathrm{idn}\,F^{(1)\mathrm{c}}(\mathrm{pr}_1\tilde X;\mathscr{A}_1,\mathscr{C}_1)
&=\mathrm{idn}\,F^{(1)}(\mathrm{pr}_1\tilde X;\mathscr{A}_1)\\
&=F^{(2)}(\mathrm{pr}_2\tilde X;\mathscr{A}_2)
=F^{(2)\mathrm{c}}(\mathrm{pr}_2\tilde X;\mathscr{A}_2,\mathscr{C}_2);
\end{align*}
that is, $F^{(1)\mathrm{c}}\sim F^{(2)\mathrm{c}}$.
Conversely, suppose that $F^{(1)\mathrm{c}}\sim F^{(2)\mathrm{c}}$, and let $\mathscr{A}_i\in\mathrm{Dm}^{\mathrm{o}}(T_i)$ have equal lifts on $N(\tilde X)$.
The pairs $(\mathscr{A}_i,0)$ lie in $\mathrm{Dm}(T_i)$ and have equal lifts on $N(\tilde X)$, and (D4) for $F^{(1)\mathrm{c}}$, $F^{(2)\mathrm{c}}$ at these pairs reads
\[
\mathrm{idn}\,F^{(1)}(\mathrm{pr}_1\tilde X;\mathscr{A}_1)=F^{(2)}(\mathrm{pr}_2\tilde X;\mathscr{A}_2).
\]

\emph{Clause (c) without slot.}
By (c), $F^{(1)\mathrm{c}}\sim F^{(2)\mathrm{c}}$ holds if and only if
\[
\mathrm{idn}\,F^{(1)\mathrm{c}}(\mathrm{pr}_1\tilde X;\mathscr{A},\mathscr{C})=F^{(2)\mathrm{c}}(\mathrm{pr}_2\tilde X;\mathrm{pj}^{*}\mathscr{A},\mathrm{pj}^{*}\mathscr{C})
\]
for every connected admissible $\tilde X$ and all $(\mathscr{A},\mathscr{C})\in\mathrm{Dm}(T_1)$.
Neither side contains $\mathscr{C}$, and $\mathfrak{X}_k(T_1)$ is not empty.
So this is the identity
\[
\mathrm{idn}\,F^{(1)}(\mathrm{pr}_1\tilde X;\mathscr{A})=F^{(2)}(\mathrm{pr}_2\tilde X;\mathrm{pj}^{*}\mathscr{A})\qquad\text{for all }\mathscr{A}\in\mathrm{Dm}^{\mathrm{o}}(T_1).
\]
If $\mathscr{A}_i\in\mathrm{Dm}^{\mathrm{o}}(T_i)$ have equal lifts on the bonds that meet $\tilde X$, then so have the pairs $(\mathscr{A}_i,0)$, and the last assertion of (c) for $F^{(1)\mathrm{c}}$, $F^{(2)\mathrm{c}}$ gives
\[
\mathrm{idn}\,F^{(1)}(\mathrm{pr}_1\tilde X;\mathscr{A}_1)=F^{(2)}(\mathrm{pr}_2\tilde X;\mathscr{A}_2).
\]
\end{proof}

\paragraph{Level 0.}
In (a) the generators $W$ carry the label: they are the $W(x,\beta,\omega)$ with $x=(c,g)$ and $c\in N(\tilde X)$, that is, $2\cdot4\cdot48$ generators over each photon site, and there is no family $\Psi$.
The weight $\nu=\frac{1}{48}$ of these indices does not enter (I.5.1).
In (b) every pair has $d_X\ge1$, and the pairs inside a cluster over $\tilde X$ are of kind (I).
All other arguments apply without change.

\paragraph{Uniformity.}
Proposition~27.2 introduces no constant and proves no bound: its only standing hypothesis is (TT), and its clauses hold for all data (TT) and at every level $k\le k_m$.
The constants $L$, $M$ and $d_U$ that occur in its clauses and in the proof above are those of the lines (C3), (C4) and (C5) of Section~I.8.2, where they are chosen before $\bar{m}$, $e$, the torus and $N$;
they do not depend on the data (TT), and the conditions on them on which the proof rests are those listed at the end of Chapter~27a.

\subsection*{3. Proposition 27.3}

\subsubsection*{Proof of Proposition 27.3}

\begin{proof}
\emph{The two systems.}
By (TT) one has $k<k_m\le k_{\mathrm{top}}$.
The system $\mathrm{Op}^{(2)}$ is the system of Proposition~26.5 on the torus $T_2=T_k(R_{\mathrm{top}})$.
According to Proposition~26.5 it has the properties W1--W17$^{\mathrm{ch}}$ of Section~G4.0.4 for $\mathrm{par}^{0}$; among them are W8$^{\mathrm{ch}}$, W12$^{\mathrm{ch}}$ and W14$^{\mathrm{ch}}$.
The system $\mathrm{Op}^{(1)}$ is the projection of $\mathrm{Op}^{(2)}$ to $T_1$ in the sense of Proposition~26.1.
Proposition~26.3\,(a) expresses its kernels by the formula (8.0), and Proposition~26.3\,(b) shows that it is a system of primitives of level $k$ on $T_1$ with W1--W17$^{\mathrm{ch}}$ for $\mathrm{par}$, in particular with W8$^{\mathrm{ch}}$ and W12$^{\mathrm{ch}}$.
By Proposition~26.6\,(d), the data $(k,\ell_i,\mu_k,e_k)$ are (SF) in the sense of Section~I.1.8, for $i=1$ and for $i=2$.
Proposition~26.4 therefore holds on $T_1$ for $\mathrm{Op}^{(1)}$ and on $T_2$ for $\mathrm{Op}^{(2)}$.
Kernels are taken at indices of the covering lattice (Section~G4.0.4): for an operator $O$ on $T_i$,
\[
O\left(t+\ell_i\cdot e_{\mu},t'\right)=O\left(t,t'+\ell_i\cdot e_{\mu}\right)=\boldsymbol{\chi}(e_{\mu})\cdot O(t,t'),
\]
so that whether a kernel vanishes at $(t,t')$ depends on $\mathrm{pr}_i(t)$ and $\mathrm{pr}_i(t')$ alone.

\paragraph{Step 1 (geometry).}
Put
\[
\widetilde{\mathrm{Nb}}:=\tilde{Y}^{\sim(j+c_{l})},\qquad\mathrm{Nb}_i:=\mathrm{pr}_i\left(\widetilde{\mathrm{Nb}}\right).
\]
As $c_{l}+1=d_U$ (recall that $c_{l}=2L$ and $d_U=2L+1$), (0.3b) gives $\widetilde{\mathrm{Nb}}^{\sim1}=\tilde{Y}^{\sim(j+d_U)}$.
The last sentence of Proposition~27.1\,(d) yields
\[
\mathrm{pr}_i\left(\tilde{Y}^{\sim j}\right)=Y_i^{\sim j},\qquad\mathrm{Nb}_i=Y_i^{\sim(j+c_{l})},\qquad\mathrm{pr}_i\left(\widetilde{\mathrm{Nb}}^{\sim1}\right)=Y_i^{\sim(j+d_U)}.
\]
Since $j\le5\cdot d_U$, one has $j+d_U\le6\cdot d_U$, and (0.4c) shows that
\begin{equation}\tag{3.1}
\begin{gathered}
\tilde{Y}^{\sim j},\quad\widetilde{\mathrm{Nb}}\quad\text{and}\quad\widetilde{\mathrm{Nb}}^{\sim1}=\tilde{Y}^{\sim(j+d_U)}\\
\text{are contained in}\quad\mathrm{box}(\tilde{Y})^{\sim6d_U}=N(\tilde{Y}).
\end{gathered}
\end{equation}

\emph{The lift $\tilde{t}$.}
The row index $t$ is contained in $Y_1^{\sim j}=\mathrm{pr}_1\left(\tilde{Y}^{\sim j}\right)$.
In view of (3.1), Proposition~27.1\,(d) applies with $\tilde{S}=\tilde{Y}^{\sim j}$; consequently $t$ has exactly one lift $\tilde{t}$ that is contained in $\tilde{Y}^{\sim j}$.

\emph{The set $\square(\tilde{t})$.}
Recall that $\square(\tilde{t})$ is the union of the $M$-cubes that $\tilde{t}$ meets (Section~G4.0.4).
It is a box consisting either of one $M$-cube or of two $M$-cubes with a common face.
Indeed, a unit site and a fermion site have one support point each.
Let $\tilde{t}$ be a bond with the ends $c$ and $c'=c+a\cdot e_{\mu}$, where $a\in\{\pm\eta_k,\pm1\}$, and let $z$, $z'$ be the centres of the $M$-cubes of $c$ and $c'$.
For $\nu\ne\mu$ the coordinates of index $\nu$ of $c$ and $c'$ are equal, and therefore $z'_{\nu}=z_{\nu}$.
Moreover $|z'_{\mu}-z_{\mu}|<\frac{M}{2}+1+\frac{M}{2}$; as $z'_{\mu}-z_{\mu}$ is a multiple of $M$, it follows that $|z'_{\mu}-z_{\mu}|\le M$.
The support of a basis element either lies in one $L$-block, which lies in one $M$-cube by (0.2b), or consists of the two ends of one bond, which is the case just treated.
Each cube of $\square(\tilde{t})$ contains a site of $\tilde{Y}^{\sim j}$; by (0.3a), $\square(\tilde{t})\subset\tilde{Y}^{\sim j}$.
It follows that $\square(\tilde{t})^{\sim c_{l}}$ is a box (0.4b), that it is connected through faces (0.4d), and that it is contained in $\left(\tilde{Y}^{\sim j}\right)^{\sim c_{l}}=\widetilde{\mathrm{Nb}}$ (0.3b).
Define
\[
\widetilde{\mathrm{Cp}}:=\text{the component of }\widetilde{\mathrm{Nb}}\text{ that contains }\square(\tilde{t})^{\sim c_{l}}.
\]
By (0.5d) the $M$-cubes met by $\mathrm{pr}_i(\tilde{t})$ are the images of the $M$-cubes met by $\tilde{t}$; that is, $\square\left(\mathrm{pr}_i\tilde{t}\right)=\mathrm{pr}_i\left(\square(\tilde{t})\right)$.
The set $\square(\tilde{t})$ is admissible:
\[
\mathrm{diam}\left(\square(\tilde{t})\right)\le2M\le\ell_1-14\cdot d_U\cdot M,
\]
because $14\cdot d_U\cdot M=(28L+14)\cdot M$, while $\ell_1>(31L+16)\cdot M\ge(28L+16)\cdot M$ by (0.1b).
The last sentence of Proposition~27.1\,(d), with $\square(\tilde{t})$ in the place of $\tilde{X}$, therefore gives
\begin{equation}\tag{3.2}
\square\left(\mathrm{pr}_i\tilde{t}\right)^{\sim c_{l}}=\mathrm{pr}_i\left(\square(\tilde{t})^{\sim c_{l}}\right)\subset\mathrm{pr}_i\left(\widetilde{\mathrm{Cp}}\right).
\end{equation}

\emph{Components.}
The following assertion holds.
\begin{itemize}
\item[(3.3)] Let $\tilde{\square}$ be an $M$-cube of $\widetilde{\mathrm{Cp}}$, and let $\square'$ be an $M$-cube of $T_i$ that has a common face with $\mathrm{pr}_i(\tilde{\square})$ and all of whose sites are of the form $\mathrm{pr}_i(c)$ with $c\in\widetilde{\mathrm{Nb}}+v$, $v\in\ell_1\mathbb{Z}^3$.
Then $\square'=\mathrm{pr}_i(\tilde{\square}')$ for an $M$-cube $\tilde{\square}'$ of $\widetilde{\mathrm{Cp}}$.
\end{itemize}
For the proof, recall relation no.~9 of the list in the proof of Proposition~27.1\,(f), the relation ``two $M$-cubes have a common face'': the two cubes $\mathrm{pr}_i(\tilde{\square})$ and $\square'$ have lifts with a common face.
After a translation by a vector of $\ell_i\mathbb{Z}^3$, the lift of $\mathrm{pr}_i(\tilde{\square})$ is $\tilde{\square}$ itself.
Thus $\square'=\mathrm{pr}_i(\tilde{\square}')$ for an $M$-cube $\tilde{\square}'$ that has a common face with $\tilde{\square}$.
The centres of $\tilde{\square}$ and $\tilde{\square}'$ differ by $M$ in one coordinate; so $\tilde{\square}'$ is a cube of $\widetilde{\mathrm{Nb}}^{\sim1}$, and it is contained in $N(\tilde{Y})$ by (3.1).
Each site of $\tilde{\square}'$ is congruent modulo $\ell_i\mathbb{Z}^3$ to a site of some set $\widetilde{\mathrm{Nb}}+v$ with $v\in\ell_1\mathbb{Z}^3$.
As $\ell_i\mathbb{Z}^3\subset\ell_1\mathbb{Z}^3$, the cube $\tilde{\square}'$ has a site in $\widetilde{\mathrm{Nb}}+v'$ for some $v'\in\ell_1\mathbb{Z}^3$, and $\widetilde{\mathrm{Nb}}+v'$ lies in $N(\tilde{Y})+v'$.
Proposition~27.1\,(b) gives $v'=0$.
By (0.3a), $\tilde{\square}'$ is then a cube of $\widetilde{\mathrm{Nb}}$; since it has a common face with the cube $\tilde{\square}$ of the component $\widetilde{\mathrm{Cp}}$, it is a cube of $\widetilde{\mathrm{Cp}}$.
This proves (3.3).

The map $\mathrm{pr}_i$ takes cubes with a common face to cubes with a common face; hence $\mathrm{pr}_i\left(\widetilde{\mathrm{Cp}}\right)$ is connected through faces.
The set $\left\{\mathrm{pr}_1(c):\ c\in\widetilde{\mathrm{Nb}}+\ell_1\mathbb{Z}^3\right\}$ is $\mathrm{Nb}_1$.
The set $\left\{\mathrm{pr}_2(c):\ c\in\widetilde{\mathrm{Nb}}+\ell_1\mathbb{Z}^3\right\}$ is $\mathrm{pj}^{-1}(\mathrm{Nb}_1)$, because $\mathrm{pr}_2(c)$ lies in $\mathrm{pj}^{-1}(\mathrm{Nb}_1)$ if and only if $\mathrm{pr}_1(c)=\mathrm{pj}\left(\mathrm{pr}_2(c)\right)$ lies in $\mathrm{pr}_1\left(\widetilde{\mathrm{Nb}}\right)$; this set contains $\mathrm{Nb}_2$.
In these terms (3.3) states:
\begin{equation}\tag{3.4}
\begin{aligned}
&\mathrm{pr}_1\left(\widetilde{\mathrm{Cp}}\right)\ \text{is a component of}\ \mathrm{Nb}_1;\\
&\mathrm{pr}_2\left(\widetilde{\mathrm{Cp}}\right)\ \text{is a component of}\ \mathrm{pj}^{-1}(\mathrm{Nb}_1),\ \text{and, being contained in}\ \mathrm{Nb}_2,\\
&\qquad\text{it is also a component of}\ \mathrm{Nb}_2.
\end{aligned}
\end{equation}

\paragraph{Step 2 (clause (a)).}
Since $\tilde{Y}\subset\widetilde{\mathrm{Nb}}$, the set $Y_1$ is contained in $\mathrm{Nb}_1$.
A parameter cube that is not contained in $\mathrm{Nb}_1$ is therefore not contained in $Y_1$, and so $s^{(1)}$ vanishes off $\mathrm{Nb}_1$.
Proposition~26.4 on $T_1$, applied to $\mathrm{Op}^{(1)}$ with the union $\mathrm{Nb}_1$ of $M$-cubes, the family $s^{(1)}$, the component $\mathrm{pr}_1\left(\widetilde{\mathrm{Cp}}\right)$ of (3.4) and the row index $t=\mathrm{pr}_1(\tilde{t})$, for which $\square(t)^{\sim c_{l}}\subset\mathrm{pr}_1\left(\widetilde{\mathrm{Cp}}\right)$ by (3.2), gives
\begin{gather*}
\mathrm{Op}^{(1)}\left(s^{(1)},A^{(1)}\right)(t,t')=0\quad\text{unless $t'$ is contained in}\ \mathrm{pr}_1\left(\widetilde{\mathrm{Cp}}\right),\\
\mathrm{Op}^{(1)}\left(s^{(1)},A^{(1)}\right)(t,\,\cdot\,)=\mathrm{Op}^{(1)}\left(s^{(1)}\cdot\mathbf{1}_{\mathrm{pr}_1(\widetilde{\mathrm{Cp}})},A_1\right)(t,\,\cdot\,)
\end{gather*}
for every $A_1\in\tilde{G}_k(T_1)$ that coincides with $A^{(1)}$ on the bonds meeting $\mathrm{pr}_1\left(\widetilde{\mathrm{Cp}}\right)$.
As $\mathrm{pr}_1\left(\widetilde{\mathrm{Cp}}\right)$ is contained in $\mathrm{Nb}_1=Y_1^{\sim(j+c_{l})}$, the first assertion of (a) follows.
Every bond of $T_1$ that meets $\mathrm{pr}_1\left(\widetilde{\mathrm{Cp}}\right)$ is of the form $\mathrm{pr}_1(b)$ with a bond $b$ of $\mathbb{L}_k$ that meets $\widetilde{\mathrm{Cp}}$ (0.5e).
Such a bond $b$ is contained in $\widetilde{\mathrm{Cp}}^{\sim1}$ (0.3c), hence in $\widetilde{\mathrm{Nb}}^{\sim1}$, and consequently $\mathrm{pr}_1(b)$ is contained in $\mathrm{pr}_1\left(\widetilde{\mathrm{Nb}}^{\sim1}\right)=Y_1^{\sim(j+d_U)}$.
Two connections in $\tilde{G}_k(T_1)$ that coincide on the bonds contained in $Y_1^{\sim(j+d_U)}$ thus coincide on the bonds that meet $\mathrm{pr}_1\left(\widetilde{\mathrm{Cp}}\right)$.
The second of the two identities above, applied to each of them with one and the same $A_1$, shows that they give the same row at $t$.
This is the second assertion of (a).

\paragraph{Step 3 (the parameter families).}
By Proposition~27.1\,(b), the map $\mathrm{pr}_2$ is injective on the parameter cubes contained in $\tilde{Y}$; hence $s^{(2)}$ is well defined.
Each value of $s^{(2)}$ is a value of $s^{(1)}$ or $0$, and so $s^{(2)}\in\mathrm{Pd}_k(T_2)$.
For every parameter cube $\widetilde{\mathrm{pc}}$ of the kind of $\mathrm{Op}$ that is contained in $N(\tilde{Y})$ one has
\begin{equation}\tag{3.5}
s^{(2)}_{\mathrm{pr}_2(\widetilde{\mathrm{pc}})}=s^{(1)}_{\mathrm{pr}_1(\widetilde{\mathrm{pc}})}=\left(\mathrm{pj}^{*}s^{(1)}\right)_{\mathrm{pr}_2(\widetilde{\mathrm{pc}})}.
\end{equation}
The second equality holds because $\left(\mathrm{pj}^{*}s\right)_{\mathrm{pc}}=s_{\mathrm{pj}(\mathrm{pc})}$ and $\mathrm{pj}\circ\mathrm{pr}_2=\mathrm{pr}_1$ (Proposition~26.1).
As for the first equality, it is the definition of $s^{(2)}$ when $\widetilde{\mathrm{pc}}$ is contained in $\tilde{Y}$.
Suppose that $\widetilde{\mathrm{pc}}$ is not contained in $\tilde{Y}$.
If one had $\mathrm{pr}_2(\widetilde{\mathrm{pc}})=\mathrm{pr}_2(\widetilde{\mathrm{pc}}')$ for a parameter cube $\widetilde{\mathrm{pc}}'$ contained in $\tilde{Y}$, Proposition~27.1\,(b) would give $\widetilde{\mathrm{pc}}=\widetilde{\mathrm{pc}}'$, which is impossible; hence $s^{(2)}_{\mathrm{pr}_2(\widetilde{\mathrm{pc}})}=0$.
If $\mathrm{pr}_1(\widetilde{\mathrm{pc}})$ were contained in $Y_1=\mathrm{pr}_1(\tilde{Y})$, it would have a lift contained in $\tilde{Y}$ (Proposition~27.1\,(d) with $\tilde{S}=\tilde{Y}$), and by Proposition~27.1\,(b) this lift would be $\widetilde{\mathrm{pc}}$.
So $\mathrm{pr}_1(\widetilde{\mathrm{pc}})$ is not contained in $Y_1$, and $s^{(1)}_{\mathrm{pr}_1(\widetilde{\mathrm{pc}})}=0$, since $s^{(1)}$ vanishes off $Y_1$.
This proves (3.5).

A parameter cube of $T_2$ that is contained in $\mathrm{pr}_2\left(\widetilde{\mathrm{Cp}}\right)$ is of the form $\mathrm{pr}_2(\widetilde{\mathrm{pc}})$ with $\widetilde{\mathrm{pc}}$ contained in $\widetilde{\mathrm{Cp}}$ (Proposition~27.1\,(d) with $\tilde{S}=\widetilde{\mathrm{Cp}}$), and $\widetilde{\mathrm{Cp}}$ lies in $N(\tilde{Y})$ by (3.1).
Hence (3.5) implies
\begin{equation}\tag{3.6}
\left(\mathrm{pj}^{*}s^{(1)}\right)\cdot\mathbf{1}_{\mathrm{pr}_2(\widetilde{\mathrm{Cp}})}=s^{(2)}\cdot\mathbf{1}_{\mathrm{pr}_2(\widetilde{\mathrm{Cp}})}.
\end{equation}
The family $\mathrm{pj}^{*}s^{(1)}$ belongs to $\mathrm{Pd}_k(T_2)$ (Proposition~26.1), and it vanishes off $\mathrm{pj}^{-1}(\mathrm{Nb}_1)$: if a parameter cube $\mathrm{pc}$ of $T_2$ has a site outside $\mathrm{pj}^{-1}(\mathrm{Nb}_1)$, then $\mathrm{pj}(\mathrm{pc})$ has a site outside $\mathrm{Nb}_1$, and $s^{(1)}_{\mathrm{pj}(\mathrm{pc})}=0$ by Step~2.
The family $s^{(2)}$ vanishes off $\mathrm{Nb}_2$, for it can differ from $0$ only at cubes $\mathrm{pr}_2(\widetilde{\mathrm{pc}})$ with $\widetilde{\mathrm{pc}}$ contained in $\tilde{Y}$, and these lie in $Y_2\subset\mathrm{Nb}_2$.

\paragraph{Step 4 (clause (b)).}
By \cite[Lemma~17.7]{QEDI}, $\mathrm{pj}^{*}A^{(1)}$ belongs to $\tilde{G}_k(T_2)$.
Every bond of $T_2$ that meets $\mathrm{pr}_2\left(\widetilde{\mathrm{Cp}}\right)$ is of the form $\mathrm{pr}_2(b)$ with a bond $b$ that meets $\widetilde{\mathrm{Cp}}$ (0.5e); this $b$ is contained in $\widetilde{\mathrm{Cp}}^{\sim1}$ (0.3c), hence in $\widetilde{\mathrm{Nb}}^{\sim1}=\tilde{Y}^{\sim(j+d_U)}$.
The hypothesis on $A^{(2)}$ and (0.5f) give
\[
A^{(2)}(\mathrm{pr}_2b)=A^{(1)}(\mathrm{pr}_1b)=\left(\mathrm{pj}^{*}A^{(1)}\right)(\mathrm{pr}_2b).
\]
Thus $A^{(2)}$ and $\mathrm{pj}^{*}A^{(1)}$ coincide on the bonds that meet $\mathrm{pr}_2\left(\widetilde{\mathrm{Cp}}\right)$.

Proposition~26.4 on $T_2$ is now applied twice to $\mathrm{Op}^{(2)}$, each time with the component $\mathrm{pr}_2\left(\widetilde{\mathrm{Cp}}\right)$ of (3.4) and with the row index $\mathrm{pr}_2(\tilde{t})$, for which (3.2) holds.
The first application uses the union $\mathrm{pj}^{-1}(\mathrm{Nb}_1)$, the family $\mathrm{pj}^{*}s^{(1)}$ and the connections $A=\mathrm{pj}^{*}A^{(1)}$, $A_1=A^{(2)}$; the second uses the union $\mathrm{Nb}_2$, the family $s^{(2)}$ and $A=A_1=A^{(2)}$.
Together with (3.6) they give
\begin{equation}\tag{3.7}
\begin{aligned}
\mathrm{Op}^{(2)}\left(\mathrm{pj}^{*}s^{(1)},\mathrm{pj}^{*}A^{(1)}\right)(\tilde{t},\,\cdot\,)
&=\mathrm{Op}^{(2)}\left(\left(\mathrm{pj}^{*}s^{(1)}\right)\cdot\mathbf{1}_{\mathrm{pr}_2(\widetilde{\mathrm{Cp}})},A^{(2)}\right)(\tilde{t},\,\cdot\,)\\
&=\mathrm{Op}^{(2)}\left(s^{(2)}\cdot\mathbf{1}_{\mathrm{pr}_2(\widetilde{\mathrm{Cp}})},A^{(2)}\right)(\tilde{t},\,\cdot\,)\\
&=\mathrm{Op}^{(2)}\left(s^{(2)},A^{(2)}\right)(\tilde{t},\,\cdot\,),
\end{aligned}
\end{equation}
and these rows vanish at the columns that are not contained in $\mathrm{pr}_2\left(\widetilde{\mathrm{Cp}}\right)$.
As $\mathrm{pr}_2\left(\widetilde{\mathrm{Cp}}\right)$ is contained in $\mathrm{Nb}_2=Y_2^{\sim(j+c_{l})}$, the last assertion of (b) follows.

Now let the column index $\tilde{t}'$ be contained in $N(\tilde{Y})$.
By (8.0), where $\widetilde{\mathrm{Op}}\left(s^{(1)},A^{(1)}\right)=\mathrm{Op}^{(2)}\left(\mathrm{pj}^{*}s^{(1)},\mathrm{pj}^{*}A^{(1)}\right)$,
\begin{multline*}
\mathrm{Op}^{(1)}\left(s^{(1)},A^{(1)}\right)(\tilde{t},\tilde{t}')\\
=\sum_{w\in(\mathbb{Z}/\mathrm{Jr}\,\mathbb{Z})^3}\boldsymbol{\chi}(w)\cdot\mathrm{Op}^{(2)}\left(\mathrm{pj}^{*}s^{(1)},\mathrm{pj}^{*}A^{(1)}\right)\left(\tilde{t},\tilde{t}'+\ell_1\cdot w\right).
\end{multline*}
Suppose that the summand of index $w$ does not vanish.
Then, by (3.7), the index $\mathrm{pr}_2\left(\tilde{t}'+\ell_1\cdot w\right)$ is contained in $\mathrm{pr}_2\left(\widetilde{\mathrm{Cp}}\right)$, and Proposition~27.1\,(d) with $\tilde{S}=\widetilde{\mathrm{Cp}}$ provides a lift of it contained in $\widetilde{\mathrm{Cp}}$: for some $w'\in\mathbb{Z}^3$ the index
\[
\tilde{t}'+\ell_1\cdot w+\ell_2\cdot w'=\tilde{t}'+\ell_1\cdot\left(w+\mathrm{Jr}\cdot w'\right)
\]
is contained in $\widetilde{\mathrm{Cp}}$, hence in $N(\tilde{Y})$.
For a point $u$ of $\operatorname{supp}\tilde{t}'$, both $u$ and $u+\ell_1\cdot\left(w+\mathrm{Jr}\cdot w'\right)$ are then sites of $N(\tilde{Y})$, and Proposition~27.1\,(b) yields $w+\mathrm{Jr}\cdot w'=0$.
Only the class $w=0$ contributes, therefore, and for its representative $w=0$ the factor is $\boldsymbol{\chi}(0)=1$.
With (3.7) one obtains
\begin{align*}
\mathrm{Op}^{(1)}\left(s^{(1)},A^{(1)}\right)(\tilde{t},\tilde{t}')
&=\mathrm{Op}^{(2)}\left(\mathrm{pj}^{*}s^{(1)},\mathrm{pj}^{*}A^{(1)}\right)(\tilde{t},\tilde{t}')\\
&=\mathrm{Op}^{(2)}\left(s^{(2)},A^{(2)}\right)(\tilde{t},\tilde{t}').
\end{align*}

\paragraph{Step 5 (clause (c)).}
Here $\mathrm{Op}$ is $S_{k,y}$ with $y\ge0$ or, when $k\ge1$, $M_k$.
By (0.3b) and (0.4c), the set $\tilde{X}^{\sim1}$ is contained in $N(\tilde{Y})$.

\emph{The polymers over $\mathrm{pr}_1(\tilde{X})$.}
The polymers $X'$ of $T_2$ with $\mathrm{pj}(X')=\mathrm{pr}_1(\tilde{X})$ are the sets $\mathrm{pr}_2(\tilde{X}+v)$, $v\in\ell_1\mathbb{Z}^3/\ell_2\mathbb{Z}^3$, and distinct classes $v$ give distinct sets.
Indeed, $\mathrm{pr}_2(\tilde{X}+v)$ is a polymer, and
\[
\mathrm{pj}\left(\mathrm{pr}_2(\tilde{X}+v)\right)=\mathrm{pr}_1(\tilde{X}+v)=\mathrm{pr}_1(\tilde{X}).
\]
Suppose that $\mathrm{pr}_2(\tilde{X}+v)$ and $\mathrm{pr}_2(\tilde{X}+v')$ have an $M$-cube in common.
Then $\tilde{\square}+v=\tilde{\square}''+v'+\ell_2\cdot w'$ with $M$-cubes $\tilde{\square}$, $\tilde{\square}''$ of $\tilde{X}$, so that $\tilde{\square}$ lies in $N(\tilde{Y})$ and in $N(\tilde{Y})+\left(v'-v+\ell_2\cdot w'\right)$.
By Proposition~27.1\,(b), $v'-v+\ell_2\cdot w'=0$: the classes of $v$ and $v'$ coincide.
Conversely, let $X'$ be a polymer of $T_2$ with $\mathrm{pj}(X')=\mathrm{pr}_1(\tilde{X})$.
An $M$-cube of $X'$ is of the form $\mathrm{pr}_2(\tilde{\square}')$, and $\mathrm{pr}_1(\tilde{\square}')=\mathrm{pj}\left(\mathrm{pr}_2(\tilde{\square}')\right)$ is an $M$-cube of $\mathrm{pr}_1(\tilde{X})$; therefore $\tilde{\square}'=\tilde{\square}+v$ with an $M$-cube $\tilde{\square}$ of $\tilde{X}$ and some $v\in\ell_1\mathbb{Z}^3$.
Let two $M$-cubes $\mathrm{pr}_2(\tilde{\square}_1+v_1)$ and $\mathrm{pr}_2(\tilde{\square}_2+v_2)$ of $X'$ have a common face.
By relation no.~9 of the list in the proof of Proposition~27.1\,(f), there is a $w'$ such that $\tilde{\square}_1+v_1$ and $\tilde{\square}_2+v_2+\ell_2\cdot w'$ have a common face.
Put $v:=v_2+\ell_2\cdot w'-v_1$.
The cube $\tilde{\square}_2+v$ has a common face with $\tilde{\square}_1$; so it lies in $\tilde{X}^{\sim1}$, and hence in $N(\tilde{Y})$, and it lies also in $N(\tilde{Y})+v$.
By Proposition~27.1\,(b), $v=0$; in other words, $v_1$ and $v_2$ are congruent modulo $\ell_2\mathbb{Z}^3$.
Since $X'$ is connected through faces, all its cubes lie in one set $\mathrm{pr}_2(\tilde{X}+v)$.
The map $\mathrm{pj}$ is injective on the $M$-cubes of $\mathrm{pr}_2(\tilde{X}+v)$, because $\mathrm{pj}\circ\mathrm{pr}_2=\mathrm{pr}_1$ is injective on $\tilde{X}+v$, a subset of $N(\tilde{Y})+v$; and it maps these cubes onto the $M$-cubes of $\mathrm{pr}_1(\tilde{X})=\mathrm{pj}(X')$.
Consequently $X'=\mathrm{pr}_2(\tilde{X}+v)$.

\emph{The kernel.}
Proposition~26.1 and (8.0) give
\begin{multline*}
\mathrm{Op}^{(1)}\left(\mathrm{pr}_1\tilde{X},A^{(1)}\right)(\tilde{t},\tilde{t}')\\
=\sum_{w\in(\mathbb{Z}/\mathrm{Jr}\,\mathbb{Z})^3}\boldsymbol{\chi}(w)\cdot\sum_{v\in\ell_1\mathbb{Z}^3/\ell_2\mathbb{Z}^3}\mathrm{Op}^{(2)}\left(\mathrm{pr}_2(\tilde{X}+v),\mathrm{pj}^{*}A^{(1)}\right)\left(\tilde{t},\tilde{t}'+\ell_1\cdot w\right).
\end{multline*}
By W12$^{\mathrm{ch}}$ on $T_2$, the summand of index $(w,v)$ vanishes unless $\mathrm{pr}_2(\tilde{t})$ and $\mathrm{pr}_2\left(\tilde{t}'+\ell_1\cdot w\right)$ lie over $\mathrm{pr}_2(\tilde{X}+v)$.
Let $u,u'\in\tilde{X}$ be the support points of $\tilde{t}$ and of $\tilde{t}'$.
If $\mathrm{pr}_2(u)$ lies in $\mathrm{pr}_2(\tilde{X}+v)$, then $u=u''+v+\ell_2\cdot w'$ with $u''\in\tilde{X}$; as $u$ and $u''$ are sites of $N(\tilde{Y})$, Proposition~27.1\,(b) gives $v+\ell_2\cdot w'=0$, and so $v$ is the class $0$.
If $\mathrm{pr}_2\left(u'+\ell_1\cdot w\right)$ lies in $\mathrm{pr}_2(\tilde{X})$, then $u'+\ell_1\cdot w+\ell_2\cdot w'\in\tilde{X}$ for some $w'$, Proposition~27.1\,(b) gives $\ell_1\cdot\left(w+\mathrm{Jr}\cdot w'\right)=0$, and so $w$ is the class $0$.
Hence
\[
\mathrm{Op}^{(1)}\left(\mathrm{pr}_1\tilde{X},A^{(1)}\right)(\tilde{t},\tilde{t}')=\mathrm{Op}^{(2)}\left(\mathrm{pr}_2\tilde{X},\mathrm{pj}^{*}A^{(1)}\right)(\tilde{t},\tilde{t}').
\]
Every bond of $T_2$ that meets $\mathrm{pr}_2(\tilde{X})$ is of the form $\mathrm{pr}_2(b)$ with a bond $b$ that meets $\tilde{X}$ (0.5e), and on it the hypothesis and (0.5f) give
\[
A^{(2)}(\mathrm{pr}_2b)=A^{(1)}(\mathrm{pr}_1b)=\left(\mathrm{pj}^{*}A^{(1)}\right)(\mathrm{pr}_2b).
\]
The second part of W12$^{\mathrm{ch}}$ on $T_2$ allows $\mathrm{pj}^{*}A^{(1)}$ to be replaced by $A^{(2)}$.
This proves (c), for every $y\ge0$.

\paragraph{Step 6 ($\bar{H}_k$).}
On both tori
\[
\bar{H}^{(i)}_k(s'',A)=\mathsf{C}^{-1}\cdot H^{(i)}_k(s'',-A)\cdot\mathsf{C}
\]
(Section~G4.0.4; for $i=1$, Proposition~26.1), so that the kernel of $\bar{H}^{(i)}_k(s'',A)$ is $\mathsf{C}^{-1}\cdot H^{(i)}_k(s'',-A)(t,t')\cdot\mathsf{C}$.
By \cite[Lemma~17.6]{QEDI}, $-A^{(i)}$ lies in $\tilde{G}_k(T_i)$; and $-A^{(1)}$, $-A^{(2)}$ have equal lifts on a set of bonds whenever $A^{(1)}$, $A^{(2)}$ have.
Clauses (a) and (b) for $H_k$ at the connections $-A^{(1)}$, $-A^{(2)}$ are therefore clauses (a) and (b) for $\bar{H}_k$ at $A^{(1)}$, $A^{(2)}$.
\end{proof}

\paragraph{Level 0.}
At level $0$ the system consists of (p2), $\Gamma_0(s'',A)$ and $S_{0,y}(X,A)$; the families (p1) and (p3) are the diagonal families $H_0(s')$ and $H_0(s'',A)$ of Section~I.5.8, and there is no family $M_0$.
The row and column indices of (p3), $\Gamma_0$ and $S_{0,y}$ are fermion sites $x=(c,g)$ of $\mathbb{L}^{\mathrm{f}}_0$, with $\operatorname{supp}x=\{c\}$ and with $\square(x)$ the $M$-cube of $c$.
Their kernels are those of (I.1.4), with the weight $\frac{1}{48}$; by (I.ke1) the kernel of $H_0(s'',A)$ is
\[
48\cdot\mathbf{1}[x=x']\cdot\prod\left\{s''_{\mathrm{pc}}:\ \mathrm{pc}\in\pi_0,\ \mathrm{pc}\text{ meets }K_c\right\}\cdot\mathbf{1}_4,
\]
and for these families $\boldsymbol{\chi}(w)=(-1)^{w_1+w_2+w_3}$.
The diagonal families on $T_2$ have W8 (by W17); Propositions~26.1 and~26.3\,(a) hold for them, and their projections are the diagonal families on $T_1$ (Proposition~26.3\,(c)).
Steps~1--4 and~6 therefore cover them.
No step moves the label $g$: it is a spectator of $\mathrm{pr}_i$, $\mathrm{pj}$, $\mathrm{pj}^{*}$ and $T_v$.

\paragraph{Uniformity.}
In Proposition~27.3 no constant is introduced and no bound is proved: its clauses assert that certain kernels vanish, that rows depend on the connection only through its values on certain bonds, and that certain kernels are equal.
The numbers $c_{l}=2L$ and $d_U=2L+1$ that occur in these clauses are the constants of (C4) in Section~I.8.2; they depend on $L$ alone and not on the data (TT).
The clauses hold for all data with (TT) and $k<k_m$, for all admissible $\tilde{Y}$, for all integers $j$ with $0\le j\le5\cdot d_U$, and for all $y\ge0$.

\subsection*{4. Proposition 27.4}

\subsubsection*{Proof of Proposition 27.4}

By Proposition~27.1\,(a), the set $\tilde{\mathbf{Y}}$ and every non-empty union $\tilde{X}$ of $M$-cubes contained in $\tilde{Y}$ are admissible, and $N(\tilde{X})\subset N(\tilde{Y})$.
By Proposition~27.2\,(a), $\mathrm{idn}[\tilde{X}]$ is the restriction of $\mathrm{idn}=\mathrm{idn}[\tilde{Y}]$.
From (0.4b), (0.4c) and (1.1) one obtains
\begin{equation}\tag{4.0}
\begin{gathered}
\mathrm{sdiam}(\tilde{X})\le\ell_1-14\cdot d_U\cdot M-\eta_k\quad\text{for}\ \tilde{X}\subset\tilde{Y},\\
\mathrm{sdiam}\left(N(\tilde{Y})\right)\le\ell_1-2\cdot d_U\cdot M-\eta_k,\qquad
\tilde{Y}^{\sim a}\subset N(\tilde{Y})\quad\left(\text{for}\ a\le6\cdot d_U\right).
\end{gathered}
\end{equation}
Consequently the following three margins hold for $i=1,2$:
$\mathrm{sdiam}\left(N(\tilde{Y})\right)+M<\ell_i$, the condition of Section~I.1.9 under which $\mathrm{cov}_{\ell_i}$ is defined on $\mathrm{Grc}_k\left(N(\tilde{Y})\right)$;
$\mathrm{sdiam}(\tilde{X})+2\,\mathrm{rg}+M<\ell_i$ for $\mathrm{rg}\le1$, the margin in Proposition~25.2;
and $\mathrm{sdiam}(\tilde{\mathbf{Y}})+3M<\ell_i$, the margin in Proposition~25.3\,(b).

\paragraph{4.0 Rows.}
Let $O$ be a linear operator between two of the function spaces of Section~I.1.9 on $T_i$.
Its arguments $f$ are the $\ell_i\mathbb{Z}^3$-periodic functions (photon objects), respectively the antiperiodic functions (spinor objects), on a covering index set.
Let $t$ be a covering row index, and let $\mathrm{ro}(t,t')$ be a family of coefficients (numbers for photon objects, $4\times4$ matrices for spinor objects), indexed by the column indices $t'$ contained in $N(\tilde{Y})$, of which only finitely many are different from $0$.
As in the definition stated in Proposition~27.4\,(a), one says that \emph{the row of $O$ at $t$, written with all indices in $N(\tilde{Y})$, is $\mathrm{ro}(t,\,\cdot\,)$} if
\begin{equation}\tag{4.1}
(Of)(t)=\sum_{t'\ \text{contained in}\ N(\tilde{Y})}\mathrm{ro}(t,t')\cdot f(t')\qquad\text{for all $f$ in the source space on}\ T_i.
\end{equation}
The column indices contained in $N(\tilde{Y})$ are pairwise incongruent modulo $\ell_i\mathbb{Z}^3$, by Proposition~27.1\,(b).
They therefore belong to a set of representatives of the column indices.
The values of $f$ on such a set are free coordinates of the source space; hence $\mathrm{ro}(t,\,\cdot\,)$ is determined by $O$ and $t$.
By the definition of the kernels (Section~G4.0.4; (I.1.4)), (4.1) amounts to
\begin{equation}\tag{4.2}
\begin{gathered}
O(t,t')=\frac{\mathrm{ro}(t,t')}{\nu(t')}\quad\text{for}\ t'\ \text{contained in}\ N(\tilde{Y});\\
O(t,t'')=0\quad\text{for every column index}\ t''\ \text{such that no}\ t''+v\ \text{with}\ v\in\ell_i\mathbb{Z}^3\\
\text{is contained in}\ N(\tilde{Y}),
\end{gathered}
\end{equation}
where $\nu(t')=\frac{\eta_k^3}{48}$ if $t'$ is a fermion site (see (I.1.3)), and $\nu(t')=1$ if $t'$ is a unit site ($k\ge1$) or a basis element.
(For columns in $\mathrm{Pzt}_k$ only the coefficients $\mathrm{ro}$ enter; compare Section~I.5.2.)
The second half of (4.2) is the assertion that the row vanishes at the columns of $T_i$ that have no lift in $N(\tilde{Y})$.

\emph{Range.}
In the parentheses of (a-1)--(a-4) the word ``range'' has the meaning of Section~I.1.9\,(f4); the range of a relation, in Proposition~27.1, is a different notion.
A number $\mathrm{rg}\ge0$ is called \emph{a range of the row at $t$} if each of the following sites has sup-distance at most $\mathrm{rg}$ from every point of $\operatorname{supp}t$: the points of the supports of the columns at which the row does not vanish, the two ends of every bond on which the row reads the connection, and the sites through which the row reads the union.
In that case the row vanishes at the columns $t'$ with $\operatorname{dist}(\operatorname{supp}t,\operatorname{supp}t')>\mathrm{rg}$, and the bonds and the sites that it reads lie within the distance $\mathrm{rg}$ of $\operatorname{supp}t$; this is Section~I.1.9\,(f4) for this row, with the number $\mathrm{rg}$.
``Range $<a$'' means that some number $\mathrm{rg}<a$ is a range of each of the rows named, and ``range $M$'' means that $M$ is such a number.

\begin{named}{Lemma 4.1}
Let $\mathrm{Op}$ be an operator of the table of Proposition~25.1, let $\mathrm{Op}_{\mathbb{L}}$ be its operator on the covering lattice, and let $t$ be a covering row index.
Two pairs $(\tilde{B},\mathrm{Ye})$ and $(\tilde{B}',\mathrm{Ye}')$, each consisting of a complex bond function on $\mathbb{L}_k$ and of a union of $M$-cubes of $\mathbb{L}_k$, are said to \emph{agree where the row at $t$ reads} if $\tilde{B}$ and $\tilde{B}'$ coincide on the bonds on which the row of $\mathrm{Op}_{\mathbb{L}}$ at $t$ reads the connection, and $\mathrm{Ye}$, $\mathrm{Ye}'$ contain the same ones among the sites through which this row reads the union (both are named in the last entry for $\mathrm{Op}$ in the table of Proposition~25.1; these sets of bonds and of sites depend on $t$ only).
For $i=1,2$ let $B_i$ be a complex bond function and $\mathrm{Yu}_i$ a union of $M$-cubes on $T_i$, with the $\ell_i\mathbb{Z}^3$-periodic extensions
\[
\tilde{B}_i(b):=B_i(\mathrm{pr}_ib),\qquad\mathrm{Ye}_i:=\mathrm{pr}_i^{-1}(\mathrm{Yu}_i).
\]
Assume:
\begin{itemize}
\item[(col)] for $i=1,2$, every column index at which the row of $\mathrm{Op}_{\mathbb{L}}(\tilde{B}_i,\mathrm{Ye}_i)$ at $t$ does not vanish is contained in $N(\tilde{Y})$;
\item[(rd)] $(\tilde{B}_1,\mathrm{Ye}_1)$ and $(\tilde{B}_2,\mathrm{Ye}_2)$ agree where the row at $t$ reads.
\end{itemize}
Then the rows at $t$ of $\mathrm{Op}(B_1,\mathrm{Yu}_1)$ on $T_1$ and of $\mathrm{Op}(B_2,\mathrm{Yu}_2)$ on $T_2$, written with all indices in $N(\tilde{Y})$, coincide; and they are the row at $t$ of $\mathrm{Op}_{\mathbb{L}}(\tilde{B},\mathrm{Ye})$ for every pair $(\tilde{B},\mathrm{Ye})$ that agrees with $(\tilde{B}_1,\mathrm{Ye}_1)$ where the row at $t$ reads.
\end{named}

\begin{proof}
(i) Let $(\tilde{B},\mathrm{Ye})$ and $(\tilde{B}',\mathrm{Ye}')$ agree where the row at $t$ reads.
Then the rows at $t$ of $\mathrm{Op}_{\mathbb{L}}(\tilde{B},\mathrm{Ye})$ and of $\mathrm{Op}_{\mathbb{L}}(\tilde{B}',\mathrm{Ye}')$ coincide.
Indeed, the word ``reads'' in Section~I.1.9\,(f4) has the meaning of \cite[Definition~5.5(1)]{QEDI}: it compares two values of one object, all other objects being held fixed, and it states that the row takes the same value at any two admitted values of this object that coincide at the places named.
The row at $t$ of $\mathrm{Op}_{\mathbb{L}}(\tilde{B},\mathrm{Ye})$ equals that of $\mathrm{Op}_{\mathbb{L}}(\tilde{B}',\mathrm{Ye})$, since the union is fixed and $\tilde{B}=\tilde{B}'$ on the bonds that are read.
The latter equals the row at $t$ of $\mathrm{Op}_{\mathbb{L}}(\tilde{B}',\mathrm{Ye}')$, since the connection is fixed and $\mathrm{Ye}$, $\mathrm{Ye}'$ contain the same ones among the sites that are read.
The pair $(\tilde{B}',\mathrm{Ye})$ is admitted, because Proposition~25.1 defines the operator $\mathrm{Op}_{\mathbb{L}}$ for every complex bond function on $\mathbb{L}_k$ and every union of $M$-cubes of $\mathbb{L}_k$, without any condition of periodicity.

(ii) By Section~I.1.9\,(f3),
\[
\left(\mathrm{Op}(B_i,\mathrm{Yu}_i)f\right)(t)=\left(\mathrm{Op}_{\mathbb{L}}(\tilde{B}_i,\mathrm{Ye}_i)f\right)(t)
\]
for $f$ in the source space on $T_i$, where $f$ is regarded as a function on the covering index set.
By Section~I.1.9\,(f1) the right side is a finite sum $\sum_{t'}\mathrm{ro}_i(t,t')\cdot f(t')$ over covering column indices $t'$.
Because of (col) this sum has the form (4.1).
By (rd) and (i), $\mathrm{ro}_1(t,\,\cdot\,)=\mathrm{ro}_2(t,\,\cdot\,)$, and by (i) this is the row at $t$ of $\mathrm{Op}_{\mathbb{L}}(\tilde{B},\mathrm{Ye})$ for every pair $(\tilde{B},\mathrm{Ye})$ named in the last assertion.
\end{proof}

\paragraph{4.2 Dictionary.}
Let $\tilde{X}$ be admissible.
The map $\mathrm{cov}_{\ell_i}$ is the homomorphism of $\mathrm{Grc}_k\left(N(\tilde{X})\right)$ into the Grassmann algebra of $T_i$ that sends every generator to the generator of $T_i$ with the same index in the sense of (I.AP) (Section~I.1.9; Section~G3.5.0).
Let $\mathrm{Grc}'_k\left(N(\tilde{X})\right)$ be $\mathrm{Grc}_k\left(N(\tilde{X})\right)$ if $k\ge1$, and, if $k=0$, its subalgebra generated by $\psi$, $\chi$ and $W$.
(By the last paragraph of Section~I.1.9, the algebra $\mathrm{Grc}_0$ has in addition the generators $\Psi(u)$, $u=(x,\beta,\omega)$, which do not occur in Proposition~27.2\,(D2).)
The generators of $\mathrm{Grc}'_k\left(N(\tilde{X})\right)$ are indexed by the set $\mathrm{Ix}(\tilde{X})$ of the proof of Proposition~27.2\,(a).
Hence $\mathrm{cov}_{\ell_i}$ maps $\mathrm{Grc}'_k\left(N(\tilde{X})\right)$ isomorphically onto $\mathrm{Gr}_i(\tilde{X})$, and
\begin{equation}\tag{4.3}
\begin{gathered}
\mathrm{idn}[\tilde{X}]=\mathrm{cov}_{\ell_2}\circ\mathrm{cov}_{\ell_1}^{-1}\quad\text{on}\ \mathrm{Gr}_1(\tilde{X});\\
\mathrm{idn}\left(\mathrm{cov}_{\ell_1}(\mathrm{Fcv})\right)=\mathrm{cov}_{\ell_2}(\mathrm{Fcv})\quad\text{for all}\ \mathrm{Fcv}\in\mathrm{Grc}'_k\left(N(\tilde{Y})\right).
\end{gathered}
\end{equation}
For a set $\tilde{U}$ of sites contained in $N(\tilde{Y})$, the algebra $\mathrm{Grc}_k(\tilde{U})$ is the subalgebra of $\mathrm{Grc}_k\left(N(\tilde{Y})\right)$ generated by the generators all of whose photon sites lie in $\tilde{U}$, and $\mathrm{cov}_{\ell_i}$ on it is the restriction.

\paragraph{(a)}
Let $\tilde{B}_i(b):=B_i(\mathrm{pr}_ib)$.
The hypothesis of (a) states that
\begin{equation}\tag{4.4}
\tilde{B}_1(b)=\tilde{B}_2(b)\qquad\text{for every bond $b$ of $\mathbb{L}_k$ that meets}\ \tilde{\mathbf{Y}}.
\end{equation}

\emph{Row indices.}
By Proposition~27.2\,(e), applied with $\tilde{\mathbf{Y}}$ in the place of $\tilde{Y}$, the polymers of $T_i$ contained in $\mathbf{Y}_i$ are the sets $\mathrm{pr}_i(\tilde{X})$, where $\tilde{X}$ is a non-empty connected union of $M$-cubes contained in $\tilde{\mathbf{Y}}$; and the unit sites and the fermion sites of (over) $\mathrm{pr}_i(\tilde{X})$ are the images of those of (over) $\tilde{X}$.
A pair of $\mathrm{Pzt}_k$ that is admissible for $\tilde{X}$ has both photon sites in $\tilde{X}^{\sim1}$: one of them lies in $\tilde{X}$, and the other has sup-distance at most $M$ from it (0.3c).
By Proposition~27.2\,(b), every pair that is admissible for $\mathrm{pr}_i(\tilde{X})$ is therefore the image of exactly one pair of $\mathrm{Pzt}_k$ that is admissible for $\tilde{X}$, namely of its lift with both photon sites in $\tilde{X}^{\sim1}$.
Conversely, the image of a pair that is admissible for $\tilde{X}$ is admissible for $\mathrm{pr}_i(\tilde{X})$, since $\mathrm{pr}_i$ maps the sites and the paths in $\tilde{X}$ into $\mathrm{pr}_i(\tilde{X})$ and $x^{(3)}$ to $(\mathrm{pr}_ix)^{(3)}$.
This is Proposition~25.3\,(b1), according to which, moreover, a pair and its image are of the same kind.
In each of (a-1)--(a-5), therefore, the row indices on $T_1$ and on $T_2$ are the images of one and the same set of covering indices, and these covering indices are contained in $N(\tilde{Y})$.

\emph{Membership.}
For a non-empty union $\tilde{X}$ of $M$-cubes contained in $\tilde{Y}$ and a site $u$ of $\tilde{X}^{\sim1}$,
\begin{equation}\tag{4.5}
u\in\mathrm{pr}_i^{-1}(\mathrm{pr}_i\tilde{X})=\tilde{X}+\ell_i\mathbb{Z}^3\quad\text{if and only if}\quad u\in\tilde{X}.
\end{equation}
Indeed, if $u$ lies in $\tilde{X}+v$ with $v\in\ell_i\mathbb{Z}^3$, then $u$ lies in $N(\tilde{Y})$, by (4.0), and in $N(\tilde{Y})+v$; so $v=0$ by Proposition~27.1\,(b).
(When $\tilde{X}$ is connected through faces, for instance $\tilde{X}=\tilde{\mathbf{Y}}$, and $u$ has sup-distance at most $\eta_k$ from $\tilde{X}$, the equivalence (4.5) is also the last clause of Proposition~25.3\,(b1), applied with $\tilde{X}$ in the place of $\tilde{\mathbf{Y}}$; the margin $\mathrm{sdiam}(\tilde{X})+3M<\ell_i$ of Proposition~25.3\,(b) follows from (4.0).
This is the situation of (a-3) below, where $\tilde{X}$ is connected through faces and $u$ is $y$ or a site $c_3$ at sup-distance at most $\eta_k$ from $\tilde{X}$.)

\emph{(a-1).}
According to Proposition~25.1, no.~12, the map $C$ is given by one formula; see \cite[Section~5.2.5]{QEDI} and \cite[Proposition~15.21(iv)]{QEDI}.
For a unit bond $b'$ of $\mathbb{Z}^3$ that is not central, $C(b',\Upsilon)=\Upsilon(b')$.
For the central bond $b_c$ of the face between $B(y')$ and $B(y'+L\cdot e_{\mu})$,
\[
C(b_c,\Upsilon)=-\sum\left\{\Upsilon(b):\ b\ \text{a non-central face bond of this face}\right\}-\sum_b\frac{n_c(b)}{L}\cdot\Upsilon(b),
\]
where the second sum extends over the interior bonds $b$ of direction $\mu$ of the two blocks.
Here $n_c(b)\in\{1,\dots,L-1\}$ is the number of the straight paths $\Gamma_{y'',y''+L\cdot e_{\mu}}$, $y''\in B(y')$, that contain $b$ (\cite[Section~5.2.5]{QEDI}; for the values $1,\dots,L-1$, \cite[Proposition~15.21(iv)]{QEDI}), and it depends only on the position of $b$ in the two blocks.
The row reads nothing, and $C(b',\Upsilon)=0$ unless $\operatorname{supp}\Upsilon$ lies in the union of the $L$-blocks of the ends of $b'$.
Let $b'$ be contained in $\tilde{\mathbf{Y}}^{\sim(1+2d_U)}$.
The $L$-block of an end of $b'$ lies in the $M$-cube of this end (0.2b), and this $M$-cube is a cube of $\tilde{\mathbf{Y}}^{\sim(1+2d_U)}$ (0.3a).
So the columns of the row are contained in $\tilde{\mathbf{Y}}^{\sim(1+2d_U)}$, which lies in $\tilde{Y}^{\sim(1+2d_U)}$ (0.3b), hence in $N(\tilde{Y})$ by (4.0), as $1+2\cdot d_U\le6\cdot d_U$.
Thus (col) holds, while (rd) is void.
Lemma~4.1 and (4.2) give (a-1).
(Since no connection is read, (a-1) is also Proposition~25.3\,(a) with $\tilde{U}=N(\tilde{Y})$, because $\mathrm{sdiam}\left(N(\tilde{Y})\right)+L-1<\ell_i$ by (4.0).)

\emph{Range.}
The union of the $L$-blocks of the ends of $b'$ is one $L$-block or, when the ends lie in two different $L$-blocks $B(y')$ and $B(y'+L\cdot e_{\mu})$, with $\mu$ the direction of $b'$, a box of unit sites with $2L$ sites in the direction $\mu$ and $L$ sites in each of the two other directions.
The ends of $b'$ lie in this set, and so do the points of $\operatorname{supp}\Upsilon$ for every column $\Upsilon$ with $C(b',\Upsilon)\ne0$.
Every point of $\operatorname{supp}\Upsilon$ therefore has sup-distance at most $2L-1$ from each end of $b'$, and the number $2L-1<2L$ is a range of the row.
(Every point of $\operatorname{supp}\Upsilon$ lies in the $L$-block of one of the two ends, at sup-distance at most $L-1$ from this end; hence $\operatorname{dist}(\operatorname{supp}b',\operatorname{supp}\Upsilon)\le L-1$, which is the number $\mathrm{rg}$ of Proposition~25.1, no.~12.)

\emph{(a-2).}
The rows are those of Proposition~25.1, no.~11, given by the display (I.1.7): for $(x,x')\in\mathrm{Pzt}_k$,
\begin{gather*}
\left(\delta_{\alpha,\tilde{B}}f\right)(x,x')=d_X(x,x')^{-\alpha}\cdot\left(U_{\tilde{B}}(x,x')\cdot f(x')-f(x)\right),\\
U_{\tilde{B}}(x,x')=\exp\left(\mathrm{i}\cdot e_k\cdot\eta_k\cdot\frac{1}{6}\sum_{\pi}\tilde{B}\left(\boldsymbol{\Gamma}^{\pi}(c,c')\right)\right),\qquad U_{\tilde{B}}(x,x')=1\ \text{if}\ c=c'.
\end{gather*}
The columns of the row are $x$ and $x'$, with the coefficients $-d_X^{-\alpha}\cdot\mathbf{1}_4$ and $d_X^{-\alpha}\cdot U_{\tilde{B}}(x,x')\cdot\mathbf{1}_4$, and the row reads $\tilde{B}$ on the bonds of the six paths.
(Inside a cluster one has $c=c'$ and $d_X=\eta_k$; the row is then $\eta_k^{-\alpha}\cdot\left(f(x')-f(x)\right)$, and no bond is read.)
Let $(x,x')$ be admissible for $\tilde{X}$, where $\tilde{X}$ is contained in $\tilde{\mathbf{Y}}$.
Both photon sites lie in $\tilde{X}^{\sim1}$, since one of them lies in $\tilde{X}$ and $|c-c'|\le M$ (0.3c); hence they lie in $N(\tilde{Y})$, and (col) holds.
If the pair is of kind (I), the six paths lie in $\tilde{X}$; their bonds are contained in $\tilde{X}$, and so they meet $\tilde{\mathbf{Y}}$.
If the pair is of kind (II), then $x'=x^{(3)}$, the six paths are the single bond $(c,c_3)$, and $c$ or $c_3$ lies in $\tilde{X}$; so this bond meets $\tilde{\mathbf{Y}}$.
By (4.4), (rd) holds.
Lemma~4.1 gives (a-2) for $\delta_{\alpha,B_i}$.
The connections $-B_1$ and $-B_2$ satisfy (4.4) as well, and this gives (a-2) for $\delta_{\alpha,-B_i}$.
(This is Proposition~25.3\,(b2).)

\emph{Range.}
One has $\operatorname{supp}(x,x')=\{c,c'\}$ with $|c-c'|\le M$, and the two columns have the supports $\{c\}$ and $\{c'\}$.
Every coordinate of a site $u$ of a path $\boldsymbol{\Gamma}^{\pi}(c,c')$ lies between the corresponding coordinates of $c$ and $c'$ (0.2c); so $|u-c|\le|c'-c|\le M$ and $|u-c'|\le M$.
Hence $M$ is a range of the row; it is the number $\mathrm{rg}$ of Proposition~25.1, no.~11.

\emph{(a-3)} ($k\ge1$).
Let $\tilde{X}$ be $\tilde{\mathbf{Y}}$ or a non-empty connected union of $M$-cubes contained in $\tilde{\mathbf{Y}}$, and let $y$ be a unit site of $\tilde{X}$.
The rows are those of Proposition~25.1, nos.~9 and~10, given by the display (23.1), which is (I.5.4).
With $\mu=\mu_k$, $x=(c,g)$ and $c_3=c+\eta_k\cdot ge_1$ they read
\begin{multline*}
\left(T^{(1)}(\tilde{B},\tilde{X})f\right)(y)=\frac{\eta_k^3}{48}\sum_{c\in\Delta_y}\sum_{g\in O_{\mathrm{cub}}}\varphi_{\tilde{B}}(y,c)\cdot\left\{\left(1+\frac{\mu_k}{b_k}\right)\cdot f(x)\right.\\
{}-\left(b_k\cdot\eta_k\right)^{-1}\cdot\left[\sum_{j\le2}M_j(x)\cdot\left(f\left(x^{(j)}\right)-f(x)\right)\right.\\
\left.\left.{}+\mathbf{1}[c_3\in\tilde{X}]\cdot M_3(x)\cdot\left(u_{\tilde{B}}(c,c_3)\cdot f\left(x^{(3)}\right)-f(x)\right)\vphantom{\sum_{j\le2}\frac{\mu_k}{b_k}}\right]\right\},
\end{multline*}
\begin{multline*}
\left(T^{(2)}(\tilde{B},\tilde{X})F\right)(y)\\
=-b_k^{-1}\cdot\eta_k^{\alpha-1}\cdot\frac{\eta_k^3}{48}\sum_{c\in\Delta_y}\sum_{g\in O_{\mathrm{cub}}}\mathbf{1}[c_3\notin\tilde{X}]\cdot\varphi_{\tilde{B}}(y,c)\cdot M_3(x)\cdot F\left(x,x^{(3)}\right),
\end{multline*}
\[
\bar{T}^{(j)}(\tilde{B},\tilde{X})=\mathsf{C}^{-1}\cdot T^{(j)}(-\tilde{B},\tilde{X})\cdot\mathsf{C},
\]
where
\[
M_j(x)=\frac{w_j}{2}\cdot\left(\rho+\gamma\cdot\left(g\hat{n}_j\right)\right).
\]
By Proposition~23.1\,(b) and Proposition~25.1, nos.~9 and~10, the columns are the sites $x$, $x^{(0)}$, $x^{(1)}$, $x^{(2)}$ over $\Delta_y$ and the sites $x^{(3)}$ with $c_3$ in the union ($\tilde{X}$, respectively $\mathrm{Ye}_i$), respectively the pairs $(x,x^{(3)})$ with $c\in\Delta_y$ and $c_3$ not in the union.
The rows read the connection on the bonds contained in $\Delta_y$ and on the bonds $(c,c_3)$ with $c\in\Delta_y$, and they read the union at $y$ and at the sites $c_3$ with $c\in\Delta_y$.
The block $\Delta_y$ lies in the $M$-cube of $y$ (0.2b), which is a cube of $\tilde{X}$, and the sites $c_3$ lie in $\tilde{X}^{\sim1}$ (0.3c).
So all columns are contained in $\tilde{X}^{\sim1}$, hence in $N(\tilde{Y})$, and (col) holds.
All bonds named have an end in $\Delta_y$, and hence they meet $\tilde{\mathbf{Y}}$.
Together with (4.4), and with (4.5) applied to the unions $\mathrm{Ye}_i=\tilde{X}+\ell_i\mathbb{Z}^3$ (the periodic extensions of $\mathrm{Yu}_i=\mathrm{pr}_i\tilde{X}$), this gives (rd).
By Lemma~4.1, the rows at $y$ of $T^{(j)}(B_i,\mathrm{pr}_i\tilde{X})$ coincide for $i=1$ and $i=2$.
By (4.5), the pair $(\tilde{B},\tilde{X})$ agrees with $(\tilde{B}_1,\mathrm{Ye}_1)$ where the row at $y$ reads, for every $\tilde{B}$ that coincides with $\tilde{B}_1$ on the bonds that meet $\tilde{\mathbf{Y}}$.
By the last assertion of Lemma~4.1, these rows are therefore the rows displayed above, formed with $\tilde{X}$ itself and with such a $\tilde{B}$.
The same holds for $\bar{T}^{(j)}$, with $-B_i$ in the place of $B_i$.
(For $\tilde{X}=\tilde{\mathbf{Y}}$ this is Proposition~25.3\,(b3); for a smaller $\tilde{X}$ it is Proposition~25.3\,(b3) with $\tilde{X}$ in the place of $\tilde{\mathbf{Y}}$.)

\emph{Range.}
The row index $y$ has the support $\{y\}$.
Two kinds of sites occur in the rows.
The first are the sites $c$ of $\Delta_y$, for which $|c-y|\le\frac{1-\eta_k}{2}$ by (0.2a).
The second are the sites $c_3=c+\eta_k\cdot g\,e_1$ with $c\in\Delta_y$, for which
\[
|c_3-y|\le\frac{1-\eta_k}{2}+\eta_k=\frac{1+\eta_k}{2}.
\]
No other site is involved: every column has the support $\{c\}$, $\{c_3\}$ or $\{c,c_3\}$, both ends of every bond that is read are sites of these two kinds, and the union is read at $y$ and at sites $c_3$.
Consequently $\frac{1+\eta_k}{2}$ is a range of the rows (it is the number $\mathrm{rg}$ of nos.~9 and~10 of the table of Proposition~25.1).
Moreover $\frac{1+\eta_k}{2}<1$, since $\eta_k\le\frac{1}{L}<1$ for $k\ge1$.

\emph{(a-4)} ($k\ge1$).
In this clause $z_1$ and $z_2$ denote unit sites, and $y$ denotes the resolvent parameter.
Nos.~4 and~6 of the table of Proposition~25.1, together with the kernels (K1)--(K3) given there, yield
\begin{align*}
\left(Q_k(\tilde{B})J^{*}\right)(z_1,x)&=\mathbf{1}[c\in\Delta_{z_1}]\cdot\varphi_{\tilde{B}}(z_1,c)\cdot\mathbf{1}_4,\\
\left(JQ_k^{T}(-\tilde{B})\right)(x,z_1)&=\mathbf{1}[c\in\Delta_{z_1}]\cdot\varphi_{\tilde{B}}(z_1,c)^{-1}\cdot\mathbf{1}_4,\\
B_{k,y}(\tilde{B})(z_1,z_2)&=\left(b_k+i\gamma_5y\right)^{-1}\cdot\mathbf{1}[z_1=z_2]\\
&\qquad+\left[\left(b_k+b_L+i\gamma_5y\right)^{-1}-\left(b_k+i\gamma_5y\right)^{-1}\right]\cdot P(\tilde{B})(z_1,z_2),
\end{align*}
where
\[
P(\tilde{B})(z_1,z_2)=
\begin{cases}
L^{-3}\cdot U_{\tilde{B}}(y',z_1)^{-1}\cdot U_{\tilde{B}}(y',z_2)\cdot\mathbf{1}_4 & \text{if $z_1$ and $z_2$ lie in one block $B(y')$,}\\
0 & \text{otherwise.}
\end{cases}
\]
The first of these kernels is taken with respect to $\int\mathrm{d}x$, the other ones with respect to the sum over the unit sites.
Let the row index lie in $\tilde{\mathbf{Y}}$ if it is a unit site, and over $\tilde{\mathbf{Y}}$ if it is a fermion site.
The columns of the row of $Q_k(\tilde{B})J^{*}$ at $z_1$ lie over $\Delta_{z_1}$.
The row of $JQ_k^{T}(-\tilde{B})$ at $x$ has the single column $[c]$, the unit site whose unit block contains $c$.
The columns of the row of $B_{k,y}(\tilde{B})$ at $z_1$ lie in the $L$-block of $z_1$.
By (0.2b) all these columns lie in, or over, the $M$-cube of the row index, which is a cube of $\tilde{\mathbf{Y}}$; so (col) holds.
As for the connection, the row of $Q_k(\tilde{B})J^{*}$ at $z_1$ reads $\tilde{B}$ on bonds contained in $\Delta_{z_1}$, the row of $JQ_k^{T}(-\tilde{B})$ at $x$ reads it on bonds contained in $\Delta_{[c]}$ (no.~4), and the row of $B_{k,y}(\tilde{B})$ at $z_1$ reads it on the fine bonds of the paths $\boldsymbol{\Gamma}^{\pi}(y',z)$, $z\in B(y')$, where $y'$ is the centre of the $L$-block of $z_1$ (nos.~5 and~6).
The sites of these paths lie in the boxes spanned by $y'$ and $z$, hence in $\mathrm{Bl}_{k+1}(y')$ by (0.2c), and $\mathrm{Bl}_{k+1}(y')$ lies in the $M$-cube of $z_1$ by (0.2b).
Therefore all the bonds that are read are contained in $\tilde{\mathbf{Y}}$, and in particular they meet $\tilde{\mathbf{Y}}$ (this is also contained in Proposition~25.2\,(o8)); with (4.4) this gives (rd).
Lemma~4.1 and (4.2) now yield (a-4): at indices over $\tilde{\mathbf{Y}}$ the kernels on $T_1$ and on $T_2$ are the kernels displayed above.

\emph{Range.}
Consider first the row of $Q_k(\tilde{B})J^{*}$ at $z_1$ and the row of $JQ_k^{T}(-\tilde{B})$ at $x$, where in the second case $z_1:=[c]$.
For either row the support of the row index, the supports of the columns and the ends of the bonds that are read lie in the single unit block $\Delta_{z_1}$, and $\mathrm{sdiam}(\Delta_{z_1})=1-\eta_k$ by (0.2a).
So $1-\eta_k<1$ is a range of these rows.
(Between the row index and a column one has $|c-z_1|\le\frac{1-\eta_k}{2}$ by (0.2a); the number $\mathrm{rg}$ of no.~4 of the table of Proposition~25.1 is $\frac{1-\eta_k}{2}$.)
Consider next the row of $B_{k,y}(\tilde{B})$ at $z_1$, where $z_1\in B(y')$.
Its columns $z_2$ lie in $B(y')$, and every coordinate of a site $u$ of a path $\boldsymbol{\Gamma}^{\pi}(y',z)$, $z\in B(y')$, lies between the corresponding coordinates of $y'$ and $z$ (0.2c).
Hence, for every $\mu$, the coordinates $\mu$ of $z_1$, $z_2$ and $u$ lie in the interval $\left[y'_\mu-\frac{L-1}{2},\,y'_\mu+\frac{L-1}{2}\right]$, and it follows that $|z_2-z_1|\le L-1$ and $|u-z_1|\le L-1$.
Thus the number $L-1<L$ is a range of the row; it is the number $\mathrm{rg}$ of no.~6 of the table of Proposition~25.1.

\emph{(a-5)} ($k=0$; here $\eta_0=1$, $e_k=q$ and $\mu=\bar{m}\cdot\eta$).
No.~8 of the table of Proposition~25.1, with the kernels (K4) and (K2), gives for $x=(c,g)$ and $x'=(c',g')$ in $\mathbb{L}^{\mathrm{f}}_0$
\begin{multline*}
F_0(\tilde{B})(x,x')=48\cdot\left\{\mathbf{1}[x'=x]\cdot\left(\bar{m}\cdot\eta+\frac{\rho}{2}\cdot\mathrm{sw}\right)\cdot\mathbf{1}_4-\sum_{j=0}^{3}\mathbf{1}\left[x'=x^{(j)}\right]\cdot M_j(x)\cdot U_{\tilde{B}}\left(x,x^{(j)}\right)\right\}\\
+b_L\cdot P(\tilde{B})(c,c'),
\end{multline*}
where $U_{\tilde{B}}\left(x,x^{(j)}\right)=1$ for $j\le2$, $U_{\tilde{B}}\left(x,x^{(3)}\right)=\exp\left(iq\cdot\tilde{B}(c,c_3)\right)$, and
\[
P(\tilde{B})(c,c')=
\begin{cases}
L^{-3}\cdot U_{\tilde{B}}(y',c)^{-1}\cdot U_{\tilde{B}}(y',c')\cdot\mathbf{1}_4 & \text{if $c$ and $c'$ lie in one block $B(y')$,}\\
0 & \text{otherwise.}
\end{cases}
\]
This kernel vanishes unless $x'$ is one of $x$, $x^{(0)}$, $x^{(1)}$, $x^{(2)}$, $x^{(3)}$, or $c$ and $c'$ lie in one $L$-block.
The row at $x$ reads $\tilde{B}$ on the bond $(c,c_3)$ and on the bonds of the paths $\boldsymbol{\Gamma}^{\pi}(y',z)$, $z\in B(y')$, where $y'$ is the centre of the $L$-block of $c$.
Let $x$ lie over $\tilde{\mathbf{Y}}$.
Each column either lies over the $M$-cube of $c$ (0.2b) or is $x^{(3)}$, and $c_3$ lies in $\tilde{\mathbf{Y}}^{\sim1}$ (0.3c); so the columns are contained in $N(\tilde{Y})$, and (col) holds.
The bond $(c,c_3)$ meets $\tilde{\mathbf{Y}}$.
The sites of the paths $\boldsymbol{\Gamma}^{\pi}(y',z)$ lie in the boxes spanned by $y'$ and $z$, hence in $B(y')=\mathrm{Bl}_1(y')$ by (0.2c), and this block lies in the $M$-cube of $c$ by (0.2b); so the bonds of these paths are contained in $\tilde{\mathbf{Y}}$.
Now (4.4) gives (rd), and Lemma~4.1 together with (4.2) gives (a-5).
(All sites that occur lie in $B(y')$, at sup-distance at most $L-1$ from $c$, or coincide with $c_3$, and $|c_3-c|=1\le L-1$; so the number $L-1$ is a range of the row, namely the number $\mathrm{rg}$ of no.~8 of the table of Proposition~25.1.)

\paragraph{(b)}
Proposition~25.2\,(o1), applied with $\mathrm{Mx}=1$ and with $\tilde{\square}$ in the place of $\tilde{X}$, gives $\langle\bar{\psi},\psi\rangle_{\mathrm{pr}_i\tilde{\square}}=\mathrm{cov}_{\ell_i}(\mathrm{Fcv})$, where
\[
\mathrm{Fcv}:=\int_{\tilde{\square}}\mathrm{d}x\,\bar{\psi}(x)\cdot\psi(x)=\frac{\eta_k^3}{48}\cdot\sum_{c\in\tilde{\square}}\sum_{g\in O_{\mathrm{cub}}}\sum_{\beta=1}^{4}\bar{\psi}_{\beta}(c,g)\cdot\psi_{\beta}(c,g)
\]
is an element of $\mathrm{Grc}_k(\tilde{\square})$.
(On $T_i$ the summand of $\int_{\mathrm{pr}_i\tilde{\square}}\mathrm{d}x\,\bar{\psi}(x)\cdot\psi(x)$ is $\ell_i\mathbb{Z}^3$-periodic, being the product of two antiperiodic factors; therefore the sum may be formed over the lift $\tilde{\square}\times O_{\mathrm{cub}}$.)
By (4.3),
\[
\mathrm{idn}\,\langle\bar{\psi},\psi\rangle_{\mathrm{pr}_1\tilde{\square}}=\mathrm{cov}_{\ell_2}(\mathrm{Fcv})=\langle\bar{\psi},\psi\rangle_{\mathrm{pr}_2\tilde{\square}}.
\]

\emph{The case $k=0$.}
Here $G^{(1)}[\mathrm{Zf}](\square)=\mathrm{Vx}(\mathrm{Zf},\square)$ (Proposition~23.2, display (23.3)).
Put $\widetilde{\mathrm{Zf}}_i(b):=\mathrm{Zf}_i(\mathrm{pr}_i b)$.
Proposition~25.2\,(o6), with $\mathrm{rg}=1$ and the datum $\mathrm{Zf}$ on the bonds that meet $\tilde{\square}$, gives
\[
\mathrm{Vx}(\mathrm{Zf}_i,\mathrm{pr}_i\tilde{\square})=\mathrm{cov}_{\ell_i}\left(\mathrm{Vx}(\widetilde{\mathrm{Zf}}_i,\tilde{\square})\right),
\]
where
\begin{multline*}
\mathrm{Vx}(\widetilde{\mathrm{Zf}},\tilde{\square})=\frac{1}{2}\cdot\frac{1}{48}\cdot\sum_{c\in\tilde{\square}}\sum_{g\in O_{\mathrm{cub}}}\Bigl\{\bar{\psi}(x)\cdot M_3(x)\cdot\left(\exp\left(iq\cdot\widetilde{\mathrm{Zf}}(c,c_3)\right)-1\right)\cdot\left(\psi(x)+\chi\left(x,x^{(3)}\right)\right)\\
+\left(\bar{\psi}(x)+\bar{\chi}\left(x,x^{(3)}\right)\right)\cdot M_3\left(x^{(3)}\right)\cdot\left(\exp\left(-iq\cdot\widetilde{\mathrm{Zf}}(c,c_3)\right)-1\right)\cdot\psi(x)\Bigr\},
\end{multline*}
with $x=(c,g)$, $c_3=c+g\,e_1$, $M_3(x)=3\cdot\left(\rho-\gamma\cdot(g\,e_1)\right)$ and $M_3\left(x^{(3)}\right)=3\cdot\left(\rho+\gamma\cdot(g\,e_1)\right)$.
The only generators in this element are $\psi$, $\bar{\psi}$ at the sites over $\tilde{\square}$ and $\chi$, $\bar{\chi}$ at the pairs $\left(x,x^{(3)}\right)$ with $c\in\tilde{\square}$; this is read off the display, and on the torus it is Proposition~23.2\,(c).
For these pairs $c_3$ lies in $\tilde{\square}^{\sim1}$ (0.3c); hence $\mathrm{Vx}(\widetilde{\mathrm{Zf}},\tilde{\square})$ lies in $\mathrm{Grc}'_0(N(\tilde{Y}))$.
It reads $\widetilde{\mathrm{Zf}}$ on the bonds $(c,c_3)$, $c\in\tilde{\square}$, and these bonds meet $\tilde{\square}$ (Proposition~23.2\,(c)).
(On $T_i$ every summand is $\ell_i\mathbb{Z}^3$-periodic: $M_3$ depends on the label only, $\mathrm{Zf}_i$ is periodic, and every monomial contains two antiperiodic generators.)
By hypothesis $\widetilde{\mathrm{Zf}}_1=\widetilde{\mathrm{Zf}}_2$ on the bonds that meet $\tilde{\square}$, and $q$ is the same number on both tori.
Hence $\mathrm{Vx}(\widetilde{\mathrm{Zf}}_1,\tilde{\square})=\mathrm{Vx}(\widetilde{\mathrm{Zf}}_2,\tilde{\square})$, and (4.3) yields
\[
\mathrm{idn}\,G^{(1)}[\mathrm{Zf}_1](\mathrm{pr}_1\tilde{\square})=G^{(1)}[\mathrm{Zf}_2](\mathrm{pr}_2\tilde{\square}).
\]

\paragraph{(c)}
The objects $\left(F\right)_{00}$ and $\left(F\right)_{20}$ are those defined in Proposition~27.4\,(c): for an element $F$ of the algebra, $\left(F\right)_{00}$ is the coefficient of the empty monomial of $F$, and $\left(F\right)_{20}$ is the antisymmetric kernel, with respect to the weight $\frac{\eta_k^3}{48}$, of the part of $F$ that has degree $2$ in $\psi$ and degree $0$ in each of the other families.
For $F\in\mathrm{Gr}[\psi,\chi]$ there is no further family: the coefficient of the empty monomial is the part $F^{[0,0]}$, the part of degree $2$ in $\psi$ and of degree $0$ in $\chi$ is $F^{[2,0]}$, and the antisymmetric kernel of $F^{[2,0]}$ for the weight $\frac{\eta_k^3}{48}$ is $F_{20}$.
So in that case $\left(F\right)_{00}$ and $\left(F\right)_{20}$ are the objects $F_{00}$ and $F_{20}$ of Section~I.5.1.
By 4.2 and $\mathrm{idn}\,F_1=F_2$ there is a single element $\mathrm{Fcv}$ of $\mathrm{Grc}'_k(N(\tilde{X}))$ such that $F_i=\mathrm{cov}_{\ell_i}(\mathrm{Fcv})$ for both $i=1$ and $i=2$, namely $\mathrm{Fcv}:=\mathrm{cov}_{\ell_1}^{-1}(F_1)$; indeed, by (4.3),
\[
\mathrm{cov}_{\ell_2}(\mathrm{Fcv})=\mathrm{idn}[\tilde{X}]\,F_1=\mathrm{idn}\,F_1=F_2.
\]
According to 4.2, $\mathrm{cov}_{\ell_i}$ sends every monomial in the generators of $\mathrm{Grc}'_k(N(\tilde{X}))$ to the monomial in the generators of $T_i$ with the same indices.
It therefore sends the empty monomial to the empty monomial and preserves the degree in each of the families $\psi$, $\chi$, $\Psi$, $W$ (for $k=0$: $\psi$, $\chi$, $W$); moreover, by the hypothesis, $\mathrm{Fcv}$ contains $\psi(\xi)$ only for $\xi$ over $\tilde{X}$.
Let $\mathrm{Kcv}_{00}$ be the coefficient of the empty monomial of $\mathrm{Fcv}$.
Let $\mathrm{Fcv}^{[2]}$ be the part of $\mathrm{Fcv}$ of degree $(2,0,0,0)$ in $(\psi,\chi,\Psi,W)$ (for $k=0$: of degree $(2,0,0)$ in $(\psi,\chi,W)$), and let $\mathrm{Kcv}_{20}$ be its antisymmetric kernel with respect to the weight $\nu=\frac{\eta_k^3}{48}$ (Proposition~15\,(a)):
\[
\mathrm{Fcv}^{[2]}=\frac{1}{2}\cdot\sum_{\xi_1,\xi_2\text{ over }\tilde{X}}\left(\frac{\eta_k^3}{48}\right)^2\cdot\mathrm{Kcv}_{20}(\xi_1,\xi_2)\cdot\psi(\xi_1)\,\psi(\xi_2).
\]
Apply $\mathrm{cov}_{\ell_i}$.
The coefficient of the empty monomial of $F_i$ is $\mathrm{Kcv}_{00}$, and the part of $F_i$ of degree $2$ in $\psi$ and of degree $0$ in each of the other families is $\mathrm{cov}_{\ell_i}\left(\mathrm{Fcv}^{[2]}\right)$; it is given by the same expansion in the generators of $T_i$ with the same indices, and these generators belong to the system (2.0) formed with a set $\mathrm{Rep}_i$ that contains $N(\tilde{Y})$.
The antisymmetric kernels are determined by the element (Proposition~15\,(a)).
Consequently $(F_i)_{00}=\mathrm{Kcv}_{00}$; $(F_i)_{20}(\xi_1,\xi_2)=\mathrm{Kcv}_{20}(\xi_1,\xi_2)$ at arguments over the lift $\tilde{X}$; and $(F_i)_{20}$ vanishes at the arguments of the system (2.0) that do not lie over $\tilde{X}$.
In the matrix notation of Section~I.5.1 one has $(F_i)_{20}(x,x')_{\beta\beta'}=(F_i)_{20}\left((x,\beta,1),(x',\beta',0)\right)$, and, for $i=1$ and for $i=2$,
\begin{multline*}
\int_{\tilde{X}}\mathrm{d}x\int_{\tilde{X}}\mathrm{d}x'\,(F_i)_{20}(x,x')\\
=\left(\frac{\eta_k^3}{48}\right)^2\cdot\sum_{c,c'\in\tilde{X}}\sum_{g,g'\in O_{\mathrm{cub}}}\left[\mathrm{Kcv}_{20}\left(((c,g),\beta,1),((c',g'),\beta',0)\right)\right]_{\beta,\beta'}.
\end{multline*}
(For $F_i\in\mathrm{Gr}[\psi,\chi]$ this is the instance (o3) of Proposition~25.3\,(b5).)

\paragraph{(d)}
Put $\tilde{B}_i(b):=\mathscr{T}_i(\mathrm{pr}_i b)$; by hypothesis $\tilde{B}_1=\tilde{B}_2$ on the bonds contained in $\tilde{X}$.
By Proposition~27.2\,(d), $\mathrm{pr}_i(\tilde{X})$ is a small polymer with the lift $\tilde{X}$.
Proposition~25.2\,(o4), with $\mathrm{rg}=0$ and the datum $\tilde{B}$ on the bonds contained in $\tilde{X}$, gives on the subalgebra of $\mathrm{Grc}_k(\tilde{X})$ generated by the $\psi(\xi)$
\[
\boldsymbol{\Theta}_{\mathrm{pr}_i\tilde{X},\mathscr{T}_i}\circ\mathrm{cov}_{\ell_i}=\mathrm{cov}_{\ell_i}\circ\boldsymbol{\Theta}_{\tilde{X},\tilde{B}_i}.
\]
Here $\boldsymbol{\Theta}_{\tilde{X},\tilde{B}}$ is the homomorphism defined on the lift by (I.5.3): for $x=(c,g)$ over $\tilde{X}$,
\begin{multline*}
\boldsymbol{\Theta}_{\tilde{X},\tilde{B}}\,\psi(x)=\operatorname{Vol}(\tilde{X})^{-1}\cdot\frac{\eta_k^3}{48}\cdot\sum_{\tilde{c}\in\tilde{X}}\sum_{\tilde{g}\in O_{\mathrm{cub}}}\\
\operatorname{Mn}\left\{U_{n+2}\cdot\psi(\tilde{x})-\sum_{j=0}^{n+1}U_j\cdot d_X(p_j,p_{j+1})^{\alpha}\cdot\chi(p_j,p_{j+1})\right\},\qquad\tilde{x}=(\tilde{c},\tilde{g}),
\end{multline*}
with the following notation.
The symbol $\operatorname{Mn}$ denotes the mean over the shortest chains $(\square_0,\dots,\square_n)$ of $M$-cubes of $\tilde{X}$ from the cube of $c$ to the cube of $\tilde{c}$ in which any two consecutive cubes have a common face.
The points are $p_0=x$, $p_{j+1}=(z_j,g)$ with $z_j$ the centre of $\square_j$ ($0\le j\le n$), and $p_{n+2}=\tilde{x}$; further $U_j=\prod_{j'<j}U_{\tilde{B}}(p_{j'},p_{j'+1})$, and the terms with $p_j=p_{j+1}$ are omitted.
The first $n+1$ pairs have the same label $g$; the last pair $\left((z_n,g),(\tilde{c},\tilde{g})\right)$ lies inside a cluster if $\tilde{c}=z_n$ and $\tilde{g}\ne g$, and then $d_X=\eta_k$ and $U_{\tilde{B}}=1$.
For $\bar{\psi}$ the formula is the same, with $U_{\tilde{B}}^{-1}$, $\bar{\psi}$ and $\bar{\chi}$ in the places of $U_{\tilde{B}}$, $\psi$ and $\chi$.
All the pairs $(p_j,p_{j+1})$ are of kind (I) for $\tilde{X}$ (Proposition~17.1\,(c1)): their paths lie in $\tilde{X}$, so the transports read $\tilde{B}$ on bonds contained in $\tilde{X}$.
It follows that $\boldsymbol{\Theta}_{\tilde{X},\tilde{B}_1}=\boldsymbol{\Theta}_{\tilde{X},\tilde{B}_2}=:\boldsymbol{\Theta}_{\tilde{X},\tilde{B}}$; this is the consequence stated in Proposition~17.2\,(f).
The values of $\boldsymbol{\Theta}_{\tilde{X},\tilde{B}}$ lie in $\mathrm{Grc}_k(\tilde{X})$.
Now let $F=\mathrm{cov}_{\ell_1}(\mathrm{Fcv})$ lie in the subalgebra of $\mathrm{Gr}_1(\tilde{Y})$ generated by the $\psi(\xi)$, $\xi$ over $\tilde{X}$.
By (4.3),
\begin{align*}
\mathrm{idn}\,\boldsymbol{\Theta}_{\mathrm{pr}_1\tilde{X},\mathscr{T}_1}F
&=\mathrm{idn}\,\mathrm{cov}_{\ell_1}\left(\boldsymbol{\Theta}_{\tilde{X},\tilde{B}}\,\mathrm{Fcv}\right)
=\mathrm{cov}_{\ell_2}\left(\boldsymbol{\Theta}_{\tilde{X},\tilde{B}}\,\mathrm{Fcv}\right)\\
&=\boldsymbol{\Theta}_{\mathrm{pr}_2\tilde{X},\mathscr{T}_2}\,\mathrm{cov}_{\ell_2}(\mathrm{Fcv})
=\boldsymbol{\Theta}_{\mathrm{pr}_2\tilde{X},\mathscr{T}_2}\,\mathrm{idn}\,F.
\end{align*}

\paragraph{(e)}
Since $k+1\le k_m$, (TT) holds at the level $k+1$, with the sides $\ell_{k+1}(R)=\frac{\ell_1}{L}$ and $\ell_{k+1}(R_{\mathrm{top}})=\frac{\ell_2}{L}$, and Propositions~27.1 and~27.2 hold at that level.
The map $c'\mapsto Lc'$ is a bijection of $\mathbb{L}_{k+1}=\eta_{k+1}\mathbb{Z}^3$ onto $\mathbb{L}_k=\eta_k\mathbb{Z}^3$.
It carries the $M$-cube of $\mathbb{L}_{k+1}$ with the centre $z'\in M\mathbb{Z}^3$ onto
\[
\left\{c\in\mathbb{L}_k:|c-Lz'|<\frac{LM}{2}\right\}=\mathrm{Kub}_{LM}(Lz'),
\]
and by (0.2b) this set is the union of the $M$-cubes with the centres $Lz'+Mw$, $|w|\le\frac{L-1}{2}$.
So $L\tilde{U}$ is a non-empty finite union of $M$-cubes of $\mathbb{L}_k$.
Let $\mathrm{lo}'_\mu$, $\mathrm{hi}'_\mu$ and $\mathrm{sid}'_\mu=\mathrm{hi}'_\mu-\mathrm{lo}'_\mu+M$ be the numbers of subsection~0.4 formed for $\tilde{U}$ at the level $k+1$.
The numbers of subsection~0.4 formed for $L\tilde{U}$ at the level $k$ are then
\[
\mathrm{lo}_\mu=L\cdot\mathrm{lo}'_\mu-M\cdot\frac{L-1}{2},\qquad
\mathrm{hi}_\mu=L\cdot\mathrm{hi}'_\mu+M\cdot\frac{L-1}{2},
\]
\[
\mathrm{sid}_\mu=L\cdot\left(\mathrm{hi}'_\mu-\mathrm{lo}'_\mu\right)+M\cdot(L-1)+M=L\cdot\mathrm{sid}'_\mu.
\]

\emph{Admissibility.}
One has
\[
\mathrm{diam}(L\tilde{U})=L\cdot\mathrm{diam}(\tilde{U})\le L\cdot\left(\frac{\ell_1}{L}-14d_U\cdot M\right)=\ell_1-14d_U\cdot LM\le\ell_1-14d_U\cdot M.
\]

\emph{$N(L\tilde{U})$ is contained in $L\,N(\tilde{U})$.}
Since $\mathrm{lo}_\mu-\frac{M}{2}=L\cdot\mathrm{lo}'_\mu-\frac{LM}{2}$ and $\mathrm{hi}_\mu+\frac{M}{2}=L\cdot\mathrm{hi}'_\mu+\frac{LM}{2}$, (0.4b) gives
\begin{align*}
N(L\tilde{U})&=\left\{c:L\cdot\mathrm{lo}'_\mu-\frac{LM}{2}-6d_U\cdot M<c_\mu<L\cdot\mathrm{hi}'_\mu+\frac{LM}{2}+6d_U\cdot M\right\},\\
L\,N(\tilde{U})&=\left\{c:L\cdot\left(\mathrm{lo}'_\mu-\frac{M}{2}-6d_U\cdot M\right)<c_\mu<L\cdot\left(\mathrm{hi}'_\mu+\frac{M}{2}+6d_U\cdot M\right)\right\},
\end{align*}
and $6d_U\cdot M\le6d_U\cdot LM$.

\emph{The map $S_L$.}
By Section~I.3.7,
\begin{gather*}
S_L\,\psi(x,\beta,\omega)=L^{-1}\cdot\psi(x/L,\beta,\omega),\qquad S_L\,\chi(x,x',\beta,\omega)=L^{-1-\alpha}\cdot\chi(x/L,x'/L,\beta,\omega),\\
x/L=(c/L,g),
\end{gather*}
on the generators indexed by the covering lattices.
This rule is compatible with (I.AP), because a translation of $x$ by $\ell_i e_\mu$ is a translation of $x/L$ by $\frac{\ell_i}{L}e_\mu$.
Let $c$ and $c'$ lie in $N(L\tilde{U})$; then $c/L$ and $c'/L$ lie in $N(\tilde{U})$.
For $(x,x')\in\mathrm{Pzt}_k$ one has $x/L\ne x'/L$ and $d_X(x/L,x'/L)=\frac{d_X(x,x')}{L}\le\frac{M}{L}\le M$ (Section~I.3.7), so that $(x/L,x'/L)\in\mathrm{Pzt}_{k+1}$.
Therefore $S_L$ maps the generators $\psi$, $\chi$ of $\mathrm{Gr}_i(L\tilde{U})$ to multiples of generators of $\mathrm{Gr}_i(\tilde{U})$ of the level $k+1$, and the subalgebra generated by them into $\mathrm{Gr}_i(\tilde{U})$.
On such a generator of $T_1$ one has
\[
S_L\,\mathrm{idn}[L\tilde{U}]\,\psi(x,\beta,\omega)=L^{-1}\cdot\psi(x/L,\beta,\omega)=\mathrm{idn}[\tilde{U}]\,S_L\,\psi(x,\beta,\omega),
\]
where $\psi(x/L,\beta,\omega)$ in the middle is a generator of $T_{k+1}(R_{\mathrm{top}})$, and likewise
\[
S_L\,\mathrm{idn}[L\tilde{U}]\,\chi(x,x',\beta,\omega)=L^{-1-\alpha}\cdot\chi(x/L,x'/L,\beta,\omega)=\mathrm{idn}[\tilde{U}]\,S_L\,\chi(x,x',\beta,\omega).
\]
Both sides are homomorphisms of algebras.
Hence $S_L\,\mathrm{idn}[L\tilde{U}]=\mathrm{idn}[\tilde{U}]\,S_L$ on the part of $\mathrm{Gr}_1(L\tilde{U})$ generated by $\psi$ and $\chi$.
In the notation of Proposition~25.2\,(o7) and Proposition~18.5\,(f) this is the identity $\mathrm{cov}_{\ell_i/L}\circ S_L=S_L\circ\mathrm{cov}_{\ell_i}$ on the subalgebra of $\mathrm{Grc}_k(N(L\tilde{U}))$ generated by $\psi$ and $\chi$.
The margin required there is available: the sides of $N(L\tilde{U})$ are $L\cdot\mathrm{sid}'_\mu+12d_U\cdot M\le\ell_1-14d_U\cdot LM+12d_U\cdot M$, so that
\[
\mathrm{sdiam}(N(L\tilde{U}))+LM<\ell_1-\left(14d_U\cdot L-12d_U-L\right)\cdot M<\ell_i.
\]

\paragraph{(f)}
\emph{The integral over $W$.}
Let $F=\mathrm{cov}_{\ell_1}(\mathrm{Fcv})$ with $\mathrm{Fcv}\in\mathrm{Grc}'_k(N(\tilde{Y}))$.
By Proposition~25.2\,(o9),
\[
\int\mathrm{d}\mu_I(W)\,\mathrm{cov}_{\ell_i}(\mathrm{Fcv})=\mathrm{cov}_{\ell_i}\left(\int_{\mathrm{c}}\mathrm{Fcv}\right),
\]
where $\int_{\mathrm{c}}$ is the integral (I.3.5) with $\mathrm{d}\mu_I$ on the index set of the $W(u)$ of $\mathrm{Grc}_k(N(\tilde{Y}))$, with $\nu=1$ for $k\ge1$ and $\nu=\frac{1}{48}$ for $k=0$.
By (4.3),
\begin{align*}
\mathrm{idn}\int\mathrm{d}\mu_I(W)\,F
&=\mathrm{idn}\,\mathrm{cov}_{\ell_1}\left(\int_{\mathrm{c}}\mathrm{Fcv}\right)
=\mathrm{cov}_{\ell_2}\left(\int_{\mathrm{c}}\mathrm{Fcv}\right)\\
&=\int\mathrm{d}\mu_I(W)\,\mathrm{cov}_{\ell_2}(\mathrm{Fcv})
=\int\mathrm{d}\mu_I(W)\,\mathrm{idn}\,F.
\end{align*}
(The same argument in terms of monomials.
$\mathrm{Fcv}$ is a linear combination of products $\mathrm{Fmo}\cdot W(u_1)\cdots W(u_p)$, where $\mathrm{Fmo}$ is a monomial in $\psi$, $\chi$, $\Psi$, and $u_1,\dots,u_p$ are pairwise different indices of $W$ in $N(\tilde{Y})$.
On $T_i$ the integral of the image of such a product is the image of $\mathrm{Fmo}$ multiplied by the number $\int\mathrm{d}\mu_I(W)\,W(u_1)\cdots W(u_p)$.
By Proposition~15\,(e), which is (I.B4), this number vanishes unless $\{u_1,\dots,u_p\}$ is paired.
If the set is paired, the number equals $\mathrm{sn}\cdot\prod\nu^{-1}$, the product extending over the $\frac{p}{2}$ modes that occur, where $\mathrm{sn}\in\{1,-1\}$ is the sign of the permutation that brings $W(u_1)\cdots W(u_p)$ to the form $\prod W(\mathrm{mo})\,\bar{W}(\mathrm{mo})$.
The number is determined by the sequence $(u_1,\dots,u_p)$ and by $\nu$; the modes of $T_i$ that do not occur contribute the factor $\int\mathrm{d}\mu_I=1$.
For $k\ge1$ the modes are $\mathrm{mo}=(y,\beta)$, $\nu=1$, and the number is $\mathrm{sn}$.
For $k=0$ the modes are $\mathrm{mo}=(x,\beta)$ with $x=(c,g)$, $\nu=\frac{1}{48}$, and the number is $\mathrm{sn}\cdot48^{\frac{p}{2}}$; in particular $\int\mathrm{d}\mu_I\,W_{\beta}(x)\,\bar{W}_{\beta'}(x')=48\cdot\mathbf{1}[x=x',\beta=\beta']$ on both tori, by Proposition~15\,(n) and (I.L0.5).)

\emph{The measure $\hat{\mu}$.}
Consider the numbers $v_k=-\log e_k$ and $v_K=-\log\left(e\cdot h^{\frac{1}{2}}\right)$ of Section~I.6.1.
One has $v_k\ge v_K$ by (G4.Sc3) of Section~G4.0.2, because $k\le k_m=k_{*}(R)$, and $v_K\ge8$ by (G4.Sc6).
Hence $p_{0,k}=v_k^{p_0}>0$.
The measure $\mu_I$ is the standard Gaussian measure of the Euclidean space $V_k(T_i)$, equipped with the counting scalar product \cite[Section~5.2.5]{QEDI}, and $\{\Upsilon\}$ is an orthonormal basis of $V_k(T_i)$ with $d_V\cdot\ell_i^3$ elements, where $d_V=2-2L^{-3}$ \cite[Proposition~15.21(ii)]{QEDI}.
Therefore, in the coordinates $\mathsf{Z}(\Upsilon)=\langle\mathsf{Z},\Upsilon\rangle$, the measure $\mu_I$ is the product, over the basis, of copies of the standard normal law $(2\pi)^{-\frac{1}{2}}\cdot\exp\left(-\frac{z^2}{2}\right)\mathrm{d}z$, and
\[
\mathrm{Box}_k=\left\{|\mathsf{Z}(\Upsilon)|\le p_{0,k}\text{ for all }\Upsilon\right\}
\]
is the product of the intervals $[-p_{0,k},p_{0,k}]$.
It follows that
\[
\mu_I(\mathrm{Box}_k)=\left((2\pi)^{-\frac{1}{2}}\int_{-p_{0,k}}^{p_{0,k}}\mathrm{d}z\,\exp\left(-\frac{z^2}{2}\right)\right)^{d_V\cdot\ell_i^3},
\]
which is positive, since $p_{0,k}>0$ and the integrand is positive.
So $\hat{\mu}$ and $\delta\varepsilon^{0}_{k}$ are defined \cite[Section~5.2.5]{QEDI}.
The measure
\[
\mathrm{d}\hat{\mu}=\frac{\mathbf{1}_{\mathrm{Box}_k}}{\mu_I(\mathrm{Box}_k)}\,\mathrm{d}\mu_I
\]
is the product, over the basis, of copies of the standard normal law restricted to $[-p_{0,k},p_{0,k}]$ and normalised, and
\begin{align*}
\delta\varepsilon^{0}_{k}&=-n_k^{-1}\cdot\log\mu_I(\mathrm{Box}_k)\\
&=-d_V\cdot\log\left((2\pi)^{-\frac{1}{2}}\int_{-p_{0,k}}^{p_{0,k}}\mathrm{d}z\,\exp\left(-\frac{z^2}{2}\right)\right)\qquad\left(n_k=\ell_i^3\right).
\end{align*}
The number $p_{0,k}=v_k^{p_0}$, with $v_k=-\log e_k$, is the same on both tori by (TT), and $d_V$ depends on $L$ only.
(This is Proposition~25.2\,(o10), whose proviso $p_{0,k}>0$ is satisfied.)
\qed

\emph{Level $0$.}
Clauses (a-1) and (a-2) hold at $k=0$ as stated: $\eta_0=1$, every pair has $d_X\ge1$, and inside a cluster $(\delta_{\alpha,B}f)(x,x')=f(x')-f(x)$.
Clauses (a-3) and (a-4) do not occur.
Clause (a-5) and the second half of (b) occur at $k=0$ only; the factor $48$ in $F_0(\tilde{B})(x,x')$ is the inverse of the weight in (I.ke1).
In (c) and (d) the weight is $\frac{1}{48}$.
In (e), $S_L$ maps the level $0$ to the level $1$ and leaves the label unchanged.
In (f), $\nu=\frac{1}{48}$, and the moments are $\mathrm{sn}\cdot48^{\frac{p}{2}}$.

In Proposition~27.4 no constant is introduced and no bound is proved; its clauses hold for all data (TT) with $k<k_m$, for all admissible $\tilde{Y}$ and for all $y\ge0$.

\subsection*{Conditions on the constants}

Propositions~27.1--27.4 are identities and inclusions: no constant is introduced, and no bound is proved.
The arguments of this chapter rest on the following conditions on the constants of Section~I.8.
\begin{description}
\item[$L$ odd; $M$ a power of $L$ with $M\ge L^2$]
These conditions stand in Sections~I.0.5 and~I.8.2, (C3), (C5), and in Proposition~14.2\,(b), (L5), (M1).
That $L$ is odd and that $M$ is a power of $L$ make $\frac{a}{\eta_k}$ an odd integer for $a=L^j\cdot\eta_k$, $M$, $LM$, and $L$ an odd integer.
The quotient $\frac{M}{L}$ is a power of $L$ as well, and the inequality $M\ge L^2$ of (M1) gives $\frac{M}{L}\ge L$; hence $\frac{M}{L}$ is an odd integer.
For a power $M$ of $L$ the inequality $M>1$ already suffices for this.
The facts (0.2a) and (0.2b) of subsection~0.2 rest on this, and through subsection~0.2 it underlies every part of this chapter.
\item[$\eta_k\le1\le M$]
One has $M\ge L^2$ by Section~I.0.5 and by Proposition~14.2\,(b), (M1).
Consequently a bond that meets $Z$ lies in $Z^{\sim1}$; this is (0.3c).
\item[$L\ge5$]
This condition stands in Sections~I.0.5 and~I.8.2, (C3), and in Proposition~14.2\,(b), (L5).
It yields $\eta_k\le\frac{1}{L}<1$ for $k\ge1$, and $1\le L-1$; the ranges found in Proposition~27.4\,(a-3) and (a-5) rest on these two inequalities.
\item[Condition (2.1) of Section~G4.0.3, $r_*>31L+16$]
Together with (TT) and (G4.Sc3) it gives $\ell_1>(31L+16)\cdot M$.
This inequality is applied in (0.5b) and in Proposition~27.1\,(f) in the forms $L-1<\ell_i$, $M<\ell_i$ and $2M<\ell_i$, which already follow from $\ell_i\in\mathrm{Cl}$, $\ell_i\ge13LM$;
in Proposition~27.1\,(f) also in the form $(4L+2)\cdot M<\frac{\ell_i}{2}$;
in Proposition~27.2\,(d) in the forms $(L+2)\cdot M\le\ell_1-14d_U\cdot M$ and $(2L+4)\cdot M<\ell_i$;
and in Step~1 of the proof of Proposition~27.3 in the form $2M\le\ell_1-14d_U\cdot M$.
The condition is stated in Section~I.8.4 and, as (2.1), in Proposition~14.2\,(c), where it is derived from $e\cdot h^{\frac{1}{2}}\le e_1$ through the row (e1) of Proposition~14.2\,(b) and the line (C9); see also Annex~B, Section~G4.A.6.
\item[$v_K\ge8$]
By (G4.Sc3) and (G4.Sc6) one has $v_k\ge v_K\ge8$ for $k\le k_m$.
Of this only $v_k>0$ is needed: it gives $p_{0,k}=v_k^{p_0}>0$, hence $\mu_I(\mathrm{Box}_k)>0$ in Proposition~27.4\,(f), so that $\hat{\mu}$ and $\delta\varepsilon^{0}_{k}$ are defined, and it is the proviso of Proposition~25.2\,(o10).
The inequality $v_K\ge8$ is Proposition~14.2\,(f), where it is derived from $e\cdot h^{\frac{1}{2}}\le e_1$ through the first inequality of the row (e1) of Proposition~14.2\,(b) and the line (C9); see also Annex~B, Section~G4.A.6.
\item[$d_U=2L+1$ and $c_l=2L$]
These are definitions (Section~I.1.1; Section~I.8.2, (C4)).
They give
$2LM<4LM<2d_U\cdot M$, used in Proposition~27.1\,(d) and (f);
$c_l+1=d_U$, and $j+d_U\le6d_U$ for $j\le5d_U$, used in Step~1 of the proof of Proposition~27.3;
$1\le6d_U$, $L+1\le6d_U$ and $1+2d_U\le6d_U$, used in Proposition~27.2\,(b), (e) (step (e4)) and in Proposition~27.4\,(a-1);
and $12d_U+L<14d_U\cdot L$, used in Proposition~27.4\,(e).
Together with $M\ge1$ they give moreover
$M<2d_U\cdot M$, used for the pairs in Proposition~27.1\,(d) and for the first of the three margins stated after (4.0);
$2+M<14d_U\cdot M$ and $3M<14d_U\cdot M$, used for the second and the third of these margins;
and $L-1<2d_U\cdot M$, used in the proof of Proposition~27.4\,(a-1), in the remark that (a-1) is also Proposition~25.3\,(a) with $\tilde{U}=N(\tilde{Y})$.
\item[$M>4L$]
This inequality, stated in Section~I.8.2, (C5), and in Proposition~14.2\,(b), (M2), is one of the conditions stated with Proposition~26.4, and it enters Proposition~27.3 through Proposition~26.4.
\end{description}
Every inequality between constants that occurs in the arguments of this chapter follows from the conditions listed above.
The arguments of this chapter apply, moreover, the following clauses, definitions and assertions:
Proposition~15\,(a), (e) and~(n);
Proposition~17.1\,(c1) and Proposition~17.2\,(f);
Proposition~18.5\,(f);
Proposition~23.1\,(b), and Proposition~23.2\,(c) together with the display (23.3);
the rows nos.~4--6 and~8--12 of the table of Proposition~25.1, with the kernels (K1)--(K4);
Proposition~25.2\,(o1), (o3), (o4) and~(o6)--(o10);
Proposition~25.3\,(a), (b1)--(b3) and~(b5);
Propositions~26.1, 26.3\,(a)--(c), 26.4 and~26.5, and Proposition~26.6\,(d);
the Definition in Proposition~48.1, and Proposition~48.1\,(b);
the assertions (G4.Sc2), (G4.Sc3) and~(G4.Sc6) of Section~G4.0.2, which enter with the standing data (TT), and the assertions (G4.Sc1) and~(G4.Sc4) of that section, which enter through Proposition~26.6.
The conditions on the constants under which these hold enter Propositions~27.1--27.4 through them.
They are stated at the following places:
for Propositions~15, 17.1, 17.2, 18.5 and~48.1, in Annex~B, Section~G2.9;
for Propositions~23.1, 23.2 and~25.1--25.3, in Annex~B, Section~G3.6.1, in Chapter~23, Section~4, and in Chapter~25, Section~5;
for Propositions~26.3, 26.4 and~26.5, with these Propositions;
for Proposition~26.6 and for the assertions of Section~G4.0.2 named above, in Chapter~26b, Section~3;
for Propositions~26.1 and~26.3--26.6 and for these assertions, moreover, in Annex~B, Section~G4.A.6.
The complete list of the conditions on the constants is that of Proposition~14.2.
One of the conditions that enter in this way, $M>4L$, is stated with Proposition~26.4 and is listed above.
Those that are not listed above include the following, arranged by the places at which they are stated.
\begin{itemize}
\item Stated with Proposition~26.3 and in the first entry of Annex~B, Section~G4.A.6; they enter through Proposition~26.3\,(a)--(c):
$\gamma_{\mathrm{w}}\cdot M\ge1$, which follows from (M3); $M\ge M^{P}$ (M6); $0<\alpha_{0}\le\frac{1}{4}$ and $\kappa_{\mathrm{w}}\ge0$ (C4); $K_{0}\ge1$ (C0); $C_w^{(1)}\ge4\exp(2\gamma_{\mathrm{w}})\cdot M^{27\alpha_{0}}$ (C6).
\item Stated with Proposition~26.5 and in the second entry of Annex~B, Section~G4.A.6; they enter through Proposition~26.5:
$M\ge M^{F}$ at $\kappa'=\kappa_{\mathrm{w}}+\kappa_{0}$ (M5); $M\ge M^{P}$ (M6); (Z2.s1)--(Z2.s4) (C8); $0\le\alpha<1$ (C1); $0<\alpha_{0}\le\frac{1}{4}$ (C4).
\item Stated in Chapter~26b, Section~3, and in Annex~B, Section~G4.A.6; they enter through Proposition~26.6\,(d) and through the assertions of Section~G4.0.2 named above:
the third inequality of the row (e1) of Proposition~14.2\,(b), namely $\left(1+\frac{3\,\mathrm{lg}}{v}\right)^{r}<L$ for all $v\ge-\log e_1$; $L\ge13$, which follows from (L1); $e_1\le e_0$ (C9); $r>0$ (C7).
\item Stated in Annex~B, Section~G3.6.1, in its entries on Propositions~23.1\,(a)--(f) with~23.2\,(a)--(d) and on Propositions~25.1--25.3, in Chapter~23, Section~4, and in Chapter~25, Section~5; they enter through Propositions~23.1\,(b), 23.2\,(c) and~25.1--25.3:
$b>0$, which together with $L>1$ gives $b_k>0$; $L-1<M$; the closure relation (I.1.1); $0<\rho<1$ (C0), which enters the proof of Proposition~23.2\,(c).
\end{itemize}
Here (C0)--(C9) are lines of Section~I.8.2, and (L1), (M3), (M5), (M6) and (Z2.s1)--(Z2.s4) are rows of Proposition~14.2\,(b).
Of the statements of \cite{QEDI}, the arguments of this chapter apply \cite[Lemmas~17.6 and~17.7]{QEDI}, the assertions (t7) and~(t10) of \cite[Theorem~13.2]{QEDI}, which are listed under the same marks in \cite[Proposition~5.58]{QEDI}, and \cite[Propositions~5.35 and~15.21]{QEDI}.
In the forms in which they are applied, these statements require only that $L$ be odd with $L\ge3$, which follows from the first and the third of the conditions listed above, and that the sides $\ell_1$, $\ell_2$ satisfy (0.1a), which holds by Proposition~26.6\,(d).

\section*{Chapter 27b. The objects of one step, and the step itself, agree on two tori: proofs of Propositions 27.5 and 27.6}

\subsection*{0. Conventions of this chapter; facts that both proofs use}

\subsubsection*{0.0. Standing data and notation}

Throughout this chapter the standing data (TT) of the Propositions~27.m are assumed; they are stated before Proposition~27.1.
A set $\tilde{X}$ is said to be \emph{connected in} $\tilde{\mathbf{Y}}$ if $\tilde{X}$ is a non-empty union of $M$-cubes of $\mathbb{L}_k$ that is connected through faces and contained in $\tilde{\mathbf{Y}}$; the same words are used with other sets of sites of $\mathbb{L}_k$, such as $\tilde{Y}$, in the place of $\tilde{\mathbf{Y}}$.
The data $\mathrm{Bs}$, $\mathsf{Z}^{(i)}$ and $s^{(i)}$ of Proposition~27.5 depend on $\tilde{\mathbf{Y}}$, on $\mathrm{zval}$ and on $\mathrm{sval}$.
Where this dependence matters, it is shown in the notation:
$\mathrm{Bs}(\tilde{\mathbf{Y}})$, $\mathsf{Z}^{(i)}[\tilde{\mathbf{Y}};\mathrm{zval}]$, $s^{(i)}[\tilde{\mathbf{Y}};\mathrm{sval}]$.
Bare labels (a), (b), (c1)--(c7), (d), (e) refer to the clauses of the Propositions~27.5 and 27.6.

\subsubsection*{0.1. The two tori}

\begin{itemize}
\item[(TT.1)] (Both tori carry the step.)
Let (TT) hold with $0\le k<k_m$.
The quadruple $(\bar{m},e,R_{\mathrm{top}},N)$ is (FL) by Proposition~26.6\,(a), and (TT) gives $k<k_m=k_{*}(R)\le k_{\mathrm{top}}=k_{*}(R_{\mathrm{top}})$.
Proposition~26.6\,(b), applied on $R$ and on $R_{\mathrm{top}}$, shows that for $i=1$ and for $i=2$ the data $(k,\ell_i,\mu_k,e_k)$, together with the system $\mathrm{Op}^{(i)}$, are (OS).
By the Dictionary of Section~I.1.8, (OS) contains (ST1)$^{\mathrm{ch}}$.
Consequently Definition~24.D defines the step map on $T_1$ and on $T_2$, and the Propositions~24.1 and 24.3 hold on both tori.
Moreover $\ell_i$ and $\ell_i/L$ belong to $\mathrm{Cl}$, and the level $k+1$ has the sides $\ell_1/L=\ell_{k+1}(R)$ and $\ell_2/L=\ell_{k+1}(R_{\mathrm{top}})$, the mass $\mu_{k+1}$ and the charge $e_{k+1}$.
As $k+1\le k_m$, the standing data (TT) hold at the level $k+1$ as well.
For $m$ and $\mathscr{E}^{(i)}$ as in the Propositions~27.5 and 27.6, the pair $(m,\mathscr{E}^{(i)})$ lies in $\mathrm{Ball}^{x}_{k}(1)$, which is a subset of $\mathrm{Ball}^{x}_{k}(4)$.
\item[(TT.2)] (Periodic objects.)
By Proposition~26.1 an object on $T_i$ is an $\ell_i\mathbb{Z}^3$-periodic, respectively antiperiodic, object on the covering index sets, and $\mathrm{pj}^{*}$ is the inclusion.
A bond function $B$ on $T_1$ and the bond function $\mathrm{pj}^{*}B$ are therefore one and the same function on the bonds of $\mathbb{L}_k$; in particular $B$ and $\mathrm{pj}^{*}B$ have equal lifts, in the sense of Proposition~27.2\,(D1), on every set of bonds.
The map $\mathrm{pj}^{*}$ is linear, and it commutes with $B\mapsto\overline{B}$.
\item[(TT.3)] (Domains.)
By (OS), see (TT.1),
\[
e_k=L^{-\frac{1}{2}}\cdot e_{k+1}\le L^{-\frac{1}{2}}\cdot e_0<e_0 ;
\]
here $e_0$ is the ceiling of the line (C8) of Section~I.8.2, and at $k=0$ the symbol $e_k$ denotes the charge $q$ of level~$0$ (I.L0.2).
The first clause of condition (A8) of Section~I.8.4 (\cite[Section~5.6.2]{QEDI}) is $e_0^{\varepsilon}\le\frac{1}{4}$; as $e_k<e_0$, it gives $e_k^{\varepsilon}\le\frac{1}{4}$.
Hence $e_k<1$ and $0<1-e_k^{\varepsilon}<1$.
By \cite[Lemma~5.20]{QEDI}, $\mathrm{pj}^{*}$ maps each of the sets $\tilde{R}_k$, $\tilde{R}_k^{-}$, $\mathfrak{X}_k$, $\tilde{G}_k$ of $T_1$ into the set of the same name of $T_2$.
By \cite[Lemma~5.17(a)]{QEDI}, the sets $\tilde{R}_k$ and $\tilde{R}_k+\mathfrak{X}_k$ are contained in $\tilde{G}_k$.
In \cite[Definition~5.16]{QEDI} (Section~I.5.3) the set $\tilde{R}_k^{-}$ is defined in the same way as $\tilde{R}_k$, except that the summand of the splitting which for $\tilde{R}_k$ is required to lie in $R_k$ is required to lie in $\left(1-e_k^{\varepsilon}\right)\cdot R_k$.
It follows that $\tilde{R}_k^{-}\subset\tilde{R}_k$.
Moreover $0$ belongs to $\tilde{R}_k^{-}$, to $\tilde{R}_k$ and to $\mathfrak{X}_k$, since all the defining inequalities are strict inequalities between $0$ and positive numbers.
By \cite[Lemma~5.19]{QEDI}, $\tilde{G}_k$ is invariant under $B\mapsto-B$ and under $B\mapsto\overline{B}$.
For the data of Proposition~27.5 one obtains: $(A_2,\mathscr{C}_2)$ lies in $\tilde{R}_k^{-}(T_2)\times\mathfrak{X}_k(T_2)$; $\mathscr{T}_2=\mathrm{pj}^{*}\mathscr{T}_1$; the bond functions $\mathscr{T}_i$, $\mathscr{C}_i$ and $0$ lie in $\tilde{G}_k(T_i)$; and $A_1$, $A_2$ have equal lifts on all bonds, as have $\mathscr{C}_1$, $\mathscr{C}_2$ and $\mathscr{T}_1$, $\mathscr{T}_2$.
\end{itemize}

\subsubsection*{0.2. Coefficients}

Let $\tilde{Y}$ be admissible.
Denote by $\mathrm{Jx}(\tilde{Y})$ the finite set of the indices of the generators that Proposition~27.2\,(D2) names for $\tilde{Y}$:
the indices $(x,\beta,\omega)$ with $x$ over $N(\tilde{Y})$, for $\psi$;
the indices $(x,x',\beta,\omega)$ with $(x,x')\in\mathrm{Pzt}_k$ and with photon sites $c$, $c'$ in $N(\tilde{Y})$, for $\chi$;
for $k\ge1$, the indices $(y,\beta,\omega)$, $y$ a unit site of $N(\tilde{Y})$, each taken once for $\Psi$ and once for $W$;
for $k=0$, the indices $(x,\beta,\omega)$ with $x$ over $N(\tilde{Y})$, for $W$.
Fix a total order on $\mathrm{Jx}(\tilde{Y})$.
For $I\subset\mathrm{Jx}(\tilde{Y})$ let $\mathrm{mon}_i(I)$ be the product, in this order, of the generators of $T_i$ with the indices in $I$, and let $\mathrm{mon}_i(\emptyset)=1$.
By Proposition~27.2\,(a) these generators are independent on $T_i$.
Hence the elements $\mathrm{mon}_i(I)$ form a basis of $\mathrm{Gr}_i(\tilde{Y})$, every $F\in\mathrm{Gr}_i(\tilde{Y})$ has a unique expansion
\[
F=\sum_{I}c_F(I)\cdot\mathrm{mon}_i(I)
\]
with complex numbers $c_F(I)$ (I.5.1), and $\mathrm{idn}\,\mathrm{mon}_1(I)=\mathrm{mon}_2(I)$ for $\mathrm{idn}=\mathrm{idn}[\tilde{Y}]$ by the definition in Proposition~27.2\,(D2).
The following facts are consequences.
\begin{itemize}
\item[(Co1)] For $F_1\in\mathrm{Gr}_1(\tilde{Y})$ and $F_2\in\mathrm{Gr}_2(\tilde{Y})$ one has $\mathrm{idn}\,F_1=F_2$ if and only if $c_{F_1}(I)=c_{F_2}(I)$ for all $I$.
\item[(Co2)] The map $\mathrm{idn}$ is an isomorphism of algebras with unit (Proposition~27.2\,(a)), and it fixes the numbers.
Let $F$ be even, and write $F=c_F(\emptyset)+F'$.
Then $F'$ is nilpotent, the exponential
\[
\exp F=\exp\left(c_F(\emptyset)\right)\cdot\sum_{n\ge0}\frac{F'^{\,n}}{n!}
\]
is a finite sum, and $\mathrm{idn}\,\exp F=\exp\left(\mathrm{idn}\,F\right)$.
\item[(Co3)] Let $F_1(\mathrm{pm})$ and $F_2(\mathrm{pm})$ be functions of a parameter $\mathrm{pm}$ with values in $\mathrm{Gr}_1(\tilde{Y})$ and in $\mathrm{Gr}_2(\tilde{Y})$, such that $\mathrm{idn}\,F_1(\mathrm{pm})=F_2(\mathrm{pm})$ for every $\mathrm{pm}$.
By (Co1) the coefficient functions $\mathrm{pm}\mapsto c_{F_i(\mathrm{pm})}(I)$ do not depend on $i\in\{1,2\}$.
In a finite-dimensional space, the integral over $\mathrm{pm}$ with respect to a measure, the partial derivative with respect to a component of $\mathrm{pm}$, the limit of a sequence and the sum of a series are all formed coefficient by coefficient.
Each of these operations is therefore defined either for both $F_1$ and $F_2$ or for neither of them, and where it is defined the two results correspond to each other under $\mathrm{idn}$.
\item[(Co4)] Let $\tilde{X}$ be a non-empty union of $M$-cubes contained in $\tilde{Y}$.
Then $\tilde{X}$ is admissible and $N(\tilde{X})\subset N(\tilde{Y})$ (Proposition~27.1\,(a)); moreover $\mathrm{Gr}_i(\tilde{X})\subset\mathrm{Gr}_i(\tilde{Y})$, and $\mathrm{idn}[\tilde{X}]$ is the restriction of $\mathrm{idn}[\tilde{Y}]$ (Proposition~27.2\,(a)).
Accordingly the index of $\mathrm{idn}$ may be omitted.
\end{itemize}

\subsubsection*{0.3. Sums over representatives; membership}

Let $\tilde{Y}$ be admissible.
\begin{itemize}
\item[(Mb)] Let $\tilde{S}$ be a set of sites contained in $N(\tilde{Y})$, and let $u$ be a site of $N(\tilde{Y})$.
Then $\mathrm{pr}_i(u)\in\mathrm{pr}_i(\tilde{S})$ if and only if $u\in\tilde{S}$.
Equivalently, $u$ lies in the $\ell_i\mathbb{Z}^3$-periodic extension $\tilde{S}+\ell_i\mathbb{Z}^3$ of $\mathrm{pr}_i(\tilde{S})$ if and only if $u$ lies in $\tilde{S}$.
Indeed, $\mathrm{pr}_i$ is injective on the sites of $N(\tilde{Y})$ by Proposition~27.1\,(b).
\item[(Rs)] Let $\tilde{S}$ be a set of sites contained in $N(\tilde{Y})$, and let the index $t$ range over one of the covering index sets: the unit sites, the positively oriented unit bonds, the basis elements, or the fermion sites.
Let $\mathrm{smd}(t)$ take values in a vector space, and assume that $\mathrm{smd}(t+\ell_iw)=\mathrm{smd}(t)$ for all $w\in\mathbb{Z}^3$ and that $\mathrm{smd}(t)=0$ unless $\mathrm{pr}_i(t)$ is contained in $\mathrm{pr}_i(\tilde{S})$ (for fermion sites: unless $\mathrm{pr}_i(t)$ lies over $\mathrm{pr}_i(\tilde{S})$).
Then the sum of $\mathrm{smd}(t)$ over a set of representatives of the indices modulo $\ell_i\mathbb{Z}^3$ is
\[
\sum_{t\ \mathrm{mod}\ \ell_i\mathbb{Z}^3}\mathrm{smd}(t)=\sum_{t\text{ contained in (over) }\tilde{S}}\mathrm{smd}(t).
\]
Indeed, by Proposition~27.1\,(d) every index of $T_i$ that is contained in $\mathrm{pr}_i(\tilde{S})$ has exactly one lift that is contained in $\tilde{S}$.

The rule (Rs) is applied to summands $\mathrm{smd}(t)=O(t',t)\cdot\Phi(t)$, where $t'$ is a fixed row index, $O$ is a kernel in the sense of Section~G4.0.4, and $\Phi$ is a field or a family of generators.
For photon indices both factors are $\ell_i\mathbb{Z}^3$-periodic; for spinor indices both factors change their sign under $t\mapsto t+\ell_ie_{\mu}$ (Section~G4.0.4; (I.AP) in Section~I.3.6), so that their product is periodic.
\end{itemize}

\subsubsection*{0.4. The generators that belong to a set $\tilde{X}$}

Let $\tilde{Y}$ be admissible and let $\tilde{X}$ be connected in $\tilde{Y}$.
Let $\mathrm{Pad}(\tilde{X})$ be the set of those pairs $(x,x')\in\mathrm{Pzt}_k$ that are of kind (I) for $\tilde{X}$, which means that $c$ and $c'$ lie in $\tilde{X}$ and the six paths $\boldsymbol{\Gamma}^{\pi}(c,c')$ lie in $\tilde{X}$ (for $c=c'$: $c$ lies in $\tilde{X}$), or of kind (II) for $\tilde{X}$, which means that $x'=x^{(3)}$ and that $c$ or $c_3$ lies in $\tilde{X}$.
For such a pair the photon sites of both points have sup-distance at most $\eta_k$ from $\tilde{X}$.
Since $\eta_k\le1<M$, they lie in $\tilde{X}^{\sim1}$, which is a subset of $N(\tilde{X})$.
Three subalgebras of $\mathrm{Gr}_i(\tilde{X})$ are attached to $\tilde{X}$:
\begin{itemize}
\item $\mathrm{Grq}_i(\tilde{X})$ is the subalgebra generated by the $\psi(x,\beta,\omega)$ with $x$ over $\tilde{X}$ and the $\chi(x,x',\beta,\omega)$ with $(x,x')\in\mathrm{Pad}(\tilde{X})$;
\item for $k\ge1$, $\mathrm{GrqP}_i(\tilde{X})$ is the subalgebra generated by $\mathrm{Grq}_i(\tilde{X})$ and the $\Psi(y,\beta,\omega)$, $y$ a unit site of $\tilde{X}$;
\item $\mathrm{GrqW}_i(\tilde{X})$ is the subalgebra generated by $\mathrm{Grq}_i(\tilde{X})$ and the $W(y,\beta,\omega)$, $y$ a unit site of $\tilde{X}$, if $k\ge1$, respectively the $W(x,\beta,\omega)$, $x$ over $\tilde{X}$, if $k=0$.
\end{itemize}
The following facts hold.
\begin{itemize}
\item[(Gq0)] The image in $\mathrm{Pz}_k(T_i)$ of a pair $(x,x')\in\mathrm{Pad}(\tilde{X})$ is admissible for $\mathrm{pr}_i(\tilde{X})$, and it is of the same kind.
Indeed, both photon sites of $(x,x')$ lie in $N(\tilde{X})$, so that $(x,x')$ is the shortest lift of its image (Proposition~27.1\,(d)).
The six paths of the image are formed on this lift (Sections~I.1.1 and I.1.7); they are therefore the images of the six paths $\boldsymbol{\Gamma}^{\pi}(c,c')$.
Finally $(\mathrm{pr}_ix)^{(3)}=\mathrm{pr}_i(x^{(3)})$.
Combined with Proposition~27.2\,(b), this shows that $\mathrm{pr}_i$ maps $\mathrm{Pad}(\tilde{X})$ bijectively onto the set of the pairs of $\mathrm{Pz}_k(T_i)$ that are admissible for $\mathrm{pr}_i(\tilde{X})$.
\item[(Gq1)] The map $\mathrm{idn}$ carries $\mathrm{Grq}_1(\tilde{X})$, $\mathrm{GrqP}_1(\tilde{X})$, $\mathrm{GrqW}_1(\tilde{X})$ onto $\mathrm{Grq}_2(\tilde{X})$, $\mathrm{GrqP}_2(\tilde{X})$, $\mathrm{GrqW}_2(\tilde{X})$, respectively, since it sends every generator to the generator with the same index.
\item[(Gq2)] Let $\tilde{X}\subset\tilde{X}'$, where $\tilde{X}'$ is connected in $\tilde{Y}$.
Then $\mathrm{Pad}(\tilde{X})\subset\mathrm{Pad}(\tilde{X}')$, because the conditions of both kinds are monotone in the set.
Consequently $\mathrm{Grq}_i(\tilde{X})\subset\mathrm{Grq}_i(\tilde{X}')$, $\mathrm{GrqP}_i(\tilde{X})\subset\mathrm{GrqP}_i(\tilde{X}')$ and $\mathrm{GrqW}_i(\tilde{X})\subset\mathrm{GrqW}_i(\tilde{X}')$.
\item[(Gq3)] Let $F$ be an element of $\mathrm{Gr}[\psi,\chi,\Psi,W]$ on $T_i$ such that every generator $\psi(\xi)$ contained in $F$ has $\xi$ over $\mathrm{pr}_i(\tilde{X})$ and every generator $\chi(\boldsymbol{\zeta})$ contained in $F$ has $\boldsymbol{\zeta}$ admissible for $\mathrm{pr}_i(\tilde{X})$.
If $F$ contains no generator $\Psi$ and no generator $W$, then $F\in\mathrm{Grq}_i(\tilde{X})$.
If the only further generators contained in $F$ are generators $\Psi(u)$ with $u$ in $\mathrm{pr}_i(\tilde{X})$, then $F\in\mathrm{GrqP}_i(\tilde{X})$.
If the only further generators contained in $F$ are generators $W(u)$ with $u$ in (over) $\mathrm{pr}_i(\tilde{X})$, then $F\in\mathrm{GrqW}_i(\tilde{X})$.
For the proof, recall from (I.AP) that the generators at two lifts of one index of $T_i$ differ by a sign; hence the subalgebra generated by the generators at a set of indices of $T_i$ is generated by the generators at one lift of each of these indices.
By Proposition~27.2\,(b), a fermion site over $\mathrm{pr}_i(\tilde{X})$ has exactly one lift over $\tilde{X}$, and a pair that is admissible for $\mathrm{pr}_i(\tilde{X})$ has exactly one lift whose two photon sites lie in $\tilde{X}^{\sim1}$; this lift belongs to $\mathrm{Pad}(\tilde{X})$.
By Proposition~27.2\,(e), a unit site of $\mathrm{pr}_i(\tilde{X})$ has exactly one lift in $\tilde{X}$.
\end{itemize}

\subsubsection*{0.5. Adding fields}

Let $\tilde{X}$ be connected in $\tilde{Y}$.
Let $v$ be a family of odd elements $v(\xi)$, $v(\boldsymbol{\zeta})$ of $\mathrm{Gr}_i(\tilde{Y})$, indexed by the $\xi=(x,\beta,\omega)$ with $x$ over $\tilde{X}$ and by the $\boldsymbol{\zeta}=(x,x',\beta,\omega)$ with $(x,x')\in\mathrm{Pad}(\tilde{X})$.
Define
\[
\mathrm{adf}_i[v]\colon\mathrm{Grq}_i(\tilde{X})\to\mathrm{Gr}_i(\tilde{Y})
\]
as the homomorphism of algebras with $\psi(\xi)\mapsto\psi(\xi)+v(\xi)$ and $\chi(\boldsymbol{\zeta})\mapsto\chi(\boldsymbol{\zeta})+v(\boldsymbol{\zeta})$.
It exists and is unique, because the generators are independent and their images are odd.
Let $v^{\#}$ be a family on $T_i$ as in Definition~24.D\,(D1), paragraph \emph{Adding and inserting fields}, whose members at the indices just named lie in $\mathrm{Gr}_i(\tilde{Y})$, and let $v$ be the restriction of $v^{\#}$ to these indices.
Then
\[
F(\psi^{\#}+v^{\#})=\mathrm{adf}_i[v]\,F\qquad\text{for every }F\in\mathrm{Grq}_i(\tilde{X}).
\]
\begin{itemize}
\item[(Ad)] Let $v^{(1)}$ and $v^{(2)}$ be families of this kind on $T_1$ and on $T_2$, with $\mathrm{idn}\,v^{(1)}(\xi)=v^{(2)}(\xi)$ and $\mathrm{idn}\,v^{(1)}(\boldsymbol{\zeta})=v^{(2)}(\boldsymbol{\zeta})$ at all the indices named.
Then
\[
\mathrm{idn}\left(\mathrm{adf}_1[v^{(1)}]\,F\right)=\mathrm{adf}_2[v^{(2)}]\left(\mathrm{idn}\,F\right)\qquad\text{for all }F\in\mathrm{Grq}_1(\tilde{X}).
\]
Here $\mathrm{idn}\,F\in\mathrm{Grq}_2(\tilde{X})$ by (Gq1).
Both sides are homomorphisms of algebras $\mathrm{Grq}_1(\tilde{X})\to\mathrm{Gr}_2(\tilde{Y})$ in the variable $F$, and they agree on the generators: at $\psi(\xi)$ both take the value $\psi(\xi)+v^{(2)}(\xi)$, and at $\chi(\boldsymbol{\zeta})$ the value $\chi(\boldsymbol{\zeta})+v^{(2)}(\boldsymbol{\zeta})$.
Hence they coincide.
\end{itemize}

\subsubsection*{0.6. Two forms}

Let $\tilde{Y}$ be admissible.
\begin{itemize}
\item[(Fm1)] Let $\tilde{X}'$ be a non-empty union of $M$-cubes contained in $\tilde{Y}$, and let $\mathrm{Mx}$ be a complex $4\times4$ matrix.
On $T_i$,
\[
\int_{\mathrm{pr}_i\tilde{X}'}\mathrm{d}x\,\bar{\psi}(x)^T\cdot\mathrm{Mx}\cdot\psi(x)=\frac{\eta_k^3}{48}\cdot\sum_{c\in\tilde{X}'}\ \sum_{g\in O_{\mathrm{cub}}}\ \sum_{\beta,\beta'}\mathrm{Mx}_{\beta\beta'}\cdot\bar{\psi}_{\beta}(c,g)\cdot\psi_{\beta'}(c,g).
\]
In particular this element lies in $\mathrm{Gr}_i(\tilde{Y})$, and in $\mathrm{Grq}_i(\tilde{X}')$ if $\tilde{X}'$ is connected; and $\mathrm{idn}$ maps the element on $T_1$ to the element on $T_2$.
For $\mathrm{Mx}=\mathbf{1}_4$ and an $M$-cube $\tilde{X}'=\tilde{\square}$ the element is $\langle\bar{\psi},\psi\rangle_{\mathrm{pr}_i\tilde{\square}}$, and in this case the assertion is Proposition~27.4\,(b).
To prove the display, note that by (I.1.3) the left side equals $\frac{\eta_k^3}{48}$ times the sum of $\mathrm{smd}(x):=\bar{\psi}(x)^T\cdot\mathrm{Mx}\cdot\psi(x)$ over the $x=(c,g)$ with $c$ in a fundamental domain of $\ell_i\mathbb{Z}^3$ and $\mathrm{pr}_i(c)\in\mathrm{pr}_i(\tilde{X}')$.
By (I.AP), $\mathrm{smd}(x+\ell_ie_{\mu})=\mathrm{smd}(x)$, the two signs cancelling.
The rule (Rs) with $\tilde{S}=\tilde{X}'$, applied to $\mathbf{1}[\mathrm{pr}_i(c)\in\mathrm{pr}_i(\tilde{X}')]\cdot\mathrm{smd}(x)$, gives the display.
\item[(Fm2)] Let $k\ge1$, let $\tilde{X}$ be connected in $\tilde{Y}$, and let $O$ be an operator on $\mathfrak{S}(1,\ell_i)$ whose kernel $O(y,y')$ vanishes unless both unit sites lie in $\mathrm{pr}_i(\tilde{X})$.
Then, for the form of Section~I.3.6, paragraph \emph{Forms},
\[
\langle\bar{\Psi},O\Psi\rangle=\sum_{y,y'\in\tilde{X}\cap\mathbb{Z}^3}\ \sum_{\beta,\beta'}O(y,y')_{\beta\beta'}\cdot\bar{\Psi}_{\beta}(y)\cdot\Psi_{\beta'}(y').
\]
Indeed, the summand $\bar{\Psi}(y)\cdot O(y,y')\Psi(y')$ is unchanged under $y\mapsto y+\ell_ie_{\mu}$ and under $y'\mapsto y'+\ell_ie_{\mu}$, since in each case two signs cancel (Section~G4.0.4, (I.AP)).
The rule (Rs) with $\tilde{S}=\tilde{X}$, applied in $y$ and in $y'$, gives the display.
\end{itemize}

\subsubsection*{0.7. The reading (o)}

Proposition~27.6\,(e) concerns the step map without slot, Definition~24.D\,(D7).
The proof of Proposition~27.5 in Section~1 is given for the step map with slot.
The reading~(o) of Definition~24.D\,(D7) is extended to assertions and arguments on the step map with slot: such an assertion or argument is taken \emph{in the reading (o)} when the data, the objects and the results of (o-1)--(o-4) below are put in the places of the corresponding ones with slot.
\begin{itemize}
\item[(o-1)] \emph{Data.}
Let (TT) hold with $k<k_m$.
Let $m$ be real with $|m|\le w^{\mathrm{m}}_k$, and let $E_{(i)}\in K_{k}(\ell_i)$ have norm at most $w^{\mathrm{E}}_k$ ($E_{(i)}=0$ if $k=0$), with $E_{(1)}\sim E_{(2)}$ in the sense of the last sentence of Proposition~27.2\,(D4); these are the data of Proposition~27.6\,(e).
The number of the torus is written as a lower index, because the upper indices in $E^{(j)}$, $j=1,2,3$, mark the localised objects of Definition~24.D\,(D7).
Let $\tilde{Y}$ be admissible and $\mathscr{A}\in\tilde{R}_{k+1}$ on $T_{k+1}(R)$, and put
\[
A_1:=\mathscr{A}_L,\qquad A_2:=\mathrm{pj}^{*}A_1=\left(\mathrm{pj}^{*}\mathscr{A}\right)_L .
\]
By \cite[Lemma~5.18]{QEDI} and \cite[Lemma~5.20]{QEDI}, $A_i$ lies in $\tilde{R}_k^{-}(T_i)$, and $\mathrm{pj}^{*}\mathscr{A}$ lies in $\tilde{R}_{k+1}$ on $T_{k+1}(R_{\mathrm{top}})$.
By (TT.3), $A_i$ and $0$ lie in $\tilde{R}_k(T_i)$, which is a subset of $\tilde{G}_k(T_i)$; by (TT.2), $A_1$ and $A_2$ have equal lifts on all bonds.
There is no $\mathscr{C}_i$, and $\mathscr{T}_i:=A_i$.
Further $m'^{(i)}:=m+m(E_{(i)})$ and $E^{r,(i)}:=R\,E_{(i)}$.
The pair $(m,E_{(i)})$ lies in $\mathrm{Ball}_{k}(1)$, which is a subset of $\mathrm{Ball}_{k}(4)$.
The objects $\tilde{\mathbf{Y}}$, $\mathbf{Y}_i$, $\mathrm{Bs}$, $\mathrm{zval}$, $\mathsf{Z}^{(i)}$, $\mathrm{sval}$, $s^{(i)}$ are those of Proposition~27.5.
\item[(o-2)] \emph{Objects.}
Every product $\mathrm{Dm}^{\mathrm{o}}\times\mathfrak{X}_k$, where $\mathrm{Dm}^{\mathrm{o}}=\tilde{R}_k$ or $\mathrm{Dm}^{\mathrm{o}}=\tilde{R}_k^{-}$, is read $\mathrm{Dm}^{\mathrm{o}}$.
The objects $\mathscr{E}^{(i)}$, $\mathscr{E}^{r,(i)}$, $E^{(j),x}$, $E^{(\le3),x}$, $E^{(\le5),x}$, $E^{\#,x}_k$ are read $E_{(i)}$, $E^{r,(i)}$, $E^{(j)}$, $E^{(\le3)}$, $E^{(\le5)}$, $E^{\#}_k$; the arguments $\mathscr{C}$ and $\mathscr{C}_i$ are deleted; and $\mathrm{Km}$, $\mathrm{Hm}$ are formed with $E^{(\le5)}$.
These are the objects of Definition~24.D\,(D7).
\item[(o-3)] \emph{Results used.}
In the reading (o), Proposition~16.3\,(a), (d), (f) is replaced by Proposition~48.2\,(a), (d), (f); Proposition~17.2\,(c), (f) by Proposition~48.3\,(c), (f); the property (k2) by the property (k2)$^{\mathrm{o}}$ of Proposition~48.1; Proposition~24.3 by its last sentence, which asserts that the same holds in the reading (o); and Proposition~24.1\,(a) by Proposition~24.1\,(b) in the following form:
$E^{*}_k=L\,E^{\#}_k$ is a map defined on $\mathrm{Ball}_{k}(4)$, with values in families on $\tilde{R}_{k+1}$.
In other words, for $(m,E)\in\mathrm{Ball}_{k}(4)$, for every $X\in D_k$ and for every connection $A=\mathscr{A}_L$ with $\mathscr{A}\in\tilde{R}_{k+1}$, all displays of Definition~24.D in the reading (D7) that lead to $E^{\#}_k(m,E)(X;A)$ are defined, and the series (8.4) of Definition~24.D for this element converges.
Every $X\in D_k$ occurs here, since there is a $Y\in D_{k+1}$ with $\bar{X}^{L}=LY$ \cite[Theorem~13.2(t5)]{QEDI}.
The form of Proposition~24.1\,(b) just stated is the special case $A=\mathscr{A}_L$ of the first assertion of Proposition~24.1\,(b), which holds for every $A\in\tilde{R}_k^{-}$: indeed $\mathscr{A}_L\in\tilde{R}_k^{-}$ for every $\mathscr{A}\in\tilde{R}_{k+1}$ \cite[Lemma~5.18(a)]{QEDI}.
The Propositions~27.1, 27.3, 27.4, 23.1 and 23.2, the properties W7$^{\mathrm{ch}}$ and W12$^{\mathrm{ch}}$ of Section~G4.0.4, (I.Q3), \cite[Lemmas~5.17--5.20]{QEDI} and the facts of Sections~0.2--0.6 contain no slot; they hold in the reading (o) as they stand.
In Proposition~27.4\,(d) the connections are the $\mathscr{T}_i'=\mathscr{A}_i$ of (a.3), which are arbitrary elements of $\tilde{R}_k(T_i)$ with equal lifts on $N(\tilde{X})$, and not the $\mathscr{T}_i=A_i$ of (o-1); and $\tilde{R}_k(T_i)$ is a subset of $\left(\tilde{R}_k+\mathfrak{X}_k\right)(T_i)$, because $0\in\mathfrak{X}_k$.
Among the facts of Section~0.1, the slot occurs only in the final assertions of (TT.1) and of (TT.3), which concern the pairs $(m,\mathscr{E}^{(i)})$ and the data of Proposition~27.5; in the reading (o), the assertions of (o-1) on the data of (o-1) take their place.
\item[(o-4)] \emph{Local families without slot.}
They are defined in Proposition~27.2\,(D3), and their relation $\sim$ in the last sentence of Proposition~27.2\,(D4).
By Proposition~27.2\,(g), the clauses (b) and (c) of Proposition~27.2 hold for them, with $\mathscr{C}$ and $\mathscr{C}_i$ deleted and with $\mathrm{Dm}^{\mathrm{o}}$ in the place of $\mathrm{Dm}$.
\end{itemize}

\subsubsection*{0.8. Lemma 27.A (the low kernels agree)}

\begin{named}{Lemma 27.A}[The low kernels agree]
Let (TT) hold with $k<k_m$, let $\mathscr{E}^{(i)}\in K^{x}_{k}(\ell_i)$ for $i=1,2$, and let $\mathscr{E}^{(1)}\sim\mathscr{E}^{(2)}$.
\begin{itemize}
\item[(i)] For every connected non-empty union $\tilde{X}$ of $M$-cubes of $\mathbb{L}_k$ with $d_M(\tilde{X})\le L$,
\[
\mathrm{al}_{0}(\mathscr{E}^{(1)},\mathrm{pr}_1\tilde{X})=\mathrm{al}_{0}(\mathscr{E}^{(2)},\mathrm{pr}_2\tilde{X}),\qquad
\mathrm{al}_{2}(\mathscr{E}^{(1)},\mathrm{pr}_1\tilde{X})=\mathrm{al}_{2}(\mathscr{E}^{(2)},\mathrm{pr}_2\tilde{X}).
\]
\item[(ii)] $\varepsilon(\mathscr{E}^{(1)})=\varepsilon(\mathscr{E}^{(2)})$ and $m(\mathscr{E}^{(1)})=m(\mathscr{E}^{(2)})$.
\end{itemize}
The same holds for $E_{(i)}\in K_{k}(\ell_i)$ with $E_{(1)}\sim E_{(2)}$.
\end{named}

\begin{proof}
(i) By Proposition~27.2\,(d), $\tilde{X}$ is admissible, $\mathrm{pr}_i(\tilde{X})$ is a small polymer of $T_i$, and $\operatorname{Vol}(\mathrm{pr}_i\tilde{X})=\operatorname{Vol}(\tilde{X})$.
By (TT.3) the pairs $(0,0)$ belong to $\tilde{R}_k(T_i)\times\mathfrak{X}_k(T_i)$, and they have equal lifts on $N(\tilde{X})$.
Put $F_i:=\mathscr{E}^{(i)}(\mathrm{pr}_i\tilde{X};0,0)$ and $\mathrm{idn}=\mathrm{idn}[\tilde{X}]$.
Since $\mathscr{E}^{(1)}\sim\mathscr{E}^{(2)}$, Proposition~27.2\,(D4) gives $\mathrm{idn}\,F_1=F_2$.
The element $F_i$ is even.
By the locality property (k2) of Section~I.5.4, every generator $\psi(\xi)$ that $F_i$ contains has $\xi$ over $\mathrm{pr}_i(\tilde{X})$, and every generator $\chi(\boldsymbol{\zeta})$ that it contains has $\boldsymbol{\zeta}$ admissible for $\mathrm{pr}_i(\tilde{X})$; hence $F_i\in\mathrm{Grq}_i(\tilde{X})$ by (Gq3).
Apply Proposition~27.4\,(c) with the admissible set $\tilde{X}$ in the place of $\tilde{Y}$: the numbers $(F_i)_{00}$ and the matrices $\int\mathrm{d}x\,\mathrm{d}x'\,(F_i)_{20}(x,x')$, formed with all arguments over the lift $\tilde{X}$, coincide for $i=1$ and $i=2$.
By Section~I.5.6, together with the last sentence of (I.AP), $\mathrm{al}_{0}(\mathscr{E}^{(i)},\mathrm{pr}_i\tilde{X})$ and $\mathrm{al}_{2}(\mathscr{E}^{(i)},\mathrm{pr}_i\tilde{X})$ are these numbers and these matrices divided by $\operatorname{Vol}(\tilde{X})$.
This proves (i).

The same conclusion can be read off the coefficients.
By Section~I.5.1,
\begin{align*}
F_i^{[2,0]}&=\int\mathrm{d}x\,\mathrm{d}x'\,\bar{\psi}(x)^T\cdot(F_i)_{20}(x,x')\cdot\psi(x')\\
&=\left(\frac{\eta_k^3}{48}\right)^2\cdot\sum_{x,x'}\ \sum_{\beta,\beta'}(F_i)_{20}(x,x')_{\beta\beta'}\cdot\bar{\psi}_{\beta}(x)\cdot\psi_{\beta'}(x'),
\end{align*}
the sum extending over the $x=(c,g)$ and $x'=(c',g')$ with $c,c'\in\tilde{X}$ and $g,g'\in O_{\mathrm{cub}}$.
Thus $\left(\frac{\eta_k^3}{48}\right)^2\cdot(F_i)_{20}(x,x')_{\beta\beta'}$ is the coefficient of $\bar{\psi}_{\beta}(x)\cdot\psi_{\beta'}(x')$ in $F_i$.
The entry $(\beta,\beta')$ of
\[
\int\mathrm{d}x\,\mathrm{d}x'\,(F_i)_{20}(x,x')=\left(\frac{\eta_k^3}{48}\right)^2\cdot\sum_{c,c'\in\tilde{X}}\ \sum_{g,g'\in O_{\mathrm{cub}}}(F_i)_{20}(x,x')
\]
is therefore the sum of these coefficients over the sites and over both labels.
Moreover $(F_i)_{00}=c_{F_i}(\emptyset)$.
Together with (Co1) this gives the assertion once more.

(ii) Fix an $M$-cube $\tilde{\square}$ of $\mathbb{L}_k$.
By Proposition~27.2\,(d), $\tilde{X}\mapsto\mathrm{pr}_i(\tilde{X})$ is a bijection from the set of the connected $\tilde{X}$ with $d_M(\tilde{X})\le L$ that contain $\tilde{\square}$ onto the set of the small polymers of $T_i$ that contain $\mathrm{pr}_i(\tilde{\square})$.
By Section~I.5.6 and Proposition~16.3\,(a), according to which the two sums below do not depend on the $M$-cube and the second of them is a multiple of $\mathbf{1}_4$,
\[
\varepsilon(\mathscr{E}^{(i)})=-\sum_{\tilde{X}}\mathrm{al}_{0}(\mathscr{E}^{(i)},\mathrm{pr}_i\tilde{X}),\qquad
m(\mathscr{E}^{(i)})\cdot\mathbf{1}_4=-\sum_{\tilde{X}}\mathrm{al}_{2}(\mathscr{E}^{(i)},\mathrm{pr}_i\tilde{X}),
\]
where both sums extend over this set of $\tilde{X}$, which is the same for $i=1$ and for $i=2$.
By (i) the terms coincide, and (ii) follows.

Without slot the argument is the same, with $F_i:=E_{(i)}(\mathrm{pr}_i\tilde{X};0)$, with (k2)$^{\mathrm{o}}$, and with Proposition~48.2 (its definitions of $\mathrm{al}_{0}$, $\mathrm{al}_{2}$, $\varepsilon$, $m$, and its clause (a)) in the places of $\mathscr{E}^{(i)}(\mathrm{pr}_i\tilde{X};0,0)$, (k2), Section~I.5.6 and Proposition~16.3\,(a).
Proposition~27.2\,(d), Proposition~27.4\,(c) and (Gq3) contain no slot.
\end{proof}

\subsection*{1. Proposition 27.5}

\subsubsection*{Proof of Proposition 27.5}

The setting is that of Section~0.1.
From Step (b) on, $\mathrm{idn}=\mathrm{idn}[\tilde{Y}]$, see (Co4).
Step (a) concerns arbitrary connected admissible sets $\tilde{X}$, which need not be contained in $\tilde{Y}$; there $\mathrm{idn}=\mathrm{idn}[\tilde{X}]$, and the facts of Sections~0.2--0.6 are applied with $\tilde{X}$ in the place of $\tilde{Y}$.

\paragraph{Step (a): extraction, adjustment}

(a.1) By Lemma~27.A\,(ii), $m(\mathscr{E}^{(1)})=m(\mathscr{E}^{(2)})$, and therefore
\[
m'^{(1)}=m+m(\mathscr{E}^{(1)})=m+m(\mathscr{E}^{(2)})=m'^{(2)}.
\]
This common number is denoted by $m'$.

(a.2) \emph{The extraction $R$.}
By Proposition~16.3\,(d), $R\mathscr{E}^{(i)}$ lies in $K^{x}_{k}(\ell_i)$.
So it consists of even elements and has the property (k2); that is, it is a local family on $\tilde{R}_k\times\mathfrak{X}_k$ in the sense of Proposition~27.2\,(D3).
Let $\tilde{X}$ be connected and admissible, and let $(\mathscr{A}_i,\mathscr{C}_i)\in\tilde{R}_k(T_i)\times\mathfrak{X}_k(T_i)$ have equal lifts on $N(\tilde{X})$.
The hypothesis $\mathscr{E}^{(1)}\sim\mathscr{E}^{(2)}$ gives
\begin{equation}
\mathrm{idn}\,\mathscr{E}^{(1)}(\mathrm{pr}_1\tilde{X};\mathscr{A}_1,\mathscr{C}_1)=\mathscr{E}^{(2)}(\mathrm{pr}_2\tilde{X};\mathscr{A}_2,\mathscr{C}_2).\tag{27.1}
\end{equation}
By Proposition~27.2\,(d), $\mathrm{pr}_1(\tilde{X})$ is small if and only if $d_M(\tilde{X})\le L$, and this holds if and only if $\mathrm{pr}_2(\tilde{X})$ is small.
If $d_M(\tilde{X})>L$, then $\left(R\mathscr{E}^{(i)}\right)(\mathrm{pr}_i\tilde{X})=\mathscr{E}^{(i)}(\mathrm{pr}_i\tilde{X})$ for $i=1,2$ (Section~I.5.6), and (27.1) is the assertion.
Let $d_M(\tilde{X})\le L$, and write
\[
\mathrm{al}_{0,i}:=\mathrm{al}_{0}(\mathscr{E}^{(i)},\mathrm{pr}_i\tilde{X}),\qquad\mathrm{al}_{2,i}:=\mathrm{al}_{2}(\mathscr{E}^{(i)},\mathrm{pr}_i\tilde{X}).
\]
Section~I.5.6, the equality $\operatorname{Vol}(\mathrm{pr}_i\tilde{X})=\operatorname{Vol}(\tilde{X})$ of Proposition~27.2\,(d) and (Fm1) give
\begin{multline*}
\left(R\mathscr{E}^{(i)}\right)(\mathrm{pr}_i\tilde{X};\mathscr{A}_i,\mathscr{C}_i)=\mathscr{E}^{(i)}(\mathrm{pr}_i\tilde{X};\mathscr{A}_i,\mathscr{C}_i)-\mathrm{al}_{0,i}\cdot\operatorname{Vol}(\tilde{X})\\
-\frac{\eta_k^3}{48}\cdot\sum_{c\in\tilde{X}}\ \sum_{g\in O_{\mathrm{cub}}}\ \sum_{\beta,\beta'}(\mathrm{al}_{2,i})_{\beta\beta'}\cdot\bar{\psi}_{\beta}(c,g)\cdot\psi_{\beta'}(c,g).
\end{multline*}
By Lemma~27.A\,(i), $\mathrm{al}_{0,1}=\mathrm{al}_{0,2}$ and $\mathrm{al}_{2,1}=\mathrm{al}_{2,2}$.
The map $\mathrm{idn}$ fixes the numbers and sends every generator to the generator with the same index.
Together with (27.1) this yields
\[
\mathrm{idn}\,\left(R\mathscr{E}^{(1)}\right)(\mathrm{pr}_1\tilde{X};\mathscr{A}_1,\mathscr{C}_1)=\left(R\mathscr{E}^{(2)}\right)(\mathrm{pr}_2\tilde{X};\mathscr{A}_2,\mathscr{C}_2).
\]
Hence $R\mathscr{E}^{(1)}\sim R\mathscr{E}^{(2)}$.

(a.3) \emph{The adjustment.}
Recall that $\mathscr{E}^{r,(i)}=R\mathscr{E}^{(i)}$.
The family $\left(\mathscr{E}^{r,(i)}\right)^{\mathrm{nat}}$ has the property (k2), by Proposition~17.2\,(c) applied to the family $\mathscr{E}^{r,(i)}$, which has (k2).
Its elements are even.
Indeed, $\boldsymbol{\Theta}_{X,\mathscr{T}}$ is a homomorphism that maps every generator $\psi(\xi)$ to a linear combination of generators (I.5.3); so it maps the part of bidegree $(2,0)$ to an element of degree two.
Hence $\left(\mathscr{E}^{r,(i)}\right)^{\mathrm{nat}}$ is a local family on $\tilde{R}_k\times\mathfrak{X}_k$.
Let $\tilde{X}$ and $(\mathscr{A}_i,\mathscr{C}_i)$ be as in (a.2), and put $\mathscr{T}_i':=\mathscr{A}_i+\mathscr{C}_i$.
If $d_M(\tilde{X})>L$, then $\left(\mathscr{E}^{r,(i)}\right)^{\mathrm{nat}}(\mathrm{pr}_i\tilde{X})=\mathscr{E}^{r,(i)}(\mathrm{pr}_i\tilde{X})$, and (a.2) is the assertion.
Let $d_M(\tilde{X})\le L$, and let $F_i$ denote the part of bidegree $(2,0)$ of $\mathscr{E}^{r,(i)}(\mathrm{pr}_i\tilde{X};\mathscr{A}_i,\mathscr{C}_i)$.
By Section~I.5.6,
\[
\left(\mathscr{E}^{r,(i)}\right)^{\mathrm{nat}}(\mathrm{pr}_i\tilde{X};\mathscr{A}_i,\mathscr{C}_i)=\mathscr{E}^{r,(i)}(\mathrm{pr}_i\tilde{X};\mathscr{A}_i,\mathscr{C}_i)-F_i+\boldsymbol{\Theta}_{\mathrm{pr}_i\tilde{X},\mathscr{T}_i'}(F_i).
\]
Since $\mathrm{idn}$ preserves the bidegree (Proposition~27.2\,(a)), (a.2) gives $\mathrm{idn}\,F_1=F_2$.
By (k2), $F_i$ lies in the subalgebra generated by the $\psi(\xi)$ with $\xi$ over $\tilde{X}$.
The connection $\mathscr{T}_i'$ lies in $\left(\tilde{R}_k+\mathfrak{X}_k\right)(T_i)$, and $\mathscr{T}_1'$, $\mathscr{T}_2'$ have equal lifts on $N(\tilde{X})$, hence on the bonds contained in $\tilde{X}$.
Proposition~27.4\,(d), applied with the admissible set $\tilde{X}$ in the place of $\tilde{Y}$, therefore gives
\[
\mathrm{idn}\,\boldsymbol{\Theta}_{\mathrm{pr}_1\tilde{X},\mathscr{T}_1'}(F_1)=\boldsymbol{\Theta}_{\mathrm{pr}_2\tilde{X},\mathscr{T}_2'}(\mathrm{idn}\,F_1)=\boldsymbol{\Theta}_{\mathrm{pr}_2\tilde{X},\mathscr{T}_2'}(F_2).
\]
Combining this with (a.2), one obtains
\[
\mathrm{idn}\,\left(\mathscr{E}^{r,(1)}\right)^{\mathrm{nat}}(\mathrm{pr}_1\tilde{X};\mathscr{A}_1,\mathscr{C}_1)=\left(\mathscr{E}^{r,(2)}\right)^{\mathrm{nat}}(\mathrm{pr}_2\tilde{X};\mathscr{A}_2,\mathscr{C}_2).
\]
The chain of (I.5.3) from $x=(c,g)$ to $\tilde{x}=(\tilde{c},\tilde{g})$ consists of the fermion site $x$, then of the sites $(z,g)$ over the centres $z$ of the $M$-cubes of the chain, all with the label $g$ of the starting point, and finally of $\tilde{x}$.
The integral $\int_X\mathrm{d}\tilde{x}$ contains the factor $\frac{1}{48}$ and the sum over $\tilde{g}\in O_{\mathrm{cub}}$.
The last pair of a chain has the labels $(g,\tilde{g})$.
If $\tilde{c}$ is the last centre, this pair lies inside a cluster when $\tilde{g}\ne g$, while for $\tilde{g}=g$ its two points coincide and the term is omitted.
All these data are formed over the lift $\tilde{X}$, and the fact that they are the same for $i=1$ and for $i=2$ is the content of Proposition~27.4\,(d).

\emph{Level $0$.}
Here $\mathscr{E}^{(i)}=0$; hence $m'=m$, $R\mathscr{E}^{(i)}=0$ and $\left(R\mathscr{E}^{(i)}\right)^{\mathrm{nat}}=0$.

\emph{In the reading (o).}
Step (a) gives $m'^{(1)}=m'^{(2)}$, $R\,E_{(1)}\sim R\,E_{(2)}$ and $\left(R\,E_{(1)}\right)^{\mathrm{nat}}\sim\left(R\,E_{(2)}\right)^{\mathrm{nat}}$, as local families without slot on $\tilde{R}_k$.
It rests on Lemma~27.A without slot, on Proposition~48.2\,(d) ($R$ maps $K_k$ into itself), on the display of $R\,E$ in Proposition~48.2, and on Proposition~48.3\,(c) together with the display of $E^{\mathrm{nat}}$ in Proposition~48.3, formed with $\boldsymbol{\Theta}_{X,\mathscr{A}}$; here $\mathscr{T}_i'=\mathscr{A}_i$.

\paragraph{Step (b): determinant}

Let $\tilde{X}$ be connected in $\tilde{Y}$, put $X_i:=\mathrm{pr}_i(\tilde{X})$, and let $B_i\in\tilde{G}_k(T_i)$, $i=1,2$, be connections with equal lifts on the bonds that meet $\tilde{X}$.
Then Proposition~27.3\,(c) and Proposition~27.4\,(a), the latter with $\tilde{\mathbf{Y}}:=\tilde{X}$, are applicable to the pair $(B_1,B_2)$.
In (b.1), (b.2), (b.4) and in the first part of (b.6) only one torus is involved:
there $\ell$ stands for $\ell_1$ or for $\ell_2$, $X$ is a polymer of $T_k=\mathbb{L}_k/\ell\mathbb{Z}^3$, $B\in\tilde{G}_k$, and $y\ge0$.

\emph{Levels $k\ge1$.}

(b.1) \emph{The kernel of $B_{k,y}(B)$.}
Let $v\in\mathfrak{S}(1,\ell)$, and let $z$ be a unit site of the $L$-block $B(y')$.
By Section~I.3.3,
\[
\left(P(B)\cdot v\right)(z)=U_B(y',z)^{-1}\cdot L^{-3}\sum_{z'\in B(y')}U_B(y',z')\cdot v(z').
\]
The kernel $P(B)(z,z')$ is therefore a multiple of $\mathbf{1}_4$, and it vanishes unless $z$ and $z'$ belong to one $L$-block of $T^0_k$.
Moreover $Q(B)\cdot Q^T(-B)=1$, because $L^{-3}\sum_{z\in B(y')}U_B(y',z)\cdot U_B(y',z)^{-1}=1$ (Section~I.3.3); consequently
\[
P(B)^2=Q^T(-B)\cdot Q(B)\cdot Q^T(-B)\cdot Q(B)=P(B).
\]
Let $\mathrm{bv}$ denote either of the two real numbers $b_k$ and $b_k+b_L$.
One has $b_k\ge b_-=b\,(1-L^{-1})>0$ and $b_L=b/L>0$ (Sections~I.3.4 and~I.0.5), and $\gamma_5^2=1$; hence
\[
\left(\mathrm{bv}+i\gamma_5y\right)\cdot\left(\mathrm{bv}-i\gamma_5y\right)=\mathrm{bv}^2+y^2>0 .
\]
Thus the matrices $b_k+i\gamma_5y$ and $b_k+b_L+i\gamma_5y$ are invertible, and they commute with $P(B)$, whose kernel consists of multiples of $\mathbf{1}_4$.
One has
\begin{equation}
B_{k,y}(B)=\left(b_k+i\gamma_5y\right)^{-1}\cdot\left(1-P(B)\right)+\left(b_k+b_L+i\gamma_5y\right)^{-1}\cdot P(B).\tag{27.2}
\end{equation}
Indeed, since $P(B)\cdot\left(1-P(B)\right)=0$, the product of $b_k+b_L\cdot P(B)+i\gamma_5y$ with the right side of (27.2), the former standing on the left, equals $\left(1-P(B)\right)+P(B)=1$.
As $B_{k,y}(B)$ is the inverse of $b_k+b_L\cdot P(B)+i\gamma_5y$ on the finite-dimensional space $\mathfrak{S}(1,\ell)$ (Section~I.3.5), the identity (27.2) follows.
Consequently $B_{k,y}(B)(z,z')=0$ unless $z$ and $z'$ lie in one $L$-block.
Every $L$-block is contained in a single $M$-cube (\cite[Section~17.4.1, (g0)]{QEDI}).

(b.2) \emph{The trace.}
Put
\[
O_1:=B_{k,y}(B)\cdot\gamma_5\cdot B_{k,y}(B),\qquad
O_2:=Q_k(B)\cdot J^{*}\cdot S_{k,y}(X,B)\cdot J\cdot Q_k^T(-B);
\]
both are operators on $\mathfrak{S}(1,\ell)$.
The trace is cyclic, so that
\[
\operatorname{Tr}\left[\gamma_5\cdot B_{k,y}(B)\cdot O_2\cdot B_{k,y}(B)\right]=\operatorname{Tr}\left[O_1\cdot O_2\right].
\]
An element of $\mathfrak{S}(1,\ell)$ is determined by its values on a set of representatives of $\mathbb{Z}^3$ modulo $\ell\mathbb{Z}^3$.
In the basis of the delta functions at these points an operator $O$ is given by its $4\times4$ blocks $O(z,z')$, the kernels being taken as in Section~G4.0.4.
Hence
\[
\operatorname{Tr}\left[O_1\cdot O_2\right]=\sum_{z_1,z_2}\operatorname{tr}\left[O_1(z_2,z_1)\cdot O_2(z_1,z_2)\right],
\]
with $z_1$ and $z_2$ ranging over the representatives.
The summand is invariant under $z_1\mapsto z_1+\ell e_\mu$ and under $z_2\mapsto z_2+\ell e_\mu$, since in each case both factors change sign.
By Section~I.3.3, for $x=(c,g)$,
\[
\left(Q_k(B)\cdot J^{*}f\right)(z)=\int_{\Delta_z}\mathrm{d}x\,\varphi_B(z,c)\cdot f(x),\qquad
\left(J\cdot Q_k^T(-B)\,v\right)(x)=\varphi_B(z,c)^{-1}\cdot v(z)\quad\text{for }x\text{ over }\Delta_z .
\]
Combining this with (I.1.4) and (I.1.3), and taking for the fundamental domain in (I.1.4) the union of the blocks $\Delta_{z_2}$, $z_2$ ranging over the representatives, one obtains, with $x=(c,g)$ and $x'=(c',g')$,
\begin{align}
O_2(z_1,z_2)&=\int_{\Delta_{z_1}}\mathrm{d}x\int_{\Delta_{z_2}}\mathrm{d}x'\,\varphi_B(z_1,c)\cdot S_{k,y}(X,B)(x,x')\cdot\varphi_B(z_2,c')^{-1}\notag\\
&=\left(\frac{\eta_k^3}{48}\right)^2\sum_{c\in\Delta_{z_1}}\sum_{g\in O_{\mathrm{cub}}}\sum_{c'\in\Delta_{z_2}}\sum_{g'\in O_{\mathrm{cub}}}\varphi_B(z_1,c)\cdot S_{k,y}(X,B)\left((c,g),(c',g')\right)\cdot\varphi_B(z_2,c')^{-1}.\tag{27.3}
\end{align}
By W12$^{\mathrm{ch}}$ (one of the properties W1--W17$^{\mathrm{ch}}$ listed in Section~G4.0.4) the kernel of $S_{k,y}(X,B)$ vanishes unless both $x$ and $x'$ lie over $X$, and by (g0) the block $\Delta_z$ is contained in the $M$-cube of $z$.
Therefore $O_2(z_1,z_2)=0$ unless $z_1$ and $z_2$ are unit sites of $X$, and
\[
\operatorname{Tr}\left[\gamma_5\cdot B_{k,y}(B)\cdot Q_k(B)\cdot J^{*}\cdot S_{k,y}(X,B)\cdot J\cdot Q_k^T(-B)\cdot B_{k,y}(B)\right]
=\sum_{z_1,z_2}\operatorname{tr}\left[O_1(z_2,z_1)\cdot O_2(z_1,z_2)\right],
\]
where now $z_1$, $z_2$ range over the unit sites of $X$.
This is the first display in Proposition~27.5\,(b).

(b.3) \emph{Comparison.}
Consider (b.2) on $T_i$ with $X=X_i$ and $B=B_i$.
By (Rs) with $\tilde{S}=\tilde{X}$, the sum extends over the pairs $z_1,z_2\in\tilde{X}\cap\mathbb{Z}^3$.
Let $z_1$, $z_2$ be such sites.
\begin{itemize}
\item
In (27.3) the blocks $\Delta_{z_1}$, $\Delta_{z_2}$ of the covering lattice are contained in $\tilde{X}$, by (g0).
At indices over $\tilde{X}$ the kernels
\[
\left(Q_k(B_i)\cdot J^{*}\right)(z_1,x)=\varphi_{B_i}(z_1,c)\cdot\mathbf{1}_4,\qquad
\left(J\cdot Q_k^T(-B_i)\right)(x',z_2)=\varphi_{B_i}(z_2,c')^{-1}\cdot\mathbf{1}_4
\]
are the same for $i=1$ and for $i=2$ (Proposition~27.4\,(a-4)).
The same holds for the kernels $S^{(i)}_{k,y}(X_i,B_i)(x,x')$ with $x$, $x'$ over $\tilde{X}$ (Proposition~27.3\,(c)).
Hence $O_2^{(1)}(z_1,z_2)=O_2^{(2)}(z_1,z_2)$.
\item
One has
\[
O_1^{(i)}(z_2,z_1)=\sum_zB_{k,y}(B_i)(z_2,z)\cdot\gamma_5\cdot B_{k,y}(B_i)(z,z_1),
\]
$z$ ranging over the representatives.
By (b.1) the term with index $z$ vanishes unless $\mathrm{pr}_i(z)$ lies in the $L$-block of $\mathrm{pr}_i(z_2)$; this block is contained in the $M$-cube of $\mathrm{pr}_i(z_2)$, which is a cube of $X_i$.
By (Rs), the sum therefore extends over $z\in\tilde{X}\cap\mathbb{Z}^3$.
At indices in $\tilde{X}$ the kernels of $B_{k,y}(B_1)$ and of $B_{k,y}(B_2)$ coincide (Proposition~27.4\,(a-4)).
Hence $O_1^{(1)}(z_2,z_1)=O_1^{(2)}(z_2,z_1)$.
\end{itemize}
It follows that, for every $y\ge0$, the trace on $T_1$ equals the trace on $T_2$.

\emph{Level $k=0$.}

(b.4) \emph{The traces.}
For operators $O$, $O'$ on $\mathfrak{S}^{\mathrm{ch}}(1,\ell)$, by (I.1.4), the kernel of the product is
\[
(O\cdot O')(x,x'')=\int\mathrm{d}x'\,O(x,x')\cdot O'(x',x''),
\]
and the kernel of $\gamma_5\cdot O$ is $\gamma_5\cdot O(x,x')$.
By (I.ke2),
\begin{align*}
\operatorname{Tr}\left[F_0(B)\cdot S_{0,y}(X,B)\right]&=\int\mathrm{d}x\int\mathrm{d}x'\,\operatorname{tr}\left[F_0(B)(x,x')\cdot S_{0,y}(X,B)(x',x)\right],\\
\operatorname{Tr}\left[\gamma_5\cdot S_{0,y}(X,B)\right]&=\int\mathrm{d}x\,\operatorname{tr}\left[\gamma_5\cdot S_{0,y}(X,B)(x,x)\right].
\end{align*}
By W12$^{\mathrm{ch}}$ the kernel of $S_{0,y}(X,B)$ vanishes unless both of its arguments lie over $X$.
So all integrals may be restricted to $X$, and the two identities become the second and the third display in Proposition~27.5\,(b).
Here, by (I.1.3) and $\eta_0=1$,
\[
\int_X\mathrm{d}x\,(\,\cdot\,)=\frac{1}{48}\sum_{c\in X}\sum_{g\in O_{\mathrm{cub}}}(\,\cdot\,).
\]

(b.5) \emph{Comparison.}
On $T_i$, with $X=X_i$ and $B=B_i$, the integrands are $\ell_i\mathbb{Z}^3$-periodic in $x$ and in $x'$, the sign changes of the two kernels compensating each other.
By (Rs),
\begin{align*}
\operatorname{Tr}\left[F_0(B_i)\cdot S^{(i)}_{0,y}(X_i,B_i)\right]
&=\left(\frac{1}{48}\right)^2\sum_{c,c'\in\tilde{X}}\sum_{g,g'\in O_{\mathrm{cub}}}\operatorname{tr}\left[F_0(B_i)\left((c,g),(c',g')\right)\cdot S^{(i)}_{0,y}(X_i,B_i)\left((c',g'),(c,g)\right)\right],\\
\operatorname{Tr}\left[\gamma_5\cdot S^{(i)}_{0,y}(X_i,B_i)\right]
&=\frac{1}{48}\sum_{c\in\tilde{X}}\sum_{g\in O_{\mathrm{cub}}}\operatorname{tr}\left[\gamma_5\cdot S^{(i)}_{0,y}(X_i,B_i)\left((c,g),(c,g)\right)\right].
\end{align*}
At the fermion sites over $\tilde{X}$ the kernels $F_0(B_i)(x,x')$ are the same for $i=1,2$ (Proposition~27.4\,(a-5)), and so are the kernels $S^{(i)}_{0,y}(X_i,B_i)(x,x')$ (Proposition~27.3\,(c)).
Consequently, for every $y\ge0$, each of the two traces takes the same value on $T_1$ and on $T_2$.

\emph{All levels.}

(b.6) \emph{Existence of the integrals.}
Like (b.1), (b.2) and (b.4), this part is a statement on a single torus: $X$ is a polymer of $T_k$, $B$ lies in $\tilde{G}_k$, and $S_{k,y}(X,B)$ is the family (p5) of the system of (OS) on this torus (see (TT.1)).
The data (OS) contain (SF), and the family has the properties W1, W2, W7, W12--W14$^{\mathrm{ch}}$ for $\mathrm{par}=(\alpha_{0},\gamma_{\mathrm{w}},\kappa_{\mathrm{w}};C_{w})$; this is the setting in which $\acute{E}_k$ is defined in Section~I.5.7.
In particular W7$^{\mathrm{ch}}$ holds, and the entries of the kernel of $S_{k,y}(X,B)$ are Borel functions of $y\in[0,\infty)$ (Sections~I.5.8 and~G4.0.4).
Let $|\cdot|_c$ be the norm of Section~I.1.5: on $\mathbb{C}^4$ it is the largest modulus of a component, on $4\times4$ matrices the largest row sum of moduli.
It is submultiplicative, and no entry of a matrix exceeds the norm of the matrix in modulus.
Hence, for $4\times4$ matrices $\mathrm{Mx}_1$, $\mathrm{Mx}_2$,
\[
\left|\operatorname{tr}\left(\mathrm{Mx}_1\cdot\mathrm{Mx}_2\right)\right|
\le\sum_{\beta,\beta'}\left|(\mathrm{Mx}_1)_{\beta\beta'}\right|\cdot\left|(\mathrm{Mx}_2)_{\beta'\beta}\right|
\le4\,|\mathrm{Mx}_1|_c\cdot|\mathrm{Mx}_2|_c .
\]

\emph{The case $k\ge1$.}
The matrix $\gamma_5$ commutes with $P(B)$, whose kernel consists of multiples of $\mathbf{1}_4$, and with the matrices $\left(\mathrm{bv}+i\gamma_5y\right)^{-1}$ for $\mathrm{bv}=b_k$ and $\mathrm{bv}=b_k+b_L$.
From (27.2), $P(B)^2=P(B)$ and $P(B)\cdot\left(1-P(B)\right)=0$ one obtains
\[
O_1=B_{k,y}(B)\cdot\gamma_5\cdot B_{k,y}(B)
=\gamma_5\cdot\left[\left(b_k+i\gamma_5y\right)^{-2}\cdot\left(1-P(B)\right)+\left(b_k+b_L+i\gamma_5y\right)^{-2}\cdot P(B)\right].
\]
By (b.1),
\[
\left(\mathrm{bv}+i\gamma_5y\right)^{-1}=\frac{\mathrm{bv}-i\gamma_5y}{\mathrm{bv}^2+y^2}.
\]
Every row of $\gamma_5$ has exactly one non-zero entry; this entry equals $1$ and lies off the diagonal (Section~I.2.1).
So every row of $\mathrm{bv}-i\gamma_5y$ has the two non-zero entries $\mathrm{bv}$ and $-iy$, and
\[
\left|\left(\mathrm{bv}+i\gamma_5y\right)^{-1}\right|_c=\frac{\mathrm{bv}+y}{\mathrm{bv}^2+y^2}\le\sqrt{2}\cdot\left(\mathrm{bv}^2+y^2\right)^{-\frac{1}{2}},\qquad|\gamma_5|_c=1.
\]
Therefore
\[
\left|\gamma_5\cdot\left(\mathrm{bv}+i\gamma_5y\right)^{-2}\right|_c\le\frac{2}{\mathrm{bv}^2+y^2}\le\frac{2}{b_k^2+y^2}.
\]
The kernel of $1-P(B)$ at $(z_2,z_1)$ has the norm $|\cdot|_c$ at most $1+|P(B)(z_2,z_1)|_c$.
It follows that
\[
\left|O_1(z_2,z_1)\right|_c\le\frac{2\left(1+2\,|P(B)(z_2,z_1)|_c\right)}{b_k^2+y^2},
\]
where $P(B)$ does not depend on $y$.

Let $v\in\mathfrak{S}(1,\ell)$ vanish at all unit sites outside $z_2+\ell\mathbb{Z}^3$.
Then $(O_2v)(z_1)=O_2(z_1,z_2)\cdot v(z_2)$ and $|v|_\infty=|v(z_2)|_c$, and $v(z_2)$ may be any element of $\mathbb{C}^4$.
The spinor $f:=J\cdot Q_k^T(-B)\,v$ vanishes at every fermion site that does not lie over the block $\Delta_{z_2}$ of $T_k$ (Section~I.3.3).
Thus $f=\mathbf{1}_{\Delta_{z_2}}f$, and
\[
|f|_\infty\le\exp\left(\frac{3}{2}\,\mathfrak{i}(B)\right)\cdot|v(z_2)|_c
\]
by (I.Q3).
For every $f'\in\mathfrak{S}^{\mathrm{ch}}(\eta_k,\ell)$,
\[
\left(Q_k(B)\cdot J^{*}f'\right)(z_1)=\int_{\Delta_{z_1}}\mathrm{d}x\,\varphi_B(z_1,c)\cdot f'(x)=\left(Q_k(B)\cdot J^{*}\,\mathbf{1}_{\Delta_{z_1}}f'\right)(z_1),
\]
and by (I.Q3) the norm $|\cdot|_c$ of this vector is at most $\exp\left(\frac{3}{2}\,\mathfrak{i}(B)\right)\cdot\left|\mathbf{1}_{\Delta_{z_1}}f'\right|_\infty$.
Choosing $f'=S_{k,y}(X,B)\,\mathbf{1}_{\Delta_{z_2}}f$ and using W7$^{\mathrm{ch}}$, one finds
\[
\left|O_2(z_1,z_2)\cdot v(z_2)\right|_c\le\exp\left(3\,\mathfrak{i}(B)\right)\cdot C_{w}\cdot\exp\left(-\gamma_{\mathrm{w}}\,d(z_1,z_2)-\kappa_{\mathrm{w}}\,d_M(X)\right)\cdot|v(z_2)|_c .
\]
Since $\gamma_{\mathrm{w}}>0$ and $\kappa_{\mathrm{w}}>0$, this gives, for all $y\ge0$,
\[
\left|O_2(z_1,z_2)\right|_c\le\exp\left(3\,\mathfrak{i}(B)\right)\cdot C_{w}\cdot\exp\left(-\gamma_{\mathrm{w}}\,d(z_1,z_2)-\kappa_{\mathrm{w}}\,d_M(X)\right)\le\exp\left(3\,\mathfrak{i}(B)\right)\cdot C_{w}.
\]
By the last display of (b.2), the modulus of the trace is bounded by
\[
8\,C_{w}\cdot\exp\left(3\,\mathfrak{i}(B)\right)\cdot\left[\sum_{z_1,z_2}\left(1+2\,|P(B)(z_2,z_1)|_c\right)\right]\cdot\left(b_k^2+y^2\right)^{-1},
\]
the sum being over the unit sites $z_1$, $z_2$ of $X$.
The factor in front of $\left(b_k^2+y^2\right)^{-1}$ is a finite number that does not depend on $y$.
The entries of $O_1(z_2,z_1)$ are continuous functions of $y$.
By (27.3), the entries of $O_2(z_1,z_2)$ are finite linear combinations, with coefficients independent of $y$, of entries of the kernel of $S_{k,y}(X,B)$.
Hence the trace is a Borel function of $y$.
As $b_k>0$ by (b.1) and
\[
\int_0^{\infty}\frac{\mathrm{d}y}{b_k^2+y^2}=\frac{\pi}{2\,b_k},
\]
the integral in $\acute{E}_k(X,B)$ converges absolutely.

\emph{The case $k=0$.}
W7$^{\mathrm{ch}}$ in its form for $k=0$ gives
\[
\left|S_{0,y}(X,B)(x,x')\right|_c\le C_{w}\cdot(1+y)^{-1}
\]
for all fermion sites $x$, $x'$ and all $y\ge0$.
The kernel of $F_0(B)$ does not depend on $y$.
By (b.4),
\begin{align*}
\left|\operatorname{Tr}\left[F_0(B)\cdot S_{0,y}(X,B)\right]\right|&\le4\,C_{w}\cdot(1+y)^{-1}\int_X\mathrm{d}x\int_X\mathrm{d}x'\,\left|F_0(B)(x,x')\right|_c,\\
\left|\operatorname{Tr}\left[\gamma_5\cdot S_{0,y}(X,B)\right]\right|&\le4\,C_{w}\cdot(1+y)^{-1}\cdot\operatorname{Vol}(X);
\end{align*}
here the double integral is the finite sum $\left(\frac{1}{48}\right)^2\sum_{c,c'\in X}\sum_{g,g'\in O_{\mathrm{cub}}}$, and $\int_X\mathrm{d}x=\operatorname{Vol}(X)$.
Each of the two traces is a finite linear combination, with coefficients independent of $y$, of entries of the kernel of $S_{0,y}(X,B)$, and is therefore a Borel function of $y$.
Since
\[
\int_1^{\infty}\frac{\mathrm{d}y}{y\,(1+y)}=\log2,\qquad\int_0^1\frac{\mathrm{d}y}{1+y}=\log2,
\]
the two integrals in $\acute{E}_0(X,B)$ converge absolutely.

\emph{The equalities.}
By Section~I.5.7, $\acute{E}_k^{(i)}(X_i,B_i)$ arises from the trace of (b.3) by the operation $-i\,b_k^2\int_0^{\infty}\mathrm{d}y\,(\,\cdot\,)$ if $k\ge1$, and from the traces of (b.5) by the operation $\int_1^{\infty}\frac{\mathrm{d}y}{y}\,(\,\cdot\,)-i\int_0^1\mathrm{d}y\,(\,\cdot\,)$ if $k=0$; the number $b_k$ is the same on both tori.
By (b.3) and (b.5), the integrands on $T_1$ and on $T_2$ are one and the same function of $y$.
Therefore
\[
\acute{E}_k^{(1)}(X_1,B_1)=\acute{E}_k^{(2)}(X_2,B_2).
\]
The connections $-B_i$, $\overline{B_i}$ and $-\overline{B_i}$ lie in $\tilde{G}_k(T_i)$ by \cite[Lemma~17.6]{QEDI}, and each of the pairs $(-B_1,-B_2)$, $(\overline{B_1},\overline{B_2})$, $(-\overline{B_1},-\overline{B_2})$ has equal lifts on the bonds that meet $\tilde{X}$.
So the equality just proved holds for these pairs as well.
By Section~I.5.7, $E^{d}_k(X,B)$ is the mean of the four numbers $\acute{E}_k(X,B)$, $\acute{E}_k(X,-B)$, $\overline{\acute{E}_k(X,\overline{B})}$ and $\overline{\acute{E}_k(X,-\overline{B})}$; hence
\[
E^{d,(1)}_k(X_1,B_1)=E^{d,(2)}_k(X_2,B_2).
\]
By (TT.3), this applies to $(B_1,B_2)=(\mathscr{T}_1,\mathscr{T}_2)$, to $(\mathscr{C}_1,\mathscr{C}_2)$ and to $(0,0)$.
By Section~I.5.7,
\begin{align*}
{\delta E^{\mathrm{det}}}_k^{(i)}(X_i;A_i,\mathscr{C}_i)&=E^{d,(i)}_k(X_i,\mathscr{T}_i)-E^{d,(i)}_k(X_i,\mathscr{C}_i),\\
E_k^{\mathrm{det},(i)}(X_i,\mathscr{C}_i)&=E^{d,(i)}_k(X_i,\mathscr{C}_i)-E^{d,(i)}_k(X_i,0),
\end{align*}
and both quantities take the same value for $i=1$ and for $i=2$.

\emph{In the reading (o).}
No activity occurs in Step (b).
The assertion on the integrands, the three trace identities and the equalities of $\acute{E}_k$ and of $E^{d}_k$ do not involve the second slot, and they hold without change.
The equality of $E^{d}_k$, applied to $(B_1,B_2)=(A_1,A_2)$ and to $(B_1,B_2)=(0,0)$, yields
\[
E_k^{\mathrm{det},(1)}(X_1,A_1)=E_k^{\mathrm{det},(2)}(X_2,A_2).
\]

\paragraph{Step (c): preliminaries}

Fix $\tilde{\mathbf{Y}}$ connected in $\tilde{Y}$ and data $\mathrm{zval}$, $\mathrm{sval}$ as in Proposition~27.5, and write
\[
s^{(i)}=s^{(i)}[\tilde{\mathbf{Y}};\mathrm{sval}],\qquad
\mathsf{Z}^{(i)}=\mathsf{Z}^{(i)}[\tilde{\mathbf{Y}};\mathrm{zval}],\qquad
\mathscr{Z}^{(i)}:=\mathscr{Z}^{(i)}(s'^{(i)}).
\]

(c0.1) The set $\tilde{\mathbf{Y}}$ is admissible, and $N(\tilde{\mathbf{Y}})$ is contained in $N(\tilde{Y})$ (Proposition~27.1\,(a)).
As $\tilde{\mathbf{Y}}$ is contained in $\mathrm{box}(\tilde{\mathbf{Y}})$, the set $\tilde{\mathbf{Y}}^{\sim a}$ is contained in $\mathrm{box}(\tilde{\mathbf{Y}})^{\sim a}$, and hence in $N(\tilde{\mathbf{Y}})$, for $0\le a\le6d_U$.
Note that $c_l=2L=d_U-1$.
Proposition~27.3 is applied below with the admissible set $\tilde{\mathbf{Y}}$ in the place of the set $\tilde{Y}$ of its statement, so that $Y_i=\mathbf{Y}_i$; Proposition~27.4 is applied with the given sets $\tilde{Y}$ and $\tilde{\mathbf{Y}}$.

(c0.2) The map $\mathrm{pr}_i$ is injective on the indices contained in $N(\tilde{\mathbf{Y}})$ (Proposition~27.1\,(b)), and parameter cubes and basis elements are indices (Proposition~27.1, Definitions).
Hence $s^{(i)}$ and $\mathsf{Z}^{(i)}$ are well defined, and $s^{(i)}\in\mathrm{Pd}_k(T_i)$.
The family $s^{(i)}$ vanishes at every parameter cube of $T_i$ that is not contained in $\mathbf{Y}_i$, since the cubes $\mathrm{pr}_i(\tilde{\mathrm{pc}})$ with $\tilde{\mathrm{pc}}$ contained in $\tilde{\mathbf{Y}}$ are contained in $\mathbf{Y}_i$.
Thus $s^{(i)}$ vanishes off $\mathbf{Y}_i$ in the sense of Proposition~26.4 (Definitions), and it is a family as in the first sentence of Proposition~24.3\,(ga), with $\mathbf{Y}=\mathbf{Y}_i$.
The family $s^{(2)}$ is the one that Proposition~27.3 associates with $s^{(1)}$: for every parameter cube $\tilde{\mathrm{pc}}$ contained in $\tilde{\mathbf{Y}}$ its value at $\mathrm{pr}_2(\tilde{\mathrm{pc}})$ is $s^{(1)}_{\mathrm{pr}_1(\tilde{\mathrm{pc}})}=\mathrm{sval}_{\tilde{\mathrm{pc}}}$, and it vanishes at all other cubes.

(c0.3) One has $\mathsf{Z}^{(i)}\in\mathrm{Box}_{k}(T_i)$, and $\mathsf{Z}^{(1)}$, $\mathsf{Z}^{(2)}$ have equal lifts on $N(\tilde{\mathbf{Y}})$ in the sense of Proposition~27.2\,(D1).
Indeed, let $\Upsilon$ be a basis element contained in $N(\tilde{\mathbf{Y}})$.
If $\Upsilon\in\mathrm{Bs}(\tilde{\mathbf{Y}})$, both coordinates equal $\mathrm{zval}_{\Upsilon}$.
Otherwise, by the injectivity of $\mathrm{pr}_i$, the element $\mathrm{pr}_i(\Upsilon)$ differs from $\mathrm{pr}_i(\Upsilon'')$ for every $\Upsilon''\in\mathrm{Bs}(\tilde{\mathbf{Y}})$, and both coordinates vanish.

\emph{In the reading (o)} nothing changes, since neither a connection nor an activity occurs.

\paragraph{Step (c1): the field $\mathscr{Z}$}

Let $b$ be a positively oriented fine bond contained in $\tilde{\mathbf{Y}}^{\sim(1+d_U)}$; this is no restriction, because both sides of the identity in Proposition~27.5\,(c1) change sign when $b$ is reversed.
On $T_i$ write $C^{(i)}(b',\Upsilon):=(C\Upsilon)(b')$.
By Definition~24.D\,(D2) and the kernels of Section~G4.0.4,
\begin{equation}
\mathscr{Z}^{(i)}(\mathrm{pr}_ib)=\sum_{b'}\sum_{\Upsilon}\sum_{\Upsilon'}H_{k}^{(i)}(s'^{(i)})(b,b')\cdot C^{(i)}(b',\Upsilon)\cdot C_{k}^{1/2,(i)}(s'^{(i)})(\Upsilon,\Upsilon')\cdot\mathsf{Z}^{(i)}(\Upsilon'),\tag{27.4}
\end{equation}
where $b'$ ranges over the positively oriented unit bonds and $\Upsilon$, $\Upsilon'$ over the basis elements, each over a set of representatives modulo $\ell_i\mathbb{Z}^3$.
All factors are $\ell_i\mathbb{Z}^3$-periodic.

(1) \emph{Rows of $H_k(s')$.}
Put $j:=1+d_U$; then $j\le5d_U$ and $j+c_l=2d_U$.
The bond $\mathrm{pr}_1(b)$ is a row index of (p1) on $T_1$ that is contained in $\mathbf{Y}_1^{\sim j}=\mathrm{pr}_1(\tilde{\mathbf{Y}}^{\sim j})$, and $b$ is its lift contained in $\tilde{\mathbf{Y}}^{\sim j}$ (Proposition~27.1\,(d)).
By Proposition~27.3\,(a) for $i=1$ and Proposition~27.3\,(b) for $i=2$, one has $H_{k}^{(i)}(s'^{(i)})(b,b')=0$ unless $\mathrm{pr}_i(b')$ is contained in $\mathbf{Y}_i^{\sim2d_U}=\mathrm{pr}_i(\tilde{\mathbf{Y}}^{\sim2d_U})$, and
\[
H_{k}^{(1)}(s'^{(1)})(b,b')=H_{k}^{(2)}(s'^{(2)})(b,b')\qquad\text{for every }b'\text{ contained in }N(\tilde{\mathbf{Y}}).
\]
By (Rs) with $\tilde{S}=\tilde{\mathbf{Y}}^{\sim2d_U}$, the sum over $b'$ in (27.4) extends over the $b'$ contained in $\tilde{\mathbf{Y}}^{\sim2d_U}$, which is a subset of $\tilde{\mathbf{Y}}^{\sim(1+2d_U)}$.

(2) \emph{Rows of $C$.}
Let $b'$ be a unit bond contained in $\tilde{\mathbf{Y}}^{\sim(1+2d_U)}$.
By \cite[Proposition~15.21(iv)]{QEDI} (the formula for the value of $C\mathsf{Z}$ at a central bond, together with $C\mathsf{Z}=\mathsf{Z}$ off the central bonds), the number $(C\mathsf{Z})(b')$ depends on the coordinates $\mathsf{Z}(\Upsilon)$ only for the $\Upsilon$ in the $L$-blocks of the two ends of $b'$ and on their common face.
On $T_i$ this means that $C^{(i)}(b',\Upsilon)=0$ unless either $\mathrm{pr}_i(\Upsilon)$ is one of the elements $e^{y}_{n}$, with $y$ the centre of the $L$-block of an end of $\mathrm{pr}_i(b')$, or $\mathrm{pr}_i(\Upsilon)=\delta_{b''}$ for a face bond $b''$ joining these two $L$-blocks.
In both cases $\operatorname{supp}\mathrm{pr}_i(\Upsilon)$ is contained in the union of the $L$-blocks of the ends of $\mathrm{pr}_i(b')$ (\cite[Proposition~15.21(ii)]{QEDI}).
By (g0) each of these $L$-blocks lies in a single $M$-cube, namely in the $M$-cube of the corresponding end.
So $\mathrm{pr}_i(\Upsilon)$ is contained in the union of the $M$-cubes of the ends of $\mathrm{pr}_i(b')$.
These are the images of the $M$-cubes of the ends of $b'$, which are cubes of $\tilde{\mathbf{Y}}^{\sim(1+2d_U)}$.
By (Rs), with $\tilde{S}$ the union of the $M$-cubes of the ends of $b'$, the sum over $\Upsilon$ extends over the $\Upsilon$ contained in this union, a subset of $\tilde{\mathbf{Y}}^{\sim(1+2d_U)}$.
For these $\Upsilon$ one has $C^{(1)}(b',\Upsilon)=C^{(2)}(b',\Upsilon)$ (Proposition~27.4\,(a-1)).

(3) \emph{Rows of $C_{k}^{1/2}(s')$.}
Let $\Upsilon$ be contained in $\tilde{\mathbf{Y}}^{\sim j'}$, where $j':=1+2d_U\le5d_U$; then $j'+c_l=3d_U$.
By Proposition~27.3\,(a), (b) for the family (p2), $C_{k}^{1/2,(i)}(s'^{(i)})(\Upsilon,\Upsilon')=0$ unless $\mathrm{pr}_i(\Upsilon')$ is contained in $\mathrm{pr}_i(\tilde{\mathbf{Y}}^{\sim3d_U})$, and the entries for $i=1$ and $i=2$ coincide for every $\Upsilon'$ contained in $N(\tilde{\mathbf{Y}})$.
By (Rs), the sum over $\Upsilon'$ extends over the $\Upsilon'$ contained in $\tilde{\mathbf{Y}}^{\sim3d_U}$.
This set is a subset of $\tilde{\mathbf{Y}}^{\sim(1+3d_U)}$, which lies in $N(\tilde{\mathbf{Y}})$ by (c0.1), because $1+3d_U\le6d_U$.

(4) For these $\Upsilon'$ one has $\mathsf{Z}^{(1)}(\Upsilon')=\mathsf{Z}^{(2)}(\Upsilon')$ by (c0.3).

Hence the two sums (27.4) extend over the same covering indices and have the same terms, and
\[
\mathscr{Z}^{(1)}(\mathrm{pr}_1b)=\mathscr{Z}^{(2)}(\mathrm{pr}_2b).
\]

\emph{Level $0$.}
Here (p1) is the diagonal family
\[
H_0(s')(b,b')=\mathbf{1}[b=b']\cdot\prod\left\{s'_{\mathrm{pc}}:\mathrm{pc}\in\pi'_0,\ b\text{ contained in }\mathrm{pc}\right\}
\]
of Section~I.5.8.
It belongs to the system on $T_2$, and its projection is the diagonal family on $T_1$ (Section~G4.0.4; Proposition~26.3\,(c)).
Assertion (1) holds for it as a special case of Proposition~27.3, and the sum over $b'$ reduces to the single term $b'=b$.

\emph{In the reading (o)} nothing changes, as only photon objects occur.

\paragraph{Step (c2): connections}

By (OS) and the first clause of condition (A8), as in (TT.3), one has $e_k=L^{-\frac{1}{2}}e_{k+1}<e_0<1$; here and below (A1), (A5), (A8) are the conditions of Section~I.8.4, that is, of \cite[Section~5.6.2]{QEDI} with the constants of Section~I.8.2.
By the third clause of (A1),
\[
\frac{\mathrm{rad}_2}{\mathrm{rad}_1}=2e_k^{\frac{5\varepsilon}{2}}\le1 .
\]
By (A5), $e_k^{-\left(\frac{1}{4}-6\varepsilon\right)}\ge88M^3\ge2$.
Since $\varepsilon\le\frac{1}{200}$, the exponents satisfy
\[
-\frac{3}{4}+7\varepsilon\le-\frac{1}{4}+3\varepsilon\le-\frac{1}{4}+6\varepsilon ,
\]
and therefore $\mathrm{rad}_2\ge\mathrm{rad}_4\ge e_k^{-\frac{1}{4}+6\varepsilon}\ge2$.
(These are the inequalities of the first assertion of Proposition~24.3\,(q1).)

One has $s^{(i)}\in\mathrm{Pd}_k(T_i)$ and $\mathsf{Z}^{(i)}\in\mathrm{Box}_{k}(T_i)$ by (c0.2), (c0.3), and $A_i\in\tilde{R}_k^{-}(T_i)$, $\mathscr{C}_i\in\mathfrak{X}_k(T_i)$ by (TT.3).
Proposition~24.3\,(q1) on $T_i$ shows that the three norms of $\mathscr{Z}^{(i)}$ are at most $\frac{1}{2}e_k^{-\frac{\varepsilon}{2}}$, and that for $|t|\le\mathrm{rad}_1$ the field $A_i+t\cdot\mathscr{Z}^{(i)}$ lies in $\tilde{R}_k(T_i)$, while the field
\[
B_i(t):=\mathscr{T}_i+t\cdot\mathscr{Z}^{(i)}
\]
lies in $(\tilde{R}_k+\mathfrak{X}_k)(T_i)$, a subset of $\tilde{G}_k(T_i)$.
On the bonds contained in $\tilde{\mathbf{Y}}^{\sim(1+d_U)}$ the connections $A_1$, $A_2$ have equal lifts, and so have $\mathscr{T}_1$, $\mathscr{T}_2$, by (TT.3); the fields $\mathscr{Z}^{(1)}$, $\mathscr{Z}^{(2)}$ have equal lifts there by (c1).
Consequently $A_1+t\cdot\mathscr{Z}^{(1)}$ and $A_2+t\cdot\mathscr{Z}^{(2)}$ have equal lifts on these bonds, and so have $B_1(t)$ and $B_2(t)$.
A bond that meets $\tilde{\mathbf{Y}}$ is contained in $\tilde{\mathbf{Y}}^{\sim1}$, since $\eta_k<M$, and hence in $\tilde{\mathbf{Y}}^{\sim(1+d_U)}$.

\emph{In the reading (o)}, $B_i(t)=A_i+t\cdot\mathscr{Z}^{(i)}$ lies in $\tilde{R}_k(T_i)$, which is a subset of $\tilde{G}_k(T_i)$ (Proposition~24.3\,(q1) in the reading (o), and \cite[Lemma~17.4(a)]{QEDI}).

\paragraph{Step (c3): rows of the fermionic primitives}

Let $|t|\le\mathrm{rad}_1$ and $B_i:=B_i(t)$; the value $t=0$ gives $B_i=\mathscr{T}_i$.

(1) \emph{The families $H_k$ and $\bar{H}_k$.}
Let $x=(c,g)$ be a fermion site with $\operatorname{dist}(c,\tilde{\mathbf{Y}})\le\eta_k$.
Then $c$ lies in $\tilde{\mathbf{Y}}^{\sim1}$, the site $\mathrm{pr}_1(x)$ is a row index of (p3) on $T_1$ contained in $\mathbf{Y}_1^{\sim1}=\mathrm{pr}_1(\tilde{\mathbf{Y}}^{\sim1})$, and $x$ is its lift contained in $\tilde{\mathbf{Y}}^{\sim1}$ (Proposition~27.1\,(d)).
By (c2), $B_1\in\tilde{G}_k(T_1)$ and $B_2\in\tilde{G}_k(T_2)$ have equal lifts on the bonds of $\tilde{\mathbf{Y}}^{\sim(1+d_U)}$.
Proposition~27.3\,(b) with $j=1$ gives, for every column index $u$ contained in $N(\tilde{\mathbf{Y}})$ (a unit site if $k\ge1$, a fermion site if $k=0$),
\[
H_k^{(1)}(s''^{(1)},B_1)(x,u)=H_k^{(2)}(s''^{(2)},B_2)(x,u),\qquad
\bar{H}_k^{(1)}(s''^{(1)},B_1)(x,u)=\bar{H}_k^{(2)}(s''^{(2)},B_2)(x,u);
\]
for $\bar{H}_k$ this rests on the last sentence of Proposition~27.3.

(2) \emph{The family $\Gamma_k$.}
Let $u$ be a unit site of $\tilde{\mathbf{Y}}$ (for $k=0$: a fermion site over $\tilde{\mathbf{Y}}$).
Proposition~27.3\,(b) with $j=0$, applied to the connections $\mathscr{T}_1$, $\mathscr{T}_2$, which have equal lifts on all bonds, gives
\[
\Gamma_k^{(1)}(s''^{(1)},\mathscr{T}_1)(u,u')=\Gamma_k^{(2)}(s''^{(2)},\mathscr{T}_2)(u,u')
\qquad\text{for every }u'\text{ contained in }N(\tilde{\mathbf{Y}}).
\]

(3) \emph{The family $M_k$, $k\ge1$.}
Let $\tilde{X}$ be connected in $\tilde{\mathbf{Y}}$.
The connections $B_1$, $B_2$ have equal lifts on the bonds that meet $\tilde{X}$.
By Proposition~27.3\,(c),
\[
M_k^{(1)}(\mathrm{pr}_1\tilde{X},B_1)(y,y')=M_k^{(2)}(\mathrm{pr}_2\tilde{X},B_2)(y,y')
\]
for the unit sites $y$, $y'$ of $\tilde{X}$.
If $y$, $y'$ are unit sites of $N(\tilde{\mathbf{Y}})$ and one of them does not lie in $\tilde{X}$, then its image does not lie in $\mathrm{pr}_i(\tilde{X})$, by (Mb), and both kernels vanish, by W12$^{\mathrm{ch}}$ on $T_1$ and on $T_2$.

(4) \emph{Columns.}
By (c0.2), Proposition~24.3\,(ga2) applies on $T_i$ with $\mathbf{Y}=\mathbf{Y}_i$ and $s=s^{(i)}$.
Accordingly, the rows of $H_k^{(i)}(s''^{(i)},B_i)$ and $\bar{H}_k^{(i)}(s''^{(i)},B_i)$ at every fermion site, and the rows of $\Gamma_k^{(i)}(s''^{(i)},\mathscr{T}_i)$ at the unit indices in $\mathbf{Y}_i$ (over $\mathbf{Y}_i$ if $k=0$), have their columns in $\mathbf{Y}_i$ (over $\mathbf{Y}_i$ if $k=0$); moreover the rows of $H_k^{(i)}$ and $\bar{H}_k^{(i)}$ at $x$ vanish unless $\mathrm{pr}_i(c)$ lies in $\mathbf{Y}_i$.
An index of $T_i$ contained in $\mathbf{Y}_i=\mathrm{pr}_i(\tilde{\mathbf{Y}})$ has exactly one lift contained in $\tilde{\mathbf{Y}}$ (Proposition~27.1\,(d)).
Therefore these rows vanish at the columns that have no lift in $\tilde{\mathbf{Y}}$.

\emph{Level $0$.}
The operators $H_0^{(i)}(s''^{(i)},B)$ and $\bar{H}_0^{(i)}(s''^{(i)},B)=\mathsf{C}^{-1}\cdot H_0^{(i)}(s''^{(i)},-B)\cdot\mathsf{C}$ are the multiplication by the number
\[
\mathrm{dgf}_i(c):=\prod\left\{s''^{(i)}_{\mathrm{pc}}:\mathrm{pc}\in\pi_0(T_i),\ \mathrm{pc}\text{ meets }K_c\right\},
\]
which is the same at the $48$ sites over $c$ (Section~I.5.8).
By (I.ke1), their kernel (I.1.4) is $48\cdot\mathbf{1}[x=x']\cdot\mathrm{dgf}_i(c)\cdot\mathbf{1}_4$.
Assertion (1) states that $\mathrm{dgf}_1(c)=\mathrm{dgf}_2(c)$ whenever $\operatorname{dist}(c,\tilde{\mathbf{Y}})\le1$.
Assertion (4) states that $\mathrm{dgf}_i(c)=0$ unless $\mathrm{pr}_i(c)$ lies in $\mathbf{Y}_i$, that is, for a site $c$ of $N(\tilde{\mathbf{Y}})$, unless $c$ lies in $\tilde{\mathbf{Y}}$ (by (Mb)).
The kernels of $\Gamma_0^{(i)}$ are those of (I.1.4) at pairs of fermion sites.
There is no family $M_0$.

\emph{In the reading (o)} nothing changes; here $\mathscr{T}_i=A_i$.

\paragraph{Step (c4): data.}
Let $\tilde{X}$ be connected in $\tilde{\mathbf{Y}}$, and let $|t|\le\mathrm{rad}_1$.
By (c2) the pairs $\left(A_i+t\cdot\mathscr{Z}^{(i)},\mathscr{C}_i\right)$ belong to $\tilde{R}_k(T_i)\times\mathfrak{X}_k(T_i)$, and they have equal lifts on the bonds that meet $\tilde{X}$ (such bonds meet $\tilde{\mathbf{Y}}$).
The set $\tilde{X}$ is connected and admissible, and by (a) one has $R\,\mathscr{E}^{(1)}\sim R\,\mathscr{E}^{(2)}$ as local families on $\tilde{R}_k\times\mathfrak{X}_k$.
The second sentence of Proposition~27.2\,(c), combined with (Co4), therefore yields
\[
\mathrm{idn}\,\mathscr{E}^{r,(1)}\left(\mathrm{pr}_1\tilde{X};A_1+t\cdot\mathscr{Z}^{(1)},\mathscr{C}_1\right)=\mathscr{E}^{r,(2)}\left(\mathrm{pr}_2\tilde{X};A_2+t\cdot\mathscr{Z}^{(2)},\mathscr{C}_2\right).
\]
Both elements belong to $\mathrm{Grq}_i(\tilde{X})$, by the property (k2) of Section~I.5.4 and by (Gq3).
For $k=0$ both of them vanish.

\emph{Reading (o).}
By (a) in the reading (o) and by (o-4),
\[
\mathrm{idn}\,E^{r,(1)}\left(\mathrm{pr}_1\tilde{X};A_1+t\cdot\mathscr{Z}^{(1)}\right)=E^{r,(2)}\left(\mathrm{pr}_2\tilde{X};A_2+t\cdot\mathscr{Z}^{(2)}\right).
\]

\paragraph{Step (c5): first localisation, the functions $G^{(1)}$, $G^{(2)}$, $G^{(3)}$.}
In this step $\tilde{X}$ is a set connected in $\tilde{\mathbf{Y}}$.
For a connection $B$ on $T_i$ write
\[
H^{(i)}[B]:=H_k^{(i)}\left(s''^{(i)},B\right),\qquad\bar{H}^{(i)}[B]:=\bar{H}_k^{(i)}\left(s''^{(i)},B\right).
\]

\emph{Levels $k\ge1$.}

(c5.1) \emph{The fields.}
Let $x=(c,g)$ be a fermion site with $\mathrm{dist}(c,\tilde{\mathbf{Y}})\le\eta_k$, and let $|t|\le\mathrm{rad}_1$.
With the kernels taken as in Section~G4.0.4, part~(4) of Step~(c3) and (Rs) with $\tilde{S}=\tilde{\mathbf{Y}}$ give
\begin{equation}
\begin{aligned}
\left(H^{(i)}[B_i(t)]\,\Psi\right)_{\beta}(x)&=\sum_{y\in\tilde{\mathbf{Y}}\cap\mathbb{Z}^3}\ \sum_{\beta'}H^{(i)}[B_i(t)](x,y)_{\beta\beta'}\cdot\Psi(y,\beta',0),\\
\left(\bar{H}^{(i)}[B_i(t)]\,\bar{\Psi}\right)_{\beta}(x)&=\sum_{y\in\tilde{\mathbf{Y}}\cap\mathbb{Z}^3}\ \sum_{\beta'}\bar{H}^{(i)}[B_i(t)](x,y)_{\beta\beta'}\cdot\Psi(y,\beta',1).
\end{aligned}
\tag{27.5}
\end{equation}
These are odd elements of $\mathrm{Gr}_i(\tilde{\mathbf{Y}})$.
By part~(1) of Step~(c3) their coefficients do not depend on $i$; hence, by (Co1), $\mathrm{idn}$ maps the elements (27.5) on $T_1$ to the elements (27.5) on $T_2$.

Now let $(x,x')$ be a pair of $\mathrm{Pad}(\tilde{X})$.
The photon sites of $x$ and of $x'$ are at sup-distance at most $\eta_k$ from $\tilde{X}$ (Section~0.4), hence from $\tilde{\mathbf{Y}}$.
By (I.1.7) and (I.1.6), the latter with $U_{-B}=U_B^{-1}$, one has at $B=B_i(t)$
\begin{equation}
\begin{aligned}
\left(\delta_{\alpha,B}H^{(i)}[B]\,\Psi\right)(x,x')&=d_X(x,x')^{-\alpha}\cdot\left(U_B(x,x')\left(H^{(i)}[B]\,\Psi\right)(x')-\left(H^{(i)}[B]\,\Psi\right)(x)\right),\\
\left(\delta_{\alpha,-B}\bar{H}^{(i)}[B]\,\bar{\Psi}\right)(x,x')&=d_X(x,x')^{-\alpha}\cdot\left(U_B(x,x')^{-1}\left(\bar{H}^{(i)}[B]\,\bar{\Psi}\right)(x')-\left(\bar{H}^{(i)}[B]\,\bar{\Psi}\right)(x)\right).
\end{aligned}
\tag{27.6}
\end{equation}
For a pair inside a cluster ($c=c'$, $g\ne g'$) one has $d_X=\eta_k$ and $U_B=1$; for a pair $(x,x^{(3)})$ one has $d_X=\eta_k$ and $U_B(x,x^{(3)})=u_B(c,c_3)$.
The connection $B_i(t)$ lies in $\tilde{G}_k(T_i)$, and by (c2) the connections $B_1(t)$, $B_2(t)$ have equal lifts on the bonds that meet $\tilde{\mathbf{Y}}$.
By (Gq0) the pairs of $\mathrm{Pad}(\tilde{X})$ are the lifts of the pairs that are admissible for $\mathrm{pr}_i(\tilde{X})$.
Proposition~27.4\,(a-2) therefore asserts that the rows of $\delta_{\alpha,B_i(t)}$ and of $\delta_{\alpha,-B_i(t)}$ at these pairs are the same for $i=1$ and $i=2$; in other words, the numbers $d_X(x,x')$ and $U_{B_i(t)}(x,x')^{\pm1}$ that occur in (27.6) do not depend on $i$.
(Here $d_X$ is a distance of the covering lattice, (I.1.5), and
\[
U_B(x,x')=\exp\left(i\,e_k\cdot\eta_k\cdot\frac{1}{6}\sum_{\pi}B\left(\boldsymbol{\Gamma}^{\pi}(c,c')\right)\right);
\]
for a pair of kind (I) the bonds of the six paths are contained in $\tilde{X}$, and for a pair of kind (II) they reduce to the single bond $(c,c_3)$, which meets $\tilde{X}$.)
Hence $\mathrm{idn}$ maps the elements (27.6) on $T_1$ to those on $T_2$.

Since $1\le\mathrm{rad}_1$ by (c2), the identities (27.5) and (27.6) hold at $t=1$ and at $t=0$, where $B_i(1)=\mathscr{T}_i+\mathscr{Z}^{(i)}$ and $B_i(0)=\mathscr{T}_i$.
Put
\[
\mathrm{vJ}^{(i)}:=J_k^{\#,(i)}\left(s''^{(i)},\mathscr{T}_i,\mathscr{Z}^{(i)}\right)\Psi.
\]
By Definition~24.D\,(D2), the members of this family at the indices $(x,\beta,0)$, $(x,\beta,1)$ with $x$ over $\tilde{X}$, and $(x,x',\beta,0)$, $(x,x',\beta,1)$ with $(x,x')\in\mathrm{Pad}(\tilde{X})$, are the differences
\[
\bigl[\,(27.5),\ (27.6)\ \text{at}\ t=1\,\bigr]-\bigl[\,(27.5),\ (27.6)\ \text{at}\ t=0\,\bigr],
\]
taken in the component $\beta$, the first lines of the two displays for $\omega=0$ and the second lines for $\omega=1$.
Consequently the members of $\mathrm{vJ}^{(i)}$ at these indices are odd elements of $\mathrm{Gr}_i(\tilde{\mathbf{Y}})$, linear in the generators $\Psi(y,\beta',\omega)$, $y\in\tilde{\mathbf{Y}}\cap\mathbb{Z}^3$, and
\begin{equation}
\mathrm{idn}\,\mathrm{vJ}^{(1)}(\xi)=\mathrm{vJ}^{(2)}(\xi),\qquad\mathrm{idn}\,\mathrm{vJ}^{(1)}(\boldsymbol{\zeta})=\mathrm{vJ}^{(2)}(\boldsymbol{\zeta}).
\tag{27.7}
\end{equation}
The same holds for the families $t\cdot\mathrm{vJ}^{(i)}$, $t$ a complex number.

(c5.2) \emph{The function $G^{(1)}$.}
Let $|t|=\mathrm{rad}_1$.
By W12$^{\mathrm{ch}}$ (property~12 of \cite[Section~5.4.2]{QEDI}, read as in Section~I.5.8), the kernel of $M_k^{(i)}\left(\mathrm{pr}_i\tilde{X},B_i(t)\right)$ vanishes unless both of its unit sites lie in $\mathrm{pr}_i(\tilde{X})$.
So (Fm2) applies, and the integrand of $G^{(1),(i)}\left(s^{(i)},\mathrm{pr}_i\tilde{X}\right)$ is
\[
\left\langle\bar{\Psi},M_k^{(i)}\left(\mathrm{pr}_i\tilde{X},B_i(t)\right)\Psi\right\rangle_1=\sum_{y,y'\in\tilde{X}\cap\mathbb{Z}^3}\ \sum_{\beta,\beta'}M_k^{(i)}\left(\mathrm{pr}_i\tilde{X},B_i(t)\right)(y,y')_{\beta\beta'}\cdot\bar{\Psi}_{\beta}(y)\,\Psi_{\beta'}(y'),
\]
an element of $\mathrm{GrqP}_i(\tilde{X})$.
By part~(3) of Step~(c3) the coefficients are the same for $i=1,2$, and (Co1) gives $\mathrm{idn}\,(\text{integrand})^{(1)}=(\text{integrand})^{(2)}$.

(c5.3) \emph{The function $G^{(2)}$.}
Let $|t|=\mathrm{rad}_2$; then $|t|\le\mathrm{rad}_1$, by (c2).
By (c4) the element $\mathscr{E}^{r,(i)}\left(\mathrm{pr}_i\tilde{X};A_i+t\cdot\mathscr{Z}^{(i)},\mathscr{C}_i\right)$ belongs to $\mathrm{Grq}_i(\tilde{X})$.
By Section~0.5 the integrand of $G^{(2),(i)}\left(s^{(i)},\mathrm{pr}_i\tilde{X}\right)$ is
\[
\mathscr{E}^{r,(i)}\left(\mathrm{pr}_i\tilde{X};A_i+t\cdot\mathscr{Z}^{(i)},\mathscr{C}_i,\psi^{\#}+t\cdot\mathrm{vJ}^{(i)}\right)=\mathrm{adf}_i\left[t\cdot\mathrm{vJ}^{(i)}\right]\left(\mathscr{E}^{r,(i)}\left(\mathrm{pr}_i\tilde{X};A_i+t\cdot\mathscr{Z}^{(i)},\mathscr{C}_i\right)\right).
\]
It is a polynomial in the generators of $\mathrm{Grq}_i(\tilde{X})$ and in the $\Psi(y,\beta,\omega)$, $y\in\tilde{\mathbf{Y}}\cap\mathbb{Z}^3$; by (Gq2) it therefore belongs to $\mathrm{GrqP}_i(\tilde{\mathbf{Y}})$.
Applying (Ad) with (27.7), and then (c4), one obtains
\begin{align*}
&\mathrm{idn}\,\mathrm{adf}_1\left[t\cdot\mathrm{vJ}^{(1)}\right]\left(\mathscr{E}^{r,(1)}\left(\mathrm{pr}_1\tilde{X};A_1+t\cdot\mathscr{Z}^{(1)},\mathscr{C}_1\right)\right)\\
&\qquad=\mathrm{adf}_2\left[t\cdot\mathrm{vJ}^{(2)}\right]\left(\mathrm{idn}\,\mathscr{E}^{r,(1)}\left(\mathrm{pr}_1\tilde{X};A_1+t\cdot\mathscr{Z}^{(1)},\mathscr{C}_1\right)\right)\\
&\qquad=\mathrm{adf}_2\left[t\cdot\mathrm{vJ}^{(2)}\right]\left(\mathscr{E}^{r,(2)}\left(\mathrm{pr}_2\tilde{X};A_2+t\cdot\mathscr{Z}^{(2)},\mathscr{C}_2\right)\right).
\end{align*}

(c5.4) \emph{The function $G^{(3)}$.}
Let $\tilde{X}=\tilde{\square}$ be an $M$-cube of $\tilde{\mathbf{Y}}$, and let $|t|=\mathrm{rad}_2$.
Recall that $J_k$, $\bar{J}_k$ are the components of $J_k^{\#}$ without H\"older quotient.
By (Fm1) the element
\[
\langle\bar{\psi},\psi\rangle_{\mathrm{pr}_i\tilde{\square}}=\frac{\eta_k^3}{48}\sum_{c\in\tilde{\square}}\ \sum_{g\in O_{\mathrm{cub}}}\ \sum_{\beta}\bar{\psi}_{\beta}(c,g)\,\psi_{\beta}(c,g)
\]
belongs to $\mathrm{Grq}_i(\tilde{\square})$, and the integrand of $G^{(3),(i)}\left(s^{(i)},\mathrm{pr}_i\tilde{\square}\right)$ is
\[
(-m')\cdot\left\langle\bar{\psi}+t\cdot\bar{J}_k\bar{\Psi},\psi+t\cdot J_k\Psi\right\rangle_{\mathrm{pr}_i\tilde{\square}}=(-m')\cdot\mathrm{adf}_i\left[t\cdot\mathrm{vJ}^{(i)}\right]\left(\langle\bar{\psi},\psi\rangle_{\mathrm{pr}_i\tilde{\square}}\right),
\]
an element of $\mathrm{GrqP}_i(\tilde{\mathbf{Y}})$.
The number $m'$ is the same for $i=1,2$, by (a).
(Ad) with (27.7), together with (Fm1) (which in this case is Proposition~27.4\,(b)), gives $\mathrm{idn}\,(\text{integrand})^{(1)}=(\text{integrand})^{(2)}$.

(c5.5) \emph{The contour integrals.}
The functional of Section~I.5.2,
\[
\oint_{\mathrm{rad}}f=\frac{1}{2\pi i}\oint_{|t|=\mathrm{rad}}\mathrm{d}t\,\frac{f(t)}{t\cdot(t-1)},
\]
is an integral over the parameter $t$, with the same radius on both tori.
The integrals exist, by Proposition~24.1\,(a) on $T_i$.
By (Co3),
\[
\mathrm{idn}\,G^{(j),(1)}\left(s^{(1)},\mathrm{pr}_1\tilde{X}\right)=G^{(j),(2)}\left(s^{(2)},\mathrm{pr}_2\tilde{X}\right),\qquad j=1,2,3.
\]
Since $\mathrm{GrqP}_i(\tilde{\mathbf{Y}})$ is a linear subspace of finite dimension, $G^{(j),(i)}\left(s^{(i)},\mathrm{pr}_i\tilde{X}\right)$ belongs to $\mathrm{GrqP}_i(\tilde{\mathbf{Y}})$, a subalgebra of $\mathrm{Gr}_i(\tilde{\mathbf{Y}})$; thus it is a polynomial in generators $\psi$, $\chi$, $\Psi$ with indices in $N(\tilde{\mathbf{Y}})$.
On a polymer that is not an $M$-cube one has $G^{(3)}=0$ on both tori (Definition~24.D\,(D3)), and $\mathrm{pr}_i(\tilde{X})$ is an $M$-cube if and only if $\tilde{X}$ is one.

\emph{Level $k=0$.}

(c5.6)
By Definition~24.D\,(D4) and (D3), $G^{(2)}=G^{(3)}=0$, and $G^{(1)}(\,\cdot\,,X)=0$ unless $X$ is an $M$-cube.
Let $\tilde{\square}$ be an $M$-cube of $\tilde{\mathbf{Y}}$.
For $x=(c,g)$ write $c_3=c+g e_1$ and $x^{(3)}=(c_3,g s_3)$, and for spinors $v$, $v'$ and a complex $4\times4$ matrix $\mathrm{Mx}$ write $v\cdot\mathrm{Mx}\,v':=\sum_{\beta,\beta'}v_{\beta}\,\mathrm{Mx}_{\beta\beta'}\,v'_{\beta'}$.
Proposition~23.2, with the displays (23.3) and (8.3$_0$)$^{\mathrm{ch}}$, is combined with (Rs), applied to a summand that is $\ell_i\mathbb{Z}^3$-periodic in $x$.
Indeed, a translation $x\mapsto x+\ell_i e_{\mu}$, $\mu=1,2,3$, moves $c_3$ and $x^{(3)}$ by the same vector, and by (I.AP) it changes the sign of exactly two factors in each of the two terms below: of $\bar{\psi}(x)$ and of $\psi(x)+\chi(x,x^{(3)})$ in the first term, and of $\bar{\psi}(x)+\bar{\chi}(x,x^{(3)})$ and of $\psi(x)$ in the second.
The remaining factors, the matrices $M_3(x)$, $M_3(x^{(3)})$, which depend on $g$ alone, and the numbers $\exp\left(\pm i\,q\cdot\mathscr{Z}^{(i)}(c,c_3)\right)-1$, are $\ell_i\mathbb{Z}^3$-periodic.
The two signs cancel; so each term is invariant under the translations by $\ell_i e_1$, $\ell_i e_2$, $\ell_i e_3$, and therefore under every element of $\ell_i\mathbb{Z}^3$.
The result is
\begin{multline*}
G^{(1),(i)}\left(s^{(i)},\mathrm{pr}_i\tilde{\square}\right)=\frac{1}{2}\cdot\frac{1}{48}\sum_{c\in\tilde{\square}}\ \sum_{g\in O_{\mathrm{cub}}}\biggl\{\bar{\psi}(x)\cdot M_3(x)\left(\exp\left(i\,q\cdot\mathscr{Z}^{(i)}(c,c_3)\right)-1\right)\left(\psi(x)+\chi(x,x^{(3)})\right)\\
+\left(\bar{\psi}(x)+\bar{\chi}(x,x^{(3)})\right)\cdot M_3(x^{(3)})\left(\exp\left(-i\,q\cdot\mathscr{Z}^{(i)}(c,c_3)\right)-1\right)\psi(x)\biggr\},
\end{multline*}
where $M_3(x)=3\cdot\left(\rho-\gamma\cdot(g e_1)\right)$ and $M_3(x^{(3)})=3\cdot\left(\rho+\gamma\cdot(g e_1)\right)$.
The pairs $(x,x^{(3)})$, $c\in\tilde{\square}$, belong to $\mathrm{Pad}(\tilde{\square})$ (kind (II)); so this element belongs to $\mathrm{Grq}_i(\tilde{\square})$, a subalgebra of $\mathrm{Gr}_i(\tilde{\mathbf{Y}})$.
The bonds $(c,c_3)$, $c\in\tilde{\square}$, meet $\tilde{\square}$ and are contained in $\tilde{\mathbf{Y}}^{\sim1}$; by (c1), $\mathscr{Z}^{(1)}$ and $\mathscr{Z}^{(2)}$ have equal lifts on them.
Proposition~27.4\,(b) with $\mathrm{Zf}_i=\mathscr{Z}^{(i)}$ now gives
\[
\mathrm{idn}\,G^{(1),(1)}\left(s^{(1)},\mathrm{pr}_1\tilde{\square}\right)=G^{(1),(2)}\left(s^{(2)},\mathrm{pr}_2\tilde{\square}\right).
\]
At this level neither a contour nor a field $\Psi$ occurs.

\emph{Reading (o).}
In (c5.3) the activity is $E^{r,(i)}\left(\mathrm{pr}_i\tilde{X};A_i+t\cdot\mathscr{Z}^{(i)}\right)$, and (c4) in the reading (o) applies to it.
The rest of the step is unchanged, with $\mathscr{T}_i=A_i$.

\paragraph{Step (c6): first localisation, $\mathrm{Loc}$ and the substitution.}

(c6.1) \emph{The terms correspond.}
Let the kinds of the parameter cubes be fixed: here they are $\pi_k$, $\pi'_k$, $\pi''_k$, and in (c7) the kind $\pi_k$ alone; $P_0(\,\cdot\,)$ is formed with these kinds.
Let $\mathrm{Tmk}(\tilde{\mathbf{Y}})$ be the set of the pairs $(\tilde{X},\tilde{P})$ such that $\tilde{X}$ is connected in $\tilde{\mathbf{Y}}$, $\tilde{P}$ is a set of parameter cubes of the covering lattice, of these kinds, which are contained in $\tilde{\mathbf{Y}}$ and do not belong to $P_0(\tilde{X})$, and $Y'(\tilde{X},\tilde{P})=\tilde{\mathbf{Y}}$.
Then
\[
(\tilde{X},\tilde{P})\longmapsto\left(\mathrm{pr}_i\tilde{X},\mathrm{pr}_i\tilde{P}\right)
\]
is a bijection of $\mathrm{Tmk}(\tilde{\mathbf{Y}})$ onto the set of the pairs $(X,P)$ of (8.2) on $T_i$ with $Y'(X,P)=\mathbf{Y}_i$.

\begin{proof}[Proof of (c6.1)]
Proposition~27.2\,(e) is used with the admissible set $\tilde{\mathbf{Y}}$ in the place of $\tilde{Y}$.
(Since the cubes of $P_0(X)$ are contained in $X$, the set $X\cup\mathrm{Yp}(X,P)$ is the union of $X$ and of the cubes of $P$; hence $Y'(X,P)$ does not depend on the kinds that are counted in $P_0(X)$.
As $\mathrm{pr}_i$ preserves the kind of a cube, the identity $P_0(\mathrm{pr}_i\tilde{X})=\mathrm{pr}_iP_0(\tilde{X})$ holds kind by kind.
So the identities of Proposition~27.2\,(e) hold for either of the two choices of kinds.)
Let $Y'(X,P)=\mathbf{Y}_i$.
By the last sentence of Proposition~27.2\,(e) (see also Proposition~24.3\,(al)), $X$ and all cubes of $P$ are contained in $\mathbf{Y}_i$.
Hence $X=\mathrm{pr}_i(\tilde{X})$ for exactly one $\tilde{X}$ connected in $\tilde{\mathbf{Y}}$, and $P=\mathrm{pr}_i(\tilde{P})$ for exactly one set $\tilde{P}$ of parameter cubes contained in $\tilde{\mathbf{Y}}$.
Since $\mathrm{pr}_i$ is injective on these cubes and $P_0(X)=\mathrm{pr}_iP_0(\tilde{X})$, the set $\tilde{P}$ is disjoint from $P_0(\tilde{X})$.
Moreover $Y'(\tilde{X},\tilde{P})$ is contained in $\tilde{\mathbf{Y}}$ and $\mathrm{pr}_iY'(\tilde{X},\tilde{P})=Y'(X,P)=\mathrm{pr}_i(\tilde{\mathbf{Y}})$; as $\mathrm{pr}_i$ is injective on $\tilde{\mathbf{Y}}$, it follows that $Y'(\tilde{X},\tilde{P})=\tilde{\mathbf{Y}}$.
Conversely, let $(\tilde{X},\tilde{P})\in\mathrm{Tmk}(\tilde{\mathbf{Y}})$.
Then $\mathrm{pr}_i(\tilde{P})$ is disjoint from $\mathrm{pr}_iP_0(\tilde{X})=P_0(\mathrm{pr}_i\tilde{X})$, and $Y'(\mathrm{pr}_i\tilde{X},\mathrm{pr}_i\tilde{P})=\mathrm{pr}_iY'(\tilde{X},\tilde{P})=\mathbf{Y}_i$.
\end{proof}

(c6.2) \emph{$\mathrm{Loc}$ of corresponding functions.}
Let the kinds be fixed as in (c6.1).
For $i=1,2$ let $\mathrm{Gfn}^{(i)}(\,\cdot\,,X)$, $X\in D_k(T_i)$, be functions as in Definition~24.D\,(D3), with parameters of these kinds, and let $\mathrm{Vs}_i$ be a linear subspace of $\mathrm{Gr}_i(\tilde{\mathbf{Y}})$.
Assume that for every $\tilde{X}$ connected in $\tilde{\mathbf{Y}}$ and every $\mathrm{sval}$ the element $\mathrm{Gfn}^{(i)}\left(s^{(i)}[\tilde{\mathbf{Y}};\mathrm{sval}],\mathrm{pr}_i\tilde{X}\right)$ belongs to $\mathrm{Vs}_i$ and
\begin{equation}
\mathrm{idn}\,\mathrm{Gfn}^{(1)}\left(s^{(1)}[\tilde{\mathbf{Y}};\mathrm{sval}],\mathrm{pr}_1\tilde{X}\right)=\mathrm{Gfn}^{(2)}\left(s^{(2)}[\tilde{\mathbf{Y}};\mathrm{sval}],\mathrm{pr}_2\tilde{X}\right).
\tag{27.8}
\end{equation}
Then $\mathrm{Loc}[\mathrm{Gfn}^{(i)}](\mathbf{Y}_i)$ belongs to $\mathrm{Vs}_i$, and
\[
\mathrm{idn}\,\mathrm{Loc}[\mathrm{Gfn}^{(1)}](\mathbf{Y}_1)=\mathrm{Loc}[\mathrm{Gfn}^{(2)}](\mathbf{Y}_2).
\]

\begin{proof}[Proof of (c6.2)]
Let $(\tilde{X},\tilde{P})\in\mathrm{Tmk}(\tilde{\mathbf{Y}})$.
For a family $\mathrm{spar}=(\mathrm{spar}_{\tilde{\mathrm{pc}}})_{\tilde{\mathrm{pc}}\in\tilde{P}}$ of complex numbers of modulus at most $M^{\alpha_0}$, let $\mathrm{sval}[\mathrm{spar}]$ be the family with the value $1$ at the cubes of $P_0(\tilde{X})$, the value $\mathrm{spar}_{\tilde{\mathrm{pc}}}$ at $\tilde{\mathrm{pc}}\in\tilde{P}$, and the value $0$ at the other parameter cubes contained in $\tilde{\mathbf{Y}}$ (note that $1\le M^{\alpha_0}$).
Since $P_0(\mathrm{pr}_i\tilde{X})=\mathrm{pr}_iP_0(\tilde{X})$, the family $s^{(i)}[\tilde{\mathbf{Y}};\mathrm{sval}[\mathrm{spar}]]$ takes, among the cubes of the fixed kinds, the value $1$ on $P_0(\mathrm{pr}_i\tilde{X})$, the value $\mathrm{spar}_{\tilde{\mathrm{pc}}}$ at $\mathrm{pr}_i(\tilde{\mathrm{pc}})$ for $\tilde{\mathrm{pc}}\in\tilde{P}$, and the value $0$ outside $P_0(\mathrm{pr}_i\tilde{X})\cup\mathrm{pr}_i(\tilde{P})$.
By (8.2) and (c6.1), therefore,
\begin{equation}
\begin{split}
\mathrm{Loc}[\mathrm{Gfn}^{(i)}](\mathbf{Y}_i)&=\sum_{(\tilde{X},\tilde{P})\in\mathrm{Tmk}(\tilde{\mathbf{Y}})}\ \int_{[0,1]^{\tilde{P}}}\mathrm{d}\,\mathrm{spar}\\
&\qquad\left(\prod_{\tilde{\mathrm{pc}}\in\tilde{P}}\frac{\partial}{\partial\,\mathrm{spar}_{\tilde{\mathrm{pc}}}}\right)\mathrm{Gfn}^{(i)}\left(s^{(i)}\left[\tilde{\mathbf{Y}};\mathrm{sval}[\mathrm{spar}]\right],\mathrm{pr}_i\tilde{X}\right).
\end{split}
\tag{27.9}
\end{equation}
Because $M^{\alpha_0}>1$, the cube $[0,1]^{\tilde{P}}$ lies in the interior of the polydisk, and the derivatives and integrals exist by the hypothesis of Definition~24.D\,(D3).
The identity (27.8) holds at all $\mathrm{sval}[\mathrm{spar}]$ with $\mathrm{spar}$ in the polydisk.
By (Co3), applied term by term, $\mathrm{idn}$ maps the right side of (27.9) for $i=1$ to the right side for $i=2$.
Finally, $\mathrm{Vs}_i$ has finite dimension, and therefore it is closed under derivatives and integrals.
\end{proof}

(c6.2$'$) \emph{Application.}
The functions $G^{(j),(i)}$, $j=1,2,3$, taken at $\left(A_i,\mathscr{C}_i,\mathsf{Z}^{(i)}[\tilde{\mathbf{Y}};\mathrm{zval}]\right)$, satisfy the hypothesis of Definition~24.D\,(D3): the displays are defined by Proposition~24.1\,(a), and the last sentence of Proposition~24.3\,(q5) applies.
By (c5) they also satisfy the hypothesis of (c6.2) for the kinds $\pi_k$, $\pi'_k$, $\pi''_k$, with $\mathrm{Vs}_i=\mathrm{GrqP}_i(\tilde{\mathbf{Y}})$ for $k\ge1$ and $\mathrm{Vs}_i=\mathrm{Grq}_i(\tilde{\mathbf{Y}})$ for $k=0$.
Consequently
\[
\mathrm{idn}\,\mathrm{Loc}[G^{(j)}]^{(1)}(\mathbf{Y}_1)=\mathrm{Loc}[G^{(j)}]^{(2)}(\mathbf{Y}_2),\qquad j=1,2,3.
\]
These are the elements $E^{(j),*,(i)}(\mathbf{Y}_i)$; they belong to $\mathrm{GrqP}_i(\tilde{\mathbf{Y}})$ if $k\ge1$, and to $\mathrm{Grq}_i(\tilde{\mathbf{Y}})$ if $k=0$.

(c6.3) \emph{The substitution; $k\ge1$.}
Let $y$ be a unit site of $\tilde{\mathbf{Y}}$.
The block $\Delta_y$ lies in $\tilde{\mathbf{Y}}$, by \cite[Section~17.4.1, (g0)]{QEDI}.
Read (8.3a)$^{\mathrm{ch}}$ and (23.1) at the point $y$ of the covering lattice, with $B=\mathscr{T}_i$, with $x=(c,g)$, $x^{(a)}=(c,g s_a)$ for $a=0,1,2$, $c_3=c+\eta_k\cdot g e_1$, $x^{(3)}=(c_3,g s_3)$, and with the periodic extension $\tilde{\mathbf{Y}}+\ell_i\mathbb{Z}^3$ of $\mathbf{Y}_i$.
This gives
\begin{align*}
\mathrm{sub}_{\mathbf{Y}_i,\mathscr{T}_i}\Psi(y)&=\frac{\eta_k^3}{48}\sum_{c\in\Delta_y}\ \sum_{g\in O_{\mathrm{cub}}}\varphi_{\mathscr{T}_i}(y,c)\cdot\biggl\{\left(1+\frac{\mu_k}{b_k}\right)\psi(x)\\
&\qquad-\left(b_k\cdot\eta_k\right)^{-1}\biggl[\sum_{a\le2}M_a(x)\left(\psi(x^{(a)})-\psi(x)\right)\\
&\qquad\qquad+\mathbf{1}\left[c_3\in\tilde{\mathbf{Y}}+\ell_i\mathbb{Z}^3\right]M_3(x)\left(u_{\mathscr{T}_i}(c,c_3)\,\psi(x^{(3)})-\psi(x)\right)\biggr]\\
&\qquad-b_k^{-1}\cdot\eta_k^{\alpha-1}\cdot\mathbf{1}\left[c_3\notin\tilde{\mathbf{Y}}+\ell_i\mathbb{Z}^3\right]M_3(x)\,\chi(x,x^{(3)})\biggr\},
\end{align*}
and $\mathrm{sub}_{\mathbf{Y}_i,\mathscr{T}_i}\bar{\Psi}(y)$ is the same expression with $\varphi_{\mathscr{T}_i}(y,c)^{-1}$, $u_{\mathscr{T}_i}(c,c_3)^{-1}$, $\mathsf{C}^{-1}M_a(x)\,\mathsf{C}$ ($a=0,\dots,3$), $\bar{\psi}$, $\bar{\chi}$ in the places of $\varphi_{\mathscr{T}_i}(y,c)$, $u_{\mathscr{T}_i}(c,c_3)$, $M_a(x)$, $\psi$, $\chi$.
(Indeed $\bar{T}^{(n)}(B,Y)=\mathsf{C}^{-1}T^{(n)}(-B,Y)\,\mathsf{C}$ for $n=1,2$, $\varphi_{-B}=\varphi_B^{-1}$ and $u_{-B}=u_B^{-1}$; the type of a bond is indexed here by $a$, since $i$ numbers the tori.)
The site $c_3$ has sup-distance at most $\eta_k$ from $\tilde{\mathbf{Y}}$; so it lies in $N(\tilde{\mathbf{Y}})$, and by (Mb) it lies in $\tilde{\mathbf{Y}}+\ell_i\mathbb{Z}^3$ if and only if it lies in $\tilde{\mathbf{Y}}$.
The sites $x$, $x^{(a)}$ lie over $\tilde{\mathbf{Y}}$; the site $x^{(3)}$ occurs in $\psi$ only if $c_3$ lies in $\tilde{\mathbf{Y}}$; and the pairs $(x,x^{(3)})$ with $c\in\tilde{\mathbf{Y}}$ and $c_3\notin\tilde{\mathbf{Y}}$ belong to $\mathrm{Pad}(\tilde{\mathbf{Y}})$ (kind (II)).
Hence $\mathrm{sub}_{\mathbf{Y}_i,\mathscr{T}_i}$, which fixes $\psi$ and $\chi$, maps $\mathrm{GrqP}_i(\tilde{\mathbf{Y}})$ into $\mathrm{Grq}_i(\tilde{\mathbf{Y}})$.
The connections $\mathscr{T}_i$ lie in $\tilde{G}_k(T_i)$, and $\mathscr{T}_1$, $\mathscr{T}_2$ have equal lifts on all bonds.
By Proposition~27.4\,(a-3), the rows of $T^{(1)}(\mathscr{T}_i)$, $T^{(2)}(\mathscr{T}_i)$, $\bar{T}^{(1)}(\mathscr{T}_i)$, $\bar{T}^{(2)}(\mathscr{T}_i)$, formed with $\mathbf{Y}_i$, at the unit sites of $\tilde{\mathbf{Y}}$ are the same for $i=1,2$; that is, the two displays have the same coefficients for $i=1$ and for $i=2$.
So the homomorphisms $\mathrm{idn}\circ\mathrm{sub}_{\mathbf{Y}_1,\mathscr{T}_1}$ and $\mathrm{sub}_{\mathbf{Y}_2,\mathscr{T}_2}\circ\mathrm{idn}$ of $\mathrm{GrqP}_1(\tilde{\mathbf{Y}})$ into $\mathrm{Grq}_2(\tilde{\mathbf{Y}})$ agree on all generators, and hence everywhere.
By Definition~24.D\,(D4), $E^{(j),x,(i)}(\mathbf{Y}_i)=\mathrm{sub}_{\mathbf{Y}_i,\mathscr{T}_i}\left(E^{(j),*,(i)}(\mathbf{Y}_i)\right)$.
Together with (c6.2$'$) this yields
\[
\mathrm{idn}\,E^{(j),x,(1)}(\mathbf{Y}_1)=E^{(j),x,(2)}(\mathbf{Y}_2),\qquad j=1,2,3,
\]
at the arguments $\left(A_i,\mathscr{C}_i,\mathsf{Z}^{(i)}[\tilde{\mathbf{Y}};\mathrm{zval}]\right)$, and these elements belong to $\mathrm{Grq}_i(\tilde{\mathbf{Y}})$.

\emph{Level $0$.}
There is no substitution: $E^{(1),x}=E^{(1),*}=E^{(\le3),x}$, and $E^{(2),x}=E^{(3),x}=0$.

Summation over $j$ gives
\[
\mathrm{idn}\,E^{(\le3),x,(1)}(\mathbf{Y}_1)=E^{(\le3),x,(2)}(\mathbf{Y}_2).
\]
Here $\tilde{\mathbf{Y}}$ was an arbitrary set connected in $\tilde{Y}$.

\emph{Reading (o).}
The step is the same, with $\mathscr{T}_i=A_i$ and the arguments $\left(A_i,\mathsf{Z}^{(i)}\right)$.

\paragraph{Step (c7): second localisation.}

(c7.1) \emph{Change of $\mathsf{Z}$.}
Let $\tilde{X}$ be connected in $\tilde{\mathbf{Y}}$.
The basis elements of $T_i$ that are contained in $\mathrm{pr}_i(\tilde{X})$ are the $\mathrm{pr}_i(\Upsilon)$, $\Upsilon\in\mathrm{Bs}(\tilde{X})$ (Proposition~27.2\,(e), with the admissible set $\tilde{X}$ in the place of $\tilde{Y}$), and $\mathrm{Bs}(\tilde{X})$ is a subset of $\mathrm{Bs}(\tilde{\mathbf{Y}})$.
Let $\mathrm{zval}'$ be the restriction of $\mathrm{zval}$ to $\mathrm{Bs}(\tilde{X})$.
Then $\mathsf{Z}^{(i)}[\tilde{\mathbf{Y}};\mathrm{zval}]$ and $\mathsf{Z}^{(i)}[\tilde{X};\mathrm{zval}']$ have the same coordinates at the basis elements contained in $\mathrm{pr}_i(\tilde{X})$.
By the fourth sentence of Proposition~24.3\,(ze), each of the objects $E^{(j),*}$, $E^{(j),x}$, $E^{(\le3),x}$, $E^{(\le5),x}$, $\mathrm{Km}$ attached to $\mathrm{pr}_i(\tilde{X})$ takes the same value at these two elements of $\mathrm{Box}_k(T_i)$.
It follows that an identity between such objects on $T_1$ and on $T_2$ which holds at $\mathsf{Z}^{(i)}[\tilde{X};\mathrm{zval}']$ for all $\mathrm{zval}'$ also holds at $\mathsf{Z}^{(i)}[\tilde{\mathbf{Y}};\mathrm{zval}]$ for all $\mathrm{zval}$.
In particular, (c6) with $\tilde{X}$ in the place of $\tilde{\mathbf{Y}}$ gives
\begin{equation}
\mathrm{idn}\,E^{(\le3),x,(1)}(\mathrm{pr}_1\tilde{X})=E^{(\le3),x,(2)}(\mathrm{pr}_2\tilde{X})\qquad\text{at }\left(A_i,\mathscr{C}_i,\mathsf{Z}^{(i)}[\tilde{\mathbf{Y}};\mathrm{zval}]\right),
\tag{27.10}
\end{equation}
and both elements belong to $\mathrm{Grq}_i(\tilde{X})$.

(c7.2) \emph{The field $W_k^{\#}$.}
Let $x=(c,g)$ be a fermion site with $\mathrm{dist}(c,\tilde{\mathbf{Y}})\le\eta_k$.
By Definition~24.D\,(D2), with the kernels of Section~G4.0.4, by part~(4) of Step~(c3) and by (Rs) with $\tilde{S}=\tilde{\mathbf{Y}}$, applied twice, the components $\mathrm{Wc}^{(i)}$, $\overline{\mathrm{Wc}}^{(i)}$ of $W_k^{\#,(i)}\left(s''^{(i)},\mathscr{T}_i\right)$ are
\begin{equation}
\begin{aligned}
k\ge1:\quad&\mathrm{Wc}^{(i)}_{\beta}(x)=\sum_{y,y'\in\tilde{\mathbf{Y}}\cap\mathbb{Z}^3}\ \sum_{\beta',\beta''}H^{(i)}[\mathscr{T}_i](x,y)_{\beta\beta'}\cdot\Gamma_k^{(i)}\left(s''^{(i)},\mathscr{T}_i\right)(y,y')_{\beta'\beta''}\cdot W(y',\beta'',0),\\
&\overline{\mathrm{Wc}}^{(i)}_{\beta}(x)=\sum_{y\in\tilde{\mathbf{Y}}\cap\mathbb{Z}^3}\ \sum_{\beta'}\bar{H}^{(i)}[\mathscr{T}_i](x,y)_{\beta\beta'}\cdot W(y,\beta',1);\\
k=0:\quad&\mathrm{Wc}^{(i)}_{\beta}(x)=\mathrm{dgf}_i(c)\cdot\frac{1}{48}\sum_{c'\in\tilde{\mathbf{Y}}}\ \sum_{g'\in O_{\mathrm{cub}}}\ \sum_{\beta'}\Gamma_0^{(i)}\left(s''^{(i)},\mathscr{T}_i\right)\left(x,(c',g')\right)_{\beta\beta'}\cdot W\left((c',g'),\beta',0\right),\\
&\overline{\mathrm{Wc}}^{(i)}_{\beta}(x)=\mathrm{dgf}_i(c)\cdot W(x,\beta,1).
\end{aligned}
\tag{27.11}
\end{equation}
(For $k=0$ one has $\left(\Gamma_0W\right)(x)=\int\mathrm{d}x'\,\Gamma_0(x,x')\,W(x')$ with the weight $\frac{1}{48}$ of (I.1.3).
If $c$ does not lie in $\tilde{\mathbf{Y}}$, then $\mathrm{dgf}_i(c)=0$ and both members vanish; if $c$ lies in $\tilde{\mathbf{Y}}$, the row of $\Gamma_0^{(i)}$ at $x$ is one of the rows of part~(2) of Step~(c3).)
These are odd elements of $\mathrm{Gr}_i(\tilde{\mathbf{Y}})$, linear in the $W(u)$, $u$ in (over) $\tilde{\mathbf{Y}}$.
By parts~(1) and~(2) of Step~(c3) their coefficients are the same for $i=1,2$.
At a pair $(x,x')$ of $\mathrm{Pad}(\tilde{X})$, $\tilde{X}$ connected in $\tilde{\mathbf{Y}}$, one has
\begin{align*}
\left(\delta_{\alpha,\mathscr{T}_i}\mathrm{Wc}^{(i)}\right)(x,x')&=d_X(x,x')^{-\alpha}\cdot\left(U_{\mathscr{T}_i}(x,x')\,\mathrm{Wc}^{(i)}(x')-\mathrm{Wc}^{(i)}(x)\right),\\
\left(\delta_{\alpha,-\mathscr{T}_i}\overline{\mathrm{Wc}}^{(i)}\right)(x,x')&=d_X(x,x')^{-\alpha}\cdot\left(U_{\mathscr{T}_i}(x,x')^{-1}\,\overline{\mathrm{Wc}}^{(i)}(x')-\overline{\mathrm{Wc}}^{(i)}(x)\right),
\end{align*}
with numbers $d_X(x,x')$ and $U_{\mathscr{T}_i}(x,x')^{\pm1}$ that are the same for $i=1,2$ (Proposition~27.4\,(a-2) at $B_i=\mathscr{T}_i$); inside a cluster $d_X=\eta_k$ and $U_{\mathscr{T}_i}=1$, and at level $0$ one has $d_X\ge1$ for all pairs.
Let $\mathrm{wW}^{(i)}$ be the restriction of $W_k^{\#,(i)}\left(s''^{(i)},\mathscr{T}_i\right)$ to the indices $(x,\beta,\omega)$ with $x$ over $\tilde{X}$ and $(x,x',\beta,\omega)$ with $(x,x')\in\mathrm{Pad}(\tilde{X})$.
By (Co1),
\begin{equation}
\mathrm{idn}\,\mathrm{wW}^{(1)}(\xi)=\mathrm{wW}^{(2)}(\xi),\qquad\mathrm{idn}\,\mathrm{wW}^{(1)}(\boldsymbol{\zeta})=\mathrm{wW}^{(2)}(\boldsymbol{\zeta}),
\tag{27.12}
\end{equation}
and the same holds for $t\cdot\mathrm{wW}^{(i)}$, $t$ complex.

(c7.3) \emph{The three functions.}
Let $\tilde{X}$ be connected in $\tilde{\mathbf{Y}}$.
Definition~24.D\,(D4), Section~0.5, (27.10), (c4) at $t=0$ and (Fm1) give
\begin{align*}
G^{(\le3),(i)}\left(s''^{(i)},\mathrm{pr}_i\tilde{X}\right)&=\mathrm{adf}_i\left[\mathrm{wW}^{(i)}\right]\left(E^{(\le3),x,(i)}\left(\mathrm{pr}_i\tilde{X};A_i,\mathscr{C}_i,\mathsf{Z}^{(i)}[\tilde{\mathbf{Y}};\mathrm{zval}]\right)\right),\\
G^{(4),(i)}\left(s''^{(i)},\mathrm{pr}_i\tilde{X}\right)&=\oint_{\mathrm{rad}_4}\mathrm{adf}_i\left[t\cdot\mathrm{wW}^{(i)}\right]\left(\mathscr{E}^{r,(i)}\left(\mathrm{pr}_i\tilde{X};A_i,\mathscr{C}_i\right)\right)\qquad(k\ge1;\ G^{(4)}=0\text{ for }k=0),\\
G^{(5),(i)}\left(s''^{(i)},\mathrm{pr}_i\tilde{\square}\right)&=\oint_{\mathrm{rad}_4}(-m')\cdot\mathrm{adf}_i\left[t\cdot\mathrm{wW}^{(i)}\right]\left(\langle\bar{\psi},\psi\rangle_{\mathrm{pr}_i\tilde{\square}}\right)\qquad(\tilde{X}=\tilde{\square}\text{ an }M\text{-cube}),
\end{align*}
and $G^{(5)}=0$ on the polymers that are not $M$-cubes.
The arguments of $\mathrm{adf}_i$ belong to $\mathrm{Grq}_i(\tilde{X})$; the values are polynomials in the generators of $\mathrm{Grq}_i(\tilde{X})$ and in the $W(u)$, $u$ in (over) $\tilde{\mathbf{Y}}$, and so they belong to $\mathrm{GrqW}_i(\tilde{\mathbf{Y}})$.
By (Ad) with (27.12), and by (27.10), (c4) at $t=0$, (Fm1) and (a) (there is one number $m'$), the map $\mathrm{idn}$ takes the first function on $T_1$ to the first function on $T_2$, and likewise the integrands of the second and of the third function for every $t$ with $|t|=\mathrm{rad}_4$.
By (Co3), the contour integrals being defined by Proposition~24.1\,(a),
\[
\mathrm{idn}\left(G^{(\le3)}+G^{(4)}+G^{(5)}\right)^{(1)}\left(s''^{(1)},\mathrm{pr}_1\tilde{X}\right)=\left(G^{(\le3)}+G^{(4)}+G^{(5)}\right)^{(2)}\left(s''^{(2)},\mathrm{pr}_2\tilde{X}\right)
\]
for all $\mathrm{zval}$ and all $\mathrm{sval}$, and both sides belong to $\mathrm{GrqW}_i(\tilde{\mathbf{Y}})$.
Only the entries of $\mathrm{sval}$ at the cubes of the kind $\pi_k$ occur.

(c7.4) \emph{$\mathrm{Loc}$.}
By Definition~24.D\,(D4), $E^{(\le5),x,(i)}(\mathbf{Y}_i)=\mathrm{Loc}\left[G^{(\le3)}+G^{(4)}+G^{(5)}\right]^{(i)}(\mathbf{Y}_i)$, with the parameter cubes of the kind $\pi_k$.
The functions $\left(G^{(\le3)}+G^{(4)}+G^{(5)}\right)^{(i)}$, taken at $\left(A_i,\mathscr{C}_i,\mathsf{Z}^{(i)}[\tilde{\mathbf{Y}};\mathrm{zval}]\right)$, satisfy the hypothesis of Definition~24.D\,(D3) (Proposition~24.1\,(a); Proposition~24.3\,(q5), last sentence), and by (c7.3) they satisfy the hypothesis of (c6.2) for the kind $\pi_k$, with $\mathrm{Vs}_i=\mathrm{GrqW}_i(\tilde{\mathbf{Y}})$.
Therefore (c6.2) gives
\[
\mathrm{idn}\,E^{(\le5),x,(1)}(\mathbf{Y}_1)=E^{(\le5),x,(2)}(\mathbf{Y}_2)\qquad\text{at }\left(A_i,\mathscr{C}_i,\mathsf{Z}^{(i)}[\tilde{\mathbf{Y}};\mathrm{zval}]\right),\ \text{for all }\mathrm{zval},
\]
and these elements belong to $\mathrm{GrqW}_i(\tilde{\mathbf{Y}})$.

\emph{Reading (o).}
In (c7.3) the objects are $E^{(\le3),(i)}\left(\mathrm{pr}_i\tilde{X};A_i,\mathsf{Z}^{(i)}\right)$ and $E^{r,(i)}\left(\mathrm{pr}_i\tilde{X};A_i\right)$.
The rest of the step is unchanged, with $\mathscr{T}_i=A_i$.

\paragraph{Step (d): cluster expansion.}
Let $\tilde{X}$ be connected in $\tilde{Y}$.
By Proposition~27.2\,(e), applied with the admissible set $\tilde{X}$ in the place of $\tilde{Y}$, the map $\tilde{X}'\mapsto\mathrm{pr}_i(\tilde{X}')$ is a bijection of the sets $\tilde{X}'$ connected in $\tilde{X}$ onto the polymers of $T_i$ contained in $\mathrm{pr}_i(\tilde{X})$, and it preserves unions in both directions.

(d.1) \emph{The function $\mathrm{Km}$.}
The polymers of a set $\{X_a\}$ whose union is $\mathrm{pr}_i(\tilde{X})$ are contained in $\mathrm{pr}_i(\tilde{X})$.
Hence Definition~24.D\,(D5) gives
\[
\mathrm{Km}^{(i)}(\mathrm{pr}_i\tilde{X})=\sum_{\{\tilde{X}_a\}}\ \prod_a\left(\exp\left(E^{(\le5),x,(i)}(\mathrm{pr}_i\tilde{X}_a)\right)-1\right),
\]
the sum being over the sets $\{\tilde{X}_a\}$ of distinct sets connected in $\tilde{X}$ whose union is $\tilde{X}$; this index set is the same for $i=1,2$.
(The factors are even, so that their order is immaterial.)
Let all objects be taken at $\left(A_i,\mathscr{C}_i,\mathsf{Z}^{(i)}[\tilde{X};\mathrm{zval}]\right)$.
By (c7) with $\tilde{\mathbf{Y}}:=\tilde{X}_a$, and by (c7.1),
\[
\mathrm{idn}\,E^{(\le5),x,(1)}(\mathrm{pr}_1\tilde{X}_a)=E^{(\le5),x,(2)}(\mathrm{pr}_2\tilde{X}_a)
\]
at these arguments.
By (Co2),
\[
\mathrm{idn}\,\mathrm{Km}^{(1)}(\mathrm{pr}_1\tilde{X})=\mathrm{Km}^{(2)}(\mathrm{pr}_2\tilde{X})\qquad\text{at }\mathsf{Z}^{(i)}[\tilde{X};\mathrm{zval}],\ \text{for all }\mathrm{zval},
\]
and $\mathrm{Km}^{(i)}(\mathrm{pr}_i\tilde{X})$ belongs to $\mathrm{GrqW}_i(\tilde{X})$, by (Gq2).

(d.2) \emph{The function $\mathrm{Hm}$.}
Let $\mathrm{nrm}_k$ be the standard normal law on $\mathbb{R}$, restricted to $[-p_{0,k},p_{0,k}]$ and normalised.
By Proposition~27.4\,(f), in the coordinates $\mathsf{Z}(\Upsilon)$ the measure $\hat{\mu}$ on $T_i$ is the product of copies of $\mathrm{nrm}_k$ over the basis elements of $T_i$.
The element $\mathrm{Km}^{(i)}(\mathrm{pr}_i\tilde{X})$ depends on $\mathsf{Z}$ only through the coordinates at the basis elements contained in $\mathrm{pr}_i(\tilde{X})$ (Proposition~24.3\,(ze)), and these basis elements are the $\mathrm{pr}_i(\Upsilon)$, $\Upsilon\in\mathrm{Bs}(\tilde{X})$.
The remaining coordinates are integrated with probability measures (Fubini's theorem; Proposition~24.3\,(de)).
Hence
\[
\mathrm{Hm}^{(i)}(\mathrm{pr}_i\tilde{X})=\int_{[-p_{0,k},p_{0,k}]^{\mathrm{Bs}(\tilde{X})}}\ \prod_{\Upsilon\in\mathrm{Bs}(\tilde{X})}\mathrm{d}\,\mathrm{nrm}_k(\mathrm{zval}_{\Upsilon})\left(\int\mathrm{d}\mu_I(W)\left.\mathrm{Km}^{(i)}(\mathrm{pr}_i\tilde{X})\right|_{\mathsf{Z}=\mathsf{Z}^{(i)}[\tilde{X};\mathrm{zval}]}\right).
\]
Proposition~27.4\,(f) and (d.1) give, for every $\mathrm{zval}$,
\[
\mathrm{idn}\int\mathrm{d}\mu_I(W)\,\mathrm{Km}^{(1)}(\mathrm{pr}_1\tilde{X})=\int\mathrm{d}\mu_I(W)\,\mathrm{idn}\,\mathrm{Km}^{(1)}(\mathrm{pr}_1\tilde{X})=\int\mathrm{d}\mu_I(W)\,\mathrm{Km}^{(2)}(\mathrm{pr}_2\tilde{X}).
\]
(The integral of a monomial in the $W(u)$ is $0$ or $\pm\prod\nu(\mathrm{mo})^{-1}$, the product being over the modes $\mathrm{mo}$ whose pair $W(\mathrm{mo})$, $\bar{W}(\mathrm{mo})$ occurs; here $\nu=1$ for $k\ge1$, and $\nu=\frac{1}{48}$ for $k=0$, where $\int\mathrm{d}\mu_I\,W_{\beta}(x)\,\bar{W}_{\beta'}(x')=48\cdot\mathbf{1}[x=x',\beta=\beta']$.
These numbers are the same on both tori: (I.B3), (I.B4), (I.L0.5).)
The measure in $\mathrm{zval}$ is the same for $i=1,2$.
By (Co3),
\[
\mathrm{idn}\,\mathrm{Hm}^{(1)}(\mathrm{pr}_1\tilde{X})=\mathrm{Hm}^{(2)}(\mathrm{pr}_2\tilde{X}).
\]
The element $\mathrm{Hm}^{(i)}(\mathrm{pr}_i\tilde{X})$ belongs to $\mathrm{Grq}_i(\tilde{X})$ (Proposition~24.3\,(ze); (Gq3)).

(d.3) \emph{The series.}
The polymers of a family $(X_1,\dots,X_n)$ whose union is $\mathrm{pr}_i(\tilde{X})$ are contained in $\mathrm{pr}_i(\tilde{X})$.
So by (8.4)
\[
E^{\#,x,(i)}_k\left(m,\mathscr{E}^{(i)}\right)\left(\mathrm{pr}_i\tilde{X};A_i,\mathscr{C}_i\right)=\sum_{n\ge1}\frac{1}{n!}\sum_{(\tilde{X}_1,\dots,\tilde{X}_n)}\mathrm{urs}^{T}\left(\mathrm{pr}_i\tilde{X}_1,\dots,\mathrm{pr}_i\tilde{X}_n\right)\prod_{a=1}^{n}\mathrm{Hm}^{(i)}(\mathrm{pr}_i\tilde{X}_a),
\]
where the inner sum is over the finitely many families of sets connected in $\tilde{X}$ whose union is $\tilde{X}$; this index set is the same for $i=1,2$.
By its definition in \cite[Section~13.3]{QEDI}, where it is written $\rho^T$, the Ursell coefficient $\mathrm{urs}^{T}(X_1,\dots,X_n)$ is a function of the graph on $\{1,\dots,n\}$ whose edges are the pairs $\{a,j\}$, $a\ne j$, such that $X_a\cup X_j$ is connected through faces.
By Proposition~27.2\,(e), $\mathrm{pr}_i(\tilde{X}_a)\cup\mathrm{pr}_i(\tilde{X}_j)$ is connected if and only if $\tilde{X}_a\cup\tilde{X}_j$ is connected.
So the graph, and with it the coefficient $\mathrm{urs}^{T}$, is the same for $i=1,2$.
By (d.2) for the sets $\tilde{X}_a$ (which are connected in $\tilde{Y}$) and by (Co2), $\mathrm{idn}$ maps every term on $T_1$ to the term on $T_2$ with the same index.
Both series converge absolutely in the norm $\|\cdot\|_{h_k/4}$ (Proposition~24.1\,(a), first sentence, on $T_1$ and on $T_2$, by (TT.1)).
By (Co3),
\[
\mathrm{idn}\,E^{\#,x,(1)}_k\left(\mathrm{pr}_1\tilde{X};A_1,\mathscr{C}_1\right)=E^{\#,x,(2)}_k\left(\mathrm{pr}_2\tilde{X};A_2,\mathscr{C}_2\right).
\]
By (d.2) and (Gq2) every term of the series on $T_i$ belongs to $\mathrm{Grq}_i(\tilde{X})$, which is a linear subspace of finite dimension and hence closed.
Therefore $E^{\#,x,(i)}_k\left(\mathrm{pr}_i\tilde{X};A_i,\mathscr{C}_i\right)$ belongs to $\mathrm{Grq}_i(\tilde{X})$.

\emph{Reading (o).}
One has $\mathrm{idn}\,\mathrm{Km}^{(1)}=\mathrm{Km}^{(2)}$, $\mathrm{idn}\,\mathrm{Hm}^{(1)}=\mathrm{Hm}^{(2)}$ and
\[
\mathrm{idn}\,E^{\#,(1)}_k\left(m,E_{(1)}\right)\left(\mathrm{pr}_1\tilde{X};A_1\right)=E^{\#,(2)}_k\left(m,E_{(2)}\right)\left(\mathrm{pr}_2\tilde{X};A_2\right),
\]
and $E^{\#,(i)}_k\left(m,E_{(i)}\right)\left(\mathrm{pr}_i\tilde{X};A_i\right)$ belongs to $\mathrm{Grq}_i(\tilde{X})$.
The two series converge by Proposition~24.1\,(b), in the form stated in (o-3).

\paragraph{Range of validity.}
Proposition~27.5 introduces no constant and asserts no bound.
The only inequalities between numbers among its assertions are $1<\mathrm{rad}_2\le\mathrm{rad}_1$ and $\mathrm{rad}_4>1$ in (c2); by their definition in Section~I.5.7 the radii are expressions in $e_k$ and $\varepsilon$ alone.
The remaining assertions are of the following kinds.

First, there are identities between objects on $T_1$ and on $T_2$, among them the relations $\sim$ of (a); such identities are asserted in (a), (b), (c1), in the last sentence of (c2), in (c3)--(c7) and in (d).

Second, there are assertions on each torus separately.
The families of (a) are local families in the sense of Proposition~27.2\,(D3): on each torus their elements are even, and they have the property (k2) of Section~I.5.4.
In (b) the traces are given by three formulas, as a sum over the unit sites of $X$ for $k\ge1$ and as integrals over $X$ for $k=0$.
The clause (c1) says which terms of the sum that defines $\mathscr{Z}(s')$ can be different from zero, and (c3) says that the rows vanish at the columns that have no lift in $\tilde{\mathbf{Y}}$.
The clause (c2) asserts the memberships $B_i(t)\in\tilde{G}_k(T_i)$ and $A_i+t\cdot\mathscr{Z}^{(i)}\left(s'^{(i)}\right)\in\tilde{R}_k(T_i)$, and (c5) asserts that $G^{(j),(i)}\left(s^{(i)},\mathrm{pr}_i\tilde{X}\right)$ belongs to $\mathrm{Gr}_i(\tilde{\mathbf{Y}})$, that is, that it is a polynomial in the generators $\psi$, $\chi$, $\Psi$ with indices in $N(\tilde{\mathbf{Y}})$.
The clauses (c1), (c5) and (c7) write out on each torus the field $\mathscr{Z}(s')$, the functions $G^{(j)}$ and the field $W_k^{\#}$, the last one together with the indices at which it enters.
The clause (d) asserts that the series (8.4) converge in norm on both tori; this follows from the first sentence of Proposition~24.1\,(a), applied on $T_1$ and on $T_2$, and in the reading (o) from Proposition~24.1\,(b).

Third, (c6) asserts a correspondence of index sets: the pairs $(X,P)$ of $T_1$ with $Y'(X,P)=\mathbf{Y}_1$ correspond to the pairs $(X,P)$ of $T_2$ with $Y'(X,P)=\mathbf{Y}_2$.

The identities of (b) between $\acute{E}_k$, $E^{d}_k$, ${\delta E^{\mathrm{det}}}_k$ and $E_k^{\mathrm{det}}$ presuppose that the integrals over $y$ in $\acute{E}_k$ exist on each torus.
This is shown in (b.6) by means of the property W7$^{\mathrm{ch}}$ for $\mathrm{par}=(\alpha_{0},\gamma_{\mathrm{w}},\kappa_{\mathrm{w}};C_{w})$; the bounds obtained there serve only to show that the integrals exist, and their size enters nowhere.
The constants that occur in the clauses are $M$ and $d_U$; in the data there occurs moreover $M^{\alpha_0}$, which bounds the moduli of the entries of $\mathrm{sval}$ and is the radius of the polydisk in (c6).
Here $d_U$ and $\alpha_0$ belong to the line (C4) of Section~I.8.2 and $M$ to the line (C5); by the order in which the constants are chosen (Section~I.8.1), none of them depends on the data (TT).
The numbers $e_k$, $\eta_k$, $p_{0,k}$, $\mathrm{rad}_1$, $\mathrm{rad}_2$, $\mathrm{rad}_4$ take the same values on both tori.
The clauses hold for all data $(\bar{m},e,R,N)$ of (FL) with $R<R_{\mathrm{top}}$, for every level $k<k_m$, every $\mathrm{Jr}=R_{\mathrm{top}}/R$, every admissible $\tilde{Y}$, and all $(A_1,\mathscr{C}_1)$, $\mathrm{zval}$, $\mathrm{sval}$ as above.
\qed

\subsection*{2. Proposition 27.6}

The statement of Proposition~27.6 is given in Part~II.

\subsubsection*{Proof of Proposition 27.6}

Throughout, the setting is that of Section~0.1.
Objects of level $k+1$ carry a prime where this is convenient.
Thus $T_1':=T_{k+1}(R)$ and $T_2':=T_{k+1}(R_{\mathrm{top}})$ are the tori $\mathbb{L}_{k+1}/(\ell_i/L)\mathbb{Z}^3$, and $\mathrm{pr}_i'\colon\mathbb{L}_{k+1}\to T_i'$ are the covering maps.
The map $c'\mapsto Lc'$ is a bijection of $\mathbb{L}_{k+1}$ onto $\mathbb{L}_k$; it carries $(\ell_i/L)\mathbb{Z}^3$ onto $\ell_i\mathbb{Z}^3$ and the $M$-cubes of $\mathbb{L}_{k+1}$ onto the $LM$-cubes of $\mathbb{L}_k$.
Consequently
\[
L\,\mathrm{pr}_i'(Z')=\mathrm{pr}_i(LZ')
\]
for every set $Z'$ of sites of $\mathbb{L}_{k+1}$.

\paragraph{(a)}
The equalities $\varepsilon(\mathscr{E}^{(1)})=\varepsilon(\mathscr{E}^{(2)})$ and $m(\mathscr{E}^{(1)})=m(\mathscr{E}^{(2)})$ are Lemma~27.A\,(ii).
By Proposition~27.4\,(f), the number
\[
\delta\varepsilon^{0}_{k}=-d_V\cdot\log\left((2\pi)^{-\frac{1}{2}}\int_{-p_{0,k}}^{p_{0,k}}\mathrm{d}z\,\exp\left(-\frac{z^2}{2}\right)\right)
\]
is the same on both tori, since neither $d_V=2-2L^{-3}$ nor $p_{0,k}$ depends on the side.

\emph{Level $0$.}
Here $\mathscr{E}^{(i)}=0$; hence $\mathrm{al}_{0}$ and $\mathrm{al}_{2}$ vanish, and $\varepsilon(0)=0$, $m(0)=0$ on both tori.

\paragraph{(b)}
(B.0) \emph{What has to be shown.}
By (TT.1), Proposition~24.1\,(d) holds on $T_i$.
Its hypotheses are those of Proposition~24.1\,(c): $m$ is real with $|m|\le w^{\mathrm{m}}_k$, $\mathscr{E}^{(i)}\in K^{x}_{k}(\ell_i)$ has norm at most $w^{\mathrm{E}}_k$, and $\mathscr{E}^{(i)}=0$ if $k=0$.
It follows that $\mathscr{E}^{(i)+}\in K^{x}_{k+1}(\ell_i/L)$.
In particular $\mathscr{E}^{(i)+}$ consists of even elements and has the property (k2) of Section~I.5.4 at level $k+1$; in other words, it is a local family on $\tilde{R}_{k+1}\times\mathfrak{X}_{k+1}$ in the sense of Proposition~27.2\,(D3) at level $k+1$.
By Proposition~27.2\,(c) at level $k+1$ it therefore suffices to prove the following.
Let $\tilde{U}$ be a connected union of $M$-cubes of $\mathbb{L}_{k+1}$ that is admissible at level $k+1$, that is, $\mathrm{diam}(\tilde{U})\le\ell_1/L-14\,d_U\cdot M$, and let $(\mathscr{A},\mathscr{C})\in\tilde{R}_{k+1}(T_1')\times\mathfrak{X}_{k+1}(T_1')$.
Then
\begin{equation}
\mathrm{idn}[\tilde{U}]\,\mathscr{E}^{(1)+}(\mathrm{pr}_1'\tilde{U};\mathscr{A},\mathscr{C})=\mathscr{E}^{(2)+}(\mathrm{pr}_2'\tilde{U};\mathrm{pj}^{*}\mathscr{A},\mathrm{pj}^{*}\mathscr{C}).
\tag{27.13}
\end{equation}
Here $\mathrm{pj}^{*}(\mathscr{A},\mathscr{C})$ lies in $\tilde{R}_{k+1}(T_2')\times\mathfrak{X}_{k+1}(T_2')$ by \cite[Lemma~17.7]{QEDI} at level $k+1$, because $\ell_1/L$ and $\ell_2/L$ lie in $\mathrm{Cl}$ by (TT.1).
Fix $\tilde{U}$, $\mathscr{A}$ and $\mathscr{C}$, and put
\[
Y_i:=\mathrm{pr}_i'(\tilde{U}),\qquad \tilde{Y}:=L\tilde{U}.
\]

(B.1) \emph{The set $\tilde{Y}$.}
Proposition~27.4\,(e) gives the following: $\tilde{Y}$ is admissible at level $k$; $N(\tilde{Y})$ is contained in $L\,N(\tilde{U})$; the homomorphism $S_L$ maps the part of $\mathrm{Gr}_i(\tilde{Y})$ generated by $\psi$, $\chi$ into $\mathrm{Gr}_i(\tilde{U})$ of level $k+1$; and
\begin{equation}
S_L\,\mathrm{idn}[\tilde{Y}]=\mathrm{idn}[\tilde{U}]\,S_L\qquad\text{on this part.}
\tag{27.14}
\end{equation}
Moreover, $\tilde{Y}$ is a union of $LM$-cubes of $\mathbb{L}_k$ that is connected through faces, and $LY_i=\mathrm{pr}_i(\tilde{Y})$.
For the connectedness, recall that $\tilde{U}$ is connected through faces, and that $c'\mapsto Lc'$ carries two $M$-cubes of $\mathbb{L}_{k+1}$ with a common face onto two $LM$-cubes of $\mathbb{L}_k$ with centres $z$ and $z+LM\cdot e_{\mu}$ (Section~I.1.1).
By (g0) of \cite[Section~17.4.1]{QEDI}, and since $L$ is odd, the $LM$-cube with centre $z$ is the union of the $M$-cubes with the centres $z+Mw$, where $w\in\mathbb{Z}^3$ and $|w|\le\frac{L-1}{2}$.
Any two of these $M$-cubes are joined by a chain of such cubes in which consecutive centres differ by $M\cdot e_{\nu}$ or by $-M\cdot e_{\nu}$; one brings the three coordinates of $w$ to their final values one after the other.
Finally, the $M$-cubes with the centres $z+M\cdot\frac{L-1}{2}\cdot e_{\mu}$ and $z+M\cdot\frac{L+1}{2}\cdot e_{\mu}$ have a common face, and the first of them lies in the $LM$-cube with centre $z$, the second in the $LM$-cube with centre $z+LM\cdot e_{\mu}$.
The relation $LY_i=\mathrm{pr}_i(\tilde{Y})$ is the identity $L\,\mathrm{pr}_i'(Z')=\mathrm{pr}_i(LZ')$ with $Z'=\tilde{U}$.

(B.2) \emph{The connections.}
Put $(A_1,\mathscr{C}_1):=(\mathscr{A}_L,\mathscr{C}_L)$.
Since $\ell_1$ and $\ell_1/L$ lie in $\mathrm{Cl}$ and $e_{k+1}=L^{\frac{1}{2}}\cdot e_k$, \cite[Lemma~17.5]{QEDI} applies.
Its part (a), with $c=1$, gives $A_1\in L^{-\frac{3}{4}}\cdot\tilde{R}_k(T_1)$; by its part (b) this set is contained in $\tilde{R}_k^{-}(T_1)$; and its part (c) gives $\mathscr{C}_1\in\mathfrak{X}_k(T_1)$.
As functions on the bonds of the covering lattice one has $\mathrm{pj}^{*}\mathscr{A}=\mathscr{A}$ and $\mathrm{pj}^{*}\mathscr{C}=\mathscr{C}$, by (TT.2) at level $k+1$; moreover $B_L(b)=L^{-\frac{1}{2}}\cdot B(b/L)$ (Section~I.3.7).
Hence
\[
\left((\mathrm{pj}^{*}\mathscr{A})_L,(\mathrm{pj}^{*}\mathscr{C})_L\right)=\mathrm{pj}^{*}(A_1,\mathscr{C}_1)=:(A_2,\mathscr{C}_2).
\]
So $m$, $\mathscr{E}^{(i)}$, $\tilde{Y}$ and $(A_1,\mathscr{C}_1)$ are data of Proposition~27.5.

(B.3) \emph{The polymers.}
For $\tilde{X}$ connected in $\tilde{Y}$ let $\mathrm{clo}(\tilde{X})$ denote the union of the $LM$-cubes of $\mathbb{L}_k$ that contain an $M$-cube of $\tilde{X}$.
Then
\[
\left\{X\in D_k(T_i):\bar{X}^{L}=LY_i\right\}=\left\{\mathrm{pr}_i(\tilde{X}):\tilde{X}\text{ connected in }\tilde{Y},\ \mathrm{clo}(\tilde{X})=\tilde{Y}\right\},
\]
and $\tilde{X}$ is determined by $\mathrm{pr}_i(\tilde{X})$.
Indeed, let $X\in D_k(T_i)$ satisfy $\bar{X}^{L}=LY_i=\mathrm{pr}_i(\tilde{Y})$.
The set $\bar{X}^{L}$ is the union of the $LM$-cubes that contain an $M$-cube of $X$ (Section~I.1.1), and every $M$-cube lies in one $LM$-cube by (g0); so $X$ is contained in $\bar{X}^{L}$.
By Proposition~27.2\,(e), $X=\mathrm{pr}_i(\tilde{X})$ for exactly one $\tilde{X}$ connected in $\tilde{Y}$; moreover $\mathrm{clo}(\tilde{X})$ is contained in $\tilde{Y}^{\sim(L-1)}$, and $\mathrm{pr}_i(\mathrm{clo}(\tilde{X}))=\bar{X}^{L}=\mathrm{pr}_i(\tilde{Y})$.
The map $\mathrm{pr}_i$ is injective on $\tilde{Y}^{\sim(L+1)}$, which contains both $\mathrm{clo}(\tilde{X})$ and $\tilde{Y}$; therefore $\mathrm{clo}(\tilde{X})=\tilde{Y}$.
Conversely, if $\tilde{X}$ is connected in $\tilde{Y}$ and $\mathrm{clo}(\tilde{X})=\tilde{Y}$, then
\[
\overline{\mathrm{pr}_i(\tilde{X})}^{\,L}=\mathrm{pr}_i(\mathrm{clo}(\tilde{X}))=\mathrm{pr}_i(\tilde{Y})=LY_i.
\]
The index set on the right-hand side does not depend on $i$.

(B.4) \emph{The display.}
One uses the display (8.6)$^{\mathrm{ch}}$ of Definition~24.D, the relation $LR\mathscr{E}=L\left((R\mathscr{E})^{\mathrm{nat}}\right)$ of Section~I.5.6, and the formula
\[
(LE)(Y;\mathscr{A},\mathscr{C})=\sum_{X:\bar{X}^{L}=LY}S_L\left[E(X;\mathscr{A}_L,\mathscr{C}_L)\right]
\]
of Section~I.3.7.
Together with (B.2) and (B.3) they give
\begin{equation}
\begin{split}
\mathscr{E}^{(i)+}(Y_i;\,\cdot\,)=\sum_{\substack{\tilde{X}\text{ connected in }\tilde{Y},\\ \mathrm{clo}(\tilde{X})=\tilde{Y}}}S_L\Bigl[&(R\mathscr{E}^{(i)})^{\mathrm{nat}}(\mathrm{pr}_i\tilde{X};A_i,\mathscr{C}_i)+{\delta E^{\mathrm{det}}}_k^{(i)}(\mathrm{pr}_i\tilde{X};A_i,\mathscr{C}_i)\\
&+E^{\#,x,(i)}_k(m,\mathscr{E}^{(i)})(\mathrm{pr}_i\tilde{X};A_i,\mathscr{C}_i)\Bigr],
\end{split}
\tag{27.15}
\end{equation}
where the arguments are $(\mathscr{A},\mathscr{C})$ for $i=1$ and $\mathrm{pj}^{*}(\mathscr{A},\mathscr{C})$ for $i=2$.

(B.5) \emph{The three terms.}
Let $\tilde{X}$ be connected in $\tilde{Y}$.
Then $\tilde{X}$ is connected and admissible (Proposition~27.1\,(a)).
Write $\mathrm{idn}=\mathrm{idn}[\tilde{Y}]$.
\begin{itemize}
\item By (TT.3) the pairs $(A_i,\mathscr{C}_i)$ lie in $\tilde{R}_k(T_i)\times\mathfrak{X}_k(T_i)$, and the two pairs have equal lifts on all bonds.
Proposition~27.5\,(a), Proposition~27.2\,(D4) and (Co4) give
\[
\mathrm{idn}\,(R\mathscr{E}^{(1)})^{\mathrm{nat}}(\mathrm{pr}_1\tilde{X};A_1,\mathscr{C}_1)=(R\mathscr{E}^{(2)})^{\mathrm{nat}}(\mathrm{pr}_2\tilde{X};A_2,\mathscr{C}_2).
\]
\item Proposition~27.5\,(b) gives
\[
{\delta E^{\mathrm{det}}}_k^{(1)}(\mathrm{pr}_1\tilde{X};A_1,\mathscr{C}_1)={\delta E^{\mathrm{det}}}_k^{(2)}(\mathrm{pr}_2\tilde{X};A_2,\mathscr{C}_2);
\]
both sides are numbers, and numbers are fixed by $\mathrm{idn}$ and by $S_L$.
\item Proposition~27.5\,(d) gives
\[
\mathrm{idn}\,E^{\#,x,(1)}_k(\mathrm{pr}_1\tilde{X};A_1,\mathscr{C}_1)=E^{\#,x,(2)}_k(\mathrm{pr}_2\tilde{X};A_2,\mathscr{C}_2).
\]
\end{itemize}
The three elements lie in the part of $\mathrm{Gr}_i(\tilde{Y})$ generated by $\psi$, $\chi$, and the second of them is a number.
By Proposition~27.5\,(a) the family $(R\mathscr{E}^{(i)})^{\mathrm{nat}}$ is a local family; so it has (k2) (Proposition~27.2\,(D3)), and the first element lies in $\mathrm{Grq}_i(\tilde{X})$ by (Gq3).
The third element lies in $\mathrm{Grq}_i(\tilde{X})$ by the last sentence of (d.3).

(B.6) \emph{Conclusion.}
Denote by $[\,\cdots\,]^{(i)}$ the bracket in (27.15).
Apply $\mathrm{idn}[\tilde{U}]$ to (27.15) for $i=1$, use (27.14) in each term, and then (B.5):
\begin{align*}
\mathrm{idn}[\tilde{U}]\,\mathscr{E}^{(1)+}(Y_1;\mathscr{A},\mathscr{C})
&=\sum_{\tilde{X}}S_L\,\mathrm{idn}[\tilde{Y}]\,[\,\cdots\,]^{(1)}
=\sum_{\tilde{X}}S_L\,[\,\cdots\,]^{(2)}\\
&=\mathscr{E}^{(2)+}(Y_2;\mathrm{pj}^{*}\mathscr{A},\mathrm{pj}^{*}\mathscr{C}).
\end{align*}
This is (27.13).

\emph{Level $0$.}
Here $\mathscr{E}^{(i)}=0$ and $(R\mathscr{E}^{(i)})^{\mathrm{nat}}=0$, so that only ${\delta E^{\mathrm{det}}}_0^{(i)}$ and $E^{\#,x,(i)}_0(m,0)$ occur in (27.15); that is,
\[
\mathscr{E}^{(i)+}=L\,{\delta E^{\mathrm{det}}}_0^{(i)}+E^{*,x,(i)}_0(m,0).
\]
The clauses of Proposition~27.5 used above hold at $k=0$; at this level the generators $W(x,\beta,\omega)$, $x\in T^{\mathrm{f}}_0$, carry the weight $1/48$, the traces are those of (b.4), (b.5) over $\mathfrak{S}^{\mathrm{ch}}(1,\ell_i)$, and the vertex is that of (c5.6).

\paragraph{(c)}
Let $\tilde{U}$ be connected and admissible at level $k+1$, let $\mathscr{C}'\in\mathfrak{X}_{k+1}(T_1')$, and let $Y_i$, $\tilde{Y}$ be as in (B.0).
By \cite[Lemma~17.7]{QEDI} at level $k+1$, $\mathrm{pj}^{*}\mathscr{C}'$ lies in $\mathfrak{X}_{k+1}(T_2')$.
By \cite[Lemma~17.5(c)]{QEDI}, $\mathscr{C}_1:=\mathscr{C}'_L$ lies in $\mathfrak{X}_k(T_1)$; by (B.2), $(\mathrm{pj}^{*}\mathscr{C}')_L=\mathrm{pj}^{*}\mathscr{C}_1=:\mathscr{C}_2$, and $\mathscr{C}_2\in\mathfrak{X}_k(T_2)$ by (TT.3).
The definition of $\mathfrak{P}_{k+1}^{(i)}$ and (B.3) give
\begin{align*}
\mathfrak{P}_{k+1}^{(1)}(Y_1,\mathscr{C}')&=\sum_{\tilde{X}}\left(\mathfrak{P}_k^{(1)}(\mathrm{pr}_1\tilde{X},\mathscr{C}_1)+E_k^{\mathrm{det},(1)}(\mathrm{pr}_1\tilde{X},\mathscr{C}_1)\right),\\
\mathfrak{P}_{k+1}^{(2)}(Y_2,\mathrm{pj}^{*}\mathscr{C}')&=\sum_{\tilde{X}}\left(\mathfrak{P}_k^{(2)}(\mathrm{pr}_2\tilde{X},\mathrm{pj}^{*}\mathscr{C}_1)+E_k^{\mathrm{det},(2)}(\mathrm{pr}_2\tilde{X},\mathscr{C}_2)\right),
\end{align*}
both sums extending over the $\tilde{X}$ connected in $\tilde{Y}$ with $\mathrm{clo}(\tilde{X})=\tilde{Y}$.
Each such $\tilde{X}$ is connected and admissible.
Hence $\mathfrak{P}_k^{(2)}(\mathrm{pr}_2\tilde{X},\mathrm{pj}^{*}\mathscr{C}_1)=\mathfrak{P}_k^{(1)}(\mathrm{pr}_1\tilde{X},\mathscr{C}_1)$, by the hypothesis $\mathfrak{P}_k^{(1)}\sim\mathfrak{P}_k^{(2)}$ (Proposition~27.2\,(D4)).
Further, $E_k^{\mathrm{det},(1)}(\mathrm{pr}_1\tilde{X},\mathscr{C}_1)=E_k^{\mathrm{det},(2)}(\mathrm{pr}_2\tilde{X},\mathscr{C}_2)$ by Proposition~27.5\,(b), applied with the data $m$, $\mathscr{E}^{(i)}$, $\tilde{Y}$ and $(A_1,\mathscr{C}_1):=(0,\mathscr{C}_1)$; here $0\in\tilde{R}_k^{-}(T_1)$ by (TT.3).
It follows that
\[
\mathfrak{P}_{k+1}^{(2)}(\mathrm{pr}_2'\tilde{U},\mathrm{pj}^{*}\mathscr{C}')=\mathfrak{P}_{k+1}^{(1)}(\mathrm{pr}_1'\tilde{U},\mathscr{C}'),
\]
that is, $\mathfrak{P}_{k+1}^{(1)}\sim\mathfrak{P}_{k+1}^{(2)}$.

\paragraph{(d)}
Proposition~24.1\,(d) on $T_i$ gives
\begin{align*}
\mathscr{E}^{(i)+}|_0
&=L R(\mathscr{E}^{(i)}|_0)+L E_k^{\mathrm{det},(i)}+E^{*,(i)}_k(m,\mathscr{E}^{(i)}|_0)\\
&=L R E_{(i)}+L E_k^{\mathrm{det},(i)}+E^{*,(i)}_k(m,E_{(i)}).
\end{align*}
By (b), $\mathscr{E}^{(1)+}\sim\mathscr{E}^{(2)+}$ as families with slot at level $k+1$.
Proposition~27.2\,(f) at level $k+1$ now yields $\mathscr{E}^{(1)+}|_0\sim\mathscr{E}^{(2)+}|_0$.

\paragraph{(e)}
The data are $m$ and $E_{(i)}$ as in (o-1) of Section~0.7.
Put
\[
F^{(i)}:=L R E_{(i)}+L E_k^{\mathrm{det},(i)}+E^{*,(i)}_k(m,E_{(i)}),
\]
a family without slot on $\tilde{R}_{k+1}(T_i')$.

(e.1) \emph{Clause (a) without slot.}
The equalities $\varepsilon(E_{(1)})=\varepsilon(E_{(2)})$ and $m(E_{(1)})=m(E_{(2)})$ are the last sentence of Lemma~27.A.
For $\delta\varepsilon^{0}_{k}$ the argument is that of (a).

(e.2) \emph{Proposition~27.5 without slot.}
Each of the steps (a), (b), (c), (c1)--(c7) and (d) of the proof of Proposition~27.5 in Section~1 holds in the reading (o) of Section~0.7, that is, with the data (o-1) and the objects (o-2).
The reason is that every fact on which such a step rests either does not involve the slot or has the counterpart named in (o-3), (o-4).
For each step this is verified at the end of the step, in the part that begins with the words \emph{In the reading (o)} (Steps (a), (b), (c), (c1), (c2), (c3)) or \emph{Reading (o)} (Steps (c4)--(c7), (d)); the replacements, and the facts on which Section~1 rests that do not involve the slot, are listed below.
Consequently, for every admissible $\tilde{Y}$, every $\mathscr{A}\in\tilde{R}_{k+1}(T_1')$, $A_1=\mathscr{A}_L$, $A_2=\mathrm{pj}^{*}A_1$, and every $\tilde{X}$ connected in $\tilde{Y}$, one has
\begin{equation}
\begin{aligned}
&(RE_{(1)})^{\mathrm{nat}}\sim(RE_{(2)})^{\mathrm{nat}}\qquad\text{(local families without slot on }\tilde{R}_k\text{)},\\
&E_k^{\mathrm{det},(1)}(\mathrm{pr}_1\tilde{X},A_1)=E_k^{\mathrm{det},(2)}(\mathrm{pr}_2\tilde{X},A_2),\\
&\mathrm{idn}\,E^{\#,(1)}_k(m,E_{(1)})(\mathrm{pr}_1\tilde{X};A_1)=E^{\#,(2)}_k(m,E_{(2)})(\mathrm{pr}_2\tilde{X};A_2).
\end{aligned}
\tag{27.16}
\end{equation}
Propositions~24.1 and 24.3 and the parts (D1)--(D6) of Definition~24.D have $(m,\mathscr{E})$ and $\mathscr{C}$ among their data.
Without slot they are available only through Proposition~24.1\,(b), the last sentence of Proposition~24.3, and Definition~24.D\,(D7).
This holds also for those of their clauses used in Section~1 that involve neither $\mathscr{C}$ nor an activity, among them the hypothesis at the head of Proposition~24.3\,(ga), the clauses (ga2) and (de) of Proposition~24.3, and Definition~24.D\,(D2), (D3).
The part (D0) of Definition~24.D involves neither these data nor the slot: the sets $P_0(X)$, $\mathrm{Yp}(X,P)$ and $Y'(X,P)$ defined there, which enter (c6.1), (c6.2) and Proposition~27.2\,(e), are formed from a polymer $X$ and a set $P$ of parameter cubes alone.
Since (D7) gives a reading of the parts (D1)--(D6) only, the part (D0) holds unchanged in the reading (o).
Finally, where a step of Section~1 refers to (TT.1), (TT.3), to Lemma~27.A, to a clause of Proposition~27.5, or to a step or a display of Section~1, the object referred to is taken in the reading (o).
In detail, the reading (o) of (TT.1), (TT.3), of Lemma~27.A and of the steps of Section~1 is obtained by the following replacements.

\begin{description}
\item[(TT.1), (TT.3)] Their final assertions are replaced by the assertions of (o-1) on the data of (o-1).
The final assertion of (TT.1) is that $(m,\mathscr{E}^{(i)})$ lies in $\mathrm{Ball}^{x}_{k}(1)$, a subset of $\mathrm{Ball}^{x}_{k}(4)$.
The final assertion of (TT.3), on the data of Proposition~27.5, is that $(A_2,\mathscr{C}_2)$ lies in $\tilde{R}_k^{-}(T_2)\times\mathfrak{X}_k(T_2)$, that $\mathscr{T}_2=\mathrm{pj}^{*}\mathscr{T}_1$, that $\mathscr{T}_i$, $\mathscr{C}_i$ and $0$ lie in $\tilde{G}_k(T_i)$, and that $A_1$ and $A_2$, $\mathscr{C}_1$ and $\mathscr{C}_2$, $\mathscr{T}_1$ and $\mathscr{T}_2$ have equal lifts on all bonds.
The assertions of (o-1) are that $(m,E_{(i)})$ lies in $\mathrm{Ball}_{k}(1)$, a subset of $\mathrm{Ball}_{k}(4)$; that $A_i$ lies in $\tilde{R}_k^{-}(T_i)$; that $A_i$ and $0$ lie in $\tilde{R}_k(T_i)$, a subset of $\tilde{G}_k(T_i)$; and that $A_1$ and $A_2$ have equal lifts on all bonds; moreover, there is no $\mathscr{C}_i$, and $\mathscr{T}_i:=A_i$.
\item[Lemma~27.A] The property (k2), the maps $\mathrm{al}_{0}$, $\mathrm{al}_{2}$, $\varepsilon$, $m$ of Section~I.5.6, Proposition~16.3\,(a) and Proposition~27.2\,(D4) are replaced by (k2)$^{\mathrm{o}}$, the maps $\mathrm{al}_{0}$, $\mathrm{al}_{2}$, $\varepsilon$, $m$ of Proposition~48.2, Proposition~48.2\,(a) and the last sentence of Proposition~27.2\,(D4).
\item[Steps (a.1), (a.2)] Proposition~48.2\,(d), (k2)$^{\mathrm{o}}$, (o-4) and the map $R$ of Proposition~48.2 take the places of Proposition~16.3\,(d), (k2), Proposition~27.2\,(D3), (D4) and the map $R$ of Section~I.5.6.
\item[Step (a.3)] Proposition~48.3\,(c), (k2)$^{\mathrm{o}}$ and the map ${}^{\mathrm{nat}}$ of Proposition~48.3 take the places of Proposition~17.2\,(c), (k2) and the map ${}^{\mathrm{nat}}$ of Section~I.5.6.
The connection $\mathscr{T}_i'$ at which this step applies Proposition~27.4\,(d) is defined in the step as $\mathscr{T}_i':=\mathscr{A}_i+\mathscr{C}_i$; in the reading (o) it is $\mathscr{T}_i':=\mathscr{A}_i$.
Proposition~27.4\,(d) itself is the same in both readings.
\item[Step (b)] The last two assertions, on ${\delta E^{\mathrm{det}}}_k$ and on $E_k^{\mathrm{det}}$ at $\mathscr{C}_i$, are replaced by the assertion on $E_k^{\mathrm{det}}$ at $A_i$, which follows from the equality of $E^{d}_k$ at $(A_1,A_2)$ and at $(0,0)$.
\item[Steps (c), (c1)] The hypothesis at the head of Proposition~24.3\,(ga), on which the part (c0.2) of Step (c) rests, is taken in the reading (o); the definition of $\mathscr{Z}(s')$ in Definition~24.D\,(D2), on which (c1) rests, is replaced by Definition~24.D\,(D7).
\item[Step (c2)] Proposition~24.3\,(q1) is taken in the reading (o), and $\mathscr{T}_i=A_i+\mathscr{C}_i$ is replaced by $\mathscr{T}_i=A_i$.
\item[Step (c3)] Proposition~24.3\,(ga2) is taken in the reading (o).
\item[Step (c4)] (a), Proposition~27.2\,(c) and (k2) are replaced by (a) in the reading (o), (o-4) and (k2)$^{\mathrm{o}}$.
\item[Step (c5)] Definition~24.D\,(D2), (D3), (D4), Step (c4) and Proposition~24.1\,(a) are replaced by Definition~24.D\,(D7), Step (c4) in the reading (o) and Proposition~24.1\,(b).
\item[Step (c6)] Definition~24.D\,(D3), (D4), Proposition~24.3\,(al), (q5) and Proposition~24.1\,(a) are replaced by Definition~24.D\,(D7), the same clauses of Proposition~24.3 in the reading (o) and Proposition~24.1\,(b).
\item[Step (c7)] Definition~24.D\,(D2), (D3), (D4), Proposition~24.3\,(ze), (q5), Step (c4) and Proposition~24.1\,(a) are replaced by Definition~24.D\,(D7), the same clauses of Proposition~24.3 in the reading (o), Step (c4) in the reading (o) and Proposition~24.1\,(b).
\item[Step (d)] Definition~24.D\,(D5), Proposition~24.3\,(ze), (de) and Proposition~24.1\,(a) are replaced by Definition~24.D\,(D7), the same clauses of Proposition~24.3 in the reading (o) and Proposition~24.1\,(b).
\end{description}

The other facts on which Section~1 rests are the following.
None of them involves $\mathscr{C}$ or a family of activities, so that they hold unchanged in the reading (o).
\begin{itemize}
\item Proposition~26.3\,(c), the definitions in Proposition~26.4, Proposition~27.1, Proposition~27.2\,(D1), (a), (d), (e), Propositions~27.3, 27.4, 23.1 and 23.2, Sections~0.2--0.6, and Definition~24.D\,(D0); further, from Proposition~27.2\,(b), the lifts of the fermion sites and of the admissible pairs, in Section~0.4, and from Definition~24.D\,(D1), the notation $F(\psi^{\#}+v^{\#})$, in Section~0.5.
They concern index sets, parameter families, elements of $V_{k}$, elements of Grassmann algebras, single bond functions ($B_i$, $\mathscr{T}_i'$, $\mathrm{Zf}_i$), and sets formed from a polymer and a set of parameter cubes.
\item (OS), (SF), W7$^{\mathrm{ch}}$, W12$^{\mathrm{ch}}$ (in (b.6) also W1, W2, W13, W14$^{\mathrm{ch}}$; the properties W1--W17$^{\mathrm{ch}}$ of a system of primitives are those of Section~G4.0.4), Section~G4.0.4, and \cite[Lemmas~17.4, 17.6 and 17.7]{QEDI}.
They concern the primitives and their kernels at a single connection of $\tilde{G}_k$, and relations between domains of bond functions.
\item From Part~I: Sections~I.0.5, I.1.5, I.2.1, I.3.3--I.3.5, I.5.2, I.5.8, (I.1.3)--(I.1.7), (I.5.3), (I.ke1), (I.ke2), (I.Q3), (I.AP) of Section~I.3.6, (I.B3), (I.B4), (I.L0.5), and the displays of $\acute{E}_k$ and of $E^{d}_k$ in Section~I.5.7.
They concern numbers, norms, kernels, exact operators, traces and the homomorphism $\boldsymbol{\Theta}$ at a single connection, the functional $\oint_{\mathrm{rad}}$, the Grassmann integral, and the generators $\psi$, $\chi$, $\Psi$, $W$ with their signs under translations: by (I.AP), which enters (c5.6), these generators are indexed by the covering index sets and change sign under a translation by $\ell_i\cdot e_{\mu}$, for $\chi$ a translation of both points.
\item (g0) of \cite[Section~17.4.1]{QEDI}, \cite[Proposition~15.21\,(ii), (iv)]{QEDI} and \cite[Section~13.3]{QEDI}.
They concern cubes and blocks, the basis $\{\Upsilon\}$ and the map $C$, and the coefficient $\mathrm{urs}^{T}$.
\item (A1), (A5), (A8), $\varepsilon\le1/200$, $M^{\alpha_{0}}>1$, $\gamma_{\mathrm{w}}>0$ and $\kappa_{\mathrm{w}}>0$; the last two are the rows $\gamma_{\mathrm{w}}$ and $\kappa_{\mathrm{w}}$ of Proposition~14.2\,(a), and they enter (b.6) at every level.
These are inequalities between numbers.
\end{itemize}

(e.3) \emph{$F^{(i)}$ is a local family without slot at level $k+1$.}
By Proposition~48.5\,(a), $L R E_{(i)}$ lies in $K_{k+1}(\ell_i/L)$.
By Proposition~24.1\,(b), applied to the real data $(m,E_{(i)})$ and $(m,0)$ of $\mathrm{Ball}_{k}(4)$, the families $E^{*,(i)}_k(m,E_{(i)})$ and $E^{*,(i)}_k(m,0)$ lie in $K_{k+1}(\ell_i/L)$.
Next apply Proposition~24.1\,(d) with the datum $\mathscr{E}=0$, which lies in $K^{x}_{k}$ and has norm $0$.
It shows that the family $L\,{\delta E^{\mathrm{det}}}_k^{(i)}+E^{*,x,(i)}_k(m,0)$ lies in $K^{x}_{k+1}(\ell_i/L)$, and that its restriction to $\mathscr{C}=0$ is
\[
L R\,0+L E_k^{\mathrm{det},(i)}+E^{*,(i)}_k(m,0)=L E_k^{\mathrm{det},(i)}+E^{*,(i)}_k(m,0);
\]
this restriction lies in $K_{k+1}(\ell_i/L)$ by Section~I.5.4, that is, by Proposition~48.1\,(a), at level $k+1$.
Since $K_{k+1}(\ell_i/L)$ is a real vector space, $L E_k^{\mathrm{det},(i)}$ lies in it.
Hence $F^{(i)}$ lies in $K_{k+1}(\ell_i/L)$: its elements are even, and it has (k2)$^{\mathrm{o}}$.

(e.4) \emph{Clause (b) without slot.}
By (o-4) at level $k+1$, that is, by Proposition~27.2\,(c) for local families without slot, it suffices to show that
\[
\mathrm{idn}[\tilde{U}]\,F^{(1)}(\mathrm{pr}_1'\tilde{U};\mathscr{A})=F^{(2)}(\mathrm{pr}_2'\tilde{U};\mathrm{pj}^{*}\mathscr{A})
\]
for every connected $\tilde{U}$ that is admissible at level $k+1$ and all $\mathscr{A}\in\tilde{R}_{k+1}(T_1')$.
Let $Y_i$ and $\tilde{Y}$ be as in (B.0), let $A_1:=\mathscr{A}_L$, which lies in $\tilde{R}_k^{-}(T_1)$ by \cite[Lemma~17.5(a), (b)]{QEDI}, and let $A_2:=(\mathrm{pj}^{*}\mathscr{A})_L=\mathrm{pj}^{*}A_1$.
The statements (B.1) and (B.3) do not involve the slot.
By Definition~24.D\,(D6) in the reading (D7), where $E^{*}_k(m,E)=L E^{\#}_k(m,E)$, by the display of $L R E$ in Proposition~48.5, and by Section~I.3.7,
\begin{equation*}
\begin{split}
F^{(i)}(Y_i;\,\cdot\,)=\sum_{\substack{\tilde{X}\text{ connected in }\tilde{Y},\\ \mathrm{clo}(\tilde{X})=\tilde{Y}}}S_L\Bigl[&(RE_{(i)})^{\mathrm{nat}}(\mathrm{pr}_i\tilde{X};A_i)+E_k^{\mathrm{det},(i)}(\mathrm{pr}_i\tilde{X},A_i)\\
&+E^{\#,(i)}_k(m,E_{(i)})(\mathrm{pr}_i\tilde{X};A_i)\Bigr].
\end{split}
\end{equation*}
By (27.16) the three terms correspond under $\mathrm{idn}[\tilde{Y}]$; for the first term this rests on Proposition~27.2\,(D4), last sentence, since the $A_i$ lie in $\tilde{R}_k(T_i)$ and have equal lifts on all bonds.
The second term is a number.
The first term lies in $\mathrm{Grq}_i(\tilde{X})$ by (k2)$^{\mathrm{o}}$ (Proposition~48.3\,(c)) and (Gq3), and the third by the last sentence of (d.3) in the reading (o).
With (27.14) the assertion follows as in (B.6):
\[
\mathrm{idn}[\tilde{U}]\,F^{(1)}(Y_1;\mathscr{A})=\sum_{\tilde{X}}S_L\,\mathrm{idn}[\tilde{Y}]\,[\,\cdots\,]^{(1)}=\sum_{\tilde{X}}S_L\,[\,\cdots\,]^{(2)}=F^{(2)}(Y_2;\mathrm{pj}^{*}\mathscr{A}),
\]
where now $[\,\cdots\,]^{(i)}$ denotes the bracket in the last display.

\emph{Level $0$.}
Here $E_{(i)}=0$; hence $\varepsilon=0$, $m(\,\cdot\,)=0$, and $F^{(i)}=L E_0^{\mathrm{det},(i)}+E^{*,(i)}_0(m,0)$.

Every assertion of Proposition~27.6 is an identity.
Clause (a), also in its form without slot in (e), consists of equalities of numbers.
Clauses (b), (c), (d), and (b) in its form without slot in (e), are relations $\sim$ in the sense of Proposition~27.2\,(D4), that is, families of identities; in (b) and in (e) the families so related are local families at level $k+1$, by (B.0) and (e.3).
No constant is introduced in these assertions, and no bound is asserted.
They hold for all data of (TT) with $0\le k<k_m$.
\qed

\subsection*{3. Conditions on the constants}

\subsubsection*{3.1. The conditions used in the two proofs}

\begin{description}
\item[(A1), third clause] $\exp\left(-\frac{5\varepsilon v}{2}\right)\le\frac{1}{2}$ for all $v\ge-\log e_0$.
In Step (c2) it gives $\mathrm{rad}_2\le\mathrm{rad}_1$.
See Proposition~14.2\,(b), row (A1), and Section~I.8.2\,(C8).
\item[(A5)] $88M^3\cdot\exp\left(-\left(\frac{1}{4}-6\varepsilon\right)v\right)\le1$ for all $v\ge-\log e_0$.
In Step (c2) it gives $\mathrm{rad}_2\ge\mathrm{rad}_4\ge2$.
See Proposition~14.2\,(b), row (A5), and Section~I.8.2\,(C8).
\item[$0<\varepsilon\le\bar{\varepsilon}=1/200$] In Step (c2) the three exponents are ordered by
\[
-\frac{3}{4}+7\varepsilon\le-\frac{1}{4}+3\varepsilon\le-\frac{1}{4}+6\varepsilon ;
\]
the first of these inequalities is equivalent to $4\varepsilon\le\frac{1}{2}$, the second to $\varepsilon\ge0$.
The inequality $\varepsilon>0$ enters moreover (TT.3) and Step (c2), through the first clause of (A8).
Through \cite[Lemma~17.5(a), (c)]{QEDI} the bound $\varepsilon\le\frac{1}{200}$ enters (o-1) and (o-3) of Section~0.7, (B.2), (c) and (e.4).
See Proposition~14.2\,(b), row (S1), that is, $1-\alpha<\varepsilon\le\bar{\varepsilon}$, which gives $\varepsilon>0$ since $\alpha<1$; Section~I.0.5; and Section~I.8.2\,(C7).
\item[$\alpha>0$] Together with $\varepsilon\le\frac{1}{200}$ it enters \cite[Lemma~17.5(a)]{QEDI}, and through it (o-1), (o-3), (B.2) and (e.4).
See Section~I.0.5 and Section~I.8.2\,(C1).
\item[(A8), first clause] $\exp(-\varepsilon v)\le\frac{1}{4}$ for all $v\ge-\log e_0$.
At $v=-\log e_k$ and at $v=-\log e_0$ it gives $e_k^{\varepsilon}\le\frac{1}{4}$ and $e_0^{\varepsilon}\le\frac{1}{4}$; since $\varepsilon>0$, it follows that $e_0<1$, $e_k<1$ and $0<1-e_k^{\varepsilon}<1$.
In (TT.3) the inequality $e_k^{\varepsilon}\le\frac{1}{4}$ is obtained from $e_0^{\varepsilon}\le\frac{1}{4}$ and $e_k<e_0$, again because $\varepsilon>0$.
These inequalities enter (TT.3) and Step (c2).
The inequality $e_k^{\varepsilon}\le\frac{1}{4}$ enters moreover \cite[Lemma~17.4(a)]{QEDI}, which is \cite[Lemma~5.17(a)]{QEDI}, and through it (TT.3) and Step (c2) in the reading (o); and it enters \cite[Lemma~17.5(b)]{QEDI}, and through it (o-1), (o-3), (B.2) and (e.4).
See Proposition~14.2\,(b), row (A8), and Section~I.8.2\,(C8).
\item[(A8), fourth clause] $\exp\left(-\frac{v}{4}\right)+\exp(-\varepsilon v)\le1$ for all $v\ge-\log e_0$.
At $v=-\log e_k$ it gives $e_k^{\frac{1}{4}}+e_k^{\varepsilon}\le1$.
It enters \cite[Lemma~17.4(a)]{QEDI}, the inclusions of $\tilde{R}_k$ and of $\tilde{R}_k+\mathfrak{X}_k$ in $\tilde{G}_k$, and through it (TT.3) and Step (c2) in the reading (o).
See Proposition~14.2\,(b), row (A8), and Section~I.8.2\,(C8).
\item[$M^{\alpha_{0}}>1$] This follows from $M^{\alpha_{0}}\ge1+\exp(\hat{\kappa})$.
In (c6.2) it shows that the value $1$ is admitted in $\mathrm{Pd}_k$ and that $[0,1]^{\tilde{P}}$ lies in the interior of the polydisk.
See Proposition~14.2\,(b), row (M4), and Section~I.8.2\,(C5).
\item[$L\ge5$, $b=1$] In the forms $L>1$ and $b>0$ they enter (TT.3), where $e_k=L^{-\frac{1}{2}}\cdot e_{k+1}<e_0$, and (b.1), (b.6), where $b_{k}\ge b_{-}=b\cdot\left(1-\frac{1}{L}\right)>0$ and $b_{L}>0$.
The inequality $L\ge5$ itself enters \cite[Lemma~17.4(a)]{QEDI}, whose proof uses that the cube $\square^{\sim5}$ lies in $\square^{\sim L}$ for every $M$-cube $\square$, and through it (TT.3) and Step (c2) in the reading (o).
See Proposition~14.2\,(b), row (L5), and Section~I.8.2\,(C3), (C0), (C4).
\item[$L^{-\frac{3}{4}}\le\frac{1}{2}$] Together with $e_k^{\varepsilon}\le\frac{1}{4}$ it enters \cite[Lemma~17.5(b)]{QEDI}, the inclusion $L^{-\frac{3}{4}}\cdot\tilde{R}_k\subset\tilde{R}_k^{-}$, and through it (o-1), (o-3), (B.2) and (e.4); in (o-1) and (o-3) it enters through \cite[Lemma~5.18]{QEDI}, which is a special case of \cite[Lemma~17.5]{QEDI}.
It follows from row (L1) of Proposition~14.2\,(b) by Proposition~14.2\,(f); see Section~I.8.2\,(C3).
\item[$\gamma_{\mathrm{w}}>0$, $\kappa_{\mathrm{w}}=\kappa+\kappa_{0}+3>0$] In the forms $\gamma_{\mathrm{w}}\ge0$ and $\kappa_{\mathrm{w}}\ge0$ these show in (b.6) that the factor $\exp\left(-\gamma_{\mathrm{w}}\cdot d-\kappa_{\mathrm{w}}\cdot d_M(X)\right)$ of W7$^{\mathrm{ch}}$ is at most $1$; for $k\ge1$ this enters the bound of $O_2(z_1,z_2)$, for $k=0$ the bound of the kernel of $S_{0,y}(X,B)$.
Proposition~14.2\,(a) gives $\gamma_{\mathrm{w}}>0$ in the row $\gamma_{\mathrm{w}}$, and $\kappa_{0}=54\log72>0$, $\kappa\ge7\kappa_{0}+7$, $\kappa_{\mathrm{w}}=\kappa+\kappa_{0}+3$ in the rows $\kappa_{0}$, $\kappa$, $\kappa_{\mathrm{w}}$; see also Section~I.8.2\,(C4r) for $\gamma_{\mathrm{w}}$, and (C4) with (C2), (C0) for $\kappa_{\mathrm{w}}$.
\item[$L$ odd, $M$ a power of $L$] Then $M$ is odd.
Through (g0) of \cite[Section~17.4.1]{QEDI} this enters (b.1)--(b.3), (c1)\,(2), (c6.3), (B.1) and (B.3).
Through (g0) it enters also \cite[Lemma~17.5(a)]{QEDI}, and with it (o-1), (o-3), (B.2) and (e.4).
See Proposition~14.2\,(b), rows (L5), (M1), and Section~I.8.2\,(C3), (C5).
\item[$M>1$] Since $\eta_k\le1$, one has $\eta_k\le1<M$, so that the sites at distance at most $\eta_k$ from a union $Z$ of $M$-cubes lie in $Z^{\sim1}$; this enters Section~0.4, (c2), (c3)\,(1), (c5.6) and (c6.3).
See Proposition~14.2\,(b), row (M1), and Section~I.8.2\,(C5).
\item[$d_U=2L+1$, $c_{l}=2L$] These are definitions.
They are used in three forms: $c_{l}=d_U-1$, stated in (c0.1) and applied in (c1)\,(1) and (c1)\,(3); $1+d_U\le1+2d_U\le5d_U$, in (c1)\,(1) and (c1)\,(3); and $1+3d_U\le6d_U$, in (c1)\,(3), together with (c0.1).
See Section~I.8.2\,(C4).
\item[(G4.2.1) of Section~G4.0.3] $r_{*}>31L+16$.
This enters the two proofs through (TT), Proposition~26.6\,(b) (in (TT.1)), Proposition~27.1, Proposition~27.2\,(d) and Proposition~27.3.
It is a consequence of $e_K\le e_1$ (Proposition~14.2\,(b), row (e1); Section~I.8.2\,(C9)).
\item[(e1), first and third inequality, at $v=v_K$; $r\ge1$; $M$ a power of $L$] The row (e1) of Proposition~14.2\,(b) holds for all $v\ge-\log e_1$ (Proposition~14.2\,(c), Proposition~14.3\,(f)).
Since $e_K=e\cdot h^{\frac{1}{2}}\le e_1$ (Section~G4.0.2), it holds at $v=v_K=-\log e_K$, where its first and its third inequality read
\[
v_K\ge8,\qquad\left(1+\frac{3\,\mathrm{lg}}{v_K}\right)^{r}<L .
\]
The integer $r\ge2$ enters in the form $r>0$, that is, $r\ge1$.
As $\mathrm{lg}>0$ and $r>0$, the two inequalities give
\[
v_K\ge1,\qquad\left(1+\frac{\mathrm{lg}}{v_K}\right)^{r}\le\left(1+\frac{3\,\mathrm{lg}}{v_K}\right)^{r}<L ;
\]
these are the hypotheses of \cite[Proposition~19.1]{QEDI}, under which all clauses of that proposition are stated.
The proof of \cite[Proposition~19.1]{QEDI} uses, moreover, that $M$ is a power of $L$, and that $x\mapsto x^{r}$ is increasing on $(0,\infty)$, which holds since $r>0$.
The clauses (d) and (e) of \cite[Proposition~19.1]{QEDI} are the assertions (Sc2) and (Sc3) of Section~G4.0.2, and both enter the standing data (TT): by (Sc2) one has $k_m\le k_{\mathrm{top}}$, and (Sc3) together with (G4.2.1) gives $\ell_1\ge\ell_{k_m}(R)=M\cdot r_{*}>(31L+16)\cdot M$.
Hence the first and the third inequality of (e1) at $v=v_K$, the condition $r\ge1$ and the condition that $M$ is a power of $L$ enter the two proofs through (TT), that is, through (Sc2) and (Sc3).
See Proposition~14.2\,(b), rows (e1) and (M1); Section~I.8.2\,(C9), (C7), (C5); and Annex~B, Section~G4.A.6.
\end{description}

\subsubsection*{3.2. The cluster, the weights and the spin matrices in the two proofs}

The cluster, the weights and the spin matrices enter the proofs of Propositions~27.5 and 27.6 at the following places.
\begin{itemize}
\item The cluster sum $\frac{\eta_k^3}{48}\sum_{c}\sum_{g\in O_{\mathrm{cub}}}$, that is, one integral $\int\mathrm{d}x$, occurs in (Fm1), in the subtracted term of $R$ in (a.2), in the integral $\int_X\mathrm{d}\tilde{x}$ over the end points $\tilde{x}$ of the chains of (I.5.3) in (a.3), in the integrals over the blocks $\Delta_z$ in $Q_k(B)\cdot J^{*}$ in (b.2) and (b.6), in $\langle\bar{\psi},\psi\rangle_{\square}$ in (c5.4) and (c7.3), and in the rows of $T^{(1)}$, $T^{(2)}$, $\bar{T}^{(1)}$, $\bar{T}^{(2)}$ in (c6.3).
\item The factor $\left(\frac{\eta_k^3}{48}\right)^2$, times a sum over two sites and two labels, occurs in $\mathrm{al}_{2}$ in Lemma~27.A and in the kernel (27.3) of $Q_kJ^{*}S_{k,y}(X,B)JQ_k^{T}$.
\item At level $0$, the factor $\frac{1}{48}$ occurs in $\int_X\mathrm{d}x\,(\,\cdot\,)=\frac{1}{48}\sum_{c\in X}\sum_{g\in O_{\mathrm{cub}}}(\,\cdot\,)$ in (b.4), the factors $\left(\frac{1}{48}\right)^2$ and $\frac{1}{48}$ in the two traces in (b.5) and in their bounds in (b.6), and the factor $\frac{1}{2}\cdot\frac{1}{48}\sum_{c\in\tilde{\square}}\sum_{g\in O_{\mathrm{cub}}}$, with $8$ labels for each of the six directions $g e_1$, in the vertex in (c5.6).
Also at level $0$, the factor $\frac{1}{48}$ in $(\Gamma_0W)(x)$ occurs in (27.11), the factor $48$ in the kernel of $H_{0}(s'',B)$ in (c3), the factor $48$ in $\int\mathrm{d}\mu_I\,W\bar{W}$ in (d.2), and the weight $\frac{1}{48}$ of the generators $W(x,\beta,\omega)$ in the paragraph on level $0$ of the proof of Proposition~27.6\,(b).
\item The norm $|\cdot|_c$ of a $4\times4$ block, the largest row sum, occurs in (b.6): for the blocks of $O_1$, $O_2$ if $k\ge1$, and for the kernels of $F_0(B)$, $S_{0,y}(X,B)$ at pairs of fermion sites if $k=0$, where the factor $48$ of (I.ke4) is contained in $C_{w}$.
\item The pairs inside a cluster, with $d_X=\eta_k$ and without a phase, occur in (27.6), in (c7.2), in the last link of the chain in (a.3), and in the definition of $\mathrm{Pad}(\tilde{X})$, where they fall under kind (I) as the case $c=c'$.
\item The pairs $(x,x^{(3)})$ of kind (II) occur in the definition of $\mathrm{Pad}(\tilde{X})$ (Section~0.4), in (Gq0), in (27.6) within (c5.1), where $d_X=\eta_k$ and $U_B(x,x^{(3)})=u_B(c,c_{3})$, in (c5.6), and in the rows of $T^{(2)}$, $\bar{T}^{(2)}$ in (c6.3).
\item The spin matrices occur as $\mathrm{Mx}$, a general complex $4\times4$ matrix, in (Fm1); as $\mathrm{al}_{2}$, a general $4\times4$ matrix, in (a.2); as $\gamma_5$ in (b.1)--(b.6); as $\mathsf{C}$ in $\bar{H}_0(s'',B)=\mathsf{C}^{-1}\cdot H_0(s'',-B)\cdot\mathsf{C}$ in (c3), at level $0$; as $M_{a}(x)$ and $\mathsf{C}^{-1}M_{a}(x)\mathsf{C}$, $a=0,\dots,3$, in (c6.3); and as $M_{3}(x)$, $M_{3}(x^{(3)})$ in (c5.6).
\end{itemize}
No norm of an element of a Grassmann algebra is estimated in the two proofs.
Such norms, and with them the factor $\sqrt{48}$, which is the weight $\mathsf{w}_W$ of (I.5.2) at level $0$, enter only through cited statements.
In particular, by Proposition~24.1\,(a), and in the reading (o) by Proposition~24.1\,(b), the displays and the contour integrals are defined and the series (8.4) converge; and the memberships of $\mathscr{E}^{(i)+}$ in $K^{x}_{k+1}(\ell_i/L)$ in (B.0) and of $F^{(i)}$ in $K_{k+1}(\ell_i/L)$ in (e.3) are taken from Propositions~24.1\,(b), (d), 48.1\,(a) and 48.5\,(a).
The clauses (b) and (d) of Proposition~24.1 have bounds on the norms of the data among their hypotheses.

\section*{Chapter 28. The flow; the tuned pair: proofs of Propositions 28.1--28.3}

\subsection*{0. Conventions of this chapter; setting of the three proofs}

\subsubsection*{0.1. Data, scales, windows}

Throughout the chapter the flow data (FL) $=(\bar{m},e,R,N)$ of Section~G4.0.1 are fixed, together with the primitives that belong to them:
on $T_k(R_{\mathrm{top}})$, $k<k_{\mathrm{top}}$, the systems of Proposition~26.5, and on $T_k(R)$, $R<R_{\mathrm{top}}$, $k<k_{*}(R)$, their projections.
Unless another torus is named, every object is formed on the torus of side $R$; accordingly
\begin{gather*}
\ell_k:=\ell_k(R)=R/a_k,\qquad T_k:=T_k(R)=\mathbb{L}_k/\ell_k\mathbb{Z}^3,\qquad T^{\mathrm{f}}_k=T_k\times O_{\mathrm{cub}},\\
k_{*}:=k_{*}(R),\qquad K_{k}:=K_{k}(\ell_k),\qquad K^{x}_{k}:=K^{x}_{k}(\ell_k).
\end{gather*}
Section~I.6.1 gives $a_{k}=L^{k-N}$, $e_k=e\cdot a_{k}^{\frac{1}{2}}$ and $\mu_k=\bar{m}\cdot a_{k}$.
Consequently, for $0\le k\le i$,
\begin{equation}
\tag{28.1}
e_i=L^{\frac{i-k}{2}}\cdot e_k,\qquad\mu_{k+1}=L\cdot\mu_k,\qquad\ell_{k+1}=\ell_k/L.
\end{equation}
In a formula written for a general level $k$, the symbol $e_k$ stands for the running charge $q=e\cdot\eta^{\frac{1}{2}}$ at $k=0$ and for the running charge $L^{\frac{1}{2}}\cdot q$ at $k=1$;
the symbols $e_0$ and $e_1$ themselves, read as numbers, denote the ceilings of the lines (C8) and (C9) of Section~I.8.2 and never running charges (Section~I.0.5 and Section~I.2.5).
Moreover $1\le k_{*}$ (Section~I.6.2; (Sc1)).
The windows and the balls are those of Section~I.5.4:
\begin{gather*}
w^{\mathrm{m}}_k=e_k^{\frac{3}{4}-8\varepsilon},\qquad w^{\mathrm{E}}_k=e_k^{\frac{1}{4}-7\varepsilon},\\
\mathrm{Ball}_{k}(\mathrm{cf})=\left\{(m,E)\in\mathbb{C}\times K_{k}^{C}:\ |m|\le\mathrm{cf}\cdot w^{\mathrm{m}}_k,\ \|E\|_k\le\mathrm{cf}\cdot w^{\mathrm{E}}_k\right\},
\end{gather*}
and $\mathrm{Ball}^{x}_{k}(\mathrm{cf})$ is defined in the same way, with $K^{x,C}_{k}$ and $\|\cdot\|^{x}_k$ in place of $K_{k}^{C}$ and $\|\cdot\|_k$.
At $k=0$ both balls are equal to $\left\{(m,0):\ |m|\le\mathrm{cf}\cdot w^{\mathrm{m}}_0\right\}$, where $w^{\mathrm{m}}_0=q^{\frac{3}{4}-8\varepsilon}$.

\subsubsection*{0.2. The recursion (Section~I.5.7, Flows; Section~G4.0.5)}

Let $(m_{\mathrm{in}},\varepsilon^{\mathrm{v}}_{\mathrm{in}})$ be a real pair.
The data of level $k$ consist of the five entries $\left(\varepsilon^{\mathrm{v},t}_k,m_k,E_k,\mathscr{E}_k,\mathfrak{P}_k\right)$.
At level $0$ they are
\begin{equation}
\tag{rec0}
\varepsilon^{\mathrm{v},t}_0=\varepsilon^{\mathrm{v}}_{\mathrm{in}},\qquad m_{\mathrm{in}},\qquad E_0=0,\qquad\mathscr{E}_0=0,\qquad\mathfrak{P}_0=0,
\end{equation}
the second entry being the mass coupling $m_{\mathrm{in}}$, which is the meaning of the symbol $m_k$ at $k=0$.
Let $k<k_{*}$, and suppose that the data of level $k$ are defined and that $(m_k,\mathscr{E}_k)$ lies in $\mathrm{Ball}^{x}_{k}(4)$.
Then the data of level $k+1$ are
\begin{align}
\varepsilon^{\mathrm{v},t}_{k+1}&=L^3\cdot\left(\varepsilon^{\mathrm{v},t}_k+\varepsilon(E_k)+\delta\varepsilon^{0}_{k}\right),\tag{rec1}\\
m_{k+1}&=L\cdot\left(m_k+m(E_k)\right),\tag{rec2}\\
E_{k+1}&=L R E_k+L E_k^{\mathrm{det}}+E^{*}_k(m_k,E_k),\tag{rec3}\\
\mathscr{E}_{k+1}&=L R\mathscr{E}_k+L\,{\delta E^{\mathrm{det}}}_k+E^{*,x}_k(m_k,\mathscr{E}_k),\tag{rec4}\\
\mathfrak{P}_{k+1}(Y,\mathscr{C})&=\sum_{X\in D_k:\ \bar{X}^{L}=LY}\left(\mathfrak{P}_k(X,\mathscr{C}_L)+E_k^{\mathrm{det}}(X,\mathscr{C}_L)\right)\qquad\left(Y\in D_{k+1},\ \mathscr{C}\in\mathfrak{X}_{k+1}\right),\tag{rec5}
\end{align}
The right-hand sides of (rec1)--(rec5) are formed with the step map of level $k$ on $T_k$ (Definition~24.D) and with the system of primitives of (FL) at level $k$.
The phrase ``the trajectory exists up to level $k_{\mathrm{up}}$'' means that $(m_k,\mathscr{E}_k)$ lies in $\mathrm{Ball}^{x}_{k}(4)$ for every $k<k_{\mathrm{up}}$.
The pair is called \emph{controlled on $R$} if its trajectory on $R$ exists up to $k_{*}$ and $|m_k|\le w^{\mathrm{m}}_k$ for $0\le k\le k_{*}$.

\emph{The terms multiplied by the two couplings.}
In the density $\mathscr{F}^{x}_k$ of Section~I.5.7 the couplings $m_k$ and $\varepsilon^{\mathrm{v},t}_k$ occur in the factor $\exp\left\{-m_k\langle\bar{\psi},\psi\rangle-\varepsilon^{\mathrm{v},t}_k\cdot n_k\right\}$, where $n_k=\ell_k^3$.
By (I.1.3) and Section~I.3.6,
\[
\langle\bar{\psi},\psi\rangle=\int\mathrm{d}x\,\bar{\psi}(x)\cdot\psi(x)=\eta_k^3\cdot\frac{1}{48}\sum_{c\in T_k}\sum_{g\in O_{\mathrm{cub}}}\sum_{\beta=1}^{4}\bar{\psi}_{\beta}(c,g)\,\psi_{\beta}(c,g).
\]
Thus a single coefficient multiplies the $48$ sites of every cluster and the four spin components, that is, all $192$ components that belong to a photon site.
At $k=0$ one has $\psi=\Psi$ and $\eta_0=1$, and the mass term is $-m_{\mathrm{in}}\cdot\frac{1}{48}\sum_{c}\sum_{g\in O_{\mathrm{cub}}}\bar{\Psi}(c,g)\cdot\Psi(c,g)$.

\emph{The two extracted numbers} (Section~I.5.6).
For $E\in K_{k}^{C}$, a small polymer $X$ and an $M$-cube $\square$,
\begin{align*}
\mathrm{al}_{0}(E,X)&=\operatorname{Vol}(X)^{-1}\cdot E_{00}(X;0),&
\mathrm{al}_{2}(E,X)&=\operatorname{Vol}(X)^{-1}\int\mathrm{d}x\int\mathrm{d}x'\,E_{20}(X;0;x,x'),\\
\varepsilon({E})&=-\sum_{X\text{ small},\ X\supset\square}\mathrm{al}_{0}(E,X),&
m({E})\,\mathbf{1}&=-\sum_{X\text{ small},\ X\supset\square}\mathrm{al}_{2}(E,X).
\end{align*}
Here $\mathrm{al}_{2}(E,X)$ is a $4\times4$ matrix and $\mathbf{1}$ is the $4\times4$ identity matrix;
each of the two integrals in $\mathrm{al}_{2}$ equals $\eta_k^3\cdot\frac{1}{48}$ times the sum over the photon sites $c$ of one lift of $X$ and over the $48$ labels $g\in O_{\mathrm{cub}}$.
For $\mathscr{E}\in K^{x,C}_{k}$ the same definitions are made with $\mathscr{E}_{00}(X;0,0)$ and $\mathscr{E}_{20}(X;0,0;x,x')$.
That the two sums do not depend on $\square$, and that the second one is a multiple of $\mathbf{1}$, is part of Proposition~48.2\,(a) for $E\in K_{k}^{C}$ and of Proposition~16.3\,(a) for $\mathscr{E}\in K^{x,C}_{k}$.

\subsubsection*{0.3. The data of the statements applied}

By Proposition~26.6\,(b), for every $k<k_{*}$ the data $(k,\ell_k,\mu_k,e_k)$, together with the system of primitives of (FL) at level $k$, are (OS); explicitly,
\begin{equation}
\tag{28.2}
\begin{aligned}
&\ell_k\in\mathrm{Cl}\ \text{ and }\ \ell_k/L=\ell_{k+1}\in\mathrm{Cl};\qquad\mu_k\in(0,1/L];\qquad e_k>0,\quad e_{k+1}=L^{\frac{1}{2}}\cdot e_k\le e_0;\\
&\ell_k\ge M\cdot r_k;\qquad\text{the system has W1--W17$^{\mathrm{ch}}$ for $\mathrm{par}$.}
\end{aligned}
\end{equation}
Since $k_{*}\ge1$, every level $k\le k_{*}$ is either a level $k<k_{*}$ or the level $(k_{*}-1)+1$.
Therefore (28.2) yields
\begin{equation}
\tag{28.3}
\ell_k\in\mathrm{Cl}\quad\text{and}\quad0<e_k\le e_0\qquad\text{for }0\le k\le k_{*},
\end{equation}
and for $k<k_{*}$ even $e_k<e_{k+1}\le e_0$.

In the following table, the second column lists the data that the statement in the first column requires, and the third column the place from which these data are obtained.

\noindent
\begin{tabular}{p{0.26\textwidth}p{0.42\textwidth}p{0.2\textwidth}}
\hline
Statement applied & Data required & Origin of the data\\
\hline
Definition~24.D and Proposition~24.1\,(a)--(e), (i), at a level $k<k_{*}$ &
(ST1)$^{\mathrm{ch}}$, which is contained in (OS) (Section~I.1.8, Dictionary); for \cite[Lemma~17.4(a)]{QEDI} and \cite[Lemma~17.5(a)--(c)]{QEDI}, on which the closing assertions of Definition~24.D\,(D6) on $\mathscr{A}_L$ and $\mathscr{C}_L$ rest: $\ell_k,\ell_k/L\in\mathrm{Cl}$, $e_{k+1}=L^{\frac{1}{2}}\cdot e_k\le e_0$ &
(28.2)\\
\hline
Proposition~13\,(d5), (d6), at a level $k<k_{*}$ (and (d1), which requires only the first three of these conditions, that is (SF), and the family) &
(SF+): $\ell_k\in\mathrm{Cl}$, $\mu_k\in(0,1]$, $e_k\in(0,e_0]$, $\ell_k/L\in\mathrm{Cl}$, $\ell_k\ge M\cdot r_k$, $e_{k+1}\le e_0$; a family (p5) with W1, W2, W7, W12, W13, W14 for $\mathrm{par}$ &
(28.2); $(0,1/L]\subset(0,1]$; $e_k<e_{k+1}\le e_0$; the family (p5) is that of the system in (28.2)\\
\hline
Proposition~16.3\,(a) and Proposition~48.2\,(a), (b), at a level $k<k_{*}$; Proposition~48.1\,(a), (d), at a level $k\le k_{*}$ &
(S-1): $k\ge0$, $\ell_k\in\mathrm{Cl}$, $e_k\in(0,e_0]$ &
(28.3)\\
\hline
Proposition~48.5\,(a), (b), at a level $k<k_{*}$ &
(S-2): (S-1), $\ell_k/L\in\mathrm{Cl}$, $e_{k+1}\le e_0$ &
(28.2), (28.3)\\
\hline
Proposition~48.6\,(b) &
(S-4), which is (FL) without the primitives; systems of the levels $k<k_{*}$ on $T_k(R)$ with W1--W8, W10--W17 for $\mathrm{par}$ &
(FL); (28.2)\\
\hline
Propositions~27.2 and~27.6 &
(TT): (FL) with $R<R_{\mathrm{top}}$ and a level $k\le k_{*}(R)$ (for Proposition~27.6: $k<k_{*}(R)$) &
Section~3\\
\hline
\end{tabular}

The spaces of level $k+1$ that occur in these statements are formed with the side $\ell_k/L$ and the charge $L^{\frac{1}{2}}\cdot e_k$;
by (28.1) they are $K_{k+1}=K_{k+1}(\ell_{k+1})$ and $K^{x}_{k+1}=K^{x}_{k+1}(\ell_{k+1})$, formed with $e_{k+1}$.

\subsubsection*{0.4. Arithmetic of the windows}

The following conditions on the constants enter:
$0<\varepsilon\le\bar{\varepsilon}=1/200$ ((S1), (C0));
$L>1$, $L$ being an odd integer $\ge5$ (Section~I.0.5; (L5));
$c_{\mathrm{fl}}\cdot L^{-\frac{3}{8}+6\bar{\varepsilon}}\le\frac{1}{16}$ ((L2));
$L^{-\frac{5}{8}}\le\frac{1}{4}$ (Proposition~14.2\,(f));
$\exp(-\varepsilon v)\le\frac{1}{4}$ for all $v\ge-\log e_0$ (the first inequality of (A8));
$c_{\mathrm{fl}}\ge1$ ((C0));
and, for (u6), the row (e2), through (Sc6).

The first group of facts holds for $0\le k\le k_{*}$:
\begin{equation}
\tag{u0}
0<e_k\le e_0,\qquad4\,e_k^{\varepsilon}\le1,\qquad e_k<1.
\end{equation}
Indeed, the first assertion is (28.3);
hence $v_k=-\log e_k\ge-\log e_0$, so that $e_k^{\varepsilon}=\exp(-\varepsilon v_k)\le\frac{1}{4}$, and $e_k<1$ because $\varepsilon>0$.

By (28.1), $w^{\mathrm{m}}_{k+1}/w^{\mathrm{m}}_k=L^{\frac{1}{2}\left(\frac{3}{4}-8\varepsilon\right)}=L^{\frac{3}{8}-4\varepsilon}$ and $w^{\mathrm{E}}_{k+1}/w^{\mathrm{E}}_k=L^{\frac{1}{2}\left(\frac{1}{4}-7\varepsilon\right)}=L^{\frac{1}{8}-\frac{7\varepsilon}{2}}$.
It follows that, for $0\le k<k_{*}$,
\begin{align}
L^{-1}\cdot\frac{w^{\mathrm{m}}_{k+1}}{w^{\mathrm{m}}_k}&=L^{-\frac{5}{8}-4\varepsilon}\le L^{-\frac{5}{8}}\le\frac{1}{4},\tag{u1}\\
c_{\mathrm{fl}}\cdot e_k^{\frac{1}{2}}\cdot\frac{w^{\mathrm{E}}_k}{w^{\mathrm{m}}_k}&=c_{\mathrm{fl}}\cdot e_k^{\frac{1}{2}+\frac{1}{4}-7\varepsilon-\frac{3}{4}+8\varepsilon}=c_{\mathrm{fl}}\cdot e_k^{\varepsilon},\tag{u2}\\
c_{\mathrm{fl}}\cdot L^{-\frac{1}{4}+2\varepsilon}\cdot\frac{w^{\mathrm{E}}_k}{w^{\mathrm{E}}_{k+1}}&=c_{\mathrm{fl}}\cdot L^{-\frac{1}{4}+2\varepsilon-\frac{1}{8}+\frac{7\varepsilon}{2}}=c_{\mathrm{fl}}\cdot L^{-\frac{3}{8}+\frac{11\varepsilon}{2}}\le c_{\mathrm{fl}}\cdot L^{-\frac{3}{8}+6\bar{\varepsilon}}\le\frac{1}{16},\tag{u3}\\
\frac{e_{k+1}^{\frac{1}{4}-2\varepsilon}}{w^{\mathrm{E}}_{k+1}}&=e_{k+1}^{5\varepsilon}\le e_{k+1}^{\varepsilon},\tag{u4}\\
\frac{e_k^{\frac{1}{4}-6\varepsilon}}{w^{\mathrm{E}}_{k+1}}&=\left(\frac{e_k}{e_{k+1}}\right)^{\frac{1}{4}-6\varepsilon}\cdot e_{k+1}^{\frac{1}{4}-6\varepsilon-\frac{1}{4}+7\varepsilon}=L^{-\left(\frac{1}{8}-3\varepsilon\right)}\cdot e_{k+1}^{\varepsilon}\le e_{k+1}^{\varepsilon}.\tag{u5}
\end{align}
In (u3) one uses $\frac{11\varepsilon}{2}\le6\bar{\varepsilon}$ and $L\ge1$, in (u4) the inequality $e_{k+1}<1$, and in (u5) the inequality $\frac{1}{8}-3\varepsilon>0$.
Finally, for $0\le k\le k_{*}$,
\begin{equation}
\tag{u6}
\begin{aligned}
&\mathrm{sm}_k=(c_{\mathrm{fl}}+2)\cdot e_k^{\varepsilon}\le\frac{1}{16};\\
&\text{hence}\quad c_{\mathrm{fl}}\cdot e_k^{\varepsilon}\le\frac{1}{16},\qquad2\,e_k^{\varepsilon}\le\frac{1}{16},\qquad e_k\le e_k^{\varepsilon}\le\frac{1}{48}.
\end{aligned}
\end{equation}
To see this, note that by (Sc6) of Section~G4.0.2, that is, by Proposition~14.2\,(c), (f), with the row (e1) at $v=v_K$ supplied by the second assertion of Proposition~14.3\,(f), the condition (fl) of Section~G4.0.3 holds under (FL):
\[
\sum_{k\le k_{*}}\mathrm{sm}_k\le\frac{1}{16}.
\]
In detail, (28.1) gives
\[
\sum_{k\le k_{*}}e_k^{\varepsilon}=e_{*}^{\varepsilon}\sum_{i=0}^{k_{*}}L^{-\frac{i\varepsilon}{2}}\le e_{*}^{\varepsilon}\cdot\left(1-L^{-\frac{\varepsilon}{2}}\right)^{-1};
\]
by (Sc1) one has $k_{*}\le K$, so that $e_{*}\le e_K\le e_1$ and $v_{*}\ge-\log e_1$;
and the row (e2) at $v=v_{*}$ reads $(c_{\mathrm{fl}}+2)\cdot e_{*}^{\varepsilon}\le\frac{1}{16}\left(1-L^{-\frac{\varepsilon}{2}}\right)$.
All the terms $\mathrm{sm}_k$ are positive, so each of them is at most $\frac{1}{16}$.
Since $c_{\mathrm{fl}}\ge1$, one has $3\,e_k^{\varepsilon}\le\mathrm{sm}_k$, as well as $c_{\mathrm{fl}}\cdot e_k^{\varepsilon}\le\mathrm{sm}_k$ and $2\,e_k^{\varepsilon}\le\mathrm{sm}_k$.
Finally, $0<e_k<1$ and $0<\varepsilon\le1$ give $e_k\le e_k^{\varepsilon}$.
The numbers $\frac{1}{4}$, $\frac{1}{16}$, $\frac{1}{48}$ in (u0)--(u6) do not depend on $(\bar{m},e,R,N,k)$.

\subsubsection*{0.5. Notation}

The following letters have a meaning local to this chapter.
$\mathrm{Vfl}$ is a Banach space of sequences and $\mathrm{Bfl}$ a closed subset of it, $|\cdot|_{\mathrm{fl}}$ is the norm of $\mathrm{Vfl}$, $\mathrm{Ffl}$ is a map of $\mathrm{Bfl}$, and $\mathrm{zfl}$, $\widetilde{\mathrm{zfl}}$ are elements of $\mathrm{Bfl}$;
$\mathrm{qfl}^{m}_k$ and $\mathrm{qfl}^{E}_k$ are the quotients of a coupling by its window, and $\mathrm{dfl}^{m}_k$, $\mathrm{dfl}^{E}_k$ the quotients of the difference of two couplings by the window;
the upper index $\mathrm{traj}$ (in 2.6) marks the entries of a trajectory.
The mass coupling of level $0$ is never written with the letter $m$ and the bare index $0$: it is $m_{\mathrm{in}}$, or the term of index $0$ of a sequence (Section~G4.0.5).

Labels defined outside this chapter have the following meaning.
(Sc0)--(Sc6) are the facts of Section~G4.0.2;
(fl) and (T1) are the conditions of Section~G4.0.3;
(S1), (L1), (L2), (L5), (M1), (A8) and (e1)--(e8) are the rows with these labels in Proposition~14.2\,(b);
(C0)--(C11) are the lines of Section~I.8.2;
(8.6)$^{\mathrm{ch}}$ is the display of Definition~24.D\,(D6);
and (6.1) is the display of Proposition~28.3\,(a).

\subsection*{1. Proposition~28.1 (the flow on one torus)}

\begin{proof}[Proof of Proposition~28.1]
\leavevmode
\paragraph{1.1. The hypothesis.}
Fix $k$ with $0\le k\le k_{*}$, and assume the hypothesis of Proposition~28.1:
\begin{equation}
\tag{Hyp$_k$}
|m_i|\le w^{\mathrm{m}}_i\qquad\text{for every }i<k\text{ for which the data of level }i\text{ are defined},
\end{equation}
the numbers $m_i$, $i<k$, being produced by the recursion one after the other.
That all these data are in fact defined is one of the conclusions.

\paragraph{1.2. Claim.}
For every $j$ with $0\le j\le k$ the following hold:
\begin{itemize}
\item[(i)$_j$] the data of level $j$ are defined, and $(m_i,\mathscr{E}_i)$ lies in $\mathrm{Ball}^{x}_{i}(1)$ for every $i<j$;
\item[(ii)$_j$] $\varepsilon^{\mathrm{v},t}_j$ and $m_j$ are real; $\mathscr{E}_j$ lies in $K^{x}_{j}$; $E_j=\mathscr{E}_j|_0$, and $E_j$ lies in $K_{j}$; for every $X\in D_j$, $\mathfrak{P}_j(X,\cdot\,)$ is a complex function on $\mathfrak{X}_j$, and $\mathfrak{P}_j(X,0)=0$;
\item[(iii)$_j$] $\|\mathscr{E}_j\|^{x}_j\le\frac{32}{15}\,e_j^{\varepsilon}\cdot w^{\mathrm{E}}_j$.
\end{itemize}
The claim is proved by induction on $j$.

\paragraph{1.3. The case $j=0$.}
The data of level $0$ are those of (rec0), and there is no index $i<0$.
The numbers $\varepsilon^{\mathrm{v}}_{\mathrm{in}}$ and $m_{\mathrm{in}}$ are real.
The zero family lies in $K^{x}_{0}$; its restriction $0|_0=0$ lies in the space $K_{k}$ at $k=0$, and it is $E_0$.
The function $\mathfrak{P}_0(X,\cdot\,)=0$ is a complex function on $\mathfrak{X}_0$ that vanishes at $0$.
Finally $\|\mathscr{E}_0\|^{x}_0=0$.

\paragraph{1.4. The step from $j$ to $j+1$ ($0\le j<k$).}
Assume (i)$_j$, (ii)$_j$ and (iii)$_j$.

\emph{(1) The couplings of level $j$.}
By (i)$_j$ the data of level $j$ are defined, so that (Hyp$_k$) gives $|m_j|\le w^{\mathrm{m}}_j$.
By (iii)$_j$ and (u0),
\[
\|\mathscr{E}_j\|^{x}_j\le\frac{32}{15}\cdot\frac{1}{4}\,w^{\mathrm{E}}_j=\frac{8}{15}\,w^{\mathrm{E}}_j\le w^{\mathrm{E}}_j.
\]
Hence $(m_j,\mathscr{E}_j)$ lies in $\mathrm{Ball}^{x}_{j}(1)$, which is a subset of $\mathrm{Ball}^{x}_{j}(4)$, and $j<k\le k_{*}$.
By Section~0.2 the data of level $j+1$ are therefore given by (rec1)--(rec5).
Together with (i)$_j$ it follows that $(m_i,\mathscr{E}_i)$ lies in $\mathrm{Ball}^{x}_{i}(1)$ for every $i<j+1$;
thus the trajectory exists up to level $j+1$.

\emph{(2) The right-hand sides of (rec1)--(rec5) are meaningful.}
Apply Proposition~48.6\,(b) with $j+1$ as the level up to which the trajectory exists; here $j+1\le k\le k_{*}$, and the hypothesis of that statement is the last conclusion of (1).
It shows that all five entries of level $j+1$ are defined.
For the fifth entry this means, in view of the range of the variables in (rec5), that $\mathfrak{P}_{j+1}(Y,\cdot\,)$ is a complex function on $\mathfrak{X}_{j+1}$ for every $Y\in D_{j+1}$.
In detail, the entries are formed as follows.

\emph{(rec1), (rec2).}
By (ii)$_j$, $E_j=\mathscr{E}_j|_0$, that is, $E_j(X;\mathscr{A})=\mathscr{E}_j(X;\mathscr{A},0)$ (Section~I.5.4).
Hence the kernels $E_{j,00}(X;0)$ and $E_{j,20}(X;0;x,x')$ are $\mathscr{E}_{j,00}(X;0,0)$ and $\mathscr{E}_{j,20}(X;0,0;x,x')$, and
\[
\mathrm{al}_{0}(E_j,X)=\mathrm{al}_{0}(\mathscr{E}_j,X),\qquad\mathrm{al}_{2}(E_j,X)=\mathrm{al}_{2}(\mathscr{E}_j,X)
\]
for every small polymer $X$.
So $\varepsilon(E_j)$ and $m(E_j)$ are defined by Proposition~16.3\,(a): the two sums of Section~0.2 do not depend on $\square$, and the second one is a multiple of $\mathbf{1}$.
For the number $\delta\varepsilon^{0}_{j}$ see (4).

\emph{(rec3), (rec4).}
Let $(\mathscr{A},\mathscr{C})\in\tilde{R}_{j+1}\times\mathfrak{X}_{j+1}$.
The closing assertions of Definition~24.D\,(D6) on $\mathscr{A}_L$ and $\mathscr{C}_L$, which rest on \cite[Lemma~17.5(a)--(c)]{QEDI} and \cite[Lemma~17.4(a)]{QEDI}, give the following:
$\mathscr{A}_L$ lies in $\tilde{R}_j^{-}$,
$\mathscr{C}_L$ lies in $\mathfrak{X}_j$,
and $\mathscr{A}_L+\mathscr{C}_L$ and $\mathscr{C}_L$ lie in $\tilde{G}_j$.
Here $\tilde{R}_j^{-}$ is a subset of $\tilde{R}_j$.
Indeed (Section~I.5.3; \cite[Definition~5.16(c)]{QEDI}), the set $\tilde{R}_j^{-}$ is defined as $\tilde{R}_j$, except that the summand $A_1$ of the splitting is taken in $\left(1-e_j^{\varepsilon}\right)R_j$ instead of $R_j$;
the set $R_j$ is defined by strict upper bounds on the three norms, and $0<1-e_j^{\varepsilon}<1$ by (u0), so that $\left(1-e_j^{\varepsilon}\right)R_j\subset R_j$.
The case $\mathscr{C}=0$ of the membership of $\mathscr{A}_L+\mathscr{C}_L$ in $\tilde{G}_j$ shows that $\mathscr{A}_L$ lies in $\tilde{G}_j$ as well.
Consequently $L R E_j$ and $L R\mathscr{E}_j$ are formed according to Section~I.5.6.
The activities $E_j^{\mathrm{det}}$ and ${\delta E^{\mathrm{det}}}_j$ are formed according to Section~I.5.7 with the family (p5) of the system in (28.2):
$\acute{E}_j(X,\cdot\,)$, and hence $E^{d}_j(X,\cdot\,)$, are functions on $\tilde{G}_j$ (Proposition~13\,(d1)), and $L E_j^{\mathrm{det}}$, $L\,{\delta E^{\mathrm{det}}}_j$ are then formed according to Section~I.5.6.
By Proposition~48.1\,(a), (ii)$_j$ and (1),
\[
\|E_j\|_j=\left\|\mathscr{E}_j|_0\right\|_j\le\|\mathscr{E}_j\|^{x}_j\le w^{\mathrm{E}}_j;
\]
thus $(m_j,E_j)$ lies in $\mathrm{Ball}_{j}(1)$, a subset of $\mathrm{Ball}_{j}(4)$, which is the set of the data of the step map without slot (Definition~24.D\,(D1) in the reading (D7));
on this set $E^{*}_j$ is defined (Proposition~24.1\,(b), first and second sentences).
Likewise $(m_j,\mathscr{E}_j)$ lies in $\mathrm{Ball}^{x}_{j}(4)$, the set of the data of the step map (Definition~24.D\,(D1)), on which all displays of Definition~24.D are defined, $E^{*,x}_j(m_j,\mathscr{E}_j)$ among them (the first assertion of Proposition~24.1\,(a)).

\emph{(rec5).}
Let $\mathscr{C}\in\mathfrak{X}_{j+1}$.
Then $\mathscr{C}_L$ lies in $\mathfrak{X}_j$ and in $\tilde{G}_j$.
By (ii)$_j$, $\mathfrak{P}_j(X,\cdot\,)$ is a function on $\mathfrak{X}_j$;
and $E_j^{\mathrm{det}}(X,\cdot\,)=E^{d}_j(X,\cdot\,)-E^{d}_j(X,0)$ is a function on $\tilde{G}_j$, a set that contains $0$.
As $D_j$ is finite, the sum in (rec5) is finite.

\emph{(3) The recursion is the step map.}
Proposition~16.3\,(a) and $E_j=\mathscr{E}_j|_0$ give
\[
\varepsilon(\mathscr{E}_j)=\varepsilon\left(\mathscr{E}_j|_0\right)=\varepsilon(E_j),\qquad m(\mathscr{E}_j)=m\left(\mathscr{E}_j|_0\right)=m(E_j),
\]
and both numbers are real, since $\mathscr{E}_j$ lies in $K^{x}_{j}$.
Therefore (rec1), (rec2), (rec4) and (rec5) are the display (8.6)$^{\mathrm{ch}}$ of Definition~24.D\,(D6), taken at $(\varepsilon^{\mathrm{v}},m,\mathscr{E})=\left(\varepsilon^{\mathrm{v},t}_j,m_j,\mathscr{E}_j\right)$ and with the family $\mathfrak{P}_j$:
\begin{equation}
\tag{28.4}
\left(\varepsilon^{\mathrm{v},t}_{j+1},m_{j+1},\mathscr{E}_{j+1}\right)=\left(\varepsilon^{\mathrm{v},+},m^{+},\mathscr{E}^{+}\right),\qquad\mathfrak{P}_{j+1}\ \text{is the family of (8.6)$^{\mathrm{ch}}$ formed from }\mathfrak{P}_j.
\end{equation}
By (1) and (ii)$_j$ these data are in the situation of Proposition~24.1\,(c):
$\varepsilon^{\mathrm{v}}$ and $m$ are real, $|m|\le w^{\mathrm{m}}_j$, $\mathscr{E}$ lies in $K^{x}_{j}$ with $\|\mathscr{E}\|^{x}_j\le w^{\mathrm{E}}_j$ ($\mathscr{E}=0$ if $j=0$), and $\mathfrak{P}_j$ is a family of complex functions on $\mathfrak{X}_j$.
The two remaining data of Proposition~24.1\,(c), namely $A'\in\mathrm{Dom}_{j+1}$ and $\mathscr{C}'\in\mathfrak{X}_{j+1}$, occur neither in (8.6)$^{\mathrm{ch}}$ nor in the assertions of Proposition~24.1\,(d).
Such data exist: one may take $\mathscr{C}'=0$, and $A'=0$, because $H_{j+1}0=0$ ($H_{j+1}$ is linear by \cite[Proposition~15.5]{QEDI}) and $0$ lies in $\frac{1}{4}\tilde{R}_{j+1}$, while $\mathrm{Dom}_{j+1}$ is the set of the real $A'$ with $H_{j+1}A'\in\frac{1}{4}\tilde{R}_{j+1}$ (Section~I.5.3; \cite[Definition~5.16(g)]{QEDI}).

\emph{(4) Proof of (ii)$_{j+1}$.}
The number $\delta\varepsilon^{0}_{j}=-n_j^{-1}\log\mu_I\left(\mathrm{Box}_{j}\right)$ (Section~I.5.7; \cite[Section~5.2.5]{QEDI}) is real.
Indeed, $\mu_I$ is a probability measure on $V_j$ with a positive density, and
\[
\mathrm{Box}_{j}=\left\{\zeta:\ |\zeta(\Upsilon)|\le p_{0,j}\text{ for all }\Upsilon\right\},\qquad p_{0,j}=v_j^{p_0}>0,
\]
where $v_j>0$ by (u0); hence $0<\mu_I\left(\mathrm{Box}_{j}\right)\le1$.
With (3) and (ii)$_j$, the right-hand sides of (rec1) and (rec2) are real.
By Proposition~24.1\,(d), $\mathscr{E}_{j+1}=\mathscr{E}^{+}$ lies in $K^{x}_{j+1}$, and
\[
\mathscr{E}_{j+1}|_0=L R\left(\mathscr{E}_j|_0\right)+L E_j^{\mathrm{det}}+E^{*}_j\left(m_j,\mathscr{E}_j|_0\right)=L R E_j+L E_j^{\mathrm{det}}+E^{*}_j(m_j,E_j)=E_{j+1}
\]
by (ii)$_j$ and (rec3).
By Proposition~48.1\,(a), the restriction $\mathscr{E}_{j+1}|_0$ of the element $\mathscr{E}_{j+1}$ of $K^{x}_{j+1}$ lies in $K_{j+1}$.
At $\mathscr{C}=0$ one has $\mathscr{C}_L=0$ (Section~I.3.7: $A_{L}(b)=L^{-\frac{1}{2}}\cdot A(b/L)$), $\mathfrak{P}_j(X,0)=0$ by (ii)$_j$, and $E_j^{\mathrm{det}}(X,0)=E^{d}_j(X,0)-E^{d}_j(X,0)=0$ (Section~I.5.7);
so (rec5) gives $\mathfrak{P}_{j+1}(Y,0)=0$.
Together with (2) this is (ii)$_{j+1}$.
The same assertions also follow from Proposition~48.6\,(b), applied as in (2), and the assertion on $\mathfrak{P}_{j+1}(\cdot\,,0)$ also from the last sentence of Proposition~24.1\,(c).

\emph{(5) Proof of (iii)$_{j+1}$.}
By Proposition~24.1\,(d) and (28.4),
\[
\|\mathscr{E}_{j+1}\|^{x}_{j+1}\le c_{\mathrm{fl}}\cdot L^{-\frac{1}{4}+2\varepsilon}\cdot\|\mathscr{E}_j\|^{x}_j+e_{j+1}^{\frac{1}{4}-2\varepsilon}+e_j^{\frac{1}{4}-6\varepsilon}.
\]
Dividing by $w^{\mathrm{E}}_{j+1}$ and using (u3), (u4) and (u5) at $k=j$, one obtains
\begin{equation}
\tag{28.5}
\frac{\|\mathscr{E}_{j+1}\|^{x}_{j+1}}{w^{\mathrm{E}}_{j+1}}\le\frac{1}{16}\cdot\frac{\|\mathscr{E}_j\|^{x}_j}{w^{\mathrm{E}}_j}+2\,e_{j+1}^{\varepsilon}.
\end{equation}
By (iii)$_j$ and $e_j\le e_{j+1}$, the right-hand side is at most
\[
\left(\frac{1}{16}\cdot\frac{32}{15}+2\right)e_{j+1}^{\varepsilon}=\left(\frac{2}{15}+\frac{30}{15}\right)e_{j+1}^{\varepsilon}=\frac{32}{15}\,e_{j+1}^{\varepsilon}.
\]
To sum up, (1) together with (2) is (i)$_{j+1}$, (4) is (ii)$_{j+1}$, and (5) is (iii)$_{j+1}$.

\paragraph{1.5. The first step written out ($j=0$).}
The data are $(m_{\mathrm{in}},0)$, with $|m_{\mathrm{in}}|\le w^{\mathrm{m}}_0=q^{\frac{3}{4}-8\varepsilon}$.
The operations of Section~I.5.6 send the zero family to $0$:
$\mathrm{al}_{0}(0,X)=0$ and $\mathrm{al}_{2}(0,X)=0$, so that $\varepsilon(0)=0$, $m(0)=0$ and $R\,0=0$;
and since $\boldsymbol{\Theta}_{X,\mathscr{T}}$ and $S_L$ are homomorphisms of algebras, $0^{\mathrm{nat}}=0$ and $L\,0=0$.
Hence $\varepsilon(\mathscr{E}_0)=\varepsilon(E_0)=0$, $m(\mathscr{E}_0)=m(E_0)=0$, $L R\mathscr{E}_0=0$ and $L R E_0=0$ ((I.L0.21)), and (rec1)--(rec5) become ((I.L0.27); Definition~24.D\,(D6), case $k=0$)
\begin{gather*}
\varepsilon^{\mathrm{v},t}_1=L^3\cdot\left(\varepsilon^{\mathrm{v}}_{\mathrm{in}}+\delta\varepsilon^{0}_{0}\right),\qquad m_1=L\cdot m_{\mathrm{in}},\qquad E_1=L E_0^{\mathrm{det}}+E^{*}_0(m_{\mathrm{in}},0),\\
\mathscr{E}_1=L\,{\delta E^{\mathrm{det}}}_0+E^{*,x}_0(m_{\mathrm{in}},0),\qquad\mathfrak{P}_1(Y,\mathscr{C})=\sum_{X\in D_0:\ \bar{X}^{L}=LY}E_0^{\mathrm{det}}(X,\mathscr{C}_L).
\end{gather*}
The activities $E_0^{\mathrm{det}}$ and ${\delta E^{\mathrm{det}}}_0$ are formed with the function $\acute{E}_0$ of Section~I.5.7, whose traces are traces on $\mathfrak{S}^{\mathrm{ch}}(1,\ell_0)$, a space of dimension $192\,n_0$; by (I.ke2),
\[
\operatorname{Tr}O=\int\mathrm{d}x\,\operatorname{tr}O(x,x)=\frac{1}{48}\sum_{c\in T_0}\sum_{g\in O_{\mathrm{cub}}}\operatorname{tr}O\left((c,g),(c,g)\right).
\]
In the identity of Proposition~24.1\,(c) at $k=0$ the left-hand side is an integral over the $192\,n_0$ pairs of generators $\Psi$ on $T^{\mathrm{f}}_0$, of weight $\frac{1}{48}$, and the right-hand side is an element of $\mathrm{Gr}[\Psi']$ with $4\,n_1$ pairs of generators.
The bound of Proposition~24.1\,(d) at $k=0$, together with (u4) and (u5) at $k=0$, gives
\[
\|\mathscr{E}_1\|^{x}_1\le\left(L^{\frac{1}{2}}\cdot q\right)^{\frac{1}{4}-2\varepsilon}+q^{\frac{1}{4}-6\varepsilon}\le2\left(L^{\frac{1}{2}}\cdot q\right)^{\varepsilon}\cdot w^{\mathrm{E}}_1;
\]
this is (28.5) at $j=0$, where the first term on the right-hand side vanishes.
In every other respect the step $j=0$ is as in 1.4.

\paragraph{1.6. Conclusion.}
By (i)$_k$ the data of level $k$ are defined, and $(m_i,\mathscr{E}_i)$ lies in $\mathrm{Ball}^{x}_{i}(1)$, hence in $\mathrm{Ball}^{x}_{i}(4)$, for all $i<k$:
the trajectory on $R$ exists up to level $k$.
The assertions of Proposition~28.1 on $\varepsilon^{\mathrm{v},t}_k$, $m_k$, $\mathscr{E}_k$, $E_k$ and $\mathfrak{P}_k$ are contained in (ii)$_k$.
By Proposition~48.1\,(a), (ii)$_k$, (iii)$_k$, the inequality $\frac{32}{15}\le4$ and (u0),
\[
\|E_k\|_k=\left\|\mathscr{E}_k|_0\right\|_k\le\|\mathscr{E}_k\|^{x}_k\le\frac{32}{15}\,e_k^{\varepsilon}\cdot w^{\mathrm{E}}_k\le4\,e_k^{\varepsilon}\cdot w^{\mathrm{E}}_k\le w^{\mathrm{E}}_k.
\]

\paragraph{1.7. The characterisation of the controlled pairs.}
Suppose that, recursively, $|m_k|\le w^{\mathrm{m}}_k$ for $k=0,1,\dots,k_{*}$, that is, $|m_k|\le w^{\mathrm{m}}_k$ for every $k\le k_{*}$ for which the data of level $k$ are defined.
This condition contains (Hyp$_{k_{*}}$); so, by 1.2 and 1.6 at $k=k_{*}$, the trajectory exists up to level $k_{*}$, and the data of all levels $j\le k_{*}$ are defined ((i)$_j$, $0\le j\le k_{*}$).
Hence $|m_k|\le w^{\mathrm{m}}_k$ for $0\le k\le k_{*}$, and the pair is controlled on $R$ (Section~0.2).
Conversely, if the pair is controlled on $R$, then by the definition the data of all levels $k\le k_{*}$ are defined and $|m_k|\le w^{\mathrm{m}}_k$.

\paragraph{1.8. Level $0$.}
The assertion of Proposition~28.1 on level $0$ is (rec0).

\paragraph{1.9. Uniformity.}
The numbers $\frac{1}{16}$, $2$, $\frac{32}{15}$ and $4$ are absolute.
The inequalities (u0) and (u3)--(u5) hold for all data (FL) and all $k$ in the stated ranges, by conditions on $L$, $\varepsilon$ and $e_0$ alone.
The bound of Proposition~24.1\,(d) holds with the stated numbers for all data (OS) and every system with the properties named there (Proposition~24.1\,(i)).
The bounds of Proposition~28.1 contain $\varepsilon$ and the datum $e_k=e\cdot L^{\frac{k-N}{2}}$; the constants $c_{\mathrm{fl}}$ and $L$ enter the proof through (u3).
Apart from $e_k$, none of the level $k$, the side $\ell_k$, the mass $\mu_k$, $N$, $R$, $\bar{m}$ and the pair occurs in these bounds, and the recursion of Section~0.2 contains no twist.
Therefore, apart from the datum $e_k$, the bounds contain neither $N$, $R$, $\theta$ nor the pair.
\end{proof}

\subsection*{2. Proposition~28.2 (exactly one pair stays in the windows and ends at zero)}

\begin{proof}[Proof of Proposition~28.2]
The argument rests on Banach's fixed point theorem:
a map $\mathrm{Ffl}$ of a non-empty complete metric space $\mathrm{Bfl}$ into itself such that
$\mathrm{dist}(\mathrm{Ffl}(\mathrm{zfl}),\mathrm{Ffl}(\widetilde{\mathrm{zfl}}))\le\mathrm{cf}\cdot\mathrm{dist}(\mathrm{zfl},\widetilde{\mathrm{zfl}})$
for all $\mathrm{zfl},\widetilde{\mathrm{zfl}}\in\mathrm{Bfl}$, with a number $\mathrm{cf}<1$,
has exactly one fixed point.

\paragraph{2.1. The space.}
For $0\le k\le k_{*}$ the space $K_{k}$ with the norm $\|\cdot\|_k$ is a real Banach space (Proposition~48.1\,(d)),
and the windows satisfy $w^{\mathrm{m}}_k>0$, $w^{\mathrm{E}}_k>0$.
Let $\mathrm{Vfl}$ be the real vector space of the sequences $\mathrm{zfl}=(m_k,E_k)_{0\le k\le k_{*}}$ with $m_k\in\mathbb{R}$ and $E_k\in K_{k}$,
equipped with the norm
\[
|\mathrm{zfl}|_{\mathrm{fl}}:=\max_{0\le k\le k_{*}}\max\left(\frac{|m_k|}{w^{\mathrm{m}}_k},\frac{\|E_k\|_k}{w^{\mathrm{E}}_k}\right).
\]
Then $\mathrm{Vfl}$ is a Banach space:
a Cauchy sequence for $|\cdot|_{\mathrm{fl}}$ is a Cauchy sequence in each of the $2(k_{*}+1)$ components;
the limits of the components form an element of $\mathrm{Vfl}$;
and convergence in all components is convergence for $|\cdot|_{\mathrm{fl}}$.
Put
\[
\mathrm{Bfl}:=\left\{\mathrm{zfl}\in\mathrm{Vfl}:\ |m_k|\le w^{\mathrm{m}}_k\ \text{and}\ \|E_k\|_k\le w^{\mathrm{E}}_k\ \text{for}\ 0\le k\le k_{*};\ E_0=0\right\}.
\]
Each of the conditions that define $\mathrm{Bfl}$ is closed, so that $\mathrm{Bfl}$ is closed in $\mathrm{Vfl}$;
it contains the sequence $0$.
With $\mathrm{dist}(\mathrm{zfl},\widetilde{\mathrm{zfl}}):=|\mathrm{zfl}-\widetilde{\mathrm{zfl}}|_{\mathrm{fl}}$ it is therefore a non-empty complete metric space.
For $\mathrm{zfl}=(m_k,E_k)$ and $\widetilde{\mathrm{zfl}}=(\tilde{m}_k,\tilde{E}_k)$ in $\mathrm{Bfl}$ let
\[
\mathrm{dfl}^{m}_k:=\frac{|m_k-\tilde{m}_k|}{w^{\mathrm{m}}_k},\qquad
\mathrm{dfl}^{E}_k:=\frac{\|E_k-\tilde{E}_k\|_k}{w^{\mathrm{E}}_k};
\]
then
\[
\mathrm{dist}(\mathrm{zfl},\widetilde{\mathrm{zfl}})=\max_k\max\left(\mathrm{dfl}^{m}_k,\mathrm{dfl}^{E}_k\right),\qquad\mathrm{dfl}^{E}_0=0.
\]
For $\mathrm{zfl}\in\mathrm{Bfl}$ and $k<k_{*}$ one has
\begin{equation}
m_k\in\mathbb{R},\qquad E_k\in K_{k},\qquad(m_k,E_k)\in\mathrm{Ball}_{k}(1);
\tag{28.6}
\end{equation}
at $k=0$ this is the pair $(m_k,0)$ with $|m_k|\le w^{\mathrm{m}}_0$.

\paragraph{2.2. The map.}
For $\mathrm{zfl}\in\mathrm{Bfl}$ let $\mathrm{Ffl}(\mathrm{zfl}):=(m'_k,E'_k)_{0\le k\le k_{*}}$, where
\begin{equation}
\begin{aligned}
m'_k&:=L^{-1}\cdot m_{k+1}-m(E_k)\quad(0\le k<k_{*}),&m'_{k_{*}}&:=0,\\
E'_0&:=0,&E'_k&:=L R E_{k-1}+L E_{k-1}^{\mathrm{det}}+E^{*}_{k-1}(m_{k-1},E_{k-1})\quad(1\le k\le k_{*}).
\end{aligned}
\tag{28.7}
\end{equation}
The right sides are defined.
Indeed, $m(\cdot)$ is defined on $K_{k}$ (Proposition~48.2\,(a); see 0.2),
and $L R$ on $K_{k-1}$ (Section~I.5.6; Proposition~48.5\,(a));
the activity $L E_{k-1}^{\mathrm{det}}$ is formed with the family (p5) of the system of primitives of (FL) at the level $k-1$
(Sections~I.5.7, I.5.6; Proposition~13\,(d5));
and $E^{*}_{k-1}$ is defined on $\mathrm{Ball}_{k-1}(4)$ (Proposition~24.1\,(b)),
a set which contains $(m_{k-1},E_{k-1})$ by (28.6).

\paragraph{2.3. Fixed points.}
An element $\mathrm{zfl}$ of $\mathrm{Bfl}$ is a fixed point of $\mathrm{Ffl}$ if and only if
\begin{equation}
\begin{aligned}
&m_{k+1}=L\cdot(m_k+m(E_k)),\qquad E_{k+1}=L R E_k+L E_k^{\mathrm{det}}+E^{*}_k(m_k,E_k)\qquad(0\le k<k_{*}),\\
&m_{k_{*}}=0.
\end{aligned}
\tag{28.8}
\end{equation}
In fact, the equation $m_k=L^{-1}\cdot m_{k+1}-m(E_k)$ is the first relation of (28.8);
the equations $E_k=E'_k$, $1\le k\le k_{*}$, are the second relation with $k-1$ in the place of $k$;
the equation $m_{k_{*}}=m'_{k_{*}}$ is the last relation;
and $E_0=0=E'_0$ holds on $\mathrm{Bfl}$.

\paragraph{2.4. $\mathrm{Ffl}$ maps $\mathrm{Bfl}$ into $\mathrm{Bfl}$.}
Let $\mathrm{zfl}\in\mathrm{Bfl}$.

\emph{Masses.}
Let $k<k_{*}$.
Since $E_k$ lies in $K_{k}$, the number $m(E_k)$ is real (Proposition~48.2\,(a)), and so is $m'_k$.
Proposition~48.2\,(b) gives
$|m(E_k)|\le c_{\mathrm{fl}}\cdot e_k^{\frac12}\cdot\|E_k\|_k\le c_{\mathrm{fl}}\cdot e_k^{\frac12}\cdot w^{\mathrm{E}}_k$.
Together with (u1), (u2) and (u6) it follows that
\[
\frac{|m'_k|}{w^{\mathrm{m}}_k}
\le L^{-1}\cdot\frac{w^{\mathrm{m}}_{k+1}}{w^{\mathrm{m}}_k}\cdot\frac{|m_{k+1}|}{w^{\mathrm{m}}_{k+1}}+c_{\mathrm{fl}}\cdot e_k^{\frac12}\cdot\frac{w^{\mathrm{E}}_k}{w^{\mathrm{m}}_k}
\le\frac14+c_{\mathrm{fl}}\cdot e_k^{\varepsilon}
\le\frac14+\frac1{16}\le1.
\]
Moreover $m'_{k_{*}}=0$.

\emph{Activities.}
Let $1\le k\le k_{*}$.
The activity $L R E_{k-1}$ lies in $K_{k}$ by Proposition~48.5\,(a),
and $L E_{k-1}^{\mathrm{det}}$ lies in $K_{k}$ by Proposition~13\,(d5);
the data $(m_{k-1},E_{k-1})$ are real and lie in $\mathrm{Ball}_{k-1}(4)$ by (28.6),
so that $E^{*}_{k-1}(m_{k-1},E_{k-1})$ lies in $K_{k}$ by Proposition~24.1\,(b).
Hence $E'_k$ lies in $K_{k}$.
By Propositions~48.5\,(b), 13\,(d6) and 24.1\,(b),
\[
\begin{gathered}
\|L R E_{k-1}\|_k\le c_{\mathrm{fl}}\cdot L^{-\frac14+2\varepsilon}\cdot\|E_{k-1}\|_{k-1}\le c_{\mathrm{fl}}\cdot L^{-\frac14+2\varepsilon}\cdot w^{\mathrm{E}}_{k-1},\\
\|L E_{k-1}^{\mathrm{det}}\|_k\le e_k^{\frac14-2\varepsilon},\qquad
\|E^{*}_{k-1}(m_{k-1},E_{k-1})\|_k\le e_{k-1}^{\frac14-6\varepsilon}.
\end{gathered}
\]
Dividing by $w^{\mathrm{E}}_k$ and using (u3), (u4), (u5), each with $k-1$ in the place of $k$, together with (u6), one obtains
\[
\frac{\|E'_k\|_k}{w^{\mathrm{E}}_k}\le\frac1{16}+e_k^{\varepsilon}+e_k^{\varepsilon}\le\frac1{16}+\frac1{16}\le1.
\]
Moreover $E'_0=0$.

\paragraph{2.5. $\mathrm{Ffl}$ is a contraction.}
Let $\mathrm{zfl},\widetilde{\mathrm{zfl}}\in\mathrm{Bfl}$,
let $(m'_k,E'_k)=\mathrm{Ffl}(\mathrm{zfl})$ and $(\tilde{m}'_k,\tilde{E}'_k)=\mathrm{Ffl}(\widetilde{\mathrm{zfl}})$,
and let $\mathrm{dfl}^{m'}_k$, $\mathrm{dfl}^{E'}_k$ be the quotients of 2.1 formed with the two images.

\emph{Masses.}
The map $m(\cdot)$ is linear (Proposition~48.2\,(a)).
For $k<k_{*}$,
\[
m'_k-\tilde{m}'_k=L^{-1}\cdot(m_{k+1}-\tilde{m}_{k+1})-m(E_k-\tilde{E}_k),
\]
and $|m(E_k-\tilde{E}_k)|\le c_{\mathrm{fl}}\cdot e_k^{\frac12}\cdot\|E_k-\tilde{E}_k\|_k$ (Proposition~48.2\,(b)).
With (u1), (u2) and (u6),
\[
\mathrm{dfl}^{m'}_k
\le L^{-1}\cdot\frac{w^{\mathrm{m}}_{k+1}}{w^{\mathrm{m}}_k}\cdot\mathrm{dfl}^{m}_{k+1}+c_{\mathrm{fl}}\cdot e_k^{\frac12}\cdot\frac{w^{\mathrm{E}}_k}{w^{\mathrm{m}}_k}\cdot\mathrm{dfl}^{E}_k
\le\frac14\,\mathrm{dfl}^{m}_{k+1}+c_{\mathrm{fl}}\cdot e_k^{\varepsilon}\cdot\mathrm{dfl}^{E}_k
\le\left(\frac14+\frac1{16}\right)\mathrm{dist}(\mathrm{zfl},\widetilde{\mathrm{zfl}}).
\]
Moreover $\mathrm{dfl}^{m'}_{k_{*}}=0$.

\emph{Activities.}
The map $L R$ is linear (Proposition~48.5\,(a)),
and the term $L E_{k-1}^{\mathrm{det}}$ does not depend on $\mathrm{zfl}$.
For $1\le k\le k_{*}$ therefore
\[
E'_k-\tilde{E}'_k=L R(E_{k-1}-\tilde{E}_{k-1})+E^{*}_{k-1}(m_{k-1},E_{k-1})-E^{*}_{k-1}(\tilde{m}_{k-1},\tilde{E}_{k-1}).
\]
Proposition~48.5\,(b) gives
$\|L R(E_{k-1}-\tilde{E}_{k-1})\|_k\le c_{\mathrm{fl}}\cdot L^{-\frac14+2\varepsilon}\cdot\|E_{k-1}-\tilde{E}_{k-1}\|_{k-1}$,
and (28.6), together with the form of Proposition~24.1\,(e) without slot
(the bound for $E^{*}_{k-1}$ and for pairs in $\mathrm{Ball}_{k-1}(1)$, with the norms $\|\cdot\|_{k-1}$ and $\|\cdot\|_k$), gives
\[
\|E^{*}_{k-1}(m_{k-1},E_{k-1})-E^{*}_{k-1}(\tilde{m}_{k-1},\tilde{E}_{k-1})\|_k\le e_{k-1}^{\frac14-6\varepsilon}\cdot\left(\mathrm{dfl}^{m}_{k-1}+\mathrm{dfl}^{E}_{k-1}\right).
\]
Dividing by $w^{\mathrm{E}}_k$ and using (u3), (u5), with $k-1$ in the place of $k$, and (u6), one obtains
\[
\mathrm{dfl}^{E'}_k
\le\frac1{16}\,\mathrm{dfl}^{E}_{k-1}+e_k^{\varepsilon}\cdot\left(\mathrm{dfl}^{m}_{k-1}+\mathrm{dfl}^{E}_{k-1}\right)
\le\left(\frac1{16}+2e_k^{\varepsilon}\right)\mathrm{dist}(\mathrm{zfl},\widetilde{\mathrm{zfl}})
\le\left(\frac1{16}+\frac1{16}\right)\mathrm{dist}(\mathrm{zfl},\widetilde{\mathrm{zfl}}).
\]
Moreover $\mathrm{dfl}^{E'}_0=0$.

The two parts together give
\[
\mathrm{dist}(\mathrm{Ffl}(\mathrm{zfl}),\mathrm{Ffl}(\widetilde{\mathrm{zfl}}))\le\frac{5}{16}\,\mathrm{dist}(\mathrm{zfl},\widetilde{\mathrm{zfl}}).
\]
By Banach's fixed point theorem $\mathrm{Ffl}$ has exactly one fixed point in $\mathrm{Bfl}$.
In what follows $\mathrm{zfl}=(m_k,E_k)_{0\le k\le k_{*}}$ denotes this fixed point.

\paragraph{2.6. Existence.}
Put
\begin{equation}
\begin{aligned}
m_{\mathrm{in}}&:=\text{the term of index }k=0\text{ of the sequence }(m_k)_{0\le k\le k_{*}}\text{ of the fixed point},\\
\varepsilon^{\mathrm{v}}_{\mathrm{in}}&:=-\sum_{i=0}^{k_{*}-1}L^{-3i}\cdot\left(\varepsilon(E_i)+\delta\varepsilon^{0}_{i}\right).
\end{aligned}
\tag{28.9}
\end{equation}
The number $m_{\mathrm{in}}$ is real by the definition of $\mathrm{Vfl}$.
Each $\varepsilon(E_i)$ is real because $E_i$ lies in $K_{i}$ (Proposition~48.2\,(a)),
and $\delta\varepsilon^{0}_{i}$ is real (Proposition~24.1\,(a)).
Thus $(m_{\mathrm{in}},\varepsilon^{\mathrm{v}}_{\mathrm{in}})$ is a real pair.
Within this subsection the mass coupling and the activity without slot of level $k$ of its trajectory on $R$
are written $m^{\mathrm{traj}}_k$ and $E^{\mathrm{traj}}_k$.
For $0\le k\le k_{*}$ consider the assertion
\begin{itemize}
\item[$(\mathrm{Cq}_k)$] the data of the levels $i\le k$ of the trajectory are defined,
and $m^{\mathrm{traj}}_i=m_i$, $E^{\mathrm{traj}}_i=E_i$ for all $i\le k$.
\end{itemize}
The assertion $(\mathrm{Cq}_0)$ holds:
by (rec0) and (28.9) the mass coupling of level $0$ of the trajectory is $m_{\mathrm{in}}$, the term of index $0$ of the fixed point,
and the activity without slot of level $0$ of the trajectory is $0=E_0$, because $\mathrm{zfl}$ lies in $\mathrm{Bfl}$.

Let $k<k_{*}$ and assume $(\mathrm{Cq}_k)$.
By $(\mathrm{Cq}_k)$ and $\mathrm{zfl}\in\mathrm{Bfl}$, one has $|m^{\mathrm{traj}}_i|=|m_i|\le w^{\mathrm{m}}_i$ for all $i<k+1$.
Proposition~28.1, applied with $k+1$ ($\le k_{*}$) in the place of $k$,
shows that the trajectory exists up to the level $k+1$ and that the data of level $k+1$ are defined.
By (rec2), (rec3), $(\mathrm{Cq}_k)$ and (28.8),
\[
m^{\mathrm{traj}}_{k+1}=L\cdot(m_k+m(E_k))=m_{k+1},\qquad
E^{\mathrm{traj}}_{k+1}=L R E_k+L E_k^{\mathrm{det}}+E^{*}_k(m_k,E_k)=E_{k+1}.
\]
This is $(\mathrm{Cq}_{k+1})$.

By $(\mathrm{Cq}_{k_{*}})$ the trajectory exists up to the level $k_{*}$,
and $|m^{\mathrm{traj}}_k|=|m_k|\le w^{\mathrm{m}}_k$ for $0\le k\le k_{*}$;
hence the pair is controlled on $R$.
Moreover $m^{\mathrm{traj}}_{k_{*}}=m_{k_{*}}=0$ by (28.8).
From now on the upper index $\mathrm{traj}$ is omitted.

\emph{The vacuum energy.}
For $k<k_{*}$ the relation (rec1) is equivalent to
$\varepsilon^{\mathrm{v},t}_k=L^{-3}\cdot\varepsilon^{\mathrm{v},t}_{k+1}-(\varepsilon(E_k)+\delta\varepsilon^{0}_{k})$.
It follows that, for $0\le k\le k_{*}$,
\begin{equation}
\varepsilon^{\mathrm{v},t}_k=L^{-3(k_{*}-k)}\cdot\varepsilon^{\mathrm{v},t}_{k_{*}}-\sum_{i=k}^{k_{*}-1}L^{-3(i-k)}\cdot\left(\varepsilon(E_i)+\delta\varepsilon^{0}_{i}\right).
\tag{28.10}
\end{equation}
This is seen by induction on $k_{*}-k$.
For $k=k_{*}$ the sum is empty.
If (28.10) holds at $k+1$, then
\[
\varepsilon^{\mathrm{v},t}_k=L^{-3}\cdot\left\{L^{-3(k_{*}-k-1)}\cdot\varepsilon^{\mathrm{v},t}_{k_{*}}-\sum_{i=k+1}^{k_{*}-1}L^{-3(i-k-1)}\cdot\left(\varepsilon(E_i)+\delta\varepsilon^{0}_{i}\right)\right\}-\left(\varepsilon(E_k)+\delta\varepsilon^{0}_{k}\right),
\]
which is the right side of (28.10) at $k$.
At $k=0$ the left side of (28.10) is $\varepsilon^{\mathrm{v}}_{\mathrm{in}}$;
comparison with (28.9) gives $L^{-3k_{*}}\cdot\varepsilon^{\mathrm{v},t}_{k_{*}}=0$, that is, $\varepsilon^{\mathrm{v},t}_{k_{*}}=0$.
With this value (28.10) is the formula asserted in Proposition~28.2.

\paragraph{2.7. Uniqueness.}
Let $(m_{\mathrm{in}},\varepsilon^{\mathrm{v}}_{\mathrm{in}})$ be a real pair that is controlled on $R$
and satisfies $m_{k_{*}}=0$, $\varepsilon^{\mathrm{v},t}_{k_{*}}=0$,
and let $(\varepsilon^{\mathrm{v},t}_k,m_k,E_k)$ be its trajectory on $R$.
Then $|m_k|\le w^{\mathrm{m}}_k$ for $0\le k\le k_{*}$.
Proposition~28.1, applied at every $k\le k_{*}$, shows that $m_k$ is real,
that $E_k$ lies in $K_{k}$ and that $\|E_k\|_k\le w^{\mathrm{E}}_k$;
moreover $E_0=0$ by (rec0).
Hence $(m_k,E_k)_{0\le k\le k_{*}}$ lies in $\mathrm{Bfl}$.
The relations (rec2), (rec3), together with $m_{k_{*}}=0$, are (28.8).
So this sequence is a fixed point of $\mathrm{Ffl}$, and by 2.5 it is the fixed point.
The first entry $m_{\mathrm{in}}$ of the pair is the mass coupling of level $0$ of its trajectory (rec0),
hence the term of index $0$ of the fixed point.
The identity (28.10) follows from (rec1) alone;
at $k=0$, with $\varepsilon^{\mathrm{v},t}_{k_{*}}=0$, it gives
\[
\varepsilon^{\mathrm{v}}_{\mathrm{in}}=-\sum_{i=0}^{k_{*}-1}L^{-3i}\cdot\left(\varepsilon(E_i)+\delta\varepsilon^{0}_{i}\right)
\]
with the terms $E_i$ of the fixed point.
Therefore the pair is the pair (28.9).

\paragraph{2.8. Bounds.}
Let $(\varepsilon^{\mathrm{v},t}_k,m_k,E_k)$ be the trajectory on $R$ of the pair (28.9), and put
\[
\mathrm{qfl}^{m}_k:=\frac{|m_k|}{w^{\mathrm{m}}_k},\qquad\mathrm{qfl}^{E}_k:=\frac{\|E_k\|_k}{w^{\mathrm{E}}_k}\qquad(\le1).
\]

\emph{Mass.}
One has $\mathrm{qfl}^{m}_{k_{*}}=0$.
For $k<k_{*}$, (28.8) gives $m_k=L^{-1}\cdot m_{k+1}-m(E_k)$,
and $|m(E_k)|\le c_{\mathrm{fl}}\cdot e_k^{\frac12}\cdot\|E_k\|_k$ by Proposition~48.2\,(b).
Division by $w^{\mathrm{m}}_k$, together with (u1), (u2) and (u6), gives
\[
\mathrm{qfl}^{m}_k
\le L^{-1}\cdot\frac{w^{\mathrm{m}}_{k+1}}{w^{\mathrm{m}}_k}\cdot\mathrm{qfl}^{m}_{k+1}+c_{\mathrm{fl}}\cdot e_k^{\frac12}\cdot\frac{w^{\mathrm{E}}_k}{w^{\mathrm{m}}_k}\cdot\mathrm{qfl}^{E}_k
\le\frac14\,\mathrm{qfl}^{m}_{k+1}+c_{\mathrm{fl}}\cdot e_k^{\varepsilon}\cdot\mathrm{qfl}^{E}_k
\le\frac14\,\mathrm{qfl}^{m}_{k+1}+\frac1{16}.
\]
By induction on $k_{*}-k$,
\[
\mathrm{qfl}^{m}_k\le\frac1{16}\sum_{i=0}^{k_{*}-k-1}4^{-i}\le\frac1{16}\cdot\frac43=\frac1{12}\le\frac18,
\]
that is, $|m_k|\le w^{\mathrm{m}}_k/8$.

\emph{Vacuum energy.}
Let $0\le k\le i\le k_{*}-1$.
By Proposition~48.2\,(b), Proposition~28.1 and (28.1),
\[
|\varepsilon(E_i)|\le c_{\mathrm{fl}}\cdot\|E_i\|_i\le4c_{\mathrm{fl}}\cdot e_i^{\varepsilon}\cdot w^{\mathrm{E}}_i
=4c_{\mathrm{fl}}\cdot e_i^{\frac14-6\varepsilon}
=4c_{\mathrm{fl}}\cdot e_k^{\frac14-6\varepsilon}\cdot L^{(i-k)\left(\frac18-3\varepsilon\right)}.
\]
The inequality $L^{-\frac58}\le\frac14$ implies $L\ge4$.
Hence
\[
L^{-3(i-k)}\cdot L^{(i-k)\left(\frac18-3\varepsilon\right)}=L^{-(i-k)\left(\frac{23}{8}+3\varepsilon\right)}\le L^{-2(i-k)}\le16^{-(i-k)},
\]
and
\[
\sum_{i=k}^{k_{*}-1}L^{-3(i-k)}\cdot|\varepsilon(E_i)|\le4c_{\mathrm{fl}}\cdot e_k^{\frac14-6\varepsilon}\sum_{j\ge0}16^{-j}=\frac{64}{15}\,c_{\mathrm{fl}}\cdot e_k^{\frac14-6\varepsilon}.
\]
By Proposition~24.1\,(a), $0\le\delta\varepsilon^{0}_{i}\le e_i^7$, and by (28.1)
\[
L^{-3(i-k)}\cdot e_i^7=L^{-3(i-k)}\cdot\left(e_k\cdot L^{\frac{i-k}{2}}\right)^6\cdot e_i=e_k^6\cdot e_i.
\]
Further $e_i=L^{-\frac{k_{*}-i}{2}}\cdot e_{*}$ and $L^{\frac12}\ge2$.
Consequently
\[
0\le\sum_{i=k}^{k_{*}-1}L^{-3(i-k)}\cdot\delta\varepsilon^{0}_{i}
\le e_k^6\sum_{i=k}^{k_{*}-1}e_i
\le e_k^6\cdot e_{*}\sum_{j\ge1}L^{-\frac{j}{2}}
=e_k^6\cdot e_{*}\cdot\left(L^{\frac12}-1\right)^{-1}
\le e_k^6\cdot e_{*}
\le\frac1{48}\,e_k^{\frac14-6\varepsilon},
\]
where the last step uses $e_k<1$, $6\ge\frac14-6\varepsilon$ and $e_{*}\le\frac1{48}$ (u6).
Inserting the two bounds into the formula of Proposition~28.2 for $\varepsilon^{\mathrm{v},t}_k$, and using $c_{\mathrm{fl}}\ge1$, one obtains
\[
|\varepsilon^{\mathrm{v},t}_k|\le\left(\frac{64}{15}+\frac1{48}\right)c_{\mathrm{fl}}\cdot e_k^{\frac14-6\varepsilon}
=\frac{343}{80}\,c_{\mathrm{fl}}\cdot e_k^{\frac14-6\varepsilon}
\le5c_{\mathrm{fl}}\cdot e_k^{\frac14-6\varepsilon}
=5c_{\mathrm{fl}}\cdot e_k^{\varepsilon}\cdot w^{\mathrm{E}}_k.
\]
At $k=k_{*}$ both sums are empty and $\varepsilon^{\mathrm{v},t}_{k_{*}}=0$.

\paragraph{2.9. Level 0 written out.}
The component of index $0$ of an element of $\mathrm{Bfl}$ is a pair $(m,0)$ with $m$ real and $|m|\le w^{\mathrm{m}}_0=q^{\frac34-8\varepsilon}$,
that is, an element of $\mathrm{Ball}_{0}(1)$ (Section~I.5.4, (I.L0.21)).
Since $m(\cdot)$, $\varepsilon(\cdot)$ and $L R$ are linear, one has $m(E_0)=0$, $\varepsilon(E_0)=0$ and $L R E_0=0$,
and the components of (28.7) that involve the level $0$ read
\[
m'_0=L^{-1}\cdot m_1,\qquad E'_0=0,\qquad E'_1=L E_0^{\mathrm{det}}+E^{*}_0(m,0),
\]
where $m$ and $m_1$ are the terms of index $0$ and $1$ of the sequence of the masses.
Propositions~13\,(d5), (d6) and 24.1\,(b), (e) include the case $k=0$:
the activities $L E_0^{\mathrm{det}}$ and $E^{*}_0(m,0)$ lie in the space $K_{k+1}$ at $k=0$, and
\[
\|L E_0^{\mathrm{det}}\|_1\le\left(L^{\frac12}\cdot q\right)^{\frac14-2\varepsilon},\qquad
\|E^{*}_0(m,0)\|_1\le q^{\frac14-6\varepsilon},\qquad
\|E^{*}_0(m,0)-E^{*}_0(\tilde{m},0)\|_1\le q^{\frac14-6\varepsilon}\cdot\frac{|m-\tilde{m}|}{w^{\mathrm{m}}_0}
\]
(the second term of Proposition~24.1\,(e) is absent).
The activity $E_0^{\mathrm{det}}$ is formed with $\acute{E}_0$ of Section~I.5.7,
whose traces are traces on $\mathfrak{S}^{\mathrm{ch}}(1,\ell_0)$, a space of dimension $192\,n_0$,
with $\int\mathrm{d}x=\frac1{48}\sum_c\sum_g$ (see 1.5).
At these components the estimates of 2.4 and 2.5 are those written there, with $\|E_0\|_0=0$ and $\mathrm{dfl}^{E}_0=0$.
In (28.9) and (28.10) the term with $i=0$ is $\delta\varepsilon^{0}_{0}$.
By (28.8), $m_{k_{*}}=0$ and $m_k=L^{-1}\cdot m_{k+1}-m(E_k)$ for $k<k_{*}$;
induction on $k_{*}-k$ gives
\[
m_k=-\sum_{i=k}^{k_{*}-1}L^{-(i-k)}\cdot m(E_i)
\]
(if this holds at $k+1$, then $m_k=-\sum_{i=k+1}^{k_{*}-1}L^{-1-(i-k-1)}\cdot m(E_i)-m(E_k)$),
and at $k=0$, where $m(E_0)=0$,
\[
m_{\mathrm{in}}=-\sum_{i=1}^{k_{*}-1}L^{-i}\cdot m(E_i).
\]
The unknowns of the fixed point are the real numbers $m_k$ and the activities $E_k$.
The pair consists of the coupling $m_{\mathrm{in}}$ of the term
\[
-m_{\mathrm{in}}\cdot\frac1{48}\sum_{c}\sum_{g\in O_{\mathrm{cub}}}\bar{\Psi}(c,g)\cdot\Psi(c,g)
\]
and of the coupling $\varepsilon^{\mathrm{v}}_{\mathrm{in}}$ of the term $-\varepsilon^{\mathrm{v}}_{\mathrm{in}}\cdot n_0$ (see 0.2).
No other coupling occurs in (rec0)--(rec5).

\paragraph{2.10. Uniformity.}
The numbers $\frac14$, $\frac1{16}$, $\frac5{16}$, $\frac1{12}$, $\frac{64}{15}$, $\frac1{48}$, $\frac{343}{80}$ are absolute.
The inequalities (u0)--(u6) hold for all data (FL).
The bounds obtained from other Propositions hold, with the numbers stated there, for all data of 0.3
(Proposition~24.1\,(i); the constants of Propositions~13, 48.2 and 48.5 are those of Section~I.8).
The number $k_{*}+1$ of the levels depends on $N$, but it enters no estimate:
the sums over the levels are bounded by $\sum_i4^{-i}$, $\sum_j16^{-j}$ and $\sum_jL^{-\frac{j}{2}}$.
Hence the bounds of Proposition~28.2 contain $c_{\mathrm{fl}}$, $\varepsilon$ and the datum $e_k=e\cdot L^{\frac{k-N}{2}}$;
apart from $e_k$, neither $N$, $R$, $\bar{m}$ nor a twist occurs in them.
\end{proof}

\subsection*{3. Proposition~28.3 (one pair for all tori of the scheme)}

\begin{proof}[Proof of Proposition~28.3]
\leavevmode
\paragraph{3.0. The definition is meaningful.}
Let $\bar{m}\in(0,1]$ and $e>0$ with $e\cdot h^{\frac12}\le e_1$.
The side $M$ is a power of $L$ ((M1)),
and $r_K(e)=\min\{L^i\ge v_K^r:\ i\ge0\ \text{an integer}\}$ is a power $L^i$ with $i\ge0$ (Section~I.6.1).
Hence $R_{\mathrm{top}}(e)=h\cdot M\cdot r_K(e)=h\cdot L^T$ with an integer $T\ge0$, and $R_{\mathrm{top}}(e)\le R_{\mathrm{top}}(e)$;
so $R_{\mathrm{top}}(e)$ lies in $\mathrm{Tor}(e)$ (Section~I.6.2; this is part of (Sc0) of Section~G4.0.2).
From $e\cdot h^{\frac12}\le e_1$ one has $v_K=-\log\left(e\cdot h^{\frac12}\right)\ge-\log e_1$.
Since (e1)--(e8) hold for all $v\ge-\log e_1$ (Proposition~14.3\,(f)), the row (e1) at $v=v_K$ gives
\[
v_K\ge8\qquad\text{and}\qquad\left(1+\frac{\mathrm{lg}}{v_K}\right)^r\le\left(1+\frac{3\,\mathrm{lg}}{v_K}\right)^r<L;
\]
both inequalities belong to (Sc0), and neither contains $R$ or $N$.
They are the hypotheses of \cite[Proposition~19.1]{QEDI} for the torus $R_{\mathrm{top}}(e)$ in $\mathrm{Tor}(e)$,
the number $h'$ of \cite{QEDI} being read as $h$ (Section~I.6.1).
By \cite[Proposition~19.1(b), (c)]{QEDI} the number $\mathrm{js}:=\mathrm{js}(R_{\mathrm{top}}(e),e)$ exists
(it is denoted $j_{*}$ in \cite[Proposition~19.1(b)]{QEDI}), and
\[
N_{\min}(R_{\mathrm{top}}(e),e)=\max\{0,\ \mathrm{js}+1-\log_Lh\},
\]
an integer $\ge0$ that is formed with $(e\cdot h^{\frac12},h,L,M,r)$;
this is the assertion of (Sc0) on $N_{\min}(R,e)$ at $R=R_{\mathrm{top}}(e)$.
Now let $N\ge N_{\min}(R_{\mathrm{top}}(e),e)$ be an integer.
Then $(\bar{m},e,R_{\mathrm{top}}(e),N)$ has (S-4) and is therefore (FL) (Proposition~26.6\,(e); Section~G4.0.1);
the primitives that belong to it are the systems of Proposition~26.5 on $T_k(R_{\mathrm{top}})$, $k<k_{\mathrm{top}}$.
By Proposition~28.2 on the torus $R_{\mathrm{top}}(e)$ there is exactly one real pair that is controlled on $R_{\mathrm{top}}(e)$
and satisfies $m_{k_{\mathrm{top}}}=0$, $\varepsilon^{\mathrm{v},t}_{k_{\mathrm{top}}}=0$.
This pair is $P_N(e)$, and its trajectories on $R_{\mathrm{top}}(e)$ exist up to the level $k_{\mathrm{top}}$.
By Propositions~28.1 and 28.2 on $R_{\mathrm{top}}(e)$ one has, for $0\le k\le k_{\mathrm{top}}$,
\begin{equation}
\begin{gathered}
m^{a}_k,\ \varepsilon^{\mathrm{v},a}_k\ \text{are real};\qquad
|m^{a}_k|\le\frac{w^{\mathrm{m}}_k}{8};\qquad
|\varepsilon^{\mathrm{v},a}_k|\le5c_{\mathrm{fl}}\cdot e_k^{\varepsilon}\cdot w^{\mathrm{E}}_k;\\
\mathfrak{P}^{a}_k(X,\cdot)\ \text{is a complex function on}\ \mathfrak{X}_k;\qquad
\mathscr{E}^{a}_k\in K^{x}_{k}(\ell_k(R_{\mathrm{top}})),\qquad
E^{a}_k=\mathscr{E}^{a}_k|_0\in K_{k}(\ell_k(R_{\mathrm{top}}));\\
\|E^{a}_k\|_k\le\|\mathscr{E}^{a}_k\|^{x}_k\le4e_k^{\varepsilon}\cdot w^{\mathrm{E}}_k\le w^{\mathrm{E}}_k.
\end{gathered}
\tag{28.11}
\end{equation}
The numbers $e_k=e\cdot L^{\frac{k-N}{2}}$, $w^{\mathrm{m}}_k$ and $w^{\mathrm{E}}_k$ are formed with $(e,N,k)$ and $\varepsilon$ (Sections~I.6.1, I.5.4);
they are the same numbers on all tori.

\paragraph{3.1. Clause (a): $P_N(e)$ is defined.}
Let (FL) hold.
By (Sc2) with $R'=R_{\mathrm{top}}(e)$ one has $N_{\min}(R_{\mathrm{top}},e)\le N_{\min}(R,e)\le N$ and $k_{*}(R)\le k_{\mathrm{top}}$.
So the data of the Definition of $P_N(e)$ are given (this is also Proposition~26.6\,(a)),
and the right sides of the display (6.1) of Proposition~28.3\,(a) are defined for $0\le k\le k_{*}(R)$.

\paragraph{3.2. Clause (a) for $R=R_{\mathrm{top}}$.}
In this case the trajectory on $R$ of $P_N(e)$ is $(\varepsilon^{\mathrm{v},a}_k,m^{a}_k,E^{a}_k)$, with slot $(\mathscr{E}^{a}_k,\mathfrak{P}^{a}_k)$,
and (6.1) holds with equality.
The pair is controlled on $R$ by 3.0.

\paragraph{3.3. Clause (a) for $R<R_{\mathrm{top}}$.}
Put $k_m:=k_{*}(R)$.
The data (FL) with $R<R_{\mathrm{top}}$, together with a level $k\le k_m$, are the standing data (TT) of the Propositions~27.m,
with $T_1=T_k(R)$, $T_2=T_k(R_{\mathrm{top}})$, $\ell_1=\ell_k(R)$, $\ell_2=\ell_k(R_{\mathrm{top}})$.
For $k<k_m$ the system $\mathrm{Op}^{(2)}$ of (TT) is the system of Proposition~26.5 on $T_2$,
and $\mathrm{Op}^{(1)}$ is its projection to $T_1$;
these are the primitives of (FL) on $R_{\mathrm{top}}$ and on $R$ (Section~G4.0.1).
Hence the step maps on $T_1$ and $T_2$ that occur in Proposition~27.6 are the step maps of the recursion 0.2 on $R$ and on $R_{\mathrm{top}}$.

An element of $K^{x}_{k}(\ell_i)$ is a family of even elements of $\mathrm{Gr}[\psi,\chi]$ on $\tilde{R}_k\times\mathfrak{X}_k$
with the property (I.k2) of Section~I.5.4,
that is, a local family on $T_i$ in the sense of Proposition~27.2\,(D3) with $\mathrm{Dm}=\tilde{R}_k\times\mathfrak{X}_k$;
an element of $K_{k}(\ell_i)$ is a local family without slot.
So the relation $\sim$ of Proposition~27.2\,(D4) is defined for the entries of (6.1).

For $0\le k\le k_m$ consider the assertion
\begin{itemize}
\item[$(\mathrm{Ag}_k)$] the data of the levels $i\le k$ of the trajectory on $R$ of $P_N(e)$ are defined, and for all $i\le k$
\[
m_i=m^{a}_i,\qquad\varepsilon^{\mathrm{v},t}_i=\varepsilon^{\mathrm{v},a}_i,\qquad\mathscr{E}_i\sim\mathscr{E}^{a}_i,\qquad\mathfrak{P}_i\sim\mathfrak{P}^{a}_i,\qquad E_i\sim E^{a}_i\qquad\text{at the level }i.
\]
\end{itemize}

The assertion $(\mathrm{Ag}_0)$ holds.
By (rec0) the data of level $0$ are $(\varepsilon^{\mathrm{v}}_{0,N}(e),m_{0,N}(e),0,0,0)$ on both tori.
The map $\mathrm{idn}[\tilde{X}]$ is linear, so $\mathrm{idn}[\tilde{X}]\,0=0$;
hence $\mathscr{E}_0\sim\mathscr{E}^{a}_0$ and $E_0\sim E^{a}_0$.
Moreover $\mathfrak{P}^{a}_0(\mathrm{pr}_2\tilde{X},\mathrm{pj}^{*}\mathscr{C})=0=\mathfrak{P}_0(\mathrm{pr}_1\tilde{X},\mathscr{C})$
for every connected admissible $\tilde{X}$ and every $\mathscr{C}$, so that $\mathfrak{P}_0\sim\mathfrak{P}^{a}_0$.

Let $k<k_m$ and assume $(\mathrm{Ag}_k)$.
\begin{itemize}
\item[(1)]
By $(\mathrm{Ag}_k)$ and (28.11), $|m_i|=|m^{a}_i|\le w^{\mathrm{m}}_i/8\le w^{\mathrm{m}}_i$ for all $i\le k$.
Proposition~28.1 on $R$, applied with $k+1$ ($\le k_m$) and with $k$, shows:
the trajectory on $R$ exists up to the level $k+1$, and the data of level $k+1$ are defined;
$m_k$ and $\varepsilon^{\mathrm{v},t}_k$ are real;
$\mathscr{E}_k$ lies in $K^{x}_{k}(\ell_k(R))$ with $\|\mathscr{E}_k\|^{x}_k\le w^{\mathrm{E}}_k$;
$\mathscr{E}_k|_0=E_k$;
the $\mathfrak{P}_k(X,\cdot)$ are complex functions on $\mathfrak{X}_k(T_1)$.
On $R_{\mathrm{top}}$ the same holds by (28.11), since $k<k_m\le k_{\mathrm{top}}$.
\item[(2)]
Proposition~27.6 is applied with
\[
m:=m_k=m^{a}_k,\qquad\mathscr{E}^{(1)}:=\mathscr{E}_k,\qquad\mathscr{E}^{(2)}:=\mathscr{E}^{a}_k,\qquad\mathfrak{P}_k^{(1)}:=\mathfrak{P}_k,\qquad\mathfrak{P}_k^{(2)}:=\mathfrak{P}^{a}_k.
\]
Its hypotheses are satisfied:
(TT) holds with $0\le k<k_m$;
$m$ is real and $|m|\le w^{\mathrm{m}}_k$;
$\mathscr{E}^{(i)}$ lies in $K^{x}_{k}(\ell_i)$ with norm at most $w^{\mathrm{E}}_k$ (both vanish if $k=0$);
$\mathscr{E}^{(1)}\sim\mathscr{E}^{(2)}$ and $\mathfrak{P}_k^{(1)}\sim\mathfrak{P}_k^{(2)}$ by $(\mathrm{Ag}_k)$.
\item[(3)]
By Proposition~27.6\,(a), $\varepsilon(\mathscr{E}_k)=\varepsilon(\mathscr{E}^{a}_k)$ and $m(\mathscr{E}_k)=m(\mathscr{E}^{a}_k)$,
and $\delta\varepsilon^{0}_{k}$ is the same number on both tori.
Proposition~16.3\,(a) on each of the two tori, whose data (S-1) at the level $k<k_m\le k_{\mathrm{top}}$ are given by (28.3),
read for the flow data $(\bar{m},e,R,N)$ and $(\bar{m},e,R_{\mathrm{top}}(e),N)$,
yields, with $\mathscr{E}_k|_0=E_k$ and $\mathscr{E}^{a}_k|_0=E^{a}_k$,
\[
\varepsilon(E_k)=\varepsilon(\mathscr{E}_k)=\varepsilon(\mathscr{E}^{a}_k)=\varepsilon(E^{a}_k),\qquad
m(E_k)=m(\mathscr{E}_k)=m(\mathscr{E}^{a}_k)=m(E^{a}_k).
\]
With (rec2), (rec1) on the two tori and $(\mathrm{Ag}_k)$,
\[
\begin{aligned}
m_{k+1}&=L\cdot(m_k+m(E_k))=L\cdot(m^{a}_k+m(E^{a}_k))=m^{a}_{k+1},\\
\varepsilon^{\mathrm{v},t}_{k+1}&=L^3\cdot(\varepsilon^{\mathrm{v},t}_k+\varepsilon(E_k)+\delta\varepsilon^{0}_{k})=L^3\cdot(\varepsilon^{\mathrm{v},a}_k+\varepsilon(E^{a}_k)+\delta\varepsilon^{0}_{k})=\varepsilon^{\mathrm{v},a}_{k+1}.
\end{aligned}
\]
\item[(4)]
By (rec4) on the two tori, $\mathscr{E}_{k+1}=\mathscr{E}^{(1)+}$ and $\mathscr{E}^{a}_{k+1}=\mathscr{E}^{(2)+}$ in the notation of Proposition~27.6\,(b);
hence $\mathscr{E}_{k+1}\sim\mathscr{E}^{a}_{k+1}$ at the level $k+1$.
By (rec5), $\mathfrak{P}_{k+1}=\mathfrak{P}_{k+1}^{(1)}$ and $\mathfrak{P}^{a}_{k+1}=\mathfrak{P}_{k+1}^{(2)}$ in the notation of Proposition~27.6\,(c);
hence $\mathfrak{P}_{k+1}\sim\mathfrak{P}^{a}_{k+1}$.
\item[(5)]
By Proposition~28.1 on $R$ (at $k+1$) and by (28.11), $E_{k+1}=\mathscr{E}_{k+1}|_0$ and $E^{a}_{k+1}=\mathscr{E}^{a}_{k+1}|_0$.
Since $k+1\le k_m$, (TT) holds at the level $k+1$,
and Proposition~27.2\,(f) at this level turns (4) into $E_{k+1}\sim E^{a}_{k+1}$.
\end{itemize}
This proves $(\mathrm{Ag}_{k+1})$.

The assertion $(\mathrm{Ag}_{k_m})$ contains (6.1) for $0\le k\le k_m$.
By $(\mathrm{Ag}_{k_m})$ and (28.11), $|m_k|=|m^{a}_k|\le w^{\mathrm{m}}_k/8\le w^{\mathrm{m}}_k$ for $0\le k\le k_m$;
by the assertion of Proposition~28.1 that characterises the controlled pairs (see 1.7), the pair is controlled on $R$.

\paragraph{3.4. Clause (b).}
Let $k\le k_{*}(R)$.
By (6.1) and (28.11), $|m_k|=|m^{a}_k|\le w^{\mathrm{m}}_k/8$.
The pair is controlled on $R$, so Proposition~28.1 on $R$ gives $\|E_k\|_k\le\|\mathscr{E}_k\|^{x}_k\le4e_k^{\varepsilon}\cdot w^{\mathrm{E}}_k$.
At $k=k_{*}(R)$ ($\le k_{\mathrm{top}}$), (6.1) and (28.11) give
$|\varepsilon^{\mathrm{v},t}_k|=|\varepsilon^{\mathrm{v},a}_k|\le5c_{\mathrm{fl}}\cdot e_k^{\varepsilon}\cdot w^{\mathrm{E}}_k$,
and the numbers $e_k$, $w^{\mathrm{E}}_k$ at $k=k_{*}(R)$ are the numbers $e_{*}$, $w^{\mathrm{E}}_{*}$ of the torus $R$.

\paragraph{3.5. Clause (c).}
The pair $P_N(e)$ is real by 3.0.
The data of the Definition of $P_N(e)$ are $(\bar{m},e,N)$.
The following list contains every object that enters the definition of $P_N(e)$,
that is, Proposition~28.2 on the torus $R_{\mathrm{top}}(e)$, Section~G4.0.5 and the recursion 0.2.
Each entry names the place where the object is defined and the quantities with which it is formed, besides objects of earlier entries.
\begin{itemize}
\item[(i)] The constants of the lines (C0)--(C9) of Section~I.8.2.
\item[(ii)] Scales (Sections~I.0.5, I.6.1):
$\eta=L^{-N}$;
$h=L^{\left\lfloor\log_L\frac{1}{\bar{m}}\right\rfloor}$;
$a_k=L^{k-N}$;
$e_k=e\cdot a_k^{\frac12}$;
$v_k=-\log e_k$;
$\mu_k=\bar{m}\cdot a_k$;
$K=N+\log_Lh$;
$r_k$, $p_k$, $p_{0,k}$.
Formed with $(N,e,\bar{m})$ and $L$, $r$, $p$, $p_0$.
\item[(iii)] The torus (Sections~I.6.1, I.6.2):
$R_{\mathrm{top}}(e)=h\cdot M\cdot r_K(e)$, where $r_K(e)$ is formed with $v_K=-\log\left(e\cdot h^{\frac12}\right)$;
$N_{\min}(R_{\mathrm{top}}(e),e)$;
$\ell_k(R_{\mathrm{top}})=R_{\mathrm{top}}(e)/a_k$;
$n_k$;
$\mathfrak{n}_k$;
$k_{\mathrm{top}}$, the first $k$ with $\mathfrak{n}_k=0$.
Formed with (ii) and $M$.
\item[(iv)] Index sets, measure, generators (Sections~I.1.1--I.1.7, I.3.6):
$\mathbb{L}_k$, $T_k(R_{\mathrm{top}})$, $T^0_k$, $T^{\mathrm{f}}_k=T_k\times O_{\mathrm{cub}}$, $\mathrm{Pz}_k$, the $M$-cubes, $D_k$, the parameter cubes;
$\int\mathrm{d}x=\eta_k^3\cdot\frac1{48}\sum_c\sum_{g\in O_{\mathrm{cub}}}$;
$\psi$, $\chi$ with the weights $\frac{\eta_k^3}{48}$, $\left(\frac{\eta_k^3}{48}\right)^2$;
$\Psi$, $W$ with the weight $1$ on $T^0_k$ ($k\ge1$) and with the weight $\frac1{48}$ on $T^{\mathrm{f}}_0$ ($k=0$).
Formed with $k$, $\ell_k(R_{\mathrm{top}})$, $L$, $M$.
\item[(v)] Spin table, groups, the four transformations (Sections~I.2.1, I.4.1--I.4.5):
$\gamma_\mu$, $\gamma_5$, $\Sigma_3$, $\mathsf{C}$, $\varkappa$;
$\mathbb{G}_k$, $\mathbb{G}_k^{L}$;
$a_{\tilde{g}}$, $G_{\Lambda}$, $\mathrm{Cg}$, $\mathbf{K}$.
Formed with (iv) and $e_k$.
\item[(vi)] Domains (Section~I.5.3):
$\tilde{R}_k$, $\tilde{R}_k^{-}$, $\tilde{G}_k$, $\mathfrak{X}_k$.
Formed with $e_k$, $\varepsilon$, $\alpha$, $L$, $M$.
\item[(vii)] Norms, spaces, windows, balls (Sections~I.5.1, I.5.4):
$h_k=\left(e_k^{-\frac14},e_k^{-\frac14+\varepsilon}\right)$;
$K_{k}$, $K^{x}_{k}$ and their complex forms, with $\kappa$ in (I.k6);
$w^{\mathrm{m}}_k$, $w^{\mathrm{E}}_k$;
$\mathrm{Ball}_{k}$, $\mathrm{Ball}^{x}_{k}$.
\item[(viii)] Extraction, adjustment, reblocking (Sections~I.5.6, I.3.7):
$\mathrm{al}_{0}$, $\mathrm{al}_{2}$, $\varepsilon(\cdot)$, $m(\cdot)$, $R$,
$\boldsymbol{\Theta}_{X,\mathscr{T}}$ with the chain (I.SY6)$^{\mathrm{ch}}$,
the adjustment $E\mapsto E^{\mathrm{nat}}$, $S_L$, $L$.
Formed with (iv), $\alpha$, $L$, $M$ and $e_k$ (through $U_{\mathscr{T}}$).
\item[(ix)] Photon objects (Section~I.5.7; \cite[Section~5.2.5]{QEDI}):
the basis $\{e^{0}_{i}\}$, $\{\Upsilon\}$, $V_{k}$, $C$, $\mathrm{Box}_{k}$, $\mu_I$, $\hat{\mu}$, $\delta\varepsilon^{0}_{k}$.
Formed with $L$, $\ell_k(R_{\mathrm{top}})$, $p_{0,k}$.
\item[(x)] Exact fermionic operators and the forms in the displays
(Sections~I.1.8, I.2.2, I.2.3, I.3.1--I.3.6, I.3.9, I.5.9; the definitions at the heads of Propositions~23.1, 23.2):
the phases $u_B$, $U_B$, $\varphi_B$;
$M_{i}(x)$, $D_A$, $J$, $J^{*}$, $Q_k(A)$, $Q(A)$, $P(A)$, $B_{k,y}(A)$, $F_0(A)$, $\delta_{\alpha,B}$;
$T^{(1)}$, $T^{(2)}$, $\bar{T}^{(1)}$, $\bar{T}^{(2)}$ and the substitution $\mathrm{sub}_{\mathbf{Y},\mathscr{T}}$ ($k\ge1$);
the vertex $\mathrm{Vx}$, that is, the display (I.5.5) ($k=0$);
the forms $\langle\bar{\psi},\psi\rangle_{\square}$.
Formed with $\rho$, the weights $w$, $b$, $b_{k}$, $L$, $\alpha$, $\mu_k$, $e_k$.
\item[(xi)] The systems of primitives of Proposition~26.5 on $T_k(R_{\mathrm{top}})$, $k<k_{\mathrm{top}}$.
\item[(xii)] The determinant (Section~I.5.7):
$\acute{E}_k$, $E^{d}_k$, $E_k^{\mathrm{det}}$, ${\delta E^{\mathrm{det}}}_k$.
Formed with (x) and the family (p5) of (xi).
\item[(xiii)] The step map of level $k<k_{\mathrm{top}}$, with and without slot (Definition~24.D\,(D0)--(D8)):
formed with (iv)--(xii),
the contours $\mathrm{rad}_1$, $\mathrm{rad}_2$, $\mathrm{rad}_4$ (formed with $e_k$, $\varepsilon$),
the polydisk $\mathrm{Pd}_k$ (formed with $M^{\alpha_{0}}$)
and the coefficients $\mathrm{urs}^{T}$.
\item[(xiv)] The recursion (rec0)--(rec5);
the notions ``the trajectory exists up to level $k_{\mathrm{up}}$'' and ``controlled'' (Section~I.5.7, Section~G4.0.5).
\end{itemize}
Consequently every object of the list is formed with $(N,e,\bar{m})$, with constants of the lines (C0)--(C9),
with the basis $\{e^{0}_{i}\}$ and with the systems (xi), and with nothing else.
Once these are given, the set $\mathrm{Bfl}$ and the map $\mathrm{Ffl}$ of 2.1, 2.2 on the torus $R_{\mathrm{top}}(e)$
and the numbers $\delta\varepsilon^{0}_{i}$, $i<k_{\mathrm{top}}$, are given;
hence so are the fixed point of $\mathrm{Ffl}$ and, by (28.9), the pair.
This proves the second sentence of (c).

The remaining assertions of (c).
\emph{Torus:}
the data of the Definition of $P_N(e)$ are $(\bar{m},e,N)$ only, and the datum $R$ of (FL) is not among them;
no entry of the list is formed with this datum.
The only torus in the list is $R_{\mathrm{top}}(e)$, entry (iii);
it lies in $\mathrm{Tor}(e)$ by 3.0, and it is formed with $(e,\bar{m})$ and $L$, $M$, $r$ (entries (ii), (iii)).
\emph{Twist:}
$\theta$ and the connections $c^{(k)}_{\theta}$, $H_{\theta}$ of Section~I.6.6 occur in no entry;
in Proposition~24.1\,(c) a harmonic part occurs only inside the argument $A'$ of the two identities,
and the map (8.6)$^{\mathrm{ch}}$ does not contain $A'$.
\emph{External connection, probe:}
the second slot $\mathscr{C}$ of $\mathscr{E}_k$ and of $\mathfrak{P}_k$ is a variable that ranges over $\mathfrak{X}_k$;
no connection $a^{[k]}$ is inserted in (rec0)--(rec5).
A probe $\beta$, the translation vector $x_0$ and the numbers of the line (C10) that are formed with them
($D_{\mathrm{be}}$, $D_\beta$, $C_{\beta}$, $\mathbf{D}$, $\mathrm{Cb}$, $\mathrm{Cah}_i$, $\mathrm{Cpv}_1$, $\mathrm{Cpv}_2$, $\mathrm{epv}$, $\mathrm{Cpv}_\beta$,
$C_{\mathrm{ct}}$, $C_{1}^{\mathrm{N2}}$, $C_{2}^{\mathrm{N2}}$, $C_{\mathrm{fr}}$, $e_{\sharp}$, $C^{G}_{\beta}$, $s_{B}$, $C_{3\mathrm{p}}$, $C_{U}$, $C_{\bar{\kappa}}$, $C'_{\bar{\kappa}}$, $e_2$)
occur in no entry.
Of the line (C10) only $\bar{m}$ and $h$ occur, of the line (C11) only $e$ and $N$.

\paragraph{3.6. Clause (d).}
Let $\bar{m}\in(0,1]$, let $e>0$ with $e\cdot h^{\frac12}\le e_1$, and let $N\ge0$ be an integer;
put $q=e\cdot\eta^{\frac12}=e\cdot L^{-\frac{N}{2}}$.
Since $\bar{m}\le1$, one has $\left\lfloor\log_L\frac{1}{\bar{m}}\right\rfloor\ge0$;
so $h\ge1$ and $e\le e\cdot h^{\frac12}\le e_1\le\frac13$ ((C9); the inequality $e\le\frac13$ is part of (Sc0)).
Hence $0<q\le\frac13$, and, since $\frac34-8\varepsilon>0$,
\[
0<w^{\mathrm{m}}_0=q^{\frac34-8\varepsilon}<1.
\]
(For $N\ge N_{\min}(R_{\mathrm{top}}(e),e)$ this is the third clause of the condition (T1) of Section~G4.0.3 for the triple $(R_{\mathrm{top}}(e),e,N)$.)
If $N\ge N_{\min}(R_{\mathrm{top}}(e),e)$, then $m_{0,N}(e)$ is the first entry of $P_N(e)$, that is, by (rec0),
the mass coupling of level $0$ of its trajectory on $R_{\mathrm{top}}(e)$;
by (28.11) at $k=0$, $|m_{0,N}(e)|\le w^{\mathrm{m}}_0/8$.
If $N<N_{\min}(R_{\mathrm{top}}(e),e)$, then $m_{0,N}(e)=0\le w^{\mathrm{m}}_0/8$.
In both cases $w^{\mathrm{m}}_0/8<1/8$.

The number $m_{0,N}(e)$ is real: by (c), respectively because it vanishes.
The number $N_{\min}(R_{\mathrm{top}}(e),e)$ is formed with $(e,\bar{m})$ and $L$, $M$, $r$
(by the formula of 3.0; Section~I.6.2, with (ii), (iii) of 3.5).
Together with (c) it follows that $m_{0,N}(e)$ is formed with $(N,e,\bar{m})$ and the scheme.
Its data are $(\bar{m},e,N)$ only, as are those of the Definition of $P_N(e)$;
no torus is among them, in particular not the datum $R$ of (FL).
The only torus that occurs in the formation of $m_{0,N}(e)$,
namely in the Definition of $P_N(e)$ and in the threshold $N_{\min}(R_{\mathrm{top}}(e),e)$, is $R_{\mathrm{top}}(e)$.
It lies in $\mathrm{Tor}(e)$ by 3.0, and it is itself formed with $(e,\bar{m})$ and $L$, $M$, $r$:
in $R_{\mathrm{top}}(e)=h\cdot M\cdot r_K(e)$ the number $h=L^{\left\lfloor\log_L\frac{1}{\bar{m}}\right\rfloor}$ is formed with $\bar{m}$ and $L$,
and $r_K(e)$ with $v_K=-\log\left(e\cdot h^{\frac12}\right)$, $r$ and $L$ (entries (ii), (iii) of 3.5).
No other element of $\mathrm{Tor}(e)$ occurs,
and neither the even tori of Section~I.6.5, of side $8\cdot h\cdot L^{T}$, nor their index $T$ occurs.
Hence $N\mapsto m_{0,N}(e)$ is a real sequence that does not depend on $T$:
a mass assignment in the sense of Section~I.7.5.

\paragraph{3.7. Level 0 written out.}
In 3.3 at $k=0$ the data on both tori are $(m_{0,N}(e),0)$, with $\mathfrak{P}_0=\mathfrak{P}^{a}_0=0$.
At level $0$, Proposition~27.6 reads as follows.
In (a), $\varepsilon(\cdot)=0$ and $m(\cdot)=0$ on both tori,
so that $m_1=L\cdot m_{0,N}(e)=m^{a}_1$ and $\varepsilon^{\mathrm{v},t}_1=L^3\cdot(\varepsilon^{\mathrm{v}}_{0,N}(e)+\delta\varepsilon^{0}_{0})=\varepsilon^{\mathrm{v},a}_1$.
In (b) only $L\,{\delta E^{\mathrm{det}}}_0$ and $E^{*,x}_0(m_{0,N}(e),0)$ occur.
In (c),
\[
\mathfrak{P}_1(Y,\mathscr{C})=\sum_{X\in D_0:\ \bar{X}^{L}=LY}E_0^{\mathrm{det}}(X,\mathscr{C}_L).
\]
The generators $W$ of the first step are indexed by $T^{\mathrm{f}}_0$ and carry the weight $\frac1{48}$ on both tori.

\paragraph{3.8. Uniformity.}
The bounds of (b) and (d) are those of Propositions~28.1 and 28.2:
they contain $c_{\mathrm{fl}}$, $\varepsilon$ and the datum $e_k=e\cdot L^{\frac{k-N}{2}}$ (which is $q$ at $k=0$);
apart from $e_k$, neither $N$, $R$, $\bar{m}$ nor a twist occurs in them.
The induction 3.3 involves no estimate.
\end{proof}

\subsection*{4. Conditions on the constants}

The list below contains every condition on a constant that enters the arguments of Sections~0--3 directly, and not only through another Proposition.
Condition no.~7 enters through (Sc6), that is, through Proposition~14.2\,(c), (f);
the corresponding estimates are written out in (u6).
Conditions nos.~10 and 11 enter through results of \cite{QEDI};
no.~10 is also used directly in 3.0.
Each condition is stated among the Propositions~14.m, at the place named under ``stated in''.

\begin{description}
\item[No.~1]
Condition: $c_{\mathrm{fl}}\cdot L^{-\frac38+6\bar{\varepsilon}}\le\frac1{16}$.
Used in: (u3), and through it in 1.4\,(5), 2.4, 2.5.
Stated in: Proposition~14.2\,(b), row (L2); Section~I.8.2, line (C3).
\item[No.~2]
Condition: $L^{-\frac58}\le\frac14$.
Used in: (u1), and through it in 2.4, 2.5, 2.8.
Stated in: Proposition~14.2\,(f), as a consequence of (L1); Section~I.8.2, line (C3).
\item[No.~3]
Condition: $L\ge4$, in the forms $L^{-2}\le\frac1{16}$ and $L^{\frac12}-1\ge1$; and $L>1$.
Used in: $L\ge4$ in 2.8;
$L>1$ in (28.3), (u1), (u3), (u5), (u6), 1.4\,(5), 2.8, 3.0 (where $\mathrm{lg}>0$), 3.6.
Stated in: a consequence of no.~2 (which gives $L\ge4^{\frac85}$);
also Proposition~14.2\,(f) ($L\ge4^8$) and (b), row (L5);
Sections~I.0.5 and I.8.2, line (C3).
\item[No.~4]
Condition: $0<\varepsilon\le\bar{\varepsilon}=\frac1{200}$, in the forms
$\varepsilon>0$, $\frac{11\varepsilon}{2}\le6\bar{\varepsilon}$, $\frac18-3\varepsilon>0$, $\frac14-6\varepsilon>0$, $\frac34-8\varepsilon>0$, $\varepsilon\le1$.
Used in: (u1)--(u6), 1.4\,(5), 2.8, 3.6.
Stated in: Proposition~14.2\,(b), row (S1), and (a), rows $\bar{\varepsilon}$, $\varepsilon$; Section~I.8.2, lines (C0), (C7).
\item[No.~5]
Condition: $\exp(-\varepsilon v)\le\frac14$ for all $v\ge-\log e_0$, that is, $e_0^{\varepsilon}\le\frac14$; it contains $e_0<1$.
Used in: (u0), and through it in (u4), (u6), 1.4\,(1), 1.4\,(2), 1.4\,(4), 1.6, 2.8.
Stated in: Proposition~14.2\,(b), row (A8), first inequality; Section~I.8.2, line (C8).
\item[No.~6]
Condition: $c_{\mathrm{fl}}\ge1$.
Used in: (u6), 2.8.
Stated in: Proposition~14.2\,(a), row $c_{\mathrm{fl}}=61c_{X}\cdot K_{0}$, with $c_{X}=54$ and $K_{0}=\frac{72^{54}}{36}$; Section~I.8.2, line (C0).
\item[No.~7]
Condition: $(c_{\mathrm{fl}}+2)\exp(-\varepsilon v)\le\frac1{16}\left(1-L^{-\frac{\varepsilon}{2}}\right)$ for all $v\ge-\log e_1$.
Used in: (u6), through (fl) and (Sc6).
Stated in: Proposition~14.2\,(b), row (e2), and (c), condition (fl); Proposition~14.3\,(f); Section~I.8.2, line (C9).
\item[No.~8]
Condition: $e_1\le\frac13$.
Used in: 3.6 ($e\le\frac13$, hence $w^{\mathrm{m}}_0<1$ at all $N\ge0$; this contains the third clause of (T1)).
Stated in: Proposition~14.2\,(a), row $e_1$, and (f); Section~I.8.2, line (C9).
\item[No.~9]
Condition: $M$ is a power of $L$.
Used in: 3.0 ($R_{\mathrm{top}}(e)\in\mathrm{Tor}(e)$).
Stated in: Proposition~14.2\,(b), row (M1); Section~I.8.2, line (C5).
\item[No.~10]
Condition: $v\ge8$ and $\left(1+\frac{3\,\mathrm{lg}}{v}\right)^r<L$ for all $v\ge-\log e_1$.
Used in: 3.0 (at $v=v_K$);
through \cite[Proposition~19.1(b), (c)]{QEDI} in 3.0;
and through (Sc1) ($1\le k_{*}\le K$) and (Sc2), that is, \cite[Proposition~19.1(c), (d)]{QEDI},
whose hypotheses $v_K\ge1$, $\left(1+\frac{\mathrm{lg}}{v_K}\right)^r\le L$ it provides.
Stated in: Proposition~14.2\,(b), row (e1); Proposition~14.3\,(f); Section~I.8.2, line (C9).
\item[No.~11]
Condition: (A8); $L\ge5$, $L^{-\frac34}\le\frac12$; $\varepsilon\le\frac1{200}$.
Used in: 1.4\,(2), through \cite[Lemmas~17.4(a) and 17.5(a)--(c)]{QEDI}, on which the last sentence of Definition~24.D\,(D6) rests.
Stated in: Proposition~14.2\,(b), rows (A8), (L5), (S1), and (f); Section~I.8.2, lines (C8), (C3), (C7).
\end{description}

The condition $e_1\le e_0$ of line (C9) is used in the proofs of Propositions~26.6 and 48.6\,(b)
(Annex~B, Sections~G4.A.6 and G2.9, rows $e_1\le e_0$).
It enters the arguments of Sections~0--3 only through other Propositions,
in particular through the bounds in terms of $e_0$ that are taken from Proposition~26.6\,(b) (see (28.2), (28.3)).
Condition no.~3 enters Proposition~28.2 as a consequence of no.~2,
and Propositions~28.1 and 28.3 through the row (L5) of Proposition~14.2\,(b).
Each of the conditions nos.~1--11 is a row of Proposition~14.2\,(a), (b), or a consequence of these rows stated in Proposition~14.2\,(f).

\section*{Chapter 29. The terminal integral: proofs of Proposition~29}

\subsection*{Proof of Proposition~29}

\subsubsection*{P0. Conventions of this chapter}

Throughout the chapter the data (FL), the real pair $(m_{\mathrm{in}},\varepsilon^{\mathrm{v}}_{\mathrm{in}})$, its trajectories $(\varepsilon^{\mathrm{v},t}_k,m_k,E_k)$ and $(\mathscr{E}_k,\mathfrak{P}_k)$ on $R$ (Section~G4.0.5; at $k=0$ the symbol $m_k$ denotes $m_{\mathrm{in}}$) and the external connection $a$ are those of Proposition~29.
All lattices are those of the torus $R$:
\[
\ell_k=\ell_k(R),\qquad T_k=\mathbb{L}_k/\ell_k\mathbb{Z}^3,\qquad T^0_k=\mathbb{Z}^3/\ell_k\mathbb{Z}^3,\qquad T^{\mathrm{f}}_k=T_k\times O_{\mathrm{cub}},\qquad n_k=\ell_k^3;
\]
moreover $k_{*}=k_{*}(R)$ and $n=R/\eta=\ell_0$.
At $k=0$ the symbol $e_k$ denotes $q$ (Section~I.2.5).

The lattice mass is $\boldsymbol{\mu}:=\bar{m}\cdot\eta+m_{\mathrm{in}}$, and $M_{\mathrm{ch}}:=\boldsymbol{\mu}+\frac{\rho}{2}\cdot\mathrm{sw}$ is the mass parameter (Section~I.2.5).
No condition is imposed on the sign or on the size of $\boldsymbol{\mu}$ (see (I.K10)).
Without an argument, $D_0$ is the operator $D_A$ of (I.2.3) of level $0$ at $A=0$; with an argument, $D_0(B)=D_B+\bar{m}\cdot\eta$ is the operator of level $0$ of Section~I.3.9.

The spaces of the real functions on the sites and of the real bond functions of $T^0_k$ are denoted $C^0(T^0_k)$ and $C^1(T^0_k)$, respectively; $C^1(T^0_0)=C^1(n)$.
A constant bond function is named by its three components; in this way a constant bond function on $T^0_k$ is also one on $\mathbb{Z}^3/n\mathbb{Z}^3$.
For $B\in C^1(T^0_k)$, $\mathrm{mean}(B)$ is the constant bond function with the components
\[
\ell_k^{-3}\sum_{y\in T^0_k}B_{\mu}(y).
\]
The twist connection $c_{\theta}$ is the constant bond function on the unit lattice $T^0_{k_{*}}$ with the components $\frac{\theta_{\mu}}{e_{*}\cdot\ell_*}$, and $H_{\theta}$ is the constant one-cochain on $\mathbb{Z}^3/n\mathbb{Z}^3$ with the components $\frac{\theta_{\mu}}{q\cdot n}$ (Section~I.6.6).

In this chapter $\mathrm{Jm}$ is the map of \cite[Theorem~14.17\,(b)]{QEDI}, denoted $\mathsf{J}$ there and written out in (Q7.4)\,(b) below, and $\mathrm{Dm}$ denotes a matrix, as in Proposition~46\,(e).
Behind a number, as in $L\mathbb{N}$, $L^k\mathbb{N}$, $L^{k_{*}}\mathbb{N}$ and $L^ja\mathbb{N}$, the symbol $\mathbb{N}$ is the set $\{1,2,3,\dots\}$ of the positive integers, as in \cite[Section~14.1.1]{QEDI}; the letter $N$ is the index of the spacing.

\subsubsection*{P1. Prerequisites}

The proof rests on the following facts.

\paragraph{(Q1) Proposition~26.6\,(b), (c).}
It is applied to the data (FL): clause (b) at every level $k<k_{*}$, clause (c) at $k=k_{*}$.
The data (OS) contain (ST1)$^{\mathrm{ch}}$ (Section~I.1.8).
One reads off the following.
\begin{itemize}
\item[(q1)] $\ell_k\in\mathrm{Cl}$ for $k\le k_{*}$, and $\ell_k/L=\ell_{k+1}\in\mathrm{Cl}$ for $k<k_{*}$.
In particular $n=\ell_0\in\mathrm{Cl}$, that is, $n/(LM)$ is an odd integer $\ge13$ (Section~I.1.1).
Hence $n$ is odd, because $L$ is odd and $M$ is a power of $L$, and $n\ge13LM\ge4$; this is the setting of Sections~I.7.1 and I.7.3.
\item[(q2)] $0<e_k\le e_0$ for $k\le k_{*}$.
For $k<k_{*}$ the data (OS) give $0<e_k<L^{\frac{1}{2}}\cdot e_k\le e_0$; for $k=k_{*}$ the bound is part of (SF).
\item[(q3)] The data of every level $k\le k_{*}$ are (SF): for $k<k_{*}$ the data (OS) contain $\ell_k\in\mathrm{Cl}$ and $\mu_k\in(0,1/L]$, and (q2) holds; for $k=k_{*}$ this is clause (c).
\item[(q4)] For $k\le k_{*}$ one has $n=\ell_0=L^k\cdot\ell_k$ with the integer $\ell_k$, so that $n\in L^k\mathbb{N}$.
\end{itemize}

\paragraph{(Q2) Proposition~28.1.}
The pair is controlled on $R$, that is, $|m_k|\le w^{\mathrm{m}}_k$ for $0\le k\le k_{*}$ (Section~G4.0.5).
Therefore Proposition~28.1 applies at every $k\le k_{*}$.

\paragraph{(Q3) Proposition~24.1\,(c).}
It is applied in the form of its display (2.4x)$^{\mathrm{ch}}$, together with the display (8.6)$^{\mathrm{ch}}$ of Definition~24.D\,(D6) (Section~G3.4.1).
The setting is (ST1)$^{\mathrm{ch}}$ (Section~G3.4.0) at a level $k\ge0$.

\paragraph{(Q4) Proposition~22.2\,(c).}
It is applied in the flow, that is, with $\ell=n$, $\mu=\bar{m}\cdot\eta$ and $\mathfrak{P}_0=0$.
It gives $\det\left(D_0+\bar{m}\cdot\eta\right)\ne0$, and the formula
\[
y^{x}_{0}[\varepsilon^{\mathrm{v}}_{\mathrm{in}},m_{\mathrm{in}},0](A_0;a)=\exp\left(-\varepsilon^{\mathrm{v}}_{\mathrm{in}}\cdot n^3\right)\cdot\frac{\det\left(D_{A_0+a}+\bar{m}\cdot\eta+m_{\mathrm{in}}\right)}{\det\left(D_0+\bar{m}\cdot\eta\right)}
\]
at $A_0\in\mathrm{Dom}_0$ and $a\in\mathfrak{X}_0$.
In detail, by (I.L0.22), (I.L0.4), (I.L0.5) and (I.B1),
\begin{multline*}
\mathscr{F}^{x}_{0}[\varepsilon^{\mathrm{v}}_{\mathrm{in}},m_{\mathrm{in}},0](A_0,a;\Psi)
=\left(\det\left(D_0+\bar{m}\cdot\eta\right)\right)^{-1}\\
\cdot\exp\left\{-\left\langle\bar{\Psi},\left(D_{A_0+a}+\bar{m}\cdot\eta\right)\Psi\right\rangle^{\mathrm{ch}}-m_{\mathrm{in}}\cdot\left\langle\bar{\Psi},\Psi\right\rangle^{\mathrm{ch}}-\varepsilon^{\mathrm{v}}_{\mathrm{in}}\cdot n^3\right\},
\end{multline*}
where
\[
\left\langle\bar{\Psi},O\Psi\right\rangle^{\mathrm{ch}}=\frac{1}{48}\sum_{c\in\mathbb{Z}^3/n\mathbb{Z}^3}\ \sum_{g\in O_{\mathrm{cub}}}\ \sum_{\beta=1}^{4}\bar{\Psi}_{\beta}(c,g)\cdot\left(O\Psi\right)_{\beta}(c,g);
\]
$\int\mathrm{D}\Psi$ is the integral (I.3.5) with the weight $\nu=\frac{1}{48}$ over the $192n^3$ pairs of generators $\Psi_{\beta}(c,g)$, $\bar{\Psi}_{\beta}(c,g)$; and
\[
\int\mathrm{D}\Psi\,\exp\left(-\left\langle\bar{\Psi},O\Psi\right\rangle^{\mathrm{ch}}\right)=\det O
\]
at $O=D_{A_0+a}+\bar{m}\cdot\eta+m_{\mathrm{in}}$, a linear map of the space $\mathfrak{S}^{\mathrm{ch}}(1,n)$, whose dimension is $192n^3$.

\paragraph{(Q5) Proposition~16.3\,(a).}
The two equalities $\varepsilon(\mathscr{E})=\varepsilon(\mathscr{E}|_0)$ and $m(\mathscr{E})=m(\mathscr{E}|_0)$ of this clause hold for the data (SF) without the mass and for $\mathscr{E}\in K^{x,C}_{k}$.

\paragraph{(Q6) Fields in the domains.}
This is \cite[Proposition~5.38\,(i), (ii)]{QEDI}, in the form of \cite[Proposition~15.30\,(i), (ii)]{QEDI}.
Setting: an integer $k\ge0$, a side $\ell\in\mathrm{Cl}$ and a number $e_k>0$.
Hypotheses:
\begin{gather*}
8\,C_{B}\cdot p_k\le e_k^{-\frac{3}{4}+3\varepsilon},\qquad 8\,e_k^{\frac{3}{4}-\frac{7\varepsilon}{2}}\le1,\qquad e_k<1,\\
p_{0,k}\le e_k^{-\frac{\varepsilon}{4}},\qquad 2\,C_{w}^2\cdot c_{\gamma}^2\cdot c_{C}\cdot e_k^{\frac{\varepsilon}{4}}\le1,
\end{gather*}
with numbers $C_{w}\ge C_{P}$ and $0<\gamma_{\mathrm{w}}\le\gamma_{P}$, the number $c_{\gamma}$ being formed with $\gamma_{\mathrm{w}}$.
Then:
\begin{itemize}
\item[(i)] $\left\{A\in C^1(T^0_k):|dA|_{\infty}\le p_k\right\}\subset\mathrm{Dom}_k$;
\item[(ii)] if moreover $\ell/L\in\mathrm{Cl}$, and $\mathrm{Dom}_{k+1}$ is formed with the side $\ell/L$ and with $e_{k+1}=L^{\frac{1}{2}}\cdot e_k$, then $A'\in\mathrm{Dom}_{k+1}$ and $\zeta\in\mathrm{Box}_{k}$ imply $A^{\min}_{k}(A')+\mathrm{Sq}_{k}\zeta\in\mathrm{Dom}_k$.
\end{itemize}
In the application, $C_{w}$, $\gamma_{\mathrm{w}}$, $c_{\gamma}$, $c_{C}$, $C_{B}$, $C_{P}$, $\gamma_{P}$ are the numbers of Section~I.8.2.
Their definitions are rows of the table of Proposition~14.2\,(a); the conditions (A1), (A2), (A8), (A10), (L5), (M1), (M3), (S1) are rows of the table of Proposition~14.2\,(b); and (C1) is a line of Section~I.8.2.
\begin{itemize}
\item[(6a)] \emph{$C_{P}\le C_{w}$.}
By the definition of $C_{P}$ \cite[Section~15.7]{QEDI} one has $C_{P}\ge4\exp(2)>0$.
Next, $C^{P}=116\,C_{P}\cdot M^{402\alpha_{0}}\ge C_{P}$, because $M\ge L^2\ge1$ by (M1) and $\alpha_{0}=\frac{1}{54}>0$.
Further $C_{w}^{(0)}=\max\left(C^{F},C^{P}\right)\ge C^{P}$ and $C_{w}^{(1)}=\max\left(C_{w}^{(0)},4\exp(2\gamma_{\mathrm{w}})\cdot M^{27\alpha_{0}}\right)\ge C_{w}^{(0)}$.
Finally $C_{w}=64\,K_{0}\cdot C_{w}^{(1)}\ge C_{w}^{(1)}$, because $C_{w}^{(1)}>0$ and
\[
64\,K_{0}=\frac{64}{36}\exp(\kappa_{0})=\frac{16}{9}\cdot72^{54}\ge1\qquad\left(K_{0}=\frac{\exp(\kappa_{0})}{36},\ \kappa_{0}=54\log72\right).
\]
\item[(6b)] \emph{$0<\gamma_{\mathrm{w}}\le\gamma_{P}$.}
The inequality $\gamma_{\mathrm{w}}\cdot M\ge16$ of (M3) and $M>0$ give $\gamma_{\mathrm{w}}>0$.
Moreover $\gamma_{\mathrm{w}}=\min\left(\gamma^{F},\gamma^{P}\right)\le\gamma^{P}=\frac{\gamma_{P}}{2}$.
So $\gamma_{P}\ge2\gamma_{\mathrm{w}}>0$, and $\frac{\gamma_{P}}{2}\le\gamma_{P}$.
\item[(6c)] \emph{The number $c_{\gamma}$.}
The number $c_{\gamma}$ of Section~I.5.8 is the supremum of the sums
\[
\sum_{t'}\nu(t')\exp\left(-\gamma_{\mathrm{w}}\,d(t,t')\right)
\]
over all indices $t$ (fine bonds, unit bonds and basis elements are indices, Section~G4.0.4) and, among others, over the two index sets for $t'$ that consist of the positively oriented unit bonds and of the basis $\{\Upsilon\}$, with $\nu=1$.
These are the two sums that the proof of \cite[Proposition~15.30\,(iii)]{QEDI} bounds by $c_{\gamma}$: the one with $t=\Upsilon$ and $t'=\Upsilon'$, and the one with $t$ a fine bond $b$ and $t'$ the unit bonds $b'$.
\item[(6d)] \emph{The five hypotheses.}
Let $v=v_k=-\log e_k$, so that $p_k=v^p$ and $p_{0,k}=v^{p_0}$ (Section~I.6.1).
The first, the second, the fourth and the fifth hypothesis are, in this order, the third and the fourth inequality of (A10), the first inequality of (A1), and (A2).
These conditions are imposed for all $v\ge-\log e_0$ (Proposition~14.2\,(b)), and $v_k\ge-\log e_0$ at every level $k\le k_{*}$ by (q2).
As for the third hypothesis, the first inequality of (A8), $\exp(-\varepsilon v)\le\frac{1}{4}$, reads $e_0^{\varepsilon}\le\frac{1}{4}$ at $v=-\log e_0$, and $\varepsilon>1-\alpha>0$ by (S1) and (C1).
Hence $e_0<1$, and $e_k\le e_0<1$ by (q2).
\item[(6e)] \emph{The parameters in \cite[Propositions~15.29 and 15.30]{QEDI}.}
The proofs of these Propositions use the following four facts on $\varepsilon$, $L$, $M$ and $\alpha$, and each of them holds here.
First, $\frac{3\varepsilon}{2}\le\frac{1}{8}$, since $\varepsilon\le\bar{\varepsilon}=\frac{1}{200}$ by (S1).
Second, $L^{\frac{3}{4}}\ge\frac{8}{3}$: by (L5), $L\ge5$, so that $L^3\ge125>\frac{4096}{81}=\left(\frac{8}{3}\right)^4$.
Third, the $M$-cubes nest in the $LM$-cubes: $L$ is odd by (L5) and $M$ is a power of $L$ by (M1); so $M$ is odd, the $LM$-cube with the centre $z\in LM\mathbb{Z}^3$ is the union of the $L^3$ $M$-cubes with the centres $z+My$, $y\in\left\{-\frac{L-1}{2},\dots,\frac{L-1}{2}\right\}^3$, and no site lies on a face (Section~I.1.1).
Fourth, $0\le\alpha<1$, which is the range of \cite[Theorem~15.14]{QEDI} (see also \cite[Proposition~5.32]{QEDI}); that theorem enters Step~3 of the proof of \cite[Proposition~15.29]{QEDI} and the proof of \cite[Proposition~15.30\,(iii)]{QEDI}, and the range holds by (C1).
\item[(6f)] \emph{Two steps in the proof of \cite[Proposition~15.29]{QEDI}.}
Clause~(i) rests on \cite[Proposition~15.29]{QEDI}, which is also stated as \cite[Proposition~5.37]{QEDI}.
In the present setting two steps of its proof take the following form.
Step~3 applies the bound of the paragraph `Kernels' of \cite[Section~15.1.1]{QEDI} on the set $\square^{\dagger}$, with the supremum over the bonds of $\square^{\dagger}$.
In this form the bound reads as follows.
Let $O$ map the bond functions on $T^0_k$ to the bond functions on $T_k$,
\[
(OA)(b)=\sum_{b'}O(b,b')\cdot A(b'),
\]
where $b$ is a positively oriented fine bond and $b'$ ranges over the positively oriented unit bonds, and let the moduli of the three kernels $O(b,b')$, $(\partial_{\mathfrak{l}}O)(b,b')$ and $(\delta_{\alpha}\partial_{\mathfrak{l}}O)((b,\tilde{b});b')$ be at most $\mathrm{maj}(b,b')$.
Then, for every set $Z$ of sites of $T_k$, each of the three norms of $OA$ on $Z$ is at most
\[
\sup_{b}\ \sum_{b'}\mathrm{maj}(b,b')\cdot|A(b')|,
\]
where $b$ ranges over the fine bonds contained in $Z$.
Indeed, let $b=(c,c+\eta_ke_{\mu})$ and $\tilde{b}=(\tilde{c},\tilde{c}+\eta_ke_{\mu})$ with $0<|c-\tilde{c}|\le1$.
By the definitions of the three kernels,
\begin{align*}
(OA)_{\mu}(c)&=\sum_{b'}O(b,b')\cdot A(b'),\\
\left(\partial_{\mathfrak{l}}(OA)_{\mu}\right)(c)&=\sum_{b'}(\partial_{\mathfrak{l}}O)(b,b')\cdot A(b'),\\
|c-\tilde{c}|^{-\alpha}\left(\left(\partial_{\mathfrak{l}}(OA)_{\mu}\right)(\tilde{c})-\left(\partial_{\mathfrak{l}}(OA)_{\mu}\right)(c)\right)&=\sum_{b'}(\delta_{\alpha}\partial_{\mathfrak{l}}O)((b,\tilde{b});b')\cdot A(b').
\end{align*}
The three norms of $OA$ on $Z$ are the suprema of the moduli of these three quantities over the bonds and the differences inside $Z$ \cite[Section~5.3.1]{QEDI}, and in each of them $b$ is contained in $Z$.
Step~1 uses that $\square^{\dagger}$ does not wrap, in the sense of Section~I.1.1: the cube of $\mathbb{L}_k$ of the side $s_d=(2L+1)M$ is a lift of $\square^{\dagger}$, its sup-diameter is $s_d-\eta_k$, and
\[
(s_d-\eta_k)+\eta_k=s_d\le3LM<13LM\le\ell.
\]
\end{itemize}
By (q1) and (Q1), (i) holds at the level $k_{*}$, and (ii) holds at every level $k<k_{*}$, with $\ell=\ell_k$, $\ell/L=\ell_{k+1}$ and $e_{k+1}=L^{\frac{1}{2}}\cdot e_k$.
\emph{Level $0$.}
At $k=0$, (ii) is applied with $e_k=q$ and $H_{0}=1$, so that
\[
\mathrm{Dom}_0=\left\{A\in C^1(n):A\in\tfrac{1}{4}\tilde{R}_0\right\};
\]
the photon objects $\mathrm{Box}_{0}$, $A^{\min}_{0}$, $\mathrm{Sq}_{0}$ and the domains are formed as at the levels $k\ge1$ (Section~I.5.10, part~H).

\paragraph{(Q7) Block fields and their coordinates.}
The facts (Q7.0)--(Q7.5) concern \cite[Sections~14.1--14.5]{QEDI}, whose standing datum is an odd integer $L\ge3$; here it is given by (L5).

\paragraph{(Q7.0) One set of photon objects.}
The objects $V_j$, $\{\Upsilon\}$, $C$, $H_{j}$, $\Delta_j$, $\mathrm{Dt}_{j}=C^{T}\Delta_jC$, $C_{j}=\mathrm{Dt}_{j}^{-1}$, $C_{j}^{1/2}$, $\mathrm{Sq}_{j}=C\,C_{j}^{1/2}$, $A^{\min}_{j}$, $\mu_I$, $\mathrm{Box}_{j}$, $\mathrm{Sl}$, $A_{0}(F^{\mathrm{u}})$, $A_{j}$, $A_{*}$, $\zeta_j$, $\mathrm{Af}_0$ of Section~I.5.7, of Section~G4.0.6 and of (Q3) are the objects of \cite[Section~5.2.5 and Definition~5.11]{QEDI}; there $\mathrm{Af}_0$ is denoted $A^{\downarrow}_0$.
Those among them that occur in \cite[Section~14]{QEDI} are the objects of the same names there.
For the sides of (Q1) ($\ell_j$ a positive integer for $j\le k_{*}$, and $\ell_j=L\cdot\ell_{j+1}\in L\mathbb{N}$ for $j<k_{*}$) they exist, and those other than $\{\Upsilon\}$ and $\mathrm{Box}_{j}$ are unique, by the following results.
\begin{itemize}
\item $H_{j}$ and $\Delta_j$: \cite[Proposition~14.8\,(ii)]{QEDI} for the tower $(\eta_j,j,\ell_j)$, and the sentence after that Proposition.
\item $V_j$: \cite[Section~14.4.1]{QEDI}, where it is also shown that $\{\Upsilon\}$ is an orthonormal basis of $V_j$.
\item $\{\Upsilon\}$ is formed, by the display that defines it in \cite[Section~5.2.5]{QEDI}, with the basis $\{e^{0}_{i}\}$ of that section.
A basis of this kind exists by \cite[Proposition~15.21\,(i)]{QEDI} and is not unique; $\{e^{0}_{i}\}$ is the one fixed with the scheme (Section~I.6.2), the same at all levels, on all tori and in all chapters.
The basis $\{\Upsilon\}$, the coordinates $\zeta(\Upsilon)$ and $\mathrm{Box}_{j}$ are determined by $\{e^{0}_{i}\}$.
\item $C$: \cite[Lemma~14.13\,(b)]{QEDI} and the paragraph `The map $C$' below.
\item $\mathrm{Dt}_{j}$, $C_{j}$, $C_{j}^{1/2}$ and $\mathrm{Sq}_{j}$: \cite[Section~14.4.2]{QEDI} with \cite[Proposition~14.14\,(a), (b)]{QEDI}; these are the displays of \cite[Section~5.2.5]{QEDI}, formed with $\Delta_j$ and $C$.
\item $A^{\min}_{j}$: \cite[Theorem~14.15\,(a)]{QEDI}.
\item $A_{j}$ and $\mathrm{Af}_0$: \cite[Section~14.5.1]{QEDI}, namely the display that defines $A_{j}$ and the recursion displayed before \cite[Theorem~14.16]{QEDI}.
\item $\mathrm{Sl}$ and $A_{*}:=A_{k_{*}}$: the first paragraph of \cite[Section~14.5.2]{QEDI}.
\item $A_{0}(F^{\mathrm{u}})$: (Q7.4)\,(a).
\item $\zeta_j$: (Q7.3) and the first paragraph of the proof of (t7).
\end{itemize}

\emph{The map $C$.}
The map $C$ is defined in \cite[Section~5.2.5]{QEDI} and once more, in other terms, in \cite[Lemma~14.13\,(b)]{QEDI}; the two definitions give the same map.
Indeed, in \cite[Section~5.2.5]{QEDI} as in \cite[Section~14.4.1]{QEDI} a bond function on $T^0_j$ is read as an $\ell_j\mathbb{Z}^3$-periodic bond function on $\mathbb{Z}^3$, and the $L$-blocks $B(y)$, $y\in L\mathbb{Z}^3$, the bonds, the faces and the paths are those of $\mathbb{Z}^3$, counted modulo $\ell_j\mathbb{Z}^3$; in \cite[Section~14.4.1]{QEDI} the blocks and the paths are those of \cite[Section~14.1.2]{QEDI} at the mesh~$1$.
In both sections the faces are indexed by the pairs $(y,\mu)$ of a point $y$ and a direction $\mu$, and the face bonds of $(y,\mu)$ are the bonds, not interior, that join $B(y)$ and $B(y+Le_{\mu})$.
Hence the interior bonds, the faces $\langle y,y+Le_{\mu}\rangle$ and their $L^2$ face bonds of \cite[Section~5.2.5]{QEDI} are the inner bonds, the pairs $(y,\mu)$ and their $L^2$ face bonds of \cite[Section~14.4.1]{QEDI}.
For $j<k_{*}$ these objects are moreover, modulo $\ell_j\mathbb{Z}^3$, those formed on the torus $T^0_j$ itself: \cite[Section~5.2.5]{QEDI} states this for every side $\ell\ge3L$, and $\ell_j\ge13L^2M\ge3L$ by (q1).
The same is seen directly: the set $B(y)\cup B(y+Le_{\mu})$ has the sup-diameter $2L-1$, and $(2L-1)+1<13L^2M\le\ell_j$; so this set does not wrap (Section~I.1.1, read on the unit lattice, with the mesh $1$), and the notions on $T^0_j$ and on $\mathbb{Z}^3$ correspond.
Off the central bonds both maps are the identity: in \cite[Section~5.2.5]{QEDI}, as in \cite[Lemma~14.13\,(b)]{QEDI}, the field $C\mathsf{Z}$ coincides with $\mathsf{Z}$ on every bond that is not central.
The central bond of the face $\langle y,y+Le_{\mu}\rangle$ in \cite[Section~5.2.5]{QEDI} is the face bond on the line through the two centres $y$ and $y+Le_{\mu}$.
On this line the last site of $B(y)$ is $y+\frac{L-1}{2}e_{\mu}$, and the first site of $B(y+Le_{\mu})$ is $y+\frac{L+1}{2}e_{\mu}$.
So the central bond is the bond
\[
b(y,\mu)=\left(y+\frac{L-1}{2}e_{\mu},\ y+\frac{L+1}{2}e_{\mu}\right)
\]
of \cite[Section~14.4.1]{QEDI}.
On this bond the definition of \cite[Section~5.2.5]{QEDI} gives
\[
(C\mathsf{Z})(b(y,\mu))=-\sum_{\substack{b\text{ face bond of this face}\\ b\text{ not central}}}\mathsf{Z}(b)-\sum_{b}\frac{\mathrm{npt}(b)}{L}\cdot\mathsf{Z}(b),
\]
the second sum being over the interior bonds $b$ of the direction $\mu$ in $B(y)$ and in $B(y+Le_{\mu})$, and $\mathrm{npt}(b)$ being the number of the paths $\boldsymbol{\Gamma}_{y'',y''+Le_{\mu}}$, $y''\in B(y)$, that contain $b$ (there it is denoted $n_c(b)$, and the starting point $y''$ is written $x$).
The display of $(C\mathsf{Z})(b(y,\mu))$ in \cite[Lemma~14.13\,(b)]{QEDI} is the same, with the weight $w(b)$ in place of $\mathrm{npt}(b)$; here $w(b)=i$ for the bond from the $i$-th to the $(i+1)$-th site of its line in $B(y)$, and $w(b)=L-i$ for the bond from the $i$-th to the $(i+1)$-th site of its line in $B(y+Le_{\mu})$.
The path that starts at the $i'$-th site of a line in $B(y)$, $1\le i'\le L$, contains the bonds $(i,i+1)$, $i'\le i\le L-1$, of this line in $B(y)$, one face bond, and the bonds $(i,i+1)$, $1\le i\le i'-1$, of the continuation of the line in $B(y+Le_{\mu})$.
So the bond $(i,i+1)$ in $B(y)$ lies on the $i$ paths with $i'\le i$, and the bond $(i,i+1)$ in $B(y+Le_{\mu})$ lies on the $L-i$ paths with $i'\ge i+1$:
\[
\mathrm{npt}(b)=w(b).
\]
Hence $C$, $\mathrm{Dt}_{j}$, $C_{j}^{1/2}$ and $\mathrm{Sq}_{j}$ are the same operators in \cite[Section~5.2.5]{QEDI} and in \cite[Section~14.4]{QEDI}.

\paragraph{(Q7.1) The Hodge decomposition.}
This is the display (Hodge) of \cite[Section~14.1.1]{QEDI}.
On each $T^0_k$ one has the orthogonal decomposition
\[
C^1=dC^0\oplus\{\text{constants}\}\oplus{V_{\mathrm{t}}},
\]
and $\ker d=dC^0\oplus\{\text{constants}\}$ on $C^1$; the map $d$ is injective on ${V_{\mathrm{t}}}$; $\mathrm{mean}(d\boldsymbol{\lambda})=0$, and the mean of a constant is the constant.
It follows that the elements of ${V_{\mathrm{t}}}$ have mean zero, that $d$ maps ${V_{\mathrm{t}}}$ bijectively onto $dC^1$, and that a closed bond function with mean zero is exact.
At $k=0$ these are the spaces ${V_{\mathrm{g}}}$, ${V_{\mathrm{h}}}$, ${V_{\mathrm{t}}}$ of Section~I.7.1; $d$ is the coboundary of Section~I.1.1, and the pairing is the counting pairing.

\paragraph{(Q7.2) \cite[Proposition~14.4\,(c)]{QEDI}, with $a=1$, $j=k_{*}$, $\ell=n$.}
Let $Q_{k_{*}}$ be the $k_{*}$-fold composition of $Q$ on $\mathbb{Z}^3/n\mathbb{Z}^3$, a map of $C^1(n)$ into the real bond functions on $L^{k_{*}}\mathbb{Z}^3/n\mathbb{Z}^3$.
Then
\[
\mathrm{mean}\left(Q_{k_{*}}A_0\right)=\mathrm{mean}(A_0)\qquad\left(A_0\in C^1(n)\right);
\]
here the mean of a bond function $B$ on $L^{k_{*}}\mathbb{Z}^3/n\mathbb{Z}^3$ is the constant with the components
\[
\left(\frac{L^{k_{*}}}{n}\right)^3\sum_{z}B_{\mu}(z),
\]
$z$ ranging over the sites of this lattice.
The hypothesis $\ell\in L^ja\mathbb{N}$ of that Proposition reads $n\in L^{k_{*}}\mathbb{N}$; it holds by (q4).

\paragraph{(Q7.3) \cite[Section~14.5.1, Theorem~14.16\,(a) and Proposition~14.14\,(b)]{QEDI}.}
Data: an integer $k\ge1$, $n\in L^k\mathbb{N}$, and $\ell_j=n/L^j$.
For $A_0\in C^1(n)$ let
\[
A_{j}:=\left(Q_jA_0\right)^{[j]}\in C^1(T^0_j)\qquad(0\le j\le k),
\]
where $Q_j$ is the $j$-fold composition of $Q$ on $\mathbb{Z}^3/n\mathbb{Z}^3$, $Q_0=1$, and $B^{[j]}(b)=L^{\frac{j}{2}}B(L^jb)$; these are the fields $A_{j}$ of \cite[Definition~5.11]{QEDI}.
\begin{itemize}
\item By the first sentence of \cite[Theorem~14.16\,(a)]{QEDI}, the map $(A,\zeta)\mapsto\mathrm{Af}_0(A,\zeta)$, defined by
\[
\mathrm{Af}_k:=A,\qquad\mathrm{Af}_j:=A^{\min}_{j}\left(\mathrm{Af}_{j+1}\right)+\mathrm{Sq}_{j}\zeta_j\qquad(j=k-1,\dots,0),
\]
is linear on $C^1(T^0_k)\times\prod_{j<k}V_j$.
\item By the second sentence, if $A_0=\mathrm{Af}_0(A,\zeta)$, then $A_{j}=\mathrm{Af}_j$ for $0\le j\le k$; in particular $A=A_{k}$, and $\mathrm{Sq}_{j}\zeta_j=A_{j}-A^{\min}_{j}(A_{j+1})$ for $j<k$.
\item By \cite[Proposition~14.14\,(b)]{QEDI}, $\mathrm{Sq}_{j}$ is a linear bijection of $V_j$ onto $\left\{Z\in C^1(T^0_j):QZ=0,\ \tau Z=0\right\}$; so it is injective on $V_j$.
\end{itemize}
The first sentence is applied for every $1\le k\le k_{*}$, the second one at $k=k_{*}$; the data hold by (Q1) and (q4).

\paragraph{(Q7.4) \cite[Theorem~14.17\,(a), (b), (e)]{QEDI}.}
This theorem is also stated as \cite[Proposition~5.53]{QEDI}.
The data are those of (Q7.3), at $k=k_{*}\ge1$.
The space $F_N=dC^1(n)$ and its standard Gaussian measure $\gamma_N$ (Section~I.7.3) are those of the first paragraph of \cite[Section~14.5.2]{QEDI}: $d$ is the coboundary of Section~I.1.1, which is the coboundary of \cite[Section~14.1.1]{QEDI} at the mesh $1$, and the inner product of $C^2(n)$ is the counting pairing of Section~I.7.1, which is the pairing of \cite[Section~14.1.1]{QEDI} at the mesh $1$.
\begin{itemize}
\item[(a)] $d$ maps $\mathrm{Sl}$ bijectively onto $F_N$.
\item[(b)] For $F^{\mathrm{u}}\in F_N$ let $A_{0}(F^{\mathrm{u}})\in\mathrm{Sl}$ be the potential with $dA_{0}(F^{\mathrm{u}})=F^{\mathrm{u}}$, and let $(A_{*};\zeta_0,\dots,\zeta_{k_{*}-1})$ be its coordinates:
\[
A_{0}(F^{\mathrm{u}})=\mathrm{Af}_0(A_{*},\zeta).
\]
The map $\mathrm{Jm}:F^{\mathrm{u}}\mapsto(A_{*};\zeta)$ is a linear bijection of $F_N$ onto ${V_{\mathrm{t}}}(T^0_{k_{*}})\times\prod_{j<k_{*}}V_j$, with the inverse $(A,\zeta)\mapsto d\,\mathrm{Af}_0(A,\zeta)$.
\item[(e)] Under $\gamma_N$ the variables $\zeta_j(\Upsilon)$ are independent standard normal variables, independent of $A_{*}$; the law of $A_{*}$ is $\mathrm{d}\mathsf{m}_{*}$.
\end{itemize}
In (b) and (e) the letters $A_{*}$, $\zeta_j$ denote the entries of $\mathrm{Jm}\,F^{\mathrm{u}}$.
That these are the variables $A_{*}$, $\zeta_j$ of \cite[Definition~5.11]{QEDI} is the first assertion of (t7).

\paragraph{(Q7.5) \cite[Corollary~14.18]{QEDI} without its last sentence, the one that begins with ``In particular''.}
This is also stated as \cite[Proposition~5.54]{QEDI}, likewise without its last sentence, which also begins with ``In particular''.
Data as in (Q7.4).
If $A^{c}\in C^1(T^0_{k_{*}})$ is closed, then for all $(A,\zeta)\in C^1(T^0_{k_{*}})\times\prod_jV_j$
\[
\mathrm{Af}_0(A+A^{c},\zeta)=\mathrm{Af}_0(A,\zeta)+L^{-\frac{k_{*}}{2}}\cdot\mathrm{mean}(A^{c})+d\boldsymbol{\lambda}
\]
with a real function $\boldsymbol{\lambda}$ on $\mathbb{Z}^3/n\mathbb{Z}^3$.

\paragraph{(Q8) Classical facts.}
Fubini's theorem, \cite[Theorem~8.8]{Rudin1987}.
Image measures: for a Borel map between finite-dimensional real vector spaces and a finite Borel measure on the first space, the integral of a bounded Borel function $G$ with respect to the image measure is the integral, with respect to the given measure, of the composition of $G$ with the map.
For indicators this is the definition of the image measure; the general case follows by linearity and monotone limits \cite[Theorems~1.17, 1.26]{Rudin1987}.
A continuous function on a compact metric space is bounded and uniformly continuous \cite[Theorems~4.15, 4.19]{Rudin1976}.

\subsubsection*{P2. The objects of the data; the two displays}

\paragraph{(P2.1) The numbers $y^{x}_{k}$ are defined.}
Let $k\le k_{*}$.
By (q3) the data of level $k$ are (SF), so the paragraph `Settings' of Section~I.5.7 applies.
\begin{itemize}
\item For $(\mathscr{A},\mathscr{C})\in\tilde{R}_k\times\mathfrak{X}_k$ the total connection $\mathscr{T}=\mathscr{A}+\mathscr{C}$ lies in $\tilde{G}_k$ \cite[Lemma~17.4\,(a)]{QEDI}, and so does $-\mathscr{T}$ \cite[Lemma~17.6]{QEDI}.
\item The operators $H_{k}(\mathscr{T})$, $\bar{H}_{k}(\mathscr{T})=\mathsf{C}^{-1}H_{k}(-\mathscr{T})\mathsf{C}$ and $D_k(\mathscr{T})$ exist: for $k\ge1$ by Proposition~14.1\,(i), with the displays of Section~I.3.5; for $k=0$ they are $1$, $1$ and $D_{\mathscr{T}}+\bar{m}\cdot\eta$ by Section~I.3.9.
\item $\det D_k(0)\ne0$: for $k\ge1$ by Proposition~19\,(d); for $k=0$ by Proposition~19\,(a), which is also contained in (Q4).
\item For $A\in\mathrm{Dom}_k$ the field $H_{k}A$ lies in $\frac{1}{4}\tilde{R}_k$; this is the definition
\[
\mathrm{Dom}_k=\left\{A\in C^1(T^0_k):H_{k}A\in\tfrac{1}{4}\tilde{R}_k\right\}
\]
(\cite[Definition~5.16]{QEDI}; Section~I.5.3).
The set $\frac{1}{4}\tilde{R}_k$ is contained in $\tilde{R}_k$, because its thresholds are one quarter of those of $\tilde{R}_k$ (Section~I.5.3).
\end{itemize}
By (Q2), $\mathscr{E}_k$ lies in $K^{x}_{k}$, a subset of $K^{x,C}_{k}$, and the $\mathfrak{P}_k(X,\cdot)$ are complex functions on $\mathfrak{X}_k$.
Consequently, for $A\in\mathrm{Dom}_k$ and $\mathscr{C}\in\mathfrak{X}_k$, the element $\mathscr{F}^{x}_{k}[\varepsilon^{\mathrm{v},t}_k,m_k,\mathscr{E}_k](A,\mathscr{C};\Psi)$ of Section~I.5.7, formed with $\mathfrak{P}_k$, is an even element of $\mathrm{Gr}[\Psi]$, and
\[
y^{x}_{k}[\varepsilon^{\mathrm{v},t}_k,m_k,\mathscr{E}_k](A;\mathscr{C})=\int\mathrm{D}\Psi\,\mathscr{F}^{x}_{k}[\varepsilon^{\mathrm{v},t}_k,m_k,\mathscr{E}_k](A,\mathscr{C};\Psi)
\]
is a complex number.
At $k=k_{*}$ this number is $Y^{x}_{*}(A;\mathscr{C})$.

\paragraph{(P2.2) The display of $Y^{x}_{*}$.}
Let $A\in\mathrm{Dom}_{k_{*}}$ and $\mathscr{C}\in\mathfrak{X}_{k_{*}}$, and put $\mathscr{A}:=H_{k_{*}}A$ and $\mathscr{T}:=\mathscr{A}+\mathscr{C}$.
The three definitions of $I^{x}_{k}$, $\mathrm{Act}^{x}_{k}$ and $\mathscr{F}^{x}_{k}$ in Section~I.5.7, at $k=k_{*}$ and with $n_{k_{*}}=\ell_*^3$, give
\begin{multline*}
\mathscr{F}^{x}_{k_{*}}[\varepsilon^{\mathrm{v},t}_{*},m_{*},\mathscr{E}_{*}](A,\mathscr{C};\Psi)
=\left(\det D_{*}(0)\right)^{-1}\cdot\exp\left\{-\left\langle\bar{\Psi},D_{*}(\mathscr{T})\Psi\right\rangle-m_{*}\cdot\left\langle\bar{\psi},\psi\right\rangle\right.\\
\left.{}+\sum_{X}\mathscr{E}_{*}(X;\mathscr{A},\mathscr{C},\psi,\chi)-\varepsilon^{\mathrm{v},t}_{*}\cdot\ell_*^3+\sum_{X}\mathfrak{P}_{*}(X,\mathscr{C})\right\},
\end{multline*}
where the fields
\[
\psi=H_{k_{*}}(\mathscr{T})\Psi,\qquad\chi=\delta_{\alpha,\mathscr{T}}\psi,\qquad\bar{\psi}=\bar{H}_{k_{*}}(\mathscr{T})\bar{\Psi},\qquad\bar{\chi}=\delta_{\alpha,-\mathscr{T}}\bar{\psi}
\]
are inserted (Section~I.5.4), and $X$ ranges over $D_{k_{*}}$.
The last two terms of the exponent are complex numbers.
A number commutes with every element of $\mathrm{Gr}[\Psi]$; hence the exponential of the sum of an even element $F$ and a number is the exponential of the number times $\exp\{F\}$.
Moreover $\int\mathrm{D}\Psi$ is linear.
So the factor $\exp\left\{-\varepsilon^{\mathrm{v},t}_{*}\cdot\ell_*^3+\sum_{X}\mathfrak{P}_{*}(X,\mathscr{C})\right\}$ stands in front of the integral:
\begin{multline*}
Y^{x}_{*}(A;\mathscr{C})=\left(\det D_{*}(0)\right)^{-1}\cdot\exp\left\{-\varepsilon^{\mathrm{v},t}_{*}\cdot\ell_*^3+\sum_{X}\mathfrak{P}_{*}(X,\mathscr{C})\right\}\\
\cdot\int\mathrm{D}\Psi\,\exp\left\{-\left\langle\bar{\Psi},D_{*}(\mathscr{T})\Psi\right\rangle-m_{*}\cdot\left\langle\bar{\psi},\psi\right\rangle+\sum_{X}\mathscr{E}_{*}(X;\mathscr{A},\mathscr{C},\psi,\chi)\right\}.
\end{multline*}
This is the display of $Y^{x}_{*}(A;\mathscr{C})$ among the definitions of Proposition~29.

Since $k_{*}\ge1$ by (Q1), the table of Section~I.3.6 shows that the generators are $\Psi_{\beta}(y)$, $\bar{\Psi}_{\beta}(y)$, $y\in T^0_{k_{*}}$, $\beta\in\{1,\dots,4\}$: these are $4\ell_*^3$ pairs, with the weight $\nu=1$.
Further,
\[
\left\langle\bar{\Psi},D_{*}(\mathscr{T})\Psi\right\rangle=\sum_{y,y'\in T^0_{k_{*}}}\bar{\Psi}(y)\cdot D_{*}(\mathscr{T})(y,y')\,\Psi(y'),
\]
and, by (I.1.3) and Section~I.3.6,
\[
\left\langle\bar{\psi},\psi\right\rangle=\int\mathrm{d}x\,\bar{\psi}(x)\cdot\psi(x)=\eta_{k_{*}}^3\cdot\frac{1}{48}\sum_{c\in T_{k_{*}}}\ \sum_{g\in O_{\mathrm{cub}}}\ \sum_{\beta=1}^{4}\bar{\psi}_{\beta}(c,g)\cdot\psi_{\beta}(c,g).
\]

\paragraph{(P2.3) The display (Y).}
The point $0$ lies in $\mathfrak{X}_{k_{*}}$.
At $\mathscr{C}=0$ one has $\mathscr{T}=\mathscr{A}$, and $\mathfrak{P}_{*}(X,0)=0$ by (Q2).
Further,
\[
\mathscr{E}_{*}(X;\mathscr{A},0,\psi,\chi)=\left(\mathscr{E}_{*}|_0\right)(X;\mathscr{A},\psi,\chi)
\]
by the definition of $\mathscr{E}|_0$ (Section~I.5.4), and $\mathscr{E}_{*}|_0=E_{*}$ by (Q2).
Insert this into the display of (P2.2), and move the number $\exp\left(-\varepsilon^{\mathrm{v},t}_{*}\cdot\ell_*^3\right)$ under the integral and into the exponential.
The result is the display (Y) of Proposition~29:
\[
Y_{*}(A)=\left(\det D_{*}(0)\right)^{-1}\int\mathrm{D}\Psi\,\exp\left\{-\left\langle\bar{\Psi},D_{*}(\mathscr{A})\Psi\right\rangle-m_{*}\cdot\left\langle\bar{\psi},\psi\right\rangle-\varepsilon^{\mathrm{v},t}_{*}\cdot\ell_*^3+\sum_{X}E_{*}(X;\mathscr{A},\psi,\chi)\right\}.
\]

\subsubsection*{P3. Level 0: the weight as a function of a complex one-cochain}

Let $B$ be a complex one-cochain on $\mathbb{Z}^3/n\mathbb{Z}^3$, that is, a complex bond function on $T_0=T^0_0$.
Put
\begin{equation}
\mathrm{wdn}(B):=\exp\left(-\varepsilon^{\mathrm{v}}_{\mathrm{in}}\cdot n^3\right)\cdot\frac{\det\left(D_B+\boldsymbol{\mu}\right)}{\det\left(D_0+\bar{m}\eta\right)},
\tag{29.1}
\end{equation}
where both determinants are determinants of linear maps of $\mathfrak{S}^{\mathrm{ch}}(1,n)$, a space of dimension $192n^3$.

\begin{named}{Lemma 1}
\begin{itemize}
\item[(i)] One has $\det\left(D_0+\bar{m}\eta\right)\neq0$, and $\mathrm{wdn}$ is an entire function of the $3n^3$ complex numbers $B_\mu(c)$, $c\in\mathbb{Z}^3/n\mathbb{Z}^3$, $\mu\in\{1,2,3\}$.
\item[(ii)] $\mathrm{wdn}(B+d\boldsymbol{\lambda})=\mathrm{wdn}(B)$ for every complex one-cochain $B$ and every real function $\boldsymbol{\lambda}$ on $\mathbb{Z}^3/n\mathbb{Z}^3$.
\item[(iii)] For every complex one-cochain $B$,
\[
\left|\mathrm{wdn}(B)\right|\le\exp\left(-\varepsilon^{\mathrm{v}}_{\mathrm{in}}\cdot n^3\right)\cdot\left|\det\left(D_0+\bar{m}\eta\right)\right|^{-1}\cdot\mathrm{rsm}(B)^{192n^3},
\]
where
\begin{multline*}
\mathrm{rsm}(B):=\left|M_{\mathrm{ch}}\right|+3(1+\rho)+6\sqrt{2}\left(\rho+\sqrt{2}\right)+3(1+\rho)\cdot\exp\left(q\cdot\left|\operatorname{Im}B\right|_\infty\right)\\
\le\left|\boldsymbol{\mu}\right|+41+6\exp\left(q\cdot\left|\operatorname{Im}B\right|_\infty\right).
\end{multline*}
In particular, for every real one-cochain $B_0$,
\begin{multline*}
\left|\mathrm{wdn}(B_0+a)\right|\le\mathrm{Bdn}\\
:=\exp\left(-\varepsilon^{\mathrm{v}}_{\mathrm{in}}\cdot n^3\right)\cdot\left|\det\left(D_0+\bar{m}\eta\right)\right|^{-1}\cdot\left(\left|\boldsymbol{\mu}\right|+41+6\exp\left(q\cdot\left|\operatorname{Im}a\right|_\infty\right)\right)^{192n^3}.
\end{multline*}
\item[(iv)] $y^{x}_{0}[\varepsilon^{\mathrm{v}}_{\mathrm{in}},m_{\mathrm{in}},0](A_0;a)=\mathrm{wdn}(A_0+a)$ for every $A_0\in\mathrm{Dom}_0$.
\item[(v)] $\mathrm{dens}^{\theta,a}_N(F^{\mathrm{u}})=\mathrm{wdn}\left(A(F^{\mathrm{u}})+H_\theta+a\right)$ for $\theta\in\mathbb{R}^3$ and $F^{\mathrm{u}}\in F_N$.
This is a continuous function of $(\theta,F^{\mathrm{u}})$ on $\mathbb{R}^3\times F_N$, and its modulus does not exceed $\mathrm{Bdn}$.
\end{itemize}
\end{named}

\begin{proof}
\emph{The matrix of $D_B+\boldsymbol{\mu}$.}
Let $\mathrm{Fdm}$ be a fundamental domain of $n\mathbb{Z}^3$ in $\mathbb{Z}^3$.
By Section~I.1.5, an element $f$ of $\mathfrak{S}^{\mathrm{ch}}(1,n)$ is determined by the $192n^3$ numbers $f_\beta(c,g)$ with $c\in\mathrm{Fdm}$, $g\in O_{\mathrm{cub}}$, $\beta\in\{1,\dots,4\}$; these numbers can be prescribed arbitrarily, and
\[
f(c+nv,g)=\boldsymbol{\chi}(v)\cdot f(c,g)\qquad(v\in\mathbb{Z}^3),
\]
where $\boldsymbol{\chi}(v)=1$ if the sum of the three components of $v$ is even and $\boldsymbol{\chi}(v)=-1$ if it is odd (Sections~I.1.5 and I.6.3).
At level $0$ the phases in (I.2.3) are $U_B(x,x^{(i)})=1$ for $i=0,1,2$ and $U_B(x,x^{(3)})=\exp\left\{iqB(c,c_3)\right\}$.
Together with (I.2.5), (I.2.2) and (I.1.2) this gives, at the fermion site $x=(c,g)$,
\begin{equation}
\begin{split}
\left(\left(D_B+\boldsymbol{\mu}\right)f\right)(c,g)&=M_{\mathrm{ch}}\cdot f(c,g)-\sum_{i=0}^{2}M_i(c,g)\cdot f(c,gs_i)\\
&\quad-M_3(c,g)\cdot\exp\left\{iqB(c,c+ge_1)\right\}\cdot f(c+ge_1,gs_3),
\end{split}
\tag{29.2}
\end{equation}
where $M_i(c,g)=\frac{w_i}{2}\left(\rho+\gamma\cdot(g\hat{n}_i)\right)$ are the $4\times4$ matrices of (I.2.1).
In the first four terms $f$ is evaluated at sites of the cluster of $c$, since the hops of the types $0$, $1$, $2$ stay over $c$.
The last term involves a value of $f$ over the nearest neighbour $c_3=c+ge_1$, reached by a hop of type $3$, and it is the only term in which $B$ occurs.
Write $c+ge_1=c'+nv$ with $c'\in\mathrm{Fdm}$ and $v\in\mathbb{Z}^3$.
Then $f(c+ge_1,gs_3)=\boldsymbol{\chi}(v)\cdot f(c',gs_3)$.
Consequently, in the coordinates just described, $D_B+\boldsymbol{\mu}$ is represented by a matrix $\mathrm{Dm}(B)$ whose row $(c,g,\beta)$ has the following entries, where $\beta'$ ranges over $\{1,\dots,4\}$:
\begin{equation}
\begin{aligned}
&M_{\mathrm{ch}}&&\text{in the column }(c,g,\beta);\\
&-M_i(c,g)_{\beta\beta'}&&\text{in the columns }(c,gs_i,\beta'),\ i=0,1,2;\\
&-\boldsymbol{\chi}(v)\cdot M_3(c,g)_{\beta\beta'}\cdot\exp\left\{iqB(c,c+ge_1)\right\}&&\text{in the columns }(c',gs_3,\beta');\\
&0&&\text{in all other columns}.
\end{aligned}
\tag{29.3}
\end{equation}
The five matrices $1$, $s_0$, $s_1$, $s_2$, $s_3$ are pairwise distinct (Section~I.1.2).
Hence $g$, $gs_0$, $gs_1$, $gs_2$, $gs_3$ are five distinct labels, and (29.3) concerns $1+4\cdot4=17$ distinct columns.

\emph{(i)}
The first assertion is part of (Q4).
By Section~I.1.1, $B(c,c+e_\mu)=B_\mu(c)$ and $B(c,c-e_\mu)=-B_\mu(c-e_\mu)$.
Every entry in (29.3) is therefore a constant, or a constant multiple of $\exp\left\{iqB_\mu(c)\right\}$ or of $\exp\left\{-iqB_\mu(c-e_\mu)\right\}$; thus it is an entire function of the numbers $B_\mu(\cdot)$ (I.O1).
Since $\det\left(D_B+\boldsymbol{\mu}\right)=\det\mathrm{Dm}(B)$ is a polynomial in the entries, the second assertion follows.

\emph{(ii)}
The function $\boldsymbol{\lambda}$ is real and $n\mathbb{Z}^3$-periodic in $c$.
The multiplication operators
\[
\left(\mathrm{e}^{\pm iq\boldsymbol{\lambda}}f\right)(c,g):=\mathrm{e}^{\pm iq\boldsymbol{\lambda}(c)}\cdot f(c,g)
\]
act by one and the same factor on the $48$ sites over $c$ (I.POS); they map $\mathfrak{S}^{\mathrm{ch}}(1,n)$ into itself and are inverse to each other.
For every bond $(c,c'')$ one has $(d\boldsymbol{\lambda})(c,c'')=\boldsymbol{\lambda}(c'')-\boldsymbol{\lambda}(c)$, by Section~I.1.1 at $\eta_0=1$, for both orientations.
Apply (29.2), with $B+d\boldsymbol{\lambda}$ in place of $B$, to $\mathrm{e}^{-iq\boldsymbol{\lambda}}f$.
In the first four terms all sites lie over $c$, and the factor is $\mathrm{e}^{-iq\boldsymbol{\lambda}(c)}$.
In the last term one has, with $c_3=c+ge_1$,
\begin{multline*}
\exp\left\{iqB(c,c_3)+iq\left(\boldsymbol{\lambda}(c_3)-\boldsymbol{\lambda}(c)\right)\right\}\cdot\mathrm{e}^{-iq\boldsymbol{\lambda}(c_3)}\cdot f(c_3,gs_3)\\
=\mathrm{e}^{-iq\boldsymbol{\lambda}(c)}\cdot\exp\left\{iqB(c,c_3)\right\}\cdot f(c_3,gs_3).
\end{multline*}
It follows that
\[
\left(D_{B+d\boldsymbol{\lambda}}+\boldsymbol{\mu}\right)\mathrm{e}^{-iq\boldsymbol{\lambda}}=\mathrm{e}^{-iq\boldsymbol{\lambda}}\left(D_B+\boldsymbol{\mu}\right),
\]
which is (I.O4) with $-\boldsymbol{\lambda}$ in place of $\boldsymbol{\lambda}$, the scalar $\boldsymbol{\mu}$ being added on both sides.
So the maps $D_{B+d\boldsymbol{\lambda}}+\boldsymbol{\mu}$ and $D_B+\boldsymbol{\mu}$ are similar, and their determinants coincide.

\emph{(iii)}
Let $\mathrm{Dm}$ be a square matrix whose rows and columns are indexed by the same finite set of indices $t$.
Then $\det\mathrm{Dm}=\sum_{\mathrm{prm}}(\pm1)\prod_t\mathrm{Dm}_{t,\mathrm{prm}(t)}$, where $\mathrm{prm}$ ranges over the permutations of the index set and each term carries the sign of $\mathrm{prm}$.
Consequently
\begin{equation}
\left|\det\mathrm{Dm}\right|\le\sum_{\mathrm{prm}}\prod_t\left|\mathrm{Dm}_{t,\mathrm{prm}(t)}\right|\le\sum_{t'(\cdot)}\prod_t\left|\mathrm{Dm}_{t,t'(t)}\right|=\prod_t\left(\sum_{t'}\left|\mathrm{Dm}_{t,t'}\right|\right),
\tag{29.3a}
\end{equation}
where the middle sum extends over all maps $t\mapsto t'(t)$ of the index set into itself.
By (29.3), the sum of the moduli of the entries in the row $(c,g,\beta)$ of $\mathrm{Dm}(B)$ is
\begin{multline*}
\left|M_{\mathrm{ch}}\right|+\sum_{i=0}^{2}\sum_{\beta'}\left|M_i(c,g)_{\beta\beta'}\right|+\left|\exp\left\{iqB(c,c_3)\right\}\right|\cdot\sum_{\beta'}\left|M_3(c,g)_{\beta\beta'}\right|\\
\le\left|M_{\mathrm{ch}}\right|+\left|M_0(x)\right|_c+\left|M_1(x)\right|_c+\left|M_2(x)\right|_c+\left|M_3(x)\right|_c\cdot\exp\left(q\cdot\left|\operatorname{Im}B\right|_\infty\right),
\end{multline*}
because $|m|_c$ is the largest row sum of the moduli of the entries of a $4\times4$ block $m$ (Section~I.1.5) and $\left|\exp\left\{iqB(b)\right\}\right|=\exp\left(-q\cdot\operatorname{Im}B(b)\right)$.
By (I.m7), $\left|M_0(x)\right|_c=\left|M_3(x)\right|_c=3(1+\rho)$ and $\left|M_1(x)\right|_c,\ \left|M_2(x)\right|_c\le3\sqrt{2}\left(\rho+\sqrt{2}\right)$.
Hence every row sum is at most $\mathrm{rsm}(B)$.
There are $192n^3$ rows, and (29.3a) gives
\[
\left|\det\left(D_B+\boldsymbol{\mu}\right)\right|=\left|\det\mathrm{Dm}(B)\right|\le\mathrm{rsm}(B)^{192n^3}.
\]
For the bound on $\mathrm{rsm}(B)$, note that $\left|M_{\mathrm{ch}}\right|\le|\boldsymbol{\mu}|+\frac{\rho}{2}\,\mathrm{sw}=|\boldsymbol{\mu}|+\rho\left(6+6\sqrt{2}\right)$, and that for $0<\rho<1$
\[
\rho\left(6+6\sqrt{2}\right)+3(1+\rho)+6\sqrt{2}\left(\rho+\sqrt{2}\right)<\left(6+6\sqrt{2}\right)+6+\left(6\sqrt{2}+12\right)=24+12\sqrt{2}<41,
\]
the last inequality because $288<289$; moreover $3(1+\rho)<6$.
If $B=B_0+a$ with $B_0$ real, then $\operatorname{Im}B=\operatorname{Im}a$, and the bound by $\mathrm{Bdn}$ follows.

\emph{(iv)}
By Section~I.6.6, $a^{[0]}(b)=L^0a(L^0b)=a(b)$; hence $a$ lies in $\mathfrak{X}_0$ by the data.
The identity is then (Q4), because $\bar{m}\eta+m_{\mathrm{in}}=\boldsymbol{\mu}$.

\emph{(v)}
The identity follows from (I.7.2) and (29.1).
By (Q7.1) at $k=0$, $d$ maps ${V_{\mathrm{t}}}$ bijectively onto $F_N=dC^1(n)$; the map $F^{\mathrm{u}}\mapsto A(F^{\mathrm{u}})$ is the inverse map, which is linear, and $\theta\mapsto H_\theta$ is linear as well.
Therefore $(\theta,F^{\mathrm{u}})\mapsto A(F^{\mathrm{u}})+H_\theta+a$ is an affine map of $\mathbb{R}^3\times F_N$ into the finite-dimensional space of the complex one-cochains, hence continuous; and $\mathrm{wdn}$ is continuous by (i).
Since $A(F^{\mathrm{u}})+H_\theta$ is real, the bound is a consequence of (iii).
\end{proof}

\subsubsection*{P4. The tower of the block fields}

Let $0\le k\le k_*$, let $A\in C^1(T^0_k)$, and let $\zeta_j\in V_j$ for $j<k$, the space $V_j$ being formed on $T^0_j$.
Define recursively
\begin{equation}
\mathrm{Af}^{(k)}_k:=A,\qquad\mathrm{Af}^{(k)}_j:=A^{\min}_j\left(\mathrm{Af}^{(k)}_{j+1}\right)+\mathrm{Sq}_j\zeta_j\quad(j=k-1,\dots,0),
\tag{29.4}
\end{equation}
and denote the last field of this recursion by $\mathrm{Af}^{(k)}_0(A;\zeta_0,\dots,\zeta_{k-1})\in C^1(n)$.
Then $\mathrm{Af}^{(0)}_0(A)=A$, and $\mathrm{Af}^{(k_*)}_0(A;\zeta)$ is the field $\mathrm{Af}_0(A,\zeta)$ of Section~G4.0.6.
All these fields are real.

Put
\[
\mathrm{Bxp}_k:=\prod_{j<k}\mathrm{Box}_j,
\]
and let $\mathrm{mup}_k$ be the product of the measures $\mathbf{1}_{\mathrm{Box}_j}\mu_I$, $j<k$.
It is a positive Borel measure on $\prod_{j<k}V_j$, concentrated on $\mathrm{Bxp}_k$, of total mass $\prod_{j<k}\mu_I(\mathrm{Box}_j)\le1$.
For $k=0$ the product $\prod_{j<0}V_j$ is a space with one point, and $\mathrm{mup}_0$ gives this point the mass $1$.
With the convention of Section~I.0.4 on products with indicators, one has, for every bounded Borel function $G$ on $\mathrm{Bxp}_k$,
\[
\int\prod_{j<k}\mathrm{d}\mu_I(\zeta_j)\,\mathbf{1}_{\mathrm{Box}_j}(\zeta_j)\;G(\zeta)=\int_{\mathrm{Bxp}_k}\mathrm{d}\,\mathrm{mup}_k(\zeta)\,G(\zeta).
\]

\begin{named}{Lemma 2}
\begin{itemize}
\item[(t1)] Let $k<k_*$ and $A\in C^1(T^0_{k+1})$. Then
\[
\mathrm{Af}^{(k+1)}_0(A;\zeta_0,\dots,\zeta_k)=\mathrm{Af}^{(k)}_0\left(A^{\min}_k(A)+\mathrm{Sq}_k\zeta_k;\zeta_0,\dots,\zeta_{k-1}\right).
\]
\item[(t2)] The map $(A;\zeta_0,\dots,\zeta_{k-1})\mapsto\mathrm{Af}^{(k)}_0(A;\zeta)$ of $C^1(T^0_k)\times\prod_{j<k}V_j$ into $C^1(n)$ is linear, hence continuous.
\item[(t3)] The sets $\mathrm{Box}_j$, $j<k_*$, and $\mathrm{Bxp}_k$, $k\le k_*$, are compact.
\item[(t4)] If $A\in\mathrm{Dom}_k$ and $(\zeta_j)_{j<k}\in\mathrm{Bxp}_k$, then $\mathrm{Af}^{(k)}_j\in\mathrm{Dom}_j$ for $0\le j\le k$.
\item[(t5)] For every real function $\boldsymbol{\lambda}$ on $T^0_{k_*}$ there is a real function $\boldsymbol{\lambda}'$ on $\mathbb{Z}^3/n\mathbb{Z}^3$ such that
\[
\mathrm{Af}_0(A+d\boldsymbol{\lambda},\zeta)=\mathrm{Af}_0(A,\zeta)+d\boldsymbol{\lambda}'\qquad\text{for all }(A,\zeta)\in C^1(T^0_{k_*})\times\prod_jV_j.
\]
\item[(t6)] For every $\theta\in\mathbb{R}^3$ there is a real function $\boldsymbol{\lambda}'$ on $\mathbb{Z}^3/n\mathbb{Z}^3$ such that
\[
\mathrm{Af}_0(A+c_\theta,\zeta)=\mathrm{Af}_0(A,\zeta)+H_\theta+d\boldsymbol{\lambda}'\qquad\text{for all }(A,\zeta)\in C^1(T^0_{k_*})\times\prod_jV_j.
\]
\item[(t7)] For every $F^{\mathrm{u}}\in F_N$ the entries of $\mathrm{Jm}\,F^{\mathrm{u}}$ are the variables $A_*,\zeta_0,\dots,\zeta_{k_*-1}$ of Section~G4.0.6 (\cite[Definition~5.11]{QEDI}) at $F^{\mathrm{u}}$.
The image of $\gamma_N$ under $\mathrm{Jm}$ is the product measure $\mathrm{d}\mathsf{m}_*\times\prod_{j<k_*}\mu_I$ on ${V_{\mathrm{t}}}(T^0_{k_*})\times\prod_{j<k_*}V_j$.
For $F^{\mathrm{u}}\in F_N$ with $\mathrm{Jm}\,F^{\mathrm{u}}=(A_*;\zeta)$,
\[
\mathbf{1}_{\mathfrak{G}}(F^{\mathrm{u}})=\chi_*(A_*)\cdot\prod_{j<k_*}\mathbf{1}_{\mathrm{Box}_j}(\zeta_j).
\]
\item[(t8)] Let $(A,\zeta)\in{V_{\mathrm{t}}}(T^0_{k_*})\times\prod_{j<k_*}V_j$ and $F^{\mathrm{u}}:=d\,\mathrm{Af}_0(A,\zeta)$.
There is a real function $\boldsymbol{\lambda}''$ on $\mathbb{Z}^3/n\mathbb{Z}^3$ such that $\mathrm{Af}_0(A,\zeta)=A(F^{\mathrm{u}})+d\boldsymbol{\lambda}''$.
\end{itemize}
\end{named}

\begin{proof}
\emph{(t1)}
By (29.4) with $k+1$ in place of $k$, one has $\mathrm{Af}^{(k+1)}_{k+1}=A$ and $\mathrm{Af}^{(k+1)}_k=A^{\min}_k(A)+\mathrm{Sq}_k\zeta_k$.
The fields $\mathrm{Af}^{(k+1)}_j$, $j<k$, are obtained from $\mathrm{Af}^{(k+1)}_k$ by the same recursion that produces the fields $\mathrm{Af}^{(k)}_j$ from $\mathrm{Af}^{(k)}_k$.

\emph{(t2)}
For $k=0$ the map is the identity.
For $k\ge1$ the linearity is the first sentence of (Q7.3).
A linear map between finite-dimensional real vector spaces is continuous.

\emph{(t3)}
By (Q7.0), $\{\Upsilon\}$ is an orthonormal basis of the finite-dimensional space $V_j$.
In the coordinates $\zeta\mapsto\left(\zeta(\Upsilon)\right)_\Upsilon$ the set $\mathrm{Box}_j=\left\{\zeta:|\zeta(\Upsilon)|\le p_{0,j}\text{ for all }\Upsilon\right\}$ is the closed cube of side $2p_{0,j}$.
A finite product of compact sets is compact.

\emph{(t4)}
The proof is by downward induction on $j$.
The field $\mathrm{Af}^{(k)}_k=A$ lies in $\mathrm{Dom}_k$.
Suppose that $\mathrm{Af}^{(k)}_{j+1}\in\mathrm{Dom}_{j+1}$ for some $j<k$.
Then (Q6)\,(ii) at the level $j<k_*$, applied with $A'=\mathrm{Af}^{(k)}_{j+1}$ and $\zeta=\zeta_j\in\mathrm{Box}_j$, shows that $\mathrm{Af}^{(k)}_j\in\mathrm{Dom}_j$.

\emph{(t5), (t6)}
Let $A^c\in C^1(T^0_{k_*})$ be closed.
By (t2), the difference
\[
\mathrm{Af}_0(A+A^c,\zeta)-\mathrm{Af}_0(A,\zeta)=\mathrm{Af}_0(A^c,0)
\]
does not depend on $(A,\zeta)$, and by (Q7.5) it equals $L^{-\frac{k_*}{2}}\cdot\mathrm{mean}(A^c)+d\boldsymbol{\lambda}'$ with a real function $\boldsymbol{\lambda}'$ on $\mathbb{Z}^3/n\mathbb{Z}^3$.
For $A^c=d\boldsymbol{\lambda}$: the field $d\boldsymbol{\lambda}$ is closed and $\mathrm{mean}(d\boldsymbol{\lambda})=0$ (Q7.1); this gives (t5).
For $A^c=c_\theta$: a constant bond function is closed, and $\mathrm{mean}(c_\theta)=c_\theta$.
By Section~I.6.1, $e_k\cdot\ell_k=e\cdot a_k^{\frac{1}{2}}\cdot R/a_k=e\cdot R\cdot a_k^{-\frac{1}{2}}$ and $a_k=L^{k-N}=L^k\eta$; hence
\[
e_*\cdot\ell_*=L^{-\frac{k_*}{2}}\cdot e\cdot R\cdot\eta^{-\frac{1}{2}}=L^{-\frac{k_*}{2}}\cdot q\cdot n,
\]
because $q\cdot n=e\cdot\eta^{\frac{1}{2}}\cdot R/\eta$.
Therefore $L^{-\frac{k_*}{2}}\cdot\theta_\mu/(e_*\ell_*)=\theta_\mu/(q\cdot n)$, that is, $L^{-\frac{k_*}{2}}\cdot c_\theta=H_\theta$, and (t6) follows.

\emph{(t7)}
\emph{The variables.}
Let $F^{\mathrm{u}}\in F_N$, put $A_0:=A_0(F^{\mathrm{u}})$ and $A_j:=\left(Q_jA_0\right)^{[j]}$ for $0\le j\le k_*$, and let $(A;\zeta):=\mathrm{Jm}\,F^{\mathrm{u}}$, that is, $A_0=\mathrm{Af}_0(A,\zeta)$ ((Q7.4)\,(b)).
By the second sentence of (Q7.3) at $k=k_*$, one has $A=A_{k_*}$, which is the variable $A_*$ of \cite[Definition~5.11]{QEDI}, and $A_j=A^{\min}_j(A_{j+1})+\mathrm{Sq}_j\zeta_j$ for $j<k_*$.
In \cite[Definition~5.11]{QEDI} the variable $\zeta_j(F^{\mathrm{u}})$ is defined as the element of $V_j$ with this property; there is only one such element, because $\mathrm{Sq}_j$ is injective on $V_j$ (Q7.3).
Hence the entries of $\mathrm{Jm}\,F^{\mathrm{u}}$ are the variables $A_*$, $\zeta_j$ of Section~G4.0.6, with which $\mathfrak{G}$, $\mathrm{d}\mathsf{m}_*$ and the left side of the identity in Proposition~29\,(c) are formed.

\emph{The image measure.}
The measure $\mu_I$ is the standard Gaussian measure of the Euclidean space $V_j$.
Since $\{\Upsilon\}$ is an orthonormal basis (Q7.0), $\mu_I$ is the law of a random element $\zeta$ of $V_j$ whose coordinates $\zeta(\Upsilon)$ are independent standard normal variables.
So (Q7.4)\,(e) states that under $\gamma_N$ the variables $A_*,\zeta_0,\dots,\zeta_{k_*-1}$ are independent and that $\zeta_j$ has the law $\mu_I$.
The variable $A_*$ has the law $\mathrm{d}\mathsf{m}_*$: in view of the first assertion, this is the definition of $\mathrm{d}\mathsf{m}_*$ (Section~G4.0.6).
The joint law of independent variables is the product of their laws.

\emph{The indicator.}
Given the first assertion, the third one follows by comparison of the definitions
\begin{gather*}
\mathfrak{G}=\left\{F^{\mathrm{u}}\in F_N:\ |\zeta_j(\Upsilon)|\le p_{0,j}\text{ for all }\Upsilon\text{ and all }j<k_*;\ |dA_*|_\infty\le p_*\right\},\\
\mathrm{Box}_j=\left\{\zeta\in V_j:\ |\zeta(\Upsilon)|\le p_{0,j}\text{ for all }\Upsilon\right\},\qquad\chi_*(A)=\mathbf{1}\left[|dA|_\infty\le p_*\right]
\end{gather*}
(\cite[Section~5.2.5 and Definition~5.11]{QEDI}; Section~G4.0.6).

\emph{(t8)}
Put $A_0:=\mathrm{Af}_0(A,\zeta)$.
By (Q7.4)\,(b), $\mathrm{Jm}\,F^{\mathrm{u}}=(A,\zeta)$, and $A_0=A_0(F^{\mathrm{u}})$ lies in $\mathrm{Sl}$.
By the definition of $\mathrm{Sl}$ (\cite[Definition~5.11]{QEDI}) the bond function $Q_{k_*}A_0$ has mean zero, and by (Q7.2), $\mathrm{mean}(A_0)=\mathrm{mean}(Q_{k_*}A_0)=0$.
The field $A(F^{\mathrm{u}})$ lies in ${V_{\mathrm{t}}}$, so it has mean zero (Q7.1), and $dA(F^{\mathrm{u}})=F^{\mathrm{u}}=dA_0$.
Thus $A_0-A(F^{\mathrm{u}})$ is closed and has mean zero, and it is therefore exact (Q7.1).
\end{proof}

\subsubsection*{P5. One step along the trajectory}

\begin{named}{Lemma 3}
Let $0\le k<k_*$ and $A'\in\mathrm{Dom}_{k+1}$, and let $y^{x}_{k}$ and $y^{x}_{k+1}$ be formed with the families $\mathfrak{P}_k$ and $\mathfrak{P}_{k+1}$ of the trajectory with slot.
Then
\begin{equation}
\begin{split}
&y^{x}_{k+1}[\varepsilon^{\mathrm{v},t}_{k+1},m_{k+1},\mathscr{E}_{k+1}]\left(A';a^{[k+1]}\right)\\
&\qquad=\int_{\mathrm{Box}_k}\mathrm{d}\mu_I(\zeta)\,y^{x}_{k}[\varepsilon^{\mathrm{v},t}_k,m_k,\mathscr{E}_k]\left(A^{\min}_k(A')+\mathrm{Sq}_k\zeta;a^{[k]}\right).
\end{split}
\tag{29.5}
\end{equation}
\end{named}

\begin{proof}
\emph{(1) The setting.}
By (Q1), the data of level $k$, together with the primitives of (FL), are (ST1)$^{\mathrm{ch}}$.
By Section~G4.0.5, the step map in the recursion of Section~I.5.7 is the step map of level $k$ on $T_k(R)$ formed with these primitives, that is, the step map of Definition~24.D for these data (Section~I.5.7, paragraph ``Step map''; Definition~24.D of Section~G3.4.1 is the counterpart of \cite[Definition~5.25]{QEDI} and of the displays of \cite[Section~17.4]{QEDI}).
The activities $E_k^{\mathrm{det}}$ and ${\delta E^{\mathrm{det}}}_k$ are formed with the family (p5) of the same system.

\emph{(2) The hypotheses of (Q3)} at $(\varepsilon^{\mathrm{v}},m,\mathscr{E})=(\varepsilon^{\mathrm{v},t}_k,m_k,\mathscr{E}_k)$, with the family $\mathfrak{P}_k$ of the trajectory, are provided by (Q2): the numbers $\varepsilon^{\mathrm{v},t}_k$ and $m_k$ are real, $|m_k|\le w^{\mathrm{m}}_k$, $\mathscr{E}_k\in K^{x}_{k}$ with $\|\mathscr{E}_k\|^{x}_{k}\le w^{\mathrm{E}}_k$, $\mathscr{E}_0=0$, and the functions $\mathfrak{P}_k(X,\cdot\,)$ are complex functions on $\mathfrak{X}_k$.
By the data, $\mathscr{C}':=a^{[k+1]}$ lies in $\mathfrak{X}_{k+1}$.

\emph{(3) The image.}
Consider first the levels $1\le k<k_*$.
By (Q5), whose data hold by (q3), and by $\mathscr{E}_k|_0=E_k$ (Q2),
\[
\varepsilon(\mathscr{E}_k)=\varepsilon(E_k),\qquad m(\mathscr{E}_k)=m(E_k).
\]
(The same can be read off the definitions: $\mathrm{al}_{0}(\mathscr{E},X)$ and $\mathrm{al}_{2}(\mathscr{E},X)$ of Section~I.5.6 are formed with the kernels $\mathscr{E}_{00}(X;0,0)$ and $\mathscr{E}_{20}(X;0,0;x,x')$; the kernels of $\mathscr{E}|_0=\mathscr{E}(\,\cdot\,;\,\cdot\,,0,\,\cdot\,)$ at $\mathscr{A}=0$ are these kernels (Section~I.5.4); and $\varepsilon$ and $m\,1$ are the negatives of the sums of $\mathrm{al}_{0}$ and $\mathrm{al}_{2}$ over the small polymers $X$ that contain a given $M$-cube.)
Hence the equations of Section~I.5.7 (paragraph ``Flows'')
\begin{gather*}
\varepsilon^{\mathrm{v},t}_{k+1}=L^3\left(\varepsilon^{\mathrm{v},t}_k+\varepsilon(E_k)+\delta\varepsilon^{0}_{k}\right),\qquad m_{k+1}=L\left(m_k+m(E_k)\right),\\
\mathscr{E}_{k+1}=LR\,\mathscr{E}_k+L\,{\delta E^{\mathrm{det}}}_k+E^{*,x}_k(m_k,\mathscr{E}_k),\\
\mathfrak{P}_{k+1}(Y,\mathscr{C})=\sum_{X:\bar{X}^{L}=LY}\left(\mathfrak{P}_k(X,\mathscr{C}_L)+E_k^{\mathrm{det}}(X,\mathscr{C}_L)\right)
\end{gather*}
state that $(\varepsilon^{\mathrm{v},t}_{k+1},m_{k+1},\mathscr{E}_{k+1})=(\varepsilon^{\mathrm{v},+},m^+,\mathscr{E}^+)$ and that $\mathfrak{P}_{k+1}$ is the family of the display (8.6)$^{\mathrm{ch}}$ of Definition~24.D.

\emph{The level $k=0$.}
Here $\mathscr{E}_0=E_0=0$ and $\mathfrak{P}_0=0$, and $\varepsilon(0)=0$, $m(0)=0$, $LR\,0=0$ (I.L0.21).
The same four equations become
\begin{gather*}
\varepsilon^{\mathrm{v},t}_1=L^3\left(\varepsilon^{\mathrm{v}}_{\mathrm{in}}+\delta\varepsilon^{0}_{0}\right),\qquad m_1=L\,m_{\mathrm{in}},\qquad\mathscr{E}_1=L\,{\delta E^{\mathrm{det}}}_0+E^{*,x}_0(m_{\mathrm{in}},0),\\
\mathfrak{P}_1(Y,\mathscr{C})=\sum_{X:\bar{X}^{L}=LY}E_0^{\mathrm{det}}(X,\mathscr{C}_L)
\end{gather*}
(I.L0.27), which is the form of (8.6)$^{\mathrm{ch}}$ at $k=0$ with $(\varepsilon^{\mathrm{v}},m)=(\varepsilon^{\mathrm{v}}_{\mathrm{in}},m_{\mathrm{in}})$ and $\mathfrak{P}_0=0$.

\emph{(4) The slot.}
For a bond $b$ of $T_k$, Sections~I.3.7 and I.6.6 give
\begin{align*}
\left(a^{[k+1]}\right)_L(b)&=L^{-\frac{1}{2}}\cdot a^{[k+1]}(b/L)=L^{-\frac{1}{2}}\cdot L^{\frac{k+1}{2}}\cdot a\left(L^{k+1}b/L\right)\\
&=L^{\frac{k}{2}}\cdot a\left(L^kb\right)=a^{[k]}(b).
\end{align*}

\emph{(5)}
By (1)--(4), the display (2.4x)$^{\mathrm{ch}}$ of Proposition~24.1\,(c), taken at $A'$ and $\mathscr{C}'=a^{[k+1]}$, is (29.5).
\end{proof}

\emph{The first step, written out.}
At $k=0$, (29.5) reads
\begin{multline*}
y^{x}_{1}\left[L^3\left(\varepsilon^{\mathrm{v}}_{\mathrm{in}}+\delta\varepsilon^{0}_{0}\right),\,L\,m_{\mathrm{in}},\,L\,{\delta E^{\mathrm{det}}}_0+E^{*,x}_0(m_{\mathrm{in}},0)\right]\left(A';a^{[1]}\right)\\
=\int_{\mathrm{Box}_0}\mathrm{d}\mu_I(\zeta)\,y^{x}_{0}[\varepsilon^{\mathrm{v}}_{\mathrm{in}},m_{\mathrm{in}},0]\left(A^{\min}_0(A')+\mathrm{Sq}_0\zeta;a\right).
\end{multline*}
The left side is a Berezin integral over the $4n_1$ pairs of generators $\Psi_\beta(y)$, $\bar{\Psi}_\beta(y)$, $y\in T^0_1$, with weight $1$.
On the right side the integrand of the integral over $\mathrm{Box}_0$ is the Berezin integral of (Q4) over the $192n^3$ pairs of generators on $T^{\mathrm{f}}_0$, with weight $\frac{1}{48}$.
By (Q6)\,(ii) its first argument lies in $\mathrm{Dom}_0$, and by Lemma~1\,(iv) its value is $\mathrm{wdn}\left(A^{\min}_0(A')+\mathrm{Sq}_0\zeta+a\right)$, a quotient of two determinants of size $192n^3$.

\subsubsection*{P6. Proof of (a)}

\begin{proof}[Proof of Proposition~29\,(a)]
\emph{Continuity of the integrands.}
Let $0\le k\le k_*$ and $A\in C^1(T^0_k)$.
By (t2), the map $\zeta\mapsto\mathrm{Af}^{(k)}_0(A;\zeta)+a$ of $\prod_{j<k}V_j$ into the space of the complex one-cochains is affine, hence continuous; and $\mathrm{wdn}$ is continuous by Lemma~1\,(i).
So $\zeta\mapsto\mathrm{wdn}\left(\mathrm{Af}^{(k)}_0(A;\zeta)+a\right)$ is continuous on $\prod_{j<k}V_j$.
On the compact set $\mathrm{Bxp}_k$ (see (t3)) this function is bounded (Q8) and Borel, and therefore its integral with respect to the finite measure $\mathrm{mup}_k$ exists.

\emph{Claim (I$_k$), $0\le k\le k_*$.}
For every $A\in\mathrm{Dom}_k$,
\begin{equation}
y^{x}_{k}[\varepsilon^{\mathrm{v},t}_k,m_k,\mathscr{E}_k]\left(A;a^{[k]}\right)=\int_{\mathrm{Bxp}_k}\mathrm{d}\,\mathrm{mup}_k(\zeta)\,\mathrm{wdn}\left(\mathrm{Af}^{(k)}_0(A;\zeta)+a\right).
\tag{29.6}
\end{equation}

\emph{Proof of (I$_0$).}
By (Q2), the data of level $0$ are $(\varepsilon^{\mathrm{v}}_{\mathrm{in}},m_{\mathrm{in}},0)$, and $\mathfrak{P}_0=0$.
Moreover $a^{[0]}=a$, $\mathrm{Af}^{(0)}_0(A)=A$, and $\mathrm{mup}_0$ is the measure of mass $1$ on a space with one point.
Hence (29.6) at $k=0$ is Lemma~1\,(iv).

\emph{(I$_k$) implies (I$_{k+1}$) for $k<k_*$.}
Let $A\in\mathrm{Dom}_{k+1}$.
For every $\zeta_k\in\mathrm{Box}_k$ the field $A^{\min}_k(A)+\mathrm{Sq}_k\zeta_k$ lies in $\mathrm{Dom}_k$, by (Q6)\,(ii).
Hence
\begin{align*}
&y^{x}_{k+1}[\varepsilon^{\mathrm{v},t}_{k+1},m_{k+1},\mathscr{E}_{k+1}]\left(A;a^{[k+1]}\right)\\
&\quad=\int_{\mathrm{Box}_k}\mathrm{d}\mu_I(\zeta_k)\,y^{x}_{k}[\varepsilon^{\mathrm{v},t}_k,m_k,\mathscr{E}_k]\left(A^{\min}_k(A)+\mathrm{Sq}_k\zeta_k;a^{[k]}\right)\\
&\quad=\int_{\mathrm{Box}_k}\mathrm{d}\mu_I(\zeta_k)\int_{\mathrm{Bxp}_k}\mathrm{d}\,\mathrm{mup}_k(\zeta_0,\dots,\zeta_{k-1})\\
&\qquad\qquad\mathrm{wdn}\left(\mathrm{Af}^{(k)}_0\left(A^{\min}_k(A)+\mathrm{Sq}_k\zeta_k;\zeta_0,\dots,\zeta_{k-1}\right)+a\right)\\
&\quad=\int_{\mathrm{Box}_k}\mathrm{d}\mu_I(\zeta_k)\int_{\mathrm{Bxp}_k}\mathrm{d}\,\mathrm{mup}_k(\zeta_0,\dots,\zeta_{k-1})\,\mathrm{wdn}\left(\mathrm{Af}^{(k+1)}_0(A;\zeta_0,\dots,\zeta_k)+a\right)\\
&\quad=\int_{\mathrm{Bxp}_{k+1}}\mathrm{d}\,\mathrm{mup}_{k+1}(\zeta)\,\mathrm{wdn}\left(\mathrm{Af}^{(k+1)}_0(A;\zeta)+a\right).
\end{align*}
The first equality is Lemma~3; the second is (I$_k$), applied to the field $A^{\min}_k(A)+\mathrm{Sq}_k\zeta_k\in\mathrm{Dom}_k$; the third is (t1); the fourth is Fubini's theorem.
For the last step note that $\mathrm{Bxp}_{k+1}=\mathrm{Bxp}_k\times\mathrm{Box}_k$, that $\mathrm{mup}_{k+1}$ is the product of the finite measures $\mathrm{mup}_k$ and $\mathbf{1}_{\mathrm{Box}_k}\mu_I$, and that the integrand is bounded and Borel on $\mathrm{Bxp}_{k+1}$ by the first paragraph, with $k+1$ in place of $k$.

\emph{Conclusion.}
Let $A\in\mathrm{Dom}_{k_*}$.
Since $Y^{x}_{*}(A;\mathscr{C}_*)=y^{x}_{k_*}[\varepsilon^{\mathrm{v},t}_*,m_*,\mathscr{E}_*]\left(A;a^{[k_*]}\right)$, the claim (I$_{k_*}$) reads
\[
Y^{x}_{*}(A;\mathscr{C}_*)=\int_{\mathrm{Bxp}_{k_*}}\mathrm{d}\,\mathrm{mup}_{k_*}(\zeta)\,\mathrm{wdn}\left(\mathrm{Af}_0(A,\zeta)+a\right).
\]
For $\zeta\in\mathrm{Bxp}_{k_*}$ the field $\mathrm{Af}_0(A,\zeta)$ lies in $\mathrm{Dom}_0$ ((t4) at $k=k_*$, $j=0$), and Lemma~1\,(iv) gives
\[
\mathrm{wdn}\left(\mathrm{Af}_0(A,\zeta)+a\right)=y^{x}_{0}[\varepsilon^{\mathrm{v}}_{\mathrm{in}},m_{\mathrm{in}},0]\left(\mathrm{Af}_0(A,\zeta);a\right).
\]
Together with the definition of $\mathrm{mup}_{k_*}$ and Section~I.0.4, this is the identity of (a).
The formula for $y^{x}_{0}$ that stands in (a) is that of (Q4).
By the first paragraph, the integrand is continuous on the compact set $\prod_{j<k_*}\mathrm{Box}_j$.
\end{proof}

\subsubsection*{P7. Proof of (b)}

\begin{proof}[Proof of Proposition~29\,(b)]
Let $A$ and $A+d\boldsymbol{\lambda}$ lie in $\mathrm{Dom}_{k_*}$, where $\boldsymbol{\lambda}$ is a real function on $T^0_{k_*}$.
By (t5) there is a real function $\boldsymbol{\lambda}'$ on $\mathbb{Z}^3/n\mathbb{Z}^3$ with $\mathrm{Af}_0(A+d\boldsymbol{\lambda},\zeta)=\mathrm{Af}_0(A,\zeta)+d\boldsymbol{\lambda}'$ for all $\zeta$.
Lemma~1\,(ii), applied with $B=\mathrm{Af}_0(A,\zeta)+a$, yields
\[
\mathrm{wdn}\left(\mathrm{Af}_0(A+d\boldsymbol{\lambda},\zeta)+a\right)=\mathrm{wdn}\left(\mathrm{Af}_0(A,\zeta)+a+d\boldsymbol{\lambda}'\right)=\mathrm{wdn}\left(\mathrm{Af}_0(A,\zeta)+a\right)
\]
for all $\zeta\in\prod_{j<k_*}V_j$.
It remains to integrate over $\mathrm{Bxp}_{k_*}$ with respect to $\mathrm{mup}_{k_*}$ and to apply (I$_{k_*}$) of P6 at $A+d\boldsymbol{\lambda}$ and at $A$.
\end{proof}

\subsubsection*{P8. Proof of (c)}

\begin{proof}[Proof of Proposition~29\,(c)]
\emph{Step 1 (the density in the coordinates of (Q7.4)).}
Let $\theta\in\mathbb{R}^3$, $F^{\mathrm{u}}\in F_N$ and $\mathrm{Jm}\,F^{\mathrm{u}}=(A_*;\zeta)$.
Then
\begin{equation}
\mathrm{dens}^{\theta,a}_N(F^{\mathrm{u}})=\mathrm{wdn}\left(\mathrm{Af}_0(A_*+c_\theta,\zeta)+a\right).
\tag{29.7}
\end{equation}
Indeed, $F^{\mathrm{u}}=d\,\mathrm{Af}_0(A_*,\zeta)$ by (Q7.4)\,(b).
By (t6) and (t8) there are real functions $\boldsymbol{\lambda}'$, $\boldsymbol{\lambda}''$ on $\mathbb{Z}^3/n\mathbb{Z}^3$ with
\[
\mathrm{Af}_0(A_*+c_\theta,\zeta)=\mathrm{Af}_0(A_*,\zeta)+H_\theta+d\boldsymbol{\lambda}'=A(F^{\mathrm{u}})+H_\theta+d\left(\boldsymbol{\lambda}'+\boldsymbol{\lambda}''\right).
\]
Lemma~1\,(ii), with $B=A(F^{\mathrm{u}})+H_\theta+a$, and Lemma~1\,(v) now give
\[
\mathrm{wdn}\left(\mathrm{Af}_0(A_*+c_\theta,\zeta)+a\right)=\mathrm{wdn}\left(A(F^{\mathrm{u}})+H_\theta+a\right)=\mathrm{dens}^{\theta,a}_N(F^{\mathrm{u}}).
\]

\emph{Step 2 (the function $\mathrm{Ith}$).}
For $\theta\in\mathbb{R}^3$ and $A\in C^1(T^0_{k_*})$ put
\[
\mathrm{Ith}(\theta,A):=\int_{\mathrm{Bxp}_{k_*}}\mathrm{d}\,\mathrm{mup}_{k_*}(\zeta)\,\mathrm{wdn}\left(\mathrm{Af}_0(A+c_\theta,\zeta)+a\right).
\]
The map $\theta\mapsto c_\theta$ is linear; so, by (t2), $(\theta,A,\zeta)\mapsto\mathrm{Af}_0(A+c_\theta,\zeta)$ is linear, and $(\theta,A,\zeta)\mapsto\mathrm{wdn}\left(\mathrm{Af}_0(A+c_\theta,\zeta)+a\right)$ is continuous on $\mathbb{R}^3\times C^1(T^0_{k_*})\times\prod_jV_j$ (Lemma~1\,(i)).
In particular the integral exists (first paragraph of P6, with $A+c_\theta$ in place of $A$).

(2a)
$\left|\mathrm{Ith}(\theta,A)\right|\le\mathrm{Bdn}$.
Indeed, $\mathrm{Af}_0(A+c_\theta,\zeta)$ is a real one-cochain, so that Lemma~1\,(iii) applies, and the total mass of $\mathrm{mup}_{k_*}$ is at most $1$.
If moreover $A\in{V_{\mathrm{t}}}(T^0_{k_*})$, then
\[
\left|\mathrm{Ith}(\theta,A)\right|\le\sup\left\{\left|\mathrm{dens}^{\theta,a}_N(F^{\mathrm{u}})\right|:F^{\mathrm{u}}\in F_N\right\},
\]
because by (29.7) the integrand at $(A,\zeta)$ equals $\mathrm{dens}^{\theta,a}_N\left(d\,\mathrm{Af}_0(A,\zeta)\right)$ (for $F^{\mathrm{u}}=d\,\mathrm{Af}_0(A,\zeta)$ one has $\mathrm{Jm}\,F^{\mathrm{u}}=(A,\zeta)$ by (Q7.4)\,(b)).

(2b)
$\mathrm{Ith}$ is continuous on $\mathbb{R}^3\times C^1(T^0_{k_*})$.
Let $(\theta',A')$ be a point of this space.
The product of a compact neighbourhood of $(\theta',A')$ with $\mathrm{Bxp}_{k_*}$ is compact (t3), and on it the integrand is uniformly continuous (Q8).
Hence
\[
\sup_{\zeta\in\mathrm{Bxp}_{k_*}}\left|\mathrm{wdn}\left(\mathrm{Af}_0(A+c_\theta,\zeta)+a\right)-\mathrm{wdn}\left(\mathrm{Af}_0(A'+c_{\theta'},\zeta)+a\right)\right|\longrightarrow0
\]
as $(\theta,A)\to(\theta',A')$.
Since the total mass of $\mathrm{mup}_{k_*}$ is at most $1$, the difference $\left|\mathrm{Ith}(\theta,A)-\mathrm{Ith}(\theta',A')\right|$ does not exceed this supremum.

(2c)
If $|dA|_\infty\le p_*$, then $A+c_\theta\in\mathrm{Dom}_{k_*}$, and
\[
\mathrm{Ith}(\theta,A)=Y^{x}_{*}\left(A+c_\theta;\mathscr{C}_*\right).
\]
Indeed, $c_\theta$ is constant, so $d(A+c_\theta)=dA$; thus $A+c_\theta$ is a real bond function on $T^0_{k_*}$ with $\left|d(A+c_\theta)\right|_\infty\le p_*=p_{k_*}$, and (Q6)\,(i) at the level $k_*$ shows that it lies in $\mathrm{Dom}_{k_*}$.
The identity is then (I$_{k_*}$) of P6 at $A+c_\theta$.

\emph{Step 3 (integration).}
Let $\Phi$ be a bounded Borel function on ${V_{\mathrm{t}}}(T^0_{k_*})$.
On ${V_{\mathrm{t}}}(T^0_{k_*})\times\prod_{j<k_*}V_j$ define
\[
G(A,\zeta):=\Phi(A)\cdot\chi_*(A)\cdot\prod_{j<k_*}\mathbf{1}_{\mathrm{Box}_j}(\zeta_j)\cdot\mathrm{wdn}\left(\mathrm{Af}_0(A+c_\theta,\zeta)+a\right).
\]
The function $G$ is Borel: $\Phi$ is Borel, $\chi_*$ and $\mathbf{1}_{\mathrm{Box}_j}$ are indicators of closed sets (the maps $A\mapsto dA$ and $\zeta_j\mapsto\zeta_j(\Upsilon)$ are linear), and the last factor is continuous.
Moreover $|G|\le\sup|\Phi|\cdot\mathrm{Bdn}$ (Lemma~1\,(iii)).
Since $A_*$ is a linear function of $F^{\mathrm{u}}$ ((Q7.4)\,(b)), the function $F^{\mathrm{u}}\mapsto\Phi(A_*)$ is bounded and Borel on $F_N$.
By (t7) and (29.7), for every $F^{\mathrm{u}}\in F_N$,
\[
\mathbf{1}_{\mathfrak{G}}(F^{\mathrm{u}})\cdot\Phi(A_*)\cdot\mathrm{dens}^{\theta,a}_N(F^{\mathrm{u}})=G\left(\mathrm{Jm}\,F^{\mathrm{u}}\right).
\]
Consequently
\begin{align*}
\nu^{\theta,a}\left(\mathbf{1}_{\mathfrak{G}}\,\Phi(A_*)\right)
&=\int\mathrm{d}\gamma_N(F^{\mathrm{u}})\,\mathbf{1}_{\mathfrak{G}}(F^{\mathrm{u}})\cdot\Phi(A_*)\cdot\mathrm{dens}^{\theta,a}_N(F^{\mathrm{u}})\\
&=\int\mathrm{d}\gamma_N(F^{\mathrm{u}})\,G\left(\mathrm{Jm}\,F^{\mathrm{u}}\right)\\
&=\int\mathrm{d}\left(\mathrm{d}\mathsf{m}_*\times\prod_{j<k_*}\mu_I\right)\,G\\
&=\int\mathrm{d}\mathsf{m}_*(A)\,\Phi(A)\cdot\chi_*(A)\int_{\mathrm{Bxp}_{k_*}}\mathrm{d}\,\mathrm{mup}_{k_*}(\zeta)\,\mathrm{wdn}\left(\mathrm{Af}_0(A+c_\theta,\zeta)+a\right)\\
&=\int\mathrm{d}\mathsf{m}_*(A)\,\Phi(A)\cdot\chi_*(A)\cdot\mathrm{Ith}(\theta,A)\\
&=\int\mathrm{d}\mathsf{m}_*(A)\,\Phi(A)\cdot\chi_*(A)\cdot Y^{x}_{*}\left(A+c_\theta;\mathscr{C}_*\right).
\end{align*}
The first equality holds because $\nu^{\theta,a}=\mathrm{dens}^{\theta,a}_N\,\gamma_N$ (Section~I.7.3); the second is the preceding identity; the third is the rule for integrals with respect to an image measure (Q8), together with (t7); the fourth is Fubini's theorem, which applies because $G$ is bounded and Borel and the measures are probability measures; the fifth is the definition of $\mathrm{Ith}$; and the sixth follows from (2c) and Section~I.0.4.
This is the identity of (c).

\emph{Step 4 (the density of the right side).}
By (2c) and Section~I.0.4,
\[
\chi_*(A)\cdot Y^{x}_{*}\left(A+c_\theta;\mathscr{C}_*\right)=\chi_*(A)\cdot\mathrm{Ith}(\theta,A)\qquad\text{for all }A\in{V_{\mathrm{t}}}(T^0_{k_*})\text{ and }\theta\in\mathbb{R}^3.
\]
\emph{Boundedness.}
By (2a) the modulus of this function is at most $\mathrm{Bdn}$, a number that contains neither $\theta$ nor $A$.

\emph{Measurability.}
The set
\[
\left\{A\in{V_{\mathrm{t}}}(T^0_{k_*}):\chi_*(A)=1\right\}=\left\{A:\left|(dA)(p)\right|\le p_*\text{ for all plaquettes }p\text{ of }T^0_{k_*}\right\}
\]
is closed, and $\mathrm{Ith}(\theta,\cdot\,)$ is continuous by (2b); the product of the indicator of a closed set with a continuous function is Borel.

\emph{Continuity in $(\theta,A)$ on $\{\chi_*=1\}$.}
On $\mathbb{R}^3\times\left\{A\in{V_{\mathrm{t}}}(T^0_{k_*}):\chi_*(A)=1\right\}$ the function $(\theta,A)\mapsto Y^{x}_{*}\left(A+c_\theta;\mathscr{C}_*\right)$ is the restriction of $\mathrm{Ith}$, which is continuous by (2b).
\end{proof}

\subsubsection*{P9. The two sentences that follow Proposition~29; range of validity}

\emph{No constant enters Proposition~29.}
The assertions (a) and (b) and the first assertion of (c) are identities between objects of the data; the remaining assertions concern continuity, boundedness and measurability.
No inequality involving a constant is asserted.
Every constant of Section~I.8 that occurs at all occurs inside one of the objects named in the statement:
$L$, $M$ and $r$ in the lattices and in the stop level ($\ell_j$, $r_k$, $k_*$, $\ell_*$);
$p_0$ in $\mathrm{Box}_j$, and $p$ in $\chi_*$;
$\varepsilon$ in the domains $\mathrm{Dom}_j$ and $\mathfrak{X}_j$;
$\alpha$, and again $L$ and $M$, in $\mathrm{Dom}_j$ (through the third norm and the cubes $\square^{\dagger}$), and $\alpha$ also in $\chi=\delta_{\alpha,\mathscr{T}}\psi$;
$b$, through $b_{k_*}$, in $D_*(\mathscr{T})$ and $H_{k_*}^{\#}(\mathscr{T})$;
$\rho$ and the weights in $D_B$, $D_*(\mathscr{T})$ and $H_{k_*}^{\#}(\mathscr{T})$;
the constants of the lines (C0)--(C8), through the step maps and the primitives of (FL), in $\mathscr{E}_*$, $\mathfrak{P}_*$, $\varepsilon^{\mathrm{v},t}_*$ and $m_*$;
and $e_1$ in the range of the data (FL).
In Step~4 of P8 the bound on the function $A\mapsto\chi_*(A)\cdot Y^{x}_{*}\left(A+c_\theta;\mathscr{C}_*\right)$ is the number $\mathrm{Bdn}$ of Lemma~1\,(iii), which is formed with $n$, $q$, $\bar{m}\eta$, $m_{\mathrm{in}}$, $\varepsilon^{\mathrm{v}}_{\mathrm{in}}$, $|\operatorname{Im}a|_\infty$ and $\det\left(D_0+\bar{m}\eta\right)$.

\emph{The twist enters only through the harmonic part $c_\theta$ of the first argument of $Y^{x}_{*}$, and $a$ enters only through $\mathscr{C}_*$.}
This sentence concerns the right side of (c).
The function $\Phi$ is a datum.
The function $\chi_*$ is the indicator of $\left\{|dA|_\infty\le p_*\right\}$, and $\mathrm{d}\mathsf{m}_*$ is the law of $A_*$ under $\gamma_N$ (Section~G4.0.6); both are photon objects, formed without $\theta$ and without $a$.
The numbers $\varepsilon^{\mathrm{v},t}_*$, $m_*$ and the families $\mathscr{E}_*$ (on $\tilde{R}_{k_*}\times\mathfrak{X}_{k_*}$) and $\mathfrak{P}_*$ (on $\mathfrak{X}_{k_*}$) are defined by the recursion of Section~I.5.7 from $(m_{\mathrm{in}},\varepsilon^{\mathrm{v}}_{\mathrm{in}})$ and the primitives of (FL) (Section~G4.0.5).
This recursion contains neither $\theta$ nor $a$: the two connections are variables of the families, and no particular connection is substituted for them.
Hence, in the display of (P2.2) at $(A+c_\theta;\mathscr{C}_*)$, the twist $\theta$ occurs only in $\mathscr{A}=H_{k_*}(A+c_\theta)$, and $a$ occurs only in $\mathscr{C}=\mathscr{C}_*=a^{[k_*]}$.
For $A\in{V_{\mathrm{t}}}(T^0_{k_*})$ one has $\mathrm{mean}(A)=0$ (Q7.1).
So $c_\theta=\mathrm{mean}(A+c_\theta)$ is the constant part of $A+c_\theta$ in the decomposition $C^1=dC^0\oplus\{\text{constants}\}\oplus{V_{\mathrm{t}}}$, and $A$ is its part in ${V_{\mathrm{t}}}$.

\emph{Range of validity.}
The arguments of P3--P8 prove (a), (b) and (c) for all data that satisfy (Q1) and (Q2) and for which $a^{[j]}\in\mathfrak{X}_j$ for $j\le k_*$.
Hence (a), (b) and (c) hold for every $\bar{m}\in(0,1]$, every $e>0$ with $e\cdot h^{\frac{1}{2}}\le e_1$, every $R\in\mathrm{Tor}(e)$, every integer $N\ge N_{\min}(R,e)$, every real pair that is controlled on $R$, every external connection of the data (including $a=0$, since $0$ lies in every $\mathfrak{X}_j$), every $\theta\in\mathbb{R}^3$ and every bounded Borel function $\Phi$.
The induction of P6 extends over all levels $0\le k<k_*(R)$, and its first step, $k=0$, is the display at the end of P5.

\subsection*{Conditions on the constants}

The conditions (A1), (A2), (A8), (A10), (L5), (M1), (M3), (S1) below are rows of the table of Proposition~14.2\,(b), and (C0)--(C9) are the lines of Section~I.8.2.

The conditions on the constants enter the proof of Proposition~29 directly in three ways: through the hypotheses and the proofs of \cite[Propositions~15.29 and 15.30]{QEDI}, by way of (Q6), which is used in (t4) of P4, in the last display of P5, in P6 and in (2c) of P8; through \cite[Lemma~17.4(a)]{QEDI}, in (P2.1); and through the standing hypothesis of \cite[Section~14]{QEDI}, by way of (Q7).
The single exception is the range $0<\rho<1$, (K13), which is used in the proof of Lemma~1\,(iii).
In the conditions (K1)--(K9), (K11)--(K13) listed below, $v=v_k=-\log e_k$ ranges over the values at the levels $k\le k_*(R)$, and $e_k\le e_0$ by (q2).

The conditions (K1)--(K9), (K11) and (K12) are those that are named with Proposition~29 as hypotheses of the results on which its proof rests, in the following way: (A1), first clause: (K4); (A2): (K5); (A8), first and fourth clause: (K3), (K12); (A10), third and fourth clause: (K1), (K2); (M3): (K6); (L5): (K7), (K8), (K12); (M1): (K9); (S1): (K7); (C1): (K7), (K11).
For each of these conditions, the part headed \emph{Used for} names the hypothesis of \cite[Proposition~15.30(i),~(ii)]{QEDI} or of \cite[Lemma~17.4(a)]{QEDI}, or the place in the proofs of \cite[Propositions~15.29 and 15.30]{QEDI} (which rest on \cite[Theorem~15.14]{QEDI}) or in the standing hypotheses of \cite[Sections~14 and 15]{QEDI} at which the condition enters; for (K13) it names the step of the proof above.
The part headed \emph{Place} gives the line of Section~I.8.2 and the row of Proposition~14.2.

\begin{description}
\item[(K1)]
\emph{Condition:} $8C_B\cdot v^{p}\le\exp\left(\left(\frac{3}{4}-3\varepsilon\right)v\right)$ for all $v\ge-\log e_0$.
\emph{Used for:} \cite[Proposition~15.30]{QEDI}, hypothesis $8C_B\cdot p_k\le e_k^{-\frac{3}{4}+3\varepsilon}$; see (6d).
\emph{Place:} (C8); Proposition~14.2\,(b), (A10), third inequality.
\item[(K2)]
\emph{Condition:} $8\exp\left(-\left(\frac{3}{4}-\frac{7\varepsilon}{2}\right)v\right)\le1$ for all $v\ge-\log e_0$.
\emph{Used for:} \cite[Proposition~15.30]{QEDI}, hypothesis $8e_k^{\frac{3}{4}-\frac{7\varepsilon}{2}}\le1$; see (6d).
\emph{Place:} (C8); Proposition~14.2\,(b), (A10), fourth inequality.
\item[(K3)]
\emph{Condition:} $\exp(-\varepsilon v)\le\frac{1}{4}$ for all $v\ge-\log e_0$; it is used at $v=-\log e_0$, where it reads $e_0^{\varepsilon}\le\frac{1}{4}$, whence $e_0<1$.
\emph{Used for:} \cite[Proposition~15.30]{QEDI}, hypothesis $e_k<1$; see (6d).
Also the inequality $e_k^{-\frac{3}{4}+\varepsilon}\le e_k^{-1+\varepsilon}$ in the proof of \cite[Lemma~17.4(a)]{QEDI}.
\emph{Place:} (C8); Proposition~14.2\,(b), (A8), first inequality.
\item[(K4)]
\emph{Condition:} $v^{p_0}\le\exp\left(\frac{\varepsilon v}{4}\right)$ for all $v\ge-\log e_0$.
\emph{Used for:} \cite[Proposition~15.30]{QEDI}, hypothesis $p_{0,k}\le e_k^{-\frac{\varepsilon}{4}}$; see (6d).
\cite[Proposition~15.30]{QEDI} assumes that the whole of (A1) and (A2) holds at $e_k$; of the five inequalities of (A1), its proof draws only on the first one, $p_{0,k}\le e_k^{-\frac{\varepsilon}{4}}$, which the statement displays after the words ``In particular''.
The other four inequalities of (A1) hold at $e_k$ as well, by the row (A1) of Proposition~14.2\,(b), since $v_k\ge-\log e_0$ by (q2).
\emph{Place:} (C8); Proposition~14.2\,(b), (A1), first inequality.
\item[(K5)]
\emph{Condition:} $2C_w^2\cdot c_{\gamma}^2\cdot c_{C}\cdot\exp\left(-\frac{\varepsilon v}{4}\right)\le1$ for all $v\ge-\log e_0$.
\emph{Used for:} \cite[Proposition~15.30]{QEDI}, hypothesis (A2); see (6d).
\emph{Place:} (C8); Proposition~14.2\,(b), (A2).
\item[(K6)]
\emph{Condition:} $\gamma_{\mathrm{w}}\cdot M\ge16$, needed in the form $\gamma_{\mathrm{w}}>0$.
The inequalities $C_P\le C_w$ and $\gamma_{\mathrm{w}}\le\gamma_P$ are not conditions: they follow from the definitions, see (6a) and (6b).
\emph{Used for:} \cite[Proposition~15.30]{QEDI}: the numbers $C_w$, $\gamma_{\mathrm{w}}$, $c_{\gamma}$ in (A2).
\emph{Place:} (C5); Proposition~14.2\,(b), (M3).
Definitions, Proposition~14.2\,(a): $\kappa_0$, $K_0$ (C0); $\alpha_0$, $\gamma^{P}$ (C4); $\gamma_{\mathrm{w}}$ (C4r); $C^{P}$, $C_{w}^{(0)}$, $C_{w}^{(1)}$, $C_w$ (C6); $C_P$: \cite[Section~15.7]{QEDI}.
\item[(K7)]
\emph{Condition:} $\varepsilon\le\bar{\varepsilon}=\frac{1}{200}$, needed in the form $\frac{3\varepsilon}{2}\le\frac{1}{8}$; $\varepsilon>0$; $L\ge5$, needed in the form $L^3\ge125>\left(\frac{8}{3}\right)^4$.
\emph{Used for:} the proof of \cite[Proposition~15.30(ii)]{QEDI}, see (6e); $\varepsilon>0$: (6d).
\emph{Place:} (C7): (S1), together with (C1); (C3): (L5).
\item[(K8)]
\emph{Condition:} $L$ odd, $L\ge3$.
\emph{Used for:} \cite[Sections~14.1.1 and 15]{QEDI}.
\emph{Place:} (C3): (L5).
\item[(K9)]
\emph{Condition:} $M$ a power of $L$, $M\ge L^2$; needed, together with (K8), in the form: $M$ is odd, $M\ge1$, and the $M$-cubes are nested in the $LM$-cubes.
\emph{Used for:} the proof of \cite[Proposition~15.30(ii)]{QEDI}, see (6e); (6a).
\emph{Place:} (C5): (M1).
\item[(K11)]
\emph{Condition:} $0\le\alpha<1$.
\emph{Used for:} \cite[Theorem~15.14]{QEDI}, in Step~3 of the proof of \cite[Proposition~15.29]{QEDI} and in the proof of \cite[Proposition~15.30(iii)]{QEDI}; see (6e).
\emph{Place:} (C1).
\item[(K12)]
\emph{Condition:} $\exp\left(-\frac{v}{4}\right)+\exp(-\varepsilon v)\le1$ for all $v\ge-\log e_0$; $L\ge5$.
\emph{Used for:} \cite[Lemma~17.4(a)]{QEDI}, the inclusion $\tilde{R}_k+\mathfrak{X}_k\subset\tilde{G}_k$, through $e_k^{-\frac{3}{4}+\varepsilon}+e_k^{-1+2\varepsilon}\le e_k^{-1+\varepsilon}$ and $\square^{\sim5}\subset\square^{\sim L}$; see (P2.1).
\emph{Place:} (C8); Proposition~14.2\,(b), (A8), fourth inequality; (C3): (L5).
\item[(K13)]
\emph{Condition:} $0<\rho<1$.
\emph{Used for:} the proof of Lemma~1\,(iii): (I.m7), and the numerals $41$ and $6$ in the bound on $\mathrm{rsm}(B)$, hence in $\mathrm{Bdn}$:
\[
\rho\left(6+6\sqrt{2}\right)+3(1+\rho)+6\sqrt{2}\left(\rho+\sqrt{2}\right)<24+12\sqrt{2}<41,\qquad3(1+\rho)<6.
\]
Of $\mathrm{Bdn}$ only its finiteness is needed.
\emph{Place:} (C0); Proposition~14.2\,(a), row $\rho$.
\end{description}

The conditions (K8), (K9), (K11), (K13), and the parts $L\ge5$ and $\varepsilon>0$ of (K7), are ranges fixed in Section~I.0.5.
All the conditions (K1)--(K9), (K11)--(K13) are among the ranges of Section~I.0.5 and the conditions of Section~I.8.2.

Propositions~26.6, 28.1, 24.1, 22.2, 16.3 and, through Section~I.5.7, Propositions~14.1\,(i) and 19 hold under the conditions on the constants that their proofs require.
For Proposition~28.1 these are $c_{\mathrm{fl}}\cdot L^{-\frac{3}{8}+6\bar{\varepsilon}}\le\frac{1}{16}$ (C3), $e_0^{\varepsilon}\le\frac{1}{4}$ (A8), (S1) in the form $0<\varepsilon\le\bar{\varepsilon}$, (L5) in the form $L>1$, and, through \cite[Lemmas~17.4 and 17.5]{QEDI} and (Sc1), the conditions (A8), $L\ge5$, $L^{-\frac{3}{4}}\le\frac{1}{2}$ (from (L1)), $\varepsilon\le\frac{1}{200}$ and the row (e1) at $v=v_K$.
For Proposition~26.6 they are listed in Chapter~26b, Section~3, and in Annex~B, Section~G4.A.6, and they fall into two groups.
The first group consists of $e_1\le e_0$ (C9); the first three inequalities of the row (e1), through (Sc1)--(Sc4) and (Sc6) of Section~G4.0.2; $M$ a power of $L$ (M1); $L$ odd (L5); $L\ge13$, a consequence of (L1), in (Sc4); and $r\ge2$ (C7), needed in the form $r>0$.
The second group enters through Propositions~26.5, 26.3\,(b) and 14.1\,(ii), for the primitives in the clauses (a), (b) and (e) of Proposition~26.6:
$M\ge M^{F}$ (M5), which contains $M>4L$ (M2);
$M\ge L^2$ (M1);
$\gamma_{\mathrm{w}}\cdot M\ge1$, from (M3);
$M\ge M^{P}$ (M6);
(Z2.s1)--(Z2.s4);
$0\le\alpha<1$ (C1);
$0<\alpha_{0}\le\frac{1}{4}$ (C4);
and the inequalities $\kappa_{\mathrm{w}}\ge0$, $K_{0}\ge1$ and $C_{w}^{(1)}\ge4\exp\left(2\gamma_{\mathrm{w}}\right)\cdot M^{27\alpha_{0}}$, which follow from the definitions.
Of all these conditions, only those of the list (K1)--(K9), (K11)--(K13) enter the proof above again; among the ones just named these are
(A8), in (K3) and (K12);
(S1), in (K7);
(L5), in (K7), (K8) and (K12);
(M1), in (K9);
(M3), in (K6);
(C1), in (K7) and (K11);
and the first two inequalities of (Z2.s4), which are the first and the fourth inequality of (A8), in (K3) and (K12).

\section*{Chapter 38. The determinant at complex connections; the zero-flux sector; the shift identity: proofs of Propositions~38.1--38.3}

\subsection*{0. Conventions of this chapter}

\paragraph{0.1. Units.}
Throughout the chapter the level is~$0$, and lattice units (Section~I.2.5) are used:
$\eta_0=1$, the charge that enters the phases is $q>0$, and $r_q=2\pi/q$.
Both $\exp(\cdot)$ and a power $\mathrm{e}^{\pm i(\cdot)}$ denote the exponential function (Section~I.0.4).
The theorems of Fubini and Tonelli are applied in the form of \cite[Theorem~8.8]{Rudin1987}.

\paragraph{0.2. Period cell.}
Let $n=(n_1,n_2,n_3)$ with integers $n_\mu\ge1$, let $|n|=n_1n_2n_3$, and put
\[
\mathrm{Bn}:=\left\{c\in\mathbb{Z}^3:\ 0\le c_\mu\le n_\mu-1,\ \mu=1,2,3\right\}.
\]
Every $c\in\mathbb{Z}^3$ admits exactly one decomposition
\[
c=\mathrm{red}(c)+\sum_{\mu}\mathrm{kq}_\mu(c)\cdot n_\mu e_\mu,
\qquad \mathrm{red}(c)\in\mathrm{Bn},\quad \mathrm{kq}(c)\in\mathbb{Z}^3,
\]
and one sets $\mathrm{sn}(c):=(-1)^{\mathrm{kq}_1(c)+\mathrm{kq}_2(c)+\mathrm{kq}_3(c)}$.
The set $\mathrm{Xn}:=\mathrm{Bn}\times O_{\mathrm{cub}}$ consists of $48|n|$ fermion sites.
For $x=(c,g)\in\mathbb{Z}^3\times O_{\mathrm{cub}}$ let $\mathrm{red}(x):=(\mathrm{red}(c),g)$ and $\mathrm{sn}(x):=\mathrm{sn}(c)$.
Finally, $\mathrm{Ix}:=\mathrm{Xn}\times\{1,2,3,4\}$, so that $\#|\mathrm{Ix}|=4\cdot48\cdot|n|=192|n|$.

\paragraph{0.3. Cochains; Lebesgue measures.}
A $j$-cochain of the torus is a function on the positively oriented $j$-cells of $\mathbb{Z}^3$ that is periodic under $\sum_\mu n_\mu\mathbb{Z}e_\mu$ (Section~G6.A1).
The real, the integer valued and the complex $j$-cochains form the spaces $C^{j}(n)$, $C^{j}(n;\mathbb{Z})$ and $C^{j}(n)\otimes\mathbb{C}$, respectively.
The orbits of cells are represented by the cells whose base point $c$ lies in $\mathrm{Bn}$:
these are $|n|$ sites, $3|n|$ links $(c,\mu)$, $3|n|$ plaquettes $(c;\mu\nu)$ with $\mu<\nu$, and $|n|$ cubes.
The representatives of the links form the set
\[
\mathrm{Ln}(n):=\left\{(c,\mu):\ c\in\mathrm{Bn},\ \mu=1,2,3\right\}.
\]
A cochain is determined by its values on the representatives, and these values may be prescribed at will;
hence $C^{1}(n)\otimes\mathbb{C}=\mathbb{C}^{\mathrm{Ln}(n)}=\mathbb{C}^{3|n|}$.
For two cochains $u$, $v$ of the same degree, $\langle u,v\rangle$ is the sum of $u\cdot v$ over the representatives;
this form is bilinear on complex cochains, and $\|u\|^2:=\langle u,u\rangle$ for real $u$.
The coboundary $d$ is that of Section~I.1.1, written in lattice units:
\begin{gather*}
(d\boldsymbol{\lambda})_\mu(c)=\boldsymbol{\lambda}(c+e_\mu)-\boldsymbol{\lambda}(c),\qquad
(dA)_{\mu\nu}(c)=A_\mu(c)+A_\nu(c+e_\mu)-A_\mu(c+e_\nu)-A_\nu(c),\\
(du)(c)=\left[u_{23}(c+e_1)-u_{23}(c)\right]-\left[u_{13}(c+e_2)-u_{13}(c)\right]+\left[u_{12}(c+e_3)-u_{12}(c)\right].
\end{gather*}
A cochain $u$ is called closed if $du=0$;
$d^{*}$ denotes the adjoint of $d$ with respect to $\langle\cdot,\cdot\rangle$, and $\Delta_0=d^{*}d$ on $C^{0}(n)$.
A one-cochain $B$ is also regarded as the odd function on oriented pairs of nearest neighbours with $B(c,c+e_\mu)=B_\mu(c)=-B(c+e_\mu,c)$ (Section~I.1.1);
for $b=(c,\mu)\in\mathrm{Ln}(n)$ put $B(b):=B_\mu(c)$.
The cochains $\mathbf{1}_c\in C^{0}(n;\mathbb{Z})$ and $\mathbf{1}_b\in C^{1}(n;\mathbb{Z})$ take the value $1$ on the orbit of the site $c$, respectively of the link $b$, and the value $0$ elsewhere;
moreover $\mathbf{1}_\mu:=\sum_{c\in\mathrm{Bn}}\mathbf{1}_{(c,\mu)}$.
For $\theta\in\mathbb{C}^3$ the twist connection is (Section~I.6.6)
\[
H_{\theta}:=\sum_{\mu}\frac{\theta_\mu}{q\,n_\mu}\,\mathbf{1}_\mu .
\]

The families $(\mathbf{1}_c)_{c\in\mathrm{Bn}}$ and $(\mathbf{1}_b)_{b\in\mathrm{Ln}(n)}$ are orthonormal bases of $C^{0}(n)$ and of $C^{1}(n)$.
The Lebesgue measure of a Euclidean space of dimension $\mathrm{Dq}$ is the Lebesgue measure of $\mathbb{R}^{\mathrm{Dq}}$ transported by the coordinates with respect to an orthonormal basis.
It does not depend on the basis, since two orthonormal bases differ by an orthogonal matrix, whose determinant is $\pm1$ (\cite[Theorem~2.20(e) with~2.23]{Rudin1987});
and it is invariant under translations (\cite[Theorem~2.20(c)]{Rudin1987}).
Consequently $\mathrm{d}A=\prod_{b\in\mathrm{Ln}(n)}\mathrm{d}A(b)$ is the Lebesgue measure of $\left(C^{1}(n),\langle\cdot,\cdot\rangle\right)$,
and on an orthogonal sum of subspaces the Lebesgue measure is the product of the Lebesgue measures of the summands
(one takes an orthonormal basis that is the union of orthonormal bases of the summands).
The letter $\sigma$ denotes the Lebesgue measure of $\left({V_{\mathrm{t}}},\langle\cdot,\cdot\rangle\right)$.

\begin{named}{(Lin)}
Let $\mathrm{Tl}\colon\mathrm{Vq}\to\mathrm{Wv}$ be a linear bijection between Euclidean spaces of the same dimension, whose Lebesgue measures are $\mathrm{d}A$ and $\mathrm{d}u$,
and let $|\det\mathrm{Tl}|$ be the modulus of the determinant of the matrix of $\mathrm{Tl}$ with respect to orthonormal bases of $\mathrm{Vq}$ and of $\mathrm{Wv}$.
Then every Borel function $\mathrm{Gv}\ge0$ on $\mathrm{Wv}$ satisfies
\[
\int_{\mathrm{Vq}}\mathrm{d}A\,\mathrm{Gv}(\mathrm{Tl}\,A)=|\det\mathrm{Tl}|^{-1}\int_{\mathrm{Wv}}\mathrm{d}u\,\mathrm{Gv}(u).
\]
\end{named}

\begin{proof}[Proof of (Lin)]
When $\mathrm{Gv}$ is the indicator function of a Borel set $\mathrm{Ew}\subset\mathrm{Wv}$, the assertion reads
$\operatorname{vol}\left(\mathrm{Tl}^{-1}\mathrm{Ew}\right)=\left|\det\mathrm{Tl}^{-1}\right|\cdot\operatorname{vol}(\mathrm{Ew})$,
and this is \cite[Theorem~2.20(e) with~2.23]{Rudin1987}, applied to the matrix of $\mathrm{Tl}^{-1}$.
Both sides are additive in $\mathrm{Gv}$ and homogeneous with respect to factors $\ge0$;
hence the identity holds for all finite linear combinations of indicator functions with coefficients $\ge0$.
Every Borel function $\mathrm{Gv}\ge0$ is the limit of a non-decreasing sequence of such combinations (\cite[Theorem~1.17]{Rudin1987}),
and along this sequence both sides converge by the monotone convergence theorem (\cite[Theorem~1.26]{Rudin1987}).
\end{proof}

\paragraph{0.4. Spinors, operator, determinant.}
According to the data of Proposition~46,
\[
\mathfrak{S}^{\mathrm{ch}}(1;n)=\left\{f\colon\mathbb{Z}^3\times O_{\mathrm{cub}}\to\mathbb{C}^4:\ f(c+n_\mu e_\mu,g)=-f(c,g)\right\}.
\]
Therefore $f(x)=\mathrm{sn}(x)\cdot f(\mathrm{red}(x))$.
Let $\mathrm{ev}_1,\dots,\mathrm{ev}_4$ be the unit vectors of $\mathbb{C}^4$, and for $(x',\beta)\in\mathrm{Ix}$ let
\[
\delta_{x',\beta}(x):=\mathrm{sn}(x)\cdot\mathbf{1}[\mathrm{red}(x)=x']\cdot\mathrm{ev}_\beta .
\]
These $192|n|$ functions form a basis of $\mathfrak{S}^{\mathrm{ch}}(1;n)$, in which
$f=\sum_{(x',\beta)\in\mathrm{Ix}}f_\beta(x')\cdot\delta_{x',\beta}$.

Let $B\in C^{1}(n)\otimes\mathbb{C}$ and $x=(c,g)$.
By (I.1.2), (I.1.6) and (I.2.1)--(I.2.3) at $k=0$,
\begin{gather*}
x^{(i)}=(c,g s_i)\quad(i=0,1,2),\qquad x^{(3)}=(c_3,g s_3),\qquad c_3=c+g e_1,\\
U_B(x,x^{(i)})=1\quad(i=0,1,2),\qquad U_B(x,x^{(3)})=\mathrm{e}^{iqB(c,c_3)},\\
M_i(x)=\frac{w_i}{2}\left(\rho+\gamma\cdot(g\,\hat n_i)\right),\qquad
\left(R_i(B)f\right)(x)=U_B(x,x^{(i)})\cdot f(x^{(i)}),\\
(D_Bf)(x)=-\sum_{i=0}^{3}M_i(x)\left[U_B(x,x^{(i)})\cdot f(x^{(i)})-f(x)\right].
\end{gather*}
By (I.s1) one has $(x+n_\mu e_\mu)^{(i)}=x^{(i)}+n_\mu e_\mu$;
moreover $B$ is periodic, and $M_i(x)$ depends on $g$ alone.
So $R_i(B)$, the multiplication by $M_i(\cdot)$ and $D_B$ map $\mathfrak{S}^{\mathrm{ch}}(1;n)$ into itself.
Let $\det\nolimits_n(D_B+\boldsymbol{\mu})$ denote the determinant of the linear map $D_B+\boldsymbol{\mu}$ of $\mathfrak{S}^{\mathrm{ch}}(1;n)$ into itself.
It depends neither on a basis nor on a pairing (Proposition~46\,(e)); in particular the weight $1/48$ of (I.1.3) plays no role in it.

\paragraph{0.5. Entire functions.}
Let $\mathrm{Nq}\ge1$ be an integer.
For $\mathrm{kx}\in\{0,1,2,\dots\}^{\mathrm{Nq}}$ and $a\in\mathbb{C}^{\mathrm{Nq}}$ put
\[
a^{\mathrm{kx}}:=\prod_j a_j^{\mathrm{kx}_j},\qquad
\mathrm{kx}!:=\prod_j\mathrm{kx}_j!,\qquad
|\mathrm{kx}|:=\sum_j\mathrm{kx}_j,\qquad
\mathrm{d}_a^{\mathrm{kx}}:=\prod_j\left(\frac{\mathrm{d}}{\mathrm{d}a_j}\right)^{\mathrm{kx}_j},
\]
where the derivatives are complex derivatives.
A function $\mathrm{fe}\colon\mathbb{C}^{\mathrm{Nq}}\to\mathbb{C}$ is called entire if there are complex numbers $\mathrm{tay}_{\mathrm{kx}}$ such that
$\mathrm{fe}(a)=\sum_{\mathrm{kx}}\mathrm{tay}_{\mathrm{kx}}\cdot a^{\mathrm{kx}}$ for every $a\in\mathbb{C}^{\mathrm{Nq}}$, the series being absolutely convergent at every $a$.
Taken at the point $a=(\mathrm{rd},\dots,\mathrm{rd})$, this requirement gives
\begin{equation}
\sum_{\mathrm{kx}}|\mathrm{tay}_{\mathrm{kx}}|\cdot\mathrm{rd}^{|\mathrm{kx}|}<\infty\qquad\text{for every }\mathrm{rd}\in(0,\infty).
\tag{38.0a}
\end{equation}
The terms of an absolutely convergent series may be collected into groups, and the groups may be summed in any order (Fubini's theorem for counting measures, \cite[Theorem~8.8]{Rudin1987}).

\begin{named}{(ef1)}[Derivatives]
Let $\mathrm{fe}$ be entire.
For every $j$ the derivative $\frac{\mathrm{d}}{\mathrm{d}a_j}\mathrm{fe}$ exists at every point of $\mathbb{C}^{\mathrm{Nq}}$, it is entire, and its series is
\[
\sum_{\mathrm{kx}:\ \mathrm{kx}_j\ge1}\mathrm{kx}_j\cdot\mathrm{tay}_{\mathrm{kx}}\cdot a_j^{\mathrm{kx}_j-1}\prod_{j'\ne j}a_{j'}^{\mathrm{kx}_{j'}} .
\]
Consequently all the derivatives $\mathrm{d}_a^{\mathrm{kx}}\mathrm{fe}$ exist and are entire, they do not depend on the order in which the differentiations are performed, and
\begin{equation}
\mathrm{d}_a^{\mathrm{kx}}\mathrm{fe}(0)=\mathrm{kx}!\cdot\mathrm{tay}_{\mathrm{kx}} .
\tag{38.0b}
\end{equation}
In particular, $\mathrm{fe}$ determines the coefficients $\mathrm{tay}_{\mathrm{kx}}$.
\end{named}

\begin{proof}[Proof of (ef1)]
Since $\mathrm{kx}_j\le2^{|\mathrm{kx}|}$, the bound (38.0a) with $2\,\mathrm{rd}$ in place of $\mathrm{rd}$ yields
\[
\sum_{\mathrm{kx}}\mathrm{kx}_j\cdot|\mathrm{tay}_{\mathrm{kx}}|\cdot\mathrm{rd}^{|\mathrm{kx}|-1}\le\mathrm{rd}^{-1}\sum_{\mathrm{kx}}|\mathrm{tay}_{\mathrm{kx}}|\cdot(2\,\mathrm{rd})^{|\mathrm{kx}|}<\infty .
\]
Hence the series displayed in (ef1) converges absolutely at every point.
Now fix the variables $a_{j'}$ with $j'\ne j$.
If the terms of the series of $\mathrm{fe}$ are grouped according to the value $i$ of $\mathrm{kx}_j$, there results a power series in the one variable $a_j$ that converges for every $a_j\in\mathbb{C}$;
its coefficient of $a_j^i$ is the sum of $\mathrm{tay}_{\mathrm{kx}}\prod_{j'\ne j}a_{j'}^{\mathrm{kx}_{j'}}$ over the $\mathrm{kx}$ with $\mathrm{kx}_j=i$.
By \cite[Theorem~10.6]{Rudin1987} the sum of this power series is differentiable in $a_j$, and its derivative is the power series differentiated term by term.
After the same grouping, this is the series displayed in (ef1).
By induction on $|\mathrm{kx}|$ it follows that $\mathrm{d}_a^{\mathrm{kx}}\mathrm{fe}$ is entire and that its series is the series of $\mathrm{fe}$ differentiated term by term.
On a monomial the differentiations commute, and at $a=0$ the derivative $\mathrm{d}_a^{\mathrm{kx}}a^{\mathrm{kx}'}$ equals $\mathrm{kx}!$ if $\mathrm{kx}'=\mathrm{kx}$ and $0$ otherwise.
\end{proof}

By (38.0a), the series of an entire function converges uniformly on every polydisc $\left\{\max_j|a_j|\le\mathrm{rd}\right\}$.
An entire function in the present sense is therefore continuous on $\mathbb{C}^{\mathrm{Nq}}$;
by (ef1) it has complex partial derivatives of all orders, so that it is holomorphic on $\mathbb{C}^{\mathrm{Nq}}$.

\begin{named}{(ef2)}[Affine substitution]
Let $\mathrm{fe}$ be entire on $\mathbb{C}^{\mathrm{Nq}}$, let $a^0\in\mathbb{C}^{\mathrm{Nq}}$, and let $(\mathrm{lm}_{ji})$ be a complex matrix with $\mathrm{Nq}$ rows and $\mathrm{Nq}'$ columns;
put $(\mathrm{lm}\,t)_j:=\sum_i\mathrm{lm}_{ji}t_i$.
Then $t\mapsto\mathrm{fe}(a^0+\mathrm{lm}\,t)$ is entire on $\mathbb{C}^{\mathrm{Nq}'}$.
\end{named}

\begin{proof}[Proof of (ef2)]
One has
\[
\mathrm{fe}(a^0+\mathrm{lm}\,t)=\sum_{\mathrm{kx}}\mathrm{tay}_{\mathrm{kx}}\prod_j\left(a^0_j+\sum_i\mathrm{lm}_{ji}t_i\right)^{\mathrm{kx}_j}.
\]
Expand the products.
The term of index $\mathrm{kx}$ becomes a finite sum of complex numbers times monomials in $t$, and the moduli of these summands add up to
\[
|\mathrm{tay}_{\mathrm{kx}}|\prod_j\left(|a^0_j|+\sum_i|\mathrm{lm}_{ji}|\cdot|t_i|\right)^{\mathrm{kx}_j}\le|\mathrm{tay}_{\mathrm{kx}}|\cdot\mathrm{rd}^{|\mathrm{kx}|},
\qquad
\mathrm{rd}:=1+\max_j\left(|a^0_j|+\sum_i|\mathrm{lm}_{ji}|\cdot|t_i|\right).
\]
By (38.0a), the family of all these summands, for all $\mathrm{kx}$, is absolutely summable.
Grouped according to the monomials in $t$, it becomes a power series in $t$ which converges absolutely at the point $t$ and whose sum is $\mathrm{fe}(a^0+\mathrm{lm}\,t)$.
The coefficient of a monomial is a series of numbers that do not contain $t$, and this series converges absolutely, as the point $t=(1,\dots,1)$ shows.
\end{proof}

\begin{named}{(ef3)}[Sums, products, exponentials]
Finite linear combinations and products of entire functions are entire.
If $\mathrm{Pol}$ is a polynomial in $a_1,\dots,a_{\mathrm{Nq}}$ with complex coefficients, then $\exp\mathrm{Pol}(a)$ is entire.
\end{named}

\begin{proof}[Proof of (ef3)]
For linear combinations one argues term by term.
Let $\mathrm{fe}$ and $\mathrm{fe}'$ be entire, with the coefficients $\mathrm{tay}_{\mathrm{kx}}$ and $\mathrm{tay}'_{\mathrm{kx}}$.
The family $\mathrm{tay}_{\mathrm{kx}}\cdot\mathrm{tay}'_{\mathrm{kx}'}\cdot a^{\mathrm{kx}}a^{\mathrm{kx}'}$ is absolutely summable, because the sum of its moduli equals
$\left(\sum_{\mathrm{kx}}|\mathrm{tay}_{\mathrm{kx}}|\cdot|a^{\mathrm{kx}}|\right)\left(\sum_{\mathrm{kx}'}|\mathrm{tay}'_{\mathrm{kx}'}|\cdot|a^{\mathrm{kx}'}|\right)$;
grouped according to $\mathrm{kx}+\mathrm{kx}'$, it is a power series whose sum is $\mathrm{fe}(a)\cdot\mathrm{fe}'(a)$.
For the exponential, write $\exp\mathrm{Pol}(a)=\sum_{i\ge0}\frac{\mathrm{Pol}(a)^i}{i!}$ and expand the powers.
Let $\mathrm{Pol}^+$ be the polynomial whose coefficients are the moduli of the coefficients of $\mathrm{Pol}$.
The moduli of all the summands add up to
\[
\sum_{i\ge0}\frac{\mathrm{Pol}^+\left(|a_1|,\dots,|a_{\mathrm{Nq}}|\right)^i}{i!}=\exp\mathrm{Pol}^+\left(|a_1|,\dots,|a_{\mathrm{Nq}}|\right)<\infty .
\]
Grouped according to the monomials in $a$, the family is a power series whose sum is $\exp\mathrm{Pol}(a)$;
the coefficient of a monomial is a series of numbers that do not contain $a$, and it converges absolutely, as the point $a=(1,\dots,1)$ shows.
\end{proof}

\paragraph{0.6. Locally uniform absolute convergence.}
Let an integral, a series, or a series of integrals depend on a parameter $a\in\mathbb{C}^{\mathrm{Nq}}$.
It is said to converge absolutely, locally uniformly in $a$, if for every $\mathrm{rd}\in(0,\infty)$ the moduli of the integrands (of the terms) have an integrable (a summable) majorant which is the same for all $a$ in $\left\{\max_j|a_j|\le\mathrm{rd}\right\}$.
In Proposition~38.3\,(a) and~(a$'$) the words are meant in this sense, with $a\in\mathbb{C}^{\mathrm{Ln}(n)}$.

\paragraph{0.7. Letters.}
The notation of Part~I and of Section~G6.A3 is in force;
in particular $\mathrm{wq}$, $Z^{\mathrm{ph}}_n$, $\mathrm{Ld}(n)$ and the class $\mathfrak{F}(n)$ with its constant $\mathrm{cf}$ have the meanings given there,
and $\mathrm{zV}_n$, $\mathrm{Ev}$ and $a(t)$ have the meanings given in Propositions~38.1--38.3.
In this chapter $\mathrm{Pol}$ is the arbitrary polynomial in $a_1,\dots,a_{\mathrm{Nq}}$ with complex coefficients of (ef3),
and $\mathrm{Pol}^+$ is the polynomial whose coefficients are the moduli of the coefficients of $\mathrm{Pol}$.
The set $\mathrm{Ln}(n)$ is the one defined in 0.3, formed by the representatives $(c,\mu)$, $c\in\mathrm{Bn}$, of the $3|n|$ orbits of links;
it is the index set that Part~I denotes by the word ``links'', as in $[0,r_q)^{\mathrm{links}}$ (Sections~I.7.1, I.7.2 and~I.7.4).
As in Section~I.1.1, $\mathfrak{l}$ denotes the third direction;
$\Pi_{\mu\nu}$ is the set of the plaquettes of one plane (Section~I.7.1).
The letter $y$ denotes a real variable and $z$ a complex variable:
$z\in\mathbb{C}^r$ in 3.6 and $z\in\mathbb{C}^{N_{*}}$ in 3.7, as in Proposition~38.3\,(d), (e), and $z\in\mathbb{C}$ in the paragraph of 3.7 on the Hermite polynomials;
in this chapter neither letter denotes a unit site or a site of the $L$-lattice.
In 0.4, 1.1 and 1.3 the letter $\beta\in\{1,2,3,4\}$ is a spinor index, and so is $\alpha\in\{1,2,3,4\}$ in 1.1 and 1.3;
elsewhere in this chapter $\alpha$ denotes the H\"older order.
The letters $i$, $j$, $j'$ and $k_\mu$ denote integers.
For $i\ge0$, $H_i(y)$ is the polynomial defined in 3.7; $H_{2}$ and $H_{4}$ are those of Section~I.7.6.
In the second half of 1.3 the symbols $\mathbb{T}$, $\chi$, $\mathbb{X}$, $\mathbb{X}_{\mathbb{T}}$, $\mathbb{F}$, $r^{(i)}$, $\widehat{r^{(i)}}$, $g_i(r)$, $w(r,r^{(i)})$, $y(r)$, $v_1$, $u_b$, $\mathfrak{s}$, $n^{\uparrow}_b(\mathfrak{s})$, $n^{\downarrow}_b(\mathfrak{s})$ and $P_{\mathbb{T}}$
have the meanings of \cite[Definitions~2.4, 2.10, 4.1, 4.3 and~4.4]{QEDII}, a site of $\mathbb{X}$ being written $r$ instead of $x$.
Here \cite[Definition~2.4]{QEDII} fixes the base point $v_0=(v_1,v_2,v_3)\in C_0$, so that $0<v_3<v_2<v_1<\tfrac12$,
and defines the set $\mathbb{X}=Wv_0$ of the fermion sites, the neighbours $r^{(i)}$ of a site $r$, the map $\pi$ that assigns to a site its nearest photon site, and $y(r)=r-\pi(r)$.

\subsection*{1. Proposition~38.1. The determinant at complex connections}

\begin{proof}[Proof of Proposition~38.1]
The notation is that of Section~0.

\paragraph{1.1. The matrix.}
The formula for $D_B$ in 0.4 can be written
\[
(D_Bf)(x)=-\sum_iM_i(x)\cdot\left(R_i(B)f\right)(x)+\left(\sum_iM_i(x)\right)f(x),
\]
and the closure relation $\sum_iw_i\hat n_i=0$ ((I.1.1); this is (I.m1)) gives
\[
\sum_iM_i(x)=\frac{\rho}{2}\,\mathrm{sw}+\frac12\,\gamma\cdot\left(g\sum_iw_i\hat n_i\right)=\frac{\rho}{2}\,\mathrm{sw}.
\]
Since $M_{\mathrm{ch}}=\boldsymbol{\mu}+\frac{\rho}{2}\,\mathrm{sw}$, this is (I.2.5):
\begin{equation}
D_B+\boldsymbol{\mu}=M_{\mathrm{ch}}-\sum_{i=0}^{3}M_i(\cdot)\,R_i(B).
\tag{38.1}
\end{equation}
For $(x,\alpha),(x',\beta)\in\mathrm{Ix}$ let $\mathrm{Dm}[(x,\alpha),(x',\beta)]$ denote the component $\alpha$ of $\left((D_B+\boldsymbol{\mu})\,\delta_{x',\beta}\right)(x)$.
By 0.4,
\[
(D_B+\boldsymbol{\mu})\,\delta_{x',\beta}=\sum_{(x,\alpha)\in\mathrm{Ix}}\mathrm{Dm}[(x,\alpha),(x',\beta)]\cdot\delta_{x,\alpha}.
\]
Thus $\mathrm{Dm}$ is the matrix of $D_B+\boldsymbol{\mu}$ in the basis $(\delta_{x',\beta})$, and $\det\nolimits_n(D_B+\boldsymbol{\mu})=\det\mathrm{Dm}$.
As $\delta_{x',\beta}(x^{(i)})=\mathrm{sn}(x^{(i)})\cdot\mathbf{1}[\mathrm{red}(x^{(i)})=x']\cdot\mathrm{ev}_\beta$, the identity (38.1) gives
\begin{equation}
\begin{split}
\mathrm{Dm}[(x,\alpha),(x',\beta)]={}&M_{\mathrm{ch}}\cdot\mathbf{1}[x=x']\cdot\mathbf{1}[\alpha=\beta]\\
&-\sum_{i=0}^{3}\mathbf{1}[\mathrm{red}(x^{(i)})=x']\cdot\mathrm{sn}(x^{(i)})\cdot U_B(x,x^{(i)})\cdot\left(M_i(x)\right)_{\alpha\beta},\\
M_i(x)={}&\frac{w_i}{2}\left(\rho+\gamma\cdot(g\,\hat n_i)\right).
\end{split}
\tag{38.2}
\end{equation}
By (I.1.2) the labels of $x,x^{(0)},x^{(1)},x^{(2)},x^{(3)}$ are $g,gs_0,gs_1,gs_2,gs_3$,
and $1,s_0,s_1,s_2,s_3$ are five distinct elements of $O_{\mathrm{cub}}$ (Section~I.1.2).
Hence $x,\mathrm{red}(x^{(0)}),\mathrm{red}(x^{(1)}),\mathrm{red}(x^{(2)}),\mathrm{red}(x^{(3)})$ are five distinct sites of $\mathrm{Xn}$, whatever the sides $n_\mu\ge1$;
for given $x$ and $x'$, at most one of the five terms of (38.2) is different from $0$.
For $i=0,1,2$ the site $x^{(i)}=(c,gs_i)$ belongs to the cluster of $c$;
therefore $\mathrm{red}(x^{(i)})=x^{(i)}$ and $\mathrm{sn}(x^{(i)})=1$, and the $4\times4$ block of $\mathrm{Dm}$ at $(x,x^{(i)})$ is $-M_i(x)$, which carries no phase.
The block at $(x,\mathrm{red}(x^{(3)}))$ is $-\mathrm{sn}(x^{(3)})\cdot\mathrm{e}^{iqB(c,c_3)}\cdot M_3(x)$, where $M_3(x)=3\left(\rho-\gamma\cdot(g e_1)\right)$;
it connects the cluster of $c$ with the cluster of $\mathrm{red}(c_3)$, and the two clusters coincide if the side in the direction of $g e_1$ equals $1$.

\paragraph{1.2. The variable of a site.}
Let $x=(c,g)\in\mathrm{Xn}$.
The vector $g e_1$ is one of the six vectors $\pm e_\mu$.
\begin{itemize}
\item
If $g e_1=e_\mu$, then $c_3=c+e_\mu$ and $B(c,c_3)=B_\mu(c)=B(b)$ with $b=(c,\mu)\in\mathrm{Ln}(n)$, so that $U_B(x,x^{(3)})=\mathrm{e}^{iqB(b)}$.
In this case $x$ is said to carry $(b,+)$.
\item
If $g e_1=-e_\mu$, then $c_3=c-e_\mu$ and, by the periodicity of $B$, $B(c,c_3)=-B_\mu(c-e_\mu)=-B(b)$ with $b=(\mathrm{red}(c-e_\mu),\mu)\in\mathrm{Ln}(n)$, so that $U_B(x,x^{(3)})=\mathrm{e}^{-iqB(b)}$.
In this case $x$ is said to carry $(b,-)$.
\end{itemize}
Every site of $\mathrm{Xn}$ carries exactly one of the $6|n|$ pairs $(b,+)$, $(b,-)$ with $b\in\mathrm{Ln}(n)$.
Conversely, let $b=(c',\mu)\in\mathrm{Ln}(n)$.
The sites of $\mathrm{Xn}$ that carry $(b,+)$ are the sites $(c',g)$ with $g e_1=e_\mu$.
The sites that carry $(b,-)$ are the sites $(c,g)$ with $g e_1=-e_\mu$ and $\mathrm{red}(c-e_\mu)=c'$, that is, with $c=\mathrm{red}(c'+e_\mu)$.
In each of the two cases there are $8$ such sites, by (I.s4):
a signed permutation matrix with a prescribed first column is determined by the order in which the two remaining unit vectors occupy its second and third columns and by their two signs, which leaves $2\cdot2\cdot2=8$ possibilities.

\paragraph{1.3. Proof of (a).}
Let $u_b^+$, $u_b^-$ ($b\in\mathrm{Ln}(n)$) be $6|n|$ indeterminates.
Let $\mathrm{Dmp}$ be the matrix with entries in the polynomial ring $\mathbb{C}\left[u_b^+,u_b^-:\ b\in\mathrm{Ln}(n)\right]$ that is obtained from (38.2) when $U_B(x,x^{(i)})$ is replaced by $1$ for $i=0,1,2$,
and $U_B(x,x^{(3)})$ by $u_b^+$ if $x$ carries $(b,+)$ and by $u_b^-$ if $x$ carries $(b,-)$.
Each entry of the four rows $(x,1),\dots,(x,4)$ of $\mathrm{Dmp}$ is either a complex number or a complex number times the one indeterminate that belongs to $x$;
these complex numbers are $M_{\mathrm{ch}}$ and the numbers $-\mathrm{sn}(x^{(i)})\cdot\left(M_i(x)\right)_{\alpha\beta}$, and they contain neither $B$ nor $q$.
Put
\begin{equation}
\mathrm{Pt}_n:=\sum_{\mathrm{pe}\in\mathrm{Sym}(\mathrm{Ix})}\operatorname{sgn}(\mathrm{pe})\prod_{\mathrm{jx}\in\mathrm{Ix}}\mathrm{Dmp}[\mathrm{jx},\mathrm{pe}(\mathrm{jx})]
\qquad\text{(Leibniz formula)}.
\tag{38.3}
\end{equation}
Each of the products contains exactly one entry of every row.
By 1.2 a given indeterminate occurs in the entries of only $8\cdot4=32$ rows, and in each such entry with degree at most $1$.
Hence $\mathrm{Pt}_n$ is a polynomial in the $6|n|$ indeterminates, of degree at most $32$ in each of them.
The substitution $u_b^+\mapsto\mathrm{e}^{iqB(b)}$, $u_b^-\mapsto\mathrm{e}^{-iqB(b)}$ turns $\mathrm{Dmp}$ into $\mathrm{Dm}$; therefore
\begin{equation}
\det\nolimits_n(D_B+\boldsymbol{\mu})=\mathrm{Pt}_n\left(\left(\mathrm{e}^{iqB(b)}\right)_b,\left(\mathrm{e}^{-iqB(b)}\right)_b\right)\qquad\text{for every }B\in C^{1}(n)\otimes\mathbb{C}.
\tag{38.4}
\end{equation}
Under this substitution a monomial $\prod_b(u_b^+)^{\mathrm{kx}_b}(u_b^-)^{\mathrm{kx}'_b}$ becomes $\exp\left(iq\sum_b(\mathrm{kx}_b-\mathrm{kx}'_b)B(b)\right)$, the exponential of a linear form in the $3|n|$ numbers $B(b)$.
Thus $\det\nolimits_n(D_B+\boldsymbol{\mu})$ is a finite linear combination of exponentials of linear forms in these numbers.
By (ef3) it is an entire function of $B\in\mathbb{C}^{\mathrm{Ln}(n)}=\mathbb{C}^{3|n|}$ in the sense of 0.5;
and it is continuous there, since the exponential function and linear forms are continuous.

\emph{The parenthesis in (a).}
Let $\mathbb{T}:=\left(\sum_\mu n_\mu\mathbb{Z}e_\mu,\ \chi\right)$ with $\chi(n_\mu e_\mu)=-1$;
\cite[Definition~4.1]{QEDII} admits every subgroup of $\mathbb{Z}^3$ of rank three, hence all $n_\mu\ge1$.
By Proposition~46\,(a), (b), under the identification given by $\mathrm{geo}$, the following holds.
The spinors of $\mathbb{T}$ are the elements of $\mathfrak{S}^{\mathrm{ch}}(1;n)$.
For complex $B$, the chamber operator $D_B$ of \cite[Definition~2.10]{QEDII} with the mass parameter $M_{\mathrm{ch}}$ is $D_B+\boldsymbol{\mu}$.
The set $\mathrm{geo}(\mathrm{Xn})$ is a set of representatives of $\mathbb{X}_{\mathbb{T}}$ in $\mathbb{X}$, and it is taken as the set $\mathbb{F}$ of representatives of \cite[Definition~4.3]{QEDII}.
For $r=\mathrm{geo}(x)$ with $x\in\mathrm{Xn}$, and with the hopping matrix $w$ of \cite[Definition~2.10]{QEDII},
\[
\widehat{r^{(i)}}=\mathrm{geo}\left(\mathrm{red}(x^{(i)})\right),\qquad
\chi\left(g_i(r)\right)=\mathrm{sn}(x^{(i)}),\qquad
w\left(r,r^{(i)}\right)=M_i(x),\qquad
U_B\left(r,r^{(i)}\right)=U_B\left(x,x^{(i)}\right).
\]
So (38.2) is the display of \cite[Definition~4.3]{QEDII}.
Moreover $y(r)=g v_0$, and the coordinate of $g v_0$ of modulus $v_1$ (the first coordinate of $v_0$) is the coordinate $\nu$ with $g e_1=\pm e_\nu$, and it has the same sign.
Consequently the lower ends of the bonds over $b$ in the sense of \cite[Definition~4.4]{QEDII} are the sites that carry $(b,+)$, and the upper ends are the sites that carry $(b,-)$ (\cite[Lemma~4.6(d)]{QEDII}).
Hence the determinant polynomial $P_{\mathbb{T}}$ of \cite[Definition~4.4]{QEDII} is the image of $\mathrm{Pt}_n$ under $u_b^+\mapsto u_b$, $u_b^-\mapsto u_b^{-1}$;
the identity (38.4) takes the form $\det\nolimits_n(D_B+\boldsymbol{\mu})=P_{\mathbb{T}}\left(\left(\mathrm{e}^{iqB(b)}\right)_b\right)$;
and the bound on the degrees is \cite[Lemma~4.7]{QEDII}: $0\le n^{\uparrow}_b(\mathfrak{s})\le32$ and $0\le n^{\downarrow}_b(\mathfrak{s})\le32$ for every admissible permutation $\mathfrak{s}$.

\paragraph{1.4. Proof of (b).}
Let $m\in C^{1}(n;\mathbb{Z})$ and $x=(c,g)$.
As $q\,r_q=2\pi$ and $m(c,c_3)$ is an integer,
\[
\mathrm{e}^{iq(B+r_qm)(c,c_3)}=\mathrm{e}^{iqB(c,c_3)}\cdot\exp\left(2\pi i\,m(c,c_3)\right)=\mathrm{e}^{iqB(c,c_3)}.
\]
Hence $U_{B+r_qm}=U_B$, and $D_{B+r_qm}=D_B$ ((I.O4), last sentence).

\paragraph{1.5. Proof of (c).}
Let $\boldsymbol{\lambda}\in C^{0}(n)\otimes\mathbb{C}$.
The multiplication operator $\left(\mathrm{e}^{iq\boldsymbol{\lambda}}f\right)(x):=\mathrm{e}^{iq\boldsymbol{\lambda}(c)}f(x)$ ((I.POS)) maps $\mathfrak{S}^{\mathrm{ch}}(1;n)$ into itself, because $\boldsymbol{\lambda}$ is periodic;
its inverse is the multiplication by $\mathrm{e}^{-iq\boldsymbol{\lambda}}$.
Let $f\in\mathfrak{S}^{\mathrm{ch}}(1;n)$ and $x=(c,g)$.
For $i=0,1,2$ the site $x^{(i)}$ lies over $c$ and the phases equal $1$, so that
\[
U_{B-d\boldsymbol{\lambda}}(x,x^{(i)})\cdot\mathrm{e}^{iq\boldsymbol{\lambda}(c)}f(x^{(i)})=\mathrm{e}^{iq\boldsymbol{\lambda}(c)}\cdot U_B(x,x^{(i)})\cdot f(x^{(i)}).
\]
For $i=3$ one has $(d\boldsymbol{\lambda})(c,c_3)=\boldsymbol{\lambda}(c_3)-\boldsymbol{\lambda}(c)$ for both orientations of the link, whence
\[
U_{B-d\boldsymbol{\lambda}}(x,x^{(3)})\cdot\mathrm{e}^{iq\boldsymbol{\lambda}(c_3)}f(x^{(3)})
=\mathrm{e}^{iqB(c,c_3)}\cdot\mathrm{e}^{-iq\left(\boldsymbol{\lambda}(c_3)-\boldsymbol{\lambda}(c)\right)}\cdot\mathrm{e}^{iq\boldsymbol{\lambda}(c_3)}f(x^{(3)})
=\mathrm{e}^{iq\boldsymbol{\lambda}(c)}\cdot U_B(x,x^{(3)})\cdot f(x^{(3)}).
\]
Inserting these identities into the formula for $D_B$ in 0.4 one obtains
\[
\left(D_{B-d\boldsymbol{\lambda}}\,\mathrm{e}^{iq\boldsymbol{\lambda}}f\right)(x)
=-\sum_iM_i(x)\left[\mathrm{e}^{iq\boldsymbol{\lambda}(c)}\cdot U_B(x,x^{(i)})\cdot f(x^{(i)})-\mathrm{e}^{iq\boldsymbol{\lambda}(c)}f(x)\right]
=\mathrm{e}^{iq\boldsymbol{\lambda}(c)}(D_Bf)(x),
\]
which is (I.O4).
Consequently
\[
D_{B-d\boldsymbol{\lambda}}+\boldsymbol{\mu}=\mathrm{e}^{iq\boldsymbol{\lambda}}\left(D_B+\boldsymbol{\mu}\right)\mathrm{e}^{-iq\boldsymbol{\lambda}},
\]
and conjugate linear maps have the same determinant.
Statement (c) is this identity with $-\boldsymbol{\lambda}$ in place of $\boldsymbol{\lambda}$.

\paragraph{1.6. Proof of (d).}
On $\mathfrak{S}^{\mathrm{ch}}(1;n)$ let
\[
\|f\|^2:=\sum_{x\in\mathrm{Xn}}|f(x)|_2^2=\sum_{c\in\mathrm{Bn}}\sum_{g\in O_{\mathrm{cub}}}|f(c,g)|_2^2 ;
\]
by (I.1.3) at $\eta_0=1$ this is $48\,(f,f)^{\mathrm{ch}}$, and the two norms give the same operator norm.
The function $x\mapsto|f(x)|_2$ is periodic.
By (I.s1) the map $x\mapsto x^{(i)}$ is a bijection of $\mathbb{Z}^3\times O_{\mathrm{cub}}$ that commutes with the translations, so $x\mapsto\mathrm{red}(x^{(i)})$ is a bijection of $\mathrm{Xn}$.
By (I.m6), $|M_i(x)|_2=\frac{w_i}{2}(1+\rho)$.
For $i=0,\dots,3$ it follows that
\begin{align*}
\left\|M_i(\cdot)\,R_i(B)f\right\|^2&=\sum_{x\in\mathrm{Xn}}\left|M_i(x)\cdot U_B(x,x^{(i)})\cdot f(x^{(i)})\right|_2^2\\
&\le\left(\frac{w_i}{2}(1+\rho)\right)^2\sup_x\left|U_B(x,x^{(i)})\right|^2\sum_{x\in\mathrm{Xn}}\left|f\left(\mathrm{red}(x^{(i)})\right)\right|_2^2\\
&=\left(\frac{w_i}{2}(1+\rho)\cdot\sup_x\left|U_B(x,x^{(i)})\right|\right)^2\|f\|^2 .
\end{align*}
Here $|U_B(x,x^{(i)})|=1$ for $i=0,1,2$, while $|U_B(x,x^{(3)})|=\exp\left(-q\operatorname{Im}B(c,c_3)\right)\le\exp\left(q|\operatorname{Im}B|_\infty\right)$.
Together with (38.1), $w_3=6$ and $w_0+w_1+w_2=\mathrm{sw}-6$ (Section~I.0.5), this gives
\begin{equation}
\left\|D_B+\boldsymbol{\mu}\right\|\le|M_{\mathrm{ch}}|+\frac{1+\rho}{2}\left(\mathrm{sw}-6\right)+\frac{1+\rho}{2}\cdot6\cdot\exp\left(q|\operatorname{Im}B|_\infty\right).
\tag{38.5}
\end{equation}
For a linear map $O$ of a normed complex vector space of dimension $\mathrm{Dq}$, every eigenvalue has modulus at most $\|O\|$,
and $\det O$ is the product of the $\mathrm{Dq}$ eigenvalues, counted with their algebraic multiplicities;
hence $|\det O|\le\|O\|^{\mathrm{Dq}}$.
With $\mathrm{Dq}=192|n|$ and (38.5) this is (d).

\paragraph{1.7. Proof of (e).}
Fix $\mu$.
Let $\boldsymbol{\lambda}^{(\mu)}$ be the element of $C^{0}(n)$ with $\boldsymbol{\lambda}^{(\mu)}(c)=-\frac{r_q}{n_\mu}\,\mathrm{red}(c)_\mu$,
and let $m^{(\mu)}:=\sum_{c\in\mathrm{Bn},\ c_\mu=n_\mu-1}\mathbf{1}_{(c,\mu)}\in C^{1}(n;\mathbb{Z})$.
By 0.3, $H_{2\pi e_\mu}=\frac{2\pi}{q\,n_\mu}\,\mathbf{1}_\mu=\frac{r_q}{n_\mu}\,\mathbf{1}_\mu$.
Let $(c,\nu)\in\mathrm{Ln}(n)$.
\begin{itemize}
\item
If $\nu\ne\mu$, then $\mathrm{red}(c+e_\nu)_\mu=c_\mu$, hence $\left(d\boldsymbol{\lambda}^{(\mu)}\right)_\nu(c)=0$;
also $m^{(\mu)}_\nu(c)=0$ and $\left(H_{2\pi e_\mu}\right)_\nu(c)=0$.
\item
If $\nu=\mu$ and $c_\mu\le n_\mu-2$, then $\mathrm{red}(c+e_\mu)_\mu=c_\mu+1$, hence $\left(d\boldsymbol{\lambda}^{(\mu)}\right)_\mu(c)=-\frac{r_q}{n_\mu}$;
also $m^{(\mu)}_\mu(c)=0$, and $r_q\,m^{(\mu)}_\mu(c)-\left(d\boldsymbol{\lambda}^{(\mu)}\right)_\mu(c)=\frac{r_q}{n_\mu}$.
\item
If $\nu=\mu$ and $c_\mu=n_\mu-1$, then $\mathrm{red}(c+e_\mu)_\mu=0$, hence $\left(d\boldsymbol{\lambda}^{(\mu)}\right)_\mu(c)=0+\frac{r_q}{n_\mu}(n_\mu-1)=r_q-\frac{r_q}{n_\mu}$;
also $m^{(\mu)}_\mu(c)=1$, and $r_q\,m^{(\mu)}_\mu(c)-\left(d\boldsymbol{\lambda}^{(\mu)}\right)_\mu(c)=\frac{r_q}{n_\mu}$.
\end{itemize}
(For $n_\mu=1$ only the third case occurs in the direction $\mu$; then $\boldsymbol{\lambda}^{(\mu)}=0$ and $m^{(\mu)}=\mathbf{1}_\mu$.)
This proves
\[
H_{2\pi e_\mu}=r_q\,m^{(\mu)}-d\boldsymbol{\lambda}^{(\mu)}.
\]
The map $\theta\mapsto H_\theta$ is linear.
Therefore, for $\theta\in\mathbb{C}^3$, part (b) with $B+H_\theta-d\boldsymbol{\lambda}^{(\mu)}$ in place of $B$ and part (c) with $-\boldsymbol{\lambda}^{(\mu)}$ in place of $\boldsymbol{\lambda}$ give
\begin{align*}
\det\nolimits_n\left(D_{B+H_{\theta+2\pi e_\mu}}+\boldsymbol{\mu}\right)
&=\det\nolimits_n\left(D_{B+H_\theta-d\boldsymbol{\lambda}^{(\mu)}+r_qm^{(\mu)}}+\boldsymbol{\mu}\right)\\
&=\det\nolimits_n\left(D_{B+H_\theta-d\boldsymbol{\lambda}^{(\mu)}}+\boldsymbol{\mu}\right)
=\det\nolimits_n\left(D_{B+H_\theta}+\boldsymbol{\mu}\right).
\end{align*}
\end{proof}

\paragraph{1.8. Dependence on the parameters.}
Parts (a)--(e) of Proposition~38.1 hold for all integers $n_\mu\ge1$, all $q>0$, all real $\boldsymbol{\mu}$ and all complex $B$, $\boldsymbol{\lambda}$, $\theta$.
The numbers $\rho$, $\boldsymbol{\mu}$ and the weights $w_i$ enter the proofs of (a)--(c) and (e) only through the identity (I.2.5).
In (d) the numbers $\rho$, $\mathrm{sw}$ and $6=w_3$ are those of the line (C0) of Section~I.8.2;
the base of the power in (d) depends only on $(\boldsymbol{\mu},\rho,q|\operatorname{Im}B|_\infty)$, and the torus enters (d) only through the exponent $192|n|$.

\subsection*{2. Proposition~38.2. The zero-flux sector is the harmonic average}

\begin{proof}[Proof of Proposition~38.2]
The notation is that of Section~0, and throughout $n_\mu\ge4$ for $\mu=1,2,3$.
Put
\[
\mathrm{Bo}:=[0,r_q)^{\mathrm{Ln}(n)};
\]
this is a Borel subset of $C^{1}(n)=\mathbb{R}^{\mathrm{Ln}(n)}$, and its volume is $r_q^{3|n|}$.
For $w\in\mathbb{Z}^3$ let $\mathrm{Mz}_w$ denote the set of all $\mathfrak{z}\in C^{2}(n;\mathbb{Z})$ with $d\mathfrak{z}=0$ whose flux is $w$.
For $y\ge0$ let
\[
\mathrm{Kd}(y):=\left(|M_{\mathrm{ch}}|+\frac{1+\rho}{2}\left(\mathrm{sw}-6+6\exp(qy)\right)\right)^{192|n|};
\]
this is the right side of the inequality of Proposition~38.1\,(d) at $|\operatorname{Im}B|_\infty=y$.

Two facts on cochains will be needed.
They are proved in \cite[Section~10.1.1]{QEDI} for all sides $n_\mu\ge2$, and no fermion enters them.
\begin{itemize}
\item[(H1)]
One has the orthogonal decomposition $C^{1}(n)={V_{\mathrm{g}}}\oplus{V_{\mathrm{h}}}\oplus{V_{\mathrm{t}}}$,
where ${V_{\mathrm{g}}}=dC^{0}(n)$, the space ${V_{\mathrm{h}}}$ is spanned by $\mathbf{1}_1,\mathbf{1}_2,\mathbf{1}_3$, and ${V_{\mathrm{t}}}=\left({V_{\mathrm{g}}}+{V_{\mathrm{h}}}\right)^{\perp}$.
The kernel of $d$ on $C^{1}(n)$ is ${V_{\mathrm{g}}}+{V_{\mathrm{h}}}$, and $dC^{1}(n)$ is contained in the kernel of $d$ on $C^{2}(n)$.
\item[(H2)]
Let $\mathfrak{z}\in C^{2}(n;\mathbb{Z})$ satisfy $d\mathfrak{z}=0$.
The sums $w_{\mu\nu}$ of Section~I.7.1 do not depend on the plane over which they are taken,
and $\mathfrak{z}=dm$ for some $m\in C^{1}(n;\mathbb{Z})$ if and only if the flux of $\mathfrak{z}$ is $0$.
\end{itemize}

\paragraph{2.1. Majorants.}

\begin{named}{Lemma 2.A}
For $\epsilon>0$ put
\[
\mathrm{Cth}(\epsilon):=\sup_{y\in\mathbb{R}}\ \sum_{j\in\mathbb{Z}}\exp\left(-\epsilon\left(y+r_q j\right)^2\right),
\qquad
\mathrm{Cth}:=\mathrm{Cth}\left(\tfrac12\right).
\]
Then $\mathrm{Cth}(\epsilon)<\infty$, and
\begin{equation}
\sum_{\mathfrak{z}\in C^{2}(n;\mathbb{Z})}\ \int_{\mathrm{Bo}}\mathrm{d}A\,\exp\left(-\epsilon\left\|dA+r_q\mathfrak{z}\right\|^2\right)\le\left(r_q\cdot\mathrm{Cth}(\epsilon)\right)^{3|n|}.
\tag{38.6}
\end{equation}
If a function $\Phi$ on $C^{2}(n)$ satisfies $|\Phi(u)|\le\mathrm{cf}\cdot\exp\left(\mathrm{cf}\cdot\|u\|\right)$ for all $u\in C^{2}(n)$, then
\begin{equation}
|\Phi(u)|\cdot\exp\left(-\frac{1}{2}\|u\|^2\right)\le\mathrm{cf}\cdot\exp\left(\mathrm{cf}^2\right)\cdot\exp\left(-\frac{1}{4}\|u\|^2\right)
\qquad\left(u\in C^{2}(n)\right).
\tag{38.7}
\end{equation}
\end{named}

\begin{proof}
The estimate (38.6) is the chain of inequalities of \cite[Lemma~3.13]{QEDII} (where $\epsilon=\frac12$) and of the proof of \cite[Lemma~3.14]{QEDII} (for every $\epsilon>0$),
from the point at which the condition $d\mathfrak{z}=0$ is dropped onwards,
written for the period lattice $\sum_\mu n_\mu\mathbb{Z}e_\mu$.
The shortest non-zero vectors of this lattice have the sup norm $\min_\mu n_\mu\ge4$, so that it is a period lattice in the sense of \cite[Definition~2.6]{QEDII}.
In detail, the argument is the following.

The sum over $j$ in the definition of $\mathrm{Cth}(\epsilon)$ is a periodic function of $y$ with the period $r_q$.
Let $0\le y<r_q$.
Then $(y+r_q j)^2\ge r_q^2 j^2$ for $j\ge0$, and $|y+r_q j|\ge r_q\left(|j|-1\right)$ for $j\le-1$.
Consequently
\[
\mathrm{Cth}(\epsilon)\le2\sum_{j\ge0}\exp\left(-\epsilon r_q^2 j^2\right)<\infty.
\]
An element $\mathfrak{z}$ of $C^{2}(n;\mathbb{Z})$ is an arbitrary family of integers $\mathfrak{z}_p$ indexed by the $3|n|$ plaquettes $p$ of the torus,
and $\|dA+r_q\mathfrak{z}\|^2=\sum_p\left((dA)_p+r_q\mathfrak{z}_p\right)^2$.
Hence, for fixed $A$,
\[
\sum_{\mathfrak{z}\in C^{2}(n;\mathbb{Z})}\exp\left(-\epsilon\left\|dA+r_q\mathfrak{z}\right\|^2\right)
=\prod_p\ \sum_{j\in\mathbb{Z}}\exp\left(-\epsilon\left((dA)_p+r_q j\right)^2\right)
\le\mathrm{Cth}(\epsilon)^{3|n|}.
\]
All terms are non-negative and continuous in $A$.
Therefore the sum over $\mathfrak{z}$ may be taken under the integral over $\mathrm{Bo}$, by the termwise integration of a series of non-negative measurable functions \cite[Theorem~1.27]{Rudin1987};
since $\mathrm{Bo}$ has the volume $r_q^{3|n|}$, this gives (38.6).

The inequality (38.7) follows from $\mathrm{cf}\cdot\|u\|\le\frac14\|u\|^2+\mathrm{cf}^2$, which holds because
\[
\frac{1}{4}\|u\|^2-\mathrm{cf}\cdot\|u\|+\mathrm{cf}^2=\left(\frac{1}{2}\|u\|-\mathrm{cf}\right)^2\ge0.
\]
\end{proof}

\paragraph{2.2. Two fundamental domains.}

\begin{named}{Lemma 2.B}
Let $\mathrm{Vq}$ be a real vector space of finite dimension with a Lebesgue measure $\mathrm{d}A$ (a translation invariant Borel measure), and let $\mathrm{Gq}$ be a countable subgroup of $\mathrm{Vq}$.
A Borel set $\mathrm{Da}\subset\mathrm{Vq}$ is called a fundamental domain of $\mathrm{Gq}$ if every $A\in\mathrm{Vq}$ can be written as $A=A'+\mathrm{tv}$ with exactly one pair $(A',\mathrm{tv})\in\mathrm{Da}\times\mathrm{Gq}$.
Let $\mathrm{Da}$ and $\mathrm{Db}$ be fundamental domains of $\mathrm{Gq}$,
and let $\mathrm{Gv}$ be a Borel function on $\mathrm{Vq}$ with $\mathrm{Gv}(A+\mathrm{tv})=\mathrm{Gv}(A)$ for all $A\in\mathrm{Vq}$ and all $\mathrm{tv}\in\mathrm{Gq}$.
If $\mathrm{Gv}\ge0$, then
\[
\int_{\mathrm{Da}}\mathrm{d}A\,\mathrm{Gv}(A)=\int_{\mathrm{Db}}\mathrm{d}A\,\mathrm{Gv}(A).
\]
If $\mathrm{Gv}$ is complex valued and $\int_{\mathrm{Da}}\mathrm{d}A\,|\mathrm{Gv}(A)|<\infty$,
then $\int_{\mathrm{Db}}\mathrm{d}A\,|\mathrm{Gv}(A)|<\infty$, and the integrals of $\mathrm{Gv}$ over $\mathrm{Da}$ and over $\mathrm{Db}$ are equal.
\end{named}

\begin{proof}
The space $\mathrm{Vq}$ is the disjoint union of the sets $\mathrm{Db}+\mathrm{tv}$, $\mathrm{tv}\in\mathrm{Gq}$, and it is also the disjoint union of the sets $\mathrm{Da}-\mathrm{tv}$, $\mathrm{tv}\in\mathrm{Gq}$.
Let $\mathrm{Gv}\ge0$.
The countable additivity of the integral (\cite[Theorem~1.27]{Rudin1987}, applied to the functions $\mathbf{1}_{\mathrm{Db}+\mathrm{tv}}\cdot\mathbf{1}_{\mathrm{Da}}\cdot\mathrm{Gv}$, respectively $\mathbf{1}_{\mathrm{Da}-\mathrm{tv}}\cdot\mathbf{1}_{\mathrm{Db}}\cdot\mathrm{Gv}$),
the invariance of $\mathrm{d}A$ under translations and the invariance of $\mathrm{Gv}$ give
\[
\begin{aligned}
\int_{\mathrm{Da}}\mathrm{d}A\,\mathrm{Gv}(A)
&=\sum_{\mathrm{tv}\in\mathrm{Gq}}\ \int_{\mathrm{Da}\cap(\mathrm{Db}+\mathrm{tv})}\mathrm{d}A\,\mathrm{Gv}(A)
=\sum_{\mathrm{tv}\in\mathrm{Gq}}\ \int_{(\mathrm{Da}-\mathrm{tv})\cap\mathrm{Db}}\mathrm{d}A\,\mathrm{Gv}(A+\mathrm{tv})\\
&=\sum_{\mathrm{tv}\in\mathrm{Gq}}\ \int_{(\mathrm{Da}-\mathrm{tv})\cap\mathrm{Db}}\mathrm{d}A\,\mathrm{Gv}(A)
=\int_{\mathrm{Db}}\mathrm{d}A\,\mathrm{Gv}(A).
\end{aligned}
\]
For complex valued $\mathrm{Gv}$, apply this first to $|\mathrm{Gv}|$, and then to the positive and the negative parts of $\operatorname{Re}\mathrm{Gv}$ and of $\operatorname{Im}\mathrm{Gv}$.
These four functions are non-negative Borel functions, invariant under $\mathrm{Gq}$ and bounded by $|\mathrm{Gv}|$, so that their integrals over $\mathrm{Da}$ and over $\mathrm{Db}$ are finite;
and $\mathrm{Gv}$ is a linear combination of them.
\end{proof}

\paragraph{2.3. Closed integer one-cochains.}

\begin{named}{Lemma 2.C}
Let $\mathrm{Zcl}:=\left\{m\in C^{1}(n;\mathbb{Z}):\ dm=0\right\}$, and let $m^{(\mu)}$ and $\boldsymbol{\lambda}^{(\mu)}$ be as in Proposition~38.1\,(e).
\begin{itemize}
\item[(i)]
If $\boldsymbol{\lambda}\in C^{0}(n)$ and $d\boldsymbol{\lambda}$ is integer valued, then $\boldsymbol{\lambda}-\boldsymbol{\lambda}(0)\in C^{0}(n;\mathbb{Z})$;
if $d\boldsymbol{\lambda}=0$, then $\boldsymbol{\lambda}$ is constant.
Consequently $\mathrm{Zcl}\cap{V_{\mathrm{g}}}=dC^{0}(n;\mathbb{Z})$.
\item[(ii)]
$m^{(\mu)}\in\mathrm{Zcl}$ for $\mu=1,2,3$.
Every $\mathrm{mc}\in\mathrm{Zcl}$ can be written as
\[
\mathrm{mc}=d\boldsymbol{\lambda}'+\sum_\mu\mathrm{kz}_\mu\cdot m^{(\mu)}
\qquad\text{with }\boldsymbol{\lambda}'\in C^{0}(n;\mathbb{Z})\text{ and }\mathrm{kz}\in\mathbb{Z}^3,
\]
and $\mathrm{kz}$ and $d\boldsymbol{\lambda}'$ are determined by $\mathrm{mc}$.
\item[(iii)]
For $\mathrm{mc}$ as in (ii),
\[
r_q\cdot\mathrm{mc}=\left[r_q\,d\boldsymbol{\lambda}'+\sum_\mu\mathrm{kz}_\mu\cdot d\boldsymbol{\lambda}^{(\mu)}\right]+H_{2\pi\mathrm{kz}},
\]
where the bracket lies in ${V_{\mathrm{g}}}$ and $H_{2\pi\mathrm{kz}}$ lies in ${V_{\mathrm{h}}}$.
\end{itemize}
\end{named}

\begin{proof}
(i)
Every $c\in\mathbb{Z}^3$ is joined to $0$ by a lattice path, and $\boldsymbol{\lambda}(c)-\boldsymbol{\lambda}(0)$ is the sum of the values of $d\boldsymbol{\lambda}$ on the oriented links of this path.
This gives the first two assertions.
If $\mathrm{mc}\in\mathrm{Zcl}\cap{V_{\mathrm{g}}}$, then $\mathrm{mc}=d\boldsymbol{\lambda}=d\left(\boldsymbol{\lambda}-\boldsymbol{\lambda}(0)\right)$ for some $\boldsymbol{\lambda}\in C^{0}(n)$, and $\boldsymbol{\lambda}-\boldsymbol{\lambda}(0)\in C^{0}(n;\mathbb{Z})$.
Conversely, the elements of $dC^{0}(n;\mathbb{Z})$ are integer valued and lie in ${V_{\mathrm{g}}}$, which is contained in the kernel of $d$ by (H1).

(ii), (iii)
By Proposition~38.1\,(e) and the definition of $H_\theta$ in 0.3,
\[
r_q\,m^{(\mu)}=H_{2\pi e_\mu}+d\boldsymbol{\lambda}^{(\mu)}=\frac{r_q}{n_\mu}\,\mathbf{1}_\mu+d\boldsymbol{\lambda}^{(\mu)}.
\]
The right side lies in ${V_{\mathrm{h}}}+{V_{\mathrm{g}}}$, which is the kernel of $d$ by (H1); hence $m^{(\mu)}\in\mathrm{Zcl}$.
Let $\mathrm{mc}\in\mathrm{Zcl}$.
By (H1) one can write $\mathrm{mc}=d\boldsymbol{\lambda}+\sum_\mu\mathrm{tq}_\mu\cdot\mathbf{1}_\mu$ with $\boldsymbol{\lambda}\in C^{0}(n)$ and $\mathrm{tq}\in\mathbb{R}^3$.
For $c\in\mathbb{Z}^3$, the sum of $\mathrm{mc}$ along the closed loop of the direction $\mu$ through $c$ is
\[
\sum_{j=0}^{n_\mu-1}\mathrm{mc}\left(c+je_\mu,\,c+(j+1)e_\mu\right)=\boldsymbol{\lambda}(c+n_\mu e_\mu)-\boldsymbol{\lambda}(c)+n_\mu\mathrm{tq}_\mu=n_\mu\mathrm{tq}_\mu .
\]
This number is an integer; denote it by $\mathrm{kz}_\mu$.
Thus $\mathrm{mc}=d\boldsymbol{\lambda}+\sum_\mu\frac{\mathrm{kz}_\mu}{n_\mu}\,\mathbf{1}_\mu$, and
\[
\mathrm{mc}-\sum_\mu\mathrm{kz}_\mu\cdot m^{(\mu)}=d\left(\boldsymbol{\lambda}-r_q^{-1}\sum_\mu\mathrm{kz}_\mu\cdot\boldsymbol{\lambda}^{(\mu)}\right)
\]
is integer valued and lies in ${V_{\mathrm{g}}}$.
By (i) it equals $d\boldsymbol{\lambda}'$ for some $\boldsymbol{\lambda}'\in C^{0}(n;\mathbb{Z})$.
This proves the representation in (ii), and (iii) follows from it, since $\sum_\mu\mathrm{kz}_\mu H_{2\pi e_\mu}=H_{2\pi\mathrm{kz}}$.
As to uniqueness, the component of $\mathrm{mc}$ in ${V_{\mathrm{h}}}$ is $\sum_\mu\frac{\mathrm{kz}_\mu}{n_\mu}\,\mathbf{1}_\mu$;
it determines $\mathrm{kz}$, and then $d\boldsymbol{\lambda}'=\mathrm{mc}-\sum_\mu\mathrm{kz}_\mu\cdot m^{(\mu)}$.
\end{proof}

\paragraph{2.4. Modes.}

\begin{named}{Lemma 2.D}
Let
\[
\mathrm{Pm}:=\left\{p\in\mathbb{R}^3:\ p_\mu=\frac{2\pi k_\mu}{n_\mu},\ k_\mu\in\{0,\dots,n_\mu-1\}\right\},
\qquad
\mathrm{em}_p(c):=|n|^{-\frac12}\,\mathrm{e}^{ip\cdot c},
\]
\[
\mathrm{xq}_\mu(p):=\mathrm{e}^{ip_\mu}-1,
\qquad
|\mathrm{xq}(p)|^2:=\sum_\mu|\mathrm{xq}_\mu(p)|^2=\sum_\mu\left(2-2\cos p_\mu\right),
\]
and let $(\,\cdot\,,\,\cdot\,)$ be the Hermitian extension of $\langle\,\cdot\,,\,\cdot\,\rangle$ to complex cochains, antilinear in the first argument.
\begin{itemize}
\item[(i)]
$\Delta_0\mathrm{em}_p=|\mathrm{xq}(p)|^2\cdot\mathrm{em}_p$, and $|\mathrm{xq}(p)|^2>0$ for $p\in\mathrm{Pm}\setminus\{0\}$.
Hence $0$ is a simple eigenvalue of $\Delta_0$, and
\[
\operatorname{det}'\Delta_0=\prod_{p\in\mathrm{Pm},\ p\ne0}\ \sum_\mu\left(2-2\cos p_\mu\right).
\]
\item[(ii)]
The map $d^{*}d$ sends ${V_{\mathrm{t}}}$ into itself and is positive definite on ${V_{\mathrm{t}}}$;
$\dim{V_{\mathrm{t}}}=2|n|-2$, and $\det\left(d^{*}d|_{{V_{\mathrm{t}}}}\right)=\left(\operatorname{det}'\Delta_0\right)^2$.
\item[(iii)]
\[
\mathrm{zV}_n:=\int_{{V_{\mathrm{t}}}}\mathrm{d}\sigma(A)\,\exp\left(-\frac{1}{2}\|dA\|^2\right)=\frac{(2\pi)^{|n|-1}}{\operatorname{det}'\Delta_0}.
\]
For $|n|=n^3$ this is the value of $\mathrm{zV}_n$ in Proposition~38.2\,(c).
\end{itemize}
\end{named}

\begin{proof}
For $p,p'\in\mathrm{Pm}$,
\[
\sum_{c\in\mathrm{Bn}}\mathrm{e}^{i(p-p')\cdot c}=\prod_\mu\ \sum_{j=0}^{n_\mu-1}\exp\left(\frac{2\pi i\left(k_\mu-k'_\mu\right)j}{n_\mu}\right)=|n|\cdot\mathbf{1}[p=p'];
\]
so $(\mathrm{em}_p)_{p\in\mathrm{Pm}}$ is an orthonormal basis of $C^{0}(n)\otimes\mathbb{C}$.
For $\mathrm{za}\in\mathbb{C}^3$ let $\mathrm{za}\cdot\mathrm{em}_p$ denote the one-cochain with the components $\mathrm{za}_\mu\cdot\mathrm{em}_p$.
These cochains span $C^{1}(n)\otimes\mathbb{C}$, and
\[
\left(\mathrm{za}\cdot\mathrm{em}_p,\ \mathrm{za}'\cdot\mathrm{em}_{p'}\right)=\mathbf{1}[p=p']\sum_\mu\overline{\mathrm{za}_\mu}\cdot\mathrm{za}'_\mu .
\]
The formulas for $d$ in 0.3 give
\begin{equation}
\begin{aligned}
d\,\mathrm{em}_p&=\mathrm{xq}(p)\cdot\mathrm{em}_p,\\
\left(d(\mathrm{za}\cdot\mathrm{em}_p)\right)_{\mu\nu}
&=\left(\mathrm{za}_\mu+\mathrm{za}_\nu\,\mathrm{e}^{ip_\mu}-\mathrm{za}_\mu\,\mathrm{e}^{ip_\nu}-\mathrm{za}_\nu\right)\mathrm{em}_p
=\left(\mathrm{xq}_\mu(p)\cdot\mathrm{za}_\nu-\mathrm{xq}_\nu(p)\cdot\mathrm{za}_\mu\right)\mathrm{em}_p .
\end{aligned}
\tag{38.8}
\end{equation}

(i)
By (38.8), $(\mathrm{em}_{p'},\Delta_0\mathrm{em}_p)=(d\,\mathrm{em}_{p'},d\,\mathrm{em}_p)=\mathbf{1}[p=p']\cdot|\mathrm{xq}(p)|^2$;
so $\Delta_0$ is diagonal in the basis $(\mathrm{em}_p)$, with the eigenvalues $|\mathrm{xq}(p)|^2$.
For $0\le p_\mu<2\pi$ one has $2-2\cos p_\mu=0$ only if $p_\mu=0$.

(ii)
The range of $d^{*}$ is the orthogonal complement of the kernel of $d$, that is, ${V_{\mathrm{t}}}$ by (H1);
hence $d^{*}d$ maps $C^{1}(n)$ into ${V_{\mathrm{t}}}$.
By (38.8), ${V_{\mathrm{g}}}\otimes\mathbb{C}$ is spanned by the cochains $\mathrm{xq}(p)\cdot\mathrm{em}_p$ with $p\ne0$, and ${V_{\mathrm{h}}}\otimes\mathbb{C}$ is spanned by the cochains $\mathrm{za}\cdot\mathrm{em}_0$.
Therefore ${V_{\mathrm{t}}}\otimes\mathbb{C}$ is the orthogonal sum, over the $|n|-1$ elements $p\ne0$ of $\mathrm{Pm}$, of the spaces
\[
\mathrm{Wm}_p:=\left\{\mathrm{za}\cdot\mathrm{em}_p:\ \sum_\mu\overline{\mathrm{xq}_\mu(p)}\cdot\mathrm{za}_\mu=0\right\},
\]
each of which has the complex dimension $2$.
For $\mathrm{zb},\mathrm{zc}\in\mathbb{C}^3$ one has
\[
\begin{aligned}
\sum_{\mu<\nu}\left|\mathrm{zb}_\mu\cdot\mathrm{zc}_\nu-\mathrm{zb}_\nu\cdot\mathrm{zc}_\mu\right|^2
&=\frac{1}{2}\sum_{\mu,\nu}\left(|\mathrm{zb}_\mu|^2|\mathrm{zc}_\nu|^2+|\mathrm{zb}_\nu|^2|\mathrm{zc}_\mu|^2-2\operatorname{Re}\left[\mathrm{zb}_\mu\cdot\overline{\mathrm{zc}_\mu}\cdot\overline{\mathrm{zb}_\nu}\cdot\mathrm{zc}_\nu\right]\right)\\
&=|\mathrm{zb}|_2^2\cdot|\mathrm{zc}|_2^2-\left|\sum_\mu\overline{\mathrm{zb}_\mu}\cdot\mathrm{zc}_\mu\right|^2 .
\end{aligned}
\]
Let $A=\sum_{p\ne0}\mathrm{za}(p)\cdot\mathrm{em}_p$ with $\mathrm{za}(p)\cdot\mathrm{em}_p\in\mathrm{Wm}_p$.
By (38.8) the two-cochains $d\left(\mathrm{za}(p)\cdot\mathrm{em}_p\right)$ are pairwise orthogonal,
and the identity above with $\mathrm{zb}=\mathrm{xq}(p)$ and $\mathrm{zc}=\mathrm{za}(p)$ gives
\[
\left(A,\,d^{*}dA\right)=(dA,dA)=\sum_{p\ne0}|\mathrm{xq}(p)|^2\cdot|\mathrm{za}(p)|_2^2 .
\]
A Hermitian form determines its operator; so $d^{*}d$ acts on $\mathrm{Wm}_p$ as the multiplication by $|\mathrm{xq}(p)|^2>0$.
The determinant of a real linear map equals the determinant of its complex extension, and therefore
\[
\det\left(d^{*}d|_{{V_{\mathrm{t}}}}\right)=\prod_{p\ne0}\left(|\mathrm{xq}(p)|^2\right)^2=\left(\operatorname{det}'\Delta_0\right)^2 .
\]

(iii)
One has $\int_{\mathbb{R}}\mathrm{d}\mathrm{tq}\,\exp\left(-\mathrm{tq}^2\right)=\pi^{\frac12}$ \cite[8.21]{Rudin1976}.
For $\mathrm{lq}>0$, (Lin) for the map $\mathrm{tq}\mapsto\left(\frac{\mathrm{lq}}{2}\right)^{\frac12}\mathrm{tq}$ of $\mathbb{R}$ and the function $\exp\left(-\mathrm{tq}^2\right)$ gives
\[
\int_{\mathbb{R}}\mathrm{d}\mathrm{tq}\,\exp\left(-\frac{\mathrm{lq}}{2}\,\mathrm{tq}^2\right)=\left(\frac{2}{\mathrm{lq}}\right)^{\frac12}\pi^{\frac12}=\left(\frac{2\pi}{\mathrm{lq}}\right)^{\frac12}.
\]
Let $\mathrm{Sy}$ be a positive definite symmetric map of a Euclidean space of dimension $\mathrm{Dq}$, with the eigenvalues $\mathrm{lq}_1,\dots,\mathrm{lq}_{\mathrm{Dq}}$.
In the coordinates $\mathrm{tq}_1,\dots,\mathrm{tq}_{\mathrm{Dq}}$ with respect to an orthonormal basis of eigenvectors (see 0.3), Tonelli's theorem gives
\[
\int\mathrm{d}A\,\exp\left(-\frac{1}{2}\langle A,\mathrm{Sy}\,A\rangle\right)
=\int_{\mathbb{R}^{\mathrm{Dq}}}\mathrm{d}\mathrm{tq}\,\exp\left(-\frac{1}{2}\sum_j\mathrm{lq}_j\,\mathrm{tq}_j^2\right)
=\prod_j\left(\frac{2\pi}{\mathrm{lq}_j}\right)^{\frac12}
=(2\pi)^{\frac{\mathrm{Dq}}{2}}\left(\det\mathrm{Sy}\right)^{-\frac12}.
\]
For $\mathrm{Sy}=d^{*}d|_{{V_{\mathrm{t}}}}$ one has $\langle A,\mathrm{Sy}\,A\rangle=\|dA\|^2$ and $\mathrm{Dq}=2|n|-2$, and together with (ii) this is (iii).
\end{proof}

\paragraph{2.5. The reduced Laplacian.}

\begin{named}{Lemma 2.E}
The cochains $d\mathbf{1}_c$, $c\in\mathrm{Bn}\setminus\{0\}$, form a basis of the real vector space ${V_{\mathrm{g}}}$ and a basis of the group $dC^{0}(n;\mathbb{Z})$, and
\[
\det\left(\langle d\mathbf{1}_c,d\mathbf{1}_{c'}\rangle\right)_{c,c'\in\mathrm{Bn}\setminus\{0\}}=\frac{\operatorname{det}'\Delta_0}{|n|}.
\]
\end{named}

\begin{proof}
The sum $\sum_{c\in\mathrm{Bn}}\mathbf{1}_c$ is the constant function $1$; hence $d\mathbf{1}_0=-\sum_{c\ne0}d\mathbf{1}_c$, and
\[
d\left(\sum_{c\in\mathrm{Bn}}\mathrm{tq}_c\cdot\mathbf{1}_c\right)=\sum_{c\ne0}\left(\mathrm{tq}_c-\mathrm{tq}_0\right)d\mathbf{1}_c .
\]
So the $d\mathbf{1}_c$ with $c\ne0$ span ${V_{\mathrm{g}}}$ over $\mathbb{R}$ and $dC^{0}(n;\mathbb{Z})$ over $\mathbb{Z}$.
If $\sum_{c\ne0}\mathrm{tq}_c\cdot d\mathbf{1}_c=0$, then $\boldsymbol{\lambda}:=\sum_{c\ne0}\mathrm{tq}_c\cdot\mathbf{1}_c$ is constant by Lemma~2.C\,(i), and $\boldsymbol{\lambda}(0)=0$;
hence all $\mathrm{tq}_c$ vanish.

For $c,c'\in\mathrm{Bn}$ put $\mathrm{Sy}_{cc'}:=\langle d\mathbf{1}_c,d\mathbf{1}_{c'}\rangle=\langle\mathbf{1}_c,\Delta_0\mathbf{1}_{c'}\rangle$;
this is the matrix of $\Delta_0$ in the orthonormal basis $(\mathbf{1}_c)_{c\in\mathrm{Bn}}$.
Let $\mathrm{Sy}^{(c)}$ be the matrix obtained from $\mathrm{Sy}$ by deleting the row $c$ and the column $c$,
let $\mathrm{tq}$ be a real variable, and let $\mathrm{tq}+\mathrm{Sy}$ denote the sum of $\mathrm{Sy}$ and $\mathrm{tq}$ times the identity matrix.
\begin{itemize}
\item[(1)]
By Lemma~2.D\,(i), $\det(\mathrm{tq}+\mathrm{Sy})=\mathrm{tq}\cdot\prod_{p\ne0}\left(\mathrm{tq}+|\mathrm{xq}(p)|^2\right)$; its coefficient of $\mathrm{tq}^1$ is $\operatorname{det}'\Delta_0$.
\item[(2)]
The determinant is multilinear in the columns, and the column $c$ of $\mathrm{tq}+\mathrm{Sy}$ is the sum of $\mathrm{tq}$ times the standard basis vector of index $c$ and the column $c$ of $\mathrm{Sy}$.
Hence the coefficient of $\mathrm{tq}^1$ is the sum over $c\in\mathrm{Bn}$ of the determinants of the matrices that arise from $\mathrm{Sy}$ when its column $c$ is replaced by this basis vector, that is, $\sum_{c\in\mathrm{Bn}}\det\mathrm{Sy}^{(c)}$.
\item[(3)]
For $\mathrm{tv}\in\mathbb{Z}^3$ the map $c\mapsto\mathrm{red}(c+\mathrm{tv})$ is a bijection of $\mathrm{Bn}$, and $\mathrm{Sy}_{\mathrm{red}(c+\mathrm{tv})\,\mathrm{red}(c'+\mathrm{tv})}=\mathrm{Sy}_{cc'}$, because the translations commute with $d$ and preserve $\langle\,\cdot\,,\,\cdot\,\rangle$.
So $\mathrm{Sy}^{(\mathrm{red}(c+\mathrm{tv}))}$ arises from $\mathrm{Sy}^{(c)}$ by one and the same bijection of the rows and of the columns, and $\det\mathrm{Sy}^{(\mathrm{red}(c+\mathrm{tv}))}=\det\mathrm{Sy}^{(c)}$.
\end{itemize}
By (1)--(3), $|n|\cdot\det\mathrm{Sy}^{(0)}=\operatorname{det}'\Delta_0$, and $\mathrm{Sy}^{(0)}$ is the matrix of the assertion.
\end{proof}

\paragraph{2.6. Proof of (a).}
Let $\Phi\in\mathfrak{F}(n)$ with $|\Phi(u)|\le\mathrm{cf}\cdot\exp\left(\mathrm{cf}\cdot\|u\|\right)$, and put
\[
\mathrm{Ps}(A):=\Phi(dA)\cdot\exp\left(-\frac{1}{2}\|dA\|^2\right)\cdot\mathrm{wq}(A),\qquad A\in C^{1}(n).
\]
The function $\mathrm{Ps}$ is Borel, since $\Phi$ is Borel, $d$ is linear and $\mathrm{wq}$ is continuous.

\emph{Step 1 (the left side converges absolutely).}
By (38.7), the modulus of the integrand on the left side of (a) is at most
$\sup|\mathrm{wq}|\cdot\mathrm{cf}\cdot\exp\left(\mathrm{cf}^2\right)\cdot\exp\left(-\frac14\|dA+r_q\mathfrak{z}\|^2\right)$.
By (38.6) with $\epsilon=\frac14$, the sum over all $\mathfrak{z}\in C^{2}(n;\mathbb{Z})$ of the integrals of this majorant over $\mathrm{Bo}$ is at most
$\sup|\mathrm{wq}|\cdot\mathrm{cf}\cdot\exp\left(\mathrm{cf}^2\right)\cdot\left(r_q\cdot\mathrm{Cth}\left(\tfrac14\right)\right)^{3|n|}$.

\emph{Step 2 (unfolding).}
Let $\mathrm{Rep}\subset C^{1}(n;\mathbb{Z})$ contain exactly one element of every class of $C^{1}(n;\mathbb{Z})/\mathrm{Zcl}$.
For $m\in C^{1}(n;\mathbb{Z})$ the cochain $dm$ is integer valued and closed, by (H1); hence $\mathrm{Mz}_0=dC^{1}(n;\mathbb{Z})$ by (H2).
Consequently $\mathrm{mr}\mapsto d\,\mathrm{mr}$ maps $\mathrm{Rep}$ into $\mathrm{Mz}_0$.
This map is onto, because $dm=d\,\mathrm{mr}$ when $\mathrm{mr}$ is the representative of the class of $m$;
and it is injective, because $d\,\mathrm{mr}=d\,\mathrm{mr}'$ means $\mathrm{mr}-\mathrm{mr}'\in\mathrm{Zcl}$.
For $m\in C^{1}(n;\mathbb{Z})$ the invariance $\mathrm{wq}(A+r_q m)=\mathrm{wq}(A)$ gives
\[
\Phi\left(dA+r_q\,dm\right)\cdot\exp\left(-\frac{1}{2}\left\|dA+r_q\,dm\right\|^2\right)\cdot\mathrm{wq}(A)=\mathrm{Ps}(A+r_q m).
\]
The sets $\mathrm{Bo}+r_q m$, $m\in C^{1}(n;\mathbb{Z})$, are pairwise disjoint and cover $C^{1}(n)$, by division with remainder by $r_q$ in each coordinate $A(b)$.
Therefore the left side of (a) equals
\begin{equation}
\sum_{\mathrm{mr}\in\mathrm{Rep}}\ \int_{\mathrm{Bo}}\mathrm{d}A\,\mathrm{Ps}(A+r_q\,\mathrm{mr})
=\sum_{\mathrm{mr}\in\mathrm{Rep}}\ \int_{\mathrm{Bo}+r_q\mathrm{mr}}\mathrm{d}A\,\mathrm{Ps}(A)
=\int_{\mathrm{Da}}\mathrm{d}A\,\mathrm{Ps}(A),
\qquad
\mathrm{Da}:=\bigcup_{\mathrm{mr}\in\mathrm{Rep}}\left(\mathrm{Bo}+r_q\,\mathrm{mr}\right),
\tag{38.9}
\end{equation}
and $\int_{\mathrm{Da}}\mathrm{d}A\,|\mathrm{Ps}(A)|<\infty$ by Step~1.
The set $\mathrm{Da}$ is Borel, being a countable union of boxes.

\emph{The set $\mathrm{Da}$ is a fundamental domain of $\mathrm{Gq}:=r_q\,\mathrm{Zcl}$.}
Let $A\in C^{1}(n)$.
Then $A=A'+r_q m$ with a unique $A'\in\mathrm{Bo}$ and a unique $m\in C^{1}(n;\mathbb{Z})$,
and $m=\mathrm{mr}+\mathrm{mc}$ with a unique $\mathrm{mr}\in\mathrm{Rep}$ and a unique $\mathrm{mc}\in\mathrm{Zcl}$;
so $A=(A'+r_q\,\mathrm{mr})+r_q\,\mathrm{mc}$ with $A'+r_q\,\mathrm{mr}\in\mathrm{Da}$.
If also $A=(A''+r_q\,\mathrm{mr}')+r_q\,\mathrm{mc}'$ with $A''\in\mathrm{Bo}$, $\mathrm{mr}'\in\mathrm{Rep}$ and $\mathrm{mc}'\in\mathrm{Zcl}$,
then $A''=A'$ and $\mathrm{mr}'+\mathrm{mc}'=m$, whence $\mathrm{mr}'=\mathrm{mr}$ and $\mathrm{mc}'=\mathrm{mc}$.

\emph{The function $\mathrm{Ps}$ is invariant under $\mathrm{Gq}$.}
For $\mathrm{mc}\in\mathrm{Zcl}$ one has $d(A+r_q\,\mathrm{mc})=dA$ and $\mathrm{wq}(A+r_q\,\mathrm{mc})=\mathrm{wq}(A)$.

\emph{Step 3 (a second fundamental domain).}
Put
\[
\mathrm{Fg}:=\left\{r_q\sum_{c\in\mathrm{Bn}\setminus\{0\}}\mathrm{tq}_c\cdot d\mathbf{1}_c:\ 0\le\mathrm{tq}_c<1\right\}\subset{V_{\mathrm{g}}},
\qquad
\mathrm{Fh}:=\left\{H_\theta:\ \theta\in[-\pi,\pi)^3\right\}\subset{V_{\mathrm{h}}},
\qquad
\mathrm{Db}:=\mathrm{Fg}+\mathrm{Fh}+{V_{\mathrm{t}}}.
\]
The set $\mathrm{Db}$ is the image of the Borel set $\mathrm{Fg}\times\mathrm{Fh}\times{V_{\mathrm{t}}}$ under the linear bijection $(A_{\mathrm{g}},A_{\mathrm{h}},A_{\mathrm{t}})\mapsto A_{\mathrm{g}}+A_{\mathrm{h}}+A_{\mathrm{t}}$ of ${V_{\mathrm{g}}}\times{V_{\mathrm{h}}}\times{V_{\mathrm{t}}}$ onto $C^{1}(n)$ (see (H1));
hence it is a Borel set.
By Lemma~2.E, every element of ${V_{\mathrm{g}}}$ is, in exactly one way, the sum of an element of $\mathrm{Fg}$ and an element of $r_q\,dC^{0}(n;\mathbb{Z})$:
these are given by the fractional and the integer parts of its coordinates in the basis $\left(r_q\,d\mathbf{1}_c\right)_{c\ne0}$.
The map $\theta\mapsto H_\theta$ is a linear bijection of $\mathbb{R}^3$ onto ${V_{\mathrm{h}}}$;
hence every element of ${V_{\mathrm{h}}}$ is, in exactly one way, the sum of an element of $\mathrm{Fh}$ and of some $H_{2\pi\mathrm{kz}}$ with $\mathrm{kz}\in\mathbb{Z}^3$.
Let $A=A_{\mathrm{g}}+A_{\mathrm{h}}+A_{\mathrm{t}}\in C^{1}(n)$, and let $\mathrm{mc}=d\boldsymbol{\lambda}'+\sum_\mu\mathrm{kz}_\mu\cdot m^{(\mu)}\in\mathrm{Zcl}$ (Lemma~2.C\,(ii)).
By Lemma~2.C\,(iii), the three components of $A-r_q\,\mathrm{mc}$ are
\[
A_{\mathrm{g}}-\sum_\mu\mathrm{kz}_\mu\cdot d\boldsymbol{\lambda}^{(\mu)}-r_q\,d\boldsymbol{\lambda}'\in{V_{\mathrm{g}}},
\qquad
A_{\mathrm{h}}-H_{2\pi\mathrm{kz}}\in{V_{\mathrm{h}}},
\qquad
A_{\mathrm{t}}\in{V_{\mathrm{t}}}.
\]
There is exactly one $\mathrm{kz}\in\mathbb{Z}^3$ for which the second component lies in $\mathrm{Fh}$,
and for this $\mathrm{kz}$ there is exactly one $d\boldsymbol{\lambda}'\in dC^{0}(n;\mathbb{Z})$ for which the first component lies in $\mathrm{Fg}$.
By Lemma~2.C\,(ii), the pairs $(\mathrm{kz},d\boldsymbol{\lambda}')$ correspond bijectively to the elements $\mathrm{mc}$ of $\mathrm{Zcl}$.
Hence $\mathrm{Db}$ is a fundamental domain of $\mathrm{Gq}$.

\emph{Step 4 (volumes).}
Let $\mathrm{bv}_1,\dots,\mathrm{bv}_{\mathrm{Dq}}$ be a basis of a Euclidean space $\mathrm{Vq}$.
In the standard basis of $\mathbb{R}^{\mathrm{Dq}}$ and an orthonormal basis of $\mathrm{Vq}$, the linear bijection $\mathrm{Tl}\colon\mathrm{tq}\mapsto\sum_j\mathrm{tq}_j\cdot\mathrm{bv}_j$ of $\mathbb{R}^{\mathrm{Dq}}$ onto $\mathrm{Vq}$ has a matrix $\mathrm{Vm}$ with $\mathrm{Vm}^{T}\mathrm{Vm}=\left(\langle\mathrm{bv}_i,\mathrm{bv}_j\rangle\right)_{i,j}$.
For the indicator function of the Borel set
$\mathrm{Tl}\left([0,1)^{\mathrm{Dq}}\right)=\left\{\sum_j\mathrm{tq}_j\cdot\mathrm{bv}_j:\ 0\le\mathrm{tq}_j<1\right\}$,
the identity (Lin) reads $1=|\det\mathrm{Vm}|^{-1}\operatorname{vol}\left(\mathrm{Tl}\left([0,1)^{\mathrm{Dq}}\right)\right)$.
So this set has the volume $|\det\mathrm{Vm}|=\det\left(\langle\mathrm{bv}_i,\mathrm{bv}_j\rangle\right)^{\frac12}$.
With $\mathrm{bv}_c=r_q\,d\mathbf{1}_c$, $c\in\mathrm{Bn}\setminus\{0\}$, and Lemma~2.E (so that $\dim{V_{\mathrm{g}}}=|n|-1$) this gives
\begin{equation}
\operatorname{vol}(\mathrm{Fg})=r_q^{|n|-1}\left(\frac{\operatorname{det}'\Delta_0}{|n|}\right)^{\frac12}.
\tag{38.10}
\end{equation}
The cochains $\mathrm{eh}_\mu:=|n|^{-\frac12}\mathbf{1}_\mu$, $\mu=1,2,3$, form an orthonormal basis of ${V_{\mathrm{h}}}$, and
\[
H_\theta=\sum_\mu\frac{\theta_\mu|n|^{\frac12}}{q\,n_\mu}\,\mathrm{eh}_\mu .
\]
So the matrix of the linear bijection $\theta\mapsto H_\theta$ of $\mathbb{R}^3$ onto ${V_{\mathrm{h}}}$ is diagonal, with the determinant
$\prod_\mu\frac{|n|^{\frac12}}{q\,n_\mu}=|n|^{\frac12}q^{-3}$;
moreover $q^{-1}=\frac{r_q}{2\pi}$.
By (Lin) for this map and the function $\mathbf{1}_{\mathrm{Fh}}\cdot\mathrm{Gv}$, every non-negative Borel function $\mathrm{Gv}$ on ${V_{\mathrm{h}}}$ satisfies
\begin{equation}
\int_{\mathrm{Fh}}\mathrm{d}A_{\mathrm{h}}\,\mathrm{Gv}(A_{\mathrm{h}})
=|n|^{\frac12}q^{-3}\int_{[-\pi,\pi)^3}\mathrm{d}\theta\,\mathrm{Gv}(H_\theta)
=r_q^3\,|n|^{\frac12}\int_{[-\pi,\pi)^3}\frac{\mathrm{d}\theta}{(2\pi)^3}\,\mathrm{Gv}(H_\theta).
\tag{38.11}
\end{equation}

\emph{Step 5 (conclusion).}
Lemma~2.B, applied with Steps~2 and~3, gives
\[
\int_{\mathrm{Db}}\mathrm{d}A\,|\mathrm{Ps}(A)|=\int_{\mathrm{Da}}\mathrm{d}A\,|\mathrm{Ps}(A)|<\infty
\qquad\text{and}\qquad
\int_{\mathrm{Db}}\mathrm{d}A\,\mathrm{Ps}(A)=\int_{\mathrm{Da}}\mathrm{d}A\,\mathrm{Ps}(A).
\]
By 0.3, $\mathrm{d}A$ is the product of the Lebesgue measures of ${V_{\mathrm{g}}}$, ${V_{\mathrm{h}}}$ and ${V_{\mathrm{t}}}$.
For $A_{\mathrm{g}}=d\boldsymbol{\lambda}\in{V_{\mathrm{g}}}$ with real $\boldsymbol{\lambda}$, for $A_{\mathrm{h}}\in{V_{\mathrm{h}}}$ and for $A_{\mathrm{t}}\in{V_{\mathrm{t}}}$,
one has $d(A_{\mathrm{g}}+A_{\mathrm{h}}+A_{\mathrm{t}})=dA_{\mathrm{t}}$ by (H1), and $\mathrm{wq}(A_{\mathrm{g}}+A_{\mathrm{h}}+A_{\mathrm{t}})=\mathrm{wq}(A_{\mathrm{h}}+A_{\mathrm{t}})$ by the first invariance of $\mathrm{wq}$.
Hence
\[
\mathrm{Ps}(A_{\mathrm{g}}+A_{\mathrm{h}}+A_{\mathrm{t}})=\Phi(dA_{\mathrm{t}})\cdot\exp\left(-\frac{1}{2}\|dA_{\mathrm{t}}\|^2\right)\cdot\mathrm{wq}(A_{\mathrm{t}}+A_{\mathrm{h}})
\]
does not depend on $A_{\mathrm{g}}$.
Let $\mathrm{Gv}\ge0$ be a Borel function on $C^{1}(n)$ with $\mathrm{Gv}(A_{\mathrm{g}}+A_{\mathrm{h}}+A_{\mathrm{t}})=\mathrm{Gv}(A_{\mathrm{h}}+A_{\mathrm{t}})$.
By Tonelli's theorem, (38.10) and (38.11),
\[
\begin{aligned}
\int_{\mathrm{Db}}\mathrm{d}A\,\mathrm{Gv}(A)
&=\operatorname{vol}(\mathrm{Fg})\int_{\mathrm{Fh}}\mathrm{d}A_{\mathrm{h}}\int_{{V_{\mathrm{t}}}}\mathrm{d}\sigma(A_{\mathrm{t}})\,\mathrm{Gv}(A_{\mathrm{h}}+A_{\mathrm{t}})\\
&=r_q^{|n|-1}\left(\frac{\operatorname{det}'\Delta_0}{|n|}\right)^{\frac12}r_q^3\,|n|^{\frac12}\int_{[-\pi,\pi)^3}\frac{\mathrm{d}\theta}{(2\pi)^3}\int_{{V_{\mathrm{t}}}}\mathrm{d}\sigma(A)\,\mathrm{Gv}(A+H_\theta).
\end{aligned}
\]
The constant in front of the integrals is $r_q^{|n|+2}\left(\operatorname{det}'\Delta_0\right)^{\frac12}=\mathrm{Nc}_n$.
For $\mathrm{Gv}=|\mathrm{Ps}|$ the left side is finite; therefore
\[
\int_{[-\pi,\pi)^3}\frac{\mathrm{d}\theta}{(2\pi)^3}\int_{{V_{\mathrm{t}}}}\mathrm{d}\sigma(A)\,|\Phi(dA)|\cdot\exp\left(-\frac{1}{2}\|dA\|^2\right)\cdot|\mathrm{wq}(A+H_\theta)|
=\mathrm{Nc}_n^{-1}\int_{\mathrm{Db}}\mathrm{d}A\,|\mathrm{Ps}(A)|<\infty,
\]
that is, the right side of (a) converges absolutely.
The positive and the negative parts of $\operatorname{Re}\mathrm{Ps}$ and of $\operatorname{Im}\mathrm{Ps}$ are non-negative Borel functions which do not depend on $A_{\mathrm{g}}$ and are bounded by $|\mathrm{Ps}|$.
For each of these four functions the identity just proved for $\mathrm{Gv}$ has finite numbers on both sides, and by linearity
\[
\int_{\mathrm{Db}}\mathrm{d}A\,\mathrm{Ps}(A)
=\mathrm{Nc}_n\int_{[-\pi,\pi)^3}\frac{\mathrm{d}\theta}{(2\pi)^3}\int_{{V_{\mathrm{t}}}}\mathrm{d}\sigma(A)\,\Phi(dA)\cdot\exp\left(-\frac{1}{2}\|dA\|^2\right)\cdot\mathrm{wq}(A+H_\theta).
\]
The set $[-\pi,\pi]^3\setminus[-\pi,\pi)^3$ is a null set.
Together with (38.9) this is (a).
Of the function $\mathrm{wq}$ only the following properties have entered: it is a bounded Borel function with complex values, and it has the two invariances.

\paragraph{2.7. Proof of (b).}

\begin{named}{Lemma 2.F}[The flux is a function of the curvature]
For $\mu<\nu$ let $\mathfrak{l}$ be the third direction, and let
\[
\Pi_{\mu\nu}:=\left\{(c;\mu\nu):\ c\in\mathrm{Bn},\ c_{\mathfrak{l}}=0\right\}.
\]
\begin{itemize}
\item[(i)]
$\sum_{p\in\Pi_{\mu\nu}}(dA)_p=0$ for every $A\in C^{1}(n)\otimes\mathbb{C}$.
Consequently, for $m\in C^{1}(n;\mathbb{Z})$ the cochain $dm$, which is closed by (H1), has the flux $0$.
The flux is additive in $\mathfrak{z}$.
\item[(ii)]
If $(A,\mathfrak{z})$ is a configuration (Section~I.7.1) and $\mathfrak{z}$ has the flux $w$, then $\sum_{p\in\Pi_{\mu\nu}}F^{\mathrm{u}}_p=r_q w_{\mu\nu}$.
\item[(iii)]
Let
\[
\mathrm{If}_w(u):=\prod_{\mu<\nu}\mathbf{1}\left[\sum_{p\in\Pi_{\mu\nu}}u_p=r_q w_{\mu\nu}\right];
\]
this is a bounded Borel function on $C^{2}(n)$.
Then $\mathbf{1}[\text{flux}=w]=\mathrm{If}_w(F^{\mathrm{u}})$; for $\Phi\in\mathfrak{F}(n)$ also $\mathrm{If}_w\cdot\Phi\in\mathfrak{F}(n)$, and
\begin{equation}
Z^{(w)}_n[\Phi]=Z_n\left[\mathrm{If}_w\cdot\Phi\right]
=\sum_{\mathfrak{z}\in\mathrm{Mz}_w}\ \int_{\mathrm{Bo}}\mathrm{d}A\,\exp\left(-\frac{1}{2}\|F^{\mathrm{u}}\|^2\right)\cdot\det\nolimits_n\left(D_A+\boldsymbol{\mu}\right)\cdot\Phi(F^{\mathrm{u}}),
\qquad F^{\mathrm{u}}=dA+r_q\mathfrak{z}.
\tag{38.12}
\end{equation}
\end{itemize}
\end{named}

\begin{proof}
(i)
One has $(dA)_{\mu\nu}(c)=\left[A_\mu(c)-A_\mu(c+e_\nu)\right]+\left[A_\nu(c+e_\mu)-A_\nu(c)\right]$.
The maps $c\mapsto\mathrm{red}(c+e_\nu)$ and $c\mapsto\mathrm{red}(c+e_\mu)$ are bijections of the set $\{c\in\mathrm{Bn}:\ c_{\mathfrak{l}}=0\}$, and $A$ is periodic;
so each of the two brackets has the sum $0$ over this set.
(ii)
This follows from (i) and from the definition of $w_{\mu\nu}$ in Section~I.7.1; by (H2) the choice of the plane is immaterial.
(iii)
By (ii), $\mathrm{If}_w(F^{\mathrm{u}})=1$ if $\mathfrak{z}$ has the flux $w$, and $\mathrm{If}_w(F^{\mathrm{u}})=0$ if it has another flux.
The remaining assertions follow from (I.7.1) and from the definition of $Z^{(w)}_n$ in Section~I.7.3.
\end{proof}

\emph{The first assertion of (b), with the convergence.}
Let $\boldsymbol{\mu}$ be real, and put $\mathrm{wq}(A):=\det\nolimits_n\left(D_A+\boldsymbol{\mu}\right)$ for $A\in C^{1}(n)$.
This function is continuous by Proposition~38.1\,(a): it is a polynomial in the numbers $\mathrm{e}^{\pm iqA(b)}$, which are continuous functions of $A$.
Moreover $|\mathrm{wq}(A)|\le\mathrm{Kd}(0)$ by Proposition~38.1\,(d) with $\operatorname{Im}B=0$;
$\mathrm{wq}(A+d\boldsymbol{\lambda})=\mathrm{wq}(A)$ for real $\boldsymbol{\lambda}$ by Proposition~38.1\,(c);
and $\mathrm{wq}(A+r_q m)=\mathrm{wq}(A)$ for $m\in C^{1}(n;\mathbb{Z})$ by Proposition~38.1\,(b).
By (38.12) with $w=0$, the left side of (a) is $Z^{(0)}_n[\Phi]$.
Hence (a), applied to this $\mathrm{wq}$, is the first assertion of (b), together with the absolute convergence of both sides.

\emph{The law.}
For $A\in{V_{\mathrm{t}}}$ let
\[
\mathrm{wd}(A):=\exp\left(-\frac{1}{2}\|dA\|^2\right)\int_{[-\pi,\pi]^3}\mathrm{d}\theta\,\det\nolimits_n\left(D_{A+H_\theta}+\boldsymbol{\mu}\right).
\]
The connection $A+H_\theta$ is real, so that the determinant is non-negative by Proposition~50\,(a), and $\mathrm{wd}$ is a non-negative Borel function by Tonelli's theorem.
The first assertion with $\Phi=1$ gives
\[
\int_{{V_{\mathrm{t}}}}\mathrm{d}\sigma\,\mathrm{wd}=\frac{(2\pi)^3\,Z^{(0)}_n[1]}{\mathrm{Nc}_n},
\]
which is a number in $(0,\infty)$ by Proposition~50\,(e).
For every bounded Borel function $\Phi$ on $C^{2}(n)$, the first assertion, Fubini's theorem and the definition of $E^{\mathrm{av}}_n$ (Section~I.7.3) give
\[
E^{\mathrm{av}}_n\,\Phi(F^{\mathrm{u}})=\frac{Z^{(0)}_n[\Phi]}{Z^{(0)}_n[1]}
=\frac{\int_{{V_{\mathrm{t}}}}\mathrm{d}\sigma(A)\,\Phi(dA)\cdot\mathrm{wd}(A)}{\int_{{V_{\mathrm{t}}}}\mathrm{d}\sigma\,\mathrm{wd}}.
\]
The right side is the integral of $\Phi$ with respect to the law of $dA$ under the probability measure $\left(\int_{{V_{\mathrm{t}}}}\mathrm{d}\sigma\,\mathrm{wd}\right)^{-1}\mathrm{wd}\cdot\sigma$ on ${V_{\mathrm{t}}}$.

\paragraph{2.8. Proof of (c).}
The torus is now cubic: $n_1=n_2=n_3=n\ge4$, and $|n|=n^3$.
By (H1), $F_N=dC^{1}(n)=d{V_{\mathrm{t}}}$, and $d_{\mathrm{t}}:=d|_{{V_{\mathrm{t}}}}$ is a linear bijection of ${V_{\mathrm{t}}}$ onto $F_N$, its kernel being ${V_{\mathrm{t}}}\cap\left({V_{\mathrm{g}}}+{V_{\mathrm{h}}}\right)=\{0\}$.
Hence $A(\,\cdot\,)=d_{\mathrm{t}}^{-1}$, and $A(dA)=A$ for $A\in{V_{\mathrm{t}}}$.
With respect to the Lebesgue measure $\mathrm{d}u$ of the Euclidean space $\left(F_N,\langle\,\cdot\,,\,\cdot\,\rangle\right)$, the measure $\gamma_N$ has the density $(2\pi)^{-\frac{\mathrm{Dq}}{2}}\exp\left(-\frac12\|u\|^2\right)$,
where $\mathrm{Dq}:=\dim F_N=\dim{V_{\mathrm{t}}}=2n^3-2$.

\begin{named}{Lemma 2.G}[Change of variables; the identity of Proposition~38.2\,(c) with $\mathrm{zV}_n$ and $\gamma_N$]
Let $G$ be a Borel function on $F_N$ which is either non-negative, or complex valued with $|G(u)|\le\mathrm{cf}\cdot\exp\left(\mathrm{cf}\cdot\|u\|\right)$ for some $\mathrm{cf}<\infty$.
Then
\[
\int_{{V_{\mathrm{t}}}}\mathrm{d}\sigma(A)\,G(dA)\cdot\exp\left(-\frac{1}{2}\|dA\|^2\right)=\mathrm{zV}_n\int\mathrm{d}\gamma_N\,G .
\]
\end{named}

\begin{proof}
By (Lin) for $\mathrm{Tl}=d_{\mathrm{t}}\colon{V_{\mathrm{t}}}\to F_N$, every non-negative Borel function $\mathrm{Gv}$ on $F_N$ satisfies
\[
\int_{{V_{\mathrm{t}}}}\mathrm{d}\sigma(A)\,\mathrm{Gv}(d_{\mathrm{t}}A)=|\det d_{\mathrm{t}}|^{-1}\int_{F_N}\mathrm{d}u\,\mathrm{Gv}(u).
\]
For $\mathrm{Gv}(u)=G(u)\exp\left(-\frac12\|u\|^2\right)$ with $G\ge0$, the left side of the Lemma thus equals $|\det d_{\mathrm{t}}|^{-1}(2\pi)^{\frac{\mathrm{Dq}}{2}}\int\mathrm{d}\gamma_N\,G$;
the choice $G=1$ shows that the constant is $\mathrm{zV}_n$.
Let now $G$ be complex valued with the stated bound.
By the computation that proves (38.7), $|G(u)|\exp\left(-\frac12\|u\|^2\right)\le\mathrm{cf}\cdot\exp\left(\mathrm{cf}^2\right)\cdot\exp\left(-\frac14\|u\|^2\right)$,
and $\int_{F_N}\mathrm{d}u\,\exp\left(-\frac14\|u\|^2\right)=(4\pi)^{\frac{\mathrm{Dq}}{2}}$ by the formula in the proof of Lemma~2.D\,(iii), with $\mathrm{Sy}$ equal to $\frac12$ times the identity.
Hence $\int\mathrm{d}\gamma_N\,|G|<\infty$, and the case $G\ge0$, applied to $|G|$, shows that both sides converge absolutely.
The identity for $G$ follows from the case $G\ge0$, applied to the positive and the negative parts of $\operatorname{Re}G$ and of $\operatorname{Im}G$.
\end{proof}

\emph{The numbers $\mathrm{zV}_n$ and $Z^{\mathrm{ph}}_n$.}
The value $\mathrm{zV}_n=\frac{(2\pi)^{n^3-1}}{\operatorname{det}'\Delta_0}$ is Lemma~2.D\,(iii) for $|n|=n^3$, and the change of variables asserted in (c) is Lemma~2.G.
By the definition of $Z^{\mathrm{ph}}_n$ and by Lemma~2.D\,(iii) and~(i),
\[
Z^{\mathrm{ph}}_n=\mathrm{Nc}_n\cdot\mathrm{zV}_n
=r_q^{n^3+2}\left(\operatorname{det}'\Delta_0\right)^{\frac12}(2\pi)^{n^3-1}\left(\operatorname{det}'\Delta_0\right)^{-1}
=r_q^{n^3+2}(2\pi)^{n^3-1}\exp\left(-\frac{1}{2}\mathrm{Ld}(n)\right),
\]
where
\[
\mathrm{Ld}(n)=\log\operatorname{det}'\Delta_0=\sum_{k\in(\mathbb{Z}/n)^3,\ k\ne0}\log\sum_\mu\left(2-2\cos\frac{2\pi k_\mu}{n}\right).
\]

\emph{The density.}
Let $\boldsymbol{\mu}=\bar{m}\eta+m_{\mathrm{in}}$.
Here $\bar{m}>0$ (Section~I.0.5) and $\eta=L^{-N}>0$.
Proposition~50\,(a) with the lattice mass $\bar{m}\eta>0$ and $B=0$ gives
\[
\det\nolimits_n\left(D_0+\bar{m}\eta\right)\ge\left(\bar{m}\eta\right)^{192n^3}>0.
\]
By (I.7.2),
\[
\mathrm{dens}^{\theta,0}_N(F^{\mathrm{u}})=\exp\left(-\varepsilon^{\mathrm{v}}_{\mathrm{in}}n^3\right)\frac{\det\nolimits_n\left(D_{A(F^{\mathrm{u}})+H_\theta}+\boldsymbol{\mu}\right)}{\det\nolimits_n\left(D_0+\bar{m}\eta\right)}.
\]
As a function of $(\theta,F^{\mathrm{u}})\in\mathbb{R}^3\times F_N$ it is continuous (Proposition~38.1\,(a); the maps $A(\,\cdot\,)$ and $\theta\mapsto H_\theta$ are linear),
non-negative (Proposition~50\,(a)), and bounded by $\exp\left(-\varepsilon^{\mathrm{v}}_{\mathrm{in}}n^3\right)\mathrm{Kd}(0)/\det\nolimits_n\left(D_0+\bar{m}\eta\right)$ (Proposition~38.1\,(d)).

\emph{The identities.}
Let $G$ be a Borel function on $C^{2}(n)$ with $|G(u)|\le\mathrm{cf}\cdot\exp\left(\mathrm{cf}\cdot\|u\|\right)$.
For every $\theta$, apply Lemma~2.G to the function $u\mapsto G(u)\cdot\det\nolimits_n\left(D_{A(u)+H_\theta}+\boldsymbol{\mu}\right)$ on $F_N$.
Its modulus is at most $\mathrm{Kd}(0)\cdot\mathrm{cf}\cdot\exp\left(\mathrm{cf}\cdot\|u\|\right)\le\mathrm{cf}'\cdot\exp\left(\mathrm{cf}'\cdot\|u\|\right)$ with $\mathrm{cf}':=\max\left(1,\mathrm{Kd}(0)\right)\cdot\mathrm{cf}$,
and the integrands are Borel functions of $(\theta,u)$, respectively of $(\theta,A)$, so that the integrals over $\theta$ exist.
The definitions of Section~I.7.3, these applications of Lemma~2.G and part (b) give
\begin{equation}
\begin{aligned}
\bar{\nu}(G)
&=\int_{[-\pi,\pi]^3}\frac{\mathrm{d}\theta}{(2\pi)^3}\int\mathrm{d}\gamma_N\,G\cdot\mathrm{dens}^{\theta,0}_N\\
&=\frac{\exp\left(-\varepsilon^{\mathrm{v}}_{\mathrm{in}}n^3\right)}{\det\nolimits_n\left(D_0+\bar{m}\eta\right)\cdot\mathrm{zV}_n}
\int_{[-\pi,\pi]^3}\frac{\mathrm{d}\theta}{(2\pi)^3}\int_{{V_{\mathrm{t}}}}\mathrm{d}\sigma(A)\,G(dA)\cdot\exp\left(-\frac{1}{2}\|dA\|^2\right)\cdot\det\nolimits_n\left(D_{A+H_\theta}+\boldsymbol{\mu}\right)\\
&=\frac{\exp\left(-\varepsilon^{\mathrm{v}}_{\mathrm{in}}n^3\right)}{\det\nolimits_n\left(D_0+\bar{m}\eta\right)\cdot\mathrm{zV}_n\cdot\mathrm{Nc}_n}\,Z^{(0)}_n[G].
\end{aligned}
\tag{38.13}
\end{equation}
The definition of $\mathrm{zr}_n$ in Section~I.7.3, Lemma~2.G for every $\theta$, applied to the function $u\mapsto\det\nolimits_n\left(D_{A(u)+H_\theta}+\boldsymbol{\mu}\right)$ on $F_N$ (which is continuous and of modulus at most $\mathrm{Kd}(0)$), and part (b) with $\Phi=1$ give
\[
\mathrm{zr}_n(q,\boldsymbol{\mu})
=\mathrm{zV}_n^{-1}\int_{[-\pi,\pi]^3}\frac{\mathrm{d}\theta}{(2\pi)^3}\int_{{V_{\mathrm{t}}}}\mathrm{d}\sigma(A)\,\exp\left(-\frac{1}{2}\|dA\|^2\right)\cdot\det\nolimits_n\left(D_{A+H_\theta}+\boldsymbol{\mu}\right)
=\frac{Z^{(0)}_n[1]}{\mathrm{Nc}_n\cdot\mathrm{zV}_n}.
\]
So $Z^{(0)}_n[1]=Z^{\mathrm{ph}}_n\cdot\mathrm{zr}_n(q,\boldsymbol{\mu})$.
Since $Z^{(0)}_n[1]>0$ (Proposition~50\,(e)) and $Z^{\mathrm{ph}}_n>0$, it follows that $\mathrm{zr}_n(q,\boldsymbol{\mu})>0$;
and $\mathrm{zr}_n(q,\boldsymbol{\mu})\le\mathrm{Kd}(0)$, because $\gamma_N$ and $\frac{\mathrm{d}\theta}{(2\pi)^3}$ on $[-\pi,\pi]^3$ are probability measures.
For $G=1$, (38.13) gives
\[
\bar{\nu}(1)=\frac{\exp\left(-\varepsilon^{\mathrm{v}}_{\mathrm{in}}n^3\right)Z^{(0)}_n[1]}{\det\nolimits_n\left(D_0+\bar{m}\eta\right)\cdot Z^{\mathrm{ph}}_n}
=\frac{\exp\left(-\varepsilon^{\mathrm{v}}_{\mathrm{in}}n^3\right)\mathrm{zr}_n(q,\boldsymbol{\mu})}{\det\nolimits_n\left(D_0+\bar{m}\eta\right)},
\]
a number in $(0,\infty)$.
Division of (38.13) by this number gives
\[
\bar{\mu}(G)=\frac{\bar{\nu}(G)}{\bar{\nu}(1)}=\frac{Z^{(0)}_n[G]}{Z^{(0)}_n[1]}=E^{\mathrm{av}}_n\,G(F^{\mathrm{u}})
\]
for every bounded Borel function $G$ on $C^{2}(n)$.
So $\bar{\mu}$, a measure on $F_N$, is the law of $F^{\mathrm{u}}$ under $E^{\mathrm{av}}_n$.

\emph{The last sentence of (c).}
In 2.1--2.8 the number $q>0$ is arbitrary, and no relation between $q$ and $\eta$ has been assumed.
For a real number $\boldsymbol{\mu}$ put $m_{\mathrm{in}}:=\boldsymbol{\mu}-\bar{m}\eta$;
it is real, and $\boldsymbol{\mu}=\bar{m}\eta+m_{\mathrm{in}}$, so that every real $\boldsymbol{\mu}$ occurs.
Apart from $\rho$ and the weights, the formulas that define the following numbers contain only $(n,q,\boldsymbol{\mu})$:
\begin{itemize}
\item
$Z^{(0)}_n[1]$ ((I.7.1) and Section~I.7.3: it contains $n$, $r_q$, and $\det\nolimits_n\left(D_A+\boldsymbol{\mu}\right)$, in whose phases the charge $q$ enters);
\item
$\mathrm{Nc}_n$ (it contains $n$ and $r_q$);
\item
$\mathrm{zV}_n$ (it contains $n$);
\item
$Z^{\mathrm{ph}}_n=\mathrm{Nc}_n\cdot\mathrm{zV}_n$;
\item
$\mathrm{zr}_n(q,\boldsymbol{\mu})$ (Section~I.7.3: $F_N=dC^{1}(n)$, $\gamma_N$ and $A(\,\cdot\,)$ contain only $n$, $H_\theta$ contains $n$ and $q$, and the determinant is $\det\nolimits_n\left(D_{A(F^{\mathrm{u}})+H_\theta}+\boldsymbol{\mu}\right)$).
\end{itemize}
Hence $\bar{m}$, $\eta$ and $m_{\mathrm{in}}$ enter the assertions of (c) on these numbers only through $\boldsymbol{\mu}$, and $\varepsilon^{\mathrm{v}}_{\mathrm{in}}$ does not enter them.

\paragraph{2.9. Range.}
Parts (a) and (b) hold for all sides $n_\mu\ge4$, even or odd, for all $q>0$ and for all real $\boldsymbol{\mu}$, with the single constant $\mathrm{Nc}_n$, which depends only on $q$ and $n$.
No assumption has been made on the sign or the size of $\boldsymbol{\mu}$ (compare (I.K10)), on the parity of a side, or on a sign of $\mathrm{wq}$.
Part (c) holds for every cubic torus with $n\ge4$, every real pair $(m_{\mathrm{in}},\varepsilon^{\mathrm{v}}_{\mathrm{in}})$, every $\bar{m}\in(0,1]$ and every $\eta=L^{-N}$.
\end{proof}

\subsection*{3. Proposition~38.3. The shift identity}

This section contains the proof of Proposition~38.3.
Throughout, $n_\mu\ge4$.
The notation is that of Section~0, in particular the entire functions of Section~0.5 with (38.0a), (38.0b) and (ef1)--(ef3), the complex derivatives being written $\partial_a^{\mathrm{kx}}=\prod_j\left(\partial/\partial a_j\right)^{\mathrm{kx}_j}$;
from Section~2 we keep $\mathrm{Bo}$, $\mathrm{Mz}_w$, $\mathrm{Kd}$, $\mathrm{Cth}$, $\mathrm{zV}_n$, the two facts (H1), (H2) on cochains \cite[Section~10.1.1]{QEDI}, and Lemmas~2.A--2.G.

\subsubsection*{3.1. Integrals that depend on complex parameters}

\begin{named}{Lemma 3.A}
Let $(\mathrm{Yq},\mathrm{vq})$ be a $\sigma$-finite measure space and $\mathrm{Nq}\ge1$ an integer.
Let $\mathrm{vph}\colon\mathrm{Yq}\times\mathbb{C}^{\mathrm{Nq}}\to\mathbb{C}$ be measurable with respect to the product of the $\sigma$-algebra of $\mathrm{Yq}$ and the Borel $\sigma$-algebra of $\mathbb{C}^{\mathrm{Nq}}$, and assume:
\begin{itemize}
\item[(i)] for every $\mathrm{yq}\in\mathrm{Yq}$ the function $a\mapsto\mathrm{vph}(\mathrm{yq},a)$ is entire;
\item[(ii)] for every $\mathrm{rd}\in(0,\infty)$ there is a $\mathrm{vq}$-integrable function $\mathrm{Mj}_{\mathrm{rd}}\ge0$ on $\mathrm{Yq}$ such that $|\mathrm{vph}(\mathrm{yq},a)|\le\mathrm{Mj}_{\mathrm{rd}}(\mathrm{yq})$ for every $\mathrm{yq}\in\mathrm{Yq}$ and every $a$ with $\max_j|a_j|\le\mathrm{rd}$.
\end{itemize}
Then the integral
\[
\mathrm{Vph}(a):=\int\mathrm{d}\mathrm{vq}(\mathrm{yq})\,\mathrm{vph}(\mathrm{yq},a)
\]
converges absolutely, and $\int\mathrm{d}\mathrm{vq}(\mathrm{yq})\,|\mathrm{vph}(\mathrm{yq},a)|\le\int\mathrm{d}\mathrm{vq}\,\mathrm{Mj}_{\mathrm{rd}}$ whenever $\max_j|a_j|\le\mathrm{rd}$.
The function $\mathrm{Vph}$ is entire.
For every multi-index $\mathrm{kx}$ and every $a\in\mathbb{C}^{\mathrm{Nq}}$,
\[
\partial_a^{\mathrm{kx}}\mathrm{Vph}(a)=\int\mathrm{d}\mathrm{vq}(\mathrm{yq})\,\partial_a^{\mathrm{kx}}\mathrm{vph}(\mathrm{yq},a),
\]
where the integral on the right converges absolutely.
\end{named}

\begin{proof}
The first assertion is a restatement of (ii).
For $\mathrm{yq}\in\mathrm{Yq}$ let $\mathrm{tay}_{\mathrm{kx}}(\mathrm{yq})$ denote the coefficients of the entire function $\mathrm{vph}(\mathrm{yq},\,\cdot\,)$ (Section~0.5).
Fix $\mathrm{rd}\in(0,\infty)$, and for $\mathrm{ang}\in[0,2\pi]^{\mathrm{Nq}}$ put
\[
\mathrm{zt}(\mathrm{ang}):=\left(\mathrm{rd}\cdot\exp(i\,\mathrm{ang}_1),\dots,\mathrm{rd}\cdot\exp(i\,\mathrm{ang}_{\mathrm{Nq}})\right),\qquad
\mathrm{kx}\cdot\mathrm{ang}:=\sum_j\mathrm{kx}_j\,\mathrm{ang}_j .
\]
Then
\[
\mathrm{vph}\left(\mathrm{yq},\mathrm{zt}(\mathrm{ang})\right)=\sum_{\mathrm{kx}'}\mathrm{tay}_{\mathrm{kx}'}(\mathrm{yq})\cdot\mathrm{rd}^{|\mathrm{kx}'|}\cdot\exp\left(i\,\mathrm{kx}'\cdot\mathrm{ang}\right),
\qquad
\sum_{\mathrm{kx}'}\left|\mathrm{tay}_{\mathrm{kx}'}(\mathrm{yq})\right|\cdot\mathrm{rd}^{|\mathrm{kx}'|}<\infty,
\]
the latter by (38.0a).
Consequently the series, multiplied by $\exp(-i\,\mathrm{kx}\cdot\mathrm{ang})$, may be integrated over $[0,2\pi]^{\mathrm{Nq}}$ term by term (Fubini's theorem for the product of the counting measure and $\mathrm{d}\,\mathrm{ang}$).
Since
\[
\int_{[0,2\pi]^{\mathrm{Nq}}}\mathrm{d}\,\mathrm{ang}\,\exp\left(i(\mathrm{kx}'-\mathrm{kx})\cdot\mathrm{ang}\right)=(2\pi)^{\mathrm{Nq}}\cdot\mathbf{1}[\mathrm{kx}'=\mathrm{kx}],
\]
one obtains
\begin{equation}\tag{38.14}
\mathrm{tay}_{\mathrm{kx}}(\mathrm{yq})=(2\pi)^{-\mathrm{Nq}}\cdot\mathrm{rd}^{-|\mathrm{kx}|}\int_{[0,2\pi]^{\mathrm{Nq}}}\mathrm{d}\,\mathrm{ang}\;\mathrm{vph}\left(\mathrm{yq},\mathrm{zt}(\mathrm{ang})\right)\cdot\exp\left(-i\,\mathrm{kx}\cdot\mathrm{ang}\right).
\end{equation}
The integrand is a measurable function of $(\mathrm{yq},\mathrm{ang})$, because the map $(\mathrm{yq},\mathrm{ang})\mapsto\left(\mathrm{yq},\mathrm{zt}(\mathrm{ang})\right)$ is measurable.
By (ii) its modulus does not exceed $\mathrm{Mj}_{\mathrm{rd}}(\mathrm{yq})$, a function that is integrable for the product of $\mathrm{vq}$ and $\mathrm{d}\,\mathrm{ang}$.
Fubini's theorem shows that $\mathrm{yq}\mapsto\mathrm{tay}_{\mathrm{kx}}(\mathrm{yq})$ is measurable, and (38.14) together with (ii) yields
\begin{equation}\tag{38.15}
\left|\mathrm{tay}_{\mathrm{kx}}(\mathrm{yq})\right|\le\mathrm{rd}^{-|\mathrm{kx}|}\cdot\mathrm{Mj}_{\mathrm{rd}}(\mathrm{yq})\qquad\text{for all }\mathrm{yq}\in\mathrm{Yq},\text{ all }\mathrm{kx}\text{ and every }\mathrm{rd}\in(0,\infty).
\end{equation}
Now let $\max_j|a_j|\le\mathrm{rd}$.
Applying (38.15) with $2\,\mathrm{rd}$ in place of $\mathrm{rd}$, and using
$\sum_{\mathrm{kx}}2^{-|\mathrm{kx}|}=\left(\sum_{i\ge0}2^{-i}\right)^{\mathrm{Nq}}=2^{\mathrm{Nq}}$, one finds
\[
\sum_{\mathrm{kx}}\int\mathrm{d}\mathrm{vq}(\mathrm{yq})\,\left|\mathrm{tay}_{\mathrm{kx}}(\mathrm{yq})\right|\cdot\left|a^{\mathrm{kx}}\right|
\le\sum_{\mathrm{kx}}(2\,\mathrm{rd})^{-|\mathrm{kx}|}\cdot\mathrm{rd}^{|\mathrm{kx}|}\int\mathrm{d}\mathrm{vq}\,\mathrm{Mj}_{2\mathrm{rd}}
=2^{\mathrm{Nq}}\int\mathrm{d}\mathrm{vq}\,\mathrm{Mj}_{2\mathrm{rd}}<\infty .
\]
Fubini's theorem for the product of $\mathrm{vq}$ and the counting measure on the set of the multi-indices therefore gives
\[
\mathrm{Vph}(a)=\int\mathrm{d}\mathrm{vq}(\mathrm{yq})\sum_{\mathrm{kx}}\mathrm{tay}_{\mathrm{kx}}(\mathrm{yq})\cdot a^{\mathrm{kx}}
=\sum_{\mathrm{kx}}\left(\int\mathrm{d}\mathrm{vq}(\mathrm{yq})\,\mathrm{tay}_{\mathrm{kx}}(\mathrm{yq})\right)\cdot a^{\mathrm{kx}},
\]
and the series on the right converges absolutely.
As $\mathrm{rd}$ is arbitrary, $\mathrm{Vph}$ is entire, its coefficients being the numbers $\int\mathrm{d}\mathrm{vq}(\mathrm{yq})\,\mathrm{tay}_{\mathrm{kx}}(\mathrm{yq})$.
By (38.0b), applied to $\mathrm{Vph}$ and to $\mathrm{vph}(\mathrm{yq},\,\cdot\,)$,
\[
\partial_a^{\mathrm{kx}}\mathrm{Vph}(0)=\mathrm{kx}!\int\mathrm{d}\mathrm{vq}(\mathrm{yq})\,\mathrm{tay}_{\mathrm{kx}}(\mathrm{yq})=\int\mathrm{d}\mathrm{vq}(\mathrm{yq})\,\partial_a^{\mathrm{kx}}\mathrm{vph}(\mathrm{yq},0),
\]
and the last integral converges absolutely by (38.15).
For the derivatives at an arbitrary point $a^0\in\mathbb{C}^{\mathrm{Nq}}$, apply what has been shown to the function $(\mathrm{yq},a)\mapsto\mathrm{vph}(\mathrm{yq},a^0+a)$.
It is measurable; it satisfies (i) by (ef2), and (ii) with $\mathrm{Mj}_{\mathrm{rd}+\max_j|a^0_j|}$ in place of $\mathrm{Mj}_{\mathrm{rd}}$; its integral is $\mathrm{Vph}(a^0+a)$; and the derivatives at $a=0$ of $a\mapsto\mathrm{vph}(\mathrm{yq},a^0+a)$ and of $a\mapsto\mathrm{Vph}(a^0+a)$ are the derivatives of $\mathrm{vph}(\mathrm{yq},\,\cdot\,)$ and of $\mathrm{Vph}$ at $a^0$.
\end{proof}

\subsubsection*{3.2. The identity theorem}

\begin{named}{Lemma 3.B}
Two entire functions on $\mathbb{C}^{\mathrm{Nq}}$ that coincide on $\mathbb{R}^{\mathrm{Nq}}$ are equal.
\end{named}

\begin{proof}
Let $\mathrm{fe}$ be the difference of the two functions.
It is entire by (ef3), and it vanishes on $\mathbb{R}^{\mathrm{Nq}}$.
Whenever an entire function vanishes on $\mathbb{R}^{\mathrm{Nq}}$, each of its first partial derivatives vanishes there as well: by (ef1) the complex derivative with respect to $a_j$ exists, hence at a real point it is the limit of the difference quotients with real increments, and these are $0$.
By (ef1) the partial derivatives are again entire.
Induction on $|\mathrm{kx}|$ shows that every $\partial_a^{\mathrm{kx}}\mathrm{fe}$ vanishes on $\mathbb{R}^{\mathrm{Nq}}$, in particular at $0$.
By (38.0b) all coefficients of $\mathrm{fe}$ are $0$.
\end{proof}

\subsubsection*{3.3. Proof of (a)}

\begin{proof}[Proof of Proposition~38.3\,(a)]
Let $a\in C^{1}(n)\otimes\mathbb{C}$, $A\in C^{1}(n)$ and $\theta\in\mathbb{R}^3$.
Then $\operatorname{Im}(A+a)=\operatorname{Im}(A+H_{\theta}+a)=\operatorname{Im}a$, and Proposition~38.1\,(d) gives
\begin{equation}\tag{38.16}
\left|\operatorname{det}_n\left(D_{A+a}+\boldsymbol{\mu}\right)\right|\le\mathrm{Kd}\left(|\operatorname{Im}a|_{\infty}\right),\qquad
\left|\operatorname{det}_n\left(D_{A+H_{\theta}+a}+\boldsymbol{\mu}\right)\right|\le\mathrm{Kd}\left(|\operatorname{Im}a|_{\infty}\right);
\end{equation}
the function $\mathrm{Kd}$ is non-decreasing.
By Proposition~38.1\,(a), $B\mapsto\operatorname{det}_n(D_B+\boldsymbol{\mu})$ is a polynomial in the numbers $\exp\left(\pm iqB(b)\right)$, and
$\exp\left(\pm iq(A+a)(b)\right)=\exp\left(\pm iqA(b)\right)\cdot\exp\left(\pm iqa(b)\right)$.
For fixed $A$ and $\theta$ the two determinants are therefore finite linear combinations of exponentials of linear forms in $a$, and hence entire functions of $a\in\mathbb{C}^{\mathrm{Ln}(n)}=\mathbb{C}^{3|n|}$ by (ef3);
moreover they are continuous in $(A,a)$, resp.\ in $(\theta,A,a)$.
If $\max_b|a(b)|\le\mathrm{rd}$, then $|\operatorname{Im}a|_{\infty}\le\mathrm{rd}$, and the right sides of (38.16) are at most $\mathrm{Kd}(\mathrm{rd})$.

\emph{The functions $Z^{\mathrm{x},(w)}_n$.}
We apply Lemma~3.A with $\mathrm{Nq}=3|n|$, $\mathrm{Yq}=\mathrm{Mz}_w\times\mathrm{Bo}$, and $\mathrm{vq}$ the product of the counting measure and $\mathrm{d}A$; this measure is $\sigma$-finite, since $\mathrm{Mz}_w$ is countable and $\mathrm{Bo}$ has finite volume.
Put
\[
\mathrm{vph}\left((\mathfrak{z},A),a\right):=\exp\left(-\frac{1}{2}\|dA+r_q\mathfrak{z}\|^2\right)\cdot\operatorname{det}_n\left(D_{A+a}+\boldsymbol{\mu}\right),
\qquad
\mathrm{Mj}_{\mathrm{rd}}(\mathfrak{z},A):=\mathrm{Kd}(\mathrm{rd})\cdot\exp\left(-\frac{1}{2}\|dA+r_q\mathfrak{z}\|^2\right).
\]
For every $\mathfrak{z}$ the function $\mathrm{vph}$ is continuous in $(A,a)$, hence it is measurable.
The sets $\mathrm{Mz}_w$, $w\in\mathbb{Z}^3$, are pairwise disjoint subsets of $C^{2}(n;\mathbb{Z})$; so (38.6) with $\epsilon=\frac{1}{2}$ gives
\begin{equation}\tag{38.17}
\sum_{w\in\mathbb{Z}^3}\ \sum_{\mathfrak{z}\in\mathrm{Mz}_w}\ \int_{\mathrm{Bo}}\mathrm{d}A\;\mathrm{Kd}(\mathrm{rd})\cdot\exp\left(-\frac{1}{2}\|dA+r_q\mathfrak{z}\|^2\right)\le\mathrm{Kd}(\mathrm{rd})\cdot\left(r_q\cdot\mathrm{Cth}\right)^{3|n|}.
\end{equation}
Hence, for every $w$, the sum and the integral that define $Z^{\mathrm{x},(w)}_n(a)$ converge absolutely, with a majorant that is the same for all $a$ in $\left\{\max_b|a(b)|\le\mathrm{rd}\right\}$, and $Z^{\mathrm{x},(w)}_n$ is entire.
For every $a$ with $\max_b|a(b)|\le\mathrm{rd}$ the term of index $w$ of the series in (38.17) is a bound of $\left|Z^{\mathrm{x},(w)}_n(a)\right|$.
Therefore $\sum_w\left|Z^{\mathrm{x},(w)}_n(a)\right|\le\mathrm{Kd}(\mathrm{rd})\cdot\left(r_q\cdot\mathrm{Cth}\right)^{3|n|}$ on this set, and the series converges uniformly there.

\emph{The function $\mathbf{Z}^{\mathrm{x}}_n$.}
We apply Lemma~3.A with $\mathrm{Nq}=3|n|$, $\mathrm{Yq}=[-\pi,\pi]^3\times{V_{\mathrm{t}}}$, $\mathrm{vq}$ the product of $\mathrm{d}\theta/(2\pi)^3$ and $\sigma$, and
\[
\mathrm{vph}\left((\theta,A),a\right):=\exp\left(-\frac{1}{2}\|dA\|^2\right)\cdot\operatorname{det}_n\left(D_{A+H_{\theta}+a}+\boldsymbol{\mu}\right),
\qquad
\mathrm{Mj}_{\mathrm{rd}}(\theta,A):=\mathrm{Kd}(\mathrm{rd})\cdot\exp\left(-\frac{1}{2}\|dA\|^2\right).
\]
Here $\mathrm{vph}$ is continuous in $(\theta,A,a)$, and $\int\mathrm{d}\mathrm{vq}\,\mathrm{Mj}_{\mathrm{rd}}=\mathrm{Kd}(\mathrm{rd})\cdot\mathrm{zV}_n<\infty$ by Lemma~2.D\,(iii); the coboundary $d$ is injective on ${V_{\mathrm{t}}}$, because by (H1) its kernel on $C^{1}(n)$ is ${V_{\mathrm{g}}}+{V_{\mathrm{h}}}$.
Lemma~3.A now gives the absolute convergence of $\mathbf{Z}^{\mathrm{x}}_n(a)$, locally uniformly in $a$, and shows that $\mathbf{Z}^{\mathrm{x}}_n$ is entire.
\end{proof}

\subsubsection*{3.4. Proof of (b)}

\begin{proof}[Proof of Proposition~38.3\,(b)]
\emph{Step 0 (the left side is defined).}
Let $a\in C^{1}(n)\otimes\mathbb{C}$.
The function $u\mapsto\exp\langle da,u\rangle$ on $C^{2}(n)$ is continuous, and
\[
\left|\exp\langle da,u\rangle\right|=\exp\langle d\operatorname{Re}a,u\rangle\le\exp\left(\|d\operatorname{Re}a\|\cdot\|u\|\right).
\]
Hence it belongs to $\mathfrak{F}(n)$, with $\mathrm{cf}=\max\left(1,\|d\operatorname{Re}a\|\right)$.
By Proposition~50\,(e) the expressions $Z_n\left[\exp\langle da,F^{\mathrm{u}}\rangle\right]$ and $Z^{(w)}_n\left[\exp\langle da,F^{\mathrm{u}}\rangle\right]$ converge absolutely, and the second one is given by (38.12).

\emph{Step 1 (real $a$).}
Let $a\in C^{1}(n)$ and $w\in\mathbb{Z}^3$, and define, for $A\in C^{1}(n)$,
\[
\mathrm{Gf}(A):=\sum_{\mathfrak{z}\in\mathrm{Mz}_w}\exp\left(-\frac{1}{2}\|dA+r_q\mathfrak{z}\|^2+\langle da,dA+r_q\mathfrak{z}\rangle\right)\cdot\operatorname{det}_n\left(D_A+\boldsymbol{\mu}\right).
\]
Every term is continuous in $A$ and non-negative, the latter by Proposition~50\,(a), since $A$ is real.
Thus $\mathrm{Gf}$ is a Borel function with values in $[0,\infty]$.
Let $b\in\mathrm{Ln}(n)$.
The integer two-cochain $d\mathbf{1}_b$ is closed by (H1), and its flux is $0$ by Lemma~2.F\,(i); since the flux is additive, $\mathfrak{z}\mapsto\mathfrak{z}+d\mathbf{1}_b$ is a bijection of $\mathrm{Mz}_w$ onto itself, with the inverse $\mathfrak{z}\mapsto\mathfrak{z}-d\mathbf{1}_b$.
As $d(A+r_q\mathbf{1}_b)+r_q\mathfrak{z}=dA+r_q(\mathfrak{z}+d\mathbf{1}_b)$, and as $\operatorname{det}_n(D_{A+r_q\mathbf{1}_b}+\boldsymbol{\mu})=\operatorname{det}_n(D_A+\boldsymbol{\mu})$ by Proposition~38.1\,(b), it follows that
$\mathrm{Gf}(A+r_q\mathbf{1}_b)=\mathrm{Gf}(A)$.
So $\mathrm{Gf}$ is invariant under the group $\mathrm{Gq}:=r_q\,C^{1}(n;\mathbb{Z})$.
The set $\mathrm{Bo}$ is a fundamental domain of this group (division with remainder in each coordinate), and so is its translate $\mathrm{Bo}+a$.
By Lemma~2.B,
\begin{equation}\tag{38.18}
\int_{\mathrm{Bo}}\mathrm{d}A\,\mathrm{Gf}(A)=\int_{\mathrm{Bo}+a}\mathrm{d}A\,\mathrm{Gf}(A)=\int_{\mathrm{Bo}}\mathrm{d}A\,\mathrm{Gf}(A+a).
\end{equation}
By Tonelli's theorem and (38.12), the left side is $Z^{(w)}_n\left[\exp\langle da,F^{\mathrm{u}}\rangle\right]$.
On the right side, with $F^{\mathrm{u}}=dA+r_q\mathfrak{z}$,
\[
-\frac{1}{2}\|F^{\mathrm{u}}+da\|^2+\langle da,F^{\mathrm{u}}+da\rangle
=-\frac{1}{2}\|F^{\mathrm{u}}\|^2-\langle da,F^{\mathrm{u}}\rangle-\frac{1}{2}\|da\|^2+\langle da,F^{\mathrm{u}}\rangle+\|da\|^2
=-\frac{1}{2}\|F^{\mathrm{u}}\|^2+\frac{1}{2}\|da\|^2 .
\]
Therefore
\[
\mathrm{Gf}(A+a)=\exp\left(\frac{1}{2}\|da\|^2\right)\sum_{\mathfrak{z}\in\mathrm{Mz}_w}\exp\left(-\frac{1}{2}\|F^{\mathrm{u}}\|^2\right)\cdot\operatorname{det}_n\left(D_{A+a}+\boldsymbol{\mu}\right),
\]
with non-negative terms, because $A+a$ is real.
By Tonelli's theorem the right side of (38.18) equals $\exp\left(\frac{1}{2}\langle da,da\rangle\right)\cdot Z^{\mathrm{x},(w)}_n(a)$.
This is (b) for real $a$.

\emph{Step 2 (complex $a$).}
The right side of (b) is an entire function of $a\in\mathbb{C}^{\mathrm{Ln}(n)}$: the function $Z^{\mathrm{x},(w)}_n$ is entire by (a); $\langle da,da\rangle$ is a polynomial in the numbers $a(b)$, so that $\exp\left(\frac{1}{2}\langle da,da\rangle\right)$ is entire by (ef3); and the product of the two is entire by (ef3).
By (38.12) the left side of (b) is $\int\mathrm{d}\mathrm{vq}\,\mathrm{vph}$, where $\mathrm{Yq}=\mathrm{Mz}_w\times\mathrm{Bo}$, $\mathrm{vq}$ is the product of the counting measure and $\mathrm{d}A$, and
\[
\mathrm{vph}\left((\mathfrak{z},A),a\right):=\exp\left(-\frac{1}{2}\|F^{\mathrm{u}}\|^2\right)\cdot\operatorname{det}_n\left(D_A+\boldsymbol{\mu}\right)\cdot\exp\langle da,F^{\mathrm{u}}\rangle .
\]
For every $\mathfrak{z}$ this function is continuous in $(A,a)$; it is the exponential of a linear form in $a$ times a factor that does not contain $a$, hence entire in $a$ by (ef3).
Every component $(d\operatorname{Re}a)_p$ is a sum of four numbers $\pm\operatorname{Re}a(b)$; so, for $\max_b|a(b)|\le\mathrm{rd}$,
\[
\|d\operatorname{Re}a\|\le\mathrm{cf}_{\mathrm{rd}}:=4\,\mathrm{rd}\cdot(3|n|)^{\frac{1}{2}} .
\]
By Step 0, by the inequality $\mathrm{cf}_{\mathrm{rd}}\|F^{\mathrm{u}}\|\le\frac{1}{4}\|F^{\mathrm{u}}\|^2+\mathrm{cf}_{\mathrm{rd}}^2$ and by (38.16) at $a=0$,
\[
|\mathrm{vph}|\le\mathrm{Mj}_{\mathrm{rd}}:=\mathrm{Kd}(0)\cdot\exp\left(\mathrm{cf}_{\mathrm{rd}}^2\right)\cdot\exp\left(-\frac{1}{4}\|F^{\mathrm{u}}\|^2\right),
\]
and $\mathrm{Mj}_{\mathrm{rd}}$ is integrable by (38.6) with $\epsilon=\frac{1}{4}$.
By Lemma~3.A the left side of (b) is entire in $a$.
By Step 1 the two sides of (b) agree on $\mathbb{R}^{\mathrm{Ln}(n)}$, and by Lemma~3.B they agree on $\mathbb{C}^{\mathrm{Ln}(n)}$.

\emph{Step 3 (the sum over $w$).}
Each closed $\mathfrak{z}\in C^{2}(n;\mathbb{Z})$ belongs to exactly one of the sets $\mathrm{Mz}_w$.
The sum and the integral in (I.7.1) for $G=\exp\langle da,\,\cdot\,\rangle$ converge absolutely by Step 0.
Hence
\[
Z_n\left[\exp\langle da,F^{\mathrm{u}}\rangle\right]=\sum_wZ^{(w)}_n\left[\exp\langle da,F^{\mathrm{u}}\rangle\right]=\exp\left(\frac{1}{2}\langle da,da\rangle\right)\sum_wZ^{\mathrm{x},(w)}_n(a),
\]
and the last series converges absolutely by (a).
\end{proof}

\subsubsection*{3.5. Proof of (c)}

\begin{proof}[Proof of Proposition~38.3\,(c)]
Fix $a\in C^{1}(n)\otimes\mathbb{C}$ and put $\mathrm{wq}(A):=\operatorname{det}_n(D_{A+a}+\boldsymbol{\mu})$ for $A\in C^{1}(n)$.
The function $\mathrm{wq}$ is continuous by Proposition~38.1\,(a), and bounded by $\mathrm{Kd}\left(|\operatorname{Im}a|_{\infty}\right)$ by (38.16).
Moreover $\mathrm{wq}(A+d\boldsymbol{\lambda})=\mathrm{wq}(A)$ for every real $\boldsymbol{\lambda}$ (Proposition~38.1\,(c) with $B=A+a$), and $\mathrm{wq}(A+r_qm)=\mathrm{wq}(A)$ for every $m\in C^{1}(n;\mathbb{Z})$ (Proposition~38.1\,(b) with $B=A+a$).
Proposition~38.2\,(a), applied to this $\mathrm{wq}$ and to $\Phi=1$, has the left side $Z^{\mathrm{x},(0)}_n(a)$ and the right side $\mathrm{Nc}_n\cdot\mathbf{Z}^{\mathrm{x}}_n(a)$.
At $a=0$ one has $Z^{\mathrm{x},(0)}_n(0)=Z^{(0)}_n[1]$ by (38.12), and $Z^{(0)}_n[1]>0$ by Proposition~50\,(e).
Since $\mathrm{Nc}_n>0$, it follows that $\mathbf{Z}^{\mathrm{x}}_n(0)=Z^{(0)}_n[1]/\mathrm{Nc}_n>0$.
\end{proof}

\subsubsection*{3.6. Proof of (d)}

\begin{proof}[Proof of Proposition~38.3\,(d)]
Let $z\in\mathbb{C}^r$ and $a:=\sum_jz_ja^{(j)}\in C^{1}(n)\otimes\mathbb{C}$.
By bilinearity,
\[
\langle da,F^{\mathrm{u}}\rangle=\sum_jz_jX^{(j)},\qquad\langle da,da\rangle=\sum_{i,j}z_iz_j\,\langle da^{(i)},da^{(j)}\rangle .
\]
By (b) at $w=0$ and by (c),
\[
Z^{(0)}_n\left[\exp\langle da,F^{\mathrm{u}}\rangle\right]=\exp\left(\frac{1}{2}\langle da,da\rangle\right)\cdot\mathrm{Nc}_n\cdot\mathbf{Z}^{\mathrm{x}}_n(a).
\]
Divide by $Z^{(0)}_n[1]=\mathrm{Nc}_n\cdot\mathbf{Z}^{\mathrm{x}}_n(0)>0$ and by $\exp\left(\frac{1}{2}\langle da,da\rangle\right)$, a number that does not depend on $F^{\mathrm{u}}$:
\[
E^{\mathrm{av}}_n\exp\left(\langle da,F^{\mathrm{u}}\rangle-\frac{1}{2}\langle da,da\rangle\right)=\frac{\mathbf{Z}^{\mathrm{x}}_n(a)}{\mathbf{Z}^{\mathrm{x}}_n(0)} .
\]
With the two identities obtained from bilinearity this is (d).
\end{proof}

\subsubsection*{3.7. Proof of (e)}

\begin{proof}[Proof of Proposition~38.3\,(e)]
\emph{The setting} (Sections~I.6.2, I.6.4, I.7.6).
Let $e>0$ with $eh^{\frac{1}{2}}\le e_1$, let $N\ge N_0(e)$, $\eta=L^{-N}$ and $q=e\eta^{\frac{1}{2}}$, and consider the cubic torus with $n_*=hL^{T^{*}(e)+N}$ sites per side.
One has $h\ge1$ (because $\bar{m}\le1$), $L\ge5$ and $N\ge0$ (Section~I.0.5), and $T^{*}(e)\ge1$ (Section~I.6.2); hence $n_*\ge5$, and (a)--(d) hold for $n=n_*$.
The probe $\beta$ fits $(N,T^{*}(e))$, $v$ is a placement, $\hat{\mathfrak{a}}\in C^{1}(n_*)$ is real, and $\boldsymbol{\mu}$ is a real number.
$\bar{Z}^{\mathrm{ext}}$ and $\mathring{\Phi}$ are the displays of Section~I.7.6 formed with this $\boldsymbol{\mu}$, and $\Theta^{(0)}_m$ is the display of Section~I.7.6 formed with this $\mathring{\Phi}$ and with the real number $c_m$ of the clause.

\emph{The cochains of the cells.}
The torus has $n_*/\mathrm{hc}=2L$ cells per side.
$\mathrm{Cel}$ is the set of the classes of $\mathfrak{c}\in\mathbb{Z}^3$ modulo $2L\mathbb{Z}^3$, and $N_{*}=(2L)^3$.
By Section~I.6.4, in the coordinate $\mu$ the map $\mathsf{g}_{\mathfrak{c}+2Le_{\mu}}$ arises from $\mathsf{g}_{\mathfrak{c}}$ when $i$ is replaced by $i+L$ (where $\mathfrak{c}_\mu=2i$ or $\mathfrak{c}_\mu=2i+1$); since $L\cdot\mathrm{per}=n_*$, the map $\mathsf{g}_{\mathfrak{c}+2Le_{\mu}}$ is $\mathsf{g}_{\mathfrak{c}}$ followed by the translation by $n_*e_{\mu}$.
The linear part of $\mathsf{g}_{\mathfrak{c}}$ is a diagonal matrix of signs, so $(\mathsf{g}_{\mathfrak{c}})_*$ maps $C^{1}(n_*)$ into itself (Section~I.4.7; Section~G6.A1).
The rules $(g_{\mathrm{w}})_*=(g_{\mathrm{w}}^{-1})^{*}$ and $(g_{\mathrm{w}}g_{\mathrm{w}}')^{*}=(g_{\mathrm{w}}')^{*}g_{\mathrm{w}}^{*}$ of Section~I.4.7 give
\[
(g_{\mathrm{w}}g_{\mathrm{w}}')_*=\left((g_{\mathrm{w}}')^{-1}g_{\mathrm{w}}^{-1}\right)^{*}=\left(g_{\mathrm{w}}^{-1}\right)^{*}\left((g_{\mathrm{w}}')^{-1}\right)^{*}=(g_{\mathrm{w}})_*(g_{\mathrm{w}}')_* .
\]
A cochain of the torus is invariant under the translation by $n_*e_{\mu}$.
Consequently $\hat{\mathfrak{a}}_{\mathfrak{c}}=(\mathsf{g}_{\mathfrak{c}})_*\hat{\mathfrak{a}}\in C^{1}(n_*)$ depends only on the class of $\mathfrak{c}$ in $\mathrm{Cel}$, and it is real.
By Proposition~41\,(c), with the number $\boldsymbol{\sigma}=\|d\hat{\mathfrak{a}}\|$ of Section~I.7.6,
\begin{equation}\tag{38.19}
\langle d\hat{\mathfrak{a}}_{\mathfrak{c}},d\hat{\mathfrak{a}}_{\mathfrak{c}'}\rangle=\boldsymbol{\sigma}^2\cdot\mathbf{1}[\mathfrak{c}=\mathfrak{c}']\qquad(\mathfrak{c},\mathfrak{c}'\in\mathrm{Cel}).
\end{equation}

\emph{$\bar{Z}^{\mathrm{ext}}$ and $\mathring{\Phi}$: the first three assertions of (e).}
For $t\in\mathbb{C}^{\mathrm{Cel}}$ let
\[
a(t):=e^{-1}\sum_{\mathfrak{c}\in\mathrm{Cel}}t_{\mathfrak{c}}\hat{\mathfrak{a}}_{\mathfrak{c}}\in C^{1}(n_*)\otimes\mathbb{C}.
\]
For every $\theta\in\mathbb{R}^3$ the function $u\mapsto\operatorname{det}_{n_*}\left(D_{A(u)+H_{\theta}+a(t)}+\boldsymbol{\mu}\right)$ on $F_N$ is continuous and bounded by $\mathrm{Kd}\left(|\operatorname{Im}a(t)|_{\infty}\right)$, by (38.16); so it satisfies the bound of Lemma~2.G with $\mathrm{cf}:=\max\left(1,\mathrm{Kd}\left(|\operatorname{Im}a(t)|_{\infty}\right)\right)$.
Lemma~2.G, applied for every $\theta$, together with $A(dA)=A$ on ${V_{\mathrm{t}}}$, gives
\[
\bar{Z}^{\mathrm{ext}}(t)=\mathrm{zV}_{n_*}^{-1}\cdot\mathbf{Z}^{\mathrm{x}}_{n_*}\left(a(t)\right).
\]
By (c), $\bar{Z}^{\mathrm{ext}}(0)=\mathbf{Z}^{\mathrm{x}}_{n_*}(0)/\mathrm{zV}_{n_*}>0$.
Hence
\[
\mathring{\Phi}(t)=\frac{\mathbf{Z}^{\mathrm{x}}_{n_*}\left(a(t)\right)}{\mathbf{Z}^{\mathrm{x}}_{n_*}(0)} .
\]
The map $t\mapsto a(t)$ is linear, and $\mathbf{Z}^{\mathrm{x}}_{n_*}$ is entire by (a); so $\mathring{\Phi}$ is entire on $\mathbb{C}^{\mathrm{Cel}}=\mathbb{C}^{N_{*}}$ by (ef2).

\emph{The generating identity.}
We apply (d) to the $N_{*}$ cochains $\hat{\mathfrak{a}}_{\mathfrak{c}}$, $\mathfrak{c}\in\mathrm{Cel}$.
By (38.19),
$\sum_{\mathfrak{c},\mathfrak{c}'}z_{\mathfrak{c}}z_{\mathfrak{c}'}\langle d\hat{\mathfrak{a}}_{\mathfrak{c}},d\hat{\mathfrak{a}}_{\mathfrak{c}'}\rangle=\boldsymbol{\sigma}^2\sum_{\mathfrak{c}}z_{\mathfrak{c}}^2$,
and $\sum_{\mathfrak{c}}z_{\mathfrak{c}}\hat{\mathfrak{a}}_{\mathfrak{c}}=a(ez)$.
So, for $z\in\mathbb{C}^{N_{*}}$,
\begin{equation}\tag{38.20}
E^{\mathrm{av}}_{n_*}\exp\sum_{\mathfrak{c}}\left(z_{\mathfrak{c}}X_{\mathfrak{c}}-\frac{1}{2}z_{\mathfrak{c}}^2\boldsymbol{\sigma}^2\right)=\frac{\mathbf{Z}^{\mathrm{x}}_{n_*}\left(a(ez)\right)}{\mathbf{Z}^{\mathrm{x}}_{n_*}(0)}=\mathring{\Phi}(ez).
\end{equation}

\emph{Hermite polynomials.}
For a real number $y$ the function $z\mapsto\exp\left(zy-\frac{1}{2}z^2\boldsymbol{\sigma}^2\right)$ is entire on $\mathbb{C}$ by (ef3).
Let
\[
H_i(y):=\left.\left(\frac{\mathrm{d}}{\mathrm{d}z}\right)^i\exp\left(zy-\frac{1}{2}z^2\boldsymbol{\sigma}^2\right)\right|_{z=0}\qquad(i\ge0);
\]
by (38.0b), $\exp\left(zy-\frac{1}{2}z^2\boldsymbol{\sigma}^2\right)=\sum_{i\ge0}H_i(y)\,z^i/i!$.
The product of the two exponential series is
\[
\sum_{j,j'\ge0}\frac{y^j}{j!}\cdot\frac{1}{j'!}\left(-\frac{\boldsymbol{\sigma}^2}{2}\right)^{j'}z^{j+2j'} ,
\]
and its coefficients of $z^0$, $z^2$, $z^4$ are
\[
1,\qquad\frac{y^2}{2}-\frac{\boldsymbol{\sigma}^2}{2},\qquad\frac{y^4}{24}-\frac{\boldsymbol{\sigma}^2}{2}\cdot\frac{y^2}{2}+\frac{1}{2}\cdot\frac{\boldsymbol{\sigma}^4}{4} .
\]
Multiplying them by $0!$, $2!$, $4!$ one obtains
\[
H_0=1,\qquad H_2(y)=y^2-\boldsymbol{\sigma}^2,\qquad H_4(y)=y^4-6\boldsymbol{\sigma}^2y^2+3\boldsymbol{\sigma}^4 ,
\]
which are the polynomials $H_2$, $H_4$ of Section~I.7.6.

\emph{Differentiation under the expectation.}
We apply Lemma~3.A with $\mathrm{Nq}=N_{*}$, $\mathrm{Yq}=\mathrm{Mz}_0\times\mathrm{Bo}$, the measure
\[
\mathrm{d}\mathrm{vq}:=Z^{(0)}_{n_*}[1]^{-1}\cdot\exp\left(-\frac{1}{2}\|F^{\mathrm{u}}\|^2\right)\cdot\operatorname{det}_{n_*}\left(D_A+\boldsymbol{\mu}\right)\ \text{times the product of the counting measure and }\mathrm{d}A,
\]
which is a probability measure (its density is non-negative by Proposition~50\,(a), and (38.12) holds), and
\[
\mathrm{vph}\left((\mathfrak{z},A),z\right):=\exp\sum_{\mathfrak{c}}\left(z_{\mathfrak{c}}X_{\mathfrak{c}}-\frac{1}{2}z_{\mathfrak{c}}^2\boldsymbol{\sigma}^2\right)=\prod_{\mathfrak{c}}\exp\left(z_{\mathfrak{c}}X_{\mathfrak{c}}-\frac{1}{2}z_{\mathfrak{c}}^2\boldsymbol{\sigma}^2\right),\qquad X_{\mathfrak{c}}=\langle d\hat{\mathfrak{a}}_{\mathfrak{c}},F^{\mathrm{u}}\rangle .
\]
For every $\mathfrak{z}$ the function $\mathrm{vph}$ is continuous in $(A,z)$, and it is entire in $z$ by (ef3).
By (38.12), $\int\mathrm{d}\mathrm{vq}\,G(F^{\mathrm{u}})=E^{\mathrm{av}}_{n_*}G$ for $G\in\mathfrak{F}(n_*)$.
The Cauchy--Schwarz inequality and (38.19) give $|X_{\mathfrak{c}}|\le\boldsymbol{\sigma}\|F^{\mathrm{u}}\|$.
Hence, for $\max_{\mathfrak{c}}|z_{\mathfrak{c}}|\le\mathrm{rd}$,
\[
|\mathrm{vph}|\le\mathrm{Mj}_{\mathrm{rd}}:=\exp\left(\mathrm{rd}\cdot N_{*}\boldsymbol{\sigma}\|F^{\mathrm{u}}\|+\frac{1}{2}N_{*}\cdot\mathrm{rd}^2\boldsymbol{\sigma}^2\right),
\]
a function of $F^{\mathrm{u}}$ that belongs to $\mathfrak{F}(n_*)$, with $\mathrm{cf}:=\max\left(\mathrm{rd}\cdot N_{*}\boldsymbol{\sigma},\exp\left(\frac{1}{2}N_{*}\cdot\mathrm{rd}^2\boldsymbol{\sigma}^2\right)\right)$; it is therefore $\mathrm{vq}$-integrable (Proposition~50\,(e) together with (38.12)).
Let $m\in\{2,4\}$, and let $\mathrm{Sb}$ be a subset of $\mathrm{Cel}$.
The derivative $\prod_{\mathfrak{c}\in\mathrm{Sb}}\left(\partial/\partial z_{\mathfrak{c}}\right)^m$ of $\mathrm{vph}$ at $z=0$ is
\[
\prod_{\mathfrak{c}\in\mathrm{Sb}}H_m(X_{\mathfrak{c}})\cdot\prod_{\mathfrak{c}\notin\mathrm{Sb}}H_0(X_{\mathfrak{c}})=\prod_{\mathfrak{c}\in\mathrm{Sb}}H_m(X_{\mathfrak{c}}).
\]
For real $y$,
\[
|H_2(y)|\le y^2+\boldsymbol{\sigma}^2,\qquad|H_4(y)|\le y^4+6\boldsymbol{\sigma}^2y^2+3\boldsymbol{\sigma}^4\le3\left(y^2+\boldsymbol{\sigma}^2\right)^2 .
\]
Together with $|X_{\mathfrak{c}}|\le\boldsymbol{\sigma}\|F^{\mathrm{u}}\|$ this gives $|H_m(X_{\mathfrak{c}})|\le3\boldsymbol{\sigma}^m\left(1+\|F^{\mathrm{u}}\|\right)^m$.
Using $(1+y)^i\le i!\cdot\exp(1+y)$ for $y\ge0$, one sees that the function $u\mapsto\prod_{\mathfrak{c}\in\mathrm{Sb}}H_m\left(\langle d\hat{\mathfrak{a}}_{\mathfrak{c}},u\rangle\right)$ on $C^{2}(n_*)$ is continuous and has modulus at most
\[
3^{\#|\mathrm{Sb}|}\cdot\boldsymbol{\sigma}^{m\#|\mathrm{Sb}|}\cdot\left(m\#|\mathrm{Sb}|\right)!\cdot\mathrm{e}\cdot\exp\left(\|u\|\right).
\]
It therefore lies in $\mathfrak{F}(n_*)$, and so does every finite linear combination of such functions; $E^{\mathrm{av}}_{n_*}$ is defined on all of them.
If $\mathrm{tay}_{\mathrm{kx}}$ denote the coefficients of $\mathring{\Phi}$, then $\mathring{\Phi}(ez)=\sum_{\mathrm{kx}}\mathrm{tay}_{\mathrm{kx}}\cdot e^{|\mathrm{kx}|}z^{\mathrm{kx}}$.
By (38.0b), $\partial_z^{\mathrm{kx}}\mathring{\Phi}(ez)$ at $z=0$ is therefore $e^{|\mathrm{kx}|}\cdot\partial_t^{\mathrm{kx}}\mathring{\Phi}(t)$ at $t=0$.
For $\mathrm{kx}_{\mathfrak{c}}=m$ ($\mathfrak{c}\in\mathrm{Sb}$) and $\mathrm{kx}_{\mathfrak{c}}=0$ ($\mathfrak{c}\notin\mathrm{Sb}$) one has $|\mathrm{kx}|=m\#|\mathrm{Sb}|$.
Lemma~3.A and (38.20) now yield
\begin{equation}\tag{38.21}
E^{\mathrm{av}}_{n_*}\prod_{\mathfrak{c}\in\mathrm{Sb}}H_m(X_{\mathfrak{c}})=\left.\prod_{\mathfrak{c}\in\mathrm{Sb}}\left(\frac{\partial}{\partial z_{\mathfrak{c}}}\right)^m\mathring{\Phi}(ez)\right|_{z=0}=e^{m\#|\mathrm{Sb}|}\left.\prod_{\mathfrak{c}\in\mathrm{Sb}}\partial_{t_{\mathfrak{c}}}^m\mathring{\Phi}(t)\right|_{t=0}.
\end{equation}

\emph{Expansion.}
With the sum over all subsets $\mathrm{Sb}$ of $\mathrm{Cel}$,
\[
\prod_{\mathfrak{c}\in\mathrm{Cel}}\left(e^{-m}H_m(X_{\mathfrak{c}})-c_m\right)=\sum_{\mathrm{Sb}}\left(-c_m\right)^{N_{*}-\#|\mathrm{Sb}|}\prod_{\mathfrak{c}\in\mathrm{Sb}}e^{-m}H_m(X_{\mathfrak{c}}).
\]
The operators $\partial_{t_{\mathfrak{c}}}$ commute on entire functions by (ef1); hence
\[
\prod_{\mathfrak{c}\in\mathrm{Cel}}\left(\partial_{t_{\mathfrak{c}}}^m-c_m\right)=\sum_{\mathrm{Sb}}\left(-c_m\right)^{N_{*}-\#|\mathrm{Sb}|}\prod_{\mathfrak{c}\in\mathrm{Sb}}\partial_{t_{\mathfrak{c}}}^m .
\]
Using (38.21), divided by $e^{m\#|\mathrm{Sb}|}$, and the linearity of $E^{\mathrm{av}}_{n_*}$, we obtain
\[
E^{\mathrm{av}}_{n_*}\prod_{\mathfrak{c}\in\mathrm{Cel}}\left(e^{-m}H_m(X_{\mathfrak{c}})-c_m\right)=\left.\prod_{\mathfrak{c}\in\mathrm{Cel}}\left(\partial_{t_{\mathfrak{c}}}^m-c_m\right)\mathring{\Phi}(t)\right|_{t=0}.
\]
By the definition of $\Theta^{(0)}_m$ in Section~I.7.6, the modulus of the right side is $\left(\Theta^{(0)}_m\right)^{N_{*}}$.
\end{proof}

\subsubsection*{3.8. Proof of (a$'$)}

The expressions $Z^{\mathrm{x},(w)}_n(a;\Phi)$ and $\mathbf{Z}^{\mathrm{x}}_n(a;\Phi)$ of (a$'$) are the displays (38.22) below.
The first half of (a$'$) is Lemma~3.C: the convergence in the sense of Section~0.6, the fact that both expressions are entire, the relation (38.23), and the value at $a=0$.
The second half, which concerns the cubic torus and the measures $\bar{\nu}^a$, is Corollary~3.D.
There the power that appears in (a$'$),
\[
\left(|M_{\mathrm{ch}}|+\frac{1+\rho}{2}\left(\mathrm{sw}-6+6\exp\left(q|\operatorname{Im}a|_{\infty}\right)\right)\right)^{192n^3},
\]
is $\mathrm{Kd}\left(|\operatorname{Im}a|_{\infty}\right)$, by the definition of $\mathrm{Kd}$ at the beginning of Section~2 with $|n|=n^3$.

\begin{named}{Lemma 3.C}[A further factor: the first half of (a$'$)]
Let $\Phi\in\mathfrak{F}(n)$ with $|\Phi(u)|\le\mathrm{cf}\cdot\exp\left(\mathrm{cf}\|u\|\right)$.
For $a\in C^{1}(n)\otimes\mathbb{C}$ and $w\in\mathbb{Z}^3$ put, with $F^{\mathrm{u}}=dA+r_q\mathfrak{z}$,
\begin{equation}\tag{38.22}
\begin{aligned}
Z^{\mathrm{x},(w)}_n(a;\Phi)&:=\sum_{\mathfrak{z}\in\mathrm{Mz}_w}\int_{\mathrm{Bo}}\mathrm{d}A\;\Phi(F^{\mathrm{u}})\cdot\exp\left(-\frac{1}{2}\|F^{\mathrm{u}}\|^2\right)\cdot\operatorname{det}_n\left(D_{A+a}+\boldsymbol{\mu}\right),\\
\mathbf{Z}^{\mathrm{x}}_n(a;\Phi)&:=\int_{[-\pi,\pi]^3}\frac{\mathrm{d}\theta}{(2\pi)^3}\int_{{V_{\mathrm{t}}}}\mathrm{d}\sigma(A)\;\Phi(dA)\cdot\exp\left(-\frac{1}{2}\|dA\|^2\right)\cdot\operatorname{det}_n\left(D_{A+H_{\theta}+a}+\boldsymbol{\mu}\right).
\end{aligned}
\end{equation}
Both expressions converge absolutely; for every $\mathrm{rd}\in(0,\infty)$ the integrands have integrable majorants that are the same for all $a$ in $\left\{\max_b|a(b)|\le\mathrm{rd}\right\}$; both are entire functions of $a\in\mathbb{C}^{3|n|}$; and
\begin{equation}\tag{38.23}
Z^{\mathrm{x},(0)}_n(a;\Phi)=\mathrm{Nc}_n\cdot\mathbf{Z}^{\mathrm{x}}_n(a;\Phi)\qquad\text{for every }a\in C^{1}(n)\otimes\mathbb{C}.
\end{equation}
For $\Phi=1$ these expressions are $Z^{\mathrm{x},(w)}_n(a)$ and $\mathbf{Z}^{\mathrm{x}}_n(a)$ of Section~I.7.4.
At $a=0$, $Z^{\mathrm{x},(w)}_n(0;\Phi)=Z^{(w)}_n[\Phi]$ by (38.12).
\end{named}

\begin{proof}
Lemma~3.A is applied twice, with $\mathrm{Nq}=3|n|$.

\emph{The first expression.}
Take $\mathrm{Yq}=\mathrm{Mz}_w\times\mathrm{Bo}$, $\mathrm{vq}$ the product of the counting measure and $\mathrm{d}A$, and
\[
\mathrm{vph}\left((\mathfrak{z},A),a\right):=\Phi(dA+r_q\mathfrak{z})\cdot\exp\left(-\frac{1}{2}\|dA+r_q\mathfrak{z}\|^2\right)\cdot\operatorname{det}_n\left(D_{A+a}+\boldsymbol{\mu}\right).
\]
For every $\mathfrak{z}$ the function $A\mapsto\Phi(dA+r_q\mathfrak{z})$ is a Borel function ($\Phi$ is a Borel function and $A\mapsto dA+r_q\mathfrak{z}$ is continuous), and the two other factors are continuous in $(A,a)$ (Section~3.3); so $\mathrm{vph}$ is measurable.
It is the determinant times a factor that does not contain $a$, hence entire in $a$ (Section~3.3).
For $\max_b|a(b)|\le\mathrm{rd}$, (38.16) and (38.7) give
\[
|\mathrm{vph}|\le\mathrm{Mj}_{\mathrm{rd}}:=\mathrm{Kd}(\mathrm{rd})\cdot\mathrm{cf}\cdot\exp\left(\mathrm{cf}^2\right)\cdot\exp\left(-\frac{1}{4}\|dA+r_q\mathfrak{z}\|^2\right),
\]
and (38.6) with $\epsilon=\frac{1}{4}$ gives
\[
\int\mathrm{d}\mathrm{vq}\,\mathrm{Mj}_{\mathrm{rd}}\le\mathrm{Kd}(\mathrm{rd})\cdot\mathrm{cf}\cdot\exp\left(\mathrm{cf}^2\right)\cdot\left(r_q\cdot\mathrm{Cth}\left(\frac{1}{4}\right)\right)^{3|n|}.
\]

\emph{The second expression.}
Take $\mathrm{Yq}=[-\pi,\pi]^3\times{V_{\mathrm{t}}}$, $\mathrm{vq}$ the product of $\mathrm{d}\theta/(2\pi)^3$ and $\sigma$, and
\[
\mathrm{vph}\left((\theta,A),a\right):=\Phi(dA)\cdot\exp\left(-\frac{1}{2}\|dA\|^2\right)\cdot\operatorname{det}_n\left(D_{A+H_{\theta}+a}+\boldsymbol{\mu}\right).
\]
The function $A\mapsto\Phi(dA)$ is a Borel function, and the two other factors are continuous in $(\theta,A,a)$ (Section~3.3); so $\mathrm{vph}$ is measurable.
It is the determinant times a factor that does not contain $a$, hence entire in $a$ (Section~3.3).
For $\max_b|a(b)|\le\mathrm{rd}$, (38.16) and (38.7) give
\[
|\mathrm{vph}|\le\mathrm{Mj}_{\mathrm{rd}}:=\mathrm{Kd}(\mathrm{rd})\cdot\mathrm{cf}\cdot\exp\left(\mathrm{cf}^2\right)\cdot\exp\left(-\frac{1}{4}\|dA\|^2\right).
\]
Apply (Lin) to the map $A\mapsto2^{-\frac{1}{2}}A$ of ${V_{\mathrm{t}}}$, whose determinant is $2^{-(|n|-1)}$ because $\dim{V_{\mathrm{t}}}=2|n|-2$, and to the function $\exp\left(-\frac{1}{2}\|dA\|^2\right)$:
\[
\int_{{V_{\mathrm{t}}}}\mathrm{d}\sigma(A)\,\exp\left(-\frac{1}{4}\|dA\|^2\right)=2^{|n|-1}\cdot\mathrm{zV}_n .
\]
Hence $\int\mathrm{d}\mathrm{vq}\,\mathrm{Mj}_{\mathrm{rd}}=\mathrm{Kd}(\mathrm{rd})\cdot\mathrm{cf}\cdot\exp\left(\mathrm{cf}^2\right)\cdot2^{|n|-1}\cdot\mathrm{zV}_n<\infty$.

\emph{The relation (38.23).}
Apply Proposition~38.2\,(a) to this $\Phi$ and to $\mathrm{wq}(A):=\operatorname{det}_n(D_{A+a}+\boldsymbol{\mu})$; the hypotheses on $\mathrm{wq}$ were checked in Section~3.5.
The left side of Proposition~38.2\,(a) is $Z^{\mathrm{x},(0)}_n(a;\Phi)$, and its right side is $\mathrm{Nc}_n\cdot\mathbf{Z}^{\mathrm{x}}_n(a;\Phi)$.
\end{proof}

\begin{named}{Corollary 3.D}[The measures $\bar{\nu}^a$: the second half of (a$'$)]
Let the torus be cubic with $n\ge4$ sites per side, with the objects of Section~I.7.3 for a real pair $(m_{\mathrm{in}},\varepsilon^{\mathrm{v}}_{\mathrm{in}})$, and let $\boldsymbol{\mu}=\bar{m}\eta+m_{\mathrm{in}}$.
For $a\in C^{1}(n)\otimes\mathbb{C}$ the displays of Section~I.7.3 define complex measures $\nu^{\theta,a}$ ($\theta\in\mathbb{R}^3$) and $\bar{\nu}^a$ on the Borel sets of $F_N$.
For $\Phi\in\mathfrak{F}(n)$ the integral
\[
\bar{\nu}^a(\Phi):=\int_{[-\pi,\pi]^3}\frac{\mathrm{d}\theta}{(2\pi)^3}\int\mathrm{d}\gamma_N\,\Phi\cdot\mathrm{dens}^{\theta,a}_N
\]
converges absolutely, and
\begin{equation}\tag{38.24}
\bar{\nu}^a(\Phi)=\frac{\exp\left(-\varepsilon^{\mathrm{v}}_{\mathrm{in}}n^3\right)}{\operatorname{det}_n\left(D_0+\bar{m}\eta\right)\cdot\mathrm{zV}_n}\cdot\mathbf{Z}^{\mathrm{x}}_n(a;\Phi)
=\frac{\exp\left(-\varepsilon^{\mathrm{v}}_{\mathrm{in}}n^3\right)}{\operatorname{det}_n\left(D_0+\bar{m}\eta\right)\cdot\mathrm{zV}_n\cdot\mathrm{Nc}_n}\cdot Z^{\mathrm{x},(0)}_n(a;\Phi).
\end{equation}
So $a\mapsto\bar{\nu}^a(\Phi)$ is an entire function on $\mathbb{C}^{3n^3}$.
In particular, for every Borel subset $\mathrm{Ev}$ of $F_N$, the function $a\mapsto\bar{\nu}^a(\mathrm{Ev})$ is entire on $\mathbb{C}^{3n^3}$,
\[
\left|\bar{\nu}^a(\mathrm{Ev})\right|\le\frac{\exp\left(-\varepsilon^{\mathrm{v}}_{\mathrm{in}}n^3\right)\cdot\mathrm{Kd}\left(|\operatorname{Im}a|_{\infty}\right)}{\operatorname{det}_n\left(D_0+\bar{m}\eta\right)},
\]
and for every linear map $t\mapsto a(t)$ of $\mathbb{C}^{\mathrm{Nq}'}$ into $C^{1}(n)\otimes\mathbb{C}$ the function $t\mapsto\bar{\nu}^{a(t)}(\mathrm{Ev})$ is entire on $\mathbb{C}^{\mathrm{Nq}'}$.
On the torus $n=n_*$ this applies to $a(t)=e^{-1}\sum_{\mathfrak{c}\in\mathrm{Cel}}t_{\mathfrak{c}}\hat{\mathfrak{a}}_{\mathfrak{c}}$ of Section~3.7.
\end{named}

\begin{proof}
By (I.7.2),
\[
\mathrm{dens}^{\theta,a}_N(u)=\exp\left(-\varepsilon^{\mathrm{v}}_{\mathrm{in}}n^3\right)\cdot\frac{\operatorname{det}_n\left(D_{A(u)+H_{\theta}+a}+\boldsymbol{\mu}\right)}{\operatorname{det}_n\left(D_0+\bar{m}\eta\right)},
\]
and the denominator is positive (Section~2.8).
Fix $\theta\in\mathbb{R}^3$, and let
\[
G(u):=\Phi(u)\cdot\operatorname{det}_n\left(D_{A(u)+H_{\theta}+a}+\boldsymbol{\mu}\right),\qquad u\in F_N .
\]
$G$ is a Borel function (the determinant is continuous in $u$, because $A(\,\cdot\,)$ is linear), and by (38.16)
\[
|G(u)|\le\mathrm{Kd}\left(|\operatorname{Im}a|_{\infty}\right)\cdot\mathrm{cf}\cdot\exp\left(\mathrm{cf}\|u\|\right)\le\mathrm{cf}'\cdot\exp\left(\mathrm{cf}'\|u\|\right),\qquad\mathrm{cf}':=\max\left(1,\mathrm{Kd}\left(|\operatorname{Im}a|_{\infty}\right)\right)\cdot\mathrm{cf}.
\]
By Lemma~2.G, and since $A(dA)=A$ for $A\in{V_{\mathrm{t}}}$,
\[
\int\mathrm{d}\gamma_N\,G=\mathrm{zV}_n^{-1}\int_{{V_{\mathrm{t}}}}\mathrm{d}\sigma(A)\;\Phi(dA)\cdot\exp\left(-\frac{1}{2}\|dA\|^2\right)\cdot\operatorname{det}_n\left(D_{A+H_{\theta}+a}+\boldsymbol{\mu}\right),
\]
both sides converging absolutely.
The integrands are Borel functions of $(\theta,u)$, resp.\ of $(\theta,A)$, and the majorant
$\mathrm{Kd}\left(|\operatorname{Im}a|_{\infty}\right)\cdot\mathrm{cf}\cdot\exp\left(\mathrm{cf}^2\right)\cdot\exp\left(-\frac{1}{4}\|dA\|^2\right)$
of the integrand on the right does not contain $\theta$.
Integrating over $\theta$ with respect to $\mathrm{d}\theta/(2\pi)^3$ one obtains the absolute convergence and the first equality of (38.24); the second equality is (38.23).
By Lemma~3.C the right sides of (38.24) are entire in $a$.

By (38.16),
\[
\left|\mathrm{dens}^{\theta,a}_N(u)\right|\le\frac{\exp\left(-\varepsilon^{\mathrm{v}}_{\mathrm{in}}n^3\right)\cdot\mathrm{Kd}\left(|\operatorname{Im}a|_{\infty}\right)}{\operatorname{det}_n\left(D_0+\bar{m}\eta\right)}\qquad\text{for all }\theta\in\mathbb{R}^3\text{ and }u\in F_N,
\]
and $\mathrm{dens}^{\theta,a}_N(u)$ is continuous in $(\theta,u)$ (Section~3.3; $A(\,\cdot\,)$ is linear).
If $\mathrm{Gv}$ is a bounded Borel function on $F_N$, then $\mathrm{Ev}\mapsto\int_{\mathrm{Ev}}\mathrm{d}\gamma_N\,\mathrm{Gv}$ is a complex measure on the Borel sets of $F_N$.
Indeed, $\gamma_N$ is a probability measure, and for pairwise disjoint Borel sets $\mathrm{Ev}_i$ with the union $\mathrm{Ev}$ the series $\sum_i\mathbf{1}_{\mathrm{Ev}_i}\mathrm{Gv}=\mathbf{1}_{\mathrm{Ev}}\mathrm{Gv}$ may be integrated term by term (Fubini's theorem for the product of the counting measure and $\gamma_N$; the moduli have the sum $\int_{\mathrm{Ev}}\mathrm{d}\gamma_N\,|\mathrm{Gv}|<\infty$).
With $\mathrm{Gv}=\mathrm{dens}^{\theta,a}_N$ this shows that $\nu^{\theta,a}$ is a complex measure.
By Fubini's theorem,
\[
\bar{\nu}^a(\mathrm{Ev})=\int_{[-\pi,\pi]^3}\frac{\mathrm{d}\theta}{(2\pi)^3}\,\nu^{\theta,a}(\mathrm{Ev})=\int_{\mathrm{Ev}}\mathrm{d}\gamma_N(u)\left(\int_{[-\pi,\pi]^3}\frac{\mathrm{d}\theta}{(2\pi)^3}\,\mathrm{dens}^{\theta,a}_N(u)\right),
\]
and the inner integral is a bounded Borel function of $u$.
So $\bar{\nu}^a$ is a complex measure.

$F_N$ is a linear subspace of $C^{2}(n)$, hence a closed set.
Therefore a Borel subset $\mathrm{Ev}$ of $F_N$ is a Borel subset of $C^{2}(n)$, the function $\Phi$ that equals $1$ on $\mathrm{Ev}$ and $0$ on $C^{2}(n)\setminus\mathrm{Ev}$ lies in $\mathfrak{F}(n)$ (with $\mathrm{cf}=1$), and $\bar{\nu}^a(\mathrm{Ev})=\bar{\nu}^a(\Phi)$ by the definition of $\bar{\nu}^a$ in Section~I.7.3.
The bound on $\left|\bar{\nu}^a(\mathrm{Ev})\right|$ follows from (38.16), because $\gamma_N$ and $\mathrm{d}\theta/(2\pi)^3$ on $[-\pi,\pi]^3$ are probability measures.
The assertion on $t\mapsto\bar{\nu}^{a(t)}(\mathrm{Ev})$ follows by (ef2).
\end{proof}

\subsubsection*{3.9. Range}

The assertions (a)--(d) hold for all $n_\mu\ge4$ of either parity, all $q>0$ and all real $\boldsymbol{\mu}$.
The bounds in (a) are
\[
\sum_w\left|Z^{\mathrm{x},(w)}_n(a)\right|\le\mathrm{Kd}\left(|\operatorname{Im}a|_{\infty}\right)\cdot\left(r_q\cdot\mathrm{Cth}\right)^{3|n|},\qquad
\left|\mathbf{Z}^{\mathrm{x}}_n(a)\right|\le\mathrm{Kd}\left(|\operatorname{Im}a|_{\infty}\right)\cdot\mathrm{zV}_n .
\]
Assertion (e) holds for every $e$ with $eh^{\frac{1}{2}}\le e_1$, every $N\ge N_0(e)$, every probe that fits, every placement, every real $\boldsymbol{\mu}$ and all real $c_m$; the case $\boldsymbol{\sigma}=0$ is included.
As to (a$'$): Lemma~3.C holds for all $n_\mu\ge4$ of either parity, all $q>0$, all real $\boldsymbol{\mu}$ and all $\Phi\in\mathfrak{F}(n)$; Corollary~3.D holds for every cubic torus with $n\ge4$, every real pair $(m_{\mathrm{in}},\varepsilon^{\mathrm{v}}_{\mathrm{in}})$, every $\bar{m}\in(0,1]$ and every $\eta=L^{-N}$.

\subsection*{4. Conditions on the constants}

The following conditions on the constants enter the proofs of Propositions~38.1--38.3.
The last column names the line of the ordered list of constants of Section~I.8.2 that contains the condition.

\begin{tabular}{p{0.28\textwidth}p{0.36\textwidth}p{0.26\textwidth}}
\hline
condition & where it is needed & line of the list of constants \\
\hline
$\rho\ge0$ (part of $\rho\in(0,1)$), in the form $|M_i(x)|_2=\frac{1}{2}w_i(1+\rho)$ of (I.m6)
&
Proposition~38.1\,(d), (38.5); through $\mathrm{Kd}$ in Proposition~38.2\,(b), (c) and in Proposition~38.3\,(a), (a$'$), (b)--(e)
&
(C0) of Section~I.8.2, that is, Proposition~14.2\,(a), row $\rho$ \\
\hline
the standard weights: $w_3=6$, $\mathrm{sw}=w_0+w_1+w_2+w_3$, and $\sum_iw_i\hat n_i=0$ ((I.1.1), in the form (I.2.5))
&
(38.1), (38.2) (all clauses of Proposition~38.1); Proposition~38.1\,(d) (the numbers $|M_{\mathrm{ch}}|$, $6$, $\mathrm{sw}-6$)
&
(C0), row $w$ \\
\hline
$\bar{m}\in(0,1]$
&
Proposition~38.2\,(c) and Proposition~38.3\,(a$'$) (Corollary~3.D): $\bar{m}\eta>0$, so that the denominator $\operatorname{det}_n(D_0+\bar{m}\eta)$ is positive (Section~2.8); Proposition~38.3\,(e): $\bar{m}\le1$ gives $h\ge1$, which yields $n_*\ge5$ (Section~3.7)
&
(C10), rows $\bar{m}$, $h$; Section~I.0.5 \\
\hline
$L\ge5$; $M$ a power of $L$ with $M\ge L^2$ (through $T^{*}(e)\ge1$, Section~I.6.2)
&
Proposition~38.3\,(e): $n_*\ge5\ge4$ (Section~3.7)
&
(C3), label (L5) of Proposition~14.2\,(b); (C5), label (M1) of Proposition~14.2\,(b) (also Annex~B, Section~G6.Z4, row $M\ge L^2$) \\
\hline
the conditions behind the clauses Proposition~50\,(a), (e) ($\rho\in(0,1)$: Annex~B, Section~G6.Z4, last paragraph) and Proposition~41\,(c), which are applied
&
Proposition~38.2\,(b), (c); Proposition~38.3\,(b)--(e); Proposition~38.3\,(a$'$) (Corollary~3.D: Proposition~50\,(a), through Section~2.8)
&
Proposition~50\,(a), (e): (C0); Proposition~41\,(c): (C3) with (L5), (C5) with (M1), (C10), rows $\bar{m}$, $h$ (the three rows of the table of conditions of Chapter~41) \\
\hline
\end{tabular}

The sides $n_\mu\ge1$ (Proposition~38.1) and $n_\mu\ge4$ (Propositions~38.2, 38.3), the real number $\boldsymbol{\mu}$, the inequalities $eh^{\frac{1}{2}}\le e_1$ and $N\ge N_0(e)$, and the requirement that $\beta$ fits are data of the statements.

\section*{Chapter 50. The determinant and the invariances of the functionals on rectangular tori of either parity: proof of Proposition~50}

This chapter contains the proof of Proposition~50.
The proof is arranged as follows.
Section~0 fixes the setting and proves Lemmas~1--6; Lemma~1 contains the last sentence of the preamble of Proposition~50.
The clauses are then proved in the order (a); (b); (c); (d); the first part of (e) (convergence, the sum over the sectors); (f) and (g); (g$'$); the second part of (e) (positivity, the order of the sectors, the laws); (g$''$); (h).
The inequality $Z^{(w)}_n[1]\le Z^{(0)}_n[1]$ of (e) is proved after (g), because its derivation in (e), Step~2, (2a)--(2e), which starts from the identity of \cite[Lemma~13.8(a)]{QEDII} with $e_1$ as the time, applies (g) to single sectors in (2e).
Of (e), the clauses (f), (g) and (g$'$) use only the first part.

Outside the third step (h3) of the proof of (h), the only constants of Section~I.8 on which the assertions and their proofs depend are $\rho$, the weights $w_i$ and $\mathrm{sw}$ (line~(C0)).
The step (h3) depends in addition, through the data of the last sentence of (h), on the numbers $e$, $\eta$, $\bar{m}$ and, for the mirror copies, on $N$, $L$, $h$, $e_1$, $N_0(e)$, $n_*$, $\mathrm{per}$ (Sections~I.6.2 and~I.6.4).
The conditions on the constants under which the assertions hold are those listed in the table at the end of this chapter.
Every identity and every inequality below holds for every $q>0$, for every real $\boldsymbol{\mu}$ and for all sides admitted in the respective clause; none of them depends on a level or on the parity of a side, and $N$ enters only through the data of (h3) listed above.

\subsection*{0. Conventions of this chapter; setting and six lemmas}

\subsubsection*{0.0. Notation; classical theorems}

In this chapter $\mathrm{fl}$ denotes the flux of Section~I.7.1; its definition is recalled in Lemma~4\,(i).

In the steps (2b) and (2c) of the proof of (e) the following objects of \cite{QEDII} occur, all of them formed on the torus with $e_1$ as the time direction:
the Hilbert spaces $\mathrm{Hs}^{\mathrm{lk}}$ and $\mathrm{Hs}^{\mathrm{st}}$ of the link plane and of the site plane, the electric field $E=(E_b)$, the projection $P_G$ onto the gauge invariant vectors, the operators $K_\beta$ and $K_{\beta,S}$, and the slice operator $M_0$, which is formed with a plaquette $p_0$ and with the matrix function $\mathrm{Qs}$ \cite[Definitions~2.26 and~2.28]{QEDII};
the functions $u_t$, $\mathrm{ve}_t$ and the operators $\mathrm{Ko}_t$ of the word formula \cite[Proposition~5.43]{QEDII};
the projection $\mathrm{Pr}_{\mathrm{jz}}$ of $\mathrm{Hs}^{\mathrm{lk}}$, and of $\mathrm{Hs}^{\mathrm{st}}$, onto the summand with the flux index $\mathrm{jz}\in\mathbb{Z}$ \cite[Lemma~13.8]{QEDII}.
There $M_0$ denotes the slice operator of \cite[Definition~2.28]{QEDII}; it is not the hopping matrix $M_i(x)$ at $i=0$.
Let $L_2$ and $L_3$ be the two closed lines of links of the slice torus through a fixed site that are chosen in \cite[Lemma~13.8]{QEDII}; the vector $q\cdot\left(\mathrm{bl}_2\cdot\mathbf{1}_{L_2}+\mathrm{bl}_3\cdot\mathbf{1}_{L_3}\right)$ of that lemma is denoted $\beta(\mathrm{bl})$.
The number of sites of the slice torus \cite[Definition~2.26]{QEDII} is written out as $n_2n_3$.

The following classical results are applied.
The theorems of Fubini and Tonelli \cite[Theorem~8.8]{Rudin1987} enter in Lemma~5, in (e.3), (e.4) and in the sums over the sectors.
The formula for the integral with respect to an image measure \cite[Theorem~3.6.1]{Bogachev2007} enters in (e), Step~4, and in (g$''$).
The identities $\det O^{T}=\det O$ and $\det O^{*}=\overline{\det O}$ and the inequality $|\det O|\le\|O\|^{\dim}$ are (0.5) and (0.6) below, where they are proved.
In (e), Step~2, (2c), the following facts on operators on Hilbert spaces are applied:
the product of an operator of trace class and a bounded operator is of trace class \cite[Theorem~VI.19(b)]{ReedSimon1980};
$\operatorname{Tr}(OO')=\operatorname{Tr}(O'O)$ if $O$ is bounded and $O'$ is of trace class \cite[Theorem~VI.25(c)]{ReedSimon1980};
the trace of a non-negative operator of trace class is non-negative;
for the operator $\mathrm{Yr}:\mathrm{Hs}^{\mathrm{st}}\to\mathrm{Hs}^{\mathrm{lk}}$ of (2c) one has $\mathrm{Yr}\cdot\mathrm{Yr}^{*}\ge0$, and $\operatorname{Tr}\left(\mathrm{Yr}\cdot\mathrm{Yr}^{*}\right)$ is the square of the Hilbert--Schmidt norm of $\mathrm{Yr}^{*}$, which is finite, so that $\mathrm{Yr}\cdot\mathrm{Yr}^{*}$ is of trace class by the definition of the trace class \cite[Section~VI.6]{ReedSimon1980}.

\subsubsection*{0.1. Lattice units}

Lattice units are those of Section~I.2.5 and of Proposition~46: $\eta_0=1$, and the charge $q$ takes the place of $e_k$.
The numbers $\rho$, $w=(w_0,w_1,w_2,w_3)$ and $\mathrm{sw}$ are those of Section~I.0.5.

\emph{Period lattices.}
Let $\mathrm{Lt}(n):=\sum_{\mu}n_{\mu}\mathbb{Z}\,e_{\mu}$, let $\mathrm{Zd}(n):=\prod_\mu\{0,\dots,n_\mu-1\}$, a set of representatives of $\mathbb{Z}^3/\mathrm{Lt}(n)$, and let $|n|=n_1n_2n_3$.
For $r\in O_{\mathrm{cub}}$ write $re_\mu=\epsilon_\mu\cdot e_{\mu'}$, with $\epsilon_\mu=\pm1$ and a permutation $\mu\mapsto\mu'$, as in Proposition~50.
Then $r\,\mathrm{Lt}(n)=\sum_\mu n_\mu\mathbb{Z}\,e_{\mu'}$.
This group equals $\mathrm{Lt}(n')=\sum_\nu n'_\nu\mathbb{Z}\,e_\nu$ if and only if $n'_{\mu'}=n_\mu$ for $\mu=1,2,3$, since both numbers are the positive generator of the intersection of the respective group with $\mathbb{Z}\,e_{\mu'}$.
In this case $|n'|=|n|$.

\emph{Fermion sites} (Section~I.1.3).
A fermion site is a pair $x=(c,g)\in\mathbb{Z}^3\times O_{\mathrm{cub}}$; one writes $\pi(x)=c$, and $x+v=(c+v,g)$ for $v\in\mathbb{Z}^3$.
By (I.1.2) with $\eta_0=1$,
\[
x^{(i)}=(c,gs_i)\quad(i=0,1,2),\qquad x^{(3)}=(c+ge_1,gs_3),\qquad\hat d_i(x)=-g\hat n_i\quad(i=0,\dots,3).
\]
The facts (I.s1)--(I.s3) read as follows.
By (I.s1), $(x^{(i)})^{(i)}=x$, the map $x\mapsto x^{(i)}$ is a bijection of $\mathbb{Z}^3\times O_{\mathrm{cub}}$, and $(x+v)^{(i)}=x^{(i)}+v$.
By (I.s2), $\hat d_i(x^{(i)})=-\hat d_i(x)$.
By (I.s3), $\sum_iw_i\hat d_i(x)=0$; this is the closure relation (I.1.1).

\emph{Phases.}
A cochain $B\in C^1(n)\otimes\mathbb{C}$ is regarded as a function of the ordered pairs $(c,\tilde c)$ of nearest neighbours of $\mathbb{Z}^3$, with $B(c,c+e_\mu)=B_\mu(c)$, $B(\tilde c,c)=-B(c,\tilde c)$ and $B(c+v,\tilde c+v)=B(c,\tilde c)$ for $v\in\mathrm{Lt}(n)$; moreover $B(c,c):=0$.
As in (I.1.6), $U_B(x,x^{(i)}):=\exp\left(iqB(\pi x,\pi x^{(i)})\right)$.
This number is $1$ for $i=0,1,2$ and $\exp\left(iqB(c,c+ge_1)\right)$ for $i=3$.
Since $B$ is odd,
\begin{equation}\tag{0.1}
U_B(x^{(i)},x)=U_B(x,x^{(i)})^{-1}=U_{-B}(x,x^{(i)});\qquad\text{for real }B:\quad U_B(x^{(i)},x)=\overline{U_B(x,x^{(i)})},\quad\left|U_B(x,x^{(i)})\right|=1.
\end{equation}

\emph{Hopping matrices.}
By (I.2.1), $M_i(x)=\frac{w_i}{2}\cdot\left(\rho-\gamma\cdot\hat d_i(x)\right)$; this matrix depends only on the label $g$.
By (I.s2) and (I.s3),
\begin{equation}\tag{0.2}
M_i(x^{(i)})=\frac{w_i}{2}\cdot\left(\rho+\gamma\cdot\hat d_i(x)\right),\qquad\sum_{i=0}^{3}M_i(x)=\frac{\rho}{2}\cdot\mathrm{sw}\cdot1_4;
\end{equation}
the second identity is (I.m1).

\emph{The operator.}
Let $(R_i(B)f)(x):=U_B(x,x^{(i)})\cdot f(x^{(i)})$, as in (I.2.2).
The letters $M_i$ and $\gamma\cdot\hat d_i$ also denote the multiplication operators $f\mapsto M_i(x)f(x)$ and $f\mapsto\left(\gamma\cdot\hat d_i(x)\right)f(x)$.
In this notation (I.2.3), (I.2.5) and (I.2.4) read
\begin{equation}\tag{0.3}
\begin{aligned}
(D_Bf)(x)&=-\sum_{i=0}^{3}M_i(x)\left[U_B(x,x^{(i)})\cdot f(x^{(i)})-f(x)\right],\\
D_B+\boldsymbol{\mu}&=M_{\mathrm{ch}}-\sum_iM_i\cdot R_i(B),\qquad D_B=\frac{\rho}{2}\cdot\nabla_B^{*}\nabla_B+\frac{1}{2}\cdot J_B,\\
\nabla_B^{*}\nabla_B&:=\sum_iw_i\cdot\left(1-R_i(B)\right),\qquad J_B:=\sum_iw_i\cdot(\gamma\cdot\hat d_i)\cdot R_i(B),\qquad M_{\mathrm{ch}}=\boldsymbol{\mu}+\frac{\rho}{2}\cdot\mathrm{sw}.
\end{aligned}
\end{equation}
The second and the third form follow from the first by (0.2), because
\[
\sum_iM_i(x)f(x)=\frac{\rho}{2}\cdot\mathrm{sw}\cdot f(x)=\frac{\rho}{2}\sum_iw_i\cdot f(x),\qquad -M_i(x)=-\frac{\rho}{2}\cdot w_i+\frac{1}{2}\cdot w_i\cdot\left(\gamma\cdot\hat d_i(x)\right).
\]

\emph{Spinors.}
Recall that
\[
\mathfrak{S}^{\mathrm{ch}}(1;n)=\left\{f:\mathbb{Z}^3\times O_{\mathrm{cub}}\to\mathbb{C}^4:\ f(c+n_\mu e_\mu,g)=-f(c,g),\ \mu=1,2,3\right\}.
\]
An element is determined by its restriction to $\mathrm{Zd}(n)\times O_{\mathrm{cub}}$, and this restriction is arbitrary: there are $48|n|$ sites with four components each, and the dimension is $192|n|$.
The operators $R_i(B)$, $M_i$ and $\gamma\cdot\hat d_i$ map $\mathfrak{S}^{\mathrm{ch}}(1;n)$ into itself, because $(x+v)^{(i)}=x^{(i)}+v$, a translation does not change the label, and $B$ is $\mathrm{Lt}(n)$-periodic.
Hence $D_B$ maps $\mathfrak{S}^{\mathrm{ch}}(1;n)$ into itself.

\emph{Pairings} (Section~I.1.6 with $\eta_0=1$).
\[
\langle f',f\rangle^{\mathrm{ch}}=\frac{1}{48}\sum_{c\in\mathrm{Zd}(n)}\sum_{g\in O_{\mathrm{cub}}}f'(c,g)\cdot f(c,g),\qquad
(f',f)^{\mathrm{ch}}=\frac{1}{48}\sum_{c\in\mathrm{Zd}(n)}\sum_{g\in O_{\mathrm{cub}}}\overline{f'(c,g)}\cdot f(c,g).
\]
The summands are $\mathrm{Lt}(n)$-periodic in $c$, since the two factors $-1$ cancel.
A bijection $\mathrm{bj}$ of $\mathbb{Z}^3\times O_{\mathrm{cub}}$ with $\mathrm{bj}(x+v)=\mathrm{bj}(x)+v$ for $v\in\mathrm{Lt}(n)$ induces a bijection of $(\mathbb{Z}^3\times O_{\mathrm{cub}})/\mathrm{Lt}(n)$.
Hence every $\mathrm{Lt}(n)$-periodic function $\mathrm{Fs}$ on $\mathbb{Z}^3\times O_{\mathrm{cub}}$ satisfies
\begin{equation}\tag{0.4}
\sum_{x\in\mathrm{Zd}(n)\times O_{\mathrm{cub}}}\mathrm{Fs}(\mathrm{bj}\,x)=\sum_{x\in\mathrm{Zd}(n)\times O_{\mathrm{cub}}}\mathrm{Fs}(x).
\end{equation}
The adjoint $O^{*}$ and the transpose $O^{T}$ of an operator $O$ on $\mathfrak{S}^{\mathrm{ch}}(1;n)$ are formed with these pairings: $(f',Of)^{\mathrm{ch}}=(O^{*}f',f)^{\mathrm{ch}}$ and $\langle f',Of\rangle^{\mathrm{ch}}=\langle O^{T}f',f\rangle^{\mathrm{ch}}$.
The weight $\frac{1}{48}$ stands on both sides of these relations; so $O^{*}$ and $O^{T}$ are also the adjoint and the transpose for the counting pairings.
For $x\in\mathrm{Zd}(n)\times O_{\mathrm{cub}}$ and $\beta=1,\dots,4$ let $\delta_{x,\beta}$ be the element of $\mathfrak{S}^{\mathrm{ch}}(1;n)$ whose restriction to $\mathrm{Zd}(n)\times O_{\mathrm{cub}}$ is the $\beta$-th unit vector of $\mathbb{C}^4$ at $x$ and $0$ at the other sites.
In this basis the two pairings are $\frac{1}{48}$ times the standard bilinear form and the standard sesquilinear form.
Therefore the matrix of $O^{T}$ is the transpose of the matrix of $O$, and the matrix of $O^{*}$ is its conjugate transpose, so that
\begin{equation}\tag{0.5}
\det O^{T}=\det O,\qquad\det O^{*}=\overline{\det O}.
\end{equation}
Let $\|f\|:=\left((f,f)^{\mathrm{ch}}\right)^{\frac{1}{2}}$ be the $\ell^2$ norm, and $\|O\|$ the operator norm for it.
For every linear map $O$ of a space of finite dimension $\dim$ with a scalar product, each eigenvalue has modulus at most $\|O\|$, and $\det O$ is the product of the $\dim$ eigenvalues, counted with their algebraic multiplicities; hence
\begin{equation}\tag{0.6}
|\det O|\le\|O\|^{\dim}.
\end{equation}

\subsubsection*{0.2. Cells, orientations, pull-backs; the last sentence of the preamble of Proposition~50}

\emph{Cells.}
For $j\in\{0,1,2,3\}$ a $j$-cell of $\mathbb{Z}^3$ is a pair $\boldsymbol{\sigma}=(c;\mathrm{Iv})$ with $c\in\mathbb{Z}^3$ and a subset $\mathrm{Iv}$ of $\{1,2,3\}$ with $j$ elements $i_1<\dots<i_j$.
It stands for the closed set $c+\sum_{i\in\mathrm{Iv}}[0,1]\,e_i$ and carries the orientation $\mathrm{or}(\boldsymbol{\sigma}):=e_{i_1}\wedge\dots\wedge e_{i_j}\in\Lambda^j\mathbb{R}^3$, with $\mathrm{or}(\boldsymbol{\sigma}):=1$ for $j=0$.
For $j=0,1,2,3$ the $j$-cells are the sites, the links $(c;\{\mu\})$, the plaquettes $(c;\{\mu,\nu\})=(c;\mu\nu)$ with $\mu<\nu$, and the cubes.
Put $\boldsymbol{\sigma}+v:=(c+v;\mathrm{Iv})$.
A $j$-cochain of $\mathbb{Z}^3$ is a real function of the $j$-cells.
The space $C^j(n)$ consists of the $j$-cochains $u$ with $u(\boldsymbol{\sigma}+v)=u(\boldsymbol{\sigma})$ for $v\in\mathrm{Lt}(n)$ (Section~G6.A1), and $\langle u,\tilde u\rangle$ is the sum of $u(\boldsymbol{\sigma})\cdot\tilde u(\boldsymbol{\sigma})$ over the classes of the $j$-cells modulo $\mathrm{Lt}(n)$.
For all $n_\mu\ge1$ there are $3|n|$ classes of links and $3|n|$ classes of plaquettes.

\emph{Faces, coboundary.}
Let $\boldsymbol{\sigma}^{+}=(c;\mathrm{Iv}^{+})$ be a $(j+1)$-cell.
For $i\in\mathrm{Iv}^{+}$ let $\#(i)$ be the number of the elements of $\mathrm{Iv}^{+}$ that are smaller than $i$.
The faces of $\boldsymbol{\sigma}^{+}$ are the $2(j+1)$ cells $(c;\mathrm{Iv}^{+}\setminus\{i\})$, with the outward unit normal $\hat n=-e_i$, and $(c+e_i;\mathrm{Iv}^{+}\setminus\{i\})$, with the outward unit normal $\hat n=+e_i$, where $i\in\mathrm{Iv}^{+}$.
For a face $\boldsymbol{\sigma}$ of $\boldsymbol{\sigma}^{+}$ the number $[\boldsymbol{\sigma}^{+}:\boldsymbol{\sigma}]\in\{1,-1\}$ is defined by
\begin{equation}\tag{0.7}
\hat n\wedge\mathrm{or}(\boldsymbol{\sigma})=[\boldsymbol{\sigma}^{+}:\boldsymbol{\sigma}]\cdot\mathrm{or}(\boldsymbol{\sigma}^{+}).
\end{equation}
Since $e_i\wedge\mathrm{or}(c;\mathrm{Iv}^{+}\setminus\{i\})=(-1)^{\#(i)}\cdot\mathrm{or}(\boldsymbol{\sigma}^{+})$, one has $[\boldsymbol{\sigma}^{+}:\boldsymbol{\sigma}]=(-1)^{\#(i)}$ for the face at $c+e_i$ and $[\boldsymbol{\sigma}^{+}:\boldsymbol{\sigma}]=-(-1)^{\#(i)}$ for the face at $c$.
Hence
\[
(du)(\boldsymbol{\sigma}^{+}):=\sum_{\boldsymbol{\sigma}\text{ a face of }\boldsymbol{\sigma}^{+}}[\boldsymbol{\sigma}^{+}:\boldsymbol{\sigma}]\cdot u(\boldsymbol{\sigma})=\sum_{i\in\mathrm{Iv}^{+}}(-1)^{\#(i)}\left[u(c+e_i;\mathrm{Iv}^{+}\setminus\{i\})-u(c;\mathrm{Iv}^{+}\setminus\{i\})\right].
\]
For $\mathrm{Iv}^{+}=\{\mu\}$ and a zero-cochain $u=\boldsymbol{\lambda}$ this is $\boldsymbol{\lambda}(c+e_\mu)-\boldsymbol{\lambda}(c)$.
For $\mathrm{Iv}^{+}=\{\mu,\nu\}$, $\mu<\nu$, and a one-cochain $u=A$ it is $\left[A_\nu(c+e_\mu)-A_\nu(c)\right]-\left[A_\mu(c+e_\nu)-A_\mu(c)\right]$.
For $\mathrm{Iv}^{+}=\{1,2,3\}$ and a two-cochain $u=F$ it is
\[
\left[F_{23}(c+e_1)-F_{23}(c)\right]-\left[F_{13}(c+e_2)-F_{13}(c)\right]+\left[F_{12}(c+e_3)-F_{12}(c)\right].
\]
These are the three formulas for $d$ in Section~I.1.1.

\emph{The identity $dd=0$.}
For a function $\boldsymbol{\lambda}$ of $c\in\mathbb{Z}^3$ let $(\partial_\mu\boldsymbol{\lambda})(c)=\boldsymbol{\lambda}(c+e_\mu)-\boldsymbol{\lambda}(c)$ (Section~I.1.1 with $\eta_0=1$).
The function
\[
(\partial_\mu\partial_\nu\boldsymbol{\lambda})(c)=\boldsymbol{\lambda}(c+e_\mu+e_\nu)-\boldsymbol{\lambda}(c+e_\mu)-\boldsymbol{\lambda}(c+e_\nu)+\boldsymbol{\lambda}(c)
\]
is symmetric in $\mu,\nu$; hence $\partial_\mu\partial_\nu=\partial_\nu\partial_\mu$.
The three formulas read $(d\boldsymbol{\lambda})_\mu=\partial_\mu\boldsymbol{\lambda}$, $(dA)_{\mu\nu}=\partial_\mu A_\nu-\partial_\nu A_\mu$ ($\mu<\nu$) and $dF=\partial_1F_{23}-\partial_2F_{13}+\partial_3F_{12}$.
Consequently $(dd\boldsymbol{\lambda})_{\mu\nu}=\partial_\mu\partial_\nu\boldsymbol{\lambda}-\partial_\nu\partial_\mu\boldsymbol{\lambda}=0$, and
\[
ddA=\partial_1\left(\partial_2A_3-\partial_3A_2\right)-\partial_2\left(\partial_1A_3-\partial_3A_1\right)+\partial_3\left(\partial_1A_2-\partial_2A_1\right)=0,
\]
the six terms cancelling in pairs.
Thus $dd=0$ on the cochains of $\mathbb{Z}^3$, and therefore on $C^j(n)$ for all $n_\mu\ge1$.

\emph{Action of $g_{\mathrm{w}}=(a,r)$.}
The map $y\mapsto ry+a$ of $\mathbb{R}^3$ maps the cell $\boldsymbol{\sigma}=(c;\mathrm{Iv})$ onto the set $rc+a+\sum_{i\in\mathrm{Iv}}[0,1]\,\epsilon_i\cdot e_{i'}$, which is the cell
\[
g_{\mathrm{w}}\boldsymbol{\sigma}:=\left(rc+a-\sum_{i\in\mathrm{Iv},\ \epsilon_i=-1}e_{i'}\,;\ \mathrm{Iv}'\right),\qquad\mathrm{Iv}':=\{i':i\in\mathrm{Iv}\}.
\]
The sign $\epsilon(g_{\mathrm{w}},\boldsymbol{\sigma})\in\{1,-1\}$ is defined by $(\Lambda^jr)\,\mathrm{or}(\boldsymbol{\sigma})=\epsilon(g_{\mathrm{w}},\boldsymbol{\sigma})\cdot\mathrm{or}(g_{\mathrm{w}}\boldsymbol{\sigma})$.
It equals $\prod_{i\in\mathrm{Iv}}\epsilon_i$ times the sign of the permutation that sorts $(i_1',\dots,i_j')$; in particular it depends only on $r$ and $\mathrm{Iv}$.
As in Section~I.4.7, $(g_{\mathrm{w}}^{*}u)(\boldsymbol{\sigma}):=\epsilon(g_{\mathrm{w}},\boldsymbol{\sigma})\cdot u(g_{\mathrm{w}}\boldsymbol{\sigma})$.
Moreover $g_{\mathrm{w}}(\boldsymbol{\sigma}+v)=g_{\mathrm{w}}\boldsymbol{\sigma}+rv$.

\begin{named}{Lemma 1}
Let $g_{\mathrm{w}}=(a,r)$ and $\tilde g_{\mathrm{w}}=(\tilde a,\tilde r)$ be elements of $W_{\mathrm{af}}$, and let $j\in\{0,1,2,3\}$.
\begin{itemize}
\item[(i)] $\epsilon(g_{\mathrm{w}}\tilde g_{\mathrm{w}},\boldsymbol{\sigma})=\epsilon(g_{\mathrm{w}},\tilde g_{\mathrm{w}}\boldsymbol{\sigma})\cdot\epsilon(\tilde g_{\mathrm{w}},\boldsymbol{\sigma})$ and $(g_{\mathrm{w}}\tilde g_{\mathrm{w}})^{*}=\tilde g_{\mathrm{w}}^{*}g_{\mathrm{w}}^{*}$.
Consequently $g_{\mathrm{w}}^{*}$ is a bijection of the $j$-cochains of $\mathbb{Z}^3$ with the inverse $g_{\mathrm{w}*}=\left(g_{\mathrm{w}}^{-1}\right)^{*}$, and $\epsilon(g_{\mathrm{w}}^{-1},g_{\mathrm{w}}\boldsymbol{\sigma})=\epsilon(g_{\mathrm{w}},\boldsymbol{\sigma})$.
\item[(ii)] $d\,g_{\mathrm{w}}^{*}=g_{\mathrm{w}}^{*}d$ and $d\,g_{\mathrm{w}*}=g_{\mathrm{w}*}d$ on the cochains of $\mathbb{Z}^3$.
\item[(iii)] If $r\,\mathrm{Lt}(n)=\mathrm{Lt}(n')$, then $g_{\mathrm{w}}^{*}$ maps $C^j(n')$ onto $C^j(n)$ and $g_{\mathrm{w}*}$ maps $C^j(n)$ onto $C^j(n')$; the two maps are inverse to each other, and they map integer valued cochains onto integer valued cochains.
\item[(iv)] In the setting of (iii), for $u,\tilde u\in C^j(n)$ and $u',\tilde u'\in C^j(n')$,
\[
\langle g_{\mathrm{w}}^{*}u',g_{\mathrm{w}}^{*}\tilde u'\rangle\ (\text{on }n)=\langle u',\tilde u'\rangle\ (\text{on }n'),\qquad
\langle g_{\mathrm{w}*}u,g_{\mathrm{w}*}\tilde u\rangle\ (\text{on }n')=\langle u,\tilde u\rangle\ (\text{on }n),
\]
\[
\langle u,g_{\mathrm{w}}^{*}u'\rangle\ (\text{on }n)=\langle g_{\mathrm{w}*}u,u'\rangle\ (\text{on }n').
\]
\item[(v)] $d^{*}g_{\mathrm{w}*}=g_{\mathrm{w}*}d^{*}$, where $d^{*}$ is the adjoint of $d$ for the counting pairings on $n$ and on $n'$.
\item[(vi)] For a one-cochain $A'$ and nearest neighbours $c,\tilde c$ of $\mathbb{Z}^3$ one has $(g_{\mathrm{w}}^{*}A')(c,\tilde c)=A'(g_{\mathrm{w}}c,g_{\mathrm{w}}\tilde c)$; hence $(g_{\mathrm{w}*}A)(c,\tilde c)=A(g_{\mathrm{w}}^{-1}c,g_{\mathrm{w}}^{-1}\tilde c)$.
\end{itemize}
All maps are extended complex linearly to complex cochains; (i)--(iii), (v) and (vi) hold for the extensions, and (iv) holds for the bilinear extension of $\langle\,\cdot\,,\,\cdot\,\rangle$.
\end{named}

\begin{proof}
(i) As sets, $(g_{\mathrm{w}}\tilde g_{\mathrm{w}})\boldsymbol{\sigma}=g_{\mathrm{w}}(\tilde g_{\mathrm{w}}\boldsymbol{\sigma})$, and
\[
\left(\Lambda^j(r\tilde r)\right)\mathrm{or}(\boldsymbol{\sigma})=(\Lambda^jr)\left(\epsilon(\tilde g_{\mathrm{w}},\boldsymbol{\sigma})\cdot\mathrm{or}(\tilde g_{\mathrm{w}}\boldsymbol{\sigma})\right)=\epsilon(\tilde g_{\mathrm{w}},\boldsymbol{\sigma})\cdot\epsilon(g_{\mathrm{w}},\tilde g_{\mathrm{w}}\boldsymbol{\sigma})\cdot\mathrm{or}(g_{\mathrm{w}}\tilde g_{\mathrm{w}}\boldsymbol{\sigma}).
\]
Therefore
\[
(\tilde g_{\mathrm{w}}^{*}g_{\mathrm{w}}^{*}u)(\boldsymbol{\sigma})=\epsilon(\tilde g_{\mathrm{w}},\boldsymbol{\sigma})\cdot(g_{\mathrm{w}}^{*}u)(\tilde g_{\mathrm{w}}\boldsymbol{\sigma})=\epsilon(\tilde g_{\mathrm{w}},\boldsymbol{\sigma})\cdot\epsilon(g_{\mathrm{w}},\tilde g_{\mathrm{w}}\boldsymbol{\sigma})\cdot u(g_{\mathrm{w}}\tilde g_{\mathrm{w}}\boldsymbol{\sigma})=\left((g_{\mathrm{w}}\tilde g_{\mathrm{w}})^{*}u\right)(\boldsymbol{\sigma}).
\]
For the unit element one has $\epsilon=1$, so that $g_{\mathrm{w}}^{*}g_{\mathrm{w}*}$ and $g_{\mathrm{w}*}g_{\mathrm{w}}^{*}$ are the identity.
With $g_{\mathrm{w}}^{-1}$ in the place of $g_{\mathrm{w}}$ and $g_{\mathrm{w}}$ in the place of $\tilde g_{\mathrm{w}}$ the first identity gives $1=\epsilon(g_{\mathrm{w}}^{-1},g_{\mathrm{w}}\boldsymbol{\sigma})\cdot\epsilon(g_{\mathrm{w}},\boldsymbol{\sigma})$.

(ii) Let $\boldsymbol{\sigma}$ be a face of $\boldsymbol{\sigma}^{+}$ with the outward unit normal $\hat n$.
The vector $\hat n$ is the unit vector that is parallel to $\boldsymbol{\sigma}^{+}$, orthogonal to $\boldsymbol{\sigma}$, and such that the open half-line $y+(0,\infty)\,\hat n$ does not meet $\boldsymbol{\sigma}^{+}$ for any $y\in\boldsymbol{\sigma}$.
The isometry $y\mapsto ry+a$ maps $\boldsymbol{\sigma}^{+}$ onto $g_{\mathrm{w}}\boldsymbol{\sigma}^{+}$, $\boldsymbol{\sigma}$ onto $g_{\mathrm{w}}\boldsymbol{\sigma}$, and $y+(0,\infty)\,\hat n$ onto $(ry+a)+(0,\infty)\,r\hat n$.
Hence $g_{\mathrm{w}}\boldsymbol{\sigma}$ is a face of $g_{\mathrm{w}}\boldsymbol{\sigma}^{+}$ with the outward unit normal $r\hat n$, and $\boldsymbol{\sigma}\mapsto g_{\mathrm{w}}\boldsymbol{\sigma}$ is a bijection of the faces of $\boldsymbol{\sigma}^{+}$ onto the faces of $g_{\mathrm{w}}\boldsymbol{\sigma}^{+}$.
In terms of the list of the faces given above this reads as follows.
Let $\boldsymbol{\sigma}^{+}=(c;\mathrm{Iv}^{+})$, $g_{\mathrm{w}}\boldsymbol{\sigma}^{+}=(c^{*};\mathrm{Iv}^{+\prime})$ and $\boldsymbol{\sigma}=(c+\mathrm{ti}\cdot e_{\mathfrak{l}};\mathrm{Iv}^{+}\setminus\{\mathfrak{l}\})$ with $\mathfrak{l}\in\mathrm{Iv}^{+}$ and $\mathrm{ti}\in\{0,1\}$.
By the formula for $g_{\mathrm{w}}\boldsymbol{\sigma}$, the first entry of $g_{\mathrm{w}}\boldsymbol{\sigma}$ is
\[
r(c+\mathrm{ti}\cdot e_{\mathfrak{l}})+a-\sum_{i\in\mathrm{Iv}^{+}\setminus\{\mathfrak{l}\},\ \epsilon_i=-1}e_{i'}=c^{*}+\left(\mathbf{1}[\epsilon_{\mathfrak{l}}=-1]+\mathrm{ti}\cdot\epsilon_{\mathfrak{l}}\right)\cdot e_{\mathfrak{l}'},
\]
so that $g_{\mathrm{w}}\boldsymbol{\sigma}=(c^{*}+\mathrm{ti}'\cdot e_{\mathfrak{l}'};\mathrm{Iv}^{+\prime}\setminus\{\mathfrak{l}'\})$ with $\mathrm{ti}'=\mathrm{ti}$ if $\epsilon_{\mathfrak{l}}=1$ and $\mathrm{ti}'=1-\mathrm{ti}$ if $\epsilon_{\mathfrak{l}}=-1$.
Thus $\boldsymbol{\sigma}\mapsto g_{\mathrm{w}}\boldsymbol{\sigma}$ maps the $2(j+1)$ faces of $\boldsymbol{\sigma}^{+}$ bijectively onto those of $g_{\mathrm{w}}\boldsymbol{\sigma}^{+}$, and it maps the normal $(2\,\mathrm{ti}-1)\cdot e_{\mathfrak{l}}$ to $(2\,\mathrm{ti}'-1)\cdot e_{\mathfrak{l}'}=r(2\,\mathrm{ti}-1)\cdot e_{\mathfrak{l}}$.
Applying $\Lambda^{j+1}r$ to (0.7) one obtains
\begin{align*}
[\boldsymbol{\sigma}^{+}:\boldsymbol{\sigma}]\cdot\epsilon(g_{\mathrm{w}},\boldsymbol{\sigma}^{+})\cdot\mathrm{or}(g_{\mathrm{w}}\boldsymbol{\sigma}^{+})
&=(r\hat n)\wedge(\Lambda^jr)\,\mathrm{or}(\boldsymbol{\sigma})
=\epsilon(g_{\mathrm{w}},\boldsymbol{\sigma})\cdot(r\hat n)\wedge\mathrm{or}(g_{\mathrm{w}}\boldsymbol{\sigma})\\
&=\epsilon(g_{\mathrm{w}},\boldsymbol{\sigma})\cdot[g_{\mathrm{w}}\boldsymbol{\sigma}^{+}:g_{\mathrm{w}}\boldsymbol{\sigma}]\cdot\mathrm{or}(g_{\mathrm{w}}\boldsymbol{\sigma}^{+}).
\end{align*}
Hence $[\boldsymbol{\sigma}^{+}:\boldsymbol{\sigma}]\cdot\epsilon(g_{\mathrm{w}},\boldsymbol{\sigma})=\epsilon(g_{\mathrm{w}},\boldsymbol{\sigma}^{+})\cdot[g_{\mathrm{w}}\boldsymbol{\sigma}^{+}:g_{\mathrm{w}}\boldsymbol{\sigma}]$, and, the sums being over the faces $\boldsymbol{\sigma}$ of $\boldsymbol{\sigma}^{+}$,
\begin{align*}
(d\,g_{\mathrm{w}}^{*}u)(\boldsymbol{\sigma}^{+})
&=\sum_{\boldsymbol{\sigma}}[\boldsymbol{\sigma}^{+}:\boldsymbol{\sigma}]\cdot\epsilon(g_{\mathrm{w}},\boldsymbol{\sigma})\cdot u(g_{\mathrm{w}}\boldsymbol{\sigma})
=\epsilon(g_{\mathrm{w}},\boldsymbol{\sigma}^{+})\sum_{\boldsymbol{\sigma}}[g_{\mathrm{w}}\boldsymbol{\sigma}^{+}:g_{\mathrm{w}}\boldsymbol{\sigma}]\cdot u(g_{\mathrm{w}}\boldsymbol{\sigma})\\
&=\epsilon(g_{\mathrm{w}},\boldsymbol{\sigma}^{+})\cdot(du)(g_{\mathrm{w}}\boldsymbol{\sigma}^{+})
=(g_{\mathrm{w}}^{*}du)(\boldsymbol{\sigma}^{+}).
\end{align*}
The second identity is the first one for $g_{\mathrm{w}}^{-1}$.

(iii) Let $u'$ be $\mathrm{Lt}(n')$-periodic and $v\in\mathrm{Lt}(n)$.
Then $(g_{\mathrm{w}}^{*}u')(\boldsymbol{\sigma}+v)=\epsilon(g_{\mathrm{w}},\boldsymbol{\sigma})\cdot u'(g_{\mathrm{w}}\boldsymbol{\sigma}+rv)=(g_{\mathrm{w}}^{*}u')(\boldsymbol{\sigma})$, because $\epsilon(g_{\mathrm{w}},\boldsymbol{\sigma}+v)=\epsilon(g_{\mathrm{w}},\boldsymbol{\sigma})$ and $rv\in\mathrm{Lt}(n')$.
Further $g_{\mathrm{w}}^{-1}=(-r^{-1}a,r^{-1})$ and $r^{-1}\mathrm{Lt}(n')=\mathrm{Lt}(n)$; so $g_{\mathrm{w}*}$ maps $C^j(n)$ into $C^j(n')$.
By (i) the two maps are inverse to each other; hence both are onto.
The values of $g_{\mathrm{w}}^{*}u'$ are values of $u'$ up to signs.

(iv) By $g_{\mathrm{w}}(\boldsymbol{\sigma}+v)=g_{\mathrm{w}}\boldsymbol{\sigma}+rv$ and $r\,\mathrm{Lt}(n)=\mathrm{Lt}(n')$, the map $\boldsymbol{\sigma}\mapsto g_{\mathrm{w}}\boldsymbol{\sigma}$ induces a bijection of the classes of the $j$-cells modulo $\mathrm{Lt}(n)$ onto the classes modulo $\mathrm{Lt}(n')$.
Hence
\[
\langle g_{\mathrm{w}}^{*}u',g_{\mathrm{w}}^{*}\tilde u'\rangle=\sum_{\boldsymbol{\sigma}}\epsilon(g_{\mathrm{w}},\boldsymbol{\sigma})^2\cdot u'(g_{\mathrm{w}}\boldsymbol{\sigma})\cdot\tilde u'(g_{\mathrm{w}}\boldsymbol{\sigma})=\langle u',\tilde u'\rangle.
\]
The second identity is the first one for $u'=g_{\mathrm{w}*}u$ and $\tilde u'=g_{\mathrm{w}*}\tilde u$.
The third one is the first one for $\tilde u':=g_{\mathrm{w}*}u$, read from right to left: $\langle g_{\mathrm{w}*}u,u'\rangle=\langle g_{\mathrm{w}}^{*}g_{\mathrm{w}*}u,g_{\mathrm{w}}^{*}u'\rangle=\langle u,g_{\mathrm{w}}^{*}u'\rangle$.

(v) For $u\in C^{j+1}(n)$ and $u'\in C^j(n')$, by (iv) and (ii),
\[
\langle d^{*}g_{\mathrm{w}*}u,u'\rangle=\langle g_{\mathrm{w}*}u,du'\rangle=\langle u,g_{\mathrm{w}}^{*}du'\rangle=\langle u,d\,g_{\mathrm{w}}^{*}u'\rangle=\langle d^{*}u,g_{\mathrm{w}}^{*}u'\rangle=\langle g_{\mathrm{w}*}d^{*}u,u'\rangle.
\]

(vi) Let $\tilde c=c+e_\mu$.
If $\epsilon_\mu=1$, then $g_{\mathrm{w}}(c;\{\mu\})=(g_{\mathrm{w}}c;\{\mu'\})$ and $\epsilon(g_{\mathrm{w}},(c;\{\mu\}))=1$, so that $(g_{\mathrm{w}}^{*}A')_\mu(c)=A'_{\mu'}(g_{\mathrm{w}}c)=A'(g_{\mathrm{w}}c,g_{\mathrm{w}}c+e_{\mu'})$.
If $\epsilon_\mu=-1$, then $g_{\mathrm{w}}(c;\{\mu\})=(g_{\mathrm{w}}c-e_{\mu'};\{\mu'\})$ and $\epsilon(g_{\mathrm{w}},(c;\{\mu\}))=-1$, so that $(g_{\mathrm{w}}^{*}A')_\mu(c)=-A'(g_{\mathrm{w}}c-e_{\mu'},g_{\mathrm{w}}c)=A'(g_{\mathrm{w}}c,g_{\mathrm{w}}c-e_{\mu'})$.
In both cases the second argument is $g_{\mathrm{w}}c+re_\mu=g_{\mathrm{w}}\tilde c$.
Both sides of the identity are odd under the exchange of $c$ and $\tilde c$, which settles the case $\tilde c=c-e_\mu$.
\end{proof}

Parts (iii), (iv) and (ii) of Lemma~1 are the assertions of the last sentence of the preamble of Proposition~50.

\subsubsection*{0.3. Fermion sites; spin matrices}

\begin{named}{Lemma 2}[Section~I.4.7]
The map $g_{\mathrm{w}}(c,g):=(rc+a,rg)$ is a bijection of $\mathbb{Z}^3\times O_{\mathrm{cub}}$ with the inverse $g_{\mathrm{w}}^{-1}(c,g)=\left(r^{-1}(c-a),r^{-1}g\right)$, and for $i=0,\dots,3$ and $v\in\mathbb{Z}^3$
\[
\pi(g_{\mathrm{w}}x)=g_{\mathrm{w}}\pi(x),\qquad g_{\mathrm{w}}(x+v)=g_{\mathrm{w}}x+rv,\qquad(g_{\mathrm{w}}x)^{(i)}=g_{\mathrm{w}}(x^{(i)}),\qquad\hat d_i(g_{\mathrm{w}}x)=r\,\hat d_i(x).
\]
\end{named}

\begin{proof}
Let $x=(c,g)$.
For $i\le2$ one has $(g_{\mathrm{w}}x)^{(i)}=(rc+a,rgs_i)=g_{\mathrm{w}}(c,gs_i)$.
Further $(g_{\mathrm{w}}x)^{(3)}=(rc+a+(rg)e_1,rgs_3)=g_{\mathrm{w}}(c+ge_1,gs_3)$, and $\hat d_i(g_{\mathrm{w}}x)=-rg\hat n_i=r\,\hat d_i(x)$.
The two remaining relations hold by definition.
\end{proof}

\begin{named}{Lemma 3}[Sections~I.4.1 and~I.4.7]
For a unit vector $\hat n$ of $\mathbb{R}^3$ let $r_{\hat n}:=1-2\hat n\hat n^{T}$ and $S_{\mathrm{nv}}:=(\gamma\cdot\hat n)\cdot\gamma_5$.
Then $S_{\mathrm{nv}}$ is invertible, and $S_{\mathrm{nv}}\cdot(\gamma\cdot v)\cdot S_{\mathrm{nv}}^{-1}=\gamma\cdot(r_{\hat n}v)$ for $v\in\mathbb{C}^3$.
For $r\in O_{\mathrm{cub}}$ let $r=r_{\hat n_{(1)}}\cdots r_{\hat n_{(\mathrm{wl})}}$, $\mathrm{wl}\ge0$, be a representation of $r$ as a product of reflections $r_{\hat n}$; such a representation exists, because $s_0=r_{\hat n_0}$, $s_1=r_{\hat n_1}$ and $s_2=r_{\hat n_2}$ generate $O_{\mathrm{cub}}$ (Section~I.1.2).
Let $S_{\mathrm{gw}}$ be the product, in this order, of the matrices $S_{\mathrm{nv}}$ formed with $\hat n=\hat n_{(1)},\dots,\hat n_{(\mathrm{wl})}$, and let $S_{\mathrm{gw}}:=1$ for $\mathrm{wl}=0$.
Then
\begin{equation}\tag{0.8}
S_{\mathrm{gw}}\cdot(\gamma\cdot v)\cdot S_{\mathrm{gw}}^{-1}=\gamma\cdot(rv)\qquad\text{for all }v\in\mathbb{C}^3.
\end{equation}
\end{named}

\begin{proof}
By (I.c1), $(\gamma\cdot\hat n)^2=1=\gamma_5^2$; hence $S_{\mathrm{nv}}^{-1}=\gamma_5\cdot(\gamma\cdot\hat n)$.
Moreover $\gamma_5\cdot(\gamma\cdot v)\cdot\gamma_5=-\gamma\cdot v$ and $(\gamma\cdot v)\cdot(\gamma\cdot\hat n)=2(v\cdot\hat n)-(\gamma\cdot\hat n)\cdot(\gamma\cdot v)$, where $v\cdot\hat n$ denotes the bilinear product.
Therefore
\[
S_{\mathrm{nv}}\cdot(\gamma\cdot v)\cdot S_{\mathrm{nv}}^{-1}=-(\gamma\cdot\hat n)\cdot(\gamma\cdot v)\cdot(\gamma\cdot\hat n)=-(\gamma\cdot\hat n)\cdot\left[2(v\cdot\hat n)-(\gamma\cdot\hat n)\cdot(\gamma\cdot v)\right]=\gamma\cdot v-2(v\cdot\hat n)\cdot(\gamma\cdot\hat n)=\gamma\cdot(r_{\hat n}v).
\]
The identity (0.8) follows by induction on $\mathrm{wl}$: if $S'$ is the product of the matrices $S_{\mathrm{nv}}$ formed with $\hat n_{(2)},\dots,\hat n_{(\mathrm{wl})}$, then $S_{\mathrm{gw}}=S_{\mathrm{nv}}\cdot S'$ with $\hat n=\hat n_{(1)}$, and
\[
S_{\mathrm{nv}}\cdot\left(\gamma\cdot\left(r_{\hat n_{(2)}}\cdots r_{\hat n_{(\mathrm{wl})}}v\right)\right)\cdot S_{\mathrm{nv}}^{-1}=\gamma\cdot\left(r_{\hat n_{(1)}}r_{\hat n_{(2)}}\cdots r_{\hat n_{(\mathrm{wl})}}v\right).
\]
\end{proof}

Of the matrix $S_{\mathrm{gw}}$, the proof of (d) uses only (0.8) and the invertibility.
Hence (d) holds for every representation of $r$ of this kind.

\subsubsection*{0.4. The flux}

Let $\mathrm{Mz}(n):=\left\{\mathfrak{z}\in C^2(n;\mathbb{Z}):d\mathfrak{z}=0\right\}$.
For $\mu<\nu$ let $\mathfrak{l}$ be the third direction, let $\mathbf{1}_{\mu\nu}\in C^2(n)$ be the cochain with the value $1$ on the plaquettes $(c;\mu\nu)$ and $0$ on the other plaquettes, and put $\mathbf{1}_{\nu\mu}:=-\mathbf{1}_{\mu\nu}$.
For $\mathrm{jz}\in\mathbb{Z}$ let $\Pi_{\mu\nu}(\mathrm{jz})$ be the set of the $n_\mu n_\nu$ classes modulo $\mathrm{Lt}(n)$ that contain a plaquette $(c;\mu\nu)$ with $c_{\mathfrak{l}}=\mathrm{jz}$; this is a plane in the sense of Section~I.7.1.
One has $\Pi_{\mu\nu}(\mathrm{jz}+n_{\mathfrak{l}})=\Pi_{\mu\nu}(\mathrm{jz})$.

\begin{named}{Lemma 4}
Let $n_\mu,n'_\mu\ge1$, and let $g_{\mathrm{w}}=(a,r)\in W_{\mathrm{af}}$ satisfy $r\,\mathrm{Lt}(n)=\mathrm{Lt}(n')$.
\begin{itemize}
\item[(i)] For a closed $u\in C^2(n)$ the number $\sum_{p\in\Pi_{\mu\nu}(\mathrm{jz})}u_p$ does not depend on $\mathrm{jz}$.
Put $\mathrm{fl}_{\mu\nu}(u):=r_q^{-1}\sum_{p\in\Pi_{\mu\nu}(\mathrm{jz})}u_p$ for $\mu<\nu$, $\mathrm{fl}_{\nu\mu}(u):=-\mathrm{fl}_{\mu\nu}(u)$, and $\mathrm{fl}(u):=\left(\mathrm{fl}_{23}(u),\mathrm{fl}_{13}(u),\mathrm{fl}_{12}(u)\right)\in\mathbb{R}^3$ (Section~I.7.1).
Then $\langle\mathbf{1}_{\mu\nu},u\rangle=r_q\cdot n_{\mathfrak{l}}\cdot\mathrm{fl}_{\mu\nu}(u)$ for all $\mu\ne\nu$.
\item[(ii)] $\mathrm{fl}$ is linear; $\mathrm{fl}(dA)=0$ for $A\in C^1(n)$; for $\mathfrak{z}\in\mathrm{Mz}(n)$ the vector $\mathrm{fl}(r_q\cdot\mathfrak{z})$ is the flux $w$ of $\mathfrak{z}$ (Section~I.7.1).
Hence $\mathrm{fl}(dA+r_q\cdot\mathfrak{z})=w$.
\item[(iii)] $g_{\mathrm{w}*}\mathbf{1}_{\mu\nu}=\epsilon_\mu\epsilon_\nu\cdot\mathbf{1}_{\mu'\nu'}$ for all $\mu\ne\nu$, the right side being formed on $n'$.
\item[(iv)] For a closed $u'\in C^2(n')$ the cochain $g_{\mathrm{w}}^{*}u'$ is closed, and $\mathrm{fl}_{\mu\nu}(g_{\mathrm{w}}^{*}u')=\epsilon_\mu\epsilon_\nu\cdot\mathrm{fl}_{\mu'\nu'}(u')$ for all $\mu\ne\nu$.
\item[(v)] The rule $(g_{\mathrm{w}}\cdot w)_{\mu'\nu'}:=\epsilon_\mu\epsilon_\nu\cdot w_{\mu\nu}$ ($\mu\ne\nu$; $w_{\nu\mu}=-w_{\mu\nu}$) defines a bijection $w\mapsto g_{\mathrm{w}}\cdot w$ of $\mathbb{Z}^3$ with $g_{\mathrm{w}}\cdot0=0$, and $\mathrm{fl}(g_{\mathrm{w}}^{*}u')=w$ if and only if $\mathrm{fl}(u')=g_{\mathrm{w}}\cdot w$.
\item[(vi)] $g_{\mathrm{w}*}h_w=h_{g_{\mathrm{w}}\cdot w}$, the right side being formed on $n'$, and $\left\|h_{g_{\mathrm{w}}\cdot w}\right\|=\left\|h_w\right\|$.
\item[(vii)] For $w\in\mathbb{Z}^3$ let $\mathfrak{z}^{(w)}_{\mu\nu}(c):=w_{\mu\nu}\cdot\mathbf{1}\left[c_\mu\in n_\mu\mathbb{Z},\ c_\nu\in n_\nu\mathbb{Z}\right]$ ($\mu<\nu$).
Then $\mathfrak{z}^{(w)}$ lies in $\mathrm{Mz}(n)$, and its flux is $w$.
\end{itemize}
\end{named}

\begin{proof}
(i) By Section~0.2,
\[
(du)(c;\{1,2,3\})=\sum_{\mu<\nu}\mathrm{sn}_{\mu\nu}\cdot\left[u_{\mu\nu}(c+e_{\mathfrak{l}})-u_{\mu\nu}(c)\right],\qquad\mathrm{sn}_{23}=\mathrm{sn}_{12}=1,\quad\mathrm{sn}_{13}=-1,
\]
where $\mathfrak{l}$ is the third direction of the pair $\mu,\nu$.
Fix a pair $\mu<\nu$ with the third direction $\mathfrak{l}$.
Since $du=0$, the sum of $(du)(c;\{1,2,3\})$ over the $n_\mu n_\nu$ classes of the sites $c$ with $c_{\mathfrak{l}}=\mathrm{jz}$ vanishes.
The bracket of each of the two other pairs is a difference in the direction $\mu$ or in the direction $\nu$, and $c\mapsto c+e_\mu$, $c\mapsto c+e_\nu$ are bijections of the set of these classes; so the sum of such a bracket vanishes.
There remains
\[
0=\mathrm{sn}_{\mu\nu}\cdot\left[\sum_{p\in\Pi_{\mu\nu}(\mathrm{jz}+1)}u_p-\sum_{p\in\Pi_{\mu\nu}(\mathrm{jz})}u_p\right].
\]
The plaquettes $(c;\mu\nu)$ fall into the $n_{\mathfrak{l}}$ planes $\mathrm{jz}=0,\dots,n_{\mathfrak{l}}-1$; hence $\langle\mathbf{1}_{\mu\nu},u\rangle=n_{\mathfrak{l}}\cdot r_q\cdot\mathrm{fl}_{\mu\nu}(u)$.
For $\mu>\nu$ both sides change their sign.

(ii) The number $\mathrm{fl}_{\mu\nu}(u)$ is $r_q^{-1}$ times a sum of values of $u$; so $\mathrm{fl}$ is linear on the space of the closed cochains.
With the sums over the classes of the sites $c$ with $c_{\mathfrak{l}}=\mathrm{jz}$,
\[
\sum_{c}(dA)_{\mu\nu}(c)=\sum_c\left[A_\mu(c)-A_\mu(c+e_\nu)\right]+\sum_c\left[A_\nu(c+e_\mu)-A_\nu(c)\right]=0,
\]
because $c\mapsto c+e_\nu$ and $c\mapsto c+e_\mu$ are bijections of the set of these classes.
Finally, $r_q^{-1}\sum_{p\in\Pi_{\mu\nu}(\mathrm{jz})}r_q\cdot\mathfrak{z}_p=w_{\mu\nu}$ is the definition of Section~I.7.1.

(iii) Let $\mu<\nu$ and let $\boldsymbol{\sigma}'$ be a plaquette.
Then $(g_{\mathrm{w}*}\mathbf{1}_{\mu\nu})(\boldsymbol{\sigma}')=\epsilon(g_{\mathrm{w}}^{-1},\boldsymbol{\sigma}')\cdot\mathbf{1}_{\mu\nu}(g_{\mathrm{w}}^{-1}\boldsymbol{\sigma}')$.
This vanishes unless $\boldsymbol{\sigma}'=g_{\mathrm{w}}\boldsymbol{\sigma}$ with $\boldsymbol{\sigma}=(c;\mu\nu)$, that is, unless $\boldsymbol{\sigma}'$ has the directions $\mu',\nu'$.
In this case $\epsilon(g_{\mathrm{w}}^{-1},g_{\mathrm{w}}\boldsymbol{\sigma})=\epsilon(g_{\mathrm{w}},\boldsymbol{\sigma})$ (Lemma~1\,(i)), and $(\Lambda^2r)(e_\mu\wedge e_\nu)=\epsilon_\mu\epsilon_\nu\cdot e_{\mu'}\wedge e_{\nu'}$, which is $\epsilon_\mu\epsilon_\nu\cdot\mathrm{or}(g_{\mathrm{w}}\boldsymbol{\sigma})$ if $\mu'<\nu'$ and $-\epsilon_\mu\epsilon_\nu\cdot\mathrm{or}(g_{\mathrm{w}}\boldsymbol{\sigma})$ if $\mu'>\nu'$.
By the convention $\mathbf{1}_{\nu'\mu'}=-\mathbf{1}_{\mu'\nu'}$ this is the assertion.
For $\mu>\nu$ both sides change their sign.

(iv) One has $d\,g_{\mathrm{w}}^{*}u'=g_{\mathrm{w}}^{*}du'=0$ by Lemma~1\,(ii).
Since $\mu\mapsto\mu'$ is a permutation, $\mathfrak{l}'$ is the third direction of $\mu',\nu'$, and $n'_{\mathfrak{l}'}=n_{\mathfrak{l}}$.
By (i), Lemma~1\,(iv), (iii) and (i),
\[
r_q\cdot n_{\mathfrak{l}}\cdot\mathrm{fl}_{\mu\nu}(g_{\mathrm{w}}^{*}u')=\langle\mathbf{1}_{\mu\nu},g_{\mathrm{w}}^{*}u'\rangle=\langle g_{\mathrm{w}*}\mathbf{1}_{\mu\nu},u'\rangle=\epsilon_\mu\epsilon_\nu\cdot\langle\mathbf{1}_{\mu'\nu'},u'\rangle=\epsilon_\mu\epsilon_\nu\cdot r_q\cdot n'_{\mathfrak{l}'}\cdot\mathrm{fl}_{\mu'\nu'}(u').
\]

(v) The map $(\mu,\nu)\mapsto(\mu',\nu')$ is a bijection of the six ordered pairs, and the rule is compatible with the antisymmetry: $(g_{\mathrm{w}}\cdot w)_{\nu'\mu'}=\epsilon_\nu\epsilon_\mu\cdot w_{\nu\mu}=-(g_{\mathrm{w}}\cdot w)_{\mu'\nu'}$.
So $g_{\mathrm{w}}\cdot w$ is well defined, and $w\mapsto g_{\mathrm{w}}\cdot w$ permutes the three components and changes some of their signs.
By (iv), $\mathrm{fl}(g_{\mathrm{w}}^{*}u')=w$ holds if and only if $\epsilon_\mu\epsilon_\nu\cdot\mathrm{fl}_{\mu'\nu'}(u')=w_{\mu\nu}$ for all $\mu\ne\nu$, that is, if and only if $\mathrm{fl}_{\mu'\nu'}(u')=\epsilon_\mu\epsilon_\nu\cdot w_{\mu\nu}=(g_{\mathrm{w}}\cdot w)_{\mu'\nu'}$ for all $\mu\ne\nu$, that is, if and only if $\mathrm{fl}(u')=g_{\mathrm{w}}\cdot w$.

(vi) By Section~I.7.1, $h_w=\sum_{\mu<\nu}\frac{r_q\cdot w_{\mu\nu}}{n_\mu n_\nu}\cdot\mathbf{1}_{\mu\nu}$.
By (iii) and $n_\mu=n'_{\mu'}$,
\[
g_{\mathrm{w}*}h_w=\sum_{\mu<\nu}\frac{r_q\cdot\epsilon_\mu\epsilon_\nu\cdot w_{\mu\nu}}{n'_{\mu'}n'_{\nu'}}\cdot\mathbf{1}_{\mu'\nu'}=\sum_{\mu<\nu}\frac{r_q\cdot(g_{\mathrm{w}}\cdot w)_{\mu'\nu'}}{n'_{\mu'}n'_{\nu'}}\cdot\mathbf{1}_{\mu'\nu'}.
\]
The summand does not change if $\mu'$ and $\nu'$ are exchanged, and $\{\mu',\nu'\}$ ranges over the three pairs; so the sum is $h_{g_{\mathrm{w}}\cdot w}$.
Since $g_{\mathrm{w}*}$ is orthogonal (Lemma~1\,(iv)), $\left\|h_{g_{\mathrm{w}}\cdot w}\right\|=\left\|h_w\right\|$.

(vii) The cochain $\mathfrak{z}^{(w)}_{\mu\nu}$ is $\mathrm{Lt}(n)$-periodic and integer valued, and it does not depend on $c_{\mathfrak{l}}$; so every bracket in the formula for $d$ in the proof of (i) vanishes.
Exactly one class of $\Pi_{\mu\nu}(\mathrm{jz})$ has $c_\mu\in n_\mu\mathbb{Z}$ and $c_\nu\in n_\nu\mathbb{Z}$.
\end{proof}

As in Section~I.7.1, the statement that the flux of the closed cochain $u$ is $w$ means $\mathrm{fl}(u)=w$.
By (ii), together with (i), of Lemma~4, the function $\mathbf{1}[\text{flux}=w]$ of a configuration $(A,\mathfrak{z})$ (Section~I.7.3) is $\mathrm{chf}_w(F^{\mathrm{u}})$, where
\[
\mathrm{chf}_w(u):=\prod_{\mu<\nu}\mathbf{1}\left[\langle\mathbf{1}_{\mu\nu},u\rangle=r_q\cdot n_{\mathfrak{l}}\cdot w_{\mu\nu}\right]
\]
is a Borel function on $C^2(n)$ with values in $\{0,1\}$.
Let $\mathrm{Mz}_w(n):=\left\{\mathfrak{z}\in\mathrm{Mz}(n):\mathrm{fl}(r_q\cdot\mathfrak{z})=w\right\}$.
Then $\mathrm{Mz}(n)$ is the disjoint union of the sets $\mathrm{Mz}_w(n)$, $w\in\mathbb{Z}^3$.

\subsubsection*{0.5. Periodic integrands; theta sums}

The coordinates on $C^{1}(n)=\mathbb{R}^{3|n|}$ are the numbers $A(b)$, where $b$ ranges over the classes of links,
and $[0,r_q)^{\mathrm{links}}:=\{A\in C^{1}(n):0\le A(b)<r_q\text{ for every }b\}$ is the box in $C^{1}(n)$.
When two tori occur, the box of the torus $n$ is written $[0,r_q)^{\mathrm{links}(n)}$.
The Lebesgue measure on $C^{1}(n)$ is denoted by $\mathrm{d}A$.

\begin{named}{Lemma 5}
Let $g_{\mathrm{w}}=(a,r)\in W_{\mathrm{af}}$ satisfy $r\,\mathrm{Lt}(n)=\mathrm{Lt}(n')$,
and let $\mathrm{Fn}\colon C^{1}(n')\to\mathbb{C}$ be a bounded Borel function such that $\mathrm{Fn}(A'+r_qm')=\mathrm{Fn}(A')$ for all $m'\in C^{1}(n';\mathbb{Z})$.
Then
\begin{equation*}
\int_{[0,r_q)^{\mathrm{links}(n)}}\mathrm{d}A\,\mathrm{Fn}(g_{\mathrm{w}*}A)=\int_{[0,r_q)^{\mathrm{links}(n')}}\mathrm{d}A'\,\mathrm{Fn}(A').
\end{equation*}
Moreover, if $\mathrm{Fn}$ is a bounded Borel function on $C^{1}(n)$ with $\mathrm{Fn}(A+r_qm)=\mathrm{Fn}(A)$ for all $m\in C^{1}(n;\mathbb{Z})$, then
\begin{equation*}
\int_{[0,r_q)^{\mathrm{links}}}\mathrm{d}A\,\mathrm{Fn}(-A)=\int_{[0,r_q)^{\mathrm{links}}}\mathrm{d}A\,\mathrm{Fn}(A).
\end{equation*}
\end{named}

\begin{proof}
By Lemma~1\,(i), (iv), the map $b\mapsto g_{\mathrm{w}}b$ is a bijection from the classes of links of $n$ onto those of $n'$, and
\begin{equation*}
(g_{\mathrm{w}*}A)(g_{\mathrm{w}}b)=\epsilon(g_{\mathrm{w}}^{-1},g_{\mathrm{w}}b)\cdot A(b)=\epsilon(g_{\mathrm{w}},b)\cdot A(b).
\end{equation*}
Thus $A\mapsto g_{\mathrm{w}*}A$ permutes the coordinates and then changes the signs of some of them.
Such a map preserves the Lebesgue measure, and it carries the box onto a product $\prod_{b'}\mathrm{Int}_{b'}$ in which every factor $\mathrm{Int}_{b'}$ is either $[0,r_q)$ or $(-r_q,0]$.
The left side of the first identity is therefore the integral of $\mathrm{Fn}$ over $\prod_{b'}\mathrm{Int}_{b'}$.
If $\mathrm{fq}$ is a bounded Borel function of one real variable with the period $r_q$, then
\begin{equation*}
\int_{(-r_q,0]}\mathrm{d}s\,\mathrm{fq}(s)=\int_{(0,r_q]}\mathrm{d}s\,\mathrm{fq}(s)=\int_{[0,r_q)}\mathrm{d}s\,\mathrm{fq}(s).
\end{equation*}
By Fubini's theorem this identity may be applied to the coordinates $b'$ with $\mathrm{Int}_{b'}=(-r_q,0]$, one after another; this gives the first assertion.
As to the second, the map $A\mapsto-A$ preserves the Lebesgue measure and carries the box onto $(-r_q,0]^{\mathrm{links}}$.
Hence the left side is the integral of $\mathrm{Fn}$ over $(-r_q,0]^{\mathrm{links}}$, and the identity for $\mathrm{fq}$ is applied to each of the $3|n|$ coordinates.
\end{proof}

\begin{named}{Lemma 6}[{cf.\ \cite[Lemmas~3.13 and~3.14]{QEDII}}]
For $\mathrm{aq}>0$ put
\begin{equation*}
\mathrm{Cth}(\mathrm{aq}):=\sup_{\mathrm{tq}\in\mathbb{R}}\ \sum_{\mathrm{jz}\in\mathbb{Z}}\exp\left(-\mathrm{aq}\cdot(\mathrm{tq}+r_q\cdot\mathrm{jz})^2\right).
\end{equation*}
Then
\begin{equation*}
\mathrm{Cth}(\mathrm{aq})\le2+\frac{2}{\exp(\mathrm{aq}\cdot r_q^2)-1},
\end{equation*}
and for every $A\in C^{1}(n)$
\begin{equation*}
\sum_{\mathfrak{z}\in C^{2}(n;\mathbb{Z})}\exp\left(-\mathrm{aq}\cdot\|dA+r_q\mathfrak{z}\|^2\right)\le\mathrm{Cth}(\mathrm{aq})^{3|n|}.
\end{equation*}
\end{named}

\begin{proof}
The sum over $\mathrm{jz}$ is periodic in $\mathrm{tq}$ with the period $r_q$, so that one may assume $0\le\mathrm{tq}<r_q$.
For $\mathrm{jz}\ge0$ one has $(\mathrm{tq}+r_q\cdot\mathrm{jz})^2\ge r_q^2\cdot\mathrm{jz}^2$;
for $\mathrm{jz}\le-1$ one has $|\mathrm{tq}+r_q\cdot\mathrm{jz}|\ge r_q\cdot(|\mathrm{jz}|-1)$.
The sum is therefore at most
\begin{equation*}
2\sum_{\mathrm{jz}\ge0}\exp\left(-\mathrm{aq}\cdot r_q^2\cdot\mathrm{jz}^2\right)\le2\sum_{\mathrm{jz}\ge0}\exp\left(-\mathrm{aq}\cdot r_q^2\cdot\mathrm{jz}\right)=2+\frac{2}{\exp(\mathrm{aq}\cdot r_q^2)-1}.
\end{equation*}
The number $\|dA+r_q\mathfrak{z}\|^2$ is the sum of $\left((dA)_p+r_q\mathfrak{z}_p\right)^2$ over the $3|n|$ classes $p$ of plaquettes, and the $\mathfrak{z}_p$ are independent integers.
So the sum over $\mathfrak{z}$ is the product over $p$ of the sums
\begin{equation*}
\sum_{\mathrm{jz}\in\mathbb{Z}}\exp\left(-\mathrm{aq}\cdot\left((dA)_p+r_q\cdot\mathrm{jz}\right)^2\right),
\end{equation*}
each of which is at most $\mathrm{Cth}(\mathrm{aq})$.
\end{proof}

\subsection*{(a)}

\begin{proof}[Proof of Proposition~50\,(a)]
Let $B$ be real and $n_\mu\ge1$.
In (a1)--(a3) the pair $(\mathrm{sp},\mathrm{gq})$ is either $(4,\gamma)$ or $(2,\sigma)$.
Let $\mathfrak{S}^{\mathrm{ch},\mathrm{sp}}(1;n)$ be the space of the functions $f\colon\mathbb{Z}^3\times O_{\mathrm{cub}}\to\mathbb{C}^{\mathrm{sp}}$ with $f(c+n_\mu e_\mu,g)=-f(c,g)$;
its dimension is $48\cdot\mathrm{sp}\cdot|n|$.
It carries the scalar product
\begin{equation*}
(f',f):=\frac{1}{48}\sum_{x\in\mathrm{Zd}(n)\times O_{\mathrm{cub}}}\overline{f'(x)}\cdot f(x)
\end{equation*}
and the corresponding norm $\|f\|$;
for $\mathrm{sp}=4$ these are $\mathfrak{S}^{\mathrm{ch}}(1;n)$ and $(\,\cdot\,,\cdot\,)^{\mathrm{ch}}$.
For a vector $v$ put $\mathrm{gq}\cdot v:=\sum_\mu v_\mu\,\mathrm{gq}_\mu$;
for real $v$ this is a Hermitian $\mathrm{sp}\times\mathrm{sp}$ matrix, by (I.c1) and because the Pauli matrices are Hermitian.
On $\mathfrak{S}^{\mathrm{ch},\mathrm{sp}}(1;n)$ define
\begin{equation*}
(R_i(B)f)(x):=U_B(x,x^{(i)})\,f(x^{(i)}),\qquad
\nabla_B^{*}\nabla_B:=\sum_iw_i\left(1-R_i(B)\right),\qquad
J_B^{\mathrm{gq}}:=\sum_iw_i\,(\mathrm{gq}\cdot\hat d_i)\,R_i(B).
\end{equation*}
The summands of $(\,\cdot\,,\cdot\,)$ are $\mathrm{Lt}(n)$-periodic, so that (0.4) applies.
In the basis $\delta_{x,\beta}$, $\beta=1,\dots,\mathrm{sp}$, the scalar product is $1/48$ times the standard one;
hence $\det O^*=\overline{\det O}$, as in (0.5).

\emph{(a1) $R_i(B)$ is a unitary involution.}
By (0.4) for the bijection $x\mapsto x^{(i)}$, by (0.1) and by (I.s1), all sums being over $x\in\mathrm{Zd}(n)\times O_{\mathrm{cub}}$,
\begin{align*}
(f',R_i(B)f)&=\frac{1}{48}\sum_x\overline{f'(x)}\cdot U_B(x,x^{(i)})\,f(x^{(i)})=\frac{1}{48}\sum_x\overline{f'(x^{(i)})}\cdot U_B(x^{(i)},x)\,f(x)\\
&=\frac{1}{48}\sum_x\overline{U_B(x,x^{(i)})\,f'(x^{(i)})}\cdot f(x)=(R_i(B)f',f).
\end{align*}
Moreover $(R_i(B)^2f)(x)=U_B(x,x^{(i)})\,U_B(x^{(i)},x)\,f(x)=f(x)$.
Hence $R_i(B)^*=R_i(B)=R_i(B)^{-1}$, and $\|R_i(B)\|=1$.

\emph{(a2) $\nabla_B^{*}\nabla_B=\left(\nabla_B^{*}\nabla_B\right)^{*}\ge0$.}
The operator $\nabla_B^{*}\nabla_B$ is Hermitian by (a1), and
\begin{equation*}
(f,(1-R_i(B))f)\ge\|f\|^2-\|f\|\cdot\|R_i(B)f\|=0;
\end{equation*}
moreover $w_i>0$.

\emph{(a3) $(J_B^{\mathrm{gq}})^*=-J_B^{\mathrm{gq}}$.}
Since $\hat d_i(x)$ is real, the multiplication operator $\mathrm{gq}\cdot\hat d_i$ is Hermitian, and $\left((\mathrm{gq}\cdot\hat d_i)\,R_i(B)\right)^*=R_i(B)\,(\mathrm{gq}\cdot\hat d_i)$.
By (I.s2),
\begin{align*}
\left(R_i(B)\,(\mathrm{gq}\cdot\hat d_i)f\right)(x)&=U_B(x,x^{(i)})\left(\mathrm{gq}\cdot\hat d_i(x^{(i)})\right)f(x^{(i)})\\
&=-\left(\mathrm{gq}\cdot\hat d_i(x)\right)U_B(x,x^{(i)})\,f(x^{(i)})=-\left((\mathrm{gq}\cdot\hat d_i)\,R_i(B)f\right)(x).
\end{align*}

\emph{(a4) Real part.}
The operator $J_B$ of (0.3) is $J_B^{\gamma}$.
By (0.3), by (a2) and (a3) for $\mathrm{sp}=4$, and since $\boldsymbol{\mu}$ is real,
\begin{equation*}
\operatorname{Re}(D_B+\boldsymbol{\mu}):=\frac{1}{2}\left[(D_B+\boldsymbol{\mu})+(D_B+\boldsymbol{\mu})^*\right]=\boldsymbol{\mu}+\frac{\rho}{2}\,\nabla_B^{*}\nabla_B\ge\boldsymbol{\mu},
\end{equation*}
because $\rho\ge0$.

\emph{(a5) Norm.}
By (0.3), $D_B+\boldsymbol{\mu}=M_{\mathrm{ch}}-\sum_iM_iR_i(B)$.
The multiplication operator $M_i$ has the norm $\max_x|M_i(x)|_2=\frac{w_i}{2}(1+\rho)$:
by (I.m6) the eigenvalues of the matrix $M_i(x)$, which is Hermitian by (I.c1), are $\frac{w_i}{2}(\rho-1)$ and $\frac{w_i}{2}(\rho+1)$,
and the larger of their moduli is $\frac{w_i}{2}(1+\rho)$, since $w_i>0$ and $\rho\ge0$.
Together with $\|R_i(B)\|=1$ from (a1) this gives
\begin{equation*}
\|D_B+\boldsymbol{\mu}\|\le|M_{\mathrm{ch}}|+\sum_i\frac{w_i}{2}(1+\rho)=|M_{\mathrm{ch}}|+\frac{1+\rho}{2}\,\mathrm{sw}.
\end{equation*}

\emph{(a6) Blocks.}
One has $\Sigma_3=\operatorname{diag}(1,1,-1,-1)$ and $\gamma_\mu=\operatorname{diag}(\sigma_\mu,-\sigma_\mu)$ (Section~I.2.1).
For $f\in\mathfrak{S}^{\mathrm{ch}}(1;n)$ write $f=(f^+,f^-)$, where $f^+$ consists of the first two and $f^-$ of the last two components of $f$;
then $f^+$ and $f^-$ range over $\mathfrak{S}^{\mathrm{ch},2}(1;n)$.
The eigenspaces of $\Sigma_3$ for the eigenvalues $+1$ and $-1$ are $\{(f^+,0)\}$ and $\{(0,f^-)\}$;
they are orthogonal, and $(f',f)^{\mathrm{ch}}=(f'^+,f^+)+(f'^-,f^-)$.
Moreover
\begin{gather*}
R_i(B)(f^+,f^-)=\left(R_i(B)f^+,R_i(B)f^-\right),\qquad
\nabla_B^{*}\nabla_B(f^+,f^-)=\left(\nabla_B^{*}\nabla_Bf^+,\nabla_B^{*}\nabla_Bf^-\right),\\
J_B(f^+,f^-)=\left(J_B^{\sigma}f^+,-J_B^{\sigma}f^-\right).
\end{gather*}
Put
\begin{equation*}
D_B^{\boldsymbol{\sigma}}:=\frac{\rho}{2}\,\nabla_B^{*}\nabla_B+\frac{1}{2}\,J_B^{\sigma}\qquad\text{on }\mathfrak{S}^{\mathrm{ch},2}(1;n).
\end{equation*}
By (0.3),
\begin{equation*}
D_B(f^+,f^-)=\left(D_B^{\boldsymbol{\sigma}}f^+,\left[\frac{\rho}{2}\,\nabla_B^{*}\nabla_B-\frac{1}{2}\,J_B^{\sigma}\right]f^-\right),
\end{equation*}
and (a2), (a3) for $\mathrm{sp}=2$ give $\frac{\rho}{2}\,\nabla_B^{*}\nabla_B-\frac{1}{2}\,J_B^{\sigma}=(D_B^{\boldsymbol{\sigma}})^*$.
As $\boldsymbol{\mu}$ is real, it follows that
\begin{gather*}
D_B+\boldsymbol{\mu}=\operatorname{diag}\left(D_B^{\boldsymbol{\sigma}}+\boldsymbol{\mu},(D_B^{\boldsymbol{\sigma}}+\boldsymbol{\mu})^*\right),\\
\det\nolimits_n(D_B+\boldsymbol{\mu})=\det(D_B^{\boldsymbol{\sigma}}+\boldsymbol{\mu})\cdot\overline{\det(D_B^{\boldsymbol{\sigma}}+\boldsymbol{\mu})}=\left|\det(D_B^{\boldsymbol{\sigma}}+\boldsymbol{\mu})\right|^2\ge0.
\end{gather*}

\emph{(a7) The case $\boldsymbol{\mu}>0$.}
Let $\mathrm{zl}$ be an eigenvalue of $D_B+\boldsymbol{\mu}$ and $f\ne0$ an eigenvector.
By (a4),
\begin{equation*}
(\operatorname{Re}\mathrm{zl})\cdot\|f\|^2=\operatorname{Re}\,(f,(D_B+\boldsymbol{\mu})f)^{\mathrm{ch}}=\left(f,\operatorname{Re}(D_B+\boldsymbol{\mu})f\right)^{\mathrm{ch}}\ge\boldsymbol{\mu}\cdot\|f\|^2.
\end{equation*}
So $|\mathrm{zl}|\ge\operatorname{Re}\mathrm{zl}\ge\boldsymbol{\mu}$.
The number $\det\nolimits_n(D_B+\boldsymbol{\mu})$ is the product of the $192|n|$ eigenvalues, and it is non-negative by (a6);
it is therefore equal to its modulus, which is at least $\boldsymbol{\mu}^{192|n|}$.
\end{proof}

From (a5), (a6) and (0.6) one obtains, for every real $B$ and every real $\boldsymbol{\mu}$,
\begin{equation}
0\le\det\nolimits_n(D_B+\boldsymbol{\mu})\le\mathrm{Dmax}:=\left(|M_{\mathrm{ch}}|+\frac{1+\rho}{2}\,\mathrm{sw}\right)^{192|n|}.\tag{a.1}
\end{equation}

\subsection*{(b)}

\begin{proof}[Proof of Proposition~50\,(b)]
By Proposition~38.1\,(a), $\det\nolimits_n(D_B+\boldsymbol{\mu})$ is a polynomial in the numbers $\exp(iqB(b))$ and $\exp(-iqB(b))$;
it is therefore an entire function of $B\in\mathbb{C}^{3|n|}$, and in particular it is continuous on $C^{1}(n)$.
Proposition~38.1\,(c) is the invariance under $B\mapsto B+d\boldsymbol{\lambda}$, $\boldsymbol{\lambda}$ a complex zero-cochain of the torus,
and Proposition~38.1\,(b) is the invariance under $B\mapsto B+r_qm$, $m\in C^{1}(n;\mathbb{Z})$.
By Proposition~38.1\,(d), on the set $\{B:\max_b|\operatorname{Im}B(b)|\le\mathrm{cf}\}$ the modulus of the determinant does not exceed
\begin{equation*}
\left(|M_{\mathrm{ch}}|+\frac{1+\rho}{2}\left(\mathrm{sw}-6+6\exp(q\cdot\mathrm{cf})\right)\right)^{192|n|};
\end{equation*}
for $\mathrm{cf}=0$ this is the bound $\mathrm{Dmax}$ on $C^{1}(n)$.
\end{proof}

\subsection*{(c)}

\begin{proof}[Proof of Proposition~50\,(c)]
Let $B$ be complex, $n_\mu\ge1$, and $f,f'\in\mathfrak{S}^{\mathrm{ch}}(1;n)$.
For $v,v'\in\mathbb{C}^4$ one has $v'\cdot\left(M_i(x^{(i)})v\right)=\left(M_i(x^{(i)})^Tv'\right)\cdot v$, where $M_i(x^{(i)})^T$ denotes the transposed $4\times4$ matrix.
By (0.3), by (0.4) for the bijection $x\mapsto x^{(i)}$, by (0.1) and by (I.s1), with all sums over $x\in\mathrm{Zd}(n)\times O_{\mathrm{cub}}$,
\begin{align*}
\langle f',(D_B+\boldsymbol{\mu})f\rangle^{\mathrm{ch}}&=\frac{1}{48}\sum_xM_{\mathrm{ch}}\,f'(x)\cdot f(x)-\frac{1}{48}\sum_i\sum_xf'(x)\cdot M_i(x)\,U_B(x,x^{(i)})\,f(x^{(i)})\\
&=\frac{1}{48}\sum_xM_{\mathrm{ch}}\,f'(x)\cdot f(x)-\frac{1}{48}\sum_i\sum_xf'(x^{(i)})\cdot M_i(x^{(i)})\,U_B(x^{(i)},x)\,f(x)\\
&=\frac{1}{48}\sum_x\left[M_{\mathrm{ch}}\,f'(x)-\sum_iM_i(x^{(i)})^T\,U_{-B}(x,x^{(i)})\,f'(x^{(i)})\right]\cdot f(x).
\end{align*}
Consequently
\begin{equation*}
\left((D_B+\boldsymbol{\mu})^Tf'\right)(x)=M_{\mathrm{ch}}\,f'(x)-\sum_iM_i(x^{(i)})^T\,U_{-B}(x,x^{(i)})\,f'(x^{(i)}).
\end{equation*}
By (I.c3), $\gamma_\mu^T=-\mathsf{C}^{-1}\gamma_\mu\mathsf{C}$ and $\mathsf{C}^{-1}=-\mathsf{C}$;
hence $\gamma_\mu^T=-\mathsf{C}\gamma_\mu\mathsf{C}^{-1}$, and (0.2) yields
\begin{equation*}
M_i(x^{(i)})^T=\frac{w_i}{2}\left(\rho+\left(\gamma\cdot\hat d_i(x)\right)^T\right)=\mathsf{C}\,\frac{w_i}{2}\left(\rho-\gamma\cdot\hat d_i(x)\right)\mathsf{C}^{-1}=\mathsf{C}\,M_i(x)\,\mathsf{C}^{-1}
\end{equation*}
(this is (I.m2)).
Let $\mathsf{C}$ act on $\mathfrak{S}^{\mathrm{ch}}(1;n)$ site by site, $(\mathsf{C}f)(x):=\mathsf{C}f(x)$;
this operator is invertible and commutes with $R_i(-B)$.
Therefore
\begin{equation*}
(D_B+\boldsymbol{\mu})^T=\mathsf{C}\left[M_{\mathrm{ch}}-\sum_iM_iR_i(-B)\right]\mathsf{C}^{-1}=\mathsf{C}\,(D_{-B}+\boldsymbol{\mu})\,\mathsf{C}^{-1}
\end{equation*}
(this is (I.O5), since $\mathsf{C}^{-1}=-\mathsf{C}$), and (0.5) gives
\begin{equation*}
\det\nolimits_n(D_B+\boldsymbol{\mu})=\det\nolimits_n\left((D_B+\boldsymbol{\mu})^T\right)=\det\nolimits_n(D_{-B}+\boldsymbol{\mu}).
\end{equation*}
\end{proof}

\subsection*{(d)}

\begin{proof}[Proof of Proposition~50\,(d)]
Let $B$ be complex, $n_\mu,n'_\mu\ge1$, $r\,\mathrm{Lt}(n)=\mathrm{Lt}(n')$, and let $S_{\mathrm{gw}}$ be as in Lemma~3.

\emph{(d1) $U_{\mathrm{gw}}$ maps $\mathfrak{S}^{\mathrm{ch}}(1;n)$ onto $\mathfrak{S}^{\mathrm{ch}}(1;n')$.}
Let $f\in\mathfrak{S}^{\mathrm{ch}}(1;n)$.
Lemma~2 for $g_{\mathrm{w}}^{-1}$ gives
\begin{equation*}
g_{\mathrm{w}}^{-1}(x+n'_{\mu'}e_{\mu'})=g_{\mathrm{w}}^{-1}x+n'_{\mu'}\,r^{-1}e_{\mu'}=g_{\mathrm{w}}^{-1}x+\epsilon_\mu n_\mu e_\mu.
\end{equation*}
For every fermion site $\tilde x$ one has $f(\tilde x+n_\mu e_\mu)=-f(\tilde x)$ and $f(\tilde x-n_\mu e_\mu)=-f(\tilde x)$.
Hence
\begin{equation*}
(U_{\mathrm{gw}}f)(x+n'_{\mu'}e_{\mu'})=S_{\mathrm{gw}}f(g_{\mathrm{w}}^{-1}x+\epsilon_\mu n_\mu e_\mu)=-(U_{\mathrm{gw}}f)(x)\qquad(\mu=1,2,3),
\end{equation*}
and $\mu'$ ranges over the three directions.
For $f'\in\mathfrak{S}^{\mathrm{ch}}(1;n')$ put $(\mathrm{Vs}f')(\tilde x):=S_{\mathrm{gw}}^{-1}f'(g_{\mathrm{w}}\tilde x)$.
By Lemma~2, $g_{\mathrm{w}}(\tilde x+n_\mu e_\mu)=g_{\mathrm{w}}\tilde x+n_\mu re_\mu=g_{\mathrm{w}}\tilde x+\epsilon_\mu n'_{\mu'}e_{\mu'}$, so that
\begin{equation*}
(\mathrm{Vs}f')(\tilde x+n_\mu e_\mu)=S_{\mathrm{gw}}^{-1}f'(g_{\mathrm{w}}\tilde x+\epsilon_\mu n'_{\mu'}e_{\mu'})=-(\mathrm{Vs}f')(\tilde x);
\end{equation*}
thus $\mathrm{Vs}$ maps $\mathfrak{S}^{\mathrm{ch}}(1;n')$ into $\mathfrak{S}^{\mathrm{ch}}(1;n)$.
Finally,
\begin{equation*}
(\mathrm{Vs}\,U_{\mathrm{gw}}f)(\tilde x)=S_{\mathrm{gw}}^{-1}S_{\mathrm{gw}}f(g_{\mathrm{w}}^{-1}g_{\mathrm{w}}\tilde x)=f(\tilde x),\qquad
(U_{\mathrm{gw}}\mathrm{Vs}f')(x)=S_{\mathrm{gw}}S_{\mathrm{gw}}^{-1}f'(g_{\mathrm{w}}g_{\mathrm{w}}^{-1}x)=f'(x).
\end{equation*}
So $U_{\mathrm{gw}}$ is a linear bijection, and $U_{\mathrm{gw}}^{-1}=\mathrm{Vs}$.

\emph{(d2) The similarity.}
Let $f'\in\mathfrak{S}^{\mathrm{ch}}(1;n')$, $x\in\mathbb{Z}^3\times O_{\mathrm{cub}}$, $\tilde x:=g_{\mathrm{w}}^{-1}x$ and $\tilde x^{(i)}:=(\tilde x)^{(i)}$.
By Lemma~2, $g_{\mathrm{w}}(\tilde x^{(i)})=(g_{\mathrm{w}}\tilde x)^{(i)}=x^{(i)}$.
Hence
\begin{align*}
\left(U_{\mathrm{gw}}D_BU_{\mathrm{gw}}^{-1}f'\right)(x)&=S_{\mathrm{gw}}\left(D_BU_{\mathrm{gw}}^{-1}f'\right)(\tilde x)\\
&=-\sum_iS_{\mathrm{gw}}M_i(\tilde x)\left[U_B(\tilde x,\tilde x^{(i)})\,S_{\mathrm{gw}}^{-1}f'(g_{\mathrm{w}}\tilde x^{(i)})-S_{\mathrm{gw}}^{-1}f'(g_{\mathrm{w}}\tilde x)\right]\\
&=-\sum_iS_{\mathrm{gw}}M_i(\tilde x)S_{\mathrm{gw}}^{-1}\left[U_B(\tilde x,\tilde x^{(i)})\,f'(x^{(i)})-f'(x)\right].
\end{align*}
By (0.8) and Lemma~2,
\begin{equation*}
S_{\mathrm{gw}}M_i(\tilde x)S_{\mathrm{gw}}^{-1}=\frac{w_i}{2}\left(\rho-\gamma\cdot\left(r\hat d_i(\tilde x)\right)\right)=\frac{w_i}{2}\left(\rho-\gamma\cdot\hat d_i(g_{\mathrm{w}}\tilde x)\right)=M_i(x).
\end{equation*}
Lemma~2 gives $\pi\tilde x=g_{\mathrm{w}}^{-1}\pi x$ and $\pi\tilde x^{(i)}=g_{\mathrm{w}}^{-1}\pi x^{(i)}$; together with Lemma~1\,(vi) this yields
\begin{equation*}
U_B(\tilde x,\tilde x^{(i)})=\exp\left(iqB(g_{\mathrm{w}}^{-1}\pi x,g_{\mathrm{w}}^{-1}\pi x^{(i)})\right)=\exp\left(iq\,(g_{\mathrm{w}*}B)(\pi x,\pi x^{(i)})\right)=U_{g_{\mathrm{w}*}B}(x,x^{(i)}),
\end{equation*}
where for $i\le2$ both sides equal $1$.
The cochain $g_{\mathrm{w}*}B$ lies in $C^{1}(n')\otimes\mathbb{C}$ (Lemma~1\,(iii)).
Hence $U_{\mathrm{gw}}D_BU_{\mathrm{gw}}^{-1}=D_{g_{\mathrm{w}*}B}$ on $\mathfrak{S}^{\mathrm{ch}}(1;n')$, and $U_{\mathrm{gw}}\boldsymbol{\mu}U_{\mathrm{gw}}^{-1}=\boldsymbol{\mu}$.

\emph{(d3) Determinants.}
If $(f_1,\dots,f_{192|n|})$ is a basis of $\mathfrak{S}^{\mathrm{ch}}(1;n)$, then $(U_{\mathrm{gw}}f_1,\dots,U_{\mathrm{gw}}f_{192|n|})$ is a basis of $\mathfrak{S}^{\mathrm{ch}}(1;n')$,
and in this basis $D_{g_{\mathrm{w}*}B}+\boldsymbol{\mu}$ has the matrix that $D_B+\boldsymbol{\mu}$ has in the first one.
Therefore $\det\nolimits_{n'}(D_{g_{\mathrm{w}*}B}+\boldsymbol{\mu})=\det\nolimits_n(D_B+\boldsymbol{\mu})$.
\end{proof}

\subsection*{(e), first part: convergence; the sum over the sectors}

Let $\Phi\in\mathfrak{F}(n)$ (Section~G6.A3);
thus $\Phi$ is a Borel function on $C^{2}(n)$ with $|\Phi(u)|\le\mathrm{cf}\cdot\exp(\mathrm{cf}\cdot\|u\|)$.
Since $\left(\frac{1}{2}\|u\|-\mathrm{cf}\right)^2\ge0$, one has $\mathrm{cf}\cdot\|u\|\le\mathrm{cf}^2+\frac{1}{4}\|u\|^2$, and therefore
\begin{equation}
|\Phi(u)|\cdot\exp\left(-\frac{1}{2}\|u\|^2\right)\le\mathrm{cf}\cdot\exp(\mathrm{cf}^2)\cdot\exp\left(-\frac{1}{4}\|u\|^2\right)\qquad\text{for all }u\in C^{2}(n).\tag{e.1}
\end{equation}
For a subset $\mathrm{Smz}$ of $\mathrm{Mz}(n)$ and for $A\in C^{1}(n)$ put
\begin{equation*}
\mathrm{Psw}[\Phi,\mathrm{Smz}](A):=\det\nolimits_n(D_A+\boldsymbol{\mu})\cdot\sum_{\mathfrak{z}\in\mathrm{Smz}}\Phi(dA+r_q\mathfrak{z})\cdot\exp\left(-\frac{1}{2}\|dA+r_q\mathfrak{z}\|^2\right).
\end{equation*}
By (a.1), (e.1) and Lemma~6 with $\mathrm{aq}=1/4$ (the sum being enlarged by dropping the condition $d\mathfrak{z}=0$), the series converges absolutely for every $A$, and
\begin{align}
&\det\nolimits_n(D_A+\boldsymbol{\mu})\cdot\sum_{\mathfrak{z}\in\mathrm{Mz}(n)}|\Phi(dA+r_q\mathfrak{z})|\cdot\exp\left(-\frac{1}{2}\|dA+r_q\mathfrak{z}\|^2\right)\notag\\
&\qquad\le\mathrm{Dmax}\cdot\mathrm{cf}\cdot\exp(\mathrm{cf}^2)\cdot\mathrm{Cth}(1/4)^{3|n|}\qquad\text{for all }A\in C^{1}(n).\tag{e.2}
\end{align}
Every term is a Borel function of $A$:
$\Phi$ is Borel, $A\mapsto dA+r_q\mathfrak{z}$ is continuous, and the determinant is continuous by (b).
The box has the volume $r_q^{3|n|}$.
Tonelli's theorem, together with the non-negativity of the determinant from (a), therefore gives
\begin{align}
Z_n[|\Phi|]&=\sum_{\mathfrak{z}\in\mathrm{Mz}(n)}\int_{[0,r_q)^{\mathrm{links}}}\mathrm{d}A\,|\Phi(F^{\mathrm{u}})|\cdot\exp\left(-\frac{1}{2}\|F^{\mathrm{u}}\|^2\right)\cdot\det\nolimits_n(D_A+\boldsymbol{\mu})\notag\\
&\le\left(r_q\cdot\mathrm{Cth}(1/4)\right)^{3|n|}\cdot\mathrm{Dmax}\cdot\mathrm{cf}\cdot\exp(\mathrm{cf}^2)<\infty.\tag{e.3}
\end{align}
This is the absolute convergence of the sum and the integral in (I.7.1).
By Fubini's theorem, by Lemma~4\,(ii) and by the definition of $\mathrm{chf}_w$,
\begin{align}
Z_n[\Phi]&=\int_{[0,r_q)^{\mathrm{links}}}\mathrm{d}A\,\mathrm{Psw}[\Phi,\mathrm{Mz}(n)](A),\notag\\
Z^{(w)}_n[\Phi]&=Z_n[\mathrm{chf}_w\Phi]=\int_{[0,r_q)^{\mathrm{links}}}\mathrm{d}A\,\mathrm{Psw}[\Phi,\mathrm{Mz}_w(n)](A)\qquad(w\in\mathbb{Z}^3);\tag{e.4}
\end{align}
here $\mathrm{chf}_w\Phi$ lies in $\mathfrak{F}(n)$.
As the determinant is non-negative, $|Z^{(w)}_n[\Phi]|\le Z^{(w)}_n[|\Phi|]$.
The set $\mathrm{Mz}(n)$ is the disjoint union of the sets $\mathrm{Mz}_w(n)$.
Hence $\sum_wZ^{(w)}_n[|\Phi|]=Z_n[|\Phi|]$ by Tonelli's theorem, so that
\begin{equation*}
\sum_w\left|Z^{(w)}_n[\Phi]\right|\le Z_n[|\Phi|]<\infty,\qquad\sum_wZ^{(w)}_n[\Phi]=Z_n[\Phi].
\end{equation*}
For $\Phi=1$, the bound (a.1) and Lemma~6 with $\mathrm{aq}=1/2$ give directly
\begin{equation*}
Z_n[1]\le\left(r_q\cdot\mathrm{Cth}(1/2)\right)^{3|n|}\cdot\mathrm{Dmax};
\end{equation*}
this is the upper bound of \cite[Theorem~4.13\,(c)]{QEDII}, whose constant $C_q$ is $\mathrm{Cth}(1/2)$.

\emph{Periodicity.}
Let $\mathrm{Smz}=\mathrm{Mz}(n)$ or $\mathrm{Smz}=\mathrm{Mz}_w(n)$, and let $m\in C^{1}(n;\mathbb{Z})$.
The map $\mathfrak{z}\mapsto\mathfrak{z}+dm$ is a bijection of $\mathrm{Smz}$, because $d\,dm=0$ (Section~0.2) and $\mathrm{fl}(r_q\,dm)=0$ (Lemma~4\,(ii)).
With Proposition~38.1\,(b) it follows that
\begin{align}
\mathrm{Psw}[\Phi,\mathrm{Smz}](A+r_qm)&=\det\nolimits_n(D_A+\boldsymbol{\mu})\cdot\sum_{\mathfrak{z}\in\mathrm{Smz}}\Phi\left(dA+r_q(\mathfrak{z}+dm)\right)\cdot\exp\left(-\frac{1}{2}\left\|dA+r_q(\mathfrak{z}+dm)\right\|^2\right)\notag\\
&=\mathrm{Psw}[\Phi,\mathrm{Smz}](A).\tag{e.5}
\end{align}
By (e.2), $\mathrm{Psw}[\Phi,\mathrm{Smz}]$ is a bounded Borel function on $C^{1}(n)$.

\subsection*{(f) and (g)}

\begin{proof}[Proof of Proposition~50\,(f) and (g)]
Let $r\,\mathrm{Lt}(n)=\mathrm{Lt}(n')$ and $\Phi\in\mathfrak{F}(n)$.
The function $\Phi\circ g_{\mathrm{w}}^{*}$ is a Borel function on $C^{2}(n')$, and by Lemma~1\,(iv)
\begin{equation*}
|\Phi(g_{\mathrm{w}}^{*}u')|\le\mathrm{cf}\cdot\exp\left(\mathrm{cf}\cdot\|g_{\mathrm{w}}^{*}u'\|\right)=\mathrm{cf}\cdot\exp\left(\mathrm{cf}\cdot\|u'\|\right);
\end{equation*}
so it lies in $\mathfrak{F}(n')$.
The assertions of (g) concerning the map $w\mapsto g_{\mathrm{w}}\cdot w$, the norm $\|h_{g_{\mathrm{w}}\cdot w}\|$ and the flux of $g_{\mathrm{w}}^{*}u'$ are Lemma~4\,(v), (vi).

\emph{(f1) Sectors.}
By Lemma~1\,(ii), (iii), $g_{\mathrm{w}*}$ maps $\mathrm{Mz}(n)$ onto $\mathrm{Mz}(n')$.
For $\mathfrak{z}\in\mathrm{Mz}(n)$ and $\mathfrak{z}':=g_{\mathrm{w}*}\mathfrak{z}$ one has $g_{\mathrm{w}}^{*}(r_q\mathfrak{z}')=r_q\mathfrak{z}$;
by Lemma~4\,(v), therefore, $\mathrm{fl}(r_q\mathfrak{z})=w$ holds if and only if $\mathrm{fl}(r_q\mathfrak{z}')=g_{\mathrm{w}}\cdot w$.
Hence $g_{\mathrm{w}*}\mathrm{Mz}_w(n)=\mathrm{Mz}_{g_{\mathrm{w}}\cdot w}(n')$.

\emph{(f2) Integrands.}
Let either $\mathrm{Smz}=\mathrm{Mz}(n)$ and $\mathrm{Smz}'=\mathrm{Mz}(n')$, or $\mathrm{Smz}=\mathrm{Mz}_w(n)$ and $\mathrm{Smz}'=\mathrm{Mz}_{g_{\mathrm{w}}\cdot w}(n')$;
in both cases $\mathrm{Smz}'=g_{\mathrm{w}*}\mathrm{Smz}$ by (f1).
Let $\mathrm{Psw}'$ denote the function $\mathrm{Psw}$ formed on $n'$.
For $A\in C^{1}(n)$, writing the sum over $\mathrm{Smz}'$ as a sum over $\mathfrak{z}'=g_{\mathrm{w}*}\mathfrak{z}$ with $\mathfrak{z}\in\mathrm{Smz}$, one obtains
\begin{multline*}
\mathrm{Psw}'[\Phi\circ g_{\mathrm{w}}^{*},\mathrm{Smz}'](g_{\mathrm{w}*}A)=\det\nolimits_{n'}(D_{g_{\mathrm{w}*}A}+\boldsymbol{\mu})\\
\cdot\sum_{\mathfrak{z}\in\mathrm{Smz}}\Phi\left(g_{\mathrm{w}}^{*}\left(d\,g_{\mathrm{w}*}A+r_q\,g_{\mathrm{w}*}\mathfrak{z}\right)\right)\cdot\exp\left(-\frac{1}{2}\left\|d\,g_{\mathrm{w}*}A+r_q\,g_{\mathrm{w}*}\mathfrak{z}\right\|^2\right).
\end{multline*}
Here $d\,g_{\mathrm{w}*}A+r_q\,g_{\mathrm{w}*}\mathfrak{z}=g_{\mathrm{w}*}(dA+r_q\mathfrak{z})$ by Lemma~1\,(ii),
$g_{\mathrm{w}}^{*}g_{\mathrm{w}*}=1$ by Lemma~1\,(i),
$\|g_{\mathrm{w}*}u\|=\|u\|$ by Lemma~1\,(iv),
and $\det\nolimits_{n'}(D_{g_{\mathrm{w}*}A}+\boldsymbol{\mu})=\det\nolimits_n(D_A+\boldsymbol{\mu})$ by (d).
Thus
\begin{equation*}
\mathrm{Psw}'[\Phi\circ g_{\mathrm{w}}^{*},\mathrm{Smz}'](g_{\mathrm{w}*}A)=\mathrm{Psw}[\Phi,\mathrm{Smz}](A).
\end{equation*}

\emph{(f3) Integration.}
By (e.2) and (e.5) on $n'$, the function $\mathrm{Fn}:=\mathrm{Psw}'[\Phi\circ g_{\mathrm{w}}^{*},\mathrm{Smz}']$ is bounded, Borel and invariant under $A'\mapsto A'+r_qm'$, $m'\in C^{1}(n';\mathbb{Z})$.
By (e.4) on $n'$, Lemma~5, (f2) and (e.4) on $n$,
\begin{align*}
\int_{[0,r_q)^{\mathrm{links}(n')}}\mathrm{d}A'\,\mathrm{Psw}'[\Phi\circ g_{\mathrm{w}}^{*},\mathrm{Smz}'](A')&=\int_{[0,r_q)^{\mathrm{links}(n)}}\mathrm{d}A\,\mathrm{Psw}'[\Phi\circ g_{\mathrm{w}}^{*},\mathrm{Smz}'](g_{\mathrm{w}*}A)\\
&=\int_{[0,r_q)^{\mathrm{links}(n)}}\mathrm{d}A\,\mathrm{Psw}[\Phi,\mathrm{Smz}](A).
\end{align*}
For $\mathrm{Smz}=\mathrm{Mz}(n)$ this identity reads $Z_{n'}[\Phi\circ g_{\mathrm{w}}^{*}]=Z_n[\Phi]$;
for $\mathrm{Smz}=\mathrm{Mz}_w(n)$ it reads $Z^{(g_{\mathrm{w}}\cdot w)}_{n'}[\Phi\circ g_{\mathrm{w}}^{*}]=Z^{(w)}_n[\Phi]$;
and for $w=0$, since $g_{\mathrm{w}}\cdot0=0$, it reads $Z^{(0)}_{n'}[\Phi\circ g_{\mathrm{w}}^{*}]=Z^{(0)}_n[\Phi]$.
\end{proof}

\subsection*{(g$'$)}

\begin{proof}[Proof of Proposition~50\,(g$'$)]
Let $\Phi\in\mathfrak{F}(n)$ and $\Phi^{-}(u):=\Phi(-u)$; then $\Phi^{-}\in\mathfrak{F}(n)$.
Since $\mathrm{fl}$ is linear, the map $\mathfrak{z}\mapsto-\mathfrak{z}$ carries $\mathrm{Mz}_w(n)$ onto $\mathrm{Mz}_{-w}(n)$.
For $A\in C^{1}(n)$, writing the sum over $\mathrm{Mz}_{-w}(n)$ as a sum over $-\mathfrak{z}$ with $\mathfrak{z}\in\mathrm{Mz}_w(n)$, and using (c), one finds
\begin{align*}
\mathrm{Psw}[\Phi^{-},\mathrm{Mz}_{-w}(n)](-A)&=\det\nolimits_n(D_{-A}+\boldsymbol{\mu})\cdot\sum_{\mathfrak{z}\in\mathrm{Mz}_w(n)}\Phi\left(-\left(d(-A)-r_q\mathfrak{z}\right)\right)\cdot\exp\left(-\frac{1}{2}\left\|d(-A)-r_q\mathfrak{z}\right\|^2\right)\\
&=\mathrm{Psw}[\Phi,\mathrm{Mz}_w(n)](A).
\end{align*}
The function $\mathrm{Psw}[\Phi^{-},\mathrm{Mz}_{-w}(n)]$ satisfies the hypotheses of the second assertion of Lemma~5, by (e.2) and (e.5).
By (e.4), that assertion and (e.4) once more, with all integrals extended over $[0,r_q)^{\mathrm{links}}$,
\begin{align*}
Z^{(-w)}_n[\Phi^{-}]&=\int\mathrm{d}A\,\mathrm{Psw}[\Phi^{-},\mathrm{Mz}_{-w}(n)](A)=\int\mathrm{d}A\,\mathrm{Psw}[\Phi^{-},\mathrm{Mz}_{-w}(n)](-A)\\
&=\int\mathrm{d}A\,\mathrm{Psw}[\Phi,\mathrm{Mz}_w(n)](A)=Z^{(w)}_n[\Phi].
\end{align*}
At $w=0$ this is $Z^{(0)}_n[\Phi^{-}]=Z^{(0)}_n[\Phi]$.
Summation over $w$, which converges absolutely by the first part of (e), gives $Z_n[\Phi^{-}]=Z_n[\Phi]$.
\end{proof}

\subsection*{(e), second part: positivity, order of the sectors, laws}

\begin{proof}[Proof of the remaining assertions of Proposition~50\,(e)]
Let $n_\mu\ge4$ and $w\in\mathbb{Z}^3$.

\emph{Step~1: $Z^{(w)}_n[1]>0$.}
The assertion \cite[Theorem~4.13\,(d)]{QEDII} holds for every real $\rho$ with $\rho^2\ne1$, every real value of the mass parameter of \cite[Definition~2.10]{QEDII}, all weights $w_i>0$, every $q>0$, every subgroup of $\mathbb{Z}^3$ of rank three and every unitary character of it.
It is applied with $0<\rho<1$, the standard weights, the mass parameter $M_{\mathrm{ch}}$, the group $\mathrm{Lt}(n)$ and the character that takes the value $-1$ at each $n_\mu e_\mu$, that is, on the torus $C(n)$.
By Proposition~46\,(e), the determinant of \cite[Definition~2.11]{QEDII} on $C(n)$ is $\det\nolimits_n(D_A+\boldsymbol{\mu})$.
The theorem thus asserts that for every $\mathfrak{z}\in\mathrm{Mz}(n)$
\begin{equation*}
\int_{[0,r_q)^{\mathrm{links}}}\mathrm{d}A\,\exp\left(-\frac{1}{2}\|dA+r_q\mathfrak{z}\|^2\right)\cdot\det\nolimits_n(D_A+\boldsymbol{\mu})>0.
\end{equation*}
By (e.4), $Z^{(w)}_n[1]$ is the sum of these integrals over $\mathfrak{z}\in\mathrm{Mz}_w(n)$.
All of them are non-negative by (a), and $\mathrm{Mz}_w(n)$ contains $\mathfrak{z}^{(w)}$ (Lemma~4\,(vii); the same cochain is used in \cite[Remark~13.13]{QEDII}).
Hence $Z^{(w)}_n[1]>0$.

\emph{Step~2: $Z^{(w)}_n[1]\le Z^{(0)}_n[1]$.}
The assertion \cite[Lemma~13.12\,(b)]{QEDII} holds on every torus $C(n)$ with $n_\mu\ge4$, of either parity, for $-1\le\rho\le1$, every real value of the mass parameter, all weights $w_i>0$ and every $q>0$:
the number $Z^{(w)}[1]$ of \cite[Definition~13.7]{QEDII} on $C(n)$ does not exceed the number $Z^{(0)}[1]$, for every $w$.
By Proposition~46\,(f+), at the mass parameter $M_{\mathrm{ch}}$ these numbers are $Z^{(w)}_n[1]$ and $Z^{(0)}_n[1]$.
The inequality also follows directly from the identity of \cite[Lemma~13.8\,(a)]{QEDII}, taken with $e_1$ as the time, together with (g), as follows.
In (2a)--(2d), $n$ is an arbitrary torus with $n_\mu\ge4$, and $w_{23}$ is an arbitrary integer.

\emph{(2a) The series.}
For $\mathrm{bl}=(\mathrm{bl}_2,\mathrm{bl}_3)\in[0,1]^2$ put
\begin{equation*}
\mathrm{gf}(\mathrm{bl}):=\sum_{w_{12},w_{13}\in\mathbb{Z}}\exp\left(2\pi i\left(\mathrm{bl}_2\cdot w_{12}+\mathrm{bl}_3\cdot w_{13}\right)\right)\cdot Z^{(w_{23},w_{13},w_{12})}_n[1].
\end{equation*}
The coefficients are non-negative, and their sum is at most $Z_n[1]$ by the first part of (e).
Hence the series converges absolutely and uniformly in $\mathrm{bl}$, the function $\mathrm{gf}$ is continuous, and the series may be integrated term by term.

\emph{(2b) The trace.}
The objects introduced below are those of \cite[Definitions~2.26 and~2.28 and Lemma~13.8]{QEDII} on $C(n)$, with $e_1$ as the time, at the parameters of Step~1 ($0<\rho<1$, the standard weights, the mass parameter $M_{\mathrm{ch}}$) and with the charge $q$.
The statements \cite[Definitions~2.26, 2.28, Lemmas~5.35, 5.40, Corollary~5.41, Proposition~5.43 and Lemma~13.8]{QEDII} are formulated for the tori $C(n)$ with $n_\mu\ge4$ of either parity (\cite[Definition~5.32]{QEDII}),
for $-1\le\rho\le1$ (the standing assumptions at the beginning of \cite[Section~5]{QEDII} and of \cite[Section~13.3]{QEDII}, and \cite[Definition~2.28]{QEDII}),
for every real value of the mass parameter, all weights $w_i>0$ and every $q>0$;
so they hold at these parameters.
Let $\mathrm{Lk}_s$ be the set of the links of the torus $\mathbb{Z}^2/(n_2\mathbb{Z}\times n_3\mathbb{Z})$.
The spaces $\mathrm{Hs}^{\mathrm{lk}}$ and $\mathrm{Hs}^{\mathrm{st}}$ are Hilbert spaces, each of them a direct sum of spaces indexed by the integers;
$\mathrm{Pr}_{\mathrm{jz}}$ denotes the projection of $\mathrm{Hs}^{\mathrm{lk}}$, and of $\mathrm{Hs}^{\mathrm{st}}$, onto the summand with the index $\mathrm{jz}\in\mathbb{Z}$.
Further, $E=(E_b)_{b\in\mathrm{Lk}_s}$ is a family of commuting self-adjoint operators, $P_G$ is a projection on $\mathrm{Hs}^{\mathrm{lk}}$, and
\begin{equation*}
K_\beta:=r_q^{n_2n_3}\cdot\exp\left(-\frac{1}{2}|E-\beta|^2\right)P_G\qquad\text{for }\beta\in\mathbb{R}^{\mathrm{Lk}_s}.
\end{equation*}
The operator $M_0\colon\mathrm{Hs}^{\mathrm{st}}\to\mathrm{Hs}^{\mathrm{lk}}$ is a multiplication operator, and it is bounded:
by \cite[Definition~2.28]{QEDII} it multiplies by the product of a number in $(0,1]$ and the matrix function $\mathrm{Qs}$, and $\mathrm{Qs}$ is bounded by \cite[Lemma~5.35\,(a)]{QEDII}.
It is formed with a plaquette $p_0$, which is arbitrary here (\cite[Proposition~5.43\,(c)(iii)]{QEDII}).
The sets $L_2$ and $L_3$ are the two subsets of $\mathrm{Lk}_s$ of \cite[Lemma~13.8]{QEDII}, $\mathbf{1}_{L_2}$ and $\mathbf{1}_{L_3}$ are their indicator functions, and
\begin{equation*}
\beta(\mathrm{bl}):=q\left(\mathrm{bl}_2\cdot\mathbf{1}_{L_2}+\mathrm{bl}_3\cdot\mathbf{1}_{L_3}\right)\in\mathbb{R}^{\mathrm{Lk}_s}.
\end{equation*}
In \cite{QEDII} the objects $\mathrm{Lk}_s$, $\mathrm{Hs}^{\mathrm{lk}}$, $\mathrm{Hs}^{\mathrm{st}}$, $\mathrm{Pr}_{\mathrm{jz}}$, $M_0$, $\mathrm{Qs}$, $p_0$, $L_2$, $L_3$, $\mathrm{bl}$ and $\beta(\mathrm{bl})$ are written
$\Lambda^{(1)}_{\mathrm{s}}$, $\mathcal{H}^{\mathrm{l}}$, $\mathcal{H}^{\mathrm{s}}$, $P_{\mathrm{jz}}$, $\mathfrak{M}$, $\widehat{\mathfrak{N}}$, $\sigma_0$, $S_2$, $S_3$, $\vartheta$ and $\beta_\vartheta$;
the operator $M_0$ is not one of the hopping matrices $M_i$.

The identity of \cite[Lemma~13.8\,(a)]{QEDII} is now taken at $u_t=\mathrm{ve}_t=1$ for all $t$;
there these two functions are written $u_j$ and $v_j$, and the constant $1$ is admitted for both of them in \cite[Proposition~5.43\,(b)]{QEDII}.
By Proposition~46\,(f+), the left side of that identity is $\mathrm{gf}(\mathrm{bl})$.
On its right side $u_t$ acts as the identity, and each of the $n_1$ factors is $r_q^{n_2n_3}\cdot\mathrm{Ko}^{\mathrm{bl}}P_GM_0M_0^*$;
here $\mathrm{Ko}^{\mathrm{bl}}$ is the operator $\mathrm{Ko}_t$ of \cite[Proposition~5.43\,(b)]{QEDII} (written $K_j$ there), without local step data, formed with the function $\exp\left(i\langle\beta(\mathrm{bl}),\cdot\,\rangle\right)$ on $\mathbb{R}^{\mathrm{Lk}_s}$.
By \cite[Proposition~5.43\,(c)(i)]{QEDII}, $r_q^{n_2n_3}\cdot\mathrm{Ko}^{\mathrm{bl}}P_G=K_{\beta(\mathrm{bl})}$.
With $\beta:=\beta(\mathrm{bl})$, and with $\mathrm{Tr}$ the trace on $\mathrm{Hs}^{\mathrm{lk}}$, it follows that
\begin{equation}
\mathrm{gf}(\mathrm{bl})=(2\pi)^{|n|}\cdot\mathrm{Tr}\left(\mathrm{Pr}_{w_{23}}\cdot\left(K_\beta\cdot M_0\cdot M_0^*\right)^{n_1}\right).\tag{e.6}
\end{equation}

\emph{(2c) $\mathrm{gf}(\mathrm{bl})\ge0$.}
\begin{itemize}
\item[(i)] By \cite[Lemma~5.40\,(c)]{QEDII}, taken at the empty subset of $\mathrm{Lk}_s$ (for which the operator $K_{\beta,S}$ of \cite[Definition~2.28]{QEDII} is $K_\beta$):
the bounded Borel functions of $E$ commute with $P_G$;
$K_\beta\ge0$;
and $K_\beta^{\frac{1}{2}}=r_q^{\frac{n_2n_3}{2}}\cdot\exp\left(-\frac{1}{4}|E-\beta|^2\right)P_G$, a bounded self-adjoint operator.
\item[(ii)] By \cite[Corollary~5.41]{QEDII}, taken at the empty subset, the operator $\mathrm{Yr}:=K_\beta^{\frac{1}{2}}\cdot M_0\colon\mathrm{Hs}^{\mathrm{st}}\to\mathrm{Hs}^{\mathrm{lk}}$ is a Hilbert--Schmidt operator, and $\mathrm{Yr}^*=M_0^*\cdot K_\beta^{\frac{1}{2}}$.
So $\mathrm{Yr}\cdot\mathrm{Yr}^*$ is non-negative and of trace class.
\item[(iii)] All operators of \cite[Definition~2.28]{QEDII} are diagonal with respect to the two direct sums, as stated in that definition.
Hence $\mathrm{Pr}_{w_{23}}\cdot K_\beta^{\frac{1}{2}}=K_\beta^{\frac{1}{2}}\cdot\mathrm{Pr}_{w_{23}}$ and $\mathrm{Pr}_{w_{23}}\cdot M_0=M_0\cdot\mathrm{Pr}_{w_{23}}$, and the adjoint of the second relation is $M_0^*\cdot\mathrm{Pr}_{w_{23}}=\mathrm{Pr}_{w_{23}}\cdot M_0^*$.
Consequently $\mathrm{Pr}_{w_{23}}\cdot\mathrm{Yr}=\mathrm{Yr}\cdot\mathrm{Pr}_{w_{23}}$ and $\mathrm{Yr}^*\cdot\mathrm{Pr}_{w_{23}}=\mathrm{Pr}_{w_{23}}\cdot\mathrm{Yr}^*$, and $\mathrm{Pr}_{w_{23}}$ commutes with $\mathrm{Yr}\cdot\mathrm{Yr}^*$.
\end{itemize}
By (i) and (ii),
\begin{align*}
\left(K_\beta\cdot M_0\cdot M_0^*\right)^{n_1}&=\left(K_\beta^{\frac{1}{2}}\cdot\mathrm{Yr}\cdot M_0^*\right)^{n_1}=K_\beta^{\frac{1}{2}}\cdot\mathrm{Yr}\cdot\left(M_0^*\cdot K_\beta^{\frac{1}{2}}\cdot\mathrm{Yr}\right)^{n_1-1}\cdot M_0^*\\
&=K_\beta^{\frac{1}{2}}\cdot\mathrm{Yr}\cdot\left(\mathrm{Yr}^*\cdot\mathrm{Yr}\right)^{n_1-1}\cdot M_0^*=K_\beta^{\frac{1}{2}}\cdot\left(\mathrm{Yr}\cdot\mathrm{Yr}^*\right)^{n_1-1}\cdot\mathrm{Yr}\cdot M_0^*.
\end{align*}
The operator $\left(\mathrm{Yr}\cdot\mathrm{Yr}^*\right)^{n_1-1}\cdot\mathrm{Yr}\cdot M_0^*$ is of trace class, because $n_1-1\ge1$ and $\mathrm{Yr}\cdot M_0^*$ is bounded \cite[Theorem~VI.19\,(b)]{ReedSimon1980};
the operator $\mathrm{Pr}_{w_{23}}\cdot K_\beta^{\frac{1}{2}}$ is bounded.
If $O$ is bounded and $O'$ is of trace class, then $\mathrm{Tr}(OO')=\mathrm{Tr}(O'O)$ \cite[Theorem~VI.25\,(c)]{ReedSimon1980}.
This rule and (iii) give
\begin{align*}
\mathrm{Tr}\left(\mathrm{Pr}_{w_{23}}\cdot\left(K_\beta\cdot M_0\cdot M_0^*\right)^{n_1}\right)&=\mathrm{Tr}\left(\left(\mathrm{Yr}\cdot\mathrm{Yr}^*\right)^{n_1-1}\cdot\mathrm{Yr}\cdot M_0^*\cdot\mathrm{Pr}_{w_{23}}\cdot K_\beta^{\frac{1}{2}}\right)\\
&=\mathrm{Tr}\left(\left(\mathrm{Yr}\cdot\mathrm{Yr}^*\right)^{n_1-1}\cdot\mathrm{Yr}\cdot M_0^*\cdot K_\beta^{\frac{1}{2}}\cdot\mathrm{Pr}_{w_{23}}\right)=\mathrm{Tr}\left(\left(\mathrm{Yr}\cdot\mathrm{Yr}^*\right)^{n_1}\cdot\mathrm{Pr}_{w_{23}}\right).
\end{align*}
Since $\mathrm{Pr}_{w_{23}}$ is a projection that commutes with $\mathrm{Yr}\cdot\mathrm{Yr}^*$, one has
\begin{equation*}
\left(\mathrm{Yr}\cdot\mathrm{Yr}^*\right)^{n_1}\cdot\mathrm{Pr}_{w_{23}}=\left(\mathrm{Pr}_{w_{23}}\cdot\mathrm{Yr}\cdot\mathrm{Yr}^*\cdot\mathrm{Pr}_{w_{23}}\right)^{n_1}.
\end{equation*}
The operator $\mathrm{Pr}_{w_{23}}\cdot\mathrm{Yr}\cdot\mathrm{Yr}^*\cdot\mathrm{Pr}_{w_{23}}=\left(\mathrm{Pr}_{w_{23}}\cdot\mathrm{Yr}\right)\cdot\left(\mathrm{Pr}_{w_{23}}\cdot\mathrm{Yr}\right)^*$ is non-negative and of trace class, and so is its $n_1$-th power;
the trace of a non-negative operator of trace class is non-negative.
Together with (e.6) this gives
\begin{equation}
\mathrm{gf}(\mathrm{bl})=(2\pi)^{|n|}\cdot\mathrm{Tr}\left(\left(\mathrm{Pr}_{w_{23}}\cdot\mathrm{Yr}\cdot\mathrm{Yr}^*\cdot\mathrm{Pr}_{w_{23}}\right)^{n_1}\right)\ge0\qquad\text{for all }\mathrm{bl}\in[0,1]^2.\tag{e.7}
\end{equation}

\emph{(2d) Fourier coefficients.}
By (2a) and (e.7), for all $w_{12},w_{13}\in\mathbb{Z}$,
\begin{align}
Z^{(w_{23},w_{13},w_{12})}_n[1]&=\int_{[0,1]^2}\mathrm{d}\,\mathrm{bl}\,\mathrm{gf}(\mathrm{bl})\cdot\exp\left(-2\pi i\left(\mathrm{bl}_2\cdot w_{12}+\mathrm{bl}_3\cdot w_{13}\right)\right)\notag\\
&\le\int_{[0,1]^2}\mathrm{d}\,\mathrm{bl}\,|\mathrm{gf}(\mathrm{bl})|=\int_{[0,1]^2}\mathrm{d}\,\mathrm{bl}\,\mathrm{gf}(\mathrm{bl})=Z^{(w_{23},0,0)}_n[1].\tag{e.8}
\end{align}

\emph{(2e) The second direction.}
Let $g_{\mathrm{w}}:=(0,s_2)$, where $s_2$ is the transposition $e_1\leftrightarrow e_2$.
Then $\epsilon_\mu=1$ for all $\mu$, $1'=2$, $2'=1$, $3'=3$, and $n'=(n_2,n_1,n_3)$.
The rule of (g) gives $(g_{\mathrm{w}}\cdot w)_{21}=w_{12}$, $(g_{\mathrm{w}}\cdot w)_{23}=w_{13}$ and $(g_{\mathrm{w}}\cdot w)_{13}=w_{23}$, that is,
\begin{equation*}
g_{\mathrm{w}}\cdot(w_{23},w_{13},w_{12})=(w_{13},w_{23},-w_{12}).
\end{equation*}
By (g) with $\Phi=1$, by (e.8) on the torus $n'$, and by (g) at $w=0$,
\begin{equation*}
Z^{(w_{23},0,0)}_n[1]=Z^{(0,w_{23},0)}_{n'}[1]\le Z^{(0,0,0)}_{n'}[1]=Z^{(0,0,0)}_n[1].
\end{equation*}
Combined with (e.8) on $n$, this yields $Z^{(w)}_n[1]\le Z^{(0)}_n[1]$.

\emph{Step~3.}
All terms being non-negative, the first part of (e) gives
\begin{equation*}
Z^{(0)}_n[1]\le\sum_wZ^{(w)}_n[1]=Z_n[1]<\infty.
\end{equation*}

\emph{Step~4: laws.}
Let $\mathrm{mq}$ be the measure on $[0,r_q)^{\mathrm{links}}\times\mathrm{Mz}(n)$ that arises from the product of the Lebesgue measure and the counting measure by multiplication with the density
\begin{equation*}
\exp\left(-\frac{1}{2}\|dA+r_q\mathfrak{z}\|^2\right)\cdot\det\nolimits_n(D_A+\boldsymbol{\mu}).
\end{equation*}
The density is non-negative by (a), and the total mass of $\mathrm{mq}$ is $Z_n[1]$, a number in $(0,\infty)$ by Steps~1 and~3.
The map $(A,\mathfrak{z})\mapsto F^{\mathrm{u}}=dA+r_q\mathfrak{z}$ is a Borel map into $C^{2}(n)$.
The image of $Z_n[1]^{-1}\cdot\mathrm{mq}$ under this map is a Borel probability measure on $C^{2}(n)$, the law of $F^{\mathrm{u}}$;
by the formula for the integral with respect to an image measure, $E^{\mathrm{af}}_n\Phi=Z_n[\Phi]/Z_n[1]$ is the integral of $\Phi$ with respect to it, for every $\Phi\in\mathfrak{F}(n)$.
Let $\mathrm{mq}_0$ be the restriction of $\mathrm{mq}$ to $[0,r_q)^{\mathrm{links}}\times\mathrm{Mz}_0(n)$.
By (e.4) its total mass is $Z^{(0)}_n[1]$, a number in $(0,\infty)$ by Steps~1 and~3.
The image of $Z^{(0)}_n[1]^{-1}\cdot\mathrm{mq}_0$ under $(A,\mathfrak{z})\mapsto F^{\mathrm{u}}$ is a Borel probability measure on $C^{2}(n)$,
and by (e.4) $E^{\mathrm{av}}_n\Phi=Z^{(0)}_n[\Phi]/Z^{(0)}_n[1]$ is the integral of $\Phi$ with respect to it, for every $\Phi\in\mathfrak{F}(n)$.
In all four steps $\boldsymbol{\mu}$ is an arbitrary real number.
\end{proof}

\subsection*{(g$''$)}

\begin{proof}[Proof of Proposition~50\,(g$''$)]
Let $n'=n=(n,n,n)$ with an integer $n\ge4$, so that $\mathrm{Lt}(n)=n\mathbb{Z}^3$.
Every $r\in O_{\mathrm{cub}}$ maps $n\mathbb{Z}^3$ onto itself.
Hence clause (f) applies, with $n'=n$, to every $g_{\mathrm{w}}\in W_{\mathrm{af}}$, and it gives
\begin{equation*}
Z_n\left[\Phi\circ g_{\mathrm{w}}^{*}\right]=Z_n[\Phi],\qquad Z^{(0)}_n\left[\Phi\circ g_{\mathrm{w}}^{*}\right]=Z^{(0)}_n[\Phi]\qquad\left(\Phi\in\mathfrak{F}(n)\right).
\end{equation*}
By (e) one has $Z_n[1]>0$ and $Z^{(0)}_n[1]>0$.
Dividing by these two numbers, one obtains
\begin{equation*}
E^{\mathrm{af}}_n\left[\Phi\circ g_{\mathrm{w}}^{*}\right]=E^{\mathrm{af}}_n[\Phi],\qquad E^{\mathrm{av}}_n\left[\Phi\circ g_{\mathrm{w}}^{*}\right]=E^{\mathrm{av}}_n[\Phi]\qquad\left(\Phi\in\mathfrak{F}(n)\right).
\end{equation*}
Let $\Phi$ range over the indicator functions of the Borel subsets of $C^2(n)$.
The last two identities then state that each of the two laws of $F^{\mathrm{u}}$ of (e), Step~4, is its own image under the map $u\mapsto g_{\mathrm{w}}^{*}u$ of $C^2(n)$.

In the same manner clause (g$'$) yields
\begin{equation*}
Z_n\left[\Phi(-\,\cdot\,)\right]=Z_n[\Phi],\qquad Z^{(0)}_n\left[\Phi(-\,\cdot\,)\right]=Z^{(0)}_n[\Phi],
\end{equation*}
and division by $Z_n[1]$ and by $Z^{(0)}_n[1]$ gives
\begin{equation*}
E^{\mathrm{af}}_n\left[\Phi(-\,\cdot\,)\right]=E^{\mathrm{af}}_n[\Phi],\qquad E^{\mathrm{av}}_n\left[\Phi(-\,\cdot\,)\right]=E^{\mathrm{av}}_n[\Phi]\qquad\left(\Phi\in\mathfrak{F}(n)\right).
\end{equation*}
With the indicator functions of the Borel sets in the place of $\Phi$, it follows that each of the two laws is also its own image under $u\mapsto-u$.

For the converse, suppose that one of the two laws is its own image under $u\mapsto g_{\mathrm{w}}^{*}u$, or under $u\mapsto-u$, and let $E$ be the corresponding expectation, $E=E^{\mathrm{af}}_n$ or $E=E^{\mathrm{av}}_n$ as in the clause.
By (e), Step~4, $E[\Phi]$ is the integral of $\Phi$ with respect to this law.
The formula for the integral with respect to an image measure therefore returns
\begin{equation*}
E\left[\Phi\circ g_{\mathrm{w}}^{*}\right]=E[\Phi],\qquad\text{respectively}\qquad E\left[\Phi(-\,\cdot\,)\right]=E[\Phi],
\end{equation*}
for every $\Phi\in\mathfrak{F}(n)$.
Consequently the two formulations of the invariance in (g$''$), through the expectations and through the image laws, are equivalent.
\end{proof}

\subsection*{(h)}

\begin{proof}[Proof of Proposition~50\,(h)]
Let $n_\mu\ge1$, let $\mathbf{m}>0$, and let $a\in C^1(n)$ be real.
In this proof the letter $a$ denotes this cochain, and the translation part of an element of $W_{\mathrm{af}}$ is not written.

\emph{(h1) The assertions of the second sentence of (h) up to $N_a(0)=1$.}
They are: the denominator of $N_a$ is positive; $N_a$ is entire and even; for real $t$ the value $N_a(t)$ is real and
\begin{equation*}
N_a(t)\ge\frac{\mathbf{m}^{192|n|}}{\det_n\left(D_0+\mathbf{m}\right)}>0 .
\end{equation*}
Clause (a), taken with $B=0$ and $\boldsymbol{\mu}=\mathbf{m}$, gives $\det_n\left(D_0+\mathbf{m}\right)\ge\mathbf{m}^{192|n|}>0$.
The function $t\mapsto\det_n\left(D_{t\cdot a}+\mathbf{m}\right)$ is the composition of the linear map $\mathbb{C}\to C^1(n)\otimes\mathbb{C}$, $t\mapsto t\cdot a$, with the entire function of (b).
Hence $N_a$ is entire.
Since $D_{0\cdot a}=D_0$, one has $N_a(0)=1$.
Clause (c), taken with $B=t\cdot a$, gives $N_a(-t)=N_a(t)$.
For real $t$ the cochain $t\cdot a$ is real.
Then clause (a), taken with $B=t\cdot a$ and $\boldsymbol{\mu}=\mathbf{m}$, shows that $\det_n\left(D_{t\cdot a}+\mathbf{m}\right)$ is real and not smaller than $\mathbf{m}^{192|n|}$.
It follows that $N_a(t)$ is real for real $t$ and satisfies the lower bound displayed above.

\emph{(h2)}
Let $g_{\mathrm{w}}\in W_{\mathrm{af}}$ have the linear part $r$, with $r\,\mathrm{Lt}(n)=\mathrm{Lt}(n')$.
The map $(g_{\mathrm{w}})_{*}$ is linear: $(g_{\mathrm{w}})_{*}(t\cdot a)=t\cdot(g_{\mathrm{w}})_{*}a$ and $(g_{\mathrm{w}})_{*}0=0$.
By Lemma~1\,(iii), $(g_{\mathrm{w}})_{*}a$ is a real element of $C^1(n')$.
Clause (d), taken with $B=t\cdot a$ and with $B=0$, and with $\boldsymbol{\mu}=\mathbf{m}$, gives for the function $N_{(g_{\mathrm{w}})_{*}a}$, formed on $n'$,
\begin{equation*}
N_{(g_{\mathrm{w}})_{*}a}(t)=\frac{\det_{n'}\left(D_{t\cdot(g_{\mathrm{w}})_{*}a}+\mathbf{m}\right)}{\det_{n'}\left(D_0+\mathbf{m}\right)}=\frac{\det_n\left(D_{t\cdot a}+\mathbf{m}\right)}{\det_n\left(D_0+\mathbf{m}\right)}=N_a(t)\qquad\left(t\in\mathbb{C}\right).
\end{equation*}

\emph{(h3) The numbers of Section~I.7.6.}
Let $0<\eta\le1$, let $R=n\cdot\eta$ with an integer $n\ge1$, let $e>0$, and let $q=e\cdot\eta^{\frac{1}{2}}$ (Section~I.2.5); the torus is $(n,n,n)$.
For a probe $\beta'$ let
\begin{equation*}
u[\beta'],\qquad a_\eta[\beta']=d_\eta^{*}u[\beta'],\qquad\hat{\mathfrak{a}}[\beta']:=\eta^{\frac{1}{2}}\cdot a_\eta[\beta']
\end{equation*}
be the cochains of Section~I.7.6 formed with $\beta'$, written in the form of Section~G5.1.3; for $n=1$ the same displays are used, as the last sentence of (h) prescribes.
Thus
\begin{equation*}
u[\beta']_{\mu\nu}(c)=\sum_{\mathrm{kz}\in\mathbb{Z}^3}\left(J_\eta\beta'_{\mu\nu}\right)\left(c+\frac{\eta}{2}\cdot\left(e_\mu+e_\nu\right)+R\cdot\mathrm{kz}\right)\qquad\left(c\in\eta\mathbb{Z}^3\right),
\end{equation*}
where only finitely many terms of the sum are non-zero.
Further $d_\eta=\eta^{-1}\cdot d$, and $d_\eta^{*}$ is the adjoint of $d_\eta$ with respect to $\langle\cdot,\cdot\rangle_\eta$ (Section~I.7.5; Section~G5.1.3; \cite[Section~7.1.1]{QEDI}).
Cochains on $\Lambda_{\eta,R}$ are read in lattice units through $c\mapsto c/\eta$ (Section~I.7.5).
Put
\begin{equation*}
a:=\frac{\hat{\mathfrak{a}}[\beta']}{e},\qquad\mathbf{m}:=\bar{m}\cdot\eta>0 .
\end{equation*}
The free ratio formed with $\hat{\mathfrak{a}}[\beta']$ is then
\begin{equation*}
N_0(t)=\frac{\det_n\left(D_{(t/e)\cdot\hat{\mathfrak{a}}[\beta']}+\bar{m}\cdot\eta\right)}{\det_n\left(D_0+\bar{m}\cdot\eta\right)}=N_a(t).
\end{equation*}
The numbers
\begin{equation*}
c_m=N_0^{(m)}(0),\qquad a_{\mathrm{ch}}=c_2,\qquad b_{\mathrm{ch}}=\left.\frac{\mathrm{d}^4}{\mathrm{d}t^4}\log N_0(t)\right|_{t=0}
\end{equation*}
are determined by the function $N_0$ alone; here $\log N_0$ is holomorphic near $0$, because $N_0(0)=1$.
Therefore these numbers do not change when $\hat{\mathfrak{a}}[\beta']$ is replaced by a cochain with the same function $N_a$.

\emph{Mirror copies} (the setting of Section~I.6.4).
Here $\eta=L^{-N}$, $e\cdot h^{\frac{1}{2}}\le e_1$ and $N\ge N_0(e)$, so that $\mathrm{per}=n_*/L$ is defined (Section~I.6.4).
The side $n$ of the torus remains arbitrary.
In particular $\mathrm{per}$ need not divide $n$: the argument needs only $r_{\mathfrak{c}}(n\mathbb{Z}^3)=n\mathbb{Z}^3$ and $a_{\mathfrak{c}}\in\mathbb{Z}^3$.
By Section~I.6.4 the map $\mathsf{g}_{\mathfrak{c}}$ acts on the coordinate $\mu$ either by $y_\mu\mapsto i\cdot\mathrm{per}+y_\mu$ (if $\mathfrak{c}_\mu=2i$) or by $y_\mu\mapsto(i+1)\cdot\mathrm{per}-y_\mu$ (if $\mathfrak{c}_\mu=2i+1$), and $\mathrm{per}$ is an integer (last row of the table at the end of this chapter).
Hence $\mathsf{g}_{\mathfrak{c}}=(a_{\mathfrak{c}},r_{\mathfrak{c}})$, where $r_{\mathfrak{c}}$ lies in the subgroup $\{\pm1\}^3$ of $O_{\mathrm{cub}}$ and $a_{\mathfrak{c}}\in\mathbb{Z}^3$.
So $\mathsf{g}_{\mathfrak{c}}\in W_{\mathrm{af}}$, and $r_{\mathfrak{c}}(n\mathbb{Z}^3)=n\mathbb{Z}^3$.
One has $\hat{\mathfrak{a}}_{\mathfrak{c}}=(\mathsf{g}_{\mathfrak{c}})_{*}\hat{\mathfrak{a}}$ and $(\mathsf{g}_{\mathfrak{c}})_{*}\left(\hat{\mathfrak{a}}/e\right)=\hat{\mathfrak{a}}_{\mathfrak{c}}/e$.
By (h2) with $n'=n$, the function $N_0$ formed with $\hat{\mathfrak{a}}_{\mathfrak{c}}$ coincides with the function $N_0$ formed with $\hat{\mathfrak{a}}$.

\emph{Placements.}
Let $\mathrm{vd}\in\eta\mathbb{Z}^3$, and let $g_{\mathrm{w},\mathrm{vd}}:=(\mathrm{vd}/\eta,1)\in W_{\mathrm{af}}$; in lattice units it is the translation $c\mapsto c+\mathrm{vd}/\eta$.
The form $\tau_{\mathrm{vd}}\beta'$, given by $(\tau_{\mathrm{vd}}\beta')(y)=\beta'(y-\mathrm{vd})$, is a probe (Section~I.7.5; \cite[Section~7.1.3]{QEDI}).
For a function $\mathrm{fs}$ on $\mathbb{R}^3$ one has
\begin{equation*}
\left(J_\eta\tau_{\mathrm{vd}}\mathrm{fs}\right)(y)=\eta^{-3}\int_{\left[-\frac{\eta}{2},\frac{\eta}{2}\right]^3}\mathrm{d}y'\,\mathrm{fs}\left(y+y'-\mathrm{vd}\right)=\left(J_\eta\mathrm{fs}\right)(y-\mathrm{vd}).
\end{equation*}
Consequently
\begin{equation*}
u[\tau_{\mathrm{vd}}\beta']_{\mu\nu}(c)=u[\beta']_{\mu\nu}(c-\mathrm{vd}),\qquad\text{that is,}\qquad u[\tau_{\mathrm{vd}}\beta']=(g_{\mathrm{w},\mathrm{vd}})_{*}u[\beta'];
\end{equation*}
indeed, for a translation the sign $\epsilon\left(g_{\mathrm{w},\mathrm{vd}}^{-1},\sigma\right)$ equals $1$ for every cell $\sigma$, and $g_{\mathrm{w},\mathrm{vd}}^{-1}$ maps the plaquette $(c;\mu\nu)$ of $\eta\mathbb{Z}^3$ onto the plaquette $(c-\mathrm{vd};\mu\nu)$.
Next, $d_\eta=\eta^{-1}\cdot d$ (Section~I.7.5; Section~G5.1.3), and in every degree the pairing $\langle\cdot,\cdot\rangle_\eta$ is $\eta^3$ times the counting pairing (Section~I.7.5).
Hence $d_\eta^{*}=\eta^{-1}\cdot d^{*}$, and by Lemma~1\,(v) this operator commutes with $(g_{\mathrm{w},\mathrm{vd}})_{*}$.
It follows that
\begin{equation*}
\hat{\mathfrak{a}}[\tau_{\mathrm{vd}}\beta']=(g_{\mathrm{w},\mathrm{vd}})_{*}\hat{\mathfrak{a}}[\beta'].
\end{equation*}
Now let $v$ and $\tilde{v}$ be two vectors of $\eta\mathbb{Z}^3$.
Since $\tau_{\tilde{v}}\beta=\tau_{\tilde{v}-v}\tau_v\beta$, the last identity gives
\begin{equation*}
\hat{\mathfrak{a}}[\tau_{\tilde{v}}\beta]=(g_{\mathrm{w},\tilde{v}-v})_{*}\hat{\mathfrak{a}}[\tau_v\beta],
\end{equation*}
and by (h2) with $n'=n$ the two cochains have the same function $N_0$.

By (h1) the function $N_0$ is even, so that $c_1=c_3=0$ and, as $t\to0$,
\begin{align*}
N_0(t)&=1+\frac{c_2}{2}\cdot t^2+\frac{c_4}{24}\cdot t^4+O\left(t^6\right),\\
\log N_0(t)&=\frac{c_2}{2}\cdot t^2+\left(\frac{c_4}{24}-\frac{c_2^2}{8}\right)\cdot t^4+O\left(t^6\right).
\end{align*}
Hence $b_{\mathrm{ch}}=c_4-3c_2^2$, in agreement with Section~I.7.6.
Since $N_0$ is real on the real axis, all the numbers $c_m$ are real.
\end{proof}

\subsection*{Conditions on the constants}

The proofs of this chapter rest on the following conditions on the constants.

\begin{center}
\begin{tabular}{p{0.19\textwidth}p{0.48\textwidth}p{0.21\textwidth}}
\hline
condition & where it is used & where the constant is fixed\\
\hline
$\rho\ge0$
& (a4): $\mathrm{Re}\left(D_B+\boldsymbol{\mu}\right)=\boldsymbol{\mu}+\frac{\rho}{2}\cdot\nabla_B^{*}\nabla_B\ge\boldsymbol{\mu}$; through (a4) also (a7) and (h1).
(a5), through (I.m6): the norm of $M_i(x)$ is $\frac{w_i}{2}\cdot(1+\rho)$ (for $\rho<0$ it is $\frac{w_i}{2}\cdot(1-\rho)$); through (a5) also $\mathrm{Dmax}$, (a.1), (e.2), (e.3) and the upper bound for $Z_n[1]$ in (e).
The same number enters the bound of Proposition~38.1\,(d), on which (b) rests.
& Section~I.8.2, line (C0); Proposition~14.2\,(a), line (C0): $\rho\in(0,1)$\\
\hline
$\rho^2\neq1$
& (e), Step~1: \cite[Theorem~4.13(d)]{QEDII}
& the same line\\
\hline
$-1\le\rho\le1$
& (e), Step~2: \cite[Lemma~13.12(b)]{QEDII}; in (2b) and (2c): \cite[Lemma~13.8]{QEDII}, \cite[Lemma~5.35]{QEDII}, \cite[Corollary~5.41]{QEDII}, \cite[Proposition~5.43]{QEDII}, \cite[Lemma~5.40]{QEDII}
& the same line\\
\hline
the standard weights satisfy $\sum_i w_i\cdot\hat{n}_i=0$, the relation (I.1.1)
& (0.2), hence (0.3): the diagonal term of $D_B$ is the scalar $\frac{\rho}{2}\cdot\mathrm{sw}$
& Section~I.8.2, line (C0); Proposition~14.2\,(a), line (C0): $w$\\
\hline
the standard weights satisfy $w_i>0$
& (a2); (a5), through (I.m6), and through (a5) what the first row lists; (e), Steps~1 and~2 (a hypothesis of \cite[Theorem~4.13(d)]{QEDII} and of \cite[Lemma~13.12(b)]{QEDII})
& the same line\\
\hline
$M$ is a power of $L$, and $M\ge L^2$ (here $L$ is an integer, and $h$ is a power of $L$ with an exponent $\ge0$: Section~I.0.5)
& (h3), Mirror copies: $T^{*}(e)=\log_L\left(M\cdot r_K(e)\right)-1$ is an integer $\ge1$ (Section~I.6.2), and $N\ge0$; hence $\mathrm{per}=n_*/L=h\cdot L^{T^{*}(e)+N-1}$ is an integer (Section~I.6.4), $a_{\mathfrak{c}}\in\mathbb{Z}^3$, and $\mathsf{g}_{\mathfrak{c}}\in W_{\mathrm{af}}$
& Section~I.8.2, line (C5); Proposition~14.2\,(b), row (M1)\\
\hline
\end{tabular}
\end{center}

The three conditions of the first three rows follow from the range $\rho\in(0,1)$ of line (C0).
The closure relation (I.1.1) and the positivity $w_i>0$, the two properties of the standard weights in the fourth and the fifth row, are properties of the numbers fixed in line (C0).
Among these numbers is the value $w_3=6$, which occurs in the bound of Proposition~38.1\,(d), on which (b) rests.
Since $0<\bar{m}\le1$, the number $\mathbf{m}=\bar{m}\cdot\eta$ of (h3) is positive, and $\log_L\left(1/\bar{m}\right)\ge0$;
hence the exponent $\left\lfloor\log_L\left(1/\bar{m}\right)\right\rfloor$ of the power $h$ of $L$ (Section~I.0.5) is $\ge0$, as required in the last row.

The sides ($n_\mu\ge1$ in (a)--(d) and (h), $n_\mu\ge4$ in (e)--(g$''$)), the numbers $q>0$ and $\mathbf{m}>0$, and the real number $\boldsymbol{\mu}$ are data of the statements.
So are, in the last sentence of (h), the numbers $e>0$, $0<\eta\le1$ and $\bar{m}$ (with $\mathbf{m}=\bar{m}\cdot\eta$), and, for the mirror copies, $\eta=L^{-N}$, $e\cdot h^{\frac{1}{2}}\le e_1$ and $N\ge N_0(e)$ (the setting of Section~I.6.4).

\section*{Chapter 47. Lattice probes on every torus; the cells on the even tori and the translation form: proofs of Propositions 47.1 and 47.2}

\subsection*{0. Conventions of this chapter; five lemmas}

\paragraph{Notation.}
In this chapter $\mathrm{sn}_\mu(p)$ denotes a Fourier multiplier, $\mathrm{Stp}(\cdot)$ a step form, $\mathrm{tq}$ and $\mathrm{tq}'$ real variables, $\mathrm{mi}$ a multi-index, $\mathrm{Gq}$ a function of one real variable, $\mathrm{Bq}$ a Borel set, $\mathrm{Dq}$ the side of a cube, $\mathrm{lo}_\mu$ and $\mathrm{hi}_\mu$ end points of intervals, $\mathfrak{l}_3$ the direction different from $\mu$ and $\nu$, $\mathrm{dg}$ a degree, and $\mathrm{cf}_1$, $\mathrm{cf}_2$ numbers.
As in Section~I.7.6, $z$ and $z'$ are vectors of $\mathbb{Z}^3$;
$\mathrm{gc}$ is a two-cochain, as in Proposition~31.1\,(d), and $\mathrm{cf}$ is a number, as in Section~G6.A3;
$y$, $y'$, $y_0$, $y_1$ are points of $\mathbb{R}^3$.
The symbol $a'$ stands for the vector field $(a'_\nu)$ that Proposition~47.1 forms with a probe (Subsection~0.2);
the constant $a'$ of (C4r) in Section~I.8.2 occurs neither in Propositions~47.1 and~47.2 nor in this chapter.
The symbol $\boldsymbol{\sigma}$ denotes a plaquette $\boldsymbol{\sigma}=(c;\mu\nu)$ (Sections~I.4.7 and~I.7.5) as well as the number $\boldsymbol{\sigma}=\|d\hat{\mathfrak{a}}\|$ (Section~I.7.6);
the form of the expression shows which of the two is meant.
Both meanings occur in (1.6) and in the paragraph \emph{Support} that follows it.

\paragraph{Cells of a torus.}
For Proposition~47.2\,(a), (c), the cells of a torus with the period lattice $n\mathbb{Z}^3$ are the orbits of $n\mathbb{Z}^3$ on the set $\left\{Q^{\circ}_{\mathfrak{c}}:\mathfrak{c}\in\mathbb{Z}^3\right\}$ of the open cells.
This is the notion of the cells of a torus of \cite[Definition~2.18]{QEDII}, where the width of a cell structure is a positive integer;
here it is applied with the width $\mathrm{hc}\in\frac{1}{2}\mathbb{Z}$, which (I.h1) in Section~I.6.4 admits for the coordinate frame.

\subsubsection*{0.1. Lattice, cells, cochains}

The conventions of this subsection follow Sections~I.1.1, I.7.5, G6.A1 and \cite[Section~7.1.1]{QEDI}.
Let $\eta>0$ and $R>0$, and let $n:=R/\eta$ be an integer $\ge1$;
in the statements one has $0<\eta\le1$ and $n\ge2$.
Put $\Lambda_{\eta,R}=\eta\mathbb{Z}^3/R\mathbb{Z}^3$, and let $c$, $c'$ denote sites of $\eta\mathbb{Z}^3$.
A $j$-cochain on $\Lambda_{\eta,R}$ is an $R\mathbb{Z}^3$-periodic function on the $j$-cells of $\eta\mathbb{Z}^3$;
it is real unless it is called complex.
The bond $(c,c+\eta e_\nu)$ has the centre $c+\frac{\eta}{2}e_\nu$, and for a one-cochain $A$ one writes $A_\nu(c):=A(c,c+\eta e_\nu)$.
The plaquette $\boldsymbol{\sigma}=(c;\mu\nu)$, $\mu<\nu$, has the centre $\mathrm{ce}(\boldsymbol{\sigma})=c+\frac{\eta}{2}\left(e_\mu+e_\nu\right)$, and
\[
Q(\boldsymbol{\sigma})=\mathrm{ce}(\boldsymbol{\sigma})+\left[-\frac{\eta}{2},\frac{\eta}{2}\right)^3.
\]
A sum over $c\in\Lambda_{\eta,R}$ is the sum over a set of representatives of $\eta\mathbb{Z}^3$ modulo $R\mathbb{Z}^3$, a set with $n^3$ elements;
for a periodic summand the sum does not depend on the choice of the representatives, and it is invariant under the substitution $c\mapsto c+c'$.
For a zero-cochain $\boldsymbol{\lambda}$, one-cochains $A$, $A'$ and two-cochains $u$, $u'$ put
\begin{equation}\tag{0.1}
\begin{aligned}
(d\boldsymbol{\lambda})_\nu(c)&=\boldsymbol{\lambda}(c+\eta e_\nu)-\boldsymbol{\lambda}(c),\\
(dA)_{\mu\nu}(c)&=A_\mu(c)+A_\nu(c+\eta e_\mu)-A_\mu(c+\eta e_\nu)-A_\nu(c)\qquad(\mu<\nu),\\
d_\eta&:=\eta^{-1}\cdot d,
\end{aligned}
\end{equation}
and
\begin{gather*}
\langle A,A'\rangle_\eta=\eta^3\sum_{c\in\Lambda_{\eta,R}}\sum_{\nu}A_\nu(c)\cdot A'_\nu(c),
\qquad
\langle u,u'\rangle_\eta=\eta^3\sum_{c\in\Lambda_{\eta,R}}\sum_{\mu<\nu}u_{\mu\nu}(c)\cdot u'_{\mu\nu}(c),\\
|u|_\eta^2=\eta^3\sum_{c\in\Lambda_{\eta,R}}\sum_{\mu<\nu}\left|u_{\mu\nu}(c)\right|^2.
\end{gather*}
The operator $d_\eta^*$ is the adjoint of $d_\eta$ for these pairings:
$\langle d_\eta A,u\rangle_\eta=\langle A,d_\eta^*u\rangle_\eta$.
The bijection $c\mapsto c/\eta$ of $\Lambda_{\eta,R}$ onto $\mathbb{Z}^3/n\mathbb{Z}^3$ identifies the cochains of the two lattices (Section~I.7.5);
under this identification (0.1) is the coboundary of Section~I.1.1 in lattice units.
The counting pairing and the counting norm are written $\langle\cdot,\cdot\rangle$ and $\|\cdot\|$;
thus $\langle u,u'\rangle_\eta=\eta^3\cdot\langle u,u'\rangle$.

\subsubsection*{0.2. Forms}

A two-form $u$ on $\mathbb{R}^3$ has the components $u_{\mu\nu}$, $\mu<\nu$, and one sets $u_{\nu\mu}:=-u_{\mu\nu}$ and $u_{\mu\mu}:=0$;
the same convention applies to two-cochains.
The support $\operatorname{supp}u$ is the union of the supports of the components.
For a square integrable two-form $u$,
\[
\|u\|_2^2=\sum_{\mu<\nu}\int_{\mathbb{R}^3}\mathrm{d}y\,\left|u_{\mu\nu}(y)\right|^2;
\]
for $u\in\mathcal{S}$ this equals $\|u\|_{H^0}^2$, by Plancherel's theorem.
With this convention on the indices one has, for a probe $\beta$ (Section~I.7.6) and $0<\eta\le1$,
\[
a_\nu=-\sum_{\mu\ne\nu}\partial_\mu\beta_{\mu\nu},
\qquad
a'_\nu=-\sum_{\mu\ne\nu}\delta_{\eta,\mu}\beta_{\mu\nu},
\qquad
f_{\mu\nu}=\partial_\mu a_\nu-\partial_\nu a_\mu
\]
(the sign convention for $d$ is that of Section~I.9.7), and
\begin{equation}\tag{0.2}
\begin{aligned}
f_{\mu\nu}&=-\sum_{\mathfrak{l}\ne\nu}\partial_\mu\partial_{\mathfrak{l}}\beta_{\mathfrak{l}\nu}+\sum_{\mathfrak{l}\ne\mu}\partial_\nu\partial_{\mathfrak{l}}\beta_{\mathfrak{l}\mu},\\
f^{(\eta)}_{\mu\nu}&=-\sum_{\mathfrak{l}\ne\nu}\delta_{\eta,\mu}\delta_{\eta,\mathfrak{l}}\beta_{\mathfrak{l}\nu}+\sum_{\mathfrak{l}\ne\mu}\delta_{\eta,\nu}\delta_{\eta,\mathfrak{l}}\beta_{\mathfrak{l}\mu}.
\end{aligned}
\end{equation}
Each of the two right sides consists of four terms.
For $\mu<\nu$, with $\mathfrak{l}_3$ the third direction, these terms contain the components $\beta_{\mu\nu}$ (the term $\mathfrak{l}=\mu$ of the first sum), $\beta_{\mathfrak{l}_3\nu}$, $\beta_{\nu\mu}$ (the term $\mathfrak{l}=\nu$ of the second sum) and $\beta_{\mathfrak{l}_3\mu}$.
When the probe has to be named, one writes $a'[\beta]$, $\tilde{a}_\eta[\beta]$, $f^{(\eta)}[\beta]$, $u[\beta]$, $a_\eta[\beta]$, $\hat{\mathfrak{a}}[\beta]$.

\subsubsection*{0.3. Translations}

For $b\in\mathbb{R}^3$ let $(\tau_b\mathrm{fs})(y):=\mathrm{fs}(y-b)$ (Section~I.7.5);
on forms $\tau_b$ acts componentwise.
For $v\in\eta\mathbb{Z}^3$ let $g_{\mathrm{w}}=\left(\frac{v}{\eta},1\right)\in W_{\mathrm{af}}$ be the translation by $v/\eta$ (Section~I.4.7; lattice units).
In the labelling of the cells by $\eta\mathbb{Z}^3$ (Subsection~0.1) it maps a cell $\boldsymbol{\sigma}$ to the cell $\boldsymbol{\sigma}+v$ with the same orientation.
Hence $\epsilon(g_{\mathrm{w}},\boldsymbol{\sigma})=1$, and
\[
(g_{\mathrm{w}}^*u)(\boldsymbol{\sigma})=u(\boldsymbol{\sigma}+v),
\qquad
(g_{\mathrm{w}*}u)(\boldsymbol{\sigma})=\left((g_{\mathrm{w}}^{-1})^*u\right)(\boldsymbol{\sigma})=u(\boldsymbol{\sigma}-v)
\]
for cochains $u$ of every degree.
By (0.1), $g_{\mathrm{w}*}d=d\,g_{\mathrm{w}*}$.
The substitution $c\mapsto c+v$ in the sums yields the first two of the following formulas, and the definition of $g_{\mathrm{w}*}$ yields the third:
\begin{equation}\tag{0.3}
\begin{aligned}
&\langle g_{\mathrm{w}*}u,u'\rangle_\eta=\langle u,g_{\mathrm{w}}^*u'\rangle_\eta,
\qquad
|g_{\mathrm{w}*}u|_\eta=|u|_\eta,\\
&\left(\frac{v}{\eta},1\right)_*\left(\frac{v'}{\eta},1\right)_*=\left(\frac{v+v'}{\eta},1\right)_*\qquad(v,v'\in\eta\mathbb{Z}^3).
\end{aligned}
\end{equation}
Since $\langle\cdot,\cdot\rangle_\eta=\eta^3\cdot\langle\cdot,\cdot\rangle$ (Subsection~0.1), the first two formulas of (0.3) hold for the counting pairing and the counting norm as well.
Moreover $g_{\mathrm{w}}^*=(g_{\mathrm{w}}^{-1})_*$ and $g_{\mathrm{w}}^{-1}=\left(-\frac{v}{\eta},1\right)$;
therefore $\|g_{\mathrm{w}}^*u\|=\|u\|$.

\subsubsection*{0.4. Periodisation}

For a bounded Borel function $\mathrm{fs}$ on $\mathbb{R}^3$ with compact support put
\[
\mathrm{fs}^{R}(y):=\sum_{z\in\mathbb{Z}^3}\mathrm{fs}(y+R\cdot z).
\]
For $y$ in a bounded set only finitely many terms differ from $0$, and $\mathrm{fs}^{R}$ is $R\mathbb{Z}^3$-periodic.
The value of $\mathrm{fs}$ at a point of $\mathbb{R}^3/R\mathbb{Z}^3$ is that of $\mathrm{fs}^{R}$ (the rule of Section~I.7.6).
Accordingly, the definitions of Sections~I.7.6 and~I.7.5 read
\begin{equation}\tag{0.4}
\begin{aligned}
&u_{\mu\nu}(c)=\left(J_{\eta}\beta_{\mu\nu}\right)^{R}\left(c+\frac{\eta}{2}\left(e_\mu+e_\nu\right)\right),
\qquad
a_\eta=d_\eta^*u,\\
&f'_\eta(\boldsymbol{\sigma})=\eta^{-3}\int_{Q(\boldsymbol{\sigma})}\mathrm{d}y\,\left(f'_{\mu\nu}\right)^{R}(y)\qquad\left(\boldsymbol{\sigma}=(c;\mu\nu)\right).
\end{aligned}
\end{equation}
For an $R\mathbb{Z}^3$-periodic locally square integrable two-form $u'$ put
\[
\|u'\|_{L^2(\mathbb{R}^3/R\mathbb{Z}^3)}^2:=\sum_{\mu<\nu}\int_{y_1+[0,R)^3}\mathrm{d}y\,\left|u'_{\mu\nu}(y)\right|^2 .
\]
The right side does not depend on $y_1\in\mathbb{R}^3$:
the integral of an $R$-periodic function of one variable over $[\mathrm{tq},\mathrm{tq}+R)$ does not depend on $\mathrm{tq}$, and this is applied in each coordinate.

\subsubsection*{0.5. The cube}

Let $[\mathrm{lo}_\mu,\mathrm{hi}_\mu]$, $\mu=1,2,3$, be the smallest intervals with $\operatorname{supp}\beta\subset\prod_\mu[\mathrm{lo}_\mu,\mathrm{hi}_\mu]$.
The projection of $\operatorname{supp}\beta+[-2,2]^3$ to the axis $\mu$ has the least element $\mathrm{lo}_\mu-2$ and the largest element $\mathrm{hi}_\mu+2$.
Hence a closed coordinate cube $y_1+[0,\mathrm{Dq}]^3$ contains $\operatorname{supp}\beta+[-2,2]^3$ if and only if
\[
y_{1,\mu}\le\mathrm{lo}_\mu-2
\quad\text{and}\quad
\mathrm{hi}_\mu+2\le y_{1,\mu}+\mathrm{Dq}
\qquad\text{for all }\mu,
\]
and such a $y_1$ exists if and only if $\mathrm{Dq}\ge\max_\mu\left(\mathrm{hi}_\mu-\mathrm{lo}_\mu\right)+4$.
The same reasoning with $0$ in the place of $2$ gives $D_{\mathrm{be}}=\max_\mu\left(\mathrm{hi}_\mu-\mathrm{lo}_\mu\right)$.
Therefore $D_\beta=D_{\mathrm{be}}+4$.
Since $\beta$ is continuous and $\beta\ne0$, the set $\operatorname{supp}\beta$ contains a ball;
so $D_{\mathrm{be}}>0$ and $D_\beta>4$.
Of the cube $Q_\beta=y_\beta+\left[-\frac{D_\beta}{2},\frac{D_\beta}{2}\right]^3$ only the following two properties are needed:
\begin{itemize}
\item[(cu1)] $Q_\beta$ is a closed coordinate cube of side $D_\beta$;
\item[(cu2)] $\operatorname{supp}\beta+[-2,2]^3$ is contained in $Q_\beta$.
\end{itemize}

\begin{named}{Lemma 1}[The operators $\delta_{\eta,\mu}$ and $J_\eta$]
Let $\eta>0$ and $\mathrm{fs}\in C^\infty(\mathbb{R}^3)$.
\begin{itemize}
\item[(a)] The operators $\delta_{\eta,\mu}$, $\delta_{\eta,\nu}$, $J_\eta$, $\partial_{\mathfrak{l}}$ and $\tau_b$ ($b\in\mathbb{R}^3$) commute pairwise on $C^\infty(\mathbb{R}^3)$.
They map real functions to real functions, and $C_c^\infty(\mathbb{R}^3)$ into itself.
\item[(b)] If $\mathrm{fs}$ has compact support, then
\[
\operatorname{supp}\delta_{\eta,\mu}\mathrm{fs}\subset\operatorname{supp}\mathrm{fs}+\left[-\frac{\eta}{2},\frac{\eta}{2}\right]e_\mu
\qquad\text{and}\qquad
\operatorname{supp}J_\eta\mathrm{fs}\subset\operatorname{supp}\mathrm{fs}+\left[-\frac{\eta}{2},\frac{\eta}{2}\right]^3.
\]
\item[(c)] One has
\[
\delta_{\eta,\mu}\mathrm{fs}(y)=\int_{-\frac{1}{2}}^{\frac{1}{2}}\mathrm{d}\mathrm{tq}\,\left(\partial_\mu\mathrm{fs}\right)\left(y+\mathrm{tq}\cdot\eta e_\mu\right).
\]
If the derivatives that occur are bounded, then $|\delta_{\eta,\mu}\mathrm{fs}|_\infty\le|\partial_\mu\mathrm{fs}|_\infty$ and $|J_\eta\mathrm{fs}|_\infty\le|\mathrm{fs}|_\infty$.
\item[(d)] One has
\[
\left|\delta_{\eta,\mu}\mathrm{fs}-\partial_\mu\mathrm{fs}\right|_\infty\le\frac{\eta^2}{24}\left|\partial_\mu^3\mathrm{fs}\right|_\infty
\qquad\text{and}\qquad
\left|J_\eta\mathrm{fs}-\mathrm{fs}\right|_\infty\le\frac{9}{8}\eta^2\left|D^2\mathrm{fs}\right|_\infty.
\]
\item[(e)] For $\mathrm{fs}\in C_c^\infty(\mathbb{R}^3)$,
\[
\widehat{\partial_\mu\mathrm{fs}}(p)=ip_\mu\,\widehat{\mathrm{fs}}(p),
\qquad
\widehat{\delta_{\eta,\mu}\mathrm{fs}}(p)=i\,\mathrm{sn}_\mu(p)\,\widehat{\mathrm{fs}}(p),
\qquad\text{where}\quad
\mathrm{sn}_\mu(p):=\frac{2}{\eta}\sin\left(\frac{\eta p_\mu}{2}\right);
\]
moreover $|\mathrm{sn}_\mu(p)|\le|p_\mu|$ and $|\mathrm{sn}_\mu(p)-p_\mu|\le\frac{\eta^2|p_\mu|^3}{24}$, and for all directions $\mu$, $\mathfrak{l}$, equal or not,
\begin{equation}\tag{0.5}
\left|\mathrm{sn}_\mu(p)\cdot\mathrm{sn}_{\mathfrak{l}}(p)-p_\mu p_{\mathfrak{l}}\right|\le\frac{\eta^2}{12}|p|_2^4.
\end{equation}
\end{itemize}
\end{named}

\begin{proof}
(a) One has $\delta_{\eta,\mu}=\eta^{-1}\cdot\left(\tau_{-\eta e_\mu/2}-\tau_{\eta e_\mu/2}\right)$.
Translations commute with one another and with $\partial_{\mathfrak{l}}$;
hence the operators $\delta_{\eta,\mu}$, $\tau_b$ and $\partial_{\mathfrak{l}}$ commute.
Further,
\[
(J_\eta\tau_b\mathrm{fs})(y)=\eta^{-3}\int_{\left[-\frac{\eta}{2},\frac{\eta}{2}\right]^3}\mathrm{d}y'\,\mathrm{fs}(y+y'-b)=(J_\eta\mathrm{fs})(y-b),
\]
so $J_\eta$ commutes with the translations and therefore with the $\delta_{\eta,\mu}$.
The integrand $\mathrm{fs}(y+y')$ is of class $C^\infty$ in $(y,y')$, and $y'$ ranges over a compact set;
so $J_\eta\mathrm{fs}$ is differentiable with $\partial_{\mathfrak{l}}J_\eta\mathrm{fs}=J_\eta\partial_{\mathfrak{l}}\mathrm{fs}$, and by induction $J_\eta\mathrm{fs}\in C^\infty(\mathbb{R}^3)$.
If $\mathrm{fs}$ is real, so are $\tau_b\mathrm{fs}$ and $\partial_{\mathfrak{l}}\mathrm{fs}$;
$\delta_{\eta,\mu}$ is a real linear combination of two translations, and $J_\eta$ is an average with the real weight $\eta^{-3}\,\mathrm{d}y'$.
If $\operatorname{supp}\mathrm{fs}$ is compact, then $\operatorname{supp}\tau_b\mathrm{fs}=b+\operatorname{supp}\mathrm{fs}$ is compact, so is $\operatorname{supp}\partial_{\mathfrak{l}}\mathrm{fs}$, a closed subset of $\operatorname{supp}\mathrm{fs}$, and $\operatorname{supp}\delta_{\eta,\mu}\mathrm{fs}$, $\operatorname{supp}J_\eta\mathrm{fs}$ are compact by (b), whose proof below does not rely on (a).

(b) If $\delta_{\eta,\mu}\mathrm{fs}(y)\ne0$, then one of the points $y+\frac{\eta}{2}e_\mu$, $y-\frac{\eta}{2}e_\mu$ lies in $\operatorname{supp}\mathrm{fs}$, and $y$ lies in $\operatorname{supp}\mathrm{fs}+\left[-\frac{\eta}{2},\frac{\eta}{2}\right]e_\mu$.
If $J_\eta\mathrm{fs}(y)\ne0$, then there is a $y'\in\left[-\frac{\eta}{2},\frac{\eta}{2}\right]^3$ with $y+y'\in\operatorname{supp}\mathrm{fs}$, and $y$ lies in $\operatorname{supp}\mathrm{fs}+\left[-\frac{\eta}{2},\frac{\eta}{2}\right]^3$.
Both sets are sums of two compact sets, hence closed;
therefore they contain the closures of $\{\delta_{\eta,\mu}\mathrm{fs}\ne0\}$ and of $\{J_\eta\mathrm{fs}\ne0\}$, respectively.

(c) One has $\frac{\mathrm{d}}{\mathrm{d}\mathrm{tq}}\,\mathrm{fs}(y+\mathrm{tq}\cdot\eta e_\mu)=\eta\cdot(\partial_\mu\mathrm{fs})(y+\mathrm{tq}\cdot\eta e_\mu)$;
integration over $\left[-\frac{1}{2},\frac{1}{2}\right]$ and division by $\eta$ give the formula.
The two bounds follow because the interval $\left[-\frac{1}{2},\frac{1}{2}\right]$ has length $1$ and the cube $\left[-\frac{\eta}{2},\frac{\eta}{2}\right]^3$ has volume $\eta^3$.

(d) For $\mathrm{Gq}\in C^2(\mathbb{R})$ Taylor's formula reads
\[
\mathrm{Gq}(\mathrm{tq})=\mathrm{Gq}(0)+\mathrm{tq}\cdot\mathrm{Gq}'(0)+\int_0^{\mathrm{tq}}\mathrm{d}\mathrm{tq}'\,\left(\mathrm{tq}-\mathrm{tq}'\right)\cdot\mathrm{Gq}''(\mathrm{tq}');
\]
indeed, both sides have the same value and the same first derivative at $0$, and they have the same second derivative.
Consequently
\[
\left|\mathrm{Gq}(\mathrm{tq})-\mathrm{Gq}(0)-\mathrm{tq}\cdot\mathrm{Gq}'(0)\right|\le\frac{\mathrm{tq}^2}{2}\sup\left|\mathrm{Gq}''\right|.
\]
For the first bound put $\mathrm{Gq}(\mathrm{tq}):=(\partial_\mu\mathrm{fs})(y+\mathrm{tq}\cdot\eta e_\mu)$, so that $\mathrm{Gq}''(\mathrm{tq})=\eta^2\cdot(\partial_\mu^3\mathrm{fs})(y+\mathrm{tq}\cdot\eta e_\mu)$.
By (c), and since $\int_{-\frac{1}{2}}^{\frac{1}{2}}\mathrm{d}\mathrm{tq}\,\mathrm{tq}=0$ and $\int_{-\frac{1}{2}}^{\frac{1}{2}}\mathrm{d}\mathrm{tq}\,\mathrm{tq}^2=\frac{1}{12}$,
\[
\left|\delta_{\eta,\mu}\mathrm{fs}(y)-\partial_\mu\mathrm{fs}(y)\right|
=\left|\int_{-\frac{1}{2}}^{\frac{1}{2}}\mathrm{d}\mathrm{tq}\,\left(\mathrm{Gq}(\mathrm{tq})-\mathrm{Gq}(0)-\mathrm{tq}\cdot\mathrm{Gq}'(0)\right)\right|
\le\frac{1}{2}\cdot\frac{1}{12}\cdot\eta^2\left|\partial_\mu^3\mathrm{fs}\right|_\infty.
\]
For the second bound let $y'\in\left[-\frac{\eta}{2},\frac{\eta}{2}\right]^3$ and put $\mathrm{Gq}(\mathrm{tq}):=\mathrm{fs}(y+\mathrm{tq}\cdot y')$.
Then
\[
\mathrm{Gq}'(0)=\sum_jy'_j\,\partial_j\mathrm{fs}(y),
\qquad
\mathrm{Gq}''(\mathrm{tq})=\sum_{j,j'}y'_jy'_{j'}\cdot(\partial_j\partial_{j'}\mathrm{fs})(y+\mathrm{tq}\cdot y'),
\]
and
\[
\left|\mathrm{Gq}''\right|\le\left(\sum_j|y'_j|\right)^2\left|D^2\mathrm{fs}\right|_\infty\le\left(\frac{3\eta}{2}\right)^2\left|D^2\mathrm{fs}\right|_\infty.
\]
At $\mathrm{tq}=1$ this gives
\[
\left|\mathrm{fs}(y+y')-\mathrm{fs}(y)-\sum_jy'_j\,\partial_j\mathrm{fs}(y)\right|\le\frac{9}{8}\eta^2\left|D^2\mathrm{fs}\right|_\infty.
\]
The integral of $y'_j$ over $\left[-\frac{\eta}{2},\frac{\eta}{2}\right]^3$ vanishes;
averaging over $y'$ yields the second bound.

(e) With the Fourier transform of Section~I.9.7,
\[
\widehat{\tau_b\mathrm{fs}}(p)=\int\mathrm{d}y\,\exp(-ip\cdot y)\,\mathrm{fs}(y-b)=\exp(-ip\cdot b)\,\widehat{\mathrm{fs}}(p).
\]
Hence
\[
\widehat{\delta_{\eta,\mu}\mathrm{fs}}(p)
=\eta^{-1}\cdot\left(\exp\left(\frac{i\eta p_\mu}{2}\right)-\exp\left(-\frac{i\eta p_\mu}{2}\right)\right)\cdot\widehat{\mathrm{fs}}(p)
=i\,\frac{2}{\eta}\sin\left(\frac{\eta p_\mu}{2}\right)\cdot\widehat{\mathrm{fs}}(p).
\]
The formula for $\widehat{\partial_\mu\mathrm{fs}}$ is an integration by parts.
With $\mathrm{tq}:=\eta p_\mu/2$ one has
\[
|\mathrm{sn}_\mu|=\frac{2}{\eta}|\sin\mathrm{tq}|\le\frac{2}{\eta}|\mathrm{tq}|=|p_\mu|,
\qquad
|\mathrm{sn}_\mu-p_\mu|=\frac{2}{\eta}\left|\sin\mathrm{tq}-\mathrm{tq}\right|\le\frac{2}{\eta}\cdot\frac{|\mathrm{tq}|^3}{6}=\frac{\eta^2|p_\mu|^3}{24}.
\]
Therefore
\[
\left|\mathrm{sn}_\mu\cdot\mathrm{sn}_{\mathfrak{l}}-p_\mu p_{\mathfrak{l}}\right|
\le\left|\mathrm{sn}_\mu-p_\mu\right|\cdot\left|\mathrm{sn}_{\mathfrak{l}}\right|+|p_\mu|\cdot\left|\mathrm{sn}_{\mathfrak{l}}-p_{\mathfrak{l}}\right|
\le\frac{\eta^2}{24}|p_\mu|\cdot|p_{\mathfrak{l}}|\cdot\left(p_\mu^2+p_{\mathfrak{l}}^2\right).
\]
For $\mu=\mathfrak{l}$ the right side equals $\frac{\eta^2}{12}p_\mu^4\le\frac{\eta^2}{12}|p|_2^4$.
For $\mu\ne\mathfrak{l}$ one has $|p_\mu|\cdot|p_{\mathfrak{l}}|\le\frac{1}{2}\left(p_\mu^2+p_{\mathfrak{l}}^2\right)$ and $p_\mu^2+p_{\mathfrak{l}}^2\le|p|_2^2$;
so the right side is at most $\frac{\eta^2}{48}|p|_2^4$.
\end{proof}

\begin{named}{Lemma 2}[Periodisation]
Let $R>0$ and $\eta>0$, and let $\mathrm{fs}$ be a bounded Borel function on $\mathbb{R}^3$ with compact support.
On locally bounded Borel functions the operators $\tau_b$, $\delta_{\eta,\mu}$ and $J_\eta$ act by the formulas of Subsection~0.3, of Proposition~47.1 and of Section~I.7.6, respectively.
\begin{itemize}
\item[(a)] For $b\in\mathbb{R}^3$,
\[
(\tau_b\mathrm{fs})^{R}=\tau_b\left(\mathrm{fs}^{R}\right),
\qquad
(\delta_{\eta,\mu}\mathrm{fs})^{R}=\delta_{\eta,\mu}\left(\mathrm{fs}^{R}\right),
\qquad
(J_\eta\mathrm{fs})^{R}=J_\eta\left(\mathrm{fs}^{R}\right).
\]
If $\mathrm{fs}$ is smooth, then $\mathrm{fs}^{R}$ is smooth and $\partial_{\mathfrak{l}}\left(\mathrm{fs}^{R}\right)=\left(\partial_{\mathfrak{l}}\mathrm{fs}\right)^{R}$.
\item[(b)] Let $\mathrm{fs}$ vanish outside $y_0+\left(-\frac{R}{2},\frac{R}{2}\right)^3$ for some $y_0\in\mathbb{R}^3$.
Then for every $y\in\mathbb{R}^3$ at most one of the points $y+R\cdot z$, $z\in\mathbb{Z}^3$, lies in $y_0+\left(-\frac{R}{2},\frac{R}{2}\right)^3$.
Consequently the following hold.
\begin{itemize}
\item[(b-i)] $\left|\mathrm{fs}^{R}(y)\right|^2=\sum_z\left|\mathrm{fs}(y+R\cdot z)\right|^2$ for all $y\in\mathbb{R}^3$, and $\sup\left|\mathrm{fs}^{R}\right|=\sup|\mathrm{fs}|$.
\item[(b-ii)] For every $y_1\in\mathbb{R}^3$,
\[
\int_{y_1+[0,R)^3}\mathrm{d}y\,\left|\mathrm{fs}^{R}(y)\right|^2=\int_{\mathbb{R}^3}\mathrm{d}y\,\left|\mathrm{fs}(y)\right|^2.
\]
\item[(b-iii)] If $R/\eta$ is an integer, then for every $y_1\in\mathbb{R}^3$
\[
\sum_{c\in\Lambda_{\eta,R}}\left|\mathrm{fs}^{R}(c+y_1)\right|^2=\sum_{c\in\eta\mathbb{Z}^3}\left|\mathrm{fs}(c+y_1)\right|^2.
\]
\item[(b-iv)] For every Borel subset $\mathrm{Bq}$ of $y_0+\left(-\frac{R}{2},\frac{R}{2}\right)^3$ and every $y_1\in\mathbb{R}^3$ the Lebesgue measure of $\left(\mathrm{Bq}+R\mathbb{Z}^3\right)\cap\left(y_1+[0,R)^3\right)$ equals that of $\mathrm{Bq}$.
\end{itemize}
\end{itemize}
\end{named}

\begin{proof}
(a) One has $(\tau_b\mathrm{fs})^{R}(y)=\sum_z\mathrm{fs}(y-b+R\cdot z)=\mathrm{fs}^{R}(y-b)$, and $\delta_{\eta,\mu}$ is a linear combination of two translations.
For $y$ in a bounded set the sum over $z$ has only finitely many non-zero terms;
so it may be interchanged with the integral over $y'\in\left[-\frac{\eta}{2},\frac{\eta}{2}\right]^3$ and with $\partial_{\mathfrak{l}}$.

(b) If $y+R\cdot z$ and $y+R\cdot z'$ both lie in $y_0+\left(-\frac{R}{2},\frac{R}{2}\right)^3$, then $R\cdot|z-z'|_\infty<R$, and so $z=z'$.
(b-i) The sum defining $\mathrm{fs}^{R}(y)$ has at most one non-zero term;
this gives the first formula and $\sup\left|\mathrm{fs}^{R}\right|\le\sup|\mathrm{fs}|$.
For $y\in y_0+\left(-\frac{R}{2},\frac{R}{2}\right)^3$ the terms with $z\ne0$ vanish, so that $\mathrm{fs}^{R}(y)=\mathrm{fs}(y)$;
this gives $\sup\left|\mathrm{fs}^{R}\right|\ge\sup|\mathrm{fs}|$.
(b-ii) By (b-i) and Tonelli's theorem the left side equals
\[
\sum_{z\in\mathbb{Z}^3}\int_{y_1+R\cdot z+[0,R)^3}\mathrm{d}y\,\left|\mathrm{fs}(y)\right|^2,
\]
and the cubes $y_1+R\cdot z+[0,R)^3$, $z\in\mathbb{Z}^3$, partition $\mathbb{R}^3$.
(b-iii) By (b-i) the left side equals $\sum_{c\in\Lambda_{\eta,R}}\sum_{z\in\mathbb{Z}^3}\left|\mathrm{fs}(c+R\cdot z+y_1)\right|^2$, and $\eta\mathbb{Z}^3$ is the disjoint union of the classes $c+R\mathbb{Z}^3$.
(b-iv) This is (b-ii) for the indicator function of $\mathrm{Bq}$, whose periodisation is the indicator function of $\mathrm{Bq}+R\mathbb{Z}^3$, because at most one term equals $1$.
\end{proof}

\begin{named}{Lemma 3}[Tiling; norm of the test map]
Let $\eta>0$, and let $R=n\eta$ with an integer $n\ge1$.
\begin{itemize}
\item[(a)] For fixed $\mu<\nu$ the cubes $Q(c;\mu\nu)$, $c\in\eta\mathbb{Z}^3$, form a partition of $\mathbb{R}^3$, and $Q(c+R\cdot z;\mu\nu)=Q(c;\mu\nu)+R\cdot z$ for $z\in\mathbb{Z}^3$.
The union of the $n^3$ cubes $Q(c;\mu\nu)$, $c\in\eta\cdot\{0,\dots,n-1\}^3$ (a set of representatives of $\eta\mathbb{Z}^3$ modulo $R\mathbb{Z}^3$), is the cube $y_1+[0,R)^3$, where
\[
y_1:=\frac{\eta}{2}\left(e_\mu+e_\nu\right)-\frac{\eta}{2}(1,1,1).
\]
\item[(b)] For every bounded Borel two-form $f'$, real or complex, with compact support,
\[
|f'_\eta|_\eta\le\left\|f'^{R}\right\|_{L^2(\mathbb{R}^3/R\mathbb{Z}^3)}.
\]
\item[(c)] If moreover $f'$ vanishes outside $y_0+\left(-\frac{R}{2},\frac{R}{2}\right)^3$ for some $y_0\in\mathbb{R}^3$, then $|f'_\eta|_\eta\le\|f'\|_2$.
\end{itemize}
\end{named}

\begin{proof}
(a) Fix a coordinate $j$.
The half-open intervals $\left[\mathrm{tq}-\frac{\eta}{2},\mathrm{tq}+\frac{\eta}{2}\right)$, where $\mathrm{tq}$ ranges over $\eta\mathbb{Z}+\frac{\eta}{2}$ if $j\in\{\mu,\nu\}$ and over $\eta\mathbb{Z}$ if $j$ is the third direction, form a partition of $\mathbb{R}$.
The cube $Q(c;\mu\nu)$ is the product of three such intervals, and every such product arises from exactly one $c\in\eta\mathbb{Z}^3$.
Moreover $R\mathbb{Z}^3=n\eta\mathbb{Z}^3$ is a subgroup of $\eta\mathbb{Z}^3$, and $\mathrm{ce}(c+R\cdot z;\mu\nu)=\mathrm{ce}(c;\mu\nu)+R\cdot z$.
In the coordinate $j$ the $n$ intervals with $\mathrm{tq}=\eta i+y_{1,j}+\frac{\eta}{2}$, $i=0,\dots,n-1$, have the union $\left[y_{1,j},y_{1,j}+n\eta\right)$.

(b) The Cauchy--Schwarz inequality on the cube $Q(\boldsymbol{\sigma})$, whose volume is $\eta^3$, gives
\[
\left|f'_\eta(\boldsymbol{\sigma})\right|^2=\eta^{-6}\left|\int_{Q(\boldsymbol{\sigma})}\mathrm{d}y\,\left(f'_{\mu\nu}\right)^{R}(y)\right|^2\le\eta^{-3}\int_{Q(\boldsymbol{\sigma})}\mathrm{d}y\,\left|\left(f'_{\mu\nu}\right)^{R}(y)\right|^2.
\]
Multiply by $\eta^3$ and sum over $c\in\eta\cdot\{0,\dots,n-1\}^3$;
by (a) the right side becomes the integral of $\left|\left(f'_{\mu\nu}\right)^{R}\right|^2$ over the cube $y_1+[0,R)^3$.
Summation over $\mu<\nu$ turns the left side into $|f'_\eta|_\eta^2$ (Subsection~0.1) and the right side into $\left\|f'^{R}\right\|_{L^2(\mathbb{R}^3/R\mathbb{Z}^3)}^2$ (Subsection~0.4);
here $y_1$ depends on the pair $(\mu,\nu)$, but the integral of a periodic function over a period cube does not depend on the corner of the cube.

(c) Lemma~2\,(b-ii), applied to the three components, gives $\left\|f'^{R}\right\|_{L^2(\mathbb{R}^3/R\mathbb{Z}^3)}=\|f'\|_2$.
\end{proof}

\begin{named}{Lemma 4}[Translations]
Let $\beta$ be a probe, $0<\eta\le1$ and $b\in\mathbb{R}^3$.
\begin{itemize}
\item[(a)] $\tau_b\beta$ is a probe, and
\[
a'[\tau_b\beta]=\tau_ba'[\beta],
\qquad
\tilde{a}_\eta[\tau_b\beta]=\tau_b\tilde{a}_\eta[\beta],
\qquad
f^{(\eta)}[\tau_b\beta]=\tau_bf^{(\eta)}[\beta],
\qquad
dd^*\tau_b\beta=\tau_b\,dd^*\beta.
\]
The norms $\|\cdot\|_{H^s}$, $\|\cdot\|_2$, $|\cdot|_{C^i}$ and $|D^i\cdot|_\infty$ of $\tau_b\beta$, $\tau_bf$, $\tau_bf^{(\eta)}$ coincide with those of $\beta$, $f$, $f^{(\eta)}$.
\item[(b)] $D_{\tau_b\beta}=D_\beta$, and $b+Q_\beta$ has the properties (cu1), (cu2) with respect to $\tau_b\beta$.
\item[(c)] Let $R=n\eta$ with an integer $n\ge1$, let $v\in\eta\mathbb{Z}^3$, and let $f'$ be a bounded Borel two-form with compact support.
Then
\[
\left(\tau_vf'\right)_\eta=\left(\frac{v}{\eta},1\right)_*f'_\eta.
\]
\end{itemize}
\end{named}

\begin{proof}
(a) The support $\operatorname{supp}\tau_b\beta=b+\operatorname{supp}\beta$ is compact, and $\tau_b\beta$ is real, smooth and $\ne0$.
The four identities hold because $\tau_b$ commutes with $\delta_{\eta,\mu}$, $J_\eta$ and $\partial_\mu$ (Lemma~1\,(a)).
Moreover $\left|\widehat{\tau_b\mathrm{fs}}(p)\right|=\left|\widehat{\mathrm{fs}}(p)\right|$ (proof of Lemma~1\,(e)), sup norms are invariant under translations, and $\partial^{\mathrm{mi}}\tau_b=\tau_b\partial^{\mathrm{mi}}$.

(b) A cube contains $b+\operatorname{supp}\beta+[-2,2]^3$ if and only if its translate by $-b$ contains $\operatorname{supp}\beta+[-2,2]^3$.

(c) Since $\mathrm{ce}(\boldsymbol{\sigma}-v)=\mathrm{ce}(\boldsymbol{\sigma})-v$, one has $Q(\boldsymbol{\sigma})-v=Q(\boldsymbol{\sigma}-v)$.
With Lemma~2\,(a) and (0.4) it follows that
\begin{align*}
\left(\tau_vf'\right)_\eta(\boldsymbol{\sigma})
&=\eta^{-3}\int_{Q(\boldsymbol{\sigma})}\mathrm{d}y\,\left(f'_{\mu\nu}\right)^{R}(y-v)
=\eta^{-3}\int_{Q(\boldsymbol{\sigma}-v)}\mathrm{d}y\,\left(f'_{\mu\nu}\right)^{R}(y)\\
&=f'_\eta(\boldsymbol{\sigma}-v)
=\left(\left(\frac{v}{\eta},1\right)_*f'_\eta\right)(\boldsymbol{\sigma}),
\end{align*}
the last equality by Subsection~0.3.
\end{proof}

\begin{named}{Lemma 5}[The condition $R>D_\beta$]
Let $\beta$ be a probe and $R>0$.
\begin{itemize}
\item[(a)] $Q_\beta$ lies in the interior of a translate of $[0,R)^3$ if and only if $R>D_\beta$.
In this case $Q_\beta\subset y_\beta+\left(-\frac{R}{2},\frac{R}{2}\right)^3$.
\item[(b)] The condition $R>D_\beta$ holds for $\tau_b\beta$, $b\in\mathbb{R}^3$, if and only if it holds for $\beta$.
\item[(c)] If $R>D_\beta$, $0<\eta\le1$ and $n=R/\eta$ is an integer, then $n\ge5$.
\end{itemize}
\end{named}

\begin{proof}
(a) The interior of $y_1+[0,R)^3$ is $y_1+(0,R)^3$.
It contains $y_\beta+\left[-\frac{D_\beta}{2},\frac{D_\beta}{2}\right]^3$ if and only if
\[
y_{1,\mu}<y_{\beta,\mu}-\frac{D_\beta}{2}
\quad\text{and}\quad
y_{\beta,\mu}+\frac{D_\beta}{2}<y_{1,\mu}+R
\qquad\text{for all }\mu .
\]
Such a $y_1$ exists if and only if $D_\beta<R$, and in that case $y_1=y_\beta-\frac{R}{2}(1,1,1)$ is one of them.
(b) This is Lemma~4\,(b).
(c) One has $n=R/\eta\ge R>D_\beta>4$ (Subsection~0.5), and $n$ is an integer.
\end{proof}

\subsection*{1. Proposition~47.1}

\subsubsection*{1.2. Proof of Proposition~47.1}

Throughout, the definitions in the statement, in Sections~I.7.1, I.7.5, I.7.6 and in Section~G5.1.3 (the display that defines $\mathrm{gs}_\nu[\beta']$) are read with the conventions of Subsections~0.1--0.5.

\paragraph{The data.}
\begin{itemize}
\item[(1)] By Lemma~1\,(a),
\[
\tilde{a}_{\eta,\nu}=J_\eta a'_\nu=-\sum_{\mu\ne\nu}J_\eta\delta_{\eta,\mu}\beta_{\mu\nu}=-\sum_{\mu\ne\nu}\delta_{\eta,\mu}J_\eta\beta_{\mu\nu},
\]
and the right side is $\mathrm{gs}_\nu[\beta]$ of Section~G5.1.3.
\item[(2)] By Plancherel's theorem $\|u\|_{H^0}^2=\sum_{\mu<\nu}\int\mathrm{d}y\,\left|u_{\mu\nu}(y)\right|^2$ for $u\in\mathcal{S}$;
so $\|\cdot\|_2$ is the norm of $L^2(\mathbb{R}^3;\Lambda^2)$.
The forms $\beta$, $f$ and $f^{(\eta)}$ lie in $C_c^\infty(\mathbb{R}^3;\Lambda^2)$ (Lemma~1\,(a)), which is contained in $\mathcal{S}$;
hence all norms that occur in the statement are finite.
\item[(3)] Grade~0 of the statement contains three assertions on the condition $R>D_\beta$:
it means that $Q_\beta$ lies in the interior of a translate of $[0,R)^3$;
it is invariant under $\beta\mapsto\tau_b\beta$, $b\in\mathbb{R}^3$;
and, because it holds, at most one term of the periodising sum differs from $0$.
The first two are Lemma~5\,(a) and~(b);
the third is proved in the paragraph following (1.0).
By Lemma~5\,(c), $n\ge5$;
hence the standing condition $n_\mu\ge4$ of Section~I.7.1 is satisfied wherever $C^2(n)$ occurs in Proposition~47.1.
\end{itemize}

\paragraph{Grade 0.}

\textbf{(b2), first sentence.}
The probe $\beta$ is real and smooth and has compact support;
by Lemma~1\,(a) the same is true of $a'$, of $\tilde a_\eta=J_\eta a'$ and of $f^{(\eta)}$.
Lemma~1\,(b) yields the inclusions
\begin{align*}
\operatorname{supp}a'_\nu
&\subset\bigcup_{\mu\ne\nu}\left(\operatorname{supp}\beta+[-\eta/2,\eta/2]\cdot e_\mu\right)
\subset\operatorname{supp}\beta+[-\eta/2,\eta/2]^3,\\
\operatorname{supp}\tilde a_{\eta,\nu}=\operatorname{supp}J_\eta a'_\nu
&\subset\operatorname{supp}a'_\nu+[-\eta/2,\eta/2]^3,\\
\operatorname{supp}f^{(\eta)}_{\mu\nu}
&\subset\left(\operatorname{supp}a'_\nu+[-\eta/2,\eta/2]\cdot e_\mu\right)\cup\left(\operatorname{supp}a'_\mu+[-\eta/2,\eta/2]\cdot e_\nu\right).
\end{align*}
Since $[-\eta/2,\eta/2]^3+[-\eta/2,\eta/2]^3=[-\eta,\eta]^3$, it follows that
\begin{equation}\tag{1.0}
\operatorname{supp}\tilde a_\eta,\ \operatorname{supp}f^{(\eta)}\subset\operatorname{supp}\beta+[-\eta,\eta]^3,
\qquad
\operatorname{supp}J_\eta f^{(\eta)}_{\mu\nu}\subset\operatorname{supp}\beta+[-3\eta/2,3\eta/2]^3.
\end{equation}
As $\eta\le1$, one has $3\eta/2\le2$.
Hence the three supports in (1.0) lie in $\operatorname{supp}\beta+[-2,2]^3$, therefore in $Q_\beta$ by (cu2), and therefore in $y_\beta+(-R/2,R/2)^3$ by Lemma~5\,(a).
The support of $J_\eta\beta_{\mu\nu}$, which by Lemma~1\,(b) is contained in $\operatorname{supp}\beta+[-\eta/2,\eta/2]^3$, lies in the same set.
Consequently Lemma~2\,(b), with $y_0=y_\beta$, applies to the components of $J_\eta\beta$, $\tilde a_\eta$, $f^{(\eta)}$, $J_\eta f^{(\eta)}$ and to all their derivatives:
for each such function $\mathrm{fs}$ and each $y\in\mathbb R^3$, at most one term of the sum $\mathrm{fs}^R(y)$ is different from $0$.
In particular, every such $\mathrm{fs}$ satisfies
\begin{equation}\tag{1.1}
\mathrm{fs}^R(y)=\mathrm{fs}(y)\qquad\text{for all }y\in y_\beta+[-R/2,R/2)^3,
\end{equation}
since for $z\ne0$ the point $y+Rz$ lies outside the fundamental domain $y_\beta+[-R/2,R/2)^3$ of $R\mathbb Z^3$, and therefore outside $Q_\beta$.

\textbf{(b1), the adjoint.}
For every two-cochain $u$ on $\Lambda_{\eta,R}$,
\begin{equation}\tag{1.2}
\left(d_\eta^*u\right)_\nu(c)=-\eta^{-1}\sum_{\mu\ne\nu}\left[u_{\mu\nu}(c)-u_{\mu\nu}(c-\eta e_\mu)\right].
\end{equation}
To prove this, let $A$ be a one-cochain.
By (0.1) and by the definition of $\langle\cdot,\cdot\rangle_\eta$ on two-cochains in Subsection~0.1,
\[
\langle d_\eta A,u\rangle_\eta=\eta^3\sum_{c\in\Lambda_{\eta,R}}\sum_{\mu<\nu}\eta^{-1}\left[A_\mu(c)+A_\nu(c+\eta e_\mu)-A_\mu(c+\eta e_\nu)-A_\nu(c)\right]\cdot u_{\mu\nu}(c).
\]
After the substitution $c\to c-\eta e_\mu$ in the second term and $c\to c-\eta e_\nu$ in the third, this becomes
\[
\langle d_\eta A,u\rangle_\eta=\eta^3\sum_{c}\sum_{\mu<\nu}\eta^{-1}\left\{A_\mu(c)\left[u_{\mu\nu}(c)-u_{\mu\nu}(c-\eta e_\nu)\right]-A_\nu(c)\left[u_{\mu\nu}(c)-u_{\mu\nu}(c-\eta e_\mu)\right]\right\}.
\]
In the first group of terms write $u_{\mu\nu}=-u_{\nu\mu}$ and interchange the names $\mu$ and $\nu$;
the group turns into
\[
-\sum_{\nu<\mu}A_\nu(c)\left[u_{\mu\nu}(c)-u_{\mu\nu}(c-\eta e_\mu)\right].
\]
The second group is the same expression summed over the pairs $\mu<\nu$, so that the two groups together form a sum over all ordered pairs $\mu\ne\nu$:
\[
\langle d_\eta A,u\rangle_\eta=\eta^3\sum_{c}\sum_{\nu}A_\nu(c)\left(-\eta^{-1}\sum_{\mu\ne\nu}\left[u_{\mu\nu}(c)-u_{\mu\nu}(c-\eta e_\mu)\right]\right).
\]
This is (1.2).

\textbf{(b1), first identity.}
Let $u$ be the cochain of (0.4).
The first formula of (0.4) holds for all ordered pairs $\mu\ne\nu$, its two sides being antisymmetric in $(\mu,\nu)$.
Let $y:=c+(\eta/2)e_\nu$ be the centre of the bond $(c,c+\eta e_\nu)$.
In $u_{\mu\nu}(c)$ and in $u_{\mu\nu}(c-\eta e_\mu)$ the function $\left(J_\eta\beta_{\mu\nu}\right)^R$ is evaluated at $y+(\eta/2)e_\mu$ and at $y-(\eta/2)e_\mu$, respectively.
The definition of $\delta_{\eta,\mu}$ and Lemma~2\,(a) therefore give
\[
u_{\mu\nu}(c)-u_{\mu\nu}(c-\eta e_\mu)=\eta\cdot\left(\delta_{\eta,\mu}\left(J_\eta\beta_{\mu\nu}\right)^R\right)(y)=\eta\cdot\left(\delta_{\eta,\mu}J_\eta\beta_{\mu\nu}\right)^R(y).
\]
Inserting this into (1.2) and using (1) of \emph{The data}, one obtains, for all $c\in\eta\mathbb Z^3$ and $\nu=1,2,3$,
\begin{equation}\tag{1.3}
a_{\eta,\nu}(c)=\left(\tilde a_{\eta,\nu}\right)^R\left(c+\tfrac{\eta}{2}e_\nu\right)=\sum_{z\in\mathbb Z^3}\tilde a_{\eta,\nu}\left(c+\tfrac{\eta}{2}e_\nu+Rz\right).
\end{equation}
This is the first formula of (b1): its right side is the value of $\tilde a_{\eta,\nu}$ at a point of $\mathbb R^3/R\mathbb Z^3$, in the sense of Subsection~0.4.
By Lemma~2\,(b) at most one term of the sum is non-zero, and by (1.1) one has $a_{\eta,\nu}(c)=\tilde a_{\eta,\nu}\left(c+(\eta/2)e_\nu\right)$ for that representative $c$ of its class modulo $R\mathbb Z^3$ for which $c+(\eta/2)e_\nu\in y_\beta+[-R/2,R/2)^3$.

\textbf{(b1), second identity.}
Let $\boldsymbol{\sigma}=(c;\mu\nu)$ with $\mu<\nu$, and put $y:=\mathrm{ce}(\boldsymbol{\sigma})=c+(\eta/2)(e_\mu+e_\nu)$.
From (1.3),
\begin{align*}
a_{\eta,\nu}(c+\eta e_\mu)-a_{\eta,\nu}(c)
&=\left(\tilde a_{\eta,\nu}\right)^R\left(y+\tfrac{\eta}{2}e_\mu\right)-\left(\tilde a_{\eta,\nu}\right)^R\left(y-\tfrac{\eta}{2}e_\mu\right)
=\eta\cdot\left(\delta_{\eta,\mu}\tilde a_{\eta,\nu}\right)^R(y),\\
a_{\eta,\mu}(c)-a_{\eta,\mu}(c+\eta e_\nu)
&=\left(\tilde a_{\eta,\mu}\right)^R\left(y-\tfrac{\eta}{2}e_\nu\right)-\left(\tilde a_{\eta,\mu}\right)^R\left(y+\tfrac{\eta}{2}e_\nu\right)
=-\eta\cdot\left(\delta_{\eta,\nu}\tilde a_{\eta,\mu}\right)^R(y).
\end{align*}
Together with (0.1), with $\tilde a_\eta=J_\eta a'$ and with Lemma~1\,(a), this yields
\begin{equation}\tag{1.4}
\left(d_\eta a_\eta\right)_{\mu\nu}(c)=\left(\delta_{\eta,\mu}J_\eta a'_\nu-\delta_{\eta,\nu}J_\eta a'_\mu\right)^R(y)=\left(J_\eta f^{(\eta)}_{\mu\nu}\right)^R\left(\mathrm{ce}(\boldsymbol{\sigma})\right).
\end{equation}
On the other hand, the sets $Q(\boldsymbol{\sigma})$ and $\mathrm{ce}(\boldsymbol{\sigma})+[-\eta/2,\eta/2]^3$ differ by a null set; so (0.4) and Lemma~2\,(a) give
\[
\left(f^{(\eta)}\right)_\eta(\boldsymbol{\sigma})
=\eta^{-3}\int_{[-\eta/2,\eta/2]^3}\mathrm dy'\,\left(f^{(\eta)}_{\mu\nu}\right)^R\left(\mathrm{ce}(\boldsymbol{\sigma})+y'\right)
=\left(J_\eta\left(f^{(\eta)}_{\mu\nu}\right)^R\right)\left(\mathrm{ce}(\boldsymbol{\sigma})\right)
=\left(J_\eta f^{(\eta)}_{\mu\nu}\right)^R\left(\mathrm{ce}(\boldsymbol{\sigma})\right).
\]
Hence $\left(f^{(\eta)}\right)_\eta=d_\eta a_\eta$.
The relations (1.2)--(1.4) hold without the hypothesis $R>D_\beta$.

\textbf{(b1), the form on $\eta\mathbb Z^3$} (clause (b2$^{\mathrm{o}}$), first and third assertions).
Of the data of grade~0, only $0<\eta\le1$ enters here: neither $R$ nor $n$ appears in (1) of \emph{The data}, in the derivation of (1.0), or in the following lines up to (1.3$'$).
On the cochains of $\eta\mathbb Z^3$ let $\langle\cdot,\cdot\rangle_\eta$ be the pairings of Subsection~0.1 with the sums extended over all $c\in\eta\mathbb Z^3$;
they are defined whenever one of the two factors has finite support.
For a two-cochain $u$ of finite support, $d_\eta^*u$ denotes the one-cochain with $\langle d_\eta A,u\rangle_\eta=\langle A,d_\eta^*u\rangle_\eta$ for all one-cochains $A$;
there is at most one such cochain, as one sees by taking for $A$ the indicators of the bonds.
The computation leading to (1.2) remains valid for any sum over $c$ that has finitely many non-zero terms and is invariant under the substitutions $c\to c-\eta e_\mu$ and $c\to c-\eta e_\nu$.
Therefore (1.2) holds on $\eta\mathbb Z^3$ for every $u$ of finite support.
The cochain $u^{\mathrm{o}}$ of (b2$^{\mathrm{o}}$) has finite support, because $\operatorname{supp}J_\eta\beta_{\mu\nu}$ is compact (Lemma~1\,(b)).
With $y=c+(\eta/2)e_\nu$, the computation that precedes (1.3) gives
\[
u^{\mathrm{o}}_{\mu\nu}(c)-u^{\mathrm{o}}_{\mu\nu}(c-\eta e_\mu)=\eta\cdot\left(\delta_{\eta,\mu}J_\eta\beta_{\mu\nu}\right)(y),
\]
and from (1.2) and (1) of \emph{The data} it follows that the cochain $a^{\mathrm{o}}_\eta=d_\eta^*u^{\mathrm{o}}$ has the components
\begin{equation}\tag{1.3$'$}
a^{\mathrm{o}}_{\eta,\nu}(c)=\tilde a_{\eta,\nu}\left(c+\tfrac{\eta}{2}e_\nu\right)\qquad\left(c\in\eta\mathbb Z^3,\ \nu=1,2,3\right).
\end{equation}
The right side contains neither $R$ nor $n$.
Now let $R>0$ be such that $R/\eta$ is an integer $\ge2$.
Since (1.3) holds without the hypothesis $R>D_\beta$, (1.3) and (1.3$'$) give
\[
a_{\eta,\nu}(c)=\sum_{z\in\mathbb Z^3}a^{\mathrm{o}}_{\eta,\nu}(c+Rz);
\]
in other words, $a_\eta$ is the $R\mathbb Z^3$-periodisation of the cochain $a^{\mathrm{o}}_\eta$.

\textbf{(b1), third assertion.}
By its definition, $a_\eta=d_\eta^*u$ lies in the range of $d_\eta^*$.
If $A$ is a one-cochain with $dA=0$, then $\langle a_\eta,A\rangle_\eta=\langle u,d_\eta A\rangle_\eta=0$.
For $A=d\boldsymbol\lambda$, (0.1) gives
\begin{multline*}
\left(dd\boldsymbol\lambda\right)_{\mu\nu}(c)
=\left[\boldsymbol\lambda(c+\eta e_\mu)-\boldsymbol\lambda(c)\right]+\left[\boldsymbol\lambda(c+\eta e_\mu+\eta e_\nu)-\boldsymbol\lambda(c+\eta e_\mu)\right]\\
-\left[\boldsymbol\lambda(c+\eta e_\nu+\eta e_\mu)-\boldsymbol\lambda(c+\eta e_\nu)\right]-\left[\boldsymbol\lambda(c+\eta e_\nu)-\boldsymbol\lambda(c)\right]=0.
\end{multline*}
For a constant one-cochain, that is, one for which $A_\nu(c)$ does not depend on $c$, the four terms of $(dA)_{\mu\nu}(c)$ cancel in pairs.
So $a_\eta$ is orthogonal to the exact and to the constant one-cochains, which means that it lies in $V_{\mathrm t}$ (Section~I.7.1).

\textbf{(b2), second sentence.}
Let the bond $(c,c+\eta e_\nu)$ belong to $\operatorname{supp}a_\eta$, so that $a_{\eta,\nu}(c)\ne0$.
By (1.3) there is $z\in\mathbb Z^3$ such that $y:=c+(\eta/2)e_\nu+Rz$ lies in $\operatorname{supp}\tilde a_{\eta,\nu}$.
By (1.0), the points $c+Rz=y-(\eta/2)e_\nu$ and $c+\eta e_\nu+Rz=y+(\eta/2)e_\nu$ belong to
\[
\operatorname{supp}\beta+[-\eta,\eta]^3+[-\eta/2,\eta/2]\cdot e_\nu\subset\operatorname{supp}\beta+[-3\eta/2,3\eta/2]^3,
\]
and one and the same vector $Rz$ serves for both end points.
As was shown after (1.0), the set on the right lies in $Q_\beta$.

\emph{The cochain $a^{\mathrm{o}}_\eta$} (clause (b2$^{\mathrm{o}}$), second assertion).
Let $c\in\eta\mathbb Z^3$ with $a^{\mathrm{o}}_{\eta,\nu}(c)\ne0$.
By (1.3$'$) the point $y:=c+(\eta/2)e_\nu$ lies in $\operatorname{supp}\tilde a_{\eta,\nu}$.
By (1.0), therefore, $c=y-(\eta/2)e_\nu$ and $c+\eta e_\nu=y+(\eta/2)e_\nu$ belong to $\operatorname{supp}\beta+[-\eta,\eta]^3+[-\eta/2,\eta/2]\cdot e_\nu$, a subset of $\operatorname{supp}\beta+[-3\eta/2,3\eta/2]^3$.
Thus both end points of every bond of $\operatorname{supp}a^{\mathrm{o}}_\eta$ lie in this set, as points of $\mathbb R^3$.

\textbf{(b3).}
By (0.2) and Lemma~1\,(e), where $(i\,p_\mu)(i\,p_{\mathfrak l})=-p_\mu p_{\mathfrak l}$ and $(i\,\mathrm{sn}_\mu)(i\,\mathrm{sn}_{\mathfrak l})=-\mathrm{sn}_\mu\cdot\mathrm{sn}_{\mathfrak l}$,
\[
\widehat{f^{(\eta)}_{\mu\nu}}(p)-\hat f_{\mu\nu}(p)
=\sum_{\mathfrak l\ne\nu}\left(\mathrm{sn}_\mu\cdot\mathrm{sn}_{\mathfrak l}-p_\mu p_{\mathfrak l}\right)\hat\beta_{\mathfrak l\nu}(p)
-\sum_{\mathfrak l\ne\mu}\left(\mathrm{sn}_\nu\cdot\mathrm{sn}_{\mathfrak l}-p_\nu p_{\mathfrak l}\right)\hat\beta_{\mathfrak l\mu}(p).
\]
There are four terms on the right, and by (0.5) each coefficient has modulus at most $(\eta^2/12)\cdot|p|_2^4$.
Let $\mathfrak l_3$ denote the third direction (Subsection~0.2, after (0.2)).
The Cauchy--Schwarz inequality for four terms gives
\[
\left|\widehat{f^{(\eta)}_{\mu\nu}}-\hat f_{\mu\nu}\right|^2
\le4\left(\frac{\eta^2}{12}\right)^2|p|_2^8\cdot\left(2\left|\hat\beta_{\mu\nu}\right|^2+\left|\hat\beta_{\mathfrak l_3\nu}\right|^2+\left|\hat\beta_{\mathfrak l_3\mu}\right|^2\right).
\]
Sum over the three pairs $\mu<\nu$.
For a given pair $\{\mu,\nu\}$, the square of the modulus of the component of $\hat\beta$ on this pair occurs twice in the bracket of the pair itself, once in the bracket of the pair $\{\mu,\mathfrak l_3\}$, whose third direction is $\nu$, and once in that of the pair $\{\nu,\mathfrak l_3\}$, whose third direction is $\mu$: four times in all.
Therefore
\[
\sum_{\mu<\nu}\left|\widehat{f^{(\eta)}_{\mu\nu}}(p)-\hat f_{\mu\nu}(p)\right|^2
\le16\left(\frac{\eta^2}{12}\right)^2|p|_2^8\sum_{\mu<\nu}\left|\hat\beta_{\mu\nu}(p)\right|^2
=\left(\frac{\eta^2}{3}\right)^2|p|_2^8\sum_{\mu<\nu}\left|\hat\beta_{\mu\nu}(p)\right|^2.
\]
Multiply by $\left(1+|p|_2^2\right)^s$, use $|p|_2^8\le\left(1+|p|_2^2\right)^4$, and integrate with respect to $\mathrm dp/(2\pi)^3$; the result is
\[
\left\|f^{(\eta)}-f\right\|_{H^s}^2\le\left(\frac{\eta^2}{3}\right)^2\|\beta\|_{H^{s+4}}^2,
\]
valid for every real number $s$.
Since $\eta^2/3\le\eta^2$, this contains (b3).

\textbf{(b4).}
Let $\mathrm{mi}$ be a multi-index of order $|\mathrm{mi}|=j$.
By Lemma~1\,(a),
\[
\partial^{\mathrm{mi}}\tilde a_{\eta,\nu}=-\sum_{\mu\ne\nu}J_\eta\delta_{\eta,\mu}\partial^{\mathrm{mi}}\beta_{\mu\nu},
\]
and by Lemma~1\,(c),
\[
\left|J_\eta\delta_{\eta,\mu}\partial^{\mathrm{mi}}\beta_{\mu\nu}\right|_\infty\le\left|\partial_\mu\partial^{\mathrm{mi}}\beta_{\mu\nu}\right|_\infty\le\left|D^{j+1}\beta\right|_\infty.
\]
There are two values $\mu\ne\nu$; taking $j=i$, one gets $\left|D^i\tilde a_\eta\right|_\infty\le2\left|D^{i+1}\beta\right|_\infty\le2|\beta|_{C^{i+1}}$.

Next,
\[
\partial^{\mathrm{mi}}\left(\tilde a_{\eta,\nu}-a_\nu\right)
=-\sum_{\mu\ne\nu}\left[J_\eta\left(\delta_{\eta,\mu}-\partial_\mu\right)\partial^{\mathrm{mi}}\beta_{\mu\nu}+\left(J_\eta-1\right)\partial_\mu\partial^{\mathrm{mi}}\beta_{\mu\nu}\right].
\]
By Lemma~1\,(c), (d), the first term in the bracket has sup norm at most
\[
\frac{\eta^2}{24}\left|\partial_\mu^3\partial^{\mathrm{mi}}\beta_{\mu\nu}\right|_\infty\le\frac{\eta^2}{24}\left|D^{j+3}\beta\right|_\infty,
\]
and the second at most
\[
\frac98\eta^2\left|D^2\partial_\mu\partial^{\mathrm{mi}}\beta_{\mu\nu}\right|_\infty\le\frac98\eta^2\left|D^{j+3}\beta\right|_\infty.
\]
Hence
\[
\left|D^j\left(\tilde a_\eta-a\right)\right|_\infty\le2\left(\frac1{24}+\frac98\right)\eta^2\left|D^{j+3}\beta\right|_\infty=\frac73\eta^2\left|D^{j+3}\beta\right|_\infty,
\]
and the maximum over $j\le i$ gives $\left|\tilde a_\eta-a\right|_{C^i}\le\frac73\eta^2|\beta|_{C^{i+3}}\le3\eta^2|\beta|_{C^{i+3}}$.

\textbf{(b0).}
A test $f'\in\mathcal S$ whose support is contained in the bounded set $y_0+(-R/2,R/2)^3$ is bounded and continuous and has compact support.
The assertion is therefore Lemma~3\,(c).

\paragraph{Grade 1.}
Let $v\in\mathbb R^3$.
By Lemma~4\,(a), (b) and Lemma~5\,(b), $\tau_v\beta$ is a probe, $D_{\tau_v\beta}=D_\beta<R$, and the cube $v+Q_\beta$ has the properties (cu1), (cu2) with respect to $\tau_v\beta$.
Hence everything proved in grade~0 holds for $\tau_v\beta$, with $v+Q_\beta$ and $v+y_\beta$ in place of $Q_\beta$ and $y_\beta$.
Put
\[
f':=\tau_vf^{(\eta)}=f^{(\eta)}[\tau_v\beta]\quad\text{(Lemma~4\,(a))},\qquad a_\eta:=a_\eta[\tau_v\beta],\qquad\hat{\mathfrak a}:=\eta^{\frac12}\cdot a_\eta.
\]
The second identity of (b1) for $\tau_v\beta$, together with $d=\eta\cdot d_\eta$, gives
\begin{equation}\tag{1.5}
d\hat{\mathfrak a}=\eta^{\frac12}\cdot\eta\cdot d_\eta a_\eta=\eta^{\frac32}\cdot f'_\eta.
\end{equation}

\textbf{(b2$'$).}
Let $F^{\mathrm u}\in C^2(n)$ and $F=\eta^{-\frac32}\cdot F^{\mathrm u}$ (Section~I.7.1).
From the definition of $X_{\eta,R}$ (Section~I.7.5), from $\langle u,u'\rangle_\eta=\eta^3\cdot\langle u,u'\rangle$ (Subsection~0.1) and from (1.5),
\[
X_{\eta,R}(f')=\langle f'_\eta,F\rangle_\eta=\eta^3\cdot\left\langle f'_\eta,\eta^{-\frac32}F^{\mathrm u}\right\rangle=\left\langle\eta^{\frac32}f'_\eta,F^{\mathrm u}\right\rangle=\left\langle d\hat{\mathfrak a},F^{\mathrm u}\right\rangle=X.
\]

\textbf{(b5), first display.}
By (1.5), $\boldsymbol{\sigma}^2=\|d\hat{\mathfrak a}\|^2=\eta^3\cdot\|f'_\eta\|^2=|f'_\eta|_\eta^2$.
By (1.0) for $\tau_v\beta$, the support of $f'$ lies in $v+Q_\beta$, which by Lemma~5\,(a) is a subset of $v+y_\beta+(-R/2,R/2)^3$.
So (b0) with $y_0=v+y_\beta$, together with Lemma~4\,(a), gives
\[
\boldsymbol{\sigma}=|f'_\eta|_\eta\le\|f'\|_2=\left\|f^{(\eta)}\right\|_2,
\]
and (b3) at $s=0$, together with $\eta\le1$, gives
\[
\left\|f^{(\eta)}\right\|_2\le\|f\|_2+\left\|f^{(\eta)}-f\right\|_2\le\|f\|_2+\eta^2\|\beta\|_{H^4}\le\|f\|_2+\|\beta\|_{H^4}.
\]

\textbf{(b5), second display.}
Let $\mathrm{Stp}(f')$ be the $R\mathbb Z^3$-periodic two-form given by
\[
\mathrm{Stp}(f')_{\mu\nu}(y):=f'_\eta(\boldsymbol{\sigma})\qquad\text{for }y\in Q(\boldsymbol{\sigma})+R\mathbb Z^3,\ \boldsymbol{\sigma}=(c;\mu\nu).
\]
By Lemma~3\,(a) it is defined at every point, and the integral of $\left|\mathrm{Stp}(f')_{\mu\nu}\right|^2$ over the cube $y_1+[0,R)^3$ of Lemma~3\,(a) equals
\[
\sum_{c\in\eta\{0,\dots,n-1\}^3}\eta^3\cdot\left|f'_\eta(c;\mu\nu)\right|^2.
\]
Consequently, with Lemma~2\,(b-ii) for the second relation and the triangle inequality for the third,
\begin{equation}\tag{1.6}
\begin{aligned}
&\left\|\mathrm{Stp}(f')\right\|_{L^2(\mathbb R^3/R\mathbb Z^3)}=|f'_\eta|_\eta=\boldsymbol{\sigma},\qquad
\left\|{f'}^R\right\|_{L^2(\mathbb R^3/R\mathbb Z^3)}=\|f'\|_2=\left\|f^{(\eta)}\right\|_2,\\
&\left|\boldsymbol{\sigma}-\left\|f^{(\eta)}\right\|_2\right|\le\left\|\mathrm{Stp}(f')-{f'}^R\right\|_{L^2(\mathbb R^3/R\mathbb Z^3)}.
\end{aligned}
\end{equation}

\emph{Pointwise bound.}
The two-form ${f'}^R$ is smooth, $\partial_j\left(f'_{\mu\nu}\right)^R=\left(\partial_jf'_{\mu\nu}\right)^R$ by Lemma~2\,(a), and by Lemma~2\,(b-i)
\[
\left|\left(\partial_jf'_{\mu\nu}\right)^R\right|_\infty=\left|\partial_jf'_{\mu\nu}\right|_\infty=\left|\partial_jf^{(\eta)}_{\mu\nu}\right|_\infty.
\]
By (0.2) and Lemma~1\,(a), (c), the derivative $\partial_jf^{(\eta)}_{\mu\nu}$ is a sum of four terms $\pm\delta_{\eta,\mu'}\delta_{\eta,\mathfrak l}\partial_j\beta_{\mathfrak l\nu'}$, with $(\mu',\nu')=(\mu,\nu)$ or $(\nu,\mu)$, each of sup norm at most $\left|\partial_{\mu'}\partial_{\mathfrak l}\partial_j\beta_{\mathfrak l\nu'}\right|_\infty\le\left|D^3\beta\right|_\infty$.
So $\left|\partial_j\left(f'_{\mu\nu}\right)^R\right|_\infty\le4|\beta|_{C^3}$.
Let $y$ and $y'$ lie in one cube $Q(\boldsymbol{\sigma})$, which is a convex subset of $\mathbb R^3$ of sup-diameter at most $\eta$.
Then
\begin{multline*}
\left|\left(f'_{\mu\nu}\right)^R(y')-\left(f'_{\mu\nu}\right)^R(y)\right|
=\left|\int_0^1\mathrm d\mathrm{tq}\,\sum_j\left(y'_j-y_j\right)\cdot\left(\partial_j\left(f'_{\mu\nu}\right)^R\right)\left(y+\mathrm{tq}\cdot(y'-y)\right)\right|\\
\le3\eta\cdot4|\beta|_{C^3}=12\eta\cdot|\beta|_{C^3}.
\end{multline*}
For $y\in Q(\boldsymbol{\sigma})$ one has
\[
\mathrm{Stp}(f')_{\mu\nu}(y)-\left(f'_{\mu\nu}\right)^R(y)=\eta^{-3}\int_{Q(\boldsymbol{\sigma})}\mathrm dy'\,\left[\left(f'_{\mu\nu}\right)^R(y')-\left(f'_{\mu\nu}\right)^R(y)\right],
\]
and therefore $\left|\mathrm{Stp}(f')_{\mu\nu}-\left(f'_{\mu\nu}\right)^R\right|\le12\eta\cdot|\beta|_{C^3}$ everywhere.

\emph{Support.}
Let $y\in Q(\boldsymbol{\sigma})$, $\boldsymbol{\sigma}=(c;\mu\nu)$, be a point with $\mathrm{Stp}(f')_{\mu\nu}(y)\ne\left(f'_{\mu\nu}\right)^R(y)$.
Then $\left(f'_{\mu\nu}\right)^R$ does not vanish identically on $Q(\boldsymbol{\sigma})$: there is $y'\in Q(\boldsymbol{\sigma})$ with $\left(f'_{\mu\nu}\right)^R(y')\ne0$.
This point $y'$ lies in $\operatorname{supp}f'+R\mathbb Z^3$, and $|y-y'|_\infty\le\eta$.
By (1.0) for $\tau_v\beta$, by $2\eta\le2$ and by (cu2), the point $y$ belongs to
\[
v+\operatorname{supp}\beta+[-2\eta,2\eta]^3+R\mathbb Z^3\subset v+Q_\beta+R\mathbb Z^3.
\]
By Lemma~2\,(b-iv), applied with $y_0=v+y_\beta$ and $\mathrm{Bq}=v+Q_\beta$, the part of this set that lies in a cube $y_1+[0,R)^3$ has measure $D_\beta^3$.

\emph{Conclusion.}
It follows that
\[
\left\|\mathrm{Stp}(f')-{f'}^R\right\|_{L^2(\mathbb R^3/R\mathbb Z^3)}^2
\le3\cdot\left(12\eta\cdot|\beta|_{C^3}\right)^2\cdot D_\beta^3
=432\eta^2\cdot|\beta|_{C^3}^2\cdot D_\beta^3
\le\left(21\eta\cdot|\beta|_{C^3}\cdot D_\beta^{\frac32}\right)^2,
\]
because $432\le441$.
Combined with (1.6) and with (b3) at $s=0$, this gives
\[
\left|\boldsymbol{\sigma}-\|f\|_2\right|
\le\left|\boldsymbol{\sigma}-\left\|f^{(\eta)}\right\|_2\right|+\left|\left\|f^{(\eta)}\right\|_2-\|f\|_2\right|
\le21\eta\cdot|\beta|_{C^3}\cdot D_\beta^{\frac32}+\eta^2\|\beta\|_{H^4}.
\]

\textbf{(b5$'$).}
\emph{A formula for $\boldsymbol{\sigma}^2$, valid for every $v\in\mathbb R^3$.}
By (1.5), by (1.4) for $\tau_v\beta$, and since $J_\eta\tau_v=\tau_vJ_\eta$,
\[
\boldsymbol{\sigma}^2=|f'_\eta|_\eta^2=\eta^3\sum_{\mu<\nu}\sum_{c\in\Lambda_{\eta,R}}\left|\left(\tau_vJ_\eta f^{(\eta)}_{\mu\nu}\right)^R\left(c+\tfrac{\eta}{2}(e_\mu+e_\nu)\right)\right|^2.
\]
By (1.0) for $\tau_v\beta$, the function $\tau_vJ_\eta f^{(\eta)}_{\mu\nu}$ vanishes outside $v+Q_\beta$, a subset of $v+y_\beta+(-R/2,R/2)^3$.
Lemma~2\,(b-iii) then gives
\begin{equation}\tag{1.7}
\boldsymbol{\sigma}^2=\eta^3\sum_{\mu<\nu}\sum_{c\in\eta\mathbb Z^3}\left|\left(J_\eta f^{(\eta)}_{\mu\nu}\right)\left(c+\tfrac{\eta}{2}(e_\mu+e_\nu)-v\right)\right|^2,
\end{equation}
a sum with finitely many non-zero terms.
Neither $R$ nor $n$ occurs on the right side; hence $\boldsymbol{\sigma}^2$ does not depend on $R$.

\emph{The case $v\in\eta\mathbb Z^3$.}
The substitution $c\to c+v$ in (1.7) removes $v$, so that $\boldsymbol{\sigma}^2$ is the same for all $v\in\eta\mathbb Z^3$.
By (1.3) for $\tau_v\beta$, by Lemma~4\,(a) and Lemma~2\,(a), and by (1.3) for $\beta$ at the site $c-v$ of $\eta\mathbb Z^3$,
\[
a_\eta[\tau_v\beta]_\nu(c)
=\left(\tau_v\tilde a_{\eta,\nu}[\beta]\right)^R\left(c+\tfrac{\eta}{2}e_\nu\right)
=\left(\tilde a_{\eta,\nu}[\beta]\right)^R\left(c-v+\tfrac{\eta}{2}e_\nu\right)
=a_\eta[\beta]_\nu(c-v).
\]
After multiplication by $\eta^{\frac12}$ this reads, in the notation of Subsection~0.3,
\begin{equation}\tag{1.8}
\hat{\mathfrak a}[\tau_v\beta]=(v/\eta,1)_*\,\hat{\mathfrak a}[\beta]\qquad\left(v\in\eta\mathbb Z^3\right):
\end{equation}
the cochain formed with $\tau_v\beta$ is the translate by $v/\eta$, measured in lattice spacings, of the cochain formed with $\beta$.

\paragraph{Grade 2.}
The assertions of grade~2 are those of Proposition~41.

\paragraph{Uniformity.}
The numbers $1$, $2$, $3$ and $21$ in (b3), (b4), (b5) are absolute.
Besides $\eta$, the right sides contain only the four quantities $\|\beta\|_{H^{s+4}}$, $|\beta|_{C^i}$, $D_\beta$, $\|f\|_2$, and these are invariant under $\beta\to\tau_b\beta$ (Lemma~4).
No bound depends on $R$, on $n$ or its parity, or on $v$; no charge, no mass and no constant of Section~I.8 enters.
\qed

\subsection*{2. Proposition 47.2}

\subsubsection*{2.2. Proof of Proposition 47.2}

\begin{proof}
\textbf{The numbers.}
The integers $N$ and $T$ are $\ge0$, and $\eta=L^{-N}$ (Sections~I.0.5 and I.6.5);
the last identity can also be read off the data of the statement, since
\[
\eta=\frac{R}{n}=\frac{8h\cdot L^{T}}{8h\cdot L^{T+N}}=L^{-N}.
\]
As $\bar m\in(0,1]$, the number $i'':=\left\lfloor\log_L(1/\bar m)\right\rfloor$ is an integer $\ge0$, and $h=L^{i''}$ (Section~I.0.5).
Moreover $r_K(e)=L^i$ for an integer $i\ge0$ (Section~I.6.1), and $M=L^{i'}$ for an integer $i'\ge2$ by condition (M1) of Proposition~14.2\,(b).
Hence $T^*(e)=\log_L\left(M\cdot r_K(e)\right)-1=i'+i-1$ is an integer $\ge1$ (Section~I.6.2), and Sections~I.6.2 and I.6.4 give
\begin{equation}\tag{2.1}
\begin{aligned}
R^*(e)&=h\cdot L^{T^*(e)},&\qquad n_*&=\frac{R^*(e)}{\eta}=h\cdot L^{T^*(e)+N},\\
\mathrm{per}&=\frac{n_*}{L}=h\cdot L^{T^*(e)+N-1},&\qquad\frac{R^*(e)}{L}&=h\cdot L^{T^*(e)-1}=\mathrm{per}\cdot\eta.
\end{aligned}
\end{equation}
Thus $\mathrm{per}=L^{i''+T^*(e)+N-1}$, where the exponent is an integer $\ge0+1+0-1=0$.
The number $L$ is odd by condition (L5) of Proposition~14.2\,(b).
So $\mathrm{per}$ is an odd integer $\ge1$, $m_{\mathrm p}=(\mathrm{per}-1)/2$ is an integer $\ge0$, and $\mathrm{hc}=\mathrm{per}/2=m_{\mathrm p}+1/2$.

\textbf{(a).}
The numbers $T$ and $T^*(e)$ are integers with $T\ge T^*(e)$; hence $T-T^*(e)$ is an integer $\ge0$, and $\mathrm{kt}:=8L^{T-T^*(e)}$ is an integer $\ge8$.
By (2.1),
\[
\mathrm{kt}\cdot n_*=8L^{T-T^*(e)}\cdot h\cdot L^{T^*(e)+N}=8h\cdot L^{T+N}=n,\qquad
\frac{n}{\mathrm{per}}=\frac{\mathrm{kt}\cdot n_*}{\mathrm{per}}=\mathrm{kt}\cdot L,
\]
which is an integer.
So $\mathrm{per}$ divides $n$, and $C(n,n,n)$ carries $\mathcal Q_*$, because by Section~I.6.4 a torus $C(n_1,n_2,n_3)$ carries $\mathcal Q_*$ if and only if $\mathrm{per}$ divides every $n_\mu$.
Moreover $n=8L^{i''+T+N}$ with a non-negative exponent; so $n$ is an integer $\ge8\ge4$.

\emph{Cells.}
For $z\in\mathbb Z^3$, the translation by $nz$ maps the open cell
\[
Q^\circ_{\mathfrak c}=\left\{y:\ \mathfrak c_\mu\cdot\mathrm{hc}<y_\mu<(\mathfrak c_\mu+1)\cdot\mathrm{hc},\ \mu=1,2,3\right\}
\]
onto the set $\left\{y:\ \mathfrak c_\mu\cdot\mathrm{hc}+nz_\mu<y_\mu<(\mathfrak c_\mu+1)\cdot\mathrm{hc}+nz_\mu,\ \mu=1,2,3\right\}$, and $n=(n/\mathrm{hc})\cdot\mathrm{hc}=2\mathrm{kt}\cdot L\cdot\mathrm{hc}$.
Therefore
\begin{equation}\tag{2.2}
Q^\circ_{\mathfrak c}+nz=Q^\circ_{\mathfrak c+2\mathrm{kt}\cdot L\cdot z}\qquad\left(\mathfrak c,z\in\mathbb Z^3\right).
\end{equation}
The cells of the torus are the orbits of $n\mathbb Z^3$ on the set of the open cells (Section~0, paragraph \emph{Cells of a torus}).
By (2.2) they are indexed by $\mathfrak c\in\left(\mathbb Z/(2\mathrm{kt}\cdot L)\mathbb Z\right)^3$: there are $2\mathrm{kt}\cdot L$ cells per side and $(2\mathrm{kt}\cdot L)^3$ cells in all.

\emph{The inequality $\mathrm{kt}\ge L$.}
Condition (L1) of Proposition~14.2\,(b) reads $L^{\frac18-\frac{3\bar\varepsilon}{2}}\ge2400$, and $\bar\varepsilon=1/200$ (Section~I.8.2, line (C0));
so the exponent is $\frac18-\frac{3}{400}$, a number in $(0,1)$.
Since $L\ge1$, it follows that $L\ge L^{\frac18-\frac{3}{400}}\ge2400>8$.
If $T=T^*(e)$, then $\mathrm{kt}=8<L$; if $T\ge T^*(e)+1$, then $\mathrm{kt}\ge8L\ge L$.

\emph{Walls.}
In this paragraph and in (b), $i$ stands for an arbitrary integer.
For $j=2i$ one has $j\cdot\mathrm{hc}=i\cdot\mathrm{per}\in\mathbb Z$, so that the wall $\{y_\mu=j\cdot\mathrm{hc}\}$ is a site plane (Section~I.4.8).
For $j=2i+1$ one has $j\cdot\mathrm{hc}=i\cdot\mathrm{per}+\mathrm{per}/2=i\cdot\mathrm{per}+m_{\mathrm p}+1/2\in\mathbb Z+1/2$, and the wall is a link plane.

\emph{The two mirrors.}
By Section~I.4.8, $\theta_{y_\mu,\mathrm{po}}(y)=y-2\left(y_\mu-\mathrm{po}\right)\cdot e_\mu$.
Hence $\theta_{y_\mu,n/2}(y)=\theta_{y_\mu,0}(y)+n\cdot e_\mu$, and the two reflections induce one and the same map of $\mathbb R^3/n\mathbb Z^3$.
A point $y$ is fixed by this map if and only if $2y_\mu\in n\mathbb Z$, that is, if and only if it lies on one of the planes $\{y_\mu=0\}$, $\{y_\mu=n/2\}$ modulo $n$.
The complement of these two planes in the torus consists of the two sets $\{0<y_\mu<n/2\}$ and $\{n/2<y_\mu<n\}$ (modulo $n$), which the map interchanges.
Now $0=0\cdot\mathrm{hc}$, and $n/2=\mathrm{kt}\cdot L\cdot\mathrm{per}/2=\mathrm{kt}\cdot L\cdot\mathrm{hc}$, where $\mathrm{kt}\cdot L=8L^{T-T^*(e)+1}$ is even.
So both planes are walls $\{y_\mu=j\cdot\mathrm{hc}\}$ with even $j$, and by the preceding paragraph both are site planes; indeed $n/2=4h\cdot L^{T+N}$ is an integer.

\textbf{The hypotheses of Proposition~47.1.}
From now on $\beta$ fits $(N,T^*(e))$, that is, $D_\beta+3\eta\le\ell_c$, and $v$ is a placement; in particular $v\in\eta\mathbb Z^3$ (Section~I.7.6).
One has $0<\eta=L^{-N}\le1$ because $N\ge0$, and $n=R/\eta$ is an integer $\ge8$ by (a).
By Section~I.6.4 and $T\ge T^*(e)$,
\[
D_\beta<\ell_c=\frac{R^*(e)}{2L}<R^*(e)=h\cdot L^{T^*(e)}\le h\cdot L^{T}<R.
\]
Hence the data of the grades~0 and~1 of Proposition~47.1 are given on $\Lambda_{\eta,R}$, with an arbitrary vector of $\mathbb R^3$ in place of $v$.
On this torus, $\hat{\mathfrak a}=\hat{\mathfrak a}[\tau_v\beta]$, $\boldsymbol{\sigma}^2$, $X$, $\hat{\mathfrak a}_{\mathfrak c}$ and $X_{\mathfrak c}$ are the objects of Section~I.7.6.

\textbf{(b).}
By (2.1), $\left(R^*(e)/L\right)\mathbb Z^3=\mathrm{per}\cdot\eta\,\mathbb Z^3$, which is a subset of $\eta\mathbb Z^3$.
Let $y=\mathrm{per}\cdot\eta\cdot z$ with $z\in\mathbb Z^3$; then $\mathfrak c(y)=2y/(\mathrm{per}\cdot\eta)=2z\in2\mathbb Z^3$.
According to Section~I.6.4, $\mathsf g_{\mathfrak c}$ acts in the coordinate $\mu$ by $y'_\mu\mapsto i\cdot\mathrm{per}+y'_\mu$ if $\mathfrak c_\mu=2i$;
for $\mathfrak c=2z$ this is $y'_\mu\mapsto z_\mu\cdot\mathrm{per}+y'_\mu$.
So $\mathsf g_{\mathfrak c(y)}$ is the translation by $\mathrm{per}\cdot z=y/\eta=\mathrm{hc}\cdot\mathfrak c(y)$, in agreement with Proposition~41\,(e):
\[
\mathsf g_{\mathfrak c(y)}=(y/\eta,1)\in W_{\mathrm{af}}.
\]
The sum $\sum_\mu\mathfrak c(y)_\mu=2\sum_\mu z_\mu$ is even.
Hence, by the definition of the mirror copies in Section~I.6.4, every $O\in\mathcal A(Q^\circ_0)$ that is a function of two-cochains $u$ satisfies
\[
O^{[\mathfrak c(y)]}(u)=O\left(\mathsf g_{\mathfrak c(y)}^*u\right),
\]
without complex conjugation.

The vectors $v$ and $v+y$ belong to $\eta\mathbb Z^3$.
By (1.8) for $v+y$ and for $v$, and by the third formula of (0.3),
\begin{equation}\tag{2.3}
\begin{split}
\hat{\mathfrak a}[\tau_{v+y}\beta]
&=\left((v+y)/\eta,1\right)_*\hat{\mathfrak a}[\beta]
=(y/\eta,1)_*(v/\eta,1)_*\hat{\mathfrak a}[\beta]\\
&=(y/\eta,1)_*\hat{\mathfrak a}[\tau_v\beta]
=\left(\mathsf g_{\mathfrak c(y)}\right)_*\hat{\mathfrak a}
=\hat{\mathfrak a}_{\mathfrak c(y)}.
\end{split}
\end{equation}
The definition of $X_{\mathfrak c}$ (Section~I.7.6), (2.3), and Proposition~47.1\,(b2$'$) with the vector $v+y$ in place of $v$ yield
\[
X_{\mathfrak c(y)}=\left\langle d\hat{\mathfrak a}_{\mathfrak c(y)},F^{\mathrm u}\right\rangle=\left\langle d\hat{\mathfrak a}[\tau_{v+y}\beta],F^{\mathrm u}\right\rangle=X_{\eta,R}\left(\tau_{v+y}f^{(\eta)}\right).
\]

\textbf{(c).}
By (2.2), two indices $\mathfrak c,\mathfrak c'\in\mathbb Z^3$ give the same cell of the torus if and only if $\mathfrak c-\mathfrak c'\in2\mathrm{kt}\cdot L\,\mathbb Z^3=(2n/\mathrm{per})\,\mathbb Z^3$.
For $y,y'\in\mathrm{per}\cdot\eta\,\mathbb Z^3$,
\[
\mathfrak c(y)-\mathfrak c(y')=\frac{2(y-y')}{\mathrm{per}\cdot\eta},
\]
and this vector lies in $(2n/\mathrm{per})\,\mathbb Z^3$ if and only if $(y-y')/\eta\in n\mathbb Z^3$, that is, if and only if $y-y'\in n\eta\,\mathbb Z^3=R\mathbb Z^3$.

\emph{The sufficient condition.}
Suppose that $R>2\max_{y\in\mathrm{Ypt}}|y|_\infty$, and let $y,y'\in\mathrm{Ypt}$ with $y-y'\in R\mathbb Z^3$.
Then $|y-y'|_\infty\le|y|_\infty+|y'|_\infty<R$.
Since $0$ is the only vector of $R\mathbb Z^3$ of sup norm less than $R$, it follows that $y=y'$.

\textbf{(d).}
Let $g_{\mathrm w}:=(v/\eta,1)\in W_{\mathrm{af}}$, and for $y\in\mathrm{Ypt}$ put $\mathrm{gc}_y:=\eta^{\frac32}\cdot\left(\tau_yf^{(\eta)}\right)_\eta\in C^2(n)$.
By Section~I.7.5 ($X_{\eta,R}(f')=\langle f'_\eta,F\rangle_\eta$), by Section~I.7.1 ($F=\eta^{-\frac32}F^{\mathrm u}$) and by $\langle u,u'\rangle_\eta=\eta^3\cdot\langle u,u'\rangle$ (Subsection~0.1), every test $f'$ satisfies
\[
X_{\eta,R}(f')=\eta^3\cdot\left\langle f'_\eta,\eta^{-\frac32}F^{\mathrm u}\right\rangle=\left\langle\eta^{\frac32}f'_\eta,F^{\mathrm u}\right\rangle.
\]
In particular $X_{\eta,R}\left(\tau_yf^{(\eta)}\right)=\left\langle\mathrm{gc}_y,F^{\mathrm u}\right\rangle$.
Using $\tau_{v+y}=\tau_v\tau_y$, Lemma~4\,(c) for the test $\tau_yf^{(\eta)}$ and the vector $v\in\eta\mathbb Z^3$, and the first formula of (0.3) for the counting pairing (Subsection~0.3), one obtains
\begin{equation}\tag{2.4}
\begin{split}
X_{\eta,R}\left(\tau_{v+y}f^{(\eta)}\right)\left(F^{\mathrm u}\right)
&=\left\langle\left(g_{\mathrm w}\right)_*\mathrm{gc}_y,F^{\mathrm u}\right\rangle
=\left\langle\mathrm{gc}_y,g_{\mathrm w}^*F^{\mathrm u}\right\rangle\\
&=X_{\eta,R}\left(\tau_yf^{(\eta)}\right)\left(g_{\mathrm w}^*F^{\mathrm u}\right).
\end{split}
\end{equation}
The identity $(\tau_vf')_\eta=(v/\eta,1)_*f'_\eta$ of Lemma~4\,(c) is the first identity of \cite[Lemma~7.6]{QEDI}, here for $n$ of either parity.
Put
\[
\Phi(u):=\mathrm{Pol}\left(\left(\left\langle\mathrm{gc}_y,u\right\rangle\right)_{y\in\mathrm{Ypt}}\right)\qquad\left(u\in C^2(n)\right).
\]
The function $\Phi$ is continuous.
Let $\mathrm{dg}$ be the degree of $\mathrm{Pol}$, $\mathrm{cf}_1$ the sum of the moduli of its coefficients, and $\mathrm{cf}_2:=\max_y\|\mathrm{gc}_y\|$, with $\mathrm{cf}_2:=0$ if $\mathrm{Ypt}$ is empty.
Then $\left|\left\langle\mathrm{gc}_y,u\right\rangle\right|\le\mathrm{cf}_2\cdot\|u\|$, and since $1+\mathrm{tq}\le\exp(\mathrm{tq})$ for real $\mathrm{tq}\ge0$, every monomial of degree at most $\mathrm{dg}$ in these numbers has modulus at most
\[
\left(1+\mathrm{cf}_2\cdot\|u\|\right)^{\mathrm{dg}}\le\exp\left(\mathrm{dg}\cdot\mathrm{cf}_2\cdot\|u\|\right).
\]
Therefore $|\Phi(u)|\le\mathrm{cf}_1\cdot\exp\left(\mathrm{dg}\cdot\mathrm{cf}_2\cdot\|u\|\right)$, and $\Phi$ lies in $\mathfrak F(n)$ (Section~G6.A3), with $\mathrm{cf}:=\max\left(\mathrm{cf}_1,\mathrm{dg}\cdot\mathrm{cf}_2\right)$.
Because $\left\|g_{\mathrm w}^*u\right\|=\|u\|$ (Subsection~0.3), the function $\Phi\circ g_{\mathrm w}^*$ satisfies the same bound and lies in $\mathfrak F(n)$ as well.
Consequently $E^{\mathrm{af}}_n\Phi$ and $E^{\mathrm{af}}_n\left(\Phi\circ g_{\mathrm w}^*\right)$ are defined (Section~I.7.2).
By (2.4), the left side of (d) is $E^{\mathrm{af}}_n\left(\Phi\circ g_{\mathrm w}^*\right)$, and the right side is $E^{\mathrm{af}}_n\Phi$.
Since $n\ge4$ by (a), the two are equal by Proposition~50\,(g$''$).

\textbf{(e).}
By Proposition~47.1\,(b2), $g_N=f^{(\eta)}$ is real and smooth, and its support lies in the compact set $\operatorname{supp}\beta+[-\eta,\eta]^3$;
as $\eta\le1<2$, it lies in $Q_\beta$ by (cu2).
Every $y\in Q_\beta=y_\beta+[-D_\beta/2,D_\beta/2]^3$ satisfies
\[
|y|_\infty\le|y_\beta|_\infty+D_\beta/2=r_g.
\]
The number $r_g$ depends only on the probe, with its cube, and not on $N$, $e$ or $T$.
By Proposition~47.1\,(b3) at $s=0$, $\|g_N-f\|_2\le\eta^2\|\beta\|_{H^4}$.
By Proposition~47.1\,(b5), and since $\boldsymbol{\sigma}\ge0$,
\begin{multline*}
\left|\boldsymbol{\sigma}^2-\|f\|_2^2\right|=\left|\boldsymbol{\sigma}-\|f\|_2\right|\cdot\left(\boldsymbol{\sigma}+\|f\|_2\right)\\
\le\left(\eta^2\|\beta\|_{H^4}+21\eta\cdot|\beta|_{C^3}\cdot D_\beta^{\frac32}\right)\cdot\left(2\|f\|_2+\|\beta\|_{H^4}\right),\qquad\eta=L^{-N}.
\end{multline*}
The width $\ell_c=R^*(e)/(2L)$ does not involve $N$.
Hence a probe that fits $(N,T^*(e))$ also fits $(N',T^*(e))$ for every $N'\ge N$:
\[
D_\beta+3L^{-N'}\le D_\beta+3L^{-N}\le\ell_c.
\]
At every cutoff a placement exists: for each $\mu$ there is a point $v_\mu\in\eta\mathbb Z$ at distance at most $\eta/2$ from the real number $\ell_c/2-y_{\beta,\mu}$, and the modulus in the definition of a placement (Section~I.7.6) is the sup norm (Section~I.1.1).
So $\boldsymbol{\sigma}^2$ is defined for every $N'\ge N$, the preceding bound holds for each such $N'$, and its right side tends to $0$ as $N\to\infty$.
By Proposition~47.1\,(b5$'$), $\boldsymbol{\sigma}^2$ depends only on $(\beta,\eta)$; it depends neither on the placement nor on $T$.
Finally, $R^*(e)/(L\eta)=\mathrm{per}$ by (2.1), and $\mathrm{per}$ is an integer.

\textbf{Uniformity.}
The assertions (a)--(c) are identities between integers and between sets; they hold for all data of the statement.
In (d), the charge, the mass $\boldsymbol\mu$ and the parity of $n$ enter only through Proposition~50\,(g$''$), which is valid for every real $\boldsymbol\mu$, every $q>0$ and every $n\ge4$.
The bounds in (e) involve $\eta$ and the probe only.
\end{proof}

\subsection*{3. Conditions on the constants}

\begin{center}
\begin{tabular}{p{0.15\textwidth}p{0.4\textwidth}p{0.33\textwidth}}
condition & on which constant; where it stands & used in\\
\hline
$L$ is odd
& $L$; Section~I.0.5; Section~I.8.2, line (C3); Proposition~14.2\,(b), (L5)
& proof of Proposition~47.2, \emph{The numbers}: $\mathrm{per}$ is odd, $\mathrm{hc}=m_{\mathrm p}+1/2$; proof of (a): the walls with odd $j$ are link planes\\
\hline
$L>8$
& $L$; Section~I.8.2, line (C3); Proposition~14.2\,(b), (L1) (the inequality $L^{\frac18-\frac{3\bar\varepsilon}{2}}\ge2400$, with $\bar\varepsilon=1/200$ of the line (C0), gives $L\ge2400$); Proposition~14.2\,(f); Annex~B, Section~G6.Z4, first row
& proof of Proposition~47.2\,(a): $\mathrm{kt}\ge L$ if and only if $T\ge T^*(e)+1$\\
\hline
$M$ is a power of $L$, $M\ge L^2$
& $M$; Section~I.0.5; Section~I.8.2, line (C5); Proposition~14.2\,(b), (M1)
& proof of Proposition~47.2, \emph{The numbers}: $T^*(e)$ is an integer $\ge1$, so that $\mathrm{per}=n_*/L$ is an integer
\end{tabular}
\end{center}

All three conditions are implied by the lines (C3), (C5) of Section~I.8.2, more precisely by the conditions (L1), (L5), (M1) of Proposition~14.2\,(b).
Proposition~47.1 holds without any condition on a constant.

The remaining hypotheses in Section~2 are restrictions on the data of the statements;
they impose no condition on the constants, and they serve as follows.
\begin{itemize}
\item $N$ and $T$ are integers $\ge0$ (Sections~I.0.5 and I.6.5); hence the exponent of $\mathrm{per}$ is $\ge0$, $\eta\le1$, and $\mathrm{kt}$ is an integer.
\item $\bar m\in(0,1]$; hence $h$ is a power of $L$ with an exponent $\ge0$.
\item $e\cdot h^{\frac12}\le e_1$ and $N\ge N_0(e)$: this is the range in which Section~I.6.4 defines $n_*$, $\mathrm{per}$, $\mathrm{hc}$, $\ell_c$, $\mathcal Q_*$; the arithmetic of Section~2 remains valid without these two inequalities.
\item $T\ge T^*(e)$; hence $\mathrm{kt}$ is an integer and $R>D_\beta$.
\item $\beta$ fits $(N,T^*(e))$; hence $D_\beta<\ell_c$, and a placement is defined.
\end{itemize}

\section*{Chapter 41. The probes live in the open cell: proofs of Proposition~41}

\subsection*{Proof}

All objects of Proposition~41 are cochains of the photon lattice in lattice units, that is, of the lattice of level~$0$, together with functions of such cochains.

\subsubsection*{0. Conventions of this chapter}

\paragraph{0.1. Numbers, units.}
By Section~I.0.5, $L$ is an odd integer with $L\ge5$ and $N\ge0$ is an integer;
hence $0<\eta=L^{-N}\le1$ and $2/\eta\ge2$.
Since $0<\bar{m}\le1$ (Section~I.0.5), the number $h=L^{\lfloor\log_L(1/\bar{m})\rfloor}$ is a power of $L$ with a non-negative integer exponent.
The same is true of $r_K(e)$ (Section~I.6.1),
and $M$ is a power of $L$ with $M\ge L^2$ (Section~I.0.5).
Consequently $M\cdot r_K(e)$ is a power of $L$ which is at least $L^2$,
and $T^{*}(e)=\log_L\left(M\cdot r_K(e)\right)-1$ is an integer $\ge1$ (Section~I.6.2).
Section~I.6.4 gives $n_*=h\cdot L^{T^{*}(e)+N}$, so that
\[
\mathrm{per}=\frac{n_*}{L}=h\cdot L^{T^{*}(e)+N-1}
\]
is a power of the odd integer $L$ with a non-negative integer exponent, and therefore an odd integer $\ge1$.
Thus $\mathrm{per}=2m_{\mathrm{p}}+1$ with an integer $m_{\mathrm{p}}\ge0$,
$\mathrm{hc}=\mathrm{per}/2=m_{\mathrm{p}}+\frac12$,
and $\ell_c=\mathrm{hc}\cdot\eta$.
On $\mathbb{R}^3$, $|\cdot|$ denotes the sup norm (Section~I.1.1).
Physical coordinates are $\eta$ times lattice coordinates:
the map $c\mapsto\eta\cdot c$ sends the cells of $\mathbb{Z}^3$ onto the cells of $\eta\mathbb{Z}^3$,
and a cochain on $\eta\mathbb{Z}^3$ is identified with the cochain on $\mathbb{Z}^3$ which takes the same values on corresponding cells (Section~I.7.5).
The walls of $\mathcal{Q}_{*}$ are the planes $\{y_\mu=j\cdot\mathrm{hc}\}$ in lattice coordinates and $\{y_\mu=j\cdot\ell_c\}$ in physical coordinates, where $j\in\mathbb{Z}$ and $\mu=1,2,3$ (Section~I.6.4).

\paragraph{0.2. Cochains on $\mathbb{Z}^3$.}
For $j=1,2$ let $C^{j}$ be the space of the real functions on the $j$-cells of $\mathbb{Z}^3$ (the bonds for $j=1$, the plaquettes for $j=2$),
and let $C^{j}_{\mathrm{c}}\subset C^{j}$ be the subspace of the finitely supported functions \cite[Definition~2.1]{QEDII}.
For $u\in C^{j}_{\mathrm{c}}$ and $u'\in C^{j}$ set
\[
\langle u,u'\rangle=\langle u',u\rangle:=\sum_{\boldsymbol{\sigma}}u(\boldsymbol{\sigma})\cdot u'(\boldsymbol{\sigma}),\qquad\|u\|^2:=\langle u,u\rangle ,
\]
where $\boldsymbol{\sigma}$ ranges over the $j$-cells; only finitely many terms of the sum are non-zero.
An element $A$ of $C^1$ is regarded as a function on the ordered pairs of nearest neighbours,
with $A(c,c+e_\nu)=A_\nu(c)$ and $A(c',c)=-A(c,c')$ (Section~I.1.1).
The product $A(c,c')\cdot A'(c,c')$ does not change when the two points are exchanged,
and $\langle A,A'\rangle$ is the sum of these products over the unordered pairs $\{c,c'\}$.

The plaquette $p=(c;\mu\nu)$, $\mu<\nu$, has the corners
\[
\mathrm{vx}_0:=c,\qquad\mathrm{vx}_1:=c+e_\mu,\qquad\mathrm{vx}_2:=c+e_\mu+e_\nu,\qquad\mathrm{vx}_3:=c+e_\nu ,
\]
taken in this order (Section~I.1.1);
the index of a corner is read modulo~$4$.
Regarded as a closed subset of $\mathbb{R}^3$, the plaquette is the convex hull of its corners,
and the bond $\{c,c'\}$ is the segment with the end points $c$ and $c'$ \cite[Definition~2.1]{QEDII}.
For $u\in C^2$ write $u_{\mu\nu}(c):=u(c;\mu\nu)$.
In lattice units the coboundary of Section~I.1.1 reads, because $A(\mathrm{vx}_2,\mathrm{vx}_3)=-A_\mu(c+e_\nu)$ and $A(\mathrm{vx}_3,\mathrm{vx}_0)=-A_\nu(c)$,
\begin{equation}
(dA)(p)=A_\mu(c)+A_\nu(c+e_\mu)-A_\mu(c+e_\nu)-A_\nu(c)=\sum_{i\bmod 4}A(\mathrm{vx}_i,\mathrm{vx}_{i+1})\qquad(A\in C^1).\tag{0.1}
\end{equation}
If $(dA)(p)\neq0$, then one of the four bonds $\{\mathrm{vx}_i,\mathrm{vx}_{i+1}\}$ of $p$ belongs to $\operatorname{supp}A$.
Each bond is a bond of four plaquettes.
It follows that $d$ maps $C^1_{\mathrm{c}}$ into $C^2_{\mathrm{c}}$.

\paragraph{0.3. Cochains on the torus; periodisation.}
Let $n$ be an integer in $\mathrm{per}\,\mathbb{Z}$ with $n\ge n_*$.
Then $n\ge15$:
this inequality is part of display~(6.2),
and that part of (6.2) follows from the hypotheses of Proposition~41 and from the inequality $m_{\mathrm{p}}\ge7$ established in Section~3.
In particular $n$ is a side in the sense of Section~I.7.1, where $n\ge4$ is required.
Let $C^{j}(n)$ be the space of the $u\in C^{j}$ with $u(\boldsymbol{\sigma}+n\cdot z)=u(\boldsymbol{\sigma})$ for all $z\in\mathbb{Z}^3$ (Section~G6.A1; Section~I.7.1).
Choose a set $\mathrm{Rep}^j$ of representatives of the $j$-cells modulo $n\mathbb{Z}^3$, and put
\[
\langle u,u'\rangle_n:=\sum_{\boldsymbol{\sigma}\in\mathrm{Rep}^j}u(\boldsymbol{\sigma})\cdot u'(\boldsymbol{\sigma}),\qquad\|u\|_n^2:=\langle u,u\rangle_n\qquad\left(u,u'\in C^{j}(n)\right).
\]
This is the counting pairing of Section~I.7.1, which is written $\langle\,,\rangle$ in Proposition~41.
Since the summand is $n\mathbb{Z}^3$-periodic, the value does not depend on the choice of $\mathrm{Rep}^j$.
By (0.1) the coboundary $d$ maps $C^1(n)$ into $C^2(n)$;
let $d^{*}_{n}\colon C^2(n)\to C^1(n)$ denote the adjoint of this map with respect to $\langle\,,\rangle_n$.
The periodisation is the map
\[
\mathrm{Prd}_n\colon C^{j}_{\mathrm{c}}\to C^{j}(n),\qquad(\mathrm{Prd}_n u)(\boldsymbol{\sigma}):=\sum_{z\in\mathbb{Z}^3}u(\boldsymbol{\sigma}+n\cdot z),
\]
in which only finitely many terms are different from zero.
In this chapter $z$ and $z'$ denote vectors of $\mathbb{Z}^3$;
only as the argument of the polynomials $\mathrm{Pol}$, $H_2$, $H_4$ does $z$ denote the variable of a polynomial.

\paragraph{0.4. The action of $W_{\mathrm{af}}$.}
An element $g_{\mathrm{w}}=(a,r)$ of $W_{\mathrm{af}}$, with $a\in\mathbb{Z}^3$ and $r\in O_{\mathrm{cub}}$, is the map $y\mapsto ry+a$ of $\mathbb{R}^3$ (Section~I.4.7).
It is an isometry with $g_{\mathrm{w}}\mathbb{Z}^3=\mathbb{Z}^3$;
hence it sends pairs of nearest neighbours to pairs of nearest neighbours, and plaquettes onto plaquettes.
As in Section~I.4.7, put
\[
(g_{\mathrm{w}}^{*}A)(c,c'):=A(g_{\mathrm{w}}c,g_{\mathrm{w}}c')\quad\text{on }C^1,\qquad
(g_{\mathrm{w}}^{*}u)(p):=\epsilon(g_{\mathrm{w}},p)\cdot u(g_{\mathrm{w}}p)\quad\text{on }C^2,
\]
where $g_{\mathrm{w}}p$ is the plaquette with the corners $g_{\mathrm{w}}\mathrm{vx}_0,\dots,g_{\mathrm{w}}\mathrm{vx}_3$, and where the sign $\epsilon(g_{\mathrm{w}},p)$ is defined as follows.
With respect to the relation of being nearest neighbours, the four corners of a plaquette form a cycle of length four.
A sequence of the four corners in which any two consecutive members (indices modulo~$4$) are nearest neighbours is therefore determined by its first member and by the choice of one of the two neighbours of this member as the second member.
Hence there are exactly eight such sequences:
the four cyclic shifts of the sequence of the corners in their order, and the four cyclic shifts of the reversed sequence.
The sequence $(g_{\mathrm{w}}\mathrm{vx}_0,g_{\mathrm{w}}\mathrm{vx}_1,g_{\mathrm{w}}\mathrm{vx}_2,g_{\mathrm{w}}\mathrm{vx}_3)$ is one of the eight sequences of the plaquette $g_{\mathrm{w}}p$.
Set $\epsilon(g_{\mathrm{w}},p):=+1$ if it is a cyclic shift of the corners of $g_{\mathrm{w}}p$ in their order,
and $\epsilon(g_{\mathrm{w}},p):=-1$ if it is a cyclic shift of the reversed sequence.
This sign is the comparison of orientations of Section~I.4.7.
Indeed, the first two steps of the sequence are $g_{\mathrm{w}}\mathrm{vx}_1-g_{\mathrm{w}}\mathrm{vx}_0=re_\mu$ and $g_{\mathrm{w}}\mathrm{vx}_2-g_{\mathrm{w}}\mathrm{vx}_1=re_\nu$.
Let $\mu'<\nu'$ be the directions of $g_{\mathrm{w}}p$.
The corners of $g_{\mathrm{w}}p$ in their order have the steps $e_{\mu'},e_{\nu'},-e_{\mu'},-e_{\nu'}$,
and in each of the four cases the wedge product of a step with the following one is $e_{\mu'}\wedge e_{\nu'}$;
the reversed sequence has the steps $e_{\mu'},-e_{\nu'},-e_{\mu'},e_{\nu'}$,
and in each of the four cases the wedge product of a step with the following one is $-e_{\mu'}\wedge e_{\nu'}$.
Therefore $\epsilon(g_{\mathrm{w}},p)$ is the sign in
\[
re_\mu\wedge re_\nu=\epsilon(g_{\mathrm{w}},p)\cdot e_{\mu'}\wedge e_{\nu'} .
\]

Further, $(g_{\mathrm{w}})_{*}:=(g_{\mathrm{w}}^{-1})^{*}$, and $(g_{\mathrm{w}}g_{\mathrm{w}}')^{*}=(g_{\mathrm{w}}')^{*}g_{\mathrm{w}}^{*}$ (Section~I.4.7).
On $C^1$ the last identity is the computation
\[
\left((g_{\mathrm{w}}g_{\mathrm{w}}')^{*}A\right)(c,c')=A(g_{\mathrm{w}}g_{\mathrm{w}}'c,g_{\mathrm{w}}g_{\mathrm{w}}'c')=(g_{\mathrm{w}}^{*}A)(g_{\mathrm{w}}'c,g_{\mathrm{w}}'c')=\left((g_{\mathrm{w}}')^{*}g_{\mathrm{w}}^{*}A\right)(c,c').
\]
On $C^2$ let $r$, $r'$ be the linear parts of $g_{\mathrm{w}}$, $g_{\mathrm{w}}'$,
let $\mu'<\nu'$ be the directions of $g_{\mathrm{w}}'p$, and $\mu''<\nu''$ those of $g_{\mathrm{w}}g_{\mathrm{w}}'p$.
Then
\[
rr'e_\mu\wedge rr'e_\nu=\epsilon(g_{\mathrm{w}}',p)\cdot re_{\mu'}\wedge re_{\nu'}=\epsilon(g_{\mathrm{w}}',p)\cdot\epsilon(g_{\mathrm{w}},g_{\mathrm{w}}'p)\cdot e_{\mu''}\wedge e_{\nu''} ,
\]
so that $\epsilon(g_{\mathrm{w}}g_{\mathrm{w}}',p)=\epsilon(g_{\mathrm{w}}',p)\cdot\epsilon(g_{\mathrm{w}},g_{\mathrm{w}}'p)$, and
\[
\left((g_{\mathrm{w}}g_{\mathrm{w}}')^{*}u\right)(p)=\epsilon(g_{\mathrm{w}}',p)\cdot\epsilon(g_{\mathrm{w}},g_{\mathrm{w}}'p)\cdot u(g_{\mathrm{w}}g_{\mathrm{w}}'p)=\epsilon(g_{\mathrm{w}}',p)\cdot(g_{\mathrm{w}}^{*}u)(g_{\mathrm{w}}'p)=\left((g_{\mathrm{w}}')^{*}g_{\mathrm{w}}^{*}u\right)(p).
\]
It follows that
\[
(g_{\mathrm{w}})_{*}g_{\mathrm{w}}^{*}=(g_{\mathrm{w}}g_{\mathrm{w}}^{-1})^{*}=1,\qquad
g_{\mathrm{w}}^{*}(g_{\mathrm{w}})_{*}=1,\qquad
(g_{\mathrm{w}}g_{\mathrm{w}}')_{*}=\left((g_{\mathrm{w}}')^{-1}g_{\mathrm{w}}^{-1}\right)^{*}=(g_{\mathrm{w}})_{*}(g_{\mathrm{w}}')_{*}.
\]
For $z\in\mathbb{Z}^3$ let $\tau_z:=(z,1)$ be the translation $y\mapsto y+z$.
It sends the corners of $p$, in their order, to the corners of $p+z$, in their order;
hence $\epsilon(\tau_z,p)=1$, and
\begin{equation}
\left((\tau_z)_{*}u\right)(\boldsymbol{\sigma})=u(\boldsymbol{\sigma}-z)\qquad(u\in C^{j},\ j=1,2).\tag{0.2}
\end{equation}

\paragraph{0.5. Letters.}
Let $c\in\eta\mathbb{Z}^3$ and $\nu\in\{1,2,3\}$.
The centre of the bond $(c,c+\eta\cdot e_\nu)$ of $\eta\mathbb{Z}^3$ is the point $c+\frac{\eta}{2}\cdot e_\nu$, in physical coordinates;
it is the point denoted $c(\sigma)$ in \cite[Section~7.1.1]{QEDI}, for the $1$-cell $\sigma=(c;\{\nu\})$ of that section.
In this chapter the letter $\mathrm{ce}$, written without an argument, denotes this point for the bond under consideration;
in Section~I.7.5 the same letter denotes the centre of a plaquette.
In the sense of \cite[Definition~2.12]{QEDII}, $G$ and $G_1$ denote integrands, $\Phi$ the bounded factor of a product,
and $\mathrm{Cq}$, $\varepsilon_G$ the two constants of the growth condition (the first of them is written $C$ there);
the letter $G$ is also used as in Section~I.4.8.
The translation part and the linear part of $\mathsf{g}_{\mathfrak{c}}$ are $a_{\mathfrak{c}}$ and $r_{\mathfrak{c}}$, as in Proposition~49 and in Proposition~35.2\,(a).
$\mathrm{Pol}$ and $\hat{\mathfrak{a}}^{\mathrm{o}}$ have the meaning of Proposition~41, and $\mathrm{cf}_j$ denotes the coefficients of $\mathrm{Pol}$.
The cochains $u^{\mathrm{o}}$ and $a^{\mathrm{o}}_{\eta}$ of the displays (2.2) and (2.3) are the cochains $u^{\mathrm{o}}$, $a^{\mathrm{o}}_{\eta}$ of Proposition~47.1\,(b2$^{\mathrm{o}}$) for the probe $\beta^{v}$.
Likewise $\delta_{\eta,\mu}$ is the difference quotient of Proposition~47.1,
and the functions $a'$, $\tilde{a}_{\eta}$ of display (2.1) are the functions $a'$, $\tilde{a}_{\eta}$ of Proposition~47.1 for the probe $\beta^{v}$.
In this chapter $a'$ always stands for the vector field $(a'_\nu)$ of (2.1), and never for the constant $a'$ of line~(C4r) in Section~I.8.2.
A cell is denoted $\boldsymbol{\sigma}$, as in Sections~I.4.7 and~I.7.5;
the number of Section~I.7.6 is written only in the forms $\boldsymbol{\sigma}^2$ and $\boldsymbol{\sigma}^4=(\boldsymbol{\sigma}^2)^2$.

\subsubsection*{1. Four lemmas on cochains}

\begin{named}{Lemma 1}[The adjoint of $d$ on $\mathbb{Z}^3$]
For $u\in C^2_{\mathrm{c}}$ set $u_{\nu\mu}:=-u_{\mu\nu}$ and
\begin{equation}
(d^{*}u)_\nu(c):=-\sum_{\mu\neq\nu}\left[u_{\mu\nu}(c)-u_{\mu\nu}(c-e_\mu)\right]\qquad(c\in\mathbb{Z}^3,\ \nu=1,2,3).\tag{1.1}
\end{equation}
Then $d^{*}u\in C^1_{\mathrm{c}}$, and $\langle dA,u\rangle=\langle A,d^{*}u\rangle$ for every $A\in C^1$, finitely supported or not.
In particular, $d^{*}$ is the adjoint of $d$ on the finitely supported cochains in the sense of \cite[Definition~2.1]{QEDII}.
\end{named}

\begin{proof}
The right side of (1.1) vanishes unless $\operatorname{supp}u$ contains a plaquette whose base point is $c$ or $c-e_\mu$ for some $\mu\neq\nu$;
therefore $d^{*}u$ has finite support.
By (0.1),
\[
\langle dA,u\rangle=\sum_{\mu<\nu}\sum_{c}\left[A_\mu(c)+A_\nu(c+e_\mu)-A_\mu(c+e_\nu)-A_\nu(c)\right]\cdot u_{\mu\nu}(c),
\]
and only finitely many terms of this sum are different from zero.
Replace $c$ by $c-e_\mu$ in the second term and by $c-e_\nu$ in the third term; this gives
\[
\langle dA,u\rangle=\sum_{\mu<\nu}\sum_{c}A_\mu(c)\cdot\left[u_{\mu\nu}(c)-u_{\mu\nu}(c-e_\nu)\right]-\sum_{\mu<\nu}\sum_{c}A_\nu(c)\cdot\left[u_{\mu\nu}(c)-u_{\mu\nu}(c-e_\mu)\right].
\]
After the two indices of the first double sum are renamed ($\mu\to\nu$, $\nu\to\mu$), it becomes
\[
\sum_{\nu<\mu}\sum_{c}A_\nu(c)\cdot\left[u_{\nu\mu}(c)-u_{\nu\mu}(c-e_\mu)\right]=-\sum_{\nu<\mu}\sum_{c}A_\nu(c)\cdot\left[u_{\mu\nu}(c)-u_{\mu\nu}(c-e_\mu)\right].
\]
Combining it with the second double sum one obtains
\[
\langle dA,u\rangle=-\sum_{\nu}\sum_{c}A_\nu(c)\cdot\sum_{\mu\neq\nu}\left[u_{\mu\nu}(c)-u_{\mu\nu}(c-e_\mu)\right]=\langle A,d^{*}u\rangle .
\]
\end{proof}

\begin{named}{Lemma 2}[Periodisation]
\begin{itemize}
\item[(a)] $\langle\mathrm{Prd}_n u,u'\rangle_n=\langle u,u'\rangle$ for $u\in C^{j}_{\mathrm{c}}$ and $u'\in C^{j}(n)$.
\item[(b)] $d\,\mathrm{Prd}_n A=\mathrm{Prd}_n\,dA$ for $A\in C^1_{\mathrm{c}}$.
\item[(c)] $d^{*}_{n}\,\mathrm{Prd}_n u=\mathrm{Prd}_n\,d^{*}u$ for $u\in C^2_{\mathrm{c}}$.
\end{itemize}
\end{named}

\begin{proof}
\begin{itemize}
\item[(a)] The map $(\boldsymbol{\sigma},z)\mapsto\boldsymbol{\sigma}+n\cdot z$ is a bijection of $\mathrm{Rep}^j\times\mathbb{Z}^3$ onto the set of the $j$-cells:
each cell is congruent modulo $n\mathbb{Z}^3$ to exactly one element of $\mathrm{Rep}^j$,
and no translation by a vector $n\cdot z\neq0$ fixes a cell.
Moreover $u'(\boldsymbol{\sigma})=u'(\boldsymbol{\sigma}+n\cdot z)$.
Hence
\[
\langle\mathrm{Prd}_n u,u'\rangle_n=\sum_{\boldsymbol{\sigma}\in\mathrm{Rep}^j}\sum_{z}u(\boldsymbol{\sigma}+n\cdot z)\cdot u'(\boldsymbol{\sigma}+n\cdot z)=\langle u,u'\rangle .
\]
\item[(b)] The corners of $p+n\cdot z$, in their order, are the points $\mathrm{vx}_i+n\cdot z$.
By (0.1),
\[
(d\,\mathrm{Prd}_n A)(p)=\sum_{i\bmod 4}\sum_{z}A(\mathrm{vx}_i+n\cdot z,\mathrm{vx}_{i+1}+n\cdot z)=\sum_{z}(dA)(p+n\cdot z)=(\mathrm{Prd}_n\,dA)(p).
\]
\item[(c)] Let $B\in C^1(n)$.
Then $dB\in C^2(n)$, and
\[
\langle B,d^{*}_{n}\mathrm{Prd}_n u\rangle_n=\langle dB,\mathrm{Prd}_n u\rangle_n=\langle dB,u\rangle=\langle B,d^{*}u\rangle=\langle B,\mathrm{Prd}_n d^{*}u\rangle_n ;
\]
the four equalities hold by the definition of $d^{*}_{n}$, by (a), by Lemma~1 with $A=B$, and by (a) again.
Since $B$ is an arbitrary element of $C^1(n)$ and $\langle\,,\rangle_n$ is positive definite on $C^1(n)$, the assertion follows.
\end{itemize}
\end{proof}

\begin{named}{Lemma 3}[$W_{\mathrm{af}}$ on cochains]
Let $g_{\mathrm{w}}=(a,r)\in W_{\mathrm{af}}$ and $j\in\{1,2\}$.
\begin{itemize}
\item[(a)] $d\,g_{\mathrm{w}}^{*}A=g_{\mathrm{w}}^{*}\,dA$ for every $A\in C^1$.
\item[(b)] $g_{\mathrm{w}}^{*}$ maps $C^{j}_{\mathrm{c}}$ into $C^{j}_{\mathrm{c}}$,
and $\langle g_{\mathrm{w}}^{*}u,g_{\mathrm{w}}^{*}u'\rangle=\langle u,u'\rangle$ for $u\in C^{j}_{\mathrm{c}}$ and $u'\in C^{j}$.
\item[(c)] $g_{\mathrm{w}}^{*}$ maps $C^{j}(n)$ into $C^{j}(n)$;
$\langle g_{\mathrm{w}}^{*}u,g_{\mathrm{w}}^{*}u'\rangle_n=\langle u,u'\rangle_n$ for $u,u'\in C^{j}(n)$;
and $g_{\mathrm{w}}^{*}\,\mathrm{Prd}_n u=\mathrm{Prd}_n\,g_{\mathrm{w}}^{*}u$ for $u\in C^{j}_{\mathrm{c}}$.
\item[(d)] The assertions (a), (b), (c) remain true with $(g_{\mathrm{w}})_{*}$ in place of $g_{\mathrm{w}}^{*}$.
Moreover $\langle u,g_{\mathrm{w}}^{*}u'\rangle=\langle(g_{\mathrm{w}})_{*}u,u'\rangle$ for $u\in C^{j}_{\mathrm{c}}$ and $u'\in C^{j}$,
and $\langle u,g_{\mathrm{w}}^{*}u'\rangle_n=\langle(g_{\mathrm{w}})_{*}u,u'\rangle_n$ for $u,u'\in C^{j}(n)$.
\item[(e)] $\operatorname{supp}(g_{\mathrm{w}})_{*}u=\{g_{\mathrm{w}}\boldsymbol{\sigma}:\boldsymbol{\sigma}\in\operatorname{supp}u\}$ for $u\in C^{j}$.
\item[(f)] $(\tau_{n\cdot z})_{*}u=u$ for $u\in C^{j}(n)$ and $z\in\mathbb{Z}^3$.
\end{itemize}
\end{named}

\begin{proof}
\begin{itemize}
\item[(a)] Let $p$ be a plaquette with the corners $\mathrm{vx}_i$,
and let $\mathrm{vx}'_0,\dots,\mathrm{vx}'_3$ be the corners of $g_{\mathrm{w}}p$ in their order.
By (0.1), $(d\,g_{\mathrm{w}}^{*}A)(p)=\sum_{i\bmod 4}A(g_{\mathrm{w}}\mathrm{vx}_i,g_{\mathrm{w}}\mathrm{vx}_{i+1})$.
If $\epsilon(g_{\mathrm{w}},p)=1$, there is an $i_0$ such that $g_{\mathrm{w}}\mathrm{vx}_i=\mathrm{vx}'_{i+i_0}$ for all $i$,
and the sum equals $\sum_iA(\mathrm{vx}'_i,\mathrm{vx}'_{i+1})=(dA)(g_{\mathrm{w}}p)$.
If $\epsilon(g_{\mathrm{w}},p)=-1$, there is an $i_0$ such that $g_{\mathrm{w}}\mathrm{vx}_i=\mathrm{vx}'_{i_0-i}$ for all $i$, and the sum equals
\[
\sum_iA(\mathrm{vx}'_{i_0-i},\mathrm{vx}'_{i_0-i-1})=\sum_iA(\mathrm{vx}'_{i+1},\mathrm{vx}'_i)=-(dA)(g_{\mathrm{w}}p).
\]
In both cases $(d\,g_{\mathrm{w}}^{*}A)(p)=\epsilon(g_{\mathrm{w}},p)\cdot(dA)(g_{\mathrm{w}}p)=(g_{\mathrm{w}}^{*}dA)(p)$.
\item[(b)] The map $\boldsymbol{\sigma}\mapsto g_{\mathrm{w}}\boldsymbol{\sigma}$ is a bijection of the set of the $j$-cells.
Let $j=2$.
Then $(g_{\mathrm{w}}^{*}u)(p)\neq0$ if and only if $g_{\mathrm{w}}p\in\operatorname{supp}u$, and
\[
\langle g_{\mathrm{w}}^{*}u,g_{\mathrm{w}}^{*}u'\rangle=\sum_p\epsilon(g_{\mathrm{w}},p)^2\cdot u(g_{\mathrm{w}}p)\cdot u'(g_{\mathrm{w}}p)=\sum_pu(p)\cdot u'(p).
\]
Let $j=1$.
Then $(g_{\mathrm{w}}^{*}A)(c,c')\neq0$ if and only if $\{g_{\mathrm{w}}c,g_{\mathrm{w}}c'\}\in\operatorname{supp}A$,
and $\langle g_{\mathrm{w}}^{*}A,g_{\mathrm{w}}^{*}A'\rangle$ is the sum of the products $A(g_{\mathrm{w}}c,g_{\mathrm{w}}c')\cdot A'(g_{\mathrm{w}}c,g_{\mathrm{w}}c')$ over the unordered pairs $\{c,c'\}$, which is $\langle A,A'\rangle$.
\item[(c)] For $z\in\mathbb{Z}^3$ one has $g_{\mathrm{w}}(y+n\cdot z)=g_{\mathrm{w}}y+n\cdot rz$, and $rz\in\mathbb{Z}^3$;
hence $g_{\mathrm{w}}(\boldsymbol{\sigma}+n\cdot z)=g_{\mathrm{w}}\boldsymbol{\sigma}+n\cdot rz$.
The corners of $p+n\cdot z$ and of $g_{\mathrm{w}}p+n\cdot rz$, in their order, are the translates of the corners of $p$ and of $g_{\mathrm{w}}p$, in their order;
therefore $\epsilon(g_{\mathrm{w}},p+n\cdot z)=\epsilon(g_{\mathrm{w}},p)$.
Consequently
\[
(g_{\mathrm{w}}^{*}u)(p+n\cdot z)=\epsilon(g_{\mathrm{w}},p)\cdot u(g_{\mathrm{w}}p+n\cdot rz)=(g_{\mathrm{w}}^{*}u)(p)\qquad\left(u\in C^2(n)\right),
\]
\[
(g_{\mathrm{w}}^{*}A)(c+n\cdot z,c'+n\cdot z)=A(g_{\mathrm{w}}c+n\cdot rz,g_{\mathrm{w}}c'+n\cdot rz)=(g_{\mathrm{w}}^{*}A)(c,c')\qquad\left(A\in C^1(n)\right).
\]
As $r\mathbb{Z}^3=\mathbb{Z}^3$, the bijection $\boldsymbol{\sigma}\mapsto g_{\mathrm{w}}\boldsymbol{\sigma}$ maps the classes modulo $n\mathbb{Z}^3$ onto the classes modulo $n\mathbb{Z}^3$;
hence $g_{\mathrm{w}}\mathrm{Rep}^j:=\{g_{\mathrm{w}}\boldsymbol{\sigma}:\boldsymbol{\sigma}\in\mathrm{Rep}^j\}$ is again a set of representatives.
For $j=2$ this gives
\[
\langle g_{\mathrm{w}}^{*}u,g_{\mathrm{w}}^{*}u'\rangle_n=\sum_{p\in\mathrm{Rep}^2}\epsilon(g_{\mathrm{w}},p)^2\cdot u(g_{\mathrm{w}}p)\cdot u'(g_{\mathrm{w}}p)=\sum_{p\in g_{\mathrm{w}}\mathrm{Rep}^2}u(p)\cdot u'(p)=\langle u,u'\rangle_n ,
\]
because the value of $\langle\,,\rangle_n$ does not depend on the set of representatives (Subsection~0.3);
for $j=1$ it gives
\[
\langle g_{\mathrm{w}}^{*}A,g_{\mathrm{w}}^{*}A'\rangle_n=\sum_{\{c,c'\}\in\mathrm{Rep}^1}A(g_{\mathrm{w}}c,g_{\mathrm{w}}c')\cdot A'(g_{\mathrm{w}}c,g_{\mathrm{w}}c')=\sum_{\{c,c'\}\in g_{\mathrm{w}}\mathrm{Rep}^1}A(c,c')\cdot A'(c,c')=\langle A,A'\rangle_n .
\]
Let $u\in C^2_{\mathrm{c}}$.
Since $z\mapsto rz$ is a bijection of $\mathbb{Z}^3$,
\[
(\mathrm{Prd}_n\,g_{\mathrm{w}}^{*}u)(p)=\sum_z\epsilon(g_{\mathrm{w}},p+n\cdot z)\cdot u(g_{\mathrm{w}}p+n\cdot rz)=\epsilon(g_{\mathrm{w}},p)\cdot\sum_zu(g_{\mathrm{w}}p+n\cdot z)=(g_{\mathrm{w}}^{*}\,\mathrm{Prd}_n u)(p).
\]
Let $A\in C^1_{\mathrm{c}}$.
The same bijection gives
\[
(\mathrm{Prd}_n\,g_{\mathrm{w}}^{*}A)(c,c')=\sum_zA(g_{\mathrm{w}}c+n\cdot rz,g_{\mathrm{w}}c'+n\cdot rz)=\sum_zA(g_{\mathrm{w}}c+n\cdot z,g_{\mathrm{w}}c'+n\cdot z)=(g_{\mathrm{w}}^{*}\,\mathrm{Prd}_n A)(c,c').
\]
\item[(d)] One has $(g_{\mathrm{w}})_{*}=(g_{\mathrm{w}}^{-1})^{*}$ with $g_{\mathrm{w}}^{-1}\in W_{\mathrm{af}}$;
so the assertions (a), (b), (c) for $g_{\mathrm{w}}^{-1}$ are the assertions for $(g_{\mathrm{w}})_{*}$.
By (b) for $(g_{\mathrm{w}})_{*}$ and by $(g_{\mathrm{w}})_{*}g_{\mathrm{w}}^{*}=1$ (Subsection~0.4),
\[
\langle u,g_{\mathrm{w}}^{*}u'\rangle=\langle(g_{\mathrm{w}})_{*}u,(g_{\mathrm{w}})_{*}g_{\mathrm{w}}^{*}u'\rangle=\langle(g_{\mathrm{w}})_{*}u,u'\rangle .
\]
By (c) for $(g_{\mathrm{w}})_{*}$ and by $(g_{\mathrm{w}})_{*}g_{\mathrm{w}}^{*}=1$,
\[
\langle u,g_{\mathrm{w}}^{*}u'\rangle_n=\langle(g_{\mathrm{w}})_{*}u,(g_{\mathrm{w}})_{*}g_{\mathrm{w}}^{*}u'\rangle_n=\langle(g_{\mathrm{w}})_{*}u,u'\rangle_n .
\]
\item[(e)] One has $\left((g_{\mathrm{w}})_{*}u\right)(p)=\epsilon(g_{\mathrm{w}}^{-1},p)\cdot u(g_{\mathrm{w}}^{-1}p)$ and $\left((g_{\mathrm{w}})_{*}A\right)(c,c')=A(g_{\mathrm{w}}^{-1}c,g_{\mathrm{w}}^{-1}c')$;
the value at a cell $\boldsymbol{\sigma}$ is different from zero if and only if $g_{\mathrm{w}}^{-1}\boldsymbol{\sigma}\in\operatorname{supp}u$.
\item[(f)] This follows from (0.2) and from the periodicity of $u$.
\end{itemize}
\end{proof}

\begin{named}{Lemma 4}[The maps $\mathsf{g}_{\mathfrak{c}}$]
For $\mathfrak{c}\in\mathbb{Z}^3$, $y\in\mathbb{R}^3$ and $\mu=1,2,3$,
\begin{equation}
(\mathsf{g}_{\mathfrak{c}}y)_\mu=\left(\mathfrak{c}_\mu+\frac12\right)\cdot\mathrm{hc}+(-1)^{\mathfrak{c}_\mu}\cdot\left(y_\mu-\frac{\mathrm{hc}}{2}\right).\tag{1.2}
\end{equation}
\begin{itemize}
\item[(a)] $\mathsf{g}_{\mathfrak{c}}=(a_{\mathfrak{c}},r_{\mathfrak{c}})$ with
$r_{\mathfrak{c}}=\operatorname{diag}\left((-1)^{\mathfrak{c}_1},(-1)^{\mathfrak{c}_2},(-1)^{\mathfrak{c}_3}\right)\in\{\pm1\}^3$ and $a_{\mathfrak{c}}\in\mathrm{per}\,\mathbb{Z}^3$.
Hence $\mathsf{g}_{\mathfrak{c}}\in W_{\mathrm{af}}$ and $\det r_{\mathfrak{c}}=(-1)^{\mathfrak{c}_1+\mathfrak{c}_2+\mathfrak{c}_3}$;
thus $\mathsf{g}_{\mathfrak{c}}$ is even (Section~I.4.7) if and only if $\sum_\mu\mathfrak{c}_\mu$ is even.
The map $\mathsf{g}_0$ is the identity.
Each $\mathsf{g}_{\mathfrak{c}}$ is a product of reflections in walls of $\mathcal{Q}_{*}$, and therefore belongs to $W(\mathcal{Q}_{*})$ (Section~I.6.4).
\item[(b)] $\left|(\mathsf{g}_{\mathfrak{c}}y)_\mu-\left(\mathfrak{c}_\mu+\frac12\right)\cdot\mathrm{hc}\right|=\left|y_\mu-\frac{\mathrm{hc}}{2}\right|$.
Consequently $\mathsf{g}_{\mathfrak{c}}$ maps $Q^{\circ}_{0}$ onto $Q^{\circ}_{\mathfrak{c}}$,
and it maps the cube $\mathrm{Bx}_0=\frac{\mathrm{hc}}{2}\cdot(1,1,1)+[-r_0,r_0]^3$ of Proposition~41\,(b) onto the cube $\left(\left(\mathfrak{c}_\mu+\frac12\right)\cdot\mathrm{hc}\right)_\mu+[-r_0,r_0]^3$.
\item[(c)] The number $2n/\mathrm{per}$ is an even integer, and for $z\in\mathbb{Z}^3$
\[
\mathsf{g}_{\mathfrak{c}+(2n/\mathrm{per})\cdot z}=\tau_{n\cdot z}\,\mathsf{g}_{\mathfrak{c}},\qquad Q^{\circ}_{\mathfrak{c}+(2n/\mathrm{per})\cdot z}=Q^{\circ}_{\mathfrak{c}}+n\cdot z .
\]
Hence two indices $\mathfrak{c},\mathfrak{c}'\in\mathbb{Z}^3$ denote the same cell of the torus $n$ if and only if $\mathfrak{c}-\mathfrak{c}'\in\frac{2n}{\mathrm{per}}\mathbb{Z}^3$.
Here the cells of a torus are the orbits of its period lattice on the open cells, as in \cite[Definition~2.18]{QEDII};
there the width is a positive integer, and by (I.h1) of Section~I.6.4 that definition is applied with the width $\mathrm{hc}\in\frac12\mathbb{Z}$.
\item[(d)] The open cells $Q^{\circ}_{\mathfrak{c}}$, $\mathfrak{c}\in\mathbb{Z}^3$, are pairwise disjoint.
\end{itemize}
\end{named}

\begin{proof}
Recall from Section~I.6.4 the description of $\mathsf{g}_{\mathfrak{c}}$ in each coordinate, and that $\mathrm{per}=2\,\mathrm{hc}$.
If $\mathfrak{c}_\mu=2i$, then
\[
(\mathsf{g}_{\mathfrak{c}}y)_\mu=i\cdot\mathrm{per}+y_\mu=2i\cdot\mathrm{hc}+y_\mu=\left(2i+\frac12\right)\cdot\mathrm{hc}+\left(y_\mu-\frac{\mathrm{hc}}{2}\right);
\]
if $\mathfrak{c}_\mu=2i+1$, then
\[
(\mathsf{g}_{\mathfrak{c}}y)_\mu=(i+1)\cdot\mathrm{per}-y_\mu=(2i+2)\cdot\mathrm{hc}-y_\mu=\left(2i+1+\frac12\right)\cdot\mathrm{hc}-\left(y_\mu-\frac{\mathrm{hc}}{2}\right).
\]
Together these two cases are (1.2).
\begin{itemize}
\item[(a)] The two formulas of Section~I.6.4 show that $(r_{\mathfrak{c}})_{\mu\mu}=(-1)^{\mathfrak{c}_\mu}$ and that $(a_{\mathfrak{c}})_\mu$ equals $i\cdot\mathrm{per}$, respectively $(i+1)\cdot\mathrm{per}$.
The reflection in the wall $\{y_\mu=j\cdot\mathrm{hc}\}$ is $\theta_{y_\mu,\mathrm{po}}$ with $\mathrm{po}=j\cdot\mathrm{hc}$ (Section~I.4.8);
it sends $y_\mu$ to $2j\cdot\mathrm{hc}-y_\mu=j\cdot\mathrm{per}-y_\mu$ and leaves the other two coordinates fixed.
If $\mathfrak{c}_\mu=2i+1$, the map $y_\mu\mapsto(i+1)\cdot\mathrm{per}-y_\mu$ is the reflection in the wall with $j=i+1$.
If $\mathfrak{c}_\mu=2i$, the map $y_\mu\mapsto i\cdot\mathrm{per}+y_\mu$ is the reflection in the wall with $j=0$ followed by the reflection in the wall with $j=i$:
$y_\mu\mapsto-y_\mu\mapsto i\cdot\mathrm{per}+y_\mu$.
The map $\mathsf{g}_{\mathfrak{c}}$ is the product of these maps for $\mu=1,2,3$.
\item[(b)] The identity is read off from (1.2).
The map $\mathsf{g}_{\mathfrak{c}}$ is a bijection of $\mathbb{R}^3$, and
\begin{gather*}
Q^{\circ}_{\mathfrak{c}}=\left\{y:\left|y_\mu-\left(\mathfrak{c}_\mu+\frac12\right)\cdot\mathrm{hc}\right|<\frac{\mathrm{hc}}{2},\ \mu=1,2,3\right\}\quad\text{(Section~I.6.4)},\\
\mathrm{Bx}_0=\left\{y:\left|y_\mu-\frac{\mathrm{hc}}{2}\right|\le r_0,\ \mu=1,2,3\right\}.
\end{gather*}
\item[(d)] Let $\mathfrak{c}\neq\mathfrak{c}'$.
There is a $\mu$ with $\mathfrak{c}_\mu\neq\mathfrak{c}'_\mu$, say with $\mathfrak{c}_\mu+1\le\mathfrak{c}'_\mu$.
Then $y_\mu<(\mathfrak{c}_\mu+1)\cdot\mathrm{hc}\le\mathfrak{c}'_\mu\cdot\mathrm{hc}$ on $Q^{\circ}_{\mathfrak{c}}$, whereas $y_\mu>\mathfrak{c}'_\mu\cdot\mathrm{hc}$ on $Q^{\circ}_{\mathfrak{c}'}$.
\item[(c)] Since $\mathrm{per}$ divides $n$, the number $2n/\mathrm{per}=2\cdot(n/\mathrm{per})$ is even;
hence $(-1)^{\mathfrak{c}_\mu+\frac{2n}{\mathrm{per}}z_\mu}=(-1)^{\mathfrak{c}_\mu}$, and moreover $(2n/\mathrm{per})\cdot\mathrm{hc}=n$.
Therefore the right side of (1.2) for $\mathfrak{c}+(2n/\mathrm{per})\cdot z$ equals the right side of (1.2) for $\mathfrak{c}$ plus $n\cdot z_\mu$,
and the inequalities defining $Q^{\circ}_{\mathfrak{c}+(2n/\mathrm{per})\cdot z}$ are the inequalities defining $Q^{\circ}_{\mathfrak{c}}$, written for the point $y-n\cdot z$.
So the orbit of the cell $Q^{\circ}_{\mathfrak{c}}$ under $n\mathbb{Z}^3$ consists of the cells $Q^{\circ}_{\mathfrak{c}}+n\cdot z=Q^{\circ}_{\mathfrak{c}+(2n/\mathrm{per})\cdot z}$.
The cells are non-empty;
by (d), therefore, $Q^{\circ}_{\mathfrak{c}'}$ is one of them if and only if $\mathfrak{c}'=\mathfrak{c}+(2n/\mathrm{per})\cdot z$ for some $z\in\mathbb{Z}^3$.
\end{itemize}
\end{proof}

\subsubsection*{2. The lattice probe without periodisation}

For the placement $v$ of Proposition~41 let $\beta^{v}:=\tau_v\beta$, that is, $\beta^{v}(y)=\beta(y-v)$ (\cite[Section~7.1.3]{QEDI}; Section~I.7.5).
This form is real and smooth, $\beta^{v}_{\nu\mu}=-\beta^{v}_{\mu\nu}$, and $\operatorname{supp}\beta^{v}=v+\operatorname{supp}\beta$ is compact.
With the difference quotients $\delta_{\eta,\mu}$ of Proposition~47.1 and the average $J_\eta$ of Section~I.7.6 define, in physical coordinates,
\begin{align}
a'_\nu&:=-\sum_{\mu\neq\nu}\delta_{\eta,\mu}\beta^{v}_{\mu\nu},\qquad\tilde{a}_{\eta,\nu}:=J_\eta a'_\nu ,\tag{2.1}\\
u^{\mathrm{o}}_{\mu\nu}(c)&:=\left(J_\eta\beta^{v}_{\mu\nu}\right)\left(c+\frac{\eta}{2}\cdot(e_\mu+e_\nu)\right)\qquad\left(c\in\eta\mathbb{Z}^3,\ \mu<\nu\right),\tag{2.2}\\
a^{\mathrm{o}}_{\eta}&:=d_\eta^{*}u^{\mathrm{o}},\qquad\hat{\mathfrak{a}}^{\mathrm{o}}:=\eta^{\frac12}\cdot a^{\mathrm{o}}_{\eta};\qquad\text{in lattice coordinates}\quad\hat{\mathfrak{a}}^{\mathrm{o}}_{\nu}(c)=\eta^{\frac12}\cdot a^{\mathrm{o}}_{\eta,\nu}(\eta\cdot c)\quad\left(c\in\mathbb{Z}^3\right).\tag{2.3}
\end{align}
The functions (2.1) are the functions $a'$, $\tilde{a}_\eta$ of Proposition~47.1 for the probe $\beta^{v}$:
the sum in Proposition~47.1 contains in addition the term $\mu=\nu$, and this term vanishes, because $\beta^{v}_{\nu\nu}=-\beta^{v}_{\nu\nu}=0$.
The cochains (2.2) and (2.3) are those of the hypotheses of Proposition~41,
namely the cochains of Section~I.7.6 for the form $\tau_v\beta$, formed without the periodisation of $u$.
Here $d_\eta=\eta^{-1}\cdot d$, and $d_\eta^{*}$ is the adjoint of $d_\eta$ for the pairing $\langle u,u'\rangle_\eta=\eta^3\cdot\langle u,u'\rangle$ (Section~I.7.5; \cite[Section~7.1.1]{QEDI}).
For $A\in C^1_{\mathrm{c}}$ and $u\in C^2_{\mathrm{c}}$, Lemma~1 gives
\[
\langle d_\eta A,u\rangle_\eta=\eta^2\cdot\langle dA,u\rangle=\eta^2\cdot\langle A,d^{*}u\rangle=\langle A,\eta^{-1}\cdot d^{*}u\rangle_\eta ;
\]
hence $d_\eta^{*}=\eta^{-1}\cdot d^{*}$ on $C^2_{\mathrm{c}}$.
On the torus $n$ one has $\langle u,u'\rangle_\eta=\eta^3\cdot\langle u,u'\rangle_n$ (Section~I.7.5), and for $A\in C^1(n)$, $u\in C^2(n)$,
\[
\langle d_\eta A,u\rangle_\eta=\eta^2\cdot\langle dA,u\rangle_n=\eta^2\cdot\langle A,d^{*}_{n}u\rangle_n=\langle A,\eta^{-1}\cdot d^{*}_{n}u\rangle_\eta ;
\]
hence $d_\eta^{*}=\eta^{-1}\cdot d^{*}_{n}$ on $C^2(n)$.

\begin{named}{Lemma 5}
\begin{itemize}
\item[(a)] $u^{\mathrm{o}}\in C^2_{\mathrm{c}}$ and $\hat{\mathfrak{a}}^{\mathrm{o}}\in C^1_{\mathrm{c}}$.
\item[(b)] $a^{\mathrm{o}}_{\eta,\nu}(c)=\tilde{a}_{\eta,\nu}\left(c+\frac{\eta}{2}\cdot e_\nu\right)$ for $c\in\eta\mathbb{Z}^3$ and $\nu=1,2,3$.
\item[(c)] If $\tilde{a}_{\eta,\nu}(y)\neq0$, then $y\in\operatorname{supp}\beta^{v}+[-\eta,\eta]^3$.
For the probe $\beta^{v}$ this is the assertion on $\tilde{a}_\eta$ in the first sentence of Proposition~47.1\,(b2).
\item[(d)] On the torus $n$ the cochain $\hat{\mathfrak{a}}$ of Section~I.7.6, formed with $\tau_v\beta$, equals $\mathrm{Prd}_n\hat{\mathfrak{a}}^{\mathrm{o}}$, and $d\hat{\mathfrak{a}}=\mathrm{Prd}_n\,d\hat{\mathfrak{a}}^{\mathrm{o}}$.
\item[(e)] Let $v'=v+\eta\cdot z$ with $z\in\mathbb{Z}^3$.
The cochain (2.3) formed with $v'$ in place of $v$ is $(\tau_z)_{*}\hat{\mathfrak{a}}^{\mathrm{o}}$.
\end{itemize}
\end{named}

\begin{proof}
\begin{itemize}
\item[(a)] For a function $\mathrm{fs}$ on $\mathbb{R}^3$ the value $(J_\eta\mathrm{fs})(y)$ is $\eta^{-3}$ times the integral of $\mathrm{fs}$ over $y+[-\eta/2,\eta/2]^3$.
Hence $u^{\mathrm{o}}_{\mu\nu}(c)\neq0$ only if $c+\frac{\eta}{2}\cdot(e_\mu+e_\nu)$ lies in the bounded set $\operatorname{supp}\beta^{v}+[-\eta/2,\eta/2]^3$.
Moreover $d^{*}$ maps $C^2_{\mathrm{c}}$ into $C^1_{\mathrm{c}}$ (Lemma~1).
\item[(b)] From $d_\eta^{*}=\eta^{-1}\cdot d^{*}$ and from (1.1), written in physical coordinates,
\[
a^{\mathrm{o}}_{\eta,\nu}(c)=-\eta^{-1}\cdot\sum_{\mu\neq\nu}\left[u^{\mathrm{o}}_{\mu\nu}(c)-u^{\mathrm{o}}_{\mu\nu}(c-\eta\cdot e_\mu)\right],\qquad u^{\mathrm{o}}_{\nu\mu}:=-u^{\mathrm{o}}_{\mu\nu}.
\]
Formula (2.2) holds for $\mu>\nu$ as well, since $\beta^{v}_{\mu\nu}=-\beta^{v}_{\nu\mu}$ and the point $c+\frac{\eta}{2}\cdot(e_\mu+e_\nu)$ is symmetric in $(\mu,\nu)$.
Let $\mathrm{ce}:=c+\frac{\eta}{2}\cdot e_\nu$ be the centre of the bond $(c,c+\eta\cdot e_\nu)$.
Then $c+\frac{\eta}{2}\cdot(e_\mu+e_\nu)=\mathrm{ce}+\frac{\eta}{2}\cdot e_\mu$ and $(c-\eta\cdot e_\mu)+\frac{\eta}{2}\cdot(e_\mu+e_\nu)=\mathrm{ce}-\frac{\eta}{2}\cdot e_\mu$, so that
\[
u^{\mathrm{o}}_{\mu\nu}(c)-u^{\mathrm{o}}_{\mu\nu}(c-\eta\cdot e_\mu)=\eta\cdot\left(\delta_{\eta,\mu}J_\eta\beta^{v}_{\mu\nu}\right)(\mathrm{ce}).
\]
For a continuous function $\mathrm{fs}$,
\begin{multline*}
\left(J_\eta\delta_{\eta,\mu}\mathrm{fs}\right)(y)=\eta^{-3}\int_{[-\eta/2,\eta/2]^3}\mathrm{d}y'\,\eta^{-1}\cdot\left[\mathrm{fs}\left(y+y'+\frac{\eta}{2}\cdot e_\mu\right)-\mathrm{fs}\left(y+y'-\frac{\eta}{2}\cdot e_\mu\right)\right]\\
=\eta^{-1}\cdot\left[(J_\eta\mathrm{fs})\left(y+\frac{\eta}{2}\cdot e_\mu\right)-(J_\eta\mathrm{fs})\left(y-\frac{\eta}{2}\cdot e_\mu\right)\right]=\left(\delta_{\eta,\mu}J_\eta\mathrm{fs}\right)(y).
\end{multline*}
Therefore
\[
a^{\mathrm{o}}_{\eta,\nu}(c)=-\sum_{\mu\neq\nu}\left(J_\eta\delta_{\eta,\mu}\beta^{v}_{\mu\nu}\right)(\mathrm{ce})=\left(J_\eta a'_\nu\right)(\mathrm{ce})=\tilde{a}_{\eta,\nu}(\mathrm{ce}).
\]
\item[(c)] By the definitions,
\[
\tilde{a}_{\eta,\nu}(y)=-\sum_{\mu\neq\nu}\eta^{-4}\int_{[-\eta/2,\eta/2]^3}\mathrm{d}y'\,\left[\beta^{v}_{\mu\nu}\left(y+y'+\frac{\eta}{2}\cdot e_\mu\right)-\beta^{v}_{\mu\nu}\left(y+y'-\frac{\eta}{2}\cdot e_\mu\right)\right].
\]
If this number is different from zero, there are a $\mu\neq\nu$, a sign $\pm$ and a point $y'\in[-\eta/2,\eta/2]^3$ such that $y+y'\pm\frac{\eta}{2}\cdot e_\mu\in\operatorname{supp}\beta^{v}$;
and $\left|y'\pm\frac{\eta}{2}\cdot e_\mu\right|\le\eta$.
\item[(d)] In Section~I.7.6 the two-cochain $u$ of the torus with $n$ sites per side (of period $n\cdot\eta$ in physical coordinates) is the periodisation of (2.2), that is, $u=\mathrm{Prd}_n u^{\mathrm{o}}$.
By $d_\eta^{*}=\eta^{-1}\cdot d^{*}_{n}$ on $C^2(n)$, by Lemma~2\,(c), and by $d_\eta^{*}=\eta^{-1}\cdot d^{*}$ on $C^2_{\mathrm{c}}$,
\[
a_\eta=d_\eta^{*}u=\eta^{-1}\cdot d^{*}_{n}\,\mathrm{Prd}_n u^{\mathrm{o}}=\mathrm{Prd}_n\left(\eta^{-1}\cdot d^{*}u^{\mathrm{o}}\right)=\mathrm{Prd}_n a^{\mathrm{o}}_{\eta}.
\]
Hence $\hat{\mathfrak{a}}=\eta^{\frac12}\cdot a_\eta=\mathrm{Prd}_n\hat{\mathfrak{a}}^{\mathrm{o}}$, and Lemma~2\,(b) gives $d\hat{\mathfrak{a}}=\mathrm{Prd}_n\,d\hat{\mathfrak{a}}^{\mathrm{o}}$.
\item[(e)] One has $\beta^{v'}(y)=\beta(y-v-\eta\cdot z)=\beta^{v}(y-\eta\cdot z)$.
The operators $\delta_{\eta,\mu}$ and $J_\eta$ are formed with translates of their argument, so they commute with the translation by $\eta\cdot z$;
hence the function (2.1) formed with $v'$ is $y\mapsto\tilde{a}_{\eta,\nu}(y-\eta\cdot z)$.
By (b) and (2.3), the cochain formed with $v'$ takes, in lattice coordinates, on the bond $(c,c+e_\nu)$ the value
\[
\eta^{\frac12}\cdot\tilde{a}_{\eta,\nu}\left(\eta\cdot(c-z)+\frac{\eta}{2}\cdot e_\nu\right)=\hat{\mathfrak{a}}^{\mathrm{o}}_{\nu}(c-z).
\]
By (0.2) this is the cochain $(\tau_z)_{*}\hat{\mathfrak{a}}^{\mathrm{o}}$.
\end{itemize}
\end{proof}

\subsubsection*{3. Proof of (a)}

\begin{proof}[Proof of Proposition~41\,(a)]
\emph{The cube $v+Q_\beta$.}
Let $y\in v+Q_\beta=v+y_\beta+\left[-\frac{D_\beta}{2},\frac{D_\beta}{2}\right]^3$ and $\mu\in\{1,2,3\}$.
The inequality $\left|v+y_\beta-\frac{\ell_c}{2}(1,1,1)\right|\le\frac{\eta}{2}$ that defines a placement, together with $D_\beta\le\ell_c-3\eta$, gives
\[
\left|y_\mu-\frac{\ell_c}{2}\right|
\le\left|y_\mu-v_\mu-(y_\beta)_\mu\right|+\left|v_\mu+(y_\beta)_\mu-\frac{\ell_c}{2}\right|
\le\frac{D_\beta}{2}+\frac{\eta}{2}
\le\frac{\ell_c-3\eta}{2}+\frac{\eta}{2}
=\frac{\ell_c}{2}-\eta .
\]
Hence $\eta\le y_\mu\le\ell_c-\eta$, and
\begin{equation}
v+Q_\beta\subset[\eta,\ell_c-\eta]^3 .
\tag{3.1}
\end{equation}

\emph{The bound $m_{\mathrm{p}}\ge7$.}
As $\beta\neq0$, the set $\mathrm{supp}\,\beta$ contains a point $y$.
By Section~I.7.6 the cube $y+[-2,2]^3$, of side $4$, is contained in $Q_\beta$; therefore $D_\beta\ge4$.
It follows that
\[
\left(m_{\mathrm{p}}+\frac12\right)\cdot\eta=\ell_c\ge D_\beta+3\eta\ge4+3\eta ,
\]
that is, $m_{\mathrm{p}}+\frac12\ge\frac{4}{\eta}+3\ge7$, the last step because $\eta\le1$.
Thus $m_{\mathrm{p}}\ge\frac{13}{2}$, and since $m_{\mathrm{p}}$ is an integer, $m_{\mathrm{p}}\ge7$.

\emph{Finite support; independence of $n$.}
By Lemma~5\,(a), $\hat{\mathfrak{a}}^{\mathrm{o}}$ lies in $C^{1}_{\mathrm{c}}$.
The definitions (2.2) and (2.3) involve $\beta$, $v$ and $\eta$, and they do not involve $n$.

\emph{End points of the bonds.}
Let $y\in\mathrm{supp}\,\beta^{v}=v+\mathrm{supp}\,\beta$.
Then $y-v\in\mathrm{supp}\,\beta$, so that $y-v+[-2,2]^3\subset Q_\beta$ by Section~I.7.6, and consequently $y+[-2,2]^3\subset v+Q_\beta$.
By (3.1) one has $\eta\le y_\mu-2$ and $y_\mu+2\le\ell_c-\eta$ for $\mu=1,2,3$.
Therefore
\begin{equation}
\mathrm{supp}\,\beta^{v}\subset[2+\eta,\ell_c-2-\eta]^3 .
\tag{3.2}
\end{equation}
Now let $(c,c+\eta e_\nu)$ be a bond of $\eta\mathbb{Z}^3$ with $a^{\mathrm{o}}_{\eta,\nu}(c)\neq0$, and put $\mathrm{ce}:=c+\frac{\eta}{2}e_\nu$.
Lemma~5\,(b), (c) provides a point $y\in\mathrm{supp}\,\beta^{v}$ with $|\mathrm{ce}-y|\le\eta$.
The end points of the bond are $\mathrm{ce}-\frac{\eta}{2}e_\nu$ and $\mathrm{ce}+\frac{\eta}{2}e_\nu$, and the sup distance of each of them from $y$ is at most $\frac{3\eta}{2}$.
By (3.2), their physical coordinates lie in
\[
\left[2+\eta-\frac{3\eta}{2},\ \ell_c-2-\eta+\frac{3\eta}{2}\right]=\left[2-\frac{\eta}{2},\ \ell_c-2+\frac{\eta}{2}\right].
\]
Dividing by $\eta$ and using $\ell_c/\eta=\mathrm{hc}=m_{\mathrm{p}}+\frac12$, one finds that the lattice coordinates of both end points lie in
\begin{equation}
\left[\frac{2}{\eta}-\frac12,\ \mathrm{hc}-\frac{2}{\eta}+\frac12\right]=\left[\frac{2}{\eta}-\frac12,\ m_{\mathrm{p}}+1-\frac{2}{\eta}\right].
\tag{3.3}
\end{equation}
Since $\frac{2}{\eta}\ge2$, this interval is contained in $\left[\frac32,m_{\mathrm{p}}-1\right]$.
Lattice coordinates are integers; so they belong to $\{2,\dots,m_{\mathrm{p}}-1\}$.
\end{proof}

\subsubsection*{4. Proof of (b)}

\begin{proof}[Proof of Proposition~41\,(b)]
\emph{Positivity of $r_0$.}
The inequality $D_\beta+3\eta\le\ell_c$ gives $\mathrm{hc}=\ell_c/\eta\ge D_\beta/\eta+3$, and $D_\beta\ge4$ was shown in Section~3.
Consequently
\begin{equation}
r_0=\frac{\mathrm{hc}}{2}-\frac{2}{\eta}+\frac32
\ge\frac{D_\beta}{2\eta}+\frac32-\frac{2}{\eta}+\frac32
=\frac{D_\beta-4}{2\eta}+3\ge3 .
\tag{4.1}
\end{equation}

\emph{The box.}
One has $\frac{\mathrm{hc}}{2}-r_0=\frac{1}{\eta}\left(2-\frac{3\eta}{2}\right)=\frac{2}{\eta}-\frac32$ and $\frac{\mathrm{hc}}{2}+r_0=\mathrm{hc}-\frac{2}{\eta}+\frac32=m_{\mathrm{p}}+2-\frac{2}{\eta}$.
Hence
\begin{equation}
\mathrm{Bx}_0=\left[\frac{2}{\eta}-\frac32,\ m_{\mathrm{p}}+2-\frac{2}{\eta}\right]^3 .
\tag{4.2}
\end{equation}

\emph{Plaquettes.}
Let $p=(c;\mu\nu)$, $\mu<\nu$, be a plaquette with $(d\hat{\mathfrak{a}}^{\mathrm{o}})(p)\neq0$, and let $\mathrm{vx}_0,\dots,\mathrm{vx}_3$ be its corners in the order of Subsection~0.2.
By (0.1) there is an index $i$ with $\hat{\mathfrak{a}}^{\mathrm{o}}(\mathrm{vx}_i,\mathrm{vx}_{i+1})\neq0$.
The list of the corners in Subsection~0.2 shows that $\mathrm{vx}_{i+2}-\mathrm{vx}_{i+1}=\mathrm{vx}_{i+3}-\mathrm{vx}_i$, and that this vector is $+e_{\mathfrak{l}}$ or $-e_{\mathfrak{l}}$, where $\mathfrak{l}$ is the direction of $p$ other than the direction of the bond $\{\mathrm{vx}_i,\mathrm{vx}_{i+1}\}$:
it equals $e_\nu$ for $i=0$, $-e_\mu$ for $i=1$, $-e_\nu$ for $i=2$, and $e_\mu$ for $i=3$.
By (3.3), the coordinates of $\mathrm{vx}_i$ and of $\mathrm{vx}_{i+1}$ lie in $\left[\frac{2}{\eta}-\frac12,m_{\mathrm{p}}+1-\frac{2}{\eta}\right]$.
So the coordinates of all four corners lie in $\left[\frac{2}{\eta}-\frac32,m_{\mathrm{p}}+2-\frac{2}{\eta}\right]$, and by (4.2) the four corners belong to $\mathrm{Bx}_0$.
The box $\mathrm{Bx}_0$ is convex, and $p$ is the convex hull of its corners; hence $p\subset\mathrm{Bx}_0$.
As $\frac{2}{\eta}\ge2$, the coordinates of the corners lie in $\left[\frac12,m_{\mathrm{p}}\right]$; they are integers, so the corners lie in $\{1,\dots,m_{\mathrm{p}}\}^3$.

\emph{The distances from the walls; the inclusion $\mathrm{Bx}_0\subset Q^{\circ}_{0}$.}
Let $y\in\mathrm{Bx}_0$, $\mu\in\{1,2,3\}$ and $j\in\mathbb{Z}$.
If $j\le0$, then (4.2) gives $y_\mu-j\cdot\mathrm{hc}\ge\frac{2}{\eta}-\frac32+(-j)\cdot\mathrm{hc}$, with equality for $y_\mu=\frac{\mathrm{hc}}{2}-r_0$.
If $j\ge1$, then (4.2) gives
\[
j\cdot\mathrm{hc}-y_\mu\ge j\cdot\mathrm{hc}-\left(\mathrm{hc}-\frac{2}{\eta}+\frac32\right)=\frac{2}{\eta}-\frac32+(j-1)\cdot\mathrm{hc},
\]
with equality for $y_\mu=\frac{\mathrm{hc}}{2}+r_0$.
The distance of a point $y$ from the plane $\{y_\mu=j\cdot\mathrm{hc}\}$ is $|y_\mu-j\cdot\mathrm{hc}|$.
Therefore
\begin{equation}
\mathrm{dist}\left(\mathrm{Bx}_0,\{y_\mu=j\cdot\mathrm{hc}\}\right)=\frac{2}{\eta}-\frac32+\max(-j,j-1)\cdot\mathrm{hc}\qquad(j\in\mathbb{Z},\ \mu=1,2,3).
\tag{4.3}
\end{equation}
For $j\in\{0,1\}$, that is, for the six walls that bound $Q^{\circ}_{0}$, this distance equals $\frac{2}{\eta}-\frac32$; for every other wall it is larger.
Moreover $\frac{2}{\eta}-\frac32\ge\frac12$, because $\eta\le1$.
In particular $0<y_\mu<\mathrm{hc}$ for every $y\in\mathrm{Bx}_0$ and every $\mu$, so that $\mathrm{Bx}_0\subset Q^{\circ}_{0}=(0,\mathrm{hc})^3$.
\end{proof}

\subsubsection*{5. Proof of (c)}

\begin{proof}[Proof of Proposition~41\,(c)]
For $\mathfrak{c}\in\mathbb{Z}^3$ put
\begin{equation}
\hat{\mathfrak{a}}^{\mathrm{o}}_{\mathfrak{c}}:=(\mathsf{g}_{\mathfrak{c}})_{*}\hat{\mathfrak{a}}^{\mathrm{o}}\in C^{1}_{\mathrm{c}},\qquad
\mathrm{Bx}_{\mathfrak{c}}:=\mathsf{g}_{\mathfrak{c}}\mathrm{Bx}_0=\left(\left(\mathfrak{c}_\mu+\frac12\right)\cdot\mathrm{hc}\right)_\mu+[-r_0,r_0]^3 .
\tag{5.1}
\end{equation}
The first cochain lies in $C^{1}_{\mathrm{c}}$ by Lemma~4\,(a) and Lemma~3\,(b), (d); the second equality is Lemma~4\,(b).
Since $r_0>0$ by (4.1), $\mathrm{Bx}_{\mathfrak{c}}$ is a box in the sense of \cite[Definition~2.12]{QEDII}.
Clause (b) gives $\mathrm{Bx}_0\subset Q^{\circ}_{0}$, and $\mathsf{g}_{\mathfrak{c}}Q^{\circ}_{0}=Q^{\circ}_{\mathfrak{c}}$ by Lemma~4\,(b); hence $\mathrm{Bx}_{\mathfrak{c}}\subset Q^{\circ}_{\mathfrak{c}}$.

\emph{The cochains $\hat{\mathfrak{a}}_{\mathfrak{c}}$ as periodisations.}
By Section~I.7.6, $\hat{\mathfrak{a}}_{\mathfrak{c}}=(\mathsf{g}_{\mathfrak{c}})_{*}\hat{\mathfrak{a}}$, and $(\mathsf{g}_{\mathfrak{c}})_{*}$ maps $C^{1}(n)$ into itself by Lemma~3\,(c), (d).
From Lemma~5\,(d), Lemma~3\,(c), (d), Lemma~2\,(b) and Lemma~3\,(a), (d) one obtains
\[
\hat{\mathfrak{a}}_{\mathfrak{c}}=(\mathsf{g}_{\mathfrak{c}})_{*}\mathrm{Prd}_n\hat{\mathfrak{a}}^{\mathrm{o}}=\mathrm{Prd}_n\hat{\mathfrak{a}}^{\mathrm{o}}_{\mathfrak{c}},\qquad
d\hat{\mathfrak{a}}_{\mathfrak{c}}=\mathrm{Prd}_n\,d\hat{\mathfrak{a}}^{\mathrm{o}}_{\mathfrak{c}},\qquad
d\hat{\mathfrak{a}}^{\mathrm{o}}_{\mathfrak{c}}=(\mathsf{g}_{\mathfrak{c}})_{*}\,d\hat{\mathfrak{a}}^{\mathrm{o}}.
\]

\emph{The supports lie in the boxes; uniqueness of $z$.}
Lemma~3\,(e) gives $\mathrm{supp}\,d\hat{\mathfrak{a}}^{\mathrm{o}}_{\mathfrak{c}}=\{\mathsf{g}_{\mathfrak{c}}p':\ p'\in\mathrm{supp}\,d\hat{\mathfrak{a}}^{\mathrm{o}}\}$.
By clause (b), every such $p'$ is contained in $\mathrm{Bx}_0$; hence every plaquette of $\mathrm{supp}\,d\hat{\mathfrak{a}}^{\mathrm{o}}_{\mathfrak{c}}$ is contained in $\mathrm{Bx}_{\mathfrak{c}}$.
Let $p$ be a plaquette of $\mathbb{Z}^3$ with
\[
(d\hat{\mathfrak{a}}_{\mathfrak{c}})(p)=\sum_{z\in\mathbb{Z}^3}(d\hat{\mathfrak{a}}^{\mathrm{o}}_{\mathfrak{c}})(p+nz)\neq0 .
\]
Then $p+nz\subset\mathrm{Bx}_{\mathfrak{c}}$ for some $z\in\mathbb{Z}^3$.
By (5.1) and $\frac{2n}{\mathrm{per}}\cdot\mathrm{hc}=n$, this means
\begin{equation}
p\subset\mathrm{Bx}_{\mathfrak{c}}-nz=\mathrm{Bx}_{\mathfrak{c}-\frac{2n}{\mathrm{per}}z},\qquad\text{a box in}\qquad Q^{\circ}_{\mathfrak{c}-\frac{2n}{\mathrm{per}}z}=Q^{\circ}_{\mathfrak{c}}-nz .
\tag{5.2}
\end{equation}
By Lemma~4\,(d) only one $z$ has this property: for two distinct such vectors the non-empty set $p$ would lie in two open cells with distinct indices.
By Lemma~4\,(c), the index $\mathfrak{c}-\frac{2n}{\mathrm{per}}z$ denotes the same cell of the torus $n$ as $\mathfrak{c}$.

\emph{The pairings.}
Let $\mathfrak{c},\mathfrak{c}'\in\mathbb{Z}^3$.
Lemma~2\,(a), applied to $u=d\hat{\mathfrak{a}}^{\mathrm{o}}_{\mathfrak{c}}\in C^{2}_{\mathrm{c}}$ and $u'=\mathrm{Prd}_n\,d\hat{\mathfrak{a}}^{\mathrm{o}}_{\mathfrak{c}'}\in C^{2}(n)$, yields
\begin{equation}
\begin{split}
\left\langle d\hat{\mathfrak{a}}_{\mathfrak{c}},d\hat{\mathfrak{a}}_{\mathfrak{c}'}\right\rangle_n
&=\left\langle\mathrm{Prd}_n\,d\hat{\mathfrak{a}}^{\mathrm{o}}_{\mathfrak{c}},\mathrm{Prd}_n\,d\hat{\mathfrak{a}}^{\mathrm{o}}_{\mathfrak{c}'}\right\rangle_n
=\left\langle d\hat{\mathfrak{a}}^{\mathrm{o}}_{\mathfrak{c}},\mathrm{Prd}_n\,d\hat{\mathfrak{a}}^{\mathrm{o}}_{\mathfrak{c}'}\right\rangle\\
&=\sum_p\sum_{z\in\mathbb{Z}^3}(d\hat{\mathfrak{a}}^{\mathrm{o}}_{\mathfrak{c}})(p)\cdot(d\hat{\mathfrak{a}}^{\mathrm{o}}_{\mathfrak{c}'})(p+nz).
\end{split}
\tag{5.3}
\end{equation}
Suppose that the term with the index $(p,z)$ is different from $0$.
Then $p\subset\mathrm{Bx}_{\mathfrak{c}}\subset Q^{\circ}_{\mathfrak{c}}$ and $p+nz\subset\mathrm{Bx}_{\mathfrak{c}'}\subset Q^{\circ}_{\mathfrak{c}'}$.
Thus $p$ is contained both in $Q^{\circ}_{\mathfrak{c}}$ and in $Q^{\circ}_{\mathfrak{c}'}-nz=Q^{\circ}_{\mathfrak{c}'-\frac{2n}{\mathrm{per}}z}$ (Lemma~4\,(c)), and Lemma~4\,(d) gives $\mathfrak{c}=\mathfrak{c}'-\frac{2n}{\mathrm{per}}z$.

\emph{Different cells of the torus.}
Let $\mathfrak{c}$ and $\mathfrak{c}'$ denote different cells of the torus $n$, that is (Lemma~4\,(c)), $\mathfrak{c}-\mathfrak{c}'\notin\frac{2n}{\mathrm{per}}\mathbb{Z}^3$.
Then every term of (5.3) vanishes, and $\left\langle d\hat{\mathfrak{a}}_{\mathfrak{c}},d\hat{\mathfrak{a}}_{\mathfrak{c}'}\right\rangle_n=0$.

\emph{The same cell of the torus.}
Let $\mathfrak{c}'=\mathfrak{c}+\frac{2n}{\mathrm{per}}z'$ with $z'\in\mathbb{Z}^3$.
Then only the terms of (5.3) with $z=z'$ can be different from $0$.
Lemma~4\,(c) and Subsection~0.4 give $\mathsf{g}_{\mathfrak{c}'}=\tau_{nz'}\mathsf{g}_{\mathfrak{c}}$ and $\hat{\mathfrak{a}}^{\mathrm{o}}_{\mathfrak{c}'}=(\tau_{nz'})_{*}\hat{\mathfrak{a}}^{\mathrm{o}}_{\mathfrak{c}}$; by Lemma~3\,(a), (d) and (0.2),
\[
(d\hat{\mathfrak{a}}^{\mathrm{o}}_{\mathfrak{c}'})(p+nz')=\left((\tau_{nz'})_{*}\,d\hat{\mathfrak{a}}^{\mathrm{o}}_{\mathfrak{c}}\right)(p+nz')=(d\hat{\mathfrak{a}}^{\mathrm{o}}_{\mathfrak{c}})(p).
\]
Hence, by Lemma~3\,(b), (d),
\begin{equation}
\left\langle d\hat{\mathfrak{a}}_{\mathfrak{c}},d\hat{\mathfrak{a}}_{\mathfrak{c}'}\right\rangle_n=\sum_p(d\hat{\mathfrak{a}}^{\mathrm{o}}_{\mathfrak{c}})(p)^2=\left\|(\mathsf{g}_{\mathfrak{c}})_{*}\,d\hat{\mathfrak{a}}^{\mathrm{o}}\right\|^2=\left\|d\hat{\mathfrak{a}}^{\mathrm{o}}\right\|^2 .
\tag{5.4}
\end{equation}

\emph{Dependence on the cell only.}
For indices of the same cell the two cochains coincide: $\hat{\mathfrak{a}}_{\mathfrak{c}'}=(\tau_{nz'})_{*}(\mathsf{g}_{\mathfrak{c}})_{*}\hat{\mathfrak{a}}=\hat{\mathfrak{a}}_{\mathfrak{c}}$ by Lemma~3\,(f).
Thus $\hat{\mathfrak{a}}_{\mathfrak{c}}$ depends only on the cell of the torus, that is (Lemma~4\,(c)), only on the class of $\mathfrak{c}$ modulo $\frac{2n}{\mathrm{per}}\mathbb{Z}^3$.

\emph{The number $\boldsymbol{\sigma}^2$.}
The map $\mathsf{g}_0$ is the identity, so $\hat{\mathfrak{a}}_0=\hat{\mathfrak{a}}$.
For $\mathfrak{c}=\mathfrak{c}'=0$ the identity (5.4) reads $\boldsymbol{\sigma}^2=\|d\hat{\mathfrak{a}}\|_n^2=\|d\hat{\mathfrak{a}}^{\mathrm{o}}\|^2$, where $\boldsymbol{\sigma}^2$ is the number of Section~I.7.6.
With this, (5.4) states that $\left\langle d\hat{\mathfrak{a}}_{\mathfrak{c}},d\hat{\mathfrak{a}}_{\mathfrak{c}'}\right\rangle_n=\boldsymbol{\sigma}^2$ whenever $\mathfrak{c}$ and $\mathfrak{c}'$ are the same cell of the torus $n$.

\emph{Independence.}
The number $\|d\hat{\mathfrak{a}}^{\mathrm{o}}\|^2$ involves neither $n$ nor $\mathfrak{c}$ (clause (a)).
Let $v'$ be a second placement.
Both $v$ and $v'$ lie in $\eta\mathbb{Z}^3$; so $v'=v+\eta z$ with $z\in\mathbb{Z}^3$.
By Lemma~5\,(e) the cochain formed with $v'$ is $(\tau_z)_{*}\hat{\mathfrak{a}}^{\mathrm{o}}$, and Lemma~3\,(a), (b), (d) gives
\[
\left\|d(\tau_z)_{*}\hat{\mathfrak{a}}^{\mathrm{o}}\right\|^2=\left\|(\tau_z)_{*}\,d\hat{\mathfrak{a}}^{\mathrm{o}}\right\|^2=\left\|d\hat{\mathfrak{a}}^{\mathrm{o}}\right\|^2 .
\]
\end{proof}

\subsubsection*{6. Proof of (d)}

\begin{proof}[Proof of Proposition~41\,(d)]
For $\mathfrak{c}\in\mathbb{Z}^3$ and $u\in C^{2}$ put
\begin{equation}
X^{\mathrm{o}}_{\mathfrak{c}}(u):=\left\langle d\hat{\mathfrak{a}}^{\mathrm{o}}_{\mathfrak{c}},u\right\rangle=\sum_{p\subset\mathrm{Bx}_{\mathfrak{c}}}(d\hat{\mathfrak{a}}^{\mathrm{o}}_{\mathfrak{c}})(p)\cdot u(p);
\tag{6.1}
\end{equation}
the second equality holds because every plaquette of $\mathrm{supp}\,d\hat{\mathfrak{a}}^{\mathrm{o}}_{\mathfrak{c}}$ is contained in $\mathrm{Bx}_{\mathfrak{c}}$ (Section~5).
Let $\mathrm{Pol}(z)=\sum_j\mathrm{cf}_j z^j$ be a polynomial with complex coefficients $\mathrm{cf}_j$, where $j$ ranges over the integers $\ge0$ and only finitely many $\mathrm{cf}_j$ are different from $0$.

\textbf{6.1. $\mathrm{Pol}(X^{\mathrm{o}}_{\mathfrak{c}})$ is a bosonic element of $\mathrm{Ad}(\mathrm{Bx}_{\mathfrak{c}})$ and an element of $\mathcal{A}(Q^{\circ}_{\mathfrak{c}})$.}
Let $\mathrm{Ln}(\mathrm{Bx}_{\mathfrak{c}})$ and $\mathrm{Pq}(\mathrm{Bx}_{\mathfrak{c}})$ be the sets of the bonds and of the plaquettes of $\mathbb{Z}^3$ that are contained in $\mathrm{Bx}_{\mathfrak{c}}$;
in \cite[Definition~2.12]{QEDII} they are the sets $B^{(1)}$ and $B^{(2)}$ of the box $B=\mathrm{Bx}_{\mathfrak{c}}$.
The four bonds of a plaquette contained in $\mathrm{Bx}_{\mathfrak{c}}$ are contained in $\mathrm{Bx}_{\mathfrak{c}}$; so (0.1) defines $(d\hat A)(p)$ for $\hat A\in\mathbb{R}^{\mathrm{Ln}(\mathrm{Bx}_{\mathfrak{c}})}$ and $p\in\mathrm{Pq}(\mathrm{Bx}_{\mathfrak{c}})$.
Put
\[
G(\hat A):=\mathrm{Pol}\left(\sum_{p\in\mathrm{Pq}(\mathrm{Bx}_{\mathfrak{c}})}(d\hat{\mathfrak{a}}^{\mathrm{o}}_{\mathfrak{c}})(p)\cdot(d\hat A)(p)\right),\qquad\hat A\in\mathbb{R}^{\mathrm{Ln}(\mathrm{Bx}_{\mathfrak{c}})}.
\]
For a one-cochain $\hat A$ on $\mathbb{Z}^3$ this is $\mathrm{Pol}\left(X^{\mathrm{o}}_{\mathfrak{c}}(d\hat A)\right)$, and it depends only on the restriction of $\hat A$ to $\mathrm{Ln}(\mathrm{Bx}_{\mathfrak{c}})$.
With $B=\mathrm{Bx}_{\mathfrak{c}}$, the set $\mathrm{Ad}(\mathrm{Bx}_{\mathfrak{c}})$ is the set $\mathcal{G}(B)$ of \cite[Definition~2.12]{QEDII}, and $\mathrm{Ad}^{\mathrm{ev}}(\mathrm{Bx}_{\mathfrak{c}})$ is the set $\mathcal{G}^{\mathrm{ev}}(B)$ of the even elements of $\mathcal{G}(B)$ defined there.
By that definition, an element of $\mathrm{Ad}(\mathrm{Bx}_{\mathfrak{c}})$ is subject to the requirements (i), (ii), (iii) below, and (iv) is the definition of a bosonic element.
\begin{itemize}
\item[(i)] \emph{Borel map.}
$G$ is a polynomial in the finitely many coordinates of $\hat A$, and its values are scalars, that is, elements of the part of degree $0$ of the Grassmann algebra of \cite[Definition~2.12]{QEDII}.
It is continuous.
\item[(ii)] \emph{Gauge invariance.}
Let $\boldsymbol{\lambda}$ be a real function on the sites of $\mathrm{Bx}_{\mathfrak{c}}$.
By (0.1),
\[
(dd\boldsymbol{\lambda})(p)=\sum_{i\bmod4}\left(\boldsymbol{\lambda}(\mathrm{vx}_{i+1})-\boldsymbol{\lambda}(\mathrm{vx}_i)\right)=0\qquad\left(p\in\mathrm{Pq}(\mathrm{Bx}_{\mathfrak{c}})\right).
\]
Therefore $d(\hat A-d\boldsymbol{\lambda})=d\hat A$ on $\mathrm{Pq}(\mathrm{Bx}_{\mathfrak{c}})$, and $G(\hat A-d\boldsymbol{\lambda})=G(\hat A)$.
No generator $\psi$, $\bar\psi$ occurs in $G$.
\item[(iii)] \emph{Growth.}
Put
\[
\mathrm{nr}:=\left(\sum_{p\in\mathrm{Pq}(\mathrm{Bx}_{\mathfrak{c}})}(d\hat A)(p)^2\right)^{\frac12}.
\]
The Cauchy--Schwarz inequality, (5.4) and the identity $\boldsymbol{\sigma}^2=\|d\hat{\mathfrak{a}}^{\mathrm{o}}\|^2$ of Section~5 give
\[
\left|\sum_p(d\hat{\mathfrak{a}}^{\mathrm{o}}_{\mathfrak{c}})(p)\cdot(d\hat A)(p)\right|\le\left\|d\hat{\mathfrak{a}}^{\mathrm{o}}_{\mathfrak{c}}\right\|\cdot\mathrm{nr}=(\boldsymbol{\sigma}^2)^{\frac12}\cdot\mathrm{nr}.
\]
For an integer $j\ge1$, the function $\mathrm{nr}\mapsto\mathrm{nr}^j\exp\left(-\frac{\mathrm{nr}^2}{4}\right)$ on $[0,\infty)$ has the logarithmic derivative $\frac{j}{\mathrm{nr}}-\frac{\mathrm{nr}}{2}$.
Its maximum is therefore attained at $\mathrm{nr}^2=2j$, and its value there is $\left(\frac{2j}{\mathrm{e}}\right)^{\frac{j}{2}}$; for $j=0$ the maximum is $1$.
With $\left(\frac{2j\boldsymbol{\sigma}^2}{\mathrm{e}}\right)^{\frac{j}{2}}$ read as $1$ for $j=0$, it follows that
\[
|G(\hat A)|\le\sum_j|\mathrm{cf}_j|\cdot(\boldsymbol{\sigma}^2)^{\frac{j}{2}}\cdot\mathrm{nr}^j\le\mathrm{Cq}\cdot\exp\left(\frac{\mathrm{nr}^2}{4}\right),\qquad
\mathrm{Cq}:=\sum_j|\mathrm{cf}_j|\cdot\left(\frac{2j\boldsymbol{\sigma}^2}{\mathrm{e}}\right)^{\frac{j}{2}}.
\]
This is the growth bound of \cite[Definition~2.12]{QEDII} with the constants $\mathrm{Cq}$ and $\varepsilon_G=\frac14$:
\[
|G(\hat A)|\le\mathrm{Cq}\cdot\exp\left(\left(\frac12-\varepsilon_G\right)\sum_{p\in\mathrm{Pq}(\mathrm{Bx}_{\mathfrak{c}})}(d\hat A)(p)^2\right).
\]
\item[(iv)] \emph{Scalar values.}
Every value $G(\hat A)$ is a scalar.
Hence $G$ is bosonic, and it is even.
\end{itemize}
By (i)--(iv), $G$ lies in $\mathrm{Ad}(\mathrm{Bx}_{\mathfrak{c}})$, is bosonic, and lies in $\mathrm{Ad}^{\mathrm{ev}}(\mathrm{Bx}_{\mathfrak{c}})$.
The box $\mathrm{Bx}_{\mathfrak{c}}$ is contained in the open set $Q^{\circ}_{\mathfrak{c}}$ (Section~5).
The product $\Phi\cdot G_1$ with the bounded factor $\Phi=1$ and the single factor $G_1=G\in\mathrm{Ad}^{\mathrm{ev}}(\mathrm{Bx}_{\mathfrak{c}})$ is one of the products whose linear span is $\mathcal{A}(Q^{\circ}_{\mathfrak{c}})$ (\cite[Definition~2.12]{QEDII}; Section~I.6.4).
Consequently $G\in\mathcal{A}(Q^{\circ}_{\mathfrak{c}})$.
Finally, $G$ is built from $\mathrm{Pol}$, from $\hat{\mathfrak{a}}^{\mathrm{o}}_{\mathfrak{c}}=(\mathsf{g}_{\mathfrak{c}})_{*}\hat{\mathfrak{a}}^{\mathrm{o}}$ and from $\mathrm{Bx}_{\mathfrak{c}}$, and none of these involves $n$ (clause (a); (1.2); (5.1)); so $G$ does not involve $n$.

\textbf{6.2. The embedding in $n\mathbb{Z}^3$; the reading on the torus $n$.}
The side $n$ lies in $\mathrm{per}\,\mathbb{Z}$ and satisfies $n\ge n_*>0$; hence $n$ is a positive multiple of $\mathrm{per}$, and $\mathrm{per}\le n$.
From $r_0=\frac{\mathrm{hc}}{2}-\frac{2}{\eta}+\frac32$, from $\eta\le1$ and from $m_{\mathrm{p}}\ge7$ (Section~3) one obtains
\begin{equation}
2r_0+2=\mathrm{hc}-\frac{4}{\eta}+5\le\mathrm{hc}+1=m_{\mathrm{p}}+\frac32<2m_{\mathrm{p}}+1=\mathrm{per}\le n,\qquad n\ge\mathrm{per}=2m_{\mathrm{p}}+1\ge15\ge4 .
\tag{6.2}
\end{equation}
Since $n\ge4$, the lattice $n\mathbb{Z}^3$ is a period lattice in the sense of \cite[Definition~2.6]{QEDII}, its systole being $\mathrm{sys}(n\mathbb{Z}^3)=n\ge4$.
So $C(n,n,n)$ is a torus in the sense of Section~I.7.1 and of \cite[Definition~2.11]{QEDII}, and it carries $\mathcal{Q}_{*}$ because $\mathrm{per}$ divides $n$ (Section~I.6.4).
The sets $\mathrm{Bx}_{\mathfrak{c}}+[-1,1]^3+nz$, $z\in\mathbb{Z}^3$, are closed cubes of side $2r_0+2<n$, and their centres differ by vectors of $n\mathbb{Z}^3$.
For $z\neq z'$ there is a $\mu$ with $|nz_\mu-nz'_\mu|\ge n>2r_0+2$; two closed intervals of length $2r_0+2$ whose centres are more than $2r_0+2$ apart are disjoint; hence these cubes are pairwise disjoint.
Thus $\mathrm{Bx}_{\mathfrak{c}}$ embeds in $n\mathbb{Z}^3$, the sets $\mathrm{Bx}_{\mathfrak{c}}+nz$ are pairwise disjoint, and the product $1\cdot G$ of Subsection~6.1 embeds in $n\mathbb{Z}^3$ (\cite[Definition~2.12]{QEDII}).
Let $(A,\mathfrak{z})$ be a configuration on the torus $n$, let $F^{\mathrm{u}}=dA+r_q\cdot\mathfrak{z}\in C^{2}(n)$ (Section~I.7.1), and let $\hat A$ be a local lift on $\mathrm{Bx}_{\mathfrak{c}}$ in the sense of \cite[Definition~2.12]{QEDII}, so that $d\hat A=F^{\mathrm{u}}$ on $\mathrm{Pq}(\mathrm{Bx}_{\mathfrak{c}})$.
By (6.1), Lemma~2\,(a) and Section~5,
\begin{equation}
\begin{split}
\sum_{p\in\mathrm{Pq}(\mathrm{Bx}_{\mathfrak{c}})}(d\hat{\mathfrak{a}}^{\mathrm{o}}_{\mathfrak{c}})(p)\cdot(d\hat A)(p)
&=X^{\mathrm{o}}_{\mathfrak{c}}(F^{\mathrm{u}})=\left\langle d\hat{\mathfrak{a}}^{\mathrm{o}}_{\mathfrak{c}},F^{\mathrm{u}}\right\rangle\\
&=\left\langle\mathrm{Prd}_n\,d\hat{\mathfrak{a}}^{\mathrm{o}}_{\mathfrak{c}},F^{\mathrm{u}}\right\rangle_n=\left\langle d\hat{\mathfrak{a}}_{\mathfrak{c}},F^{\mathrm{u}}\right\rangle_n=X_{\mathfrak{c}}.
\end{split}
\tag{6.3}
\end{equation}
Hence, on every torus $n$ admitted in the statement, the element $G$ of Subsection~6.1 is the function $\mathrm{Pol}(X_{\mathfrak{c}})$ of $F^{\mathrm{u}}$.

Consider the index $\mathfrak{c}=0$.
Then $\mathsf{g}_0$ is the identity (Lemma~4\,(a)), $\hat{\mathfrak{a}}^{\mathrm{o}}_{0}=\hat{\mathfrak{a}}^{\mathrm{o}}$, the box (5.1) is $\mathrm{Bx}_0$, the cell is $Q^{\circ}_{0}$, and $X_0=X$.
In this case Subsection~6.1 asserts that the map
\[
G\colon\ \hat A\mapsto\mathrm{Pol}\left(\left\langle d\hat{\mathfrak{a}}^{\mathrm{o}},d\hat A\right\rangle\right)\qquad\text{on }\mathbb{R}^{\mathrm{Ln}(\mathrm{Bx}_0)},
\]
in which $\left\langle d\hat{\mathfrak{a}}^{\mathrm{o}},d\hat A\right\rangle$ stands for the sum over $\mathrm{Pq}(\mathrm{Bx}_0)$ in the definition of $G$, is a bosonic element of $\mathrm{Ad}(\mathrm{Bx}_0)$, hence an element of $\mathcal{A}(Q^{\circ}_{0})$, and that it does not involve $n$.
By (6.2) and the argument that follows it, $\mathrm{Bx}_0$ embeds in $n\mathbb{Z}^3$, with $2r_0+2<\mathrm{per}\le n$.
By (6.3), the reading of $G$ on the torus $n$ (\cite[Definition~2.12]{QEDII}, local lifts) is the function $\mathrm{Pol}(X)$, $X=\left\langle d\hat{\mathfrak{a}},F^{\mathrm{u}}\right\rangle$, of $F^{\mathrm{u}}\in C^{2}(n)$.
In this sense $X$ (the case $\mathrm{Pol}(z)=z$) and every polynomial in $X$ with complex coefficients are bosonic elements of $\mathrm{Ad}(\mathrm{Bx}_0)$, each of them one and the same element for all sides $n$ admitted in the statement.

\textbf{6.3. $O_m$ is real.}
The number $\boldsymbol{\sigma}^2$ is real; hence the polynomials $H_2(z)=z^2-\boldsymbol{\sigma}^2$ and $H_4(z)=z^4-6\boldsymbol{\sigma}^2z^2+3\boldsymbol{\sigma}^4$ of Section~I.7.6 have real coefficients.
Moreover $e>0$, the numbers $c_m$ are real, and $X^{\mathrm{o}}_{0}(u)$ is real for every $u\in C^{2}$.
Therefore
\[
O_m(u):=e^{-m}\cdot H_m\left(X^{\mathrm{o}}_{0}(u)\right)-c_m
\]
is a real number for every $u\in C^{2}$; for $u=F^{\mathrm{u}}\in C^{2}(n)$ it equals $e^{-m}\cdot H_m(X)-c_m$ by (6.3).

\textbf{6.4. Mirror copies.}
By Subsection~6.1, applied with $\mathrm{Pol}=e^{-m}\cdot H_m-c_m$ and $\mathfrak{c}=0$, the function $O_m$ lies in $\mathcal{A}(Q^{\circ}_{0})$.
For a function $O$ of two-cochains $u$, Section~I.6.4 (mirror copies; compare the action of $W$ on bosonic integrands in \cite[Section~2.3]{QEDII} and \cite[Definition~2.19]{QEDII}) gives
\[
O^{[\mathfrak{c}]}(u)=O(\mathsf{g}_{\mathfrak{c}}^{*}u)\quad\text{if }\sum_\mu\mathfrak{c}_\mu\text{ is even},\qquad
O^{[\mathfrak{c}]}(u)=\overline{O(\mathsf{g}_{\mathfrak{c}}^{*}u)}\quad\text{if }\sum_\mu\mathfrak{c}_\mu\text{ is odd}.
\]
As $O_m$ is real (Subsection~6.3), for every $\mathfrak{c}\in\mathbb{Z}^3$ and every $u\in C^{2}$
\[
O_m^{[\mathfrak{c}]}(u)=O_m(\mathsf{g}_{\mathfrak{c}}^{*}u)=e^{-m}\cdot H_m\left(X^{\mathrm{o}}_{0}(\mathsf{g}_{\mathfrak{c}}^{*}u)\right)-c_m .
\]
Lemma~3\,(d), (a) and (5.1) give
\[
X^{\mathrm{o}}_{0}(\mathsf{g}_{\mathfrak{c}}^{*}u)=\left\langle d\hat{\mathfrak{a}}^{\mathrm{o}},\mathsf{g}_{\mathfrak{c}}^{*}u\right\rangle=\left\langle(\mathsf{g}_{\mathfrak{c}})_{*}\,d\hat{\mathfrak{a}}^{\mathrm{o}},u\right\rangle=\left\langle d(\mathsf{g}_{\mathfrak{c}})_{*}\hat{\mathfrak{a}}^{\mathrm{o}},u\right\rangle=\left\langle d\hat{\mathfrak{a}}^{\mathrm{o}}_{\mathfrak{c}},u\right\rangle=X^{\mathrm{o}}_{\mathfrak{c}}(u).
\]
For $u=F^{\mathrm{u}}\in C^{2}(n)$ the cochain $\mathsf{g}_{\mathfrak{c}}^{*}F^{\mathrm{u}}$ lies in $C^{2}(n)$ (Lemma~3\,(c)), and $X^{\mathrm{o}}_{\mathfrak{c}}(F^{\mathrm{u}})=X_{\mathfrak{c}}$ by (6.3).
Hence $O_m^{[\mathfrak{c}]}=e^{-m}\cdot H_m(X_{\mathfrak{c}})-c_m$ for every $\mathfrak{c}\in\mathbb{Z}^3$.
\end{proof}

\subsubsection*{7. Proof of (e)}

\begin{proof}[Proof of Proposition~41\,(e)]
Let $\mathfrak{c}\in2\mathbb{Z}^3$.
Then $(-1)^{\mathfrak{c}_\mu}=1$ for $\mu=1,2,3$, and (1.2) becomes
\[
(\mathsf{g}_{\mathfrak{c}}y)_\mu=\left(\mathfrak{c}_\mu+\frac12\right)\cdot\mathrm{hc}+y_\mu-\frac{\mathrm{hc}}{2}=y_\mu+\mathrm{hc}\cdot\mathfrak{c}_\mu .
\]
So $\mathsf{g}_{\mathfrak{c}}$ is the translation by $\mathrm{hc}\cdot\mathfrak{c}$, and $\mathrm{hc}\cdot\mathfrak{c}_\mu=\frac{\mathrm{per}}{2}\cdot\mathfrak{c}_\mu=\mathrm{per}\cdot\frac{\mathfrak{c}_\mu}{2}$ lies in $\mathrm{per}\,\mathbb{Z}$.
\end{proof}

\subsubsection*{8. Uniformity}

\begin{named}{Remark}
Let $q=e\cdot\eta^{\frac12}$ be the charge of level~$0$ (Section~I.2.5) and $r_q=2\pi/q$, so that $F^{\mathrm{u}}=dA+r_q\cdot\mathfrak{z}$ for a configuration $(A,\mathfrak{z})$ on the torus $n$ (Section~I.7.1).
The assertions (a)--(e) hold for every $q>0$ in the place of $e\cdot\eta^{\frac12}$, and $G$, $\varepsilon_G=\frac14$ and $\mathrm{Cq}$ do not involve $q$.
Indeed, in the condition of gauge invariance in the definition of $\mathrm{Ad}(\mathrm{Bx}_{\mathfrak{c}})$, \cite[Definition~2.12]{QEDII} attaches to a monomial $\Psi$ in the generators the factor $\exp\left(iq\left\langle\boldsymbol{\lambda},\mathfrak{q}(\Psi)\right\rangle\right)$;
as $G$ is scalar, no generator occurs in it, $\Psi=1$ is the only monomial, its charge density vanishes, and this factor equals $1$.
Moreover, (6.3) holds for every $q>0$, since of a local lift $\hat A$ on $\mathrm{Bx}_{\mathfrak{c}}$ it requires only the identity $d\hat A=F^{\mathrm{u}}$ on $\mathrm{Pq}(\mathrm{Bx}_{\mathfrak{c}})$.

The constants $\varepsilon_G=\frac14$ and $\mathrm{Cq}=\sum_j|\mathrm{cf}_j|\cdot\left(\frac{2j\boldsymbol{\sigma}^2}{\mathrm{e}}\right)^{\frac{j}{2}}$ of Subsection~6.1 depend only on $\mathrm{Pol}$ and on $\boldsymbol{\sigma}^2$;
they depend neither on $n$ nor on $\mathfrak{c}$.

The sides admitted in Proposition~41 are all multiples $n\ge n_*$ of $\mathrm{per}$, of either parity.
A torus of the flow, $R=h\cdot L^{T}\in\mathrm{Tor}(e)$ (Section~I.6.2), has $R/\eta=h\cdot L^{T+N}$ sites per side,
and among the admitted sides only $n_*=R^{*}(e)/\eta$ and $L\cdot n_*=R_{\mathrm{top}}(e)/\eta$ are of this form.
\end{named}

\subsection*{Conditions on the constants}
\mbox{}

\medskip
\noindent
\begin{tabular}{|p{2.4cm}|p{4.7cm}|p{4.1cm}|}
\hline
condition & where it is used & line of the list of constants \\
\hline
$L$ is an odd integer $\ge5$
&
$\mathrm{per}=n_*/L$ is odd, and $\mathrm{hc}=m_{\mathrm{p}}+\frac12$ with an integer $m_{\mathrm{p}}$ (Subsection~0.1, through Section~I.6.4); $0<\eta=L^{-N}\le1$ (Subsection~0.1)
&
Section~I.8.2, line (C3), which is Proposition~14.2\,(a), line (C3); with Proposition~14.2\,(b), condition (L5) \\
\hline
$M$ is a power of $L$, and $M\ge L^2$
&
$T^{*}(e)=\log_L\left(M\cdot r_K(e)\right)-1$ is an integer $\ge1$, so that $\mathrm{per}=h\cdot L^{T^{*}(e)+N-1}$ is an integer; here $h=L^i$ with an integer $i\ge0$, because $\bar{m}\le1$ (Section~I.0.5) (Subsection~0.1, through Section~I.6.2)
&
Section~I.8.2, line (C5), which is Proposition~14.2\,(a), line (C5); with Proposition~14.2\,(b), condition (M1) \\
\hline
$\bar{m}\in(0,1]$
&
$h=L^{\left\lfloor\log_L\frac{1}{\bar{m}}\right\rfloor}$ is a power of $L$ with an integer exponent $\ge0$ (Subsection~0.1)
&
Section~I.0.5; Section~I.8.2, line (C10), which is Proposition~14.2\,(a), the lines (C10) for $\bar{m}$ and $h$ (the range of a standing datum: Section~G6.A1) \\
\hline
\end{tabular}

\medskip
By (6.2), every side $n$ admitted in Proposition~41 satisfies $n\ge15$, and hence the requirement $n\ge4$ on the sides in Annex~B, Section~G6.Z4.

\section*{Chapter 31. The zero-flux sector and the complement of the small-field set: proofs of Propositions 31.1, 31.2}

\subsection*{0. Conventions of this chapter}

In this chapter, (9.2) denotes the display of Proposition~31.2.
In Section~1 the labels (s1), (s2), (s3) denote the facts (I.s1)--(I.s3), written in lattice units in the form stated in 1.1.1.

The symbols $\mathrm{kc}$, $\mathrm{Mo}_i$, $\mathrm{La}$ and $\delta_{\mathrm{xo},\beta}$ have meanings that are local to this chapter.
The symbol $\mathrm{cf}'$ denotes a second constant of the same kind as $\mathrm{cf}$.
A plaquette is denoted by $\mathrm{pl}$; the letter $p$ is the integer of (C7).

The set $T^{\mathrm{f}}_0(n)$ is the set $T^{\mathrm{f}}_0$ of Section~I.1.3 formed with the side $n$, which need not lie in $\mathrm{Cl}$.
The operator $D_a^{\boldsymbol{\sigma}}$ is the operator $D_A^{\boldsymbol{\sigma}}$ of (I.O3) at $A=a$.
The letter $\mathrm{Ev}$ denotes a Borel set, as in Proposition~35.7\,(a).
The letters $u$ and $u'$ denote two-cochains, as in Section~I.6.4, Section~G6.A3 and Proposition~50\,(g).
The letters $m$ and $m^{(\mu)}$ denote integer-valued one-cochains, as in Proposition~38.1\,(b), (e).
In the proof of Lemma~3 the letter $t$ denotes a real number $t\ge0$.
The symbols $\operatorname{covol}$ and $\Gamma_n$ are those of \cite[Lemma~21.32]{QEDI} and of its proof.
The operator $\Delta_0$ is $d^{*}d$ on $C^0(n)$, and $\operatorname{det}'\Delta_0$ is the product of its non-zero eigenvalues, as in Proposition~38.2 and in the definition of $\mathrm{Nc}_n$ in Section~I.7.3.
This $\Delta_0$ acts on the functions of the photon sites; it is not the operator $\Delta_k=(dH_k)^T dH_k$ of \cite[Section~5.2.5]{QEDI} at $k=0$, which acts on one-cochains.

In Section~1 the letters $T$ and $T'$ denote coordinate tori $C(n_1,n_2,n_3)$ in the sense of \cite[Definition~2.11]{QEDII}; in Section~2 the letter $T$ is the integer with $R=hL^T$.
The symbols $\mathrm{hz}$, $\mathrm{hz}'$ and $\mathrm{hz}^{(t)}$ denote two-cochains on such tori; they play the role of the cochains $\beta$ of \cite[Proposition~13.11]{QEDII}.
In Step~D2 the index $t$ labels the planes $\{c_1=t\}$, and the number $\mathrm{Z0}_T(\mathrm{hz})$ defined there is the number $\mathcal{Z}(\beta)$ of the proof of \cite[Proposition~13.11]{QEDII} at $\beta=\mathrm{hz}$.
The function $\Phi_0$ is the indicator $\Phi_{k'}$ of \cite[Definition~13.7(a)]{QEDII} at $k'=0$.
An element of the group $W$ of \cite[Definition~2.2]{QEDII}, denoted $g$ there, is written $g_{\mathrm{w}}\in W_{\mathrm{af}}$.
In Step~D2, where such an element occurs, the letter $r$ denotes its linear part, as in Section~I.4.7.
In Sections~2 and~3 the letter $r$ is the integer $r$ of (C7).
The curvature $F$ of \cite[Definition~2.6]{QEDII} is written $F^{\mathrm{u}}$.
In this chapter the letter $g$ itself denotes the label of a fermion site and nothing else.

\subsection*{1. Proof of Proposition 31.1}

The statement proved in this section is Proposition~31.1 of Part~II.

\subsubsection*{1.1. Setting in lattice units; three lemmas}

Throughout Section~1, $n\ge4$ is the integer of the data of Proposition~31.1, $\rho\in(0,1)$, $w=(w_0,w_1,w_2,w_3)=(6,6\sqrt{2},6\sqrt{2},6)$ and $\mathrm{sw}=w_0+w_1+w_2+w_3$ (Section~I.0.5).
The conventions are those of Sections~I.1 and~I.2 at $k=0$, with $\eta_0=1$, with $q$ in the place of $e_k$ and with the period $n$ in the place of $\ell$; this is the setting of Proposition~46 at $(n,n,n)$.
All objects of Section~1 are objects of level~$0$ in lattice units.
Over every photon site the spinors have $192$ components, four at each of the $48$ fermion sites over it, and every fermion site carries the weight $\frac{1}{48}$.
The counting norm and the counting pairing of cochains (Section~I.7.1) are denoted by $\|\cdot\|$ and $\langle\cdot,\cdot\rangle$.

\paragraph{1.1.1. Sites.}
A fermion site is a pair $x=(c,g)$ with $c\in\mathbb{Z}^3$ and $g\in O_{\mathrm{cub}}$; the $48$ sites of the cluster $\{c\}\times O_{\mathrm{cub}}$ lie over the photon site $c$.
At $\eta_0=1$ the display (I.1.2) reads
\[
x^{(i)}=(c,gs_i)\quad(i=0,1,2),\qquad x^{(3)}=(c+ge_1,gs_3),\qquad c_3=c+ge_1,\qquad\hat d_i(x)=-g\hat n_i\quad(i=0,1,2,3).
\]
The following three facts of Section~I.1.3 are used.
\begin{itemize}
\item[(s1)] $\left(x^{(i)}\right)^{(i)}=x$, and $(x+v)^{(i)}=x^{(i)}+v$ for $v\in\mathbb{Z}^3$.
For $i\le2$ the first identity follows from $s_i^2=1$.
For $i=3$ one has $\left(x^{(3)}\right)^{(3)}=\left(c+ge_1+gs_3e_1,gs_3^2\right)=(c,g)$, because $s_3e_1=-e_1$.
The translations act on $c$ only.
\item[(s2)] $\hat d_i\left(x^{(i)}\right)=-\hat d_i(x)$.
Indeed $\hat d_i\left(x^{(i)}\right)=-gs_i\hat n_i=g\hat n_i$, because $s_i\hat n_i=-\hat n_i$.
\item[(s3)] $\sum_{i=0}^{3}w_i\hat d_i(x)=0$.
Indeed the sum equals $-g\sum_iw_i\hat n_i$, and $\sum_iw_i\hat n_i=6e_3+6(e_2-e_3)+6(e_1-e_2)-6e_1=0$ by (I.1.1).
\end{itemize}
Let
\[
T^{\mathrm{f}}_0(n):=\left(\mathbb{Z}^3/n\mathbb{Z}^3\right)\times O_{\mathrm{cub}}
\]
be the set of the fermion sites of the torus; it is the set $T^{\mathrm{f}}_0$ of Section~I.1.3 formed with the side $n$, and $\#|T^{\mathrm{f}}_0(n)|=48n^3$.
By (s1), for each $i$ the map $x\mapsto x^{(i)}$ induces a map of $T^{\mathrm{f}}_0(n)$ into itself whose square is the identity.
Consequently it is a bijection of $T^{\mathrm{f}}_0(n)$.

\paragraph{1.1.2. Spinors with $\mathrm{kc}$ components.}
Let $\mathrm{kc}\in\{2,4\}$, and put
\begin{gather*}
\mathfrak{S}^{\mathrm{ch}}_{\mathrm{kc}}(n):=\left\{f\colon\mathbb{Z}^3\times O_{\mathrm{cub}}\to\mathbb{C}^{\mathrm{kc}}\ :\ f(c+ne_\mu,g)=-f(c,g)\ \text{for}\ \mu=1,2,3\right\},\\
(f',f)_{\mathrm{kc}}:=\frac{1}{48}\sum_{c\in\{0,\dots,n-1\}^3}\ \sum_{g\in O_{\mathrm{cub}}}\ \sum_{\beta=1}^{\mathrm{kc}}\overline{f'_\beta(c,g)}\cdot f_\beta(c,g).
\end{gather*}
An element of $\mathfrak{S}^{\mathrm{ch}}_{\mathrm{kc}}(n)$ is determined by its values on $\{0,\dots,n-1\}^3\times O_{\mathrm{cub}}$, and these values can be prescribed arbitrarily.
Hence the complex dimension of $\mathfrak{S}^{\mathrm{ch}}_{\mathrm{kc}}(n)$ is $48\cdot\mathrm{kc}\cdot n^3$.
One has $\mathfrak{S}^{\mathrm{ch}}_4(n)=\mathfrak{S}^{\mathrm{ch}}(1;n)$, a space of dimension $192n^3$, and $(\cdot,\cdot)_4=(\cdot,\cdot)^{\mathrm{ch}}$ by (I.1.3) and Section~I.1.6 at $\eta_0=1$, the weight of a fermion site being $\frac{1}{48}$.
For $f',f\in\mathfrak{S}^{\mathrm{ch}}_{\mathrm{kc}}(n)$ the function $x\mapsto\overline{f'(x)}\cdot f(x)$, where $v'\cdot v=\sum_\beta v'_\beta v_\beta$ (Section~I.1.6), is $n\mathbb{Z}^3$-periodic, since the two signs $-1$ cancel.
Therefore
\[
(f',f)_{\mathrm{kc}}=\frac{1}{48}\sum_{x\in T^{\mathrm{f}}_0(n)}\overline{f'(x)}\cdot f(x),
\]
and this sum may be re-indexed by any bijection of $T^{\mathrm{f}}_0(n)$.
Let $|f|:=(f,f)_{\mathrm{kc}}^{\frac{1}{2}}$.
For a linear map $O$ of $\mathfrak{S}^{\mathrm{ch}}_{\mathrm{kc}}(n)$ into itself, $|O|$ denotes its operator norm for this norm, $O^{*}$ its adjoint with respect to $(\cdot,\cdot)_{\mathrm{kc}}$, and $\operatorname{Re}O:=\frac{1}{2}\left(O+O^{*}\right)$.
For a real number $z$, the relation $\operatorname{Re}O\ge z$ means that $\operatorname{Re}(f,Of)_{\mathrm{kc}}\ge z\cdot|f|^2$ for every $f$.
The objects $O^{*}$, $|O|$, $\operatorname{Re}O$ and the relation $\operatorname{Re}O\ge z$ do not change when $(\cdot,\cdot)_{\mathrm{kc}}$ is multiplied by a positive number.
They are therefore the same for $(\cdot,\cdot)_{\mathrm{kc}}$ and for the counting scalar product $\sum_{x\in T^{\mathrm{f}}_0(n)}\overline{f'(x)}\cdot f(x)=48\cdot(f',f)_{\mathrm{kc}}$ of \cite[Definition~3.9]{QEDII}.
A determinant does not depend on a scalar product.

For $v\in\mathbb{R}^3$ put
\[
\mathrm{Gv}_4(v):=\gamma\cdot v=\sum_\mu v_\mu\gamma_\mu,\qquad\mathrm{Gv}_2(v):=\sigma\cdot v=\sum_\mu v_\mu\sigma_\mu.
\]
By Section~I.2.1 one has $\gamma_\mu=\operatorname{diag}(\sigma_\mu,-\sigma_\mu)$, and hence $\mathrm{Gv}_4(v)=\operatorname{diag}\left(\mathrm{Gv}_2(v),-\mathrm{Gv}_2(v)\right)$.
The Pauli matrices are Hermitian and satisfy $\sigma_\mu\sigma_\nu+\sigma_\nu\sigma_\mu=2\delta_{\mu\nu}$.
It follows that, for $\mathrm{kc}=2$ and for $\mathrm{kc}=4$,
\[
\mathrm{Gv}_{\mathrm{kc}}(v)^{*}=\mathrm{Gv}_{\mathrm{kc}}(v),\qquad\mathrm{Gv}_{\mathrm{kc}}(v)^2=|v|_2^2\cdot1_{\mathrm{kc}},
\]
where $1_{\mathrm{kc}}$ is the identity matrix of size $\mathrm{kc}$.
So for $|v|_2=1$ the matrix $\mathrm{Gv}_{\mathrm{kc}}(v)$ is Hermitian and unitary, and its eigenvalues lie in $\{1,-1\}$.
For $\mathrm{kc}=4$ this is (I.c1).

\paragraph{1.1.3. Cochains, phases.}
A one-cochain $B$ on $\mathbb{Z}^3/n\mathbb{Z}^3$, real or complex, is an $n\mathbb{Z}^3$-periodic function on the ordered pairs of nearest neighbours of $\mathbb{Z}^3$ such that $B(c',c)=-B(c,c')$ (Section~I.1.1 and Section~G6.A1).
With the charge $q$, (I.1.6) and (I.2.3) give
\[
U_B\left(x,x^{(i)}\right)=1\quad(i=0,1,2),\qquad U_B\left(x,x^{(3)}\right)=\exp\left(\mathrm{i}q\cdot B(c,c_3)\right).
\]
By (s1), the photon site of $x^{(3)}$ is $c_3$, and the photon site of $\left(x^{(3)}\right)^{(3)}$ is $c$.
Hence, for $i=0,\dots,3$,
\begin{equation}\tag{31.1}
\begin{aligned}
&U_B\left(x^{(i)},x\right)\cdot U_B\left(x,x^{(i)}\right)=1,\qquad U_B\left(x+v,x^{(i)}+v\right)=U_B\left(x,x^{(i)}\right)\quad\text{for}\ v\in n\mathbb{Z}^3,\\
&\text{and, for real}\ B,\qquad\left|U_B\left(x,x^{(i)}\right)\right|=1.
\end{aligned}
\end{equation}

\paragraph{1.1.4. The operator.}
For a complex one-cochain $B$ on $\mathbb{Z}^3/n\mathbb{Z}^3$ and $f\in\mathfrak{S}^{\mathrm{ch}}(1;n)$, the displays (I.2.1) and (I.2.3) at $\eta_0=1$ give
\begin{equation}\tag{31.2}
\begin{aligned}
(D_Bf)(x)&=-\sum_{i=0}^{3}M_i(x)\left[U_B\left(x,x^{(i)}\right)\cdot f\left(x^{(i)}\right)-f(x)\right],\\
M_i(x)&=\frac{w_i}{2}\cdot\left(\rho-\gamma\cdot\hat d_i(x)\right).
\end{aligned}
\end{equation}

\paragraph{1.1.5. Hodge decomposition.}
The space $C^1(n)$ is the orthogonal sum ${V_{\mathrm{g}}}\oplus{V_{\mathrm{h}}}\oplus{V_{\mathrm{t}}}$, where ${V_{\mathrm{g}}}=dC^0(n)$, ${V_{\mathrm{h}}}$ is the space of the constant one-cochains, and ${V_{\mathrm{t}}}$ is the orthogonal complement of ${V_{\mathrm{g}}}\oplus{V_{\mathrm{h}}}$ (Section~I.7.1).
By \cite[Section~10.1.1]{QEDI}, where the identity is stated for all sides $n_\mu\ge2$,
\[
\ker\left(d\colon C^1(n)\to C^2(n)\right)={V_{\mathrm{g}}}\oplus{V_{\mathrm{h}}}.
\]
Hence $dC^1(n)=d{V_{\mathrm{t}}}$, and $d$ is injective on ${V_{\mathrm{t}}}$.
Consequently the restriction $\mathrm{dVt}$ of $d$ to ${V_{\mathrm{t}}}$ is a linear bijection of ${V_{\mathrm{t}}}$ onto $F_N=dC^1(n)$.
The field $A(F^{\mathrm{u}})\in{V_{\mathrm{t}}}$ with $dA(F^{\mathrm{u}})=F^{\mathrm{u}}$ (Section~I.7.3) is therefore
\[
A(F^{\mathrm{u}})=\mathrm{dVt}^{-1}F^{\mathrm{u}};
\]
it is linear in $F^{\mathrm{u}}$, and $A(dA)=A$ for $A\in{V_{\mathrm{t}}}$.
Let $\mathrm{dimF}:=\dim F_N=\dim{V_{\mathrm{t}}}$.
The spaces ${V_{\mathrm{t}}}$ and $F_N$ carry the scalar products of $C^1(n)$ and of $C^2(n)$.
Their Lebesgue measures are denoted by $\sigma$ and $\mathrm{lbF}$; each is the image of the Lebesgue measure of $\mathbb{R}^{\mathrm{dimF}}$ under the map that sends a vector of coordinates to the corresponding linear combination of the vectors of an orthonormal basis.
One has
\[
\mathrm{d}\gamma_N(F^{\mathrm{u}})=(2\pi)^{-\frac{\mathrm{dimF}}{2}}\cdot\exp\left(-\frac{1}{2}\|F^{\mathrm{u}}\|^2\right)\mathrm{d}\,\mathrm{lbF}(F^{\mathrm{u}}).
\]
Let $\mathrm{Jac}$ be the modulus of the determinant of the matrix of $\mathrm{dVt}$ with respect to orthonormal bases of ${V_{\mathrm{t}}}$ and of $F_N$.
Then $\mathrm{Jac}>0$, and $\mathrm{Jac}$ does not depend on the choice of the bases.

\begin{named}{Lemma 1}[Image of the measure; the identity stated with the definition of $\mathrm{zV}_n$ in Proposition~31.1]
The number $\mathrm{zV}_n=\int_{{V_{\mathrm{t}}}}\mathrm{d}\sigma(A)\,\exp\left(-\frac{1}{2}\|dA\|^2\right)$ of Proposition~31.1 is
\[
\mathrm{zV}_n=\frac{(2\pi)^{\frac{\mathrm{dimF}}{2}}}{\mathrm{Jac}}\in(0,\infty).
\]
Let $G$ be a Borel function on $F_N$ that is $\ge0$, or that satisfies $|G(F^{\mathrm{u}})|\le\mathrm{cf}\cdot\exp\left(\mathrm{cf}\cdot\|F^{\mathrm{u}}\|\right)$ for some $\mathrm{cf}$.
Then
\[
\int_{{V_{\mathrm{t}}}}\mathrm{d}\sigma(A)\,G(dA)\cdot\exp\left(-\frac{1}{2}\|dA\|^2\right)=\mathrm{zV}_n\cdot\int_{F_N}\mathrm{d}\gamma_N\,G.
\]
In the second case both integrals converge absolutely, and
\[
\int\mathrm{d}\gamma_N(F^{\mathrm{u}})\,\exp\left(\mathrm{cf}\cdot\|F^{\mathrm{u}}\|\right)\le\exp\left(\mathrm{cf}^2\right)\cdot2^{\frac{\mathrm{dimF}}{2}}.
\]
\end{named}

\begin{proof}
In the coordinates of orthonormal bases of ${V_{\mathrm{t}}}$ and of $F_N$, the map $\mathrm{dVt}$ becomes an invertible linear map of $\mathbb{R}^{\mathrm{dimF}}$ whose determinant has the modulus $\mathrm{Jac}$.
Applied to the inverse map, \cite[Theorem~2.20(e) and 2.23]{Rudin1987} gives
\[
\sigma\left(\mathrm{dVt}^{-1}\mathrm{Ev}\right)=\mathrm{Jac}^{-1}\cdot\mathrm{lbF}(\mathrm{Ev})
\]
for every Borel subset $\mathrm{Ev}$ of $F_N$.
This is the identity
\[
\int_{{V_{\mathrm{t}}}}\mathrm{d}\sigma(A)\,G_1(\mathrm{dVt}\,A)=\mathrm{Jac}^{-1}\cdot\int_{F_N}\mathrm{d}\,\mathrm{lbF}\,G_1
\]
for $G_1=\mathbf{1}_{\mathrm{Ev}}$.
By linearity the identity holds for the finite linear combinations of such indicator functions with coefficients $\ge0$.
Every Borel function $G_1\ge0$ on $F_N$ is the pointwise limit of an increasing sequence of such combinations; so the identity holds for it by monotone convergence \cite[Theorems~1.17 and 1.26]{Rudin1987}.
By linearity once more, it holds for every Borel function $G_1$ that is integrable with respect to $\mathrm{lbF}$.
For $G_1:=G\cdot\exp\left(-\frac{1}{2}\|\cdot\|^2\right)$ the left side is the left side of the assertion, since $\mathrm{dVt}\,A=dA$, and the right side equals $\mathrm{Jac}^{-1}\cdot(2\pi)^{\frac{\mathrm{dimF}}{2}}\cdot\int\mathrm{d}\gamma_N\,G$.
The choice $G=1$ gives the value of $\mathrm{zV}_n$.
For $y\ge0$ one has $\mathrm{cf}\cdot y\le\mathrm{cf}^2+\frac{y^2}{4}$, because $\left(\frac{y}{2}-\mathrm{cf}\right)^2\ge0$.
Hence
\[
\exp\left(\mathrm{cf}\cdot\|F^{\mathrm{u}}\|-\frac{1}{2}\|F^{\mathrm{u}}\|^2\right)\le\exp\left(\mathrm{cf}^2\right)\cdot\exp\left(-\frac{1}{4}\|F^{\mathrm{u}}\|^2\right),
\]
and
\[
(2\pi)^{-\frac{\mathrm{dimF}}{2}}\cdot\int\mathrm{d}\,\mathrm{lbF}(F^{\mathrm{u}})\,\exp\left(-\frac{1}{4}\|F^{\mathrm{u}}\|^2\right)=(2\pi)^{-\frac{\mathrm{dimF}}{2}}\cdot(4\pi)^{\frac{\mathrm{dimF}}{2}}=2^{\frac{\mathrm{dimF}}{2}}.
\]
In the second case this bound makes $G_1$ integrable with respect to $\mathrm{lbF}$, and the identity for $|G|\ge0$ shows that both integrals converge absolutely.
\end{proof}

\begin{named}{Lemma 2}[The flux is a function of the curvature]
For $\mu<\nu$ let $\Pi_{\mu\nu}$ be the set of the $n^2$ plaquettes $(c;\mu\nu)$ with $c_{\mathfrak{l}}=0$, where $\mathfrak{l}$ is the third direction (Section~I.7.1).
For $u\in C^2(n)$ put
\[
\Phi_0(u):=\prod_{\mu<\nu}\mathbf{1}\left[\sum_{\mathrm{pl}\in\Pi_{\mu\nu}}u(\mathrm{pl})=0\right],
\]
a Borel function on $C^2(n)$ with the values $0$ and $1$.
For every configuration $(A,\mathfrak{z})$ (Section~I.7.1) one has
\[
\sum_{\mathrm{pl}\in\Pi_{\mu\nu}}F^{\mathrm{u}}(\mathrm{pl})=r_q\cdot w_{\mu\nu}.
\]
Consequently $\mathbf{1}[\text{flux}=0]=\Phi_0(F^{\mathrm{u}})$, and $Z^{(0)}_n[G]=Z_n[\Phi_0\cdot G]$ for every Borel function $G$ with $|G|\le\mathrm{cf}\cdot\exp\left(\mathrm{cf}\cdot\|F^{\mathrm{u}}\|\right)$; the function $\Phi_0\cdot G$ satisfies the same bound.
(The two sides denote the same sum of integrals; its absolute convergence is shown in Step~B4.)
For a closed $u\in C^2(n)$ the sum of $u$ over the plaquettes $(c;\mu\nu)$ with $c_{\mathfrak{l}}=j$ does not depend on $j\in\mathbb{Z}/n\mathbb{Z}$.
Hence, on closed cochains, in particular at $u=F^{\mathrm{u}}$ and at $u=\mathfrak{z}$, the choice of the planes in the definitions of $\Phi_0$ and of the flux is immaterial.
\end{named}

\begin{proof}
By Section~I.1.1,
\[
(dA)_{\mu\nu}(c)=A_\mu(c)-A_\mu(c+e_\nu)+A_\nu(c+e_\mu)-A_\nu(c).
\]
In the sum over the $c$ with $c_{\mathfrak{l}}=0$, that is over $(c_\mu,c_\nu)\in(\mathbb{Z}/n\mathbb{Z})^2$, the translations $c\mapsto c+e_\nu$ and $c\mapsto c+e_\mu$ are bijections of the set of summation.
So the sum of the first two terms vanishes, and so does the sum of the last two.
Since $F^{\mathrm{u}}=dA+r_q\cdot\mathfrak{z}$, it follows that
\[
\sum_{\mathrm{pl}\in\Pi_{\mu\nu}}F^{\mathrm{u}}(\mathrm{pl})=r_q\cdot\sum_{\mathrm{pl}\in\Pi_{\mu\nu}}\mathfrak{z}(\mathrm{pl})=r_q\cdot w_{\mu\nu},
\]
and $r_q>0$.
The identity for $Z^{(0)}_n$ is the definition $Z^{(0)}_n[G]=Z_n\left[\mathbf{1}[\text{flux}=0]\cdot G\right]$ of Section~I.7.3.
For the planes, let $\mathrm{Su}_{\mu\nu}(j)$ be the sum of $u$ over the plaquettes $(c;\mu\nu)$ with $c_{\mathfrak{l}}=j$.
By Section~I.1.1,
\[
(du)(c)=\left[u_{23}(c+e_1)-u_{23}(c)\right]-\left[u_{13}(c+e_2)-u_{13}(c)\right]+\left[u_{12}(c+e_3)-u_{12}(c)\right].
\]
When $(du)(c)$ is summed over the $c$ with $c_{\mathfrak{l}}=j$, the two brackets that contain a translation within the plane contribute zero, and the third bracket contributes $\pm\left(\mathrm{Su}_{\mu\nu}(j+1)-\mathrm{Su}_{\mu\nu}(j)\right)$.
Hence $du=0$ implies $\mathrm{Su}_{\mu\nu}(j+1)=\mathrm{Su}_{\mu\nu}(j)$.
\end{proof}

\begin{named}{Lemma 3}[From the exponential moment to the tails]
Let $\mathrm{prm}$ be a Borel probability measure on $F_N$ such that
\[
\int\mathrm{d}\,\mathrm{prm}(F^{\mathrm{u}})\,\exp\langle\phi,F^{\mathrm{u}}\rangle\le\exp\left(\frac{1}{2}\|\phi\|^2\right)\qquad\text{for every}\ \phi\in F_N.
\]
Then the following hold.
\begin{itemize}
\item[(i)] For $0\ne\phi\in F_N$ and $y\ge0$,
\[
\mathrm{prm}\left\{\langle\phi,F^{\mathrm{u}}\rangle\ge y\right\}\le\exp\left(-\frac{y^2}{2\|\phi\|^2}\right).
\]
\item[(ii)] For $\phi\in F_N$ with $\|\phi\|\le1$ and $y\ge0$,
\[
\mathrm{prm}\left\{\left|\langle\phi,F^{\mathrm{u}}\rangle\right|>y\right\}\le2\exp\left(-\frac{y^2}{2}\right).
\]
\end{itemize}
The measure $\gamma_N$ satisfies the hypothesis, with equality.
\end{named}

\begin{proof}
(i) For $t\ge0$ one has $\mathbf{1}\left[\langle\phi,F^{\mathrm{u}}\rangle\ge y\right]\le\exp\left(t\cdot\left(\langle\phi,F^{\mathrm{u}}\rangle-y\right)\right)$.
Integration of this inequality (Markov's inequality) and the hypothesis for $t\phi$ give
\[
\mathrm{prm}\left\{\langle\phi,F^{\mathrm{u}}\rangle\ge y\right\}\le\exp\left(-t\cdot y+\frac{t^2}{2}\|\phi\|^2\right).
\]
At $t=y/\|\phi\|^2$ the exponent equals $-\frac{y^2}{\|\phi\|^2}+\frac{y^2}{2\|\phi\|^2}=-\frac{y^2}{2\|\phi\|^2}$.

(ii) If $\phi=0$, the set $\{0>y\}$ is empty.
If $\phi\ne0$, the set $\left\{\left|\langle\phi,F^{\mathrm{u}}\rangle\right|>y\right\}$ is contained in the union of $\left\{\langle\phi,F^{\mathrm{u}}\rangle\ge y\right\}$ and $\left\{\langle-\phi,F^{\mathrm{u}}\rangle\ge y\right\}$.
By (i), applied to $\phi$ and to $-\phi$, its measure is at most
\[
2\exp\left(-\frac{y^2}{2\|\phi\|^2}\right)\le2\exp\left(-\frac{y^2}{2}\right),
\]
because $\|\phi\|\le1$.

For the measure $\gamma_N$ and $\phi\in F_N$ one has
\[
\langle\phi,F^{\mathrm{u}}\rangle-\frac{1}{2}\|F^{\mathrm{u}}\|^2=\frac{1}{2}\|\phi\|^2-\frac{1}{2}\|F^{\mathrm{u}}-\phi\|^2,
\]
and $\mathrm{lbF}$ is invariant under the translation by $\phi$.
Therefore $\int\mathrm{d}\gamma_N(F^{\mathrm{u}})\,\exp\langle\phi,F^{\mathrm{u}}\rangle=\exp\left(\frac{1}{2}\|\phi\|^2\right)$.
\end{proof}

\subsubsection*{1.2. Proof of (a)}

\begin{proof}[Proof of Proposition~31.1\,(a)]
In Steps A1--A3, $\mathrm{kc}\in\{2,4\}$, and $a$ is a real one-cochain on $\mathbb{Z}^3/n\mathbb{Z}^3$.

\textbf{A1 (the hops).}
For $i=0,\dots,3$ and $f\in\mathfrak{S}^{\mathrm{ch}}_{\mathrm{kc}}(n)$ let $\left(R_i(a)f\right)(x):=U_a\left(x,x^{(i)}\right)\cdot f\left(x^{(i)}\right)$, as in (I.2.2).
At $x=(c,g)$ this reads
\begin{align*}
\left(R_i(a)f\right)(c,g)&=f(c,gs_i)\qquad(i=0,1,2),\\
\left(R_3(a)f\right)(c,g)&=\exp\left(\mathrm{i}q\cdot a(c,c+ge_1)\right)\cdot f(c+ge_1,gs_3).
\end{align*}
For $i\le2$ this is a hop within the cluster of $c$, without phase; for $i=3$ it is a hop to the cluster of the nearest neighbour $c+ge_1$, with the phase of one link.
\begin{itemize}
\item[(i)] $R_i(a)f\in\mathfrak{S}^{\mathrm{ch}}_{\mathrm{kc}}(n)$.
Indeed, by (s1) and (31.1),
\[
\left(R_i(a)f\right)(x+ne_\mu)=U_a\left(x,x^{(i)}\right)\cdot f\left(x^{(i)}+ne_\mu\right)=-\left(R_i(a)f\right)(x).
\]
\item[(ii)] $R_i(a)^2=1$.
Indeed, by (s1) and (31.1),
\[
\left(R_i(a)^2f\right)(x)=U_a\left(x,x^{(i)}\right)\cdot U_a\left(x^{(i)},\left(x^{(i)}\right)^{(i)}\right)\cdot f\left(\left(x^{(i)}\right)^{(i)}\right)=U_a\left(x,x^{(i)}\right)\cdot U_a\left(x^{(i)},x\right)\cdot f(x)=f(x).
\]
\item[(iii)] $R_i(a)$ is isometric.
Indeed, by (31.1), $a$ being real, and since $x\mapsto x^{(i)}$ is a bijection of $T^{\mathrm{f}}_0(n)$ (1.1.1, 1.1.2),
\[
\left(R_i(a)f',R_i(a)f\right)_{\mathrm{kc}}=\frac{1}{48}\sum_{x\in T^{\mathrm{f}}_0(n)}\left|U_a\left(x,x^{(i)}\right)\right|^2\cdot\overline{f'\left(x^{(i)}\right)}\cdot f\left(x^{(i)}\right)=\frac{1}{48}\sum_{x\in T^{\mathrm{f}}_0(n)}\overline{f'(x)}\cdot f(x)=(f',f)_{\mathrm{kc}}.
\]
\end{itemize}
By (iii), $R_i(a)^{*}R_i(a)=1$; together with (ii) this gives
\[
R_i(a)^{*}=R_i(a)^{-1}=R_i(a),\qquad\left|R_i(a)f\right|=|f|.
\]
The assertions (i) and (ii) do not use that $a$ is real.

\textbf{A2 (the Hermitian part).}
On $\mathfrak{S}^{\mathrm{ch}}_{\mathrm{kc}}(n)$ let
\[
\nabla_a^{*}\nabla_a:=\sum_{i=0}^{3}w_i\cdot\left(1-R_i(a)\right).
\]
By A1, $\left(\nabla_a^{*}\nabla_a\right)^{*}=\nabla_a^{*}\nabla_a$ and $\left(1-R_i(a)\right)^2=2\left(1-R_i(a)\right)$, so that $\left(f,\left(1-R_i(a)\right)f\right)_{\mathrm{kc}}=\frac{1}{2}\left|\left(1-R_i(a)\right)f\right|^2$.
Hence
\[
\left(f,\nabla_a^{*}\nabla_af\right)_{\mathrm{kc}}=\frac{1}{2}\sum_{i=0}^{3}w_i\cdot\frac{1}{48}\sum_{x\in T^{\mathrm{f}}_0(n)}\left|U_a\left(x,x^{(i)}\right)\cdot f\left(x^{(i)}\right)-f(x)\right|_2^2\ \ge0,
\]
because $w_i>0$.
In the inner sum the terms of $x$ and of $x^{(i)}$ are equal, because
\[
U_a\left(x^{(i)},x\right)\cdot f(x)-f\left(x^{(i)}\right)=-U_a\left(x^{(i)},x\right)\cdot\left(U_a\left(x,x^{(i)}\right)\cdot f\left(x^{(i)}\right)-f(x)\right)
\]
and $\left|U_a\left(x^{(i)},x\right)\right|=1$.
So the right side is the sum, over the unoriented bonds $\left\{x,x^{(i)}\right\}$, of $\frac{1}{48}w_i\cdot\left|U_a\left(x,x^{(i)}\right)\cdot f\left(x^{(i)}\right)-f(x)\right|_2^2$; this is the form given in (I.O2).

\textbf{A3 (the anti-Hermitian part).}
Let $\left(\mathrm{Gd}_if\right)(x):=\mathrm{Gv}_{\mathrm{kc}}\left(\hat d_i(x)\right)f(x)$.
The vector $\hat d_i(x)=-g\hat n_i$ is real, has Euclidean length one and depends on the label $g$ only.
Hence $\mathrm{Gd}_i$ maps $\mathfrak{S}^{\mathrm{ch}}_{\mathrm{kc}}(n)$ into itself, and by 1.1.2 one has $\mathrm{Gd}_i^{*}=\mathrm{Gd}_i$ and $\left|\mathrm{Gd}_if\right|=|f|$.
By (s2),
\[
\left(R_i(a)\mathrm{Gd}_if\right)(x)=U_a\left(x,x^{(i)}\right)\cdot\mathrm{Gv}_{\mathrm{kc}}\left(\hat d_i\left(x^{(i)}\right)\right)f\left(x^{(i)}\right)=-\mathrm{Gv}_{\mathrm{kc}}\left(\hat d_i(x)\right)U_a\left(x,x^{(i)}\right)\cdot f\left(x^{(i)}\right)=-\left(\mathrm{Gd}_iR_i(a)f\right)(x).
\]
So $\left(\mathrm{Gd}_iR_i(a)\right)^{*}=R_i(a)^{*}\mathrm{Gd}_i^{*}=R_i(a)\mathrm{Gd}_i=-\mathrm{Gd}_iR_i(a)$.
Therefore the operator
\[
J_a:=\sum_{i=0}^{3}w_i\cdot\mathrm{Gd}_iR_i(a),\qquad\text{that is,}\qquad\left(J_af\right)(x)=\sum_iw_i\cdot\mathrm{Gv}_{\mathrm{kc}}\left(\hat d_i(x)\right)U_a\left(x,x^{(i)}\right)\cdot f\left(x^{(i)}\right),
\]
satisfies $J_a^{*}=-J_a$.
For $\mathrm{kc}=4$ the operators $\nabla_a^{*}\nabla_a$ and $J_a$ are those of (I.2.4); for $\mathrm{kc}=2$ they are denoted by $\mathrm{Lw2}_a$ and $\mathrm{Ghs}_a$.

\textbf{A4 (two forms of the operator).}
Let $\mathrm{kc}=4$, and let $B$ be a complex one-cochain on $\mathbb{Z}^3/n\mathbb{Z}^3$.
Let $R_i(B)$, $\nabla_B^{*}\nabla_B$ and $J_B$ be the operators given by the formulas of A1--A3 (the assertion A1\,(i) holds for complex $B$), and let $\left(\mathrm{Mo}_if\right)(x):=M_i(x)f(x)$.
Since $M_i(x)$ depends on the label $g$ only, $\mathrm{Mo}_i$ maps $\mathfrak{S}^{\mathrm{ch}}(1;n)$ into itself.
Inserting $M_i(x)=\frac{w_i}{2}\cdot\left(\rho-\gamma\cdot\hat d_i(x)\right)$ in (31.2) one obtains
\[
(D_Bf)(x)=\frac{\rho}{2}\sum_iw_i\cdot\left[f(x)-\left(R_i(B)f\right)(x)\right]+\frac{1}{2}\sum_iw_i\left(\gamma\cdot\hat d_i(x)\right)\left(R_i(B)f\right)(x)-\frac{1}{2}\left(\gamma\cdot\sum_iw_i\hat d_i(x)\right)f(x).
\]
The last term vanishes by (s3).
Therefore
\begin{equation}\tag{31.3}
D_B=\frac{\rho}{2}\cdot\nabla_B^{*}\nabla_B+\frac{1}{2}\cdot J_B.
\end{equation}
By (s3) also
\[
\sum_iM_i(x)=\frac{\rho}{2}\cdot\mathrm{sw}-\frac{1}{2}\left(\gamma\cdot\sum_iw_i\hat d_i(x)\right)=\frac{\rho}{2}\cdot\mathrm{sw}.
\]
So (31.2) can be written, for every number $z$, as
\begin{equation}\tag{31.4}
D_B+z=\left(z+\frac{\rho}{2}\cdot\mathrm{sw}\right)-\sum_{i=0}^{3}\mathrm{Mo}_iR_i(B).
\end{equation}
The identities (31.3) and (31.4) are (I.2.4) and (I.2.5).

\textbf{A5 (real part).}
Let $a$ and $z$ be real.
By (31.3), and by A2 and A3 at $\mathrm{kc}=4$, one has $\left(D_a+z\right)^{*}=\frac{\rho}{2}\cdot\nabla_a^{*}\nabla_a-\frac{1}{2}\cdot J_a+z$.
Hence
\[
\operatorname{Re}\left(D_a+z\right)=z+\frac{\rho}{2}\cdot\nabla_a^{*}\nabla_a\ge z\qquad\text{for}\ (\cdot,\cdot)^{\mathrm{ch}},
\]
because $\rho>0$ and $\nabla_a^{*}\nabla_a\ge0$ (Step~A2).
This is the first assertion of (a).

\textbf{A6 (blocks; sign of the determinant).}
Let $a$ and $z$ be real.
Write $f\in\mathfrak{S}^{\mathrm{ch}}(1;n)$ as $f=(\mathrm{fup},\mathrm{flo})$ with $\mathrm{fup}:=(f_1,f_2)$ and $\mathrm{flo}:=(f_3,f_4)$.
Then $\mathrm{fup}$ and $\mathrm{flo}$ range independently over $\mathfrak{S}^{\mathrm{ch}}_2(n)$, and $(f',f)^{\mathrm{ch}}=(\mathrm{fup}',\mathrm{fup})_2+(\mathrm{flo}',\mathrm{flo})_2$ for $f'=(\mathrm{fup}',\mathrm{flo}')$.
The subspaces $\{\mathrm{flo}=0\}$ and $\{\mathrm{fup}=0\}$ are the eigenspaces of $\Sigma_3=\operatorname{diag}(1,1,-1,-1)$, acting at every site, for the eigenvalues $1$ and $-1$; each of them has the dimension $96n^3$.
The operator $R_i(a)$ acts on $\mathrm{fup}$ and on $\mathrm{flo}$ by the formula of A1, and $\mathrm{Gv}_4(v)=\operatorname{diag}\left(\mathrm{Gv}_2(v),-\mathrm{Gv}_2(v)\right)$ by 1.1.2.
Hence
\[
\nabla_a^{*}\nabla_a(\mathrm{fup},\mathrm{flo})=\left(\mathrm{Lw2}_a\mathrm{fup},\mathrm{Lw2}_a\mathrm{flo}\right),\qquad J_a(\mathrm{fup},\mathrm{flo})=\left(\mathrm{Ghs}_a\mathrm{fup},-\mathrm{Ghs}_a\mathrm{flo}\right),
\]
and by (31.3)
\begin{gather*}
\left(D_a+z\right)(\mathrm{fup},\mathrm{flo})=\left(\left(D_a^{\boldsymbol{\sigma}}+z\right)\mathrm{fup},\left(D_a^{-\boldsymbol{\sigma}}+z\right)\mathrm{flo}\right),\\
D_a^{\boldsymbol{\sigma}}:=\frac{\rho}{2}\cdot\mathrm{Lw2}_a+\frac{1}{2}\cdot\mathrm{Ghs}_a,\qquad D_a^{-\boldsymbol{\sigma}}:=\frac{\rho}{2}\cdot\mathrm{Lw2}_a-\frac{1}{2}\cdot\mathrm{Ghs}_a.
\end{gather*}
By A2 and A3 at $\mathrm{kc}=2$, and because $\rho$ and $z$ are real,
\[
\left(D_a^{\boldsymbol{\sigma}}+z\right)^{*}=\frac{\rho}{2}\cdot\mathrm{Lw2}_a-\frac{1}{2}\cdot\mathrm{Ghs}_a+z=D_a^{-\boldsymbol{\sigma}}+z.
\]
The determinant of a direct sum of two linear maps is the product of their determinants \cite[Sections~0.3 and 0.9]{HornJohnson2013}.
For a linear map $O$ of a finite-dimensional space with a scalar product one has $\det\left(O^{*}\right)=\overline{\det O}$: in an orthonormal basis the matrix of $O^{*}$ is the transpose of the complex conjugate of the matrix of $O$, a determinant does not change under transposition, and a determinant is a polynomial with integer coefficients in the entries.
Hence
\[
\det\left(D_a+z\right)=\det\left(D_a^{\boldsymbol{\sigma}}+z\right)\cdot\overline{\det\left(D_a^{\boldsymbol{\sigma}}+z\right)}=\left|\det\left(D_a^{\boldsymbol{\sigma}}+z\right)\right|^2\ge0.
\]

\textbf{A7 ($z>0$).}
Let $a$ be real and $z>0$.
If $\left(D_a+z\right)f=0$, then by A5
\[
0=\operatorname{Re}\left(f,\left(D_a+z\right)f\right)^{\mathrm{ch}}\ge z\cdot|f|^2,
\]
so that $f=0$.
Thus $D_a+z$ is injective on a space of finite dimension, $\det\left(D_a+z\right)\ne0$, and $\det\left(D_a+z\right)>0$ by A6.
At $a=0$ this is the third assertion of (a).

\textbf{A8 (bound).}
Let $B$ and $z$ be real.
The matrix $\gamma\cdot\hat d_i(x)$ is Hermitian and its square is $1_4$ (1.1.2); so its eigenvalues lie in $\{1,-1\}$, and the eigenvalues of the Hermitian matrix $M_i(x)$ lie in $\left\{\frac{w_i}{2}\cdot(\rho-1),\frac{w_i}{2}\cdot(\rho+1)\right\}$.
The norm $|\cdot|_2$ of a Hermitian matrix (Section~I.1.5: the operator norm for the Euclidean norm of $\mathbb{C}^4$) is the largest modulus of its eigenvalues, by the spectral theorem \cite[Section~2.5]{HornJohnson2013}, and $|\rho-1|\le1+\rho$.
Hence $\left|M_i(x)\right|_2\le\frac{w_i}{2}\cdot(1+\rho)$, and therefore $\left|\mathrm{Mo}_if\right|\le\frac{w_i}{2}\cdot(1+\rho)\cdot|f|$.
Since $\left|R_i(B)f\right|=|f|$ (Step~A1), (31.4) gives
\[
\left|\left(D_B+z\right)f\right|\le\mathrm{bdD}(z)\cdot|f|,\qquad\mathrm{bdD}(z):=\left|z+\frac{\rho}{2}\cdot\mathrm{sw}\right|+\frac{1+\rho}{2}\cdot\mathrm{sw}.
\]
Let $O$ be a linear map of a space of the finite dimension $\mathrm{dimS}$ with a scalar product.
By Step~A6, $|\det O|^2=\det\left(O^{*}\right)\cdot\det O=\det\left(O^{*}O\right)$, and this is the product of the $\mathrm{dimS}$ eigenvalues of the Hermitian map $O^{*}O$.
An eigenvalue with the eigenvector $f$ equals $|Of|^2/|f|^2$, a number in $\left[0,|O|^2\right]$.
Hence $|\det O|\le|O|^{\mathrm{dimS}}$.
Together with A6 this gives
\begin{equation}\tag{31.5}
0\le\det\left(D_B+z\right)\le\mathrm{bdD}(z)^{192n^3}\qquad\text{for every real one-cochain}\ B\ \text{on}\ \mathbb{Z}^3/n\mathbb{Z}^3.
\end{equation}

\textbf{A9 (continuity).}
Let $\mathrm{Fdm}(n):=\{0,\dots,n-1\}^3\times O_{\mathrm{cub}}$.
For $\mathrm{xo}\in\mathrm{Fdm}(n)$ and $\beta\in\{1,\dots,4\}$ let $\delta_{\mathrm{xo},\beta}$ be the element of $\mathfrak{S}^{\mathrm{ch}}(1;n)$ whose value at $\mathrm{xo}$ is the $\beta$-th standard basis vector of $\mathbb{C}^4$ and which vanishes at the other points of $\mathrm{Fdm}(n)$.
These $192n^3$ elements form a basis of $\mathfrak{S}^{\mathrm{ch}}(1;n)$ (1.1.2).
For $x\in\mathrm{Fdm}(n)$ and $i\in\{0,\dots,3\}$ write $x^{(i)}=\check x^{(i)}+n\cdot\mathrm{vz}_i(x)$ with $\check x^{(i)}\in\mathrm{Fdm}(n)$ and $\mathrm{vz}_i(x)\in\mathbb{Z}^3$; both are unique, and $\mathrm{vz}_i(x)=0$ for $i\le2$.
Then $f\left(x^{(i)}\right)=(-1)^{|\mathrm{vz}_i(x)|_1}f\left(\check x^{(i)}\right)$.
By (31.4) the matrix of $D_B+z$ in this basis consists of $4\times4$ blocks indexed by $\mathrm{Fdm}(n)\times\mathrm{Fdm}(n)$.
It is the sum, over all $x\in\mathrm{Fdm}(n)$, of the block $\left(z+\frac{\rho}{2}\cdot\mathrm{sw}\right)\cdot1_4$ at the position $(x,x)$ and of the four blocks $-(-1)^{|\mathrm{vz}_i(x)|_1}U_B\left(x,x^{(i)}\right)\cdot M_i(x)$ at the positions $\left(x,\check x^{(i)}\right)$, $i=0,\dots,3$.
Each $U_B\left(x,x^{(i)}\right)$ equals $1$ or $\exp\left(\pm\mathrm{i}q\cdot B(b)\right)$ for one link $b$ of the torus.
So every entry of the matrix is a finite sum of constants and of constants times numbers $\exp\left(\pm\mathrm{i}q\cdot B(b)\right)$, and the determinant is a polynomial in the entries.
Consequently the map $B\mapsto\det\left(D_B+z\right)$ is continuous on $C^1(n)$, and it is entire on $C^1(n)\otimes\mathbb{C}$.

The map $(\theta,F^{\mathrm{u}})\mapsto A(F^{\mathrm{u}})+H_\theta+a$ from $\mathbb{R}^3\times F_N$ to $C^1(n)$ is affine: $A(\cdot)$ is linear (1.1.5), and $H_\theta(c,c+e_\mu)=\theta_\mu/(q\cdot n)$ is linear in $\theta$ (Section~I.6.6).
Therefore the function $(\theta,F^{\mathrm{u}})\mapsto\det\left(D_{A(F^{\mathrm{u}})+H_\theta+a}+z\right)$ is continuous on $\mathbb{R}^3\times F_N$.

\textbf{A10 (density and measure).}
Let $a$ be real, $\theta\in\mathbb{R}^3$ and $F^{\mathrm{u}}\in F_N$.
Then $B:=A(F^{\mathrm{u}})+H_\theta+a$ is a real one-cochain on $\mathbb{Z}^3/n\mathbb{Z}^3$, the constant cochain $H_\theta$ being periodic, and $\boldsymbol{\mu}=\bar{m}\cdot\eta+m_{\mathrm{in}}$ is real.
Since $\bar{m}\cdot\eta>0$, Step~A7 gives $\det\left(D_0+\bar{m}\cdot\eta\right)>0$.
By (I.7.2) and by (31.5) at $z=\boldsymbol{\mu}$,
\begin{equation}\tag{31.6}
\begin{aligned}
&0\le\mathrm{dens}^{\theta,a}_N(F^{\mathrm{u}})\le\mathrm{dsup}:=\frac{\exp\left(-\varepsilon^{\mathrm{v}}_{\mathrm{in}}\cdot n^3\right)\cdot\mathrm{bdD}(\boldsymbol{\mu})^{192n^3}}{\det\left(D_0+\bar{m}\cdot\eta\right)}<\infty,\\
&\mathrm{bdD}(\boldsymbol{\mu})=\left|M_{\mathrm{ch}}\right|+\frac{1+\rho}{2}\cdot\mathrm{sw}.
\end{aligned}
\end{equation}
The number $\mathrm{dsup}$ does not depend on $(\theta,F^{\mathrm{u}},a)$.
By A9 the function $(\theta,F^{\mathrm{u}})\mapsto\mathrm{dens}^{\theta,a}_N(F^{\mathrm{u}})$ is continuous.
For a Borel subset $\mathrm{Ev}$ of $F_N$, the number
\[
\nu^{\theta,a}(\mathrm{Ev})=\int_{\mathrm{Ev}}\mathrm{d}\gamma_N\,\mathrm{dens}^{\theta,a}_N
\]
is the integral of a bounded Borel function $\ge0$ with respect to the probability measure $\gamma_N$.
Hence $\nu^{\theta,a}$ is countably additive \cite[Theorem~1.29]{Rudin1987}, and its values lie in $[0,\mathrm{dsup}]$.
So $\nu^{\theta,a}$ is a finite positive measure on $F_N$.
This proves (a).
\end{proof}

\subsubsection*{1.3. Proof of (b)}

\begin{proof}[Proof of Proposition~31.1\,(b)]
\textbf{B1 (the measure $\bar{\nu}$).}
Step~A10 at $a=0$ shows that $(\theta,F^{\mathrm{u}})\mapsto\mathrm{dens}^{\theta,0}_N(F^{\mathrm{u}})$ is a continuous function on $\mathbb{R}^3\times F_N$ with values in $[0,\mathrm{dsup}]$.
Let $\mathrm{Ev}$ be a Borel subset of $F_N$.
The function $(\theta,F^{\mathrm{u}})\mapsto\mathbf{1}_{\mathrm{Ev}}(F^{\mathrm{u}})\cdot\mathrm{dens}^{\theta,0}_N(F^{\mathrm{u}})$ is then Borel and non-negative on $[-\pi,\pi]^3\times F_N$.
By Tonelli's theorem \cite[Theorem~8.8(a)]{Rudin1987}, the function $\theta\mapsto\nu^{\theta}(\mathrm{Ev})$ is Borel on $[-\pi,\pi]^3$, and
\begin{gather*}
\bar{\nu}(\mathrm{Ev})=\frac{1}{(2\pi)^3}\int_{[-\pi,\pi]^3}\mathrm{d}\theta\,\nu^{\theta}(\mathrm{Ev})=\int_{\mathrm{Ev}}\mathrm{d}\gamma_N(F^{\mathrm{u}})\,\mathrm{densb}(F^{\mathrm{u}}),\\
\mathrm{densb}(F^{\mathrm{u}}):=\frac{1}{(2\pi)^3}\int_{[-\pi,\pi]^3}\mathrm{d}\theta\,\mathrm{dens}^{\theta,0}_N(F^{\mathrm{u}})\in[0,\mathrm{dsup}].
\end{gather*}
Consequently $\bar{\nu}=\mathrm{densb}\cdot\gamma_N$ is a non-negative measure on $F_N$ \cite[Theorem~1.29]{Rudin1987}, and $\bar{\nu}(1)=\bar{\nu}(F_N)\le\mathrm{dsup}$.

\textbf{B2 (integrability).}
Let $G$ be a Borel function on $C^2(n)$ with $|G(F^{\mathrm{u}})|\le\mathrm{cf}\cdot\exp\left(\mathrm{cf}\cdot\|F^{\mathrm{u}}\|\right)$.
The number $\bar{\nu}(G)$ is the integral, with respect to $\bar{\nu}$, of the restriction of $G$ to $F_N$.
From B1 and Lemma~1 one obtains
\[
\int\mathrm{d}\bar{\nu}\,|G|\le\mathrm{dsup}\cdot\mathrm{cf}\cdot\exp\left(\mathrm{cf}^2\right)\cdot2^{\frac{\mathrm{dimF}}{2}}<\infty,
\]
and the same number bounds
\[
\frac{1}{(2\pi)^3}\int_{[-\pi,\pi]^3}\mathrm{d}\theta\int\mathrm{d}\gamma_N\,|G|\cdot\mathrm{dens}^{\theta,0}_N .
\]
Fubini's theorem \cite[Theorem~8.8(b), (c)]{Rudin1987} and (I.7.2) at $a=0$ therefore give
\begin{equation}
\begin{split}
\bar{\nu}(G)={}&\exp\left(-\varepsilon^{\mathrm{v}}_{\mathrm{in}}\cdot n^3\right)\cdot\det\left(D_0+\bar{m}\cdot\eta\right)^{-1}\\
&\cdot\frac{1}{(2\pi)^3}\int_{[-\pi,\pi]^3}\mathrm{d}\theta\int_{F_N}\mathrm{d}\gamma_N(F^{\mathrm{u}})\,G(F^{\mathrm{u}})\cdot\det\left(D_{A(F^{\mathrm{u}})+H_{\theta}}+\boldsymbol{\mu}\right).
\end{split}
\tag{31.7}
\end{equation}

\textbf{B3 (from $F_N$ to ${V_{\mathrm{t}}}$).}
Fix $\theta$.
By Step~A9 the function $F^{\mathrm{u}}\mapsto G(F^{\mathrm{u}})\cdot\det\left(D_{A(F^{\mathrm{u}})+H_{\theta}}+\boldsymbol{\mu}\right)$ is Borel on $F_N$.
By (31.5) its modulus does not exceed $\mathrm{cf}'\cdot\exp\left(\mathrm{cf}'\cdot\|F^{\mathrm{u}}\|\right)$, where
\[
\mathrm{cf}':=\mathrm{cf}\cdot\max\left(1,\mathrm{bdD}(\boldsymbol{\mu})^{192n^3}\right).
\]
Lemma~1, together with the identity $A(dA)=A$ for $A\in{V_{\mathrm{t}}}$ (1.1.5), yields
\begin{equation}
\begin{split}
&\int_{F_N}\mathrm{d}\gamma_N(F^{\mathrm{u}})\,G(F^{\mathrm{u}})\cdot\det\left(D_{A(F^{\mathrm{u}})+H_{\theta}}+\boldsymbol{\mu}\right)\\
&\qquad=\mathrm{zV}_n^{-1}\cdot\int_{{V_{\mathrm{t}}}}\mathrm{d}\sigma(A)\,G(dA)\cdot\det\left(D_{A+H_{\theta}}+\boldsymbol{\mu}\right)\cdot\exp\left(-\frac{1}{2}\|dA\|^2\right).
\end{split}
\tag{31.8}
\end{equation}

\textbf{B4 (the zero-flux sector).}
Apply Proposition~38.2\,(b) on the torus $(n,n,n)$, $n\ge4$, with the real number $\boldsymbol{\mu}$ and with $\Phi=G$, a Borel function with $|G|\le\mathrm{cf}\cdot\exp\left(\mathrm{cf}\cdot\|F^{\mathrm{u}}\|\right)$.
It states that both sides of
\begin{equation}
Z^{(0)}_n[G]=\mathrm{Nc}_n\cdot\frac{1}{(2\pi)^3}\int_{[-\pi,\pi]^3}\mathrm{d}\theta\int_{{V_{\mathrm{t}}}}\mathrm{d}\sigma(A)\,G(dA)\cdot\exp\left(-\frac{1}{2}\|dA\|^2\right)\cdot\det\left(D_{A+H_{\theta}}+\boldsymbol{\mu}\right)
\tag{31.9}
\end{equation}
converge absolutely and that they coincide.
On the left side this says that the sum over the closed $\mathfrak{z}\in C^2(n;\mathbb{Z})$ of flux $0$ and the integrals over $[0,r_q)^{\text{links}}$ of (I.7.1) and Section~I.7.3 converge absolutely.
Hence $Z^{(0)}_n[G]$ is a complex number; this applies in particular where it occurs below, in (31.10), C5 and D1.

The absolute convergence can also be read off from the equality alone, as follows.
The function $|G|$ is Borel and obeys the bound of $G$.
Every term of $Z^{(0)}_n\left[|G|\right]$ is the integral of a non-negative function, since $\det\left(D_A+\boldsymbol{\mu}\right)\ge0$ for the real $A$ of the box (Step~A6).
So the left side of (31.9) for $|G|$ is defined as an element of $[0,\infty]$.
By (31.5) and Lemma~1 the right side of (31.9) for $|G|$ is at most
\[
\begin{aligned}
&\mathrm{Nc}_n\cdot\mathrm{bdD}(\boldsymbol{\mu})^{192n^3}\cdot\int_{{V_{\mathrm{t}}}}\mathrm{d}\sigma(A)\,|G(dA)|\cdot\exp\left(-\frac{1}{2}\|dA\|^2\right)=\mathrm{Nc}_n\cdot\mathrm{bdD}(\boldsymbol{\mu})^{192n^3}\cdot\mathrm{zV}_n\cdot\int\mathrm{d}\gamma_N\,|G|\\
&\qquad\le\mathrm{Nc}_n\cdot\mathrm{bdD}(\boldsymbol{\mu})^{192n^3}\cdot\mathrm{zV}_n\cdot\mathrm{cf}\cdot\exp\left(\mathrm{cf}^2\right)\cdot2^{\frac{\mathrm{dimF}}{2}}<\infty .
\end{aligned}
\]
It follows that $Z^{(0)}_n\left[|G|\right]<\infty$, and this is the absolute convergence of the sum and of the integrals in $Z^{(0)}_n[G]$.

Inserting (31.8) into (31.7), and then (31.9), one finds
\[
\bar{\nu}(G)=\exp\left(-\varepsilon^{\mathrm{v}}_{\mathrm{in}}\cdot n^3\right)\cdot\det\left(D_0+\bar{m}\cdot\eta\right)^{-1}\cdot\mathrm{zV}_n^{-1}\cdot\mathrm{Nc}_n^{-1}\cdot Z^{(0)}_n[G].
\]
The three inverted numbers belong to $(0,\infty)$.
For $\det\left(D_0+\bar{m}\cdot\eta\right)$ this is Step~A7, and for $\mathrm{zV}_n$ it is Lemma~1.
For $\mathrm{Nc}_n=r_q^{n^3+2}\cdot\left(\operatorname{det}'\Delta_0\right)^{\frac{1}{2}}$ it holds because $r_q>0$ and $\operatorname{det}'\Delta_0$ is the product of the non-zero eigenvalues of the non-negative operator $d^{*}d$ on $C^0(n)$, all of which are positive.
The last display is the first display of Proposition~31.1\,(b).

\textbf{B5 ($G=1$).}
The definition of $\mathrm{zr}_n$ (Section~I.7.3) and (31.7) at $G=1$ give
\[
\mathrm{zr}_n(q,\boldsymbol{\mu})=\frac{1}{(2\pi)^3}\int_{[-\pi,\pi]^3}\mathrm{d}\theta\int\mathrm{d}\gamma_N(F^{\mathrm{u}})\,\det\left(D_{A(F^{\mathrm{u}})+H_{\theta}}+\boldsymbol{\mu}\right)=\exp\left(\varepsilon^{\mathrm{v}}_{\mathrm{in}}\cdot n^3\right)\cdot\det\left(D_0+\bar{m}\cdot\eta\right)\cdot\bar{\nu}(1).
\]
By the first display of Proposition~31.1\,(b) at $G=1$, the last member equals $Z^{(0)}_n[1]/\left(\mathrm{Nc}_n\cdot\mathrm{zV}_n\right)$.
This is the second display of Proposition~31.1\,(b).
\end{proof}

\subsubsection*{1.4. Proof of (c)}

\begin{proof}[Proof of Proposition~31.1\,(c)]
\textbf{C1 (finiteness).}
By B1 and (31.6), $\bar{\nu}(1)\le\mathrm{dsup}<\infty$.

\textbf{C2 ($Z^{(0)}_n[1]>0$).}
By (I.7.1) and Section~I.7.3,
\[
Z^{(0)}_n[1]=\sum_{\mathfrak{z}}\mathrm{Tm}(\mathfrak{z}),\qquad\mathrm{Tm}(\mathfrak{z}):=\int_{[0,r_q)^{\text{links}}}\mathrm{d}A\,\exp\left(-\frac{1}{2}\left\|dA+r_q\cdot\mathfrak{z}\right\|^2\right)\cdot\det\left(D_A+\boldsymbol{\mu}\right),
\]
where $\mathfrak{z}$ ranges over the elements of $C^2(n;\mathbb{Z})$ with $d\mathfrak{z}=0$ and flux $0$.
Each $\mathrm{Tm}(\mathfrak{z})$ is non-negative, by Step~A6.
The cochain $\mathfrak{z}=0$ is closed and has flux $0$.
By Proposition~46\,(b), (e), $\det\left(D_A+\boldsymbol{\mu}\right)$ is the determinant $\det D^{T}_A$ of \cite[Definition~2.11]{QEDII} for the torus $T=C(n,n,n)$ and the parameters $(\rho,w,M=M_{\mathrm{ch}},q)$ of \cite[Definition~2.10]{QEDII}.
The hypotheses of \cite[Theorem~4.13]{QEDII} are satisfied.
Indeed, $\rho$ is real with $\rho^2\ne1$, as $0<\rho<1$; the parameter $M$ of \cite[Definition~2.10]{QEDII} is real, being $M_{\mathrm{ch}}=\boldsymbol{\mu}+\frac{\rho}{2}\cdot\mathrm{sw}$; $w_i>0$; $q>0$; and the period lattice $n\mathbb{Z}^3$ is a subgroup of $\mathbb{Z}^3$ of rank three, carrying a unitary character, namely the one that takes the value $-1$ at $ne_1$, $ne_2$ and $ne_3$.
At $\mathfrak{z}=0$, \cite[Theorem~4.13(d)]{QEDII} gives $\mathrm{Tm}(0)>0$.
Hence $Z^{(0)}_n[1]\ge\mathrm{Tm}(0)>0$.

\textbf{C3 (positivity; the probability measure).}
By the first display of Proposition~31.1\,(b) at $G=1$, the number $\bar{\nu}(1)$ is the product of $Z^{(0)}_n[1]>0$ (Step~C2) and of
\[
\exp\left(-\varepsilon^{\mathrm{v}}_{\mathrm{in}}\cdot n^3\right)\cdot\left[\det\left(D_0+\bar{m}\cdot\eta\right)\cdot\mathrm{Nc}_n\cdot\mathrm{zV}_n\right]^{-1},
\]
a number in $(0,\infty)$ (Step~B4).
With C1 it follows that $0<\bar{\nu}(1)<\infty$, and also that $Z^{(0)}_n[1]<\infty$.
Thus $\bar{\mu}:=\bar{\nu}/\bar{\nu}(1)$ is a non-negative measure on $F_N$ with $\bar{\mu}(F_N)=1$.

\textbf{C4 (the law under $E^{\mathrm{av}}_n$).}
For every $G$ as in (b) the constant of (b) cancels in the quotient:
\begin{equation}
\bar{\mu}(G)=\frac{\bar{\nu}(G)}{\bar{\nu}(1)}=\frac{Z^{(0)}_n[G]}{Z^{(0)}_n[1]}=E^{\mathrm{av}}_n(G).
\tag{31.10}
\end{equation}
For $G=\mathbf{1}_{\mathrm{Ev}}$, where $\mathrm{Ev}$ is a Borel subset of $C^2(n)$, this becomes
\[
E^{\mathrm{av}}_n\left(\mathbf{1}\left[F^{\mathrm{u}}\in\mathrm{Ev}\right]\right)=\bar{\mu}\left(\mathrm{Ev}\cap F_N\right).
\]
So the law of $F^{\mathrm{u}}$ under $E^{\mathrm{av}}_n$, a Borel probability measure on $C^2(n)$, is the image of $\bar{\mu}$ under the inclusion of $F_N$ into $C^2(n)$.

\textbf{C5 ($E^{\mathrm{av}}_n=E^{\mathrm{af}}_n(\,\cdot\mid\text{flux}=0)$).}
\cite[Theorem~4.26]{QEDII} is stated for real $\rho$, real values of the parameter $M$ of \cite[Definition~2.10]{QEDII}, $w_i>0$, $q>0$ and every torus in the sense of \cite[Definition~2.11]{QEDII}; the torus $C(n,n,n)$ with $n\ge4$ is such a torus.
Part~(a) of that theorem and Proposition~46\,(f) at $\Phi=1$ give $0<Z_n[1]<\infty$.
Hence $E^{\mathrm{af}}_n=Z_n[\,\cdot\,]/Z_n[1]$ is defined (Section~I.7.2), and $E^{\mathrm{af}}_n\left(\mathbf{1}[\text{flux}=0]\right)=Z^{(0)}_n[1]/Z_n[1]>0$ by C2.
For $G$ as in (b) it follows that
\[
E^{\mathrm{af}}_n(G\mid\text{flux}=0)=\frac{E^{\mathrm{af}}_n\left(\mathbf{1}[\text{flux}=0]\cdot G\right)}{E^{\mathrm{af}}_n\left(\mathbf{1}[\text{flux}=0]\right)}=\frac{Z^{(0)}_n[G]}{Z^{(0)}_n[1]}=E^{\mathrm{av}}_n(G).
\]

\textbf{C6 (what $\bar{\mu}$ depends on).}
By B1, (I.7.2) and B5, the factor $\exp\left(-\varepsilon^{\mathrm{v}}_{\mathrm{in}}\cdot n^3\right)/\det\left(D_0+\bar{m}\cdot\eta\right)$ appears both in $\bar{\nu}$ and in $\bar{\nu}(1)$ and cancels in the quotient:
\[
\mathrm{d}\bar{\mu}(F^{\mathrm{u}})=\mathrm{zr}_n(q,\boldsymbol{\mu})^{-1}\cdot\left[\frac{1}{(2\pi)^3}\int_{[-\pi,\pi]^3}\mathrm{d}\theta\,\det\left(D_{A(F^{\mathrm{u}})+H_{\theta}}+\boldsymbol{\mu}\right)\right]\mathrm{d}\gamma_N(F^{\mathrm{u}}).
\]
The right side is built from $F_N$, $\gamma_N$ and $A(\cdot)$, which are determined by $n$; from $H_\theta$ and the operator (31.2), which are determined by $n$, $q$ and the fixed $\rho$ and $w$; and from the number $\boldsymbol{\mu}$.
Neither $\varepsilon^{\mathrm{v}}_{\mathrm{in}}$ nor the way in which $\boldsymbol{\mu}$ is split into $\bar{m}\cdot\eta$ and $m_{\mathrm{in}}$ enters it.
\end{proof}

\subsubsection*{1.5. Proof of (d)}

\begin{proof}[Proof of Proposition~31.1\,(d)]
\textbf{D1 (reduction to $F_N$).}
Let $\mathrm{gc}$ be a real two-cochain on $\mathbb{Z}^3/n\mathbb{Z}^3$, and put $\phi:=\mathrm{PrF}\,\mathrm{gc}\in F_N$.
Since $\mathrm{gc}-\phi$ is orthogonal to $F_N$, one has $\langle\mathrm{gc},F^{\mathrm{u}}\rangle=\langle\phi,F^{\mathrm{u}}\rangle$ for $F^{\mathrm{u}}\in F_N$.
As $\bar{\mu}$ is a measure on $F_N$, this gives $\bar{\mu}\left(\exp\langle\mathrm{gc},F^{\mathrm{u}}\rangle\right)=\bar{\mu}\left(\exp\langle\phi,F^{\mathrm{u}}\rangle\right)$.
The function $G_{\phi}(u):=\exp\langle\phi,u\rangle$, $u\in C^2(n)$, is continuous, and
\[
0<G_{\phi}(u)\le\exp\left(\|\phi\|\cdot\|u\|\right)\le\mathrm{cf}\cdot\exp\left(\mathrm{cf}\cdot\|u\|\right),\qquad\mathrm{cf}:=\max\left(1,\|\phi\|\right).
\]
By (31.10),
\[
\bar{\mu}\left(\exp\langle\mathrm{gc},F^{\mathrm{u}}\rangle\right)=\frac{Z^{(0)}_n[G_{\phi}]}{Z^{(0)}_n[1]}.
\]

\textbf{D2 (Gaussian domination on the torus).}

\emph{The dictionary.}
Let $T:=C(n,n,n)$, with the parameters $(\rho,w,M=M_{\mathrm{ch}},q)$ of \cite[Definition~2.10]{QEDII} and with an arbitrary $v_0\in\mathrm{Chm}_{0}$.
By \cite[Definition~13.7(a)]{QEDII}, $Z^{(0)}_T[G]:=Z_T[\Phi_0\cdot G]$, where $\Phi_0$ is the function of Lemma~2.
In \cite[Definition~13.7(a)]{QEDII} this function is the indicator $\Phi_{k'}$ of the flux sector $k'=0$, formed with one coordinate plane $\{c_{\mathfrak{l}}=\mathrm{const}\}$ for each pair $\mu<\nu$; by Lemma~2 the choice of these planes does not affect $\Phi_0(F^{\mathrm{u}})$, because $F^{\mathrm{u}}$ is closed.
The curvature $F$ of \cite[Definition~2.6]{QEDII} is $F^{\mathrm{u}}$.
The functions $\Phi_0\cdot G_{\phi}$ and $\Phi_0$ are Borel, non-negative and bounded by $\mathrm{cf}\cdot\exp\left(\mathrm{cf}\cdot\|u\|\right)$.
Hence $Z_T$ of \cite[Definition~2.13]{QEDII} is defined on them, both through the class~(b) of that definition and through its definition on non-negative Borel functions of the curvature, and the two definitions give the same number.
Proposition~46\,(f), applicable since $n\ge4$, and Lemma~2 give
\[
Z^{(0)}_T[G_{\phi}]=Z_T[\Phi_0\cdot G_{\phi}]=Z_n[\Phi_0\cdot G_{\phi}]=Z^{(0)}_n[G_{\phi}],\qquad Z^{(0)}_T[1]=Z_T[\Phi_0]=Z_n[\Phi_0]=Z^{(0)}_n[1].
\]

\emph{The statement on the torus.}
\cite[Proposition~13.11]{QEDII} asserts that for every coordinate torus $C(n_1,n_2,n_3)$ and every two-cochain $\mathrm{hz}$ on it
\[
0<Z^{(0)}\left[\exp\left(-\langle\mathrm{hz},F^{\mathrm{u}}\rangle-\frac{1}{2}\|\mathrm{hz}\|^2\right)\right]\le Z^{(0)}[1],
\]
where $Z^{(0)}$ is the functional of this torus.
Its hypotheses are those of \cite[Section~13.3]{QEDII}, read with \cite[Definitions~2.6, 2.10 and 2.11]{QEDII}: $w_i>0$, the parameter $M$ of \cite[Definition~2.10]{QEDII} real, $q>0$, $-1\le\rho\le1$, and $n_\mu\ge4$.
Here $0<\rho<1$, the weights are the standard weights, the parameter is the real number $M_{\mathrm{ch}}$, and $n\ge4$.
No condition is imposed on the sign or the modulus of the parameter, and none on the parity of the $n_\mu$:
in the proof of \cite[Proposition~13.11]{QEDII} the side $n_1$ enters only through the number of the factors of a trace and through the common exponent in H\"older's inequality, both for the trace and for the integral over the two Bloch angles.

\emph{The second time direction by a rotation.}
In the proof of \cite[Proposition~13.11]{QEDII} the inequality with $e_1$ as time is followed by a second inequality, which is derived from \cite[Lemma~13.8(c)]{QEDII}, with $e_2$ as time.
The second inequality can also be reduced to the first one by a rotation of the torus, through \cite[Lemma~6.26]{QEDII};
in this form \cite[Lemma~13.8]{QEDII} is needed with $e_1$ as time only.
The reduction is as follows.
In this paragraph $T$ is an arbitrary coordinate torus $C(n_1,n_2,n_3)$ with $n_\mu\ge4$, $\Phi_0$ is the function of \cite[Definition~13.7(a)]{QEDII} on $T$, and for a two-cochain $\mathrm{hz}$ on $T$
\[
\mathrm{Z0}_T(\mathrm{hz}):=Z^{(0)}_T\left[\exp\left(-\langle\mathrm{hz},F^{\mathrm{u}}\rangle-\frac{1}{2}\|\mathrm{hz}\|^2\right)\right],
\]
the middle term of the assertion of \cite[Proposition~13.11]{QEDII}.
The first half of the proof of \cite[Proposition~13.11]{QEDII}, the part that takes $e_1$ as time, rests on the Bloch representation \cite[Lemma~13.8(b)]{QEDII} with $e_1$ as time and yields
\[
\mathrm{Z0}_T(\mathrm{hz})\le\prod_{t\in\mathbb{Z}/n_1}\mathrm{Z0}_T\left(\mathrm{hz}^{(t)}\right)^{\frac{1}{n_1}}.
\]
Here $\mathrm{hz}^{(t)}$ is the two-cochain whose $(23)$-components in every plane $\{c_1=\mathrm{const}\}$ are the $(23)$-components of $\mathrm{hz}$ in the plane $\{c_1=t\}$, and whose $(12)$- and $(13)$-components vanish.
In particular $\mathrm{Z0}_T(\mathrm{hz})\le\mathrm{Z0}_T(0)$ whenever $\mathrm{hz}$ has no $(23)$-components.
The second half requires the bound $\mathrm{Z0}_T\left(\mathrm{hz}^{(t)}\right)\le\mathrm{Z0}_T(0)$, which follows from the first half by a rotation.
Let $g_{\mathrm{w}}\in W_{\mathrm{af}}$ be the rotation $y\mapsto ry$, where $r\in O_{\mathrm{cub}}$ is the permutation matrix with $re_1=e_3$, $re_2=e_1$, $re_3=e_2$.
Thus $r$ is the linear part of $g_{\mathrm{w}}$, and $g_{\mathrm{w}}^{-1}$ is the map $y\mapsto r^{-1}y$, where $r^{-1}$ is the cyclic permutation matrix with $r^{-1}e_1=e_2$, $r^{-1}e_2=e_3$, $r^{-1}e_3=e_1$.
Let $T':=g_{\mathrm{w}}T=C(n_2,n_3,n_1)$ be the image torus of \cite[Definition~2.11]{QEDII}; for the cubic torus $T'=T$.
For a two-cochain $\mathrm{hz}'$ on $T$ the function $G:=\Phi_0\cdot\exp\left(-\langle\mathrm{hz}',\cdot\,\rangle-\frac{1}{2}\|\mathrm{hz}'\|^2\right)$ is Borel and non-negative, and \cite[Lemma~6.26]{QEDII}, applied to $g_{\mathrm{w}}$, $T$ and $G$, gives
\[
Z_{T'}\left[G\circ g_{\mathrm{w}}^{*}\right]=Z_T[G]=\mathrm{Z0}_T(\mathrm{hz}').
\]
By \cite[Definition~2.1]{QEDII} the map $g_{\mathrm{w}}^{*}$, which takes the two-cochains on $T'$ to those on $T$, is orthogonal; its inverse, which is also its adjoint, is the pull-back $\left(g_{\mathrm{w}}^{-1}\right)^{*}$ by $g_{\mathrm{w}}^{-1}$.
Therefore $\left\langle\mathrm{hz}',g_{\mathrm{w}}^{*}u'\right\rangle=\left\langle\left(g_{\mathrm{w}}^{-1}\right)^{*}\mathrm{hz}',u'\right\rangle$ for the two-cochains $u'$ on $T'$, and $\left\|\left(g_{\mathrm{w}}^{-1}\right)^{*}\mathrm{hz}'\right\|=\|\mathrm{hz}'\|$.
By \cite[Definition~13.7(e)]{QEDII}, applied to the element $g_{\mathrm{w}}$, one has $(\Phi_0)^{g_{\mathrm{w}}}=\Phi_0$ on the closed two-cochains of $g_{\mathrm{w}}T=T'$.
An even element of $W_{\mathrm{af}}$ acts on functions of the curvature by composition with its pull-back, and an odd element by the same composition followed by complex conjugation \cite[Section~2.3]{QEDII}.
The element $g_{\mathrm{w}}$ is even, since $\det r=1$, and $\Phi_0$ is real; so $(\Phi_0)^{g_{\mathrm{w}}}(u')=\Phi_0\left(g_{\mathrm{w}}^{*}u'\right)$ for the two-cochains $u'$ on $T'$.
Hence $\Phi_0\left(g_{\mathrm{w}}^{*}u'\right)=\Phi_0(u')$ for the closed two-cochains $u'$ on $T'$, where $\Phi_0$ on the left is the function of $T$ and $\Phi_0$ on the right is that of $T'$.
So, at every closed two-cochain $u'$ on $T'$,
\[
\left(G\circ g_{\mathrm{w}}^{*}\right)(u')=\Phi_0(u')\cdot\exp\left(-\left\langle\left(g_{\mathrm{w}}^{-1}\right)^{*}\mathrm{hz}',u'\right\rangle-\frac{1}{2}\left\|\left(g_{\mathrm{w}}^{-1}\right)^{*}\mathrm{hz}'\right\|^2\right).
\]
Moreover $Z_{T'}$ evaluates its argument only at curvatures, and these are closed.
Consequently
\[
\mathrm{Z0}_T(\mathrm{hz}')=\mathrm{Z0}_{T'}\left(\left(g_{\mathrm{w}}^{-1}\right)^{*}\mathrm{hz}'\right)\quad\text{for every two-cochain}\ \mathrm{hz}'\ \text{on}\ T,\qquad\text{in particular}\quad\mathrm{Z0}_T(0)=\mathrm{Z0}_{T'}(0).
\]
Now take $\mathrm{hz}'=\mathrm{hz}^{(t)}$.
By \cite[Definition~2.1]{QEDII}, $\left(\left(g_{\mathrm{w}}^{-1}\right)^{*}\mathrm{hz}^{(t)}\right)(\mathrm{pl})=\pm\mathrm{hz}^{(t)}\left(g_{\mathrm{w}}^{-1}\mathrm{pl}\right)$ for the plaquettes $\mathrm{pl}$ of $T'$, and $g_{\mathrm{w}}^{-1}$ maps the plaquettes of the directions $(12)$, $(13)$, $(23)$ onto plaquettes of the directions $(23)$, $(12)$, $(13)$, respectively.
Hence $\left(g_{\mathrm{w}}^{-1}\right)^{*}\mathrm{hz}^{(t)}$ has $(12)$-components only, and the first half, applied on $T'$, gives $\mathrm{Z0}_{T'}\left(\left(g_{\mathrm{w}}^{-1}\right)^{*}\mathrm{hz}^{(t)}\right)\le\mathrm{Z0}_{T'}(0)$.
Combining,
\[
\mathrm{Z0}_T\left(\mathrm{hz}^{(t)}\right)=\mathrm{Z0}_{T'}\left(\left(g_{\mathrm{w}}^{-1}\right)^{*}\mathrm{hz}^{(t)}\right)\le\mathrm{Z0}_{T'}(0)=\mathrm{Z0}_T(0)
\]
for every $t$, and therefore $\mathrm{Z0}_T(\mathrm{hz})\le\mathrm{Z0}_T(0)=Z^{(0)}_T[1]$.

\emph{Application.}
Let again $T=C(n,n,n)$.
The two-cochains on $T$ are the elements of $C^2(n)$, with the counting pairing \cite[Definitions~2.1 and 2.6]{QEDII}.
\cite[Proposition~13.11]{QEDII} at $\mathrm{hz}:=-\phi$ reads
\[
Z^{(0)}_T\left[\exp\left(\langle\phi,F^{\mathrm{u}}\rangle-\frac{1}{2}\|\phi\|^2\right)\right]\le Z^{(0)}_T[1],\qquad\text{that is,}\qquad Z^{(0)}_n[G_{\phi}]\le\exp\left(\frac{1}{2}\|\phi\|^2\right)\cdot Z^{(0)}_n[1].
\]

\textbf{D3 (the first display of (d)).}
Dividing by $Z^{(0)}_n[1]>0$ (Step~C2) and using D1, one obtains
\[
\bar{\mu}\left(\exp\langle\mathrm{gc},F^{\mathrm{u}}\rangle\right)\le\exp\left(\frac{1}{2}\|\phi\|^2\right)=\exp\left(\frac{1}{2}\left\|\mathrm{PrF}\,\mathrm{gc}\right\|^2\right).
\]

\textbf{D4 (the second and the third display of (d)).}
For $\phi\in F_N$ one has $\mathrm{PrF}\,\phi=\phi$.
So, by D3, the probability measure $\bar{\mu}$ satisfies the hypothesis of Lemma~3.
Lemma~3\,(i) is the second display of (d), and Lemma~3\,(ii) is the third.
\end{proof}

\subsubsection*{1.6. The last paragraph of the statement}

\begin{itemize}
\item[(i)] \emph{No constant.}
Of the constants of Section~I.8.2, only $\rho$ and the weights (line (C0)) occur in 1.2--1.5.
All identities and inequalities of (a)--(d) hold for all data $(\bar{m},e,N,n,m_{\mathrm{in}},\varepsilon^{\mathrm{v}}_{\mathrm{in}})$, and the inequalities of (d) hold with the constant one for every $n\ge4$, every $q>0$ and every real $\boldsymbol{\mu}$.
\item[(ii)] \emph{Every real mass} (condition (I.K10)).
The number $\boldsymbol{\mu}$ is inserted, in the place of $z$, in Steps~A6 (the sign of the determinant), A8 (that is, (31.5)) and A9 (continuity), each of which is stated for every real $z$, as is Step~A5.
It enters further in Proposition~38.2\,(b), which holds for every real $\boldsymbol{\mu}$, and in \cite[Theorems~4.13(d) and 4.26, Proposition~13.11]{QEDII}, which hold for every real value of the parameter $M$ of \cite[Definition~2.10]{QEDII}, here $M_{\mathrm{ch}}=\boldsymbol{\mu}+\frac{\rho}{2}\cdot\mathrm{sw}$.
Step~A7 is stated for $z>0$ only, and it is applied only at $a=0$ and $z=\bar{m}\cdot\eta$ (Steps~A10 and B4).
This is the only positivity of a mass that enters: $\bar{m}\cdot\eta>0$, in the denominator $\det\left(D_0+\bar{m}\cdot\eta\right)$ (Step~A10).
\item[(iv)] \emph{The number $Z^{\mathrm{ph}}_n$.}
In Proposition~38.2\,(c) the number $Z^{\mathrm{ph}}_n$ is defined as $\mathrm{Nc}_n\cdot\int_{{V_{\mathrm{t}}}}\mathrm{d}\sigma(A)\,\exp\left(-\frac{1}{2}\|dA\|^2\right)$, which is $\mathrm{Nc}_n\cdot\mathrm{zV}_n$ by the definition of $\mathrm{zV}_n$.
In \cite[Section~21.8.1 and Lemma~21.32]{QEDI} the same letter denotes
\[
\int_{C^1(n)/\Gamma_n}\mathrm{d}A\,\exp\left(-\frac{1}{2}\|dA\|^2\right),\qquad\Gamma_n:=r_q\cdot\mathrm{Zcl},\qquad\mathrm{Zcl}:=\left\{m\in C^1(n;\mathbb{Z}):\ dm=0\right\},
\]
that is, the integral of the $\Gamma_n$-periodic integrand over a fundamental domain of $\Gamma_n$, with respect to the Lebesgue measure of $C^1(n)$.
Two Borel fundamental domains give the same integral of a non-negative $\Gamma_n$-periodic Borel function: one cuts one of them along the translates of the other and uses the invariance of Lebesgue measure under translations; \cite[Definition~2.6]{QEDII} contains the same fact.
For a lattice $\mathrm{La}$ of full rank in a Euclidean space, $\operatorname{covol}(\mathrm{La})$ denotes the modulus of the determinant of the matrix of a basis of $\mathrm{La}$ with respect to an orthonormal basis of the space.
It is the Lebesgue measure of the parallelepiped spanned by that basis \cite[Theorem~2.20(e) and 2.23]{Rudin1987}, and this parallelepiped is a fundamental domain of $\mathrm{La}$.
The number of \cite[Lemma~21.32]{QEDI} also equals $\mathrm{Nc}_n\cdot\mathrm{zV}_n$; the proof occupies the steps (1)--(6).
\begin{itemize}
\item[(1)] By \cite[Section~10.1.1]{QEDI}, $\ker\left(d\colon C^1(n)\to C^2(n)\right)={V_{\mathrm{g}}}\oplus{V_{\mathrm{h}}}$, so that $\mathrm{Zcl}$ is contained in ${V_{\mathrm{g}}}\oplus{V_{\mathrm{h}}}$.
The Lebesgue measure of $C^1(n)$ is the product of the Lebesgue measures of ${V_{\mathrm{g}}}\oplus{V_{\mathrm{h}}}$ and of ${V_{\mathrm{t}}}$, and the integrand depends only on the component of $A$ in ${V_{\mathrm{t}}}$.
Hence, if $\Gamma_n$ is a lattice of full rank in ${V_{\mathrm{g}}}\oplus{V_{\mathrm{h}}}$, the product of a fundamental domain of $\Gamma_n$ in ${V_{\mathrm{g}}}\oplus{V_{\mathrm{h}}}$ with ${V_{\mathrm{t}}}$ is a fundamental domain in $C^1(n)$, and the integral equals $\operatorname{covol}(\Gamma_n)\cdot\mathrm{zV}_n$.
\item[(2)] One has $\mathrm{Zcl}\cap{V_{\mathrm{g}}}=dC^0(n;\mathbb{Z})$.
Indeed, let $m=d\boldsymbol{\lambda}$ with a real $\boldsymbol{\lambda}$, and let $m$ be integer valued.
Then $\boldsymbol{\lambda}(c')-\boldsymbol{\lambda}(c)$ is an integer for all nearest neighbours $c$, $c'$.
As the torus is connected, $\boldsymbol{\lambda}-\boldsymbol{\lambda}(0)$ is integer valued, and $m=d\left(\boldsymbol{\lambda}-\boldsymbol{\lambda}(0)\right)$.
Conversely, $dd=0$, so that $dC^0(n;\mathbb{Z})$ consists of closed integer cochains, and it lies in ${V_{\mathrm{g}}}$.
\item[(3)] The orthogonal projection of $\mathrm{Zcl}$ onto ${V_{\mathrm{h}}}$ is the set of the constant cochains $\mathrm{cm}$ with $n\cdot\mathrm{cm}_\mu\in\mathbb{Z}$ for $\mu=1,2,3$.
Indeed, let $m\in\mathrm{Zcl}$.
By \cite[Section~10.1.1]{QEDI}, $m=d\boldsymbol{\lambda}+\mathrm{cm}$ with $\mathrm{cm}\in{V_{\mathrm{h}}}$, and $\mathrm{cm}$ is the projection of $m$.
The sum of $m$ along the closed path $c,c+e_\mu,\dots,c+ne_\mu$ is an integer, and it equals $n\cdot\mathrm{cm}_\mu$, because the terms of $d\boldsymbol{\lambda}$ telescope.
Conversely, for $\mu\in\{1,2,3\}$ the one-cochain
\[
m^{(\mu)}(c,c+e_\nu):=\delta_{\mu\nu}\cdot\mathbf{1}\left[c_\mu\equiv n-1\ \mathrm{mod}\ n\right]\qquad(\nu=1,2,3)
\]
is integer valued.
It is closed: its only non-vanishing component depends on $c_\mu$ alone, so that both differences in
\[
\left(dm^{(\mu)}\right)_{\alpha\beta}(c)=m^{(\mu)}_{\alpha}(c)-m^{(\mu)}_{\alpha}(c+e_\beta)+m^{(\mu)}_{\beta}(c+e_\alpha)-m^{(\mu)}_{\beta}(c)
\]
vanish.
Its sums along the three closed paths in the directions $\nu=1,2,3$ are $\delta_{\mu\nu}$.
The three constant cochains with the components $\delta_{\mu\nu}$, $\nu=1,2,3$, are pairwise orthogonal, and each has the norm $n^{\frac{3}{2}}$.
So the projection is a lattice of full rank in ${V_{\mathrm{h}}}$, with the covolume $\left(n^{\frac{3}{2}}/n\right)^3=n^{\frac{3}{2}}$.
\item[(4)] Let $\mathrm{La}$ be a lattice of full rank in an orthogonal sum $\mathrm{Eu}_1\oplus\mathrm{Eu}_2$ of Euclidean spaces such that $\mathrm{La}\cap\mathrm{Eu}_1$ has full rank in $\mathrm{Eu}_1$.
Then the orthogonal projection of $\mathrm{La}$ onto $\mathrm{Eu}_2$ is a lattice of full rank in $\mathrm{Eu}_2$, and
\[
\operatorname{covol}(\mathrm{La})=\operatorname{covol}\left(\mathrm{La}\cap\mathrm{Eu}_1\right)\cdot\operatorname{covol}(\text{projection}).
\]
Indeed, the group $\mathrm{La}/\left(\mathrm{La}\cap\mathrm{Eu}_1\right)$ is torsion free: if a non-zero multiple of an element of $\mathrm{La}$ lies in $\mathrm{Eu}_1$, the element itself lies in $\mathrm{Eu}_1$.
As a quotient of $\mathrm{La}$, this group has a finite system of generators; it is therefore free, and a basis of $\mathrm{La}\cap\mathrm{Eu}_1$ extends to a basis of $\mathrm{La}$.
Here one uses that an abelian group with a finite system of generators and without torsion is free, and that a homomorphism onto a free abelian group splits \cite[Chapter~I, Sections~7 and 8]{Lang2002}.
In orthonormal coordinates adapted to $\mathrm{Eu}_1\oplus\mathrm{Eu}_2$ the matrix of this basis is block triangular.
Its diagonal blocks are the matrix of the basis of $\mathrm{La}\cap\mathrm{Eu}_1$ and the matrix of the projections onto $\mathrm{Eu}_2$ of the remaining basis vectors.
The determinant is the product of the determinants of the two blocks, and it is not zero.
Hence these projections are linearly independent, and the projection of $\mathrm{La}$ is the group generated by them.
\item[(5)] The covolume of $dC^0(n;\mathbb{Z})$ in ${V_{\mathrm{g}}}$ is $\left(\operatorname{det}'\Delta_0\right)^{\frac{1}{2}}\cdot n^{-\frac{3}{2}}$.
Indeed, let $\mathrm{Ko}$ be the orthogonal complement of the constants in $C^0(n)$.
The kernel of $d$ on $C^0(n)$ is the space of the constants; so $d$ maps $\mathrm{Ko}$ bijectively onto ${V_{\mathrm{g}}}$, and in orthonormal bases the modulus of its determinant is $\det\left(d^{*}d\ \text{on}\ \mathrm{Ko}\right)^{\frac{1}{2}}=\left(\operatorname{det}'\Delta_0\right)^{\frac{1}{2}}$.
The group $dC^0(n;\mathbb{Z})$ is the image under $d$ of the orthogonal projection of $C^0(n;\mathbb{Z})=\mathbb{Z}^{n^3}$ onto $\mathrm{Ko}$.
By (4), with $\mathrm{La}=\mathbb{Z}^{n^3}$ (covolume $1$), $\mathrm{Eu}_1$ the constants, $\mathrm{La}\cap\mathrm{Eu}_1=\mathbb{Z}\cdot(1,\dots,1)$ (covolume $n^{\frac{3}{2}}$) and $\mathrm{Eu}_2=\mathrm{Ko}$, this projection has the covolume $n^{-\frac{3}{2}}$.
Finally, the image of a lattice under a linear bijection has the covolume of the lattice times the modulus of the determinant.
\item[(6)] The group $\mathrm{Zcl}$ is a discrete subgroup of ${V_{\mathrm{g}}}\oplus{V_{\mathrm{h}}}$.
By (2), (3) and (5) it contains a lattice of full rank in ${V_{\mathrm{g}}}$, and its projection spans ${V_{\mathrm{h}}}$.
So $\mathrm{Zcl}$ is a lattice of full rank in ${V_{\mathrm{g}}}\oplus{V_{\mathrm{h}}}$, a space of the dimension $(n^3-1)+3$.
By (4) with $\mathrm{Eu}_1={V_{\mathrm{g}}}$ and $\mathrm{Eu}_2={V_{\mathrm{h}}}$, and by (2), (3) and (5),
\[
\operatorname{covol}(\mathrm{Zcl})=\left(\operatorname{det}'\Delta_0\right)^{\frac{1}{2}}\cdot n^{-\frac{3}{2}}\cdot n^{\frac{3}{2}}=\left(\operatorname{det}'\Delta_0\right)^{\frac{1}{2}}.
\]
Hence $\operatorname{covol}(\Gamma_n)=r_q^{n^3+2}\cdot\left(\operatorname{det}'\Delta_0\right)^{\frac{1}{2}}=\mathrm{Nc}_n$, and by (1) the number of \cite[Lemma~21.32]{QEDI} is $\mathrm{Nc}_n\cdot\mathrm{zV}_n$.
\end{itemize}
\end{itemize}

\subsection*{2. Proof of Proposition~31.2}

\subsubsection*{2.1. Scales}

Let (FL) hold (Section~G4.0.1): $\bar{m}\in(0,1]$; $e>0$ with $e\cdot h^{\frac{1}{2}}\le e_1$; $R=h\cdot L^T\in\mathrm{Tor}(e)$; and $N\ge N_{\min}(R,e)$.
Put $\eta=L^{-N}$ and $q=e\cdot\eta^{\frac{1}{2}}$, and for integers $j\ge0$
\begin{gather*}
a_j=L^{j-N},\qquad e_j=e\cdot a_j^{\frac{1}{2}},\qquad v_j=-\log e_j,\qquad\ell_j=\frac{R}{a_j},\qquad n_j=\ell_j^3,\\
r_j=\min\left\{L^i:\ i\ge0\ \text{an integer},\ L^i\ge v_j^r\right\},\qquad p_{0,j}=v_j^{p_0},\qquad p_j=v_j^p,
\end{gather*}
where $\mathrm{lg}=\frac{1}{2}\log L$ and $r$, $p_0$, $p$ are the integers of (C7); for $j=0$ the number $e_j$ is $q$.
Further $K=N+\log_Lh$, so that $a_K=h$ and $e_K=e\cdot h^{\frac{1}{2}}$, and the index $*$ refers to $j=k_{*}=k_{*}(R)$ (Sections~I.6.1 and I.6.2).
\begin{itemize}
\item[(i)] By (Sc1) of Section~G4.0.2, that is by \cite[Proposition~19.1(c)]{QEDI}, $1\le k_{*}\le K$ and $\ell_*=M\cdot r_{*}$.
(The hypotheses $v_K\ge1$ and $\left(1+\frac{\mathrm{lg}}{v_K}\right)^r\le L$ of \cite[Proposition~19.1]{QEDI}, read with $h'=h$ (Section~I.6.1), hold by (Sc6) of Section~G4.0.2: $v_K\ge8\ge1$, and $\left(1+\frac{\mathrm{lg}}{v_K}\right)^r\le\left(1+\frac{3\,\mathrm{lg}}{v_K}\right)^r<L$, because $\mathrm{lg}>0$ and $r>0$.)
Since $M$ is a power of $L$ and $r_{*}$ is a power of $L$ with a non-negative exponent, $\ell_*$ is a positive integer.
For $0\le j\le k_{*}$ one has $\ell_j/\ell_*=a_{*}/a_j=L^{k_{*}-j}$.
Hence
\[
\ell_j=L^{k_{*}-j}\cdot\ell_*,\qquad n_j=\ell_*^3\cdot L^{3(k_{*}-j)},\qquad n=\frac{R}{\eta}=\ell_0=L^{k_{*}}\cdot\ell_*,
\]
and $\ell_j$ is a multiple of $L$ for $j<k_{*}$.
\item[(ii)] By the definition of $N_{\min}(R,e)$ (Section~I.6.2), $n=R\cdot L^N\ge L\cdot M\cdot r_0$.
The number $r_0$ is a power $L^i$ with $i\ge0$; so $r_0\ge1$, and $L\cdot M\cdot r_0\ge L\cdot M\ge4$, because $L\ge5$ and $M\ge L^2$ (Section~I.0.5).
Moreover $n=h\cdot L^{T+N}$, and $h=L^{\lfloor\log_L\frac{1}{\bar{m}}\rfloor}$ is a power of $L$ with a non-negative exponent.
As $L$ is odd, $n$ is odd.
Hence $(\bar{m},e,N,n)$, together with the real pair $(m_{\mathrm{in}},\varepsilon^{\mathrm{v}}_{\mathrm{in}})$, are data of Proposition~31.1, and $\bar{\mu}$ is the probability measure on $F_N=dC^1(n)$ of Proposition~31.1\,(c).
\item[(iii)] For $0\le j\le k_{*}$,
\[
e_j=e_{*}\cdot\left(\frac{a_j}{a_{*}}\right)^{\frac{1}{2}}=e_{*}\cdot L^{-\frac{k_{*}-j}{2}},\qquad\text{hence}\qquad v_j=v_{*}+(k_{*}-j)\cdot\mathrm{lg}.
\]
\item[(iv)] By (Sc6) the conditions of Section~G4.0.3 hold under (FL).
Of (T1), the proof of Proposition~31.2 needs the first two clauses:
\[
v_{*}\ge8,\qquad p_0\cdot v_{*}^{p_0-1}\ge28.
\]
\end{itemize}

\subsubsection*{2.2. The functionals}

The objects introduced below are the photon objects of Section~G4.0.6 and Section~I.5.7; their definitions are those of \cite[Section~5.2.5]{QEDI} and \cite[Definition~5.11]{QEDI}.
Of all objects of Section~2, the measure $\bar{\mu}$ is the only one that is formed with the fermion operator, and it enters through Proposition~31.1\,(d).

For $0\le j\le k_{*}$ put $T^0_j:=\mathbb{Z}^3/\ell_j\mathbb{Z}^3$; in particular $T^0_0=\mathbb{Z}^3/n\mathbb{Z}^3$, since $\ell_0=n$ by 2.1\,(i).
For $0\le j<k_{*}$, the space $V_j$ consists of the real bond functions $\mathsf{Z}$ on $T^0_j$ with $\tau\mathsf{Z}=0$ that vanish on the central bonds; it carries the counting scalar product.
Its basis $\{\Upsilon\}$ is the one fixed in \cite[Section~5.2.5]{QEDI}, and $\mathsf{Z}(\Upsilon):=\langle\mathsf{Z},\Upsilon\rangle$.
Let $F^{\mathrm{u}}\in F_N$.
Then $A_0=A_0(F^{\mathrm{u}})$ denotes the potential of $F^{\mathrm{u}}$ in the slice $\mathrm{Sl}$, the bond function $A_j:=(Q_jA_0)^{[j]}$ lives on $T^0_j$, and $A_{*}:=A_{k_{*}}$.
For $j<k_{*}$ the element $\zeta_j\in V_j$ is defined by
\[
A_j=A^{\min}_{j}(A_{j+1})+\mathrm{Sq}_{j}\zeta_j .
\]
The law of $A_{*}$ under $\gamma_N$ is denoted by $\mathrm{d}\mathsf{m}_{*}$, and $\chi_{*}(A):=\mathbf{1}\left[\,|dA|_{\infty}\le p_{*}\,\right]$.
The small-field set is
\[
\mathfrak{G}=\left\{F^{\mathrm{u}}\in F_N:\ |\zeta_j(\Upsilon)|\le p_{0,j}\text{ for all }\Upsilon\text{ and all }j<k_{*};\ |dA_{*}|_{\infty}\le p_{*}\right\}.
\]
Here $|dA_{*}|_{\infty}$ is the largest of the numbers $\left|(dA_{*})(\mathrm{pl})\right|$, where $\mathrm{pl}=(y;\mu\nu)$ ranges over the plaquettes with $y\in T^0_{k_{*}}$ and $\mu<\nu$; there are $3\ell_*^3$ such plaquettes.
A plaquette is denoted by $\mathrm{pl}$, the letter $p$ being the integer of (C7).

\emph{The setting of the hierarchical coordinates.}
The results of \cite[Section~14.5.1]{QEDI} and \cite[Section~14.5.2]{QEDI} are formulated under the following hypotheses: $L$ is an odd integer with $L\ge3$ \cite[Section~14.1.1]{QEDI}, $k\ge1$ is an integer, $n$ is a positive multiple of $L^k$, and $\ell_j:=\frac{n}{L^j}$.
By 2.1\,(i) these hypotheses are met with $k=k_{*}$, and the numbers $\ell_j$ are then the sides of 2.1.
At the level $j$, the setting of \cite[Section~15.6]{QEDI} requires a side that is a multiple of $L$; by 2.1\,(i) this requirement is met for each $j<k_{*}$.
Moreover, for $j<k_{*}$,
\[
\ell_j=L^{k_{*}-j}\cdot M\cdot r_{*}\ge L^3\ge3L,
\]
because $M\ge L^2$ and $r_{*}\ge1$.
In \cite[Section~5.2.5]{QEDI} and in \cite[Section~14.4.1]{QEDI} the interior bonds (called inner bonds in the latter), the faces, the face bonds and the number of the elements of the basis $\{\Upsilon\}$ are described in the same way, for every side $\ell\in L\mathbb{N}$: a bond function on the torus of side $\ell$ is read as an $\ell\mathbb{Z}^3$-periodic bond function on $\mathbb{Z}^3$, the $L$-blocks, the bonds and the faces are those of $\mathbb{Z}^3$, counted modulo $\ell\mathbb{Z}^3$, and the basis has $d_V\cdot\ell^3$ elements, the dimension of the fluctuation space by \cite[Lemma~14.13(c)]{QEDI}.
For the side $\ell_j$ this number is $d_V\cdot n_j=\dim V_j$; \cite[Section~15.6]{QEDI} prescribes the periodic lift for the side $L$.
Hence the two descriptions agree for every side in $L\mathbb{N}$, the sides $L$ and $2L$ included.
For every side $\ell\ge3L$ the $L$-blocks, the bonds and the faces so defined, hence also the interior bonds and the face bonds, are moreover, modulo $\ell\mathbb{Z}^3$, those formed on the torus itself \cite[Section~5.2.5]{QEDI}; by the last display this is the case for the torus $T^0_j$ of side $\ell_j$, for every $j<k_{*}$.
The objects are taken on $\mathbb{Z}^3$ modulo $\ell\mathbb{Z}^3$, and not on the torus, because of the sides $L$ and $2L$: on the torus of side $L$ there is a single block, and the bonds that wrap around the torus have both ends in it; on the torus of side $2L$ two neighbouring blocks are joined by two faces.
Neither of these two sides occurs here.

\emph{The objects of \cite[Definition~5.11]{QEDI} are the coordinates of \cite[Section~14.5.1]{QEDI} and \cite[Section~14.5.2]{QEDI}.}
Let $A\in C^1(T^0_{k_{*}})$ and $\zeta=(\zeta_j)_{j<k_{*}}$ with $\zeta_j\in V_j$.
Put
\[
\mathrm{Af}_{k_{*}}:=A,\qquad\mathrm{Af}_j:=A^{\min}_{j}(\mathrm{Af}_{j+1})+\mathrm{Sq}_{j}\zeta_j\quad(j=k_{*}-1,\dots,0),
\]
and let $\mathrm{Af}_0(A,\zeta)$ denote the field reached at the end of this recursion (\cite[Definition~5.11]{QEDI}; \cite[Section~14.5.1]{QEDI}).
According to \cite[Section~14.5.2]{QEDI}, the slice $\mathrm{Sl}$ of \cite[Definition~5.11]{QEDI} is the set
\[
\mathrm{Sl}=\left\{A_0\in C^1(n):\ \tau A_j=0\ \ (0\le j<k_{*}),\ \ A_{k_{*}}\in{V_{\mathrm{t}}}(T^0_{k_{*}})\right\},\qquad A_j=(Q_jA_0)^{[j]},
\]
where ${V_{\mathrm{t}}}(T^0_{k_{*}})$ is the space of the coclosed bond functions on $T^0_{k_{*}}$ with mean zero.
The coboundary $d$ maps $\mathrm{Sl}$ isomorphically onto $F_N$ \cite[Theorem~14.17(a)]{QEDI}.
Therefore each $F^{\mathrm{u}}\in F_N$ has exactly one potential $A_0(F^{\mathrm{u}})$ in $\mathrm{Sl}$, and $A_0(F^{\mathrm{u}})$ is linear in $F^{\mathrm{u}}$.
By \cite[Theorem~14.16(a)]{QEDI}, every $A_0$ with $\tau A_j=0$ for all $j<k_{*}$ equals $\mathrm{Af}_0(A,\zeta)$ for exactly one pair $(A,\zeta)$, and for this pair one has $A=A_{k_{*}}$ and $\mathrm{Sq}_{j}\zeta_j=A_j-A^{\min}_{j}(A_{j+1})$ for $j<k_{*}$.
Since $\mathrm{Sq}_{j}$ is injective on $V_j$ \cite[Proposition~14.14(b)]{QEDI}, the equation $A_j=A^{\min}_{j}(A_{j+1})+\mathrm{Sq}_{j}\zeta_j$ of \cite[Definition~5.11]{QEDI} possesses exactly one solution $\zeta_j\in V_j$.
Consequently the coordinates $(A_{*};\zeta_0,\dots,\zeta_{k_{*}-1})$ of $F^{\mathrm{u}}$ in \cite[Theorem~14.17(b)]{QEDI}, to which the assertions of \cite[Theorem~14.17(c), (d)]{QEDI} refer, are the field $A_{k_{*}}$ and the elements $\zeta_j$ of \cite[Definition~5.11]{QEDI}.
Each of them is a linear function of $F^{\mathrm{u}}$.
\begin{itemize}
\item[(I1)] By \cite[Theorem~14.17(c)]{QEDI} there is an orthonormal family $\{\phi_{j,\Upsilon}:\ j<k_{*},\ \Upsilon\}$ in $F_N$ such that $\zeta_j(\Upsilon)=\langle\phi_{j,\Upsilon},F^{\mathrm{u}}\rangle$ for all $F^{\mathrm{u}}\in F_N$.
In particular $\|\phi_{j,\Upsilon}\|=1$.
\item[(I2)] By \cite[Theorem~14.17(d)]{QEDI}, to every plaquette $\mathrm{pl}$ of $T^0_{k_{*}}$ there corresponds an element $\phi_{\mathrm{pl}}\in F_N$ with $\|\phi_{\mathrm{pl}}\|\le1$ and $(dA_{*})(\mathrm{pl})=\langle\phi_{\mathrm{pl}},F^{\mathrm{u}}\rangle$ for all $F^{\mathrm{u}}\in F_N$.
\item[(I3)] By \cite[Proposition~15.21(ii)]{QEDI}, for $j<k_{*}$ the basis $\{\Upsilon\}$ of $V_j$ has $d_V\cdot n_j$ elements, where $d_V=2-2L^{-3}<2$.
\end{itemize}
From (I1), (I2) and the definition of $\mathfrak{G}$ one obtains
\begin{equation}
\begin{split}
F_N\setminus\mathfrak{G}={}&\bigcup_{j<k_{*}}\bigcup_{\Upsilon}\left\{\left|\langle\phi_{j,\Upsilon},F^{\mathrm{u}}\rangle\right|>p_{0,j}\right\}\\
&\cup\bigcup_{\mathrm{pl}}\left\{\left|\langle\phi_{\mathrm{pl}},F^{\mathrm{u}}\rangle\right|>p_{*}\right\},
\end{split}
\tag{31.11}
\end{equation}
a finite union of open subsets of $F_N$.

\subsubsection*{2.3. The first bound of (9.2)}

\begin{proof}[Proof of the first inequality of (9.2)]
\textbf{Step 1} (union bound).
The thresholds $p_{0,j}$ and $p_{*}$ are nonnegative, because $v_j\ge v_{*}\ge8$ by 2.1\,(iii), (iv).
All the elements $\phi_{j,\Upsilon}$ and $\phi_{\mathrm{pl}}$ belong to $F_N$ and have norm at most one.
The third display of Proposition~31.1\,(d), applied with $\phi=\phi_{j,\Upsilon}$, $y=p_{0,j}$ and with $\phi=\phi_{\mathrm{pl}}$, $y=p_{*}$, together with (31.11), the subadditivity of $\bar{\mu}$, the count of (I3) and the count $3\ell_*^3$ of the plaquettes, gives
\begin{equation}
\begin{split}
\bar{\mu}\left(F_N\setminus\mathfrak{G}\right)&\le\sum_{j=0}^{k_{*}-1}d_V\cdot n_j\cdot2\exp\left(-\frac{p_{0,j}^2}{2}\right)+3\ell_*^3\cdot2\exp\left(-\frac{p_{*}^2}{2}\right)\\
&\le\sum_{j=0}^{k_{*}-1}4\,n_j\cdot\exp\left(-\frac{p_{0,j}^2}{2}\right)+6\,\ell_*^3\cdot\exp\left(-\frac{p_{*}^2}{2}\right).
\end{split}
\tag{31.12}
\end{equation}

\textbf{Step 2} (one level).
Fix $j<k_{*}$ and set $i:=k_{*}-j$; this is an integer with $1\le i\le k_{*}$.
By 2.1\,(i),
\[
n_j=\ell_*^3\cdot L^{3i}=\ell_*^3\cdot\exp\left(6\,i\cdot\mathrm{lg}\right),
\]
and by 2.1\,(iii), $p_{0,j}=\left(v_{*}+i\cdot\mathrm{lg}\right)^{p_0}$.
For real numbers $\mathrm{ya},\mathrm{yb}\ge0$ and an integer $p_0\ge1$, every term of the binomial expansion of $(\mathrm{ya}+\mathrm{yb})^{p_0}$ is nonnegative, and retaining the first two terms yields
\[
(\mathrm{ya}+\mathrm{yb})^{p_0}\ge\mathrm{ya}^{p_0}+p_0\cdot\mathrm{ya}^{p_0-1}\cdot\mathrm{yb}.
\]
With $\mathrm{ya}=v_{*}$ and $\mathrm{yb}=i\cdot\mathrm{lg}$, both nonnegative, and with 2.1\,(iv) it follows that
\[
p_{0,j}\ge v_{*}^{p_0}+p_0\cdot v_{*}^{p_0-1}\cdot i\cdot\mathrm{lg}\ge p_{0*}+28\,i\cdot\mathrm{lg}.
\]
As $v_j\ge8$, one has $p_{0,j}\ge1$, and therefore $p_{0,j}^2\ge p_{0,j}$.
Hence
\begin{align*}
n_j\cdot\exp\left(-\frac{p_{0,j}^2}{2}\right)&\le\ell_*^3\cdot\exp\left(6\,i\cdot\mathrm{lg}\right)\cdot\exp\left(-\frac{p_{0,j}}{2}\right)\\
&\le\ell_*^3\cdot\exp\left(-\frac{p_{0*}}{2}\right)\cdot\exp\left((6-14)\,i\cdot\mathrm{lg}\right)=\ell_*^3\cdot\exp\left(-\frac{p_{0*}}{2}\right)\cdot L^{-4i}.
\end{align*}

\textbf{Step 3} (sum).
Since $L\ge5$,
\[
\sum_{i=1}^{k_{*}}L^{-4i}\le\sum_{i\ge1}L^{-4i}=\frac{1}{L^4-1}\le1.
\]
Inserting the bound of Step~2 into (31.12), one finds
\[
\bar{\mu}\left(F_N\setminus\mathfrak{G}\right)\le4\,\ell_*^3\cdot\exp\left(-\frac{p_{0*}}{2}\right)+6\,\ell_*^3\cdot\exp\left(-\frac{p_{*}^2}{2}\right).
\]
By Proposition~14.2\,(f) one has $C_0\ge1$, whence $\exp\left(-\frac{p_{*}^2}{2}\right)\le\exp\left(-\frac{p_{*}^2}{2C_0^2}\right)$.
The right side is therefore at most
\[
6\,\ell_*^3\cdot\exp\left(-\frac{p_{0*}}{2}\right)+6\,\ell_*^3\cdot\exp\left(-\frac{p_{*}^2}{2C_0^2}\right)=\varepsilon_3,
\]
by the definition of $\varepsilon_3$ in Section~G4.0.3 (Proposition~14.2\,(c)).
\end{proof}

\subsubsection*{2.4. The second bound of (9.2)}

\begin{proof}[Proof of the second inequality of (9.2)]
By Section~2.2, $\mathrm{d}\mathsf{m}_{*}$ is the image of $\gamma_N$ under the linear map $F^{\mathrm{u}}\mapsto A_{*}$.
Hence (I2) gives
\begin{align*}
\mathrm{d}\mathsf{m}_{*}\left(\chi_{*}=0\right)&=\gamma_N\left\{F^{\mathrm{u}}\in F_N:\ |dA_{*}|_{\infty}>p_{*}\right\}=\gamma_N\left(\bigcup_{\mathrm{pl}}\left\{\left|\langle\phi_{\mathrm{pl}},F^{\mathrm{u}}\rangle\right|>p_{*}\right\}\right)\\
&\le\sum_{\mathrm{pl}}\gamma_N\left\{\left|\langle\phi_{\mathrm{pl}},F^{\mathrm{u}}\rangle\right|>p_{*}\right\}.
\end{align*}
By Lemma~3 of Section~1.1 the measure $\gamma_N$ satisfies
\[
\int\mathrm{d}\gamma_N(F^{\mathrm{u}})\,\exp\langle\phi,F^{\mathrm{u}}\rangle=\exp\left(\frac{\|\phi\|^2}{2}\right)\qquad(\phi\in F_N),
\]
so that under $\gamma_N$ the variable $\langle\phi_{\mathrm{pl}},F^{\mathrm{u}}\rangle$ is centred Gaussian with variance $\|\phi_{\mathrm{pl}}\|^2\le1$.
Lemma~3\,(ii), applied with $\phi=\phi_{\mathrm{pl}}$ and $y=p_{*}\ge0$, yields
\[
\gamma_N\left\{\left|\langle\phi_{\mathrm{pl}},F^{\mathrm{u}}\rangle\right|>p_{*}\right\}\le2\exp\left(-\frac{p_{*}^2}{2}\right).
\]
As there are $3\ell_*^3$ plaquettes, it follows that
\[
\mathrm{d}\mathsf{m}_{*}\left(\chi_{*}=0\right)\le6\,\ell_*^3\cdot\exp\left(-\frac{p_{*}^2}{2}\right)\le6\,\ell_*^3\cdot\exp\left(-\frac{p_{*}^2}{2C_0^2}\right)=\varepsilon_4;
\]
here $C_0\ge1$ by Proposition~14.2\,(f), and $\varepsilon_4$ is the number defined in Section~G4.0.3.
\end{proof}

\subsubsection*{2.5. What the bounds depend on}

The numbers $\varepsilon_3$ and $\varepsilon_4$ are formed with $\ell_*$, with $v_{*}=-\log e_{*}$, with the integers $p_0$, $p$ and with $C_0$.
Both bounds hold for every $R\in\mathrm{Tor}(e)$, every $N\ge N_{\min}(R,e)$ and every real pair $(m_{\mathrm{in}},\varepsilon^{\mathrm{v}}_{\mathrm{in}})$, with $\varepsilon_3$ and $\varepsilon_4$ formed on the given $R$; indeed, the pair enters the argument only through Proposition~31.1\,(d), and the inequalities of Proposition~31.1\,(d) involve neither the pair nor $q$, $n$, $\boldsymbol{\mu}$.
By (Sc1), for fixed $R$ neither $\ell_*$ nor $e_{*}$ depends on $N\ge N_{\min}(R,e)$.
Hence, for fixed $(\bar{m},e,R)$, the numbers $\varepsilon_3$ and $\varepsilon_4$ depend neither on $N$ nor on the pair.
They depend on $R$ through $\ell_*$ and $e_{*}$, which by (Sc3) satisfy $e_{*}^2\cdot\ell_*=e^2\cdot R$.
The second bound involves no fermion object.

\subsection*{3. Conditions on the constants}
\leavevmode\par\noindent
\begin{tabular}{p{0.26\textwidth}p{0.36\textwidth}p{0.26\textwidth}}
\hline
condition & used in & line of Section~I.8.2 / Proposition~14.2\\
\hline
$\rho>0$; only $\rho\ge0$ is needed & A5 ($\operatorname{Re}D_a=\frac{\rho}{2}\nabla_a^{*}\nabla_a\ge0$); A8 (the modulus of $\rho-1$ is at most $1+\rho$) & (C0): $\rho\in(0,1)$\\
\hline
$\rho$ real, $\rho^2\ne1$ & C2, through \cite[Theorem~4.13(d)]{QEDII} & (C0): $\rho\in(0,1)$\\
\hline
$-1\le\rho\le1$ & D2, through \cite[Proposition~13.11]{QEDII} & (C0): $\rho\in(0,1)$\\
\hline
$w_i>0$ & A2; A8 (the moduli of the eigenvalues $\frac{w_i}{2}(\rho\mp1)$ of $M_{i}(x)$ are at most $\frac{w_i}{2}(1+\rho)$); \cite[Theorem~4.13]{QEDII}, \cite[Theorem~4.26]{QEDII}, \cite[Proposition~13.11]{QEDII} & (C0): the standard weights\\
\hline
$\sum_iw_i\cdot\hat n_i=0$ (I.1.1) & (s3), hence A4: (31.3), (31.4) & (C0): the standard weights\\
\hline
$L$ an odd integer $\ge5$; $M$ a power of $L$, $M\ge L^2$ & 2.1\,(i), (ii) ($\ell_*$ an integer; $L\cdot M\ge4$; $n$ odd); 2.2 ($L$ odd and $L\ge3$, the setting of \cite[Section~14.1.1]{QEDI}; $\ell_j$ a multiple of $L$; $\ell_j\ge L^3\ge3L$ for $j<k_{*}$); 2.3, Steps~1 and~2 ($\mathrm{lg}\ge0$); 2.3, Step~3 ($\frac{1}{L^4-1}\le1$) & Section~I.0.5; (C3), (C5)\\
\hline
$p_0$ an integer $\ge1$ & 2.3, Step~2 (binomial expansion) & (C7): $p_0>3r+2$, $r\ge2$\\
\hline
$C_0\ge1$ & 2.3, Step~3; 2.4 & (C6); Proposition~14.2\,(f): $C_0\ge\max(1,C^{F})$\\
\hline
(T1), first clause: $v_{*}\ge8$; only $v_j\ge0$ and $p_{0,j}\ge1$ are needed & 2.3, Steps~1, 2; 2.4 ($y=p_{*}\ge0$ in Lemma~3\,(ii)) & Proposition~14.2\,(c); it follows from $e\cdot h^{\frac{1}{2}}\le e_1$ and (e1), line (C9), by (Sc6)\\
\hline
(T1), second clause: $p_0\cdot v_{*}^{p_0-1}\ge28$ & 2.3, Step~2 & Proposition~14.2\,(c); from (e1), line (C9), by (Sc6)\\
\hline
$v_K\ge1$, $\left(1+\frac{\mathrm{lg}}{v_K}\right)^r\le L$ (the hypotheses of \cite[Proposition~19.1]{QEDI}) & 2.1\,(i), through (Sc1) & (e1), line (C9), by (Sc6): $v_K\ge8$, $\left(1+\frac{3\,\mathrm{lg}}{v_K}\right)^r<L$; with (C7): $r$ an integer $\ge2$, of which only $r>0$ is needed, and $\mathrm{lg}>0$ ($L\ge5$)\\
\hline
\end{tabular}

The proofs of Propositions~31.1 and~31.2 require no condition on the constants other than those of the lines named in the table.
None of these conditions involves $N$, $\theta$, $n$, the pair $(m_{\mathrm{in}},\varepsilon^{\mathrm{v}}_{\mathrm{in}})$ or the sign of $\boldsymbol{\mu}$.
The proof of Proposition~31.1 uses no condition on, and no property of, a constant of a line of Section~I.8.2 after the line (C0).
Of the constants of these lines, only $L$ (line (C3)) and the data $\bar{m}$ (line (C10)), $e$ and $N$ (line (C11)) occur in the statement of Proposition~31.1, and they enter it only through the positive numbers $\eta=L^{-N}$, $q=e\cdot\eta^{\frac{1}{2}}$ and $\bar{m}\cdot\eta$.

\section*{Chapter 32. Two-sided bound on the partition function: proofs of Propositions 32.1--32.3}

\subsection*{0. Conventions of this chapter; setting; facts for the three proofs}

\subsubsection*{0.0. Notation}

In this chapter $F_\infty(B)$ denotes the operator $D_B+\mu_*+b_{k_*}\,JP_{k_*}(B)J^*$ of Section~I.3.5 at the level $k_*$.
The letters $z$, $z'$ always denote real numbers; unit sites are denoted by $y$, $y'$, and fermion sites by $x$, $x'$.
Accordingly, the variable of the function $\varepsilon_5$ of Section~G4.0.3 is written $z$.
The terminal exponent is denoted by $\mathrm{Etm}$, and its image under the substitution of Step~6 of Section~1 by $\mathrm{Etm}'$.
The letter $t$ is the contour variable of Section~I.5.2; $\mathrm{fD}(t)$ is a quadratic form in the generators that depends on $t$, and $\mathscr{A}^{(t)}:=t\cdot\mathscr{A}^0+c_\theta$.
Further, $\mathrm{Dtw}:=D_*(c_\theta)$, with inverse $\mathrm{Cov}_\theta$; $\mathrm{Osub}$ is a substitution operator and $\mathrm{crw}$ the supremum of its row sums; $\mathrm{Stw}$ is the set on which the function $\mathrm{vrem}$ is defined; $\mathrm{Ety}(B)$ is the exponent in (Y), with the two parts $\mathrm{Ety}_0$, $\mathrm{Ety}_1$; and $\mathscr{A}[\theta,A]$ is the background connection attached to a point $(\theta,A)$ of $\mathrm{Stw}$.
A plaquette of $T^0_{k_*}$ is denoted by $\mathrm{pl}$.

Besides the displays $(32.1)$, $(32.2)$, $\dots$ of this chapter, displays of other Propositions occur:
(9.1) is the display of Proposition~32.1;
(7.1)--(7.4) are the displays of Proposition~32.2;
(10.1)--(10.3) are the displays of Proposition~32.3;
(Y) is the display of Proposition~29;
(8.1) is the display of Proposition~30.1;
(8.5) is the display of Proposition~30.3;
(9.2) is the display of Proposition~31.2.

\subsubsection*{0.1. Data, scales, level 0}

Throughout Sections~0--2 the quadruple $(\bar{m},e,R,N)$ satisfies (FL) of Sections~I.1.8 and~G4.0.1, and $\eta$, $q$, $n$ are the numbers that it determines.
The scales $a_k$, $e_k$, $v_k$, $\mu_k$, $K$, $\ell_k$, $r_k$, $p_k$, $p_{0,k}$ are those of Section~I.6.1.
The stop level $k_*=k_*(R)$ and the index $*$ are those of Section~I.6.2, and $p_{0*}$ is the number of Section~G4.0.3.
The weights $h_k$ of the norms are those of Section~I.5.1, and $h_*:=h_{k_*}$; the windows $w^{\mathrm{m}}_k$, $w^{\mathrm{E}}_k$ are those of Section~I.5.4.
Section~G4.0.2 contains, as (Sc1)--(Sc5), the assertions of \cite[Proposition~19.1(c)--(g)]{QEDI}, read with $h'=h$, $m_0=0$, and with the value $1$ for the upper end $\mathfrak{m}_0$ of the mass range used there.

One has
\[
e_*\le e_K\le e_1\le e_0<1,
\]
by (Sc3), by (FL), by (C9) and by Proposition~14.2\,(f), in this order.
By Proposition~26.6\,(c) the quadruple $(k_*,\ell_*,\mu_*,e_*)$ is (SF), with $k_*\ge1$.
Hence every statement whose data are (SF), or (S-3), or the data of Proposition~19\,(b)--(f), holds at the level $k_*$.
The definitions give
\begin{equation}\tag{32.1}
L^{-k_*}\cdot\mu_*=\bar{m}\cdot L^{-N}=\bar{m}\cdot\eta,\qquad \ell_*\cdot L^{k_*}=R\cdot L^{N}=n,\qquad n=\ell_0\ge L\cdot M\cdot r_0\ge4,
\end{equation}
where the inequality $\ell_0\ge L\cdot M\cdot r_0$ is the definition of $N_{\min}(R,e)$.

\emph{The levels $k_*$ and $0$.}
By (Sc1), $k_*\ge1$.
The terminal integral of Section~0.3 is therefore an integral at a level $\ge1$.
Its generators $\Psi_\beta(y)$, $\bar{\Psi}_\beta(y)$ carry four components at each unit site and the weight $\nu=1$, so that $\mathsf{w}_W=1$;
at the level $0$, by contrast, $\nu=\frac{1}{48}$ and $\mathsf{w}_W=\sqrt{48}$ (Sections~I.3.6 and~I.5.2).
In the terminal integral the factor $\frac{1}{48}$ comes from the fine lattice $T^{\mathrm{f}}_{k_*}$ alone:
from $\int\mathrm{d}x$ (in the form $\left\langle\bar{\psi},\psi\right\rangle$ and in the integrals $\int\mathrm{d}\xi$, $\int\mathrm{d}\boldsymbol{\zeta}$ of the activities) and from the cluster mean $J^*$ contained in $D_*(B)$ and in $H_{k_*}(B)$.
Besides objects of the level $k_*$, the Propositions~32.m contain the determinants $\det\left(D_B+z\right)$ of the level $0$, where $B$ is a real one-cochain on $\mathbb{Z}^3/n\mathbb{Z}^3$ and $z$ is real.
These are determinants of linear maps of $\mathfrak{S}^{\mathrm{ch}}(1;n)$, a space of dimension $192\,n^3$ (Proposition~46\,(e), first sentence), and they stand in $\mathrm{dens}^{\theta,0}_N$, $\mathrm{zr}_n$ and $\mathrm{zr}^{\mathrm{free}}_n$ (Section~0.5).

\subsubsection*{0.2. The conditions (T2)--(T5)}

The conditions (T2)--(T5) are those of Section~G4.0.3.
Let $\mathrm{lh}(v)=L\cdot M\cdot v^r$ be the function of Proposition~14.2\,(a).
By Proposition~14.3\,(f) the inequalities (e3)--(e6) of Proposition~14.2\,(b) hold for every $v\ge-\log e_1$, and $v_*\ge v_K=-\log e_K\ge-\log e_1$.
At $v=v_*$ one has $\exp(-v_*)=e_*$, $v_*^p=p_*$ and $v_*^{p_0}=p_{0*}$, and (Sc3) gives
\[
\ell_*<\mathrm{lh}(v_*),\qquad \mu_*^{-1}<L\cdot\ell_*<L\cdot\mathrm{lh}(v_*),\qquad \ell_*\ge M\cdot v_*^r .
\]
It follows that:
\begin{itemize}
\item[(T2)] by (e3),
\begin{align*}
8C_0^2\cdot\ell_*^6\cdot\left(b_-^{-1}+\mu_*^{-1}\right)&\le 8C_0^2\cdot\mathrm{lh}(v_*)^6\cdot\left(b_-^{-1}+L\cdot\mathrm{lh}(v_*)\right)\\
&\le\exp\left(\frac{\varepsilon\cdot v_*}{2}\right)=e_*^{-\frac{\varepsilon}{2}},\\
4\left(2C_0+c_\kappa+5c_{\mathrm{fl}}+6\right)\cdot\ell_*^3&\le 4\left(2C_0+c_\kappa+5c_{\mathrm{fl}}+6\right)\cdot\mathrm{lh}(v_*)^3\le e_*^{-\varepsilon},\\
e_*^{\frac14-9\varepsilon}&=\exp\left(-\left(\frac14-9\varepsilon\right)\cdot v_*\right)\le\frac18 ;
\end{align*}
\item[(T3)] by (e4),
\[
4C_0\cdot\ell_*^2\cdot p_*\le 4C_0\cdot\mathrm{lh}(v_*)^2\cdot v_*^p<e_*^{-\varepsilon};
\]
\item[(T4)] by (e5),
\begin{multline*}
12\ell_*^3\cdot\left(\exp\left(-\frac{p_{0*}}{2}\right)+\exp\left(-\frac{p_*^2}{2C_0^2}\right)\right)\\
\le 12\,\mathrm{lh}(v_*)^3\cdot\left(\exp\left(-\frac{v_*^{p_0}}{2}\right)+\exp\left(-\frac{v_*^{2p}}{2C_0^2}\right)\right)\le e_*;
\end{multline*}
\item[(T5)] by (e6),
\[
C_\Sigma\cdot\exp\left(-\gamma_\Sigma\cdot\ell_*\right)\le C_\Sigma\cdot\exp\left(-\gamma_\Sigma\cdot M\cdot v_*^r\right)\le e_* .
\]
\end{itemize}
For these four conditions this is the assertion (Sc6).
The first inequality of each chain uses $C_0,c_\kappa,c_{\mathrm{fl}},C_\Sigma\ge0$, $b_->0$ and $\gamma_\Sigma>0$.
At $v=v_K$ the third inequality of (e3) reads
\begin{equation}\tag{32.2}
e_K^{\frac14-9\varepsilon}\le\frac18 .
\end{equation}
Since $0<\varepsilon\le\frac{1}{200}$ by (C7), $C_0\ge1$ by Proposition~14.2\,(f), and $\ell_*\ge1$, one obtains the consequences
\begin{equation}\tag{32.a}
\begin{gathered}
e_*^{z}\le e_*^{z'}\le1\quad\text{for real numbers }z\ge z'\ge0;\qquad e_*^{\frac14}\le e_*^{\frac14-9\varepsilon}\le\frac18;\\
e_*^{\frac34-4\varepsilon}\le e_*^{\frac14-9\varepsilon}\le\frac18;\qquad e_*^{\frac34}\le8^{-3};
\end{gathered}
\end{equation}
\begin{equation}\tag{32.b}
\begin{gathered}
\varepsilon_1=\frac12\,e_*^{\frac14-9\varepsilon}\le\frac{1}{16};\\
\varepsilon_3=6\ell_*^3\cdot\left(\exp\left(-\frac{p_{0*}}{2}\right)+\exp\left(-\frac{p_*^2}{2C_0^2}\right)\right)\le\frac{e_*}{2}\quad\text{by (T4)};\\
\varepsilon_4\le\varepsilon_3\le\frac{e_*}{2};\qquad
\varepsilon_2=\ell_*^{-1}\cdot C_\Sigma\cdot\exp\left(-\gamma_\Sigma\cdot\ell_*\right)\le e_*\quad\text{by (T5)}.
\end{gathered}
\end{equation}

\subsubsection*{0.3. The terminal integral, written out}

\emph{Modes.}
Let
\[
\mathrm{Mo}:=\left\{(y,\beta):\ y\in T^0_{k_*},\ \beta\in\{1,2,3,4\}\right\};
\]
these are $4\ell_*^3$ modes, each with the weight $\nu=1$ (Section~I.3.6; $k_*\ge1$).
Write $\Psi_\beta(y):=\Psi(y,\beta,0)$ and $\bar{\Psi}_\beta(y):=\Psi(y,\beta,1)$.
By (I.AP) the site $y$ ranges over a fundamental domain of $\ell_*\mathbb{Z}^3$ in $\mathbb{Z}^3$, and the restriction to this domain identifies $\mathfrak{S}(1,\ell_*)$ with the space of the functions on $\mathrm{Mo}$.
A linear map $O$ of $\mathfrak{S}(1,\ell_*)$ thus has the matrix $O(y,y')_{\beta\beta'}$, defined by
\[
(Ov)_\beta(y)=\sum_{y',\beta'}O(y,y')_{\beta\beta'}\,v_{\beta'}(y'),
\]
and this matrix is the kernel of $O$ in the sense of Proposition~15\,(D2) at $\nu=1$.
A different fundamental domain conjugates the matrix by a diagonal matrix of signs; this leaves unchanged $\det O$, the moduli of the entries, the sums of the moduli along the rows, and the operator norm for the counting scalar product.
The form of Proposition~15\,(D2) is
\[
\left\langle\bar{\Psi},O\Psi\right\rangle=\sum_{y,y'}\sum_{\beta,\beta'}\bar{\Psi}_\beta(y)\,O(y,y')_{\beta\beta'}\,\Psi_{\beta'}(y')=\left\langle\bar{\Psi},O\Psi\right\rangle_\nu .
\]
The Berezin integral $\int\mathrm{D}\Psi$ is that of (I.3.5) with $\nu=1$.
For $\mathrm{aw}>0$ and $F\in\mathrm{Gr}[\Psi]$, $|F|_{\mathrm{aw}}$ denotes the norm (I.5.1) of $F$ with the single weight $\mathrm{aw}$ on all $8\ell_*^3$ generators.

\emph{Fields.}
Let $B\in\tilde{G}_{k_*}$, and let $x=(c,g)$, $x'=(c',g')$ be fermion sites of $T^{\mathrm{f}}_{k_*}$ (Sections~I.3.5 and~I.5.4, (I.1.6), (I.1.7)).
The fields are
\begin{equation}\tag{32.3}
\begin{aligned}
\psi_\beta(c,g)&=\sum_{y'}\sum_{\beta'}H_{k_*}(B)\left((c,g),y'\right)_{\beta\beta'}\,\Psi_{\beta'}(y'),\\
\bar{\psi}_\beta(c,g)&=\sum_{y'}\sum_{\beta'}\bar{H}_{k_*}(B)\left((c,g),y'\right)_{\beta\beta'}\,\bar{\Psi}_{\beta'}(y'),\\
H_{k_*}(B)&=b_{k_*}\,S_{k_*}(B)\,J\,Q^{T}_{k_*}(-B),\qquad S_{k_*}(B)=F_\infty(B)^{-1},\\
F_\infty(B)&=D_B+\mu_*+b_{k_*}\,JP_{k_*}(B)J^*,\qquad \bar{H}_{k_*}(B)=\mathsf{C}^{-1}H_{k_*}(-B)\,\mathsf{C},\\
\chi_\beta(x,x')&=d_X(x,x')^{-\alpha}\left(U_B(x,x')\,\psi_\beta(x')-\psi_\beta(x)\right),\\
\bar{\chi}_\beta(x,x')&=d_X(x,x')^{-\alpha}\left(U_B(x,x')^{-1}\,\bar{\psi}_\beta(x')-\bar{\psi}_\beta(x)\right)\qquad\left((x,x')\in\mathrm{Pz}_{k_*}\right).
\end{aligned}
\end{equation}
Here the pairs range over the whole of $\mathrm{Pz}_{k_*}$.
Over every photon site $c$ there are $48\cdot47$ ordered pairs inside the cluster ($c=c'$, $g\ne g'$), and for them $d_X=\eta_{k_*}$ and $U_B=1$.
By (I.1.6), $U_{-B}=U_B^{-1}$.
The substitution (32.3) is the homomorphism $F\mapsto F\left(H^{\#}_{k_*}(B)\Psi\right)$ of Proposition~20.1.
Finally, by (I.1.3) and Section~I.3.6,
\[
\left\langle\bar{\psi},\psi\right\rangle=\int\mathrm{d}x\,\bar{\psi}(x)\cdot\psi(x)=\frac{\eta_{k_*}^3}{48}\sum_{c\in T_{k_*}}\sum_{g\in O_{\mathrm{cub}}}\sum_{\beta=1}^{4}\bar{\psi}_\beta(c,g)\,\psi_\beta(c,g).
\]

\emph{The display (Y).}
This is the display in the data and definitions of Proposition~29.
It is applied to a real pair that is controlled on $R$, at a point $\tilde{A}\in\mathrm{Dom}_{k_*}$ in the place of $A$, with $B:=H_{k_*}\tilde{A}$ in the place of $\mathscr{A}$, and with the fields (32.3); that is,
\begin{multline*}
Y_*(\tilde{A})=\det D_*(0)^{-1}\int\mathrm{D}\Psi\,\exp\biggl\{-\left\langle\bar{\Psi},D_*(B)\Psi\right\rangle-m_*\cdot\left\langle\bar{\psi},\psi\right\rangle\\
-\varepsilon^{\mathrm{v},t}_*\cdot\ell_*^3+\sum_{X\in D_{k_*}}E_*(X;B,\psi,\chi)\biggr\}.
\end{multline*}
By the definition of $\mathrm{Dom}_{k_*}$ the bond function $B$ lies in $\frac14\tilde{R}_{k_*}$, hence in $\tilde{R}_{k_*}$, and hence in $\tilde{G}_{k_*}$ by \cite[Lemma~5.17(a)]{QEDI}.

\subsubsection*{0.4. Statements on photon objects and on domains}

\begin{description}
\item[The background field.]
By \cite[Proposition~5.31]{QEDI}, that is \cite[Proposition~15.5]{QEDI}, the map $H_k$ is linear and real, and it sends a constant bond function on $T^0_k$ to the same constant on $T_k$.

\item[Kernel bounds.]
Let $k\ge0$, let $\ell\ge1$ be an integer, and let $A'$ be a real bond function on $T^0_k$.
Each of the three norms of $H_kA'$ (Section~I.5.3) is at most the first member of
\[
\sup_b\sum_{b'}C_P\cdot\exp\left(-\gamma_P\cdot d(b,b')\right)\cdot|A'(b')|\le C_P\cdot c_{\mathrm{bd}}(\gamma_P)\cdot|A'|_\infty\le C_0\cdot|A'|_\infty ,
\]
where $C_0$ is the constant of (C6).
This follows from \cite[Theorem~15.14]{QEDI} (see \cite[Proposition~5.32]{QEDI}), with the constants $C_P$, $\gamma_P$ fixed in \cite[Section~15.7]{QEDI}, with the kernels of \cite[Section~15.1.1]{QEDI}, and with the bound
\[
\sum_{b'}\exp\left(-\gamma_P\cdot d(b,b')\right)\le c_{\mathrm{bd}}(\gamma_P)
\]
of \cite[Lemma~15.2]{QEDI} at the rate $\gamma_P\in(0,1]$.
The hypotheses of \cite[Theorem~15.14]{QEDI}, namely $0\le\alpha<1$, $L$ odd and $L\ge3$, hold by (C1) and (C3).

\item[Fields in the domains.]
This is \cite[Proposition~5.38(i)]{QEDI}, that is \cite[Proposition~15.30(i)]{QEDI}.
Let $\ell\in\mathrm{Cl}$ and $e_k<1$, and assume
\begin{gather*}
8C_B\cdot p_k\le e_k^{-\frac34+3\varepsilon},\qquad 8e_k^{\frac34-\frac{7\varepsilon}{2}}\le1,\\
p_{0,k}\le e_k^{-\frac{\varepsilon}{4}},\qquad 2C_w^2\cdot c_\gamma^2\cdot c_C\cdot e_k^{\frac{\varepsilon}{4}}\le1,
\end{gather*}
as well as $C_w\ge C_P$ and $\gamma_{\mathrm{w}}\le\gamma_P$.
Then every real bond function $\tilde{A}$ on $T^0_k$ with $|d\tilde{A}|_\infty\le p_k$, whatever its harmonic part, lies in $\mathrm{Dom}_k$.
At $k=k_*$ the four displayed inequalities are, at $v=v_*\ge-\log e_0$, the third and the fourth inequality of (A10), the first inequality of (A1), and (A2); these are conditions of Proposition~14.2\,(b).
The two relations $C_w\ge C_P$ and $\gamma_{\mathrm{w}}\le\gamma_P$ are contained in Proposition~14.2\,(f).

\item[Gauge fixing.]
By \cite[Corollary~15.20]{QEDI} together with \cite[Lemma~15.19]{QEDI} (see \cite[Proposition~5.40]{QEDI}), for every integer $\ell\ge2$ and every real bond function $A$ on $\mathbb{Z}^3/\ell\mathbb{Z}^3$ with mean zero there is a real function $\boldsymbol{\lambda}$ of the sites with
\[
\left|A-d\boldsymbol{\lambda}\right|_\infty\le2\ell^2\cdot|dA|_\infty .
\]

\item[Domains.]
These are \cite[Lemmas~17.3, 17.4(a) and~17.6]{QEDI} (see \cite[Lemmas~5.17(a) and~5.19]{QEDI}).
By the hypotheses stated at the head of \cite[Section~17.3]{QEDI} the three lemmas hold for every level $k\ge0$, every side $\ell\in\mathrm{Cl}$ and every $e_k\in(0,e_0]$, the third one without its assertion on $\mathrm{Dom}_k$.
They assert, in this order:
the sets $R_k$, $\mathrm{cf}\cdot\tilde{R}_k$ with $0<\mathrm{cf}\le1$, and $\tilde{G}_k$ are open in the space of the complex bond functions on $T_k$;
$\tilde{R}_k\subset\tilde{G}_k$;
the sets $\mathrm{cf}\cdot\tilde{R}_k$ and $\tilde{G}_k$ are invariant under the addition of real constant bond functions and under $B\mapsto-B$.
By the definitions of Section~I.5.3, that is of \cite[Section~5.3.1 and Definition~5.16]{QEDI}, one has moreover $R_k\subset\tilde{R}_k$ (take $A_0=0$) and $\frac14\tilde{R}_k\subset\tilde{R}_k$.

\item[Coordinates on $F_N$.]
By \cite[Theorem~14.17(b)--(e)]{QEDI} (see \cite[Proposition~5.53]{QEDI}), the map $F^{\mathrm{u}}\mapsto\left(A_*;\zeta_0,\dots,\zeta_{k_*-1}\right)$ is a linear bijection of $F_N$ onto $V_{\mathrm{t}}(T^0_{k_*})\times\prod_{j<k_*}V_j$.
The law of $A_*$ under $\gamma_N$ is $\mathrm{d}\mathsf{m}_*$, a probability measure on $V_{\mathrm{t}}(T^0_{k_*})$; the elements of $V_{\mathrm{t}}(T^0_{k_*})$ are coclosed and have mean zero (Section~G4.0.6).
\end{description}

\subsubsection*{0.5. The measures and the numbers of Section~I.7.3}

Let $F_N$, $\gamma_N$, $A(F^{\mathrm{u}})$, $\mathrm{dens}^{\theta,0}_N$, $\nu^\theta$, $\bar{\nu}$, $\bar{\mu}$, $\mathrm{zr}_n(q,\mathbf{m})$ and $\mathrm{zr}^{\mathrm{free}}_n$ be the objects of Section~I.7.3 at $a=0$ (with $H_\theta$ of Section~I.6.6), formed on the cubic torus with $n$ sites per side, for a real pair $(m_{\mathrm{in}},\varepsilon^{\mathrm{v}}_{\mathrm{in}})$ and for $\theta\in\mathbb{R}^3$.
Inserting the definition of $\mathrm{dens}^{\theta,0}_N$ into
\[
\bar{\nu}(1)=\int_{[-\pi,\pi]^3}\frac{\mathrm{d}\theta}{(2\pi)^3}\int\mathrm{d}\gamma_N\,\mathrm{dens}^{\theta,0}_N
\]
one finds
\begin{equation}\tag{32.4}
\mathrm{zr}_n\left(q,\bar{m}\eta+m_{\mathrm{in}}\right)=\exp\left(\varepsilon^{\mathrm{v}}_{\mathrm{in}}\cdot n^3\right)\cdot\det\left(D_0+\bar{m}\eta\right)\cdot\bar{\nu}(1).
\end{equation}

\subsubsection*{0.6. Elementary inequalities}

Here $z$ is a real variable.
\begin{itemize}
\item[(el1)] $\exp(z)-1\le2z$ for $0\le z\le1$.
Indeed, $\exp$ is convex, so that $\exp(z)\le1+(\mathrm{e}-1)\,z$ on $[0,1]$, and $\mathrm{e}-1<2$.
\item[(el2)] $-\log(1-z)\le2z$ for $0\le z\le\frac12$, and $\log(1+z)\le z$ for $z\ge0$.
Indeed,
\[
-\log(1-z)=\sum_{j\ge1}\frac{z^j}{j}\le\frac{z}{1-z}\le2z ,
\]
and $1+z\le\exp(z)$.
\item[(el3)] Every entry of a square matrix has modulus at most the operator norm of the matrix for the Euclidean norm.
Indeed, the entry in the row $i$ and the column $j$ is a coordinate of the image of the $j$-th vector of the standard basis, and the Euclidean norm of this image does not exceed the operator norm.
\item[(el4)] Let $t\mapsto F(t)$ be a continuous map of a circle into $\mathrm{Gr}[\Psi]$.
Then
\[
\left|\oint\mathrm{d}t\,F(t)\right|_{\mathrm{aw}}\le\oint|\mathrm{d}t|\,\left|F(t)\right|_{\mathrm{aw}} .
\]
Indeed, $|\cdot|_{\mathrm{aw}}$ is a sum of moduli of coefficients with positive weights.
\end{itemize}

\subsection*{1. Proposition 32.1 (terminal estimate)}

\begin{proof}[Proof of Proposition 32.1]
Let $\mathrm{Stw}$ be the set of the pairs $(\theta,A)$ in which $\theta\in\mathbb{R}^3$ and $A$ is a real bond function on $T^0_{k_*}$ with mean zero and $|dA|_\infty\le p_*$.
Throughout Steps~0--6 a point $(\theta,A)$ of $\mathrm{Stw}$ is fixed.

\emph{Step 0 (the data of the level $k_*$).}
The pair is controlled on $R$: its trajectory on $R$ exists up to $k_*$, and $|m_k|\le w^{\mathrm{m}}_k$ for $0\le k\le k_*$ (Section~G4.0.5).
Proposition~28.1 therefore applies at $k=k_*$, and gives
\[
\varepsilon^{\mathrm{v},t}_*,\ m_*\in\mathbb{R};\qquad \mathscr{E}_*\in K^{x}_{k_*}(\ell_*),\qquad E_*=\mathscr{E}_*|_0\in K_{k_*}(\ell_*),\qquad \|E_*\|_{k_*}\le w^{\mathrm{E}}_* .
\]
The case $k=k_*$ of the bounds on the $m_k$ gives in addition $|m_*|\le w^{\mathrm{m}}_*$.
By the definition of the norm of $K_{k_*}$, that is by (I.k6) of Section~I.5.4 in the form without the second slot,
\begin{equation}\tag{32.5}
\left\|E_*(X;B)\right\|_{h_*}\le\exp\left(-\kappa\cdot d_M(X)\right)\cdot w^{\mathrm{E}}_*\qquad\text{for all }X\in D_{k_*}\text{ and }B\in\tilde{R}_{k_*}.
\end{equation}
The data $(\ell,\mathfrak{N},\mu,\mathsf{b})=(\ell_*,L^{k_*},\mu_*,b_{k_*})$ are admitted in Proposition~30.3: $L^{k_*}$ is odd and $\ge3$, because $L$ is odd and $k_*\ge1$; $0<\mu_*\le1$ by (Sc3); and $b_{k_*}>0$.
Moreover $\ell_*\in\mathrm{Cl}$ by (Sc4), and $b_-\le b_{k_*}\le b_+$ (Section~I.3.4).
Let $c_\theta$ be the constant real bond function with the components $\frac{\theta_\mu}{e_*\cdot\ell_*}$ (Section~I.6.6); it is considered both on the unit lattice $T^0_{k_*}$, where it is added to the block field, and on the fine lattice $T_{k_*}$, where it is a connection.
Proposition~30.3\,(c), (f) with $\mathrm{ec}=e_*$ gives
\begin{equation}\tag{32.6}
\mathrm{Rtw}_\theta=\frac{\det D_*(c_\theta)}{\det D_*(0)}>0 .
\end{equation}

\emph{Step 1 (a representative of $A$; the block fields lie in $\mathrm{Dom}_{k_*}$).}
Since $\ell_*\in\mathrm{Cl}$, one has $\ell_*\ge13\,L\cdot M\ge2$; and $A$ has mean zero.
By the gauge fixing of Section~0.4 (\cite[Corollary~15.20]{QEDI}) there is a real function $\boldsymbol{\lambda}$ on the sites of $T^0_{k_*}$ such that $A':=A-d\boldsymbol{\lambda}$ satisfies
\[
|A'|_\infty\le2\ell_*^2\cdot|dA|_\infty\le2\ell_*^2\cdot p_* .
\]
As $dc_\theta=0$ and $dd\boldsymbol{\lambda}=0$,
\[
d\left(A+c_\theta\right)=d\left(A'+c_\theta\right)=dA ,
\]
and the sup norm of this cochain is at most $p_*$.
By \cite[Proposition~5.38(i)]{QEDI} (Section~0.4) both $A+c_\theta$ and $A'+c_\theta$ lie in $\mathrm{Dom}_{k_*}$.
This proves the first assertion of Proposition~32.1.

The cochain $a=0$ is an external connection in the sense of Proposition~29, since $0\in\mathfrak{X}_j$ for every $j\le k_*$, and for it $\mathscr{C}_*=0$.
Proposition~29\,(b) at $a=0$, applied to the two elements $A'+c_\theta$ and $\left(A'+c_\theta\right)+d\boldsymbol{\lambda}=A+c_\theta$ of $\mathrm{Dom}_{k_*}$, yields
\begin{equation}\tag{32.7}
Y_*\left(A+c_\theta\right)=Y_*\left(A'+c_\theta\right).
\end{equation}
Put $\mathscr{A}^0:=H_{k_*}A'$; it is real by \cite[Proposition~5.31]{QEDI}.
The kernel bounds of Section~0.4 and (T3) give
\begin{equation}\tag{32.8}
\text{each of the three norms of }\mathscr{A}^0\text{ is at most }C_0\cdot|A'|_\infty\le2C_0\cdot\ell_*^2\cdot p_*<\frac12\,e_*^{-\varepsilon}.
\end{equation}
Since $H_{k_*}$ is linear and sends constants to the same constants (\cite[Proposition~5.31]{QEDI}),
\[
\mathscr{A}:=H_{k_*}\left(A'+c_\theta\right)=\mathscr{A}^0+c_\theta .
\]

\emph{Step 2 (a disc in $\tilde{R}_{k_*}$).}
Put $\mathrm{rad}:=e_*^{-\frac34+4\varepsilon}$; by (32.a), $\mathrm{rad}\ge8$.
Let $t$ be a complex number with $|t|<2\,\mathrm{rad}$.
By (32.8) each of the three norms of $t\cdot\mathscr{A}^0$ is smaller than
\[
2\,\mathrm{rad}\cdot\frac12\,e_*^{-\varepsilon}=e_*^{-\frac34+3\varepsilon}.
\]
As $e_*<1$, this number is the smallest of the three thresholds $e_*^{-\frac34+\varepsilon}\ge e_*^{-\frac34+2\varepsilon}\ge e_*^{-\frac34+3\varepsilon}$ of $R_{k_*}$.
Hence $t\cdot\mathscr{A}^0\in R_{k_*}\subset\tilde{R}_{k_*}$, and \cite[Lemmas~17.6 and~17.4(a)]{QEDI} give
\begin{equation}\tag{32.9}
\begin{gathered}
\mathscr{A}^{(t)}:=t\cdot\mathscr{A}^0+c_\theta\in\tilde{R}_{k_*}\subset\tilde{G}_{k_*}\qquad\text{for }|t|<2\,\mathrm{rad};\\
\mathscr{A}^{(1)}=\mathscr{A},\qquad \mathscr{A}^{(0)}=c_\theta .
\end{gathered}
\end{equation}

\emph{Step 3 (change of the reference operator).}
The bond function $c_\theta$ on $T_{k_*}$ is real.
Proposition~19\,(d) at the level $k_*$ gives
\begin{equation}\tag{32.10}
\begin{gathered}
\mathrm{Dtw}:=D_*(c_\theta)\text{ is invertible},\qquad \mathrm{Cov}_\theta:=\mathrm{Dtw}^{-1},\\
|\mathrm{Cov}_\theta|_{\mathrm{op}}\le b_{k_*}^{-1}+\mu_*^{-1}\le b_-^{-1}+\mu_*^{-1}.
\end{gathered}
\end{equation}
For $|t|<2\,\mathrm{rad}$ let
\[
\mathrm{fD}(t):=\left\langle\bar{\Psi},D_*\left(\mathscr{A}^{(t)}\right)\Psi\right\rangle .
\]
Then
\[
\left\langle\bar{\Psi},D_*(\mathscr{A})\Psi\right\rangle=\left\langle\bar{\Psi},\mathrm{Dtw}\,\Psi\right\rangle+\left(\mathrm{fD}(1)-\mathrm{fD}(0)\right).
\]
With the fields (32.3) at $B=\mathscr{A}$ define the element
\begin{equation}\tag{32.11}
\mathrm{Etm}:=-\left(\mathrm{fD}(1)-\mathrm{fD}(0)\right)-m_*\cdot\left\langle\bar{\psi},\psi\right\rangle+\sum_{X\in D_{k_*}}E_*(X;\mathscr{A},\psi,\chi)
\end{equation}
of $\mathrm{Gr}[\Psi]$.
The element $\mathrm{Etm}$ is even.
Indeed, $\mathrm{fD}(t)$ and $\left\langle\bar{\psi},\psi\right\rangle$ are homogeneous of degree two in the generators $\Psi(u)$; the $E_*(X;\mathscr{A})$ are even elements of $\mathrm{Gr}[\psi,\chi]$ (Section~I.5.4); and the substitution (32.3) replaces each generator by a linear combination of generators.
Even elements commute with each other.
Consequently the exponential in (Y) at $\tilde{A}=A'+c_\theta$ is the product
\[
\exp\left(-\varepsilon^{\mathrm{v},t}_*\cdot\ell_*^3\right)\cdot\exp\left(-\left\langle\bar{\Psi},\mathrm{Dtw}\,\Psi\right\rangle\right)\cdot\exp(\mathrm{Etm}).
\]
The last identity of Proposition~15\,(l)\,(iii), with $\mathrm{Cov}=\mathrm{Cov}_\theta$ and $\mathrm{Cov}^{-1}=\mathrm{Dtw}$, reads
\[
\int\mathrm{D}\Psi\,\exp\left(-\left\langle\bar{\Psi},\mathrm{Dtw}\,\Psi\right\rangle\right)\cdot\exp(\mathrm{Etm})=\det(\mathrm{Dtw})\cdot\int\mathrm{d}\mu_{\mathrm{Cov}_\theta}(\Psi)\,\exp(\mathrm{Etm}).
\]
Combined with (Y), (32.6) and (32.7) this gives
\begin{equation}\tag{32.12}
\begin{gathered}
Y_*\left(A+c_\theta\right)=\exp\left(-\varepsilon^{\mathrm{v},t}_*\cdot\ell_*^3\right)\cdot\mathrm{Rtw}_\theta\cdot\int\mathrm{d}\mu_{\mathrm{Cov}_\theta}(\Psi)\,\exp(\mathrm{Etm}),\\
\text{that is}\qquad \mathrm{vrem}(\theta,A)=\int\mathrm{d}\mu_{\mathrm{Cov}_\theta}(\Psi)\,\exp(\mathrm{Etm})-1 .
\end{gathered}
\end{equation}

\emph{Step 4 (the difference of the two quadratic forms).}
Put $\mathrm{aw}_1:=e_*^{-\frac14}$.
The coefficients of $\mathrm{fD}(t)$ are the numbers $D_*\left(\mathscr{A}^{(t)}\right)(y,y')_{\beta\beta'}$.
By Proposition~20.1\,(a) the entries of $D_*(\cdot)$ are analytic on the open set $\tilde{G}_{k_*}$, and $t\mapsto\mathscr{A}^{(t)}$ is affine.
By (32.9) the function $\mathrm{fD}$ is therefore analytic on $\left\{|t|<2\,\mathrm{rad}\right\}$, and on this disc Proposition~20.1\,(a) gives
\[
\left|\mathrm{fD}(t)\right|_{\mathrm{aw}_1}\le C_0\cdot e_*^{-\frac12}\cdot\ell_*^3 .
\]
The points $0$ and $1$ lie inside the circle $|t|=\mathrm{rad}$, and $\frac{1}{t(t-1)}=\frac{1}{t-1}-\frac{1}{t}$.
Cauchy's integral formula, applied to each coefficient, gives
\[
\mathrm{fD}(1)-\mathrm{fD}(0)=\frac{1}{2\pi i}\oint_{|t|=\mathrm{rad}}\frac{\mathrm{d}t}{t(t-1)}\,\mathrm{fD}(t).
\]
On the circle $|t(t-1)|\ge\mathrm{rad}\cdot(\mathrm{rad}-1)$, the length of the circle is $2\pi\,\mathrm{rad}$, and $\mathrm{rad}-1\ge\frac{\mathrm{rad}}{2}$.
With (el4) it follows that
\begin{equation}\tag{32.13}
\begin{aligned}
\left|\mathrm{fD}(1)-\mathrm{fD}(0)\right|_{\mathrm{aw}_1}&\le\frac{C_0\cdot e_*^{-\frac12}\cdot\ell_*^3}{\mathrm{rad}-1}\le2C_0\cdot e_*^{-\frac12}\cdot\ell_*^3\cdot e_*^{\frac34-4\varepsilon}\\
&=2C_0\cdot e_*^{\frac14-4\varepsilon}\cdot\ell_*^3 .
\end{aligned}
\end{equation}

\emph{Step 5 (the norm of $\mathrm{Etm}$).}
Put $\mathrm{aw}_2:=C_0^{-1}\cdot e_*^{-\frac14+\varepsilon}$; then $\mathrm{aw}_2\le\mathrm{aw}_1$, since $C_0\ge1$ and $e_*<1$.
\begin{itemize}
\item[(i)] By Proposition~15\,(k),
\[
\left|\mathrm{fD}(1)-\mathrm{fD}(0)\right|_{\mathrm{aw}_2}\le\left|\mathrm{fD}(1)-\mathrm{fD}(0)\right|_{\mathrm{aw}_1}.
\]
\item[(ii)] Proposition~20.1\,(b) at the connection $\mathscr{A}\in\tilde{G}_{k_*}$, together with $|m_*|\le w^{\mathrm{m}}_*=e_*^{\frac34-8\varepsilon}$, gives
\[
\left|m_*\cdot\left\langle\bar{\psi},\psi\right\rangle\right|_{\mathrm{aw}_2}\le e_*^{\frac34-8\varepsilon}\cdot4e_*^{-\frac12}\cdot\ell_*^3=4e_*^{\frac14-8\varepsilon}\cdot\ell_*^3 .
\]
\item[(iii)] The pair $(\mathscr{A},0)$ lies in $\tilde{R}_{k_*}\times\mathfrak{X}_{k_*}$.
Proposition~20.1\,(c) at $(\mathscr{A},\mathscr{C})=(\mathscr{A},0)$, $\mathscr{T}=\mathscr{A}$, applied to the element $E_*(X;\mathscr{A})=\mathscr{E}_*(X;\mathscr{A},0)$ of $\mathrm{Gr}[\psi,\chi]$, and (32.5) give
\[
\left|E_*(X;\mathscr{A},\psi,\chi)\right|_{\mathrm{aw}_2}\le\left\|E_*(X;\mathscr{A})\right\|_{h_*}\le\exp\left(-\kappa\cdot d_M(X)\right)\cdot w^{\mathrm{E}}_* .
\]
Every polymer contains an $M$-cube, and $T_{k_*}$ contains $\left(\frac{\ell_*}{M}\right)^3$ $M$-cubes $\square$.
The definition of $c_\kappa$ in (C6) therefore gives
\begin{align*}
\sum_{X\in D_{k_*}}\exp\left(-\kappa\cdot d_M(X)\right)&\le\sum_{\square}\ \sum_{X\in D_{k_*},\,X\supset\square}\exp\left(-\kappa\cdot d_M(X)\right)\\
&\le\left(\frac{\ell_*}{M}\right)^3\cdot c_\kappa\le c_\kappa\cdot\ell_*^3 .
\end{align*}
\end{itemize}
From (32.11), the triangle inequality, (i)--(iii), (32.13), $w^{\mathrm{E}}_*=e_*^{\frac14-7\varepsilon}$ and (32.a) one obtains
\begin{equation}\tag{32.14}
\begin{aligned}
\left|\mathrm{Etm}\right|_{\mathrm{aw}_2}&\le\left(2C_0\cdot e_*^{\frac14-4\varepsilon}+4e_*^{\frac14-8\varepsilon}+c_\kappa\cdot e_*^{\frac14-7\varepsilon}\right)\cdot\ell_*^3\\
&\le\left(2C_0+c_\kappa+4\right)\cdot e_*^{\frac14-8\varepsilon}\cdot\ell_*^3 .
\end{aligned}
\end{equation}

\emph{Step 6 (integration).}
By the first identity of Proposition~15\,(l)\,(iii),
\[
\int\mathrm{d}\mu_{\mathrm{Cov}_\theta}(\Psi)\,\exp(\mathrm{Etm})=\int\mathrm{d}\mu_I(\Psi)\,\left[\exp(\mathrm{Etm})\right]\left(\bar{\Psi},\mathrm{Cov}_\theta\Psi\right),
\]
where $F\mapsto F\left(\bar{\Psi},\mathrm{Cov}_\theta\Psi\right)$ is the homomorphism of $\mathrm{Gr}[\Psi]$ with
\begin{gather*}
\Psi(\mathrm{mo})\mapsto\left(\mathrm{Cov}_\theta\Psi\right)(\mathrm{mo})=\sum_{\mathrm{mo}'\in\mathrm{Mo}}\mathrm{Cov}_\theta(\mathrm{mo},\mathrm{mo}')\,\Psi(\mathrm{mo}'),\\
\bar{\Psi}(\mathrm{mo})\mapsto\bar{\Psi}(\mathrm{mo}).
\end{gather*}
Let $\mathrm{Etm}'$ be the image of $\mathrm{Etm}$.
Then the image of $\exp(\mathrm{Etm})$ is $\exp(\mathrm{Etm}')$, and $\mathrm{Etm}'$ is even.
This homomorphism is the linear substitution of Proposition~15\,(d) with a single family on each side, $\phi_1=\theta_1=\Psi$ on the index set $\mathrm{Mo}\times\{0,1\}$ with $\nu=1$, and with the operator $O_{11}=\mathrm{Osub}$ whose kernel is
\[
\mathrm{Osub}\left((\mathrm{mo},0),(\mathrm{mo}',0)\right)=\mathrm{Cov}_\theta(\mathrm{mo},\mathrm{mo}'),\qquad \mathrm{Osub}\left((\mathrm{mo},1),(\mathrm{mo}',1)\right)=\mathbf{1}[\mathrm{mo}=\mathrm{mo}'],
\]
and $\mathrm{Osub}=0$ at all other pairs of indices; thus
\[
|\mathrm{Osub}|_{1,\infty}=\max\left\{1,\ \max_{\mathrm{mo}}\sum_{\mathrm{mo}'}\left|\mathrm{Cov}_\theta(\mathrm{mo},\mathrm{mo}')\right|\right\}=:\mathrm{crw}.
\]
A row of $\mathrm{Cov}_\theta$ has $4\ell_*^3$ entries, each of modulus at most $|\mathrm{Cov}_\theta|_{\mathrm{op}}$ by (el3).
By (32.10) every row sum is therefore at most $4\ell_*^3\cdot\left(b_-^{-1}+\mu_*^{-1}\right)$, and this number is $\ge1$ because $\mu_*\le1$.
The first inequality of (T2) now gives
\[
\mathrm{crw}\le4\ell_*^3\cdot\left(b_-^{-1}+\mu_*^{-1}\right)\le\frac{e_*^{-\frac{\varepsilon}{2}}}{2C_0^2\cdot\ell_*^3}=\mathrm{aw}_2\cdot\frac{e_*^{\frac14-\frac{3\varepsilon}{2}}}{2C_0\cdot\ell_*^3}\le\mathrm{aw}_2 .
\]
So Proposition~15\,(d) applies with $C_{11}=\mathrm{crw}$, $\mathsf{w}'_1=1$ and $\mathsf{w}_1=\mathrm{aw}_2$, and yields $\left|\mathrm{Etm}'\right|_1\le\left|\mathrm{Etm}\right|_{\mathrm{aw}_2}$.
Since $c_{\mathrm{fl}}>0$, the second inequality of (T2) gives
\[
\left(2C_0+c_\kappa+4\right)\cdot\ell_*^3\le\left(2C_0+c_\kappa+5c_{\mathrm{fl}}+6\right)\cdot\ell_*^3\le\frac14\,e_*^{-\varepsilon}.
\]
Together with (32.14) this proves
\begin{equation}\tag{32.15}
\left|\mathrm{Etm}'\right|_1\le\frac14\,e_*^{\frac14-9\varepsilon}\le\frac{1}{32}.
\end{equation}
By Proposition~15\,(l)\,(iii) at $\mathrm{Cov}=1$ one has $\int\mathrm{d}\mu_I=1$.
Now (32.12), Proposition~15\,(l)\,(ii) with $\mathsf{w}_W=1$ (as $\nu=1$), Proposition~15\,(b) for the even element $\mathrm{Etm}'$, and (el1) give
\begin{align*}
\left|\mathrm{vrem}(\theta,A)\right|&=\left|\int\mathrm{d}\mu_I(\Psi)\,\left(\exp(\mathrm{Etm}')-1\right)\right|\le\left|\exp(\mathrm{Etm}')-1\right|_1\le\exp\left(\left|\mathrm{Etm}'\right|_1\right)-1\\
&\le2\left|\mathrm{Etm}'\right|_1\le\frac12\,e_*^{\frac14-9\varepsilon}=\varepsilon_1 .
\end{align*}
This is (9.1).

\emph{Step 7 (continuity).}
(a)
For $(\theta,A)\in\mathrm{Stw}$ put
\[
\mathscr{A}[\theta,A]:=H_{k_*}\left(A+c_\theta\right)=H_{k_*}A+c_\theta
\]
(\cite[Proposition~5.31]{QEDI}).
The maps $\theta\mapsto c_\theta$ and $A\mapsto H_{k_*}A$ are linear maps between finite-dimensional spaces; hence $(\theta,A)\mapsto\mathscr{A}[\theta,A]$ is continuous.
By Step~1 and the definition of $\mathrm{Dom}_{k_*}$ its values lie in $\frac14\tilde{R}_{k_*}$, a subset of $\tilde{R}_{k_*}$.
By (Y), $Y_*\left(A+c_\theta\right)$ is the right side of (Y) at $B=\mathscr{A}[\theta,A]$.

(b)
The right side of (Y) is a continuous function of $B$ on $\tilde{R}_{k_*}$, an open subset of $\tilde{G}_{k_*}$ by \cite[Lemmas~17.3 and~17.4(a)]{QEDI}.
To see this, consider the following four facts.
\begin{itemize}
\item[(b1)] The entries of $D_*(B)$ are analytic on $\tilde{G}_{k_*}$ (Proposition~20.1\,(a)).
\item[(b2)] The quadruple $(k_*,\ell_*,\mu_*,e_*)$ is (SF) with $k_*\ge1$ (Section~0.1).
By Proposition~14.1\,(i) (FE1$^{\mathrm{ch}}$, first clause), for every $B\in\tilde{G}_{k_*}$ the operator $F_\infty(B)=D_B+\mu_*+b_{k_*}\,JP_{k_*}(B)J^*$ is bijective on $\mathfrak{S}^{\mathrm{ch}}(\eta_{k_*},\ell_*)$.
(The same follows from Proposition~20.1\,(a): the operator $D_{k_*}(B)$, which by Section~I.3.5 is formed with $S_{k_*}(B)=F_\infty(B)^{-1}$, is defined for every $B\in\tilde{G}_{k_*}$.)
In the basis of the delta functions at the points $(x,\beta)$, $x$ ranging over a fundamental domain, the entries of $F_\infty(B)$ are entire functions of the numbers $B(b)$.
For $D_B$ this is (I.O1).
For the remaining term, Section~I.3.3 gives, for $c\in\Delta_y$,
\begin{gather*}
\left(JP_{k_*}(B)J^*f\right)(c,g)=\varphi_B(y,c)^{-1}\,\frac{\eta_{k_*}^3}{48}\sum_{c'\in\Delta_y}\sum_{g'\in O_{\mathrm{cub}}}\varphi_B(y,c')\,f(c',g'),\\
\varphi_B(y,c)=\exp\left\{i\,e_*\cdot\eta_{k_*}\cdot\left(\tau_{k_*}B\right)(y,c)\right\},
\end{gather*}
where $\left(\tau_{k_*}B\right)(y,c)$ is a linear form in $B$ ((I.3.3)).
By Cramer's rule the entries of $S_{k_*}(B)=F_\infty(B)^{-1}$ are polynomials in the entries of $F_\infty(B)$, divided by $\det F_\infty(B)$, which does not vanish on $\tilde{G}_{k_*}$.
They are therefore continuous on $\tilde{G}_{k_*}$.
\item[(b3)] For $c\in\Delta_y$ one has $\left(JQ^{T}_{k_*}(-B)v\right)(c,g)=\varphi_B(y,c)^{-1}v(y)$, and $U_B(x,x')^{\pm1}=\exp\left\{\pm i\,e_*\cdot\eta_{k_*}\cdot(\tau B)(c,c')\right\}$ (Section~I.3.3, (I.1.6)); moreover $\tilde{G}_{k_*}$ is invariant under $B\mapsto-B$ (\cite[Lemma~17.6]{QEDI}).
With (b2) it follows that the entries of $H_{k_*}(B)$ and of $\bar{H}_{k_*}(B)=\mathsf{C}^{-1}H_{k_*}(-B)\,\mathsf{C}$, and hence all coefficients in (32.3), are continuous on $\tilde{G}_{k_*}$.
\item[(b4)] The kernels of $E_*(X;B)$ are analytic in $B\in\tilde{R}_{k_*}$, by (I.k1) of Section~I.5.4, since $E_*\in K_{k_*}$.
\end{itemize}
All index sets are finite.
The exponent in (Y) is therefore an even element $\mathrm{Ety}(B)$ of $\mathrm{Gr}[\Psi]$ whose coefficients are finite sums of products of the functions in (b1)--(b4), of the real numbers $m_*$ and $\varepsilon^{\mathrm{v},t}_*\cdot\ell_*^3$, and of the weights of Section~I.3.6; these coefficients are continuous on $\tilde{R}_{k_*}$.
Write $\mathrm{Ety}(B)=\mathrm{Ety}_0+\mathrm{Ety}_1$, where $\mathrm{Ety}_0$ is the coefficient of the unit.
Every monomial of $\mathrm{Ety}_1$ has degree $\ge2$, and $\mathrm{Gr}[\Psi]$ has $8\ell_*^3$ generators; hence
\[
\exp\mathrm{Ety}(B)=\exp\left(\mathrm{Ety}_0\right)\cdot\sum_{j=0}^{4\ell_*^3}\frac{\mathrm{Ety}_1^j}{j!},
\]
and the coefficients of $\exp\mathrm{Ety}(B)$ are continuous functions of those of $\mathrm{Ety}(B)$.
Finally, $\int\mathrm{D}\Psi$ is the coefficient of $\mathrm{Top}$ ((I.3.5) with $\nu=1$).

(c)
The bond function $0$ lies in $\tilde{G}_{k_*}$, and so does $c_\theta$ for every $\theta\in\mathbb{R}^3$, by \cite[Lemma~17.6]{QEDI}.
By (32.6), (b1) and the linearity of $\theta\mapsto c_\theta$, the function $\theta\mapsto\mathrm{Rtw}_\theta$ is continuous, and it is positive.

From (a), (b) and (c) it follows that the function
\[
\mathrm{vrem}(\theta,A)=Y_*\left(A+c_\theta\right)\cdot\frac{\exp\left(\varepsilon^{\mathrm{v},t}_*\cdot\ell_*^3\right)}{\mathrm{Rtw}_\theta}-1
\]
is continuous on $\mathrm{Stw}$.
\end{proof}

\emph{Uniformity.}
The right side of (9.1) is a function of $e_*$ and $\varepsilon$ alone, and by (Sc1) the number $e_*$ does not depend on $N\ge N_{\min}(R,e)$.
The constants $C_0$, $c_\kappa$, $C_B$, $C_w$, $c_\gamma$ of (C6), $b_-$, $c_C$, $C_P$, $\gamma_P$ of (C4), $\gamma_{\mathrm{w}}$ of (C4r), $c_{\mathrm{fl}}$ of (C0), $\alpha$, $\kappa$, $L$, $M$ of (C1), (C2), (C3), (C5), $\varepsilon$, $r$, $p$, $p_0$ of (C7), and $e_0$, $e_1$ of (C8), (C9),
on which the argument depends, are constants of the lines (C0)--(C9);
by (I.U4)$^{\mathrm{ch}}$ (Proposition~14.4\,(a)) they depend neither on $N$, nor on a level or a side, nor on $\theta$, on the running mass, on the pair or on $\bar{m}$.
The numbers $\ell_*$, $\mu_*$, $p_*$, which do depend on $(e,R,\bar{m})$, are absorbed by (T2) and (T3).
Hence (9.1) holds for all data (FL), for every real pair that is controlled on $R$, for every $\theta\in\mathbb{R}^3$ and for every $A$ with $(\theta,A)\in\mathrm{Stw}$.
Since the masses in the argument are the running mass $\mu_*=\bar{m}\cdot a_*\in(0,1]$ and the coupling $m_*$, which enters through $|m_*|\le w^{\mathrm{m}}_*$,
the estimate (9.1) needs no condition on the lattice mass $\boldsymbol{\mu}=\bar{m}\cdot\eta+m_{\mathrm{in}}$ besides the control of the pair on $R$;
in particular it holds whatever the sign of $\boldsymbol{\mu}$ (I.K10).

\subsection*{2. Proposition~32.2 (the twist-averaged partition function on one torus)}

\begin{proof}[Proof of Proposition~32.2]
Throughout this proof the objects of Section~0.5 are formed with the pair $\left(m_{\mathrm{in}},\varepsilon^{\mathrm{v}}_{\mathrm{in}}\right)$.
The numbers $\varepsilon_1,\dots,\varepsilon_4$ and the function $\varepsilon_5$ are those of Section~G4.0.3.

\emph{Step~0 (measures; the set $\mathfrak{G}$; signs).}
By (32.1) one has $n\ge4$.
Proposition~31.1\,(a) gives four facts:
$\mathrm{dens}^{\theta,0}_N$ is a bounded continuous function of $\left(\theta,F^{\mathrm{u}}\right)$ with values in $[0,\infty)$;
for every $\theta\in\mathbb{R}^3$, $\nu^{\theta}$ is a finite positive measure on $F_N$;
$\det\left(D_0+\bar{m}\cdot\eta\right)>0$;
and $\det\left(D_{H_{\theta}}+\bar{m}\cdot\eta\right)\ge0$.
By Proposition~31.1\,(c), $0<\bar{\nu}(1)<\infty$, and $\bar{\mu}$ is a probability measure on $F_N$.
It follows from (32.4) that $\mathrm{zr}_n>0$.

According to \cite[Theorem~14.17(b)--(e)]{QEDI}, each of the numbers $\zeta_j(\Upsilon)$, $j<k_{*}$, and $(dA_{*})(\mathrm{pl})$ is a linear functional of $F^{\mathrm{u}}\in F_N$.
Consequently
\[
\mathfrak{G}=\left\{F^{\mathrm{u}}\in F_N:\ \left|\zeta_j(\Upsilon)\right|\le p_{0,j}\ \text{for all }\Upsilon\text{ and all }j<k_{*};\ \left|dA_{*}\right|_{\infty}\le p_{*}\right\}
\]
is a closed subset of $F_N$.
The measure $\gamma_N$ is a probability measure.
By the dominated convergence theorem the map
\[
\theta\mapsto\nu^{\theta}(\mathfrak{G})=\int_{\mathfrak{G}}\mathrm{d}\gamma_N\,\mathrm{dens}^{\theta,0}_N
\]
is continuous on $\mathbb{R}^3$, and
\[
\bar{\nu}(\mathfrak{G})=\int_{[-\pi,\pi]^3}\frac{\mathrm{d}\theta}{(2\pi)^3}\,\nu^{\theta}(\mathfrak{G}).
\]

The number $\varepsilon^{\mathrm{v},t}_{*}$ is real.
Indeed, the pair is controlled on $R$, so that $\left|m_k\right|\le w^{\mathrm{m}}_k$ for $k<k_{*}$ (Section~G4.0.5), and Proposition~28.1 applies at $k=k_{*}$.

Let $\mathrm{Rtw}_{\theta}$ denote the number of Proposition~30.3 for the data
\[
\left(\ell,\mathfrak{N},\mu,\mathsf{b}\right)=\left(\ell_{*},L^{k_{*}},\mu_{*},b_{k_{*}}\right).
\]
These data are admissible in Proposition~30.3:
$L$ is odd and, by (Sc1), $k_{*}\ge1$, so that $L^{k_{*}}$ is odd and at least $3$;
$0<\mu_{*}\le1$ by (Sc3);
and $b_{k_{*}}>0$ by Section~I.3.4.
Moreover $\mathrm{Rtw}_{\theta}>0$ by Proposition~30.3\,(c).

\emph{Step~1 (the inequalities (7.1)).}
Let $\theta\in\mathbb{R}^3$.
\begin{itemize}
\item[(i)] By Step~0, the number $\exp\left(-\varepsilon^{\mathrm{v},t}_{*}\cdot\ell_{*}^3\right)\cdot\mathrm{Rtw}_{\theta}$ is real and positive.
\item[(ii)] The elements of ${V_{\mathrm{t}}}(T^0_{k_{*}})$ are real bond functions on $T^0_{k_{*}}$ with mean zero (\cite[Section~14.5.2]{QEDI}, \cite[Theorem~14.17(b)--(e)]{QEDI}).
Let $A\in{V_{\mathrm{t}}}(T^0_{k_{*}})$ satisfy $\chi_{*}(A)=1$, that is, $\left|dA\right|_{\infty}\le p_{*}$.
Then $(\theta,A)$ belongs to the set on which Proposition~32.1 is stated, and the definition of $\mathrm{vrem}$ in that proposition gives
\[
Y_{*}\left(A+c_{\theta}\right)=\exp\left(-\varepsilon^{\mathrm{v},t}_{*}\cdot\ell_{*}^3\right)\cdot\mathrm{Rtw}_{\theta}\cdot\left(1+\mathrm{vrem}(\theta,A)\right).
\]
\item[(iii)] The set $\left\{\chi_{*}=1\right\}$ is closed in ${V_{\mathrm{t}}}(T^0_{k_{*}})$.
On this set the function $\mathrm{vrem}(\theta,\cdot\,)$ is continuous and satisfies $\left|\mathrm{vrem}(\theta,\cdot\,)\right|\le\varepsilon_1$ (Proposition~32.1).
Therefore the function $A\mapsto\chi_{*}(A)\left(1+\mathrm{vrem}(\theta,A)\right)$, given the value $0$ where $\chi_{*}(A)=0$, is bounded and Borel on ${V_{\mathrm{t}}}(T^0_{k_{*}})$.
By (ii), it coincides with the function $A\mapsto\chi_{*}(A)\,Y_{*}\left(A+c_{\theta}\right)$ of Proposition~29\,(c), divided by the constant of (i).
\end{itemize}
Now apply Proposition~29\,(c) with $\Phi=1$ and at $a=0$; for $a=0$ one has $\mathscr{C}_{*}=0$, $Y^{x}_{*}(\,\cdot\,;0)=Y_{*}$ and $\nu^{\theta,0}=\nu^{\theta}$.
The result is
\begin{equation}
\begin{split}
\nu^{\theta}(\mathfrak{G})&=\int\mathrm{d}\mathsf{m}_{*}(A)\,\chi_{*}(A)\,Y_{*}\left(A+c_{\theta}\right)\\
&=\exp\left(-\varepsilon^{\mathrm{v},t}_{*}\cdot\ell_{*}^3\right)\cdot\mathrm{Rtw}_{\theta}\int\mathrm{d}\mathsf{m}_{*}(A)\,\chi_{*}(A)\left(1+\mathrm{vrem}(\theta,A)\right).
\end{split}
\tag{32.16}
\end{equation}
The left-hand side is a real number, and the factor in front of the last integral is real and positive.
Hence the last integral is real, and it equals
\[
\int\mathrm{d}\mathsf{m}_{*}\,\chi_{*}\left(1+\operatorname{Re}\mathrm{vrem}(\theta,\cdot\,)\right).
\]
As $\left|\operatorname{Re}\mathrm{vrem}\right|\le\varepsilon_1$, $\chi_{*}\ge0$ and $\mathrm{d}\mathsf{m}_{*}$ is a positive measure, one has
\[
\left(1-\varepsilon_1\right)\int\mathrm{d}\mathsf{m}_{*}\,\chi_{*}\le\int\mathrm{d}\mathsf{m}_{*}\,\chi_{*}\left(1+\operatorname{Re}\mathrm{vrem}(\theta,\cdot\,)\right)\le\left(1+\varepsilon_1\right)\int\mathrm{d}\mathsf{m}_{*}\,\chi_{*}.
\]
The measure $\mathrm{d}\mathsf{m}_{*}$ is a probability measure \cite[Theorem~14.17(e)]{QEDI}, and by the second bound of (9.2) in Proposition~31.2 one has $\mathrm{d}\mathsf{m}_{*}\left(\chi_{*}=0\right)\le\varepsilon_4$.
Therefore
\[
1-\varepsilon_4\le\int\mathrm{d}\mathsf{m}_{*}\,\chi_{*}\le1.
\]
Finally, $1-\varepsilon_1>0$ by (32.b).
Combined with (32.16) and (i), these inequalities give (7.1).

\emph{Step~2 (the inequalities (7.2)).}
Since $\nu^{\theta}\ge0$ for every $\theta$, the measure $\bar{\nu}$ is positive, and $\bar{\nu}(\mathfrak{G})\le\bar{\nu}(F_N)=\bar{\nu}(1)$.
The first bound of (9.2) in Proposition~31.2 states that $\bar{\mu}\left(F_N\setminus\mathfrak{G}\right)\le\varepsilon_3$.
It follows that
\[
\bar{\nu}(1)-\bar{\nu}(\mathfrak{G})=\bar{\nu}\left(F_N\setminus\mathfrak{G}\right)=\bar{\mu}\left(F_N\setminus\mathfrak{G}\right)\cdot\bar{\nu}(1)\le\varepsilon_3\cdot\bar{\nu}(1),
\]
that is, $\left(1-\varepsilon_3\right)\bar{\nu}(1)\le\bar{\nu}(\mathfrak{G})$.
Here $1-\varepsilon_3>0$ by (32.b).
This proves (7.2).

\emph{Step~3 (the inequalities (7.3)).}
By (32.1), the data $\left(\ell,\mathfrak{N},\mu,\mathsf{b}\right)=\left(\ell_{*},L^{k_{*}},\mu_{*},b_{k_{*}}\right)$ of Proposition~30.3 satisfy
\[
\mathrm{et}=L^{-k_{*}},\qquad\mathrm{et}\cdot\mu=\bar{m}\cdot\eta,\qquad\ell\cdot\mathfrak{N}=n.
\]
Accordingly, Proposition~30.3\,(d) takes the form
\begin{equation}
\det\left(D_0+\bar{m}\cdot\eta\right)\cdot\mathrm{Rtw}_{\theta}=\det\left(D_{H_{\theta}}+\bar{m}\cdot\eta\right)\cdot\frac{\det T_{\theta=0}}{\det T_{\theta}},
\tag{32.17}
\end{equation}
in which $\det\left(D_0+\bar{m}\cdot\eta\right)$ and $\det\left(D_{H_{\theta}}+\bar{m}\cdot\eta\right)$ are the determinants of Proposition~30.1 on the torus with $n$ sites per side at the mass $\mathbf{m}=\bar{m}\cdot\eta\in(0,1]$.
They are determinants on $\mathfrak{S}^{\mathrm{ch}}(1;n)$, of size $192\,n^3$, with
\[
H_{\theta}(c,c+e_{\mu})=\frac{\theta_{\mu}}{q\cdot n},
\]
and they are the determinants of Section~0.5.
One has $b_{-}\le b_{k_{*}}\le b_{+}$ (Section~I.3.4).
The first inequality of (8.5) in Proposition~30.3\,(e), taken at $\ell=\ell_{*}$, therefore gives
\[
\exp\left(-\varepsilon_2\right)\le\frac{\det T_{\theta=0}}{\det T_{\theta}}\le\exp\left(\varepsilon_2\right).
\]
Multiply (7.1) by $\det\left(D_0+\bar{m}\cdot\eta\right)>0$, insert (32.17), and use $\det\left(D_{H_{\theta}}+\bar{m}\cdot\eta\right)\ge0$.
One obtains
\begin{equation}
\begin{split}
&\exp\left(-\varepsilon^{\mathrm{v},t}_{*}\cdot\ell_{*}^3\right)\cdot\det\left(D_{H_{\theta}}+\bar{m}\cdot\eta\right)\cdot\exp\left(-\varepsilon_2\right)\left(1-\varepsilon_1\right)\left(1-\varepsilon_4\right)\\
&\qquad\le\det\left(D_0+\bar{m}\cdot\eta\right)\cdot\nu^{\theta}(\mathfrak{G})\\
&\qquad\le\exp\left(-\varepsilon^{\mathrm{v},t}_{*}\cdot\ell_{*}^3\right)\cdot\det\left(D_{H_{\theta}}+\bar{m}\cdot\eta\right)\cdot\exp\left(\varepsilon_2\right)\left(1+\varepsilon_1\right).
\end{split}
\tag{32.18}
\end{equation}
Each of the three members of (32.18) is a continuous function of $\theta$:
the middle one by Step~0;
the two others because $\det\left(D_{H_{\theta}}+\bar{m}\cdot\eta\right)$ is an entire function of the numbers $H_{\theta}(b)$ (I.O1), and these numbers are linear in $\theta$.
Integrating over $[-\pi,\pi]^3$ with respect to $\frac{\mathrm{d}\theta}{(2\pi)^3}$, one finds
\begin{equation}
\begin{split}
&\exp\left(-\varepsilon^{\mathrm{v},t}_{*}\cdot\ell_{*}^3\right)\cdot\mathrm{zr}^{\mathrm{free}}_n\cdot\exp\left(-\varepsilon_2\right)\left(1-\varepsilon_1\right)\left(1-\varepsilon_4\right)\\
&\qquad\le\det\left(D_0+\bar{m}\cdot\eta\right)\cdot\bar{\nu}(\mathfrak{G})\\
&\qquad\le\exp\left(-\varepsilon^{\mathrm{v},t}_{*}\cdot\ell_{*}^3\right)\cdot\mathrm{zr}^{\mathrm{free}}_n\cdot\exp\left(\varepsilon_2\right)\left(1+\varepsilon_1\right).
\end{split}
\tag{32.19}
\end{equation}
Here $\mathrm{zr}^{\mathrm{free}}_n>0$ by Proposition~30.2\,(b).
From (32.4) and (7.2),
\[
\exp\left(\varepsilon^{\mathrm{v}}_{\mathrm{in}}\cdot n^3\right)\cdot\det\left(D_0+\bar{m}\cdot\eta\right)\cdot\bar{\nu}(\mathfrak{G})\le\mathrm{zr}_n\le\frac{\exp\left(\varepsilon^{\mathrm{v}}_{\mathrm{in}}\cdot n^3\right)\cdot\det\left(D_0+\bar{m}\cdot\eta\right)\cdot\bar{\nu}(\mathfrak{G})}{1-\varepsilon_3},
\]
and all members of these inequalities are positive.
Inserting (32.19) and taking logarithms gives
\begin{multline*}
-\varepsilon_2+\log\left(1-\varepsilon_1\right)+\log\left(1-\varepsilon_4\right)\le\log\mathrm{zr}_n-\log\mathrm{zr}^{\mathrm{free}}_n-\varepsilon^{\mathrm{v}}_{\mathrm{in}}\cdot n^3+\varepsilon^{\mathrm{v},t}_{*}\cdot\ell_{*}^3\\
\le\varepsilon_2+\log\left(1+\varepsilon_1\right)-\log\left(1-\varepsilon_3\right).
\end{multline*}
By (32.b), the numbers $\varepsilon_1$, $\varepsilon_3$, $\varepsilon_4$ lie in $\left[0,\frac12\right]$.
Hence (el2) shows that the left member is at least $-\varepsilon_2-2\varepsilon_1-2\varepsilon_4$ and that the right member is at most $\varepsilon_2+\varepsilon_1+2\varepsilon_3$.
This is the first inequality of (7.3).
As to the second, (32.b) gives
\[
2\varepsilon_1+\varepsilon_2+2\varepsilon_3+2\varepsilon_4\le e_{*}^{\frac{1}{4}-9\varepsilon}+3\,e_{*},
\]
while (32.a) gives
\[
3\,e_{*}=3\,e_{*}^{\frac{3}{4}}\cdot e_{*}^{\frac{1}{4}}\le e_{*}^{\frac{1}{4}}\le e_{*}^{\frac{1}{4}-9\varepsilon}.
\]
The two bounds together give the second inequality of (7.3).

\emph{Step~4 (the inequality (7.4)).}
The number $n=R/\eta$ is an integer $\ge1$, and Proposition~30.2\,(b) gives
\[
\left|\log\mathrm{zr}^{\mathrm{free}}_n-p^{0}_{\mathrm{ch}}\left(\bar{m}\cdot\eta\right)\cdot n^3\right|\le\varepsilon_5\left(\bar{m}\cdot R\right).
\]
By (7.3) and the triangle inequality,
\begin{align*}
&\left|\log\mathrm{zr}_n-\left(p^{0}_{\mathrm{ch}}\left(\bar{m}\cdot\eta\right)+\varepsilon^{\mathrm{v}}_{\mathrm{in}}\right)n^3+\varepsilon^{\mathrm{v},t}_{*}\cdot\ell_{*}^3\right|\\
&\qquad\le\left|\log\mathrm{zr}_n-\log\mathrm{zr}^{\mathrm{free}}_n-\varepsilon^{\mathrm{v}}_{\mathrm{in}}\cdot n^3+\varepsilon^{\mathrm{v},t}_{*}\cdot\ell_{*}^3\right|+\left|\log\mathrm{zr}^{\mathrm{free}}_n-p^{0}_{\mathrm{ch}}\left(\bar{m}\cdot\eta\right)\cdot n^3\right|\\
&\qquad\le2\,e_{*}^{\frac{1}{4}-9\varepsilon}+\varepsilon_5\left(\bar{m}\cdot R\right).
\end{align*}
\end{proof}

\paragraph{Uniformity.}
The numbers $\varepsilon_1,\dots,\varepsilon_4$ are functions of $\left(e_{*},\ell_{*},p_{*}=v_{*}^{p},p_{0*}=v_{*}^{p_0}\right)$ and of the constants $C_0$, $C_{\Sigma}$, $\gamma_{\Sigma}$ (line (C6)) and $\varepsilon$, $p$, $p_0$ (line (C7)).
By (Sc1) and (I.U4)$^{\mathrm{ch}}$ (Proposition~14.4\,(a)) they do not depend on $N\ge N_{\min}(R,e)$.
The last member of (7.3) is a function of $\left(e_{*},\varepsilon\right)$.
The right-hand side of (7.4) is a function of $\left(e_{*},\varepsilon,\bar{m}\cdot R\right)$ and of $C_6$, $\kappa_{(0)}$ (line (C0r)).
The inequalities (7.1) hold for every $\theta\in\mathbb{R}^3$, and (7.1)--(7.4) hold for all data (FL) and for every real pair that is controlled on $R$.
The lattice mass $\boldsymbol{\mu}=\bar{m}\cdot\eta+m_{\mathrm{in}}$ appears only inside $\mathrm{dens}^{\theta,0}_N$ and $\mathrm{zr}_n$.
It enters the argument only through Proposition~29\,(c), for the controlled pair, and through Propositions~31.1 and~31.2, which hold for every real pair.
Hence (7.1)--(7.4) hold for either sign of $\boldsymbol{\mu}$ (I.K10).

\subsection*{3. Proposition~32.3 (all tori of the scheme, one $p_F$, one mass, uniformly in $N$; the reference datum (R2))}

\begin{proof}[Proof of Proposition~32.3]
Let $\bar{m}\in(0,1]$, and let $e>0$ satisfy $e_K=e\cdot h^{\frac{1}{2}}\le e_1$.

\emph{Step~0 (the data; the pair; $p_F$).}
Let $R\in\mathrm{Tor}(e)$ and $N\ge N_{\min}(R,e)$.
Then the quadruple $\left(\bar{m},e,R,N\right)$ has (S-4); hence it is (FL) (Proposition~26.6\,(e); Section~G4.0.1).
So the facts of Section~0 are available, and Proposition~32.2 applies to every real pair that is controlled on $R$.
By (Sc2) of Section~G4.0.2, which is \cite[Proposition~19.1(d)]{QEDI}, one has $N_{\min}(R',e)\le N_{\min}(R,e)$ whenever $R\le R'$ lie in $\mathrm{Tor}(e)$; hence $N\ge N_{\min}\left(R_{\mathrm{top}}(e),e\right)$.
Consequently the pair
\[
P_N(e)=\left(m_{0,N}(e),\varepsilon^{\mathrm{v}}_{0,N}(e)\right)
\]
is defined: by the definition in Proposition~28.3, it is the pair of Proposition~28.2 on the torus $R_{\mathrm{top}}(e)$.
By Proposition~28.3\,(c) it is real, and it is a function of $\left(N,e,\bar{m}\right)$ and of the scheme; $R$ does not enter its definition.
By Proposition~28.3\,(a) it is controlled on $R$.
Thus Proposition~32.2 applies with $\left(m_{\mathrm{in}},\varepsilon^{\mathrm{v}}_{\mathrm{in}}\right)=P_N(e)$.
The number $\mathrm{zr}_n\left(q,\bar{m}\cdot\eta+m_{\mathrm{in}}\right)$ of Proposition~32.2 is then the value
\[
\mathrm{zr}_n=\mathrm{zr}_n\left(e\cdot\eta^{\frac{1}{2}},\bar{m}\cdot\eta+m_{0,N}(e)\right)
\]
of the statement.
It is positive by (32.4), by Proposition~31.1\,(a) ($\det\left(D_0+\bar{m}\cdot\eta\right)>0$) and by Proposition~31.1\,(c) ($\bar{\nu}(1)>0$), where $n\ge4$ by (32.1).
Moreover $\mathrm{zr}^{\mathrm{free}}_n>0$ by Proposition~30.2\,(b).
Hence the logarithms in (10.1)--(10.3) are defined.

One has $0<\bar{m}\cdot\eta\le1$.
The number
\[
p^{0}_{\mathrm{ch}}\left(\bar{m}\cdot\eta\right)=\int_{[-\pi,\pi]^3}\frac{\mathrm{d}^3p}{(2\pi)^3}\,\log\mathfrak{d}_{\bar{m}\cdot\eta}(p)
\]
is real by Proposition~30.1\,(a), and by its definition it is a function of $\bar{m}\cdot\eta=\bar{m}\cdot L^{-N}$ and of $(\rho,w)$.
Therefore $p_F=p^{0}_{\mathrm{ch}}\left(\bar{m}\cdot\eta\right)+\varepsilon^{\mathrm{v}}_{0,N}(e)$ is a real number which depends on $\left(N,e,\bar{m}\right)$ and on the scheme, and not on $R$.

\emph{Step~1 (the vacuum energy at the stop level).}
For the trajectory of $P_N(e)$ on $R$, Proposition~28.3\,(b) gives
\[
\left|\varepsilon^{\mathrm{v},t}_{*}\right|\le5\,c_{\mathrm{fl}}\cdot e_{*}^{\varepsilon}\cdot w^{\mathrm{E}}_{*}=5\,c_{\mathrm{fl}}\cdot e_{*}^{\frac{1}{4}-6\varepsilon}.
\]
Since $C_0\ge0$ and $c_{\kappa}\ge0$, the second part of (T2) yields
\[
5\,c_{\mathrm{fl}}\cdot\ell_{*}^3\le\left(2C_0+c_{\kappa}+5\,c_{\mathrm{fl}}+6\right)\ell_{*}^3\le\frac{1}{4}\,e_{*}^{-\varepsilon}.
\]
Combining the two bounds with (32.a), one obtains
\begin{equation}
\left|\varepsilon^{\mathrm{v},t}_{*}\right|\cdot\ell_{*}^3\le\frac{1}{4}\,e_{*}^{\frac{1}{4}-7\varepsilon}\le e_{*}^{\frac{1}{4}-9\varepsilon}.
\tag{32.20}
\end{equation}

\emph{Step~2 (the inequalities (10.1) and (10.2)).}
From (7.3) and (32.20),
\begin{align*}
\left|\log\mathrm{zr}_n-\log\mathrm{zr}^{\mathrm{free}}_n-\varepsilon^{\mathrm{v}}_{0,N}(e)\cdot n^3\right|
&\le\left|\log\mathrm{zr}_n-\log\mathrm{zr}^{\mathrm{free}}_n-\varepsilon^{\mathrm{v}}_{0,N}(e)\cdot n^3+\varepsilon^{\mathrm{v},t}_{*}\cdot\ell_{*}^3\right|+\left|\varepsilon^{\mathrm{v},t}_{*}\right|\cdot\ell_{*}^3\\
&\le2\,e_{*}^{\frac{1}{4}-9\varepsilon}+e_{*}^{\frac{1}{4}-9\varepsilon},
\end{align*}
which is (10.1).
From (7.4) and (32.20),
\begin{align*}
\left|\log\mathrm{zr}_n-p_F\cdot n^3\right|
&\le\left|\log\mathrm{zr}_n-p_F\cdot n^3+\varepsilon^{\mathrm{v},t}_{*}\cdot\ell_{*}^3\right|+\left|\varepsilon^{\mathrm{v},t}_{*}\right|\cdot\ell_{*}^3\\
&\le\varepsilon_5\left(\bar{m}\cdot R\right)+2\,e_{*}^{\frac{1}{4}-9\varepsilon}+e_{*}^{\frac{1}{4}-9\varepsilon}=\varepsilon_P(R),
\end{align*}
which is the first inequality of (10.2).

The function $\varepsilon_5$ of Section~G4.0.3,
\[
\varepsilon_5(z)=C_6\cdot\log\left(\frac{1}{1-\exp\left(-\kappa_{(0)}\cdot z\right)}\right)\qquad(z\ \text{real},\ z>0),
\]
is decreasing, that is, $\varepsilon_5(z)\le\varepsilon_5(z')$ for $z\ge z'>0$.
Indeed, $\kappa_{(0)}>0$ by the description of the constants in Proposition~30.1, so that $z\mapsto\exp\left(-\kappa_{(0)}\cdot z\right)$ is decreasing on $(0,\infty)$, with values in $(0,1)$.
The function $\log\left(\frac{1}{1-(\,\cdot\,)}\right)$ is increasing on $(0,1)$, with positive values.
Finally $C_6\ge0$, because in (8.1) the left-hand side is nonnegative and the logarithm on the right-hand side is positive.
By (Sc3), $\bar{m}\cdot R\ge\mu_K>\frac{1}{L}$ and $e_{*}\le e_K$; moreover $\frac{1}{4}-9\varepsilon>0$.
Hence
\[
\varepsilon_5\left(\bar{m}\cdot R\right)\le\varepsilon_5\left(\frac{1}{L}\right),\qquad e_{*}^{\frac{1}{4}-9\varepsilon}\le e_K^{\frac{1}{4}-9\varepsilon},
\]
and the second inequality of (10.2) follows.

\emph{Step~3 (the first inequality of (10.3)).}
Let $N\ge N_0(e)=N_{\min}\left(R^{*}(e),e\right)$, and write $R^{*}=R^{*}(e)$.
Both $R^{*}$ and $L\cdot R^{*}=R_{\mathrm{top}}(e)$ belong to $\mathrm{Tor}(e)$.
Indeed, by Section~I.6.2,
\[
R^{*}=h\cdot L^{T^{*}(e)},\qquad R_{\mathrm{top}}(e)=h\cdot L^{T^{*}(e)+1},
\]
with the integer $T^{*}(e)=\log_L\left(M\cdot r_K(e)\right)-1\ge1$ (recall that $M$ is a power of $L$ with $M\ge L^2$, and that $r_K(e)$ is a power of $L$ with a nonnegative exponent); and both numbers are at most $R_{\mathrm{top}}(e)$.
This is the last sentence of (Sc5); it is shown at the beginning of the proof of \cite[Proposition~19.1(g)]{QEDI}, and $N$ does not enter it.
Therefore $\left(\bar{m},e,R^{*},N\right)$ has (S-4) and is (FL) (Proposition~26.6\,(e)), and (Sc2) gives
\[
N_{\min}\left(L\cdot R^{*},e\right)\le N_{\min}\left(R^{*},e\right)\le N.
\]
Consequently Steps~0--2 apply at this $N$ on the torus $R=R^{*}$ and on the torus $R=L\cdot R^{*}$, with the same $\eta$ and $q$ and with the same pair $P_N(e)$; hence with the same lattice mass $\boldsymbol{\mu}_N=\bar{m}\cdot\eta+m_{0,N}(e)$ and the same number $p_F$.
The numbers of sites per side are $n_*=R^{*}/\eta$ and $L\cdot R^{*}/\eta=L\cdot n_*$.
The numbers $\mathrm{zr}_{n_*}$ and $\mathrm{zr}_{L n_*}$ of Section~I.7.6, where $\mathrm{zr}_n:=\mathrm{zr}_n\left(q,\boldsymbol{\mu}_N\right)$, are the values $\mathrm{zr}_n$ of the statement on these two tori.
Since $L^3\cdot n_*^3=\left(L\cdot n_*\right)^3$,
\[
D^{\mathrm{mat}}=L^3\cdot\log\mathrm{zr}_{n_*}-\log\mathrm{zr}_{L n_*}=L^3\left(\log\mathrm{zr}_{n_*}-p_F\cdot n_*^3\right)-\left(\log\mathrm{zr}_{L n_*}-p_F\cdot\left(L\cdot n_*\right)^3\right).
\]
The first inequality of (10.2), applied on both tori, now gives
\[
\left|D^{\mathrm{mat}}\right|\le L^3\cdot\varepsilon_P\left(R^{*}\right)+\varepsilon_P\left(L\cdot R^{*}\right).
\]

\emph{Step~4 (the second and the third inequality of (10.3)).}
Let $R$ be $R^{*}$ or $L\cdot R^{*}$, and let $e_{*}=e_{*}(R)$.
Then $\bar{m}\cdot R\ge\bar{m}\cdot R^{*}$, the function $\varepsilon_5$ is decreasing, and $e_{*}\le e_K$ by (Sc3).
So on both tori
\[
\varepsilon_P(R)\le\varepsilon_5\left(\bar{m}\cdot R^{*}\right)+3\,e_K^{\frac{1}{4}-9\varepsilon},
\]
which gives the second inequality.
By (Sc3) on $R^{*}$ one has $\bar{m}\cdot R^{*}\ge\mu_K>\frac{1}{L}$, and hence $\varepsilon_5\left(\bar{m}\cdot R^{*}\right)\le\varepsilon_5\left(\frac{1}{L}\right)$.
By (32.2), $3\,e_K^{\frac{1}{4}-9\varepsilon}\le\frac{3}{8}\le1$.
With $C_{\mathrm{mat}}=\left(L^3+1\right)\left(\varepsilon_5\left(\frac{1}{L}\right)+1\right)$ (line (C10)) this is the third inequality.

\emph{Step~5 (the dependence of the right-hand sides; (R2)).}
One has $\varepsilon_P(R)=\varepsilon_5\left(\bar{m}\cdot R\right)+3\,e_{*}^{\frac{1}{4}-9\varepsilon}$, and $e_{*}=e_{*}(R)$ does not depend on $N\ge N_{\min}(R,e)$ by (Sc1).
Next, $R^{*}(e)=h\cdot M\cdot r_K(e)/L$, where $r_K(e)$ is formed with $v_K=-\log\left(e\cdot h^{\frac{1}{2}}\right)$ (Sections~I.6.1 and~I.6.2); thus $R^{*}(e)$ does not contain $N$.
Neither does $e_K=e\cdot h^{\frac{1}{2}}$.
Consequently no member to the right of the first inequality sign in (10.2) and in (10.3) depends on $N$.
The last member of (10.2), namely $\varepsilon_5\left(\frac{1}{L}\right)+3\,e_K^{\frac{1}{4}-9\varepsilon}$, is a function of $L$, of $C_6$ and $\kappa_{(0)}$ (functions of $\rho$ and of the weights, line (C0r); by the convention of Section~I.8.1 the dependence on the fixed weights is not indicated), of the exponent $\varepsilon$ (line (C7)) and of $e_K$.
The last member of (10.3), $C_{\mathrm{mat}}$, is a function of $L$ and of $C_6$, $\kappa_{(0)}$, that is, of $(L,\rho)$.
Finally, for all $N\ge N_0(e)$,
\[
D^{\mathrm{mat}}\le\left|D^{\mathrm{mat}}\right|\le\left(L^3+1\right)\left(\varepsilon_5\left(\bar{m}\cdot R^{*}(e)\right)+3\left(e\cdot h^{\frac{1}{2}}\right)^{\frac{1}{4}-9\varepsilon}\right)\le C_{\mathrm{mat}}.
\]
This is (R2), and its middle member is the form (G6.T-32) of Section~G6.Z1.
\end{proof}

\paragraph{Uniformity.}
The inequalities (10.1) and (10.2) hold for every $\bar{m}\in(0,1]$, every $e$ with $e\cdot h^{\frac{1}{2}}\le e_1$, every $R\in\mathrm{Tor}(e)$ and all $N\ge N_{\min}(R,e)$.
Let $\bar{m}\in(0,1]$ and $e$ with $e\cdot h^{\frac{1}{2}}\le e_1$ be given, and let $N\ge N_{\min}\left(R_{\mathrm{top}}(e),e\right)$.
Then (10.1) and (10.2) hold with one real number $p_F$ and one lattice mass $\bar{m}\cdot\eta+m_{0,N}(e)$ for all $R\in\mathrm{Tor}(e)$ with $N_{\min}(R,e)\le N$.
This set of tori contains $R_{\mathrm{top}}(e)$, and by (Sc2) it contains, together with any of its elements $R$, every $R'\in\mathrm{Tor}(e)$ with $R'\ge R$.
The inequalities (10.3) hold for every $\bar{m}\in(0,1]$, every $e$ with $e\cdot h^{\frac{1}{2}}\le e_1$ and all $N\ge N_0(e)=N_{\min}\left(R^{*}(e),e\right)$.
For these $N$ both $R^{*}(e)$ and $L\cdot R^{*}(e)=R_{\mathrm{top}}(e)$ belong to the set of all $R\in\mathrm{Tor}(e)$ with $N_{\min}(R,e)\le N$: both lie in $\mathrm{Tor}(e)$, and $N_{\min}\left(L\cdot R^{*}(e),e\right)\le N_{\min}\left(R^{*}(e),e\right)\le N$ by (Sc2) (Step~3).
The bound $\varepsilon_5\left(\frac{1}{L}\right)+3\,e_K^{\frac{1}{4}-9\varepsilon}$ is valid for all the data above simultaneously, and the bound $C_{\mathrm{mat}}$ on $\left|D^{\mathrm{mat}}\right|$ is valid simultaneously for all $\bar{m}$ and $e$ as above and all $N\ge N_0(e)$.
The first of them contains neither $N$ nor $R$, and it depends on $\left(e,\bar{m}\right)$ only through $e_K=e\cdot h^{\frac{1}{2}}\le e_1$; the constant $C_{\mathrm{mat}}$ contains none of $N$, $R$, $e$, $\bar{m}$.
The constants in the bounds are $C_6$, $\kappa_{(0)}$ (line (C0r)), $L$ (line (C3)) and $\varepsilon$ (line (C7)).
In addition, the proof involves $c_{\mathrm{fl}}$ (line (C0)), $C_0$, $c_{\kappa}$ (line (C6)), $M$, $r$ (lines (C5), (C7)) and $e_1$ (line (C9)); they enter through the second clause of (T2), through (32.2) and through (Sc1)--(Sc3), (Sc5).
All of them are constants of the lines (C0)--(C9), and by (I.U4)$^{\mathrm{ch}}$ (Proposition~14.4\,(a)) none of them depends on $N$, on a level or a side, on $\theta$, on the pair or on $\bar{m}$.
The constant $C_{\mathrm{mat}}$ belongs to the line (C10) and is formed from $L$, $C_6$, $\kappa_{(0)}$ alone.
The inequalities (10.1)--(10.3) hold whatever the sign of the lattice mass $\bar{m}\cdot\eta+m_{0,N}(e)$ (I.K10).

\subsection*{4. Conditions on the constants}

The arguments of this chapter use the following conditions on the constants.
The conditions on $e_1$ are imposed for all $v\ge-\log e_1$, those on $e_0$ for all $v\ge-\log e_0$ (Proposition~14.2\,(b)); here $\mathrm{lh}(v)=L\cdot M\cdot v^r$.
\begin{itemize}
\item
$8\,C_0^2\cdot\mathrm{lh}(v)^6\left(b_{-}^{-1}+L\cdot\mathrm{lh}(v)\right)\le\exp\left(\frac{\varepsilon\cdot v}{2}\right)$.
\emph{Imposed in:} Proposition~14.2\,(b), first inequality of (e3); a condition on $e_1$, line (C9).
\emph{Used in:} the first clause of (T2): Section~1, Step~6 ($\mathrm{crw}\le\mathrm{aw}_2$).
\item
$4\left(2C_0+c_{\kappa}+5\,c_{\mathrm{fl}}+6\right)\mathrm{lh}(v)^3\le\exp\left(\varepsilon\cdot v\right)$.
\emph{Imposed in:} Proposition~14.2\,(b), second inequality of (e3); line (C9).
\emph{Used in:} the second clause of (T2): Section~1, Step~6 (32.15); Section~3, Step~1 (32.20).
\item
$\exp\left(-\left(\frac{1}{4}-9\varepsilon\right)v\right)\le\frac{1}{8}$.
\emph{Imposed in:} Proposition~14.2\,(b), third inequality of (e3); line (C9).
\emph{Used in:} at $v=v_{*}$: the third clause of (T2), (32.a), (32.b): Section~1, Steps~2, 4, 6; Section~2, Steps~1, 3.
At $v=v_K$: (32.2): Section~3, Step~4.
\item
$4\,C_0\cdot\mathrm{lh}(v)^2\cdot v^p<\exp\left(\varepsilon\cdot v\right)$.
\emph{Imposed in:} Proposition~14.2\,(b), (e4); line (C9).
\emph{Used in:} (T3): Section~1, Step~1 (32.8).
\item
$12\,\mathrm{lh}(v)^3\left(\exp\left(-\frac{v^{p_0}}{2}\right)+\exp\left(-\frac{v^{2p}}{2C_0^2}\right)\right)\le\exp(-v)$.
\emph{Imposed in:} Proposition~14.2\,(b), (e5); line (C9).
\emph{Used in:} (T4), (32.b): Section~2, Steps~2, 3.
\item
$C_{\Sigma}\cdot\exp\left(-\gamma_{\Sigma}\cdot M\cdot v^r\right)\le\exp(-v)$.
\emph{Imposed in:} Proposition~14.2\,(b), (e6); line (C9).
\emph{Used in:} (T5), (32.b): Section~2, Step~3.
\item
$v\ge8$; $v^r>31L+16$; $\left(1+\frac{3\,\mathrm{lg}}{v}\right)^r<L$.
\emph{Imposed in:} Proposition~14.2\,(b), (e1); line (C9).
\emph{Used in:} through (Sc1)--(Sc5), that is, through the hypotheses $v_K\ge1$, $\left(1+\frac{\mathrm{lg}}{v_K}\right)^r\le L$, $r_{*}>31L+16$ (Section~G4.0.3) and $\left(1+\frac{3\,\mathrm{lg}}{v_K}\right)^r<L$ of \cite[Proposition~19.1]{QEDI}: Section~0.1; Section~3, Steps~0, 3.
\item
$L\ge33$ (hence $L\ge13$).
\emph{Imposed in:} Proposition~14.2\,(f) ($L\ge2400^8$, from (L1)); line (C3).
\emph{Used in:} through (Sc4) (in the proof of \cite[Proposition~19.1(f)]{QEDI}: $31L+16<L^2$; $L\ge13$): Section~0.1.
\item
$e_1\le e_0$; $e_0<1$.
\emph{Imposed in:} line (C9); Proposition~14.2\,(a), (f).
\emph{Used in:} $e_{*}\le e_0<1$: the data (SF) at $k_{*}$; \cite[Lemma~17.3]{QEDI}, \cite[Lemma~17.4(a)]{QEDI}, \cite[Lemma~17.6]{QEDI}, \cite[Proposition~15.30(i)]{QEDI}.
\item
$v\ge1$, that is, $e_0\le\exp(-1)$.
\emph{Imposed in:} Proposition~14.2\,(b), (Z3.s1); a condition on $e_0$, line (C8).
\emph{Used in:} the standing condition $e_0\le\exp(-1)$ of \cite[Section~15]{QEDI}, the section of \cite[Theorem~15.14]{QEDI} and \cite[Proposition~15.30]{QEDI}: Section~0.4; Section~1, Step~1.
\item
$8\,C_B\cdot v^p\le\exp\left(\left(\frac{3}{4}-3\varepsilon\right)v\right)$; $8\exp\left(-\left(\frac{3}{4}-\frac{7\varepsilon}{2}\right)v\right)\le1$.
\emph{Imposed in:} Proposition~14.2\,(b), third and fourth inequality of (A10); conditions on $e_0$, line (C8).
\emph{Used in:} the hypotheses of \cite[Proposition~15.30(i)]{QEDI}: Section~1, Step~1.
\item
$v^{p_0}\le\exp\left(\frac{\varepsilon\cdot v}{4}\right)$; $2\,C_w^2\cdot c_{\gamma}^2\cdot c_{C}\cdot\exp\left(-\frac{\varepsilon\cdot v}{4}\right)\le1$.
\emph{Imposed in:} Proposition~14.2\,(b), first inequality of (A1), and (A2); line (C8).
\emph{Used in:} the hypotheses of \cite[Proposition~15.30(i)]{QEDI}.
\item
$12\exp\left(-\frac{v^{2p_0}}{2}\right)\le\exp(-7v)$; $\exp\left(-\frac{5\varepsilon\cdot v}{2}\right)\le\frac{1}{2}$; $c_{\mathrm{fl}}\cdot\exp\left(-\varepsilon\cdot v\right)\le1$; $3M\cdot\exp\left(-\varepsilon\cdot v\right)\le1$.
\emph{Imposed in:} Proposition~14.2\,(b), second to fifth inequality of (A1); conditions on $e_0$, line (C8).
\emph{Used in:} with the first inequality of (A1), these inequalities form condition (A1) of \cite[Section~5.6.2]{QEDI}, which is a hypothesis of \cite[Proposition~15.30(i)]{QEDI}; written at $e_k$, it is among the conditions named at the beginning of \cite[Section~17.3]{QEDI}, the section of \cite[Lemma~17.3]{QEDI}, \cite[Lemma~17.4(a)]{QEDI} and \cite[Lemma~17.6]{QEDI}.
\item
$C_{P}\le C_{w}$; $\gamma_{\mathrm{w}}\le\gamma_{P}$.
\emph{Imposed in:} Proposition~14.2\,(f) (definitions of the lines (C4r), (C6)).
\emph{Used in:} the hypotheses of \cite[Proposition~15.30(i)]{QEDI}.
\item
$\gamma_{\mathrm{w}}>0$.
\emph{Imposed in:} Proposition~14.2\,(a), definition of $\gamma_{\mathrm{w}}$; it also follows from (M3) of Proposition~14.2\,(b), $\gamma_{\mathrm{w}}\cdot M\ge16$, as $M>0$.
\emph{Used in:} the hypotheses of \cite[Proposition~15.30(i)]{QEDI}, through the constant $c_{\gamma}$ of (A2), which is formed with $\gamma_{\mathrm{w}}$.
\item
$L^{\frac{3}{4}}\ge\frac{8}{3}$.
\emph{Imposed in:} Proposition~14.2\,(f) ($L^{\frac{3}{4}}\ge110$, from (L1)); line (C3).
\emph{Used in:} the hypotheses of \cite[Proposition~15.30(i)]{QEDI}.
\item
$\exp\left(-\varepsilon\cdot v\right)\le\frac{1}{4}$; $\exp\left(-\frac{v}{4}\right)+\exp\left(-\varepsilon\cdot v\right)\le1$; $L\ge5$.
\emph{Imposed in:} Proposition~14.2\,(b), first and fourth inequality of (A8); (L5).
\emph{Used in:} \cite[Lemma~17.4(a)]{QEDI}: Section~1, Steps~2, 7.
\item
$C_0\ge C_{P}\cdot c_{\mathrm{bd}}(\gamma_{P})$; $C_0\ge1$.
\emph{Imposed in:} the definition of $C_0$, line (C6); Proposition~14.2\,(f).
\emph{Used in:} \cite[Proposition~5.32]{QEDI} in the form (32.8); $\mathrm{aw}_2\le\mathrm{aw}_1$; $\mathrm{crw}\le\mathrm{aw}_2$.
\item
$c_{\kappa}<\infty$ ($c_{\kappa}\le K_{0}$).
\emph{Imposed in:} line (C6); Proposition~14.2\,(f).
\emph{Used in:} Section~1, Step~5\,(iii).
\item
$c_{\mathrm{fl}}=61\,c_{X}\cdot K_{0}>0$; $c_{\kappa}\ge0$.
\emph{Imposed in:} the definitions, lines (C0), (C6) (Proposition~14.2\,(a)).
\emph{Used in:} Section~0.2, second clause of (T2); $2C_0+c_{\kappa}+4\le2C_0+c_{\kappa}+5\,c_{\mathrm{fl}}+6$ in (32.15); $5\,c_{\mathrm{fl}}\le2C_0+c_{\kappa}+5\,c_{\mathrm{fl}}+6$ in (32.20).
\item
$\gamma_{\Sigma}\in(0,1]$, $0\le C_{\Sigma}<\infty$; $\kappa_{(0)}>0$, $0\le C_6<\infty$.
\emph{Imposed in:} Proposition~30.3\,(e) with (8.5); the constants of Proposition~30.1, with (8.1); lines (C6), (C0r).
\emph{Used in:} (T5) in Section~0.2; Section~3, Step~2 ($\varepsilon_5$ is decreasing).
\item
$0<\alpha<1$; $0<\gamma_{P}\le1$.
\emph{Imposed in:} line (C1), Proposition~14.2\,(f); \cite[Section~15.7]{QEDI}, Proposition~14.2\,(f).
\emph{Used in:} the kernel bounds of Section~0.4, which follow from \cite[Theorem~15.14]{QEDI} (see \cite[Proposition~5.32]{QEDI}) together with the lattice sum bound of \cite[Lemma~15.2]{QEDI}, applied at the rate $\gamma_{P}$ to sum the kernel bound over the unit bonds $b'$.
The theorem requires $0\le\alpha<1$, and the lattice sum bound requires a rate in $(0,1]$.
\item
$0<\varepsilon\le\frac{1}{200}$.
\emph{Imposed in:} line (C7) and Proposition~14.2\,(b), (S1), which give $1-\alpha<\varepsilon\le\bar{\varepsilon}$;
line (C1), which gives $\alpha<1$, so that $\varepsilon>1-\alpha>0$;
line (C0), in which $\bar{\varepsilon}:=\frac{1}{200}$ is defined.
\emph{Used in:} $\frac{1}{4}-9\varepsilon>0$; $\frac{1}{4}-\frac{3\varepsilon}{2}>0$; $\frac{3}{4}-4\varepsilon\ge\frac{1}{4}-9\varepsilon$; the hypothesis $\varepsilon\in\left(0,\bar{\varepsilon}\right]$ of \cite[Proposition~15.30(i)]{QEDI}.
\item
$0<b_{-}\le b_{k}\le b_{+}$ for $k\ge1$.
\emph{Imposed in:} Section~I.3.4; line (C4).
\emph{Used in:} (32.10); Proposition~30.3\,(e) at $\mathsf{b}=b_{k_{*}}$.
\item
$L$ odd.
\emph{Imposed in:} Proposition~14.2\,(b), (L5); line (C3).
\emph{Used in:} $\mathfrak{N}=L^{k_{*}}$ is odd (Proposition~30.3); (Sc4) ($\ell_{*}/(L\cdot M)$ is an odd integer); \cite[Propositions~15.5 and~5.32]{QEDI} (by \cite[Section~15.1.1 and Theorem~15.14]{QEDI}: $L$ odd, $L\ge3$); \cite[Theorem~14.17(b)--(e)]{QEDI} (blocks of the odd side $L^j$ in \cite[Section~14]{QEDI}).
\item
$M$ a power of $L$, $M\ge L^2$.
\emph{Imposed in:} Proposition~14.2\,(b), (M1); line (C5).
\emph{Used in:} through (Sc1)--(Sc4) ($\mathfrak{n}_k$ is an integer); $T^{*}(e)\ge1$ (Section~I.6.2, (Sc5)): Section~3, Step~3; $L\cdot M\cdot r_0\ge4$ in (32.1).
\item
$C_{\mathrm{mat}}=\left(L^3+1\right)\left(\varepsilon_5\left(\frac{1}{L}\right)+1\right)$.
\emph{Imposed in:} line (C10) (a definition; it contains $L$, $C_6$, $\kappa_{(0)}$ only).
\emph{Used in:} Section~3, Step~4.
\end{itemize}
\emph{Order.}
Besides the constant on which it is imposed and the variable $v$, every condition of the list contains only constants of earlier lines of Section~I.8.2 (Proposition~14.2\,(e)); the constant $C_{\mathrm{mat}}$ is formed from constants of the lines (C0r) and (C3).
No condition contains $N$, $\theta$, $R$, $\bar{m}$ or the pair.

\section*{Chapter 49. Symmetry of the odd reference torus under the reflections in its walls: proof of Proposition~49}

This chapter contains the proof of Proposition~49.
Section~0 fixes the setting and establishes five lemmas on sites, cochains, spinors and the operator;
Sections~1--6 prove the clauses (a)--(f) in turn, and Section~7 lists the conditions on the constants.
Sections~0, 1 and 3--6 are valid for every integer $\mathrm{kt}\ge1$, even or odd, whereas Section~2 concerns $\mathrm{kt}=1$.

The clauses of Proposition~49 assert identities, memberships, facts on group actions, signs of numbers and the analyticity of functions of $t$;
none of them contains a bound or a constant, and each of them is proved for one fixed tuple $(\bar{m},e,N,\mathrm{kt})$.
The only majorant of the proof, the bound on $\left|\mathrm{Igd}(F^{\mathrm{u}},\theta,t)\right|$ displayed in 5.2, is used at fixed $(\bar{m},e,N,\mathrm{kt},\boldsymbol{\mu})$, locally uniformly in $t$, to secure the convergence of the integrals.

\subsection*{0. Conventions of this chapter; five lemmas}

\subsubsection*{0.1. Setting (lattice units)}

All objects are those of Sections~I.1 and I.2 at $k=0$, with $\eta_0=1$ and with $q$ in the place of $e_k$ (Section~I.2.5; Section~I.2.9, the first paragraph of Proposition~46).
They are formed on the cubic torus with $n=\mathrm{kt}\cdot n_*$ photon sites per side.

\emph{Scales.}
Since $\bar{m}\le1$ (Section~I.0.5), the number $h=L^{\left\lfloor\log_L\frac{1}{\bar{m}}\right\rfloor}$ is a power of $L$ with a non-negative exponent.
By Section~I.0.5 and the line (C5) of Section~I.8.2, $M$ is a power of $L$ with $M\ge L^2$;
by Section~I.6.1, $r_K=L^i$ for an integer $i\ge0$.
Hence
\[
T^{*}(e)=\log_L\left(M\cdot r_K\right)-1
\]
is an integer, not smaller than $2+0-1=1$ (Section~I.6.2), and $n_*=h\cdot L^{T^{*}(e)+N}$ (Section~I.6.4) is a power of $L$ with an exponent at least $1$.
It follows that $\mathrm{per}=n_*/L$ is a power of the odd integer $L$ with a non-negative exponent, that is, a positive odd integer, and that
\begin{equation}\tag{0.1}
n=\mathrm{kt}\cdot n_*=\mathrm{kt}\cdot L\cdot\mathrm{per},\qquad n\mathbb{Z}^3\subset\mathrm{per}\cdot\mathbb{Z}^3,\qquad n\ge L\ge5.
\end{equation}

\emph{Fermion sites.}
A fermion site is a pair $x=(c,g)\in\mathbb{Z}^3\times O_{\mathrm{cub}}$;
one writes $\pi(x)=c$, and $x+v=(c+v,g)$ for $v\in\mathbb{Z}^3$.
At $\eta_0=1$ the display (I.1.2) reads
\[
x^{(i)}=(c,gs_i)\quad(i=0,1,2),\qquad x^{(3)}=\left(c+ge_1,\,gs_3\right),\qquad\hat{d}_i(x)=-g\hat n_i.
\]
The relations $s_i^2=1$, $s_i\hat n_i=-\hat n_i$, $s_3e_1=-e_1$ and (I.1.1) yield
\begin{align}
&\left(x^{(i)}\right)^{(i)}=x\quad\left(\text{for $i=3$: }\left(x^{(3)}\right)^{(3)}=\left(c+ge_1+gs_3e_1,\,gs_3^2\right)=x\right),\ \text{and}\notag\\
&\qquad x\mapsto x^{(i)}\ \text{is a bijection of $\mathbb{Z}^3\times O_{\mathrm{cub}}$ with}\ (x+v)^{(i)}=x^{(i)}+v;\tag{I.s1}\\
&\hat{d}_i\left(x^{(i)}\right)=-\hat{d}_i(x);\tag{I.s2}\\
&\sum_{i=0}^{3}w_i\cdot\hat{d}_i(x)=-g\sum_{i}w_i\cdot\hat n_i=0.\tag{I.s3}
\end{align}

\emph{Cochains.}
A $j$-cochain on $\mathbb{Z}^3/n\mathbb{Z}^3$ is an $n\mathbb{Z}^3$-periodic $j$-cochain on $\mathbb{Z}^3$, and $C^{j}(n)$ is the space of the real ones.
The number $\langle u,v\rangle$ is the sum of $u(\boldsymbol{\sigma})\cdot v(\boldsymbol{\sigma})$ over the positively oriented $j$-cells $\boldsymbol{\sigma}$ modulo $n\mathbb{Z}^3$ (Section~I.7.1);
this pairing is extended bilinearly to $C^{j}(n)\otimes\mathbb{C}$.

A one-cochain $A$ is a function on the oriented bonds $(c,c')$, $|c-c'|_2=1$, with $A(c',c)=-A(c,c')$;
one writes $A_\mu(c)=A(c,c+e_\mu)$ and puts $A(c,c):=0$ (Section~I.1.1).

An oriented plaquette is a quadruple $\mathrm{pl}=(c^1,c^2,c^3,c^4)$ of points of $\mathbb{Z}^3$ such that
\[
c^2-c^1=c^3-c^4\in\{\pm e_\mu\},\qquad c^4-c^1=c^3-c^2\in\{\pm e_\nu\},\qquad\mu\ne\nu;
\]
these are the corners of a unit square, listed in cyclic order.
The cyclic rotations of $\mathrm{pl}$ are $\mathrm{pl}$ itself, $(c^2,c^3,c^4,c^1)$, $(c^3,c^4,c^1,c^2)$ and $(c^4,c^1,c^2,c^3)$;
its reversal is $(c^1,c^4,c^3,c^2)$.
All these quadruples satisfy the two conditions again, with $\mu$ and $\nu$ interchanged for the rotations by one and by three places and for the reversal.
A unit square carries eight such quadruples, namely the four cyclic rotations of
\[
(c;\mu\nu):=\left(c,\,c+e_\mu,\,c+e_\mu+e_\nu,\,c+e_\nu\right),\qquad\mu<\nu,
\]
where $c$ is the corner with the least coordinates (Section~I.1.1), and the four cyclic rotations of the reversed quadruple $\left(c,\,c+e_\nu,\,c+e_\mu+e_\nu,\,c+e_\mu\right)$.
A two-cochain $u$ is regarded as the function with $u(\mathrm{pl}):=u_{\mu\nu}(c)$ on the first four and $u(\mathrm{pl}):=-u_{\mu\nu}(c)$ on the last four.

The coboundary is given by
\begin{equation}\tag{0.2}
\begin{aligned}
(d\boldsymbol{\lambda})(c,c')&=\boldsymbol{\lambda}(c')-\boldsymbol{\lambda}(c),\\
(dA)(c^1,c^2,c^3,c^4)&=A(c^1,c^2)+A(c^2,c^3)+A(c^3,c^4)+A(c^4,c^1).
\end{aligned}
\end{equation}
Since $A$ is odd, the second expression is invariant under cyclic rotations and changes its sign under reversal.
At $\mathrm{pl}=(c;\mu\nu)$ it equals $A_\mu(c)+A_\nu(c+e_\mu)-A_\mu(c+e_\nu)-A_\nu(c)$, which is the display of Section~I.1.1;
thus (0.2) is the coboundary of Section~I.1.1.

\emph{Phases.}
Let $A$ be a complex one-cochain on $\mathbb{Z}^3/n\mathbb{Z}^3$ and $i\in\{0,\dots,3\}$.
By Section~I.1.8 (the display (I.1.6) and the sentence following it: for nearest neighbours the six paths reduce to the bond, and $U_A(c,c)=1$),
\[
U_A\left(x,x^{(i)}\right)=\exp\left(\mathrm{i}\,q\cdot A\left(\pi x,\pi x^{(i)}\right)\right).
\]
For $i\le2$ this number is $1$, because then $\pi x^{(i)}=\pi x$.
As $A$ is odd and $n\mathbb{Z}^3$-periodic,
\begin{equation}\tag{0.3}
U_A\left(x^{(i)},x\right)=U_{-A}\left(x,x^{(i)}\right),\qquad U_A\left(x+ne_\mu,x^{(i)}+ne_\mu\right)=U_A\left(x,x^{(i)}\right).
\end{equation}

\emph{Spinors} (Sections~I.1.5 and I.2.9).
The space
\[
\mathfrak{S}^{\mathrm{ch}}(1;n)=\left\{f:\mathbb{Z}^3\times O_{\mathrm{cub}}\to\mathbb{C}^4\ :\ f(c+ne_\mu,g)=-f(c,g),\ \mu=1,2,3\right\}
\]
has the dimension $192\,n^3=4\cdot48\cdot n^3$.

\emph{Operator.}
By (I.2.1), $M_i(x)=\frac{w_i}{2}\left(\rho-\gamma\cdot\hat{d}_i(x)\right)$; this matrix depends on the label $g$ only.
By (I.s3),
\[
\sum_{i=0}^{3}M_i(x)=\frac{\rho}{2}\cdot\mathrm{sw}-\frac{1}{2}\,\gamma\cdot\left(\sum_iw_i\cdot\hat{d}_i(x)\right)=\frac{\rho}{2}\cdot\mathrm{sw},
\]
so that (I.2.3) takes the form
\begin{equation}\tag{0.4}
\begin{split}
(D_Af)(x)&=-\sum_{i=0}^{3}M_i(x)\cdot\left[U_A\left(x,x^{(i)}\right)\cdot f\left(x^{(i)}\right)-f(x)\right]\\
&=\frac{\rho}{2}\cdot\mathrm{sw}\cdot f(x)-\sum_{i=0}^{3}M_i(x)\cdot U_A\left(x,x^{(i)}\right)\cdot f\left(x^{(i)}\right).
\end{split}
\end{equation}
The operator $D_A$ maps $\mathfrak{S}^{\mathrm{ch}}(1;n)$ into itself:
$(x+ne_\mu)^{(i)}=x^{(i)}+ne_\mu$ by (I.s1), $M_i$ does not depend on $c$, and (0.3) holds.

\emph{The group.}
Let $\mathrm{Gc}=\mathbb{Z}^3\rtimes\{\pm1\}^3$ (Section~I.4.7).
An element $g_{\mathrm{w}}=(a,r)$, with $a\in\mathbb{Z}^3$ and $r\in\{\pm1\}^3$, is the map $y\mapsto ry+a$ of $\mathbb{R}^3$.
One writes $re_\mu=\epsilon_\mu e_\mu$ with $\epsilon_\mu\in\{1,-1\}$.
Since $r^2=1$,
\[
(a,r)(a',r')=\left(a+ra',\,rr'\right),\qquad(a,r)^{-1}=(-ra,\,r).
\]

\subsubsection*{0.2. Notation}

In this chapter $\epsilon_\mu$ is always the sign of 0.1, $re_\mu=\epsilon_\mu e_\mu$.
Further,
\[
\mathrm{rf}_\mu:=1-2e_\mu e_\mu^T\in\{\pm1\}^3
\]
(Section~I.4.8 denotes this matrix by $r_\mu$), and $\mathrm{thw}_{\mu,j}$ is the reflection in the wall $\{y_\mu=j\cdot\mathrm{hc}\}$.
The letter $\mathrm{jv}$ stands for a vector of $\mathbb{Z}^3$, and $\mathrm{pl}$ for an oriented plaquette.
One puts
\[
\mathrm{Gp}:=\left\{(a,r)\ :\ a\in\mathrm{per}\cdot\mathbb{Z}^3,\ r\in\{\pm1\}^3\right\},\qquad\mathrm{Yr}:=\{0,\dots,n-1\}^3\times O_{\mathrm{cub}};
\]
thus $\mathrm{Yr}$ is a set of $48\,n^3$ points of $\mathbb{Z}^3\times O_{\mathrm{cub}}$, and $\sum_{x\bmod n}$ denotes the sum over $x\in\mathrm{Yr}$, a sum of $48\,n^3$ terms.
For a cell $\mathfrak{c}$ of the torus, $\mathrm{ue}_{\mathfrak{c}}$ is the unit vector of this cell, as in Section~G5.0.3:
the vector with one entry for each cell, equal to $1$ at $\mathfrak{c}$ and to $0$ at all other cells.
In Section~5, $a(t)$ is an external connection, and $B(F^{\mathrm{u}},\theta,t)$ is a connection that serves as the argument of the operator.

\begin{named}{Lemma 0.1}[fermion sites]
For $g_{\mathrm{w}}=(a,r)\in\mathrm{Gc}$ put $g_{\mathrm{w}}(c,g):=(rc+a,\,rg)$ (Section~I.4.7).
Then the following hold.
\begin{itemize}
\item[(i)] $g_{\mathrm{w}}\left(g'_{\mathrm{w}}x\right)=\left(g_{\mathrm{w}}g'_{\mathrm{w}}\right)x$;
$g_{\mathrm{w}}$ is a bijection of $\mathbb{Z}^3\times O_{\mathrm{cub}}$;
$\pi(g_{\mathrm{w}}x)=g_{\mathrm{w}}\pi(x)$.
\item[(ii)] $g_{\mathrm{w}}(x+v)=g_{\mathrm{w}}x+rv$ for $v\in\mathbb{Z}^3$;
in particular $g_{\mathrm{w}}(x+ne_\mu)=g_{\mathrm{w}}x+n\,re_\mu=g_{\mathrm{w}}x+n\epsilon_\mu e_\mu$.
\item[(iii)] $\left(g_{\mathrm{w}}x\right)^{(i)}=g_{\mathrm{w}}\left(x^{(i)}\right)$ and $\hat{d}_i(g_{\mathrm{w}}x)=r\hat{d}_i(x)$ for $i=0,\dots,3$.
\item[(iv)] If $r\ne1$, then $g_{\mathrm{w}}x$ and $x$ carry different labels;
hence $g_{\mathrm{w}}$ fixes no point of $\mathbb{Z}^3\times O_{\mathrm{cub}}$ and no point of $\left(\mathbb{Z}^3\times O_{\mathrm{cub}}\right)/n\mathbb{Z}^3$.
\end{itemize}
\end{named}

\begin{proof}
As $\{\pm1\}^3$ is a subgroup of $O_{\mathrm{cub}}$ (Section~I.1.2), $rg$ lies in $O_{\mathrm{cub}}$.

(i) For $g'_{\mathrm{w}}=(a',r')$ one computes
\[
g_{\mathrm{w}}\left(g'_{\mathrm{w}}(c,g)\right)=\left(r(r'c+a')+a,\,rr'g\right)=\left((rr')c+(a+ra'),\,(rr')g\right);
\]
the element $(0,1)$ acts as the identity, and therefore $g_{\mathrm{w}}^{-1}$ inverts $g_{\mathrm{w}}$.

(ii) $\left(r(c+v)+a,\,rg\right)=(rc+a,\,rg)+rv$.

(iii) For $i\le2$ one has $\left(g_{\mathrm{w}}x\right)^{(i)}=\left(rc+a,\,rgs_i\right)=g_{\mathrm{w}}(c,gs_i)$.
For $i=3$,
\[
\left(g_{\mathrm{w}}x\right)^{(3)}=\left(rc+a+(rg)e_1,\,rgs_3\right)=\left(r(c+ge_1)+a,\,r(gs_3)\right)=g_{\mathrm{w}}\left(x^{(3)}\right).
\]
Finally $\hat{d}_i(g_{\mathrm{w}}x)=-(rg)\hat n_i=r\left(-g\hat n_i\right)$.

(iv) Since $g$ is invertible, $rg=g$ forces $r=1$;
and translations by $n\mathbb{Z}^3$ leave the label unchanged.
\end{proof}

\begin{named}{Lemma 0.2}[cochains]
Let $g_{\mathrm{w}}=(a,r)\in\mathrm{Gc}$.
For a cell $\boldsymbol{\sigma}$ (a site $c$, an oriented bond $(c,c')$, or an oriented plaquette $(c^1,\dots,c^4)$) let $g_{\mathrm{w}}\boldsymbol{\sigma}$ be the tuple of the images of its points.
It is a cell of the same kind, because $g_{\mathrm{w}}$ is affine with $re_\mu=\pm e_\mu$.
The map $\boldsymbol{\sigma}\mapsto g_{\mathrm{w}}\boldsymbol{\sigma}$ commutes with the reversal of a bond, and with the cyclic rotations and the reversal of a quadruple.
Put
\[
\left(g_{\mathrm{w}}^*u\right)(\boldsymbol{\sigma}):=u(g_{\mathrm{w}}\boldsymbol{\sigma}),\qquad g_{\mathrm{w}*}:=\left(g_{\mathrm{w}}^{-1}\right)^*.
\]
By the preceding sentence $g_{\mathrm{w}}^*u$ is odd under reversal and, for $j=2$, invariant under cyclic rotations;
so it is a $j$-cochain.
(This is the pull-back of Section~I.4.7: for a positively oriented $\boldsymbol{\sigma}$ the sign $\epsilon(g_{\mathrm{w}},\boldsymbol{\sigma})$ is $+1$ or $-1$ according as the tuple $g_{\mathrm{w}}\boldsymbol{\sigma}$ is positively or negatively oriented.)
Then, for $j=0,1,2$:
\begin{itemize}
\item[(i)] $g_{\mathrm{w}}^*$ and $g_{\mathrm{w}*}$ map $n\mathbb{Z}^3$-periodic cochains to $n\mathbb{Z}^3$-periodic cochains and real ones to real ones, and they commute with complex conjugation.
\item[(ii)] $\left(g_{\mathrm{w}}g'_{\mathrm{w}}\right)^*=\left(g'_{\mathrm{w}}\right)^*g_{\mathrm{w}}^*$;
$\left(g_{\mathrm{w}}g'_{\mathrm{w}}\right)_*=g_{\mathrm{w}*}\,g'_{\mathrm{w}*}$;
$g_{\mathrm{w}*}\,g_{\mathrm{w}}^*=g_{\mathrm{w}}^*\,g_{\mathrm{w}*}=1$.
\item[(iii)] $d\,g_{\mathrm{w}}^*=g_{\mathrm{w}}^*\,d$ and $d\,g_{\mathrm{w}*}=g_{\mathrm{w}*}\,d$ on zero- and one-cochains.
\item[(iv)] $\left\langle g_{\mathrm{w}}^*u,g_{\mathrm{w}}^*v\right\rangle=\langle u,v\rangle=\left\langle g_{\mathrm{w}*}u,g_{\mathrm{w}*}v\right\rangle$ for $u,v\in C^{j}(n)\otimes\mathbb{C}$.
Hence $g_{\mathrm{w}*}$ is orthogonal on $C^{j}(n)$, and $\left\langle g_{\mathrm{w}*}u,v\right\rangle=\left\langle u,g_{\mathrm{w}}^*v\right\rangle$.
\item[(v)] Let $\mathrm{Hcn}$ be the constant one-cochain with $\mathrm{Hcn}(c,c+e_\mu)=\mathrm{hv}_\mu$, where $\mathrm{hv}\in\mathbb{C}^3$.
Then $g_{\mathrm{w}*}\mathrm{Hcn}$ is the constant one-cochain with the components $\epsilon_\mu\cdot\mathrm{hv}_\mu=(r\,\mathrm{hv})_\mu$.
In particular $g_{\mathrm{w}*}H_{\theta}=H_{r\theta}$ for $\theta\in\mathbb{R}^3$ (Section~I.6.6).
Moreover
\[
g_{\mathrm{w}*}{V_{\mathrm{g}}}={V_{\mathrm{g}}},\qquad g_{\mathrm{w}*}{V_{\mathrm{h}}}={V_{\mathrm{h}}},\qquad g_{\mathrm{w}*}{V_{\mathrm{t}}}={V_{\mathrm{t}}},\qquad g_{\mathrm{w}*}F_N=F_N;
\]
$A\left(g_{\mathrm{w}*}F^{\mathrm{u}}\right)=g_{\mathrm{w}*}A(F^{\mathrm{u}})$ and $A\left(-F^{\mathrm{u}}\right)=-A(F^{\mathrm{u}})$ for $F^{\mathrm{u}}\in F_N$;
$d:{V_{\mathrm{t}}}\to F_N$ is a linear bijection whose inverse is $F^{\mathrm{u}}\mapsto A(F^{\mathrm{u}})$, which is therefore linear, and $A(dA)=A$ for $A\in{V_{\mathrm{t}}}$;
the measure $\gamma_N$ is invariant under $F^{\mathrm{u}}\mapsto g_{\mathrm{w}*}F^{\mathrm{u}}$ and under $F^{\mathrm{u}}\mapsto-F^{\mathrm{u}}$.
\item[(vi)] In components,
\begin{gather*}
\left(g_{\mathrm{w}}^*A\right)_\mu(c)=\epsilon_\mu\cdot A_\mu\left(g_{\mathrm{w}}c-\mathbf{1}[\epsilon_\mu=-1]\,e_\mu\right),\\
\left(g_{\mathrm{w}}^*u\right)_{\mu\nu}(c)=\epsilon_\mu\epsilon_\nu\cdot u_{\mu\nu}\left(g_{\mathrm{w}}c-\mathbf{1}[\epsilon_\mu=-1]\,e_\mu-\mathbf{1}[\epsilon_\nu=-1]\,e_\nu\right)\qquad(\mu<\nu).
\end{gather*}
\end{itemize}
\end{named}

\begin{proof}
(i) One has $g_{\mathrm{w}}(\boldsymbol{\sigma}+n\,\mathrm{jv})=g_{\mathrm{w}}\boldsymbol{\sigma}+n\,r\,\mathrm{jv}$ with $r\,\mathrm{jv}\in\mathbb{Z}^3$;
therefore
\[
\left(g_{\mathrm{w}}^*u\right)(\boldsymbol{\sigma}+n\,\mathrm{jv})=u\left(g_{\mathrm{w}}\boldsymbol{\sigma}+n\,r\,\mathrm{jv}\right)=u(g_{\mathrm{w}}\boldsymbol{\sigma}).
\]
For real $u$ the number $\left(g_{\mathrm{w}}^*u\right)(\boldsymbol{\sigma})=u(g_{\mathrm{w}}\boldsymbol{\sigma})$ is real.
Further, $\overline{\left(g_{\mathrm{w}}^*u\right)(\boldsymbol{\sigma})}=\overline{u}(g_{\mathrm{w}}\boldsymbol{\sigma})=\left(g_{\mathrm{w}}^*\overline{u}\right)(\boldsymbol{\sigma})$.
The map $g_{\mathrm{w}*}$ is the map $\left(g'_{\mathrm{w}}\right)^*$ for the element $g'_{\mathrm{w}}=g_{\mathrm{w}}^{-1}$ of $\mathrm{Gc}$.

(ii) One has $\left(\left(g_{\mathrm{w}}g'_{\mathrm{w}}\right)^*u\right)(\boldsymbol{\sigma})=u\left(g_{\mathrm{w}}(g'_{\mathrm{w}}\boldsymbol{\sigma})\right)=\left(g_{\mathrm{w}}^*u\right)(g'_{\mathrm{w}}\boldsymbol{\sigma})=\left(\left(g'_{\mathrm{w}}\right)^*g_{\mathrm{w}}^*u\right)(\boldsymbol{\sigma})$.
Consequently
\[
\left(g_{\mathrm{w}}g'_{\mathrm{w}}\right)_*=\left(\left(g'_{\mathrm{w}}\right)^{-1}g_{\mathrm{w}}^{-1}\right)^*=\left(g_{\mathrm{w}}^{-1}\right)^*\left(\left(g'_{\mathrm{w}}\right)^{-1}\right)^*=g_{\mathrm{w}*}\,g'_{\mathrm{w}*},
\]
and $g_{\mathrm{w}*}\,g_{\mathrm{w}}^*=\left(g_{\mathrm{w}}^{-1}\right)^*g_{\mathrm{w}}^*=\left(g_{\mathrm{w}}g_{\mathrm{w}}^{-1}\right)^*=1$, $g_{\mathrm{w}}^*\,g_{\mathrm{w}*}=\left(g_{\mathrm{w}}^{-1}g_{\mathrm{w}}\right)^*=1$.

(iii) By (0.2),
\begin{align*}
\left(g_{\mathrm{w}}^*d\boldsymbol{\lambda}\right)(c,c')&=\boldsymbol{\lambda}(g_{\mathrm{w}}c')-\boldsymbol{\lambda}(g_{\mathrm{w}}c)=\left(d\,g_{\mathrm{w}}^*\boldsymbol{\lambda}\right)(c,c'),\\
\left(g_{\mathrm{w}}^*dA\right)(c^1,\dots,c^4)&=A\left(g_{\mathrm{w}}c^1,g_{\mathrm{w}}c^2\right)+A\left(g_{\mathrm{w}}c^2,g_{\mathrm{w}}c^3\right)+A\left(g_{\mathrm{w}}c^3,g_{\mathrm{w}}c^4\right)+A\left(g_{\mathrm{w}}c^4,g_{\mathrm{w}}c^1\right)\\
&=\left(d\,g_{\mathrm{w}}^*A\right)(c^1,\dots,c^4).
\end{align*}
Since $g_{\mathrm{w}}^{-1}$ is an element of $\mathrm{Gc}$, the same two identities with $g_{\mathrm{w}}^{-1}$ in the place of $g_{\mathrm{w}}$ read $g_{\mathrm{w}*}d\boldsymbol{\lambda}=d\,g_{\mathrm{w}*}\boldsymbol{\lambda}$ and $g_{\mathrm{w}*}dA=d\,g_{\mathrm{w}*}A$.

(iv) The map $\boldsymbol{\sigma}\mapsto g_{\mathrm{w}}\boldsymbol{\sigma}$ takes the two orientations of a cell to the two orientations of its image.
It therefore induces a map of the unoriented $j$-cells of $\mathbb{Z}^3$;
by $g_{\mathrm{w}}(\boldsymbol{\sigma}+n\,\mathrm{jv})=g_{\mathrm{w}}\boldsymbol{\sigma}+n\,r\,\mathrm{jv}$ this map is compatible with congruence modulo $n\mathbb{Z}^3$ and so acts on the classes modulo $n\mathbb{Z}^3$, and the map induced by $g_{\mathrm{w}}^{-1}$ is its inverse.
For $j=1,2$ both factors of $u(\boldsymbol{\sigma})\cdot v(\boldsymbol{\sigma})$ are odd, so that this product does not depend on the orientation of $\boldsymbol{\sigma}$.
Hence $\left\langle g_{\mathrm{w}}^*u,g_{\mathrm{w}}^*v\right\rangle=\sum_{\boldsymbol{\sigma}}u(g_{\mathrm{w}}\boldsymbol{\sigma})\cdot v(g_{\mathrm{w}}\boldsymbol{\sigma})$ is the sum of $u\cdot v$ over all unoriented $j$-cells modulo $n\mathbb{Z}^3$, that is, $\langle u,v\rangle$.
With $g_{\mathrm{w}}^{-1}$ in the place of $g_{\mathrm{w}}$ this gives the second equality.
Finally, by (ii), $\left\langle g_{\mathrm{w}*}u,v\right\rangle=\left\langle g_{\mathrm{w}}^*g_{\mathrm{w}*}u,g_{\mathrm{w}}^*v\right\rangle=\left\langle u,g_{\mathrm{w}}^*v\right\rangle$.

(v) One has $g_{\mathrm{w}}^{-1}(c+e_\mu)=r(c+e_\mu-a)=g_{\mathrm{w}}^{-1}c+\epsilon_\mu e_\mu$.
Therefore $\left(g_{\mathrm{w}*}\mathrm{Hcn}\right)(c,c+e_\mu)=\mathrm{Hcn}\left(c',c'+\epsilon_\mu e_\mu\right)$ with $c'=g_{\mathrm{w}}^{-1}c$.
This number is $\mathrm{hv}_\mu$ if $\epsilon_\mu=1$, and it is $\mathrm{Hcn}(c',c'-e_\mu)=-\mathrm{Hcn}(c'-e_\mu,c')=-\mathrm{hv}_\mu$ if $\epsilon_\mu=-1$.
The cochain $H_{\theta}$ has $\mathrm{hv}_\mu=\theta_\mu/(qn)$.
The same computation shows that $g_{\mathrm{w}*}{V_{\mathrm{h}}}$ is contained in ${V_{\mathrm{h}}}$.
By (iii), $g_{\mathrm{w}*}d\boldsymbol{\lambda}=d\,g_{\mathrm{w}*}\boldsymbol{\lambda}$ and $g_{\mathrm{w}*}dA=d\,g_{\mathrm{w}*}A$;
hence $g_{\mathrm{w}*}{V_{\mathrm{g}}}\subset{V_{\mathrm{g}}}$, as ${V_{\mathrm{g}}}=dC^{0}(n)$, and $g_{\mathrm{w}*}F_N\subset F_N$, as $F_N=dC^{1}(n)$.
Since $g_{\mathrm{w}}^{-1}$ is an element of $\mathrm{Gc}$, also $g_{\mathrm{w}}^*=\left(g_{\mathrm{w}}^{-1}\right)_*$ maps ${V_{\mathrm{h}}}$ into ${V_{\mathrm{h}}}$, ${V_{\mathrm{g}}}$ into ${V_{\mathrm{g}}}$ and $F_N$ into $F_N$.
By (ii) the spaces ${V_{\mathrm{h}}}=g_{\mathrm{w}*}g_{\mathrm{w}}^*{V_{\mathrm{h}}}$, ${V_{\mathrm{g}}}=g_{\mathrm{w}*}g_{\mathrm{w}}^*{V_{\mathrm{g}}}$ and $F_N=g_{\mathrm{w}*}g_{\mathrm{w}}^*F_N$ are contained in $g_{\mathrm{w}*}{V_{\mathrm{h}}}$, $g_{\mathrm{w}*}{V_{\mathrm{g}}}$ and $g_{\mathrm{w}*}F_N$, respectively.
Thus $g_{\mathrm{w}*}{V_{\mathrm{h}}}={V_{\mathrm{h}}}$, $g_{\mathrm{w}*}{V_{\mathrm{g}}}={V_{\mathrm{g}}}$ and $g_{\mathrm{w}*}F_N=F_N$.
The map $g_{\mathrm{w}*}$ is orthogonal and maps ${V_{\mathrm{g}}}+{V_{\mathrm{h}}}$ onto itself;
so it maps the orthogonal complement ${V_{\mathrm{t}}}$ onto itself.

By \cite[Section~10.1.1]{QEDI},
\[
\ker\left(d|_{C^{1}(n)}\right)={V_{\mathrm{g}}}+{V_{\mathrm{h}}}.
\]
So $d$ is injective on ${V_{\mathrm{t}}}$, and $d{V_{\mathrm{t}}}=d\left({V_{\mathrm{g}}}+{V_{\mathrm{h}}}+{V_{\mathrm{t}}}\right)=dC^{1}(n)=F_N$.
Consequently, for $F^{\mathrm{u}}\in F_N$, the element $A(F^{\mathrm{u}})$ of Section~I.7.3 is the only element of ${V_{\mathrm{t}}}$ with $dA(F^{\mathrm{u}})=F^{\mathrm{u}}$, and $F^{\mathrm{u}}\mapsto A(F^{\mathrm{u}})$ is the inverse of the linear bijection $d:{V_{\mathrm{t}}}\to F_N$.
The cochain $g_{\mathrm{w}*}A(F^{\mathrm{u}})$ lies in ${V_{\mathrm{t}}}$, and $d\,g_{\mathrm{w}*}A(F^{\mathrm{u}})=g_{\mathrm{w}*}F^{\mathrm{u}}$;
hence it is $A\left(g_{\mathrm{w}*}F^{\mathrm{u}}\right)$.

The measure $\gamma_N$ is the standard Gaussian measure of the Euclidean space $\left(F_N,\|\cdot\|\right)$:
with respect to the Lebesgue measure that assigns the volume one to the cube spanned by an orthonormal basis, it has the density
\[
(2\pi)^{-\frac{\dim F_N}{2}}\cdot\exp\left(-\frac{\left\|F^{\mathrm{u}}\right\|^2}{2}\right).
\]
The density and this Lebesgue measure are invariant under every orthogonal map of $F_N$;
the restriction of $g_{\mathrm{w}*}$ to $F_N$ (by (iv)) and the map $-1$ are such maps.

(vi) One has $g_{\mathrm{w}}(c,c+e_\mu)=\left(g_{\mathrm{w}}c,\,g_{\mathrm{w}}c+\epsilon_\mu e_\mu\right)$;
for $\epsilon_\mu=-1$ this is the reversed bond $\left(c''+e_\mu,\,c''\right)$ with $c''=g_{\mathrm{w}}c-e_\mu$.
Next,
\[
g_{\mathrm{w}}(c;\mu\nu)=\left(g_{\mathrm{w}}c,\,g_{\mathrm{w}}c+\epsilon_\mu e_\mu,\,g_{\mathrm{w}}c+\epsilon_\mu e_\mu+\epsilon_\nu e_\nu,\,g_{\mathrm{w}}c+\epsilon_\nu e_\nu\right)
\]
is a quadruple on the unit square whose least corner is $c''=g_{\mathrm{w}}c-\mathbf{1}[\epsilon_\mu=-1]\,e_\mu-\mathbf{1}[\epsilon_\nu=-1]\,e_\nu$.
For $(\epsilon_\mu,\epsilon_\nu)=(1,1)$ it is $(c'';\mu\nu)$.
For $(-1,-1)$ it is $\left(c''+e_\mu+e_\nu,\,c''+e_\nu,\,c'',\,c''+e_\mu\right)$, a cyclic rotation of $(c'';\mu\nu)$.
For $(-1,1)$ it is $\left(c''+e_\mu,\,c'',\,c''+e_\nu,\,c''+e_\mu+e_\nu\right)$, and for $(1,-1)$ it is $\left(c''+e_\nu,\,c''+e_\mu+e_\nu,\,c''+e_\mu,\,c''\right)$;
both are cyclic rotations of the reversed quadruple $\left(c'',\,c''+e_\nu,\,c''+e_\mu+e_\nu,\,c''+e_\mu\right)$.
\end{proof}

\begin{named}{Lemma 0.3}[sums over the sites of the torus; pairings; determinants; the maps $\tilde{g}$]
\begin{itemize}
\item[(i)] Let $\mathrm{ks}$ be an $n\mathbb{Z}^3$-periodic complex function on $\mathbb{Z}^3\times O_{\mathrm{cub}}$, and let $\mathrm{bj}$ be a bijection of $\mathbb{Z}^3\times O_{\mathrm{cub}}$ for which there is an $r\in\{\pm1\}^3$ with $\mathrm{bj}(x+n\,\mathrm{jv})=\mathrm{bj}(x)+n\,r\,\mathrm{jv}$ for all $x$ and all $\mathrm{jv}\in\mathbb{Z}^3$.
Then
\[
\sum_{x\bmod n}\mathrm{ks}(\mathrm{bj}\,x)=\sum_{x\bmod n}\mathrm{ks}(x).
\]
This applies to $\mathrm{bj}=\left(x\mapsto x^{(i)}\right)$, $i=0,\dots,3$ (with $r=1$, by (I.s1)), and to $\mathrm{bj}=g_{\mathrm{w}}\in\mathrm{Gc}$ (Lemma~0.1\,(i), (ii)).
\item[(ii)] Let $\mathrm{nsp}$ be $4$ or $2$, and let $\mathrm{Sfn}$ be the space of the functions $f:\mathbb{Z}^3\times O_{\mathrm{cub}}\to\mathbb{C}^{\mathrm{nsp}}$ with $f(c+ne_\mu,g)=-f(c,g)$ for $\mu=1,2,3$;
for $\mathrm{nsp}=4$ it is $\mathfrak{S}^{\mathrm{ch}}(1;n)$.
For $f',f\in\mathrm{Sfn}$ put
\[
f'(x)\cdot f(x):=\sum_{\beta=1}^{\mathrm{nsp}}f'_\beta(x)\cdot f_\beta(x),\qquad\langle f',f\rangle:=\frac{1}{48}\sum_{x\bmod n}f'(x)\cdot f(x),\qquad(f',f):=\frac{1}{48}\sum_{x\bmod n}\overline{f'(x)}\cdot f(x);
\]
the summands are $n\mathbb{Z}^3$-periodic, since each of their two factors changes its sign.
For $\mathrm{nsp}=4$ these are $\langle f',f\rangle^{\mathrm{ch}}$ and $(f',f)^{\mathrm{ch}}$, by (I.1.3) and Section~I.1.6 at $\eta_0=1$.
For $x\in\mathrm{Yr}$ and $\beta\in\{1,\dots,\mathrm{nsp}\}$ let $\mathrm{ebs}_{x,\beta}$ be the element of $\mathrm{Sfn}$ whose value at a point $x'$ of $\mathrm{Yr}$ is $\mathbf{1}[x'=x]$ times the unit vector number $\beta$ of $\mathbb{C}^{\mathrm{nsp}}$.
These $48\cdot\mathrm{nsp}\cdot n^3$ functions ($192\,n^3$ for $\mathrm{nsp}=4$, $96\,n^3$ for $\mathrm{nsp}=2$) form a basis of $\mathrm{Sfn}$, and
\[
\overline{\mathrm{ebs}_{x,\beta}}=\mathrm{ebs}_{x,\beta},\qquad\left\langle\mathrm{ebs}_{x,\beta},\mathrm{ebs}_{x',\beta'}\right\rangle=\left(\mathrm{ebs}_{x,\beta},\mathrm{ebs}_{x',\beta'}\right)=\frac{1}{48}\cdot\mathbf{1}[x=x',\ \beta=\beta'].
\]
For a linear map $O$ of $\mathrm{Sfn}$ let $\mathrm{Mat}(O)$ be its matrix in this basis;
then $\det O=\det\mathrm{Mat}(O)$.
There is exactly one linear map $O^T$ and exactly one linear map $O^*$ of $\mathrm{Sfn}$ such that $\langle O^Tf',f\rangle=\langle f',Of\rangle$ and $(O^*f',f)=(f',Of)$ for all $f',f\in\mathrm{Sfn}$ (the transpose and the adjoint for the two pairings);
and $f\mapsto\overline{O\overline{f}}$ is a linear map of $\mathrm{Sfn}$, denoted $\overline{O}$.
Then
\[
\mathrm{Mat}\left(O^T\right)=\mathrm{Mat}(O)^T,\qquad\mathrm{Mat}\left(O^*\right)=\overline{\mathrm{Mat}(O)}^T,\qquad\mathrm{Mat}\left(\overline{O}\right)=\overline{\mathrm{Mat}(O)};
\]
hence $\det O^T=\det O$ and $\det O^*=\det\overline{O}=\overline{\det O}$.
\item[(iii)] For $S\in\mathrm{Sp}$ and $a\in\mathbb{Z}^3$ put $\tilde{g}:=(S,a)$, $g_{\mathrm{w}}:=(a,r(S))$ and $(\tilde{g}f)(x):=S\cdot f\left(g_{\mathrm{w}}^{-1}x\right)$ (Section~I.4.3; on sites $\tilde{g}$ acts as $g_{\mathrm{w}}$).
Then $\tilde{g}$ is a linear bijection of $\mathfrak{S}^{\mathrm{ch}}(1;n)$ onto itself, and its inverse is $f\mapsto S^{-1}\cdot f(g_{\mathrm{w}}\,\cdot\,)$.
\end{itemize}
\end{named}

\begin{proof}
(i) Every class of $\left(\mathbb{Z}^3\times O_{\mathrm{cub}}\right)/n\mathbb{Z}^3$ contains exactly one point of $\mathrm{Yr}$;
so there are $48\,n^3$ classes.
Let $\mathrm{bj}(x)=\mathrm{bj}(x')+n\,\mathrm{jv}$ with $x,x'\in\mathrm{Yr}$ and $\mathrm{jv}\in\mathbb{Z}^3$.
Because $r^2=1$, this means $\mathrm{bj}(x)=\mathrm{bj}\left(x'+n\,r\,\mathrm{jv}\right)$;
hence $x=x'+n\,r\,\mathrm{jv}$, and so $x=x'$.
The $48\,n^3$ points of $\mathrm{bj}(\mathrm{Yr})$ are therefore pairwise incongruent modulo $n\mathbb{Z}^3$, that is, $\mathrm{bj}(\mathrm{Yr})$ is a set of representatives of $\left(\mathbb{Z}^3\times O_{\mathrm{cub}}\right)/n\mathbb{Z}^3$;
the sum of the periodic function $\mathrm{ks}$ over this set equals its sum over $\mathrm{Yr}$.

(ii) A function on $\mathrm{Yr}$ with values in $\mathbb{C}^{\mathrm{nsp}}$ has exactly one extension in $\mathrm{Sfn}$, namely
\[
f(c+n\,\mathrm{jv},g):=(-1)^{\mathrm{jv}_1+\mathrm{jv}_2+\mathrm{jv}_3}\cdot f(c,g)\qquad\left((c,g)\in\mathrm{Yr},\ \mathrm{jv}\in\mathbb{Z}^3\right).
\]
So an element of $\mathrm{Sfn}$ is determined by its values on $\mathrm{Yr}$, and these values are arbitrary;
this gives the basis.
The values of $\mathrm{ebs}_{x,\beta}$ are real.
In $\left\langle\mathrm{ebs}_{x,\beta},\mathrm{ebs}_{x',\beta'}\right\rangle$ and in $\left(\mathrm{ebs}_{x,\beta},\mathrm{ebs}_{x',\beta'}\right)$ the sum over $\mathrm{Yr}$ equals $\mathbf{1}[x=x']\cdot\mathbf{1}[\beta=\beta']$.
Write $\mathrm{tb}=(x,\beta)$ and $\mathrm{tb}'=(x',\beta')$ for the indices of the basis, and
\[
O\,\mathrm{ebs}_{\mathrm{tb}'}=\sum_{\mathrm{tb}}\mathrm{Mat}(O)_{\mathrm{tb}\,\mathrm{tb}'}\cdot\mathrm{ebs}_{\mathrm{tb}}.
\]
Then
\[
\left\langle\mathrm{ebs}_{\mathrm{tb}},O\,\mathrm{ebs}_{\mathrm{tb}'}\right\rangle=\left(\mathrm{ebs}_{\mathrm{tb}},O\,\mathrm{ebs}_{\mathrm{tb}'}\right)=\frac{1}{48}\cdot\mathrm{Mat}(O)_{\mathrm{tb}\,\mathrm{tb}'},
\]
and, for a linear map $O_1$ of $\mathrm{Sfn}$,
\[
\left\langle O_1\mathrm{ebs}_{\mathrm{tb}},\mathrm{ebs}_{\mathrm{tb}'}\right\rangle=\frac{1}{48}\cdot\mathrm{Mat}(O_1)_{\mathrm{tb}'\,\mathrm{tb}},\qquad\left(O_1\mathrm{ebs}_{\mathrm{tb}},\mathrm{ebs}_{\mathrm{tb}'}\right)=\frac{1}{48}\cdot\overline{\mathrm{Mat}(O_1)_{\mathrm{tb}'\,\mathrm{tb}}}.
\]
The first pairing is bilinear, and the second is antilinear in its first and linear in its second argument;
so both are determined by their values on the basis.
Hence $\langle O_1f',f\rangle=\langle f',Of\rangle$ holds for all $f',f$ if and only if $\mathrm{Mat}(O_1)_{\mathrm{tb}'\,\mathrm{tb}}=\mathrm{Mat}(O)_{\mathrm{tb}\,\mathrm{tb}'}$ for all $\mathrm{tb},\mathrm{tb}'$,
and $(O_1f',f)=(f',Of)$ holds for all $f',f$ if and only if $\overline{\mathrm{Mat}(O_1)_{\mathrm{tb}'\,\mathrm{tb}}}=\mathrm{Mat}(O)_{\mathrm{tb}\,\mathrm{tb}'}$ for all $\mathrm{tb},\mathrm{tb}'$.
This proves the existence and the uniqueness of $O^T$ and $O^*$, and gives their matrices;
the weight $\frac{1}{48}$ is the same on both sides and drops out.
The map $f\mapsto\overline{f}$ maps $\mathrm{Sfn}$ onto itself, because the sign $-1$ is real;
moreover $\overline{O\,\overline{zf}}=z\cdot\overline{O\overline{f}}$ for every complex number $z$.
So $\overline{O}$ is a linear map of $\mathrm{Sfn}$, and
\[
\overline{O}\,\mathrm{ebs}_{\mathrm{tb}'}=\overline{O\,\mathrm{ebs}_{\mathrm{tb}'}}=\sum_{\mathrm{tb}}\overline{\mathrm{Mat}(O)_{\mathrm{tb}\,\mathrm{tb}'}}\cdot\mathrm{ebs}_{\mathrm{tb}}.
\]
Finally, a square matrix and its transpose have the same determinant, and the determinant of the complex conjugate of a square matrix is the complex conjugate of its determinant.

(iii) Here $g_{\mathrm{w}}^{-1}=(-ra,r)$ with $r=r(S)$.
By Lemma~0.1\,(ii) for $g_{\mathrm{w}}^{-1}$,
\[
(\tilde{g}f)(x+ne_\mu)=S\cdot f\left(g_{\mathrm{w}}^{-1}x+n\epsilon_\mu e_\mu\right)=-S\cdot f\left(g_{\mathrm{w}}^{-1}x\right)=-(\tilde{g}f)(x),
\]
because $f(x'+ne_\mu)=-f(x')$ and, with $x'-ne_\mu$ in the place of $x'$, $f(x'-ne_\mu)=-f(x')$.
Thus the map that belongs to a pair $(S,a)$, $S\in\mathrm{Sp}$, $a\in\mathbb{Z}^3$, maps $\mathfrak{S}^{\mathrm{ch}}(1;n)$ into itself.
The matrix $S^{-1}$ lies in $\mathrm{Sp}$.
Multiplying $\gamma_\nu\cdot S=r_\nu(S)\cdot S\cdot\gamma_\nu$ by $S^{-1}$ from the left and from the right gives
\[
S^{-1}\cdot\gamma_\nu=r_\nu(S)\cdot\gamma_\nu\cdot S^{-1},\quad\text{that is,}\quad\gamma_\nu\cdot S^{-1}=r_\nu(S)\cdot S^{-1}\cdot\gamma_\nu\qquad\left(r_\nu(S)=\pm1\right);
\]
so $r(S^{-1})=r$.
Moreover $(-ra,r)^{-1}=(a,r)=g_{\mathrm{w}}$.
Hence $f\mapsto S^{-1}\cdot f(g_{\mathrm{w}}\,\cdot\,)$ is the map that belongs to the pair $(S^{-1},-ra)$, and by what has been shown it maps $\mathfrak{S}^{\mathrm{ch}}(1;n)$ into itself.
By Lemma~0.1\,(i) the two maps are inverse to each other:
\[
S^{-1}\cdot(\tilde{g}f)(g_{\mathrm{w}}x)=S^{-1}\cdot S\cdot f\left(g_{\mathrm{w}}^{-1}g_{\mathrm{w}}x\right)=f(x),\qquad S\cdot\left[S^{-1}\cdot f(g_{\mathrm{w}}\,\cdot\,)\right]\left(g_{\mathrm{w}}^{-1}x\right)=f\left(g_{\mathrm{w}}g_{\mathrm{w}}^{-1}x\right)=f(x).
\]
\end{proof}

\begin{named}{Lemma 0.4}[three identities; complex $A$]
Let $A$ be a complex one-cochain on $\mathbb{Z}^3/n\mathbb{Z}^3$.
On $\mathfrak{S}^{\mathrm{ch}}(1;n)$ the following identities hold, constant $4\times4$ matrices acting at each site.
\begin{itemize}
\item[(i)] $\tilde{g}\,D_A=D_{g_{\mathrm{w}*}A}\,\tilde{g}$ for $\tilde{g}=(S,a)$, $g_{\mathrm{w}}=(a,r(S))$, $S\in\mathrm{Sp}$, $a\in\mathbb{Z}^3$;
here $g_{\mathrm{w}*}A$ is a complex one-cochain on $\mathbb{Z}^3/n\mathbb{Z}^3$ by Lemma~0.2\,(i).
\item[(ii)] $D_A^T=\mathsf{C}^{-1}\cdot D_{-A}\cdot\mathsf{C}$.
\item[(iii)] $\overline{D_{\overline{A}}}=\varkappa^{-1}\cdot D_{-A}\cdot\varkappa$.
\end{itemize}
\end{named}

\begin{proof}
\emph{The matrices.}
By the definition of $r(S)$ in Section~I.4.1, $\gamma_\nu\cdot S=r_\nu(S)\cdot S\cdot\gamma_\nu$ with $r_\nu(S)=\pm1$;
hence $S\cdot\gamma_\nu\cdot S^{-1}=r_\nu(S)\cdot\gamma_\nu$, that is, $S\cdot(\gamma\cdot v)\cdot S^{-1}=\gamma\cdot\left(r(S)v\right)$ (I.4.1).
By (I.c3), $\mathsf{C}^{-1}=-\mathsf{C}$, so that $\mathsf{C}^2=-1$, and $\gamma_\mu^T=-\mathsf{C}^{-1}\cdot\gamma_\mu\cdot\mathsf{C}$;
consequently
\[
\mathsf{C}^{-1}\cdot\gamma_\mu^T\cdot\mathsf{C}=-\mathsf{C}^{-2}\cdot\gamma_\mu\cdot\mathsf{C}^2=-\gamma_\mu.
\]
By (I.c4), $\overline{\gamma_\mu}=\varkappa^{-1}\cdot\gamma_\mu\cdot\varkappa$.
The numbers $\rho$, $w_i$ and the vectors $\hat{d}_i(x)$ are real.
With $r=r(S)$, Lemma~0.1\,(iii) and (I.s2) therefore give
\begin{align}
S\cdot M_i(x)\cdot S^{-1}&=\frac{w_i}{2}\left(\rho-\gamma\cdot\left(r\hat{d}_i(x)\right)\right)=\frac{w_i}{2}\left(\rho-\gamma\cdot\hat{d}_i(g_{\mathrm{w}}x)\right)=M_i(g_{\mathrm{w}}x);\tag{I.m4}\\
\mathsf{C}^{-1}\cdot M_i(x)^T\cdot\mathsf{C}&=\frac{w_i}{2}\left(\rho+\gamma\cdot\hat{d}_i(x)\right)=\frac{w_i}{2}\left(\rho-\gamma\cdot\hat{d}_i\left(x^{(i)}\right)\right)=M_i\left(x^{(i)}\right).\tag{I.m2}
\end{align}
Read at $x^{(i)}$, with $\left(x^{(i)}\right)^{(i)}=x$, the identity (I.m2) becomes
\[
M_i\left(x^{(i)}\right)^T=\mathsf{C}\cdot M_i(x)\cdot\mathsf{C}^{-1}=\left(-\mathsf{C}^{-1}\right)\cdot M_i(x)\cdot\left(-\mathsf{C}\right)=\mathsf{C}^{-1}\cdot M_i(x)\cdot\mathsf{C}.
\]
Finally,
\begin{equation}\tag{I.m3}
\overline{M_i(x)}=\frac{w_i}{2}\left(\rho-\overline{\gamma}\cdot\hat{d}_i(x)\right)=\varkappa^{-1}\cdot M_i(x)\cdot\varkappa.
\end{equation}

(i) Let $x':=g_{\mathrm{w}}^{-1}x$.
By Lemma~0.1\,(iii) for $g_{\mathrm{w}}^{-1}$, $\left(x'\right)^{(i)}=g_{\mathrm{w}}^{-1}\left(x^{(i)}\right)$.
By (0.4),
\[
(\tilde{g}D_Af)(x)=S\cdot(D_Af)(x')=\frac{\rho}{2}\cdot\mathrm{sw}\cdot S\cdot f(x')-\sum_{i=0}^{3}S\cdot M_i(x')\cdot U_A\left(x',g_{\mathrm{w}}^{-1}x^{(i)}\right)\cdot f\left(g_{\mathrm{w}}^{-1}x^{(i)}\right).
\]
Here $S\cdot M_i(x')=M_i(g_{\mathrm{w}}x')\cdot S=M_i(x)\cdot S$ by (I.m4);
further $S\cdot f(x')=(\tilde{g}f)(x)$ and $S\cdot f\left(g_{\mathrm{w}}^{-1}x^{(i)}\right)=(\tilde{g}f)\left(x^{(i)}\right)$;
and, by Lemma~0.1\,(i) and the definition of $g_{\mathrm{w}*}$,
\[
U_A\left(g_{\mathrm{w}}^{-1}x,g_{\mathrm{w}}^{-1}x^{(i)}\right)=\exp\left(\mathrm{i}\,q\cdot A\left(g_{\mathrm{w}}^{-1}\pi x,g_{\mathrm{w}}^{-1}\pi x^{(i)}\right)\right)=\exp\left(\mathrm{i}\,q\cdot\left(g_{\mathrm{w}*}A\right)\left(\pi x,\pi x^{(i)}\right)\right)=U_{g_{\mathrm{w}*}A}\left(x,x^{(i)}\right),
\]
both sides being $1$ for $i\le2$.
It follows that
\[
(\tilde{g}D_Af)(x)=\frac{\rho}{2}\cdot\mathrm{sw}\cdot(\tilde{g}f)(x)-\sum_{i=0}^{3}M_i(x)\cdot U_{g_{\mathrm{w}*}A}\left(x,x^{(i)}\right)\cdot(\tilde{g}f)\left(x^{(i)}\right)=\left(D_{g_{\mathrm{w}*}A}\tilde{g}f\right)(x).
\]

(ii) Let $f',f\in\mathfrak{S}^{\mathrm{ch}}(1;n)$.
By (0.4),
\[
\langle f',D_Af\rangle^{\mathrm{ch}}=\frac{\rho}{2}\cdot\mathrm{sw}\cdot\langle f',f\rangle^{\mathrm{ch}}-\sum_{i=0}^{3}\frac{1}{48}\sum_{x\bmod n}\mathrm{ks}_i(x),\qquad\mathrm{ks}_i(x):=f'(x)\cdot M_i(x)\cdot U_A\left(x,x^{(i)}\right)\cdot f\left(x^{(i)}\right).
\]
The function $\mathrm{ks}_i$ is $n\mathbb{Z}^3$-periodic:
$f'$ and $f\left((\,\cdot\,)^{(i)}\right)$ change their signs under $x\mapsto x+ne_\mu$ (by (I.s1)), $M_i$ does not depend on $c$, and (0.3) holds.
Lemma~0.3\,(i) with $\mathrm{bj}=\left(x\mapsto x^{(i)}\right)$, together with (0.3) and (I.s1), gives
\begin{align*}
\sum_{x\bmod n}\mathrm{ks}_i(x)&=\sum_{x\bmod n}\mathrm{ks}_i\left(x^{(i)}\right)=\sum_{x\bmod n}f'\left(x^{(i)}\right)\cdot M_i\left(x^{(i)}\right)\cdot U_A\left(x^{(i)},x\right)\cdot f(x)\\
&=\sum_{x\bmod n}\left[M_i\left(x^{(i)}\right)^T\cdot U_{-A}\left(x,x^{(i)}\right)\cdot f'\left(x^{(i)}\right)\right]\cdot f(x).
\end{align*}
Therefore $\langle f',D_Af\rangle^{\mathrm{ch}}=\langle O'f',f\rangle^{\mathrm{ch}}$ for all $f$, where
\[
(O'f')(x):=\frac{\rho}{2}\cdot\mathrm{sw}\cdot f'(x)-\sum_{i=0}^{3}M_i\left(x^{(i)}\right)^T\cdot U_{-A}\left(x,x^{(i)}\right)\cdot f'\left(x^{(i)}\right).
\]
By the last expression for $M_i\left(x^{(i)}\right)^T$ obtained above from (I.m2), and by (0.4),
\[
(O'f')(x)=\frac{\rho}{2}\cdot\mathrm{sw}\cdot f'(x)-\sum_{i=0}^{3}\mathsf{C}^{-1}\cdot M_i(x)\cdot\mathsf{C}\cdot U_{-A}\left(x,x^{(i)}\right)\cdot f'\left(x^{(i)}\right)=\left(\mathsf{C}^{-1}D_{-A}\mathsf{C}f'\right)(x).
\]
So $O'=\mathsf{C}^{-1}\cdot D_{-A}\cdot\mathsf{C}$ maps $\mathfrak{S}^{\mathrm{ch}}(1;n)$ into itself, and it is the transpose of $D_A$ (Lemma~0.3\,(ii) with $\mathrm{nsp}=4$).

(iii) The number $q$ is real;
hence
\[
\overline{U_{\overline{A}}\left(x,x^{(i)}\right)}=\overline{\exp\left(\mathrm{i}\,q\cdot\overline{A}\left(\pi x,\pi x^{(i)}\right)\right)}=\exp\left(-\mathrm{i}\,q\cdot A\left(\pi x,\pi x^{(i)}\right)\right)=U_{-A}\left(x,x^{(i)}\right).
\]
By (0.4) and (I.m3),
\begin{align*}
\overline{\left(D_{\overline{A}}\,\overline{f}\right)(x)}&=\frac{\rho}{2}\cdot\mathrm{sw}\cdot f(x)-\sum_{i=0}^{3}\overline{M_i(x)}\cdot U_{-A}\left(x,x^{(i)}\right)\cdot f\left(x^{(i)}\right)\\
&=\frac{\rho}{2}\cdot\mathrm{sw}\cdot f(x)-\sum_{i=0}^{3}\varkappa^{-1}\cdot M_i(x)\cdot\varkappa\cdot U_{-A}\left(x,x^{(i)}\right)\cdot f\left(x^{(i)}\right)=\left(\varkappa^{-1}D_{-A}\varkappa f\right)(x).
\end{align*}
\end{proof}

\begin{named}{Lemma 0.5}[real $A$]
Let $A\in C^{1}(n)$, and let $z$ be a real number.
\begin{itemize}
\item[(i)] $\det\left(D_A+z\right)\ge0$.
\item[(ii)] $\operatorname{Re}\left(f,(D_A+z)f\right)^{\mathrm{ch}}\ge z\cdot(f,f)^{\mathrm{ch}}$ for all $f\in\mathfrak{S}^{\mathrm{ch}}(1;n)$.
If $z>0$, then $\det\left(D_A+z\right)>0$.
\end{itemize}
\end{named}

\begin{proof}
(i) By Section~I.2.1, $\gamma_\mu=\operatorname{diag}\left(\sigma_\mu,-\sigma_\mu\right)$.
Hence $M_i(x)=\operatorname{diag}\left(M^{\boldsymbol{\sigma}}_i(x),M^{-\boldsymbol{\sigma}}_i(x)\right)$ with the $2\times2$ matrices
\[
M^{\boldsymbol{\sigma}}_i(x):=\frac{w_i}{2}\left(\rho-\sigma\cdot\hat{d}_i(x)\right),\qquad M^{-\boldsymbol{\sigma}}_i(x):=\frac{w_i}{2}\left(\rho+\sigma\cdot\hat{d}_i(x)\right).
\]
Let $\mathrm{Sfn}$ be the space of Lemma~0.3\,(ii) with $\mathrm{nsp}=2$, and $(\,\cdot\,,\,\cdot\,)$ its sesquilinear pairing;
the elements of $\mathrm{Sfn}$ are written $\mathrm{us}$.
The map that sends $f$ to the pair formed by its first two and by its last two components is a linear bijection of $\mathfrak{S}^{\mathrm{ch}}(1;n)$ onto $\mathrm{Sfn}\oplus\mathrm{Sfn}$;
it corresponds to the decomposition $\mathbb{C}^4=\mathbb{C}^2\oplus\mathbb{C}^2$ into the eigenspaces of the matrix $\operatorname{diag}(1,1,-1,-1)$.
By (0.4) this map takes $D_A$ to $\operatorname{diag}\left(D^{\boldsymbol{\sigma}}_A,D^{-\boldsymbol{\sigma}}_A\right)$, where, for $\mathrm{us}\in\mathrm{Sfn}$,
\[
\left(D^{\pm\boldsymbol{\sigma}}_A\mathrm{us}\right)(x):=\frac{\rho}{2}\cdot\mathrm{sw}\cdot\mathrm{us}(x)-\sum_{i=0}^{3}M^{\pm\boldsymbol{\sigma}}_i(x)\cdot U_A\left(x,x^{(i)}\right)\cdot\mathrm{us}\left(x^{(i)}\right).
\]
(The operator $D^{\boldsymbol{\sigma}}_A$ is the block $D_A^{\boldsymbol{\sigma}}$ of (I.O3).)
The operator $D^{\pm\boldsymbol{\sigma}}_A$ maps $\mathrm{Sfn}$ into itself:
$(x+ne_\mu)^{(i)}=x^{(i)}+ne_\mu$ by (I.s1), $M^{\pm\boldsymbol{\sigma}}_i(x)$ depends on $g$ only, and (0.3) holds.
The matrices $\sigma_\mu$ are Hermitian, and $\rho$, $w_i$, $\hat{d}_i(x)$ are real;
so, by (I.s2),
\[
M^{\boldsymbol{\sigma}}_i\left(x^{(i)}\right)^*=M^{\boldsymbol{\sigma}}_i\left(x^{(i)}\right)=\frac{w_i}{2}\left(\rho+\sigma\cdot\hat{d}_i(x)\right)=M^{-\boldsymbol{\sigma}}_i(x).
\]
As $A$ is real, (0.3) gives $\overline{U_A\left(x^{(i)},x\right)}=U_{-A}\left(x^{(i)},x\right)=U_A\left(x,x^{(i)}\right)$.
The function $x\mapsto\overline{\mathrm{us}'(x)}\cdot M^{\boldsymbol{\sigma}}_i(x)\cdot U_A\left(x,x^{(i)}\right)\cdot\mathrm{us}\left(x^{(i)}\right)$ is $n\mathbb{Z}^3$-periodic, since two of its factors change their signs ((0.3), (I.s1)).
Lemma~0.3\,(i) with $\mathrm{bj}=\left(x\mapsto x^{(i)}\right)$ and (I.s1) give
\begin{align*}
\sum_{x\bmod n}\overline{\mathrm{us}'(x)}\cdot M^{\boldsymbol{\sigma}}_i(x)\cdot U_A\left(x,x^{(i)}\right)\cdot\mathrm{us}\left(x^{(i)}\right)&=\sum_{x\bmod n}\overline{\mathrm{us}'\left(x^{(i)}\right)}\cdot M^{\boldsymbol{\sigma}}_i\left(x^{(i)}\right)\cdot U_A\left(x^{(i)},x\right)\cdot\mathrm{us}(x)\\
&=\sum_{x\bmod n}\overline{\left[M^{-\boldsymbol{\sigma}}_i(x)\cdot U_A\left(x,x^{(i)}\right)\cdot\mathrm{us}'\left(x^{(i)}\right)\right]}\cdot\mathrm{us}(x).
\end{align*}
The number $\frac{\rho}{2}\cdot\mathrm{sw}$ is real.
Hence $\left(\mathrm{us}',D^{\boldsymbol{\sigma}}_A\mathrm{us}\right)=\left(D^{-\boldsymbol{\sigma}}_A\mathrm{us}',\mathrm{us}\right)$ for all $\mathrm{us}',\mathrm{us}\in\mathrm{Sfn}$, that is, $\left(D^{\boldsymbol{\sigma}}_A\right)^*=D^{-\boldsymbol{\sigma}}_A$;
and $\left(D^{\boldsymbol{\sigma}}_A+z\right)^*=D^{-\boldsymbol{\sigma}}_A+z$, because $z$ is real.
The determinant of $\operatorname{diag}(O_1,O_2)$ on $\mathrm{Sfn}\oplus\mathrm{Sfn}$ is $\det O_1\cdot\det O_2$.
Therefore Lemma~0.3\,(ii) with $\mathrm{nsp}=2$ yields
\[
\det\left(D_A+z\right)=\det\left(D^{\boldsymbol{\sigma}}_A+z\right)\cdot\det\left(D^{-\boldsymbol{\sigma}}_A+z\right)=\det\left(D^{\boldsymbol{\sigma}}_A+z\right)\cdot\det\left(\left(D^{\boldsymbol{\sigma}}_A+z\right)^*\right)=\left|\det\left(D^{\boldsymbol{\sigma}}_A+z\right)\right|^2\ge0.
\]

(ii) Let $\left(R_i(A)f\right)(x):=U_A\left(x,x^{(i)}\right)\cdot f\left(x^{(i)}\right)$ (I.2.2) and $\left(\mathrm{Gd}_if\right)(x):=\gamma\cdot\hat{d}_i(x)\,f(x)$.
By (0.3) and (I.s1),
\[
\left(R_i(A)f\right)(x+ne_\mu)=U_A\left(x+ne_\mu,x^{(i)}+ne_\mu\right)\cdot f\left(x^{(i)}+ne_\mu\right)=-\left(R_i(A)f\right)(x),
\]
and $\hat{d}_i(x)$ depends on $g$ only.
So $R_i(A)$ and $\mathrm{Gd}_i$ map $\mathfrak{S}^{\mathrm{ch}}(1;n)$ into itself, and all summands below are products of two antiperiodic functions, hence $n\mathbb{Z}^3$-periodic.
Inserting $M_i(x)=\frac{w_i}{2}\left(\rho-\gamma\cdot\hat{d}_i(x)\right)$ into (0.4) one obtains
\[
D_A=\frac{\rho}{2}\sum_{i=0}^{3}w_i\cdot\left(1-R_i(A)\right)+\frac{1}{2}\sum_{i=0}^{3}w_i\cdot\mathrm{Gd}_i\,R_i(A),
\]
which is (I.2.4).
For real $A$, Lemma~0.3\,(i) with $\mathrm{bj}=\left(x\mapsto x^{(i)}\right)$, (I.s1) and $\overline{U_A\left(x^{(i)},x\right)}=U_A\left(x,x^{(i)}\right)$ give
\begin{align*}
\left(f',R_i(A)f\right)^{\mathrm{ch}}&=\frac{1}{48}\sum_{x\bmod n}\overline{f'(x)}\cdot U_A\left(x,x^{(i)}\right)\cdot f\left(x^{(i)}\right)=\frac{1}{48}\sum_{x\bmod n}\overline{f'\left(x^{(i)}\right)}\cdot U_A\left(x^{(i)},x\right)\cdot f(x)\\
&=\frac{1}{48}\sum_{x\bmod n}\overline{\left[U_A\left(x,x^{(i)}\right)\cdot f'\left(x^{(i)}\right)\right]}\cdot f(x)=\left(R_i(A)f',f\right)^{\mathrm{ch}},
\end{align*}
and
\[
\left(R_i(A)^2f\right)(x)=U_A\left(x,x^{(i)}\right)\cdot U_A\left(x^{(i)},x\right)\cdot f(x)=f(x).
\]
Thus $R_i(A)$ is self-adjoint with $R_i(A)^2=1$;
it preserves the norm, the number $\left(f,R_i(A)f\right)^{\mathrm{ch}}$ is real, and $\left(f,R_i(A)f\right)^{\mathrm{ch}}\le(f,f)^{\mathrm{ch}}$ by the Cauchy--Schwarz inequality.
The operator $\mathrm{Gd}_i$ is self-adjoint ($\gamma_\mu$ is Hermitian by (I.c1), and $\hat{d}_i(x)$ is real), and $R_i(A)\,\mathrm{Gd}_i=-\mathrm{Gd}_i\,R_i(A)$ by (I.s2).
Hence $\left(\mathrm{Gd}_i\,R_i(A)\right)^*=R_i(A)\,\mathrm{Gd}_i=-\mathrm{Gd}_i\,R_i(A)$, and $\operatorname{Re}\left(f,\mathrm{Gd}_i\,R_i(A)f\right)^{\mathrm{ch}}=0$.
Since $\rho>0$ and $w_i>0$,
\[
\operatorname{Re}\left(f,(D_A+z)f\right)^{\mathrm{ch}}=z\cdot(f,f)^{\mathrm{ch}}+\frac{\rho}{2}\sum_{i=0}^{3}w_i\cdot\left[(f,f)^{\mathrm{ch}}-\left(f,R_i(A)f\right)^{\mathrm{ch}}\right]\ge z\cdot(f,f)^{\mathrm{ch}}.
\]
If $z>0$ and $(D_A+z)f=0$, then $0\ge z\cdot(f,f)^{\mathrm{ch}}$, so that $f=0$.
Thus $D_A+z$ is injective, $\det\left(D_A+z\right)\ne0$, and (i) gives the sign.
\end{proof}

\subsection*{1. Proof of (a)}

\begin{proof}[Proof of Proposition~49\,(a)]
\textbf{1.1. The reflections.}
The reflection formula of Section~I.4.8, in the case $\ell = y_\mu$, $n_\ell = e_\mu$, $|n_\ell|^{2} = 1$, shows that the reflection in the plane $\{y_\mu = \mathrm{po}\}$ is the map $y \mapsto y - 2\left(y_\mu - \mathrm{po}\right)e_\mu$.
Let $\mathrm{po} = j\,\mathrm{hc}$.
Since $2\,\mathrm{hc} = \mathrm{per}$, the reflection $\mathrm{thw}_{\mu,j}$ in the wall $\{y_\mu = j\,\mathrm{hc}\}$ acts by
\[
\mathrm{thw}_{\mu,j}\colon\ y_\mu \mapsto 2j\,\mathrm{hc} - y_\mu = j\cdot\mathrm{per} - y_\mu ,
\]
and it leaves the two other coordinates fixed.
Thus $\mathrm{thw}_{\mu,j} = \left(j\cdot\mathrm{per}\,e_\mu, \mathrm{rf}_\mu\right)$, where $\mathrm{rf}_\mu = 1 - 2e_\mu e_\mu^{T} \in \{\pm1\}^3$.
As $\mathrm{per}$ is an integer, the vector $j\cdot\mathrm{per}\,e_\mu$ lies in $\mathbb{Z}^3$ for every integer $j$, even or odd; so $\mathrm{thw}_{\mu,j}$ belongs to $\mathrm{Gc}$, and in fact to the set $\mathrm{Gp}$ of the pairs $(a,r)$ with $a \in \mathrm{per}\,\mathbb{Z}^3$ and $r \in \{\pm1\}^3$.
(For even $j$ the number $j\,\mathrm{hc} = (j/2)\cdot\mathrm{per}$ is an integer, and the wall is a site plane.
For odd $j$ the integer $j\cdot\mathrm{per}$ is odd, so that $j\,\mathrm{hc} = j\cdot\mathrm{per}/2$ lies in $\mathbb{Z} + \frac{1}{2}$, and the wall is a link plane.)

\textbf{1.2. $\mathrm{Gp}$ is a group, and $\mathfrak{c} \mapsto (a_{\mathfrak{c}}, r_{\mathfrak{c}})$ is a bijection of $\mathbb{Z}^3$ onto $\mathrm{Gp}$.}
For $(a,r)$ and $(a',r')$ in $\mathrm{Gp}$, the product $(a,r)(a',r') = (a + ra', rr')$ and the inverse $(a,r)^{-1} = (-ra, r)$ again lie in $\mathrm{Gp}$, because $r\,\mathrm{per}\,\mathbb{Z}^3 = \mathrm{per}\,\mathbb{Z}^3$.
Let $r_{\mathfrak{c}}$ and $a_{\mathfrak{c}}$ be given by the formulas displayed in (a).
Every entry $(a_{\mathfrak{c}})_\mu$ is $\mathrm{per}$ times the integer $\mathfrak{c}_\mu/2$, respectively $(\mathfrak{c}_\mu + 1)/2$; hence $a_{\mathfrak{c}} \in \mathrm{per}\,\mathbb{Z}^3$ and $(a_{\mathfrak{c}}, r_{\mathfrak{c}}) \in \mathrm{Gp}$.
Conversely, let $(a,r) \in \mathrm{Gp}$, and write $a = \mathrm{per}\,\mathrm{jv}$ with $\mathrm{jv} \in \mathbb{Z}^3$ and $re_\mu = \epsilon_\mu e_\mu$.
Then $(a,r) = (a_{\mathfrak{c}}, r_{\mathfrak{c}})$ holds if and only if, for each $\mu$, the entry $\mathfrak{c}_\mu$ is even with $\mathfrak{c}_\mu/2 = \mathrm{jv}_\mu$ in the case $\epsilon_\mu = 1$, and odd with $(\mathfrak{c}_\mu + 1)/2 = \mathrm{jv}_\mu$ in the case $\epsilon_\mu = -1$; that is, if and only if
\begin{equation}
\mathfrak{c}_\mu = 2\,\mathrm{jv}_\mu \quad (\epsilon_\mu = 1), \qquad \mathfrak{c}_\mu = 2\,\mathrm{jv}_\mu - 1 \quad (\epsilon_\mu = -1).
\tag{a.1}
\end{equation}
Hence there is exactly one such $\mathfrak{c}$.

\textbf{1.3. $W(\mathcal{Q}_{*}) = \mathrm{Gp}$.}
The group $W(\mathcal{Q}_{*})$ is generated by the reflections $\mathrm{thw}_{\mu,j}$ (Section~I.6.4), and these belong to the group $\mathrm{Gp}$ by 1.1 and 1.2; therefore $W(\mathcal{Q}_{*}) \subseteq \mathrm{Gp}$.
For the opposite inclusion observe that $(a_{\mathfrak{c}}, r_{\mathfrak{c}})$ acts on each coordinate separately: on $y_\mu$ it acts by $y_\mu \mapsto i\cdot\mathrm{per} + y_\mu$ if $\mathfrak{c}_\mu = 2i$, and by $y_\mu \mapsto (i+1)\cdot\mathrm{per} - y_\mu$ if $\mathfrak{c}_\mu = 2i + 1$.
The first of these maps is $\mathrm{thw}_{\mu,i}\cdot\mathrm{thw}_{\mu,0}$, which sends $y_\mu$ first to $-y_\mu$ and then to $i\cdot\mathrm{per} + y_\mu$; the second is $\mathrm{thw}_{\mu,i+1}$; both fix the two other coordinates.
Consequently $(a_{\mathfrak{c}}, r_{\mathfrak{c}})$ is the product, over $\mu = 1,2,3$, of these elements of $W(\mathcal{Q}_{*})$.
Together with 1.2 this gives
\[
W(\mathcal{Q}_{*}) = \mathrm{Gp} = \left\{(a_{\mathfrak{c}}, r_{\mathfrak{c}}) : \mathfrak{c} \in \mathbb{Z}^3\right\}.
\]

\textbf{1.4. Cells; simple transitivity; $\mathsf{g}_{\mathfrak{c}}$.}
One has $Q^{\circ}_{0} = (0,\mathrm{hc})^3$.
If $\mathfrak{c}_\mu = 2i$, the map $y_\mu \mapsto i\cdot\mathrm{per} + y_\mu$ carries $(0,\mathrm{hc})$ onto $\left(2i\,\mathrm{hc}, (2i+1)\,\mathrm{hc}\right) = \left(\mathfrak{c}_\mu\mathrm{hc}, (\mathfrak{c}_\mu + 1)\,\mathrm{hc}\right)$.
If $\mathfrak{c}_\mu = 2i+1$, the map $y_\mu \mapsto (i+1)\cdot\mathrm{per} - y_\mu$ carries $(0,\mathrm{hc})$ onto $\left((2i+2)\,\mathrm{hc} - \mathrm{hc}, (2i+2)\,\mathrm{hc}\right) = \left(\mathfrak{c}_\mu\mathrm{hc}, (\mathfrak{c}_\mu + 1)\,\mathrm{hc}\right)$.
Therefore
\begin{equation}
(a_{\mathfrak{c}}, r_{\mathfrak{c}})\,Q^{\circ}_{0} = Q^{\circ}_{\mathfrak{c}} \qquad \text{for every } \mathfrak{c} \in \mathbb{Z}^3.
\tag{a.2}
\end{equation}
The cells $Q^{\circ}_{\mathfrak{c}}$, $\mathfrak{c} \in \mathbb{Z}^3$, are non-empty and pairwise disjoint, hence pairwise distinct.
By 1.2, 1.3 and (a.2) the map $g_{\mathrm{w}} \mapsto g_{\mathrm{w}}Q^{\circ}_{0}$ is thus a bijection of $W(\mathcal{Q}_{*})$ onto the set of the cells; in other words, $(a_{\mathfrak{c}}, r_{\mathfrak{c}})$ is the only element of $W(\mathcal{Q}_{*})$ that carries $Q^{\circ}_{0}$ onto $Q^{\circ}_{\mathfrak{c}}$.
It is the element $\mathsf{g}_{\mathfrak{c}}$ of Section~I.6.4:
\[
\mathsf{g}_{\mathfrak{c}}\,y = r_{\mathfrak{c}}\,y + a_{\mathfrak{c}}.
\]

Let $g_{\mathrm{w}} \in W(\mathcal{Q}_{*})$ and $\mathfrak{c} \in \mathbb{Z}^3$.
Since $g_{\mathrm{w}}\mathsf{g}_{\mathfrak{c}}$ lies in $W(\mathcal{Q}_{*})$, one has $g_{\mathrm{w}}\mathsf{g}_{\mathfrak{c}} = \mathsf{g}_{\mathfrak{c}'}$ for exactly one $\mathfrak{c}'$, and then $g_{\mathrm{w}}Q^{\circ}_{\mathfrak{c}} = g_{\mathrm{w}}\mathsf{g}_{\mathfrak{c}}Q^{\circ}_{0} = Q^{\circ}_{\mathfrak{c}'}$.
So $g_{\mathrm{w}}$ maps cells onto cells, the index $g_{\mathrm{w}}.\mathfrak{c}$ equals $\mathfrak{c}'$, and
\begin{equation}
g_{\mathrm{w}}\mathsf{g}_{\mathfrak{c}} = \mathsf{g}_{g_{\mathrm{w}}.\mathfrak{c}}.
\tag{a.3}
\end{equation}
The rule $\mathfrak{c} \mapsto g_{\mathrm{w}}.\mathfrak{c}$ is an action of $W(\mathcal{Q}_{*})$ on $\mathbb{Z}^3$, since $g_{\mathrm{w}}g_{\mathrm{w}}'Q^{\circ}_{\mathfrak{c}} = g_{\mathrm{w}}Q^{\circ}_{g_{\mathrm{w}}'.\mathfrak{c}} = Q^{\circ}_{g_{\mathrm{w}}.(g_{\mathrm{w}}'.\mathfrak{c})}$.
This action is simply transitive.
Indeed, $\mathsf{g}_{\mathfrak{c}'}\mathsf{g}_{\mathfrak{c}}^{-1}$ carries $Q^{\circ}_{\mathfrak{c}}$ onto $Q^{\circ}_{\mathfrak{c}'}$; and if $g_{\mathrm{w}}.\mathfrak{c} = \mathfrak{c}$, then $g_{\mathrm{w}}\mathsf{g}_{\mathfrak{c}} = \mathsf{g}_{\mathfrak{c}}$ by (a.3), whence $g_{\mathrm{w}} = 1$.

\emph{Formula.}
Let $g_{\mathrm{w}} = (\mathrm{per}\,\mathrm{jv}, r)$ with $re_\mu = \epsilon_\mu e_\mu$.
The map $y_\mu \mapsto \epsilon_\mu y_\mu + \mathrm{jv}_\mu\cdot\mathrm{per}$ carries $\left(\mathfrak{c}_\mu\mathrm{hc}, (\mathfrak{c}_\mu + 1)\,\mathrm{hc}\right)$ onto $\left((\mathfrak{c}_\mu + 2\,\mathrm{jv}_\mu)\,\mathrm{hc}, (\mathfrak{c}_\mu + 2\,\mathrm{jv}_\mu + 1)\,\mathrm{hc}\right)$ if $\epsilon_\mu = 1$, and onto $\left((2\,\mathrm{jv}_\mu - \mathfrak{c}_\mu - 1)\,\mathrm{hc}, (2\,\mathrm{jv}_\mu - \mathfrak{c}_\mu)\,\mathrm{hc}\right)$ if $\epsilon_\mu = -1$.
Hence
\begin{equation}
(g_{\mathrm{w}}.\mathfrak{c})_\mu = \mathfrak{c}_\mu + 2\,\mathrm{jv}_\mu \quad (\epsilon_\mu = 1), \qquad (g_{\mathrm{w}}.\mathfrak{c})_\mu = 2\,\mathrm{jv}_\mu - 1 - \mathfrak{c}_\mu \quad (\epsilon_\mu = -1).
\tag{a.4}
\end{equation}
For $g_{\mathrm{w}} = \mathrm{thw}_{\mu,j}$ one has $\mathrm{jv} = je_\mu$, $\epsilon_\mu = -1$, and $\epsilon_\nu = 1$, $\mathrm{jv}_\nu = 0$ for $\nu \ne \mu$.
So $(g_{\mathrm{w}}.\mathfrak{c})_\mu = 2j - 1 - \mathfrak{c}_\mu$, while the other entries of $\mathfrak{c}$ are unchanged.

\textbf{1.5. The torus $\mathrm{kt}$.}
By (0.1) the translations $(n\,\mathrm{jv}, 1)$, $\mathrm{jv} \in \mathbb{Z}^3$, belong to $\mathrm{Gp} = W(\mathcal{Q}_{*})$.
They form a normal subgroup, denoted $n\mathbb{Z}^3$, because
\[
(a,r)(n\,\mathrm{jv},1)(a,r)^{-1} = (n\,r\,\mathrm{jv}, 1).
\]
The classes of $\mathrm{Gp}$ modulo $n\mathbb{Z}^3$ are given by the pairs formed by $a$ modulo $n\mathbb{Z}^3$ and by $r$; the group $\mathrm{per}\,\mathbb{Z}^3/n\mathbb{Z}^3$ has $(n/\mathrm{per})^3 = (\mathrm{kt}\cdot L)^3$ elements, and $\{\pm1\}^3$ has $8$.
Hence $W(\mathcal{Q}_{*})/n\mathbb{Z}^3$ has $8\,(\mathrm{kt}\cdot L)^3 = (2\,\mathrm{kt}\cdot L)^3$ elements.

Since $g_{\mathrm{w}}(c + n\,\mathrm{jv}) = g_{\mathrm{w}}c + n\,r\,\mathrm{jv}$, every $g_{\mathrm{w}}$ induces a map of $\mathbb{Z}^3/n\mathbb{Z}^3$, and this map is the identity for $g_{\mathrm{w}} \in n\mathbb{Z}^3$; thus $W(\mathcal{Q}_{*})/n\mathbb{Z}^3$ acts on $\mathbb{Z}^3/n\mathbb{Z}^3$.

By (a.4), $(n\,\mathrm{jv},1).\mathfrak{c} = \mathfrak{c} + 2\,(n/\mathrm{per})\,\mathrm{jv} = \mathfrak{c} + 2\,\mathrm{kt}\cdot L\,\mathrm{jv}$; that is, $Q^{\circ}_{\mathfrak{c}} + n\,\mathrm{jv} = Q^{\circ}_{\mathfrak{c} + 2\mathrm{kt}\cdot L\,\mathrm{jv}}$.
As the cells are pairwise disjoint, the images of $Q^{\circ}_{\mathfrak{c}}$ and of $Q^{\circ}_{\mathfrak{c}'}$ in $\mathbb{R}^3/n\mathbb{Z}^3$ coincide if $\mathfrak{c}' - \mathfrak{c} \in 2\,\mathrm{kt}\cdot L\,\mathbb{Z}^3$, and are disjoint otherwise.
So the cells of the torus $\mathrm{kt}$ are indexed by $\left(\mathbb{Z}/2\,\mathrm{kt}\cdot L\,\mathbb{Z}\right)^3$, and there are $(2\,\mathrm{kt}\cdot L)^3$ of them.
From $g_{\mathrm{w}}(n\,\mathrm{jv},1) = (n\,r\,\mathrm{jv},1)\,g_{\mathrm{w}}$ one obtains
\[
g_{\mathrm{w}}.(\mathfrak{c} + 2\,\mathrm{kt}\cdot L\,\mathrm{jv}) = g_{\mathrm{w}}.\mathfrak{c} + 2\,\mathrm{kt}\cdot L\,r\,\mathrm{jv}.
\]
Hence $\mathfrak{c} \mapsto g_{\mathrm{w}}.\mathfrak{c}$ descends to $\left(\mathbb{Z}/2\,\mathrm{kt}\cdot L\,\mathbb{Z}\right)^3$, on which the elements of $n\mathbb{Z}^3$ act as the identity: $W(\mathcal{Q}_{*})/n\mathbb{Z}^3$ acts on the cells of the torus $\mathrm{kt}$.
This action is transitive by 1.4.
It is simply transitive: if $g_{\mathrm{w}}.\mathfrak{c} = \mathfrak{c} + 2\,\mathrm{kt}\cdot L\,\mathrm{jv} = (n\,\mathrm{jv},1).\mathfrak{c}$, then $(n\,\mathrm{jv},1)^{-1}g_{\mathrm{w}}$ fixes $\mathfrak{c}$, hence equals $1$ by 1.4, and $g_{\mathrm{w}} \in n\mathbb{Z}^3$.
Formula (a.3) with $g_{\mathrm{w}} = (n\,\mathrm{jv},1)$ reads
\begin{equation}
\mathsf{g}_{\mathfrak{c} + 2\mathrm{kt}\cdot L\,\mathrm{jv}} = (n\,\mathrm{jv},1)\,\mathsf{g}_{\mathfrak{c}}.
\tag{a.5}
\end{equation}
Therefore the class of $\mathsf{g}_{\mathfrak{c}}$ modulo $n\mathbb{Z}^3$ depends only on the class of $\mathfrak{c}$ modulo $2\,\mathrm{kt}\cdot L\,\mathbb{Z}^3$, and (a.3) holds modulo $n\mathbb{Z}^3$ for the cells of the torus.

\textbf{1.6. Fermion sites; spinors.}
The group $W(\mathcal{Q}_{*})$ is a subgroup of $\mathrm{Gc}$.
For $g_{\mathrm{w}} = (a,r)$, Lemma~0.1\,(i)--(iii) shows that $g_{\mathrm{w}}(c,g) = (rc + a, rg)$ defines a bijection of $\mathbb{Z}^3 \times O_{\mathrm{cub}}$ with $g_{\mathrm{w}}(x + ne_\mu) = g_{\mathrm{w}}x + n\,re_\mu$ and $(g_{\mathrm{w}}x)^{(i)} = g_{\mathrm{w}}(x^{(i)})$.
A reflection $\mathrm{thw}_{\mu,j}$ has the linear part $r = \mathrm{rf}_\mu \ne 1$, and by Lemma~0.1\,(iv) it fixes no fermion site; explicitly, as in Section~I.4.8,
\[
\mathrm{thw}_{\mu,j}(c,g) = \left(c - (2c_\mu - j\cdot\mathrm{per})\,e_\mu,\ \mathrm{rf}_\mu g\right).
\]
The last assertion of (a) is Lemma~0.3\,(iii).
\end{proof}

\subsection*{2. Proof of (b) ($\mathrm{kt} = 1$, $n = n_*$)}

\begin{proof}[Proof of Proposition~49\,(b)]
\textbf{2.1. The data are flow data.}
By the paragraph \emph{Scales} of 0.1, $R^{*}(e) = h L^{T^{*}(e)}$ with the integer $T^{*}(e) \ge 1$; moreover $R^{*}(e) = R_{\mathrm{top}}(e)/L \le R_{\mathrm{top}}(e)$ (Section~I.6.2).
By the definition of $\mathrm{Tor}(e)$ in Section~I.6.2, $R^{*}(e)$ is therefore an element of $\mathrm{Tor}(e)$.
Further, $N \ge N_0(e) = N_{\min}(R^{*}(e),e)$ (Section~I.6.2), the number $N_{\min}(R^{*}(e),e)$ being well defined by the assertion (Sc0) of Section~G4.0.2.
Hence the quadruple $(\bar{m}, e, R^{*}(e), N)$ satisfies the four conditions of (FL) in Section~G4.0.1, which constitute the data of Section~G4.0.2:
\[
\bar{m} \in (0,1], \qquad e\,h^{\frac{1}{2}} \le e_1, \qquad R^{*}(e) \in \mathrm{Tor}(e), \qquad N \ge N_{\min}(R^{*}(e),e);
\]
and the assertions (Sc1), (Sc4) and (Sc6) of Section~G4.0.2 hold with $R = R^{*}(e)$.

In detail, (Sc1) and (Sc4) are \cite[Proposition~19.1(c), (f)]{QEDI}, read with $h' = h$ and $m_0 = 0$ (Section~G4.0.2).
The hypothesis $R \in \mathrm{Tor}(e)$ of that proposition has just been established.
Its hypotheses $v_K \ge 1$ and $\left(1 + \mathrm{lg}/v_K\right)^r \le L$ hold by (Sc0).
The additional hypothesis of its clause (f) is the condition (2.1) of Section~G4.0.3, namely $r_{*} > 31L + 16$; it is one of the assertions of (Sc6), that is, of Proposition~14.2\,(c), (f) together with the row (e1) at $v = v_K$, and this row comes from the definition of $e_1$.
Indeed, by Proposition~14.2\,(a), line (C9), the number $e_1$ is one half of the supremum of the $e' \in \left(0, \min(e_0, 1/3)\right]$ for which the rows (e1)--(e8) of Proposition~14.2\,(b) hold for all $v \ge -\log e'$.
The data give $e_1 \ge e\,h^{\frac{1}{2}} > 0$; so this supremum, which is $2e_1$, exceeds $e_1$, and some such $e'$ is larger than $e_1$.
Consequently the rows (e1)--(e8) hold for all $v \ge -\log e_1$ (this is the first assertion of Proposition~14.2\,(c)), and $v_K = -\log\left(e\,h^{\frac{1}{2}}\right) \ge -\log e_1$.
The only hypothesis of (Sc0) is $e\,h^{\frac{1}{2}} \le e_1$, and the inequalities of the first display of (Sc0),
\[
e \le \frac{1}{3}, \qquad v_K \ge 8, \qquad \left(1 + \frac{\mathrm{lg}}{v_K}\right)^{r} \le \left(1 + \frac{3\,\mathrm{lg}}{v_K}\right)^{r} < L,
\]
are contained in (Sc6); in addition to this hypothesis, (Sc6) requires $R \in \mathrm{Tor}(e)$ and $N \ge N_{\min}(R,e)$.

\textbf{2.2. Sides.}
By Section~I.6.1, $a_j = L^{j-N}$ and $\eta = L^{-N}$.
Hence, with $k_{*} = k_{*}(R^{*}(e))$,
\begin{equation}
\ell_j := \ell_j(R^{*}(e)) = \frac{R^{*}(e)}{a_j} = L^{-j}\,\frac{R^{*}(e)}{\eta} = L^{-j}n_* \qquad (0 \le j \le k_{*}).
\tag{b.1}
\end{equation}
By (Sc4), $\ell_j$ belongs to $\mathrm{Cl}$: the number $\ell_j/(LM)$ is an odd integer $\ge 13$.
Hence $\ell_j\mathbb{Z}^3 \subseteq LM\,\mathbb{Z}^3$, and the groups $\mathbb{G}_j$, $\mathbb{G}_j^{L}$, $\mathbb{G}^{L}_{j,3}$ of Section~I.4.2 with the side $\ell_j$ are defined.

\textbf{2.3. Arithmetic.}
By (Sc1), $k_{*} \ge 1$ and $\ell_{k_{*}} = \ell_* = M r_{*}$.
Taking $j = k_{*}$ in (b.1), one finds
\[
n_* = L^{k_{*}}\cdot M\cdot r_{*}, \qquad \mathrm{per} = \frac{n_*}{L} = L^{k_{*}-1}\cdot M\cdot r_{*}.
\]
The number $r_{*} = r_{k_{*}}$ is a power of $L$ (Section~I.6.1: $r_k$ is the least $L^i \ge v_k^r$).
By (Sc6) the condition (2.1) of Section~G4.0.3 holds: $r_{*} > 31L + 16 > L$.
A power of $L$ that exceeds $L$ is at least $L^2$; therefore $r_{*}/L^2$ is a positive integer, and for $0 \le j \le k_{*}$
\begin{equation}
L^{-j}\cdot\mathrm{per} = L^{k_{*}-1-j}\cdot M\cdot r_{*} = LM\cdot L^{k_{*}-j}\cdot\frac{r_{*}}{L^2}, \qquad \text{with the integral exponent } k_{*} - j \ge 0.
\tag{b.2}
\end{equation}
So $L^{-j}\cdot\mathrm{per} \in LM\,\mathbb{Z}$; and as $a_{\mathfrak{c}} \in \mathrm{per}\,\mathbb{Z}^3$ by (a), the vector $L^{-j}a_{\mathfrak{c}}$ lies in $LM\,\mathbb{Z}^3$.

\textbf{2.4. Coordinates of level $j$.}
One unit of length of level $j$ equals $a_j/\eta = L^j$ lattice units (Sections~I.3.8, I.6.1).
The photon site with the coordinates $c^0 \in \mathbb{Z}^3$ in lattice units has the coordinates $c = L^{-j}c^0 \in \mathbb{L}_j = L^{-j}\mathbb{Z}^3$ at level $j$ (Section~I.6.1; the definition $a^{[j]}(b) = L^{\frac{j}{2}}a(L^jb)$ of Section~I.6.6 rests on this identification).
By (b.1) the map $c^0 \mapsto L^{-j}c^0$ carries $n_*\mathbb{Z}^3$ onto $\ell_j\mathbb{Z}^3$, and hence $\mathbb{Z}^3/n_*\mathbb{Z}^3$ onto $T_j = \mathbb{L}_j/\ell_j\mathbb{Z}^3$.
In the coordinates of level $j$ the map $\mathsf{g}_{\mathfrak{c}}$ becomes
\[
c \mapsto L^{-j}\mathsf{g}_{\mathfrak{c}}(L^jc) = L^{-j}\left(r_{\mathfrak{c}}L^jc + a_{\mathfrak{c}}\right) = r_{\mathfrak{c}}c + L^{-j}a_{\mathfrak{c}}.
\]

\textbf{2.5. The matrices $S$ with $r(S) = r$.}
The map $r\colon \mathrm{Sp} \to \{\pm1\}^3$ is a homomorphism, since
\[
\gamma_\nu SS' = r_\nu(S)\,S\gamma_\nu S' = r_\nu(S)\,r_\nu(S')\,SS'\gamma_\nu.
\]
By (I.c1), for $\nu \ne \mu$,
\[
\gamma_\nu(\gamma_\mu\gamma_5) = -\gamma_\mu\gamma_\nu\gamma_5 = (\gamma_\mu\gamma_5)\gamma_\nu,
\]
while
\[
\gamma_\mu(\gamma_\mu\gamma_5) = -\gamma_\mu\gamma_5\gamma_\mu = -(\gamma_\mu\gamma_5)\gamma_\mu.
\]
Hence $r(\gamma_\mu\gamma_5) = \mathrm{rf}_\mu$.
Moreover $(\gamma_\mu\gamma_5)\gamma_5 = \gamma_\mu = -\gamma_5(\gamma_\mu\gamma_5)$.
By (I.c2) the matrix $i\Sigma_3 = \gamma_1\gamma_2\gamma_3$ belongs to $\mathrm{Sp}$, commutes with the three matrices $\gamma_\nu$ (so that $r(i\Sigma_3) = 1$), and anticommutes with $\gamma_5$.

For $r \in \{\pm1\}^3$ with $re_\mu = \epsilon_\mu e_\mu$, let $\mathrm{nr}$ be the number of the indices $\mu$ with $\epsilon_\mu = -1$, and let $S_r$ be the product of the matrices $\gamma_\mu\gamma_5$ over these $\mu$, taken in increasing order of $\mu$ ($S_r := 1$ for $r = 1$).
Then $r(S_r) = \prod_{\mu\colon\epsilon_\mu = -1}\mathrm{rf}_\mu = r$.
Hence $r$ maps $\mathrm{Sp}$ onto $\{\pm1\}^3$, a group with $8$ elements; as $\mathrm{Sp}$ has $32$ elements (Section~I.4.1), the kernel of $r$ has $4$ elements, and $1$, $-1$, $i\Sigma_3$, $-i\Sigma_3$ are four distinct elements of it.
Therefore $\ker r = \{1, -1, i\Sigma_3, -i\Sigma_3\}$ (Section~I.4.1), and the matrices $S \in \mathrm{Sp}$ with $r(S) = r$ are the four matrices $S_r$, $-S_r$, $i\Sigma_3S_r$, $-i\Sigma_3S_r$; they satisfy
\[
S_r\gamma_5 = (-1)^{\mathrm{nr}}\gamma_5S_r, \qquad \left(i\Sigma_3S_r\right)\gamma_5 = (-1)^{\mathrm{nr}+1}\gamma_5\left(i\Sigma_3S_r\right).
\]
Consequently $\pm S_r$ lie in $\mathrm{Sp}_3$ if $\mathrm{nr}$ is even, and $\pm i\Sigma_3S_r$ lie in $\mathrm{Sp}_3$ if $\mathrm{nr}$ is odd; so for every $r \in \{\pm1\}^3$ there is an $S \in \mathrm{Sp}_3$ with $r(S) = r$ (Section~I.4.1: $r(\mathrm{Sp}_3) = \{\pm1\}^3$).
For $r = r_{\mathfrak{c}}$ the number $\mathrm{nr}$ is the number of the odd entries of $\mathfrak{c}$.

\textbf{2.6. The elements $\tilde{g}^{(j)}_{\mathfrak{c}}$.}
Let $S \in \mathrm{Sp}$ with $r(S) = r_{\mathfrak{c}}$.
By 2.3 the vector $L^{-j}a_{\mathfrak{c}}$ lies in $LM\,\mathbb{Z}^3$; therefore $\tilde{g}^{(j)}_{\mathfrak{c}} = (S, L^{-j}a_{\mathfrak{c}})$, the second entry taken modulo $\ell_j\mathbb{Z}^3$, is an element of $\mathbb{G}_j^{L}$ (Section~I.4.2).
In Section~I.4.3 an element $(S,a)$ acts on photon sites by $c \mapsto r(S)c + a$; so $\tilde{g}^{(j)}_{\mathfrak{c}}$ acts on the photon sites of $T_j$ by $c \mapsto r(S)c + L^{-j}a_{\mathfrak{c}} = r_{\mathfrak{c}}c + L^{-j}a_{\mathfrak{c}}$, and by 2.4 this is $\mathsf{g}_{\mathfrak{c}}$ in the coordinates of level $j$: the element $\tilde{g}^{(j)}_{\mathfrak{c}}$ induces $\mathsf{g}_{\mathfrak{c}}$.
On the fermion sites of $T^{\mathrm{f}}_j = T_j \times O_{\mathrm{cub}}$ it acts by (I.4.2):
\[
\tilde{g}^{(j)}_{\mathfrak{c}}(c,g) = \left(r_{\mathfrak{c}}c + L^{-j}a_{\mathfrak{c}},\ r_{\mathfrak{c}}g\right).
\]
Writing $x/L^j := (c^0/L^j, g)$ for $x = (c^0,g) \in \mathbb{Z}^3 \times O_{\mathrm{cub}}$ (Section~I.3.7, applied $j$ times), one may express this as $\tilde{g}^{(j)}_{\mathfrak{c}}(x/L^j) = (\mathsf{g}_{\mathfrak{c}}x)/L^j$, where $\mathsf{g}_{\mathfrak{c}}x$ is the action of (a).
If $S$ is chosen in $\mathrm{Sp}_3$ (2.5), then $\tilde{g}^{(j)}_{\mathfrak{c}}$ belongs to $\mathbb{G}^{L}_{j,3}$, by the definition of this group.
If the same $S$ is used at all levels, then for $j < k_{*}$ the element $\tilde{g}^{(j+1)}_{\mathfrak{c}} = (S, L^{-j}a_{\mathfrak{c}}/L)$ is the element $(\tilde{g}^{(j)}_{\mathfrak{c}})'$ of (I.a4).

In particular, for $j = 0$ one has $\mathbb{L}_0 = \mathbb{Z}^3$ and $\ell_0 = n_*$, and $\tilde{g}^{(0)}_{\mathfrak{c}} = (S, a_{\mathfrak{c}})$ is the pair $\tilde{g}$ of (a) for $g_{\mathrm{w}} = \mathsf{g}_{\mathfrak{c}}$.

\textbf{2.7. The identity for $a^{[j]}$.}
Let $a$ be a complex one-cochain on $\mathbb{Z}^3/n_*\mathbb{Z}^3$.
For a bond $(c,c')$ of $\mathbb{L}_j$ one has $a^{[j]}(c,c') = L^{\frac{j}{2}}a(L^jc, L^jc')$ (Section~I.6.6).
This function is $\ell_j\mathbb{Z}^3$-periodic, since $L^j\ell_j = n_*$ by (b.1); thus it is a bond function on $T_j$.
Put $\tilde{g} := \tilde{g}^{(j)}_{\mathfrak{c}}$.
On photon sites $\tilde{g}^{-1}c = r_{\mathfrak{c}}\left(c - L^{-j}a_{\mathfrak{c}}\right)$, and so $L^j\tilde{g}^{-1}c = r_{\mathfrak{c}}\left(L^jc - a_{\mathfrak{c}}\right) = \mathsf{g}_{\mathfrak{c}}^{-1}(L^jc)$.
The action on bond functions in Section~I.4.3, $(\tilde{g}A)(c,c') = A\left(\tilde{g}^{-1}c, \tilde{g}^{-1}c'\right)$, and the definition of $(\mathsf{g}_{\mathfrak{c}})_{*}$ give
\begin{align*}
\left(\tilde{g}\,a^{[j]}\right)(c,c') &= a^{[j]}\left(\tilde{g}^{-1}c, \tilde{g}^{-1}c'\right) = L^{\frac{j}{2}}a\left(\mathsf{g}_{\mathfrak{c}}^{-1}L^jc, \mathsf{g}_{\mathfrak{c}}^{-1}L^jc'\right)\\
&= L^{\frac{j}{2}}\left((\mathsf{g}_{\mathfrak{c}})_{*}a\right)(L^jc, L^jc') = \left((\mathsf{g}_{\mathfrak{c}})_{*}a\right)^{[j]}(c,c').
\end{align*}
\end{proof}

\subsection*{3. Proof of (c)}

\begin{proof}[Proof of Proposition~49\,(c)]
Let $A$ be a complex one-cochain on $\mathbb{Z}^3/n\mathbb{Z}^3$ and $z$ a complex number.

\emph{First identity.}
Let $g_{\mathrm{w}} = (a,r) \in \mathrm{Gc}$.
Choose $S \in \mathrm{Sp}$ with $r(S) = r$ (for instance $S = S_r$ of 2.5; that construction involves $4 \times 4$ matrices only and does not depend on $\mathrm{kt}$), and put $\tilde{g} := (S,a)$.
By Lemma~0.3\,(iii), $\tilde{g}$ is a linear bijection of $\mathfrak{S}^{\mathrm{ch}}(1;n)$; by Lemma~0.4\,(i), and since $\tilde{g}$ commutes with multiplication by the number $z$,
\[
\tilde{g}\left(D_A + z\right)\tilde{g}^{-1} = D_{g_{\mathrm{w}*}A} + z.
\]
Similar linear maps have the same determinant, so that $\det\left(D_{g_{\mathrm{w}*}A} + z\right) = \det\left(D_A + z\right)$.

\emph{Second identity.}
Lemma~0.4\,(ii) gives
\[
\left(D_A + z\right)^{T} = D_A^{T} + z = \mathsf{C}^{-1}\left(D_{-A} + z\right)\mathsf{C}.
\]
By Lemma~0.3\,(ii) with $\mathrm{nsp} = 4$, one has $\det\left(D_A + z\right) = \det\left(\left(D_A + z\right)^{T}\right) = \det\left(D_{-A} + z\right)$.

\emph{Third identity.}
For a linear map $O$ and a number $z'$ one has $\overline{O + z'} = \overline{O} + \overline{z'}$, as the definition of the conjugate map, $f \mapsto \overline{(O + z')\overline{f}}$, shows.
Lemma~0.4\,(iii) gives
\[
\overline{D_{\overline{A}} + \overline{z}} = \overline{D_{\overline{A}}} + z = \varkappa^{-1}\left(D_{-A} + z\right)\varkappa.
\]
By Lemma~0.3\,(ii) with $\mathrm{nsp} = 4$ and by the second identity,
\[
\overline{\det\left(D_{\overline{A}} + \overline{z}\right)} = \det\left(\overline{D_{\overline{A}} + \overline{z}}\right) = \det\left(D_{-A} + z\right) = \det\left(D_A + z\right).
\]
\end{proof}

\subsection*{4. Proof of (d)}

\begin{proof}[Proof of Proposition~49\,(d)]
Let $\mathrm{ah} \in C^{1}(n)$ and $\mathbf{m} \in (0,1]$, and set $\mathrm{Nm}(t) := \det\left(D_{(t/e)\cdot\mathrm{ah}} + \mathbf{m}\right)$ for $t \in \mathbb{C}$; then $N[\mathrm{ah}](t) = \mathrm{Nm}(t)/\mathrm{Nm}(0)$.

\textbf{4.1.}
By Lemma~0.5\,(ii), applied with $A = 0$ and $z = \mathbf{m} > 0$, one has $\mathrm{Nm}(0) = \det\left(D_0 + \mathbf{m}\right) > 0$.
Thus $N[\mathrm{ah}]$ is defined, and $N[\mathrm{ah}](0) = 1$.

\textbf{4.2 (entire).}
Let $\mathrm{ebs}_{\mathrm{tb}'}$ and $\mathrm{Mat}(\,\cdot\,)$ be as in Lemma~0.3\,(ii) with $\mathrm{nsp} = 4$.
Let $A$ be a complex one-cochain and $\mathrm{tb}' = (x',\beta')$ an index.
By (0.4) the value of $\left(D_A + \mathbf{m}\right)\mathrm{ebs}_{\mathrm{tb}'}$ at a point $x$ of $\mathrm{Yr} = \{0,\dots,n-1\}^3 \times O_{\mathrm{cub}}$ is
\[
\left(\frac{\rho}{2}\,\mathrm{sw} + \mathbf{m}\right)\mathrm{ebs}_{\mathrm{tb}'}(x) - \sum_{i=0}^{3}M_i(x)\,U_A\left(x,x^{(i)}\right)\mathrm{ebs}_{\mathrm{tb}'}\left(x^{(i)}\right),
\]
and each of the vectors $\mathrm{ebs}_{\mathrm{tb}'}(x)$, $\mathrm{ebs}_{\mathrm{tb}'}\left(x^{(i)}\right)$ is $0$ or $\pm$ a unit vector of $\mathbb{C}^4$, whatever $A$ is.
Hence every entry of $\mathrm{Mat}\left(D_A + \mathbf{m}\right)$ is the sum of a number independent of $A$ and of a linear combination of the four numbers $U_A\left(x,x^{(i)}\right)$ with coefficients independent of $A$; this is (I.O1).
For $A = (t/e)\cdot\mathrm{ah}$, with $c = \pi(x)$ and $c_3 = \pi\left(x^{(3)}\right)$,
\[
U_A\left(x,x^{(i)}\right) = 1 \quad (i \le 2), \qquad U_A\left(x,x^{(3)}\right) = \exp\left(iq\,\frac{t}{e}\,\mathrm{ah}(c,c_3)\right),
\]
and these are entire functions of $t$.
The determinant is a polynomial in the entries of the matrix.
Therefore $\mathrm{Nm}$ and $N[\mathrm{ah}]$ are entire functions.

\textbf{4.3 (even).}
One has $-(t/e)\cdot\mathrm{ah} = \left((-t)/e\right)\cdot\mathrm{ah}$.
The second identity of (c), with $A = (t/e)\cdot\mathrm{ah}$ and $z = \mathbf{m}$, gives $\mathrm{Nm}(-t) = \mathrm{Nm}(t)$, and hence $N[\mathrm{ah}](-t) = N[\mathrm{ah}](t)$.

\textbf{4.4 (real; $N[\mathrm{ah}](t) > 0$ on the real axis).}
The cochain $\mathrm{ah}$ and the numbers $e$, $\mathbf{m}$ are real; hence $\overline{(t/e)\cdot\mathrm{ah}} = (\overline{t}/e)\cdot\mathrm{ah}$ and $\overline{\mathbf{m}} = \mathbf{m}$.
The third identity of (c), with $A = (t/e)\cdot\mathrm{ah}$ and $z = \mathbf{m}$, gives
\[
\overline{\mathrm{Nm}(\overline{t})} = \overline{\det\left(D_{\overline{A}} + \overline{z}\right)} = \det\left(D_A + z\right) = \mathrm{Nm}(t).
\]
Since $\mathrm{Nm}(0)$ is real, it follows that $\overline{N[\mathrm{ah}](\overline{t})} = N[\mathrm{ah}](t)$; for real $t$ this says that $N[\mathrm{ah}](t)$ is real.
For real $t$ the cochain $(t/e)\cdot\mathrm{ah}$ is real and $\mathbf{m} > 0$; so $\mathrm{Nm}(t) > 0$ by Lemma~0.5\,(ii), and $N[\mathrm{ah}](t) > 0$ by 4.1.

\textbf{4.5 (covariance).}
Let $g_{\mathrm{w}} \in \mathrm{Gc}$.
The cochain $g_{\mathrm{w}*}\mathrm{ah}$ lies in $C^{1}(n)$ (Lemma~0.2\,(i)), and $(t/e)\cdot g_{\mathrm{w}*}\mathrm{ah} = g_{\mathrm{w}*}\left((t/e)\cdot\mathrm{ah}\right)$, since $g_{\mathrm{w}*}$ is linear.
By the first identity of (c), $\det\left(D_{(t/e)\cdot g_{\mathrm{w}*}\mathrm{ah}} + \mathbf{m}\right) = \mathrm{Nm}(t)$.
The denominator $\det\left(D_0 + \mathbf{m}\right)$ does not involve $\mathrm{ah}$.
Hence $N[g_{\mathrm{w}*}\mathrm{ah}] = N[\mathrm{ah}]$.

\textbf{4.6 (the probe).}
The probe $\beta$ is real; therefore $u$ (the periodised means of the components of $\tau_v\beta$), $a_\eta = d_\eta^{*}u$ and $\hat{\mathfrak{a}} = \eta^{\frac{1}{2}}a_\eta$ are real (Section~I.7.6), and $\hat{\mathfrak{a}} \in C^{1}(n)$.
The number $\mathbf{m} := \bar{m}\cdot\eta$ lies in $(0,1]$, because $\bar{m} \in (0,1]$ and $\eta = L^{-N} \le 1$.
By Section~I.7.6, $N_0 = N[\hat{\mathfrak{a}}]$ with this $\mathbf{m}$.
As $\mathsf{g}_{\mathfrak{c}}$ lies in $W(\mathcal{Q}_{*})$, a subgroup of $\mathrm{Gc}$, 4.5 yields
\[
N[\hat{\mathfrak{a}}_{\mathfrak{c}}] = N\left[(\mathsf{g}_{\mathfrak{c}})_{*}\hat{\mathfrak{a}}\right] = N[\hat{\mathfrak{a}}] = N_0.
\]
By 4.3 and 4.4, $N_0$ is even and satisfies $\overline{N_0(\overline{t})} = N_0(t)$.
Since $N_0$ is entire,
\[
N_0(t) = \sum_{m \ge 0}c_m\,\frac{t^m}{m!} \quad \text{for all } t, \qquad c_m = N_0^{(m)}(0).
\]
The relation $N_0(-t) = N_0(t)$ gives $(-1)^mc_m = c_m$, so that $c_m = 0$ for odd $m$, in particular for $m = 1$ and $m = 3$.
The relation $\overline{N_0(\overline{t})} = \sum_m\overline{c_m}\,t^m/m! = N_0(t)$ gives $\overline{c_m} = c_m$: every $c_m$ is real, in particular $c_2$ and $c_4$.

\textbf{4.7 (the phases).}
One has
\[
q\cdot\frac{t}{e}\cdot\hat{\mathfrak{a}}(c,c_3) = e\,\eta^{\frac{1}{2}}\cdot\frac{t}{e}\cdot\eta^{\frac{1}{2}}a_\eta(c,c_3) = t\cdot\eta\cdot a_\eta(c,c_3).
\]
Hence $U_{(t/e)\cdot\hat{\mathfrak{a}}}\left(x,x^{(3)}\right) = \exp\left(it\cdot\eta\cdot a_\eta(c,c_3)\right)$, where $a_\eta(c,c_3) = \pm a_\eta(b)$ for the positively oriented bond $b$ with the end points $c$ and $c_3$: the factor $e$ has cancelled.
By Section~I.7.6 the cochain $a_\eta$ is formed with the period $n = \mathrm{kt}\cdot R^{*}(e)/\eta$ and with $\tau_v\beta$, where the placement $v$ is tied to $\ell_c = R^{*}(e)/(2L)$; both of these data depend on $e$.
\end{proof}

\subsection*{5. Proof of (e)}

\begin{proof}[Proof of Proposition~49\,(e)]
\textbf{5.1. Notation; three rules.}
For $\mathfrak{c} \in \mathbb{Z}^3$ the cochain $\hat{\mathfrak{a}}_{\mathfrak{c}} = (\mathsf{g}_{\mathfrak{c}})_{*}\hat{\mathfrak{a}}$ is a real $n\mathbb{Z}^3$-periodic one-cochain (4.6, Lemma~0.2\,(i)).
By (a.5) and Lemma~0.2\,(ii),
\[
\hat{\mathfrak{a}}_{\mathfrak{c} + 2\mathrm{kt}\cdot L\,\mathrm{jv}} = (n\,\mathrm{jv},1)_{*}\hat{\mathfrak{a}}_{\mathfrak{c}} = \hat{\mathfrak{a}}_{\mathfrak{c}},
\]
where the last equality holds because $\hat{\mathfrak{a}}_{\mathfrak{c}}$ is $n\mathbb{Z}^3$-periodic.
Thus $\hat{\mathfrak{a}}_{\mathfrak{c}}$ depends only on the cell of the torus $\mathrm{kt}$.
By (a.3) and Lemma~0.2\,(ii), for $g_{\mathrm{w}} \in W(\mathcal{Q}_{*})$,
\begin{equation}
g_{\mathrm{w}*}\hat{\mathfrak{a}}_{\mathfrak{c}} = \left(g_{\mathrm{w}}\mathsf{g}_{\mathfrak{c}}\right)_{*}\hat{\mathfrak{a}} = \hat{\mathfrak{a}}_{g_{\mathrm{w}}.\mathfrak{c}}.
\tag{e.1}
\end{equation}
Let $\mathbb{C}^{\mathrm{cells}}$ denote the space of the complex vectors $t = (t_{\mathfrak{c}})$ indexed by the cells of the torus $\mathrm{kt}$, and $\mathbb{R}^{\mathrm{cells}}$ the subspace of the real ones.
For $t \in \mathbb{C}^{\mathrm{cells}}$, $\theta \in \mathbb{R}^3$ and $F^{\mathrm{u}} \in F_N$ put
\begin{align*}
a(t) &:= e^{-1}\sum_{\mathfrak{c}}t_{\mathfrak{c}}\,\hat{\mathfrak{a}}_{\mathfrak{c}}, \qquad B(F^{\mathrm{u}},\theta,t) := A(F^{\mathrm{u}}) + H_{\theta} + a(t),\\
\mathrm{Igd}(F^{\mathrm{u}},\theta,t) &:= \det\left(D_{B(F^{\mathrm{u}},\theta,t)} + \boldsymbol{\mu}\right).
\end{align*}
Since $A(F^{\mathrm{u}})$, $H_{\theta}$, $\hat{\mathfrak{a}}_{\mathfrak{c}}$ and $e$ are real, one has $\mathrm{Im}\,B(F^{\mathrm{u}},\theta,t) = \mathrm{Im}\,a(t) = a(\mathrm{Im}\,t)$.
For $g_{\mathrm{w}} \in W(\mathcal{Q}_{*})$ with the linear part $r$,
\begin{equation}
\begin{aligned}
&g_{\mathrm{w}*}B(F^{\mathrm{u}},\theta,t) = B\left(g_{\mathrm{w}*}F^{\mathrm{u}}, r\theta, g_{\mathrm{w}}t\right),\\
&-B(F^{\mathrm{u}},\theta,t) = B\left(-F^{\mathrm{u}}, -\theta, -t\right), \qquad \overline{B(F^{\mathrm{u}},\theta,t)} = B\left(F^{\mathrm{u}}, \theta, \overline{t}\right).
\end{aligned}
\tag{e.2}
\end{equation}

\emph{Proof of (e.2).}
By Lemma~0.2\,(v), $g_{\mathrm{w}*}A(F^{\mathrm{u}}) = A(g_{\mathrm{w}*}F^{\mathrm{u}})$ and $g_{\mathrm{w}*}H_{\theta} = H_{r\theta}$.
By 1.5 the map $\mathfrak{c} \mapsto g_{\mathrm{w}}.\mathfrak{c}$ is a bijection of the set of the cells of the torus $\mathrm{kt}$; so the map $t \mapsto g_{\mathrm{w}}t$, given by $(g_{\mathrm{w}}t)_{g_{\mathrm{w}}.\mathfrak{c}} := t_{\mathfrak{c}}$, is well defined, and (e.1) gives
\[
a(g_{\mathrm{w}}t) = e^{-1}\sum_{\mathfrak{c}}(g_{\mathrm{w}}t)_{g_{\mathrm{w}}.\mathfrak{c}}\,\hat{\mathfrak{a}}_{g_{\mathrm{w}}.\mathfrak{c}} = e^{-1}\sum_{\mathfrak{c}}t_{\mathfrak{c}}\,g_{\mathrm{w}*}\hat{\mathfrak{a}}_{\mathfrak{c}} = g_{\mathrm{w}*}a(t).
\]
The second rule follows from $A(-F^{\mathrm{u}}) = -A(F^{\mathrm{u}})$ (Lemma~0.2\,(v)), $H_{-\theta} = -H_{\theta}$ (Section~I.6.6) and $a(-t) = -a(t)$.
The third rule is the reality noted above.

From (c) with the real number $z = \boldsymbol{\mu}$ and from (e.2) one obtains
\begin{equation}
\begin{aligned}
&\mathrm{Igd}\left(g_{\mathrm{w}*}F^{\mathrm{u}}, r\theta, g_{\mathrm{w}}t\right) = \mathrm{Igd}(F^{\mathrm{u}},\theta,t), \qquad \mathrm{Igd}\left(-F^{\mathrm{u}}, -\theta, -t\right) = \mathrm{Igd}(F^{\mathrm{u}},\theta,t),\\
&\overline{\mathrm{Igd}\left(F^{\mathrm{u}}, \theta, \overline{t}\right)} = \mathrm{Igd}(F^{\mathrm{u}},\theta,t).
\end{aligned}
\tag{e.3}
\end{equation}
(For the third identity: $\overline{\det\left(D_{\overline{B}} + \overline{\boldsymbol{\mu}}\right)} = \det\left(D_B + \boldsymbol{\mu}\right)$ with $B = B(F^{\mathrm{u}},\theta,t)$.)

\textbf{5.2. Convergence; $\bar{Z}^{\mathrm{ext}}$ is entire.}
By Proposition~38.1\,(a), $B \mapsto \det\left(D_B + \boldsymbol{\mu}\right)$ is an entire function on $C^{1}(n) \otimes \mathbb{C}$.
The map $(F^{\mathrm{u}},\theta,t) \mapsto B(F^{\mathrm{u}},\theta,t)$ is real linear, because $A(\,\cdot\,)$ is linear (Lemma~0.2\,(v)).
Hence $\mathrm{Igd}$ is continuous in $(F^{\mathrm{u}},\theta,t)$.
By Proposition~38.1\,(d),
\[
\left|\mathrm{Igd}(F^{\mathrm{u}},\theta,t)\right| \le \left(|M_{\mathrm{ch}}| + \frac{1+\rho}{2}\left(\mathrm{sw} - 6 + 6\exp\left(q\,|\mathrm{Im}\,a(t)|_{\infty}\right)\right)\right)^{192n^3},
\]
a bound that does not involve $(F^{\mathrm{u}},\theta)$.
The measures $\gamma_N$ and $\mathrm{d}\theta/(2\pi)^3$ on $[-\pi,\pi]^3$ are probability measures.
Therefore the integral
\[
\mathrm{Zmk}^{\theta}(t) = \int\mathrm{d}\gamma_N(F^{\mathrm{u}})\,\mathrm{Igd}(F^{\mathrm{u}},\theta,t)
\]
converges absolutely, is bounded by the same number, and is continuous in $\theta$ (dominated convergence); and
\[
\bar{Z}^{\mathrm{ext}}(t) = \int_{[-\pi,\pi]^3}\frac{\mathrm{d}\theta}{(2\pi)^3}\,\mathrm{Zmk}^{\theta}(t)
\]
is defined.
This is the display of Section~I.7.6, where $\mathrm{kt} = 1$ and $\boldsymbol{\mu} = \boldsymbol{\mu}_N$.

The coboundary $d\colon {V_{\mathrm{t}}} \to F_N$ is a linear bijection (Lemma~0.2\,(v)).
Hence the image of $\sigma$ under $d$ is a constant multiple of the Lebesgue measure of $F_N$, and the image of $\exp\left(-\frac{\|dA\|^2}{2}\right)\mathrm{d}\sigma(A)$ is a constant multiple of $\gamma_N$, the constant being the total mass $\mathrm{zV}_n \in (0,\infty)$ (Proposition~31.1, Definitions).
Moreover $A(dA) = A$ for $A \in {V_{\mathrm{t}}}$.
Therefore, for every complex one-cochain $a$ on $\mathbb{Z}^3/n\mathbb{Z}^3$, with $\mathbf{Z}^{\mathrm{x}}_n$ as in Section~I.7.4,
\begin{align*}
\mathbf{Z}^{\mathrm{x}}_n(a) &= \int\frac{\mathrm{d}\theta}{(2\pi)^3}\int_{{V_{\mathrm{t}}}}\mathrm{d}\sigma(A)\,\exp\left(-\frac{\|dA\|^2}{2}\right)\det\left(D_{A + H_{\theta} + a} + \boldsymbol{\mu}\right)\\
&= \mathrm{zV}_n\int\frac{\mathrm{d}\theta}{(2\pi)^3}\int\mathrm{d}\gamma_N(F^{\mathrm{u}})\,\det\left(D_{A(F^{\mathrm{u}}) + H_{\theta} + a} + \boldsymbol{\mu}\right),
\end{align*}
that is, $\bar{Z}^{\mathrm{ext}}(t) = \mathrm{zV}_n^{-1}\,\mathbf{Z}^{\mathrm{x}}_n\left(a(t)\right)$ (this is the form (ii) in Annex~A, Section~G4.A.4).
By Proposition~38.3\,(a), which applies because $n \ge 4$ by (0.1), $\mathbf{Z}^{\mathrm{x}}_n$ is an entire function on $C^{1}(n) \otimes \mathbb{C}$; and $t \mapsto a(t)$ is complex linear.
Hence $\bar{Z}^{\mathrm{ext}}$ is an entire function of $t$.

\textbf{5.3. The identities for $\mathrm{Zmk}^{\theta}$.}
Let $g_{\mathrm{w}} \in W(\mathcal{Q}_{*})$ have the linear part $r$.
The invariance of $\gamma_N$ under $F^{\mathrm{u}} \mapsto g_{\mathrm{w}*}F^{\mathrm{u}}$ (Lemma~0.2\,(v)) and the first identity of (e.3) give
\begin{align*}
\mathrm{Zmk}^{r\theta}(g_{\mathrm{w}}t) &= \int\mathrm{d}\gamma_N(F^{\mathrm{u}})\,\mathrm{Igd}\left(F^{\mathrm{u}}, r\theta, g_{\mathrm{w}}t\right) = \int\mathrm{d}\gamma_N(F^{\mathrm{u}})\,\mathrm{Igd}\left(g_{\mathrm{w}*}F^{\mathrm{u}}, r\theta, g_{\mathrm{w}}t\right)\\
&= \int\mathrm{d}\gamma_N(F^{\mathrm{u}})\,\mathrm{Igd}(F^{\mathrm{u}},\theta,t) = \mathrm{Zmk}^{\theta}(t).
\end{align*}
The invariance of $\gamma_N$ under $F^{\mathrm{u}} \mapsto -F^{\mathrm{u}}$ and the second identity of (e.3) give
\[
\mathrm{Zmk}^{-\theta}(-t) = \int\mathrm{d}\gamma_N(F^{\mathrm{u}})\,\mathrm{Igd}\left(-F^{\mathrm{u}}, -\theta, -t\right) = \mathrm{Zmk}^{\theta}(t).
\]
As $\gamma_N$ is a positive measure, the third identity of (e.3) gives
\[
\overline{\mathrm{Zmk}^{\theta}(\overline{t})} = \int\mathrm{d}\gamma_N(F^{\mathrm{u}})\,\overline{\mathrm{Igd}\left(F^{\mathrm{u}}, \theta, \overline{t}\right)} = \mathrm{Zmk}^{\theta}(t).
\]

\textbf{5.4. The identities for $\bar{Z}^{\mathrm{ext}}$.}
The maps $\theta \mapsto r\theta$ ($r \in \{\pm1\}^3$) and $\theta \mapsto -\theta$ carry $[-\pi,\pi]^3$ onto itself and preserve $\mathrm{d}\theta$.
Hence, by 5.3,
\begin{align*}
\bar{Z}^{\mathrm{ext}}(g_{\mathrm{w}}t) &= \int\frac{\mathrm{d}\theta}{(2\pi)^3}\,\mathrm{Zmk}^{\theta}(g_{\mathrm{w}}t) = \int\frac{\mathrm{d}\theta}{(2\pi)^3}\,\mathrm{Zmk}^{r\theta}(g_{\mathrm{w}}t)\\
&= \int\frac{\mathrm{d}\theta}{(2\pi)^3}\,\mathrm{Zmk}^{\theta}(t) = \bar{Z}^{\mathrm{ext}}(t),\\
\bar{Z}^{\mathrm{ext}}(-t) &= \int\frac{\mathrm{d}\theta}{(2\pi)^3}\,\mathrm{Zmk}^{-\theta}(-t) = \bar{Z}^{\mathrm{ext}}(t),\\
\overline{\bar{Z}^{\mathrm{ext}}(\overline{t})} &= \int\frac{\mathrm{d}\theta}{(2\pi)^3}\,\overline{\mathrm{Zmk}^{\theta}(\overline{t})} = \bar{Z}^{\mathrm{ext}}(t).
\end{align*}

\textbf{5.5. Sign; $\bar{Z}^{\mathrm{ext}}(0) > 0$.}
For real $t$ the cochain $B(F^{\mathrm{u}},\theta,t)$ is real, and $\boldsymbol{\mu}$ is real; so $\mathrm{Igd} \ge 0$ by Lemma~0.5\,(i), and hence $\mathrm{Zmk}^{\theta}(t) \ge 0$ and $\bar{Z}^{\mathrm{ext}}(t) \ge 0$.
Consider the real pair $(m_{\mathrm{in}}, \varepsilon^{\mathrm{v}}_{\mathrm{in}}) := \left(\boldsymbol{\mu} - \bar{m}\cdot\eta, 0\right)$ on the torus with $n \ge 4$ sites per side.
By (I.7.2),
\[
\mathrm{dens}^{\theta,0}_N(F^{\mathrm{u}}) = \frac{\mathrm{Igd}(F^{\mathrm{u}},\theta,0)}{\det\left(D_0 + \bar{m}\cdot\eta\right)}, \qquad \text{where } \det\left(D_0 + \bar{m}\cdot\eta\right) > 0 \text{ (Lemma~0.5\,(ii))}.
\]
Hence $\bar{\nu}(1) = \bar{Z}^{\mathrm{ext}}(0)/\det\left(D_0 + \bar{m}\cdot\eta\right)$.
By Proposition~31.1\,(c), $\bar{\nu}(1) > 0$; therefore $\bar{Z}^{\mathrm{ext}}(0) > 0$.

\textbf{5.6. $\mathring{\Phi}$.}
The number $\bar{Z}^{\mathrm{ext}}(0)$ is positive.
Dividing the three identities of 5.4 by it, one obtains
\[
\mathring{\Phi}(g_{\mathrm{w}}t) = \mathring{\Phi}(t), \qquad \mathring{\Phi}(-t) = \mathring{\Phi}(t), \qquad \overline{\mathring{\Phi}(\overline{t})} = \mathring{\Phi}(t),
\]
and $\mathring{\Phi}$ is entire by 5.2.

\emph{One mark.}
Let $\mathfrak{c}$ and $\mathfrak{c}'$ be cells of the torus $\mathrm{kt}$, and let $z \in \mathbb{C}$.
By 1.5 there is a $g_{\mathrm{w}} \in W(\mathcal{Q}_{*})$ with $g_{\mathrm{w}}.\mathfrak{c} = \mathfrak{c}'$.
Let $\mathrm{ue}_{\mathfrak{c}} \in \mathbb{C}^{\mathrm{cells}}$ be the unit vector of the cell $\mathfrak{c}$, whose entry at $\mathfrak{c}$ is $1$ and whose other entries are $0$.
By the definition of $g_{\mathrm{w}}t$, one has $g_{\mathrm{w}}\left(z\,\mathrm{ue}_{\mathfrak{c}}\right) = z\,\mathrm{ue}_{g_{\mathrm{w}}.\mathfrak{c}} = z\,\mathrm{ue}_{\mathfrak{c}'}$.
So $\mathring{\Phi}\left(z\,\mathrm{ue}_{\mathfrak{c}'}\right) = \mathring{\Phi}\left(z\,\mathrm{ue}_{\mathfrak{c}}\right)$: the function $z \mapsto \mathring{\Phi}\left(z\,\mathrm{ue}_{\mathfrak{c}}\right)$ is the same for all cells.

\emph{Reality (every integer $m \ge 0$, every real $c_m$).}
The third identity, taken at real $t$, shows that $\mathring{\Phi}$ is real on $\mathbb{R}^{\mathrm{cells}}$.
If an entire function on $\mathbb{C}^{\mathrm{cells}}$ is real on $\mathbb{R}^{\mathrm{cells}}$, then so is each of its complex partial derivatives $\partial_{t_{\mathfrak{c}}}$: such a derivative is again entire, and at a real point it is the limit of the difference quotients with real increments of $t_{\mathfrak{c}}$, which are real.
By induction on the order, all partial derivatives of $\mathring{\Phi}$ are real on $\mathbb{R}^{\mathrm{cells}}$, in particular at $t = 0$.
Let $m \ge 0$ be an integer and $c_m$ a real number (the numbers $c_m$ of Section~I.7.6 are real by 4.6, on whichever torus $\mathrm{kt}$ they are formed).
The operators $\partial_{t_{\mathfrak{c}}}$ commute; hence
\[
\prod_{\mathfrak{c}}\left(\partial_{t_{\mathfrak{c}}}^m - c_m\right) = \sum_{\mathrm{Yc}}\left(-c_m\right)^{(2\mathrm{kt}\cdot L)^3 - \#|\mathrm{Yc}|}\prod_{\mathfrak{c} \in \mathrm{Yc}}\partial_{t_{\mathfrak{c}}}^m,
\]
where $\mathrm{Yc}$ ranges over the subsets of the set of the $(2\,\mathrm{kt}\cdot L)^3$ cells; this is a linear combination of partial derivatives with real coefficients.
Applied to $\mathring{\Phi}$ at $t = 0$, it yields a real number.
\end{proof}

\subsection*{6. Proof of (f)}

\begin{proof}[Proof of Proposition~49\,(f)]
Let $(m_{\mathrm{in}}, \varepsilon^{\mathrm{v}}_{\mathrm{in}})$ be a real pair and $\boldsymbol{\mu} := \bar{m}\cdot\eta + m_{\mathrm{in}}$.
By Proposition~31.1\,(a), (c), which apply because $n \ge 4$ by (0.1), the measure $\nu^{\theta} = \mathrm{dens}^{\theta,0}_N\,\gamma_N$ is a finite positive measure on $F_N$ with a bounded continuous density, $0 < \bar{\nu}(1) < \infty$, and $\bar{\mu} = \bar{\nu}/\bar{\nu}(1)$.
By (I.7.2),
\[
\mathrm{dens}^{\theta,0}_N(F^{\mathrm{u}}) = \exp\left(-\varepsilon^{\mathrm{v}}_{\mathrm{in}}n^3\right)\frac{\det\left(D_{A(F^{\mathrm{u}}) + H_{\theta}} + \boldsymbol{\mu}\right)}{\det\left(D_0 + \bar{m}\cdot\eta\right)}.
\]

\textbf{6.1.}
Let $g_{\mathrm{w}} = (a,r) \in \mathrm{Gc}$.
By Lemma~0.2\,(v),
\[
A(g_{\mathrm{w}*}F^{\mathrm{u}}) + H_{r\theta} = g_{\mathrm{w}*}\left(A(F^{\mathrm{u}}) + H_{\theta}\right), \qquad A(-F^{\mathrm{u}}) + H_{-\theta} = -\left(A(F^{\mathrm{u}}) + H_{\theta}\right).
\]
The first two identities of (c), with $z = \boldsymbol{\mu}$, give
\begin{equation}
\mathrm{dens}^{r\theta,0}_N\left(g_{\mathrm{w}*}F^{\mathrm{u}}\right) = \mathrm{dens}^{\theta,0}_N(F^{\mathrm{u}}), \qquad \mathrm{dens}^{-\theta,0}_N\left(-F^{\mathrm{u}}\right) = \mathrm{dens}^{\theta,0}_N(F^{\mathrm{u}}).
\tag{f.1}
\end{equation}

\textbf{6.2 (images of $\nu^{\theta}$).}
Let $G$ be a bounded Borel function on $F_N$.
By (f.1) and the invariance of $\gamma_N$ (Lemma~0.2\,(v)),
\begin{align*}
\int\mathrm{d}\nu^{\theta}(F^{\mathrm{u}})\,G\left(g_{\mathrm{w}*}F^{\mathrm{u}}\right) &= \int\mathrm{d}\gamma_N(F^{\mathrm{u}})\,G\left(g_{\mathrm{w}*}F^{\mathrm{u}}\right)\mathrm{dens}^{r\theta,0}_N\left(g_{\mathrm{w}*}F^{\mathrm{u}}\right)\\
&= \int\mathrm{d}\gamma_N(F^{\mathrm{u}})\,G(F^{\mathrm{u}})\,\mathrm{dens}^{r\theta,0}_N(F^{\mathrm{u}}) = \int\mathrm{d}\nu^{r\theta}\,G,\\
\int\mathrm{d}\nu^{\theta}(F^{\mathrm{u}})\,G\left(-F^{\mathrm{u}}\right) &= \int\mathrm{d}\gamma_N(F^{\mathrm{u}})\,G\left(-F^{\mathrm{u}}\right)\mathrm{dens}^{-\theta,0}_N\left(-F^{\mathrm{u}}\right) = \int\mathrm{d}\nu^{-\theta}\,G.
\end{align*}
Thus the image of $\nu^{\theta}$ under $F^{\mathrm{u}} \mapsto g_{\mathrm{w}*}F^{\mathrm{u}}$ is $\nu^{r\theta}$, and its image under $F^{\mathrm{u}} \mapsto -F^{\mathrm{u}}$ is $\nu^{-\theta}$.

\textbf{6.3 ($\bar{\nu}$, $\bar{\mu}$).}
One has
\[
\bar{\nu}(G) = \int_{[-\pi,\pi]^3}\frac{\mathrm{d}\theta}{(2\pi)^3}\,\nu^{\theta}(G),
\]
and the maps $\theta \mapsto r\theta$ and $\theta \mapsto -\theta$ preserve $[-\pi,\pi]^3$ and $\mathrm{d}\theta$.
Hence, by 6.2,
\[
\int\mathrm{d}\bar{\nu}\,G\left(g_{\mathrm{w}*}F^{\mathrm{u}}\right) = \int\frac{\mathrm{d}\theta}{(2\pi)^3}\,\nu^{r\theta}(G) = \bar{\nu}(G), \qquad \int\mathrm{d}\bar{\nu}\,G\left(-F^{\mathrm{u}}\right) = \int\frac{\mathrm{d}\theta}{(2\pi)^3}\,\nu^{-\theta}(G) = \bar{\nu}(G).
\]
After division by $\bar{\nu}(1) > 0$,
\[
\int\mathrm{d}\bar{\mu}\,G\left(g_{\mathrm{w}*}F^{\mathrm{u}}\right) = \bar{\mu}(G), \qquad \int\mathrm{d}\bar{\mu}\,G\left(-F^{\mathrm{u}}\right) = \bar{\mu}(G).
\]

\textbf{6.4 ($X_{\mathfrak{c}}$).}
By Lemma~0.2\,(iii), (iv),
\[
X_{\mathfrak{c}} = \left\langle d\hat{\mathfrak{a}}_{\mathfrak{c}}, F^{\mathrm{u}}\right\rangle = \left\langle d(\mathsf{g}_{\mathfrak{c}})_{*}\hat{\mathfrak{a}}, F^{\mathrm{u}}\right\rangle = \left\langle(\mathsf{g}_{\mathfrak{c}})_{*}d\hat{\mathfrak{a}}, F^{\mathrm{u}}\right\rangle = \left\langle d\hat{\mathfrak{a}}, \mathsf{g}_{\mathfrak{c}}^{*}F^{\mathrm{u}}\right\rangle.
\]
With $X(F^{\mathrm{u}}) = \langle d\hat{\mathfrak{a}}, F^{\mathrm{u}}\rangle$ this means $X_{\mathfrak{c}}(F^{\mathrm{u}}) = X\left(g_{\mathrm{w}*}F^{\mathrm{u}}\right)$ for $g_{\mathrm{w}} := \mathsf{g}_{\mathfrak{c}}^{-1} \in \mathrm{Gc}$, for which $g_{\mathrm{w}*} = \mathsf{g}_{\mathfrak{c}}^{*}$.
By 6.3, for every bounded Borel function $G$ on $\mathbb{R}$,
\[
\int\mathrm{d}\bar{\mu}\,G(X_{\mathfrak{c}}) = \int\mathrm{d}\bar{\mu}(F^{\mathrm{u}})\,G\left(X\left(g_{\mathrm{w}*}F^{\mathrm{u}}\right)\right) = \int\mathrm{d}\bar{\mu}\,G(X).
\]
So the law of $X_{\mathfrak{c}}$ under $\bar{\mu}$ is the law of $X$, for every $\mathfrak{c}$.
Finally $X_{\mathfrak{c}}(-F^{\mathrm{u}}) = -X_{\mathfrak{c}}(F^{\mathrm{u}})$; so 6.3 gives $\int\mathrm{d}\bar{\mu}\,G(-X_{\mathfrak{c}}) = \int\mathrm{d}\bar{\mu}\,G(X_{\mathfrak{c}})$, and the law is symmetric.
\end{proof}

\subsection*{7. Conditions on the constants}

\begin{description}
\item[$\rho \in (0,1)$.]
This condition is imposed in Section~I.8.2, line (C0), and in Proposition~14.2\,(a).
Within the five lemmas of Section~0 it enters only through the reality of $\rho$ (Lemmas~0.4 and 0.5\,(i)) and through $\rho \ge 0$ (Lemma~0.5\,(ii)).
Proposition~38.1\,(d), applied in 5.2, rests on $\rho \ge 0$, by (I.m6); Proposition~31.1\,(c), applied in 5.5 and in Section~6, rests on $\rho^2 \ne 1$, by \cite[Theorem~4.13(d)]{QEDII}.
In this way the clauses (c)--(f) depend on it.
\item[The standard weights.]
They are fixed in Sections~I.0.5 and I.1.2 and in the line (C0) of Section~I.8.2.
Within the five lemmas of Section~0 only two of their properties enter.
The first is $w_i > 0$: the reality of the $w_i$ is used in Lemmas~0.4\,(iii) and 0.5\,(i), and their positivity in Lemma~0.5\,(ii).
The second is the relation $\sum_i w_i\hat n_i = 0$, which is (I.1.1) and enters (0.4) through (I.s3).
Proposition~38.1\,(d), applied in 5.2, contains the value $w_3 = 6$, together with $w_0 + w_1 + w_2 = \mathrm{sw} - 6$, and rests on $w_i \ge 0$, by (I.m6).
The same holds for Proposition~38.3\,(a), applied in 5.2, through the function $\mathrm{Kd}$ of Chapter~38, Section~2, which is the right side of the inequality of Proposition~38.1\,(d).
Of the bound of Proposition~38.1\,(d), 5.2 uses only its finiteness and its independence of $(F^{\mathrm{u}},\theta)$.
Proposition~31.1\,(a), (c), applied in 5.5 and in Section~6, rests on $w_i > 0$, which is a hypothesis of \cite[Theorem~4.13]{QEDII}, and on (I.1.1).
In this way the clauses (c)--(f) depend on $w_i > 0$ and on (I.1.1), and the clause (e) depends in addition on the value $w_3 = 6$.
\item[$L$ an odd integer $\ge 5$.]
Section~I.0.5, the line (C3) of Section~I.8.2 and the row (L5) of Proposition~14.2\,(b) contain this condition.
By 0.1 it makes $\mathrm{per}$ an odd integer and gives $n \ge 5$, hence $n \ge 4$; every clause depends on it.
\item[$L \ge 13$.]
This bound is contained in the line (C3) of Section~I.8.2, which gives $L \ge 4^8$, and in the row (L1) of Proposition~14.2\,(b).
Clause (b) requires it through (Sc4): a power $L^i$ with $i \ge 1$ is an odd integer $\ge 13$.
The proof of \cite[Proposition~19.1(f)]{QEDI} also uses $L \ge 33$.
\item[$M$ a power of $L$, $M \ge L^2$.]
It is imposed in Section~I.0.5, in the line (C5) of Section~I.8.2 and in the row (M1) of Proposition~14.2\,(b).
In 0.1 it shows that $T^{*}(e)$ is an integer $\ge 1$, and in 2.1 that $R^{*}(e)$ lies in $\mathrm{Tor}(e)$; thus 0.1 and (b) rest on it.
\item[The number $e_1$.]
Its definition is given in the line (C9) of Section~I.8.2 and in Proposition~14.2, namely in (a), line (C9), entry~[118], and in (b), rows (e1)--(e8); see 2.1.
Of these rows only (e1) is used in clause (b), and of (e1) only its first three inequalities.
For all $v \ge -\log e_1$ one has $v \ge 8$, $v^r > 31L + 16$ and $\left(1 + 3\,\mathrm{lg}/v\right)^r < L$.
Clause (b) depends on these inequalities through (Sc0), (Sc1), (Sc4) and (Sc6): the first and the third, at $v = v_K$, give the hypotheses $v_K \ge 1$ and $\left(1 + \mathrm{lg}/v_K\right)^r \le L$ of \cite[Proposition~19.1]{QEDI}; the second, at $v = v_{*}$, gives the condition (2.1), $r_{*} > 31L + 16$, which 2.3 applies in the form $r_{*} \ge L^2$.
The condition (2.1) is part of Proposition~14.2\,(c) and has its own entry in Annex~B, Section~G4.A.6.
\item[The integer $r$ of the line (C7).]
Its range $r \ge 2$ is fixed in the line (C7) of Section~I.8.2 and in Proposition~14.2\,(a).
It enters through $r_k$ (Section~I.6.1) and through the inequalities listed under $e_1$ above.
Inside (Sc0), (Sc1) and (Sc4), namely in the step $\left(1 + \mathrm{lg}/v_K\right)^r \le \left(1 + 3\,\mathrm{lg}/v_K\right)^r$ and in \cite[Proposition~19.1(a), (b)]{QEDI}, it suffices that $r > 0$, that is, $r \ge 1$.
\end{description}

In (e) and (f) the lattice mass $\boldsymbol{\mu}$ is an arbitrary real number (I.K10).

\section*{Chapter 35a. Scales, cells, probes and the inputs at the stop level: proofs of Propositions 35.1--35.4}

\subsection*{0. Conventions of this chapter}

Letters, lattices and norms are those of Part~I and of Sections~G5.0.3--G5.1.9.
As in Proposition~14.2\,(a), $\mathrm{lh}(v) := L M v^r$; as in Proposition~6\,(f), $\mathrm{vv}$ denotes a vector of $\mathbb{C}^4$.
In Sections~1 and~5 the bare letter $v$ is the real variable of the conditions of Proposition~14.2\,(b), and in Sections~3 and~4 it is the placement of a probe.
In the paragraph on the torus in the proof of Proposition~35.2\,(a), $w$ is a vector of $\mathbb{Z}^3$ with the components $w_\mu$;
it enters through $n_* w$, in the same way as the vector $w$ of Sections~I.1.9 and~I.6.3, $n_* = \ell_0$ being the side.
The weights of Section~I.0.5 and the flux, for which the letter $w$ is also used, do not occur in this chapter.
The displays of this chapter are numbered (Eq1.1), (Eq2.1), \dots.
The labels (ps1)--(ps6) mark the preliminaries of Section~3.0, and (Pr-a)--(Pr-c), (Pr-e)--(Pr-g) are the clauses (a)--(c), (e)--(g) of \cite[Proposition~19.1]{QEDI} in the form stated in Step~0 of Section~1.
The labels (2.1)--(2.5) denote the displays of Proposition~35.4, except in Section~1, where (2.1) is the condition $r_{*} > 31L + 16$ of Section~G5.1.9;
(T1) and (T3) are the conditions of the same names in Section~G5.1.9.
The labels (e1)--(e8), (A1), (A2), (A8), (A10) denote rows of the table of Proposition~14.2\,(b), whereas (A6) and (A7) are clauses of Proposition~35.4;
(C0)--(C10) are the lines of the list of constants of Section~I.8.2, which is the table of Proposition~14.2\,(a).

\subsection*{1. Proposition 35.1 (scales on the reference torus)}

\begin{proof}[Proof of Proposition~35.1]
\emph{Step 0.}
In \cite[Proposition~19.1]{QEDI} and in the notation of \cite[Section~19.2]{QEDI} put $h' = h$ and $m_0 = 0$, and let the upper end of the mass range be $1$ (Section~I.6.1).
The numbers written $\hat{v}_j$, $\hat{r}_j$, $\hat{\mathfrak{n}}_j$, $j_*$ in \cite[Section~19.2]{QEDI} are denoted here by $\mathrm{vh}_j$, $\mathrm{rh}_j$, $\mathrm{nch}_j$, $\mathrm{js}$.
For $e > 0$, an integer $T \ge 0$, $R = h L^T$ and integers $j \ge 0$ they are
\begin{gather*}
\mathrm{vh}_j := v_K + j \cdot \mathrm{lg}, \qquad \mathrm{rh}_j := \min\left\{ L^i : i \ge 0 \text{ an integer},\ L^i \ge \mathrm{vh}_j^r \right\},\\
\mathrm{nch}_j := T + j - \log_L\left( M \cdot \mathrm{rh}_j \right).
\end{gather*}
These numbers are functions of $\left( e\,h^{\frac{1}{2}}, L, M, r, T, j \right)$; the integer $N$ does not enter them.
Let $0 \le k \le K$.
Section~I.6.1 gives $a_k = L^{k-N}$, $a_K = h$, $e_k = e\,a_k^{\frac{1}{2}}$ and $\ell_k = R/a_k$; hence
\begin{equation}
\begin{gathered}
a_k = h L^{-(K-k)}, \qquad v_k = v_K + (K-k) \cdot \mathrm{lg} = \mathrm{vh}_{K-k}, \qquad r_k = \mathrm{rh}_{K-k},\\
\ell_k = L^{T+K-k}, \qquad \mathfrak{n}_k = \mathrm{nch}_{K-k}.
\end{gathered}
\tag{Eq1.1}
\end{equation}
In this notation the clauses (a)--(c), (e)--(g) of \cite[Proposition~19.1]{QEDI} read as follows.

\begin{named}{Proposition}[{\cite[Proposition~19.1]{QEDI}}]
Let $v_K \ge 1$ and $\left( 1 + \frac{\mathrm{lg}}{v_K} \right)^r \le L$, and let $R = h L^T$ be an element of $\mathrm{Tor}(e)$.
Then the following hold.
\begin{itemize}
\item[(Pr-a)] For every $j \ge 0$ one has $\mathrm{rh}_{j+1} \in \left\{ \mathrm{rh}_j, L \cdot \mathrm{rh}_j \right\}$ and $\mathrm{nch}_{j+1} - \mathrm{nch}_j \in \{0, 1\}$.
\item[(Pr-b)] The number $\mathrm{js} = \mathrm{js}(R,e) := \max\left\{ j \ge 0 : \mathrm{nch}_j = 0 \right\}$ exists;
$\mathrm{nch}_j \le 0$ for $j \le \mathrm{js}$, and $\mathrm{nch}_j \ge 1$ for $j > \mathrm{js}$.
\item[(Pr-c)] One has $N_{\min}(R,e) = \max\left\{ 0, \mathrm{js} + 1 - \log_L h \right\}$.
For $N \ge N_{\min}(R,e)$ there is a first $k$ with $\mathfrak{n}_k = 0$;
it is $k_{*}(R) = K - \mathrm{js}$, and $1 \le k_{*} \le K$.
Moreover $a_{*} = h L^{-\mathrm{js}}$ and $\ell_* = M r_{*} = L^{T + \mathrm{js}}$.
\item[(Pr-e)] Let $N \ge N_{\min}(R,e)$ and $k \le k_{*}$.
Then $\mathfrak{n}_k \ge 0$, that is $\ell_k \ge M r_k$;
further $v_k \ge v_{*} \ge v_K$ and $e_k \le e_{*} \le e_K$.
Moreover $1/L < \bar{m} \cdot h \le 1$.
\item[(Pr-f)] Let $N \ge N_{\min}(R,e)$.
If $r_{*} > 31L + 16$, then $\ell_k \in \mathrm{Cl}$ for every $k \le k_{*}$.
\item[(Pr-g)] If $\left( 1 + \frac{3\,\mathrm{lg}}{v_K} \right)^r < L$, then $\mathrm{js}\left( R^{*}(e), e \right) \in \{1, 2\}$.
Consequently, on $R^{*}(e)$ and for $N \ge N_{\min}\left( R^{*}(e), e \right)$, one has $k_{*} \in \{K-1, K-2\}$, $a_{*} \in \left\{ h/L, h/L^2 \right\}$ and $r_{*} \in \left\{ r_K, L r_K \right\}$.
\end{itemize}
\end{named}

\begin{proof}
By (C5), $\log_L M$ is an integer $\ge 2$, and every $\mathrm{rh}_j$ is a power of $L$;
so all the numbers $\mathrm{nch}_j$ are integers.

\emph{(Pr-a)}
From $\mathrm{vh}_{j+1} > \mathrm{vh}_j$ one obtains $\mathrm{rh}_{j+1} \ge \mathrm{rh}_j$.
Since $\mathrm{vh}_j \ge v_K$,
\[
\mathrm{vh}_{j+1}^r = \mathrm{vh}_j^r \left( 1 + \frac{\mathrm{lg}}{\mathrm{vh}_j} \right)^r \le \mathrm{vh}_j^r \left( 1 + \frac{\mathrm{lg}}{v_K} \right)^r \le L \cdot \mathrm{vh}_j^r \le L \cdot \mathrm{rh}_j ,
\]
and $L \cdot \mathrm{rh}_j$ is a power of $L$;
therefore $\mathrm{rh}_{j+1} \le L \cdot \mathrm{rh}_j$.
Both numbers are powers of $L$, so that $\mathrm{rh}_{j+1} \in \left\{ \mathrm{rh}_j, L \cdot \mathrm{rh}_j \right\}$ and
\[
\mathrm{nch}_{j+1} - \mathrm{nch}_j = 1 - \log_L \frac{\mathrm{rh}_{j+1}}{\mathrm{rh}_j} \in \{0, 1\}.
\]

\emph{(Pr-b)}
One has $\mathrm{rh}_0 = r_K$, and the inequality $R \le R_{\mathrm{top}}(e) = h M r_K$ says $L^T \le M \cdot \mathrm{rh}_0$;
so $\mathrm{nch}_0 \le 0$.
Next, $\mathrm{rh}_j < L \cdot \mathrm{vh}_j^r$:
if $\mathrm{rh}_j = L^i$ with $i \ge 1$, then $L^{i-1} < \mathrm{vh}_j^r$ by minimality;
if $\mathrm{rh}_j = 1$, then $1 < L \le L \cdot \mathrm{vh}_j^r$, because $\mathrm{vh}_j \ge v_K \ge 1$.
Hence
\[
\mathrm{nch}_j > T + j - \log_L\left( L M \left( v_K + j \cdot \mathrm{lg} \right)^r \right),
\]
and the right side tends to infinity with $j$.
A sequence of integers that begins with a value $\le 0$, tends to infinity and increases by steps $0$ or $1$ takes the value $0$.
Let $\mathrm{js}$ be the last index at which it does;
then $\mathrm{nch}_j \le 0$ for $j \le \mathrm{js}$ and $\mathrm{nch}_j \ge 1$ for $j > \mathrm{js}$.

\emph{(Pr-c)}
By (Eq1.1), $R L^N = R/\eta = \ell_0$.
So the inequality $R L^N \ge L M r_0$ says $\mathfrak{n}_0 \ge 1$, that is $\mathrm{nch}_K \ge 1$, that is $K > \mathrm{js}$, by (Pr-b).
The number $K = N + \log_L h$ increases with $N$.
Hence, for an integer $N_1 \ge 0$, the inequality holds for all $N \ge N_1$ if and only if $N_1 + \log_L h \ge \mathrm{js} + 1$.
The least such $N_1$ is $\max\left\{ 0, \mathrm{js} + 1 - \log_L h \right\}$, and this is the value of $N_{\min}(R,e)$.
Let now $N \ge N_{\min}(R,e)$;
then $K \ge \mathrm{js} + 1$.
For $0 \le k < K - \mathrm{js}$ one has $\mathfrak{n}_k = \mathrm{nch}_{K-k} \ge 1$, because $K - k > \mathrm{js}$;
and $\mathfrak{n}_{K - \mathrm{js}} = \mathrm{nch}_{\mathrm{js}} = 0$.
So the first $k$ with $\mathfrak{n}_k = 0$ is $k_{*} = K - \mathrm{js}$, and $1 \le k_{*} \le K$ because $0 \le \mathrm{js} \le K - 1$.
The equation $\mathfrak{n}_{k_{*}} = 0$ says $\ell_* = M r_{*}$.
By (Eq1.1), $a_{*} = h L^{-\mathrm{js}}$ and $\ell_* = L^{T + \mathrm{js}}$.

\emph{(Pr-e)}
Let $k \le k_{*}$, that is $K - k \ge \mathrm{js}$.
Then $\mathfrak{n}_k = \mathrm{nch}_{K-k} \ge 0$ by (Pr-b), and $\ell_k = M r_k L^{\mathfrak{n}_k} \ge M r_k$.
The numbers $\mathrm{vh}_j$ increase with $j$;
so $v_k = \mathrm{vh}_{K-k} \ge \mathrm{vh}_{\mathrm{js}} = v_{*} \ge \mathrm{vh}_0 = v_K$, and, since $e_k = \exp(-v_k)$, also $e_k \le e_{*} \le e_K$.
Put $\mathrm{tl} := \log_L(1/\bar{m}) \ge 0$.
Then $h = L^{\lfloor \mathrm{tl} \rfloor}$, so that $\bar{m} \cdot h = L^{\lfloor \mathrm{tl} \rfloor - \mathrm{tl}}$ lies in $(1/L, 1]$.

\emph{(Pr-f)}
The number $r_{*}$ is a power of $L$, and $r_{*} > 31L + 16 > L$;
so $r_{*} \ge L^2$.
Let $k \le k_{*}$.
Then $r_k = \mathrm{rh}_{K-k} \ge \mathrm{rh}_{\mathrm{js}} = r_{*}$ by (Pr-a), and $\mathfrak{n}_k \ge 0$ by (Pr-e);
hence
\[
\frac{\ell_k}{LM} = r_k L^{\mathfrak{n}_k - 1}
\]
is a power $L^i$ of the odd number $L$ with $i \ge 1$, and $L \ge 13$ by (C3).
So $\ell_k/(LM)$ is an odd integer $\ge 13$, that is, $\ell_k \in \mathrm{Cl}$ (Section~I.1.1).

\emph{(Pr-g)}
For $i = 1, 2, 3$ one has
\[
\mathrm{vh}_i^r = v_K^r \left( 1 + \frac{i \cdot \mathrm{lg}}{v_K} \right)^r < L v_K^r \le L \cdot \mathrm{rh}_0 ;
\]
since $L \cdot \mathrm{rh}_0$ is a power of $L$ and $\mathrm{rh}_i \ge \mathrm{rh}_0$ by (Pr-a), it follows that $\mathrm{rh}_i \in \left\{ \mathrm{rh}_0, L \cdot \mathrm{rh}_0 \right\}$.
The side $R^{*}(e) = h M r_K/L$ equals $h L^T$ with the integer $T = \log_L M - 1 + \log_L r_K \ge 1$, and $R^{*}(e) \le R_{\mathrm{top}}(e)$;
so $R^{*}(e)$ lies in $\mathrm{Tor}(e)$.
On $R^{*}(e)$ one has $L^T = M \cdot \mathrm{rh}_0/L$, hence
\[
\mathrm{nch}_j = j - 1 - \log_L \frac{\mathrm{rh}_j}{\mathrm{rh}_0},
\]
and therefore $\mathrm{nch}_0 = -1$, $\mathrm{nch}_1 \in \{0, -1\}$, $\mathrm{nch}_2 \in \{1, 0\}$, $\mathrm{nch}_3 \in \{2, 1\}$.
By (Pr-a) the sequence $(\mathrm{nch}_j)$ does not decrease, and $\mathrm{nch}_3 \ge 1$ gives $\mathrm{js} \le 2$.
If $\mathrm{nch}_2 = 0$, then $\mathrm{js} = 2$ and $\mathrm{rh}_2 = L \cdot \mathrm{rh}_0$.
Otherwise $\mathrm{nch}_2 = 1$;
then $\mathrm{nch}_1$ lies in $\{0, 1\}$ by (Pr-a) and in $\{0, -1\}$, so that $\mathrm{nch}_1 = 0$, $\mathrm{js} = 1$ and $\mathrm{rh}_1 = \mathrm{rh}_0$.
The remaining assertions are (Pr-c) with $r_{*} = \mathrm{rh}_{\mathrm{js}}$:
one has $(a_{*}, r_{*}) = (h/L, r_K)$ if $\mathrm{js} = 1$, and $(a_{*}, r_{*}) = (h/L^2, L r_K)$ if $\mathrm{js} = 2$.
\end{proof}

\emph{Step 1 (the number $e_1$; the conditions at $v_K$ and at $v_{*}$).}
By Proposition~14.2\,(a), line (C9), $e_1$ is one half of the supremum of the set $\mathrm{Se}$ of the numbers $e' \in \left( 0, \min(e_0, 1/3) \right]$ such that (e1)--(e8) hold for all $v \ge -\log e'$.
The data contain a number $e > 0$ with $e\,h^{\frac{1}{2}} \le e_1$;
so $e_1 > 0$, $\sup \mathrm{Se} = 2 e_1 > e_1$, and some $e' \in \mathrm{Se}$ satisfies $e' > e_1$.
Hence
\begin{equation}
e_1 \le \frac{1}{2} \min\left( e_0, \frac{1}{3} \right), \qquad \text{and (e1)--(e8) hold for all } v \ge -\log e_1 . \tag{Eq1.2}
\end{equation}
In particular, the rows (e1) and (e4) read
\begin{align*}
&\text{(e1)} && v \ge 8, \quad v^r > 31L + 16, \quad \left( 1 + \frac{3\,\mathrm{lg}}{v} \right)^r < L, \quad p_0\, v^{p_0 - 1} \ge 28;\\
&\text{(e4)} && 4 C_0 \cdot \mathrm{lh}(v)^2 \cdot v^p < \exp(\varepsilon v), \quad \mathrm{lh}(v) = L M v^r .
\end{align*}
Since $v_K = -\log\left( e\,h^{\frac{1}{2}} \right) \ge -\log e_1$, (e1) holds at $v = v_K$:
\[
v_K \ge 8, \qquad v_K^r > 31L + 16, \qquad \left( 1 + \frac{\mathrm{lg}}{v_K} \right)^r \le \left( 1 + \frac{3\,\mathrm{lg}}{v_K} \right)^r < L .
\]
These are the hypotheses on $v_K$ of the proposition of Step~0 and of its clause (Pr-g).
By (C5), $M$ is a power of $L$ with $M \ge L^2$, so that $\log_L M$ is an integer $\ge 2$;
and $r_K = L^i$ with an integer $i \ge 0$, by the definition of $r_k$ in Section~I.6.1.
So $T^{*}(e) = \log_L\left( M r_K \right) - 1 = \log_L M + i - 1$ is an integer $\ge 1$, $R^{*}(e) = h L^{T^{*}(e)}$, and $R^{*}(e) = R_{\mathrm{top}}(e)/L \le R_{\mathrm{top}}(e)$;
thus $R^{*}(e)$ lies in $\mathrm{Tor}(e)$ (Section~I.6.2).
The datum $N \ge N_0(e)$ means $N \ge N_{\min}\left( R^{*}(e), e \right)$ (Section~I.6.2).
Therefore (Pr-a)--(Pr-c) and (Pr-e)--(Pr-g) hold for $R = R^{*}(e)$ and $T = T^{*}(e)$, with $\mathrm{js} := \mathrm{js}\left( R^{*}(e), e \right)$.
By (Pr-e), $v_{*} \ge v_K \ge -\log e_1$;
so (e1) and (e4) hold at $v = v_{*}$ as well.

\emph{The two assertions that precede the datum $N$.}
That $R^{*}(e)$ lies in $\mathrm{Tor}(e)$ was shown in Step~1.
The hypotheses of the proposition of Step~0 at $R = R^{*}(e)$, $T = T^{*}(e)$ are conditions on $v_K$ and on $R^{*}(e)$;
Step~1 verifies them without the datum $N$, and the numbers $\mathrm{vh}_j$, $\mathrm{rh}_j$, $\mathrm{nch}_j$, $\mathrm{js}$ do not involve $N$ (Step~0).
By (Pr-b) the number $\mathrm{js} = \mathrm{js}\left( R^{*}(e), e \right)$ exists.
By the proof of (Pr-c), an integer $N_1 \ge 0$ has the property required in Section~I.6.2, namely $R^{*}(e) L^N \ge L M r_0$ for all $N \ge N_1$, if and only if $N_1 + \log_L h \ge \mathrm{js} + 1$.
Such integers exist.
So the set of the integers $N_1$ of Section~I.6.2 is not empty, and
\[
N_0(e) = N_{\min}\left( R^{*}(e), e \right) = \max\left\{ 0, \mathrm{js} + 1 - \log_L h \right\}
\]
is finite.

\emph{(a)}
Since $\bar{m} \le 1$, one has $\log_L(1/\bar{m}) \ge 0$;
hence $h = L^{\lfloor \log_L(1/\bar{m}) \rfloor} \ge 1$ (Section~I.0.5), and $e = e_K h^{-\frac{1}{2}} \le e_K \le e_1 \le 1/3$ by (Eq1.2).
The inequality $v_K \ge 8$ and the fact that $T^{*}(e)$ is an integer $\ge 1$ were shown in Step~1.
By (Pr-g), $k_{*} \in \{K-1, K-2\}$, $r_{*} \in \left\{ r_K, L r_K \right\}$ and $a_{*} \in \left\{ h/L, h/L^2 \right\}$;
by (Pr-c), $k_{*} \ge 1$ and $\ell_* = M r_{*}$.
The mass $\mu_* = \bar{m} \cdot a_{*}$ (Section~I.6.1) is one of the numbers $\bar{m} \cdot h/L$, $\bar{m} \cdot h/L^2$, and $1/L < \bar{m} \cdot h \le 1$ by (Pr-e);
so $L^{-3} < \mu_* \le 1/L$, while $\mu_- = L^{-3}$ (Section~G5.1.1).
Next, $e_{*} = e\,a_{*}^{\frac{1}{2}}$ (Section~I.6.1) with $h/L^2 \le a_{*} \le h/L \le h$;
hence
\[
e \left( \frac{h}{L^2} \right)^{\frac{1}{2}} \le e_{*} \le e\,h^{\frac{1}{2}} = e_K .
\]
Here $e_K \le e_1$ is one of the data, and $e_1 \le e_0$ by (Eq1.2).
By Section~I.3.4 and line (C0), $b_k = b \left( 1 - L^{-1} \right)/\left( 1 - L^{-k} \right)$ with $b = 1$.
For $k \ge 1$ one has $1 - L^{-1} \le 1 - L^{-k} < 1$, so that $b_{-} = 1 - 1/L < b_{k} \le 1$;
at $k = k_{*} \ge 1$ this gives $b_{-} \le b_{*} \le 1$.
Finally $\eta_* = L^{-k_{*}} \le 1/L$, because $k_{*} \ge 1$.

\emph{(b)}
One has $v_{*} = -\log e - \frac{1}{2} \log a_{*} = v_K + \frac{1}{2} \log\left( h/a_{*} \right)$, and $h/a_{*} \in \left\{ L, L^2 \right\}$ by (a);
so $v_K \le v_{*} \le v_K + \log L$.
The number $r_{*}$ is the least power $L^i$, $i \ge 0$, with $L^i \ge v_{*}^r$;
so $r_{*} \ge v_{*}^r$.
If $r_{*} = L^i$ with $i \ge 1$, then $L^{i-1} < v_{*}^r$ by minimality, that is $r_{*} < L v_{*}^r$;
if $r_{*} = 1$, then $r_{*} < L \le L v_{*}^r$, because $v_{*} \ge 8 \ge 1$.
By (e1) at $v = v_{*}$ one has $r_{*} \ge v_{*}^r > 31L + 16 > L$, and a power of $L$ that exceeds $L$ is at least $L^2$.
Since $\ell_* = M r_{*}$ by (a), the inequalities $v_{*}^r \le r_{*} < L v_{*}^r$ become $M v_{*}^r \le \ell_* < L M v_{*}^r$.

\emph{(c)}
The number $r_K$ is a function of $v_K = -\log\left( e\,h^{\frac{1}{2}} \right)$ and of $(L, r)$.
So $T^{*}(e)$, $R^{*}(e) = h M r_K/L$ and $\ell_c = R^{*}(e)/(2L)$ (Sections~I.6.2 and~I.6.4) are functions of $(e, h; L, M, r)$.
The numbers $\mathrm{nch}_j$ formed with $T = T^{*}(e)$, and hence $\mathrm{js}$, are functions of $\left( e\,h^{\frac{1}{2}}; L, M, r \right)$.
By (Pr-c), $K - k_{*} = \mathrm{js}$ and $a_{*} = h L^{-\mathrm{js}}$;
the number $h L^{-\mathrm{js}}$ is defined without $N$, and in the equivalence below $a_{*}$ denotes this number.
Then the numbers
\begin{gather*}
e_{*} = e\,a_{*}^{\frac{1}{2}}, \qquad v_{*} = -\log e_{*},\\
r_{*} = \min\left\{ L^i : i \ge 0 \text{ an integer},\ L^i \ge v_{*}^r \right\}, \qquad \ell_* = M r_{*},\\
\mu_* = \bar{m} \cdot a_{*}, \qquad p_{*} = v_{*}^p, \qquad p_{0*} = v_{*}^{p_0}
\end{gather*}
(Section~I.6.1) are functions of $(e, \bar{m}; L, M, r, p, p_0)$.
None of these numbers involves $N$.
Moreover $\log_L a_{*} = \log_L h - \mathrm{js}$, and (Pr-c) gives
\[
N_0(e) = \max\left\{ 0, \mathrm{js} + 1 - \log_L h \right\} = \max\left\{ 0, 1 - \log_L a_{*} \right\}.
\]
So, for an integer $N$, one has $N \ge N_0(e)$ if and only if $N \ge 0$ and $N + \log_L a_{*} \ge 1$.
Since $K = N + \log_L h$ (Section~I.6.1), $k_{*} = K - \mathrm{js} = N + \log_L h - \mathrm{js} = N + \log_L a_{*}$.
The numbers $k_{*}$, $\eta_*$, $b_{*}$, $n_*$ and $\mathrm{per}$ depend on $N$:
$k_{*} = N + \log_L a_{*}$ increases strictly with $N$;
$\eta_* = L^{-k_{*}}$ and $b_{*} = \left( 1 - L^{-1} \right)/\left( 1 - L^{-k_{*}} \right)$ are strictly decreasing functions of $k_{*}$;
and $n_* = R^{*}(e) L^N$ and $\mathrm{per} = n_*/L$ increase strictly with $N$.

\emph{(d)}
The inequality $r_{*} > 31L + 16$ was shown in (b);
so (Pr-f) gives $\ell_j \in \mathrm{Cl}$ for $0 \le j \le k_{*}$.
In the quadruple $(k_{*}, \ell_*, \mu_*, e_{*})$, $k_{*}$ is an integer $\ge 1$, $\ell_* \in \mathrm{Cl}$, $\mu_* \in \left( L^{-3}, 1/L \right] \subset (0, 1]$ and $e_{*} \in (0, e_0]$, by (a).
These are data (SF) of Section~I.1.8.
Let $0 \le j < k_{*}$.
Since $a_{j+1} = L a_j$, one has $\ell_{j+1} = R^{*}(e)/a_{j+1} = \ell_j/L$, and both $\ell_j$ and $\ell_{j+1}$ lie in $\mathrm{Cl}$, because $j + 1 \le k_{*}$.
By (Pr-e), $\ell_j \ge M r_j$.
The mass $\bar{m} \cdot a_j = \mu_* L^{-(k_{*} - j)}$ lies in $\left( 0, \mu_*/L \right] \subset (0, 1/L]$.
Further $e_{j+1} = e\,a_{j+1}^{\frac{1}{2}} = L^{\frac{1}{2}} e_j$, and $e_{j+1} \le e_{*} \le e_0$ by (Pr-e) (as $j + 1 \le k_{*}$) and by (a).
At $j = 0$ one has $a_0 = L^{-N} = \eta$;
so the mass is $\bar{m} \cdot \eta$, and the charge $e\,a_0^{\frac{1}{2}} = e\,\eta^{\frac{1}{2}}$ is $q$ (Section~I.2.5).

\emph{(e)}
Condition (2.1), that is $r_{*} > 31L + 16$, was shown in (b).
The first two clauses of (T1), namely $v_{*} \ge 8$ and $p_0\, v_{*}^{p_0 - 1} \ge 28$, are the first and the fourth inequality of (e1) at $v = v_{*}$.
For (T3): by (b), $\ell_* < L M v_{*}^r = \mathrm{lh}(v_{*})$;
moreover $p_{*} = v_{*}^p$ and $e_{*}^{-\varepsilon} = \exp(\varepsilon v_{*})$.
Hence, by (e4) at $v = v_{*}$,
\[
4 C_0 \cdot \ell_*^2 \cdot p_{*} \le 4 C_0 \cdot \mathrm{lh}(v_{*})^2 \cdot v_{*}^p < \exp(\varepsilon v_{*}) = e_{*}^{-\varepsilon} .
\]

\emph{(f)}
One has $\ell_c = R^{*}(e)/(2L) = h M r_K/(2 L^2)$, and $r_K \ge v_K^r$ by the definition of $r_K$.
Next, $n_* = R^{*}(e)/\eta$ and $\ell_* = R^{*}(e)/a_{*}$ with $a_{*} = L^{k_{*} - N} = L^{k_{*}} \eta$;
so $n_* = L^{k_{*}} \ell_*$, and
\[
\mathrm{per} = \frac{n_*}{L} = L^{k_{*} - 1} \ell_* = L^{k_{*} - 1} M r_{*} .
\]
As $k_{*} \ge 1$, $M \ge L^2$ and $r_{*} \ge L^2$, it follows that $\mathrm{per} \ge L^4$.
Further $R^{*}(e) = h M r_K/L \ge h L^2/L = h L$, because $M \ge L^2$ and $r_K \ge 1$.
Since $\bar{m} \cdot h > 1/L$ by (Pr-e) and $M \ge L^2$,
\[
\bar{m} \cdot R^{*}(e) = \left( \bar{m} \cdot h \right) \frac{M r_K}{L} > \frac{1}{L} \cdot \frac{L^2 r_K}{L} = r_K .
\]
Moreover $r_K \ge v_K^r \ge v_K$, because $v_K \ge 1$ and $r \ge 1$.
Hence $\exp\left( -\bar{m} \cdot R^{*}(e) \right) < \exp(-v_K) = e_K = e\,h^{\frac{1}{2}}$.
\end{proof}

\begin{named}{Remark}[Uniformity]
The assertions of Proposition~35.1 hold for every $\bar{m} \in (0, 1]$, every $e$ with $e\,h^{\frac{1}{2}} \le e_1$ and every $N \ge N_0(e)$.
Every bound in them is explicit, and the only constants that enter them are $b$ (line (C0)), $L$ (line (C3)), $b_{-}$ and $\mu_-$ (line (C4)), and $M$, $r$, $p$, $p_0$, $\varepsilon$, $C_{0}$, $e_0$, $e_1$ (lines (C5)--(C9)).
\end{named}

\subsection*{2. Proposition 35.2 (cells and the groups of the scheme)}

\begin{proof}[Proof of Proposition~35.2]
\emph{Step 0 (the group of the walls).}
The coordinates are lattice units at the cutoff $N$ (Section~I.6.4):
$\mathrm{per} = n_*/L$ is an odd integer, $\mathrm{hc} = \mathrm{per}/2$, and the walls are the planes $\left\{ y_\mu = j \cdot \mathrm{hc} \right\}$, $j \in \mathbb{Z}$, $\mu = 1, 2, 3$.
The reflection in the wall $\left\{ y_\mu = j \cdot \mathrm{hc} \right\}$ is $y_\mu \mapsto 2j \cdot \mathrm{hc} - y_\mu = j \cdot \mathrm{per} - y_\mu$, and it fixes the two other coordinates.
Consider the maps
\begin{equation}
y \mapsto r y + a, \qquad r \in \{\pm1\}^3, \quad a \in \mathrm{per}\,\mathbb{Z}^3 . \tag{Eq2.1}
\end{equation}
They form a group under composition:
inserting $y' = r' y + a'$ into $r y' + a$ gives $r r' y + (a + r a')$, and $r a' \in \mathrm{per}\,\mathbb{Z}^3$;
the inverse of (Eq2.1) is $y \mapsto r y - r a$.
The reflections in the walls are of the form (Eq2.1).
So every element of $W(\mathcal{Q}_{*})$, the group generated by these reflections, is of the form (Eq2.1).
Conversely, the reflection in $\left\{ y_\mu = 0 \right\}$ followed by the reflection in $\left\{ y_\mu = j \cdot \mathrm{hc} \right\}$ is $y_\mu \mapsto j \cdot \mathrm{per} + y_\mu$.
Hence $W(\mathcal{Q}_{*})$ contains the translations by $\mathrm{per}\,\mathbb{Z}^3$ and the three maps $y_\mu \mapsto -y_\mu$, and therefore all maps (Eq2.1).
A map (Eq2.1) that maps $Q^{\circ}_{0} = (0, \mathrm{hc})^3$ onto itself is the identity.
Indeed, in the coordinate $\mu$ it is $y_\mu \mapsto y_\mu + a_\mu$ or $y_\mu \mapsto -y_\mu + a_\mu$;
it maps the interval $(0, \mathrm{hc})$ onto itself only if $a_\mu = 0$ in the first case and $a_\mu = \mathrm{hc}$ in the second;
and $\mathrm{hc} = \mathrm{per}/2$ does not lie in $\mathrm{per}\,\mathbb{Z}$.
Consequently, for every $\mathfrak{c} \in \mathbb{Z}^3$, at most one element of $W(\mathcal{Q}_{*})$ maps $Q^{\circ}_{0}$ onto $Q^{\circ}_{\mathfrak{c}}$.

\emph{(a)}
Let $\mathfrak{c} \in \mathbb{Z}^3$, and let $r_{\mathfrak{c}}$ and $a_{\mathfrak{c}}$ be as in Proposition~35.2\,(a).
In the coordinate $\mu$ the map $y \mapsto r_{\mathfrak{c}} y + a_{\mathfrak{c}}$ is $y_\mu \mapsto i \cdot \mathrm{per} + y_\mu$ if $\mathfrak{c}_\mu = 2i$, and $y_\mu \mapsto (i+1) \cdot \mathrm{per} - y_\mu$ if $\mathfrak{c}_\mu = 2i + 1$;
this is the map of Section~I.6.4, and it is of the form (Eq2.1).
In the first case it maps $(0, \mathrm{hc})$ onto $\left( i \cdot \mathrm{per}, i \cdot \mathrm{per} + \mathrm{hc} \right) = \left( 2i \cdot \mathrm{hc}, (2i+1) \cdot \mathrm{hc} \right)$;
in the second case onto $\left( (i+1) \cdot \mathrm{per} - \mathrm{hc}, (i+1) \cdot \mathrm{per} \right) = \left( (2i+1) \cdot \mathrm{hc}, (2i+2) \cdot \mathrm{hc} \right)$.
In both cases the image is $\left( \mathfrak{c}_\mu \cdot \mathrm{hc}, (\mathfrak{c}_\mu + 1) \cdot \mathrm{hc} \right)$.
So the map carries $Q^{\circ}_{0}$ onto $Q^{\circ}_{\mathfrak{c}}$, and by Step~0 it is the element $\mathsf{g}_{\mathfrak{c}}$ of $W(\mathcal{Q}_{*})$.

\emph{The torus.}
One has $n_* = L \cdot \mathrm{per}$.
If $\mathfrak{c}$ is replaced by $\mathfrak{c} + 2L w$, $w \in \mathbb{Z}^3$, then $(-1)^{\mathfrak{c}_\mu}$ does not change and $i$ is replaced by $i + L w_\mu$;
so $r_{\mathfrak{c}}$ does not change, and $a_{\mathfrak{c}}$ is replaced by $a_{\mathfrak{c}} + n_* w$.
Hence, as a map of $\mathbb{Z}^3/n_*\mathbb{Z}^3$, $\mathsf{g}_{\mathfrak{c}}$ depends only on the class of $\mathfrak{c}$ in $(\mathbb{Z}/2L)^3$, and $a_{\mathfrak{c}}$ is a class in $\mathrm{per}\,\mathbb{Z}^3/n_*\mathbb{Z}^3$.
This proves the last sentence of (a) in its parts on $r_{\mathfrak{c}}$ and $a_{\mathfrak{c}}$;
the class of $L^{-j} a_{\mathfrak{c}}$ is treated in the proof of (b).

\emph{(b)}
\emph{Coordinates.}
The site $c^0 \in \mathbb{Z}^3$ (lattice units, level $0$) has the coordinates $c = L^{-j} c^0$ in $\mathbb{L}_j = L^{-j}\mathbb{Z}^3$ of level $j$;
here one unit of level $j$ is $a_j/\eta = L^j$ lattice units (Section~I.6.1), and $a^{[j]}(b) = L^{\frac{j}{2}} a\left( L^j b \right)$ (Section~I.6.6).
In these coordinates $\mathbb{Z}^3/n_*\mathbb{Z}^3$ becomes $T_j = \mathbb{L}_j/\ell_j\mathbb{Z}^3$, with $\ell_j = R^{*}(e)/a_j = n_*/L^j$, and $\mathsf{g}_{\mathfrak{c}}$ is the map
\[
c \mapsto L^{-j} \left( r_{\mathfrak{c}} L^j c + a_{\mathfrak{c}} \right) = r_{\mathfrak{c}} c + L^{-j} a_{\mathfrak{c}} .
\]

\emph{The translation part.}
The relations $n_* = R^{*}(e)/\eta$, $\ell_* = R^{*}(e)/a_{*}$, $a_{*} = L^{k_{*}} \eta$ (Section~I.6.1) and $\ell_* = M r_{*}$ (Proposition~35.1\,(a)) give $\mathrm{per} = n_*/L = L^{k_{*} - 1} M r_{*}$.
So $L^{-j} \cdot \mathrm{per} = L^{k_{*} - 1 - j} M r_{*}$.
By Proposition~35.1\,(b), $r_{*} = L^i$ with an integer $i \ge 2$.
For $0 \le j \le k_{*} - 1$ this gives $L^{-j} \cdot \mathrm{per} = LM \cdot L^{k_{*} - 1 - j} L^{i-1}$, and for $j = k_{*}$ it gives $L^{-j} \cdot \mathrm{per} = LM \cdot L^{i-2}$.
In both cases $L^{-j} \cdot \mathrm{per}$ is an integer multiple of $LM$.
Since $a_{\mathfrak{c}} \in \mathrm{per}\,\mathbb{Z}^3$, it follows that $L^{-j} a_{\mathfrak{c}} \in LM\,\mathbb{Z}^3$;
this vector is determined modulo $L^{-j} n_*\mathbb{Z}^3 = \ell_j\mathbb{Z}^3$.

\emph{The matrices $S$.}
Let $\mathrm{Rf} := \left\{ \mu : \mathfrak{c}_\mu \text{ odd} \right\}$, and let $S_{\mathrm{Rf}}$ be the product of the matrices $\gamma_\mu\gamma_5$, $\mu \in \mathrm{Rf}$, in increasing order of $\mu$ ($S_{\mathrm{Rf}} := 1$ if $\mathrm{Rf}$ is empty).
By the table of Section~I.4.1, $r(\gamma_\mu\gamma_5)$ has the entry $-1$ at $\mu$ and $+1$ at the two other places, and $r$ is a homomorphism;
so $r(S_{\mathrm{Rf}}) = r_{\mathfrak{c}}$.
An element $S$ of $\mathrm{Sp}$ satisfies $r(S) = r_{\mathfrak{c}}$ if and only if $S S_{\mathrm{Rf}}^{-1}$ lies in $\ker r = \left\{ 1, -1, i\Sigma_3, -i\Sigma_3 \right\}$ (Section~I.4.1).
So there are exactly four such $S$, namely
\[
S_{\mathrm{Rf}}, \qquad -S_{\mathrm{Rf}}, \qquad i\Sigma_3 S_{\mathrm{Rf}}, \qquad -i\Sigma_3 S_{\mathrm{Rf}} .
\]
By (I.c1), $\gamma_5\gamma_\mu = -\gamma_\mu\gamma_5$ and $\gamma_5^2 = 1$;
hence $\gamma_5 \left( \gamma_\mu\gamma_5 \right) = -\gamma_\mu = -\left( \gamma_\mu\gamma_5 \right) \gamma_5$, and $\gamma_5 S_{\mathrm{Rf}} = (-1)^{\#|\mathrm{Rf}|} S_{\mathrm{Rf}} \gamma_5$.
By (I.c2), $\gamma_5\Sigma_3 = -\Sigma_3\gamma_5$.
Consequently, if $\#|\mathrm{Rf}|$ is even, then $S_{\mathrm{Rf}}$ and $-S_{\mathrm{Rf}}$ commute with $\gamma_5$ and the two others anticommute with it;
if $\#|\mathrm{Rf}|$ is odd, then $i\Sigma_3 S_{\mathrm{Rf}}$ and $-i\Sigma_3 S_{\mathrm{Rf}}$ commute with $\gamma_5$ and the two others anticommute with it.
In both cases exactly two of the four lie in $\mathrm{Sp}_3 = \left\{ S \in \mathrm{Sp} : S\gamma_5 = \gamma_5 S \right\}$.

\emph{The element $\tilde{g}$.}
By Proposition~35.1\,(d), $\ell_j \in \mathrm{Cl}$ for $0 \le j \le k_{*}$;
so $\ell_j$ is a multiple of $LM$, and the groups $\mathbb{G}_j$, $\mathbb{G}_j^{L}$ of the side $\ell_j$ are defined (Section~I.4.2).
The pair $\tilde{g} = \left( S, L^{-j} a_{\mathfrak{c}} \right)$ has $S \in \mathrm{Sp}$ and its translation part in $LM\,\mathbb{Z}^3/\ell_j\mathbb{Z}^3$;
so it is an element of $\mathbb{G}_j^{L}$.

\emph{Actions} (Section~I.4.3, with $r(S) = r_{\mathfrak{c}}$).
On photon sites, $\tilde{g} c = r_{\mathfrak{c}} c + L^{-j} a_{\mathfrak{c}}$;
this is $\mathsf{g}_{\mathfrak{c}}$ in the coordinates of level $j$.
The images of cells, sets and polymers are formed with this map of the sites, so they are the images under $\mathsf{g}_{\mathfrak{c}}$.
On fermion sites, by (I.4.2), $\tilde{g}(c, g) = \left( r_{\mathfrak{c}} c + L^{-j} a_{\mathfrak{c}}, r_{\mathfrak{c}} g \right)$:
the cluster $\{c\} \times O_{\mathrm{cub}}$ ($48$ sites) is mapped onto the cluster of $\tilde{g} c$, and the label is transformed by $g \mapsto r_{\mathfrak{c}} g$, the left multiplication by the element $r_{\mathfrak{c}}$ of $\{\pm1\}^3 \subset O_{\mathrm{cub}}$;
this map is a bijection of $O_{\mathrm{cub}}$.
On bond functions, $(\tilde{g} A)(c, c') = A\left( \tilde{g}^{-1} c, \tilde{g}^{-1} c' \right)$ (Section~I.4.3) and $\left( (\mathsf{g}_{\mathfrak{c}})_* A \right)(c, c') = A\left( \mathsf{g}_{\mathfrak{c}}^{-1} c, \mathsf{g}_{\mathfrak{c}}^{-1} c' \right)$ (Section~G5.1.2);
the two maps of the sites coincide, and therefore $\tilde{g} A = (\mathsf{g}_{\mathfrak{c}})_* A$.

\emph{Level $0$.}
At $j = 0$ the coordinates are the lattice units, $T_0 = \mathbb{Z}^3/n_*\mathbb{Z}^3$, $T^{\mathrm{f}}_0 = T_0 \times O_{\mathrm{cub}}$, $\tilde{g} = \left( S, a_{\mathfrak{c}} \right)$, and $a_{\mathfrak{c}}$ lies in $\mathrm{per}\,\mathbb{Z}^3 \subset LM\,\mathbb{Z}^3$.
Nothing else changes.
\end{proof}

\subsection*{3. Proposition 35.3 (probes)}

\subsubsection*{Proof of Proposition 35.3}

\paragraph{3.0. Preliminaries}

\emph{(ps1) Translates and reflections of a probe.}
The functions $\mathrm{gs}_\nu[\beta']$ and the numbers $\left| D^i \mathrm{fs} \right|_\infty$ are those of Section~G5.1.3.
One has $(\tau_v\beta')(y) = \beta'(y - v)$ and, for $r = \mathrm{diag}(r_1, r_2, r_3) \in \{\pm1\}^3$, $(U_r\beta')_{\mu\nu}(y) = r_\mu r_\nu \beta'_{\mu\nu}\left( r^{-1} y \right)$;
this is the rule $(U_r f)(y) = \left( \wedge^2 r \right) f\left( r^{-1} y \right)$ of Section~I.7.5.
So $\tau_v\beta'$ and $U_r\beta'$ are probes, and
\begin{equation}
\begin{gathered}
\left| \tau_v\beta' \right|_{C^i} = \left| U_r\beta' \right|_{C^i} = \left| \beta' \right|_{C^i},\\
\mathrm{supp}\,\tau_v\beta' = v + \mathrm{supp}\,\beta', \qquad \mathrm{supp}\,U_r\beta' = r\,\mathrm{supp}\,\beta'
\end{gathered}
\tag{Eq3.1}
\end{equation}
(a derivative of $U_r\beta'$ is $\pm$ a derivative of the same order of $\beta'$, taken at the point $r^{-1} y$).
The cube $[-2, 2]^3$ is invariant under $r$, and the maps $y \mapsto v + y$ and $y \mapsto r y$, as well as their inverses, carry closed coordinate cubes to closed coordinate cubes of the same side.
Let $Q_{\beta'}$ be a closed coordinate cube of the side $D_{\beta'}$ that contains $\mathrm{supp}\,\beta' + [-2, 2]^3$ (Section~I.7.6).
Then
\begin{equation}
\begin{gathered}
\mathrm{supp}\,\tau_v\beta' + [-2, 2]^3 \subset v + Q_{\beta'}, \qquad \mathrm{supp}\,U_r\beta' + [-2, 2]^3 = r \left( \mathrm{supp}\,\beta' + [-2, 2]^3 \right) \subset r Q_{\beta'},\\
D_{\tau_v\beta'} = D_{U_r\beta'} = D_{\beta'} .
\end{gathered}
\tag{Eq3.1$'$}
\end{equation}

\emph{(ps2) The operators $\delta_{\eta,\mu}$ and $J_\eta$.}
For a smooth function $\mathrm{fs}$ on $\mathbb{R}^3$ with compact support and $0 < \eta \le 1$ one has
\begin{gather*}
(\delta_{\eta,\mu}\mathrm{fs})(y) = \int_{-\frac{1}{2}}^{\frac{1}{2}} \mathrm{d}\,\mathrm{sv}\, (\partial_\mu\mathrm{fs})\left( y + \mathrm{sv}\,\eta e_\mu \right),\\
(J_\eta\mathrm{fs})(y) = \int_{[-\frac{1}{2}, \frac{1}{2}]^3} \mathrm{d}y'\, \mathrm{fs}\left( y + \eta y' \right).
\end{gather*}
Both are means, with total weight one, of translates of $\partial_\mu\mathrm{fs}$, respectively of $\mathrm{fs}$.
Hence:
\begin{itemize}
\item[(i)] they map $C_c^\infty(\mathbb{R}^3)$ into itself and commute with translations and with partial derivatives;
\item[(ii)] for every multi-index $\mathrm{ak}$ one has $\left| \partial^{\mathrm{ak}} \delta_{\eta,\mu} J_\eta \mathrm{fs} \right|_\infty \le \left| \partial^{\mathrm{ak}} \partial_\mu \mathrm{fs} \right|_\infty$;
\item[(iii)] if $(J_\eta\mathrm{fs})(y) \ne 0$, then $y + [-\eta/2, \eta/2]^3$ meets $\mathrm{supp}\,\mathrm{fs}$;
if $(\delta_{\eta,\mu}\mathrm{fs})(y) \ne 0$, then one of the points $y \pm (\eta/2) e_\mu$ lies in $\mathrm{supp}\,\mathrm{fs}$;
consequently $(\delta_{\eta,\mu} J_\eta \mathrm{fs})(y) \ne 0$ only if $y$ has sup distance at most $\eta/2 + \eta/2 = \eta$ from $\mathrm{supp}\,\mathrm{fs}$.
\end{itemize}
By (i),
\begin{equation}
\mathrm{gs}_\nu[\tau_v\beta'](y) = \mathrm{gs}_\nu[\beta'](y - v) \qquad \text{for all } v \in \mathbb{R}^3 . \tag{Eq3.2}
\end{equation}

\emph{(ps3) The three norms as suprema; separated supports.}
Let $0 < \eta'' \le 1$, and let the lattice be $\eta''\mathbb{Z}^3$ or a torus $\eta''\mathbb{Z}^3/\ell\mathbb{Z}^3$ with $\ell/\eta''$ an integer and $\ell > 6$.
For a bond function $A$ let $\mathrm{Edp}(A)$ be the set of the end points of the bonds $b$ with $A(b) \ne 0$.
By Section~G5.1.3 the three norms of $A$ are the suprema of the moduli of the following linear functionals $\mathrm{fn}$ of $A$, where $\mu, \nu \in \{1, 2, 3\}$ and $\mu = \nu$ is admitted
(Section~G5.1.3 writes $\partial_\nu A_\mu$; here the roles of the two letters are exchanged):
\begin{align*}
&\text{kind 1:} && A_\nu(c);\\
&\text{kind 2:} && (\partial_\mu A_\nu)(c) = (\eta'')^{-1} \left( A_\nu\left( c + \eta'' e_\mu \right) - A_\nu(c) \right);\\
&\text{kind 3:} && |c - c'|^{-\alpha} \left( (\partial_\mu A_\nu)(c') - (\partial_\mu A_\nu)(c) \right), \qquad 0 < |c - c'| \le 1 .
\end{align*}
A functional of kind 1, 2, 3 reads $A$ on one, two, four bonds, respectively:
on the bond $\left( c, c + \eta'' e_\nu \right)$;
on this bond and its translate by $\eta'' e_\mu$;
on these two and on the two bonds formed in the same way with $c'$.
The end points of the bonds read by one functional lie in $c + \left\{ 0, \eta'' e_\nu, \eta'' e_\mu, \eta'' (e_\mu + e_\nu) \right\}$ and, for kind 3, also in $c'$ plus the same set;
so any two of them have sup distance at most $1 + 2\eta''$ (the number $2\eta''$ occurs for $\mu = \nu$).

\begin{named}{Lemma P}
Let $\left( A^{\mathrm{ia}} \right)_{\mathrm{ia} \in \mathrm{Ia}}$ be a family of bond functions such that, for $\mathrm{ia} \ne \mathrm{ia}'$, every point of $\mathrm{Edp}\left( A^{\mathrm{ia}} \right)$ has sup distance $> 1 + 2\eta''$ from every point of $\mathrm{Edp}\left( A^{\mathrm{ia}'} \right)$.
Then every bond lies in the support of at most one $A^{\mathrm{ia}}$, so that $\sum_{\mathrm{ia}} t_{\mathrm{ia}} A^{\mathrm{ia}}$ is defined for all complex numbers $t_{\mathrm{ia}}$, and each of the three norms $\mathrm{nrm}(\cdot)$ satisfies
\[
\mathrm{nrm}\left( \sum_{\mathrm{ia}} t_{\mathrm{ia}} A^{\mathrm{ia}} \right) = \sup_{\mathrm{ia}} \left| t_{\mathrm{ia}} \right| \cdot \mathrm{nrm}\left( A^{\mathrm{ia}} \right).
\]
\end{named}

\begin{proof}
A bond in the supports of $A^{\mathrm{ia}}$ and of $A^{\mathrm{ia}'}$, $\mathrm{ia} \ne \mathrm{ia}'$, would have its end points both in $\mathrm{Edp}\left( A^{\mathrm{ia}} \right)$ and in $\mathrm{Edp}\left( A^{\mathrm{ia}'} \right)$;
so the sum is defined.
Write $\sum t A := \sum_{\mathrm{ia}} t_{\mathrm{ia}} A^{\mathrm{ia}}$, and let $\mathrm{fn}$ be one of the functionals.
If $A^{\mathrm{ia}}$ and $A^{\mathrm{ia}'}$, $\mathrm{ia} \ne \mathrm{ia}'$, did not vanish on bonds $b$, $b'$ read by $\mathrm{fn}$, then an end point of $b$ and an end point of $b'$ would have distance at most $1 + 2\eta''$.
So there is at most one index $\mathrm{ia}(\mathrm{fn})$ whose member does not vanish on all bonds read by $\mathrm{fn}$.
If there is one, then $\mathrm{fn}\left( \sum t A \right) = t_{\mathrm{ia}(\mathrm{fn})} \cdot \mathrm{fn}\left( A^{\mathrm{ia}(\mathrm{fn})} \right)$;
if there is none, then $\mathrm{fn}\left( \sum t A \right) = 0$.
This gives $\left| \mathrm{fn}\left( \sum t A \right) \right| \le \sup_{\mathrm{ia}} \left| t_{\mathrm{ia}} \right| \cdot \mathrm{nrm}\left( A^{\mathrm{ia}} \right)$.
Conversely, let $\mathrm{ia} \in \mathrm{Ia}$ and $\mathrm{fn}\left( A^{\mathrm{ia}} \right) \ne 0$.
Then $\mathrm{ia}(\mathrm{fn}) = \mathrm{ia}$, and $\left| t_{\mathrm{ia}} \right| \cdot \left| \mathrm{fn}\left( A^{\mathrm{ia}} \right) \right| = \left| \mathrm{fn}\left( \sum t A \right) \right| \le \mathrm{nrm}\left( \sum t A \right)$.
The supremum over the functionals $\mathrm{fn}$ of the kind in question and over $\mathrm{ia}$ gives the assertion.
\end{proof}

\emph{(ps4) Torus and covering lattice.}
Let $A$ be a bond function on the torus $\eta''\mathbb{Z}^3/\ell\mathbb{Z}^3$, $\ell > 2$, and let $\tilde{A}$ be its $\ell\mathbb{Z}^3$-periodic extension to $\eta''\mathbb{Z}^3$.
The three norms of $A$ and of $\tilde{A}$ coincide.
Indeed, a functional of kind 1 or 2 of $\tilde{A}$ at the site $\tilde{c}$ is the functional of $A$ at the image of $\tilde{c}$;
two sites $c$, $c'$ of the torus with $0 < d(c, c') \le 1$ have lifts with $\left| \tilde{c} - \tilde{c}' \right| = d(c, c')$;
and two sites $\tilde{c}$, $\tilde{c}'$ of $\eta''\mathbb{Z}^3$ with $0 < \left| \tilde{c} - \tilde{c}' \right| \le 1 < \ell/2$ have images at the distance $\left| \tilde{c} - \tilde{c}' \right|$.

\emph{(ps5) Lattice isometries with diagonal linear part.}
Let $g_{\mathrm{w}} c = r c + a$ with $r = \mathrm{diag}(r_1, r_2, r_3) \in \{\pm1\}^3$ and $a \in \eta''\mathbb{Z}^3$, on $\eta''\mathbb{Z}^3$ or on a torus $\eta''\mathbb{Z}^3/\ell\mathbb{Z}^3$, and let $\left( (g_{\mathrm{w}})_* A \right)(c, c') = A\left( g_{\mathrm{w}}^{-1} c, g_{\mathrm{w}}^{-1} c' \right)$, where $g_{\mathrm{w}}^{-1} c = r (c - a)$.
Put $\mathrm{sh}_\lambda := 0$ if $r_\lambda = 1$, and $\mathrm{sh}_\lambda := -\eta'' e_\lambda$ if $r_\lambda = -1$.
Then, for all $\mu$, $\nu$ ($\mu = \nu$ included) and all sites $c$,
\begin{equation}
\begin{aligned}
\left( (g_{\mathrm{w}})_* A \right)_\nu(c) &= r_\nu A_\nu\left( g_{\mathrm{w}}^{-1} c + \mathrm{sh}_\nu \right),\\
\left( \partial_\mu \left( (g_{\mathrm{w}})_* A \right)_\nu \right)(c) &= r_\mu r_\nu (\partial_\mu A_\nu)\left( g_{\mathrm{w}}^{-1} c + \mathrm{sh}_\mu + \mathrm{sh}_\nu \right).
\end{aligned}
\tag{Eq3.3}
\end{equation}

\begin{proof}
One has $\left( (g_{\mathrm{w}})_* A \right)_\nu(c) = A\left( g_{\mathrm{w}}^{-1} c, g_{\mathrm{w}}^{-1} c + r_\nu \eta'' e_\nu \right)$.
For $r_\nu = 1$ this is $A_\nu\left( g_{\mathrm{w}}^{-1} c \right)$.
For $r_\nu = -1$ it is
\[
A\left( g_{\mathrm{w}}^{-1} c, g_{\mathrm{w}}^{-1} c - \eta'' e_\nu \right) = -A\left( g_{\mathrm{w}}^{-1} c - \eta'' e_\nu, g_{\mathrm{w}}^{-1} c \right) = -A_\nu\left( g_{\mathrm{w}}^{-1} c - \eta'' e_\nu \right).
\]
This is the first formula.
Let $\tilde{c} := g_{\mathrm{w}}^{-1} c + \mathrm{sh}_\nu$.
Since $g_{\mathrm{w}}^{-1}\left( c + \eta'' e_\mu \right) = g_{\mathrm{w}}^{-1} c + r_\mu \eta'' e_\mu$, the first formula gives
\[
\left( \partial_\mu \left( (g_{\mathrm{w}})_* A \right)_\nu \right)(c) = (\eta'')^{-1} r_\nu \left( A_\nu\left( \tilde{c} + r_\mu \eta'' e_\mu \right) - A_\nu(\tilde{c}) \right).
\]
For $r_\mu = 1$ this is $r_\nu (\partial_\mu A_\nu)(\tilde{c})$.
For $r_\mu = -1$ it is $-r_\nu (\eta'')^{-1} \left( A_\nu(\tilde{c}) - A_\nu\left( \tilde{c} - \eta'' e_\mu \right) \right) = -r_\nu (\partial_\mu A_\nu)\left( \tilde{c} - \eta'' e_\mu \right)$.
\end{proof}

\emph{Consequences.}
For fixed $(\mu, \nu)$ the map $c \mapsto g_{\mathrm{w}}^{-1} c + \mathrm{sh}_\mu + \mathrm{sh}_\nu$ is a bijection of the sites that preserves sup distances.
Hence:
\begin{itemize}
\item[(i)] the three norms of $(g_{\mathrm{w}})_* A$ are those of $A$;
\item[(ii)] $(g_{\mathrm{w}})_*$ is a signed permutation of the coordinates $A_\nu(c)$, and so it is orthogonal for the counting pairing;
\item[(iii)] $\mathrm{Edp}\left( (g_{\mathrm{w}})_* A \right) = g_{\mathrm{w}}\, \mathrm{Edp}(A)$;
\item[(iv)] at the mesh one, where $(dA)_{\mu\nu} = \partial_\mu A_\nu - \partial_\nu A_\mu$ (Section~I.1.1),
\begin{equation}
\left( d\left( (g_{\mathrm{w}})_* A \right) \right)_{\mu\nu}(c) = r_\mu r_\nu (dA)_{\mu\nu}\left( g_{\mathrm{w}}^{-1} c + \mathrm{sh}_\mu + \mathrm{sh}_\nu \right), \tag{Eq3.4}
\end{equation}
and the plaquette $\left( g_{\mathrm{w}}^{-1} c + \mathrm{sh}_\mu + \mathrm{sh}_\nu; \mu\nu \right)$ is the one whose set of vertices is the image under $g_{\mathrm{w}}^{-1}$ of the set $c + \left\{ 0, e_\mu, e_\nu, e_\mu + e_\nu \right\}$ of the vertices of $(c; \mu\nu)$.
\end{itemize}
So $(g_{\mathrm{w}})_*$ maps $\ker d$ onto $\ker d$, $\left\| d\left( (g_{\mathrm{w}})_* A \right) \right\| = \| dA \|$, and a plaquette $p$ lies in the support of $d\left( (g_{\mathrm{w}})_* A \right)$ if and only if $g_{\mathrm{w}}^{-1} p$ lies in the support of $dA$.

\emph{(ps6) Three systems of coordinates} (setting (ST*)$^{\mathrm{ch}}$).
The site $c^0 \in \mathbb{Z}^3$ (lattice units) has the physical coordinates $\eta c^0$ (Section~G5.1.2) and the coordinates $L^{-k_{*}} c^0$ of level $k_{*}$ (Section~I.6.1, the paragraph on the coordinates of level $k$, at $k = k_{*}$, together with Sections~I.3.8 and~I.6.6; at $j = k_{*}$ this is also the first sentence of the proof of Proposition~35.2\,(b)).
Since $a_{*} = L^{k_{*}} \eta$ (Section~I.6.1), the physical coordinates are $a_{*}$ times the coordinates of level $k_{*}$, and
\begin{gather*}
\frac{\eta}{a_{*}} = \eta_*, \qquad \frac{R^{*}(e)}{a_{*}} = \ell_*, \qquad \frac{\ell_c}{a_{*}} = \frac{R^{*}(e)}{2L a_{*}} = \frac{\ell_*}{2L},\\
\frac{\ell_c}{\eta} = \mathrm{hc}, \qquad \frac{R^{*}(e)}{\eta} = n_* = 2L \cdot \mathrm{hc} .
\end{gather*}
For a one-cochain $a$ on $\mathbb{Z}^3/n_*\mathbb{Z}^3$ and a bond $\left( c, c + \eta_k e_\nu \right)$ of $T_k$ one has $a^{[k]}_\nu(c) = L^{\frac{k}{2}} a_\nu\left( L^k c \right)$ (Section~I.6.6).
By Proposition~35.1\,(a), (b), by $M \ge L^2$ (line (C5)) and by $L \ge 5$ (line (C3)),
\begin{equation}
\begin{gathered}
\ell_* = M r_{*} \ge L^4, \qquad \eta_* \le \frac{1}{L}, \qquad 1 + 2\eta_* < 3,\\
\frac{\ell_*}{8L} \ge \frac{L^3}{8} > 15, \qquad \mathrm{hc} = \frac{L^{k_{*} - 1} \ell_*}{2} \ge \frac{L^4}{2} .
\end{gathered}
\tag{Eq3.5}
\end{equation}

\paragraph{3.1. The clauses}

\begin{proof}[Proof of Proposition~35.3\,(a)]
\emph{The formula for $d_\eta^{*}u$.}
Let $n=R/\eta$, and let $u$ be a real two-cochain on $\mathbb{Z}^3/n\mathbb{Z}^3$ with the components $u_{\mu\nu}$, $\mu<\nu$; for $\mu<\nu$ set $u_{\nu\mu}:=-u_{\mu\nu}$.
For the cochain $u$ of Section~G5.1.3 this convention agrees with the defining formula, because $\beta'_{\nu\mu}=-\beta'_{\mu\nu}$ and the point $c+\frac{\eta}{2}(e_\mu+e_\nu)$ is symmetric in $\mu$ and $\nu$.
For a one-cochain $A$ one has, by Section~I.1.1,
\[
\langle u,dA\rangle=\sum_{c}\sum_{\mu<\nu}u_{\mu\nu}(c)\cdot\left[A_\mu(c)+A_\nu(c+e_\mu)-A_\mu(c+e_\nu)-A_\nu(c)\right].
\]
In this sum the coefficient of $A_\nu(c)$ collects the terms
\[
u_{\mu\nu}(c-e_\mu)-u_{\mu\nu}(c)
\]
coming from the pairs $(\mu,\nu)$ with $\mu<\nu$, and the terms
\[
u_{\nu\mathfrak{l}}(c)-u_{\nu\mathfrak{l}}(c-e_{\mathfrak{l}})=-u_{\mathfrak{l}\nu}(c)+u_{\mathfrak{l}\nu}(c-e_{\mathfrak{l}})
\]
coming from the pairs $(\nu,\mathfrak{l})$ with $\nu<\mathfrak{l}$.
Consequently the adjoint $d^{*}$ of $d$ with respect to the counting pairings is given by
\[
(d^{*}u)_\nu(c)=-\sum_{\mu\ne\nu}\left(u_{\mu\nu}(c)-u_{\mu\nu}(c-e_\mu)\right).
\]
One has $d_\eta=\eta^{-1}d$, and $\langle\,,\rangle_\eta=\eta^3\langle\,,\rangle$ in both degrees (Sections~G5.1.3 and~I.7.5); hence $d_\eta^{*}=\eta^{-1}d^{*}$, and this is the formula displayed in Section~G5.1.3.

\emph{The assertion.}
Take $n=n_*$ and form $u$ with $\tau_v\beta$.
Then
\[
\hat{\mathfrak{a}}=\eta^{\frac{1}{2}}\cdot a_\eta=d^{*}\left(\eta^{-\frac{1}{2}}u\right)
\]
is a real cochain in the range of $d^{*}$.
In the Euclidean space $C^1(n_*)$ the range of $d^{*}$ is the orthogonal complement of $\ker d$.
By the Hodge decomposition on the discrete torus \cite[Lemma~15.1]{QEDI}, a closed bond function is the sum of a constant one and an exact one; therefore $\ker d={V_{\mathrm{g}}}+{V_{\mathrm{h}}}$, and $(\ker d)^{\perp}={V_{\mathrm{t}}}$ (Section~I.7.1).
By Proposition~35.2\,(a), $\mathsf{g}_{\mathfrak{c}}$ is a map of the type considered in (ps5), on $\mathbb{Z}^3/n_*\mathbb{Z}^3$.
By (ps5)\,(ii) and (iv) the map $(\mathsf{g}_{\mathfrak{c}})_*$ is orthogonal and maps $\ker d$ onto itself; so it also maps $(\ker d)^{\perp}$ onto itself.
It follows that $\hat{\mathfrak{a}}_{\mathfrak{c}}=(\mathsf{g}_{\mathfrak{c}})_*\hat{\mathfrak{a}}$ lies in ${V_{\mathrm{t}}}$.
The argument is valid for every $v\in\mathbb{R}^3$.
\end{proof}

\begin{proof}[Proof of Proposition~35.3\,(b1)]
Let $c\in\Lambda_{\eta,R}$, and let $\mathrm{bm}:=c+\frac{\eta}{2}e_\nu$ be the midpoint of the bond $(c,c+\eta e_\nu)$.
For $\mu\ne\nu$ the points $c+\frac{\eta}{2}(e_\mu+e_\nu)$ and $c-\eta e_\mu+\frac{\eta}{2}(e_\mu+e_\nu)$ are the points $\mathrm{bm}+\frac{\eta}{2}e_\mu$ and $\mathrm{bm}-\frac{\eta}{2}e_\mu$.
Hence
\begin{align*}
u_{\mu\nu}(c)-u_{\mu\nu}(c-\eta e_\mu)
&=\sum_{z}\left[(J_\eta\beta'_{\mu\nu})\left(\mathrm{bm}+\frac{\eta}{2}e_\mu+Rz\right)-(J_\eta\beta'_{\mu\nu})\left(\mathrm{bm}-\frac{\eta}{2}e_\mu+Rz\right)\right]\\
&=\eta\sum_{z}\left(\delta_{\eta,\mu}J_\eta\beta'_{\mu\nu}\right)(\mathrm{bm}+Rz),
\end{align*}
and the formula for $a_\eta[\beta']$ yields
\[
a_\eta[\beta']_\nu(c)=-\sum_{\mu\ne\nu}\sum_{z}\left(\delta_{\eta,\mu}J_\eta\beta'_{\mu\nu}\right)(\mathrm{bm}+Rz)=\sum_{z}\mathrm{gs}_\nu[\beta'](\mathrm{bm}+Rz).
\]
These sums are finite, the functions involved having compact support ((ps2)\,(i)); the right-hand side does not change when $c$ is replaced by $c+Rz'$.

In the setting $(\mathrm{ST}^{*})^{\mathrm{ch}}$ one has $\eta=L^{-N}\le1$, and $R^{*}(e)/\eta=n_*=h\cdot L^{T^{*}(e)+N}$ is an integer (Section~I.6.4) with $n_*=2L\cdot\mathrm{hc}\ge2$ ((ps6), (Eq3.5)).
At $R=R^{*}(e)$ and $\beta'=\tau_v\beta$ the formula says that $a_{\eta,\nu}(c)$ is the value at $c+\frac{\eta}{2}e_\nu$ of the $R^{*}(e)\mathbb{Z}^3$-periodisation of $\mathrm{gs}_\nu[\tau_v\beta]$; by the convention stated at the beginning of Proposition~35.3 this value is written $\mathrm{gs}_\nu[\tau_v\beta]\left(c+\frac{\eta}{2}e_\nu\right)$.

\emph{Bounds.}
Let $\mathrm{ak}$ be a multi-index of order $i$.
By (ps2)\,(ii),
\[
\left|\partial^{\mathrm{ak}}\delta_{\eta,\mu}J_\eta\beta'_{\mu\nu}\right|_\infty\le\left|\partial^{\mathrm{ak}}\partial_\mu\beta'_{\mu\nu}\right|_\infty\le|\beta'|_{C^{i+1}},
\]
and $\mathrm{gs}_\nu[\beta']$ is a sum of two terms of this kind, one for each of the two values $\mu\ne\nu$.
Therefore $|D^i\mathrm{gs}_\nu[\beta']|_\infty\le2|\beta'|_{C^{i+1}}$.
If $\mathrm{gs}_\nu[\beta'](y)\ne0$, then one of the two terms does not vanish at $y$, and (ps2)\,(iii) shows that $y$ has sup distance at most $\eta$ from $\operatorname{supp}\beta'_{\mu\nu}$, which is a subset of $\operatorname{supp}\beta'$.
\end{proof}

\begin{proof}[Proof of Proposition~35.3\,(b2)]
\emph{A fact valid for every probe $\beta'$, every $\eta\in(0,1]$ and every point $\mathrm{bm}$ of $\mathbb{R}^3$:}
\begin{equation}\tag{Eq3.6}
\mathrm{gs}_\nu[\beta'](\mathrm{bm})\ne0\quad\Longrightarrow\quad\text{the two points }\mathrm{bm}\pm\frac{\eta}{2}e_\nu\text{ lie in }\operatorname{supp}\beta'+[-2,2]^3.
\end{equation}
Indeed, by (b1) the point $\mathrm{bm}$ has sup distance at most $\eta$ from $\operatorname{supp}\beta'$, so that the two points have sup distance at most $\frac{3\eta}{2}\le2$ from the compact set $\operatorname{supp}\beta'$.

Now let $(c,c+\eta e_\nu)$ be a bond of $\Lambda_{\eta,R^{*}(e)}$ on which $\hat{\mathfrak{a}}$ does not vanish, and let $\mathrm{bm}=c+\frac{\eta}{2}e_\nu$.
By (b1) there is $z\in\mathbb{Z}^3$ with $\mathrm{gs}_\nu[\tau_v\beta](\mathrm{bm}+R^{*}(e)z)\ne0$.
By (Eq3.6) and (Eq3.1') the points $c+R^{*}(e)z$ and $c+\eta e_\nu+R^{*}(e)z$ lie in $\operatorname{supp}\tau_v\beta+[-2,2]^3$, a subset of $v+Q_\beta=\mathrm{Qb}$.
These two points are lifts of the two end points of the bond.
(The cube $\mathrm{Qb}$ is contained in $[\ell_c/8,7\ell_c/8]^3$, hence in the period cube $[0,R^{*}(e))^3$.)
\end{proof}

\begin{proof}[Proof of Proposition~35.3\,(c1)]
Let $(N',\mathrm{asc},r)$ be a triple of $\mathrm{Ind}_\beta$, let $\eta'=L^{-N'}\le1$ and $\eta''=\eta'/\mathrm{asc}$, and let $\mathrm{gs}^0_\nu=\mathrm{gs}_\nu[U_r\beta]$ be formed with $\eta'$.
Then $\mathrm{asc}\cdot\eta''=\eta'$ and $\mathrm{asc}\le h/L\le h$, and (b1) together with (Eq3.1) gives
\[
|D^i\mathrm{gs}^0_\nu|_\infty\le2|\beta|_{C^{i+1}}\qquad(i=0,1,2).
\]

\emph{The first norm.}
For $c\in\eta''\mathbb{Z}^3$,
\[
|\mathrm{gsc}_\nu(c)|=\mathrm{asc}^{\frac{1}{2}}\cdot\left|\mathrm{gs}^0_\nu\left(\mathrm{asc}\cdot c+\frac{\eta'}{2}e_\nu\right)\right|\le2\,\mathrm{asc}^{\frac{1}{2}}\cdot|\beta|_{C^1}.
\]

\emph{The second norm.}
For $c\in\eta''\mathbb{Z}^3$ and all $\mu,\nu$ one has, since ${\eta''}^{-1}\eta'=\mathrm{asc}$,
\[
(\partial_\mu\mathrm{gsc}_\nu)(c)={\eta''}^{-1}\mathrm{asc}^{\frac{1}{2}}\cdot\left[\mathrm{gs}^0_\nu\left(\mathrm{asc}\cdot c+\eta'e_\mu+\frac{\eta'}{2}e_\nu\right)-\mathrm{gs}^0_\nu\left(\mathrm{asc}\cdot c+\frac{\eta'}{2}e_\nu\right)\right]=\mathrm{asc}^{\frac{3}{2}}\cdot\mathrm{Gd}(\mathrm{asc}\cdot c),
\]
where
\[
\mathrm{Gd}(y):=\int_0^1\mathrm{d}\mathrm{sv}\,(\partial_\mu\mathrm{gs}^0_\nu)\left(y+\mathrm{sv}\,\eta'e_\mu+\frac{\eta'}{2}e_\nu\right).
\]
As $|\mathrm{Gd}|_\infty\le|D\mathrm{gs}^0_\nu|_\infty\le2|\beta|_{C^2}$, it follows that $|\partial\mathrm{gsc}|_\infty\le2\,\mathrm{asc}^{\frac{3}{2}}\cdot|\beta|_{C^2}$.

\emph{The third norm.}
For $0<|c-c'|\le1$,
\[
\left|\mathrm{Gd}(\mathrm{asc}\cdot c')-\mathrm{Gd}(\mathrm{asc}\cdot c)\right|\le\left|\mathrm{asc}\cdot(c'-c)\right|_1\cdot|D^2\mathrm{gs}^0_\nu|_\infty\le3\,\mathrm{asc}\cdot|c-c'|\cdot2|\beta|_{C^3}\le6\,\mathrm{asc}\cdot|c-c'|^{\alpha}\cdot|\beta|_{C^3},
\]
because $|c-c'|\le1$ and $\alpha<1$ (line (C1) of Section~I.8.2).
Hence $|\delta_\alpha\partial\mathrm{gsc}|_\infty\le6\,\mathrm{asc}^{\frac{5}{2}}\cdot|\beta|_{C^3}$.

The sum of the three norms of $\mathrm{gsc}$ is therefore at most
\[
2\,\mathrm{asc}^{\frac{1}{2}}\cdot\left(1+\mathrm{asc}+3\,\mathrm{asc}^2\right)\cdot|\beta|_{C^3}\le2h^{\frac{1}{2}}\cdot\left(1+h+3h^2\right)\cdot|\beta|_{C^3}.
\]
The triple does not occur in this bound; so the bound holds for the supremum $C_\beta$.

\emph{The number $\mathbf{D}$.}
For $c\in\mathbb{Z}^3$ and $\mu<\nu$ one has (Section~I.1.1)
\[
(d\,\mathrm{w0})_{\mu\nu}(c)=\left[\mathrm{w0}_\nu(c+e_\mu)-\mathrm{w0}_\nu(c)\right]-\left[\mathrm{w0}_\mu(c+e_\nu)-\mathrm{w0}_\mu(c)\right].
\]
The first bracket equals
\[
{\eta'}^{\frac{1}{2}}\cdot\left[\mathrm{gs}^0_\nu\left(\eta'c+\eta'e_\mu+\frac{\eta'}{2}e_\nu\right)-\mathrm{gs}^0_\nu\left(\eta'c+\frac{\eta'}{2}e_\nu\right)\right],
\]
and its modulus is at most ${\eta'}^{\frac{1}{2}}\eta'\cdot|D\mathrm{gs}^0_\nu|_\infty\le2{\eta'}^{\frac{3}{2}}\cdot|\beta|_{C^2}$.
The second bracket equals
\[
{\eta'}^{\frac{1}{2}}\cdot\left[\mathrm{gs}^0_\mu\left(\eta'c+\eta'e_\nu+\frac{\eta'}{2}e_\mu\right)-\mathrm{gs}^0_\mu\left(\eta'c+\frac{\eta'}{2}e_\mu\right)\right],
\]
and its modulus is at most ${\eta'}^{\frac{1}{2}}\eta'\cdot|D\mathrm{gs}^0_\mu|_\infty\le2{\eta'}^{\frac{3}{2}}\cdot|\beta|_{C^2}$.
Thus $|(d\,\mathrm{w0})(p)|\le4{\eta'}^{\frac{3}{2}}\cdot|\beta|_{C^2}$ for every plaquette $p$.

Let $p=(c;\mu\nu)$ be a plaquette with $(d\,\mathrm{w0})(p)\ne0$.
Then $\mathrm{w0}$ does not vanish on one of the four bonds of $p$.
By (Eq3.6) for $\beta'=U_r\beta$ and the mesh $\eta'$, together with (Eq3.1'), the end points of this bond, multiplied by $\eta'$, lie in $rQ_\beta$, a closed coordinate cube $\prod_{\mathfrak{l}}[\mathrm{lo}_{\mathfrak{l}},\mathrm{lo}_{\mathfrak{l}}+D_\beta]$.
These end points are vertices of $p$, that is, points $c+\epsilon e_\mu+\epsilon'e_\nu$ with $\epsilon,\epsilon'\in\{0,1\}$.
Hence $\eta'c_{\mathfrak{l}}\in[\mathrm{lo}_{\mathfrak{l}}-\eta',\mathrm{lo}_{\mathfrak{l}}+D_\beta]$ for $\mathfrak{l}=1,2,3$.
A closed interval of length $D_\beta+\eta'$ contains at most $(D_\beta+\eta')/\eta'+1=D_\beta/\eta'+2$ points of $\eta'\mathbb{Z}$.
There are therefore at most $(D_\beta/\eta'+2)^3$ such $c$, and there are three pairs $\mu<\nu$; it follows that
\[
\|d\,\mathrm{w0}\|^2\le3\left(\frac{D_\beta}{\eta'}+2\right)^3\cdot16{\eta'}^3\cdot|\beta|_{C^2}^2=48|\beta|_{C^2}^2\cdot\left(D_\beta+2\eta'\right)^3\le48|\beta|_{C^2}^2\cdot\left(D_\beta+3\right)^3.
\]
The triple does not occur in this bound; so the bound holds for $\mathbf{D}$.

\emph{Dependence.}
The set $\mathrm{Ind}_\beta$ is determined by $(h,L)$.
For a triple of $\mathrm{Ind}_\beta$ the functions $\mathrm{gsc}$, $\mathrm{w0}$ and the mesh $\eta''$ are determined by $\beta$ and the triple, and the third norm involves $\alpha$.
No other quantity enters the two suprema.
(The number $h$ is a function of $(\bar{m},L)$: Section~I.0.5.)
\end{proof}

\begin{proof}[Proof of Proposition~35.3\,(c2)]
\emph{The triple.}
One has $N\ge0$; the number $a_*$ is one of $h/L$, $h/L^2$ (Proposition~35.1\,(a)); $N+\log_La_*=k_*\ge1$ (Proposition~35.1\,(a), (c)); and $1\in\{\pm1\}^3$.
Hence $(N,a_*,1)$ belongs to $\mathrm{Ind}_\beta$.
At this triple $\eta'=\eta$, $\eta''=\eta/a_*=\eta_*$ and $\mathrm{gs}^0_\nu=\mathrm{gs}_\nu[U_1\beta]=\mathrm{gs}_\nu[\beta]$.

\emph{Two formulas.}
The vector $v$ lies in $\eta\mathbb{Z}^3$; so $v/\eta\in\mathbb{Z}^3$ and $v/a_*\in(\eta/a_*)\mathbb{Z}^3=\eta_*\mathbb{Z}^3$.
By (b1), (Eq3.2) and (ps6) one obtains, for $c\in\mathbb{Z}^3$ and for $c\in\eta_*\mathbb{Z}^3$ respectively (periodic extensions),
\begin{align}
\hat{\mathfrak{a}}_\nu(c)&=\eta^{\frac{1}{2}}\sum_{z}\mathrm{gs}_\nu[\beta]\left(\eta c+\frac{\eta}{2}e_\nu-v+R^{*}(e)z\right)=\sum_{z}\mathrm{w0}_\nu\left(c-\frac{v}{\eta}+n_*z\right),\tag{Eq3.7}\\
\hat{\mathfrak{a}}^{[k_*]}_\nu(c)&=L^{\frac{k_*}{2}}\cdot\hat{\mathfrak{a}}_\nu(L^{k_*}c)=a_*^{\frac{1}{2}}\sum_{z}\mathrm{gs}_\nu[\beta]\left(a_*c+\frac{\eta}{2}e_\nu-v+a_*\ell_*z\right)\notag\\
&=\sum_{z}\mathrm{gsc}_\nu\left(c-\frac{v}{a_*}+\ell_*z\right).\tag{Eq3.8}
\end{align}

\emph{The norms.}
For $z\in\mathbb{Z}^3$ let $A^z(c):=\mathrm{gsc}\left(c-\frac{v}{a_*}+\ell_*z\right)$, a bond function on $\eta_*\mathbb{Z}^3$.
If $A^0_\nu(c)\ne0$, then $\mathrm{gs}_\nu[\tau_v\beta]\left(a_*c+\frac{\eta}{2}e_\nu\right)\ne0$, and by (Eq3.6) and (Eq3.1') the points $a_*c$ and $a_*c+\eta e_\nu$ lie in $\mathrm{Qb}$.
So $\mathrm{Edp}(A^0)$ is contained in $a_*^{-1}\mathrm{Qb}$, a closed cube of side $D_\beta/a_*$.
Since the placement is admitted, $\mathrm{Qb}$ lies in a cube of side $\frac{3}{4}\ell_c$; hence
\[
\frac{D_\beta}{a_*}\le\frac{3\ell_c}{4a_*}=\frac{3\ell_*}{8L}.
\]
Moreover $\mathrm{Edp}(A^z)=\mathrm{Edp}(A^0)-\ell_*z$.
For $z\ne z'$, a point of $\mathrm{Edp}(A^z)$ and a point of $\mathrm{Edp}(A^{z'})$ have sup distance at least
\[
\ell_*-\frac{3\ell_*}{8L}>\frac{\ell_*}{2}>3>1+2\eta_*
\]
by (Eq3.5).
Lemma~P, applied with all $t$ equal to $1$, shows that each of the three norms of $\sum_zA^z$ is the supremum over $z$ of the same norm of $A^z$.
A translation by a vector of $\eta_*\mathbb{Z}^3$ permutes the functionals of each kind; so the norms of $A^z$ coincide with those of $\mathrm{gsc}$.
By (Eq3.8) and (ps4), the three norms of $\hat{\mathfrak{a}}^{[k_*]}$ on $T_*$ are those of $\mathrm{gsc}$ on $\eta_*\mathbb{Z}^3$.
Their sum is $|\hat{\mathfrak{a}}|_{\mathrm{ext}}$ (Section~G5.1.3).

\emph{The number $\boldsymbol{\sigma}^2$.}
For the plaquettes $p$ of $\mathbb{Z}^3$ put $F^0(p):=(d\,\mathrm{w0})(p-v/\eta)$.
The coboundary $d$ commutes with translations and with locally finite sums; so, by (Eq3.7), the periodic extension of $d\hat{\mathfrak{a}}$ is $\sum_zF^0(\cdot+n_*z)$.
If $F^0(p)\ne0$, some bond of $p$ carries a non-zero value of $\mathrm{w0}(\cdot-v/\eta)$, and by (Eq3.6) and (Eq3.1') its end points, multiplied by $\eta$, lie in $\mathrm{Qb}$.
Therefore all vertices of $p$ lie in the cube $\eta^{-1}\mathrm{Qb}$ enlarged by $1$ on all sides, whose side is
\[
\frac{D_\beta}{\eta}+2\le\frac{3}{4}\mathrm{hc}+2<2L\cdot\mathrm{hc}=n_*
\]
by (Eq3.5).
Hence $p$ and $p+n_*z$ with $z\ne0$ never lie both in $\operatorname{supp}F^0$: for every $p$ at most one term of $\sum_zF^0(p+n_*z)$ differs from $0$.
Letting $p$ range over a set of representatives of the plaquettes modulo $n_*\mathbb{Z}^3$, one finds
\[
\boldsymbol{\sigma}^2=\|d\hat{\mathfrak{a}}\|^2=\sum_{p}\left|\sum_{z}F^0(p+n_*z)\right|^2=\sum_{p}\sum_{z}\left|F^0(p+n_*z)\right|^2=\sum_{p'}\left|F^0(p')\right|^2=\|d\,\mathrm{w0}\|^2,
\]
where $p'$ ranges over all plaquettes of $\mathbb{Z}^3$.

\emph{The inequalities.}
The numbers $C_\beta$ and $\mathbf{D}$ are suprema over $\mathrm{Ind}_\beta$, and the triple $(N,a_*,1)$ belongs to $\mathrm{Ind}_\beta$.
\end{proof}

\paragraph{Separation of the cells}
In physical coordinates $\mathsf{g}_{\mathfrak{c}}$ is the map $y\mapsto r_{\mathfrak{c}}y+\eta\cdot a_{\mathfrak{c}}$, and $\eta\cdot\mathrm{per}=2\ell_c$.
By the formula for $\mathsf{g}_{\mathfrak{c}}$ in Section~I.6.4, that is, by Proposition~35.2\,(a), in lattice units and in the coordinate $\mu$ the map $\mathsf{g}_{\mathfrak{c}}$ is
\begin{itemize}
\item the translation $y_\mu\mapsto i\cdot\mathrm{per}+y_\mu=2i\cdot\mathrm{hc}+y_\mu$ if $\mathfrak{c}_\mu=2i$, which maps $[0,\mathrm{hc}]$ onto $[2i\cdot\mathrm{hc},(2i+1)\mathrm{hc}]$;
\item the reflection $y_\mu\mapsto(i+1)\mathrm{per}-y_\mu=(2i+2)\mathrm{hc}-y_\mu$ if $\mathfrak{c}_\mu=2i+1$, which maps $[0,\mathrm{hc}]$ onto $[(2i+1)\mathrm{hc},(2i+2)\mathrm{hc}]$.
\end{itemize}
As $\ell_c=\mathrm{hc}\cdot\eta$, in physical coordinates and in the coordinate $\mu$ the map sends $[0,\ell_c]$ onto $[\mathfrak{c}_\mu\ell_c,(\mathfrak{c}_\mu+1)\ell_c]$, by a translation or by a reflection.
The subinterval $[\ell_c/8,7\ell_c/8]$ is symmetric about the midpoint of $[0,\ell_c]$; so its image is
\[
\mathrm{Iv}_\mu(\mathfrak{c}):=\left[\mathfrak{c}_\mu\ell_c+\frac{\ell_c}{8},\,(\mathfrak{c}_\mu+1)\ell_c-\frac{\ell_c}{8}\right].
\]
From (ps5)\,(iii) and (b2) one obtains
\begin{equation}\tag{Eq3.9}
\begin{aligned}
&\text{the end points of the bonds in the support of }\hat{\mathfrak{a}}_{\mathfrak{c}}\text{ lie in }\mathsf{g}_{\mathfrak{c}}\mathrm{Qb},\text{ a closed cube of side }D_\beta\\
&\text{contained in }\prod_{\mu}\mathrm{Iv}_\mu(\mathfrak{c})\quad(\text{modulo }R^{*}(e)\mathbb{Z}^3).
\end{aligned}
\end{equation}
Let $\mathfrak{c}\ne\mathfrak{c}'$ in $\mathrm{Cel}=(\mathbb{Z}/2L)^3$, and let $\mu$ be a direction with $\mathfrak{c}_\mu\ne\mathfrak{c}'_\mu$ modulo $2L$.
The circle $\mathbb{R}/R^{*}(e)\mathbb{Z}$ is the union of the $2L$ arcs $[j\ell_c,(j+1)\ell_c]$, $j\in\mathbb{Z}/2L$, whose interiors are pairwise disjoint.
Let $y_\mu\in\mathrm{Iv}_\mu(\mathfrak{c})$ and $y'_\mu\in\mathrm{Iv}_\mu(\mathfrak{c}')$.
Every arc of the circle from $y_\mu$ to $y'_\mu$ contains an end point of the arc of index $\mathfrak{c}_\mu$ and an end point of the arc of index $\mathfrak{c}'_\mu$.
The distance of $y_\mu$ from the first is at least $\ell_c/8$, the distance of $y'_\mu$ from the second is at least $\ell_c/8$, and the two pieces (from $y_\mu$ to the first end point met, and from the last end point met to $y'_\mu$) do not overlap.
Hence the distance of $y_\mu$ and $y'_\mu$ on the circle is at least $\ell_c/4$, and
\begin{equation}\tag{Eq3.10}
\begin{aligned}
&\text{for }\mathfrak{c}\ne\mathfrak{c}'\text{, a point of }\mathsf{g}_{\mathfrak{c}}\mathrm{Qb}\text{ and a point of }\mathsf{g}_{\mathfrak{c}'}\mathrm{Qb}\text{ have on the torus}\\
&\text{the physical sup distance}\ge\frac{\ell_c}{4},\text{ that is, the sup distance}\ge\frac{\ell_c}{4a_*}=\frac{\ell_*}{8L}\text{ in the coordinates of level }k_*.
\end{aligned}
\end{equation}

\begin{proof}[Proof of Proposition~35.3\,(d)]
\emph{Levels.}
Let $a$ be a complex one-cochain on $\mathbb{Z}^3/n_*\mathbb{Z}^3$, let $0\le j\le k_*$, and let $\tilde{g}=(S,L^{-j}a_{\mathfrak{c}})$ be as in Proposition~35.2\,(b).
Since $\tilde{g}$ is $\mathsf{g}_{\mathfrak{c}}$ in the coordinates of level $j$, every bond $b$ of $T_j$ satisfies
\[
\left((\mathsf{g}_{\mathfrak{c}})_*a\right)^{[j]}(b)=L^{\frac{j}{2}}\cdot a\left(\mathsf{g}_{\mathfrak{c}}^{-1}L^jb\right)=L^{\frac{j}{2}}\cdot a\left(L^j\tilde{g}^{-1}b\right)=a^{[j]}\left(\tilde{g}^{-1}b\right)=\left(\tilde{g}\,a^{[j]}\right)(b).
\]
Therefore
\begin{equation}\tag{Eq3.11}
\left((\mathsf{g}_{\mathfrak{c}})_*a\right)^{[j]}=\tilde{g}\left(a^{[j]}\right)=\left(c\mapsto r_{\mathfrak{c}}c+L^{-j}a_{\mathfrak{c}}\right)_*a^{[j]}.
\end{equation}

\emph{Norms.}
By (Eq3.11) at $j=k_*$ and $a=\hat{\mathfrak{a}}$, together with (ps5)\,(i) (mesh $\eta_*$, translation part in $LM\mathbb{Z}^3$, a subset of $\eta_*\mathbb{Z}^3$), the three norms of $\hat{\mathfrak{a}}_{\mathfrak{c}}^{[k_*]}$ are those of $\hat{\mathfrak{a}}^{[k_*]}$.
Consequently
\[
\epsilon_{\mathfrak{c}}=e_*\cdot|\hat{c}_{\mathfrak{c}}|_\infty=\frac{e_*}{e}\cdot\left|\hat{\mathfrak{a}}_{\mathfrak{c}}^{[k_*]}\right|_\infty=\frac{e_*}{e}\cdot\left|\hat{\mathfrak{a}}^{[k_*]}\right|_\infty
\]
does not depend on $\mathfrak{c}$.

\emph{Supports.}
Suppose that a plaquette $p$ of $\mathbb{Z}^3/n_*\mathbb{Z}^3$ lies in the support of $d\hat{\mathfrak{a}}_{\mathfrak{c}}$ and in the support of $d\hat{\mathfrak{a}}_{\mathfrak{c}'}$, where $\mathfrak{c}\ne\mathfrak{c}'$.
Then $p$ has a bond in the support of $\hat{\mathfrak{a}}_{\mathfrak{c}}$ and a bond in the support of $\hat{\mathfrak{a}}_{\mathfrak{c}'}$.
By (Eq3.9), $p$ has a vertex in $\eta^{-1}\mathsf{g}_{\mathfrak{c}}\mathrm{Qb}$ and a vertex in $\eta^{-1}\mathsf{g}_{\mathfrak{c}'}\mathrm{Qb}$.
Two vertices of a plaquette have physical sup distance at most $\eta\le1$, whereas by (Eq3.10) this distance is at least $\ell_c/4$.
But $\mathrm{Qb}$, a cube of side $D_\beta=D_{\mathrm{be}}+4\ge4$, lies in a cube of side $\frac{3}{4}\ell_c$; so $\ell_c\ge\frac{16}{3}$ and $\frac{\ell_c}{4}\ge\frac{4}{3}>1$, a contradiction.

\emph{The norm of the curvature.}
By (ps5)\,(iv),
\[
\|d\hat{\mathfrak{a}}_{\mathfrak{c}}\|=\|d(\mathsf{g}_{\mathfrak{c}})_*\hat{\mathfrak{a}}\|=\|d\hat{\mathfrak{a}}\|;
\]
so $\|d\hat{\mathfrak{a}}_{\mathfrak{c}}\|^2=\boldsymbol{\sigma}^2$.
\end{proof}

\begin{proof}[Proof of Proposition~35.3\,(e)]
One has $\hat{c}_{\mathfrak{c}}(b)=\frac{1}{e}L^{\frac{k_*}{2}}\cdot\hat{\mathfrak{a}}_{\mathfrak{c}}(L^{k_*}b)$.
So $b$ belongs to $\mathrm{Bd}_{\mathfrak{c}}$ if and only if $L^{k_*}b$ lies in the support of $\hat{\mathfrak{a}}_{\mathfrak{c}}$, and by (Eq3.9) and (ps6) the end points of the bonds of $\mathrm{Bd}_{\mathfrak{c}}$ lie in $a_*^{-1}\mathsf{g}_{\mathfrak{c}}\mathrm{Qb}$, which has a lift $\prod_\mu[\mathrm{lo}_\mu,\mathrm{lo}_\mu+D_\beta/a_*]$ (coordinates of level $k_*$).

The last assertion of (e) is (Eq3.10).

\emph{The number $d_{\mathrm{sep}}$.}
Let $y\in\mathrm{Ypr}_{\mathfrak{c}}$ and $y'\in\mathrm{Ypr}_{\mathfrak{c}'}$ with $\mathfrak{c}\ne\mathfrak{c}'$.
There are end points $c\in\Delta_y$ and $c'\in\Delta_{y'}$ of bonds of $\mathrm{Bd}_{\mathfrak{c}}$ and of $\mathrm{Bd}_{\mathfrak{c}'}$, respectively.
Since $\Delta_y=\{c:|c-y|<\frac{1}{2}\}$ (Section~I.1.1), one has $d(y,c)<\frac{1}{2}$ and $d(y',c')<\frac{1}{2}$, while $d(c,c')\ge\ell_*/(8L)$.
Hence $d(y,y')>\ell_*/(8L)-1$, and so $d_{\mathrm{sep}}\ge\ell_*/(8L)-1$ (if the sets are empty, then $d_{\mathrm{sep}}=+\infty$ by Section~G5.1.4).
Finally $\ell_*=Mr_*$ (Proposition~35.1\,(a)), and by (Eq3.5)
\[
\frac{Mr_*}{8L}-1\ge\frac{L^3}{8}-1>14\ge4.
\]

\emph{The number $\#|\mathrm{Ypr}_{\mathfrak{c}}^{+}|$.}
A site $y$ of $\mathrm{Ypr}_{\mathfrak{c}}$ has a lift with $y_\mu\in\left(\mathrm{lo}_\mu-\frac{1}{2},\,\mathrm{lo}_\mu+\frac{D_\beta}{a_*}+\frac{1}{2}\right)$ for $\mu=1,2,3$.
A site $y'$ of $\mathrm{Ypr}_{\mathfrak{c}}^{+}$ has distance at most $M+1$ from such a $y$; so it has a lift with
\[
y'_\mu\in\left(\mathrm{lo}_\mu-M-\frac{3}{2},\,\mathrm{lo}_\mu+\frac{D_\beta}{a_*}+M+\frac{3}{2}\right),
\]
an open interval of length $D_\beta/a_*+2M+3$.
An open interval of length $\mathrm{len}$ contains fewer than $\mathrm{len}+1$ integers, since the first and the last of them differ by less than $\mathrm{len}$.
Therefore
\[
\#|\mathrm{Ypr}_{\mathfrak{c}}^{+}|\le\left(\frac{D_\beta}{a_*}+2M+4\right)^3=\mathrm{pp}.
\]
(If $\mathrm{Ypr}_{\mathfrak{c}}$ is empty, so is $\mathrm{Ypr}_{\mathfrak{c}}^{+}$.)
\end{proof}

\begin{proof}[Proof of Proposition~35.3\,(f)]
The bond functions $\hat{c}_{\mathfrak{c}}$, $\mathfrak{c}\in\mathrm{Cel}$, on $T_*$ (mesh $\eta_*$, side $\ell_*>6$) satisfy the hypothesis of Lemma~P: by (e), for $\mathfrak{c}\ne\mathfrak{c}'$ the points of $\mathrm{Edp}(\hat{c}_{\mathfrak{c}})$ and of $\mathrm{Edp}(\hat{c}_{\mathfrak{c}'})$ have distance at least
\[
\frac{\ell_*}{8L}>15>1+2\eta_*
\]
by (Eq3.5).
Hence each of the three norms $\mathrm{nrm}(\cdot)$ satisfies
\[
\mathrm{nrm}\left(\mathscr{C}(t)\right)=\mathrm{nrm}\left(\sum_{\mathfrak{c}}t_{\mathfrak{c}}\hat{c}_{\mathfrak{c}}\right)=\max_{\mathfrak{c}}|t_{\mathfrak{c}}|\cdot\mathrm{nrm}(\hat{c}_{\mathfrak{c}}).
\]
By (d), $\mathrm{nrm}(\hat{c}_{\mathfrak{c}})=\mathrm{nrm}(\hat{\mathfrak{a}}^{[k_*]})/e$ for all $\mathfrak{c}$.
Adding the three norms and using (c2) one obtains
\[
|a(t)|_{\mathrm{ext}}=\max_{\mathfrak{c}}|t_{\mathfrak{c}}|\cdot\frac{|\hat{\mathfrak{a}}|_{\mathrm{ext}}}{e}\le\max_{\mathfrak{c}}|t_{\mathfrak{c}}|\cdot\frac{C_\beta}{e}.
\]

Let $t\in\mathrm{DDo}(t_x^{+})$.
If $C_\beta=0$, then $|a(t)|_{\mathrm{ext}}=0<e_*^{-1+3\varepsilon}$.
If $C_\beta>0$, then $\max_{\mathfrak{c}}|t_{\mathfrak{c}}|<t_x^{+}=e_*^{3\varepsilon}\cdot h^{-\frac{1}{2}}\cdot C_\beta^{-1}$, and
\[
|a(t)|_{\mathrm{ext}}<e_*^{3\varepsilon}\cdot\left(h^{\frac{1}{2}}\cdot e\right)^{-1}=e_*^{3\varepsilon}\cdot e_K^{-1}\le e_*^{3\varepsilon}\cdot e_*^{-1}=e_*^{-1+3\varepsilon},
\]
because $e_*\le e_K$ (Proposition~35.1\,(a)).

Since $\mathrm{Cb}=1+h^{\frac{1}{2}}\cdot C_\beta>h^{\frac{1}{2}}\cdot C_\beta$, one has $t_x=e_*^{3\varepsilon}/\mathrm{Cb}<t_x^{+}$, so that $t_x\mathrm{ue}_{\mathfrak{c}}\in\mathrm{DDo}(t_x^{+})$.
Moreover $\mathscr{C}(t_x\mathrm{ue}_{\mathfrak{c}})=t_x\hat{c}_{\mathfrak{c}}$.
Hence
\[
t_x\cdot\epsilon_{\mathfrak{c}}=e_*\cdot\left|\mathscr{C}(t_x\mathrm{ue}_{\mathfrak{c}})\right|_\infty\le e_*\cdot\left|a(t_x\mathrm{ue}_{\mathfrak{c}})\right|_{\mathrm{ext}}<e_*^{3\varepsilon}.
\]
As $e_*\le e_1\le\frac{1}{3}$ (Proposition~35.1\,(a); line (C9) of Section~I.8.2), it follows that
\[
\epsilon_{\mathfrak{c}}<\frac{e_*^{3\varepsilon}}{t_x}\le\frac{1}{t_x}\le\frac{2}{t_x}=q_x.
\]
\end{proof}

\begin{proof}[Proof of Proposition~35.3\,(g)]
By (Eq3.11) and (Eq3.3), $\hat{c}_{\mathfrak{c}}=\frac{1}{e}\,\tilde{g}\left(\hat{\mathfrak{a}}^{[k_*]}\right)$, and each of the maps $a\mapsto a^{[k_*]}$ and $A\mapsto\tilde{g}A$ is a relabelling of the bonds followed by a multiplication with non-zero numbers.
Hence $\hat{c}_{\mathfrak{c}}=0$ if and only if $\hat{\mathfrak{a}}=0$.

Let $\hat{\mathfrak{a}}\ne0$.
Then no $\mathrm{Bd}_{\mathfrak{c}}$ is empty.
Every site of $T_*$ lies in exactly one block $\Delta_y$ (Section~I.1.1); so an end point of a bond of $\mathrm{Bd}_{\mathfrak{c}}$ furnishes a site of $\mathrm{Ypr}_{\mathfrak{c}}$, and $\mathrm{Ypr}_{\mathfrak{c}}$ is contained in $\mathrm{Ypr}_{\mathfrak{c}}^{+}$.
The set $\mathrm{Cel}$ has $N_{*}=8L^3\ge2$ elements, and $d_{\mathrm{sep}}$ is the minimum of finitely many distances of non-empty finite sets; hence $d_{\mathrm{sep}}<\infty$.

Let $\hat{\mathfrak{a}}=0$.
Then all $\hat{c}_{\mathfrak{c}}$ vanish; so all $\mathrm{Bd}_{\mathfrak{c}}$, hence all $\mathrm{Ypr}_{\mathfrak{c}}$, and (the distance to the empty set being $+\infty$: Section~G5.1.4) all $\mathrm{Ypr}_{\mathfrak{c}}^{+}$ are empty.
\end{proof}

\paragraph{Uniformity}
The bounds of (b1) and (c1) involve $\beta$ and $h$ only.
The assertions (c2)--(f) hold for every $e$ with $e\cdot h^{\frac{1}{2}}\le e_1$, every $N\ge N_0(e)$ and every admitted placement.
The numbers $C_\beta$, $\mathbf{D}$, $D_\beta$ occurring in them depend neither on $e$, nor on $N$, nor on the placement, nor on a replacement of $\beta$ by $U_r\beta$ with $r\in\{\pm1\}^3$; indeed $\mathrm{Ind}_\beta$ is invariant under $r'\mapsto r'r$, one has $U_{r'}U_r=U_{r'r}$, and (Eq3.1') applies.
The number $\mathrm{pp}$ depends on $(e,\beta)$ through $D_\beta/a_*$, where $a_*\in\{h/L,h/L^2\}$.

\subsection*{4. Proposition~35.4 (what holds at the stop level and below it)}

\subsubsection*{Proof}

\paragraph{4.0. The data of the levels}
The setting $(\mathrm{ST}^{*})^{\mathrm{ch}}$ is the setting (FL) at $R=R^{*}(e)$ and $N\ge N_0(e)=N_{\min}(R^{*}(e),e)$, with the pair $P_N(e)$ and the primitives of (FL) (Sections~I.1.8, G5.1.1 and~G4.0.1).
Moreover $R^{*}(e)\in\mathrm{Tor}(e)$ and $R^{*}(e)=R_{\mathrm{top}}(e)/L<R_{\mathrm{top}}(e)$ (Section~I.6.2; the opening assertion of Proposition~35.1).
For $0\le j\le k_*$ the data of level $j$ are $(j,\ell_j,\bar{m}a_j,e_j)$, where
\[
e_j=e\cdot a_j^{\frac{1}{2}}=e_*\cdot L^{-\frac{k_*-j}{2}}.
\]
Proposition~35.1\,(a), (d) gives
\begin{equation}\tag{Eq4.1}
\ell_j\in\mathrm{Cl};\qquad0<\bar{m}a_j\le\frac{1}{L};\qquad0<e_j\le e_*\le e_0\qquad(0\le j\le k_*).
\end{equation}
Hence the data of every level $j\le k_*$ are data (SF) (Section~I.1.8).
At $j=0$ they are $(0,n_*,\bar{m}\eta,q)$: the side is $\ell_0=n_*$, the mass is $\bar{m}\eta$, and the symbol $e_j$ stands for $q$ (Section~I.2.5).
At $j=k_*$ they are $(k_*,\ell_*,\mu_*,e_*)$ with $k_*\ge1$; this is the setting (G2.S-3) (Section~I.1.8, paragraph Dictionary).

Since $v_j=-\log e_j\ge-\log e_0$, every condition of Proposition~14.2\,(b) on $e_0$ holds at $v_j$, $0\le j\le k_*$.
In particular, the row (A8) of Proposition~14.2\,(b) gives
\begin{equation}\tag{Eq4.2}
e_j^{\varepsilon}\le\frac{1}{4},\qquad3M\cdot e_j^{\varepsilon}\le1,\qquad e_j^{\frac{1}{4}}+e_j^{\varepsilon}\le1,\qquad e_j\le e_0<1.
\end{equation}

\paragraph{4.1. (SF1)}

\begin{proof}[Proof of Proposition~35.4\,(SF1)]
The trajectories in question are $(\varepsilon^{\mathrm{v},t}_k,m_k,E_k)$ and $(\mathscr{E}_k,\mathfrak{P}_k)$, $0\le k\le k_*$, the trajectories of $P_N(e)$ on $R^{*}(e)$ (Section~G5.1.1; Section~I.5.7, paragraph Flows; Section~G4.0.5).
Proposition~28.3 at the data (FL) of 4.0 states:
\begin{itemize}
\item[(a)] $P_N(e)$ is defined and is controlled on $R^{*}(e)$;
\item[(b)] $|m_k|\le w^{\mathrm{m}}_k/8$ for $k\le k_*(R^{*}(e))=k_*$;
\item[(c)] $P_N(e)$ is real.
\end{itemize}
By Section~G4.0.5, the words ``controlled on $R^{*}(e)$'' mean that the trajectory on $R^{*}(e)$ exists up to the level $k_*$ and that $|m_k|\le w^{\mathrm{m}}_k$ for $0\le k\le k_*$.
Since $|m_i|\le w^{\mathrm{m}}_i$ for all $i<k$, Proposition~28.1 applies at every $k\le k_*$: the data of level $k$ are defined, and $\varepsilon^{\mathrm{v},t}_k$ and $m_k$ are real.
At $k=k_*$ one has, since $w^{\mathrm{m}}_k=e_k^{\frac{3}{4}-8\varepsilon}$ (Section~I.5.4),
\[
|m_*|\le\frac{w^{\mathrm{m}}_*}{8}=\frac{e_*^{\frac{3}{4}-8\varepsilon}}{8}.
\]
\end{proof}

\paragraph{4.2. (SF2)}

\begin{proof}[Proof of Proposition~35.4\,(SF2)]
\emph{Domains.}
The sets $G_j$, $\tilde{G}_j$, $\mathrm{cf}\,\tilde{R}_j$, $\mathfrak{X}_j$ are those of Section~I.5.3, that is, of \cite[Definition~5.16]{QEDI}, formed with $(j,\ell_j,e_j)$; at $j=0$ they are formed with $q$.
In \cite[Section~17.3]{QEDI} the Lemmas~17.3 and~17.4, and Lemma~17.6 with the exception of its assertion on $\mathrm{Dom}_k$, are stated for every level $k\ge0$, every side $\ell\in\mathrm{Cl}$ and every $e_k\in(0,e_0]$; neither a mass nor a primitive occurs in them.
By (Eq4.1) they hold at every level $j\le k_*$:
\begin{itemize}
\item by \cite[Lemma~17.3]{QEDI}, the sets $G_j$, $\mathfrak{X}_j$, $\mathrm{cf}\,\tilde{R}_j$, $\tilde{G}_j$ are open in the space of the complex bond functions on $T_j$;
\item by \cite[Lemma~17.4(a)]{QEDI}, the sets $\tilde{R}_j$ and $\tilde{R}_j+\mathfrak{X}_j$ are contained in $\tilde{G}_j$;
\item by \cite[Lemma~17.6]{QEDI}, the sets $\mathrm{cf}\,\tilde{R}_j$, $\tilde{G}_j$, $\mathfrak{X}_j$ are invariant under $A\mapsto-A$, and $\mathrm{cf}\,\tilde{R}_j$, $\tilde{G}_j$ are invariant under $A\mapsto A-d\boldsymbol{\lambda}$, $\boldsymbol{\lambda}$ a real function on the sites of $T_j$, and under the addition of real constant bond functions.
\end{itemize}
The set $G_j$ is contained in $\tilde{G}_j$: an element $A$ of $G_j$ is $A_0+A_1$ with $A_0=0$ (representatives $A'=0$, $\boldsymbol{\lambda}=0$ on every $\square^{\sim5}$) and $A_1=A$.
Further, $0\in\tilde{R}_j$ (with $A_0=0$, $A_1=0$).
Hence $\mathfrak{X}_j=0+\mathfrak{X}_j$ is contained in $\tilde{R}_j+\mathfrak{X}_j$, and therefore in $\tilde{G}_j$.

\emph{The inverse; the four majorants.}
The data $(k_*,\ell_*,\mu_*,e_*)$ are data (SF) with $k_*\ge1$ (4.0).
The operator $\mathrm{Fst}(A)=D_A+\mu_*+b_*JP_*(A)J^{*}$ is the operator of the property FE1$^{\mathrm{ch}}$ of Proposition~14.1\,(i), which is Proposition~2.8\,(a)--(c) at $y=\infty$, in the norm $|\cdot|_c$, with the constant $C^{F}$.
By Proposition~14.1\,(i), for $A\in\tilde{G}_*$ the operator $\mathrm{Fst}(A)$ is bijective on $\mathfrak{S}^{\mathrm{ch}}(\eta_*,\ell_*)$; the map $S_*(\cdot)=S_{k_*}(\cdot)$ is analytic on $\tilde{G}_*$, which is open (\cite[Lemma~17.3]{QEDI}); and for all unit sites $z,z'$ and all $f\in\mathfrak{S}^{\mathrm{ch}}(\eta_*,\ell_*)$
\begin{align*}
\left|\mathbf{1}_{\Delta_z}S_*(A)\mathbf{1}_{\Delta_{z'}}f\right|_\infty&\le C^{F}\cdot\exp\left(-\gamma^{F}d(z,z')\right)\cdot|f|_\infty,\\
\left|\mathbf{1}_{\Delta_z}\delta_{\alpha,A}S_*(A)\mathbf{1}_{\Delta_{z'}}f\right|_\infty&\le C^{F}\cdot\exp\left(-\gamma^{F}d(z,z')\right)\cdot|f|_\infty,
\end{align*}
the second bound for the short pairs.
Here, on spinors, $\mathbf{1}_{\Delta_z}$ is the indicator of $\Delta_z\times O_{\mathrm{cub}}$ ($48\eta_*^{-3}$ fermion sites: the photon block and the $48$ labels of each of its clusters), and on functions of pairs it restricts to the pairs whose first point lies over $\Delta_z$ (Section~G1.3).
These are the indicators $\mathbf{1}_y$ of Section~G5.1.7.
The short pairs are the pairs of $\mathrm{Pz}_*$ with $d_X<1$; since $\eta_*\le1/L<1$ (Proposition~35.1\,(a)), the $48\cdot47$ ordered pairs $(c,g)$, $(c,g')$, $g\ne g'$, inside each cluster, whose distance is $d_X=\eta_*$, are among them (Section~I.1.7).
Taking the supremum over $|f|_\infty\le1$, one obtains
\[
|S_*(A)|(y,y')\le C^{F}\cdot\mathrm{Edy}(y,y')\qquad\text{and}\qquad|\mathrm{Sa}_*(A)|(y,y')\le C^{F}\cdot\mathrm{Edy}(y,y').
\]
By Proposition~6\,(f) (the property FE4$^{\mathrm{ch}}$, at $k=k_*\ge1$), for $A\in\tilde{G}_*$, unit sites $z,z'$ and $\mathrm{vv}\in\mathbb{C}^4$,
\begin{align*}
\left|\mathbf{1}_{\Delta_z}H_*(A)\mathbf{1}_{z'}\mathrm{vv}\right|_\infty&\le C^{F}\cdot\exp\left(-\gamma^{F}d(z,z')\right)\cdot|\mathrm{vv}|_c,\\
\left|\mathbf{1}_{\Delta_z}\delta_{\alpha,A}H_*(A)\mathbf{1}_{z'}\mathrm{vv}\right|_\infty&\le C^{F}\cdot\exp\left(-\gamma^{F}d(z,z')\right)\cdot|\mathrm{vv}|_c
\end{align*}
(short pairs, that is $d_X<1$, the pairs inside a cluster included: Section~I.5.8).
For $f\in\mathfrak{S}(1,\ell_*)$ with $|f|_\infty\le1$ the spinor $\mathbf{1}_{y'}f$ equals $\mathbf{1}_{y'}\mathrm{vv}$ with $\mathrm{vv}=f(y')$, and $|\mathrm{vv}|_c\le1$.
Hence $|H_*(A)|(y,y')$ and $|\delta_{\alpha,A}H_*(A)|(y,y')$ are at most $C^{F}\cdot\mathrm{Edy}(y,y')$.

\emph{Long pairs.}
Proposition~6\,(b), which is (I.Lc2)$^{\mathrm{ch}}$, gives for $A\in\tilde{G}_*$ and a pair with $1\le d_X(x,x')\le M$, in the form with $|\cdot|_c$,
\[
\left|(\delta_{\alpha,A}f)(x,x')\right|_c\le\mathrm{e}\cdot|f(x')|_c+|f(x)|_c;
\]
indeed, by (I.1.7), $d_X^{-\alpha}\le1$, and $|U_A(x,x')|\le\mathrm{e}$.

\emph{The set $\mathrm{Gwk}$.}
By its definition (Section~G5.1.7) the set $\mathrm{Gwk}$ is contained in $\tilde{R}_*+\mathfrak{X}_*$, which lies in $\tilde{G}_*$.
If $A=A_r+\mathscr{C}$ with $A_r$ real in $\tilde{R}_*$ and $\mathscr{C}\in\mathfrak{X}_*$, then $-A=(-A_r)+(-\mathscr{C})$, where $-A_r$ is real and lies in $\tilde{R}_*$ and $-\mathscr{C}\in\mathfrak{X}_*$ (\cite[Lemma~17.6]{QEDI}).
So $\mathrm{Gwk}$ is invariant under $A\mapsto-A$.

For $A=A_r+\mathscr{C}\in\mathrm{Gwk}$ one has $\operatorname{Im}A=\operatorname{Im}\mathscr{C}$; hence, by the definition of $\mathfrak{X}_*$,
\[
\mathfrak{i}(A)=e_*\cdot|\operatorname{Im}\mathscr{C}|_\infty\le e_*\cdot|\mathscr{C}|_\infty<e_*\cdot e_*^{-1+2\varepsilon}=e_*^{2\varepsilon}.
\]
By (Eq4.2), $e_*^{2\varepsilon}\le\frac{1}{16}$ and $3M\cdot\mathfrak{i}(A)<\left(3M\cdot e_*^{\varepsilon}\right)\cdot e_*^{\varepsilon}\le\frac{1}{4}$.

\emph{The number $\mathrm{ssum}$.}
Let $0<\gamma\le6$, and let $\ell\ge1$ be an integer.
Every class $y'$ of $(\mathbb{Z}/\ell)^3$ has a lift with $|y'-y|=d(y,y')$, and $|z|\ge|z|_1/3$ for $z\in\mathbb{Z}^3$.
Therefore
\[
\sum_{y'}\exp\left(-\gamma\cdot d(y,y')\right)\le\sum_{z\in\mathbb{Z}^3}\exp\left(-\frac{\gamma\cdot|z|_1}{3}\right)=\left(\sum_{i\in\mathbb{Z}}\exp\left(-\frac{\gamma\cdot|i|}{3}\right)\right)^3=\coth\left(\frac{\gamma}{6}\right)^3.
\]
Since $\coth t=1+\frac{2}{\exp(2t)-1}\le1+\frac{1}{t}$ for $t>0$, one has $\coth(\gamma/6)\le1+6/\gamma\le12/\gamma$.
Consequently
\[
\sum_{y'}\exp\left(-\gamma\cdot d(y,y')\right)\le\left(\frac{12}{\gamma}\right)^3=\mathrm{ssum}(\gamma).
\]
By Proposition~14.2\,(f), $0<\gamma^{F}\le1/(6L)$; so $\gamma=\gamma^{F}/2$ lies in $(0,6]$.
For $O=S_*(A)$ or $O=\mathrm{Sa}_*(A)$, where $A\in\tilde{G}_*$ (in particular, $A\in\mathrm{Gwk}$),
\[
\sum_{y'}|O|(y,y')\cdot\exp\left(\frac{\gamma^{F}}{2}d(y,y')\right)\le C^{F}\sum_{y'}\exp\left(-\frac{\gamma^{F}}{2}d(y,y')\right)\le C^{F}\cdot\mathrm{ssum}\left(\frac{\gamma^{F}}{2}\right)=\mathbf{S},
\]
and the same holds for the sum over $y$ at fixed $y'$, because the bound is symmetric in $(y,y')$.
Hence $\|S_*(A)\|_{\frac{\gamma^{F}}{2}}\le\mathbf{S}$ and $\|\mathrm{Sa}_*(A)\|_{\frac{\gamma^{F}}{2}}\le\mathbf{S}$.

\emph{The set $\mathrm{Gwk}^{+}$.}
For $A+A_1$ with $A\in\mathrm{Gwk}$ and $e_*\cdot|A_1|_\infty\le c_{\tau}$ one has
\[
\mathfrak{i}(A+A_1)\le\mathfrak{i}(A)+e_*\cdot|A_1|_\infty<\frac{1}{16}+c_{\tau}\le\frac{1}{16}+\frac{1}{8}<\frac{1}{4},
\]
because $c_{\tau}\le\frac{1}{8}$ (Section~G5.1.8).
\end{proof}

\paragraph{4.3. (SF3)}
By (SF1), $|m_i|\le w^{\mathrm{m}}_i$ for every $i<k_*$.
Proposition~28.1 therefore applies at $k=k_*$ (see also Proposition~28.3\,(b)), and it yields that $\mathscr{E}_*=\mathscr{E}_{k_*}$ belongs to $K^{x}_{k_*}(\ell_*)=K^{x}_*$, with
\[
\left\|\mathscr{E}_*\right\|^{x}_*\le 4e_*^{\varepsilon}\cdot w^{\mathrm{E}}_*.
\]
Membership in this space means that the clauses of its definition in Section~I.5.4 hold, read on $T^{\mathrm{f}}_*$.
Thus $\mathscr{E}_*(X;\mathscr{A},\mathscr{C},\psi,\chi)$ is an even element of $\mathrm{Gr}[\psi,\chi]$.
Its kernels $\mathscr{E}_{nm}=\mathscr{E}_{nm}(X;\mathscr{A},\mathscr{C})$ are taken antisymmetric in the $\xi_i$ and, separately, in the $\boldsymbol{\zeta}_j$; such kernels exist and are determined by the element (Proposition~15\,(a)), and the clauses (I.k1), (I.k2) below and the display (2.1) refer to them.
Each $\mathscr{E}_{nm}$ has as many barred as unbarred legs, and:
\begin{itemize}
\item[(I.k1)] the kernels are analytic in $(\mathscr{A},\mathscr{C})\in\tilde{R}_*\times\mathfrak{X}_*$;
\item[(I.k2)] they vanish unless every $\xi_i$ lies over $X$ and every pair $\boldsymbol{\zeta}_j$ is admissible for $X$ (Section~I.1.7: either of kind (I), to which every pair inside a cluster over $X$ belongs, or of kind (II), with $x'=x^{(3)}$ and $c\in X$ or $c_3\in X$), and they depend on $\mathscr{A}(b)$, $\mathscr{C}(b)$ only for the bonds $b$ that meet $X$;
\item[(I.k6)] for all $X\in D_*$ and all $(\mathscr{A},\mathscr{C})\in\tilde{R}_*\times\mathfrak{X}_*$,
\[
\exp\left(\kappa\cdot d_M(X)\right)\cdot\left\|\mathscr{E}_*(X;\mathscr{A},\mathscr{C})\right\|_{h_*}\le\left\|\mathscr{E}_*\right\|^{x}_*.
\]
\end{itemize}
For these antisymmetric kernels, Section~I.5.1 gives
\[
\left\|\mathscr{E}_*(X;\mathscr{A},\mathscr{C})\right\|_{h_*}=\sum_{n,m}\frac{h_{*,1}^{n}\cdot h_{*,2}^{m}}{n!\,m!}\left\|\mathscr{E}_{nm}(X;\mathscr{A},\mathscr{C})\right\|,
\qquad
\left\|\mathscr{E}_{nm}\right\|=\int\mathrm{d}\xi\,\mathrm{d}\boldsymbol{\zeta}\,\left|\mathscr{E}_{nm}\right|,
\]
where, by Section~I.3.6, $\int\mathrm{d}\xi$ denotes the sum over all $\xi=(c,g,\beta,\omega)$ with $c\in T_*$, $g\in O_{\mathrm{cub}}$ (the $48$ labels), $\beta\in\{1,\dots,4\}$, $\omega\in\{0,1\}$, each term with the weight $\eta_*^3/48$, and $\int\mathrm{d}\boldsymbol{\zeta}$ denotes the sum over all $\boldsymbol{\zeta}=(x,x',\beta,\omega)$ with $(x,x')\in\mathrm{Pz}_*$, each term with the weight $(\eta_*^3/48)^2$.
Combined with (I.k6) and the bound on $\|\mathscr{E}_*\|^{x}_*$, this is (2.1).

\paragraph{4.4. (SF4)}
\textbf{The primitives.}
Recall that $k_*=k_*(R^*(e))$.
By Proposition~26.6\,(b), at every level $j<k_*$ the data $(j,\ell_j,\bar{m}\cdot a_j,e_j)$, together with the primitives of (FL), are (OS).
According to the definition of (OS) (Section~I.1.8; Section~G4.0.1), this includes that the primitives form a system of primitives of level $j$ on $T_j(R^*(e))$ having the properties W1--W17$^{\mathrm{ch}}$ for $\mathrm{par}$.
Here W1--W17 are the properties of systems of primitives listed in \cite[Section~5.4.2]{QEDI}, and the mark ch means that they are read as prescribed in Section~I.5.8.
In more detail, $R^*(e)<R_{\mathrm{top}}(e)$, so that the primitives on $T_j(R^*(e))$ are the projections (Proposition~26.1), with $\mathrm{Jr}=R_{\mathrm{top}}(e)/R^*(e)=L$, of the systems of Proposition~26.5 on $T_j(R_{\mathrm{top}}(e))$.
The data of Proposition~26.1 are given at every level $j\le k_*$ by Proposition~26.6\,(d).
By Proposition~26.5 the systems on $T_j(R_{\mathrm{top}}(e))$ have W1--W17$^{\mathrm{ch}}$ for $\mathrm{par}^0$.
By Proposition~26.3\,(b) their projection is a system of primitives of level $j$ on $T_j(R^*(e))$ with W1--W17$^{\mathrm{ch}}$ for $\mathrm{par}=(\alpha_0,\gamma_{\mathrm{w}},\kappa_{\mathrm{w}};C_w)$.

At the level $j=0$ (Section~G4.0.4; Section~I.5.8) the system consists of (p2), $\Gamma_0(s'',A)$ and $S_{0,y}(X,A)$ on $\mathfrak{S}^{\mathrm{ch}}(1,n_*)$, together with the two diagonal families, and $S_{0,y}(X,A)$ is a Borel function of $y$.
There W7 is the following bound on the kernel (I.1.4) at pairs of fermion sites (Proposition~26.3\,(b), second sentence):
\[
\left|S_{0,y}(X,A)(x,x')\right|_c\le C_w\cdot(1+y)^{-1}\cdot\exp\left(-\gamma_{\mathrm{w}}\cdot d(x,x')-\kappa_{\mathrm{w}}\cdot d_M(X)\right).
\]

\textbf{(i)--(iv).}
Let $0\le j<k_*$.
By 4.0 the data of level $j$ are (SF).
The family (p5) of the system, that is, $S_{j,y}(X,A)$ with $X\in D_j$, $A\in\tilde{G}_j$ and $y\ge0$, is a Borel function of $y$ and has the properties W1, W2, W7, W12, W13, W14 for $\mathrm{par}$; here W7 is meant, at $j=0$, in its form for the level $0$, and W14 refers to the group $\mathbb{G}^{L}_{j,3}$.
For the family (p5) the displays of these six properties in Section~G4.0.4 coincide with those of Section~G1.11, which is the form in which Proposition~13 takes them.
So the data of Proposition~13 are given.
The functions $\acute{E}_j$ and $E^{d}_j$ are formed with this family by the displays of Section~I.5.7.
For $j\ge1$,
\[
\acute{E}_j(X,A)=-i\cdot b_j^2\cdot\int_0^{\infty}\mathrm{d}y\,\mathrm{Tr}\left[\gamma_5\,B_{j,y}(A)\,Q_j(A)\,J^{*}\,S_{j,y}(X,A)\,J\,Q_j^{T}(-A)\,B_{j,y}(A)\right],
\]
a trace on $\mathfrak{S}(1,\ell_j)$, which has four components per unit site.
At $j=0$,
\[
\acute{E}_0(X,A)=\int_1^{\infty}\frac{\mathrm{d}y}{y}\,\mathrm{Tr}\left[F_0(A)\,S_{0,y}(X,A)\right]-i\int_0^1\mathrm{d}y\,\mathrm{Tr}\left[\gamma_5\,S_{0,y}(X,A)\right],
\]
with traces on $\mathfrak{S}^{\mathrm{ch}}(1,n_*)$, which has $192$ components per photon site; here $\mathrm{Tr}\,O=\int\mathrm{d}x\,\mathrm{tr}\,O(x,x)$ by (I.ke2).
Finally, $E^{d}_j(X,A)$ is the mean of the four numbers
\[
\acute{E}_j(X,A),\qquad\acute{E}_j(X,-A),\qquad\overline{\acute{E}_j\left(X,\overline{A}\right)},\qquad\overline{\acute{E}_j\left(X,-\overline{A}\right)}.
\]
Proposition~13\,(d1) states that $E^{d}_j(X,\cdot\,)$ is analytic on $\tilde{G}_j$ and that
\[
\left|E^{d}_j(X,A)\right|\le C_d\cdot M^3\cdot\exp\left(-(\kappa+1)\cdot d_M(X)\right);
\]
these are the analyticity and (i).
Proposition~13\,(d2) states that $E^{d}_j(X,A)=E^{d}_j(X,\tilde{A})$ whenever $A,\tilde{A}\in\tilde{G}_j$ agree on the bonds that meet $X$; this is (ii).
Proposition~13\,(d3) states that $E^{d}_j(X,A-d\boldsymbol{\lambda})=E^{d}_j(X,A)$ for every real function $\boldsymbol{\lambda}$ of the photon sites $c$ (by (SF2), $A-d\boldsymbol{\lambda}$ lies in $\tilde{G}_j$), and that $E^{d}_j(X,-A)=E^{d}_j(X,A)$; these are (iii) and (iv).
The relation $E_j^{\mathrm{det}}(X,A)=E^{d}_j(X,A)-E^{d}_j(X,0)$ is the definition of Section~I.5.7; since $0\in\tilde{G}_j$, the map $E_j^{\mathrm{det}}(X,\cdot\,)$ is a function on $\tilde{G}_j$.
The equation $\mathfrak{P}_0=0$ and the recursion for $\mathfrak{P}_{j+1}$ constitute the definition of the trajectory with slot (Section~I.5.7, paragraph \emph{Flows}; Section~G4.0.5).

\emph{The right side of the recursion is defined.}
Let $0\le j<k_*$ and $\mathscr{C}\in\mathfrak{X}_{j+1}$.
By Section~I.3.7, $\mathscr{C}_L(b)=L^{-\frac{1}{2}}\cdot\mathscr{C}(b/L)$ is a bond function on $T_j$; indeed, $\eta_j=L\cdot\eta_{j+1}$ and $\ell_j=L\cdot\ell_{j+1}$, so that $b\mapsto b/L$ maps the bonds of $T_j$ bijectively onto the bonds of $T_{j+1}$.
Moreover $\mathscr{C}_L\in\mathfrak{X}_j$.
This is \cite[Lemma~17.5(c)]{QEDI}, whose hypotheses hold here: $\ell_j$ and $\ell_j/L=\ell_{j+1}$ lie in $\mathrm{Cl}$, $e_{j+1}=L^{\frac{1}{2}}\cdot e_j$, and $e_j\le e_0$ (Proposition~35.1\,(d), (Eq4.1)).
In detail,
\begin{align*}
|\mathscr{C}_L|_{\infty}&=L^{-\frac{1}{2}}\cdot|\mathscr{C}|_{\infty}<L^{-\frac{1}{2}}\cdot e_{j+1}^{-1+2\varepsilon}=L^{-1+\varepsilon}\cdot e_j^{-1+2\varepsilon}\le e_j^{-1+2\varepsilon},\\
\left(\partial_\mu(\mathscr{C}_L)_\nu\right)(c)&=\eta_j^{-1}\cdot L^{-\frac{1}{2}}\cdot\left(\mathscr{C}_\nu\left(\tfrac{c}{L}+\eta_{j+1}\cdot e_\mu\right)-\mathscr{C}_\nu\left(\tfrac{c}{L}\right)\right)=L^{-\frac{3}{2}}\cdot\left(\partial_\mu\mathscr{C}_\nu\right)\left(\tfrac{c}{L}\right),\\
|\partial\mathscr{C}_L|_{\infty}&=L^{-\frac{3}{2}}\cdot|\partial\mathscr{C}|_{\infty}<L^{-\frac{3}{2}}\cdot e_{j+1}^{-1+3\varepsilon}=L^{-2+\frac{3\varepsilon}{2}}\cdot e_j^{-1+3\varepsilon}\le e_j^{-1+3\varepsilon},
\end{align*}
where the exponents $-1+\varepsilon$ and $-2+\frac{3\varepsilon}{2}$ of $L$ are negative because $\varepsilon\le1/200$.
By (SF2), $\mathfrak{X}_j$ is contained in $\tilde{G}_j$.
The assertion now follows by induction on $j$.
For every $X\in D_0$ the function $\mathfrak{P}_0(X,\cdot\,)=0$ is a complex function on $\mathfrak{X}_0$.
Assume that $\mathfrak{P}_j(X,\cdot\,)$ is a complex function on $\mathfrak{X}_j$ for every $X\in D_j$.
For $\mathscr{C}\in\mathfrak{X}_{j+1}$ the numbers $\mathfrak{P}_j(X,\mathscr{C}_L)$ and $E_j^{\mathrm{det}}(X,\mathscr{C}_L)$ are then defined, and $D_j$ is a finite set.
Hence the right side of the recursion is a finite sum of complex numbers, and $\mathfrak{P}_{j+1}(Y,\cdot\,)$ is a complex function on $\mathfrak{X}_{j+1}$ for every $Y\in D_{j+1}$.
(For the trajectory, this assertion is also one of the clauses of Proposition~28.1.)
At $j=0$ the domain $\mathfrak{X}_0$ is formed with $q$, one has $e_{j+1}=L^{\frac{1}{2}}\cdot q$, and the right side is the sum of the numbers $E_0^{\mathrm{det}}(X,\mathscr{C}_L)$, formed with $\acute{E}_0$.

\textbf{(v).}
The estimates on polymers of \cite[Theorem~13.2]{QEDI} are stated for every $k\ge0$ and every $\ell\in\mathrm{Cl}$; by (Eq4.1) they apply at every level $j\le k_*$.
Its assertions bearing the marks (t2) and (t7) give the following.
\begin{itemize}
\item[(t2)] For every $M$-cube $\square$,
\[
\sum_{X\in D_j,\ X\supset\square}\exp\left(-\kappa_0\cdot d_M(X)\right)\le K_0.
\]
\item[(t7)] Any two points of a polymer $X$ have sup distance at most $M\cdot\left(d_M(X)+2\right)$.
\end{itemize}
By \cite[Section~17.4.1, (g0)]{QEDI}, which requires $L$ odd, $M$ a power of $L$ and $\ell_j\in\mathrm{Cl}$, every block $B_j(y)=\Delta_y$ is contained in one $M$-cube.
For $j<k_*$ both $\ell_j$ and $\ell_j/L=\ell_{j+1}$ lie in $\mathrm{Cl}$ (Proposition~35.1\,(d)), and the assertions of \cite[Theorem~13.2]{QEDI} bearing the marks (t5) and (t8) give the following.
\begin{itemize}
\item[(t5)] For every $X\in D_j$ there is exactly one $Y\in D_{j+1}$ with $\bar{X}^{L}=LY$, and $L\cdot d_M(Y)\le d_M(X)$.
\item[(t8)] At the rate $\kappa$, which is admissible since $\kappa\ge\kappa_0$ by (C2),
\[
\sum_{X:\ \bar{X}^{L}=LY}\exp\left(-\kappa\cdot d_M(X)\right)\le c_{\mathrm{rb}}\cdot L^3\cdot\exp\left(-L\cdot(\kappa-\kappa_0-1)\cdot d_M(Y)\right).
\]
\end{itemize}

\textbf{(vi).}
Let $a$ be a complex one-cochain on $\mathbb{Z}^3/n_*\mathbb{Z}^3$ such that $a^{[j]}\in\mathfrak{X}_j$ for all $j\le k_*$.
Then $a^{[k_*]}$ lies in $\mathfrak{X}_*$, which is a subset of $\tilde{G}_*$ by (SF2).
Consequently $\mathrm{Fst}\left(a^{[k_*]}\right)$ is bijective (SF2), and $S_*\left(a^{[k_*]}\right)$, $D_*\left(a^{[k_*]}\right)$ are defined (Section~G5.1.7).
The bond function $0$ is real, so that (A7) (see 4.9 below; Proposition~19\,(d)) shows that $D_*(0)$ is invertible, that is, $\det D_*(0)\ne0$.

Proposition~20.3 is now applied at $k=k_*$.
Its data are the following:
the setting (S-4), that is, (FL) without the primitives (Section~I.1.8), which is given by 4.0;
the integer $k_*$, which satisfies $0\le k_*\le k_*(R^*(e))$;
for every $j<k_*$ a family (p5) of level $j$ on $T_j(R^*(e))$ with W1, W2, W7$^{\mathrm{ch}}$, W12, W13, W14 for $\mathrm{par}$, which is given by the first paragraph of 4.4;
the functions $\acute{E}_j$, $E^{d}_j$, $E_j^{\mathrm{det}}$ formed with these families;
the family determined by $\mathfrak{P}_0=0$ and
\[
\mathfrak{P}_{j+1}(Y,\mathscr{C})=\sum_{X:\ \bar{X}^{L}=LY}\left(\mathfrak{P}_j(X,\mathscr{C}_L)+E_j^{\mathrm{det}}(X,\mathscr{C}_L)\right),
\]
which is the family $\mathfrak{P}_j$ of the trajectory;
and the cochain $a$.
Clause (c) of Proposition~20.3 at $k=k_*$ reads as follows, and it is (2.2):
\begin{gather*}
\frac{\det\left(D_a+\bar{m}\cdot\eta\right)}{\det\left(D_0+\bar{m}\cdot\eta\right)}=\exp\left\{\sum_{X\in D_*}\mathfrak{P}_*\left(X,a^{[k_*]}\right)\right\}\cdot\frac{\det D_*\left(a^{[k_*]}\right)}{\det D_*(0)},\\
\sum_{X\in D_*}\mathfrak{P}_*\left(X,a^{[k_*]}\right)=\sum_{j<k_*}\ \sum_{X\in D_j}E_j^{\mathrm{det}}\left(X,a^{[j]}\right).
\end{gather*}

\emph{Sizes.}
The operators $D_a+\bar{m}\cdot\eta$ and $D_0+\bar{m}\cdot\eta$ are linear maps of $\mathfrak{S}^{\mathrm{ch}}(1,n_*)$, whose dimension is $4\cdot48\cdot n_*^3=192\,n_*^3$ (Section~I.1.5).
The operator $D_*(\cdot\,)$ is a linear map of $\mathfrak{S}(1,\ell_*)$, whose dimension is $4\ell_*^3$.
The index $j$ in the last sum ranges over $0\le j<k_*$.
As $k_*\ge1$, this range contains $j=0$, where $E_0^{\mathrm{det}}$ is formed with $\acute{E}_0$ and $a^{[0]}=a$.

\paragraph{4.5. (SF5)}
All objects of this paragraph are the photon objects of \cite[Section~5.2.5 and Definition~5.11]{QEDI} (Section~I.5.7; Section~G5.1.5), on the tower of the sides $\ell_j=n_*/L^j$, $0\le j\le k_*$, where $n_*=L^{k_*}\cdot\ell_*$ and $k_*\ge1$.
This is the setting of \cite[Sections~14.5.1 and 14.5.2]{QEDI}, where $k\ge1$ and $n\in L^k\mathbb{N}$, taken at $k=k_*$ and $n=n_*$.
The slice $\mathrm{Sl}$ and the coordinates $(A_*;\zeta_j)$ of \cite[Section~14.5.2]{QEDI} are those of \cite[Definition~5.11]{QEDI}.
Indeed, $A_j=(Q_jA_0)^{[j]}$ by \cite[Section~14.5.1]{QEDI}; $\mathrm{Sl}$ is the set of \cite[Definition~5.11]{QEDI} by the second paragraph of \cite[Section~14.5.2]{QEDI}, the one that follows the paragraph in which $\mathrm{Sl}$ is defined; and, by \cite[Theorem~14.16(a)]{QEDI}, for $A_0=\mathrm{Af}_0(A,\zeta)$, where $\mathrm{Af}_0$ is the map denoted $A^{\downarrow}_0$ in \cite[Definition~5.11]{QEDI} and in \cite[Section~14.5.1]{QEDI}, one has $A=A_k$ and $\mathrm{Sq}_j\zeta_j=A_j-A^{\min}_j(A_{j+1})$, which is the definition of $A_*$ and $\zeta_j$ in \cite[Definition~5.11]{QEDI}.

\textbf{Coordinates.}
The following assertions are part of \cite[Theorem~14.17]{QEDI}.
\begin{itemize}
\item[(b)] The map $F^{\mathrm{u}}\mapsto(A_*;\zeta_0,\dots,\zeta_{k_*-1})$ is a linear bijection of $F_N$ onto ${V_{\mathrm{t}}}(T^0_*)\times\prod_{j<k_*}V_j$.
\item[(c)] There is an orthonormal family $\{\phi_{j,\Upsilon}\}$ in $F_N$ such that $\zeta_j(\Upsilon)=\langle\phi_{j,\Upsilon},F^{\mathrm{u}}\rangle$.
\item[(d)] For every plaquette $p$ of $T^0_*$ one has $(dA_*)(p)=\langle\phi_p,F^{\mathrm{u}}\rangle$ with some $\phi_p\in F_N$ such that $\|\phi_p\|\le1$.
\item[(e)] The law $\mathrm{d}\mathsf{m}_*$ of $A_*$ under $\gamma_N$ is the centred Gaussian measure on ${V_{\mathrm{t}}}(T^0_*)$ whose density with respect to the Lebesgue measure of ${V_{\mathrm{t}}}(T^0_*)$ is proportional to $\exp\left(-\frac{1}{2}|dH_*A|_2^2\right)$.
\end{itemize}
A Gaussian measure that has a density is non-degenerate.
The torus $T^0_*$ has $\ell_*^3$ sites, and at each site there are three plaquettes $(y;\mu\nu)$, $\mu<\nu$; so there are $3\ell_*^3$ plaquettes.
By \cite[Proposition~15.21(ii)]{QEDI}, which requires that $\ell_j$ be a multiple of $L$, the basis $\{\Upsilon\}$ of $V_j$ has $(2-2L^{-3})\cdot n_j\le2n_j$ elements, where $n_j=\ell_j^3$.
Since $\ell_j=R^*(e)/a_j=L^{k_*-j}\cdot\ell_*$, one has $n_j=\ell_*^3\cdot L^{3(k_*-j)}$.

\textbf{The set $\{\chi_*=1\}$.}
This set equals $\left\{A\in{V_{\mathrm{t}}}(T^0_*):\ |dA|_{\infty}\le p_*\right\}$.
Let $A\in{V_{\mathrm{t}}}(T^0_*)$ satisfy $dA=0$.
Then $A$ is closed and coclosed, hence constant by the Hodge decomposition on the discrete torus \cite[Lemma~15.1]{QEDI}; as its mean vanishes, $A=0$.
Therefore $A\mapsto|dA|_{\infty}$ is a norm on the finite-dimensional real vector space ${V_{\mathrm{t}}}(T^0_*)$, and the set in question is the closed ball of radius $p_*=v_*^{p}>0$ for this norm.
It is compact and convex, and $0$ is an interior point of it.

\textbf{The representative.}
Let $A\in{V_{\mathrm{t}}}(T^0_*)$ with $|dA|_{\infty}\le p_*$.
Then $A$ is a real one-cochain with mean zero on $(\mathbb{Z}/\ell_*\mathbb{Z})^3$, and $\ell_*\ge2$ because $\ell_*\in\mathrm{Cl}$ by (Eq4.1).
By \cite[Corollary~15.20]{QEDI} there is a function $\boldsymbol{\lambda}$ on $T^0_*$ such that $|A-d\boldsymbol{\lambda}|_{\infty}\le2\ell_*^2\cdot|dA|_{\infty}$.
Since $A$ is real,
\[
\left|A-d\operatorname{Re}\boldsymbol{\lambda}\right|_{\infty}=\left|\operatorname{Re}\left(A-d\boldsymbol{\lambda}\right)\right|_{\infty}\le\left|A-d\boldsymbol{\lambda}\right|_{\infty},
\]
and $\boldsymbol{\lambda}$ may be taken real.
Put $A^{g}:=A-d\boldsymbol{\lambda}$; then $|A^{g}|_{\infty}\le2\ell_*^2\cdot p_*$.

\emph{The minimiser.}
Here $H_*[\,\cdot\,]$ is the photon minimiser $H_{k_*}$ (see (2.3)).
By \cite[Proposition~15.5]{QEDI}, which holds for every $k\ge0$ and every positive integer $\ell$, it is linear and real, and $H_*[\mathrm{cst}]=\mathrm{cst}$ for every constant $\mathrm{cst}$: the constant bond function on $T^0_*$ is mapped to the constant bond function on $T_*$ with the same components.
Consequently $\mathscr{A}^0=H_*[A^{g}]$ is real, and
\[
H_*\left[\mathrm{sv}\cdot A^{g}+c_\theta\right]=\mathrm{sv}\cdot\mathscr{A}^0+c_\theta\qquad\text{for all real }\mathrm{sv}\text{ and all }\theta\in\mathbb{R}^3.
\]

\emph{The sup norm.}
Write $(H_*[A'])(b)=\sum_{b'}H_*(b,b')\,A'(b')$, where $b'$ ranges over the positively oriented unit bonds.
By \cite[Theorem~15.14]{QEDI}, which holds for $0\le\alpha<1$, every $k\ge0$ and every positive integer $\ell$, together with the choice of $C_P$ and $\gamma_P$ in \cite[Section~15.7]{QEDI}, which gives $C_P\ge C_{\mathrm{Hk}}(\alpha)$ and $\gamma_P\le\gamma_H$ (see the closing paragraph of that section, on the constants $C_P$, $\gamma_P$, $c_P$), the kernel satisfies
\[
\left|H_*(b,b')\right|\le C_P\cdot\exp\left(-\gamma_P\cdot d(b,b')\right).
\]
Since $0<\gamma_P\le1$, \cite[Lemma~15.2]{QEDI} gives
\[
\sum_{b'}\exp\left(-\gamma_P\cdot d(b,b')\right)\le c_{\mathrm{bd}}(\gamma_P).
\]
Hence $|H_*[A']|_{\infty}\le C_P\cdot c_{\mathrm{bd}}(\gamma_P)\cdot|A'|_{\infty}\le C_0\cdot|A'|_{\infty}$, by the definition of $C_0$ in (C6):
\[
C_0=\max\left\{C_P\cdot c_{\mathrm{bd}}(\gamma_P),\ C^{F},\ 4b_++16C^{F}\cdot c(\gamma^{F})\right\}.
\]
It follows, by (T3) (Proposition~35.1\,(e)), that
\[
|\mathscr{A}^0|_{\infty}\le C_0\cdot|A^{g}|_{\infty}\le2C_0\cdot\ell_*^2\cdot p_*<\frac{1}{2}e_*^{-\varepsilon}.
\]

\emph{The domain $\mathrm{Dom}_*$.}
By \cite[Proposition~15.30(i)]{QEDI}, applied at the level $k_*$, with the side $\ell_*\in\mathrm{Cl}$ and the charge $e_*$, the set $\left\{A'\text{ real on }T^0_*:\ |dA'|_{\infty}\le p_*\right\}$ is contained in $\mathrm{Dom}_*$.
The hypotheses of that proposition hold.
Indeed, $\varepsilon\in(0,\bar{\varepsilon}]$ with $\bar{\varepsilon}=1/200$ by (C7) together with (C1) (Proposition~14.2\,(b), row (S1)), and $L^{\frac{3}{4}}\ge\frac{8}{3}$ because Proposition~14.2\,(f) gives $L^{\frac{3}{4}}\ge110$.
Since $v_*\ge-\log e_0$, the inequalities $8C_B\cdot p_*\le e_*^{-\frac{3}{4}+3\varepsilon}$ and $8e_*^{\frac{3}{4}-\frac{7\varepsilon}{2}}\le1$ are the third and the fourth inequality of (A10); $e_*<1$ by (Eq4.2); (A1) and (A2) are rows of Proposition~14.2\,(b); and $C_w\ge C_P$, $\gamma_{\mathrm{w}}\le\gamma_P$ hold by Proposition~14.2\,(f).
Let $\theta\in\mathbb{R}^3$ and $\mathrm{sv}\in[0,1]$.
The fields $A+c_\theta$ and $\mathrm{sv}\cdot A^{g}+c_\theta$ are real, and
\[
d\left(A+c_\theta\right)=dA,\qquad d\left(\mathrm{sv}\cdot A^{g}+c_\theta\right)=\mathrm{sv}\cdot dA,
\]
because $dd\boldsymbol{\lambda}=0$ and $d$ annihilates constants; so the coboundaries of both fields have sup norm at most $p_*$.
Hence both fields lie in $\mathrm{Dom}_*$.
By definition (Section~I.5.3, that is, \cite[Definition~5.16]{QEDI}),
\[
\mathrm{Dom}_*=\left\{A'\text{ real on the bonds of }T^0_*:\ H_*[A']\in\tfrac{1}{4}\tilde{R}_*\right\};
\]
so $H_*\left[\mathrm{sv}\cdot A^{g}+c_\theta\right]=\mathrm{sv}\cdot\mathscr{A}^0+c_\theta$ lies in $\frac{1}{4}\tilde{R}_*$.
This proves (2.3).

\paragraph{4.6. (SF6)}
\textbf{(2.4).}
Let $A\in\mathrm{Dom}_*$ and $\mathscr{C}\in\mathfrak{X}_*$.
The field $\mathscr{A}=H_*[A]$ lies in $\frac{1}{4}\tilde{R}_*$, a subset of $\tilde{R}_*$, because the thresholds of $\mathrm{cf}\,\tilde{R}_*$ increase with $\mathrm{cf}$ (Section~I.5.3, that is, \cite[Definition~5.16]{QEDI}).
Hence $\mathscr{T}=\mathscr{A}+\mathscr{C}$ belongs to $\tilde{R}_*+\mathfrak{X}_*$, which is contained in $\tilde{G}_*$, and so does $-\mathscr{T}$ (SF2).
Therefore $H_*(\mathscr{T})$, $\bar{H}_*(\mathscr{T})=\mathsf{C}^{-1}H_*(-\mathscr{T})\,\mathsf{C}$ and $D_*(\mathscr{T})$ are defined (SF2), the kernels of $\mathscr{E}_*$ are defined at $(\mathscr{A},\mathscr{C})$ (SF3), and $\det D_*(0)\ne0$ (4.4\,(vi)).
By Section~G5.1.5 and Section~I.5.7, with $n_{k_*}=\ell_*^3$,
\begin{align*}
Y^{x}_*(A;\mathscr{C})&=y^{x}_{k_*}\left[\varepsilon^{\mathrm{v},t}_*,m_*,\mathscr{E}_*\right](A;\mathscr{C})=\int\mathrm{D}\Psi\,\mathscr{F}^{x}_{k_*}\left[\varepsilon^{\mathrm{v},t}_*,m_*,\mathscr{E}_*\right](A,\mathscr{C};\Psi),\\
\mathscr{F}^{x}_{k_*}\left[\varepsilon^{\mathrm{v},t}_*,m_*,\mathscr{E}_*\right](A,\mathscr{C};\Psi)&=\det D_*(0)^{-1}\cdot\exp\left\{-\left\langle\bar{\Psi},D_*(\mathscr{T})\Psi\right\rangle-m_*\cdot\left\langle\bar{\psi},\psi\right\rangle+\sum_{X\in D_*}\mathscr{E}_*(X;\mathscr{A},\mathscr{C},\psi,\chi)\right.\\
&\qquad\left.-\varepsilon^{\mathrm{v},t}_*\cdot\ell_*^3+\sum_{X\in D_*}\mathfrak{P}_*(X,\mathscr{C})\right\},
\end{align*}
where the fields $\psi=H_*(\mathscr{T})\Psi$, $\chi=\delta_{\alpha,\mathscr{T}}\psi$, $\bar{\psi}=\bar{H}_*(\mathscr{T})\bar{\Psi}$, $\bar{\chi}=\delta_{\alpha,-\mathscr{T}}\bar{\psi}$ are inserted (Section~I.5.4).
The quantity $-\varepsilon^{\mathrm{v},t}_*\cdot\ell_*^3+\sum_X\mathfrak{P}_*(X,\mathscr{C})$ is a complex number, so its exponential is a scalar factor of the even element, and $\int\mathrm{D}\Psi$ is linear.
This gives (2.4).
The forms are those of Section~I.3.6:
\begin{gather*}
\left\langle\bar{\psi},\psi\right\rangle=\int\mathrm{d}x\,\bar{\psi}(x)\cdot\psi(x)=\eta_*^3\cdot\frac{1}{48}\sum_{c\in T_*}\ \sum_{g\in O_{\mathrm{cub}}}\ \sum_{\beta}\bar{\psi}_\beta(c,g)\,\psi_\beta(c,g),\\
\left\langle\bar{\Psi},D_*(\mathscr{T})\Psi\right\rangle=\sum_{y,y'\in T^0_*}\bar{\Psi}(y)\cdot D_*(\mathscr{T})(y,y')\,\Psi(y').
\end{gather*}
The integral $\int\mathrm{D}\Psi$ is the integral (I.3.5) with $\nu=1$ over the generators $\Psi_\beta(y)$, $\bar{\Psi}_\beta(y)$, $y\in T^0_*$, $\beta\in\{1,\dots,4\}$, that is, over $4\ell_*^3$ pairs.
Since $k_*\ge1$, the generators $\Psi$, $\bar{\Psi}$ have four components per unit site and the weight $\nu=1$; the weight $1/48$ appears only in the forms and kernels on $T^{\mathrm{f}}_*$.

\textbf{Gauge invariance; (2.5).}
The pair $P_N(e)$ is real and controlled on $R^*(e)$ by (SF1).
So the data of Proposition~29 are given, and the function $Y^{x}_*$ of Proposition~29 is the function just described.
Let $a$ be a complex one-cochain on $\mathbb{Z}^3/n_*\mathbb{Z}^3$ such that $a^{[j]}\in\mathfrak{X}_j$ for all $j\le k_*$.
By Proposition~29\,(b),
\[
Y^{x}_*\left(A+d\boldsymbol{\lambda};a^{[k_*]}\right)=Y^{x}_*\left(A;a^{[k_*]}\right)
\]
whenever $A$ and $A+d\boldsymbol{\lambda}$ lie in $\mathrm{Dom}_*$ and $\boldsymbol{\lambda}$ is a real function on $T^0_*$.
By Proposition~29\,(c), for every $\theta\in\mathbb{R}^3$ and every bounded Borel function $\Phi$ on ${V_{\mathrm{t}}}(T^0_*)$,
\[
\nu^{\theta,a}\left(\mathbf{1}_{\mathfrak{G}}\,\Phi(A_*)\right)=\int\mathrm{d}\mathsf{m}_*(A)\,\Phi(A)\cdot\chi_*(A)\cdot Y^{x}_*\left(A+c_\theta;a^{[k_*]}\right).
\]
This is (2.5).
Here $A+c_\theta$ lies in $\mathrm{Dom}_*$ wherever $\chi_*(A)=1$, by (SF5).

\paragraph{4.7. (DF5)}
\textbf{The scaling relations.}
Let $a$ be a complex one-cochain on $\mathbb{Z}^3/n_*\mathbb{Z}^3$, let $0\le k\le k_*$ and put $i:=k_*-k$.
No bound on $a$ is needed here.
By Section~I.6.6, for every bond $b$ of $T_k$,
\[
a^{[k]}(b)=L^{\frac{k}{2}}\cdot a\left(L^kb\right)=L^{-\frac{i}{2}}\cdot L^{\frac{k_*}{2}}\cdot a\left(L^{k_*}\left(b/L^i\right)\right)=L^{-\frac{i}{2}}\cdot a^{[k_*]}\left(b/L^i\right),
\]
and $b\mapsto b/L^i$ is a bijection of the bonds of $T_k=\mathbb{L}_k/\ell_k\mathbb{Z}^3$ onto the bonds of $T_*$, since $\eta_k=L^i\cdot\eta_*$ and $\ell_k=L^i\cdot\ell_*$.
Thus $|a^{[k]}|_{\infty}=L^{-\frac{i}{2}}\cdot|a^{[k_*]}|_{\infty}$.
For all $\mu$, $\nu$ and every site $c$ of $T_k$,
\begin{align*}
\left(\partial_\mu a^{[k]}_\nu\right)(c)&=\eta_k^{-1}\cdot\left(a^{[k]}_\nu\left(c+\eta_k\cdot e_\mu\right)-a^{[k]}_\nu(c)\right)\\
&=L^{-i}\cdot\eta_*^{-1}\cdot L^{-\frac{i}{2}}\cdot\left(a^{[k_*]}_\nu\left(c/L^i+\eta_*\cdot e_\mu\right)-a^{[k_*]}_\nu\left(c/L^i\right)\right)=L^{-\frac{3i}{2}}\cdot\left(\partial_\mu a^{[k_*]}_\nu\right)\left(c/L^i\right).
\end{align*}
Thus $|\partial a^{[k]}|_{\infty}=L^{-\frac{3i}{2}}\cdot|\partial a^{[k_*]}|_{\infty}$.
For $k<k_*$, Section~I.3.7, where $A_L(b)=L^{-\frac{1}{2}}\cdot A(b/L)$, gives
\[
\left(a^{[k+1]}\right)_L(b)=L^{-\frac{1}{2}}\cdot L^{\frac{k+1}{2}}\cdot a\left(L^{k+1}b/L\right)=L^{\frac{k}{2}}\cdot a\left(L^kb\right)=a^{[k]}(b).
\]

\textbf{The domains $\mathfrak{X}_k$.}
Assume $|a|_{\mathrm{ext}}<e_*^{-1+3\varepsilon}$.
Two of the three summands of $|a|_{\mathrm{ext}}$ are $|a^{[k_*]}|_{\infty}$ and $|\partial a^{[k_*]}|_{\infty}$.
Hence
\[
|a^{[k]}|_{\infty}\le|a^{[k_*]}|_{\infty}<e_*^{-1+3\varepsilon}\le e_*^{-1+2\varepsilon},\qquad|\partial a^{[k]}|_{\infty}\le|\partial a^{[k_*]}|_{\infty}<e_*^{-1+3\varepsilon},
\]
where $e_*<1$ was used.
By (Eq4.1) and (Eq4.2), $e_k\le e_*<1$, and the exponents $-1+2\varepsilon$, $-1+3\varepsilon$ are negative because $\varepsilon\le1/200$; so $e_*^{-1+2\varepsilon}\le e_k^{-1+2\varepsilon}$ and $e_*^{-1+3\varepsilon}\le e_k^{-1+3\varepsilon}$.
The right sides are the thresholds of $\mathfrak{X}_k$ (Section~I.5.3), and therefore $a^{[k]}\in\mathfrak{X}_k$.
At $k=0$ one has $a^{[0]}=a$, and $\mathfrak{X}_0$ is formed with $q=e_*\cdot L^{-\frac{k_*}{2}}$.

\textbf{Probes.}
Let $\beta$ be a probe, $v$ an admitted placement and $t\in\mathrm{DDo}(t_x^{+})$.
By Proposition~35.3\,(f), $|a(t)|_{\mathrm{ext}}<e_*^{-1+3\varepsilon}$.
Consequently $a(t)^{[j]}\in\mathfrak{X}_j$ for all $j\le k_*$, and
\[
e_*\cdot|\mathscr{C}(t)|_{\infty}=e_*\cdot\left|a(t)^{[k_*]}\right|_{\infty}\le e_*\cdot|a(t)|_{\mathrm{ext}}<e_*^{3\varepsilon}.
\]

\paragraph{4.8. (A6)}
By 4.0 the data $(k_*,\ell_*,\mu_*,e_*)$ are (S-3).
In Proposition~20.1 one has $\mathrm{aw}_2=C_0^{-1}\cdot e_*^{-\frac{1}{4}+\varepsilon}=C_0^{-1}\cdot h_{*,2}=\mathrm{aw}$.
For $A\in\tilde{G}_*$ and $\psi=H_*(A)\Psi$, $\bar{\psi}=\bar{H}_*(A)\bar{\Psi}$, Proposition~20.1\,(b) gives
\[
\left|\left\langle\bar{\psi},\psi\right\rangle\right|_{\mathrm{aw}}\le4e_*^{-\frac{1}{2}+2\varepsilon}\cdot\ell_*^3;
\]
here $\langle\bar{\psi},\psi\rangle$ is the form displayed in 4.6, in which the sum extends over the photon sites of $T_*$, the $48$ labels and the four components, with the weight $\eta_*^3/48$.
Proposition~20.1\,(c), applied to $\mathscr{E}=\mathscr{E}_*$, which lies in $K^{x}_*\subset K^{x,C}_*$ by (SF3), to $(\mathscr{A},\mathscr{C})\in\tilde{R}_*\times\mathfrak{X}_*$ and to $X\in D_*$, gives
\[
\left|\mathscr{E}_*\left(X;\mathscr{A},\mathscr{C},H_*^{\#}(\mathscr{T})\Psi\right)\right|_{\mathrm{aw}}\le\left\|\mathscr{E}_*(X;\mathscr{A},\mathscr{C})\right\|_{h_*},
\]
and the right side is the left side of (2.1) (see 4.3).

\paragraph{4.9. (A7)}
Proposition~19 is applied with the following data: the side $\ell_*\in\mathrm{Cl}$, the number $e_*>0$, the mass $\mu_*>0$, the level $k_*\ge1$, and a real bond function $A$ on $T_*$, on which no bound is imposed.
It yields the following.
\begin{itemize}
\item[(a)] $\operatorname{Re}\left(f,(D_A+\mu_*)f\right)^{\mathrm{ch}}=\mu_*\cdot(f,f)^{\mathrm{ch}}+\frac{\rho}{2\eta_*}\cdot\left(f,\nabla_A^{*}\nabla_Af\right)^{\mathrm{ch}}\ge\mu_*\cdot(f,f)^{\mathrm{ch}}$.
\item[(b)] $\left(Q_*(A)J^{*}\right)^{*}=J\,Q_*^{T}(-A)$.
\item[(c)] The operator $\mathrm{Fst}(A)=D_A+\mu_*+b_*\cdot J\,P_*(A)\,J^{*}$ has real part at least $\mu_*$ and is invertible; hence $S_*(A)$ and $D_*(A)$ are defined.
\item[(d)] $D_*(A)$ is invertible, and $|D_*(A)^{-1}|_{\mathrm{op}}\le b_*^{-1}+\mu_*^{-1}$, where the operator norm is the one for the counting scalar product on $\mathfrak{S}(1,\ell_*)$, that is, for $|\cdot|_{\ell^2}$.
\item[(e)] Since $\mu_*\le1/L\le12\mathrm{e}$ (Proposition~35.1\,(a)),
\[
\left|D_*(A)^{-1}(y,y')\right|_2\le b_*^{-1}\cdot\mathbf{1}[y=y']+\frac{2\mathrm{e}}{\mu_*}\cdot\exp\left(-\frac{\mu_*}{12\mathrm{e}}\cdot d(y,y')\right).
\]
\end{itemize}
This completes the proof of Proposition~35.4.\hfill$\square$

\paragraph{4.10. Uniformity}
The constants that occur in the bounds are $C^{F}$, $\gamma^{F}$, $\mathbf{S}$, $C_d$, $K_0$, $c_{\mathrm{rb}}$, $\kappa$, $\kappa_0$, $C_0$, $M$, $L$, the exponent $\varepsilon$, the integers $r$ and $p$ (through $\ell_*=M\cdot r_*$ and $p_*=v_*^{p}$), and the explicit numerals.
They belong to the lines (C0)--(C7), and none of them depends on $\bar{m}$, $e$, $N$, the level, the side, $\theta$, $A$, $\mathscr{A}$, $\mathscr{C}$, $a$, the pair, the probe or the placement ((I.U4)$^{\mathrm{ch}}$ in Section~I.8.1; Proposition~14.1\,(iv)).
The propositions applied above hold for all data (SF), (OS), (FL) of the kinds named; by 4.0 such data are given for every $\bar{m}\in(0,1]$, every $e$ with $e\cdot h^{\frac{1}{2}}\le e_1$ and every $N\ge N_0(e)$.
The bounds of (SF2) hold on all of $\tilde{G}_*$, respectively on all of $\mathrm{Gwk}$ and of $\mathrm{Gwk}^{+}$, as stated there, and the first bound of (A6) holds on all of $\tilde{G}_*$; the bound (2.1) and the second bound of (A6) hold on all of $\tilde{R}_*\times\mathfrak{X}_*$.
The sets $\tilde{G}_*$ and $\tilde{R}_*$ are invariant under the addition of real constants (SF2), and hence so are $\mathrm{Gwk}$ and $\mathrm{Gwk}^{+}$ (Section~G5.1.7: the constant is added to $A_r$).
The twist connection $c_\theta$ is a real constant, and it is counted with $\mathscr{A}$, in the first slot (Section~I.6.6; Proposition~14.1\,(iv)).
Therefore these bounds are uniform in the twist.
(A7) holds for every real bond function.
The numbers $e_*$, $\ell_*$, $\mu_*$, $p_*$ that enter the bounds do not depend on $N$ (Proposition~35.1\,(c)), and $b_*\in[b_-,1]$, $\mu_*\in(\mu_-,1/L]$ for all $N$ (Proposition~35.1\,(a)).

\subsection*{5. Conditions on the constants}
The proofs of Propositions~35.1--35.4 use the following conditions on the constants (the lines are those of Section~I.8.2).
\begin{description}
\item[The number $e_1$.]
$e_1$ is one half of $\sup\mathrm{Se}$, and $e_1\le\frac{1}{2}\min\left(e_0,\frac{1}{3}\right)$.
\emph{Stands in:} Proposition~14.2\,(a), line (C9).
\emph{Enters:} the proof of Proposition~35.1 (Step~1, and (a)); Proposition~35.3\,(f) ($e_*\le\frac{1}{3}$).
\item[Row (e1).]
$v\ge8$; $v^{r}>31L+16$; $\left(1+3\,\mathrm{lg}/v\right)^{r}<L$; $p_0\cdot v^{p_0-1}\ge28$; for all $v\ge-\log e_1$.
\emph{Stands in:} Proposition~14.2\,(b), row (e1).
\emph{Enters:} the proof of Proposition~35.1 (Step~1, and (a), (b), (d), (e), (f)); through Proposition~35.1\,(e) ((T1): $v_*\ge8$), Proposition~35.4\,(SF5) ($p_*=v_*^{p}>0$).
\item[Row (e4).]
$4C_0\cdot\mathrm{lh}(v)^2\cdot v^{p}<\exp(\varepsilon\cdot v)$ for all $v\ge-\log e_1$.
\emph{Stands in:} Proposition~14.2\,(b), row (e4).
\emph{Enters:} Proposition~35.1\,(e) ((T3)); through it, Proposition~35.4\,(SF5).
\item[The number $M$.]
$M$ is a power of $L$, and $M\ge L^2$.
\emph{Stands in:} (C5); Proposition~14.2\,(b), row (M1).
\emph{Enters:} the proof of Proposition~35.1 (Step~1, and (f)); the proof of Proposition~35.3, (Eq3.5); \cite[Section~17.4.1, (g0)]{QEDI}.
\item[The number $L$.]
$L$ is odd and $L\ge13$; hence $L\ge5$ and $L>1$.
Moreover $L^{\frac{3}{4}}\ge110$; hence $L^{\frac{3}{4}}\ge\frac{8}{3}$, that is, $L\ge4$.
\emph{Stands in:} (C3); Proposition~14.2\,(b), rows (L1), (L5); Proposition~14.2\,(f) ($L\ge2400^8$; $L^{\frac{3}{4}}\ge110$).
\emph{Enters:} \cite[Proposition~19.1(f)]{QEDI} ($L$ odd, $L\ge13$); the proof of Proposition~35.3, (Eq3.5) ($L\ge5$), and (g) ($L$ odd: every site of $T_*$ lies in exactly one block $\Delta_y$, Section~I.1.1); \cite[Lemma~17.4]{QEDI} ($L\ge5$) and \cite[Section~17.4.1, (g0)]{QEDI} ($L$ odd); \cite[Theorem~13.2, (t5), (t8)]{QEDI} in Proposition~35.4\,(SF4)\,(v) ($L$ odd); \cite[Section~14.5]{QEDI} and \cite[Propositions~15.5, 15.21(ii) and Theorem~15.14]{QEDI} in Proposition~35.4\,(SF5) ($L$ odd, $L\ge3$); \cite[Proposition~15.30(i)]{QEDI} in Proposition~35.4\,(SF5) ($L^{\frac{3}{4}}\ge\frac{8}{3}$); through $L>1$: the proof of Proposition~35.1 (Step~0, (a), (b), (d), (f)), Proposition~35.4\,(SF4) (the powers $L^{-1+\varepsilon}$ and $L^{-2+\frac{3\varepsilon}{2}}$ in the scaling of $\mathscr{C}_L$), (DF5) ($L^{-\frac{i}{2}}\le1$ and $L^{-\frac{3i}{2}}\le1$) and (A7)\,(e) ($1/L\le12\mathrm{e}$).
\item[The number $r$.]
$r$ is an integer, $r\ge2$.
\emph{Stands in:} (C7).
\emph{Enters:} Proposition~35.1\,(f), through $r\ge1$ ($v_K^r\ge v_K$ for $v_K\ge1$); the proof of Proposition~35.1, Step~0 ((Pr-a), (Pr-b), (Pr-g)), Step~1 and (b), through $r\ge0$ (the map $t\mapsto t^r$ is non-decreasing on $(0,\infty)$, and $t^r\ge1$ for $t\ge1$).
\item[The number $b$.]
$b=1$.
\emph{Stands in:} (C0).
\emph{Enters:} Proposition~35.1\,(a).
\item[The number $\alpha$.]
$0\le\alpha<1$.
\emph{Stands in:} (C1) ($\alpha\in(1-\bar{\varepsilon},1)$, $\bar{\varepsilon}=1/200$).
\emph{Enters:} Proposition~35.3\,(c1) ($\alpha<1$); \cite[Theorem~15.14]{QEDI} in Proposition~35.4\,(SF5) ($0\le\alpha<1$).
\item[The number $\varepsilon$.]
$0<\varepsilon\le\bar{\varepsilon}$, where $\bar{\varepsilon}=1/200$.
\emph{Stands in:} (C7) ($\varepsilon\in(1-\alpha,\bar{\varepsilon}]$, where $1-\alpha>0$ by (C1)); Proposition~14.2\,(b), row (S1).
\emph{Enters:} through $0<\varepsilon<1/3$: Proposition~35.3\,(f) ($e_*^{3\varepsilon}\le1$) and, in the proof of Proposition~35.4, (Eq4.2) ($e_0^{\varepsilon}\le\frac{1}{4}$ gives $e_0<1$), (SF4) (\cite[Lemma~17.5(c)]{QEDI}: the exponents $-1+\varepsilon$ and $-2+\frac{3\varepsilon}{2}$ are negative) and (DF5) ($e_*^{3\varepsilon}\le e_*^{2\varepsilon}$; the exponents $-1+2\varepsilon$ and $-1+3\varepsilon$ are negative); further \cite[Lemma~17.4]{QEDI}; and, through the whole range $0<\varepsilon\le\bar{\varepsilon}$, \cite[Proposition~15.30]{QEDI} in Proposition~35.4\,(SF5), of which $\varepsilon\in(0,\bar{\varepsilon}]$ is a hypothesis.
\item[The number $\kappa$.]
$\kappa\ge\kappa_0$.
\emph{Stands in:} (C2); Proposition~14.2\,(f).
\emph{Enters:} Proposition~35.4\,(SF4)\,(v), in (t8).
\item[The number $\gamma^{F}$.]
$0<\gamma^{F}$ and $\gamma^{F}\cdot L\le\frac{1}{6}$.
\emph{Stands in:} (C4r); Proposition~14.2\,(f).
\emph{Enters:} Proposition~35.4\,(SF2) (the number $\mathrm{ssum}(\gamma)$, whose bound holds for $0<\gamma\le6$ and is taken at $\gamma=\gamma^{F}/2$, which lies in $(0,6]$).
\item[The number $c_\tau$.]
$c_\tau\le\frac{1}{8}$.
\emph{Stands in:} (C6'), definition of $c_\tau$.
\emph{Enters:} Proposition~35.4\,(SF2) (the set $\mathrm{Gwk}^{+}$).
\item[The number $C_0$.]
$C_0\ge C_P\cdot c_{\mathrm{bd}}(\gamma_P)$.
\emph{Stands in:} (C6), definition of $C_0$.
\emph{Enters:} Proposition~35.4\,(SF5).
\item[The numbers $C_P$, $\gamma_P$.]
$C_P\le C_w$; $\gamma_{\mathrm{w}}\le\gamma_P$; $0<\gamma_P\le1$.
\emph{Stands in:} Proposition~14.2\,(f) ($\gamma_P>0$: paragraph \emph{Ranges}).
\emph{Enters:} Proposition~35.4\,(SF5): the hypotheses $C_P\le C_w$, $\gamma_{\mathrm{w}}\le\gamma_P$ of \cite[Proposition~15.30]{QEDI}, and the hypothesis $0<\gamma_P\le1$ of \cite[Lemma~15.2]{QEDI}.
\item[Row (A8).]
$\exp(-\varepsilon\cdot v)\le\frac{1}{4}$; $3M\cdot\exp(-\varepsilon\cdot v)\le1$; $\exp\left(-\frac{v}{4}\right)+\exp(-\varepsilon\cdot v)\le1$; for all $v\ge-\log e_0$.
\emph{Stands in:} Proposition~14.2\,(b), row (A8).
\emph{Enters:} the proof of Proposition~35.4: (Eq4.2), (SF2) ($\mathfrak{i}$ on $\mathrm{Gwk}$), (DF5) ($e_*<1$); \cite[Lemma~17.4(a)]{QEDI}.
\item[Rows (A1), (A2).]
The conditions (A1) and (A2).
\emph{Stand in:} Proposition~14.2\,(b), rows (A1), (A2).
\emph{Enter:} Proposition~35.4\,(SF5), as hypotheses of \cite[Proposition~15.30]{QEDI}.
\item[Row (A10), third and fourth inequality.]
$8C_B\cdot v^{p}\le\exp\left(\left(\frac{3}{4}-3\varepsilon\right)\cdot v\right)$; $8\exp\left(-\left(\frac{3}{4}-\frac{7\varepsilon}{2}\right)\cdot v\right)\le1$.
\emph{Stands in:} Proposition~14.2\,(b), row (A10).
\emph{Enters:} Proposition~35.4\,(SF5), as hypotheses of \cite[Proposition~15.30]{QEDI}.
\end{description}

\section*{Chapter 33. The fermionic operators at the stop level as functions of the connection: proofs of Propositions 33.1--33.4}

\subsection*{0. Conventions of this chapter; setting, preliminaries}

\subsubsection*{0.0. Letters, labels, displays}

In this chapter $\operatorname{adj}O$ denotes the adjugate of a square matrix $O$.
Every other letter that is particular to the chapter is defined at the place where it first occurs.
The letter $\mathrm{Rop}$ has the meaning of the operator $\mathbf{R}$ in the proof of \cite[Theorem~13.18]{QEDI}.
In Step~3 of the proof of Proposition~33.4\,(b), and only there, $V_1$, $V_2$, $Q$, $Q'$, $b$ are the letters of the Woodbury identity in the proof of \cite[Theorem~13.18]{QEDI}, with the meaning they have there.
In (Bk3), (Bk4) and in the proof of Proposition~33.1\,(b), $a$ is a vector of $\{\pm e_\mu\}$, as in (I.s4).
The labels (M1), (M2) designate the two facts on majorants of Section~G5.1.7;
they are not to be confused with the conditions (M1)--(M6) of Proposition~14.2\,(b).
The displays numbered (0.n), (1.n), (2.n), (3.n) and the display (4.4) belong to this chapter;
(4.1), (4.2), (4.3) are the displays of Propositions~33.1 and~33.2.

\subsubsection*{0.1. Setting}

The setting is (ST*)$^{\mathrm{ch}}$ of Section~G5.1.1, with the notation of Section~G5.1.7 and the constants of Section~G5.1.8.
All objects live at the level $k_*$:
the fine photon torus $T_*=\eta_*\mathbb{Z}^3/\ell_*\mathbb{Z}^3$,
the set of fermion sites $T^{\mathrm{f}}_*=T_*\times O_{\mathrm{cub}}$, with $48$ sites $x=(c,g)$ over every photon site $c$,
and the unit lattice $T^0_*=\mathbb{Z}^3/\ell_*\mathbb{Z}^3$.
The side $\ell_*$ lies in $\mathrm{Cl}$ (Proposition~35.1\,(d), first sentence; $\mathrm{Cl}$ is defined in Section~I.1.1), so that $\ell_*\ge13\,L\cdot M>2M+2$.
Proposition~35.1\,(a), together with $L\ge5$ (Section~I.0.5) for the two inequalities $\frac{1}{L}<1$, gives
\begin{equation}\tag{Dat}
k_*\ge1;\qquad\eta_*=L^{-k_*}\le\frac{1}{L}<1;\qquad b_-\le b_*\le1;\qquad\mu_-<\mu_*\le\frac{1}{L}<1.
\end{equation}

Since $k_*\ge1$, the level $0$ does not occur.
The block spinors are the elements of $\mathfrak{S}(1,\ell_*)$, with $4$ components at every unit site and the weight $1$;
the fine spinors are the elements of $\mathfrak{S}^{\mathrm{ch}}(\eta_*,\ell_*)$, with $192=48\cdot4$ components at every photon site,
and they are integrated with the measure $\int\mathrm{d}x=\eta_*^3\cdot\frac{1}{48}\sum_c\sum_{g\in O_{\mathrm{cub}}}$ of (I.1.3);
the averaging operators are those of the first line of (I.3.4).
The weight $\sqrt{48}$ of the integration at level $0$ (Section~I.5.2) does not enter.

By Section~I.0.5 and Section~G5.1.8 (see also Section~I.8.2 and Proposition~14.2\,(f)),
\begin{gather*}
0<\rho<1;\qquad0<\alpha<1;\qquad L\text{ is an odd integer }\ge5;\qquad M\ge1\text{ is an integer};\\
C^{F}\ge1;\qquad0<\gamma^{F}\le1;\qquad b_-=1-\frac{1}{L}>0;\qquad\mu_-=L^{-3}\le1.
\end{gather*}

In this chapter a bond function is a complex bond function on $T_*$.
The bond functions form a complex vector space of dimension $3\left(\frac{\ell_*}{\eta_*}\right)^3$,
and the words ``open'' and ``analytic'' refer to this space.
All spaces of spinors and of functions of pairs that occur have finite dimension.

\begin{description}
\item[(Cv1)]
For a unit site $y$ and sets $Z$, $Z'$ of unit sites of $T^0_*$, the distances $d(y,Z)=d(Z,y)$ and $d(Z,Z')$, together with the conventions on empty sets and empty sums, are those of Section~G5.1.4.
For an index $u=(y',\beta',\omega')$ the distances $d(Z,u)=d(u,Z)=d(y',Z)$ are those of Section~G5.1.6.
The set $\mathrm{Ypr}_B$ is empty if and only if $B=0$, and then $\mathrm{Ypr}_B^{+}$ is empty as well.
\item[(Cv2)]
For an analytic function $O(\cdot)$ of the connection with values in the operators, the second derivative $\partial_B\partial_{B'}O$ is the iterated derivative of Section~G5.1.7;
it is symmetric in $(B,B')$.
\item[(Cv3)]
For a linear operator $O$ with values in the functions of pairs, the majorant of $O$ for the short pairs is the one defined in Section~G5.1.7.
With this definition $\mathrm{Sa}_*(A)$ is the restriction of $\delta_{\alpha,A}S_*(A)$ to the short pairs,
and the restriction of $\delta_{\alpha,A}S_*(A)\cdot O'$ to the short pairs is $\mathrm{Sa}_*(A)\cdot O'$, for every operator $O'$ with values in $\mathfrak{S}^{\mathrm{ch}}(\eta_*,\ell_*)$.
\item[(Cv4)]
The pairs inside a cluster, $x=(c,g)$ and $x'=(c,g')$ with $g\ne g'$ (there are $48\cdot47$ ordered pairs of this kind over every photon site), satisfy $d_X(x,x')=\max(0,\eta_*)=\eta_*$ by (I.1.5), and $\eta_*<1$ by (Dat).
Hence they are short pairs, $[x]=[x']$, and
\[
\left(\delta_{\alpha,A}f\right)(x,x')=\eta_*^{-\alpha}\cdot\left(f(c,g')-f(c,g)\right)
\]
contains no phase (Section~I.1.8).
\end{description}

\subsubsection*{0.2. Numbers}

\begin{description}
\item[(Nu0)]
$2<\mathrm{e}<2.7183$.
Indeed, $\mathrm{e}=\sum_{n\ge0}\frac{1}{n!}>1+1+\frac{1}{2}$.
For $n\ge9$ one has $n!\ge9!\cdot10^{n-9}$; therefore
\[
\sum_{n\ge9}\frac{1}{n!}\le\frac{1}{9!}\sum_{j\ge0}10^{-j}=\frac{10}{9}\cdot\frac{1}{9!}<\frac{2}{9!},
\]
and
\[
\mathrm{e}<\sum_{n\le8}\frac{1}{n!}+\frac{2}{9!}=\frac{109601}{40320}+\frac{1}{181440}=\frac{986411}{362880}<2.7183,
\]
because $2.7183\cdot362880=986416.704>986411$.
\end{description}
It follows that:
\begin{description}
\item[(Nu1)]
$\exp\left(\frac{3}{8}\right)\le\frac{3}{2}$, since $\mathrm{e}^3<20.09<25.62<\frac{6561}{256}=\left(\frac{3}{2}\right)^8$;
\item[(Nu2)]
$\exp\left(\frac{1}{4}\right)\le\frac{4}{3}$, since $\mathrm{e}<3.16<\frac{256}{81}=\left(\frac{4}{3}\right)^4$;
\item[(Nu3)]
$15+48\,\mathrm{e}<15+130.48<146$;
\item[(Nu4)]
$4\,\mathrm{e}<10.88<11$;
\item[(Nu5)]
$\mathrm{e}<3$;
\item[(Nu6)]
$\exp\left(\frac{1}{16}\right)\le2^{\frac{1}{2}}$, since $\mathrm{e}<2^8$, and $\exp\left(\frac{1}{24}\right)\le2$, since $\mathrm{e}<2^{24}$;
\item[(Nu7)]
$\left|\exp(\mathrm{tq})-1\right|\le|\mathrm{tq}|\cdot\exp|\mathrm{tq}|$ for every complex number $\mathrm{tq}$, because
\[
\left|\sum_{n\ge1}\frac{\mathrm{tq}^n}{n!}\right|\le|\mathrm{tq}|\cdot\sum_{n\ge0}\frac{|\mathrm{tq}|^n}{(n+1)!}\le|\mathrm{tq}|\cdot\exp|\mathrm{tq}|.
\]
\end{description}

\subsubsection*{0.3. Blocks (Sections~I.1.1, I.1.4)}

\begin{description}
\item[(Bk1)]
The number $\eta_*^{-1}=L^{k_*}$ is an odd integer.
For $y\in\mathbb{Z}^3$,
\[
\Delta_y=\left\{c\in\mathbb{L}_*:|c-y|<\frac{1}{2}\right\}=y+\eta_*\cdot\left\{-\frac{\eta_*^{-1}-1}{2},\dots,\frac{\eta_*^{-1}-1}{2}\right\}^3.
\]
These sets partition $\mathbb{L}_*$, every $\Delta_y$ has $\eta_*^{-3}$ points, and $\bigl|c-[c]\bigr|\le\frac{1-\eta_*}{2}$.
Consequently $\Delta_y\times O_{\mathrm{cub}}$ consists of $48\,\eta_*^{-3}$ sites, each of the weight $\frac{\eta_*^3}{48}$, and
\[
\int_{\Delta_y}\mathrm{d}x=\eta_*^3\cdot\frac{1}{48}\sum_{c\in\Delta_y}\sum_{g\in O_{\mathrm{cub}}}1=1.
\]
\item[(Bk2)]
For $c,c'\in\mathbb{L}_*$,
\[
\bigl|[c]-[c']\bigr|\le\bigl|[c]-c\bigr|+|c-c'|+\bigl|c'-[c']\bigr|\le|c-c'|+1-\eta_*<|c-c'|+1,
\]
and the left side is an integer.
Hence $\bigl|[c]-[c']\bigr|\le M$ if $|c-c'|\le M$, and $\bigl|[c]-[c']\bigr|\le1$ if $|c-c'|<1$.
If $[c]=[c']$, then $|c-c'|\le1-\eta_*<1$.
Therefore $[c]\ne[c']$ whenever $|c-c'|\ge1$;
if moreover $|c-c'|\le M$, then $0<\bigl|[c]-[c']\bigr|\le M<\ell_*$, and the classes of $[c]$ and $[c']$ in $T^0_*$ are different.
\item[(Bk3)]
If $c'=c+\eta_*\cdot a$ with $a\in\{\pm e_\mu\}$, then $[c']=[c]$ or $[c']=[c]+a$.
Indeed, $c-[c]=\eta_*\cdot n$ with integers $n_\nu$ such that $|n_\nu|\le\frac{\eta_*^{-1}-1}{2}$.
If $|n_\mu+a_\mu|\le\frac{\eta_*^{-1}-1}{2}$, then $c'$ lies in $\Delta_{[c]}$.
Otherwise $n_\mu+a_\mu=a_\mu\cdot\frac{\eta_*^{-1}+1}{2}$;
then the coordinate $\mu$ of $c'-([c]+a)$ is $a_\mu\cdot\frac{1+\eta_*}{2}-a_\mu=-a_\mu\cdot\frac{1-\eta_*}{2}$, while its two other coordinates are those of $c-[c]$;
so $c'$ lies in $\Delta_{[c]+a}$.
\item[(Bk4)]
The distance $d$ of two points of $T^0_*$ is at most the sup norm of the difference of any two of their representatives in $\mathbb{Z}^3$.
Two unit sites $z$, $z'$ of $T^0_*$ are nearest neighbours if $z'=z+a$ modulo $\ell_*\mathbb{Z}^3$ for some $a\in\{\pm e_\mu\}$.
A unit site has at most six nearest neighbours, and each of them has the distance $1$ from it,
because $|a+\ell_*\cdot n|\ge\ell_*-1>1$ for $n\in\mathbb{Z}^3$, $n\ne0$, as $\ell_*>2$.
\end{description}

\subsubsection*{0.4. Phases (photon objects: Section~I.1.1, (I.1.6), Section~I.3.2 with (I.3.3); Section~G5.1.7)}

The periodic extension of a bond function to the bonds of $\mathbb{L}_*$ is denoted by the same letter.

\begin{description}
\item[(Ph1)]
Let $y\in\mathbb{Z}^3$ and $c\in\Delta_y$, and let $c=y_0,y_1,\dots,y_{k_*}=y$ be the centres of the blocks of the sides $\eta_*,L\eta_*,\dots,1$ that contain $c$ (Section~I.3.2).
Since $L$ is odd, $\mathrm{Bl}_{j+1}(z)$ is the union of the $L^3$ blocks $\mathrm{Bl}_j(z+L^j\eta_*\cdot n)$, $n\in\left\{-\frac{L-1}{2},\dots,\frac{L-1}{2}\right\}^3$.
Hence $y_j=y_{j+1}+L^j\eta_*\cdot n$ with such an $n$, and
\[
|y_{j+1}-y_j|_1\le\frac{3\,(L-1)\cdot L^j\eta_*}{2}.
\]
Each of the six paths $\boldsymbol{\Gamma}^{\pi}(y_{j+1},y_j)$ has $\frac{|y_{j+1}-y_j|_1}{\eta_*}$ bonds, and the end points of these bonds lie in the lattice box spanned by $y_{j+1}$ and $y_j$.
This box lies in $\mathrm{Bl}_{j+1}(y_{j+1})$, a product of three intervals of $\mathbb{L}_*$ that contains both points;
and $\mathrm{Bl}_{j+1}(y_{j+1})$ lies in $\Delta_y=\mathrm{Bl}_{k_*}(y)$, because the blocks that contain $c$ are nested, by the union property.
Therefore
\[
\mathrm{phx}_A(y,c)=e_*\cdot\eta_*\cdot\sum_{j<k_*}\frac{1}{6}\sum_\pi A\left(\boldsymbol{\Gamma}^{\pi}(y_{j+1},y_j)\right)
\]
is a linear combination, with real coefficients, of the numbers $A(b)$, where $b$ ranges over the bonds with both end points in $\Delta_y$;
and the moduli of the coefficients add up to at most
\[
e_*\cdot\eta_*\cdot\sum_{j<k_*}\frac{3\,(L-1)\cdot L^j}{2}=\frac{3}{2}\,e_*\cdot\eta_*\cdot\left(\eta_*^{-1}-1\right)\le\frac{3}{2}\,e_*.
\]
Hence
\begin{gather*}
\left|\mathrm{phx}_A(y,c)\right|\le\frac{3}{2}\,e_*\cdot|A|_\infty,\qquad\operatorname{Im}\mathrm{phx}_A=\mathrm{phx}_{\operatorname{Im}A},\\
\left|\varphi_A(y,c)^{\pm1}\right|=\exp\left(\mp\operatorname{Im}\mathrm{phx}_A(y,c)\right)\le\exp\left(\frac{3}{2}\,\mathfrak{i}(A)\right),\\
\mathrm{phx}_B(y,\cdot)=0\ \text{ if $y$ does not lie in $\mathrm{Ypr}_B$}.
\end{gather*}
For the last assertion: if the class of $y$ does not lie in $\mathrm{Ypr}_B$, then no bond $b$ with $B(b)\ne0$ has an end point in $\Delta_y$.
\item[(Ph2)]
Let $c,c'\in\mathbb{L}_*$ with $|c-c'|\le M$.
Each of the six paths $\boldsymbol{\Gamma}^{\pi}(c,c')$ has $\frac{|c-c'|_1}{\eta_*}\le\frac{3\,|c-c'|}{\eta_*}$ bonds, with end points in the lattice box spanned by $c$ and $c'$.
So
\[
\mathrm{ups}_A(c,c')=e_*\cdot\eta_*\cdot\frac{1}{6}\sum_\pi A\left(\boldsymbol{\Gamma}^{\pi}(c,c')\right)
\]
is a linear combination with real coefficients of the numbers $A(b)$, $b$ in this box, and
\begin{gather*}
|\mathrm{ups}_A(c,c')|\le3\,|c-c'|\cdot e_*\cdot|A|_\infty,\qquad\mathrm{ups}_A(c,c)=0,\\
\operatorname{Im}\mathrm{ups}_A=\mathrm{ups}_{\operatorname{Im}A},\qquad\left|U_A(c,c')^{\pm1}\right|\le\exp\left(3\,|c-c'|\cdot\mathfrak{i}(A)\right).
\end{gather*}
If $\mathrm{ups}_B(c,c')\ne0$, then one of the six paths contains a bond $b$ with $B(b)\ne0$;
an end point $\tilde{c}$ of $b$ lies in the lattice box spanned by $c$ and $c'$, and the class of $[\tilde{c}]$ lies in $\mathrm{Ypr}_B$.
Every point $\tilde{c}$ of this box satisfies $|c-\tilde{c}|\le|c-c'|$ and $|c'-\tilde{c}|\le|c-c'|$,
and any two points $\tilde{c}$, $\tilde{c}'$ of the box satisfy $|\tilde{c}-\tilde{c}'|\le|c-c'|$.
\item[(Ph3)]
For a bond $b$,
\[
\left|u_A(b)^{\pm1}\right|=\exp\left(\mp e_*\cdot\eta_*\cdot\operatorname{Im}A(b)\right)\le\exp\left(\eta_*\cdot\mathfrak{i}(A)\right).
\]
\end{description}
The facts (Ph1)--(Ph3) are the facts on phases recorded in Section~G5.1.7;
compare the definition of the phases in \cite[Section~13.10]{QEDI} and \cite[Lemma~16.2(a)]{QEDI}.

\subsubsection*{0.5. The operators, written out ((I.1.2), (I.2.1), (I.2.3), Section~I.3.3, (I.1.7); Section~G5.1.7)}

Let $f\in\mathfrak{S}^{\mathrm{ch}}(\eta_*,\ell_*)$, $v\in\mathfrak{S}(1,\ell_*)$, $x=(c,g)$, $x'=(c',g')$, and $y=[c]=[x]$ (see (I.POS)).
Then
\begin{equation}\tag{0.1}
\begin{split}
(D_Af)(x)={}&-\eta_*^{-1}\sum_{i=0}^{2}M_i(x)\cdot\left[f(c,gs_i)-f(c,g)\right]\\
&-\eta_*^{-1}\cdot M_3(x)\cdot\left[u_A(c,c_3)\cdot f(c_3,gs_3)-f(c,g)\right],\\
&c_3=c+\eta_*\cdot ge_1,\qquad u_A(c,c_3)=\exp\left(i\cdot e_*\cdot\eta_*\cdot A(c,c_3)\right);
\end{split}
\end{equation}
\begin{equation}\tag{0.2}
\left(Q_*(A)J^{*}f\right)(y)=\eta_*^3\sum_{c'\in\Delta_y}\varphi_A(y,c')\cdot\frac{1}{48}\sum_{g'\in O_{\mathrm{cub}}}f(c',g')=\int_{\Delta_y}\mathrm{d}x'\,\varphi_A(y,c')\cdot f(x');
\end{equation}
\begin{equation}\tag{0.3}
\left(JQ_*^{T}(-A)v\right)(c,g)=\varphi_A(y,c)^{-1}\cdot v(y)\qquad\text{(the same value at the $48$ labels $g$)};
\end{equation}
\begin{equation}\tag{0.4}
\begin{split}
&JP_*(A)J^{*}=\left(JQ_*^{T}(-A)\right)\cdot\left(Q_*(A)J^{*}\right),\\
&\left(JP_*(A)J^{*}f\right)(c,g)=\varphi_A(y,c)^{-1}\int_{\Delta_y}\mathrm{d}x'\,\varphi_A(y,c')\cdot f(x');
\end{split}
\end{equation}
\begin{equation}\tag{0.5}
\begin{split}
&\mathrm{Fst}(A)=D_A+\mu_*+b_*\cdot JP_*(A)J^{*},\\
&\left(\delta_{\alpha,A}f\right)(x,x')=d_X(x,x')^{-\alpha}\cdot\left(U_A(c,c')\cdot f(x')-f(x)\right)\quad\left((x,x')\in\mathrm{Pzt}_*\right),\\
&U_A(c,c')=\exp\left(i\cdot\mathrm{ups}_A(c,c')\right),\qquad U_A(c,c)=1.
\end{split}
\end{equation}
The three hops $i=0,1,2$ stay in the cluster of $c$ and carry no phase;
the hop $i=3$ leads to the cluster of the nearest neighbour $c_3$ and carries the phase of the single bond $(c,c_3)$ ((I.O1), (I.s4)).
Its matrix is $M_3(x)=3\,\bigl(\rho-\gamma\cdot(ge_1)\bigr)$ with $ge_1\in\{\pm e_\mu\}$, and, by (I.m6) and (I.m7),
\begin{equation}\tag{0.6}
|M_3(x)\cdot v|_c\le3\,(1+\rho)\cdot|v|_c\le6\,|v|_c,\qquad|M_3(x)\cdot v|_2\le3\,(1+\rho)\cdot|v|_2\le6\,|v|_2\qquad\left(v\in\mathbb{C}^4\right).
\end{equation}
Indeed, $\gamma\cdot(ge_1)=\pm\gamma_\mu$.
The matrices $\gamma_1$ and $\gamma_2$ have in each row one entry of modulus one, off the diagonal, and $\gamma_3=\mathrm{diag}(1,-1,-1,1)$ (Section~I.2.1, (I.c6)).
So a row of $M_3(x)$ consists of $3\rho$ on the diagonal and one entry of modulus $3$ off it (if $\mu=1,2$), or of the one entry $3\,(\rho\mp1)$ (if $\mu=3$);
its $\ell^1$ norm is at most $3\,(1+\rho)$.
Moreover $M_3(x)$ is Hermitian and $\bigl(\gamma\cdot(ge_1)\bigr)^2=1$ by (I.c1), so that the eigenvalues of $M_3(x)$ are $3\,(\rho\mp1)$.

\subsubsection*{0.6. Majorants}

\begin{description}
\item[(Mj0)]
(Localisation.)
For a unit site $y$ of $T^0_*$ let $\mathbf{1}_y$ be the projection given by
\begin{align*}
(\mathbf{1}_yv)(y'')&:=v(y'')\ \text{ if $y''$ is congruent to $y$ modulo $\ell_*\mathbb{Z}^3$},\quad:=0\ \text{ otherwise}&&\left(v\in\mathfrak{S}(1,\ell_*)\right),\\
(\mathbf{1}_yf)(x)&:=f(x)\ \text{ if $[x]$ is congruent to $y$},\quad:=0\ \text{ otherwise}&&\left(f\in\mathfrak{S}^{\mathrm{ch}}(\eta_*,\ell_*)\right),\\
(\mathbf{1}_yF)(x,x')&:=F(x,x')\ \text{ if $[x]$ is congruent to $y$},\quad:=0\ \text{ otherwise}&&\left(F\text{ a function of pairs}\right).
\end{align*}
In each of the three spaces one has $\sum_y\mathbf{1}_y=1$, $\mathbf{1}_y\cdot\mathbf{1}_y=\mathbf{1}_y$ and $|\mathbf{1}_yf|_\infty\le|f|_\infty$.
The moduli $|f(x)|_c$ and $|F(x,x')|_c$ are $\ell_*\mathbb{Z}^3$-periodic.
By homogeneity, $|\mathbf{1}_y\cdot O\cdot\mathbf{1}_{y'}f|_\infty\le|O|(y,y')\cdot|f|_\infty$ for every $f$.
Moreover $|O_1+O_2|(y,y')\le|O_1|(y,y')+|O_2|(y,y')$, and $|\mathrm{cf}\cdot O|(y,y')=|\mathrm{cf}|\cdot|O|(y,y')$ for complex numbers $\mathrm{cf}$.
Finally $\|O\|_\gamma=0$ only for $O=0$; so $\|\cdot\|_\gamma$ is a norm.
\end{description}

\begin{named}{(M1)}[Products; Section~G5.1.7]
Let $O_1$, $O_2$ be linear operators of the kinds considered in Section~G5.1.7 such that $O_1\cdot O_2$ is defined, and let $\gamma\ge0$.
Then, for all unit sites $y$, $y'$,
\[
|O_1\cdot O_2|(y,y')\le\sum_z|O_1|(y,z)\cdot|O_2|(z,y'),\qquad\|O_1\cdot O_2\|_\gamma\le\|O_1\|_\gamma\cdot\|O_2\|_\gamma.
\]
\end{named}
\begin{proof}[Proof of (M1)]
Let $f$ lie in the space on which $O_2$ acts.
By (Mj0),
\[
\mathbf{1}_y\cdot O_1\cdot O_2\cdot\mathbf{1}_{y'}f=\sum_z\left(\mathbf{1}_y\cdot O_1\cdot\mathbf{1}_z\right)\left(\mathbf{1}_z\cdot O_2\cdot\mathbf{1}_{y'}f\right),
\]
and $\left|\left(\mathbf{1}_y\cdot O_1\cdot\mathbf{1}_z\right)f'\right|_\infty\le|O_1|(y,z)\cdot|f'|_\infty$ with $f':=\mathbf{1}_z\cdot O_2\cdot\mathbf{1}_{y'}f$, where $|f'|_\infty\le|O_2|(z,y')\cdot|f|_\infty$.
This gives the first inequality.
Then, by $d(y,y')\le d(y,z)+d(z,y')$,
\begin{multline*}
\sum_{y'}|O_1\cdot O_2|(y,y')\cdot\exp\left(\gamma\cdot d(y,y')\right)\\
\le\sum_z|O_1|(y,z)\cdot\exp\left(\gamma\cdot d(y,z)\right)\sum_{y'}|O_2|(z,y')\cdot\exp\left(\gamma\cdot d(z,y')\right)\le\|O_1\|_\gamma\cdot\|O_2\|_\gamma,
\end{multline*}
and
\begin{multline*}
\sum_y|O_1\cdot O_2|(y,y')\cdot\exp\left(\gamma\cdot d(y,y')\right)\\
\le\sum_z|O_2|(z,y')\cdot\exp\left(\gamma\cdot d(z,y')\right)\sum_y|O_1|(y,z)\cdot\exp\left(\gamma\cdot d(y,z)\right)\le\|O_1\|_\gamma\cdot\|O_2\|_\gamma.
\end{multline*}
\end{proof}

\begin{named}{(M2)}[Operators on $\mathfrak{S}(1,\ell_*)$: blocks, $\ell^2$ norm, trace norm; Section~G5.1.7]
Let $O$ be an operator on $\mathfrak{S}(1,\ell_*)$ with the $4\times4$ blocks $O(y,y')$, let $|O|_{\ell^2}$ be its $\ell^2$ operator norm and $|O|_1$ its trace norm.
Then
\[
|O(y,y')|_c\le|O|(y,y'),\qquad|O|_{\ell^2}\le2\,\|O\|_0,\qquad|O|_1\le4\sum_{y,y'}|O|(y,y').
\]
\end{named}
\begin{proof}[Proof of (M2)]
One has $(\mathbf{1}_y\cdot O\cdot\mathbf{1}_{y'}v)(y)=O(y,y')\cdot v(y')$, and the supremum of $|O(y,y')\cdot v(y')|_c$ over $|v(y')|_c\le1$ is $|O(y,y')|_c$.
The row with the index $(y,\beta)$ of the matrix $O$ has the $\ell^1$ norm
\[
\sum_{y'}\sum_{\beta'}|O(y,y')_{\beta\beta'}|\le\sum_{y'}|O(y,y')|_c\le\|O\|_0.
\]
The column with the index $(y',\beta')$ has the $\ell^1$ norm
\[
\sum_y\sum_\beta|O(y,y')_{\beta\beta'}|\le4\sum_y|O(y,y')|_c\le4\,\|O\|_0,
\]
since every entry has modulus at most $|O(y,y')|_c$.
Schur's test: for $v,v'\in\mathfrak{S}(1,\ell_*)$ the $\ell^2$ scalar product of $v'$ with $Ov$ is
\[
\sum\overline{v'_\beta(y)}\cdot O(y,y')_{\beta\beta'}\cdot v_{\beta'}(y'),
\]
where this sum and the following ones extend over the pairs of indices $(y,\beta)$, $(y',\beta')$ in $T^0_*\times\{1,\dots,4\}$.
Its modulus is at most $\sum|O(y,y')_{\beta\beta'}|\cdot|v'_\beta(y)|\cdot|v_{\beta'}(y')|$, and by the Cauchy--Schwarz inequality this is at most
\begin{multline*}
\left(\sum|O(y,y')_{\beta\beta'}|\cdot|v'_\beta(y)|^2\right)^{\frac{1}{2}}\cdot\left(\sum|O(y,y')_{\beta\beta'}|\cdot|v_{\beta'}(y')|^2\right)^{\frac{1}{2}}\\
\le\left(\|O\|_0\sum_{(y,\beta)}|v'_\beta(y)|^2\right)^{\frac{1}{2}}\cdot\left(4\,\|O\|_0\sum_{(y',\beta')}|v_{\beta'}(y')|^2\right)^{\frac{1}{2}};
\end{multline*}
in the first factor the sum over $(y',\beta')$ is the $\ell^1$ norm of a row, in the second the sum over $(y,\beta)$ is the $\ell^1$ norm of a column.
Hence $|O|_{\ell^2}\le\left(\|O\|_0\cdot4\,\|O\|_0\right)^{\frac{1}{2}}=2\,\|O\|_0$.
For the trace norm, $|O|_1=\operatorname{Tr}\left(O^{*}O\right)^{\frac{1}{2}}$.
By the polar decomposition \cite[Section~7.3]{HornJohnson2013}, $O=\mathrm{Wpd}\cdot\left(O^{*}O\right)^{\frac{1}{2}}$ with a unitary matrix $\mathrm{Wpd}$.
So
\[
|O|_1=\operatorname{Tr}\left(\mathrm{Wpd}^{*}\cdot O\right)=\sum\overline{\mathrm{Wpd}(y,y')_{\beta\beta'}}\cdot O(y,y')_{\beta\beta'},
\]
and an entry of a unitary matrix has modulus at most one.
Hence $|O|_1$ is at most the sum of the moduli of the entries of $O$, and the $16$ entries of a block add up to at most $4\,|O(y,y')|_c$.
\end{proof}

\begin{description}
\item[(Mj3)]
(Rows.)
Let $O$ be a linear operator from $\mathfrak{S}(1,\ell_*)$ to $\mathfrak{S}^{\mathrm{ch}}(\eta_*,\ell_*)$ or to the functions of pairs, let $t=x$, respectively $t=(x,x')$, be a point of the target, and let $O(t,y')$ be the $4\times4$ blocks:
$(Ov)(t)=\sum_{y'}O(t,y')\cdot v(y')$, with $y'$ ranging over a fundamental domain of $\ell_*\mathbb{Z}^3$ in $\mathbb{Z}^3$ (Section~I.1.6).
Then for every $\beta$
\[
\sum_{\beta'}|O(t,y')_{\beta\beta'}|\le|O(t,y')|_c=\sup\left\{\left|\left(O\cdot\mathbf{1}_{y'}v\right)(t)\right|_c:|v|_\infty\le1\right\}\le|O|([x],y').
\]
Indeed, $\left(O\cdot\mathbf{1}_{y'}v\right)(t)=O(t,y')\cdot v(y')$;
$|m|_c$ is the largest row sum of the moduli of the entries of a matrix $m$ (Section~I.1.5);
and $t$ is localised at $[x]$.
\item[(Mj4)]
(Integrals.)
If $\mathrm{sv}\mapsto O(\mathrm{sv})$ is continuous on $[0,1]$, then
\[
\left|\int_0^1\mathrm{d}\mathrm{sv}\,O(\mathrm{sv})\right|(y,y')\le\int_0^1\mathrm{d}\mathrm{sv}\,|O(\mathrm{sv})|(y,y').
\]
Hence
\begin{gather*}
\left\|\int_0^1\mathrm{d}\mathrm{sv}\,O(\mathrm{sv})\right\|_\gamma\le\sup_{\mathrm{sv}}\|O(\mathrm{sv})\|_\gamma,\\
\sum_{y,y'}\left|\int_0^1\mathrm{d}\mathrm{sv}\,O(\mathrm{sv})\right|(y,y')\le\sup_{\mathrm{sv}}\sum_{y,y'}|O(\mathrm{sv})|(y,y').
\end{gather*}
Indeed,
\[
\left|\mathbf{1}_y\cdot\int\mathrm{d}\mathrm{sv}\,O(\mathrm{sv})\cdot\mathbf{1}_{y'}f\right|_\infty\le\int\mathrm{d}\mathrm{sv}\,\left|\mathbf{1}_y\cdot O(\mathrm{sv})\cdot\mathbf{1}_{y'}f\right|_\infty,
\]
and the integrand $\mathrm{sv}\mapsto|O(\mathrm{sv})|(y,y')$ is continuous, being a seminorm of a continuous function with values in a space of finite dimension.
\end{description}

\subsubsection*{0.7. Exponential sums}

\begin{description}
\item[(Su)]
For $0<\gamma\le6$ and every unit site $y$ of $T^0_*$,
\[
\sum_{y'\in T^0_*}\exp\left(-\gamma\cdot d(y,y')\right)\le\mathrm{ssum}(\gamma)=\left(\frac{12}{\gamma}\right)^3
\]
(Section~G5.1.8; \cite[Section~13.1.6]{QEDI}).
Indeed, every $y'$ has a representative with $|y'-y|=d(y,y')$;
so the sum is at most
\[
\sum_{n\in\mathbb{Z}^3}\exp\left(-\gamma\cdot|n|\right)\le\left(\sum_{j\in\mathbb{Z}}\exp\left(-\frac{\gamma\cdot|j|}{3}\right)\right)^3=\coth\left(\frac{\gamma}{6}\right)^3
\]
by $|n|\ge\frac{|n|_1}{3}$, and $\coth(\mathrm{tq})=1+\frac{2}{\exp(2\,\mathrm{tq})-1}\le1+\frac{1}{\mathrm{tq}}\le\frac{2}{\mathrm{tq}}$ for $0<\mathrm{tq}\le1$.
\item[(Su1)]
$\mathrm{ssum}\left(\frac{\gamma^{F}}{2}\right)=\left(\frac{24}{\gamma^{F}}\right)^3\ge1$;
so $\mathbf{S}=C^{F}\cdot\mathrm{ssum}\left(\frac{\gamma^{F}}{2}\right)\ge C^{F}\ge1$ and $\mathrm{hL}=6\,\mathbf{S}\ge6\,C^{F}$.
\item[(Su2)]
For $0\le\gamma\le\frac{\gamma^{F}}{2}$ and unit sites $y$, $z$, $y'$,
\[
\exp\left(\gamma\cdot d(y,y')\right)\cdot\mathrm{Edy}(y,z)\cdot\mathrm{Edy}(z,y')\le\exp\left(-\frac{\gamma^{F}}{2}\,d(y,z)\right)\cdot\exp\left(-\frac{\gamma^{F}}{2}\,d(z,y')\right),
\]
because $d(y,y')\le d(y,z)+d(z,y')$.
\item[(Su3)]
If $d(z,z')\le1$, then $\mathrm{Edy}(z',y')\le\mathrm{e}\cdot\mathrm{Edy}(z,y')$ for every $y'$, because $d(z,y')\le1+d(z',y')$ and $\gamma^{F}\le1$.
\end{description}

\subsubsection*{0.8. Derivatives}

\begin{description}
\item[(De1)]
If $O_1$, $O_2$ are analytic near $A$, then $\partial_B\left(O_1\cdot O_2\right)=\left(\partial_BO_1\right)\cdot O_2+O_1\cdot\partial_BO_2$.
\item[(De2)]
If $O$ is analytic near $A$ and has invertible values, then $\partial_B\left(O^{-1}\right)=-O^{-1}\cdot\left(\partial_BO\right)\cdot O^{-1}$.
Indeed, $O(\cdot)^{-1}=\frac{\operatorname{adj}O(\cdot)}{\det O(\cdot)}$ is analytic near $A$ (by the argument of 3.2);
it remains to differentiate $O(A+\mathrm{sv}\cdot B)\cdot O(A+\mathrm{sv}\cdot B)^{-1}=1$.
\item[(De3)]
If $O$ is analytic in a neighbourhood of the segment $\{A+\mathrm{sv}\cdot B:0\le\mathrm{sv}\le1\}$, then
\[
O(A+B)-O(A)=\int_0^1\mathrm{d}\mathrm{sv}\,\left(\partial_BO\right)(A+\mathrm{sv}\cdot B).
\]
\item[(De4)]
For a linear form $\mathrm{lfm}(\cdot)$ on the bond functions,
\begin{gather*}
\partial_B\exp\left(i\cdot\mathrm{lfm}(A)\right)=i\cdot\mathrm{lfm}(B)\cdot\exp\left(i\cdot\mathrm{lfm}(A)\right),\\
\partial_B\partial_{B'}\exp\left(i\cdot\mathrm{lfm}(A)\right)=-\mathrm{lfm}(B)\cdot\mathrm{lfm}(B')\cdot\exp\left(i\cdot\mathrm{lfm}(A)\right).
\end{gather*}
This applies to $u_A(b)^{\pm1}$, $\varphi_A(y,c)^{\pm1}$, $U_A(c,c')^{\pm1}$, with $\mathrm{lfm}(A)=\pm e_*\cdot\eta_*\cdot A(b)$, $\pm\mathrm{phx}_A(y,c)$, $\pm\mathrm{ups}_A(c,c')$, respectively.
\end{description}

\subsection*{1. Proposition 33.1}

The setting of Propositions~33.1--33.4 is the one fixed in front of Proposition~33.1.

\begin{proof}[Proof of Proposition~33.1]
\textbf{First sentence (entire functions).}
In the bases of the delta functions, and up to the signs of the antiperiodic continuation, the entries of the six operators are the following:
by (0.1), numbers that do not contain $A$, and the numbers $-\eta_*^{-1}\cdot M_3(x)_{\beta\beta'}\cdot u_A(c,c_3)$;
by (0.2)--(0.4), the numbers $\frac{\eta_*^3}{48}\,\varphi_A(y,c')$, $\varphi_A(y,c)^{-1}$ and $\frac{\eta_*^3}{48}\,\varphi_A(y,c)^{-1}\cdot\varphi_A(y,c')$;
by (0.5), the numbers $d_X(x,x')^{-\alpha}\cdot U_A(c,c')$ and $-d_X(x,x')^{-\alpha}$.
By 0.4, $u_A(b)$, $\varphi_A(y,c)^{\pm1}$ and $U_A(c,c')$ are exponentials of linear forms in $A$.
Hence the entries are entire functions of the numbers $A(b)$, and $\mathrm{Fst}(A)$ is a sum of such operators and of the constant $\mu_*$.
This holds for an arbitrary complex bond function $A$ on $T_*$: the bound $\mathfrak{i}(A)\le\frac{1}{4}$ is not needed for it.

\textbf{(a)}
Let $|f|_\infty\le1$ and $|v|_\infty\le1$.
By (0.2),
\[
\left(Q_*(A)J^{*}\cdot\mathbf{1}_{y'}f\right)(y)=\int_{\Delta_y}\mathrm{d}x'\,\varphi_A(y,c')\cdot\left(\mathbf{1}_{y'}f\right)(x').
\]
For $y\ne y'$ the function $\mathbf{1}_{y'}f$ vanishes over $\Delta_y$, and the integral is $0$.
For $y=y'$, (Ph1), (Bk1), $\mathfrak{i}(A)\le\frac{1}{4}$ and (Nu1) give
\[
\left|\left(Q_*(A)J^{*}\cdot\mathbf{1}_yf\right)(y)\right|_c\le\sup_{c'\in\Delta_y}|\varphi_A(y,c')|\cdot\int_{\Delta_y}\mathrm{d}x'\le\exp\left(\frac{3}{2}\,\mathfrak{i}(A)\right)\le\exp\left(\frac{3}{8}\right)\le\frac{3}{2}.
\]
By (0.3),
\[
\left(JQ_*^{T}(-A)\cdot\mathbf{1}_{y'}v\right)(c,g)=\varphi_A([c],c)^{-1}\cdot\left(\mathbf{1}_{y'}v\right)([c]);
\]
this vanishes unless $[c]=y'$, and its norm $|\cdot|_c$ is at most $\exp\left(\frac{3}{2}\,\mathfrak{i}(A)\right)\le\frac{3}{2}$.
Applying (De4) under the finite sums of (0.2) and in (0.3), one finds
\begin{gather*}
\left(\partial_B\bigl(Q_*(A)J^{*}\bigr)f\right)(y)=\int_{\Delta_y}\mathrm{d}x'\,i\cdot\mathrm{phx}_B(y,c')\cdot\varphi_A(y,c')\cdot f(x'),\\
\left(\partial_B\bigl(JQ_*^{T}(-A)\bigr)v\right)(c,g)=-i\cdot\mathrm{phx}_B(y,c)\cdot\varphi_A(y,c)^{-1}\cdot v(y)\qquad(y=[c]).
\end{gather*}
Both operators map functions localised at $y'$ to functions localised at $y'$.
By (Ph1), $|\mathrm{phx}_B(y,\cdot)|\le\frac{3}{2}\,\epsilon_B$, and $\mathrm{phx}_B(y,\cdot)=0$ for $y$ outside $\mathrm{Ypr}_B$.
With the bound $\frac{3}{2}$ for $|\varphi_A^{\pm1}|$ and $\int_{\Delta_y}\mathrm{d}x'=1$, the two majorants are at most $\frac{3}{2}\,\epsilon_B\cdot\frac{3}{2}=\frac{9}{4}\,\epsilon_B$ for $y=y'\in\mathrm{Ypr}_B$, and they vanish otherwise.

\textbf{(b)}
\emph{The formula.}
In (0.1) only the factor $u_A(c,c_3)$ contains $A$.
By (De4), $\partial_Bu_A(c,c_3)=i\cdot e_*\cdot\eta_*\cdot B(c,c_3)\cdot u_A(c,c_3)$.
So
\begin{align*}
(\partial_BD_Af)(x)&=-\eta_*^{-1}\cdot M_3(x)\cdot i\cdot e_*\cdot\eta_*\cdot B(c,c_3)\cdot u_A(c,c_3)\cdot f(x^{(3)})\\
&=-i\cdot e_*\cdot B(c,c_3)\cdot M_3(x)\cdot u_A(c,c_3)\cdot f(x^{(3)}).
\end{align*}

\emph{Majorant of $\partial_BD_A$.}
Let $|f|_\infty\le1$, and let $x=(c,g)$ lie over $\Delta_z$.
Then
\[
\left(\left(\partial_BD_A\right)\cdot\mathbf{1}_{z'}f\right)(x)=-i\cdot e_*\cdot B(c,c_3)\cdot M_3(x)\cdot u_A(c,c_3)\cdot\left(\mathbf{1}_{z'}f\right)(x^{(3)}),
\]
and $\left(\mathbf{1}_{z'}f\right)(x^{(3)})=0$ unless $[c_3]=z'$.
By (Bk3) with $a=ge_1$, either $[c_3]=z$ or $[c_3]=z+ge_1$, a nearest neighbour of $z$.
So the majorant vanishes unless $z'=z$ or $z$, $z'$ are nearest neighbours.
By (Ph3), $\eta_*\le1$, $\mathfrak{i}(A)\le\frac{1}{4}$ and (Nu2),
\[
|u_A(c,c_3)|\le\exp\left(\eta_*\cdot\mathfrak{i}(A)\right)\le\exp\left(\frac{1}{4}\right)\le\frac{4}{3}.
\]
By (0.6) and $e_*\cdot|B(c,c_3)|\le\epsilon_B$,
\[
\left|\left(\left(\partial_BD_A\right)\cdot\mathbf{1}_{z'}f\right)(x)\right|_c\le\epsilon_B\cdot3\,(1+\rho)\cdot\frac{4}{3}=4\,(1+\rho)\cdot\epsilon_B\le8\,\epsilon_B.
\]
If the left side does not vanish, then $B(c,c_3)\ne0$;
the end points $c$ and $c_3$ of this bond lie in $\Delta_z$ and $\Delta_{z'}$, so that $z$ and $z'$ lie in $\mathrm{Ypr}_B$.

\emph{Majorant of $\partial_B\bigl(JP_*(A)J^{*}\bigr)$.}
By (0.4) and (De1),
\[
\partial_B\bigl(JP_*(A)J^{*}\bigr)=\Bigl(\partial_B\bigl(JQ_*^{T}(-A)\bigr)\Bigr)\cdot\bigl(Q_*(A)J^{*}\bigr)+\bigl(JQ_*^{T}(-A)\bigr)\cdot\partial_B\bigl(Q_*(A)J^{*}\bigr).
\]
By (M1) and (a), all four majorants being diagonal,
\begin{align*}
\left|\partial_B\bigl(JP_*(A)J^{*}\bigr)\right|(y,y')&\le\mathbf{1}[y=y'\in\mathrm{Ypr}_B]\cdot\left(\frac{9}{4}\cdot\frac{3}{2}+\frac{3}{2}\cdot\frac{9}{4}\right)\epsilon_B\\
&=\frac{27}{4}\,\epsilon_B\cdot\mathbf{1}[y=y'\in\mathrm{Ypr}_B]\le7\,\epsilon_B\cdot\mathbf{1}[y=y'\in\mathrm{Ypr}_B].
\end{align*}

\emph{Majorant of $\partial_B\mathrm{Fst}(A)$.}
One has $\partial_B\mathrm{Fst}(A)=\partial_BD_A+b_*\cdot\partial_B\bigl(JP_*(A)J^{*}\bigr)$, the term $\mu_*$ being constant, and $b_*\le1$ by (Dat).
Hence $|\partial_B\mathrm{Fst}(A)|(z,z)\le(8+7)\,\epsilon_B=15\,\epsilon_B$,
$|\partial_B\mathrm{Fst}(A)|(z,z')\le8\,\epsilon_B$ for nearest neighbours $z$, $z'$,
$|\partial_B\mathrm{Fst}(A)|(z,z')=0$ for all other pairs,
and the majorant vanishes unless $z$ and $z'$ lie in $\mathrm{Ypr}_B$.
This is the asserted bound in terms of $\mathfrak{d}$.

\emph{The norm.}
Let $0\le\gamma\le1$.
The function $\mathfrak{d}$ is symmetric, and a unit site has at most six nearest neighbours, each at the distance $1$ (Bk4).
So every row sum and every column sum of $|\partial_B\mathrm{Fst}(A)|(z,z')\cdot\exp\left(\gamma\cdot d(z,z')\right)$ is at most
\[
\epsilon_B\cdot\left(15+6\cdot8\cdot\exp(\gamma)\right)\le\epsilon_B\cdot\left(15+48\,\mathrm{e}\right)\le146\,\epsilon_B
\]
by (Nu3).

\emph{The bound (4.1).}
Only $z\in\mathrm{Ypr}_B$ contribute, and only $z'$ with $d(z,z')\le1$.
For these, $\mathrm{Edy}(z',y')\le\mathrm{e}\cdot\mathrm{Edy}(z,y')$ by (Su3), and for $z'=z$ the factor is $\mathrm{Edy}(z,y')$ itself.
So
\begin{multline*}
\sum_{z,z'}\mathrm{Edy}(y,z)\cdot|\partial_B\mathrm{Fst}(A)|(z,z')\cdot\mathrm{Edy}(z',y')\\
\le\epsilon_B\sum_{z\in\mathrm{Ypr}_B}\mathrm{Edy}(y,z)\cdot\left(15+6\cdot8\cdot\mathrm{e}\right)\cdot\mathrm{Edy}(z,y')\le146\,\epsilon_B\sum_{z\in\mathrm{Ypr}_B}\mathrm{Edy}(y,z)\cdot\mathrm{Edy}(z,y').
\end{multline*}

\textbf{(c)}
Let $(x,x')\in\mathrm{Pzt}_*$, $x=(c,g)$, $x'=(c',g')$.
Then $|c-c'|\le d_X(x,x')\le M$, by (I.1.5) and the definition of $\mathrm{Pzt}_*$ in Section~I.1.7.
By (Ph2), $3M\cdot\mathfrak{i}(A)\le\frac{1}{4}$ and (Nu2),
\[
\left|U_A(x,x')^{\pm1}\right|\le\exp\left(3\,|c-c'|\cdot\mathfrak{i}(A)\right)\le\exp\left(\frac{1}{4}\right)\le\frac{4}{3}.
\]
In (0.5) only $U_A(c,c')$ contains $A$, and by (De4) $\partial_BU_A(c,c')=i\cdot\mathrm{ups}_B(c,c')\cdot U_A(c,c')$.
So
\begin{equation}\tag{1.1}
\left(\partial_B\delta_{\alpha,A}f\right)(x,x')=d_X(x,x')^{-\alpha}\cdot i\cdot\mathrm{ups}_B(c,c')\cdot U_A(c,c')\cdot f(x').
\end{equation}
By (Ph2), $|\mathrm{ups}_B(c,c')|\le3\,|c-c'|\cdot\epsilon_B\le3\,d_X(x,x')\cdot\epsilon_B$.
Therefore the norm $|\cdot|_c$ of (1.1) is at most
\[
d_X^{-\alpha}\cdot3\,d_X\cdot\epsilon_B\cdot\frac{4}{3}\,|f(x')|_c=4\,d_X(x,x')^{1-\alpha}\cdot\epsilon_B\cdot|f(x')|_c.
\]
Since $0<1-\alpha<1$, one has $d_X^{1-\alpha}\le1$ if $d_X<1$, and $d_X^{1-\alpha}\le M^{1-\alpha}\le M$ if $1\le d_X\le M$, as $M\ge1$.

For a pair inside a cluster, $c=c'$, one has $\mathrm{ups}_B(c,c)=0$, and (1.1) vanishes.

Let (1.1) be different from $0$.
Then $\mathrm{ups}_B(c,c')\ne0$, and (Ph2) provides a photon site $\tilde{c}$ with $[\tilde{c}]\in\mathrm{Ypr}_B$, $|c-\tilde{c}|\le|c-c'|$ and $|c'-\tilde{c}|\le|c-c'|$.
Since $|c-c'|\le M$, (Bk2) and (Bk4) give
\[
d\left([x],\mathrm{Ypr}_B\right)\le\bigl|[c]-[\tilde{c}]\bigr|\le M,\qquad d\left([x'],\mathrm{Ypr}_B\right)\le\bigl|[c']-[\tilde{c}]\bigr|\le M;
\]
and if $d_X(x,x')<1$, then $|c-\tilde{c}|\le|c-c'|<1$, whence $d\left([x],\mathrm{Ypr}_B\right)\le\bigl|[c]-[\tilde{c}]\bigr|\le1$.

\textbf{(d)}
Let $d\left(\mathrm{Ypr}_B,\mathrm{Ypr}_{B'}\right)\ge M+2$.
(If $B=0$ or $B'=0$, the derivatives vanish.)
Then $\mathrm{Ypr}_B$ and $\mathrm{Ypr}_{B'}$ are disjoint.
\begin{description}
\item[(d1)]
No bond $b$ has $B(b)\ne0$ and $B'(b)\ne0$: the unit site of an end point of $b$ would lie in both sets.
By (De4),
\[
\partial_B\partial_{B'}u_A(b)=-e_*^2\cdot\eta_*^2\cdot B(b)\cdot B'(b)\cdot u_A(b)=0
\]
for every bond $b$.
By (0.1), $\partial_B\partial_{B'}D_A=0$.
\item[(d2)]
By (De4), $\partial_B\partial_{B'}\varphi_A(y,c)^{\pm1}=-\mathrm{phx}_B(y,c)\cdot\mathrm{phx}_{B'}(y,c)\cdot\varphi_A(y,c)^{\pm1}$.
By (Ph1) the first factor vanishes for $y$ outside $\mathrm{Ypr}_B$, and the second for $y$ outside $\mathrm{Ypr}_{B'}$.
So the product vanishes for every $y$, and (0.2), (0.3) give
\[
\partial_B\partial_{B'}\bigl(Q_*(A)J^{*}\bigr)=0,\qquad\partial_B\partial_{B'}\bigl(JQ_*^{T}(-A)\bigr)=0.
\]
\item[(d3)]
By (0.4), (De1) and (d2),
\[
\partial_B\partial_{B'}\bigl(JP_*(A)J^{*}\bigr)=\Bigl(\partial_B\bigl(JQ_*^{T}(-A)\bigr)\Bigr)\cdot\partial_{B'}\bigl(Q_*(A)J^{*}\bigr)+\Bigl(\partial_{B'}\bigl(JQ_*^{T}(-A)\bigr)\Bigr)\cdot\partial_B\bigl(Q_*(A)J^{*}\bigr).
\]
By (a) each of the four factors is diagonal in the unit sites, and its block at $y$ vanishes for $y$ outside $\mathrm{Ypr}_B$, respectively outside $\mathrm{Ypr}_{B'}$.
So both products vanish.
Together with (d1) this gives $\partial_B\partial_{B'}\mathrm{Fst}(A)=0$.
\item[(d4)]
By (De4) and (0.5),
\[
\left(\partial_B\partial_{B'}\delta_{\alpha,A}f\right)(x,x')=-d_X(x,x')^{-\alpha}\cdot\mathrm{ups}_B(c,c')\cdot\mathrm{ups}_{B'}(c,c')\cdot U_A(c,c')\cdot f(x').
\]
If both factors $\mathrm{ups}_B(c,c')$, $\mathrm{ups}_{B'}(c,c')$ were different from $0$, then (Ph2) would give photon sites $\tilde{c}$, $\tilde{c}'$ in the lattice box spanned by $c$ and $c'$ with $[\tilde{c}]\in\mathrm{Ypr}_B$ and $[\tilde{c}']\in\mathrm{Ypr}_{B'}$;
then $|\tilde{c}-\tilde{c}'|\le|c-c'|\le M$, and (Bk2), (Bk4) would give
\[
d\left(\mathrm{Ypr}_B,\mathrm{Ypr}_{B'}\right)\le\bigl|[\tilde{c}]-[\tilde{c}']\bigr|\le M<M+2,
\]
which is excluded.
Hence $\partial_B\partial_{B'}\delta_{\alpha,A}=0$.
\end{description}

\textbf{Uniformity.}
The bounds contain only numerals, $\rho$, $M$, $\alpha$, the constant $\gamma^{F}$ (through $\mathrm{Edy}$ in (4.1)), the parameter $\gamma\in[0,1]$ and the data $\epsilon_B$, $\mathrm{Ypr}_B$, $d_X$.
Apart from the hypothesis $\mathfrak{i}(A)\le\frac{1}{4}$ (in (c), $3M\cdot\mathfrak{i}(A)\le\frac{1}{4}$), their derivation rests on two kinds of facts of the setting (ST*)$^{\mathrm{ch}}$, properties of the data and conditions on the constants, all of which hold for every $N$, every $b_*\in[b_-,1]$, every $\mu_*\in\left(\mu_-,\frac{1}{L}\right]$ and all bond functions;
both kinds are collected in Section~5 of this chapter.
The properties of the data that enter include
$k_*\ge1$, $\eta_*\le1$, $0\le b_*\le1$ (since $b_-\le b_*\le1$ and $b_->0$),
and $\ell_*>2$, which follows from $\ell_*\in\mathrm{Cl}$ (through (Bk4)).
The conditions on the constants that enter include
$0\le\gamma^{F}\le1$ (through (Su3)), $0<\alpha<1$, $0\le\rho<1$ (through (0.6)),
$M\ge1$ together with the integrality of $M$ (through (Bk2), in (c) and in (d4)),
the oddness of $L$, and $b_->0$.
The numbers $\mu_*$, $N$, $\ell_*$, $e_*$ do not enter the numerals.
\end{proof}

\subsection*{2. Proposition~33.2}

\begin{proof}[Proof of Proposition~33.2]
\textbf{2.0. Facts from Proposition~35.4; notation.}
Proposition~35.4\,(SF2) supplies the following three facts.
\begin{itemize}
\item[(i)] The set $\mathrm{Gwk}$ is contained in the open set $\tilde{G}_*$ and is invariant under $A\mapsto -A$.
Every $A\in\mathrm{Gwk}$ satisfies $\mathfrak{i}(A)<1/16$ and $3M\,\mathfrak{i}(A)<1/4$, so that the clauses (a)--(d) of Proposition~33.1 hold at every point of $\mathrm{Gwk}$.
\item[(ii)] For $A\in\tilde{G}_*$ the operator $\mathrm{Fst}(A)$ is bijective, $S_*(\cdot)$ is analytic on $\tilde{G}_*$, and for all unit sites $y,y'$
\begin{equation*}
|S_*(A)|(y,y'),\quad|\mathrm{Sa}_*(A)|(y,y'),\quad|H_*(A)|(y,y'),\quad|\delta_{\alpha,A}H_*(A)|(y,y')\;\le\;C^{F}\,\mathrm{Edy}(y,y'),
\end{equation*}
where the second and the fourth majorant are those for the short pairs, the pairs inside a cluster included.
\item[(iii)] For $A\in\tilde{G}_*$, for every $(x,x')\in\mathrm{Pz}_*$ with $d_X(x,x')\ge1$ and for every $f\in\mathfrak{S}^{\mathrm{ch}}(\eta_*,\ell_*)$,
\begin{equation*}
|(\delta_{\alpha,A}f)(x,x')|_c\le\mathrm{e}\,|f(x')|_c+|f(x)|_c .
\end{equation*}
\end{itemize}
By (ii) and the first sentence of Proposition~33.1, the functions
\begin{equation*}
H_*(\cdot)=b_*\,S_*(\cdot)\,JQ_*^{T}(-\,\cdot\,),\qquad
D_*(\cdot)=b_*-b_*\,Q_*(\cdot)J^{*}H_*(\cdot),\qquad
\delta_{\alpha,\cdot}\,H_*(\cdot)
\end{equation*}
are analytic on $\tilde{G}_*$.
Fix $A\in\mathrm{Gwk}$.
For bond functions $B,B'$ and unit sites $y,z,z',y'$ put
\begin{equation*}
\mathrm{Fpr}_B:=\partial_B\mathrm{Fst}(A),\qquad
\mathrm{Fpr}_{B'}:=\partial_{B'}\mathrm{Fst}(A),\qquad
\mathrm{Ech}(y;z,z';y'):=\mathrm{Edy}(y,z)\cdot\mathrm{Edy}(z,z')\cdot\mathrm{Edy}(z',y');
\end{equation*}
Here $\mathrm{Fpr}_B$ and $\mathrm{Fpr}_{B'}$ are operators on $\mathfrak{S}^{\mathrm{ch}}(\eta_*,\ell_*)$.
By (Dat), $b_*\le1$.
In this notation the display (4.1) of Proposition~33.1\,(b) reads: for all unit sites $y,y'$,
\begin{equation*}
\sum_{z_1,z_2}\mathrm{Edy}(y,z_1)\cdot|\mathrm{Fpr}_B|(z_1,z_2)\cdot\mathrm{Edy}(z_2,y')
\le146\,\epsilon_B\sum_{z\in\mathrm{Ypr}_B}\mathrm{Edy}(y,z)\cdot\mathrm{Edy}(z,y');
\end{equation*}
the same inequality holds with $B'$ in place of $B$, Proposition~33.1\,(b) being valid for every bond function.

\textbf{2.1. (H1).}
By (De2), $\partial_B S_*(A)=-S_*(A)\cdot\mathrm{Fpr}_B\cdot S_*(A)$; hence, by (De1),
\begin{equation}
\begin{aligned}
\partial_B H_*(A)
&=b_*\bigl(\partial_B S_*(A)\bigr)\cdot JQ_*^{T}(-A)+b_*\,S_*(A)\cdot\partial_B\bigl(JQ_*^{T}(-A)\bigr)\\
&=-S_*(A)\cdot\mathrm{Fpr}_B\cdot H_*(A)+b_*\,S_*(A)\cdot\partial_B\bigl(JQ_*^{T}(-A)\bigr).
\end{aligned}
\tag{2.1}
\end{equation}
The first term of (2.1) is estimated by (M1) of Section~G5.1.7, by (ii) and by (4.1):
\begin{equation*}
\begin{aligned}
\bigl|S_*(A)\cdot\mathrm{Fpr}_B\cdot H_*(A)\bigr|(y,y')
&\le\sum_{z_1,z_2}C^{F}\,\mathrm{Edy}(y,z_1)\cdot|\mathrm{Fpr}_B|(z_1,z_2)\cdot C^{F}\,\mathrm{Edy}(z_2,y')\\
&\le146\,(C^{F})^2\epsilon_B\sum_{z\in\mathrm{Ypr}_B}\mathrm{Edy}(y,z)\cdot\mathrm{Edy}(z,y').
\end{aligned}
\end{equation*}
For the second term one uses (M1), (ii) and Proposition~33.1\,(a), according to which the majorant of $\partial_B\bigl(JQ_*^{T}(-A)\bigr)$ is diagonal:
\begin{equation*}
\begin{aligned}
\bigl|b_*\,S_*(A)\cdot\partial_B\bigl(JQ_*^{T}(-A)\bigr)\bigr|(y,y')
&\le C^{F}\,\mathrm{Edy}(y,y')\cdot\frac{9}{4}\,\epsilon_B\,\mathbf{1}[y'\in\mathrm{Ypr}_B]\\
&=\frac{9}{4}\,C^{F}\epsilon_B\,\mathrm{Edy}(y,y')\cdot\mathrm{Edy}(y',y')\cdot\mathbf{1}[y'\in\mathrm{Ypr}_B]\\
&\le\frac{9}{4}\,C^{F}\epsilon_B\sum_{z\in\mathrm{Ypr}_B}\mathrm{Edy}(y,z)\cdot\mathrm{Edy}(z,y');
\end{aligned}
\end{equation*}
the expression in the second line is the term $z=y'$ of the sum in the third.
As $C^{F}\ge1$, the two coefficients add up to
\begin{equation*}
146\,(C^{F})^2+\frac{9}{4}\,C^{F}\le\left(146+\frac{9}{4}\right)(C^{F})^2\le149\,(C^{F})^2 ,
\end{equation*}
and the first inequality of (4.2) is proved.

For the short pairs, (2.1) and (Cv3) give
\begin{equation*}
\delta_{\alpha,A}\partial_B H_*(A)=-\mathrm{Sa}_*(A)\cdot\mathrm{Fpr}_B\cdot H_*(A)+b_*\,\mathrm{Sa}_*(A)\cdot\partial_B\bigl(JQ_*^{T}(-A)\bigr).
\end{equation*}
In the two estimates above the first factor $S_*(A)$ entered only through the bound $|S_*(A)|(y,z)\le C^{F}\,\mathrm{Edy}(y,z)$, and by (ii) the operator $\mathrm{Sa}_*(A)$ satisfies the same bound on all short pairs, the pairs inside a cluster included.
Therefore the second inequality of (4.2) holds with the same constant $149\,(C^{F})^2$.

\textbf{2.2. (H2).}
Let $d(\mathrm{Ypr}_B,\mathrm{Ypr}_{B'})\ge M+2$.
By Proposition~33.1\,(d),
\begin{equation*}
\partial_B\partial_{B'}\mathrm{Fst}(A)=0,\qquad\partial_B\partial_{B'}\bigl(JQ_*^{T}(-A)\bigr)=0 .
\end{equation*}
Differentiating the identity $\partial_{B'}S_*(\cdot)=-S_*(\cdot)\cdot\bigl(\partial_{B'}\mathrm{Fst}(\cdot)\bigr)\cdot S_*(\cdot)$ in the direction $B$ by means of (De1) and (De2), one obtains
\begin{equation*}
\partial_B\partial_{B'}S_*(A)=S_*(A)\cdot\mathrm{Fpr}_B\cdot S_*(A)\cdot\mathrm{Fpr}_{B'}\cdot S_*(A)+S_*(A)\cdot\mathrm{Fpr}_{B'}\cdot S_*(A)\cdot\mathrm{Fpr}_B\cdot S_*(A).
\end{equation*}
Consequently, by $H_*(\cdot)=b_*\,S_*(\cdot)\,JQ_*^{T}(-\,\cdot\,)$ and (De1), the term containing $\partial_B\partial_{B'}\bigl(JQ_*^{T}(-A)\bigr)$ being zero,
\begin{equation}
\begin{aligned}
\partial_B\partial_{B'}H_*(A)
&=b_*\bigl(\partial_B\partial_{B'}S_*(A)\bigr)\cdot JQ_*^{T}(-A)+b_*\bigl(\partial_B S_*(A)\bigr)\cdot\partial_{B'}\bigl(JQ_*^{T}(-A)\bigr)\\
&\quad+b_*\bigl(\partial_{B'}S_*(A)\bigr)\cdot\partial_B\bigl(JQ_*^{T}(-A)\bigr)\\
&=S_*(A)\cdot\mathrm{Fpr}_B\cdot S_*(A)\cdot\mathrm{Fpr}_{B'}\cdot H_*(A)+S_*(A)\cdot\mathrm{Fpr}_{B'}\cdot S_*(A)\cdot\mathrm{Fpr}_B\cdot H_*(A)\\
&\quad-b_*\,S_*(A)\cdot\mathrm{Fpr}_B\cdot S_*(A)\cdot\partial_{B'}\bigl(JQ_*^{T}(-A)\bigr)\\
&\quad-b_*\,S_*(A)\cdot\mathrm{Fpr}_{B'}\cdot S_*(A)\cdot\partial_B\bigl(JQ_*^{T}(-A)\bigr).
\end{aligned}
\tag{2.2}
\end{equation}
For the first term of (2.2) one uses (M1) and (ii), and then (4.1) twice: first for the sum over $(z_1,z_2)$ with the outer points $(y,z_3)$, then for the sum over $(z_3,z_4)$ with the outer points $(z,y')$.
This gives
\begin{equation*}
\begin{aligned}
&\bigl|S_*(A)\cdot\mathrm{Fpr}_B\cdot S_*(A)\cdot\mathrm{Fpr}_{B'}\cdot H_*(A)\bigr|(y,y')\\
&\qquad\le(C^{F})^3\sum_{z_3,z_4}\left[\sum_{z_1,z_2}\mathrm{Edy}(y,z_1)\cdot|\mathrm{Fpr}_B|(z_1,z_2)\cdot\mathrm{Edy}(z_2,z_3)\right]|\mathrm{Fpr}_{B'}|(z_3,z_4)\cdot\mathrm{Edy}(z_4,y')\\
&\qquad\le146\,(C^{F})^3\epsilon_B\sum_{z\in\mathrm{Ypr}_B}\mathrm{Edy}(y,z)\sum_{z_3,z_4}\mathrm{Edy}(z,z_3)\cdot|\mathrm{Fpr}_{B'}|(z_3,z_4)\cdot\mathrm{Edy}(z_4,y')\\
&\qquad\le146^2\,(C^{F})^3\epsilon_B\epsilon_{B'}\sum_{z\in\mathrm{Ypr}_B}\sum_{z'\in\mathrm{Ypr}_{B'}}\mathrm{Ech}(y;z,z';y').
\end{aligned}
\end{equation*}
For the third term of (2.2), (M1), (ii), (4.1) and Proposition~33.1\,(a) give
\begin{equation*}
\begin{aligned}
&\bigl|b_*\,S_*(A)\cdot\mathrm{Fpr}_B\cdot S_*(A)\cdot\partial_{B'}\bigl(JQ_*^{T}(-A)\bigr)\bigr|(y,y')\\
&\qquad\le(C^{F})^2\sum_{z_1,z_2}\mathrm{Edy}(y,z_1)\cdot|\mathrm{Fpr}_B|(z_1,z_2)\cdot\mathrm{Edy}(z_2,y')\cdot\frac{9}{4}\,\epsilon_{B'}\,\mathbf{1}[y'\in\mathrm{Ypr}_{B'}]\\
&\qquad\le146\cdot\frac{9}{4}\,(C^{F})^2\epsilon_B\epsilon_{B'}\sum_{z\in\mathrm{Ypr}_B}\mathrm{Edy}(y,z)\cdot\mathrm{Edy}(z,y')\cdot\mathrm{Edy}(y',y')\cdot\mathbf{1}[y'\in\mathrm{Ypr}_{B'}]\\
&\qquad\le\frac{657}{2}\,(C^{F})^2\epsilon_B\epsilon_{B'}\sum_{z\in\mathrm{Ypr}_B}\sum_{z'\in\mathrm{Ypr}_{B'}}\mathrm{Ech}(y;z,z';y'),
\end{aligned}
\end{equation*}
the sum in the third line consisting of the terms $z'=y'$ of the sum in the last line.
Both inequalities hold for an arbitrary ordered pair of bond functions.
Applied to the pair $(B',B)$, they bound the second term of (2.2) by
\begin{equation*}
146^2\,(C^{F})^3\epsilon_B\epsilon_{B'}\sum_{z\in\mathrm{Ypr}_B}\sum_{z'\in\mathrm{Ypr}_{B'}}\mathrm{Ech}(y;z',z;y')
\end{equation*}
and the fourth term by
\begin{equation*}
\frac{657}{2}\,(C^{F})^2\epsilon_B\epsilon_{B'}\sum_{z\in\mathrm{Ypr}_B}\sum_{z'\in\mathrm{Ypr}_{B'}}\mathrm{Ech}(y;z',z;y').
\end{equation*}
Thus each of the two chains $\mathrm{Ech}(y;z,z';y')$ and $\mathrm{Ech}(y;z',z;y')$ enters with a coefficient at most
\begin{equation*}
146^2\,(C^{F})^3+\frac{657}{2}\,(C^{F})^2\le(21316+329)\,(C^{F})^3=21645\,(C^{F})^3 ,
\end{equation*}
which is the first inequality of (4.3).

For the short pairs, (2.2) and (Cv3) show that $\delta_{\alpha,A}\partial_B\partial_{B'}H_*(A)$ is the right-hand side of (2.2) with the first factor $S_*(A)$ of each of the four terms replaced by $\mathrm{Sa}_*(A)$.
Of this factor the estimates above used only the bound $|S_*(A)|(y,z_1)\le C^{F}\,\mathrm{Edy}(y,z_1)$, which $\mathrm{Sa}_*(A)$ satisfies as well, by (ii).
Hence the second inequality of (4.3) holds.

\textbf{2.3. Legs: blocks, rows, barred legs.}
Let $\mathrm{up}$ be a leg index with $\omega=0$.
Put $t:=x$ and $\mathrm{Tl}_A:=H_*(A)$ if $\mathrm{up}=\xi=(x,\beta,0)$, and $t:=(x,x')$ and $\mathrm{Tl}_A:=\delta_{\alpha,A}H_*(A)$ if $\mathrm{up}=\boldsymbol{\zeta}=(x,x',\beta,0)$.
In terms of the blocks of (Mj3),
\begin{equation*}
\psi_\beta(x)=\sum_{y'}\sum_{\beta'}H_*(A)(x,y')_{\beta\beta'}\,\Psi_{\beta'}(y'),\qquad
\chi_\beta(x,x')=\sum_{y'}\sum_{\beta'}\bigl(\delta_{\alpha,A}H_*(A)\bigr)\bigl((x,x'),y'\bigr)_{\beta\beta'}\,\Psi_{\beta'}(y').
\end{equation*}
Hence $\mathrm{Leg}_A\bigl(\mathrm{up};(y',\beta',0)\bigr)=\mathrm{Tl}_A(t,y')_{\beta\beta'}$.
The blocks of a derivative of $\mathrm{Tl}_A$ are the derivatives of its blocks.
Therefore (Mj3) yields, for $\mathrm{pd}$ equal to the identity, to $\partial_B$ or to $\partial_B\partial_{B'}$,
\begin{equation}
|\mathrm{pd}\,\mathrm{Leg}_A|(\mathrm{up};y')=\sum_{\beta'}\bigl|(\mathrm{pd}\,\mathrm{Tl}_A)(t,y')_{\beta\beta'}\bigr|\le\sup\left\{\bigl|\bigl((\mathrm{pd}\,\mathrm{Tl}_A)\mathbf{1}_{y'}v\bigr)(t)\bigr|_c:\ v\in\mathfrak{S}(1,\ell_*),\ |v|_\infty\le1\right\}.
\tag{2.3}
\end{equation}
By (I.AP), another choice of the representative of $y'$ or of $t$ changes signs only.

\emph{Barred legs.}
Let $\mathrm{ub}_1,\dots,\mathrm{ub}_4$ be the standard basis of $\mathbb{C}^4$.
The entries of $\mathsf{C}$ listed in Section~I.2.1 give
\begin{equation*}
\mathsf{C}\,\mathrm{ub}_1=-\mathrm{ub}_2,\qquad\mathsf{C}\,\mathrm{ub}_2=\mathrm{ub}_1,\qquad\mathsf{C}\,\mathrm{ub}_3=\mathrm{ub}_4,\qquad\mathsf{C}\,\mathrm{ub}_4=-\mathrm{ub}_3
\end{equation*}
($\mathsf{C}$ is monomial by (I.c6)), and $\mathsf{C}^{-1}=\mathsf{C}^{T}$ by (I.c3).
Let $\beta\mapsto\tilde\beta$ be the permutation that exchanges $1$ with $2$ and $3$ with $4$.
Then, for every $4\times4$ matrix $m$,
\begin{equation*}
(\mathsf{C}^{-1}m\,\mathsf{C})_{\beta\beta'}=\pm\,m_{\tilde\beta\tilde\beta'} .
\end{equation*}
One has $\bar\psi=\bar{H}_*(A)\bar\Psi=\mathsf{C}^{-1}H_*(-A)\,\mathsf{C}\,\bar\Psi$, and, $U_{-A}(x,x')$ being a scalar,
\begin{equation*}
\bar\chi(x,x')=(\delta_{\alpha,-A}\bar\psi)(x,x')=d_X(x,x')^{-\alpha}\bigl(U_{-A}(x,x')\bar\psi(x')-\bar\psi(x)\bigr)=\mathsf{C}^{-1}\bigl(\delta_{\alpha,-A}H_*(-A)\,\mathsf{C}\,\bar\Psi\bigr)(x,x').
\end{equation*}
Hence, for $\mathrm{up}=(x,\beta,1)$ or $\mathrm{up}=(x,x',\beta,1)$, and with $\widetilde{\mathrm{up}}:=(x,\tilde\beta,0)$, resp.\ $\widetilde{\mathrm{up}}:=(x,x',\tilde\beta,0)$,
\begin{equation*}
\mathrm{Leg}_A\bigl(\mathrm{up};(y',\beta',1)\bigr)=\pm\,\mathrm{Leg}_{-A}\bigl(\widetilde{\mathrm{up}};(y',\tilde\beta',0)\bigr),
\end{equation*}
with signs that do not depend on $A$.
For a function $\mathrm{Gf}$ of the connection,
\begin{equation*}
\partial_B\bigl[\mathrm{Gf}(-\,\cdot\,)\bigr](A)=-(\partial_B\mathrm{Gf})(-A),\qquad\partial_B\partial_{B'}\bigl[\mathrm{Gf}(-\,\cdot\,)\bigr](A)=(\partial_B\partial_{B'}\mathrm{Gf})(-A).
\end{equation*}
Summing the moduli over $\beta'$, one finds, for $\mathrm{pd}$ equal to the identity, to $\partial_B$ or to $\partial_B\partial_{B'}$,
\begin{equation*}
|\mathrm{pd}\,\mathrm{Leg}_A|(\mathrm{up};y')=|\mathrm{pd}\,\mathrm{Leg}_{A'}|(\widetilde{\mathrm{up}};y')\quad\text{at }A'=-A .
\end{equation*}
By (i), the point $-A$ lies in $\mathrm{Gwk}$; moreover $\mathrm{pos}(\widetilde{\mathrm{up}})=\mathrm{pos}(\mathrm{up})$, and the right-hand sides of (L1)--(L3) contain neither $A$ nor $\beta$.
So (L1)--(L3) for barred legs follow from (L1)--(L3) for unbarred legs at the point $-A$ of $\mathrm{Gwk}$.
In the rest of the proof $\omega=0$.

\emph{Positions.}
Write $x=(c,g)$ and $x'=(c',g')$.
For a short pair $(x,x')$ one has $|c-c'|\le d_X(x,x')<1$; hence $d([x],[x'])\le1$ by (Bk2) and (Bk4), and
\begin{equation*}
\mathrm{Edy}([x'],y'')\le\mathrm{e}\,\mathrm{Edy}([x],y'')\qquad\text{for every unit site }y''
\end{equation*}
by (Su3).
(For a pair inside a cluster, $[x']=[x]$.)
For a long pair, $d_X(x,x')=|c-c'|\ge1$, because $\eta_*<1$ by (Dat); so $[x]\ne[x']$ in $T^0_*$ by (Bk2), the set $\mathrm{pos}(\boldsymbol{\zeta})$ has two elements, and $d_X(x,x')^{-\alpha}\le1$.

\textbf{2.4. (L1).}
Let $|v|_\infty\le1$ and $\mathrm{gq}:=H_*(A)\mathbf{1}_{y'}v$.
By (ii), $|\mathrm{gq}(x)|_c\le|H_*(A)|([x],y')\le C^{F}\,\mathrm{Edy}([x],y')$ for every $x$.
If $\mathrm{up}=\xi$, then (2.3) gives $|\mathrm{Leg}_A|(\xi;y')\le C^{F}\,\mathrm{Edy}([x],y')$.
If $\mathrm{up}=\boldsymbol{\zeta}$ and the pair is short, then (2.3) and (ii) give
\begin{equation*}
|\mathrm{Leg}_A|(\boldsymbol{\zeta};y')\le|\delta_{\alpha,A}H_*(A)|([x],y')\le C^{F}\,\mathrm{Edy}([x],y').
\end{equation*}
If $\mathrm{up}=\boldsymbol{\zeta}$ and the pair is long, then (2.3), (iii) and (Nu5) give
\begin{equation*}
\begin{aligned}
|\mathrm{Leg}_A|(\boldsymbol{\zeta};y')
&\le\sup_v\bigl(\mathrm{e}\,|\mathrm{gq}(x')|_c+|\mathrm{gq}(x)|_c\bigr)
\le\mathrm{e}\,C^{F}\,\mathrm{Edy}([x'],y')+C^{F}\,\mathrm{Edy}([x],y')\\
&\le3\,C^{F}\bigl(\mathrm{Edy}([x'],y')+\mathrm{Edy}([x],y')\bigr).
\end{aligned}
\end{equation*}
In each of the three cases the bound is at most $3\,C^{F}\sum_{w\in\mathrm{pos}(\mathrm{up})}\mathrm{Edy}(w,y')$.

\emph{Consequences.}
The set $\mathrm{pos}(\mathrm{up})$ has at most two elements, and $\mathrm{Edy}(w,y')\le\exp\bigl(-\gamma^{F}d(\mathrm{pos}(\mathrm{up}),y')\bigr)$ for $w\in\mathrm{pos}(\mathrm{up})$.
Hence, by (Su1),
\begin{equation*}
|\mathrm{Leg}_A|(\mathrm{up};y')\le6\,C^{F}\exp\bigl(-\gamma^{F}d(\mathrm{pos}(\mathrm{up}),y')\bigr)\le\mathrm{hL}\,\exp\bigl(-\gamma^{F}d(\mathrm{pos}(\mathrm{up}),y')\bigr).
\end{equation*}
Let $0\le\gamma\le\gamma^{F}/2$.
The coefficient $\mathrm{Leg}_A(\mathrm{up};u)$ vanishes for $\omega'\ne\omega$, and $d(\mathrm{pos}(\mathrm{up}),u)=d(\mathrm{pos}(\mathrm{up}),y')\le d(w,y')$ for $w\in\mathrm{pos}(\mathrm{up})$, by (Cv1).
Since $\gamma^{F}-\gamma\ge\gamma^{F}/2$, (Su) yields
\begin{equation*}
\begin{aligned}
\sum_u|\mathrm{Leg}_A(\mathrm{up};u)|\exp\bigl(\gamma\,d(\mathrm{pos}(\mathrm{up}),u)\bigr)
&=\sum_{y'}|\mathrm{Leg}_A|(\mathrm{up};y')\exp\bigl(\gamma\,d(\mathrm{pos}(\mathrm{up}),y')\bigr)\\
&\le3\,C^{F}\sum_{w\in\mathrm{pos}(\mathrm{up})}\sum_{y'}\exp\bigl(-(\gamma^{F}-\gamma)\,d(w,y')\bigr)\\
&\le3\,C^{F}\cdot2\cdot\mathrm{ssum}(\gamma^{F}/2)=6\,\mathbf{S}=\mathrm{hL}.
\end{aligned}
\end{equation*}

\textbf{2.5. (L2).}
Let $|v|_\infty\le1$, let $\mathrm{gq}:=H_*(A)\mathbf{1}_{y'}v$ be as in 2.4, and put $\mathrm{gq}_B:=\bigl(\partial_B H_*(A)\bigr)\mathbf{1}_{y'}v$.
By (ii), as in 2.4, and by (H1), one has for every $x$
\begin{equation}
|\mathrm{gq}(x)|_c\le C^{F}\,\mathrm{Edy}([x],y'),\qquad
|\mathrm{gq}_B(x)|_c\le|\partial_B H_*(A)|([x],y')\le149\,(C^{F})^2\epsilon_B\sum_{z\in\mathrm{Ypr}_B}\mathrm{Edy}([x],z)\cdot\mathrm{Edy}(z,y').
\tag{2.4}
\end{equation}
Since $\mathrm{Ypr}_B\subset\mathrm{Ypr}_B^{+}$ and all terms are nonnegative, a sum over $z\in\mathrm{Ypr}_B$ is at most the same sum over $z\in\mathrm{Ypr}_B^{+}$.

\emph{The case $\mathrm{up}=\xi$.}
By (2.3) and (2.4),
\begin{equation*}
|\partial_B\mathrm{Leg}_A|(\xi;y')\le149\,(C^{F})^2\epsilon_B\sum_{z\in\mathrm{Ypr}_B^{+}}\mathrm{Edy}([x],z)\cdot\mathrm{Edy}(z,y'),
\end{equation*}
and $\mathrm{pos}(\xi)=\{[x]\}$.

\emph{The case $\mathrm{up}=\boldsymbol{\zeta}$.}
By (De1), $\partial_B\bigl(\delta_{\alpha,A}H_*(A)\bigr)=(\partial_B\delta_{\alpha,A})H_*(A)+\delta_{\alpha,A}\partial_B H_*(A)$; hence
\begin{equation}
\bigl((\partial_B\mathrm{Tl}_A)\mathbf{1}_{y'}v\bigr)(x,x')=\bigl((\partial_B\delta_{\alpha,A})\mathrm{gq}\bigr)(x,x')+\bigl(\delta_{\alpha,A}\mathrm{gq}_B\bigr)(x,x').
\tag{2.5}
\end{equation}

\emph{Short pairs.}
By Proposition~33.1\,(c), the paragraph \emph{Positions} of 2.3 and (Nu4), the first term of (2.5) has the norm at most
\begin{equation*}
4\,\epsilon_B|\mathrm{gq}(x')|_c\le4\,\epsilon_B C^{F}\,\mathrm{Edy}([x'],y')\le4\mathrm{e}\,C^{F}\epsilon_B\,\mathrm{Edy}([x],y')\le11\,C^{F}\epsilon_B\,\mathrm{Edy}([x],y');
\end{equation*}
it vanishes unless $d([x],\mathrm{Ypr}_B)\le1$, hence unless $[x]\in\mathrm{Ypr}_B^{+}$ (note $1\le M+1$), and it vanishes for a pair inside a cluster.
So it is at most
\begin{equation*}
11\,C^{F}\epsilon_B\,\mathrm{Edy}([x],[x])\cdot\mathrm{Edy}([x],y')\cdot\mathbf{1}\bigl[[x]\in\mathrm{Ypr}_B^{+}\bigr],
\end{equation*}
a term ($w=z=[x]$) of the sum in (L2).
By the second inequality of (H1), the second term of (2.5) is at most
\begin{equation*}
|\delta_{\alpha,A}\partial_B H_*(A)|([x],y')\le149\,(C^{F})^2\epsilon_B\sum_{z\in\mathrm{Ypr}_B}\mathrm{Edy}([x],z)\cdot\mathrm{Edy}(z,y').
\end{equation*}
The coefficients add up to $11\,C^{F}+149\,(C^{F})^2\le160\,(C^{F})^2$.

\emph{Long pairs.}
By Proposition~33.1\,(c), the first term of (2.5) has the norm at most $4M\epsilon_B|\mathrm{gq}(x')|_c\le4M\,C^{F}\epsilon_B\,\mathrm{Edy}([x'],y')$, and it vanishes unless $d([x'],\mathrm{Ypr}_B)\le M$, hence unless $[x']\in\mathrm{Ypr}_B^{+}$.
So it is at most
\begin{equation*}
4M\,C^{F}\epsilon_B\,\mathrm{Edy}([x'],[x'])\cdot\mathrm{Edy}([x'],y')\cdot\mathbf{1}\bigl[[x']\in\mathrm{Ypr}_B^{+}\bigr],
\end{equation*}
a term ($w=z=[x']$) of the sum in (L2); here $[x']\in\mathrm{pos}(\boldsymbol{\zeta})$.
The second term of (2.5) equals $d_X(x,x')^{-\alpha}\bigl(U_A(x,x')\,\mathrm{gq}_B(x')-\mathrm{gq}_B(x)\bigr)$.
By $d_X(x,x')^{-\alpha}\le1$, by $|U_A(x,x')|\le4/3$ (Proposition~33.1\,(c)) and by (2.4), its norm is at most
\begin{equation*}
\frac{4}{3}\cdot149\,(C^{F})^2\epsilon_B\sum_{z\in\mathrm{Ypr}_B}\mathrm{Edy}([x'],z)\cdot\mathrm{Edy}(z,y')+149\,(C^{F})^2\epsilon_B\sum_{z\in\mathrm{Ypr}_B}\mathrm{Edy}([x],z)\cdot\mathrm{Edy}(z,y').
\end{equation*}
For $w=[x']$ the sum carries the coefficient
\begin{equation*}
4M\,C^{F}+\frac{596}{3}\,(C^{F})^2\le(4M+199)\,(C^{F})^2\le203\,M(C^{F})^2\qquad(M\ge1,\ C^{F}\ge1),
\end{equation*}
and for $w=[x]$ the coefficient $149\,(C^{F})^2\le203\,M(C^{F})^2$.

In every case the coefficient is therefore at most $203\,M(C^{F})^2=C_H$, since $149\le160\le203\,M$.

\textbf{2.6. (L3).}
Let $d(\mathrm{Ypr}_B,\mathrm{Ypr}_{B'})\ge M+2$ and $|v|_\infty\le1$, let $\mathrm{gq}_B$ be as in 2.5, and put
\begin{equation*}
\mathrm{gq}_{B'}:=\bigl(\partial_{B'}H_*(A)\bigr)\mathbf{1}_{y'}v,\qquad\mathrm{gq}_{BB'}:=\bigl(\partial_B\partial_{B'}H_*(A)\bigr)\mathbf{1}_{y'}v .
\end{equation*}
By (H1), for $B$ and for $B'$, and by (H2), one has for every $x$
\begin{equation}
\begin{aligned}
|\mathrm{gq}_B(x)|_c&\le149\,(C^{F})^2\epsilon_B\sum_{z\in\mathrm{Ypr}_B}\mathrm{Edy}([x],z)\cdot\mathrm{Edy}(z,y'),\\
|\mathrm{gq}_{B'}(x)|_c&\le149\,(C^{F})^2\epsilon_{B'}\sum_{z'\in\mathrm{Ypr}_{B'}}\mathrm{Edy}([x],z')\cdot\mathrm{Edy}(z',y'),\\
|\mathrm{gq}_{BB'}(x)|_c&\le21645\,(C^{F})^3\epsilon_B\epsilon_{B'}\sum_{z\in\mathrm{Ypr}_B}\sum_{z'\in\mathrm{Ypr}_{B'}}\bigl[\mathrm{Ech}([x];z,z';y')+\mathrm{Ech}([x];z',z;y')\bigr].
\end{aligned}
\tag{2.6}
\end{equation}
For $z\in\mathrm{Ypr}_B^{+}$ and $z'\in\mathrm{Ypr}_{B'}^{+}$, $\mathrm{Ech}(w;z,z';y')$ is called the first chain and $\mathrm{Ech}(w;z',z;y')$ the second chain.
All terms are nonnegative, and $\mathrm{Ypr}_B\subset\mathrm{Ypr}_B^{+}$, $\mathrm{Ypr}_{B'}\subset\mathrm{Ypr}_{B'}^{+}$.

\emph{The case $\mathrm{up}=\xi$.}
By (2.3) and (2.6), each chain has the coefficient $21645\,(C^{F})^3$.

\emph{The case $\mathrm{up}=\boldsymbol{\zeta}$.}
By Proposition~33.1\,(d), $\partial_B\partial_{B'}\delta_{\alpha,A}=0$.
Applying (De1) twice, one finds
\begin{equation}
\begin{aligned}
\bigl((\partial_B\partial_{B'}\mathrm{Tl}_A)\mathbf{1}_{y'}v\bigr)(x,x')
&=\bigl((\partial_B\delta_{\alpha,A})\mathrm{gq}_{B'}\bigr)(x,x')+\bigl((\partial_{B'}\delta_{\alpha,A})\mathrm{gq}_B\bigr)(x,x')\\
&\quad+\bigl(\delta_{\alpha,A}\mathrm{gq}_{BB'}\bigr)(x,x').
\end{aligned}
\tag{2.7}
\end{equation}

\emph{Short pairs.}
By Proposition~33.1\,(c), (2.6) and $\mathrm{Edy}([x'],z')\le\mathrm{e}\,\mathrm{Edy}([x],z')$, the first term of (2.7) has the norm at most
\begin{equation*}
\begin{aligned}
&4\,\epsilon_B\,\mathbf{1}\bigl[[x]\in\mathrm{Ypr}_B^{+}\bigr]\,|\mathrm{gq}_{B'}(x')|_c\\
&\qquad\le4\mathrm{e}\cdot149\,(C^{F})^2\epsilon_B\epsilon_{B'}\,\mathbf{1}\bigl[[x]\in\mathrm{Ypr}_B^{+}\bigr]\sum_{z'\in\mathrm{Ypr}_{B'}}\mathrm{Edy}([x],[x])\cdot\mathrm{Edy}([x],z')\cdot\mathrm{Edy}(z',y')\\
&\qquad\le11\cdot149\,(C^{F})^2\epsilon_B\epsilon_{B'}\sum_{z\in\mathrm{Ypr}_B^{+}}\sum_{z'\in\mathrm{Ypr}_{B'}^{+}}\mathrm{Ech}([x];z,z';y')
\end{aligned}
\end{equation*}
(the terms $z=[x]$ of the first chain).
By Proposition~33.1\,(c) for $B'$, (2.6) and $\mathrm{Edy}([x'],z)\le\mathrm{e}\,\mathrm{Edy}([x],z)$, the second term of (2.7) has the norm at most
\begin{equation*}
\begin{aligned}
&4\,\epsilon_{B'}\,\mathbf{1}\bigl[[x]\in\mathrm{Ypr}_{B'}^{+}\bigr]\,|\mathrm{gq}_B(x')|_c\\
&\qquad\le4\mathrm{e}\cdot149\,(C^{F})^2\epsilon_B\epsilon_{B'}\,\mathbf{1}\bigl[[x]\in\mathrm{Ypr}_{B'}^{+}\bigr]\sum_{z\in\mathrm{Ypr}_B}\mathrm{Edy}([x],[x])\cdot\mathrm{Edy}([x],z)\cdot\mathrm{Edy}(z,y')\\
&\qquad\le11\cdot149\,(C^{F})^2\epsilon_B\epsilon_{B'}\sum_{z\in\mathrm{Ypr}_B^{+}}\sum_{z'\in\mathrm{Ypr}_{B'}^{+}}\mathrm{Ech}([x];z',z;y')
\end{aligned}
\end{equation*}
(the terms $z'=[x]$ of the second chain).
For a pair inside a cluster these two terms vanish (Proposition~33.1\,(c)).
The third term of (2.7) is at most $|\delta_{\alpha,A}\partial_B\partial_{B'}H_*(A)|([x],y')$, which the second inequality of (H2) bounds with the coefficient $21645\,(C^{F})^3$ for each chain.
Each chain therefore has the coefficient
\begin{equation*}
1639\,(C^{F})^2+21645\,(C^{F})^3\le23284\,(C^{F})^3 .
\end{equation*}

\emph{Long pairs.}
By Proposition~33.1\,(c) and (2.6), the first term of (2.7) has the norm at most
\begin{equation*}
\begin{aligned}
&4M\epsilon_B\,\mathbf{1}\bigl[[x']\in\mathrm{Ypr}_B^{+}\bigr]\,|\mathrm{gq}_{B'}(x')|_c\\
&\qquad\le4M\cdot149\,(C^{F})^2\epsilon_B\epsilon_{B'}\,\mathbf{1}\bigl[[x']\in\mathrm{Ypr}_B^{+}\bigr]\sum_{z'\in\mathrm{Ypr}_{B'}}\mathrm{Edy}([x'],[x'])\cdot\mathrm{Edy}([x'],z')\cdot\mathrm{Edy}(z',y')
\end{aligned}
\end{equation*}
(terms $w=z=[x']$ of the first chain).
By Proposition~33.1\,(c) for $B'$ and (2.6), the second term of (2.7) has the norm at most
\begin{equation*}
\begin{aligned}
&4M\epsilon_{B'}\,\mathbf{1}\bigl[[x']\in\mathrm{Ypr}_{B'}^{+}\bigr]\,|\mathrm{gq}_B(x')|_c\\
&\qquad\le4M\cdot149\,(C^{F})^2\epsilon_B\epsilon_{B'}\,\mathbf{1}\bigl[[x']\in\mathrm{Ypr}_{B'}^{+}\bigr]\sum_{z\in\mathrm{Ypr}_B}\mathrm{Edy}([x'],[x'])\cdot\mathrm{Edy}([x'],z)\cdot\mathrm{Edy}(z,y')
\end{aligned}
\end{equation*}
(terms $w=z'=[x']$ of the second chain).
For the third term, $d_X(x,x')^{-\alpha}\le1$ and $|U_A(x,x')|\le4/3$ give
\begin{equation*}
d_X(x,x')^{-\alpha}\bigl|U_A(x,x')\,\mathrm{gq}_{BB'}(x')-\mathrm{gq}_{BB'}(x)\bigr|_c\le\frac{4}{3}\,|\mathrm{gq}_{BB'}(x')|_c+|\mathrm{gq}_{BB'}(x)|_c ,
\end{equation*}
and the right-hand side is estimated by (2.6).
At $w=[x']$ each chain has the coefficient
\begin{equation*}
596\,M(C^{F})^2+\frac{4}{3}\cdot21645\,(C^{F})^3=596\,M(C^{F})^2+28860\,(C^{F})^3\le29456\,M(C^{F})^3 ,
\end{equation*}
and at $w=[x]$ the coefficient $21645\,(C^{F})^3$.

In every case the coefficient of each chain is at most
\begin{equation*}
(596\,M+1639+28860)\,(C^{F})^3\le31095\,M(C^{F})^3\le41209\,M^2(C^{F})^4=C_H^2 .
\end{equation*}

\textbf{2.7. (D1).}
Recall $D_*(A)=b_*-b_*\,Q_*(A)J^{*}H_*(A)$.
By (De1),
\begin{equation}
\partial_B D_*(A)=-b_*\bigl(\partial_B(Q_*(A)J^{*})\bigr)\cdot H_*(A)-b_*\,Q_*(A)J^{*}\cdot\partial_B H_*(A).
\tag{2.8}
\end{equation}
By (M1), Proposition~33.1\,(a) (whose majorants are diagonal), (ii), (H1) and $b_*\le1$,
\begin{equation*}
|\partial_B D_*(A)|(y,y')\le\frac{9}{4}\,\epsilon_B\,\mathbf{1}[y\in\mathrm{Ypr}_B]\cdot C^{F}\,\mathrm{Edy}(y,y')+\frac{3}{2}\cdot149\,(C^{F})^2\epsilon_B\sum_{z\in\mathrm{Ypr}_B}\mathrm{Edy}(y,z)\cdot\mathrm{Edy}(z,y').
\end{equation*}
The first term equals $\frac{9}{4}\,C^{F}\epsilon_B\,\mathrm{Edy}(y,y)\cdot\mathrm{Edy}(y,y')\cdot\mathbf{1}[y\in\mathrm{Ypr}_B]$, which is the term $z=y$ of the sum.
Hence the coefficient is at most
\begin{equation*}
\frac{9}{4}\,C^{F}+\frac{447}{2}\,(C^{F})^2\le\frac{903}{4}\,(C^{F})^2\le226\,(C^{F})^2 .
\end{equation*}

\emph{Norm.}
Let $\gamma=\gamma^{F}/2$.
By (Su2) and (Su), after the sum over $z$ has been extended to all unit sites,
\begin{equation*}
\begin{aligned}
\sum_{y'}|\partial_B D_*(A)|(y,y')\exp\bigl(\gamma\,d(y,y')\bigr)
&\le226\,(C^{F})^2\epsilon_B\sum_z\exp\left(-\frac{\gamma^{F}}{2}\,d(y,z)\right)\sum_{y'}\exp\left(-\frac{\gamma^{F}}{2}\,d(z,y')\right)\\
&\le226\,(C^{F})^2\,\mathrm{ssum}(\gamma^{F}/2)^2\,\epsilon_B=226\,\mathbf{S}^2\epsilon_B=C_K\epsilon_B ,
\end{aligned}
\end{equation*}
and in the same way
\begin{equation*}
\sum_y|\partial_B D_*(A)|(y,y')\exp\bigl(\gamma\,d(y,y')\bigr)\le226\,(C^{F})^2\epsilon_B\sum_z\exp\left(-\frac{\gamma^{F}}{2}\,d(z,y')\right)\sum_y\exp\left(-\frac{\gamma^{F}}{2}\,d(y,z)\right)\le C_K\epsilon_B .
\end{equation*}

\emph{Sum.}
Since $\mathrm{Edy}(y,z)\le\exp\left(-\frac{\gamma^{F}}{2}\,d(y,z)\right)$, (Su) gives
\begin{equation*}
\begin{aligned}
\sum_{y,y'}|\partial_B D_*(A)|(y,y')
&\le226\,(C^{F})^2\epsilon_B\sum_{z\in\mathrm{Ypr}_B}\left(\sum_y\mathrm{Edy}(y,z)\right)\left(\sum_{y'}\mathrm{Edy}(z,y')\right)\\
&\le226\,(C^{F})^2\,\mathrm{ssum}(\gamma^{F}/2)^2\,\epsilon_B\,\#|\mathrm{Ypr}_B|=C_K\epsilon_B\,\#|\mathrm{Ypr}_B| .
\end{aligned}
\end{equation*}

\textbf{2.8. (D2).}
Let $d(\mathrm{Ypr}_B,\mathrm{Ypr}_{B'})\ge M+2$.
By Proposition~33.1\,(d), $\partial_B\partial_{B'}\bigl(Q_*(A)J^{*}\bigr)=0$.
Differentiating (2.8), written for $B'$, in the direction $B$, one obtains
\begin{equation*}
\begin{aligned}
\partial_B\partial_{B'}D_*(A)=-b_*\Bigl[&\bigl(\partial_B(Q_*(A)J^{*})\bigr)\cdot\partial_{B'}H_*(A)+\bigl(\partial_{B'}(Q_*(A)J^{*})\bigr)\cdot\partial_B H_*(A)\\
&+Q_*(A)J^{*}\cdot\partial_B\partial_{B'}H_*(A)\Bigr].
\end{aligned}
\end{equation*}
By (M1), Proposition~33.1\,(a), (H1), (H2) and $b_*\le1$, the majorant of the first term is at most
\begin{equation*}
\frac{9}{4}\,\epsilon_B\,\mathbf{1}[y\in\mathrm{Ypr}_B]\cdot149\,(C^{F})^2\epsilon_{B'}\sum_{z'\in\mathrm{Ypr}_{B'}}\mathrm{Edy}(y,y)\cdot\mathrm{Edy}(y,z')\cdot\mathrm{Edy}(z',y')
\end{equation*}
(the terms $z=y$ of the first chain $\mathrm{Ech}(y;z,z';y')$), that of the second term is at most
\begin{equation*}
\frac{9}{4}\,\epsilon_{B'}\,\mathbf{1}[y\in\mathrm{Ypr}_{B'}]\cdot149\,(C^{F})^2\epsilon_B\sum_{z\in\mathrm{Ypr}_B}\mathrm{Edy}(y,y)\cdot\mathrm{Edy}(y,z)\cdot\mathrm{Edy}(z,y')
\end{equation*}
(the terms $z'=y$ of the second chain $\mathrm{Ech}(y;z',z;y')$), and that of the third term is at most
\begin{equation*}
\frac{3}{2}\cdot21645\,(C^{F})^3\epsilon_B\epsilon_{B'}\sum_{z\in\mathrm{Ypr}_B}\sum_{z'\in\mathrm{Ypr}_{B'}}\bigl[\mathrm{Ech}(y;z,z';y')+\mathrm{Ech}(y;z',z;y')\bigr].
\end{equation*}
The coefficient of each chain is therefore at most
\begin{equation*}
\frac{1341}{4}\,(C^{F})^2+\frac{64935}{2}\,(C^{F})^3\le\frac{131211}{4}\,(C^{F})^3\le32803\,(C^{F})^3\le41209\,M^2(C^{F})^4=C_H^2 .
\end{equation*}

\textbf{2.9. (D3).}
The segment $\{A_0+\mathrm{sv}\,B:\ 0\le\mathrm{sv}\le1\}$ lies in $\mathrm{Gwk}$, hence in the open set $\tilde{G}_*$, on which $D_*(\cdot)$ is analytic.
By (De3),
\begin{equation*}
D_*(A_0+B)-D_*(A_0)=\int_0^1\mathrm{d}\mathrm{sv}\,(\partial_B D_*)(A_0+\mathrm{sv}\,B).
\end{equation*}
Both bounds of (D3) follow from (D1) at the points $A_0+\mathrm{sv}\,B$ of $\mathrm{Gwk}$ and from (Mj4).

\emph{Uniformity.}
The right-hand sides of all bounds are formed from numerals, from $M$, $C^{F}$, $\gamma^{F}$ (the last through $\mathrm{Edy}$ and $\mathrm{ssum}$) and from the data $\epsilon_B$, $\epsilon_{B'}$, $\mathrm{Ypr}_B$, $\mathrm{Ypr}_{B'}$, $\mathrm{pos}(\mathrm{up})$ only.
Of the setting, the argument requires only Proposition~35.4\,(SF2), Proposition~33.1, and the inequalities $b_*\le1$, $\eta_*<1$, $M<\ell_*$, $\alpha>0$, $M\ge1$, $C^{F}\ge1$, $\gamma^{F}\le1$; the constants $C^{F}$, $\gamma^{F}$ of (SF2) do not depend on $N$, $b_*$, $\mu_*$, $A$ (Section~I.8.1).
\end{proof}

\subsection*{3. Proposition~33.3}

\begin{proof}[Proof of Proposition~33.3]
\textbf{3.0. Facts from Proposition~35.4 and Section~G5.1.8.}
By Proposition~35.4\,(SF2), the set $\mathrm{Gwk}$ is contained in $\tilde{G}_*$, at every point of which $\mathrm{Fst}(\cdot)$ is bijective; $\|S_*(A)\|_{\gamma^{F}/2}\le\mathbf{S}$ for $A\in\mathrm{Gwk}$; and $\mathfrak{i}(A'')<1/4$ for $A''\in\mathrm{Gwk}^{+}$.
Hence the clauses (a) and (b) of Proposition~33.1 hold at every point of $\mathrm{Gwk}^{+}$.
By Section~G5.1.8, $c_{\tau}\le1/(292\,\mathbf{S})$ and $C_K'=1328\,\mathbf{S}^2$.
Moreover $\gamma^{F}/2\le1$, $\mathbf{S}\ge1$ by (Su1), and $b_*\le1$ by (Dat).

\textbf{3.1. Invertibility and the norm of $S_*$.}
Let $A\in\mathrm{Gwk}$ and $e_{*}|A_1|_\infty\le c_{\tau}$.
For $0\le\mathrm{sv}\le1$ the point $A+\mathrm{sv}\,A_1$ lies in $\mathrm{Gwk}^{+}$, because $e_{*}|\mathrm{sv}\,A_1|_\infty\le c_{\tau}$.
Since $\mathrm{Fst}(\cdot)$ is entire (first sentence of Proposition~33.1), (De3) gives
\begin{equation*}
\mathrm{dD}:=\mathrm{Fst}(A+A_1)-\mathrm{Fst}(A)=\int_0^1\mathrm{d}\mathrm{sv}\,(\partial_{A_1}\mathrm{Fst})(A+\mathrm{sv}\,A_1).
\end{equation*}
Proposition~33.1\,(b) at the points $A+\mathrm{sv}\,A_1$, with the direction $A_1$ and with $\gamma=\gamma^{F}/2\in[0,1]$, together with (Mj4), yields
\begin{equation*}
\|\mathrm{dD}\|_{\gamma^{F}/2}\le146\,\epsilon_{A_1}=146\,e_{*}|A_1|_\infty\le146\,c_{\tau},
\end{equation*}
and then (M1) gives
\begin{equation*}
\|\mathrm{dD}\cdot S_*(A)\|_{\gamma^{F}/2}\le146\,c_{\tau}\mathbf{S}\le\frac{146}{292}=\frac{1}{2}.
\end{equation*}
By (Mj0), $\|\cdot\|_{\gamma^{F}/2}$ is a norm on the operators of the finite-dimensional space $\mathfrak{S}^{\mathrm{ch}}(\eta_*,\ell_*)$; by (M1) it is submultiplicative, and the identity has the norm $1$.
Hence the Neumann series $\sum_{n\ge0}\bigl(-\mathrm{dD}\cdot S_*(A)\bigr)^n$ converges, its sum has the norm at most $\sum_{n\ge0}2^{-n}=2$, and
\begin{equation*}
\bigl(1+\mathrm{dD}\cdot S_*(A)\bigr)\sum_{n\ge0}\bigl(-\mathrm{dD}\cdot S_*(A)\bigr)^n=1 .
\end{equation*}
Now $\mathrm{Fst}(A+A_1)=\mathrm{Fst}(A)+\mathrm{dD}=\bigl(1+\mathrm{dD}\cdot S_*(A)\bigr)\mathrm{Fst}(A)$, and both factors are invertible.
Therefore $\mathrm{Fst}(A+A_1)$ is invertible, and
\begin{equation}
S_*(A+A_1)=S_*(A)\sum_{n\ge0}\bigl(-\mathrm{dD}\cdot S_*(A)\bigr)^n,\qquad\|S_*(A+A_1)\|_{\gamma^{F}/2}\le\mathbf{S}\cdot2=2\,\mathbf{S}.
\tag{3.1}
\end{equation}
In particular, since $\exp\left(\frac{\gamma^{F}}{2}\,d(y,y')\right)\ge1$,
\begin{equation*}
\sup_y\sum_{y'}|S_*(A+A_1)|(y,y')\le2\,\mathbf{S},\qquad\sup_{y'}\sum_y|S_*(A+A_1)|(y,y')\le2\,\mathbf{S}.
\end{equation*}

\textbf{3.2. Analyticity.}
By the first sentence of Proposition~33.1, $\det\mathrm{Fst}(\cdot)$ is an entire function of the connection, and by 3.1 it does not vanish at the points of $\mathrm{Gwk}^{+}$.
Hence every point of $\mathrm{Gwk}^{+}$ has a neighbourhood on which $\det\mathrm{Fst}(\cdot)\ne0$.
On such a neighbourhood $S_*(\cdot)=\mathrm{Fst}(\cdot)^{-1}=\operatorname{adj}\mathrm{Fst}(\cdot)/\det\mathrm{Fst}(\cdot)$ is analytic, the entries of the adjugate being polynomials in the entries of $\mathrm{Fst}(\cdot)$ \cite[Section~0.8.2]{HornJohnson2013}.
So is
\begin{equation*}
D_*(\cdot)=b_*-b_*^2\,Q_*(\cdot)J^{*}S_*(\cdot)JQ_*^{T}(-\,\cdot\,),
\end{equation*}
whose two outer factors are entire (first sentence of Proposition~33.1).

\textbf{3.3. The derivative of $D_*(\cdot)$.}
Let $A''\in\mathrm{Gwk}^{+}$.
By (De1) and (De2),
\begin{equation}
\begin{aligned}
\partial_B D_*(A'')=-b_*^2\Bigl[&\bigl(\partial_B(Q_*(A'')J^{*})\bigr)\cdot S_*(A'')\cdot JQ_*^{T}(-A'')\\
&-Q_*(A'')J^{*}\cdot S_*(A'')\cdot\bigl(\partial_B\mathrm{Fst}(A'')\bigr)\cdot S_*(A'')\cdot JQ_*^{T}(-A'')\\
&+Q_*(A'')J^{*}\cdot S_*(A'')\cdot\partial_B\bigl(JQ_*^{T}(-A'')\bigr)\Bigr].
\end{aligned}
\tag{3.2}
\end{equation}
By Proposition~33.1\,(a), the majorants of $Q_*(A'')J^{*}$, of $JQ_*^{T}(-A'')$ and of their derivatives are diagonal; so for every $\gamma\ge0$ their norms $\|\cdot\|_\gamma$ are at most $3/2$, $3/2$, $\frac{9}{4}\,\epsilon_B$ and $\frac{9}{4}\,\epsilon_B$, respectively.
By (M1), (3.1), Proposition~33.1\,(b) and $b_*\le1$,
\begin{equation*}
\begin{aligned}
\|\partial_B D_*(A'')\|_{\gamma^{F}/2}
&\le\frac{9}{4}\,\epsilon_B\,(2\mathbf{S})\,\frac{3}{2}+\frac{3}{2}\,(2\mathbf{S})\,(146\,\epsilon_B)\,(2\mathbf{S})\,\frac{3}{2}+\frac{3}{2}\,(2\mathbf{S})\,\frac{9}{4}\,\epsilon_B\\
&=\left(\frac{27}{2}\,\mathbf{S}+1314\,\mathbf{S}^2\right)\epsilon_B\le\left(\frac{27}{2}+1314\right)\mathbf{S}^2\epsilon_B\le1328\,\mathbf{S}^2\epsilon_B=C_K'\epsilon_B .
\end{aligned}
\end{equation*}

\emph{Sum over $(y,y')$.}
By (M1) and Proposition~33.1\,(a), (b), together with the row and column sums of 3.1, the three terms of (3.2) are estimated as follows.
For the first term,
\begin{equation*}
\begin{aligned}
\sum_{y,y'}\frac{9}{4}\,\epsilon_B\,\mathbf{1}[y\in\mathrm{Ypr}_B]\cdot|S_*(A'')|(y,y')\cdot\frac{3}{2}
&\le\frac{27}{8}\,\epsilon_B\sum_{y\in\mathrm{Ypr}_B}\sum_{y'}|S_*(A'')|(y,y')\\
&\le\frac{27}{8}\,(2\mathbf{S})\,\epsilon_B\,\#|\mathrm{Ypr}_B| .
\end{aligned}
\end{equation*}
For the third term,
\begin{equation*}
\begin{aligned}
\sum_{y,y'}\frac{3}{2}\,|S_*(A'')|(y,y')\cdot\frac{9}{4}\,\epsilon_B\,\mathbf{1}[y'\in\mathrm{Ypr}_B]
&\le\frac{27}{8}\,\epsilon_B\sum_{y'\in\mathrm{Ypr}_B}\sum_y|S_*(A'')|(y,y')\\
&\le\frac{27}{8}\,(2\mathbf{S})\,\epsilon_B\,\#|\mathrm{Ypr}_B| .
\end{aligned}
\end{equation*}
For the second term,
\begin{equation*}
\begin{aligned}
&\sum_{y,y'}\sum_{z,z'}\frac{3}{2}\,|S_*(A'')|(y,z)\cdot|\partial_B\mathrm{Fst}(A'')|(z,z')\cdot|S_*(A'')|(z',y')\cdot\frac{3}{2}\\
&\qquad=\frac{9}{4}\sum_{z,z'}|\partial_B\mathrm{Fst}(A'')|(z,z')\left(\sum_y|S_*(A'')|(y,z)\right)\left(\sum_{y'}|S_*(A'')|(z',y')\right)\\
&\qquad\le\frac{9}{4}\,(2\mathbf{S})^2\sum_{z\in\mathrm{Ypr}_B}\sum_{z'}|\partial_B\mathrm{Fst}(A'')|(z,z')\le\frac{9}{4}\,(2\mathbf{S})^2\cdot146\,\epsilon_B\,\#|\mathrm{Ypr}_B| ;
\end{aligned}
\end{equation*}
in the last step Proposition~33.1\,(b) was used: the majorant of $\partial_B\mathrm{Fst}(A'')$ vanishes for $z\notin\mathrm{Ypr}_B$, and its row sums are at most $\|\partial_B\mathrm{Fst}(A'')\|_0\le146\,\epsilon_B$.
Together,
\begin{equation*}
\sum_{y,y'}|\partial_B D_*(A'')|(y,y')\le\left(\frac{27}{2}\,\mathbf{S}+1314\,\mathbf{S}^2\right)\epsilon_B\,\#|\mathrm{Ypr}_B|\le C_K'\epsilon_B\,\#|\mathrm{Ypr}_B| .
\end{equation*}

\textbf{3.4. The difference.}
Let $A'$, $A'+B$ and the segment between them lie in $\mathrm{Gwk}^{+}$.
By 3.2, $D_*(\cdot)$ is analytic in a neighbourhood of the segment, and (De3) gives
\begin{equation*}
D_*(A'+B)-D_*(A')=\int_0^1\mathrm{d}\mathrm{sv}\,(\partial_B D_*)(A'+\mathrm{sv}\,B).
\end{equation*}
Both bounds follow from 3.3 at the points $A''=A'+\mathrm{sv}\,B$ and from (Mj4).

\emph{Uniformity.}
The bounds contain only $\mathbf{S}$ and the data $\epsilon_B$, $\#|\mathrm{Ypr}_B|$.
Of the setting, the argument requires only Proposition~35.4\,(SF2), the first sentence and the clauses (a), (b) of Proposition~33.1, the inequalities $b_*\le1$, $\gamma^{F}\le1$, $C^{F}\ge1$, and the definitions of $c_{\tau}$ and $C_K'$.
\end{proof}

\subsection*{4. Proposition~33.4}

\begin{proof}[Proof of Proposition~33.4]
\emph{Norms.}
On $\mathfrak{S}^{\mathrm{ch}}(\eta_*,\ell_*)$ and on $\mathfrak{S}(1,\ell_*)$ the norms
\[
\|f\|:=\left((f,f)^{\mathrm{ch}}\right)^{\frac{1}{2}}=\left(\eta_*^3\cdot\frac{1}{48}\sum_{c}\sum_{g\in O_{\mathrm{cub}}}\left|f(c,g)\right|_2^2\right)^{\frac{1}{2}},
\qquad
\|v\|:=\left(\sum_{y}\left|v(y)\right|_2^2\right)^{\frac{1}{2}}
\]
are used, where $c$ ranges over a fundamental domain.
The operator norms for these norms are written $|\cdot|_{L^2}$ for operators on $\mathfrak{S}^{\mathrm{ch}}(\eta_*,\ell_*)$ and $|\cdot|_{\ell^2}$ for operators on $\mathfrak{S}(1,\ell_*)$;
an operator from one of the two spaces to the other is measured by the operator norm for the same two norms (Sections~I.1.5 and~I.1.6).

\textup{(Rw)} For every $4\times4$ matrix $m$, each row and each column of $m$ has $\ell^1$ norm at most $2|m|_2$.
Indeed, let $\mathrm{ub}_1,\dots,\mathrm{ub}_4$ be the standard basis of $\mathbb{C}^4$.
The Cauchy--Schwarz inequality for four terms gives
\[
\sum_{\beta'}\left|m_{\beta\beta'}\right|\le2\left(\sum_{\beta'}\left|m_{\beta\beta'}\right|^2\right)^{\frac{1}{2}}=2\left|m^{*}\,\mathrm{ub}_{\beta}\right|_2\le2|m|_2,
\qquad
\sum_{\beta}\left|m_{\beta\beta'}\right|\le2\left|m\,\mathrm{ub}_{\beta'}\right|_2\le2|m|_2.
\]

\emph{Proof of (a).}
Let $A$ be a real bond function on $T_*$.
By Proposition~35.4\,(A7), $\mathrm{Fst}(A)$ is invertible, so that $S_*(A)$ and $D_*(A)$ are defined.
Moreover $D_*(A)$ is invertible, $\left|D_*(A)^{-1}\right|_{\ell^2}\le b_*^{-1}+\mu_*^{-1}$, and the blocks of $\mathrm{Gcv}=D_*(A)^{-1}$ satisfy
\begin{equation}
\left|\mathrm{Gcv}(y,y')\right|_2\le b_*^{-1}\cdot\mathbf{1}[y=y']+\frac{2\mathrm{e}}{\mu_*}\cdot\exp\left(-\frac{\mu_*}{12\mathrm{e}}\cdot d(y,y')\right).
\tag{4.4}
\end{equation}
From (Dat), together with $b_->0$ and $\mu_->0$ (Section~0.1), one obtains $b_*^{-1}\le b_-^{-1}$ and $\mu_*^{-1}<\mu_-^{-1}$.
Hence
\[
\left|\mathrm{Gcv}\right|_{\ell^2}\le b_-^{-1}+\mu_-^{-1}=\mathrm{gGa}^2.
\]
By (Dat) and Section~G5.1.8 one has $\frac{\mu_*}{12\mathrm{e}}>\frac{\mu_-}{12\mathrm{e}}=\theta_{\mathrm{CT}}$ and $\mathrm{vt}\le\frac{\theta_{\mathrm{CT}}}{2}$.
It follows that
\[
\exp\left(\mathrm{vt}\cdot d(y,y')\right)\cdot\exp\left(-\frac{\mu_*}{12\mathrm{e}}\cdot d(y,y')\right)\le\exp\left(-\left(\theta_{\mathrm{CT}}-\mathrm{vt}\right)\cdot d(y,y')\right)\le\exp\left(-\frac{\theta_{\mathrm{CT}}}{2}\cdot d(y,y')\right).
\]
Further $0<\frac{\theta_{\mathrm{CT}}}{2}\le6$, because $\mu_-=L^{-3}\le1$ and $\mathrm{e}>1$.
Let $(y,\beta)$ index a row.
By (Rw), (4.4) and (Su),
\begin{align*}
\sum_{(y',\beta')}\left|\mathrm{Gcv}(y,\beta;y',\beta')\right|\cdot\exp\left(\mathrm{vt}\cdot d(y,y')\right)
&\le\sum_{y'}2\left|\mathrm{Gcv}(y,y')\right|_2\cdot\exp\left(\mathrm{vt}\cdot d(y,y')\right)\\
&\le2b_*^{-1}+\frac{4\mathrm{e}}{\mu_*}\sum_{y'}\exp\left(-\frac{\theta_{\mathrm{CT}}}{2}\cdot d(y,y')\right)\\
&\le2b_-^{-1}+4\mathrm{e}\cdot\mu_-^{-1}\cdot\mathrm{ssum}\left(\frac{\theta_{\mathrm{CT}}}{2}\right)=\mathrm{alG}.
\end{align*}
Let $(y',\beta')$ index a column.
By (Rw) for columns, by (4.4), and by (Su) with the centre $y'$,
\begin{align*}
\sum_{(y,\beta)}\left|\mathrm{Gcv}(y,\beta;y',\beta')\right|\cdot\exp\left(\mathrm{vt}\cdot d(y,y')\right)
&\le\sum_{y}2\left|\mathrm{Gcv}(y,y')\right|_2\cdot\exp\left(\mathrm{vt}\cdot d(y,y')\right)\\
&\le2b_*^{-1}+\frac{4\mathrm{e}}{\mu_*}\sum_{y}\exp\left(-\frac{\theta_{\mathrm{CT}}}{2}\cdot d(y,y')\right)\le\mathrm{alG}.
\end{align*}

\emph{Proof of (b).}
Let $A=A_r+A_1$ be as in (b): $A_r$ is real, $e_*\cdot|A_1|_\infty\le\frac{\mu_*}{24}$, and $\mathrm{Fst}(A)$ is invertible.

\emph{Step~1: the difference $D_A-D_{A_r}$.}
Let $f\in\mathfrak{S}^{\mathrm{ch}}(\eta_*,\ell_*)$ and $x=(c,g)$, and let $c_3=c+\eta_*\cdot g e_1$ be as in (0.1), so that $x^{(3)}=(c_3,g s_3)$.
Of the terms of (0.1), only the one that contains $u_A$ depends on the connection (I.O6), and $u_A(b)=u_{A_r}(b)\cdot u_{A_1}(b)$ for every bond $b$.
Hence the difference consists of a single term:
\[
\left(\left(D_A-D_{A_r}\right)f\right)(x)=-\eta_*^{-1}\cdot M_{3}(x)\cdot\left(u_{A_1}(c,c_3)-1\right)\cdot u_{A_r}(c,c_3)\cdot f\left(x^{(3)}\right).
\]
As $A_r$ is real, $\left|u_{A_r}(c,c_3)\right|=1$.
Next, $u_{A_1}(c,c_3)=\exp(\mathrm{tq})$ with $\mathrm{tq}=i\,e_*\cdot\eta_*\cdot A_1(c,c_3)$, and, by the hypothesis on $A_1$ and by (Dat),
\[
|\mathrm{tq}|\le\eta_*\cdot e_*\cdot|A_1|_\infty\le\frac{\mu_*}{24}\le\frac{1}{24}.
\]
By (Nu7) and (Nu6),
\[
\left|u_{A_1}(c,c_3)-1\right|\le|\mathrm{tq}|\cdot\exp|\mathrm{tq}|\le2\eta_*\cdot e_*\cdot|A_1|_\infty.
\]
With (0.6) this yields
\[
\left|\left(\left(D_A-D_{A_r}\right)f\right)(x)\right|_2\le\eta_*^{-1}\cdot6\cdot2\eta_*\cdot e_*\cdot|A_1|_\infty\cdot\left|f\left(x^{(3)}\right)\right|_2=12e_*\cdot|A_1|_\infty\cdot\left|f\left(x^{(3)}\right)\right|_2.
\]
In the sum over $T^{\mathrm{f}}_*$ that gives $\|f\|^2$, the site $x$ may be replaced by $x^{(3)}$:
by (I.s1) the substitution $x\mapsto x^{(3)}$ permutes the sites of $T^{\mathrm{f}}_*$,
the modulus $|f(\cdot)|_2$ is a function on $T^{\mathrm{f}}_*$,
and each site enters with one and the same weight $\frac{\eta_*^3}{48}$.
Therefore
\[
\eta_*^3\cdot\frac{1}{48}\sum_{x\in T^{\mathrm{f}}_*}\left|f\left(x^{(3)}\right)\right|_2^2=\|f\|^2,
\qquad
\left\|\left(D_A-D_{A_r}\right)f\right\|\le12e_*\cdot|A_1|_\infty\cdot\|f\|\le\frac{\mu_*}{2}\cdot\|f\|.
\]

\emph{Step~2: the operator $D_A+\mu_*$.}
Put $\mathrm{Rop}:=D_A+\mu_*$.
For the real connection $A_r$, (I.O2) gives
\[
\operatorname{Re}\left(f,D_{A_r}f\right)^{\mathrm{ch}}=\frac{\rho}{2\eta_*}\cdot\left(f,\nabla_{A_r}^{*}\nabla_{A_r}f\right)^{\mathrm{ch}}\ge0.
\]
By Step~1 and the Cauchy--Schwarz inequality,
\begin{align*}
\operatorname{Re}\left(f,\mathrm{Rop}\,f\right)^{\mathrm{ch}}
&=\operatorname{Re}\left(f,D_{A_r}f\right)^{\mathrm{ch}}+\mu_*\cdot\|f\|^2+\operatorname{Re}\left(f,\left(D_A-D_{A_r}\right)f\right)^{\mathrm{ch}}\\
&\ge\mu_*\cdot\|f\|^2-\frac{\mu_*}{2}\cdot\|f\|^2=\frac{\mu_*}{2}\cdot\|f\|^2.
\end{align*}
Consequently
\[
\left\|\mathrm{Rop}\,f\right\|\cdot\|f\|\ge\left|\left(f,\mathrm{Rop}\,f\right)^{\mathrm{ch}}\right|\ge\frac{\mu_*}{2}\cdot\|f\|^2 .
\]
Thus $\mathrm{Rop}$ is injective, hence invertible, the space being of finite dimension, and $\left|\mathrm{Rop}^{-1}\right|_{L^2}\le\frac{2}{\mu_*}$
(compare \cite[proof of Theorem~13.18]{QEDI}).

\emph{Step~3: the formula for $D_*(A)^{-1}$.}
The argument is that of the identity in \cite[proof of Theorem~13.18]{QEDI}, with $V_1=\mathfrak{S}^{\mathrm{ch}}(\eta_*,\ell_*)$, $V_2=\mathfrak{S}(1,\ell_*)$, $Q=Q_*(A)\cdot J^{*}$, $Q'=J\cdot Q_*^{T}(-A)$ and $b=b_*$.
Put
\[
\mathrm{cq}:=Q_*(A)\cdot J^{*}\cdot\mathrm{Rop}^{-1}\cdot J\cdot Q_*^{T}(-A),
\]
an operator on $\mathfrak{S}(1,\ell_*)$.
By (0.4),
\[
\mathrm{Fst}(A)=\mathrm{Rop}+b_*\cdot\left(J\cdot Q_*^{T}(-A)\right)\cdot\left(Q_*(A)\cdot J^{*}\right)=\mathrm{Rop}\cdot\left(1+b_*\cdot\mathrm{Rop}^{-1}\cdot J\cdot Q_*^{T}(-A)\cdot Q_*(A)\cdot J^{*}\right).
\]
Recall Sylvester's identity \cite[Section~0.8.5]{HornJohnson2013}: for linear maps $O_1\colon\mathfrak{S}(1,\ell_*)\to\mathfrak{S}^{\mathrm{ch}}(\eta_*,\ell_*)$ and $O_2\colon\mathfrak{S}^{\mathrm{ch}}(\eta_*,\ell_*)\to\mathfrak{S}(1,\ell_*)$ one has
\[
\det\left(1+O_1\cdot O_2\right)=\det\left(1+O_2\cdot O_1\right).
\]
Indeed, on $\mathfrak{S}^{\mathrm{ch}}(\eta_*,\ell_*)\oplus\mathfrak{S}(1,\ell_*)$ one has, in block form,
\[
\left(\begin{array}{cc}1&-O_1\\O_2&1\end{array}\right)
=\left(\begin{array}{cc}1&0\\O_2&1\end{array}\right)\left(\begin{array}{cc}1&-O_1\\0&1+O_2\cdot O_1\end{array}\right)
=\left(\begin{array}{cc}1+O_1\cdot O_2&-O_1\\0&1\end{array}\right)\left(\begin{array}{cc}1&0\\O_2&1\end{array}\right).
\]
Every factor on the right-hand sides is block triangular,
so that its determinant equals $\det$ of its upper left block times $\det$ of its lower right block.
Taking determinants in the two factorisations therefore gives
\[
\det\left(1+O_2\cdot O_1\right)=\det\left(\begin{array}{cc}1&-O_1\\O_2&1\end{array}\right)=\det\left(1+O_1\cdot O_2\right).
\]
With $O_1=b_*\cdot\mathrm{Rop}^{-1}\cdot J\cdot Q_*^{T}(-A)$ and $O_2=Q_*(A)\cdot J^{*}$ one obtains
\begin{align*}
\frac{\det\mathrm{Fst}(A)}{\det\mathrm{Rop}}
&=\det\left(1+b_*\cdot\mathrm{Rop}^{-1}\cdot J\cdot Q_*^{T}(-A)\cdot Q_*(A)\cdot J^{*}\right)\\
&=\det\left(1+b_*\cdot Q_*(A)\cdot J^{*}\cdot\mathrm{Rop}^{-1}\cdot J\cdot Q_*^{T}(-A)\right)=\det\left(1+b_*\cdot\mathrm{cq}\right).
\end{align*}
By hypothesis $\mathrm{Fst}(A)$ is invertible.
Hence $1+b_*\cdot\mathrm{cq}$ is invertible, and so is $b_*^{-1}+\mathrm{cq}$.
Put $\mathrm{Yq}:=\left(b_*^{-1}+\mathrm{cq}\right)^{-1}$.
Then
\[
\mathrm{cq}\cdot\mathrm{Yq}=\mathrm{Yq}\cdot\mathrm{cq}=1-b_*^{-1}\cdot\mathrm{Yq}.
\]
Put
\[
\mathrm{Swd}:=\mathrm{Rop}^{-1}-\mathrm{Rop}^{-1}\cdot J\cdot Q_*^{T}(-A)\cdot\mathrm{Yq}\cdot Q_*(A)\cdot J^{*}\cdot\mathrm{Rop}^{-1}.
\]
Then
\begin{align*}
\mathrm{Fst}(A)\cdot\mathrm{Swd}
&=1-J\cdot Q_*^{T}(-A)\cdot\mathrm{Yq}\cdot Q_*(A)\cdot J^{*}\cdot\mathrm{Rop}^{-1}+b_*\cdot J\cdot Q_*^{T}(-A)\cdot Q_*(A)\cdot J^{*}\cdot\mathrm{Rop}^{-1}\\
&\qquad-b_*\cdot J\cdot Q_*^{T}(-A)\cdot\mathrm{cq}\cdot\mathrm{Yq}\cdot Q_*(A)\cdot J^{*}\cdot\mathrm{Rop}^{-1}\\
&=1+J\cdot Q_*^{T}(-A)\cdot\left[-\mathrm{Yq}+b_*-b_*\cdot\left(1-b_*^{-1}\cdot\mathrm{Yq}\right)\right]\cdot Q_*(A)\cdot J^{*}\cdot\mathrm{Rop}^{-1}=1.
\end{align*}
So $S_*(A)=\mathrm{Swd}$, and
\begin{align*}
Q_*(A)\cdot J^{*}\cdot S_*(A)\cdot J\cdot Q_*^{T}(-A)
&=\mathrm{cq}-\mathrm{cq}\cdot\mathrm{Yq}\cdot\mathrm{cq}=\mathrm{cq}-\left(1-b_*^{-1}\cdot\mathrm{Yq}\right)\cdot\mathrm{cq}\\
&=b_*^{-1}\cdot\mathrm{Yq}\cdot\mathrm{cq}=b_*^{-1}\cdot\left(1-b_*^{-1}\cdot\mathrm{Yq}\right).
\end{align*}
Therefore
\[
D_*(A)=b_*-b_*^2\cdot Q_*(A)\cdot J^{*}\cdot S_*(A)\cdot J\cdot Q_*^{T}(-A)=b_*-b_*\cdot\left(1-b_*^{-1}\cdot\mathrm{Yq}\right)=\mathrm{Yq}=\left(b_*^{-1}+\mathrm{cq}\right)^{-1}.
\]
It follows that $D_*(A)$ is invertible and
\[
D_*(A)^{-1}=b_*^{-1}+\mathrm{cq}=b_*^{-1}+Q_*(A)\cdot J^{*}\cdot\left(D_A+\mu_*\right)^{-1}\cdot J\cdot Q_*^{T}(-A).
\]

\emph{Step~4: the averaging operators between $L^2$ and $\ell^2$.}
Since $A_r$ is real, $\mathfrak{i}(A)=e_*\cdot\left|\operatorname{Im}A_1\right|_\infty\le\frac{\mu_*}{24}\le\frac{1}{24}$.
By (Ph1) and (Nu6),
\[
\left|\varphi_A(y,c)^{\pm1}\right|\le\exp\left(\frac{3}{2}\cdot\mathfrak{i}(A)\right)\le\exp\left(\frac{1}{16}\right)\le2^{\frac{1}{2}}.
\]
Let $x'=(c',g')$.
By (0.2) and the Cauchy--Schwarz inequality on $\Delta_y\times O_{\mathrm{cub}}$, a set of measure $\int_{\Delta_y}\mathrm{d}x=1$ (Bk1),
\[
\left|\left(Q_*(A)\cdot J^{*}f\right)(y)\right|_2\le\int_{\Delta_y}\mathrm{d}x'\,\left|\varphi_A(y,c')\right|\cdot\left|f(x')\right|_2\le2^{\frac{1}{2}}\cdot\left(\int_{\Delta_y}\mathrm{d}x'\,\left|f(x')\right|_2^2\right)^{\frac{1}{2}}.
\]
Summing the squares over $y$ gives
\[
\left\|Q_*(A)\cdot J^{*}f\right\|\le2^{\frac{1}{2}}\cdot\|f\|.
\]
By (0.3) and (Bk1), with $x=(c,g)$,
\[
\left\|J\cdot Q_*^{T}(-A)\,v\right\|^2=\sum_{y}\int_{\Delta_y}\mathrm{d}x\,\left|\varphi_A(y,c)^{-1}\right|^2\cdot\left|v(y)\right|_2^2\le2\sum_{y}\left|v(y)\right|_2^2=2\|v\|^2.
\]

\emph{Step~5: conclusion.}
By Steps~2 and~4,
\[
\left|\mathrm{cq}\right|_{\ell^2}\le2^{\frac{1}{2}}\cdot\frac{2}{\mu_*}\cdot2^{\frac{1}{2}}=\frac{4}{\mu_*}.
\]
By Step~3 and (Dat),
\[
\left|D_*(A)^{-1}\right|_{\ell^2}\le b_*^{-1}+\frac{4}{\mu_*}\le b_-^{-1}+\frac{4}{\mu_-}=\mathrm{gGa}_1.
\]
For $(y,\beta)\in\mathrm{Ix}$ let $v_{(y,\beta)}$ be the element of the standard basis of $\mathfrak{S}(1,\ell_*)$ attached to $(y,\beta)$;
its norm is $1$.
Let $(\cdot,\cdot)_1$ be the scalar product of $\mathfrak{S}(1,\ell_*)$ whose norm is $\|\cdot\|$.
For $(y,\beta),(y',\beta')\in\mathrm{Ix}$ one has
\[
D_*(A)^{-1}(y,\beta;y',\beta')=\left(v_{(y,\beta)},D_*(A)^{-1}\,v_{(y',\beta')}\right)_1,
\qquad\text{whence}\qquad
\left|D_*(A)^{-1}(y,\beta;y',\beta')\right|\le\left|D_*(A)^{-1}\right|_{\ell^2}\le\mathrm{gGa}_1.
\]
A row of the matrix $D_*(A)^{-1}$ has $4\ell_*^3$ entries, one for each element of the index set $\mathrm{Ix}$.
Hence the $\ell^1$ norm of every row is at most $4\ell_*^3\cdot\mathrm{gGa}_1$.
\end{proof}

\begin{named}{Remark}
The bounds $\mathrm{gGa}^2$ and $\mathrm{alG}$ of (a) are functions of $b_-$ and $\mu_-$, that is, of $L$ alone.
The bounds in the proof of (b) are formed from the numerals $12$, $2$, $4$ and from $\mathrm{gGa}_1$;
they rest on $\eta_*\le1$, $\mu_-<\mu_*\le1$, $b_-\le b_*$ and $0\le\rho\le1$, the last through $\rho\ge0$ in (I.O2) and $3(1+\rho)\le6$ in (0.6).
They hold uniformly in the real part $A_r$, whatever its size, and also for $A$ outside $\tilde{G}_*$.
\end{named}

\subsection*{5. Conditions on the constants}

The proofs of Propositions~33.1--33.4 rest on the following conditions on the constants.

\begin{description}
\item[1.]
$0<\rho<1$ (Section~I.8.2, line~(C0); Proposition~14.2\,(a), row~$\rho$).
It enters as $3(1+\rho)\le6$ through (0.6), in Proposition~33.1\,(b) and in Step~1 of the proof of Proposition~33.4\,(b),
and as $\rho\ge0$ through (I.O2), in Step~2 of the same proof.
\item[2.]
$0<\alpha<1$ (Section~I.8.2, line~(C1); Proposition~14.2\,(f)).
It enters Proposition~33.1\,(c), through $d_X^{1-\alpha}\le1$, resp.\ $d_X^{1-\alpha}\le M$,
and 2.3, 2.5 and 2.6 in the proof of Proposition~33.2, through $d_X^{-\alpha}\le1$ for long pairs.
\item[3.]
$L$ is odd and $L\ge5$ (Section~I.0.5; Section~I.8.2, line~(C3); Proposition~14.2\,(b), (L5)).
It enters in two forms.
That $L$ is odd enters (Bk1), (Bk3) and (Ph1).
The inequality $L>1$, that is, $\eta_*<1$ and $b_->0$, enters (Dat), (Cv4) and 2.3 in the proof of Proposition~33.2, and gives $b_-=1-\frac{1}{L}>0$.
\item[4.]
$M$ is an integer and $M\ge1$ (Section~I.8.2, line~(C5); Proposition~14.2\,(b), (M1)).
It enters (Bk2), Proposition~33.1\,(c), and Proposition~33.2\,(L2), (L3), (D2).
\item[5.]
$0<\gamma^{F}\le1$ (Section~I.8.2, line~(C4r); Proposition~14.2\,(a), row~$\gamma^{F}$, in the form $\gamma^{F}\le\frac{1}{6L}$; Proposition~14.2\,(f)).
It enters (Su1)--(Su3), Proposition~33.1\,(b) in the display (4.1), and the choice $\gamma=\frac{\gamma^{F}}{2}\le1$ in the proof of Proposition~33.3;
as $\frac{\gamma^{F}}{2}\le6$ it enters (Su).
\item[6.]
$C^{F}\ge1$ (Section~I.8.2, line~(C6); Proposition~14.2\,(f), in the form $4\le C^{F}$).
It enters Proposition~33.2 in the sums of coefficients, (Su1), and Proposition~33.3.
\item[7.]
$b=1$, $b_-=1-\frac{1}{L}>0$ and $b_+=1$ (Section~I.8.2, lines~(C0) and~(C4); Section~I.3.4).
They enter through $b_-\le b_*\le1$ (Proposition~35.1\,(a)), and as $b_*^{-1}\le b_-^{-1}$ in Proposition~33.4\,(a) and~(b).
\item[8.]
$\mu_-=L^{-3}\le1$ (Section~I.8.2, line~(C4); Proposition~14.2\,(a), row~$\mu_-$).
It enters Proposition~33.4\,(a), through $\frac{\theta_{\mathrm{CT}}}{2}\le6$.
\item[9.]
$\mathbf{S}=C^{F}\cdot\mathrm{ssum}\left(\frac{\gamma^{F}}{2}\right)$ and $\mathrm{hL}=6\mathbf{S}$, with the inequalities $6C^{F}\le\mathrm{hL}$ and $\mathbf{S}\ge1$ (line~(C6$'$); Section~G5.9.3).
They enter Proposition~33.2\,(L1) and Proposition~33.3.
\item[10.]
$C_K=226\,\mathbf{S}^2=226\left(C^{F}\right)^2\cdot\mathrm{ssum}\left(\frac{\gamma^{F}}{2}\right)^2$ (line~(C6$'$)).
It enters Proposition~33.2\,(D1) and~(D3).
\item[11.]
$C_H=203M\cdot\left(C^{F}\right)^2$, with the inequalities $149\left(C^{F}\right)^2\le C_H$, $160\left(C^{F}\right)^2\le C_H$, and
\[
21645\left(C^{F}\right)^3,\quad31095M\cdot\left(C^{F}\right)^3,\quad32803\left(C^{F}\right)^3\;\le\;41209M^2\cdot\left(C^{F}\right)^4=C_H^2,
\]
all stated in the line~(C6$'$) and in Section~G5.9.3.
They enter Proposition~33.2\,(L2), (L3) and~(D2).
\item[12.]
$C_K'=1328\,\mathbf{S}^2\ge\frac{27}{2}\mathbf{S}+1314\,\mathbf{S}^2$ (line~(C6$'$); Section~G5.9.3).
It enters Proposition~33.3.
\item[13.]
$c_{\tau}\le\frac{1}{292\,\mathbf{S}}$, that is, $146\,c_{\tau}\cdot\mathbf{S}\le\frac{1}{2}$ (line~(C6$'$); Section~G5.9.3).
It enters 3.1 in the proof of Proposition~33.3.
\item[14.]
$\theta_{\mathrm{CT}}=\frac{\mu_-}{12\mathrm{e}}$ and $\mathrm{vt}\le\frac{\theta_{\mathrm{CT}}}{2}$ (line~(C6$'$); Section~G5.9.3).
They enter Proposition~33.4\,(a).
\item[15.]
$\mathrm{gGa}^2=b_-^{-1}+\mu_-^{-1}$,
$\mathrm{alG}=2b_-^{-1}+4\mathrm{e}\cdot\mu_-^{-1}\cdot\mathrm{ssum}\left(\frac{\theta_{\mathrm{CT}}}{2}\right)$ and
$\mathrm{gGa}_1=b_-^{-1}+\frac{4}{\mu_-}$ (line~(C6$'$)).
They enter Proposition~33.4\,(a) and~(b).
\end{description}

The conditions 9--15 consist of definitions made in the line (C6$'$) and of inequalities that follow from these definitions and from the conditions 4--6;
they are therefore met by the choice of the constants of the line (C6$'$).

Two further conditions enter the proofs above through the bounds
$\mathfrak{i}(A)<e_*^{2\varepsilon}\le\frac{1}{16}$ and $3M\cdot\mathfrak{i}(A)<\frac{1}{4}$ for $A\in\mathrm{Gwk}$, and $\mathfrak{i}(A)<\frac{1}{4}$ for $A\in\mathrm{Gwk}^{+}$,
of Proposition~35.4\,(SF2), which rest on them:
the condition (A8) of Proposition~14.2\,(b), which gives $e_*^{2\varepsilon}\le\frac{1}{16}$ and $3M\cdot e_*^{2\varepsilon}\le\frac{1}{4}$, and the condition $c_{\tau}\le\frac{1}{8}$.

Besides these conditions, the proofs rest on the following properties of the data:
$k_*\ge1$, $\eta_*\le\frac{1}{L}$, $b_*\in[b_-,1]$, $\mu_*\in\left(\mu_-,\frac{1}{L}\right]$ (Proposition~35.1\,(a));
$\ell_*\in\mathrm{Cl}$, hence $\ell_*>2M+2$ (Proposition~35.1\,(d) and Section~I.1.1);
$\mathfrak{i}(A)\le\frac{1}{4}$ and $3M\cdot\mathfrak{i}(A)\le\frac{1}{4}$ (Proposition~33.1);
$e_*\cdot|A_1|_\infty\le\frac{\mu_*}{24}$ and the invertibility of $\mathrm{Fst}(A)$ (Proposition~33.4\,(b)).

\section*{Chapter 34. The interacting factor, the functions $\mathfrak{P}_j$, the free ratio, one cell and the twist: proofs of Propositions 34.1--34.4}

\subsection*{0. Conventions of this chapter; preliminaries (P1)--(P7)}

\subsubsection*{0.1. Setting}

Throughout this chapter the setting is (ST*)$^{\mathrm{ch}}$ of Section~G5.1.1, and from Section~2 onwards the data are those of (DAT) in Section~G5.4.0.
The constants are those of Section~I.8, the line (C6$'$) included, together with the quantities of Section~G5.1.8.
For $0<e\le e_{\sharp}(\beta)$ the conditions (g1)--(g9) of Section~G5.1.9 hold, by the definition of $e_{\sharp}(\beta)$.

By Proposition~35.1\,(a) the stop level satisfies $k_{*}\ge1$.
Consequently, by Sections~I.3.6 and~I.5.2, the fields of the stop level are the following.
\begin{itemize}
\item The block spinors $\Psi(u)$, $u=(y,\beta,\omega)$ with $y\in T^0_*$, have four components per unit site and the weight $\nu=1$.
Their modes are the pairs $\mathrm{mo}=(y,\beta)$ in $\mathrm{Mo}:=T^0_*\times\{1,\dots,4\}$, whose number is $4\ell_*^3$, and $\mathsf{w}_W=1$.
\item The fine fields $\psi(\xi)$, $\xi=(x,\beta,\omega)$ with $x=(c,g)\in T^{\mathrm{f}}_*=T_*\times O_{\mathrm{cub}}$, have $192$ components per photon site $c$ and the weight $\eta_*^3/48$.
\item The H\"older fields $\chi(\boldsymbol{\zeta})$, $\boldsymbol{\zeta}=(x,x',\beta,\omega)$ with $(x,x')\in\mathrm{Pz}_*$, have the weight $\left(\eta_*^3/48\right)^2$.
The set $\mathrm{Pz}_*$ contains the $48\cdot47$ ordered pairs inside each cluster; for these pairs $c=c'$ and $d_X=\eta_*$, and no phase occurs.
\end{itemize}
The weight $\nu=1/48$ of the block spinors of level~$0$ and the number $\mathsf{w}_W=\sqrt{48}$ belong to the step from level~$0$ to level~$1$.
In this chapter the levels $j<k_{*}$, level~$0$ among them, occur only through the functions $E^{d}_j$ of Proposition~35.4\,(SF4), namely in Sections~1 and~3.

In this chapter (SF1)--(SF6) and (DF5) denote the clauses of Proposition~35.4, and (2.1)--(2.5) denote its displays.
(M2) is the fact so named in Section~G5.1.7.
The display (G5.5.1) of Section~G5.4.0 defines $\mathrm{Ztm}(t)$.
The display (6.1) is that of Proposition~34.1, and (6.4) is a display of Proposition~34.3.
(A8), (g3) and (g4) are the conditions listed under these names in Section~G5.1.9, and the condition (2.1) of Section~G5.1.9 is the inequality $r_{*}>31L+16$.
Finally, (C0)--(C11) and (C6$'$) are the lines of Section~I.8.2.

\subsubsection*{0.2. Notation}

In this chapter $\mathrm{sv}$ denotes a real parameter in $[0,1]$ and $\mathrm{zs}$ a complex variable; for a level $j$ one puts $i:=k_{*}-j$.
Moreover $\varepsilon_1=\frac{1}{2}\cdot e_{*}^{\frac{1}{4}-9\varepsilon}$ (Section~I.8.4).
The letter $\omega$ stands only for the bar index of Section~I.3.6; the real constant connection of Section~4 is denoted $\mathrm{Acl}$, as in Proposition~10\,(a).

\subsubsection*{0.3. (P1) Scales, charges, domains}

\begin{itemize}
\item[(a)]
Let $0\le j\le k_{*}$ and $i:=k_{*}-j$.
By Section~I.6.1,
\[
a_j=L^{-i}\cdot a_{*},\qquad e_j=e\cdot a_j^{\frac{1}{2}}=L^{-\frac{i}{2}}\cdot e_{*},\qquad \ell_j=L^{i}\cdot\ell_*,\qquad \eta_j=L^{i}\cdot\eta_*,\qquad a_j\cdot\eta_j=\eta.
\]
At $j=0$ the number $e_j$ is the charge $q=e\cdot\eta^{\frac{1}{2}}$ of level~$0$ (Section~I.2.5), and $a_0=\eta$, $\eta_0=1$, $\ell_0=n_*$.
(By Section~I.2.5 the numbers $e_0$ and $e_1$ are the ceilings of (C8) and (C9); the charges of the levels $0$ and $1$ are written $q$ and $L^{\frac{1}{2}}\cdot q$, and in a formula written for a general level $j$, such as (34.1) below, $e_j$ is the charge of level~$j$ for every $j$, so that $e_j=q$ and $e_{j+1}=L^{\frac{1}{2}}\cdot q$ when $j=0$.)
By (C7) and (C1), $0<1-\alpha<\varepsilon\le\bar{\varepsilon}=1/200$.
Proposition~35.1\,(a) gives $e_{*}\le e_0$, and the first condition of (A8) is $e_0^{\varepsilon}\le\frac{1}{4}$; hence
\begin{equation}\tag{34.1}
0<e_j\le e_{*}\le e_0<1,\qquad e_j=L^{-\frac{1}{2}}\cdot e_{j+1}<e_{j+1}\quad(j<k_{*}),\qquad e_{*}^{\varepsilon}\le\frac{1}{4}.
\end{equation}
Here $e_0<1$ holds because $e_0^{\varepsilon}\le\frac{1}{4}<1$ and $\varepsilon>0$, and $e_{*}^{\varepsilon}\le e_0^{\varepsilon}$ holds because $e_{*}\le e_0$ and $\varepsilon>0$.
Two consequences follow.
\begin{itemize}
\item[(a1)] For $0\le j\le k_{*}$ a power of $e_j$ with a real exponent decreases strictly when the exponent is increased, since $0<e_j<1$.
As $\varepsilon>0$, this gives $e_j^{3\varepsilon}<e_j^{2\varepsilon}<e_j^{\varepsilon}<1$.
\item[(a2)] $-1+2\varepsilon<-1+3\varepsilon<0$, since $0<\varepsilon$ and $3\varepsilon<1$.
\end{itemize}
Finally, $\ell_j\in\mathrm{Cl}$ for $j\le k_{*}$, and $\ell_{j+1}=\ell_j/L$ (Proposition~35.1\,(d)).
\item[(b)]
(Section~I.5.3, that is, \cite[Definition~5.16]{QEDI}.)
Let $0<\mathrm{cf}\le1$.
The conditions that define $\mathrm{cf}\,\tilde{R}_j$ are those that define $\tilde{R}_j$, with every threshold multiplied by $\mathrm{cf}$.
Therefore $\mathrm{cf}\,\tilde{R}_j\subset\tilde{R}_j$, and $0\in\mathrm{cf}\,\tilde{R}_j$.
The set
\[
\mathfrak{X}_j=\left\{\mathscr{C}:\ |\mathscr{C}|_{\infty}<e_j^{-1+2\varepsilon},\ |\partial\mathscr{C}|_{\infty}<e_j^{-1+3\varepsilon}\right\}
\]
contains $0$, and together with $\mathscr{C}$ it contains $\mathrm{zs}\cdot\mathscr{C}$ for every $\mathrm{zs}$ with $|\mathrm{zs}|\le1$.
The sets $G_j$, $\mathfrak{X}_j$, $\mathrm{cf}\,\tilde{R}_j$ and $\tilde{G}_j$ are open subsets of the space of the complex bond functions on $T_j$ \cite[Lemma~17.3]{QEDI}; that lemma holds at every level, for every side in $\mathrm{Cl}$ and for every charge in $(0,e_0]$.
Here $|\partial\mathscr{C}|_{\infty}$ is the supremum of $|(\partial_\nu\mathscr{C}_\mu)(c)|$ over all $\mu,\nu\in\{1,2,3\}$ and all sites $c$.
\item[(c)]
Let $\theta\in\mathbb{R}^3$.
The constant bond function $c_{\theta}$ is real, and it lies in $\frac{1}{4}\tilde{R}_*$: the function $0$ does by (b), and $\frac{1}{4}\tilde{R}_*$ is invariant under the addition of real constants (SF2).
By (b) and the definition of $\mathrm{Gwk}$ it follows that $c_{\theta}+\mathscr{C}\in\mathrm{Gwk}$ for every $\theta\in\mathbb{R}^3$ and every $\mathscr{C}\in\mathfrak{X}_*$.
\item[(d)]
The set $\mathrm{Gwk}$ is contained in $\tilde{G}_*$; the set $\tilde{G}_*$ is open and invariant under $A\mapsto-A$; and for $A\in\tilde{G}_*$ the operator $\mathrm{Fst}(A)$ is a bijection of $\mathfrak{S}^{\mathrm{ch}}(\eta_*,\ell_*)$ (SF2).
\end{itemize}

\subsubsection*{0.4. (P2) The operators as functions of the connection; the function $\mathrm{Zin}$}

Let $A$ be a complex bond function on $T_*$.
All spaces that occur have finite dimension, and ``entries'' are the matrix entries with respect to the bases of delta functions on a fundamental domain.
\begin{itemize}
\item[(a)]
Let $f\in\mathfrak{S}^{\mathrm{ch}}(\eta_*,\ell_*)$, $v\in\mathfrak{S}(1,\ell_*)$ and $x=(c,g)$.
By (I.2.3), (I.3.4), (I.1.7) and Section~G5.1.7 one has
\begin{align*}
(D_Af)(x)={}&-\eta_*^{-1}\sum_{i=0}^{2}M_{i}(x)\cdot\left[f(c,gs_i)-f(c,g)\right]\\
&-\eta_*^{-1}\cdot M_{3}(x)\cdot\left[\exp\left\{i\cdot e_{*}\cdot\eta_*\cdot A(c,c_3)\right\}\cdot f(c_3,gs_3)-f(c,g)\right],
\end{align*}
where $c_3=c+\eta_*\cdot ge_1$ and the vector $e_1$ is the first unit vector of $\mathbb{R}^3$ (Section~I.1.1);
\begin{align*}
\left(Q_*(A)J^{*}f\right)(y)&=\int_{\Delta_y}\mathrm{d}x\,\varphi_A(y,c)\cdot f(x)\\
&=\frac{\eta_*^3}{48}\sum_{c\in\Delta_y}\ \sum_{g\in O_{\mathrm{cub}}}\exp\left\{i\cdot\mathrm{phx}_A(y,c)\right\}\cdot f(c,g);
\end{align*}
for $c\in\Delta_y$
\[
\left(JQ_*^{T}(-A)v\right)(c,g)=\exp\left\{-i\cdot\mathrm{phx}_A(y,c)\right\}\cdot v(y),
\]
a value that is the same for the $48$ labels $g$; and for $(x,x')\in\mathrm{Pz}_*$
\[
\left(\delta_{\alpha,A}f\right)(x,x')=d_X(x,x')^{-\alpha}\cdot\left(\exp\left\{i\cdot\mathrm{ups}_A(c,c')\right\}\cdot f(x')-f(x)\right).
\]
Here $\mathrm{ups}_A(c,c)=0$, so that inside a cluster the quotient is $\eta_*^{-\alpha}\cdot\left(f(c,g')-f(c,g)\right)$.
The matrices $M_{i}(x)=\frac{w_i}{2}\left(\rho+\gamma\cdot(g\hat{n}_i)\right)$ and $M_{3}(x)=3\left(\rho-\gamma\cdot(ge_1)\right)$ of (I.2.1) contain no connection, and $\mathrm{phx}_A(y,c)$ and $\mathrm{ups}_A(c,c')$ are linear in $A$ (Section~G5.1.7).
Hence the entries of $D_A$, $Q_*(A)J^{*}$, $JQ_*^{T}(-A)$, $JP_*(A)J^{*}=\left(JQ_*^{T}(-A)\right)\cdot\left(Q_*(A)J^{*}\right)$ and $\mathrm{Fst}(A)=D_A+\mu_*+b_*\cdot JP_*(A)J^{*}$, and the coefficients of $\delta_{\alpha,A}$, are entire functions of the numbers $A(b)$.
\item[(b)]
Let $\mathrm{Inv}:=\left\{A:\ \mathrm{Fst}(A)\ \text{is invertible}\right\}$.
This set is open, because $\det\mathrm{Fst}$ is entire.
By Cramer's rule the entries of $S_*(A)=\mathrm{Fst}(A)^{-1}$ are analytic on $\mathrm{Inv}$, and therefore so are the entries of $H_{*}(A)=b_*\cdot S_*(A)\cdot JQ_*^{T}(-A)$, of $\delta_{\alpha,A}H_{*}(A)$ and of
\[
D_*(A)=b_*-b_*^2\cdot Q_*(A)J^{*}\cdot S_*(A)\cdot JQ_*^{T}(-A).
\]
By (P1)\,(d), $\tilde{G}_*\subset\mathrm{Inv}$.
\item[(c)]
If $A\in\tilde{G}_*$, then also $-A\in\tilde{G}_*$.
Therefore $\bar{H}_{*}(A)=\mathsf{C}^{-1}\cdot H_{*}(-A)\cdot\mathsf{C}$, where $\mathsf{C}=i\gamma_2$ is a constant matrix, and $\delta_{\alpha,-A}\bar{H}_{*}(A)$ are analytic on $\tilde{G}_*$.
Consequently every number $\mathrm{Leg}_A(\mathrm{up};u)$ of Section~G5.1.7 is an analytic function of $A$ on $\tilde{G}_*$; here $\mathrm{up}=\xi$ or $\mathrm{up}=\boldsymbol{\zeta}$, barred or unbarred, the pairs inside a cluster and the long pairs included.
\item[(d)]
Let $(\mathscr{A},\mathscr{C})\in\tilde{R}_*\times\mathfrak{X}_*$ and $\mathscr{T}:=\mathscr{A}+\mathscr{C}$; then $\mathscr{T}\in\tilde{G}_*$ (SF2).
Put
\begin{equation}\tag{34.2}
\begin{aligned}
\mathrm{Aex}(\mathscr{A},\mathscr{C})&:=-\left\langle\bar{\Psi},D_*(\mathscr{T})\Psi\right\rangle-m_*\cdot\left\langle\bar{\psi},\psi\right\rangle+\sum_{X\in D_*}\mathscr{E}_*(X;\mathscr{A},\mathscr{C},\psi,\chi)\ \in\mathrm{Gr}[\Psi],\\
\mathrm{Zin}(\mathscr{A},\mathscr{C})&:=\int\mathrm{D}\Psi\,\exp\mathrm{Aex}(\mathscr{A},\mathscr{C}),
\end{aligned}
\end{equation}
where the inserted fields are $\psi=H_{*}(\mathscr{T})\Psi$, $\bar{\psi}=\bar{H}_{*}(\mathscr{T})\bar{\Psi}$, $\chi=\delta_{\alpha,\mathscr{T}}\psi$ and $\bar{\chi}=\delta_{\alpha,-\mathscr{T}}\bar{\psi}$; that is, $\psi(\xi)=\sum_u\mathrm{Leg}_{\mathscr{T}}(\xi;u)\cdot\Psi(u)$ and $\chi(\boldsymbol{\zeta})=\sum_u\mathrm{Leg}_{\mathscr{T}}(\boldsymbol{\zeta};u)\cdot\Psi(u)$.
Written out (Sections~I.3.6 and~I.5.4),
\begin{align*}
\left\langle\bar{\Psi},D_*(\mathscr{T})\Psi\right\rangle&=\sum_{y,y'\in T^0_*}\ \sum_{\beta,\beta'}D_*(\mathscr{T})(y,y')_{\beta\beta'}\cdot\bar{\Psi}_{\beta}(y)\cdot\Psi_{\beta'}(y'),\\
\left\langle\bar{\psi},\psi\right\rangle&=\int\mathrm{d}x\,\bar{\psi}(x)\cdot\psi(x)=\frac{\eta_*^3}{48}\sum_{c\in T_*}\ \sum_{g\in O_{\mathrm{cub}}}\ \sum_{\beta=1}^{4}\bar{\psi}_{\beta}(c,g)\cdot\psi_{\beta}(c,g),\\
\mathscr{E}_*(X;\mathscr{A},\mathscr{C},\psi,\chi)&=\sum_{n,m}\frac{1}{n!\,m!}\sum_{\xi_1,\dots,\xi_n}\ \sum_{\boldsymbol{\zeta}_1,\dots,\boldsymbol{\zeta}_m}\left(\frac{\eta_*^3}{48}\right)^{n+2m}\\
&\qquad\cdot\mathscr{E}_{nm}(X;\mathscr{A},\mathscr{C};\xi_1,\dots,\xi_n;\boldsymbol{\zeta}_1,\dots,\boldsymbol{\zeta}_m)\\
&\qquad\cdot\psi(\xi_1)\cdots\psi(\xi_n)\cdot\chi(\boldsymbol{\zeta}_1)\cdots\chi(\boldsymbol{\zeta}_m).
\end{align*}
All these sums are finite.
The element $\mathrm{Aex}(\mathscr{A},\mathscr{C})$ of $\mathrm{Gr}[\Psi]$ is even: its first two terms are quadratic in $\Psi$, each $\mathscr{E}_*(X;\dots)$ is an even element of $\mathrm{Gr}[\psi,\chi]$, and the insertion is a homomorphism that maps generators to linear combinations of generators.
Each coefficient of $\mathrm{Aex}(\mathscr{A},\mathscr{C})$ is a finite sum of products of $m_*$, of the weights, of entries of $D_*(\mathscr{T})$, of numbers $\mathrm{Leg}_{\mathscr{T}}(\mathrm{up};u)$ and of kernels $\mathscr{E}_{nm}(X;\mathscr{A},\mathscr{C};\,\cdot\,)$.
By (b), (c) and (SF3) it is an analytic function of $(\mathscr{A},\mathscr{C})$ on the open set $\tilde{R}_*\times\mathfrak{X}_*$.

Write $\mathrm{Aex}=\mathrm{Aex}_0+\mathrm{Aex}_+$, where $\mathrm{Aex}_0$ is the coefficient of $1$.
The element $\mathrm{Aex}_+$ is even and has no constant term, and $\mathrm{Gr}[\Psi]$ has $8\ell_*^3$ generators; so $\mathrm{Aex}_+^n=0$ for $n>4\ell_*^3$, and
\[
\exp\mathrm{Aex}=\exp\left(\mathrm{Aex}_0\right)\cdot\sum_{n\le4\ell_*^3}\frac{\mathrm{Aex}_+^n}{n!}.
\]
The integral $\int\mathrm{D}\Psi$ is the coefficient of $\mathrm{Top}=\prod_{(y,\beta)}\Psi_{\beta}(y)\bar{\Psi}_{\beta}(y)$ ((I.3.5) with $\nu=1$).
Hence $\mathrm{Zin}$ is $\exp\left(\mathrm{Aex}_0\right)$ times a polynomial in the coefficients of $\mathrm{Aex}_+$, and $\mathrm{Zin}$ is analytic on $\tilde{R}_*\times\mathfrak{X}_*$.
\item[(e)]
With $\mathrm{Zin}$ as in (34.2), the definition (G5.5.1) of $\mathrm{Ztm}(t)$ and the representation (2.4) give
\begin{equation}\tag{34.3}
\begin{aligned}
\mathrm{Ztm}(t)&=\mathrm{Zin}\left(\mathscr{A},\mathscr{C}(t)\right)\qquad\left(t\in\mathrm{DDo}(t_x^{+})\right),\\
Y^{x}_{*}(A';\mathscr{C})&=\det D_*(0)^{-1}\cdot\exp\left\{-\varepsilon^{\mathrm{v},t}_*\cdot\ell_*^3+\sum_{X\in D_*}\mathfrak{P}_*(X,\mathscr{C})\right\}\cdot\mathrm{Zin}\left(H_{*}A',\mathscr{C}\right)\\
&\qquad\left(A'\in\mathrm{Dom}_*,\ \mathscr{C}\in\mathfrak{X}_*\right).
\end{aligned}
\end{equation}
In the first line $\mathscr{A}\in\frac{1}{4}\tilde{R}_*$ by (2.3) and $\mathscr{C}(t)\in\mathfrak{X}_*$ by (DF5); in the second line $H_{*}A'\in\frac{1}{4}\tilde{R}_*$ by the definition of $\mathrm{Dom}_*$.
\end{itemize}

\subsubsection*{0.5. (P3) Two rules in $\mathrm{Gr}[\Psi]$}

\begin{itemize}
\item[(a)]
Let $O$ be a linear map of $\mathfrak{S}(1,\ell_*)$ with the $4\times4$ blocks $O(y,y')$, and let $\mathrm{aw}>0$.
Then
\begin{equation}\tag{34.4}
\left|\left\langle\bar{\Psi},O\Psi\right\rangle\right|_{\mathrm{aw}}=\mathrm{aw}^2\sum_{y,y'}\ \sum_{\beta,\beta'}\left|O(y,y')_{\beta\beta'}\right|\le4\,\mathrm{aw}^2\sum_{y,y'}|O|(y,y').
\end{equation}
\begin{proof}
For different quadruples $(y,\beta,y',\beta')$ the products $\bar{\Psi}_{\beta}(y)\Psi_{\beta'}(y')$ are, up to signs, different ordered monomials of degree two; by (I.5.1) this gives the equality.
The number $|O(y,y')|_c$ is the largest of the four row sums of moduli of the block (Section~I.1.5); hence the sixteen moduli of a block add up to at most $4\,|O(y,y')|_c$.
It remains to show that $|O(y,y')|_c=|O|(y,y')$; in particular this gives the first clause of (M2).
The function $\mathbf{1}_yO\mathbf{1}_{y'}f$ vanishes outside the site $y$ and takes the value $O(y,y')f(y')$ at $y$; and when $f$ ranges over the spinors with $|f|_{\infty}\le1$, the vector $f(y')$ ranges over all $v\in\mathbb{C}^4$ with $|v|_c\le1$.
So $|O|(y,y')$ (Section~G5.1.7) is the operator norm of the block for $|\cdot|_c$, that is, $|O(y,y')|_c$.
\end{proof}
\item[(b)]
Let $\mathrm{aw}\ge1$, let $\mathrm{Cov}$ be an invertible linear map of the functions on $\mathrm{Mo}$ whose matrix has rows of $\ell^1$ norm at most $\mathrm{aw}$, and let $F$ be an even element of $\mathrm{Gr}[\Psi]$.
Then
\begin{equation}\tag{34.5}
\begin{aligned}
&\int\mathrm{D}\Psi\,\exp\left(-\left\langle\bar{\Psi},\mathrm{Cov}^{-1}\Psi\right\rangle+F\right)=\det\left(\mathrm{Cov}^{-1}\right)\cdot\left(1+\mathrm{Rq}\right)\\
&\qquad\text{with}\quad|\mathrm{Rq}|\le\exp\left(|F|_{\mathrm{aw}}\right)-1.
\end{aligned}
\end{equation}
\begin{proof}
Even elements commute; therefore the integrand equals $\exp\left(-\left\langle\bar{\Psi},\mathrm{Cov}^{-1}\Psi\right\rangle\right)\cdot\exp F$.
By Proposition~15\,(l)\,(iii), applied with the weight $\nu=1$, the left side of (34.5) equals
\[
\det\left(\mathrm{Cov}^{-1}\right)\int\mathrm{d}\mu_{\mathrm{Cov}}\,\exp F=\det\left(\mathrm{Cov}^{-1}\right)\int\mathrm{d}\mu_I\,s_{\mathrm{Cov}}(\exp F),
\]
where $s_{\mathrm{Cov}}$ is the homomorphism of $\mathrm{Gr}[\Psi]$ with $\Psi(\mathrm{mo})\mapsto(\mathrm{Cov}\,\Psi)(\mathrm{mo})=\sum_{\mathrm{mo}'}\mathrm{Cov}(\mathrm{mo},\mathrm{mo}')\cdot\Psi(\mathrm{mo}')$ and $\bar{\Psi}(\mathrm{mo})\mapsto\bar{\Psi}(\mathrm{mo})$.
Since $s_{\mathrm{Cov}}(1)=1$ and $\int\mathrm{d}\mu_I=1$ (Proposition~15\,(e)), one has $\mathrm{Rq}=\int\mathrm{d}\mu_I\,s_{\mathrm{Cov}}(\exp F-1)$, and $|\mathrm{Rq}|\le\left|s_{\mathrm{Cov}}(\exp F-1)\right|_1$ by Proposition~15\,(l)\,(ii), because $\mathsf{w}_W=1$.
Apply Proposition~15\,(d) with one family on each side, both equal to $\Psi$ on the index set $\mathrm{Mo}\times\{0,1\}$ with the weights one, and with the operator $O_{11}$ given by
\[
(O_{11}f)(\mathrm{mo},0)=\sum_{\mathrm{mo}'}\mathrm{Cov}(\mathrm{mo},\mathrm{mo}')\cdot f(\mathrm{mo}',0),\qquad(O_{11}f)(\mathrm{mo},1)=f(\mathrm{mo},1).
\]
Its rows have $\ell^1$ norm at most $\max(\mathrm{aw},1)=\mathrm{aw}=:C_{11}$, and with $\mathsf{w}'=1$, $\mathsf{w}=\mathrm{aw}$ the condition $C_{11}\cdot\mathsf{w}'\le\mathsf{w}$ is satisfied.
Hence $\left|s_{\mathrm{Cov}}(F')\right|_1\le|F'|_{\mathrm{aw}}$ for all $F'\in\mathrm{Gr}[\Psi]$.
Finally, $|\exp F-1|_{\mathrm{aw}}\le\exp\left(|F|_{\mathrm{aw}}\right)-1$ by Proposition~15\,(b).
\end{proof}
\end{itemize}

\subsubsection*{0.6. (P4) Matrices on $\mathfrak{S}(1,\ell_*)$}

The space $\mathfrak{S}(1,\ell_*)$ has the dimension $4\ell_*^3$ and the norm $\left(\sum_y|v(y)|_2^2\right)^{\frac{1}{2}}$.
For a linear map $Y$ of this space, $|Y|_{\ell^2}$ denotes its operator norm and $|Y|_1$ its trace norm, the sum of its singular values.
\begin{itemize}
\item[(a)] $|\operatorname{Tr}Y|\le|Y|_1$ and $|YY'|_1\le|Y|_{\ell^2}\cdot|Y'|_1$.
Indeed, by the polar decomposition in finite dimension, $Y=\mathrm{Uq}\,(Y^{*}Y)^{\frac{1}{2}}$ with a unitary $\mathrm{Uq}$.
Computed in an orthonormal eigenbasis of $(Y^{*}Y)^{\frac{1}{2}}$, the trace of $Y$ is the sum of the singular values, each multiplied by a diagonal entry of $\mathrm{Uq}$ in this basis, which is a number of modulus at most one.
By the min-max principle every singular value of $YY'$ is at most $|Y|_{\ell^2}$ times the singular value of $Y'$ with the same index.
\item[(b)] If $|Y|_{\ell^2}<1$, the series $\log(1+Y):=\sum_{n\ge1}(-1)^{n-1}Y^n/n$ converges in norm, and $\det(1+Y)=\exp\operatorname{Tr}\log(1+Y)$.
Indeed, in a basis in which $Y$ is upper triangular with the diagonal $(\mathrm{zq}_i)$ one has $|\mathrm{zq}_i|\le|Y|_{\ell^2}<1$, and $Y^n$ is upper triangular with the diagonal $(\mathrm{zq}_i^n)$.
Hence $\operatorname{Tr}\log(1+Y)=\sum_i\log(1+\mathrm{zq}_i)$ with the principal branch, and the exponential of this number is $\prod_i(1+\mathrm{zq}_i)=\det(1+Y)$.
\item[(c)] If $|Y|_{\ell^2}\le\frac{1}{2}$, then $|\operatorname{Tr}\log(1+Y)|\le2|Y|_1$.
Indeed, $|\operatorname{Tr}Y^n|\le|Y|_{\ell^2}^{n-1}\cdot|Y|_1$ by (a), and $\sum_{n\ge1}2^{-(n-1)}/n=2\log2\le2$.
\item[(d)] If the entries of $Y$ are analytic functions on an open subset of some $\mathbb{C}^n$ on which $|Y|_{\ell^2}<1$, then $\operatorname{Tr}\log(1+Y)$ is analytic there.
Indeed, by continuity $|Y|_{\ell^2}$ is bounded on every compact subset by a number smaller than one; there the series converges uniformly, and its terms are polynomials in the entries.
\item[(e)] (The second and the third clause of (M2).)
Let $O$ be a linear map of $\mathfrak{S}(1,\ell_*)$ with the $4\times4$ blocks $O(y,y')$, let $O(\mathrm{mo},\mathrm{mo}')=O(y,y')_{\beta\beta'}$ be its matrix on $\mathrm{Mo}$, where $\mathrm{mo}=(y,\beta)$ and $\mathrm{mo}'=(y',\beta')$, and let
\[
\|O\|_0=\max\left\{\sup_y\sum_{y'}|O|(y,y'),\ \sup_{y'}\sum_y|O|(y,y')\right\}
\]
(Section~G5.1.7).
Then
\[
|O|_{\ell^2}\le2\|O\|_0,\qquad|O|_1\le4\sum_{y,y'}|O|(y,y').
\]
\begin{proof}
By (P3)\,(a), $|O|(y,y')=|O(y,y')|_c$.

\emph{Operator norm.}
The row $(y,\beta)$ of the matrix has the $\ell^1$ norm $\sum_{y'}\sum_{\beta'}|O(y,y')_{\beta\beta'}|\le\sum_{y'}|O(y,y')|_c\le\|O\|_0$.
Every entry of a block has modulus at most $|O(y,y')|_c$; so the column $(y',\beta')$ has the $\ell^1$ norm $\sum_y\sum_{\beta}|O(y,y')_{\beta\beta'}|\le4\sum_y|O(y,y')|_c\le4\|O\|_0$.
Write $v,v'\in\mathfrak{S}(1,\ell_*)$ as functions on $\mathrm{Mo}$, $v(\mathrm{mo})=v_{\beta}(y)$.
The Cauchy--Schwarz inequality for the sum over the pairs $(\mathrm{mo},\mathrm{mo}')$ with the weights $|O(\mathrm{mo},\mathrm{mo}')|$ gives
\begin{align*}
&\left|\sum_{\mathrm{mo},\mathrm{mo}'}\overline{v'(\mathrm{mo})}\cdot O(\mathrm{mo},\mathrm{mo}')\cdot v(\mathrm{mo}')\right|\\
&\qquad\le\left(\sum_{\mathrm{mo},\mathrm{mo}'}|O(\mathrm{mo},\mathrm{mo}')|\cdot|v'(\mathrm{mo})|^2\right)^{\frac{1}{2}}\cdot\left(\sum_{\mathrm{mo},\mathrm{mo}'}|O(\mathrm{mo},\mathrm{mo}')|\cdot|v(\mathrm{mo}')|^2\right)^{\frac{1}{2}}\\
&\qquad\le\left(\|O\|_0\sum_{\mathrm{mo}}|v'(\mathrm{mo})|^2\right)^{\frac{1}{2}}\cdot\left(4\|O\|_0\sum_{\mathrm{mo}'}|v(\mathrm{mo}')|^2\right)^{\frac{1}{2}},
\end{align*}
and the right side is $2\|O\|_0$ times the product of the norms of $v'$ and $v$.

\emph{Trace norm.}
For every linear map $Y$,
\[
|Y|_1=\max\left\{|\operatorname{Tr}(\mathrm{Vq}\,Y)|:\ \mathrm{Vq}\ \text{a linear map with}\ |\mathrm{Vq}|_{\ell^2}\le1\right\}.
\]
Indeed, $|\operatorname{Tr}(\mathrm{Vq}\,Y)|\le|\mathrm{Vq}\,Y|_1\le|Y|_1$ by (a), and $\mathrm{Vq}=\mathrm{Uq}^{*}$ gives $\operatorname{Tr}(\mathrm{Uq}^{*}Y)=\operatorname{Tr}(Y^{*}Y)^{\frac{1}{2}}=|Y|_1$.
A maximum of moduli of linear functionals is subadditive; hence $|Y+Y'|_1\le|Y|_1+|Y'|_1$.
Now $O$ is the sum over $(\mathrm{mo},\mathrm{mo}')$ of the number $O(\mathrm{mo},\mathrm{mo}')$ times the matrix unit that maps the basis vector of $\mathrm{mo}'$ to the basis vector of $\mathrm{mo}$ and the other basis vectors to $0$, and a matrix unit has the singular values $1,0,\dots,0$.
Therefore
\[
|O|_1\le\sum_{\mathrm{mo},\mathrm{mo}'}|O(\mathrm{mo},\mathrm{mo}')|=\sum_{y,y'}\ \sum_{\beta,\beta'}|O(y,y')_{\beta\beta'}|\le4\sum_{y,y'}|O(y,y')|_c.
\]
\end{proof}
\end{itemize}

\subsubsection*{0.7. (P5) Continuous logarithms}

Let $f_1$ and $f_2$ be continuous complex functions on a connected topological space such that $\exp f_1=\exp f_2$, and assume that $f_1$ and $f_2$ take the same value at one point.
Then $f_1=f_2$.
Indeed, the function $(f_1-f_2)/(2\pi i)$ is continuous, takes its values in $\mathbb{Z}$, and vanishes at one point.

\subsubsection*{0.8. (P6) Elementary inequalities}

\begin{itemize}
\item[(Ei1)] $\exp(x)-1\le2x$ for $0\le x\le1$; indeed, by convexity $\exp(x)-1\le(\mathrm{e}-1)\cdot x$ on this interval.
\item[(Ei2)] $|\log(1+z)|\le|z|/(1-|z|)$ for complex $z$ with $|z|<1$ (principal branch), by the power series of the logarithm.
\item[(Ei3)] $(x+2)^3\le4\left(x^3+8\right)$ for $x\ge0$, by the convexity of the cube: $\left(\frac{x+2}{2}\right)^3\le\frac{x^3+8}{2}$.
\item[(Ei4)] $(x+3)^2\exp(-2x)\le9$ for $x\ge0$; indeed, the derivative of the left side is $-2(x+3)(x+2)\exp(-2x)<0$, and its value at $x=0$ is $9$.
\item[(Ei5)] If $L\ge3$ and $r_{*}\ge12$, then for every integer $i\ge1$
\[
L^{3i}\cdot\exp\left(-r_{*}\cdot L^{i}\right)\le\exp(-r_{*})\cdot2^{-i}.
\]
Indeed, let $x=L^{i}$; then $x\ge3^{i}\ge3i$.
Hence $x-1\ge2x/3$, so that $r_{*}(x-1)\ge8x$ and $x^3\exp(-r_{*}x)\le\exp(-r_{*})\cdot x^3\exp(-8x)$.
The function $x\mapsto x^3\exp(-8x)$ decreases for $x\ge3/8$; so its value at $x=L^{i}$ is at most its value at $3^{i}$, which is $27^{i}\exp\left(-8\cdot3^{i}\right)\le\left(27\exp(-24)\right)^{i}\le2^{-i}$.
\end{itemize}

\subsubsection*{0.9. (P7) Uniformity}

The numbers $e_{*}$, $a_{*}$, $\ell_*$, $r_{*}$, $\mu_*$ do not depend on $N\ge N_0(e)$ (Proposition~35.1\,(c)).
Hence $\mathrm{pp}$, $t_x$, $t_x^{+}$, $\tau_s$, $\varepsilon_1$, $B_1$, $B_{\mathrm{wr}}$, $B_0$ and the truth of (g1)--(g9) are functions of $(e;\bar{m};C_{\beta},\mathbf{D},D_\beta)$ and of the constants (C0)--(C9), (C6$'$) (Sections~G5.1.8 and~G5.1.9).
They involve neither $N$, nor $\theta$, nor $A$, nor the choice of $(A^{g},\mathscr{A}^{0})$, nor the placement.

Let $r'\in\{\pm1\}^3$.
By Section~I.7.5, for $r\in\{\pm1\}^3$ the operator $U_r$ is the map $f\mapsto\mathrm{Gi}\cdot f$ for $\mathrm{Gi}=r$, that is, $(U_rf)(y)=(\wedge^2r)\,f(r^{-1}y)$.
Since $\wedge^2(rr')=(\wedge^2r)(\wedge^2r')$, it follows that $U_rU_{r'}=U_{rr'}$, and $r\mapsto rr'$ is a bijection of $\{\pm1\}^3$.
Therefore the suprema over $\mathrm{Ind}_\beta$ that define $C_{\beta}$ and $\mathbf{D}$ (Section~G5.1.3) are the same for $\beta$ and for $U_{r'}\beta$.
Moreover $\operatorname{supp}U_{r'}\beta=r'\operatorname{supp}\beta$ and $r'[-2,2]^3=[-2,2]^3$, and $r'$ maps every closed coordinate cube onto a closed coordinate cube of the same side; so $D_{\mathrm{be}}$ and $D_\beta$ are the same for $\beta$ and for $U_{r'}\beta$.
Hence $\mathrm{Cb}$, $C_{\mathrm{fr}}$, $e_{\sharp}$ and all numbers of the first paragraph of (P7) are the same for $\beta$ and for $U_{r'}\beta$.

Consequently a bound in terms of the numbers of the first paragraph of (P7) and of constants is uniform in the sense of Section~G5.4.0.

\subsection*{1. The functions $\mathfrak{P}_j$}

\subsubsection*{Proof of Proposition 34.2}

\begin{proof}
Recall that the constants satisfy the inequalities
\[
L(\kappa-\kappa_0-1)\ge\kappa,\qquad c_{\mathrm{rb}}\cdot L^3\cdot\exp\left(-\left(L(\kappa-\kappa_0-1)-\kappa\right)\right)\le\frac{1}{2},
\]
which follow from the lines (C2), (C3) and stand in Section~G5.1.9;
they also satisfy the first condition of (A8), the inequalities $0<\varepsilon<1/3$, which hold by (C7), and $L>1$.

\emph{Step 0: the map $\mathscr{C}\mapsto\mathscr{C}_L$.}
Let $0\le j<k_{*}$.
Division by $L$ maps $\mathbb{L}_j=\eta_j\mathbb{Z}^3$ onto $\mathbb{L}_{j+1}=\eta_{j+1}\mathbb{Z}^3$, because $\eta_j/L=\eta_{j+1}$, and it maps $\ell_j\mathbb{Z}^3$ onto $\ell_{j+1}\mathbb{Z}^3$.
Hence $c\mapsto c/L$ is a bijection of the sites, and $b\mapsto b/L$ a bijection of the bonds, of $T_j$ onto those of $T_{j+1}$.
For a complex bond function $\mathscr{C}$ on $T_{j+1}$ one has $\mathscr{C}_L(b)=L^{-\frac{1}{2}}\cdot\mathscr{C}(b/L)$ (Section~I.3.7); in components $(\mathscr{C}_L)_\mu(c)=L^{-\frac{1}{2}}\cdot\mathscr{C}_\mu(c/L)$, and
\[
\left(\partial_\nu(\mathscr{C}_L)_\mu\right)(c)=\eta_j^{-1}\cdot L^{-\frac{1}{2}}\cdot\left[\mathscr{C}_\mu\left(c/L+\eta_{j+1}e_\nu\right)-\mathscr{C}_\mu(c/L)\right]=L^{-\frac{3}{2}}\cdot\left(\partial_\nu\mathscr{C}_\mu\right)(c/L).
\]
Consequently $|\mathscr{C}_L|_{\infty}=L^{-\frac{1}{2}}\cdot|\mathscr{C}|_{\infty}\le|\mathscr{C}|_{\infty}$ and $|\partial\mathscr{C}_L|_{\infty}=L^{-\frac{3}{2}}\cdot|\partial\mathscr{C}|_{\infty}\le|\partial\mathscr{C}|_{\infty}$, as $L\ge1$.
By (34.1), $0<e_j<e_{j+1}$, and by (P1)\,(a2) the exponents $-1+2\varepsilon$ and $-1+3\varepsilon$ are negative; hence $e_{j+1}^{-1+2\varepsilon}<e_j^{-1+2\varepsilon}$ and $e_{j+1}^{-1+3\varepsilon}<e_j^{-1+3\varepsilon}$.
Therefore
\begin{equation}\tag{34.6}
\begin{aligned}
&\text{$\mathscr{C}\mapsto\mathscr{C}_L$ is a linear map of the bond functions on $T_{j+1}$}\\
&\text{to the bond functions on $T_j$, and it maps $\mathfrak{X}_{j+1}$ into $\mathfrak{X}_j$.}
\end{aligned}
\end{equation}
(At level~$0$: for $j=0$ the set $\mathfrak{X}_0$ is formed with $q$ and the set $\mathfrak{X}_1$ with $L^{\frac{1}{2}}\cdot q$, and $T_0$ is the lattice $\mathbb{Z}^3/n_*\mathbb{Z}^3$.)

\emph{Step 1: analyticity, the value at $0$, locality.}
The three assertions are proved by induction on $j$.
For $j=0$ one has $\mathfrak{P}_0=0$.
Let $0\le j<k_{*}$, assume the three assertions for all $X\in D_j$, and let $Y\in D_{j+1}$.
By (SF4),
\begin{equation}\tag{34.7}
\mathfrak{P}_{j+1}(Y,\mathscr{C})=\sum_{X\in D_j:\ \bar{X}^{L}=LY}\left(\mathfrak{P}_j(X,\mathscr{C}_L)+E_j^{\mathrm{det}}(X,\mathscr{C}_L)\right),\qquad E_j^{\mathrm{det}}(X,A)=E^{d}_j(X,A)-E^{d}_j(X,0),
\end{equation}
and the sum is finite.
The set $\mathfrak{X}_j$ is contained in $\tilde{G}_j$ (SF2), on which $E^{d}_j(X,\cdot)$ is analytic (SF4), and $\mathfrak{P}_j(X,\cdot)$ is analytic on $\mathfrak{X}_j$.
Together with (34.6) this shows that $\mathfrak{P}_{j+1}(Y,\cdot)$ is analytic on $\mathfrak{X}_{j+1}$.
At $\mathscr{C}=0$ one has $\mathscr{C}_L=0$, $\mathfrak{P}_j(X,0)=0$ and $E_j^{\mathrm{det}}(X,0)=0$; so $\mathfrak{P}_{j+1}(Y,0)=0$.

Let $\mathscr{C},\tilde{\mathscr{C}}\in\mathfrak{X}_{j+1}$ coincide on the bonds that meet $Y$.
Let $X\in D_j$ with $\bar{X}^{L}=LY$, and let $b$ be a bond of $T_j$ that meets $X$, that is, a bond with an end point $c$ in $X$ (Section~I.1.1).
The set $\bar{X}^{L}$ is the union of the $LM$-cubes that contain an $M$-cube of $X$; so $X\subset\bar{X}^{L}=LY$, the site $c/L$ lies in $Y$, and $c/L$ is an end point of $b/L$.
Hence $b/L$ meets $Y$, and $\mathscr{C}_L(b)=L^{-\frac{1}{2}}\cdot\mathscr{C}(b/L)=L^{-\frac{1}{2}}\cdot\tilde{\mathscr{C}}(b/L)=\tilde{\mathscr{C}}_L(b)$.
Thus $\mathscr{C}_L$ and $\tilde{\mathscr{C}}_L$, both in $\mathfrak{X}_j$, coincide on the bonds that meet $X$.
By the induction hypothesis $\mathfrak{P}_j(X,\mathscr{C}_L)=\mathfrak{P}_j(X,\tilde{\mathscr{C}}_L)$, and by (SF4)\,(ii) $E^{d}_j(X,\mathscr{C}_L)=E^{d}_j(X,\tilde{\mathscr{C}}_L)$.
By (34.7), $\mathfrak{P}_{j+1}(Y,\mathscr{C})=\mathfrak{P}_{j+1}(Y,\tilde{\mathscr{C}})$.

\emph{Step 2: the bound.}
For $0\le j\le k_{*}$ put
\[
\mathrm{qP}_j:=\sup\left\{\exp\left(\kappa\cdot d_M(X)\right)\cdot\left|\mathfrak{P}_j(X,\mathscr{C})\right|:\ X\in D_j,\ d_M(X)\ge1,\ \mathscr{C}\in\mathfrak{X}_j\right\},
\]
with $\mathrm{qP}_j:=0$ if there is no such $X$; in particular $\mathrm{qP}_0=0$.
Let $0\le j<k_{*}$, let $Y\in D_{j+1}$ with $d_M(Y)\ge1$, and let $\mathscr{C}\in\mathfrak{X}_{j+1}$.
By (SF4)\,(v), every $X\in D_j$ with $\bar{X}^{L}=LY$ satisfies $d_M(X)\ge L\cdot d_M(Y)\ge1$.
Hence $\left|\mathfrak{P}_j(X,\mathscr{C}_L)\right|\le\mathrm{qP}_j\cdot\exp\left(-\kappa\cdot d_M(X)\right)$, and (SF4)\,(i), applied at $\mathscr{C}_L$ and at $0$, both of which lie in $\tilde{G}_j$, gives
\[
\left|E_j^{\mathrm{det}}(X,\mathscr{C}_L)\right|\le2C_{d}\cdot M^3\cdot\exp\left(-(\kappa+1)\cdot d_M(X)\right)\le2C_{d}\cdot M^3\cdot\exp\left(-\kappa\cdot d_M(X)\right).
\]
By (34.7) and the last inequality of (SF4)\,(v),
\begin{align*}
\left|\mathfrak{P}_{j+1}(Y,\mathscr{C})\right|&\le\left(\mathrm{qP}_j+2C_{d}\cdot M^3\right)\sum_{X:\ \bar{X}^{L}=LY}\exp\left(-\kappa\cdot d_M(X)\right)\\
&\le\left(\mathrm{qP}_j+2C_{d}\cdot M^3\right)\cdot c_{\mathrm{rb}}\cdot L^3\cdot\exp\left(-L(\kappa-\kappa_0-1)\cdot d_M(Y)\right).
\end{align*}
Since $L(\kappa-\kappa_0-1)-\kappa\ge0$ and $d_M(Y)\ge1$,
\[
\exp\left(-L(\kappa-\kappa_0-1)\cdot d_M(Y)\right)\le\exp\left(-\left(L(\kappa-\kappa_0-1)-\kappa\right)\right)\cdot\exp\left(-\kappa\cdot d_M(Y)\right).
\]
Together with $c_{\mathrm{rb}}\cdot L^3\cdot\exp\left(-\left(L(\kappa-\kappa_0-1)-\kappa\right)\right)\le\frac{1}{2}$ this yields
\[
\mathrm{qP}_{j+1}\le\frac{1}{2}\left(\mathrm{qP}_j+2C_{d}\cdot M^3\right).
\]
Starting from $\mathrm{qP}_0=0$, induction gives $\mathrm{qP}_j\le2C_{d}\cdot M^3$ for all $j\le k_{*}$, which is the asserted bound.

\emph{Level~$0$.}
For $j=0$ the recursion (34.7) reads
\[
\mathfrak{P}_1(Y,\mathscr{C})=\sum_{X\in D_0:\ \bar{X}^{L}=LY}E_0^{\mathrm{det}}(X,\mathscr{C}_L).
\]
Here $E^{d}_0$ is formed with $\acute{E}_0$, whose traces are taken on $\mathfrak{S}^{\mathrm{ch}}(1,\ell_0)$, a space of the dimension $192\,n_0$ ($48$ sites with four components over each photon site of $T_0$).
At this level Steps~0--2 require the following from Proposition~35.4:
from (SF2), the inclusion $\mathfrak{X}_0\subset\tilde{G}_0$;
from (SF4), the analyticity of $E^{d}_0(X,\cdot)$ on $\tilde{G}_0$, the relation $E_0^{\mathrm{det}}(X,A)=E^{d}_0(X,A)-E^{d}_0(X,0)$ together with $\mathfrak{P}_0=0$ and the recursion, that is, (34.7) at $j=0$, and the clauses (i), (ii) and (v).
Proposition~35.4 asserts all of these for every $j$ with $0\le j<k_{*}$, the case $j=0$ included, with the same constants $C_{d}$, $c_{\mathrm{rb}}$, $\kappa_0$.
Step~2 gives $\mathrm{qP}_1\le\frac{1}{2}\left(\mathrm{qP}_0+2C_{d}\cdot M^3\right)=C_{d}\cdot M^3$, since $\mathrm{qP}_0=0$.

\emph{Uniformity.}
The bound contains only $C_{d}$, $M$ and $\kappa$, which are constants of the lines (C2), (C5) and (C6).
It does not depend on $j$, $\ell_j$, $N$, $e$ or $\bar{m}$.
\end{proof}

\subsection*{2. The interacting factor; the functions $\mathrm{Rt}$ and $\mathrm{Lg}$}

This section contains the proof of Proposition~34.1.
Throughout it the data are those of (DAT), and the objects $\mathscr{A}$, $\mathscr{T}(t)$, $\mathrm{Ztm}(t)$ are those defined in Section~G5.4.0.

\subsubsection*{Proof of Proposition 34.1}

The conditions on the constants that enter this proof are the three conditions (g3) and the first condition of (g4), which stand in Section~G5.1.9;
the first condition of (A8), $e_0^{\varepsilon}\le\frac{1}{4}$;
the inequality $0<\varepsilon$ of line (C7) of Section~I.8.2;
and $\kappa\ge\kappa_0$, $C_0\ge1$, $M\ge1$, $\mathrm{gGa}_1\ge1$.

\begin{proof}[Proof of Proposition~34.1\,(a)]
Let $t\in\mathrm{DDo}(t_x^{+})$.
Proposition~35.4\,(DF5) shows that the bond function $\mathscr{C}(t)=a(t)^{[k_*]}$ belongs to $\mathfrak{X}_*$ and that $e_*\cdot|\mathscr{C}(t)|_\infty<e_*^{3\varepsilon}$.
If $0\le\mathrm{sv}\le1$, the point $\mathrm{sv}\cdot t$ again lies in $\mathrm{DDo}(t_x^{+})$, and $\mathrm{sv}\cdot\mathscr{C}(t)=\mathscr{C}(\mathrm{sv}\cdot t)$;
therefore $\mathrm{sv}\cdot\mathscr{C}(t)\in\mathfrak{X}_*$ as well.
By the display (2.3) of Proposition~35.4 the bond functions $\mathrm{sv}\cdot\mathscr{A}^0+c_\theta$, $0\le\mathrm{sv}\le1$, are real and belong to $\frac{1}{4}\tilde{R}_*$, hence to $\tilde{R}_*$ by (P1)\,(b);
for $\mathrm{sv}=1$ this bond function is $\mathscr{A}$, and for $\mathrm{sv}=0$ it is $c_\theta$.
According to the definition of $\mathrm{Gwk}$ in Section~G5.1.7, the sum of a real element of $\tilde{R}_*$ and of an element of $\mathfrak{X}_*$ belongs to $\mathrm{Gwk}$.
It follows that, for $0\le\mathrm{sv}\le1$, the four connections
\[
\mathscr{T}(t)=\mathscr{A}+\mathscr{C}(t),\qquad
c_\theta+\mathscr{C}(t),\qquad
\mathscr{A}+\mathrm{sv}\cdot\mathscr{C}(t),\qquad
c_\theta+\mathrm{sv}\cdot\mathscr{A}^0+\mathscr{C}(t)
\]
are elements of $\mathrm{Gwk}$.
\end{proof}

\begin{proof}[Proof of Proposition~34.1\,(b)]
Let $t\in\mathrm{DDo}(t_x^{+})$, and put
\[
\mathrm{aw}:=C_0^{-1}\cdot h_{*,2}=C_0^{-1}\cdot e_*^{-\frac{1}{4}+\varepsilon}.
\]

\emph{Step 1: the covariance.}
By (a) the connection $c_\theta+\mathscr{C}(t)$ belongs to $\mathrm{Gwk}$, hence to $\tilde{G}_*$;
so $\mathrm{Fst}\left(c_\theta+\mathscr{C}(t)\right)$ is invertible by (P1)\,(d).
This connection is the sum of the real bond function $A_r:=c_\theta$ and of $A_1:=\mathscr{C}(t)$.
From (a), (P1)\,(a1), the first condition of (g4) and the inequality $\mu_-<\mu_*$ of Proposition~35.1\,(a) one obtains
\[
e_*\cdot|A_1|_\infty<e_*^{3\varepsilon}<e_*^{2\varepsilon}\le\frac{\mu_-}{48}<\frac{\mu_*}{24}.
\]
Proposition~33.4\,(b) therefore applies:
$D_*\left(c_\theta+\mathscr{C}(t)\right)$ is invertible, and every row of the matrix of
\[
\mathrm{Cvt}:=D_*\left(c_\theta+\mathscr{C}(t)\right)^{-1}
\]
has $\ell^1$ norm at most $4\ell_*^3\cdot\mathrm{gGa}_1$.
The first condition of (g3), $4C_0\cdot\ell_*^3\cdot\mathrm{gGa}_1\le e_*^{-\frac{1}{4}+\varepsilon}$, says that $4\ell_*^3\cdot\mathrm{gGa}_1\le\mathrm{aw}$.
Moreover $\mathrm{aw}\ge4\ell_*^3\cdot\mathrm{gGa}_1\ge1$, since $\mathrm{gGa}_1\ge b_-^{-1}\ge1$.

\emph{Step 2: the integral.}
In the definition (G5.5.1) of $\mathrm{Ztm}(t)$ in Section~G5.4.0 the inserted fields are formed with $\mathscr{T}(t)$.
Splitting $D_*\left(\mathscr{T}(t)\right)$ as $\mathrm{Cvt}^{-1}+\left[D_*\left(\mathscr{T}(t)\right)-D_*\left(c_\theta+\mathscr{C}(t)\right)\right]$, one obtains
\begin{gather*}
\mathrm{Ztm}(t)=\int\mathrm{D}\Psi\,\exp\left(-\left\langle\bar{\Psi},\mathrm{Cvt}^{-1}\Psi\right\rangle+\mathrm{Ef}\right),\\
\mathrm{Ef}:=-\left\langle\bar{\Psi},\left[D_*\left(\mathscr{T}(t)\right)-D_*\left(c_\theta+\mathscr{C}(t)\right)\right]\Psi\right\rangle-m_*\cdot\langle\bar{\psi},\psi\rangle+\sum_{X\in D_*}\mathscr{E}_*\left(X;\mathscr{A},\mathscr{C}(t),\psi,\chi\right).
\end{gather*}
By (P2)\,(d), $\mathrm{Ef}$ is an even element of $\mathrm{Gr}[\Psi]$.
Hence (P3)\,(b), applied with $\mathrm{Cov}=\mathrm{Cvt}$ and $F=\mathrm{Ef}$, gives, since $\det\left(\mathrm{Cvt}^{-1}\right)=\det D_*\left(c_\theta+\mathscr{C}(t)\right)$,
\begin{equation}
\mathrm{ff}(t)=\frac{\mathrm{Ztm}(t)}{\det D_*\left(c_\theta+\mathscr{C}(t)\right)}=1+\mathrm{Rq},\qquad
\left|\mathrm{ff}(t)-1\right|\le\exp\left(|\mathrm{Ef}|_{\mathrm{aw}}\right)-1.
\tag{34.8}
\end{equation}

\emph{Step 3: the three terms of $\mathrm{Ef}$.}
\begin{itemize}
\item[(i)]
Write $\mathscr{T}(t)=A_0+B$ with $A_0:=c_\theta+\mathscr{C}(t)$ and $B:=\mathscr{A}^0$.
By (a) the segment $A_0+\mathrm{sv}\cdot B$, $0\le\mathrm{sv}\le1$, is contained in $\mathrm{Gwk}$.
The display (2.3) of Proposition~35.4 gives $\epsilon_B=e_*\cdot|\mathscr{A}^0|_\infty<\frac{1}{2}e_*^{1-\varepsilon}$, and $\#|\mathrm{Ypr}_B|\le\ell_*^3$, the number of sites of $T^0_*$.
Proposition~33.2\,(D3) yields
\[
\sum_{y,y'}\left|D_*\left(\mathscr{T}(t)\right)-D_*\left(c_\theta+\mathscr{C}(t)\right)\right|(y,y')\le C_K\cdot\epsilon_B\cdot\#|\mathrm{Ypr}_B|\le\frac{1}{2}C_K\cdot e_*^{1-\varepsilon}\cdot\ell_*^3.
\]
Together with (34.4) and $C_0\ge1$ this bounds the norm $|\cdot|_{\mathrm{aw}}$ of the first term of $\mathrm{Ef}$ by
\[
4\mathrm{aw}^2\cdot\frac{1}{2}C_K\cdot e_*^{1-\varepsilon}\cdot\ell_*^3=2C_0^{-2}\cdot C_K\cdot\ell_*^3\cdot e_*^{-\frac{1}{2}+2\varepsilon}\cdot e_*^{1-\varepsilon}\le2C_K\cdot\ell_*^3\cdot e_*^{\frac{1}{2}+\varepsilon}.
\]
\item[(ii)]
The connection $\mathscr{T}(t)$ lies in $\tilde{G}_*$.
Proposition~35.4\,(A6) gives $\left|\langle\bar{\psi},\psi\rangle\right|_{\mathrm{aw}}\le4e_*^{-\frac{1}{2}+2\varepsilon}\cdot\ell_*^3$;
here
\[
\langle\bar{\psi},\psi\rangle=\eta_*^3\cdot\frac{1}{48}\sum_{c,g,\beta}\bar{\psi}_{\beta}(c,g)\cdot\psi_{\beta}(c,g)
\]
is the form of (P2)\,(d), into which the fields $H_*\left(\mathscr{T}(t)\right)\Psi$ and $\bar{H}_*\left(\mathscr{T}(t)\right)\bar{\Psi}$ are inserted.
By Proposition~35.4\,(SF1), $|m_*|\le\frac{1}{8}e_*^{\frac{3}{4}-8\varepsilon}$.
Consequently
\[
\left|m_*\cdot\langle\bar{\psi},\psi\rangle\right|_{\mathrm{aw}}\le\frac{1}{8}e_*^{\frac{3}{4}-8\varepsilon}\cdot4e_*^{-\frac{1}{2}+2\varepsilon}\cdot\ell_*^3=\frac{1}{2}\ell_*^3\cdot e_*^{\frac{1}{4}-6\varepsilon}.
\]
\item[(iii)]
The pair $\left(\mathscr{A},\mathscr{C}(t)\right)$ lies in $\tilde{R}_*\times\mathfrak{X}_*$.
For every $X\in D_*$, Proposition~35.4\,(A6), together with the display (2.1) of the same proposition, gives
\[
\left|\mathscr{E}_*\left(X;\mathscr{A},\mathscr{C}(t),H_*^{\#}\left(\mathscr{T}(t)\right)\Psi\right)\right|_{\mathrm{aw}}
\le\sum_{n,m}\frac{h_{*,1}^n\cdot h_{*,2}^m}{n!\,m!}\left\|\mathscr{E}_{nm}\left(X;\mathscr{A},\mathscr{C}(t)\right)\right\|
\le4e_*^{\varepsilon}\cdot w^{\mathrm{E}}_*\cdot\exp\left(-\kappa\cdot d_M(X)\right);
\]
in the norms $\|\mathscr{E}_{nm}\|=\int\mathrm{d}\xi\,\mathrm{d}\boldsymbol{\zeta}\,|\mathscr{E}_{nm}|$ each $\xi$ carries the weight $\eta_*^3/48$ and each $\boldsymbol{\zeta}$ the weight $(\eta_*^3/48)^2$, and $\boldsymbol{\zeta}$ ranges over all pairs admissible for $X$, the pairs inside a cluster included.
The torus $T_*$ is the union of $(\ell_*/M)^3$ $M$-cubes $\square$ (recall $\ell_*=M\cdot r_*$ from Proposition~35.1\,(a)), every $X\in D_*$ contains one of them, and $\kappa\ge\kappa_0$.
Proposition~35.4\,(SF4)\,(v) therefore gives
\[
\sum_{X\in D_*}\exp\left(-\kappa\cdot d_M(X)\right)\le\sum_{\square}\sum_{X\supseteq\square}\exp\left(-\kappa_0\cdot d_M(X)\right)\le K_0\cdot\left(\frac{\ell_*}{M}\right)^3\le K_0\cdot\ell_*^3.
\]
Since $4e_*^{\varepsilon}\le1$ by (34.1) and $w^{\mathrm{E}}_*=e_*^{\frac{1}{4}-7\varepsilon}$, the norm of the third term of $\mathrm{Ef}$ does not exceed $K_0\cdot\ell_*^3\cdot e_*^{\frac{1}{4}-7\varepsilon}$.
\end{itemize}

\emph{Step 4: the sum.}
As $\varepsilon>0$, both exponents $\frac{1}{2}+\varepsilon$ and $\frac{1}{4}-6\varepsilon$ are larger than $\frac{1}{4}-7\varepsilon$, and (P1)\,(a1) gives
\[
e_*^{\frac{1}{2}+\varepsilon}\le e_*^{\frac{1}{4}-7\varepsilon},\qquad e_*^{\frac{1}{4}-6\varepsilon}\le e_*^{\frac{1}{4}-7\varepsilon}.
\]
Adding the three bounds of Step 3 and using the second condition of (g3), one finds
\[
|\mathrm{Ef}|_{\mathrm{aw}}\le\ell_*^3\cdot\left(2C_K+\frac{1}{2}+K_0\right)\cdot e_*^{\frac{1}{4}-7\varepsilon}\le\ell_*^3\cdot\left(2C_K+K_0+4\right)\cdot e_*^{\frac{1}{4}-7\varepsilon}\le\frac{1}{4}e_*^{-2\varepsilon}\cdot e_*^{\frac{1}{4}-7\varepsilon}=\frac{\varepsilon_1}{2}.
\]
The third condition of (g3) gives $\varepsilon_1\le\frac{1}{16}$.
With (34.8) and (Ei1) it follows that
\[
\left|\mathrm{ff}(t)-1\right|\le\exp\left(\frac{\varepsilon_1}{2}\right)-1\le\varepsilon_1=\frac{1}{2}e_*^{\frac{1}{4}-9\varepsilon}\le\frac{1}{16}.
\]
In particular $\mathrm{ff}(t)\neq0$, and therefore $\mathrm{Ztm}(t)\neq0$.

\emph{Step 5: analyticity and uniformity.}
The map $t\mapsto\mathscr{C}(t)=\sum_{\mathfrak{c}}t_{\mathfrak{c}}\cdot\hat{c}_{\mathfrak{c}}$ is linear and sends $\mathrm{DDo}(t_x^{+})$ into $\mathfrak{X}_*$.
By (34.3) and (P2)\,(d) the function $\mathrm{Ztm}(t)=\mathrm{Zin}\left(\mathscr{A},\mathscr{C}(t)\right)$ is analytic on $\mathrm{DDo}(t_x^{+})$.
The map $t\mapsto c_\theta+\mathscr{C}(t)$ is affine with values in $\tilde{G}_*$;
by (P2)\,(b) the entries of $D_*\left(c_\theta+\mathscr{C}(t)\right)$, and with them its determinant, are analytic on $\mathrm{DDo}(t_x^{+})$, and by Step 1 this determinant does not vanish.
Hence $\mathrm{ff}$ is analytic.
The bound $\varepsilon_1$ is a function of $e_*$ alone;
by (P7) it is uniform.
\end{proof}

\begin{proof}[Proof of Proposition~34.1\,(c)]
Let $t\in\mathrm{DDo}(t_x^{+})$.
By Proposition~35.4\,(DF5), $a(t)^{[j]}\in\mathfrak{X}_j$ for every $j\le k_*$, and $a(t)^{[k_*]}=\mathscr{C}(t)$.
The display (2.3) of Proposition~35.4 at $\mathrm{sv}=1$ shows that both $A+c_\theta$ and $A^{g}+c_\theta=(A+c_\theta)-d\boldsymbol{\lambda}$ lie in $\mathrm{Dom}_*$, where $\boldsymbol{\lambda}$ is the real function on $T^0_*$ provided by Proposition~35.4\,(SF5).
The gauge invariance contained in Proposition~35.4\,(SF6) gives
\[
Y^{x}_*\left(A+c_\theta;\mathscr{C}(t)\right)=Y^{x}_*\left(A^{g}+c_\theta;\mathscr{C}(t)\right).
\]
Again by (2.3), $H_*\left[A^{g}+c_\theta\right]=\mathscr{A}^0+c_\theta=\mathscr{A}$;
so (34.3) yields
\[
Y^{x}_*\left(A^{g}+c_\theta;\mathscr{C}(t)\right)=\left(\det D_*(0)\right)^{-1}\cdot\exp\left\{-\varepsilon^{\mathrm{v},t}_*\cdot\ell_*^3+\sum_{X}\mathfrak{P}_*\left(X,\mathscr{C}(t)\right)\right\}\cdot\mathrm{Ztm}(t),
\]
a number different from zero, because $\mathrm{Ztm}(t)\neq0$ by (b).
For $t=0$ one has $\mathscr{C}(0)=0$ and $\mathfrak{P}_*(X,0)=0$ (Proposition~34.2).
Dividing the expression at $t$ by the expression at $0$ gives the identity (6.1) of Proposition~34.1, namely
\[
\mathrm{Rt}(t;\theta,A)=\frac{Y^{x}_*\left(A+c_\theta;\mathscr{C}(t)\right)}{Y^{x}_*\left(A+c_\theta;0\right)}=\exp\left\{\sum_{X\in D_*}\mathfrak{P}_*\left(X,\mathscr{C}(t)\right)\right\}\cdot\frac{\mathrm{Ztm}(t)}{\mathrm{Ztm}(0)}.
\]
The quotient $Y^{x}_*\left(A+c_\theta;\mathscr{C}(t)\right)/Y^{x}_*\left(A+c_\theta;0\right)$ does not involve the pair $(A^{g},\mathscr{A}^0)$;
hence $\mathrm{Rt}$ is independent of the choice of this pair.
Being a quotient of two non-zero numbers, $\mathrm{Rt}(t;\theta,A)$ is never zero.

\emph{Continuity and analyticity.}
Let $\mathscr{A}^{u}:=H_*\left[A+c_\theta\right]$ be the photon minimiser of $A+c_\theta$, as in (2.3).
As $A+c_\theta\in\mathrm{Dom}_*$, the definition of $\mathrm{Dom}_*$ in Section~I.5.3 gives $\mathscr{A}^{u}\in\frac{1}{4}\tilde{R}_*$, a subset of $\tilde{R}_*$.
The photon minimiser $H_*[\,\cdot\,]$ is a linear map from the real bond functions on $T^0_*$ to the real bond functions on $T_*$ (\cite[Proposition~15.5]{QEDI}, at the level $k_*$ and the side $\ell_*$), and $\theta\mapsto c_\theta$ is linear.
So $(\theta,A)\mapsto\mathscr{A}^{u}$ is linear, and in particular continuous, on $\mathbb{R}^3\times V_{\mathrm{t}}(T^0_*)$.
By (34.3) with $A'=A+c_\theta$,
\[
Y^{x}_*\left(A+c_\theta;\mathscr{C}(t)\right)=\left(\det D_*(0)\right)^{-1}\cdot\exp\left\{-\varepsilon^{\mathrm{v},t}_*\cdot\ell_*^3+\sum_{X\in D_*}\mathfrak{P}_*\left(X,\mathscr{C}(t)\right)\right\}\cdot\mathrm{Zin}\left(\mathscr{A}^{u},\mathscr{C}(t)\right).
\]
Here $\mathrm{Zin}$ is analytic on the open set $\tilde{R}_*\times\mathfrak{X}_*$ by (P2)\,(d), $\mathfrak{P}_*(X,\cdot)$ is analytic on $\mathfrak{X}_*$ by Proposition~34.2, and the map $(t,\theta,A)\mapsto\left(\mathscr{A}^{u},\mathscr{C}(t)\right)$ is continuous with values in $\tilde{R}_*\times\mathfrak{X}_*$ and affine in $t$ for fixed $(\theta,A)$.
It follows that $Y^{x}_*\left(A+c_\theta;\mathscr{C}(t)\right)$ is continuous in $(t,\theta,A)$ on $\mathrm{DDo}(t_x^{+})\times\mathbb{R}^3\times\{\chi_*=1\}$ and analytic in $t$.
The denominator $Y^{x}_*\left(A+c_\theta;0\right)$ depends continuously on $(\theta,A)$ and has no zero.
Therefore $\mathrm{Rt}$ is continuous in $(t,\theta,A)$ and analytic in $t$.
\end{proof}

\begin{proof}[Proof of Proposition~34.1\,(d)]
Write $\partial_{\mathfrak{c}}:=\partial/\partial t_{\mathfrak{c}}$.
The polydisk $\mathrm{DDo}(t_x^{+})$ is convex and contains $0$.

\emph{Step 1: $\partial_{\mathfrak{c}}\mathrm{Rt}$ is continuous in $(t,\theta,A)$.}
Let $t^0\in\mathrm{DDo}(t_x^{+})$, and choose $\mathrm{rad}>0$ such that $|t^0_{\mathfrak{c}'}|+2\,\mathrm{rad}<t_x^{+}$ for every $\mathfrak{c}'$.
For $|t_{\mathfrak{c}'}-t^0_{\mathfrak{c}'}|\le\mathrm{rad}$ (all $\mathfrak{c}'$), Cauchy's formula in the variable $t_{\mathfrak{c}}$ gives
\[
\left(\partial_{\mathfrak{c}}\mathrm{Rt}\right)(t;\theta,A)=\frac{1}{2\pi i}\oint_{|\mathrm{zs}|=\mathrm{rad}}\mathrm{d}\mathrm{zs}\,\frac{\mathrm{Rt}\left(t+\mathrm{zs}\cdot\mathrm{ue}_{\mathfrak{c}};\theta,A\right)}{\mathrm{zs}^2}.
\]
By (c) the integrand is a continuous function of $(\mathrm{zs},t,\theta,A)$, and the circle is compact;
hence the integral depends continuously on $(t,\theta,A)$.

\emph{Step 2: continuity of $\mathrm{Lg}$.}
By Step 1, and because $\mathrm{Rt}$ has no zero by (c), the function
\[
\mathrm{gq}(\mathrm{sv},t;\theta,A):=\sum_{\mathfrak{c}}t_{\mathfrak{c}}\cdot\left(\frac{\partial_{\mathfrak{c}}\mathrm{Rt}}{\mathrm{Rt}}\right)(\mathrm{sv}\cdot t;\theta,A)
\]
is continuous on $[0,1]\times\mathrm{DDo}(t_x^{+})\times\mathbb{R}^3\times\{\chi_*=1\}$.
Since $[0,1]$ is compact, $\mathrm{Lg}=\int_0^1\mathrm{d}\mathrm{sv}\,\mathrm{gq}$ is continuous in $(t,\theta,A)$.
Moreover $\mathrm{Lg}(0;\theta,A)=0$, because $\mathrm{gq}(\mathrm{sv},0;\cdot)=0$.

\emph{Step 3: analyticity in $t$.}
Fix $(\theta,A)$.
For every $\mathrm{sv}$ the function $\mathrm{gq}(\mathrm{sv},\cdot)$ is analytic on $\mathrm{DDo}(t_x^{+})$.
On a closed polydisk contained in $\mathrm{DDo}(t_x^{+})$, Cauchy's formula in the $N_*$ variables $t_{\mathfrak{c}}$ represents $\mathrm{gq}(\mathrm{sv},t)$ as the integral of $\mathrm{gq}(\mathrm{sv},\cdot)$ against the Cauchy kernel over the distinguished boundary.
Integrating over $\mathrm{sv}\in[0,1]$ and interchanging the two integrations, which is legitimate because the integrand is continuous on a compact set, one obtains the same representation for $\mathrm{Lg}$.
The Cauchy kernel is analytic in $t$ in the open polydisk;
hence $\mathrm{Lg}$ is analytic in $t$.

\emph{Step 4: $\exp\mathrm{Lg}=\mathrm{Rt}$.}
Fix $(t,\theta,A)$ and let $\mathrm{uq}(\mathrm{sv}):=\mathrm{Rt}(\mathrm{sv}\cdot t;\theta,A)$ for $0\le\mathrm{sv}\le1$.
This function is continuously differentiable and has no zero;
by (6.1), $\mathrm{uq}(0)=\mathrm{Rt}(0;\theta,A)=1$;
and $\mathrm{uq}'(\mathrm{sv})=\sum_{\mathfrak{c}}t_{\mathfrak{c}}\cdot\left(\partial_{\mathfrak{c}}\mathrm{Rt}\right)(\mathrm{sv}\cdot t)$, so that $\mathrm{gq}(\mathrm{sv},t)=\mathrm{uq}'(\mathrm{sv})/\mathrm{uq}(\mathrm{sv})$.
Consequently the derivative of
\[
\mathrm{sv}\mapsto\mathrm{uq}(\mathrm{sv})\cdot\exp\left(-\int_0^{\mathrm{sv}}\mathrm{d}\mathrm{sv}'\,\mathrm{gq}(\mathrm{sv}',t)\right)
\]
vanishes identically.
Comparing the values at $\mathrm{sv}=1$ and at $\mathrm{sv}=0$, one gets $\mathrm{Rt}(t;\theta,A)\cdot\exp\left(-\mathrm{Lg}(t;\theta,A)\right)=1$.

\emph{Step 5: uniqueness.}
Let two continuous functions of $t$ on the connected set $\mathrm{DDo}(t_x^{+})$ vanish at $t=0$ and have the exponential $\mathrm{Rt}$.
By (P5) they coincide.
\end{proof}

\subsection*{3. The free ratio on the disk of radius $t_x$}

This section contains the proof of Proposition~34.3.

\subsubsection*{Proof of Proposition 34.3}

The conditions on the constants that enter this proof are the second and the fifth condition of (g4) and the condition (2.1) of Section~G5.1.9, that is $r_*>31L+16$;
the first condition of (A8), $e_0^{\varepsilon}\le\frac{1}{4}$;
the inequalities $0<\varepsilon\le\frac{1}{200}$ of line (C7) of Section~I.8.2;
the inequality $\gamma^{F}\ge0$, a consequence of $0<\gamma^{F}\le1$ (Proposition~14.2\,(f); line (C4r) of Section~I.8.2 defines $\gamma^{F}=\frac{\gamma_{z}}{2}\le\frac{1}{6L}$),
which enters the proof of Proposition~34.3\,(c) through $\|\mathrm{Kf}\|_0\le\|\mathrm{Kf}\|_{\gamma^{F}/2}$, since the norms $\|\cdot\|_\gamma$ of Section~G5.1.7 are defined for $\gamma\ge0$ and carry the factor $\exp\left(\gamma\cdot d(y,y')\right)\ge1$;
and $\kappa-1\ge\kappa_0$, $L\ge3$, $L^{1-\varepsilon}\ge2$, $M\ge2$, $C_K\le C_K'$, $\mathrm{gGa}^2\le\mathrm{gGa}_1$, $\mathrm{Cb}\ge1$.

In the whole proof $t$ denotes a complex number with $|t|<t_x^{+}$.
Then $t\cdot\mathrm{ue}_{\mathfrak{c}}\in\mathrm{DDo}(t_x^{+})$, $a(t\cdot\mathrm{ue}_{\mathfrak{c}})=\frac{t}{e}\hat{\mathfrak{a}}_{\mathfrak{c}}$ and $\mathscr{C}(t\cdot\mathrm{ue}_{\mathfrak{c}})=t\cdot\hat{c}_{\mathfrak{c}}$.
Proposition~35.4\,(DF5) and (P1)\,(a1) give
\begin{equation}
\left(\frac{t}{e}\hat{\mathfrak{a}}_{\mathfrak{c}}\right)^{[j]}\in\mathfrak{X}_j\ \ \text{for all}\ j\le k_*,\qquad
t\cdot\hat{c}_{\mathfrak{c}}\in\mathfrak{X}_*,\qquad
e_*\cdot\left|t\cdot\hat{c}_{\mathfrak{c}}\right|_\infty<e_*^{3\varepsilon}<e_*^{2\varepsilon}.
\tag{34.9}
\end{equation}
The clauses will be established in the order (c), (a), (b), (d), (e).

\begin{proof}[Proof of Proposition~34.3\,(c)]
Let $\theta\in\mathbb{R}^3$.
By (P1)\,(c), (b) and (34.9) the connections $c_\theta+\mathrm{sv}\cdot t\cdot\hat{c}_{\mathfrak{c}}$, $0\le\mathrm{sv}\le1$, belong to $\mathrm{Gwk}$, so that $D_*(\cdot)$ is defined at each of them by (P1)\,(d).
Let
\[
\mathrm{Kf}:=D_*\left(c_\theta+t\cdot\hat{c}_{\mathfrak{c}}\right)-D_*(c_\theta),\qquad
\mathrm{Yf}:=D_*(c_\theta)^{-1}\cdot\mathrm{Kf}.
\]
Since $c_\theta$ is real, Proposition~33.4\,(a) shows that $D_*(c_\theta)^{-1}$ exists and that $\left|D_*(c_\theta)^{-1}\right|_{\ell^2}\le\mathrm{gGa}^2\le\mathrm{gGa}_1$;
indeed $\mathrm{gGa}^2=b_-^{-1}+\mu_-^{-1}\le b_-^{-1}+\frac{4}{\mu_-}=\mathrm{gGa}_1$.
Apply Proposition~33.2\,(D3) with $A_0=c_\theta$ and $B=t\cdot\hat{c}_{\mathfrak{c}}$.
By (34.9), $\epsilon_B<e_*^{2\varepsilon}$.
The set $\mathrm{Ypr}_B$ is empty if $t=0$ and equals $\mathrm{Ypr}_{\mathfrak{c}}$, a subset of $\mathrm{Ypr}^{+}_{\mathfrak{c}}$, if $t\neq0$;
so $\#|\mathrm{Ypr}_B|\le\mathrm{pp}$ by Proposition~35.3\,(e).
It follows that
\[
\|\mathrm{Kf}\|_0\le\|\mathrm{Kf}\|_{\gamma^{F}/2}\le C_K\cdot e_*^{2\varepsilon},\qquad
\sum_{y,y'}|\mathrm{Kf}|(y,y')\le C_K\cdot e_*^{2\varepsilon}\cdot\mathrm{pp}.
\]
The fact (M2) of Section~G5.1.7, in the form (P4)\,(e), now gives
\[
|\mathrm{Kf}|_{\ell^2}\le2C_K\cdot e_*^{2\varepsilon},\qquad
|\mathrm{Kf}|_1\le4C_K\cdot e_*^{2\varepsilon}\cdot\mathrm{pp}.
\]
With $C_K\le C_K'$, the second condition of (g4) and (P4)\,(a) one obtains
\[
|\mathrm{Yf}|_{\ell^2}\le2\mathrm{gGa}_1\cdot C_K\cdot e_*^{2\varepsilon}\le2\mathrm{gGa}_1\cdot C_K'\cdot e_*^{2\varepsilon}\le\frac{1}{2},\qquad
|\mathrm{Yf}|_1\le4\mathrm{gGa}_1\cdot C_K\cdot\mathrm{pp}\cdot e_*^{2\varepsilon}.
\]
Hence the logarithm in $\mathrm{gD}^{\theta}(t)=\operatorname{Tr}\log\left(1+\mathrm{Yf}\right)$ is given by its power series, as in (P4)\,(b).
By (P4)\,(c),
\[
\left|\mathrm{gD}^{\theta}(t)\right|\le2|\mathrm{Yf}|_1\le8\mathrm{gGa}_1\cdot C_K\cdot\mathrm{pp}\cdot e_*^{2\varepsilon}\le16\mathrm{gGa}_1\cdot C_K'\cdot\mathrm{pp}\cdot e_*^{2\varepsilon}=B_1,
\]
and by (P4)\,(b)
\[
\exp\mathrm{gD}^{\theta}(t)=\det\left(1+\mathrm{Yf}\right)=\det\left[D_*(c_\theta)^{-1}\cdot D_*\left(c_\theta+t\cdot\hat{c}_{\mathfrak{c}}\right)\right]=\frac{\det D_*\left(c_\theta+t\cdot\hat{c}_{\mathfrak{c}}\right)}{\det D_*(c_\theta)}.
\]

\emph{Analyticity in $t$ and the value at $t=0$ (the last sentence of (c)).}
The map $t\mapsto c_\theta+t\cdot\hat{c}_{\mathfrak{c}}$ is affine, and for $|t|<t_x^{+}$ its values lie in $\tilde{G}_*$.
By (P2)\,(b) the entries of $\mathrm{Yf}$ are analytic in $t$ on this disk, on which moreover $|\mathrm{Yf}|_{\ell^2}\le\frac{1}{2}$;
so $\mathrm{gD}^{\theta}$ is analytic on $|t|<t_x^{+}$ by (P4)\,(d).
For $t=0$ one has $\mathrm{Kf}=0$, and hence $\mathrm{gD}^{\theta}(0)=0$.
\end{proof}

\begin{proof}[Proof of Proposition~34.3\,(a)]
By (34.9) the cochain $a:=\frac{t}{e}\hat{\mathfrak{a}}_{\mathfrak{c}}$ satisfies the hypothesis of Proposition~35.4\,(SF4)\,(vi), and $a^{[k_*]}=t\cdot\hat{c}_{\mathfrak{c}}$.
The first identity of the display (2.2) of Proposition~35.4, combined with (c) at $\theta=0$ (where $c_\theta$ is the zero connection), gives
\begin{align*}
\frac{\det\left(D_{(t/e)\hat{\mathfrak{a}}_{\mathfrak{c}}}+\bar{m}\eta\right)}{\det\left(D_0+\bar{m}\eta\right)}
&=\exp\left\{\sum_{X\in D_*}\mathfrak{P}_*\left(X,t\cdot\hat{c}_{\mathfrak{c}}\right)\right\}\cdot\frac{\det D_*\left(t\cdot\hat{c}_{\mathfrak{c}}\right)}{\det D_*(0)}\\
&=\exp\left\{\sum_{X}\mathfrak{P}_*\left(X,t\cdot\hat{c}_{\mathfrak{c}}\right)+\mathrm{gD}^{0}(t)\right\}=\exp\mathrm{lgN}(t);
\end{align*}
the determinants on the left are those of linear maps of $\mathfrak{S}^{\mathrm{ch}}(1;n_*)$, a space of dimension $192n_*^3$, and those in the middle are those of linear maps of $\mathfrak{S}(1,\ell_*)$, a space of dimension $4\ell_*^3$.
In the notation of Proposition~49\,(d), with $\mathrm{kt}=1$ and $\mathbf{m}=\bar{m}\eta\in(0,1]$, the left side is $N[\hat{\mathfrak{a}}_{\mathfrak{c}}](t)$, and $N_0=N[\hat{\mathfrak{a}}]$.
Here $\hat{\mathfrak{a}}$ is a real one-cochain, $\hat{\mathfrak{a}}_{\mathfrak{c}}=(\mathsf{g}_{\mathfrak{c}})_*\hat{\mathfrak{a}}$, and $\mathsf{g}_{\mathfrak{c}}$ is an element of $\mathrm{Gc}=\mathbb{Z}^3\rtimes\{\pm1\}^3$
(Sections~I.4.7 and I.6.4: in each coordinate $\mathsf{g}_{\mathfrak{c}}$ acts as $y_\mu\mapsto i\cdot\mathrm{per}+y_\mu$ or as $y_\mu\mapsto(i+1)\cdot\mathrm{per}-y_\mu$, with an integer $i$, where $\mathrm{per}$ is an integer).
Proposition~49\,(d) gives $N[\hat{\mathfrak{a}}_{\mathfrak{c}}]=N[\hat{\mathfrak{a}}]$.
Therefore $N_0(t)=\exp\mathrm{lgN}(t)$.
\end{proof}

\begin{proof}[Proof of Proposition~34.3\,(b)]
\emph{Step 1: reduction to the levels $j<k_*$.}
For $0\le j\le k_*$ let
\[
\mathrm{ch}_j:=\frac{1}{e}\hat{\mathfrak{a}}_{\mathfrak{c}}^{[j]},
\]
a real bond function on $T_j$;
one has $\mathrm{ch}_{k_*}=\hat{c}_{\mathfrak{c}}$, and $t\cdot\mathrm{ch}_j\in\mathfrak{X}_j$ by (34.9).
For $a=\frac{t}{e}\hat{\mathfrak{a}}_{\mathfrak{c}}$, whose cochains of level $j$ are $a^{[j]}=t\cdot\mathrm{ch}_j$, the second identity of (2.2) reads
\begin{equation}
\sum_{X\in D_*}\mathfrak{P}_*\left(X,t\cdot\hat{c}_{\mathfrak{c}}\right)=\sum_{j=0}^{k_*-1}\sum_{X\in D_j}E_j^{\mathrm{det}}\left(X,t\cdot\mathrm{ch}_j\right).
\tag{34.10}
\end{equation}
Fix $j<k_*$ and put $i:=k_*-j\ge1$.
Suppose that no bond $b$ with $\mathrm{ch}_j(b)\neq0$ meets $X$.
Then $t\cdot\mathrm{ch}_j$ and $0$, which both belong to $\mathfrak{X}_j$ and therefore to $\tilde{G}_j$ (Proposition~35.4\,(SF2)), agree on the bonds that meet $X$;
by Proposition~35.4\,(SF4)\,(ii), $E^{d}_j\left(X,t\cdot\mathrm{ch}_j\right)=E^{d}_j(X,0)$, that is, $E_j^{\mathrm{det}}\left(X,t\cdot\mathrm{ch}_j\right)=0$.
Thus only those $X$ that are met by a bond of the support of $\mathrm{ch}_j$ contribute to (34.10).

\emph{Step 2: the derivative of $\mathrm{ch}_j$.}
Proposition~35.4\,(DF5), applied to the cochain $\frac{1}{e}\hat{\mathfrak{a}}_{\mathfrak{c}}$ (the scaling relations stated there hold for every cochain), gives
\[
|\partial\,\mathrm{ch}_j|_\infty=L^{-\frac{3i}{2}}\cdot|\partial\hat{c}_{\mathfrak{c}}|_\infty.
\]
The number $|\partial\hat{c}_{\mathfrak{c}}|_\infty$ is one of the three norms whose sum is $\left|\frac{1}{e}\hat{\mathfrak{a}}_{\mathfrak{c}}\right|_{\mathrm{ext}}=\left|a(\mathrm{ue}_{\mathfrak{c}})\right|_{\mathrm{ext}}$, and $\left|a(\mathrm{ue}_{\mathfrak{c}})\right|_{\mathrm{ext}}\le C_\beta/e$ by Proposition~35.3\,(f) at the point $\mathrm{ue}_{\mathfrak{c}}$ of $\mathbb{C}^{\mathrm{Cel}}$.
Using $e_j=e\cdot a_j^{\frac{1}{2}}$ and $L^{-i}=a_j/a_*$ from (P1)\,(a), one obtains
\begin{equation}
e_j\cdot|\partial\,\mathrm{ch}_j|_\infty\le e\cdot a_j^{\frac{1}{2}}\cdot\left(\frac{a_j}{a_*}\right)^{\frac{3}{2}}\cdot\frac{C_\beta}{e}=a_j^2\cdot a_*^{-\frac{3}{2}}\cdot C_\beta.
\tag{34.11}
\end{equation}

\emph{Step 3: the cubes.}
Let $\mathrm{Cu}_j$ be the set of the $M$-cubes of $T_j$ that contain an end point of a bond $b$ with $\mathrm{ch}_j(b)\neq0$.
By Section~I.6.6, $\mathrm{ch}_j(b)=\frac{1}{e}L^{\frac{j}{2}}\cdot\hat{\mathfrak{a}}_{\mathfrak{c}}(L^jb)$, and $\left((g_{\mathrm{w}})_*A\right)(c,c')=A\left(g_{\mathrm{w}}^{-1}c,g_{\mathrm{w}}^{-1}c'\right)$ by Section~G5.1.2;
hence $\mathrm{ch}_j(b)\neq0$ holds if and only if $\hat{\mathfrak{a}}_{\mathfrak{c}}(L^jb)\neq0$, that is, if and only if $\hat{\mathfrak{a}}\left(\mathsf{g}_{\mathfrak{c}}^{-1}L^jb\right)\neq0$.
By Proposition~35.3\,(b2) the end points of the bonds of the support of $\hat{\mathfrak{a}}$ lie, in physical coordinates, in the closed coordinate cube $\mathrm{Qb}$ of side $D_\beta$, and $\mathsf{g}_{\mathfrak{c}}$ maps this cube onto a closed coordinate cube of side $D_\beta$.
A point with the lattice coordinates $c$ has the physical coordinates $\eta c$ and, at level $j$, the coordinates $L^{-j}c=\eta c/a_j$.
So the end points of the bonds of the support of $\mathrm{ch}_j$ lie in the image in $T_j$ of a closed coordinate cube of side $D_\beta/a_j$.
In each coordinate direction a closed interval of length $D_\beta/a_j$ meets at most $D_\beta/(Ma_j)+2$ of the intervals of length $M$ into which the $M$-cubes divide the line.
From (Ei3), $a_j=L^{-i}a_*$ and $M\ge1$ one concludes
\begin{equation}
\begin{aligned}
\#|\mathrm{Cu}_j|&\le\left(\frac{D_\beta}{Ma_j}+2\right)^3\le4\left(D_\beta^3\cdot M^{-3}\cdot a_j^{-3}+8\right),\\
\#|\mathrm{Cu}_j|&\le\left(\frac{L^iD_\beta}{Ma_*}+2\right)^3\le L^{3i}\cdot\left(\frac{D_\beta}{a_*}+2\right)^3\le L^{3i}\cdot\mathrm{pp}.
\end{aligned}
\tag{34.12}
\end{equation}
If a bond of the support of $\mathrm{ch}_j$ meets $X$, then one of its end points lies in $X$, hence in an $M$-cube of $X$, and this cube belongs to $\mathrm{Cu}_j$.
Consequently every $X$ that contributes to (34.10) at the level $j$ contains a cube of $\mathrm{Cu}_j$.
For $X\in D_j$ put $\mathrm{dq}:=d_M(X)$ and $s_X:=M\cdot(\mathrm{dq}+3)$.

\emph{Step 4: the polymers with $2s_X<\ell_j$.}
Let $\tilde{c}$ be a site of $X$ and $\mathrm{rb}:=M\cdot(\mathrm{dq}+2)+1$;
since $M\ge2$, one has $\mathrm{rb}\le s_X-1$.
By Proposition~35.4\,(SF4)\,(v) every point of $X$ is at distance at most $M\cdot(\mathrm{dq}+2)$ from $\tilde{c}$.
A bond that meets $X$ has an end point in $X$, and its other end point is at distance $\eta_j\le1$ from that one.
Hence both end points of every bond that meets $X$ lie in
\[
\mathrm{Kb}:=\left\{c\in T_j:\ d(c,\tilde{c})\le\mathrm{rb}\right\}.
\]

\emph{(4a) $\mathrm{Kb}$ does not wrap around the torus.}
Let $\underline{\tilde{c}}\in\mathbb{L}_j$ be a point over $\tilde{c}$.
The cube $\left\{\underline{c}\in\mathbb{L}_j:\ |\underline{c}-\underline{\tilde{c}}|\le\mathrm{rb}\right\}$ (sup norm) is mapped onto $\mathrm{Kb}$, its sup-diameter is at most $2\,\mathrm{rb}$, and $2\,\mathrm{rb}+\eta_j\le2(s_X-1)+1<2s_X<\ell_j$.
It is therefore a lift of $\mathrm{Kb}$ in the sense of Section~I.1.1:
the covering map is injective on it, and the bonds with both end points in $\mathrm{Kb}$ are the images of the bonds of the cube.
For $c\in\mathrm{Kb}$ let $\underline{c}$ be the point of the cube over $c$, and write $c-\tilde{c}$ for the vector $\underline{c}-\underline{\tilde{c}}$, whose sup norm is at most $\mathrm{rb}$.
Define
\[
\boldsymbol{\lambda}_0(c):=\sum_\nu\mathrm{ch}_{j,\nu}(\tilde{c})\cdot(c-\tilde{c})_\nu\quad(c\in\mathrm{Kb}),\qquad
\boldsymbol{\lambda}_0(c):=0\quad(c\notin\mathrm{Kb}).
\]
Then $\boldsymbol{\lambda}_0$ is a real function of the photon sites of $T_j$, and $(d\boldsymbol{\lambda}_0)_\nu(c)=\mathrm{ch}_{j,\nu}(\tilde{c})$ for every bond $(c,c+\eta_je_\nu)$ with both end points in $\mathrm{Kb}$.

\emph{(4b) The field $\mathrm{Bq}$.}
For a bond function $A$ on $T_j$ let $[A]_X$ denote the bond function that coincides with $A$ on the bonds meeting $X$ and vanishes on all other bonds.
Put $\mathrm{Bq}:=\left[\mathrm{ch}_j-d\boldsymbol{\lambda}_0\right]_X$, a real bond function.
For a bond $(c,c+\eta_je_\nu)$ that meets $X$ one has $\mathrm{Bq}_\nu(c)=\mathrm{ch}_{j,\nu}(c)-\mathrm{ch}_{j,\nu}(\tilde{c})$.
Join $\underline{\tilde{c}}$ to $\underline{c}$ inside the cube by a rectilinear lattice path;
it has $|\underline{c}-\underline{\tilde{c}}|_1/\eta_j\le3\,\mathrm{rb}/\eta_j$ steps, and along a step in the direction $\pm e_\mu$ the number $\mathrm{ch}_{j,\nu}$ changes by $\pm\eta_j$ times a value of $\partial_\mu\mathrm{ch}_{j,\nu}$.
Therefore
\begin{equation}
|\mathrm{Bq}|_\infty\le3\,\mathrm{rb}\cdot|\partial\,\mathrm{ch}_j|_\infty\le3s_X\cdot|\partial\,\mathrm{ch}_j|_\infty.
\tag{34.13}
\end{equation}

\emph{(4c) Removal of the constant part.}
Let
\begin{gather*}
\mathrm{Dq}:=\left\{\mathrm{zs}\in\mathbb{C}:\ \left|\left[t\cdot\mathrm{ch}_j-\mathrm{zs}\cdot d\boldsymbol{\lambda}_0\right]_X\right|_\infty<e_j^{-1+\varepsilon}\right\},\\
\mathrm{fq}(\mathrm{zs}):=E^{d}_j\left(X,\left[t\cdot\mathrm{ch}_j-\mathrm{zs}\cdot d\boldsymbol{\lambda}_0\right]_X\right)\quad(\mathrm{zs}\in\mathrm{Dq}).
\end{gather*}
The set $\mathrm{Dq}$ is the intersection, over the finitely many bonds $b$ that meet $X$, of the sets $\left\{\mathrm{zs}:\ \left|t\cdot\mathrm{ch}_j(b)-\mathrm{zs}\cdot(d\boldsymbol{\lambda}_0)(b)\right|<e_j^{-1+\varepsilon}\right\}$.
Each of these sets contains $0$, because $|t\cdot\mathrm{ch}_j(b)|<e_j^{-1+2\varepsilon}<e_j^{-1+\varepsilon}$ (recall $t\cdot\mathrm{ch}_j\in\mathfrak{X}_j$, and use (P1)\,(a1));
it is an open disk if $(d\boldsymbol{\lambda}_0)(b)\neq0$ and the whole plane if $(d\boldsymbol{\lambda}_0)(b)=0$.
So $\mathrm{Dq}$ is open and convex, and $0\in\mathrm{Dq}$.
For $\mathrm{zs}\in\mathrm{Dq}$ the field $\left[t\cdot\mathrm{ch}_j-\mathrm{zs}\cdot d\boldsymbol{\lambda}_0\right]_X$ belongs to $G_j$, a subset of $\tilde{G}_j$ (Proposition~35.4\,(SF2)), and it depends affinely on $\mathrm{zs}$;
so $\mathrm{fq}$ is analytic on $\mathrm{Dq}$ by Proposition~35.4\,(SF4).

Now let $\mathrm{zs}\in\mathrm{Dq}$ be real.
The field $t\cdot\mathrm{ch}_j-d(\mathrm{zs}\cdot\boldsymbol{\lambda}_0)$ lies in $\tilde{G}_j$, because $t\cdot\mathrm{ch}_j$ does and $\tilde{G}_j$ is invariant under real gauge transformations (Proposition~35.4\,(SF2));
on the bonds that meet $X$ it agrees with $\left[t\cdot\mathrm{ch}_j-\mathrm{zs}\cdot d\boldsymbol{\lambda}_0\right]_X$.
Proposition~35.4\,(SF4)\,(ii) and (iii) give
\[
\mathrm{fq}(\mathrm{zs})=E^{d}_j\left(X,t\cdot\mathrm{ch}_j-d(\mathrm{zs}\cdot\boldsymbol{\lambda}_0)\right)=E^{d}_j\left(X,t\cdot\mathrm{ch}_j\right).
\]
The real points of $\mathrm{Dq}$ form an open interval containing $0$, and an analytic function on a connected open subset of $\mathbb{C}$ that is constant on an interval is constant (identity theorem).
Hence
\begin{equation}
\mathrm{fq}(\mathrm{zs})=E^{d}_j\left(X,t\cdot\mathrm{ch}_j\right)\qquad\text{for all}\ \mathrm{zs}\in\mathrm{Dq}.
\tag{34.14}
\end{equation}

\emph{(4d) Schwarz's lemma.}
Let
\begin{gather*}
\mathrm{rad}_X:=\frac{e_j^{-1+\varepsilon}}{3s_X\cdot|\partial\,\mathrm{ch}_j|_\infty}\quad\left(\mathrm{rad}_X:=+\infty\ \text{if}\ \partial\,\mathrm{ch}_j=0\right),\\
\mathrm{vq}(\mathrm{zs}):=E^{d}_j\left(X,\mathrm{zs}\cdot\mathrm{Bq}\right)-E^{d}_j(X,0)\quad\left(|\mathrm{zs}|<\mathrm{rad}_X\right).
\end{gather*}
For $|\mathrm{zs}|<\mathrm{rad}_X$, (34.13) gives $|\mathrm{zs}\cdot\mathrm{Bq}|_\infty<e_j^{-1+\varepsilon}$;
hence $\mathrm{zs}\cdot\mathrm{Bq}\in G_j$, the function $\mathrm{vq}$ is analytic on $|\mathrm{zs}|<\mathrm{rad}_X$, and
\[
|\mathrm{vq}(\mathrm{zs})|\le2C_d\cdot M^3\cdot\exp\left(-(\kappa+1)\cdot\mathrm{dq}\right)
\]
by Proposition~35.4\,(SF4)\,(i).
Moreover $\mathrm{vq}(0)=0$, and $\mathrm{vq}(-\mathrm{zs})=\mathrm{vq}(\mathrm{zs})$ by Proposition~35.4\,(SF4)\,(iv);
so $\mathrm{vq}$ has a zero of order at least two at $0$, and $\mathrm{vq}(\mathrm{zs})/\mathrm{zs}^2$ is analytic on $|\mathrm{zs}|<\mathrm{rad}_X$.
The maximum principle on the disks $|\mathrm{zs}|\le\mathrm{rad}'$ with $\mathrm{rad}'<\mathrm{rad}_X$, followed by the limit $\mathrm{rad}'\to\mathrm{rad}_X$, gives
\[
|\mathrm{vq}(\mathrm{zs})|\le2C_d\cdot M^3\cdot\exp\left(-(\kappa+1)\cdot\mathrm{dq}\right)\cdot\left(\frac{|\mathrm{zs}|}{\mathrm{rad}_X}\right)^2\qquad\left(|\mathrm{zs}|<\mathrm{rad}_X\right).
\]
(If $\mathrm{rad}_X=+\infty$, then $\mathrm{Bq}=0$ by (34.13) and $\mathrm{vq}=0$, and the quotient $|\mathrm{zs}|/\mathrm{rad}_X$ is read as $0$.)

Suppose first that $|t|<\mathrm{rad}_X$.
Then $\left[t\cdot\mathrm{ch}_j-t\cdot d\boldsymbol{\lambda}_0\right]_X=t\cdot\mathrm{Bq}$ has sup norm less than $e_j^{-1+\varepsilon}$, so $t\in\mathrm{Dq}$, and (34.14) gives
\[
E_j^{\mathrm{det}}\left(X,t\cdot\mathrm{ch}_j\right)=\mathrm{fq}(t)-E^{d}_j(X,0)=\mathrm{vq}(t).
\]
If instead $|t|\ge\mathrm{rad}_X$, then $\left|E_j^{\mathrm{det}}\left(X,t\cdot\mathrm{ch}_j\right)\right|\le2C_d\cdot M^3\cdot\exp\left(-(\kappa+1)\cdot\mathrm{dq}\right)$ by Proposition~35.4\,(SF4)\,(i), while $\left(|t|/\mathrm{rad}_X\right)^2\ge1$.
In both cases, using (34.11),
\begin{equation}
\begin{aligned}
\left|E_j^{\mathrm{det}}\left(X,t\cdot\mathrm{ch}_j\right)\right|&\le2C_d\cdot M^3\cdot\exp\left(-(\kappa+1)\cdot\mathrm{dq}\right)\cdot\left(\frac{|t|}{\mathrm{rad}_X}\right)^2,\\
\left(\frac{|t|}{\mathrm{rad}_X}\right)^2&=9M^2\cdot(\mathrm{dq}+3)^2\cdot|t|^2\cdot e_j^{-2\varepsilon}\cdot\left(e_j\cdot|\partial\,\mathrm{ch}_j|_\infty\right)^2\\
&\le9M^2\cdot(\mathrm{dq}+3)^2\cdot|t|^2\cdot e_j^{-2\varepsilon}\cdot a_j^4\cdot a_*^{-3}\cdot C_\beta^2.
\end{aligned}
\tag{34.15}
\end{equation}

\emph{(4e) The sum over these polymers.}
By (Ei4) and $\kappa-1\ge\kappa_0$,
\[
(\mathrm{dq}+3)^2\cdot\exp\left(-(\kappa+1)\cdot\mathrm{dq}\right)\le9\exp\left(-(\kappa-1)\cdot\mathrm{dq}\right)\le9\exp\left(-\kappa_0\cdot\mathrm{dq}\right).
\]
Proposition~35.4\,(SF4)\,(v) therefore gives, for every $M$-cube $\square$,
\[
\sum_{X\supseteq\square}\left(d_M(X)+3\right)^2\cdot\exp\left(-(\kappa+1)\cdot d_M(X)\right)\le9K_0.
\]
By Step 3, (34.15) and the first bound of (34.12), the polymers $X\in D_j$ with $2s_X<\ell_j$ contribute to (34.10) at most
\begin{multline*}
2C_d\cdot M^3\cdot9M^2\cdot|t|^2\cdot e_j^{-2\varepsilon}\cdot a_j^4\cdot a_*^{-3}\cdot C_\beta^2\cdot9K_0\cdot4\left(D_\beta^3\cdot M^{-3}\cdot a_j^{-3}+8\right)\\
=648C_d\cdot K_0\cdot C_\beta^2\cdot|t|^2\cdot e_j^{-2\varepsilon}\cdot a_*^{-3}\cdot\left(M^2\cdot D_\beta^3\cdot a_j+8M^5\cdot a_j^4\right).
\end{multline*}

\emph{Step 5: the polymers with $2s_X\ge\ell_j$.}
For these polymers $\mathrm{dq}\ge\ell_j/(2M)-3=r_*\cdot L^i/2-3$, because $\ell_j=L^i\cdot\ell_*$ and $\ell_*=M\cdot r_*$ (Proposition~35.1\,(a)).
By Proposition~35.4\,(SF4)\,(i),
\begin{multline*}
\left|E_j^{\mathrm{det}}\left(X,t\cdot\mathrm{ch}_j\right)\right|\le2C_d\cdot M^3\cdot\exp\left(-(\kappa+1)\cdot\mathrm{dq}\right)\le2C_d\cdot M^3\cdot\exp\left(-2\,\mathrm{dq}\right)\cdot\exp\left(-\kappa_0\cdot\mathrm{dq}\right)\\
\le2C_d\cdot M^3\cdot\exp\left(6-r_*\cdot L^i\right)\cdot\exp\left(-\kappa_0\cdot\mathrm{dq}\right).
\end{multline*}
By Step 3, Proposition~35.4\,(SF4)\,(v) and the second bound of (34.12), these polymers contribute to (34.10) at most
\[
2\mathrm{e}^6\cdot C_d\cdot M^3\cdot K_0\cdot\mathrm{pp}\cdot L^{3i}\cdot\exp\left(-r_*\cdot L^i\right).
\]

\emph{Step 6: the sum over the levels.}
By (P1)\,(a), $a_j=a_*\cdot L^{-i}$ and $e_j^{-2\varepsilon}=e_*^{-2\varepsilon}\cdot L^{i\varepsilon}$.
Hence
\[
e_j^{-2\varepsilon}\cdot a_*^{-3}\cdot a_j=e_*^{-2\varepsilon}\cdot a_*^{-2}\cdot L^{-i(1-\varepsilon)},\qquad
e_j^{-2\varepsilon}\cdot a_*^{-3}\cdot a_j^4=e_*^{-2\varepsilon}\cdot a_*\cdot L^{-i(4-\varepsilon)}.
\]
The inequality $L^{1-\varepsilon}\ge2$ (a consequence of $L\ge5$ and $\varepsilon\le\frac{1}{200}$) gives
\[
\sum_{i\ge1}L^{-i(4-\varepsilon)}\le\sum_{i\ge1}L^{-i(1-\varepsilon)}\le\sum_{i\ge1}2^{-i}=1.
\]
Since $a_*\in\{h/L,h/L^2\}$ (Proposition~35.1\,(a)), $a_*^{-2}\le L^4\cdot h^{-2}$ and $a_*\le h$.
So the contributions of Step 4, summed over $i=1,\dots,k_*$, are at most
\[
648C_d\cdot K_0\cdot C_\beta^2\cdot\left(M^2\cdot D_\beta^3\cdot L^4\cdot h^{-2}+8M^5\cdot h\right)\cdot e_*^{-2\varepsilon}\cdot|t|^2=\frac{1}{4}C_{\mathrm{fr}}\cdot e_*^{-2\varepsilon}\cdot|t|^2.
\]
The condition (2.1) of Section~G5.1.9 holds by Proposition~35.1\,(e);
together with $L\ge3$ it gives $r_*>31L+16\ge12$.
By (Ei5) the contributions of Step 5 are at most
\[
2\mathrm{e}^6\cdot C_d\cdot M^3\cdot K_0\cdot\mathrm{pp}\cdot\exp(-r_*)\cdot\sum_{i\ge1}2^{-i}=\frac{1}{8}B_{\mathrm{wr}}.
\]
Adding the two contributions one finds
\[
\left|\sum_{X\in D_*}\mathfrak{P}_*\left(X,t\cdot\hat{c}_{\mathfrak{c}}\right)\right|\le C_{\mathrm{fr}}\cdot e_*^{-2\varepsilon}\cdot|t|^2+B_{\mathrm{wr}}.
\]

\emph{The term $j=0$ of (34.10).}
For this term $i=k_*$, and the objects of level $0$ take the following form:
$\eta_0=1$, $a_0=\eta$, the charge $e_j$ of the level $j=0$ is $q$ (see (P1)\,(a)), $\ell_0=n_*$, $\mathrm{ch}_0=\frac{1}{e}\hat{\mathfrak{a}}_{\mathfrak{c}}$ and $T_0=\mathbb{Z}^3/n_*\mathbb{Z}^3$;
the sets $G_0$, $\mathfrak{X}_0$, $\tilde{G}_0$ are formed with $q$;
and $E^{d}_0(X,\cdot)$ is formed with $\acute{E}_0$, that is, with traces over $\mathfrak{S}^{\mathrm{ch}}(1,\ell_0)$ ($192$ components per photon site).
In Steps~1--5 the function $E^{d}_j(X,\cdot)$ enters through (SF2) and (SF4) of Proposition~35.4, in particular through the analyticity of $E^{d}_j(X,\cdot)$ on $\tilde{G}_j$ and through the clauses (i)--(v) of (SF4);
these statements hold for $j=0$ with the same constants $C_d$, $K_0$, $\kappa_0$.
For $j=0$ the function $\boldsymbol{\lambda}_0$ of (4a) is a function of the photon sites of $T_0$, so that (SF4)\,(iii) applies to it,
and in Step~4 the second end point of a bond that meets $X$ is at distance $\eta_0=1$ from the first.
Consequently Steps~1--5 apply to the term $j=0$ as well.
\end{proof}

\begin{proof}[Proof of Proposition~34.3\,(d)]
Let $|t|\le t_x$;
then $|t|<t_x^{+}$.
Since $\mathrm{Cb}\ge1$, one has $t_x=e_*^{3\varepsilon}/\mathrm{Cb}\le e_*^{3\varepsilon}$, and therefore $C_{\mathrm{fr}}\cdot e_*^{-2\varepsilon}\cdot|t|^2\le C_{\mathrm{fr}}\cdot e_*^{4\varepsilon}$.
From (b), from (c) at $\theta=0$ and from the fifth condition of (g4) it follows that
\[
\left|\mathrm{lgN}(t)\right|\le C_{\mathrm{fr}}\cdot e_*^{4\varepsilon}+B_{\mathrm{wr}}+B_1=B_0\le\frac{1}{4}.
\]

\emph{Analyticity and the value at $0$.}
The map $t\mapsto t\cdot\hat{c}_{\mathfrak{c}}$ is linear with values in $\mathfrak{X}_*$ by (34.9), and $\mathfrak{P}_*(X,\cdot)$ is analytic on $\mathfrak{X}_*$ with $\mathfrak{P}_*(X,0)=0$ (Proposition~34.2);
by (c), $\mathrm{gD}^{0}$ is analytic on $|t|<t_x^{+}$, and $\mathrm{gD}^{0}(0)=0$.
Hence $\mathrm{lgN}$ is analytic on $|t|<t_x^{+}$, and $\mathrm{lgN}(0)=0$.

\emph{Evenness.}
Proposition~49\,(d) gives $N_0(-t)=N_0(t)$.
By (a), the continuous functions $t\mapsto\mathrm{lgN}(-t)$ and $\mathrm{lgN}$ on the disk $|t|<t_x^{+}$ have the same exponential, and both vanish at $0$.
By (P5), $\mathrm{lgN}(-t)=\mathrm{lgN}(t)$.

\emph{Reality on the real axis.}
For real $t\in(-t_x^{+},t_x^{+})$ the number $N_0(t)=\exp\mathrm{lgN}(t)$ is real by Proposition~49\,(d), and it is different from zero;
as $N_0(0)=1$, continuity gives $N_0(t)>0$ on this interval.
Hence $\exp\left(i\operatorname{Im}\mathrm{lgN}(t)\right)=N_0(t)/|N_0(t)|=1$, so that the continuous function $\operatorname{Im}\mathrm{lgN}/(2\pi)$ takes integer values on the interval and vanishes at $0$.
Thus $\operatorname{Im}\mathrm{lgN}(t)=0$.

\emph{Independence of $\mathfrak{c}$.}
The function $N_0$ is formed with $\hat{\mathfrak{a}}$ and does not involve $\mathfrak{c}$.
By (a), for every $\mathfrak{c}$ the function $\mathrm{lgN}$ formed with $\mathfrak{c}$ is continuous on the disk $|t|<t_x^{+}$, has the exponential $N_0$, and vanishes at $0$.
By (P5) all these functions coincide.
\end{proof}

\begin{proof}[Proof of Proposition~34.3\,(e)]
By Proposition~49\,(d), $N_0$ is entire, $N_0(0)=1$, $N_0(-t)=N_0(t)$ and $\overline{N_0(\overline{t})}=N_0(t)$.
Hence $\mathrm{nj}_0=1$ and $\mathrm{nj}_1=\mathrm{nj}_3=0$, and all $\mathrm{nj}_i$ are real.
For $m\le4$ the letter $c_m$ denotes here the number $N_0^{(m)}(0)$, the $m$-th derivative of $N_0$ at $0$, and not a lattice site;
one has $c_m=\phi_0^{(m)}(0)=m!\cdot\mathrm{nj}_m$, and in particular $c_1=c_3=0$.
On $|t|\le t_x$, (a), (d) and (Ei1) give
\[
\left|N_0(t)-1\right|=\left|\exp\mathrm{lgN}(t)-1\right|\le\exp(B_0)-1\le2B_0.
\]
Cauchy's estimate on the circle $|t|=t_x$ yields $|\mathrm{nj}_i|\le2B_0\cdot t_x^{-i}$ for every $i\ge1$.
With $t_x=4\mathbf{r}$ and $B_0\le\frac{1}{4}$ one obtains
\begin{gather*}
|c_2|=2|\mathrm{nj}_2|\le\frac{4B_0}{16\mathbf{r}^2}\le\frac{\mathbf{r}^{-2}}{16},\qquad
|c_4|=24|\mathrm{nj}_4|\le\frac{48B_0}{256\mathbf{r}^4}\le\frac{3}{64}\mathbf{r}^{-4}\le\frac{\mathbf{r}^{-4}}{16},\\
\|\phi_0-1\|_{\mathbf{r}}=\sum_{i=1}^{4}|\mathrm{nj}_i|\cdot\mathbf{r}^i\le2B_0\sum_{i\ge1}4^{-i}=\frac{2}{3}B_0\le\frac{1}{6},\qquad
\|\phi_0\|_{\mathbf{r}}\le1+\frac{1}{6}\le2.
\end{gather*}

\emph{The last relation of (6.4) in Proposition~34.3.}
Near $t=0$,
\[
N_0(t)=1+\frac{c_2}{2}t^2+\frac{c_4}{24}t^4+O(t^6),\qquad
\log N_0(t)=\frac{c_2}{2}t^2+\left(\frac{c_4}{24}-\frac{c_2^2}{8}\right)t^4+O(t^6),
\]
so that
\[
\left.\left(\frac{\mathrm{d}}{\mathrm{d}t}\right)^4\log N_0(t)\right|_{t=0}=c_4-3c_2^2.
\]
By Section~I.7.6 this derivative is $b_{\mathrm{ch}}(\eta,R^{*}(e))$.

\emph{The function $N_0$ does not depend on the vector $v\in\eta\mathbb{Z}^3$ with which $\hat{\mathfrak{a}}$ is formed.}
In particular the number $c_4-3c_2^2$ is the same for the placement at hand, for the centred placements of Section~I.7.6 and for $v=0$.
Let $v,v'\in\eta\mathbb{Z}^3$, let $g_{\mathrm{w}}$ be the translation of $\mathbb{Z}^3$ by $(v'-v)/\eta$, an element of $\mathrm{Gc}$, and let $\hat{\mathfrak{a}}$, $\hat{\mathfrak{a}}'$ be the lattice probes formed with $\tau_v\beta$ and $\tau_{v'}\beta$.
One has $(\tau_v\beta)(y)=\beta(y-v)$ (Section~I.7.5), and $J_\eta$ commutes with translations.
In the definitions of Section~G5.1.3 this gives, for $c\in\Lambda_{\eta,R^{*}(e)}$,
\[
u_{\mu\nu}[\tau_{v'}\beta](c)=u_{\mu\nu}[\tau_v\beta]\left(c-(v'-v)\right),\qquad
a_\eta[\tau_{v'}\beta]_\nu(c)=a_\eta[\tau_v\beta]_\nu\left(c-(v'-v)\right).
\]
On the bonds of $\mathbb{Z}^3/n_*\mathbb{Z}^3$ this reads $\hat{\mathfrak{a}}'(c,c+e_\nu)=\hat{\mathfrak{a}}\left(g_{\mathrm{w}}^{-1}c,g_{\mathrm{w}}^{-1}c+e_\nu\right)$, that is, $\hat{\mathfrak{a}}'=(g_{\mathrm{w}})_*\hat{\mathfrak{a}}$.
Proposition~49\,(d) gives $N[\hat{\mathfrak{a}}']=N[\hat{\mathfrak{a}}]$.
\end{proof}

\begin{proof}[Uniformity of the bounds of Proposition~34.3]
The bounds in (b)--(e) are expressed through $C_{\mathrm{fr}}$, $B_{\mathrm{wr}}$, $B_1$, $B_0$, $t_x$, $\mathbf{r}$ and $e_{*}$ alone.
By (P7) they are therefore uniform.
\end{proof}

\subsection*{4. One cell; the twisted free ratio}

This section contains the proof of Proposition~34.4.

\subsubsection*{Proof of Proposition 34.4}

Of the conditions on the constants, the proof uses the third condition of (g3) and the first four conditions of (g4), which stand in Section~G5.1.9;
the inequality $c_{\tau}\le\mu_-/48$;
the inequality $\gamma^{F}\ge0$, which holds since $0<\gamma^{F}\le1$ (line (C4r) of Section~I.8.2);
and, through (P1)\,(a1), the first condition of (A8), $e_0^{\varepsilon}\le\frac{1}{4}$, together with $0<\varepsilon$ (line (C7) of Section~I.8.2).
Through Proposition~34.1 it rests in addition on (g3) and on the first condition of (g4).

\begin{proof}[Proof of Proposition~34.4\,(a)]
Fix $t\in\mathbb{C}$ with $|t|<t_x^{+}$; then (34.9) holds.
For $\theta\in\mathbb{C}^3$ the twist connection $c_{\theta}$ is the constant bond function on $T_*$ whose components are $\theta_\mu/(e_{*}\cdot\ell_*)$ (Section~I.6.6).
Hence $\theta\mapsto c_{\theta}$ is complex linear, $c_{\theta}=c_{\operatorname{Re}\theta}+i\cdot c_{\operatorname{Im}\theta}$,
\[
\left|c_{\theta}\right|_{\infty}=\max_{\mu}\frac{|\theta_\mu|}{e_{*}\cdot\ell_*},
\]
and every difference quotient $\partial_\nu\left(c_{\theta}\right)_\mu$ vanishes.
Let
\[
\mathrm{wd}_s:=\ell_*\cdot\left(e_{*}^{2\varepsilon}-e_{*}^{3\varepsilon}\right),
\]
a positive number by (P1)\,(a1), and define the following subsets of $\mathbb{C}^3$:
\begin{align*}
\mathrm{Stp}&:=\left\{\theta\in\mathbb{C}^3:\ \left|\operatorname{Im}\theta_\mu\right|\le\tau_s\ \text{for}\ \mu=1,2,3\right\},\\
\mathrm{Stp}^{+}&:=\left\{\theta\in\mathbb{C}^3:\ \left|\operatorname{Im}\theta_\mu\right|\le\tau_s+\mathrm{wd}_s\ \text{for}\ \mu=1,2,3\right\},\\
\mathrm{Osp}&:=\left\{\theta\in\mathbb{C}^3:\ \left|\operatorname{Im}\theta_\mu\right|<\tau_s+\mathrm{wd}_s\ \text{for}\ \mu=1,2,3\right\}.
\end{align*}
The closed strip of the statement is $\mathrm{Stp}$.
The set $\mathrm{Osp}$ is open and convex;
it contains $\mathrm{Stp}$ and $\mathbb{R}^3$, it is contained in $\mathrm{Stp}^{+}$, and each translation $\theta\mapsto\theta+2\pi e_\mu$ maps it onto itself.
In Steps~1--3, $\theta$ denotes an arbitrary point of $\mathrm{Stp}^{+}$.

\emph{Step~1: the connections.}
Let $\theta\in\mathrm{Stp}^{+}$ and $0\le\mathrm{sv}\le1$.
Define $\mathrm{thi}\in\mathbb{R}^3$ by
\[
\mathrm{thi}_\mu:=\max\left(-\tau_s,\min\left(\tau_s,\operatorname{Im}\theta_\mu\right)\right)\qquad(\mu=1,2,3),
\]
and put $\mathrm{tho}:=\operatorname{Im}\theta-\mathrm{thi}$.
Then $\left|\mathrm{thi}_\mu\right|\le\tau_s$ and $\left|\mathrm{tho}_\mu\right|\le\mathrm{wd}_s$;
for $\theta\in\mathrm{Stp}$ one has $\mathrm{thi}=\operatorname{Im}\theta$ and $\mathrm{tho}=0$.
The linearity of $\theta\mapsto c_{\theta}$ gives the decomposition
\[
c_{\theta}+\mathrm{sv}\cdot t\cdot\hat{c}_{\mathfrak{c}}=c_{\operatorname{Re}\theta}+\left(\mathrm{sv}\cdot t\cdot\hat{c}_{\mathfrak{c}}+i\cdot c_{\mathrm{tho}}\right)+i\cdot c_{\mathrm{thi}},
\]
whose three terms have the following properties.
\begin{itemize}
\item[(1)] The bond function $c_{\operatorname{Re}\theta}$ is real, and it lies in $\frac{1}{4}\tilde{R}_*$, a subset of $\tilde{R}_*$ ((P1)\,(c), (b)).
\item[(2)] The bond function $\mathrm{sv}\cdot t\cdot\hat{c}_{\mathfrak{c}}+i\cdot c_{\mathrm{tho}}$ lies in $\mathfrak{X}_*$.
Indeed, by (34.9) its sup norm is at most
\[
\left|t\cdot\hat{c}_{\mathfrak{c}}\right|_{\infty}+\max_{\mu}\frac{\left|\mathrm{tho}_\mu\right|}{e_{*}\cdot\ell_*}<e_{*}^{-1+3\varepsilon}+\frac{\mathrm{wd}_s}{e_{*}\cdot\ell_*}=e_{*}^{-1+3\varepsilon}+e_{*}^{-1}\cdot\left(e_{*}^{2\varepsilon}-e_{*}^{3\varepsilon}\right)=e_{*}^{-1+2\varepsilon}.
\]
Its difference quotients are $\mathrm{sv}$ times those of $t\cdot\hat{c}_{\mathfrak{c}}$, and $t\cdot\hat{c}_{\mathfrak{c}}$ lies in $\mathfrak{X}_*$ by (34.9);
consequently
\[
\left|\partial\left(\mathrm{sv}\cdot t\cdot\hat{c}_{\mathfrak{c}}+i\cdot c_{\mathrm{tho}}\right)\right|_{\infty}\le\left|\partial\left(t\cdot\hat{c}_{\mathfrak{c}}\right)\right|_{\infty}<e_{*}^{-1+3\varepsilon}.
\]
\item[(3)] Since $\tau_s=c_{\tau}\cdot\ell_*$ (Section~G5.1.8),
\[
e_{*}\cdot\left|i\cdot c_{\mathrm{thi}}\right|_{\infty}=\max_{\mu}\frac{\left|\mathrm{thi}_\mu\right|}{\ell_*}\le\frac{\tau_s}{\ell_*}=c_{\tau}.
\]
\end{itemize}
By (1), (2) and the definition of $\mathrm{Gwk}$ (Section~G5.1.7), the bond function $c_{\operatorname{Re}\theta}+\left(\mathrm{sv}\cdot t\cdot\hat{c}_{\mathfrak{c}}+i\cdot c_{\mathrm{tho}}\right)$ belongs to $\mathrm{Gwk}$.
By (3) and the definition of $\mathrm{Gwk}^{+}$ it follows that
\[
c_{\theta}+\mathrm{sv}\cdot t\cdot\hat{c}_{\mathfrak{c}}\in\mathrm{Gwk}^{+}\qquad\left(\theta\in\mathrm{Stp}^{+},\ 0\le\mathrm{sv}\le1\right).
\]
By Proposition~33.3 the operator $\mathrm{Fst}\left(c_{\theta}+\mathrm{sv}\cdot t\cdot\hat{c}_{\mathfrak{c}}\right)$ is invertible.
Thus these connections belong to $\mathrm{Inv}$, and at them $D_*(\,\cdot\,)$ is defined by the formula of Section~G5.1.7.

\emph{Step~2: the inverse of $D_*(c_{\theta})$.}
Let $\theta\in\mathrm{Stp}^{+}$.
Write $c_{\theta}=A_r+A_1$ with the real bond function $A_r:=c_{\operatorname{Re}\theta}$ and with $A_1:=i\cdot c_{\operatorname{Im}\theta}$.
The definition of $c_{\tau}$ gives $c_{\tau}\le\mu_-/48$, the first condition of (g4) states $e_{*}^{2\varepsilon}\le\mu_-/48$, and $\mu_-<\mu_*$ by Proposition~35.1\,(a).
Therefore
\[
e_{*}\cdot\left|A_1\right|_{\infty}=\max_{\mu}\frac{\left|\operatorname{Im}\theta_\mu\right|}{\ell_*}\le\frac{\tau_s+\mathrm{wd}_s}{\ell_*}=c_{\tau}+e_{*}^{2\varepsilon}-e_{*}^{3\varepsilon}<\frac{\mu_-}{48}+\frac{\mu_-}{48}=\frac{\mu_-}{24}<\frac{\mu_*}{24}.
\]
(For $\theta\in\mathrm{Stp}$ one has directly $e_{*}\cdot\left|A_1\right|_{\infty}\le\tau_s/\ell_*=c_{\tau}\le\mu_-/48<\mu_*/24$.)
By Step~1, taken at $\mathrm{sv}=0$, the operator $\mathrm{Fst}(c_{\theta})$ is invertible.
Proposition~33.4\,(b) now shows that $D_*(c_{\theta})$ is invertible and that
\[
\left|D_*(c_{\theta})^{-1}\right|_{\ell^2}\le\mathrm{gGa}_{1}.
\]

\emph{Step~3: the bound.}
Let $\theta\in\mathrm{Stp}^{+}$, and put
\[
\mathrm{Kf}:=D_*\left(c_{\theta}+t\cdot\hat{c}_{\mathfrak{c}}\right)-D_*(c_{\theta}),\qquad\mathrm{Yf}:=D_*(c_{\theta})^{-1}\cdot\mathrm{Kf}.
\]
Proposition~33.3 applies with $A'=c_{\theta}$ and $B=t\cdot\hat{c}_{\mathfrak{c}}$:
by Step~1 the segment from $A'$ to $A'+B$ lies in $\mathrm{Gwk}^{+}$;
$\epsilon_B<e_{*}^{2\varepsilon}$ by (34.9);
and $\#\left|\mathrm{Ypr}_B\right|\le\mathrm{pp}$, because $\mathrm{Ypr}_B$ is empty for $t=0$ and equals $\mathrm{Ypr}_{\mathfrak{c}}$, a subset of $\mathrm{Ypr}_{\mathfrak{c}}^{+}$, for $t\ne0$ (Proposition~35.3\,(e)).
It yields
\[
\left\|\mathrm{Kf}\right\|_0\le\left\|\mathrm{Kf}\right\|_{\gamma^{F}/2}\le C_K'\cdot e_{*}^{2\varepsilon},\qquad\sum_{y,y'}\left|\mathrm{Kf}\right|(y,y')\le C_K'\cdot e_{*}^{2\varepsilon}\cdot\mathrm{pp}.
\]
From (M2) of Section~G5.1.7, in the form (P4)\,(e), from (P4)\,(a), from Step~2 and from the second condition of (g4) one obtains
\begin{gather*}
\left|\mathrm{Kf}\right|_{\ell^2}\le2C_K'\cdot e_{*}^{2\varepsilon},\qquad\left|\mathrm{Kf}\right|_1\le4C_K'\cdot e_{*}^{2\varepsilon}\cdot\mathrm{pp},\\
\left|\mathrm{Yf}\right|_{\ell^2}\le2\,\mathrm{gGa}_{1}\cdot C_K'\cdot e_{*}^{2\varepsilon}\le\frac{1}{2},\qquad\left|\mathrm{Yf}\right|_1\le4\,\mathrm{gGa}_{1}\cdot C_K'\cdot\mathrm{pp}\cdot e_{*}^{2\varepsilon}.
\end{gather*}
Consequently $\mathrm{gD}^{\theta}(t)=\operatorname{Tr}\log\left(1+\mathrm{Yf}\right)$ is defined by the power series of the logarithm, and (P4)\,(c) gives
\[
\left|\mathrm{gD}^{\theta}(t)\right|\le2\left|\mathrm{Yf}\right|_1\le8\,\mathrm{gGa}_{1}\cdot C_K'\cdot\mathrm{pp}\cdot e_{*}^{2\varepsilon}=\frac{B_1}{2}\le B_1\qquad\left(\theta\in\mathrm{Stp}^{+}\right).
\]

\emph{Step~4: analyticity on $\mathrm{Osp}$.}
Let $\theta\in\mathrm{Osp}$.
By Step~1 the connections $c_{\theta}$ and $c_{\theta}+t\cdot\hat{c}_{\mathfrak{c}}$ belong to $\mathrm{Inv}$, and by Step~2 the operator $D_*(c_{\theta})$ is invertible.
The map $\theta\mapsto c_{\theta}$ is linear, the set $\mathrm{Inv}$ is open, the entries of $D_*(\,\cdot\,)$ are analytic on $\mathrm{Inv}$ ((P2)\,(b)), and, by Cramer's rule, the entries of the inverse of an invertible matrix are analytic functions of its entries.
Hence the entries of $\mathrm{Yf}$ are analytic functions of $\theta$ on the open set $\mathrm{Osp}$, on which $\left|\mathrm{Yf}\right|_{\ell^2}\le\frac{1}{2}$ by Step~3.
By (P4)\,(d) the function $\theta\mapsto\mathrm{gD}^{\theta}(t)$ is analytic on $\mathrm{Osp}$.

\emph{Step~5: periodicity.}
Let $\mu\in\{1,2,3\}$, and let $\mathrm{Acl}$ be the real constant bond function on $T_*$ whose component in the direction $\mu$ is $-2\pi/(e_{*}\cdot\ell_*)$ and whose two other components vanish.
Then
\[
c_{\theta+2\pi e_\mu}=c_{\theta}-\mathrm{Acl}\quad\left(\theta\in\mathbb{C}^3\right),\qquad e_{*}\cdot\ell_*\cdot\mathrm{Acl}_\nu\in2\pi\mathbb{Z}\quad(\nu=1,2,3).
\]
By Proposition~35.1\,(d), $(k_{*},\ell_*,\mu_*,e_{*})$ are scale-free data (SF) in the sense of Section~I.1.8, with $k_{*}\ge1$.
Proposition~10\,(a), in the second of its examples, taken with $\mathrm{uf}(c)=\exp\left(i\cdot e_{*}\cdot\mathrm{Acl}\cdot c\right)$ and with the restriction $\mathrm{uf}^{0}$ of $\mathrm{uf}$ to $T^0_*$, gives
\[
D_*\left(A'-\mathrm{Acl}\right)=\mathrm{uf}^{0}\cdot D_*(A')\cdot\left(\mathrm{uf}^{0}\right)^{-1}\qquad\left(A'\in\tilde{G}_*\right),
\]
where $\mathrm{uf}^{0}$ does not depend on $A'$.
For real $\theta$ the connections $c_{\theta}$ and $c_{\theta}+t\cdot\hat{c}_{\mathfrak{c}}$ lie in $\mathrm{Gwk}$, which is a subset of $\tilde{G}_*$ ((P1)\,(c), (d), (34.9)).
Therefore the operators $D_*(c_{\theta})$, $\mathrm{Kf}$ and $\mathrm{Yf}$ formed at the point $\theta+2\pi e_\mu$ are the conjugates by $\mathrm{uf}^{0}$ of the same operators formed at $\theta$.
The traces $\operatorname{Tr}\mathrm{Yf}^n$ do not change under this conjugation, and hence
\[
\mathrm{gD}^{\theta+2\pi e_\mu}(t)=\mathrm{gD}^{\theta}(t)\qquad\left(\theta\in\mathbb{R}^3\right).
\]
The function $\theta\mapsto\mathrm{gD}^{\theta+2\pi e_\mu}(t)-\mathrm{gD}^{\theta}(t)$ is analytic on $\mathrm{Osp}$, by Step~4 and the invariance of $\mathrm{Osp}$ under the translation, and it vanishes on $\mathbb{R}^3$.
At a real point, a complex partial derivative of an analytic function is a limit of difference quotients taken in a real direction;
by induction on the order, every partial derivative of this function therefore vanishes on $\mathbb{R}^3$.
So the Taylor series of the function at a real point is zero, and by the identity theorem the function vanishes on the connected open set $\mathrm{Osp}$.

\emph{Outcome of Steps~1--5.}
For $\theta\in\mathrm{Stp}$ the connections $c_{\theta}+\mathrm{sv}\cdot t\cdot\hat{c}_{\mathfrak{c}}$, $0\le\mathrm{sv}\le1$, lie in $\mathrm{Gwk}^{+}$, and $D_*(c_{\theta})$ is invertible with $\left|D_*(c_{\theta})^{-1}\right|_{\ell^2}\le\mathrm{gGa}_{1}$.
On the open set $\mathrm{Osp}$, which contains the closed strip $\mathrm{Stp}$, the function $\theta\mapsto\mathrm{gD}^{\theta}(t)$ is analytic and $2\pi$-periodic in each $\theta_\mu$, and $\left|\mathrm{gD}^{\theta}(t)\right|\le B_1$.

\emph{Step~6: the dependence on a real twist.}
By the third condition of (g4), $\tau_s\ge1$.
By Steps~3--5 the function $\theta\mapsto\mathrm{gD}^{\theta}(t)$ is continuous on $\mathrm{Stp}$, holomorphic in the interior of $\mathrm{Stp}$, $2\pi$-periodic in each variable, and bounded in modulus by $B_1$ on $\mathrm{Stp}$.
These are the hypotheses of \cite[Lemma~20.13]{QEDI}, with $\tau_s$ as the width of the closed tube, $B_1$ as the bound, and the three variables numbered $1,2,3$ instead of $0,1,2$;
that lemma bounds the difference of the values of such a function at two real points by $52$ times the bound times the exponential of minus the width.
Taking the second real point to be $0$ and using the fourth condition of (g4), one obtains for every real $\theta$, with $\varepsilon_1=\frac{1}{2}e_{*}^{\frac{1}{4}-9\varepsilon}$ (Section~I.8.4),
\begin{equation}
\left|\mathrm{gD}^{\theta}(t)-\mathrm{gD}^{0}(t)\right|\le52B_1\cdot\exp\left(-\tau_s\right)\le e_{*}^{\frac{1}{4}-9\varepsilon}=2\varepsilon_1.
\tag{34.16}
\end{equation}
The numbers $\mathrm{gGa}_{1}$, $B_1$, $\tau_s$ and $e_{*}$ that enter the bounds of (a) are constants or numbers of (P7).
Hence these bounds depend neither on $t$, $\theta$, $\mathfrak{c}$ or $N$, nor on the placement, nor on a replacement of $\beta$ by $U_r\beta$.
\end{proof}

\begin{proof}[Proof of Proposition~34.4\,(b)]
Let $\theta\in\mathbb{R}^3$, let $A$ satisfy $\chi_{*}(A)=1$, let $\left(A^{g},\mathscr{A}^{0}\right)$ be as in Proposition~35.4\,(SF5), and let $|t|<t_x^{+}$.
By the definition of $\mathrm{ff}$ in Proposition~34.1\,(b), at the points $t\cdot\mathrm{ue}_{\mathfrak{c}}$ and $0$ of $\mathrm{DDo}\left(t_x^{+}\right)$ one has
\[
\mathrm{Ztm}\left(t\cdot\mathrm{ue}_{\mathfrak{c}}\right)=\mathrm{ff}\left(t\cdot\mathrm{ue}_{\mathfrak{c}}\right)\cdot\det D_*\left(c_{\theta}+t\cdot\hat{c}_{\mathfrak{c}}\right),\qquad\mathrm{Ztm}(0)=\mathrm{ff}(0)\cdot\det D_*(c_{\theta}).
\]
Together with the display (6.1) of Proposition~34.1\,(c) and with Proposition~34.3\,(c), this gives
\[
\mathrm{Rt}\left(t\cdot\mathrm{ue}_{\mathfrak{c}};\theta,A\right)=\exp\left\{\sum_{X\in D_*}\mathfrak{P}_*\left(X,t\cdot\hat{c}_{\mathfrak{c}}\right)\right\}\cdot\exp\left(\mathrm{gD}^{\theta}(t)\right)\cdot\frac{\mathrm{ff}\left(t\cdot\mathrm{ue}_{\mathfrak{c}}\right)}{\mathrm{ff}(0)}.
\]
By Proposition~34.1\,(b), $\mathrm{ff}$ is analytic on $\mathrm{DDo}\left(t_x^{+}\right)$, and $\left|\mathrm{ff}-1\right|\le\varepsilon_1\le\frac{1}{16}$ there.
So the principal branch $\log\mathrm{ff}$ is analytic on $\mathrm{DDo}\left(t_x^{+}\right)$, and (Ei2) gives
\[
\left|\log\mathrm{ff}\right|\le\frac{\varepsilon_1}{1-\varepsilon_1}\le\frac{16}{15}\cdot\varepsilon_1\le2\varepsilon_1.
\]
Consider on the disk $|t|<t_x^{+}$ the two functions
\[
t\mapsto\mathrm{Lg}\left(t\cdot\mathrm{ue}_{\mathfrak{c}};\theta,A\right)
\qquad\text{and}\qquad
t\mapsto\sum_{X\in D_*}\mathfrak{P}_*\left(X,t\cdot\hat{c}_{\mathfrak{c}}\right)+\mathrm{gD}^{\theta}(t)+\log\mathrm{ff}\left(t\cdot\mathrm{ue}_{\mathfrak{c}}\right)-\log\mathrm{ff}(0).
\]
Both are analytic on this disk.
For the first, this is Proposition~34.1\,(d).
In the second, each $\mathfrak{P}_*\left(X,t\cdot\hat{c}_{\mathfrak{c}}\right)$ is analytic by Proposition~34.2 and (34.9);
and $\mathrm{gD}^{\theta}$, at the fixed real $\theta$, is analytic by (P4)\,(d), because the entries of
\[
\mathrm{Yf}=D_*(c_{\theta})^{-1}\cdot\left(D_*\left(c_{\theta}+t\cdot\hat{c}_{\mathfrak{c}}\right)-D_*(c_{\theta})\right)
\]
are analytic in $t$ ((P2)\,(b); the connection $c_{\theta}+t\cdot\hat{c}_{\mathfrak{c}}$ lies in $\tilde{G}_*$) and $\left|\mathrm{Yf}\right|_{\ell^2}\le\frac{1}{2}$ (Proposition~34.3\,(c)).
Both functions vanish at $t=0$:
$\mathrm{Lg}(0;\theta,A)=0$, $\mathfrak{P}_*(X,0)=0$, and $\mathrm{Yf}=0$ at $t=0$, so that $\mathrm{gD}^{\theta}(0)=0$.
Both have the exponential $\mathrm{Rt}\left(t\cdot\mathrm{ue}_{\mathfrak{c}};\theta,A\right)$:
the first by Proposition~34.1\,(d), the second by the factorization of $\mathrm{Rt}$ obtained above.
By (P5) the two functions coincide.
Subtracting $\mathrm{lgN}(t)=\sum_{X\in D_*}\mathfrak{P}_*\left(X,t\cdot\hat{c}_{\mathfrak{c}}\right)+\mathrm{gD}^{0}(t)$, one finds
\[
\mathrm{wdf}_{\mathfrak{c}}(t)=\mathrm{gD}^{\theta}(t)-\mathrm{gD}^{0}(t)+\log\mathrm{ff}\left(t\cdot\mathrm{ue}_{\mathfrak{c}}\right)-\log\mathrm{ff}(0).
\]
For the same reasons, applied at $\theta$ and at $0$, the right side is analytic on $|t|<t_x^{+}$, and it vanishes at $t=0$.
By (34.16) and the bound for $\log\mathrm{ff}$,
\[
\left|\mathrm{wdf}_{\mathfrak{c}}(t)\right|\le2\varepsilon_1+2\varepsilon_1+2\varepsilon_1=3e_{*}^{\frac{1}{4}-9\varepsilon}.
\]
This bound is a function of $e_{*}$, and by (P7) it is uniform.
\end{proof}

\subsection*{5. Conditions on the constants}

The proofs of Sections~0--4 rest on the following conditions on the constants, which stand in the lines of Section~I.8.2 and in Proposition~14.2 as indicated.

\begin{description}
\item[$\kappa\ge7\kappa_{0}+7$]
(line (C2); Proposition~14.2\,(a), entry~$\kappa$, and~(f)).
In the form $\kappa\ge\kappa_{0}$ it enters Step~3\,(iii) of the proof of Proposition~34.1\,(b);
in the form $\kappa-1\ge\kappa_{0}$ it enters (4e) and Step~5 of the proof of Proposition~34.3\,(b).
\item[$L\cdot\left(\kappa-\kappa_{0}-1\right)\ge\kappa$]
(line (C3); condition (L3) of Proposition~14.2\,(b); Proposition~14.2\,(f)).
Step~2 of the proof of Proposition~34.2 rests on it.
\item[$c_{\mathrm{rb}}\cdot L^3\cdot\exp\left(-\left(L\cdot\left(\kappa-\kappa_{0}-1\right)-\kappa\right)\right)\le\frac{1}{2}$]
(line (C3); condition (L4) of Proposition~14.2\,(b)).
It is used in Step~2 of the proof of Proposition~34.2 as well.
\item[$L\ge5$]
(line (C3); condition (L5) of Proposition~14.2\,(b)).
Step~6 of the proof of Proposition~34.3\,(b) uses it, through (Ei5), in the form $L\ge3$.
In the form $L>1$ it gives $e_j<e_{j+1}$ in (34.1), and it is used in Steps~0 and~2 of the proof of Proposition~34.2 and in the proof of Proposition~34.3\,(b), namely in Steps~3 and~4 ($L^i\ge1$, $\eta_j\le1$) and in Step~6 ($a_{*}\le h$).
\item[$L^{1-\varepsilon}\ge2$]
This is a consequence of $L\ge5$ and $\varepsilon\le1/200$ (lines (C3) and (C7); conditions (L5) and (S1) of Proposition~14.2\,(b)).
It is used in Step~6 of the proof of Proposition~34.3\,(b).
\item[$0<1-\alpha<\varepsilon\le\bar{\varepsilon}=1/200$]
(line (C7), with line (C1) and the number $\bar{\varepsilon}$ of line (C0); Proposition~14.2\,(a), entry~$\varepsilon$; condition (S1) of Proposition~14.2\,(b); Proposition~14.2\,(f)).
It enters in the forms $\varepsilon>0$ and $3\varepsilon<1$, at the following places:
(P1)\,(a), namely (34.1) ($e_0<1$, $e_{*}^{\varepsilon}\le e_0^{\varepsilon}$), (a1) and (a2);
Step~0 of the proof of Proposition~34.2, where the exponents $-1+2\varepsilon$ and $-1+3\varepsilon$ are negative;
Step~1 of the proof of Proposition~34.1\,(b), and (34.9), through $e_{*}^{3\varepsilon}<e_{*}^{2\varepsilon}$;
Step~4 of the proof of Proposition~34.1\,(b), through $e_{*}^{\frac{1}{2}+\varepsilon}\le e_{*}^{\frac{1}{4}-7\varepsilon}$ and $e_{*}^{\frac{1}{4}-6\varepsilon}\le e_{*}^{\frac{1}{4}-7\varepsilon}$;
(4c) of the proof of Proposition~34.3\,(b), through $e_j^{-1+2\varepsilon}<e_j^{-1+\varepsilon}$;
the proof of Proposition~34.4\,(a), through $\mathrm{wd}_s>0$.
\item[$0<\gamma^{F}\le1$]
(line (C4r), which gives $\gamma^{F}=\gamma_{z}/2\le1/(6L)$; Proposition~14.2\,(a), entry~$\gamma^{F}$; Proposition~14.2\,(f)).
It enters in the form $\gamma^{F}\ge0$.
Then $\exp\left(\frac{\gamma^{F}}{2}\cdot d(y,y')\right)\ge1$ for all $y,y'$, so that the norms of Section~G5.1.7 satisfy
\[
\left\|O\right\|_0\le\left\|O\right\|_{\gamma^{F}/2}
\]
for every operator $O$ on which they are defined.
This inequality is used in the proof of Proposition~34.3\,(c) and in Step~3 of the proof of Proposition~34.4\,(a).
\item[$M\ge L^2$]
(line (C5); condition (M1) of Proposition~14.2\,(b)).
It is used in the form $M\ge2$ in Step~4 of the proof of Proposition~34.3\,(b) ($\mathrm{rb}\le s_X-1$), and in the form $M\ge1$ in Step~3 of that proof and in Step~3\,(iii) of the proof of Proposition~34.1\,(b).
\item[$C_{0}\ge1$]
(line (C6); Proposition~14.2\,(f)).
Step~3\,(i) of the proof of Proposition~34.1\,(b) uses it.
\item[$C_K\le C_K'$; $\mathrm{gGa}^2\le\mathrm{gGa}_{1}$; $\mathrm{gGa}_{1}\ge1$; $c_{\tau}\le\mu_-/48$]
(line (C6$'$), that is, the definitions in Proposition~14.2\,(a); all four inequalities are stated in Annex~B, Section~G5.9.3, and in Proposition~14.2\,(f)).
The first two are used in the proof of Proposition~34.3\,(c), the third in Step~1 of the proof of Proposition~34.1\,(b), and the fourth in Step~2 of the proof of Proposition~34.4\,(a).
The third follows also from the definitions alone:
\[
\mathrm{gGa}_{1}=b_{-}^{-1}+\frac{4}{\mu_-}\ge b_{-}^{-1}\ge1,
\]
because $\mu_-=L^{-3}>0$ and $0<b_{-}=1-1/L\le1$ (Section~G5.1.8; lines (C0), (C4), (C6$'$)).
\item[$e_0^{\varepsilon}\le\frac{1}{4}$]
(line (C8); the first inequality of condition (A8) of Proposition~14.2\,(b)).
It implies $e_0<1$, $e_j<1$ and $4e_{*}^{\varepsilon}\le1$.
It is used in (34.1) and in (P1)\,(a1), and through them in Step~0 of the proof of Proposition~34.2, in Steps~1, 3\,(iii) and~4 of the proof of Proposition~34.1\,(b), in (34.9) and in (4c) of the proof of Proposition~34.3, and in the proof of Proposition~34.4\,(a) ($\mathrm{wd}_s>0$).
\item[$r_{*}>31L+16$]
This is the condition (2.1) of Section~G5.1.9; it enters in the form $r_{*}\ge12$, in Step~6 of the proof of Proposition~34.3\,(b).
It is a consequence of $e\cdot h^{\frac{1}{2}}\le e_1$ (Proposition~14.2\,(c), from (e1) of line (C9); Proposition~35.1\,(e)).
\item[$\mathrm{Cb}\ge1$]
(line (C10), namely the definition $\mathrm{Cb}=1+h^{\frac{1}{2}}\cdot C_{\beta}$).
It is used in the proof of Proposition~34.3\,(d).
\end{description}

The three conditions of (g3) and the five conditions of (g4) of Section~G5.1.9 stand in line (C10), through the definition of $e_{\sharp}(\beta)$, and in Proposition~14.2\,(d).
They enter the proofs as follows.
\begin{description}
\item[(g3), first:] $4C_{0}\cdot\ell_*^3\cdot\mathrm{gGa}_{1}\le e_{*}^{-\frac{1}{4}+\varepsilon}$, in Step~1 of the proof of Proposition~34.1\,(b).
\item[(g3), second:] $4\ell_*^3\cdot\left(2C_K+K_{0}+4\right)\le e_{*}^{-2\varepsilon}$, in Step~4 of the proof of Proposition~34.1\,(b).
\item[(g3), third:] $e_{*}^{\frac{1}{4}-9\varepsilon}\le\frac{1}{8}$, in Step~4 of the proof of Proposition~34.1\,(b) and in the proof of Proposition~34.4\,(b).
\item[(g4), first:] $48e_{*}^{2\varepsilon}\le\mu_-$, in Step~1 of the proof of Proposition~34.1\,(b), and in Step~2 of the proof of Proposition~34.4\,(a) for $\theta\in\mathrm{Stp}^{+}$ outside $\mathrm{Stp}$.
\item[(g4), second:] $2\,\mathrm{gGa}_{1}\cdot C_K'\cdot e_{*}^{2\varepsilon}\le\frac{1}{2}$, in the proof of Proposition~34.3\,(c) and in Step~3 of the proof of Proposition~34.4\,(a).
\item[(g4), third:] $\tau_s\ge1$, in Step~6 of the proof of Proposition~34.4\,(a).
\item[(g4), fourth:] $52B_1\cdot\exp\left(-\tau_s\right)\le e_{*}^{\frac{1}{4}-9\varepsilon}$, in Step~6 of the proof of Proposition~34.4\,(a).
\item[(g4), fifth:] $B_0\le\frac{1}{4}$, in the proofs of Proposition~34.3\,(d) and~(e).
\end{description}

The proofs also use the following facts on the scales, which are not conditions on constants (Proposition~35.1\,(a), (c), (d)):
$k_{*}\ge1$;
$\ell_*=M\cdot r_{*}$, hence $\ell_*\ge1$;
$a_{*}\in\left\{h/L,\,h/L^2\right\}$;
$\mu_-<\mu_*$;
$e_{*}\le e_0$;
$\ell_j\in\mathrm{Cl}$;
the independence of the scale numbers of $N$.

\section*{Chapter 35b. Several cells, the complement of the small-field set, the four-jet representation: proofs of Propositions 35.5--35.8}

\subsection*{0. Conventions of this chapter: setting, letters, facts used in the four proofs, results of \cite{QEDI}}

\subsubsection*{0.1. Setting}

The four statements proved in this chapter, Propositions~35.5--35.8, have the data (DAT) of Section~G5.4.0.
The letters $\mathscr{A}$, $\mathscr{T}(t)$ and $\mathrm{Gcv}$, the function $\mathrm{Ztm}(t)$ together with its defining display (G5.5.1), and the meaning of the word ``uniform'' are those fixed in that section.
By the definition of $e_{\sharp}(\beta)$ in Section~G5.1.9 one has $e\cdot h^{\frac{1}{2}}\le e_1$, and the conditions (g1)--(g9) are satisfied.

\subsubsection*{0.2. Letters}

Put
\begin{gather*}
\mathrm{sF}:=\mathrm{ssum}\left(\frac{\gamma^{F}}{2}\right)=\left(\frac{24}{\gamma^{F}}\right)^{3},\qquad
\mathrm{Ksf}:=\left\{A\in{V_{\mathrm{t}}}(T^0_{*}):\ \chi_{*}(A)=1\right\},\\
X^{+}:=\left\{y\in T^0_{*}:\ d\left(y,X\cap T^0_{*}\right)\le1\right\}\qquad\text{for a union $X$ of $M$-cubes of $T_{*}$};
\end{gather*}
when $X$ is a single $M$-cube $\square$, this set is denoted by $\square^{+}$.
Thus
\[
\mathbf{S}=C^{F}\cdot\mathrm{sF},\qquad C_K=226\left(C^{F}\right)^{2}\cdot\mathrm{sF}^{2},\qquad\mathrm{CVj}=10^{3}\,C_H^{2}\cdot\mathrm{sF}^{4}\cdot K_{0}.
\]
A sum over $u=(y,\beta,\omega)$ extends over $y\in T^0_{*}$, $\beta\in\{1,\dots,4\}$ and $\omega\in\{0,1\}$; these indices carry the weight $\nu=1$ (Section~I.3.6).
Sums over the sites and over the pairs of the tori extend over a fundamental domain of $\ell_*\mathbb{Z}^3$ ((I.AP)).
The moduli of the kernels and of the numbers $\mathrm{Leg}_A(\mathrm{up};u)$, as well as all norms, do not depend on the choice of this domain.

In this chapter $\mathrm{zs}$ denotes a complex variable of integration, $\mathrm{sv}$ a real parameter, $\mathrm{pl}$ a plaquette of $T^0_{*}$, and $w$ a unit site.
Kernels are denoted by $\mathrm{Kv}$, and a slot $\mathrm{sl}$ of a kernel $\mathrm{Kv}_p(u_1,\dots,u_p)$ is one of the positions $1,\dots,p$ of its arguments, as in the bounds (3.3) of (Q2) in 0.6.
In 1.8 and 1.9, $\mathrm{Leg}$ is short for $\mathrm{Leg}_{\mathscr{T}(t)}$ of Section~G5.1.7.
The letter $p$ denotes a degree (in Section~1, $p=n+m$), except in 3.2, 4.6 and Section~5, where it is the integer of the line (C7).
In 3.6 and 3.7, $F'$ denotes a point of $F_N$, in lattice units, and in (n1), 3.2, 4.4 and 4.5 the letter $y$ denotes a real number, as in Proposition~31.1\,(d).

The conditions (A8) and (T1) are those of Section~I.8.4, and (C0)--(C10) are the lines of constants of Section~I.8.2.
By Proposition~35.1\,(e) the first two clauses of (T1) hold.
The fact (M2) is the one so labelled in Section~G5.1.7, and (M1), (M3), (L5), (e1) are the conditions so labelled in Proposition~14.2\,(b).
The displays (2.1)--(2.5) are those of Proposition~35.4; (G5.5.1) is the display of Section~G5.4.0 that defines $\mathrm{Ztm}(t)$, and (5.2), (5.3) are those of Proposition~35.5; (6.1) and (6.2) are those of Proposition~34.1, (6.4) is that of Proposition~34.3, and (8.1) that of Proposition~35.8.
All other numbered displays belong to this chapter.

\subsubsection*{0.3. Numbers}

\begin{itemize}
\item[(n1)] By Proposition~35.1\,(a), $k_*\ge1$, $\eta_*\le1/L<1$ and $0<e_*\le e_0$; moreover $e_0<1$ by (A8).
Hence the function $y\mapsto e_*^{y}$ is non-increasing on $\mathbb{R}$.
\item[(n2)] The third condition of (g3) gives $w^{\mathrm{E}}_{*}=e_*^{\frac{1}{4}-7\varepsilon}\le e_*^{\frac{1}{4}-9\varepsilon}\le\frac{1}{8}$.
Consequently $\bar{\kappa}_1=6\,e_*^{\frac{1}{4}-9\varepsilon}\le\frac{3}{4}$.
\item[(n3)] By Section~G5.1.8 and (G5.T-1), $C^{F}\ge1$, $0<\gamma^{F}\le1$, $\gamma^{F}\cdot M\ge16$, $M\ge L^2\ge2$ and $K_0\ge1$.
It follows that $\mathrm{sF}\ge1$, $\mathbf{S}\ge1$, $\mathrm{hL}=6\,\mathbf{S}\ge1$, $C_H=203\,M\cdot\left(C^{F}\right)^{2}\ge2$ and $h_G=1+\mathrm{gGa}\ge1$.
\item[(n4)] By Section~G5.1.8, $\mathrm{vt}\le\gamma^{F}/8$ and $\mathrm{vt}\le1/M$.
Therefore $\mathrm{vt}\le\gamma^{F}/2$, $\gamma^{F}-\mathrm{vt}\ge7\gamma^{F}/8$ and $\mathrm{vt}\cdot(M+1)\le2$.
\item[(n5)] Let $0<\gamma\le6$, let $\ell\ge1$ be an integer, and let $y\in(\mathbb{Z}/\ell\mathbb{Z})^3$.
Then (Section~G5.1.8)
\[
\sum_{y'}\exp\left(-\gamma\cdot d(y,y')\right)\le\mathrm{ssum}(\gamma).
\]
\begin{proof}
Every class $y'$ has a lift whose sup distance from a fixed lift of $y$ equals $d(y,y')$, and distinct classes have distinct lifts.
As $|z|\ge|z|_1/3$ for $z\in\mathbb{Z}^3$, the sum is therefore bounded by
\[
\left(\sum_{j\in\mathbb{Z}}\exp\left(-\frac{\gamma\cdot|j|}{3}\right)\right)^{3}=\coth\left(\frac{\gamma}{6}\right)^{3},
\]
and $\coth(\gamma/6)\le1+6/\gamma\le12/\gamma$.
\end{proof}
The left side does not increase when $\gamma$ increases.
Hence it is at most $\mathrm{sF}$ for every $\gamma\ge\gamma^{F}/2$.
\item[(n6)] For a non-empty set $P$ of unit sites and a unit site $w$,
\[
\sum_{z\in P}\mathrm{Edy}(w,z)\le\mathrm{sF}\cdot\exp\left(-\frac{\gamma^{F}}{2}\cdot d(w,P)\right).
\]
\begin{proof}
For $z\in P$ one has
\[
\mathrm{Edy}(w,z)=\exp\left(-\frac{\gamma^{F}}{2}\cdot d(w,z)\right)\cdot\exp\left(-\frac{\gamma^{F}}{2}\cdot d(w,z)\right)\le\exp\left(-\frac{\gamma^{F}}{2}\cdot d(w,P)\right)\cdot\exp\left(-\frac{\gamma^{F}}{2}\cdot d(w,z)\right),
\]
and by (n5) the sum over $z$ of the last factor is at most $\mathrm{sF}$.
\end{proof}
\item[(n7)] $\epsilon_{\mathfrak{c}}\le q_x$ (Proposition~35.3\,(f)).
\item[(n8)] $\#|\mathrm{Ypr}_{\mathfrak{c}}|\le\#|\mathrm{Ypr}_{\mathfrak{c}}^{+}|\le\mathrm{pp}$ (Proposition~35.3\,(e); $\mathrm{Ypr}_{\mathfrak{c}}$ is a subset of $\mathrm{Ypr}_{\mathfrak{c}}^{+}$).
\item[(n9)] For $\mathfrak{c}\ne\mathfrak{c}'$,
\[
d\left(\mathrm{Ypr}_{\mathfrak{c}},\mathrm{Ypr}_{\mathfrak{c}'}\right)\ge d_{\mathrm{sep}}\ge\frac{M\cdot r_{*}}{8L}-1\ge6M+8,
\]
by the definition of $d_{\mathrm{sep}}$ in Section~G5.1.4, by Proposition~35.3\,(e) and by (g2).
In particular $d\left(\mathrm{Ypr}_{\mathfrak{c}},\mathrm{Ypr}_{\mathfrak{c}'}\right)\ge M+2$.
\end{itemize}

\subsubsection*{0.4. Fields}

\begin{itemize}
\item[(f1)] For $t\in\mathrm{DDo}(t_x^{+})$ the bond function $\mathscr{C}(t)=a(t)^{[k_*]}$ belongs to $\mathfrak{X}_{*}$, and $e_*\cdot\left|\mathscr{C}(t)\right|_{\infty}<e_*^{3\varepsilon}$ (Proposition~35.4\,(DF5) at $j=k_*$).
For $0\le\mathrm{sv}\le1$ one has $\mathrm{sv}\cdot\mathscr{C}(t)=\mathscr{C}(\mathrm{sv}\cdot t)$, and $\mathrm{sv}\cdot t\in\mathrm{DDo}(t_x^{+})$.
\item[(f2)] The bond function $\mathscr{A}$ is real and belongs to $\frac{1}{4}\tilde{R}_{*}$ (Proposition~35.4\,(SF5), (2.3) at $\mathrm{sv}=1$).
The set $\frac{1}{4}\tilde{R}_{*}$ is contained in $\tilde{R}_{*}$, because the thresholds of $\mathrm{cf}\cdot\tilde{R}_{*}$ are $\mathrm{cf}$ times those of $\tilde{R}_{*}$ (Section~I.5.3; \cite[Definition~5.16]{QEDI}).
Hence, for $t\in\mathrm{DDo}(t_x^{+})$ and $0\le\mathrm{sv}\le1$, the pair $\left(\mathscr{A},\mathscr{C}(t)\right)$ lies in $\tilde{R}_{*}\times\mathfrak{X}_{*}$, and $\mathscr{A}+\mathrm{sv}\cdot\mathscr{C}(t)$ lies in the set $\mathrm{Gwk}$ defined in Section~G5.1.7.
By Proposition~35.4\,(SF2), $\mathrm{Gwk}$ is a subset of the open set $\tilde{G}_{*}$.
\item[(f3)] By Proposition~35.4\,(SF2), $S_{*}(\cdot)$ is analytic on $\tilde{G}_{*}$, and $\tilde{G}_{*}$ is invariant under $A\mapsto-A$; by the first sentence of Proposition~33.1, $Q_{*}(A)J^{*}$, $J\,Q_{*}^{T}(-A)$ and $\delta_{\alpha,A}$ are entire functions of $A$.
In view of the formulas of Section~G5.1.7,
\[
H_{*}(A)=b_{*}\cdot S_{*}(A)\,J\,Q_{*}^{T}(-A),\qquad\bar{H}_{*}(A)=\mathsf{C}^{-1}H_{*}(-A)\,\mathsf{C},\qquad D_{*}(A)=b_{*}-b_{*}\cdot Q_{*}(A)\,J^{*}H_{*}(A),
\]
the entries $D_{*}(A)(y,y')_{\beta\beta'}$ and all numbers $\mathrm{Leg}_A(\mathrm{up};u)$, which are the coefficients of $H_{*}(A)\Psi$, $\bar{H}_{*}(A)\bar{\Psi}$, $\delta_{\alpha,A}H_{*}(A)\Psi$ and $\delta_{\alpha,-A}\bar{H}_{*}(A)\bar{\Psi}$, are analytic functions of $A$ on $\tilde{G}_{*}$.
The map $t\mapsto\mathscr{T}(t)=\mathscr{A}+\sum_{\mathfrak{c}}t_{\mathfrak{c}}\cdot\hat{c}_{\mathfrak{c}}$ is affine, and by (f2) it maps $\mathrm{DDo}(t_x^{+})$ into $\tilde{G}_{*}$.
Therefore, for each function $O$ among those just named, $t\mapsto O\left(\mathscr{T}(t)\right)$ is analytic on $\mathrm{DDo}(t_x^{+})$, and the chain rule gives, with $\partial_B$ as in Section~G5.1.7,
\[
\partial_{\mathfrak{c}}\left[O\left(\mathscr{T}(t)\right)\right]=\left(\partial_{\hat{c}_{\mathfrak{c}}}O\right)\left(\mathscr{T}(t)\right),\qquad
\partial_{\mathfrak{c}}\partial_{\mathfrak{c}'}\left[O\left(\mathscr{T}(t)\right)\right]=\left(\partial_{\hat{c}_{\mathfrak{c}}}\partial_{\hat{c}_{\mathfrak{c}'}}O\right)\left(\mathscr{T}(t)\right).
\]
\end{itemize}

\subsubsection*{0.5. The marks are local}

Let $X\in D_{*}$ and $\mathfrak{c}\in\mathrm{Cel}$.

\begin{itemize}
\item[(m1)] If $c$ is a photon site in $X$, then the unit site $[c]$ lies in $X$.
\begin{proof}
By Proposition~35.4\,(SF4)\,(v) the block $\Delta_{[c]}$ is contained in one $M$-cube, and this cube contains $c$.
Distinct $M$-cubes are disjoint (Section~I.1.4), and $X$ is a union of $M$-cubes; so this cube is one of the cubes of $X$.
Since $[c]\in\Delta_{[c]}$, the claim follows.
\end{proof}
\item[(m2)] If a bond $b$ of $\mathrm{Bd}_{\mathfrak{c}}$ meets $X$, then $X$ contains a point of $\mathrm{Ypr}_{\mathfrak{c}}$.
\begin{proof}
The bond $b$ has an end point $c$ in $X$ (Section~I.1.1).
As $c\in\Delta_{[c]}$, the definition of $\mathrm{Ypr}_{\mathfrak{c}}$ in Section~G5.1.4 gives $[c]\in\mathrm{Ypr}_{\mathfrak{c}}$, and $[c]\in X$ by (m1).
\end{proof}
\item[(m3)] Let $\mathrm{gf}$ be a function on $\mathfrak{X}_{*}$ such that $\mathrm{gf}(\mathscr{C})=\mathrm{gf}(\tilde{\mathscr{C}})$ whenever $\mathscr{C},\tilde{\mathscr{C}}\in\mathfrak{X}_{*}$ agree on the bonds that meet $X$.
If $X$ contains no point of $\mathrm{Ypr}_{\mathfrak{c}}$, then the function $t\mapsto\mathrm{gf}\left(\mathscr{C}(t)\right)$ on $\mathrm{DDo}(t_x^{+})$ does not depend on $t_{\mathfrak{c}}$.
\begin{proof}
By (m2), $\hat{c}_{\mathfrak{c}}(b)=0$ for every bond $b$ that meets $X$.
Consequently, if $t,t'\in\mathrm{DDo}(t_x^{+})$ differ in the component $\mathfrak{c}$ only, then $\mathscr{C}(t)$ and $\mathscr{C}(t')$ agree on these bonds, and both belong to $\mathfrak{X}_{*}$ by (f1).
\end{proof}
\end{itemize}

\subsubsection*{0.6. Results of \cite[Section~20.3]{QEDI}}

\paragraph{(Q1)}
The following holds by \cite[Lemma~20.8]{QEDI}; here $\mathrm{gf}$ is a complex function on $\mathrm{DD}(\mathrm{rad})$, and the functions $\mathrm{gf}_{\mathrm{Scl}}$ are those of Section~G5.1.6.
\begin{itemize}
\item[(a)] $\mathrm{gf}=\sum_{\mathrm{Scl}\subset\mathrm{Cel}}\mathrm{gf}_{\mathrm{Scl}}$ and $\mathrm{gf}_{\emptyset}=\mathrm{gf}(0)$.
The function $\mathrm{gf}_{\mathrm{Scl}}$ depends only on the variables $t_{\mathfrak{c}}$, $\mathfrak{c}\in\mathrm{Scl}$, and it vanishes at every point at which $t_{\mathfrak{c}}=0$ for some $\mathfrak{c}\in\mathrm{Scl}$.
If $\mathrm{gf}$ does not depend on $t_{\mathfrak{c}}$, then $\mathrm{gf}_{\mathrm{Scl}}=0$ for every $\mathrm{Scl}$ that contains $\mathfrak{c}$.
\item[(b)] $\sup_{\mathrm{DD}(\mathrm{rad})}\left|\mathrm{gf}_{\mathrm{Scl}}\right|\le2^{\#|\mathrm{Scl}|}\cdot\sup_{\mathrm{DD}(\mathrm{rad})}\left|\mathrm{gf}\right|$.
\item[(c)] If $\mathrm{gf}$ is analytic in a neighbourhood of $\mathrm{DD}(\mathrm{rad})$ and $\mathfrak{c}\ne\mathfrak{c}'$ are elements of $\mathrm{Scl}$, then
\[
\sup_{\mathrm{DD}(\mathrm{rad})}\left|\mathrm{gf}_{\mathrm{Scl}}\right|\le2^{\#|\mathrm{Scl}|-2}\cdot\mathrm{rad}^{2}\cdot\sup_{\mathrm{DD}(\mathrm{rad})}\left|\partial_{\mathfrak{c}}\partial_{\mathfrak{c}'}\mathrm{gf}\right|.
\]
\item[(d)] If $\mathrm{gf}$ is analytic in a neighbourhood of $\mathrm{DD}(\mathrm{rad})$, $\mathrm{Scl}$ is not empty and $0<\mathrm{rad}'\le\mathrm{rad}/2$, then every monomial of the Taylor series of $\mathrm{gf}_{\mathrm{Scl}}$ at $0$ contains every variable $t_{\mathfrak{c}}$ with $\mathfrak{c}\in\mathrm{Scl}$ and no other variable, and the image of $\mathrm{gf}_{\mathrm{Scl}}$ in $\mathcal{J}$ satisfies
\[
\left\|\mathrm{gf}_{\mathrm{Scl}}\right\|_{\mathrm{rad}'}\le\sup_{\mathrm{DD}(\mathrm{rad})}\left|\mathrm{gf}_{\mathrm{Scl}}\right|.
\]
\end{itemize}
The map $\mathrm{gf}\mapsto\mathrm{gf}_{\mathrm{Scl}}$ is linear by its definition, and by the last sentence of (a) one has $\mathrm{gf}_{\mathrm{Scl}}=0$ whenever $\mathrm{gf}$ is constant and $\mathrm{Scl}\ne\emptyset$.

\paragraph{(Q2)}
By \cite[Section~20.3.4]{QEDI}, the functionals $\mathrm{rw}_G(\cdot\,;\mathrm{wg})$ and $\mathrm{Nw}_G(\cdot\,;\mathrm{wg})$ are subadditive.
Let
\[
\mathrm{Fv}=\sum_{p\ge2}\frac{1}{p!}\sum_{u_1,\dots,u_p}\mathrm{Kv}_p(u_1,\dots,u_p)\,\Psi(u_1)\cdots\Psi(u_p),
\]
where the kernels $\mathrm{Kv}_p$ need not be antisymmetric.
Then, by the same section, one has for every weight $\mathrm{wg}$ the two bounds (3.3) of $\mathrm{Nw}_G(\mathrm{Fv};\mathrm{wg})$ and $\mathrm{rw}_G(\mathrm{Fv};\mathrm{wg})$ by the kernels:
\begin{equation}\tag{3.3}
\begin{aligned}
\mathrm{Nw}_G(\mathrm{Fv};\mathrm{wg})&\le\sum_{p}\frac{h_G^{p}}{(p-1)!}\sum_{u_1,\dots,u_p}\left|\mathrm{Kv}_p(u_1,\dots,u_p)\right|\cdot\mathrm{wg}(u_1,\dots,u_p),\\
\mathrm{rw}_G(\mathrm{Fv};\mathrm{wg})&\le\sum_{p}\frac{h_G^{p}}{(p-1)!}\max_{\mathrm{sl}}\sup_{u}\sum_{\substack{u_1,\dots,u_p\\ u_{\mathrm{sl}}=u}}\left|\mathrm{Kv}_p(u_1,\dots,u_p)\right|\cdot\mathrm{wg}(u_1,\dots,u_p).
\end{aligned}
\end{equation}

\paragraph{(Q3)}
The assertions (a) and (c) below are \cite[Lemma~20.11\,(a), (c)]{QEDI}.
The data of that lemma are: a finite set $\mathrm{Ix}$ with a pseudo-metric; the Grassmann algebra with one pair of generators for each element of $\mathrm{Ix}$; an invertible complex matrix on $\mathrm{Ix}$; a rate $\ge0$; a number whose square bounds the $\ell^2$ operator norm of the matrix; a number which is at least one plus the preceding one; a bound for the row sums and for the column sums of the moduli of the entries of the matrix, each entry being multiplied by the exponential of the rate times the distance; and a bound, not larger than the preceding one, for the same sums without the exponential.

\emph{The setting of (Q2) and (Q3).}
One takes $\mathrm{Ix}=T^0_{*}\times\{1,\dots,4\}$, with $d\left((y,\beta),(y',\beta')\right):=d(y,y')$.
The pair of generators that belongs to $(y,\beta)$ consists of $\Psi(u)$, $u=(y,\beta,0)$, and $\Psi(u)$, $u=(y,\beta,1)$.
Thus the algebra is $\mathrm{Gr}[\Psi]$, and the functionals and the weights of \cite[Section~20.3.4]{QEDI} are $\mathrm{rw}_G$, $\mathrm{Nw}_G$, $\mathrm{wg}_{\delta}$ and $1$ of Section~G5.1.6.
The weight of \cite[Section~20.3.4]{QEDI} that belongs to the subset $P\times\{1,\dots,4\}$ of $\mathrm{Ix}$, $P$ a non-empty set of unit sites, is $\mathrm{wg}_P$, and the distance of $P\times\{1,\dots,4\}$ and $P'\times\{1,\dots,4\}$ is $d(P,P')$.
The matrix is $\mathrm{Gcv}$, the rate is $\mathrm{vt}$, and the four numbers are $\mathrm{gGa}$, $h_G$, $\mathrm{alG}$, $\mathrm{alG}$.
The conditions on the four numbers are
\[
\left|\mathrm{Gcv}\right|_{\ell^2}\le\mathrm{gGa}^{2},\qquad h_G\ge1+\mathrm{gGa},
\]
together with the inequality of Proposition~33.4\,(a),
\begin{align*}
\sup_{(y,\beta)}\sum_{(y',\beta')}\left|\mathrm{Gcv}(y,\beta;y',\beta')\right|\cdot\exp\left(\mathrm{vt}\cdot d(y,y')\right)&\le\mathrm{alG},\\
\sup_{(y',\beta')}\sum_{(y,\beta)}\left|\mathrm{Gcv}(y,\beta;y',\beta')\right|\cdot\exp\left(\mathrm{vt}\cdot d(y,y')\right)&\le\mathrm{alG}.
\end{align*}
They hold by Proposition~33.4\,(a), which applies because $\mathscr{A}$ is a real bond function on $T_{*}$ by (f2), and by $h_G=1+\mathrm{gGa}$.
For a matrix $\mathrm{Cov}$ the Gaussian functional of \cite[Section~20.3.2]{QEDI} is
\[
\int\mathrm{d}\mu_{\mathrm{Cov}}\;F=(\det\mathrm{Cov})\int\mathrm{D}\Psi\;\exp\left(-\left\langle\bar{\Psi},\mathrm{Cov}^{-1}\Psi\right\rangle\right)F,
\]
where $\int\mathrm{D}\Psi$ is normalised by $\int\mathrm{D}\Psi\;\exp\left(-\left\langle\bar{\Psi},O\Psi\right\rangle\right)=\det O$.
This is $\mathrm{d}\mu_{\mathrm{Cov}}$ of Proposition~15\,(D3) at $\nu=1$ (Proposition~15\,(e), last sentence; (I.B1)).
For $i\ge1$, $\tilde{E}_i\left(\mathrm{Fv}^{(1)},\dots,\mathrm{Fv}^{(i)}\right)$ denotes the function of $i$ even elements of $\mathrm{Gr}[\Psi]$ without constant term that is defined in \cite[Section~20.3.4]{QEDI}, formed with $\mathrm{Gcv}$.
It is multilinear and symmetric, and it is a polynomial in the coefficients of its arguments.
In this setting the lemma reads as follows.
\begin{itemize}
\item[(a)] Let $\mathrm{Fv}$ be even and without constant term, and let $\mathrm{alG}\cdot\mathrm{rw}_G(\mathrm{Fv};1)\le1$.
Then the series $\sum_{i\ge1}\frac{1}{i!}\tilde{E}_i(\mathrm{Fv},\dots,\mathrm{Fv})$ converges absolutely, and $\int\mathrm{d}\mu_{\mathrm{Gcv}}\;\exp(\mathrm{Fv})$ is the exponential of its sum.
\item[(c)] Let $\mathrm{Uo}$ be an open subset of $\mathbb{C}^{\mathrm{Cel}}$, and for $t\in\mathrm{Uo}$ let $\mathrm{Fv}(t)$ be even and without constant term, with coefficients that are analytic functions of $t$ on $\mathrm{Uo}$, such that $\mathrm{alG}\cdot\mathrm{rw}_G\left(\mathrm{Fv}(t);\mathrm{wg}_{\delta}\right)\le\frac{1}{2}$ on $\mathrm{Uo}$.
Then
\[
\mathrm{gf}(t):=\sum_{i\ge1}\frac{1}{i!}\tilde{E}_i\left(\mathrm{Fv}(t),\dots,\mathrm{Fv}(t)\right)
\]
is analytic on $\mathrm{Uo}$, and for $\mathfrak{c}\ne\mathfrak{c}'$ and non-empty sets $P$, $P'$ of unit sites its mixed second derivatives satisfy the bound (3.5):
\begin{equation}\tag{3.5}
\begin{split}
\left|\partial_{\mathfrak{c}}\partial_{\mathfrak{c}'}\mathrm{gf}(t)\right|\le{}&2\,\mathrm{Nw}_G\left(\partial_{\mathfrak{c}}\partial_{\mathfrak{c}'}\mathrm{Fv}(t);1\right)\\
&+2\,\mathrm{alG}\cdot\mathrm{Nw}_G\left(\partial_{\mathfrak{c}}\mathrm{Fv}(t);\mathrm{wg}_P\right)\cdot\mathrm{rw}_G\left(\partial_{\mathfrak{c}'}\mathrm{Fv}(t);\mathrm{wg}_{P'}\right)\cdot\exp\left(-\mathrm{vt}\cdot d(P,P')\right).
\end{split}
\end{equation}
\end{itemize}

\paragraph{(Q4)}
The following is \cite[Lemma~20.12]{QEDI}.
Let $\mathrm{gf}$ be holomorphic on $\left\{t\in\mathbb{C}^n:\ |t_i|<\mathrm{rad}\right\}$, with $|\mathrm{gf}|\le\mathrm{Mg}$ there and $|\mathrm{gf}(t)|\le\delta\le\mathrm{Mg}$ at the real points $t$ of this polydisc.
Then $|\mathrm{gf}(t)|\le\delta^{\frac{1}{3}}\cdot\mathrm{Mg}^{\frac{2}{3}}$ whenever $|t_i|\le\mathrm{rad}/(2n)$ for $i=1,\dots,n$.

\subsubsection*{0.7. Integrals of holomorphic functions}

\begin{itemize}
\item[(h1)] Let $n\ge1$ and $0<\mathrm{rad}_0\le\infty$, and put $\mathrm{Po}:=\left\{t\in\mathbb{C}^n:\ |t_i|<\mathrm{rad}_0\right\}$.
Let $(\mathrm{Ms},\mathrm{mq})$ be a measure space with a $\sigma$-finite positive measure, and let $\mathrm{gf}$ be a measurable complex function on $\mathrm{Po}\times\mathrm{Ms}$ which is holomorphic in $t$ for every point $\mathrm{oq}$ of $\mathrm{Ms}$ and has the following property: for every $\mathrm{rad}<\mathrm{rad}_0$ there is an $\mathrm{mq}$-integrable function $\mathrm{maj}_{\mathrm{rad}}\ge0$ with $|\mathrm{gf}(t,\mathrm{oq})|\le\mathrm{maj}_{\mathrm{rad}}(\mathrm{oq})$ whenever $|t_i|\le\mathrm{rad}$ for all $i$.
Then the function
\[
\mathrm{Gf}(t):=\int\mathrm{d}\,\mathrm{mq}(\mathrm{oq})\;\mathrm{gf}(t,\mathrm{oq})
\]
is holomorphic on $\mathrm{Po}$, and each Taylor coefficient of $\mathrm{Gf}$ at $0$ is the integral of the same Taylor coefficient of $\mathrm{gf}(\cdot\,,\mathrm{oq})$.
\begin{proof}
Let $|t_i|<\mathrm{rad}<\mathrm{rad}_0$ for all $i$, and write $\mathrm{d}\,\mathrm{zs}:=\mathrm{d}\,\mathrm{zs}_1\cdots\mathrm{d}\,\mathrm{zs}_n$.
Cauchy's formula in $n$ variables gives
\[
\mathrm{gf}(t,\mathrm{oq})=(2\pi i)^{-n}\oint_{|\mathrm{zs}_1|=\mathrm{rad}}\cdots\oint_{|\mathrm{zs}_n|=\mathrm{rad}}\mathrm{d}\,\mathrm{zs}\;\mathrm{gf}(\mathrm{zs},\mathrm{oq})\prod_{i}\left(\mathrm{zs}_i-t_i\right)^{-1}.
\]
The integrand is a measurable function on the product of the torus $\left\{|\mathrm{zs}_i|=\mathrm{rad}\ \text{for all}\ i\right\}$ and $\mathrm{Ms}$, and its modulus is at most $\mathrm{maj}_{\mathrm{rad}}(\mathrm{oq})\prod_i\left(\mathrm{rad}-|t_i|\right)^{-1}$, an integrable function on this product.
By Fubini's theorem the function $\mathrm{zs}\mapsto\mathrm{Gf}(\mathrm{zs})$ is integrable on the torus, and
\[
\mathrm{Gf}(t)=(2\pi i)^{-n}\oint\cdots\oint\mathrm{d}\,\mathrm{zs}\;\mathrm{Gf}(\mathrm{zs})\prod_{i}\left(\mathrm{zs}_i-t_i\right)^{-1}.
\]
For $t$ in a compact subset of $\left\{|t_i|<\mathrm{rad}\ \text{for all}\ i\right\}$ the expansion
\[
\prod_{i}\left(\mathrm{zs}_i-t_i\right)^{-1}=\sum_{\mathrm{ak}}\prod_{i}t_i^{\mathrm{ak}_i}\,\mathrm{zs}_i^{-\mathrm{ak}_i-1},
\]
the sum being over the multi-indices $\mathrm{ak}\in\{0,1,2,\dots\}^n$, converges uniformly in $\mathrm{zs}$ on the torus.
Hence $\mathrm{Gf}(t)=\sum_{\mathrm{ak}}\mathrm{Gf}_{\mathrm{ak}}\,t^{\mathrm{ak}}$ for these $t$, where, by Fubini's theorem once more,
\begin{align*}
\mathrm{Gf}_{\mathrm{ak}}&=(2\pi i)^{-n}\oint\cdots\oint\mathrm{d}\,\mathrm{zs}\;\mathrm{Gf}(\mathrm{zs})\prod_{i}\mathrm{zs}_i^{-\mathrm{ak}_i-1}\\
&=\int\mathrm{d}\,\mathrm{mq}(\mathrm{oq})\left[(2\pi i)^{-n}\oint\cdots\oint\mathrm{d}\,\mathrm{zs}\;\mathrm{gf}(\mathrm{zs},\mathrm{oq})\prod_{i}\mathrm{zs}_i^{-\mathrm{ak}_i-1}\right].
\end{align*}
The bracket is the Taylor coefficient of $\mathrm{gf}(\cdot\,,\mathrm{oq})$ with the multi-index $\mathrm{ak}$, and a convergent power series is holomorphic.
\end{proof}
\item[(h2)] Two entire functions on $\mathbb{C}^n$ that agree on $\mathbb{R}^n$ agree everywhere.
\begin{proof}
Their Taylor coefficients at $0$ are partial derivatives taken along real directions.
\end{proof}
If a function is holomorphic near $0$ and real at the real points near $0$, then its Taylor coefficients at $0$ are real.
\end{itemize}

\subsection*{1. Proposition 35.5}

\subsubsection*{Proof of Proposition 35.5}

\paragraph{1.1. The kernels. Clause (a).}
For $(\mathscr{A},\mathscr{C})\in\tilde{R}_{*}\times\mathfrak{X}_{*}$, Sections~I.5.4 and~I.3.6 give the expansion and the norm
\begin{align*}
\mathscr{E}_{*}(X;\mathscr{A},\mathscr{C},\psi,\chi)={}&\sum_{n,m}\frac{1}{n!\,m!}\sum_{\xi_1,\dots,\xi_n}\sum_{\boldsymbol{\zeta}_1,\dots,\boldsymbol{\zeta}_m}\left(\frac{\eta_*^{3}}{48}\right)^{n}\left(\frac{\eta_*^{3}}{48}\right)^{2m}\mathscr{E}_{nm}(X;\mathscr{A},\mathscr{C};\xi_1,\dots,\xi_n;\boldsymbol{\zeta}_1,\dots,\boldsymbol{\zeta}_m)\\
&\cdot\psi(\xi_1)\cdots\psi(\xi_n)\,\chi(\boldsymbol{\zeta}_1)\cdots\chi(\boldsymbol{\zeta}_m),\\
\left\|\mathscr{E}_{nm}(X;\mathscr{A},\mathscr{C})\right\|={}&\sum_{\xi_1,\dots,\xi_n}\sum_{\boldsymbol{\zeta}_1,\dots,\boldsymbol{\zeta}_m}\left(\frac{\eta_*^{3}}{48}\right)^{n+2m}\left|\mathscr{E}_{nm}(X;\mathscr{A},\mathscr{C};\xi_1,\dots,\xi_n;\boldsymbol{\zeta}_1,\dots,\boldsymbol{\zeta}_m)\right|;
\end{align*}
here $\xi=(c,g,\beta,\omega)$ ranges over $(c,g)\in T^{\mathrm{f}}_{*}$, $\beta\in\{1,\dots,4\}$ and $\omega\in\{0,1\}$, and $\boldsymbol{\zeta}=(x,x',\beta,\omega)$ over $(x,x')\in\mathrm{Pz}_{*}$, $\beta$ and $\omega$, the $48\cdot47$ ordered pairs inside each cluster being included.

Now let $X$ be an $M$-cube.
It contains $(M/\eta_*)^{3}$ photon sites, and Section~I.3.6 together with (I.1.3) gives
\[
\langle\bar{\psi},\psi\rangle_X=\int_X\mathrm{d}x\;\bar{\psi}(x)\cdot\psi(x)=\sum_{c\in X}\sum_{g\in O_{\mathrm{cub}}}\sum_{\beta=1}^{4}\frac{\eta_*^{3}}{48}\,\psi(c,g,\beta,1)\,\psi(c,g,\beta,0).
\]
This element is of bidegree $(2,0)$.
Define $\mathrm{Km}_{20}$ by
\[
\mathrm{Km}_{20}\big((x,\beta,1),(x',\beta',0)\big):=48\,\eta_*^{-3}\cdot\mathbf{1}[x=x']\cdot\mathbf{1}[\beta=\beta']\qquad(x\text{ over }X),
\]
by $\mathrm{Km}_{20}\big((x',\beta',0),(x,\beta,1)\big):=-\mathrm{Km}_{20}\big((x,\beta,1),(x',\beta',0)\big)$, and by $\mathrm{Km}_{20}:=0$ at all other arguments.
Then $\mathrm{Km}_{20}$ is antisymmetric, and
\begin{align*}
\frac{1}{2!}\sum_{\xi_1,\xi_2}\left(\frac{\eta_*^{3}}{48}\right)^{2}\mathrm{Km}_{20}(\xi_1,\xi_2)\,\psi(\xi_1)\,\psi(\xi_2)&=\frac{1}{2}\sum_{x\text{ over }X}\sum_{\beta}\frac{\eta_*^{3}}{48}\left[\psi(x,\beta,1)\,\psi(x,\beta,0)-\psi(x,\beta,0)\,\psi(x,\beta,1)\right]=\langle\bar{\psi},\psi\rangle_X,\\
\left\|\mathrm{Km}_{20}\right\|&=2\sum_{c\in X}\sum_{g\in O_{\mathrm{cub}}}\sum_{\beta}\left(\frac{\eta_*^{3}}{48}\right)^{2}\cdot48\,\eta_*^{-3}=2\cdot\left(\frac{M}{\eta_*}\right)^{3}\cdot48\cdot4\cdot\frac{\eta_*^{3}}{48}=8M^{3}.
\end{align*}
At pairwise distinct arguments an antisymmetric kernel coincides, up to a sign and the weights, with the coefficient of the corresponding ordered monomial; hence an element determines its antisymmetric kernels.
Consequently, for $t\in\mathrm{DDo}(t_x^{+})$, where $\left(\mathscr{A},\mathscr{C}(t)\right)\in\tilde{R}_{*}\times\mathfrak{X}_{*}$ by (f2),
\begin{equation}\tag{1.1}
\begin{aligned}
\mathrm{Ea}_{nm}(X;t)&=\mathscr{E}_{nm}\left(X;\mathscr{A},\mathscr{C}(t)\right)\qquad\text{for }(n,m)\ne(2,0),\\
\mathrm{Ea}_{20}(X;t)&=\mathscr{E}_{20}\left(X;\mathscr{A},\mathscr{C}(t)\right)-\mathbf{1}[X\text{ is an }M\text{-cube}]\cdot m_{*}\cdot\mathrm{Km}_{20}.
\end{aligned}
\end{equation}

Let $0<\mathrm{aw}\le e_*^{-\frac{1}{4}+\varepsilon}$ and $t\in\mathrm{DD}(t_x)$.
As $t_x<t_x^{+}$, the point $t$ lies in $\mathrm{DDo}(t_x^{+})$.
By (n1), $\mathrm{aw}\le h_{*,2}=e_*^{-\frac{1}{4}+\varepsilon}\le e_*^{-\frac{1}{4}}=h_{*,1}$, and therefore (2.1) of Proposition~35.4\,(SF3) yields
\[
\sum_{n,m}\frac{\mathrm{aw}^{n+m}}{n!\,m!}\left\|\mathscr{E}_{nm}\left(X;\mathscr{A},\mathscr{C}(t)\right)\right\|\le4\,e_*^{\varepsilon}\cdot w^{\mathrm{E}}_{*}\cdot\exp\left(-\kappa\cdot d_M(X)\right).
\]
If $X$ is an $M$-cube, then $d_M(X)=0$: by Section~I.1.1 and \cite[Section~5.2.3]{QEDI}, $M\cdot d_M(X)$ is the length of a shortest tree that meets all $M$-cubes of $X$, and for a single cube a point is such a tree.
Moreover, Proposition~35.4\,(SF1) gives
\[
\frac{\mathrm{aw}^{2}}{2!}\,|m_{*}|\cdot\left\|\mathrm{Km}_{20}\right\|=4M^{3}\cdot|m_{*}|\cdot\mathrm{aw}^{2}\le4M^{3}\cdot\frac{e_*^{\frac{3}{4}-8\varepsilon}}{8}\cdot e_*^{-\frac{1}{2}+2\varepsilon}=\frac{M^{3}}{2}\,e_*^{\frac{1}{4}-6\varepsilon}=\frac{M^{3}}{2}\,e_*^{\varepsilon}\cdot w^{\mathrm{E}}_{*}.
\]
From (1.1), the triangle inequality for $\|\cdot\|$ and the second condition of (g1) one obtains
\[
\sum_{n,m}\frac{\mathrm{aw}^{n+m}}{n!\,m!}\left\|\mathrm{Ea}_{nm}(X;t)\right\|\le\left(4+\frac{M^{3}}{2}\right)e_*^{\varepsilon}\cdot w^{\mathrm{E}}_{*}\cdot\exp\left(-\kappa\cdot d_M(X)\right)\le w^{\mathrm{E}}_{*}\cdot\exp\left(-\kappa\cdot d_M(X)\right).
\]

\paragraph{1.2. Clause (b).}
The set $\mathfrak{X}_{*}$ is open, since it is defined by strict inequalities (Section~I.5.3).
By Proposition~35.4\,(SF3), for fixed $\mathscr{A}$ each value $\mathscr{E}_{nm}(X;\mathscr{A},\mathscr{C};\xi_1,\dots,\xi_n;\boldsymbol{\zeta}_1,\dots,\boldsymbol{\zeta}_m)$ is an analytic function of $\mathscr{C}\in\mathfrak{X}_{*}$.
The map $t\mapsto\mathscr{C}(t)$ is linear, and by (f1) it maps $\mathrm{DDo}(t_x^{+})$ into $\mathfrak{X}_{*}$; the term $m_{*}\cdot\mathrm{Km}_{20}$ does not contain $t$.
By (1.1), every value of $\mathrm{Ea}_{nm}(X;t;\cdot\,)$ is therefore an analytic function of $t$ on $\mathrm{DDo}(t_x^{+})$.

Now let $t\in\mathrm{DD}(t_x/2)$ and $\mathfrak{c}\in\mathrm{Cel}$.
For $|\mathrm{zs}|=t_x/2$ the point $t+\mathrm{zs}\cdot\mathrm{ue}_{\mathfrak{c}}$ belongs to $\mathrm{DD}(t_x)$, and Cauchy's formula gives, at every argument,
\[
\partial_{\mathfrak{c}}\mathrm{Ea}_{nm}(X;t;\cdot\,)=(2\pi i)^{-1}\oint_{|\mathrm{zs}|=t_x/2}\mathrm{d}\,\mathrm{zs}\;\mathrm{Ea}_{nm}\left(X;t+\mathrm{zs}\cdot\mathrm{ue}_{\mathfrak{c}};\cdot\,\right)\mathrm{zs}^{-2}.
\]
Hence $\left|\partial_{\mathfrak{c}}\mathrm{Ea}_{nm}(X;t;\cdot\,)\right|$ is at most $2/t_x=q_x$ times the mean value of $\left|\mathrm{Ea}_{nm}\left(X;t+\mathrm{zs}\cdot\mathrm{ue}_{\mathfrak{c}};\cdot\,\right)\right|$ over the circle $|\mathrm{zs}|=t_x/2$.
One multiplies by $\left(\eta_*^{3}/48\right)^{n+2m}\mathrm{aw}^{n+m}/(n!\,m!)$ and sums over the arguments and over $(n,m)$ under the mean value; all these sums are finite.
Clause (a) at the points $t+\mathrm{zs}\cdot\mathrm{ue}_{\mathfrak{c}}$ then gives
\[
\sum_{n,m}\frac{\mathrm{aw}^{n+m}}{n!\,m!}\left\|\partial_{\mathfrak{c}}\mathrm{Ea}_{nm}(X;t)\right\|\le q_x\cdot w^{\mathrm{E}}_{*}\cdot\exp\left(-\kappa\cdot d_M(X)\right).
\]
Let $\mathfrak{c}\ne\mathfrak{c}'$.
For $|\mathrm{zs}|=|\mathrm{zs}'|=t_x/2$ the point $t+\mathrm{zs}\cdot\mathrm{ue}_{\mathfrak{c}}+\mathrm{zs}'\cdot\mathrm{ue}_{\mathfrak{c}'}$ belongs to $\mathrm{DD}(t_x)$, and
\[
\partial_{\mathfrak{c}}\partial_{\mathfrak{c}'}\mathrm{Ea}_{nm}(X;t;\cdot\,)=(2\pi i)^{-2}\oint_{|\mathrm{zs}|=t_x/2}\oint_{|\mathrm{zs}'|=t_x/2}\mathrm{d}\,\mathrm{zs}\,\mathrm{d}\,\mathrm{zs}'\;\mathrm{Ea}_{nm}\left(X;t+\mathrm{zs}\cdot\mathrm{ue}_{\mathfrak{c}}+\mathrm{zs}'\cdot\mathrm{ue}_{\mathfrak{c}'};\cdot\,\right)\mathrm{zs}^{-2}\,\mathrm{zs}'^{-2}.
\]
Its modulus is at most $q_x^{2}$ times the mean value over the two circles of $\left|\mathrm{Ea}_{nm}\left(X;t+\mathrm{zs}\cdot\mathrm{ue}_{\mathfrak{c}}+\mathrm{zs}'\cdot\mathrm{ue}_{\mathfrak{c}'};\cdot\,\right)\right|$.
Summing under the mean value as before and applying (a), one arrives at the second inequality of (b).

\paragraph{1.3. Clause (c).}
By Proposition~35.4\,(SF3), the function $\mathscr{C}\mapsto\mathscr{E}_{nm}(X;\mathscr{A},\mathscr{C};\xi_1,\dots,\xi_n;\boldsymbol{\zeta}_1,\dots,\boldsymbol{\zeta}_m)$ on $\mathfrak{X}_{*}$ depends on the values $\mathscr{C}(b)$ for the bonds $b$ that meet $X$ only.
If no point of $\mathrm{Ypr}_{\mathfrak{c}}$ lies in $X$, its value at $\mathscr{C}(t)$ is therefore independent of $t_{\mathfrak{c}}$, by (m3).
The second term in (1.1) does not contain $t$.
Hence $\partial_{\mathfrak{c}}\mathrm{Ea}_{nm}(X;t)=0$.

\paragraph{1.4. Clause (d).}
The sets $\square\times O_{\mathrm{cub}}$, where $\square$ ranges over the $M$-cubes of $T_{*}$, form a partition of $T^{\mathrm{f}}_{*}$ (Section~I.1.4).
Therefore
\[
\langle\bar{\psi},\psi\rangle=\int\mathrm{d}x\;\bar{\psi}(x)\cdot\psi(x)=\sum_{\square}\langle\bar{\psi},\psi\rangle_{\square}.
\]
Every $M$-cube is a polymer (Section~I.1.1).
It follows that, in $\mathrm{Gr}[\psi,\chi]$,
\[
-m_{*}\cdot\langle\bar{\psi},\psi\rangle+\sum_{X\in D_{*}}\mathscr{E}_{*}\left(X;\mathscr{A},\mathscr{C}(t),\psi,\chi\right)=\sum_{X\in D_{*}}\mathrm{Ea}(X;t).
\]
The insertion of the fields formed with $\mathscr{T}(t)$ is the homomorphism of $\mathrm{Gr}[\psi,\chi]$ into $\mathrm{Gr}[\Psi]$ given by (Section~G5.1.7, paragraph ``Legs'')
\begin{equation}\tag{1.2}
\psi(\xi)\longmapsto\sum_{u}\mathrm{Leg}_{\mathscr{T}(t)}(\xi;u)\,\Psi(u),\qquad\chi(\boldsymbol{\zeta})\longmapsto\sum_{u}\mathrm{Leg}_{\mathscr{T}(t)}(\boldsymbol{\zeta};u)\,\Psi(u).
\end{equation}
Writing $D_{*}\left(\mathscr{T}(t)\right)=D_{*}(\mathscr{A})+\mathrm{Kq}(t)$, one sees that the exponent in (G5.5.1) equals $-\left\langle\bar{\Psi},D_{*}(\mathscr{A})\Psi\right\rangle+\mathrm{Vj}(t)$.
Each monomial of $\mathrm{Ea}(X;t)$ has as many barred as unbarred factors (Section~I.5.4; (1.1)), and (1.2) maps a monomial of bidegree $(n,m)$ to a homogeneous element of degree $n+m$.
Hence $\mathrm{Vj}(t)$ is even; its component of degree zero comes from $(n,m)=(0,0)$ alone and equals
\[
\mathrm{Vj}_0(t)=\sum_{X}\mathrm{Ea}_{00}(X;t)=\sum_{X}\mathscr{E}_{00}\left(X;\mathscr{A},\mathscr{C}(t)\right),
\]
and $\mathrm{Vj}_{+}(t)=\mathrm{Vj}(t)-\mathrm{Vj}_0(t)$ is even and has no constant term.
Since even elements commute, the integrand of (G5.5.1) equals
\[
\exp\left(-\left\langle\bar{\Psi},D_{*}(\mathscr{A})\Psi\right\rangle\right)\cdot\exp\left(\mathrm{Vj}_0(t)\right)\cdot\exp\left(\mathrm{Vj}_{+}(t)\right).
\]
Proposition~15\,(l)\,(iii) at $\nu=1$, applied with $\mathrm{Cov}=\mathrm{Gcv}$ and $\mathrm{Cov}^{-1}=D_{*}(\mathscr{A})$, gives
\[
\int\mathrm{D}\Psi\;\exp\left(-\left\langle\bar{\Psi},\mathrm{Gcv}^{-1}\Psi\right\rangle\right)\cdot\exp\left(\mathrm{Vj}_{+}(t)\right)=\det\left(\mathrm{Gcv}^{-1}\right)\int\mathrm{d}\mu_{\mathrm{Gcv}}\;\exp\left(\mathrm{Vj}_{+}(t)\right).
\]
This proves (5.3).
The coefficients of $\mathrm{Vj}(t)$ are finite sums of products of the weights $\eta_*^{3}/48$ and of the numbers
\[
D_{*}\left(\mathscr{T}(t)\right)(y,y')_{\beta\beta'}-D_{*}(\mathscr{A})(y,y')_{\beta\beta'},\qquad\mathrm{Ea}_{nm}(X;t;\cdot\,),\qquad\mathrm{Leg}_{\mathscr{T}(t)}(\mathrm{up};u).
\]
By (f3) and 1.2 these numbers are analytic functions of $t$ on $\mathrm{DDo}(t_x^{+})$.

\paragraph{1.5. Clause (N0).}
Let $\hat{\mathfrak{a}}=0$.
Then $\hat{c}_{\mathfrak{c}}=0$ for every $\mathfrak{c}$, and $\mathscr{C}(t)=0$ for every $t$ (Section~G5.1.4).
Hence $\mathscr{T}(t)=\mathscr{A}$, $\mathrm{Kq}(t)=D_{*}(\mathscr{A})-D_{*}(\mathscr{A})=0$,
\[
\mathrm{Ea}(X;t)=\mathscr{E}_{*}(X;\mathscr{A},0,\cdot\,)-\mathbf{1}[X\text{ is an }M\text{-cube}]\cdot m_{*}\cdot\langle\bar{\psi},\psi\rangle_X,
\]
and the homomorphism (1.2) is formed with $\mathscr{A}$.
None of these objects contains $t$, and so $\mathrm{Vj}(t)$, $\mathrm{Vj}_0(t)$ and $\mathrm{Vj}_{+}(t)$ do not depend on $t$.

\paragraph{1.6. Step 0: preliminaries for (N1)--(N3).}

\emph{(i) Quadratic forms.}
Let $\mathrm{Kf}$ be a linear operator on $\mathfrak{S}(1,\ell_*)$ with the $4\times4$ blocks $\mathrm{Kf}(y,y')$, and put (Section~I.3.6)
\[
\mathrm{Fv}:=-\left\langle\bar{\Psi},\mathrm{Kf}\,\Psi\right\rangle=-\sum_{y,y'}\sum_{\beta,\beta'}\Psi(y,\beta,1)\,\mathrm{Kf}(y,y')_{\beta\beta'}\,\Psi(y',\beta',0).
\]
Its antisymmetric kernel is given by
\[
\mathrm{Fv}_2\big((y,\beta,1),(y',\beta',0)\big)=-\mathrm{Kf}(y,y')_{\beta\beta'}=-\mathrm{Fv}_2\big((y',\beta',0),(y,\beta,1)\big),
\]
and $\mathrm{Fv}_2=0$ at all other arguments; indeed
\begin{align*}
\frac{1}{2!}\sum_{u_1,u_2}\mathrm{Fv}_2(u_1,u_2)\,\Psi(u_1)\,\Psi(u_2)
&=\frac{1}{2}\sum_{y,y'}\sum_{\beta,\beta'}\Bigl[-\mathrm{Kf}(y,y')_{\beta\beta'}\,\Psi(y,\beta,1)\,\Psi(y',\beta',0)\\
&\qquad\qquad+\mathrm{Kf}(y,y')_{\beta\beta'}\,\Psi(y',\beta',0)\,\Psi(y,\beta,1)\Bigr]\\
&=\mathrm{Fv}.
\end{align*}
The weights $\mathrm{wg}_{\delta}$, $\mathrm{wg}_P$ and $1$ depend on the sites of their arguments only; for them one writes $\mathrm{wg}(y,y')$.
By (M2) of Section~G5.1.7, $|\mathrm{Kf}(y,y')|_c\le|\mathrm{Kf}|(y,y')$.
Since $|\mathrm{Kf}(y,y')|_c$ is the largest row sum of the moduli of the entries of the block (Section~I.1.5), every row of the block has $\ell^1$ norm at most $|\mathrm{Kf}|(y,y')$, and every column has $\ell^1$ norm at most $4\,|\mathrm{Kf}|(y,y')$.
For $u_1=(y,\beta,1)$ the sum over $u_2$ is a sum over $y'$ of row sums, and for $u_1=(y',\beta',0)$ it is a sum over $y$ of column sums.
The definitions of Section~G5.1.6, with $m=2$, therefore give
\begin{align*}
\mathrm{rw}_G(\mathrm{Fv};\mathrm{wg})&\le4\,h_G^{2}\cdot\max\left\{\sup_{y}\sum_{y'}|\mathrm{Kf}|(y,y')\cdot\mathrm{wg}(y,y'),\ \sup_{y'}\sum_{y}|\mathrm{Kf}|(y,y')\cdot\mathrm{wg}(y,y')\right\},\\
\mathrm{Nw}_G(\mathrm{Fv};\mathrm{wg})&\le8\,h_G^{2}\sum_{y,y'}|\mathrm{Kf}|(y,y')\cdot\mathrm{wg}(y,y').
\end{align*}

\emph{(ii) Kernels after the insertion.}
The letter $\mathrm{up}$ denotes a leg index.
For $n+m=p$, the first $n$ of the indices $\mathrm{up}_1,\dots,\mathrm{up}_p$ are indices $\xi$ and the last $m$ are indices $\boldsymbol{\zeta}$, and $\int\mathrm{d}\,\mathrm{up}$ denotes the weighted sum
\[
\int\mathrm{d}\,\mathrm{up}:=\sum_{\mathrm{up}_1,\dots,\mathrm{up}_p}\left(\frac{\eta_*^{3}}{48}\right)^{n+2m}.
\]
By (1.2) the image of $\mathrm{Ea}(X;t)$ in $\mathrm{Gr}[\Psi]$ is
\begin{gather*}
\sum_{n,m}\frac{1}{n!\,m!}\sum_{u_1,\dots,u_p}\mathrm{Kv}_{nm}(X;t;u_1,\dots,u_p)\,\Psi(u_1)\cdots\Psi(u_p),\\
\mathrm{Kv}_{nm}(X;t;u_1,\dots,u_p):=\int\mathrm{d}\,\mathrm{up}\;\mathrm{Ea}_{nm}(X;t;\mathrm{up}_1,\dots,\mathrm{up}_p)\prod_{i=1}^{p}\mathrm{Leg}_{\mathscr{T}(t)}(\mathrm{up}_i;u_i).
\end{gather*}
Its component of degree $p$ is
\[
\frac{1}{p!}\sum_{u_1,\dots,u_p}\left[\sum_{n+m=p}\frac{p!}{n!\,m!}\,\mathrm{Kv}_{nm}(X;t;u_1,\dots,u_p)\right]\Psi(u_1)\cdots\Psi(u_p),
\]
and
\[
\frac{h_G^{p}}{(p-1)!}\cdot\frac{p!}{n!\,m!}=\frac{p\cdot h_G^{p}}{n!\,m!}.
\]

\emph{(iii) Positions.}
If $\mathrm{Ea}_{nm}(X;t;\mathrm{up}_1,\dots,\mathrm{up}_p)\ne0$, then $\mathrm{pos}(\mathrm{up}_i)\subset X^{+}$ for every $i$.
Indeed, by Proposition~35.4\,(SF3) and (1.1), each $\xi_i$ lies over $X$, and each $\boldsymbol{\zeta}_j$ is admissible for $X$ (Section~I.1.7).
Three cases occur.
If $\mathrm{up}=\xi=(x,\beta,\omega)$ with $x=(c,g)$ and $c\in X$, then $\mathrm{pos}(\mathrm{up})=\{[c]\}$, and $[c]\in X$ by (m1).
If $\mathrm{up}=\boldsymbol{\zeta}$ is of kind (I), with $x=(c,g)$ and $x'=(c',g')$, then $c$ and $c'$ lie in $X$, so that $[c],[c']\in X$ by (m1); this case contains the pairs inside a cluster ($c=c'$, $g\ne g'$), for which $d_X=\eta_*<1$ and $\mathrm{pos}(\mathrm{up})=\{[c]\}$.
If $\mathrm{up}=\boldsymbol{\zeta}$ is of kind (II), with $x'=x^{(3)}$, then $d_X\left(x,x^{(3)}\right)=\max\left(|c-c_3|,\eta_*\right)=\eta_*<1$ by (n1), whence $\mathrm{pos}(\mathrm{up})=\{[c]\}$.
In this case $c$ or $c_3$ lies in $X$; if $c_3$ does, then $[c_3]\in X\cap T^0_{*}$ and $d\left([c],[c_3]\right)<\frac{1}{2}+\eta_*+\frac{1}{2}<2$, and since this distance is an integer, $d\left([c],[c_3]\right)\le1$, so that $[c]\in X^{+}$.

By Proposition~35.4\,(SF4)\,(v), any two points of $X$ have distance at most $M\cdot\left(d_M(X)+2\right)$.
Hence $\operatorname{diam}(X^{+})\le M\cdot\left(d_M(X)+2\right)+2$, and since $\mathrm{vt}\le1/M$ and $M\ge2$, one obtains the bound (5.4) on the diameter of $X^{+}$:
\begin{equation}\tag{5.4}
\mathrm{vt}\cdot\operatorname{diam}(X^{+})\le d_M(X)+3.
\end{equation}

\emph{(iv) Numbers.}
For integers $p\ge2$ one has $p\le5^{p}$ and $p\,(p+1)^{2}\le5^{p}$; the latter holds for $p=2$, where it reads $18\le25$, and $(p+1)(p+2)^{2}\le5p\,(p+1)^{2}$ for $p\ge2$.
Moreover
\[
\sum_{p\ge2}p\cdot5^{-p}=\frac{5}{16}-\frac{1}{5}=\frac{9}{80},\qquad\sum_{p\ge2}p^{2}\cdot5^{-p}=\frac{15}{32}-\frac{1}{5}=\frac{43}{160}.
\]
By the first condition of (g1), the number $\mathrm{aw}:=5\,\mathrm{hL}\cdot h_G$ satisfies $\mathrm{aw}\le e_*^{-\frac{1}{4}+\varepsilon}$; therefore (a) and (b) hold with this value of $\mathrm{aw}$.

\emph{(v) Sums over cubes and polymers.}
Let $z\in T^0_{*}$ and $\gamma\ge\gamma^{F}/4$.
Then
\[
\sum_{\square}\exp\left(-\gamma\cdot d\left(z,\square^{+}\right)\right)\le32,
\]
the sum being over the $M$-cubes $\square$ of $T_{*}$.
\begin{proof}
Let $\square_0$ be the $M$-cube that contains $z$, and let $z_0\in M\mathbb{Z}^3$ be its centre.
Every $M$-cube has a centre of the form $z_0+M\cdot i$ with an $i\in\mathbb{Z}^3$ whose sup norm $|i|$ is the least in its class modulo $(\ell_*/M)\mathbb{Z}^3$, and at most $(2j+1)^{3}$ cubes have $|i|=j$.
Since $M$ is odd (Section~I.1.1), on suitable lifts the unit sites $y$ of such a cube and the site $z$ satisfy $\left|y-z_0-M\cdot i\right|\le(M-1)/2$ and $|z-z_0|\le(M-1)/2$.
Hence $d(z,y)\ge M\cdot|i|-(M-1)$, and $d\left(z,\square^{+}\right)\ge\left(|i|-1\right)\cdot M$.
Moreover $\gamma\cdot M\ge\gamma^{F}\cdot M/4\ge4$.
As $27$ cubes have $|i|\le1$, the sum is at most
\[
27+\sum_{j\ge2}(2j+1)^{3}\exp\left(-4(j-1)\right).
\]
Here $125\,\mathrm{e}^{-4}\le2.29$ and $343\,\mathrm{e}^{-8}\le0.12$, and for $j\ge4$ the quotient of two consecutive terms is at most $(11/9)^{3}\,\mathrm{e}^{-4}\le1/25$, so that
\[
\sum_{j\ge4}(2j+1)^{3}\,\mathrm{e}^{-4(j-1)}\le729\,\mathrm{e}^{-12}\cdot\frac{25}{24}\le0.01.
\]
The sum is therefore at most $27+2.29+0.12+0.01\le32$.
\end{proof}
The set $X^{+}$ is the union of the sets $\square^{+}$, $\square$ ranging over the cubes of $X$; hence
\[
\exp\left(-\gamma\cdot d\left(z,X^{+}\right)\right)\le\sum_{\square\subset X}\exp\left(-\gamma\cdot d\left(z,\square^{+}\right)\right).
\]
Together with $\kappa-1\ge\kappa_0$ and Proposition~35.4\,(SF4)\,(v), this yields the bound (5.5) for sums over the polymers:
\begin{equation}\tag{5.5}
\sum_{X\in D_{*}}\exp\left(-\gamma\cdot d\left(z,X^{+}\right)\right)\cdot\exp\left(-(\kappa-1)\cdot d_M(X)\right)\le\sum_{\square}\exp\left(-\gamma\cdot d\left(z,\square^{+}\right)\right)\sum_{X\supset\square}\exp\left(-\kappa_0\cdot d_M(X)\right)\le32\,K_0.
\end{equation}

\emph{(vi) Legs without derivative.}
Let $\mathrm{Aq}\in\mathrm{Gwk}$, $X\in D_{*}$ and $n+m=p\ge2$.
Let $\mathrm{Gk}(\mathrm{up}_1,\dots,\mathrm{up}_p)$ be a complex function of $n$ indices $\xi$ and $m$ indices $\boldsymbol{\zeta}$ such that $\mathrm{pos}(\mathrm{up}_i)\subset X^{+}$ for all $i$ wherever $\mathrm{Gk}\ne0$, and put
\[
\left\|\mathrm{Gk}\right\|:=\int\mathrm{d}\,\mathrm{up}\;\left|\mathrm{Gk}\right|,\qquad\mathrm{Kv}(u_1,\dots,u_p):=\int\mathrm{d}\,\mathrm{up}\;\mathrm{Gk}(\mathrm{up}_1,\dots,\mathrm{up}_p)\prod_{i}\mathrm{Leg}_{\mathrm{Aq}}(\mathrm{up}_i;u_i).
\]
Let $\mathrm{wg}$ be a weight and $\mathrm{cw}$ a number such that
\begin{equation}\tag{1.3}
\mathrm{wg}(u_1,\dots,u_p)\le\mathrm{cw}\cdot\prod_{i=1}^{p}\exp\left(\mathrm{vt}\cdot d\left(u_i,\mathrm{pos}(\mathrm{up}_i)\right)\right)
\end{equation}
for all $u_1,\dots,u_p$ and all $(\mathrm{up}_1,\dots,\mathrm{up}_p)$ with $\mathrm{Gk}(\mathrm{up}_1,\dots,\mathrm{up}_p)\ne0$.
Then
\[
\sum_{u_1,\dots,u_p}|\mathrm{Kv}|\cdot\mathrm{wg}\le\mathrm{cw}\cdot\mathrm{hL}^{p}\cdot\left\|\mathrm{Gk}\right\|,
\]
and for every slot $\mathrm{sl}$ and every index $u$
\[
\sum_{\substack{u_1,\dots,u_p\\ u_{\mathrm{sl}}=u}}\left|\mathrm{Kv}(u_1,\dots,u_p)\right|\cdot\mathrm{wg}(u_1,\dots,u_p)\le\mathrm{cw}\cdot\mathrm{hL}^{p}\cdot\exp\left(-\frac{\gamma^{F}}{4}\cdot d\left(u,X^{+}\right)\right)\left\|\mathrm{Gk}\right\|.
\]
\begin{proof}
By (1.3),
\[
\left|\mathrm{Kv}(u_1,\dots,u_p)\right|\cdot\mathrm{wg}(u_1,\dots,u_p)\le\mathrm{cw}\cdot\int\mathrm{d}\,\mathrm{up}\;\left|\mathrm{Gk}(\mathrm{up}_1,\dots,\mathrm{up}_p)\right|\prod_{i}\left\{\left|\mathrm{Leg}_{\mathrm{Aq}}(\mathrm{up}_i;u_i)\right|\exp\left(\mathrm{vt}\cdot d\left(u_i,\mathrm{pos}(\mathrm{up}_i)\right)\right)\right\}.
\]
Proposition~33.2\,(L1) with $\gamma=\mathrm{vt}\le\gamma^{F}/2$ gives
\[
\sum_{u_i}\left|\mathrm{Leg}_{\mathrm{Aq}}(\mathrm{up}_i;u_i)\right|\exp\left(\mathrm{vt}\cdot d\left(\mathrm{pos}(\mathrm{up}_i),u_i\right)\right)\le\mathrm{hL}.
\]
For $u=(y',\beta',\omega')$, Proposition~33.2\,(L1) gives moreover
\begin{align*}
\left|\mathrm{Leg}_{\mathrm{Aq}}(\mathrm{up}_{\mathrm{sl}};u)\right|\exp\left(\mathrm{vt}\cdot d\left(\mathrm{pos}(\mathrm{up}_{\mathrm{sl}}),u\right)\right)&\le\left|\mathrm{Leg}_{\mathrm{Aq}}\right|(\mathrm{up}_{\mathrm{sl}};y')\exp\left(\mathrm{vt}\cdot d\left(\mathrm{pos}(\mathrm{up}_{\mathrm{sl}}),y'\right)\right)\\
&\le\mathrm{hL}\cdot\exp\left(-\left(\gamma^{F}-\mathrm{vt}\right)d\left(\mathrm{pos}(\mathrm{up}_{\mathrm{sl}}),y'\right)\right)\\
&\le\mathrm{hL}\cdot\exp\left(-\frac{\gamma^{F}}{4}\cdot d\left(u,X^{+}\right)\right),
\end{align*}
because $\mathrm{pos}(\mathrm{up}_{\mathrm{sl}})\subset X^{+}$ and $\gamma^{F}-\mathrm{vt}\ge\gamma^{F}/4$.
\end{proof}

\emph{(vii) Summation.}
Let $\mathrm{Aq}\in\mathrm{Gwk}$, let $\mathrm{Dsub}$ be a subset of $D_{*}$, and for $X\in\mathrm{Dsub}$ and $n+m\ge2$ let $\mathrm{Gk}_{nm}(X;\cdot\,)$ be functions as in (vi), with the kernels $\mathrm{Kv}_{nm}(X;\cdot\,)$.
Let $\mathcal{G}\ge0$ be a number with

\[
\sum_{n+m\ge2}\frac{\left(5\,\mathrm{hL}\cdot h_G\right)^{n+m}}{n!\,m!}\cdot\left\|\mathrm{Gk}_{nm}(X)\right\|\le\mathcal{G}\cdot\exp\left(-\kappa\cdot d_M(X)\right)\qquad\left(X\in\mathrm{Dsub}\right),
\]
and let $\mathrm{wg}$ be a weight which satisfies (1.3) with $\mathrm{cw}=\exp\left(d_M(X)+3\right)$ for every $X\in\mathrm{Dsub}$ and every $(n,m)$.
Put
\[
\mathrm{Fv}:=\sum_{X\in\mathrm{Dsub}}\sum_{n+m\ge2}\frac{1}{n!\,m!}\sum_{u_1,\dots,u_p}\mathrm{Kv}_{nm}(X;u_1,\dots,u_p)\cdot\Psi(u_1)\cdots\Psi(u_p)\qquad\left(p=n+m\right).
\]
In the two applications below one has $\mathrm{Gk}_{nm}=0$ whenever $n+m$ is odd, so that $\mathrm{Fv}$ is even.
Then
\[
\mathrm{rw}_G\left(\mathrm{Fv};\mathrm{wg}\right)\le650\,K_0\cdot\mathcal{G},\qquad
\mathrm{Nw}_G\left(\mathrm{Fv};\mathrm{wg}\right)\le\mathrm{e}^3\cdot\mathcal{G}\cdot\sum_{X\in\mathrm{Dsub}}\exp\left(-(\kappa-1)\cdot d_M(X)\right).
\]

\emph{Proof of (vii).}
Every term on the left side of the hypothesis is non-negative.
Hence, for every $p\ge2$,
\[
\sum_{n+m=p}\frac{\left(\mathrm{hL}\cdot h_G\right)^p}{n!\,m!}\cdot\left\|\mathrm{Gk}_{nm}(X)\right\|\le5^{-p}\cdot\mathcal{G}\cdot\exp\left(-\kappa\cdot d_M(X)\right).
\]
We combine (3.3), (ii), (vi), and (5.5) with $\gamma=\gamma^F/4$; since $\mathrm{Dsub}\subseteq D_*$, the sum over $\mathrm{Dsub}$ is at most the sum over $D_*$ in (5.5).
This gives
\begin{align*}
\mathrm{rw}_G\left(\mathrm{Fv};\mathrm{wg}\right)
&\le\sum_{p\ge2}\sup_u\sum_{X\in\mathrm{Dsub}}\sum_{n+m=p}\frac{p\cdot h_G^p}{n!\,m!}\cdot\exp\left(d_M(X)+3\right)\cdot\mathrm{hL}^p\cdot\exp\left(-\frac{\gamma^F}{4}\cdot d\left(u,X^+\right)\right)\cdot\left\|\mathrm{Gk}_{nm}(X)\right\|\\
&\le\sum_{p\ge2}p\cdot5^{-p}\cdot\mathrm{e}^3\cdot\mathcal{G}\cdot\sup_u\sum_{X\in\mathrm{Dsub}}\exp\left(-\frac{\gamma^F}{4}\cdot d\left(u,X^+\right)\right)\cdot\exp\left(-(\kappa-1)\cdot d_M(X)\right)\\
&\le\frac{9}{80}\cdot32\,\mathrm{e}^3\cdot K_0\cdot\mathcal{G}\le650\,K_0\cdot\mathcal{G},
\end{align*}
where $32\,\mathrm{e}^3\le643$ has been used in the last step.
In the same way, using $p\le5^p$,
\begin{align*}
\mathrm{Nw}_G\left(\mathrm{Fv};\mathrm{wg}\right)
&\le\sum_{p\ge2}\sum_{X\in\mathrm{Dsub}}\sum_{n+m=p}\frac{p\cdot h_G^p}{n!\,m!}\cdot\exp\left(d_M(X)+3\right)\cdot\mathrm{hL}^p\cdot\left\|\mathrm{Gk}_{nm}(X)\right\|\\
&\le\mathrm{e}^3\cdot\mathcal{G}\cdot\sum_{X\in\mathrm{Dsub}}\exp\left(-(\kappa-1)\cdot d_M(X)\right).
\end{align*}
This proves (vii).

\paragraph{1.7. Step 1: (N1).}
Let $t\in\mathrm{DD}(t_x)$.
By the definition (5.2) in Proposition~35.5 and by (ii),
\[
\mathrm{Vj}_+(t)=-\left\langle\bar{\Psi},\mathrm{Kq}(t)\Psi\right\rangle+\mathrm{Fv},
\]
where $\mathrm{Fv}$ now denotes the sum over $X\in D_*$ of the parts of degree at least $2$ of the image of $\mathrm{Ea}(X;t)$.
By (Q2), $\mathrm{rw}_G$ is subadditive, and we estimate the two terms separately.

\emph{The quadratic term.}
By (f2), $\mathscr{A}$ and $\mathscr{A}+\mathscr{C}(t)$ lie in $\mathrm{Gwk}$, together with the segment joining them.
We apply Proposition~33.2\,(D3) with $A_0=\mathscr{A}$ and $B=\mathscr{C}(t)$; by (f1) and (n1),
\[
\epsilon_B=e_*\cdot\left|\mathscr{C}(t)\right|_\infty<e_*^{3\varepsilon}\le e_*^{2\varepsilon},
\]
and therefore
\[
\left\|\mathrm{Kq}(t)\right\|_{\gamma^F/2}\le C_K\cdot e_*^{2\varepsilon}.
\]
For a pair of indices $u_1$, $u_2$ with the sites $y$, $y'$ one has $\mathrm{wg}_\delta(u_1,u_2)=\exp\left(\mathrm{vt}\cdot d(y,y')\right)$, and $\mathrm{vt}\le\gamma^F/2$.
Consequently the row sums and the column sums that occur in (i) do not exceed $\left\|\mathrm{Kq}(t)\right\|_{\gamma^F/2}$, and
\[
\mathrm{rw}_G\left(-\left\langle\bar{\Psi},\mathrm{Kq}(t)\Psi\right\rangle;\mathrm{wg}_\delta\right)\le4\,h_G^2\cdot C_K\cdot e_*^{2\varepsilon}.
\]

\emph{The activities.}
We apply (vii) with
\[
\mathrm{Aq}=\mathscr{T}(t),\qquad\mathrm{Dsub}=D_*,\qquad\mathrm{Gk}_{nm}(X;\cdot\,)=\mathrm{Ea}_{nm}(X;t;\cdot\,),\qquad\mathcal{G}=w^{\mathrm{E}}_*,\qquad\mathrm{wg}=\mathrm{wg}_\delta.
\]
The condition on the positions required in (vi) is (iii).
The hypothesis on $\mathcal{G}$ is clause (a) with $\mathrm{aw}=5\,\mathrm{hL}\cdot h_G$, which is admissible by (iv).
It remains to verify (1.3).
For each $i$ let $\mathrm{zn}_i\in\mathrm{pos}(\mathrm{up}_i)$ be a point nearest to $u_i$.
For $i\ne j$,
\[
d(u_i,u_j)\le d(u_i,\mathrm{zn}_i)+d(\mathrm{zn}_i,\mathrm{zn}_j)+d(\mathrm{zn}_j,u_j)\le d\left(u_i,\mathrm{pos}(\mathrm{up}_i)\right)+\operatorname{diam}\left(X^+\right)+d\left(u_j,\mathrm{pos}(\mathrm{up}_j)\right),
\]
since $\mathrm{zn}_i$ and $\mathrm{zn}_j$ lie in $X^+$ (by (iii)).
Thus
\[
\mathrm{vt}\cdot\max_{i,j}d(u_i,u_j)\le\mathrm{vt}\cdot\operatorname{diam}\left(X^+\right)+\sum_i\mathrm{vt}\cdot d\left(u_i,\mathrm{pos}(\mathrm{up}_i)\right),
\]
and (5.4) yields
\[
\mathrm{wg}_\delta(u_1,\dots,u_p)\le\exp\left(d_M(X)+3\right)\cdot\prod_i\exp\left(\mathrm{vt}\cdot d\left(u_i,\mathrm{pos}(\mathrm{up}_i)\right)\right).
\]
Hence (vii) gives $\mathrm{rw}_G\left(\mathrm{Fv};\mathrm{wg}_\delta\right)\le650\,K_0\cdot w^{\mathrm{E}}_*$.

\emph{Conclusion of Step 1.}
Adding the two bounds, one obtains
\[
\mathrm{rw}_G\left(\mathrm{Vj}_+(t);\mathrm{wg}_\delta\right)\le4\,h_G^2\cdot C_K\cdot e_*^{2\varepsilon}+650\,K_0\cdot e_*^{\frac{1}{4}-7\varepsilon}=\bar{r}.
\]
The sets $\mathrm{Ypr}_{\mathfrak{c}}$ do not enter this step.

\paragraph{1.8. Step 2: (N2).}
Let $\hat{\mathfrak{a}}\ne0$.
Then every $\mathrm{Ypr}_{\mathfrak{c}}$ and every $\mathrm{Ypr}_{\mathfrak{c}}^+$ is non-empty (Proposition~35.3\,(g)), so that the weight $\mathrm{wg}_{\mathrm{Ypr}_{\mathfrak{c}}}$ is defined.
Let $t\in\mathrm{DD}(t_x/2)$ and $\mathfrak{c}\in\mathrm{Cel}$, and write
\begin{gather*}
\mathrm{Leg}:=\mathrm{Leg}_{\mathscr{T}(t)},\qquad
\mathrm{Leg}_{\mathfrak{c}}:=\partial_{\mathfrak{c}}\mathrm{Leg}_{\mathscr{T}(t)}=\left(\partial_{\hat{c}_{\mathfrak{c}}}\mathrm{Leg}_A\right)\Big|_{A=\mathscr{T}(t)}\quad\text{(by (f3))},\\
\left|\mathrm{Leg}_{\mathfrak{c}}\right|(\mathrm{up};y'):=\sum_{\beta'}\left|\mathrm{Leg}_{\mathfrak{c}}\left(\mathrm{up};(y',\beta',\omega)\right)\right|,
\end{gather*}
where $\omega$ is the bar index of $\mathrm{up}$.
By (f2), $\mathscr{T}(t)$ lies in $\mathrm{Gwk}$.
By (5.2), (ii) and the product rule,
\[
\partial_{\mathfrak{c}}\mathrm{Vj}_+(t)=-\left\langle\bar{\Psi},\partial_{\mathfrak{c}}\mathrm{Kq}(t)\Psi\right\rangle+\mathrm{Fv}^A+\sum_j\mathrm{Fv}^{B,j},
\]
where $\mathrm{Fv}^A$ has the kernels
\[
\int\mathrm{d}\,\mathrm{up}\;\left(\partial_{\mathfrak{c}}\mathrm{Ea}_{nm}\right)(X;t;\mathrm{up}_1,\dots,\mathrm{up}_p)\cdot\prod_i\mathrm{Leg}(\mathrm{up}_i;u_i),
\]
and where, for $j\ge1$, $\mathrm{Fv}^{B,j}$ has the kernels, defined for $p=n+m\ge\max(2,j)$,
\[
\mathrm{Kv}^{(j)}_{nm}(X;u_1,\dots,u_p):=\int\mathrm{d}\,\mathrm{up}\;\mathrm{Ea}_{nm}(X;t;\mathrm{up}_1,\dots,\mathrm{up}_p)\cdot\mathrm{Leg}_{\mathfrak{c}}(\mathrm{up}_j;u_j)\cdot\prod_{i\ne j}\mathrm{Leg}(\mathrm{up}_i;u_i).
\]

\emph{The quadratic term.}
By (f3), $\partial_{\mathfrak{c}}\mathrm{Kq}(t)=\left(\partial_{\hat{c}_{\mathfrak{c}}}D_*\right)\left(\mathscr{T}(t)\right)$.
For $z\in\mathrm{Ypr}_{\mathfrak{c}}$ one has $\max\left(d(y,\mathrm{Ypr}_{\mathfrak{c}}),d(y',\mathrm{Ypr}_{\mathfrak{c}})\right)\le d(y,z)+d(z,y')$, and hence
\[
\mathrm{wg}_{\mathrm{Ypr}_{\mathfrak{c}}}(y,y')\le\exp\left(\mathrm{vt}\cdot\left(d(y,z)+d(z,y')\right)\right).
\]
Proposition~33.2\,(D1), applied with $B=\hat{c}_{\mathfrak{c}}$ (so that $\epsilon_B=\epsilon_{\mathfrak{c}}$ and $\mathrm{Ypr}_B=\mathrm{Ypr}_{\mathfrak{c}}$), gives
\[
\left|\partial_{\mathfrak{c}}\mathrm{Kq}(t)\right|(y,y')\cdot\mathrm{wg}_{\mathrm{Ypr}_{\mathfrak{c}}}(y,y')\le226\,\left(C^F\right)^2\cdot\epsilon_{\mathfrak{c}}\cdot\sum_{z\in\mathrm{Ypr}_{\mathfrak{c}}}\exp\left(-\left(\gamma^F-\mathrm{vt}\right)\cdot d(y,z)\right)\cdot\exp\left(-\left(\gamma^F-\mathrm{vt}\right)\cdot d(z,y')\right).
\]
Since $\gamma^F-\mathrm{vt}\ge\gamma^F/2$, (n5) bounds the sum over $y$ and the sum over $y'$ by $\mathrm{sF}$ each.
As $226\,\left(C^F\right)^2\cdot\mathrm{sF}^2=C_K$, (n8) and (i) give
\[
\mathrm{Nw}_G\left(-\left\langle\bar{\Psi},\partial_{\mathfrak{c}}\mathrm{Kq}(t)\Psi\right\rangle;\mathrm{wg}_{\mathrm{Ypr}_{\mathfrak{c}}}\right)\le8\,h_G^2\cdot C_K\cdot\epsilon_{\mathfrak{c}}\cdot\mathrm{pp}.
\]
The sum over $(z,y')$ at fixed $y$, and the sum over $(z,y)$ at fixed $y'$, are at most $\mathrm{sF}^2$; therefore
\[
\mathrm{rw}_G\left(-\left\langle\bar{\Psi},\partial_{\mathfrak{c}}\mathrm{Kq}(t)\Psi\right\rangle;\mathrm{wg}_{\mathrm{Ypr}_{\mathfrak{c}}}\right)\le4\,h_G^2\cdot C_K\cdot\epsilon_{\mathfrak{c}}.
\]

\emph{(A) The derivative acts on the kernel.}
We apply (vii) with $\mathrm{Aq}=\mathscr{T}(t)$, with $\mathrm{Dsub}$ the set of those $X\in D_*$ that contain a point of $\mathrm{Ypr}_{\mathfrak{c}}$, and with
\[
\mathrm{Gk}_{nm}(X;\cdot\,)=\partial_{\mathfrak{c}}\mathrm{Ea}_{nm}(X;t;\cdot\,),\qquad\mathcal{G}=q_x\cdot w^{\mathrm{E}}_*,\qquad\mathrm{wg}=\mathrm{wg}_{\mathrm{Ypr}_{\mathfrak{c}}}.
\]
By clause (c), $\partial_{\mathfrak{c}}\mathrm{Ea}_{nm}(X;t)=0$ for $X\notin\mathrm{Dsub}$.
The function $\partial_{\mathfrak{c}}\mathrm{Ea}_{nm}(X;t;\cdot\,)$ vanishes at every argument at which $\mathrm{Ea}_{nm}(X;t';\cdot\,)$ vanishes for all $t'$; so the condition on the positions holds by (iii).
The hypothesis on $\mathcal{G}$ is clause (b) with $\mathrm{aw}=5\,\mathrm{hL}\cdot h_G$.
To verify (1.3), let $z\in\mathrm{Ypr}_{\mathfrak{c}}\cap X$; then $z\in X^+$.
With $\mathrm{zn}_i\in\mathrm{pos}(\mathrm{up}_i)$ nearest to $u_i$,
\[
d\left(u_i,\mathrm{Ypr}_{\mathfrak{c}}\right)\le d(u_i,\mathrm{zn}_i)+d(\mathrm{zn}_i,z)\le d\left(u_i,\mathrm{pos}(\mathrm{up}_i)\right)+\operatorname{diam}\left(X^+\right),
\]
and by (5.4)
\[
\mathrm{wg}_{\mathrm{Ypr}_{\mathfrak{c}}}(u_1,\dots,u_p)\le\exp\left(d_M(X)+3\right)\cdot\prod_i\exp\left(\mathrm{vt}\cdot d\left(u_i,\mathrm{pos}(\mathrm{up}_i)\right)\right).
\]
The number of $M$-cubes that contain a point of $\mathrm{Ypr}_{\mathfrak{c}}$ is at most $\#|\mathrm{Ypr}_{\mathfrak{c}}|\le\mathrm{pp}$, every $X\in\mathrm{Dsub}$ contains one of these cubes, and $\kappa-1\ge\kappa_0$.
Proposition~35.4\,(SF4)\,(v) therefore yields
\[
\sum_{X\in\mathrm{Dsub}}\exp\left(-(\kappa-1)\cdot d_M(X)\right)\le\mathrm{pp}\cdot K_0,
\]
and (vii) gives
\[
\mathrm{Nw}_G\left(\mathrm{Fv}^A;\mathrm{wg}_{\mathrm{Ypr}_{\mathfrak{c}}}\right)\le\mathrm{e}^3\cdot K_0\cdot\mathrm{pp}\cdot q_x\cdot w^{\mathrm{E}}_*,\qquad
\mathrm{rw}_G\left(\mathrm{Fv}^A;\mathrm{wg}_{\mathrm{Ypr}_{\mathfrak{c}}}\right)\le650\,K_0\cdot q_x\cdot w^{\mathrm{E}}_*.
\]

\emph{(B) The derivative acts on the leg $j$.}
By Proposition~33.2\,(L2) at $A=\mathscr{T}(t)$ and $B=\hat{c}_{\mathfrak{c}}$,
\begin{equation}
\left|\mathrm{Leg}_{\mathfrak{c}}\right|(\mathrm{up};y')\le C_H\cdot\epsilon_{\mathfrak{c}}\cdot\sum_{w\in\mathrm{pos}(\mathrm{up})}\sum_{z\in\mathrm{Ypr}_{\mathfrak{c}}^+}\mathrm{Edy}(w,z)\cdot\mathrm{Edy}(z,y').\tag{1.4}
\end{equation}
For $u=(y',\beta',\omega')$, the number $\mathrm{Leg}_{\mathfrak{c}}(\mathrm{up};u)$ vanishes if $\omega'\ne\omega$; hence, for every non-negative function $\mathrm{fq}$ on $T^0_*$,
\[
\sum_u\left|\mathrm{Leg}_{\mathfrak{c}}(\mathrm{up};u)\right|\cdot\mathrm{fq}(y')=\sum_{y'}\left|\mathrm{Leg}_{\mathfrak{c}}\right|(\mathrm{up};y')\cdot\mathrm{fq}(y').
\]
Let $\mathrm{Ea}_{nm}(X;t;\mathrm{up}_1,\dots,\mathrm{up}_p)\ne0$, $w\in\mathrm{pos}(\mathrm{up}_j)$ and $z\in\mathrm{Ypr}_{\mathfrak{c}}^+$.
Then $d(z,\mathrm{Ypr}_{\mathfrak{c}})\le M+1$, so that $d(u_j,\mathrm{Ypr}_{\mathfrak{c}})\le d(u_j,z)+M+1$.
For $i\ne j$, with $\mathrm{zn}_i\in\mathrm{pos}(\mathrm{up}_i)$ nearest to $u_i$,
\begin{align*}
d\left(u_i,\mathrm{Ypr}_{\mathfrak{c}}\right)&\le d(u_i,\mathrm{zn}_i)+d(\mathrm{zn}_i,w)+d(w,z)+M+1\\
&\le d\left(u_i,\mathrm{pos}(\mathrm{up}_i)\right)+\operatorname{diam}\left(X^+\right)+d(w,z)+M+1,
\end{align*}
because $\mathrm{zn}_i$ and $w$ lie in $X^+$ by (iii).
Since $\mathrm{vt}\cdot(M+1)\le2$, (5.4) yields
\begin{equation}
\mathrm{wg}_{\mathrm{Ypr}_{\mathfrak{c}}}(u_1,\dots,u_p)\le\exp\left(d_M(X)+5\right)\cdot\exp\left(\mathrm{vt}\cdot\left(d(w,z)+d(z,u_j)\right)\right)\cdot\prod_{i\ne j}\exp\left(\mathrm{vt}\cdot d\left(u_i,\mathrm{pos}(\mathrm{up}_i)\right)\right).\tag{1.5}
\end{equation}

\emph{The bound for $\mathrm{Nw}_G$.}
We insert (1.4) for the leg $j$ and use (1.5) in each term $(w,z)$.
Let $y_j$ denote the site of $u_j$.
By (n5),
\[
\sum_{y_j}\mathrm{Edy}(z,y_j)\cdot\exp\left(\mathrm{vt}\cdot d(z,y_j)\right)\le\mathrm{sF}.
\]
The set $\mathrm{pos}(\mathrm{up}_j)$ has at most two points, and they lie in $X^+$; hence
\[
\sum_{w\in\mathrm{pos}(\mathrm{up}_j)}\mathrm{Edy}(w,z)\cdot\exp\left(\mathrm{vt}\cdot d(w,z)\right)\le2\exp\left(-\frac{\gamma^F}{2}\cdot d\left(z,X^+\right)\right).
\]
Each leg $i\ne j$ contributes a factor $\mathrm{hL}$, by Proposition~33.2\,(L1) with $\gamma=\mathrm{vt}$.
So
\[
\sum_{u_1,\dots,u_p}\left|\mathrm{Kv}^{(j)}_{nm}(X;\cdot\,)\right|\cdot\mathrm{wg}_{\mathrm{Ypr}_{\mathfrak{c}}}\le2\exp\left(d_M(X)+5\right)\cdot\mathrm{hL}^{p-1}\cdot C_H\cdot\epsilon_{\mathfrak{c}}\cdot\mathrm{sF}\cdot\sum_{z\in\mathrm{Ypr}_{\mathfrak{c}}^+}\exp\left(-\frac{\gamma^F}{2}\cdot d\left(z,X^+\right)\right)\cdot\left\|\mathrm{Ea}_{nm}(X;t)\right\|.
\]
There are $p$ values of $j$.
We use (3.3), (ii), the inequalities $p^2\le p\cdot(p+1)^2\le5^p$ and $\mathrm{hL}\ge1$, clause (a) with $\mathrm{aw}=5\,\mathrm{hL}\cdot h_G$, (5.5) with $\gamma=\gamma^F/2$, and (n8):
\begin{align*}
&\mathrm{Nw}_G\Bigl(\sum_j\mathrm{Fv}^{B,j};\mathrm{wg}_{\mathrm{Ypr}_{\mathfrak{c}}}\Bigr)\\
&\quad\le\sum_X\sum_{n+m\ge2}\frac{p^2\cdot\left(\mathrm{hL}\cdot h_G\right)^p}{n!\,m!}\cdot\left\|\mathrm{Ea}_{nm}(X;t)\right\|\cdot2\exp\left(d_M(X)+5\right)\cdot C_H\cdot\epsilon_{\mathfrak{c}}\cdot\mathrm{sF}\cdot\sum_{z\in\mathrm{Ypr}_{\mathfrak{c}}^+}\exp\left(-\frac{\gamma^F}{2}\cdot d\left(z,X^+\right)\right)\\
&\quad\le2\,\mathrm{e}^5\cdot C_H\cdot\epsilon_{\mathfrak{c}}\cdot\mathrm{sF}\cdot w^{\mathrm{E}}_*\cdot\sum_{z\in\mathrm{Ypr}_{\mathfrak{c}}^+}\sum_X\exp\left(-\frac{\gamma^F}{2}\cdot d\left(z,X^+\right)\right)\cdot\exp\left(-(\kappa-1)\cdot d_M(X)\right)\\
&\quad\le2\,\mathrm{e}^5\cdot32\,C_H\cdot\mathrm{sF}\cdot K_0\cdot w^{\mathrm{E}}_*\cdot\mathrm{pp}\cdot\epsilon_{\mathfrak{c}}\le9500\,C_H\cdot\mathrm{sF}\cdot K_0\cdot w^{\mathrm{E}}_*\cdot\mathrm{pp}\cdot\epsilon_{\mathfrak{c}}.
\end{align*}

\emph{The bound for $\mathrm{rw}_G$.}
Fix a slot $\mathrm{sl}$ and the index $u=(y',\beta',\omega')$ at this slot.

\emph{Case $j=\mathrm{sl}$.}
By (1.4) and (1.5), the factor of the leg $j$ is at most
\[
C_H\cdot\epsilon_{\mathfrak{c}}\cdot\sum_{w\in\mathrm{pos}(\mathrm{up}_j)}\sum_{z\in\mathrm{Ypr}_{\mathfrak{c}}^+}\exp\left(-\frac{7\gamma^F}{8}\cdot\left(d(w,z)+d(z,y')\right)\right).
\]
Since $\frac{7}{8}=\frac{3}{8}+\frac{1}{2}$ and $d(w,z)+d(z,y')\ge d(w,y')$,
\begin{align*}
&\sum_z\exp\left(-\frac{7\gamma^F}{8}\cdot\left(d(w,z)+d(z,y')\right)\right)\\
&\quad\le\exp\left(-\frac{3\gamma^F}{8}\cdot d(w,y')\right)\cdot\sum_z\exp\left(-\frac{\gamma^F}{2}\cdot d(w,z)\right)\\
&\quad\le\mathrm{sF}\cdot\exp\left(-\frac{3\gamma^F}{8}\cdot d(w,y')\right).
\end{align*}
The points $w$ are at most two and lie in $X^+$; so the sum over $w$ is at most $2\,\mathrm{sF}\cdot\exp\left(-\frac{\gamma^F}{4}\cdot d\left(u,X^+\right)\right)$.
The other $p-1$ legs, summed over their indices, contribute $\mathrm{hL}$ each.

\emph{Case $j\ne\mathrm{sl}$.}
By the last estimate in the proof of (vi), the leg $\mathrm{sl}$, together with its factor $\exp\left(\mathrm{vt}\cdot d\left(u,\mathrm{pos}(\mathrm{up}_{\mathrm{sl}})\right)\right)$ from (1.5), contributes at most $\mathrm{hL}\cdot\exp\left(-\frac{\gamma^F}{4}\cdot d\left(u,X^+\right)\right)$.
The leg $j$, summed over $u_j$, $z$ and $w$, contributes at most $C_H\cdot\epsilon_{\mathfrak{c}}\cdot2\cdot\mathrm{sF}\cdot\mathrm{sF}$.
The remaining $p-2$ legs contribute $\mathrm{hL}$ each.

In both cases, since $\mathrm{sF}\ge1$,
\begin{align*}
&\sum_{u_i,\,i\ne\mathrm{sl}}\left|\mathrm{Kv}^{(j)}_{nm}(X;\dots,u\text{ at }\mathrm{sl},\dots)\right|\cdot\mathrm{wg}_{\mathrm{Ypr}_{\mathfrak{c}}}\\
&\quad\le2\exp\left(d_M(X)+5\right)\cdot\mathrm{hL}^{p-1}\cdot C_H\cdot\epsilon_{\mathfrak{c}}\cdot\mathrm{sF}^2\\
&\qquad\cdot\exp\left(-\frac{\gamma^F}{4}\cdot d\left(u,X^+\right)\right)\cdot\left\|\mathrm{Ea}_{nm}(X;t)\right\|.
\end{align*}
By clause (a) with $\mathrm{aw}=5\,\mathrm{hL}\cdot h_G$,
\[
\sum_{n+m=p}\frac{\left(\mathrm{hL}\cdot h_G\right)^p}{n!\,m!}\cdot\left\|\mathrm{Ea}_{nm}(X;t)\right\|\le5^{-p}\cdot w^{\mathrm{E}}_*\cdot\exp\left(-\kappa\cdot d_M(X)\right).
\]
Taking into account the $p$ values of $j$, and using (3.3), (ii), $\mathrm{hL}\ge1$ and (5.5) with $\gamma=\gamma^F/4$, one obtains
\begin{align*}
&\mathrm{rw}_G\Bigl(\sum_j\mathrm{Fv}^{B,j};\mathrm{wg}_{\mathrm{Ypr}_{\mathfrak{c}}}\Bigr)\\
&\quad\le\sum_{p\ge2}p^2\cdot5^{-p}\cdot2\,\mathrm{e}^5\cdot C_H\cdot\epsilon_{\mathfrak{c}}\cdot\mathrm{sF}^2\cdot w^{\mathrm{E}}_*\cdot\sup_u\sum_X\exp\left(-\frac{\gamma^F}{4}\cdot d\left(u,X^+\right)\right)\cdot\exp\left(-(\kappa-1)\cdot d_M(X)\right)\\
&\quad\le\frac{43}{160}\cdot2\,\mathrm{e}^5\cdot32\,C_H\cdot\mathrm{sF}^2\cdot K_0\cdot w^{\mathrm{E}}_*\cdot\epsilon_{\mathfrak{c}}\le9500\,C_H\cdot\mathrm{sF}^2\cdot K_0\cdot w^{\mathrm{E}}_*\cdot\epsilon_{\mathfrak{c}}.
\end{align*}

\emph{Conclusion of Step 2.}
One has $\epsilon_{\mathfrak{c}}\le q_x$ by (n7), $w^{\mathrm{E}}_*\le1/8$ by (n2), $h_G\ge1$, and
\[
\frac{\mathrm{e}^3}{8}\le3,\qquad\frac{650}{8}\le82,\qquad\frac{9500}{8}\le1200.
\]
Moreover, since $M\ge2$,
\[
C_K=226\,\left(C^F\right)^2\cdot\mathrm{sF}^2\le203\,M\cdot\left(C^F\right)^2\cdot\mathrm{sF}^2=C_H\cdot\mathrm{sF}^2.
\]
By the subadditivity of $\mathrm{Nw}_G$ and $\mathrm{rw}_G$ (Q2) it follows that
\begin{align*}
\mathrm{Nw}_G\left(\partial_{\mathfrak{c}}\mathrm{Vj}_+(t);\mathrm{wg}_{\mathrm{Ypr}_{\mathfrak{c}}}\right)
&\le q_x\cdot\mathrm{pp}\cdot h_G^2\cdot\left(8\,C_K+3\,K_0+1200\,C_H\cdot\mathrm{sF}\cdot K_0\right)\\
&\le q_x\cdot\mathrm{pp}\cdot h_G^2\cdot1211\,C_H\cdot\mathrm{sF}^2\cdot K_0\le q_x\cdot\mathrm{CVj}\cdot h_G^2\cdot\mathrm{pp},\\
\mathrm{rw}_G\left(\partial_{\mathfrak{c}}\mathrm{Vj}_+(t);\mathrm{wg}_{\mathrm{Ypr}_{\mathfrak{c}}}\right)
&\le q_x\cdot h_G^2\cdot\left(4\,C_K+82\,K_0+1200\,C_H\cdot\mathrm{sF}^2\cdot K_0\right)\\
&\le q_x\cdot h_G^2\cdot1286\,C_H\cdot\mathrm{sF}^2\cdot K_0\le q_x\cdot\mathrm{CVj}\cdot h_G^2,
\end{align*}
because
\[
\mathrm{CVj}=10^3\,C_H^2\cdot\mathrm{sF}^4\cdot K_0\ge2000\,C_H\cdot\mathrm{sF}^2\cdot K_0\qquad\left(C_H\ge2,\ \mathrm{sF}\ge1,\ K_0\ge1\right).
\]

\paragraph{1.9. Step 3: (N3).}
Let $\hat{\mathfrak{a}}\ne0$, $t\in\mathrm{DD}(t_x/2)$ and $\mathfrak{c}\ne\mathfrak{c}'$.
Put $B:=\hat{c}_{\mathfrak{c}}$ and $B':=\hat{c}_{\mathfrak{c}'}$.
By (n9),
\[
d\left(\mathrm{Ypr}_B,\mathrm{Ypr}_{B'}\right)=d\left(\mathrm{Ypr}_{\mathfrak{c}},\mathrm{Ypr}_{\mathfrak{c}'}\right)\ge M+2,
\]
so that Proposition~33.2\,(L3) and (D2) apply at $A=\mathscr{T}(t)$.
Define
\[
\mathrm{dsp}_0:=d\left(\mathrm{Ypr}_{\mathfrak{c}}^+,\mathrm{Ypr}_{\mathfrak{c}'}^+\right)\ge d_{\mathrm{sep}}-2M-2,
\]
and, for $X\in D_*$,
\[
\mathrm{dsp}_1:=d\left(X^+,\mathrm{Ypr}_{\mathfrak{c}}^+\right),\qquad\mathrm{dsp}_2:=d\left(X^+,\mathrm{Ypr}_{\mathfrak{c}'}^+\right).
\]
The lower bound for $\mathrm{dsp}_0$ holds because every point of $\mathrm{Ypr}_{\mathfrak{c}}^+$ has distance at most $M+1$ from $\mathrm{Ypr}_{\mathfrak{c}}$; all the sets involved are non-empty.

\emph{The quadratic term.}
By (f3), $\partial_{\mathfrak{c}}\partial_{\mathfrak{c}'}\mathrm{Kq}(t)=\left(\partial_B\partial_{B'}D_*\right)\left(\mathscr{T}(t)\right)$.
Proposition~33.2\,(D2), (n5) in the form $\sum_y\mathrm{Edy}(y,z)\le\mathrm{sF}$, and (i) with the weight $1$ give
\[
\mathrm{Nw}_G\left(-\left\langle\bar{\Psi},\partial_{\mathfrak{c}}\partial_{\mathfrak{c}'}\mathrm{Kq}(t)\Psi\right\rangle;1\right)\le8\,h_G^2\cdot C_H^2\cdot\epsilon_{\mathfrak{c}}\cdot\epsilon_{\mathfrak{c}'}\cdot\sum_{z\in\mathrm{Ypr}_{\mathfrak{c}}}\sum_{z'\in\mathrm{Ypr}_{\mathfrak{c}'}}2\,\mathrm{sF}^2\cdot\mathrm{Edy}(z,z').
\]
For $z\in\mathrm{Ypr}_{\mathfrak{c}}$, (n6) and (n9) give
\[
\sum_{z'\in\mathrm{Ypr}_{\mathfrak{c}'}}\mathrm{Edy}(z,z')\le\mathrm{sF}\cdot\exp\left(-\frac{\gamma^F}{2}\cdot d_{\mathrm{sep}}\right).
\]
Together with (n7) and (n8) this yields
\[
\mathrm{Nw}_G\left(-\left\langle\bar{\Psi},\partial_{\mathfrak{c}}\partial_{\mathfrak{c}'}\mathrm{Kq}(t)\Psi\right\rangle;1\right)\le16\,h_G^2\cdot C_H^2\cdot\mathrm{sF}^3\cdot\mathrm{pp}\cdot q_x^2\cdot\exp\left(-\frac{\gamma^F}{2}\cdot d_{\mathrm{sep}}\right).
\]

\emph{The activities, the degree $p=0$ included.}
Fix $X$ and $(n,m)$, let $p=n+m\ge0$, and write the integrand of $\mathrm{Kv}_{nm}(X;t;\cdot\,)$ as the product of the $p+1$ factors
\[
\mathrm{Fac}_0:=\mathrm{Ea}_{nm}(X;t;\mathrm{up}_1,\dots,\mathrm{up}_p),\qquad\mathrm{Fac}_i:=\mathrm{Leg}(\mathrm{up}_i;u_i)\quad(1\le i\le p).
\]
By the product rule, $\partial_{\mathfrak{c}}\partial_{\mathfrak{c}'}\prod_i\mathrm{Fac}_i$ is a sum of $(p+1)^2$ products: $\partial_{\mathfrak{c}}$ acts on one of the $p+1$ factors, and $\partial_{\mathfrak{c}'}$ acts on one of the $p+1$ factors.
Let $\mathrm{Leg}_{\mathfrak{c}}$ and $\mathrm{Leg}_{\mathfrak{c}'}$ be the first derivatives, defined as in 1.8, and put
\[
\mathrm{Leg}_{\mathfrak{c}\mathfrak{c}'}:=\partial_{\mathfrak{c}}\partial_{\mathfrak{c}'}\mathrm{Leg}_{\mathscr{T}(t)}=\left(\partial_B\partial_{B'}\mathrm{Leg}_A\right)\Big|_{A=\mathscr{T}(t)}.
\]
At every argument that is not of the form described in (iii), $\mathrm{Ea}_{nm}(X;t';\cdot\,)$ vanishes for all $t'$.
Hence $\mathrm{Fac}_0$ and all its derivatives in $t$ vanish unless $\mathrm{pos}(\mathrm{up}_i)\subseteq X^+$ for every $i$.
For a leg index $\mathrm{up}$ with $\mathrm{pos}(\mathrm{up})\subseteq X^+$, Proposition~33.2\,(L1)--(L3), (n5) and (n6) give
\begin{align*}
\sum_u\left|\mathrm{Leg}(\mathrm{up};u)\right|&\le\mathrm{hL},\\
\sum_u\left|\mathrm{Leg}_{\mathfrak{c}}(\mathrm{up};u)\right|&\le C_H\cdot\epsilon_{\mathfrak{c}}\cdot\sum_{w\in\mathrm{pos}(\mathrm{up})}\sum_{z\in\mathrm{Ypr}_{\mathfrak{c}}^+}\mathrm{Edy}(w,z)\cdot\mathrm{sF}\le2\,C_H\cdot\epsilon_{\mathfrak{c}}\cdot\mathrm{sF}^2\cdot\exp\left(-\frac{\gamma^F}{2}\cdot\mathrm{dsp}_1\right),\\
\sum_u\left|\mathrm{Leg}_{\mathfrak{c}'}(\mathrm{up};u)\right|&\le2\,C_H\cdot\epsilon_{\mathfrak{c}'}\cdot\mathrm{sF}^2\cdot\exp\left(-\frac{\gamma^F}{2}\cdot\mathrm{dsp}_2\right),\\
\sum_u\left|\mathrm{Leg}_{\mathfrak{c}\mathfrak{c}'}(\mathrm{up};u)\right|&\le C_H^2\cdot\epsilon_{\mathfrak{c}}\cdot\epsilon_{\mathfrak{c}'}\cdot\mathrm{sF}\cdot\sum_{w\in\mathrm{pos}(\mathrm{up})}\sum_{z\in\mathrm{Ypr}_{\mathfrak{c}}^+}\sum_{z'\in\mathrm{Ypr}_{\mathfrak{c}'}^+}\left[\mathrm{Edy}(w,z)\cdot\mathrm{Edy}(z,z')+\mathrm{Edy}(w,z')\cdot\mathrm{Edy}(z',z)\right]\\
&\le2\,C_H^2\cdot\epsilon_{\mathfrak{c}}\cdot\epsilon_{\mathfrak{c}'}\cdot\mathrm{sF}^3\cdot\exp\left(-\frac{\gamma^F}{2}\cdot\mathrm{dsp}_0\right)\cdot\left[\exp\left(-\frac{\gamma^F}{2}\cdot\mathrm{dsp}_1\right)+\exp\left(-\frac{\gamma^F}{2}\cdot\mathrm{dsp}_2\right)\right].
\end{align*}
The last inequality rests on (n6), used twice.
For $z\in\mathrm{Ypr}_{\mathfrak{c}}^+$ one has $\sum_{z'\in\mathrm{Ypr}_{\mathfrak{c}'}^+}\mathrm{Edy}(z,z')\le\mathrm{sF}\cdot\exp\left(-\frac{\gamma^F}{2}\cdot\mathrm{dsp}_0\right)$, and then $\sum_{z\in\mathrm{Ypr}_{\mathfrak{c}}^+}\mathrm{Edy}(w,z)\le\mathrm{sF}\cdot\exp\left(-\frac{\gamma^F}{2}\cdot\mathrm{dsp}_1\right)$.
Symmetrically, for $z'\in\mathrm{Ypr}_{\mathfrak{c}'}^+$ one has $\sum_{z\in\mathrm{Ypr}_{\mathfrak{c}}^+}\mathrm{Edy}(z',z)\le\mathrm{sF}\cdot\exp\left(-\frac{\gamma^F}{2}\cdot\mathrm{dsp}_0\right)$, and then $\sum_{z'\in\mathrm{Ypr}_{\mathfrak{c}'}^+}\mathrm{Edy}(w,z')\le\mathrm{sF}\cdot\exp\left(-\frac{\gamma^F}{2}\cdot\mathrm{dsp}_2\right)$.
Finally, there are at most two points $w$, and they lie in $X^+$.

By clause (c), $\partial_{\mathfrak{c}}\mathrm{Fac}_0=0$ unless $X$ contains a point of $\mathrm{Ypr}_{\mathfrak{c}}$; such a point lies in $X^+$ and in $\mathrm{Ypr}_{\mathfrak{c}}^+$, and then $\mathrm{dsp}_1=0$.
By clause (c) for the cell $\mathfrak{c}'$, $\partial_{\mathfrak{c}'}\mathrm{Fac}_0=0$ unless $X$ contains a point of $\mathrm{Ypr}_{\mathfrak{c}'}$, and then $\mathrm{dsp}_2=0$.
Clause (c) holds for all $t\in\mathrm{DDo}(t_x^+)$.
Therefore $\partial_{\mathfrak{c}}\partial_{\mathfrak{c}'}\mathrm{Fac}_0$, being a derivative both of $\partial_{\mathfrak{c}}\mathrm{Fac}_0$ and of $\partial_{\mathfrak{c}'}\mathrm{Fac}_0$, vanishes unless $\mathrm{dsp}_1=\mathrm{dsp}_2=0$.

Let $\mathrm{Ty}$ denote any one of the four kernels
\[
\mathrm{Ea}_{nm}(X;t),\qquad q_x^{-1}\cdot\partial_{\mathfrak{c}}\mathrm{Ea}_{nm}(X;t),\qquad q_x^{-1}\cdot\partial_{\mathfrak{c}'}\mathrm{Ea}_{nm}(X;t),\qquad q_x^{-2}\cdot\partial_{\mathfrak{c}}\partial_{\mathfrak{c}'}\mathrm{Ea}_{nm}(X;t),
\]
which are the four types of the factor $\mathrm{Fac}_0$.
Since $\epsilon_{\mathfrak{c}},\epsilon_{\mathfrak{c}'}\le q_x$, $\mathrm{hL}\ge1$, $C_H\ge1$ and $\mathrm{sF}\ge1$, each of the $(p+1)^2$ products, integrated over $\mathrm{up}$ and summed over $u_1,\dots,u_p$, is at most
\[
q_x^2\cdot4\,C_H^2\cdot\mathrm{sF}^4\cdot\mathrm{hL}^p\cdot\left\|\mathrm{Ty}\right\|\cdot\mathrm{De}(X),
\]
where $\mathrm{Ty}$ is the type of the factor $\mathrm{Fac}_0$ of the product, and
\[
\mathrm{De}(X):=\exp\left(-\frac{\gamma^F}{2}\cdot\left(\mathrm{dsp}_1+\mathrm{dsp}_2\right)\right)+\exp\left(-\frac{\gamma^F}{2}\cdot\mathrm{dsp}_0\right)\cdot\left[\exp\left(-\frac{\gamma^F}{2}\cdot\mathrm{dsp}_1\right)+\exp\left(-\frac{\gamma^F}{2}\cdot\mathrm{dsp}_2\right)\right].
\]
Indeed, according to the place of the two derivatives, the product is bounded as follows.
\begin{itemize}
\item[--] The derivatives act on two different legs:
$\mathrm{hL}^{p-2}\cdot2\,C_H\cdot\epsilon_{\mathfrak{c}}\cdot\mathrm{sF}^2\cdot2\,C_H\cdot\epsilon_{\mathfrak{c}'}\cdot\mathrm{sF}^2\cdot\exp\left(-\frac{\gamma^F}{2}\cdot\left(\mathrm{dsp}_1+\mathrm{dsp}_2\right)\right)\cdot\left\|\mathrm{Ea}_{nm}\right\|$.
\item[--] Both derivatives act on the same leg:
$\mathrm{hL}^{p-1}$ times the bound for $\mathrm{Leg}_{\mathfrak{c}\mathfrak{c}'}$ times $\left\|\mathrm{Ea}_{nm}\right\|$.
\item[--] $\partial_{\mathfrak{c}}$ acts on $\mathrm{Fac}_0$ and $\partial_{\mathfrak{c}'}$ on a leg:
$\left\|\partial_{\mathfrak{c}}\mathrm{Ea}_{nm}\right\|\cdot\mathrm{hL}^{p-1}\cdot2\,C_H\cdot\epsilon_{\mathfrak{c}'}\cdot\mathrm{sF}^2\cdot\exp\left(-\frac{\gamma^F}{2}\cdot\mathrm{dsp}_2\right)$, with $\mathrm{dsp}_1=0$.
\item[--] $\partial_{\mathfrak{c}'}$ acts on $\mathrm{Fac}_0$ and $\partial_{\mathfrak{c}}$ on a leg:
$\left\|\partial_{\mathfrak{c}'}\mathrm{Ea}_{nm}\right\|\cdot\mathrm{hL}^{p-1}\cdot2\,C_H\cdot\epsilon_{\mathfrak{c}}\cdot\mathrm{sF}^2\cdot\exp\left(-\frac{\gamma^F}{2}\cdot\mathrm{dsp}_1\right)$, with $\mathrm{dsp}_2=0$.
\item[--] Both derivatives act on $\mathrm{Fac}_0$:
$\left\|\partial_{\mathfrak{c}}\partial_{\mathfrak{c}'}\mathrm{Ea}_{nm}\right\|\cdot\mathrm{hL}^p$, with $\mathrm{dsp}_1=\mathrm{dsp}_2=0$.
\end{itemize}

There are points of $X^+$ at the distance $\mathrm{dsp}_1$ from $\mathrm{Ypr}_{\mathfrak{c}}^+$ and at the distance $\mathrm{dsp}_2$ from $\mathrm{Ypr}_{\mathfrak{c}'}^+$; hence
\[
\mathrm{dsp}_0\le\mathrm{dsp}_1+\operatorname{diam}\left(X^+\right)+\mathrm{dsp}_2\le\mathrm{dsp}_1+\mathrm{dsp}_2+M\cdot d_M(X)+2M+2.
\]
As $\mathrm{vt}\le\gamma^F/8$ and $\mathrm{vt}\cdot M\le1$, it follows that
\[
\mathrm{vt}\cdot\left(\mathrm{dsp}_0-2M-2\right)\le\frac{\gamma^F}{8}\cdot\left(\mathrm{dsp}_1+\mathrm{dsp}_2\right)+d_M(X),
\]
and moreover $\frac{\gamma^F}{8}\cdot\left(\mathrm{dsp}_1+\mathrm{dsp}_2\right)+\frac{\gamma^F}{4}\cdot\mathrm{dsp}_1\le\frac{\gamma^F}{2}\cdot\left(\mathrm{dsp}_1+\mathrm{dsp}_2\right)$.
Hence
\begin{align*}
&\exp\left(-\frac{\gamma^F}{2}\cdot\left(\mathrm{dsp}_1+\mathrm{dsp}_2\right)\right)\cdot\exp\left(-\kappa\cdot d_M(X)\right)\\
&\quad\le\exp\left(-\mathrm{vt}\cdot\left(\mathrm{dsp}_0-2M-2\right)\right)\cdot\exp\left(-\frac{\gamma^F}{4}\cdot\mathrm{dsp}_1\right)\cdot\exp\left(-(\kappa-1)\cdot d_M(X)\right).
\end{align*}
Since $\mathrm{dsp}_0,\mathrm{dsp}_1,\mathrm{dsp}_2,d_M(X)\ge0$, one also has
\begin{gather*}
\exp\left(-\frac{\gamma^F}{2}\cdot\mathrm{dsp}_0\right)\le\exp\left(-\mathrm{vt}\cdot\mathrm{dsp}_0\right)\le\exp\left(-\mathrm{vt}\cdot\left(\mathrm{dsp}_0-2M-2\right)\right),\\
\exp\left(-\frac{\gamma^F}{2}\cdot\mathrm{dsp}_i\right)\le\exp\left(-\frac{\gamma^F}{4}\cdot\mathrm{dsp}_i\right)\quad(i=1,2),
\end{gather*}
and $\exp\left(-\kappa\cdot d_M(X)\right)\le\exp\left(-(\kappa-1)\cdot d_M(X)\right)$.
Therefore
\begin{equation}
\begin{split}
&\mathrm{De}(X)\cdot\exp\left(-\kappa\cdot d_M(X)\right)\\
&\quad\le\exp\left(-\mathrm{vt}\cdot\left(\mathrm{dsp}_0-2M-2\right)\right)\cdot2\left[\exp\left(-\frac{\gamma^F}{4}\cdot\mathrm{dsp}_1\right)+\exp\left(-\frac{\gamma^F}{4}\cdot\mathrm{dsp}_2\right)\right]\cdot\exp\left(-(\kappa-1)\cdot d_M(X)\right).
\end{split}\tag{1.6}
\end{equation}

For $p\ge2$ the factor in (ii) is $p\cdot h_G^p/(n!\,m!)$, and $p\cdot(p+1)^2\le5^p$.
For $p=0$ there is a single product, and its factor is $1=5^0$; odd $p$ do not occur, and the part of degree $p=0$ is
\[
\partial_{\mathfrak{c}}\partial_{\mathfrak{c}'}\mathrm{Vj}_0(t)=\sum_X\partial_{\mathfrak{c}}\partial_{\mathfrak{c}'}\mathrm{Ea}_{00}(X;t).
\]
For each of the four types, clauses (a) and (b) with $\mathrm{aw}=5\,\mathrm{hL}\cdot h_G$ give
\[
\sum_{n,m}\frac{\left(5\,\mathrm{hL}\cdot h_G\right)^{n+m}}{n!\,m!}\cdot\left\|\mathrm{Ty}\right\|\le w^{\mathrm{E}}_*\cdot\exp\left(-\kappa\cdot d_M(X)\right).
\]
At most $(p+1)^2$ products are of a given type.
Let $\mathrm{Fv}$ again denote the sum over $X$ of the parts of degree at least $2$ of the image of $\mathrm{Ea}(X;t)$.
By (3.3),
\[
\mathrm{Nw}_G\left(\partial_{\mathfrak{c}}\partial_{\mathfrak{c}'}\mathrm{Fv};1\right)+\left|\partial_{\mathfrak{c}}\partial_{\mathfrak{c}'}\mathrm{Vj}_0(t)\right|\le4\cdot q_x^2\cdot4\,C_H^2\cdot\mathrm{sF}^4\cdot w^{\mathrm{E}}_*\cdot\sum_{X\in D_*}\mathrm{De}(X)\cdot\exp\left(-\kappa\cdot d_M(X)\right).
\]
Now $\mathrm{dsp}_1$ is the least of the numbers $d(z,X^+)$ with $z\in\mathrm{Ypr}_{\mathfrak{c}}^+$, and $\mathrm{dsp}_2$ is the least of the numbers $d(z,X^+)$ with $z\in\mathrm{Ypr}_{\mathfrak{c}'}^+$.
Consequently
\[
\exp\left(-\frac{\gamma^F}{4}\cdot\mathrm{dsp}_1\right)\le\sum_{z\in\mathrm{Ypr}_{\mathfrak{c}}^+}\exp\left(-\frac{\gamma^F}{4}\cdot d\left(z,X^+\right)\right),\qquad\exp\left(-\frac{\gamma^F}{4}\cdot\mathrm{dsp}_2\right)\le\sum_{z\in\mathrm{Ypr}_{\mathfrak{c}'}^+}\exp\left(-\frac{\gamma^F}{4}\cdot d\left(z,X^+\right)\right).
\]
By (1.6), (5.5) with $\gamma=\gamma^F/4$, and (n8),
\[
\sum_X\mathrm{De}(X)\cdot\exp\left(-\kappa\cdot d_M(X)\right)\le\exp\left(-\mathrm{vt}\cdot\left(\mathrm{dsp}_0-2M-2\right)\right)\cdot2\cdot2\cdot32\,K_0\cdot\mathrm{pp}.
\]
Since $w^{\mathrm{E}}_*\le1/8$ and $\mathrm{dsp}_0-2M-2\ge d_{\mathrm{sep}}-4M-4$, the right side of the preceding bound for $\mathrm{Nw}_G\left(\partial_{\mathfrak{c}}\partial_{\mathfrak{c}'}\mathrm{Fv};1\right)+\left|\partial_{\mathfrak{c}}\partial_{\mathfrak{c}'}\mathrm{Vj}_0(t)\right|$ is at most
\[
256\,C_H^2\cdot\mathrm{sF}^4\cdot K_0\cdot\mathrm{pp}\cdot q_x^2\cdot\exp\left(-\mathrm{vt}\cdot\left(d_{\mathrm{sep}}-4M-4\right)\right).
\]

\emph{Conclusion of Step 3.}
One has $\exp\left(-\frac{\gamma^F}{2}\cdot d_{\mathrm{sep}}\right)\le\exp\left(-\mathrm{vt}\cdot\left(d_{\mathrm{sep}}-4M-4\right)\right)$, $\mathrm{sF}^3\le\mathrm{sF}^4\cdot K_0$, and $16\,h_G^2+256\le272\,h_G^2\le10^3\,h_G^2$.
Adding the two bounds, one obtains
\begin{align*}
&\mathrm{Nw}_G\left(\partial_{\mathfrak{c}}\partial_{\mathfrak{c}'}\mathrm{Vj}_+(t);1\right)+\left|\partial_{\mathfrak{c}}\partial_{\mathfrak{c}'}\mathrm{Vj}_0(t)\right|\\
&\quad\le\left(16\,h_G^2+256\right)\cdot C_H^2\cdot\mathrm{sF}^4\cdot K_0\cdot\mathrm{pp}\cdot q_x^2\cdot\exp\left(-\mathrm{vt}\cdot\left(d_{\mathrm{sep}}-4M-4\right)\right)\\
&\quad\le q_x^2\cdot\mathrm{CVj}\cdot h_G^2\cdot\mathrm{pp}\cdot\exp\left(-\mathrm{vt}\cdot\left(d_{\mathrm{sep}}-4M-6\right)\right).
\end{align*}

\paragraph{1.10. Uniformity.}
The right sides of (a), (b), (N1) and (N2) are built from $\kappa$, $d_M(X)$, $K_0$, $h_G$, $C_K$, $\mathrm{CVj}$ and $M$, which are constants of the lines (C0)--(C6$'$) or are determined by the polymer, and from the numbers
\[
w^{\mathrm{E}}_*=e_*^{\frac{1}{4}-7\varepsilon},\qquad q_x=\frac{2}{t_x}=2\,\mathrm{Cb}\cdot e_*^{-3\varepsilon},\qquad\bar{r}=4\,h_G^2\cdot C_K\cdot e_*^{2\varepsilon}+650\,K_0\cdot e_*^{\frac{1}{4}-7\varepsilon},\qquad\mathrm{pp}=\left(\frac{D_\beta}{a_*}+2M+4\right)^3.
\]
By Proposition~35.1\,(c), $e_*$ and $a_*$ do not depend on $N\ge N_0(e)$.
None of these numbers depends on $\theta$, $A$, $(A^{g},\mathscr{A}^0)$ or on the placement.
The numbers $C_\beta$, $\mathbf{D}$ (hence also $\mathrm{Cb}=1+h^{\frac{1}{2}}\cdot C_\beta$) and $D_\beta$ do not change if $\beta$ is replaced by $U_{r'}\beta$ with $r'\in\{\pm1\}^3$.
Indeed, $U_rU_{r'}\beta=U_{rr'}\beta$ (Section~I.7.5), the map $r\mapsto rr'$ is a bijection of $\{\pm1\}^3$, and the conditions on $(N',\mathrm{asc})$ in the definition of $\mathrm{Ind}_\beta$ do not involve the probe; so each of the suprema $C_\beta$, $\mathbf{D}$ of Section~G5.1.3 is the supremum of the same set of numbers for the two probes.
Moreover $\operatorname{supp}U_{r'}\beta=r'\operatorname{supp}\beta$, and $r'$ maps the closed coordinate cubes onto the closed coordinate cubes of the same side, and $[-2,2]^3$ onto itself.
The conditions (g1)--(g9) involve the probe only through $(C_\beta,\mathbf{D},D_\beta)$ (Section~G5.1.9).
Hence $e_\sharp(U_{r'}\beta)=e_\sharp(\beta)$, and (DAT) holds for $U_{r'}\beta$ with every placement that is admitted for it.
The right side of (N3) contains, in addition, the constant $\mathrm{vt}$ of the line (C6$'$) and the number $d_{\mathrm{sep}}$.
It decreases as $d_{\mathrm{sep}}$ increases, and $d_{\mathrm{sep}}\ge M\cdot r_*/(8L)-1$ for all these choices (Proposition~35.3\,(e)).
So it is at most
\[
q_x^2\cdot\mathrm{CVj}\cdot h_G^2\cdot\mathrm{pp}\cdot\exp\left(-\mathrm{vt}\cdot\left(\frac{M\cdot r_*}{8L}-4M-7\right)\right),
\]
and $r_*$ does not depend on $N$ (Proposition~35.1\,(c)).
\qed

\subsection*{2. Proposition 35.6}

\subsubsection*{Proof of Proposition 35.6}

Throughout the proof, $\theta\in\mathbb{R}^3$, $A\in\mathrm{Ksf}$ and $(A^{g},\mathscr{A}^0)$ are fixed, and we write $\mathrm{Lg}(t):=\mathrm{Lg}(t;\theta,A)$ and $\mathrm{Rt}(t):=\mathrm{Rt}(t;\theta,A)$.

\paragraph{2.1. The case $\hat{\mathfrak{a}}=0$.}
In this case $\mathscr{C}(t)=0$ for all $t$ (Section~G5.1.4).
By (6.1) in Proposition~34.1\,(c),
\[
\mathrm{Rt}(t)=\frac{Y^{x}_{*}\left(A+c_\theta;0\right)}{Y^{x}_{*}\left(A+c_\theta;0\right)}=1\qquad\text{on }\mathrm{DDo}(t_x^+);
\]
hence $\partial_{\mathfrak{c}}\mathrm{Rt}=0$, and (6.2) in Proposition~34.1\,(d) gives $\mathrm{Lg}=0$.
Consequently $\mathrm{Lg}_{\mathrm{Scl}}=0$ for every $\mathrm{Scl}$.
Moreover $\varpi_{\mathrm{sep}}=0$ (Section~G5.1.8) and $\varpi_{\mathrm{sep}}^+\ge0$, so that the two inequalities of the statement become $0\le0\le\varpi_{\mathrm{sep}}^+$.

From now on let $\hat{\mathfrak{a}}\ne0$.
Then every $\mathrm{Ypr}_{\mathfrak{c}}$ is non-empty and $d_{\mathrm{sep}}<\infty$ (Proposition~35.3\,(g)).

\paragraph{2.2. The tree expansion.}
By Proposition~35.5\,(N1) and the third condition of (g1), and since $\mathrm{wg}_\delta\ge1$,
\[
\mathrm{alG}\cdot\mathrm{rw}_G\left(\mathrm{Vj}_+(t);1\right)\le\mathrm{alG}\cdot\mathrm{rw}_G\left(\mathrm{Vj}_+(t);\mathrm{wg}_\delta\right)\le\mathrm{alG}\cdot\bar{r}\le\frac{1}{2}\qquad\left(t\in\mathrm{DD}(t_x)\right).
\]
The coefficients of $\mathrm{Vj}_+(t)$ are analytic on $\mathrm{DDo}(t_x^+)$ (Proposition~35.5\,(d)), and $\mathrm{Vj}_+(t)$ is even and has no constant term.
Hence (Q3)\,(a) and (Q3)\,(c) apply with $\mathrm{Fv}(t)=\mathrm{Vj}_+(t)$ and $\mathrm{Uo}=\mathrm{DDo}(t_x)$.
Put
\[
\mathrm{Wtr}(t):=\mathrm{Vj}_0(t)+\sum_{i\ge1}\frac{1}{i!}\,\tilde{E}_i\left(\mathrm{Vj}_+(t),\dots,\mathrm{Vj}_+(t)\right)\qquad\left(t\in\mathrm{DDo}(t_x)\right).
\]
The function $\mathrm{Wtr}$ is analytic on $\mathrm{DDo}(t_x)$, by (Q3)\,(c) and, for $\mathrm{Vj}_0$, by Proposition~35.5\,(d).
By (Q3)\,(a) and (5.3) in Proposition~35.5\,(d),
\[
\mathrm{Ztm}(t)=\det D_*(\mathscr{A})\cdot\exp\mathrm{Wtr}(t).
\]
By (6.1) in Proposition~34.1\,(c), for $t\in\mathrm{DDo}(t_x)$,
\[
\exp\mathrm{Lg}(t)=\mathrm{Rt}(t)=\exp\left\{\sum_{X\in D_*}\mathfrak{P}_*\left(X,\mathscr{C}(t)\right)\right\}\cdot\frac{\mathrm{Ztm}(t)}{\mathrm{Ztm}(0)}=\exp\left\{\sum_{X\in D_*}\mathfrak{P}_*\left(X,\mathscr{C}(t)\right)+\mathrm{Wtr}(t)-\mathrm{Wtr}(0)\right\}.
\]
The function $t\mapsto\mathfrak{P}_*\left(X,\mathscr{C}(t)\right)$ is analytic on $\mathrm{DDo}(t_x^+)$, by Proposition~34.2 and (f1), since $t\mapsto\mathscr{C}(t)$ is linear; it vanishes at $t=0$, because $\mathfrak{P}_*(X,0)=0$ (Proposition~34.2); and $D_*$ is a finite set.
So $\mathrm{Lg}(t)$ and the exponent on the right are continuous functions on the connected set $\mathrm{DDo}(t_x)$ with the same exponential.
Their difference is continuous with values in $2\pi i\,\mathbb{Z}$, hence constant, and both functions vanish at $t=0$ (Proposition~34.1\,(d)).
Therefore
\begin{equation}
\mathrm{Lg}(t)=\sum_{X\in D_*}\mathfrak{P}_*\left(X,\mathscr{C}(t)\right)+\mathrm{Wtr}(t)-\mathrm{Wtr}(0)\qquad\left(t\in\mathrm{DDo}(t_x)\right).\tag{6.3}
\end{equation}
Let $\mathrm{Scl}$ be a subset of $\mathrm{Cel}$ with $\#|\mathrm{Scl}|\ge2$.
The M\"obius components are formed on $\mathrm{DD}(t_x/2)$, which is contained in $\mathrm{DDo}(t_x)$.
They depend linearly on the function, and the component of the constant $\mathrm{Wtr}(0)$ is $0$ by (Q1).
Hence (6.3) gives
\[
\mathrm{Lg}_{\mathrm{Scl}}=\sum_{X\in D_*}\left(\mathfrak{P}_*\left(X,\mathscr{C}(\cdot\,)\right)\right)_{\mathrm{Scl}}+\mathrm{Wtr}_{\mathrm{Scl}}\qquad\text{on }\mathrm{DD}(t_x/2).
\]

\paragraph{2.3. The functions $\mathfrak{P}_*$.}
Recall that $\mathfrak{P}_*$ collects the determinants of the steps $j=0,\dots,k_*-1$, including the first step with its determinant of size $192\,n_0$.
By Proposition~34.2 the function $\mathscr{C}\mapsto\mathfrak{P}_*(X,\mathscr{C})$ on $\mathfrak{X}_*$ has the property required in (m3).
Hence $t\mapsto\mathfrak{P}_*\left(X,\mathscr{C}(t)\right)$ does not depend on $t_{\mathfrak{c}}$ unless $X$ contains a point of $\mathrm{Ypr}_{\mathfrak{c}}$.
By the last sentence of (Q1)\,(a), it follows that $\left(\mathfrak{P}_*\left(X,\mathscr{C}(\cdot\,)\right)\right)_{\mathrm{Scl}}=0$ unless $X$ contains a point of $\mathrm{Ypr}_{\mathfrak{c}}$ for every $\mathfrak{c}\in\mathrm{Scl}$.
Let $X$ be such a polymer, and let $\mathfrak{c}\ne\mathfrak{c}'$ be elements of $\mathrm{Scl}$.
Then $X$ contains a point of $\mathrm{Ypr}_{\mathfrak{c}}$ and a point of $\mathrm{Ypr}_{\mathfrak{c}'}$, and the distance of these two points is at least $d_{\mathrm{sep}}$.
By Proposition~35.4\,(SF4)\,(v), $M\cdot\left(d_M(X)+2\right)\ge d_{\mathrm{sep}}$.
Together with (n9) this gives
\[
d_M(X)\ge\frac{d_{\mathrm{sep}}}{M}-2\ge\frac{6M+8}{M}-2\ge1.
\]
Since $d_M(X)\ge1$, Proposition~34.2 gives $\left|\mathfrak{P}_*(X,\mathscr{C})\right|\le2\,C_d\cdot M^3\cdot\exp\left(-\kappa\cdot d_M(X)\right)$ for all $\mathscr{C}\in\mathfrak{X}_*$, and (Q1)\,(b) gives
\[
\sup_{\mathrm{DD}(t_x/2)}\left|\left(\mathfrak{P}_*\left(X,\mathscr{C}(\cdot\,)\right)\right)_{\mathrm{Scl}}\right|\le2^{\#|\mathrm{Scl}|}\cdot2\,C_d\cdot M^3\cdot\exp\left(-\kappa\cdot d_M(X)\right).
\]
Fix $\mathfrak{c}\in\mathrm{Scl}$.
The number of $M$-cubes that contain a point of $\mathrm{Ypr}_{\mathfrak{c}}$ is at most $\#|\mathrm{Ypr}_{\mathfrak{c}}|\le\mathrm{pp}$, and every polymer $X$ under consideration contains one of them.
As $\kappa-\kappa_0\ge2$, Proposition~35.4\,(SF4)\,(v) yields
\begin{align*}
\sum_X\sup\left|\left(\mathfrak{P}_*\left(X,\mathscr{C}(\cdot\,)\right)\right)_{\mathrm{Scl}}\right|
&\le2^{\#|\mathrm{Scl}|}\cdot2\,C_d\cdot M^3\cdot\mathrm{pp}\cdot\sup_{\square}\sum_{\substack{X\supseteq\square\\ d_M(X)\ge\frac{d_{\mathrm{sep}}}{M}-2}}\exp\left(-\kappa_0\cdot d_M(X)\right)\cdot\exp\left(-2\,d_M(X)\right)\\
&\le2^{\#|\mathrm{Scl}|}\cdot2\,C_d\cdot M^3\cdot\mathrm{pp}\cdot K_0\cdot\exp\left(-2\left(\frac{d_{\mathrm{sep}}}{M}-2\right)\right),
\end{align*}
where $\square$ ranges over the $M$-cubes.
Since $\mathrm{vt}\le1/M\le2/M$,
\[
\exp\left(-\frac{2\,d_{\mathrm{sep}}}{M}+4\right)\le\mathrm{e}^4\cdot\exp\left(-\mathrm{vt}\cdot d_{\mathrm{sep}}\right).
\]
Hence the left side of the preceding estimate is at most
\[
2^{\#|\mathrm{Scl}|}\cdot2\,\mathrm{e}^4\cdot C_d\cdot M^3\cdot K_0\cdot\mathrm{pp}\cdot\exp\left(-\mathrm{vt}\cdot d_{\mathrm{sep}}\right).
\]

\paragraph{2.4. The function $\mathrm{Wtr}$.}
Choose $\mathfrak{c}\ne\mathfrak{c}'$ in $\mathrm{Scl}$.
The function $\mathrm{Wtr}$ is analytic on $\mathrm{DDo}(t_x)$, which is a neighbourhood of $\mathrm{DD}(t_x/2)$.
By (Q1)\,(c) with $\mathrm{rad}=t_x/2$,
\[
\sup_{\mathrm{DD}(t_x/2)}\left|\mathrm{Wtr}_{\mathrm{Scl}}\right|\le2^{\#|\mathrm{Scl}|-2}\cdot\left(\frac{t_x}{2}\right)^2\cdot\sup_{\mathrm{DD}(t_x/2)}\left|\partial_{\mathfrak{c}}\partial_{\mathfrak{c}'}\mathrm{Wtr}\right|.
\]
For $t\in\mathrm{DD}(t_x/2)$, the bound (3.5) of (Q3)\,(c), applied to $\mathrm{Wtr}-\mathrm{Vj}_0$ with the non-empty sets $P=\mathrm{Ypr}_{\mathfrak{c}}$ and $P'=\mathrm{Ypr}_{\mathfrak{c}'}$, gives
\begin{align*}
\left|\partial_{\mathfrak{c}}\partial_{\mathfrak{c}'}\mathrm{Wtr}(t)\right|
&\le\left|\partial_{\mathfrak{c}}\partial_{\mathfrak{c}'}\mathrm{Vj}_0(t)\right|+2\,\mathrm{Nw}_G\left(\partial_{\mathfrak{c}}\partial_{\mathfrak{c}'}\mathrm{Vj}_+(t);1\right)\\
&\quad+2\,\mathrm{alG}\cdot\mathrm{Nw}_G\left(\partial_{\mathfrak{c}}\mathrm{Vj}_+(t);\mathrm{wg}_{\mathrm{Ypr}_{\mathfrak{c}}}\right)\cdot\mathrm{rw}_G\left(\partial_{\mathfrak{c}'}\mathrm{Vj}_+(t);\mathrm{wg}_{\mathrm{Ypr}_{\mathfrak{c}'}}\right)\cdot\exp\left(-\mathrm{vt}\cdot d\left(\mathrm{Ypr}_{\mathfrak{c}},\mathrm{Ypr}_{\mathfrak{c}'}\right)\right).
\end{align*}
By Proposition~35.5\,(N3), the first two terms together are at most
\[
2\,q_x^2\cdot\mathrm{CVj}\cdot h_G^2\cdot\mathrm{pp}\cdot\exp\left(-\mathrm{vt}\cdot\left(d_{\mathrm{sep}}-4M-6\right)\right).
\]
By Proposition~35.5\,(N2), applied to $\mathfrak{c}$ and to $\mathfrak{c}'$, and since $d\left(\mathrm{Ypr}_{\mathfrak{c}},\mathrm{Ypr}_{\mathfrak{c}'}\right)\ge d_{\mathrm{sep}}$, the third term is at most
\[
2\,\mathrm{alG}\cdot q_x^2\cdot\mathrm{CVj}^2\cdot h_G^4\cdot\mathrm{pp}\cdot\exp\left(-\mathrm{vt}\cdot d_{\mathrm{sep}}\right).
\]
As $q_x\cdot t_x/2=1$, it follows that
\[
\left(\frac{t_x}{2}\right)^2\cdot\left|\partial_{\mathfrak{c}}\partial_{\mathfrak{c}'}\mathrm{Wtr}(t)\right|\le2\,\mathrm{CVj}\cdot h_G^2\cdot\mathrm{pp}\cdot\exp\left(-\mathrm{vt}\cdot\left(d_{\mathrm{sep}}-4M-6\right)\right)+2\,\mathrm{alG}\cdot\mathrm{CVj}^2\cdot h_G^4\cdot\mathrm{pp}\cdot\exp\left(-\mathrm{vt}\cdot d_{\mathrm{sep}}\right).
\]

\paragraph{2.5. Summation.}
One has $\#|\mathrm{Scl}|\le N_*$ and $\exp\left(-\mathrm{vt}\cdot d_{\mathrm{sep}}\right)\le\exp\left(-\mathrm{vt}\cdot\left(d_{\mathrm{sep}}-4M-6\right)\right)$.
By 2.2--2.4, on $\mathrm{DD}(t_x/2)$,
\begin{align*}
\left|\mathrm{Lg}_{\mathrm{Scl}}\right|
&\le2^{N_*}\cdot\left(2\,\mathrm{CVj}\cdot h_G^2+2\,\mathrm{alG}\cdot\mathrm{CVj}^2\cdot h_G^4+2\,\mathrm{e}^4\cdot C_d\cdot M^3\cdot K_0\right)\cdot\mathrm{pp}\cdot\exp\left(-\mathrm{vt}\cdot\left(d_{\mathrm{sep}}-4M-6\right)\right)\\
&\le2^{N_*}\cdot\mathrm{CWj}\cdot\mathrm{pp}\cdot\exp\left(-\mathrm{vt}\cdot\left(d_{\mathrm{sep}}-4M-6\right)\right)=\varpi_{\mathrm{sep}},
\end{align*}
because $\mathrm{CWj}=3\,\mathrm{CVj}\cdot h_G^2+2\,\mathrm{alG}\cdot\mathrm{CVj}^2\cdot h_G^4+4\,\mathrm{e}^6\cdot C_d\cdot M^3\cdot K_0$.
By Proposition~35.3\,(e), $d_{\mathrm{sep}}-4M-6\ge M\cdot r_*/(8L)-4M-7$; hence $\varpi_{\mathrm{sep}}\le\varpi_{\mathrm{sep}}^+$.

\paragraph{2.6. Uniformity.}
The number $\varpi_{\mathrm{sep}}$ is built from $N_*$, $\mathrm{CWj}$, $\mathrm{vt}$ and $M$, which are constants of the lines (C4)--(C6$'$), from $\mathrm{pp}$ and from $d_{\mathrm{sep}}$; the number $\varpi_{\mathrm{sep}}^+$ is built in the same way, with $r_*$ and $L$ in the place of $d_{\mathrm{sep}}$.
The numbers $\mathrm{pp}$ and $r_*$ do not depend on $N$, $\theta$, $A$, $(A^{g},\mathscr{A}^0)$, the placement, or on a replacement of $\beta$ by $U_r\beta$ with $r\in\{\pm1\}^3$ (see 1.10 and Proposition~35.1\,(c)).
\qed

\subsection*{3. Proposition 35.7}

\subsubsection*{Proof of Proposition 35.7}

\paragraph{3.1. Clause (a).}
\emph{The operator of level~$0$.}
Lattice units are used (Section~I.2.5): $\eta_0=1$, the charge of level~$0$ is $q=e\,\eta^{\frac{1}{2}}$, and the spinors are the elements of $\mathfrak{S}^{\mathrm{ch}}(1;n_*)$, a space of dimension $192\,n_*^3$ (there are $n_*^3$ photon sites, $48$ labels and $4$ components).
One has $n_*=h\,L^{T^{*}(e)+N}\ge L\ge5$ (Section~I.6.4; $h\ge1$ by Section~I.0.5, because $\bar{m}\le1$; $T^{*}(e)\ge1$ by Section~I.6.2).
Hence Proposition~31.1 applies with $n=n_*$.
The pair $P_N(e)=\left(m_{0,N}(e),\varepsilon^{\mathrm{v}}_{0,N}(e)\right)$ of Section~I.6.2 is real (Proposition~35.4\,(SF1)).
Put
\[
\boldsymbol{\mu}_N=\bar{m}\,\eta+m_{0,N}(e),\qquad M_{\mathrm{ch}}:=\boldsymbol{\mu}_N+\frac{\rho}{2}\,\mathrm{sw}.
\]
For a complex one-cochain $B$ on $\mathbb{Z}^3/n_*\mathbb{Z}^3$, (I.2.3), (I.2.5) and (I.1.2) give
\begin{equation}
\tag{7.0}
\left((D_B+\boldsymbol{\mu}_N)f\right)(x)=M_{\mathrm{ch}}\,f(x)-\sum_{i=0}^{2}M_{i}(x)\,f\left(x^{(i)}\right)-M_{3}(x)\exp\left\{i\,q\,B(c,c_{3})\right\}f\left(x^{(3)}\right),
\end{equation}
where $x=(c,g)$, $x^{(i)}=(c,g\,s_i)$ for $i\le2$, $x^{(3)}=(c_{3},g\,s_3)$ and $c_{3}=c+g\,e_1$.
In the basis of the delta functions at the points $(x,\beta)$, with $x$ in a fundamental domain of $n_*\mathbb{Z}^3$ and $\beta\in\{1,\dots,4\}$, every entry of the matrix of $D_B+\boldsymbol{\mu}_N$ is either a constant or an entry of $-M_{3}(x)$ multiplied by $\exp\{i\,q\,B(c,c_{3})\}$ and by a sign; the sign is $-1$ if $x^{(3)}$ leaves the fundamental domain (Section~I.1.5).
This is (I.O1).
The determinant is a polynomial in the entries.

\emph{Entire functions; continuity.}
Let $B=A(F^{\mathrm{u}})+H_{\theta}+a(t)$.
Since $q/e=\eta^{\frac{1}{2}}$, one has for every bond $b$
\[
\exp\left\{i\,q\,B(b)\right\}=\exp\left\{i\,q\,\left(A(F^{\mathrm{u}})+H_{\theta}\right)(b)\right\}\cdot\exp\left\{i\,\eta^{\frac{1}{2}}\sum_{\mathfrak{c}}t_{\mathfrak{c}}\,\hat{\mathfrak{a}}_{\mathfrak{c}}(b)\right\}.
\]
The map $F^{\mathrm{u}}\mapsto A(F^{\mathrm{u}})$ is linear (it is the inverse of the bijection $d\colon V_{\mathrm{t}}\to F_N$), the map $\theta\mapsto H_{\theta}$ is linear (Section~I.6.6), and both are real.
Hence the first factor has modulus one and is continuous in $(F^{\mathrm{u}},\theta)$, while the second factor is an entire function of $t$.
Moreover $\det(D_0+\bar{m}\,\eta)>0$ by Proposition~31.1\,(a) with $z=\bar{m}\,\eta>0$.
By (I.7.2), $\mathrm{dens}^{\theta,a(t)}_N(F^{\mathrm{u}})$ is a polynomial in finitely many functions, each of which is entire in $t$ and continuous in $(t,F^{\mathrm{u}},\theta)$.
Consequently $\mathrm{dens}^{\theta,a(t)}_N(F^{\mathrm{u}})$ is entire in $t$ and continuous in $(t,F^{\mathrm{u}},\theta)$.

\emph{A bound.}
Operator norms on $\mathfrak{S}^{\mathrm{ch}}(1;n_*)$ are taken with respect to $(\cdot,\cdot)^{\mathrm{ch}}$.
One has $|M_{i}(x)|_2=\frac{w_i}{2}(1+\rho)$ by (I.m6); the map $x\mapsto x^{(i)}$ is a bijection of $T^{\mathrm{f}}_0$ which preserves the measure $\frac{1}{48}\sum_c\sum_g$ (I.s1); and $|\exp\{i\,q\,B(b)\}|\le\exp\left(q\,|\operatorname{Im}B|_\infty\right)$.
By (7.0) the norm of $D_B+\boldsymbol{\mu}_N$ is at most
\[
|M_{\mathrm{ch}}|+\frac{1+\rho}{2}\left(w_0+w_1+w_2+w_3\exp\left(q\,|\operatorname{Im}B|_\infty\right)\right),
\]
and $w_3=6$.
Every eigenvalue of a linear map has modulus at most its norm; hence the modulus of the determinant is at most the norm raised to the power $192\,n_*^3$.
For $B=A(F^{\mathrm{u}})+H_{\theta}+a(t)$ one has $\operatorname{Im}B=\operatorname{Im}a(t)$ and
\[
q\,|\operatorname{Im}a(t)|_\infty\le\eta^{\frac{1}{2}}\sum_{\mathfrak{c}}|\operatorname{Im}t_{\mathfrak{c}}|\cdot|\hat{\mathfrak{a}}_{\mathfrak{c}}|_\infty .
\]
Therefore $\left|\mathrm{dens}^{\theta,a(t)}_N(F^{\mathrm{u}})\right|\le\mathrm{bnd}(t)$, where
\begin{multline*}
\mathrm{bnd}(t):=\exp\left(-\varepsilon^{\mathrm{v}}_{0,N}(e)\,n_*^3\right)\det(D_0+\bar{m}\,\eta)^{-1}\\
\times\left[|M_{\mathrm{ch}}|+\frac{1+\rho}{2}\left(\mathrm{sw}-6+6\exp\left(\eta^{\frac{1}{2}}\sum_{\mathfrak{c}}|\operatorname{Im}t_{\mathfrak{c}}|\cdot|\hat{\mathfrak{a}}_{\mathfrak{c}}|_\infty\right)\right)\right]^{192\,n_*^3}.
\end{multline*}
The function $\mathrm{bnd}$ is continuous in $t$ and contains neither $F^{\mathrm{u}}$ nor $\theta$.
Hence $\mathrm{dens}^{\theta,a(t)}_N(F^{\mathrm{u}})$ is bounded on compact sets of $t$, uniformly in $(F^{\mathrm{u}},\theta)$.

\emph{The function $N_0$.}
Apply (7.0) with $\bar{m}\,\eta$ in place of $\boldsymbol{\mu}_N$ (also inside $M_{\mathrm{ch}}$) and with $B=(t/e)\,\hat{\mathfrak{a}}$, $t\in\mathbb{C}$.
The entries of the matrix of $D_{(t/e)\hat{\mathfrak{a}}}+\bar{m}\,\eta$ are then constants or constants times $\exp\{i\,t\,\eta^{\frac{1}{2}}\hat{\mathfrak{a}}(b)\}$.
Hence $N_0$ is entire on $\mathbb{C}$.

\emph{The functions $t\mapsto\bar{\nu}^{a(t)}(\mathrm{Ev})$.}
Let $\mathrm{Ev}$ be a Borel subset of $F_N$.
By Section~I.7.3,
\[
\bar{\nu}^{a(t)}(\mathrm{Ev})=\int_{[-\pi,\pi]^3}\frac{\mathrm{d}\theta}{(2\pi)^3}\int_{\mathrm{Ev}}\mathrm{d}\gamma_N(F^{\mathrm{u}})\,\mathrm{dens}^{\theta,a(t)}_N(F^{\mathrm{u}}).
\]
Apply (h1) with $n=N_{*}$, $\mathrm{rad}_0=\infty$, $\mathrm{Ms}=[-\pi,\pi]^3\times\mathrm{Ev}$, $\mathrm{mq}=(2\pi)^{-3}\,\mathrm{d}\theta\times\gamma_N$ (a measure of mass at most $1$), $\mathrm{gf}=\mathrm{dens}$ (continuous, hence measurable), and with $\mathrm{maj}_{\mathrm{rad}}$ the constant $\sup_{\mathrm{DD}(\mathrm{rad})}\mathrm{bnd}$.
It follows that $t\mapsto\bar{\nu}^{a(t)}(\mathrm{Ev})$ is entire on $\mathbb{C}^{\mathrm{Cel}}$.

\emph{The number $\bar{\nu}(F_N)$.}
The inequalities $0<\bar{\nu}(F_N)<\infty$ are Proposition~31.1\,(c).

\paragraph{3.2. Clause (b).}
By \cite[Definition~5.11]{QEDI} (see Section~I.5.7),
\[
\mathfrak{G}=\left\{F^{\mathrm{u}}\in F_N:\ |\zeta_j(\Upsilon)|\le p_{0,j}\ \text{for all }\Upsilon\text{ and all }j<k_{*},\ \ |dA_{*}|_\infty\le p_{*}\right\},
\]
and $\mathfrak{G}^{-}$ is the same set with $p_{0,j}-1$ and $p_{*}-1$ in place of $p_{0,j}$ and $p_{*}$.
Let $\mathrm{pl}$ denote a plaquette of $T^0_*$.
The set $F_N\setminus\mathfrak{G}^{-}$ is the union of the events $\{|\zeta_j(\Upsilon)|>p_{0,j}-1\}$, where $j<k_{*}$ and $\Upsilon$ is a basis element of $V_{j}$, and of the events $\{|(dA_{*})(\mathrm{pl})|>p_{*}-1\}$.
By Proposition~35.4\,(SF5), $\zeta_j(\Upsilon)=\langle\phi_{j,\Upsilon},F^{\mathrm{u}}\rangle$ with $\|\phi_{j,\Upsilon}\|=1$, and $(dA_{*})(\mathrm{pl})=\langle\phi_{\mathrm{pl}},F^{\mathrm{u}}\rangle$ with $\|\phi_{\mathrm{pl}}\|\le1$, all these $\phi$ lying in $F_N$.

Let $j<k_{*}$ and $i:=k_{*}-j\ge1$.
By Section~I.6.1, $a_j=a_{*}L^{-i}$, $e_j=e_{*}L^{-\frac{i}{2}}$ and $v_j=v_{*}+i\,\mathrm{lg}$.
The first two clauses of (T1), which hold by Proposition~35.1\,(e), give $v_{*}\ge8$ and $p_0\cdot v_{*}^{p_0-1}\ge28$.
The numbers $p_0$ and $p$ are integers $\ge1$ (line (C7)).
Hence $p_{*}=v_{*}^{p}\ge8$, $p_{0,j}=v_j^{p_0}\ge v_{*}^{p_0}=p_{0*}\ge8$, and by the binomial theorem
\[
p_{0,j}=\left(v_{*}+i\,\mathrm{lg}\right)^{p_0}\ge v_{*}^{p_0}+p_0\cdot v_{*}^{p_0-1}\cdot i\cdot\mathrm{lg}\ge p_{0*}+28\,i\cdot\mathrm{lg}=p_{0*}+14\,i\log L .
\]
Use the third inequality of Proposition~31.1\,(d) at the levels $p_{0,j}-1\ge0$ and $p_{*}-1\ge0$ respectively, the inequality $(y-1)^2\ge y$ for real $y\ge3$ (at $y=p_{0,j}$), and $C_{0}\ge1$.
This gives
\begin{align*}
\bar{\mu}\left\{|\zeta_j(\Upsilon)|>p_{0,j}-1\right\}&\le2\exp\left(-\frac{(p_{0,j}-1)^2}{2}\right)\le2\exp\left(-\frac{p_{0,j}}{2}\right)\le2\exp\left(-\frac{p_{0*}}{2}\right)L^{-7i},\\
\bar{\mu}\left\{|(dA_{*})(\mathrm{pl})|>p_{*}-1\right\}&\le2\exp\left(-\frac{(p_{*}-1)^2}{2}\right)\le2\exp\left(-\frac{(p_{*}-1)^2}{2C_{0}^2}\right).
\end{align*}
The space $V_{j}$ has at most $2n_j=2\ell_*^3L^{3i}$ basis elements, and $T^0_*$ has $3\ell_*^3$ plaquettes (Proposition~35.4\,(SF5)).
Since $\sum_{i\ge1}L^{-4i}\le1$, it follows that
\begin{align*}
\bar{\mu}\left(F_N\setminus\mathfrak{G}^{-}\right)&\le\sum_{i\ge1}2\ell_*^3L^{3i}\cdot2\exp\left(-\frac{p_{0*}}{2}\right)L^{-7i}+3\ell_*^3\cdot2\exp\left(-\frac{(p_{*}-1)^2}{2C_{0}^2}\right)\\
&\le4\ell_*^3\exp\left(-\frac{p_{0*}}{2}\right)+6\ell_*^3\exp\left(-\frac{(p_{*}-1)^2}{2C_{0}^2}\right)\le\varpi_0 .
\end{align*}

\paragraph{3.3. Clause (c): continuity, boundedness.}
Let $t\in\mathrm{DDo}(t_x^{+})$; by (f1), $\mathscr{C}:=\mathscr{C}(t)$ lies in $\mathfrak{X}_*$.
Let $(\theta,A)\in\mathbb{R}^3\times\mathrm{Ksf}$.
By (2.3) in Proposition~35.4\,(SF5), $A+c_{\theta}$ lies in $\mathrm{Dom}_*$; hence, by the definition of $\mathrm{Dom}_*$ (Section~I.5.3), $\mathscr{A}_{\mathrm{u}}:=H_{*}[A+c_{\theta}]$ lies in $\frac14\tilde{R}_*$, a subset of $\tilde{R}_*$.
By \cite[Proposition~15.5]{QEDI} the photon minimiser $H_{*}$ is linear and $H_{*}c_{\theta}=c_{\theta}$; that proposition is formulated for an integer $k\ge0$ and a side in $\mathrm{Cl}$, and it is used here with $k=k_{*}$ and the side $\ell_*$, which lies in $\mathrm{Cl}$ by Proposition~35.1\,(d).
Hence $\mathscr{A}_{\mathrm{u}}=H_{*}A+c_{\theta}$, and $(\theta,A)\mapsto\mathscr{A}_{\mathrm{u}}$ is the restriction of a linear map between spaces of finite dimension ($c_{\theta}$ is linear in $\theta$ by Section~I.6.6); in particular it is continuous.
The field $\mathscr{T}_{\mathrm{u}}:=\mathscr{A}_{\mathrm{u}}+\mathscr{C}$ lies in $\tilde{R}_*+\mathfrak{X}_*$, a subset of $\tilde{G}_*$ (Proposition~35.4\,(SF2)).
By (2.4) in Proposition~35.4\,(SF6),
\begin{multline*}
Y^{x}_{*}(A+c_{\theta};\mathscr{C})=\det D_{*}(0)^{-1}\exp\left\{-\varepsilon^{\mathrm{v},t}_*\,\ell_*^3+\sum_{X\in D_{*}}\mathfrak{P}_*(X,\mathscr{C})\right\}\\
\times\int\mathrm{D}\Psi\,\exp\left\{-\langle\bar{\Psi},D_{*}(\mathscr{T}_{\mathrm{u}})\Psi\rangle-m_{*}\langle\bar{\psi},\psi\rangle+\sum_{X\in D_{*}}\mathscr{E}_*(X;\mathscr{A}_{\mathrm{u}},\mathscr{C},\psi,\chi)\right\},
\end{multline*}
the fields being inserted by
\[
\psi(\xi)\mapsto\sum_u\mathrm{Leg}_{\mathscr{T}_{\mathrm{u}}}(\xi;u)\,\Psi(u),\qquad\chi(\boldsymbol{\zeta})\mapsto\sum_u\mathrm{Leg}_{\mathscr{T}_{\mathrm{u}}}(\boldsymbol{\zeta};u)\,\Psi(u).
\]
The coefficients of the exponent, which is an even element of $\mathrm{Gr}[\Psi]$, are finite sums of products of the weights $\eta_*^3/48$, of $m_{*}$, and of the numbers $D_{*}(\mathscr{T}_{\mathrm{u}})(y,y')_{\beta\beta'}$, $\mathrm{Leg}_{\mathscr{T}_{\mathrm{u}}}(\mathrm{up};u)$ and $\mathscr{E}_{nm}(X;\mathscr{A}_{\mathrm{u}},\mathscr{C};\,\cdot\,)$.
These numbers are analytic functions of $\mathscr{T}_{\mathrm{u}}$ on $\tilde{G}_*$ (by (f3)), respectively of $\mathscr{A}_{\mathrm{u}}$ on $\tilde{R}_*$ (by Proposition~35.4\,(SF3)); hence they are continuous in $\mathscr{A}_{\mathrm{u}}$.
The algebra $\mathrm{Gr}[\Psi]$ has finite dimension; the exponential of an even element is the exponential of its constant term times a polynomial in its other coefficients; and $\int\mathrm{D}\Psi$ is linear.
The factor in front of the integral does not contain $(\theta,A)$.
Therefore $(\theta,A)\mapsto Y^{x}_{*}(A+c_{\theta};\mathscr{C}(t))$ is continuous on $\mathbb{R}^3\times\mathrm{Ksf}$.

The set $\mathrm{Ksf}$ is compact (Proposition~35.4\,(SF5)); so this function is bounded on $[-\pi,\pi]^3\times\mathrm{Ksf}$.
The function on $[-\pi,\pi]^3\times V_{\mathrm{t}}(T^0_*)$ which equals it on the closed set $[-\pi,\pi]^3\times\mathrm{Ksf}$ and vanishes elsewhere, written $\chi_{*}(A)\,Y^{x}_{*}(A+c_{\theta};\mathscr{C}(t))$, is a bounded Borel function.
At $t=0$, divided by $\bar{\nu}(F_N)\in(0,\infty)$ (see 3.1), it is the density of $\mathrm{vm}$ with respect to the finite measure $(2\pi)^{-3}\,\mathrm{d}\theta\times\mathrm{d}\mathsf{m}_{*}$.

\paragraph{3.4. Positivity.}
\textup{(P)}\enspace\emph{For every real $t\in\mathrm{DDo}(t_x^{+})$, every $\theta\in\mathbb{R}^3$ and every $A\in\mathrm{Ksf}$ one has $Y^{x}_{*}(A+c_{\theta};\mathscr{C}(t))>0$.}

\emph{Proof of} (P).
Fix $t$ and $\theta$, and put $\mathrm{gY}(A):=Y^{x}_{*}(A+c_{\theta};\mathscr{C}(t))$; by 3.3 this is a continuous complex function on $\mathrm{Ksf}$.
The cochain $a(t)$ is real; hence $\nu^{\theta,a(t)}$ is a positive measure (Proposition~31.1\,(a)).
The cochain $a(t)^{[j]}$ lies in $\mathfrak{X}_j$ for all $j\le k_{*}$ (Proposition~35.4\,(DF5)), and $a(t)^{[k_{*}]}=\mathscr{C}(t)$.
By (2.5) in Proposition~35.4\,(SF6), for every bounded Borel function $\Phi\ge0$ on $V_{\mathrm{t}}(T^0_*)$,
\[
\int\mathrm{d}\mathsf{m}_{*}(A)\,\Phi(A)\,\chi_{*}(A)\,\mathrm{gY}(A)=\nu^{\theta,a(t)}\left(\mathbf{1}_{\mathfrak{G}}\,\Phi(A_{*})\right),
\]
and the right side is a real number $\ge0$.
The measure $\mathrm{d}\mathsf{m}_{*}$ is a non-degenerate centred Gaussian measure on $V_{\mathrm{t}}(T^0_*)$ (Proposition~35.4\,(SF5)); hence it has a positive density with respect to Lebesgue measure, and $\mathrm{d}\mathsf{m}_{*}(\mathrm{Bo})>0$ for every open ball $\mathrm{Bo}$.
Let $\mathrm{Ain}$ be an interior point of $\mathrm{Ksf}$.
If $\operatorname{Im}\mathrm{gY}(\mathrm{Ain})\neq0$, there is an open ball $\mathrm{Bo}$ around $\mathrm{Ain}$ inside $\mathrm{Ksf}$ on which $\operatorname{Im}\mathrm{gY}$ has the sign of $\operatorname{Im}\mathrm{gY}(\mathrm{Ain})$; for $\Phi=\mathbf{1}_{\mathrm{Bo}}$ the left side of the last display has the imaginary part $\int_{\mathrm{Bo}}\mathrm{d}\mathsf{m}_{*}\,\operatorname{Im}\mathrm{gY}\neq0$.
If $\operatorname{Im}\mathrm{gY}(\mathrm{Ain})=0$ and $\operatorname{Re}\mathrm{gY}(\mathrm{Ain})<0$, there is such a ball on which $\operatorname{Re}\mathrm{gY}<0$; for $\Phi=\mathbf{1}_{\mathrm{Bo}}$ the left side has a negative real part.
Both cases contradict the display.
Hence $\mathrm{gY}(\mathrm{Ain})\ge0$.
The set $\mathrm{Ksf}$ is convex and contains a ball $\mathrm{Bo}_0$ around $0$ (Proposition~35.4\,(SF5)).
For $A\in\mathrm{Ksf}$ and $0\le\mathrm{sv}<1$ the set $\mathrm{sv}\,A+(1-\mathrm{sv})\,\mathrm{Bo}_0$ is a ball around $\mathrm{sv}\,A$ inside $\mathrm{Ksf}$; hence $\mathrm{sv}\,A$ is an interior point, $\mathrm{gY}(\mathrm{sv}\,A)\ge0$, and
\[
\mathrm{gY}(A)=\lim_{\mathrm{sv}\to1}\mathrm{gY}(\mathrm{sv}\,A)\ge0 .
\]
By Proposition~34.1\,(c), $\mathrm{gY}(A)\neq0$.
This proves (P).

\paragraph{3.5. Clauses (c) (positivity, mass), (d), (e).}
By (P) at $t=0$ the density of $\mathrm{vm}$ is $\ge0$.
Hence $\mathrm{vm}$ is a positive finite measure, and the complement of $[-\pi,\pi]^3\times\mathrm{Ksf}$ is a null set for it.

\emph{Clause} (d).
Let $t\in\mathrm{DDo}(t_x^{+})$.
By (2.5) with $\Phi=1$ and $a=a(t)$ (which is admitted by Proposition~35.4\,(DF5)),
\[
\nu^{\theta,a(t)}(\mathfrak{G})=\int\mathrm{d}\mathsf{m}_{*}(A)\,\chi_{*}(A)\,Y^{x}_{*}(A+c_{\theta};\mathscr{C}(t))\qquad\text{for every }\theta .
\]
Integration over $\theta$ (Section~I.7.3) gives
\[
\mathrm{Ysf}(t)=\bar{\nu}^{a(t)}(\mathfrak{G})=\int_{[-\pi,\pi]^3}\frac{\mathrm{d}\theta}{(2\pi)^3}\int\mathrm{d}\mathsf{m}_{*}(A)\,\chi_{*}(A)\,Y^{x}_{*}(A+c_{\theta};\mathscr{C}(t)).
\]
The integrand is a bounded Borel function on the product (by 3.3), and the two measures are finite; by Fubini's theorem the iterated integral is the integral with respect to the product measure.
For $A\in\mathrm{Ksf}$, (6.1) and (6.2) in Proposition~34.1\,(c), (d) give
\[
Y^{x}_{*}(A+c_{\theta};\mathscr{C}(t))=Y^{x}_{*}(A+c_{\theta};0)\,\mathrm{Rt}(t;\theta,A)=Y^{x}_{*}(A+c_{\theta};0)\,\exp\mathrm{Lg}(t;\theta,A).
\]
By the definition of $\mathrm{vm}$ it follows that
\[
\mathrm{Ysf}(t)=\bar{\nu}(F_N)\int\mathrm{vm}(\mathrm{d}\theta\,\mathrm{d}A)\,\exp\mathrm{Lg}(t;\theta,A).
\]

\emph{Mass.}
At $t=0$ one has $\mathrm{Lg}(0;\theta,A)=0$ (Proposition~34.1\,(d)) and $a(0)=0$.
Hence the total mass of $\mathrm{vm}$ is
\[
\frac{\mathrm{Ysf}(0)}{\bar{\nu}(F_N)}=\frac{\bar{\nu}(\mathfrak{G})}{\bar{\nu}(F_N)}=\bar{\mu}(\mathfrak{G})\le1 .
\]
The set $\mathfrak{G}^{-}$ is a subset of $\mathfrak{G}$; by (b), $\bar{\mu}(\mathfrak{G})\ge\bar{\mu}(\mathfrak{G}^{-})\ge1-\varpi_0$.

\emph{Clause} (e).
Let $t\in\mathrm{DDo}(t_x^{+})$ be real, $\theta\in\mathbb{R}^3$ and $A\in\mathrm{Ksf}$.
For $0\le\mathrm{sv}\le1$ the point $\mathrm{sv}\,t$ is real and lies in $\mathrm{DDo}(t_x^{+})$.
By (P) at $\mathrm{sv}\,t$ and at $0$,
\[
\exp\mathrm{Lg}(\mathrm{sv}\,t;\theta,A)=\mathrm{Rt}(\mathrm{sv}\,t;\theta,A)=\frac{Y^{x}_{*}(A+c_{\theta};\mathscr{C}(\mathrm{sv}\,t))}{Y^{x}_{*}(A+c_{\theta};0)}>0 .
\]
Hence $\mathrm{sv}\mapsto\operatorname{Im}\mathrm{Lg}(\mathrm{sv}\,t;\theta,A)$ is a continuous function on $[0,1]$ (Proposition~34.1\,(d)) with values in $2\pi\mathbb{Z}$, which vanishes at $\mathrm{sv}=0$.
Being continuous with values in a discrete set on an interval, it is constant; so it vanishes at $\mathrm{sv}=1$.

\paragraph{3.6. The shift.}
Let $z\in\mathbb{R}^{\mathrm{Cel}}$, put $|z|_2^2:=\sum_{\mathfrak{c}}z_{\mathfrak{c}}^2$, and
\[
a(ez)=\sum_{\mathfrak{c}}z_{\mathfrak{c}}\,\hat{\mathfrak{a}}_{\mathfrak{c}},\qquad\mathrm{Gd}:=d\,a(ez)=\sum_{\mathfrak{c}}z_{\mathfrak{c}}\,d\hat{\mathfrak{a}}_{\mathfrak{c}},\qquad\langle\mathrm{Gd},F^{\mathrm{u}}\rangle=\sum_{\mathfrak{c}}z_{\mathfrak{c}}X_{\mathfrak{c}};
\]
here $\mathrm{Gd}$ is an element of $F_N$.
The supports of the $d\hat{\mathfrak{a}}_{\mathfrak{c}}$ are pairwise disjoint and $\|d\hat{\mathfrak{a}}_{\mathfrak{c}}\|^2=\boldsymbol{\sigma}^2$ (Proposition~35.3\,(d)); hence $\|\mathrm{Gd}\|^2=|z|_2^2\cdot\boldsymbol{\sigma}^2$.
The cochain $a(ez)$ lies in $V_{\mathrm{t}}$ (Proposition~35.3\,(a)).
So $A(F^{\mathrm{u}})+a(ez)$ lies in $V_{\mathrm{t}}$ and has the coboundary $F^{\mathrm{u}}+\mathrm{Gd}$; since $d$ is injective on $V_{\mathrm{t}}$, one has $A(F^{\mathrm{u}})+a(ez)=A(F^{\mathrm{u}}+\mathrm{Gd})$, and (I.7.2) gives
\[
\mathrm{dens}^{\theta,a(ez)}_N(F^{\mathrm{u}})=\mathrm{dens}^{\theta,0}_N(F^{\mathrm{u}}+\mathrm{Gd}).
\]
The measure $\gamma_N$ has the density $(2\pi)^{-\frac{\dim F_N}{2}}\exp\left(-\frac{\|F^{\mathrm{u}}\|^2}{2}\right)$ with respect to Lebesgue measure on the Euclidean space $F_N$.
The substitution $F'=F^{\mathrm{u}}+\mathrm{Gd}$ and the identity
\[
-\frac{\|F'-\mathrm{Gd}\|^2}{2}=-\frac{\|F'\|^2}{2}+\langle F',\mathrm{Gd}\rangle-\frac{\|\mathrm{Gd}\|^2}{2}
\]
give, for every $\theta\in\mathbb{R}^3$ and every Borel set $\mathrm{Ev}\subset F_N$,
\begin{equation}
\tag{7.1}
\begin{split}
\nu^{\theta,a(ez)}(\mathrm{Ev})&=\int\mathrm{d}\gamma_N(F')\,\mathbf{1}_{\mathrm{Ev}}(F'-\mathrm{Gd})\,\mathrm{dens}^{\theta,0}_N(F')\,\exp\left(\langle F',\mathrm{Gd}\rangle-\frac{\|\mathrm{Gd}\|^2}{2}\right)\\
&=\exp\left(-\frac{|z|_2^2\cdot\boldsymbol{\sigma}^2}{2}\right)\int_{\mathrm{Ev}+\mathrm{Gd}}\mathrm{d}\nu^{\theta}\,\exp\left(\sum_{\mathfrak{c}}z_{\mathfrak{c}}X_{\mathfrak{c}}\right).
\end{split}
\end{equation}
All integrands are $\ge0$ by Proposition~31.1\,(a); moreover $\mathrm{dens}^{\theta,0}_N$ is bounded and $\exp\langle F',\mathrm{Gd}\rangle$ is $\gamma_N$-integrable.

\paragraph{3.7. Clause (f).}
The numbers $C_{\beta}$ and $\mathbf{D}$ are finite (Proposition~35.3\,(c1)).
Hence
\[
\mathrm{zq}:=\frac{t_2}{e}=\min\left\{N_{*}^{-\frac{1}{2}},\ C_{\beta}^{-1},\ \frac{1}{2}\left(N_{*}\cdot\mathbf{D}\right)^{-\frac{1}{2}}\right\}>0,
\]
and $N_{*}\cdot\mathrm{zq}^2\le1$, $N_{*}\cdot\mathrm{zq}^2\cdot\mathbf{D}\le\frac14$.
Moreover $\boldsymbol{\sigma}^2\le\mathbf{D}$ (Proposition~35.3\,(c2)).
Let $t=ez$ be a real point of $\mathrm{DD}(t_2)$, that is, $z\in\mathbb{R}^{\mathrm{Cel}}$ and $|z_{\mathfrak{c}}|\le\mathrm{zq}$ for all $\mathfrak{c}$.
Then $|z|_2^2\le N_{*}\cdot\mathrm{zq}^2\le1$ and $\|\mathrm{Gd}\|^2=|z|_2^2\cdot\boldsymbol{\sigma}^2\le N_{*}\cdot\mathrm{zq}^2\cdot\mathbf{D}\le\frac14$.
The functionals $\zeta_j(\Upsilon)$ and $(dA_{*})(\mathrm{pl})$ are linear in $F^{\mathrm{u}}$ (Proposition~35.4\,(SF5)), and by the Cauchy--Schwarz inequality
\[
|\langle\phi_{j,\Upsilon},\mathrm{Gd}\rangle|\le\|\mathrm{Gd}\|\le1,\qquad|\langle\phi_{\mathrm{pl}},\mathrm{Gd}\rangle|\le\|\mathrm{Gd}\|\le1 .
\]
Hence $F'\in\mathfrak{G}^{-}$ implies $F'-\mathrm{Gd}\in\mathfrak{G}$; in other words, $(F_N\setminus\mathfrak{G})+\mathrm{Gd}$ is a subset of $F_N\setminus\mathfrak{G}^{-}$.
Take (7.1) with $\mathrm{Ev}=F_N\setminus\mathfrak{G}$, integrate over $\theta$ and divide by $\bar{\nu}(F_N)$.
Then use $\nu^{\theta}\ge0$, the Cauchy--Schwarz inequality for $\bar{\mu}$, clause (b), and the first inequality of Proposition~31.1\,(d) with $\mathrm{gc}=2\,\mathrm{Gd}\in F_N$.
This gives
\begin{align*}
0\le\frac{\mathrm{zL}(t)}{\bar{\nu}(F_N)}&=\exp\left(-\frac{|z|_2^2\cdot\boldsymbol{\sigma}^2}{2}\right)\bar{\mu}\left(\mathbf{1}\left[F^{\mathrm{u}}\in(F_N\setminus\mathfrak{G})+\mathrm{Gd}\right]\exp\langle\mathrm{Gd},F^{\mathrm{u}}\rangle\right)\\
&\le\exp\left(-\frac{|z|_2^2\cdot\boldsymbol{\sigma}^2}{2}\right)\bar{\mu}\left(F_N\setminus\mathfrak{G}^{-}\right)^{\frac{1}{2}}\,\bar{\mu}\left(\exp\langle2\,\mathrm{Gd},F^{\mathrm{u}}\rangle\right)^{\frac{1}{2}}\\
&\le\exp\left(-\frac{|z|_2^2\cdot\boldsymbol{\sigma}^2}{2}\right)\varpi_0^{\frac{1}{2}}\exp\left(\|\mathrm{Gd}\|^2\right)=\varpi_0^{\frac{1}{2}}\exp\left(\frac{|z|_2^2\cdot\boldsymbol{\sigma}^2}{2}\right)\\
&\le\varpi_0^{\frac{1}{2}}\exp\left(\frac{\mathbf{D}}{2}\right)=\varpi_L .
\end{align*}

\paragraph{3.8. Clause (g).}
\emph{The identity.}
Take (7.1) with $\mathrm{Ev}=F_N$ (note that $F_N+\mathrm{Gd}=F_N$), integrate over $\theta$ and divide by $\bar{\nu}(F_N)$.
This gives
\[
\mathring{\Phi}(ez)=\bar{\mu}\left(\exp\sum_{\mathfrak{c}}\left(z_{\mathfrak{c}}X_{\mathfrak{c}}-\frac{z_{\mathfrak{c}}^2\,\boldsymbol{\sigma}^2}{2}\right)\right)\qquad\text{for all }z\in\mathbb{R}^{\mathrm{Cel}}.
\]
The left side is an entire function of $z\in\mathbb{C}^{\mathrm{Cel}}$ by (a).
For the right side apply (h1) with $n=N_{*}$, $\mathrm{rad}_0=\infty$, $(\mathrm{Ms},\mathrm{mq})=(F_N,\bar{\mu})$ and
\[
\mathrm{gf}(z,F^{\mathrm{u}})=\exp\sum_{\mathfrak{c}}\left(z_{\mathfrak{c}}X_{\mathfrak{c}}-\frac{z_{\mathfrak{c}}^2\,\boldsymbol{\sigma}^2}{2}\right),
\]
which is continuous, and entire in $z$.
For $|z_{\mathfrak{c}}|\le\mathrm{rad}$ one has
\begin{gather*}
\left|\exp\left(-\frac{z_{\mathfrak{c}}^2\,\boldsymbol{\sigma}^2}{2}\right)\right|\le\exp\left(\frac{\mathrm{rad}^2\cdot\boldsymbol{\sigma}^2}{2}\right),\\
\left|\exp\left(z_{\mathfrak{c}}X_{\mathfrak{c}}\right)\right|=\exp\left((\operatorname{Re}z_{\mathfrak{c}})\,X_{\mathfrak{c}}\right)\le\exp\left(\mathrm{rad}\cdot X_{\mathfrak{c}}\right)+\exp\left(-\mathrm{rad}\cdot X_{\mathfrak{c}}\right).
\end{gather*}
Hence $|\mathrm{gf}(z,F^{\mathrm{u}})|\le\mathrm{maj}_{\mathrm{rad}}(F^{\mathrm{u}})$, where
\[
\mathrm{maj}_{\mathrm{rad}}(F^{\mathrm{u}}):=\exp\left(\frac{N_{*}\cdot\mathrm{rad}^2\cdot\boldsymbol{\sigma}^2}{2}\right)\sum_{\mathrm{sgn}}\exp\left\langle\mathrm{rad}\sum_{\mathfrak{c}}\mathrm{sgn}_{\mathfrak{c}}\,d\hat{\mathfrak{a}}_{\mathfrak{c}},F^{\mathrm{u}}\right\rangle,
\]
the sum extending over $\mathrm{sgn}\in\{1,-1\}^{\mathrm{Cel}}$.
By Proposition~31.1\,(d) and Proposition~35.3\,(d),
\[
\bar{\mu}\left(\mathrm{maj}_{\mathrm{rad}}\right)\le\exp\left(\frac{N_{*}\cdot\mathrm{rad}^2\cdot\boldsymbol{\sigma}^2}{2}\right)2^{N_{*}}\exp\left(\frac{N_{*}\cdot\mathrm{rad}^2\cdot\boldsymbol{\sigma}^2}{2}\right)<\infty .
\]
Hence the right side is entire as well.
By (h2) the identity holds for all $z\in\mathbb{C}^{\mathrm{Cel}}$.

\emph{The bound of $\mathring{\Phi}$.}
Let $t=ez\in\mathbb{C}^{\mathrm{Cel}}$.
The measure $\bar{\mu}$ is positive and of mass one, and $\operatorname{Re}(z_{\mathfrak{c}}^2)=(\operatorname{Re}z_{\mathfrak{c}})^2-(\operatorname{Im}z_{\mathfrak{c}})^2$.
Apply Proposition~31.1\,(d) with $\mathrm{gc}=\sum_{\mathfrak{c}}(\operatorname{Re}z_{\mathfrak{c}})\,d\hat{\mathfrak{a}}_{\mathfrak{c}}\in F_N$, for which $\|\mathrm{gc}\|^2=\sum_{\mathfrak{c}}(\operatorname{Re}z_{\mathfrak{c}})^2\boldsymbol{\sigma}^2$:
\begin{align*}
|\mathring{\Phi}(ez)|&\le\exp\left(-\sum_{\mathfrak{c}}\frac{\operatorname{Re}(z_{\mathfrak{c}}^2)\,\boldsymbol{\sigma}^2}{2}\right)\bar{\mu}\left(\exp\sum_{\mathfrak{c}}(\operatorname{Re}z_{\mathfrak{c}})\,X_{\mathfrak{c}}\right)\\
&\le\exp\left(-\sum_{\mathfrak{c}}\frac{\operatorname{Re}(z_{\mathfrak{c}}^2)\,\boldsymbol{\sigma}^2}{2}+\sum_{\mathfrak{c}}\frac{(\operatorname{Re}z_{\mathfrak{c}})^2\boldsymbol{\sigma}^2}{2}\right)\\
&=\exp\left(\sum_{\mathfrak{c}}\frac{(\operatorname{Im}t_{\mathfrak{c}})^2\boldsymbol{\sigma}^2}{2e^2}\right).
\end{align*}

\emph{The bound of $\exp\mathrm{Lg}$.}
Let $t\in\mathrm{DD}(t_2)$, $\theta\in\mathbb{R}^3$, $A\in\mathrm{Ksf}$, and write $\mathrm{Lg}(t):=\mathrm{Lg}(t;\theta,A)$.
By (g7), $t_2\le t_x/(8N_{*})\le t_x/2$.
By (Q1)\,(a) on $\mathrm{DD}(t_x/2)$, with $\mathrm{Lg}_{\emptyset}=\mathrm{Lg}(0)=0$, one has $\mathrm{Lg}=\sum_{\mathrm{Scl}\neq\emptyset}\mathrm{Lg}_{\mathrm{Scl}}$.
By the definition in Section~G5.1.6 and by Proposition~34.4\,(b),
\[
\mathrm{Lg}_{\{\mathfrak{c}\}}(t)=\mathrm{Lg}(t_{\mathfrak{c}}\,\mathrm{ue}_{\mathfrak{c}})-\mathrm{Lg}(0)=\mathrm{lgN}(t_{\mathfrak{c}})+\mathrm{wdf}_{\mathfrak{c}}(t_{\mathfrak{c}}).
\]
The function $\mathrm{lgN}$ is analytic on $|t|<t_x^{+}$ and even, $\mathrm{lgN}(0)=0$, and $|\mathrm{lgN}|\le B_0\le\frac14$ on $|t|\le t_x$ (Proposition~34.3\,(d)).
So $\mathrm{lgN}(t)/t^2$ is analytic, and the maximum principle on $|t|\le t_x$ gives $|\mathrm{lgN}(t)|\le B_0\,|t|^2/t_x^2$.
The function $\mathrm{wdf}_{\mathfrak{c}}$ is analytic on $|t|<t_x^{+}$, $\mathrm{wdf}_{\mathfrak{c}}(0)=0$, and $|\mathrm{wdf}_{\mathfrak{c}}|\le3e_{*}^{\frac{1}{4}-9\varepsilon}\le\frac38$ (Proposition~34.4\,(b); the third condition of (g3)).
Schwarz's lemma on $|t|\le t_x$ gives $|\mathrm{wdf}_{\mathfrak{c}}(t)|\le\frac38\,|t|/t_x$.
Since $|t_{\mathfrak{c}}|\le t_2\le t_x/(8N_{*})$,
\[
\sum_{\mathfrak{c}}\left|\mathrm{Lg}_{\{\mathfrak{c}\}}(t)\right|\le N_{*}\left[\frac{1}{4}\,(8N_{*})^{-2}+\frac{3}{8}\,(8N_{*})^{-1}\right]=\frac{1}{256\,N_{*}}+\frac{3}{64}.
\]
By (g6), $\bar{\kappa}_1^{1-s}\le1/(2N_{*})<1$; since $0<s<1$, it follows that $\bar{\kappa}_1<1$ and $\bar{\kappa}_1\le\bar{\kappa}_1^{1-s}\le1/(2N_{*})$.
There are fewer than $2^{N_{*}}$ sets $\mathrm{Scl}$.
By Proposition~35.6 and (g5),
\[
\sum_{\#|\mathrm{Scl}|\ge2}\left|\mathrm{Lg}_{\mathrm{Scl}}(t)\right|\le2^{N_{*}}\varpi_{\mathrm{sep}}^{+}\le2^{N_{*}}\,\frac{\bar{\kappa}_1^{N_{*}}}{2}=\frac{(2\bar{\kappa}_1)^{N_{*}}}{2}\le\frac{N_{*}^{-N_{*}}}{2}.
\]
Since $N_{*}=8L^3\ge8$, it follows that
\[
|\mathrm{Lg}(t)|\le\frac{1}{256\,N_{*}}+\frac{3}{64}+\frac{N_{*}^{-N_{*}}}{2}\le\frac12,\qquad|\exp\mathrm{Lg}(t)|\le\exp\left(\frac12\right)\le2 .
\]

\emph{The bound of $\mathrm{Ysf}$.}
By (d), (c) and the last bound, $|\mathrm{Ysf}(t)|\le\bar{\nu}(F_N)\cdot2\cdot\bar{\mu}(\mathfrak{G})\le2\,\bar{\nu}(F_N)$ for $t\in\mathrm{DD}(t_2)$.

\paragraph{3.9. Clause (h).}
The set $F_N$ is the disjoint union of $\mathfrak{G}$ and $F_N\setminus\mathfrak{G}$; hence
\[
\frac{\mathrm{zL}(t)}{\bar{\nu}(F_N)}=\mathring{\Phi}(t)-\frac{\mathrm{Ysf}(t)}{\bar{\nu}(F_N)},
\]
an entire function by (a).
On $\mathrm{DD}(t_2)$ one has
\[
\sum_{\mathfrak{c}}\frac{(\operatorname{Im}t_{\mathfrak{c}})^2\boldsymbol{\sigma}^2}{2e^2}\le\frac{N_{*}\cdot\mathrm{zq}^2\cdot\mathbf{D}}{2}\le\frac{\mathbf{D}}{2},
\]
and by (g), $|\mathrm{zL}(t)/\bar{\nu}(F_N)|\le\exp(\mathbf{D}/2)+2=\mathrm{mxL}$.
At the real points of $\mathrm{DD}(t_2)$, by (f), $|\mathrm{zL}(t)/\bar{\nu}(F_N)|\le\varpi_L$, and $\varpi_L=\varpi_0^{\frac{1}{2}}\exp(\mathbf{D}/2)\le\mathrm{mxL}$ because $\varpi_0\le\frac14$ by (g7).
Now (Q4) with $n=N_{*}$, $\mathrm{rad}=t_2$, $\mathrm{Mg}=\mathrm{mxL}$ and $\delta=\varpi_L$ gives
\[
\left|\frac{\mathrm{zL}(t)}{\bar{\nu}(F_N)}\right|\le\varpi_L^{\frac{1}{3}}\cdot\mathrm{mxL}^{\frac{2}{3}}\qquad\text{for }t\in\mathrm{DD}\left(\frac{t_2}{2N_{*}}\right).
\]

\paragraph{3.10. Clause (i).}
Write $\mathrm{Lr}=\sum_{\mathrm{ak}}\mathrm{Lr}_{\mathrm{ak}}\,t^{\mathrm{ak}}$, the sum extending over the multi-indices with $\mathrm{ak}_{\mathfrak{c}}\le4$, where the $\mathrm{Lr}_{\mathrm{ak}}$ are the Taylor coefficients at $0$ of $\mathrm{zL}/\bar{\nu}(F_N)$.
By Cauchy's estimates on $\mathrm{DD}\left(t_2/(2N_{*})\right)$ and by (h),
\[
|\mathrm{Lr}_{\mathrm{ak}}|\le\varpi_L^{\frac{1}{3}}\cdot\mathrm{mxL}^{\frac{2}{3}}\left(\frac{2N_{*}}{t_2}\right)^{|\mathrm{ak}|}.
\]
Put $\mathrm{yq}:=2N_{*}\mathbf{r}/t_2$.
By (g7), $\mathbf{r}/t_2=t_x/(4t_2)\ge2N_{*}$; hence $\mathrm{yq}\ge1$ and $\sum_{i=0}^{4}\mathrm{yq}^i\le5\,\mathrm{yq}^4$.
Therefore
\begin{align*}
\|\mathrm{Lr}\|_{\mathbf{r}}=\sum_{\mathrm{ak}}|\mathrm{Lr}_{\mathrm{ak}}|\,\mathbf{r}^{|\mathrm{ak}|}&\le\varpi_L^{\frac{1}{3}}\cdot\mathrm{mxL}^{\frac{2}{3}}\left(\sum_{i=0}^{4}\mathrm{yq}^i\right)^{N_{*}}\\
&\le\varpi_L^{\frac{1}{3}}\cdot\mathrm{mxL}^{\frac{2}{3}}\left(5\left(\frac{2N_{*}\mathbf{r}}{t_2}\right)^4\right)^{N_{*}}=\lambda_L^{N_{*}} .
\end{align*}
For real $t$ the cochain $a(t)$ is real, $\mathrm{dens}^{\theta,a(t)}_N\ge0$ (Proposition~31.1\,(a)), and $\mathrm{zL}(t)=\bar{\nu}^{a(t)}(F_N\setminus\mathfrak{G})$ is real.
By (h2) the coefficients $\mathrm{Lr}_{\mathrm{ak}}$ are real.
\qed

\subsection*{4. Proposition 35.8}

\subsubsection*{Proof of Proposition 35.8}

\begin{proof}
For $(\theta,A)\in\mathbb{R}^3\times\mathrm{Ksf}$ write $\mathrm{Lg}(t):=\mathrm{Lg}(t;\theta,A)$; this function is analytic in $t$ on $\mathrm{DDo}(t_x^{+})$ and continuous in $(t,\theta,A)$ (Proposition~34.1\,(d)).
By (g7), $t_2\le t_x/(8N_{*})<t_x^{+}$; so $\mathrm{DD}(t_2)$ is a compact subset of $\mathrm{DDo}(t_x^{+})$.

\paragraph{4.1. The identity (8.1).}
The set $F_N$ is the disjoint union of $\mathfrak{G}$ and $F_N\setminus\mathfrak{G}$.
Hence, by the definitions in Section~G5.1.5 and in Proposition~35.7,
\[
\mathring{\Phi}(t)=\frac{\mathrm{Ysf}(t)}{\bar{\nu}(F_N)}+\frac{\mathrm{zL}(t)}{\bar{\nu}(F_N)}\qquad\left(t\in\mathbb{C}^{\mathrm{Cel}}\right),
\]
and the three functions are entire (Proposition~35.7\,(a)).
The map which sends a function analytic near $0$ to its image in $\mathcal{J}$ is linear, and the image of $\mathrm{zL}/\bar{\nu}(F_N)$ is $\mathrm{Lr}$ (Proposition~35.7\,(i)).

\emph{The image of $\mathrm{Ysf}/\bar{\nu}(F_N)$.}
By Proposition~35.7\,(d),
\[
\frac{\mathrm{Ysf}(t)}{\bar{\nu}(F_N)}=\int\mathrm{vm}(\mathrm{d}\theta\,\mathrm{d}A)\,\exp\mathrm{Lg}(t;\theta,A)\qquad\text{on }\mathrm{DDo}(t_x^{+}).
\]
Apply (h1) with $n=N_{*}$, with $\mathrm{Po}$ the open polydisk of radius $\mathrm{rad}_0:=t_2$, with $\mathrm{Ms}=[-\pi,\pi]^3\times\mathrm{Ksf}$, $\mathrm{mq}=\mathrm{vm}$ (a finite positive measure by Proposition~35.7\,(c)), $\mathrm{gf}=\exp\mathrm{Lg}$ (continuous in $(t,\theta,A)$, holomorphic in $t$), and $\mathrm{maj}_{\mathrm{rad}}:=2$ (Proposition~35.7\,(g)).
Hence every Taylor coefficient at $0$ of $\mathrm{Ysf}/\bar{\nu}(F_N)$ is the integral with respect to $\mathrm{vm}$ of the same Taylor coefficient of $\exp\mathrm{Lg}(\,\cdot\,;\theta,A)$: the image of $\mathrm{Ysf}/\bar{\nu}(F_N)$ in $\mathcal{J}$ is the integral, coefficient by coefficient, of the image of $\exp\mathrm{Lg}(\,\cdot\,;\theta,A)$.

\emph{The image of $\exp\mathrm{Lg}$.}
By (Q1)\,(a) on $\mathrm{DD}(t_x/2)$, by $\mathrm{Lg}(0)=0$, by the definition of $\mathrm{gf}_{\mathrm{Scl}}$ in Section~G5.1.6 and by Proposition~34.4\,(b),
\begin{gather*}
\mathrm{Lg}(t)=\sum_{\mathfrak{c}\in\mathrm{Cel}}\left(\mathrm{lgN}(t_{\mathfrak{c}})+\mathrm{wdf}_{\mathfrak{c}}(t_{\mathfrak{c}})\right)+\sum_{\#|\mathrm{Scl}|\ge2}\mathrm{Lg}_{\mathrm{Scl}}(t),\\
\mathrm{Lg}_{\{\mathfrak{c}\}}(t)=\mathrm{Lg}(t_{\mathfrak{c}}\,\mathrm{ue}_{\mathfrak{c}})=\mathrm{lgN}(t_{\mathfrak{c}})+\mathrm{wdf}_{\mathfrak{c}}(t_{\mathfrak{c}}).
\end{gather*}
All summands are analytic near $t=0$.
Hence
\[
\exp\mathrm{Lg}(t)=\prod_{\mathfrak{c}}\exp\mathrm{lgN}(t_{\mathfrak{c}})\cdot\prod_{\mathfrak{c}}\exp\mathrm{wdf}_{\mathfrak{c}}(t_{\mathfrak{c}})\cdot\prod_{\#|\mathrm{Scl}|\ge2}\exp\mathrm{Lg}_{\mathrm{Scl}}(t).
\]
By Proposition~34.3\,(a), $\exp\mathrm{lgN}(t_{\mathfrak{c}})=N_0(t_{\mathfrak{c}})$; its image in $\mathcal{J}$ is its Taylor polynomial of degree $\le4$ in the single variable $t_{\mathfrak{c}}$, that is, $\phi_0(t_{\mathfrak{c}})$ (Proposition~34.3\,(e)).
Taking the image in $\mathcal{J}$ is a ring homomorphism (Section~G5.1.6).
With the definitions of the elements $\mathrm{kj}_{\mathrm{Scl}}$,
\[
\exp\mathrm{Lg}=\prod_{\mathfrak{c}\in\mathrm{Cel}}\phi_0(t_{\mathfrak{c}})\prod_{\mathrm{Scl}\neq\emptyset}\left(1+\mathrm{kj}_{\mathrm{Scl}}(\theta,A;\,\cdot\,)\right)\qquad\text{in }\mathcal{J}.
\]
The factor $\prod_{\mathfrak{c}}\phi_0(t_{\mathfrak{c}})$ does not contain $(\theta,A)$, and multiplication by a fixed element of $\mathcal{J}$ is linear in the coefficients; so this factor can be taken out of the integral.
This is (8.1).
The coefficients of the $\mathrm{kj}_{\mathrm{Scl}}$ are continuous on $\mathbb{R}^3\times\mathrm{Ksf}$ by clause (c) (see 4.3 below), and the coefficient at $t^{\mathrm{ak}}$ has modulus at most $\mathbf{r}^{-|\mathrm{ak}|}\,\|\mathrm{kj}_{\mathrm{Scl}}\|_{\mathbf{r}}\le\mathbf{r}^{-|\mathrm{ak}|}$, by the definition of the norm (Section~G5.1.6), by clauses (d) and (e) (see 4.4 and 4.5 below) and by $\bar{\kappa}_1\le\frac34$ (see (n2)).
The coefficients of $\prod_{\mathrm{Scl}}\left(1+\mathrm{kj}_{\mathrm{Scl}}\right)$ are polynomials in these numbers.
Hence they are bounded Borel functions on $[-\pi,\pi]^3\times\mathrm{Ksf}$, which carries $\mathrm{vm}$.

\paragraph{4.2. Clauses (a), (b), (f), (g).}
Clause (a) is contained in Proposition~34.3\,(e): $\phi_0$ is the $4$-jet of $N_0$, $\mathrm{nj}_0=1$, $\mathrm{nj}_1=\mathrm{nj}_3=0$, the $\mathrm{nj}_i$ are real, and (6.4) holds.
Clause (b) is Proposition~35.7\,(c).
Clause (f) is Proposition~35.7\,(i).
Clause (g) consists of the definitions $\mathbf{r}=t_x/4$ and $t_x=e_{*}^{3\varepsilon}/\mathrm{Cb}$ of Section~G5.1.4.

\paragraph{4.3. Clause (c).}
\emph{Monomials.}
Let $\#|\mathrm{Scl}|\ge2$.
The function $\mathrm{Lg}$ is analytic in the neighbourhood $\mathrm{DDo}(t_x^{+})$ of $\mathrm{DD}(t_x/2)$.
By (Q1)\,(d), every monomial of the Taylor series of $\mathrm{Lg}_{\mathrm{Scl}}$ at $0$ contains every $t_{\mathfrak{c}}$ with $\mathfrak{c}\in\mathrm{Scl}$, and no other variable.
The function $\mathrm{wdf}_{\mathfrak{c}}(t_{\mathfrak{c}})$ is an analytic function of the single variable $t_{\mathfrak{c}}$ which vanishes at $0$; so every monomial of its Taylor series is a power $t_{\mathfrak{c}}^i$ with $i\ge1$.
A product of monomials each of which contains every variable of $\mathrm{Scl}$ and no other variable has the same property.
Now
\[
\mathrm{kj}_{\mathrm{Scl}}=\sum_{i\ge1}\frac{1}{i!}\left(\text{image of }\mathrm{Lg}_{\mathrm{Scl}}\right)^i,
\]
respectively, for $\mathrm{Scl}=\{\mathfrak{c}\}$, the same expression with $\mathrm{wdf}_{\mathfrak{c}}$, computed in $\mathcal{J}$ (Section~G5.1.6), where the monomials with an exponent $\ge5$ are dropped.
Hence $\mathrm{kj}_{\mathrm{Scl}}$ has the property.

\emph{Continuity.}
The Taylor coefficient of index $\mathrm{ak}$ of $\mathrm{Lg}(\,\cdot\,;\theta,A)$ at $0$ is
\[
(2\pi i)^{-N_{*}}\oint\cdots\oint_{|\mathrm{zs}_{\mathfrak{c}}|=t_2}\mathrm{d}\mathrm{zs}\,\mathrm{Lg}(\mathrm{zs};\theta,A)\prod_{\mathfrak{c}}\mathrm{zs}_{\mathfrak{c}}^{-\mathrm{ak}_{\mathfrak{c}}-1}.
\]
The function $\mathrm{Lg}$ is continuous in $(\mathrm{zs},\theta,A)$, hence uniformly continuous on the product of the torus $\{|\mathrm{zs}_{\mathfrak{c}}|=t_2\}$ with a compact subset of $\mathbb{R}^3\times\mathrm{Ksf}$.
So the coefficient is continuous in $(\theta,A)$.
The Taylor coefficients of $\mathrm{Lg}_{\mathrm{Scl}}$ are those of $\mathrm{Lg}$ at the multi-indices that contain exactly the variables of $\mathrm{Scl}$, and $0$ at the others: by (Q1)\,(a), (d) the Taylor series of $\mathrm{Lg}$ is the sum of those of the $\mathrm{Lg}_{\mathrm{Scl}}$, and these contain pairwise different monomials.
The coefficients of $\mathrm{wdf}_{\mathfrak{c}}$ are those of $\mathrm{Lg}_{\{\mathfrak{c}\}}$ minus those of $\mathrm{lgN}$, which do not contain $(\theta,A)$.
The coefficients of $\mathrm{kj}_{\mathrm{Scl}}$ are polynomials in these numbers.

\emph{Reality.}
For real $t\in\mathrm{DDo}(t_x^{+})$ the number $\mathrm{Lg}(t;\theta,A)$ is real (Proposition~35.7\,(e)); so the Taylor coefficients of $\mathrm{Lg}$ at $0$ are real (by (h2)), and hence so are those of every $\mathrm{Lg}_{\mathrm{Scl}}$.
The function $\mathrm{lgN}$ is real for real $t$ (Proposition~34.3\,(d)); so its Taylor coefficients, and those of $\mathrm{wdf}_{\mathfrak{c}}$, are real.
The coefficients of $\mathrm{kj}_{\mathrm{Scl}}$ are polynomials with rational coefficients in these numbers.

\paragraph{4.4. Clause (d).}
The function $\mathrm{wdf}_{\mathfrak{c}}$ is analytic on $|t|<t_x^{+}$, a neighbourhood of $|t|\le t_x$, vanishes at $0$, and satisfies $|\mathrm{wdf}_{\mathfrak{c}}|\le3e_{*}^{\frac{1}{4}-9\varepsilon}$ (Proposition~34.4\,(b)).
By Cauchy's estimates on $|t|=t_x$, its coefficient of $t^i$ has modulus at most $3e_{*}^{\frac{1}{4}-9\varepsilon}\cdot t_x^{-i}$.
Since $\mathbf{r}/t_x=\frac14$, the third condition of (g3) gives
\[
\|\mathrm{wdf}_{\mathfrak{c}}\|_{\mathbf{r}}\le3e_{*}^{\frac{1}{4}-9\varepsilon}\sum_{i=1}^{4}4^{-i}\le3e_{*}^{\frac{1}{4}-9\varepsilon}\le\frac38 .
\]
The norm $\|\cdot\|_{\mathbf{r}}$ is submultiplicative (Section~G5.1.6), and $\exp(y)-1\le y\exp(y)\le2y$ for $0\le y\le\frac12$.
Hence
\[
\|\mathrm{kj}_{\{\mathfrak{c}\}}\|_{\mathbf{r}}\le\sum_{i\ge1}\frac{\|\mathrm{wdf}_{\mathfrak{c}}\|_{\mathbf{r}}^{i}}{i!}=\exp\left(\|\mathrm{wdf}_{\mathfrak{c}}\|_{\mathbf{r}}\right)-1\le2\,\|\mathrm{wdf}_{\mathfrak{c}}\|_{\mathbf{r}}\le6e_{*}^{\frac{1}{4}-9\varepsilon}=\bar{\kappa}_1 .
\]
The bound does not contain $(\theta,A)$.

\paragraph{4.5. Clause (e).}
Let $\#|\mathrm{Scl}|\ge2$.
By (Q1)\,(d) with $\mathrm{rad}=t_x/2$ and $\mathrm{rad}'=\mathbf{r}=t_x/4=\mathrm{rad}/2$, and by Proposition~35.6,
\[
\|\mathrm{Lg}_{\mathrm{Scl}}\|_{\mathbf{r}}\le\sup_{\mathrm{DD}(t_x/2)}|\mathrm{Lg}_{\mathrm{Scl}}|\le\varpi_{\mathrm{sep}}\le\varpi_{\mathrm{sep}}^{+}.
\]
By (g5) and (n2), $\varpi_{\mathrm{sep}}^{+}\le\bar{\kappa}_1^{N_{*}}/2\le\frac12$.
The norm is submultiplicative, and $\exp(y)-1\le2y$ for $0\le y\le\frac12$.
Hence
\[
\|\mathrm{kj}_{\mathrm{Scl}}\|_{\mathbf{r}}\le\sum_{i\ge1}\frac{\|\mathrm{Lg}_{\mathrm{Scl}}\|_{\mathbf{r}}^{i}}{i!}=\exp\left(\|\mathrm{Lg}_{\mathrm{Scl}}\|_{\mathbf{r}}\right)-1\le2\,\|\mathrm{Lg}_{\mathrm{Scl}}\|_{\mathbf{r}}\le2\varpi_{\mathrm{sep}}\le2\varpi_{\mathrm{sep}}^{+}\le\bar{\kappa}_1^{N_{*}} .
\]
The bounds do not contain $(\theta,A)$.

\paragraph{4.6. The last sentence.}
The numbers $a_{*}$, $e_{*}$, $v_{*}$, $r_{*}$, $\ell_*$, $p_{*}=v_{*}^{p}$ and $p_{0*}=v_{*}^{p_0}$ do not depend on $N\ge N_0(e)$ (Proposition~35.1\,(c)).
The numbers $C_{\beta}$ and $\mathbf{D}$ are suprema over $\mathrm{Ind}_{\beta}$, and $D_\beta$ is a number attached to $\operatorname{supp}\beta$ (Section~G5.1.3); they contain no $N$, no torus, no placement and no $\theta$, and they do not change if $\beta$ is replaced by $U_{r'}\beta$ with $r'\in\{\pm1\}^3$ (see 1.10).
By Sections~G5.1.4 and G5.1.8,
\begin{gather*}
\mathbf{r}=\frac{e_{*}^{3\varepsilon}}{4\,\mathrm{Cb}},\qquad\mathrm{Cb}=1+h^{\frac{1}{2}}C_{\beta},\qquad\bar{\kappa}_1=6e_{*}^{\frac{1}{4}-9\varepsilon},\\
\varpi_0=18\,\ell_*^3\exp\left(-\frac{p_{0*}}{2}\right)+6\,\ell_*^3\exp\left(-\frac{(p_{*}-1)^2}{2C_{0}^2}\right),\\
\varpi_{\mathrm{sep}}^{+}=2^{N_{*}}\,\mathrm{CWj}\cdot\mathrm{pp}\cdot\exp\left(-\mathrm{vt}\cdot\left(\frac{M\,r_{*}}{8L}-4M-7\right)\right),\\
\mathrm{pp}=\left(\frac{D_\beta}{a_{*}}+2M+4\right)^3,\qquad\lambda_L^{N_{*}}=\varpi_L^{\frac{1}{3}}\cdot\mathrm{mxL}^{\frac{2}{3}}\left(5\left(\frac{2N_{*}\mathbf{r}}{t_2}\right)^4\right)^{N_{*}},\\
\varpi_L=\varpi_0^{\frac{1}{2}}\exp\left(\frac{\mathbf{D}}{2}\right),\qquad\mathrm{mxL}=2+\exp\left(\frac{\mathbf{D}}{2}\right),\\
t_2=e\,\min\left\{N_{*}^{-\frac{1}{2}},\ C_{\beta}^{-1},\ \frac{1}{2}\left(N_{*}\cdot\mathbf{D}\right)^{-\frac{1}{2}}\right\}.
\end{gather*}
The bounds in (6.4) are $\mathbf{r}^{-2}/16$, $\mathbf{r}^{-4}/16$, $2$ and $1/6$.
All these numbers are functions of $e$, of $\bar{m}$, of $(C_{\beta},\mathbf{D},D_\beta)$, of the numbers named in the first sentence of 4.6 and of constants of the lines (C0)--(C7).
Hence they do not depend on $N$, on $\theta$, on the admitted placement, or on a replacement of $\beta$ by $U_r\beta$, $r\in\{\pm1\}^3$.
\end{proof}

\subsection*{5. Conditions on the constants}

The proofs of Propositions~35.5--35.8 use the following conditions on the constants.
All of them hold by the choice of the constants made in Section~I.8.2 (the lines (C0)--(C10)), in Proposition~14.2 and Section~G5.1.9, and in Section~G5.9.3 of Annex~B.

\begin{description}
\item[Conditions of line (C10).]
The following conditions are part of Proposition~14.2\,(d), line (C10) (the number $e_{\sharp}$), and of Section~G5.1.9.
\begin{itemize}
\item The first condition of (g1), $5\,\mathrm{hL}\cdot h_G\le e_{*}^{-\frac{1}{4}+\varepsilon}$, enters in 1.6\,(iv), and hence in 1.7--1.9.
\item The second condition of (g1), $\left(4+\frac{M^3}{2}\right)e_{*}^{\varepsilon}\le1$, is the condition (Z9.s1); it is used in 1.1.
\item The third condition of (g1), $\mathrm{alG}\cdot\bar{r}\le\frac{1}{2}$, is needed in 2.2.
\item The condition (g2), $\frac{M\,r_{*}}{8L}-1\ge6M+8$, enters in (n9) and in 1.9 and 2.3.
\item The third condition of (g3), $e_{*}^{\frac{1}{4}-9\varepsilon}\le\frac{1}{8}$, enters in (n2) and in 1.8, 1.9, 3.8, 4.1, 4.4 and 4.5.
\item The condition (g5), $2\varpi_{\mathrm{sep}}^{+}\le\bar{\kappa}_1^{N_{*}}$, is used in 3.8 and 4.5.
\item The condition (g6), $N_{*}\cdot\bar{\kappa}_1^{1-s}\le\frac{1}{2}$, is used in 3.8.
\item The condition (g7), that is, $\varpi_0\le\frac{1}{4}$ and $8N_{*}t_2\le t_x$, enters in 3.8, 3.9, 3.10 and Section~4.
\end{itemize}
\item[The first two clauses of (T1).]
They are the inequalities $v_{*}\ge8$ and $p_0\cdot v_{*}^{p_0-1}\ge28$ of Proposition~14.2\,(c); they follow from (e1), line (C9) (the number $e_1$), and hold by Proposition~35.1\,(e).
Both are needed in 3.2.
\item[The first condition of (A8).]
This is the inequality $e_0^{\varepsilon}\le\frac{1}{4}$ of Proposition~14.2\,(b), line (C8) (the number $e_0$).
Since $\varepsilon>0$, it yields $e_0<1$; together with $0<e_{*}\le e_0$ (Proposition~35.1\,(a)) this gives $e_{*}<1$, the form in which it is used in (n1).
\item[$\gamma^{F}M\ge16$.]
This inequality is obtained from (M3) together with $\gamma_{\mathrm{w}}\le\gamma^{F}$ (Proposition~14.2\,(b), (f), line (C5)); it is needed in 1.6\,(v).
\item[$M\ge2$; $M$ odd.]
The first follows from $M\ge L^2$; the second holds because $M$ is a power of the odd integer $L$ (Proposition~14.2\,(b): (M1), line (C5), and (L5), line (C3)).
The first is used in (5.4) and 1.8, the second in 1.6\,(v).
\item[$\kappa-1\ge\kappa_{0}$; $\kappa-\kappa_{0}\ge2$.]
Both follow from line (C2), $\kappa\ge7\kappa_{0}+7$.
The first is used in (5.5) and 1.8\,(A), the second in 2.3.
\item[$K_{0}\ge1$.]
By line (C0), $K_{0}=\exp(\kappa_{0})/36$ with $\kappa_{0}=54\log72$.
The inequality is used in 1.8 and 1.9.
\item[$C^{F}\ge1$; $0<\gamma^{F}\le1$.]
See Proposition~14.2\,(f) ($4\le C^{F}$; $\gamma^{F}L\le\frac{1}{6}$) and the lines (C6), (C4r).
These bounds enter in (n3) and (n5).
\item[$\mathrm{vt}\le\gamma^{F}/8$; $\mathrm{vt}\le1/M$.]
Both hold by the definition of $\mathrm{vt}$, line (C6$'$); they enter in (n4).
\item[Conditions of Section~G5.9.3.]
The following inequalities stand in Section~G5.9.3 of Annex~B, line (C6$'$).
\begin{itemize}
\item The inequality $C_K=226\,\mathbf{S}^2\le C_H\cdot\mathrm{sF}^2$ (with $226\le203M$) is needed in 1.8.
\item The two bounds
\[
8C_K+3K_{0}+1200\,C_H\cdot\mathrm{sF}\cdot K_{0}\le\mathrm{CVj},\qquad4C_K+82K_{0}+1200\,C_H\cdot\mathrm{sF}^2\cdot K_{0}\le\mathrm{CVj}
\]
are needed in 1.8.
\item The bound $\left(16h_G^2+256\right)\cdot C_H^2\cdot\mathrm{sF}^4\cdot K_{0}\le\mathrm{CVj}\cdot h_G^2$, with $\mathrm{CVj}=10^3\,C_H^2\cdot\mathrm{sF}^4\cdot K_{0}$ and $h_G\ge1$, is needed in 1.9.
\item The bound
\[
2\,\mathrm{CVj}\cdot h_G^2+2\,\mathrm{alG}\cdot\mathrm{CVj}^2\cdot h_G^4+2\,\mathrm{e}^4\cdot C_{d}\cdot M^3\cdot K_{0}\le\mathrm{CWj}
\]
is needed in 2.5.
\end{itemize}
\item[$C_{0}\ge1$.]
This is part of Proposition~14.2\,(f), line (C6); it is used in 3.2.
\item[$p_0$ and $p$ are integers $\ge1$; $0<s<1$.]
Both are imposed in line (C7); the first is used in 3.2, the second in 3.8.
\item[$\rho\ge0$.]
This is part of $\rho\in(0,1)$, the row $\rho$ of the line (C0) in Proposition~14.2\,(a) (Section~I.0.5).
It is used in 3.1 through (I.m6), in the form $|M_{i}(x)|_2=\frac{w_i}{2}(1+\rho)$.
\item[The standard weights.]
They are the row $w$ of the line (C0) in Proposition~14.2\,(a): $w=(w_0,w_1,w_2,w_3)=\left(6,6\sqrt{2},6\sqrt{2},6\right)$ and $\mathrm{sw}=w_0+w_1+w_2+w_3$ (Section~I.0.5).
They are used in 3.1.
The closure relation (I.1.1) enters (I.m1), from which the second identity of (I.2.5) follows, and this identity gives (7.0).
The inequalities $w_i>0$ enter through (I.m6).
The values $w_3=6$ and $\mathrm{sw}=\sum_iw_i$ enter the bound of the norm of $D_B+\boldsymbol{\mu}_N$ and the function $\mathrm{bnd}(t)$.
\item[{$\bar{m}\in(0,1]$; $h\ge1$.}]
The range of $\bar{m}$ and the definition $h=L^{\left\lfloor\log_L\frac{1}{\bar{m}}\right\rfloor}$ are those of Section~I.0.5; they are the rows $\bar{m}$ and $h$ of the line (C10) in Proposition~14.2\,(a), and $h\ge1$ stands in the paragraph ``Ranges'' of Proposition~14.2\,(f).
The mass $\bar{m}$ is a datum of the setting (ST*)$^{\mathrm{ch}}$ of Section~G5.1.1.
Both bounds on $\bar{m}$ are used in 3.1.
Since $\bar{m}\le1$, the exponent $\left\lfloor\log_L\frac{1}{\bar{m}}\right\rfloor$ is an integer $\ge0$; hence $h$ is an integral power of $L$, and $h\ge1$.
Since $\bar{m}>0$, one has $\bar{m}\,\eta>0$, which is used for $\det(D_0+\bar{m}\,\eta)>0$.
\item[$L\ge5$; $N_{*}=8L^3\ge8$.]
See Proposition~14.2\,(b): (L5), line (C3), and Proposition~14.2\,(a), line (C4).
The first enters in 3.2, through $\sum_{i\ge1}L^{-4i}\le1$, and in 3.1, through $n_*\ge5$: the number $n_*=h\,L^{T^{*}(e)+N}$ is an integer $\ge L\ge5$, by $L\ge5$ together with $h\ge1$ (an integral power of $L$, see the preceding entry), $T^{*}(e)\ge1$ and $N\ge0$.
That $T^{*}(e)$ and $N$ are integers with $T^{*}(e)\ge1$ and $N\ge0$ is a fact on the scales, not a condition on the constants (Section~I.6.2; Proposition~35.1\,(a), (c)).
The second enters in 3.8.
\item[$\varepsilon>0$.]
This is imposed in line (C7).
It enters in (n1), through $e_0<1$ (see the first condition of (A8)), and in (n2), 1.1 and 1.7, through the inequalities $e_{*}^{\frac{1}{4}-7\varepsilon}\le e_{*}^{\frac{1}{4}-9\varepsilon}$, $e_{*}^{-\frac{1}{4}+\varepsilon}\le e_{*}^{-\frac{1}{4}}$ and $e_{*}^{3\varepsilon}\le e_{*}^{2\varepsilon}$.
\end{description}

\section*{Chapter 35c. The bounds on the numbers $\Theta^{(0)}_m$ and the positivity of $e_{\sharp}(\beta)$: proofs of Propositions 35.9 and 35.10}

\subsection*{0. Conventions of this chapter}

Unless stated otherwise, the notation is that of Sections~I.6--I.8 and of Sections~G5.0.3, G5.1.1--G5.1.9 and G5.4.0.
On $\mathbb{R}^3$, $|\cdot|$ is the sup norm (Section~I.1.1).
Let $\mathrm{e}=\exp(1)$.
Since $i!\ge2\cdot3^{i-2}$ for $i\ge2$,
\[
2<\mathrm{e}=\sum_{i\ge0}\frac{1}{i!}\le2+\frac{1}{2}\sum_{i\ge0}3^{-i}=\frac{11}{4},
\]
so that $2<\mathrm{e}<3$.
For a real number $\mathrm{sv}\ge0$ put $\log^{+}\mathrm{sv}:=\max\{0,\log\mathrm{sv}\}$ if $\mathrm{sv}>0$, and $\log^{+}0:=0$; then $\mathrm{sv}\le\exp\left(\log^{+}\mathrm{sv}\right)$.
For every real $\mathrm{sv}$ one has $1+\mathrm{sv}\le\exp(\mathrm{sv})$; consequently $\log\mathrm{sv}\le\mathrm{sv}-1\le\mathrm{sv}$ for $\mathrm{sv}>0$.
The label (G5.1.1) refers to the inclusion $\mathrm{Qb}=v+Q_\beta\subset\left[\frac{\ell_c}{8},\frac{7\ell_c}{8}\right]^3$ for a placement $v\in\eta\mathbb{Z}^3$, displayed in Section~G5.1.3;
the labels (6.4) and (8.1) refer to the displays so labelled in Proposition~34.3 and in Proposition~35.8, respectively.

In this chapter $P$, $P'$, $P''$ denote elements of $\mathcal{J}$ (Section~G5.1.6), $\mathrm{sv}$ and $\mathrm{sv}'$ are real variables (as in Section~G5.0.3), and $\gamma$ is a rate (as in Section~G5.1.8); $\mathrm{cf}$, $\mathrm{cf}_0$, $\mathrm{cf}_1$, $\mathrm{cf}_2$ are real numbers and $\mathrm{dg}$ is an exponent.
The integers $N_1$, $N_2$, $i_1$, $i_2$ occur in 3.1 only, and $N_1$ has the meaning of Section~I.6.2.
The linear functionals $\mathrm{ell}_{\mathfrak{c}}$, $\mathrm{ell}$ and the families $\mathrm{Yf}$ of 2.1--2.5 play the roles of the functionals and of the families of sets of cells in the proof of \cite[Proposition~20.31]{QEDI}.
The labels (F1)--(F4) refer to the four facts established in 2.1, and not to the one-step bounds bearing the same labels in Section~I.8.5.

\subsection*{1. The range $0<e\le e_{\sharp}(\beta)$; centred placements; the two definitions of $\mathring{\Phi}$}

\begin{named}{Lemma 0}
Let the constants be those of Section~I.8, let $\bar{m}\in(0,1]$, let $\beta$ be a probe, and let $0<e\le e_{\sharp}(\beta)$.
Then $e\cdot h^{\frac{1}{2}}\le\frac{e_1}{2}$, and the conditions (g1)--(g9) hold at the charge $e$.
\end{named}

By 3.0 and 3.1, each of the conditions (g1)--(g9) is a condition on the number $e$ alone, so that the set in Section~G5.1.9 is well defined.

\begin{proof}
According to Section~G5.1.9, $e_{\sharp}(\beta)$ is one half of the supremum of the set of those $e'\in\left(0,e_1\cdot h^{-\frac{1}{2}}\right]$ for which (g1)--(g9) hold at every charge in $(0,e']$, and $e_{\sharp}(\beta)=0$ if this set is empty.
Since $e_{\sharp}(\beta)\ge e>0$, the set is not empty.
Its supremum $2\,e_{\sharp}(\beta)$ is at most $e_1\cdot h^{-\frac{1}{2}}$, whence $e\cdot h^{\frac{1}{2}}\le e_{\sharp}(\beta)\cdot h^{\frac{1}{2}}\le\frac{e_1}{2}$.
As $e_{\sharp}(\beta)>0$, this supremum exceeds $e_{\sharp}(\beta)$; hence the set contains a number $e'>e_{\sharp}(\beta)\ge e$.
The conditions (g1)--(g9) therefore hold at every charge in $(0,e']$, and in particular at $e$.
\end{proof}

\begin{named}{Lemma 1}
Assume the setting $(\mathrm{ST}^{*})^{\mathrm{ch}}$, let $\beta$ be a probe, and let $Q_\beta=y_\beta+\left[-\frac{D_\beta}{2},\frac{D_\beta}{2}\right]^3$ be its cube (Section~I.7.6).
\begin{itemize}
\item[(i)] There is a centred placement.
\item[(ii)] If $\frac{D_\beta}{2}+\frac{1}{2}\le\frac{3}{8}\,\ell_c$ (the first condition of (g9)), then every centred placement is admitted.
\item[(iii)] If $D_\beta+3\le\ell_c$ (the second condition of (g9)), then $D_\beta+3\eta\le\ell_c$ and $R^{*}(e)\ge D_{\mathrm{be}}+1$.
\end{itemize}
\end{named}

\begin{proof}
Since $N\ge0$ (Section~I.0.5) and $L\ge1$, one has $0<\eta=L^{-N}\le1$.

(i) Every real number lies at distance at most $\frac{\eta}{2}$ from $\eta\mathbb{Z}$.
Applying this to each of the three coordinates of $\frac{\ell_c}{2}(1,1,1)-y_\beta$, one finds $v\in\eta\mathbb{Z}^3$ with $\left|v+y_\beta-\frac{\ell_c}{2}(1,1,1)\right|\le\frac{\eta}{2}$.

(ii) Let $v$ be a centred placement, and put $y':=v+y_\beta$; then $\left|y'_\mu-\frac{\ell_c}{2}\right|\le\frac{\eta}{2}$ for $\mu=1,2,3$.
By Sections~G5.1.3 and I.7.6, $\mathrm{Qb}=v+Q_\beta=y'+\left[-\frac{D_\beta}{2},\frac{D_\beta}{2}\right]^3$.
For $y\in\mathrm{Qb}$ and $\mu=1,2,3$ it follows that
\[
\left|y_\mu-\frac{\ell_c}{2}\right|\le\left|y_\mu-y'_\mu\right|+\left|y'_\mu-\frac{\ell_c}{2}\right|\le\frac{D_\beta}{2}+\frac{\eta}{2}\le\frac{D_\beta}{2}+\frac{1}{2}\le\frac{3}{8}\,\ell_c,
\]
that is, $\frac{\ell_c}{8}\le y_\mu\le\frac{7\ell_c}{8}$.
This is the inclusion (G5.1.1) of Section~G5.1.3.

(iii) One has $D_\beta+3\eta\le D_\beta+3\le\ell_c$.
By Section~I.6.4, $\ell_c=\frac{R^{*}(e)}{2L}$; hence, as $D_\beta=D_{\mathrm{be}}+4$ (Section~G5.1.3),
\[
R^{*}(e)=2L\cdot\ell_c\ge\ell_c\ge D_\beta+3=D_{\mathrm{be}}+7\ge D_{\mathrm{be}}+1.
\]
\end{proof}

\begin{named}{Lemma 2}
Let data (DAT) be given (without a block field), the placement $v$ being centred.
The objects $a_\eta$, $\hat{\mathfrak{a}}$, $\boldsymbol{\sigma}^2$,$\hat{\mathfrak{a}}_{\mathfrak{c}}$, $X_{\mathfrak{c}}$, $N_0$, $c_m$ ($m=1,\dots,4$), $\mathring{\Phi}$ and $\Theta^{(0)}_m$ ($m=2,4$) of Sections~G5.1.3 and G5.1.5, formed with $v$, coincide with the objects bearing the same names in Section~I.7.6, formed with the placement $v$ on $C(n_*)$ at the lattice mass $\boldsymbol{\mu}_N$.
\end{named}

\begin{proof}
The condition $\left|v+y_\beta-\frac{\ell_c}{2}(1,1,1)\right|\le\frac{\eta}{2}$, $v\in\eta\mathbb{Z}^3$, defines the placements of Section~I.7.6 as well as the centred placements of Section~G5.1.3.
In both places $u_{\mu\nu}$ is the $R^{*}(e)\mathbb{Z}^3$-periodisation of the function $c\mapsto\left(J_\eta(\tau_v\beta)_{\mu\nu}\right)\left(c+\frac{\eta}{2}\left(e_\mu+e_\nu\right)\right)$, and
\[
a_\eta=d_\eta^{*}u,\qquad\hat{\mathfrak{a}}=\eta^{\frac{1}{2}}\cdot a_\eta,\qquad\boldsymbol{\sigma}^2=\|d\hat{\mathfrak{a}}\|^2,\qquad\hat{\mathfrak{a}}_{\mathfrak{c}}=(\mathsf{g}_{\mathfrak{c}})_{*}\hat{\mathfrak{a}},\qquad X_{\mathfrak{c}}=\langle d\hat{\mathfrak{a}}_{\mathfrak{c}},F^{\mathrm{u}}\rangle;
\]
the function $N_0$ and the numbers $c_m=N_0^{(m)}(0)$ are given by the same formula in both places.
Let $t\in\mathbb{C}^{\mathrm{Cel}}$.
By Section~G5.1.4, $a(t)=e^{-1}\sum_{\mathfrak{c}}t_{\mathfrak{c}}\cdot\hat{\mathfrak{a}}_{\mathfrak{c}}$; by $(\mathrm{ST}^{*})^{\mathrm{ch}}$, $(m_{\mathrm{in}},\varepsilon^{\mathrm{v}}_{\mathrm{in}})=P_N(e)=\left(m_{0,N}(e),\varepsilon^{\mathrm{v}}_{0,N}(e)\right)$ and $\bar{m}\cdot\eta+m_{0,N}(e)=\boldsymbol{\mu}_N$.
Therefore (I.7.2), together with the definitions of $\nu^{\theta,a}$ and $\bar{\nu}^{a}$ in Section~I.7.3 at $n=n_*$, gives
\begin{align*}
\bar{\nu}^{a(t)}(F_N)&=\int_{[-\pi,\pi]^3}\frac{\mathrm{d}\theta}{(2\pi)^3}\int\mathrm{d}\gamma_N(F^{\mathrm{u}})\,\mathrm{dens}^{\theta,a(t)}_N(F^{\mathrm{u}})\\
&=\exp\left(-\varepsilon^{\mathrm{v}}_{0,N}(e)\cdot n_*^3\right)\cdot\det\left(D_0+\bar{m}\cdot\eta\right)^{-1}\cdot\bar{Z}^{\mathrm{ext}}(t).
\end{align*}
The factor in front of $\bar{Z}^{\mathrm{ext}}(t)$ does not depend on $t$, and it is not zero, being the reciprocal of the number $\exp\left(\varepsilon^{\mathrm{v}}_{0,N}(e)\cdot n_*^3\right)\cdot\det\left(D_0+\bar{m}\cdot\eta\right)$.
At $t=0$ the left side is $\bar{\nu}(F_N)$, which lies in $(0,\infty)$ by Proposition~35.7\,(a), whose data (DAT) are given.
Hence $\bar{Z}^{\mathrm{ext}}(0)\ne0$, and
\[
\frac{\bar{\nu}^{a(t)}(F_N)}{\bar{\nu}(F_N)}=\frac{\bar{Z}^{\mathrm{ext}}(t)}{\bar{Z}^{\mathrm{ext}}(0)}.
\]
Finally, in both places $\Theta^{(0)}_m$ is obtained from $\mathring{\Phi}$ and $c_m$ by the same formula.
\end{proof}

\subsection*{2. Proposition 35.9}

Proposition~35.9 is stated in Part~II for the data (DAT) of Section~G5.4.0; the numbers $c_m$ and $\Theta^{(0)}_m$ are those of Section~G5.1.5.

\subsubsection*{Proof of Proposition 35.9}

Recall that $0<s<1$, $\varepsilon>0$ and $\sigma>0$ by (C7), that $e_1\le\frac{1}{3}$ by (C9), and that $L\ge1$ and $N\ge0$ (Section~I.0.5).

\paragraph{2.0. The data; four numbers.}
Let data (DAT) be given; a block field is not needed.
By Lemma~0 the conditions (g1)--(g9) hold at the charge $e$.
The numbers $\bar{\kappa}_1$, $\lambda_L$, $\ell_c$ occurring in these conditions are those formed with the $N$ of the data: $\bar{\kappa}_1$ and $\lambda_L^{N_*}$ do not depend on $N\ge N_0(e)$ by the last sentence of Proposition~35.8, and $\ell_c=\frac{R^{*}(e)}{2L}=\frac{h\cdot M\cdot r_K}{2L^2}$ does not contain $N$ (Sections~I.6.1, I.6.2, I.6.4).
The following facts hold throughout the proof.
First,
\begin{equation}
0<e_*\le e\cdot h^{\frac{1}{2}}\le e_1\le\frac{1}{3}
\tag{35.9.1}
\end{equation}
by Proposition~35.1\,(a) and (C9).
Next,
\[
\mathrm{Cb}=1+h^{\frac{1}{2}}\cdot C_\beta\in[1,\infty).
\]
Indeed, $C_\beta\ge0$, being a supremum of sums of norms over the set $\mathrm{Ind}_\beta$, which contains the triples $\left(N',\frac{h}{L},1\right)$ for all $N'\ge2$ (Section~G5.1.3); moreover $C_\beta\le2\,h^{\frac{1}{2}}\cdot\left(1+h+3h^2\right)\cdot|\beta|_{C^3}$ by Proposition~35.3\,(c1), which requires only $\bar{m}$ and $\beta$; and $|\beta|_{C^3}<\infty$, because $\beta$ is smooth with compact support.
Further,
\begin{equation}
\mathbf{r}=\frac{e_*^{3\varepsilon}}{4\,\mathrm{Cb}},\qquad\text{so that}\qquad0<\mathbf{r}\le\frac{1}{4}\qquad\text{and}\qquad\mathbf{r}^{-1}=4\,\mathrm{Cb}\cdot e_*^{-3\varepsilon},
\tag{35.9.2}
\end{equation}
by Proposition~35.8\,(g), (35.9.1), $\varepsilon>0$ and $1\le\mathrm{Cb}<\infty$.
Finally,
\begin{equation}
0<\bar{\kappa}_1<1,
\tag{35.9.3}
\end{equation}
since $\bar{\kappa}_1=6\,e_*^{\frac{1}{4}-9\varepsilon}>0$, while (g6) gives $\bar{\kappa}_1^{1-s}\le\frac{1}{2N_*}<1$ and $1-s>0$.

By Proposition~35.8\,(a), $\phi_0(t)=\sum_{i\le4}\mathrm{nj}_i\cdot t^i$ is the $4$-jet of $N_0$ at $0$; hence $\phi_0^{(m)}(0)=N_0^{(m)}(0)=c_m=m!\cdot\mathrm{nj}_m$ for $m=1,\dots,4$, and $\mathrm{nj}_0=\phi_0(0)=1$.
As $\phi_0$ is even, $c_1=c_3=0$ (Proposition~35.8\,(a)).
By Proposition~35.7\,(a) the function $t\mapsto\bar{\nu}^{a(t)}(F_N)$ is entire and $0<\bar{\nu}(F_N)<\infty$; so $\mathring{\Phi}$ is entire.
From now on $\mathring{\Phi}$ also denotes its image in $\mathcal{J}$ (Section~G5.1.6).

\paragraph{2.1. The functional $\mathrm{ell}$.}
Fix $m\in\{1,2,3,4\}$.
For a cell $\mathfrak{c}$ let $\mathrm{ell}_{\mathfrak{c}}$ be the linear functional on the span of $1,t_{\mathfrak{c}},\dots,t_{\mathfrak{c}}^4$ with
\[
\mathrm{ell}_{\mathfrak{c}}\left(t_{\mathfrak{c}}^i\right):=m!\cdot\mathbf{1}[i=m]-c_m\cdot\mathbf{1}[i=0]\qquad(0\le i\le4),
\]
and let $\mathrm{ell}$ be the linear functional on $\mathcal{J}$ with
\[
\mathrm{ell}\left(\prod_{\mathfrak{c}}t_{\mathfrak{c}}^{\mathrm{ak}_{\mathfrak{c}}}\right):=\prod_{\mathfrak{c}}\mathrm{ell}_{\mathfrak{c}}\left(t_{\mathfrak{c}}^{\mathrm{ak}_{\mathfrak{c}}}\right)
\]
for all multi-indices $\mathrm{ak}$ with $\mathrm{ak}_{\mathfrak{c}}\le4$; these monomials form a basis of $\mathcal{J}$.
\begin{itemize}
\item[(F1)] \emph{Let $\mathfrak{c}'$ be a cell, $P''$ an element of the span of $1,t_{\mathfrak{c}'},\dots,t_{\mathfrak{c}'}^4$, and $P'$ an element of $\mathcal{J}$ in which $t_{\mathfrak{c}'}$ does not occur.
Then $\mathrm{ell}(P''P')=\mathrm{ell}_{\mathfrak{c}'}(P'')\cdot\mathrm{ell}'(P')$, where $\mathrm{ell}'$ denotes the linear functional on the span of the monomials without $t_{\mathfrak{c}'}$ with $\mathrm{ell}'\left(\prod_{\mathfrak{c}\ne\mathfrak{c}'}t_{\mathfrak{c}}^{\mathrm{ak}_{\mathfrak{c}}}\right):=\prod_{\mathfrak{c}\ne\mathfrak{c}'}\mathrm{ell}_{\mathfrak{c}}\left(t_{\mathfrak{c}}^{\mathrm{ak}_{\mathfrak{c}}}\right)$.}
Both sides are bilinear in $(P'',P')$.
For $P''=t_{\mathfrak{c}'}^i$ and $P'=\prod_{\mathfrak{c}\ne\mathfrak{c}'}t_{\mathfrak{c}}^{\mathrm{ak}_{\mathfrak{c}}}$ the product $P''P'$ is the monomial with the exponent $i$ at $\mathfrak{c}'$ and the exponent $\mathrm{ak}_{\mathfrak{c}}$ at each $\mathfrak{c}\ne\mathfrak{c}'$; no exponent exceeds $4$, so that nothing is cut off, and the assertion is the definition of $\mathrm{ell}$.
\item[(F2)] \emph{If $\mathrm{gf}$ is analytic near $t=0$ and $P$ is its image in $\mathcal{J}$, then $\left.\prod_{\mathfrak{c}}\left(\partial_{t_{\mathfrak{c}}}^m-c_m\right)\mathrm{gf}(t)\right|_{t=0}=\mathrm{ell}(P)$.}
The operators $\partial_{t_{\mathfrak{c}}}$ commute with one another; hence
\[
\prod_{\mathfrak{c}}\left(\partial_{t_{\mathfrak{c}}}^m-c_m\right)=\sum_{\mathrm{Scl}\subset\mathrm{Cel}}\left(-c_m\right)^{N_*-\#|\mathrm{Scl}|}\prod_{\mathfrak{c}\in\mathrm{Scl}}\partial_{t_{\mathfrak{c}}}^m.
\]
Let $\mathrm{ak}(\mathrm{Scl})$ be the multi-index whose entries are $m$ for $\mathfrak{c}\in\mathrm{Scl}$ and $0$ for $\mathfrak{c}\notin\mathrm{Scl}$.
If $\mathrm{gf}_{\mathrm{ak}}$ are the Taylor coefficients of $\mathrm{gf}$ at $0$, then $\prod_{\mathfrak{c}\in\mathrm{Scl}}\partial_{t_{\mathfrak{c}}}^m\mathrm{gf}$ takes at $t=0$ the value $(m!)^{\#|\mathrm{Scl}|}\cdot\mathrm{gf}_{\mathrm{ak}(\mathrm{Scl})}$, and $\mathrm{gf}_{\mathrm{ak}(\mathrm{Scl})}=P_{\mathrm{ak}(\mathrm{Scl})}$, because $m\le4$.
On the other hand, since $m\ge1$, $\mathrm{ell}\left(t^{\mathrm{ak}}\right)$ vanishes unless every $\mathrm{ak}_{\mathfrak{c}}$ lies in $\{0,m\}$, that is, unless $\mathrm{ak}=\mathrm{ak}(\mathrm{Scl})$ with $\mathrm{Scl}=\{\mathfrak{c}:\mathrm{ak}_{\mathfrak{c}}=m\}$; and $\mathrm{ell}\left(t^{\mathrm{ak}(\mathrm{Scl})}\right)=(m!)^{\#|\mathrm{Scl}|}\cdot\left(-c_m\right)^{N_*-\#|\mathrm{Scl}|}$.
So both numbers are equal to
\[
\sum_{\mathrm{Scl}\subset\mathrm{Cel}}(m!)^{\#|\mathrm{Scl}|}\cdot\left(-c_m\right)^{N_*-\#|\mathrm{Scl}|}\cdot P_{\mathrm{ak}(\mathrm{Scl})}.
\]
\item[(F3)] \emph{For all $P\in\mathcal{J}$ one has $\left|\mathrm{ell}(P)\right|\le\left(m!\cdot\mathbf{r}^{-m}\right)^{N_*}\cdot\|P\|_{\mathbf{r}}$.}
For $0\le i\le4$, $\left|\mathrm{ell}_{\mathfrak{c}}\left(t_{\mathfrak{c}}^i\right)\right|\le m!\cdot\mathbf{r}^{-m}\cdot\mathbf{r}^i$.
Indeed, for $i=m$ both sides are equal to $m!$.
For $i=0$ the assertion reads $|c_m|\le m!\cdot\mathbf{r}^{-m}$: here $c_1=c_3=0$, and by (6.4) $|c_2|\le\frac{\mathbf{r}^{-2}}{16}\le2\,\mathbf{r}^{-2}$ and $|c_4|\le\frac{\mathbf{r}^{-4}}{16}\le24\,\mathbf{r}^{-4}$.
For the remaining $i$ the left side is $0$.
The product over the $N_*$ cells yields $\left|\mathrm{ell}\left(t^{\mathrm{ak}}\right)\right|\le\left(m!\cdot\mathbf{r}^{-m}\right)^{N_*}\cdot\mathbf{r}^{|\mathrm{ak}|}$; it remains to multiply by $|P_{\mathrm{ak}}|$ and to sum over $\mathrm{ak}$, since $\|P\|_{\mathbf{r}}=\sum_{\mathrm{ak}}|P_{\mathrm{ak}}|\cdot\mathbf{r}^{|\mathrm{ak}|}$ (Section~G5.1.6).
\item[(F4)] \emph{For every cell $\mathfrak{c}$ one has $\mathrm{ell}_{\mathfrak{c}}\left(\phi_0(t_{\mathfrak{c}})\right)=0$.}
Indeed, $\mathrm{ell}_{\mathfrak{c}}\left(\phi_0(t_{\mathfrak{c}})\right)=m!\cdot\mathrm{nj}_m-c_m\cdot\mathrm{nj}_0=c_m-c_m$.
\end{itemize}
Applying (F2) to $\mathrm{gf}=\mathring{\Phi}$ and recalling the definition of $\Theta^{(0)}_m$ (Section~G5.1.5), one obtains
\begin{equation}
\left(\Theta^{(0)}_m\right)^{N_*}=\left|\mathrm{ell}(\mathring{\Phi})\right|.
\tag{35.9.4}
\end{equation}

\paragraph{2.2. Expansion of (8.1).}
Call a family any set of non-empty subsets of $\mathrm{Cel}$; there are finitely many families, and the empty family is one of them.
For $(\theta,A)\in\mathbb{R}^3\times\{\chi_*=1\}$ one has, in $\mathcal{J}$,
\[
\prod_{\emptyset\ne\mathrm{Scl}\subset\mathrm{Cel}}\left(1+\mathrm{kj}_{\mathrm{Scl}}(\theta,A;\cdot)\right)=\sum_{\mathrm{Yf}}\prod_{\mathrm{Scl}\in\mathrm{Yf}}\mathrm{kj}_{\mathrm{Scl}}(\theta,A;\cdot),
\]
where $\mathrm{Yf}$ ranges over all families, the product being $1$ for $\mathrm{Yf}=\emptyset$.
The measure $\mathrm{vm}$ is finite and positive (Proposition~35.8\,(b)), and it is concentrated on $[-\pi,\pi]^3\times\{\chi_*=1\}$, because its density contains the factor $\chi_*(A)$ (the Definitions in Proposition~35.7).
The coefficients of $\mathrm{kj}_{\mathrm{Scl}}$ are continuous on $\mathbb{R}^3\times\{\chi_*=1\}$ (Proposition~35.8\,(c)), and the coefficient at $t^{\mathrm{ak}}$ has modulus at most $\mathbf{r}^{-|\mathrm{ak}|}\cdot\|\mathrm{kj}_{\mathrm{Scl}}\|_{\mathbf{r}}\le\mathbf{r}^{-|\mathrm{ak}|}$ (Proposition~35.8\,(d), (e), and (35.9.3)).
Hence the coefficients of all products of the $\mathrm{kj}_{\mathrm{Scl}}$ are bounded Borel functions on the set on which $\mathrm{vm}$ is concentrated.
The integral of a function with values in the finite-dimensional space $\mathcal{J}$ is formed coefficient by coefficient, and the map $P\mapsto\mathrm{ell}\left(\prod_{\mathfrak{c}}\phi_0(t_{\mathfrak{c}})\cdot P\right)$ is a linear functional on $\mathcal{J}$, that is, a finite linear combination of the coefficients of $P$.
Consequently (8.1) gives
\begin{align}
\mathrm{ell}(\mathring{\Phi})&=\sum_{\mathrm{Yf}}\int\mathrm{vm}(\mathrm{d}\theta\,\mathrm{d}A)\,\mathrm{ell}\left(\prod_{\mathfrak{c}}\phi_0(t_{\mathfrak{c}})\prod_{\mathrm{Scl}\in\mathrm{Yf}}\mathrm{kj}_{\mathrm{Scl}}(\theta,A;\cdot)\right)\notag\\
&\quad+\mathrm{ell}(\mathrm{Lr}).
\tag{35.9.5}
\end{align}

\paragraph{2.3. Families that do not cover $\mathrm{Cel}$.}
Let $\mathrm{Yf}$ be a family such that some cell $\mathfrak{c}'$ belongs to no member of $\mathrm{Yf}$; the empty family is of this kind.
By Proposition~35.8\,(c) the variable $t_{\mathfrak{c}'}$ occurs in no $\mathrm{kj}_{\mathrm{Scl}}$ with $\mathrm{Scl}\in\mathrm{Yf}$.
As a product in $\mathcal{J}$ of elements without $t_{\mathfrak{c}'}$ is an element without $t_{\mathfrak{c}'}$, the variable $t_{\mathfrak{c}'}$ does not occur in
\[
P':=\prod_{\mathfrak{c}\ne\mathfrak{c}'}\phi_0(t_{\mathfrak{c}})\prod_{\mathrm{Scl}\in\mathrm{Yf}}\mathrm{kj}_{\mathrm{Scl}}.
\]
By (F1) and (F4),
\[
\mathrm{ell}\left(\prod_{\mathfrak{c}}\phi_0(t_{\mathfrak{c}})\prod_{\mathrm{Scl}\in\mathrm{Yf}}\mathrm{kj}_{\mathrm{Scl}}\right)=\mathrm{ell}\left(\phi_0(t_{\mathfrak{c}'})\cdot P'\right)=\mathrm{ell}_{\mathfrak{c}'}\left(\phi_0(t_{\mathfrak{c}'})\right)\cdot\mathrm{ell}'(P')=0.
\]

\paragraph{2.4. Families that cover $\mathrm{Cel}$.}
For a family $\mathrm{Yf}$ let $\mathrm{Yf}_1$ be the set of its members with one element and $\mathrm{Yf}_2$ the set of its members with at least two elements, and put
\[
\mathrm{cnt}(\mathrm{Yf}):=\#|\mathrm{Yf}_1|+N_*\cdot\#|\mathrm{Yf}_2|.
\]
By Proposition~35.8\,(d), (e), for all $(\theta,A)\in\mathbb{R}^3\times\{\chi_*=1\}$ one has $\|\mathrm{kj}_{\mathrm{Scl}}\|_{\mathbf{r}}\le\bar{\kappa}_1$ if $\#|\mathrm{Scl}|=1$, and $\|\mathrm{kj}_{\mathrm{Scl}}\|_{\mathbf{r}}\le\bar{\kappa}_1^{N_*}$ if $\#|\mathrm{Scl}|\ge2$; the last inequality in Proposition~35.8\,(e) is the condition (g5).
The element $\phi_0(t_{\mathfrak{c}})$ of $\mathcal{J}$ has the norm $\sum_i|\mathrm{nj}_i|\cdot\mathbf{r}^i=\|\phi_0\|_{\mathbf{r}}\le2$ by (6.4), and the norm is submultiplicative (Section~G5.1.6).
Together with (F3) this gives, for every family $\mathrm{Yf}$ and all $(\theta,A)$,
\[
\left|\mathrm{ell}\left(\prod_{\mathfrak{c}}\phi_0(t_{\mathfrak{c}})\prod_{\mathrm{Scl}\in\mathrm{Yf}}\mathrm{kj}_{\mathrm{Scl}}\right)\right|\le\left(m!\cdot\mathbf{r}^{-m}\right)^{N_*}\cdot2^{N_*}\cdot\bar{\kappa}_1^{\mathrm{cnt}(\mathrm{Yf})}.
\]
Now let the union of the members of $\mathrm{Yf}$ be $\mathrm{Cel}$.
Then $\mathrm{cnt}(\mathrm{Yf})\ge N_*$.
Indeed, if $\mathrm{Yf}_2$ is not empty, then $\mathrm{cnt}(\mathrm{Yf})\ge N_*\cdot\#|\mathrm{Yf}_2|\ge N_*$; if $\mathrm{Yf}_2$ is empty, all members of $\mathrm{Yf}$ have one element and their union is $\mathrm{Cel}$, so that $\mathrm{Yf}_1$ consists of all the $N_*$ one-element subsets of $\mathrm{Cel}$.
By (35.9.3) and $0<s<1$ it follows that
\[
\bar{\kappa}_1^{\mathrm{cnt}(\mathrm{Yf})}=\bar{\kappa}_1^{s\cdot\mathrm{cnt}(\mathrm{Yf})}\cdot\bar{\kappa}_1^{(1-s)\cdot\mathrm{cnt}(\mathrm{Yf})}\le\bar{\kappa}_1^{s\cdot N_*}\cdot\bar{\kappa}_1^{(1-s)\cdot\mathrm{cnt}(\mathrm{Yf})}.
\]

\paragraph{2.5. The sum over the families.}
The number $\mathrm{cnt}$ is additive: $\mathrm{cnt}(\mathrm{Yf})=\sum_{\mathrm{Scl}\in\mathrm{Yf}}\mathrm{cnt}(\{\mathrm{Scl}\})$, where $\mathrm{cnt}(\{\mathrm{Scl}\})=1$ if $\#|\mathrm{Scl}|=1$ and $\mathrm{cnt}(\{\mathrm{Scl}\})=N_*$ if $\#|\mathrm{Scl}|\ge2$.
Multiplying out the product $\prod_{\emptyset\ne\mathrm{Scl}\subset\mathrm{Cel}}\left(1+\bar{\kappa}_1^{(1-s)\cdot\mathrm{cnt}(\{\mathrm{Scl}\})}\right)$, one therefore obtains
\[
\sum_{\mathrm{Yf}}\prod_{\mathrm{Scl}\in\mathrm{Yf}}\bar{\kappa}_1^{(1-s)\cdot\mathrm{cnt}(\{\mathrm{Scl}\})}=\sum_{\mathrm{Yf}}\bar{\kappa}_1^{(1-s)\cdot\mathrm{cnt}(\mathrm{Yf})},
\]
the sum over all families.
There are $N_*$ subsets of $\mathrm{Cel}$ with one element and $2^{N_*}-1-N_*$ subsets with at least two elements.
Hence
\begin{align}
\sum_{\mathrm{Yf}}\bar{\kappa}_1^{(1-s)\cdot\mathrm{cnt}(\mathrm{Yf})}&=\left(1+\bar{\kappa}_1^{1-s}\right)^{N_*}\cdot\left(1+\bar{\kappa}_1^{(1-s)\cdot N_*}\right)^{2^{N_*}-1-N_*}\notag\\
&\le\exp\left(N_*\cdot\bar{\kappa}_1^{1-s}+2^{N_*}\cdot\bar{\kappa}_1^{(1-s)\cdot N_*}\right)\le\mathrm{e}.
\tag{35.9.6}
\end{align}
For the last step note that $N_*\cdot\bar{\kappa}_1^{1-s}\le\frac{1}{2}$ is (g6), and that it implies $2^{N_*}\cdot\left(\bar{\kappa}_1^{1-s}\right)^{N_*}\le2^{N_*}\cdot\left(2N_*\right)^{-N_*}=N_*^{-N_*}\le\frac{1}{2}$, because $N_*=8L^3\ge2$.

\paragraph{2.6. Proof of (a).}
The total mass of $\mathrm{vm}$ is at most $1$ (Proposition~35.8\,(b)), and by (F3) and Proposition~35.8\,(f)
\[
\left|\mathrm{ell}(\mathrm{Lr})\right|\le\left(m!\cdot\mathbf{r}^{-m}\right)^{N_*}\cdot\|\mathrm{Lr}\|_{\mathbf{r}}\le\left(m!\cdot\mathbf{r}^{-m}\right)^{N_*}\cdot\lambda_L^{N_*}.
\]
By (35.9.5) and 2.3 only the families that cover $\mathrm{Cel}$ contribute to $\mathrm{ell}(\mathring{\Phi})$, and with 2.4, (35.9.6) and (35.9.4) it follows that
\[
\left(\Theta^{(0)}_m\right)^{N_*}\le\left(m!\cdot\mathbf{r}^{-m}\right)^{N_*}\cdot\left[\mathrm{e}\cdot\left(2\,\bar{\kappa}_1^{s}\right)^{N_*}+\lambda_L^{N_*}\right].
\]
For real numbers $\mathrm{sv},\mathrm{sv}'\ge0$ one has $\left(\mathrm{sv}+\mathrm{sv}'\right)^{\frac{1}{N_*}}\le\mathrm{sv}^{\frac{1}{N_*}}+\mathrm{sv}'^{\frac{1}{N_*}}$: the $N_*$-th power of the right side, expanded by the binomial formula, contains the terms $\mathrm{sv}$ and $\mathrm{sv}'$, and its other terms are $\ge0$.
Therefore
\[
\Theta^{(0)}_m\le m!\cdot\mathbf{r}^{-m}\cdot\left(2\,\mathrm{e}^{\frac{1}{N_*}}\cdot\bar{\kappa}_1^{s}+\lambda_L\right)\le m!\cdot\mathbf{r}^{-m}\cdot\left(2\,\mathrm{e}\cdot\bar{\kappa}_1^{s}+\lambda_L\right),
\]
which is the first inequality of (a).
By (g8), $\lambda_L\le\bar{\kappa}_1^{s}$, and $2\,\mathrm{e}+1<7\le8$; this is the second inequality of (a).

\paragraph{2.7. Proof of (b).}
By (35.9.2), $\mathbf{r}^{-4}=256\,\mathrm{Cb}^4\cdot e_*^{-12\varepsilon}$, and since $0<s<1$, $\bar{\kappa}_1^{s}=6^s\cdot e_*^{s\cdot\left(\frac{1}{4}-9\varepsilon\right)}\le6\,e_*^{s\cdot\left(\frac{1}{4}-9\varepsilon\right)}$.
With $\mathfrak{e}=s\cdot\left(\frac{1}{4}-9\varepsilon\right)-12\varepsilon$ (Section~I.8.5) this gives
\begin{equation}
\mathbf{r}^{-4}\cdot\bar{\kappa}_1^{s}\le256\cdot6\,\mathrm{Cb}^4\cdot e_*^{s\cdot\left(\frac{1}{4}-9\varepsilon\right)-12\varepsilon}=1536\,\mathrm{Cb}^4\cdot e_*^{\mathfrak{e}}.
\tag{35.9.7}
\end{equation}
At $m=4$, (a) yields
\[
\Theta^{(0)}_4\le8\cdot24\,\mathbf{r}^{-4}\cdot\bar{\kappa}_1^{s}\le192\cdot1536\,\mathrm{Cb}^4\cdot e_*^{\mathfrak{e}}=294912\,\mathrm{Cb}^4\cdot e_*^{\mathfrak{e}}\le3\cdot10^5\,\mathrm{Cb}^4\cdot e_*^{\mathfrak{e}}.
\]
At $m=2$, (a) yields $\Theta^{(0)}_2\le8\cdot2\,\mathbf{r}^{-2}\cdot\bar{\kappa}_1^{s}=16\,\mathbf{r}^{-2}\cdot\bar{\kappa}_1^{s}$.
By (35.9.2), $1\le\mathbf{r}^{-2}$, and by (6.4), $|c_2|\le\frac{\mathbf{r}^{-2}}{16}$; so $1+|c_2|\le\frac{17}{16}\,\mathbf{r}^{-2}$.
Multiplying the two bounds, one finds
\[
\left(1+|c_2|\right)\cdot\Theta^{(0)}_2\le17\,\mathbf{r}^{-4}\cdot\bar{\kappa}_1^{s}\le17\cdot1536\,\mathrm{Cb}^4\cdot e_*^{\mathfrak{e}}=26112\,\mathrm{Cb}^4\cdot e_*^{\mathfrak{e}}\le3\cdot10^4\,\mathrm{Cb}^4\cdot e_*^{\mathfrak{e}}.
\]

\emph{From $e_*$ to $e$.}
By (C7), $\sigma=\mathfrak{e}-\frac{2a_L}{N_*}>0$, and $a_L=\frac{L^3-1}{(L-1)^3}>0$; hence $\mathfrak{e}>0$.
Moreover $\mathfrak{e}=s\cdot\left(\frac{1}{4}-9\varepsilon\right)-12\varepsilon<\frac{1}{4}$, because $0<s<1$ and $\varepsilon>0$.
Since $\bar{m}\le1$ (Section~I.0.5), $h=L^{\lfloor\log_L\frac{1}{\bar{m}}\rfloor}\ge1$.
By (35.9.1), $0<e_*\le e\cdot h^{\frac{1}{2}}$.
Hence
\[
e_*^{\mathfrak{e}}\le\left(e\cdot h^{\frac{1}{2}}\right)^{\mathfrak{e}}=h^{\frac{\mathfrak{e}}{2}}\cdot e^{\mathfrak{e}}\le h^{\frac{1}{8}}\cdot e^{\mathfrak{e}}.
\]

\emph{An exponent.}
One has $s\cdot\left(\frac{1}{4}-9\varepsilon\right)=\mathfrak{e}+12\varepsilon>0$, because $\mathfrak{e}>0$ and $\varepsilon>0$; as $s>0$, it follows that
\begin{equation}
\frac{1}{4}-9\varepsilon>0.
\tag{35.9.8}
\end{equation}

\paragraph{2.8. Proof of (c).}
By (6.4) and (35.9.2),
\[
|c_2|+|c_4|^{\frac{1}{2}}\le\frac{\mathbf{r}^{-2}}{16}+\frac{\mathbf{r}^{-2}}{4}=\frac{5}{16}\,\mathbf{r}^{-2}\le\frac{\mathbf{r}^{-2}}{2}=8\,\mathrm{Cb}^2\cdot e_*^{-6\varepsilon}.
\]
By Proposition~35.1\,(a), $e_*\ge e\cdot\left(\frac{h}{L^2}\right)^{\frac{1}{2}}>0$, and the exponent $-6\varepsilon$ is negative; therefore
\[
e_*^{-6\varepsilon}\le\left(e\cdot\left(\frac{h}{L^2}\right)^{\frac{1}{2}}\right)^{-6\varepsilon}=\left(\frac{L^2}{h}\right)^{3\varepsilon}\cdot e^{-6\varepsilon}.
\]

\paragraph{2.9. The last sentence of the statement.}
Let $0<e\le e_{\sharp}(\beta)$ and $N\ge N_0(e)$.
By Lemma~0, $e\cdot h^{\frac{1}{2}}\le e_1$, so that the setting $(\mathrm{ST}^{*})^{\mathrm{ch}}$ is given.
By Lemma~0 the first condition of (g9) holds; recall from 2.0 that $\ell_c$ does not contain $N$.
By Lemma~1\,(i), (ii), centred placements exist, and every centred placement $v$ is admitted.
Thus $(\mathrm{ST}^{*})^{\mathrm{ch}}$, $\beta$, $e$, $v$ are data (DAT), and (a), (b), (c) hold for them.
By Lemma~2, applied to these data, the objects $\mathring{\Phi}$, $N_0$, $c_m$ and $\Theta^{(0)}_m$ for $m=2,4$ are then those of Section~I.7.6.

\paragraph{2.10. Uniformity.}
By the last sentence of Proposition~35.8, the numbers $\mathbf{r}$, $\bar{\kappa}_1$ and $\lambda_L^{N_*}$ depend neither on $N\ge N_0(e)$, nor on $\theta$, nor on the admitted placement, nor on a replacement of $\beta$ by $U_r\beta$, $r\in\{\pm1\}^3$.
The same is true of $\lambda_L=\left(\lambda_L^{N_*}\right)^{\frac{1}{N_*}}$ (recall that $\lambda_L\ge0$), of $e_*=\left(\frac{\bar{\kappa}_1}{6}\right)^{\left(\frac{1}{4}-9\varepsilon\right)^{-1}}$ (the exponent exists, since $\frac{1}{4}-9\varepsilon\ne0$ by (35.9.8)), and of $\mathrm{Cb}=\frac{e_*^{3\varepsilon}}{4\,\mathbf{r}}$.
The right sides of (a) are formed from $m$, $s$, $\mathbf{r}$, $\bar{\kappa}_1$, $\lambda_L$, and those of (b) and (c) from $\mathrm{Cb}$, $h$, $L$, $\varepsilon$, $\mathfrak{e}$ and $e_*$, respectively $e$.
Consequently all the bounds of Proposition~35.9 hold with the same right sides for all $N\ge N_0(e)$ and all admitted placements; neither a twist nor a block field enters them.
The factors $3\cdot10^5\,\mathrm{Cb}^4\cdot h^{\frac{1}{8}}$, $3\cdot10^4\,\mathrm{Cb}^4\cdot h^{\frac{1}{8}}$ and $8\,\mathrm{Cb}^2\cdot\left(\frac{L^2}{h}\right)^{3\varepsilon}$ in front of $e^{\mathfrak{e}}$ and $e^{-6\varepsilon}$ are finite, because $1\le\mathrm{Cb}<\infty$ (2.0).
They contain neither $e$ nor $N$, because $\mathrm{Cb}=1+h^{\frac{1}{2}}\cdot C_\beta$ and $C_\beta$ depends on $(\beta,h,L,\alpha)$ only (Proposition~35.3\,(c1)).
Hence they are the same for all $e\in\left(0,e_{\sharp}(\beta)\right]$.
\qed

\subsection*{3. Proposition 35.10}

In Proposition~35.10 the numbers $C_\beta$, $\mathbf{D}$, $\mathrm{Cb}$ are those of Section~G5.1.3, the numbers $t_x$, $t_x^{+}$, $\mathbf{r}$, $q_x$, $\mathrm{pp}$ those of Section~G5.1.4, the line (C10) and the quantities that depend on $e$ those of Section~G5.1.8, and the conditions (g1)--(g9) and the number $e_{\sharp}(\beta)$ those of Section~G5.1.9.

\subsubsection*{Proof of Proposition 35.10}

Recall that $r\ge2$, $p_0\ge2$, $p>2p_0$, $0<\varepsilon\le\frac{1}{200}$ and $0<s<1$ by (C0) and (C7), and that $0<e_1\le\frac{1}{3}$ by (C9), the positivity of $e_1$ being Proposition~14.3\,(f).

\paragraph{3.0. The constants; the numbers of the line (C10).}
\emph{The constants of the lines (C0)--(C6$'$) that occur in (g1)--(g9).}
The number $L$ is an odd integer $\ge5$, and $M$ is a power of $L$ with $M\ge L^2$ (Section~I.0.5).
The numbers $K_0=\frac{\exp(\kappa_0)}{36}$ with $\kappa_0=54\log72$, $N_*=8L^3$, $\mu_-=L^{-3}$ and $b_-=1-\frac{1}{L}$ (Section~I.8.2, lines (C0), (C4); Section~G5.1.8) are finite and $>0$.
The numbers $C^F$, $C_0$, $C_d$ belong to the line (C6), and $\gamma^F$ belongs to the line (C4r) (Section~I.8.2).
By (G5.T-1) in Section~G5.7 (Annex~A), $C^F\ge1$, $0<\gamma^F\le1$, $C_0\ge\max(1,C^F)$ and $C_d<\infty$.
By the paragraph ``Ranges'' of Proposition~14.2\,(f), every entry of the rows $\rho,\dots,\sigma$ of the table in Proposition~14.2\,(a), apart from $C_E$, $\mathrm{Csy}_{\mathrm{ak}}$ and $\mathrm{Cpv}_0$, is defined and finite; in particular $C^F<\infty$, $C_0<\infty$ and $C_w<\infty$.
Moreover $C_d>0$: by Proposition~14.2\,(a),
\[
C_d=16\,c_X\cdot C_w\cdot\max\left\{\frac{\pi}{2}\,b\cdot\mathrm{e}\cdot\left(1+2\,\mathrm{e}\right),\ 2+58\,\mathrm{e}+b\cdot\mathrm{e}\right\},
\]
a product of numbers $>0$, since $c_X=54$, $b=1$ and $C_w\ge C^F$ (Proposition~14.2\,(f)).
For $\gamma>0$ the number $\mathrm{ssum}(\gamma)=\left(\frac{12}{\gamma}\right)^3$ lies in $(0,\infty)$.
Hence, by the formulas of the line (C6$'$) (Section~G5.1.8), read in their order, the numbers
\begin{gather*}
\mathbf{S}=C^F\cdot\mathrm{ssum}\left(\frac{\gamma^F}{2}\right),\qquad\mathrm{hL}=6\,\mathbf{S},\qquad C_K=226\,\mathbf{S}^2,\qquad C_K'=1328\,\mathbf{S}^2,\\
C_H=203\,M\cdot\left(C^F\right)^2,\qquad\mathrm{CVj}=10^3\,C_H^2\cdot\mathrm{ssum}\left(\frac{\gamma^F}{2}\right)^4\cdot K_0,\\
\theta_{\mathrm{CT}}=\frac{\mu_-}{12\,\mathrm{e}},\qquad\mathrm{vt}=\min\left\{\frac{\theta_{\mathrm{CT}}}{2},\frac{\gamma^F}{8},\frac{1}{M}\right\},\\
\mathrm{gGa}=\left(b_-^{-1}+\mu_-^{-1}\right)^{\frac{1}{2}},\qquad h_{G}=1+\mathrm{gGa},\qquad\mathrm{gGa}_{1}=b_-^{-1}+\frac{4}{\mu_-},\\
\mathrm{alG}=2\,b_-^{-1}+4\,\mathrm{e}\cdot\mu_-^{-1}\cdot\mathrm{ssum}\left(\frac{\theta_{\mathrm{CT}}}{2}\right),\qquad c_{\tau}=\min\left\{\frac{1}{8},\frac{\mu_-}{48},\frac{1}{292\,\mathbf{S}}\right\},\\
\mathrm{CWj}=3\,\mathrm{CVj}\cdot h_{G}^2+2\,\mathrm{alG}\cdot\mathrm{CVj}^2\cdot h_{G}^4+4\,\mathrm{e}^6\cdot C_d\cdot M^3\cdot K_0
\end{gather*}
are finite and $>0$.

\emph{The numbers of the line (C10).}
One has $h=L^{\lfloor\log_L\frac{1}{\bar{m}}\rfloor}$, and $\log_L h=\lfloor\log_L\frac{1}{\bar{m}}\rfloor$ is an integer $\ge0$, because $\bar{m}\le1$ (Section~I.0.5); so $h\ge1$.
The support of $\beta$ is compact and not empty; so $0\le D_{\mathrm{be}}<\infty$, and $D_\beta=D_{\mathrm{be}}+4$ lies in $[4,\infty)$ (Section~G5.1.3).
The numbers $C_\beta$ and $\mathbf{D}$ are suprema of numbers $\ge0$ over the set $\mathrm{Ind}_\beta$, which is not empty: since $h\ge1$, it contains $\left(N',\frac{h}{L},1\right)$ for all $N'\ge2$.
By Proposition~35.3\,(c1), which requires only $\bar{m}$ and $\beta$, they are finite: $0\le C_\beta<\infty$ and $0\le\mathbf{D}<\infty$.
Hence their formulas (Sections~G5.1.3, G5.1.8) give
\begin{gather*}
\mathrm{Cb}=1+h^{\frac{1}{2}}\cdot C_\beta\in[1,\infty),\qquad\mathrm{mxL}=2+\exp\left(\frac{\mathbf{D}}{2}\right)\in[3,\infty),\\
C_{\mathrm{fr}}=2592\,C_d\cdot K_0\cdot C_\beta^2\cdot\left(M^2\cdot D_\beta^3\cdot L^4\cdot h^{-2}+8\,M^5\cdot h\right)\in[0,\infty),
\end{gather*}
and
\[
\frac{t_2}{e}=\min\left\{N_*^{-\frac{1}{2}},\ \frac{1}{C_\beta},\ \frac{1}{2\sqrt{N_*\cdot\mathbf{D}}}\right\},
\]
where an entry with the denominator $0$ is read as $+\infty$; this number does not contain $e$, and $0<\frac{t_2}{e}\le N_*^{-\frac{1}{2}}$.
Finally put
\[
\mathrm{pp}^{+}:=\left(\frac{D_\beta\cdot L^2}{h}+2M+4\right)^3\in[1,\infty).
\]

\paragraph{3.1. The scales as functions of $e$; six bounds.}
Let $e>0$ with $e\cdot h^{\frac{1}{2}}\le e_1$.
By Sections~I.6.1, I.6.2 and I.6.4, the numbers
\begin{gather*}
v_K=-\log\left(e\cdot h^{\frac{1}{2}}\right),\qquad r_K=\min\left\{L^i:\ i\ge0\text{ an integer},\ L^i\ge v_K^r\right\},\\
R^{*}(e)=\frac{h\cdot M\cdot r_K}{L},\qquad\ell_c=\frac{R^{*}(e)}{2L}
\end{gather*}
are formed from $(e,h;L,M,r)$.
For an integer $N\ge0$ and a level $k$, the numbers
\begin{gather*}
K=N+\log_L h,\qquad a_k=L^{k-N},\qquad e_k=e\cdot a_k^{\frac{1}{2}},\qquad v_k=-\log e_k,\\
r_k=\min\left\{L^i:\ i\ge0\text{ an integer},\ L^i\ge v_k^r\right\},\qquad\ell_k=\frac{R^{*}(e)}{a_k},\\
\mathfrak{n}_k=\log_L\frac{\ell_k}{M\cdot r_k},\qquad p_k=v_k^p,\qquad p_{0,k}=v_k^{p_0}
\end{gather*}
are formed from $(e,h,N,k;L,M,r,p_0,p)$.
The level $k_{*}$ is the first $k$ with $\mathfrak{n}_k=0$, and $N_0(e)=N_{\min}\left(R^{*}(e),e\right)$ is the least integer $N_1\ge0$ such that $R^{*}(e)\cdot L^N\ge L\cdot M\cdot r_0$ for all integers $N\ge N_1$; here $r_0$ is the number $r_k$ at $k=0$, which is formed with $N$.
No other quantity enters these definitions, and $\bar{m}$ enters through $h$ only.

\emph{The reference torus lies in $\mathrm{Tor}(e)$, and $N_0(e)$ is finite.}
By Section~I.0.5, $M=L^{i_1}$ with an integer $i_1\ge2$, and $r_K=L^{i_2}$ with an integer $i_2\ge0$.
Hence $R^{*}(e)=h\cdot L^{i_1+i_2-1}$, where $i_1+i_2-1\ge1$ is an integer, and $R^{*}(e)=\frac{R_{\mathrm{top}}(e)}{L}\le R_{\mathrm{top}}(e)$; thus $R^{*}(e)$ lies in $\mathrm{Tor}(e)$, and $N_{\min}\left(R^{*}(e),e\right)$ is formed according to Section~I.6.2.
By (C9), $e\cdot h^{\frac{1}{2}}\le e_1\le\frac{1}{3}$; so $v_K\ge\log3>1$, as $\mathrm{e}<3$.
Let $N\ge0$ be an integer.
With $\mathrm{lg}=\frac{1}{2}\log L>0$ (Section~I.6.1) one has $v_K=-\log e-\frac{1}{2}\log h=-\log e-\left(\log_L h\right)\cdot\mathrm{lg}$, and $\log_L h\ge0$ (3.0).
At $k=0$ the numbers $e_k$ and $v_k$ are $q=e\cdot\eta^{\frac{1}{2}}=e\cdot L^{-\frac{N}{2}}$ and $-\log q$ (Section~I.6.1).
Therefore
\[
-\log q=-\log\left(e\cdot L^{-\frac{N}{2}}\right)=-\log e+N\cdot\mathrm{lg}=v_K+\left(N+\log_L h\right)\cdot\mathrm{lg}\ge v_K>1.
\]
Consequently $\left(-\log q\right)^r>1=L^0$.
Hence $r_0=L^i$ with an integer $i\ge1$, and $L^{i-1}<\left(-\log q\right)^r$, because $i$ is the least exponent with $L^i\ge\left(-\log q\right)^r$.
This gives
\[
r_0<L\cdot\left(-\log q\right)^r=L\cdot\left(v_K+\left(N+\log_L h\right)\cdot\mathrm{lg}\right)^r.
\]
The numbers $L^{3-N}\cdot\left(v_K+\left(N+\log_L h\right)\cdot\mathrm{lg}\right)^r$, $N=0,1,2,\dots$, are $>0$.
The quotient of the number at $N+1$ by the number at $N$ equals
\[
L^{-1}\cdot\left(1+\frac{\mathrm{lg}}{v_K+\left(N+\log_L h\right)\cdot\mathrm{lg}}\right)^r;
\]
it tends to $\frac{1}{L}$ as $N\to\infty$, and $\frac{1}{L}<\frac{1}{2}\left(1+\frac{1}{L}\right)<1$, because $L>1$.
So there is an integer $N_2\ge0$ such that this quotient is at most $\frac{1}{2}\left(1+\frac{1}{L}\right)$ for all $N\ge N_2$, and for $N\ge N_2$ the number at $N$ is then at most $\left(\frac{1}{2}\left(1+\frac{1}{L}\right)\right)^{N-N_2}$ times the number at $N_2$.
These numbers therefore tend to $0$, and there is an integer $N_1\ge0$ with $\left(v_K+\left(N+\log_L h\right)\cdot\mathrm{lg}\right)^r\le L^{N-3}$ for all $N\ge N_1$.
For all $N\ge N_1$, since $h\ge1$ and $r_K\ge1$,
\[
L\cdot M\cdot r_0<L^2\cdot M\cdot\left(v_K+\left(N+\log_L h\right)\cdot\mathrm{lg}\right)^r\le M\cdot L^{N-1}\le h\cdot M\cdot r_K\cdot L^{N-1}=R^{*}(e)\cdot L^N.
\]
So the set of the integers $N_1\ge0$ in the definition of $N_{\min}\left(R^{*}(e),e\right)$ is not empty, and its least element $N_0(e)$ is a finite number, formed from $(e,h;L,M,r)$.

\emph{The scales at the stop level.}
Integers $N\ge N_0(e)$ therefore exist, and for each of them $\bar{m}$, $e$, $N$ and $R=R^{*}(e)$ satisfy the hypotheses of Proposition~35.1.
By Proposition~35.1\,(a) the level $k_{*}$ exists; it equals $K-1$ or $K-2$.
By Proposition~35.1\,(c) the numbers $a_*$, $e_*$, $v_*$, $r_*$, $\ell_*$, $p_*$, $p_{0*}$ do not depend on $N\ge N_0(e)$; so they coincide with the numbers formed with $N=N_0(e)$.
Hence
\begin{equation}
v_K,\ r_K,\ R^{*}(e),\ \ell_c,\ N_0(e),\ a_*,\ e_*,\ v_*,\ r_*,\ \ell_*,\ p_*,\ p_{0*}\quad\text{are functions of}\quad(e,h;L,M,r,p_0,p).
\tag{35.10.1}
\end{equation}
The quantities $\mathrm{pp}$, $t_x$, $\mathbf{r}$ (Section~G5.1.4) and $\tau_s$, $\bar{\kappa}_1$, $\bar{r}$, $B_1$, $B_{\mathrm{wr}}$, $B_0$, $\varpi_{\mathrm{sep}}^{+}$, $\varpi_0$, $t_2$, $\varpi_L$, $\lambda_L$ (Section~G5.1.8) are formed from the numbers (35.10.1), from $e$, from the numbers of 3.0 and from constants of the lines (C0)--(C7), (C6$'$).
Hence, once the constants, $\bar{m}$ and $\beta$ are fixed, each of the nineteen inequalities (g1)--(g9) is a condition on the number $e\in\left(0,e_1\cdot h^{-\frac{1}{2}}\right]$ alone, and $e_{\sharp}(\beta)$ is defined by Section~G5.1.9.
The following bounds hold for every $e>0$ with $e\cdot h^{\frac{1}{2}}\le e_1$; in them Proposition~35.1 is applied with an arbitrary integer $N\ge N_0(e)$.
By Proposition~35.1\,(a), (b),
\begin{equation}
8\le v_K\le v_*,\qquad e_*=\exp(-v_*)\in(0,1).
\tag{35.10.2}
\end{equation}
By Proposition~35.1\,(b), and since $r\ge2$, $p_0\ge2$ (C7) and $v_*\ge1$,
\begin{equation}
v_*^2\le v_*^r\le r_*,\qquad M\cdot v_*^2\le M\cdot v_*^r\le\ell_*<L\cdot M\cdot v_*^r,\qquad v_*^2\le v_*^{p_0}=p_{0*}.
\tag{35.10.3}
\end{equation}
By Proposition~35.1\,(a), $e_*=e\cdot a_*^{\frac{1}{2}}$ with $a_*\in\left\{\frac{h}{L},\frac{h}{L^2}\right\}$; hence
\begin{equation}
\left(\frac{L}{h}\right)^{\frac{1}{2}}\cdot e_*\le e\le L\cdot h^{-\frac{1}{2}}\cdot e_*.
\tag{35.10.4}
\end{equation}
Since $\mathrm{pp}=\left(\frac{D_\beta}{a_*}+2M+4\right)^3$ and $a_*\ge\frac{h}{L^2}$,
\begin{equation}
\mathrm{pp}\le\mathrm{pp}^{+}.
\tag{35.10.5}
\end{equation}
By Proposition~35.1\,(f), and since $r\ge2$ and $v_K\ge1$,
\begin{equation}
\ell_c\ge\frac{h\cdot M\cdot v_K^r}{2L^2}\ge\frac{h\cdot M\cdot v_K^2}{2L^2}.
\tag{35.10.6}
\end{equation}
Finally,
\begin{equation}
\exp\left(-\frac{v_*}{4}\right)\le\bar{\kappa}_1=6\,\exp\left(-\left(\frac{1}{4}-9\varepsilon\right)\cdot v_*\right).
\tag{35.10.7}
\end{equation}

\emph{Exponents.}
The lines (C0) and (C7) give $0<\varepsilon\le\bar{\varepsilon}=\frac{1}{200}$ and $0<s<1$.
Hence
\begin{gather*}
\frac{1}{4}-9\varepsilon\ge\frac{41}{200}>0,\qquad
\frac{1}{4}-\varepsilon>0,\qquad
\frac{1}{4}-7\varepsilon\ge\frac{43}{200}>2\varepsilon,\\
1-3\varepsilon>0,\qquad
(1-s)\left(\frac{1}{4}-9\varepsilon\right)>0.
\end{gather*}
The left inequality in (35.10.7) is a consequence of $6\ge1$ and $0<\frac{1}{4}-9\varepsilon\le\frac{1}{4}$.

\paragraph{3.2. Three elementary inequalities.}

\begin{named}{Lemma A}
For real numbers $\mathrm{cf}\ge0$, $\mathrm{dg}\ge0$ and $\gamma>0$ put
\[
\mathrm{thr}(\mathrm{cf},\mathrm{dg},\gamma):=
\begin{cases}
\max\left\{0,\ \dfrac{2}{\gamma}\cdot\log\left(\mathrm{cf}\cdot\left(\dfrac{2\,\mathrm{dg}}{\gamma\cdot\mathrm{e}}\right)^{\mathrm{dg}}\right)\right\} & \text{if }\mathrm{cf}>0,\\[2mm]
0 & \text{if }\mathrm{cf}=0,
\end{cases}
\]
where $\left(\frac{2\,\mathrm{dg}}{\gamma\cdot\mathrm{e}}\right)^{\mathrm{dg}}:=1$ for $\mathrm{dg}=0$.
\begin{itemize}
\item[(i)] $\mathrm{sv}^{\mathrm{dg}}\cdot\exp(-\gamma\cdot\mathrm{sv})\le\left(\frac{\mathrm{dg}}{\gamma\cdot\mathrm{e}}\right)^{\mathrm{dg}}$ for every $\mathrm{sv}>0$; for $\mathrm{dg}=0$ the right side is $1$.
\item[(ii)] $\mathrm{cf}\cdot\mathrm{sv}^{\mathrm{dg}}\cdot\exp(-\gamma\cdot\mathrm{sv})\le1$ for every $\mathrm{sv}>0$ with $\mathrm{sv}\ge\mathrm{thr}(\mathrm{cf},\mathrm{dg},\gamma)$.
\item[(iii)] If $\mathrm{cf}_2>0$, $\mathrm{cf}_1\ge0$, $\mathrm{cf}_0\ge0$ and $\mathrm{sv}\ge\max\left\{1,\frac{\mathrm{cf}_1+\mathrm{cf}_0}{\mathrm{cf}_2}\right\}$, then $\mathrm{cf}_2\cdot\mathrm{sv}^2\ge\mathrm{cf}_1\cdot\mathrm{sv}+\mathrm{cf}_0$.
\end{itemize}
\end{named}

\begin{proof}[Proof of Lemma A]
(i) For $\mathrm{dg}=0$ the left side equals $\exp(-\gamma\cdot\mathrm{sv})\le1$.
Let $\mathrm{dg}>0$.
Multiplying the inequality $\log\left(\frac{\gamma\cdot\mathrm{sv}}{\mathrm{dg}}\right)\le\frac{\gamma\cdot\mathrm{sv}}{\mathrm{dg}}-1$ by $\mathrm{dg}$, one finds
\[
\mathrm{dg}\cdot\log\mathrm{sv}\le\mathrm{dg}\cdot\log\frac{\mathrm{dg}}{\gamma}+\gamma\cdot\mathrm{sv}-\mathrm{dg},
\]
and the exponential of both sides gives (i).

(ii) For $\mathrm{cf}=0$ the left side vanishes.
Let $\mathrm{cf}>0$.
Part (i), with $\gamma/2$ in place of $\gamma$, gives
\begin{align*}
\mathrm{cf}\cdot\mathrm{sv}^{\mathrm{dg}}\cdot\exp(-\gamma\cdot\mathrm{sv})
&=\mathrm{cf}\cdot\left[\mathrm{sv}^{\mathrm{dg}}\cdot\exp\left(-\frac{\gamma}{2}\cdot\mathrm{sv}\right)\right]\cdot\exp\left(-\frac{\gamma}{2}\cdot\mathrm{sv}\right)\\
&\le\mathrm{cf}\cdot\left(\frac{2\,\mathrm{dg}}{\gamma\cdot\mathrm{e}}\right)^{\mathrm{dg}}\cdot\exp\left(-\frac{\gamma}{2}\cdot\mathrm{sv}\right).
\end{align*}
The right side is at most $1$, because $\mathrm{sv}\ge\mathrm{thr}(\mathrm{cf},\mathrm{dg},\gamma)$ implies $\frac{\gamma}{2}\cdot\mathrm{sv}\ge\log\left(\mathrm{cf}\cdot\left(\frac{2\,\mathrm{dg}}{\gamma\cdot\mathrm{e}}\right)^{\mathrm{dg}}\right)$.

(iii) The hypotheses give $\mathrm{cf}_2\cdot\mathrm{sv}\ge\mathrm{cf}_1+\mathrm{cf}_0$ and $\mathrm{sv}\ge1$; therefore
\[
\mathrm{cf}_2\cdot\mathrm{sv}^2=(\mathrm{cf}_2\cdot\mathrm{sv})\cdot\mathrm{sv}\ge(\mathrm{cf}_1+\mathrm{cf}_0)\cdot\mathrm{sv}\ge\mathrm{cf}_1\cdot\mathrm{sv}+\mathrm{cf}_0.
\]
\end{proof}

\paragraph{3.3. The conditions, one by one.}
Let $e>0$ with $e\cdot h^{\frac{1}{2}}\le e_1$.
The 21 numbers $\mathrm{vq}_{1a},\dots,\mathrm{vq}_9$ introduced below are formed only from the constants of the lines (C0)--(C7), (C6$'$) and from $h$, $D_\beta$, $C_{\beta}$, $\mathbf{D}$.
Each of them is finite: the first arguments of $\mathrm{thr}$ are finite and $\ge0$ by 3.0, the third arguments are $>0$ by the paragraph \emph{Exponents} of 3.1, and $\mu_-$, $c_{\tau}$, $\mathrm{vt}$, $M$, $C_{0}$, $t_2/e$ are $>0$ by 3.0.
Lemma~A is always applied with $\mathrm{sv}=v_{*}$, which is positive by (35.10.2).

\emph{(g1), first inequality.}
Put $\mathrm{vq}_{1a}:=\mathrm{thr}\left(5\,\mathrm{hL}\cdot h_{G},\,0,\,\frac{1}{4}-\varepsilon\right)$.
If $v_{*}\ge\mathrm{vq}_{1a}$, Lemma~A\,(ii) yields
\[
5\,\mathrm{hL}\cdot h_{G}\cdot e_{*}^{\frac{1}{4}-\varepsilon}=5\,\mathrm{hL}\cdot h_{G}\cdot\exp\left(-\left(\frac{1}{4}-\varepsilon\right)v_{*}\right)\le1.
\]

\emph{(g1), second inequality.}
With $\mathrm{vq}_{1b}:=\mathrm{thr}\left(4+\frac{M^3}{2},\,0,\,\varepsilon\right)$ and $v_{*}\ge\mathrm{vq}_{1b}$, Lemma~A\,(ii) gives
\[
\left(4+\frac{M^3}{2}\right)e_{*}^{\varepsilon}=\left(4+\frac{M^3}{2}\right)\exp(-\varepsilon\cdot v_{*})\le1.
\]

\emph{(g1), third inequality.}
Let $\mathrm{vq}_{1c}:=\mathrm{thr}\left(2\,\mathrm{alG}\cdot\left(4h_{G}^2\cdot C_K+650\,K_{0}\right),\,0,\,2\varepsilon\right)$.
Because $0<e_{*}<1$ and $\frac{1}{4}-7\varepsilon>2\varepsilon$, one has $e_{*}^{\frac{1}{4}-7\varepsilon}\le e_{*}^{2\varepsilon}$; so the definition of $\bar{r}$ (Section~G5.1.8) gives
\begin{align*}
\bar{r}&=4h_{G}^2\cdot C_K\cdot e_{*}^{2\varepsilon}+650\,K_{0}\cdot e_{*}^{\frac{1}{4}-7\varepsilon}\\
&\le\left(4h_{G}^2\cdot C_K+650\,K_{0}\right)\exp(-2\varepsilon\cdot v_{*}).
\end{align*}
For $v_{*}\ge\mathrm{vq}_{1c}$ it follows from Lemma~A\,(ii) that $2\,\mathrm{alG}\cdot\bar{r}\le1$.

\emph{(g2).}
Put $\mathrm{vq}_2:=(120L)^{\frac{1}{2}}$ and let $v_{*}\ge\mathrm{vq}_2$.
By (35.10.3) and $M\ge1$,
\[
\frac{M\cdot r_{*}}{8L}\ge\frac{M\cdot v_{*}^2}{8L}\ge15M\ge6M+9,
\]
which is the inequality $\frac{M\cdot r_{*}}{8L}-1\ge6M+8$.

\emph{(g3), first inequality.}
Let $\mathrm{vq}_{3a}:=\mathrm{thr}\left(4C_{0}\cdot\mathrm{gGa}_{1}\cdot(L\cdot M)^3,\,3r,\,\frac{1}{4}-\varepsilon\right)$.
By (35.10.3), $\ell_*^3<(L\cdot M)^3\cdot v_{*}^{3r}$.
If $v_{*}\ge\mathrm{vq}_{3a}$, Lemma~A\,(ii) therefore gives
\begin{align*}
&4C_{0}\cdot\ell_*^3\cdot\mathrm{gGa}_{1}\cdot e_{*}^{\frac{1}{4}-\varepsilon}\\
&\qquad\le4C_{0}\cdot\mathrm{gGa}_{1}\cdot(L\cdot M)^3\cdot v_{*}^{3r}\cdot\exp\left(-\left(\frac{1}{4}-\varepsilon\right)v_{*}\right)\le1.
\end{align*}

\emph{(g3), second inequality.}
Let $\mathrm{vq}_{3b}:=\mathrm{thr}\left(4\,(2C_K+K_{0}+4)\cdot(L\cdot M)^3,\,3r,\,2\varepsilon\right)$.
For $v_{*}\ge\mathrm{vq}_{3b}$ the same bound on $\ell_*^3$ and Lemma~A\,(ii) give
\begin{align*}
&4\,\ell_*^3\cdot(2C_K+K_{0}+4)\cdot e_{*}^{2\varepsilon}\\
&\qquad\le4\,(2C_K+K_{0}+4)\cdot(L\cdot M)^3\cdot v_{*}^{3r}\cdot\exp(-2\varepsilon\cdot v_{*})\le1.
\end{align*}

\emph{(g3), third inequality.}
For $\mathrm{vq}_{3c}:=\mathrm{thr}\left(8,\,0,\,\frac{1}{4}-9\varepsilon\right)$ and $v_{*}\ge\mathrm{vq}_{3c}$, Lemma~A\,(ii) gives
\[
8\,e_{*}^{\frac{1}{4}-9\varepsilon}=8\exp\left(-\left(\frac{1}{4}-9\varepsilon\right)v_{*}\right)\le1.
\]

\emph{(g4), first inequality.}
Let $\mathrm{vq}_{4a}:=\mathrm{thr}\left(\frac{48}{\mu_-},\,0,\,2\varepsilon\right)$.
If $v_{*}\ge\mathrm{vq}_{4a}$, then $\frac{48}{\mu_-}\cdot e_{*}^{2\varepsilon}\le1$ by Lemma~A\,(ii).

\emph{(g4), second inequality.}
Let $\mathrm{vq}_{4b}:=\mathrm{thr}\left(4\,\mathrm{gGa}_{1}\cdot C_K',\,0,\,2\varepsilon\right)$.
If $v_{*}\ge\mathrm{vq}_{4b}$, Lemma~A\,(ii) gives $4\,\mathrm{gGa}_{1}\cdot C_K'\cdot e_{*}^{2\varepsilon}\le1$, that is, $2\,\mathrm{gGa}_{1}\cdot C_K'\cdot e_{*}^{2\varepsilon}\le\frac{1}{2}$.

\emph{(g4), third inequality.}
Let $\mathrm{vq}_{4c}:=(c_{\tau}\cdot M)^{-\frac{1}{2}}$.
If $v_{*}\ge\mathrm{vq}_{4c}$, then (35.10.3) gives $\tau_s=c_{\tau}\cdot\ell_*\ge c_{\tau}\cdot M\cdot v_{*}^2\ge1$.

\emph{(g4), fourth inequality.}
Let
\[
\mathrm{vq}_{4d}:=\max\left\{1,\ \frac{\frac{1}{4}+\log^{+}\left(832\,\mathrm{gGa}_{1}\cdot C_K'\cdot\mathrm{pp}^{+}\right)}{c_{\tau}\cdot M}\right\}.
\]
By (35.10.5) and $e_{*}<1$,
\[
52\,B_1=832\,\mathrm{gGa}_{1}\cdot C_K'\cdot\mathrm{pp}\cdot e_{*}^{2\varepsilon}\le832\,\mathrm{gGa}_{1}\cdot C_K'\cdot\mathrm{pp}^{+},
\]
and by (35.10.3), $\tau_s\ge c_{\tau}\cdot M\cdot v_{*}^2$.
Let $v_{*}\ge\mathrm{vq}_{4d}$.
Lemma~A\,(iii), applied with $\mathrm{cf}_2=c_{\tau}\cdot M$, $\mathrm{cf}_1=\frac{1}{4}$ and $\mathrm{cf}_0=\log^{+}\left(832\,\mathrm{gGa}_{1}\cdot C_K'\cdot\mathrm{pp}^{+}\right)$, gives
\[
c_{\tau}\cdot M\cdot v_{*}^2\ge\frac{v_{*}}{4}+\log^{+}\left(832\,\mathrm{gGa}_{1}\cdot C_K'\cdot\mathrm{pp}^{+}\right).
\]
Since $\mathrm{sv}\le\exp\left(\log^{+}\mathrm{sv}\right)$ for every $\mathrm{sv}\ge0$, it follows that
\begin{align*}
52\,B_1\cdot\exp(-\tau_s)
&\le832\,\mathrm{gGa}_{1}\cdot C_K'\cdot\mathrm{pp}^{+}\cdot\exp\left(-c_{\tau}\cdot M\cdot v_{*}^2\right)\\
&\le\exp\left(-\frac{v_{*}}{4}\right)\le\exp\left(-\left(\frac{1}{4}-9\varepsilon\right)v_{*}\right)=e_{*}^{\frac{1}{4}-9\varepsilon}.
\end{align*}

\emph{(g4), fifth inequality.}
Put
\begin{align*}
\mathrm{vq}_{4e}&:=\mathrm{thr}\left(12\,C_{\mathrm{fr}},\,0,\,4\varepsilon\right),\\
\mathrm{vq}_{4f}&:=\mathrm{thr}\left(192\,\mathrm{e}^6\cdot C_{d}\cdot M^3\cdot K_{0}\cdot\mathrm{pp}^{+},\,0,\,1\right),\\
\mathrm{vq}_{4g}&:=\mathrm{thr}\left(192\,\mathrm{gGa}_{1}\cdot C_K'\cdot\mathrm{pp}^{+},\,0,\,2\varepsilon\right),
\end{align*}
and let $v_{*}\ge\max\left\{\mathrm{vq}_{4e},\mathrm{vq}_{4f},\mathrm{vq}_{4g}\right\}$.
Lemma~A\,(ii) gives $C_{\mathrm{fr}}\cdot e_{*}^{4\varepsilon}\le\frac{1}{12}$.
By (35.10.3) and $v_{*}\ge1$ one has $r_{*}\ge v_{*}^2\ge v_{*}$; together with (35.10.5) and Lemma~A\,(ii) this gives
\[
B_{\mathrm{wr}}\le16\,\mathrm{e}^6\cdot C_{d}\cdot M^3\cdot K_{0}\cdot\mathrm{pp}^{+}\cdot\exp(-v_{*})\le\frac{1}{12}.
\]
Again by (35.10.5) and Lemma~A\,(ii),
\[
B_1\le16\,\mathrm{gGa}_{1}\cdot C_K'\cdot\mathrm{pp}^{+}\cdot\exp(-2\varepsilon\cdot v_{*})\le\frac{1}{12}.
\]
Adding the three bounds one obtains $B_0\le\frac{1}{4}$.

\emph{(g5).}
Let
\begin{multline*}
\mathrm{vq}_5:=\max\Biggl\{1,\ \frac{8L}{\mathrm{vt}\cdot M}\\
\cdot\left(\frac{N_{*}}{4}+\log^{+}\left(2^{N_{*}+1}\cdot\mathrm{CWj}\cdot\mathrm{pp}^{+}\right)+\mathrm{vt}\cdot(4M+7)\right)\Biggr\}.
\end{multline*}
By (35.10.5) and (35.10.3),
\begin{align*}
2\varpi_{\mathrm{sep}}^{+}
&=2^{N_{*}+1}\cdot\mathrm{CWj}\cdot\mathrm{pp}\cdot\exp\left(-\mathrm{vt}\cdot\left(\frac{M\cdot r_{*}}{8L}-4M-7\right)\right)\\
&\le2^{N_{*}+1}\cdot\mathrm{CWj}\cdot\mathrm{pp}^{+}\cdot\exp\left(\mathrm{vt}\cdot(4M+7)\right)\\
&\qquad\cdot\exp\left(-\frac{\mathrm{vt}\cdot M\cdot v_{*}^2}{8L}\right).
\end{align*}
Let $v_{*}\ge\mathrm{vq}_5$.
Lemma~A\,(iii), applied with $\mathrm{cf}_2=\frac{\mathrm{vt}\cdot M}{8L}$, $\mathrm{cf}_1=\frac{N_{*}}{4}$ and $\mathrm{cf}_0=\log^{+}\left(2^{N_{*}+1}\cdot\mathrm{CWj}\cdot\mathrm{pp}^{+}\right)+\mathrm{vt}\cdot(4M+7)$, gives
\begin{align*}
\frac{\mathrm{vt}\cdot M\cdot v_{*}^2}{8L}
&\ge\frac{N_{*}\cdot v_{*}}{4}+\log^{+}\left(2^{N_{*}+1}\cdot\mathrm{CWj}\cdot\mathrm{pp}^{+}\right)\\
&\qquad+\mathrm{vt}\cdot(4M+7).
\end{align*}
Together with (35.10.7) this yields
\[
2\varpi_{\mathrm{sep}}^{+}\le\exp\left(-\frac{N_{*}\cdot v_{*}}{4}\right)\le\bar{\kappa}_1^{N_{*}}.
\]

\emph{(g6).}
Let $\mathrm{vq}_6:=\mathrm{thr}\left(12N_{*},\,0,\,(1-s)\left(\frac{1}{4}-9\varepsilon\right)\right)$.
By (35.10.7) and $0<1-s<1$,
\begin{align*}
\bar{\kappa}_1^{1-s}
&=6^{1-s}\cdot\exp\left(-(1-s)\left(\frac{1}{4}-9\varepsilon\right)v_{*}\right)\\
&\le6\exp\left(-(1-s)\left(\frac{1}{4}-9\varepsilon\right)v_{*}\right).
\end{align*}
If $v_{*}\ge\mathrm{vq}_6$, Lemma~A\,(ii) gives
\[
2N_{*}\cdot\bar{\kappa}_1^{1-s}\le12N_{*}\cdot\exp\left(-(1-s)\left(\frac{1}{4}-9\varepsilon\right)v_{*}\right)\le1.
\]

\emph{(g7), first inequality.}
Let $\mathrm{vq}_{7a}:=4C_{0}^2$ and $\mathrm{vq}_{7b}:=\mathrm{thr}\left(96\,(L\cdot M)^3,\,3r,\,1\right)$, and suppose $v_{*}\ge\mathrm{vq}_{7a}$.
As $v_{*}\ge8$ by (35.10.2) and $p\ge1$, the number $p_{*}=v_{*}^{p}$ is at least $8$; hence $p_{*}-1\ge\frac{p_{*}}{2}$ and $(p_{*}-1)^2\ge\frac{v_{*}^{2p}}{4}$.
The integer $2p-p_0$ is at least $1$, since $p>2p_0$ and $p_0>0$; together with $v_{*}\ge1$ this gives $v_{*}^{2p-p_0}\ge v_{*}\ge4C_{0}^2$.
Therefore
\[
(p_{*}-1)^2\ge C_{0}^2\cdot v_{*}^{p_0}=C_{0}^2\cdot p_{0*},
\]
that is, $\frac{(p_{*}-1)^2}{2C_{0}^2}\ge\frac{p_{0*}}{2}$.
In combination with (35.10.3) this gives
\begin{equation}
\begin{split}
\varpi_0&=18\,\ell_*^3\cdot\exp\left(-\frac{p_{0*}}{2}\right)+6\,\ell_*^3\cdot\exp\left(-\frac{(p_{*}-1)^2}{2C_{0}^2}\right)\\
&\le24\,\ell_*^3\cdot\exp\left(-\frac{p_{0*}}{2}\right)\\
&\le24\,(L\cdot M)^3\cdot v_{*}^{3r}\cdot\exp\left(-\frac{v_{*}^2}{2}\right)\qquad\left(v_{*}\ge\mathrm{vq}_{7a}\right).
\end{split}
\tag{35.10.8}
\end{equation}
Since $v_{*}\ge2$, one has $\frac{v_{*}^2}{2}\ge v_{*}$.
If moreover $v_{*}\ge\mathrm{vq}_{7b}$, Lemma~A\,(ii) gives
\[
4\varpi_0\le96\,(L\cdot M)^3\cdot v_{*}^{3r}\cdot\exp(-v_{*})\le1.
\]

\emph{(g7), second inequality.}
Let $\mathrm{vq}_{7c}:=\mathrm{thr}\left(8N_{*}^{\frac{1}{2}}\cdot L\cdot h^{-\frac{1}{2}}\cdot\mathrm{Cb},\,0,\,1-3\varepsilon\right)$.
By 3.0 and (35.10.4),
\[
t_2=e\cdot\frac{t_2}{e}\le N_{*}^{-\frac{1}{2}}\cdot e\le N_{*}^{-\frac{1}{2}}\cdot L\cdot h^{-\frac{1}{2}}\cdot e_{*},
\]
and $t_x=\frac{e_{*}^{3\varepsilon}}{\mathrm{Cb}}$.
Hence, if $v_{*}\ge\mathrm{vq}_{7c}$, Lemma~A\,(ii) gives
\[
\frac{8N_{*}\cdot t_2}{t_x}\le8N_{*}^{\frac{1}{2}}\cdot L\cdot h^{-\frac{1}{2}}\cdot\mathrm{Cb}\cdot\exp\left(-(1-3\varepsilon)v_{*}\right)\le1.
\]

\emph{(g8).}
Put
\begin{align*}
\mathrm{cf}_8&:=\left(24\,(L\cdot M)^3\right)^{\frac{1}{6N_{*}}}\cdot\exp\left(\frac{\mathbf{D}}{6N_{*}}\right)\cdot\mathrm{mxL}^{\frac{2}{3N_{*}}}\\
&\qquad\cdot5\cdot\left(\frac{N_{*}\cdot h^{\frac{1}{2}}\cdot L^{-\frac{1}{2}}}{2\,\mathrm{Cb}}\cdot\left(\frac{t_2}{e}\right)^{-1}\right)^4,\\
\mathrm{vq}_8&:=\max\left\{1,\ 12N_{*}\cdot\left(\frac{r}{2N_{*}}+\frac{17}{4}+\log^{+}\mathrm{cf}_8\right)\right\};
\end{align*}
here $\mathrm{cf}_8>0$.
Let $v_{*}\ge\max\left\{\mathrm{vq}_{7a},\mathrm{vq}_8\right\}$.
By the definition of $\varpi_L$ in Section~G5.1.8, $\varpi_L^{\frac{1}{3}}=\varpi_0^{\frac{1}{6}}\cdot\exp\left(\frac{\mathbf{D}}{6}\right)$.
From (35.10.8) and from $v_{*}^{\frac{r}{2N_{*}}}\le\exp\left(\frac{r\cdot v_{*}}{2N_{*}}\right)$, which holds because $\log v_{*}\le v_{*}$, one obtains
\begin{align*}
\left(\varpi_L^{\frac{1}{3}}\cdot\mathrm{mxL}^{\frac{2}{3}}\right)^{\frac{1}{N_{*}}}
&=\varpi_0^{\frac{1}{6N_{*}}}\cdot\exp\left(\frac{\mathbf{D}}{6N_{*}}\right)\cdot\mathrm{mxL}^{\frac{2}{3N_{*}}}\\
&\le\left(24\,(L\cdot M)^3\right)^{\frac{1}{6N_{*}}}\cdot\exp\left(\frac{\mathbf{D}}{6N_{*}}\right)\cdot\mathrm{mxL}^{\frac{2}{3N_{*}}}\\
&\qquad\cdot\exp\left(\frac{r\cdot v_{*}}{2N_{*}}-\frac{v_{*}^2}{12N_{*}}\right).
\end{align*}
By Section~G5.1.4, $\mathbf{r}=\frac{e_{*}^{3\varepsilon}}{4\,\mathrm{Cb}}$; by (35.10.4), $\frac{1}{e}\le\left(\frac{h}{L}\right)^{\frac{1}{2}}\cdot\frac{1}{e_{*}}$; and $e_{*}^{3\varepsilon}\le1$.
Consequently
\begin{align*}
\frac{2N_{*}\cdot\mathbf{r}}{t_2}
&=\frac{2N_{*}\cdot e_{*}^{3\varepsilon}}{4\,\mathrm{Cb}\cdot e}\cdot\left(\frac{t_2}{e}\right)^{-1}\\
&\le\frac{N_{*}\cdot h^{\frac{1}{2}}\cdot L^{-\frac{1}{2}}}{2\,\mathrm{Cb}}\cdot\left(\frac{t_2}{e}\right)^{-1}\cdot\exp(v_{*}).
\end{align*}
By the definition of $\lambda_L$ in Section~G5.1.8, the last two estimates give
\[
\lambda_L\le\mathrm{cf}_8\cdot\exp\left(-\frac{v_{*}^2}{12N_{*}}+\left(\frac{r}{2N_{*}}+4\right)v_{*}\right).
\]
Lemma~A\,(iii), applied with $\mathrm{cf}_2=\frac{1}{12N_{*}}$, $\mathrm{cf}_1=\frac{r}{2N_{*}}+\frac{17}{4}$ and $\mathrm{cf}_0=\log^{+}\mathrm{cf}_8$, gives
\[
\frac{v_{*}^2}{12N_{*}}\ge\left(\frac{r}{2N_{*}}+\frac{17}{4}\right)v_{*}+\log^{+}\mathrm{cf}_8.
\]
Hence, by (35.10.7) and $0<s<1$,
\[
\lambda_L\le\exp\left(-\frac{v_{*}}{4}\right)\le\exp\left(-\frac{s\cdot v_{*}}{4}\right)\le\bar{\kappa}_1^{s}.
\]

\emph{(g9).}
Let $\mathrm{vq}_9:=\left(\frac{2L^2}{h\cdot M}\cdot\left(\frac{4}{3}\,D_\beta+3\right)\right)^{\frac{1}{2}}$.
If $v_K\ge\mathrm{vq}_9$, then by (35.10.6)
\[
\ell_c\ge\frac{h\cdot M\cdot v_K^2}{2L^2}\ge\frac{4}{3}\,D_\beta+3;
\]
consequently $\frac{3}{8}\,\ell_c\ge\frac{D_\beta}{2}+\frac{9}{8}\ge\frac{D_\beta}{2}+\frac{1}{2}$, and $\ell_c\ge D_\beta+3$.

\paragraph{3.4. $e_{\sharp}(\beta)>0$.}
Let $\mathrm{vq}$ be the largest of the 21 numbers $\mathrm{vq}_{1a}$, $\mathrm{vq}_{1b}$, $\mathrm{vq}_{1c}$, $\mathrm{vq}_2$, $\mathrm{vq}_{3a}$, $\mathrm{vq}_{3b}$, $\mathrm{vq}_{3c}$, $\mathrm{vq}_{4a},\dots,\mathrm{vq}_{4g}$, $\mathrm{vq}_5$, $\mathrm{vq}_6$, $\mathrm{vq}_{7a}$, $\mathrm{vq}_{7b}$, $\mathrm{vq}_{7c}$, $\mathrm{vq}_8$, $\mathrm{vq}_9$, and put
\[
e_{\mathrm{vq}}:=h^{-\frac{1}{2}}\cdot\min\left\{e_1,\exp(-\mathrm{vq})\right\}.
\]
Since $e_1>0$ (line (C9); Proposition~14.3\,(f)) and $\mathrm{vq}<\infty$, the number $e_{\mathrm{vq}}$ is positive; moreover $e_{\mathrm{vq}}\le e_1\cdot h^{-\frac{1}{2}}$.
Let $0<e\le e_{\mathrm{vq}}$.
Then $e\cdot h^{\frac{1}{2}}\le e_1$ and $v_K=-\log\left(e\cdot h^{\frac{1}{2}}\right)\ge\mathrm{vq}$, so that $v_{*}\ge v_K\ge\mathrm{vq}$ by (35.10.2).
By 3.3 all the inequalities (g1)--(g9) hold at $e$.
Thus $e_{\mathrm{vq}}$ belongs to the set of Section~G5.1.9, whose supremum is therefore at least $e_{\mathrm{vq}}$, and
\begin{equation}
e_{\sharp}(\beta)\ge\frac{e_{\mathrm{vq}}}{2}=\frac{1}{2}\,h^{-\frac{1}{2}}\cdot\min\left\{e_1,\exp(-\mathrm{vq})\right\}>0.
\tag{35.10.9}
\end{equation}

\paragraph{3.5. What $e_{\sharp}(\beta)$ depends on.}
By Section~G5.1.9 the number $e_{\sharp}(\beta)$ is determined by the interval $\left(0,e_1\cdot h^{-\frac{1}{2}}\right]$ together with the set of those $e$ in this interval at which (g1)--(g9) hold.
Besides $e$ and the numbers (35.10.1), the conditions contain the following letters.

\medskip
\noindent
\begin{tabular}{lp{0.3\textwidth}p{0.45\textwidth}}
condition & through & letters\\
\hline
(g1) & $\bar{r}$ & $\mathrm{hL}$, $h_{G}$, $\varepsilon$; $M$; $\mathrm{alG}$, $C_K$, $K_{0}$\\
(g2) & -- & $M$, $L$\\
(g3) & -- & $C_{0}$, $\mathrm{gGa}_{1}$, $\varepsilon$; $C_K$, $K_{0}$\\
(g4) & $\tau_s$, $B_1$, $B_{\mathrm{wr}}$, $B_0$, $\mathrm{pp}$ & $\mu_-$, $\varepsilon$; $\mathrm{gGa}_{1}$, $C_K'$; $c_{\tau}$; $D_\beta$, $M$; $C_{\mathrm{fr}}$; $C_{d}$, $K_{0}$\\
(g5) & $\varpi_{\mathrm{sep}}^{+}$, $\bar{\kappa}_1$, $\mathrm{pp}$ & $N_{*}$, $\mathrm{CWj}$, $\mathrm{vt}$, $M$, $L$, $\varepsilon$; $D_\beta$\\
(g6) & $\bar{\kappa}_1$ & $N_{*}$, $s$, $\varepsilon$\\
(g7) & $\varpi_0$, $t_2$, $t_x$ & $C_{0}$; $N_{*}$, $C_{\beta}$, $\mathbf{D}$; $\mathrm{Cb}$, $\varepsilon$\\
(g8) & $\lambda_L$, $\varpi_L$, $\mathrm{mxL}$, $\varpi_0$, $\mathbf{r}$, $t_2$, $\bar{\kappa}_1$ & $N_{*}$, $\mathbf{D}$, $C_{\beta}$, $\mathrm{Cb}$, $C_{0}$, $s$, $\varepsilon$\\
(g9) & $\ell_c$ & $D_\beta$
\end{tabular}
\medskip

Here $\mathrm{Cb}$ is formed from $(h,C_{\beta})$, and $C_{\mathrm{fr}}$ from $(C_{d},K_{0},M,L,h,C_{\beta},D_\beta)$.
By (35.10.1) the scales are functions of $(e,h;L,M,r,p_0,p)$, and $h$ is a function of $(\bar{m},L)$.
Consequently $e_{\sharp}(\beta)$ is a function of the following quantities, and of no others:
\begin{itemize}
\item $e_1$; $L$, $M$, $r$, $p_0$, $p$, $\varepsilon$, $s$, $N_{*}$, $K_{0}$, $\mu_-$, $C_{0}$, $C_{d}$, $\mathrm{hL}$, $h_{G}$, $\mathrm{alG}$, $\mathrm{gGa}_{1}$, $C_K$, $C_K'$, $c_{\tau}$, $\mathrm{vt}$, $\mathrm{CWj}$, which are constants of the lines (C0)--(C9), (C6$'$);
\item $\bar{m}$, through $h$;
\item $C_{\beta}$, $\mathbf{D}$, $D_\beta$.
\end{itemize}
The parameter $\rho$ stands in the line (C0), and the constants $M$, $C_{0}$, $C_{d}$, $\mathrm{hL}$, $C_K$, $C_K'$, $c_{\tau}$, $\mathrm{vt}$, $\mathrm{CWj}$, $e_1$ may depend on it (Sections~I.8.1 and I.8.2).
The probe enters only through the three numbers $(C_{\beta},\mathbf{D},D_\beta)$: two probes with the same three numbers have, for the same $\bar{m}$, the same value of $e_{\sharp}$.
The list contains no $N$ (by (35.10.1)), no torus (the only torus of the setting is $R^{*}(e)$, which is a function of $e$ by (35.10.1)), no placement and no $\theta$.

\paragraph{3.6. The assertions for $0<e\le e_{\sharp}(\beta)$.}
Let $0<e\le e_{\sharp}(\beta)$ and $N\ge N_0(e)$.
Lemma~0 gives $e\cdot h^{\frac{1}{2}}\le e_1$, so that the setting $(\mathrm{ST}^{*})^{\mathrm{ch}}$ is given, and it gives that (g1)--(g9) hold at $e$.
By 3.1 these conditions do not contain $N$.
The first condition of (g9) and Lemma~1\,(ii) show that every centred placement is admitted.
The second condition of (g9) and Lemma~1\,(iii) give $D_\beta+3\eta\le\ell_c$, which by Section~I.7.6 means that $\beta$ fits $(N,T^{*}(e))$, and they give $R^{*}(e)\ge D_{\mathrm{be}}+1$.
This completes the proof of Proposition~35.10.\qed

\subsection*{4. Conditions on the constants}

Each condition on the constants on which the proofs of Propositions~35.9 and 35.10 rest is one of the conditions of Section~I.0.5, of Section~I.8 or of Proposition~14.2.
Hence both propositions hold for all values of $\rho$, $\alpha$, $\kappa$, $L$, $M$, $\varepsilon$, $s$, $e_0$, $e_1$ that satisfy those conditions.
The conditions, with the steps in which they enter, are the following.

\subsubsection*{4.1. Conditions in the proof of Proposition~35.9}

\begin{description}
\item[(g5): $2\varpi_{\mathrm{sep}}^{+}\le\bar{\kappa}_1^{N_{*}}$.]
This inequality enters 2.2, where together with (35.9.3) it gives $\|\mathrm{kj}_{\mathrm{Scl}}\|_{\mathbf{r}}\le\bar{\kappa}_1^{N_{*}}<1$ for $\#|\mathrm{Scl}|\ge2$, and 2.4; at both places it enters only through the last inequality of Proposition~35.8\,(e).
It is a condition on $e_{\sharp}(\beta)$ in the line (C10) (Section~G5.1.9; Proposition~14.2\,(d)).
\item[(g6): $N_{*}\cdot\bar{\kappa}_1^{1-s}\le\frac{1}{2}$.]
It gives $\bar{\kappa}_1<1$ in (35.9.3) and enters (35.9.6).
Like (g5), it is a condition on $e_{\sharp}(\beta)$ in the line (C10).
\item[(g8): $\lambda_L\le\bar{\kappa}_1^{s}$.]
It yields the second inequality of (a) in 2.6.
Like (g5), it is a condition on $e_{\sharp}(\beta)$ in the line (C10).
\item[(g9), first condition: $\frac{D_\beta}{2}+\frac{1}{2}\le\frac{3}{8}\,\ell_c$.]
Through Lemma~1\,(ii) it enters 2.9, the proof of the last sentence of Proposition~35.9.
Like (g5), it is a condition on $e_{\sharp}(\beta)$ in the line (C10).
\item[$0<s<1$.]
This enters (35.9.3), 2.4 and 2.7.
It belongs to the line (C7); see Proposition~14.2\,(a), row $s$.
\item[$\sigma=\mathfrak{e}-\frac{2a_{L}}{N_{*}}>0$.]
It gives $\mathfrak{e}>0$ in 2.7.
It belongs to the line (C7); see (S2) of Proposition~14.2\,(b).
\item[$\varepsilon>0$.]
It enters (35.9.2), 2.7 (for $\mathfrak{e}<\frac{1}{4}$ and for (35.9.8)) and 2.8.
It belongs to the line (C7); see (S1) of Proposition~14.2\,(b).
\item[$\frac{1}{4}-9\varepsilon>0$.]
It enters 2.10, where $e_{*}$ is written as a function of $\bar{\kappa}_1$.
It is (35.9.8), a consequence of $\sigma>0$, $\varepsilon>0$ and $s>0$; it also follows from $\varepsilon\le\bar{\varepsilon}=\frac{1}{200}$ (lines (C0), (C7); (S1) of Proposition~14.2\,(b)).
\item[$e_1\le\frac{1}{3}$.]
It gives $e_{*}<1$ in (35.9.1).
It belongs to the line (C9); see Proposition~14.2\,(a), row $e_1$, and Proposition~14.2\,(f).
\item[$N_{*}=8L^3\ge2$; $a_{L}>0$.]
These enter (35.9.6) and 2.7.
They follow from the definitions in the line (C4), since $L\ge5$ by (L5) of Proposition~14.2\,(b).
\item[$L\ge1$; $N\ge0$.]
Together they give $\eta\le1$, which enters 2.9 through Lemma~1.
They hold by Section~I.0.5 and by (L5) of Proposition~14.2\,(b).
\item[$h\ge1$.]
It enters 2.0: in the definition of $\mathrm{Ind}_\beta$ (Section~G5.1.3) the condition $N'+\log_L\frac{h}{L}\ge1$ means $N'\ge2-\log_L h$, so that $h\ge1$ gives $\left(N',\frac{h}{L},1\right)\in\mathrm{Ind}_\beta$ for all $N'\ge2$.
In 2.7 it gives $h^{\frac{\mathfrak{e}}{2}}\le h^{\frac{1}{8}}$.
It holds by Section~I.0.5 (the definition of $h$, and $\bar{m}\le1$).
\item[$C_{\beta}<\infty$.]
The number $C_{\beta}$ of the line (C10) is finite by Proposition~35.3\,(c1).
This enters (35.9.2) (for $\mathbf{r}>0$) and 2.10 (finiteness of the right sides).
\end{description}

\subsubsection*{4.2. Conditions in the proof of Proposition~35.10}

The proof of Proposition~35.10 establishes the inequalities (g1)--(g9): for $0<e\le e_{\mathrm{vq}}$ in 3.3 and 3.4, and for $0<e\le e_{\sharp}(\beta)$ through the definition of $e_{\sharp}(\beta)$ (Lemma~0); in this second form the two conditions (g9) enter 3.6.
The conditions on the constants on which this proof rests are the following.

\begin{description}
\item[$r\ge2$.]
It gives $(-\log q)^{r}>1$ in 3.1, whence the finiteness of $N_0(e)$; it also enters (35.10.3) and (35.10.6), and through them (g2), the third to the fifth inequality of (g4), (g5) and (g9).
It belongs to the line (C7); see Proposition~14.2\,(a), row $r$.
\item[$p_0\ge2$.]
It enters (35.10.3), and through it the first inequality of (g7) and (g8).
It follows from $p_0>3r+2$ in the line (C7); see Proposition~14.2\,(f).
\item[$p>2p_0$.]
In the form $2p-p_0\ge1$, $p\ge1$ it enters (35.10.8).
It belongs to the line (C7); see Proposition~14.2\,(a), row $p$.
\item[$0<\varepsilon\le\bar{\varepsilon}=\frac{1}{200}$.]
It gives the inequalities of the paragraph \emph{Exponents} of 3.1.
See the lines (C0), (C7) and (S1) of Proposition~14.2\,(b).
\item[$0<s<1$.]
It enters the proofs of (g6) and (g8) in 3.3.
It belongs to the line (C7).
\item[$M\ge1$; $L\ge1$.]
These enter (g2), Lemma~1 and (35.10.4).
They follow from (M1) and (L5) of Proposition~14.2\,(b).
\item[$L>1$.]
It gives $b_{-}>0$ in 3.0, and $\mathrm{lg}>0$ and $\frac{1}{L}<1$ in 3.1.
By Section~I.0.5, $L$ is an odd integer $\ge5$; see also (L5) of Proposition~14.2\,(b).
\item[$M$ a power of $L$ with $M\ge L^2$.]
It gives in 3.1 that $R^{*}(e)$ lies in $\mathrm{Tor}(e)$.
See Section~I.0.5, the line (C5) and (M1) of Proposition~14.2\,(b).
\item[$e_1\le\frac{1}{3}$.]
It gives $v_K>1$ in 3.1, whence the finiteness of $N_0(e)$.
It belongs to the line (C9); see Proposition~14.2\,(a), row $e_1$, and Proposition~14.2\,(f).
\item[$e_1>0$.]
It enters 3.4.
It belongs to the line (C9) and holds by Proposition~14.3\,(f), in the form (G5.T-1) of Section~G5.7.
\item[Finiteness and positivity of constants.]
The numbers $\mathrm{hL}$, $h_{G}$, $\mathrm{alG}$, $C_K$, $C_K'$, $K_{0}$, $C_{0}$, $\mathrm{gGa}_{1}$, $C_{d}$, $\mathrm{CWj}$, $N_{*}$, $\mu_-$, $c_{\tau}$, $\mathrm{vt}$ are finite and $>0$; this enters the first paragraph of 3.0 and, in 3.3, the finiteness of the numbers $\mathrm{vq}_{1a},\dots,\mathrm{vq}_9$.
It follows from the formulas of the lines (C0), (C4), (C6$'$) (Section~I.8.2; Section~G5.1.8; Proposition~14.2\,(a)) together with (G5.T-1) of Section~G5.7, which follows from Propositions~14.1 and 14.2\,(a), (f): $C^{F}\ge1$, $0<\gamma^{F}\le1$, $C_{0}\ge\max\left(1,C^{F}\right)$, $M\ge L^2$, $K_{0}\ge1$, $C_{d}<\infty$; further $C_{w}\ge C^{F}$ (Proposition~14.2\,(f)), and $C^{F},C_{0},C_{w}<\infty$ (the paragraph ``Ranges'' of Proposition~14.2\,(f)).
\item[$h\ge1$; $\log_L h$ is an integer.]
These enter 3.0 and 3.1.
They hold by Section~I.0.5 (the definition of $h$, and $\bar{m}\le1$).
\item[$C_{\beta}<\infty$, $\mathbf{D}<\infty$.]
The numbers $C_{\beta}$ and $\mathbf{D}$ of the line (C10) are finite by Proposition~35.3\,(c1); this enters 3.0.
\item[The conditions (e1) of Proposition~14.2\,(b) on $e_1$.]
They enter the proof only through Proposition~35.1, on which (35.10.2)--(35.10.6) rest.
\end{description}

\section*{Chapter 37. The Bloch symbol and its inverse, the four-photon coefficient, the free coefficients on the reference torus: proofs of Propositions 37.1--37.3}

\subsection*{0. Conventions of this chapter}

\subsubsection*{0.1. Setting}

Notation and conventions are those of Part~I and of Section~G5.0.3.
The setting (FF), together with its paragraphs on the fibre and on the continuum, is the one of Section~G5.6.0;
Propositions~37.1, 37.2 and~37.3 are formulated in it.

\subsubsection*{0.2. Lattice units, the fibre, uniformity}

Proposition~37.1 concerns the $192\times192$ matrices $D^{(p)}$ of (I.2.6).
Propositions~37.2 and~37.3 are formulated in lattice units (Section~I.2.5, $\eta_0=1$).
In these units the index set is $\mathbb{Z}^3\times O_{\mathrm{cub}}$:
over every photon site $c$ stand the $48$ fermion sites $(c,g)$, $g\in O_{\mathrm{cub}}$,
and every pairing carries the weight $\frac{1}{48}$ of (I.1.3);
these are the conventions of level~$0$.
The uniformity asserted in the three statements is uniformity in $(p,\mathbf{m})$ for Proposition~37.1;
in $(\eta,R)$, hence in $n=R/\eta$, and in the position of $\mathrm{supp}\,\beta$ for Proposition~37.2;
and in $(e,N,v,\mathfrak{c})$ for Proposition~37.3, the only one of the three statements that involves the charge $e$.

\emph{The fibre.}
The map $J^{*}$ is the mean over the $48$ labels $g\in O_{\mathrm{cub}}$.
The space ${V_{\mathrm{v}}}$ has dimension $192=48\cdot4$;
consequently $|\mathrm{tr}\,O|\le192\,|O|_2$ and $\|\mathrm{Mx}\|_4^4\le192\,|\mathrm{Mx}|_2^4$ for all matrices $O$ and $\mathrm{Mx}$ on ${V_{\mathrm{v}}}$.
The space $\mathfrak{S}^{\mathrm{ch}}(1;n)$ has dimension $192\,n^3$.

\emph{Norms.}
On $\mathbb{R}^3$ and on $\mathbb{C}^3$ the norms $|y|_2$, $|y|_\infty$, $|y|_1$ are those of Section~I.1.1.
On $\mathbb{C}^4$ and on ${V_{\mathrm{v}}}$ all norms are Euclidean:
on $\mathbb{C}^4$ the norm is $|\xi|_2$, and on ${V_{\mathrm{v}}}$
\[
(u',u):=\frac{1}{48}\sum_{g\in O_{\mathrm{cub}}}\overline{u'(g)}\cdot u(g),\qquad |u|_2:=(u,u)^{\frac{1}{2}}.
\]
For a matrix acting on $\mathbb{C}^4$ or on ${V_{\mathrm{v}}}$, and for a bounded operator on a Hilbert space, $|O|_2$ is the operator norm;
adjoints $O^{*}$ are taken with respect to these inner products.
If the inner product of a space is multiplied by a positive number, then the adjoint, the operator norm, the singular values, the trace and the determinant of a linear map of this space into itself remain unchanged.

\emph{Integration theory.}
For the theorems of Fubini and Tonelli, the monotone and the dominated convergence theorems, H\"older's inequality (whose case of two exponents equal to $2$ is the Cauchy--Schwarz inequality) and Plancherel's theorem, in the form that the Fourier transform is injective on $L^2(\mathbb{R}^3)$,
see \cite[Theorems~1.12, 1.6, 1.8, 2.3, 5.3]{LiebLoss2001}.

\subsubsection*{0.3. Notation}

In this chapter $\mathrm{Plt}$ and $\mathrm{Phv}$ are the projections of (I.J4), formed in the fibre,
and $\mathrm{Mx}$ denotes a matrix, or a matrix-valued function, as in (I.ke1).
The number $\delta(\rho)$ that is named in Proposition~37.1 after clause~(h) is fixed in [1.10].
In Section~2, $\mathrm{fs}$ denotes a scalar function, as in Section~I.7.6, $\psi$ a spinor function, and $\mathrm{cf}$ a number that is fixed where it appears;
$I_1$ and $I_2$ are the Schatten classes of order one and two (the trace class and the Hilbert--Schmidt class), in the same way as $I_4$ is the Schatten class of order four in the definitions that precede Proposition~37.2;
$P$ denotes an element of $\mathrm{Kf}_R$, a momentum;
and $N_1$ denotes a value of the index $N$ of the spacing (in 2.12\,(f): $RL^N$ is an integer for all $N\ge N_1$).
Further, $\mathrm{Ran}\,O$ is the range of $O$;
$\mathrm{End}(\mathrm{Fb})$ is the space of the linear maps of $\mathrm{Fb}$ into itself;
for multi-indices, $\binom{\mathrm{ak}}{\mathrm{bi}}:=\prod_\nu\binom{\mathrm{ak}_\nu}{\mathrm{bi}_\nu}$;
$\mathrm{d}S$ is the surface measure of a sphere;
$\ell^1$, $\ell^2$, $\ell^\infty$ are the sequence spaces, and $\|\cdot\|_{\ell^4}$ is the norm defined in 2.1.
An index or a prime attached to a letter ($u_i$, $\mathrm{hm}_i$, $\mathrm{fn}_i$, $O_1$, $t_r'$, $p'$, $y''$) denotes a further object of the same kind.

\subsection*{1. Proposition 37.1}

\subsubsection*{Proof of Proposition 37.1}

\paragraph{1.0. Notation and eight facts.}
Throughout, $\rho\in(0,1)$, $\mathbf{m}\ge0$, and $p\in\mathbb{C}^3$ unless $p$ is said to be real.
We write $R_3$ for $R_3(p)$.
For $i=0,\dots,3$ put
\[
(\mathrm{Gd}_i u)(g):=\gamma\cdot\hat{d}_i(g)\,u(g),\qquad \hat{d}_i(g)=-g\hat n_i
\]
(in particular $\hat{d}_3(g)=ge_1$).
Further $\mathrm{Plt}:=JJ^{*}$ and $\mathrm{Phv}:=1-\mathrm{Plt}$.
For real $p$ and $\mathbf{m}>0$ let $\mathrm{Sdr}(p,\mathbf{m}):=(\mathbf{m}+i\gamma\cdot p)^{-1}$, a $4\times4$ matrix.
By (I.2.1), $M_i(g)=\frac{w_i}{2}\left(\rho-\gamma\cdot\hat{d}_i(g)\right)$.
\begin{itemize}
\item[(w1)] By Section~I.1.2, $s_i^2=1$.
Hence $g\mapsto gs_i$ is a bijection of $O_{\mathrm{cub}}$ onto itself, and $(gs_i)s_i=g$;
this is (I.s1), read on the labels.
\item[(w2)] Since $s_i\hat n_i=-\hat n_i$, one has $\hat{d}_i(gs_i)=-gs_i\hat n_i=g\hat n_i=-\hat{d}_i(g)$;
this is (I.s2).
As $g$ is orthogonal, $|\hat{d}_i(g)|_2=|\hat n_i|_2=1$.
\item[(w3)] By (I.1.1), $\sum_{i=0}^{3}w_i\hat{d}_i(g)=-g\sum_iw_i\hat n_i=0$;
this is (I.s3).
\item[(w4)] Each of the six vectors $\pm e_\mu$ is the first column $ge_1$ of exactly $8$ of the $48$ labels $g$;
compare (I.s4).
\item[(w5)] One has $J^{*}J\xi=\frac{1}{48}\sum_g\xi=\xi$ and $(u,J\xi)=\frac{1}{48}\sum_g\overline{u(g)}\cdot\xi=\overline{J^{*}u}\cdot\xi$.
Hence $J^{*}$ is the adjoint of $J$, and $J$ is an isometric embedding of $\mathbb{C}^4$ in ${V_{\mathrm{v}}}$;
the operator $\mathrm{Plt}$ is the orthogonal projection of ${V_{\mathrm{v}}}$ onto $J\mathbb{C}^4$, which consists of the functions that are constant on $O_{\mathrm{cub}}$;
one has $J^{*}\cdot\mathrm{Phv}=0$ and $\mathrm{Phv}\cdot J=0$;
and $|JOJ^{*}|_2=|O|_2$ holds for all $4\times4$ matrices $O$.
\item[(w6)] By (I.c1), for every $y\in\mathbb{R}^3$ the matrix $\gamma\cdot y$ is Hermitian and $(\gamma\cdot y)^2=|y|_2^2$.
For a unit vector $y$ the eigenvalues of $\gamma\cdot y$ are therefore in $\{1,-1\}$, and $|\rho+\gamma\cdot y|_2\le1+\rho$.
For real $p$ and $\mathbf{m}>0$,
\[
(\mathbf{m}+i\gamma\cdot p)^{*}(\mathbf{m}+i\gamma\cdot p)=(\mathbf{m}-i\gamma\cdot p)(\mathbf{m}+i\gamma\cdot p)=\mathbf{m}^2+|p|_2^2 .
\]
Therefore $\left(\mathbf{m}^2+|p|_2^2\right)^{\frac{1}{2}}\mathrm{Sdr}(p,\mathbf{m})$ is unitary,
$|\mathrm{Sdr}(p,\mathbf{m})|_2=\left(\mathbf{m}^2+|p|_2^2\right)^{-\frac{1}{2}}$,
and $|p|_2\,|\mathrm{Sdr}(p,\mathbf{m})|_2\le1$.
\item[(w7)] Given $4\times4$ matrices $\mathrm{Mx}(g)$, $g\in O_{\mathrm{cub}}$, and an index $i\in\{0,\dots,3\}$,
the operator $O$ on ${V_{\mathrm{v}}}$ defined by $(Ou)(g):=\mathrm{Mx}(g)\,u(gs_i)$ has the norm $|O|_2\le\max_g|\mathrm{Mx}(g)|_2$;
indeed, by (w1),
\begin{align*}
|Ou|_2^2=\frac{1}{48}\sum_g|\mathrm{Mx}(g)\,u(gs_i)|_2^2
&\le\max_g|\mathrm{Mx}(g)|_2^2\cdot\frac{1}{48}\sum_g|u(gs_i)|_2^2\\
&=\max_g|\mathrm{Mx}(g)|_2^2\,|u|_2^2 .
\end{align*}
\item[(w8)] Suppose that $O$ is a linear map of a finite-dimensional inner product space into itself and that, for some $\mathrm{cw}>0$, $\mathrm{Re}\,(u,Ou)\ge\mathrm{cw}\,(u,u)$ for every $u$.
Then $O$ is invertible with $|O^{-1}|_2\le\frac{1}{\mathrm{cw}}$;
indeed $\mathrm{cw}\,|u|_2^2\le|(u,Ou)|\le|u|_2\,|Ou|_2$.
\end{itemize}

\paragraph{1.1. Clause (a).}
\begin{proof}[Proof of Proposition~37.1\,(a)]
\emph{The operators $R_i$.}
For $i\le2$ one has $(R_i^2u)(g)=u(gs_is_i)=u(g)$.
For $i=3$ and every $p\in\mathbb{C}^3$, since $s_3e_1=-e_1$,
\[
(R_3(p)^2u)(g)=\exp\left(ip\cdot(ge_1)\right)\exp\left(ip\cdot(gs_3e_1)\right)u(gs_3s_3)=u(g).
\]
If $p$ is real, then (w1) and $\left|\exp\left(ip\cdot(ge_1)\right)\right|=1$ show that every $R_i$ preserves the norm.
Thus for real $p$ the four operators $R_0$, $R_1$, $R_2$, $R_3(p)$ are unitary involutions, hence Hermitian;
in particular $(u,R_iu)$ is real, and
\begin{equation*}
\left(u,(1-R_i)u\right)=\frac{1}{2}\left(|u|_2^2-2(u,R_iu)+|R_iu|_2^2\right)=\frac{1}{2}|u-R_iu|_2^2\ge0. \tag*{[1.1]}
\end{equation*}

\emph{The operators $\mathrm{Gd}_i$.}
By (w6) and (w2), $\mathrm{Gd}_i^{*}=\mathrm{Gd}_i$.
For every $p$, (w2) yields
\begin{align*}
(R_i\cdot\mathrm{Gd}_iu)(g)&=\gamma\cdot\hat{d}_i(gs_i)\,u(gs_i)\\
&=-\gamma\cdot\hat{d}_i(g)\,u(gs_i)=-(\mathrm{Gd}_i\cdot R_iu)(g)\qquad(i\le2),\\
(R_3(p)\cdot\mathrm{Gd}_3u)(g)&=\exp\left(ip\cdot(ge_1)\right)\gamma\cdot\hat{d}_3(gs_3)\,u(gs_3)\\
&=-\gamma\cdot\hat{d}_3(g)\exp\left(ip\cdot(ge_1)\right)u(gs_3)=-\left(\mathrm{Gd}_3\cdot R_3(p)u\right)(g).
\end{align*}
Hence
\begin{equation*}
R_i\cdot\mathrm{Gd}_i=-\mathrm{Gd}_i\cdot R_i\qquad(i=0,\dots,3). \tag*{[1.2]}
\end{equation*}

\emph{The decomposition, for every $p\in\mathbb{C}^3$.}
We use (I.2.6) at $\boldsymbol{\mu}=0$, (I.m1) in the form $\sum_iM_i(g)=\frac{\rho}{2}\,\mathrm{sw}$, and $\mathrm{sw}=\sum_iw_i$:
\begin{align*}
(D_0^{(p)}u)(g)&=\sum_{i\le3}M_i(g)\,u(g)-\sum_{i\le3}M_i(g)\,(R_iu)(g)\\
&=\frac{\rho}{2}\sum_iw_i\,u(g)-\sum_i\frac{w_i}{2}\left(\rho-\gamma\cdot\hat{d}_i(g)\right)(R_iu)(g)\\
&=\frac{\rho}{2}\sum_iw_i\,(u-R_iu)(g)+\frac{1}{2}\sum_iw_i\,\gamma\cdot\hat{d}_i(g)\,(R_iu)(g).
\end{align*}
In the last line the first sum equals $\frac{\rho}{2}\left((\nabla^{*}\nabla)(p)u\right)(g)$,
while the second equals $\frac{1}{2}(\mathrm{Gh}(p)u)(g)$, because $\mathrm{Gh}(p)=\sum_{i\le3}w_i\,\mathrm{Gd}_i\cdot R_i$.
We have shown $D_0^{(p)}=\frac{\rho}{2}(\nabla^{*}\nabla)(p)+\frac{1}{2}\mathrm{Gh}(p)$.

\emph{Real $p$.}
Being a sum of Hermitian operators, $(\nabla^{*}\nabla)(p)$ is Hermitian, and by [1.1]
\begin{equation*}
\left(u,(\nabla^{*}\nabla)(p)u\right)=\sum_{i=0}^{3}\frac{w_i}{2}|u-R_iu|_2^2\ge0. \tag*{[1.3]}
\end{equation*}
Moreover $(\mathrm{Gd}_i\cdot R_i)^{*}=R_i\cdot\mathrm{Gd}_i$, which equals $-\mathrm{Gd}_i\cdot R_i$ by [1.2].
Therefore $\mathrm{Gh}(p)^{*}=-\mathrm{Gh}(p)$, so that $(u,\mathrm{Gh}(p)u)$ is purely imaginary.
Since $D^{(p)}=\mathbf{m}+D_0^{(p)}$, it follows that
$\mathrm{Re}\,(u,D^{(p)}u)=\mathbf{m}\,(u,u)+\frac{\rho}{2}\left(u,(\nabla^{*}\nabla)(p)u\right)\ge\mathbf{m}\,(u,u)$.
\end{proof}

\paragraph{1.2. Clause (b).}
\begin{proof}[Proof of Proposition~37.1\,(b)]
\emph{Generators.}
The transpositions $(23)$ and $(12)$, whose permutation matrices are $s_1$ and $s_2$, generate the symmetric group of $\{1,2,3\}$.
Conjugating $s_0=\mathrm{diag}(1,1,-1)$ by permutation matrices one obtains the three diagonal matrices with one entry $-1$;
they generate $\{\pm1\}^3$.
Each element of $O_{\mathrm{cub}}$ is a product of an element of $\{\pm1\}^3$ and a permutation matrix.
Therefore $s_0$, $s_1$, $s_2$ generate $O_{\mathrm{cub}}$.

\emph{Kernel.}
Let $p$ be real and $(\nabla^{*}\nabla)(p)u=0$.
Then [1.3] forces $u=R_iu$ for every $i\in\{0,\dots,3\}$.
For $i\le2$ this states $u(g)=u(gs_i)$ for all $g$;
as $s_0$, $s_1$, $s_2$ generate $O_{\mathrm{cub}}$, one gets $u(g)=u(1)=:\xi$ for all $g$, that is, $u=J\xi$.
For $i=3$ it then states $\left(1-\exp\left(ip\cdot(ge_1)\right)\right)\xi=0$ for all $g$, that is, by (w4), $\left(1-\exp(\pm ip_\mu)\right)\xi=0$ for $\mu=1,2,3$.
When $p\notin2\pi\mathbb{Z}^3$, at least one of these scalar factors does not vanish, whence $\xi=0$:
$\ker(\nabla^{*}\nabla)(p)=\{0\}$.
For $p=0$, on the other hand, $R_iJ\xi=J\xi$ for all $i$ and all $\xi$, so that $(\nabla^{*}\nabla)(0)J\xi=0$.
Hence $\ker(\nabla^{*}\nabla)(0)=J\mathbb{C}^4$.

\emph{Reduction of $(\nabla^{*}\nabla)(0)$ to eight pieces.}
For $\mathrm{dn}\in\{\pm1\}^3$ let $(\mathrm{Lt}_{\mathrm{dn}}u)(g):=u(\mathrm{dn}\cdot g)$.
For a subset $\mathrm{Ia}$ of $\{1,2,3\}$ put
\begin{gather*}
\mathrm{sna}_{\mathrm{Ia}}(\mathrm{dn}):=\prod_{\mu\in\mathrm{Ia}}\mathrm{dn}_{\mu\mu},\\
{V_{\mathrm{v}}}_{\mathrm{Ia}}:=\left\{u\in{V_{\mathrm{v}}}:\ \mathrm{Lt}_{\mathrm{dn}}u=\mathrm{sna}_{\mathrm{Ia}}(\mathrm{dn})\,u\ \text{for all }\mathrm{dn}\in\{\pm1\}^3\right\}.
\end{gather*}
Since $g\mapsto\mathrm{dn}\cdot g$ is a bijection of $O_{\mathrm{cub}}$, each $\mathrm{Lt}_{\mathrm{dn}}$ is unitary;
as moreover $\mathrm{Lt}_{\mathrm{dn}}^2=1$, it is Hermitian.
Further, $\mathrm{Lt}_{\mathrm{dn}}\cdot\mathrm{Lt}_{\mathrm{dn}'}=\mathrm{Lt}_{\mathrm{dn}\cdot\mathrm{dn}'}$, as $\{\pm1\}^3$ is abelian;
$\mathrm{sna}_{\mathrm{Ia}}(\mathrm{dn})\cdot\mathrm{sna}_{\mathrm{Ia}}(\mathrm{dn}')=\mathrm{sna}_{\mathrm{Ia}}(\mathrm{dn}\cdot\mathrm{dn}')$;
$\mathrm{sna}_{\mathrm{Ia}}(\mathrm{dn})^2=1$;
and
\[
\sum_{\mathrm{Ia}\subset\{1,2,3\}}\mathrm{sna}_{\mathrm{Ia}}(\mathrm{dn})=\prod_\mu(1+\mathrm{dn}_{\mu\mu})=8\cdot\mathbf{1}[\mathrm{dn}=1].
\]
It follows that the eight operators $\frac{1}{8}\sum_{\mathrm{dn}\in\{\pm1\}^3}\mathrm{sna}_{\mathrm{Ia}}(\mathrm{dn})\cdot\mathrm{Lt}_{\mathrm{dn}}$ are Hermitian idempotents whose ranges are the spaces ${V_{\mathrm{v}}}_{\mathrm{Ia}}$ and whose sum is $1$;
so ${V_{\mathrm{v}}}$ is the orthogonal direct sum of the eight spaces ${V_{\mathrm{v}}}_{\mathrm{Ia}}$.
Since left and right multiplications in $O_{\mathrm{cub}}$ commute, $\mathrm{Lt}_{\mathrm{dn}}$ commutes with $R_0$, $R_1$, $R_2$ and $R_3(0)$;
consequently $(\nabla^{*}\nabla)(0)$ maps every ${V_{\mathrm{v}}}_{\mathrm{Ia}}$ into itself.
Every eigenvalue of $(\nabla^{*}\nabla)(0)$ is therefore an eigenvalue of its restriction to one of the spaces ${V_{\mathrm{v}}}_{\mathrm{Ia}}$, and conversely.

\emph{Coordinates on ${V_{\mathrm{v}}}_{\mathrm{Ia}}$.}
To a permutation $\mathrm{pe}$ of $\{1,2,3\}$ associate the matrix $\mathrm{Pe}(\mathrm{pe})$ defined by $\mathrm{Pe}(\mathrm{pe})e_\mu=e_{\mathrm{pe}(\mu)}$;
then $\mathrm{Pe}(\mathrm{pe})\cdot\mathrm{Pe}(\mathrm{pe}')=\mathrm{Pe}(\mathrm{pe}\cdot\mathrm{pe}')$.
Let $\mathrm{dg}_\mu:=1-2e_\mu e_\mu^{T}$, $\mathrm{tp}_1:=(23)$, $\mathrm{tp}_2:=(12)$.
Then $s_1=\mathrm{Pe}(\mathrm{tp}_1)$, $s_2=\mathrm{Pe}(\mathrm{tp}_2)$, $s_0=\mathrm{dg}_3$, $s_3=\mathrm{dg}_1$, and $\mathrm{Pe}(\mathrm{pe})\cdot\mathrm{dg}_\mu=\mathrm{dg}_{\mathrm{pe}(\mu)}\cdot\mathrm{Pe}(\mathrm{pe})$.
Every $g\in O_{\mathrm{cub}}$ can be written uniquely as $g=\mathrm{dn}\cdot\mathrm{Pe}(\mathrm{pe})$ with $\mathrm{dn}\in\{\pm1\}^3$ and a permutation $\mathrm{pe}$.
Consequently an element $u$ of ${V_{\mathrm{v}}}_{\mathrm{Ia}}$ is determined by its six values $u(\mathrm{Pe}(\mathrm{pe}))$, which may be prescribed arbitrarily in $\mathbb{C}^4$;
one has $u(\mathrm{dn}\cdot\mathrm{Pe}(\mathrm{pe}))=\mathrm{sna}_{\mathrm{Ia}}(\mathrm{dn})\,u(\mathrm{Pe}(\mathrm{pe}))$ and $(u,u)=\frac{8}{48}\sum_{\mathrm{pe}}|u(\mathrm{Pe}(\mathrm{pe}))|_2^2$.
For $u\in{V_{\mathrm{v}}}_{\mathrm{Ia}}$,
\begin{align*}
(R_1u)(\mathrm{Pe}(\mathrm{pe}))&=u(\mathrm{Pe}(\mathrm{pe}\cdot\mathrm{tp}_1)),\\
(R_2u)(\mathrm{Pe}(\mathrm{pe}))&=u(\mathrm{Pe}(\mathrm{pe}\cdot\mathrm{tp}_2)),\\
(R_0u)(\mathrm{Pe}(\mathrm{pe}))&=u\left(\mathrm{dg}_{\mathrm{pe}(3)}\cdot\mathrm{Pe}(\mathrm{pe})\right)=(-1)^{\mathbf{1}[\mathrm{pe}(3)\in\mathrm{Ia}]}\,u(\mathrm{Pe}(\mathrm{pe})),\\
(R_3(0)u)(\mathrm{Pe}(\mathrm{pe}))&=u\left(\mathrm{dg}_{\mathrm{pe}(1)}\cdot\mathrm{Pe}(\mathrm{pe})\right)=(-1)^{\mathbf{1}[\mathrm{pe}(1)\in\mathrm{Ia}]}\,u(\mathrm{Pe}(\mathrm{pe})).
\end{align*}
With $w_0=w_3=6$, $w_1=w_2=6\sqrt{2}$ and $\mathrm{pot}_{\mathrm{Ia}}(\mathrm{pe}):=12\left(\mathbf{1}[\mathrm{pe}(3)\in\mathrm{Ia}]+\mathbf{1}[\mathrm{pe}(1)\in\mathrm{Ia}]\right)$ one obtains
\begin{equation*}
\begin{aligned}
\left((\nabla^{*}\nabla)(0)u\right)(\mathrm{Pe}(\mathrm{pe}))
&=6\sqrt{2}\,\left(2u(\mathrm{Pe}(\mathrm{pe}))-u(\mathrm{Pe}(\mathrm{pe}\cdot\mathrm{tp}_1))-u(\mathrm{Pe}(\mathrm{pe}\cdot\mathrm{tp}_2))\right)\\
&\quad+\mathrm{pot}_{\mathrm{Ia}}(\mathrm{pe})\cdot u(\mathrm{Pe}(\mathrm{pe})).
\end{aligned}
\tag*{[1.4]}
\end{equation*}

\emph{The six permutations form a cycle.}
Starting from a permutation $\mathrm{pe}_0$, define recursively $\mathrm{pe}_{j+1}:=\mathrm{pe}_j\cdot\mathrm{tp}_1$ for even $j$ and $\mathrm{pe}_{j+1}:=\mathrm{pe}_j\cdot\mathrm{tp}_2$ for odd $j$.
The product $\mathrm{tp}_1\cdot\mathrm{tp}_2$ has order three.
Hence $\mathrm{pe}_6=\mathrm{pe}_0$;
$\mathrm{pe}_0$, $\mathrm{pe}_2$, $\mathrm{pe}_4$ are the three permutations with the sign of $\mathrm{pe}_0$, and $\mathrm{pe}_1$, $\mathrm{pe}_3$, $\mathrm{pe}_5$ are the three others.
All indices $j$ are taken modulo $6$.
If $j$ is even, then $\mathrm{pe}_j\cdot\mathrm{tp}_1=\mathrm{pe}_{j+1}$, and from $\mathrm{pe}_j=\mathrm{pe}_{j-1}\cdot\mathrm{tp}_2$ one gets $\mathrm{pe}_j\cdot\mathrm{tp}_2=\mathrm{pe}_{j-1}$.
If $j$ is odd, then $\mathrm{pe}_j\cdot\mathrm{tp}_2=\mathrm{pe}_{j+1}$, and from $\mathrm{pe}_j=\mathrm{pe}_{j-1}\cdot\mathrm{tp}_1$ one gets $\mathrm{pe}_j\cdot\mathrm{tp}_1=\mathrm{pe}_{j-1}$.
In terms of $\mathrm{yc}_j:=u(\mathrm{Pe}(\mathrm{pe}_j))\in\mathbb{C}^4$ the identity [1.4] reads
\begin{equation*}
\begin{aligned}
\left((\nabla^{*}\nabla)(0)u\right)(\mathrm{Pe}(\mathrm{pe}_j))
&=6\sqrt{2}\,(2\mathrm{yc}_j-\mathrm{yc}_{j+1}-\mathrm{yc}_{j-1})+\mathrm{pot}_{\mathrm{Ia}}(\mathrm{pe}_j)\cdot\mathrm{yc}_j\\
&\qquad(j\ \text{modulo}\ 6).
\end{aligned}
\tag*{[1.5]}
\end{equation*}
The operator defined by the first term is non-negative:
\[
\sum_j\overline{\mathrm{yc}_j}\cdot(2\mathrm{yc}_j-\mathrm{yc}_{j+1}-\mathrm{yc}_{j-1})=\sum_j|\mathrm{yc}_j-\mathrm{yc}_{j+1}|_2^2 .
\]

\emph{The case $\mathrm{Ia}=\emptyset$.}
Here $\mathrm{pot}_{\mathrm{Ia}}=0$.
For $\mathrm{jj}\in\{0,\dots,5\}$ and $\xi\in\mathbb{C}^4$ consider $\mathrm{yc}_j=\exp\left(\frac{i\pi\,\mathrm{jj}\,j}{3}\right)\xi$.
These sequences span the space of all sequences $(\mathrm{yc}_j)_j$, and each of them is an eigenvector of [1.5] with the eigenvalue $6\sqrt{2}\left(2-2\cos\frac{\pi\,\mathrm{jj}}{3}\right)$;
these eigenvalues are $0$ for $\mathrm{jj}=0$, $6\sqrt{2}$ for $\mathrm{jj}=1,5$, $18\sqrt{2}$ for $\mathrm{jj}=2,4$, and $24\sqrt{2}$ for $\mathrm{jj}=3$.

\emph{The case $\mathrm{Ia}=\{1,2,3\}$.}
Here $\mathrm{pot}_{\mathrm{Ia}}=24$, and every eigenvalue is at least $24$.

\emph{The case where $\mathrm{Ia}$ is $\{1,2,3\}$ with one element $\mathfrak{l}$ removed.}
Here $\mathrm{pot}_{\mathrm{Ia}}(\mathrm{pe})=12\left(2-\mathbf{1}[\mathrm{pe}(3)=\mathfrak{l}]-\mathbf{1}[\mathrm{pe}(1)=\mathfrak{l}]\right)\ge12$, because $\mathrm{pe}(3)\ne\mathrm{pe}(1)$;
every eigenvalue is at least $12$.

\emph{The case $\mathrm{Ia}=\{\mathfrak{l}\}$.}
Let $\mathrm{pe}_0$ be chosen with $\mathrm{pe}_0(2)=\mathfrak{l}$.
Then
\begin{align*}
\mathrm{pe}_1(3)&=\mathrm{pe}_0(\mathrm{tp}_1(3))=\mathrm{pe}_0(2)=\mathfrak{l}, &
\mathrm{pe}_2(3)&=\mathrm{pe}_1(\mathrm{tp}_2(3))=\mathrm{pe}_1(3)=\mathfrak{l},\\
\mathrm{pe}_3(2)&=\mathrm{pe}_2(\mathrm{tp}_1(2))=\mathrm{pe}_2(3)=\mathfrak{l}, &
\mathrm{pe}_4(1)&=\mathrm{pe}_3(\mathrm{tp}_2(1))=\mathrm{pe}_3(2)=\mathfrak{l},\\
\mathrm{pe}_5(1)&=\mathrm{pe}_4(\mathrm{tp}_1(1))=\mathrm{pe}_4(1)=\mathfrak{l}.
\end{align*}
Hence $\left(\mathrm{pot}_{\mathrm{Ia}}(\mathrm{pe}_j)\right)_{j=0,\dots,5}=(0,12,12,0,12,12)$, and [1.5] is $\mathrm{Hx}\otimes1_4$ with the real symmetric $6\times6$ matrix $\mathrm{Hx}$ given by
\begin{gather*}
(\mathrm{Hx}\cdot\mathrm{yc})_j:=6\sqrt{2}\,(2\mathrm{yc}_j-\mathrm{yc}_{j+1}-\mathrm{yc}_{j-1})+\mathrm{pot}_j\cdot\mathrm{yc}_j\qquad(\mathrm{yc}\in\mathbb{C}^6),\\
(\mathrm{pot}_j)_{j=0,\dots,5}=(0,12,12,0,12,12).
\end{gather*}
Since $\mathrm{pot}_{j+3}=\mathrm{pot}_j=\mathrm{pot}_{-j}$, the matrix $\mathrm{Hx}$ commutes with the two commuting involutions $(\mathrm{ro}\cdot\mathrm{yc})_j:=\mathrm{yc}_{j+3}$ and $(\mathrm{rf}\cdot\mathrm{yc})_j:=\mathrm{yc}_{-j}$.
Their four joint eigenspaces, listed for the pairs of eigenvalues $(\mathrm{ro},\mathrm{rf})=(+,+),(+,-),(-,+),(-,-)$ in this order, are spanned by
\begin{gather*}
\mathrm{hb}_1=(1,0,0,1,0,0),\quad\mathrm{hb}_2=(0,1,1,0,1,1);\qquad\mathrm{hb}_5=(0,1,-1,0,1,-1);\\
\mathrm{hb}_3=(1,0,0,-1,0,0),\quad\mathrm{hb}_4=(0,1,-1,0,-1,1);\qquad\mathrm{hb}_6=(0,1,1,0,-1,-1)
\end{gather*}
(the components are indexed by $j=0,\dots,5$; the dimensions add up to $6$).
$\mathrm{Hx}$ maps each of the four spaces into itself, and the components $j=0$ and $j=1$ determine the images:
\begin{align*}
\mathrm{Hx}\cdot\mathrm{hb}_1&=12\sqrt{2}\,\mathrm{hb}_1-6\sqrt{2}\,\mathrm{hb}_2, &
\mathrm{Hx}\cdot\mathrm{hb}_2&=-12\sqrt{2}\,\mathrm{hb}_1+(12+6\sqrt{2})\,\mathrm{hb}_2,\\
\mathrm{Hx}\cdot\mathrm{hb}_5&=(12+18\sqrt{2})\,\mathrm{hb}_5;\\
\mathrm{Hx}\cdot\mathrm{hb}_3&=12\sqrt{2}\,\mathrm{hb}_3-6\sqrt{2}\,\mathrm{hb}_4, &
\mathrm{Hx}\cdot\mathrm{hb}_4&=-12\sqrt{2}\,\mathrm{hb}_3+(12+18\sqrt{2})\,\mathrm{hb}_4,\\
\mathrm{Hx}\cdot\mathrm{hb}_6&=(12+6\sqrt{2})\,\mathrm{hb}_6.
\end{align*}
Indeed, the components $0$ and $1$ of the left sides are:
for $\mathrm{hb}_1$ and for $\mathrm{hb}_3$, $6\sqrt{2}\cdot2$ and $6\sqrt{2}\,(0-0-1)$;
for $\mathrm{hb}_2$, $6\sqrt{2}\,(0-1-1)$ and $6\sqrt{2}\,(2-1-0)+12$;
for $\mathrm{hb}_4$, $6\sqrt{2}\,(0-1-1)$ and $6\sqrt{2}\,(2+1-0)+12$;
for $\mathrm{hb}_5$, $6\sqrt{2}\,(0-1+1)=0$ and $6\sqrt{2}\,(2+1-0)+12$;
for $\mathrm{hb}_6$, $6\sqrt{2}\,(0-1+1)=0$ and $6\sqrt{2}\,(2-1-0)+12$.
Hence $\mathrm{Hx}$ has the eigenvalues $12+18\sqrt{2}$ and $12+6\sqrt{2}$, together with the two roots of $\mathrm{qp}_1$ and the two roots of $\mathrm{qp}_3$, the characteristic polynomials of the two $2\times2$ blocks, which, computed from their traces and determinants, are
\begin{align*}
\mathrm{qp}_1(\mathrm{xv})&:=\mathrm{xv}^2-(12+18\sqrt{2})\,\mathrm{xv}+144\sqrt{2},\\
\mathrm{qp}_3(\mathrm{xv})&:=\mathrm{xv}^2-(12+30\sqrt{2})\,\mathrm{xv}+288+144\sqrt{2};
\end{align*}
the determinants are
\begin{align*}
12\sqrt{2}\,(12+6\sqrt{2})-(12\sqrt{2})(6\sqrt{2})&=144\sqrt{2},\\
12\sqrt{2}\,(12+18\sqrt{2})-(12\sqrt{2})(6\sqrt{2})&=288+144\sqrt{2}.
\end{align*}
The polynomial $\mathrm{qp}_1$ has the roots
\[
6+9\sqrt{2}\pm\left((6+9\sqrt{2})^2-144\sqrt{2}\right)^{\frac{1}{2}}=6+9\sqrt{2}\pm\sqrt{198-36\sqrt{2}},
\]
the smaller of which is the number $\lambda_1$ of (C0r).

\emph{Comparison.}
The polynomial $\mathrm{qp}_1$ is decreasing on $(-\infty,6+9\sqrt{2}]$, $\mathrm{qp}_1(6)=36(\sqrt{2}-1)>0$, and $\mathrm{qp}_1(7)=18\sqrt{2}-35<0$ (because $648<1225$).
Therefore $6<\lambda_1<7$, and the other root of $\mathrm{qp}_1$ is larger than $7$.
The polynomial $\mathrm{qp}_3$ is decreasing on $(-\infty,6+15\sqrt{2}]$, and $\mathrm{qp}_3(7)=253-66\sqrt{2}>0$ (because $8712<64009$);
so both roots of $\mathrm{qp}_3$ are larger than $7$.
All other non-zero eigenvalues obtained in the four cases are $6\sqrt{2}=\sqrt{72}>7$, $18\sqrt{2}$, $24\sqrt{2}$, $12+6\sqrt{2}$, $12+18\sqrt{2}$, and numbers $\ge12$.
It follows that the smallest non-zero eigenvalue of $(\nabla^{*}\nabla)(0)$ is $\lambda_1$.
As $(\nabla^{*}\nabla)(0)$ is Hermitian and its kernel is $J\mathbb{C}^4=\mathrm{Ran}\,\mathrm{Plt}$,
\begin{equation*}
\left(u,(\nabla^{*}\nabla)(0)u\right)\ge\lambda_1\,(u,u)\qquad\text{for }u\in\mathrm{Ran}\,\mathrm{Phv}. \tag*{[1.6]}
\end{equation*}
\end{proof}

\paragraph{1.3. Clause (c).}
\begin{proof}[Proof of Proposition~37.1\,(c)]
Let $\xi\in\mathbb{C}^4$.
Since $R_iJ\xi=J\xi$ for $i\le2$ and $R_3(0)J\xi=J\xi$, every difference in (I.2.6), taken at $\boldsymbol{\mu}=0$ and $p=0$, vanishes on $J\xi$;
that is, $D_0^{(0)}J=0$.
For $u\in{V_{\mathrm{v}}}$, clause (a) together with (w1), (w2), (w3) gives
\begin{align*}
48\,J^{*}(\nabla^{*}\nabla)(0)u&=\sum_iw_i\sum_g\left(u(g)-u(gs_i)\right)=0,\\
48\,J^{*}\mathrm{Gh}(0)u&=\sum_iw_i\sum_g\gamma\cdot\hat{d}_i(g)\,u(gs_i)=\sum_iw_i\sum_g\gamma\cdot\hat{d}_i(gs_i)\,u(g)\\
&=-\sum_g\gamma\cdot\left(\sum_iw_i\hat{d}_i(g)\right)u(g)=0.
\end{align*}
Hence $J^{*}D_0^{(0)}=\frac{\rho}{2}J^{*}(\nabla^{*}\nabla)(0)+\frac{1}{2}J^{*}\mathrm{Gh}(0)=0$.
For every $p\in\mathbb{C}^3$, (I.2.6) gives
\begin{align*}
(D_0^{(p)}J\xi)(g)&=-M_3(g)\left(\exp\left(ip\cdot(ge_1)\right)-1\right)\xi\\
&=3\left(\rho-\gamma\cdot(ge_1)\right)\left(1-\exp\left(ip\cdot(ge_1)\right)\right)\xi .
\end{align*}
By (w4), in the mean over the $48$ labels each of the six values $\pm e_\mu$ of $ge_1$ carries the weight $3\cdot\frac{8}{48}=\frac{1}{2}$;
therefore
\begin{align*}
J^{*}D_0^{(p)}J&=\frac{1}{2}\sum_\mu\left[(\rho-\gamma_\mu)\left(1-\exp(ip_\mu)\right)+(\rho+\gamma_\mu)\left(1-\exp(-ip_\mu)\right)\right]\\
&=\sum_\mu\left[\rho\,(1-\cos p_\mu)+i\gamma_\mu\sin p_\mu\right].
\end{align*}
At $\boldsymbol{\mu}=0$ this is the formula for $J^{*}D^{(p)}J$ at the end of (I.O8).
\end{proof}

\paragraph{1.4. Clause (d).}
\begin{proof}[Proof of Proposition~37.1\,(d)]
If $ge_1=\pm e_\nu$, then $M_3(g)=3(\rho\mp\gamma_\nu)$ and $\exp\left(ip\cdot(ge_1)\right)=\exp(\pm ip_\nu)$.
Every $g$ satisfies exactly one of the six conditions $ge_1=\pm e_\nu$.
Hence
\[
M_3(g)\exp\left(ip\cdot(ge_1)\right)u(gs_3)=\sum_\nu\left(\left(\exp(ip_\nu)\,\mathrm{Ehp}^{+}_{\nu}+\exp(-ip_\nu)\,\mathrm{Ehp}^{-}_{\nu}\right)u\right)(g),
\]
and (I.2.6) at $\boldsymbol{\mu}=0$ takes the form
\begin{equation*}
\begin{aligned}
&D_0^{(p)}=D_{0,\mathrm{loc}}-\sum_\nu\left(\exp(ip_\nu)\,\mathrm{Ehp}^{+}_{\nu}+\exp(-ip_\nu)\,\mathrm{Ehp}^{-}_{\nu}\right),\\
&(D_{0,\mathrm{loc}}u)(g):=-\sum_{i\le2}M_i(g)\left(u(gs_i)-u(g)\right)+M_3(g)\,u(g).
\end{aligned}
\tag*{[1.7]}
\end{equation*}
Given complex numbers $z^{\pm}_{\nu}$, the operator $\sum_\nu\left(z^{+}_{\nu}\mathrm{Ehp}^{+}_{\nu}+z^{-}_{\nu}\mathrm{Ehp}^{-}_{\nu}\right)$ acts as $u\mapsto\mathrm{Mx}(g)\,u(gs_3)$ with the $4\times4$ matrices $\mathrm{Mx}(g)=3z^{\pm}_{\nu}(\rho\mp\gamma_\nu)$ for $ge_1=\pm e_\nu$.
Now (w6) gives $|\mathrm{Mx}(g)|_2\le3(1+\rho)\max|z^{\pm}_{\nu}|$, and (w7) turns this into the same bound for the norm of the operator.
By (w4), $J^{*}\mathrm{Ehp}^{\pm}_{\nu}J=\frac{8}{48}\cdot3(\rho\mp\gamma_\nu)=\frac{1}{2}(\rho\mp\gamma_\nu)$.
As $D^{(p)}=\mathbf{m}+D_0^{(p)}$, differentiation of [1.7] gives, for $j\ge1$,
\begin{equation*}
\partial_\nu^jD^{(p)}=-\left(i^j\exp(ip_\nu)\,\mathrm{Ehp}^{+}_{\nu}+(-i)^j\exp(-ip_\nu)\,\mathrm{Ehp}^{-}_{\nu}\right); \tag*{[1.8]}
\end{equation*}
for $j=1$ and $j=2$ these are the two formulas of clause (d), for $\mathrm{Vx}_\nu(p)$ and for $\partial_\nu^2D^{(p)}$.
Since every term of [1.7] involves a single coordinate of $p$, mixed derivatives are zero.
Further, $J^{*}\mathrm{Vx}_\nu(0)J=-i\left(\frac{1}{2}(\rho-\gamma_\nu)-\frac{1}{2}(\rho+\gamma_\nu)\right)=i\gamma_\nu$.
When $p$ is real, the coefficients in [1.8] are of modulus one, whence $|\partial^{\mathrm{ak}}D^{(p)}|_2\le3(1+\rho)\le6$ for $\mathrm{ak}\ne0$.
For real $p$ and $p'$,
\begin{align*}
D^{(p)}-D^{(p')}&=-\sum_\nu\left(\exp(ip_\nu)-\exp(ip'_\nu)\right)\mathrm{Ehp}^{+}_{\nu}\\
&\quad-\sum_\nu\left(\exp(-ip_\nu)-\exp(-ip'_\nu)\right)\mathrm{Ehp}^{-}_{\nu},
\end{align*}
and every coefficient is bounded in modulus by $|p_\nu-p'_\nu|$;
whence $|D^{(p)}-D^{(p')}|_2\le3(1+\rho)\,|p-p'|_\infty\le6\,|p-p'|_\infty$.
Moreover $(\nabla^{*}\nabla)(p)-(\nabla^{*}\nabla)(p')=-w_3\left(R_3(p)-R_3(p')\right)$ with $w_3=6$,
\[
\left(\left(R_3(p)-R_3(p')\right)u\right)(g)=\left(\exp\left(ip\cdot(ge_1)\right)-\exp\left(ip'\cdot(ge_1)\right)\right)u(gs_3),
\]
and $\left|\exp\left(ip\cdot(ge_1)\right)-\exp\left(ip'\cdot(ge_1)\right)\right|\le|(p-p')\cdot(ge_1)|\le|p-p'|_\infty$.
By (w7),
\begin{equation*}
|(\nabla^{*}\nabla)(p)-(\nabla^{*}\nabla)(p')|_2\le6\,|p-p'|_\infty\qquad(p,p'\ \text{real}). \tag*{[1.9]}
\end{equation*}
\end{proof}

\paragraph{1.5. Clause (e).}
\begin{proof}[Proof of Proposition~37.1\,(e)]
The matrix $-1$ belongs to $O_{\mathrm{cub}}$, so that $g\mapsto-g$ is a bijection of $O_{\mathrm{cub}}$;
moreover $-1$ commutes with every element of $O_{\mathrm{cub}}$.
By (I.c1), $\gamma_5$ is Hermitian, $\gamma_5^2=1$ and $\gamma_5\gamma_\mu\gamma_5=-\gamma_\mu$.
Consequently $\mathrm{Pa}$ preserves the norm, and $(\mathrm{Pa}^2u)(g)=\gamma_5\gamma_5\,u(g)=u(g)$;
so $\mathrm{Pa}$ is a unitary involution.
Since $M_i(-g)=\frac{w_i}{2}\left(\rho+\gamma\cdot(-g\hat n_i)\right)$, one has
\[
\gamma_5M_i(-g)\gamma_5=\frac{w_i}{2}\left(\rho+\gamma\cdot(g\hat n_i)\right)=M_i(g),
\]
and $(\mathrm{Pa}\,u)(-gs_i)=\gamma_5u(gs_i)$.
Together with (I.2.6) this gives, for every $p\in\mathbb{C}^3$,
\begin{align*}
\left(\mathrm{Pa}\cdot D^{(p)}\cdot\mathrm{Pa}\,u\right)(g)
&=\gamma_5\left(D^{(p)}\cdot\mathrm{Pa}\,u\right)(-g)\\
&=\gamma_5\left[\mathbf{m}\gamma_5u(g)-\sum_{i\le2}M_i(-g)\gamma_5\left(u(gs_i)-u(g)\right)\right.\\
&\qquad\left.{}-M_3(-g)\gamma_5\left(\exp\left(ip\cdot(-ge_1)\right)u(gs_3)-u(g)\right)\right]\\
&=\mathbf{m}u(g)-\sum_{i\le2}M_i(g)\left(u(gs_i)-u(g)\right)\\
&\qquad{}-M_3(g)\left(\exp\left(-ip\cdot(ge_1)\right)u(gs_3)-u(g)\right)\\
&=\left(D^{(-p)}u\right)(g).
\end{align*}
Differentiating with respect to $p_\nu$ one obtains $\mathrm{Pa}\cdot\mathrm{Vx}_\nu(p)\cdot\mathrm{Pa}=-\mathrm{Vx}_\nu(-p)$.
\end{proof}

\paragraph{1.6. Clause (f).}
\begin{proof}[Proof of Proposition~37.1\,(f)]
\emph{Constants.}
Let
\begin{equation*}
\begin{gathered}
\mathrm{kah}:=\frac{\rho\lambda_1}{4},\qquad
\mathrm{Crs}:=1+\frac{36}{\mathrm{kah}},\qquad
\delta(\rho):=\min\left\{\frac{\lambda_1}{12},\frac{1}{2\,\mathrm{Crs}}\right\}\quad\left(\le\frac{1}{2}\right),\\
\mathrm{lmn}(p):=\min\left\{\left(u,(\nabla^{*}\nabla)(p)u\right):|u|_2=1\right\}\quad(p\in\mathbb{R}^3),\\
\mathrm{lmin}:=\min\left\{\mathrm{lmn}(p):p\in[-\pi,\pi]^3,\ |p|_2\ge\delta(\rho)\right\},\\
C_E:=\max\left\{2\,\mathrm{Crs}+\frac{25}{\mathrm{kah}}+\frac{36}{\mathrm{kah}^2},\ \frac{2}{\rho\cdot\mathrm{lmin}}+\frac{1}{\delta(\rho)}\right\}.
\end{gathered}
\tag*{[1.10]}
\end{equation*}
For real $p$ the number $\mathrm{lmn}(p)$ is the lowest eigenvalue of $(\nabla^{*}\nabla)(p)$.

The number $\mathrm{lmin}$ is positive.
Indeed, the lowest eigenvalue of a Hermitian matrix moves by at most the norm of a Hermitian perturbation, so that [1.9] gives $|\mathrm{lmn}(p)-\mathrm{lmn}(p')|\le6\,|p-p'|_\infty$;
thus $\mathrm{lmn}$ is continuous.
The set $\left\{p\in[-\pi,\pi]^3:|p|_2\ge\delta(\rho)\right\}$ is compact and non-empty, and it contains no point of $2\pi\mathbb{Z}^3$, since $0$ is the only point of $2\pi\mathbb{Z}^3$ in $[-\pi,\pi]^3$;
by (b), $\mathrm{lmn}>0$ on this set.
The numbers $\mathrm{kah}$, $\mathrm{Crs}$, $\delta(\rho)$, $\mathrm{lmin}$ and $C_E$ depend on $\rho$ alone.

Now let $\mathbf{m}>0$ and $p\in[-\pi,\pi]^3$.
By (a) and (w8), $D^{(p)}$ is invertible.
Abbreviate $\mathrm{Sdr}:=\mathrm{Sdr}(p,\mathbf{m})$, and put $\mathrm{Er}:=\mathrm{Rsy}(p)-J\cdot\mathrm{Sdr}\cdot J^{*}$.

\emph{First case: $|p|_2\le\delta(\rho)$.}
Set
\[
\mathrm{mm}_{\mathrm{ll}}:=J^{*}D^{(p)}J,\qquad
\mathrm{mm}_{\mathrm{lh}}:=J^{*}D^{(p)}\cdot\mathrm{Phv},\qquad
\mathrm{mm}_{\mathrm{hl}}:=\mathrm{Phv}\cdot D^{(p)}J,
\]
and let $\mathrm{mm}_{\mathrm{hh}}$ be the restriction of $\mathrm{Phv}\cdot D^{(p)}\cdot\mathrm{Phv}$ to $\mathrm{Ran}\,\mathrm{Phv}$.
\begin{itemize}
\item[(1)] By (c), $J^{*}D^{(0)}=\mathbf{m}J^{*}$ and $D^{(0)}J=\mathbf{m}J$.
As $J^{*}\cdot\mathrm{Phv}=0=\mathrm{Phv}\cdot J$ by (w5), it follows that
\[
\mathrm{mm}_{\mathrm{lh}}=J^{*}\left(D^{(p)}-D^{(0)}\right)\mathrm{Phv},\qquad
\mathrm{mm}_{\mathrm{hl}}=\mathrm{Phv}\left(D^{(p)}-D^{(0)}\right)J.
\]
By (d), $|\mathrm{mm}_{\mathrm{lh}}|_2\le6\,|p|_\infty\le6\,|p|_2$ and $|\mathrm{mm}_{\mathrm{hl}}|_2\le6\,|p|_\infty\le6\,|p|_2$.
\item[(2)] Let $u\in\mathrm{Ran}\,\mathrm{Phv}$.
From (a), [1.6], [1.9] and $|p|_\infty\le|p|_2\le\lambda_1/12$ one gets
\begin{align*}
\mathrm{Re}\,(u,\mathrm{mm}_{\mathrm{hh}}u)&=\mathbf{m}\,(u,u)+\frac{\rho}{2}\left(u,(\nabla^{*}\nabla)(p)u\right)\\
&\ge\frac{\rho}{2}\left(\lambda_1-6\,|p|_\infty\right)(u,u)\ge\frac{\rho}{2}\cdot\frac{\lambda_1}{2}\,(u,u)=\mathrm{kah}\,(u,u).
\end{align*}
By (w8), $\mathrm{mm}_{\mathrm{hh}}$ is invertible on $\mathrm{Ran}\,\mathrm{Phv}$, and $|\mathrm{mm}_{\mathrm{hh}}^{-1}|_2\le1/\mathrm{kah}$.
\item[(3)] By (c), $\mathrm{mm}_{\mathrm{ll}}=\mathbf{m}+i\gamma\cdot\sin p+\rho\sum_\mu(1-\cos p_\mu)$.
Define
\begin{gather*}
\mathrm{ssc}:=\mathrm{mm}_{\mathrm{ll}}-\mathrm{mm}_{\mathrm{lh}}\cdot\mathrm{mm}_{\mathrm{hh}}^{-1}\cdot\mathrm{mm}_{\mathrm{hl}}=\mathbf{m}+i\gamma\cdot p+\mathrm{rsc},\\
\mathrm{rsc}:=i\gamma\cdot(\sin p-p)+\rho\sum_\mu(1-\cos p_\mu)-\mathrm{mm}_{\mathrm{lh}}\cdot\mathrm{mm}_{\mathrm{hh}}^{-1}\cdot\mathrm{mm}_{\mathrm{hl}}.
\end{gather*}
By (w6), $|\gamma\cdot(\sin p-p)|_2=|\sin p-p|_2\le\sum_\mu|p_\mu|^3/6\le|p|_2^3/6$;
further $0\le\rho\sum_\mu(1-\cos p_\mu)\le|p|_2^2/2$.
With $|p|_2\le1/2$ and with (1), (2) it follows that
\[
|\mathrm{rsc}|_2\le\left(\frac{1}{12}+\frac{1}{2}\right)|p|_2^2+\frac{36\,|p|_2^2}{\mathrm{kah}}\le\mathrm{Crs}\cdot|p|_2^2.
\]
\item[(4)] By (w6), $|\mathrm{rsc}\cdot\mathrm{Sdr}|_2\le\mathrm{Crs}\cdot|p|_2\left(|p|_2\cdot|\mathrm{Sdr}|_2\right)\le\mathrm{Crs}\cdot|p|_2\le1/2$, because $|p|_2\le\delta(\rho)\le1/(2\,\mathrm{Crs})$.
Hence $1+\mathrm{rsc}\cdot\mathrm{Sdr}$ is invertible and $\left|(1+\mathrm{rsc}\cdot\mathrm{Sdr})^{-1}\right|_2\le2$;
the matrix $\mathrm{ssc}=(1+\mathrm{rsc}\cdot\mathrm{Sdr})(\mathbf{m}+i\gamma\cdot p)$ is invertible, $\mathrm{ssc}^{-1}=\mathrm{Sdr}\cdot(1+\mathrm{rsc}\cdot\mathrm{Sdr})^{-1}$, and
\begin{gather*}
|\mathrm{ssc}^{-1}|_2\le2\,|\mathrm{Sdr}|_2,\qquad|p|_2\cdot|\mathrm{ssc}^{-1}|_2\le2,\\
|\mathrm{ssc}^{-1}-\mathrm{Sdr}|_2=\left|\mathrm{Sdr}\cdot(1+\mathrm{rsc}\cdot\mathrm{Sdr})^{-1}\cdot\mathrm{rsc}\cdot\mathrm{Sdr}\right|_2\le2\,\mathrm{Crs}\cdot|p|_2^2\cdot|\mathrm{Sdr}|_2^2\le2\,\mathrm{Crs}.
\end{gather*}
\item[(5)] (Block inversion.)
With respect to the orthogonal decomposition ${V_{\mathrm{v}}}=J\mathbb{C}^4\oplus\mathrm{Ran}\,\mathrm{Phv}$ one has
\[
D^{(p)}=J\cdot\mathrm{mm}_{\mathrm{ll}}\cdot J^{*}+J\cdot\mathrm{mm}_{\mathrm{lh}}+\mathrm{mm}_{\mathrm{hl}}\cdot J^{*}+\mathrm{mm}_{\mathrm{hh}}\cdot\mathrm{Phv}.
\]
Since $\mathrm{mm}_{\mathrm{hh}}$ and $\mathrm{ssc}$ are invertible,
\begin{equation*}
\begin{aligned}
\mathrm{Rsy}(p)={}&J\cdot\mathrm{ssc}^{-1}\cdot J^{*}-J\cdot\mathrm{ssc}^{-1}\cdot\mathrm{mm}_{\mathrm{lh}}\cdot\mathrm{mm}_{\mathrm{hh}}^{-1}\cdot\mathrm{Phv}-\mathrm{mm}_{\mathrm{hh}}^{-1}\cdot\mathrm{mm}_{\mathrm{hl}}\cdot\mathrm{ssc}^{-1}\cdot J^{*}\\
&+\mathrm{mm}_{\mathrm{hh}}^{-1}\cdot\mathrm{Phv}+\mathrm{mm}_{\mathrm{hh}}^{-1}\cdot\mathrm{mm}_{\mathrm{hl}}\cdot\mathrm{ssc}^{-1}\cdot\mathrm{mm}_{\mathrm{lh}}\cdot\mathrm{mm}_{\mathrm{hh}}^{-1}\cdot\mathrm{Phv}.
\end{aligned}
\tag*{[1.11]}
\end{equation*}
To verify [1.11], apply $D^{(p)}$ from the left to its right-hand side and use $J^{*}J=1$, $J^{*}\cdot\mathrm{mm}_{\mathrm{hh}}^{-1}=0$, $\mathrm{Phv}\cdot J=0$ and $\mathrm{mm}_{\mathrm{lh}}=\mathrm{mm}_{\mathrm{lh}}\cdot\mathrm{Phv}$.
The terms ending with $J^{*}$ add up to
\[
J\left(\mathrm{mm}_{\mathrm{ll}}-\mathrm{mm}_{\mathrm{lh}}\cdot\mathrm{mm}_{\mathrm{hh}}^{-1}\cdot\mathrm{mm}_{\mathrm{hl}}\right)\mathrm{ssc}^{-1}J^{*}+\left(\mathrm{mm}_{\mathrm{hl}}-\mathrm{mm}_{\mathrm{hl}}\right)\mathrm{ssc}^{-1}J^{*}=JJ^{*},
\]
and the terms ending with $\mathrm{Phv}$ add up to
\begin{multline*}
J\left(-\mathrm{mm}_{\mathrm{ll}}+\mathrm{mm}_{\mathrm{lh}}\cdot\mathrm{mm}_{\mathrm{hh}}^{-1}\cdot\mathrm{mm}_{\mathrm{hl}}\right)\mathrm{ssc}^{-1}\cdot\mathrm{mm}_{\mathrm{lh}}\cdot\mathrm{mm}_{\mathrm{hh}}^{-1}\cdot\mathrm{Phv}+J\cdot\mathrm{mm}_{\mathrm{lh}}\cdot\mathrm{mm}_{\mathrm{hh}}^{-1}\cdot\mathrm{Phv}\\
+\left(-\mathrm{mm}_{\mathrm{hl}}+\mathrm{mm}_{\mathrm{hl}}\right)\mathrm{ssc}^{-1}\cdot\mathrm{mm}_{\mathrm{lh}}\cdot\mathrm{mm}_{\mathrm{hh}}^{-1}\cdot\mathrm{Phv}+\mathrm{Phv}=\mathrm{Phv}.
\end{multline*}
The total is $JJ^{*}+\mathrm{Phv}=1$.
\item[(6)] By [1.11], $\mathrm{Er}$ is the sum of $J\left(\mathrm{ssc}^{-1}-\mathrm{Sdr}\right)J^{*}$ and of the last four terms of [1.11].
By (w5), (1), (2), (4) and $|p|_2\le1/2$,
\begin{align*}
|\mathrm{Er}|_2&\le2\,\mathrm{Crs}+\frac{2\left(2\,|\mathrm{Sdr}|_2\right)\left(6\,|p|_2\right)}{\mathrm{kah}}+\frac{1}{\mathrm{kah}}+\frac{\left(6\,|p|_2\right)\left(2\,|\mathrm{Sdr}|_2\right)\left(6\,|p|_2\right)}{\mathrm{kah}^2}\\
&\le2\,\mathrm{Crs}+\frac{24}{\mathrm{kah}}+\frac{1}{\mathrm{kah}}+\frac{72\,|p|_2}{\mathrm{kah}^2}\le2\,\mathrm{Crs}+\frac{25}{\mathrm{kah}}+\frac{36}{\mathrm{kah}^2}\le C_E.
\end{align*}
\end{itemize}

\emph{Second case: $|p|_2\ge\delta(\rho)$.}
By (a), $\mathrm{Re}\,(u,D^{(p)}u)\ge\frac{\rho}{2}\,\mathrm{lmn}(p)\,(u,u)\ge\frac{\rho\cdot\mathrm{lmin}}{2}\,(u,u)$, and (w8) yields $|\mathrm{Rsy}(p)|_2\le2/(\rho\cdot\mathrm{lmin})$.
By (w5) and (w6), $|J\cdot\mathrm{Sdr}\cdot J^{*}|_2=|\mathrm{Sdr}|_2\le1/|p|_2\le1/\delta(\rho)$.
Hence $|\mathrm{Er}|_2\le2/(\rho\cdot\mathrm{lmin})+1/\delta(\rho)\le C_E$.

\emph{The two consequences.}
In both cases
\[
|\mathrm{Rsy}(p)|_2\le|\mathrm{Er}|_2+|J\cdot\mathrm{Sdr}\cdot J^{*}|_2\le C_E+\left(\mathbf{m}^2+|p|_2^2\right)^{-\frac{1}{2}}.
\]
The inequality $\left(|p|_2+\mathbf{m}\right)^2\le2\left(|p|_2^2+\mathbf{m}^2\right)$ gives $\left(\mathbf{m}^2+|p|_2^2\right)^{-\frac{1}{2}}\le\sqrt{2}/\left(|p|_2+\mathbf{m}\right)$.
If $\mathbf{m}\le1$, then $|p|_2+\mathbf{m}\le\sqrt{3}\,\pi+1<7$, so that $C_E\le7C_E/\left(|p|_2+\mathbf{m}\right)$;
altogether $|\mathrm{Rsy}(p)|_2\le\left(\sqrt{2}+7C_E\right)/\left(|p|_2+\mathbf{m}\right)$.
\end{proof}

\paragraph{1.7. Clause (g).}
\begin{proof}[Proof of Proposition~37.1\,(g)]
Let $0<\mathbf{m}\le1$, and for real $p$ put $\mathrm{xd}(p):=\mathrm{dist}_2(p,2\pi\mathbb{Z}^3)+\mathbf{m}$.
One has $\mathrm{dist}_2(p,2\pi\mathbb{Z}^3)^2=\sum_\nu\mathrm{dist}(p_\nu,2\pi\mathbb{Z})^2$, and $\mathrm{dist}(p_\nu,2\pi\mathbb{Z})=|p_\nu|$ whenever $|p_\nu|\le\pi$.
Hence $\mathrm{xd}(p)=|p|_2+\mathbf{m}$ on $[-\pi,\pi]^3$;
the function $\mathrm{xd}$ is $2\pi\mathbb{Z}^3$-periodic, and $\mathbf{m}\le\mathrm{xd}(p)\le\sqrt{3}\,\pi+1<7$.
The entries of $D^{(p)}$ are entire functions of $p$ (Section~I.2.6), and $D^{(p)}$ is invertible for real $p$;
therefore $\mathrm{Rsy}$ is a smooth $2\pi\mathbb{Z}^3$-periodic function on $\mathbb{R}^3$, and so are all its derivatives.
Define, by recursion on $|\mathrm{ak}|$,
\begin{equation*}
\mathrm{Csy}_0:=\sqrt{2}+7C_E,\qquad
\mathrm{Csy}_{\mathrm{ak}}:=6\,\mathrm{Csy}_0\sum_{0<\mathrm{bi}\le\mathrm{ak}}\binom{\mathrm{ak}}{\mathrm{bi}}\cdot7^{|\mathrm{bi}|-1}\cdot\mathrm{Csy}_{\mathrm{ak}-\mathrm{bi}}\quad(\mathrm{ak}\ne0),
\tag*{[1.12]}
\end{equation*}
where $\mathrm{bi}$ ranges over the multi-indices with $\mathrm{bi}\le\mathrm{ak}$ componentwise and $\mathrm{bi}\ne0$.

It is claimed that
\[
\left|\partial^{\mathrm{ak}}\mathrm{Rsy}(p)\right|_2\le\mathrm{Csy}_{\mathrm{ak}}\cdot\mathrm{xd}(p)^{-1-|\mathrm{ak}|}\qquad\text{for all real }p.
\]
For $\mathrm{ak}=0$ and $p\in[-\pi,\pi]^3$ this is the last bound of (f), and by periodicity it holds for all real $p$.
Let $\mathrm{ak}\ne0$, and assume the claim for all orders smaller than $|\mathrm{ak}|$.
Apply $\partial^{\mathrm{ak}}$ to the identity $\mathrm{Rsy}(p)D^{(p)}=1$ by Leibniz' rule, and multiply by $\mathrm{Rsy}(p)$ from the right;
this gives
\begin{equation*}
\partial^{\mathrm{ak}}\mathrm{Rsy}=-\sum_{0<\mathrm{bi}\le\mathrm{ak}}\binom{\mathrm{ak}}{\mathrm{bi}}\left(\partial^{\mathrm{ak}-\mathrm{bi}}\mathrm{Rsy}\right)\left(\partial^{\mathrm{bi}}D^{(p)}\right)\mathrm{Rsy}.
\tag*{[1.13]}
\end{equation*}
By (d) and by the claim for $\mathrm{ak}-\mathrm{bi}$ and for $0$, the term with index $\mathrm{bi}$ has norm at most
\begin{align*}
&\binom{\mathrm{ak}}{\mathrm{bi}}\cdot\mathrm{Csy}_{\mathrm{ak}-\mathrm{bi}}\cdot\mathrm{xd}^{-1-|\mathrm{ak}|+|\mathrm{bi}|}\cdot6\,\mathrm{Csy}_0\cdot\mathrm{xd}^{-1}\\
&\qquad=6\binom{\mathrm{ak}}{\mathrm{bi}}\cdot\mathrm{Csy}_{\mathrm{ak}-\mathrm{bi}}\cdot\mathrm{Csy}_0\cdot\mathrm{xd}^{|\mathrm{bi}|-1}\cdot\mathrm{xd}^{-1-|\mathrm{ak}|},
\end{align*}
and $\mathrm{xd}^{|\mathrm{bi}|-1}\le7^{|\mathrm{bi}|-1}$.
Summation over $\mathrm{bi}$ gives the claim.
The numbers $\mathrm{Csy}_{\mathrm{ak}}$ depend on $\rho$ alone.
\end{proof}

\paragraph{1.8. Clause (h).}
\begin{proof}[Proof of Proposition~37.1\,(h)]
Let $\bar m\in(0,1]$, let $\eta_j\in(0,1]$ with $\eta_j\to0$, and let $\mathrm{kph}_j\to\mathrm{kph}$ in $\mathbb{R}^3$.
At the index $j$ the matrices $D^{(p)}$, $\mathrm{Rsy}(p)$, $\mathrm{Er}(p,\mathbf{m})$ are formed with $\mathbf{m}=\bar m\eta_j$, a number in $(0,1]$;
the matrices $\mathrm{Vx}_\mu(p)$ and $\partial_\mu^2D^{(p)}$ do not contain $\mathbf{m}$.
For all sufficiently large $j$ the point $\eta_j\mathrm{kph}_j$ lies in $[-\pi,\pi]^3$.
For $\eta>0$ and $\mathrm{kph}'\in\mathbb{R}^3$,
\[
\eta\cdot\mathrm{Sdr}(\eta\,\mathrm{kph}',\bar m\eta)=\eta\left(\bar m\eta+i\eta\gamma\cdot\mathrm{kph}'\right)^{-1}=\mathrm{Sdr}(\mathrm{kph}',\bar m).
\]
Therefore (f) gives
\[
\left|\eta_j\mathrm{Rsy}(\eta_j\mathrm{kph}_j)-J\cdot\mathrm{Sdr}(\mathrm{kph}_j,\bar m)\cdot J^{*}\right|_2=\eta_j\left|\mathrm{Er}(\eta_j\mathrm{kph}_j,\bar m\eta_j)\right|_2\le C_E\eta_j\to0.
\]
The map $\mathrm{kph}'\mapsto\mathrm{Sdr}(\mathrm{kph}',\bar m)$ is continuous on $\mathbb{R}^3$.
Hence $\eta_j\mathrm{Rsy}(\eta_j\mathrm{kph}_j)\to J\cdot\mathrm{Sdr}(\mathrm{kph},\bar m)\cdot J^{*}$.
For $|\mathrm{ak}|=1,2$, formula [1.13] together with (d), according to which the mixed second derivatives of $D^{(p)}$ vanish, reads
\begin{align*}
\partial_\mu\mathrm{Rsy}&=-\mathrm{Rsy}\cdot\mathrm{Vx}_\mu\cdot\mathrm{Rsy},\\
\partial_\mu\partial_\nu\mathrm{Rsy}&=\mathrm{Rsy}\cdot\mathrm{Vx}_\mu\cdot\mathrm{Rsy}\cdot\mathrm{Vx}_\nu\cdot\mathrm{Rsy}+\mathrm{Rsy}\cdot\mathrm{Vx}_\nu\cdot\mathrm{Rsy}\cdot\mathrm{Vx}_\mu\cdot\mathrm{Rsy}\\
&\quad-\delta_{\mu\nu}\,\mathrm{Rsy}\cdot\left(\partial_\mu^2D^{(p)}\right)\cdot\mathrm{Rsy}.
\end{align*}
As $\mathrm{Vx}_\mu$ is continuous, $\mathrm{Vx}_\mu(\eta_j\mathrm{kph}_j)\to\mathrm{Vx}_\mu(0)$;
moreover $J^{*}\mathrm{Vx}_\mu(0)J=i\gamma_\mu$, and $\left|\partial_\mu^2D^{(p)}\right|_2\le6$ by (d).
Write $\mathrm{Sdr}$ for $\mathrm{Sdr}(\mathrm{kph},\bar m)$, and evaluate $\mathrm{Rsy}$, $\mathrm{Vx}_\mu$, $\mathrm{Vx}_\nu$ and $\partial_\mu^2D^{(p)}$ at $\eta_j\mathrm{kph}_j$.
Then
\begin{align*}
\eta_j^2\left(\partial_\mu\mathrm{Rsy}\right)(\eta_j\mathrm{kph}_j)&=-\left(\eta_j\mathrm{Rsy}\right)\cdot\mathrm{Vx}_\mu\cdot\left(\eta_j\mathrm{Rsy}\right)\\
&\to-J\cdot\mathrm{Sdr}\cdot J^{*}\,\mathrm{Vx}_\mu(0)\,J\cdot\mathrm{Sdr}\cdot J^{*}=-J\cdot\mathrm{Sdr}\cdot i\gamma_\mu\cdot\mathrm{Sdr}\cdot J^{*},\\
\eta_j^3\left(\partial_\mu\partial_\nu\mathrm{Rsy}\right)(\eta_j\mathrm{kph}_j)&\to J\left(\mathrm{Sdr}\cdot i\gamma_\mu\cdot\mathrm{Sdr}\cdot i\gamma_\nu\cdot\mathrm{Sdr}+\mathrm{Sdr}\cdot i\gamma_\nu\cdot\mathrm{Sdr}\cdot i\gamma_\mu\cdot\mathrm{Sdr}\right)J^{*},
\end{align*}
since the third term equals $\eta_j\left(\eta_j\mathrm{Rsy}\right)\left(\partial_\mu^2D^{(p)}\right)\left(\eta_j\mathrm{Rsy}\right)$, whose norm is at most $6\,\eta_j\left|\eta_j\mathrm{Rsy}\right|_2^2\to0$.
Differentiation of $\mathrm{Sdr}(\mathrm{kph},\bar m)\left(\bar m+i\gamma\cdot\mathrm{kph}\right)=1$ with respect to $\mathrm{kph}$ gives
\begin{align*}
\partial_\mu\mathrm{Sdr}&=-\mathrm{Sdr}\cdot i\gamma_\mu\cdot\mathrm{Sdr},\\
\partial_\mu\partial_\nu\mathrm{Sdr}&=\mathrm{Sdr}\cdot i\gamma_\mu\cdot\mathrm{Sdr}\cdot i\gamma_\nu\cdot\mathrm{Sdr}+\mathrm{Sdr}\cdot i\gamma_\nu\cdot\mathrm{Sdr}\cdot i\gamma_\mu\cdot\mathrm{Sdr};
\end{align*}
these are the limits found above.
\end{proof}

\paragraph{1.9. The two sentences after (h).}
\begin{proof}[Proof of the two sentences that follow Proposition~37.1\,(h)]
The constant $C_E$ is the number defined in [1.10];
it involves $\mathrm{lmin}$, the minimum of a continuous positive function over a compact set.
The numbers $\mathrm{Csy}_{\mathrm{ak}}$ are those of [1.12];
they are built from $C_E$, and hence they involve $\mathrm{lmin}$ as well.
\begin{itemize}
\item[(i)] \emph{Every number $C_E$ for which (f) holds satisfies $C_E\ge1/(\rho\cdot\mathrm{sw})$.}
Let $\xi\in\mathbb{C}^4$, $\xi\ne0$, and $u(g):=\det(g)\,\xi$.
By (I.s5), $\det(gs_i)=-\det(g)$.
Hence $R_iu=-u$ for $i\le2$, and $R_3(0)u=-u$.
It follows that $(\nabla^{*}\nabla)(0)u=2\,\mathrm{sw}\cdot u$, that
\[
(\mathrm{Gh}(0)u)(g)=-\det(g)\,\gamma\cdot\left(\sum_iw_i\hat d_i(g)\right)\xi=0
\]
by (w3), and that $J^{*}u=\frac{1}{48}\sum_g\det(g)\,\xi=0$, because $g\mapsto gs_0$ reverses the sign of the determinant.
By (a), $D^{(0)}u=\left(\mathbf{m}+\rho\cdot\mathrm{sw}\right)u$.
Consequently
\[
\mathrm{Er}(0,\mathbf{m})u=\mathrm{Rsy}(0)u-J\cdot\mathrm{Sdr}(0,\mathbf{m})\cdot J^{*}u=\frac{u}{\mathbf{m}+\rho\cdot\mathrm{sw}},
\]
so that $|\mathrm{Er}(0,\mathbf{m})|_2\ge1/\left(\mathbf{m}+\rho\cdot\mathrm{sw}\right)$ for every $\mathbf{m}>0$;
the supremum of the right-hand side over $\mathbf{m}>0$ equals $1/(\rho\cdot\mathrm{sw})$.
\item[(ii)] \emph{Every number $C_E$ for which (f) holds satisfies $C_E\ge2/(\rho\lambda_1)$};
this bound exceeds $1/(\rho\cdot\mathrm{sw})$, because $\lambda_1<7<2\,\mathrm{sw}$.
Let
\[
\mathrm{Tm}:=\sum_{i=0}^{3}w_i(1-s_i)=2\sum_iw_i\hat n_i\hat n_i^{T},
\]
a real symmetric $3\times3$ matrix.
With the vectors $\hat n_i$ of Section~I.1.2 and $w=(6,6\sqrt{2},6\sqrt{2},6)$,
\[
\mathrm{Tm}=12\,e_3e_3^{T}+6\sqrt{2}\,(e_2-e_3)(e_2-e_3)^{T}+6\sqrt{2}\,(e_1-e_2)(e_1-e_2)^{T}+12\,e_1e_1^{T},
\]
that is,
\[
\mathrm{Tm}=\left(\begin{array}{ccc}12+6\sqrt{2}&-6\sqrt{2}&0\\-6\sqrt{2}&12\sqrt{2}&-6\sqrt{2}\\0&-6\sqrt{2}&12+6\sqrt{2}\end{array}\right).
\]
Put $\tilde y:=e_1+\tilde y_2e_2+e_3$ with $\tilde y_2:=\left(12+6\sqrt{2}-\lambda_1\right)/\left(6\sqrt{2}\right)$.
The first and the third component of $\mathrm{Tm}\,\tilde y$ are equal to $12+6\sqrt{2}-6\sqrt{2}\,\tilde y_2=\lambda_1$.
The second component is $12\sqrt{2}\,(\tilde y_2-1)$, and, by 1.2,
\begin{align*}
12\sqrt{2}\,(\tilde y_2-1)-\lambda_1\tilde y_2&=\frac{\left(12\sqrt{2}-\lambda_1\right)\left(12+6\sqrt{2}-\lambda_1\right)-144}{6\sqrt{2}}\\
&=\frac{\lambda_1^2-\left(12+18\sqrt{2}\right)\lambda_1+144\sqrt{2}}{6\sqrt{2}}=\frac{\mathrm{qp}_1(\lambda_1)}{6\sqrt{2}}=0.
\end{align*}
Hence $\mathrm{Tm}\,\tilde y=\lambda_1\tilde y$.

Let $\xi\in\mathbb{C}^4$, $\xi\ne0$, and $u(g):=\gamma\cdot(g\tilde y)\,\xi$.
By (I.c1), $\left(\gamma\cdot(g\tilde y)\right)^2=|\tilde y|_2^2\ne0$, so $u\ne0$.
For $i\le2$, and for $R_3(0)$ in the place of $R_3$, one has $(R_iu)(g)=\gamma\cdot(gs_i\tilde y)\,\xi$.
Therefore
\[
\left((\nabla^{*}\nabla)(0)u\right)(g)=\gamma\cdot\left(g\sum_iw_i(1-s_i)\tilde y\right)\xi=\gamma\cdot\left(g\,\mathrm{Tm}\,\tilde y\right)\xi=\lambda_1u(g).
\]
From $s_i\tilde y=\tilde y-2(\hat n_i\cdot\tilde y)\hat n_i$, $\hat d_i(g)=-g\hat n_i$ and $\left(\gamma\cdot(g\hat n_i)\right)^2=1$ one obtains
\[
\gamma\cdot\hat d_i(g)\;\gamma\cdot(gs_i\tilde y)=-\gamma\cdot(g\hat n_i)\;\gamma\cdot(g\tilde y)+2\,(\hat n_i\cdot\tilde y).
\]
Hence, by (I.1.1) applied to both terms,
\begin{align*}
(\mathrm{Gh}(0)u)(g)&=\sum_iw_i\,\gamma\cdot\hat d_i(g)\;\gamma\cdot(gs_i\tilde y)\,\xi\\
&=-\gamma\cdot\left(g\sum_iw_i\hat n_i\right)\gamma\cdot(g\tilde y)\,\xi+2\left(\left(\sum_iw_i\hat n_i\right)\cdot\tilde y\right)\xi=0.
\end{align*}
Moreover $J^{*}u=\gamma\cdot\left(\frac{1}{48}\sum_gg\tilde y\right)\xi=0$, since $g\mapsto-g$ is a bijection of $O_{\mathrm{cub}}$.
By (a), $D^{(0)}u=\left(\mathbf{m}+\frac{1}{2}\rho\lambda_1\right)u$.
Consequently
\[
\mathrm{Er}(0,\mathbf{m})u=\mathrm{Rsy}(0)u-J\cdot\mathrm{Sdr}(0,\mathbf{m})\cdot J^{*}u=\left(\mathbf{m}+\frac{1}{2}\rho\lambda_1\right)^{-1}u,
\]
and $|\mathrm{Er}(0,\mathbf{m})|_2\ge\left(\mathbf{m}+\frac{1}{2}\rho\lambda_1\right)^{-1}$ for every $\mathbf{m}>0$;
the supremum of the right-hand side over $\mathbf{m}>0$ is $2/(\rho\lambda_1)$.
\item[(iii)] By [1.12], $\mathrm{Csy}_0\ge7C_E$, and for $\mathrm{ak}\ne0$ the term $\mathrm{bi}=\mathrm{ak}$ gives $\mathrm{Csy}_{\mathrm{ak}}\ge6\cdot7^{|\mathrm{ak}|-1}\cdot\mathrm{Csy}_0^2$.
Together with (i) or (ii) it follows that $C_E$ and all the numbers $\mathrm{Csy}_{\mathrm{ak}}$ of [1.12] tend to infinity as $\rho\to0$.
\end{itemize}
\end{proof}

\subsection*{2. Proposition 37.2}

\subsubsection*{Proof of Proposition 37.2}

\paragraph{2.0. Standing data; the numbers of the proof.}
Throughout this proof the data are those of the setting (FF) of Section~G5.6.0, with $R\ge D_{\mathrm{be}}+1$.
Let $\mathbf{m}:=\bar m\eta\in(0,1]$ and $n=R/\eta$;
the matrices $D^{(p)}$, $\mathrm{Rsy}(p)$ and $\mathrm{Er}(p,\mathbf{m})$ are formed with this $\mathbf{m}$.
For real $p$ and $\mathbf{m}'>0$ let $\mathrm{Sdr}(p,\mathbf{m}'):=\left(\mathbf{m}'+i\gamma\cdot p\right)^{-1}$.
By (I.c1), $\left(\mathbf{m}'-i\gamma\cdot p\right)\left(\mathbf{m}'+i\gamma\cdot p\right)=\mathbf{m}'^2+|p|_2^2$, so that $\left(\mathbf{m}'^2+|p|_2^2\right)^{\frac{1}{2}}\mathrm{Sdr}(p,\mathbf{m}')$ is unitary.
Let $Q_{\mathrm{be}}$ be a closed coordinate cube of side $D_{\mathrm{be}}$ that contains $\mathrm{supp}\,\beta$.
Since $\beta$ is continuous and $\beta\ne0$, one has $D_{\mathrm{be}}>0$;
hence $R>1$ and $R>D_{\mathrm{be}}$, and the sets $Q_{\mathrm{be}}+Rz$, $z\in\mathbb{Z}^3$, are pairwise disjoint.
Put $\beta^R:=\sum_{z\in\mathbb{Z}^3}\beta(\,\cdot+Rz)$;
then $a^R_\nu=-\sum_{\mu\ne\nu}\partial_\mu\beta^R_{\mu\nu}$.
The functions $\mathrm{va}:=i\gamma\cdot a$ and $\mathrm{va}^R:=i\gamma\cdot a^R$ act as multiplication operators.

\emph{Absolute numbers.}
Exactly $(2j+1)^3-(2j-1)^3=24j^2+2$ points $z\in\mathbb{Z}^3$ satisfy $|z|_\infty=j$;
hence, by $\sum_{j\ge1}j^{-2}=\pi^2/6$ and $\sum_{j\ge1}j^{-4}=\pi^4/90$,
\[
\mathrm{zs}_4:=\sum_{z\in\mathbb{Z}^3,\,z\ne0}|z|_\infty^{-4}=\sum_{j\ge1}\left(24j^2+2\right)j^{-4}=4\pi^2+\frac{\pi^4}{45}<42,
\]
and $\mathrm{zs}_4/\pi^4<1/2$.
Further
\[
\mathrm{cs}_1:=\frac{\left(1+\sqrt{3}\,\pi\right)^4}{6\pi^2}<30,\qquad\mathrm{Lz}:=\frac{3\sqrt{3}\,\pi}{4}<5;
\]
the first bound holds because $1+\sqrt{3}\,\pi<6.45$, $6.45^4<1731$ and $6\pi^2>59$.

\emph{Numbers that depend on $\beta$ only.}
\begin{gather*}
\mathrm{Ab}_1:=\max_\nu\sum_{\mu\ne\nu}|\partial_\mu\beta_{\mu\nu}|_\infty,\qquad
\mathrm{Ab}_2:=\max_\nu\sum_{\mu\ne\nu}\max_j|\partial_j\partial_\mu\beta_{\mu\nu}|_\infty,\\
\mathrm{Ab}_5:=\max_\nu\sum_{\mu\ne\nu}\max_j\left\|\partial_j^4\partial_\mu\beta_{\mu\nu}\right\|_{L^1(\mathbb{R}^3)},\qquad
\mathrm{Vbe}:=\left(D_{\mathrm{be}}+3\right)^3,\qquad
\mathrm{kpr}:=\mathrm{Ab}_1\cdot\mathrm{Vbe}^{\frac{1}{4}},\\
\mathrm{bq}(\mathrm{kph}):=\sum_\nu\sum_{\mu\ne\nu}|\mathrm{kph}_\mu|\cdot|\hat\beta_{\mu\nu}(\mathrm{kph})|\quad(\mathrm{kph}\in\mathbb{R}^3),\qquad
\mathrm{bs}_4:=\sup_{\mathrm{kph}\in\mathbb{R}^3}\left(1+|\mathrm{kph}|_2\right)^4\cdot\mathrm{bq}(\mathrm{kph}),\\
\mathrm{Afr}_1:=\mathrm{cs}_1\cdot\mathrm{bs}_4,\qquad
\mathrm{Afr}_0:=2\,\mathrm{bs}_4.
\end{gather*}
Here $\hat\beta_{\mu\nu}$ is the Fourier transform of Section~I.2.6;
it is a Schwartz function, and therefore $\mathrm{bs}_4<\infty$.

\emph{Numbers that depend on $(\rho,\bar m,\beta)$; $\rho$ enters through $C_E$ only.}
\begin{gather*}
\mathrm{Cf}:=\sqrt{2}+7C_E,\qquad
\mathrm{kn}_0:=3\left(\frac{4}{\pi\bar m}\right)^{\frac{1}{4}}\cdot\mathrm{kpr},\qquad
\mathrm{kn}_1:=36\left(1536\left(C_E^4+\frac{1}{\pi\bar m}\right)\right)^{\frac{1}{4}}\cdot\mathrm{kpr},\\
\mathrm{kn}_2:=\left(\frac{9\,\mathrm{Ab}_1^4D_{\mathrm{be}}^3}{2\pi\bar m}\right)^{\frac{1}{4}},\qquad
\mathrm{rn}_0:=\frac{3}{2}\left[2\cdot64^{\frac{1}{4}}\cdot\mathrm{kpr}+\left(\frac{128}{\pi\bar m}\right)^{\frac{1}{4}}\cdot\mathrm{Ab}_2\cdot\mathrm{Vbe}^{\frac{1}{4}}\right],\\
\mathrm{lnk}_1:=4\max\left(\mathrm{kn}_1,12\,\mathrm{kn}_0\right)^3\cdot36\cdot192^{\frac{1}{4}}\cdot\left(C_E+1\right)\cdot\mathrm{kpr},\qquad
\mathrm{lnk}_2:=4\left(2\,\mathrm{kn}_0\right)^3\cdot\mathrm{rn}_0,\\
\mathrm{lnk}_3:=4\,\mathrm{kn}_0^3\cdot\left(\frac{24\sqrt{3}}{\pi^9}\right)^{\frac{1}{2}}\cdot\mathrm{zs}_4\cdot\mathrm{Ab}_5,\qquad
\mathrm{lnk}_4:=36\,\mathrm{kn}_0^3\cdot\left(\frac{4}{\pi\bar m}\right)^{\frac{1}{4}}\cdot\mathrm{Ab}_2\cdot\mathrm{Vbe}^{\frac{1}{4}},\\
\mathrm{lnk}_5:=12\,\mathrm{kn}_0^3\cdot\left(\frac{128}{3\pi^3}\right)^{\frac{1}{4}}\cdot\mathrm{kpr},\qquad
\mathrm{Sgm}:=\sum_{j\ge1}\left(24j^2+2\right)\exp\left(-\bar m(j-1)\right),\\
\mathrm{wim}_0:=\frac{3\sqrt{3}}{2\pi}\cdot\mathrm{Sgm}\cdot\mathrm{Ab}_1\cdot D_{\mathrm{be}}^3,\qquad
\mathrm{lnk}_6:=4\left(\mathrm{kn}_2+\mathrm{wim}_0\right)^3\cdot\mathrm{wim}_0.
\end{gather*}

\emph{The three constants of the statement.}
\begin{equation*}
\begin{gathered}
C_{\mathrm{ct}}:=3456\left(6\,\mathrm{Cf}\right)^3\cdot\mathrm{Afr}_1^3\cdot\mathrm{Afr}_0,\\
C_{1}^{\mathrm{N2}}:=C_{\mathrm{ct}}\left(1+\frac{4}{3}\left(\frac{5}{\bar m}\right)^{\frac{3}{4}}\right)+6\left(\mathrm{lnk}_1+\mathrm{lnk}_2+\mathrm{lnk}_3+\mathrm{lnk}_4+\mathrm{lnk}_5\right),\qquad
C_{2}^{\mathrm{N2}}:=6\,\mathrm{lnk}_6.
\end{gathered}
\tag*{[2.1]}
\end{equation*}

\paragraph{2.1. Spaces, momenta, three kinds of operators, the propagator, the vertices.}
\emph{Spaces.}
Let $\mathrm{Sfc}:=\mathfrak{S}^{\mathrm{ch}}(1;n)$ (Proposition~46), with the inner product
\[
(f',f)^{\mathrm{ch}}=\frac{1}{48}\sum_{c\in\{0,\dots,n-1\}^3}\ \sum_{g\in O_{\mathrm{cub}}}\overline{f'(c,g)}\cdot f(c,g)
\]
(Section~I.1.6 with $\eta_0=1$; the summand is $n\mathbb{Z}^3$-periodic in $c$);
its dimension is $192\,n^3$.
Let $\mathrm{Sfl}:=\mathfrak{S}(1;n)$ (Proposition~46, with the sides $(n,n,n)$), that is,
\[
\mathrm{Sfl}=\left\{\xi:\mathbb{Z}^3\to\mathbb{C}^4:\ \xi(c+ne_\mu)=-\xi(c)\right\},\qquad
(\xi',\xi):=\sum_{c\in\{0,\dots,n-1\}^3}\overline{\xi'(c)}\cdot\xi(c);
\]
its dimension is $4n^3$.
In what follows $(\mathrm{Fb},\mathrm{Sfb})$ stands for $({V_{\mathrm{v}}},\mathrm{Sfc})$ or for $(\mathbb{C}^4,\mathrm{Sfl})$;
an element of $\mathrm{Sfb}$ is a function $\psi:\mathbb{Z}^3\to\mathrm{Fb}$ with $\psi(c+ne_\mu)=-\psi(c)$.

\emph{The cluster maps.}
For $\xi\in\mathrm{Sfl}$ and $f\in\mathrm{Sfc}$ let
\[
(J\xi)(c,g):=\xi(c),\qquad(J^{*}f)(c):=\frac{1}{48}\sum_{g\in O_{\mathrm{cub}}}f(c,g);
\]
these are the two displays (I.3.1), read on $\mathrm{Sfl}$ and on $\mathrm{Sfc}$ as in the preamble of Proposition~46.
The function $J\xi$ is antiperiodic in $c$ for every label, and $J^{*}f$ is antiperiodic;
thus $J:\mathrm{Sfl}\to\mathrm{Sfc}$ and $J^{*}:\mathrm{Sfc}\to\mathrm{Sfl}$.
For $\xi\in\mathrm{Sfl}$ and $f\in\mathrm{Sfc}$, with sums over $c\in\{0,\dots,n-1\}^3$ and $g\in O_{\mathrm{cub}}$,
\begin{gather*}
(J^{*}J\xi)(c)=\frac{1}{48}\sum_g\xi(c)=\xi(c),\qquad
(J\xi,J\xi)^{\mathrm{ch}}=(J^{*}J\xi,\xi)=(\xi,\xi),\\
(f,J\xi)^{\mathrm{ch}}=\frac{1}{48}\sum_c\sum_g\overline{f(c,g)}\cdot\xi(c)=\sum_c\overline{(J^{*}f)(c)}\cdot\xi(c)=(J^{*}f,\xi).
\end{gather*}
Hence $J^{*}J=1$, $J^{*}$ is the adjoint of $J$, and $J$ is an isometry of $\mathrm{Sfl}$ into $\mathrm{Sfc}$.

\emph{Momenta.}
Let
\[
\mathrm{Pb}_n:=\frac{2\pi}{n}\left(\mathbb{Z}+\frac{1}{2}\right)^3\cap(-\pi,\pi]^3.
\]
This set has $n^3$ elements, and reduction modulo $2\pi\mathbb{Z}^3$ maps the set $P_{n,0}$ of Section~I.2.6 bijectively onto it.
For $p\in\mathrm{Pb}_n$ let $\mathrm{ew}_p(c):=n^{-\frac{3}{2}}\exp(ip\cdot c)$, $c\in\mathbb{Z}^3$;
since $\exp(inp_\mu)=-1$, one has $\mathrm{ew}_p(c+ne_\mu)=-\mathrm{ew}_p(c)$.
For $p,p'\in\mathrm{Pb}_n$ the vector $p'-p$ belongs to $\frac{2\pi}{n}\mathbb{Z}^3\cap(-2\pi,2\pi)^3$, and
\[
n^{-3}\sum_{c\in\{0,\dots,n-1\}^3}\exp\left(i(p'-p)\cdot c\right)=\prod_\mu\left(n^{-1}\sum_{j=0}^{n-1}\exp\left(i(p'-p)_\mu j\right)\right)=\mathbf{1}[p=p'].
\]
Consequently, for $u,u'\in\mathrm{Fb}$,
\[
\left(\mathrm{ew}_p\otimes u',\mathrm{ew}_{p'}\otimes u\right)=\mathbf{1}[p=p']\,(u',u),\qquad\text{where }(\mathrm{ew}_p\otimes u)(c):=\mathrm{ew}_p(c)\,u.
\]
Counting dimensions, one sees that $\mathrm{Sfb}$ is the orthogonal direct sum of the $n^3$ subspaces $\mathrm{ew}_p\otimes\mathrm{Fb}$, $p\in\mathrm{Pb}_n$.

\emph{Blocks.}
Every linear map $O$ of $\mathrm{Sfb}$ into itself has blocks $O(p,p')\in\mathrm{End}(\mathrm{Fb})$, determined by
\[
O\left(\mathrm{ew}_{p'}\otimes u\right)=\sum_{p\in\mathrm{Pb}_n}\mathrm{ew}_p\otimes O(p,p')u.
\]
One has $(O_1O_2)(p,p'')=\sum_{p'}O_1(p,p')O_2(p',p'')$ and $\mathrm{Tr}\,O=\sum_p\mathrm{tr}\,O(p,p)$.

\emph{Multipliers.}
For $\mathrm{hm}:\mathrm{Pb}_n\to\mathrm{End}(\mathrm{Fb})$ let $\mathrm{Mu}[\mathrm{hm}]\left(\mathrm{ew}_p\otimes u\right):=\mathrm{ew}_p\otimes\mathrm{hm}(p)u$.
Then $\mathrm{Mu}[\mathrm{hm}]^{*}=\mathrm{Mu}[\mathrm{hm}^{*}]$ and $\mathrm{Mu}[\mathrm{hm}]\cdot\mathrm{Mu}[\mathrm{hm}']=\mathrm{Mu}[\mathrm{hm}\cdot\mathrm{hm}']$.
A matrix $O\in\mathrm{End}(\mathrm{Fb})$ acts on $\mathrm{Sfb}$ as a constant multiplier: $(O\psi)(c)=O\left(\psi(c)\right)$.

\emph{Multiplication operators.}
An $n\mathbb{Z}^3$-periodic function $\mathrm{fs}:\mathbb{Z}^3\to\mathbb{C}$ acts by $(\mathrm{fs}\cdot\psi)(c):=\mathrm{fs}(c)\psi(c)$.
Its blocks are the multiples
\[
\mathrm{fs}(p,p')=n^{-3}\sum_{c\in\{0,\dots,n-1\}^3}\mathrm{fs}(c)\exp\left(-i(p-p')\cdot c\right)
\]
of the identity;
in particular $\left(\mathrm{ew}_p,\mathrm{fs}\cdot\mathrm{ew}_p\right)=n^{-3}\sum_c\mathrm{fs}(c)$ for every $p$.
Here, for scalar functions $\mathrm{fn},\mathrm{fn}'$ on $\mathbb{Z}^3$ whose product $\overline{\mathrm{fn}'}\cdot\mathrm{fn}$ is $n\mathbb{Z}^3$-periodic, $(\mathrm{fn}',\mathrm{fn}):=\sum_{c\in\{0,\dots,n-1\}^3}\overline{\mathrm{fn}'(c)}\cdot\mathrm{fn}(c)$.
Multiplication operators commute with constant multipliers.
Finally,
\[
\|\mathrm{fs}\|_{\ell^4}:=\left(\sum_{c\in\{0,\dots,n-1\}^3}|\mathrm{fs}(c)|^4\right)^{\frac{1}{4}}.
\]

\emph{Shifts.}
Let $(\mathrm{Sh}^{\pm}_\nu\psi)(c):=\psi(c\pm e_\nu)$.
Thus $\mathrm{Sh}^{\pm}_\nu=\mathrm{Mu}\left[\exp(\pm ip_\nu)\right]$, and these operators are unitary.

\emph{The cluster maps and the three kinds of operators.}
The maps $J$ and $J^{*}$ commute with multiplication operators and with shifts:
for an $n\mathbb{Z}^3$-periodic scalar function $\mathrm{fs}$, $\xi\in\mathrm{Sfl}$ and $f\in\mathrm{Sfc}$,
\begin{gather*}
\left(J(\mathrm{fs}\cdot\xi)\right)(c,g)=\mathrm{fs}(c)\xi(c)=(\mathrm{fs}\cdot J\xi)(c,g),\qquad
\left(J\cdot\mathrm{Sh}^{\pm}_\nu\xi\right)(c,g)=\xi(c\pm e_\nu)=\left(\mathrm{Sh}^{\pm}_\nu J\xi\right)(c,g),\\
\left(J^{*}(\mathrm{fs}\cdot f)\right)(c)=\frac{1}{48}\sum_g\mathrm{fs}(c)f(c,g)=\mathrm{fs}(c)(J^{*}f)(c),\\
\left(J^{*}\mathrm{Sh}^{\pm}_\nu f\right)(c)=\frac{1}{48}\sum_gf(c\pm e_\nu,g)=\left(\mathrm{Sh}^{\pm}_\nu J^{*}f\right)(c).
\end{gather*}
Let $J:\mathbb{C}^4\to{V_{\mathrm{v}}}$ and $J^{*}:{V_{\mathrm{v}}}\to\mathbb{C}^4$ denote also the fibre maps of Section~G5.6.0;
they occur on the right-hand sides below.
Then $J\left(\mathrm{ew}_p\otimes\xi\right)=\mathrm{ew}_p\otimes J\xi$ for $\xi\in\mathbb{C}^4$, both sides taking the value $\mathrm{ew}_p(c)\xi$ at $(c,g)$;
and $J^{*}\left(\mathrm{ew}_p\otimes u\right)=\mathrm{ew}_p\otimes J^{*}u$ for $u\in{V_{\mathrm{v}}}$, both sides taking the value $\mathrm{ew}_p(c)\frac{1}{48}\sum_gu(g)$ at $c$.
Hence, for $\mathrm{hm}:\mathrm{Pb}_n\to\mathrm{End}(\mathbb{C}^4)$ and for a matrix $O\in\mathrm{End}({V_{\mathrm{v}}})$ acting as a constant multiplier,
\[
J\cdot\mathrm{Mu}[\mathrm{hm}]\cdot J^{*}\left(\mathrm{ew}_p\otimes u\right)=J\left(\mathrm{ew}_p\otimes\mathrm{hm}(p)J^{*}u\right)=\mathrm{ew}_p\otimes J\cdot\mathrm{hm}(p)\cdot J^{*}u,
\]
that is, $J\cdot\mathrm{Mu}[\mathrm{hm}]\cdot J^{*}=\mathrm{Mu}\left[J\cdot\mathrm{hm}\cdot J^{*}\right]$;
and
\[
J^{*}OJ\left(\mathrm{ew}_p\otimes\xi\right)=\mathrm{ew}_p\otimes J^{*}OJ\xi,
\]
that is, $J^{*}OJ$ is the constant multiplier with the $4\times4$ matrix $J^{*}OJ$.

\emph{The operator.}
By Proposition~46\,(b), $D^{\mathrm{ch}}[B]$ is the operator (I.2.3) at the mesh $\eta_0=1$ of Section~I.2.5, with the phase $\exp(iB)$ on the bonds of type $3$:
\begin{equation*}
\begin{aligned}
\left(D^{\mathrm{ch}}[B]f\right)(c,g)={}&-\sum_{i\le2}M_i(g)\left[f(c,gs_i)-f(c,g)\right]\\
&-M_3(g)\left[\exp\left(iB(c,c+ge_1)\right)f(c+ge_1,gs_3)-f(c,g)\right].
\end{aligned}
\tag*{[2.2]}
\end{equation*}
By (I.2.6), $D^{\mathrm{ch}}[0]+\mathbf{m}=\mathrm{Mu}\left[D^{(p)}\right]$, where $D^{(p)}$ is $2\pi\mathbb{Z}^3$-periodic in $p$.
By Proposition~37.1\,(f), every $D^{(p)}$ with $p\in\mathrm{Pb}_n$ is invertible.
Therefore $D^{\mathrm{ch}}[0]+\mathbf{m}$ is invertible, $\det\left(D^{\mathrm{ch}}[0]+\mathbf{m}\right)\ne0$, and, by Proposition~37.1\,(f) and $\mathbf{m}\le1$,
\begin{equation*}
\begin{gathered}
\mathrm{G0}=\mathrm{Mu}[\mathrm{Rsy}],\\
|\mathrm{Rsy}(p)|_2\le C_E+\left(\mathbf{m}^2+|p|_2^2\right)^{-\frac{1}{2}},\qquad
|\mathrm{Rsy}(p)|_2\le\frac{\mathrm{Cf}}{|p|_2+\mathbf{m}}\qquad(p\in\mathrm{Pb}_n).
\end{gathered}
\tag*{[2.3]}
\end{equation*}

\emph{The vertices.}
For $c\in\mathbb{Z}^3$ let $\mathrm{alat}_\nu(c):=a_\eta[\beta]_\nu(\eta c)$, a real $n\mathbb{Z}^3$-periodic function, and $\mathrm{alat}^{-}_\nu(c):=\mathrm{alat}_\nu(c-e_\nu)$;
further let
\[
\mathrm{bph}^{+}_\nu(c):=B_1(c,c+e_\nu)=\eta\cdot\mathrm{alat}_\nu(c),\qquad
\mathrm{bph}^{-}_\nu(c):=B_1(c,c-e_\nu)=-B_1(c-e_\nu,c)=-\eta\cdot\mathrm{alat}^{-}_\nu(c).
\]
One has $B_t=tB_1$.
If $ge_1=\pm e_\nu$, then $M_3(g)=3(\rho\mp\gamma_\nu)$, $B_t(c,c+ge_1)=t\cdot\mathrm{bph}^{\pm}_\nu(c)$ and $f(c+ge_1,gs_3)=\left(\mathrm{Sh}^{\pm}_\nu f\right)(c,gs_3)$.
With the fibre matrices $\mathrm{Ehp}^{\pm}_\nu$ of Section~G5.6.0, [2.2] therefore yields, for every $t\in\mathbb{C}$,
\begin{equation*}
\begin{gathered}
\mathrm{Dw}(t):=D^{\mathrm{ch}}[B_t]-D^{\mathrm{ch}}[0]=-\sum_\nu\sum_{\pm}\left(\exp\left(it\cdot\mathrm{bph}^{\pm}_\nu\right)-1\right)\mathrm{Ehp}^{\pm}_\nu\cdot\mathrm{Sh}^{\pm}_\nu,\\
\mathrm{Vd}_j=-\sum_\nu\sum_{\pm}\left(i\,\mathrm{bph}^{\pm}_\nu\right)^j\mathrm{Ehp}^{\pm}_\nu\cdot\mathrm{Sh}^{\pm}_\nu\quad(j\ge1),\qquad
\mathrm{Dw}(t)=\sum_{j\ge1}\frac{t^j}{j!}\,\mathrm{Vd}_j;
\end{gathered}
\tag*{[2.4]}
\end{equation*}
in each product the shift acts first, then the fibre matrix, then the multiplication operator, and the series converges for every $t$.
The second formula of [2.4] coincides with the formula for $\mathrm{Vd}_j$ in the definitions of Proposition~37.2.
By Proposition~37.1\,(d), $\left|\mathrm{Ehp}^{\pm}_\nu\cdot\mathrm{Sh}^{\pm}_\nu\right|_2\le3(1+\rho)\le6$.
Finally, let $\mathrm{Yv}_j:=\mathrm{G0}\cdot\mathrm{Vd}_j$.

\paragraph{2.2. Clause (b)}

\begin{proof}[Proof of Proposition~37.2\,(b)]
Recall that $\mathrm{Dw}(t)=D^{\mathrm{ch}}[B_t]-D^{\mathrm{ch}}[0]$ and $\mathbf{m}=\bar{m}\eta$.
Hence
\[
N_0(t)=\frac{\det\left(D^{\mathrm{ch}}[0]+\mathbf{m}+\mathrm{Dw}(t)\right)}{\det\left(D^{\mathrm{ch}}[0]+\mathbf{m}\right)}=\det\left(1+\mathrm{G0}\cdot\mathrm{Dw}(t)\right),
\]
an entire function of $t$ with $N_0(0)=1$.

\emph{A formula.}
Let $\mathrm{Yq}$ be a linear map of $\mathrm{Sfc}$ into itself with $|\mathrm{Yq}|_2<1$.
Then
\[
\det(1+\mathrm{Yq})=\exp\left(\sum_{\mathrm{lf}\ge1}\frac{(-1)^{\mathrm{lf}+1}}{\mathrm{lf}}\operatorname{Tr}\mathrm{Yq}^{\mathrm{lf}}\right),
\qquad
\left|\operatorname{Tr}\mathrm{Yq}^{\mathrm{lf}}\right|\le192\,n^3\cdot|\mathrm{Yq}|_2^{\mathrm{lf}}\quad(\mathrm{lf}\ge1).
\]
To see this, denote by $\mathrm{evq}_i$, $i=1,\dots,192\,n^3$, the eigenvalues of $\mathrm{Yq}$, each listed as often as its algebraic multiplicity indicates.
Choose a basis of $\mathrm{Sfc}$ that makes the matrix of $\mathrm{Yq}$ upper triangular;
its diagonal then consists of the numbers $\mathrm{evq}_i$, and the matrices of $1+\mathrm{Yq}$ and of $\mathrm{Yq}^{\mathrm{lf}}$, again upper triangular, carry on their diagonals the numbers $1+\mathrm{evq}_i$ and $\mathrm{evq}_i^{\mathrm{lf}}$.
Consequently
\[
\det(1+\mathrm{Yq})=\prod_i\left(1+\mathrm{evq}_i\right),\qquad\operatorname{Tr}\mathrm{Yq}^{\mathrm{lf}}=\sum_i\mathrm{evq}_i^{\mathrm{lf}}.
\]
If $f\ne0$ and $\mathrm{Yq}f=\mathrm{evq}_if$, then
\[
|\mathrm{evq}_i|\cdot\left((f,f)^{\mathrm{ch}}\right)^{\frac12}=\left((\mathrm{Yq}f,\mathrm{Yq}f)^{\mathrm{ch}}\right)^{\frac12}\le|\mathrm{Yq}|_2\cdot\left((f,f)^{\mathrm{ch}}\right)^{\frac12},
\]
so that $|\mathrm{evq}_i|\le|\mathrm{Yq}|_2<1$;
the bound on $\operatorname{Tr}\mathrm{Yq}^{\mathrm{lf}}$ follows.
Every complex number $\mathrm{evq}$ of modulus less than $1$ satisfies
\[
1+\mathrm{evq}=\exp\left(\sum_{\mathrm{lf}\ge1}\frac{(-1)^{\mathrm{lf}+1}}{\mathrm{lf}}\,\mathrm{evq}^{\mathrm{lf}}\right).
\]
Indeed, on the disk $|\mathrm{evq}|<1$ the power series in the exponent has the derivative $\sum_{\mathrm{lf}\ge1}(-\mathrm{evq})^{\mathrm{lf}-1}=(1+\mathrm{evq})^{-1}$;
therefore the function
\[
(1+\mathrm{evq})\exp\left(-\sum_{\mathrm{lf}\ge1}\frac{(-1)^{\mathrm{lf}+1}}{\mathrm{lf}}\,\mathrm{evq}^{\mathrm{lf}}\right)
\]
has the derivative $0$, and its value at $\mathrm{evq}=0$ is $1$.
It follows that
\[
\prod_i\left(1+\mathrm{evq}_i\right)=\exp\left(\sum_i\sum_{\mathrm{lf}\ge1}\frac{(-1)^{\mathrm{lf}+1}}{\mathrm{lf}}\,\mathrm{evq}_i^{\mathrm{lf}}\right)=\exp\left(\sum_{\mathrm{lf}\ge1}\frac{(-1)^{\mathrm{lf}+1}}{\mathrm{lf}}\operatorname{Tr}\mathrm{Yq}^{\mathrm{lf}}\right),
\]
where the sum over $i$ is a finite sum of absolutely convergent series.

\emph{The logarithm near $t=0$.}
By [2.4],
\[
\mathrm{Yq}(t):=\mathrm{G0}\cdot\mathrm{Dw}(t)=\sum_{j\ge1}\frac{t^j}{j!}\,\mathrm{Yv}_j,
\]
the series converging in norm for every $t$.
Let
\[
\mathrm{cf}:=\sum_{j\ge1}\frac{|\mathrm{Yv}_j|_2}{j!}.
\]
This number is finite: by [2.4], $\mathrm{Vd}_j$ is a sum of six operators, each of norm at most $6\max_{\nu,\pm}|\mathrm{bph}^{\pm}_{\nu}|_\infty^j$, and $|\mathrm{Yv}_j|_2\le|\mathrm{G0}|_2\cdot|\mathrm{Vd}_j|_2$.
Put $t_r:=\min\left(1,\frac{1}{2\,\mathrm{cf}}\right)$, and $t_r:=1$ if $\mathrm{cf}=0$.
For $|t|\le t_r$ (so that $|t|\le1$)
\[
|\mathrm{Yq}(t)|_2\le\sum_{j\ge1}|t|^j\,\frac{|\mathrm{Yv}_j|_2}{j!}\le\mathrm{cf}\cdot|t|\le\frac12 ,
\]
and for these $t$ one defines
\begin{align*}
\mathrm{lgq}(t)&:=\sum_{\mathrm{lf}\ge1}\frac{(-1)^{\mathrm{lf}+1}}{\mathrm{lf}}\operatorname{Tr}\mathrm{Yq}(t)^{\mathrm{lf}},\\
\mathrm{rmq}(t)&:=\sum_{\mathrm{lf}\ge5}\frac{(-1)^{\mathrm{lf}+1}}{\mathrm{lf}}\operatorname{Tr}\mathrm{Yq}(t)^{\mathrm{lf}}=\mathrm{lgq}(t)-\sum_{\mathrm{lf}=1}^{4}\frac{(-1)^{\mathrm{lf}+1}}{\mathrm{lf}}\operatorname{Tr}\mathrm{Yq}(t)^{\mathrm{lf}}.
\end{align*}
As $\mathrm{cf}\cdot|t|\le\frac12$, one has $\sum_{\mathrm{lf}\ge i}(\mathrm{cf}\cdot|t|)^{\mathrm{lf}}\le2\,(\mathrm{cf}\cdot|t|)^i$.
Together with the formula this gives, for $|t|\le t_r$,
\[
\exp\left(\mathrm{lgq}(t)\right)=N_0(t),\qquad|\mathrm{lgq}(t)|\le384\,n^3\cdot\mathrm{cf}\cdot|t|,\qquad|\mathrm{rmq}(t)|\le384\,n^3\cdot\mathrm{cf}^5\cdot|t|^5 .
\]
In particular $N_0$ vanishes nowhere on the disk $|t|\le t_r$.
Therefore
\[
\log N_0(t):=\int_0^1\mathrm{d}\,\mathrm{sv}\;\frac{t\,N_0'(\mathrm{sv}\cdot t)}{N_0(\mathrm{sv}\cdot t)}
\]
is a holomorphic function on $|t|<t_r$ with $\log N_0(0)=0$ and with the derivative $N_0'/N_0$;
for the derivative one uses the identity $\frac{\mathrm{d}}{\mathrm{d}t}\left(t\cdot\mathrm{fs}(\mathrm{sv}\cdot t)\right)=\frac{\mathrm{d}}{\mathrm{d}\,\mathrm{sv}}\left(\mathrm{sv}\cdot\mathrm{fs}(\mathrm{sv}\cdot t)\right)$, valid for every holomorphic function $\mathrm{fs}$ on the disk.
Since $N_0\exp(-\log N_0)$ has a vanishing derivative and equals $1$ at $t=0$, one has $\exp\left(\log N_0(t)\right)=N_0(t)$.
This is the logarithm of $N_0$ that enters the definition of $b_{\mathrm{ch}}(\eta,R)$ in Section~G5.6.0;
any other continuous logarithm of $N_0$ near $t=0$ differs from it by a constant, which does not change the derivatives.
Since $\exp\left(\mathrm{lgq}(t)-\log N_0(t)\right)=1$, the number $\mathrm{lgq}(t)-\log N_0(t)$ lies in $2\pi i\mathbb{Z}$.
As $t\to0$, both $\mathrm{lgq}(t)$ and $\log N_0(t)$ converge to $0$.
Consequently there is $t_r'\in(0,t_r]$ such that $|\mathrm{lgq}(t)-\log N_0(t)|<2\pi$ whenever $|t|<t_r'$;
for such $t$ the difference, an element of $2\pi i\mathbb{Z}$, is $0$, that is, $\mathrm{lgq}(t)=\log N_0(t)$.
For $|t|<t_r'$ therefore
\[
\log N_0(t)=\sum_{\mathrm{lf}=1}^{4}\frac{(-1)^{\mathrm{lf}+1}}{\mathrm{lf}}\operatorname{Tr}\left(\sum_{j\ge1}\frac{t^j}{j!}\,\mathrm{Yv}_j\right)^{\mathrm{lf}}+\mathrm{rmq}(t).
\]
The terms with $\mathrm{lf}\le4$ form an entire function of $t$.
Hence $\mathrm{rmq}$ is holomorphic on $|t|<t_r'$, where moreover $|\mathrm{rmq}(t)|\le384\,n^3\cdot\mathrm{cf}^5\cdot|t|^5$.
It follows that the coefficients of $t^0,\dots,t^4$ in the Taylor expansion of $\mathrm{rmq}$ at $t=0$ are zero:
if the first non-vanishing one had an order $i\le4$, then $|\mathrm{rmq}(t)|/|t|^i$ would tend, as $t\to0$, to the modulus of this coefficient, which is not $0$.
Thus $b_{\mathrm{ch}}(\eta,R)=\left(\frac{\mathrm{d}}{\mathrm{d}t}\right)^4\log N_0(t)\big|_{t=0}$ equals $4!$ times the coefficient of $t^4$ in the sum over $\mathrm{lf}\le4$.
Using the cyclicity of the trace, one finds the following contributions of the four values of $\mathrm{lf}$ to this coefficient:
\begin{align*}
\mathrm{lf}=1:&\quad\frac{1}{24}\operatorname{Tr}\mathrm{Yv}_4;\\
\mathrm{lf}=2:&\quad-\frac12\left[\frac{2\operatorname{Tr}(\mathrm{Yv}_1\cdot\mathrm{Yv}_3)}{1!\,3!}+\frac{\operatorname{Tr}(\mathrm{Yv}_2^2)}{2!\,2!}\right]=-\frac16\operatorname{Tr}(\mathrm{Yv}_1\cdot\mathrm{Yv}_3)-\frac18\operatorname{Tr}(\mathrm{Yv}_2^2);\\
\mathrm{lf}=3:&\quad\frac13\left[\frac{3\operatorname{Tr}(\mathrm{Yv}_1^2\cdot\mathrm{Yv}_2)}{1!\,1!\,2!}\right]=\frac12\operatorname{Tr}(\mathrm{Yv}_1^2\cdot\mathrm{Yv}_2)\\
&\quad\text{(the three orders $(1,1,2)$, $(1,2,1)$, $(2,1,1)$ give equal traces)};\\
\mathrm{lf}=4:&\quad-\frac14\operatorname{Tr}(\mathrm{Yv}_1^4).
\end{align*}
Multiplying by $24$ one finds
\[
b_{\mathrm{ch}}(\eta,R)=\operatorname{Tr}\mathrm{Yv}_4-4\operatorname{Tr}(\mathrm{Yv}_1\cdot\mathrm{Yv}_3)-3\operatorname{Tr}(\mathrm{Yv}_2^2)+12\operatorname{Tr}(\mathrm{Yv}_1^2\cdot\mathrm{Yv}_2)-6\operatorname{Tr}(\mathrm{Yv}_1^4),
\]
and, as $\mathrm{Yv}_j=\mathrm{G0}\cdot\mathrm{Vd}_j$, this is the identity (b).
\end{proof}

The identity (b) holds for every $R$ admitted in Setting (FF);
the condition $R\ge D_{\mathrm{be}}+1$ is not required for it.

\paragraph{2.3. Trace ideals}

Let $\mathrm{Aq}$ be a compact operator on a separable Hilbert space, and let $|\mathrm{Aq}|:=(\mathrm{Aq}^*\cdot\mathrm{Aq})^{\frac12}$.
Define
\[
\|\mathrm{Aq}\|_1:=\operatorname{Tr}|\mathrm{Aq}|,\qquad
\|\mathrm{Aq}\|_2:=\left(\operatorname{Tr}\mathrm{Aq}^*\cdot\mathrm{Aq}\right)^{\frac12},\qquad
\|\mathrm{Aq}\|_4:=\left(\operatorname{Tr}(\mathrm{Aq}^*\cdot\mathrm{Aq})^2\right)^{\frac14},
\]
with values in $[0,\infty]$.
The trace class $I_{1}$, the Hilbert--Schmidt class $I_{2}$ and the class $I_{4}$ consist of those $\mathrm{Aq}$ for which the respective number is finite.

\begin{itemize}
\item[(Tr1)]
The following facts are found in \cite[Chapters~1--3]{Simon2005} and \cite[Section~VI.6]{ReedSimon1980}.
With these norms $I_{1}$, $I_{2}$, $I_{4}$ are Banach spaces, and
\[
|\mathrm{Aq}|_2\le\|\mathrm{Aq}\|_4\le\|\mathrm{Aq}\|_2\le\|\mathrm{Aq}\|_1 ,
\]
where $|\cdot|_2$ is the operator norm.
In each of the three norms $\|\mathrm{Aq}^*\|=\|\mathrm{Aq}\|$, and $\|\mathrm{Bq}\cdot\mathrm{Aq}\cdot\mathrm{Bq}'\|\le|\mathrm{Bq}|_2\cdot\|\mathrm{Aq}\|\cdot|\mathrm{Bq}'|_2$ for bounded operators $\mathrm{Bq}$, $\mathrm{Bq}'$.
One has $\|\mathrm{Aq}\|_4^4=\|\mathrm{Aq}^*\cdot\mathrm{Aq}\|_2^2$, and a bounded operator $\mathrm{Aq}$ lies in $I_{4}$ if $\mathrm{Aq}^*\cdot\mathrm{Aq}$ lies in $I_{2}$.
H\"older's inequality asserts
\[
\|\mathrm{Aq}_1\cdot\mathrm{Aq}_2\|_2\le\|\mathrm{Aq}_1\|_4\cdot\|\mathrm{Aq}_2\|_4,\qquad
\|\mathrm{Aq}_1\cdot\mathrm{Aq}_2\|_1\le\|\mathrm{Aq}_1\|_2\cdot\|\mathrm{Aq}_2\|_2,
\]
and hence $\|\mathrm{Aq}_1\cdot\mathrm{Aq}_2\cdot\mathrm{Aq}_3\cdot\mathrm{Aq}_4\|_1\le\prod_i\|\mathrm{Aq}_i\|_4$.
Further, $|\operatorname{Tr}\mathrm{Aq}|\le\|\mathrm{Aq}\|_1$;
the trace of an operator of trace class is the sum of its diagonal matrix elements in any orthonormal basis;
$\operatorname{Tr}\mathrm{Aq}^*=\overline{\operatorname{Tr}\mathrm{Aq}}$;
and $\operatorname{Tr}(\mathrm{Aq}\cdot\mathrm{Bq})=\operatorname{Tr}(\mathrm{Bq}\cdot\mathrm{Aq})$ if $\mathrm{Aq}$ is of trace class and $\mathrm{Bq}$ is bounded.
Let $\mathrm{Dq}$ be a measurable subset of $\mathbb{R}^3$, and consider an integral operator on $L^2(\mathrm{Dq};\mathbb{C}^4)$ whose kernel is a measurable function on $\mathrm{Dq}\times\mathrm{Dq}$ with values in the $4\times4$ matrices.
If the Hilbert--Schmidt norm $\|\cdot\|_2$ of the value of the kernel at $(y,y')$, as a function of $(y,y')$, lies in $L^2(\mathrm{Dq}\times\mathrm{Dq})$, then the operator is a Hilbert--Schmidt operator, and its norm $\|\cdot\|_2$ is the $L^2$ norm of this function.
For a matrix $\mathrm{Mx}$ on $\mathbb{C}^4$, resp. on ${V_{\mathrm{v}}}$, one has $\|\mathrm{Mx}\|_4^4\le4\,|\mathrm{Mx}|_2^4$, resp. $\|\mathrm{Mx}\|_4^4\le192\,|\mathrm{Mx}|_2^4$.
For a vector $y\in\mathbb{C}^3$, by (I.c1),
\[
\|\gamma\cdot y\|_2^2=\operatorname{tr}\left((\gamma\cdot y)^*\,\gamma\cdot y\right)=\sum_{\mu,\nu}\overline{y_\mu}\,y_\nu\operatorname{tr}(\gamma_\mu\gamma_\nu)=4\,|y|_2^2 .
\]

\item[(Tr2)]
For $\mathrm{Aq},\mathrm{Bq}\in I_{4}$,
\[
\left|\operatorname{Tr}\mathrm{Aq}^4-\operatorname{Tr}\mathrm{Bq}^4\right|\le4\max\left(\|\mathrm{Aq}\|_4,\|\mathrm{Bq}\|_4\right)^3\cdot\|\mathrm{Aq}-\mathrm{Bq}\|_4 .
\]
\begin{proof}
One has $\mathrm{Aq}^4-\mathrm{Bq}^4=\sum_{j=0}^{3}\mathrm{Aq}^j\,(\mathrm{Aq}-\mathrm{Bq})\,\mathrm{Bq}^{3-j}$;
H\"older's inequality for four factors, applied to each term, gives the claim.
\end{proof}

\item[(Tr3)]
Let $\mathrm{Aq}\in I_{4}$, and let $\mathrm{Prq}$ be an orthogonal projection such that $\mathrm{Aq}=\mathrm{Prq}\cdot\mathrm{Aq}$ or $\mathrm{Aq}=\mathrm{Aq}\cdot\mathrm{Prq}$.
Then $\operatorname{Tr}\mathrm{Aq}^4=\operatorname{Tr}(\mathrm{Prq}\cdot\mathrm{Aq}\cdot\mathrm{Prq})^4$, where the right side may be computed in $\operatorname{Ran}\mathrm{Prq}$, and $\|\mathrm{Prq}\cdot\mathrm{Aq}\cdot\mathrm{Prq}\|_4\le\|\mathrm{Aq}\|_4$.
\begin{proof}
Suppose first that $\mathrm{Aq}=\mathrm{Prq}\cdot\mathrm{Aq}$.
Then $\mathrm{Aq}^4=\mathrm{Prq}\cdot\mathrm{Aq}^4$, hence $\operatorname{Tr}\mathrm{Aq}^4=\operatorname{Tr}\left(\mathrm{Prq}\cdot(\mathrm{Prq}\cdot\mathrm{Aq}^4)\right)=\operatorname{Tr}(\mathrm{Prq}\cdot\mathrm{Aq}^4\cdot\mathrm{Prq})$, and
\[
\mathrm{Prq}\cdot\mathrm{Aq}^4\cdot\mathrm{Prq}=\mathrm{Prq}\cdot(\mathrm{Prq}\cdot\mathrm{Aq})^4\cdot\mathrm{Prq}=(\mathrm{Prq}\cdot\mathrm{Aq}\cdot\mathrm{Prq})^4 .
\]
Suppose next that $\mathrm{Aq}=\mathrm{Aq}\cdot\mathrm{Prq}$.
Then $\mathrm{Aq}^4=\mathrm{Aq}^4\cdot\mathrm{Prq}$, hence $\operatorname{Tr}\mathrm{Aq}^4=\operatorname{Tr}\left((\mathrm{Aq}^4\cdot\mathrm{Prq})\cdot\mathrm{Prq}\right)=\operatorname{Tr}(\mathrm{Prq}\cdot\mathrm{Aq}^4\cdot\mathrm{Prq})$, and
\[
\mathrm{Prq}\cdot\mathrm{Aq}^4\cdot\mathrm{Prq}=\mathrm{Prq}\cdot(\mathrm{Aq}\cdot\mathrm{Prq})^4=(\mathrm{Prq}\cdot\mathrm{Aq}\cdot\mathrm{Prq})^4 .
\]
The bound on the norm is contained in (Tr1).
\end{proof}

\item[(Tr4)]
Let $\mathrm{Uq}$ and $\mathrm{Vq}$ be bounded Hermitian operators, and let $\mathrm{Uq}$ have finite rank.
Then $\operatorname{Tr}(\mathrm{Uq}\cdot\mathrm{Vq}\cdot\mathrm{Uq}\cdot\mathrm{Vq})\le\operatorname{Tr}(\mathrm{Uq}^2\cdot\mathrm{Vq}^2)$.
\begin{proof}
The operator $\mathrm{Zq}:=\mathrm{Uq}\cdot\mathrm{Vq}-\mathrm{Vq}\cdot\mathrm{Uq}$ has finite rank, and $\mathrm{Zq}^*=-\mathrm{Zq}$.
By the cyclicity of the trace,
\begin{align*}
0\le\operatorname{Tr}\mathrm{Zq}^*\cdot\mathrm{Zq}=-\operatorname{Tr}\mathrm{Zq}^2
&=-\operatorname{Tr}(\mathrm{Uq}\cdot\mathrm{Vq}\cdot\mathrm{Uq}\cdot\mathrm{Vq})+\operatorname{Tr}(\mathrm{Uq}\cdot\mathrm{Vq}^2\cdot\mathrm{Uq})\\
&\quad+\operatorname{Tr}(\mathrm{Vq}\cdot\mathrm{Uq}^2\cdot\mathrm{Vq})-\operatorname{Tr}(\mathrm{Vq}\cdot\mathrm{Uq}\cdot\mathrm{Vq}\cdot\mathrm{Uq})\\
&=2\operatorname{Tr}(\mathrm{Uq}^2\cdot\mathrm{Vq}^2)-2\operatorname{Tr}(\mathrm{Uq}\cdot\mathrm{Vq}\cdot\mathrm{Uq}\cdot\mathrm{Vq}).
\end{align*}
\end{proof}

\item[(Tr5)]
(The lattice.)
For $\mathrm{hm}:\mathrm{Pb}_n\to\operatorname{End}(\mathrm{Fb})$ put
\[
\mathrm{nrm}_4(\mathrm{hm}):=\left(n^{-3}\sum_{p\in\mathrm{Pb}_n}\|\mathrm{hm}(p)\|_4^4\right)^{\frac14}.
\]
Then, for every $n\mathbb{Z}^3$-periodic scalar function $\mathrm{fs}$,
\[
\left\|\mathrm{Mu}[\mathrm{hm}]\cdot\mathrm{fs}\right\|_4\le\mathrm{nrm}_4(\mathrm{hm})\cdot\|\mathrm{fs}\|_{\ell^4}.
\]
\begin{proof}
Let $\mathrm{Aq}:=\mathrm{Mu}[\mathrm{hm}]\cdot\mathrm{fs}$, $\mathrm{Uq}:=\mathrm{Mu}[\mathrm{hm}^*\cdot\mathrm{hm}]$ and $\mathrm{Vq}:=|\mathrm{fs}|^2$.
By (Tr4),
\[
\|\mathrm{Aq}\|_4^4=\operatorname{Tr}(\mathrm{Aq}^*\cdot\mathrm{Aq})^2=\operatorname{Tr}\left(\overline{\mathrm{fs}}\cdot\mathrm{Uq}\cdot\mathrm{fs}\cdot\overline{\mathrm{fs}}\cdot\mathrm{Uq}\cdot\mathrm{fs}\right)=\operatorname{Tr}(\mathrm{Uq}\cdot\mathrm{Vq}\cdot\mathrm{Uq}\cdot\mathrm{Vq})\le\operatorname{Tr}(\mathrm{Uq}^2\cdot\mathrm{Vq}^2).
\]
With an orthonormal basis $(u_i)$ of $\mathrm{Fb}$,
\begin{align*}
\operatorname{Tr}(\mathrm{Uq}^2\cdot\mathrm{Vq}^2)
&=\sum_{p\in\mathrm{Pb}_n}\sum_i\left(\mathrm{Uq}^2(\mathrm{ew}_p\otimes u_i),|\mathrm{fs}|^4\cdot\mathrm{ew}_p\otimes u_i\right)
=\sum_{p\in\mathrm{Pb}_n}\operatorname{tr}\left((\mathrm{hm}^*\cdot\mathrm{hm})(p)^2\right)\cdot\left(\mathrm{ew}_p,|\mathrm{fs}|^4\cdot\mathrm{ew}_p\right)\\
&=\sum_{p\in\mathrm{Pb}_n}\|\mathrm{hm}(p)\|_4^4\cdot n^{-3}\sum_{c\in\{0,\dots,n-1\}^3}|\mathrm{fs}(c)|^4
=\mathrm{nrm}_4(\mathrm{hm})^4\cdot\|\mathrm{fs}\|_{\ell^4}^4 .
\end{align*}
\end{proof}

\item[(Tr6)]
(The continuum torus.)
Let
\[
\mathrm{Kf}_R:=\frac{2\pi}{R}\left(\mathbb{Z}+\frac12\right)^3,\qquad
\mathrm{Kb}_R:=\frac{2\pi}{R}\,\mathbb{Z}^3,\qquad
\mathrm{Ew}_P(y):=R^{-\frac32}\,\mathrm{e}^{iP\cdot y}\quad(P\in\mathrm{Kf}_R,\ y\in\mathbb{R}^3),
\]
and let $\mathrm{bv}_1,\dots,\mathrm{bv}_4$ be the standard basis of $\mathbb{C}^4$.
Here and in the rest of this chapter the subscript $\mathrm{be}$ of $\mathrm{bv}_{\mathrm{be}}$ is a variable index with the range $\{1,2,3,4\}$;
it is the spinor index, which Section~I.1.9 denotes by $\beta$, a letter that stands for the probe in this chapter.
This index is unrelated to the subscript of $D_{\mathrm{be}}$ and $Q_{\mathrm{be}}$, which refers to the probe $\beta$.
The functions $\mathrm{Ew}_P\cdot\mathrm{bv}_{\mathrm{be}}$, $P\in\mathrm{Kf}_R$, $\mathrm{be}\in\{1,2,3,4\}$, form an orthonormal basis of $L^2_{\mathrm{ap}}(T_R;\mathbb{C}^4)$.
Indeed, multiplication by $\mathrm{e}^{-\frac{i\pi(y_1+y_2+y_3)}{R}}$ is a unitary map onto the $R\mathbb{Z}^3$-periodic functions, and it carries the functions $\mathrm{Ew}_P\cdot\mathrm{bv}_{\mathrm{be}}$ to the exponentials $R^{-\frac32}\mathrm{e}^{i\,\mathrm{kph}\cdot y}\,\mathrm{bv}_{\mathrm{be}}$, $\mathrm{kph}\in\mathrm{Kb}_R$, $\mathrm{be}\in\{1,2,3,4\}$;
and the functions $R^{-\frac32}\mathrm{e}^{i\,\mathrm{kph}\cdot y}$, $\mathrm{kph}\in\mathrm{Kb}_R$, form an orthonormal basis of $L^2$ of a period cube, by the completeness of the exponentials \cite[Chapter~VII, Section~1]{SteinWeiss1971}, stated there for the period $1$, to which the substitution $y\mapsto y/R$ reduces the period $R$.
For a bounded map $\mathrm{hm}:\mathrm{Kf}_R\to\operatorname{End}(\mathbb{C}^4)$ let $\mathrm{Mu}[\mathrm{hm}](\mathrm{Ew}_P\,\xi):=\mathrm{Ew}_P\cdot\mathrm{hm}(P)\xi$ for $\xi\in\mathbb{C}^4$.
Then $S^R=\mathrm{Mu}[\mathrm{Sdr}(\,\cdot\,,\bar{m})]$.
For a bounded measurable $R\mathbb{Z}^3$-periodic scalar function $\mathrm{fs}$,
\begin{align*}
(\mathrm{Ew}_P,\mathrm{fs}\cdot\mathrm{Ew}_{P'})&:=\int_{\text{period cube}}\mathrm{d}y\;\overline{\mathrm{Ew}_P(y)}\cdot\mathrm{fs}(y)\cdot\mathrm{Ew}_{P'}(y)=\mathrm{fc}[\mathrm{fs}](P-P'),\\
\mathrm{fc}[\mathrm{fs}](\mathrm{kph})&:=R^{-3}\int_{\text{period cube}}\mathrm{d}y\;\mathrm{fs}(y)\,\mathrm{e}^{-i\,\mathrm{kph}\cdot y}\qquad(\mathrm{kph}\in\mathrm{Kb}_R).
\end{align*}
Claim: if $\sum_{P\in\mathrm{Kf}_R}\|\mathrm{hm}(P)\|_4^4<\infty$, then $\mathrm{Mu}[\mathrm{hm}]\cdot\mathrm{fs}$ lies in $I_{4}$, and
\[
\left\|\mathrm{Mu}[\mathrm{hm}]\cdot\mathrm{fs}\right\|_4\le\left(R^{-3}\sum_{P\in\mathrm{Kf}_R}\|\mathrm{hm}(P)\|_4^4\right)^{\frac14}\cdot\|\mathrm{fs}\|_{L^4(\text{period cube})}.
\]
\begin{proof}
Let first $\mathrm{hm}$ have finite support.
Then $\mathrm{Aq}:=\mathrm{Mu}[\mathrm{hm}]\cdot\mathrm{fs}$ and $\mathrm{Uq}:=\mathrm{Mu}[\mathrm{hm}^*\cdot\mathrm{hm}]$ have finite rank;
let $\mathrm{Vq}:=|\mathrm{fs}|^2$.
By (Tr4),
\[
\|\mathrm{Aq}\|_4^4=\operatorname{Tr}\left(\overline{\mathrm{fs}}\cdot\mathrm{Uq}\cdot\mathrm{fs}\cdot\overline{\mathrm{fs}}\cdot\mathrm{Uq}\cdot\mathrm{fs}\right)=\operatorname{Tr}(\mathrm{Uq}\cdot\mathrm{Vq}\cdot\mathrm{Uq}\cdot\mathrm{Vq})\le\operatorname{Tr}(\mathrm{Uq}^2\cdot\mathrm{Vq}^2),
\]
and, $\mathrm{hm}$ having finite support, only finitely many terms of the following sums are non-zero:
\begin{align*}
\operatorname{Tr}(\mathrm{Uq}^2\cdot\mathrm{Vq}^2)
&=\sum_{P\in\mathrm{Kf}_R}\sum_{\mathrm{be}=1}^{4}\left(\mathrm{Uq}^2(\mathrm{Ew}_P\cdot\mathrm{bv}_{\mathrm{be}}),|\mathrm{fs}|^4\cdot\mathrm{Ew}_P\cdot\mathrm{bv}_{\mathrm{be}}\right)\\
&=\sum_{P\in\mathrm{Kf}_R}\operatorname{tr}\left((\mathrm{hm}^*\cdot\mathrm{hm})(P)^2\right)\cdot\left(\mathrm{Ew}_P,|\mathrm{fs}|^4\cdot\mathrm{Ew}_P\right)\\
&=\sum_{P\in\mathrm{Kf}_R}\|\mathrm{hm}(P)\|_4^4\cdot R^{-3}\int_{\text{period cube}}\mathrm{d}y\;|\mathrm{fs}(y)|^4 .
\end{align*}
For general $\mathrm{hm}$ put $\mathrm{hm}_i:=\mathrm{hm}\cdot\mathbf{1}[|P|_2\le i]$.
Applying the case of finite support to the differences $\mathrm{hm}_i-\mathrm{hm}_{i'}$, one sees that $\left(\mathrm{Mu}[\mathrm{hm}_i]\cdot\mathrm{fs}\right)_i$ is a Cauchy sequence in $I_{4}$;
it converges in $I_{4}$, and therefore in operator norm.
Since $\sup_{|P|_2>i}|\mathrm{hm}(P)|_2\to0$, its limit in operator norm is $\mathrm{Mu}[\mathrm{hm}]\cdot\mathrm{fs}$.
Letting $i\to\infty$ in the bound for $\mathrm{hm}_i$ gives the bound for $\mathrm{hm}$.
\end{proof}
\end{itemize}

\paragraph{2.4. Sums over momenta}

\begin{named}{Lemma 2.1}
Let $R'>0$, and let $\mathrm{fn}$ be a non-increasing function from $(0,\infty)$ to $[0,\infty)$.
Then
\[
{R'}^{-3}\sum_{P\in\frac{2\pi}{R'}(\mathbb{Z}+\frac12)^3}\mathrm{fn}(|P|_2)\le\pi^{-3}\int_{\mathbb{R}^3}\mathrm{d}^3y\;\mathrm{fn}(|y|_2).
\]
If $\mathrm{Ksub}$ is a subset of $\frac{2\pi}{R'}(\mathbb{Z}+\frac12)^3$ such that all cubes $P+[-\pi/R',\pi/R')^3$, $P\in\mathrm{Ksub}$, lie in the ball $\{|y|_2\le2\,\mathrm{rba}\}$, then
\[
{R'}^{-3}\sum_{P\in\mathrm{Ksub}}\mathrm{fn}(|P|_2)\le\pi^{-3}\int_{|y|_2\le\mathrm{rba}}\mathrm{d}^3y\;\mathrm{fn}(|y|_2).
\]
\end{named}

\begin{proof}
The cubes $P+[-\pi/R',\pi/R')^3$ are pairwise disjoint, and each has the volume $(2\pi/R')^3$.
Let $y$ lie in the cube of $P$.
Then $|y_j-P_j|\le\pi/R'\le|P_j|$, whence $|y_j|\le2|P_j|$, $|y|_2\le2|P|_2$ and $\mathrm{fn}(|P|_2)\le\mathrm{fn}(|y|_2/2)$.
Integration over the cube and summation over $P$ give
\[
(2\pi)^3{R'}^{-3}\sum_P\mathrm{fn}(|P|_2)\le\int\mathrm{d}^3y\;\mathrm{fn}(|y|_2/2),
\]
where the integration extends over $\mathbb{R}^3$, resp. over $\{|y|_2\le2\,\mathrm{rba}\}$.
The substitution $y\mapsto2y$ produces the factor $8$, and $8\,(2\pi)^{-3}=\pi^{-3}$.
\end{proof}

\emph{Consequences.}
Take $R'=n$;
the set $\mathrm{Pb}_n$ is contained in $\frac{2\pi}{n}(\mathbb{Z}+\frac12)^3$.
Since
\[
\int_{\mathbb{R}^3}\mathrm{d}^3y\;\left(|y|_2^2+\mathbf{m}^2\right)^{-2}=4\pi\int_0^\infty\mathrm{d}\,\mathrm{rd}\;\mathrm{rd}^2\left(\mathrm{rd}^2+\mathbf{m}^2\right)^{-2}=\frac{\pi^2}{\mathbf{m}},
\]
one obtains
\begin{equation*}
n^{-3}\sum_{p\in\mathrm{Pb}_n}\left(|p|_2^2+\mathbf{m}^2\right)^{-2}\le\pi^{-3}\,\frac{\pi^2}{\mathbf{m}}=\frac{1}{\pi\mathbf{m}}.
\tag*{[2.5]}
\end{equation*}
Let $p\in\mathrm{Pb}_n$ and $y\in p+[-\pi/n,\pi/n)^3$.
Then $|y|_\infty\le\pi+\pi/n\le3\pi/2$, because $n\ge2$;
so $|y|_2\le3\sqrt3\,\pi/2=2\,\mathrm{Lz}$.
For $\mathrm{lf}=1,2,3$ put
\[
\mathrm{Il}_{\mathrm{lf}}:=n^{-3}\sum_{p\in\mathrm{Pb}_n}\left(|p|_2+\mathbf{m}\right)^{-\mathrm{lf}}.
\]
The second part of Lemma~2.1 gives
\[
\mathrm{Il}_{\mathrm{lf}}\le\pi^{-3}\cdot4\pi\int_0^{\mathrm{Lz}}\mathrm{d}\,\mathrm{rd}\;\mathrm{rd}^2(\mathrm{rd}+\mathbf{m})^{-\mathrm{lf}}=\frac{4}{\pi^2}\int_0^{\mathrm{Lz}}\mathrm{d}\,\mathrm{rd}\;\mathrm{rd}^2(\mathrm{rd}+\mathbf{m})^{-\mathrm{lf}},
\]
and $\mathrm{rd}^2(\mathrm{rd}+\mathbf{m})^{-\mathrm{lf}}$ does not exceed $\mathrm{rd}$, $1$, $(\mathrm{rd}+\mathbf{m})^{-1}$ for $\mathrm{lf}=1,2,3$ respectively.
Hence
\begin{equation*}
\begin{aligned}
&\mathrm{Il}_1\le\frac{4}{\pi^2}\cdot\frac{\mathrm{Lz}^2}{2}=\frac{27}{8}\le4,\qquad
\mathrm{Il}_2\le\frac{4}{\pi^2}\cdot\mathrm{Lz}=\frac{3\sqrt3}{\pi}\le2,\\
&\mathrm{Il}_3\le\frac{4}{\pi^2}\log\left(1+\frac{\mathrm{Lz}}{\mathbf{m}}\right)\le\log\left(1+\frac{5}{\eta\bar{m}}\right).
\end{aligned}
\tag*{[2.6]}
\end{equation*}
Take now $R'=R$.
Then
\begin{align*}
&R^{-3}\sum_{P\in\mathrm{Kf}_R}\left(|P|_2^2+\bar{m}^2\right)^{-2}\le\frac{1}{\pi\bar{m}},\tag*{[2.7]}\\
&R^{-3}\sum_{P\in\mathrm{Kf}_R,\ |P|_\infty\le\frac{\pi}{2\eta}}|P|_2^{-2}\le\pi^{-3}\cdot4\pi\cdot\frac{\sqrt3\,\pi}{2\eta}=\frac{2\sqrt3}{\pi\eta},\tag*{[2.8]}\\
&R^{-3}\sum_{P\in\mathrm{Kf}_R,\ |P|_\infty>\frac{\pi}{2\eta}}|P|_2^{-4}\le\frac{32\,\eta}{3\pi^3}.\tag*{[2.9]}
\end{align*}
For [2.8] the first part of Lemma~2.1 is applied with $\mathrm{fn}(\mathrm{rd})=\mathrm{rd}^{-2}\cdot\mathbf{1}[\mathrm{rd}\le\sqrt3\,\pi/(2\eta)]$.
For [2.9] it is applied with $\mathrm{fn}(\mathrm{rd})=\min\left(\mathrm{rd}_0^{-4},\mathrm{rd}^{-4}\right)$, $\mathrm{rd}_0:=\pi/(2\eta)$, for which
\[
\int\mathrm{d}^3y\;\mathrm{fn}(|y|_2)=\frac{4\pi}{3}\,\mathrm{rd}_0^{-1}+4\pi\,\mathrm{rd}_0^{-1}=\frac{32\,\eta}{3}.
\]

\begin{named}{Lemma 2.2}
For $R\ge1$,
\[
R^{-3}\sum_{\mathrm{kph}\in\mathrm{Kb}_R}\left(1+|\mathrm{kph}|_2\right)^{-4}\le\mathrm{cs}_1 .
\]
\end{named}

\begin{proof}
For $y\in\mathrm{kph}+[-\pi/R,\pi/R)^3$ one has $1+|y|_2\le1+|\mathrm{kph}|_2+\sqrt3\,\pi/R\le\left(1+\sqrt3\,\pi\right)(1+|\mathrm{kph}|_2)$.
Therefore
\[
\left(\frac{2\pi}{R}\right)^3\left(1+|\mathrm{kph}|_2\right)^{-4}\le\left(1+\sqrt3\,\pi\right)^4\int_{\mathrm{kph}+[-\pi/R,\pi/R)^3}\mathrm{d}^3y\;(1+|y|_2)^{-4}.
\]
The cubes $\mathrm{kph}+[-\pi/R,\pi/R)^3$, $\mathrm{kph}\in\mathrm{Kb}_R$, form a partition of $\mathbb{R}^3$, and
\[
\int_{\mathbb{R}^3}\mathrm{d}^3y\;(1+|y|_2)^{-4}=4\pi\int_0^\infty\mathrm{d}\,\mathrm{rd}\;\mathrm{rd}^2(1+\mathrm{rd})^{-4}=\frac{4\pi}{3}.
\]
Finally $\left(1+\sqrt3\,\pi\right)^4\,(2\pi)^{-3}\cdot\frac{4\pi}{3}=\mathrm{cs}_1$.
\end{proof}

\paragraph{2.5. The probe in position space}

\emph{Representation.}
By the definitions of Section~G5.1.3 and Section~I.7.6, every smooth function $\mathrm{fs}$ satisfies
\[
(\delta_{\eta,\mu}\mathrm{fs})(y)=\int_{-\frac12}^{\frac12}\mathrm{d}\,\mathrm{sv}\;(\partial_\mu\mathrm{fs})(y+\mathrm{sv}\cdot\eta\,e_{\mu}),
\]
and $\partial_\mu J_{\eta}=J_{\eta}\partial_\mu$.
Hence
\begin{equation*}
\mathrm{gs}_\nu[\beta](y)=-\sum_{\mu\ne\nu}\int_{-\frac12}^{\frac12}\mathrm{d}\,\mathrm{sv}\;\eta^{-3}\int_{[-\frac{\eta}{2},\frac{\eta}{2}]^3}\mathrm{d}y'\;(\partial_\mu\beta_{\mu\nu})(y+\mathrm{sv}\cdot\eta\,e_{\mu}+y').
\tag*{[2.10]}
\end{equation*}

\emph{The interpolating field.}
For $y\in\mathbb{R}^3$ let
\[
\mathrm{ait}_\nu(y):=\sum_{z\in\mathbb{Z}^3}\mathrm{gs}_\nu[\beta]\left(y+\frac{\eta}{2}\,e_{\nu}+Rz\right);
\]
the sum is finite, and $\mathrm{ait}_\nu$ is smooth and $R\mathbb{Z}^3$-periodic.
Proposition~35.3\,(b1) gives
\begin{equation*}
\mathrm{alat}_\nu(c)=a_{\eta}[\beta]_\nu(\eta c)=\mathrm{ait}_\nu(\eta c)\qquad(c\in\mathbb{Z}^3).
\tag*{[2.11]}
\end{equation*}
By [2.10],
\[
\mathrm{ait}_\nu(y)=-\sum_{\mu\ne\nu}\int_{-\frac12}^{\frac12}\mathrm{d}\,\mathrm{sv}\;\eta^{-3}\int_{[-\frac{\eta}{2},\frac{\eta}{2}]^3}\mathrm{d}y'\;(\partial_\mu\beta^R_{\mu\nu})(y+\mathrm{dsh}),\qquad
\mathrm{dsh}:=\frac{\eta}{2}\,e_{\nu}+\mathrm{sv}\cdot\eta\,e_{\mu}+y'.
\]
Because $\mu\ne\nu$, the $\nu$-th coordinate of $\mathrm{dsh}$ ranges over $[0,\eta]$, the $\mu$-th over $[-\eta,\eta]$ and the third over $[-\eta/2,\eta/2]$;
thus $|\mathrm{dsh}|_\infty\le\eta$ and $|\mathrm{dsh}|_1\le5\eta/2$.
So $\mathrm{ait}_\nu$ arises by averaging the translates of $-\partial_\mu\beta^R_{\mu\nu}$ by vectors $\mathrm{dsh}$ with $|\mathrm{dsh}|_\infty\le\eta$ and summing over $\mu\ne\nu$, while $a^R_\nu$ is the same sum without translation, i.e. with $\mathrm{dsh}=0$.
Since the sets $\operatorname{supp}\beta+Rz$, $z\in\mathbb{Z}^3$, are pairwise disjoint, every derivative of $\beta^R$ has the same sup norm as the corresponding derivative of $\beta$, and its $L^1$ norm over a period cube equals the $L^1(\mathbb{R}^3)$ norm of the corresponding derivative of $\beta$.
It follows that
\begin{equation*}
\begin{aligned}
&|\mathrm{ait}_\nu|_\infty\le\mathrm{Ab}_1,\qquad|a^R_\nu|_\infty=|a_\nu|_\infty\le\mathrm{Ab}_1,\\
&|\partial_j\mathrm{ait}_\nu|_\infty\le\mathrm{Ab}_2,\qquad\left\|\partial_j^4\mathrm{ait}_\nu\right\|_{L^1(\text{period cube})}\le\mathrm{Ab}_5,
\end{aligned}
\tag*{[2.12]}
\end{equation*}
\begin{equation*}
\left|\mathrm{ait}_\nu(y)-a^R_\nu(y)\right|\le\sup|\mathrm{dsh}|_1\cdot\mathrm{Ab}_2\le3\eta\cdot\mathrm{Ab}_2,
\tag*{[2.13]}
\end{equation*}
\begin{equation*}
\mathrm{ait}_\nu(y)\ne0\ \text{or}\ a^R_\nu(y)\ne0\quad\text{only if}\quad\inf\left\{|y-x|_\infty:\ x\in Q_{\mathrm{be}}+R\mathbb{Z}^3\right\}\le\eta .
\tag*{[2.14]}
\end{equation*}
The set appearing in [2.14] consists of the $R\mathbb{Z}^3$-translates of one cube of side $D_{\mathrm{be}}+2\eta$;
so its measure in a period cube does not exceed $(D_{\mathrm{be}}+2\eta)^3\le\mathrm{Vbe}$.
If a class in $\mathbb{Z}^3/n\mathbb{Z}^3$ has a representative $c$ with $\eta c$ in this set, then it has a representative $c'$ with $\eta c'$ in that cube; and a closed cube of side $D_{\mathrm{be}}+2\eta$ contains at most $\left((D_{\mathrm{be}}+2\eta)/\eta+1\right)^3$ points of $\eta\mathbb{Z}^3$.
Let
\[
N_s:=\#\left|\left\{c\in\mathbb{Z}^3/n\mathbb{Z}^3:\ \mathrm{alat}_\nu(c)\ne0\text{ for some }\nu\right\}\right|.
\]
Then
\begin{equation*}
\eta^3N_s\le(D_{\mathrm{be}}+3\eta)^3\le\mathrm{Vbe},\qquad|\mathrm{alat}_\nu|_\infty\le\mathrm{Ab}_1,\qquad\left|\mathrm{alat}^-_\nu(c)-\mathrm{alat}_\nu(c)\right|\le\eta\cdot\mathrm{Ab}_2,
\tag*{[2.15]}
\end{equation*}
\begin{equation*}
\begin{aligned}
&\|\mathrm{ait}_\nu\|_{L^4(\text{period cube})},\ \|a^R_\nu\|_{L^4(\text{period cube})}\le\mathrm{Ab}_1\cdot\mathrm{Vbe}^{\frac14}=\mathrm{kpr},\\
&\|\mathrm{ait}_\nu-a^R_\nu\|_{L^4(\text{period cube})}\le3\eta\cdot\mathrm{Ab}_2\cdot\mathrm{Vbe}^{\frac14},
\end{aligned}
\tag*{[2.16]}
\end{equation*}
\begin{equation*}
\begin{aligned}
&\|\mathrm{bph}^{\pm}_{\nu}\|_{\ell^4}=\|\eta\cdot\mathrm{alat}_\nu\|_{\ell^4}\le\left(\eta^4\cdot\mathrm{Ab}_1^4\cdot N_s\right)^{\frac14}\le\eta^{\frac14}\cdot\mathrm{kpr},\\
&\left\|\eta\left(\mathrm{alat}^-_\nu-\mathrm{alat}_\nu\right)\right\|_{\ell^4}\le\left(\eta^8\cdot\mathrm{Ab}_2^4\cdot2N_s\right)^{\frac14}\le\left(2\eta^5\cdot\mathrm{Ab}_2^4\cdot\mathrm{Vbe}\right)^{\frac14}.
\end{aligned}
\tag*{[2.17]}
\end{equation*}
Here one uses that $\mathrm{bph}^-_{\nu}$ is a translate of $-\mathrm{bph}^+_{\nu}$, and that $\mathrm{alat}^-_\nu-\mathrm{alat}_\nu$ vanishes outside at most $2N_s$ classes.

\paragraph{2.6. The continuum propagator; clause (a)}

\begin{itemize}
\item[(K-0)]
Let $\mathrm{Hq}$ denote one of the spaces $L^2(\mathbb{R}^3;\mathbb{C}^4)$, $L^2_{\mathrm{ap}}(T_R;\mathbb{C}^4)$, with norm $\|\cdot\|$, and let $\mathrm{Dq}$ be $\mathbb{R}^3$ in the first case and a period cube in the second.
Let $\mathrm{krn}$ be a measurable function on $\mathbb{R}^3$, with values in the $4\times4$ matrices, such that
\[
\|\mathrm{krn}\|_{L^1}:=\int_{\mathbb{R}^3}\mathrm{d}y''\;|\mathrm{krn}(y'')|_2<\infty .
\]
Then, for each $\psi\in\mathrm{Hq}$, the integral
\[
(\mathrm{krn}*\psi)(y):=\int_{\mathbb{R}^3}\mathrm{d}y''\;\mathrm{krn}(y'')\,\psi(y-y'')
\]
is absolutely convergent for almost all $y$;
the function $\mathrm{krn}*\psi$ lies in $\mathrm{Hq}$ and satisfies
\[
\|\mathrm{krn}*\psi\|\le\|\mathrm{krn}\|_{L^1}\cdot\|\psi\| .
\]
This is Minkowski's integral inequality \cite[Theorem~2.4]{LiebLoss2001};
for $\mathrm{Hq}=L^2(\mathbb{R}^3;\mathbb{C}^4)$ it is Young's inequality with the exponents $1,2,2$ \cite[Theorem~4.2]{LiebLoss2001}.
\begin{proof}
Let
\[
\mathrm{Fq}(y):=\int_{\mathbb{R}^3}\mathrm{d}y''\;|\mathrm{krn}(y'')|_2\cdot|\psi(y-y'')|_2,
\]
with values in $[0,\infty]$;
if $\mathrm{Hq}=L^2_{\mathrm{ap}}(T_R;\mathbb{C}^4)$, this function is $R\mathbb{Z}^3$-periodic.
For every $y''$ one has $\int_{\mathrm{Dq}}\mathrm{d}y\;|\psi(y-y'')|_2^2=\|\psi\|^2$.
Let $\mathrm{fn}\ge0$ be measurable on $\mathrm{Dq}$.
By Tonelli's theorem and the Cauchy--Schwarz inequality in the inner integral,
\[
\int_{\mathrm{Dq}}\mathrm{d}y\;\mathrm{fn}(y)\cdot\mathrm{Fq}(y)=\int_{\mathbb{R}^3}\mathrm{d}y''\;|\mathrm{krn}(y'')|_2\left(\int_{\mathrm{Dq}}\mathrm{d}y\;\mathrm{fn}(y)\cdot|\psi(y-y'')|_2\right)\le\|\mathrm{krn}\|_{L^1}\cdot\|\mathrm{fn}\|_{L^2(\mathrm{Dq})}\cdot\|\psi\| .
\]
For $i=1,2,\dots$ the function $\mathrm{fn}_i:=\min(\mathrm{Fq},i)\cdot\mathbf{1}[|y|_2\le i]$ lies in $L^2(\mathrm{Dq})$, and the last inequality, applied to it, gives
\[
\|\mathrm{fn}_i\|_{L^2(\mathrm{Dq})}^2\le\int_{\mathrm{Dq}}\mathrm{d}y\;\mathrm{fn}_i(y)\cdot\mathrm{Fq}(y)\le\|\mathrm{krn}\|_{L^1}\cdot\|\mathrm{fn}_i\|_{L^2(\mathrm{Dq})}\cdot\|\psi\| ,
\]
so that $\|\mathrm{fn}_i\|_{L^2(\mathrm{Dq})}\le\|\mathrm{krn}\|_{L^1}\cdot\|\psi\|$.
As $i\to\infty$, $\mathrm{fn}_i$ increases to $\mathrm{Fq}$;
by monotone convergence, $\|\mathrm{Fq}\|_{L^2(\mathrm{Dq})}\le\|\mathrm{krn}\|_{L^1}\cdot\|\psi\|$.
In particular $\mathrm{Fq}$ is finite almost everywhere.
Wherever $\mathrm{Fq}(y)$ is finite, the integral $(\mathrm{krn}*\psi)(y)$ converges absolutely and $|(\mathrm{krn}*\psi)(y)|_2\le\mathrm{Fq}(y)$.
Finally, $\mathrm{krn}*\psi$ is measurable by Fubini's theorem, and it is antiperiodic if $\psi$ is.
\end{proof}
\end{itemize}

\emph{Kernels.}
For $y\in\mathbb{R}^3$, $y\ne0$, write $\mathrm{rd}=|y|_2$ and put
\[
\mathrm{Yk}(y):=\frac{\exp(-\bar{m}\cdot\mathrm{rd})}{4\pi\,\mathrm{rd}},\qquad\mathrm{Sx}(y):=(\bar{m}-\gamma\cdot\partial)\,\mathrm{Yk}(y).
\]

\begin{itemize}
\item[(K-1)]
One has $\partial_j\mathrm{Yk}(y)=-\left(\bar{m}+\frac{1}{\mathrm{rd}}\right)\mathrm{Yk}(y)\,\frac{y_j}{\mathrm{rd}}$.
Hence
\[
\mathrm{Sx}(y)=\left[\bar{m}+\left(\bar{m}+\frac{1}{\mathrm{rd}}\right)\frac{\gamma\cdot y}{\mathrm{rd}}\right]\mathrm{Yk}(y),
\]
and since $|\gamma\cdot y/\mathrm{rd}|_2=1$ by (I.c1),
\begin{equation*}
\begin{aligned}
&|\mathrm{Sx}(y)|_2\le\left(2\bar{m}+\frac{1}{\mathrm{rd}}\right)\frac{\exp(-\bar{m}\cdot\mathrm{rd})}{4\pi\,\mathrm{rd}},\qquad
\|\mathrm{Sx}\|_{L^1(\mathbb{R}^3)}\le\int_0^\infty\mathrm{d}\,\mathrm{rd}\;(2\bar{m}\cdot\mathrm{rd}+1)\exp(-\bar{m}\cdot\mathrm{rd})=\frac{3}{\bar{m}},\\
&|\mathrm{Sx}(y)|_2\le\frac{3}{4\pi}\exp(-\bar{m}\cdot|y|_2)\quad\text{for }|y|_2\ge1\quad(\bar{m}\le1).
\end{aligned}
\tag*{[2.18]}
\end{equation*}

\item[(K-2)]
The function $\mathrm{Yk}$ is integrable, and for $p\in\mathbb{R}^3$, $p\ne0$,
\[
\int\mathrm{d}y\;\mathrm{Yk}(y)\,\mathrm{e}^{-ip\cdot y}=\frac{4\pi}{|p|_2}\int_0^\infty\mathrm{d}\,\mathrm{rd}\;\mathrm{rd}\,\sin(|p|_2\cdot\mathrm{rd})\,\frac{\exp(-\bar{m}\cdot\mathrm{rd})}{4\pi\,\mathrm{rd}}=\left(\bar{m}^2+|p|_2^2\right)^{-1};
\]
by continuity this also holds at $p=0$.
Moreover
\[
\|\mathrm{Yk}\|_{L^2}^2=\int_0^\infty\mathrm{d}\,\mathrm{rd}\;\frac{\exp(-2\bar{m}\cdot\mathrm{rd})}{4\pi}=\frac{1}{8\pi\bar{m}}.
\]
The gradient of $\mathrm{Yk}$, computed pointwise for $y\ne0$, is integrable and coincides with the gradient of $\mathrm{Yk}$ in the sense of distributions.
Indeed, let $\mathrm{tf}$ be a test function.
On $\{|y|_2>\mathrm{rd}_0\}$ the field $\mathrm{Yk}\cdot\mathrm{tf}\cdot e_j$ is smooth up to the sphere $\{|y|_2=\mathrm{rd}_0\}$ and has compact support, and the outer normal on this sphere is $-y/\mathrm{rd}_0$.
The divergence theorem gives
\[
\int\mathrm{d}y\;\mathrm{Yk}\cdot\partial_j\mathrm{tf}=\lim_{\mathrm{rd}_0\to0}\int_{|y|_2>\mathrm{rd}_0}\mathrm{d}y\;\mathrm{Yk}\cdot\partial_j\mathrm{tf}
=\lim_{\mathrm{rd}_0\to0}\left[-\int_{|y|_2>\mathrm{rd}_0}\mathrm{d}y\;(\partial_j\mathrm{Yk})\cdot\mathrm{tf}-\int_{|y|_2=\mathrm{rd}_0}\mathrm{d}S\;\mathrm{Yk}\cdot\mathrm{tf}\cdot\frac{y_j}{\mathrm{rd}_0}\right],
\]
and the modulus of the surface term is at most $|\mathrm{tf}|_\infty\cdot4\pi\,\mathrm{rd}_0^2/(4\pi\,\mathrm{rd}_0)=|\mathrm{tf}|_\infty\cdot\mathrm{rd}_0$.
The resulting identity $\int\mathrm{d}y\;\mathrm{Yk}\cdot\partial_j\mathrm{tf}=-\int\mathrm{d}y\;(\partial_j\mathrm{Yk})\cdot\mathrm{tf}$ extends to the function $\mathrm{e}^{-ip\cdot y}$ in the place of $\mathrm{tf}$.
To see this, fix a test function $\mathrm{tf}_0$ that equals $1$ in a neighbourhood of $0$, insert $\mathrm{tf}(y)=\mathrm{e}^{-ip\cdot y}\,\mathrm{tf}_0(y/\mathrm{xq})$ into the identity, and let $\mathrm{xq}\to\infty$:
the functions $\mathrm{tf}$ and
\[
\partial_j\mathrm{tf}=\left(-ip_j\,\mathrm{tf}_0(y/\mathrm{xq})+\mathrm{xq}^{-1}\cdot(\partial_j\mathrm{tf}_0)(y/\mathrm{xq})\right)\mathrm{e}^{-ip\cdot y}
\]
are bounded uniformly in $\mathrm{xq}\ge1$ and converge pointwise to $\mathrm{e}^{-ip\cdot y}$ and to $-ip_j\,\mathrm{e}^{-ip\cdot y}$;
as $\mathrm{Yk}$ and $\partial_j\mathrm{Yk}$ are integrable, dominated convergence applies.
The identity obtained in this way reads
\[
\int\mathrm{d}y\;(\partial_j\mathrm{Yk})(y)\,\mathrm{e}^{-ip\cdot y}=ip_j\int\mathrm{d}y\;\mathrm{Yk}(y)\,\mathrm{e}^{-ip\cdot y}.
\]
Consequently
\[
\int\mathrm{d}y\;\mathrm{Sx}(y)\,\mathrm{e}^{-ip\cdot y}=\frac{\bar{m}-i\gamma\cdot p}{\bar{m}^2+|p|_2^2}=\mathrm{Sdr}(p,\bar{m}).
\]

\item[(K-3)]
For $\psi\in L^2(\mathbb{R}^3;\mathbb{C}^4)$ one has $S_{\bar{m}}\psi=\mathrm{Sx}*\psi$, and the multiplier
\[
\left(\bar{m}^2+|p|_2^2\right)^{-1}=\mathrm{Sdr}(p,\bar{m})^*\cdot\mathrm{Sdr}(p,\bar{m})=\mathrm{Sdr}(p,\bar{m})\cdot\mathrm{Sdr}(p,\bar{m})^*
\]
is the convolution with $\mathrm{Yk}\,1_4$.
This follows from (K-2) and the convolution theorem for an integrable function and a function of $L^2$ \cite[Theorem~5.8]{LiebLoss2001}.
\begin{proof}
Let $\mathrm{krn}$ be $\mathrm{Sx}$ or $\mathrm{Yk}\,1_4$, and let $\widehat{\mathrm{krn}}(p):=\int\mathrm{d}y\;\mathrm{krn}(y)\,\mathrm{e}^{-ip\cdot y}$.
By (K-2), $\widehat{\mathrm{krn}}(p)$ is $\mathrm{Sdr}(p,\bar{m})$, resp. $\left(\bar{m}^2+|p|_2^2\right)^{-1}1_4$, and $|\widehat{\mathrm{krn}}(p)|_2\le\|\mathrm{krn}\|_{L^1}$.
Let $\psi\in L^1\cap L^2$.
Then $\mathrm{krn}*\psi$ lies in $L^2$ by (K-0), and in $L^1$, because Tonelli's theorem gives $\int\mathrm{d}y\;|(\mathrm{krn}*\psi)(y)|_2\le\|\mathrm{krn}\|_{L^1}\int\mathrm{d}y\;|\psi(y)|_2$.
By Fubini's theorem
\[
\int\mathrm{d}y\;(\mathrm{krn}*\psi)(y)\,\mathrm{e}^{-ip\cdot y}=\int\mathrm{d}y''\;\mathrm{krn}(y'')\,\mathrm{e}^{-ip\cdot y''}\left(\int\mathrm{d}y\;\psi(y-y'')\,\mathrm{e}^{-ip\cdot(y-y'')}\right)=\widehat{\mathrm{krn}}(p)\cdot\hat{\psi}(p).
\]
Hence $\mathrm{krn}*\psi$ coincides with the image of $\psi$ under the multiplier $\widehat{\mathrm{krn}}$, the two functions having the same Fourier transform.
Both $\psi\mapsto\mathrm{krn}*\psi$ (by (K-0)) and the multiplier $\widehat{\mathrm{krn}}$ (as $\widehat{\mathrm{krn}}$ is bounded) are bounded operators on $L^2(\mathbb{R}^3;\mathbb{C}^4)$, and they agree on the dense subspace $L^1\cap L^2$;
so they are equal.
\end{proof}
\end{itemize}

\begin{named}{Lemma 2.3}
Let $\mathrm{fs}$ be a scalar function on $\mathbb{R}^3$ that is bounded, measurable and of compact support.
Then $S_{\bar{m}}\cdot\mathrm{fs}$ and $S_{\bar{m}}^*\cdot\mathrm{fs}$ lie in $I_{4}$ on $L^2(\mathbb{R}^3;\mathbb{C}^4)$, and
\[
\|S_{\bar{m}}\cdot\mathrm{fs}\|_4^4=\|S_{\bar{m}}^*\cdot\mathrm{fs}\|_4^4\le\frac{\|\mathrm{fs}\|_{L^4}^4}{2\pi\bar{m}}.
\]
Moreover $S_{\bar{m}}\cdot\mathrm{va}$ lies in $I_{4}$, and
\[
\|S_{\bar{m}}\cdot\mathrm{va}\|_4^4\le\frac{1}{2\pi\bar{m}}\int\mathrm{d}y\;|a(y)|_2^4\le\mathrm{kn}_2^4 .
\]
\end{named}

\begin{proof}
The operator $\mathrm{Aq}:=S_{\bar{m}}\cdot\mathrm{fs}$ is bounded.
By (K-3), $\mathrm{Aq}^*\cdot\mathrm{Aq}=\overline{\mathrm{fs}}\cdot S_{\bar{m}}^*S_{\bar{m}}\cdot\mathrm{fs}$ is the integral operator with the kernel $\overline{\mathrm{fs}(y)}\cdot\mathrm{Yk}(y-y')\cdot\mathrm{fs}(y')\,1_4$, and $\|1_4\|_2^2=4$.
If $\mathrm{fn}$ and $\mathrm{krn}$ are non-negative measurable scalar functions on $\mathbb{R}^3$, then
\[
\int\mathrm{d}y\int\mathrm{d}y'\;\mathrm{fn}(y)\cdot\mathrm{krn}(y-y')\cdot\mathrm{fn}(y')=\int\mathrm{d}y''\;\mathrm{krn}(y'')\left(\int\mathrm{d}y\;\mathrm{fn}(y)\cdot\mathrm{fn}(y-y'')\right)\le\|\mathrm{krn}\|_{L^1}\cdot\|\mathrm{fn}\|_{L^2}^2 ,
\]
by Tonelli's theorem with $y''=y-y'$, the Cauchy--Schwarz inequality in the inner integral, and the invariance of the norm of $L^2$ under translations.
Hence
\[
\|\mathrm{Aq}^*\cdot\mathrm{Aq}\|_2^2=4\int\mathrm{d}y\int\mathrm{d}y'\;|\mathrm{fs}(y)|^2\cdot\mathrm{Yk}(y-y')^2\cdot|\mathrm{fs}(y')|^2\le4\,\|\mathrm{Yk}\|_{L^2}^2\cdot\left\||\mathrm{fs}|^2\right\|_{L^2}^2=\frac{\|\mathrm{fs}\|_{L^4}^4}{2\pi\bar{m}}.
\]
By (Tr1), $\mathrm{Aq}$ lies in $I_{4}$, and $\|\mathrm{Aq}\|_4^4=\|\mathrm{Aq}^*\cdot\mathrm{Aq}\|_2^2$.
The operator $(S_{\bar{m}}^*\cdot\mathrm{fs})^*(S_{\bar{m}}^*\cdot\mathrm{fs})=\overline{\mathrm{fs}}\cdot S_{\bar{m}}S_{\bar{m}}^*\cdot\mathrm{fs}$ has the same kernel, by (K-3);
so $S_{\bar{m}}^*\cdot\mathrm{fs}$ lies in $I_{4}$ and has the same norm.
As to $\mathrm{va}$: one has $\mathrm{va}^*=-\mathrm{va}=-i\gamma\cdot a$, and $(S_{\bar{m}}\cdot\mathrm{va})^*(S_{\bar{m}}\cdot\mathrm{va})$ has the kernel $(\gamma\cdot a(y))\,\mathrm{Yk}(y-y')\,(\gamma\cdot a(y'))$.
By (I.c1),
\[
\left\|(\gamma\cdot a(y))(\gamma\cdot a(y'))\right\|_2^2=\operatorname{tr}\left((\gamma\cdot a(y'))(\gamma\cdot a(y))^2(\gamma\cdot a(y'))\right)=4\,|a(y)|_2^2\,|a(y')|_2^2 ,
\]
and therefore $\|S_{\bar{m}}\cdot\mathrm{va}\|_4^4\le4\,\|\mathrm{Yk}\|_{L^2}^2\cdot\left\||a|_2^2\right\|_{L^2}^2$.
It remains to note that $|a|_2^2\le3\,\mathrm{Ab}_1^2$ and that $a=0$ outside $Q_{\mathrm{be}}$;
hence $\int\mathrm{d}y\;|a(y)|_2^4\le9\,\mathrm{Ab}_1^4\cdot D_{\mathrm{be}}^3$.
\end{proof}

\begin{named}{Lemma 2.4}
Let $\mathrm{Mx}$ be a $4\times4$ matrix, and let $\mathrm{fs}$ be a scalar function on $\mathbb{R}^3$ that is bounded, measurable and $R\mathbb{Z}^3$-periodic.
Then $S^R\cdot\mathrm{Mx}\cdot\mathrm{fs}$ and $(S^R)^*\cdot\mathrm{Mx}\cdot\mathrm{fs}$ lie in $I_{4}$ on $L^2_{\mathrm{ap}}(T_R;\mathbb{C}^4)$, and their norms $\|\cdot\|_4$ are at most
\[
|\mathrm{Mx}|_2\cdot\left(\frac{4}{\pi\bar{m}}\right)^{\frac14}\cdot\|\mathrm{fs}\|_{L^4(\text{period cube})}.
\]
Moreover $S^R\cdot\mathrm{va}^R$ lies in $I_{4}$, and $\|S^R\cdot\mathrm{va}^R\|_4\le\mathrm{kn}_0$.
\end{named}

\begin{proof}
One has $S^R\cdot\mathrm{Mx}=\mathrm{Mu}[\mathrm{Sdr}(\,\cdot\,,\bar{m})\cdot\mathrm{Mx}]$ and
\[
\left\|\mathrm{Sdr}(P,\bar{m})\cdot\mathrm{Mx}\right\|_4^4\le4\,|\mathrm{Mx}|_2^4\cdot\left(|P|_2^2+\bar{m}^2\right)^{-2};
\]
the same bound holds with $\mathrm{Sdr}(P,\bar{m})^*$ in the place of $\mathrm{Sdr}(P,\bar{m})$.
The first assertion follows from (Tr6) and [2.7].
Further, $S^R\cdot\mathrm{va}^R=\sum_\nu S^R\,i\gamma_\nu\,a^R_\nu$, $|\gamma_\nu|_2=1$, and [2.16] gives
\[
\|S^R\cdot\mathrm{va}^R\|_4\le3\left(\frac{4}{\pi\bar{m}}\right)^{\frac14}\cdot\mathrm{kpr}=\mathrm{kn}_0 .
\]
\end{proof}

\begin{proof}[Proof of Proposition~37.2\,(a)]
Lemmas~2.3 and~2.4 place the operators $S_{\bar{m}}\cdot\mathrm{va}$ and $S^R\cdot\mathrm{va}^R$ in $I_{4}$.
Their fourth powers are then of trace class by H\"older's inequality, and therefore the numbers
\[
b(f)=-6\operatorname{Tr}\left(S_{\bar{m}}\cdot\mathrm{va}\right)^4\qquad\text{and}\qquad b_{R}(f)=-6\operatorname{Tr}\left(S^R\cdot\mathrm{va}^R\right)^4
\]
are well defined.

\emph{Reality.}
The two triples $(\mathrm{Sz},\mathrm{vz},\mathrm{az}):=(S_{\bar{m}},\mathrm{va},a)$ on $L^2(\mathbb{R}^3;\mathbb{C}^4)$ and $(\mathrm{Sz},\mathrm{vz},\mathrm{az}):=(S^R,\mathrm{va}^R,a^R)$ on $L^2_{\mathrm{ap}}(T_R;\mathbb{C}^4)$ are treated together.
Let $\mathrm{fa}:=|\mathrm{az}|_2^{\frac12}$, a bounded measurable function with compact support, resp. $R\mathbb{Z}^3$-periodic;
let $\mathrm{unv}:=\mathrm{az}/|\mathrm{az}|_2$ where $\mathrm{az}\ne0$, and $\mathrm{unv}:=0$ elsewhere.
Then $\mathrm{vz}=\mathrm{fa}\cdot(i\gamma\cdot\mathrm{unv})\cdot\mathrm{fa}$, where $i\gamma\cdot\mathrm{unv}$ is a bounded multiplication operator.
By Lemma~2.3, resp. by Lemma~2.4 with $\mathrm{Mx}=1$, the operators $\mathrm{Sz}\cdot\mathrm{fa}$ and $\mathrm{Sz}^*\cdot\mathrm{fa}$ lie in $I_{4}$;
hence so does $\mathrm{fa}\cdot\mathrm{Sz}=(\mathrm{Sz}^*\cdot\mathrm{fa})^*$, and H\"older's inequality gives
\[
\mathrm{Gsv}:=\mathrm{Sz}\cdot\mathrm{vz}\cdot\mathrm{Sz}=(\mathrm{Sz}\cdot\mathrm{fa})\,(i\gamma\cdot\mathrm{unv})\,(\mathrm{fa}\cdot\mathrm{Sz})\in I_{2}.
\]
Moreover $(\mathrm{Sz}\cdot\mathrm{vz})^4=\mathrm{Gsv}\cdot\mathrm{vz}\cdot\mathrm{Gsv}\cdot\mathrm{vz}$.
Multiplication by the constant matrix $\gamma_5$ is a unitary involution.
By (I.c1), $\gamma_5\,\mathrm{Sdr}(p,\bar{m})\,\gamma_5=(\bar{m}-i\gamma\cdot p)^{-1}=\mathrm{Sdr}(p,\bar{m})^*$, so that $\gamma_5\,\mathrm{Sz}\,\gamma_5=\mathrm{Sz}^*$;
and $\gamma_5\,\mathrm{vz}\,\gamma_5=-\mathrm{vz}=\mathrm{vz}^*$.
Consequently $\gamma_5\,\mathrm{Gsv}\,\gamma_5=\mathrm{Sz}^*\cdot\mathrm{vz}^*\cdot\mathrm{Sz}^*=\mathrm{Gsv}^*$, and
\begin{align*}
\operatorname{Tr}(\mathrm{Gsv}\cdot\mathrm{vz}\cdot\mathrm{Gsv}\cdot\mathrm{vz})
&=\operatorname{Tr}(\gamma_5\,\mathrm{Gsv}\cdot\mathrm{vz}\cdot\mathrm{Gsv}\cdot\mathrm{vz}\,\gamma_5)
=\operatorname{Tr}(\mathrm{Gsv}^*\cdot\mathrm{vz}^*\cdot\mathrm{Gsv}^*\cdot\mathrm{vz}^*)
=\operatorname{Tr}\left((\mathrm{vz}\cdot\mathrm{Gsv}\cdot\mathrm{vz}\cdot\mathrm{Gsv})^*\right)\\
&=\overline{\operatorname{Tr}\left(\mathrm{vz}\cdot(\mathrm{Gsv}\cdot\mathrm{vz}\cdot\mathrm{Gsv})\right)}
=\overline{\operatorname{Tr}\left((\mathrm{Gsv}\cdot\mathrm{vz}\cdot\mathrm{Gsv})\cdot\mathrm{vz}\right)} .
\end{align*}
The last equality is $\operatorname{Tr}(\mathrm{Bq}\cdot\mathrm{Aq})=\operatorname{Tr}(\mathrm{Aq}\cdot\mathrm{Bq})$ for the bounded operator $\mathrm{Bq}=\mathrm{vz}$ and the operator $\mathrm{Aq}=\mathrm{Gsv}\cdot\mathrm{vz}\cdot\mathrm{Gsv}$, which is of trace class as the product of the two Hilbert--Schmidt operators $\mathrm{Gsv}$ and $\mathrm{vz}\cdot\mathrm{Gsv}$.
Hence $\operatorname{Tr}(\mathrm{Sz}\cdot\mathrm{vz})^4$ is real.
\end{proof}

\paragraph{2.7. The probe in momentum space.}
Let $\mathrm{Gq}_n:=\frac{2\pi}{n}\mathbb{Z}^3/2\pi\mathbb{Z}^3$;
this is a group with $n^3$ elements, whose elements are denoted $\mathrm{pd}$.
It acts simply transitively on $\frac{2\pi}{n}\left(\mathbb{Z}+\frac12\right)^3/2\pi\mathbb{Z}^3$,
and $\mathrm{Pb}_n$ is a set of representatives of the latter quotient;
accordingly, for $p,p'\in\mathrm{Pb}_n$ the symbol $p-p'$ stands for the class of the difference in $\mathrm{Gq}_n$.
For $j\ge1$ and $\mathrm{pd}\in\frac{2\pi}{n}\mathbb{Z}^3$ define
\[
\mathrm{alc}^{(j)}_\nu(\mathrm{pd}):=n^{-3}\sum_{c\in\{0,\dots,n-1\}^3}\mathrm{alat}_\nu(c)^j\,\mathrm{e}^{-i\,\mathrm{pd}\cdot c},
\qquad
\mathrm{fv}^{(j)}(\mathrm{pd}):=\sum_\nu\left|\mathrm{alc}^{(j)}_\nu(\mathrm{pd})\right| .
\]
The number $\mathrm{alc}^{(j)}_\nu(\mathrm{pd})$ depends only on the class of $\mathrm{pd}$ in $\mathrm{Gq}_n$, and
\[
\mathrm{alat}_\nu(c)^j=\sum_{\mathrm{pd}\in\mathrm{Gq}_n}\mathrm{alc}^{(j)}_\nu(\mathrm{pd})\,\mathrm{e}^{i\,\mathrm{pd}\cdot c}.
\]
The norms of $\ell^1(\mathrm{Gq}_n)$ and of $\ell^\infty(\mathrm{Gq}_n)$, with respect to the counting measure, are denoted $\|\cdot\|_1$ and $\|\cdot\|_\infty$;
the symbol $*$ denotes the convolution on $\mathrm{Gq}_n$.

\begin{named}{Lemma 2.5}
For all $j\ge1$,
\[
\left\|\mathrm{fv}^{(j)}\right\|_1\le\mathrm{Afr}_1^{\,j}
\qquad\text{and}\qquad
R^3\left\|\mathrm{fv}^{(j)}\right\|_\infty\le\mathrm{Afr}_1^{\,j-1}\cdot\mathrm{Afr}_0 .
\]
\end{named}

\begin{proof}
\begin{itemize}
\item[(1)]
\emph{The Fourier transform of $\mathrm{gs}_\nu$.}
Let $\mathrm{fs}$ be a Schwartz function and $y'\in\mathbb{R}^3$.
The transform of $\mathrm{fs}(\cdot+y')$ is $\mathrm{e}^{i\,\mathrm{kph}\cdot y'}\,\widehat{\mathrm{fs}}(\mathrm{kph})$.
Put $\mathrm{snc}(\mathrm{xq}):=\sin(\mathrm{xq})/\mathrm{xq}$ and $\mathrm{snc}(0):=1$.
Since
\[
\eta^{-1}\int_{-\frac{\eta}{2}}^{\frac{\eta}{2}}\mathrm{d}y'_{\mathfrak{l}}\,\mathrm{e}^{i\,\mathrm{kph}_{\mathfrak{l}}y'_{\mathfrak{l}}}=\mathrm{snc}\left(\frac{\eta\,\mathrm{kph}_{\mathfrak{l}}}{2}\right),
\]
the transform of $J_{\eta}\mathrm{fs}$ is $\widehat{\mathrm{fs}}(\mathrm{kph})\prod_{\mathfrak{l}}\mathrm{snc}\left(\frac{\eta\,\mathrm{kph}_{\mathfrak{l}}}{2}\right)$,
and the transform of $\delta_{\eta,\mu}\mathrm{fs}$ is
\[
\eta^{-1}\left(\mathrm{e}^{\frac{i\eta\,\mathrm{kph}_\mu}{2}}-\mathrm{e}^{-\frac{i\eta\,\mathrm{kph}_\mu}{2}}\right)\widehat{\mathrm{fs}}(\mathrm{kph})
=i\,\mathrm{kph}_\mu\,\mathrm{snc}\left(\frac{\eta\,\mathrm{kph}_\mu}{2}\right)\widehat{\mathrm{fs}}(\mathrm{kph}).
\]
Let $\mathrm{ao}_\nu(y):=\mathrm{gs}_\nu[\beta]\left(y+\frac{\eta}{2}e_\nu\right)$, a function in $C_c^\infty(\mathbb{R}^3)$.
Then
\[
\widehat{\mathrm{ao}}_\nu(\mathrm{kph})=-\mathrm{e}^{\frac{i\eta\,\mathrm{kph}_\nu}{2}}\sum_{\mu\ne\nu}i\,\mathrm{kph}_\mu\,\mathrm{snc}\left(\frac{\eta\,\mathrm{kph}_\mu}{2}\right)\prod_{\mathfrak{l}}\mathrm{snc}\left(\frac{\eta\,\mathrm{kph}_{\mathfrak{l}}}{2}\right)\hat\beta_{\mu\nu}(\mathrm{kph}),
\]
and, since $|\mathrm{snc}|\le1$,
\[
\sum_\nu\left|\widehat{\mathrm{ao}}_\nu(\mathrm{kph})\right|\le\mathrm{bq}(\mathrm{kph})\le\mathrm{bs}_4\left(1+|\mathrm{kph}|_2\right)^{-4}.
\]
\item[(2)]
\emph{Poisson summation.}
By [2.11] and the definition of $\mathrm{ait}_\nu$,
$\mathrm{alat}_\nu(c)=\sum_z\mathrm{ao}_\nu(\eta c+Rz)=\sum_z\mathrm{ao}_\nu\left(\eta(c+nz)\right)$.
For $\mathrm{pd}\in\frac{2\pi}{n}\mathbb{Z}^3$ one has $\mathrm{e}^{-i\,\mathrm{pd}\cdot(c+nz)}=\mathrm{e}^{-i\,\mathrm{pd}\cdot c}$.
Consequently
\begin{align*}
\mathrm{alc}^{(1)}_\nu(\mathrm{pd})
&=n^{-3}\sum_{c\in\mathbb{Z}^3}\mathrm{ao}_\nu(\eta c)\,\mathrm{e}^{-i\,\mathrm{pd}\cdot c}
=n^{-3}\sum_{y\in\eta\mathbb{Z}^3}\mathrm{ao}_\nu(y)\,\mathrm{e}^{-i\frac{\mathrm{pd}}{\eta}\cdot y}\\
&=n^{-3}\eta^{-3}\sum_{z\in\mathbb{Z}^3}\widehat{\mathrm{ao}}_\nu\left(\frac{\mathrm{pd}+2\pi z}{\eta}\right)
=R^{-3}\sum_{z\in\mathbb{Z}^3}\widehat{\mathrm{ao}}_\nu\left(\frac{\mathrm{pd}+2\pi z}{\eta}\right).
\end{align*}
The third equality is Poisson's summation formula
\[
\eta^3\sum_{y\in\eta\mathbb{Z}^3}\mathrm{fs}(y)=\sum_{\mathrm{kph}\in\frac{2\pi}{\eta}\mathbb{Z}^3}\widehat{\mathrm{fs}}(\mathrm{kph})
\]
\cite[Chapter~VII, Corollary~2.6]{SteinWeiss1971},
applied to the Schwartz function $\mathrm{fs}(y)=\mathrm{ao}_\nu(y)\,\mathrm{e}^{-i\frac{\mathrm{pd}}{\eta}\cdot y}$,
whose transform is $\widehat{\mathrm{ao}}_\nu\left(\cdot+\frac{\mathrm{pd}}{\eta}\right)$.
\item[(3)]
\emph{The norm in $\ell^1$.}
When $\mathrm{pd}$ ranges over a set of representatives of $\mathrm{Gq}_n$ and $z$ over $\mathbb{Z}^3$,
the map $(\mathrm{pd},z)\mapsto(\mathrm{pd}+2\pi z)/\eta$ is a bijection onto $\frac{2\pi}{n\eta}\mathbb{Z}^3=\mathrm{Kb}_R$.
By (1), (2) and Lemma~2.2 (applicable since $R\ge1$),
\[
\left\|\mathrm{fv}^{(1)}\right\|_1\le R^{-3}\sum_{\mathrm{kph}\in\mathrm{Kb}_R}\mathrm{bq}(\mathrm{kph})
\le\mathrm{bs}_4\cdot R^{-3}\sum_{\mathrm{kph}\in\mathrm{Kb}_R}\left(1+|\mathrm{kph}|_2\right)^{-4}
\le\mathrm{cs}_1\cdot\mathrm{bs}_4=\mathrm{Afr}_1 .
\]
\item[(4)]
\emph{The sup norm.}
Every class of $\mathrm{Gq}_n$ has a representative $\mathrm{pd}$ with $|\mathrm{pd}|_\infty\le\pi$.
For such a representative and for $z\ne0$,
\[
\left|\frac{\mathrm{pd}+2\pi z}{\eta}\right|_2\ge\frac{2\pi|z|_\infty-\pi}{\eta}\ge\frac{\pi|z|_\infty}{\eta}.
\]
Hence
\[
R^3\,\mathrm{fv}^{(1)}(\mathrm{pd})\le\mathrm{bq}\left(\frac{\mathrm{pd}}{\eta}\right)+\sum_{z\ne0}\mathrm{bs}_4\left(\frac{\eta}{\pi|z|_\infty}\right)^4
\le\mathrm{bs}_4\left(1+\frac{\eta^4\,\mathrm{zs}_4}{\pi^4}\right)\le2\,\mathrm{bs}_4=\mathrm{Afr}_0 .
\]
\item[(5)]
\emph{Powers.}
With sums taken in $\mathrm{Gq}_n$,
\[
\mathrm{alc}^{(j)}_\nu(\mathrm{pd})=\sum_{\mathrm{pd}_1+\dots+\mathrm{pd}_j=\mathrm{pd}}\ \prod_i\mathrm{alc}^{(1)}_\nu(\mathrm{pd}_i).
\]
All terms in the following inequalities are $\ge0$, so
\[
\mathrm{fv}^{(j)}\le\sum_\nu\left|\mathrm{alc}^{(1)}_\nu\right|^{*j}\le\left(\mathrm{fv}^{(1)}\right)^{*j},
\]
the $j$-fold convolution.
For functions $\mathrm{fn},\mathrm{fn}'\ge0$ on $\mathrm{Gq}_n$ one has $\|\mathrm{fn}*\mathrm{fn}'\|_1\le\|\mathrm{fn}\|_1\,\|\mathrm{fn}'\|_1$ and $\|\mathrm{fn}*\mathrm{fn}'\|_\infty\le\|\mathrm{fn}\|_1\,\|\mathrm{fn}'\|_\infty$;
together with (3) and (4) these give the two bounds.
\end{itemize}
\end{proof}

\paragraph{2.8. A cyclic inequality.}

\begin{named}{Lemma 2.6}
Let $\mathrm{lf}\ge1$, let $\mathrm{ufn}:\mathrm{Pb}_n\to[0,\infty)$ and $\mathrm{wfn}_1,\dots,\mathrm{wfn}_{\mathrm{lf}}:\mathrm{Gq}_n\to[0,\infty)$, and put $p_{\mathrm{lf}+1}:=p_1$.
Then
\[
\sum_{p_1,\dots,p_{\mathrm{lf}}\in\mathrm{Pb}_n}\ \prod_{i=1}^{\mathrm{lf}}\mathrm{ufn}(p_i)\,\mathrm{wfn}_i(p_i-p_{i+1})
\le\left(\sum_{p\in\mathrm{Pb}_n}\mathrm{ufn}(p)^{\mathrm{lf}}\right)\left\|\mathrm{wfn}_{\mathrm{lf}}\right\|_\infty\prod_{i<\mathrm{lf}}\left\|\mathrm{wfn}_i\right\|_1 .
\]
\end{named}

\begin{proof}
Identify $\mathrm{Pb}_n$ with $\frac{2\pi}{n}\left(\mathbb{Z}+\frac12\right)^3/2\pi\mathbb{Z}^3$.
With $\mathrm{pd}_i:=p_i-p_{i+1}$, the map $(p_1,\dots,p_{\mathrm{lf}})\mapsto(p_1,\mathrm{pd}_1,\dots,\mathrm{pd}_{\mathrm{lf}-1})$ is a bijection of $\mathrm{Pb}_n^{\mathrm{lf}}$ onto $\mathrm{Pb}_n\times\mathrm{Gq}_n^{\mathrm{lf}-1}$;
indeed $p_i=p_1-\mathrm{smq}_i$ with $\mathrm{smq}_i:=\mathrm{pd}_1+\dots+\mathrm{pd}_{i-1}$ (so $\mathrm{smq}_1=0$),
and $p_{\mathrm{lf}}-p_1=-\mathrm{smq}_{\mathrm{lf}}$.
The left side is therefore equal to
\[
\sum_{\mathrm{pd}_1,\dots,\mathrm{pd}_{\mathrm{lf}-1}}\ \prod_{i<\mathrm{lf}}\mathrm{wfn}_i(\mathrm{pd}_i)\cdot\mathrm{wfn}_{\mathrm{lf}}(-\mathrm{smq}_{\mathrm{lf}})\sum_{p_1}\prod_{i=1}^{\mathrm{lf}}\mathrm{ufn}(p_1-\mathrm{smq}_i).
\]
H\"older's inequality for $\mathrm{lf}$ factors, each with the exponent $\mathrm{lf}$, and the fact that $p_1\mapsto p_1-\mathrm{smq}_i$ is a bijection give
\[
\sum_{p_1}\prod_i\mathrm{ufn}(p_1-\mathrm{smq}_i)\le\prod_i\left(\sum_{p_1}\mathrm{ufn}(p_1-\mathrm{smq}_i)^{\mathrm{lf}}\right)^{\frac{1}{\mathrm{lf}}}=\sum_p\mathrm{ufn}(p)^{\mathrm{lf}} .
\]
Finally, $\mathrm{wfn}_{\mathrm{lf}}(-\mathrm{smq}_{\mathrm{lf}})\le\|\mathrm{wfn}_{\mathrm{lf}}\|_\infty$.
(For $\mathrm{lf}=1$ the assertion reads $\sum_p\mathrm{ufn}(p)\,\mathrm{wfn}_1(0)\le\sum_p\mathrm{ufn}(p)\,\|\mathrm{wfn}_1\|_\infty$.)
\end{proof}

\paragraph{2.9. Clause (c).}

\begin{proof}[Proof of Proposition~37.2\,(c)]
\emph{The blocks of the vertices.}
By 2.1, the block $(p,p')$ of the multiplication by $\mathrm{alat}_\nu^j$ is $\mathrm{alc}^{(j)}_\nu(p-p')$;
the block $(p,p')$ of the multiplication by $(\mathrm{alat}^-_\nu)^j$ is $\mathrm{e}^{-i(p-p')_\nu}\,\mathrm{alc}^{(j)}_\nu(p-p')$;
and the block $(p,p')$ of $\mathrm{Sh}^{\pm}_\nu$ is $\mathbf{1}[p=p']\,\mathrm{e}^{\pm ip_\nu}$.
Insert these blocks, together with
$(i\,\mathrm{bph}^+_\nu)^j=i^j\eta^j\,\mathrm{alat}_\nu^j$,
$(i\,\mathrm{bph}^-_\nu)^j=(-i)^j\eta^j\,(\mathrm{alat}^-_\nu)^j$
and $\mathrm{e}^{-i(p-p')_\nu}\,\mathrm{e}^{-ip'_\nu}=\mathrm{e}^{-ip_\nu}$, into [2.4].
The result is
\begin{align*}
\mathrm{Vd}_j(p,p')&=-\eta^j\sum_\nu\mathrm{alc}^{(j)}_\nu(p-p')\left[i^j\,\mathrm{e}^{ip'_\nu}\,\mathrm{Ehp}^+_\nu+(-i)^j\,\mathrm{e}^{-ip_\nu}\,\mathrm{Ehp}^-_\nu\right],\\
\left|\mathrm{Vd}_j(p,p')\right|_2&\le3(1+\rho)\,\eta^j\cdot\mathrm{fv}^{(j)}(p-p')\le6\,\eta^j\cdot\mathrm{fv}^{(j)}(p-p'),\tag*{[2.19]}
\end{align*}
where the bound comes from Proposition~37.1\,(d), applied to each term of the sum over $\nu$ (the coefficients have modulus one).

\emph{One trace.}
Let $\mathrm{lf}\in\{1,2,3\}$, and let $j_1+\dots+j_{\mathrm{lf}}=4$ with $j_i\ge1$; put $p_{\mathrm{lf}+1}:=p_1$.
Use 2.1 (blocks), $\mathrm{G0}=\mathrm{Mu}[\mathrm{Rsy}]$, the bound $|\operatorname{tr}O|\le192\,|O|_2$ on ${V_{\mathrm{v}}}$, [2.3] and [2.19];
then apply Lemma~2.6 with $\mathrm{ufn}(p)=\left(|p|_2+\mathbf{m}\right)^{-1}$ and $\mathrm{wfn}_i=\mathrm{fv}^{(j_i)}$, and Lemma~2.5.
This yields
\begin{align*}
\left|\operatorname{Tr}\left(\mathrm{Yv}_{j_1}\cdots\mathrm{Yv}_{j_{\mathrm{lf}}}\right)\right|
&=\left|\sum_{p_1,\dots,p_{\mathrm{lf}}}\operatorname{tr}\prod_{i=1}^{\mathrm{lf}}\left[\mathrm{Rsy}(p_i)\,\mathrm{Vd}_{j_i}(p_i,p_{i+1})\right]\right|\\
&\le192\,(6\,\mathrm{Cf})^{\mathrm{lf}}\cdot\eta^4\sum_{p_1,\dots,p_{\mathrm{lf}}}\ \prod_i\left(|p_i|_2+\mathbf{m}\right)^{-1}\mathrm{fv}^{(j_i)}(p_i-p_{i+1})\\
&\le192\,(6\,\mathrm{Cf})^{\mathrm{lf}}\cdot\eta^4\cdot n^3\cdot\mathrm{Il}_{\mathrm{lf}}\cdot\mathrm{Afr}_1^{\,j_{\mathrm{lf}}-1}\cdot\mathrm{Afr}_0\cdot R^{-3}\prod_{i<\mathrm{lf}}\mathrm{Afr}_1^{\,j_i}\\
&=192\,(6\,\mathrm{Cf})^{\mathrm{lf}}\cdot\mathrm{Afr}_1^3\cdot\mathrm{Afr}_0\cdot\eta\cdot\mathrm{Il}_{\mathrm{lf}},\tag*{[2.20]}
\end{align*}
since $\eta^4n^3R^{-3}=\eta$.
Now [2.6] and $6\,\mathrm{Cf}\ge1$ give
\begin{align*}
\left|\operatorname{Tr}\mathrm{Yv}_4\right|&\le192\,(6\,\mathrm{Cf})^3\cdot\mathrm{Afr}_1^3\cdot\mathrm{Afr}_0\cdot4\eta,\\
4\left|\operatorname{Tr}\left(\mathrm{Yv}_1\mathrm{Yv}_3\right)\right|&\le192\,(6\,\mathrm{Cf})^3\cdot\mathrm{Afr}_1^3\cdot\mathrm{Afr}_0\cdot8\eta,\\
3\left|\operatorname{Tr}\left(\mathrm{Yv}_2^2\right)\right|&\le192\,(6\,\mathrm{Cf})^3\cdot\mathrm{Afr}_1^3\cdot\mathrm{Afr}_0\cdot6\eta,\\
12\left|\operatorname{Tr}\left(\mathrm{Yv}_1^2\mathrm{Yv}_2\right)\right|&\le192\,(6\,\mathrm{Cf})^3\cdot\mathrm{Afr}_1^3\cdot\mathrm{Afr}_0\cdot12\eta\log\left(1+\frac{5}{\eta\bar{m}}\right).
\end{align*}
The sum of the four left sides is therefore at most
\[
192\cdot18\,(6\,\mathrm{Cf})^3\cdot\mathrm{Afr}_1^3\cdot\mathrm{Afr}_0\cdot\eta\left(1+\log\left(1+\frac{5}{\eta\bar{m}}\right)\right)
=C_{\mathrm{ct}}\cdot\eta\left(1+\log\left(1+\frac{5}{\eta\bar{m}}\right)\right).
\]
\end{proof}

\emph{A consequence.}
For $\mathrm{xq}>0$,
\[
\frac34\log(1+\mathrm{xq})=\log\left((1+\mathrm{xq})^{\frac34}\right)\le\log\left(1+\mathrm{xq}^{\frac34}\right)\le\mathrm{xq}^{\frac34}.
\]
With $\mathrm{xq}=5/(\eta\bar{m})$, and since $\eta\le\eta^{\frac14}$, it follows that
\[
C_{\mathrm{ct}}\cdot\eta\left(1+\log\left(1+\frac{5}{\eta\bar{m}}\right)\right)
\le C_{\mathrm{ct}}\left(\eta+\frac43\left(\frac{5}{\bar{m}}\right)^{\frac34}\eta^{\frac14}\right)
\le C_{\mathrm{ct}}\left(1+\frac43\left(\frac{5}{\bar{m}}\right)^{\frac34}\right)\eta^{\frac14}.\tag*{[2.21]}
\]

\paragraph{2.10. Clause (d).}

\begin{proof}[Proof of Proposition~37.2\,(d)]
By (b),
\begin{multline*}
b_{\mathrm{ch}}(\eta,R)-b_R(f)=\left[\operatorname{Tr}\mathrm{Yv}_4-4\operatorname{Tr}\left(\mathrm{Yv}_1\mathrm{Yv}_3\right)-3\operatorname{Tr}\left(\mathrm{Yv}_2^2\right)+12\operatorname{Tr}\left(\mathrm{Yv}_1^2\mathrm{Yv}_2\right)\right]\\
-6\left[\operatorname{Tr}\mathrm{Yv}_1^4-\operatorname{Tr}_{L^2_{\mathrm{ap}}(T_R;\mathbb{C}^4)}\left(S^R\,\mathrm{va}^R\right)^4\right].\tag*{[2.22]}
\end{multline*}
The first bracket is bounded by (c) and [2.21].
The second bracket is bounded in five steps, which are preceded by a bound on $\mathrm{Yv}_1$.

\emph{The norm of $\mathrm{Yv}_1$.}
By [2.3], by $\|\mathrm{Mx}\|_4^4\le192\,|\mathrm{Mx}|_2^4$ on ${V_{\mathrm{v}}}$, by $(\mathrm{xq}+\mathrm{xq}')^4\le8\left(\mathrm{xq}^4+\mathrm{xq}'^4\right)$ for $\mathrm{xq},\mathrm{xq}'\ge0$, and by [2.5],
\[
\mathrm{nrm}_4(\mathrm{Rsy})^4\le192\cdot8\,n^{-3}\sum_{p\in\mathrm{Pb}_n}\left(C_E^4+\left(\mathbf{m}^2+|p|_2^2\right)^{-2}\right)\le1536\left(C_E^4+\frac{1}{\pi\mathbf{m}}\right).
\]
By [2.4], $\mathrm{Yv}_1$ is minus the sum of the six operators $\left(\mathrm{G0}\,i\,\mathrm{bph}^{\pm}_\nu\right)\left(\mathrm{Ehp}^{\pm}_\nu\mathrm{Sh}^{\pm}_\nu\right)$.
By (Tr1), (Tr5), [2.17] and $\mathbf{m}=\bar{m}\eta$,
\[
\left\|\mathrm{Yv}_1\right\|_4\le6\cdot6\left(1536\left(C_E^4+\frac{1}{\pi\mathbf{m}}\right)\right)^{\frac14}\eta^{\frac14}\cdot\mathrm{kpr}
=36\left(1536\left(C_E^4\cdot\eta+\frac{1}{\pi\bar{m}}\right)\right)^{\frac14}\cdot\mathrm{kpr}\le\mathrm{kn}_1.\tag*{[2.23]}
\]

\emph{Step 1 (the heavy components).}
For $p\in\mathrm{Pb}_n$ define
\[
\mathrm{cut}(p):=\mathbf{1}\left[|p|_\infty\le\frac{\pi}{2}\right],\qquad
\mathrm{Sct}(p):=\mathrm{cut}(p)\,\mathrm{Sdr}(p,\mathbf{m}),\qquad
\mathrm{Erc}(p):=\mathrm{Rsy}(p)-J\,\mathrm{Sct}(p)\,J^{*};
\]
here $\mathrm{Sct}(p)$ is a $4\times4$ matrix.
By Proposition~37.1\,(f), $\mathrm{Erc}(p)=\mathrm{Er}(p,\mathbf{m})+\left(1-\mathrm{cut}(p)\right)J\,\mathrm{Sdr}(p,\mathbf{m})\,J^{*}$.
At the points where $\mathrm{cut}(p)=0$ one has $|\mathrm{Sdr}(p,\mathbf{m})|_2\le1/|p|_2\le1/|p|_\infty<2/\pi$.
Therefore $|\mathrm{Erc}(p)|_2\le C_E+1$, and $\mathrm{nrm}_4(\mathrm{Erc})\le192^{\frac14}(C_E+1)$.
For a $4\times4$ matrix $\mathrm{Mx}$ one has $\|\mathrm{Sct}(p)\,\mathrm{Mx}\|_4^4\le4\,|\mathrm{Mx}|_2^4\left(|p|_2^2+\mathbf{m}^2\right)^{-2}$.
By [2.5] and (Tr5) it follows that, for every $n\mathbb{Z}^3$-periodic $\mathrm{fs}$ with $\|\mathrm{fs}\|_{\ell^4}\le\eta^{\frac14}\cdot\mathrm{kpr}$,
\[
\mathrm{nrm}_4(\mathrm{Sct}\,\mathrm{Mx})\le|\mathrm{Mx}|_2\left(\frac{4}{\pi\mathbf{m}}\right)^{\frac14},\qquad
\left\|\mathrm{Mu}[\mathrm{Sct}\,\mathrm{Mx}]\,\mathrm{fs}\right\|_4\le|\mathrm{Mx}|_2\left(\frac{4}{\pi\bar{m}}\right)^{\frac14}\cdot\mathrm{kpr}=|\mathrm{Mx}|_2\,\frac{\mathrm{kn}_0}{3}.\tag*{[2.24]}
\]
Let
\[
\mathrm{Xl}:=J\,\mathrm{Mu}[\mathrm{Sct}]\,J^{*}\,\mathrm{Vd}_1=\mathrm{Mu}\left[J\,\mathrm{Sct}\,J^{*}\right]\mathrm{Vd}_1
\]
(the equality by 2.1).
The fibre map $J$ is an isometry of $\mathbb{C}^4$ into ${V_{\mathrm{v}}}$ with the adjoint $J^{*}$ ((w5));
hence the non-zero singular values of $J\,\mathrm{Sct}(p)\,J^{*}$ are those of $\mathrm{Sct}(p)$,
and [2.24] holds for $\mathrm{Mu}\left[J\,\mathrm{Sct}\,J^{*}\right]\mathrm{fs}$ with $\mathrm{Mx}=1$.
By [2.4], $\mathrm{Xl}$ is minus the sum of the six operators $\left(\mathrm{Mu}\left[J\,\mathrm{Sct}\,J^{*}\right]i\,\mathrm{bph}^{\pm}_\nu\right)\left(\mathrm{Ehp}^{\pm}_\nu\mathrm{Sh}^{\pm}_\nu\right)$;
by [2.17], [2.24] and $\left|\mathrm{Ehp}^{\pm}_\nu\mathrm{Sh}^{\pm}_\nu\right|_2\le6$,
\[
\|\mathrm{Xl}\|_4\le6\cdot\frac{\mathrm{kn}_0}{3}\cdot6=12\,\mathrm{kn}_0 .
\]
The difference $\mathrm{Yv}_1-\mathrm{Xl}=\mathrm{Mu}[\mathrm{Erc}]\,\mathrm{Vd}_1$ is minus the sum of the six operators $\left(\mathrm{Mu}[\mathrm{Erc}]\,i\,\mathrm{bph}^{\pm}_\nu\right)\left(\mathrm{Ehp}^{\pm}_\nu\mathrm{Sh}^{\pm}_\nu\right)$,
and (Tr5), [2.17] give
\[
\left\|\mathrm{Yv}_1-\mathrm{Xl}\right\|_4\le6\cdot192^{\frac14}(C_E+1)\,\eta^{\frac14}\cdot\mathrm{kpr}\cdot6=36\cdot192^{\frac14}(C_E+1)\,\mathrm{kpr}\cdot\eta^{\frac14}.
\]
By (Tr2), therefore,
\[
\left|\operatorname{Tr}\mathrm{Yv}_1^4-\operatorname{Tr}\mathrm{Xl}^4\right|\le\mathrm{lnk}_1\cdot\eta^{\frac14}.
\]
Since the spaces are finite-dimensional, the trace is cyclic, and
\[
\operatorname{Tr}_{\mathrm{Sfc}}\mathrm{Xl}^4=\operatorname{Tr}_{\mathrm{Sfc}}\left(J\cdot\mathrm{Mu}[\mathrm{Sct}]\,J^{*}\,\mathrm{Vd}_1\right)^4=\operatorname{Tr}_{\mathrm{Sfl}}\left(\mathrm{Mu}[\mathrm{Sct}]\,\mathrm{vw}_1\right)^4,
\qquad\mathrm{vw}_1:=J^{*}\,\mathrm{Vd}_1\,J.
\]
By 2.1, $J$ and $J^{*}$ commute with multiplication operators and with shifts,
and $J^{*}\,\mathrm{Ehp}^{\pm}_\nu\,J$ is the constant multiplier with the matrix $\frac12(\rho\mp\gamma_\nu)$
(Proposition~37.1\,(d): $8$ of the $48$ labels contribute, with the weight $3\cdot8/48$).
Consequently [2.4] gives
\begin{align*}
\mathrm{vw}_1&=-\frac{i}{2}\sum_\nu\left[\mathrm{bph}^+_\nu(\rho-\gamma_\nu)\,\mathrm{Sh}^+_\nu+\mathrm{bph}^-_\nu(\rho+\gamma_\nu)\,\mathrm{Sh}^-_\nu\right]\\
&=-\frac{i\eta}{2}\sum_\nu\left[\mathrm{alat}_\nu(\rho-\gamma_\nu)\,\mathrm{Sh}^+_\nu-\mathrm{alat}^-_\nu(\rho+\gamma_\nu)\,\mathrm{Sh}^-_\nu\right].\tag*{[2.25]}
\end{align*}

\emph{Step 2 (the vertex).}
As operators on $\mathrm{Sfl}$ one has
$\mathrm{alat}_\nu\,\mathrm{Sh}^+_\nu=\mathrm{Sh}^+_\nu\,\mathrm{alat}^-_\nu$ and $\mathrm{alat}^-_\nu\,\mathrm{Sh}^-_\nu=\mathrm{Sh}^-_\nu\,\mathrm{alat}_\nu$,
because
\[
\left(\mathrm{Sh}^+_\nu(\mathrm{alat}^-_\nu\xi)\right)(c)=\mathrm{alat}^-_\nu(c+e_\nu)\,\xi(c+e_\nu)=\mathrm{alat}_\nu(c)\,\xi(c+e_\nu),\qquad
\left(\mathrm{Sh}^-_\nu(\mathrm{alat}_\nu\xi)\right)(c)=\mathrm{alat}_\nu(c-e_\nu)\,\xi(c-e_\nu).
\]
Hence
\begin{align*}
\mathrm{vw}_1&=-\frac{i\eta}{2}\sum_\nu\left[(\rho-\gamma_\nu)\,\mathrm{Sh}^+_\nu\,\mathrm{alat}^-_\nu-(\rho+\gamma_\nu)\,\mathrm{Sh}^-_\nu\,\mathrm{alat}_\nu\right]
=i\eta\,\gamma\cdot\mathrm{alat}+\mathrm{rv},\qquad\gamma\cdot\mathrm{alat}:=\sum_\nu\gamma_\nu\,\mathrm{alat}_\nu,\tag*{[2.26]}\\
\mathrm{rv}&:=-\frac{i\eta}{2}\sum_\nu\left[(\rho-\gamma_\nu)\left(\mathrm{Sh}^+_\nu-1\right)\mathrm{alat}^-_\nu-(\rho+\gamma_\nu)\left(\mathrm{Sh}^-_\nu-1\right)\mathrm{alat}_\nu+(\rho-\gamma_\nu)\left(\mathrm{alat}^-_\nu-\mathrm{alat}_\nu\right)\right];
\end{align*}
indeed, the terms that contain no factor $\mathrm{Sh}-1$ are
\[
-\frac{i\eta}{2}\left((\rho-\gamma_\nu)\,\mathrm{alat}^-_\nu-(\rho+\gamma_\nu)\,\mathrm{alat}_\nu\right)
=-\frac{i\eta}{2}(\rho-\gamma_\nu)\left(\mathrm{alat}^-_\nu-\mathrm{alat}_\nu\right)+i\eta\,\gamma_\nu\,\mathrm{alat}_\nu .
\]
Next, $\mathrm{Mu}[\mathrm{Sct}](\rho\mp\gamma_\nu)\left(\mathrm{Sh}^{\pm}_\nu-1\right)=\mathrm{Mu}\left[\mathrm{hm}^{\pm}_\nu\right]$ with
\[
\mathrm{hm}^{\pm}_\nu(p):=\mathrm{Sct}(p)(\rho\mp\gamma_\nu)\left(\mathrm{e}^{\pm ip_\nu}-1\right),\qquad
\left|\mathrm{hm}^{\pm}_\nu(p)\right|_2\le(1+\rho)\,|p_\nu|\left(|p|_2^2+\mathbf{m}^2\right)^{-\frac12}\le2;
\]
so $\mathrm{nrm}_4\left(\mathrm{hm}^{\pm}_\nu\right)\le\left(4\cdot2^4\right)^{\frac14}=64^{\frac14}$.
By (Tr5), [2.17] and the first bound of [2.24] (with $|\rho-\gamma_\nu|_2\le2$),
\begin{align*}
\left\|\mathrm{Mu}[\mathrm{Sct}]\,\mathrm{rv}\right\|_4
&\le\frac12\sum_\nu\left[2\cdot64^{\frac14}\,\eta^{\frac14}\cdot\mathrm{kpr}+2\left(\frac{4}{\pi\mathbf{m}}\right)^{\frac14}\left(2\eta^5\cdot\mathrm{Ab}_2^4\cdot\mathrm{Vbe}\right)^{\frac14}\right]\\
&=\frac32\left[2\cdot64^{\frac14}\,\mathrm{kpr}\cdot\eta^{\frac14}+\left(\frac{128}{\pi\bar{m}}\right)^{\frac14}\cdot\mathrm{Ab}_2\cdot\mathrm{Vbe}^{\frac14}\cdot\eta\right]\le\mathrm{rn}_0\cdot\eta^{\frac14}.
\end{align*}
By the first expression for $\mathrm{vw}_1$ in [2.26], $\mathrm{Mu}[\mathrm{Sct}]\,\mathrm{vw}_1$ is $-i/2$ times a sum of six operators
$\mathrm{Mu}\left[\mathrm{Sct}\,(\rho\mp\gamma_\nu)\,\mathrm{e}^{\pm ip_\nu}\right]\mathrm{fs}$ with $\mathrm{fs}=\eta\,\mathrm{alat}^-_\nu$ or $\mathrm{fs}=-\eta\,\mathrm{alat}_\nu$.
The factor $\mathrm{e}^{\pm ip_\nu}$ has modulus one;
so (Tr5), [2.17] and [2.24] bound the norm $\|\cdot\|_4$ of each of these six operators by $2\,\mathrm{kn}_0/3$,
and $\left\|\mathrm{Mu}[\mathrm{Sct}]\,\mathrm{vw}_1\right\|_4\le2\,\mathrm{kn}_0$.
The operator $\mathrm{Mu}[\mathrm{Sct}]\,i\eta\,\gamma\cdot\mathrm{alat}=\sum_\nu\mathrm{Mu}\left[\mathrm{Sct}\,i\gamma_\nu\right](\eta\,\mathrm{alat}_\nu)$ is a sum of three operators,
each of which has the norm $\|\cdot\|_4\le\mathrm{kn}_0/3$ by [2.17] and [2.24];
so $\left\|\mathrm{Mu}[\mathrm{Sct}]\,i\eta\,\gamma\cdot\mathrm{alat}\right\|_4\le\mathrm{kn}_0$.
Now (Tr2) gives
\[
\left|\operatorname{Tr}\left(\mathrm{Mu}[\mathrm{Sct}]\,\mathrm{vw}_1\right)^4-\operatorname{Tr}\left(\mathrm{Mu}[\mathrm{Sct}]\,i\eta\,\gamma\cdot\mathrm{alat}\right)^4\right|\le4\left(2\,\mathrm{kn}_0\right)^3\cdot\mathrm{rn}_0\cdot\eta^{\frac14}=\mathrm{lnk}_2\cdot\eta^{\frac14}.
\]

\emph{Step 3 (from the lattice to the continuum torus).}
Since $2\pi/(n\eta)=2\pi/R$, the map $p\mapsto P:=p/\eta$ is a bijection of $\mathrm{Pb}_n$ onto $\mathrm{Kf}_R\cap(-\pi/\eta,\pi/\eta]^3$.
Let
\[
\mathrm{Kf}_R^{<}:=\left\{P\in\mathrm{Kf}_R:\ |P|_\infty\le\frac{\pi}{2\eta}\right\};
\]
this set has at most $n^3$ elements, and $\mathrm{cut}(p)=1$ if and only if $p/\eta\in\mathrm{Kf}_R^{<}$.
Moreover $\eta\,\mathrm{Sct}(p)=\mathrm{cut}(p)\,\mathrm{Sdr}(p/\eta,\bar{m})$.
The operator $\mathrm{Aq}:=\mathrm{Mu}[\mathrm{Sct}]\,i\eta\,\gamma\cdot\mathrm{alat}$ satisfies $\mathrm{Aq}=\mathrm{Mu}[\mathrm{cut}]\,\mathrm{Aq}$,
and $\mathrm{Mu}[\mathrm{cut}]$ is the orthogonal projection of $\mathrm{Sfl}$ onto the span of the vectors $\mathrm{ew}_p\otimes\mathrm{bv}_{\mathrm{be}}$ with $\mathrm{cut}(p)=1$.
Computing in this orthonormal system, one obtains from (Tr3)
\[
\operatorname{Tr}_{\mathrm{Sfl}}\left(\mathrm{Mu}[\mathrm{Sct}]\,i\eta\,\gamma\cdot\mathrm{alat}\right)^4=\operatorname{Tr}\mathrm{Tq}_{\mathrm{lat}}^4,\qquad
\left\|\mathrm{Tq}_{\mathrm{lat}}\right\|_4\le\mathrm{kn}_0,
\]
where $\mathrm{Tq}_{\mathrm{lat}}$ is the matrix of $4\times4$ blocks
\[
\mathrm{Tq}_{\mathrm{lat}}(P,P'):=\mathrm{Sdr}(P,\bar{m})\,i\sum_\nu\gamma_\nu\left(\mathrm{ew}_{\eta P},\mathrm{alat}_\nu\,\mathrm{ew}_{\eta P'}\right)\qquad\left(P,P'\in\mathrm{Kf}_R^{<}\right).
\]

\emph{Matrix elements.}
The function $\mathrm{ait}_\nu$ is smooth and $R\mathbb{Z}^3$-periodic.
Four integrations by parts in the variable $y_j$ give
$|\mathrm{kph}_j|^4\left|\mathrm{fc}[\mathrm{ait}_\nu](\mathrm{kph})\right|\le R^{-3}\left\|\partial_j^4\mathrm{ait}_\nu\right\|_{L^1(\text{period cube})}$.
Choosing $j$ with $|\mathrm{kph}_j|=|\mathrm{kph}|_\infty$ and using [2.12], one obtains
\[
\left|\mathrm{fc}[\mathrm{ait}_\nu](\mathrm{kph})\right|\le R^{-3}\cdot\mathrm{Ab}_5\cdot|\mathrm{kph}|_\infty^{-4}\qquad\left(\mathrm{kph}\in\mathrm{Kb}_R,\ \mathrm{kph}\ne0\right).\tag*{[2.27]}
\]
Since $\sum_{\mathrm{kph}\in\mathrm{Kb}_R,\,\mathrm{kph}\ne0}|\mathrm{kph}|_\infty^{-4}=\left(\frac{R}{2\pi}\right)^4\mathrm{zs}_4$,
the series $\sum_{\mathrm{kph}\in\mathrm{Kb}_R}\mathrm{fc}[\mathrm{ait}_\nu](\mathrm{kph})\,\mathrm{e}^{i\,\mathrm{kph}\cdot y}$ converges absolutely and uniformly in $y$.
Its sum is a continuous $R\mathbb{Z}^3$-periodic function, and integration term by term shows that its coefficients are the numbers $\mathrm{fc}[\mathrm{ait}_\nu](\mathrm{kph})$.
By the completeness of the exponentials ((Tr6)), a function in $L^2$ of a period cube all of whose coefficients $\mathrm{fc}[\,\cdot\,](\mathrm{kph})$ vanish is $0$ almost everywhere.
Hence the sum of the series coincides almost everywhere with the continuous function $\mathrm{ait}_\nu$, and therefore everywhere:
\[
\mathrm{ait}_\nu(y)=\sum_{\mathrm{kph}\in\mathrm{Kb}_R}\mathrm{fc}[\mathrm{ait}_\nu](\mathrm{kph})\,\mathrm{e}^{i\,\mathrm{kph}\cdot y}.
\]
For $p=\eta P$ and $p'=\eta P'$ in $\mathrm{Pb}_n$, [2.11] gives
\[
\left(\mathrm{ew}_p,\mathrm{alat}_\nu\,\mathrm{ew}_{p'}\right)=n^{-3}\sum_{c\in\{0,\dots,n-1\}^3}\mathrm{ait}_\nu(\eta c)\,\mathrm{e}^{-i(p-p')\cdot c}
=\sum_{\mathrm{kph}\in\mathrm{Kb}_R}\mathrm{fc}[\mathrm{ait}_\nu](\mathrm{kph})\,n^{-3}\sum_c\mathrm{e}^{i(\eta\,\mathrm{kph}-p+p')\cdot c}.
\]
The vector $\eta\,\mathrm{kph}-p+p'$ lies in $\frac{2\pi}{n}\mathbb{Z}^3$;
so the inner mean is $1$ if $\eta\,\mathrm{kph}-p+p'\in2\pi\mathbb{Z}^3$, and $0$ otherwise.
Thus
\[
\left(\mathrm{ew}_{\eta P},\mathrm{alat}_\nu\,\mathrm{ew}_{\eta P'}\right)=\mathrm{fc}[\mathrm{ait}_\nu](P-P')+\mathrm{ali}_\nu(P-P'),\qquad
\mathrm{ali}_\nu(\mathrm{kph}):=\sum_{z\in\mathbb{Z}^3,\,z\ne0}\mathrm{fc}[\mathrm{ait}_\nu]\left(\mathrm{kph}+\frac{2\pi z}{\eta}\right).\tag*{[2.28]}
\]
For $P,P'\in\mathrm{Kf}_R^{<}$ the vector $\mathrm{kph}=P-P'$ satisfies $|\mathrm{kph}|_\infty\le\pi/\eta$;
so, for $z\ne0$,
\[
\left|\mathrm{kph}+\frac{2\pi z}{\eta}\right|_\infty\ge\frac{2\pi|z|_\infty-\pi}{\eta}\ge\frac{\pi|z|_\infty}{\eta},
\]
and [2.27] yields
\[
\left|\mathrm{ali}_\nu(P-P')\right|\le R^{-3}\cdot\mathrm{Ab}_5\cdot\mathrm{zs}_4\left(\frac{\eta}{\pi}\right)^4.\tag*{[2.29]}
\]
Define
\[
\mathrm{Tq}_{\mathrm{it}}(P,P'):=\mathrm{Sdr}(P,\bar{m})\,i\sum_\nu\gamma_\nu\,\mathrm{fc}[\mathrm{ait}_\nu](P-P'),\qquad P,P'\in\mathrm{Kf}_R^{<}.
\]
By (Tr6), this is the matrix, in the orthonormal system $\mathrm{Ew}_P\,\mathrm{bv}_{\mathrm{be}}$, $P\in\mathrm{Kf}_R^{<}$,
of the compression to the span of this system of the operator
$\mathrm{Mu}\left[\mathbf{1}\left[P\in\mathrm{Kf}_R^{<}\right]\mathrm{Sdr}(P,\bar{m})\right]i\gamma\cdot\mathrm{ait}$ on $L^2_{\mathrm{ap}}(T_R;\mathbb{C}^4)$,
where $\gamma\cdot\mathrm{ait}:=\sum_\nu\gamma_\nu\,\mathrm{ait}_\nu$.
By (Tr1), (Tr6), [2.7] and [2.16],
\[
\left\|\mathrm{Tq}_{\mathrm{it}}\right\|_4\le3\left(\frac{4}{\pi\bar{m}}\right)^{\frac14}\cdot\mathrm{kpr}=\mathrm{kn}_0 .
\]
Use now $|\mathrm{Sdr}(P,\bar{m})|_2^2\le|P|_2^{-2}$, the identity $\|\gamma\cdot y\|_2^2=4|y|_2^2$ for $y\in\mathbb{C}^3$, [2.29], [2.8],
and $\#\left|\mathrm{Kf}_R^{<}\right|\le n^3=R^3\eta^{-3}$ for the number of the points $P'$.
They give
\begin{align*}
\left\|\mathrm{Tq}_{\mathrm{lat}}-\mathrm{Tq}_{\mathrm{it}}\right\|_2^2
&=\sum_{P,P'\in\mathrm{Kf}_R^{<}}\left\|\mathrm{Sdr}(P,\bar{m})\,\gamma\cdot\mathrm{ali}(P-P')\right\|_2^2\\
&\le\left[\frac{2\sqrt{3}\,R^3}{\pi\eta}\right]\left[R^3\eta^{-3}\right]4\cdot3\left[R^{-3}\cdot\mathrm{Ab}_5\cdot\mathrm{zs}_4\left(\frac{\eta}{\pi}\right)^4\right]^2
=\frac{24\sqrt{3}}{\pi^9}\,\mathrm{zs}_4^2\cdot\mathrm{Ab}_5^2\cdot\eta^4,
\end{align*}
where $\mathrm{ali}:=(\mathrm{ali}_1,\mathrm{ali}_2,\mathrm{ali}_3)$; the powers of $R$ cancel.
By (Tr2) and $\|\cdot\|_4\le\|\cdot\|_2$,
\[
\left|\operatorname{Tr}\mathrm{Tq}_{\mathrm{lat}}^4-\operatorname{Tr}\mathrm{Tq}_{\mathrm{it}}^4\right|\le4\,\mathrm{kn}_0^3\left(\frac{24\sqrt{3}}{\pi^9}\right)^{\frac12}\cdot\mathrm{zs}_4\cdot\mathrm{Ab}_5\cdot\eta^2=\mathrm{lnk}_3\cdot\eta^2.
\]

\emph{Step 4 ($\mathrm{ait}$ against $a^R$).}
On $L^2_{\mathrm{ap}}(T_R;\mathbb{C}^4)$ the symbol $\mathrm{cut}$ also denotes the multiplier $\mathrm{Mu}\left[\mathbf{1}\left[P\in\mathrm{Kf}_R^{<}\right]\right]$, which is an orthogonal projection.
Put
\[
\mathrm{Tq}_1:=\mathrm{cut}\,S^R\,\mathrm{va}^R,\qquad\mathrm{Tq}_0:=S^R\,\mathrm{va}^R,
\]
and let $\mathrm{Tq}_a$ be the compression of $\mathrm{Tq}_1$ to $\mathrm{Ran}\,\mathrm{cut}$.
By Lemma~2.4 and (Tr1), $\|\mathrm{Tq}_a\|_4\le\|\mathrm{Tq}_1\|_4\le\|\mathrm{Tq}_0\|_4\le\mathrm{kn}_0$.
The difference $\mathrm{Tq}_{\mathrm{it}}-\mathrm{Tq}_a$ is the compression of $\mathrm{cut}\,S^R\,i\gamma\cdot\left(\mathrm{ait}-a^R\right)$;
by Lemma~2.4 (with $\mathrm{Mx}=\gamma_\nu$) and [2.16],
\[
\left\|\mathrm{Tq}_{\mathrm{it}}-\mathrm{Tq}_a\right\|_4\le3\left(\frac{4}{\pi\bar{m}}\right)^{\frac14}\cdot3\eta\cdot\mathrm{Ab}_2\cdot\mathrm{Vbe}^{\frac14}.
\]
Then (Tr2) gives
\[
\left|\operatorname{Tr}\mathrm{Tq}_{\mathrm{it}}^4-\operatorname{Tr}\mathrm{Tq}_a^4\right|\le4\,\mathrm{kn}_0^3\cdot9\left(\frac{4}{\pi\bar{m}}\right)^{\frac14}\cdot\mathrm{Ab}_2\cdot\mathrm{Vbe}^{\frac14}\cdot\eta=\mathrm{lnk}_4\cdot\eta .
\]
Since $\mathrm{Tq}_1=\mathrm{cut}\,\mathrm{Tq}_1$, (Tr3) yields $\operatorname{Tr}\mathrm{Tq}_a^4=\operatorname{Tr}\mathrm{Tq}_1^4$.

\emph{Step 5 (removal of the cut).}
One has
\[
\mathrm{Tq}_0-\mathrm{Tq}_1=\sum_\nu\mathrm{Mu}\left[\mathbf{1}\left[P\notin\mathrm{Kf}_R^{<}\right]\mathrm{Sdr}(P,\bar{m})\,i\gamma_\nu\right]a^R_\nu,
\qquad
\left\|\mathrm{Sdr}(P,\bar{m})\,\gamma_\nu\right\|_4^4=4\left(|P|_2^2+\bar{m}^2\right)^{-2}\le4\,|P|_2^{-4}.
\]
By (Tr6), [2.9] and [2.16],
\[
\left\|\mathrm{Tq}_0-\mathrm{Tq}_1\right\|_4\le3\left(\frac{4\cdot32\,\eta}{3\pi^3}\right)^{\frac14}\cdot\mathrm{kpr},
\]
and (Tr2) gives
\[
\left|\operatorname{Tr}\mathrm{Tq}_1^4-\operatorname{Tr}\mathrm{Tq}_0^4\right|\le12\,\mathrm{kn}_0^3\left(\frac{128}{3\pi^3}\right)^{\frac14}\cdot\mathrm{kpr}\cdot\eta^{\frac14}=\mathrm{lnk}_5\cdot\eta^{\frac14}.
\]

\emph{Summation.}
Steps 1--5, together with $\eta^2\le\eta^{\frac14}$ and $\eta\le\eta^{\frac14}$, give
\[
\left|\operatorname{Tr}\mathrm{Yv}_1^4-\operatorname{Tr}_{L^2_{\mathrm{ap}}(T_R;\mathbb{C}^4)}\left(S^R\,\mathrm{va}^R\right)^4\right|\le\left(\mathrm{lnk}_1+\mathrm{lnk}_2+\mathrm{lnk}_3+\mathrm{lnk}_4+\mathrm{lnk}_5\right)\eta^{\frac14}.\tag*{[2.30]}
\]
Together with [2.22], (c), [2.21] and [2.1] this yields $\left|b_{\mathrm{ch}}(\eta,R)-b_R(f)\right|\le C_{1}^{\mathrm{N2}}\cdot\eta^{\frac14}$.
\end{proof}

\paragraph{2.11. Clause (e).}

\begin{proof}[Proof of Proposition~37.2\,(e)]
\emph{Localisation.}
Let $\mathrm{inj}_\infty:L^2(Q_{\mathrm{be}};\mathbb{C}^4)\to L^2(\mathbb{R}^3;\mathbb{C}^4)$ be the extension by zero.
Let $\mathrm{inj}_R:L^2(Q_{\mathrm{be}};\mathbb{C}^4)\to L^2_{\mathrm{ap}}(T_R;\mathbb{C}^4)$ be given by
\[
\left(\mathrm{inj}_R\,\psi_0\right)(y+Rz):=(-1)^{z_1+z_2+z_3}\,\psi_0(y)\quad\left(y\in Q_{\mathrm{be}},\ z\in\mathbb{Z}^3\right),
\qquad
\mathrm{inj}_R\,\psi_0:=0\ \text{ outside }Q_{\mathrm{be}}+R\mathbb{Z}^3;
\]
the definition is consistent because the sets $Q_{\mathrm{be}}+Rz$ are pairwise disjoint.
Both maps are isometries:
the modulus of $\mathrm{inj}_R\,\psi_0$ is periodic, and a period cube meets $Q_{\mathrm{be}}+R\mathbb{Z}^3$ in a set that is congruent to $Q_{\mathrm{be}}$ modulo $R\mathbb{Z}^3$.
Let $\mathrm{va}_Q$ denote the multiplication by $i\gamma\cdot a$ on $L^2(Q_{\mathrm{be}};\mathbb{C}^4)$.
As $a$ vanishes outside $Q_{\mathrm{be}}$, one has
$\mathrm{va}=\mathrm{inj}_\infty\,\mathrm{va}_Q\,\mathrm{inj}_\infty^{*}$ and $\mathrm{va}^R=\mathrm{inj}_R\,\mathrm{va}_Q\,\mathrm{inj}_R^{*}$.
Define the operators
\[
\mathrm{Tq}^{Q}_{\infty}:=\mathrm{inj}_\infty^{*}\,S_{\bar{m}}\,\mathrm{inj}_\infty\,\mathrm{va}_Q,\qquad
\mathrm{Tq}^{Q}_{R}:=\mathrm{inj}_R^{*}\,S^R\,\mathrm{inj}_R\,\mathrm{va}_Q
\]
on $L^2(Q_{\mathrm{be}};\mathbb{C}^4)$.
Let $(\mathrm{Sz},\mathrm{inj})$ be one of the pairs $(S_{\bar{m}},\mathrm{inj}_\infty)$ and $(S^R,\mathrm{inj}_R)$, and let $\mathrm{Prq}:=\mathrm{inj}\,\mathrm{inj}^{*}$.
Because $\mathrm{inj}^{*}\,\mathrm{inj}=1$, the operator $\mathrm{Aq}:=\mathrm{Sz}\,\mathrm{inj}\,\mathrm{va}_Q\,\mathrm{inj}^{*}$ satisfies
$\mathrm{Aq}=\mathrm{Aq}\,\mathrm{Prq}$ and $\mathrm{Prq}\,\mathrm{Aq}\,\mathrm{Prq}=\mathrm{inj}\left(\mathrm{inj}^{*}\,\mathrm{Sz}\,\mathrm{inj}\,\mathrm{va}_Q\right)\mathrm{inj}^{*}$.
From (a), from (Tr3), and from the fact that $\mathrm{inj}$ is an isometry it follows that
\[
\operatorname{Tr}\left(S_{\bar{m}}\,\mathrm{va}\right)^4=\operatorname{Tr}\left(\mathrm{Tq}^{Q}_{\infty}\right)^4,\qquad
\operatorname{Tr}\left(S^R\,\mathrm{va}^R\right)^4=\operatorname{Tr}\left(\mathrm{Tq}^{Q}_{R}\right)^4,\qquad
\left\|\mathrm{Tq}^{Q}_{\infty}\right\|_4\le\left\|S_{\bar{m}}\,\mathrm{va}\right\|_4\le\mathrm{kn}_2,\tag*{[2.31]}
\]
the last inequality by Lemma~2.3.

\emph{Images.}
For $\psi\in L^2_{\mathrm{ap}}(T_R;\mathbb{C}^4)$ let
\[
(\mathrm{Sx}*\psi)(y)=\int_{\mathbb{R}^3}\mathrm{d}y''\,\mathrm{Sx}(y'')\,\psi(y-y'').
\]
Apply (K-0) with $\mathrm{Hq}=L^2_{\mathrm{ap}}(T_R;\mathbb{C}^4)$ and $\mathrm{krn}=\mathrm{Sx}$; here $\|\mathrm{Sx}\|_{L^1}\le3/\bar{m}$ by [2.18].
It follows that the integral converges absolutely for almost every $y$, that $\mathrm{Sx}*\psi$ lies in $L^2_{\mathrm{ap}}(T_R;\mathbb{C}^4)$, and that $\|\mathrm{Sx}*\psi\|\le\|\mathrm{Sx}\|_{L^1}\,\|\psi\|$.
For $\psi=\mathrm{Ew}_P\,\xi$, (K-2) gives
\[
(\mathrm{Sx}*\psi)(y)=\left(\int\mathrm{d}y''\,\mathrm{Sx}(y'')\,\mathrm{e}^{-iP\cdot y''}\right)\mathrm{Ew}_P(y)\,\xi=\mathrm{Sdr}(P,\bar{m})\,\mathrm{Ew}_P(y)\,\xi=\left(S^R\psi\right)(y).
\]
Two bounded operators that agree on an orthonormal basis are equal; hence $S^R\psi=\mathrm{Sx}*\psi$.
Now let $\psi_0\in L^2(Q_{\mathrm{be}};\mathbb{C}^4)$.
Splitting $\mathbb{R}^3$ into the sets $Q_{\mathrm{be}}+Rz$ and their complement, one finds, for almost every $y\in Q_{\mathrm{be}}$,
\begin{align*}
\left(\mathrm{inj}_R^{*}\,S^R\,\mathrm{inj}_R\,\psi_0\right)(y)
&=\sum_{z\in\mathbb{Z}^3}(-1)^{z_1+z_2+z_3}\int_{Q_{\mathrm{be}}}\mathrm{d}y'\,\mathrm{Sx}(y-y'-Rz)\,\psi_0(y')\\
&=\left(\mathrm{inj}_\infty^{*}\,S_{\bar{m}}\,\mathrm{inj}_\infty\,\psi_0\right)(y)+\int_{Q_{\mathrm{be}}}\mathrm{d}y'\,\mathrm{Wim}(y,y')\,\psi_0(y'),\\
\mathrm{Wim}(y,y')&:=\sum_{z\in\mathbb{Z}^3,\,z\ne0}(-1)^{z_1+z_2+z_3}\,\mathrm{Sx}(y-y'-Rz);
\end{align*}
the term $z=0$ is identified by (K-3).
Therefore $\mathrm{Tq}^{Q}_{R}-\mathrm{Tq}^{Q}_{\infty}$ is the integral operator on $L^2(Q_{\mathrm{be}};\mathbb{C}^4)$ with the kernel $\mathrm{Wim}(y,y')\,i\gamma\cdot a(y')$.

\emph{Size.}
Let $y,y'\in Q_{\mathrm{be}}$ and $z\ne0$.
Since $R\ge1$ and $R-D_{\mathrm{be}}\ge1$,
\[
|y-y'-Rz|_2\ge|y-y'-Rz|_\infty\ge R|z|_\infty-D_{\mathrm{be}}=(R-D_{\mathrm{be}})+R\left(|z|_\infty-1\right)\ge(R-D_{\mathrm{be}})+\left(|z|_\infty-1\right)\ge1 .
\]
By [2.18], and because exactly $24j^2+2$ points $z$ satisfy $|z|_\infty=j$,
\[
\left|\mathrm{Wim}(y,y')\right|_2\le\frac{3}{4\pi}\sum_{j\ge1}\left(24j^2+2\right)\exp\left(-\bar{m}(R-D_{\mathrm{be}})-\bar{m}(j-1)\right)=\frac{3}{4\pi}\,\mathrm{Sgm}\cdot\exp\left(-\bar{m}(R-D_{\mathrm{be}})\right).
\]
One has $\left\|\mathrm{Wim}(y,y')\,\gamma\cdot a(y')\right\|_2^2\le\left|\mathrm{Wim}(y,y')\right|_2^2\,4\,|a(y')|_2^2\le\left|\mathrm{Wim}(y,y')\right|_2^2\,12\,\mathrm{Ab}_1^2$, and the volume of $Q_{\mathrm{be}}$ is $D_{\mathrm{be}}^3$;
hence
\[
\left\|\mathrm{Tq}^{Q}_{R}-\mathrm{Tq}^{Q}_{\infty}\right\|_2\le\frac{3}{4\pi}\,\mathrm{Sgm}\left(12\,\mathrm{Ab}_1^2\cdot D_{\mathrm{be}}^6\right)^{\frac12}\exp\left(-\bar{m}(R-D_{\mathrm{be}})\right)=\mathrm{wim}_0\cdot\exp\left(-\bar{m}(R-D_{\mathrm{be}})\right).
\]
Consequently $\left\|\mathrm{Tq}^{Q}_{R}\right\|_4\le\mathrm{kn}_2+\mathrm{wim}_0$
(by [2.31], $\|\cdot\|_4\le\|\cdot\|_2$ and $R-D_{\mathrm{be}}\ge0$),
and (Tr2) together with [2.31] gives
\[
\left|\operatorname{Tr}\left(S^R\,\mathrm{va}^R\right)^4-\operatorname{Tr}\left(S_{\bar{m}}\,\mathrm{va}\right)^4\right|\le4\left(\mathrm{kn}_2+\mathrm{wim}_0\right)^3\cdot\mathrm{wim}_0\cdot\exp\left(-\bar{m}(R-D_{\mathrm{be}})\right)=\mathrm{lnk}_6\cdot\exp\left(-\bar{m}(R-D_{\mathrm{be}})\right).
\]
Multiplication by $6$ yields (e) with $C_{2}^{\mathrm{N2}}=6\,\mathrm{lnk}_6$.
\end{proof}

\paragraph{2.12. Clauses (f), (g).}

\begin{proof}[Proof of Proposition~37.2\,(f), (g)]
\begin{itemize}
\item[(f)]
The first inequality is the triangle inequality combined with (d) and (e).
The second follows from the first, since $\eta\le1$ and $R\ge D_{\mathrm{be}}$.
The right side of the first inequality tends to $0$ as $\eta\to0$ and $R\to\infty$, in any manner.
Let $RL^N$ be an integer for all $N\ge N_1$.
Since $R>1$ and $L\ge5$, one has $RL^N\ge2$ for $N\ge1$;
so (FF) holds with $\eta=L^{-N}$ for $N\ge\max(N_1,1)$, and (d) gives
\[
\left|b_{\mathrm{ch}}\left(L^{-N},R\right)-b_R(f)\right|\le C_{1}^{\mathrm{N2}}\cdot L^{-\frac{N}{4}}\to0 .
\]
\item[(g)]
The numbers of 2.0 are formed from absolute numbers, from $\bar{m}$, from $C_E$, from $D_{\mathrm{be}}$, from sup norms and $L^1(\mathbb{R}^3)$ norms of derivatives of the components of $\beta$, and from the moduli $\left|\hat\beta_{\mu\nu}(\mathrm{kph})\right|$.
None of them contains $\eta$, $R$ or $n$.
For $v\in\mathbb{R}^3$ the probe $\tau_v\beta$ has the same $D_{\mathrm{be}}$ and the same norms of derivatives,
and $\left|\widehat{(\tau_v\beta)}_{\mu\nu}(\mathrm{kph})\right|=\left|\hat\beta_{\mu\nu}(\mathrm{kph})\right|$, since a translation multiplies the Fourier transform by a phase.
Hence $C_{\mathrm{ct}}$, $C_{1}^{\mathrm{N2}}$ and $C_{2}^{\mathrm{N2}}$ are the same for all $\tau_v\beta$.
In all bounds of Section~2 the factor $1+\rho$ is estimated by $2$, and $\rho$ occurs nowhere else.
So $\rho$ enters only through $C_E=C_E(\rho)$:
it enters $C_{\mathrm{ct}}$ through $\mathrm{Cf}$, and $C_{1}^{\mathrm{N2}}$ through $C_{\mathrm{ct}}$, $\mathrm{kn}_1$ and $\mathrm{lnk}_1$, while $C_{2}^{\mathrm{N2}}$ does not contain $C_E$.
Finally, $C_{\mathrm{ct}}$ does not contain $\bar{m}$.
\end{itemize}
\end{proof}

\subsection*{3. Proposition 37.3}

\subsubsection*{Proof of Proposition 37.3}

\paragraph{3.1. The data are data of (FF).}
The number $\eta=L^{-N}$ lies in $(0,1]$.
By Section~I.6.4, $n_*=R^{*}(e)/\eta=L\,\mathrm{per}$ with the integer $\mathrm{per}$, and $\mathrm{per}\ge L^4$ by Proposition~35.1\,(f);
hence $n_*$ is an integer $\ge2$.
By Proposition~35.1\,(f), together with $h\ge1$ and $L\ge5$ (Section~I.0.5), one has $R^{*}(e)\ge hL\ge1$.
The inequality $R^{*}(e)\ge D_{\mathrm{be}}+1$ belongs to the data.
For $v\in\mathbb{R}^3$, $\tau_v\beta$ is a probe with the same number $D_{\mathrm{be}}$.
Consequently the pair $(\eta,R)=(\eta,R^{*}(e))$, formed with $\beta$ or with $\tau_v\beta$, is admissible in the setting (FF),
in Proposition~36\,(f), (g) (which require $0<\eta\le1$, $R\ge1$ and that $R/\eta$ be an integer $\ge2$)
and in Proposition~37.2 (which requires $R\ge D_{\mathrm{be}}+1$).
For a probe $\beta'$, the symbols $N_0[\beta']$, $a_{\mathrm{ch}}(\eta,R^{*}(e))[\beta']$ and $b_{\mathrm{ch}}(\eta,R^{*}(e))[\beta']$ denote the objects of Section~G5.6.0 on $C(n_*,n_*,n_*)$.

\paragraph{3.2. The free ratio.}
For a complex one-cochain $A'$ on $\mathbb{Z}^3/n_*\mathbb{Z}^3$, ``$N_0$ formed with $A'$'' denotes the function
\[
t\mapsto\frac{\det\left(D_{(t/e)A'}+\bar{m}\eta\right)}{\det\left(D_0+\bar{m}\eta\right)}\qquad\text{on }C(n_*)
\]
(Section~I.7.6; here $D_A$ is (I.2.3) in lattice units, with the charge $q=e\,\eta^{\frac12}$).
By Proposition~46\,(b), every complex one-cochain $A$ satisfies $D_A+\bar{m}\eta=D^{\mathrm{ch}}[qA]+\bar{m}\eta$ on $\mathfrak{S}^{\mathrm{ch}}(1;n)$ with $n=(n_*,n_*,n_*)$.
Take $A'=\hat{\mathfrak{a}}=\eta^{\frac12}a_{\eta}[\tau_v\beta]$ and $A=(t/e)\,\hat{\mathfrak{a}}$.
Then
\[
qA=e\,\eta^{\frac12}\cdot\frac{t}{e}\cdot\eta^{\frac12}a_{\eta}[\tau_v\beta]=t\,\eta\,a_{\eta}[\tau_v\beta],
\]
which is the cochain $B_{t}$ of Section~G5.6.0 formed with $\tau_v\beta$; for $A=0$ one has $qA=0$.
Therefore $N_0$ formed with $\hat{\mathfrak{a}}$ is the function
\[
t\mapsto\frac{\det\left(D^{\mathrm{ch}}[B_{t}]+\bar{m}\eta\right)}{\det\left(D^{\mathrm{ch}}[0]+\bar{m}\eta\right)}=N_0[\tau_v\beta](t).
\]

\paragraph{3.3. Clause (a).}
\begin{itemize}
\item[(1)]
Let $v\in\eta\mathbb{Z}^3$.
By the last sentence of Proposition~36\,(h), $N_0$ does not change when $\beta$ is replaced by $\tau_v\beta$:
$N_0[\tau_v\beta]=N_0[\beta]$.
\item[(2)]
Let $\mathfrak{c}\in\mathrm{Cel}$.
By Proposition~35.2\,(a), $\mathsf{g}_{\mathfrak{c}}y=r_{\mathfrak{c}}y+a_{\mathfrak{c}}$,
where $r_{\mathfrak{c}}$ lies in $\{\pm1\}^3$, a subgroup of $O_{\mathrm{cub}}$, and $a_{\mathfrak{c}}$ lies in $\mathrm{per}\,\mathbb{Z}^3$, a subset of $\mathbb{Z}^3$.
Thus $\mathsf{g}_{\mathfrak{c}}$ is the element $g_{\mathrm{w}}=(a_{\mathfrak{c}},r_{\mathfrak{c}})$ of $W_{\mathrm{af}}$.
Since $r_{\mathfrak{c}}\,n_*\mathbb{Z}^3=n_*\mathbb{Z}^3$, the map $(\mathsf{g}_{\mathfrak{c}})_*$ sends one-cochains on $\mathbb{Z}^3/n_*\mathbb{Z}^3$ to one-cochains on $\mathbb{Z}^3/n_*\mathbb{Z}^3$, and it is linear.
Moreover $\hat{\mathfrak{a}}_{\mathfrak{c}}=(\mathsf{g}_{\mathfrak{c}})_*\hat{\mathfrak{a}}$.
The first sentence of Proposition~36\,(h), applied with $n=n_*$, $\boldsymbol{\mu}=\bar{m}\eta$ and the complex one-cochain $A=(t/e)\,\hat{\mathfrak{a}}$, gives for every $t\in\mathbb{C}$
\[
\det\left(D_{(t/e)\hat{\mathfrak{a}}_{\mathfrak{c}}}+\bar{m}\eta\right)=\det\left(D_{(\mathsf{g}_{\mathfrak{c}})_*\left((t/e)\hat{\mathfrak{a}}\right)}+\bar{m}\eta\right)=\det\left(D_{(t/e)\hat{\mathfrak{a}}}+\bar{m}\eta\right).
\]
The denominators coincide.
Hence $N_0$ formed with $\hat{\mathfrak{a}}_{\mathfrak{c}}$ is equal to $N_0$ formed with $\hat{\mathfrak{a}}$.
\item[(3)]
By (1), (2) and 3.2, $N_0=N_0[\beta]$ in all cases;
it depends neither on $v$ nor on $\mathfrak{c}$, and $c_{m}=N_0[\beta]^{(m)}(0)$.
Proposition~36\,(a) now gives
\[
c_{1}=c_{3}=0,\qquad c_{2}=a_{\mathrm{ch}}(\eta,R^{*}(e))[\beta],\qquad c_4-3c_{2}^2=b_{\mathrm{ch}}(\eta,R^{*}(e))[\beta].
\]
\end{itemize}

\paragraph{3.4. Clause (b).}
By (a), 3.1 and Proposition~36\,(g), $|c_{2}|=\left|a_{\mathrm{ch}}(\eta,R^{*}(e))\right|\le\mathrm{Cpv}_\beta$.
By (a) and the second inequality of Proposition~37.2\,(f),
\[
|c_4|\le\left|b_{\mathrm{ch}}(\eta,R^{*}(e))\right|+3c_{2}^2\le|b(f)|+C_{1}^{\mathrm{N2}}+C_{2}^{\mathrm{N2}}+3\,\mathrm{Cpv}_\beta^2 .
\]
The display of (b) is the sum of the first bound and of the square root of the second.
Its right side contains neither $e$ nor $N$.

\paragraph{3.5. Clause (c).}
By (a) and the first inequality of Proposition~37.2\,(f) at $R=R^{*}(e)$,
\[
\left|c_4-3c_{2}^2-b(f)\right|\le C_{1}^{\mathrm{N2}}\cdot\eta^{\frac14}+C_{2}^{\mathrm{N2}}\cdot\exp\left(-\bar{m}\left(R^{*}(e)-D_{\mathrm{be}}\right)\right).
\]
By Proposition~35.1\,(f), $\exp\left(-\bar{m}\cdot R^{*}(e)\right)<e\cdot h^{\frac12}$; hence
\[
\exp\left(-\bar{m}\left(R^{*}(e)-D_{\mathrm{be}}\right)\right)=\exp\left(\bar{m}\cdot D_{\mathrm{be}}\right)\exp\left(-\bar{m}\cdot R^{*}(e)\right)\le\exp\left(\bar{m}\cdot D_{\mathrm{be}}\right)\cdot h^{\frac12}\cdot e .
\]

\paragraph{3.6. Clause (d).}
It follows from (a) and from Proposition~36\,(f) at $R=R^{*}(e)$.
\qed

\begin{named}{Remark}
Clause~(c), and the bound $\left|c_4-3c_{2}^2\right|\le|b(f)|+C_{1}^{\mathrm{N2}}+C_{2}^{\mathrm{N2}}$ that enters 3.4, can also be derived by applying Proposition~37.2\,(f) to the probe $\tau_v\beta$ in place of $\beta$.
Indeed, by Proposition~37.2\,(g) the constants $C_{1}^{\mathrm{N2}}$, $C_{2}^{\mathrm{N2}}$ formed with $\tau_v\beta$ coincide with those formed with $\beta$;
and $b\left(d\,d^{*}\tau_v\beta\right)=b(f)$, because the translation by $v$ is a unitary map of $L^2(\mathbb{R}^3;\mathbb{C}^4)$ that commutes with $S_{\bar{m}}$ and conjugates $i\gamma\cdot a$ into $i\gamma\cdot\tau_v a$.
\end{named}

\subsection*{4. Conditions on the constants}

The proofs of Propositions~37.1--37.3 use the following conditions on the constants.

\medskip
\noindent
\begin{tabular}{p{0.48\textwidth}p{0.16\textwidth}p{0.2\textwidth}}
\hline
condition or constant & where it enters & its line in Section~I.8.2 / Proposition~14.2\,(a)\\
\hline
$0<\rho<1$; for instance, $\rho\ge0$ enters the last inequality of 1.1; $\rho>0$ enters $\mathrm{kah}$, the second case of 1.6 and 1.9; $\rho\le1$ enters $1+\rho\le2$ and $\rho\sum_\mu(1-\cos p_\mu)\le|p|_2^2/2$
& Sections 1, 2 & (C0)\\
the standard weights $w=(6,6\sqrt{2},6\sqrt{2},6)$, with (I.1.1) and $w_3=6$
& 1.1--1.4, 1.9, 2.1, 2.10 (Step 1) & (C0)\\
$\lambda_1=6+9\sqrt{2}-\sqrt{198-36\sqrt{2}}$; that this number is the gap of $(\nabla^{*}\nabla)(0)$ is proved in 1.2
& 1.2, 1.6, 1.9 & (C0r)\\
$C_E:=$ the number [1.10] (it may be replaced by any larger number)
& 1.6 & (C0r)\\
$\mathrm{Csy}_{\mathrm{ak}}:=$ the numbers [1.12], for all multi-indices $\mathrm{ak}$
& 1.7 & (C0r); Proposition~14.2\,(a) involves these numbers only for $|\mathrm{ak}|\le4$\\
$\bar{m}\in(0,1]$; for instance, $\bar{m}>0$ enters $\mathbf{m}>0$ for Proposition~37.1\,(f) ([2.3]), [2.5], the third bound of [2.6], [2.7], the bound $3/\bar{m}$ of [2.18], (K-2), Lemmas~2.3 and 2.4, and $\mathrm{Sgm}$; $\bar{m}\le1$ enters $\mathbf{m}\le1$ and [2.18]
& 1.8, Section 2 & (C10)\\
$C_{\mathrm{ct}}$, $C_{1}^{\mathrm{N2}}$, $C_{2}^{\mathrm{N2}}:=$ the numbers [2.1]; they depend on $(\rho,\beta)$, $(\rho,\bar{m},\beta)$, $(\bar{m},\beta)$ respectively, and on $\rho$ through $C_E$
& 2.0 & (C10)\\
$L\ge5$ (only $L\ge2$ enters the arguments), $h\ge1$
& 2.12\,(f), 3.1 & (C3), (C10); Section~I.0.5\\
$e\cdot h^{\frac12}\le e_1$, $N\ge N_0(e)$ and $R^{*}(e)\ge D_{\mathrm{be}}+1$ are data of Proposition~37.3; the conditions behind Propositions~35.1\,(f), 35.2\,(a) and~36 are those of these statements
& Section 3 & --\\
\hline
\end{tabular}
\medskip

Each of these conditions, apart from those of the last row, belongs to a line of Section~I.8.2 or of Proposition~14.2\,(a) named in the third column.
Hence Propositions~37.1--37.3 hold for all values of $L$, $M$, $e_0$, $e_1$, $e_2$ and $e_{\sharp}$ that Section~I.8.2 and Proposition~14.2\,(a) admit.

\section*{Chapter 36. The two-photon coefficient of the free ratio, its limit and its uniform bound: proof of Proposition 36}

Throughout this chapter the setting (FF) of Section~G5.6.0 is in force, together with the free ratio $N_0$, the fibre ${V_{\mathrm{v}}}$ and the continuum objects defined there.
The functions $\mathrm{a0}_{\eta,\nu}$, $\mathrm{ahl}_{\eta,\nu}$, $\mathrm{lp}_{\mu\nu}$, $\mathrm{Vpl}^{\eta,R}_{\mu\nu}$, $\mathrm{Vpl}^{\eta}_{\mu\nu}$ and the set $\mathrm{Kx}_{\eta,R}$ are those defined in Proposition~36.

\subsection*{Proof of Proposition 36}

The clauses are proved in the order (a), (h), (b), (c), (d), (e), (f), (g).
Clauses (c) and (d) concern the fibre ${V_{\mathrm{v}}}$ alone, and clause (e) concerns the probe alone.

\subsubsection*{0. Conventions of this chapter: setting, notation, elementary facts}

\paragraph{0.1. Where the statement lives.}
All objects are taken in lattice units in the sense of the preamble of Proposition~46.
Thus Sections~I.1 and~I.2 are read at the single level $k=0$, with $\eta_0=1$ and with a number $q>0$ in the place of $e_k$, on the cubic torus $n=(n,n,n)$, $n\ge2$; the number $q$ drops out (subsection~0.2).
A fermion site is a pair $x=(c,g)$ with $c\in\mathbb{Z}^3$ and $g\in O_{\mathrm{cub}}$.
The cluster of $c$ consists of the $48$ sites $(c,g)$, $g\in O_{\mathrm{cub}}$, and the neighbours $x^{(i)}$ are those of (I.1.2) with $\eta_0=1$.
A one-cochain on $\mathbb{Z}^3/n\mathbb{Z}^3$ is a complex function $A$ on the oriented nearest-neighbour bonds of $\mathbb{Z}^3$ which is $n\mathbb{Z}^3$-periodic and satisfies $A(c',c)=-A(c,c')$ ((FF); preamble of Proposition~46).
Every bond is a bond of $\mathbb{Z}^3$; in particular, for $n=2$ the bonds $(c,c+e_\nu)$ and $(c,c-e_\nu)$ are different bonds.
Let
\begin{equation*}
\mathrm{Hch}:=\mathfrak{S}^{\mathrm{ch}}(1;n)=\left\{f\colon\mathbb{Z}^3\times O_{\mathrm{cub}}\to\mathbb{C}^4\ :\ f(c+n\cdot e_\mu,g)=-f(c,g),\ \mu=1,2,3\right\},
\end{equation*}
so that $\dim\mathrm{Hch}=4\cdot48\cdot n^3=192\cdot n^3$, and put
\begin{equation*}
\mathbf{m}:=\bar{m}\cdot\eta\in(0,1].
\end{equation*}
The pairings on $\mathrm{Hch}$ are those of Section~I.1.6 at $\eta_0=1$,
\begin{gather*}
\langle f',f\rangle^{\mathrm{ch}}=\frac{1}{48}\sum_{c\in\{0,\dots,n-1\}^3}\ \sum_{g\in O_{\mathrm{cub}}}f'(c,g)\cdot f(c,g),\\
(f',f)^{\mathrm{ch}}=\frac{1}{48}\sum_{c\in\{0,\dots,n-1\}^3}\ \sum_{g\in O_{\mathrm{cub}}}\overline{f'(c,g)}\cdot f(c,g),
\end{gather*}
and their summands are $n\mathbb{Z}^3$-periodic in $c$.
The trace on $\mathrm{Hch}$ is denoted $\mathrm{Tr}$, the trace on ${V_{\mathrm{v}}}$ is denoted $\mathrm{tr}$, and the trace on $\mathbb{C}^4$ is denoted $\mathrm{tr}_4$.
Traces and determinants do not depend on a pairing.

\paragraph{0.2. The operator.}
Let $A$ be a one-cochain and $B:=q\cdot A$.
By Proposition~46\,(b), the operator $D^{\mathrm{ch}}[B]$ is the operator $D_A$ of (I.2.3) in lattice units.
Written out by means of (I.1.2) and (I.2.1), it acts on functions $f\colon\mathbb{Z}^3\times O_{\mathrm{cub}}\to\mathbb{C}^4$ by
\begin{equation}\tag{0.1}
\begin{aligned}
\left(D^{\mathrm{ch}}[B]f\right)(c,g)={}&-\sum_{i\le2}M_i(g)\cdot\left[f(c,gs_i)-f(c,g)\right]\\
&-M_3(g)\cdot\left[\exp\left(\mathrm{i}\,B(c,c+ge_1)\right)\cdot f(c+ge_1,gs_3)-f(c,g)\right],\\
M_i(g)={}&\frac{w_i}{2}\cdot\left(\rho+\gamma\cdot(g\hat{n}_i)\right)\quad(i=0,1,2),\qquad M_3(g)=3\cdot\left(\rho-\gamma\cdot(ge_1)\right).
\end{aligned}
\end{equation}
The number $q$ enters only through $B$.
Since $x\mapsto x^{(i)}$ commutes with the translations by $n\mathbb{Z}^3$ (I.s1) and $B$ is periodic, $D^{\mathrm{ch}}[B]$ maps $\mathrm{Hch}$ into itself.

Let $\hat{\mathfrak{a}}:=\eta^{\frac{1}{2}}\cdot a_\eta[\beta]$ be the one-cochain of Section~G5.1.3 in lattice units, let $e>0$ be an arbitrary number, $q:=e\cdot\eta^{\frac{1}{2}}$, and $t\in\mathbb{C}$.
On the bond $(c,c+e_\nu)$ the cochain $q\cdot\frac{t}{e}\cdot\hat{\mathfrak{a}}$ has the value
\begin{equation*}
e\cdot\eta^{\frac{1}{2}}\cdot\frac{t}{e}\cdot\eta^{\frac{1}{2}}\cdot a_\eta[\beta]_\nu(\eta\cdot c)=t\cdot\eta\cdot a_\eta[\beta]_\nu(\eta\cdot c),
\end{equation*}
which is $B_t(c,c+e_\nu)$ by (FF).
Hence $B_t=q\cdot\frac{t}{e}\cdot\hat{\mathfrak{a}}$, Proposition~46\,(b) yields $D^{\mathrm{ch}}[B_t]=D_{(t/e)\cdot\hat{\mathfrak{a}}}$, and on $\mathrm{Hch}$
\begin{equation*}
N_0(t)=\frac{\det\left(D_{(t/e)\cdot\hat{\mathfrak{a}}}+\mathbf{m}\right)}{\det\left(D_0+\mathbf{m}\right)}.
\end{equation*}
So $N_0$ is the free ratio of Section~I.7.6 on the torus $C(n,n,n)$, formed with $\beta$.
This is Proposition~46\,(g+) at $v=0$; that clause holds for all $\eta\in(0,1]$ and all $R$ for which $R/\eta$ is an integer $\ge2$, which is the range of (FF).

\paragraph{0.3. Notation.}
The facts of subsection~0.4 are labelled (Ef0)--(Ef5).
The letters $\mathrm{G0}$, $\mathrm{Vd}_1$, $\mathrm{Vd}_2$ have the meaning of Proposition~37.2.
A function on $\mathbb{R}^3$ or on $\mathbb{Z}^3$ is denoted $\mathrm{fs}$, and $\mathrm{fh}$ is its Fourier coefficient (Sections~I.7.6 and~I.2.6).
The letter $\mathrm{Mx}$ denotes a $4\times4$ matrix (Section~I.1.6), the letters $\mathrm{sv},\mathrm{sv}',\mathrm{sv}_1,\dots$ denote real variables, and $\mathrm{rad}$ denotes a radius.
The letters $\mathrm{lv},\mathrm{lv}'$ denote vectors of $\mathbb{R}^3$ or of $\mathbb{C}^3$, as stated at each occurrence, and $O,O',O_1,\dots$ denote arbitrary linear maps, as in Section~I.1.6.
As indices, $\alpha,\beta,\mathfrak{l},\mu,\nu$ denote directions in $\{1,2,3\}$; the index $\beta\in\{1,\dots,4\}$ of the standard basis vectors $\mathrm{ue}_\beta$ of $\mathbb{C}^4$, however, is the spinor index of Section~I.1.6, and the index of the constant $\mathrm{Cpv}_\beta$ of Proposition~36\,(g) is the probe.
The probe itself is denoted $\beta$, as in $a_\eta[\beta]$, $\tau_v\beta$, $\mathrm{Cah}_i(\beta)$ and $|\beta|_{C^i}$, and its components are $\beta_{\mu\nu}$.

Besides the imaginary unit $\mathrm{i}$, the letter $i$ denotes an index of one of two kinds: the type $i\in\{0,\dots,3\}$ of a bond (I.1.2), as in $M_i$, $R_i(A)$ and $x^{(i)}$, or an integer $i\ge0$, such as an order of decay or a number of differences or of derivatives, as in $\mathrm{Cah}_i$, $\mathrm{Cahd}_i$ and $|\beta|_{C^i}$.

\paragraph{0.4. Elementary facts.}
\begin{description}
\item[(Ef0)] \emph{Numerals.}
One has $3.1415<\pi<3.1416$; this is classical and follows, for instance, from Machin's formula $\pi=16\cdot\arctan\frac{1}{5}-4\cdot\arctan\frac{1}{239}$ and the alternating series of the arc tangent.
From $3.1415^2=9.86902225$ and $3.1416^2=9.86965056$ one obtains $9.8690<\pi^2<9.8697$.
From $1.41421^2=1.9999899241<2$ one obtains $\sqrt{2}>1.41421$, and from $1.732^2=2.999824<3<3.00017041=1.7321^2$ one obtains $1.732<\sqrt{3}<1.7321$.
These bounds imply the following inequalities between numerals.
\begin{itemize}
\item[(i)] $\pi\cdot\sqrt{3}<3.1416\cdot1.7321=5.44156536<6$.
\item[(ii)] $4\pi^2+\frac{\pi^4}{45}<4\cdot9.8697+\frac{98}{45}<39.4788+2.18<42$, because $9.8697^2=97.41097809<98$.
\item[(iii)] $\frac{42}{30\pi^2}<\frac{1}{7}$, because $30\pi^2>30\cdot9.8690=296.07>294=7\cdot42$.
\item[(iv)] $\frac{3456}{\pi^2}<351$, because $351\cdot9.8690=3464.019>3456$.
\item[(v)] $\frac{6912}{\pi^2}<701$, because $701\cdot9.8690=6918.169>6912$.
\item[(vi)] $\frac{2\sqrt{2}}{3\pi}>\frac{3}{10}$, because $20\sqrt{2}>28.2842>28.2744=9\cdot3.1416>9\pi$.
\item[(vii)] $\frac{20\sqrt{3}}{3}<12$, because $\sqrt{3}<1.8$.
\item[(viii)] $\frac{9\cdot7^4}{6\pi^2}<365$, because $9\cdot7^4=21609<21613.11=2190\cdot9.8690<6\cdot365\cdot\pi^2$.
\item[(ix)] $\frac{\sqrt{3}\cdot7^4}{6\pi}<221$, because $7^4\cdot\sqrt{3}<2401\cdot1.7321=4158.7721<4165.629=1326\cdot3.1415<6\cdot221\cdot\pi$.
\item[(x)] $192\cdot\left(6\cdot7\cdot8^4+36\cdot16^4\right)=192\cdot(172032+2359296)=192\cdot2531328=486014976<49\cdot10^7$.
\end{itemize}
\item[(Ef1)] \emph{The distance to $2\pi\mathbb{Z}^3$.}
For $p\in\mathbb{R}^3$ let
\begin{equation*}
\mathrm{xd}(p):=\mathrm{dist}_2\left(p,2\pi\mathbb{Z}^3\right)+\mathbf{m}.
\end{equation*}
The minimum of $\sum_\nu(p_\nu-2\pi z_\nu)^2$ over $z\in\mathbb{Z}^3$ is taken coordinate by coordinate, so that $\mathrm{dist}_2\left(p,2\pi\mathbb{Z}^3\right)^2=\sum_\nu\mathrm{dist}\left(p_\nu,2\pi\mathbb{Z}\right)^2$.
Hence $\mathrm{xd}$ is continuous and $2\pi\mathbb{Z}^3$-periodic, $\mathrm{xd}(p)=|p|_2+\mathbf{m}$ for $p\in[-\pi,\pi]^3$, and, by (Ef0)\,(i), $\mathbf{m}\le\mathrm{xd}(p)\le\pi\cdot\sqrt{3}+1<7$ for all $p$.
For a real number $\mathrm{sv}$ one has $\mathrm{dist}(\mathrm{sv},2\pi\mathbb{Z})=|\mathrm{sv}|$ if $|\mathrm{sv}|\le\pi$, and $\mathrm{dist}(\mathrm{sv},2\pi\mathbb{Z})=2\pi-|\mathrm{sv}|\ge\frac{\pi}{2}\ge\frac{|\mathrm{sv}|}{3}$ if $\pi\le|\mathrm{sv}|\le\frac{3\pi}{2}$.
It follows that $\mathrm{dist}_2\left(p,2\pi\mathbb{Z}^3\right)\ge\frac{|p|_2}{3}$ whenever $|p_\nu|\le\frac{3\pi}{2}$ for $\nu=1,2,3$, and
\begin{equation}\tag{0.2}
\mathrm{xd}(p)\ge\frac{|p|_2+\mathbf{m}}{3}\qquad\text{if }|p|_\infty\le\frac{3\pi}{2}.
\end{equation}
\item[(Ef2)] \emph{Two integrals.}
For an integer $j\ge4$, Euler's Beta integral $\mathrm{B}(3,j-3)$ gives
\begin{equation*}
\int_0^\infty\mathrm{d}\mathrm{sv}\,\mathrm{sv}^2\cdot(\mathrm{sv}+1)^{-j}=\frac{2}{(j-1)(j-2)(j-3)}.
\end{equation*}
By periodicity, (Ef1) and polar coordinates one has, for every $p''\in\mathbb{R}^3$,
\begin{equation}\tag{0.3}
\begin{aligned}
\int_{[-\pi,\pi]^3}\mathrm{d}^3p\,\mathrm{xd}(p+p'')^{-j}&=\int_{[-\pi,\pi]^3}\mathrm{d}^3p\,\left(|p|_2+\mathbf{m}\right)^{-j}\le4\pi\int_0^\infty\mathrm{d}\mathrm{sv}\,\mathrm{sv}^2\cdot(\mathrm{sv}+\mathbf{m})^{-j}\\
&=\frac{8\pi\cdot\mathbf{m}^{3-j}}{(j-1)(j-2)(j-3)}\quad
\begin{cases}
\le\frac{4\pi}{3}\cdot\mathbf{m}^{-1}&\text{for }j=4,\\
\le\frac{2\pi}{15}\cdot\mathbf{m}^{-3}&\text{for }j=6.
\end{cases}
\end{aligned}
\end{equation}
The invariance of the Lebesgue measure under translations, polar coordinates and the Beta integral at $j=4$ give, for every number $\mathrm{mq}>0$ and every $\mathrm{lv}'\in\mathbb{R}^3$,
\begin{equation}\tag{0.4}
\int_{\mathbb{R}^3}\mathrm{d}^3\mathrm{lv}\,\left(|\mathrm{lv}+\mathrm{lv}'|_2+\mathrm{mq}\right)^{-4}=4\pi\int_0^\infty\mathrm{d}\mathrm{sv}\,\mathrm{sv}^2\cdot(\mathrm{sv}+\mathrm{mq})^{-4}=\frac{4\pi}{3\cdot\mathrm{mq}}.
\end{equation}
\item[(Ef3)] For positive numbers $\mathrm{xd}_1,\mathrm{xd}_2$ and for $j_1,j_2\ge0$,
\begin{equation*}
\mathrm{xd}_1^{-j_1}\cdot\mathrm{xd}_2^{-j_2}\le\max\left(\mathrm{xd}_1^{-1},\mathrm{xd}_2^{-1}\right)^{j_1+j_2}\le\mathrm{xd}_1^{-j_1-j_2}+\mathrm{xd}_2^{-j_1-j_2}.
\end{equation*}
\item[(Ef4)] \emph{Traces; the weight $\frac{1}{48}$.}
Let $\mathrm{ue}_\beta$, $\beta=1,\dots,4$, be the standard basis of $\mathbb{C}^4$.
For the inner product $(u',u)=\frac{1}{48}\sum_g\overline{u'(g)}\cdot u(g)$ the $192$ vectors $\sqrt{48}\cdot\mathbf{1}[g=g']\cdot\mathrm{ue}_\beta$, $g'\in O_{\mathrm{cub}}$, $\beta=1,\dots,4$, form an orthonormal basis of ${V_{\mathrm{v}}}$.
Therefore $|\mathrm{tr}\,O|\le192\cdot|O|_2$ for every linear map $O$ of ${V_{\mathrm{v}}}$.
Let $\mathrm{Mx}$ be a $4\times4$ matrix and $O$ a linear map of ${V_{\mathrm{v}}}$, and let $(J\xi)(g)=\xi$ and $J^{*}u=\frac{1}{48}\sum_{g\in O_{\mathrm{cub}}}u(g)$, as in (I.3.1).
The cyclicity of the trace for the maps $J\cdot\mathrm{Mx}\colon\mathbb{C}^4\to{V_{\mathrm{v}}}$ and $J^{*}\cdot O\colon{V_{\mathrm{v}}}\to\mathbb{C}^4$ gives
\begin{equation}\tag{0.5}
\mathrm{tr}\left[J\cdot\mathrm{Mx}\cdot J^{*}\cdot O\right]=\mathrm{tr}_4\left[\mathrm{Mx}\cdot J^{*}\cdot O\cdot J\right].
\end{equation}
\item[(Ef5)] \emph{The resolvent of the symbol.}
Let $D^{(p)}$ be the symbol (I.2.6) at $\boldsymbol{\mu}=\mathbf{m}$; it is an entire and $2\pi\mathbb{Z}^3$-periodic function of $p$.
For real $p$, Proposition~37.1\,(a) gives $\mathrm{Re}\left(u,D^{(p)}u\right)\ge\mathbf{m}\cdot(u,u)$.
So $D^{(p)}$ is invertible for real $p$, the function $\mathrm{Rsy}(p)=\left(D^{(p)}\right)^{-1}$ of Section~G5.6.0 is real analytic and $2\pi\mathbb{Z}^3$-periodic on $\mathbb{R}^3$, and differentiation of $\mathrm{Rsy}(p)\cdot D^{(p)}=1$ gives, with $\mathrm{Vx}_\alpha(p)=\frac{\partial D^{(p)}}{\partial p_\alpha}$ as in Section~G5.6.0,
\begin{equation}\tag{0.6}
\partial_\alpha\mathrm{Rsy}=-\mathrm{Rsy}\cdot\mathrm{Vx}_\alpha\cdot\mathrm{Rsy}.
\end{equation}
Let
\begin{equation*}
\mathrm{Ksy}:=\max\left\{1,\ \max\left\{\mathrm{Csy}_{\mathrm{ak}}(\rho):|\mathrm{ak}|\le4\right\}\right\},
\end{equation*}
with the numbers $\mathrm{Csy}_{\mathrm{ak}}(\rho)$ of Proposition~37.1\,(g); there are $35$ multi-indices $\mathrm{ak}$ with $|\mathrm{ak}|\le4$, and $\mathrm{Ksy}$ is a function of $\rho$.
By Proposition~37.1\,(g) and~(d), for all real $p$ and all $0<\mathbf{m}\le1$,
\begin{equation}\tag{0.7}
\begin{gathered}
\left|\partial^{\mathrm{ak}}\mathrm{Rsy}(p)\right|_2\le\mathrm{Ksy}\cdot\mathrm{xd}(p)^{-1-|\mathrm{ak}|}\quad(|\mathrm{ak}|\le4),\qquad
\left|\partial^{\mathrm{ak}}\mathrm{Vx}_\nu(p)\right|_2\le6\quad(\text{all }\mathrm{ak}),\\
\left|\partial^{\mathrm{ak}}\partial_\nu^2D^{(p)}\right|_2\le6\quad(\text{all }\mathrm{ak}),\qquad
\partial_\mu\mathrm{Vx}_\nu=\delta_{\mu\nu}\cdot\partial_\nu^2D^{(p)}.
\end{gathered}
\end{equation}
\end{description}

\subsubsection*{1. Clause (a)}

\paragraph{1.1. The facts (I.2.4), (I.O2), (I.O3) and (I.O5) on $\mathrm{Hch}$.}
In lattice units (subsection~0.1) these facts take the following form on $\mathrm{Hch}$, for every integer $n\ge1$.
Let $q>0$ and let $A$ be a one-cochain.
For $i=0,\dots,3$ put $\left(R_i(A)f\right)(x):=U_A\left(x,x^{(i)}\right)\cdot f\left(x^{(i)}\right)$, where $U_A\left(x,x^{(i)}\right)=1$ for $i=0,1,2$ and $U_A\left(x,x^{(3)}\right)=\exp\left(\mathrm{i}\,q\cdot A(c,c_3)\right)$ with $c_3=c+ge_1$ ((I.2.2), (I.1.6)).
By (I.s1) the map $x\mapsto x^{(i)}$ is an involution of $\mathbb{Z}^3\times O_{\mathrm{cub}}$ which commutes with the translations by $n\mathbb{Z}^3$.
Hence it is a bijection of the $48\cdot n^3$ sites modulo $n\mathbb{Z}^3$, $R_i(A)$ maps $\mathrm{Hch}$ into itself, and $R_i(A)^2=1$, because $U_A\left(x,x^{(i)}\right)\cdot U_A\left(x^{(i)},x\right)=1$; for $i=3$ this is the relation $A(c,c_3)+A(c_3,c)=0$.
\begin{itemize}
\item[(i)] One has $M_i(x)=\frac{w_i}{2}\cdot\left(\rho-\gamma\cdot\hat{d}_i(x)\right)$, and $\sum_iw_i\cdot\left(\gamma\cdot\hat{d}_i(x)\right)=0$ by (I.s3).
Therefore (0.1) takes the form
\begin{equation}\tag{1.1}
\begin{gathered}
D_A=\sum_{i=0}^{3}M_i(x)\cdot\left(1-R_i(A)\right)=\frac{\rho}{2}\cdot\nabla_A^{*}\nabla_A+\frac{1}{2}\cdot J_A,\qquad\nabla_A^{*}\nabla_A=\sum_iw_i\cdot\left(1-R_i(A)\right),\\
(J_Af)(x)=\sum_iw_i\cdot\left(\gamma\cdot\hat{d}_i(x)\right)\cdot U_A\left(x,x^{(i)}\right)\cdot f\left(x^{(i)}\right).
\end{gathered}
\end{equation}
\item[(ii)] Let $A$ be real.
Then $U_A\left(x^{(i)},x\right)=\overline{U_A\left(x,x^{(i)}\right)}$, and the substitution $x\mapsto x^{(i)}$ in the sum over the sites modulo $n\mathbb{Z}^3$ gives
\begin{equation*}
\left(f',R_i(A)f\right)^{\mathrm{ch}}=\frac{1}{48}\sum_x\overline{f'(x)}\cdot U_A\left(x,x^{(i)}\right)\cdot f\left(x^{(i)}\right)=\frac{1}{48}\sum_x\overline{U_A\left(x,x^{(i)}\right)\cdot f'\left(x^{(i)}\right)}\cdot f(x)=\left(R_i(A)f',f\right)^{\mathrm{ch}}.
\end{equation*}
Consequently $R_i(A)=R_i(A)^{*}=R_i(A)^{-1}$,
\begin{equation*}
\left(f,\left(1-R_i(A)\right)f\right)^{\mathrm{ch}}=\frac{1}{2}\cdot\left(f-R_i(A)f,\,f-R_i(A)f\right)^{\mathrm{ch}}\ge0,
\end{equation*}
and $\nabla_A^{*}\nabla_A=\left(\nabla_A^{*}\nabla_A\right)^{*}\ge0$.
The multiplication by $\gamma\cdot\hat{d}_i(x)$ is Hermitian by (I.c1), and it anticommutes with $R_i(A)$, because $\hat{d}_i\left(x^{(i)}\right)=-\hat{d}_i(x)$ by (I.s2).
So $\left(\left(\gamma\cdot\hat{d}_i\right)\cdot R_i(A)\right)^{*}=R_i(A)\cdot\left(\gamma\cdot\hat{d}_i\right)=-\left(\gamma\cdot\hat{d}_i\right)\cdot R_i(A)$, that is, $J_A^{*}=-J_A$.
Together with (1.1) and $\rho>0$ this gives
\begin{equation}\tag{1.2}
\mathrm{Re}\left(f,(D_A+\mathbf{m})f\right)^{\mathrm{ch}}=\mathbf{m}\cdot(f,f)^{\mathrm{ch}}+\frac{\rho}{2}\cdot\left(f,\nabla_A^{*}\nabla_Af\right)^{\mathrm{ch}}\ge\mathbf{m}\cdot(f,f)^{\mathrm{ch}}\qquad(A\text{ real},\ f\in\mathrm{Hch}).
\end{equation}
\item[(iii)] Let $A$ be real and let $t'$ be a real number.
By Section~I.2.1, $\gamma_\mu=\mathrm{diag}(\sigma_\mu,-\sigma_\mu)$ and $\Sigma_3=\sigma_3\otimes1_2=\mathrm{diag}(1,1,-1,-1)$.
Hence $D_A$ maps the spinors with values in $\mathbb{C}^2\oplus0$ into themselves, and likewise the spinors with values in $0\oplus\mathbb{C}^2$; for $\mathrm{lv}\in\mathbb{C}^3$ the matrix $\gamma\cdot\mathrm{lv}$ acts on the first space as $\sigma\cdot\mathrm{lv}$ and on the second as $-\sigma\cdot\mathrm{lv}$.
By (1.1) the two restrictions of $D_A$ are the operators
\begin{equation*}
D_A^{\boldsymbol{\sigma}}:=\frac{\rho}{2}\cdot\nabla_A^{*}\nabla_A+\frac{1}{2}\cdot J_A^{\boldsymbol{\sigma}}\qquad\text{and}\qquad\frac{\rho}{2}\cdot\nabla_A^{*}\nabla_A-\frac{1}{2}\cdot J_A^{\boldsymbol{\sigma}}
\end{equation*}
on two-component spinors, where $J_A^{\boldsymbol{\sigma}}$ is formed as $J_A$ with $\sigma$ in the place of $\gamma$.
Since $\sigma_\mu$ is Hermitian, the computation of (ii) gives $\left(J_A^{\boldsymbol{\sigma}}\right)^{*}=-J_A^{\boldsymbol{\sigma}}$.
So the second restriction plus $t'$ is the adjoint of $D_A^{\boldsymbol{\sigma}}+t'$, and
\begin{equation}\tag{1.3}
\det(D_A+t')=\det\left(D_A^{\boldsymbol{\sigma}}+t'\right)\cdot\overline{\det\left(D_A^{\boldsymbol{\sigma}}+t'\right)}=\left|\det\left(D_A^{\boldsymbol{\sigma}}+t'\right)\right|^2\ge0.
\end{equation}
\item[(iv)] Let $A$ be complex.
For $f',f\in\mathrm{Hch}$ the substitution $x\mapsto x^{(i)}$ and the relation $U_A\left(x^{(i)},x\right)=U_{-A}\left(x,x^{(i)}\right)$ give
\begin{equation*}
\begin{aligned}
\left\langle f',M_i\cdot R_i(A)f\right\rangle^{\mathrm{ch}}&=\frac{1}{48}\sum_x\left(M_i(x)^T\cdot f'(x)\right)\cdot f\left(x^{(i)}\right)\cdot U_A\left(x,x^{(i)}\right)\\
&=\frac{1}{48}\sum_x\left(M_i\left(x^{(i)}\right)^T\cdot U_{-A}\left(x,x^{(i)}\right)\cdot f'\left(x^{(i)}\right)\right)\cdot f(x).
\end{aligned}
\end{equation*}
By (I.m2) at the site $x^{(i)}$ and by (I.s1), $\mathsf{C}^{-1}\cdot M_i\left(x^{(i)}\right)^T\cdot\mathsf{C}=M_i(x)$.
Hence $\left(M_i\cdot R_i(A)\right)^T=\mathsf{C}\cdot M_i\cdot R_i(-A)\cdot\mathsf{C}^{-1}$.
By (I.m1), $D_A=\frac{\rho}{2}\cdot\mathrm{sw}-\sum_iM_i\cdot R_i(A)$, where $\frac{\rho}{2}\cdot\mathrm{sw}$ is a scalar.
It follows that $D_A^T=\mathsf{C}\cdot D_{-A}\cdot\mathsf{C}^{-1}=\mathsf{C}^{-1}\cdot D_{-A}\cdot\mathsf{C}$, since $\mathsf{C}^{-1}=-\mathsf{C}$ by (I.c3).
In the basis of subsection~1.2 the bilinear form $\langle\,\cdot\,,\cdot\,\rangle^{\mathrm{ch}}$ has the matrix $\frac{1}{48}$ times the unit matrix; so the matrix of $O^T$ is the transpose of the matrix of $O$, and $\det O^T=\det O$.
Therefore, for every complex number $t'$,
\begin{equation}\tag{1.4}
\det\left(D_{-A}+t'\right)=\det\left(\mathsf{C}\cdot(D_A+t')^T\cdot\mathsf{C}^{-1}\right)=\det\left(D_A+t'\right)\qquad(A\text{ complex}).
\end{equation}
\end{itemize}

\paragraph{1.2. Proof of (a).}
\begin{proof}[Proof of Proposition~36\,(a)]
\emph{$N_0$ is entire.}
A basis of $\mathrm{Hch}$ is indexed by $c'\in\{0,\dots,n-1\}^3$, $g'\in O_{\mathrm{cub}}$ and $\beta\in\{1,\dots,4\}$: its element with these indices is the function with the value $(-1)^{z_1+z_2+z_3}\cdot\mathrm{ue}_\beta$ at the sites $(c'+n\cdot z,g')$, $z\in\mathbb{Z}^3$, and the value $0$ at all other sites.
By (0.1) with $B=B_t$, every entry of the matrix of $D^{\mathrm{ch}}[B_t]+\mathbf{m}$ in this basis is a finite sum of numbers that do not depend on $t$ and of products of a sign, an entry of $M_3(g)$ and a factor $\exp\left(\mathrm{i}\,t\cdot\mathrm{sv}\right)$ with a real number $\mathrm{sv}$, namely one of the numbers $\pm\eta\cdot a_\eta[\beta]_\nu(\eta\cdot c)$.
The determinant is a polynomial in the entries.
Hence $t\mapsto\det\left(D^{\mathrm{ch}}[B_t]+\mathbf{m}\right)$ is an entire function.

\emph{The denominator; positivity.}
For real $t$ the cochain $B_t$ is real.
By (1.2), $\left(D^{\mathrm{ch}}[B_t]+\mathbf{m}\right)f=0$ implies $f=0$; so $\det\left(D^{\mathrm{ch}}[B_t]+\mathbf{m}\right)\ne0$, and by (1.3) this determinant is $>0$.
In particular $\det\left(D^{\mathrm{ch}}[0]+\mathbf{m}\right)>0$.
Consequently $N_0$ is entire, $N_0(t)>0$ for real $t$, and $N_0(0)=1$ because $B_0=0$.

\emph{Evenness.}
One has $B_{-t}=-B_t$, and (1.4) gives $N_0(-t)=N_0(t)$ for all complex $t$.

\emph{The coefficients.}
As $N_0$ is even, $c_m=N_0^{(m)}(0)=0$ for odd $m$; as $N_0$ is real on the real axis, all $c_m$ are real.
Near $t=0$ one has $N_0(t)=1+\frac{c_2}{2}\cdot t^2+\frac{c_4}{24}\cdot t^4+O\left(t^6\right)$, and $\log(1+\mathrm{sv})=\mathrm{sv}-\frac{\mathrm{sv}^2}{2}+O\left(\mathrm{sv}^3\right)$ gives, for the branch with $\log1=0$,
\begin{equation*}
\log N_0(t)=\frac{c_2}{2}\cdot t^2+\left(\frac{c_4}{24}-\frac{c_2^2}{8}\right)\cdot t^4+O\left(t^6\right).
\end{equation*}
Hence $a_{\mathrm{ch}}(\eta,R)=c_2$ and $b_{\mathrm{ch}}(\eta,R)=24\cdot\left(\frac{c_4}{24}-\frac{c_2^2}{8}\right)=c_4-3\cdot c_2^2$, and both numbers are real.
\end{proof}

\subsubsection*{2. Clause (h)}

\paragraph{2.1. Lifts.}
Let $\hat{n}\in\mathbb{R}^3$ be a unit vector, $\mathrm{lv}\in\mathbb{C}^3$, and $S_{\mathrm{nv}}=(\gamma\cdot\hat{n})\cdot\gamma_5$ the matrix of Section~I.4.1.
By (I.c1), $(\gamma\cdot\hat{n})^2=1=\gamma_5^2$, so that $S_{\mathrm{nv}}^{-1}=\gamma_5\cdot(\gamma\cdot\hat{n})$.
As $\gamma_5$ anticommutes with $\gamma\cdot\mathrm{lv}$, it follows that
\begin{equation*}
\begin{aligned}
S_{\mathrm{nv}}\cdot(\gamma\cdot\mathrm{lv})\cdot S_{\mathrm{nv}}^{-1}&=-(\gamma\cdot\hat{n})\cdot(\gamma\cdot\mathrm{lv})\cdot(\gamma\cdot\hat{n})=-\left(2\,(\mathrm{lv}\cdot\hat{n})-(\gamma\cdot\mathrm{lv})\cdot(\gamma\cdot\hat{n})\right)\cdot(\gamma\cdot\hat{n})\\
&=\gamma\cdot\left(\mathrm{lv}-2\,(\mathrm{lv}\cdot\hat{n})\cdot\hat{n}\right).
\end{aligned}
\end{equation*}
For $\hat{n}=\hat{n}_i$ the right side is $\gamma\cdot(s_i\mathrm{lv})$ (Section~I.1.2).
Now let $g_{\mathrm{w}}=(a,r)\in W_{\mathrm{af}}$.
Since $s_0,s_1,s_2$ generate $O_{\mathrm{cub}}$ (Section~I.1.2), one can write $r=s_{i_1}s_{i_2}\cdots s_{i_j}$ with $i_1,\dots,i_j\in\{0,1,2\}$.
Let $S_{\mathrm{gw}}$ be the product, in this order, of the matrices $S_{\mathrm{nv}}$ formed with $\hat{n}=\hat{n}_{i_1},\hat{n}_{i_2},\dots,\hat{n}_{i_j}$ (Section~I.4.7), and $S_{\mathrm{gw}}:=1$ for the empty word.
Then
\begin{equation}\tag{2.1}
S_{\mathrm{gw}}\cdot(\gamma\cdot\mathrm{lv})\cdot S_{\mathrm{gw}}^{-1}=\gamma\cdot(r\,\mathrm{lv})\qquad\text{for all }\mathrm{lv}\in\mathbb{C}^3.
\end{equation}

\paragraph{2.2. Sites.}
Let $g_{\mathrm{w}}=(a,r)$ act on $\mathbb{Z}^3$ by $c\mapsto rc+a$.
The map $g_{\mathrm{w}}(c,g)=(rc+a,rg)$ is a bijection of $\mathbb{Z}^3\times O_{\mathrm{cub}}$ with the inverse $(c,g)\mapsto\left(r^{-1}(c-a),r^{-1}g\right)$, and $\pi(g_{\mathrm{w}}x)=g_{\mathrm{w}}\pi(x)$.
By (I.1.2), $g_{\mathrm{w}}(c,gs_i)=\left(rc+a,(rg)s_i\right)$ for $i\le2$ and $g_{\mathrm{w}}(c+ge_1,gs_3)=\left(rc+a+(rg)e_1,(rg)s_3\right)$; therefore $\left(g_{\mathrm{w}}x\right)^{(i)}=g_{\mathrm{w}}\left(x^{(i)}\right)$ for $i=0,\dots,3$.
Moreover $\hat{d}_i(g_{\mathrm{w}}x)=-(rg)\hat{n}_i=r\hat{d}_i(x)$.
Together with (2.1) this gives
\begin{equation}\tag{2.2}
S_{\mathrm{gw}}\cdot M_i(x)\cdot S_{\mathrm{gw}}^{-1}=\frac{w_i}{2}\cdot\left(\rho-\gamma\cdot\left(r\hat{d}_i(x)\right)\right)=M_i(g_{\mathrm{w}}x).
\end{equation}

\paragraph{2.3. Covariance.}
Let $A$ be a complex one-cochain.
For $f\colon\mathbb{Z}^3\times O_{\mathrm{cub}}\to\mathbb{C}^4$ put $\left(U_{\mathrm{gw}}f\right)(x):=S_{\mathrm{gw}}\cdot f\left(g_{\mathrm{w}}^{-1}x\right)$; then $\left(U_{\mathrm{gw}}^{-1}f\right)(x')=S_{\mathrm{gw}}^{-1}\cdot f(g_{\mathrm{w}}x')$.
Fix $x$ and let $x':=g_{\mathrm{w}}^{-1}x$; by subsection~2.2, $g_{\mathrm{w}}\left({x'}^{(i)}\right)=x^{(i)}$.
By (I.2.3) in lattice units,
\begin{equation*}
\begin{aligned}
\left(D_A\cdot U_{\mathrm{gw}}^{-1}f\right)(x')&=-\sum_iM_i(x')\cdot\left[U_A\left(x',{x'}^{(i)}\right)\cdot S_{\mathrm{gw}}^{-1}\cdot f\left(x^{(i)}\right)-S_{\mathrm{gw}}^{-1}\cdot f(x)\right],\\
\left(U_{\mathrm{gw}}\cdot D_A\cdot U_{\mathrm{gw}}^{-1}f\right)(x)&=-\sum_iS_{\mathrm{gw}}\cdot M_i(x')\cdot S_{\mathrm{gw}}^{-1}\cdot\left[U_A\left(x',{x'}^{(i)}\right)\cdot f\left(x^{(i)}\right)-f(x)\right]\\
&=-\sum_iM_i(x)\cdot\left[U_A\left(g_{\mathrm{w}}^{-1}x,g_{\mathrm{w}}^{-1}x^{(i)}\right)\cdot f\left(x^{(i)}\right)-f(x)\right],
\end{aligned}
\end{equation*}
the last equality following from (2.2).
For $i\le2$ one has $U_A\left(g_{\mathrm{w}}^{-1}x,g_{\mathrm{w}}^{-1}x^{(i)}\right)=1$.
For $i=3$, subsection~2.2 applied to $g_{\mathrm{w}}^{-1}$ gives $\pi\left(g_{\mathrm{w}}^{-1}x\right)=g_{\mathrm{w}}^{-1}c$ and $\pi\left(g_{\mathrm{w}}^{-1}x^{(3)}\right)=g_{\mathrm{w}}^{-1}c_3$, so that
\begin{equation*}
U_A\left(g_{\mathrm{w}}^{-1}x,g_{\mathrm{w}}^{-1}x^{(3)}\right)=\exp\left(\mathrm{i}\,q\cdot A\left(g_{\mathrm{w}}^{-1}c,g_{\mathrm{w}}^{-1}c_3\right)\right)=\exp\left(\mathrm{i}\,q\cdot\left((g_{\mathrm{w}})_{*}A\right)(c,c_3)\right)=U_{(g_{\mathrm{w}})_{*}A}\left(x,x^{(3)}\right),
\end{equation*}
with $\left((g_{\mathrm{w}})_{*}A\right)(c,c')=A\left(g_{\mathrm{w}}^{-1}c,g_{\mathrm{w}}^{-1}c'\right)$ as in Section~I.4.7.
Hence
\begin{equation}\tag{2.3}
U_{\mathrm{gw}}\cdot D_A\cdot U_{\mathrm{gw}}^{-1}=D_{(g_{\mathrm{w}})_{*}A}\qquad\text{on the functions }\mathbb{Z}^3\times O_{\mathrm{cub}}\to\mathbb{C}^4,\ \text{for every complex }A.
\end{equation}

\paragraph{2.4. The torus.}
The matrix $r$ is a signed permutation matrix: $r^{-1}e_\mu=\pm e_\nu$ for some $\nu$, and $r\left(n\mathbb{Z}^3\right)=n\mathbb{Z}^3$.
So $(g_{\mathrm{w}})_{*}A$ is $n\mathbb{Z}^3$-periodic if $A$ is, because
\begin{equation*}
\left((g_{\mathrm{w}})_{*}A\right)(c+n\cdot z,c'+n\cdot z)=A\left(g_{\mathrm{w}}^{-1}c+n\cdot r^{-1}z,\ g_{\mathrm{w}}^{-1}c'+n\cdot r^{-1}z\right).
\end{equation*}
Every $f\in\mathrm{Hch}$ satisfies also $f(c-n\cdot e_\nu,g)=-f(c,g)$, and so
\begin{equation*}
\left(U_{\mathrm{gw}}f\right)(c+n\cdot e_\mu,g)=S_{\mathrm{gw}}\cdot f\left(r^{-1}(c-a)+n\cdot r^{-1}e_\mu,\ r^{-1}g\right)=-\left(U_{\mathrm{gw}}f\right)(c,g).
\end{equation*}
Thus $U_{\mathrm{gw}}$ maps $\mathrm{Hch}$ into $\mathrm{Hch}$.
The same holds for $U_{\mathrm{gw}}^{-1}$, since $\left(U_{\mathrm{gw}}^{-1}f\right)(c+n\cdot e_\mu,g)=S_{\mathrm{gw}}^{-1}\cdot f\left(rc+a+n\cdot re_\mu,\ rg\right)$; so $U_{\mathrm{gw}}$ is a linear bijection of $\mathrm{Hch}$.
By (2.3), $D_{(g_{\mathrm{w}})_{*}A}+\boldsymbol{\mu}=U_{\mathrm{gw}}\cdot\left(D_A+\boldsymbol{\mu}\right)\cdot U_{\mathrm{gw}}^{-1}$ on $\mathrm{Hch}$ for every number $\boldsymbol{\mu}$, and the linear maps $O$ and $U_{\mathrm{gw}}\cdot O\cdot U_{\mathrm{gw}}^{-1}$ have the same determinant.
Therefore
\begin{equation}\tag{2.4}
\det\left(D_{(g_{\mathrm{w}})_{*}A}+\boldsymbol{\mu}\right)=\det\left(D_A+\boldsymbol{\mu}\right)\qquad\text{on }\mathfrak{S}^{\mathrm{ch}}(1;n).
\end{equation}
This is the first assertion of (h).

\paragraph{2.5. Translates of the probe.}
Let $v\in\eta\mathbb{Z}^3$.
Then $(\tau_v\beta)(y)=\beta(y-v)$ (Section~I.7.5; \cite[Section~7.1.3]{QEDI}), and $\tau_v\beta$ is a probe.
The operators $J_\eta$ and $\delta_{\eta,\mu}$ of Sections~I.7.6 and~G5.1.3 commute with translations; hence $\mathrm{gs}_\nu[\tau_v\beta](y)=\mathrm{gs}_\nu[\beta](y-v)$.
By the first assertion of Proposition~35.3\,(b1), for $c\in\eta\mathbb{Z}^3$,
\begin{equation*}
a_\eta[\tau_v\beta]_\nu(c)=\sum_{z\in\mathbb{Z}^3}\mathrm{gs}_\nu[\beta]\left(c-v+\frac{\eta}{2}\cdot e_\nu+R\cdot z\right)=a_\eta[\beta]_\nu(c-v).
\end{equation*}
In lattice units (Section~G5.1.3), for a probe $\beta'$ the cochain $\eta^{\frac{1}{2}}\cdot a_\eta[\beta']$ has the value $\eta^{\frac{1}{2}}\cdot a_\eta[\beta']_\nu(\eta\cdot c)$ on the bond $(c,c+e_\nu)$, $c\in\mathbb{Z}^3$.
Let $g_{\mathrm{w}}=(v/\eta,1)$ be the translation by $v/\eta$ in $\mathbb{Z}^3$.
Then
\begin{equation*}
\left((g_{\mathrm{w}})_{*}\left(\eta^{\frac{1}{2}}\cdot a_\eta[\beta]\right)\right)(c,c+e_\nu)=\eta^{\frac{1}{2}}\cdot a_\eta[\beta]_\nu\left(\eta\cdot\left(c-\frac{v}{\eta}\right)\right)=\eta^{\frac{1}{2}}\cdot a_\eta[\tau_v\beta]_\nu(\eta\cdot c).
\end{equation*}
This is the second assertion of (h).

\paragraph{2.6. Conclusion.}
\begin{proof}[Proof of Proposition~36\,(h)]
The first two assertions of (h) are (2.4) and the identity of subsection~2.5.
Let $g_{\mathrm{w}}\in W_{\mathrm{af}}$, $\hat{\mathfrak{a}}=\eta^{\frac{1}{2}}\cdot a_\eta[\beta]$, $e>0$, $q=e\cdot\eta^{\frac{1}{2}}$ and $t\in\mathbb{C}$.
The map $(g_{\mathrm{w}})_{*}$ is linear, and (2.4) with $A=\frac{t}{e}\cdot\hat{\mathfrak{a}}$ and $\boldsymbol{\mu}=\mathbf{m}$ gives
\begin{equation*}
\det\left(D_{(t/e)\cdot(g_{\mathrm{w}})_{*}\hat{\mathfrak{a}}}+\mathbf{m}\right)=\det\left(D_{(t/e)\cdot\hat{\mathfrak{a}}}+\mathbf{m}\right).
\end{equation*}
The denominator of $N_0$ does not contain the probe.
Hence $N_0$ (subsection~0.2), and with it $c_m=N_0^{(m)}(0)$, $a_{\mathrm{ch}}(\eta,R)$ and $b_{\mathrm{ch}}(\eta,R)$, remain unchanged if $\hat{\mathfrak{a}}$ is replaced by $(g_{\mathrm{w}})_{*}\hat{\mathfrak{a}}$.
By subsection~2.5, the replacement of $\beta$ by $\tau_v\beta$ with $v\in\eta\mathbb{Z}^3$ is a special case of this.
\end{proof}

\subsubsection*{3. Clause (b)}

\paragraph{3.1. The derivatives of $D^{\mathrm{ch}}[B_t]$.}
For $c\in\mathbb{Z}^3$ and $\nu=1,2,3$ let
\[
\mathrm{bq}_\nu(c):=\eta\,a_{\eta}[\beta]_\nu(\eta c).
\]
Since $a_{\eta}[\beta]$ is a real function on $\Lambda_{\eta,R}=\eta\mathbb{Z}^3/R\mathbb{Z}^3$ and $R=n\eta$, the function $\mathrm{bq}_\nu$ is real and $n\mathbb{Z}^3$-periodic.
The definition of $B_t$ in the setting (FF) of Section~G5.6.0, together with $B_t(c',c)=-B_t(c,c')$, gives
\[
B_t(c,c+e_\nu)=t\,\mathrm{bq}_\nu(c),\qquad B_t(c,c-e_\nu)=-t\,\mathrm{bq}_\nu(c-e_\nu).
\]
For every $g\in O_{\mathrm{cub}}$ the vector $ge_1$ is exactly one of the six vectors $\pm e_\nu$, and $M_3(g)=3(\rho\mp\gamma_\nu)$ if $ge_1=\pm e_\nu$.
On $\mathrm{Hch}$ define
\begin{align*}
\left(\mathrm{Hhp}^{\pm}_\nu f\right)(c,g)&:=\mathbf{1}[ge_1=\pm e_\nu]\,3(\rho\mp\gamma_\nu)\,f(c\pm e_\nu,gs_3),\\
(\mathrm{Dloc}\,f)(c,g)&:=\mathbf{m}\,f(c,g)-\sum_{i\le2}M_i(g)\left[f(c,gs_i)-f(c,g)\right]+M_3(g)\,f(c,g);
\end{align*}
a function of $c$ acts on $\mathrm{Hch}$ by multiplication, in accordance with (I.POS).
In particular
\[
\left(\mathrm{Hhp}^{-}_\nu\exp(-it\,\mathrm{bq}_\nu)f\right)(c,g)=\mathbf{1}[ge_1=-e_\nu]\,3(\rho+\gamma_\nu)\exp\left(-it\,\mathrm{bq}_\nu(c-e_\nu)\right)f(c-e_\nu,gs_3),
\]
and (0.1) takes the form
\begin{equation}
\mathrm{Dtt}(t):=D^{\mathrm{ch}}[B_t]+\mathbf{m}=\mathrm{Dloc}-\sum_\nu\left(\exp(it\,\mathrm{bq}_\nu)\,\mathrm{Hhp}^{+}_\nu+\mathrm{Hhp}^{-}_\nu\exp(-it\,\mathrm{bq}_\nu)\right)\qquad(t\in\mathbb{C}).
\tag{3.1}
\end{equation}
The operators $\mathrm{Dloc}$ and $\mathrm{Hhp}^{\pm}_\nu$ do not depend on $t$.
By clause (a) the inverse $\mathrm{G0}:=\mathrm{Dtt}(0)^{-1}=\left(D^{\mathrm{ch}}[0]+\mathbf{m}\right)^{-1}$ exists, and differentiation of (3.1) at $t=0$ gives
\begin{equation}
\begin{split}
\mathrm{Vd}_1&:=\mathrm{Dtt}'(0)=-i\sum_\nu\left(\mathrm{bq}_\nu\cdot\mathrm{Hhp}^{+}_\nu-\mathrm{Hhp}^{-}_\nu\cdot\mathrm{bq}_\nu\right),\\
\mathrm{Vd}_2&:=\mathrm{Dtt}''(0)=\sum_\nu\left(\mathrm{bq}_\nu^2\cdot\mathrm{Hhp}^{+}_\nu+\mathrm{Hhp}^{-}_\nu\cdot\mathrm{bq}_\nu^2\right).
\end{split}
\tag{3.2}
\end{equation}
These are the operators $\mathrm{G0}$, $\mathrm{Vd}_1$, $\mathrm{Vd}_2$ of Proposition~37.2.

\paragraph{3.2. Jacobi's formula.}
Let $O(t)$ be a holomorphic family of invertible linear maps of a space of finite dimension, and let $O'(t):=\frac{\mathrm{d}}{\mathrm{d}t}O(t)$.
One has
\[
\det O(t+\mathrm{sv})=\det O(t)\,\det\left(1+O(t)^{-1}\left(O(t+\mathrm{sv})-O(t)\right)\right),
\]
and, for a linear map $O_1$,
\[
\det\left(1+\mathrm{sv}\,O_1+(\text{terms of order }\mathrm{sv}^2)\right)=1+\mathrm{sv}\operatorname{Tr}O_1+(\text{terms of order }\mathrm{sv}^2).
\]
Applied to $O_1=O(t)^{-1}O'(t)$, these two facts give
\[
\frac{\mathrm{d}}{\mathrm{d}t}\det O(t)=\det O(t)\,\operatorname{Tr}\left(O(t)^{-1}O'(t)\right),
\]
while differentiation of $OO^{-1}=1$ gives $\left(O^{-1}\right)'=-O^{-1}O'O^{-1}$.
The operator $\mathrm{Dtt}(t)$ is invertible for $t$ near $0$.
Since $N_0(t)=\det\mathrm{Dtt}(t)/\det\mathrm{Dtt}(0)$, one has there
\[
\frac{\mathrm{d}}{\mathrm{d}t}\log N_0(t)=\operatorname{Tr}\left(\mathrm{Dtt}(t)^{-1}\cdot\mathrm{Dtt}'(t)\right),\qquad
\frac{\mathrm{d}^2}{\mathrm{d}t^2}\log N_0(t)=\operatorname{Tr}\left(\mathrm{Dtt}^{-1}\cdot\mathrm{Dtt}''\right)-\operatorname{Tr}\left(\mathrm{Dtt}^{-1}\cdot\mathrm{Dtt}'\cdot\mathrm{Dtt}^{-1}\cdot\mathrm{Dtt}'\right).
\]
At $t=0$ this becomes
\begin{equation}
a_{\mathrm{ch}}(\eta,R)=\operatorname{Tr}\left[\mathrm{G0}\,\mathrm{Vd}_2\right]-\operatorname{Tr}\left[\mathrm{G0}\,\mathrm{Vd}_1\,\mathrm{G0}\,\mathrm{Vd}_1\right].
\tag{3.3}
\end{equation}

\paragraph{3.3. Fourier decomposition of $\mathrm{Hch}$.}
Let
\[
\mathrm{Pm}_n:=\frac{2\pi}{n}\left(\mathbb{Z}+\tfrac12\right)^3,\qquad
\mathrm{Kt}_n:=\frac{2\pi}{n}\mathbb{Z}^3,\qquad
\mathrm{Kl}_n:=\left\{\frac{2\pi}{n}j:\ j\in\mathbb{Z}^3,\ -\frac{n}{2}<j_\nu\le\frac{n}{2}\right\}=\eta\,\mathrm{Kx}_{\eta,R}.
\]
Then $\mathrm{Pm}_n+\mathrm{Kt}_n=\mathrm{Pm}_n$.
The set $P_{n,0}=\left\{p:\ p_\mu=(2\pi j_\mu+\pi)/n,\ j\in\{0,\dots,n-1\}^3\right\}$ of Section~I.2.6 is a set of representatives of $\mathrm{Pm}_n$ modulo $2\pi\mathbb{Z}^3$, and $\mathrm{Kl}_n$ is a set of representatives of $\mathrm{Kt}_n$ modulo $2\pi\mathbb{Z}^3$; in both cases $j_\nu$ ranges, in each coordinate, over $n$ consecutive integers.
For $p\in\mathrm{Pm}_n$ and $u\in V_{\mathrm{v}}$ let
\[
\mathrm{ef}_p(c):=\exp(ip\cdot c),\qquad\left(\mathrm{ef}_p\otimes u\right)(c,g):=\mathrm{ef}_p(c)\,u(g).
\]
As $\exp(inp_\mu)=-1$, the function $\mathrm{ef}_p\otimes u$ lies in $\mathrm{Hch}$; moreover $\mathrm{ef}_p$ depends on $p$ only modulo $2\pi\mathbb{Z}^3$.
If $p,p''\in P_{n,0}$ have the integer vectors $j,j''$, then $|j_\nu-j''_\nu|\le n-1$, and therefore
\begin{equation}
n^{-3}\sum_{c\in\{0,\dots,n-1\}^3}\exp\left(i(p-p'')\cdot c\right)=\prod_\nu n^{-1}\sum_{c_\nu=0}^{n-1}\exp\left(\frac{2\pi i\,(j_\nu-j''_\nu)\,c_\nu}{n}\right)=\mathbf{1}[p=p''].
\tag{3.4}
\end{equation}
The linear map $(u_p)_{p\in P_{n,0}}\mapsto\sum_p\mathrm{ef}_p\otimes u_p$ from the direct sum of $n^3$ copies of $V_{\mathrm{v}}$ to $\mathrm{Hch}$ is injective: if the sum vanishes, multiplication by $\exp(-ip''\cdot c)$ and summation over $c\in\{0,\dots,n-1\}^3$ give $u_{p''}=0$ by (3.4).
Both spaces have the dimension $192n^3$.
Hence
\[
\mathrm{Hch}=\bigoplus_{p\in P_{n,0}}\mathrm{ef}_p\otimes V_{\mathrm{v}}.
\]

\emph{Blocks.}
For a linear map $O$ of $\mathrm{Hch}$ and $p,p''\in P_{n,0}$ let $O[p'',p]$ be the linear map of $V_{\mathrm{v}}$ determined by
\[
O\left(\mathrm{ef}_p\otimes u\right)=\sum_{p''\in P_{n,0}}\mathrm{ef}_{p''}\otimes O[p'',p]\,u\qquad(u\in V_{\mathrm{v}}).
\]
For $p,p''\in\mathrm{Pm}_n$ the block $O[p'',p]$ is defined through the representatives in $P_{n,0}$; it is thus $2\pi\mathbb{Z}^3$-periodic in each of its two arguments.
Linear algebra gives
\begin{equation}
\operatorname{Tr}O=\sum_{p\in P_{n,0}}\operatorname{tr}O[p,p],\qquad\left(OO'\right)[p'',p]=\sum_{p'''\in P_{n,0}}O[p'',p''']\,O'[p''',p].
\tag{3.5}
\end{equation}

\emph{Fourier coefficients of the probe.}
For $p'\in\mathrm{Kt}_n$ put
\[
\mathrm{bqh}_\nu(p'):=n^{-3}\sum_{c\in\{0,\dots,n-1\}^3}\mathrm{bq}_\nu(c)\exp(-ip'\cdot c).
\]
This function is $2\pi\mathbb{Z}^3$-periodic in $p'$, and $\mathrm{bqh}_\nu(-p')=\overline{\mathrm{bqh}_\nu(p')}$ because $\mathrm{bq}_\nu$ is real.
For $p'=\frac{2\pi}{n}j$ and $p''=\frac{2\pi}{n}j''$ in $\mathrm{Kl}_n$ one has $|j_\nu-j''_\nu|\le n-1$, and the geometric sums of (3.4) give
\[
n^{-3}\sum_{c\in\{0,\dots,n-1\}^3}\exp\left(i(p'-p'')\cdot c\right)=\mathbf{1}[p'=p''].
\]
Consequently the $n^3$ functions $c\mapsto\exp(ip'\cdot c)$, $p'\in\mathrm{Kl}_n$, form a basis of the space of $n\mathbb{Z}^3$-periodic functions on $\mathbb{Z}^3$ which is orthonormal for the inner product $(f_1,f_2)\mapsto n^{-3}\sum_{c\in\{0,\dots,n-1\}^3}\overline{f_1(c)}\,f_2(c)$, and
\begin{equation}
\mathrm{bq}_\nu(c)=\sum_{p'\in\mathrm{Kl}_n}\mathrm{bqh}_\nu(p')\exp(ip'\cdot c),\qquad
n^{-3}\sum_{c\in\{0,\dots,n-1\}^3}\mathrm{bq}_\nu(c)^2=\sum_{p'\in\mathrm{Kl}_n}\left|\mathrm{bqh}_\nu(p')\right|^2.
\tag{3.6}
\end{equation}

\paragraph{3.4. The blocks of the factors.}
Let $p\in\mathrm{Pm}_n$ and $u\in V_{\mathrm{v}}$.

(i) One has $\left(D^{\mathrm{ch}}[0]+\mathbf{m}\right)\left(\mathrm{ef}_p\otimes u\right)=\mathrm{ef}_p\otimes D^{(p)}u$: this is the definition (I.2.6) of the symbol, taken at $\boldsymbol{\mu}=\mathbf{m}$.
Hence $\mathrm{G0}[p'',p]=\mathbf{1}\left[p''\equiv p\bmod 2\pi\mathbb{Z}^3\right]\mathrm{Rsy}(p)$.

(ii) With $\mathrm{Ehp}^{\pm}_\nu$ as in Section~G5.6.0,
\[
\left(\mathrm{Hhp}^{\pm}_\nu\left(\mathrm{ef}_p\otimes u\right)\right)(c,g)=\mathbf{1}[ge_1=\pm e_\nu]\,3(\rho\mp\gamma_\nu)\exp\left(ip\cdot(c\pm e_\nu)\right)u(gs_3)=\exp(\pm ip_\nu)\,\mathrm{ef}_p(c)\left(\mathrm{Ehp}^{\pm}_\nu u\right)(g).
\]
Therefore $\mathrm{Hhp}^{\pm}_\nu[p'',p]=\mathbf{1}\left[p''\equiv p\bmod 2\pi\mathbb{Z}^3\right]\exp(\pm ip_\nu)\,\mathrm{Ehp}^{\pm}_\nu$.

(iii) Let $\mathrm{fs}$ be an $n\mathbb{Z}^3$-periodic function, $\mathrm{fs}(c)=\sum_{p'\in\mathrm{Kl}_n}\mathrm{fsh}(p')\exp(ip'\cdot c)$.
Then $\mathrm{fs}\left(\mathrm{ef}_p\otimes u\right)=\sum_{p'\in\mathrm{Kl}_n}\mathrm{fsh}(p')\,\mathrm{ef}_{p+p'}\otimes u$.
So, for $p'\in\mathrm{Kt}_n$, the block $\mathrm{fs}[p+p',p]$ is $\mathrm{fsh}(p')$ times the identity of $V_{\mathrm{v}}$, where $\mathrm{fsh}$ is continued periodically.

\paragraph{3.5. The blocks of $\mathrm{Vd}_1$ and $\mathrm{Vd}_2$.}
Let $p\in\mathrm{Pm}_n$ and $p'\in\mathrm{Kt}_n$.
From (3.2), (3.5) and 3.4,
\[
\left(\mathrm{bq}_\nu\cdot\mathrm{Hhp}^{+}_\nu\right)[p+p',p]=\mathrm{bqh}_\nu(p')\exp(ip_\nu)\,\mathrm{Ehp}^{+}_\nu,\qquad
\left(\mathrm{Hhp}^{-}_\nu\cdot\mathrm{bq}_\nu\right)[p+p',p]=\exp\left(-i(p_\nu+p'_\nu)\right)\mathrm{Ehp}^{-}_\nu\,\mathrm{bqh}_\nu(p').
\]
Hence
\begin{equation}
\begin{split}
\mathrm{Vd}_1[p+p',p]&=-i\sum_\nu\mathrm{bqh}_\nu(p')\left(\exp(ip_\nu)\,\mathrm{Ehp}^{+}_\nu-\exp\left(-i(p_\nu+p'_\nu)\right)\mathrm{Ehp}^{-}_\nu\right)\\
&=\sum_\nu\mathrm{bqh}_\nu(p')\exp\left(-\tfrac{i}{2}p'_\nu\right)\mathrm{Vx}_\nu(p+p'/2)\qquad\left(p\in\mathrm{Pm}_n,\ p'\in\mathrm{Kt}_n\right),
\end{split}
\tag{3.7}
\end{equation}
where the second equality follows from Proposition~37.1\,(d), which gives
\[
\mathrm{Vx}_\nu(p+p'/2)=-i\left(\exp\left(i\left(p_\nu+\tfrac12p'_\nu\right)\right)\mathrm{Ehp}^{+}_\nu-\exp\left(-i\left(p_\nu+\tfrac12p'_\nu\right)\right)\mathrm{Ehp}^{-}_\nu\right).
\]
The middle member of (3.7) is $2\pi\mathbb{Z}^3$-periodic in $p$ and in $p'$, as the left member is by definition; in the right member the two factors $\exp\left(-\tfrac{i}{2}p'_\nu\right)$ and $\mathrm{Vx}_\nu(p+p'/2)$ both change their sign under $p'_\nu\mapsto p'_\nu+2\pi$.
Since (3.7) holds for all pairs in $\mathrm{Pm}_n\times\mathrm{Kt}_n$, it may be applied to the pair $(p+p',-p')$; with $\mathrm{bqh}_\mu(-p')=\overline{\mathrm{bqh}_\mu(p')}$ and $(p+p')-p'/2=p+p'/2$ it yields
\begin{equation}
\mathrm{Vd}_1[p,p+p']=\sum_\mu\overline{\mathrm{bqh}_\mu(p')\exp\left(-\tfrac{i}{2}p'_\mu\right)}\,\mathrm{Vx}_\mu(p+p'/2).
\tag{3.8}
\end{equation}
The Fourier coefficient of the periodic function $\mathrm{bq}_\nu^2$ at $p'=0$ is $n^{-3}\sum_{c\in\{0,\dots,n-1\}^3}\mathrm{bq}_\nu(c)^2$.
By (3.2), 3.4, (3.6) and the relation $\partial_\nu^2D^{(p)}=\exp(ip_\nu)\,\mathrm{Ehp}^{+}_\nu+\exp(-ip_\nu)\,\mathrm{Ehp}^{-}_\nu$ of Proposition~37.1\,(d),
\begin{equation}
\mathrm{Vd}_2[p,p]=\sum_\nu\left(\sum_{p'\in\mathrm{Kl}_n}\left|\mathrm{bqh}_\nu(p')\right|^2\right)\partial_\nu^2D^{(p)}.
\tag{3.9}
\end{equation}

\paragraph{3.6. The two traces.}
By (3.5), 3.4\,(i) and (3.9),
\[
\operatorname{Tr}\left[\mathrm{G0}\,\mathrm{Vd}_2\right]=\sum_{p'\in\mathrm{Kl}_n}\sum_\nu\left|\mathrm{bqh}_\nu(p')\right|^2\sum_{p\in P_{n,0}}\operatorname{tr}\left[\mathrm{Rsy}(p)\,\partial_\nu^2D^{(p)}\right],
\]
and by (3.5) and 3.4\,(i),
\[
\operatorname{Tr}\left[\mathrm{G0}\,\mathrm{Vd}_1\,\mathrm{G0}\,\mathrm{Vd}_1\right]=\sum_{p\in P_{n,0}}\sum_{p'''\in P_{n,0}}\operatorname{tr}\left[\mathrm{Rsy}(p)\cdot\mathrm{Vd}_1[p,p''']\cdot\mathrm{Rsy}(p''')\cdot\mathrm{Vd}_1[p''',p]\right].
\]
For fixed $p$, the map that sends $p'\in\mathrm{Kl}_n$ to the representative of $p+p'$ in $P_{n,0}$ is a bijection of $\mathrm{Kl}_n$ onto $P_{n,0}$, and $\mathrm{Rsy}$ and the blocks are periodic.
Together with (3.7) and (3.8) this gives
\begin{multline*}
\operatorname{Tr}\left[\mathrm{G0}\,\mathrm{Vd}_1\,\mathrm{G0}\,\mathrm{Vd}_1\right]=\sum_{p'\in\mathrm{Kl}_n}\sum_{\mu,\nu}\overline{\mathrm{bqh}_\mu(p')\exp\left(-\tfrac{i}{2}p'_\mu\right)}\,\mathrm{bqh}_\nu(p')\exp\left(-\tfrac{i}{2}p'_\nu\right)\\
\times\sum_{p\in P_{n,0}}\operatorname{tr}\left[\mathrm{Rsy}(p)\cdot\mathrm{Vx}_\mu(p+p'/2)\cdot\mathrm{Rsy}(p+p')\cdot\mathrm{Vx}_\nu(p+p'/2)\right].
\end{multline*}
Recall from the definitions in Proposition~36 that
\[
\mathrm{lp}_{\mu\nu}(p,p')=\delta_{\mu\nu}\operatorname{tr}\left[\mathrm{Rsy}(p)\,\partial_\nu^2D^{(p)}\right]-\operatorname{tr}\left[\mathrm{Rsy}(p)\cdot\mathrm{Vx}_\mu(p+p'/2)\cdot\mathrm{Rsy}(p+p')\cdot\mathrm{Vx}_\nu(p+p'/2)\right].
\]
Since $\left|\mathrm{bqh}_\nu(p')\exp\left(-\tfrac{i}{2}p'_\nu\right)\right|^2=\left|\mathrm{bqh}_\nu(p')\right|^2$, it follows from (3.3) that
\begin{equation}
a_{\mathrm{ch}}(\eta,R)=\sum_{p'\in\mathrm{Kl}_n}\sum_{\mu,\nu}\overline{\mathrm{bqh}_\mu(p')\exp\left(-\tfrac{i}{2}p'_\mu\right)}\,\mathrm{bqh}_\nu(p')\exp\left(-\tfrac{i}{2}p'_\nu\right)\sum_{p\in P_{n,0}}\mathrm{lp}_{\mu\nu}(p,p').
\tag{3.10}
\end{equation}

\paragraph{3.7. Physical units.}
By the first sentence of Proposition~35.3\,(b1) and the definition of $\mathrm{a0}_{\eta,\nu}$ in Proposition~36,
\[
a_{\eta}[\beta]_\nu(y)=\sum_{z\in\mathbb{Z}^3}\mathrm{a0}_{\eta,\nu}(y+Rz)\qquad(y\in\eta\mathbb{Z}^3);
\]
the sum is finite, and no condition on $R$ is needed.
Let $p'\in\mathrm{Kt}_n$.
The map $(c,z)\mapsto c+nz$ is a bijection of $\{0,\dots,n-1\}^3\times\mathbb{Z}^3$ onto $\mathbb{Z}^3$; moreover $\eta c+Rz=\eta(c+nz)$ and $\exp\left(-ip'\cdot(c+nz)\right)=\exp(-ip'\cdot c)$.
So
\begin{align*}
\mathrm{bqh}_\nu(p')&=n^{-3}\eta\sum_{c\in\{0,\dots,n-1\}^3}\sum_{z\in\mathbb{Z}^3}\mathrm{a0}_{\eta,\nu}\left(\eta(c+nz)\right)\exp(-ip'\cdot c)\\
&=n^{-3}\eta\sum_{c\in\mathbb{Z}^3}\mathrm{a0}_{\eta,\nu}(\eta c)\exp(-ip'\cdot c).
\end{align*}
Write $y=\eta c$ and $\mathrm{kph}:=p'/\eta$.
Then $p'\cdot c=\mathrm{kph}\cdot y$, $\tfrac12p'_\nu=\mathrm{kph}\cdot\left(\tfrac{\eta}{2}e_\nu\right)$ and $n^{-3}\eta^{-3}=R^{-3}$, and the definition of $\mathrm{ahl}_{\eta,\nu}$ gives
\begin{equation}
\begin{split}
\mathrm{bqh}_\nu(p')\exp\left(-\tfrac{i}{2}p'_\nu\right)&=n^{-3}\eta\sum_{y\in\eta\mathbb{Z}^3}\mathrm{a0}_{\eta,\nu}(y)\exp\left(-i\,\mathrm{kph}\cdot\left(y+\tfrac{\eta}{2}e_\nu\right)\right)\\
&=n^{-3}\eta^{-2}\,\mathrm{ahl}_{\eta,\nu}(\mathrm{kph})=\eta R^{-3}\,\mathrm{ahl}_{\eta,\nu}(\mathrm{kph}).
\end{split}
\tag{3.11}
\end{equation}
The vector $p'$ ranges over $\mathrm{Kl}_n$ if and only if $\mathrm{kph}=p'/\eta$ ranges over $\mathrm{Kx}_{\eta,R}$.
Inserting (3.11) into (3.10) and using $\eta^2R^{-6}=R^{-3}\eta^{-1}(\eta/R)^3=R^{-3}\eta^{-1}n^{-3}$, one obtains
\[
a_{\mathrm{ch}}(\eta,R)=R^{-3}\sum_{\mathrm{kph}\in\mathrm{Kx}_{\eta,R}}\sum_{\mu,\nu}\overline{\mathrm{ahl}_{\eta,\mu}(\mathrm{kph})}\cdot\mathrm{ahl}_{\eta,\nu}(\mathrm{kph})\,\eta^{-1}n^{-3}\sum_{p\in P_{n,0}}\mathrm{lp}_{\mu\nu}(p,\eta\,\mathrm{kph}).
\]
The product of the last three factors is $\mathrm{Vpl}^{\eta,R}_{\mu\nu}(\mathrm{kph})$, and the identity obtained is clause (b).
It holds for every $\eta\in(0,1]$ and every $R>0$ for which $n=R/\eta$ is an integer $\ge2$.\qed

\subsubsection*{4. Clause (c)}

In this section $\bar m\in(0,1]$, $\eta\in(0,1]$, $\mathbf{m}=\bar m\eta$, $\mu,\nu\in\{1,2,3\}$ and a vector $p'\in\mathbb{R}^3$ are fixed unless the contrary is stated, and all bounds are uniform in them.
The number $R$ appears only in 4.1 and 4.3, through $n=R/\eta$.

\paragraph{4.1. A mean over $P_{n,0}$ and the integral.}
Let $\mathrm{fs}$ be a $2\pi\mathbb{Z}^3$-periodic function of class $C^4$ on $\mathbb{R}^3$, and let
\[
\mathrm{fh}(z)=\int_{[-\pi,\pi]^3}\frac{\mathrm{d}^3p}{(2\pi)^3}\,\mathrm{fs}(p)\exp(-ip\cdot z)\qquad(z\in\mathbb{Z}^3),
\]
as in Section~I.2.6.
Let $z\neq0$, and choose $\mathfrak{l}$ with $|z_{\mathfrak{l}}|=|z|_\infty$.
Four integrations by parts in the variable $p_{\mathfrak{l}}$, without boundary terms because of the periodicity, give
\[
z_{\mathfrak{l}}^4\,\mathrm{fh}(z)=\int_{[-\pi,\pi]^3}\frac{\mathrm{d}^3p}{(2\pi)^3}\left(\partial_{\mathfrak{l}}^4\mathrm{fs}\right)(p)\exp(-ip\cdot z).
\]
Hence
\begin{equation}
|\mathrm{fh}(z)|\le|z|_\infty^{-4}\max_{\mathfrak{l}}\int_{[-\pi,\pi]^3}\frac{\mathrm{d}^3p}{(2\pi)^3}\left|\partial_{\mathfrak{l}}^4\mathrm{fs}\right|\qquad(z\neq0).
\tag{4.1}
\end{equation}
For $j\ge1$ one has $\#\left|\left\{z\in\mathbb{Z}^3:\ |z|_\infty=j\right\}\right|=(2j+1)^3-(2j-1)^3=24j^2+2$, and therefore
\begin{equation}
\sum_{z\neq0}|z|_\infty^{-4}=\sum_{j\ge1}\left(24j^{-2}+2j^{-4}\right)=4\pi^2+\frac{\pi^4}{45}<42;
\tag{4.2}
\end{equation}
here Euler's values $\sum_{j\ge1}j^{-2}=\pi^2/6$ and $\sum_{j\ge1}j^{-4}=\pi^4/90$ are used, and the last inequality is (Ef0)\,(ii).
Consequently $\sum_z|\mathrm{fh}(z)|<\infty$.
The continuous periodic function $\sum_z\mathrm{fh}(z)\exp(ip\cdot z)$ has the same Fourier coefficients as $\mathrm{fs}$.
The difference of the two functions is a continuous periodic function that is orthogonal to all trigonometric polynomials; these are uniformly dense in the continuous periodic functions by the Stone--Weierstrass theorem, so the difference vanishes, that is, $\mathrm{fs}(p)=\sum_z\mathrm{fh}(z)\exp(ip\cdot z)$ for all $p$.
For $z\in\mathbb{Z}^3$,
\[
n^{-3}\sum_{p\in P_{n,0}}\exp(ip\cdot z)=\prod_\nu\exp\left(\frac{i\pi z_\nu}{n}\right)n^{-1}\sum_{j=0}^{n-1}\exp\left(\frac{2\pi i\,jz_\nu}{n}\right)=\prod_\nu\mathbf{1}[z_\nu\in n\mathbb{Z}]\,(-1)^{\frac{z_\nu}{n}}.
\]
Therefore $n^{-3}\sum_{p\in P_{n,0}}\mathrm{fs}(p)=\sum_{z\in\mathbb{Z}^3}(-1)^{z_1+z_2+z_3}\,\mathrm{fh}(nz)$, and the term $z=0$ is the integral of $\mathrm{fs}$.
By (4.1) and (4.2),
\begin{equation}
\left|n^{-3}\sum_{p\in P_{n,0}}\mathrm{fs}(p)-\int_{[-\pi,\pi]^3}\frac{\mathrm{d}^3p}{(2\pi)^3}\,\mathrm{fs}(p)\right|\le42\,n^{-4}\max_{\mathfrak{l}}\int_{[-\pi,\pi]^3}\frac{\mathrm{d}^3p}{(2\pi)^3}\left|\partial_{\mathfrak{l}}^4\mathrm{fs}\right|.
\tag{4.3}
\end{equation}

\paragraph{4.2. The fourth derivatives of $\mathrm{lp}$.}
By (Ef5) the function $p\mapsto\mathrm{lp}_{\mu\nu}(p,p')$ is real analytic and $2\pi\mathbb{Z}^3$-periodic.
In 4.2 all derivatives are taken with respect to $p$.
The Leibniz rule, (Ef4), (0.7) and the inequality $\mathrm{xd}(p)^{-1-j}\le7^{5-j}\,\mathrm{xd}(p)^{-6}$ for $0\le j\le4$, which holds because $\mathrm{xd}<7$, give
\begin{equation*}
\begin{split}
\left|\partial_{\mathfrak{l}}^4\operatorname{tr}\left[\mathrm{Rsy}(p)\,\partial_\nu^2D^{(p)}\right]\right|
&\le192\sum_{j=0}^{4}\binom{4}{j}\mathrm{Ksy}\cdot\mathrm{xd}(p)^{-1-j}\cdot6
\le192\cdot6\,\mathrm{Ksy}\cdot\mathrm{xd}(p)^{-6}\sum_{j=0}^{4}\binom{4}{j}7^{5-j}\\
&=192\cdot6\cdot7\cdot8^4\,\mathrm{Ksy}\cdot\mathrm{xd}(p)^{-6}.
\end{split}
\end{equation*}
For four factors,
\[
\partial_{\mathfrak{l}}^4\left(O_1O_2O_3O_4\right)=\sum_{j_1+j_2+j_3+j_4=4}\frac{4!}{j_1!\,j_2!\,j_3!\,j_4!}\left(\partial_{\mathfrak{l}}^{j_1}O_1\right)\left(\partial_{\mathfrak{l}}^{j_2}O_2\right)\left(\partial_{\mathfrak{l}}^{j_3}O_3\right)\left(\partial_{\mathfrak{l}}^{j_4}O_4\right).
\]
Take $O_1=\mathrm{Rsy}(p)$, $O_2=\mathrm{Vx}_\mu(p+p'/2)$, $O_3=\mathrm{Rsy}(p+p')$ and $O_4=\mathrm{Vx}_\nu(p+p'/2)$.
By (0.7) the norm $|\cdot|_2$ of the product $\left(\partial_{\mathfrak{l}}^{j_1}O_1\right)\left(\partial_{\mathfrak{l}}^{j_2}O_2\right)\left(\partial_{\mathfrak{l}}^{j_3}O_3\right)\left(\partial_{\mathfrak{l}}^{j_4}O_4\right)$ is at most $36\,\mathrm{Ksy}^2\cdot\mathrm{xd}(p)^{-1-j_1}\cdot\mathrm{xd}(p+p')^{-1-j_3}$; by (Ef3), $\mathrm{xd}<7$ and $4-j_1-j_3=j_2+j_4$, it is at most $36\,\mathrm{Ksy}^2\,7^{j_2+j_4}\left[\mathrm{xd}(p)^{-6}+\mathrm{xd}(p+p')^{-6}\right]$.
By the multinomial theorem the sum of $\frac{4!}{j_1!\,j_2!\,j_3!\,j_4!}\,7^{j_2+j_4}$ over $j_1+j_2+j_3+j_4=4$ equals $(1+7+1+7)^4=16^4$.
With (Ef4), $\mathrm{Ksy}\ge1$ and (Ef0)\,(x) it follows that
\begin{equation}
\begin{split}
\left|\partial_{\mathfrak{l}}^4\mathrm{lp}_{\mu\nu}(p,p')\right|&\le192\left(6\cdot7\cdot8^4\,\mathrm{Ksy}+36\cdot16^4\,\mathrm{Ksy}^2\right)\left[\mathrm{xd}(p)^{-6}+\mathrm{xd}(p+p')^{-6}\right]\\
&\le486014976\,\mathrm{Ksy}^2\left[\mathrm{xd}(p)^{-6}+\mathrm{xd}(p+p')^{-6}\right].
\end{split}
\tag{4.4}
\end{equation}

\paragraph{4.3. The first bound.}
Apply (4.3) to $\mathrm{fs}=\mathrm{lp}_{\mu\nu}(\,\cdot\,,\eta\,\mathrm{kph})$ with $\mathrm{kph}\in\mathbb{R}^3$; by the definitions in Proposition~36, $\mathrm{Vpl}^{\eta,R}_{\mu\nu}(\mathrm{kph})-\mathrm{Vpl}^{\eta}_{\mu\nu}(\mathrm{kph})$ is $\eta^{-1}$ times the quantity estimated in (4.3).
By (4.4) and by (0.3) with $j=6$,
\[
\max_{\mathfrak{l}}\int_{[-\pi,\pi]^3}\frac{\mathrm{d}^3p}{(2\pi)^3}\left|\partial_{\mathfrak{l}}^4\mathrm{fs}\right|\le486014976\,\mathrm{Ksy}^2\cdot\frac{2\cdot2\pi\,\mathbf{m}^{-3}}{15\,(2\pi)^3}=\frac{486014976\,\mathrm{Ksy}^2}{30\pi^2\,\mathbf{m}^3}.
\]
Multiply by $\eta^{-1}$.
Since $\eta^{-1}n^{-4}\mathbf{m}^{-3}=\eta^{-1}(\eta/R)^4(\bar m\eta)^{-3}=\bar m^{-3}R^{-4}$, and since $42/(30\pi^2)<1/7$ and $486014976/7<7\cdot10^7$ by (Ef0)\,(iii), (x), one obtains
\begin{equation}
\left|\mathrm{Vpl}^{\eta,R}_{\mu\nu}(\mathrm{kph})-\mathrm{Vpl}^{\eta}_{\mu\nu}(\mathrm{kph})\right|\le\frac{42}{30\pi^2}\cdot486014976\,\mathrm{Ksy}^2\,\bar m^{-3}R^{-4}\le7\cdot10^7\,\mathrm{Ksy}^2\,\bar m^{-3}R^{-4}
\tag{4.5}
\end{equation}
for all real $\mathrm{kph}$, all $\eta\in(0,1]$ and all $R$ as in (FF).

\paragraph{4.4. Three identities.}
Put
\[
\mathrm{Fb}_{\mu\nu}(p,p'):=\operatorname{tr}\left[\mathrm{Rsy}(p-p'/2)\cdot\mathrm{Vx}_\mu(p)\cdot\mathrm{Rsy}(p+p'/2)\cdot\mathrm{Vx}_\nu(p)\right]\qquad(p,p'\in\mathbb{R}^3);
\]
this function is real analytic in $(p,p')$ and $2\pi\mathbb{Z}^3$-periodic in $p$.

(i) \emph{Translation.}
The function $p\mapsto\operatorname{tr}\left[\mathrm{Rsy}(p)\cdot\mathrm{Vx}_\mu(p+p'/2)\cdot\mathrm{Rsy}(p+p')\cdot\mathrm{Vx}_\nu(p+p'/2)\right]$ is continuous and $2\pi\mathbb{Z}^3$-periodic.
Its integral over $[-\pi,\pi]^3$ is therefore equal to its integral over the translated cube $[-\pi,\pi]^3-p'/2$, and the substitution $p\mapsto p-p'/2$ turns the latter into $\int_{[-\pi,\pi]^3}\mathrm{d}^3p\,\mathrm{Fb}_{\mu\nu}(p,p')$.

(ii) \emph{A total derivative.}
By (0.6) and the last relation in (0.7), all functions being taken at $p$,
\[
\partial_\mu\operatorname{tr}\left[\mathrm{Rsy}\cdot\mathrm{Vx}_\nu\right]=-\operatorname{tr}\left[\mathrm{Rsy}\cdot\mathrm{Vx}_\mu\cdot\mathrm{Rsy}\cdot\mathrm{Vx}_\nu\right]+\delta_{\mu\nu}\operatorname{tr}\left[\mathrm{Rsy}\,\partial_\nu^2D^{(p)}\right]=\mathrm{lp}_{\mu\nu}(p,0).
\]
The integral over a period cube of a derivative of a periodic $C^1$ function vanishes; hence $\int_{[-\pi,\pi]^3}\mathrm{d}^3p\,\mathrm{lp}_{\mu\nu}(p,0)=0$.
The first term of $\mathrm{lp}_{\mu\nu}(p,p')$ does not depend on $p'$.
Together with (i) this gives
\begin{equation}
\begin{split}
\mathrm{Vpl}^{\eta}_{\mu\nu}(\mathrm{kph})&=\eta^{-1}\int_{[-\pi,\pi]^3}\frac{\mathrm{d}^3p}{(2\pi)^3}\left[\mathrm{lp}_{\mu\nu}(p,\eta\,\mathrm{kph})-\mathrm{lp}_{\mu\nu}(p,0)\right]\\
&=-\eta^{-1}\int_{[-\pi,\pi]^3}\frac{\mathrm{d}^3p}{(2\pi)^3}\left[\mathrm{Fb}_{\mu\nu}(p,\eta\,\mathrm{kph})-\mathrm{Fb}_{\mu\nu}(p,0)\right].
\end{split}
\tag{4.6}
\end{equation}

(iii) \emph{Parity.}
By Proposition~37.1\,(e) one has $\mathrm{Pa}^2=1$, $\mathrm{Pa}\,D^{(p)}\,\mathrm{Pa}=D^{(-p)}$ and $\mathrm{Pa}\cdot\mathrm{Vx}_\nu(p)\cdot\mathrm{Pa}=-\mathrm{Vx}_\nu(-p)$; the second relation gives $\mathrm{Pa}\cdot\mathrm{Rsy}(p)\cdot\mathrm{Pa}=\mathrm{Rsy}(-p)$.
Inserting $\mathrm{Rsy}(-p\pm p'/2)=\mathrm{Pa}\cdot\mathrm{Rsy}(p\mp p'/2)\cdot\mathrm{Pa}$ and, for each direction, $\mathrm{Vx}_\nu(-p)=-\mathrm{Pa}\cdot\mathrm{Vx}_\nu(p)\cdot\mathrm{Pa}$, one finds
\begin{align*}
\mathrm{Fb}_{\mu\nu}(-p,-p')&=\operatorname{tr}\left[\mathrm{Rsy}(-p+p'/2)\cdot\mathrm{Vx}_\mu(-p)\cdot\mathrm{Rsy}(-p-p'/2)\cdot\mathrm{Vx}_\nu(-p)\right]\\
&=\operatorname{tr}\left[\mathrm{Pa}\cdot\mathrm{Rsy}(p-p'/2)\cdot\mathrm{Vx}_\mu(p)\cdot\mathrm{Rsy}(p+p'/2)\cdot\mathrm{Vx}_\nu(p)\cdot\mathrm{Pa}\right]=\mathrm{Fb}_{\mu\nu}(p,p');
\end{align*}
the two signs cancel, and the last equality uses $\mathrm{Pa}^2=1$ and the cyclicity of the trace.
Differentiation in $p'_\alpha$ at $p'=0$ gives $\left(\partial_{p'_\alpha}\mathrm{Fb}_{\mu\nu}\right)(-p,0)=-\left(\partial_{p'_\alpha}\mathrm{Fb}_{\mu\nu}\right)(p,0)$.
As the cube $[-\pi,\pi]^3$ is symmetric,
\begin{equation}
\int_{[-\pi,\pi]^3}\mathrm{d}^3p\,\left(\partial_{p'_\alpha}\mathrm{Fb}_{\mu\nu}\right)(p,0)=0.
\tag{4.7}
\end{equation}

\paragraph{4.5. The second derivatives of $\mathrm{Fb}$ in $p'$.}
Write $p_\pm:=p\pm p'/2$ and $\mathrm{Vx}=\mathrm{Vx}(p)$; the vertices do not depend on $p'$.
By the chain rule,
\begin{equation}
\begin{split}
\partial_{p'_\alpha}\partial_{p'_\beta}\mathrm{Fb}_{\mu\nu}(p,p')
&=\frac14\left\{\operatorname{tr}\left[(\partial_\alpha\partial_\beta\mathrm{Rsy})(p_-)\cdot\mathrm{Vx}_\mu\cdot\mathrm{Rsy}(p_+)\cdot\mathrm{Vx}_\nu\right]\right.\\
&\qquad-\operatorname{tr}\left[(\partial_\alpha\mathrm{Rsy})(p_-)\cdot\mathrm{Vx}_\mu\cdot(\partial_\beta\mathrm{Rsy})(p_+)\cdot\mathrm{Vx}_\nu\right]\\
&\qquad-\operatorname{tr}\left[(\partial_\beta\mathrm{Rsy})(p_-)\cdot\mathrm{Vx}_\mu\cdot(\partial_\alpha\mathrm{Rsy})(p_+)\cdot\mathrm{Vx}_\nu\right]\\
&\qquad\left.{}+\operatorname{tr}\left[\mathrm{Rsy}(p_-)\cdot\mathrm{Vx}_\mu\cdot(\partial_\alpha\partial_\beta\mathrm{Rsy})(p_+)\cdot\mathrm{Vx}_\nu\right]\right\}.
\end{split}
\tag{4.8}
\end{equation}
By (Ef4) and (0.7),
\[
\left|\partial_{p'_\alpha}\partial_{p'_\beta}\mathrm{Fb}_{\mu\nu}(p,p')\right|\le\frac14\cdot192\cdot36\,\mathrm{Ksy}^2\left[\mathrm{xd}(p_-)^{-3}\cdot\mathrm{xd}(p_+)^{-1}+2\,\mathrm{xd}(p_-)^{-2}\cdot\mathrm{xd}(p_+)^{-2}+\mathrm{xd}(p_-)^{-1}\cdot\mathrm{xd}(p_+)^{-3}\right],
\]
and by (Ef3) each of the four products in the bracket is at most $\mathrm{xd}(p_-)^{-4}+\mathrm{xd}(p_+)^{-4}$; that is,
\begin{align}
\left|\partial_{p'_\alpha}\partial_{p'_\beta}\mathrm{Fb}_{\mu\nu}(p,p')\right|&\le6912\,\mathrm{Ksy}^2\left[\mathrm{xd}(p-p'/2)^{-4}+\mathrm{xd}(p+p'/2)^{-4}\right],\tag{4.9}\\
\int_{[-\pi,\pi]^3}\frac{\mathrm{d}^3p}{(2\pi)^3}\left|\partial_{p'_\alpha}\partial_{p'_\beta}\mathrm{Fb}_{\mu\nu}(p,p')\right|&\le6912\,\mathrm{Ksy}^2\cdot\frac{2\cdot4\pi\,\mathbf{m}^{-1}}{3\,(2\pi)^3}=\frac{2304\,\mathrm{Ksy}^2}{\pi^2\,\mathbf{m}}\qquad\text{for all real }p',\tag{4.10}
\end{align}
the integral bound by (0.3) with $j=4$.

\paragraph{4.6. Taylor's formula; the second and the third bound.}
For fixed $p$, Taylor's formula with integral remainder gives
\[
\mathrm{Fb}_{\mu\nu}(p,p')-\mathrm{Fb}_{\mu\nu}(p,0)=\sum_\alpha p'_\alpha\left(\partial_{p'_\alpha}\mathrm{Fb}_{\mu\nu}\right)(p,0)+\int_0^1\mathrm{d}\,\mathrm{sv}\,(1-\mathrm{sv})\sum_{\alpha,\beta}p'_\alpha p'_\beta\left(\partial_{p'_\alpha}\partial_{p'_\beta}\mathrm{Fb}_{\mu\nu}\right)(p,\mathrm{sv}\,p').
\]
Put $p'=\eta\,\mathrm{kph}$, integrate over $p$ (the integrands are continuous on compact sets), and use (4.6) and (4.7):
\begin{equation}
\mathrm{Vpl}^{\eta}_{\mu\nu}(\mathrm{kph})=-\eta\int_{[-\pi,\pi]^3}\frac{\mathrm{d}^3p}{(2\pi)^3}\int_0^1\mathrm{d}\,\mathrm{sv}\,(1-\mathrm{sv})\sum_{\alpha,\beta}\mathrm{kph}_\alpha\cdot\mathrm{kph}_\beta\left(\partial_{p'_\alpha}\partial_{p'_\beta}\mathrm{Fb}_{\mu\nu}\right)(p,\mathrm{sv}\,\eta\,\mathrm{kph}).
\tag{4.11}
\end{equation}
Differentiation of (4.6) under the integral over the compact cube, followed by (4.7) and the identity
\[
\left(\partial_{p'_\alpha}\mathrm{Fb}_{\mu\nu}\right)(p,p')-\left(\partial_{p'_\alpha}\mathrm{Fb}_{\mu\nu}\right)(p,0)=\int_0^1\mathrm{d}\,\mathrm{sv}\sum_\beta p'_\beta\left(\partial_{p'_\alpha}\partial_{p'_\beta}\mathrm{Fb}_{\mu\nu}\right)(p,\mathrm{sv}\,p'),
\]
yields
\begin{equation}
\begin{split}
\partial_{\mathrm{kph}_\alpha}\mathrm{Vpl}^{\eta}_{\mu\nu}(\mathrm{kph})&=-\int_{[-\pi,\pi]^3}\frac{\mathrm{d}^3p}{(2\pi)^3}\left(\partial_{p'_\alpha}\mathrm{Fb}_{\mu\nu}\right)(p,\eta\,\mathrm{kph})\\
&=-\eta\int_{[-\pi,\pi]^3}\frac{\mathrm{d}^3p}{(2\pi)^3}\int_0^1\mathrm{d}\,\mathrm{sv}\sum_\beta\mathrm{kph}_\beta\left(\partial_{p'_\alpha}\partial_{p'_\beta}\mathrm{Fb}_{\mu\nu}\right)(p,\mathrm{sv}\,\eta\,\mathrm{kph}).
\end{split}
\tag{4.12}
\end{equation}
Now use (4.10), $\int_0^1\mathrm{d}\,\mathrm{sv}\,(1-\mathrm{sv})=\frac12$, $\sum_{\alpha,\beta}|\mathrm{kph}_\alpha|\,|\mathrm{kph}_\beta|=|\mathrm{kph}|_1^2\le3\,|\mathrm{kph}|_2^2$, $|\mathrm{kph}|_1\le\sqrt3\,|\mathrm{kph}|_2$, $\eta/\mathbf{m}=1/\bar m$ and (Ef0)\,(iv), (v):
\begin{equation}
\begin{split}
\left|\mathrm{Vpl}^{\eta}_{\mu\nu}(\mathrm{kph})\right|&\le\frac12\cdot3\,|\mathrm{kph}|_2^2\cdot\frac{2304\,\mathrm{Ksy}^2}{\pi^2\bar m}\le\frac{351\,\mathrm{Ksy}^2\,|\mathrm{kph}|_2^2}{\bar m},\\
\left|\partial_{\mathrm{kph}_\alpha}\mathrm{Vpl}^{\eta}_{\mu\nu}(\mathrm{kph})\right|&\le\sqrt3\,|\mathrm{kph}|_2\cdot\frac{2304\,\mathrm{Ksy}^2}{\pi^2\bar m},\qquad
\left|\operatorname{grad}\mathrm{Vpl}^{\eta}_{\mu\nu}(\mathrm{kph})\right|_2\le\frac{3\cdot2304\,\mathrm{Ksy}^2\,|\mathrm{kph}|_2}{\pi^2\bar m}\le\frac{701\,\mathrm{Ksy}^2\,|\mathrm{kph}|_2}{\bar m},
\end{split}
\tag{4.13}
\end{equation}
for all real $\mathrm{kph}$ and all $\eta\in(0,1]$.

\paragraph{4.7. The constant.}
Let
\[
\mathrm{Cpv}_0:=10^8\,\mathrm{Ksy}^2,
\]
with $\mathrm{Ksy}$ of (Ef5).
This number is a function of $\rho$ alone; it contains none of $\bar m$, $\eta$, $R$, $\mathrm{kph}$.
Indeed, by Proposition~37.1\,(g) the bounds (0.7) hold with the single number $\mathrm{Ksy}$ for all $\mathbf{m}\in(0,1]$; $\bar m$ enters 4.1--4.6 only through $\mathbf{m}=\bar m\eta$; and the powers of $\mathbf{m}$ combine with those of $\eta$ and $n$ through the identities $\eta^{-1}n^{-4}\mathbf{m}^{-3}=\bar m^{-3}R^{-4}$ in 4.3 and $\eta/\mathbf{m}=1/\bar m$ in 4.6.
Therefore (4.5) and (4.13) hold with $\mathrm{Cpv}_0$ in the place of their constants, for all $\bar m\in(0,1]$, all real $\mathrm{kph}$, all $\eta\in(0,1]$ and all $R$ as in (FF).
These are the three bounds of clause (c).\qed

\subsubsection*{5. Clause (d)}

\paragraph{5.1. The continuum integrand.}
For $\mathrm{lv}\in\mathbb{R}^3$ put
\[
\mathrm{den}(\mathrm{lv}):=\bar m^2+|\mathrm{lv}|_2^2,\qquad\mathrm{Sct}(\mathrm{lv}):=\left(\bar m+i\gamma\cdot\mathrm{lv}\right)^{-1}.
\]
In the notation of Chapter~37 one has $\mathrm{Sct}(\mathrm{lv})=\mathrm{Sdr}(\mathrm{lv},\bar m)$; in this chapter the letter $\mathrm{Sct}$ always denotes this function, which carries no cut-off.
By (I.c1), $\left(\bar m+i\gamma\cdot\mathrm{lv}\right)\left(\bar m-i\gamma\cdot\mathrm{lv}\right)=\bar m^2+(\gamma\cdot\mathrm{lv})^2=\mathrm{den}(\mathrm{lv})$.
Hence
\begin{equation}
\mathrm{Sct}(\mathrm{lv})=\frac{\bar m-i\gamma\cdot\mathrm{lv}}{\mathrm{den}(\mathrm{lv})},\qquad\partial_\alpha\mathrm{Sct}=-\mathrm{Sct}\,i\gamma_\alpha\,\mathrm{Sct}.
\tag{5.1}
\end{equation}
For $\mathrm{lph},\mathrm{kph}\in\mathbb{R}^3$ let
\[
\mathrm{Fq}_{\mu\nu}(\mathrm{lph},\mathrm{kph}):=\operatorname{tr}_4\left[\mathrm{Sct}(\mathrm{lph}-\mathrm{kph}/2)\,i\gamma_\mu\,\mathrm{Sct}(\mathrm{lph}+\mathrm{kph}/2)\,i\gamma_\nu\right].
\]
This is a smooth function on $\mathbb{R}^6$, and the chain rule gives, with $\mathrm{lph}_\pm:=\mathrm{lph}\pm\mathrm{kph}/2$,
\begin{equation}
\begin{split}
\partial_{\mathrm{kph}_\alpha}\partial_{\mathrm{kph}_\beta}\mathrm{Fq}_{\mu\nu}(\mathrm{lph},\mathrm{kph})
&=\frac14\left\{\operatorname{tr}_4\left[(\partial_\alpha\partial_\beta\mathrm{Sct})(\mathrm{lph}_-)\,i\gamma_\mu\,\mathrm{Sct}(\mathrm{lph}_+)\,i\gamma_\nu\right]\right.\\
&\qquad-\operatorname{tr}_4\left[(\partial_\alpha\mathrm{Sct})(\mathrm{lph}_-)\,i\gamma_\mu\,(\partial_\beta\mathrm{Sct})(\mathrm{lph}_+)\,i\gamma_\nu\right]\\
&\qquad-\operatorname{tr}_4\left[(\partial_\beta\mathrm{Sct})(\mathrm{lph}_-)\,i\gamma_\mu\,(\partial_\alpha\mathrm{Sct})(\mathrm{lph}_+)\,i\gamma_\nu\right]\\
&\qquad\left.{}+\operatorname{tr}_4\left[\mathrm{Sct}(\mathrm{lph}_-)\,i\gamma_\mu\,(\partial_\alpha\partial_\beta\mathrm{Sct})(\mathrm{lph}_+)\,i\gamma_\nu\right]\right\}.
\end{split}
\tag{5.2}
\end{equation}
Define
\begin{equation}
\mathrm{Vpc}_{\mu\nu}(\mathrm{kph}):=-\int_{\mathbb{R}^3}\frac{\mathrm{d}^3\mathrm{lph}}{(2\pi)^3}\int_0^1\mathrm{d}\,\mathrm{sv}\,(1-\mathrm{sv})\sum_{\alpha,\beta}\mathrm{kph}_\alpha\cdot\mathrm{kph}_\beta\left(\partial_{\mathrm{kph}_\alpha}\partial_{\mathrm{kph}_\beta}\mathrm{Fq}_{\mu\nu}\right)(\mathrm{lph},\mathrm{sv}\cdot\mathrm{kph}).
\tag{5.3}
\end{equation}

\paragraph{5.2. Pointwise convergence.}
\begin{named}{Claim}
For every $\mathrm{kph}\in\mathbb{R}^3$ the integral (5.3) converges absolutely, and $\mathrm{Vpl}^{\eta}_{\mu\nu}(\mathrm{kph})\to\mathrm{Vpc}_{\mu\nu}(\mathrm{kph})$ as $\eta\to0$.
\end{named}
\begin{proof}
Let $\eta_j\in(0,1]$ with $\eta_j\to0$.
From some index on, $\eta_j|\mathrm{kph}|_\infty\le\pi$; the finitely many earlier terms are discarded.
In (4.11) at $\eta=\eta_j$, where $\mathrm{Rsy}$ is formed with $\mathbf{m}=\bar m\eta_j$, substitute $p=\eta_j\,\mathrm{lph}$, $\mathrm{d}^3p=\eta_j^3\,\mathrm{d}^3\mathrm{lph}$:
\begin{equation}
\begin{split}
\mathrm{Vpl}^{\eta_j}_{\mu\nu}(\mathrm{kph})=-\int_{\mathbb{R}^3}\frac{\mathrm{d}^3\mathrm{lph}}{(2\pi)^3}\int_0^1\mathrm{d}\,\mathrm{sv}\,(1-\mathrm{sv})\sum_{\alpha,\beta}&\mathrm{kph}_\alpha\cdot\mathrm{kph}_\beta\,\mathbf{1}\left[|\mathrm{lph}|_\infty\le\frac{\pi}{\eta_j}\right]\\
&\times\eta_j^4\left(\partial_{p'_\alpha}\partial_{p'_\beta}\mathrm{Fb}_{\mu\nu}\right)(\eta_j\,\mathrm{lph},\mathrm{sv}\,\eta_j\,\mathrm{kph}).
\end{split}
\tag{5.4}
\end{equation}

(1) \emph{Limit of the integrand.}
Fix $\mathrm{lph}$ and $\mathrm{sv}$, and put $\mathrm{lph}_\pm:=\mathrm{lph}\pm\mathrm{sv}\cdot\mathrm{kph}/2$.
In (4.8) at $p=\eta_j\,\mathrm{lph}$ and $p'=\mathrm{sv}\,\eta_j\,\mathrm{kph}$ one has $p_\pm=\eta_j\,\mathrm{lph}_\pm$.
Each of the four traces in (4.8) has the form $\operatorname{tr}\left[(\partial^{\mathrm{ak}}\mathrm{Rsy})(p_-)\cdot\mathrm{Vx}_\mu(p)\cdot(\partial^{\mathrm{ak}'}\mathrm{Rsy})(p_+)\cdot\mathrm{Vx}_\nu(p)\right]$ with multi-indices satisfying $|\mathrm{ak}|+|\mathrm{ak}'|=2$, namely
\[
(\mathrm{ak},\mathrm{ak}')=(e_\alpha+e_\beta,0),\quad(e_\alpha,e_\beta),\quad(e_\beta,e_\alpha),\quad(0,e_\alpha+e_\beta).
\]
Write $\eta_j^4=\eta_j^{1+|\mathrm{ak}|}\,\eta_j^{1+|\mathrm{ak}'|}$.
Proposition~37.1\,(h), applied to the constant sequences $\mathrm{lph}_-$ and $\mathrm{lph}_+$, gives
\[
\eta_j^{1+|\mathrm{ak}|}\left(\partial^{\mathrm{ak}}\mathrm{Rsy}\right)(\eta_j\,\mathrm{lph}_\pm)\to J\left(\partial^{\mathrm{ak}}\mathrm{Sct}\right)(\mathrm{lph}_\pm)J^{*}\qquad(|\mathrm{ak}|\le2).
\]
By the last relation in (0.7), $\mathrm{Vx}_\mu(p)-\mathrm{Vx}_\mu(0)=\int_0^1\mathrm{d}\,\mathrm{sv}'\,p_\mu\left(\partial_\mu\mathrm{Vx}_\mu\right)(\mathrm{sv}'p)$, so that $\left|\mathrm{Vx}_\mu(\eta_j\,\mathrm{lph})-\mathrm{Vx}_\mu(0)\right|_2\le6\,\eta_j\,|\mathrm{lph}_\mu|\to0$.
Products and traces are continuous.
Hence, for each of the four pairs $(\mathrm{ak},\mathrm{ak}')$,
\begin{align*}
&\eta_j^4\operatorname{tr}\left[(\partial^{\mathrm{ak}}\mathrm{Rsy})(\eta_j\,\mathrm{lph}_-)\cdot\mathrm{Vx}_\mu(\eta_j\,\mathrm{lph})\cdot(\partial^{\mathrm{ak}'}\mathrm{Rsy})(\eta_j\,\mathrm{lph}_+)\cdot\mathrm{Vx}_\nu(\eta_j\,\mathrm{lph})\right]\\
&\qquad\to\operatorname{tr}\left[J(\partial^{\mathrm{ak}}\mathrm{Sct})(\mathrm{lph}_-)J^{*}\,\mathrm{Vx}_\mu(0)\,J(\partial^{\mathrm{ak}'}\mathrm{Sct})(\mathrm{lph}_+)J^{*}\,\mathrm{Vx}_\nu(0)\right]\\
&\qquad=\operatorname{tr}_4\left[(\partial^{\mathrm{ak}}\mathrm{Sct})(\mathrm{lph}_-)\left(J^{*}\mathrm{Vx}_\mu(0)J\right)(\partial^{\mathrm{ak}'}\mathrm{Sct})(\mathrm{lph}_+)\left(J^{*}\mathrm{Vx}_\nu(0)J\right)\right]\\
&\qquad=\operatorname{tr}_4\left[(\partial^{\mathrm{ak}}\mathrm{Sct})(\mathrm{lph}_-)\,i\gamma_\mu\,(\partial^{\mathrm{ak}'}\mathrm{Sct})(\mathrm{lph}_+)\,i\gamma_\nu\right];
\end{align*}
the first equality is (0.5) with $\mathrm{Mx}=(\partial^{\mathrm{ak}}\mathrm{Sct})(\mathrm{lph}_-)$ and $O=\mathrm{Vx}_\mu(0)\,J(\partial^{\mathrm{ak}'}\mathrm{Sct})(\mathrm{lph}_+)J^{*}\,\mathrm{Vx}_\nu(0)$, where $J^{*}u=\frac{1}{48}\sum_{g\in O_{\mathrm{cub}}}u(g)$ is the sum over the $48$ labels with the weight $\frac{1}{48}$, as in (Ef4);
the second one is the relation $J^{*}\mathrm{Vx}_\nu(0)J=i\gamma_\nu$ of Proposition~37.1\,(d), which rests on the fact that, for each sign, exactly $8$ of the $48$ labels $g$ satisfy $ge_1=\pm e_\nu$.
Since (4.8) and (5.2) combine the four traces with the same coefficients, namely $\frac14$, $-\frac14$, $-\frac14$, $\frac14$ in the order of the pairs $(\mathrm{ak},\mathrm{ak}')$ listed above, it follows that
\[
\eta_j^4\left(\partial_{p'_\alpha}\partial_{p'_\beta}\mathrm{Fb}_{\mu\nu}\right)(\eta_j\,\mathrm{lph},\mathrm{sv}\,\eta_j\,\mathrm{kph})\to\left(\partial_{\mathrm{kph}_\alpha}\partial_{\mathrm{kph}_\beta}\mathrm{Fq}_{\mu\nu}\right)(\mathrm{lph},\mathrm{sv}\cdot\mathrm{kph}),
\]
and the indicator in (5.4) tends to $1$.

(2) \emph{Majorant.}
Let $|\mathrm{lph}|_\infty\le\pi/\eta_j$ and $0\le\mathrm{sv}\le1$.
Every coordinate of $\eta_j\,\mathrm{lph}\pm\mathrm{sv}\,\eta_j\,\mathrm{kph}/2$ has modulus at most $\pi+\pi/2$.
By (0.2),
\[
\mathrm{xd}\left(\eta_j\,\mathrm{lph}\pm\mathrm{sv}\,\eta_j\,\mathrm{kph}/2\right)\ge\frac{\eta_j}{3}\left(|\mathrm{lph}\pm\mathrm{sv}\cdot\mathrm{kph}/2|_2+\bar m\right).
\]
With (4.9) this gives
\begin{equation}
\begin{split}
&\mathbf{1}\left[|\mathrm{lph}|_\infty\le\frac{\pi}{\eta_j}\right]\eta_j^4\left|\left(\partial_{p'_\alpha}\partial_{p'_\beta}\mathrm{Fb}_{\mu\nu}\right)(\eta_j\,\mathrm{lph},\mathrm{sv}\,\eta_j\,\mathrm{kph})\right|\\
&\qquad\le81\cdot6912\,\mathrm{Ksy}^2\left[\left(|\mathrm{lph}-\mathrm{sv}\cdot\mathrm{kph}/2|_2+\bar m\right)^{-4}+\left(|\mathrm{lph}+\mathrm{sv}\cdot\mathrm{kph}/2|_2+\bar m\right)^{-4}\right].
\end{split}
\tag{5.5}
\end{equation}
The right side does not depend on $j$, and by (0.4) its integral over $\mathbb{R}^3\times[0,1]$ equals $81\cdot6912\,\mathrm{Ksy}^2\cdot\frac{8\pi}{3\bar m}$.

By (1), (2) and the dominated convergence theorem the right side of (5.4) tends to the right side of (5.3); by (1) the integrand of (5.3) obeys the bound (5.5), so that the integral (5.3) converges absolutely.
The sequence $(\eta_j)$ was arbitrary.
\end{proof}

\paragraph{5.3. Local uniformity.}
Let $\mathrm{rad}>0$; the ball $\left\{|\mathrm{kph}|_2\le\mathrm{rad}\right\}$ is convex.
By (4.13) and 4.7,
\[
\left|\mathrm{Vpl}^{\eta}_{\mu\nu}(\mathrm{kph})-\mathrm{Vpl}^{\eta}_{\mu\nu}(\mathrm{kph}')\right|\le\frac{\mathrm{Cpv}_0\cdot\mathrm{rad}}{\bar m}\,|\mathrm{kph}-\mathrm{kph}'|_2
\]
for $\mathrm{kph},\mathrm{kph}'$ in the ball and all $\eta\in(0,1]$; by 5.2 the same bound holds for $\mathrm{Vpc}_{\mu\nu}$.
Let $\delta>0$.
Choose finitely many points of the ball such that every point of the ball has Euclidean distance at most $\frac{\delta\bar m}{3\,\mathrm{Cpv}_0\cdot\mathrm{rad}}$ from one of them, and choose, by 5.2, a number $\eta'>0$ such that $\left|\mathrm{Vpl}^{\eta}_{\mu\nu}-\mathrm{Vpc}_{\mu\nu}\right|\le\delta/3$ at these points for all $\eta\le\eta'$.
Then $\left|\mathrm{Vpl}^{\eta}_{\mu\nu}(\mathrm{kph})-\mathrm{Vpc}_{\mu\nu}(\mathrm{kph})\right|\le\delta$ for all $\mathrm{kph}$ in the ball and all $\eta\le\eta'$.
Thus $\mathrm{Vpl}^{\eta}_{\mu\nu}\to\mathrm{Vpc}_{\mu\nu}$ uniformly on every ball.

\paragraph{5.4. Evaluation of $\mathrm{Vpc}$.}
\begin{named}{Claim}
With $\mathrm{vpo}_{\bar m}$ as in Section~G5.6.0,
\begin{equation}
\mathrm{Vpc}_{\mu\nu}(\mathrm{kph})=-\left(|\mathrm{kph}|_2^2\,\delta_{\mu\nu}-\mathrm{kph}_\mu\cdot\mathrm{kph}_\nu\right)\mathrm{vpo}_{\bar m}\left(|\mathrm{kph}|_2^2\right)\qquad\text{for all }\mathrm{kph}\in\mathbb{R}^3.
\tag{5.6}
\end{equation}
\end{named}
For $\mathrm{kph}=0$ both sides vanish.
Let $\mathrm{kph}\neq0$; in the rest of 5.4, $\mathrm{lph}_\pm:=\mathrm{lph}\pm\mathrm{kph}/2$.

(i) \emph{The integrand.}
Traces of products of the matrices $\gamma_\mu$ are obtained from (I.c1) and (I.c2).
First, $\operatorname{tr}_4\gamma_\mu=0$, because $\gamma_\mu=\mathrm{diag}(\sigma_\mu,-\sigma_\mu)$.
Next, $\operatorname{tr}_4[\gamma_\alpha\gamma_\beta]=\frac12\operatorname{tr}_4[\gamma_\alpha\gamma_\beta+\gamma_\beta\gamma_\alpha]=4\,\delta_{\alpha\beta}$.
Further, $\operatorname{tr}_4[\gamma_\alpha\gamma_\beta\gamma_{\mathfrak{l}}]=0$: if two of the indices coincide, the product is $\pm$ one matrix $\gamma$ by (I.c1); if they are distinct, it equals $\pm\gamma_1\gamma_2\gamma_3=\pm i\Sigma_3$ by (I.c1), (I.c2), and $\operatorname{tr}_4\Sigma_3=0$.
Moving $\gamma_\alpha$ to the right by means of (I.c1) gives
\[
\gamma_\alpha\gamma_\mu\gamma_\beta\gamma_\nu=2\delta_{\alpha\mu}\gamma_\beta\gamma_\nu-2\delta_{\alpha\beta}\gamma_\mu\gamma_\nu+2\delta_{\alpha\nu}\gamma_\mu\gamma_\beta-\gamma_\mu\gamma_\beta\gamma_\nu\gamma_\alpha,
\]
and by cyclicity $\operatorname{tr}_4[\gamma_\alpha\gamma_\mu\gamma_\beta\gamma_\nu]=4\left(\delta_{\alpha\mu}\delta_{\beta\nu}-\delta_{\alpha\beta}\delta_{\mu\nu}+\delta_{\alpha\nu}\delta_{\mu\beta}\right)$.
Consequently, for $\mathrm{lv},\mathrm{lv}'\in\mathbb{R}^3$,
\begin{equation}
\begin{split}
&\operatorname{tr}_4\left[\left(\bar m-i\gamma\cdot\mathrm{lv}\right)\gamma_\mu\left(\bar m-i\gamma\cdot\mathrm{lv}'\right)\gamma_\nu\right]=4\bar m^2\delta_{\mu\nu}-4\left(\mathrm{lv}_\mu\cdot\mathrm{lv}'_\nu-\mathrm{lv}\cdot\mathrm{lv}'\,\delta_{\mu\nu}+\mathrm{lv}_\nu\cdot\mathrm{lv}'_\mu\right)=4\,\mathrm{num}_{\mu\nu}(\mathrm{lv},\mathrm{lv}'),\\
&\mathrm{num}_{\mu\nu}(\mathrm{lv},\mathrm{lv}'):=\left(\bar m^2+\mathrm{lv}\cdot\mathrm{lv}'\right)\delta_{\mu\nu}-\mathrm{lv}_\mu\cdot\mathrm{lv}'_\nu-\mathrm{lv}_\nu\cdot\mathrm{lv}'_\mu,\qquad
\mathrm{Fq}_{\mu\nu}(\mathrm{lph},\mathrm{kph})=-\frac{4\,\mathrm{num}_{\mu\nu}(\mathrm{lph}_-,\mathrm{lph}_+)}{\mathrm{den}(\mathrm{lph}_-)\cdot\mathrm{den}(\mathrm{lph}_+)};
\end{split}
\tag{5.7}
\end{equation}
the last relation follows from (5.1) and $i\cdot i=-1$.
The number $\mathrm{num}_{\mu\nu}(\mathrm{lv},\mathrm{lv}')$ is real, symmetric in $(\mu,\nu)$, and symmetric under $\mathrm{lv}\leftrightarrow\mathrm{lv}'$.

(ii) \emph{Removal of the Taylor form.}
The substitution $\mathrm{kph}\mapsto-\mathrm{kph}$ exchanges $\mathrm{lph}_-$ and $\mathrm{lph}_+$.
By (5.7), $\mathrm{Fq}_{\mu\nu}(\mathrm{lph},-\mathrm{kph})=\mathrm{Fq}_{\mu\nu}(\mathrm{lph},\mathrm{kph})$; so the first derivatives of $\mathrm{Fq}_{\mu\nu}(\mathrm{lph},\,\cdot\,)$ vanish at $0$, and Taylor's formula gives
\[
\mathrm{Fq}_{\mu\nu}(\mathrm{lph},\mathrm{kph})-\mathrm{Fq}_{\mu\nu}(\mathrm{lph},0)=\int_0^1\mathrm{d}\,\mathrm{sv}\,(1-\mathrm{sv})\sum_{\alpha,\beta}\mathrm{kph}_\alpha\cdot\mathrm{kph}_\beta\left(\partial_{\mathrm{kph}_\alpha}\partial_{\mathrm{kph}_\beta}\mathrm{Fq}_{\mu\nu}\right)(\mathrm{lph},\mathrm{sv}\cdot\mathrm{kph}).
\]
The absolute convergence of (5.3) and Fubini's theorem give
\begin{equation}
\mathrm{Vpc}_{\mu\nu}(\mathrm{kph})=-\int_{\mathbb{R}^3}\frac{\mathrm{d}^3\mathrm{lph}}{(2\pi)^3}\left[\mathrm{Fq}_{\mu\nu}(\mathrm{lph},\mathrm{kph})-\mathrm{Fq}_{\mu\nu}(\mathrm{lph},0)\right],
\tag{5.8}
\end{equation}
where the integral converges absolutely and is therefore the limit, as $\mathrm{rad}\to\infty$, of the integrals over $|\mathrm{lph}|_2<\mathrm{rad}$.
By (5.7) the $3\times3$ matrix $\mathrm{Vpc}(\mathrm{kph})=\left(\mathrm{Vpc}_{\mu\nu}(\mathrm{kph})\right)$ is real and symmetric.

(iii) \emph{Covariance.}
Let $\mathrm{Ro}\in\mathrm{SO}(3)$.
As $3\times3$ matrices, $\mathrm{num}(\mathrm{Ro}\cdot\mathrm{lv},\mathrm{Ro}\cdot\mathrm{lv}')=\mathrm{Ro}\cdot\mathrm{num}(\mathrm{lv},\mathrm{lv}')\cdot\mathrm{Ro}^{T}$; moreover $\mathrm{den}(\mathrm{Ro}\cdot\mathrm{lv})=\mathrm{den}(\mathrm{lv})$ and $\mathrm{Ro}\cdot\mathrm{lph}\pm\mathrm{Ro}\cdot\mathrm{kph}/2=\mathrm{Ro}\cdot\mathrm{lph}_\pm$.
Hence $\mathrm{Fq}(\mathrm{Ro}\cdot\mathrm{lph},\mathrm{Ro}\cdot\mathrm{kph})=\mathrm{Ro}\cdot\mathrm{Fq}(\mathrm{lph},\mathrm{kph})\cdot\mathrm{Ro}^{T}$.
The substitution $\mathrm{lph}\mapsto\mathrm{Ro}\cdot\mathrm{lph}$ in (5.8) preserves the Lebesgue measure, and so $\mathrm{Vpc}(\mathrm{Ro}\cdot\mathrm{kph})=\mathrm{Ro}\cdot\mathrm{Vpc}(\mathrm{kph})\cdot\mathrm{Ro}^{T}$.
Let $\mathrm{uv}:=\mathrm{kph}/|\mathrm{kph}|_2$.
Every rotation $\mathrm{Ro}$ about the axis $\mathrm{uv}$ satisfies $\mathrm{Ro}\cdot\mathrm{kph}=\mathrm{kph}$, so that $\mathrm{Vpc}(\mathrm{kph})$ commutes with $\mathrm{Ro}$.
The vector $\mathrm{Vpc}(\mathrm{kph})\cdot\mathrm{uv}$ is then fixed by all rotations about $\mathrm{uv}$, whence $\mathrm{Vpc}(\mathrm{kph})\cdot\mathrm{uv}=\mathrm{zl}\cdot\mathrm{uv}$ with a real number $\mathrm{zl}$.
Since $\mathrm{Vpc}(\mathrm{kph})$ is symmetric, it maps the plane orthogonal to $\mathrm{uv}$ into itself, and its restriction to this plane is a symmetric $2\times2$ matrix $\left(\begin{array}{cc}\mathrm{zq}_1&\mathrm{zq}_2\\ \mathrm{zq}_2&\mathrm{zq}_3\end{array}\right)$ that commutes with the rotation by a right angle $\left(\begin{array}{cc}0&-1\\ 1&0\end{array}\right)$.
The two products are
\[
\left(\begin{array}{cc}\mathrm{zq}_1&\mathrm{zq}_2\\ \mathrm{zq}_2&\mathrm{zq}_3\end{array}\right)\left(\begin{array}{cc}0&-1\\ 1&0\end{array}\right)=\left(\begin{array}{cc}\mathrm{zq}_2&-\mathrm{zq}_1\\ \mathrm{zq}_3&-\mathrm{zq}_2\end{array}\right),\qquad
\left(\begin{array}{cc}0&-1\\ 1&0\end{array}\right)\left(\begin{array}{cc}\mathrm{zq}_1&\mathrm{zq}_2\\ \mathrm{zq}_2&\mathrm{zq}_3\end{array}\right)=\left(\begin{array}{cc}-\mathrm{zq}_2&-\mathrm{zq}_3\\ \mathrm{zq}_1&\mathrm{zq}_2\end{array}\right),
\]
and their equality forces $\mathrm{zq}_2=0$ and $\mathrm{zq}_1=\mathrm{zq}_3$.
It follows that

\begin{equation}
\begin{aligned}
&\mathrm{Vpc}(\mathrm{kph})=\mathrm{zt}\cdot\left(1-\mathrm{uv}\,\mathrm{uv}^{T}\right)+\mathrm{zl}\cdot\mathrm{uv}\,\mathrm{uv}^{T}\quad\text{with real numbers }\mathrm{zt},\ \mathrm{zl};\\
&\mathrm{tr}_3\,\mathrm{Vpc}(\mathrm{kph}):=\sum_{\mu}\mathrm{Vpc}_{\mu\mu}(\mathrm{kph})=2\,\mathrm{zt}+\mathrm{zl}.
\end{aligned}
\tag{5.9}
\end{equation}

(iv) \emph{The number $\mathrm{zl}$ vanishes.}
For $\mathrm{lv}\in\mathbb{R}^3$ let $\mathrm{jv}_\nu(\mathrm{lv}):=\mathrm{tr}_4\left[\mathrm{Sct}(\mathrm{lv})\,\mathrm{i}\gamma_\nu\right]$.
By (5.1) and (i), $\mathrm{jv}_\nu(\mathrm{lv})=\mathrm{tr}_4\left[(\gamma\cdot\mathrm{lv})\,\gamma_\nu\right]/\mathrm{den}(\mathrm{lv})=4\,\mathrm{lv}_\nu/\mathrm{den}(\mathrm{lv})$.
Since $\mathrm{lph}_+-\mathrm{lph}_-=\mathrm{kph}$, one has $\mathrm{i}\gamma\cdot\mathrm{kph}=\mathrm{Sct}(\mathrm{lph}_+)^{-1}-\mathrm{Sct}(\mathrm{lph}_-)^{-1}$.
It follows that
\[
\sum_\mu\mathrm{kph}_\mu\cdot\mathrm{Fq}_{\mu\nu}(\mathrm{lph},\mathrm{kph})
=\mathrm{tr}_4\left[\mathrm{Sct}(\mathrm{lph}_-)\left(\mathrm{Sct}(\mathrm{lph}_+)^{-1}-\mathrm{Sct}(\mathrm{lph}_-)^{-1}\right)\mathrm{Sct}(\mathrm{lph}_+)\,\mathrm{i}\gamma_\nu\right]
=\mathrm{jv}_\nu(\mathrm{lph}_-)-\mathrm{jv}_\nu(\mathrm{lph}_+),
\]
while the second relation of (5.1) gives
\[
\sum_\mu\mathrm{kph}_\mu\cdot\mathrm{Fq}_{\mu\nu}(\mathrm{lph},0)
=\mathrm{tr}_4\left[\mathrm{Sct}(\mathrm{lph})\,(\mathrm{i}\gamma\cdot\mathrm{kph})\,\mathrm{Sct}(\mathrm{lph})\,\mathrm{i}\gamma_\nu\right]
=-\sum_\alpha\mathrm{kph}_\alpha\cdot(\partial_\alpha\mathrm{jv}_\nu)(\mathrm{lph}).
\]
Consequently
\begin{align*}
\sum_\mu\mathrm{kph}_\mu\cdot\left[\mathrm{Fq}_{\mu\nu}(\mathrm{lph},\mathrm{kph})-\mathrm{Fq}_{\mu\nu}(\mathrm{lph},0)\right]
&=\mathrm{jv}_\nu(\mathrm{lph}-\mathrm{kph}/2)-\mathrm{jv}_\nu(\mathrm{lph}+\mathrm{kph}/2)+\sum_\alpha\mathrm{kph}_\alpha\cdot(\partial_\alpha\mathrm{jv}_\nu)(\mathrm{lph})\\
&=-\sum_\alpha\mathrm{kph}_\alpha\cdot(\partial_\alpha\mathrm{Wj})(\mathrm{lph}),
\end{align*}
where, for fixed $\nu$ and $\mathrm{kph}$,
\[
\mathrm{Wj}(\mathrm{lph}):=\int_0^{1/2}\mathrm{d}\,\mathrm{sv}\,\left[\mathrm{jv}_\nu(\mathrm{lph}+\mathrm{sv}\cdot\mathrm{kph})+\mathrm{jv}_\nu(\mathrm{lph}-\mathrm{sv}\cdot\mathrm{kph})-2\,\mathrm{jv}_\nu(\mathrm{lph})\right].
\]
Indeed, for either choice of the sign,
\[
\sum_\alpha\mathrm{kph}_\alpha\,\partial_\alpha\int_0^{1/2}\mathrm{d}\,\mathrm{sv}\,\mathrm{jv}_\nu(\mathrm{lph}\pm\mathrm{sv}\cdot\mathrm{kph})
=\pm\int_0^{1/2}\mathrm{d}\,\mathrm{sv}\,\frac{\mathrm{d}}{\mathrm{d}\,\mathrm{sv}}\,\mathrm{jv}_\nu(\mathrm{lph}\pm\mathrm{sv}\cdot\mathrm{kph})
=\pm\left(\mathrm{jv}_\nu(\mathrm{lph}\pm\mathrm{kph}/2)-\mathrm{jv}_\nu(\mathrm{lph})\right).
\]
The divergence theorem on the ball of radius $\mathrm{rad}$ gives
\[
\int_{|\mathrm{lph}|_2<\mathrm{rad}}\mathrm{d}^3\mathrm{lph}\,\sum_\alpha\mathrm{kph}_\alpha\cdot(\partial_\alpha\mathrm{Wj})(\mathrm{lph})
=\int_{|\mathrm{lph}|_2=\mathrm{rad}}\mathrm{d}S\,\frac{\mathrm{kph}\cdot\mathrm{lph}}{\mathrm{rad}}\,\mathrm{Wj}(\mathrm{lph}).
\]
To bound $\mathrm{Wj}$, note that
\[
\partial_\beta\partial_\alpha\mathrm{jv}_\nu(\mathrm{lv})
=-\frac{8\left(\delta_{\alpha\nu}\mathrm{lv}_\beta+\delta_{\beta\nu}\mathrm{lv}_\alpha+\delta_{\alpha\beta}\mathrm{lv}_\nu\right)}{\mathrm{den}(\mathrm{lv})^2}
+\frac{32\,\mathrm{lv}_\nu\cdot\mathrm{lv}_\alpha\cdot\mathrm{lv}_\beta}{\mathrm{den}(\mathrm{lv})^3};
\]
hence, for every $\mathrm{uv}'\in\mathbb{R}^3$ with $|\mathrm{uv}'|_2=1$,
\[
\left|\sum_{\alpha,\beta}\mathrm{uv}'_\alpha\cdot\mathrm{uv}'_\beta\cdot\partial_\alpha\partial_\beta\mathrm{jv}_\nu(\mathrm{lv})\right|
\le\frac{24\,|\mathrm{lv}|_2}{\mathrm{den}(\mathrm{lv})^2}+\frac{32\,|\mathrm{lv}|_2^3}{\mathrm{den}(\mathrm{lv})^3}
\le56\,|\mathrm{lv}|_2^{-3}.
\]
For a function $\mathrm{phs}$ of class $C^2$ of one real variable and for $\mathrm{sv}\ge0$,
\[
\mathrm{phs}(\mathrm{sv})+\mathrm{phs}(-\mathrm{sv})-2\,\mathrm{phs}(0)=\int_{-\mathrm{sv}}^{\mathrm{sv}}\mathrm{d}\,\mathrm{sv}'\,\left(\mathrm{sv}-|\mathrm{sv}'|\right)\mathrm{phs}''(\mathrm{sv}').
\]
Let $|\mathrm{lph}|_2=\mathrm{rad}>|\mathrm{kph}|_2$ and $0\le\mathrm{sv}\le1/2$, and take $\mathrm{phs}(\mathrm{sv}'):=\mathrm{jv}_\nu(\mathrm{lph}+\mathrm{sv}'\cdot\mathrm{kph})$.
For $|\mathrm{sv}'|\le1/2$ one has $|\mathrm{lph}+\mathrm{sv}'\cdot\mathrm{kph}|_2\ge\mathrm{rad}-|\mathrm{kph}|_2/2>0$, and the bound above, applied with $\mathrm{uv}'=\mathrm{kph}/|\mathrm{kph}|_2$, yields
\[
\left|\mathrm{jv}_\nu(\mathrm{lph}+\mathrm{sv}\cdot\mathrm{kph})+\mathrm{jv}_\nu(\mathrm{lph}-\mathrm{sv}\cdot\mathrm{kph})-2\,\mathrm{jv}_\nu(\mathrm{lph})\right|
\le56\,\mathrm{sv}^2\cdot|\mathrm{kph}|_2^2\cdot(\mathrm{rad}-|\mathrm{kph}|_2/2)^{-3}.
\]
As $\int_0^{1/2}\mathrm{d}\,\mathrm{sv}\,\mathrm{sv}^2=1/24$, it follows that $|\mathrm{Wj}(\mathrm{lph})|\le\frac{56}{24}\,|\mathrm{kph}|_2^2\,(\mathrm{rad}-|\mathrm{kph}|_2/2)^{-3}$ on the sphere $|\mathrm{lph}|_2=\mathrm{rad}$.
So the surface integral above has modulus at most
\[
|\mathrm{kph}|_2\cdot4\pi\,\mathrm{rad}^2\cdot\frac{56}{24}\,|\mathrm{kph}|_2^2\,(\mathrm{rad}-|\mathrm{kph}|_2/2)^{-3},
\]
which tends to $0$ as $\mathrm{rad}\to\infty$.
By (5.8), $\sum_\mu\mathrm{kph}_\mu\cdot\mathrm{Vpc}_{\mu\nu}(\mathrm{kph})$ is the limit, as $\mathrm{rad}\to\infty$, of $(2\pi)^{-3}$ times this surface integral.
Therefore $\sum_\mu\mathrm{kph}_\mu\cdot\mathrm{Vpc}_{\mu\nu}(\mathrm{kph})=0$ for $\nu=1,2,3$.
Since the matrix $\mathrm{Vpc}(\mathrm{kph})$ is symmetric, (5.9) gives $\sum_\mu\mathrm{kph}_\mu\cdot\mathrm{Vpc}_{\mu\nu}(\mathrm{kph})=\mathrm{zl}\cdot\mathrm{kph}_\nu$.
Hence $\mathrm{zl}=0$, and
\begin{equation}
\mathrm{Vpc}_{\mu\nu}(\mathrm{kph})=\frac12\left(\mathrm{tr}_3\,\mathrm{Vpc}(\mathrm{kph})\right)\left(\delta_{\mu\nu}-\frac{\mathrm{kph}_\mu\cdot\mathrm{kph}_\nu}{|\mathrm{kph}|_2^2}\right).
\tag{5.10}
\end{equation}

(v) \emph{The trace.}
The definition of $\mathrm{num}_{\mu\nu}$ in (5.7) gives $\sum_\mu\mathrm{num}_{\mu\mu}(\mathrm{lv},\mathrm{lv}')=3\bar{m}^2+\mathrm{lv}\cdot\mathrm{lv}'$.
Moreover
\[
\mathrm{lph}_-\cdot\mathrm{lph}_+=\frac12\left(|\mathrm{lph}_-|_2^2+|\mathrm{lph}_+|_2^2-|\mathrm{kph}|_2^2\right)
=\frac12\left(\mathrm{den}(\mathrm{lph}_-)+\mathrm{den}(\mathrm{lph}_+)\right)-\bar{m}^2-\frac12\,|\mathrm{kph}|_2^2.
\]
Therefore
\[
\frac{\sum_\mu\mathrm{num}_{\mu\mu}(\mathrm{lph}_-,\mathrm{lph}_+)}{\mathrm{den}(\mathrm{lph}_-)\cdot\mathrm{den}(\mathrm{lph}_+)}
=\frac{2\bar{m}^2-|\mathrm{kph}|_2^2/2}{\mathrm{den}(\mathrm{lph}_-)\cdot\mathrm{den}(\mathrm{lph}_+)}
+\frac12\left(\frac{1}{\mathrm{den}(\mathrm{lph}_-)}+\frac{1}{\mathrm{den}(\mathrm{lph}_+)}\right),
\]
and at $\mathrm{kph}=0$ the right side equals $2\bar{m}^2/\mathrm{den}(\mathrm{lph})^2+1/\mathrm{den}(\mathrm{lph})$.
By (5.7) and (5.8), $\mathrm{tr}_3\,\mathrm{Vpc}(\mathrm{kph})=4\,\mathrm{Iq}_1+4\,\mathrm{Iq}_2$, where $\mathrm{lph}_\pm=\mathrm{lph}\pm\mathrm{kph}/2$ and
\begin{align*}
\mathrm{Iq}_1&:=\int_{\mathbb{R}^3}\frac{\mathrm{d}^3\mathrm{lph}}{(2\pi)^3}\left\{\frac{2\bar{m}^2-|\mathrm{kph}|_2^2/2}{\mathrm{den}(\mathrm{lph}_-)\cdot\mathrm{den}(\mathrm{lph}_+)}-\frac{2\bar{m}^2}{\mathrm{den}(\mathrm{lph})^2}\right\},\\
\mathrm{Iq}_2&:=\int_{\mathbb{R}^3}\frac{\mathrm{d}^3\mathrm{lph}}{(2\pi)^3}\left\{\frac12\left(\frac{1}{\mathrm{den}(\mathrm{lph}_-)}+\frac{1}{\mathrm{den}(\mathrm{lph}_+)}\right)-\frac{1}{\mathrm{den}(\mathrm{lph})}\right\}.
\end{align*}
Both integrals converge absolutely.
Each of the two terms in the integrand of $\mathrm{Iq}_1$ is bounded and is $O(|\mathrm{lph}|_2^{-4})$ as $|\mathrm{lph}|_2\to\infty$; the integrand of $\mathrm{Iq}_2$ is integrable by the bound that follows.

\emph{Proof that $\mathrm{Iq}_2=0$.}
One has $\partial_\beta(1/\mathrm{den})(\mathrm{lv})=-2\,\mathrm{lv}_\beta/\mathrm{den}(\mathrm{lv})^2$, of modulus at most $2\,|\mathrm{lv}|_2^{-3}$, and
\[
\partial_\alpha\partial_\beta(1/\mathrm{den})(\mathrm{lv})=-\frac{2\,\delta_{\alpha\beta}}{\mathrm{den}(\mathrm{lv})^2}+\frac{8\,\mathrm{lv}_\alpha\cdot\mathrm{lv}_\beta}{\mathrm{den}(\mathrm{lv})^3},
\]
of modulus at most $10/\mathrm{den}(\mathrm{lv})^2$.
The formula for second differences in (iv) shows that the integrand of $\mathrm{Iq}_2$ equals
\[
\frac18\int_{-1}^{1}\mathrm{d}\,\mathrm{sv}\,(1-|\mathrm{sv}|)\sum_{\alpha,\beta}\mathrm{kph}_\alpha\cdot\mathrm{kph}_\beta\cdot\left(\partial_\alpha\partial_\beta(1/\mathrm{den})\right)(\mathrm{lph}+\mathrm{sv}\cdot\mathrm{kph}/2).
\]
It is bounded, and for $|\mathrm{lph}|_2>|\mathrm{kph}|_2/2$ its modulus is at most $\frac{10}{8}\,|\mathrm{kph}|_1^2\,(|\mathrm{lph}|_2-|\mathrm{kph}|_2/2)^{-4}$; so it is integrable.
By Fubini's theorem on the compact set $\{|\mathrm{lph}|_2\le\mathrm{rad}\}\times[-1,1]$ and by the divergence theorem, its integral over $|\mathrm{lph}|_2<\mathrm{rad}$ equals
\[
\frac18\int_{-1}^{1}\mathrm{d}\,\mathrm{sv}\,(1-|\mathrm{sv}|)\sum_{\alpha,\beta}\mathrm{kph}_\alpha\cdot\mathrm{kph}_\beta\int_{|\mathrm{lph}|_2=\mathrm{rad}}\mathrm{d}S\,\frac{\mathrm{lph}_\alpha}{\mathrm{rad}}\,\left(\partial_\beta(1/\mathrm{den})\right)(\mathrm{lph}+\mathrm{sv}\cdot\mathrm{kph}/2),
\]
whose modulus is at most $\frac18\,|\mathrm{kph}|_1^2\cdot4\pi\,\mathrm{rad}^2\cdot2\,(\mathrm{rad}-|\mathrm{kph}|_2/2)^{-3}$.
This tends to $0$ as $\mathrm{rad}\to\infty$; the integrand being integrable, $\mathrm{Iq}_2=0$.

\emph{Computation of $\mathrm{Iq}_1$.}
For positive numbers $\mathrm{zq}_1,\mathrm{zq}_2$ one has
\[
\frac{1}{\mathrm{zq}_1\cdot\mathrm{zq}_2}=\int_0^1\mathrm{d}\,\mathrm{sv}\,\left[(1-\mathrm{sv})\,\mathrm{zq}_1+\mathrm{sv}\cdot\mathrm{zq}_2\right]^{-2}.
\]
Moreover
\[
(1-\mathrm{sv})\,\mathrm{den}(\mathrm{lph}_-)+\mathrm{sv}\cdot\mathrm{den}(\mathrm{lph}_+)
=\bar{m}^2+|\mathrm{lph}|_2^2+|\mathrm{kph}|_2^2/4+(2\,\mathrm{sv}-1)\,\mathrm{lph}\cdot\mathrm{kph}
=\left|\mathrm{lph}+(\mathrm{sv}-1/2)\,\mathrm{kph}\right|_2^2+\mathrm{De}(\mathrm{sv}),
\]
with $\mathrm{De}(\mathrm{sv}):=\bar{m}^2+\mathrm{sv}\cdot(1-\mathrm{sv})\cdot|\mathrm{kph}|_2^2$.
The substitution $\mathrm{sv}'=\mathrm{De}^{\frac12}\tan\mathrm{sv}''$ turns $\int_0^\infty\mathrm{d}\,\mathrm{sv}'\,\mathrm{sv}'^2\,(\mathrm{sv}'^2+\mathrm{De})^{-2}$ into $\mathrm{De}^{-\frac12}\int_0^{\pi/2}\mathrm{d}\,\mathrm{sv}''\,\sin^2\mathrm{sv}''=\mathrm{De}^{-\frac12}\,\pi/4$, so that
\[
\int_0^\infty\mathrm{d}\,\mathrm{sv}'\,\mathrm{sv}'^2\,(\mathrm{sv}'^2+\mathrm{De})^{-2}=\frac{\pi}{4\,\mathrm{De}^{\frac12}}.
\]
With Tonelli's theorem, the invariance of the Lebesgue measure under translations and polar coordinates one obtains
\[
\mathrm{Ib}(\mathrm{kph}):=\int_{\mathbb{R}^3}\frac{\mathrm{d}^3\mathrm{lph}}{(2\pi)^3}\,\frac{1}{\mathrm{den}(\mathrm{lph}_-)\cdot\mathrm{den}(\mathrm{lph}_+)}
=\int_0^1\mathrm{d}\,\mathrm{sv}\,\frac{4\pi}{(2\pi)^3}\cdot\frac{\pi}{4\,\mathrm{De}(\mathrm{sv})^{\frac12}}
=\int_0^1\frac{\mathrm{d}\,\mathrm{sv}}{8\pi\,\mathrm{De}(\mathrm{sv})^{\frac12}},
\qquad\mathrm{Ib}(0)=\frac{1}{8\pi\bar{m}}.
\]
Since $\mathrm{Iq}_2=0$, it follows that
\[
\mathrm{tr}_3\,\mathrm{Vpc}(\mathrm{kph})=4\left[\left(2\bar{m}^2-|\mathrm{kph}|_2^2/2\right)\mathrm{Ib}(\mathrm{kph})-2\bar{m}^2\,\mathrm{Ib}(0)\right]
=\frac{1}{2\pi}\int_0^1\mathrm{d}\,\mathrm{sv}\,\left[\left(2\bar{m}^2-|\mathrm{kph}|_2^2/2\right)\mathrm{De}(\mathrm{sv})^{-\frac12}-2\bar{m}\right].
\]
Since $\frac{\mathrm{d}}{\mathrm{d}\,\mathrm{sv}}\mathrm{De}(\mathrm{sv})=(1-2\,\mathrm{sv})\,|\mathrm{kph}|_2^2$ and $(\mathrm{sv}-1/2)^2\,|\mathrm{kph}|_2^2=|\mathrm{kph}|_2^2/4+\bar{m}^2-\mathrm{De}(\mathrm{sv})$,
\[
\frac{\mathrm{d}}{\mathrm{d}\,\mathrm{sv}}\left[(\mathrm{sv}-1/2)\,\mathrm{De}(\mathrm{sv})^{\frac12}\right]
=\mathrm{De}(\mathrm{sv})^{\frac12}-(\mathrm{sv}-1/2)^2\,|\mathrm{kph}|_2^2\,\mathrm{De}(\mathrm{sv})^{-\frac12}
=2\,\mathrm{De}(\mathrm{sv})^{\frac12}-\left(\bar{m}^2+|\mathrm{kph}|_2^2/4\right)\mathrm{De}(\mathrm{sv})^{-\frac12}.
\]
Integration over $[0,1]$, with $\mathrm{De}(0)=\mathrm{De}(1)=\bar{m}^2$, gives
\[
\bar{m}=\int_0^1\mathrm{d}\,\mathrm{sv}\,\left[2\,\mathrm{De}(\mathrm{sv})^{\frac12}-\left(\bar{m}^2+|\mathrm{kph}|_2^2/4\right)\mathrm{De}(\mathrm{sv})^{-\frac12}\right].
\]
Inserting $2\bar{m}=\int_0^1\mathrm{d}\,\mathrm{sv}\,\left[4\,\mathrm{De}^{\frac12}-\left(2\bar{m}^2+|\mathrm{kph}|_2^2/2\right)\mathrm{De}^{-\frac12}\right]$ into the last expression for the trace, one finds
\begin{equation}
\begin{aligned}
\mathrm{tr}_3\,\mathrm{Vpc}(\mathrm{kph})
&=\frac{1}{2\pi}\int_0^1\mathrm{d}\,\mathrm{sv}\,\left(4\bar{m}^2-4\,\mathrm{De}(\mathrm{sv})\right)\mathrm{De}(\mathrm{sv})^{-\frac12}\\
&=-\frac{2\,|\mathrm{kph}|_2^2}{\pi}\int_0^1\mathrm{d}\,\mathrm{sv}\,\mathrm{sv}\cdot(1-\mathrm{sv})\cdot\left(\bar{m}^2+\mathrm{sv}\cdot(1-\mathrm{sv})\cdot|\mathrm{kph}|_2^2\right)^{-\frac12}
=-2\,|\mathrm{kph}|_2^2\,\mathrm{vpo}_{\bar{m}}(|\mathrm{kph}|_2^2).
\end{aligned}
\tag{5.11}
\end{equation}

(vi) The identities (5.10) and (5.11) give (5.6).
\qed

\paragraph{5.5. Conclusion.}
By 5.2 and 5.3, $\mathrm{Vpl}^{\eta}_{\mu\nu}\to\mathrm{Vpc}_{\mu\nu}$ as $\eta\to0$, uniformly on every ball, and by (5.6) the limit is the function in clause (d).
This proves (d).
The limit is reached without subtracting any term from $\mathrm{Vpl}^{\eta}$, and the mass in $D^{(p)}$ is $\mathbf{m}=\bar{m}\eta$.
\qed

\subsubsection*{6. Clause (e)}

\begin{named}{Lemma 6.1}[lattice transform of a smooth function]
Let $\eta\in(0,1]$, let $\mathrm{Lsh}$ be a translate of $\eta\mathbb{Z}^3$ in $\mathbb{R}^3$, let $y^0\in\mathbb{R}^3$ and $\mathrm{sl}\ge0$, and let $\mathrm{fs}\in C_c^\infty(\mathbb{R}^3)$ be a complex function that vanishes outside a closed coordinate cube of side $\mathrm{sl}$.
Put
\[
\mathrm{wh}(\mathrm{kph}):=\eta^3\sum_{y\in\mathrm{Lsh}}\mathrm{fs}(y)\exp\left(-\mathrm{i}\,\mathrm{kph}\cdot(y-y^0)\right)\qquad(\text{a finite sum}).
\]
Then for every integer $i\ge0$ and every $\mathrm{kph}\in\mathbb{R}^3$ with $|\mathrm{kph}|_\infty\le3\pi/(2\eta)$
\begin{equation}
|\mathrm{wh}(\mathrm{kph})|\le12^i\,(\mathrm{sl}+i+1)^3\max\left\{|\mathrm{fs}|_\infty,\ \max_\mu|\partial_\mu^i\mathrm{fs}|_\infty\right\}(1+|\mathrm{kph}|_2)^{-i}.
\tag{6.1}
\end{equation}
\end{named}

\begin{proof}
(1) A closed interval of length $\mathrm{sv}$ contains at most $\mathrm{sv}/\eta+1$ points of a translate of $\eta\mathbb{Z}$.
Hence a closed box with the sides $\mathrm{sl}+i\eta$, $\mathrm{sl}$, $\mathrm{sl}$ contains at most $(\mathrm{sl}/\eta+i+1)(\mathrm{sl}/\eta+1)^2$ points of $\mathrm{Lsh}$, and $\eta^3$ times this number is at most $(\mathrm{sl}+(i+1)\eta)(\mathrm{sl}+\eta)^2\le(\mathrm{sl}+i+1)^3$.

(2) For a function $\mathrm{fw}$ on $\mathbb{R}^3$ let $(\mathrm{Db}_\mu\mathrm{fw})(y):=\left(\mathrm{fw}(y)-\mathrm{fw}(y-\eta e_\mu)\right)/\eta$.
For smooth $\mathrm{fw}$ one has $(\mathrm{Db}_\mu\mathrm{fw})(y)=\int_0^1\mathrm{d}\,\mathrm{sv}\,(\partial_\mu\mathrm{fw})(y-\eta\,\mathrm{sv}\,e_\mu)$, and $\mathrm{Db}_\mu$ commutes with $\partial_\mu$ and with translations.
By induction,
\[
(\mathrm{Db}_\mu^i\mathrm{fs})(y)=\int_{[0,1]^i}\mathrm{d}\,\mathrm{sv}_1\cdots\mathrm{d}\,\mathrm{sv}_i\,(\partial_\mu^i\mathrm{fs})\left(y-\eta\,(\mathrm{sv}_1+\dots+\mathrm{sv}_i)\,e_\mu\right).
\]
Therefore $|\mathrm{Db}_\mu^i\mathrm{fs}|_\infty\le|\partial_\mu^i\mathrm{fs}|_\infty$, and $(\mathrm{Db}_\mu^i\mathrm{fs})(y)\neq0$ only if $y$ lies in the box obtained from the cube of $\mathrm{fs}$ by lengthening it by $i\eta$ in the direction $+e_\mu$.

(3) The set $\mathrm{Lsh}$ is invariant under $y\mapsto y+\eta e_\mu$.
For a function $\mathrm{fw}$ that vanishes outside a bounded set this gives
\[
\sum_{y\in\mathrm{Lsh}}\mathrm{fw}(y-\eta e_\mu)\exp\left(-\mathrm{i}\,\mathrm{kph}\cdot(y-y^0)\right)
=\exp(-\mathrm{i}\eta\,\mathrm{kph}_\mu)\sum_{y\in\mathrm{Lsh}}\mathrm{fw}(y)\exp\left(-\mathrm{i}\,\mathrm{kph}\cdot(y-y^0)\right),
\]
and hence
\[
\eta^3\sum_{y\in\mathrm{Lsh}}(\mathrm{Db}_\mu\mathrm{fw})(y)\exp\left(-\mathrm{i}\,\mathrm{kph}\cdot(y-y^0)\right)
=\eta^{-1}\left(1-\exp(-\mathrm{i}\eta\,\mathrm{kph}_\mu)\right)\eta^3\sum_{y\in\mathrm{Lsh}}\mathrm{fw}(y)\exp\left(-\mathrm{i}\,\mathrm{kph}\cdot(y-y^0)\right).
\]
Applying this identity $i$ times and then (1) and (2), one obtains
\begin{equation}
\left|\eta^{-1}\left(1-\exp(-\mathrm{i}\eta\,\mathrm{kph}_\mu)\right)\right|^i|\mathrm{wh}(\mathrm{kph})|
=\left|\eta^3\sum_{y\in\mathrm{Lsh}}(\mathrm{Db}_\mu^i\mathrm{fs})(y)\exp\left(-\mathrm{i}\,\mathrm{kph}\cdot(y-y^0)\right)\right|
\le(\mathrm{sl}+i+1)^3\,|\partial_\mu^i\mathrm{fs}|_\infty.
\tag{6.2}
\end{equation}

(4) One has $|1-\exp(-\mathrm{i}\,\mathrm{sv})|=2\,|\sin(\mathrm{sv}/2)|$.
The sine is concave on $[0,\pi]$; therefore $\sin\mathrm{sv}'\ge\mathrm{sv}'\cdot\sin(3\pi/4)/(3\pi/4)$ for $0\le\mathrm{sv}'\le3\pi/4$, and $\sin(3\pi/4)/(3\pi/4)=2\sqrt2/(3\pi)>3/10$ by (Ef0)\,(vi).
Hence $\left|\eta^{-1}\left(1-\exp(-\mathrm{i}\eta\,\mathrm{kph}_\mu)\right)\right|\ge\frac{3}{10}\,|\mathrm{kph}_\mu|$ whenever $\eta\,|\mathrm{kph}_\mu|\le3\pi/2$.

(5) Let $|\mathrm{kph}|_2\le1$.
Then (6.2) with $i=0$ gives $|\mathrm{wh}(\mathrm{kph})|\le(\mathrm{sl}+1)^3\,|\mathrm{fs}|_\infty$, and $1\le2^i\,(1+|\mathrm{kph}|_2)^{-i}$; this implies (6.1).
Let $|\mathrm{kph}|_2\ge1$.
Choose $\mu$ with $|\mathrm{kph}_\mu|=|\mathrm{kph}|_\infty$; then $|\mathrm{kph}_\mu|\ge|\mathrm{kph}|_2/\sqrt3\ge(1+|\mathrm{kph}|_2)/(2\sqrt3)$, and $\eta\,|\mathrm{kph}_\mu|\le3\pi/2$ by hypothesis.
By (4),
\[
\left|\eta^{-1}\left(1-\exp(-\mathrm{i}\eta\,\mathrm{kph}_\mu)\right)\right|\ge\frac{3}{20\sqrt3}\,(1+|\mathrm{kph}|_2)\ge\frac{1+|\mathrm{kph}|_2}{12},
\]
because $20\sqrt3/3<12$ by (Ef0)\,(vii); now (6.2) gives (6.1).
\end{proof}

\paragraph{6.2. The decay.}
Write $\mathrm{gs}_\nu:=\mathrm{gs}_\nu[\beta]$, formed with $\eta$.
Replacing $y+(\eta/2)e_\nu$ by $y$ in the definitions of $\mathrm{a0}_{\eta,\nu}$ and $\mathrm{ahl}_{\eta,\nu}$ in Proposition~36, one obtains
\begin{equation}
\mathrm{ahl}_{\eta,\nu}(\mathrm{kph})=\eta^3\sum_{y\in\eta\mathbb{Z}^3+(\eta/2)e_\nu}\mathrm{gs}_\nu(y)\exp(-\mathrm{i}\,\mathrm{kph}\cdot y).
\tag{6.3}
\end{equation}
By the definition of $D_{\mathrm{be}}$ there is a point $y^0\in\mathbb{R}^3$ with $\operatorname{supp}\beta\subset y^0+[-D_{\mathrm{be}}/2,D_{\mathrm{be}}/2]^3$; it is fixed from now on.
By the last sentence of Proposition~35.3\,(b1), the norm $|D^i\,\cdot\,|_\infty$ being that of Section~G5.1.3, one has $|\partial_\mu^i\mathrm{gs}_\nu|_\infty\le|D^i\mathrm{gs}_\nu|_\infty\le2\,|\beta|_{C^{i+1}}$ for $i\ge0$.
By the same sentence, $\mathrm{gs}_\nu$ vanishes outside $y^0+[-D_{\mathrm{be}}/2-\eta,D_{\mathrm{be}}/2+\eta]^3$, a closed coordinate cube of side $D_{\mathrm{be}}+2\eta\le D_{\mathrm{be}}+2$, hence also outside the concentric closed cube of side $D_{\mathrm{be}}+2$.
Moreover $|\beta|_{C^1}\le|\beta|_{C^{i+1}}$.
Lemma~6.1, applied with $\mathrm{Lsh}=\eta\mathbb{Z}^3+(\eta/2)e_\nu$, $\mathrm{sl}=D_{\mathrm{be}}+2$, $\mathrm{fs}=\mathrm{gs}_\nu$ and the point $0$ in place of $y^0$, gives for all $\eta\in(0,1]$, all $i\ge0$ and all $\mathrm{kph}$ with $|\mathrm{kph}|_\infty\le3\pi/(2\eta)$
\begin{equation}
|\mathrm{ahl}_{\eta,\nu}(\mathrm{kph})|\le\mathrm{Cah}_i(\beta)\cdot(1+|\mathrm{kph}|_2)^{-i},\qquad
\mathrm{Cah}_i(\beta):=2\cdot12^i\,(D_{\mathrm{be}}+i+3)^3\,|\beta|_{C^{i+1}}.
\tag{6.4}
\end{equation}

\paragraph{6.3. The centred transform and its gradient.}
Put
\[
\mathrm{ahc}_\nu(\mathrm{kph}):=\exp(\mathrm{i}\,\mathrm{kph}\cdot y^0)\,\mathrm{ahl}_{\eta,\nu}(\mathrm{kph})
=\eta^3\sum_{y\in\eta\mathbb{Z}^3+(\eta/2)e_\nu}\mathrm{gs}_\nu(y)\exp\left(-\mathrm{i}\,\mathrm{kph}\cdot(y-y^0)\right).
\]
Then
\begin{equation}
\overline{\mathrm{ahc}_\mu(\mathrm{kph})}\,\mathrm{ahc}_\nu(\mathrm{kph})=\overline{\mathrm{ahl}_{\eta,\mu}(\mathrm{kph})}\,\mathrm{ahl}_{\eta,\nu}(\mathrm{kph}),\qquad
|\mathrm{ahc}_\nu(\mathrm{kph})|=|\mathrm{ahl}_{\eta,\nu}(\mathrm{kph})|,
\tag{6.5}
\end{equation}
and
\[
\partial_{\mathrm{kph}_{\mathfrak{l}}}\mathrm{ahc}_\nu(\mathrm{kph})=-\mathrm{i}\,\eta^3\sum_y(y-y^0)_{\mathfrak{l}}\,\mathrm{gs}_\nu(y)\exp\left(-\mathrm{i}\,\mathrm{kph}\cdot(y-y^0)\right)\qquad(\text{a finite sum}).
\]
The function $\mathrm{fs}(y):=(y-y^0)_{\mathfrak{l}}\,\mathrm{gs}_\nu(y)$ vanishes outside the same cube, on this cube $|(y-y^0)_{\mathfrak{l}}|\le D_{\mathrm{be}}/2+1$, and the Leibniz rule gives
\begin{equation}
\begin{aligned}
\partial_\mu^i\mathrm{fs}&=(y-y^0)_{\mathfrak{l}}\,\partial_\mu^i\mathrm{gs}_\nu+i\,\delta_{\mathfrak{l}\mu}\,\partial_\mu^{i-1}\mathrm{gs}_\nu,\\
|\partial_\mu^i\mathrm{fs}|_\infty&\le(D_{\mathrm{be}}/2+1)\cdot2\,|\beta|_{C^{i+1}}+2i\,|\beta|_{C^i}\le(D_{\mathrm{be}}+2+2i)\,|\beta|_{C^{i+1}}
\end{aligned}
\tag{6.6}
\end{equation}
(for $i=0$ the second term is absent).
Moreover $|\mathrm{fs}|_\infty\le(D_{\mathrm{be}}+2)\,|\beta|_{C^1}$.
Lemma~6.1 gives, in the same range of $\eta$, $i$ and $\mathrm{kph}$,
\begin{equation}
|\partial_{\mathrm{kph}_{\mathfrak{l}}}\mathrm{ahc}_\nu(\mathrm{kph})|\le\mathrm{Cahd}_i(\beta)\cdot(1+|\mathrm{kph}|_2)^{-i},\qquad
\mathrm{Cahd}_i(\beta):=12^i\,(D_{\mathrm{be}}+i+3)^3\,(D_{\mathrm{be}}+2+2i)\,|\beta|_{C^{i+1}}.
\tag{6.7}
\end{equation}
The numbers $\mathrm{Cah}_i(\beta)$ and $\mathrm{Cahd}_i(\beta)$ depend on $\beta$ only through $D_{\mathrm{be}}$ and the norms $|\beta|_{C^j}$.

\paragraph{6.4. The limit.}
By the definitions in Section~I.7.6 and Section~G5.1.3,
\begin{gather*}
(J_\eta\mathrm{fs})(y)=\int_{[-1/2,1/2]^3}\mathrm{d}^3\mathrm{sv}'\,\mathrm{fs}(y+\eta\,\mathrm{sv}')\quad(\mathrm{sv}'\text{ a vector in }[-1/2,1/2]^3),\\
(\delta_{\eta,\mu}\mathrm{fs})(y)=\int_{-1/2}^{1/2}\mathrm{d}\,\mathrm{sv}\,(\partial_\mu\mathrm{fs})(y+\eta\,\mathrm{sv}\,e_\mu).
\end{gather*}
Hence
\[
\mathrm{gs}_\nu(y)=-\sum_{\mu\neq\nu}\int_{-1/2}^{1/2}\mathrm{d}\,\mathrm{sv}\int_{[-1/2,1/2]^3}\mathrm{d}^3\mathrm{sv}'\,(\partial_\mu\beta_{\mu\nu})(y+\eta\,\mathrm{sv}\,e_\mu+\eta\,\mathrm{sv}'),
\]
whereas $a_\nu(y)=-\sum_{\mu\neq\nu}(\partial_\mu\beta_{\mu\nu})(y)$, as $\beta_{\nu\nu}=0$.
For a function $\mathrm{fs}$ of class $C^1$ one has $|\mathrm{fs}(y+y')-\mathrm{fs}(y)|\le|y'|_1\max_{\mathfrak{l}}|\partial_{\mathfrak{l}}\mathrm{fs}|_\infty$, and $|\eta\,\mathrm{sv}\,e_\mu+\eta\,\mathrm{sv}'|_1\le2\eta$.
Two values of $\mu$ occur; therefore
\begin{equation}
|\mathrm{gs}_\nu(y)-a_\nu(y)|\le4\eta\,|\beta|_{C^2}\qquad\text{for all }y\in\mathbb{R}^3.
\tag{6.8}
\end{equation}
Let $\mathrm{Lsh}=\eta\mathbb{Z}^3+(\eta/2)e_\nu$ and $\mathrm{kph}\in\mathbb{R}^3$.
The function $\mathrm{gs}_\nu-a_\nu$ vanishes outside the cube of 6.2; by step (1) of the proof of Lemma~6.1, with $i=0$ and $\mathrm{sl}=D_{\mathrm{be}}+2\eta$, and by (6.8),
\[
\left|\eta^3\sum_{y\in\mathrm{Lsh}}\left(\mathrm{gs}_\nu(y)-a_\nu(y)\right)\exp(-\mathrm{i}\,\mathrm{kph}\cdot y)\right|\le4\eta\,|\beta|_{C^2}\,(D_{\mathrm{be}}+3)^3.
\]
Let $\mathrm{fs}(y):=a_\nu(y)\exp(-\mathrm{i}\,\mathrm{kph}\cdot y)$; then $|\partial_{\mathfrak{l}}\mathrm{fs}|_\infty\le2\,|\beta|_{C^2}+2\,|\mathrm{kph}|_\infty\,|\beta|_{C^1}$.
The cells $y+[-\eta/2,\eta/2)^3$, $y\in\mathrm{Lsh}$, tile $\mathbb{R}^3$, and
\[
\left|\eta^3\,\mathrm{fs}(y)-\int_{y+[-\eta/2,\eta/2)^3}\mathrm{d}y'\,\mathrm{fs}(y')\right|\le\eta^3\,(3\eta/2)\max_{\mathfrak{l}}|\partial_{\mathfrak{l}}\mathrm{fs}|_\infty.
\]
The cell of $y$ meets $\operatorname{supp}a$ only if $y\in y^0+[-(D_{\mathrm{be}}+\eta)/2,(D_{\mathrm{be}}+\eta)/2]^3$, and $\eta^3$ times the number of these $y$ is at most $(D_{\mathrm{be}}+2\eta)^3$.
Hence
\[
\left|\eta^3\sum_{y\in\mathrm{Lsh}}\mathrm{fs}(y)-\hat{a}_\nu(\mathrm{kph})\right|\le3\eta\,(D_{\mathrm{be}}+2)^3\left(|\beta|_{C^2}+|\mathrm{kph}|_\infty\,|\beta|_{C^1}\right).
\]
Together with (6.3) this gives
\begin{equation}
|\mathrm{ahl}_{\eta,\nu}(\mathrm{kph})-\hat{a}_\nu(\mathrm{kph})|\le7\eta\,(D_{\mathrm{be}}+3)^3\,(1+|\mathrm{kph}|_\infty)\,|\beta|_{C^2}\qquad\text{for all }\eta\in(0,1],\ \mathrm{kph}\in\mathbb{R}^3.
\tag{6.9}
\end{equation}
In particular $\mathrm{ahl}_{\eta,\nu}(\mathrm{kph})\to\hat{a}_\nu(\mathrm{kph})$ as $\eta\to0$.

\paragraph{6.5. Transversality.}
Since $\beta$ has compact support, integration by parts gives
\begin{gather*}
\hat{a}_\nu(\mathrm{kph})=-\sum_\mu\int\mathrm{d}y\,(\partial_\mu\beta_{\mu\nu})(y)\exp(-\mathrm{i}\,\mathrm{kph}\cdot y)=-\mathrm{i}\sum_\mu\mathrm{kph}_\mu\,\hat{\beta}_{\mu\nu}(\mathrm{kph}),\\
\hat{\beta}_{\mu\nu}(\mathrm{kph}):=\int\mathrm{d}y\,\beta_{\mu\nu}(y)\exp(-\mathrm{i}\,\mathrm{kph}\cdot y).
\end{gather*}
Consequently $\sum_\nu\mathrm{kph}_\nu\,\hat{a}_\nu(\mathrm{kph})=-\mathrm{i}\sum_{\mu,\nu}\mathrm{kph}_\mu\cdot\mathrm{kph}_\nu\,\hat{\beta}_{\mu\nu}(\mathrm{kph})=0$, because $\beta_{\nu\mu}=-\beta_{\mu\nu}$.
The bounds (6.4), (6.9) and this identity prove clause (e).
\qed

\subsubsection*{7. Clause (f)}

Throughout this section $\eta\in(0,1]$ and $R\ge1$, and $n=R/\eta$ is an integer $\ge2$.

\paragraph{7.1. The summand and its limit.}
For $\mathrm{kph}\in\mathbb{R}^3$ put
\[
\mathrm{Ipv}_\eta(\mathrm{kph}):=\sum_{\mu,\nu}\overline{\mathrm{ahl}_{\eta,\mu}(\mathrm{kph})}\,\mathrm{ahl}_{\eta,\nu}(\mathrm{kph})\,\mathrm{Vpl}^{\eta}_{\mu\nu}(\mathrm{kph}),\qquad
\mathrm{Ipv}_0(\mathrm{kph}):=\sum_{\mu,\nu}\overline{\hat{a}_\mu(\mathrm{kph})}\,\hat{a}_\nu(\mathrm{kph})\,\mathrm{Vpc}_{\mu\nu}(\mathrm{kph}).
\]
By clauses (d) and (e), $\mathrm{Ipv}_\eta(\mathrm{kph})\to\mathrm{Ipv}_0(\mathrm{kph})$ as $\eta\to0$, for every $\mathrm{kph}$.
By (5.6), by the identity $\sum_{\mu,\nu}\overline{\hat{a}_\mu}\,\hat{a}_\nu\,\mathrm{kph}_\mu\cdot\mathrm{kph}_\nu=\left|\sum_\nu\mathrm{kph}_\nu\,\hat{a}_\nu\right|^2$ and by 6.5,
\begin{equation}
\mathrm{Ipv}_0(\mathrm{kph})=-\mathrm{vpo}_{\bar{m}}(|\mathrm{kph}|_2^2)\left(|\mathrm{kph}|_2^2\,|\hat{a}(\mathrm{kph})|_2^2-\left|\sum_\nu\mathrm{kph}_\nu\,\hat{a}_\nu(\mathrm{kph})\right|^2\right)
=-|\mathrm{kph}|_2^2\,\mathrm{vpo}_{\bar{m}}(|\mathrm{kph}|_2^2)\,|\hat{a}(\mathrm{kph})|_2^2.
\tag{7.1}
\end{equation}

\paragraph{7.2. The majorant.}
Put
\begin{equation}
\begin{aligned}
&\mathrm{Cmj}:=9\sqrt3\,\mathrm{Cpv}_0\cdot\bar{m}^{-1}\cdot\mathrm{Cah}_3(\beta)\cdot\left(2\,\mathrm{Cahd}_3(\beta)+\mathrm{Cah}_3(\beta)\right),\qquad
\mathrm{Mpv}(\mathrm{kph}):=\mathrm{Cmj}\cdot(1+|\mathrm{kph}|_2)^{-4};\\
&\int_{\mathbb{R}^3}\frac{\mathrm{d}^3\mathrm{kph}}{(2\pi)^3}\,\mathrm{Mpv}(\mathrm{kph})=\mathrm{Cmj}\cdot\frac{4\pi}{3}\cdot\frac{1}{(2\pi)^3}=\frac{\mathrm{Cmj}}{6\pi^2}
\end{aligned}
\tag{7.2}
\end{equation}
by (0.4) with $\mathrm{mq}=1$.

\begin{named}{Claim}
For all $\mathrm{kph}$ with $|\mathrm{kph}|_\infty\le3\pi/(2\eta)$ one has $|\mathrm{Ipv}_\eta(\mathrm{kph})|\le\mathrm{Mpv}(\mathrm{kph})$ and $|\mathrm{grad}\,\mathrm{Ipv}_\eta(\mathrm{kph})|_2\le\mathrm{Mpv}(\mathrm{kph})$.
Moreover $|\mathrm{Ipv}_0(\mathrm{kph})|\le\mathrm{Mpv}(\mathrm{kph})$ for all $\mathrm{kph}\in\mathbb{R}^3$.
\end{named}

\begin{proof}
By (6.5), $\mathrm{Ipv}_\eta=\sum_{\mu,\nu}\overline{\mathrm{ahc}_\mu}\,\mathrm{ahc}_\nu\,\mathrm{Vpl}^{\eta}_{\mu\nu}$, a sum of $9$ terms.
By (6.4) with $i=3$ and by clause (c),
\[
|\mathrm{Ipv}_\eta(\mathrm{kph})|\le9\,\mathrm{Cah}_3^2\cdot(1+|\mathrm{kph}|_2)^{-6}\cdot\mathrm{Cpv}_0\cdot|\mathrm{kph}|_2^2/\bar{m}\le9\,\mathrm{Cah}_3^2\cdot\mathrm{Cpv}_0\cdot\bar{m}^{-1}\cdot(1+|\mathrm{kph}|_2)^{-4}.
\]
As $\mathrm{kph}$ is real, $\partial_{\mathrm{kph}_{\mathfrak{l}}}\overline{\mathrm{ahc}_\mu}=\overline{\partial_{\mathrm{kph}_{\mathfrak{l}}}\mathrm{ahc}_\mu}$.
The product rule, (6.4) and (6.7) with $i=3$, and clause (c) give
\begin{align*}
|\partial_{\mathrm{kph}_{\mathfrak{l}}}\mathrm{Ipv}_\eta(\mathrm{kph})|
&\le9\,(1+|\mathrm{kph}|_2)^{-6}\cdot\mathrm{Cpv}_0\cdot\bar{m}^{-1}\left[2\,\mathrm{Cah}_3\cdot\mathrm{Cahd}_3\cdot|\mathrm{kph}|_2^2+\mathrm{Cah}_3^2\cdot|\mathrm{kph}|_2\right]\\
&\le9\,\mathrm{Cpv}_0\cdot\bar{m}^{-1}\cdot\mathrm{Cah}_3\cdot\left(2\,\mathrm{Cahd}_3+\mathrm{Cah}_3\right)\cdot(1+|\mathrm{kph}|_2)^{-4};
\end{align*}
the Euclidean norm of the gradient, that is, the square root of the sum of the squared moduli of the three partial derivatives, is at most $\sqrt3$ times the largest of them.
This proves the first two bounds.
For fixed $\mathrm{kph}$ the first bound holds for all $\eta\in(0,1]$ with $\eta\,|\mathrm{kph}|_\infty\le3\pi/2$; letting $\eta\to0$ and using 7.1, one obtains the third.
\end{proof}

From (7.1), (7.2), the third bound of the Claim and the definition of $\mathrm{pv}(f)$ it follows that
\begin{equation}
\mathrm{Ipv}_0\ \text{is integrable},\qquad
\int_{\mathbb{R}^3}\frac{\mathrm{d}^3\mathrm{kph}}{(2\pi)^3}\,\mathrm{Ipv}_0(\mathrm{kph})=\mathrm{pv}(f),\qquad
|\mathrm{pv}(f)|\le\frac{\mathrm{Cmj}}{6\pi^2}.
\tag{7.3}
\end{equation}

\paragraph{7.3. Cubes.}
For $\mathrm{kph}_0\in(2\pi/R)\mathbb{Z}^3$ let $\mathrm{Cu}(\mathrm{kph}_0):=\mathrm{kph}_0+[-\pi/R,\pi/R)^3$.
These cubes tile $\mathbb{R}^3$, and $R^{-3}$ is $(2\pi)^{-3}$ times the volume of each of them.
Let $\mathrm{kph}$ lie in the closure of $\mathrm{Cu}(\mathrm{kph}_0)$, and let $\mathrm{kph}'$ lie on the segment from $\mathrm{kph}_0$ to $\mathrm{kph}$, the point $\mathrm{kph}'=\mathrm{kph}_0$ being admitted.
Then $|\mathrm{kph}'-\mathrm{kph}|_2\le|\mathrm{kph}_0-\mathrm{kph}|_2\le\sqrt3\,\pi/R\le\sqrt3\,\pi<6$, by $R\ge1$ and (Ef0)\,(i); hence $1+|\mathrm{kph}|_2\le1+|\mathrm{kph}'|_2+6\le7\,(1+|\mathrm{kph}'|_2)$.
Therefore, integrating the first of the following inequalities with $\mathrm{kph}'=\mathrm{kph}_0$ over $\mathrm{Cu}(\mathrm{kph}_0)$ and summing over $\mathrm{kph}_0$,
\begin{equation}
\begin{aligned}
&(1+|\mathrm{kph}'|_2)^{-4}\le7^4\,(1+|\mathrm{kph}|_2)^{-4}\quad\text{for such }\mathrm{kph},\ \mathrm{kph}',\\
&R^{-3}\sum_{\mathrm{kph}_0\in(2\pi/R)\mathbb{Z}^3}(1+|\mathrm{kph}_0|_2)^{-4}\le7^4\int_{\mathbb{R}^3}\frac{\mathrm{d}^3\mathrm{kph}}{(2\pi)^3}\,(1+|\mathrm{kph}|_2)^{-4}=\frac{7^4}{6\pi^2}.
\end{aligned}
\tag{7.4}
\end{equation}
Let $\mathrm{Ucu}$ be the union of the cubes $\mathrm{Cu}(\mathrm{kph}_0)$, $\mathrm{kph}_0\in\mathrm{Kx}_{\eta,R}$, and let $\mathrm{Cq}:=[-\pi/(2\eta),\pi/(2\eta)]^3$.
The set $\mathrm{Ucu}$ is the product of three copies of the interval $\left[(2\pi/R)(j_{\min}-1/2),(2\pi/R)(j_{\max}+1/2)\right)$, where $j_{\min}$ and $j_{\max}$ are the least and the largest integer in $(-n/2,n/2]$.
For even $n$ this interval is $[-\pi/\eta+\pi/R,\pi/\eta+\pi/R)$, for odd $n$ it is $[-\pi/\eta,\pi/\eta)$.
As $n\ge2$, one has $\pi/R\le\pi/(2\eta)$.
Hence
\begin{equation}
\mathrm{Cq}\subset\mathrm{Ucu},\qquad
\text{the closure of }\mathrm{Ucu}\text{ lies in }[-3\pi/(2\eta),3\pi/(2\eta)]^3,\qquad
|\mathrm{kph}|_\infty\le\pi/\eta\ \text{ for }\mathrm{kph}\in\mathrm{Kx}_{\eta,R}.
\tag{7.5}
\end{equation}

\paragraph{7.4. Step 1: $\mathrm{Vpl}^{\eta,R}$ is replaced by $\mathrm{Vpl}^{\eta}$.}
By (b), the first bound of (c), (6.4) with $i=2$ (admissible by (7.5)) and (7.4),
\begin{equation}
\begin{aligned}
\left|a_{\mathrm{ch}}(\eta,R)-R^{-3}\sum_{\mathrm{kph}\in\mathrm{Kx}_{\eta,R}}\mathrm{Ipv}_\eta(\mathrm{kph})\right|
&\le\mathrm{Cpv}_0\cdot\bar{m}^{-3}R^{-4}\cdot R^{-3}\sum_{\mathrm{kph}\in\mathrm{Kx}_{\eta,R}}9\,\mathrm{Cah}_2^2\cdot(1+|\mathrm{kph}|_2)^{-4}
\le\mathrm{Cpv}_1\cdot\bar{m}^{-3}R^{-4},\\
\mathrm{Cpv}_1&:=365\,\mathrm{Cpv}_0\cdot\mathrm{Cah}_2(\beta)^2,
\end{aligned}
\tag{7.6}
\end{equation}
since $9\cdot7^4/(6\pi^2)<365$ by (Ef0)\,(viii).

\paragraph{7.5. Step 2: the sum over $\mathrm{Kx}_{\eta,R}$ is replaced by an integral.}
Let $\mathrm{kph}_0\in\mathrm{Kx}_{\eta,R}$ and $\mathrm{kph}\in\mathrm{Cu}(\mathrm{kph}_0)$.
The segment from $\mathrm{kph}_0$ to $\mathrm{kph}$ lies in the closure of $\mathrm{Ucu}$, a box, where 7.2 applies by (7.5).
One has
\[
\mathrm{Ipv}_\eta(\mathrm{kph})-\mathrm{Ipv}_\eta(\mathrm{kph}_0)=\int_0^1\mathrm{d}\,\mathrm{sv}\,\sum_{\mathfrak{l}}(\mathrm{kph}-\mathrm{kph}_0)_{\mathfrak{l}}\,(\partial_{\mathrm{kph}_{\mathfrak{l}}}\mathrm{Ipv}_\eta)\left(\mathrm{kph}_0+\mathrm{sv}\cdot(\mathrm{kph}-\mathrm{kph}_0)\right);
\]
by the Cauchy--Schwarz inequality, 7.2 and (7.4), $|\mathrm{Ipv}_\eta(\mathrm{kph})-\mathrm{Ipv}_\eta(\mathrm{kph}_0)|\le(\sqrt3\,\pi/R)\cdot7^4\cdot\mathrm{Mpv}(\mathrm{kph})$.
Together with (7.2) this yields
\begin{equation}
\begin{aligned}
&\left|R^{-3}\sum_{\mathrm{kph}_0\in\mathrm{Kx}_{\eta,R}}\mathrm{Ipv}_\eta(\mathrm{kph}_0)-\int_{\mathrm{Ucu}}\frac{\mathrm{d}^3\mathrm{kph}}{(2\pi)^3}\,\mathrm{Ipv}_\eta(\mathrm{kph})\right|
\le\sum_{\mathrm{kph}_0\in\mathrm{Kx}_{\eta,R}}\int_{\mathrm{Cu}(\mathrm{kph}_0)}\frac{\mathrm{d}^3\mathrm{kph}}{(2\pi)^3}\,\left|\mathrm{Ipv}_\eta(\mathrm{kph}_0)-\mathrm{Ipv}_\eta(\mathrm{kph})\right|\\
&\qquad\le\frac{\sqrt3\,\pi}{R}\cdot7^4\cdot\frac{\mathrm{Cmj}}{6\pi^2}\le\mathrm{Cpv}_2\cdot R^{-1},\qquad\mathrm{Cpv}_2:=221\,\mathrm{Cmj},
\end{aligned}
\tag{7.7}
\end{equation}
since $\sqrt3\cdot7^4/(6\pi)<221$ by (Ef0)\,(ix).

\paragraph{7.6. Step 3: $\eta\to0$.}
For $\eta\in(0,1]$ define
\begin{equation}
\begin{aligned}
\mathrm{epv}(\eta)&:=(2\pi)^{-3}\left[\int_{\mathrm{Cq}}\mathrm{d}^3\mathrm{kph}\,\left|\mathrm{Ipv}_\eta(\mathrm{kph})-\mathrm{Ipv}_0(\mathrm{kph})\right|+2\int_{\mathbb{R}^3\setminus\mathrm{Cq}}\mathrm{d}^3\mathrm{kph}\,\mathrm{Mpv}(\mathrm{kph})\right],\\
\mathrm{Cq}&=[-\pi/(2\eta),\pi/(2\eta)]^3.
\end{aligned}
\tag{7.8}
\end{equation}
It is formed with $\rho$, $\bar{m}$, $\beta$ and $\eta$; $R$ does not enter.
By (7.5) and 7.2, which give $|\mathrm{Ipv}_\eta|\le\mathrm{Mpv}$ on $\mathrm{Ucu}\setminus\mathrm{Cq}$ and $|\mathrm{Ipv}_0|\le\mathrm{Mpv}$,
\begin{equation}
\begin{aligned}
&\left|\int_{\mathrm{Ucu}}\frac{\mathrm{d}^3\mathrm{kph}}{(2\pi)^3}\,\mathrm{Ipv}_\eta-\int_{\mathbb{R}^3}\frac{\mathrm{d}^3\mathrm{kph}}{(2\pi)^3}\,\mathrm{Ipv}_0\right|\\
&\qquad\le(2\pi)^{-3}\left[\int_{\mathrm{Cq}}\mathrm{d}^3\mathrm{kph}\,|\mathrm{Ipv}_\eta-\mathrm{Ipv}_0|+\int_{\mathrm{Ucu}\setminus\mathrm{Cq}}\mathrm{d}^3\mathrm{kph}\,|\mathrm{Ipv}_\eta|+\int_{\mathbb{R}^3\setminus\mathrm{Cq}}\mathrm{d}^3\mathrm{kph}\,|\mathrm{Ipv}_0|\right]
\le\mathrm{epv}(\eta).
\end{aligned}
\tag{7.9}
\end{equation}
\emph{Boundedness.}
On $\mathrm{Cq}$ one has $|\mathrm{kph}|_\infty\le\pi/(2\eta)$, so that $|\mathrm{Ipv}_\eta-\mathrm{Ipv}_0|\le2\,\mathrm{Mpv}$ there by 7.2.
Hence, by (7.2),
\[
0\le\mathrm{epv}(\eta)\le(2\pi)^{-3}\left[2\int_{\mathrm{Cq}}\mathrm{d}^3\mathrm{kph}\,\mathrm{Mpv}+2\int_{\mathbb{R}^3\setminus\mathrm{Cq}}\mathrm{d}^3\mathrm{kph}\,\mathrm{Mpv}\right]
=2\int_{\mathbb{R}^3}\frac{\mathrm{d}^3\mathrm{kph}}{(2\pi)^3}\,\mathrm{Mpv}=\frac{\mathrm{Cmj}}{3\pi^2}.
\]
\emph{Limit.}
Let $\eta_j\to0$, and let $\mathrm{Cq}(\eta_j)$ denote the cube $\mathrm{Cq}$ formed with $\eta_j$.
The functions $\mathbf{1}[\mathrm{kph}\in\mathrm{Cq}(\eta_j)]\,|\mathrm{Ipv}_{\eta_j}(\mathrm{kph})-\mathrm{Ipv}_0(\mathrm{kph})|$ tend to $0$ pointwise by 7.1, and they are bounded by the integrable function $2\,\mathrm{Mpv}$ by 7.2.
The functions $\mathbf{1}[\mathrm{kph}\notin\mathrm{Cq}(\eta_j)]\,\mathrm{Mpv}(\mathrm{kph})$ tend to $0$ pointwise and are bounded by $\mathrm{Mpv}$.
By dominated convergence, $\mathrm{epv}(\eta_j)\to0$.

\paragraph{7.7. Conclusion.}
The bounds (7.6), (7.7), (7.9), the identity (7.3) and the triangle inequality give, for all $\eta\in(0,1]$ and all $R\ge1$ with $R/\eta$ an integer $\ge2$,
\[
|a_{\mathrm{ch}}(\eta,R)-\mathrm{pv}(f)|\le\mathrm{Cpv}_1\cdot\bar{m}^{-3}R^{-4}+\mathrm{Cpv}_2\cdot R^{-1}+\mathrm{epv}(\eta).
\]
Here $\mathrm{Cpv}_1=365\,\mathrm{Cpv}_0\cdot\mathrm{Cah}_2(\beta)^2$ is formed with $\rho$ and $\beta$, and $\mathrm{Cpv}_2=221\,\mathrm{Cmj}$ with $\rho$, $\bar{m}$ and $\beta$.
The three terms on the right tend to $0$ as $R\to\infty$, as $R\to\infty$ and as $\eta\to0$, respectively, each uniformly in the other variable.
Hence $a_{\mathrm{ch}}(\eta,R)\to\mathrm{pv}(f)$ as $\eta\to0$ and $R\to\infty$ in any manner.
No relation between $R$ and the support of $\beta$ is required; the periodisation enters only through the identity of 3.7.
\qed

\subsubsection*{8. Clause (g)}

\paragraph{8.1. Proof of (g).}
Let $0<\eta\le1$ and $R\ge1$ with $R/\eta$ an integer $\ge2$.
Clause (f), together with $R^{-4}\le1$ and $R^{-1}\le1$, gives
\[
|a_{\mathrm{ch}}(\eta,R)|\le|\mathrm{pv}(f)|+\mathrm{Cpv}_1\cdot\bar{m}^{-3}+\mathrm{Cpv}_2+\mathrm{epv}(\eta)\le|\mathrm{pv}(f)|+\mathrm{Cpv}_1\cdot\bar{m}^{-3}+\mathrm{Cpv}_2+\sup\mathrm{epv}.
\]
The right side is finite by (7.3) and 7.6, and it is formed with $\rho$, $\bar{m}$ and $\beta$.
Therefore $\mathrm{Cpv}_\beta$ is finite and satisfies the bound of (g).
With the values of Section 7,
\[
\mathrm{Cpv}_\beta\le\frac{\mathrm{Cmj}}{6\pi^2}+365\,\mathrm{Cpv}_0\cdot\mathrm{Cah}_2(\beta)^2\cdot\bar{m}^{-3}+221\,\mathrm{Cmj}+\frac{\mathrm{Cmj}}{3\pi^2}.
\]
\qed

\paragraph{8.2. A bound on $\mathrm{Cpv}_\beta$ from (b), (c) and (e).}
Let $0<\eta\le1$ and $R\ge1$ with $R/\eta$ an integer $\ge2$, and let $\mathrm{kph}\in\mathrm{Kx}_{\eta,R}$.
By the first two bounds of (c),
\[
|\mathrm{Vpl}^{\eta,R}_{\mu\nu}(\mathrm{kph})|\le\mathrm{Cpv}_0\cdot\bar{m}^{-3}R^{-4}+\mathrm{Cpv}_0\cdot|\mathrm{kph}|_2^2/\bar{m}.
\]
By (7.5), the bound (6.4) is admissible at $\mathrm{kph}$; with $i=2$ and with $i=3$ it gives
\[
|\mathrm{ahl}_{\eta,\mu}(\mathrm{kph})|\,|\mathrm{ahl}_{\eta,\nu}(\mathrm{kph})|\le\mathrm{Cah}_2(\beta)^2\cdot(1+|\mathrm{kph}|_2)^{-4},\qquad
|\mathrm{ahl}_{\eta,\mu}(\mathrm{kph})|\,|\mathrm{ahl}_{\eta,\nu}(\mathrm{kph})|\le\mathrm{Cah}_3(\beta)^2\cdot(1+|\mathrm{kph}|_2)^{-6}.
\]
The first of these bounds is combined with the first term, the second with the second term, and $|\mathrm{kph}|_2^2\,(1+|\mathrm{kph}|_2)^{-6}\le(1+|\mathrm{kph}|_2)^{-4}$ is used.
There are nine pairs $(\mu,\nu)$.
By (b), $R^{-4}\le1$, (7.4) and (Ef0)\,(viii),
\begin{equation}
\begin{aligned}
|a_{\mathrm{ch}}(\eta,R)|
&\le9\,\mathrm{Cpv}_0\cdot\left[\mathrm{Cah}_2(\beta)^2\cdot\bar{m}^{-3}+\mathrm{Cah}_3(\beta)^2\cdot\bar{m}^{-1}\right]R^{-3}\sum_{\mathrm{kph}\in\mathrm{Kx}_{\eta,R}}(1+|\mathrm{kph}|_2)^{-4}\\
&\le365\,\mathrm{Cpv}_0\cdot\left[\mathrm{Cah}_2(\beta)^2\cdot\bar{m}^{-3}+\mathrm{Cah}_3(\beta)^2\cdot\bar{m}^{-1}\right].
\end{aligned}
\tag{8.1}
\end{equation}
Hence $\mathrm{Cpv}_\beta$ is at most the right side of (8.1).
\qed

\subsection*{Uniformity}

No level, no side $\ell\in\mathrm{Cl}$, no charge $e$, no cutoff index $N$ and no scale factor $L$ occurs in the statement of Proposition~36, so that its constants are independent of all of these quantities.
The number $q>0$ in the first sentence of (h) is arbitrary (subsection~0.1); the operator depends on it only through $B=qA$ (subsection~0.2), and $N_0$ does not depend on it.
The variables are $\eta\in(0,1]$, which takes every real value in this interval, and $R$, which takes every value with $R/\eta$ an integer $\ge2$; in (f) and (g), moreover, $R\ge1$.

Clauses (a), (b) and (h) are identities for each pair $(\eta,R)$.

In (c), the constant $\mathrm{Cpv}_0=10^8\,\mathrm{Ksy}^2$ depends on $\rho$ only.
The bounds hold for all $\mathrm{kph}\in\mathbb{R}^3$, all $\eta$, all $R$ and all $\bar{m}\in(0,1]$, because Proposition~37.1\,(g) holds for all $\mathbf{m}\in(0,1]$ with constants that depend on $\rho$ only, and because the powers of $\mathbf{m}$ are collected exactly: $\eta^{-1}n^{-4}\mathbf{m}^{-3}=\bar{m}^{-3}R^{-4}$ in 4.3, and $\eta/\mathbf{m}=1/\bar{m}$ in 4.6.

In (e), $\mathrm{Cah}_i(\beta)$ depends only on $D_{\mathrm{be}}$ and $|\beta|_{C^{i+1}}$; it does not depend on $\eta$, $R$, $\rho$ or $\bar{m}$, and it is the same for all translates $\tau_v\beta$, $v\in\mathbb{R}^3$.

In (f) and (g), the numbers $\mathrm{Cpv}_1$, $\mathrm{Cpv}_2$, the bound $\mathrm{Cmj}/(3\pi^2)$ on $\mathrm{epv}$ and the right side of (8.1) depend neither on $\eta$ nor on $R$; the function $\mathrm{epv}$ depends on $\eta$ and not on $R$.
These numbers depend on $\beta$ only through $D_{\mathrm{be}}$, $|\beta|_{C^3}$ and $|\beta|_{C^4}$; the function $\mathrm{epv}$ itself depends on $\beta$.

\subsection*{Conditions on the constants}

Proposition~36 holds without any of the inequalities between constants imposed in Section~I.8.
Its proof uses three data, namely $\rho$, the standard weights and $\bar{m}$, through their ranges and defining relations, and the numbers $\mathrm{Csy}_{\mathrm{ak}}(\rho)$, $|\mathrm{ak}|\le4$, of the line (C0r); for each of these, the first table gives the place where it enters the proof and the line of Section~I.8.2, or of Proposition~14.2, in which it is fixed.
The second table defines the constants of the statement.

\noindent
\begin{tabular}{|p{0.36\textwidth}|p{0.26\textwidth}|p{0.26\textwidth}|}
\hline
datum or constant & where it enters the proof & line of Section~I.8.2 / Proposition~14.2\\
\hline
$\rho>0$ & (1.2): $\operatorname{Re}(D_A+\mathbf{m})\ge\mathbf{m}$ & (C0): $\rho\in(0,1)$\\
\hline
$\rho\in(0,1)$ & only through Proposition~37.1\,(a), (d), (e), (g), (h) & (C0)\\
\hline
the standard weights: $\sum_i w_i\hat{n}_i=0$ ((I.1.1), (I.s3)); $\sum_i M_i=(\rho/2)\,\mathrm{sw}$ ((I.m1)) & (1.1), 1.1\,(iv) & (C0)\\
\hline
$\bar{m}\in(0,1]$, so that $\mathbf{m}=\bar{m}\eta$ lies in $(0,1]$, the range of Proposition~37.1\,(g) & (0.7) & (C10)\\
\hline
$\mathrm{Csy}_{\mathrm{ak}}(\rho)$, $|\mathrm{ak}|\le4$ & (0.7) & (C0r), row $\mathrm{Csy}_{\mathrm{ak}}$, $\mathrm{Cpv}_0$ of Proposition~14.2\,(a)\\
\hline
\end{tabular}

\medskip

\noindent
\begin{tabular}{|p{0.19\textwidth}|p{0.34\textwidth}|p{0.11\textwidth}|p{0.2\textwidth}|}
\hline
constant & value & depends on & line\\
\hline
$\mathrm{Ksy}$ (auxiliary) & $\max\{1,\max\{\mathrm{Csy}_{\mathrm{ak}}(\rho):|\mathrm{ak}|\le4\}\}$ & $\rho$ & (C0r), after $\mathrm{Csy}_{\mathrm{ak}}$\\
\hline
$\mathrm{Cpv}_0$ & $10^8\,\mathrm{Ksy}^2$ & $\rho$ & (C0r), row $\mathrm{Csy}_{\mathrm{ak}}$, $\mathrm{Cpv}_0$\\
\hline
$\mathrm{Cah}_i(\beta)$ & $2\cdot12^i\,(D_{\mathrm{be}}+i+3)^3\,|\beta|_{C^{i+1}}$ & $\beta$ & (C10), row $\mathrm{Cah}_i$, $\mathrm{Cpv}_1$, $\mathrm{Cpv}_2$, $\mathrm{epv}$, $\mathrm{Cpv}_\beta$, after $D_{\mathrm{be}}$\\
\hline
$\mathrm{Cahd}_i(\beta)$ (auxiliary) & $12^i\,(D_{\mathrm{be}}+i+3)^3\,(D_{\mathrm{be}}+2+2i)\,|\beta|_{C^{i+1}}$ & $\beta$ & (C10)\\
\hline
$\mathrm{Cmj}$ (auxiliary) & $9\sqrt3\,\mathrm{Cpv}_0\cdot\bar{m}^{-1}\cdot\mathrm{Cah}_3\cdot(2\,\mathrm{Cahd}_3+\mathrm{Cah}_3)$ & $\rho$, $\bar{m}$, $\beta$ & (C10)\\
\hline
$\mathrm{Cpv}_1$ & $365\,\mathrm{Cpv}_0\cdot\mathrm{Cah}_2^2$ & $\rho$, $\beta$ & (C10), same row\\
\hline
$\mathrm{Cpv}_2$ & $221\,\mathrm{Cmj}$ & $\rho$, $\bar{m}$, $\beta$ & (C10), same row\\
\hline
$\mathrm{epv}$ & the function (7.8); $0\le\mathrm{epv}\le\mathrm{Cmj}/(3\pi^2)$ & $\rho$, $\bar{m}$, $\beta$ & (C10), same row\\
\hline
$\mathrm{Cpv}_\beta$ & the supremum in (g); at most the right side of (8.1) & $\rho$, $\bar{m}$, $\beta$ & (C10), same row\\
\hline
\end{tabular}

\medskip

Accordingly, Proposition~36 requires no condition on the constants beyond those of Section~I.8.
In the order of Section~I.8.2, $\mathrm{Cpv}_0$ is formed after the numbers $\mathrm{Csy}_{\mathrm{ak}}$ within the line (C0r), and the constants of the line (C10) are formed after $\bar{m}$, $\beta$, $D_{\mathrm{be}}$ and $\mathrm{Cpv}_0$; thus each constant of Proposition~36 is formed from quantities that precede it.

\section*{Chapter 39. Chessboard estimate and lengthening for coordinate cells of half-integer width: proofs of Proposition~39}

This chapter contains the proofs of Proposition~39 of Part~II, with Definition~39.0 and the list of the places at which the width of the cells enters the proofs, and of Remark~39.R.

\subsection*{0. Conventions of this chapter}

The conventions of Part~I are in force.
The notation of \cite{QEDII} is related to the notation of this paper in Section~1.1 for Sections~1--7, and in Section~A.0 for Appendix~A and Section~8.

Several letters that carry a fixed meaning in Part~I are given a local meaning in this chapter; these local meanings are listed below.
Where a letter keeps, in addition, a meaning that it has elsewhere in the paper, this other meaning is named after the word ``besides'', and the context distinguishes the two.

\emph{Letters for objects of \cite{QEDII}.}
The following letters denote objects of \cite[Section~5, Proposition~3.20 and Definition~2.27]{QEDII}, where these objects are written with the letters given in parentheses; each of them is introduced where it first occurs (Lemma~2.4\,(a), A.2, A.6, A.8, A.9).
Here $\mathbf{K}$ is a product of three intervals of $\mathbb{Z}$, regarded as a complex of cells (in the role of the complex $\mathcal{K}$ of \cite[Lemma~5.13]{QEDII});
$\omega$, $\omega'$, $\omega_k$ are configurations $(A,\mathfrak{z})$ of a closed slab ($\varpi$ in \cite[Definition~5.2]{QEDII});
$\xi=(\mathrm{mi},\boldsymbol{\zeta},I)$ is a mirror variable, with $\boldsymbol{\zeta}=(\varphi,\mathrm{kz})$ ($\xi=(\mathfrak{m},\vartheta,k,I)$ in \cite[Definition~5.27]{QEDII}, the pair $(\vartheta,k)$ standing in the place of $\boldsymbol{\zeta}$);
$\mathrm{Ee}_i=\Theta_{\mathrm{po}}(C_i)\cdot C_i$ ($\Theta_\Pi(\zeta_i)\,\zeta_i$ in \cite[Proposition~3.20]{QEDII});
$C$ and $C_i$ are modes ($\zeta_i$ in \cite[Definition~2.27]{QEDII}), and $C_I$, $C^{(\mathrm{po})}_I$ denote the ordered monomial of the modes of $I$ ($\zeta_I$ and $\zeta^{(c)}_I$ in \cite[Definitions~2.27 and~5.20]{QEDII}); this $C_I$ is not the constant $C_I$ of the line (C4r) of Section~I.8.2;
$e_I$, $e^{+}_{k'}$, $e^{-}_{k'}$ are basis vectors ($e_I$ and $e^{\pm}_{n,k}$ with $n=e_\mu$ in \cite[Definition~2.27]{QEDII}).
In \cite{QEDII} the letter $\omega_i$ denotes the crossing coefficient of a mode $i$, which is written $c_i$ in A.6, and $\zeta_i$ denotes the mode itself; the letters $\omega$ and $\boldsymbol{\zeta}$ of this chapter have the meanings listed above.

\emph{The letters $k$, $s$, $t$.}
In this chapter these letters denote neither a level, nor an exponent of the line (C7) of Section~I.8.2, nor a parameter family, nor a contour variable; their meanings are the following.
The letter $k$ denotes: in 1.7, the integer of \cite[Lemma~10.6]{QEDI}; in Lemma~2.1 and in Sections~3--7, the integer $k\ge1$ for which $2k$ is the number of the operators, respectively the number $n_\mu/\mathrm{hw}$ of the layers of $C(n)$ in one direction, and $2k$ is then also the Schatten exponent; in Theorem~SO, in Appendix~A and in Section~8, the index of a plane $\mathrm{po}_k$ of a chain and of the objects attached to it ($G_k$, $\Phi_k$, $\mathrm{mt}_k$, $\mathrm{Ms}_k$, $\omega_k$); in (Be1), the index of a factor.
The letter $s$ denotes: the number of the factors of a product, or of the letters of a word (1.3, 1.4, Lemmas~2.3 and~2.4); the number of the slabs of a chain $\mathrm{po}_0<\dots<\mathrm{po}_s$ (1.6, Lemma~2.1, Appendix~A); an integer multiple of $\tau_\mu$ (Section~3, 7.1, A.11); a real variable (1.7\,(c), the proofs of Proposition~39\,(e) and of Proposition~7.2); and, in Remark~39.R and in Section~8, the function $s=1$ of the slice data.
The letter $t$ denotes an integer or a real variable, declared where it occurs; in A.9 and in Section~8 it is the position of a site plane.
The dimension of a cell is denoted by $\mathrm{dg}$ (A.1).

\emph{Further letters.}
Besides the translation part of $g_{\mathrm{w}}=(a,r)$ and a link $b$, the letters $a$ and $b$ denote integers (1.1, Lemma~2.4\,(a), A.0, Corollary~A.6) and spinor indices ((S3), (S4), (Ap9)); the letter $a$ is also the index of a member of a frame (7.2), and $a_\beta$, $b_\beta$ are anticommuting elements ((Be2)).
In Section~8, as in Remark~39.R, the letters $a$, $a'$, $a_t$, $a_y$, $a_b$ denote real functions on $\mathrm{Cu}^1(1/2)$ and their values; here $a'$ is not the constant of the line (C4r).
The letter $B$ denotes a positive operator of trace class (1.7\,(c), the proof of Proposition~39\,(e), 7.2\,(h4)), an operator or a matrix (the proofs of Lemma~2.1 and of Lemma~A.3\,(F5)), and a real function on $\mathrm{Mi}^1(0)$ (Section~8).
The letters $C'$, $C_n$, $C_{\mathrm{po},\mathrm{po}'}$, and $C$ in (SO3), denote constants (Lemma~2.4\,(d), A.10).
The coordinates of a photon site $c$ are $c_1$, $c_2$, $c_3$, $c_\mu$; in particular, $c_3$ denotes a coordinate and not the photon site of $x^{(3)}$.
Besides the coboundary, the letter $d$ denotes an element of a group $\mathrm{Dw}_\mu$ ($d$, $d_\mu$, $d_a$: Sections~4, 5 and~7.2), and a degree or the number of the elements of $I$ (A.6, A.7); the numbers $d_j\in\{1,2,3\}$ are the directions of a cell (A.1).
The letter $f$ denotes a function on $\mathbb{Z}$ (Lemma~2.4\,(a)), a Schwartz function ((K1)), and an element of $\mathrm{Hs}(\mathrm{po})$ (A.9, A.10).
The letters $G_j$, $G_k$ denote the factors of a product (1.4), and the element of $\mathcal{A}_{\mu,n}[\,\cdot\,,\cdot\,]$ or of $\mathrm{Ig}[\,\cdot\,,\cdot\,]$ that belongs to the layer $j$, respectively to the slab $k$.
The sets $I_1$, $I_2$, $I_3$, $I_\mu$, $I_\nu$, $I_{\mathfrak{l}}$ are intervals of $\mathbb{Z}$ (Lemma~2.4\,(a), Lemma~A.2), not the intervals of a cell of Section~I.6.4; $I_b$ is an interval of integration (Lemma~A.1, A.10); $I_0$ and $I_1$ are subsets of $\mathrm{Ri}(\mathrm{po}_0)$ and of $\mathrm{Ri}(\mathrm{po}_1)$ (A.7, A.10).
The index of a cell of $\mathcal{Q}'$ is denoted by $\varkappa$ (7.2\,(h1)); it is the index written $\kappa$ in \cite[Definition~2.18]{QEDII}.
Besides a mirror $m$, the letters $m$, $m'$ denote integer one-cochains; in 3.1 the letter $m$ is an integer.
The letters $P$, $P'$, $P_i$, $P_{\mathfrak{c}}$ denote elements of $\mathrm{Fn}$ (1.3) and products (1.4), and $P^{+}$, $P^{-}$ are the spectral projections of $\gamma_\mu$ (A.6).
Besides a plaquette $p$, the letters $p$, $p'$, $p_i$ denote Schatten exponents (1.7, Lemma~2.1, the proof of Proposition~39\,(e)).
Besides the linear part of $g_{\mathrm{w}}$, the letter $r$ denotes the number of the factors (Lemma~2.4\,(b), 3.1, (5.b)).
The letters $S$, $S_i$ denote the content of a layer (the proof of Proposition~39\,(e), Lemma~5.1); $S_p$, $S_2$, $S_{2k}$ are Schatten classes (1.7); $S_n(A)$, $S_{(\mathrm{po},\mathrm{po}')}(A)$, and in A.10 also $S$, are actions (1.5, A.3).
The number $t_0$ lies in $[0,r_q)$ (Lemma~A.1); it is not the constant of the line (C0r).
The letter $u$ denotes a cochain (as in Proposition~39\,(f)), the vector $n_\mu e_\mu$ (Lemma~5.1), and a variable of integration ((K1)).
The letters $v$, $v'$ denote vectors of $\mathbb{Z}^3$ or of $\mathbb{R}^3$.
The letter $X$ denotes an operator ((K3)), and $X(\,\cdot\,)$ denotes sets of fermion sites (1.4, A.3).
Besides a fermion site $x=(c,g)$, the letter $x$ denotes a real variable in (K1).
In Lemma~2.2 the letter $z$ denotes a complex number, not a unit site or a site of the $L$-lattice as in Part~I.
The letters $Y$, $Y_i$, $Y_j$, $Y_k$ denote operators (1.7, Lemma~2.1, Sections~3 and~6, (K3)), a $4\times4$ matrix ((S3)), and finite sets of sites or of classes of sites ((S1), A.4).
Besides a point of $\mathbb{R}^3$, the letter $y$ denotes a real number (Lemma~2.4\,(d), (K1), (K2)) and a site of the slice (Section~8, as in Remark~39.R).

\emph{Values of the curvature.}
In this chapter $F_p$, $\tilde F$, $F(\omega)$ and $F_p(\omega)$ are values of the curvature in lattice units: of a two-cochain in the place of $F^{\mathrm{u}}$ (1.4), of the periodic extension of $F^{\mathrm{u}}$ (1.5), and of $dA+r_q\cdot\mathfrak{z}$ for a configuration of a slab (A.2).
They are not values of the physical curvature $F$ of Section~I.7.1.

\emph{Labels.}
The labels (M1)--(M5), (S1)--(S6), (SO0)--(SO3), (Tq1)--(Tq6), (B1)--(B3), (C1), (C2), (C4) and (F1)--(F5) are those of this chapter; (B1)--(B3) and (C1), (C2), (C4) are the clauses of Lemmas~2.2 and~2.3, and (F1)--(F5) are the clauses of Lemma~A.3; the latter are not the properties (F1)--(F3) of \cite[Section~10.2.1]{QEDI} that occur in Remark~39.R.
The identities (Tq1)--(Tq6) of A.4 and A.6 correspond, in the setting of Section~A.0, to those of \cite[Lemma~5.8]{QEDII} for (Tq1)--(Tq3), of \cite[Lemma~5.19]{QEDII} for (Tq4), and of \cite[Definition~5.20]{QEDII} for (Tq5) and (Tq6).
The labels (h1)--(h5) are those of Section~I.6.4.

\subsection*{1. Letters, definitions and recalled results}

Lattice units are used throughout (Sections~I.2.5 and~I.7.1).
Fix $\rho\in(0,1)$, the standard weights $w=(6,6\sqrt{2},6\sqrt{2},6)$, a number $q>0$ with $r_q=2\pi/q$, and a real lattice mass $\boldsymbol{\mu}$, and put $M_{\mathrm{ch}}=\boldsymbol{\mu}+\frac{\rho}{2}\cdot\mathrm{sw}$.
Directions are denoted by $\mu,\nu,\mathfrak{l}\in\{1,2,3\}$, and $|v|$ is the sup norm of $v\in\mathbb{R}^3$.

\subsubsection*{1.1. Letters}

For $\mathrm{po}\in\tfrac12\mathbb{Z}$ put
\begin{gather*}
\theta_{\mu,\mathrm{po}}:=\theta_{y_\mu,\mathrm{po}}\colon y\mapsto y-2\cdot\left(y_\mu-\mathrm{po}\right)\cdot e_\mu,\qquad
\Theta_{\mu,\mathrm{po}}:=\Theta_{(y_\mu,\mathrm{po})},\\
r_\mu:=1-2\,e_\mu e_\mu^{T},\qquad
\mathrm{Sl}_\mu(\mathrm{po},\mathrm{po}'):=\left\{y\in\mathbb{R}^3:\ \mathrm{po}<y_\mu<\mathrm{po}'\right\}
\end{gather*}
(Section~I.4.8; the unit normal is $\hat{n}=e_\mu$, so that $\gamma\cdot\hat{n}=\gamma_\mu$).
For $v\in\mathbb{Z}^3$, $\tau_v$ denotes the translation $y\mapsto y+v$.
For $n=(n_1,n_2,n_3)$ with integers $n_\mu\ge4$ put
\begin{gather*}
\Gamma(n):=\sum_{\mu} n_{\mu}\mathbb{Z}\,e_{\mu},\qquad
\Gamma_\mu(n):=\sum_{\nu\ne\mu}n_\nu\mathbb{Z}\,e_\nu,\qquad
|n|:=n_1n_2n_3,\\
\chi_n\left(\sum_\mu a_\mu n_\mu e_\mu\right):=(-1)^{a_1+a_2+a_3}\qquad(a_\mu\in\mathbb{Z}).
\end{gather*}
The torus $(\Gamma(n),\chi_n)$ of \cite[Definition~2.11]{QEDII} is denoted by $C(n)$.
Put
\[
\mathrm{Xs}_\mu(n):=\left(\Gamma_\mu(n),\ \chi_n|_{\Gamma_\mu(n)}\right).
\]
This is a cross-section for the family $y_\mu$ in the sense of \cite[Definition~2.17]{QEDII}: $\Gamma_\mu(n)$ is a subgroup of rank two of $\mathbb{Z}^3\cap\{y_\mu=0\}$, and its non-zero elements have sup norm at least $\min_{\nu\ne\mu}n_\nu\ge4$.
Moreover $\Gamma(n)\cap\{y_\mu=0\}=\Gamma_\mu(n)$ and $\Gamma(n)=\Gamma_\mu(n)+\mathbb{Z}n_\mu e_\mu$, and for $\ell=y_\mu$ one has $n_\ell=e_\mu$ and $|n_\ell|^{2}=1$.
Hence the cross-section $\Xi(C(n),y_\mu)$ of the torus $C(n)$ for the family $y_\mu$, as defined in \cite[Definition~2.17]{QEDII}, equals $\mathrm{Xs}_\mu(n)$; moreover the vector $n_\mu e_\mu$ is a closing vector, the height period is $n_\mu$, $\chi_n(n_\mu e_\mu)=-1$, and $C(n)$ is adapted to $y_\mu$ in the sense of the same definition.
A cross-section, which \cite{QEDII} denotes by $\Xi$, is written $\mathrm{Xs}$ in this chapter; in particular, $\mathcal{A}_{\mathrm{Xs}_\mu(n)}$ in Section~1.4 is the class $\mathcal{A}_{\Xi}$ of \cite[Definition~2.17]{QEDII} for $\Xi=\mathrm{Xs}_\mu(n)$.
The group $\Gamma_\mu(n)$, the cross-section $\mathrm{Xs}_\mu(n)$ and every object formed with them depend only on the two sides $n_\nu$, $\nu\ne\mu$.

\emph{Notation of \cite{QEDII}.}
The following objects are those of \cite{QEDII}; where \cite{QEDII} writes them with other letters, these letters are given in parentheses.
They are the Grassmann algebras $\mathrm{Gr}(Y)$ and $\mathrm{Gr}_{\mathbb{T}}$ and the Berezin integral $\int_{\mathbb{T}}$ of a torus $\mathbb{T}$;
boxes $\mathrm{Bx}$ and the sets $\mathrm{Ln}(\mathrm{Bx})$, $\mathrm{Pq}(\mathrm{Bx})$ of their links and plaquettes ($B$, $B^{(1)}$, $B^{(2)}$ in \cite[Definition~2.12]{QEDII});
the class $\mathrm{Ad}(\mathrm{Bx})$ ($\mathcal{G}(B)$ in \cite[Definition~2.12]{QEDII}) and the classes $\mathcal{A}$, $\mathcal{A}_{\Xi}$ of \cite[Definitions~2.12 and~2.17]{QEDII};
local lifts $\hat{A}$;
the slab operators $\mathrm{Cs}[G]$ ($\mathfrak{S}[G]$ in \cite[Definition~5.28 and Corollary~5.30]{QEDII});
the fold $\mathrm{fold}_\ell$ of \cite[Definition~6.7]{QEDII}, written $\mathrm{fold}_\mu$ in Section~3.1;
and the set $\mathrm{Nd}(\mathcal{Q})$ of shapes ($\mathcal{M}(\mathcal{Q})$ in \cite[Definition~6.22]{QEDII}), written $\mathrm{Nd}$ in Section~7.3.
The number $\mathrm{kk}$ is not an object of \cite{QEDII}: it is one of the data of Theorem~SO, and for the spaces and operators of \cite{QEDII} it equals $1$ (Section~1.6).
Statements of \cite{QEDII} are read with the following identifications.
A fermion site of \cite{QEDII} is the point $\mathrm{geo}(x)$ of a site $x=(c,g)\in\mathbb{L}^{\mathrm{f}}_0=\mathbb{Z}^3\times O_{\mathrm{cub}}$ (Proposition~46\,(a)).
The group $W$ of \cite{QEDII} is $W_{\mathrm{af}}$, whose elements are written $g_{\mathrm{w}}=(a,r)$.
The operator $D_A$ of \cite{QEDII} is $D_A+\boldsymbol{\mu}$, and its mass parameter $M$ is $M_{\mathrm{ch}}$.
The curvature $F$ of \cite{QEDII} is $F^{\mathrm{u}}$; the position $c$ of a mirror is $\mathrm{po}$; the width $h$ of a cell structure is $\mathrm{hw}$; the index $\kappa$ of a cell is $\mathfrak{c}$; the open cell $Q_\kappa$ is $Q^{\circ}_{\mathfrak{c}}$; and the mirror copy $G^{[\kappa]}$ is $O^{[\mathfrak{c}]}$.
For the torus $\mathbb{T}=C(n)$, the fold $\operatorname{Fold}_{\mathbb{T}}(O)$ is $O^{\sharp}_n$, the functional $Z_{\mathbb{T}}$ is $Z_n$, and the expectation $\mathbb{E}_{\mathbb{T}}$ is $E^{\mathrm{af}}_n$.

\emph{Further notation of this chapter.}
The symbols $X(\mathrm{Bx})$ (1.4), $\mathrm{Hs}(\mathrm{po})$ and $V_v$ (1.6) denote the objects that correspond to those written $\mathbb{X}(B)$, $\mathcal{H}(c)$ and $V_v$ in \cite[Definition~2.12 and Corollary~5.30]{QEDII}.
The centre and the half side of a box are $\mathrm{bc}$ and $\mathrm{rd}$.
The symbols $\mathrm{Fn}$, $\mathrm{Tg}_v$, $\mathrm{qu}_n$ are introduced in 1.3 and 1.5; $\mathrm{Hb}$, $\mathrm{em}_j$, $\hat{Y}_j$ in Lemma~2.1; $S_a^t$, $\boldsymbol{\Theta}_{\varepsilon}$, $C_n$ in Lemma~2.4; $\mathrm{po}_j=\mathrm{po}^{\mu}_j$, $\tau_\mu$, $\mathrm{Dw}_\mu$, $\mathrm{dw}^{\mu}_{i\leftarrow j}$ in 3.1 and at the beginning of Section~4; $\mathrm{sh}$, $\mathrm{GG}$, $\mathrm{Rs}$, $\mathrm{Prd}$, $\mathrm{Kb}(n)$, $\mathrm{Rm}_\mu^j$, $\mathrm{CC}(O)$ in Section~5 (for $\mathrm{CC}$ cf.\ \cite[Definition~6.13]{QEDII}); and $\mathrm{kv}$ in the display (6.1) of the proof of Proposition~39\,(e).
The notation of Appendix~A, which Section~8 also uses, is introduced in Section~A.0.

\subsubsection*{1.2. Identities for the reflections}

The identities (M1)--(M4) are those of \cite[Definition~2.2\,(a)--(d)]{QEDII} for the family $\ell=y_\mu$.
For $\mathrm{po},\mathrm{po}'\in\tfrac12\mathbb{Z}$ and $v\in\mathbb{Z}^3$ one has
\begin{align*}
&\theta_{\mu,\mathrm{po}}^2=\mathrm{id};\tag{M1}\\
&\left(\theta_{\mu,\mathrm{po}}y\right)_\mu=2\,\mathrm{po}-y_\mu,\qquad\left(\theta_{\mu,\mathrm{po}}y\right)_\nu=y_\nu\quad\text{for }\nu\ne\mu;\tag{M2}\\
&\theta_{\mu,\mathrm{po}'}\,\theta_{\mu,\mathrm{po}}=\tau_{2(\mathrm{po}'-\mathrm{po})e_\mu},\quad\text{and }2\,(\mathrm{po}'-\mathrm{po})\text{ is an integer};\tag{M3}\\
&\tau_v\,\theta_{\mu,\mathrm{po}}\,\tau_{-v}=\theta_{\mu,\mathrm{po}+v_\mu};\tag{M4}\\
&\theta_{\mu,\mathrm{po}}\,\theta_{\nu,\mathrm{po}'}=\theta_{\nu,\mathrm{po}'}\,\theta_{\mu,\mathrm{po}}\quad\text{for }\mu\ne\nu.\tag{M5}
\end{align*}

\begin{proof}
(M2) is the definition.
For (M1) and (M3): by (M2), the coordinate $\mu$ of $\theta_{\mu,\mathrm{po}'}\,\theta_{\mu,\mathrm{po}}y$ is $2\,\mathrm{po}'-(2\,\mathrm{po}-y_\mu)$, and the other coordinates are those of $y$.
For (M4): at the point $y$, the coordinate $\mu$ of the left side is $2\,\mathrm{po}-(y_\mu-v_\mu)+v_\mu$, and the coordinate $\nu\ne\mu$ is $y_\nu$.
For (M5): the two maps change different coordinates.
\end{proof}

Since $2\,\mathrm{po}$ is an integer, $\theta_{\mu,\mathrm{po}}$ maps $\mathbb{Z}^3$ onto itself; it is the element $\left(2\,\mathrm{po}\cdot e_\mu,\,r_\mu\right)$ of $W_{\mathrm{af}}$.

\subsubsection*{1.3. Grassmann algebra, $\Theta$, the action of $W_{\mathrm{af}}$ on even elements}

The definitions of this subsection are those of \cite[Definitions~2.8 and~2.9]{QEDII} and of Section~I.4.8.
Let $\mathrm{Gr}:=\mathrm{Gr}(\mathbb{L}^{\mathrm{f}}_0)$ be the complex Grassmann algebra with the generators $\psi_\beta(x)$, $\bar{\psi}_\beta(x)$, $x\in\mathbb{L}^{\mathrm{f}}_0$, $\beta=1,\dots,4$; here $\psi(x)$ is a column and $\bar{\psi}(x)$ a row.
For a set $Y$ of sites, $\mathrm{Gr}(Y)$ is the subalgebra generated by the generators at the sites of $Y$; the upper index $\mathrm{ev}$ denotes the even part.
An element of even degree commutes with all elements.
\begin{itemize}
\item[(S1)] For a finite set $Y$ of sites, $\int_Y$ is the Berezin integral normalised by
\[
\int\mathrm{d}\bar{\psi}_\beta(x)\,\mathrm{d}\psi_\beta(x)\,\psi_\beta(x)\,\bar{\psi}_\beta(x)=1;
\]
consequently $\int_Y\exp\left(-\bar{\psi}\,\mathrm{Dm}\,\psi\right)=\det\mathrm{Dm}$ for every complex matrix $\mathrm{Dm}$ whose rows and columns are indexed by $Y\times\{1,\dots,4\}$.
The pairing in the exponent is the counting pairing $\sum_x\bar{\psi}(x)\cdot(\,\cdot\,)(x)$, in which every site counts with the weight $1$ and not with the weight $1/48$ of (I.1.3).
For this pairing and the $48\,|n|$ classes of sites of $C(n)$ (the integral $\int_n$ of Section~1.5), the identity above is the first identity of Proposition~46\,(e): the Berezin integral equals $\det\mathrm{Dm}$, without the factor $48^{-192|n|}$ that the pairing $\langle\,\cdot\,,\cdot\,\rangle^{\mathrm{ch}}$ produces in the second identity there.
\item[(S2)] Let $m=(\ell,\mathrm{po})$ be a mirror, with the reflection $\theta$ and the unit normal $\hat{n}$.
Then $\Theta_m$ is the antilinear map of $\mathrm{Gr}$ with
\[
\Theta_m(GG')=\Theta_m(G')\,\Theta_m(G),\qquad
\Theta_m\psi(x)=-\bar{\psi}(\theta x)\,(\gamma\cdot\hat{n}),\qquad
\Theta_m\bar{\psi}(x)=-(\gamma\cdot\hat{n})\,\psi(\theta x).
\]
One has $\Theta_m^2=1$, and $\Theta_m$ preserves the degree.
\item[(S3)] For every complex $4\times4$ matrix $Y$ and all sites $x$, $x'$,
\[
\Theta_m\left(\bar{\psi}(x)\,Y\psi(x')\right)=\bar{\psi}(\theta x')\,(\gamma\cdot\hat{n})\,Y^{*}\,(\gamma\cdot\hat{n})\,\psi(\theta x)
\]
(\cite[Lemma~3.10\,(c)]{QEDII}).
Indeed, since $\Theta_m$ is antilinear and reverses products,
\[
\Theta_m\left(\bar{\psi}_a(x)\,Y_{ab}\,\psi_b(x')\right)
=\overline{Y_{ab}}\cdot\Theta_m\left(\psi_b(x')\right)\,\Theta_m\left(\bar{\psi}_a(x)\right)
=\overline{Y_{ab}}\cdot\left(\bar{\psi}(\theta x')\,(\gamma\cdot\hat{n})\right)_b\left((\gamma\cdot\hat{n})\,\psi(\theta x)\right)_a,
\]
the two minus signs cancelling each other; the claim follows by summation over $a$ and $b$, because $(Y^{*})_{ba}=\overline{Y_{ab}}$.
\item[(S4)] For two mirrors $m$, $m'$ with the unit normals $\hat{n}$, $\hat{n}'$,
\[
\Theta_m\Theta_{m'}\psi(x)=(\gamma\cdot\hat{n}')\,(\gamma\cdot\hat{n})\,\psi(\theta_m\theta_{m'}x),\qquad
\Theta_m\Theta_{m'}\bar{\psi}(x)=\bar{\psi}(\theta_m\theta_{m'}x)\,(\gamma\cdot\hat{n})\,(\gamma\cdot\hat{n}')
\]
(\cite[Lemma~3.10\,(d)]{QEDII}).
Indeed, as $\gamma\cdot\hat{n}'$ is Hermitian,
\begin{align*}
\Theta_m\Theta_{m'}\psi_a(x)
&=\Theta_m\left(-\sum_b\bar{\psi}_b(\theta_{m'}x)\,(\gamma\cdot\hat{n}')_{ba}\right)
=-\sum_b\overline{(\gamma\cdot\hat{n}')_{ba}}\cdot\Theta_m\bar{\psi}_b(\theta_{m'}x)\\
&=\sum_b(\gamma\cdot\hat{n}')_{ab}\cdot\left((\gamma\cdot\hat{n})\,\psi(\theta_m\theta_{m'}x)\right)_b,\\
\Theta_m\Theta_{m'}\bar{\psi}_a(x)
&=-\sum_b\overline{(\gamma\cdot\hat{n}')_{ab}}\cdot\Theta_m\psi_b(\theta_{m'}x)
=\sum_b\left(\bar{\psi}(\theta_m\theta_{m'}x)\,(\gamma\cdot\hat{n})\right)_b\cdot(\gamma\cdot\hat{n}')_{ba}.
\end{align*}
\end{itemize}
A product of two antilinear anti-automorphisms is a linear automorphism.
For $m'=m$, (S4) gives $\Theta_m^2=1$.
If the normals are orthogonal, then $(\gamma\cdot\hat{n}')\,(\gamma\cdot\hat{n})=-(\gamma\cdot\hat{n})\,(\gamma\cdot\hat{n}')$ and $\theta_m\theta_{m'}=\theta_{m'}\theta_m$; hence $\Theta_m\Theta_{m'}=(-1)^{\deg}\,\Theta_{m'}\Theta_m$.
From $\gamma_\mu^2=1$ and (M3) one obtains
\[
\Theta_{\mu,\mathrm{po}'}\,\Theta_{\mu,\mathrm{po}}=\mathrm{Tg}_{2(\mathrm{po}'-\mathrm{po})e_\mu},
\]
where, for $v\in\mathbb{Z}^3$, $\mathrm{Tg}_v$ denotes the automorphism $\psi(x)\mapsto\psi(x+v)$, $\bar{\psi}(x)\mapsto\bar{\psi}(x+v)$ of $\mathrm{Gr}$, with $x+v:=(c+v,g)$ for $x=(c,g)$.

\emph{The space $\mathrm{Fn}$} (\cite[Definition~5.12]{QEDII}).
Let $C^{1}$ be the space of the real one-cochains on $\mathbb{Z}^3$.
The vector space $\mathrm{Fn}$ consists of all maps $C^{1}\to\mathrm{Gr}$, $\hat{A}\mapsto P(\hat{A})$, with the pointwise product $(PP')(\hat{A}):=P(\hat{A})\,P'(\hat{A})$; an element $P$ is called even if its values are even.
The following two definitions are those of \cite[Section~2.3]{QEDII}.
\begin{itemize}
\item[(S5)] For $P\in\mathrm{Fn}$, $(\Theta_mP)(\hat{A}):=\Theta_m\left(P(\theta_m^{*}\hat{A})\right)$, with the map (S2) on the right side.
\item[(S6)] For $g_{\mathrm{w}}\in W_{\mathrm{af}}$ and an even $P\in\mathrm{Fn}$ put $P^{\mathrm{gw}}:=\Theta_{m_1}\cdots\Theta_{m_s}P$, where $g_{\mathrm{w}}=\theta_{m_1}\cdots\theta_{m_s}$ is a word in reflections in mirrors.
By \cite[Proposition~3.16\,(f$''$)]{QEDII}, two words for the same $g_{\mathrm{w}}$ give maps that agree on $\mathrm{Gr}$ up to the factor $(-1)^{\deg}$ and that act in the same way on the argument $\hat{A}$; hence $P^{\mathrm{gw}}$ is well defined for even $P$.
\end{itemize}

\begin{proof}[Proof of the last assertion for words in the mirrors $(y_\mu,\mathrm{po})$]
These are the only mirrors that occur in Sections~3--6, and for them the assertion can be seen directly.
It suffices to show that a word with $\theta_{m_1}\cdots\theta_{m_s}=\mathrm{id}$ gives $\Theta_{m_1}\cdots\Theta_{m_s}=\mathrm{id}$ on even elements.
By (S4) the maps $\Theta$ of mirrors of different directions commute on even elements, and by (M5) the corresponding reflections commute; so the word may be sorted by directions.
The part of the direction $\mu$ acts on the coordinate $\mu$ only, so that it is the identity by itself.
Its length is even, as the determinant shows; one has
\[
\theta_{\mu,\mathrm{po}_1}\theta_{\mu,\mathrm{po}_2}\cdots\theta_{\mu,\mathrm{po}_{2i}}=\tau_{2(\mathrm{po}_1-\mathrm{po}_2+\dots-\mathrm{po}_{2i})e_\mu},
\]
and, by (S4),
\[
\Theta_{\mu,\mathrm{po}_1}\Theta_{\mu,\mathrm{po}_2}\cdots\Theta_{\mu,\mathrm{po}_{2i}}=\mathrm{Tg}_{2(\mathrm{po}_1-\mathrm{po}_2+\dots-\mathrm{po}_{2i})e_\mu}.
\]
Both are the identity if and only if the vector vanishes.
\end{proof}

Written out, with $(g_{\mathrm{w}}g_{\mathrm{w}}')^{*}=(g_{\mathrm{w}}')^{*}g_{\mathrm{w}}^{*}$, the definition (S6) reads
\begin{equation*}
P^{\mathrm{gw}}(\hat{A})=\Theta_{m_1}\cdots\Theta_{m_s}\left(P\left(g_{\mathrm{w}}^{*}\hat{A}\right)\right).\tag{1.1}
\end{equation*}
One has $\left(P^{\mathrm{gw}}\right)^{\mathrm{gw}'}=P^{\mathrm{gw}'\mathrm{gw}}$ (juxtaposition of words); the map $P\mapsto P^{\mathrm{gw}}$ is linear for even $g_{\mathrm{w}}$ and antilinear for odd $g_{\mathrm{w}}$; and $P^{\mathrm{id}}=P$, because $\Theta_m^2=1$.
By (S4) and $\tau_{e_\mu}=\theta_{\mu,1/2}\,\theta_{\mu,0}$,
\begin{equation*}
P^{\tau_v}(\hat{A})=\mathrm{Tg}_v\left(P\left(\tau_v^{*}\hat{A}\right)\right)\qquad(v\in\mathbb{Z}^3).\tag{1.2}
\end{equation*}

\subsubsection*{1.4. Boxes, integrands, products}

The definitions of this subsection are those of \cite[Definition~2.12]{QEDII}.
A box is a set $\mathrm{Bx}=\mathrm{bc}+[-\mathrm{rd},\mathrm{rd}]^3$ with $\mathrm{bc}\in\mathbb{R}^3$ and $\mathrm{rd}>0$.
The sets of the links and of the plaquettes of $\mathbb{Z}^3$ that are contained in $\mathrm{Bx}$ are denoted by $\mathrm{Ln}(\mathrm{Bx})$ and $\mathrm{Pq}(\mathrm{Bx})$, and $X(\mathrm{Bx})$ is the set of the fermion sites $x=(c,g)$ with $c+[-\frac12,\frac12]^3\subset\mathrm{Bx}$; it contains either all $48$ sites of a cluster or none of them.

An \emph{integrand that reads in} $\mathrm{Bx}$ is a gauge invariant Borel map
\[
G\colon\mathbb{R}^{\mathrm{Ln}(\mathrm{Bx})}\to\mathrm{Gr}(X(\mathrm{Bx})),\qquad
\hat{A}\mapsto G(\hat{A})=\sum_{\mathrm{mm}}c_{\mathrm{mm}}(\hat{A})\cdot\mathrm{mm},
\]
where $\mathrm{mm}$ ranges over the monomials in the generators at the sites of $X(\mathrm{Bx})$.
Gauge invariance means: for every real function $\boldsymbol{\lambda}$ on the photon sites of $\mathrm{Bx}$, the map $G$ does not change if $\hat{A}$ is replaced by $\hat{A}-d\boldsymbol{\lambda}$, $\psi(x)$ by $\exp\left(iq\boldsymbol{\lambda}(\pi(x))\right)\cdot\psi(x)$, and $\bar{\psi}(x)$ by $\exp\left(-iq\boldsymbol{\lambda}(\pi(x))\right)\cdot\bar{\psi}(x)$.
The set $\mathrm{Ad}(\mathrm{Bx})$ consists of those integrands for which there are $\mathrm{Cq}<\infty$ and $\varepsilon_G>0$ with
\[
\left|c_{\mathrm{mm}}(\hat{A})\right|\le\mathrm{Cq}\cdot\exp\left(\left(\frac12-\varepsilon_G\right)\sum_{p\in\mathrm{Pq}(\mathrm{Bx})}(d\hat{A})_p^2\right)
\]
for all $\mathrm{mm}$ and $\hat{A}$; $\mathrm{Ad}^{\mathrm{ev}}(\mathrm{Bx})$ is the set of its even elements.

A \emph{product} is an expression $P=\Phi\cdot G_1\cdots G_s$, $s\ge0$, with $G_j\in\mathrm{Ad}^{\mathrm{ev}}(\mathrm{Bx}_j)$ for pairwise disjoint boxes $\mathrm{Bx}_j$ and with a bounded Borel function $\Phi$ of finitely many $F_p$; it stands for the element
\[
\hat{A}\mapsto\Phi(d\hat{A})\cdot G_1\left(\hat{A}|_{\mathrm{Ln}(\mathrm{Bx}_1)}\right)\cdots G_s\left(\hat{A}|_{\mathrm{Ln}(\mathrm{Bx}_s)}\right)
\]
of $\mathrm{Fn}$, which is even.
For an open set $\mathrm{Uo}\subset\mathbb{R}^3$, $\mathcal{A}(\mathrm{Uo})$ is the linear span in $\mathrm{Fn}$ of the products whose boxes are contained in $\mathrm{Uo}$ and whose factor $\Phi$ reads only plaquettes contained in $\mathrm{Uo}$.

Let $\Gamma'$ be a subgroup of $\mathbb{Z}^3$.
A box $\mathrm{Bx}$ \emph{embeds} in $\Gamma'$ if the sets $\mathrm{Bx}+[-1,1]^3+v$, $v\in\Gamma'$, are pairwise disjoint.
A product \emph{embeds} in $\Gamma'$ if each of its boxes embeds in $\Gamma'$ and the sets $\mathrm{Bx}_j+v$, $1\le j\le s$, $v\in\Gamma'$, are pairwise disjoint.
A product \emph{is read on} $C(n)$ if the sets $\mathrm{Bx}_j+v$, $1\le j\le s$, $v\in\Gamma(n)$, are pairwise disjoint (\cite[Definition~2.13\,(a) and Definition~5.11]{QEDII}).
The span of the products that are read on $C(n)$ is denoted by $\mathcal{A}(C(n))$.

For $\mathrm{po}<\mathrm{po}'$ in $\tfrac12\mathbb{Z}$ let, as in \cite[Definition~2.17]{QEDII},
\[
\mathcal{A}_{\mu,n}[\mathrm{po},\mathrm{po}']:=\mathcal{A}_{\mathrm{Xs}_\mu(n)}[\mathrm{po},\mathrm{po}']
\]
be the span of those products of $\mathcal{A}\left(\mathrm{Sl}_\mu(\mathrm{po},\mathrm{po}')\right)$ that embed in $\Gamma_\mu(n)$.

\subsubsection*{1.5. The functional}

The definitions of this subsection are those of \cite[Definitions~2.11 and~2.13\,(a)]{QEDII} and of Section~I.7.2.
Let $\mathrm{Gr}_n:=\mathrm{Gr}_{C(n)}$ be the Grassmann algebra with the generators $\psi_\beta(x)$, $\bar{\psi}_\beta(x)$, $x\in\mathbb{L}^{\mathrm{f}}_0$, and the relations
\[
\psi(x+v)=\chi_n(v)\cdot\psi(x),\qquad\bar{\psi}(x+v)=\chi_n(v)\cdot\bar{\psi}(x),\qquad v\in\Gamma(n);
\]
its independent generators are those at one representative of each of the $48\,|n|$ classes of sites.
The homomorphism of algebras $\mathrm{qu}_n\colon\mathrm{Gr}\to\mathrm{Gr}_n$ sends every generator to the generator of the same name; its behaviour under translations by elements of $\Gamma(n)$ is described in Lemma~2.2\,(B3).
The integral $\int_n$ is the integral (S1) over the $48\,|n|$ classes; each pair $\psi_\beta(x)\,\bar{\psi}_\beta(x)$ does not depend on the representative.
With $x$ ranging over a set of representatives, put
\[
S_n(A):=\sum_x\bar{\psi}(x)\cdot\left((D_A+\boldsymbol{\mu})\psi\right)(x)
=\sum_x\left[M_{\mathrm{ch}}\cdot\bar{\psi}(x)\,\psi(x)-\sum_{i=0}^{3}\bar{\psi}(x)\,M_i(x)\cdot U_A(x,x^{(i)})\cdot\psi(x^{(i)})\right],
\]
the second expression by (I.2.5); then $\int_n\exp\left(-S_n(A)\right)=\det\left(D_A+\boldsymbol{\mu}\right)$ (Proposition~46\,(e)).

A \emph{configuration} on $C(n)$ is a pair $(A,\mathfrak{z})$ with $A\in C^{1}(n)$, $\mathfrak{z}\in C^{2}(n;\mathbb{Z})$ and $d\mathfrak{z}=0$; its curvature is $F^{\mathrm{u}}=dA+r_q\cdot\mathfrak{z}$, and $\tilde{A}$, $\tilde{\mathfrak{z}}$, $\tilde{F}$ denote the periodic extensions to $\mathbb{Z}^3$.
A \emph{lift} of $(A,\mathfrak{z})$ is a real one-cochain $\hat{A}$ on $\mathbb{Z}^3$ with $d\hat{A}=\tilde{F}$ and $\exp\left(iq\hat{A}(b)\right)=\exp\left(iq\tilde{A}(b)\right)$ for all links $b$.
For $P\in\mathcal{A}(C(n))$ put
\begin{equation*}
\begin{gathered}
P(A,\mathfrak{z}):=\mathrm{qu}_n\left(P(\hat{A})\right)\qquad(\hat{A}\text{ a lift of }(A,\mathfrak{z})),\\
Z_n[P]:=\sum_{\mathfrak{z}:\,d\mathfrak{z}=0}\ \int_{[0,r_q)^{\mathrm{links}}}\mathrm{d}A\,\exp\left(-\frac12\|F^{\mathrm{u}}\|^2\right)\cdot\int_n\exp\left(-S_n(A)\right)\cdot P(A,\mathfrak{z}).
\end{gathered}\tag{1.3}
\end{equation*}
That (1.3) is meaningful is the content of Lemma~2.4.
If $P$ is a function $\Phi$ of the curvature in the class of Section~I.7.2 and lies in $\mathcal{A}(C(n))$, then (1.3) is (I.7.1), because $\int_n\exp\left(-S_n(A)\right)=\det\left(D_A+\boldsymbol{\mu}\right)$; this is the identification of Proposition~46\,(e), (f).

\subsubsection*{1.6. The slab operators of one coordinate direction}

\begin{named}{Theorem SO}[{\cite[Theorem~5.29 and Corollary~5.30]{QEDII} for the family $y_\mu$}]
Let $\mu\in\{1,2,3\}$, and let the two sides $n_\nu\ge4$, $\nu\ne\mu$, be given.
There are
a number $\mathrm{kk}>0$;
separable Hilbert spaces $\mathrm{Hs}(\mathrm{po})$, $\mathrm{po}\in\tfrac12\mathbb{Z}$;
unitary operators $V_v\colon\mathrm{Hs}(\mathrm{po})\to\mathrm{Hs}(\mathrm{po}+v_\mu)$, $v\in\mathbb{Z}e_\mu$, with $V_vV_{v'}=V_{v+v'}$;
and linear maps $\mathrm{Cs}$ from $\mathcal{A}_{\mu,n}[\mathrm{po},\mathrm{po}']$ to the Hilbert--Schmidt operators $\mathrm{Hs}(\mathrm{po})\to\mathrm{Hs}(\mathrm{po}')$, $\mathrm{po}<\mathrm{po}'$ in $\tfrac12\mathbb{Z}$;
all of them formed with $\mu$, with the sides $n_\nu$, $\nu\ne\mu$, and with the parameters $\rho$, $w$, $q$, $\boldsymbol{\mu}$, but not with $n_\mu$;
such that the following hold.
\begin{itemize}
\item[(SO1)] $\mathrm{Cs}[G']\,\mathrm{Cs}[G]=\mathrm{Cs}[GG']$ for $G\in\mathcal{A}_{\mu,n}[\mathrm{po},\mathrm{po}']$ and $G'\in\mathcal{A}_{\mu,n}[\mathrm{po}',\mathrm{po}'']$.
\item[(SO2a)] $V_v\,\mathrm{Cs}[G]\,V_v^{-1}=\mathrm{Cs}\left[G^{\tau_v}\right]$ for $G\in\mathcal{A}_{\mu,n}[\mathrm{po},\mathrm{po}']$ and $v\in\mathbb{Z}e_\mu$.
\item[(SO2b)] $\mathrm{Cs}\left[\Theta_{\mu,\mathrm{po}'}G\right]=V_\tau\,\mathrm{Cs}[G]^{*}$ for $G\in\mathcal{A}_{\mu,n}[\mathrm{po},\mathrm{po}']$, where $\tau:=2\,(\mathrm{po}'-\mathrm{po})\cdot e_\mu$; both sides map $\mathrm{Hs}(\mathrm{po}')$ to $\mathrm{Hs}(2\,\mathrm{po}'-\mathrm{po})$.
\item[(SO3)] For every integer $n_\mu\ge4$, every chain $\mathrm{po}_0<\mathrm{po}_1<\dots<\mathrm{po}_s=\mathrm{po}_0+n_\mu$ in $\tfrac12\mathbb{Z}$ with $s\ge2$, and all $G_k\in\mathcal{A}_{\mu,n}[\mathrm{po}_{k-1},\mathrm{po}_k]$, the product $G_1\cdots G_s$ lies in $\mathcal{A}(C(n))$, and
\[
Z_n\left[G_1\cdots G_s\right]=\mathrm{kk}^{|n|}\cdot\operatorname{Tr}_{\mathrm{Hs}(\mathrm{po}_0)}\left(V_{n_\mu e_\mu}^{-1}\,\mathrm{Cs}[G_s]\cdots\mathrm{Cs}[G_1]\right).
\]
\end{itemize}
In (SO2a) and (SO2b) the elements $G^{\tau_v}$ and $\Theta_{\mu,\mathrm{po}'}G$ lie in $\mathcal{A}_{\mu,n}[\mathrm{po}+v_\mu,\mathrm{po}'+v_\mu]$ and in $\mathcal{A}_{\mu,n}[\mathrm{po}',2\,\mathrm{po}'-\mathrm{po}]$, respectively (Lemma~2.2\,(B2)).
\end{named}

The planes $\mathrm{po}$, $\mathrm{po}'$, $\mathrm{po}_k$ are arbitrary elements of $\tfrac12\mathbb{Z}$: site planes ($\mathrm{po}\in\mathbb{Z}$) and link planes ($\mathrm{po}\in\mathbb{Z}+\frac12$), in any order.
Theorem~SO is \cite[Corollary~5.30]{QEDII}, together with \cite[Theorem~5.29]{QEDII}, for the family $y_\mu$, the cross-section $\mathrm{Xs}_\mu(n)$, the translations along $e_\mu$ and the closing vector $n_\mu e_\mu$ of $C(n)$ (Section~1.1), and in this form it holds with $\mathrm{kk}=1$.
Indeed, the trace formula of \cite[Theorem~5.29\,(e)]{QEDII} carries no factor $\mathrm{kk}^{|n|}$, and its factor $U_{\mathrm{f}}^{\frac{1+\chi(g)}{2}}$, formed with the parity $U_{\mathrm{f}}$ of \cite[Definition~5.27]{QEDII}, is the identity for the closing vector $g=n_\mu e_\mu$, because $\chi_n(n_\mu e_\mu)=-1$.
In Sections~3--6, Theorem~SO is applied to chains of the walls $\mathrm{po}_j=o_\mu+j\cdot\mathrm{hw}$ of Section~3.1, where $o_\mu,\mathrm{hw}\in\tfrac12\mathbb{Z}$ and $\mathrm{hw}>0$.
For $i\in\mathbb{Z}$ the difference $\mathrm{po}_{j+i}-\mathrm{po}_j=i\cdot\mathrm{hw}$ lies in $\mathbb{Z}+\tfrac12$ if and only if $i$ and $2\,\mathrm{hw}$ are both odd.
Hence two consecutive walls $\mathrm{po}_j$, $\mathrm{po}_{j+1}$ are of different kinds exactly if $2\,\mathrm{hw}$ is odd; in that case site planes and link planes alternate in the chains of Sections~3--6, whose consecutive members differ by odd multiples of $\mathrm{hw}$.
If $\mathrm{hw}$ is an integer, all walls $\mathrm{po}_j$ are of the kind of $o_\mu$.
Appendix~A proves Theorem~SO, through Theorem~A together with Corollary~A.6, for planes at arbitrary positions in $\tfrac12\mathbb{Z}$, with the spaces, the measures and the kernels at the two kinds of planes written out; the alternating chain of Section~3 is the special case treated in Section~A.11.

\subsubsection*{1.7. Schatten classes}

For $1\le p<\infty$ let $S_p$ be the Schatten class of order $p$ of a single Hilbert space, normed by $\|Y\|_p^p=\operatorname{Tr}|Y|^p$.
The following three facts are \cite[Lemma~10.6]{QEDI}.
\begin{itemize}
\item[(a)] If $Y_i\in S_{p_i}$ for $i=1,\dots,k$ and $\sum_{i=1}^{k}\frac{1}{p_i}=1$, then
\[
\left|\operatorname{Tr}\left(Y_1\cdots Y_k\right)\right|\le\prod_{i=1}^{k}\|Y_i\|_{p_i}.
\]
\item[(b)] If $p\ge p'$, then $\|Y\|_p\le\|Y\|_{p'}$.
For every unitary operator $U$ one has $\|U\cdot Y\|_p=\|Y\cdot U\|_p=\|Y\|_p$; moreover $\|Y^{*}\|_p=\|Y\|_p$, and $\|Y\|_p^p=\operatorname{Tr}\left(Y^{*}\cdot Y\right)^{\frac{p}{2}}$ if $p$ is even.
\item[(c)] Let $B\ge0$ be an operator of trace class.
Then $\operatorname{Tr}B^k\le\left(\operatorname{Tr}B\right)^k$ for every integer $k\ge1$.
If in addition $B\ne0$, the function $s\mapsto\log\operatorname{Tr}B^s$ is convex on $[1,\infty)$.
\end{itemize}

\subsubsection*{1.8. Positivity of the partition function}

For every real $\boldsymbol{\mu}$ and all sides $n_\mu\ge4$ one has
\[
0<Z_n[1]<\infty
\]
(Section~I.7.2, Proposition~50\,(e), \cite[Theorem~4.26]{QEDII}).

\subsection*{2. Four lemmas}

Let $\mathrm{Hb}$ and $\mathrm{Hb}'$ be Hilbert spaces, let $Y\colon\mathrm{Hb}\to\mathrm{Hb}'$ be a Hilbert--Schmidt operator, and let $k\ge1$ be an integer.
Put
\[
\|Y\|_{2k}:=\left(\operatorname{Tr}_{\mathrm{Hb}}\left(Y^{*}\cdot Y\right)^k\right)^{\frac{1}{2k}}.
\]
This number is finite: $Y^{*}\cdot Y$ is an operator of trace class on $\mathrm{Hb}$, and 1.7\,(b) applies to it.

\begin{named}{Lemma 2.1}[{H\"older's inequality between different spaces; cf.\ \cite[Lemma~6.3]{QEDII}}]
Let $s\ge2$, let $\mathrm{Hb}_0,\dots,\mathrm{Hb}_s$ be separable Hilbert spaces, let $Y_j\colon\mathrm{Hb}_{j-1}\to\mathrm{Hb}_j$ be Hilbert--Schmidt operators for $j=1,\dots,s$, and let $U\colon\mathrm{Hb}_s\to\mathrm{Hb}_0$ be unitary.
Then $U\cdot Y_s\cdots Y_1$ is an operator of trace class on $\mathrm{Hb}_0$.
If $s=2k$, then
\[
\left|\operatorname{Tr}_{\mathrm{Hb}_0}\left(U\cdot Y_{2k}\cdots Y_1\right)\right|\le\prod_{j=1}^{2k}\|Y_j\|_{2k}.
\]
\end{named}

\begin{proof}
Let $\mathrm{Hb}:=\mathrm{Hb}_0\oplus\cdots\oplus\mathrm{Hb}_{s-1}$, and for $0\le j\le s-1$ let $\mathrm{em}_j\colon\mathrm{Hb}_j\to\mathrm{Hb}$ be the canonical isometry, so that $\mathrm{em}_j^{*}\cdot\mathrm{em}_j=1$.
Define operators on $\mathrm{Hb}$ by
\[
\hat{Y}_j:=\mathrm{em}_j\cdot Y_j\cdot\mathrm{em}_{j-1}^{*}\quad(1\le j\le s-1),\qquad\hat{Y}_s:=\mathrm{em}_0\cdot U\cdot Y_s\cdot\mathrm{em}_{s-1}^{*}.
\]
For every $j$ one has
\[
\hat{Y}_j^{*}\cdot\hat{Y}_j=\mathrm{em}_{j-1}\cdot Y_j^{*}\cdot Y_j\cdot\mathrm{em}_{j-1}^{*};
\]
for $j=s$ this uses in addition $U^{*}\cdot U=1$.
Consequently, for every integer $k\ge1$,
\[
\left(\hat{Y}_j^{*}\cdot\hat{Y}_j\right)^k=\mathrm{em}_{j-1}\cdot\left(Y_j^{*}\cdot Y_j\right)^k\cdot\mathrm{em}_{j-1}^{*},
\qquad
\operatorname{Tr}_{\mathrm{Hb}}\left(\hat{Y}_j^{*}\cdot\hat{Y}_j\right)^k=\operatorname{Tr}_{\mathrm{Hb}_{j-1}}\left(Y_j^{*}\cdot Y_j\right)^k.
\]
Thus $\hat{Y}_j$ belongs to $S_2(\mathrm{Hb})$, which is contained in $S_{2k}(\mathrm{Hb})$, and by 1.7\,(b) its norm in $S_{2k}(\mathrm{Hb})$ is $\|Y_j\|_{2k}$.
As $\mathrm{em}_j^{*}\cdot\mathrm{em}_j=1$, the product of the $\hat{Y}_j$ telescopes:
\[
\hat{Y}_s\cdots\hat{Y}_1=\mathrm{em}_0\cdot\left(U\cdot Y_s\cdots Y_1\right)\cdot\mathrm{em}_0^{*}.
\]
A product of $s\ge2$ Hilbert--Schmidt operators is of trace class; hence so is $U\cdot Y_s\cdots Y_1$.
Moreover $\operatorname{Tr}_{\mathrm{Hb}}\left(\mathrm{em}_0\cdot B\cdot\mathrm{em}_0^{*}\right)=\operatorname{Tr}_{\mathrm{Hb}_0}B$ for every operator $B$ of trace class on $\mathrm{Hb}_0$.
Let now $s=2k$.
By 1.7\,(a), applied on $\mathrm{Hb}$ with the $2k$ exponents $p_i=2k$,
\[
\left|\operatorname{Tr}_{\mathrm{Hb}_0}\left(U\cdot Y_{2k}\cdots Y_1\right)\right|
=\left|\operatorname{Tr}_{\mathrm{Hb}}\left(\hat{Y}_{2k}\cdots\hat{Y}_1\right)\right|
\le\prod_{j=1}^{2k}\|\hat{Y}_j\|_{2k}
=\prod_{j=1}^{2k}\|Y_j\|_{2k}.
\]
\end{proof}

\begin{named}{Lemma 2.2}[{The action of $W_{\mathrm{af}}$ on $\mathcal{A}$; cf.\ \cite[Lemma~6.4]{QEDII}}]
Let $g_{\mathrm{w}}=(a,r)\in W_{\mathrm{af}}$.
\begin{itemize}
\item[(B1)] Let $P_1,P_2\in\mathrm{Fn}$ be even, and let $z\in\mathbb{C}$.
Then $(P_1\cdot P_2)^{\mathrm{gw}}=P_1^{\mathrm{gw}}\cdot P_2^{\mathrm{gw}}$ and $1^{\mathrm{gw}}=1$; further $(z\cdot P_1)^{\mathrm{gw}}=z\cdot P_1^{\mathrm{gw}}$ if $g_{\mathrm{w}}$ is even, and $(z\cdot P_1)^{\mathrm{gw}}=\overline{z}\cdot P_1^{\mathrm{gw}}$ if $g_{\mathrm{w}}$ is odd.
In a product of even elements the order of the factors is immaterial.
\item[(B2)] Let $\mathrm{Bx}=\mathrm{bc}+[-\mathrm{rd},\mathrm{rd}]^3$ be a box.
Then $g_{\mathrm{w}}\mathrm{Bx}=g_{\mathrm{w}}\mathrm{bc}+[-\mathrm{rd},\mathrm{rd}]^3$ is a box, and
\[
\mathrm{Ln}(g_{\mathrm{w}}\mathrm{Bx})=g_{\mathrm{w}}\mathrm{Ln}(\mathrm{Bx}),\qquad
\mathrm{Pq}(g_{\mathrm{w}}\mathrm{Bx})=g_{\mathrm{w}}\mathrm{Pq}(\mathrm{Bx}),\qquad
X(g_{\mathrm{w}}\mathrm{Bx})=g_{\mathrm{w}}X(\mathrm{Bx}).
\]
If $G\in\mathrm{Ad}^{\mathrm{ev}}(\mathrm{Bx})$, then $G^{\mathrm{gw}}\in\mathrm{Ad}^{\mathrm{ev}}(g_{\mathrm{w}}\mathrm{Bx})$, with the same $\varepsilon_G$.
Let $\mathrm{Uo}\subset\mathbb{R}^3$ be open, and let $\Phi$ be a bounded Borel function of finitely many $F_p$ with plaquettes $p$ contained in $\mathrm{Uo}$.
Then $\Phi^{\mathrm{gw}}$ is a bounded Borel function of finitely many $F_p$ with plaquettes $p$ contained in $g_{\mathrm{w}}\mathrm{Uo}$, it has the same supremum as $\Phi$, and
\[
\Phi^{\mathrm{gw}}(u)=\Phi(g_{\mathrm{w}}^{*}u)\ \text{ if }g_{\mathrm{w}}\text{ is even},\qquad
\Phi^{\mathrm{gw}}(u)=\overline{\Phi(g_{\mathrm{w}}^{*}u)}\ \text{ if }g_{\mathrm{w}}\text{ is odd}.
\]
Consequently, the image of a product with the boxes $\mathrm{Bx}_j$ is a product with the boxes $g_{\mathrm{w}}\mathrm{Bx}_j$; $P\in\mathcal{A}(\mathrm{Uo})$ implies $P^{\mathrm{gw}}\in\mathcal{A}(g_{\mathrm{w}}\mathrm{Uo})$; and if a product embeds in a subgroup $\Gamma'$ of $\mathbb{Z}^3$, then its image embeds in $r\Gamma'$.
\item[(B3)] If $P\in\mathcal{A}(C(n))$ and $v\in\Gamma(n)$, then $P^{\tau_v}\in\mathcal{A}(C(n))$, and $P^{\tau_v}(A,\mathfrak{z})=P(A,\mathfrak{z})$ for every configuration $(A,\mathfrak{z})$ on $C(n)$.
\end{itemize}
\end{named}

\begin{proof}
(B1) By (S6) of 1.3 it suffices to take for $g_{\mathrm{w}}$ a reflection $\theta_m$.
By (S2) and (S5), $\Theta_m(P_1\cdot P_2)=\Theta_m(P_2)\cdot\Theta_m(P_1)$.
Since $\Theta_m$ preserves the degree and even elements are central, the right side equals $\Theta_m(P_1)\cdot\Theta_m(P_2)$.
The map $\Theta_m$ is antilinear; therefore a word of even length acts linearly, and a word of odd length acts antilinearly.

(B2) The linear part $r$ is a signed permutation matrix, and it maps $[-\mathrm{rd},\mathrm{rd}]^3$ onto itself; hence $g_{\mathrm{w}}\left(\mathrm{bc}+[-\mathrm{rd},\mathrm{rd}]^3\right)=g_{\mathrm{w}}\mathrm{bc}+[-\mathrm{rd},\mathrm{rd}]^3$.
The map $g_{\mathrm{w}}$ permutes the cells of $\mathbb{Z}^3$, and a cell lies in $\mathrm{Bx}$ if and only if its image lies in $g_{\mathrm{w}}\mathrm{Bx}$.
On fermion sites one has $g_{\mathrm{w}}(c,g)=(r\cdot c+a,\,r\cdot g)$ (Section~I.4.7), and $g_{\mathrm{w}}\left(c+\left[-\frac12,\frac12\right]^3\right)=g_{\mathrm{w}}c+\left[-\frac12,\frac12\right]^3$; therefore $X(g_{\mathrm{w}}\mathrm{Bx})=g_{\mathrm{w}}X(\mathrm{Bx})$.

By (S6) it suffices to treat a reflection $g_{\mathrm{w}}=\theta=\theta_m$, with the unit normal $\hat{n}$.
Let $G(\hat{A})=\sum_{\mathrm{mm}}c_{\mathrm{mm}}(\hat{A})\cdot\mathrm{mm}$ be an element of $\mathrm{Ad}^{\mathrm{ev}}(\mathrm{Bx})$, the sum being over the monomials $\mathrm{mm}$.
By (S5),
\[
(\Theta_mG)(\hat{A})=\sum_{\mathrm{mm}}\overline{c_{\mathrm{mm}}(\theta^{*}\hat{A})}\cdot\Theta_m(\mathrm{mm}),
\]
where $\hat{A}$ is a cochain on $\mathrm{Ln}(\theta\mathrm{Bx})$ and $\theta^{*}\hat{A}$ is the cochain $(c,c')\mapsto\hat{A}(\theta c,\theta c')$ on $\mathrm{Ln}(\mathrm{Bx})$.
This map is Borel and even, and its values lie in $\mathrm{Gr}\left(\theta X(\mathrm{Bx})\right)$.

\emph{Growth.}
Each $\Theta_m(\mathrm{mm})$ is a linear combination of monomials whose coefficients are products of entries of $\gamma\cdot\hat{n}$.
So the coefficients of $\Theta_mG$ at $\hat{A}$ are bounded by a constant times
\[
\max_{\mathrm{mm}}\left|c_{\mathrm{mm}}(\theta^{*}\hat{A})\right|\le\mathrm{Cq}\cdot\exp\left(\left(\frac12-\varepsilon_G\right)\sum_{p\in\mathrm{Pq}(\mathrm{Bx})}\left(d\theta^{*}\hat{A}\right)_p^2\right).
\]
Since $d\theta^{*}=\theta^{*}d$ and $(\theta^{*}u)_p=\pm u_{\theta p}$ (Section~I.4.7), the sum in the exponent equals $\sum_{p'\in\mathrm{Pq}(\theta\mathrm{Bx})}\left(d\hat{A}\right)_{p'}^2$.

\emph{Gauge invariance.}
Let $\boldsymbol{\lambda}$ be a real function on the photon sites of $\theta\mathrm{Bx}$, and put $\boldsymbol{\lambda}':=\boldsymbol{\lambda}\circ\theta$.
Under the gauge transformation with $\boldsymbol{\lambda}$, a factor $\bar{\psi}(\theta x)$ of $\Theta_m(\mathrm{mm})$, which comes from a factor $\psi(x)$ of $\mathrm{mm}$, is multiplied by
\[
\exp\left(-iq\boldsymbol{\lambda}\left(\pi(\theta x)\right)\right)=\overline{\exp\left(iq\boldsymbol{\lambda}'\left(\pi(x)\right)\right)},
\]
because $\pi(\theta x)=\theta\pi(x)$; and a factor $\psi(\theta x)$, which comes from a factor $\bar{\psi}(x)$, is multiplied by $\overline{\exp\left(-iq\boldsymbol{\lambda}'\left(\pi(x)\right)\right)}$.
Hence $\Theta_m(\mathrm{mm})$ is multiplied by the complex conjugate of the factor $z$ which $\mathrm{mm}$ receives under the gauge transformation with $\boldsymbol{\lambda}'$.
Moreover $\theta^{*}(\hat{A}-d\boldsymbol{\lambda})=\theta^{*}\hat{A}-d\boldsymbol{\lambda}'$.
The gauge invariance of $G$ states that $c_{\mathrm{mm}}(\hat{A}'-d\boldsymbol{\lambda}')\cdot z=c_{\mathrm{mm}}(\hat{A}')$ for every $\mathrm{mm}$; the complex conjugate of this identity at $\hat{A}'=\theta^{*}\hat{A}$ is the gauge invariance of $\Theta_mG$.

\emph{Bounded factors.}
One has $(\Theta_m\Phi)(\hat{A})=\overline{\Phi(d\theta^{*}\hat{A})}=\overline{\Phi(\theta^{*}d\hat{A})}$, and $\theta^{*}u$ reads $u$ at the plaquettes $\theta p$.
For a general $g_{\mathrm{w}}$ one iterates along a word in reflections and uses $(g_{\mathrm{w}}g_{\mathrm{w}}')^{*}=g_{\mathrm{w}}'^{*}g_{\mathrm{w}}^{*}$.

\emph{Embedding.}
By (B1), the image of $\Phi\cdot G_1\cdots G_s$ is $\Phi^{\mathrm{gw}}\cdot G_1^{\mathrm{gw}}\cdots G_s^{\mathrm{gw}}$, with the boxes $g_{\mathrm{w}}\mathrm{Bx}_j$.
Furthermore
\[
g_{\mathrm{w}}\left(\mathrm{Bx}_j+[-1,1]^3+v\right)=g_{\mathrm{w}}\mathrm{Bx}_j+[-1,1]^3+r\cdot v,\qquad
g_{\mathrm{w}}\left(\mathrm{Bx}_j+v\right)=g_{\mathrm{w}}\mathrm{Bx}_j+r\cdot v,
\]
and $g_{\mathrm{w}}$ is a bijection; so disjointness is preserved.

(B3) By linearity it suffices to consider a product $P$ that is read on $C(n)$, with the boxes $\mathrm{Bx}_j$.
By (B2), $P^{\tau_v}$ is a product with the boxes $\mathrm{Bx}_j+v$.
The family of the sets $\mathrm{Bx}_j+v+v'$, $v'\in\Gamma(n)$, coincides with the family of the sets $\mathrm{Bx}_j+v'$, $v'\in\Gamma(n)$; hence $P^{\tau_v}$ is read on $C(n)$.
Let $\hat{A}$ be a lift of $(A,\mathfrak{z})$.
Then $\tau_v^{*}\hat{A}$ is also a lift of $(A,\mathfrak{z})$: since $\tilde{A}$ and $\tilde{F}$ are $\Gamma(n)$-periodic, $d\tau_v^{*}\hat{A}=\tau_v^{*}\tilde{F}=\tilde{F}$, and
\[
\exp\left(iq\left(\tau_v^{*}\hat{A}\right)(b)\right)=\exp\left(iq\hat{A}(b+v)\right)=\exp\left(iq\tilde{A}(b+v)\right)=\exp\left(iq\tilde{A}(b)\right).
\]
By the relations of $\mathrm{Gr}_n$, on every homogeneous element the map $\mathrm{qu}_n\circ\mathrm{Tg}_v$ is $\mathrm{qu}_n$ multiplied by $\chi_n(v)$ raised to the degree of the element; in particular $\mathrm{qu}_n\circ\mathrm{Tg}_v=\mathrm{qu}_n$ on even elements.
By (1.2), (1.3) and Lemma~2.4\,(b),
\[
P^{\tau_v}(A,\mathfrak{z})=\mathrm{qu}_n\left(\mathrm{Tg}_v\left(P(\tau_v^{*}\hat{A})\right)\right)=\mathrm{qu}_n\left(P(\tau_v^{*}\hat{A})\right)=P(A,\mathfrak{z}).
\]
\end{proof}

\begin{named}{Lemma 2.3}[{Boxes; cf.\ \cite[Lemma~6.5(a), (b), (d)]{QEDII}}]
Let $\mathrm{Bx}=\mathrm{bc}+[-\mathrm{rd},\mathrm{rd}]^3$ be a box.
\begin{itemize}
\item[(C1)] Let $\Gamma'$ be a subgroup of $\mathbb{Z}^3$ such that every $v\in\Gamma'\setminus\{0\}$ satisfies $|v|\ge4$ and $|v|>4\mathrm{rd}$.
Then $\mathrm{Bx}$ embeds in $\Gamma'$.
\item[(C2)] If $\mathrm{Bx}\subset\mathrm{Sl}_\mu(\mathrm{po},\mathrm{po}')$, then $2\mathrm{rd}<\mathrm{po}'-\mathrm{po}$.
\item[(C4)] Let $\mathcal{Q}=(\mathfrak{F}_{\mathrm{c}};o;\mathrm{hw})$ be as in Definition~39.0, and let $C(n)$ carry $\mathcal{Q}$.
Let $\mathfrak{c}^1,\dots,\mathfrak{c}^s\in\mathbb{Z}^3$ be such that the cells $Q^{\circ}_{\mathfrak{c}^1},\dots,Q^{\circ}_{\mathfrak{c}^s}$ lie in pairwise different orbits of $\Gamma(n)$, and for $i=1,\dots,s$ let $P_i$ be a product all of whose boxes, and all of whose plaquettes read by the bounded factor, lie in $Q^{\circ}_{\mathfrak{c}^i}$.
Then $P_1\cdots P_s$ is a product; it embeds in $\Gamma(n)$ and in every subgroup of $\Gamma(n)$; in particular, it is read on $C(n)$.
If all $\mathfrak{c}^i$ have the same entry $\mathfrak{c}^i_\mu=j$, then
\[
P_1\cdots P_s\in\mathcal{A}_{\mu,n}\left[o_\mu+j\cdot\mathrm{hw},\,o_\mu+(j+1)\cdot\mathrm{hw}\right].
\]
\end{itemize}
\end{named}

\begin{proof}
(C1) The set $\mathrm{Bx}+[-1,1]^3=\mathrm{bc}+[-\mathrm{rd}-1,\mathrm{rd}+1]^3$ meets its translate by a vector $v$ if and only if $|v|\le2\mathrm{rd}+2$.
Let $v\in\Gamma'\setminus\{0\}$.
If $\mathrm{rd}<1$, then $2\mathrm{rd}+2<4\le|v|$; if $\mathrm{rd}\ge1$, then $2\mathrm{rd}+2\le4\mathrm{rd}<|v|$.
As two of the sets $\mathrm{Bx}+[-1,1]^3+v$, $v\in\Gamma'$, differ by a translation by a non-zero element of $\Gamma'$, they are pairwise disjoint.

(C2) On $\mathrm{Bx}$ the coordinate $y_\mu$ takes exactly the values of the interval $[\mathrm{bc}_\mu-\mathrm{rd},\mathrm{bc}_\mu+\mathrm{rd}]$.
This closed interval has the length $2\mathrm{rd}$, and it is contained in the open interval $(\mathrm{po},\mathrm{po}')$.

(C4) Let $\mathrm{Bx}$ be a box contained in a cell $Q^{\circ}_{\mathfrak{c}}$.
For each $\mu$ it lies in $\mathrm{Sl}_\mu\left(o_\mu+\mathfrak{c}_\mu\cdot\mathrm{hw},\,o_\mu+(\mathfrak{c}_\mu+1)\cdot\mathrm{hw}\right)$; so $2\mathrm{rd}<\mathrm{hw}$ by (C2).
Let $v\in\Gamma(n)\setminus\{0\}$.
Some entry $v_\mu$ of $v$ is non-zero, and it lies in $n_\mu\mathbb{Z}$, a subset of $2\mathrm{hw}\mathbb{Z}$.
Hence $|v|\ge|v_\mu|\ge2\mathrm{hw}>4\mathrm{rd}$, and $|v|\ge\min_\mu n_\mu\ge4$.
By (C1), $\mathrm{Bx}$ embeds in $\Gamma(n)$, hence in every subgroup of $\Gamma(n)$.

For $v\in\Gamma(n)$ the set $\tau_vQ^{\circ}_{\mathfrak{c}}$ is the cell $Q^{\circ}_{\mathfrak{c}'}$ with $\mathfrak{c}'_\mu=\mathfrak{c}_\mu+\frac{v_\mu}{\mathrm{hw}}$, and two open cells are either equal or disjoint.
Let $\mathrm{Bx}$ be a box of $P_i$ and $\mathrm{Bx}'$ a box of $P_j$, and let $v,v'\in\Gamma(n)$ be such that $(\mathrm{Bx}+v)\cap(\mathrm{Bx}'+v')$ is not empty.
Then $\tau_vQ^{\circ}_{\mathfrak{c}^i}=\tau_{v'}Q^{\circ}_{\mathfrak{c}^j}$.
So the two cells $Q^{\circ}_{\mathfrak{c}^i}$ and $Q^{\circ}_{\mathfrak{c}^j}$ lie in one orbit, whence $i=j$; and $\tau_{v-v'}$ fixes a cell, whence $v=v'$.
It follows that $\mathrm{Bx}$ and $\mathrm{Bx}'$ are boxes of one and the same product, and therefore they are equal or disjoint.
For $v=v'=0$ this shows that boxes of different factors are disjoint.
Let $\Phi_i$ be the bounded factor of $P_i$.
Then $P_1\cdots P_s$ is the product of the bounded function $\prod_i\Phi_i$ with all the factors $G$ of $P_1,\dots,P_s$, and it is a product in the sense of 1.4; by what has been shown, it embeds in $\Gamma(n)$ and in every subgroup of $\Gamma(n)$.
As for the last assertion: all boxes and all plaquettes that are read lie in the open slab $\mathrm{Sl}_\mu\left(o_\mu+j\cdot\mathrm{hw},\,o_\mu+(j+1)\cdot\mathrm{hw}\right)$, and the product embeds in the subgroup $\Gamma_\mu(n)$.
\end{proof}

\begin{named}{Lemma 2.4}[{Lifts; reading on a torus; convergence; cf.\ \cite[Lemmas~5.13, 5.14(a), (c), Proposition~5.17(c), Lemmas~3.14, 3.15]{QEDII}}]
\leavevmode
\begin{itemize}
\item[(a)] Let $\mathbf{K}=I_1\times I_2\times I_3$ be a product of three non-empty intervals of $\mathbb{Z}$, finite or infinite, regarded as the complex of those cells of $\mathbb{Z}^3$ all of whose corners lie in $\mathbf{K}$.
Every integer two-cochain on $\mathbf{K}$ that is closed on the cubes of $\mathbf{K}$ is of the form $dm$ with an integer one-cochain $m$ on $\mathbf{K}$.
Every integer one-cochain on $\mathbf{K}$ that is closed on the plaquettes of $\mathbf{K}$ is of the form $d\boldsymbol{\lambda}$ with an integer valued function $\boldsymbol{\lambda}$ on $\mathbf{K}$.
Consequently, every configuration $(A,\mathfrak{z})$ on $C(n)$ has a lift.
\item[(b)] For $P\in\mathcal{A}(C(n))$ the element $\mathrm{qu}_n\left(P(\hat{A})\right)$ does not depend on the choice of the lift $\hat{A}$ of $(A,\mathfrak{z})$.
For $m'\in C^{1}(n;\mathbb{Z})$ the configuration $(A+r_q\cdot m',\,\mathfrak{z}-dm')$ has the same lifts as $(A,\mathfrak{z})$.
The map $P\mapsto P(\cdot,\cdot)$ is linear, and for every $r\ge2$
\[
(P_1\cdots P_r)(A,\mathfrak{z})=P_1(A,\mathfrak{z})\cdots P_r(A,\mathfrak{z})
\]
whenever $P_1,\dots,P_r$ and the product $P_1\cdots P_r$ lie in $\mathcal{A}(C(n))$.
\item[(c)] For $P\in\mathcal{A}(C(n))$ there are $\mathrm{Cq}<\infty$ and $\varepsilon\in\left(0,\frac12\right]$ such that $(A,\mathfrak{z})\mapsto P(A,\mathfrak{z})$ is a Borel map with values in $\mathrm{Gr}_n^{\mathrm{ev}}$ whose coefficients in a monomial basis have modulus at most $\mathrm{Cq}\cdot\exp\left(\left(\frac12-\varepsilon\right)\|F^{\mathrm{u}}\|^2\right)$.
\item[(d)] Let $P$ be as in (c).
Every coefficient, in a monomial basis of $\mathrm{Gr}_n$, of
\[
\exp\left(-\frac{\|F^{\mathrm{u}}\|^2}{2}\right)\cdot\exp\left(-S_n(A)\right)\cdot P(A,\mathfrak{z})
\]
is absolutely summable over the integer two-cochains $\mathfrak{z}$ with $d\mathfrak{z}=0$ and absolutely integrable over $A\in[0,r_q)^{\mathrm{links}}$, the product of the intervals $[0,r_q)$ over the links of $C(n)$; more precisely,
\[
\sum_{\mathfrak{z}}\int\mathrm{d}A\,\left|\text{coefficient}\right|\le C_n\cdot\mathrm{Cq}\cdot\left(r_q\cdot\boldsymbol{\Theta}_{\varepsilon}\right)^{3|n|},
\qquad
\boldsymbol{\Theta}_{\varepsilon}:=\sup_{y\in\mathbb{R}}\sum_{j\in\mathbb{Z}}\exp\left(-\varepsilon\left(y+r_q\cdot j\right)^2\right)\le2+r_q^{-1}\cdot\left(\frac{\pi}{\varepsilon}\right)^{\frac12},
\]
where $C_n$ depends only on $n$ and on the parameters.
Hence $Z_n[P]$ is defined by an absolutely convergent sum and integral, and it is linear in $P$.
The first assertion of (d) holds, with the same bound, for every Borel map $(A,\mathfrak{z})\mapsto P(A,\mathfrak{z})$ from the configurations on $C(n)$ to $\mathrm{Gr}_n$ whose coefficients in a monomial basis have modulus at most $\mathrm{Cq}\cdot\exp\left(\left(\frac12-\varepsilon\right)\|F^{\mathrm{u}}\|^2\right)$ for some $\mathrm{Cq}<\infty$ and $\varepsilon\in\left(0,\frac12\right]$.
\end{itemize}
\end{named}

\begin{proof}
(a) For integers $a$, $t$ and a function $f$ on $\mathbb{Z}$ put
\[
S_a^tf:=\sum_{a\le j<t}f(j)\ \text{ if }t\ge a,\qquad S_a^tf:=-\sum_{t\le j<a}f(j)\ \text{ if }t<a;
\]
thus $S_a^af=0$ and $S_a^{t+1}f-S_a^tf=f(t)$.
Fix a point $(a_1,a_2,a_3)$ of $\mathbf{K}$.
For a two-cochain $u$ put
\[
m_1:=0,\qquad m_2(c):=S_{a_1}^{c_1}u_{12}(\cdot,c_2,c_3),\qquad m_3(c):=S_{a_1}^{c_1}u_{13}(\cdot,c_2,c_3)+S_{a_2}^{c_2}u_{23}(a_1,\cdot,c_3).
\]
All cells that occur lie in $\mathbf{K}$, because the $I_\mu$ are intervals.
By Section~I.1.1,
\[
(dm)_{12}(c)=m_2(c+e_1)-m_2(c)=u_{12}(c),\qquad(dm)_{13}(c)=m_3(c+e_1)-m_3(c)=u_{13}(c),
\]
and
\begin{align*}
(dm)_{23}(c)&=m_2(c)-m_2(c+e_3)+m_3(c+e_2)-m_3(c)\\
&=u_{23}(a_1,c_2,c_3)+S_{a_1}^{c_1}\left[u_{12}(\cdot,c_2,c_3)-u_{12}(\cdot,c_2,c_3+1)+u_{13}(\cdot,c_2+1,c_3)-u_{13}(\cdot,c_2,c_3)\right].
\end{align*}
Since $(du)(j,c_2,c_3)=0$, the bracket at $j$ equals $u_{23}(j+1,c_2,c_3)-u_{23}(j,c_2,c_3)$; so the sum telescopes to $u_{23}(c)-u_{23}(a_1,c_2,c_3)$, and $(dm)_{23}(c)=u_{23}(c)$.

For a closed one-cochain $u$ put
\[
\boldsymbol{\lambda}(c):=S_{a_1}^{c_1}u_1(\cdot,a_2,a_3)+S_{a_2}^{c_2}u_2(c_1,\cdot,a_3)+S_{a_3}^{c_3}u_3(c_1,c_2,\cdot).
\]
Then $(d\boldsymbol{\lambda})_3=u_3$, and
\[
(d\boldsymbol{\lambda})_2(c)=u_2(c_1,c_2,a_3)+S_{a_3}^{c_3}\left[u_3(c_1,c_2+1,\cdot)-u_3(c_1,c_2,\cdot)\right]=u_2(c),
\]
because $(du)_{23}=0$ turns the bracket at $j$ into $u_2(c_1,c_2,j+1)-u_2(c_1,c_2,j)$.
Further
\[
(d\boldsymbol{\lambda})_1(c)=u_1(c_1,a_2,a_3)+S_{a_2}^{c_2}\left[u_2(c_1+1,\cdot,a_3)-u_2(c_1,\cdot,a_3)\right]+S_{a_3}^{c_3}\left[u_3(c_1+1,c_2,\cdot)-u_3(c_1,c_2,\cdot)\right].
\]
The identity $(du)_{12}=0$ turns the first bracket at $j$ into $u_1(c_1,j+1,a_3)-u_1(c_1,j,a_3)$, and $(du)_{13}=0$ turns the second one into $u_1(c_1,c_2,j+1)-u_1(c_1,c_2,j)$.
The two sums telescope, and
\[
(d\boldsymbol{\lambda})_1(c)=u_1(c_1,a_2,a_3)+\left[u_1(c_1,c_2,a_3)-u_1(c_1,a_2,a_3)\right]+\left[u_1(c)-u_1(c_1,c_2,a_3)\right]=u_1(c).
\]
If $u$ is integer valued, then so are $m$ and $\boldsymbol{\lambda}$.

\emph{Lift.}
The periodic extension $\tilde{\mathfrak{z}}$ is a closed integer two-cochain on $\mathbf{K}=\mathbb{Z}^3$.
Write $\tilde{\mathfrak{z}}=dm$ as above, and put $\hat{A}:=\tilde{A}+r_q\cdot m$.
Then $d\hat{A}=d\tilde{A}+r_q\cdot\tilde{\mathfrak{z}}=\tilde{F}$, and $\exp\left(iq\,r_q\cdot m(b)\right)=1$ for every link $b$.

(b) Let $\hat{A}$ and $\hat{A}'$ be lifts of $(A,\mathfrak{z})$.
The cochain $\frac{1}{r_q}\left(\hat{A}-\hat{A}'\right)$ is integer valued, because the phases are equal and $q\cdot r_q=2\pi$; and it is closed.
By (a) it is $d\boldsymbol{\lambda}$ with an integer valued function $\boldsymbol{\lambda}$.
Let $P=\Phi\cdot G_1\cdots G_s$ be a product.
Then $\Phi(d\hat{A})=\Phi(\tilde{F})=\Phi(d\hat{A}')$.
The gauge invariance of $G_j$ with the real function $r_q\cdot\boldsymbol{\lambda}$ reads
\[
G_j\left(\hat{A}-r_q\cdot d\boldsymbol{\lambda};\ \exp(2\pi i\boldsymbol{\lambda})\cdot\psi,\ \exp(-2\pi i\boldsymbol{\lambda})\cdot\bar{\psi}\right)=G_j\left(\hat{A};\psi,\bar{\psi}\right),
\]
and $\exp(2\pi i\boldsymbol{\lambda})=1$; so $G_j(\hat{A}')=G_j(\hat{A})$.
Hence $P(\hat{A}')=P(\hat{A})$ for products, and by linearity for all elements of $\mathcal{A}(\mathbb{R}^3)$.
The configurations $(A+r_q\cdot m',\,\mathfrak{z}-dm')$ and $(A,\mathfrak{z})$ have the same curvature and the same phases, hence the same lifts.
For a fixed cochain $\hat{A}$ the map $P\mapsto\mathrm{qu}_n\left(P(\hat{A})\right)$ is a homomorphism of algebras $\mathrm{Fn}\to\mathrm{Gr}_n$: the evaluation at $\hat{A}$ is a homomorphism for the pointwise operations of $\mathrm{Fn}$ defined in 1.3, and $\mathrm{qu}_n$ is a homomorphism by its definition in 1.5.
So, for a lift $\hat{A}$ of $(A,\mathfrak{z})$,
\[
\mathrm{qu}_n\left((P_1\cdots P_r)(\hat{A})\right)=\mathrm{qu}_n\left(P_1(\hat{A})\right)\cdots\mathrm{qu}_n\left(P_r(\hat{A})\right).
\]
One lift serves for all factors, and by the first assertion of (b) each of these $r+1$ expressions is the function (1.3) of the element of $\mathcal{A}(C(n))$ in question.

(c) It suffices to treat a product $P=\Phi\cdot G_1\cdots G_s$ that is read on $C(n)$.
For a sum of such products one takes the largest $\mathrm{Cq}$, multiplied by the number of terms, and the least $\varepsilon$; since $\|F^{\mathrm{u}}\|^2\ge0$, the bound is non-increasing in $\varepsilon$.
The map $X(\mathrm{Bx}_j)\to\mathbb{L}^{\mathrm{f}}_0/\Gamma(n)$ is injective, and two different plaquettes of $\bigcup_j\mathrm{Pq}(\mathrm{Bx}_j)$ are inequivalent modulo $\Gamma(n)$, because the sets $\mathrm{Bx}_j+v$ are pairwise disjoint.
Let $\varepsilon:=\min\left(\frac12,\varepsilon_1,\dots,\varepsilon_s\right)$, and $\varepsilon:=\frac12$ if $s=0$.
Then
\[
\sum_j\left(\frac12-\varepsilon_j\right)\sum_{p\in\mathrm{Pq}(\mathrm{Bx}_j)}\tilde{F}_p^2\le\left(\frac12-\varepsilon\right)\sum_j\sum_{p\in\mathrm{Pq}(\mathrm{Bx}_j)}\tilde{F}_p^2\le\left(\frac12-\varepsilon\right)\|F^{\mathrm{u}}\|^2;
\]
the second inequality enlarges the set of the plaquettes from those of the boxes to all plaquettes of $C(n)$, and since $\frac12-\varepsilon\ge0$ the bound is monotone in the set of the plaquettes.
A coefficient of a product of elements is a finite sum of products of coefficients, and $\sup|\Phi|<\infty$.
As to measurability: for fixed $\mathfrak{z}$ the integer cochain $m$ of (a) is fixed, the map $A\mapsto\hat{A}=\tilde{A}+r_q\cdot m$ is affine, and the $G_j$ and $\Phi$ are Borel.

(d) The element $\exp\left(-S_n(A)\right)$ is a polynomial in the $384|n|$ independent generators whose coefficients are polynomials in $M_{\mathrm{ch}}$, in the entries of the matrices $M_{i}(x)$ and in the phases $\exp\left(\pm iqA(b)\right)$; these coefficients are bounded uniformly in $A$.
A coefficient of the product $\exp\left(-S_n(A)\right)\cdot P(A,\mathfrak{z})$ is a sum of at most $2^{384|n|}$ products of two coefficients.
So it is bounded by $C_n\cdot\mathrm{Cq}\cdot\exp\left(\left(\frac12-\varepsilon\right)\|F^{\mathrm{u}}\|^2\right)$, and the integrand is bounded by
\[
C_n\cdot\mathrm{Cq}\cdot\prod_p\exp\left(-\varepsilon\left(F^{\mathrm{u}}_p\right)^2\right),
\]
the product extending over the $3|n|$ plaquettes $p$ of $C(n)$.
Dropping the condition $d\mathfrak{z}=0$ one gets, for every $A$,
\[
\sum_{\mathfrak{z}\in C^{2}(n;\mathbb{Z})}\prod_p\exp\left(-\varepsilon\left((dA)_p+r_q\cdot\mathfrak{z}_p\right)^2\right)\le\boldsymbol{\Theta}_{\varepsilon}^{3|n|},
\]
and the box $[0,r_q)^{\mathrm{links}}$ has the volume $r_q^{3|n|}$.
For the bound on $\boldsymbol{\Theta}_{\varepsilon}$: the two points of $y+r_q\mathbb{Z}$ next to $0$ contribute at most $2$, and every other term is at most $r_q^{-1}$ times the integral of $\exp\left(-\varepsilon t^2\right)$ over the adjacent interval of length $r_q$ on the side of $0$; these intervals are disjoint, and $\int_{\mathbb{R}}\mathrm{d}t\,\exp\left(-\varepsilon t^2\right)=\left(\frac{\pi}{\varepsilon}\right)^{\frac12}$.
Of $P(A,\mathfrak{z})$ the argument has used only that it is a Borel map with values in $\mathrm{Gr}_n$ whose coefficients obey the bound of (c); this proves the last assertion of (d).
The linearity of $Z_n[P]$ in $P$ follows from (b).
\end{proof}

\subsection*{3. One direction}

\subsubsection*{3.1. Layers, the dihedral group, copies, the fold}

Let $\mu$ be a direction, and let $o_\mu\in\frac12\mathbb{Z}$ and $\mathrm{hw}\in\frac12\mathbb{Z}$ with $\mathrm{hw}>0$.
Put
\[
\mathrm{po}_j:=\mathrm{po}^{\mu}_j:=o_\mu+j\cdot\mathrm{hw}\quad(j\in\mathbb{Z}),\qquad\tau_\mu:=2\mathrm{hw}\cdot e_\mu,
\]
and call the open slab $\mathrm{Sl}_\mu(\mathrm{po}_j,\mathrm{po}_{j+1})$ the \emph{layer}~$j$.

(w1) For every $j$ the number $\mathrm{po}_j$ lies in $\frac12\mathbb{Z}=\operatorname{Adm}(y_\mu)$; so every plane $(y_\mu,\mathrm{po}_j)$ is a mirror.
If $2\mathrm{hw}$ is odd, then $\mathrm{po}_j$ and $\mathrm{po}_{j+1}$ are of different kinds (one of them lies in $\mathbb{Z}$, the other one in $\mathbb{Z}+\frac12$), while $\mathrm{po}_j$ and $\mathrm{po}_{j+2}$ are of the same kind.

(w2) Since $2\mathrm{hw}$ is an integer, $\tau_\mu$ is a vector of $\mathbb{Z}^3$.
By (M3),
\[
\theta_{\mu,\mathrm{po}_{j+s}}\,\theta_{\mu,\mathrm{po}_j}=\tau_{s\cdot\tau_\mu},\qquad\text{that is,}\qquad\theta_{\mu,\mathrm{po}_{j+s}}=\tau_{s\cdot\tau_\mu}\,\theta_{\mu,\mathrm{po}_j}\qquad(j,s\in\mathbb{Z}).
\]

Let $\mathrm{Dw}_\mu$ be the group generated by the reflections $\theta_{\mu,\mathrm{po}_j}$, $j\in\mathbb{Z}$.
By (w2), a word of even length in these reflections is a translation $\tau_{s\cdot\tau_\mu}$, and a word of odd length is of the form $\tau_{s\cdot\tau_\mu}\,\theta_{\mu,\mathrm{po}_j}=\theta_{\mu,\mathrm{po}_{j+s}}$.
Hence
\[
\mathrm{Dw}_\mu=\left\{\tau_{s\cdot\tau_\mu}:\ s\in\mathbb{Z}\right\}\cup\left\{\theta_{\mu,\mathrm{po}_j}:\ j\in\mathbb{Z}\right\}.
\]
By (M2), the reflection $\theta_{\mu,\mathrm{po}_m}$ maps the layer $j$, that is, the set $\left\{\mathrm{po}_j<y_\mu<\mathrm{po}_{j+1}\right\}$, onto
\[
\left\{2\mathrm{po}_m-\mathrm{po}_{j+1}<y_\mu<2\mathrm{po}_m-\mathrm{po}_j\right\}=\left\{\mathrm{po}_{2m-j-1}<y_\mu<\mathrm{po}_{2m-j}\right\},
\]
which is the layer $2m-j-1$; and $\tau_{s\cdot\tau_\mu}$ maps the layer $j$ onto the layer $j+2s$.
Consequently, for $i,j\in\mathbb{Z}$ there is exactly one element $\mathrm{dw}_{i\leftarrow j}=\mathrm{dw}^{\mu}_{i\leftarrow j}$ of $\mathrm{Dw}_\mu$ that maps the layer $j$ onto the layer $i$, namely
\begin{equation*}
\mathrm{dw}_{i\leftarrow j}=
\begin{cases}
\tau_{\frac{i-j}{2}\cdot\tau_\mu}&\text{if }i-j\text{ is even},\\
\theta_{\mu,\mathrm{po}_{(i+j+1)/2}}&\text{if }i-j\text{ is odd}.
\end{cases}
\tag{3.0}
\end{equation*}
In particular $\mathrm{dw}_{j\leftarrow j}=\mathrm{id}$, and $\mathrm{dw}_{j+1\leftarrow j}=\theta_{\mu,\mathrm{po}_{j+1}}$ is the reflection in the upper wall of the layer $j$.
By the uniqueness,
\begin{equation*}
\mathrm{dw}_{i'\leftarrow i}\,\mathrm{dw}_{i\leftarrow j}=\mathrm{dw}_{i'\leftarrow j},\qquad\tau_{s\cdot\tau_\mu}\,\mathrm{dw}_{i\leftarrow j}=\mathrm{dw}_{i+2s\leftarrow j}.
\tag{$3.0'$}
\end{equation*}

Let the two sides $n_\nu\ge4$, $\nu\ne\mu$, be given.
For $G\in\mathcal{A}_{\mu,n}[\mathrm{po}_j,\mathrm{po}_{j+1}]$ the \emph{copy} of $G$ in the layer $i$ is
\[
G^{(i)}:=G^{\mathrm{dw}_{i\leftarrow j}}.
\]
It lies in $\mathcal{A}_{\mu,n}[\mathrm{po}_i,\mathrm{po}_{i+1}]$.
Indeed, by Lemma~2.2\,(B2), the image of a product of $\mathcal{A}\left(\mathrm{Sl}_\mu(\mathrm{po}_j,\mathrm{po}_{j+1})\right)$ that embeds in $\Gamma_\mu(n)$ is a product of $\mathcal{A}\left(\mathrm{Sl}_\mu(\mathrm{po}_i,\mathrm{po}_{i+1})\right)$ that embeds in $r\Gamma_\mu(n)$, where $r\in\{1,r_\mu\}$ is the linear part of $\mathrm{dw}_{i\leftarrow j}$; and $r_\mu$ fixes $\Gamma_\mu(n)$ pointwise.
By (S6), (3.0) and ($3.0'$),
\begin{equation*}
G^{(j)}=G,\qquad G^{(j+1)}=\Theta_{\mu,\mathrm{po}_{j+1}}G,\qquad G^{(i+2s)}=\left(G^{(i)}\right)^{\tau_{s\cdot\tau_\mu}}.
\tag{$3.0''$}
\end{equation*}

Let now $n_\mu\ge4$ be an integer with
\begin{equation*}
n_\mu=2k\cdot\mathrm{hw},\qquad k\ge1\text{ an integer}.
\tag{w3}
\end{equation*}
For $G\in\mathcal{A}_{\mu,n}[\mathrm{po}_j,\mathrm{po}_{j+1}]$ put
\[
\mathrm{fold}_\mu G:=\prod_{i=j}^{j+2k-1}G^{(i)}.
\]
By (SO3), applied to the chain $\mathrm{po}_j<\mathrm{po}_{j+1}<\dots<\mathrm{po}_{j+2k}=\mathrm{po}_j+n_\mu$ with $2k\ge2$ slabs, this element lies in $\mathcal{A}(C(n))$.
So does every single copy $G^{(i)}$, by (SO3) for the chain $\mathrm{po}_i<\mathrm{po}_{i+1}<\mathrm{po}_i+n_\mu$ and the elements $G^{(i)}$ and $1$.
The vector $k\cdot\tau_\mu=n_\mu\cdot e_\mu$ lies in $\Gamma(n)$.
Therefore, by ($3.0''$) and Lemma~2.2\,(B3), $G^{(i+2k)}$ and $G^{(i)}$ define the same function of the configuration on $C(n)$.
By Lemma~2.4\,(b), with $r=2k$ factors, the function
\[
(A,\mathfrak{z})\mapsto\prod_iG^{(i)}(A,\mathfrak{z}),
\]
the product extending over $2k$ consecutive layers $i$, is the same for every choice of these $2k$ consecutive layers; for the layers $j,\dots,j+2k-1$ it is $(\mathrm{fold}_\mu G)(A,\mathfrak{z})$.

\begin{named}{Remark}
The constructions of 3.1 correspond to those of \cite[Definition~6.7]{QEDII} for $\ell=y_\mu$, $c_0=o_\mu$ and $h'=\mathrm{hw}$.
The hypotheses of that definition are satisfied: the positions $c_0+j\cdot h'=\mathrm{po}_j$ lie in $\operatorname{Adm}(y_\mu)$ by (w1), and by Section~1.1 and (w3) the torus $C(n)$ is adapted to $y_\mu$, with the cross-section $\mathrm{Xs}_\mu(n)$ and the height period $n_\mu=2k\cdot h'$.
For these data the vector $v_{h'}$ of \cite[Definition~6.7]{QEDII} is $\tau_\mu$, and the group $W(\ell;c_0,h')$ is $\mathrm{Dw}_\mu$.
The numbering of the slabs, however, is shifted by one.
In \cite[Definition~6.7]{QEDII} the slab~$j$ is the set $\left\{\mathrm{po}_{j-1}<y_\mu<\mathrm{po}_j\right\}$, so that the layer~$j$ of 3.1 is the slab~$j+1$ there.
Accordingly, $\mathrm{dw}_{i\leftarrow j}$ is the element of $W(\ell;c_0,h')$ that maps the slab~$j+1$ onto the slab~$i+1$; and for $G\in\mathcal{A}_{\mu,n}[\mathrm{po}_j,\mathrm{po}_{j+1}]$, which is the class attached there to the slab~$j+1$, the copy $G^{(i)}$ of 3.1 is the copy of $G$ in the slab~$i+1$ in the sense of \cite[Definition~6.7]{QEDII}.
The two folds are formed over different sets of $2k$ slabs: $\mathrm{fold}_\mu G$ is the product of the copies of $G$ in the slabs $j+1,\dots,j+2k$, whereas $\mathrm{fold}_\ell G$ of \cite[Definition~6.7]{QEDII} is the product of the copies of $G$ in the slabs $1,\dots,2k$.
For $j=0$ the two elements coincide; for every $j$ they define the same function of the configuration on $C(n)$, as shown at the end of 3.1.
\end{named}

\subsubsection*{3.2. H\"older's inequality for $2k$ layers}

\begin{named}{Proposition 3.1}[{cf.\ \cite[Proposition~6.8]{QEDII}, with both tori equal to $C(n)$ and $\ell=y_\mu$}]
Let $\mu$, $o_\mu$, $\mathrm{hw}$ and $\mathrm{po}_j$ be as in 3.1, and let $n=(n_1,n_2,n_3)$ with $n_\nu\ge4$ for all $\nu$ and $n_\mu=2k\cdot\mathrm{hw}$, $k\ge1$ an integer.
Let $\mathrm{Hs}$, $V$, $\mathrm{Cs}$ and $\mathrm{kk}$ be the objects of Theorem~SO for $(\mu;n_\nu,\nu\ne\mu)$.
For every $j\in\mathbb{Z}$ and every $G\in\mathcal{A}_{\mu,n}[\mathrm{po}_j,\mathrm{po}_{j+1}]$,
\begin{equation}
Z_n\left[\mathrm{fold}_\mu G\right]=\mathrm{kk}^{|n|}\cdot\left\|\mathrm{Cs}[G]\right\|_{2k}^{2k}=\mathrm{kk}^{|n|}\cdot\operatorname{Tr}_{\mathrm{Hs}(\mathrm{po}_j)}\left(\mathrm{Cs}[G]^{*}\cdot\mathrm{Cs}[G]\right)^k\ge0.
\tag{3.1}
\end{equation}
For every $j_0\in\mathbb{Z}$ and all $G_j\in\mathcal{A}_{\mu,n}[\mathrm{po}_j,\mathrm{po}_{j+1}]$, $j=j_0,\dots,j_0+2k-1$,
\begin{equation}
\left|Z_n\left[G_{j_0}\cdot G_{j_0+1}\cdots G_{j_0+2k-1}\right]\right|\le\prod_{j=j_0}^{j_0+2k-1}Z_n\left[\mathrm{fold}_\mu G_j\right]^{\frac{1}{2k}}.
\tag{3.2}
\end{equation}
\end{named}

\begin{proof}
\emph{Step 1: the fold is a Schatten norm.}
Fix $j$ and $G$, and put $Y:=\mathrm{Cs}[G]\colon\mathrm{Hs}(\mathrm{po}_j)\to\mathrm{Hs}(\mathrm{po}_{j+1})$ and $\tau:=\tau_\mu$.
By (SO2b) with $\mathrm{po}=\mathrm{po}_j$ and $\mathrm{po}'=\mathrm{po}_{j+1}$, for which $2(\mathrm{po}'-\mathrm{po})\cdot e_\mu=\tau$,
\[
\mathrm{Cs}\left[\Theta_{\mu,\mathrm{po}_{j+1}}G\right]=V_\tau\cdot Y^{*}\colon\mathrm{Hs}(\mathrm{po}_{j+1})\to\mathrm{Hs}(\mathrm{po}_{j+2}).
\]
Here $Y^{*}$ maps $\mathrm{Hs}(\mathrm{po}_{j+1})$ to $\mathrm{Hs}(\mathrm{po}_j)$, and $V_\tau$ maps $\mathrm{Hs}(\mathrm{po}_j)$ to $\mathrm{Hs}(\mathrm{po}_j+2\mathrm{hw})=\mathrm{Hs}(\mathrm{po}_{j+2})$.
If $2\mathrm{hw}$ is odd, then $\mathrm{Hs}(\mathrm{po}_j)$ and $\mathrm{Hs}(\mathrm{po}_{j+1})$ are the spaces of planes of different kinds, and $\mathrm{Hs}(\mathrm{po}_j)$ and $\mathrm{Hs}(\mathrm{po}_{j+2})$ those of planes of the same kind.
By ($3.0''$), the copies of $G$ in the layers $j+2s$ and $j+2s+1$ are $G^{\tau_{s\cdot\tau}}$ and $\left(\Theta_{\mu,\mathrm{po}_{j+1}}G\right)^{\tau_{s\cdot\tau}}$.
By (SO2a) and $V_v\cdot V_{v'}=V_{v+v'}$,
\begin{align*}
\mathrm{Cs}\left[G^{(j+2s)}\right]&=V_{s\cdot\tau}\cdot Y\cdot V_{s\cdot\tau}^{-1},\\
\mathrm{Cs}\left[G^{(j+2s+1)}\right]&=V_{s\cdot\tau}\cdot V_\tau\cdot Y^{*}\cdot V_{s\cdot\tau}^{-1}=V_{(s+1)\cdot\tau}\cdot Y^{*}\cdot V_{s\cdot\tau}^{-1},
\end{align*}
and therefore
\[
\mathrm{Cs}\left[G^{(j+2s+1)}\right]\cdot\mathrm{Cs}\left[G^{(j+2s)}\right]=V_{(s+1)\cdot\tau}\cdot Y^{*}\cdot Y\cdot V_{s\cdot\tau}^{-1}.
\]
Now apply (SO3) with the base plane $\mathrm{po}_j$, the chain $\mathrm{po}_j<\mathrm{po}_{j+1}<\dots<\mathrm{po}_{j+2k}=\mathrm{po}_j+n_\mu$ and the elements $G^{(j)},\dots,G^{(j+2k-1)}$; note that $n_\mu\cdot e_\mu=k\cdot\tau$.
With the later layers written to the left,
\begin{align*}
Z_n\left[\mathrm{fold}_\mu G\right]&=\mathrm{kk}^{|n|}\cdot\operatorname{Tr}_{\mathrm{Hs}(\mathrm{po}_j)}\left(V_{k\cdot\tau}^{-1}\cdot\prod_{s=k-1,\dots,0}\left[V_{(s+1)\cdot\tau}\cdot Y^{*}\cdot Y\cdot V_{s\cdot\tau}^{-1}\right]\right)\\
&=\mathrm{kk}^{|n|}\cdot\operatorname{Tr}\left(V_{k\cdot\tau}^{-1}\cdot V_{k\cdot\tau}\cdot\left(Y^{*}\cdot Y\right)^k\right)=\mathrm{kk}^{|n|}\cdot\operatorname{Tr}_{\mathrm{Hs}(\mathrm{po}_j)}\left(Y^{*}\cdot Y\right)^k;
\end{align*}
between consecutive brackets $V_{s\cdot\tau}^{-1}\cdot V_{s\cdot\tau}=1$, and at the right end $V_0=1$.
The operator $\left(Y^{*}\cdot Y\right)^k\ge0$ is of trace class.
This proves (3.1).

\emph{Step 2: H\"older's inequality.}
By (SO3) with the base plane $\mathrm{po}_{j_0}$,
\[
Z_n\left[G_{j_0}\cdots G_{j_0+2k-1}\right]=\mathrm{kk}^{|n|}\cdot\operatorname{Tr}_{\mathrm{Hs}(\mathrm{po}_{j_0})}\left(U\cdot Y_{j_0+2k-1}\cdots Y_{j_0}\right),
\]
where $Y_j:=\mathrm{Cs}[G_j]\colon\mathrm{Hs}(\mathrm{po}_j)\to\mathrm{Hs}(\mathrm{po}_{j+1})$ is a Hilbert--Schmidt operator, and $U:=V_{n_\mu\cdot e_\mu}^{-1}\colon\mathrm{Hs}(\mathrm{po}_{j_0}+n_\mu)\to\mathrm{Hs}(\mathrm{po}_{j_0})$ is unitary.
By Lemma~2.1 and (3.1),
\[
\left|Z_n\left[G_{j_0}\cdots G_{j_0+2k-1}\right]\right|\le\mathrm{kk}^{|n|}\cdot\prod_j\|Y_j\|_{2k}=\prod_j\left(\mathrm{kk}^{|n|}\cdot\|Y_j\|_{2k}^{2k}\right)^{\frac{1}{2k}}=\prod_jZ_n\left[\mathrm{fold}_\mu G_j\right]^{\frac{1}{2k}}.
\]
\end{proof}

For $k=1$, Lemma~2.1 is the Cauchy--Schwarz inequality for the Hilbert--Schmidt scalar product.

\subsection*{4. The cells: proof of Proposition~39\,(a)}

Let $\mathcal{Q}=(\mathfrak{F}_{\mathrm{c}};o_1,o_2,o_3;\mathrm{hw})$ be as in Definition~39.0.
For $\mu=1,2,3$ let $\mathrm{po}^\mu_j$, $\tau_\mu$, $\mathrm{Dw}_\mu$ and $\mathrm{dw}^\mu_{i\leftarrow j}$ be the objects of Section~3.1 formed for the direction $\mu$.
For $\mathfrak{c}\in\mathbb{Z}^3$ the open cell $Q^{\circ}_{\mathfrak{c}}(\mathcal{Q})$ is the intersection of the three layers $\mathfrak{c}_\mu$ of the directions $\mu=1,2,3$.
For a direction $\mu$ and $i\in\mathbb{Z}$, $\mathfrak{c}[\mu{:}i]$ denotes the point of $\mathbb{Z}^3$ that arises from $\mathfrak{c}$ when its entry $\mu$ is replaced by $i$.

\begin{proof}[Proof of Proposition~39\,(a)]
\emph{Walls.}
The numbers $o_\mu$ and $\mathrm{hw}$ lie in $\tfrac12\mathbb{Z}$;
hence $o_\mu+j\cdot\mathrm{hw}\in\tfrac12\mathbb{Z}=\operatorname{Adm}(y_\mu)$ for every $j\in\mathbb{Z}$, and this is (w1).
For $\mathrm{po}\in\tfrac12\mathbb{Z}$, the identity (M3) of Section~1.2 gives $\theta_{y_\mu,\mathrm{po}+\mathrm{hw}}\,\theta_{y_\mu,\mathrm{po}}=\tau_{2\mathrm{hw}\cdot e_\mu}$,
and $2\mathrm{hw}$ is an integer; this is (w2).

\emph{The group.}
By (w2),
\[
\theta_{\mu,\mathrm{po}_j}=\tau_{j\tau_\mu}\,\theta_{\mu,\mathrm{po}_0}=\left(\theta_{\mu,\mathrm{po}_1}\,\theta_{\mu,\mathrm{po}_0}\right)^j\,\theta_{\mu,\mathrm{po}_0}\qquad(j\in\mathbb{Z}),
\]
so that $\mathrm{Dw}_\mu$ is generated by the two involutions $\theta_{\mu,\mathrm{po}_0}$ and $\theta_{\mu,\mathrm{po}_1}$.
Their product $\tau_{\tau_\mu}$ has infinite order, since $\mathrm{hw}>0$.
Therefore $\mathrm{Dw}_\mu$ is an infinite dihedral group.
An element of $\mathrm{Dw}_\mu$ changes the coordinate $\mu$ only, through a map of $y_\mu$ alone;
consequently elements of $\mathrm{Dw}_\mu$ and of $\mathrm{Dw}_\nu$ commute when $\mu\neq\nu$
(for two reflections this is (M5), and translations commute with each other).
It follows that
\[
\mathrm{Dw}_1\times\mathrm{Dw}_2\times\mathrm{Dw}_3\to W(\mathcal{Q}),\qquad(d_1,d_2,d_3)\mapsto d_1\,d_2\,d_3,
\]
is a homomorphism onto $W(\mathcal{Q})$, the group generated by all the reflections $\theta_{\mu,\mathrm{po}^\mu_j}$.
It is injective.
Indeed, $d_1\,d_2\,d_3$ acts on the coordinate $\mu$ as $d_\mu$ does;
if the product is the identity, then each $d_\mu$ acts trivially on the coordinate $\mu$, and it acts trivially on the two other coordinates in any case.
Thus $W(\mathcal{Q})$ is the direct product of the three pairwise commuting infinite dihedral groups $\mathrm{Dw}_1$, $\mathrm{Dw}_2$, $\mathrm{Dw}_3$.

\emph{Simple transitivity.}
The element $d_1\,d_2\,d_3$ maps $Q^{\circ}_{\mathfrak{c}}$ onto the intersection of the layers $d_\mu(\text{layer }\mathfrak{c}_\mu)$, $\mu=1,2,3$, each layer being taken in its own direction.
By Section~3.1, for $\mathfrak{c},\mathfrak{c}'\in\mathbb{Z}^3$ exactly one triple $(d_1,d_2,d_3)$ satisfies $d_\mu(\text{layer }\mathfrak{c}_\mu)=\text{layer }\mathfrak{c}'_\mu$ for $\mu=1,2,3$,
namely the triple with $d_\mu=\mathrm{dw}^\mu_{\mathfrak{c}'_\mu\leftarrow\mathfrak{c}_\mu}$.
Hence there is exactly one element of $W(\mathcal{Q})$ that maps $Q^{\circ}_{\mathfrak{c}}$ onto $Q^{\circ}_{\mathfrak{c}'}$, namely
\begin{equation}
\begin{aligned}
&g_{\mathrm{w},\mathfrak{c}'\leftarrow\mathfrak{c}}:=\mathrm{dw}^3_{\mathfrak{c}'_3\leftarrow\mathfrak{c}_3}\,\mathrm{dw}^2_{\mathfrak{c}'_2\leftarrow\mathfrak{c}_2}\,\mathrm{dw}^1_{\mathfrak{c}'_1\leftarrow\mathfrak{c}_1}
\quad\text{(the three factors may be taken in any order)},\\
&g_{\mathrm{w},\mathfrak{c}''\leftarrow\mathfrak{c}'}\,g_{\mathrm{w},\mathfrak{c}'\leftarrow\mathfrak{c}}=g_{\mathrm{w},\mathfrak{c}''\leftarrow\mathfrak{c}},
\end{aligned}
\tag{4.1}
\end{equation}
the last identity because both sides map $Q^{\circ}_{\mathfrak{c}}$ onto $Q^{\circ}_{\mathfrak{c}''}$.
By (3.0), the element $g_{\mathrm{w},\mathfrak{c}'\leftarrow\mathfrak{c}}$ is even if and only if $\sum_\mu(\mathfrak{c}'_\mu-\mathfrak{c}_\mu)$ is even.
(For $o=0$, $\mathrm{hw}=\mathrm{hc}$ and $\mathfrak{c}=0$ the formulas (3.0) are those of $\mathsf{g}_{\mathfrak{c}'}$ in Section~I.6.4:
$y_\mu\mapsto i\cdot\mathrm{per}+y_\mu$ if $\mathfrak{c}'_\mu=2i$,
and $y_\mu\mapsto2\mathrm{po}_{i+1}-y_\mu=(i+1)\cdot\mathrm{per}-y_\mu$ if $\mathfrak{c}'_\mu=2i+1$.)

\emph{Translations.}
The translations in $\mathrm{Dw}_\mu$ are the $\tau_{s\tau_\mu}$, $s\in\mathbb{Z}$.
Hence the translations in $W(\mathcal{Q})$ are the $\tau_v$ with $v\in(2\mathrm{hw}\,\mathbb{Z})^3$.

\emph{Tori.}
Let $C(n)$ carry $\mathcal{Q}$, that is, $n_\mu\in2\mathrm{hw}\,\mathbb{Z}$ for $\mu=1,2,3$.
Then $\Gamma(n)\subset(2\mathrm{hw}\,\mathbb{Z})^3$, so that $\{\tau_v:\ v\in\Gamma(n)\}$ is a subgroup of $W(\mathcal{Q})$.
This subgroup is normal:
$g_{\mathrm{w}}\,\tau_v\,g_{\mathrm{w}}^{-1}=\tau_{rv}$, where the linear part $r$ of $g_{\mathrm{w}}$ is a diagonal matrix with entries $\pm1$, and $r\,\Gamma(n)=\Gamma(n)$.
Therefore $W(\mathcal{Q})$ permutes the $\Gamma(n)$-orbits of cells, and the quotient group acts on the set of these orbits.
The action is transitive, because $W(\mathcal{Q})$ acts transitively on the cells.
It is free:
if $g_{\mathrm{w}}Q^{\circ}_{\mathfrak{c}}=Q^{\circ}_{\mathfrak{c}}+v$ with $v\in\Gamma(n)$, then $\tau_{-v}\,g_{\mathrm{w}}$ fixes the cell $Q^{\circ}_{\mathfrak{c}}$, and so $g_{\mathrm{w}}=\tau_v$.
Since $\tau_vQ^{\circ}_{\mathfrak{c}}=Q^{\circ}_{\mathfrak{c}'}$ with $\mathfrak{c}'_\mu=\mathfrak{c}_\mu+v_\mu/\mathrm{hw}$,
the orbits correspond to the classes of $\mathbb{Z}^3$ modulo $\prod_\mu\left(\frac{n_\mu}{\mathrm{hw}}\mathbb{Z}\right)$, and their number is
\[
\mathrm{nc}(n)=\prod_{\mu=1}^{3}\frac{n_\mu}{\mathrm{hw}} .
\]
Each factor $\frac{n_\mu}{\mathrm{hw}}=2\cdot\frac{n_\mu}{2\mathrm{hw}}$ is an even integer.
\end{proof}

By (4.1) and the rule (S6) of Section~1.3 one has, for $O\in\mathcal{A}\left(Q^{\circ}_{\mathfrak{c}}\right)$,
\begin{equation}
\begin{aligned}
&O^{[\mathfrak{c}']}=O^{g_{\mathrm{w},\mathfrak{c}'\leftarrow\mathfrak{c}}}\in\mathcal{A}\left(Q^{\circ}_{\mathfrak{c}'}\right)\quad\text{(Lemma~2.2\,(B2))},\\
&\left(O^{[\mathfrak{c}']}\right)^{g_{\mathrm{w},\mathfrak{c}''\leftarrow\mathfrak{c}'}}=O^{g_{\mathrm{w},\mathfrak{c}''\leftarrow\mathfrak{c}'}\,g_{\mathrm{w},\mathfrak{c}'\leftarrow\mathfrak{c}}}=O^{[\mathfrak{c}'']}.
\end{aligned}
\tag{4.2}
\end{equation}

\subsection*{5. Families of cell contents; one step; three steps (\cite[Definition~6.13, Lemma~6.14, Proposition~6.15]{QEDII} for $\mathbb{T}_0=\mathbb{T}_1=\mathbb{T}_2=\mathbb{T}_3=C(n)$)}

Throughout this section the torus $C(n)$ carries $\mathcal{Q}$.

\subsubsection*{5.1. Definitions}

\emph{Shifts.}
For $v\in\mathbb{Z}^3$ with $v_\mu\in\mathrm{hw}\,\mathbb{Z}$ for $\mu=1,2,3$ put
\[
\mathrm{sh}(v):=\left(\frac{v_\mu}{\mathrm{hw}}\right)_{\mu=1,2,3}\in\mathbb{Z}^3 .
\]
Then $\tau_vQ^{\circ}_{\mathfrak{c}}=Q^{\circ}_{\mathfrak{c}+\mathrm{sh}(v)}$, and the map $\mathrm{sh}$ is additive and injective.
(If $2\mathrm{hw}$ is odd, then $\mathbb{Z}\cap\mathrm{hw}\,\mathbb{Z}=2\mathrm{hw}\,\mathbb{Z}$, and $\mathrm{sh}(v)$ has even entries.)
One has $\mathrm{sh}(\Gamma(n))=\prod_\mu\left(\frac{n_\mu}{\mathrm{hw}}\mathbb{Z}\right)$, and the cells of $C(n)$ are the classes of $\mathbb{Z}^3$ modulo $\mathrm{sh}(\Gamma(n))$.
The \emph{brick}
\[
\mathrm{Kb}(n):=\left\{\mathfrak{c}\in\mathbb{Z}^3:\ 0\le\mathfrak{c}_\mu<\frac{n_\mu}{\mathrm{hw}},\ \mu=1,2,3\right\}
\]
is a set of representatives of these classes,
and so is, for a direction $\mu$ and $j\in\mathbb{Z}$, the \emph{shifted brick}
\[
\mathrm{Kb}(n;\mu,j):=\left\{\mathfrak{c}\in\mathbb{Z}^3:\ j\le\mathfrak{c}_\mu<j+\frac{n_\mu}{\mathrm{hw}},\ 0\le\mathfrak{c}_\nu<\frac{n_\nu}{\mathrm{hw}}\text{ for }\nu\neq\mu\right\}.
\]

\emph{Families.}
A \emph{family on $(C(n),\mathcal{Q})$} is a collection $\mathrm{GG}=(G_{\mathfrak{c}})_{\mathfrak{c}\in\mathbb{Z}^3}$ with $G_{\mathfrak{c}}\in\mathcal{A}\left(Q^{\circ}_{\mathfrak{c}}\right)$ such that
\[
G_{\mathfrak{c}+\mathrm{sh}(v)}=\left(G_{\mathfrak{c}}\right)^{\tau_v}\qquad\text{for all }\mathfrak{c}\in\mathbb{Z}^3\text{ and }v\in\Gamma(n).
\]
For a set $\mathrm{Rs}$ of representatives of $\mathbb{Z}^3/\mathrm{sh}(\Gamma(n))$ put
\[
\mathrm{Prd}(\mathrm{GG};\mathrm{Rs}):=\prod_{\mathfrak{c}\in\mathrm{Rs}}G_{\mathfrak{c}} .
\]
\begin{itemize}
\item[(5.a)]
$\mathrm{Prd}(\mathrm{GG};\mathrm{Rs})$ lies in $\mathcal{A}(C(n))$.
Indeed, expand every factor into products;
by Lemma~2.3\,(C4) each term of the expansion is a product that is read on $C(n)$.
The order of the factors does not matter (Lemma~2.2\,(B1)).
\item[(5.b)]
The function $(A,\mathfrak{z})\mapsto\mathrm{Prd}(\mathrm{GG};\mathrm{Rs})(A,\mathfrak{z})$ does not depend on $\mathrm{Rs}$.
Indeed, each $G_{\mathfrak{c}}$ lies in $\mathcal{A}(C(n))$ (Lemma~2.3\,(C4) with $s=1$), and so does their product, by (5.a);
by Lemma~2.4\,(b) with $r=\mathrm{nc}(n)$ factors, the function is therefore $\prod_{\mathfrak{c}\in\mathrm{Rs}}G_{\mathfrak{c}}(A,\mathfrak{z})$.
If $\mathrm{Rs}'$ is a second set of representatives, then for every $\mathfrak{c}\in\mathrm{Rs}$ there is exactly one $v(\mathfrak{c})\in\Gamma(n)$ with $\mathfrak{c}+\mathrm{sh}(v(\mathfrak{c}))\in\mathrm{Rs}'$,
and $\mathfrak{c}\mapsto\mathfrak{c}+\mathrm{sh}(v(\mathfrak{c}))$ is a bijection $\mathrm{Rs}\to\mathrm{Rs}'$;
moreover $G_{\mathfrak{c}+\mathrm{sh}(v)}(A,\mathfrak{z})=\left(G_{\mathfrak{c}}\right)^{\tau_v}(A,\mathfrak{z})=G_{\mathfrak{c}}(A,\mathfrak{z})$ by Lemma~2.2\,(B3).
Hence
\[
Z_n[\mathrm{GG}]:=Z_n\left[\mathrm{Prd}(\mathrm{GG};\mathrm{Rs})\right]
\]
is well defined.
\item[(5.c)]
Contents $G_{\mathfrak{c}}\in\mathcal{A}\left(Q^{\circ}_{\mathfrak{c}}\right)$ prescribed for the points $\mathfrak{c}$ of a set $\mathrm{Rs}$ of representatives extend to exactly one family,
namely $G_{\mathfrak{c}+\mathrm{sh}(v)}:=\left(G_{\mathfrak{c}}\right)^{\tau_v}$ for $\mathfrak{c}\in\mathrm{Rs}$ and $v\in\Gamma(n)$.
Indeed, every point of $\mathbb{Z}^3$ is $\mathfrak{c}+\mathrm{sh}(v)$ with a unique pair $(\mathfrak{c},v)\in\mathrm{Rs}\times\Gamma(n)$, because $\mathrm{sh}$ is injective;
$\left(G^{\tau_v}\right)^{\tau_{v'}}=G^{\tau_{v+v'}}$ by (S6);
and $G^{\tau_v}$ lies in $\mathcal{A}\left(Q^{\circ}_{\mathfrak{c}+\mathrm{sh}(v)}\right)$ by Lemma~2.2\,(B2).
\end{itemize}

\emph{Definition~39.0 and Proposition~39\,(b), (c) in this notation.}
By Definition~39.0, the hypothesis ``$O_{\mathfrak{c}}\in\mathcal{A}\left(Q^{\circ}_{\mathfrak{c}}(\mathcal{Q})\right)$ for every cell $\mathfrak{c}$ of $C(n)$'' means that a set $\mathrm{Rs}$ of representatives is chosen
and that $O_{\mathfrak{c}}\in\mathcal{A}\left(Q^{\circ}_{\mathfrak{c}}\right)$ is given for $\mathfrak{c}\in\mathrm{Rs}$.
Then $\prod_{\mathfrak{c}}O_{\mathfrak{c}}$ is $\mathrm{Prd}(\mathrm{GG};\mathrm{Rs})$ for the family $\mathrm{GG}$ of (5.c),
and $O^{\sharp}_n$ is $\mathrm{Prd}(\mathrm{CC}(O);\mathrm{Rs})$ for the family $\mathrm{CC}(O)$ defined below (see (5.d)).
As an element of $\mathcal{A}$, $\mathrm{Prd}(\mathrm{GG};\mathrm{Rs})$ is formed with $\mathrm{Rs}$:
another set of representatives replaces factors $G_{\mathfrak{c}}$ by their translates $\left(G_{\mathfrak{c}}\right)^{\tau_v}$, $v\in\Gamma(n)$, which lie in $\mathcal{A}\left(Q^{\circ}_{\mathfrak{c}+\mathrm{sh}(v)}\right)$.
The function of the configuration on $C(n)$ defined by $\mathrm{Prd}(\mathrm{GG};\mathrm{Rs})$, and therefore its value under $Z_n$, does not depend on $\mathrm{Rs}$;
this assertion of Definition~39.0 is (5.b).

\emph{Fold of a family in the direction $\mu$ from the layer $j\in\mathbb{Z}$.}
Put
\[
\left(\mathrm{Rm}_\mu^j\mathrm{GG}\right)_{\mathfrak{c}}:=\left(G_{\mathfrak{c}[\mu{:}j]}\right)^{\mathrm{dw}^\mu_{\mathfrak{c}_\mu\leftarrow j}}\in\mathcal{A}\left(Q^{\circ}_{\mathfrak{c}}\right)\qquad(\mathfrak{c}\in\mathbb{Z}^3);
\]
the membership holds by Lemma~2.2\,(B2), since $\mathrm{dw}^\mu_{\mathfrak{c}_\mu\leftarrow j}$ maps $Q^{\circ}_{\mathfrak{c}[\mu{:}j]}$ onto $Q^{\circ}_{\mathfrak{c}}$.

\emph{Chessboard family.}
For $O\in\mathcal{A}\left(Q^{\circ}_{\mathfrak{c}}\right)$ put
\[
\mathrm{CC}(O):=\left(O^{[\mathfrak{c}']}\right)_{\mathfrak{c}'\in\mathbb{Z}^3}.
\]
\begin{itemize}
\item[(5.d)]
$\mathrm{CC}(O)$ is a family on $(C(n),\mathcal{Q})$,
one has $\mathrm{Prd}(\mathrm{CC}(O);\mathrm{Rs})=O^{\sharp}_n$,
and $\mathrm{Rm}_\mu^j\mathrm{CC}(O)=\mathrm{CC}(O)$ for all $\mu$ and $j$.
Indeed, for $v\in\Gamma(n)$ the translation $\tau_v$ lies in $W(\mathcal{Q})$ and maps $Q^{\circ}_{\mathfrak{c}'}$ onto $Q^{\circ}_{\mathfrak{c}'+\mathrm{sh}(v)}$;
so $\tau_v=g_{\mathrm{w},\mathfrak{c}'+\mathrm{sh}(v)\leftarrow\mathfrak{c}'}$, and (4.2) gives $\left(O^{[\mathfrak{c}']}\right)^{\tau_v}=O^{[\mathfrak{c}'+\mathrm{sh}(v)]}$.
Further, $\mathrm{dw}^\mu_{\mathfrak{c}'_\mu\leftarrow j}$ is the element $g_{\mathrm{w},\mathfrak{c}'\leftarrow\mathfrak{c}'[\mu{:}j]}$ of $W(\mathcal{Q})$,
and (4.2) gives $\left(O^{[\mathfrak{c}'[\mu{:}j]]}\right)^{\mathrm{dw}^\mu_{\mathfrak{c}'_\mu\leftarrow j}}=O^{[\mathfrak{c}']}$.
\end{itemize}

\subsubsection*{5.2. One step}

\begin{named}{Lemma 5.1}
Let $\mathrm{GG}$ be a family on $(C(n),\mathcal{Q})$, and let $\mu$ be a direction.
\begin{itemize}
\item[(i)] For every $j\in\mathbb{Z}$, $\mathrm{Rm}_\mu^j\mathrm{GG}$ is a family on $(C(n),\mathcal{Q})$.
\item[(ii)] For every $j\in\mathbb{Z}$ one has $Z_n\left[\mathrm{Rm}_\mu^j\mathrm{GG}\right]\ge0$, and
\[
\left|Z_n[\mathrm{GG}]\right|\le\prod_{j=0}^{n_\mu/\mathrm{hw}-1}Z_n\left[\mathrm{Rm}_\mu^j\mathrm{GG}\right]^{\frac{\mathrm{hw}}{n_\mu}}.
\]
\end{itemize}
\end{named}

\begin{proof}
(i)
One has $\Gamma(n)=\Gamma_\mu(n)+\mathbb{Z}u$ with $u:=n_\mu e_\mu$.
The set of the vectors $v\in\Gamma(n)$ with
\[
\left(\mathrm{Rm}_\mu^j\mathrm{GG}\right)_{\mathfrak{c}+\mathrm{sh}(v)}=\left(\left(\mathrm{Rm}_\mu^j\mathrm{GG}\right)_{\mathfrak{c}}\right)^{\tau_v}\qquad\text{for all }\mathfrak{c}\in\mathbb{Z}^3
\]
is a subgroup of $\Gamma(n)$, by (S6) in the form $\left(G^{\tau_v}\right)^{\tau_{v'}}=G^{\tau_{v+v'}}$.
It therefore suffices to check the elements of $\Gamma_\mu(n)$ and the vector $u$.

Let $v\in\Gamma_\mu(n)$.
The translation $\tau_v$ commutes with every element of $\mathrm{Dw}_\mu$:
with the translations, and with the reflections because $\tau_v\,\theta_{\mu,\mathrm{po}}\,\tau_v^{-1}=\theta_{\mu,\mathrm{po}+v_\mu}=\theta_{\mu,\mathrm{po}}$ by (M4).
The entry $\mu$ of $\mathrm{sh}(v)$ is $0$.
Put $d:=\mathrm{dw}^\mu_{\mathfrak{c}_\mu\leftarrow j}$.
Using (S6) in the form $\left(G^{g_{\mathrm{w}}}\right)^{g_{\mathrm{w}}'}=G^{g_{\mathrm{w}}'\,g_{\mathrm{w}}}$ one obtains
\begin{align*}
\left(\mathrm{Rm}_\mu^j\mathrm{GG}\right)_{\mathfrak{c}+\mathrm{sh}(v)}
&=\left(G_{\mathfrak{c}[\mu{:}j]+\mathrm{sh}(v)}\right)^{d}
=\left(\left(G_{\mathfrak{c}[\mu{:}j]}\right)^{\tau_v}\right)^{d}
=\left(G_{\mathfrak{c}[\mu{:}j]}\right)^{d\,\tau_v}\\
&=\left(G_{\mathfrak{c}[\mu{:}j]}\right)^{\tau_v\,d}
=\left(\left(\mathrm{Rm}_\mu^j\mathrm{GG}\right)_{\mathfrak{c}}\right)^{\tau_v}.
\end{align*}

Now consider the vector $u$.
One has $u=\frac{n_\mu}{2\mathrm{hw}}\cdot\tau_\mu$, and $\frac{n_\mu}{2\mathrm{hw}}$ is an integer;
so $\tau_u$ lies in $\mathrm{Dw}_\mu$, and
\[
\tau_u\,\mathrm{dw}^\mu_{\mathfrak{c}_\mu\leftarrow j}=\mathrm{dw}^\mu_{\mathfrak{c}_\mu+n_\mu/\mathrm{hw}\leftarrow j}
\]
by (3.0$'$).
Moreover $\mathrm{sh}(u)=\frac{n_\mu}{\mathrm{hw}}\,e_\mu$ and $\left(\mathfrak{c}+\mathrm{sh}(u)\right)[\mu{:}j]=\mathfrak{c}[\mu{:}j]$.
Hence
\[
\left(\mathrm{Rm}_\mu^j\mathrm{GG}\right)_{\mathfrak{c}+\mathrm{sh}(u)}=\left(G_{\mathfrak{c}[\mu{:}j]}\right)^{\tau_u\,d}=\left(\left(\mathrm{Rm}_\mu^j\mathrm{GG}\right)_{\mathfrak{c}}\right)^{\tau_u}.
\]

(ii)
Put $2k:=n_\mu/\mathrm{hw}$, an even integer by (a).
Fix $j_0\in\mathbb{Z}$, let $\mathrm{Rs}:=\mathrm{Kb}(n;\mu,j_0)$, and put
\[
S_i:=\prod_{\mathfrak{c}\in\mathrm{Rs},\ \mathfrak{c}_\mu=i}G_{\mathfrak{c}}\qquad(i=j_0,\dots,j_0+2k-1).
\]
The cells with $\mathfrak{c}_\mu=i$ lie in the layer $i$ of the direction $\mu$.
Expanded into products, $S_i$ lies in $\mathcal{A}_{\mu,n}[\mathrm{po}_i,\mathrm{po}_{i+1}]$ by Lemma~2.3\,(C4),
and $\mathrm{Prd}(\mathrm{GG};\mathrm{Rs})=S_{j_0}\cdots S_{j_0+2k-1}$.
By Lemma~2.2\,(B1) and the definition of $\mathrm{Rm}_\mu^i$,
\begin{align*}
\mathrm{fold}_\mu S_i
&=\prod_{i'=i}^{i+2k-1}\left(S_i\right)^{\mathrm{dw}^\mu_{i'\leftarrow i}}
=\prod_{i'=i}^{i+2k-1}\ \prod_{\mathfrak{c}\in\mathrm{Rs},\ \mathfrak{c}_\mu=i}\left(G_{\mathfrak{c}}\right)^{\mathrm{dw}^\mu_{i'\leftarrow i}}\\
&=\prod_{\mathfrak{c}'\in\mathrm{Kb}(n;\mu,i)}\left(\mathrm{Rm}_\mu^i\mathrm{GG}\right)_{\mathfrak{c}'}
=\mathrm{Prd}\left(\mathrm{Rm}_\mu^i\mathrm{GG};\mathrm{Kb}(n;\mu,i)\right),
\end{align*}
because $\mathfrak{c}'\mapsto\mathfrak{c}'[\mu{:}i]$ maps $\{\mathfrak{c}'\in\mathrm{Kb}(n;\mu,i):\ \mathfrak{c}'_\mu=i'\}$ bijectively onto $\{\mathfrak{c}\in\mathrm{Rs}:\ \mathfrak{c}_\mu=i\}$;
both sets are described by the conditions $0\le\mathfrak{c}_\nu<n_\nu/\mathrm{hw}$, $\nu\neq\mu$.
By (i) and (5.b), $Z_n\left[\mathrm{fold}_\mu S_i\right]=Z_n\left[\mathrm{Rm}_\mu^i\mathrm{GG}\right]$.
The identity (3.1) of Proposition~3.1 gives $Z_n\left[\mathrm{Rm}_\mu^i\mathrm{GG}\right]\ge0$ for $i=j_0$, hence for every integer, $j_0$ being arbitrary.
The inequality (3.2), applied with $j_0=0$, gives
\[
\left|Z_n[\mathrm{GG}]\right|=\left|Z_n\left[S_0\cdots S_{2k-1}\right]\right|\le\prod_{i=0}^{2k-1}Z_n\left[\mathrm{fold}_\mu S_i\right]^{\frac{1}{2k}},
\]
and $\frac{1}{2k}=\frac{\mathrm{hw}}{n_\mu}$.
\end{proof}

\subsubsection*{5.3. Three steps}

\begin{named}{Proposition 5.2}
For every family $\mathrm{GG}$ on $(C(n),\mathcal{Q})$,
\begin{align*}
\left|Z_n[\mathrm{GG}]\right|
&\le\prod_{j_1}Z_n\left[\mathrm{Rm}_1^{j_1}\mathrm{GG}\right]^{\frac{\mathrm{hw}}{n_1}}
\le\prod_{j_1,j_2}Z_n\left[\mathrm{Rm}_2^{j_2}\mathrm{Rm}_1^{j_1}\mathrm{GG}\right]^{\frac{\mathrm{hw}^2}{n_1n_2}}\\
&\le\prod_{\mathfrak{c}\in\mathrm{Kb}(n)}Z_n\left[\mathrm{CC}(G_{\mathfrak{c}})\right]^{\frac{1}{\mathrm{nc}(n)}},
\end{align*}
where $j_\mu$ ranges over the integers with $0\le j_\mu<n_\mu/\mathrm{hw}$.
All the factors are non-negative.
\end{named}

\begin{proof}
The first inequality is Lemma~5.1\,(ii) for $\mathrm{GG}$ and $\mu=1$.
Each $\mathrm{Rm}_1^{j_1}\mathrm{GG}$ is a family (Lemma~5.1\,(i)), and its value under $Z_n$ is a non-negative number, equal to its modulus.
Lemma~5.1\,(ii) for $\mathrm{Rm}_1^{j_1}\mathrm{GG}$ and $\mu=2$ bounds this number by $\prod_{j_2}Z_n\left[\mathrm{Rm}_2^{j_2}\mathrm{Rm}_1^{j_1}\mathrm{GG}\right]^{\frac{\mathrm{hw}}{n_2}}$,
and $t\mapsto t^{\frac{\mathrm{hw}}{n_1}}$ is non-decreasing on $[0,\infty)$; this gives the second inequality.
Next, $\mathrm{Rm}_2^{j_2}\mathrm{Rm}_1^{j_1}\mathrm{GG}$ is a family, its value under $Z_n$ is $\ge0$,
and Lemma~5.1\,(ii) for this family and $\mu=3$ bounds it by $\prod_{j_3}Z_n\left[\mathrm{Rm}_3^{j_3}\mathrm{Rm}_2^{j_2}\mathrm{Rm}_1^{j_1}\mathrm{GG}\right]^{\frac{\mathrm{hw}}{n_3}}$;
the map $t\mapsto t^{\frac{\mathrm{hw}^2}{n_1n_2}}$ is non-decreasing on $[0,\infty)$.
By the definition of $\mathrm{Rm}$, applied three times, and by (S6), (4.1), (4.2),
\begin{align*}
\left(\mathrm{Rm}_3^{j_3}\mathrm{Rm}_2^{j_2}\mathrm{Rm}_1^{j_1}\mathrm{GG}\right)_{\mathfrak{c}}
&=\left(\left(\mathrm{Rm}_2^{j_2}\mathrm{Rm}_1^{j_1}\mathrm{GG}\right)_{\mathfrak{c}[3{:}j_3]}\right)^{\mathrm{dw}^3_{\mathfrak{c}_3\leftarrow j_3}}\\
&=\left(\left(\left(\mathrm{Rm}_1^{j_1}\mathrm{GG}\right)_{\mathfrak{c}[3{:}j_3][2{:}j_2]}\right)^{\mathrm{dw}^2_{\mathfrak{c}_2\leftarrow j_2}}\right)^{\mathrm{dw}^3_{\mathfrak{c}_3\leftarrow j_3}}\\
&=\left(\left(\left(G_{(j_1,j_2,j_3)}\right)^{\mathrm{dw}^1_{\mathfrak{c}_1\leftarrow j_1}}\right)^{\mathrm{dw}^2_{\mathfrak{c}_2\leftarrow j_2}}\right)^{\mathrm{dw}^3_{\mathfrak{c}_3\leftarrow j_3}}\\
&=\left(G_{(j_1,j_2,j_3)}\right)^{\mathrm{dw}^3_{\mathfrak{c}_3\leftarrow j_3}\,\mathrm{dw}^2_{\mathfrak{c}_2\leftarrow j_2}\,\mathrm{dw}^1_{\mathfrak{c}_1\leftarrow j_1}}
=\left(G_{(j_1,j_2,j_3)}\right)^{[\mathfrak{c}]}.
\end{align*}
Hence $\mathrm{Rm}_3^{j_3}\mathrm{Rm}_2^{j_2}\mathrm{Rm}_1^{j_1}\mathrm{GG}=\mathrm{CC}\left(G_{(j_1,j_2,j_3)}\right)$,
the triple $(j_1,j_2,j_3)$ ranges over $\mathrm{Kb}(n)$, and the exponent is
\[
\frac{\mathrm{hw}}{n_1}\cdot\frac{\mathrm{hw}}{n_2}\cdot\frac{\mathrm{hw}}{n_3}=\frac{1}{\mathrm{nc}(n)} .
\]
\end{proof}

\subsection*{6. Proofs of Proposition~39\,(b)--(f)}

\begin{proof}[Proof of Proposition~39\,(b)]
Let $C(n)$ carry $\mathcal{Q}$, let $\mathrm{Rs}$ be a set of representatives of the cells of $C(n)$,
and let $O_{\mathfrak{c}}\in\mathcal{A}\left(Q^{\circ}_{\mathfrak{c}}\right)$ be given for $\mathfrak{c}\in\mathrm{Rs}$ (Definition~39.0; Section~5.1).
Each $O_{\mathfrak{c}}$ is a finite sum of products whose boxes, and whose plaquettes read by the bounded factor, lie in $Q^{\circ}_{\mathfrak{c}}$.
Hence $\prod_{\mathfrak{c}\in\mathrm{Rs}}O_{\mathfrak{c}}$ is a finite sum of terms $\prod_{\mathfrak{c}\in\mathrm{Rs}}P_{\mathfrak{c}}$, where each $P_{\mathfrak{c}}$ is such a product for the cell $Q^{\circ}_{\mathfrak{c}}$.
By Lemma~2.3\,(C4) each term is a product,
its boxes together with all their translates by $\Gamma(n)=\sum_\mu n_\mu\mathbb{Z}e_\mu$ are pairwise disjoint,
and each of its boxes embeds in $\Gamma(n)$.
So each term is read on $C(n)$ in the sense of \cite[Definition~2.13]{QEDII}, and $\prod_{\mathfrak{c}}O_{\mathfrak{c}}$ lies in $\mathcal{A}(C(n))$.
By Lemma~2.4\,(b)--(d), the function (1.3) depends neither on the lift nor on the way in which the element is written as a sum of products,
and $Z_n\left[\prod_{\mathfrak{c}}O_{\mathfrak{c}}\right]$ is given by a sum and an integral that converge absolutely, coefficient by coefficient.

\emph{One $O$ on all tori.}
An element $O$ of $\mathcal{A}\left(Q^{\circ}_0(\mathcal{Q})\right)$ is a map $C^1\to\mathrm{Gr}$ (Section~1.4), and no torus enters its definition.
For every torus $C(n)$ that carries $\mathcal{Q}$ it lies in $\mathcal{A}(C(n))$ (Lemma~2.3\,(C4) with $s=1$),
and its function on the configurations of $C(n)$ is $(A,\mathfrak{z})\mapsto\mathrm{qu}_n\left(O(\hat A)\right)$, formed with this one map.
The same holds for the copies $O^{[\mathfrak{c}]}=O^{g_{\mathrm{w},\mathfrak{c}\leftarrow0}}$, which are elements of $\mathcal{A}\left(Q^{\circ}_{\mathfrak{c}}\right)$ formed with $W(\mathcal{Q})$ alone.
If $O$ is a function of the curvature, that is, if $O(\hat A)$ is a function of the numbers $(d\hat A)_p$ with $p$ in a finite set of plaquettes contained in $Q^{\circ}_0$,
then its function on $C(n)$ is $F^{\mathrm{u}}\mapsto O\left((\tilde F_p)_p\right)$, where $\tilde F$ is the periodic extension of $F^{\mathrm{u}}$.
Here the plaquettes are plaquettes of $\mathbb{Z}^3$,
and two different plaquettes contained in $Q^{\circ}_0$ are inequivalent modulo $\Gamma(n)$, since $\tau_vQ^{\circ}_0$ is disjoint from $Q^{\circ}_0$ for $0\neq v\in\Gamma(n)$.
The constants $\mathrm{Cq}$ and $\varepsilon$ of Lemma~2.4\,(c) are those of the products that make up $O$; they do not depend on $n$.
To see this, let $P=\Phi\,G_1\cdots G_s$ be one of these products, and let $\mathrm{Cq}_j$, $\varepsilon_j$ be the constants $\mathrm{Cq}$, $\varepsilon_G$ of $G_j$ (Section~1.4).
The sets $X(\mathrm{Bx}_j)$ are pairwise disjoint;
so the coefficients of $P(\hat A)$ in $\mathrm{Gr}$ are the numbers $\pm\Phi(d\hat A)\cdot\prod_jc_{\mathrm{mm}_j}(\hat A)$,
where $c_{\mathrm{mm}_j}$ is the coefficient of $G_j$ at a monomial $\mathrm{mm}_j$ in the generators at the sites of $X(\mathrm{Bx}_j)$.
For every $n$ such that $C(n)$ carries $\mathcal{Q}$, the sets $\mathrm{Bx}_j+v$ (all $j$, all $v\in\Gamma(n)$) are pairwise disjoint (Lemma~2.3\,(C4) for the single cell $Q^{\circ}_0$),
and therefore the map $\bigcup_jX(\mathrm{Bx}_j)\to\mathbb{L}^{\mathrm{f}}_0/\Gamma(n)$ is injective.
Consequently $\mathrm{qu}_n$ maps the monomials in the generators at the sites of $\bigcup_jX(\mathrm{Bx}_j)$, up to signs, to pairwise different monomials of a basis of $\mathrm{Gr}_n$,
and the coefficients of $P(A,\mathfrak{z})=\mathrm{qu}_n\left(P(\hat A)\right)$ are, up to signs, those of $P(\hat A)$.
By the display in the proof of Lemma~2.4\,(c), their modulus is at most $\mathrm{Cq}\cdot\exp\left(\left(\tfrac12-\varepsilon\right)\|F^{\mathrm{u}}\|^2\right)$ with
\[
\mathrm{Cq}=\sup|\Phi|\cdot\prod_j\mathrm{Cq}_j,\qquad\varepsilon=\min\left(\tfrac12,\varepsilon_1,\dots,\varepsilon_s\right).
\]
For the sum $O$ of the products one takes the largest $\mathrm{Cq}$ times the number of the products, and the least $\varepsilon$.
\end{proof}

\begin{proof}[Proof of Proposition~39\,(c)]
Let $\mathrm{GG}$ be the family of (5.c) with $G_{\mathfrak{c}}=O_{\mathfrak{c}}$ for $\mathfrak{c}\in\mathrm{Rs}$;
thus $Z_n\left[\prod_{\mathfrak{c}}O_{\mathfrak{c}}\right]=Z_n[\mathrm{GG}]$.

\emph{Positivity.}
Let $O\in\mathcal{A}\left(Q^{\circ}_{\mathfrak{c}}\right)$.
By (5.d), $\mathrm{CC}(O)$ is a family, $Z_n\left[O^{\sharp}_n\right]=Z_n\left[\mathrm{CC}(O)\right]$, and $\mathrm{CC}(O)=\mathrm{Rm}_1^0\mathrm{CC}(O)$;
Lemma~5.1\,(ii) gives $Z_n\left[\mathrm{Rm}_1^0\mathrm{CC}(O)\right]\ge0$.

\emph{Inequality.}
By Proposition~5.2,
\[
\left|Z_n[\mathrm{GG}]\right|\le\prod_{\mathfrak{c}'\in\mathrm{Kb}(n)}Z_n\left[\mathrm{CC}(G_{\mathfrak{c}'})\right]^{\frac{1}{\mathrm{nc}(n)}}.
\]
For $\mathfrak{c}'\in\mathrm{Kb}(n)$ let $\mathfrak{c}\in\mathrm{Rs}$ be the representative of its class, $\mathfrak{c}'=\mathfrak{c}+\mathrm{sh}(v)$ with $v\in\Gamma(n)$.
Then $G_{\mathfrak{c}'}=\left(O_{\mathfrak{c}}\right)^{\tau_v}=\left(O_{\mathfrak{c}}\right)^{[\mathfrak{c}']}$, because $\tau_v$ is the element $g_{\mathrm{w},\mathfrak{c}'\leftarrow\mathfrak{c}}$ of $W(\mathcal{Q})$;
and $\mathrm{CC}\left(\left(O_{\mathfrak{c}}\right)^{[\mathfrak{c}']}\right)=\mathrm{CC}(O_{\mathfrak{c}})$ by (4.2).
Hence $Z_n\left[\mathrm{CC}(G_{\mathfrak{c}'})\right]=Z_n\left[\left(O_{\mathfrak{c}}\right)^{\sharp}_n\right]$, and $\mathfrak{c}'\mapsto\mathfrak{c}$ is a bijection $\mathrm{Kb}(n)\to\mathrm{Rs}$.
\end{proof}

\begin{proof}[Proof of Proposition~39\,(d)]
Let $O\in\mathcal{A}\left(Q^{\circ}_0(\mathcal{Q})\right)$.
The copy $O^{[\mathfrak{c}]}$ lies in $\mathcal{A}\left(Q^{\circ}_{\mathfrak{c}}\right)$,
and by Definition~39.0 its copy in the cell $\mathfrak{c}'$ is $\left(O^{[\mathfrak{c}]}\right)^{g_{\mathrm{w},\mathfrak{c}'\leftarrow\mathfrak{c}}}$, which is $O^{[\mathfrak{c}']}$ by (4.2).
So the chessboard families of $O^{[\mathfrak{c}]}$ and of $O$ coincide entry by entry,
and $\left(O^{[\mathfrak{c}]}\right)^{\sharp}_n=O^{\sharp}_n$, both products being formed with one and the same set of representatives.
\end{proof}

\begin{proof}[Proof of Proposition~39\,(e)]
The number $\mathrm{hw}\cdot\mathrm{sv}_\mu=(2\mathrm{hw})\cdot\frac{\mathrm{sv}_\mu}{2}$ is an integer, since $2\mathrm{hw}$ is an integer and $\mathrm{sv}_\mu$ is even.
So, for $\mathrm{sv}$ as in (e), the vector $n:=\mathrm{hw}\cdot\mathrm{sv}$ has integer entries $\ge4$, $n_\mu\in2\mathrm{hw}\,\mathbb{Z}$, the torus $C(n)$ carries $\mathcal{Q}$,
$n_\mu/\mathrm{hw}=\mathrm{sv}_\mu$, and $\mathrm{nc}(n)=\mathrm{sv}_1\cdot\mathrm{sv}_2\cdot\mathrm{sv}_3$.
Fix a direction $\mu$, let $\nu$ and $\mathfrak{l}$ be the two other directions, and fix $\mathrm{sv}_\nu$ and $\mathrm{sv}_{\mathfrak{l}}$.
Then $\Gamma_\mu(n)$, the class $\mathcal{A}_{\mu,n}[\,\cdot\,,\,\cdot\,]$ and the objects $\mathrm{Hs}$, $V$, $\mathrm{Cs}$, $\mathrm{kk}$ of Theorem~SO are fixed;
they do not involve $\mathrm{sv}_\mu$.
Put
\begin{gather*}
S:=\prod_{\mathfrak{c}:\ \mathfrak{c}_\mu=0,\ 0\le\mathfrak{c}_\nu<\mathrm{sv}_\nu,\ 0\le\mathfrak{c}_{\mathfrak{l}}<\mathrm{sv}_{\mathfrak{l}}}O^{[\mathfrak{c}]}\ \in\mathcal{A}_{\mu,n}[o_\mu,o_\mu+\mathrm{hw}]\qquad\text{(Lemma~2.3\,(C4))},\\
Y:=\mathrm{Cs}[S]:\ \mathrm{Hs}(o_\mu)\to\mathrm{Hs}(o_\mu+\mathrm{hw}),\qquad B:=Y^*Y\ge0,
\end{gather*}
where $B$ is of trace class on $\mathrm{Hs}(o_\mu)$.
The elements $S$, $Y$, $B$ do not depend on $\mathrm{sv}_\mu$.
(That the cells in the product lie in different orbits of $\Gamma(n)$, and the conclusion of (C4), involve the subgroup $\Gamma_\mu(n)$ only, because the cells have the same entry at $\mu$.)
For $\mathfrak{c}_\mu=0$ the element $\mathrm{dw}^\mu_{i'\leftarrow0}$ is $g_{\mathrm{w},\mathfrak{c}[\mu{:}i']\leftarrow\mathfrak{c}}\in W(\mathcal{Q})$;
so Lemma~2.2\,(B1) and (4.2) give
\begin{align*}
\mathrm{fold}_\mu S
&=\prod_{i'=0}^{\mathrm{sv}_\mu-1}S^{\mathrm{dw}^\mu_{i'\leftarrow0}}
=\prod_{i'=0}^{\mathrm{sv}_\mu-1}\ \prod_{\mathfrak{c}:\ \mathfrak{c}_\mu=0,\ \dots}\left(O^{[\mathfrak{c}]}\right)^{\mathrm{dw}^\mu_{i'\leftarrow0}}\\
&=\prod_{i'}\prod_{\mathfrak{c}}O^{[\mathfrak{c}[\mu{:}i']]}
=\prod_{\mathfrak{c}'\in\mathrm{Kb}(n)}O^{[\mathfrak{c}']}
=O^{\sharp}_n .
\end{align*}
The identity (3.1) with $2k=\mathrm{sv}_\mu$ now yields, for every even $\mathrm{sv}_\mu$ with $\mathrm{hw}\cdot\mathrm{sv}_\mu\ge4$,
\begin{equation}
\mathcal{Z}_{\mathrm{sv}}(O)=\mathrm{kk}^{|n|}\cdot\mathrm{Tr}\,B^{\frac{\mathrm{sv}_\mu}{2}}=\mathrm{kv}^{\mathrm{sv}_\mu}\cdot\mathrm{Tr}\,B^{\frac{\mathrm{sv}_\mu}{2}},\qquad\mathrm{kv}:=\mathrm{kk}^{\mathrm{hw}^3\cdot\mathrm{sv}_\nu\cdot\mathrm{sv}_{\mathfrak{l}}}>0.
\tag{6.1}
\end{equation}

\emph{Monotonicity.}
One has
\[
\mathcal{Z}_{\mathrm{sv}}(O)^{\frac{1}{\mathrm{sv}_\mu}}=\mathrm{kv}\cdot\left(\mathrm{Tr}\,B^{\frac{\mathrm{sv}_\mu}{2}}\right)^{\frac{1}{\mathrm{sv}_\mu}}=\mathrm{kv}\cdot\|B\|_{\mathrm{sv}_\mu/2}^{\frac12},
\]
and $\|B\|_p$ is non-increasing in $p\ge1$ (Section~1.7\,(b); note that $\mathrm{sv}_\mu/2\ge1$).

\emph{Positivity.}
One has $\mathrm{Tr}\,B^p>0$ for one $p\ge1$ if and only if $B\neq0$, and this holds if and only if $\mathrm{Tr}\,B^p>0$ for all $p\ge1$.

\emph{Convexity.}
Let $B\neq0$.
Then
\[
\log\mathcal{Z}_{\mathrm{sv}}(O)=\mathrm{sv}_\mu\cdot\log\mathrm{kv}+\log\mathrm{Tr}\,B^{\frac{\mathrm{sv}_\mu}{2}} ;
\]
the first term is linear in $\mathrm{sv}_\mu$,
and the second is the restriction to $\{\mathrm{sv}_\mu\text{ even},\ \mathrm{hw}\cdot\mathrm{sv}_\mu\ge4\}$ of the function $t\mapsto\log\mathrm{Tr}\,B^{\frac{t}{2}}$,
which is convex on $[2,\infty)$ by Section~1.7\,(c).
So, if $\mathcal{Z}(s)$ denotes $\mathcal{Z}_{\mathrm{sv}}(O)$ at $\mathrm{sv}_\mu=s$, then for $s<s'<s''$ in this set
\[
(s''-s)\cdot\log\mathcal{Z}(s')\le(s''-s')\cdot\log\mathcal{Z}(s)+(s'-s)\cdot\log\mathcal{Z}(s'').
\]
\end{proof}

\begin{proof}[Proof of Proposition~39\,(f)]
For $n\in2\mathrm{hw}\,\mathbb{Z}$ with $n\ge4$ the cubic torus $C(n,n,n)$ carries $\mathcal{Q}$, and the number of its cells is $\mathrm{nc}:=\mathrm{nc}(n,n,n)=\left(\frac{n}{\mathrm{hw}}\right)^3$.
One has $0<Z_n[1]<\infty$ (Section~1.8), and $E^{\mathrm{af}}_n=Z_n[\,\cdot\,]/Z_n[1]$ (Section~I.7.2; on $\mathcal{A}$ by (1.3)).
Divide the two statements of (c) by $Z_n[1]=\prod_{\mathfrak{c}}Z_n[1]^{\frac{1}{\mathrm{nc}}}$, a product of $\mathrm{nc}$ factors;
here $\frac{1}{\mathrm{nc}}=\left(\frac{\mathrm{hw}}{n}\right)^3$.

\emph{Functions of the curvature.}
Let $O\in\mathcal{A}\left(Q^{\circ}_{\mathfrak{c}}\right)$ be a function of the curvature, $O(\hat A)=O(d\hat A)$.
By (1.1), and since $\Theta_m$ acts on numbers by complex conjugation,
$O^{g_{\mathrm{w}}}(\hat A)$ equals $O\left(d\,g_{\mathrm{w}}^*\hat A\right)=O\left(g_{\mathrm{w}}^*\,d\hat A\right)$ for a word of even length
and $\overline{O\left(g_{\mathrm{w}}^*\,d\hat A\right)}$ for a word of odd length;
the parity of the length is that of $g_{\mathrm{w}}$.
For $g_{\mathrm{w}}=g_{\mathrm{w},\mathfrak{c}'\leftarrow\mathfrak{c}}$ this is the assertion.
\end{proof}

\paragraph{Uniformity.}
The clauses (a)--(f) of Proposition~39 hold for every real $\boldsymbol{\mu}$, every $q>0$, every $\rho\in(0,1)$, every coordinate cell structure $\mathcal{Q}$ of Definition~39.0 and every torus $C(n)$ that carries $\mathcal{Q}$.
The chessboard estimates of (c) and (f) hold with the factor $1$ and with the exponent $\frac{1}{\mathrm{nc}(n)}$;
in (f), on the cubic torus $C(n,n,n)$, this exponent is $\left(\frac{\mathrm{hw}}{n}\right)^3=\frac{1}{\mathrm{nc}(n,n,n)}$.
In (e) the exponent is $\frac{1}{\mathrm{sv}_\mu}$.
All objects of Proposition~39 are objects of level~$0$, written in lattice units ($\eta_0=1$, Section~I.2.5).

\subsection*{7. The width in Sections 2--6; the five replacements; three further statements of \cite[Section~6]{QEDII}}

\subsubsection*{7.1. The places at which the width enters}

In Sections~2--6 the number $\mathrm{hw}$ enters only through the following properties.
They substantiate the Remark that follows Proposition~39, which lists the places at which the width enters the proofs.
\begin{itemize}
\item[(w0)]
$\mathrm{hw}>0$ (Definition~39.0):
the walls $\mathrm{po}_j=o_\mu+j\cdot\mathrm{hw}$ are strictly increasing in $j$, so that the layers and the chains of walls are formed (Section~3.1, Proposition~3.1);
$\tau_{\tau_\mu}$ has infinite order, so that $\mathrm{Dw}_\mu$ is infinite (Section~4);
and $\mathrm{hw}$ is a divisor in (w3), (w5), (w6).
\item[(w1)]
$o_\mu+j\cdot\mathrm{hw}\in\tfrac12\mathbb{Z}$ for all $j$ (Section~3.1):
the walls are mirrors, and Theorem~SO applies to the chains of walls.
This requires $o_\mu,\mathrm{hw}\in\tfrac12\mathbb{Z}$ (point~(1) of that Remark).
\item[(w2)]
$\tau_\mu=2\mathrm{hw}\cdot e_\mu\in\mathbb{Z}^3$ (Section~3.1):
$\tau_{s\tau_\mu}$ acts on $\mathbb{Z}^3$, on cochains and on fermion sites,
and $V_{s\tau_\mu}$ is defined and maps $\mathrm{Hs}(\mathrm{po})$ to $\mathrm{Hs}(\mathrm{po}+2s\cdot\mathrm{hw})$, the space of a plane of the same kind.
This requires $2\mathrm{hw}\in\mathbb{Z}$ (point~(2) of that Remark).
\item[(w3)]
$n_\mu/(2\mathrm{hw})$ is an integer $\ge1$ (Section~3.1, Proposition~3.1, Lemma~5.1\,(i)):
the number $2k=n_\mu/\mathrm{hw}$ of the layers is even, and the translation by $n_\mu e_\mu=k\cdot\tau_\mu$ lies in $\mathrm{Dw}_\mu$.
This is the condition that $C(n)$ carries $\mathcal{Q}$ (point~(8) of that Remark).
\item[(w4)]
$2\mathrm{rd}<\mathrm{hw}$ for a box in a cell, and $|v_\mu|\ge2\mathrm{hw}$ for $0\neq v_\mu\in n_\mu\mathbb{Z}$ (Lemma~2.3\,(C4)):
inequalities between real numbers (point~(4) of that Remark).
\item[(w5)]
$\mathrm{sh}(v)=(v_\mu/\mathrm{hw})_\mu$ for $v\in\mathbb{Z}^3\cap(\mathrm{hw}\,\mathbb{Z})^3$ (Section~5.1; point~(5) of that Remark).
\item[(w6)]
The exponents $\mathrm{hw}/n_\mu$, the number $\mathrm{nc}(n)=\prod_\mu(n_\mu/\mathrm{hw})=|n|/\mathrm{hw}^3$, and $\mathrm{hw}^3$ in the exponent of $\mathrm{kk}$ in (6.1):
here $\mathrm{hw}^3$ occurs as a real number only (point~(6) of that Remark).
\item[(w7)]
$\mathrm{hw}\cdot\mathrm{sv}_\mu$ is an integer for even $\mathrm{sv}_\mu$ (proof of (e); point~(7) of that Remark).
\end{itemize}
The property (w0) is the hypothesis $\mathrm{hw}>0$ of Definition~39.0;
\cite[Definition~2.18]{QEDII} contains it in the words ``a positive integer $h$'', and the widening (I.h1) of Section~I.6.4 keeps it.
Point~(3) of the Remark that follows Proposition~39 is the generality of Proposition~3.1:
this proposition holds for every height $\mathrm{hw}$ with $\mathrm{po}_0+j\cdot\mathrm{hw}\in\tfrac12\mathbb{Z}$;
if $2\mathrm{hw}$ is odd, the operators $\mathrm{Cs}[G]:\mathrm{Hs}(\mathrm{po}_j)\to\mathrm{Hs}(\mathrm{po}_{j+1})$ act between the space of a site plane and the space of a link plane,
and Lemma~2.1 is H\"older's inequality between different spaces.
Two consecutive walls enter the proofs only through (SO2b) and (SO3) of Theorem~SO, which are statements about arbitrary pairs and chains in $\tfrac12\mathbb{Z}$.
Neither these two statements nor any of the properties (w0)--(w7) requires $\mathrm{hw}$ to be an integer.

\subsubsection*{7.2. The five replacements (h1)--(h5) of Section~I.6.4, and the cells on tori of even side}

\textup{(h1)} Besides definitions, \cite[Definition~2.18]{QEDII} contains the following assertions:
the walls are mirrors;
$W(\mathcal{Q})$ is a product of three infinite dihedral groups and acts simply transitively on the cells;
if $\mathcal{Q}$ cuts a torus $T$, then the lattice of $T$ permutes the cells, and the number $n(T,\mathcal{Q})$ of the cells of $T$ is $|T|/|Q|$;
if $T$ carries $\mathcal{Q}$, then the lattice of $T$ lies in the translation subgroup of $W(\mathcal{Q})$, and the quotient acts simply transitively on the cells of $T$;
the walls of a coarser structure $\mathcal{Q}'$ are walls of $\mathcal{Q}$, and a cell of $\mathcal{Q}'$ consists of $\left(\frac{\mathrm{hw}'}{\mathrm{hw}}\right)^3$ cells of $\mathcal{Q}$, where $\mathrm{hw}'$ denotes the width of $\mathcal{Q}'$.
For the frame $\mathfrak{F}_{\mathrm{c}}$, offsets and width in $\frac{1}{2}\mathbb{Z}$, and a torus $T=C(n)$ that carries $\mathcal{Q}$, the first four of these assertions are Proposition~39\,(a), with $|Q|=\mathrm{hw}^3$ and $|n|/\mathrm{hw}^3=\mathrm{nc}(n)$.

\emph{Tori that are cut.}
Let $T=(\Gamma,\chi)$ be a torus in the sense of \cite[Definition~2.11]{QEDII} that $\mathcal{Q}$ cuts, that is, $v_\mu\in\mathrm{hw}\,\mathbb{Z}$ for all $v\in\Gamma$ and $\mu=1,2,3$.
Then $\mathrm{sh}(v)=\left(\frac{v_\mu}{\mathrm{hw}}\right)_\mu$ is defined on $\Gamma$;
it is additive and injective, and $\tau_vQ^{\circ}_{\mathfrak{c}}=Q^{\circ}_{\mathfrak{c}+\mathrm{sh}(v)}$ (Section~5.1, paragraph ``Shifts''; no torus enters there).
Hence $\Gamma$ permutes the cells, and the orbits correspond to the classes of $\mathbb{Z}^3$ modulo $\mathrm{sh}(\Gamma)$.
The image $\mathrm{sh}(\Gamma)$ is a subgroup of rank three of $\mathbb{Z}^3$, and $\frac{1}{\mathrm{hw}}$ times a basis of $\Gamma$ is a basis of it;
the index in $\mathbb{Z}^3$ of a subgroup of rank three is the modulus of the determinant of a basis.
Therefore
\[
n(T,\mathcal{Q})=\left[\mathbb{Z}^3:\mathrm{sh}(\Gamma)\right]=\frac{\left[\mathbb{Z}^3:\Gamma\right]}{\mathrm{hw}^3}=\frac{|T|}{|Q|}.
\]

\emph{Coarser structures.}
Let $\mathcal{Q}'=(\mathfrak{F}_{\mathrm{c}};o';\mathrm{hw}')$ be coarser than $\mathcal{Q}$:
$\mathrm{hw}'=t\cdot\mathrm{hw}$ and $o'_\mu=o_\mu+i_\mu\cdot\mathrm{hw}$ with integers $t\ge1$ and $i_\mu$.
Then $o'_\mu+j\cdot\mathrm{hw}'=o_\mu+(i_\mu+jt)\cdot\mathrm{hw}$, so that every wall of $\mathcal{Q}'$ is a wall of $\mathcal{Q}$.
The interval $\left(o'_\mu+\varkappa_\mu\cdot\mathrm{hw}',\,o'_\mu+(\varkappa_\mu+1)\cdot\mathrm{hw}'\right)$ is the union of the $t$ intervals $\left(o_\mu+i\cdot\mathrm{hw},\,o_\mu+(i+1)\cdot\mathrm{hw}\right)$, $i_\mu+\varkappa_\mu t\le i<i_\mu+(\varkappa_\mu+1)t$, and of the $t-1$ points between them.
So the cell of $\mathcal{Q}'$ with the index $\varkappa$ is the union of the $t^3=\left(\frac{\mathrm{hw}'}{\mathrm{hw}}\right)^3$ cells $Q^{\circ}_{\mathfrak{c}}(\mathcal{Q})$ with $i_\mu+\varkappa_\mu t\le\mathfrak{c}_\mu<i_\mu+(\varkappa_\mu+1)t$, $\mu=1,2,3$, and of points that lie on walls of $\mathcal{Q}$.

Consequently, two of the assertions listed above hold for all offsets in $\frac{1}{2}\mathbb{Z}$ and all widths $\mathrm{hw}\in\frac{1}{2}\mathbb{Z}$ with $\mathrm{hw}>0$, whether or not $\mathrm{hw}$ is an integer.
The first is the third assertion of the list:
for every torus $T=(\Gamma,\chi)$ of \cite[Definition~2.11]{QEDII} that $\mathcal{Q}$ cuts, the lattice $\Gamma$ permutes the cells, and $n(T,\mathcal{Q})=|T|/|Q|$.
The second is the fifth:
the walls of a coarser structure $\mathcal{Q}'$ are walls of $\mathcal{Q}$, and a cell of $\mathcal{Q}'$ consists of $\left(\frac{\mathrm{hw}'}{\mathrm{hw}}\right)^3$ cells of $\mathcal{Q}$ together with points on walls of $\mathcal{Q}$.
The other statements of \cite{QEDII} that (h1) widens are, for these widths and $T=C(n)$, given by the following statements of this chapter:
\begin{itemize}
\item \cite[Definition~2.19]{QEDII}: Definition~39.0 with (4.2), Section~5.1 and Proposition~39\,(c) (positivity).
The fold of \cite[Definition~2.19]{QEDII} is here the product $O^{\sharp}_n=\mathrm{Prd}(\mathrm{CC}(O);\mathrm{Rs})$ over the cells of $C(n)$;
(5.d) gives this identity and $\left(O^{[\mathfrak{c}']}\right)^{\tau_v}=O^{[\mathfrak{c}'+\mathrm{sh}(v)]}$ for $v\in\Gamma(n)$,
(5.a) gives $O^{\sharp}_n\in\mathcal{A}(C(n))$,
and (5.b) gives that the function of $O^{\sharp}_n$ on the configurations, and $Z_n\left[O^{\sharp}_n\right]$, do not depend on the set $\mathrm{Rs}$ of representatives;
\item \cite[Definition~6.7]{QEDII}: Section~3.1;
\item \cite[Lemma~6.5(d)]{QEDII}: Lemma~2.3\,(C4);
\item \cite[Definition~6.13]{QEDII}, \cite[Lemma~6.14]{QEDII} and \cite[Proposition~6.15]{QEDII}: Section~5;
\item \cite[Theorem~6.16(a)]{QEDII}: Proposition~39\,(c);
\item \cite[Lemma~6.17]{QEDII}: Lemma~7.1;
\item \cite[Proposition~6.23(a)]{QEDII}: Proposition~39\,(e);
\item \cite[Proposition~6.23(b)]{QEDII}: Proposition~7.2\,(b$'$);
\item \cite[Proposition~12.10(f)]{QEDII}: Proposition~39\,(f).
\end{itemize}

\emph{The images of $\mathfrak{F}_{\mathrm{c}}$ under $O_{\mathrm{cub}}$.}
Such an image is a triple $\mathrm{fr}=(y_{\mu_1},y_{\mu_2},y_{\mu_3})$, where $(\mu_1,\mu_2,\mu_3)$ is a permutation of $(1,2,3)$ (the planes $\{-y_\mu=\mathrm{po}\}$ are the planes $\{y_\mu=-\mathrm{po}\}$).
Let $\mathcal{Q}'=(\mathrm{fr};o'_1,o'_2,o'_3;\mathrm{hw})$ with $o'_a,\mathrm{hw}\in\frac{1}{2}\mathbb{Z}$, and put $o_{\mu_a}:=o'_a$ and $\mathcal{Q}:=(\mathfrak{F}_{\mathrm{c}};o;\mathrm{hw})$.
The walls $(y_{\mu_a},o'_a+j\cdot\mathrm{hw})$ of $\mathcal{Q}'$ are the walls of $\mathcal{Q}$;
the cell of $\mathcal{Q}'$ with the index $\varkappa$ is $Q^{\circ}_{\mathfrak{c}}(\mathcal{Q})$ with $\mathfrak{c}_{\mu_a}=\varkappa_a$;
$W(\mathcal{Q}')=W(\mathcal{Q})$;
the rectangular torus $\mathrm{Rect}(\mathrm{fr};p_1,p_2,p_3)$ of \cite[Definition~2.18]{QEDII} is $C(n)$ with $n_{\mu_a}=p_a$, and it carries $\mathcal{Q}'$ if and only if $C(n)$ carries $\mathcal{Q}$.
The clauses (a)--(f) of Proposition~39 and the two paragraphs above, on tori that are cut and on coarser structures, do not involve the numbering of the directions;
so they hold for $\mathcal{Q}'$.
For the frame $\mathrm{fr}$, \cite[Proposition~6.15]{QEDII} is the chain of inequalities in the order $\mu_1,\mu_2,\mu_3$:
\begin{multline*}
\left|Z_n[\mathrm{GG}]\right|
\le\prod_{j'_1}Z_n\left[\mathrm{Rm}_{\mu_1}^{j'_1}\,\mathrm{GG}\right]^{\frac{\mathrm{hw}}{n_{\mu_1}}}
\le\prod_{j'_1,j'_2}Z_n\left[\mathrm{Rm}_{\mu_2}^{j'_2}\,\mathrm{Rm}_{\mu_1}^{j'_1}\,\mathrm{GG}\right]^{\frac{\mathrm{hw}^2}{n_{\mu_1}n_{\mu_2}}}\\
\le\prod_{\mathfrak{c}\in\mathrm{Kb}(n)}Z_n\left[\mathrm{CC}(G_{\mathfrak{c}})\right]^{\frac{1}{\mathrm{nc}(n)}}
\qquad\left(0\le j'_a<\frac{n_{\mu_a}}{\mathrm{hw}}\right).
\end{multline*}

\begin{proof}[Proof of the chain for the frame $\mathrm{fr}$]
Lemma~5.1 holds for each direction $\mu$.
Its part (ii) for $(\mathrm{GG},\mu_1)$ is the first inequality.
By (i) and (ii), each $\mathrm{Rm}_{\mu_1}^{j'_1}\,\mathrm{GG}$ is a family, and $Z_n$ of it is $\ge0$;
part (ii) for $(\mathrm{Rm}_{\mu_1}^{j'_1}\,\mathrm{GG},\mu_2)$ bounds it by $\prod_{j'_2}Z_n\left[\mathrm{Rm}_{\mu_2}^{j'_2}\,\mathrm{Rm}_{\mu_1}^{j'_1}\,\mathrm{GG}\right]^{\frac{\mathrm{hw}}{n_{\mu_2}}}$, and $t\mapsto t^{\frac{\mathrm{hw}}{n_{\mu_1}}}$ is non-decreasing on $[0,\infty)$;
this gives the second inequality.
By (i) and (ii), each $\mathrm{Rm}_{\mu_2}^{j'_2}\,\mathrm{Rm}_{\mu_1}^{j'_1}\,\mathrm{GG}$ is a family, and $Z_n$ of it is $\ge0$;
part (ii) for this family and $\mu_3$ bounds it by $\prod_{j'_3}Z_n\left[\mathrm{Rm}_{\mu_3}^{j'_3}\,\mathrm{Rm}_{\mu_2}^{j'_2}\,\mathrm{Rm}_{\mu_1}^{j'_1}\,\mathrm{GG}\right]^{\frac{\mathrm{hw}}{n_{\mu_3}}}$;
the map $t\mapsto t^{\frac{\mathrm{hw}^2}{n_{\mu_1}n_{\mu_2}}}$ is non-decreasing on $[0,\infty)$, and
\[
\frac{\mathrm{hw}}{n_{\mu_1}}\cdot\frac{\mathrm{hw}}{n_{\mu_2}}\cdot\frac{\mathrm{hw}}{n_{\mu_3}}=\frac{1}{\mathrm{nc}(n)}.
\]
Let $\mathrm{cj}\in\mathbb{Z}^3$ be the index with $\mathrm{cj}_{\mu_a}=j'_a$, and put $d_a:=\mathrm{dw}^{\mu_a}_{\mathfrak{c}_{\mu_a}\leftarrow j'_a}$.
By the definition of $\mathrm{Rm}$, applied three times, and by (S6),
\[
\left(\mathrm{Rm}_{\mu_3}^{j'_3}\,\mathrm{Rm}_{\mu_2}^{j'_2}\,\mathrm{Rm}_{\mu_1}^{j'_1}\,\mathrm{GG}\right)_{\mathfrak{c}}
=\left(\left(\left(G_{\mathrm{cj}}\right)^{d_1}\right)^{d_2}\right)^{d_3}
=\left(G_{\mathrm{cj}}\right)^{d_3d_2d_1}.
\]
Elements of $\mathrm{Dw}_\mu$ and of $\mathrm{Dw}_\nu$, $\mu\neq\nu$, commute (Section~4);
so $d_3d_2d_1$ is the product (4.1), that is, the element $g_{\mathrm{w},\mathfrak{c}\leftarrow\mathrm{cj}}$, and the family is $\mathrm{CC}(G_{\mathrm{cj}})$ by (4.2).
The index $\mathrm{cj}$ ranges over $\mathrm{Kb}(n)$.
\end{proof}

\textup{(h2)} In \cite[Definition~6.13]{QEDII}, the assertion that $o_a+jh$ lies in $\operatorname{Adm}(\ell^a)$ is, for these widths, the property (w1).

\textup{(h3)} In \cite[Definition~6.22]{QEDII}, the condition that the vector $\frac{p_a}{|n_a|^2}\cdot n_a$ be integral holds for these widths with $p_a=\mathrm{hw}\cdot\mathrm{sv}_\mu$, $|n_a|^2=1$ and $n_a=e_\mu$;
it is the property (w7).

\textup{(h4)} For widths in $\frac{1}{2}\mathbb{Z}$ the sides $\mathrm{hw}\cdot\mathrm{sv}_\mu$, $\mathrm{sv}_\mu$ even, need not be even:
for $\mathrm{hw}=\frac{5}{2}$ and $\mathrm{sv}_\mu=2$ the side is $5$.
For $\mathrm{hw}=\frac{1}{2}$ every $n$ with $n_\mu\ge4$ is $\mathrm{hw}\cdot\mathrm{sv}$ with $\mathrm{sv}=2n$, all of whose entries are even.
For $\mathrm{hw}=\frac{1}{2}$ and $O=1$ the identity (6.1) reads
\[
Z_n[1]=\mathrm{kk}^{|n|}\cdot\operatorname{Tr}B^{n_\mu}\qquad\text{for all }n_\mu\ge4,
\]
with $B=\mathrm{Cs}[1]^{*}\cdot\mathrm{Cs}[1]$ formed with a half step.

\textup{(h5)} \cite[Proposition~12.10(f)]{QEDII}, for widths $\mathrm{hw}\in\frac{1}{2}\mathbb{Z}$ and sides $n\in2\mathrm{hw}\,\mathbb{Z}$, is Proposition~39\,(f).

\emph{The cells on tori of even side.}
By Definition~39.0 the torus $C(n,n,n)$ carries $\mathcal{Q}_*=(\mathfrak{F}_{\mathrm{c}};0,0,0;\mathrm{hc})$ if and only if $n\in2\mathrm{hc}\,\mathbb{Z}=\mathrm{per}\,\mathbb{Z}$.
The parity of $n$ does not enter.

\subsubsection*{7.3. \cite[Lemma~6.17]{QEDII} and \cite[Proposition~6.23(b)]{QEDII} for widths in $\frac{1}{2}\mathbb{Z}$}

\begin{named}{Lemma 7.1}[{\cite[Lemma~6.17]{QEDII}}]
Let $C(n)$ carry $\mathcal{Q}$, and let $g_{\mathrm{w}}'=(a,r)\in W_{\mathrm{af}}$ satisfy $r\,\Gamma(n)=\Gamma(n)$ and map the set of the walls of $\mathcal{Q}$ onto itself.
Then, for $O\in\mathcal{A}(Q^{\circ}_{\mathfrak{c}})$, the set $g_{\mathrm{w}}'Q^{\circ}_{\mathfrak{c}}$ is a cell, $O^{g_{\mathrm{w}}'}$ lies in $\mathcal{A}(g_{\mathrm{w}}'Q^{\circ}_{\mathfrak{c}})$, and
\[
Z_n\left[\left(O^{g_{\mathrm{w}}'}\right)^{\sharp}_n\right]=Z_n\left[O^{\sharp}_n\right].
\]
This applies to $g_{\mathrm{w}}'=\tau_v$ with $v\in\mathbb{Z}^3\cap(\mathrm{hw}\,\mathbb{Z})^3$.
\end{named}

\begin{proof}
One has $g_{\mathrm{w}}'\,\theta_m\,g_{\mathrm{w}}'^{-1}=\theta_{g_{\mathrm{w}}'m}$;
so $g_{\mathrm{w}}'\,W(\mathcal{Q})\,g_{\mathrm{w}}'^{-1}=W(\mathcal{Q})$.
The cells are the connected components of the complement of the union of the walls;
so $g_{\mathrm{w}}'$ permutes them.
By Lemma~2.2\,(B2), $O^{g_{\mathrm{w}}'}$ lies in $\mathcal{A}(g_{\mathrm{w}}'Q^{\circ}_{\mathfrak{c}})$.
Let $\mathrm{Rs}$ be a set of representatives of the cells of $C(n)$, and for $\mathfrak{c}'\in\mathrm{Rs}$ let $g_{\mathrm{w},1}\in W(\mathcal{Q})$ be the element that maps $g_{\mathrm{w}}'Q^{\circ}_{\mathfrak{c}}$ onto $Q^{\circ}_{\mathfrak{c}'}$.
By (S6),
\[
\left(O^{g_{\mathrm{w}}'}\right)^{g_{\mathrm{w},1}}=O^{g_{\mathrm{w},1}g_{\mathrm{w}}'}=O^{g_{\mathrm{w}}'g_{\mathrm{w},2}}=\left(O^{g_{\mathrm{w},2}}\right)^{g_{\mathrm{w}}'},
\qquad g_{\mathrm{w},2}:=g_{\mathrm{w}}'^{-1}g_{\mathrm{w},1}g_{\mathrm{w}}'\in W(\mathcal{Q}),
\]
and $g_{\mathrm{w},2}$ maps $Q^{\circ}_{\mathfrak{c}}$ onto the cell $g_{\mathrm{w}}'^{-1}Q^{\circ}_{\mathfrak{c}'}=:Q^{\circ}_{\mathfrak{c}''}$;
so $\left(O^{g_{\mathrm{w}}'}\right)^{g_{\mathrm{w},1}}=\left(O^{[\mathfrak{c}'']}\right)^{g_{\mathrm{w}}'}$.
Let $\mathrm{Rs}''$ be the set of these $\mathfrak{c}''$.
By Lemma~2.2\,(B1),
\[
\left(O^{g_{\mathrm{w}}'}\right)^{\sharp}_n=\prod_{\mathfrak{c}'\in\mathrm{Rs}}\left(O^{g_{\mathrm{w}}'}\right)^{[\mathfrak{c}']}=\left(\mathrm{Prd}(\mathrm{CC}(O);\mathrm{Rs}'')\right)^{g_{\mathrm{w}}'}.
\]
One has $g_{\mathrm{w}}'^{-1}(Q^{\circ}+v)=g_{\mathrm{w}}'^{-1}Q^{\circ}+r^{-1}v$ for a cell $Q^{\circ}$, and $r^{-1}\Gamma(n)=\Gamma(n)$;
so $g_{\mathrm{w}}'^{-1}$ maps the orbits of cells bijectively onto orbits of cells, and $\mathrm{Rs}''$ is a set of representatives.
The element $\mathrm{Prd}(\mathrm{CC}(O);\mathrm{Rs}'')$ embeds in $\Gamma(n)$, term by term (Lemma~2.3\,(C4)).
Further $\chi_n\circ r^{-1}=\chi_n$, because $r$ permutes the vectors $\pm n_\mu e_\mu$;
so $g_{\mathrm{w}}'C(n)=C(n)$ in the sense of \cite[Definition~2.11]{QEDII}.
By \cite[Proposition~3.16(g)]{QEDII}, the number $Z_n\left[\left(\mathrm{Prd}(\mathrm{CC}(O);\mathrm{Rs}'')\right)^{g_{\mathrm{w}}'}\right]$ is $Z_n\left[\mathrm{Prd}(\mathrm{CC}(O);\mathrm{Rs}'')\right]$ or its complex conjugate;
by (5.b) and (5.d) this number is $Z_n\left[O^{\sharp}_n\right]$, which is real by Proposition~39\,(c).
For $g_{\mathrm{w}}'=\tau_v$ one has $r=1$, and $\tau_v$ maps the wall $(y_\mu,\mathrm{po})$ to $(y_\mu,\mathrm{po}+v_\mu)$ by (M4).
\end{proof}

\begin{named}{Proposition 7.2}[{clause (b$'$): \cite[Proposition~6.23(b)]{QEDII}}]
Let $O\in\mathcal{A}(Q^{\circ}_0(\mathcal{Q}))$ and
\[
\mathrm{Nd}:=\left\{\mathrm{sv}:\ \mathrm{sv}_\mu\text{ even},\ \mathrm{hw}\cdot\mathrm{sv}_\mu\ge4,\ \mu=1,2,3\right\}.
\]
\begin{itemize}
\item[(b$'$)] If $\mathcal{Z}_{\mathrm{sv}}(O)>0$ for one $\mathrm{sv}\in\mathrm{Nd}$, then $\mathcal{Z}_{\mathrm{sv}}(O)>0$ for all $\mathrm{sv}\in\mathrm{Nd}$, and $\left(\mathrm{sv}_1\cdot\mathrm{sv}_2\cdot\mathrm{sv}_3\right)^{-1}\cdot\log\mathcal{Z}_{\mathrm{sv}}(O)$ is non-increasing in each $\mathrm{sv}_\mu$, is bounded below, and converges to its infimum, which is finite, as $\min_\mu\mathrm{sv}_\mu\to\infty$ in any manner.
\item[(b)] If $O$ is a real function of the curvature with $0<\inf O\le\sup O<\infty$, then $\mathcal{Z}_{\mathrm{sv}}(O)>0$ and
\[
\left|\log\mathcal{Z}_{\mathrm{sv}}(O)-\log\mathcal{Z}_{\mathrm{sv}}(1)\right|\le\mathrm{sv}_1\cdot\mathrm{sv}_2\cdot\mathrm{sv}_3\cdot\sup|\log O|.
\]
\end{itemize}
\end{named}

\begin{proof}
(b$'$) The set $\mathrm{Nd}$ is the product of the three sets $\mathrm{Nd}_\mu=\{\mathrm{sv}_\mu\text{ even},\ \mathrm{hw}\cdot\mathrm{sv}_\mu\ge4\}$, each of which is unbounded above.
Let $\mathrm{sv}^0\in\mathrm{Nd}$ be a triple with $\mathcal{Z}_{\mathrm{sv}^0}(O)>0$, and let $\mathrm{sv}\in\mathrm{Nd}$.
By Proposition~39\,(e), positivity extends from $\mathrm{sv}^0$ to $(\mathrm{sv}_1,\mathrm{sv}^0_2,\mathrm{sv}^0_3)$, from there to $(\mathrm{sv}_1,\mathrm{sv}_2,\mathrm{sv}^0_3)$, and from there to $\mathrm{sv}$.
By Proposition~39\,(e), the function $\varphi(\mathrm{sv}):=\log\mathcal{Z}_{\mathrm{sv}}(O)$ has, in each variable, the two properties (Mon) [$s\mapsto\varphi/s$ is non-increasing] and (Cvx) [the three-point inequality of convexity] that are the hypotheses of \cite[Lemma~6.18]{QEDII}.
That lemma, a statement on real functions on a product of three subsets of $(0,\infty)$ that are unbounded above, gives the assertions through its parts (a), (c) and (d).

(b) Let $n:=\mathrm{hw}\cdot\mathrm{sv}$.
By Proposition~39\,(f), $O^{[\mathfrak{c}]}(u)=O(g_{\mathrm{w}}^*u)$, since $O$ is real.
With $s:=\sup|\log O|$ one has $\exp(-s)\le O^{[\mathfrak{c}]}\le\exp(s)$, and $O^{\sharp}_n$ is a product of $\mathrm{nc}(n)=\mathrm{sv}_1\cdot\mathrm{sv}_2\cdot\mathrm{sv}_3$ such factors.
For a function $\Phi\ge0$ of the curvature, $Z_n[\Phi]\ge0$, because $\operatorname{det}(D_A+\boldsymbol{\mu})\ge0$ (Proposition~50\,(a)).
Applied to $O^{\sharp}_n-\exp(-s\cdot\mathrm{nc}(n))$ and to $\exp(s\cdot\mathrm{nc}(n))-O^{\sharp}_n$, this gives
\[
\exp\left(-s\cdot\mathrm{nc}(n)\right)\cdot Z_n[1]\le\mathcal{Z}_{\mathrm{sv}}(O)\le\exp\left(s\cdot\mathrm{nc}(n)\right)\cdot Z_n[1],
\]
and $\mathcal{Z}_{\mathrm{sv}}(1)=Z_n[1]>0$ by Section~1.8.
\end{proof}

\subsection*{8. Remark~39.R}

In 8.1 and 8.2 the setting and the notation are those of Appendix~A (Section~A.0):
$\mu$ is a direction, the direction of time ($\mu=1$ in Remark~39.R);
$\nu<\mathfrak{l}$ are the two other directions;
$n_\nu,n_{\mathfrak{l}}\ge4$;
$\mathrm{Gmp}$, $\mathrm{chp}$ and $v_s=n_\nu n_{\mathfrak{l}}$ are those of A.0.

\subsubsection*{8.1. The objects of \cite[Definition~2.27]{QEDII}; the display for every time length}

\emph{Objects.}
Put
\begin{gather*}
\mathrm{Fk}':=\mathrm{Fk}(0),\qquad\mathrm{Fk}:=\mathrm{Fk}(\tfrac{1}{2}),\qquad\mathrm{cc}':=\mathrm{cc}_0,\qquad\mathrm{cc}:=\mathrm{cc}_{\frac{1}{2}};\\
\mathrm{Wq}(a):=\operatorname{diag}\left(\exp\left(iq\sum_{b\in\mathrm{Cu}^1(\frac{1}{2})}\mathrm{Qc}_I(b)\cdot a_b\right)\right)_{I\subset\mathrm{Ri}(\frac{1}{2})}\in\operatorname{End}\mathrm{Fk},\qquad a\in\mathbb{R}^{\mathrm{Cu}^1(\frac{1}{2})};\\
\mathrm{Nn}(B):=\mathrm{Ms}_{0,\frac{1}{2}}\left[\exp\left(-S_{(0,\frac{1}{2})}(B)\right)\right]:\mathrm{Fk}'\to\mathrm{Fk},\qquad
\mathrm{Qs}(B):=\mathrm{cc}^{\frac{1}{2}}\cdot\mathrm{Nn}(B)\cdot\mathrm{cc}'^{\frac{1}{2}}:\mathrm{Fk}'\to\mathrm{Fk},\qquad B\in\mathbb{R}^{\mathrm{Mi}^1(0)}
\end{gather*}
(A.6 for the first line, A.8 for the second, A.3 and A.7 for the third).
One has $\mathrm{Ln}[0,\frac{1}{2}]=\mathrm{Mi}^1(0)$, because $\mathrm{Sb}^1(0,\frac{1}{2})$ and $\mathrm{Mi}^1(\frac{1}{2})$ are empty (A.1, A.2).
So, by (Bd2), the element $S_{(0,\frac{1}{2})}(A)$ contains $A$ only through the phases $\exp(\pm iqA(b))$, $b\in\mathrm{Mi}^1(0)$, and $S_{(0,\frac{1}{2})}(B)$ is defined for a real function $B$ on $\mathrm{Mi}^1(0)$.
The matrices $\mathrm{cc}^{\frac{1}{2}}$ and $\mathrm{cc}'^{\frac{1}{2}}$ are diagonal, with the entries $c_I^{\frac{1}{2}}$ ($c_I>0$ by A.6).
The set $\mathrm{Mi}^1(0)$ consists of the $2v_s$ links of the plane $\{y_\mu=0\}$ of $\mathrm{Cy}$, that is, of the slice $\mathbb{Z}^2/(n_\nu\mathbb{Z}\times n_{\mathfrak{l}}\mathbb{Z})$;
$\mathrm{Cu}^1(\frac{1}{2})$ is the set of the $v_s$ links $(c;\{\mu\})$ with $c_\mu=0$, one over each site $y$ of the slice, and $a_y$, $\mathrm{Qc}_I(y)$ denote the values at this link.
All these objects are formed with $(\mu;n_\nu,n_{\mathfrak{l}})$ and the parameters $(q,\rho,M_{\mathrm{ch}})$.

For $\mu=1$ they are the objects of \cite[Definition~2.27]{QEDII} (its paragraphs ``Coordinate planes'' and ``Half slab'') for the cross-section $(n_2,n_3)$:
the four groups of modes, their coefficients and phases, the ordered sets of representatives, the monomials $C_I$ and the spaces $\mathrm{Fk}(\cdot)$ are those of A.6;
the charge of a mode over $y$ is $+1$ for the groups 1, 2 and $-1$ for the groups 3, 4 (A.8);
the half slab there is $\mathrm{hs}=X((0,\frac{1}{2}))/\mathrm{Gmp}$, and its action is $S_{\mathrm{hs}}=S_{(0,\frac{1}{2})}$;
and the entry
\[
\mathrm{Nn}(A_t)_{I,I'}=\int_{\mathrm{hs}}C'_{I'}\cdot\exp\left(-S_{\mathrm{hs}}(A_t)\right)\cdot\Theta_{t+\frac{1}{2}}(C_I)
\]
of the half-slab matrix defined there, $C'_{I'}$ being the ordered monomial of the modes $I'$ of the site plane, is, at $t=0$, the entry of $\mathrm{Ms}_{0,\frac{1}{2}}[\,\cdot\,]$ in the row $I$ (a subset of $\mathrm{Ri}(\frac{1}{2})$) and the column $I'$ (a subset of $\mathrm{Ri}(0)$) (A.7).
In the notation of \cite[Definition~2.27]{QEDII}, the objects $\mathrm{Fk}$, $\mathrm{Fk}'$, $\mathrm{cc}$, $\mathrm{cc}'$, $\mathrm{Wq}(a)$, $\mathrm{Qc}_I$, $\mathrm{Nn}(B)$, $\mathrm{Qs}(B)$, $\mathrm{Ri}(0)$, $\mathrm{Ri}(\frac{1}{2})$, $C_I$, $C'_{I'}$, $\mathrm{hs}$ and $S_{\mathrm{hs}}$ are written $\mathcal{F}$, $\mathcal{F}'$, $\omega$, $\omega'$, $u(a)$, $\mathfrak{q}(I)$, $\mathfrak{N}(B)$, $\widehat{\mathfrak{N}}(B)$, $\mathcal{J}^{\mathrm{s}}(\Lambda_{\mathrm{s}})$, $\mathcal{J}^{\mathrm{l}}(\Lambda_{\mathrm{s}})$, $\zeta_I$, $\zeta'_{I'}$, $\mathbb{X}_{\mathrm{h}}$ and $S_{\mathrm{h}}$.

\emph{Slices.}
Let $n_\mu\ge1$ be an integer, $\mathrm{gv}:=n_\mu e_\mu$, $n$ the triple of the three sides, $\Gamma(n)=\mathrm{Gmp}+\mathbb{Z}\,\mathrm{gv}$, and $A$ a real one-cochain on $\mathbb{Z}^3$ that is periodic under $\Gamma(n)$.
For $t\in\mathbb{Z}$ put
\[
A_t(b):=A(b+te_\mu),\quad b\in\mathrm{Mi}^1(0);\qquad
a_t(b):=A(b+te_\mu),\quad b\in\mathrm{Cu}^1(\tfrac{1}{2});
\]
here $a_t(b)$ is the value on the link oriented along $e_\mu$, and $A_{t+n_\mu}=A_t$, $a_{t+n_\mu}=a_t$.
Let $\operatorname{det}_n(D_A+\boldsymbol{\mu})$ be the determinant on $\mathfrak{S}^{\mathrm{ch}}(1;n)$ (Proposition~46, where $n_\mu\ge1$ is admitted).
Let $\mathrm{Gr}_n$ and $\int_n$ be those of Lemma~A.3\,(F5) (the lattice $\Gamma(n)$, the character that is $\mathrm{chp}$ on $\mathrm{Gmp}$ and $-1$ at $\mathrm{gv}$), and, as in Section~1.5,
\[
S_n(A):=\sum_x\bar{\psi}(x)\cdot\left((D_A+\boldsymbol{\mu})\psi\right)(x)\in\mathrm{Gr}_n,
\]
$x$ ranging over representatives of the sites modulo $\Gamma(n)$.
By Proposition~46\,(e),
\[
\int_n\exp\left(-S_n(A)\right)=\operatorname{det}_n(D_A+\boldsymbol{\mu}).
\]

\begin{named}{Lemma 8.1}[the display of Remark~39.R, for every time length and with each direction as time]
For every integer $n_\mu\ge1$ and every real $\Gamma(n)$-periodic one-cochain $A$,
\[
\operatorname{det}_n(D_A+\boldsymbol{\mu})=\operatorname{tr}_{\mathrm{Fk}'}\prod_{t=n_\mu-1,\dots,0}\left(\mathrm{Qs}(A_{t+1})^{*}\cdot\mathrm{Wq}(a_t)\cdot\mathrm{Qs}(A_t)\right),
\]
with the later times to the left and $A_{n_\mu}=A_0$.
\end{named}

\begin{proof}
Put $\mathrm{po}_k:=\frac{k}{2}$, $k\in\mathbb{Z}$.

\emph{Step 1 (the action).}
In $\mathrm{Gr}_n$,
\begin{equation}
S_n(A)=\sum_{k=1}^{2n_\mu}S_{(\mathrm{po}_{k-1},\mathrm{po}_k)}(A)+\sum_{k=0}^{2n_\mu-1}\mathrm{Xc}_{\mathrm{po}_k}(A),
\tag{8.1}
\end{equation}
where $\mathrm{Xc}_{\mathrm{po}_0}(A)$, an element of $\mathrm{Gr}_{(-\frac{1}{2},\frac{1}{2})}$, is read in $\mathrm{Gr}_n$ as in (F5): a field at a site $x$ with $\mathrm{geo}(x)_\mu\in(-n_\mu,0)$ is minus the field at $x+\mathrm{gv}$.
This is (Bd4) of A.7 for the chain $\mathrm{po}_0<\mathrm{po}_1<\dots<\mathrm{po}_{2n_\mu}=\mathrm{po}_0+n_\mu$, which has $s=2n_\mu$ slabs;
(Bd4) admits every integer $n_\mu\ge1$.

\emph{Step 2 (the trace).}
All terms of (8.1) are even;
so they commute, and $\exp(-S_n(A))$ is the product of their exponentials in any order.
The reading of $\mathrm{Gr}_{(-\frac{1}{2},\frac{1}{2})}$ in $\mathrm{Gr}_n$ is a homomorphism of algebras (generators go to generators or to their negatives).
By (Ap9),
\[
\exp\left(-\mathrm{Xc}_{\mathrm{po}_k}(A)\right)=\mathrm{Ins}_{\mathrm{po}_k}[\mathrm{mt}_k],\qquad
\mathrm{mt}_k:=\mathrm{cc}_{\mathrm{po}_k}\cdot\operatorname{diag}\left(u_I(A)\right)_{I\subset\mathrm{Ri}(\mathrm{po}_k)},
\]
a diagonal matrix, hence even.
The element $\Phi_k:=\exp\left(-S_{(\mathrm{po}_{k-1},\mathrm{po}_k)}(A)\right)$ lies in $\mathrm{Gr}^{\mathrm{ev}}_{(\mathrm{po}_{k-1},\mathrm{po}_k)}$.
By (F5) for the chain $\mathrm{po}_0<\mathrm{po}_1<\dots<\mathrm{po}_{2n_\mu}=\mathrm{po}_0+n_\mu$, which has $s=2n_\mu\ge2$ slabs,
\begin{equation}
\begin{aligned}
\operatorname{det}_n(D_A+\boldsymbol{\mu})
&=\int_n\mathrm{Ins}_{\mathrm{po}_0}[\mathrm{mt}_0]\cdot\Phi_1\cdot\mathrm{Ins}_{\mathrm{po}_1}[\mathrm{mt}_1]\cdot\Phi_2\cdots\mathrm{Ins}_{\mathrm{po}_{2n_\mu-1}}[\mathrm{mt}_{2n_\mu-1}]\cdot\Phi_{2n_\mu}\\
&=\operatorname{tr}_{\mathrm{Fk}(0)}\left(\mathrm{Ft}_{\mathrm{gv}}^{-1}\cdot\mathrm{Ms}_{2n_\mu}\cdot\mathrm{mt}_{2n_\mu-1}\cdot\mathrm{Ms}_{2n_\mu-1}\cdots\mathrm{mt}_1\cdot\mathrm{Ms}_1\cdot\mathrm{mt}_0\right),
\end{aligned}
\tag{8.2}
\end{equation}
where $\mathrm{Ms}_k:=\mathrm{Ms}_{\mathrm{po}_{k-1},\mathrm{po}_k}[\Phi_k]$.

\emph{Step 3 (the factors).}
Let $t$ be an integer, and $v:=te_\mu$.

(3a) For every real $\mathrm{Gmp}$-periodic cochain $A'$,
\[
S_{(t,t+\frac{1}{2})}(A')=\mathrm{Tg}_v\,S_{(0,\frac{1}{2})}(\tau_v^*A').
\]
Indeed, $x\mapsto x+v$ maps $X((0,\frac{1}{2}))$ onto $X((t,t+\frac{1}{2}))$, $(x+v)^{(i)}=x^{(i)}+v$, $M_i(x+v)=M_i(x)$, and $U_{A'}(x+v,x^{(i)}+v)=U_{\tau_v^*A'}(x,x^{(i)})$.
Since $\mathrm{Tg}_v$ is an automorphism, $\exp\left(-S_{(t,t+\frac{1}{2})}(A')\right)=\mathrm{Tg}_v\exp\left(-S_{(0,\frac{1}{2})}(\tau_v^*A')\right)$, and (F4) gives
\begin{equation}
\mathrm{Ms}_{t,t+\frac{1}{2}}\left[\exp\left(-S_{(t,t+\frac{1}{2})}(A')\right)\right]=\mathrm{Ft}_v\cdot\mathrm{Nn}\left((\tau_v^*A')|_{\mathrm{Mi}^1(0)}\right)\cdot\mathrm{Ft}_v^{-1},
\tag{8.3}
\end{equation}
where on the right $\mathrm{Ft}_v^{-1}:\mathrm{Fk}(t)\to\mathrm{Fk}'$ and $\mathrm{Ft}_v:\mathrm{Fk}\to\mathrm{Fk}(t+\frac{1}{2})$.
For $A'=A$ one has $(\tau_v^*A)(b)=A(b+v)=A_t(b)$;
so
\[
\mathrm{Ms}_{2t+1}=\mathrm{Ft}_v\cdot\mathrm{Nn}(A_t)\cdot\mathrm{Ft}_v^{-1}.
\]

(3b) Let $\mathrm{po}:=t$, $\mathrm{po}':=t+\frac{1}{2}$ and $\mathrm{po}'':=t+1$;
the vector $\tau$ of (F3) is then $\tau=e_\mu$.
Step (4) of the proof of (Ap10) is an identity for every real $\mathrm{Gmp}$-periodic cochain:
\[
S_{(\mathrm{po}',\mathrm{po}'')}(A)=\Theta_{\mathrm{po}'}\left(S_{(\mathrm{po},\mathrm{po}')}(\theta_{\mathrm{po}'}^*A)\right).
\]
Since $\Theta_{\mathrm{po}'}$ is an antilinear anti-automorphism and the coefficients $\frac{1}{j!}$ of the exponential series are real,
$\exp\left(-S_{(\mathrm{po}',\mathrm{po}'')}(A)\right)=\Theta_{\mathrm{po}'}\left(\exp\left(-S_{(\mathrm{po},\mathrm{po}')}(\theta_{\mathrm{po}'}^*A)\right)\right)$.
By (F3),
\[
\mathrm{Ms}_{2t+2}=\mathrm{Ms}_{\mathrm{po}',\mathrm{po}''}\left[\Theta_{\mathrm{po}'}\left(\exp\left(-S_{(\mathrm{po},\mathrm{po}')}(\theta_{\mathrm{po}'}^*A)\right)\right)\right]
=\mathrm{Ft}_{e_\mu}\cdot\left(\mathrm{Ms}_{\mathrm{po},\mathrm{po}'}\left[\exp\left(-S_{(\mathrm{po},\mathrm{po}')}(\theta_{\mathrm{po}'}^*A)\right)\right]\right)^{*}.
\]
By (Ap2), $\theta_{\mathrm{po}'}\boldsymbol{\sigma}=\boldsymbol{\sigma}+e_\mu$ with $\epsilon(\theta_{\mathrm{po}'},\boldsymbol{\sigma})=+1$ for $\boldsymbol{\sigma}\in\mathrm{Mi}(t)$;
so $(\tau_v^*\theta_{\mathrm{po}'}^*A)(b)=(\theta_{\mathrm{po}'}^*A)(b+v)=A(b+v+e_\mu)=A_{t+1}(b)$ for $b\in\mathrm{Mi}^1(0)$.
By (8.3) with $A'=\theta_{\mathrm{po}'}^*A$, and because the $\mathrm{Ft}$ are unitary,
\[
\mathrm{Ms}_{2t+2}=\mathrm{Ft}_{e_\mu}\cdot\left(\mathrm{Ft}_v\cdot\mathrm{Nn}(A_{t+1})\cdot\mathrm{Ft}_v^{-1}\right)^{*}
=\mathrm{Ft}_{e_\mu}\cdot\mathrm{Ft}_v\cdot\mathrm{Nn}(A_{t+1})^{*}\cdot\mathrm{Ft}_v^{-1}
=\mathrm{Ft}_{v+e_\mu}\cdot\mathrm{Nn}(A_{t+1})^{*}\cdot\mathrm{Ft}_v^{-1},
\]
with $\mathrm{Ft}_v^{-1}:\mathrm{Fk}(t+\frac{1}{2})\to\mathrm{Fk}$, $\mathrm{Nn}(A_{t+1})^{*}:\mathrm{Fk}\to\mathrm{Fk}'$ and $\mathrm{Ft}_{v+e_\mu}:\mathrm{Fk}'\to\mathrm{Fk}(t+1)$.

(3c) At the site plane $t$ one has $u_I=1$ (A.6, A.8);
since $\mathrm{Ft}_v\cdot\mathrm{cc}_{\mathrm{po}}=\mathrm{cc}_{\mathrm{po}+v_\mu}\cdot\mathrm{Ft}_v$ (A.6),
\[
\mathrm{mt}_{2t}=\mathrm{cc}_t=\mathrm{Ft}_v\cdot\mathrm{cc}'\cdot\mathrm{Ft}_v^{-1}.
\]
At the link plane $t+\frac{1}{2}$ one has, for $I\subset\mathrm{Ri}(\frac{1}{2})$, $\mathrm{Qc}_{I+v}(b+v)=\mathrm{Qc}_I(b)$ (the group of a mode and the link under its bond are transported by $i\mapsto i+v$);
so
\[
u_{I+v}(A)=\exp\left(iq\sum_{b\in\mathrm{Cu}^1(\frac{1}{2})}\mathrm{Qc}_I(b)\cdot A(b+v)\right)=\mathrm{Wq}(a_t)_{I,I}.
\]
Hence
\[
\mathrm{mt}_{2t+1}=\mathrm{cc}_{t+\frac{1}{2}}\cdot\operatorname{diag}\left(u_{I+v}(A)\right)=\mathrm{Ft}_v\cdot\mathrm{cc}\cdot\mathrm{Wq}(a_t)\cdot\mathrm{Ft}_v^{-1}.
\]

\emph{Step 4 (telescoping).}
By Step 3, for $0\le t\le n_\mu-1$,
\[
\mathrm{Ms}_{2t+2}\cdot\mathrm{mt}_{2t+1}\cdot\mathrm{Ms}_{2t+1}\cdot\mathrm{mt}_{2t}
=\mathrm{Ft}_{(t+1)e_\mu}\cdot\left[\mathrm{Nn}(A_{t+1})^{*}\cdot\mathrm{cc}\cdot\mathrm{Wq}(a_t)\cdot\mathrm{Nn}(A_t)\cdot\mathrm{cc}'\right]\cdot\mathrm{Ft}_{te_\mu}^{-1},
\]
because $\mathrm{Ft}_{te_\mu}^{-1}\cdot\mathrm{Ft}_{te_\mu}=1$ between the four factors.
In the ordered product over $t=n_\mu-1,\dots,0$ the factors $\mathrm{Ft}_{te_\mu}^{-1}\cdot\mathrm{Ft}_{te_\mu}$ between consecutive brackets cancel;
at the right end stands $\mathrm{Ft}_0=1$, and at the left end stands $\mathrm{Ft}_{n_\mu e_\mu}=\mathrm{Ft}_{\mathrm{gv}}$, which cancels against $\mathrm{Ft}_{\mathrm{gv}}^{-1}$ in (8.2).
So
\[
\operatorname{det}_n(D_A+\boldsymbol{\mu})=\operatorname{tr}_{\mathrm{Fk}'}\prod_{t=n_\mu-1,\dots,0}\left[\mathrm{Nn}(A_{t+1})^{*}\cdot\mathrm{cc}\cdot\mathrm{Wq}(a_t)\cdot\mathrm{Nn}(A_t)\cdot\mathrm{cc}'\right].
\]

\emph{Step 5 (symmetric form).}
The matrices $\mathrm{cc}$, $\mathrm{cc}^{\frac{1}{2}}$ and $\mathrm{Wq}(a_t)$ are diagonal, so that $\mathrm{cc}\cdot\mathrm{Wq}(a_t)=\mathrm{cc}^{\frac{1}{2}}\cdot\mathrm{Wq}(a_t)\cdot\mathrm{cc}^{\frac{1}{2}}$.
Write $\mathrm{cc}'=\mathrm{cc}'^{\frac{1}{2}}\cdot\mathrm{cc}'^{\frac{1}{2}}$ in every bracket, and move the factor $\mathrm{cc}'^{\frac{1}{2}}$ at the right end of the product (bracket $t=0$) to its left end by the cyclicity of the trace.
Then $\mathrm{cc}'^{\frac{1}{2}}$ stands to the left of each $\mathrm{Nn}(A_{t+1})^{*}$:
for $t\le n_\mu-2$ it comes from the bracket $t+1$, and for $t=n_\mu-1$ it is the factor that was moved.
So the product becomes
\[
\prod_{t=n_\mu-1,\dots,0}\left[\mathrm{cc}'^{\frac{1}{2}}\cdot\mathrm{Nn}(A_{t+1})^{*}\cdot\mathrm{cc}^{\frac{1}{2}}\cdot\mathrm{Wq}(a_t)\cdot\mathrm{cc}^{\frac{1}{2}}\cdot\mathrm{Nn}(A_t)\cdot\mathrm{cc}'^{\frac{1}{2}}\right],
\]
and $\mathrm{cc}'^{\frac{1}{2}}\cdot\mathrm{Nn}(B)^{*}\cdot\mathrm{cc}^{\frac{1}{2}}=\mathrm{Qs}(B)^{*}$, the diagonal matrices being real.
\end{proof}

\subsubsection*{8.2. Slice data}

\begin{named}{Corollary 8.2}[Remark~39.R up to its paragraph on closed cells]
\begin{itemize}
\item[(i)] $\#|\mathrm{Ri}(0)|=\#|\mathrm{Ri}(\frac{1}{2})|=64v_s$, and $\dim\mathrm{Fk}'=\dim\mathrm{Fk}=2^{64v_s}$.
\item[(ii)] \emph{(The display.)}
Let $\mathrm{Uf}:\mathrm{Fk}\to\mathrm{Fk}'$ be unitary, and put $\mathrm{Qf}(B):=\mathrm{Qs}(B)\cdot\mathrm{Uf}\in\operatorname{End}\mathrm{Fk}$, $s:=1$, $\mathrm{Gmf}:=\mathrm{Wq}$.
For every integer $n_\mu\ge1$ and $A$ as in Lemma~8.1,
\[
\operatorname{det}_n(D_A+\boldsymbol{\mu})=\operatorname{tr}_{\mathrm{Fk}}\prod_{t=n_\mu-1,\dots,0}\left[s\cdot\mathrm{Qf}(A_{t+1})^{*}\cdot\mathrm{Gmf}(a_t)\cdot\mathrm{Qf}(A_t)\cdot s\right].
\]
\item[(iii)] The entries of $\mathrm{Nn}(B)$, $\mathrm{Qs}(B)$ and $\mathrm{Qf}(B)$ are polynomials in the numbers $\exp(\pm iqB(b))$, $b\in\mathrm{Mi}^1(0)$.
So $\mathrm{Qf}$ is bounded and continuous on $\mathbb{R}^{\mathrm{Mi}^1(0)}$, and it depends on $B$ through the link phases only.
\item[(iv)] \emph{(The charges; $\mathrm{Gmf}$.)}
The operators $Q^{\mathrm{ch}}_y:=\operatorname{diag}\left(\mathrm{Qc}_I(y)\right)_{I\subset\mathrm{Ri}(\frac{1}{2})}$, $y$ a site of the slice, are commuting self-adjoint operators on $\mathrm{Fk}$ with spectrum in $\{-32,\dots,32\}$;
$\mathrm{Gmf}(a)=\exp\left(iq\sum_ya_y\cdot Q^{\mathrm{ch}}_y\right)$;
$\mathrm{Gmf}(a)$ is unitary, $\mathrm{Gmf}(a)\cdot\mathrm{Gmf}(a')=\mathrm{Gmf}(a+a')$, $\mathrm{Gmf}(a)^{*}=\mathrm{Gmf}(-a)$, and $\mathrm{Gmf}(a)=1$ if all $a_y$ lie in $r_q\mathbb{Z}$.
\item[(v)] \emph{(With what the data are formed; slice data.)}
The data $\mathrm{Fk}$, $s$, $\mathrm{Qf}$ and $Q^{\mathrm{ch}}_y$ are formed with $(\mu;n_\nu,n_{\mathfrak{l}})$, the parameters $(q,\rho,M_{\mathrm{ch}})$ and $\mathrm{Uf}$;
they do not contain $n_\mu$.
So they are slice data for the cross-section $(n_\nu,n_{\mathfrak{l}})$ in the sense of the definition of slice data in \cite[Section~11.2.1]{QEDI}, which is the sense of \cite[Lemma~21.4]{QEDI} and of \cite[Section~10.2.1]{QEDI}, for the functions $A\mapsto\operatorname{det}_n(D_A+\boldsymbol{\mu})$ in the place of $\mathrm{Dt}(qA)$:
the first condition, (SD1) of \cite[Section~11.2.1]{QEDI}, that is, property (a) of \cite[Lemma~21.4]{QEDI}, is (iii) together with (iv);
the second condition, (SD2) of \cite[Section~11.2.1]{QEDI}, that is, property (b) of \cite[Lemma~21.4]{QEDI}, which is required for all $n_0\ge2$, is (ii).
Of $\mathrm{Uf}$ only the relations $\mathrm{Uf}\cdot\mathrm{Uf}^{*}=1_{\mathrm{Fk}'}$ and $\mathrm{Uf}^{*}\cdot\mathrm{Uf}=1_{\mathrm{Fk}}$ enter;
neither condition involves a gauge covariance of $\mathrm{Qf}$ such as the one stated for the half-slab matrix in \cite[Lemma~5.35(b)]{QEDII}.
\end{itemize}
\end{named}

\begin{proof}
(i) By (Bd1), eight bonds cross the plane over each photon site of a site plane and over each cut link of a link plane;
each of them carries eight modes (A.6);
the slice has $v_s$ sites.
A set with $64v_s$ elements has $2^{64v_s}$ subsets.
In particular, unitaries $\mathrm{Uf}$ exist.

(ii) One has
\[
\mathrm{Qf}(A_{t+1})^{*}\cdot\mathrm{Wq}(a_t)\cdot\mathrm{Qf}(A_t)=\mathrm{Uf}^{*}\cdot\left[\mathrm{Qs}(A_{t+1})^{*}\cdot\mathrm{Wq}(a_t)\cdot\mathrm{Qs}(A_t)\right]\cdot\mathrm{Uf}.
\]
In the ordered product the factors $\mathrm{Uf}\cdot\mathrm{Uf}^{*}=1_{\mathrm{Fk}'}$ between consecutive brackets cancel;
so the product is $\mathrm{Uf}^{*}\cdot\mathrm{Yp}\cdot\mathrm{Uf}$, where $\mathrm{Yp}$ is the product of Lemma~8.1, and
\[
\operatorname{tr}_{\mathrm{Fk}}\left(\mathrm{Uf}^{*}\cdot\mathrm{Yp}\cdot\mathrm{Uf}\right)=\operatorname{tr}_{\mathrm{Fk}'}\left(\mathrm{Yp}\cdot\mathrm{Uf}\cdot\mathrm{Uf}^{*}\right)=\operatorname{tr}_{\mathrm{Fk}'}\mathrm{Yp}.
\]

(iii) By (Bd2), the coefficients of $S_{(0,\frac{1}{2})}(B)$ are polynomials in the phases.
The element $\exp\left(-S_{(0,\frac{1}{2})}(B)\right)$ is a polynomial in the finitely many generators of $\mathrm{Gr}_{(0,\frac{1}{2})}$, whose coefficients are polynomials in the coefficients of $S_{(0,\frac{1}{2})}(B)$;
the elements $C^{(0)}_{I'}$ and $\Theta_{\frac{1}{2}}\left(C^{(\frac{1}{2})}_I\right)$ have constant coefficients;
and $\int_{(0,\frac{1}{2})}$ picks the coefficient of $\mathrm{top}$.
The phases have modulus one.

(iv) Over $y$ there are $64$ modes, $32$ of the groups 1, 2 and $32$ of the groups 3, 4;
so $\mathrm{Qc}_I(y)$ is an integer in $[-32,32]$.
Diagonal real matrices commute and are self-adjoint.
The operator $\exp\left(iq\sum_ya_y\cdot Q^{\mathrm{ch}}_y\right)$ is the diagonal matrix with the entries $\exp\left(iq\sum_ya_y\cdot\mathrm{Qc}_I(y)\right)$;
this is $\mathrm{Wq}(a)$.
With the real number $\langle a,\mathrm{Qc}_I\rangle:=\sum_ya_y\cdot\mathrm{Qc}_I(y)$ one has
\begin{gather*}
\left|\exp\left(iq\langle a,\mathrm{Qc}_I\rangle\right)\right|=1,\qquad
\exp\left(iq\langle a,\mathrm{Qc}_I\rangle\right)\cdot\exp\left(iq\langle a',\mathrm{Qc}_I\rangle\right)=\exp\left(iq\langle a+a',\mathrm{Qc}_I\rangle\right),\\
\overline{\exp\left(iq\langle a,\mathrm{Qc}_I\rangle\right)}=\exp\left(iq\langle -a,\mathrm{Qc}_I\rangle\right);
\end{gather*}
and if $a_y=r_q\cdot j_y$ with integers $j_y$, then $q\langle a,\mathrm{Qc}_I\rangle=2\pi\sum_yj_y\cdot\mathrm{Qc}_I(y)$ lies in $2\pi\mathbb{Z}$, because $q\cdot r_q=2\pi$.

(v) See the paragraph ``Objects'' of 8.1.
\end{proof}

\subsubsection*{8.3. The paragraph ``Closed cells'' of the Remark}

\begin{proof}[Proof of the paragraph ``Closed cells'' of Remark~39.R]
The assertions of the paragraph follow from Section~9.2, Proposition~9.3 and Remark~9.4 of Chapter~39r;
Remark~9.4 of Chapter~39r is the form for closed cells of \cite[Proposition~21.22]{QEDI} and of \cite[Proposition~21.24(a)]{QEDI}.
The last assertion of the paragraph, which concerns the properties of the fermion that enter the proof, follows from Remark~5.4 of Chapter~39r.
\end{proof}

\subsection*{Appendix A. The slab operators of one coordinate direction: proof of Theorem SO for planes of both kinds}

This appendix gives the construction of \cite[Sections~5.1--5.7]{QEDII} and the proofs of \cite[Lemmas~5.3, 5.14, 5.4, 5.22, 5.24 and~5.25, Theorem~5.29 and Corollary~5.30]{QEDII} in the case needed in Sections~3--6:
the family $\ell=y_\mu$, the cross-section $\mathrm{Xs}_\mu(n)$, even integrands, translations along $e_\mu$, and the closing vector $n_\mu e_\mu$.
The positions $\mathrm{po},\mathrm{po}',\dots$ of the planes are arbitrary elements of $\tfrac12\mathbb{Z}$.
For each step it is stated what the step amounts to at a site plane ($\mathrm{po}\in\mathbb{Z}$) and at a link plane ($\mathrm{po}\in\mathbb{Z}+\tfrac12$).

\subsubsection*{A.0. Setting and notation}

The direction $\mu$ is fixed, $\nu$ and $\mathfrak{l}$ are the two other directions, and $n_\nu,n_{\mathfrak{l}}\ge4$.
Put
\[
\mathrm{Gmp}:=\Gamma_\mu(n)=n_\nu\mathbb{Z}\,e_\nu+n_{\mathfrak{l}}\mathbb{Z}\,e_{\mathfrak{l}},\qquad
\mathrm{chp}\left(a\,n_\nu e_\nu+b\,n_{\mathfrak{l}}e_{\mathfrak{l}}\right):=(-1)^{a+b}\quad(a,b\in\mathbb{Z}),\qquad
v_s:=n_\nu\cdot n_{\mathfrak{l}},
\]
and, for $\mathrm{po}\in\tfrac12\mathbb{Z}$,
\[
\theta_{\mathrm{po}}:=\theta_{\mu,\mathrm{po}},\qquad
\Theta_{\mathrm{po}}:=\Theta_{\mu,\mathrm{po}},\qquad
\mathrm{po}^{+}:=\mathrm{po}+\tfrac12,\qquad
\mathrm{po}^{-}:=\mathrm{po}-\tfrac12.
\]

\emph{Notation.}
Cells are denoted by $\boldsymbol{\sigma}$ and their dimension by $\mathrm{dg}$;
$b$ is a link, $p$ a plaquette and $\mathrm{qb}$ a cube.
The position of a mirror, the dimension of a cell and the integer attached to a cut plaquette, which are denoted by $c$, $k$ and $k$ in \cite[Section~5]{QEDII}, are denoted here by $\mathrm{po}$, $\mathrm{dg}$ and $\mathrm{kz}$;
$F(\omega)$ and $F_p(\omega)$ are curvatures in lattice units.
The letters $\omega$, $\xi$, $\boldsymbol{\zeta}$, $\mathrm{Ee}_i$, $C_I$ of this appendix, and $\mathbf{K}$ in Lemma~2.4, stand for objects of \cite[Definition~2.27, Proposition~3.20 and Section~5]{QEDII};
each of them is defined where it first occurs, and none of them has the meaning which the same letter carries in the other parts of this paper.
The letter $\mathrm{mt}$ denotes a matrix on one of the spaces $\mathrm{Fk}(\cdot)$ of A.6.

\emph{Classical facts.}
\begin{itemize}
\item[(K1)] (Poisson summation.)
For real $y$ and $t$,
\[
\sum_{j\in\mathbb{Z}}\exp\left(-\frac{\left(y+r_q\cdot j\right)^2}{2}+i\,j\,t\right)
=\frac{(2\pi)^{\frac12}}{r_q}\sum_{\mathrm{kz}\in\mathbb{Z}}\exp\left(-\frac{\kappa^2}{2}-i\,\kappa\,y\right),
\qquad
\kappa:=\frac{t-2\pi\,\mathrm{kz}}{r_q}.
\]
Indeed, $f(x):=\exp\left(-\frac{(y+r_qx)^2}{2}+i\,t\,x\right)$ is a Schwartz function of the real variable $x$, and the Poisson summation formula gives
\[
\sum_{j\in\mathbb{Z}}f(j)=\sum_{\mathrm{kz}\in\mathbb{Z}}\int\mathrm{d}x\,f(x)\,e^{-2\pi i\,\mathrm{kz}\,x}.
\]
The substitution $u=y+r_qx$ turns the integral into
\[
r_q^{-1}\,e^{-i\kappa y}\int\mathrm{d}u\,\exp\left(-\frac{u^2}{2}+i\,\kappa\,u\right),
\]
and the last integral equals $(2\pi)^{\frac12}\exp\left(-\frac{\kappa^2}{2}\right)$.
\item[(K2)] For $\varepsilon>0$,
\[
\boldsymbol{\Theta}_{\varepsilon}=\sup_{y\in\mathbb{R}}\sum_{j\in\mathbb{Z}}\exp\left(-\varepsilon\left(y+r_q\cdot j\right)^2\right)\le2+r_q^{-1}\cdot\left(\frac{\pi}{\varepsilon}\right)^{\frac12}
\]
(Lemma~2.4\,(d)).
\item[(K3)] Let $(\Xi,\mathrm{mq})$ and $(\Xi',\mathrm{mq}')$ be $\sigma$-finite measure spaces whose spaces $L^2$ are separable.
A kernel $\mathrm{kn}\in L^2(\mathrm{mq}'\times\mathrm{mq})$ defines a Hilbert--Schmidt operator $L^2(\mathrm{mq})\to L^2(\mathrm{mq}')$, whose Hilbert--Schmidt norm is the $L^2$-norm of the kernel.
The adjoint operator has the kernel $\overline{\mathrm{kn}(\xi',\xi)}$, regarded as a function of $(\xi,\xi')$.
If $\mathrm{kn}_1$ is the kernel of such an operator $L^2(\mathrm{mq})\to L^2(\mathrm{mq}')$ and $\mathrm{kn}_2$ the kernel of such an operator $L^2(\mathrm{mq}')\to L^2(\mathrm{mq}'')$, then the product of the two operators has the kernel
\[
\int\mathrm{d}\,\mathrm{mq}'(\xi')\,\mathrm{kn}_2(\xi'',\xi')\,\mathrm{kn}_1(\xi',\xi).
\]
If $X\colon L^2(\mathrm{mq})\to L^2(\mathrm{mq}')$ and $Y\colon L^2(\mathrm{mq}')\to L^2(\mathrm{mq})$ are Hilbert--Schmidt operators with the kernels $\mathrm{kn}_X$ and $\mathrm{kn}_Y$, then $YX$ is of trace class, and
\[
\operatorname{Tr}(YX)=\int\mathrm{d}\,\mathrm{mq}(\xi)\int\mathrm{d}\,\mathrm{mq}'(\xi')\,\mathrm{kn}_Y(\xi,\xi')\,\mathrm{kn}_X(\xi',\xi),
\]
the integral being absolutely convergent \cite[Theorems~VI.22 and~VI.23]{ReedSimon1980}.
\end{itemize}

\subsubsection*{A.1. Cells of the cylinder, relative to the planes of the family}

Let $\mathrm{dg}\in\{0,1,2,3\}$.
A cell of $\mathbb{Z}^3$ of dimension $\mathrm{dg}$, or $\mathrm{dg}$-cell, is
\[
\boldsymbol{\sigma}=(c;\mathrm{Md}):=c+\sum_{j=1}^{\mathrm{dg}}[0,1]\,e_{d_j},
\qquad c\in\mathbb{Z}^3,\quad\mathrm{Md}=\{d_1<\dots<d_{\mathrm{dg}}\}\subset\{1,2,3\},
\]
with the orientation $e_{d_1}\wedge\dots\wedge e_{d_{\mathrm{dg}}}$.
Its faces are the $2\,\mathrm{dg}$ cells $(c+e_{d_j};\mathrm{Md}\setminus\{d_j\})$ and $(c;\mathrm{Md}\setminus\{d_j\})$, $j=1,\dots,\mathrm{dg}$.
The incidence numbers are
\[
\left[\boldsymbol{\sigma}:(c+e_{d_j};\mathrm{Md}\setminus\{d_j\})\right]:=(-1)^{j-1},\qquad
\left[\boldsymbol{\sigma}:(c;\mathrm{Md}\setminus\{d_j\})\right]:=-(-1)^{j-1},
\]
and $[\boldsymbol{\sigma}:\boldsymbol{\sigma}']:=0$ if $\boldsymbol{\sigma}'$ is not a face of $\boldsymbol{\sigma}$.
With them the coboundary reads $(du)(\boldsymbol{\sigma})=\sum_{\boldsymbol{\sigma}'}[\boldsymbol{\sigma}:\boldsymbol{\sigma}']\cdot u(\boldsymbol{\sigma}')$;
for $\mathrm{dg}=1,2,3$ these are the formulas of Section~I.1.1, in lattice units.
\begin{itemize}
\item[(I1)] $\sum_{\boldsymbol{\sigma}'}[\boldsymbol{\sigma}:\boldsymbol{\sigma}']\cdot[\boldsymbol{\sigma}':\boldsymbol{\sigma}'']=0$, that is, $dd=0$.
\item[(I2)] For $g_{\mathrm{w}}\in W_{\mathrm{af}}$,
\[
\epsilon(g_{\mathrm{w}},\boldsymbol{\sigma})\cdot[g_{\mathrm{w}}\boldsymbol{\sigma}:g_{\mathrm{w}}\boldsymbol{\sigma}']=[\boldsymbol{\sigma}:\boldsymbol{\sigma}']\cdot\epsilon(g_{\mathrm{w}},\boldsymbol{\sigma}').
\]
This is the identity $g_{\mathrm{w}}^{*}d=d\,g_{\mathrm{w}}^{*}$ (Section~I.4.7 and \cite[Definition~2.1]{QEDII}), evaluated on the indicator function of $g_{\mathrm{w}}\boldsymbol{\sigma}'$.
If $g_{\mathrm{w}}$ is a translation, then $\epsilon(g_{\mathrm{w}},\boldsymbol{\sigma})=1$.
For the reflection $\theta_{\mathrm{po}}$ one has $\epsilon(\theta_{\mathrm{po}},\boldsymbol{\sigma})=-1$ if $\mu\in\mathrm{Md}$ and $\epsilon(\theta_{\mathrm{po}},\boldsymbol{\sigma})=+1$ if $\mu\notin\mathrm{Md}$, because the linear part $r_\mu$ reverses $e_\mu$ and fixes the two other unit vectors.
\end{itemize}

\emph{The cylinder.}
All cells, cochains and fermion sites of this appendix are taken modulo $\mathrm{Gmp}$.
A cell of $\mathrm{Cy}:=\mathbb{Z}^3/\mathrm{Gmp}$ is a $\mathrm{Gmp}$-orbit of cells of $\mathbb{Z}^3$, and a cochain on a set of cells of $\mathrm{Cy}$ is a $\mathrm{Gmp}$-periodic function.
The coordinate $y_\mu$, the reflections $\theta_{\mathrm{po}}$, which satisfy $\theta_{\mathrm{po}}(y+v)=\theta_{\mathrm{po}}(y)+v$ for $v\in\mathrm{Gmp}$ by (M4) of Section~1.2, and the translations $\tau_v$, $v\in\mathbb{Z}e_\mu$, act on orbits.
Every non-zero element of $\mathrm{Gmp}$ has sup norm at least $4$.
Hence two different faces of one cell, or of two cells with a common face, never lie in the same orbit;
so $[\boldsymbol{\sigma}:\boldsymbol{\sigma}']$ is defined for orbits, and (I1), (I2) hold for them.
For every bounded set $J\subset\mathbb{R}$ the set of the cells $(c;\mathrm{Md})$ of $\mathrm{Cy}$ with $c_\mu\in J$ is finite.

\emph{Kinds of cells.}
Let $\mathrm{po}<\mathrm{po}'$ be elements of $\tfrac12\mathbb{Z}$ and $\mathrm{dg}\in\{0,1,2,3\}$.
Put
\begin{align*}
\mathrm{Mi}^{\mathrm{dg}}(\mathrm{po})&:=\left\{(c;\mathrm{Md}):\ \#|\mathrm{Md}|=\mathrm{dg},\ \mu\notin\mathrm{Md},\ c_\mu=\mathrm{po}\right\},\\
\mathrm{Cu}^{\mathrm{dg}}(\mathrm{po})&:=\left\{(c;\mathrm{Md}):\ \#|\mathrm{Md}|=\mathrm{dg},\ \mu\in\mathrm{Md},\ c_\mu=\mathrm{po}^{-}\right\}.
\end{align*}
Thus $\mathrm{Mi}^{\mathrm{dg}}(\mathrm{po})$ is the set of the $\mathrm{dg}$-cells contained in the plane $\{y_\mu=\mathrm{po}\}$;
it is empty if $\mathrm{po}$ is a link plane.
The set $\mathrm{Cu}^{\mathrm{dg}}(\mathrm{po})$ consists of the $\mathrm{dg}$-cells cut by this plane, that is, of those with $\mu\in\mathrm{Md}$ and $c_\mu<\mathrm{po}<c_\mu+1$;
it is empty if $\mathrm{po}$ is a site plane, and $\mathrm{Cu}^{0}(\mathrm{po})$ is always empty.
Let $\mathrm{Sb}^{\mathrm{dg}}(\mathrm{po},\mathrm{po}')$ be the set of the $\mathrm{dg}$-cells that are contained in $\{\mathrm{po}\le y_\mu\le\mathrm{po}'\}$ and in neither of the two planes $\{y_\mu=\mathrm{po}\}$, $\{y_\mu=\mathrm{po}'\}$;
these are the slab cells, and they may touch the two planes.
The unions over $\mathrm{dg}$ are denoted by $\mathrm{Mi}(\mathrm{po})$, $\mathrm{Cu}(\mathrm{po})$ and $\mathrm{Sb}(\mathrm{po},\mathrm{po}')$.
The sets of the links and of the plaquettes contained in the closed slab are
\[
\mathrm{Ln}[\mathrm{po},\mathrm{po}']:=\mathrm{Mi}^1(\mathrm{po})\cup\mathrm{Sb}^1(\mathrm{po},\mathrm{po}')\cup\mathrm{Mi}^1(\mathrm{po}'),\qquad
\mathrm{Pq}[\mathrm{po},\mathrm{po}']:=\mathrm{Mi}^2(\mathrm{po})\cup\mathrm{Sb}^2(\mathrm{po},\mathrm{po}')\cup\mathrm{Mi}^2(\mathrm{po}').
\]
(On the cell $(c;\mathrm{Md})$ the coordinate $y_\mu$ takes the values of $[c_\mu,c_\mu]$ if $\mu\notin\mathrm{Md}$ and of $[c_\mu,c_\mu+1]$ if $\mu\in\mathrm{Md}$;
the only element of $\tfrac12\mathbb{Z}$ in the open interval $(c_\mu,c_\mu+1)$ is $c_\mu+\tfrac12$.)
\begin{itemize}
\item[(Ap1)] (Cut cells.)
Let $\boldsymbol{\sigma}=(c;\mathrm{Md})\in\mathrm{Cu}^{\mathrm{dg}}(\mathrm{po})$, and let $j$ be the index with $d_j=\mu$.
Put
\[
\boldsymbol{\sigma}^{+}:=(c+e_\mu;\mathrm{Md}\setminus\{\mu\})\in\mathrm{Mi}^{\mathrm{dg}-1}(\mathrm{po}^{+}),\qquad
\boldsymbol{\sigma}^{-}:=(c;\mathrm{Md}\setminus\{\mu\})\in\mathrm{Mi}^{\mathrm{dg}-1}(\mathrm{po}^{-}),
\]
\[
\mathrm{sn}(\boldsymbol{\sigma}):=[\boldsymbol{\sigma}:\boldsymbol{\sigma}^{+}]=(-1)^{j-1}=-[\boldsymbol{\sigma}:\boldsymbol{\sigma}^{-}].
\]
The other faces of $\boldsymbol{\sigma}$ contain the direction $\mu$, and the coordinate $\mu$ of their base points is $\mathrm{po}^{-}$;
so they belong to $\mathrm{Cu}^{\mathrm{dg}-1}(\mathrm{po})$.
Moreover $\theta_{\mathrm{po}}\boldsymbol{\sigma}=\boldsymbol{\sigma}$ with $\epsilon(\theta_{\mathrm{po}},\boldsymbol{\sigma})=-1$, and $\theta_{\mathrm{po}}\boldsymbol{\sigma}^{+}=\boldsymbol{\sigma}^{-}$ with $\epsilon(\theta_{\mathrm{po}},\boldsymbol{\sigma}^{+})=+1$:
the reflection $\theta_{\mathrm{po}}$ maps the point with $y_\mu=\mathrm{po}^{-}+t$, $0\le t\le1$, to the point with $y_\mu=\mathrm{po}^{-}+(1-t)$ and the same other coordinates.
\item[(Ap2)] (Cells in planes.)
Let $\boldsymbol{\sigma}\in\mathrm{Mi}(\mathrm{po}'')$.
Then for every $\mathrm{po}\in\tfrac12\mathbb{Z}$
\[
\theta_{\mathrm{po}}\boldsymbol{\sigma}=\boldsymbol{\sigma}+2\left(\mathrm{po}-\mathrm{po}''\right)e_\mu,\qquad\epsilon(\theta_{\mathrm{po}},\boldsymbol{\sigma})=+1.
\]
In particular, $\theta_{\mathrm{po}}$ fixes the cells of $\mathrm{Mi}(\mathrm{po})$.
\item[(Ap3)] (The coboundary at a plane.)
For $\boldsymbol{\sigma}\in\mathrm{Cu}^{\mathrm{dg}}(\mathrm{po})$ and a $(\mathrm{dg}-1)$-cochain $u$ on the faces of $\boldsymbol{\sigma}$,
\[
(du)(\boldsymbol{\sigma})=\mathrm{sn}(\boldsymbol{\sigma})\cdot\left(u(\boldsymbol{\sigma}^{+})-u(\boldsymbol{\sigma}^{-})\right)+(d_{\mathrm{po}}u)(\boldsymbol{\sigma}),\qquad
(d_{\mathrm{po}}u)(\boldsymbol{\sigma}):=\sum_{\boldsymbol{\sigma}'\in\mathrm{Cu}^{\mathrm{dg}-1}(\mathrm{po})}[\boldsymbol{\sigma}:\boldsymbol{\sigma}']\cdot u(\boldsymbol{\sigma}').
\]
For $p\in\mathrm{Cu}^2(\mathrm{po})$ and $\mathrm{qb}\in\mathrm{Cu}^3(\mathrm{po})$ put
\[
\alpha_p(A):=\mathrm{sn}(p)\cdot A(p^{+}),\qquad
a_{\mathrm{qb}}(\mathfrak{z}):=\mathrm{sn}(\mathrm{qb})\cdot\mathfrak{z}(\mathrm{qb}^{+}).
\]
By (Ap1), $(\theta_{\mathrm{po}}^{*}u)(\boldsymbol{\sigma}^{+})=u(\boldsymbol{\sigma}^{-})$.
Consequently
\[
(dA)(p)=\alpha_p(A)-\alpha_p(\theta_{\mathrm{po}}^{*}A)+(d_{\mathrm{po}}A)(p),\qquad
(d\mathfrak{z})(\mathrm{qb})=a_{\mathrm{qb}}(\mathfrak{z})-a_{\mathrm{qb}}(\theta_{\mathrm{po}}^{*}\mathfrak{z})+(d_{\mathrm{po}}\mathfrak{z})(\mathrm{qb}).
\]
\item[(Ap4)] $d_{\mathrm{po}}d_{\mathrm{po}}=0$ as a map $\mathbb{R}^{\mathrm{Cu}^1(\mathrm{po})}\to\mathbb{R}^{\mathrm{Cu}^3(\mathrm{po})}$.
Indeed, a plaquette that has the cut link $b=(c;\{\mu\})$ as a face is $(c';\{\mu,\nu'\})$ with a direction $\nu'\ne\mu$ and $c'=c$ or $c'=c-e_{\nu'}$;
hence $c'_\mu=\mathrm{po}^{-}$, and the plaquette is cut.
Therefore, for $\mathrm{qb}\in\mathrm{Cu}^3(\mathrm{po})$,
\[
\sum_{p\in\mathrm{Cu}^2(\mathrm{po})}[\mathrm{qb}:p]\cdot[p:b]=\sum_{p}[\mathrm{qb}:p]\cdot[p:b]=0
\]
by (I1), the middle sum extending over all plaquettes.
The adjoint maps with respect to the counting inner products are denoted by $d_{\mathrm{po}}^{*}$;
one has $d_{\mathrm{po}}^{*}d_{\mathrm{po}}^{*}=0$.
\item[(Ap5)] Let $m$ be a real one-cochain that vanishes on the links of $\mathrm{Cu}^1(\mathrm{po})$ and on the links contained in $\{y_\mu\le\mathrm{po}\}$.
Then, for every $\mathrm{qb}\in\mathrm{Cu}^3(\mathrm{po})$,
\[
a_{\mathrm{qb}}(dm)+\sum_{p\in\mathrm{Cu}^2(\mathrm{po})}[\mathrm{qb}:p]\cdot\alpha_p(m)=0,
\]
where $a_{\mathrm{qb}}(dm):=\mathrm{sn}(\mathrm{qb})\cdot(dm)(\mathrm{qb}^{+})$ is given by the formula of (Ap3).
To see this, apply (Ap3) to $dm$:
\[
0=(ddm)(\mathrm{qb})=\mathrm{sn}(\mathrm{qb})\cdot\left((dm)(\mathrm{qb}^{+})-(dm)(\mathrm{qb}^{-})\right)+\sum_{p\in\mathrm{Cu}^2(\mathrm{po})}[\mathrm{qb}:p]\cdot(dm)(p).
\]
The links of $\mathrm{qb}^{-}$ lie in the plane $\{y_\mu=\mathrm{po}^{-}\}$, so that $(dm)(\mathrm{qb}^{-})=0$.
For a cut plaquette $p$, (Ap3) gives
\[
(dm)(p)=\mathrm{sn}(p)\cdot\left(m(p^{+})-m(p^{-})\right)+(d_{\mathrm{po}}m)(p)=\mathrm{sn}(p)\cdot m(p^{+})=\alpha_p(m).
\]
\item[(Ap6)] (Partition.)
Let $(\mathrm{po}_k)_{k\in\mathbb{Z}}$ be a strictly increasing sequence in $\tfrac12\mathbb{Z}$, unbounded above and below.
Every cell of $\mathrm{Cy}$ belongs to exactly one of the sets $\mathrm{Mi}(\mathrm{po}_k)$, $\mathrm{Cu}(\mathrm{po}_k)$, $\mathrm{Sb}(\mathrm{po}_{k-1},\mathrm{po}_k)$, $k\in\mathbb{Z}$.
For $\mathrm{po}<\mathrm{po}'<\mathrm{po}''$ and every dimension $\mathrm{dg}$ the union
\[
\mathrm{Sb}^{\mathrm{dg}}(\mathrm{po},\mathrm{po}'')=\mathrm{Sb}^{\mathrm{dg}}(\mathrm{po},\mathrm{po}')\cup\mathrm{Mi}^{\mathrm{dg}}(\mathrm{po}')\cup\mathrm{Cu}^{\mathrm{dg}}(\mathrm{po}')\cup\mathrm{Sb}^{\mathrm{dg}}(\mathrm{po}',\mathrm{po}'')
\]
is disjoint.
For the proof, call $\{\mathrm{po}_{k-1}\le y_\mu\le\mathrm{po}_k\}$ the closed slab $k$ of the sequence, and let $\boldsymbol{\sigma}=(c;\mathrm{Md})$, $t:=c_\mu$.
Suppose first that $\mu\notin\mathrm{Md}$.
If $t=\mathrm{po}_k$, then $\boldsymbol{\sigma}\in\mathrm{Mi}(\mathrm{po}_k)$;
it lies in no set $\mathrm{Cu}(\mathrm{po}_{k'})$, and in no set $\mathrm{Sb}(\mathrm{po}_{k'-1},\mathrm{po}_{k'})$, because the closed slabs $k$ and $k+1$ are the only closed slabs of the sequence that contain it, and it lies in one of their planes.
Otherwise there is exactly one $k$ with $\mathrm{po}_{k-1}<t<\mathrm{po}_k$, and $\boldsymbol{\sigma}$ belongs to $\mathrm{Sb}(\mathrm{po}_{k-1},\mathrm{po}_k)$ and to none of the other sets.
Suppose now that $\mu\in\mathrm{Md}$.
If $t+\tfrac12=\mathrm{po}_k$, then $\boldsymbol{\sigma}\in\mathrm{Cu}(\mathrm{po}_k)$;
it lies in no set $\mathrm{Mi}(\mathrm{po}_{k'})$, and no closed slab of the sequence contains it.
Otherwise no $\mathrm{po}_k$ lies in $(t,t+1)$, and there is exactly one $k$ with $\mathrm{po}_{k-1}\le t<t+1\le\mathrm{po}_k$;
the cell $\boldsymbol{\sigma}$ lies in no plane, and it belongs to $\mathrm{Sb}(\mathrm{po}_{k-1},\mathrm{po}_k)$ only.
The second assertion follows by applying the first one to a sequence in which $\mathrm{po},\mathrm{po}',\mathrm{po}''$ are consecutive terms and to a sequence in which $\mathrm{po},\mathrm{po}''$ are consecutive terms.
\item[(Ap7)] Let $\boldsymbol{\sigma}\in\mathrm{Cu}(\mathrm{po})$.
Then $\boldsymbol{\sigma}^{+}\in\mathrm{Sb}(\mathrm{po},\mathrm{po}')\cup\mathrm{Mi}(\mathrm{po}')$ for every $\mathrm{po}'>\mathrm{po}$, and $\boldsymbol{\sigma}^{-}\in\mathrm{Mi}(\mathrm{po}')\cup\mathrm{Sb}(\mathrm{po}',\mathrm{po})$ for every $\mathrm{po}'<\mathrm{po}$.
Indeed, $\boldsymbol{\sigma}^{+}$ lies in the plane $\{y_\mu=\mathrm{po}^{+}\}$, and $\mathrm{po}^{+}\le\mathrm{po}'$;
if $\mathrm{po}'=\mathrm{po}^{+}$, then $\boldsymbol{\sigma}^{+}\in\mathrm{Mi}(\mathrm{po}')$, and otherwise $\boldsymbol{\sigma}^{+}$ is contained in the closed slab $\{\mathrm{po}\le y_\mu\le\mathrm{po}'\}$ and in neither of its two planes.
Likewise $\boldsymbol{\sigma}^{-}$ lies in the plane $\{y_\mu=\mathrm{po}^{-}\}$, and $\mathrm{po}'\le\mathrm{po}^{-}$;
if $\mathrm{po}'=\mathrm{po}^{-}$, then $\boldsymbol{\sigma}^{-}\in\mathrm{Mi}(\mathrm{po}')$, and otherwise $\boldsymbol{\sigma}^{-}$ is contained in the closed slab $\{\mathrm{po}'\le y_\mu\le\mathrm{po}\}$ and in neither of its two planes.
\end{itemize}

\subsubsection*{A.2. Configurations of a slab, weights, interior integrals}

Let $\mathrm{po}<\mathrm{po}'$.
A configuration of the closed slab is a pair
\[
\omega=(A,\mathfrak{z})\in\mathbb{R}^{\mathrm{Ln}[\mathrm{po},\mathrm{po}']}\times\mathbb{Z}^{\mathrm{Pq}[\mathrm{po},\mathrm{po}']}.
\]
A face of a cell contained in the closed slab is contained in the closed slab.
Hence the curvature $F(\omega):=dA+r_q\cdot\mathfrak{z}$ is defined on $\mathrm{Pq}[\mathrm{po},\mathrm{po}']$, and $d\mathfrak{z}$ is defined on $\mathrm{Sb}^3(\mathrm{po},\mathrm{po}')$, since no cube lies in a plane;
$F_p(\omega)$ denotes the value of $F(\omega)$ at $p$.
Put
\[
\mathrm{Cfg}[\mathrm{po},\mathrm{po}']:=\left\{\omega:\ (d\mathfrak{z})(\mathrm{qb})=0\ \text{for all}\ \mathrm{qb}\in\mathrm{Sb}^3(\mathrm{po},\mathrm{po}')\right\},
\]
\[
\omega+[m]:=\left(A+r_q\cdot m,\ \mathfrak{z}-dm\right)\quad\text{for}\ m\in\mathbb{Z}^{\mathrm{Ln}[\mathrm{po},\mathrm{po}']},\qquad
g_{\mathrm{w}}^{*}\omega:=\left(g_{\mathrm{w}}^{*}A,\ g_{\mathrm{w}}^{*}\mathfrak{z}\right).
\]
The curvature $F$, the cochain $d\mathfrak{z}$ and the phases $e^{iqA(b)}$ are invariant under $\omega\mapsto\omega+[m]$.
The configuration that consists of the part $\mathrm{mi}$ on $\mathrm{Mi}(\mathrm{po})$, the part $\mathrm{mi}'$ on $\mathrm{Mi}(\mathrm{po}')$ and the part $\mathrm{omi}$ on $\mathrm{Sb}(\mathrm{po},\mathrm{po}')$ is written $\omega=\mathrm{mi}\vee\mathrm{omi}\vee\mathrm{mi}'$;
at a link plane the part $\mathrm{mi}$ is empty.
The weight of the slab is
\[
\mathrm{wb}_{\mathrm{po},\mathrm{po}'}(\omega):=\prod_{p\in\mathrm{Sb}^2(\mathrm{po},\mathrm{po}')}\exp\left(-\frac{F_p(\omega)^2}{2}\right)\cdot\prod_{\mathrm{qb}\in\mathrm{Sb}^3(\mathrm{po},\mathrm{po}')}\mathbf{1}\left[(d\mathfrak{z})(\mathrm{qb})=0\right],
\]
and the interior integral of a function $\Phi$ of configurations is
\[
\mathrm{Ii}_{\mathrm{po},\mathrm{po}'}\Phi(\mathrm{mi},\mathrm{mi}'):=\sum_{\mathfrak{z}\in\mathbb{Z}^{\mathrm{Sb}^2(\mathrm{po},\mathrm{po}')}}\ \int_{[0,r_q)^{\mathrm{Sb}^1(\mathrm{po},\mathrm{po}')}}\mathrm{d}A\,\Phi\left(\mathrm{mi}\vee\mathrm{omi}\vee\mathrm{mi}'\right),\qquad\mathrm{omi}=(A,\mathfrak{z}).
\]
If $\mathrm{Sb}^1(\mathrm{po},\mathrm{po}')$ and $\mathrm{Sb}^2(\mathrm{po},\mathrm{po}')$ are empty, which happens for $\mathrm{po}'-\mathrm{po}=\tfrac12$, then $\mathrm{Ii}_{\mathrm{po},\mathrm{po}'}$ is the evaluation.

\begin{named}{Lemma A.1}[{The box is immaterial; cf.\ \cite[Lemma~5.3]{QEDII}}]
Let $\Phi$ be a Borel function on $\mathbb{R}^{\mathrm{Ln}[\mathrm{po},\mathrm{po}']}\times\mathbb{Z}^{\mathrm{Pq}[\mathrm{po},\mathrm{po}']}$ such that $\mathrm{Ii}_{\mathrm{po},\mathrm{po}'}|\Phi|<\infty$ and $\Phi(\omega+[m])=\Phi(\omega)$ for all $m\in\mathbb{Z}^{\mathrm{Sb}^1(\mathrm{po},\mathrm{po}')}$, extended by $0$ to $\mathrm{Ln}[\mathrm{po},\mathrm{po}']$.
Then $\mathrm{Ii}_{\mathrm{po},\mathrm{po}'}\Phi$ does not change if, for each $b\in\mathrm{Sb}^1(\mathrm{po},\mathrm{po}')$, the interval $[0,r_q)$ is replaced by an arbitrary interval $I_b$ of length $r_q$, open, closed or half open.
\end{named}

\begin{proof}
End points form a set of measure zero.
Since the intervals can be replaced one at a time, it suffices to replace the interval $[0,r_q)$ of one link $b$ by $[t,t+r_q)$.

Let $t=r_q$, and substitute $A_b=A'_b+r_q$.
By the invariance, $\Phi(A'+r_q\cdot\mathbf{1}_b,\mathfrak{z})=\Phi(A',\mathfrak{z}+d\mathbf{1}_b)$, and $\mathfrak{z}\mapsto\mathfrak{z}+(d\mathbf{1}_b)|_{\mathrm{Sb}^2(\mathrm{po},\mathrm{po}')}$ is a bijection of $\mathbb{Z}^{\mathrm{Sb}^2(\mathrm{po},\mathrm{po}')}$;
here $d\mathbf{1}_b$ vanishes on the plaquettes in the two planes, because their links lie in the planes.
For $t=-r_q$ substitute $A_b=A'_b-r_q$ and use the bijection $\mathfrak{z}\mapsto\mathfrak{z}-d\mathbf{1}_b$.
By iteration the assertion holds for every $t\in r_q\mathbb{Z}$.

Let $0<t<r_q$.
Then $\int_{[t,t+r_q)}=\int_{[t,r_q)}+\int_{[r_q,r_q+t)}$, and after summation over $\mathfrak{z}$ the second integral equals $\int_{[0,t)}$, by the substitution of the case $t=r_q$.
A general $t$ is of the form $j\cdot r_q+t_0$ with $j\in\mathbb{Z}$ and $0\le t_0<r_q$, and the two cases combine.
\end{proof}

\subsubsection*{A.3. Fermion sites and bonds, relative to the planes of the family (group coordinates)}

For $g\in O_{\mathrm{cub}}$ let $\mathrm{col}(g)\in\{1,2,3\}$ be the column and $\mathrm{sr}(g)\in\{1,-1\}$ the value of the non-zero entry in the row $\mu$ of $g$;
thus $g\,e_{\mathrm{col}(g)}=\mathrm{sr}(g)\,e_\mu$.
With $\mathrm{geo}(c,g)=c+g\,v_0$, where $v_0=(v_1,v_2,v_3)$ and $0<v_3<v_2<v_1<\tfrac12$ (Section~I.1.3 and Proposition~46\,(a)), one has
\[
\mathrm{geo}(c,g)_\mu=c_\mu+\mathrm{sr}(g)\cdot v_{\mathrm{col}(g)}.
\]
Since $0<v_j<\tfrac12$ and $c_\mu-\mathrm{po}\in\tfrac12\mathbb{Z}$, the following holds.
\begin{itemize}
\item[(Ap8)] Let $x=(c,g)$ and $\mathrm{po}\in\tfrac12\mathbb{Z}$.
Then $\mathrm{geo}(x)_\mu>\mathrm{po}$ if and only if either $c_\mu>\mathrm{po}$, or $c_\mu=\mathrm{po}$ and $\mathrm{sr}(g)=+1$.
Moreover $\mathrm{geo}(x)_\mu\ne\mathrm{po}$ for all $x$ and all $\mathrm{po}\in\tfrac12\mathbb{Z}$:
no site lies on a plane (Section~I.4.8), and the point $v_0$ does not enter.
\end{itemize}
For an open interval $\mathrm{Iv}$ with end points in $\tfrac12\mathbb{Z}\cup\{\pm\infty\}$ put
\[
X(\mathrm{Iv}):=\left\{x\in\mathbb{L}^{\mathrm{f}}_0:\ \mathrm{geo}(x)_\mu\in\mathrm{Iv}\right\}.
\]
Over a photon site $c$ of a site plane $\mathrm{po}=c_\mu$ the $48$ sites of the cluster are separated by the plane:
the $24$ sites with $\mathrm{sr}(g)=+1$ lie above it, and the $24$ sites with $\mathrm{sr}(g)=-1$ lie below it.
A link plane separates no cluster.
Furthermore
\[
\theta_{\mathrm{po}}(c,g)=\left(c-2\left(c_\mu-\mathrm{po}\right)e_\mu,\ r_\mu g\right),\qquad
\mathrm{col}(r_\mu g)=\mathrm{col}(g),\qquad
\mathrm{sr}(r_\mu g)=-\mathrm{sr}(g),
\]
so that $\mathrm{geo}(\theta_{\mathrm{po}}x)_\mu=2\,\mathrm{po}-\mathrm{geo}(x)_\mu$ and $\theta_{\mathrm{po}}X(\mathrm{Iv})=X(2\,\mathrm{po}-\mathrm{Iv})$.
\begin{itemize}
\item[(Bd0)] (Neighbours.)
Let $x=(c,g)$, $j:=\mathrm{col}(g)$ and $\mathrm{es}:=\mathrm{sr}(g)$.
The row $\mu$ of $g\,s_i$ is the row $\mu$ of $g$ multiplied by $s_i$ (Section~I.1.2 and (I.1.2)).
Therefore the neighbour $x^{(0)}=(c,g\,s_0)$ has $(\mathrm{col},\mathrm{sr})=(3,-\mathrm{es})$ if $j=3$ and $(\mathrm{col},\mathrm{sr})=(j,\mathrm{es})$ if $j\ne3$.
The neighbour $x^{(1)}=(c,g\,s_1)$ has $\mathrm{sr}=\mathrm{es}$, and its value of $\mathrm{col}$ is obtained from $j$ by exchanging $2$ and $3$.
The neighbour $x^{(2)}=(c,g\,s_2)$ has $\mathrm{sr}=\mathrm{es}$, and its value of $\mathrm{col}$ is obtained from $j$ by exchanging $1$ and $2$.
For the neighbour $x^{(3)}=(c+g\,e_1,g\,s_3)$ there are two cases:
if $j=1$, its photon site is $c+\mathrm{es}\,e_\mu$ and $(\mathrm{col},\mathrm{sr})=(1,-\mathrm{es})$;
if $j\ne1$, the coordinate $\mu$ of its photon site and the values of $\mathrm{col}$ and $\mathrm{sr}$ are those of $x$.
Accordingly, the heights $\mathrm{geo}(\cdot)_\mu$ of the two ends of a bond are as follows.
For the types $1$ and $2$ they are $c_\mu+\mathrm{es}\cdot v_j$ and $c_\mu+\mathrm{es}\cdot v_{j'}$, where $j'$ is the value of $\mathrm{col}$ at the other end, and both lie strictly between $c_\mu$ and $c_\mu+\tfrac{\mathrm{es}}{2}$.
For type $0$ with $j\ne3$, and for type $3$ with $j\ne1$, the two heights are equal.
For type $0$ with $j=3$ they are $c_\mu+\mathrm{es}\cdot v_3$ and $c_\mu-\mathrm{es}\cdot v_3$;
for type $3$ with $j=1$ they are $c_\mu+\mathrm{es}\cdot v_1$ and $c_\mu+\mathrm{es}-\mathrm{es}\cdot v_1$.
\item[(Bd1)] A bond whose end points lie on different sides of the plane $(y_\mu,\mathrm{po})$ is of the form $\{x,\theta_{\mathrm{po}}x\}$.
If $\mathrm{po}$ is a site plane, the bond is of type $0$, with $x=(c,g)$, $c_\mu=\mathrm{po}$ and $\mathrm{col}(g)=3$;
over each photon site of the plane there are eight such bonds, their lower ends being given by the eight $g$ with $g\,e_3=-e_\mu$, and these bonds carry no phase.
If $\mathrm{po}$ is a link plane, the bond is of type $3$;
its lower end is $x=(c,g)$ with $c_\mu=\mathrm{po}^{-}$ and $g\,e_1=e_\mu$, and its upper end lies over $c+e_\mu$.
Over each link $b=(c;\{\mu\})$ of $\mathrm{Cu}^1(\mathrm{po})$ there are eight such bonds, and
\[
U_A(x,\theta_{\mathrm{po}}x)=e^{iqA_b},\qquad A_b:=A(c,c+e_\mu).
\]
In both cases the two end points lie in $X((\mathrm{po}^{-},\mathrm{po}^{+}))$, the bond crosses no other plane of the family, and $\hat{d}_i(x)=e_\mu$ at the lower end $x$, where $i$ is the type of the bond.

To prove this, observe first that by (Bd0) the only bonds whose two ends are separated by an element of $\tfrac12\mathbb{Z}$ are the bonds of type $0$ with $j=3$, whose ends are separated by $c_\mu$ alone (since $v_3<\tfrac12$), and the bonds of type $3$ with $j=1$, whose ends are separated by $c_\mu+\tfrac{\mathrm{es}}{2}$ alone (since $0<v_1<\tfrac12$).
If $g\,e_3=\pm e_\mu$, then $g\,s_0\,g^{-1}=1-2\,(g\,e_3)(g\,e_3)^{T}=r_\mu$, so that $x^{(0)}=(c,r_\mu g)=\theta_{c_\mu}x$.
If $g\,e_1=\mathrm{es}\,e_\mu$, then $g\,s_3\,g^{-1}=1-2\,(g\,e_1)(g\,e_1)^{T}=r_\mu$, that is, $g\,s_3=r_\mu g$;
moreover $c+\mathrm{es}\,e_\mu=c-2\left(c_\mu-\mathrm{po}\right)e_\mu$ for $\mathrm{po}=c_\mu+\tfrac{\mathrm{es}}{2}$, and hence $x^{(3)}=\theta_{\mathrm{po}}x$.
For a given unit vector and a given prescribed image of it there are exactly $8$ elements $g$ of $O_{\mathrm{cub}}$.
Finally $\hat{d}_0(x)=-g\,\hat{n}_0=-g\,e_3$ and $\hat{d}_3(x)=-g\,\hat{n}_3=g\,e_1$ by (I.1.2) and Section~I.1.2, and the phases are those of (I.1.6) and (I.2.3).
\item[(Bd2)] For $\mathrm{po}<\mathrm{po}'$ in $\tfrac12\mathbb{Z}$ let $S_{(\mathrm{po},\mathrm{po}')}(A)$ be the sum of the terms
\[
M_{\mathrm{ch}}\cdot\bar{\psi}(x)\psi(x),\qquad x\in X((\mathrm{po},\mathrm{po}'))/\mathrm{Gmp},
\]
and of the terms
\[
-\bar{\psi}(x)M_i(x)\,U_A(x,x^{(i)})\,\psi(x^{(i)})\qquad\text{with}\ x\in X((\mathrm{po},\mathrm{po}'))\ \text{and}\ x^{(i)}\in X((\mathrm{po},\mathrm{po}')).
\]
It is an even element of $\mathrm{Gr}_{(\mathrm{po},\mathrm{po}')}$ (A.4).
Each of its terms is independent of the choice of the representative, because $\mathrm{chp}^2=1$, and its coefficients are polynomials in the phases $e^{\pm iqA(b)}$, $b\in\mathrm{Ln}[\mathrm{po},\mathrm{po}']$.

Indeed, only the bonds of type $3$ carry a phase, namely $U_A(x,x^{(3)})=\exp\left(iqA(c,c+g\,e_1)\right)$.
Let both ends of such a bond lie in the open slab.
If $\mathrm{col}(g)\ne1$, the link $(c,c+g\,e_1)$ lies in the plane $\{y_\mu=c_\mu\}$;
from $\mathrm{po}<\mathrm{geo}(x)_\mu<\mathrm{po}'$, $\left|\mathrm{geo}(x)_\mu-c_\mu\right|<\tfrac12$ and $c_\mu-\mathrm{po},\,c_\mu-\mathrm{po}'\in\tfrac12\mathbb{Z}$ it follows that $\mathrm{po}\le c_\mu\le\mathrm{po}'$.
If $\mathrm{col}(g)=1$, let $\mathrm{es}:=\mathrm{sr}(g)$ and $t:=\min(c_\mu,c_\mu+\mathrm{es})$;
the heights of the two ends are $t+v_1$ and $t+1-v_1$.
The inequality $\mathrm{po}<t+v_1$ gives $t>\mathrm{po}-\tfrac12$, hence $t\ge\mathrm{po}$;
the inequality $t+1-v_1<\mathrm{po}'$ gives $t+1<\mathrm{po}'+\tfrac12$, hence $t+1\le\mathrm{po}'$.
In both cases the link is contained in the closed slab.
\item[(Bd3)] Let $\mathrm{Xc}_{\mathrm{po}}(A)$ be the sum of the terms, for both orientations, of the bonds of $\mathrm{Cy}$ that cross $(y_\mu,\mathrm{po})$.
For $\mathrm{po}_0<\mathrm{po}_1<\dots<\mathrm{po}_s$,
\[
S_{(\mathrm{po}_0,\mathrm{po}_s)}=\sum_{k=1}^{s}S_{(\mathrm{po}_{k-1},\mathrm{po}_k)}+\sum_{k=1}^{s-1}\mathrm{Xc}_{\mathrm{po}_k}.
\]
The element $\mathrm{Xc}_{\mathrm{po}}$ depends only on the phases of the links of $\mathrm{Cu}^1(\mathrm{po})$, and at a site plane it contains no phase.
Indeed, by (Ap8) the set $X((\mathrm{po}_0,\mathrm{po}_s))$ is the disjoint union of the sets $X((\mathrm{po}_{k-1},\mathrm{po}_k))$, and by (Bd1) a bond whose ends lie in different slabs crosses some $\mathrm{po}_k$, and only one.
\end{itemize}

\subsubsection*{A.4. Grassmann algebras of slabs; automorphisms; Berezin integrals}

For an open interval $\mathrm{Iv}$ as in A.3 let $\mathrm{Gr}_{\mathrm{Iv}}$ be the Grassmann algebra generated by $\psi_\beta(x)$, $\bar{\psi}_\beta(x)$, $x\in X(\mathrm{Iv})$, $\beta=1,\dots,4$, subject to the relations
\[
\psi(x+v)=\mathrm{chp}(v)\cdot\psi(x),\qquad\bar{\psi}(x+v)=\mathrm{chp}(v)\cdot\bar{\psi}(x),\qquad v\in\mathrm{Gmp}.
\]
If $Y$ is a set of representatives of $X(\mathrm{Iv})/\mathrm{Gmp}$, then $\mathrm{Gr}_{\mathrm{Iv}}$ is the free Grassmann algebra $\mathrm{Gr}(Y)$.
The homomorphism $\mathrm{qu}\colon\mathrm{Gr}(X(\mathrm{Iv}))\to\mathrm{Gr}_{\mathrm{Iv}}$ sends every generator to the generator of the same name.
For bounded $\mathrm{Iv}$ the set $X(\mathrm{Iv})/\mathrm{Gmp}$ is finite;
put
\[
\mathrm{top}(x):=\prod_{\beta=1}^{4}\psi_\beta(x)\,\bar{\psi}_\beta(x),\qquad
\mathrm{top}(Y):=\prod_{x\in Y}\mathrm{top}(x),
\]
a product of even factors, and for $\mathrm{Ph}\in\mathrm{Gr}_{\mathrm{Iv}}$ let $\int_{\mathrm{Iv}}\mathrm{Ph}:=\int_Y\mathrm{Ph}$ be the coefficient of $\mathrm{top}(Y)$ in $\mathrm{Ph}$, in accordance with the normalisation (S1) of Section~1.3.
Since $\mathrm{top}(x+v)=\mathrm{chp}(v)^8\cdot\mathrm{top}(x)=\mathrm{top}(x)$, the integral $\int_{\mathrm{Iv}}$ does not depend on $Y$.

The automorphisms $\mathrm{Tg}_v$, $v\in\mathbb{Z}e_\mu$, and the maps $\Theta_{\mathrm{po}}$ of Section~1.3 are compatible with these relations, because $\theta_{\mathrm{po}}(x+v)=\theta_{\mathrm{po}}x+v$ for $v\in\mathrm{Gmp}$ and $\mathrm{chp}$ is real.
They induce
\[
\mathrm{Tg}_v\colon\mathrm{Gr}_{\mathrm{Iv}}\to\mathrm{Gr}_{\mathrm{Iv}+v_\mu},\qquad
\Theta_{\mathrm{po}}\colon\mathrm{Gr}_{\mathrm{Iv}}\to\mathrm{Gr}_{2\mathrm{po}-\mathrm{Iv}},
\]
and $\mathrm{qu}\,\mathrm{Tg}_v=\mathrm{Tg}_v\,\mathrm{qu}$, $\mathrm{qu}\,\Theta_{\mathrm{po}}=\Theta_{\mathrm{po}}\,\mathrm{qu}$.
\begin{itemize}
\item[(Tq1)] $\mathrm{Tg}_v\,\mathrm{Tg}_{v'}=\mathrm{Tg}_{v+v'}$.
\item[(Tq2)] $\mathrm{Tg}_v\,\Theta_{\mathrm{po}}\,\mathrm{Tg}_v^{-1}=\Theta_{\mathrm{po}+v_\mu}$.
Indeed,
\[
\mathrm{Tg}_v\Theta_{\mathrm{po}}\mathrm{Tg}_{-v}\,\psi(x)=\mathrm{Tg}_v\left(-\bar{\psi}\left(\theta_{\mathrm{po}}(x-v)\right)\gamma_\mu\right)=-\bar{\psi}\left(\theta_{\mathrm{po}}(x-v)+v\right)\gamma_\mu,
\]
where $\theta_{\mathrm{po}}(x-v)+v=\theta_{\mathrm{po}+v_\mu}x$ by (M4), and
\[
\mathrm{Tg}_v\Theta_{\mathrm{po}}\mathrm{Tg}_{-v}\,\bar{\psi}(x)=\mathrm{Tg}_v\left(-\gamma_\mu\psi\left(\theta_{\mathrm{po}}(x-v)\right)\right)=-\gamma_\mu\psi\left(\theta_{\mathrm{po}+v_\mu}x\right).
\]
Both sides of (Tq2) are antilinear anti-automorphisms, and they agree on the generators.
\item[(Tq3)] $\Theta_{\mathrm{po}}^2=1$, and $\Theta_{\mathrm{po}'}\Theta_{\mathrm{po}}=\mathrm{Tg}_{\tau}$ with $\tau:=2\left(\mathrm{po}'-\mathrm{po}\right)e_\mu$, by (S4) of Section~1.3.
Let $\mathrm{po}'':=2\mathrm{po}'-\mathrm{po}$;
it is $\mathrm{po}$ plus the coordinate $\mu$ of $\tau$.
By (Tq2),
\[
\Theta_{\mathrm{po}''}=\mathrm{Tg}_{\tau}\Theta_{\mathrm{po}}\mathrm{Tg}_{\tau}^{-1}
=\left(\Theta_{\mathrm{po}'}\Theta_{\mathrm{po}}\right)\Theta_{\mathrm{po}}\left(\Theta_{\mathrm{po}}\Theta_{\mathrm{po}'}\right)
=\Theta_{\mathrm{po}'}\Theta_{\mathrm{po}}\Theta_{\mathrm{po}'},
\]
and consequently
\[
\Theta_{\mathrm{po}''}\mathrm{Tg}_{\tau}=\left(\Theta_{\mathrm{po}'}\Theta_{\mathrm{po}}\Theta_{\mathrm{po}'}\right)\left(\Theta_{\mathrm{po}'}\Theta_{\mathrm{po}}\right)=\Theta_{\mathrm{po}'}.
\]
\item[(Be1)] (Cf.\ \cite[Proposition~5.17(a)]{QEDII}.)
Let $Y_1,\dots,Y_s$ be pairwise disjoint finite sets of classes of sites, and let $\mathrm{Ph}_k\in\mathrm{Gr}(Y_k)$, $k=1,\dots,s$, be arbitrary.
Then
\[
\int_{Y_1\cup\dots\cup Y_s}\mathrm{Ph}_1\mathrm{Ph}_2\cdots\mathrm{Ph}_s=\prod_{k=1}^{s}\int_{Y_k}\mathrm{Ph}_k ,
\]
without any sign.

For the proof put $\mathrm{Dp}(x,\beta):=\frac{\partial}{\partial\bar{\psi}_\beta(x)}\frac{\partial}{\partial\psi_\beta(x)}$, with left derivatives.
Then $\int_Y=\prod_{x\in Y,\,\beta}\mathrm{Dp}(x,\beta)$, and the operators $\mathrm{Dp}(x,\beta)$ commute with one another, each of them being a product of two anticommuting operators.
If $\mathrm{Ph}$ contains neither $\psi_\beta(x)$ nor $\bar{\psi}_\beta(x)$, then $\mathrm{Dp}(x,\beta)(\mathrm{Ph}\,\mathrm{Ph}')=\mathrm{Ph}\cdot\mathrm{Dp}(x,\beta)\mathrm{Ph}'$:
for homogeneous $\mathrm{Ph}$ the sign is $(-1)^{2\deg\mathrm{Ph}}=1$, and the general case follows by linearity.
With $\mathrm{Dp}_k:=\prod_{x\in Y_k,\,\beta}\mathrm{Dp}(x,\beta)$ one obtains
\[
\mathrm{Dp}_s\left(\mathrm{Ph}_1\cdots\mathrm{Ph}_s\right)=\mathrm{Ph}_1\cdots\mathrm{Ph}_{s-1}\cdot\mathrm{Dp}_s\mathrm{Ph}_s,
\]
where $\mathrm{Dp}_s\mathrm{Ph}_s$ is a number;
the assertion follows by induction on $s$.
\item[(Be2)] (Cf.\ \cite[Proposition~5.17(b)]{QEDII} and \cite[Proposition~3.16(e)]{QEDII}, here on the cylinder.)
For bounded $\mathrm{Iv}$ and $\mathrm{Ph}\in\mathrm{Gr}_{\mathrm{Iv}}$,
\[
\int_{2\mathrm{po}-\mathrm{Iv}}\Theta_{\mathrm{po}}\mathrm{Ph}=\overline{\int_{\mathrm{Iv}}\mathrm{Ph}}.
\]

For the proof let $Y$ be a set of representatives of $X(\mathrm{Iv})/\mathrm{Gmp}$;
then $\theta_{\mathrm{po}}Y$ is a set of representatives of $X(2\mathrm{po}-\mathrm{Iv})/\mathrm{Gmp}$.
Both sides are antilinear in $\mathrm{Ph}$, and $\Theta_{\mathrm{po}}$ preserves the degree;
so it suffices to show that $\Theta_{\mathrm{po}}(\mathrm{top}(Y))=\mathrm{top}(\theta_{\mathrm{po}}Y)$.
The map $\Theta_{\mathrm{po}}$ reverses products, and the factors $\mathrm{top}(x)$, as well as the pairs inside $\mathrm{top}(x)$, are even.
Hence, by (S2) of Section~1.3, whose two signs cancel,
\[
\Theta_{\mathrm{po}}(\mathrm{top}(Y))=\prod_{x\in Y}\prod_{\beta=1}^{4}\Theta_{\mathrm{po}}\left(\bar{\psi}_\beta(x)\right)\Theta_{\mathrm{po}}\left(\psi_\beta(x)\right)
=\prod_{x\in Y}\prod_{\beta=1}^{4}\left(\gamma_\mu\psi(\theta_{\mathrm{po}}x)\right)_\beta\left(\bar{\psi}(\theta_{\mathrm{po}}x)\gamma_\mu\right)_\beta .
\]
There is a sign $s_4$, the same for every family of eight anticommuting elements, such that $\prod_{\beta}a_\beta b_\beta=s_4\cdot(a_1a_2a_3a_4)(b_1b_2b_3b_4)$.
Moreover, with the products taken in the order $\beta=1,2,3,4$,
\[
\prod_{\beta=1}^{4}\left(\gamma_\mu\psi(\theta_{\mathrm{po}}x)\right)_\beta=\det(\gamma_\mu)\cdot\prod_{\beta=1}^{4}\psi_\beta(\theta_{\mathrm{po}}x),\qquad
\prod_{\beta=1}^{4}\left(\bar{\psi}(\theta_{\mathrm{po}}x)\gamma_\mu\right)_\beta=\det(\gamma_\mu)\cdot\prod_{\beta=1}^{4}\bar{\psi}_\beta(\theta_{\mathrm{po}}x),
\]
and $\det\gamma_\mu=1$:
the eigenvalues of $\gamma_\mu$ are $1,1,-1,-1$, because $\gamma_\mu^2=1$ and $\operatorname{tr}\gamma_\mu=0$, the latter since $\gamma_\mu=-\gamma_\nu\gamma_\mu\gamma_\nu$ for $\nu\ne\mu$ by (I.c1).
Therefore $\prod_{\beta}\left(\gamma_\mu\psi(\theta_{\mathrm{po}}x)\right)_\beta\left(\bar{\psi}(\theta_{\mathrm{po}}x)\gamma_\mu\right)_\beta=\mathrm{top}(\theta_{\mathrm{po}}x)$.
\item[(Tq4)] $\int_{\mathrm{Iv}+v_\mu}\mathrm{Tg}_v\mathrm{Ph}=\int_{\mathrm{Iv}}\mathrm{Ph}$, because $\mathrm{Tg}_v\mathrm{top}(x)=\mathrm{top}(x+v)$;
and $\int_{2\mathrm{po}-\mathrm{Iv}}\Theta_{\mathrm{po}}\mathrm{Ph}=\overline{\int_{\mathrm{Iv}}\mathrm{Ph}}$ by (Be2).
\end{itemize}

\subsubsection*{A.5. The classes $\mathrm{Ig}$; elements of $\mathcal{A}$ as elements of $\mathrm{Ig}$}

\begin{named}{Definition}
Let $\mathrm{po}<\mathrm{po}'$ be elements of $\tfrac12\mathbb{Z}$.
The class $\mathrm{Ig}[\mathrm{po},\mathrm{po}']$ consists of the Borel maps
\[
G\colon\mathrm{Cfg}[\mathrm{po},\mathrm{po}']\to\mathrm{Gr}^{\mathrm{ev}}_{(\mathrm{po},\mathrm{po}')}
\]
with the following two properties.
\begin{itemize}
\item[(Ig1)] $G(\omega+[m])=G(\omega)$ for every $m\in\mathbb{Z}^{\mathrm{Ln}[\mathrm{po},\mathrm{po}']}$.
\item[(Ig2)] There are $\mathrm{Cq}<\infty$ and $\varepsilon_G\in(0,\tfrac12]$ such that, for every $\omega\in\mathrm{Cfg}[\mathrm{po},\mathrm{po}']$, each coefficient of $G(\omega)$ in a monomial basis has modulus at most
\[
\mathrm{Cq}\cdot\exp\left(\left(\tfrac12-\varepsilon_G\right)\sum_{p\in\mathrm{Sb}^2(\mathrm{po},\mathrm{po}')}F_p(\omega)^2\right).
\]
\end{itemize}
\end{named}

Since $\sum_{p\in\mathrm{Sb}^2(\mathrm{po},\mathrm{po}')}F_p(\omega)^2\ge0$, the bound in (Ig2) is non-increasing in $\varepsilon_G$.
Consequently $\mathrm{Ig}[\mathrm{po},\mathrm{po}']$ is a linear space:
for the sum of two of its elements, (Ig2) holds with the sum of the two constants $\mathrm{Cq}$ and the smaller of the two exponents.
As $\tfrac12-\varepsilon_G\ge0$, the bound in (Ig2) is moreover monotone in the set of the plaquettes over which the sum extends.

Three operations act on these classes.

\emph{Product.}
Let $\mathrm{po}<\mathrm{po}'<\mathrm{po}''$, $G\in\mathrm{Ig}[\mathrm{po},\mathrm{po}']$ and $G'\in\mathrm{Ig}[\mathrm{po}',\mathrm{po}'']$, and put
\[
(GG')(\omega):=G\left(\omega|_{[\mathrm{po},\mathrm{po}']}\right)\,G'\left(\omega|_{[\mathrm{po}',\mathrm{po}'']}\right),\qquad\omega\in\mathrm{Cfg}[\mathrm{po},\mathrm{po}''].
\]
Then $GG'\in\mathrm{Ig}[\mathrm{po},\mathrm{po}'']$.
Indeed, by (Ap6) the sets $\mathrm{Sb}^2(\mathrm{po},\mathrm{po}')$ and $\mathrm{Sb}^2(\mathrm{po}',\mathrm{po}'')$ are disjoint subsets of $\mathrm{Sb}^2(\mathrm{po},\mathrm{po}'')$,
and $\mathrm{Sb}^3(\mathrm{po},\mathrm{po}')$ and $\mathrm{Sb}^3(\mathrm{po}',\mathrm{po}'')$ are subsets of $\mathrm{Sb}^3(\mathrm{po},\mathrm{po}'')$;
by the monotonicity just noted, (Ig2) holds for $GG'$ with the exponent $\min\{\varepsilon_G,\varepsilon_{G'}\}$.

\emph{Translation.}
For $v\in\mathbb{Z}e_\mu$ and $G\in\mathrm{Ig}[\mathrm{po},\mathrm{po}']$ put
\[
(\mathrm{Tg}_vG)(\omega'):=\mathrm{Tg}_v\left(G(\tau_v^*\omega')\right);
\]
then $\mathrm{Tg}_vG\in\mathrm{Ig}[\mathrm{po}+v_\mu,\mathrm{po}'+v_\mu]$.

\emph{Reflection.}
For $G\in\mathrm{Ig}[\mathrm{po},\mathrm{po}']$ put
\[
(\Theta_{\mathrm{po}'}G)(\omega'):=\Theta_{\mathrm{po}'}\left(G(\theta_{\mathrm{po}'}^*\omega')\right);
\]
then $\Theta_{\mathrm{po}'}G\in\mathrm{Ig}[\mathrm{po}',2\mathrm{po}'-\mathrm{po}]$.
Indeed, $\theta_{\mathrm{po}'}$ maps the closed slab $\{\mathrm{po}\le y_\mu\le\mathrm{po}'\}$ onto the closed slab $\{\mathrm{po}'\le y_\mu\le2\mathrm{po}'-\mathrm{po}\}$,
cells onto cells, and $\mathrm{Sb}(\mathrm{po},\mathrm{po}')$ onto $\mathrm{Sb}(\mathrm{po}',2\mathrm{po}'-\mathrm{po})$;
the pull-back $\theta_{\mathrm{po}'}^*$ commutes with $d$;
$\theta_{\mathrm{po}'}^*(\omega'+[m])=\theta_{\mathrm{po}'}^*\omega'+[\theta_{\mathrm{po}'}^*m]$;
and $(\theta_{\mathrm{po}'}^*u)_p^2=u_{\theta_{\mathrm{po}'}p}^2$.
Finally, $\gamma_\mu$ is monomial (I.c6),
so that $\Theta_{\mathrm{po}'}$ maps every monomial to a monomial times a number of modulus one,
and (Ig2) holds for $\Theta_{\mathrm{po}'}G$ with the same $\mathrm{Cq}$ and $\varepsilon_G$.

\begin{named}{Lemma A.2}[{cf.\ \cite[Lemma~5.14(b), (d), Lemma~5.13(b)]{QEDII}}]
Let $\mathrm{Ksl}[\mathrm{po},\mathrm{po}']$ be the complex of the cells of $\mathbb{Z}^3$ that are contained in $\{\mathrm{po}\le y_\mu\le\mathrm{po}'\}$;
it is the product of the intervals $I_\mu=\mathbb{Z}\cap[\mathrm{po},\mathrm{po}']$, which is not empty because $\mathrm{po}'-\mathrm{po}\ge\tfrac12$, and $I_\nu=I_{\mathfrak{l}}=\mathbb{Z}$.
For $\omega=(A,\mathfrak{z})\in\mathrm{Cfg}[\mathrm{po},\mathrm{po}']$ let $\tilde A$, $\tilde{\mathfrak{z}}$ and $\widetilde{F(\omega)}$ be the $\mathrm{Gmp}$-periodic extensions of $A$, $\mathfrak{z}$ and $F(\omega)$ to $\mathrm{Ksl}[\mathrm{po},\mathrm{po}']$.
A \emph{lift} of $\omega$ is a real cochain $\hat A$ on the links of $\mathrm{Ksl}[\mathrm{po},\mathrm{po}']$ such that $d\hat A=\widetilde{F(\omega)}$ on its plaquettes
and $\exp\left(iq\hat A(b)\right)=\exp\left(iq\tilde A(b)\right)$ for its links $b$.
For $P\in\mathcal{A}_{\mu,n}[\mathrm{po},\mathrm{po}']$ put
\[
P(\omega):=\mathrm{qu}\left(P(\hat A)\right),
\]
where $\hat A$ is a lift of $\omega$, extended by $0$ to the other links of $\mathbb{Z}^3$.
Then the following hold.
\begin{itemize}
\item[(i)] Lifts exist; $P(\omega)$ does not depend on the choice of the lift; the map $\omega\mapsto P(\omega)$ belongs to $\mathrm{Ig}[\mathrm{po},\mathrm{po}']$; and $P\mapsto P(\,\cdot\,)$ is linear.
\item[(ii)] Let $\mathrm{po}''>\mathrm{po}'$ be in $\tfrac12\mathbb{Z}$, $G\in\mathcal{A}_{\mu,n}[\mathrm{po},\mathrm{po}']$ and $G'\in\mathcal{A}_{\mu,n}[\mathrm{po}',\mathrm{po}'']$.
Then $GG'\in\mathcal{A}_{\mu,n}[\mathrm{po},\mathrm{po}'']$, and
\[
(GG')(\omega)=G\left(\omega|_{[\mathrm{po},\mathrm{po}']}\right)\,G'\left(\omega|_{[\mathrm{po}',\mathrm{po}'']}\right),\qquad\omega\in\mathrm{Cfg}[\mathrm{po},\mathrm{po}''].
\]
\item[(iii)] Let $G\in\mathcal{A}_{\mu,n}[\mathrm{po},\mathrm{po}']$ and $v\in\mathbb{Z}e_\mu$.
Then $G^{\tau_v}\in\mathcal{A}_{\mu,n}[\mathrm{po}+v_\mu,\mathrm{po}'+v_\mu]$, and the image of $G^{\tau_v}$ in $\mathrm{Ig}[\mathrm{po}+v_\mu,\mathrm{po}'+v_\mu]$ is $\mathrm{Tg}_v$ applied to the image of $G$.
\item[(iv)] Let $G\in\mathcal{A}_{\mu,n}[\mathrm{po},\mathrm{po}']$.
Then $\Theta_{\mu,\mathrm{po}'}G\in\mathcal{A}_{\mu,n}[\mathrm{po}',2\mathrm{po}'-\mathrm{po}]$, and the image of $\Theta_{\mu,\mathrm{po}'}G$ in $\mathrm{Ig}[\mathrm{po}',2\mathrm{po}'-\mathrm{po}]$ is $\Theta_{\mathrm{po}'}$ applied to the image of $G$.
\end{itemize}
\end{named}

\begin{proof}
(i) The cubes of $\mathrm{Ksl}[\mathrm{po},\mathrm{po}']$ are the cubes over $\mathrm{Sb}^3(\mathrm{po},\mathrm{po}')$.
So $\tilde{\mathfrak{z}}$ is closed on them, and Lemma~2.4\,(a) gives $\tilde{\mathfrak{z}}=dm$ with an integer one-cochain $m$ on $\mathrm{Ksl}[\mathrm{po},\mathrm{po}']$.
Then $\hat A:=\tilde A+r_q\cdot m$ is a lift of $\omega$:
$d\hat A=d\tilde A+r_q\cdot\tilde{\mathfrak{z}}=\widetilde{F(\omega)}$, and $\exp\left(iq\,r_q\cdot m(b)\right)=1$ because $q\,r_q=2\pi$.

A product of $\mathcal{A}(\mathrm{Sl}_\mu(\mathrm{po},\mathrm{po}'))$ reads only links and plaquettes that are contained in the open slab;
hence $P(\hat A)$ depends only on the restriction of $\hat A$ to $\mathrm{Ksl}[\mathrm{po},\mathrm{po}']$.
Its values lie in $\mathrm{Gr}(X((\mathrm{po},\mathrm{po}')))$:
if $c+[-\tfrac12,\tfrac12]^3$ lies in a box inside the open slab, then $\mathrm{po}<c_\mu-\tfrac12$ and $c_\mu+\tfrac12<\mathrm{po}'$,
so that $\mathrm{geo}(c,g)_\mu\in(\mathrm{po},\mathrm{po}')$ for all $48$ labels $g$.

Let $\hat A$ and $\hat A'$ be lifts of $\omega$.
Their phases are equal, so $r_q^{-1}(\hat A-\hat A')$ is integer valued;
it is closed on the plaquettes of $\mathrm{Ksl}[\mathrm{po},\mathrm{po}']$, and by Lemma~2.4\,(a) it equals $d\boldsymbol{\lambda}$ for an integer valued function $\boldsymbol{\lambda}$ on $\mathrm{Ksl}[\mathrm{po},\mathrm{po}']$.
For a product $\Phi\cdot G_1\cdots G_s$ one has $\Phi(d\hat A')=\Phi(d\hat A)$,
and the gauge invariance of $G_j$, applied with the real function $r_q\boldsymbol{\lambda}$ on the sites of $\mathrm{Bx}_j$, gives $G_j(\hat A')=G_j(\hat A)$, because $\exp(2\pi i\boldsymbol{\lambda})=1$.
Thus $P(\hat A')=P(\hat A)$ for products, and by linearity for all $P\in\mathcal{A}_{\mu,n}[\mathrm{po},\mathrm{po}']$.

(Ig1): the configuration $\omega+[m]$ has the same curvature and the same phases as $\omega$, hence the same lifts.
Linearity: for a fixed cochain $\hat A$ the map $P\mapsto\mathrm{qu}\left(P(\hat A)\right)$ is linear.

(Ig2), for a product $\Phi\cdot G_1\cdots G_s$ that embeds in $\mathrm{Gmp}$, where $\varepsilon_j$ denotes the exponent of the factor $G_j$:
the plaquettes contained in a box lie in the open slab, so that they belong to $\mathrm{Sb}^2(\mathrm{po},\mathrm{po}')$;
and two different plaquettes of $\bigcup_j\mathrm{Pq}(\mathrm{Bx}_j)$ are inequivalent modulo $\mathrm{Gmp}$, because the sets $\mathrm{Bx}_j+v$, $v\in\mathrm{Gmp}$, are pairwise disjoint.
With $\varepsilon:=\min\{\tfrac12,\varepsilon_1,\dots,\varepsilon_s\}$ it follows that
\[
\sum_j\left(\tfrac12-\varepsilon_j\right)\sum_{p\in\mathrm{Pq}(\mathrm{Bx}_j)}(d\hat A)_p^2\le\left(\tfrac12-\varepsilon\right)\sum_{p\in\mathrm{Sb}^2(\mathrm{po},\mathrm{po}')}F_p(\omega)^2 .
\]
Since $\sup|\Phi|<\infty$ and a coefficient of a product is a finite sum of products of coefficients of the factors, this gives (Ig2).
Borel measurability: for fixed $\mathfrak{z}$ the integer cochain $m$ is fixed, $A\mapsto\hat A=\tilde A+r_q\cdot m$ is affine, and the $G_j$ and $\Phi$ are Borel.

(ii) The boxes of $G$ and those of $G'$ lie in two disjoint open slabs, which are invariant under $\mathrm{Gmp}$;
so the products of which $GG'$ is composed embed in $\mathrm{Gmp}$.
If $\hat A$ is a lift of $\omega\in\mathrm{Cfg}[\mathrm{po},\mathrm{po}'']$, its restrictions to $\mathrm{Ksl}[\mathrm{po},\mathrm{po}']$ and to $\mathrm{Ksl}[\mathrm{po}',\mathrm{po}'']$ are lifts of the restrictions of $\omega$;
and $\mathrm{qu}$ is multiplicative.

(iii), (iv) Membership follows from Lemma~2.2\,(B2),
since $\tau_v\mathrm{Sl}_\mu(\mathrm{po},\mathrm{po}')=\mathrm{Sl}_\mu(\mathrm{po}+v_\mu,\mathrm{po}'+v_\mu)$ and $\theta_{\mathrm{po}'}\mathrm{Sl}_\mu(\mathrm{po},\mathrm{po}')=\mathrm{Sl}_\mu(\mathrm{po}',2\mathrm{po}'-\mathrm{po})$,
and since the linear parts $1$ and $r_\mu$ fix $\mathrm{Gmp}$ pointwise.
By (1.2), $G^{\tau_v}(\hat A)=\mathrm{Tg}_v\left(G(\tau_v^*\hat A)\right)$;
if $\hat A$ is a lift of $\omega'\in\mathrm{Cfg}[\mathrm{po}+v_\mu,\mathrm{po}'+v_\mu]$, then $\tau_v^*\hat A$ is a lift of $\tau_v^*\omega'$;
and $\mathrm{qu}\,\mathrm{Tg}_v=\mathrm{Tg}_v\,\mathrm{qu}$.
By (S5), $(\Theta_{\mu,\mathrm{po}'}G)(\hat A)=\Theta_{\mathrm{po}'}\left(G(\theta_{\mathrm{po}'}^*\hat A)\right)$.
Let $\hat A$ be a lift of $\omega'=(A',\mathfrak{z}')\in\mathrm{Cfg}[\mathrm{po}',2\mathrm{po}'-\mathrm{po}]$, and let $\tilde A'$ be the periodic extension of $A'$.
Then $\theta_{\mathrm{po}'}^*\hat A$ is a lift of $\theta_{\mathrm{po}'}^*\omega'$:
$d$ commutes with $\theta_{\mathrm{po}'}^*$, and
\begin{align*}
\exp\left(iq(\theta_{\mathrm{po}'}^*\hat A)(b)\right)
&=\left(\exp\left(iq\hat A(\theta_{\mathrm{po}'}b)\right)\right)^{\epsilon(\theta_{\mathrm{po}'},b)}\\
&=\left(\exp\left(iq\tilde A'(\theta_{\mathrm{po}'}b)\right)\right)^{\epsilon(\theta_{\mathrm{po}'},b)}
=\exp\left(iq(\theta_{\mathrm{po}'}^*\tilde A')(b)\right).
\end{align*}
Moreover $\mathrm{qu}\,\Theta_{\mathrm{po}'}=\Theta_{\mathrm{po}'}\,\mathrm{qu}$.
\end{proof}

\subsubsection*{A.6. Modes at a plane; the crossing expansion}

The objects of this section and the crossing expansion are those of \cite[Definition~2.27]{QEDII} and \cite[Proposition~3.20]{QEDII}, for the family $\ell=y_\mu$.

The matrices $P^{\pm}:=\tfrac12(1\pm\gamma_\mu)$ are the orthogonal projections onto the eigenspaces of the Hermitian matrix $\gamma_\mu$ (I.c1).
Each of them has rank two: by (I.c1), $\gamma_\mu=-\gamma_\nu\gamma_\mu\gamma_\nu$ for $\nu\ne\mu$, so that $\operatorname{tr}\gamma_\mu=0$.
Orthonormal bases $(e^-_1,e^-_2)$ of $\operatorname{ran}P^-$ and $(e^+_1,e^+_2)$ of $\operatorname{ran}P^+$ are fixed, the same for all planes of the family;
then $P^{\pm}=\sum_{k'}e^{\pm}_{k'}\left(e^{\pm}_{k'}\right)^*$.

Let $\mathrm{po}\in\tfrac12\mathbb{Z}$, let $\{x,\theta_{\mathrm{po}}x\}$ be a bond that crosses $(y_\mu,\mathrm{po})$, with lower end $x$, and let $i_b\in\{0,3\}$ be its type (Bd1);
recall that $w_0=w_3=6$.
By (Bd1), $\hat{d}_{i_b}(x)=e_\mu$; therefore (I.2.1) and (I.m1) give
\begin{align*}
M_{i_b}(x)&=3(\rho-\gamma_\mu)=3\left[(1+\rho)P^--(1-\rho)P^+\right],\\
M_{i_b}(\theta_{\mathrm{po}}x)&=3(\rho+\gamma_\mu)=3\left[(1+\rho)P^+-(1-\rho)P^-\right].
\end{align*}

\emph{Modes.}
Put $U:=U_A(x,\theta_{\mathrm{po}}x)$; thus $U=1$ at a site plane, and $U=\exp(iqA_b)$ at a link plane.
For $k'\in\{1,2\}$ the bond carries one mode $C$ of each of four groups, with a coefficient $c$ and a phase $u$:
\begin{align*}
&\text{group 1:}&C&=(e^-_{k'})^*\psi(\theta_{\mathrm{po}}x),&c&=3(1+\rho),&u&=U;\\
&\text{group 2:}&C&=(e^+_{k'})^*\psi(\theta_{\mathrm{po}}x),&c&=3(1-\rho),&u&=U;\\
&\text{group 3:}&C&=\bar{\psi}(\theta_{\mathrm{po}}x)\,e^+_{k'},&c&=3(1+\rho),&u&=\overline{U};\\
&\text{group 4:}&C&=\bar{\psi}(\theta_{\mathrm{po}}x)\,e^-_{k'},&c&=3(1-\rho),&u&=\overline{U}.
\end{align*}
Let $\mathrm{idx}(\mathrm{po})$ be the set of the triples $i=(\text{bond},\text{group},k')$; there are eight of them for each crossing bond.
For $i\in\mathrm{idx}(\mathrm{po})$, the mode, its coefficient and its phase are denoted $C_i$, $c_i$ and $u_i$.
At both kinds of planes, $c_i$ equals $3(1+\rho)$ or $3(1-\rho)$, and $c_i>0$ because $0<\rho<1$.
The eight $C_i$ of one bond form a basis of the span of the eight generators at its upper end, and the eight $\Theta_{\mathrm{po}}(C_i)$ a basis of the span of the eight generators at its lower end.
By (Bd1), $C_i\in\mathrm{Gr}_{(\mathrm{po},\mathrm{po}^+)}$ and $\Theta_{\mathrm{po}}(C_i)\in\mathrm{Gr}_{(\mathrm{po}^-,\mathrm{po})}$.

For $v\in\mathbb{Z}^3$ put $i+v:=(\text{bond}+v,\text{group},k')\in\mathrm{idx}(\mathrm{po}+v_\mu)$;
then $C_{i+v}=\mathrm{Tg}_vC_i$ and $c_{i+v}=c_i$.
In the algebras $\mathrm{Gr}_{\mathrm{Iv}}$ one has $C_{i+v}=\mathrm{chp}(v)\cdot C_i$ for $v\in\mathrm{Gmp}$.

\emph{Representatives and orders} (cf.\ \cite[Definition~2.27]{QEDII}).
For every $\mathrm{po}\in\tfrac12\mathbb{Z}$, $\mathrm{Ri}(\mathrm{po})$ is a set of representatives of $\mathrm{idx}(\mathrm{po})/\mathrm{Gmp}$ with a total order,
such that $\mathrm{Ri}(\mathrm{po}+1)=\mathrm{Ri}(\mathrm{po})+e_\mu$ with the transported order.
Such sets exist: choose $\mathrm{Ri}(0)$ and $\mathrm{Ri}(\tfrac12)$ and transport them.
Thus every site plane carries a translate of $\mathrm{Ri}(0)$, and every link plane a translate of $\mathrm{Ri}(\tfrac12)$.
By (Bd1), $\#|\mathrm{Ri}(\mathrm{po})|=64\,v_s$ at both kinds of planes.
For $I=\{i_1<\dots<i_d\}\subset\mathrm{Ri}(\mathrm{po})$ put
\[
C_I:=C_{i_1}\cdots C_{i_d},\qquad c_I:=\prod_{i\in I}c_i,\qquad u_I:=\prod_{i\in I}u_i ;
\]
then $\Theta_{\mathrm{po}}(C_I)=\Theta_{\mathrm{po}}(C_{i_d})\cdots\Theta_{\mathrm{po}}(C_{i_1})$.
Where the plane has to be named, an upper index is added: $C^{(\mathrm{po})}_I$.

Let $\mathrm{Fk}(\mathrm{po})$ be the Hilbert space with the orthonormal basis $(e_I)$, $I$ ranging over the subsets of $\mathrm{Ri}(\mathrm{po})$, and put
\[
\mathrm{Pf}\,e_I:=(-1)^{\#|I|}e_I,\qquad\mathrm{cc}_{\mathrm{po}}:=\operatorname{diag}(c_I)_{I\subset\mathrm{Ri}(\mathrm{po})}.
\]
A matrix on $\mathrm{Fk}(\mathrm{po})$, or from $\mathrm{Fk}(\mathrm{po})$ to $\mathrm{Fk}(\mathrm{po}')$, is \emph{even} if its entry with the indices $I$, $I'$ vanishes unless $\#|I|+\#|I'|$ is even.
For $v\in\mathbb{Z}e_\mu$ let
\[
\mathrm{Ft}_v\colon\mathrm{Fk}(\mathrm{po})\to\mathrm{Fk}(\mathrm{po}+v_\mu),\qquad\mathrm{Ft}_v\,e_I:=e_{I+v} .
\]
It is unitary, with entries $0$ and $1$, and
\[
\mathrm{Ft}_v\,\mathrm{Ft}_{v'}=\mathrm{Ft}_{v+v'},\qquad\mathrm{Ft}_v\,\mathrm{cc}_{\mathrm{po}}=\mathrm{cc}_{\mathrm{po}+v_\mu}\,\mathrm{Ft}_v .
\]

For $v\in\mathbb{Z}e_\mu$ and $I\subset\mathrm{Ri}(\mathrm{po})$,
\begin{equation}
\mathrm{Tg}_v\,C^{(\mathrm{po})}_I=C^{(\mathrm{po}+v_\mu)}_{I+v};
\tag{Tq5}
\end{equation}
no sign occurs, because the order is transported.
Let $\mathrm{po}''=2\mathrm{po}'-\mathrm{po}$ and $\tau=2(\mathrm{po}'-\mathrm{po})e_\mu$; this is an element of $\mathbb{Z}e_\mu$, whatever the kinds of $\mathrm{po}$ and $\mathrm{po}'$ are.
Then
\begin{equation}
\Theta_{\mathrm{po}''}\left(C^{(\mathrm{po}'')}_{I+\tau}\right)=\Theta_{\mathrm{po}'}\left(C^{(\mathrm{po})}_I\right).
\tag{Tq6}
\end{equation}
Indeed, $C^{(\mathrm{po}'')}_{I+\tau}=\mathrm{Tg}_\tau C^{(\mathrm{po})}_I$ by (Tq5), and $\Theta_{\mathrm{po}''}\mathrm{Tg}_\tau=\Theta_{\mathrm{po}'}$ by (Tq3).

For a matrix $\mathrm{mt}=(\mathrm{mt}_{I,\mathrm{Ia}})$ on $\mathrm{Fk}(\mathrm{po})$ put
\[
\mathrm{Ins}_{\mathrm{po}}[\mathrm{mt}]:=\sum_{I,\mathrm{Ia}}\mathrm{mt}_{I,\mathrm{Ia}}\cdot\Theta_{\mathrm{po}}(C_{\mathrm{Ia}})\,C_I\ \in\ \mathrm{Gr}_{(\mathrm{po}^-,\mathrm{po}^+)} ;
\]
the row index $I$ belongs to the upper side and the column index $\mathrm{Ia}$ to the lower side, as in \cite[Definition~2.27]{QEDII}.
If $\mathrm{mt}$ is even, then $\mathrm{Ins}_{\mathrm{po}}[\mathrm{mt}]$ is even.

The crossing expansion reads
\begin{equation}
\exp\left(-\mathrm{Xc}_{\mathrm{po}}(A)\right)=\sum_{I\subset\mathrm{Ri}(\mathrm{po})}c_I\cdot u_I(A)\cdot\Theta_{\mathrm{po}}(C_I)\,C_I=\mathrm{Ins}_{\mathrm{po}}\left[\mathrm{cc}_{\mathrm{po}}\operatorname{diag}\left(u_I(A)\right)\right].
\tag{Ap9}
\end{equation}

\begin{proof}[Proof of (Ap9)]
\emph{One bond.}
For the bond $\{x,\theta_{\mathrm{po}}x\}$,
\[
-(\text{terms of the bond})=\bar{\psi}(x)M_{i_b}(x)\,U\,\psi(\theta_{\mathrm{po}}x)+\bar{\psi}(\theta_{\mathrm{po}}x)M_{i_b}(\theta_{\mathrm{po}}x)\,\overline{U}\,\psi(x).
\]
Let $\mathrm{ev}\in\mathbb{C}^4$.
By antilinearity, (S2) and $\gamma_\mu^*=\gamma_\mu$, the element $C=\mathrm{ev}^*\psi(\theta_{\mathrm{po}}x)$ satisfies
\[
\Theta_{\mathrm{po}}(C)=\sum_a\mathrm{ev}_a\cdot\Theta_{\mathrm{po}}\psi_a(\theta_{\mathrm{po}}x)=-\bar{\psi}(x)\gamma_\mu\mathrm{ev},
\]
and the element $C=\bar{\psi}(\theta_{\mathrm{po}}x)\,\mathrm{ev}$ satisfies
\[
\Theta_{\mathrm{po}}(C)=\sum_a\overline{\mathrm{ev}_a}\cdot\Theta_{\mathrm{po}}\bar{\psi}_a(\theta_{\mathrm{po}}x)=-(\gamma_\mu\mathrm{ev})^*\psi(x).
\]
Hence, for the four groups:
\begin{description}
\item[group 1] ($\gamma_\mu\mathrm{ev}=-\mathrm{ev}$): $\Theta_{\mathrm{po}}(C)\,C=\bar{\psi}(x)\,\mathrm{ev}\,\mathrm{ev}^*\psi(\theta_{\mathrm{po}}x)$; the sum over $k'$ is $\bar{\psi}(x)P^-\psi(\theta_{\mathrm{po}}x)$;
\item[group 2] ($\gamma_\mu\mathrm{ev}=\mathrm{ev}$): $\Theta_{\mathrm{po}}(C)\,C=-\bar{\psi}(x)\,\mathrm{ev}\,\mathrm{ev}^*\psi(\theta_{\mathrm{po}}x)$; the sum over $k'$ is $-\bar{\psi}(x)P^+\psi(\theta_{\mathrm{po}}x)$;
\item[group 3] ($\gamma_\mu\mathrm{ev}=\mathrm{ev}$): $\Theta_{\mathrm{po}}(C)\,C=-\left(\mathrm{ev}^*\psi(x)\right)\left(\bar{\psi}(\theta_{\mathrm{po}}x)\,\mathrm{ev}\right)=+\bar{\psi}(\theta_{\mathrm{po}}x)\,\mathrm{ev}\,\mathrm{ev}^*\psi(x)$, since two odd elements anticommute; the sum over $k'$ is $\bar{\psi}(\theta_{\mathrm{po}}x)P^+\psi(x)$;
\item[group 4] ($\gamma_\mu\mathrm{ev}=-\mathrm{ev}$): $\Theta_{\mathrm{po}}(C)\,C=\left(\mathrm{ev}^*\psi(x)\right)\left(\bar{\psi}(\theta_{\mathrm{po}}x)\,\mathrm{ev}\right)=-\bar{\psi}(\theta_{\mathrm{po}}x)\,\mathrm{ev}\,\mathrm{ev}^*\psi(x)$; the sum over $k'$ is $-\bar{\psi}(\theta_{\mathrm{po}}x)P^-\psi(x)$.
\end{description}
With the coefficients and the phases of the four groups, the sum of $c_i\cdot u_i\cdot\Theta_{\mathrm{po}}(C_i)\,C_i$ over the eight modes of the bond is
\[
U\cdot\bar{\psi}(x)\,3\left[(1+\rho)P^--(1-\rho)P^+\right]\psi(\theta_{\mathrm{po}}x)+\overline{U}\cdot\bar{\psi}(\theta_{\mathrm{po}}x)\,3\left[(1+\rho)P^+-(1-\rho)P^-\right]\psi(x),
\]
and this is $-(\text{terms of the bond})$.
Therefore
\[
-\mathrm{Xc}_{\mathrm{po}}=\sum_{i\in\mathrm{Ri}(\mathrm{po})}c_i\cdot u_i\cdot\mathrm{Ee}_i,\qquad\mathrm{Ee}_i:=\Theta_{\mathrm{po}}(C_i)\,C_i ,
\]
and each $\mathrm{Ee}_i$ is independent of the representative, as $\mathrm{chp}^2=1$.

\emph{Products.}
Each $\mathrm{Ee}_i$ is even; $C_i^2=0$, so that $\mathrm{Ee}_i^2=0$; and the $\mathrm{Ee}_i$ commute with each other.
For $I=\{i_1<\dots<i_d\}$,
\[
\Theta_{\mathrm{po}}(C_I)\,C_I=\Theta_{\mathrm{po}}(C_{i_d})\cdots\Theta_{\mathrm{po}}(C_{i_2})\left[\Theta_{\mathrm{po}}(C_{i_1})\,C_{i_1}\right]C_{i_2}\cdots C_{i_d};
\]
the bracket is even and can be moved to the far left, and induction gives $\Theta_{\mathrm{po}}(C_I)\,C_I=\mathrm{Ee}_{i_1}\cdots\mathrm{Ee}_{i_d}$.
Hence
\[
\exp(-\mathrm{Xc}_{\mathrm{po}})=\prod_{i\in\mathrm{Ri}(\mathrm{po})}\left(1+c_i\cdot u_i\cdot\mathrm{Ee}_i\right)=\sum_{I\subset\mathrm{Ri}(\mathrm{po})}c_I\cdot u_I\cdot\Theta_{\mathrm{po}}(C_I)\,C_I .
\qedhere
\]
\end{proof}

\subsubsection*{A.7. Slab matrices}

The slab matrices and their rules are those of \cite[Definition~5.21, Lemma~5.22]{QEDII}, here for even elements.

For $\mathrm{po}<\mathrm{po}'$ in $\tfrac12\mathbb{Z}$ and $\mathrm{Ph}\in\mathrm{Gr}_{(\mathrm{po},\mathrm{po}')}$ let $\mathrm{Ms}_{\mathrm{po},\mathrm{po}'}[\mathrm{Ph}]\colon\mathrm{Fk}(\mathrm{po})\to\mathrm{Fk}(\mathrm{po}')$ be the linear map with the entries
\[
\mathrm{Ms}_{\mathrm{po},\mathrm{po}'}[\mathrm{Ph}]_{I',I}:=\int_{(\mathrm{po},\mathrm{po}')}C^{(\mathrm{po})}_I\,\mathrm{Ph}\,\Theta_{\mathrm{po}'}\left(C^{(\mathrm{po}')}_{I'}\right),
\]
the column index $I$ being a subset of $\mathrm{Ri}(\mathrm{po})$ and the row index $I'$ a subset of $\mathrm{Ri}(\mathrm{po}')$.
Here $C^{(\mathrm{po})}_I$ lies in $\mathrm{Gr}_{(\mathrm{po},\mathrm{po}^+)}$ and $\Theta_{\mathrm{po}'}\left(C^{(\mathrm{po}')}_{I'}\right)$ in $\mathrm{Gr}_{(\mathrm{po}'^-,\mathrm{po}')}$;
both algebras are subalgebras of $\mathrm{Gr}_{(\mathrm{po},\mathrm{po}')}$, because $\mathrm{po}'-\mathrm{po}\ge\tfrac12$.

\begin{named}{Lemma A.3}
Let $\mathrm{po}<\mathrm{po}'<\mathrm{po}''$ be in $\tfrac12\mathbb{Z}$, $\mathrm{Ph}\in\mathrm{Gr}^{\mathrm{ev}}_{(\mathrm{po},\mathrm{po}')}$ and $\mathrm{Ph}'\in\mathrm{Gr}^{\mathrm{ev}}_{(\mathrm{po}',\mathrm{po}'')}$.
\begin{itemize}
\item[(F1)] The matrix $\mathrm{Ms}_{\mathrm{po},\mathrm{po}'}[\mathrm{Ph}]$ is even.
\item[(F2)] (Composition.) For every matrix $\mathrm{mt}$ on $\mathrm{Fk}(\mathrm{po}')$,
\[
\mathrm{Ms}_{\mathrm{po},\mathrm{po}''}\left[\mathrm{Ph}\,\mathrm{Ins}_{\mathrm{po}'}[\mathrm{mt}]\,\mathrm{Ph}'\right]=\mathrm{Ms}_{\mathrm{po}',\mathrm{po}''}[\mathrm{Ph}']\,\mathrm{mt}\,\mathrm{Ms}_{\mathrm{po},\mathrm{po}'}[\mathrm{Ph}].
\]
\item[(F3)] (Adjoint.) If $\mathrm{po}''=2\mathrm{po}'-\mathrm{po}$ and $\tau=2(\mathrm{po}'-\mathrm{po})e_\mu$, then
\[
\mathrm{Ms}_{\mathrm{po}',\mathrm{po}''}[\Theta_{\mathrm{po}'}\mathrm{Ph}]_{I+\tau,I'}=\overline{\mathrm{Ms}_{\mathrm{po},\mathrm{po}'}[\mathrm{Ph}]_{I',I}}\qquad\left(I\subset\mathrm{Ri}(\mathrm{po}),\ I'\subset\mathrm{Ri}(\mathrm{po}')\right),
\]
that is, $\mathrm{Ms}_{\mathrm{po}',\mathrm{po}''}[\Theta_{\mathrm{po}'}\mathrm{Ph}]=\mathrm{Ft}_\tau\,\mathrm{Ms}_{\mathrm{po},\mathrm{po}'}[\mathrm{Ph}]^*$.
\item[(F4)] (Translation.) For $v\in\mathbb{Z}e_\mu$,
\[
\mathrm{Ms}_{\mathrm{po}+v_\mu,\mathrm{po}'+v_\mu}[\mathrm{Tg}_v\mathrm{Ph}]=\mathrm{Ft}_v\,\mathrm{Ms}_{\mathrm{po},\mathrm{po}'}[\mathrm{Ph}]\,\mathrm{Ft}_v^{-1}.
\]
\item[(F5)] (Closure.) Let $n_\mu\ge1$ be an integer, $\mathrm{gv}:=n_\mu e_\mu$, and $\mathrm{po}_0<\mathrm{po}_1<\dots<\mathrm{po}_s=\mathrm{po}_0+n_\mu$ in $\tfrac12\mathbb{Z}$, with $s\ge1$.
Let $\mathrm{Gr}_n$, with its integral $\int_n$, be the algebra of Section~1.5 for the lattice $\mathrm{Gmp}+\mathbb{Z}\,\mathrm{gv}$ and the character that equals $\mathrm{chp}$ on $\mathrm{Gmp}$ and $-1$ at $\mathrm{gv}$.
Let $\mathrm{Ph}_k\in\mathrm{Gr}^{\mathrm{ev}}_{(\mathrm{po}_{k-1},\mathrm{po}_k)}$ for $k=1,\dots,s$, and let $\mathrm{mt}_k$ be even matrices on $\mathrm{Fk}(\mathrm{po}_k)$ for $k=0,\dots,s-1$.
Then, with $\mathrm{Ms}_k:=\mathrm{Ms}_{\mathrm{po}_{k-1},\mathrm{po}_k}[\mathrm{Ph}_k]$,
\begin{multline*}
\int_n\mathrm{Ins}_{\mathrm{po}_0}[\mathrm{mt}_0]\,\mathrm{Ph}_1\,\mathrm{Ins}_{\mathrm{po}_1}[\mathrm{mt}_1]\,\mathrm{Ph}_2\cdots\mathrm{Ins}_{\mathrm{po}_{s-1}}[\mathrm{mt}_{s-1}]\,\mathrm{Ph}_s\\
=\operatorname{tr}_{\mathrm{Fk}(\mathrm{po}_0)}\left(\mathrm{Ft}_{\mathrm{gv}}^{-1}\,\mathrm{Ms}_s\,\mathrm{mt}_{s-1}\,\mathrm{Ms}_{s-1}\cdots\mathrm{mt}_1\,\mathrm{Ms}_1\,\mathrm{mt}_0\right),
\end{multline*}
where the elements of $\mathrm{Gr}_{(\mathrm{po}_0-n_\mu,\mathrm{po}_0+n_\mu)}$ are read in $\mathrm{Gr}_n$ as described at the beginning of the proof of (F5).
\end{itemize}
\end{named}

\begin{proof}
(F1) The integrand has the parity of $\#|I|+\#|I'|$, and $\mathrm{top}$ has even degree.

(F2) By the definition of $\mathrm{Ins}_{\mathrm{po}'}[\mathrm{mt}]$,
\begin{multline*}
C^{(\mathrm{po})}_I\,\mathrm{Ph}\,\mathrm{Ins}_{\mathrm{po}'}[\mathrm{mt}]\,\mathrm{Ph}'\,\Theta_{\mathrm{po}''}\left(C^{(\mathrm{po}'')}_{I''}\right)\\
=\sum_{\mathrm{Ib},\mathrm{Ia}}\mathrm{mt}_{\mathrm{Ib},\mathrm{Ia}}\cdot\left[C^{(\mathrm{po})}_I\,\mathrm{Ph}\,\Theta_{\mathrm{po}'}\left(C^{(\mathrm{po}')}_{\mathrm{Ia}}\right)\right]\left[C^{(\mathrm{po}')}_{\mathrm{Ib}}\,\mathrm{Ph}'\,\Theta_{\mathrm{po}''}\left(C^{(\mathrm{po}'')}_{I''}\right)\right];
\end{multline*}
no factor has been moved.
The first bracket lies in $\mathrm{Gr}_{(\mathrm{po},\mathrm{po}')}$, the second one in $\mathrm{Gr}_{(\mathrm{po}',\mathrm{po}'')}$.
By (Ap8), $X((\mathrm{po},\mathrm{po}''))/\mathrm{Gmp}$ is the disjoint union of $X((\mathrm{po},\mathrm{po}'))/\mathrm{Gmp}$ and $X((\mathrm{po}',\mathrm{po}''))/\mathrm{Gmp}$.
So, by (Be1), the integral over $(\mathrm{po},\mathrm{po}'')$ equals
\[
\sum_{\mathrm{Ib},\mathrm{Ia}}\mathrm{Ms}_{\mathrm{po}',\mathrm{po}''}[\mathrm{Ph}']_{I'',\mathrm{Ib}}\cdot\mathrm{mt}_{\mathrm{Ib},\mathrm{Ia}}\cdot\mathrm{Ms}_{\mathrm{po},\mathrm{po}'}[\mathrm{Ph}]_{\mathrm{Ia},I}.
\]

(F3) By (Tq6),
\begin{align*}
C^{(\mathrm{po}')}_{I'}\,\Theta_{\mathrm{po}'}(\mathrm{Ph})\,\Theta_{\mathrm{po}''}\left(C^{(\mathrm{po}'')}_{I+\tau}\right)
&=C^{(\mathrm{po}')}_{I'}\,\Theta_{\mathrm{po}'}(\mathrm{Ph})\,\Theta_{\mathrm{po}'}\left(C^{(\mathrm{po})}_I\right)\\
&=\Theta_{\mathrm{po}'}\left(C^{(\mathrm{po})}_I\,\mathrm{Ph}\,\Theta_{\mathrm{po}'}\left(C^{(\mathrm{po}')}_{I'}\right)\right),
\end{align*}
because $\Theta_{\mathrm{po}'}$ is an anti-automorphism and $\Theta_{\mathrm{po}'}^2=1$.
The reflection $\theta_{\mathrm{po}'}$ maps $X((\mathrm{po},\mathrm{po}'))$ onto $X((\mathrm{po}',\mathrm{po}''))$, and (Tq4) gives the first form of (F3).
The operator form follows from $(\mathrm{Ft}_\tau\,\mathrm{Ms}^*)_{I+\tau,I'}=(\mathrm{Ms}^*)_{I,I'}=\overline{\mathrm{Ms}_{I',I}}$.

(F4) By (Tq4), (Tq2) and (Tq5),
\begin{align*}
\mathrm{Ms}_{\mathrm{po},\mathrm{po}'}[\mathrm{Ph}]_{I',I}
&=\int_{(\mathrm{po}+v_\mu,\mathrm{po}'+v_\mu)}\mathrm{Tg}_v\left(C_I\,\mathrm{Ph}\,\Theta_{\mathrm{po}'}(C_{I'})\right)\\
&=\int_{(\mathrm{po}+v_\mu,\mathrm{po}'+v_\mu)}(\mathrm{Tg}_vC_I)\,(\mathrm{Tg}_v\mathrm{Ph})\,\Theta_{\mathrm{po}'+v_\mu}(\mathrm{Tg}_vC_{I'})\\
&=\mathrm{Ms}_{\mathrm{po}+v_\mu,\mathrm{po}'+v_\mu}[\mathrm{Tg}_v\mathrm{Ph}]_{I'+v,I+v},
\end{align*}
which is (F4), entry by entry.

(F5) \emph{Representatives.}
By (Ap8), no site $x$ has its height $\mathrm{geo}(x)_\mu$ in $\mathrm{po}_0+n_\mu\mathbb{Z}$.
So every class of sites modulo $\mathrm{Gmp}+\mathbb{Z}\,\mathrm{gv}$ has exactly one representative modulo $\mathrm{Gmp}$ with height in $(\mathrm{po}_0,\mathrm{po}_0+n_\mu)$, and
\[
\mathrm{Gr}_n=\mathrm{Gr}_{(\mathrm{po}_0,\mathrm{po}_s)},\qquad\int_n=\int_{(\mathrm{po}_0,\mathrm{po}_s)} .
\]
A field at a site $x$ with $\mathrm{geo}(x)_\mu\in(\mathrm{po}_0-n_\mu,\mathrm{po}_0)$ is read as $\psi(x)=-\psi(x+\mathrm{gv})$, $\bar{\psi}(x)=-\bar{\psi}(x+\mathrm{gv})$.
Hence an element of $\mathrm{Gr}_{(\mathrm{po}_0-n_\mu,\mathrm{po}_0)}$ that is homogeneous of degree $d$ is read as $(-1)^d$ times its image under $\mathrm{Tg}_{\mathrm{gv}}$.
In particular, by (Tq2) and (Tq5), $\Theta_{\mathrm{po}_0}\left(C^{(\mathrm{po}_0)}_{\mathrm{Ia}}\right)$ is read as
\[
(-1)^{\#|\mathrm{Ia}|}\cdot\Theta_{\mathrm{po}_s}\left(\mathrm{Tg}_{\mathrm{gv}}C^{(\mathrm{po}_0)}_{\mathrm{Ia}}\right)=(-1)^{\#|\mathrm{Ia}|}\cdot\Theta_{\mathrm{po}_s}\left(C^{(\mathrm{po}_s)}_{\mathrm{Ia}+\mathrm{gv}}\right).
\]

\emph{Reordering.}
The left side of (F5) is
\[
\sum_{I_0,\mathrm{Ia}_0}(\mathrm{mt}_0)_{I_0,\mathrm{Ia}_0}\cdot\int_n\Theta_{\mathrm{po}_0}\left(C_{\mathrm{Ia}_0}\right)\mathrm{Rr}_{I_0},\qquad
\mathrm{Rr}_{I_0}:=C^{(\mathrm{po}_0)}_{I_0}\,\mathrm{Ph}_1\,\mathrm{Ins}_{\mathrm{po}_1}[\mathrm{mt}_1]\cdots\mathrm{Ph}_s .
\]
The insertions and the $\mathrm{Ph}_k$ are even;
so $\mathrm{Rr}_{I_0}$ is homogeneous of the parity of $\#|I_0|$, which on the support of $\mathrm{mt}_0$ is the parity of $\#|\mathrm{Ia}_0|$.
Moving the factor $\Theta_{\mathrm{po}_0}(C_{\mathrm{Ia}_0})$, of degree $\#|\mathrm{Ia}_0|$, from the far left to the far right costs the sign $(-1)^{\#|\mathrm{Ia}_0|\cdot\#|\mathrm{Ia}_0|}=(-1)^{\#|\mathrm{Ia}_0|}$.
The two signs $(-1)^{\#|\mathrm{Ia}_0|}$, the one from the antiperiodicity and the one from the reordering, cancel.
Hence the left side is
\begin{align*}
&\sum_{I_0,\mathrm{Ia}_0}(\mathrm{mt}_0)_{I_0,\mathrm{Ia}_0}\cdot\int_{(\mathrm{po}_0,\mathrm{po}_s)}C^{(\mathrm{po}_0)}_{I_0}\,\mathrm{Ph}_1\,\mathrm{Ins}_{\mathrm{po}_1}[\mathrm{mt}_1]\cdots\mathrm{Ph}_s\,\Theta_{\mathrm{po}_s}\left(C^{(\mathrm{po}_s)}_{\mathrm{Ia}_0+\mathrm{gv}}\right)\\
&\qquad=\sum_{I_0,\mathrm{Ia}_0}(\mathrm{mt}_0)_{I_0,\mathrm{Ia}_0}\cdot\left(\mathrm{Ms}_s\,\mathrm{mt}_{s-1}\cdots\mathrm{mt}_1\,\mathrm{Ms}_1\right)_{\mathrm{Ia}_0+\mathrm{gv},I_0},
\end{align*}
by the definition of $\mathrm{Ms}_{\mathrm{po}_0,\mathrm{po}_s}$ and by (F2), applied $s-1$ times.
For a matrix $B$ from $\mathrm{Fk}(\mathrm{po}_0)$ to $\mathrm{Fk}(\mathrm{po}_s)$ one has $(\mathrm{Ft}_{\mathrm{gv}}^{-1}B)_{\mathrm{Ia}_0,I_0}=B_{\mathrm{Ia}_0+\mathrm{gv},I_0}$.
So the sum is $\operatorname{tr}\left(\mathrm{Ft}_{\mathrm{gv}}^{-1}\,\mathrm{Ms}_s\,\mathrm{mt}_{s-1}\cdots\mathrm{mt}_1\,\mathrm{Ms}_1\,\mathrm{mt}_0\right)$.
\end{proof}

\begin{named}{(Bd4)}[the action on a closed chain]
Let $n_\mu\ge1$, $\mathrm{gv}=n_\mu e_\mu$, let $\mathrm{po}_0<\mathrm{po}_1<\dots<\mathrm{po}_s=\mathrm{po}_0+n_\mu$ be a chain with $s\ge1$, and let $\mathrm{Gr}_n$ be as in (F5).
Put $\Gamma(n):=\mathrm{Gmp}+\mathbb{Z}\,\mathrm{gv}$, let $A$ be a real one-cochain on $\mathbb{Z}^3$ that is periodic under $\Gamma(n)$, and let
\[
S_n(A):=\sum_x\left[M_{\mathrm{ch}}\cdot\bar{\psi}(x)\psi(x)-\sum_{i=0}^{3}\bar{\psi}(x)M_i(x)\,U_A(x,x^{(i)})\,\psi(x^{(i)})\right]\ \in\ \mathrm{Gr}_n,
\]
the sum extending over representatives $x$ of the sites modulo $\Gamma(n)$; this is the display of Section~1.5, see (I.2.5).
Then, in $\mathrm{Gr}_n$,
\[
S_n(A)=S_{(\mathrm{po}_0,\mathrm{po}_s)}(A)+\mathrm{Xc}_{\mathrm{po}_0}(A)=\sum_{k=1}^{s}S_{(\mathrm{po}_{k-1},\mathrm{po}_k)}(A)+\sum_{k=0}^{s-1}\mathrm{Xc}_{\mathrm{po}_k}(A),
\]
where $\mathrm{Xc}_{\mathrm{po}_0}(A)$, an element of $\mathrm{Gr}_{(\mathrm{po}_0^-,\mathrm{po}_0^+)}$, is read in $\mathrm{Gr}_n$ as in (F5):
a field at a site $x$ with $\mathrm{geo}(x)_\mu\in(\mathrm{po}_0-n_\mu,\mathrm{po}_0)$ is minus the field at $x+\mathrm{gv}$.
\end{named}

\begin{proof}[Proof of (Bd4)]
Each term of $S_n(A)$ is independent of the representative.
Indeed, for $v\in\Gamma(n)$ one has $(x+v)^{(i)}=x^{(i)}+v$ (Bd0);
$M_i(x+v)=M_i(x)$, since $M_i$ depends on the label only (I.2.1);
$U_A(x+v,x^{(i)}+v)=U_A(x,x^{(i)})$, since $A$ is periodic;
and $\bar{\psi}(x+v)$, $\psi(x^{(i)}+v)$ differ from $\bar{\psi}(x)$, $\psi(x^{(i)})$ by the same sign.
By the paragraph ``Representatives'' of the proof of (F5), $X((\mathrm{po}_0,\mathrm{po}_s))/\mathrm{Gmp}$ is a set of representatives of the sites modulo $\Gamma(n)$, and $\mathrm{Gr}_n=\mathrm{Gr}_{(\mathrm{po}_0,\mathrm{po}_s)}$.
So $S_n(A)$ is the sum over $x\in X((\mathrm{po}_0,\mathrm{po}_s))/\mathrm{Gmp}$, a field $\psi(x^{(i)})$ at a site outside $X((\mathrm{po}_0,\mathrm{po}_s))$ being read by the relations of $\mathrm{Gr}_n$.
The terms $M_{\mathrm{ch}}\cdot\bar{\psi}(x)\psi(x)$ and the terms of the pairs $(x,x^{(i)})$ with $x^{(i)}\in X((\mathrm{po}_0,\mathrm{po}_s))$ are the terms of $S_{(\mathrm{po}_0,\mathrm{po}_s)}(A)$ (Bd2), each of them exactly once.

Let now $x\in X((\mathrm{po}_0,\mathrm{po}_s))$ and $x^{(i)}\notin X((\mathrm{po}_0,\mathrm{po}_s))$.
By (Ap8), the height of $x^{(i)}$ is either $<\mathrm{po}_0$ or $>\mathrm{po}_s$.

\emph{First case: the height of $x^{(i)}$ is $<\mathrm{po}_0$.}
Then $\{x,x^{(i)}\}$ crosses $(y_\mu,\mathrm{po}_0)$, and $x$ is its upper end.
The term of $(x,x^{(i)})$ is the term of $\mathrm{Xc}_{\mathrm{po}_0}(A)$ with $\bar{\psi}$ at the upper end, the field $\psi(x^{(i)})$ at the lower end being read in $\mathrm{Gr}_n$.

\emph{Second case: the height of $x^{(i)}$ is $>\mathrm{po}_s$.}
Then $\{x,x^{(i)}\}$ crosses $(y_\mu,\mathrm{po}_s)$, and $x$ is its lower end.
The pair $\{x-\mathrm{gv},x^{(i)}-\mathrm{gv}\}$ is a bond of the same type that crosses $(y_\mu,\mathrm{po}_0)$, with the lower end $x-\mathrm{gv}$;
moreover $M_i(x-\mathrm{gv})=M_i(x)$ and $U_A(x-\mathrm{gv},x^{(i)}-\mathrm{gv})=U_A(x,x^{(i)})$, and in $\mathrm{Gr}_n$
\[
\bar{\psi}(x-\mathrm{gv})=-\bar{\psi}(x),\qquad\psi(x^{(i)}-\mathrm{gv})=-\psi(x^{(i)}).
\]
So the term of $(x,x^{(i)})$ is the term of $\mathrm{Xc}_{\mathrm{po}_0}(A)$ with $\bar{\psi}$ at the lower end, read in $\mathrm{Gr}_n$.

By (Bd1), the upper ends of the bonds that cross $(y_\mu,\mathrm{po}_0)$ lie in $X((\mathrm{po}_0,\mathrm{po}_0^+))$, and the lower ends of the bonds that cross $(y_\mu,\mathrm{po}_s)$ lie in $X((\mathrm{po}_s^-,\mathrm{po}_s))$;
both sets are subsets of $X((\mathrm{po}_0,\mathrm{po}_s))$.
The map $x\mapsto x-\mathrm{gv}$ maps the lower ends of the bonds that cross $(y_\mu,\mathrm{po}_s)$ bijectively onto the lower ends of the bonds that cross $(y_\mu,\mathrm{po}_0)$.
So every term of $\mathrm{Xc}_{\mathrm{po}_0}(A)$ arises exactly once, and the first equality holds.
The second equality is (Bd3):
a bond with both ends in $X((\mathrm{po}_0,\mathrm{po}_s))$ either has both ends in one of the $s$ slabs or crosses exactly one of the inner planes $\mathrm{po}_1,\dots,\mathrm{po}_{s-1}$, and no sign occurs there.
\end{proof}

\subsubsection*{A.8. The photon at a plane}

The objects of this section are those of \cite[Section~5.6]{QEDII}.

Let $\mathrm{po}\in\tfrac12\mathbb{Z}$; all sets of cells that occur below are finite.
Put
\[
\mathrm{Ze}_{\mathrm{po}}:=[0,2\pi)^{\mathrm{Cu}^3(\mathrm{po})}\times\mathbb{Z}^{\mathrm{Cu}^2(\mathrm{po})},\qquad\boldsymbol{\zeta}=(\varphi,\mathrm{kz}),
\]
and, for $p\in\mathrm{Cu}^2(\mathrm{po})$,
\[
\kappa_p(\boldsymbol{\zeta}):=\frac{(d_{\mathrm{po}}^*\varphi)(p)-2\pi\,\mathrm{kz}_p}{r_q},\qquad(d_{\mathrm{po}}^*\varphi)(p)=\sum_{\mathrm{qb}\in\mathrm{Cu}^3(\mathrm{po})}[\mathrm{qb}:p]\,\varphi_{\mathrm{qb}} .
\]
Let $\mathrm{dn}_{\mathrm{po}}$ be the measure on $\mathrm{Ze}_{\mathrm{po}}$ given by
\[
\mathrm{dn}_{\mathrm{po}}(\boldsymbol{\zeta}):=\prod_{\mathrm{qb}\in\mathrm{Cu}^3(\mathrm{po})}\frac{\mathrm{d}\varphi_{\mathrm{qb}}}{2\pi}\cdot\prod_{p\in\mathrm{Cu}^2(\mathrm{po})}\frac{(2\pi)^{\frac12}}{r_q}\exp\left(-\frac{\kappa_p(\boldsymbol{\zeta})^2}{2}\right),
\]
with the counting measure in $\mathrm{kz}$, and put
\[
\mathrm{ef}_{\mathrm{po}}(\boldsymbol{\zeta};\omega):=\exp\left(i\sum_{\mathrm{qb}\in\mathrm{Cu}^3(\mathrm{po})}\varphi_{\mathrm{qb}}\,a_{\mathrm{qb}}(\mathfrak{z})-i\sum_{p\in\mathrm{Cu}^2(\mathrm{po})}\kappa_p(\boldsymbol{\zeta})\,\alpha_p(A)\right),
\]
with $\alpha_p$ and $a_{\mathrm{qb}}$ as in (Ap3).
Thus $\mathrm{ef}_{\mathrm{po}}$ reads $\omega$ on the cells $\boldsymbol{\sigma}^+$, $\boldsymbol{\sigma}\in\mathrm{Cu}^2(\mathrm{po})\cup\mathrm{Cu}^3(\mathrm{po})$, which lie in the plane $\mathrm{po}^+$.
For $I\subset\mathrm{Ri}(\mathrm{po})$ and $b\in\mathrm{Cu}^1(\mathrm{po})$ let
\[
\mathrm{Qc}_I(b):=\#|\{i\in I\text{ over }b\text{, of group 1 or 2}\}|-\#|\{i\in I\text{ over }b\text{, of group 3 or 4}\}|;
\]
then $u_I(A)=\exp\left(iq\sum_{b\in\mathrm{Cu}^1(\mathrm{po})}\mathrm{Qc}_I(b)\,A_b\right)$, by A.6.
Gauss' law is expressed by
\[
\mathrm{Ga}_{\mathrm{po}}(\boldsymbol{\zeta},I):=r_q^{\#|\mathrm{Cu}^1(\mathrm{po})|}\prod_{b\in\mathrm{Cu}^1(\mathrm{po})}\mathbf{1}\left[(d_{\mathrm{po}}^*\mathrm{kz})(b)+\mathrm{Qc}_I(b)=0\right].
\]

At a site plane the three sets $\mathrm{Cu}^1(\mathrm{po})$, $\mathrm{Cu}^2(\mathrm{po})$ and $\mathrm{Cu}^3(\mathrm{po})$ are empty:
$\mathrm{Ze}_{\mathrm{po}}$ is a point, $\mathrm{dn}_{\mathrm{po}}$ is the measure of mass one at this point, and $\mathrm{ef}_{\mathrm{po}}=1$, $\mathrm{Ga}_{\mathrm{po}}=1$, $u_I=1$.
At a link plane, (K1) with $y=0$ gives
\[
\frac{(2\pi)^{\frac12}}{r_q}\sum_{\mathrm{kz}_p}\exp\left(-\frac{\kappa_p^2}{2}\right)=\sum_j\exp\left(-\frac{r_q^2j^2}{2}+ij\,(d_{\mathrm{po}}^*\varphi)(p)\right)\le\mathrm{tz}:=\sum_j\exp\left(-\frac{r_q^2j^2}{2}\right).
\]
So $\mathrm{dn}_{\mathrm{po}}$ is a finite positive measure, of mass at most $\mathrm{tz}^{\#|\mathrm{Cu}^2(\mathrm{po})|}$.

\begin{named}{Lemma A.4}[{crossing; cf.\ \cite[Lemma~5.24]{QEDII}}]
Let $\mathrm{po}\in\tfrac12\mathbb{Z}$ and $I\subset\mathrm{Ri}(\mathrm{po})$, and let $\omega=(A,\mathfrak{z})$ be given on the cells $\boldsymbol{\sigma}^+$ and $\boldsymbol{\sigma}^-$, $\boldsymbol{\sigma}\in\mathrm{Cu}^2(\mathrm{po})\cup\mathrm{Cu}^3(\mathrm{po})$.
Then
\begin{multline*}
\sum_{\mathfrak{z}\in\mathbb{Z}^{\mathrm{Cu}^2(\mathrm{po})}}\int_{[0,r_q)^{\mathrm{Cu}^1(\mathrm{po})}}\mathrm{d}A\,\prod_{p\in\mathrm{Cu}^2(\mathrm{po})}\exp\left(-\frac{F_p^2}{2}\right)\cdot\prod_{\mathrm{qb}\in\mathrm{Cu}^3(\mathrm{po})}\mathbf{1}\left[(d\mathfrak{z})(\mathrm{qb})=0\right]\cdot u_I(A)\\
=\int\mathrm{dn}_{\mathrm{po}}(\boldsymbol{\zeta})\,\mathrm{Ga}_{\mathrm{po}}(\boldsymbol{\zeta},I)\cdot\mathrm{ef}_{\mathrm{po}}(\boldsymbol{\zeta};\omega)\cdot\overline{\mathrm{ef}_{\mathrm{po}}(\boldsymbol{\zeta};\theta_{\mathrm{po}}^*\omega)} .
\end{multline*}
On the left, $F_p=(dA)(p)+r_q\cdot\mathfrak{z}_p$ and $d\mathfrak{z}$ are formed with $\omega$ and with the variables of summation and integration.
Both sides converge absolutely, uniformly in $\omega$.
At a site plane both sides are equal to $1$.
\end{named}

\begin{proof}
Let $\mathrm{po}$ be a link plane, and write $\mathrm{Ac}$ and $\mathrm{mc}$ for the variables on $\mathrm{Cu}^1(\mathrm{po})$ and on $\mathrm{Cu}^2(\mathrm{po})$.
By (Ap3),
\[
F_p=y_p+(d_{\mathrm{po}}\mathrm{Ac})(p)+r_q\cdot\mathrm{mc}_p,\qquad y_p:=\alpha_p(A)-\alpha_p(\theta_{\mathrm{po}}^*A),
\]
and $(d\mathfrak{z})(\mathrm{qb})=t_{\mathrm{qb}}+(d_{\mathrm{po}}\mathrm{mc})(\mathrm{qb})$ with the integer $t_{\mathrm{qb}}:=a_{\mathrm{qb}}(\mathfrak{z})-a_{\mathrm{qb}}(\theta_{\mathrm{po}}^*\mathfrak{z})$.

(1) For integers $j$ one has $\mathbf{1}[j=0]=\int_0^{2\pi}\frac{\mathrm{d}\varphi}{2\pi}\,\exp(i\varphi j)$.
Therefore
\[
\prod_{\mathrm{qb}\in\mathrm{Cu}^3(\mathrm{po})}\mathbf{1}\left[(d\mathfrak{z})(\mathrm{qb})=0\right]=\int\prod_{\mathrm{qb}\in\mathrm{Cu}^3(\mathrm{po})}\frac{\mathrm{d}\varphi_{\mathrm{qb}}}{2\pi}\,\exp\left(i\langle\varphi,t\rangle\right)\exp\left(i\langle d_{\mathrm{po}}^*\varphi,\mathrm{mc}\rangle\right).
\]
The integrand of the threefold operation (sum over $\mathrm{mc}$, integral over $\mathrm{Ac}$, integral over $\varphi$) is bounded in modulus by $\prod_p\exp\left(-\tfrac12F_p^2\right)$,
whose sum over $\mathrm{mc}$ is at most $\boldsymbol{\Theta}_{1/2}^{\#|\mathrm{Cu}^2(\mathrm{po})|}$ by (K2);
and $\mathrm{Ac}$ and $\varphi$ range over sets of finite measure.
So Fubini's theorem applies.

(2) For fixed $\mathrm{Ac}$ and $\varphi$ the sum over $\mathrm{mc}$ factorises over $p$, and (K1) with $y=y_p+(d_{\mathrm{po}}\mathrm{Ac})(p)$ and $t=(d_{\mathrm{po}}^*\varphi)(p)$ gives
\begin{multline*}
\sum_{\mathrm{mc}_p}\exp\left(-\frac{F_p^2}{2}+i\,\mathrm{mc}_p\,(d_{\mathrm{po}}^*\varphi)(p)\right)\\
=\frac{(2\pi)^{\frac12}}{r_q}\sum_{\mathrm{kz}_p}\exp\left(-\frac{\kappa_p^2}{2}\right)\exp\left(-i\kappa_py_p\right)\exp\left(-i\kappa_p\,(d_{\mathrm{po}}\mathrm{Ac})(p)\right).
\end{multline*}

(3) The sum over $\mathrm{kz}$ converges absolutely, uniformly in $(\mathrm{Ac},\varphi)$, with the bound $\mathrm{tz}$ for each $p$; Fubini's theorem applies again.
The dependence on $\mathrm{Ac}$ is through the factor
\[
\exp\left(-i\langle d_{\mathrm{po}}^*\kappa,\mathrm{Ac}\rangle+iq\langle\mathrm{Qc}_I,\mathrm{Ac}\rangle\right).
\]
By (Ap4),
\[
d_{\mathrm{po}}^*\kappa=\frac{d_{\mathrm{po}}^*d_{\mathrm{po}}^*\varphi-2\pi\,d_{\mathrm{po}}^*\mathrm{kz}}{r_q}=-q\,d_{\mathrm{po}}^*\mathrm{kz}.
\]
Hence
\[
\int_{[0,r_q)^{\mathrm{Cu}^1(\mathrm{po})}}\mathrm{d}\mathrm{Ac}\,\exp\left(iq\langle d_{\mathrm{po}}^*\mathrm{kz}+\mathrm{Qc}_I,\mathrm{Ac}\rangle\right)=\mathrm{Ga}_{\mathrm{po}}(\boldsymbol{\zeta},I),
\]
because $d_{\mathrm{po}}^*\mathrm{kz}+\mathrm{Qc}_I$ is integer valued and $q\,r_q=2\pi$.

(4) Finally,
\begin{align*}
\exp\left(i\langle\varphi,t\rangle\right)\prod_p\exp\left(-i\kappa_py_p\right)
&=\exp\left(i\sum_{\mathrm{qb}}\varphi_{\mathrm{qb}}\,a_{\mathrm{qb}}(\mathfrak{z})-i\sum_p\kappa_p\,\alpha_p(A)\right)\\
&\qquad\cdot\exp\left(-i\sum_{\mathrm{qb}}\varphi_{\mathrm{qb}}\,a_{\mathrm{qb}}(\theta_{\mathrm{po}}^*\mathfrak{z})+i\sum_p\kappa_p\,\alpha_p(\theta_{\mathrm{po}}^*A)\right)\\
&=\mathrm{ef}_{\mathrm{po}}(\boldsymbol{\zeta};\omega)\cdot\overline{\mathrm{ef}_{\mathrm{po}}(\boldsymbol{\zeta};\theta_{\mathrm{po}}^*\omega)} .
\end{align*}
Inserting (2), (3) and (4) into (1) gives the identity.
\end{proof}

\begin{named}{Lemma A.5}[{properties of $\mathrm{ef}_{\mathrm{po}}$; cf.\ \cite[Lemma~5.25]{QEDII}}]
Let $\mathrm{po}\in\tfrac12\mathbb{Z}$.
\begin{itemize}
\item[(i)] $|\mathrm{ef}_{\mathrm{po}}|=1$, and $\mathrm{ef}_{\mathrm{po}}$ is continuous in $(\varphi,A)$ for fixed $(\mathrm{kz},\mathfrak{z})$.
\item[(ii)] $\mathrm{ef}_{\mathrm{po}}(\boldsymbol{\zeta};\omega+[m])=\mathrm{ef}_{\mathrm{po}}(\boldsymbol{\zeta};\omega)$ for every integer one-cochain $m$ that vanishes on the links contained in $\{y_\mu\le\mathrm{po}\}$ and on the links of $\mathrm{Cu}^1(\mathrm{po})$.
Consequently $\omega\mapsto\mathrm{ef}_{\mathrm{po}}(\boldsymbol{\zeta};\theta_{\mathrm{po}}^*\omega)$ is invariant under $\omega\mapsto\omega+[m]$ for every integer one-cochain $m$ that vanishes on the links contained in $\{y_\mu\ge\mathrm{po}\}$ and on the links of $\mathrm{Cu}^1(\mathrm{po})$.
\item[(iii)] Let $v\in\mathbb{Z}e_\mu$, and put $(\tau_v^*u)(\boldsymbol{\sigma}):=u(\boldsymbol{\sigma}+v)$.
Then
\[
\mathrm{ef}_{\mathrm{po}+v_\mu}(\boldsymbol{\zeta}';\omega')=\mathrm{ef}_{\mathrm{po}}(\tau_v^*\boldsymbol{\zeta}';\tau_v^*\omega'),\qquad\kappa_{p+v}(\boldsymbol{\zeta}')=\kappa_p(\tau_v^*\boldsymbol{\zeta}'),
\]
the map $\boldsymbol{\zeta}'\mapsto\tau_v^*\boldsymbol{\zeta}'$ maps $\mathrm{dn}_{\mathrm{po}+v_\mu}$ to $\mathrm{dn}_{\mathrm{po}}$, and $\mathrm{Ga}_{\mathrm{po}+v_\mu}(\boldsymbol{\zeta}',I+v)=\mathrm{Ga}_{\mathrm{po}}(\tau_v^*\boldsymbol{\zeta}',I)$.
\end{itemize}
\end{named}

\begin{proof}
(i) The exponent of $\mathrm{ef}_{\mathrm{po}}$ is $i$ times a real number.
For fixed $(\mathrm{kz},\mathfrak{z})$ it is a polynomial in the variables $\varphi_{\mathrm{qb}}$ and $A(p^+)$: $\kappa_p(\boldsymbol{\zeta})$ is affine in $\varphi$, and $\alpha_p(A)$ is linear in $A$.

(ii) One has $\alpha_p(A+r_q\cdot m)=\alpha_p(A)+r_q\cdot\alpha_p(m)$ and $r_q\kappa_p=(d_{\mathrm{po}}^*\varphi)(p)-2\pi\,\mathrm{kz}_p$, and $\alpha_p(m)$ is an integer;
so $\exp\left(-i\kappa_p\cdot r_q\cdot\alpha_p(m)\right)=\exp\left(-i\,(d_{\mathrm{po}}^*\varphi)(p)\cdot\alpha_p(m)\right)$.
Moreover $a_{\mathrm{qb}}(\mathfrak{z}-dm)=a_{\mathrm{qb}}(\mathfrak{z})-a_{\mathrm{qb}}(dm)$.
So $\mathrm{ef}_{\mathrm{po}}$ is multiplied by
\[
\exp\left(-i\sum_{\mathrm{qb}}\varphi_{\mathrm{qb}}\left[a_{\mathrm{qb}}(dm)+\sum_{p\in\mathrm{Cu}^2(\mathrm{po})}[\mathrm{qb}:p]\,\alpha_p(m)\right]\right)=1,
\]
by (Ap5).
For the second statement: $\theta_{\mathrm{po}}^*(\omega+[m])=\theta_{\mathrm{po}}^*\omega+[\theta_{\mathrm{po}}^*m]$,
and $\theta_{\mathrm{po}}^*m$ vanishes on the links contained in $\{y_\mu\le\mathrm{po}\}$ and on the cut links, since $\theta_{\mathrm{po}}$ maps $\mathrm{Cu}^1(\mathrm{po})$ to itself.

(iii) The translation $\tau_v$ maps $\mathrm{Mi}(\mathrm{po})$, $\mathrm{Cu}(\mathrm{po})$ and $\boldsymbol{\sigma}^+$ onto $\mathrm{Mi}(\mathrm{po}+v_\mu)$, $\mathrm{Cu}(\mathrm{po}+v_\mu)$ and $(\boldsymbol{\sigma}+v)^+$, with $\epsilon(\tau_v,\boldsymbol{\sigma})=1$ and with the same incidence numbers (I2).
Moreover $\mathrm{Qc}_{I+v}(b+v)=\mathrm{Qc}_I(b)$: the group of a mode and the link under its bond are transported by $i\mapsto i+v$.
\end{proof}

\subsubsection*{A.9. Mirror variables, Hilbert spaces, translations, slab operators}

The mirror spaces and the kernels below are those of \cite[Definitions~5.27 and~5.28]{QEDII}, in the setting of Section~A.0.
For $\mathrm{po}\in\tfrac12\mathbb{Z}$ put
\[
\Xi_{\mathrm{po}}:=\left\{\xi=(\mathrm{mi},\boldsymbol{\zeta},I):\ \mathrm{mi}=(A,\mathfrak{z})\in[0,r_q)^{\mathrm{Mi}^1(\mathrm{po})}\times\mathbb{Z}^{\mathrm{Mi}^2(\mathrm{po})},\ \boldsymbol{\zeta}\in\mathrm{Ze}_{\mathrm{po}},\ I\subset\mathrm{Ri}(\mathrm{po})\right\},
\]
\[
\mathrm{d}\,\mathrm{mq}_{\mathrm{po}}(\xi):=\prod_{b\in\mathrm{Mi}^1(\mathrm{po})}\mathrm{d}A_b\cdot\prod_{p\in\mathrm{Mi}^2(\mathrm{po})}\exp\left(-\frac12F_p(\mathrm{mi})^2\right)\cdot\mathrm{dn}_{\mathrm{po}}(\boldsymbol{\zeta})\cdot\mathrm{Ga}_{\mathrm{po}}(\boldsymbol{\zeta},I),
\]
with the counting measure in $\mathfrak{z}$, in $\mathrm{kz}$ and in $I$.
Here, for $\mathrm{mi}=(A,\mathfrak{z})$ with $A\in\mathbb{R}^{\mathrm{Mi}^1(\mathrm{po})}$ and $\mathfrak{z}\in\mathbb{Z}^{\mathrm{Mi}^2(\mathrm{po})}$,
\[
F(\mathrm{mi}):=dA+r_q\mathfrak{z}\qquad\text{on }\mathrm{Mi}^2(\mathrm{po}),
\]
and $F_p(\mathrm{mi})$ is its value at $p\in\mathrm{Mi}^2(\mathrm{po})$;
this is defined, because the faces of a plaquette contained in the plane $\{y_\mu=\mathrm{po}\}$ are links contained in that plane.
Consequently $F_p(\mathrm{mi})=F_p(\omega)$ for every configuration $\omega$ of a closed slab whose part on $\mathrm{Mi}(\mathrm{po})$ is $\mathrm{mi}$.
Put
\[
\mathrm{Hs}(\mathrm{po}):=L^2\left(\Xi_{\mathrm{po}},\mathrm{mq}_{\mathrm{po}}\right).
\]
For $v\in\mathbb{Z}e_\mu$ the operator $V_v\colon\mathrm{Hs}(\mathrm{po})\to\mathrm{Hs}(\mathrm{po}+v_\mu)$ is defined by
\[
(V_vf)(\mathrm{mi}',\boldsymbol{\zeta}',I+v):=f(\tau_v^*\mathrm{mi}',\tau_v^*\boldsymbol{\zeta}',I).
\]

\emph{Site plane $\mathrm{po}=t$.}
Here $\xi=(A_t,\mathfrak{z}_t,I)$, where $A_t$ denotes the link variables and $\mathfrak{z}_t$ the integers on the plaquettes of the plane.
Let $F(A_t,\mathfrak{z}_t):=F(\mathrm{mi})$ for $\mathrm{mi}=(A_t,\mathfrak{z}_t)$, so that
\[
\left\|F(A_t,\mathfrak{z}_t)\right\|^2=\sum_{p\in\mathrm{Mi}^2(t)}F_p(\mathrm{mi})^2;
\]
then
\[
\mathrm{d}\,\mathrm{mq}_{\mathrm{po}}=\mathrm{d}A_t\cdot\exp\left(-\frac12\left\|F(A_t,\mathfrak{z}_t)\right\|^2\right),
\]
with the counting measure in $\mathfrak{z}_t$ and in $I$;
the set $I$ ranges over the subsets of the set of the $64v_s$ modes of the bonds of type~$0$.

\emph{Link plane $\mathrm{po}=t+\frac12$.}
Here $\xi=(\varphi,\mathrm{kz},I)$, where $\varphi$ consists of one angle for each cut cube and $\mathrm{kz}$ of one integer for each cut plaquette, and
\[
\mathrm{d}\,\mathrm{mq}_{\mathrm{po}}=\mathrm{dn}_{\mathrm{po}}(\boldsymbol{\zeta})\cdot\mathrm{Ga}_{\mathrm{po}}(\boldsymbol{\zeta},I);
\]
the set $I$ ranges over the subsets of the set of the $64v_s$ modes of the bonds of type~$3$, restricted by Gauss' law.

The measure $\mathrm{mq}_{\mathrm{po}}$ is a finite positive Borel measure on a countable union of boxes.
Indeed, by (K2) the mass of the first two factors is at most $r_q^{\#|\mathrm{Mi}^1(\mathrm{po})|}\cdot\boldsymbol{\Theta}_{1/2}^{\#|\mathrm{Mi}^2(\mathrm{po})|}$;
the mass of $\mathrm{dn}_{\mathrm{po}}$ is at most $\mathrm{tz}^{\#|\mathrm{Cu}^2(\mathrm{po})|}$;
one has $\mathrm{Ga}_{\mathrm{po}}\le r_q^{\#|\mathrm{Cu}^1(\mathrm{po})|}$;
and there are $2^{64v_s}$ sets $I$.
Hence $\mathrm{Hs}(\mathrm{po})$ is separable.

Let $\mathrm{po}<\mathrm{po}'$ be elements of $\tfrac12\mathbb{Z}$ and let $G\in\mathrm{Ig}[\mathrm{po},\mathrm{po}']$.
For $\xi=(\mathrm{mi},\boldsymbol{\zeta},I)\in\Xi_{\mathrm{po}}$ and $\xi'=(\mathrm{mi}',\boldsymbol{\zeta}',I')\in\Xi_{\mathrm{po}'}$ write $\omega=\mathrm{mi}\vee\mathrm{omi}\vee\mathrm{mi}'$, and put
\begin{multline*}
\mathrm{Cs}[G](\xi',\xi):=\left(c_{I'}\cdot c_I\right)^{\frac12}\cdot\mathrm{Ii}_{\mathrm{po},\mathrm{po}'}\Bigl\{\mathrm{wb}_{\mathrm{po},\mathrm{po}'}(\omega)\cdot\mathrm{ef}_{\mathrm{po}}(\boldsymbol{\zeta};\omega)\cdot\overline{\mathrm{ef}_{\mathrm{po}'}\left(\boldsymbol{\zeta}';\theta_{\mathrm{po}'}^*\omega\right)}\\
\cdot\mathrm{Ms}_{\mathrm{po},\mathrm{po}'}\left[G(\omega)\cdot\exp\left(-S_{(\mathrm{po},\mathrm{po}')}(A)\right)\right]_{I',I}\Bigr\}
\end{multline*}
and
\[
(\mathrm{Cs}[G]f)(\xi'):=\int\mathrm{d}\,\mathrm{mq}_{\mathrm{po}}(\xi)\,\mathrm{Cs}[G](\xi',\xi)\cdot f(\xi).
\]
The expression in braces is read as $0$ where $\mathrm{wb}_{\mathrm{po},\mathrm{po}'}(\omega)=0$;
the function $G$ is defined on $\mathrm{Cfg}[\mathrm{po},\mathrm{po}']$ only.
By (Ap7), $\mathrm{ef}_{\mathrm{po}}(\boldsymbol{\zeta};\,\cdot\,)$ reads $\omega$ on cells of the plane $\mathrm{po}^+$, and $\mathrm{ef}_{\mathrm{po}'}(\boldsymbol{\zeta}';\theta_{\mathrm{po}'}^*\,\cdot\,)$ reads $\omega$ on cells of the plane $\mathrm{po}'^-$;
these are cells of the closed slab.
For $G\in\mathcal{A}_{\mu,n}[\mathrm{po},\mathrm{po}']$ the operator $\mathrm{Cs}[G]$ is, by definition, the operator $\mathrm{Cs}$ of the image of $G$ in $\mathrm{Ig}[\mathrm{po},\mathrm{po}']$ (Lemma~A.2).

\emph{The four kinds of pairs.}
The kind of $\mathrm{po}$ decides whether the lower end of the kernel carries the variables $(A_t,\mathfrak{z}_t)$ and the factor $\mathrm{ef}_{\mathrm{po}}=1$ (site plane), or the variables $(\varphi,\mathrm{kz})$ and the factor $\mathrm{ef}_{\mathrm{po}}(\boldsymbol{\zeta};\omega)$ (link plane).
The kind of $\mathrm{po}'$ decides the same for the upper end;
the two ends are independent of each other.
In all four cases the fermionic part is the matrix $\mathrm{Ms}_{\mathrm{po},\mathrm{po}'}[\,\cdot\,]$ from $\mathrm{Fk}(\mathrm{po})$ to $\mathrm{Fk}(\mathrm{po}')$, with the weights $\left(c_{I'}\cdot c_I\right)^{\frac12}$.

\subsubsection*{A.10. The theorem}

\begin{named}{Theorem A}
Let $0<\rho<1$ and $\mathrm{kk}:=1$, and let $\mathrm{po}<\mathrm{po}'<\mathrm{po}''$ be elements of $\tfrac12\mathbb{Z}$.
\begin{itemize}
\item[(SO0)] For $v,v'\in\mathbb{Z}e_\mu$ the operator $V_v$ is unitary, and $V_vV_{v'}=V_{v+v'}$.
The map $G\mapsto\mathrm{Cs}[G]$ is linear from $\mathrm{Ig}[\mathrm{po},\mathrm{po}']$ to the Hilbert--Schmidt operators from $\mathrm{Hs}(\mathrm{po})$ to $\mathrm{Hs}(\mathrm{po}')$, and the kernel of $\mathrm{Cs}[G]$ is bounded:
\[
\left|\mathrm{Cs}[G](\xi',\xi)\right|\le C_{\mathrm{po},\mathrm{po}'}\cdot\mathrm{Cq}\cdot\boldsymbol{\Theta}_{\varepsilon_G}^{\#|\mathrm{Sb}^2(\mathrm{po},\mathrm{po}')|},
\]
where $C_{\mathrm{po},\mathrm{po}'}$ depends only on $(\mu,n_\nu,n_{\mathfrak{l}},\mathrm{po},\mathrm{po}')$ and on the parameters.
\item[(SO1)] $\mathrm{Cs}[G']\,\mathrm{Cs}[G]=\mathrm{Cs}[GG']$ for $G\in\mathrm{Ig}[\mathrm{po},\mathrm{po}']$ and $G'\in\mathrm{Ig}[\mathrm{po}',\mathrm{po}'']$.
\item[(SO2a)] $V_v\,\mathrm{Cs}[G]\,V_v^{-1}=\mathrm{Cs}[\mathrm{Tg}_vG]$ for $v\in\mathbb{Z}e_\mu$.
\item[(SO2b)] $\mathrm{Cs}[\Theta_{\mathrm{po}'}G]=V_{\tau}\,\mathrm{Cs}[G]^*$ for $G\in\mathrm{Ig}[\mathrm{po},\mathrm{po}']$, where $\tau=2(\mathrm{po}'-\mathrm{po})e_\mu$.
\item[(SO3)] Let $n_\mu\ge4$ be an integer and $\mathrm{gv}:=n_\mu e_\mu$, let $\mathrm{po}_0<\mathrm{po}_1<\dots<\mathrm{po}_s=\mathrm{po}_0+n_\mu$ be elements of $\tfrac12\mathbb{Z}$ with $s\ge2$, and let $G_k\in\mathrm{Ig}[\mathrm{po}_{k-1},\mathrm{po}_k]$.
For a configuration $(A,\mathfrak{z})$ on $C(n)$ let $\omega_k$ be the restriction of its periodic extension to the closed slab $k$, that is, to $[\mathrm{po}_{k-1},\mathrm{po}_k]$.
Then the map $(A,\mathfrak{z})\mapsto G_1(\omega_1)\cdots G_s(\omega_s)$, its values being read in $\mathrm{Gr}_n$ as in (F5), is a Borel map with values in $\mathrm{Gr}_n^{\mathrm{ev}}$;
it is invariant under $(A,\mathfrak{z})\mapsto(A+r_qm,\mathfrak{z}-dm)$ for every integer one-cochain $m$ on $C(n)$;
its coefficients are bounded by $C\cdot\exp\left(\left(\frac12-\min_k\varepsilon_{G_k}\right)\left\|F^{\mathrm{u}}\right\|^2\right)$ with a constant $C$;
and
\begin{multline*}
\sum_{\mathfrak{z}:\,d\mathfrak{z}=0}\ \int_{[0,r_q)^{\mathrm{links}}}\mathrm{d}A\,\exp\left(-\frac12\left\|F^{\mathrm{u}}\right\|^2\right)\int_n\exp\left(-S_n(A)\right)\cdot G_1(\omega_1)\cdots G_s(\omega_s)\\
=\mathrm{Tr}_{\mathrm{Hs}(\mathrm{po}_0)}\left(V_{\mathrm{gv}}^{-1}\,\mathrm{Cs}[G_s]\cdots\mathrm{Cs}[G_1]\right).
\end{multline*}
\end{itemize}
\end{named}

\begin{named}{Remark}
Theorem~A is \cite[Theorem~5.29]{QEDII} in the setting of Section~A.0.
\end{named}

\begin{proof}[Proof of (SO0)]
\emph{The operators $V_v$.}
The map
\[
\Xi_{\mathrm{po}+v_\mu}\to\Xi_{\mathrm{po}},\qquad(\mathrm{mi}',\boldsymbol{\zeta}',I+v)\mapsto(\tau_v^*\mathrm{mi}',\tau_v^*\boldsymbol{\zeta}',I),
\]
is a bijection.
It carries $\mathrm{mq}_{\mathrm{po}+v_\mu}$ to $\mathrm{mq}_{\mathrm{po}}$, by Lemma~A.5\,(iii) and because $F_p(\tau_v^*\mathrm{mi}')=F_{p+v}(\mathrm{mi}')$;
translations preserve the orientation of links, so that the box $[0,r_q)^{\mathrm{Mi}^1}$ is preserved.
Hence $V_v$ is unitary.
For the group law let $f\in\mathrm{Hs}(\mathrm{po})$.
Then
\[
(V_vV_{v'}f)(\mathrm{mi}'',\boldsymbol{\zeta}'',I+v'+v)=(V_{v'}f)(\tau_v^*\mathrm{mi}'',\tau_v^*\boldsymbol{\zeta}'',I+v')=f(\tau_{v'}^*\tau_v^*\mathrm{mi}'',\tau_{v'}^*\tau_v^*\boldsymbol{\zeta}'',I),
\]
and $(\tau_{v'}^*\tau_v^*u)(\boldsymbol{\sigma})=u(\boldsymbol{\sigma}+v'+v)=(\tau_{v+v'}^*u)(\boldsymbol{\sigma})$.

\emph{The bound.}
One has $|\mathrm{ef}|=1$.
The coefficients of $\exp(-S_{(\mathrm{po},\mathrm{po}')}(A))$ are bounded uniformly in $A$ (Bd2), and the elements $C_I$ and $\Theta_{\mathrm{po}'}(C_{I'})$ have fixed coefficients.
By (Ig2) it follows that
\[
\left|\mathrm{Ms}_{\mathrm{po},\mathrm{po}'}\left[G(\omega)\cdot\exp\left(-S_{(\mathrm{po},\mathrm{po}')}\right)\right]_{I',I}\right|\le C'\cdot\mathrm{Cq}\cdot\exp\left(\left(\frac12-\varepsilon_G\right)\sum_{p\in\mathrm{Sb}^2(\mathrm{po},\mathrm{po}')}F_p^2\right).
\]
Hence, the constraints being dropped, the expression in braces is bounded by
\[
C'\cdot\mathrm{Cq}\cdot\prod_{p\in\mathrm{Sb}^2(\mathrm{po},\mathrm{po}')}\exp\left(-\varepsilon_GF_p^2\right),
\]
and by (K2) the image of this bound under $\mathrm{Ii}_{\mathrm{po},\mathrm{po}'}$ is at most
\[
C'\cdot\mathrm{Cq}\cdot r_q^{\#|\mathrm{Sb}^1(\mathrm{po},\mathrm{po}')|}\cdot\boldsymbol{\Theta}_{\varepsilon_G}^{\#|\mathrm{Sb}^2(\mathrm{po},\mathrm{po}')|}
\]
for all values of the mirror variables;
here $F_p=(dA)(p)+r_q\mathfrak{z}_p$, with $\mathfrak{z}_p$ a variable of summation.
Moreover $c_I\le6^{64v_s}$.
The kernel is a Borel function of $(\xi',\xi)$, by Tonelli's theorem.
A bounded kernel on a product of finite measure spaces is square integrable, so that by (K3) it defines a Hilbert--Schmidt operator.
The kernel is linear in $G$, because $\mathrm{Ms}_{\mathrm{po},\mathrm{po}'}[\,\cdot\,]$ and $\mathrm{Ii}_{\mathrm{po},\mathrm{po}'}$ are linear.
\end{proof}

\begin{proof}[Proof of (SO1)]
All integrals below converge absolutely by the bounds just proved:
the measures are finite, the kernels are bounded, and inside each kernel the interior integral $\mathrm{Ii}$ converges absolutely.
Fubini's theorem is used freely.
By (K3) the kernel of $\mathrm{Cs}[G']\,\mathrm{Cs}[G]$ at $(\xi'',\xi)$ is
\[
\int\mathrm{d}\,\mathrm{mq}_{\mathrm{po}'}(\xi')\,\mathrm{Cs}[G'](\xi'',\xi')\cdot\mathrm{Cs}[G](\xi',\xi).
\]
Write $\omega_1=\mathrm{mi}\vee\mathrm{omi}_1\vee\mathrm{mi}'$ and $\omega_2=\mathrm{mi}'\vee\mathrm{omi}_2\vee\mathrm{mi}''$.
Apart from the density $\mathrm{dn}_{\mathrm{po}'}(\boldsymbol{\zeta}')\cdot\mathrm{Ga}_{\mathrm{po}'}(\boldsymbol{\zeta}',I')$ of $\mathrm{mq}_{\mathrm{po}'}$, the integrand depends on $\boldsymbol{\zeta}'$ only through
\[
\mathrm{ef}_{\mathrm{po}'}(\boldsymbol{\zeta}';\omega_2)\cdot\overline{\mathrm{ef}_{\mathrm{po}'}\left(\boldsymbol{\zeta}';\theta_{\mathrm{po}'}^*\omega_1\right)},
\]
and on $I'$ only through $c_{I'}\cdot\mathrm{Ms}_{\mathrm{po}',\mathrm{po}''}[\ldots]_{I'',I'}\cdot\mathrm{Ms}_{\mathrm{po},\mathrm{po}'}[\ldots]_{I',I}$.
For fixed $I'$, Lemma~A.4 at $\mathrm{po}'$, with the cells $\boldsymbol{\sigma}^+$ read from $\omega_2$ and the cells $\boldsymbol{\sigma}^-$ read from $\omega_1$, gives
\begin{multline*}
\int\mathrm{dn}_{\mathrm{po}'}(\boldsymbol{\zeta}')\,\mathrm{Ga}_{\mathrm{po}'}(\boldsymbol{\zeta}',I')\cdot\mathrm{ef}_{\mathrm{po}'}(\boldsymbol{\zeta}';\omega_2)\cdot\overline{\mathrm{ef}_{\mathrm{po}'}\left(\boldsymbol{\zeta}';\theta_{\mathrm{po}'}^*\omega_1\right)}\\
=\sum_{\mathrm{mc}}\int\mathrm{d}\mathrm{Ac}\,\prod_{p\in\mathrm{Cu}^2(\mathrm{po}')}\exp\left(-\frac12F_p^2\right)\cdot\prod_{\mathrm{qb}\in\mathrm{Cu}^3(\mathrm{po}')}\mathbf{1}\left[(d\mathfrak{z})(\mathrm{qb})=0\right]\cdot u_{I'}(\mathrm{Ac}).
\end{multline*}
At a site plane $\mathrm{po}'$ both sides are equal to $1$;
in that case the measure $\mathrm{mq}_{\mathrm{po}'}$ contains instead the integral over the links and the sum over the plaquettes of $\mathrm{Mi}(\mathrm{po}')$, with the weight $\prod_{p\in\mathrm{Mi}^2(\mathrm{po}')}\exp\left(-\frac12F_p^2\right)$.
By (Ap9), (F2) and (Bd3), for fixed values of all photon variables,
\begin{align*}
&\sum_{I'}\mathrm{Ms}_{\mathrm{po}',\mathrm{po}''}\left[G'\cdot\exp\left(-S_{(\mathrm{po}',\mathrm{po}'')}\right)\right]_{I'',I'}\cdot c_{I'}\cdot u_{I'}(\mathrm{Ac})\cdot\mathrm{Ms}_{\mathrm{po},\mathrm{po}'}\left[G\cdot\exp\left(-S_{(\mathrm{po},\mathrm{po}')}\right)\right]_{I',I}\\
&\qquad=\mathrm{Ms}_{\mathrm{po},\mathrm{po}''}\left[G\cdot\exp\left(-S_{(\mathrm{po},\mathrm{po}')}\right)\cdot\exp\left(-\mathrm{Xc}_{\mathrm{po}'}\right)\cdot G'\cdot\exp\left(-S_{(\mathrm{po}',\mathrm{po}'')}\right)\right]_{I'',I}\\
&\qquad=\mathrm{Ms}_{\mathrm{po},\mathrm{po}''}\left[GG'\cdot\exp\left(-S_{(\mathrm{po},\mathrm{po}'')}\right)\right]_{I'',I},
\end{align*}
since all factors are even and commute.
By (Ap6) the variables $(\mathrm{omi}_1,\mathrm{mi}',\mathrm{Ac},\mathrm{mc},\mathrm{omi}_2)$ are the interior variables of the slab $(\mathrm{po},\mathrm{po}'')$, each link variable ranging over $[0,r_q)$ and each $\mathfrak{z}_p$ over $\mathbb{Z}$;
by the same partition,
\begin{multline*}
\mathrm{wb}_{\mathrm{po},\mathrm{po}'}(\omega_1)\cdot\prod_{p\in\mathrm{Mi}^2(\mathrm{po}')}\exp\left(-\frac12F_p^2\right)\cdot\prod_{p\in\mathrm{Cu}^2(\mathrm{po}')}\exp\left(-\frac12F_p^2\right)\cdot\prod_{\mathrm{qb}\in\mathrm{Cu}^3(\mathrm{po}')}\mathbf{1}\left[(d\mathfrak{z})(\mathrm{qb})=0\right]\\
\cdot\mathrm{wb}_{\mathrm{po}',\mathrm{po}''}(\omega_2)=\mathrm{wb}_{\mathrm{po},\mathrm{po}''}(\omega).
\end{multline*}
Moreover $\mathrm{ef}_{\mathrm{po}}(\boldsymbol{\zeta};\omega_1)=\mathrm{ef}_{\mathrm{po}}(\boldsymbol{\zeta};\omega)$ and $\mathrm{ef}_{\mathrm{po}''}(\boldsymbol{\zeta}'';\theta_{\mathrm{po}''}^*\omega_2)=\mathrm{ef}_{\mathrm{po}''}(\boldsymbol{\zeta}'';\theta_{\mathrm{po}''}^*\omega)$, since the same cells are read.
Finally
\[
\left(c_{I''}\cdot c_{I'}\right)^{\frac12}\cdot\left(c_{I'}\cdot c_I\right)^{\frac12}=\left(c_{I''}\cdot c_I\right)^{\frac12}\cdot c_{I'},
\]
because $c_{I'}\ge0$.
The result is $\mathrm{Cs}[GG'](\xi'',\xi)$.
\end{proof}

\begin{proof}[Proof of (SO2a)]
Let $\xi_1=(\mathrm{mi},\boldsymbol{\zeta},I)$ and $\xi_2=(\mathrm{mi}',\boldsymbol{\zeta}',I')$, and put
\[
\xi_1^v:=\left(\tau_{-v}^*\mathrm{mi},\tau_{-v}^*\boldsymbol{\zeta},I+v\right)\in\Xi_{\mathrm{po}+v_\mu},\qquad\xi_2^v:=\left(\tau_{-v}^*\mathrm{mi}',\tau_{-v}^*\boldsymbol{\zeta}',I'+v\right)\in\Xi_{\mathrm{po}'+v_\mu}.
\]
\emph{Claim:} $\mathrm{Cs}[\mathrm{Tg}_vG](\xi_2^v,\xi_1^v)=\mathrm{Cs}[G](\xi_2,\xi_1)$.
In $\mathrm{Ii}_{\mathrm{po}+v_\mu,\mathrm{po}'+v_\mu}$ substitute $\omega^v=\tau_{-v}^*\omega$;
this is a bijection of the interior variables that preserves the boxes and the counting measures.
The weight $\mathrm{wb}$ is translation invariant.
By Lemma~A.5\,(iii), $\mathrm{ef}_{\mathrm{po}+v_\mu}(\tau_{-v}^*\boldsymbol{\zeta};\tau_{-v}^*\omega)=\mathrm{ef}_{\mathrm{po}}(\boldsymbol{\zeta};\omega)$.
By (M4),
\[
\theta_{\mathrm{po}'+v_\mu}^*\tau_{-v}^*=\left(\tau_{-v}\theta_{\mathrm{po}'+v_\mu}\right)^*=\left(\theta_{\mathrm{po}'}\tau_{-v}\right)^*=\tau_{-v}^*\theta_{\mathrm{po}'}^*,
\]
so that, again by Lemma~A.5\,(iii),
\[
\mathrm{ef}_{\mathrm{po}'+v_\mu}\left(\tau_{-v}^*\boldsymbol{\zeta}';\theta_{\mathrm{po}'+v_\mu}^*\tau_{-v}^*\omega\right)=\mathrm{ef}_{\mathrm{po}'}\left(\boldsymbol{\zeta}';\theta_{\mathrm{po}'}^*\omega\right).
\]
Further $(\mathrm{Tg}_vG)(\omega^v)=\mathrm{Tg}_v(G(\omega))$, and
\[
S_{(\mathrm{po}+v_\mu,\mathrm{po}'+v_\mu)}(\tau_{-v}^*A)=\mathrm{Tg}_vS_{(\mathrm{po},\mathrm{po}')}(A),
\]
because $M_i(x)$ depends on the label only and $U_{\tau_{-v}^*A}(x+v,x^{(i)}+v)=U_A(x,x^{(i)})$.
Together with (F4) and $c_{I+v}=c_I$ this proves the claim.
By the definition of $V_v$ and the invariance of the measures, the claim is the identity $V_v\,\mathrm{Cs}[G]=\mathrm{Cs}[\mathrm{Tg}_vG]\,V_v$.
\end{proof}

\begin{proof}[Proof of (SO2b)]
Put $\mathrm{po}'':=2\mathrm{po}'-\mathrm{po}$ and $\tilde G:=\Theta_{\mathrm{po}'}G$.
Since $\mathrm{po}''-\mathrm{po}=2(\mathrm{po}'-\mathrm{po})$ is an integer, $\mathrm{po}''$ is of the same kind as $\mathrm{po}$, whatever the kind of $\mathrm{po}'$ is.
For $\xi=(\mathrm{mi},\boldsymbol{\zeta},I)\in\Xi_{\mathrm{po}}$ let
\[
\xi^{\tau}:=\left(\tau_{-\tau}^*\mathrm{mi},\tau_{-\tau}^*\boldsymbol{\zeta},I+\tau\right)\in\Xi_{\mathrm{po}''};
\]
then $(V_{\tau}f)(\xi^{\tau})=f(\xi)$.
\emph{Claim:}
\begin{equation}\tag{Ap10}
\mathrm{Cs}[\tilde G](\xi^{\tau},\xi')=\overline{\mathrm{Cs}[G](\xi',\xi)}\qquad\text{for all }\xi\in\Xi_{\mathrm{po}}\text{ and }\xi'\in\Xi_{\mathrm{po}'}.
\end{equation}
Granted (Ap10), (K3) gives
\[
\left(V_{\tau}\,\mathrm{Cs}[G]^*f\right)(\xi^{\tau})=\left(\mathrm{Cs}[G]^*f\right)(\xi)=\int\mathrm{d}\,\mathrm{mq}_{\mathrm{po}'}(\xi')\,\overline{\mathrm{Cs}[G](\xi',\xi)}\cdot f(\xi')=\left(\mathrm{Cs}[\tilde G]f\right)(\xi^{\tau}).
\]

\emph{Proof of (Ap10).}
For a configuration $\omega=(A,\mathfrak{z})$ of a closed slab let $A(\omega):=A$ be its one-cochain (the part on the links) and $\mathfrak{z}(\omega):=\mathfrak{z}$ its integer two-cochain (the part on the plaquettes);
$d\mathfrak{z}(\omega)$ denotes the coboundary of $\mathfrak{z}(\omega)$ on the cubes of the slab.
For configurations $\omega$ of $[\mathrm{po},\mathrm{po}']$ and $\tilde\omega$ of $[\mathrm{po}',\mathrm{po}'']$ put
\[
S_{(\mathrm{po},\mathrm{po}')}(\omega):=S_{(\mathrm{po},\mathrm{po}')}\left(A(\omega)\right),\qquad S_{(\mathrm{po}',\mathrm{po}'')}(\tilde\omega):=S_{(\mathrm{po}',\mathrm{po}'')}\left(A(\tilde\omega)\right).
\]
Let $\tilde\omega=\mathrm{mi}'\vee\widetilde{\mathrm{omi}}\vee\tau_{-\tau}^*\mathrm{mi}$ be a configuration of the closed slab $[\mathrm{po}',\mathrm{po}'']$, and let $\omega:=\theta_{\mathrm{po}'}^*\tilde\omega$, a configuration of $[\mathrm{po},\mathrm{po}']$;
that is, $A(\omega)=\theta_{\mathrm{po}'}^*A(\tilde\omega)$ and $\mathfrak{z}(\omega)=\theta_{\mathrm{po}'}^*\mathfrak{z}(\tilde\omega)$.
By (Ap2) the cells of $\mathrm{Mi}(\mathrm{po}')$ are fixed by $\theta_{\mathrm{po}'}$, with $\epsilon(\theta_{\mathrm{po}'},\boldsymbol{\sigma})=1$ (there are none if $\mathrm{po}'$ is a link plane), and $\theta_{\mathrm{po}'}\boldsymbol{\sigma}=\boldsymbol{\sigma}+\tau$, with $\epsilon(\theta_{\mathrm{po}'},\boldsymbol{\sigma})=1$, for $\boldsymbol{\sigma}\in\mathrm{Mi}(\mathrm{po})$ (there are none if $\mathrm{po}$ is a link plane).
Hence $\omega=\mathrm{mi}\vee\mathrm{omi}\vee\mathrm{mi}'$ with $\mathrm{omi}=\theta_{\mathrm{po}'}^*\widetilde{\mathrm{omi}}$, and $\widetilde{\mathrm{omi}}\mapsto\mathrm{omi}$ is the bijection
\[
\mathbb{R}^{\mathrm{Sb}^1(\mathrm{po}',\mathrm{po}'')}\times\mathbb{Z}^{\mathrm{Sb}^2(\mathrm{po}',\mathrm{po}'')}\to\mathbb{R}^{\mathrm{Sb}^1(\mathrm{po},\mathrm{po}')}\times\mathbb{Z}^{\mathrm{Sb}^2(\mathrm{po},\mathrm{po}')},\qquad\mathrm{omi}(\boldsymbol{\sigma})=\epsilon(\theta_{\mathrm{po}'},\boldsymbol{\sigma})\cdot\widetilde{\mathrm{omi}}(\theta_{\mathrm{po}'}\boldsymbol{\sigma}).
\]
It maps the box $[0,r_q)^{\mathrm{Sb}^1(\mathrm{po}',\mathrm{po}'')}$ onto $\prod_bI_b$, where $I_b=[0,r_q)$ for the links orthogonal to $e_\mu$ and $I_b=(-r_q,0]$ for the links parallel to $e_\mu$.
\begin{itemize}
\item[(1)] One has $F(\omega)=\theta_{\mathrm{po}'}^*F(\tilde\omega)$, $d\mathfrak{z}(\omega)=\theta_{\mathrm{po}'}^*d\mathfrak{z}(\tilde\omega)$ and $\epsilon(\theta_{\mathrm{po}'},\boldsymbol{\sigma})^2=1$;
hence $\mathrm{wb}_{\mathrm{po}',\mathrm{po}''}(\tilde\omega)=\mathrm{wb}_{\mathrm{po},\mathrm{po}'}(\omega)$.
\item[(2)] One has $\theta_{\mathrm{po}'}^*\omega=\tilde\omega$;
hence $\mathrm{ef}_{\mathrm{po}'}(\boldsymbol{\zeta}';\tilde\omega)=\overline{\overline{\mathrm{ef}_{\mathrm{po}'}\left(\boldsymbol{\zeta}';\theta_{\mathrm{po}'}^*\omega\right)}}$.
\item[(3)] By Lemma~A.5\,(iii) with $v=\tau$,
\[
\mathrm{ef}_{\mathrm{po}''}\left(\tau_{-\tau}^*\boldsymbol{\zeta};\theta_{\mathrm{po}''}^*\tilde\omega\right)=\mathrm{ef}_{\mathrm{po}}\left(\boldsymbol{\zeta};\tau_{\tau}^*\theta_{\mathrm{po}''}^*\tilde\omega\right),
\]
and $\tau_{\tau}^*\theta_{\mathrm{po}''}^*=(\theta_{\mathrm{po}''}\tau_{\tau})^*=\theta_{\mathrm{po}'}^*$, because $\tau_{\tau}=\theta_{\mathrm{po}''}\theta_{\mathrm{po}'}$ by (M3).
Hence $\overline{\mathrm{ef}_{\mathrm{po}''}\left(\tau_{-\tau}^*\boldsymbol{\zeta};\theta_{\mathrm{po}''}^*\tilde\omega\right)}=\overline{\mathrm{ef}_{\mathrm{po}}(\boldsymbol{\zeta};\omega)}$.
\item[(4)] One has $\tilde G(\tilde\omega)=\Theta_{\mathrm{po}'}(G(\omega))$ (Section~A.5).
Term by term, $\Theta_{\mathrm{po}'}\left(S_{(\mathrm{po},\mathrm{po}')}(A(\omega))\right)=S_{(\mathrm{po}',\mathrm{po}'')}(A(\tilde\omega))$.
Indeed, by (S3) the term $-\bar\psi(x)M_i(x)U_{A(\omega)}(x,x^{(i)})\psi(x^{(i)})$ is mapped to
\[
-\bar\psi(\theta_{\mathrm{po}'}x^{(i)})\,\gamma_\mu M_i(x)^*\gamma_\mu\,\overline{U_{A(\omega)}(x,x^{(i)})}\,\psi(\theta_{\mathrm{po}'}x).
\]
Here $M_i(x)^*=M_i(x)$ by (I.m5), and
\[
\gamma_\mu M_i(x)\gamma_\mu=\frac{w_i}{2}\left(\rho-\gamma_\mu\,\gamma\cdot\hat{d}_i(x)\,\gamma_\mu\right)=\frac{w_i}{2}\left(\rho+\gamma\cdot\left(r_\mu\hat{d}_i(x)\right)\right)=M_i(\theta_{\mathrm{po}'}x^{(i)}),
\]
because $\gamma_\mu\,\gamma\cdot v\,\gamma_\mu=-\gamma\cdot(r_\mu v)$, $(\theta_{\mathrm{po}'}x^{(i)})^{(i)}=\theta_{\mathrm{po}'}x$ and $\hat{d}_i(\theta_{\mathrm{po}'}x^{(i)})=r_\mu\hat{d}_i(x^{(i)})=-r_\mu\hat{d}_i(x)$ (Section~I.4.7, (I.s2)).
Moreover
\[
\overline{U_{\theta_{\mathrm{po}'}^*A(\tilde\omega)}(x,x^{(i)})}=\exp\left(-iq\,A(\tilde\omega)\left(\theta_{\mathrm{po}'}\pi(x),\theta_{\mathrm{po}'}\pi(x^{(i)})\right)\right)=U_{A(\tilde\omega)}(\theta_{\mathrm{po}'}x^{(i)},\theta_{\mathrm{po}'}x).
\]
So the image is the term of the oriented bond $(\theta_{\mathrm{po}'}x^{(i)},\theta_{\mathrm{po}'}x)$.
The term $M_{\mathrm{ch}}\bar\psi(x)\psi(x)$ is mapped to $M_{\mathrm{ch}}\bar\psi(\theta_{\mathrm{po}'}x)\gamma_\mu\gamma_\mu\psi(\theta_{\mathrm{po}'}x)$, since $M_{\mathrm{ch}}$ is real.
The map $(x,x^{(i)})\mapsto(\theta_{\mathrm{po}'}x^{(i)},\theta_{\mathrm{po}'}x)$ is a bijection of the oriented bonds with both ends in $X((\mathrm{po},\mathrm{po}'))$ onto those with both ends in $X((\mathrm{po}',\mathrm{po}''))$.
Hence, the factors being even,
\[
\tilde G(\tilde\omega)\cdot\exp\left(-S_{(\mathrm{po}',\mathrm{po}'')}(\tilde\omega)\right)=\Theta_{\mathrm{po}'}\left(G(\omega)\cdot\exp\left(-S_{(\mathrm{po},\mathrm{po}')}(\omega)\right)\right),
\]
and (F3) gives
\[
\mathrm{Ms}_{\mathrm{po}',\mathrm{po}''}\left[\tilde G(\tilde\omega)\cdot\exp(-S)\right]_{I+\tau,I'}=\overline{\mathrm{Ms}_{\mathrm{po},\mathrm{po}'}\left[G(\omega)\cdot\exp(-S)\right]_{I',I}}.
\]
\item[(5)] One has $c_{I+\tau}=c_I$.
\end{itemize}
Let $\mathrm{Bb}(\omega)$ denote the expression in braces of the kernel $\mathrm{Cs}[G](\xi',\xi)$ at $\omega$.
By (1)--(5), the expression in braces of the kernel $\mathrm{Cs}[\tilde G](\xi^{\tau},\xi')$ at $\tilde\omega$ is the complex conjugate of $\mathrm{Bb}(\omega)$.
Therefore
\[
\mathrm{Cs}[\tilde G](\xi^{\tau},\xi')=\left(c_I\cdot c_{I'}\right)^{\frac12}\cdot\overline{\sum_{\mathfrak{z}}\int_{\prod_bI_b}\mathrm{d}A\,\mathrm{Bb}(\omega)}.
\]
The function $\mathrm{Bb}$ is invariant under $\omega\mapsto\omega+[m]$ for $m\in\mathbb{Z}^{\mathrm{Sb}^1(\mathrm{po},\mathrm{po}')}$:
the weight $\mathrm{wb}$ depends on $F$ and $d\mathfrak{z}$ only;
the two factors $\mathrm{ef}$ are invariant by Lemma~A.5\,(ii), since a link of $\mathrm{Sb}^1(\mathrm{po},\mathrm{po}')$ is contained in the closed slab and in neither plane, so that it is not cut by $\mathrm{po}$ or by $\mathrm{po}'$ and is contained neither in $\{y_\mu\le\mathrm{po}\}$ nor in $\{y_\mu\ge\mathrm{po}'\}$;
the factor $G$ is invariant by (Ig1);
and $S_{(\mathrm{po},\mathrm{po}')}$ reads phases only.
Moreover $\mathrm{Ii}_{\mathrm{po},\mathrm{po}'}|\mathrm{Bb}|<\infty$ by (SO0).
By Lemma~A.1 the intervals $I_b$ may therefore be replaced by $[0,r_q)$.
This is (Ap10).
\end{proof}

\begin{proof}[Proof of (SO3)]
\emph{Step 0 (the left side).}
A $\Gamma(n)$-periodic configuration is $\mathrm{Gmp}$-periodic;
$\omega_k$ is its restriction to $\mathrm{Ln}[\mathrm{po}_{k-1},\mathrm{po}_k]$ and $\mathrm{Pq}[\mathrm{po}_{k-1},\mathrm{po}_k]$, and it lies in $\mathrm{Cfg}[\mathrm{po}_{k-1},\mathrm{po}_k]$.
Extend the sequence by $\mathrm{po}_{k+js}:=\mathrm{po}_k+jn_\mu$.
The translation $\tau_{\mathrm{gv}}$ maps the classes $\mathrm{Mi}(\mathrm{po}_k)$, $\mathrm{Cu}(\mathrm{po}_k)$, $\mathrm{Sb}(\mathrm{po}_{k-1},\mathrm{po}_k)$ of (Ap6) onto those with the index $k+s$, and the cells of $C(n)$ are the cells of $\mathrm{Cy}$ modulo $\tau_{\mathrm{gv}}$.
Hence:
\begin{equation}\tag{Ap11}
\begin{split}
&\text{the sets }\mathrm{Mi}(\mathrm{po}_k),\ \mathrm{Cu}(\mathrm{po}_k)\ \ (0\le k\le s-1)\ \text{ and }\ \mathrm{Sb}(\mathrm{po}_{k-1},\mathrm{po}_k)\ \ (1\le k\le s)\\
&\text{together form a set of representatives of the cells of }C(n).
\end{split}
\end{equation}
In particular the sets $\mathrm{Sb}^2(\mathrm{po}_{k-1},\mathrm{po}_k)$ are pairwise disjoint sets of plaquettes of $C(n)$;
so the coefficients of $\prod_kG_k(\omega_k)$ are bounded by $C\cdot\exp\left(\left(\frac12-\min_k\varepsilon_{G_k}\right)\left\|F^{\mathrm{u}}\right\|^2\right)$.
The invariance under the integer shifts follows from (Ig1).
The left side converges absolutely, by the last sentence of Lemma~2.4\,(d) with $\varepsilon=\min_k\varepsilon_{G_k}$, which lies in $(0,\frac12]$ by (Ig2).

\emph{Step 1 (reduction to $s=2$).}
By (SO1), $\mathrm{Cs}[G_{s-1}]\cdots\mathrm{Cs}[G_1]=\mathrm{Cs}[G_1\cdots G_{s-1}]$, where $G_1\cdots G_{s-1}\in\mathrm{Ig}[\mathrm{po}_0,\mathrm{po}_{s-1}]$ and $(G_1\cdots G_{s-1})(\omega)=\prod_{k<s}G_k(\omega_k)$.
It therefore suffices to treat the case $s=2$.

\emph{Step 2 (the trace).}
The operator $\mathrm{Yy}:=V_{\mathrm{gv}}^{-1}\,\mathrm{Cs}[G_2]\colon\mathrm{Hs}(\mathrm{po}_1)\to\mathrm{Hs}(\mathrm{po}_0)$ is a Hilbert--Schmidt operator with the kernel
\[
\mathrm{Yy}(\xi_0,\xi_1)=\mathrm{Cs}[G_2](\xi_0^{\mathrm{gv}},\xi_1),\qquad\xi_0^{\mathrm{gv}}:=\left(\tau_{-\mathrm{gv}}^*\mathrm{mi}_0,\tau_{-\mathrm{gv}}^*\boldsymbol{\zeta}_0,I_0+\mathrm{gv}\right)\in\Xi_{\mathrm{po}_2}.
\]
By (K3),
\[
\mathrm{Tr}\left(\mathrm{Yy}\,\mathrm{Cs}[G_1]\right)=\int\mathrm{d}\,\mathrm{mq}_{\mathrm{po}_0}(\xi_0)\int\mathrm{d}\,\mathrm{mq}_{\mathrm{po}_1}(\xi_1)\,\mathrm{Cs}[G_2](\xi_0^{\mathrm{gv}},\xi_1)\cdot\mathrm{Cs}[G_1](\xi_1,\xi_0),
\]
with absolute convergence, also after the kernels are written out (by the bounds of (SO0)).
Fubini's theorem is used freely.

\emph{Step 3 (photon).}
Here $\omega_1=\mathrm{mi}_0\vee\mathrm{omi}_1\vee\mathrm{mi}_1$ and $\omega_2=\mathrm{mi}_1\vee\mathrm{omi}_2\vee\tau_{-\mathrm{gv}}^*\mathrm{mi}_0$.
The integrand contains the factors
\begin{align*}
&\text{at }\mathrm{po}_1\colon&&\mathrm{Ga}_{\mathrm{po}_1}(\boldsymbol{\zeta}_1,I_1)\cdot\mathrm{ef}_{\mathrm{po}_1}(\boldsymbol{\zeta}_1;\omega_2)\cdot\overline{\mathrm{ef}_{\mathrm{po}_1}\left(\boldsymbol{\zeta}_1;\theta_{\mathrm{po}_1}^*\omega_1\right)},\\
&\text{at }\mathrm{po}_0\colon&&\mathrm{Ga}_{\mathrm{po}_0}(\boldsymbol{\zeta}_0,I_0)\cdot\mathrm{ef}_{\mathrm{po}_0}(\boldsymbol{\zeta}_0;\omega_1)\cdot\overline{\mathrm{ef}_{\mathrm{po}_2}\left(\tau_{-\mathrm{gv}}^*\boldsymbol{\zeta}_0;\theta_{\mathrm{po}_2}^*\omega_2\right)}.
\end{align*}
By Lemma~A.5\,(iii) with $v=\mathrm{gv}$,
\[
\mathrm{ef}_{\mathrm{po}_2}\left(\tau_{-\mathrm{gv}}^*\boldsymbol{\zeta}_0;\theta_{\mathrm{po}_2}^*\omega_2\right)=\mathrm{ef}_{\mathrm{po}_0}\left(\boldsymbol{\zeta}_0;\tau_{\mathrm{gv}}^*\theta_{\mathrm{po}_2}^*\omega_2\right),
\]
and by (M4)
\[
\tau_{\mathrm{gv}}^*\theta_{\mathrm{po}_2}^*=\left(\theta_{\mathrm{po}_2}\tau_{\mathrm{gv}}\right)^*=\left(\tau_{\mathrm{gv}}\theta_{\mathrm{po}_0}\right)^*=\theta_{\mathrm{po}_0}^*\tau_{\mathrm{gv}}^*.
\]
The configuration $\tau_{\mathrm{gv}}^*\omega_2$ is the periodic continuation of the configuration to the closed slab $[\mathrm{po}_1-n_\mu,\mathrm{po}_0]$.
So Lemma~A.4 applies at $\mathrm{po}_1$, with the cells $\boldsymbol{\sigma}^+$ read from $\omega_2$ and the cells $\boldsymbol{\sigma}^-$ from $\omega_1$, and at $\mathrm{po}_0$, with the cells $\boldsymbol{\sigma}^+$ read from $\omega_1$ and the cells $\boldsymbol{\sigma}^-$ from $\tau_{\mathrm{gv}}^*\omega_2$, for fixed $I_0$, $I_1$.
It replaces the integrals over $\boldsymbol{\zeta}_1$ and $\boldsymbol{\zeta}_0$ by the sums over $\mathfrak{z}$ on $\mathrm{Cu}^2(\mathrm{po}_1)$ and $\mathrm{Cu}^2(\mathrm{po}_0)$ and the integrals over the links of $\mathrm{Cu}^1(\mathrm{po}_1)$ and $\mathrm{Cu}^1(\mathrm{po}_0)$, with the weights
\[
\prod\exp\left(-\frac12F_p^2\right)\cdot\prod\mathbf{1}\left[(d\mathfrak{z})(\mathrm{qb})=0\right]\cdot u_{I_1}(A)\cdot u_{I_0}(A),
\]
$F$ and $d\mathfrak{z}$ being those of the periodic configuration.
Each of the two planes may be of either kind;
at a site plane this step is void.

\emph{Step 4 (fermions).}
Put
\[
\mathrm{Ms}_k:=\mathrm{Ms}_{\mathrm{po}_{k-1},\mathrm{po}_k}\left[G_k(\omega_k)\cdot\exp\left(-S_{(\mathrm{po}_{k-1},\mathrm{po}_k)}\right)\right],\qquad\mathrm{diag}(u^{(k)}):=\mathrm{diag}\left(u_I(A)\right)_{I\subset\mathrm{Ri}(\mathrm{po}_k)}.
\]
With $c_{I_0+\mathrm{gv}}=c_{I_0}$, the remaining dependence on the index sets is
\begin{multline*}
\sum_{I_0,I_1}(\mathrm{Ms}_2)_{I_0+\mathrm{gv},I_1}\cdot c_{I_1}\cdot u_{I_1}\cdot(\mathrm{Ms}_1)_{I_1,I_0}\cdot c_{I_0}\cdot u_{I_0}\\
=\mathrm{tr}_{\mathrm{Fk}(\mathrm{po}_0)}\left(\mathrm{Ft}_{\mathrm{gv}}^{-1}\,\mathrm{Ms}_2\,\mathrm{cc}_{\mathrm{po}_1}\mathrm{diag}(u^{(1)})\,\mathrm{Ms}_1\,\mathrm{cc}_{\mathrm{po}_0}\mathrm{diag}(u^{(0)})\right).
\end{multline*}
By (F5) with $\mathrm{mt}_k=\mathrm{cc}_{\mathrm{po}_k}\mathrm{diag}(u^{(k)})$, which is diagonal and hence even, and by (Ap9), this equals
\begin{multline*}
\int_n\exp\left(-\mathrm{Xc}_{\mathrm{po}_0}\right)\cdot G_1\cdot\exp\left(-S_{(\mathrm{po}_0,\mathrm{po}_1)}\right)\cdot\exp\left(-\mathrm{Xc}_{\mathrm{po}_1}\right)\cdot G_2\cdot\exp\left(-S_{(\mathrm{po}_1,\mathrm{po}_2)}\right)\\
=\int_n\exp\left(-S_n(A)\right)\cdot G_1(\omega_1)\cdot G_2(\omega_2),
\end{multline*}
since all factors are even and, by (Bd4) of Section~A.7 with $s=2$,
\[
S_n=S_{(\mathrm{po}_0,\mathrm{po}_1)}+S_{(\mathrm{po}_1,\mathrm{po}_2)}+\mathrm{Xc}_{\mathrm{po}_1}+\mathrm{Xc}_{\mathrm{po}_0}\qquad\text{in }\mathrm{Gr}_n.
\]
Here $\mathrm{Xc}_{\mathrm{po}_0}$, an element of $\mathrm{Gr}_{(\mathrm{po}_0^-,\mathrm{po}_0^+)}$, is read in $\mathrm{Gr}_n$ as in (F5);
this reading is a homomorphism of algebras, generators going to generators or to their negatives, so that (Ap9) at $\mathrm{po}_0$ holds in $\mathrm{Gr}_n$.

\emph{Step 5 (collecting).}
The variables are now:
$\mathrm{mi}_0$ and $\mathrm{mi}_1$, with the weights $\prod_{p\in\mathrm{Mi}^2(\mathrm{po}_k)}\exp\left(-\frac12F_p^2\right)$ from $\mathrm{mq}_{\mathrm{po}_k}$;
the cut variables at $\mathrm{po}_0$ and at $\mathrm{po}_1$;
and $\mathrm{omi}_1$ and $\mathrm{omi}_2$, with the weight $\mathrm{wb}_{\mathrm{po}_0,\mathrm{po}_1}(\omega_1)\cdot\mathrm{wb}_{\mathrm{po}_1,\mathrm{po}_2}(\omega_2)$.
By (Ap11) these are all links of $C(n)$, each with the range $[0,r_q)$, and all plaquettes of $C(n)$, each with the range $\mathbb{Z}$;
the product of the weights is
\[
\exp\left(-\frac12\left\|F^{\mathrm{u}}\right\|^2\right)\cdot\prod_{\mathrm{qb}\text{ a cube of }C(n)}\mathbf{1}\left[(d\mathfrak{z})(\mathrm{qb})=0\right].
\]
This is the left side of (SO3).
\end{proof}

\begin{named}{Corollary A.6}[Theorem SO]
The data $\mathrm{kk}=1$, $\mathrm{Hs}(\mathrm{po})$, $V_v$, and $\mathrm{Cs}$ restricted to $\mathcal{A}_{\mu,n}[\mathrm{po},\mathrm{po}']$ through Lemma~A.2, have the properties stated in Theorem~SO.
\end{named}

\begin{proof}
These objects are formed with $(\mu,n_\nu,n_{\mathfrak{l}})$ and the parameters.
Linearity follows from Lemma~A.2\,(i) and (SO0).
The properties (SO1), (SO2a), (SO2b) follow from Theorem~A together with Lemma~A.2\,(ii), (iii), (iv).
For (SO3), let $G_k\in\mathcal{A}_{\mu,n}[\mathrm{po}_{k-1},\mathrm{po}_k]$ be products with the boxes $\mathrm{Bx}_{k,j}$.
Every $v\in\Gamma(n)$ is of the form $v'+a\,\mathrm{gv}$ with $v'\in\mathrm{Gmp}$ and $a\in\mathbb{Z}$, and $\mathrm{Bx}_{k,j}+v$ lies in $\{\mathrm{po}_{k-1}+an_\mu<y_\mu<\mathrm{po}_k+an_\mu\}$.
Two such open slabs are disjoint unless the pairs $(k,a)$ coincide, and for equal pairs the disjointness is the hypothesis that the products embed in $\mathrm{Gmp}$.
So $G_1\cdots G_s$ is a product that is read on $C(n)$.
If $\hat{A}$ is a lift of $(A,\mathfrak{z})$ (Section~1.5), its restriction to $\mathrm{Ksl}[\mathrm{po}_{k-1},\mathrm{po}_k]$ is a lift of $\omega_k$;
on $\mathrm{Gr}(X((\mathrm{po}_0,\mathrm{po}_0+n_\mu)))$ the map $\mathrm{qu}_n$ is $\mathrm{qu}$ followed by the identification $\mathrm{Gr}_{(\mathrm{po}_0,\mathrm{po}_s)}=\mathrm{Gr}_n$ of (F5).
Hence
\[
(G_1\cdots G_s)(A,\mathfrak{z})=\mathrm{qu}_n\left(\prod_kG_k(\hat{A})\right)=\prod_kG_k(\omega_k),
\]
and the left side of (SO3) in Theorem~A is $Z_n[G_1\cdots G_s]$ of (1.3).
For sums of products both sides are multilinear in $(G_1,\dots,G_s)$.
\end{proof}

\subsubsection*{A.11. The alternating chain of Section 3}

Let $2\,\mathrm{hw}$ be odd and $\mathrm{po}_j=o_\mu+j\,\mathrm{hw}$.
Then the spaces $\mathrm{Hs}(\mathrm{po}_j)$ are alternately spaces of a site plane and of a link plane (Section~A.9).
Since Theorem~A and Corollary~A.6 hold for all positions in $\tfrac12\mathbb{Z}$, whatever their kinds, one has in particular:
(SO2b) for each pair $(\mathrm{po}_j,\mathrm{po}_{j+1})$, consisting of one site plane and one link plane, in either order;
(SO2a) for $v=s\,\tau_\mu$;
and (SO3) for the chain $\mathrm{po}_j<\mathrm{po}_{j+1}<\dots<\mathrm{po}_{j+2k}$.

\begin{named}{Remark}
The kind of a plane bears on the constructions of Sections~A.1--A.10 through the following three facts.
\begin{itemize}
\item[(i)] For every $\mathrm{po}$ one of the sets $\mathrm{Mi}(\mathrm{po})$, $\mathrm{Cu}(\mathrm{po})$ is empty (Section~A.1); which of the two is empty decides which variables $\xi$ contains and which factors of $\mathrm{mq}_{\mathrm{po}}$ and of the kernel are trivial (Sections~A.8, A.9).
\item[(ii)] The crossing bonds (Bd1) are of type~$0$ or of type~$3$; the coefficients are $c_i=3(1\pm\rho)$ in both cases, and $u_i=1$ or $u_i=\exp(\pm iqA_b)$.
\item[(iii)] For $\mathrm{po}<\mathrm{po}'$ one has $\tau=2(\mathrm{po}'-\mathrm{po})e_\mu\in\mathbb{Z}e_\mu$, and $\mathrm{Ri}(\mathrm{po}+1)=\mathrm{Ri}(\mathrm{po})+e_\mu$; by these two facts (Tq6), and hence (F3) and (Ap10), hold for every pair $\mathrm{po}<\mathrm{po}'$, whatever the kinds of $\mathrm{po}$ and $\mathrm{po}'$.
\end{itemize}
Each of Lemma~A.4 and (F2) concerns one plane, and in a chain they are applied at each plane separately, with the data of that plane.
\end{named}

\subsection*{Conditions on the constants}

Proposition~39 holds without any condition on the constants of the lines (C0r)--(C10) of the ordered list of constants (Section~I.8.2, Proposition~14.2): none of these constants occurs in it or in the proofs of this chapter.

Among the data of the line (C0), the hypothesis $0<\rho<1$ gives $c_i=3(1\pm\rho)\ge0$ (Section~A.6);
hence the square roots $\left(c_{I'}\cdot c_I\right)^{\frac12}$ in the kernels of Section~A.9 and the matrices $\mathrm{cc}^{\frac12}$, $\mathrm{cc}'^{\frac12}$ of Section~8.1 are real, and the identity $\left(c_{I''}\cdot c_{I'}\right)^{\frac12}\cdot\left(c_{I'}\cdot c_I\right)^{\frac12}=\left(c_{I''}\cdot c_I\right)^{\frac12}\cdot c_{I'}$ in the proof of (SO1) holds.
It also gives $c_i\le6$, whence $c_I\le6^{64v_s}$ in the proof of (SO0).
The condition $|\rho|\le1$ would suffice.
Of the weights (the row $w$ of the line (C0)), the closure relation $\sum_iw_i\hat n_i=0$ of the four standard weights (I.1.1), in the form (I.m1), (I.2.5), makes the diagonal term of $D_A+\boldsymbol{\mu}$ the scalar $M_{\mathrm{ch}}=\boldsymbol{\mu}+\frac{\rho}{2}\,\mathrm{sw}$ (Section~1.5, (Bd2), (Bd4));
step~(4) of the proof of (Ap10) rests on the fact that this term is a real scalar.
The values $w_0=w_3=6$ give $c_i=3(1\pm\rho)$ (Section~A.6) and the bound $c_I\le6^{64v_s}$ (proof of (SO0)).
Apart from the closure relation, the weights $w_1$ and $w_2$ occur only in $S_{(\mathrm{po},\mathrm{po}')}(A)$.
The number $q>0$ and the real number $\boldsymbol{\mu}$ are arbitrary;
in particular neither a sign of $\boldsymbol{\mu}$ nor a lower bound for $\boldsymbol{\mu}$ is assumed (I.K10).

\section*{Chapter 39r. Chessboard estimate and lengthening for closed cells, by transfer operators at the coordinate planes: proof of the paragraph ``Closed cells'' of Remark 39.R and of Proposition 39\,(a)}

This chapter proves the paragraph ``Closed cells'' of Remark~39.R and clause~(a) of Proposition~39;
both statements are formulated in the notions of Definition~39.0.
In Proposition~39 the clauses (b)--(f) concern contents in $\mathcal{A}\left(Q^{\circ}_{\mathfrak{c}}(\mathcal{Q})\right)$, Grassmann valued contents included, and every width $\mathrm{hw}\in\frac12\mathbb{Z}$ with $\mathrm{hw}>0$;
in this generality they are proved in Chapter~39.
Here the clauses (c)--(f) are established only for contents of the class $\mathrm{Fcl}(\mathfrak{c})$, that is, for scalar functions of the curvature on the closed cell, and only for the widths $\mathrm{hw}\ge\frac32$.

The paragraph ``Closed cells'' of Remark~39.R is proved in 9.2, Proposition~9.3 and Remark~9.4 below;
its sentence on the properties of the fermion that enter the argument is the content of Remark~5.4 below.
Section~1 proves clause~(a) of Proposition~39 in the explicit form (1.1), (1.2), on which the later sections rely.

\subsection*{0. Conventions of this chapter: setting, letters, inputs}

\subsubsection*{0.1. Setting}

Lattice units are used (Section~I.2.5, Section~I.7.1).
The photon lattice is $\mathbb{Z}^3$, with the directions $1,2,3$, and every object is of level~$0$.
Let $\rho\in(0,1)$ and $w=(6,6\sqrt{2},6\sqrt{2},6)$, let $\boldsymbol{\mu}$ be real and $q>0$, and put
\[
r_q:=\frac{2\pi}{q},
\qquad
M_{\mathrm{ch}}:=\boldsymbol{\mu}+\frac{\rho}{2}\cdot\mathrm{sw}.
\]
Statements of \cite{QEDII} are read with the following identifications.
The operator $D_A$ of \cite{QEDII} is $D_A+\boldsymbol{\mu}$, its mass parameter $M$ is $M_{\mathrm{ch}}$, its curvature $F$ is $F^{\mathrm{u}}$, and its group $W$ is $W_{\mathrm{af}}$;
the position $c$ of a mirror is $\mathrm{po}$, the width $h$ is $\mathrm{hw}$, the index $\kappa$ of a cell is $\mathfrak{c}$, and an element $g$ of $W$ is $g_{\mathrm{w}}$;
the coordinate forms $x_\mu$ are $y_\mu$, the torus $\mathbb{T}$ is $C(n)$, and $Z_{\mathbb{T}}$ is $Z_n$.
A fermion site $x=(c,g)$ stands for the point $\mathrm{geo}(x)=c+g\,v_0$ of $\mathbb{R}^3$.
Accordingly, ``$x$ lies in $\{\mathrm{po}<y_1<\mathrm{po}'\}$'' means $\mathrm{po}<\mathrm{geo}(x)_1<\mathrm{po}'$;
by Section~I.4.8 this condition does not depend on $v_0$.

All pairings are the counting pairings, and the Berezin integral is normalised as in \cite[Definition~2.8]{QEDII}.
With these normalisations, for $T=C(n)$ one has $\int_T\exp\left(-S_T(A)\right)=\det_n(D_A+\boldsymbol{\mu})$ (Proposition~46\,(e)),
and for a Borel function of $F^{\mathrm{u}}$ of the class $\mathfrak{F}(n)$ the number $Z_T[\,\cdot\,]$ of \cite[Definition~2.13]{QEDII} is the number $Z_n[\,\cdot\,]$ of (I.7.1) (Proposition~46\,(f)).
The index sets $\mathrm{Jx}(\mathrm{po})$ below, and with them the Fock spaces and the identities built on them, depend on the $48$ fermion sites over a photon site only through two numbers, both given by (In1):
over each photon site of a site plane eight bonds of type~$0$ cross the plane, and over each photon link lie eight bonds of type~$3$.
The number $\mathrm{Nx}$ of the fermion sites of a slab enters the value of the constant $C_{\mathrm{sl}}$ of (m1).

\subsubsection*{0.2. Letters}

\paragraph{Objects of \cite{QEDII}.}
The following objects of \cite{QEDII} keep the meaning they have there.
After each group follows its place in \cite{QEDII} and, where \cite{QEDII} uses other letters, these letters.
\begin{itemize}
\item[--] The lattice $\Gamma:=\sum_\mu n_\mu\mathbb{Z}\,e_\mu$ and $\mathrm{Conf}(\Gamma)$: \cite[Definition~2.6]{QEDII}.
\item[--] The lattice $\Gamma_{\mathrm{par}}:=n_2\mathbb{Z}e_2+n_3\mathbb{Z}e_3$,
the cross-section $\mathrm{Xs}:=(\Gamma_{\mathrm{par}};-1,-1)$, whose character takes the value $-1$ at $n_2e_2$ and at $n_3e_3$,
and the open slab $\mathrm{Sl}(\mathrm{po},\mathrm{po}'):=\{\mathrm{po}<y_1<\mathrm{po}'\}$:
\cite[Definition~2.17]{QEDII}; there $\Gamma_\parallel$ and $\Xi$.
\item[--] The slice torus $\Lambda_{\mathrm{s}}:=\mathbb{Z}^2/(n_2\mathbb{Z}\times n_3\mathbb{Z})$, its links $\mathrm{Lk}_{\mathrm{s}}$ and plaquettes $\mathrm{Pl}_{\mathrm{s}}$, the coboundary $d_{\mathrm{s}}$,
the slice variables $A_t$, $a_t$ and the parts $\mathrm{El}_t$, $\mathrm{Bm}_t$ of a two-cochain:
\cite[Definition~2.26]{QEDII}; there the links and the plaquettes form $\Lambda^{(1)}_{\mathrm{s}}$ and $\Lambda^{(2)}_{\mathrm{s}}$, and the parts are $F^{\mathrm{el}}_t$ and $F^{\mathrm{mag}}_t$.
Moreover $\mathrm{box}_0:=[0,r_q)^{\Lambda_{\mathrm{s}}}$.
\item[--] $\mathrm{Mi}$, $\mathrm{Cu}$, $\mathrm{Sb}$, $\mathrm{Ln}[\cdot,\cdot]$, $\mathrm{Pq}[\cdot,\cdot]$, $\mathrm{Cfg}[\cdot,\cdot]$, $\omega$, $\omega+[m]$, $\mathrm{mi}\vee\mathrm{omi}\vee\mathrm{mi}'$, $\mathrm{wb}$, $\mathrm{Ii}$:
\cite[Definition~5.2]{QEDII}; there, in the same order, $\mathcal{M}$, $\mathcal{X}$, $\mathcal{B}$, $[c,c']^{(1)}$, $[c,c']^{(2)}$, $\mathfrak{C}[c,c']$, $\varpi$, $\varpi+[k]$, $\mathfrak{m}\vee\mathring{\varpi}\vee\mathfrak{m}'$, $\mathcal{W}_{c,c'}$ and the interior integral.
\item[--] $X_{\mathrm{Xs}}(\cdot)$ and $\mathrm{Gr}_{\mathrm{Xs}}(\cdot)$, whose argument is an open interval of values of $y_1$:
\cite[Definition~5.7]{QEDII}; there $\mathbb{X}_\Xi$ and $\mathrm{Gr}_\Xi$.
$\mathrm{Ig}_{\mathrm{Xs}}[\cdot,\cdot]$ and its subspace $\mathrm{Ig}^{\mathrm{ev}}_{\mathrm{Xs}}[\cdot,\cdot]$ of the elements with even values:
\cite[Definition~5.9]{QEDII}; there $\mathcal{I}_\Xi$.
$\mathrm{Cfg}(\Gamma)$ and $\mathrm{Ig}(T)$: \cite[Definition~5.11]{QEDII}; there $\mathfrak{C}(\Gamma)$ and $\mathcal{I}(\mathbb{T})$.
$\mathrm{qu}_{\mathrm{Xs}}$: \cite[Definition~5.12]{QEDII}; there $r_\Xi$.
\item[--] The Berezin integral $\int_{(\cdot)}$, whose index is an open interval of values of $y_1$: \cite[Definition~5.16]{QEDII}.
\item[--] The translation $\mathrm{Tg}_v$ of the Grassmann variables: \cite[Lemma~5.8]{QEDII}; there $\tau_v$.
\item[--] Elements $\Phi$, $\Phi'$ of the Grassmann algebras: \cite[Definition~5.21]{QEDII}; there $G$, $G'$.
\item[--] The vector $\tau:=2(\mathrm{po}'-\mathrm{po})\,e_1$ that belongs to a slab $[\mathrm{po},\mathrm{po}']$:
\cite[Lemma~5.8\,(c) and Lemma~5.22\,(c)]{QEDII}; there $v_{c,c'}$.
\item[--] The modes, coefficients and phases $C_i$, $c_i$, $u_i$ and $C_I$, $c_I$, $u_I$ (there $\zeta_i$, $\omega_i$, $u_i$ and $\zeta_I$, $\omega_I$, $u_I$),
the insertion $\mathrm{Ins}$,
a matrix $\mathrm{mt}$ on a space $\mathrm{Fk}(\mathrm{po})$ (there $\Upsilon$),
the diagonal matrices $\mathrm{cc}$ of the coefficients (there $\omega$),
the Fock spaces $\mathrm{Fk}(\cdot)$ with the basis $(e_I)$ (there $\mathcal{F}(\cdot)$),
the charges $\mathrm{Qc}_I$ (there $\mathfrak{q}(I)$),
the unitaries $\mathrm{Wq}(a)$ (there $u(a)$),
the spaces $\mathrm{Fk}'$, $\mathrm{Fk}$ (there $\mathcal{F}'$, $\mathcal{F}$),
the half-slab matrices $\mathrm{Nn}(A)$ and $\mathrm{Qs}(A)=\mathrm{cc}^{\frac12}\,\mathrm{Nn}(A)\,\mathrm{cc}'^{\frac12}$ (there $\mathfrak{N}(A)$ and $\widehat{\mathfrak{N}}(A)$),
and the action $S_{\mathrm{hs}}$ of the half slab with the integral $\int_{\mathrm{hs}}$ over it (there $S_{\mathrm{h}}$ and $\int_{\mathbb{X}_{\mathrm{h}}}$):
\cite[Definitions~2.27, 5.20 and~5.21]{QEDII}.
\item[--] $S_{(\mathrm{po},\mathrm{po}')}(A)$ and $\mathrm{Xc}_{\mathrm{po}}(A)$: \cite[Lemma~5.18]{QEDII}; there $S_{(c,c')}(A)$ and $S^{\times}_c(A)$.
\item[--] $\mathrm{Cq}$, $\epsilon_G$, $\mathrm{Bx}$, $\mathrm{Ln}(\mathrm{Bx})$, $\mathrm{Pq}(\mathrm{Bx})$, $\mathrm{Ad}(\mathrm{Bx})$, $\mathrm{Ad}^{\mathrm{ev}}(\mathrm{Bx})$, $\mathcal{A}(\cdot)$, $\hat{A}$:
\cite[Definition~2.12]{QEDII}; there, in the same order, $C$, $\varepsilon_G$, $B$, $B^{(1)}$, $B^{(2)}$, $\mathcal{G}(B)$, $\mathcal{G}^{\mathrm{ev}}(B)$.
\item[--] $P_G$ and $\mathrm{Ug}(a)$ are the operators $P_{\mathrm{G}}$ and $U_{\mathrm{G}}(a)$ of \cite[Definition~2.28]{QEDII}, taken in the picture of 3.1;
$P_G^{+}$ is $P_G$ for the trivial representation (3.2).
\end{itemize}
Further, $\theta_{\mathrm{po}}:=\theta_{y_1,\mathrm{po}}$ and $\Theta_{\mathrm{po}}:=\Theta_{(y_1,\mathrm{po})}$ (Section~I.4.8),
$v_{\mathrm{s}}:=n_2n_3$ (Remark~39.R),
and $\mathrm{Js}_p$ denotes the Schatten classes of \cite[Section~10.3.1]{QEDI}, written $\mathscr{J}_p$ there.

In addition, the following objects of \cite{QEDII} carry other letters here:
$\mathrm{Jx}(\mathrm{po})$ is the ordered set $\mathcal{J}(c)$ of representatives of the modes \cite[Definition~5.20]{QEDII},
$\mathrm{Ml}_{\mathrm{po},\mathrm{po}'}[\Phi]$ is the slab matrix $\mathfrak{N}_{c,c'}[G]$ \cite[Definition~5.21]{QEDII},
$\mathrm{Tt}_j$ is the unitary $U_{je_1}$ \cite[Definition~5.20]{QEDII},
and $\mathrm{thq}(\epsilon)$ is the number $C_{q,\epsilon}$ \cite[Lemma~5.15\,(b)]{QEDII}.
The electric and the magnetic part of an integer two-cochain $\mathfrak{z}$ are written $\mathfrak{z}^{\mathrm{el}}_t$ and $\mathfrak{z}^{\mathrm{mag}}_t$;
in \cite[Example~5.6]{QEDII} they are $z^{\mathrm{el}}_t$ and $z^{\mathrm{mag}}_t$.
The measure spaces of \cite[Lemma~5.15\,(c)]{QEDII} are written $(\Xi,\mathrm{d}\xi)$ and $(\Xi',\mathrm{d}\xi')$, with the points $\xi$, $\xi'$;
there they are $(\mathcal{Y},\mu)$ and $(\mathcal{Y}',\mu')$.
The indicator of a site $y$ of $\Lambda_{\mathrm{s}}$ is written $\mathbf{1}_{\{y\}}$; in \cite[Definition~2.27]{QEDII} it is $\delta_y$.

\paragraph{Single letters.}
In this chapter $i,j,k,s,t$, also with primes or indices, are integers.
The letter $i$ is moreover the type of a bond, as in $M_i$, $x^{(i)}$, $w_i$, and in (In2), (In3) a mode, that is, a triple;
the letter $k$ is also the third entry of a mode, and $k_X$, $k_Y$ in (In9) are kernels.
The letters $t,t',t_0,t_1$ are times: positions of site planes, and numbers of slices and of steps.
The letter $r$ is the number of slabs, except in $r\Gamma$, $rv$, where it is the linear part of $g_{\mathrm{w}}$ (Section~I.4.7).
Real variables are written $\mathrm{va}$ (a variance) and $\mathrm{xv}$.
Elements of $\mathbb{R}^{\Lambda_{\mathrm{s}}}$, the temporal links of one step, are written $a,a',a'',a_0,a_t,\tilde a_t$;
$b$ is a link and $p$ a plaquette, while in (In9) the letters $p,p',p_i$ are exponents of Schatten classes.
An integer one-cochain is written $m$, with the components $m_b$;
$u,u'$ are real cochains, two-cochains when they are arguments of functions of the curvature.
Vectors of $\mathbb{Z}^3$ are written $v,v'$, and $y$ is a point of $\mathbb{R}^3$ or a site of $\Lambda_{\mathrm{s}}$.
An integrand is written $G$, a scalar element of $\mathrm{Ig}(T)$ or of $\mathcal{A}\left(Q^{\circ}_{\mathfrak{c}}\right)$ is written $P$,
operators and matrices are written $X,Y$, sets of modes $I,I'$, and finite-dimensional Hilbert spaces $V,V'$.
The letter $T$ always denotes the torus $C(n)$ named at the place, and $\ell=y_1$ is the form of the family of mirrors.

\subsubsection*{0.3. Inputs}

Throughout, $n_2,n_3\ge4$ are fixed and $\ell=y_1$.
All cells, cochains and fermion sites are taken modulo $\Gamma_{\mathrm{par}}$, that is, on the cylinder $\mathbb{Z}^3/\Gamma_{\mathrm{par}}$ of \cite[Section~5.1]{QEDII}.
The following facts enter in the form stated.

\paragraph{(In1) Bonds.}
(\cite[Proposition~3.6\,(b), (c)]{QEDII}, \cite[Lemma~5.18\,(a)]{QEDII}; Section~I.4.8.)
Let $\mathrm{po}\in\frac12\mathbb{Z}$.
No fermion site lies on the plane $\{y_1=\mathrm{po}\}$.
A bond whose end points lie on different sides of this plane has the form $\{x,\theta_{\mathrm{po}}x\}$, and both of its end points lie in $\{\mathrm{po}-\frac12<y_1<\mathrm{po}+\frac12\}$.
If $\mathrm{po}=t\in\mathbb{Z}$, these bonds are of type~$0$; eight of them lie over each photon site $(t,y)$, $y\in\Lambda_{\mathrm{s}}$, and they carry no phase.
If $\mathrm{po}=t+\frac12$, they are of type~$3$; they are the eight bonds over each photon link from $(t,y)$ to $(t+1,y)$, and for the lower end $x$ of such a bond
\[
U_A(x,\theta_{\mathrm{po}}x)=\exp\left(iq\,a_t(y)\right).
\]
In the coordinates $x=(c,g)$ of (I.1.2), with $c=(t,y)$, the lower ends of the bonds that cross the site plane $t$ are the $8$ among the $48$ sites over $c$ with $g\,e_3=-e_1$, and for them
\[
\hat{d}_0(x)=-g\,\hat n_0=e_1,
\qquad
\theta_t x=x^{(0)}=(c,g\,s_0),
\qquad
\mathrm{geo}(x)_1=t-(v_0)_3 .
\]
The lower ends of the bonds that cross the link plane $t+\frac12$ are the $8$ sites over $c$ with $g\,e_1=e_1$, and for them
\[
\hat{d}_3(x)=e_1,
\qquad
\theta_{t+1/2}x=x^{(3)}=(c+e_1,g\,s_3),
\qquad
\mathrm{geo}(x)_1=t+(v_0)_1 .
\]
For such a bond the hopping matrix at the lower end $x$, that is, $M_{0}(x)$ or $M_{3}(x)$ of (I.2.1), equals $3(\rho-\gamma_1)$,
and the hopping matrix at the upper end $\theta_{\mathrm{po}}x$ equals $3(\rho+\gamma_1)$; here $w_0=w_3=6$.

\paragraph{(In2) Modes and Fock spaces.}
(\cite[Definition~2.27]{QEDII} with its total orders of the modes, \cite[Definition~5.20]{QEDII}.)
Let $\mathrm{po}\in\frac12\mathbb{Z}$.
The index set $\mathrm{idx}((y_1,\mathrm{po}))$ consists of the triples $i$ formed by a bond that crosses $\{y_1=\mathrm{po}\}$, a group in $\{1,2,3,4\}$ and an index $k\in\{1,2\}$;
each crossing bond carries eight of them.
The mode $C_i$ is a linear form in the four components of $\psi$ (groups $1$, $2$), respectively of $\bar{\psi}$ (groups $3$, $4$), at the upper end $\theta_{\mathrm{po}}x$ of the bond.
With the orthonormal bases $(e^{-}_1,e^{-}_2)$ of the range of $\frac12(1-\gamma_1)$ and $(e^{+}_1,e^{+}_2)$ of the range of $\frac12(1+\gamma_1)$ fixed in \cite[Definition~2.7]{QEDII}, it is given by
\[
\begin{aligned}
&\text{group }1:\quad C_i=(e^{-}_k)^{*}\psi(\theta_{\mathrm{po}}x),
&\qquad
&\text{group }2:\quad C_i=(e^{+}_k)^{*}\psi(\theta_{\mathrm{po}}x),\\
&\text{group }3:\quad C_i=\bar{\psi}(\theta_{\mathrm{po}}x)\,e^{+}_k,
&\qquad
&\text{group }4:\quad C_i=\bar{\psi}(\theta_{\mathrm{po}}x)\,e^{-}_k .
\end{aligned}
\]
Its coefficient is
\[
c_i=\frac{w_0}{2}(1+\rho)=\frac{w_3}{2}(1+\rho)=3(1+\rho)\quad\text{for the groups }1,3,
\qquad
c_i=3(1-\rho)\quad\text{for the groups }2,4;
\]
thus $c_i>0$.
Let $\mathrm{Jx}(\mathrm{po})$ be a set of representatives of $\mathrm{idx}((y_1,\mathrm{po}))/\Gamma_{\mathrm{par}}$ with a total order,
such that $\mathrm{Jx}(\mathrm{po}+1)=\mathrm{Jx}(\mathrm{po})+e_1$ with the transported order.
Then $\#|\mathrm{Jx}(\mathrm{po})|=8\cdot8\cdot v_{\mathrm{s}}=64\,v_{\mathrm{s}}$.
For a subset $I=\{i_1<i_2<\dots\}$ of $\mathrm{Jx}(\mathrm{po})$ let $C_I:=C_{i_1}C_{i_2}\cdots$ be the product of the $C_i$, $i\in I$, taken in increasing order, and let $c_I:=\prod_{i\in I}c_i$.
The $C_i$, $i\in\mathrm{Jx}(\mathrm{po})$, form a basis of the span of the generators of $\mathrm{Gr}_{\mathrm{Xs}}$ at the upper ends of the crossing bonds, and they lie in $\mathrm{Gr}_{\mathrm{Xs}}((\mathrm{po},\mathrm{po}+\frac12))$;
the $\Theta_{\mathrm{po}}(C_i)$ form a basis of the span of the generators at the lower ends, and they lie in $\mathrm{Gr}_{\mathrm{Xs}}((\mathrm{po}-\frac12,\mathrm{po}))$.

The space $\mathrm{Fk}(\mathrm{po})$ is the Hilbert space with the orthonormal basis $(e_I)$, where $I$ ranges over the subsets of $\mathrm{Jx}(\mathrm{po})$;
thus $\dim\mathrm{Fk}(\mathrm{po})=2^{64v_{\mathrm{s}}}$.
The matrix $\mathrm{cc}_{\mathrm{po}}:=\operatorname{diag}(c_I)$ is diagonal and positive.
For $j\in\mathbb{Z}$ the map $\mathrm{Tt}_j$, $\mathrm{Tt}_je_I:=e_{I+je_1}$, is a unitary operator from $\mathrm{Fk}(\mathrm{po})$ onto $\mathrm{Fk}(\mathrm{po}+j)$, and \cite[Definition~5.20\,(a)]{QEDII}
\[
\mathrm{Tt}_j\cdot\mathrm{Tt}_{j'}=\mathrm{Tt}_{j+j'},
\qquad
\mathrm{Tt}_j\cdot\mathrm{cc}_{\mathrm{po}}=\mathrm{cc}_{\mathrm{po}+j}\cdot\mathrm{Tt}_j .
\]

Let now $\mathrm{po}=t+\frac12$.
A mode over the link from $(t,y)$ to $(t+1,y)$ has the charge $\mathrm{Qc}_i:=\mathbf{1}_{\{y\}}$ if it belongs to the group $1$ or $2$, and $\mathrm{Qc}_i:=-\mathbf{1}_{\{y\}}$ if it belongs to the group $3$ or $4$;
here $\mathbf{1}_{\{y\}}\in\mathbb{Z}^{\Lambda_{\mathrm{s}}}$ is the indicator of the site $y$.
Put $\mathrm{Qc}_I:=\sum_{i\in I}\mathrm{Qc}_i\in\mathbb{Z}^{\Lambda_{\mathrm{s}}}$ and, for $a\in\mathbb{R}^{\Lambda_{\mathrm{s}}}$,
\[
\mathrm{Wq}(a):=\operatorname{diag}\left(\exp\left(iq\,\langle a,\mathrm{Qc}_I\rangle\right)\right).
\]
Then $u_I(A)=\exp\left(iq\,\langle a_t,\mathrm{Qc}_I\rangle\right)$, that is, $\operatorname{diag}\left(u_I(A)\right)=\mathrm{Wq}(a_t)$.
Since $\mathrm{Qc}_{I+e_1}=\mathrm{Qc}_I$, one has $\mathrm{Tt}_j\cdot\mathrm{Wq}(a)=\mathrm{Wq}(a)\cdot\mathrm{Tt}_j$.
The operator $\mathrm{Wq}(a)$ is unitary,
\[
\mathrm{Wq}(a)\cdot\mathrm{Wq}(a')=\mathrm{Wq}(a+a'),
\qquad
\mathrm{Wq}(a)^{*}=\mathrm{Wq}(-a),
\]
and $\mathrm{Wq}(a)=1$ for $a\in r_q\mathbb{Z}^{\Lambda_{\mathrm{s}}}$, because $q\cdot r_q=2\pi$.
For $\mathrm{po}=t\in\mathbb{Z}$ one has $u_I=1$.

\emph{Identification.}
From now on $\mathrm{Fk}(\mathrm{po}+j)$ is identified with $\mathrm{Fk}(\mathrm{po})$ through $\mathrm{Tt}_j$, $j\in\mathbb{Z}$.
By the three relations just proved for $\mathrm{Tt}_j$, this identification is coherent and respects $\mathrm{cc}$ and $\mathrm{Wq}$.
Two spaces remain, $\mathrm{Fk}':=\mathrm{Fk}(0)$ for the site planes and $\mathrm{Fk}:=\mathrm{Fk}(\frac12)$ for the link planes,
together with two matrices $\mathrm{cc}'$, $\mathrm{cc}$ and the representation $\mathrm{Wq}$ on $\mathrm{Fk}$.
These are the objects of \cite[Definition~2.27]{QEDII}, and \cite[Definition~5.32]{QEDII} makes the same identification.

\paragraph{(In3) Insertions; the crossing expansion.}
(\cite[Definition~2.27]{QEDII}, \cite[Proposition~3.20]{QEDII}, \cite[Section~5.5]{QEDII}.)
For a matrix $\mathrm{mt}=(\mathrm{mt}_{II'})$ on $\mathrm{Fk}(\mathrm{po})$, the indices $I,I'$ being subsets of $\mathrm{Jx}(\mathrm{po})$, put
\[
\mathrm{Ins}_{\mathrm{po}}[\mathrm{mt}]:=\sum_{I,I'}\mathrm{mt}_{II'}\cdot\Theta_{\mathrm{po}}(C_{I'})\,C_I ;
\]
the row index belongs to the upper side.
This is an element of $\mathrm{Gr}_{\mathrm{Xs}}((\mathrm{po}-\frac12,\mathrm{po}+\frac12))$, and it is even if $\mathrm{mt}$ is even, that is, if $\mathrm{mt}_{II'}=0$ unless $\#|I|-\#|I'|$ is even.
Let $\mathrm{Xc}_{\mathrm{po}}(A)$ be the sum of those terms of $\bar{\psi}(D_A+\boldsymbol{\mu})\psi$ that belong to the bonds crossing $\{y_1=\mathrm{po}\}$.
A crossing bond with the lower end $x$ contributes the two terms
\[
-\bar{\psi}(x)\,3(\rho-\gamma_1)\,U_A(x,\theta_{\mathrm{po}}x)\,\psi(\theta_{\mathrm{po}}x)
-\bar{\psi}(\theta_{\mathrm{po}}x)\,3(\rho+\gamma_1)\,U_A(\theta_{\mathrm{po}}x,x)\,\psi(x).
\]
Since
\[
3(\rho-\gamma_1)=3(1+\rho)\cdot\frac{1-\gamma_1}{2}-3(1-\rho)\cdot\frac{1+\gamma_1}{2},
\qquad
3(\rho+\gamma_1)=3(1+\rho)\cdot\frac{1+\gamma_1}{2}-3(1-\rho)\cdot\frac{1-\gamma_1}{2},
\]
\cite[Proposition~3.20]{QEDII} gives
\[
-\mathrm{Xc}_{\mathrm{po}}(A)=\sum_{i\in\mathrm{Jx}(\mathrm{po})}c_i\cdot u_i(A)\cdot\Theta_{\mathrm{po}}(C_i)\,C_i
\]
and, in the form of \cite[Section~5.5]{QEDII},
\[
\exp\left(-\mathrm{Xc}_{\mathrm{po}}(A)\right)=\sum_Ic_I\cdot u_I(A)\cdot\Theta_{\mathrm{po}}(C_I)\,C_I=\mathrm{Ins}_{\mathrm{po}}\left[\mathrm{cc}_{\mathrm{po}}\cdot\operatorname{diag}\left(u_I(A)\right)\right];
\]
here $u_i=U_A(x,\theta_{\mathrm{po}}x)$ for the modes of the groups $1$, $2$ of the bond, and $u_i$ is the complex conjugate of this number for the groups $3$, $4$.

\paragraph{(In4) The action of a slab.}
(\cite[Lemma~5.18\,(b), (c)]{QEDII}.)
For $\mathrm{po}<\mathrm{po}'$ in $\frac12\mathbb{Z}$ let $S_{(\mathrm{po},\mathrm{po}')}(A)$ be the sum of the terms $\bar{\psi}(x)(\,\cdot\,)\psi(x')$ of $\bar{\psi}(D_A+\boldsymbol{\mu})\psi$ with $x,x'\in X_{\mathrm{Xs}}((\mathrm{po},\mathrm{po}'))$,
the diagonal terms $M_{\mathrm{ch}}\cdot\bar{\psi}(x)\psi(x)$ included.
It is an even element of $\mathrm{Gr}_{\mathrm{Xs}}((\mathrm{po},\mathrm{po}'))$.
Its coefficients are $M_{\mathrm{ch}}$ and the entries of the matrices $-M_{i}(x)\cdot U_A(x,x^{(i)})$;
only $i=3$ carries a phase, namely $\exp\left(\pm iqA(b)\right)$ with $b\in\mathrm{Ln}[\mathrm{po},\mathrm{po}']$.
For the last assertion, let a bond of type~$3$ over the photon link $b$ have both end points in the open slab.
If $b$ were cut by one of the two planes, the bond would cross that plane, by (In1); so $b$ is not cut.
By \cite[Proposition~3.6\,(c)]{QEDII} the bond then lies, for each of the two planes, on the side on which the open link $b$ lies, unless $b$ lies in the plane.
In all cases $b$ is contained in $\{\mathrm{po}\le y_1\le\mathrm{po}'\}$.
For $\mathrm{po}<\mathrm{po}'<\mathrm{po}''$ one has
\[
S_{(\mathrm{po},\mathrm{po}'')}=S_{(\mathrm{po},\mathrm{po}')}+\mathrm{Xc}_{\mathrm{po}'}+S_{(\mathrm{po}',\mathrm{po}'')}.
\]

\paragraph{(In5) Slab matrices.}
(\cite[Definition~5.21, Lemma~5.22]{QEDII}, read with the identification of (In2).)
For $\mathrm{po}<\mathrm{po}'$ and $\Phi\in\mathrm{Gr}_{\mathrm{Xs}}((\mathrm{po},\mathrm{po}'))$ let $\mathrm{Ml}_{\mathrm{po},\mathrm{po}'}[\Phi]$ be the matrix with the entries
\[
\mathrm{Ml}_{\mathrm{po},\mathrm{po}'}[\Phi]_{I',I}:=\int_{(\mathrm{po},\mathrm{po}')}C^{(\mathrm{po})}_I\,\Phi\,\Theta_{\mathrm{po}'}\left(C^{(\mathrm{po}')}_{I'}\right),
\]
where the upper index names the plane of the modes, the column index $I$ is a subset of $\mathrm{Jx}(\mathrm{po})$ and the row index $I'$ a subset of $\mathrm{Jx}(\mathrm{po}')$.
It is a linear map $\mathrm{Fk}(\mathrm{po})\to\mathrm{Fk}(\mathrm{po}')$, and it depends linearly on $\Phi$.
The rules (F2)--(F5) below are the rules (b)--(e) of \cite[Lemma~5.22]{QEDII}.
\begin{itemize}
\item[(F2)]
Let $\mathrm{po}<\mathrm{po}'<\mathrm{po}''$, $\Phi\in\mathrm{Gr}_{\mathrm{Xs}}((\mathrm{po},\mathrm{po}'))$ and $\Phi'\in\mathrm{Gr}_{\mathrm{Xs}}((\mathrm{po}',\mathrm{po}''))$.
For every matrix $\mathrm{mt}$ on $\mathrm{Fk}(\mathrm{po}')$,
\[
\mathrm{Ml}_{\mathrm{po},\mathrm{po}''}\left[\Phi\,\mathrm{Ins}_{\mathrm{po}'}[\mathrm{mt}]\,\Phi'\right]=\mathrm{Ml}_{\mathrm{po}',\mathrm{po}''}[\Phi']\cdot\mathrm{mt}\cdot\mathrm{Ml}_{\mathrm{po},\mathrm{po}'}[\Phi].
\]
\item[(F3)]
Let $\mathrm{po}'':=2\mathrm{po}'-\mathrm{po}$.
Then
\[
\mathrm{Ml}_{\mathrm{po}',\mathrm{po}''}\left[\Theta_{\mathrm{po}'}\Phi\right]=\mathrm{Ml}_{\mathrm{po},\mathrm{po}'}[\Phi]^{*}.
\]
In \cite[Lemma~5.22\,(c)]{QEDII} the right-hand side carries in addition, on the left, the unitary $U_\tau$ of \cite[Definition~5.20]{QEDII} for the vector $\tau=2(\mathrm{po}'-\mathrm{po})\,e_1\in\mathbb{Z}e_1$ of the slab $[\mathrm{po},\mathrm{po}']$;
this unitary is the identification of (In2).
\begin{proof}[The identity behind (F3)]
The rule (F3) rests on the identity
\[
\int_{(\mathrm{po}',\mathrm{po}'')}\Theta_{\mathrm{po}'}(\,\cdot\,)=\overline{\int_{(\mathrm{po},\mathrm{po}')}(\,\cdot\,)}
\qquad\text{on }\mathrm{Gr}_{\mathrm{Xs}}((\mathrm{po},\mathrm{po}'))
\]
of \cite[Lemma~5.19]{QEDII}.
For a finite set $Y$ of fermion sites and the free algebra $\mathrm{Gr}(Y)$ it is the identity of \cite[Proposition~5.17\,(b)]{QEDII}, and on the cylinder it follows from that case.
Indeed, let $Y$ be a set of representatives of $X_{\mathrm{Xs}}((\mathrm{po},\mathrm{po}'))$.
The quotient map $\mathrm{qu}_{\mathrm{Xs}}$ identifies $\mathrm{Gr}(Y)$ with $\mathrm{Gr}_{\mathrm{Xs}}((\mathrm{po},\mathrm{po}'))$ and $\int_Y$ with $\int_{(\mathrm{po},\mathrm{po}')}$.
The reflection $\theta_{\mathrm{po}'}$ commutes with the translations by $\Gamma_{\mathrm{par}}$, so that $\theta_{\mathrm{po}'}Y$ is a set of representatives of $X_{\mathrm{Xs}}((\mathrm{po}',\mathrm{po}''))$.
Finally $\mathrm{qu}_{\mathrm{Xs}}\circ\Theta_{\mathrm{po}'}=\Theta_{\mathrm{po}'}\circ\mathrm{qu}_{\mathrm{Xs}}$ by \cite[Definition~5.12]{QEDII}.
\end{proof}
\item[(F4)]
For $j\in\mathbb{Z}$,
\[
\mathrm{Ml}_{\mathrm{po}+j,\mathrm{po}'+j}\left[\mathrm{Tg}_{je_1}\Phi\right]=\mathrm{Ml}_{\mathrm{po},\mathrm{po}'}[\Phi].
\]
In \cite[Lemma~5.22\,(d)]{QEDII} the right-hand side is $U_v\cdot\mathrm{Ml}_{\mathrm{po},\mathrm{po}'}[\Phi]\cdot U_v^{-1}$ with $v=je_1$, and $U_v$ is the identification of (In2).
\item[(F5)]
Let $n_1\ge4$, $T=C(n_1,n_2,n_3)$, and let $\mathrm{po}_0<\mathrm{po}_1<\dots<\mathrm{po}_r=\mathrm{po}_0+n_1$ be points of $\frac12\mathbb{Z}$, where $r\ge1$.
Each class of fermion sites of $T$ has exactly one representative modulo $\Gamma_{\mathrm{par}}$ in $\{\mathrm{po}_0<y_1<\mathrm{po}_0+n_1\}$.
Hence $\mathrm{Gr}_T=\mathrm{Gr}_{\mathrm{Xs}}((\mathrm{po}_0,\mathrm{po}_r))$ and $\int_T=\int_{(\mathrm{po}_0,\mathrm{po}_r)}$,
and a field at a site $x$ with $\mathrm{po}_0-n_1<\mathrm{geo}(x)_1<\mathrm{po}_0$ is read as the negative of the field at $x+n_1e_1$.
Let $\Phi_k\in\mathrm{Gr}_{\mathrm{Xs}}((\mathrm{po}_{k-1},\mathrm{po}_k))$, $k=1,\dots,r$, be even, let $\mathrm{mt}_k$, $k=0,\dots,r-1$, be even matrices on $\mathrm{Fk}(\mathrm{po}_k)$, and put $\mathrm{Ml}_k:=\mathrm{Ml}_{\mathrm{po}_{k-1},\mathrm{po}_k}[\Phi_k]$.
Then
\[
\int_T\mathrm{Ins}_{\mathrm{po}_0}[\mathrm{mt}_0]\,\Phi_1\,\mathrm{Ins}_{\mathrm{po}_1}[\mathrm{mt}_1]\,\Phi_2\cdots\mathrm{Ins}_{\mathrm{po}_{r-1}}[\mathrm{mt}_{r-1}]\,\Phi_r
=\operatorname{tr}_{\mathrm{Fk}(\mathrm{po}_0)}\left(\mathrm{Ml}_r\cdot\mathrm{mt}_{r-1}\cdot\mathrm{Ml}_{r-1}\cdots\mathrm{mt}_1\cdot\mathrm{Ml}_1\cdot\mathrm{mt}_0\right).
\]
In \cite[Lemma~5.22\,(e)]{QEDII} two further factors stand under the trace, $U_g^{-1}$ and $U_{\mathrm{f}}^{\frac{1+\chi(g)}{2}}$, where $g$ is the closing vector and $U_{\mathrm{f}}$ the fermion parity.
Here $g=n_1e_1$ and $\chi(g)=-1$; so the second factor is the identity, and $U_g$ maps $e_I$ to $e_{I+n_1e_1}$, which is the identification of (In2).
\end{itemize}

\paragraph{(In6) Reflection and translation of the action.}
(\cite[Proposition~3.16\,(d)]{QEDII} with the rule for the coefficients in \cite[Section~2.3]{QEDII}; Section~I.1.3, (I.2.1), (I.2.5), Section~I.4.7, (I.O7).)
Let $\mathrm{po}<\mathrm{po}'$, $\mathrm{po}''=2\mathrm{po}'-\mathrm{po}$, and let $A'$ be a real one-cochain on $\mathrm{Ln}[\mathrm{po}',\mathrm{po}'']$.
Then
\[
\Theta_{\mathrm{po}'}\left(S_{(\mathrm{po},\mathrm{po}')}(\theta_{\mathrm{po}'}^{*}A')\right)=S_{(\mathrm{po}',\mathrm{po}'')}(A');
\]
the term of the oriented bond $(x,x')$ goes to the term of $(\theta_{\mathrm{po}'}x',\theta_{\mathrm{po}'}x)$, and $\Theta_{\mathrm{po}'}$ is the antilinear map of Section~I.4.8 on the Grassmann algebra.
This is \cite[Proposition~3.16\,(d)]{QEDII}.
For $j\in\mathbb{Z}$ and a one-cochain $A'$ on $\mathrm{Ln}[\mathrm{po}+j,\mathrm{po}'+j]$ one has moreover
\[
\mathrm{Tg}_{je_1}\,S_{(\mathrm{po},\mathrm{po}')}(\tau_{je_1}^{*}A')=S_{(\mathrm{po}+j,\mathrm{po}'+j)}(A').
\]
\begin{proof}[Proof of the second identity]
Let $v:=je_1$.
The automorphism $\mathrm{Tg}_v$ maps $\psi(x)$ to $\psi(x+v)$ and $\bar{\psi}(x)$ to $\bar{\psi}(x+v)$ \cite[Lemma~5.8]{QEDII}.
A translation moves the photon site and keeps the label:
for $x=(c,g)$ one has $x+v=(c+v,g)$, $\mathrm{geo}(x+v)=\mathrm{geo}(x)+v$ and $(x+v)^{(i)}=x^{(i)}+v$ (Section~I.1.3 with (I.s1)).
Further $M_{i}(x+v)=M_{i}(x)$, since $M_{i}(x)$ depends on the label only (I.2.1).
The phase of a bond of type $0$, $1$ or $2$ equals $1$ for every cochain, and
\[
U_{\tau_v^{*}A'}(x,x^{(3)})=\exp\left(iq\,A'(c+v,c_{3}+v)\right)=U_{A'}(x+v,x^{(3)}+v),
\]
because $(\tau_v^{*}A')(c,c')=A'(c+v,c'+v)$ (Section~I.4.7).
Consequently $\mathrm{Tg}_v$ maps the term $-\bar{\psi}(x)\,M_{i}(x)\,U_{\tau_v^{*}A'}(x,x^{(i)})\,\psi(x^{(i)})$ of (I.2.5) to the term of the oriented bond $(x+v,x^{(i)}+v)$ at $A'$,
and it maps $M_{\mathrm{ch}}\cdot\bar{\psi}(x)\psi(x)$ to $M_{\mathrm{ch}}\cdot\bar{\psi}(x+v)\psi(x+v)$.
Finally $x\mapsto x+v$ maps $X_{\mathrm{Xs}}((\mathrm{po},\mathrm{po}'))$ onto $X_{\mathrm{Xs}}((\mathrm{po}+j,\mathrm{po}'+j))$.
Written for operators, this is (I.O7).
\end{proof}

\paragraph{(In7) Configurations and classes of integrands.}
(\cite[Definitions~5.2, 5.7, 5.9, 5.11 and~5.12]{QEDII}, \cite[Proposition~5.17\,(c)]{QEDII}.)
\begin{itemize}
\item[(i)]
\emph{Cells.}
Let $t\in\mathbb{Z}$.
The spatial links $(t;b)$, $b\in\mathrm{Lk}_{\mathrm{s}}$, and the magnetic plaquettes $(t;p')$, $p'\in\mathrm{Pl}_{\mathrm{s}}$, lie in the plane $\{y_1=t\}$.
The temporal links $(t,y)$ from $(t,y)$ to $(t+1,y)$, the electric plaquettes $(t,b)$ over $b$ and the cubes $(t,p')$ over $p'$ lie between the planes $t$ and $t+1$;
they form the step~$t$.
The set $\mathrm{Mi}(\mathrm{po})$ consists of the cells in $\{y_1=\mathrm{po}\}$; it is empty for $\mathrm{po}\in\mathbb{Z}+\frac12$.
The set $\mathrm{Cu}(\mathrm{po})$ consists of the cells cut by this plane: for $\mathrm{po}=t+\frac12$ these are the cells of the step~$t$, and for $\mathrm{po}\in\mathbb{Z}$ there are none.
The set $\mathrm{Sb}(\mathrm{po},\mathrm{po}')$ consists of the cells that are contained in $\{\mathrm{po}\le y_1\le\mathrm{po}'\}$ and in neither of the two planes.
Finally $\mathrm{Ln}[\mathrm{po},\mathrm{po}']$ and $\mathrm{Pq}[\mathrm{po},\mathrm{po}']$ are the sets of the links and of the plaquettes contained in the closed slab.
An upper index $1$, $2$, $3$ denotes the dimension.
\item[(ii)]
\emph{Configurations.}
Let $\mathrm{tl}:=\lceil\mathrm{po}\rceil$ and $\mathrm{tu}:=\lfloor\mathrm{po}'\rfloor$.
A configuration $\omega=(A,\mathfrak{z})\in\mathbb{R}^{\mathrm{Ln}[\mathrm{po},\mathrm{po}']}\times\mathbb{Z}^{\mathrm{Pq}[\mathrm{po},\mathrm{po}']}$ consists of the slices
$\mathrm{mi}_t=(A_t,\mathfrak{z}^{\mathrm{mag}}_t)\in\mathbb{R}^{\mathrm{Lk}_{\mathrm{s}}}\times\mathbb{Z}^{\mathrm{Pl}_{\mathrm{s}}}$, $\mathrm{tl}\le t\le\mathrm{tu}$,
and of the steps $(a_t,\mathfrak{z}^{\mathrm{el}}_t)\in\mathbb{R}^{\Lambda_{\mathrm{s}}}\times\mathbb{Z}^{\mathrm{Lk}_{\mathrm{s}}}$, $\mathrm{tl}\le t<\mathrm{tu}$.
The curvature $F^{\mathrm{u}}(\omega):=dA+r_q\cdot\mathfrak{z}$ has the parts \cite[Definition~2.26]{QEDII}
\[
\mathrm{Bm}_t=d_{\mathrm{s}}A_t+r_q\cdot\mathfrak{z}^{\mathrm{mag}}_t\quad(\mathrm{tl}\le t\le\mathrm{tu}),
\qquad
\mathrm{El}_t=A_{t+1}-A_t-d_{\mathrm{s}}a_t+r_q\cdot\mathfrak{z}^{\mathrm{el}}_t\quad(\mathrm{tl}\le t<\mathrm{tu}),
\]
and \cite[Example~5.6]{QEDII}
\[
(d\mathfrak{z})(t,p')=\mathfrak{z}^{\mathrm{mag}}_{t+1}(p')-\mathfrak{z}^{\mathrm{mag}}_t(p')-(d_{\mathrm{s}}\mathfrak{z}^{\mathrm{el}}_t)(p').
\]
Put
\[
\mathrm{Cfg}[\mathrm{po},\mathrm{po}']:=\left\{\omega:\ \mathfrak{z}^{\mathrm{mag}}_{t+1}=\mathfrak{z}^{\mathrm{mag}}_t+d_{\mathrm{s}}\mathfrak{z}^{\mathrm{el}}_t,\ \mathrm{tl}\le t<\mathrm{tu}\right\}.
\]
For $m\in\mathbb{Z}^{\mathrm{Ln}[\mathrm{po},\mathrm{po}']}$ let $\omega+[m]:=(A+r_q\cdot m,\ \mathfrak{z}-dm)$.
The curvature $F^{\mathrm{u}}$, the coboundary $d\mathfrak{z}$ and the phases $\exp\left(iqA(b)\right)$ are invariant under $\omega\mapsto\omega+[m]$.
\item[(iii)]
\emph{Slab integrands.}
The class $\mathrm{Ig}_{\mathrm{Xs}}[\mathrm{po},\mathrm{po}']$ consists of the Borel maps $G:\mathrm{Cfg}[\mathrm{po},\mathrm{po}']\to\mathrm{Gr}_{\mathrm{Xs}}((\mathrm{po},\mathrm{po}'))$ with the two properties (Ig1), (Ig2), which are the properties (a), (b) of \cite[Definition~5.9]{QEDII}:
\begin{itemize}
\item[(Ig1)] $G(\omega+[m])=G(\omega)$ for all $m\in\mathbb{Z}^{\mathrm{Ln}[\mathrm{po},\mathrm{po}']}$;
\item[(Ig2)] there are $\mathrm{Cq}<\infty$ and $\epsilon_G\in(0,\frac12]$ such that every coefficient of $G(\omega)$ in a monomial basis has modulus at most
\[
\mathrm{Cq}\cdot\exp\left(\left(\frac12-\epsilon_G\right)\sum_{p\in\mathrm{Sb}^2(\mathrm{po},\mathrm{po}')}F^{\mathrm{u}}_p(\omega)^2\right).
\]
\end{itemize}
The subspace of the $G$ with even values is $\mathrm{Ig}^{\mathrm{ev}}_{\mathrm{Xs}}[\mathrm{po},\mathrm{po}']$.
On these classes act the translation and the reflection
\begin{gather*}
(\mathrm{Tg}_vG)(\omega'):=\mathrm{Tg}_v\left(G(\tau_v^{*}\omega')\right)\in\mathrm{Ig}_{\mathrm{Xs}}[\mathrm{po}+v_1,\mathrm{po}'+v_1],
\\
(\Theta_{\mathrm{po}'}G)(\omega'):=\Theta_{\mathrm{po}'}\left(G(\theta_{\mathrm{po}'}^{*}\omega')\right)\in\mathrm{Ig}_{\mathrm{Xs}}[\mathrm{po}',2\mathrm{po}'-\mathrm{po}],
\end{gather*}
and the following product.
For $\mathrm{po}<\mathrm{po}'<\mathrm{po}''$ and $G'\in\mathrm{Ig}_{\mathrm{Xs}}[\mathrm{po}',\mathrm{po}'']$ with the constants $(\mathrm{Cq}',\epsilon_{G'})$ put
\[
(GG')(\omega):=G(\omega_1)\,G'(\omega_2)\in\mathrm{Ig}_{\mathrm{Xs}}[\mathrm{po},\mathrm{po}''],
\]
where $\omega_1$ and $\omega_2$ are the restrictions of $\omega\in\mathrm{Cfg}[\mathrm{po},\mathrm{po}'']$ to the two closed slabs.
The product $GG'$ has the constants $(\mathrm{Cq}\cdot\mathrm{Cq}',\min(\epsilon_G,\epsilon_{G'}))$, and its values are even if those of $G$ and $G'$ are.
\begin{proof}[Proof of the properties of the product]
The steps of the two slabs are steps of the slab $[\mathrm{po},\mathrm{po}'']$; hence $\omega_1\in\mathrm{Cfg}[\mathrm{po},\mathrm{po}']$ and $\omega_2\in\mathrm{Cfg}[\mathrm{po}',\mathrm{po}'']$.
The restrictions are continuous maps, so $GG'$ is Borel.
The restriction of $\omega+[m]$ to either closed slab is $\omega_i+[m_i]$, $i=1,2$, with the restriction $m_i$ of $m$ to that slab; this gives (Ig1).
By (In1) no fermion site lies on the plane $\{y_1=\mathrm{po}'\}$.
Therefore $X_{\mathrm{Xs}}((\mathrm{po},\mathrm{po}''))$ is the disjoint union of $X_{\mathrm{Xs}}((\mathrm{po},\mathrm{po}'))$ and $X_{\mathrm{Xs}}((\mathrm{po}',\mathrm{po}''))$;
a monomial of $\mathrm{Gr}_{\mathrm{Xs}}((\mathrm{po},\mathrm{po}''))$ is, up to a sign, the product of a monomial of $\mathrm{Gr}_{\mathrm{Xs}}((\mathrm{po},\mathrm{po}'))$ and a monomial of $\mathrm{Gr}_{\mathrm{Xs}}((\mathrm{po}',\mathrm{po}''))$,
and its coefficient in $G(\omega_1)\,G'(\omega_2)$ is, up to the same sign, the product of the two coefficients.
The sets $\mathrm{Sb}^2(\mathrm{po},\mathrm{po}')$ and $\mathrm{Sb}^2(\mathrm{po}',\mathrm{po}'')$ are disjoint subsets of $\mathrm{Sb}^2(\mathrm{po},\mathrm{po}'')$ \cite[Lemma~5.4\,(d)]{QEDII},
$F^{\mathrm{u}}(\omega_i)$ is the restriction of $F^{\mathrm{u}}(\omega)$, and $\frac12-\min(\epsilon_G,\epsilon_{G'})\ge0$; this gives (Ig2).
\end{proof}
\item[(iv)]
\emph{Torus integrands.}
Let $T=C(n)$ with $n_\mu\ge4$, and
\[
\mathrm{Cfg}(\Gamma):=\left\{(A,\mathfrak{z})\in C^{1}(n)\times C^{2}(n;\mathbb{Z}):\ d\mathfrak{z}=0\right\}.
\]
The class $\mathrm{Ig}(T)$ consists of the Borel maps $G:\mathrm{Cfg}(\Gamma)\to\mathrm{Gr}_T$ with $G(A+r_q\cdot m,\mathfrak{z}-dm)=G(A,\mathfrak{z})$ for $m\in C^{1}(n;\mathbb{Z})$
whose coefficients are bounded by $\mathrm{Cq}\cdot\exp\left(\left(\frac12-\epsilon_G\right)\|F^{\mathrm{u}}\|^2\right)$ with some $\mathrm{Cq}<\infty$ and $\epsilon_G\in(0,\frac12]$.
For $G\in\mathrm{Ig}(T)$,
\[
Z_n[G]:=\sum_{\substack{\mathfrak{z}\in C^{2}(n;\mathbb{Z})\\ d\mathfrak{z}=0}}\ \int_{[0,r_q)^{\{\text{links of }C(n)\}}}\mathrm{d}A\,\exp\left(-\frac12\|F^{\mathrm{u}}\|^2\right)\int_T\exp\left(-S_T(A)\right)G(A,\mathfrak{z}),
\]
which converges absolutely, coefficient by coefficient \cite[Proposition~5.17\,(c)]{QEDII}.
\item[(vi)]
\emph{The box is immaterial} \cite[Lemma~5.3]{QEDII}.
Let $t<t'$ be integers and $\Phi$ a Borel function on $\mathbb{R}^{\mathrm{Ln}[t,t']}\times\mathbb{Z}^{\mathrm{Pq}[t,t']}$, and put
\[
\mathrm{Ii}_{t,t'}\Phi(\mathrm{mi},\mathrm{mi}'):=\sum_{\mathfrak{z}\in\mathbb{Z}^{\mathrm{Sb}^2(t,t')}}\ \int_{[0,r_q)^{\mathrm{Sb}^1(t,t')}}\mathrm{d}A\,\Phi(\mathrm{mi}\vee\mathrm{omi}\vee\mathrm{mi}'),
\]
where $\mathrm{mi}$, $\mathrm{mi}'$ are the parts on the planes $t$, $t'$ and $\mathrm{omi}=(A,\mathfrak{z})$ is the part on $\mathrm{Sb}(t,t')$.
If $\mathrm{Ii}_{t,t'}|\Phi|<\infty$ and $\Phi(\omega+[m])=\Phi(\omega)$ for all $m\in\mathbb{Z}^{\mathrm{Sb}^1(t,t')}$,
then $\mathrm{Ii}_{t,t'}\Phi$ does not change when, for each $b\in\mathrm{Sb}^1(t,t')$, the interval $[0,r_q)$ is replaced by an arbitrary interval of length $r_q$.
\item[(vii)]
\emph{(K2)}, that is, \cite[Lemma~5.15\,(b)]{QEDII}.
For $\epsilon>0$,
\[
\mathrm{thq}(\epsilon):=\sup_{\mathrm{xv}\in\mathbb{R}}\ \sum_{m\in\mathbb{Z}}\exp\left(-\epsilon\left(\mathrm{xv}+r_q\cdot m\right)^2\right)\le2+r_q^{-1}\left(\frac{\pi}{\epsilon}\right)^{\frac12}.
\]
\end{itemize}

\paragraph{(In9) Trace ideals.}
(\cite[Lemma~10.6]{QEDI}; \cite[Lemma~5.15\,(c)]{QEDII}.)
For a real number $p$ with $1\le p<\infty$, $\mathrm{Js}_p$ is the Schatten class of a separable Hilbert space, with $\|Y\|_p^p=\operatorname{Tr}|Y|^p$.
\begin{itemize}
\item[(a)] If $Y_i\in\mathrm{Js}_{p_i}$ and $\sum_{i=1}^k\frac{1}{p_i}=1$, then $\left|\operatorname{Tr}(Y_1\cdots Y_k)\right|\le\prod_i\|Y_i\|_{p_i}$.
\item[(b)] $\|Y\|_p\le\|Y\|_{p'}$ for $p\ge p'$; for even integers $p$, $\|Y\|_p^p=\operatorname{Tr}\left(Y^{*}\cdot Y\right)^{\frac{p}{2}}$.
\item[(c)] If $\mathrm{Bp}\ge0$ is of trace class and $\mathrm{Bp}\ne0$, then the function $\mathrm{xv}\mapsto\log\operatorname{Tr}\mathrm{Bp}^{\mathrm{xv}}$ of the real variable $\mathrm{xv}$ is convex on $[1,\infty)$.
\item[(K3)] (\cite[Lemma~5.15\,(c)]{QEDII}.)
Let $(\Xi,\mathrm{d}\xi)$ and $(\Xi',\mathrm{d}\xi')$ be $\sigma$-finite measure spaces whose $L^2$ spaces are separable.
A kernel in $L^2(\mathrm{d}\xi'\times\mathrm{d}\xi)$ with values in the linear maps between two finite-dimensional Hilbert spaces defines a Hilbert--Schmidt operator, whose Hilbert--Schmidt norm is the $L^2$ norm of the kernel.
The adjoint operator has the adjoint kernel with the two arguments exchanged, the product of two such operators has the composed kernel, and
\[
\operatorname{Tr}(Y\cdot X)=\int\mathrm{d}\xi\int\mathrm{d}\xi'\,\operatorname{tr}\left(k_Y(\xi,\xi')\,k_X(\xi',\xi)\right),
\]
the double integral converging absolutely.
For values in $\operatorname{Hom}(V,V')$ one applies the scalar statement to the measure spaces $\Xi\times\{1,\dots,\dim V\}$ and $\Xi'\times\{1,\dots,\dim V'\}$, with $\mathrm{d}\xi$ and $\mathrm{d}\xi'$ times the counting measure:
these are $\sigma$-finite with separable $L^2$ spaces, $L^2(\Xi;V)$ is the $L^2$ space of the first one, the matrix entries form the scalar kernel, and the sum over the discrete variables is $\operatorname{tr}$.
\end{itemize}
The measure spaces to which (In9) is applied are those of the spaces $L^2(\mathrm{Sx};\mathrm{Vf})$ of 3.1, with the measure $\mathrm{dmi}$, and, through Lemma~6.1, those of finite orthogonal sums of these spaces;
by 3.1\,(s0) they are $\sigma$-finite and have separable $L^2$ spaces.
Besides (K3), three classical facts are needed \cite[Theorems~VI.22 and~VI.23]{ReedSimon1980}:
the product of a Hilbert--Schmidt operator and a bounded operator is a Hilbert--Schmidt operator;
the product of two Hilbert--Schmidt operators is of trace class;
and every Hilbert--Schmidt operator between two such $L^2$ spaces has a kernel in $L^2$.

\subsection*{1. The group of the walls; the cells of a torus (clause (a) of Proposition~39)}

\paragraph{1.1. Walls are mirrors.}
Since $\tfrac12\mathbb{Z}$ is an additive group,
the position $o_\mu+j\,\mathrm{hw}$ belongs to $\tfrac12\mathbb{Z}=\operatorname{Adm}(y_\mu)$ (Section~I.4.8) for every $j\in\mathbb{Z}$.
For $\mathrm{po}\in\tfrac12\mathbb{Z}$ the reflection
\[
\theta_{y_\mu,\mathrm{po}}(y)=y-2\,(y_\mu-\mathrm{po})\,e_\mu
\]
changes the coordinate $y_\mu$ into $2\,\mathrm{po}-y_\mu$ and leaves the other two coordinates unchanged.

\paragraph{1.2. Two adjacent walls.}
The composition $\theta_{y_\mu,\mathrm{po}+\mathrm{hw}}\,\theta_{y_\mu,\mathrm{po}}$ changes $y_\mu$ into
\[
2\,(\mathrm{po}+\mathrm{hw})-(2\,\mathrm{po}-y_\mu)=y_\mu+2\,\mathrm{hw}.
\]
It is therefore the translation by $2\,\mathrm{hw}\cdot e_\mu$,
and $2\,\mathrm{hw}$ is a positive integer.

\paragraph{1.3. One direction.}
Let $\mathrm{Di}_\mu$ be the group generated by the reflections
\[
\mathrm{rf}_j:=\theta_{y_\mu,\,o_\mu+j\,\mathrm{hw}},\qquad j\in\mathbb{Z}.
\]
Its elements act on the coordinate $y_\mu$ alone.
The product $\mathrm{rf}_j\,\mathrm{rf}_i$ is the translation by $2\,(j-i)\,\mathrm{hw}\cdot e_\mu$,
and $\tau_{2s\cdot\mathrm{hw}\cdot e_\mu}\,\mathrm{rf}_j=\mathrm{rf}_{j+s}$,
because both maps change $y_\mu$ into $2\left(o_\mu+(j+s)\,\mathrm{hw}\right)-y_\mu$.
Consequently
\[
\mathrm{Di}_\mu=\left\{\tau_{2s\cdot\mathrm{hw}\cdot e_\mu}:\ s\in\mathbb{Z}\right\}\cup\left\{\mathrm{rf}_j:\ j\in\mathbb{Z}\right\}.
\]
This group is generated by the two involutions $\mathrm{rf}_0$ and $\mathrm{rf}_1$, whose product has infinite order;
it is the infinite dihedral group.
For $j\in\mathbb{Z}$ put
\[
\mathrm{In}_j:=\left(o_\mu+j\,\mathrm{hw},\ o_\mu+(j+1)\,\mathrm{hw}\right).
\]
The translation $\tau_{2s\cdot\mathrm{hw}\cdot e_\mu}$ maps the set $\{y_\mu\in\mathrm{In}_j\}$ onto $\{y_\mu\in\mathrm{In}_{j+2s}\}$,
and the reflection $\mathrm{rf}_k$ maps it onto $\{y_\mu\in\mathrm{In}_{2k-j-1}\}$.
Hence, for all $i,j\in\mathbb{Z}$, the group $\mathrm{Di}_\mu$ contains exactly one element $\mathrm{di}^{\mu}_{i,j}$
that maps $\{y_\mu\in\mathrm{In}_j\}$ onto $\{y_\mu\in\mathrm{In}_i\}$, namely
\begin{equation}
\tag{1.1}
\mathrm{di}^{\mu}_{i,j}=
\begin{cases}
\tau_{(i-j)\,\mathrm{hw}\cdot e_\mu} & \text{if } i-j \text{ is even},\\
\mathrm{rf}_{(i+j+1)/2}=\tau_{(i-j-1)\,\mathrm{hw}\cdot e_\mu}\,\mathrm{rf}_{j+1} & \text{if } i-j \text{ is odd}.
\end{cases}
\end{equation}
Indeed, a translation changes the index of the interval by an even number,
a reflection $\mathrm{rf}_k$ changes it by the odd number $2k-2j-1$,
and in each of the two cases the number $s$, respectively $k$, is determined by $i-j$.
In the first case of (1.1) the element $\mathrm{di}^{\mu}_{i,j}$ is an even element of $W_{\mathrm{af}}$, in the second case an odd one.

\paragraph{1.4. Three directions.}
Let $\mu\neq\nu$.
Elements of $\mathrm{Di}_\mu$ and of $\mathrm{Di}_\nu$ act on different coordinates;
so they commute,
and an element of $\mathrm{Di}_\mu$ that is a product of elements of the two other groups fixes every coordinate.
It follows that $W(\mathcal{Q})=\mathrm{Di}_1\,\mathrm{Di}_2\,\mathrm{Di}_3$ is the direct product of three infinite dihedral groups.
With the intervals $\mathrm{In}_j$ of 1.3 formed with $o_\mu$ in the direction $\mu$, one has
\[
Q^{\circ}_{\mathfrak{c}}=\left\{y:\ y_\mu\in\mathrm{In}_{\mathfrak{c}_\mu},\ \mu=1,2,3\right\}.
\]
By 1.3 the elements of $W(\mathcal{Q})$ that map $Q^{\circ}_{\mathfrak{c}}$ onto $Q^{\circ}_{\mathfrak{c}'}$
are the products $g_{\mathrm{w},1}\,g_{\mathrm{w},2}\,g_{\mathrm{w},3}$ with $g_{\mathrm{w},\mu}\in\mathrm{Di}_\mu$
and $g_{\mathrm{w},\mu}\{y_\mu\in\mathrm{In}_{\mathfrak{c}_\mu}\}=\{y_\mu\in\mathrm{In}_{\mathfrak{c}'_\mu}\}$ for $\mu=1,2,3$,
and there is exactly one such product:
\begin{equation}
\tag{1.2}
\mathrm{di}^{3}_{\mathfrak{c}'_3,\mathfrak{c}_3}\,\mathrm{di}^{2}_{\mathfrak{c}'_2,\mathfrak{c}_2}\,\mathrm{di}^{1}_{\mathfrak{c}'_1,\mathfrak{c}_1}.
\end{equation}
Thus $W(\mathcal{Q})$ acts simply transitively on the cells.
The translations that belong to $W(\mathcal{Q})$ are the $\tau_v$ with $v\in(2\,\mathrm{hw}\,\mathbb{Z})^3$.

\paragraph{1.5. Tori.}
Let $C(n)$ carry $\mathcal{Q}$;
then $\mathrm{sv}_\mu:=n_\mu/\mathrm{hw}$ is an even positive integer.
With $s=\mathrm{sv}_\mu/2$ one has $\tau_{n_\mu e_\mu}=\tau_{2s\cdot\mathrm{hw}\cdot e_\mu}\in\mathrm{Di}_\mu$.
Hence $\Gamma=\sum_\mu n_\mu\mathbb{Z}e_\mu$ is contained in $(2\,\mathrm{hw}\,\mathbb{Z})^3$;
in other words, $\{\tau_v:\ v\in\Gamma\}$ is a subgroup of the group of the translations of $W(\mathcal{Q})$.
The linear part $r$ of an element $g_{\mathrm{w}}$ of $W(\mathcal{Q})$ is a diagonal matrix with the entries $\pm1$;
one has $g_{\mathrm{w}}\,\tau_v\,g_{\mathrm{w}}^{-1}=\tau_{rv}$ and $r\Gamma=\Gamma$,
so that this subgroup is normal.
For $v\in\mathbb{Z}^3$ with $v_\mu\in\mathrm{hw}\,\mathbb{Z}$ for $\mu=1,2,3$ put
\[
\mathrm{sh}(v):=\left(\frac{v_\mu}{\mathrm{hw}}\right)_{\mu=1,2,3}\in\mathbb{Z}^3 .
\]
Then $\tau_vQ^{\circ}_{\mathfrak{c}}=Q^{\circ}_{\mathfrak{c}+\mathrm{sh}(v)}$,
and $\mathrm{sh}(\Gamma)=\prod_\mu\mathrm{sv}_\mu\mathbb{Z}$.
The cells of $C(n)$ are the orbits of the cells under $\Gamma$,
that is, the classes of $\mathbb{Z}^3$ modulo $\prod_\mu\mathrm{sv}_\mu\mathbb{Z}$.
Their number is
\[
\mathrm{nc}(n)=\mathrm{sv}_1\,\mathrm{sv}_2\,\mathrm{sv}_3=\prod_\mu\frac{n_\mu}{\mathrm{hw}},
\]
a product of three even integers.
The quotient $W(\mathcal{Q})/\Gamma$ acts on these cells, and by 1.4 it acts transitively.
Suppose that $g_{\mathrm{w}}Q^{\circ}_{\mathfrak{c}}=\tau_vQ^{\circ}_{\mathfrak{c}}$ with $v\in\Gamma$.
Then $\tau_{-v}\,g_{\mathrm{w}}$ fixes a cell, so that $g_{\mathrm{w}}=\tau_v$ by 1.4.
The action of the quotient is therefore free.
\qed

\subsection*{2. The brick}

\paragraph{2.4. The brick.}
Let $C(n)$ carry $\mathcal{Q}$.
The brick is the set
\[
\mathrm{Kb}:=\left\{\mathfrak{c}\in\mathbb{Z}^3:\ 0\le\mathfrak{c}_\mu<\mathrm{sv}_\mu,\ \mu=1,2,3\right\};
\]
by 1.5 it is a set of representatives of the cells of $C(n)$.

\subsection*{3. The slice space; the operator $\mathrm{Kl}$ of a link plane}

In this section the integers $n_2,n_3\ge4$ are fixed.
One has $\#|\Lambda_s|=\#|\mathrm{Pl}_s|=v_s$ and $\#|\mathrm{Lk}_s|=2v_s$.
The Euclidean norms of $\mathbb{R}^{\mathrm{Lk}_s}$, $\mathbb{R}^{\mathrm{Pl}_s}$ and $\mathbb{R}^{\Lambda_s}$,
as well as the norm of a finite-dimensional Hilbert space, are all denoted by $|\cdot|$.

\paragraph{3.1. The slice space.}
Let
\[
\mathrm{Sxc}:=\mathbb{R}^{\mathrm{Lk}_s}\times\mathbb{Z}^{\mathrm{Pl}_s},
\]
whose points are written $\mathrm{mi}=(A,\mathfrak{z})$,
and equip it with the measure $\mathrm{d}\mathrm{mi}$, the product of the Lebesgue measure and the counting measure.
Put
\[
\mathrm{Bm}(\mathrm{mi}):=d_sA+r_q\mathfrak{z}\in\mathbb{R}^{\mathrm{Pl}_s},
\qquad
\mathrm{mi}+[m]:=\left(A+r_qm,\ \mathfrak{z}-d_sm\right)\quad\text{for } m\in\mathbb{Z}^{\mathrm{Lk}_s},
\]
and
\[
\mathrm{Sx}:=[0,r_q)^{\mathrm{Lk}_s}\times\mathbb{Z}^{\mathrm{Pl}_s}.
\]
A function $f$ on $\mathrm{Sxc}$ is called \emph{invariant} if $f(\mathrm{mi}+[m])=f(\mathrm{mi})$ for every $m\in\mathbb{Z}^{\mathrm{Lk}_s}$.
The function $\mathrm{Bm}$ is invariant.
\begin{itemize}
\item[(s1)]
Every point of $\mathrm{Sxc}$ can be written as $\mathrm{mi}_0+[m]$ with exactly one pair $(\mathrm{mi}_0,m)\in\mathrm{Sx}\times\mathbb{Z}^{\mathrm{Lk}_s}$;
the component $m_b$ is the integer part of $A_b/r_q$.
Consequently every function on $\mathrm{Sx}$ has exactly one invariant extension to $\mathrm{Sxc}$,
and functions on $\mathrm{Sx}$ are identified with invariant functions on $\mathrm{Sxc}$.
\item[(s2)]
Let $f$ be an invariant Borel function such that either $f\ge0$,
or $f$ is complex valued with $\int_{\mathrm{Sx}}\mathrm{d}\mathrm{mi}\,|f(\mathrm{mi})|<\infty$,
and let $y\in\mathbb{R}^{\mathrm{Lk}_s}$.
Then
\[
\int_{\mathrm{Sx}}\mathrm{d}\mathrm{mi}\,f(A-y,\mathfrak{z})=\int_{\mathrm{Sx}}\mathrm{d}\mathrm{mi}\,f(A,\mathfrak{z}).
\]
\end{itemize}
\begin{proof}[Proof of (s2)]
Let $\mathrm{bx}:=[0,r_q)^{\mathrm{Lk}_s}$.
The translates $\mathrm{bx}+r_qm$, $m\in\mathbb{Z}^{\mathrm{Lk}_s}$, tile $\mathbb{R}^{\mathrm{Lk}_s}$.
Therefore
\begin{align*}
\int_{\mathrm{Sx}}\mathrm{d}\mathrm{mi}\,f(A-y,\mathfrak{z})
&=\sum_{\mathfrak{z}}\int_{\mathrm{bx}-y}\mathrm{d}A\,f(A,\mathfrak{z})
=\sum_{m}\sum_{\mathfrak{z}}\int_{(\mathrm{bx}-y)\cap(\mathrm{bx}+r_qm)}\mathrm{d}A\,f(A,\mathfrak{z})\\
&=\sum_{m}\sum_{\mathfrak{z}}\int_{(\mathrm{bx}-y-r_qm)\cap\mathrm{bx}}\mathrm{d}A\,f(A+r_qm,\mathfrak{z}).
\end{align*}
By the invariance, $f(A+r_qm,\mathfrak{z})=f(A,\mathfrak{z}+d_sm)$;
the map $\mathfrak{z}\mapsto\mathfrak{z}+d_sm$ is a bijection of $\mathbb{Z}^{\mathrm{Pl}_s}$;
and the sets $\mathrm{bx}-y-r_qm$, $m\in\mathbb{Z}^{\mathrm{Lk}_s}$, again tile $\mathbb{R}^{\mathrm{Lk}_s}$.
So the last double sum equals $\sum_{\mathfrak{z}}\int_{\mathrm{bx}}\mathrm{d}A\,f(A,\mathfrak{z})$.
\end{proof}
\begin{itemize}
\item[(s3)]
For every $\varepsilon>0$,
\[
\int_{\mathrm{Sx}}\mathrm{d}\mathrm{mi}\,\exp\left(-\varepsilon\,|\mathrm{Bm}(\mathrm{mi})|^2\right)\le\left(r_q^2\,\mathrm{thq}(\varepsilon)\right)^{v_s}.
\]
\end{itemize}
\begin{proof}[Proof of (s3)]
For fixed $A$ the sum over $\mathfrak{z}$ is the product over $p'\in\mathrm{Pl}_s$ of the sums over the values $\mathfrak{z}(p')$,
and
\[
\sum_{\mathfrak{z}(p')\in\mathbb{Z}}\exp\left(-\varepsilon\left((d_sA)(p')+r_q\,\mathfrak{z}(p')\right)^2\right)\le\mathrm{thq}(\varepsilon).
\]
The box $\mathrm{bx}$ has the volume $r_q^{2v_s}$.
\end{proof}
For a finite-dimensional Hilbert space $\mathrm{Vf}$ the space $L^2(\mathrm{Sx};\mathrm{Vf})$ is formed with the measure $\mathrm{d}\mathrm{mi}$;
its scalar product $\langle\cdot,\cdot\rangle$ is antilinear in the first argument.
The property (s0) below is the $\sigma$-finiteness of $\mathrm{d}\mathrm{mi}$ on $\mathrm{Sx}$ together with the separability of these spaces.
\begin{itemize}
\item[(s0)]
$\mathrm{Sx}$ is the union of the countably many boxes $[0,r_q)^{\mathrm{Lk}_s}\times\{\mathfrak{z}\}$, $\mathfrak{z}\in\mathbb{Z}^{\mathrm{Pl}_s}$,
each of finite measure $r_q^{2v_s}$;
so $\mathrm{d}\mathrm{mi}$ is $\sigma$-finite on $\mathrm{Sx}$.
The space $L^2(\mathrm{Sx};\mathrm{Vf})$ is the orthogonal sum of the countably many spaces $L^2([0,r_q)^{\mathrm{Lk}_s};\mathrm{Vf})$,
each of which is separable, being formed with the Lebesgue measure on a box of $\mathbb{R}^{\mathrm{Lk}_s}$ and with values in the finite-dimensional space $\mathrm{Vf}$.
Hence $L^2(\mathrm{Sx};\mathrm{Vf})$ is separable, and so is every finite orthogonal sum of such spaces.
\end{itemize}

\paragraph{3.2. Operators on $L^2(\mathrm{Sx};\mathrm{Vf})$.}
For a real number $\mathrm{va}>0$ let
\[
\mathrm{gd}_{\mathrm{va}}(y):=(2\pi\,\mathrm{va})^{-v_s}\exp\left(-\frac{|y|^2}{2\,\mathrm{va}}\right),\qquad y\in\mathbb{R}^{\mathrm{Lk}_s},
\]
be the centred Gaussian density of the variance $\mathrm{va}$ on $\mathbb{R}^{\mathrm{Lk}_s}$, a space of dimension $2v_s$.
It satisfies $\mathrm{gd}_{\mathrm{va}}\ge0$, $\int\mathrm{gd}_{\mathrm{va}}=1$, $\mathrm{gd}_{\mathrm{va}}(-y)=\mathrm{gd}_{\mathrm{va}}(y)$
and $\mathrm{gd}_{\mathrm{va}}*\mathrm{gd}_{\mathrm{va}'}=\mathrm{gd}_{\mathrm{va}+\mathrm{va}'}$.
For an invariant function $f$ put
\[
(\mathrm{sf}_yf)(A,\mathfrak{z}):=f(A-y,\mathfrak{z})\quad(y\in\mathbb{R}^{\mathrm{Lk}_s}),
\qquad
(\mathrm{Hg}_{\mathrm{va}}f)(A,\mathfrak{z}):=\int_{\mathbb{R}^{\mathrm{Lk}_s}}\mathrm{d}y\,\mathrm{gd}_{\mathrm{va}}(y)\cdot f(A-y,\mathfrak{z}).
\]
For $\mathrm{Vf}=\mathrm{Fk}$, and with the representation $\mathrm{Wq}$ of (In2), put
\[
(\mathrm{Ug}(a)f)(A,\mathfrak{z}):=\mathrm{Wq}(a)\cdot f(A-d_sa,\mathfrak{z})\quad(a\in\mathbb{R}^{\Lambda_s}),
\qquad
(P_Gf)(A,\mathfrak{z}):=r_q^{-v_s}\int_{\mathrm{box}_0}\mathrm{d}a\,\mathrm{Wq}(a)\cdot f(A-d_sa,\mathfrak{z}),
\]
and
\[
\mathrm{Kl}:=(2\pi r_q)^{v_s}\cdot\mathrm{Hg}_1\,P_G,
\qquad
\mathrm{Kl}^{\frac12}:=(2\pi r_q)^{\frac{v_s}{2}}\cdot\mathrm{Hg}_{1/2}\,P_G .
\]
For $\mathrm{mi}'=(A',\mathfrak{z}')$ and $\mathrm{mi}=(A,\mathfrak{z})$ in $\mathrm{Sx}$ define
\[
\mathrm{Kl}(\mathrm{mi}',\mathrm{mi}):=\sum_{m\in\mathbb{Z}^{\mathrm{Lk}_s}}\mathbf{1}\left[\mathfrak{z}'=\mathfrak{z}+d_sm\right]\int_{\mathrm{box}_0}\mathrm{d}a\,\exp\left(-\tfrac12\left|A'-A-d_sa+r_qm\right|^2\right)\mathrm{Wq}(a)\ \in\operatorname{End}\mathrm{Fk},
\]
and let $\mathrm{kl}^+(\mathrm{mi}',\mathrm{mi})$ denote the same expression with $1$ in the place of $\mathrm{Wq}(a)$.
Then, for all $\mathrm{mi}'$ and $\mathrm{mi}$,
\[
\mathrm{kl}^+(\mathrm{mi}',\mathrm{mi})\le r_q^{v_s}\cdot\mathrm{thq}(1/2)^{2v_s}.
\]
\begin{proof}
Omit the indicator and interchange the sum and the integral, by Tonelli's theorem.
For fixed $a$ the sum over $m$ is the product over $b\in\mathrm{Lk}_s$ of the sums
\[
\sum_{m_b\in\mathbb{Z}}\exp\left(-\tfrac12\left((A'-A-d_sa)(b)+r_qm_b\right)^2\right)\le\mathrm{thq}(1/2),
\]
where the bound is (In7)\,(vii) at $\varepsilon=1/2$.
There are $\#|\mathrm{Lk}_s|=2v_s$ factors, and $\mathrm{box}_0$ has the volume $r_q^{v_s}$.
\end{proof}
Since $\mathrm{Wq}(a)$ is unitary,
the sum and the integral that define $\mathrm{Kl}(\mathrm{mi}',\mathrm{mi})$ converge absolutely at every point $(\mathrm{mi}',\mathrm{mi})$.
Finally, $P_G^{+}$ and $\mathrm{Kl}^{+}:=(2\pi r_q)^{v_s}\cdot\mathrm{Hg}_1\,P_G^{+}$ denote the operators $P_G$ and $\mathrm{Kl}$
formed for $\mathrm{Vf}=\mathbb{C}$ with the trivial representation in the place of $\mathrm{Wq}$.

\begin{named}{Lemma 3.3}
\begin{itemize}
\item[(a)]
$\mathrm{sf}_y$ maps invariant functions to invariant functions and is isometric on $L^2(\mathrm{Sx};\mathrm{Vf})$;
moreover $\mathrm{sf}_y\,\mathrm{sf}_{y'}=\mathrm{sf}_{y+y'}$ and $\langle f',\mathrm{sf}_yf\rangle=\langle\mathrm{sf}_{-y}f',f\rangle$.
\item[(b)]
Let $\mathrm{va},\mathrm{va}'>0$ and $f\in L^2(\mathrm{Sx};\mathrm{Vf})$.
The integral that defines $(\mathrm{Hg}_{\mathrm{va}}f)(\mathrm{mi})$ converges absolutely for almost every $\mathrm{mi}$;
$\|\mathrm{Hg}_{\mathrm{va}}f\|\le\|f\|$;
$\mathrm{Hg}_{\mathrm{va}}^*=\mathrm{Hg}_{\mathrm{va}}$;
and $\mathrm{Hg}_{\mathrm{va}}\,\mathrm{Hg}_{\mathrm{va}'}=\mathrm{Hg}_{\mathrm{va}+\mathrm{va}'}$.
\item[(c)]
For $f\in L^2(\mathrm{Sx};\mathrm{Fk})$ the function $a\mapsto\mathrm{Wq}(a)\cdot f(A-d_sa,\mathfrak{z})$ is $r_q\mathbb{Z}^{\Lambda_s}$-periodic;
the integral that defines $P_Gf$ converges absolutely for almost every $\mathrm{mi}$;
$P_G$ is an orthogonal projection;
and $P_G\,\mathrm{Hg}_{\mathrm{va}}=\mathrm{Hg}_{\mathrm{va}}\,P_G$.
\item[(d)]
$\mathrm{Kl}^{\frac12}$ is bounded, with $\|\mathrm{Kl}^{\frac12}\|\le(2\pi r_q)^{\frac{v_s}{2}}$, it is self-adjoint, and $\left(\mathrm{Kl}^{\frac12}\right)^2=\mathrm{Kl}$.
Consequently $\mathrm{Kl}=\left(\mathrm{Kl}^{\frac12}\right)^*\mathrm{Kl}^{\frac12}\ge0$.
\item[(e)]
For $f\in L^2(\mathrm{Sx};\mathrm{Fk})$ and almost every $\mathrm{mi}'\in\mathrm{Sx}$,
\[
\int_{\mathrm{Sx}}\mathrm{d}\mathrm{mi}\,\mathrm{kl}^+(\mathrm{mi}',\mathrm{mi})\,|f(\mathrm{mi})|=\left(\mathrm{Kl}^+|f|\right)(\mathrm{mi}')<\infty
\quad\text{and}\quad
(\mathrm{Kl}f)(\mathrm{mi}')=\int_{\mathrm{Sx}}\mathrm{d}\mathrm{mi}\,\mathrm{Kl}(\mathrm{mi}',\mathrm{mi})\,f(\mathrm{mi}).
\]
Moreover $\|\mathrm{Kl}(\mathrm{mi}',\mathrm{mi})\|\le\mathrm{kl}^+(\mathrm{mi}',\mathrm{mi})$.
\end{itemize}
\end{named}

\begin{proof}
(a)
One has $(A-y,\mathfrak{z})+[m]=(A+r_qm-y,\ \mathfrak{z}-d_sm)$;
hence $\mathrm{sf}_y$ preserves invariance.
The isometry is (s2) applied to $|f|^2$.
Next, $(\mathrm{sf}_y\,\mathrm{sf}_{y'}f)(A,\mathfrak{z})=f(A-y-y',\mathfrak{z})$.
For the last identity consider the function
\[
\mathrm{fq}(A,\mathfrak{z}):=\left\langle f'(A+y,\mathfrak{z}),f(A,\mathfrak{z})\right\rangle_{\mathrm{Vf}} .
\]
It is invariant and, by the Cauchy--Schwarz inequality, integrable,
and $\mathrm{fq}(A-y,\mathfrak{z})=\langle f'(A,\mathfrak{z}),f(A-y,\mathfrak{z})\rangle_{\mathrm{Vf}}$;
so the identity follows from (s2).

(b)
By Tonelli's theorem and (a),
\[
\int_{\mathrm{Sx}}\mathrm{d}\mathrm{mi}\int\mathrm{d}y\,\mathrm{gd}_{\mathrm{va}}(y)\,|f(A-y,\mathfrak{z})|^2
=\int\mathrm{d}y\,\mathrm{gd}_{\mathrm{va}}(y)\,\|\mathrm{sf}_yf\|^2=\|f\|^2 .
\]
Hence, for almost every $\mathrm{mi}$, the function $y\mapsto|f(A-y,\mathfrak{z})|^2$ is integrable with respect to the probability measure $\mathrm{gd}_{\mathrm{va}}(y)\,\mathrm{d}y$.
For such $\mathrm{mi}$ the Cauchy--Schwarz inequality for this measure shows that the integral of $|f(A-y,\mathfrak{z})|$ is finite and that
\[
\left|(\mathrm{Hg}_{\mathrm{va}}f)(\mathrm{mi})\right|^2\le\int\mathrm{d}y\,\mathrm{gd}_{\mathrm{va}}(y)\,|f(A-y,\mathfrak{z})|^2 .
\]
Integrating over $\mathrm{Sx}$ one obtains $\|\mathrm{Hg}_{\mathrm{va}}f\|\le\|f\|$;
the function $\mathrm{Hg}_{\mathrm{va}}f$ is invariant.
By Fubini's theorem and (a), and since $\mathrm{gd}_{\mathrm{va}}$ is even,
\[
\langle f',\mathrm{Hg}_{\mathrm{va}}f\rangle
=\int\mathrm{d}y\,\mathrm{gd}_{\mathrm{va}}(y)\,\langle f',\mathrm{sf}_yf\rangle
=\int\mathrm{d}y\,\mathrm{gd}_{\mathrm{va}}(y)\,\langle\mathrm{sf}_{-y}f',f\rangle
=\langle\mathrm{Hg}_{\mathrm{va}}f',f\rangle .
\]
Finally,
\[
\left(\mathrm{Hg}_{\mathrm{va}}\,\mathrm{Hg}_{\mathrm{va}'}f\right)(A,\mathfrak{z})
=\int\mathrm{d}y\int\mathrm{d}y'\,\mathrm{gd}_{\mathrm{va}}(y)\,\mathrm{gd}_{\mathrm{va}'}(y')\,f(A-y-y',\mathfrak{z})
=\int\mathrm{d}y''\,\left(\mathrm{gd}_{\mathrm{va}}*\mathrm{gd}_{\mathrm{va}'}\right)(y'')\,f(A-y'',\mathfrak{z})
\]
by Fubini's theorem;
it applies because the same integral with $|f|$ in the place of $f$ equals $\left(\mathrm{Hg}_{\mathrm{va}+\mathrm{va}'}|f|\right)(\mathrm{mi})$,
which is finite for almost every $\mathrm{mi}$.

(c)
Let $m^0\in\mathbb{Z}^{\Lambda_s}$.
By (In2), $\mathrm{Wq}(a+r_qm^0)=\mathrm{Wq}(a)$.
Since $d_sd_sm^0=0$,
\[
\left(A-d_sa-r_q\,d_sm^0,\ \mathfrak{z}\right)=\left(A-d_sa,\ \mathfrak{z}\right)+\left[-d_sm^0\right],
\]
and $f$ takes the same value at these two points.
This is the periodicity.
As a consequence, in the definition of $P_G$ the box $\mathrm{box}_0$ may be replaced by $\mathrm{box}_0+a_0$ with $a_0\in\mathbb{R}^{\Lambda_s}$, and also by $-\mathrm{box}_0$.
By Tonelli's theorem, the unitarity of $\mathrm{Wq}(a)$ and (s2),
\[
\int_{\mathrm{Sx}}\mathrm{d}\mathrm{mi}\ r_q^{-v_s}\int_{\mathrm{box}_0}\mathrm{d}a\,\left|\mathrm{Wq}(a)\cdot f(A-d_sa,\mathfrak{z})\right|^2=\|f\|^2 .
\]
So, for almost every $\mathrm{mi}$, the function $a\mapsto|f(A-d_sa,\mathfrak{z})|^2$ is integrable on $\mathrm{box}_0$.
For such $\mathrm{mi}$ the Cauchy--Schwarz inequality for the probability measure $r_q^{-v_s}\,\mathrm{d}a$ on $\mathrm{box}_0$ shows that the integral defining $(P_Gf)(\mathrm{mi})$ converges absolutely, and
\[
\left|(P_Gf)(\mathrm{mi})\right|^2\le r_q^{-v_s}\int_{\mathrm{box}_0}\mathrm{d}a\,|f(A-d_sa,\mathfrak{z})|^2 .
\]
Integration over $\mathrm{Sx}$ gives $\|P_Gf\|\le\|f\|$.

The \emph{exceptional set} of $f$ is the set of the points $\mathrm{mi}=(A,\mathfrak{z})$ of $\mathrm{Sxc}$ with
$\int_{\mathrm{box}_0}\mathrm{d}a\,|f(A-d_sa,\mathfrak{z})|^2=\infty$.
By Tonelli's theorem it is a Borel set.
It is invariant under $\mathrm{mi}\mapsto\mathrm{mi}+[m]$, because $f$ is invariant,
and its intersection with $\mathrm{Sx}$ is a null set, as has just been shown;
by (s1) it is therefore a null set of $\mathrm{Sxc}$.
It is also invariant under $A\mapsto A-d_sa_0$, $a_0\in\mathbb{R}^{\Lambda_s}$,
since by the periodicity the integral over $\mathrm{box}_0+a_0$ equals the integral over $\mathrm{box}_0$.
Outside the exceptional set, $(P_Gf)(\mathrm{mi})$ is the absolutely convergent integral;
on it one puts $(P_Gf)(\mathrm{mi}):=0$.
By (a), $P_Gf$ is invariant:
\[
(P_Gf)(\mathrm{mi}+[m])=r_q^{-v_s}\int_{\mathrm{box}_0}\mathrm{d}a\,\mathrm{Wq}(a)\cdot f\left((A-d_sa,\mathfrak{z})+[m]\right)=(P_Gf)(\mathrm{mi}).
\]
Now let $\mathrm{mi}$ lie outside the exceptional set of $f$.
Then all the points $(A-d_sa,\mathfrak{z})$ lie outside it, and
\[
\left|(P_Gf)(A-d_sa,\mathfrak{z})\right|^2\le r_q^{-v_s}\int_{\mathrm{box}_0}\mathrm{d}a''\,|f(A-d_sa'',\mathfrak{z})|^2\qquad\text{for all } a,
\]
so that $\mathrm{mi}$ lies outside the exceptional set of $P_Gf$ as well.
Moreover
\begin{align*}
\int_{\mathrm{box}_0}\mathrm{d}a\int_{\mathrm{box}_0}\mathrm{d}a'\,\left|f\left(A-d_s(a+a'),\mathfrak{z}\right)\right|
&=\int_{\mathrm{box}_0}\mathrm{d}a\int_{\mathrm{box}_0+a}\mathrm{d}a''\,|f(A-d_sa'',\mathfrak{z})|\\
&=r_q^{v_s}\int_{\mathrm{box}_0}\mathrm{d}a''\,|f(A-d_sa'',\mathfrak{z})|<\infty .
\end{align*}
Hence Fubini's theorem and the periodicity give
\begin{align*}
(P_G\,P_Gf)(A,\mathfrak{z})
&=r_q^{-2v_s}\int_{\mathrm{box}_0}\mathrm{d}a\int_{\mathrm{box}_0}\mathrm{d}a'\,\mathrm{Wq}(a+a')\cdot f\left(A-d_s(a+a'),\mathfrak{z}\right)\\
&=r_q^{-2v_s}\int_{\mathrm{box}_0}\mathrm{d}a\int_{\mathrm{box}_0+a}\mathrm{d}a''\,\mathrm{Wq}(a'')\cdot f(A-d_sa'',\mathfrak{z})
=(P_Gf)(A,\mathfrak{z}).
\end{align*}
By Fubini's theorem, (a) and $\mathrm{Wq}(a)^*=\mathrm{Wq}(-a)$,
\begin{align*}
\langle f',P_Gf\rangle
&=r_q^{-v_s}\int_{\mathrm{box}_0}\mathrm{d}a\,\left\langle f',\mathrm{Wq}(a)\,\mathrm{sf}_{d_sa}f\right\rangle
=r_q^{-v_s}\int_{\mathrm{box}_0}\mathrm{d}a\,\left\langle\mathrm{Wq}(-a)\,\mathrm{sf}_{-d_sa}f',f\right\rangle\\
&=r_q^{-v_s}\int_{-\mathrm{box}_0}\mathrm{d}a'\,\left\langle\mathrm{Wq}(a')\,\mathrm{sf}_{d_sa'}f',f\right\rangle
=\langle P_Gf',f\rangle .
\end{align*}
Thus $P_G=P_G^2=P_G^*$.
The operators $\mathrm{Hg}_{\mathrm{va}}$ and $P_G$ commute by Fubini's theorem,
because translations commute with each other and $\mathrm{Wq}(a)$ is a constant matrix.
The majorant for these applications of Fubini's theorem is obtained as follows.
By Tonelli's theorem the integral
\[
r_q^{-v_s}\int_{\mathrm{box}_0}\mathrm{d}a\int_{\mathbb{R}^{\mathrm{Lk}_s}}\mathrm{d}y\,\mathrm{gd}_{\mathrm{va}}(y)\,|f(A-y-d_sa,\mathfrak{z})|,
\]
taken in either order, equals $\left(\mathrm{Hg}_{\mathrm{va}}\,P_G^+|f|\right)(\mathrm{mi})$ for every $\mathrm{mi}$.
The function $\mathrm{Hg}_{\mathrm{va}}\,P_G^+|f|$ belongs to $L^2(\mathrm{Sx})$,
by (b) and by the bound $\|P_G^+|f|\|\le\|f\|$ just proved for $\mathrm{Vf}=\mathbb{C}$ and the trivial representation;
so it is finite for almost every $\mathrm{mi}$.

(d)
By (b) and (c), $\mathrm{Kl}^{\frac12}$ is a positive number times the product of two commuting bounded self-adjoint operators of norm at most $1$, and
\[
\left(\mathrm{Kl}^{\frac12}\right)^2=(2\pi r_q)^{v_s}\cdot\mathrm{Hg}_{1/2}^2\,P_G^2=(2\pi r_q)^{v_s}\cdot\mathrm{Hg}_1\,P_G .
\]

(e)
One has $(2\pi)^{v_s}\,\mathrm{gd}_1(y)=\exp\left(-\tfrac12|y|^2\right)$.
For almost every $\mathrm{mi}'=(A',\mathfrak{z}')$, by (b), (c) and Fubini's theorem,
\begin{align*}
(\mathrm{Kl}f)(\mathrm{mi}')
&=\int_{\mathrm{box}_0}\mathrm{d}a\int_{\mathbb{R}^{\mathrm{Lk}_s}}\mathrm{d}y\,\exp\left(-\tfrac12|y|^2\right)\mathrm{Wq}(a)\,f(A'-y-d_sa,\mathfrak{z}')\\
&=\int_{\mathrm{box}_0}\mathrm{d}a\int_{\mathbb{R}^{\mathrm{Lk}_s}}\mathrm{d}A''\,\exp\left(-\tfrac12\left|A'-A''-d_sa\right|^2\right)\mathrm{Wq}(a)\,f(A'',\mathfrak{z}'),
\end{align*}
where the second line results from the substitution $y=A'-d_sa-A''$.
Write $A''=A-r_qm$ with $A\in[0,r_q)^{\mathrm{Lk}_s}$ and $m\in\mathbb{Z}^{\mathrm{Lk}_s}$.
Since $(A-r_qm,\mathfrak{z}')=(A,\mathfrak{z}'-d_sm)+[-m]$, one has $f(A-r_qm,\mathfrak{z}')=f(A,\mathfrak{z}'-d_sm)$, and therefore
\begin{align*}
(\mathrm{Kl}f)(\mathrm{mi}')
&=\sum_{m}\int_{[0,r_q)^{\mathrm{Lk}_s}}\mathrm{d}A\int_{\mathrm{box}_0}\mathrm{d}a\,\exp\left(-\tfrac12\left|A'-A-d_sa+r_qm\right|^2\right)\mathrm{Wq}(a)\,f(A,\mathfrak{z}'-d_sm)\\
&=\sum_{\mathfrak{z}}\int_{[0,r_q)^{\mathrm{Lk}_s}}\mathrm{d}A\,\mathrm{Kl}\left(\mathrm{mi}',(A,\mathfrak{z})\right)f(A,\mathfrak{z}).
\end{align*}
\emph{The interchanges.}
Replace $f$ by $|f|$ and $\mathrm{Wq}(a)$ by $1$ in the two displays.
All integrands are then non-negative,
so that each step is justified by Tonelli's theorem, with values in $[0,\infty]$:
the interchange of the integrals over $a$ and $y$,
the substitution,
the decomposition $A''=A-r_qm$ together with the invariance of $|f|$,
and the interchange of the sum with the integrals.
This yields
\[
\left(\mathrm{Kl}^+|f|\right)(\mathrm{mi}')=\int_{\mathrm{Sx}}\mathrm{d}\mathrm{mi}\,\mathrm{kl}^+(\mathrm{mi}',\mathrm{mi})\,|f(\mathrm{mi})|
\qquad\text{for every }\mathrm{mi}'.
\]
By (b) and (c), applied for $\mathrm{Vf}=\mathbb{C}$ and the trivial representation, $\mathrm{Kl}^+|f|$ lies in $L^2(\mathrm{Sx})$;
hence it is finite almost everywhere.
This is the first assertion of (e).
At every point $\mathrm{mi}'$ where it is finite, Fubini's theorem gives the two displays for $f$ itself,
because $|\mathrm{Wq}(a)f|=|f|$.
Finally, $\|\mathrm{Kl}(\mathrm{mi}',\mathrm{mi})\|\le\mathrm{kl}^+(\mathrm{mi}',\mathrm{mi})$ because $\mathrm{Wq}(a)$ is unitary.
\end{proof}

\emph{The kernel at a link plane.}
Let $\mathrm{mi}=\mathrm{mi}_t$ and $\mathrm{mi}'=\mathrm{mi}_{t+1}$ be the slices below and above the link plane $t+\tfrac12$,
and let $(a,m)=(a_t,\mathfrak{z}^{\mathrm{el}}_t)$ be the step $t$.
Then $A'-A-d_sa+r_qm=\mathrm{El}_t$,
and $\mathbf{1}\left[\mathfrak{z}'=\mathfrak{z}+d_sm\right]$ is the indicator of the condition $d\mathfrak{z}=0$ on the cubes of the step $t$ (see (In7)\,(ii)).
Hence
\begin{equation}
\tag{3.1}
\mathrm{Kl}(\mathrm{mi}_{t+1},\mathrm{mi}_t)
=\sum_{\mathfrak{z}^{\mathrm{el}}_t\in\mathbb{Z}^{\mathrm{Lk}_s}}\int_{\mathrm{box}_0}\mathrm{d}a_t\,
\prod_{p'\in\mathrm{Pl}_s}\mathbf{1}\left[(d\mathfrak{z})(t,p')=0\right]\cdot
\prod_{b\in\mathrm{Lk}_s}\exp\left(-\tfrac12\,\mathrm{El}_t(b)^2\right)\cdot\mathrm{Wq}(a_t).
\end{equation}
In words: the kernel is the sum and integral, over the cells cut by the plane,
of their Maxwell weight, of the constraint, and of the phases $\operatorname{diag}(u_I(A))$ of the crossing expansion (In3).

\begin{named}{Lemma 3.4}[a Hilbert--Schmidt kernel times $\mathrm{Kl}$]
Let $\mathrm{Vf}$ be a finite-dimensional Hilbert space,
let $\mathrm{Ro}\colon L^2(\mathrm{Sx};\mathrm{Fk})\to L^2(\mathrm{Sx};\mathrm{Vf})$ be a Hilbert--Schmidt operator with the kernel $\mathrm{Ro}(\mathrm{mi}',\mathrm{mi}'')$,
and let $\mathrm{ro}^+\in L^2(\mathrm{Sx}\times\mathrm{Sx})$ satisfy $\|\mathrm{Ro}(\mathrm{mi}',\mathrm{mi}'')\|\le\mathrm{ro}^+(\mathrm{mi}',\mathrm{mi}'')$.
Then the function
\[
\mathrm{zo}^+(\mathrm{mi}',\mathrm{mi}):=\int_{\mathrm{Sx}}\mathrm{d}\mathrm{mi}''\,\mathrm{ro}^+(\mathrm{mi}',\mathrm{mi}'')\,\mathrm{kl}^+(\mathrm{mi}'',\mathrm{mi})
\]
is finite for almost every $(\mathrm{mi}',\mathrm{mi})$ and lies in $L^2(\mathrm{Sx}\times\mathrm{Sx})$;
it is the kernel of $\mathrm{ro}^+\,\mathrm{Kl}^+$, where $\mathrm{ro}^+$ also denotes the operator with the kernel $\mathrm{ro}^+$.
Moreover, $\mathrm{Ro}\,\mathrm{Kl}$ is the Hilbert--Schmidt operator with the kernel
\[
\mathrm{Zo}(\mathrm{mi}',\mathrm{mi}):=\int_{\mathrm{Sx}}\mathrm{d}\mathrm{mi}''\,\mathrm{Ro}(\mathrm{mi}',\mathrm{mi}'')\,\mathrm{Kl}(\mathrm{mi}'',\mathrm{mi}),
\qquad
\|\mathrm{Zo}(\mathrm{mi}',\mathrm{mi})\|\le\mathrm{zo}^+(\mathrm{mi}',\mathrm{mi}).
\]
\end{named}

\begin{proof}
As the product of a Hilbert--Schmidt operator and a bounded operator, $\mathrm{ro}^+\,\mathrm{Kl}^+$ is a Hilbert--Schmidt operator;
let $\mathrm{zo}'\in L^2(\mathrm{Sx}\times\mathrm{Sx})$ be its kernel.
For non-negative $f,f'\in L^2(\mathrm{Sx})$, Lemma~3.3\,(e) and Tonelli's theorem give
\[
\int\mathrm{d}\mathrm{mi}'\int\mathrm{d}\mathrm{mi}\,f'(\mathrm{mi}')\,\mathrm{zo}^+(\mathrm{mi}',\mathrm{mi})\,f(\mathrm{mi})
=\langle f',\mathrm{ro}^+\,\mathrm{Kl}^+f\rangle
=\int\mathrm{d}\mathrm{mi}'\int\mathrm{d}\mathrm{mi}\,f'(\mathrm{mi}')\,\mathrm{zo}'(\mathrm{mi}',\mathrm{mi})\,f(\mathrm{mi})<\infty .
\]
Taking for $f$ and $f'$ the indicator functions of sets of finite measure,
one obtains $\mathrm{zo}^+=\mathrm{zo}'$ almost everywhere, since $\mathrm{d}\mathrm{mi}$ is $\sigma$-finite by (s0).
Now let $f\in L^2(\mathrm{Sx};\mathrm{Fk})$.
For almost every $\mathrm{mi}'$,
\[
(\mathrm{Ro}\,\mathrm{Kl}f)(\mathrm{mi}')
=\int\mathrm{d}\mathrm{mi}''\,\mathrm{Ro}(\mathrm{mi}',\mathrm{mi}'')\left(\int\mathrm{d}\mathrm{mi}\,\mathrm{Kl}(\mathrm{mi}'',\mathrm{mi})\,f(\mathrm{mi})\right)
=\int\mathrm{d}\mathrm{mi}\,\mathrm{Zo}(\mathrm{mi}',\mathrm{mi})\,f(\mathrm{mi});
\]
here Fubini's theorem applies, because
\[
\int\mathrm{d}\mathrm{mi}''\int\mathrm{d}\mathrm{mi}\,\mathrm{ro}^+(\mathrm{mi}',\mathrm{mi}'')\,\mathrm{kl}^+(\mathrm{mi}'',\mathrm{mi})\,|f(\mathrm{mi})|=\left(\mathrm{ro}^+\,\mathrm{Kl}^+|f|\right)(\mathrm{mi}')<\infty
\]
for almost every $\mathrm{mi}'$.
The bound on $\mathrm{Zo}$ follows from $\|\mathrm{Ro}(\mathrm{mi}',\mathrm{mi}'')\|\le\mathrm{ro}^+(\mathrm{mi}',\mathrm{mi}'')$ and from the last assertion of Lemma~3.3\,(e).
\end{proof}

\subsection*{4. Fermions between two coordinate planes; the slice data}

\begin{named}{Definition 4.1}
Let $\mathrm{po}<\mathrm{po}'$ be elements of $\tfrac12\mathbb{Z}$,
let $G\in\mathrm{Ig}^{\mathrm{ev}}_{\mathrm{Xs}}[\mathrm{po},\mathrm{po}']$,
and let $\omega=(A,\mathfrak{z})\in\mathrm{Cfg}[\mathrm{po},\mathrm{po}']$.
Put
\[
\mathrm{Ms}[G](\omega)=\mathrm{Ms}_{\mathrm{po},\mathrm{po}'}[G](\omega)
:=\mathrm{cc}_{\mathrm{po}'}^{\frac12}\cdot\mathrm{Ml}_{\mathrm{po},\mathrm{po}'}\left[G(\omega)\exp\left(-S_{(\mathrm{po},\mathrm{po}')}(A)\right)\right]\cdot\mathrm{cc}_{\mathrm{po}}^{\frac12}
\colon\ \mathrm{Fk}(\mathrm{po})\to\mathrm{Fk}(\mathrm{po}').
\]
For a one-cochain that is defined on the links cut by the plane $\{y_1=\mathrm{po}\}$ put
\[
\mathrm{Wx}(\mathrm{po}):=\mathrm{Wq}(a_{\mathrm{po}-1/2})\ \text{on }\mathrm{Fk}\quad\text{if }\mathrm{po}\in\mathbb{Z}+\tfrac12,
\qquad
\mathrm{Wx}(\mathrm{po}):=1\ \text{on }\mathrm{Fk}'\quad\text{if }\mathrm{po}\in\mathbb{Z}.
\]
\end{named}
By (In2) and (In3), and because diagonal matrices commute,
\[
\exp\left(-\mathrm{Xc}_{\mathrm{po}}(A)\right)=\mathrm{Ins}_{\mathrm{po}}\left[\mathrm{cc}_{\mathrm{po}}^{\frac12}\,\mathrm{Wx}(\mathrm{po})\,\mathrm{cc}_{\mathrm{po}}^{\frac12}\right].
\]

\begin{named}{Lemma 4.2}
Let $G\in\mathrm{Ig}^{\mathrm{ev}}_{\mathrm{Xs}}[\mathrm{po},\mathrm{po}']$ have the constants $(\mathrm{Cq},\varepsilon_G)$.
\begin{itemize}
\item[(m1)]
There is a constant $C_{\mathrm{sl}}<\infty$, depending on $(n_2,n_3,\mathrm{po}'-\mathrm{po})$ and on $(\rho,M_{\mathrm{ch}})$, such that
\[
\|\mathrm{Ms}[G](\omega)\|\le C_{\mathrm{sl}}\cdot\mathrm{Cq}\cdot\exp\left(\left(\tfrac12-\varepsilon_G\right)\sum_{p\in\mathrm{Sb}^2(\mathrm{po},\mathrm{po}')}F^{\mathrm{u}}_p(\omega)^2\right)
\qquad\text{for all }\omega\in\mathrm{Cfg}[\mathrm{po},\mathrm{po}'].
\]
\item[(m2)]
The map $\omega\mapsto\mathrm{Ms}[G](\omega)$ is Borel, it is linear in $G$,
and $\mathrm{Ms}[G](\omega+[m])=\mathrm{Ms}[G](\omega)$ for all $m\in\mathbb{Z}^{\mathrm{Ln}[\mathrm{po},\mathrm{po}']}$.
\item[(m3)]
Let $\mathrm{po}'':=2\,\mathrm{po}'-\mathrm{po}$.
For $\tilde{\omega}\in\mathrm{Cfg}[\mathrm{po}',\mathrm{po}'']$,
\[
\mathrm{Ms}[\Theta_{\mathrm{po}'}G](\tilde{\omega})=\mathrm{Ms}[G](\theta_{\mathrm{po}'}^{*}\tilde{\omega})^{*}.
\]
\item[(m4)]
For $j\in\mathbb{Z}$ and $\omega'\in\mathrm{Cfg}[\mathrm{po}+j,\mathrm{po}'+j]$,
\[
\mathrm{Ms}[\mathrm{Tg}_{je_1}G](\omega')=\mathrm{Ms}[G](\tau_{je_1}^{*}\omega').
\]
\item[(m5)]
Let $n_1\ge4$, $T=C(n_1,n_2,n_3)$, $r\ge1$,
and let $\mathrm{po}_0<\mathrm{po}_1<\dots<\mathrm{po}_r=\mathrm{po}_0+n_1$ be elements of $\tfrac12\mathbb{Z}$.
Let $G_k\in\mathrm{Ig}^{\mathrm{ev}}_{\mathrm{Xs}}[\mathrm{po}_{k-1},\mathrm{po}_k]$ for $k=1,\dots,r$,
let $(A,\mathfrak{z})\in\mathrm{Cfg}(\Gamma)$,
and let $\omega_k$ be the restriction of the periodic extension of $(A,\mathfrak{z})$ to the cells of the closed slab $\{\mathrm{po}_{k-1}\le y_1\le\mathrm{po}_k\}$.
Then $\omega_k\in\mathrm{Cfg}[\mathrm{po}_{k-1},\mathrm{po}_k]$,
and, the values of the $G_k$ being read in $\mathrm{Gr}_T$ as in the rule (F5) of (In5),
\begin{multline*}
\int_T\exp\left(-S_T(A)\right)G_1(\omega_1)\cdots G_r(\omega_r)\\
=\operatorname{tr}_{\mathrm{Fk}(\mathrm{po}_0)}\left(\mathrm{Ms}[G_r](\omega_r)\,\mathrm{Wx}(\mathrm{po}_{r-1})\,\mathrm{Ms}[G_{r-1}](\omega_{r-1})\cdots\mathrm{Wx}(\mathrm{po}_1)\,\mathrm{Ms}[G_1](\omega_1)\,\mathrm{Wx}(\mathrm{po}_0)\right).
\end{multline*}
\item[(m6)]
Let $\mathrm{po}<\mathrm{po}'<\mathrm{po}''$,
let $G'\in\mathrm{Ig}^{\mathrm{ev}}_{\mathrm{Xs}}[\mathrm{po}',\mathrm{po}'']$ and $\omega\in\mathrm{Cfg}[\mathrm{po},\mathrm{po}'']$,
let $\omega_1$ and $\omega_2$ be the restrictions of $\omega$ to the two closed slabs $\{\mathrm{po}\le y_1\le\mathrm{po}'\}$ and $\{\mathrm{po}'\le y_1\le\mathrm{po}''\}$,
and put $(GG')(\omega):=G(\omega_1)\,G'(\omega_2)$.
Then
\[
\mathrm{Ms}_{\mathrm{po},\mathrm{po}''}[GG'](\omega)=\mathrm{Ms}[G'](\omega_2)\,\mathrm{Wx}(\mathrm{po}')\,\mathrm{Ms}[G](\omega_1).
\]
\end{itemize}
\end{named}

\begin{proof}
(m1)
Let $\mathrm{Nx}:=\#|X_{\mathrm{Xs}}((\mathrm{po},\mathrm{po}'))|$.
By (In4), $\exp\left(-S_{(\mathrm{po},\mathrm{po}')}(A)\right)$ is a polynomial in the $8\,\mathrm{Nx}$ generators
whose coefficients are polynomials in $M_{\mathrm{ch}}$, in the entries of the matrices $M_i(x)$ and in phases of modulus one;
these coefficients are bounded uniformly in $A$.
The elements $C^{(\mathrm{po})}_I$ and $\Theta_{\mathrm{po}'}\left(C^{(\mathrm{po}')}_{I'}\right)$ have fixed coefficients.
The Berezin integral is the coefficient of the top monomial \cite[Definition~2.8]{QEDII},
and a coefficient of a product of four elements is a sum of at most $2^{24\,\mathrm{Nx}}$ products of four coefficients.
Hence every entry of $\mathrm{Ml}_{\mathrm{po},\mathrm{po}'}\left[G(\omega)\exp\left(-S_{(\mathrm{po},\mathrm{po}')}(A)\right)\right]$
is bounded by a constant times the largest modulus of a coefficient of $G(\omega)$,
and this modulus is estimated by the condition (Ig2) of (In7)\,(iii).
The matrices have $2^{64v_s}$ rows and $2^{64v_s}$ columns, and $c_I\le6^{64v_s}$.
By (F4) and (In6), the constant depends on $\mathrm{po}$ and $\mathrm{po}'$ only through $\mathrm{po}'-\mathrm{po}$
and through the kind of $\mathrm{po}$, that is, whether $\mathrm{po}$ lies in $\mathbb{Z}$ or in $\mathbb{Z}+\tfrac12$;
one takes the larger of the two constants.

(m2)
The map $G$ is Borel, and the coefficients of $\exp\left(-S_{(\mathrm{po},\mathrm{po}')}(A)\right)$ are continuous in $A$.
The invariance under $\omega\mapsto\omega+[m]$ follows from (Ig1),
because $S_{(\mathrm{po},\mathrm{po}')}$ depends on $A$ only through the phases $\exp\left(iqA(b)\right)$,
which do not change under $A\mapsto A+r_qm$.
The linearity in $G$ holds because $\mathrm{Ml}$ is linear in its argument.

(m3)
For a configuration $\omega$ let $A(\omega)$ denote its one-cochain.
Put $\omega:=\theta_{\mathrm{po}'}^{*}\tilde{\omega}$ and
\[
\Phi:=G(\omega)\exp\left(-S_{(\mathrm{po},\mathrm{po}')}(A(\omega))\right).
\]
By (In7)\,(iii), $(\Theta_{\mathrm{po}'}G)(\tilde{\omega})=\Theta_{\mathrm{po}'}\left(G(\omega)\right)$.
By (In6) with $A'=A(\tilde{\omega})$, for which $A(\omega)=\theta_{\mathrm{po}'}^{*}A'$,
\[
\Theta_{\mathrm{po}'}\left(S_{(\mathrm{po},\mathrm{po}')}(A(\omega))\right)=S_{(\mathrm{po}',\mathrm{po}'')}(A(\tilde{\omega})).
\]
Since $\Theta_{\mathrm{po}'}$ is an antilinear anti-automorphism and the exponent is even,
\[
\Theta_{\mathrm{po}'}\left(\exp\left(-S_{(\mathrm{po},\mathrm{po}')}(A(\omega))\right)\right)=\exp\left(-S_{(\mathrm{po}',\mathrm{po}'')}(A(\tilde{\omega}))\right).
\]
Even elements commute; hence
\[
(\Theta_{\mathrm{po}'}G)(\tilde{\omega})\cdot\exp\left(-S_{(\mathrm{po}',\mathrm{po}'')}(A(\tilde{\omega}))\right)=\Theta_{\mathrm{po}'}(\Phi).
\]
By (F3), $\mathrm{Ml}_{\mathrm{po}',\mathrm{po}''}[\Theta_{\mathrm{po}'}\Phi]=\mathrm{Ml}_{\mathrm{po},\mathrm{po}'}[\Phi]^{*}$.
Under the identification of (In2) one has $\mathrm{cc}_{\mathrm{po}''}=\mathrm{cc}_{\mathrm{po}}$, because $\mathrm{po}''-\mathrm{po}$ is an integer,
and the matrices $\mathrm{cc}^{\frac12}$ are real and diagonal.
Therefore
\[
\mathrm{Ms}[\Theta_{\mathrm{po}'}G](\tilde{\omega})
=\mathrm{cc}_{\mathrm{po}}^{\frac12}\,\mathrm{Ml}_{\mathrm{po},\mathrm{po}'}[\Phi]^{*}\,\mathrm{cc}_{\mathrm{po}'}^{\frac12}
=\left(\mathrm{cc}_{\mathrm{po}'}^{\frac12}\,\mathrm{Ml}_{\mathrm{po},\mathrm{po}'}[\Phi]\,\mathrm{cc}_{\mathrm{po}}^{\frac12}\right)^{*}.
\]

(m4)
Put $\omega:=\tau_{je_1}^{*}\omega'$.
Then $(\mathrm{Tg}_{je_1}G)(\omega')=\mathrm{Tg}_{je_1}\left(G(\omega)\right)$,
and $S_{(\mathrm{po}+j,\mathrm{po}'+j)}(A(\omega'))=\mathrm{Tg}_{je_1}S_{(\mathrm{po},\mathrm{po}')}(A(\omega))$ by (In6).
So the argument of $\mathrm{Ml}_{\mathrm{po}+j,\mathrm{po}'+j}$ is $\mathrm{Tg}_{je_1}\left(G(\omega)\exp\left(-S_{(\mathrm{po},\mathrm{po}')}(A(\omega))\right)\right)$,
and the assertion follows from (F4) together with $\mathrm{cc}_{\mathrm{po}+j}=\mathrm{cc}_{\mathrm{po}}$.

(m5)
A $\Gamma$-periodic configuration is $\Gamma_{\mathrm{par}}$-periodic, and $\mathfrak{z}$ is closed on all cubes;
hence $\omega_k\in\mathrm{Cfg}[\mathrm{po}_{k-1},\mathrm{po}_k]$.

\emph{Decomposition of $S_T$.}
Extend the sequence by $\mathrm{po}_{k+jr}:=\mathrm{po}_k+jn_1$ for $j\in\mathbb{Z}$.
No site lies on a plane of the sequence.
By (In1), a bond whose end points are separated by the plane $\mathrm{po}_k$ has both end points in $\{\mathrm{po}_k-\tfrac12<y_1<\mathrm{po}_k+\tfrac12\}$;
since the planes of the sequence have their positions in $\tfrac12\mathbb{Z}$, no second plane of the sequence separates the end points of such a bond.
Consequently every bond of $\mathbb{Z}^3/\Gamma_{\mathrm{par}}$ either has both end points in one open slab $\mathrm{Sl}(\mathrm{po}_{k-1},\mathrm{po}_k)$
or crosses exactly one plane of the sequence.
The translation by $n_1e_1$ shifts the index by $r$.
Hence every bond of $T$ either has both ends in exactly one of the slabs $k=1,\dots,r$
or crosses exactly one of the planes $\mathrm{po}_0,\dots,\mathrm{po}_{r-1}$,
and every site lies in exactly one of these slabs:
\[
S_T(A)=\sum_{k=1}^{r}S_{(\mathrm{po}_{k-1},\mathrm{po}_k)}(A)+\sum_{k=0}^{r-1}\mathrm{Xc}_{\mathrm{po}_k}(A)\qquad\text{in }\mathrm{Gr}_T .
\]

\emph{The seam.}
In this identity $\mathrm{Xc}_{\mathrm{po}_0}(A)$, an element of $\mathrm{Gr}_{\mathrm{Xs}}((\mathrm{po}_0-\tfrac12,\mathrm{po}_0+\tfrac12))$, is read in $\mathrm{Gr}_T$ as in (F5).
Put $v:=n_1e_1$.
Let $x$ be the lower end of a bond that crosses the plane $\mathrm{po}_r=\mathrm{po}_0+n_1$, and let $x':=\theta_{\mathrm{po}_r}x$.
Then $x$ and $x'-v$ lie in $\{\mathrm{po}_0<y_1<\mathrm{po}_r\}$, while $x'$ and $x-v$ do not,
and $\{x-v,x'-v\}$ is the bond that crosses the plane $\mathrm{po}_0$ with the lower end $x-v$.
In $S_T(A)$ the two terms of the bond $\{x,x'\}$, written as in (In3), are read with
$\psi(x')=-\psi(x'-v)$ and $\bar{\psi}(x')=-\bar{\psi}(x'-v)$;
these are the relations of $\mathrm{Gr}_T$, since $\chi(v)=-1$.
In $\mathrm{Xc}_{\mathrm{po}_0}(A)$ the two terms of the bond $\{x-v,x'-v\}$ have the same matrices $3\,(\rho\mp\gamma_1)$
and, $A$ being periodic, the same phases,
and by (F5) they are read with $\bar{\psi}(x-v)=-\bar{\psi}(x)$ and $\psi(x-v)=-\psi(x)$.
In both readings each term receives exactly one sign, and the two terms are
\[
+\bar{\psi}(x)\,3\,(\rho-\gamma_1)\,U_A(x,x')\,\psi(x'-v)+\bar{\psi}(x'-v)\,3\,(\rho+\gamma_1)\,U_A(x',x)\,\psi(x).
\]

All terms are even.
Therefore $\exp\left(-S_T(A)\right)G_1(\omega_1)\cdots G_r(\omega_r)$ is the product of the commuting even elements
\[
\exp\left(-\mathrm{Xc}_{\mathrm{po}_k}(A)\right)=\mathrm{Ins}_{\mathrm{po}_k}[\mathrm{mt}_k],
\qquad
\mathrm{mt}_k:=\mathrm{cc}_{\mathrm{po}_k}^{\frac12}\,\mathrm{Wx}(\mathrm{po}_k)\,\mathrm{cc}_{\mathrm{po}_k}^{\frac12}
\quad(k=0,\dots,r-1),
\]
where $\mathrm{mt}_k$ is an even diagonal matrix, and
\[
\Phi_k:=G_k(\omega_k)\exp\left(-S_{(\mathrm{po}_{k-1},\mathrm{po}_k)}(A)\right)\qquad(k=1,\dots,r),
\]
and it may be written in the order of (F5).
With $\mathrm{Ml}_k:=\mathrm{Ml}_{\mathrm{po}_{k-1},\mathrm{po}_k}[\Phi_k]$, the rule (F5) gives the value
\[
\operatorname{tr}\left(\mathrm{Ml}_r\,\mathrm{mt}_{r-1}\,\mathrm{Ml}_{r-1}\cdots\mathrm{mt}_1\,\mathrm{Ml}_1\,\mathrm{mt}_0\right)
\]
for the Berezin integral over $T$.
Move the last factor $\mathrm{cc}_{\mathrm{po}_0}^{\frac12}$ to the front, by the cyclicity of the trace, and use $\mathrm{cc}_{\mathrm{po}_0}=\mathrm{cc}_{\mathrm{po}_r}$:
this gives the assertion.

(m6)
By (In7)\,(iii), $GG'\in\mathrm{Ig}^{\mathrm{ev}}_{\mathrm{Xs}}[\mathrm{po},\mathrm{po}'']$.
By (In4),
\[
\exp\left(-S_{(\mathrm{po},\mathrm{po}'')}\right)=\exp\left(-S_{(\mathrm{po},\mathrm{po}')}\right)\exp\left(-\mathrm{Xc}_{\mathrm{po}'}\right)\exp\left(-S_{(\mathrm{po}',\mathrm{po}'')}\right),
\]
a product of commuting even elements.
Hence
\begin{multline*}
(GG')(\omega)\exp\left(-S_{(\mathrm{po},\mathrm{po}'')}(A)\right)\\
=\left[G(\omega_1)\exp\left(-S_{(\mathrm{po},\mathrm{po}')}\right)\right]\,
\mathrm{Ins}_{\mathrm{po}'}\left[\mathrm{cc}_{\mathrm{po}'}^{\frac12}\,\mathrm{Wx}(\mathrm{po}')\,\mathrm{cc}_{\mathrm{po}'}^{\frac12}\right]\,
\left[G'(\omega_2)\exp\left(-S_{(\mathrm{po}',\mathrm{po}'')}\right)\right],
\end{multline*}
and the assertion follows from (F2).
\end{proof}

\begin{named}{Lemma 4.3}[the slice data]
Let $\mathrm{Qs}(A):=\mathrm{Ms}_{0,1/2}[1](\omega)$, where $\omega$ is the configuration of the slab $[0,\tfrac12]$ with the spatial links $A_0=A\in\mathbb{R}^{\mathrm{Lk}_s}$;
this slab contains one slice and no step, and $\mathrm{Ms}[1]$ does not depend on $\mathfrak{z}$.
The operator so defined is the operator $\mathrm{Qs}(A)=\mathrm{cc}^{\frac12}\,\mathrm{Nn}(A)\,\mathrm{cc}'^{\frac12}\colon\mathrm{Fk}'\to\mathrm{Fk}$ of \cite[Definition~2.27]{QEDII},
where $\mathrm{Nn}(A)$ is the half-slab matrix.
\begin{itemize}
\item[(i)]
For $t\in\mathbb{Z}$,
\[
\mathrm{Ms}_{t,t+1/2}[1](\omega)=\mathrm{Qs}(A_t),
\qquad
\mathrm{Ms}_{t+1/2,t+1}[1](\omega)=\mathrm{Qs}(A_{t+1})^{*}.
\]
\item[(ii)]
For integers $t_0\le t_1$ put
\[
\mathrm{Qp}(t_1,t_0):=\prod_{t=t_1-1,\dots,t_0}\left(\mathrm{Qs}(A_{t+1})^{*}\,\mathrm{Wq}(a_t)\,\mathrm{Qs}(A_t)\right)\in\operatorname{End}\mathrm{Fk}',
\]
the later times standing to the left; the empty product is $1$.
Then
\begin{gather*}
\mathrm{Ms}_{t_0,t_1}[1]=\mathrm{Qp}(t_1,t_0)\quad(t_0<t_1),
\qquad
\mathrm{Ms}_{t_0,t_1+1/2}[1]=\mathrm{Qs}(A_{t_1})\,\mathrm{Qp}(t_1,t_0),\\
\mathrm{Ms}_{t_0-1/2,t_1}[1]=\mathrm{Qp}(t_1,t_0)\,\mathrm{Qs}(A_{t_0})^{*},
\qquad
\mathrm{Ms}_{t_0-1/2,t_1+1/2}[1]=\mathrm{Qs}(A_{t_1})\,\mathrm{Qp}(t_1,t_0)\,\mathrm{Qs}(A_{t_0})^{*}.
\end{gather*}
\item[(iii)]
For a scalar $G\in\mathrm{Ig}_{\mathrm{Xs}}[\mathrm{po},\mathrm{po}']$,
\[
\mathrm{Ms}[G](\omega)=G(\omega)\cdot\mathrm{Ms}_{\mathrm{po},\mathrm{po}'}[1](\omega).
\]
\item[(iv)]
The entries of $\mathrm{Qs}(A)$ are polynomials in the phases $\exp\left(\pm iqA(b)\right)$, $b\in\mathrm{Lk}_s$ \cite[Lemma~5.35(a)]{QEDII};
hence $\mathrm{Qs}$ is bounded, continuous and $r_q\mathbb{Z}^{\mathrm{Lk}_s}$-periodic.
\end{itemize}
\end{named}

\begin{proof}
By (In5) and \cite[Definition~2.27]{QEDII},
\[
\mathrm{Nn}(A_t)_{I,I'}=\int_{\mathrm{hs}}C'_{I'}\exp\left(-S_{\mathrm{hs}}(A_t)\right)\Theta_{t+1/2}(C_I)
=\mathrm{Ml}_{t,t+1/2}\left[\exp\left(-S_{(t,t+1/2)}(A)\right)\right]_{I,I'} .
\]

(i)
The first identity is (m4) for $G=1$, since $\mathrm{Tg}\,1=1$.
The second identity is (m3) with $\mathrm{po}=t$, $\mathrm{po}'=t+\tfrac12$ and $G=1$, since $\Theta\,1=1$:
the reflection $\theta_{t+1/2}$ maps the plane $t+1$ onto the plane $t$ and preserves the orientation of its cells \cite[Lemma~5.4(b)]{QEDII},
so that the slice of $\theta_{t+1/2}^{*}\tilde{\omega}$ at the time $t$ is the slice of $\tilde{\omega}$ at the time $t+1$.

(ii)
Cut the slab by all the planes of $\tfrac12\mathbb{Z}$ in its interior and apply (m6) repeatedly, with $G=G'=1$.
Between two adjacent half slabs there stands $\mathrm{Wx}=\mathrm{Wq}(a_t)$ if the plane is the link plane $t+\tfrac12$,
and $\mathrm{Wx}=1$ if it is a site plane.
It remains to insert (i).

(iii)
$\mathrm{Ml}$ is linear in its argument.

(iv)
By (In4), $\exp\left(-S_{\mathrm{hs}}(A)\right)$ is a polynomial in the generators whose coefficients are polynomials in the phases,
and the Berezin integral picks out one coefficient.
\end{proof}

\subsection*{5. Slab operators}

\subsubsection*{5.1. Definitions}

For $\mathrm{po}\in\tfrac12\mathbb{Z}$ let
\[
\mathrm{Hp}(\mathrm{po}):=L^2\left(\mathrm{Sx};\mathrm{Fk}(\mathrm{po})\right).
\]
Thus $\mathrm{Hp}(\mathrm{po})=L^2(\mathrm{Sx};\mathrm{Fk}')$ for $\mathrm{po}\in\mathbb{Z}$ and $\mathrm{Hp}(\mathrm{po})=L^2(\mathrm{Sx};\mathrm{Fk})$ for $\mathrm{po}\in\mathbb{Z}+\tfrac12$;
in particular $\mathrm{Hp}(\mathrm{po}+j)=\mathrm{Hp}(\mathrm{po})$ for $j\in\mathbb{Z}$,
and by (s0) of 3.1 each $\mathrm{Hp}(\mathrm{po})$ is a separable Hilbert space.
Further let
\[
\mathrm{Ke}(\mathrm{po}):=
\begin{cases}
1, & \mathrm{po}\in\mathbb{Z},\\
\mathrm{Kl}^{\frac12}, & \mathrm{po}\in\mathbb{Z}+\tfrac12,
\end{cases}
\qquad
\mathrm{Kf}(\mathrm{po}):=\mathrm{Ke}(\mathrm{po})^2,
\]
so that $\mathrm{Kf}(\mathrm{po})=1$ for $\mathrm{po}\in\mathbb{Z}$ and $\mathrm{Kf}(\mathrm{po})=\mathrm{Kl}$ for $\mathrm{po}\in\mathbb{Z}+\tfrac12$.
For $\mathrm{mi}\in\mathrm{Sx}$ put
\[
\mathrm{we}_{\mathrm{po}}(\mathrm{mi}):=
\begin{cases}
\exp\left(-\frac14|\mathrm{Bm}(\mathrm{mi})|^2\right), & \mathrm{po}\in\mathbb{Z},\\
\exp\left(-\frac12|\mathrm{Bm}(\mathrm{mi})|^2\right), & \mathrm{po}\in\mathbb{Z}+\tfrac12,
\end{cases}
\]
and, for $\varepsilon\in(0,\tfrac12]$,
\[
\mathrm{we}^{\varepsilon}_{\mathrm{po}}(\mathrm{mi}):=
\begin{cases}
\exp\left(-\frac14|\mathrm{Bm}(\mathrm{mi})|^2\right), & \mathrm{po}\in\mathbb{Z},\\
\exp\left(-\varepsilon|\mathrm{Bm}(\mathrm{mi})|^2\right), & \mathrm{po}\in\mathbb{Z}+\tfrac12.
\end{cases}
\]

Let $\mathrm{po}<\mathrm{po}'$ be elements of $\tfrac12\mathbb{Z}$, and put $\mathrm{tl}:=\lceil\mathrm{po}\rceil$, $\mathrm{tu}:=\lfloor\mathrm{po}'\rfloor$.
The slab $[\mathrm{po},\mathrm{po}']$ is called \emph{admitted} if $\mathrm{tu}-\mathrm{tl}\ge1$, that is, if it contains at least two slices.
One has $\mathrm{tu}-\mathrm{tl}\ge\mathrm{po}'-\mathrm{po}-1$, with equality only if $\mathrm{po}$ and $\mathrm{po}'$ both lie in $\mathbb{Z}+\tfrac12$ (both planes are link planes), and in that case $\mathrm{po}'-\mathrm{po}$ is an integer.
As $\mathrm{tu}-\mathrm{tl}$ is an integer, every slab with $\mathrm{po}'-\mathrm{po}\ge\tfrac32$ is admitted;
so is every slab $[t,t+1]$ with $t\in\mathbb{Z}$.

The corners of a cell are integer points.
Hence a cell is contained in $\{\mathrm{po}\le y_1\le\mathrm{po}'\}$ if and only if it is contained in $\{\mathrm{tl}\le y_1\le\mathrm{tu}\}$, and
\[
\mathrm{Ln}[\mathrm{po},\mathrm{po}']=\mathrm{Ln}[\mathrm{tl},\mathrm{tu}],\qquad
\mathrm{Pq}[\mathrm{po},\mathrm{po}']=\mathrm{Pq}[\mathrm{tl},\mathrm{tu}],\qquad
\mathrm{Cfg}[\mathrm{po},\mathrm{po}']=\mathrm{Cfg}[\mathrm{tl},\mathrm{tu}].
\]
Let now the slab $[\mathrm{po},\mathrm{po}']$ be admitted, so that $\mathrm{po}\le\mathrm{tl}<\mathrm{tu}\le\mathrm{po}'$.
A configuration is written $\omega=\mathrm{mi}\vee\mathrm{omi}\vee\mathrm{mi}'$;
here $\mathrm{mi}=\mathrm{mi}_{\mathrm{tl}}$ is its lowest slice, $\mathrm{mi}'=\mathrm{mi}_{\mathrm{tu}}$ its highest slice,
and $\mathrm{omi}$ is its part on $\mathrm{Sb}(\mathrm{tl},\mathrm{tu})$, which consists of the slices $\mathrm{mi}_t$, $\mathrm{tl}<t<\mathrm{tu}$, and of the steps $(a_t,\mathfrak{z}^{\mathrm{el}}_t)$, $\mathrm{tl}\le t<\mathrm{tu}$.
Let
\[
N_1:=\#|\mathrm{Sb}^1(\mathrm{tl},\mathrm{tu})|=v_{\mathrm{s}}\left(3(\mathrm{tu}-\mathrm{tl})-2\right),
\qquad
N_2:=\#|\mathrm{Sb}^2(\mathrm{tl},\mathrm{tu})|=v_{\mathrm{s}}\left(3(\mathrm{tu}-\mathrm{tl})-1\right).
\]
Since $\mathrm{tl}<\mathrm{tu}\le\mathrm{po}'$ and $\mathrm{po}\le\mathrm{tl}<\mathrm{tu}$, the plane $\{y_1=\mathrm{tl}\}$ is not the plane $\{y_1=\mathrm{po}'\}$, and the plane $\{y_1=\mathrm{tu}\}$ is not the plane $\{y_1=\mathrm{po}\}$.
The magnetic plaquettes of the plane $\mathrm{tl}$ lie in the plane $\{y_1=\mathrm{po}\}$ if and only if $\mathrm{po}=\mathrm{tl}$, that is, if and only if $\mathrm{po}\in\mathbb{Z}$;
likewise the magnetic plaquettes of the plane $\mathrm{tu}$ lie in the plane $\{y_1=\mathrm{po}'\}$ if and only if $\mathrm{po}'=\mathrm{tu}$, that is, if and only if $\mathrm{po}'\in\mathbb{Z}$.
Consequently
\begin{equation}
\mathrm{Sb}^2(\mathrm{po},\mathrm{po}')=\mathrm{Sb}^2(\mathrm{tl},\mathrm{tu})\ \cup\ \left\{(\mathrm{tl};p'):p'\in\mathrm{Pl}_{\mathrm{s}}\right\}\ \cup\ \left\{(\mathrm{tu};p'):p'\in\mathrm{Pl}_{\mathrm{s}}\right\},
\tag{5.1}
\end{equation}
where the second set is present if and only if $\mathrm{po}\in\mathbb{Z}+\tfrac12$, and the third if and only if $\mathrm{po}'\in\mathbb{Z}+\tfrac12$.
In the slice variables the slab weight of \cite[Definition~5.2]{QEDII} is
\[
\mathrm{wb}_{\mathrm{tl},\mathrm{tu}}(\omega)=\prod_{\mathrm{tl}<t<\mathrm{tu}}\exp\left(-\tfrac12|\mathrm{Bm}_t|^2\right)\cdot\prod_{\mathrm{tl}\le t<\mathrm{tu}}\left(\exp\left(-\tfrac12|\mathrm{El}_t|^2\right)\cdot\mathbf{1}\left[\mathfrak{z}^{\mathrm{mag}}_{t+1}=\mathfrak{z}^{\mathrm{mag}}_t+d_{\mathrm{s}}\mathfrak{z}^{\mathrm{el}}_t\right]\right).
\]

Let the slab $[\mathrm{po},\mathrm{po}']$ be admitted, let $G\in\mathrm{Ig}^{\mathrm{ev}}_{\mathrm{Xs}}[\mathrm{po},\mathrm{po}']$, and let $\mathrm{mi},\mathrm{mi}'\in\mathrm{Sx}$.
With the interior integral $\mathrm{Ii}_{\mathrm{tl},\mathrm{tu}}$ of (In7)\,(vi) define
\begin{align*}
\mathrm{Rk}[G](\mathrm{mi}',\mathrm{mi})&:=\mathrm{we}_{\mathrm{po}'}(\mathrm{mi}')\cdot\mathrm{we}_{\mathrm{po}}(\mathrm{mi})\cdot\mathrm{Ii}_{\mathrm{tl},\mathrm{tu}}\left\{\mathrm{wb}_{\mathrm{tl},\mathrm{tu}}(\omega)\cdot\mathrm{Ms}[G](\omega)\right\}(\mathrm{mi},\mathrm{mi}')\ \in\ \mathrm{Hom}\left(\mathrm{Fk}(\mathrm{po}),\mathrm{Fk}(\mathrm{po}')\right),\\
\mathrm{rk}^+[G](\mathrm{mi}',\mathrm{mi})&:=\mathrm{we}_{\mathrm{po}'}(\mathrm{mi}')\cdot\mathrm{we}_{\mathrm{po}}(\mathrm{mi})\cdot\mathrm{Ii}_{\mathrm{tl},\mathrm{tu}}\left\{\mathrm{wb}_{\mathrm{tl},\mathrm{tu}}(\omega)\cdot\|\mathrm{Ms}[G](\omega)\|\right\}(\mathrm{mi},\mathrm{mi}').
\end{align*}
Since $G$ is defined on $\mathrm{Cfg}[\mathrm{po},\mathrm{po}']$ only, the expressions in braces are given the value $0$ at the configurations with $\mathrm{wb}_{\mathrm{tl},\mathrm{tu}}(\omega)=0$.
The integral operators with the kernels $\mathrm{Rk}[G]$ and $\mathrm{rk}^+[G]$ are denoted by the same symbols, and
\[
\mathrm{Co}[G]:=\mathrm{Ke}(\mathrm{po}')\cdot\mathrm{Rk}[G]\cdot\mathrm{Ke}(\mathrm{po})\colon\ \mathrm{Hp}(\mathrm{po})\to\mathrm{Hp}(\mathrm{po}').
\]
The weights are thereby distributed as follows.
A slice lying in a site plane that bounds the slab gives one half of its magnetic weight to each of the two slabs adjacent to that plane.
A slice next to a link plane belongs to its slab with its whole magnetic weight.
The cells cut by a link plane enter through $\mathrm{Kl}$, by (3.1), and each of the two adjacent slabs receives one factor $\mathrm{Kl}^{\frac12}$.

\begin{named}{Proposition 5.2}
Let $[\mathrm{po},\mathrm{po}']$ be an admitted slab and $G\in\mathrm{Ig}^{\mathrm{ev}}_{\mathrm{Xs}}[\mathrm{po},\mathrm{po}']$, with the constants $(\mathrm{Cq},\varepsilon)$ of (Ig2) in (In7)\,(iii), $\varepsilon=\varepsilon_G$.
\begin{itemize}
\item[(Tc0)] $\mathrm{Rk}[G]$ and $\mathrm{rk}^+[G]$ are Borel functions on $\mathrm{Sx}\times\mathrm{Sx}$, and
\begin{equation}
\|\mathrm{Rk}[G](\mathrm{mi}',\mathrm{mi})\|\le\mathrm{rk}^+[G](\mathrm{mi}',\mathrm{mi})\le C_{\mathrm{sl}}\cdot\mathrm{Cq}\cdot r_q^{N_1}\cdot\mathrm{thq}(\varepsilon)^{N_2}\cdot\mathrm{we}^{\varepsilon}_{\mathrm{po}'}(\mathrm{mi}')\cdot\mathrm{we}^{\varepsilon}_{\mathrm{po}}(\mathrm{mi}).
\tag{5.2}
\end{equation}
The operators $\mathrm{Rk}[G]$, $\mathrm{rk}^+[G]$ and $\mathrm{Co}[G]$ are Hilbert--Schmidt operators, and the map $G\mapsto\mathrm{Co}[G]$ is linear.
\item[(Tc1)] $\mathrm{Co}[\mathrm{Tg}_{je_1}G]=\mathrm{Co}[G]$ for every $j\in\mathbb{Z}$.
\item[(Tc2)] $\mathrm{Co}[\Theta_{\mathrm{po}'}G]=\mathrm{Co}[G]^{*}$, both sides being operators from $\mathrm{Hp}(\mathrm{po}')$ to $\mathrm{Hp}(2\mathrm{po}'-\mathrm{po})=\mathrm{Hp}(\mathrm{po})$.
\item[(Tc3)] Let $n_1\ge4$, $n=(n_1,n_2,n_3)$, $T=C(n_1,n_2,n_3)$,
and let $\mathrm{po}_0<\mathrm{po}_1<\dots<\mathrm{po}_r=\mathrm{po}_0+n_1$ be points of $\tfrac12\mathbb{Z}$, $r\ge2$, such that all slabs $[\mathrm{po}_{k-1},\mathrm{po}_k]$ are admitted;
let $G_k\in\mathrm{Ig}^{\mathrm{ev}}_{\mathrm{Xs}}[\mathrm{po}_{k-1},\mathrm{po}_k]$, $k=1,\dots,r$.
Then the map $(A,\mathfrak{z})\mapsto G_1(\omega_1)\cdots G_r(\omega_r)$, with $\omega_k$ as in (m5), is an even element of $\mathrm{Ig}(T)$, denoted $G_1\cdots G_r$, and
\[
Z_n[G_1\cdots G_r]=\operatorname{Tr}_{\mathrm{Hp}(\mathrm{po}_0)}\left(\mathrm{Co}[G_r]\cdots\mathrm{Co}[G_1]\right).
\]
\end{itemize}
\end{named}

\begin{proof}[Proof of Proposition 5.2\,(Tc0)]
Let $\omega=\mathrm{mi}\vee\mathrm{omi}\vee\mathrm{mi}'\in\mathrm{Cfg}[\mathrm{po},\mathrm{po}']$.
On $\mathrm{Cfg}[\mathrm{po},\mathrm{po}']$ the weight is $\mathrm{wb}_{\mathrm{tl},\mathrm{tu}}(\omega)=\prod_{p\in\mathrm{Sb}^2(\mathrm{tl},\mathrm{tu})}\exp\left(-\frac12F^{\mathrm{u}}_p(\omega)^2\right)$,
and (m1) together with (5.1) yields
\[
\mathrm{we}_{\mathrm{po}'}(\mathrm{mi}')\cdot\mathrm{we}_{\mathrm{po}}(\mathrm{mi})\cdot\mathrm{wb}_{\mathrm{tl},\mathrm{tu}}(\omega)\cdot\|\mathrm{Ms}[G](\omega)\|
\le C_{\mathrm{sl}}\cdot\mathrm{Cq}\cdot\mathrm{we}^{\varepsilon}_{\mathrm{po}'}(\mathrm{mi}')\cdot\mathrm{we}^{\varepsilon}_{\mathrm{po}}(\mathrm{mi})\cdot\prod_{p\in\mathrm{Sb}^2(\mathrm{tl},\mathrm{tu})}\exp\left(-\varepsilon F^{\mathrm{u}}_p(\omega)^2\right).
\]
Indeed, consider first the lower end of the slab.
If $\mathrm{po}\in\mathbb{Z}+\tfrac12$, the factor $\exp\left(-\frac12|\mathrm{Bm}_{\mathrm{tl}}|^2\right)$ of $\mathrm{we}_{\mathrm{po}}(\mathrm{mi})$ meets the factor $\exp\left(\left(\frac12-\varepsilon\right)|\mathrm{Bm}_{\mathrm{tl}}|^2\right)$ of (m1), and their product is $\mathrm{we}^{\varepsilon}_{\mathrm{po}}(\mathrm{mi})$;
if $\mathrm{po}\in\mathbb{Z}$, the plaquettes of the plane $\mathrm{tl}$ do not occur in (m1), and $\mathrm{we}_{\mathrm{po}}=\mathrm{we}^{\varepsilon}_{\mathrm{po}}$.
At the upper end, if $\mathrm{po}'\in\mathbb{Z}+\tfrac12$, the factor $\exp\left(-\frac12|\mathrm{Bm}_{\mathrm{tu}}|^2\right)$ of $\mathrm{we}_{\mathrm{po}'}(\mathrm{mi}')$ meets the factor $\exp\left(\left(\frac12-\varepsilon\right)|\mathrm{Bm}_{\mathrm{tu}}|^2\right)$ of (m1), and their product is $\mathrm{we}^{\varepsilon}_{\mathrm{po}'}(\mathrm{mi}')$;
if $\mathrm{po}'\in\mathbb{Z}$, the plaquettes of the plane $\mathrm{tu}$ do not occur in (m1), and $\mathrm{we}_{\mathrm{po}'}=\mathrm{we}^{\varepsilon}_{\mathrm{po}'}$.

The right side of the last inequality does not involve the constraints that define $\mathrm{Cfg}[\mathrm{po},\mathrm{po}']$, and these constraints are now dropped:
$\mathrm{rk}^+[G](\mathrm{mi}',\mathrm{mi})$ is at most the sum over all $\mathfrak{z}\in\mathbb{Z}^{\mathrm{Sb}^2(\mathrm{tl},\mathrm{tu})}$ and the integral over the box $[0,r_q)^{\mathrm{Sb}^1(\mathrm{tl},\mathrm{tu})}$ of the right side.
For fixed link variables the sum over $\mathfrak{z}$ factorises over $p$;
since $F^{\mathrm{u}}_p=(dA)_p+r_q\mathfrak{z}_p$, each factor is at most $\mathrm{thq}(\varepsilon)$ by (In7)\,(vii).
The box $[0,r_q)^{\mathrm{Sb}^1(\mathrm{tl},\mathrm{tu})}$ has the volume $r_q^{N_1}$.
This proves the second inequality of (5.2);
the first one holds because the norm of an absolutely convergent integral does not exceed the integral of the norm.

By Tonelli's theorem and (m2) the kernels are Borel functions.
By (s3) the functions $\mathrm{we}^{\varepsilon}_{\mathrm{po}}$ belong to $L^2(\mathrm{Sx})$.
Hence $\mathrm{rk}^+[G]\in L^2(\mathrm{Sx}\times\mathrm{Sx})$, and the matrix entries of $\mathrm{Rk}[G]$ lie in this space as well;
by (K3) in (In9), $\mathrm{Rk}[G]$ and $\mathrm{rk}^+[G]$ are Hilbert--Schmidt operators.
The operators $\mathrm{Ke}(\mathrm{po})$ and $\mathrm{Ke}(\mathrm{po}')$ are bounded (Lemma~3.3\,(d)), so that $\mathrm{Co}[G]$ is Hilbert--Schmidt.
The linearity of $G\mapsto\mathrm{Co}[G]$ follows from (m2).
\end{proof}

\begin{proof}[Proof of Proposition 5.2\,(Tc1)]
Put $G':=\mathrm{Tg}_{je_1}G$.
It belongs to $\mathrm{Ig}^{\mathrm{ev}}_{\mathrm{Xs}}[\mathrm{po}+j,\mathrm{po}'+j]$, and the integer planes of this slab are $\mathrm{tl}+j$ and $\mathrm{tu}+j$.
In $\mathrm{Ii}_{\mathrm{tl}+j,\mathrm{tu}+j}$ substitute $\omega'=\tau_{-je_1}^{*}\omega$:
the slice (resp.\ the step) $t+j$ of $\omega'$ is the slice (resp.\ the step) $t$ of $\omega$, and the boxes and the counting measures are preserved.
One has $\mathrm{wb}_{\mathrm{tl}+j,\mathrm{tu}+j}(\omega')=\mathrm{wb}_{\mathrm{tl},\mathrm{tu}}(\omega)$, $\mathrm{Ms}[G'](\omega')=\mathrm{Ms}[G](\omega)$ by (m4), and $\mathrm{we}_{\mathrm{po}+j}=\mathrm{we}_{\mathrm{po}}$.
It follows that $\mathrm{Rk}[G']=\mathrm{Rk}[G]$;
as moreover $\mathrm{Ke}(\mathrm{po}+j)=\mathrm{Ke}(\mathrm{po})$, one obtains $\mathrm{Co}[G']=\mathrm{Co}[G]$.
\end{proof}

\begin{proof}[Proof of Proposition 5.2\,(Tc2)]
Put $\mathrm{po}'':=2\mathrm{po}'-\mathrm{po}$, $\theta:=\theta_{\mathrm{po}'}$ and $\tilde{G}:=\Theta_{\mathrm{po}'}G\in\mathrm{Ig}^{\mathrm{ev}}_{\mathrm{Xs}}[\mathrm{po}',\mathrm{po}'']$.
The number $2\mathrm{po}'$ is an integer.
Hence the integer planes of $[\mathrm{po}',\mathrm{po}'']$ are $\lceil\mathrm{po}'\rceil=2\mathrm{po}'-\mathrm{tu}$ and $\lfloor\mathrm{po}''\rfloor=2\mathrm{po}'-\mathrm{tl}$;
their distance is $\mathrm{tu}-\mathrm{tl}$, so that this slab is admitted.
Let $\mathrm{mi},\mathrm{mi}'\in\mathrm{Sx}$, let $\tilde{\omega}=\mathrm{mi}'\vee\tilde{\mathrm{omi}}\vee\mathrm{mi}$ be a configuration of $[\mathrm{po}',\mathrm{po}'']$ with the lowest slice $\mathrm{mi}'$ and the highest slice $\mathrm{mi}$, and put $\omega:=\theta^{*}\tilde{\omega}$.

The reflection $\theta$ maps the plane $\{y_1=t\}$ onto the plane $\{y_1=2\mathrm{po}'-t\}$ and preserves the orientation of the cells lying in it \cite[Lemma~5.4(b)]{QEDII}.
It maps the cells of the step $t$ onto the cells of the step $2\mathrm{po}'-1-t$ and reverses the orientation of the temporal links and of the electric plaquettes.
Indeed, $\theta=\tau_{(2\mathrm{po}'-2t-1)e_1}\theta_{t+\frac12}$;
the plane $t+\frac12$ cuts the cells of the step $t$;
the reflection $\theta_{t+\frac12}$ maps every cell that it cuts onto itself with the orientation sign $\epsilon=-1$ \cite[Lemma~5.4(b)]{QEDII};
and translations have $\epsilon=+1$.
Marking the parts of $\tilde{\omega}$ by a tilde, one obtains
\[
\mathrm{mi}_t(\omega)=\tilde{\mathrm{mi}}_{2\mathrm{po}'-t},\qquad
a_t(\omega)=-\tilde{a}_{2\mathrm{po}'-1-t},\qquad
\mathfrak{z}^{\mathrm{el}}_t(\omega)=-\tilde{\mathfrak{z}}^{\mathrm{el}}_{2\mathrm{po}'-1-t};
\]
for $\mathrm{po}'\in\mathbb{Z}+\tfrac12$ these are the formulas for the reflected slice variables in \cite[Definition~2.26]{QEDII}.
In particular $\omega=\mathrm{mi}\vee\mathrm{omi}\vee\mathrm{mi}'$ with $\mathrm{omi}=\theta^{*}\tilde{\mathrm{omi}}$.
The map $\tilde{\mathrm{omi}}\mapsto\mathrm{omi}$ is a bijection that preserves the Lebesgue and the counting measures;
it maps the box $[0,r_q)^{\mathrm{Sb}^1(2\mathrm{po}'-\mathrm{tu},2\mathrm{po}'-\mathrm{tl})}$ onto $\prod_bI_b$, where $I_b=[0,r_q)$ for the spatial links and $I_b=(-r_q,0]$ for the temporal links $b$ of $\mathrm{Sb}^1(\mathrm{tl},\mathrm{tu})$.
Moreover the following hold.
\begin{itemize}
\item[(1)] $F^{\mathrm{u}}(\omega)=\theta^{*}F^{\mathrm{u}}(\tilde{\omega})$ and $d\mathfrak{z}(\omega)=\theta^{*}d\mathfrak{z}(\tilde{\omega})$, because $d$ commutes with $\theta^{*}$;
since $(\theta^{*}u)(\boldsymbol{\sigma})=\pm u(\theta\boldsymbol{\sigma})$, it follows that $\mathrm{wb}_{2\mathrm{po}'-\mathrm{tu},2\mathrm{po}'-\mathrm{tl}}(\tilde{\omega})=\mathrm{wb}_{\mathrm{tl},\mathrm{tu}}(\omega)$.
\item[(2)] $\mathrm{Ms}[\tilde{G}](\tilde{\omega})=\mathrm{Ms}[G](\omega)^{*}$ by (m3).
\item[(3)] $\mathrm{we}_{\mathrm{po}''}=\mathrm{we}_{\mathrm{po}}$, because $\mathrm{po}''-\mathrm{po}$ is an integer; the weights are real.
\end{itemize}
Therefore
\[
\mathrm{Rk}[\tilde{G}](\mathrm{mi},\mathrm{mi}')=\mathrm{we}_{\mathrm{po}}(\mathrm{mi})\cdot\mathrm{we}_{\mathrm{po}'}(\mathrm{mi}')\cdot\sum_{\mathfrak{z}\in\mathbb{Z}^{\mathrm{Sb}^2(\mathrm{tl},\mathrm{tu})}}\int_{\prod_bI_b}\mathrm{d}A\,\mathrm{wb}_{\mathrm{tl},\mathrm{tu}}(\omega)\cdot\mathrm{Ms}[G](\omega)^{*}.
\]
The entries of $\omega\mapsto\mathrm{wb}_{\mathrm{tl},\mathrm{tu}}(\omega)\cdot\mathrm{Ms}[G](\omega)^{*}$ are invariant under $\omega\mapsto\omega+[m]$, $m\in\mathbb{Z}^{\mathrm{Sb}^1(\mathrm{tl},\mathrm{tu})}$,
since $\mathrm{wb}_{\mathrm{tl},\mathrm{tu}}$ depends on $F^{\mathrm{u}}$ and $d\mathfrak{z}$ only, and by (m2);
and $\mathrm{Ii}_{\mathrm{tl},\mathrm{tu}}$ of their moduli is finite by (5.2).
By (In7)\,(vi) the intervals $I_b$ may therefore be replaced by $[0,r_q)$, and
\[
\mathrm{Rk}[\tilde{G}](\mathrm{mi},\mathrm{mi}')=\left(\mathrm{Rk}[G](\mathrm{mi}',\mathrm{mi})\right)^{*}\qquad\text{for all }\mathrm{mi},\mathrm{mi}'\in\mathrm{Sx},
\]
that is, $\mathrm{Rk}[\tilde{G}]=\mathrm{Rk}[G]^{*}$ by (K3).
Finally $\mathrm{Ke}(\mathrm{po}'')=\mathrm{Ke}(\mathrm{po})$, and $\mathrm{Ke}(\mathrm{po})$, $\mathrm{Ke}(\mathrm{po}')$ are self-adjoint (Lemma~3.3\,(d)); hence
\[
\mathrm{Co}[\tilde{G}]=\mathrm{Ke}(\mathrm{po})\cdot\mathrm{Rk}[G]^{*}\cdot\mathrm{Ke}(\mathrm{po}')=\left(\mathrm{Ke}(\mathrm{po}')\cdot\mathrm{Rk}[G]\cdot\mathrm{Ke}(\mathrm{po})\right)^{*}=\mathrm{Co}[G]^{*}.
\]
\end{proof}

\begin{proof}[Proof of Proposition 5.2\,(Tc3)]
\emph{Step 0 (the class).}
Let $\mathrm{tl}_k:=\lceil\mathrm{po}_{k-1}\rceil$ and $\mathrm{tu}_k:=\lfloor\mathrm{po}_k\rfloor$ for $k=1,\dots,r$, and $\mathrm{tu}_0:=\mathrm{tu}_r-n_1$.
If $\mathrm{po}_{k-1}\in\mathbb{Z}$, then $\mathrm{tl}_k=\mathrm{po}_{k-1}=\mathrm{tu}_{k-1}$;
if $\mathrm{po}_{k-1}\in\mathbb{Z}+\tfrac12$, then $\mathrm{tl}_k=\mathrm{po}_{k-1}+\frac12=\mathrm{tu}_{k-1}+1$.
Consequently
\begin{equation}
\{\mathrm{tu}_0+1,\dots,\mathrm{tu}_r\}=\bigsqcup_{k=1}^{r}\{\mathrm{tu}_{k-1}+1,\dots,\mathrm{tu}_k\},
\qquad
\{\mathrm{tu}_0,\dots,\mathrm{tu}_r-1\}=\bigsqcup_{k=1}^{r}\{\mathrm{tu}_{k-1},\dots,\mathrm{tu}_k-1\},
\tag{5.3}
\end{equation}
the unions being disjoint;
each of the two sets on the left consists of $n_1$ consecutive integers, that is, of a system of representatives of $\mathbb{Z}/n_1\mathbb{Z}$.
By (5.1) a plaquette of $\mathrm{Sb}^2(\mathrm{po}_{k-1},\mathrm{po}_k)$ is either electric, with its step in $\{\mathrm{tl}_k,\dots,\mathrm{tu}_k-1\}$, or magnetic, with a time $t$ such that $\mathrm{po}_{k-1}<t<\mathrm{po}_k$.
Hence the sets $\mathrm{Sb}^2(\mathrm{po}_{k-1},\mathrm{po}_k)$, $k=1,\dots,r$, are pairwise disjoint sets of plaquettes of $C(n)$.
Together with (Ig2) this shows that every coefficient of $\prod_kG_k(\omega_k)$ is bounded by a constant times $\exp\left(\left(\frac12-\min_k\varepsilon_{G_k}\right)\|F^{\mathrm{u}}\|^2\right)$.
The restriction of $(A+r_qm,\mathfrak{z}-dm)$ to a closed slab is $\omega_k+[m']$, $m'$ the restriction of $m$, so that (Ig1) applies.
The map is Borel and even, and therefore belongs to $\mathrm{Ig}(T)$.

\emph{Step 1 (operators).}
Let $\mathrm{Rk}_k:=\mathrm{Rk}[G_k]$.
Since $\mathrm{Ke}(\mathrm{po}_k)^2=\mathrm{Kf}(\mathrm{po}_k)$ and $\mathrm{Ke}(\mathrm{po}_r)=\mathrm{Ke}(\mathrm{po}_0)$,
\[
\mathrm{Co}[G_r]\cdots\mathrm{Co}[G_1]=\mathrm{Ke}(\mathrm{po}_0)\cdot\mathrm{Wo}\cdot\mathrm{Ke}(\mathrm{po}_0),
\qquad
\mathrm{Wo}:=\mathrm{Rk}_r\cdot\mathrm{Kf}(\mathrm{po}_{r-1})\cdot\mathrm{Rk}_{r-1}\cdots\mathrm{Kf}(\mathrm{po}_1)\cdot\mathrm{Rk}_1.
\]
The operator $\mathrm{Wo}$ is a product of $r\ge2$ Hilbert--Schmidt operators and of bounded operators; it is therefore of trace class.
For $X$ bounded and $Y$ of trace class one has $\operatorname{Tr}(XY)=\operatorname{Tr}(YX)$ \cite[Theorem~VI.25]{ReedSimon1980}.
Hence
\begin{gather*}
\operatorname{Tr}\left(\mathrm{Co}[G_r]\cdots\mathrm{Co}[G_1]\right)=\operatorname{Tr}\left(\mathrm{Wo}\cdot\mathrm{Kf}(\mathrm{po}_0)\right)=\operatorname{Tr}\left(\mathrm{Zs}_r\cdots\mathrm{Zs}_1\right),\\
\mathrm{Zs}_k:=\mathrm{Rk}_k\cdot\mathrm{Kf}(\mathrm{po}_{k-1})\colon\ \mathrm{Hp}(\mathrm{po}_{k-1})\to\mathrm{Hp}(\mathrm{po}_k).
\end{gather*}

\emph{Step 2 (kernels).}
If $\mathrm{po}_{k-1}\in\mathbb{Z}$, the operator $\mathrm{Zs}_k$ has the kernel $\mathrm{Rk}_k(\mathrm{mi}',\mathrm{mi})$; put $\mathrm{zs}^+_k:=\mathrm{rk}^+[G_k]$.
If $\mathrm{po}_{k-1}\in\mathbb{Z}+\tfrac12$, then by Lemma~3.4 the operator $\mathrm{Zs}_k$ is Hilbert--Schmidt with the kernel
\[
\mathrm{Zs}_k(\mathrm{mi}',\mathrm{mi})=\int_{\mathrm{Sx}}\mathrm{d}\mathrm{mi}''\,\mathrm{Rk}_k(\mathrm{mi}',\mathrm{mi}'')\cdot\mathrm{Kl}(\mathrm{mi}'',\mathrm{mi}),
\]
and
\[
\mathrm{zs}^+_k(\mathrm{mi}',\mathrm{mi}):=\int_{\mathrm{Sx}}\mathrm{d}\mathrm{mi}''\,\mathrm{rk}^+[G_k](\mathrm{mi}',\mathrm{mi}'')\cdot\mathrm{kl}^+(\mathrm{mi}'',\mathrm{mi})
\]
belongs to $L^2(\mathrm{Sx}\times\mathrm{Sx})$.
In both cases $\|\mathrm{Zs}_k(\mathrm{mi}',\mathrm{mi})\|\le\mathrm{zs}^+_k(\mathrm{mi}',\mathrm{mi})$.

\emph{Step 3 (the trace as an integral).}
Apply (K3) to $Y=\mathrm{Zs}_r$ and $X=\mathrm{Zs}_{r-1}\cdots\mathrm{Zs}_1$, whose kernel is the composition of the kernels.
With $\mathrm{mi}^{(r)}:=\mathrm{mi}^{(0)}$ this gives
\[
\operatorname{Tr}\left(\mathrm{Zs}_r\cdots\mathrm{Zs}_1\right)=\int_{\mathrm{Sx}^r}\prod_{k=0}^{r-1}\mathrm{d}\mathrm{mi}^{(k)}\,\operatorname{tr}_{\mathrm{Fk}(\mathrm{po}_0)}\left(\mathrm{Zs}_r(\mathrm{mi}^{(r)},\mathrm{mi}^{(r-1)})\cdots\mathrm{Zs}_2(\mathrm{mi}^{(2)},\mathrm{mi}^{(1)})\cdot\mathrm{Zs}_1(\mathrm{mi}^{(1)},\mathrm{mi}^{(0)})\right).
\]
The integral converges absolutely, also after all kernels have been written out as sums and integrals.
Indeed, $|\operatorname{tr}(X_r\cdots X_1)|\le2^{64v_{\mathrm{s}}}\prod_k\|X_k\|$ for linear maps $X_k\colon\mathrm{Fk}(\mathrm{po}_{k-1})\to\mathrm{Fk}(\mathrm{po}_k)$, $k=1,\dots,r$, with $\mathrm{Fk}(\mathrm{po}_r)=\mathrm{Fk}(\mathrm{po}_0)$;
and the integral of $\prod_k\mathrm{zs}^+_k(\mathrm{mi}^{(k)},\mathrm{mi}^{(k-1)})$, with the non-negative kernels written out, equals $\operatorname{Tr}\left(\mathrm{zs}^+_r\cdots\mathrm{zs}^+_1\right)\le\prod_k\|\mathrm{zs}^+_k\|_2<\infty$, by Tonelli's theorem and (K3).
Fubini's theorem therefore justifies all exchanges of sums and integrals below.

\emph{Step 4 (the variables).}
The variable $\mathrm{mi}^{(k)}$ is the highest slice of the slab $k$, the slice of the time $\mathrm{tu}_k$.
If $\mathrm{po}_{k-1}\in\mathbb{Z}$, the lowest slice of the slab $k$ is $\mathrm{mi}^{(k-1)}$, because $\mathrm{tl}_k=\mathrm{tu}_{k-1}$.
If $\mathrm{po}_{k-1}\in\mathbb{Z}+\tfrac12$, the lowest slice of the slab $k$, of the time $\mathrm{tl}_k=\mathrm{tu}_{k-1}+1$, is the variable $\mathrm{mi}''$ of Step~2,
and by (3.1) the variables $(a,m)$ of $\mathrm{Kl}(\mathrm{mi}'',\mathrm{mi}^{(k-1)})$ form the step $(a_t,\mathfrak{z}^{\mathrm{el}}_t)$ with $t=\mathrm{tu}_{k-1}$.
The variables of $\mathrm{Ii}_{\mathrm{tl}_k,\mathrm{tu}_k}$ are the slices of the times $\mathrm{tl}_k<t<\mathrm{tu}_k$ and the steps $\mathrm{tl}_k\le t<\mathrm{tu}_k$.
By (5.3) the following holds.
Every time $t\in\{\mathrm{tu}_0+1,\dots,\mathrm{tu}_r\}$ occurs exactly once, as a slice $\mathrm{mi}_t$ with the range $\mathrm{Sx}=[0,r_q)^{\mathrm{Lk}_{\mathrm{s}}}\times\mathbb{Z}^{\mathrm{Pl}_{\mathrm{s}}}$.
Every step $t\in\{\mathrm{tu}_0,\dots,\mathrm{tu}_r-1\}$ occurs exactly once, with $(a_t,\mathfrak{z}^{\mathrm{el}}_t)\in\mathrm{box}_0\times\mathbb{Z}^{\mathrm{Lk}_{\mathrm{s}}}$.
The slice of the time $\mathrm{tu}_0$ is $\mathrm{mi}^{(0)}=\mathrm{mi}^{(r)}$, which is also the slice of the time $\mathrm{tu}_r=\mathrm{tu}_0+n_1$.
The indicators that occur, in the weights $\mathrm{wb}$ and in (3.1), are the indicators $\mathbf{1}\left[\mathfrak{z}^{\mathrm{mag}}_{t+1}=\mathfrak{z}^{\mathrm{mag}}_t+d_{\mathrm{s}}\mathfrak{z}^{\mathrm{el}}_t\right]$, one for each step.
By (In7)\,(i), (ii) the families of variables admitted by these indicators are exactly the pairs $(A,\mathfrak{z})$ with $A\in[0,r_q)^{\text{links of }C(n)}$ and $\mathfrak{z}\in C^{2}(n;\mathbb{Z})$, $d\mathfrak{z}=0$;
and the configuration $\mathrm{mi}_{\mathrm{tl}_k}\vee\mathrm{omi}\vee\mathrm{mi}_{\mathrm{tu}_k}$ of the slab $k$ is $\omega_k$.

\emph{Step 5 (the weights).}
Every step carries the factor $\exp\left(-\frac12|\mathrm{El}_t|^2\right)$ exactly once, from a weight $\mathrm{wb}$ or from (3.1).
A slice of a time $\mathrm{tl}_k<t<\mathrm{tu}_k$ carries $\exp\left(-\frac12|\mathrm{Bm}_t|^2\right)$, from $\mathrm{wb}$.
The slice of the time $\mathrm{tu}_k$ carries the factor $\mathrm{we}_{\mathrm{po}_k}(\mathrm{mi}^{(k)})$ of $\mathrm{Rk}_k$.
If $\mathrm{po}_k\in\mathbb{Z}+\tfrac12$, this factor is $\exp\left(-\frac12|\mathrm{Bm}|^2\right)$.
If $\mathrm{po}_k\in\mathbb{Z}$, it is $\exp\left(-\frac14|\mathrm{Bm}|^2\right)$;
then the same slice is the lowest slice of the slab $k+1$ (of the slab $1$ if $k=r$) and carries a second factor $\exp\left(-\frac14|\mathrm{Bm}|^2\right)$, from $\mathrm{Rk}_{k+1}$.
A lowest slice above a link plane carries $\mathrm{we}_{\mathrm{po}_{k-1}}=\exp\left(-\frac12|\mathrm{Bm}|^2\right)$.
The product of all weights is therefore $\exp\left(-\frac12\|F^{\mathrm{u}}\|^2\right)$.

\emph{Step 6 (the fermions).}
The matrices under the trace are the $\mathrm{Ms}[G_k](\omega_k)$, coming from $\mathrm{Rk}_k$, and, at the link planes, the matrices $\mathrm{Wq}(a_t)=\mathrm{Wx}(\mathrm{po}_{k-1})$, coming from $\mathrm{Kl}$;
at the site planes $\mathrm{Wx}(\mathrm{po}_{k-1})=1$.
They stand in the order of the left side of
\begin{multline*}
\operatorname{tr}_{\mathrm{Fk}(\mathrm{po}_0)}\left(\mathrm{Ms}[G_r](\omega_r)\cdot\mathrm{Wx}(\mathrm{po}_{r-1})\cdot\mathrm{Ms}[G_{r-1}](\omega_{r-1})\cdots\mathrm{Ms}[G_1](\omega_1)\cdot\mathrm{Wx}(\mathrm{po}_0)\right)\\
=\int_T\exp\left(-S_T(A)\right)\cdot G_1(\omega_1)\cdots G_r(\omega_r),
\end{multline*}
and the equality is (m5).

Steps~3--6 together give
\begin{align*}
\operatorname{Tr}\left(\mathrm{Zs}_r\cdots\mathrm{Zs}_1\right)
&=\sum_{\mathfrak{z}\in C^{2}(n;\mathbb{Z}),\ d\mathfrak{z}=0}\ \int_{[0,r_q)^{\text{links of }C(n)}}\mathrm{d}A\,\exp\left(-\tfrac12\|F^{\mathrm{u}}\|^2\right)\\
&\qquad\cdot\int_T\exp\left(-S_T(A)\right)\cdot G_1(\omega_1)\cdots G_r(\omega_r)\\
&=Z_n[G_1\cdots G_r]
\end{align*}
by (In7)\,(iv).
\end{proof}

\begin{named}{Remark 5.4}[Scalar contents; the properties of the fermion that enter]
Let all contents $G$ be scalar.
By Lemma~4.3\,(ii), (iii), $\mathrm{Ms}[G](\omega)$ is $G(\omega)$ times an ordered product of matrices $\mathrm{Qs}(A_t)$, $\mathrm{Qs}(A_t)^{*}$, $\mathrm{Wq}(a_t)$;
let $\mathrm{Ms}[G]$ be defined by this formula.
Then (m1) follows from $\sup_A\|\mathrm{Qs}(A)\|<\infty$ and from the unitarity of $\mathrm{Wq}(a)$.
Property (m2) follows from the continuity and the periodicity of $\mathrm{Qs}$ and from $\mathrm{Wq}(a)=1$ for $a\in r_q\mathbb{Z}^{\Lambda_{\mathrm{s}}}$;
these are the statements of \cite[Lemma~5.35(a)]{QEDII}, that is, Lemma~4.3\,(iv) together with (In2).
Property (m4) holds because the product of Lemma~4.3\,(ii) for the slab $[\mathrm{po}+j,\mathrm{po}'+j]$ at $\omega'$ involves the slices and the steps $t+j$ of $\omega'$, which are the slices and the steps $t$ of $\tau_{je_1}^{*}\omega'$.
Property (m3) is the rule $(X_1X_2\cdots X_j)^{*}=X_j^{*}\cdots X_2^{*}X_1^{*}$, combined with $\mathrm{Wq}(a)^{*}=\mathrm{Wq}(-a)$ and with the formulas for $\theta^{*}\tilde{\omega}$ in the proof of (Tc2).
Property (m5) is the identity
\[
\det_n\left(D_A+\boldsymbol{\mu}\right)=\operatorname{tr}_{\mathrm{Fk}'}\prod_{t=n_1-1,\dots,0}\left(\mathrm{Qs}(A_{t+1})^{*}\cdot\mathrm{Wq}(a_t)\cdot\mathrm{Qs}(A_t)\right),
\]
which is \cite[Proposition~5.33]{QEDII} with the operators $\mathrm{Wq}(a_t)$ at the link planes, after a cyclic permutation of the factors.
Of $\mathrm{Wq}$, Section~3 requires only that it is a unitary representation of $\mathbb{R}^{\Lambda_{\mathrm{s}}}$ with $\mathrm{Wq}(a)=1$ for $a\in r_q\mathbb{Z}^{\Lambda_{\mathrm{s}}}$, which is the last assertion of \cite[Lemma~5.35(a)]{QEDII}.

\emph{The slice data.}
For scalar contents the fermion therefore enters Sections~3, 5 and~6 only through the slice data for the time direction $e_1$:
the identity of \cite[Proposition~5.33]{QEDII} with the operators $\mathrm{Wq}(a_t)$ at the link planes, and the properties of $\mathrm{Qs}$ and $\mathrm{Wq}$ in \cite[Lemma~5.35(a)]{QEDII}.
These data give the absolute convergence of $Z_n[P]$ for scalar $P\in\mathrm{Ig}(T)$, $T=C(n)$:
by the two statements,
\[
\left|\det_n\left(D_A+\boldsymbol{\mu}\right)\right|\le2^{64v_{\mathrm{s}}}\cdot\left(\sup_A\|\mathrm{Qs}(A)\|\right)^{2n_1},
\]
and one concludes with (In7)\,(vii), as in the proof of (Tc0).
They also give Proposition~6.3 for scalar $G$, hence (7.1) and Proposition~9.3\,(iv) for $\mu=1$.

\emph{The two further properties used in Proposition~9.3.}
\begin{itemize}
\item[(i)] For the inequality (7.1) of one direction with $\mu=2,3$ and for Proposition~9.3\,(iv), where it enters through \cite[Lemma~5.46]{QEDII}:
the identity
\[
Z_{n^{\mathrm{pm}}}\left[P\circ g_{\mathrm{w}}^{*}\right]=Z_n[P]
\]
for $g_{\mathrm{w}}=\mathrm{wr}_\mu^{-1}$ (in the notation of 7.1) and scalar $P\in\mathrm{Ig}(T)$, $T=C(n)$;
here $P\circ g_{\mathrm{w}}^{*}$ is the function $(A',\mathfrak{z}')\mapsto P\left(g_{\mathrm{w}}^{*}(A',\mathfrak{z}')\right)$ on the configurations of $C(n^{\mathrm{pm}})$.
The only property of the fermion contained in this identity is the covariance
\[
\det_{n^{\mathrm{pm}}}\left(D_{g_{\mathrm{w}*}A}+\boldsymbol{\mu}\right)=\det_n\left(D_A+\boldsymbol{\mu}\right),\qquad g_{\mathrm{w}*}:=(g_{\mathrm{w}}^{-1})^{*},
\]
for real one-cochains $A$ on $C(n)$ (\cite[Proposition~3.16(b)]{QEDII}; it is the last identity of Proposition~50\,(d)).
Indeed, $g_{\mathrm{w}}^{*}$ maps the configurations $(A',\mathfrak{z}')$ of $C(n^{\mathrm{pm}})$ onto the configurations of $C(n)$.
Since $g_{\mathrm{w}}$ is a permutation matrix, it maps positively oriented links to positively oriented links;
so $g_{\mathrm{w}}^{*}$ permutes the link variables and maps the box $[0,r_q)^{\text{links of }C(n^{\mathrm{pm}})}$ onto the box $[0,r_q)^{\text{links of }C(n)}$.
On integer two-cochains $g_{\mathrm{w}}^{*}$ is a permutation of the coordinates with the signs $\epsilon(g_{\mathrm{w}},p)$;
it commutes with $d$ and preserves the norm $\|\cdot\|$ (Section~I.4.7).
Hence $F^{\mathrm{u}}\left(g_{\mathrm{w}}^{*}(A',\mathfrak{z}')\right)=g_{\mathrm{w}}^{*}F^{\mathrm{u}}(A',\mathfrak{z}')$ has the norm of $F^{\mathrm{u}}(A',\mathfrak{z}')$,
and the substitution $(A,\mathfrak{z})=g_{\mathrm{w}}^{*}(A',\mathfrak{z}')$, $A'=g_{\mathrm{w}*}A$, gives
\[
Z_{n^{\mathrm{pm}}}\left[P\circ g_{\mathrm{w}}^{*}\right]=\sum_{\mathfrak{z}\in C^{2}(n;\mathbb{Z}),\ d\mathfrak{z}=0}\ \int_{[0,r_q)^{\text{links of }C(n)}}\mathrm{d}A\,\exp\left(-\tfrac12\|F^{\mathrm{u}}\|^2\right)\cdot\det_{n^{\mathrm{pm}}}\left(D_{g_{\mathrm{w}*}A}+\boldsymbol{\mu}\right)\cdot P(A,\mathfrak{z}),
\]
which equals $Z_n[P]$ by the covariance.

\emph{Example.}
The slice data concern the single time direction $e_1$, and the following data satisfy them but not the covariance.
Let $\mathrm{Fk}'=\mathrm{Fk}=\mathbb{C}$, $\mathrm{Wq}:=1$,
\[
\mathrm{Qs}(A):=\prod_{b\in\mathrm{Lk}_{\mathrm{s}},\ b\parallel e_2}\left(2+\cos\left(qA(b)\right)\right),
\]
and take, in place of $\det_n\left(D_A+\boldsymbol{\mu}\right)$, the weight $\prod_b\left(2+\cos\left(qA(b)\right)\right)^2$, the product extending over the links $b$ of $C(n)$ parallel to $e_2$.
For every $n$ these data satisfy the identity of \cite[Proposition~5.33]{QEDII} with the operators $\mathrm{Wq}(a_t)$ and the statements of \cite[Lemma~5.35(a)]{QEDII};
under the trace there stands the number $\prod_t\mathrm{Qs}(A_{t+1})\cdot\mathrm{Qs}(A_t)=\prod_t\mathrm{Qs}(A_t)^2$.
For this weight the covariance fails on every $C(n)$, for $\mu=2$ as well as for $\mu=3$.
Let $A$ take the value $\pi/q$ on one link $b$ parallel to $e_2$ and the value $0$ on all other links of $C(n)$.
The right side equals $9^{n_1n_2n_3-1}$, because $2+\cos\pi=1$ and $2+\cos0=3$.
The cochain $g_{\mathrm{w}*}A$ takes the value $\pi/q$ on the link $g_{\mathrm{w}}b$ and $0$ on all other links of $C(n^{\mathrm{pm}})$;
the link $g_{\mathrm{w}}b$ is parallel to the vector $e_\nu$ with $\mathrm{pm}(\nu)=2$, and $\nu\ne2$;
so the left side equals $9^{n_1n_2n_3}$.

Given the identity of \cite[Proposition~5.33]{QEDII} with the operators $\mathrm{Wq}(a_t)$ for the time direction $e_1$, the covariance is equivalent to the same identity for the time direction $e_\mu$, with the objects transported by $\mathrm{wr}_\mu$ \cite[Proposition~5.47]{QEDII}.
\item[(ii)] For Proposition~9.3\,(v):
the positivity $Z_n[1]>0$ \cite[Theorem~4.26(a)]{QEDII}, which is contained in Proposition~50\,(e).
\emph{Example.}
With $\mathrm{Qs}:=0$ and the weight $0$ in place of the determinant, the identity of \cite[Proposition~5.33]{QEDII} with the operators $\mathrm{Wq}(a_t)$ and the statements of \cite[Lemma~5.35(a)]{QEDII} remain true, whereas $Z_n[1]=0$.
\end{itemize}

\emph{Summary.}
Proposition~9.3\,(iii) does not involve the fermion.
The slice data for the time direction $e_1$ give Proposition~9.3\,(i), the inequality (7.1) for $\mu=1$ and Proposition~9.3\,(iv) for $\mu=1$.
Together with the property (i) they give Proposition~9.3\,(ii) and (iv);
together with the properties (i) and (ii) they give Proposition~9.3\,(v).
\end{named}

\subsection*{6. One direction}

\begin{named}{Lemma 6.1}[H\"older's inequality for a closed chain of operators]
Let $k\ge1$, let $\mathrm{Hb}_0,\dots,\mathrm{Hb}_{2k}$ be separable Hilbert spaces with $\mathrm{Hb}_{2k}=\mathrm{Hb}_0$, and let $Y_j\colon\mathrm{Hb}_{j-1}\to\mathrm{Hb}_j$ be Hilbert--Schmidt operators, $j=1,\dots,2k$.
Then $Y_{2k}\cdots Y_1$ is of trace class on $\mathrm{Hb}_0$, and
\[
\left|\operatorname{Tr}_{\mathrm{Hb}_0}\left(Y_{2k}\cdots Y_1\right)\right|\le\prod_{j=1}^{2k}\left(\operatorname{Tr}_{\mathrm{Hb}_{j-1}}\left(Y_j^{*}Y_j\right)^k\right)^{\frac{1}{2k}}.
\]
\end{named}

\begin{proof}
Let $\mathrm{Hd}:=\mathrm{Hb}_0\oplus\dots\oplus\mathrm{Hb}_{2k-1}$, let $\mathrm{em}_j\colon\mathrm{Hb}_j\to\mathrm{Hd}$ be the canonical isometries, so that $\mathrm{em}_j^{*}\mathrm{em}_j=1$, put $\mathrm{em}_{2k}:=\mathrm{em}_0$ and $\mathrm{Yh}_j:=\mathrm{em}_j\cdot Y_j\cdot\mathrm{em}_{j-1}^{*}$.
Then
\[
\mathrm{Yh}_j^{*}\mathrm{Yh}_j=\mathrm{em}_{j-1}\cdot Y_j^{*}Y_j\cdot\mathrm{em}_{j-1}^{*},
\qquad
\left(\mathrm{Yh}_j^{*}\mathrm{Yh}_j\right)^k=\mathrm{em}_{j-1}\cdot\left(Y_j^{*}Y_j\right)^k\cdot\mathrm{em}_{j-1}^{*},
\]
and $\operatorname{Tr}_{\mathrm{Hd}}\left(\mathrm{Yh}_j^{*}\mathrm{Yh}_j\right)^k=\operatorname{Tr}_{\mathrm{Hb}_{j-1}}\left(Y_j^{*}Y_j\right)^k$;
this number is finite, since $Y_j^{*}Y_j$ is of trace class ((In9)\,(b)).
Moreover $\mathrm{Yh}_{2k}\cdots\mathrm{Yh}_1=\mathrm{em}_0\cdot\left(Y_{2k}\cdots Y_1\right)\cdot\mathrm{em}_0^{*}$.
A product of at least two Hilbert--Schmidt operators is of trace class, and $\operatorname{Tr}_{\mathrm{Hd}}\left(\mathrm{em}_0X\mathrm{em}_0^{*}\right)=\operatorname{Tr}_{\mathrm{Hb}_0}X$.
The inequality now follows from (In9)\,(a), applied with the $2k$ exponents $2k$, and from (In9)\,(b), which gives $\|\mathrm{Yh}_j\|_{2k}^{2k}=\operatorname{Tr}\left(\mathrm{Yh}_j^{*}\mathrm{Yh}_j\right)^k$.
\end{proof}

\subsubsection*{6.2. Copies in slabs (class $\mathrm{Ig}$)}

Let $\mathrm{po}_0\in\tfrac12\mathbb{Z}$ and $\mathrm{hsl}\in\tfrac12\mathbb{Z}$ with $\mathrm{hsl}\ge\tfrac32$;
put $\mathrm{po}_j:=\mathrm{po}_0+j\,\mathrm{hsl}$ for $j\in\mathbb{Z}$, and $\tau:=2\mathrm{hsl}\,e_1$, the vector $\tau$ of the slabs $[\mathrm{po}_j,\mathrm{po}_{j+1}]$.
It lies in $\mathbb{Z}e_1$, since $2\mathrm{hsl}$ is a positive integer.
All slabs $[\mathrm{po}_j,\mathrm{po}_{j+1}]$ are admitted by 5.1, because $\mathrm{po}_{j+1}-\mathrm{po}_j=\mathrm{hsl}\ge\tfrac32$.
For $G\in\mathrm{Ig}^{\mathrm{ev}}_{\mathrm{Xs}}[\mathrm{po}_j,\mathrm{po}_{j+1}]$ and $i\in\mathbb{Z}$ the \emph{copy of $G$ in the slab $i$} is
\[
\mathrm{cp}_iG:=
\begin{cases}
\mathrm{Tg}_{s\tau}G, & i-j=2s,\\
\mathrm{Tg}_{s\tau}\Theta_{\mathrm{po}_{j+1}}G, & i-j=2s+1,
\end{cases}
\qquad s\in\mathbb{Z}.
\]
By (In7)\,(iii), $\mathrm{cp}_iG$ belongs to $\mathrm{Ig}^{\mathrm{ev}}_{\mathrm{Xs}}[\mathrm{po}_i,\mathrm{po}_{i+1}]$:
the element $\Theta_{\mathrm{po}_{j+1}}G$ belongs to $\mathrm{Ig}^{\mathrm{ev}}_{\mathrm{Xs}}[\mathrm{po}_{j+1},\mathrm{po}_{j+2}]$, and the first component of $s\tau$ is $2s\,\mathrm{hsl}$.
For $i=j$ one has $s=0$, and $\mathrm{cp}_jG=G$.

\begin{named}{Proposition 6.3}[One direction, class $\mathrm{Ig}$]
In the situation of 6.2 let $k\ge1$ be an integer such that $n_1:=2k\,\mathrm{hsl}\ge4$, and let $T=C(n_1,n_2,n_3)$, $n=(n_1,n_2,n_3)$.
\begin{itemize}
\item[(i)] For $j\in\{0,\dots,2k-1\}$ and $G\in\mathrm{Ig}^{\mathrm{ev}}_{\mathrm{Xs}}[\mathrm{po}_j,\mathrm{po}_{j+1}]$ one has, with $Y:=\mathrm{Co}[G]$,
\begin{equation}
Z_n\left[(\mathrm{cp}_0G)(\mathrm{cp}_1G)\cdots(\mathrm{cp}_{2k-1}G)\right]=\operatorname{Tr}_{\mathrm{Hp}(\mathrm{po}_j)}\left(Y^{*}Y\right)^k\ge0.
\tag{6.1}
\end{equation}
\item[(ii)] For $G_j\in\mathrm{Ig}^{\mathrm{ev}}_{\mathrm{Xs}}[\mathrm{po}_j,\mathrm{po}_{j+1}]$, $j=0,\dots,2k-1$,
\begin{equation}
\left|Z_n\left[G_0G_1\cdots G_{2k-1}\right]\right|\le\prod_{j=0}^{2k-1}Z_n\left[(\mathrm{cp}_0G_j)(\mathrm{cp}_1G_j)\cdots(\mathrm{cp}_{2k-1}G_j)\right]^{\frac{1}{2k}}.
\tag{6.2}
\end{equation}
\end{itemize}
\end{named}

\begin{proof}
The number $n_1=k\cdot2\mathrm{hsl}$ is an integer, and $n_1e_1=k\tau$.

(i) By (Tc1) and (Tc2), $\mathrm{Co}[\mathrm{cp}_iG]=Y$ if $i-j$ is even, and $\mathrm{Co}[\mathrm{cp}_iG]=Y^{*}$ if $i-j$ is odd.
By (Tc3), applied with $r=2k\ge2$ and the planes $\mathrm{po}_0<\dots<\mathrm{po}_{2k}=\mathrm{po}_0+n_1$,
\[
Z_n\left[(\mathrm{cp}_0G)\cdots(\mathrm{cp}_{2k-1}G)\right]=\operatorname{Tr}\left(\mathrm{Co}[\mathrm{cp}_{2k-1}G]\cdots\mathrm{Co}[\mathrm{cp}_0G]\right).
\]
If $j$ is even, this is $\operatorname{Tr}\left(Y^{*}Y\,Y^{*}Y\cdots Y^{*}Y\right)=\operatorname{Tr}\left(Y^{*}Y\right)^k$.
If $j$ is odd, it is
\[
\operatorname{Tr}\left(YY^{*}\cdots YY^{*}\right)=\operatorname{Tr}\left(YY^{*}\right)^k=\operatorname{Tr}\left(Y\cdot Y^{*}(YY^{*})^{k-1}\right)=\operatorname{Tr}\left(Y^{*}(YY^{*})^{k-1}\cdot Y\right)=\operatorname{Tr}\left(Y^{*}Y\right)^k,
\]
where the trace of the product of the two Hilbert--Schmidt operators $Y$ and $Y^{*}(YY^{*})^{k-1}$ does not depend on their order.
The operator $(Y^{*}Y)^k\ge0$ is of trace class, so that its trace is non-negative.

(ii) By (Tc3), $Z_n[G_0\cdots G_{2k-1}]=\operatorname{Tr}\left(Y_{2k-1}\cdots Y_0\right)$ with $Y_j:=\mathrm{Co}[G_j]\colon\mathrm{Hp}(\mathrm{po}_j)\to\mathrm{Hp}(\mathrm{po}_{j+1})$ and $\mathrm{Hp}(\mathrm{po}_{2k})=\mathrm{Hp}(\mathrm{po}_0)$;
these spaces are separable by (s0) of 3.1.
The inequality (6.2) follows from Lemma~6.1 and (6.1).
\end{proof}

\subsection*{7. The three directions}

\paragraph{7.1. Rotation of the time direction.}
For $\mu\in\{1,2,3\}$ let $\mathrm{wr}_\mu\in O_{\mathrm{cub}}$ be the cyclic permutation matrix with $\mathrm{wr}_\mu e_1=e_\mu$, as in \cite[Definition~2.35]{QEDII};
in particular $\mathrm{wr}_1=1$.
Let $\mathrm{pm}$ be the cyclic permutation of $\{1,2,3\}$ determined by $\mathrm{wr}_\mu e_\nu=e_{\mathrm{pm}(\nu)}$;
then $\mathrm{pm}(1)=\mu$.
The element $g_{\mathrm{w}}:=\mathrm{wr}_\mu^{-1}$ of $W_{\mathrm{af}}$ is even, and for $y\in\mathbb{R}^3$
\[
\left(g_{\mathrm{w}}y\right)_\nu=e_\nu\cdot\mathrm{wr}_\mu^{-1}y=\left(\mathrm{wr}_\mu e_\nu\right)\cdot y=y_{\mathrm{pm}(\nu)} .
\]
For a triple $z=(z_1,z_2,z_3)$ write $z^{\mathrm{pm}}:=\left(z_{\mathrm{pm}(1)},z_{\mathrm{pm}(2)},z_{\mathrm{pm}(3)}\right)$;
thus $g_{\mathrm{w}}y=y^{\mathrm{pm}}$.
One has the following.
\begin{itemize}
\item[(r1)]
$g_{\mathrm{w}}$ maps $\Gamma=\sum_\nu n_\nu\mathbb{Z}e_\nu$ onto $\sum_\nu n_{\mathrm{pm}(\nu)}\mathbb{Z}e_\nu$,
and it carries the character of $C(n)$, which is $-1$ on the three generators $n_\nu e_\nu$,
to the character that is $-1$ on the three generators $n_{\mathrm{pm}(\nu)}e_\nu$.
Hence $g_{\mathrm{w}}C(n)=C(n^{\mathrm{pm}})$ in the sense of \cite[Definition~2.11]{QEDII}.
\item[(r2)]
$y_\nu\circ g_{\mathrm{w}}=y_{\mathrm{pm}(\nu)}$.
Consequently $g_{\mathrm{w}}$ maps the plane $\{y_{\mathrm{pm}(\nu)}=\mathrm{po}\}$ onto the plane $\{y_\nu=\mathrm{po}\}$, and
\[
g_{\mathrm{w}}\,\theta_{y_{\mathrm{pm}(\nu)},\mathrm{po}}\,g_{\mathrm{w}}^{-1}=\theta_{y_\nu,\mathrm{po}} .
\]
\item[(r3)]
$\mathcal{Q}^{\mathrm{pm}}:=\left(\mathfrak{F}_{\mathrm{c}};o^{\mathrm{pm}};\mathrm{hw}\right)$ is a coordinate cell structure.
The isometry $g_{\mathrm{w}}$ maps the walls of $\mathcal{Q}$ onto the walls of $\mathcal{Q}^{\mathrm{pm}}$ and $Q^{\circ}_{\mathfrak{c}}(\mathcal{Q})$ onto $Q^{\circ}_{\mathfrak{c}^{\mathrm{pm}}}(\mathcal{Q}^{\mathrm{pm}})$, and
\[
g_{\mathrm{w}}\,\mathrm{Di}_{\mathrm{pm}(\nu)}(\mathcal{Q})\,g_{\mathrm{w}}^{-1}=\mathrm{Di}_\nu(\mathcal{Q}^{\mathrm{pm}}),
\qquad
g_{\mathrm{w}}\,W(\mathcal{Q})\,g_{\mathrm{w}}^{-1}=W(\mathcal{Q}^{\mathrm{pm}}) .
\]
By the uniqueness in 1.3 it follows that
\[
g_{\mathrm{w}}\,\mathrm{di}^{\mathrm{pm}(\nu)}_{i,j}(\mathcal{Q})\,g_{\mathrm{w}}^{-1}=\mathrm{di}^{\nu}_{i,j}(\mathcal{Q}^{\mathrm{pm}}) .
\]
The torus $C(n)$ carries $\mathcal{Q}$ if and only if $C(n^{\mathrm{pm}})$ carries $\mathcal{Q}^{\mathrm{pm}}$.
In that case the numbers of cells per side are the entries of $\mathrm{sv}^{\mathrm{pm}}$,
and $\mathfrak{c}\mapsto\mathfrak{c}^{\mathrm{pm}}$ maps the brick $\mathrm{Kb}$ of $(C(n),\mathcal{Q})$ onto the brick of $(C(n^{\mathrm{pm}}),\mathcal{Q}^{\mathrm{pm}})$.
\end{itemize}

\paragraph{7.2. Notation; the inequality of one direction.}
For $\mathfrak{c}\in\mathbb{Z}^3$ let $\mathfrak{c}[\mu:j]$ denote the triple obtained from $\mathfrak{c}$ by replacing its entry $\mu$ with $j$.
The families $\mathrm{GG}$ on the brick $\mathrm{Kb}$, their products $\mathrm{Fp}(\mathrm{GG})$ and their folds $\mathrm{Rm}^j_\mu\mathrm{GG}$ are defined in the proof of Proposition~9.3.
For such families and for $\mu\in\{1,2,3\}$, the \emph{inequality of the direction $\mu$} is the assertion
\begin{equation}\tag{7.1}
Z_n\left[\mathrm{Fp}\left(\mathrm{Rm}^j_\mu\mathrm{GG}\right)\right]\ge0\quad(0\le j<\mathrm{sv}_\mu),
\qquad
\left|Z_n\left[\mathrm{Fp}(\mathrm{GG})\right]\right|\le\prod_{j=0}^{\mathrm{sv}_\mu-1}Z_n\left[\mathrm{Fp}\left(\mathrm{Rm}^j_\mu\mathrm{GG}\right)\right]^{\frac{1}{\mathrm{sv}_\mu}} .
\end{equation}

\subsection*{8. The tori of clause (e); convexity on a set}

Let $\mathrm{sv}$ be as in Proposition~39\,(e), and put $n:=\mathrm{hw}\cdot\mathrm{sv}$.
Then $n_\mu=\frac{\mathrm{sv}_\mu}{2}\cdot2\mathrm{hw}$ is an integer, $n_\mu\ge4$, and $n_\mu\in2\mathrm{hw}\mathbb{Z}$.
Hence $C(n)$ carries $\mathcal{Q}$, and its brick is $\mathrm{Kb}=\{0\le\mathfrak{c}_\mu<\mathrm{sv}_\mu\}$ (2.4).

A real function $f$ on a set of real numbers is called \emph{convex on this set} if
\[
\left(\mathrm{xv}''-\mathrm{xv}\right)f(\mathrm{xv}')\le\left(\mathrm{xv}''-\mathrm{xv}'\right)f(\mathrm{xv})+\left(\mathrm{xv}'-\mathrm{xv}\right)f(\mathrm{xv}'')
\]
for all $\mathrm{xv}<\mathrm{xv}'<\mathrm{xv}''$ in the set.
A convex function on an interval satisfies this inequality for all points of the interval,
because $\mathrm{xv}'$ is the convex combination of $\mathrm{xv}$ and $\mathrm{xv}''$ with the weights $\frac{\mathrm{xv}''-\mathrm{xv}'}{\mathrm{xv}''-\mathrm{xv}}$ and $\frac{\mathrm{xv}'-\mathrm{xv}}{\mathrm{xv}''-\mathrm{xv}}$.
The inequality is homogeneous of degree one in the three points.
Therefore, if $f$ is convex on an interval and $\mathrm{cf}>0$, then $\mathrm{xv}\mapsto f(\mathrm{cf}\cdot\mathrm{xv})$ is convex on every set that the map $\mathrm{xv}\mapsto\mathrm{cf}\cdot\mathrm{xv}$ sends into the interval.

\subsection*{9. Closed cells}

\paragraph{9.2. Closed cells.}
Let $\mathcal{Q}$ be as in Definition~39.0, with $\mathrm{hw}\ge\frac32$.
The sets $\mathrm{Pqp}(\mathfrak{c})$, the classes $\mathrm{Fcl}(\mathfrak{c})$, the copies $\Phi^{[\mathfrak{c}']}$ and the products $\Phi^{\sharp}_n$ are those of Remark~39.R, paragraph \emph{Closed cells}.
The closure $\mathrm{cl}\,Q^{\circ}_{\mathfrak{c}}$ of $Q^{\circ}_{\mathfrak{c}}$ is called the closed cell.
Let $g_{\mathrm{w}}\in W_{\mathrm{af}}$ be an isometry that maps walls of $\mathcal{Q}$ to walls of a cell structure $\mathcal{Q}'$ and $Q^{\circ}_{\mathfrak{c}}(\mathcal{Q})$ onto $Q^{\circ}_{\mathfrak{c}'}(\mathcal{Q}')$.
Then $g_{\mathrm{w}}$ maps $\mathrm{Pqp}(\mathfrak{c})$ onto the set $\mathrm{Pqp}(\mathfrak{c}')$ formed with $\mathcal{Q}'$.
For $\Phi\in\mathrm{Fcl}(\mathfrak{c})$ and $u\in\mathbb{R}^{\mathrm{Pqp}(\mathfrak{c}')}$ put
\[
\Phi^{g_{\mathrm{w}}}(u):=
\begin{cases}
\Phi\left(\left.\left(g_{\mathrm{w}}^{*}u\right)\right|_{\mathrm{Pqp}(\mathfrak{c})}\right) & \text{if $g_{\mathrm{w}}$ is even},\\[3pt]
\overline{\Phi\left(\left.\left(g_{\mathrm{w}}^{*}u\right)\right|_{\mathrm{Pqp}(\mathfrak{c})}\right)} & \text{if $g_{\mathrm{w}}$ is odd},
\end{cases}
\qquad
\left(g_{\mathrm{w}}^{*}u\right)(p)=\epsilon(g_{\mathrm{w}},p)\cdot u(g_{\mathrm{w}}p).
\]
Then $\Phi^{g_{\mathrm{w}}}$ belongs to $\mathrm{Fcl}(\mathfrak{c}')$, with the same constants, and
\[
\left(\Phi^{g_{\mathrm{w}}}\right)^{g_{\mathrm{w}}'}=\Phi^{g_{\mathrm{w}}'g_{\mathrm{w}}},
\]
because $g_{\mathrm{w}}^{*}g_{\mathrm{w}}'^{*}=\left(g_{\mathrm{w}}'g_{\mathrm{w}}\right)^{*}$ (Section~I.4.7) and the parities multiply.
The copy $\Phi^{[\mathfrak{c}']}$ is $\Phi^{g_{\mathrm{w}}}$ for the element $g_{\mathrm{w}}$ of $W(\mathcal{Q})$ with $g_{\mathrm{w}}Q^{\circ}_{\mathfrak{c}}=Q^{\circ}_{\mathfrak{c}'}$.
On a torus $C(n)$ that carries $\mathcal{Q}$, an element $\Phi\in\mathrm{Fcl}(\mathfrak{c})$ is read as the function
\[
(A,\mathfrak{z})\mapsto\Phi\left(\left(\widetilde F_p\right)_{p\in\mathrm{Pqp}(\mathfrak{c})}\right),
\]
where $\widetilde F$ is the periodic extension of $F^{\mathrm{u}}$.

\emph{Representatives} (Definition~39.0).
Let $v\in\Gamma$.
The translation $\tau_v$ belongs to $W(\mathcal{Q})$ (1.5) and is even;
moreover $\left(\tau_v^{*}u\right)(p)=u(p+v)$ and $\widetilde F_{p+v}=\widetilde F_p$.
Hence $\Phi^{\tau_v}$, an element of $\mathrm{Fcl}(\mathfrak{c}+\mathrm{sh}(v))$, and $\Phi$ are read as the same function on $C(n)$.
The rule of composition and 1.4 give
\[
\left(\Phi^{\tau_v}\right)^{[\mathfrak{c}']}=\Phi^{[\mathfrak{c}']},
\qquad
\left(\Phi^{[\mathfrak{c}']}\right)^{\tau_v}=\Phi^{[\mathfrak{c}'+\mathrm{sh}(v)]} .
\]
Consequently the function on $C(n)$ defined by $\prod_{\mathfrak{c}}\Phi_{\mathfrak{c}}$ does not change when a factor is moved, by a translation of $\Gamma$, to another representative of its cell;
the function defined by $\Phi^{\sharp}_n$ does not depend on the representatives $\mathfrak{c}'$ over which the product is formed;
and $\left(\Phi^{\tau_v}\right)^{\sharp}_n=\Phi^{\sharp}_n$.
It is therefore enough to take the brick $\mathrm{Kb}$ of 2.4 as the set of representatives:
\[
\Phi^{\sharp}_n=\prod_{\mathfrak{c}'\in\mathrm{Kb}}\Phi^{[\mathfrak{c}']} .
\]

\emph{Scalar elements of $\mathcal{A}(Q^{\circ}_{\mathfrak{c}})$.}
The plaquettes contained in the open cell belong to $\mathrm{Pqp}(\mathfrak{c})$,
and the scalar elements of $\mathcal{A}(Q^{\circ}_{\mathfrak{c}})$ lie in $\mathrm{Fcl}(\mathfrak{c})$, with the same reading on $C(n)$ and the same copies (Proposition~39\,(f), last sentence),
in the following sense:
for every $P\in\mathcal{A}(Q^{\circ}_{\mathfrak{c}})$ whose values are scalars there is $\Phi_P\in\mathrm{Fcl}(\mathfrak{c})$ with
\[
P(\hat{A})=\Phi_P\left(\left((d\hat{A})_p\right)_{p\in\mathrm{Pqp}(\mathfrak{c})}\right)
\]
for all real one-cochains $\hat{A}$ on $\mathbb{Z}^3$.
\begin{proof}
An element of $\mathcal{A}(Q^{\circ}_{\mathfrak{c}})$ is a map $\hat{A}\mapsto P(\hat{A})$ from the real one-cochains on $\mathbb{Z}^3$ to the Grassmann algebra, namely a finite sum of products
\[
\hat{A}\mapsto\Phi'(d\hat{A})\cdot\prod_jG_j(\hat{A}),
\]
in which $G_j\in\mathrm{Ad}^{\mathrm{ev}}(\mathrm{Bx}_j)$ reads $\hat{A}$ on $\mathrm{Ln}(\mathrm{Bx}_j)$,
the boxes $\mathrm{Bx}_j$ of one product are pairwise disjoint and contained in $Q^{\circ}_{\mathfrak{c}}$,
and $\Phi'$ is a bounded Borel function of finitely many $(d\hat{A})_p$ with $p\subset Q^{\circ}_{\mathfrak{c}}$, which absorbs the coefficient of the product in the linear combination
\cite[Definitions~2.12 and~5.12]{QEDII}.

(1) The coefficient of the empty monomial is a linear and multiplicative functional on the Grassmann algebra.
Let $G^0_j(\hat{A})$ denote the coefficient of the empty monomial in $G_j(\hat{A})$.
If the values of $P$ are scalars, then $P(\hat{A})$ is the sum, over the products, of the numbers $\Phi'(d\hat{A})\cdot\prod_jG^0_j(\hat{A})$.

(2) $G^0_j$ is a Borel function on $\mathbb{R}^{\mathrm{Ln}(\mathrm{Bx}_j)}$ with
\[
\left|G^0_j(\hat{A})\right|\le\mathrm{Cq}_j\cdot\exp\left(\left(\frac12-\varepsilon_j\right)\sum_{p\in\mathrm{Pq}(\mathrm{Bx}_j)}(d\hat{A})_p^2\right),
\]
where $(\mathrm{Cq}_j,\varepsilon_j)$ are the constants of $G_j$ and $\varepsilon_j>0$.
Moreover $G^0_j(\hat{A}-d\boldsymbol{\lambda})=G^0_j(\hat{A})$ for every real function $\boldsymbol{\lambda}$ on the sites of $\mathrm{Bx}_j$:
a gauge transformation multiplies each monomial by a phase, and this phase is $1$ for the empty monomial.
So $G^0_j$ is a bosonic element of $\mathrm{Ad}(\mathrm{Bx}_j)$, and $G^0_j(\hat{A})$ depends on $d\hat{A}$ only \cite[Definition~2.12]{QEDII}.
On the linear subspace $d\,\mathbb{R}^{\mathrm{Ln}(\mathrm{Bx}_j)}$ of $\mathbb{R}^{\mathrm{Pq}(\mathrm{Bx}_j)}$, the function $u\mapsto G^0_j(\hat{A})$, $d\hat{A}=u$, is the composition of $G^0_j$ with a linear right inverse of $d$;
hence it is Borel.
Extended by $0$ to $\mathbb{R}^{\mathrm{Pq}(\mathrm{Bx}_j)}$, it is again denoted by $G^0_j$;
it is a Borel function with
\[
\left|G^0_j(u)\right|\le\mathrm{Cq}_j\cdot\exp\left(\left(\frac12-\varepsilon_j\right)\sum_{p\in\mathrm{Pq}(\mathrm{Bx}_j)}u_p^2\right).
\]

(3) The sets $\mathrm{Pq}(\mathrm{Bx}_j)$ of one product are pairwise disjoint subsets of $\mathrm{Pqp}(\mathfrak{c})$, and $\Phi'$ reads plaquettes of $\mathrm{Pqp}(\mathfrak{c})$.
Put $\varepsilon:=\min\left(\frac12,\min_j\varepsilon_j\right)$, so that $\frac12-\varepsilon\ge0$.
For $u\in\mathbb{R}^{\mathrm{Pqp}(\mathfrak{c})}$, with $G^0_j$ evaluated at the restriction of $u$ to $\mathrm{Pq}(\mathrm{Bx}_j)$, one has
\[
\left|\Phi'(u)\cdot\prod_jG^0_j(u)\right|\le\sup\left|\Phi'\right|\cdot\left(\prod_j\mathrm{Cq}_j\right)\cdot\exp\left(\left(\frac12-\varepsilon\right)\sum_{p\in\mathrm{Pqp}(\mathfrak{c})}u_p^2\right).
\]
Thus $u\mapsto\Phi'(u)\cdot\prod_jG^0_j(u)$ lies in $\mathrm{Fcl}(\mathfrak{c})$.
The class $\mathrm{Fcl}(\mathfrak{c})$ is a linear space: a sum has the sum of the constants $\mathrm{Cq}$ and the smallest of the exponents $\varepsilon$.
Let $\Phi_P$ be the sum of these functions over the products;
by (1) and (2) it has the asserted property.

(4) $\Phi_P$ depends on the representation of $P$ as a sum of products and on the extensions by $0$;
its values at the points $\left((d\hat{A})_p\right)_{p\in\mathrm{Pqp}(\mathfrak{c})}$ do not.
On $C(n)$ the boxes of one product and their translates by $\Gamma$ are pairwise disjoint, since $\mathrm{Bx}_j+v$ lies in $Q^{\circ}_{\mathfrak{c}+\mathrm{sh}(v)}$ (1.5).
Accordingly $P$ is read as the function $(A,\mathfrak{z})\mapsto P(\hat{A})$, where $\hat{A}$ is a real one-cochain on $\mathbb{Z}^3$ with $d\hat{A}=\widetilde F$
\cite[Definition~2.13(a) and Lemma~5.14(c)]{QEDII};
this function is $\Phi_P\left(\left(\widetilde F_p\right)_{p\in\mathrm{Pqp}(\mathfrak{c})}\right)$, the reading of $\Phi_P$.
One and the same $\hat{A}$ serves all cells.
Hence a product over the cells of scalar elements $P$ is read as the product of the functions $\Phi_P$, and $Z_n$ takes the same value at the two.
Finally let $g_{\mathrm{w}}\in W(\mathcal{Q})$ with $g_{\mathrm{w}}Q^{\circ}_{\mathfrak{c}}=Q^{\circ}_{\mathfrak{c}'}$.
Then $P^{g_{\mathrm{w}}}(\hat{A})=P(g_{\mathrm{w}}^{*}\hat{A})$ if $g_{\mathrm{w}}$ is even and $P^{g_{\mathrm{w}}}(\hat{A})=\overline{P(g_{\mathrm{w}}^{*}\hat{A})}$ if $g_{\mathrm{w}}$ is odd \cite[Section~2.3]{QEDII},
and $d\,g_{\mathrm{w}}^{*}\hat{A}=g_{\mathrm{w}}^{*}d\hat{A}$ (Section~I.4.7).
Therefore
\[
P^{g_{\mathrm{w}}}(\hat{A})=\left(\Phi_P\right)^{g_{\mathrm{w}}}\left(\left((d\hat{A})_p\right)_{p\in\mathrm{Pqp}(\mathfrak{c}')}\right),
\]
that is, the copy of $\Phi_P$ is a function $\Phi_{P^{g_{\mathrm{w}}}}$.
\end{proof}
The set $\mathrm{Pqp}(\mathfrak{c})$ is the set $P^{+}$ of the cell with the index $\mathfrak{c}$ in \cite[Definition~21.18]{QEDI},
because a plaquette whose vertices lie in the closed cell is contained in it.
A Borel function $\Phi$ of $(u_p)_{p\in\mathrm{Pqp}(\mathfrak{c})}$ with $|\Phi(u)|\le\mathrm{cf}\cdot\exp\left(\mathrm{cf}\cdot|u|\right)$ lies in $\mathrm{Fcl}(\mathfrak{c})$,
since $\mathrm{cf}\cdot|u|\le\mathrm{cf}^2+\frac{|u|^2}{4}$.

\begin{named}{Proposition 9.3}
Let $\mathrm{hw}\ge\frac32$, and let $C(n)$ carry $\mathcal{Q}$.
\begin{itemize}
\item[(i)]
The sets $\mathrm{Pqp}(\mathfrak{c})+v$, $\mathfrak{c}\in\mathrm{Kb}$, $v\in\Gamma$, are pairwise disjoint, and $p+v=p'+v'$ with $p,p'\in\mathrm{Pqp}(\mathfrak{c})$ implies $p=p'$ and $v=v'$.
For $\Phi_{\mathfrak{c}}\in\mathrm{Fcl}(\mathfrak{c})$, $\mathfrak{c}\in\mathrm{Kb}$, the product $\prod_{\mathfrak{c}}\Phi_{\mathfrak{c}}$ is a scalar element of $\mathrm{Ig}(T)$, $T=C(n)$,
and $Z_n\left[\prod_{\mathfrak{c}}\Phi_{\mathfrak{c}}\right]$ converges absolutely.
\item[(ii)]
$Z_n\left[(\Phi_{\mathfrak{c}})^{\sharp}_n\right]\ge0$, and
\[
\left|Z_n\left[\prod_{\mathfrak{c}\in\mathrm{Kb}}\Phi_{\mathfrak{c}}\right]\right|\le\prod_{\mathfrak{c}\in\mathrm{Kb}}Z_n\left[(\Phi_{\mathfrak{c}})^{\sharp}_n\right]^{\frac{1}{\mathrm{nc}(n)}} .
\]
\item[(iii)]
For $\Phi\in\mathrm{Fcl}(\mathfrak{c})$ and $\mathfrak{c}',\mathfrak{c}''\in\mathbb{Z}^3$:
$\left(\Phi^{[\mathfrak{c}']}\right)^{[\mathfrak{c}'']}=\Phi^{[\mathfrak{c}'']}$.
\item[(iv)]
For $\Phi\in\mathrm{Fcl}(0)$ and $\mathcal{Z}_{\mathrm{sv}}(\Phi):=Z_{\mathrm{hw}\cdot\mathrm{sv}}\left[\Phi^{\sharp}_{\mathrm{hw}\cdot\mathrm{sv}}\right]$ the assertions of Proposition~39\,(e) hold.
\item[(v)]
On cubic tori, $n\in2\mathrm{hw}\mathbb{Z}$, $n\ge4$:
\[
E^{\mathrm{af}}_n\left[(\Phi_{\mathfrak{c}})^{\sharp}_n\right]\ge0
\qquad\text{and}\qquad
\left|E^{\mathrm{af}}_n\left[\prod_{\mathfrak{c}\in\mathrm{Kb}}\Phi_{\mathfrak{c}}\right]\right|\le\prod_{\mathfrak{c}\in\mathrm{Kb}}E^{\mathrm{af}}_n\left[(\Phi_{\mathfrak{c}})^{\sharp}_n\right]^{\left(\frac{\mathrm{hw}}{n}\right)^{3}} .
\]
\end{itemize}
\end{named}

\begin{proof}
\emph{(i).}
Let $p\in\mathrm{Pqp}(\mathfrak{c})$, $p'\in\mathrm{Pqp}(\mathfrak{c}')$ and $v,v'\in\Gamma$ with $p+v=p'+v'$.
This plaquette is contained in $\mathrm{cl}\,Q^{\circ}_{\mathfrak{c}+\mathrm{sh}(v)}$ and in $\mathrm{cl}\,Q^{\circ}_{\mathfrak{c}'+\mathrm{sh}(v')}$.
Suppose that the two indices differ in the entry $\mu$.
Two different closed intervals $\left[o_\mu+i\cdot\mathrm{hw},\,o_\mu+(i+1)\cdot\mathrm{hw}\right]$ have at most one common point, which is an end point;
so the plaquette lies in a wall, and then so does $p$, because $\tau_{-v}$ maps walls to walls (1.5).
This contradicts $p\in\mathrm{Pqp}(\mathfrak{c})$.
Hence $\mathfrak{c}+\mathrm{sh}(v)=\mathfrak{c}'+\mathrm{sh}(v')$;
since $\mathfrak{c}$ and $\mathfrak{c}'$ both lie in $\mathrm{Kb}$, it follows that $\mathfrak{c}=\mathfrak{c}'$, $v=v'$ and $p=p'$.
Consequently a plaquette of $C(n)$ has a representative in at most one of the sets $\mathrm{Pqp}(\mathfrak{c})$, $\mathfrak{c}\in\mathrm{Kb}$, and at most one representative there, so that
\[
\left|\prod_{\mathfrak{c}}\Phi_{\mathfrak{c}}\right|\le\left(\prod_{\mathfrak{c}}\mathrm{Cq}_{\mathfrak{c}}\right)\cdot\exp\left(\left(\frac12-\min_{\mathfrak{c}}\varepsilon_{\mathfrak{c}}\right)\left\|F^{\mathrm{u}}\right\|^2\right).
\]
The product is a Borel function of $F^{\mathrm{u}}$, hence invariant under $(A,\mathfrak{z})\mapsto(A+r_q\cdot m,\mathfrak{z}-dm)$.
By (In7)\,(iv) it is a scalar element of $\mathrm{Ig}(T)$, and $Z_n$ of it converges absolutely.

\emph{(iii).}
Let $g_{\mathrm{w},1},g_{\mathrm{w},2}\in W(\mathcal{Q})$ with $g_{\mathrm{w},1}Q^{\circ}_{\mathfrak{c}}=Q^{\circ}_{\mathfrak{c}'}$ and $g_{\mathrm{w},2}Q^{\circ}_{\mathfrak{c}'}=Q^{\circ}_{\mathfrak{c}''}$.
Then
\[
\left(\Phi^{[\mathfrak{c}']}\right)^{[\mathfrak{c}'']}=\left(\Phi^{g_{\mathrm{w},1}}\right)^{g_{\mathrm{w},2}}=\Phi^{g_{\mathrm{w},2}g_{\mathrm{w},1}},
\]
and $g_{\mathrm{w},2}g_{\mathrm{w},1}$ is the element of $W(\mathcal{Q})$ that maps $Q^{\circ}_{\mathfrak{c}}$ onto $Q^{\circ}_{\mathfrak{c}''}$ (1.4).

\emph{One direction.}
A \emph{family} is $\mathrm{GG}=(\Phi_{\mathfrak{c}})_{\mathfrak{c}\in\mathrm{Kb}}$ with $\Phi_{\mathfrak{c}}\in\mathrm{Fcl}(\mathfrak{c})$.
Its product is $\mathrm{Fp}(\mathrm{GG}):=\prod_{\mathfrak{c}}\Phi_{\mathfrak{c}}$, and its folds are the families
\[
\left(\mathrm{Rm}^j_\mu\mathrm{GG}\right)_{\mathfrak{c}}:=\left(\Phi_{\mathfrak{c}[\mu:j]}\right)^{\mathrm{di}^{\mu}_{\mathfrak{c}_\mu,j}}\in\mathrm{Fcl}(\mathfrak{c}).
\]
For these families the inequality (7.1) holds in each direction $\mu$.

\emph{The direction $\mu=1$.}
Take the notation of 6.2 with $\mathrm{po}_j:=o_1+j\cdot\mathrm{hw}$, $\mathrm{hsl}:=\mathrm{hw}$ and $2k:=\mathrm{sv}_1$;
the vector of the slabs $[\mathrm{po}_j,\mathrm{po}_{j+1}]$ is then $2\mathrm{hw}\cdot e_1$.
For $\omega\in\mathrm{Cfg}[\mathrm{po}_j,\mathrm{po}_{j+1}]$ put
\[
G_j(\omega):=\prod_{\mathfrak{c}\in\mathrm{Kb},\ \mathfrak{c}_1=j}\Phi_{\mathfrak{c}}\left(\left(F^{\mathrm{u}}_p(\omega)\right)_{p\in\mathrm{Pqp}(\mathfrak{c})}\right).
\]
A plaquette $p\in\mathrm{Pqp}(\mathfrak{c})$ with $\mathfrak{c}_1=j$ is contained in $\{\mathrm{po}_j\le y_1\le\mathrm{po}_{j+1}\}$ and in neither of the two planes, which are walls;
so its class modulo $\Gamma_{\mathrm{par}}=n_2\mathbb{Z}e_2+n_3\mathbb{Z}e_3$ lies in $\mathrm{Sb}^2(\mathrm{po}_j,\mathrm{po}_{j+1})$,
and by (i) different pairs $(\mathfrak{c},p)$ give different classes.
Therefore $G_j$ is well defined, Borel, and a function of $F^{\mathrm{u}}(\omega)$, which gives (Ig1); moreover
\[
\left|G_j(\omega)\right|\le\mathrm{const}\cdot\exp\left(\left(\frac12-\min_{\mathfrak{c}}\varepsilon_{\mathfrak{c}}\right)\sum_{p\in\mathrm{Sb}^2(\mathrm{po}_j,\mathrm{po}_{j+1})}F^{\mathrm{u}}_p(\omega)^2\right).
\]
Thus $G_j$ is a scalar element of $\mathrm{Ig}^{\mathrm{ev}}_{\mathrm{Xs}}[\mathrm{po}_j,\mathrm{po}_{j+1}]$.
Since $F^{\mathrm{u}}(\omega_j)$ is the restriction of $\widetilde F$, the product $G_0\cdots G_{2k-1}$ of (Tc3) is $\mathrm{Fp}(\mathrm{GG})$.
For a scalar $G$ one has $(\mathrm{Tg}_vG)(\omega')=G(\tau_v^{*}\omega')$ and $(\Theta_{\mathrm{po}'}G)(\omega')=\overline{G(\theta_{\mathrm{po}'}^{*}\omega')}$ ((In7)\,(iii)),
and $F^{\mathrm{u}}(g_{\mathrm{w}}^{*}\omega')=g_{\mathrm{w}}^{*}F^{\mathrm{u}}(\omega')$.
Hence, for $i-j=2s$,
\[
(\mathrm{cp}_iG_j)(\omega')=G_j\left(\tau_{2s\cdot\mathrm{hw}\cdot e_1}^{*}\omega'\right),
\]
and, for $i-j=2s+1$,
\[
(\mathrm{cp}_iG_j)(\omega')=\overline{G_j\left(\theta_{\mathrm{po}_{j+1}}^{*}\tau_{2s\cdot\mathrm{hw}\cdot e_1}^{*}\omega'\right)}=\overline{G_j\left(\left(\tau_{2s\cdot\mathrm{hw}\cdot e_1}\theta_{\mathrm{po}_{j+1}}\right)^{*}\omega'\right)} .
\]
By (1.1) this means
\[
(\mathrm{cp}_iG_j)(\omega')=\prod_{\mathfrak{c}\in\mathrm{Kb},\ \mathfrak{c}_1=j}\left(\Phi_{\mathfrak{c}}\right)^{\mathrm{di}^1_{i,j}}\left(F^{\mathrm{u}}(\omega')\right).
\]
Hence $(\mathrm{cp}_0G_j)\cdots(\mathrm{cp}_{2k-1}G_j)=\mathrm{Fp}\left(\mathrm{Rm}^j_1\mathrm{GG}\right)$, and Proposition~6.3 becomes (7.1) for $\mu=1$.

\emph{The directions $\mu=2,3$.}
Let $g_{\mathrm{w}}=\mathrm{wr}_\mu^{-1}$, in the notation of 7.1.
By (r3), $(\Phi_{\mathfrak{c}})^{g_{\mathrm{w}}}$ lies in the class $\mathrm{Fcl}(\mathfrak{c}^{\mathrm{pm}})$ formed with $\mathcal{Q}^{\mathrm{pm}}$;
as $g_{\mathrm{w}}$ is even, $(\Phi_{\mathfrak{c}})^{g_{\mathrm{w}}}(u')=\Phi_{\mathfrak{c}}(g_{\mathrm{w}}^{*}u')$.
Let $\mathrm{GG}^{\mathrm{pm}}$ be the family on $(C(n^{\mathrm{pm}}),\mathcal{Q}^{\mathrm{pm}})$ with the members $(\mathrm{GG}^{\mathrm{pm}})_{\mathfrak{c}^{\mathrm{pm}}}:=(\Phi_{\mathfrak{c}})^{g_{\mathrm{w}}}$.
Then $\mathrm{Fp}(\mathrm{GG}^{\mathrm{pm}})(u')=\mathrm{Fp}(\mathrm{GG})(g_{\mathrm{w}}^{*}u')$.
By (r3) and the rule $\left(\Phi^{g_{\mathrm{w}}}\right)^{g_{\mathrm{w}}'}=\Phi^{g_{\mathrm{w}}'g_{\mathrm{w}}}$,
\[
\left(\mathrm{Rm}^j_1\mathrm{GG}^{\mathrm{pm}}\right)_{\mathfrak{c}^{\mathrm{pm}}}
=\left(\left(\Phi_{\mathfrak{c}[\mu:j]}\right)^{g_{\mathrm{w}}}\right)^{g_{\mathrm{w}}\mathrm{di}^{\mu}_{\mathfrak{c}_\mu,j}g_{\mathrm{w}}^{-1}}
=\left(\Phi_{\mathfrak{c}[\mu:j]}\right)^{g_{\mathrm{w}}\mathrm{di}^{\mu}_{\mathfrak{c}_\mu,j}}
=\left(\left(\mathrm{Rm}^j_\mu\mathrm{GG}\right)_{\mathfrak{c}}\right)^{g_{\mathrm{w}}};
\]
so $\mathrm{Fp}\left(\mathrm{Rm}^j_1\mathrm{GG}^{\mathrm{pm}}\right)(u')=\mathrm{Fp}\left(\mathrm{Rm}^j_\mu\mathrm{GG}\right)(g_{\mathrm{w}}^{*}u')$.
The element $g_{\mathrm{w}}$ is a product $\theta_{m_1}\theta_{m_2}$ of the reflections in two diagonal mirrors through the origin \cite[Proposition~5.47]{QEDII}.
For a scalar $P\in\mathrm{Ig}(T)$ one has $(\Theta_mP)(A',\mathfrak{z}')=\overline{P\left(\theta_m^{*}(A',\mathfrak{z}')\right)}$;
hence
\[
(\Theta_{m_1}\Theta_{m_2}P)(A',\mathfrak{z}')=P\left(\theta_{m_2}^{*}\theta_{m_1}^{*}(A',\mathfrak{z}')\right)=P\left(g_{\mathrm{w}}^{*}(A',\mathfrak{z}')\right),
\]
and \cite[Lemma~5.46]{QEDII} yields
\[
Z_{n^{\mathrm{pm}}}\left[P\circ g_{\mathrm{w}}^{*}\right]=Z_n[P] .
\]
Moreover $F^{\mathrm{u}}\left(g_{\mathrm{w}}^{*}(A',\mathfrak{z}')\right)=g_{\mathrm{w}}^{*}F^{\mathrm{u}}(A',\mathfrak{z}')$.
With $P=\mathrm{Fp}(\mathrm{GG})$ and with $P=\mathrm{Fp}\left(\mathrm{Rm}^j_\mu\mathrm{GG}\right)$, regarded as functions of $F^{\mathrm{u}}$, this gives
\[
Z_{n^{\mathrm{pm}}}\left[\mathrm{Fp}(\mathrm{GG}^{\mathrm{pm}})\right]=Z_n\left[\mathrm{Fp}(\mathrm{GG})\right],
\qquad
Z_{n^{\mathrm{pm}}}\left[\mathrm{Fp}\left(\mathrm{Rm}^j_1\mathrm{GG}^{\mathrm{pm}}\right)\right]=Z_n\left[\mathrm{Fp}\left(\mathrm{Rm}^j_\mu\mathrm{GG}\right)\right].
\]
The direction $1$ for $(C(n^{\mathrm{pm}}),\mathcal{Q}^{\mathrm{pm}},\mathrm{GG}^{\mathrm{pm}})$ is therefore (7.1) for the direction $\mu$.

\emph{(ii).}
Apply (7.1) with $\mu=1$ to $\mathrm{GG}$, with $\mu=2$ to each $\mathrm{Rm}^{j_1}_1\mathrm{GG}$, and with $\mu=3$ to each $\mathrm{Rm}^{j_2}_2\mathrm{Rm}^{j_1}_1\mathrm{GG}$;
the left sides of the second and of the third application are non-negative numbers.
The result is
\[
\left|Z_n\left[\mathrm{Fp}(\mathrm{GG})\right]\right|\le\prod_{j_1,j_2,j_3}Z_n\left[\mathrm{Fp}\left(\mathrm{Rm}^{j_3}_3\mathrm{Rm}^{j_2}_2\mathrm{Rm}^{j_1}_1\mathrm{GG}\right)\right]^{\frac{1}{\mathrm{sv}_1\cdot\mathrm{sv}_2\cdot\mathrm{sv}_3}},
\]
with all factors $\ge0$.
By the rule $\left(\Phi^{g_{\mathrm{w}}}\right)^{g_{\mathrm{w}}'}=\Phi^{g_{\mathrm{w}}'g_{\mathrm{w}}}$ and (1.2), with $\mathrm{jv}=(j_1,j_2,j_3)$,
\[
\left(\mathrm{Rm}^{j_3}_3\mathrm{Rm}^{j_2}_2\mathrm{Rm}^{j_1}_1\mathrm{GG}\right)_{\mathfrak{c}}
=\left(\Phi_{\mathrm{jv}}\right)^{\mathrm{di}^3_{\mathfrak{c}_3,j_3}\mathrm{di}^2_{\mathfrak{c}_2,j_2}\mathrm{di}^1_{\mathfrak{c}_1,j_1}}
=\left(\Phi_{\mathrm{jv}}\right)^{[\mathfrak{c}]} .
\]
Hence $\mathrm{Fp}\left(\mathrm{Rm}^{j_3}_3\mathrm{Rm}^{j_2}_2\mathrm{Rm}^{j_1}_1\mathrm{GG}\right)=(\Phi_{\mathrm{jv}})^{\sharp}_n$, and $\mathrm{sv}_1\cdot\mathrm{sv}_2\cdot\mathrm{sv}_3=\mathrm{nc}(n)$.

\emph{(iv).}
The direction $\mu=1$.
Put
\[
G(\omega):=\prod_{0\le\mathfrak{c}_2<\mathrm{sv}_2,\ 0\le\mathfrak{c}_3<\mathrm{sv}_3}\Phi^{[(0,\mathfrak{c}_2,\mathfrak{c}_3)]}\left(F^{\mathrm{u}}(\omega)\right)\in\mathrm{Ig}^{\mathrm{ev}}_{\mathrm{Xs}}[o_1,o_1+\mathrm{hw}],
\]
$Y:=\mathrm{Co}[G]$ and $\mathrm{Bp}:=Y^{*}Y\ge0$, an operator of trace class.
These objects do not depend on $\mathrm{sv}_1$, because the definitions of 4.1 and 5.1 do not involve $n_1$.
By the computation for the direction $\mu=1$ above and by (iii),
\[
(\mathrm{cp}_iG)(\omega')=\prod_{0\le\mathfrak{c}_2<\mathrm{sv}_2,\ 0\le\mathfrak{c}_3<\mathrm{sv}_3}\Phi^{[(i,\mathfrak{c}_2,\mathfrak{c}_3)]}\left(F^{\mathrm{u}}(\omega')\right);
\]
so $(\mathrm{cp}_0G)\cdots(\mathrm{cp}_{\mathrm{sv}_1-1}G)=\Phi^{\sharp}_n$, and (6.1) gives
\[
\mathcal{Z}_{\mathrm{sv}}(\Phi)=\operatorname{Tr}\mathrm{Bp}^{\frac{\mathrm{sv}_1}{2}}
\]
for all even $\mathrm{sv}_1$ with $\mathrm{hw}\cdot\mathrm{sv}_1\ge4$.
Consequently $\mathcal{Z}_{\mathrm{sv}}(\Phi)^{\frac{1}{\mathrm{sv}_1}}=\left(\left\|\mathrm{Bp}\right\|_{\frac{\mathrm{sv}_1}{2}}\right)^{\frac12}$ is non-increasing in $\mathrm{sv}_1$ ((In9)\,(b));
$\mathcal{Z}_{\mathrm{sv}}(\Phi)>0$ if and only if $\mathrm{Bp}\ne0$;
and in that case $\log\mathcal{Z}_{\mathrm{sv}}(\Phi)=\log\operatorname{Tr}\mathrm{Bp}^{\frac{\mathrm{sv}_1}{2}}$ is convex in $\mathrm{sv}_1$ ((In9)\,(c)).
The directions $\mu=2,3$.
By (r3) and the rule for copies, $\Phi^{\sharp}_n(g_{\mathrm{w}}^{*}u')=\left(\Phi^{g_{\mathrm{w}}}\right)^{\sharp}_{n^{\mathrm{pm}}}(u')$, the right side formed with $\mathcal{Q}^{\mathrm{pm}}$.
By the identity $Z_{n^{\mathrm{pm}}}\left[P\circ g_{\mathrm{w}}^{*}\right]=Z_n[P]$ with $P=\Phi^{\sharp}_n$,
the number $\mathcal{Z}_{\mathrm{sv}}(\Phi)$ formed with $\mathcal{Q}$ equals the number $\mathcal{Z}_{\mathrm{sv}^{\mathrm{pm}}}(\Phi^{g_{\mathrm{w}}})$ formed with $\mathcal{Q}^{\mathrm{pm}}$,
and the first entry of $\mathrm{sv}^{\mathrm{pm}}$ is $\mathrm{sv}_\mu$.

\emph{(v).}
One has $\mathrm{nc}(n)=\left(\frac{n}{\mathrm{hw}}\right)^3$ and $0<Z_n[1]<\infty$ \cite[Theorem~4.26(a)]{QEDII}.
Divide (ii) by $Z_n[1]=\prod_{\mathfrak{c}\in\mathrm{Kb}}Z_n[1]^{\frac{1}{\mathrm{nc}(n)}}$.
\end{proof}

\begin{named}{Remark 9.4}[the closed-cell form of the chessboard estimate and of the lengthening of one side]
Let $o=0$ and $\mathrm{hw}\ge2$.
The cell structures of \cite[Definition~21.18]{QEDI} are of this form, the directions $0,1,2$ there being the directions $1,2,3$ here (Section~I.9.2).
By the last assertion of 9.2, the observables of the cell $\mathfrak{c}$ in \cite[Definition~21.18]{QEDI},
that is, the functions of the class $\mathfrak{F}(n)$ that depend only on the $F^{\mathrm{u}}_p$ with $p$ in the set $P^{+}$ of this cell, which is $\mathrm{Pqp}(\mathfrak{c})$,
lie in $\mathrm{Fcl}(\mathfrak{c})$.
By (1.1) and (1.2), the map $\mathsf{g}_{\mathfrak{c}}$ of \cite[Definition~21.18]{QEDI} is the element of $W(\mathcal{Q})$ that maps $Q^{\circ}_0$ onto $Q^{\circ}_{\mathfrak{c}}$,
and it is even if and only if $\sum_\mu\mathfrak{c}_\mu$ is even;
so the mirror copies and the product $O^{\sharp}$ of \cite[Definition~21.18]{QEDI} are those of 9.2.
For $O_{\mathfrak{c}}\in\mathrm{Fcl}(0)$ put $\Phi_{\mathfrak{c}}:=(O_{\mathfrak{c}})^{[\mathfrak{c}]}$.
By Proposition~9.3\,(iii), $(\Phi_{\mathfrak{c}})^{\sharp}_n=(O_{\mathfrak{c}})^{\sharp}_n$,
and Proposition~9.3\,(ii) is the inequality of \cite[Proposition~21.22]{QEDI} for the all-flux functional, with $Z_n[\,\cdot\,]$ in the place of $\boldsymbol{\omega}_n$.
If $\mathrm{sv}'$ arises from $\mathrm{sv}$ by multiplying the entry $\mathrm{sv}_\mu$ by an integer $j\ge1$,
the first assertion of Proposition~9.3\,(iv) states that
\[
\mathcal{Z}_{\mathrm{sv}'}(\Phi)^{\frac{1}{j\cdot\mathrm{sv}_\mu}}\le\mathcal{Z}_{\mathrm{sv}}(\Phi)^{\frac{1}{\mathrm{sv}_\mu}},
\qquad\text{that is,}\qquad
0\le\mathcal{Z}_{\mathrm{sv}'}(\Phi)\le\mathcal{Z}_{\mathrm{sv}}(\Phi)^{j} :
\]
this is the inequality of \cite[Proposition~21.24(a)]{QEDI}.
\end{named}

\subsection*{11. Conditions on the constants}

\begin{itemize}
\item[(K-1)]
$0<\rho<1$: Section~I.0.5, and the line (C0) of Section~I.8.2, that is, of Proposition~14.2\,(a).
It is used in (In2), 4.1 and (m3): $c_i=3(1\pm\rho)>0$, so that $\mathrm{cc}_{\mathrm{po}}^{\frac12}$ is a real matrix.
(The condition $-1\le\rho\le1$ suffices.)
\item[(K-2)]
The standard weights: line (C0).
Of these weights the arguments of this chapter use the following.
The closure relation $\sum_iw_i\hat n_i=0$ (I.1.1), which, through (I.2.5), makes the diagonal term of $D_A+\boldsymbol{\mu}$ the real scalar $M_{\mathrm{ch}}=\boldsymbol{\mu}+\frac{\rho}{2}\cdot\mathrm{sw}$ ((In4), (In6)).
The values $w_0=w_3=6>0$:
in the hopping matrices $3(\rho\mp\gamma_1)$ of the crossing bonds and in $c_i=\frac{w_0}{2}(1\pm\rho)=\frac{w_3}{2}(1\pm\rho)=3(1\pm\rho)$ ((In1)--(In3), 4.1, (m3), (m5)),
and, together with $\rho\le1$, in $c_I\le6^{64v_s}$ ((m1)).
Beyond this, $w_1$ and $w_2$ enter only as real factors of the hopping matrices $M_1(x)$, $M_2(x)$,
whose entries are coefficients of $S_{(\mathrm{po},\mathrm{po}')}(A)$ that do not depend on $A$ ((In4), (In6), (m1)).
The cited statements have their own hypotheses on the weights:
\cite[Definition~2.10]{QEDII} and all of \cite[Section~5]{QEDII} assume $w_i>0$, $i=0,\dots,3$,
and so does \cite[Theorem~4.26]{QEDII}, on which Proposition~9.3\,(v) rests.
\end{itemize}
Of the constants of Section~I.8, only $\rho$ and the standard weights enter the definitions, statements and proofs of this chapter,
and no condition on them beyond (K-1) and (K-2) is required.
The lattice mass $\boldsymbol{\mu}$ is an arbitrary real number (I.K10).

\section*{Chapter 40. Transfer to all tori: proofs of Proposition~40}

The statement proved in this chapter is Proposition~40 of Part~II.

\subsection*{Proof of Proposition~40}

\subsubsection*{0. Conventions of this chapter: setting, notation, the results used}

\paragraph{0.1 Setting.}
Throughout the chapter lengths are measured in lattice units (Sections~I.2.5 and~I.7.1), and the standing data of Section~G6.A1 are fixed: $\rho\in(0,1)$ and the standard weights.
The lattice mass $\boldsymbol{\mu}$ is an arbitrary real number, the charge satisfies $q>0$, and $M_{\mathrm{ch}}=\boldsymbol{\mu}+\frac{\rho}{2}\cdot\mathrm{sw}$.
Let $n=(n_1,n_2,n_3)$ with integers $n_\mu\ge4$.
Then $C(n)$ denotes the torus of Sections~I.6.6 and~I.7.1, and $Z_n[\,\cdot\,]$ the functional of Section~I.7.2.
On Borel functions of $F^{\mathrm{u}}$ this functional is given by the expression (I.7.1).
On the elements of $\mathcal{A}$ that are read on $C(n)$ it is the functional $Z_{\mathbb{T}}$ of \cite[Definition~2.13]{QEDII} for $\mathbb{T}=C(n)$, formed with $M_{\mathrm{ch}}$ as the mass parameter $M$ of \cite[Definition~2.10]{QEDII}; this is the last sentence of Definition~39.0.
On both classes one has $E^{\mathrm{af}}_n=Z_n[\,\cdot\,]/Z_n[1]$: on the first by (I.7.1), on the second by the definition $\mathbb{E}_{\mathbb{T}}=Z_{\mathbb{T}}[\,\cdot\,]/Z_{\mathbb{T}}[1]$ of the expectation in \cite[Definition~2.13]{QEDII}.
A single letter $n$ stands for the cubic torus $(n,n,n)$.

The fermion enters Proposition~40 and its proof only through the numbers $Z_n[\,\cdot\,]$, and these enter only through the facts collected in 0.3.
All objects live at level~$0$, in lattice units.

\paragraph{0.2 Notation; powers.}
In this chapter $\mathrm{Kr}$ and $\mathrm{Krv}$ denote sets of ratios, and $\mathrm{kr}_1$, $\mathrm{kr}_2$, $\mathrm{kr}_3$, $\mathrm{kr}'$, $\mathrm{kr}''$, $\mathrm{kr}'''$ denote ratios.
The letter $\mathrm{fz}$ stands for a real function of three ratios, $\mathrm{wr}$ for a real function of one ratio, and $\mathrm{wz}$ for a real function on the integers $\ge4$.
The number $(\mathrm{kr}-1)/(\check{j}-1)$ is denoted by $\alpha$.
Real numbers $\mathrm{bz},\mathrm{bz}'\ge0$ serve as bases of powers, and real numbers $\mathrm{xe},\mathrm{xe}'$ as exponents.
The maps $g_{\mathrm{w},1}$, $g_{\mathrm{w},2}$, $g_{\mathrm{w},3}$ of the three coordinates that make up an element $g_{\mathrm{w}}$ of $W(\mathcal{Q})$ are introduced in 1.5.
In Section~7 the letters $\mathrm{Og}_{\mathfrak{c}}$ denote a family of cell contents, in Remark~7.5 the letter $\mathrm{cz}$ denotes a complex number, and in Section~8 the letter $\mathrm{Sd}$ denotes a set of sides.
The letters $i$, $i'$, $i''$, $j$ denote integers.

\emph{Powers.}
For $\mathrm{bz}>0$ and every real $\mathrm{xe}$ put $\mathrm{bz}^{\mathrm{xe}}:=\exp\left(\mathrm{xe}\cdot\log\mathrm{bz}\right)$, and for $\mathrm{xe}>0$ put $0^{\mathrm{xe}}:=0$.
When $\mathrm{xe}$ is a positive integer, $\mathrm{bz}^{\mathrm{xe}}$ is the product of $\mathrm{xe}$ factors equal to $\mathrm{bz}$.
The following rules are used.
\begin{description}
\item[(pw1)] For fixed $\mathrm{xe}>0$ the function $\mathrm{bz}\mapsto\mathrm{bz}^{\mathrm{xe}}$ is non-decreasing on $[0,\infty)$.
\item[(pw2)] $\left(\mathrm{bz}^{\mathrm{xe}}\right)^{\mathrm{xe}'}=\mathrm{bz}^{\mathrm{xe}\cdot\mathrm{xe}'}$ holds for $\mathrm{bz}\ge0$ and $\mathrm{xe},\mathrm{xe}'>0$, and also for $\mathrm{bz}>0$ and arbitrary real $\mathrm{xe},\mathrm{xe}'$; moreover $\mathrm{bz}^1=\mathrm{bz}$.
\item[(pw3)] $\left(\mathrm{bz}\cdot\mathrm{bz}'\right)^{\mathrm{xe}}=\mathrm{bz}^{\mathrm{xe}}\cdot{\mathrm{bz}'}^{\mathrm{xe}}$ for $\mathrm{bz},\mathrm{bz}'\ge0$ and $\mathrm{xe}>0$.
\item[(pw4)] For $\mathrm{bz}>0$ and all real $\mathrm{xe},\mathrm{xe}'$ one has $\mathrm{bz}^0=1$, $\mathrm{bz}^{\mathrm{xe}}\cdot\mathrm{bz}^{\mathrm{xe}'}=\mathrm{bz}^{\mathrm{xe}+\mathrm{xe}'}$, $\mathrm{bz}^{-\mathrm{xe}}=\frac{1}{\mathrm{bz}^{\mathrm{xe}}}=\left(\frac{1}{\mathrm{bz}}\right)^{\mathrm{xe}}$ and $\log\left(\mathrm{bz}^{\mathrm{xe}}\right)=\mathrm{xe}\cdot\log\mathrm{bz}$.
\end{description}
For positive bases these rules are the functional equations and the monotonicity of $\exp$ and $\log$.
If a base vanishes, both sides of (pw2) and of (pw3) are equal to $0$, and (pw1) holds because $0^{\mathrm{xe}}=0\le\mathrm{bz}^{\mathrm{xe}}$.
Exponents $\le0$ occur only with the bases $Z_{[\mathrm{sv}]}$ and $Z_{[\mathrm{ncs}]}$, which are positive by 1.4.

\paragraph{0.3 The results used.}
The proof rests on the following facts.
\begin{description}
\item[{[Q1]}] Proposition~39\,(b): if $C(n)$ carries $\mathcal{Q}$ and $O_{\mathfrak{c}}\in\mathcal{A}(Q^{\circ}_{\mathfrak{c}}(\mathcal{Q}))$ for every cell $\mathfrak{c}$ of $C(n)$, then the product $\prod_{\mathfrak{c}}O_{\mathfrak{c}}$ is read on $C(n)$, and $Z_n$ of this product converges absolutely;
moreover, by the last sentence of that clause, every element of $\mathcal{A}(Q^{\circ}_0(\mathcal{Q}))$ is an integrand on $\mathbb{R}^3$, so that one and the same $O$ is read on all tori that carry $\mathcal{Q}$.
\item[{[Q2]}] Proposition~39\,(c), both assertions, with the number $\mathrm{nc}(n)$ of Proposition~39\,(a): in the situation of Proposition~39\,(b),
\[
Z_n\left[(O_{\mathfrak{c}})^{\sharp}_n\right]\ge0\quad\text{for every }\mathfrak{c},
\qquad
\left|Z_n\left[\prod_{\mathfrak{c}}O_{\mathfrak{c}}\right]\right|\le\prod_{\mathfrak{c}}Z_n\left[(O_{\mathfrak{c}})^{\sharp}_n\right]^{\frac{1}{\mathrm{nc}(n)}}.
\]
\item[{[Q3]}] Proposition~39\,(d): for $O\in\mathcal{A}(Q^{\circ}_0(\mathcal{Q}))$ and every cell $\mathfrak{c}$ one has $\left(O^{[\mathfrak{c}]}\right)^{[\mathfrak{c}']}=O^{[\mathfrak{c}']}$ for all $\mathfrak{c}'$, and hence $\left(O^{[\mathfrak{c}]}\right)^{\sharp}_n=O^{\sharp}_n$.
\item[{[Q4]}] Proposition~39\,(e), namely the definition
\[
\mathcal{Z}_{\mathrm{sv}}(O):=Z_{\mathrm{hw}\cdot\mathrm{sv}}\left[O^{\sharp}_{\mathrm{hw}\cdot\mathrm{sv}}\right]
\]
for $\mathrm{sv}=(\mathrm{sv}_1,\mathrm{sv}_2,\mathrm{sv}_3)$ with even positive integers $\mathrm{sv}_\mu$ such that $\mathrm{hw}\cdot\mathrm{sv}_\mu\ge4$ and for $O\in\mathcal{A}(Q^{\circ}_0(\mathcal{Q}))$, together with the monotonicity:
for each $\mu$, the two other entries being fixed, the map $\mathrm{sv}_\mu\mapsto\mathcal{Z}_{\mathrm{sv}}(O)^{\frac{1}{\mathrm{sv}_\mu}}$ is non-increasing on $\{\mathrm{sv}_\mu\text{ even},\ \mathrm{hw}\cdot\mathrm{sv}_\mu\ge4\}$.
\item[{[Q5]}] Proposition~50\,(e), the part $0<Z_n[1]<\infty$ of its chain of inequalities, for integers $n_\mu\ge4$ and every real $\boldsymbol{\mu}$.
\item[{[Q6]}] \cite[Proposition~6.19(a),(b)]{QEDII}: on $\{n\in\mathbb{Z}^3:\ n_\mu\ge4\}$, with $f(n)=|n|^{-1}\log Z_{C(n)}[1]$ as in \cite[Definition~2.21]{QEDII},
(a) $f$ is non-increasing in each $n_\mu$, and
(b) $n_\mu\mapsto\log Z_{C(n)}[1]$ is convex.
The standing hypotheses of \cite[Section~6]{QEDII} are $-1\le\rho\le1$, $w_i>0$, a real mass parameter, $q>0$ and $v_0\in\mathrm{Chm}_0$.
They are satisfied here with the mass parameter $M_{\mathrm{ch}}$ and with an arbitrary $v_0\in\mathrm{Chm}_0$, and $Z_{C(n)}[1]$ does not depend on $v_0$, by Proposition~46\,(e) and 1.3.
The integers $n_\mu\ge4$ may be of either parity, and $C(n)$ has $\chi=-1$ on the three generators, as in Section~I.6.6.
\item[{[Q7]}] \cite[Definition~2.12]{QEDII}: the constant $1$ lies in $\mathcal{A}(U)$ for every open subset $U$ of $\mathbb{R}^3$, being the product $\Phi\cdot G_1\cdots G_r$ with $r=0$ factors $G_j$ and $\Phi=1$.
The last sentence of the paragraph \emph{Theta} of Section~I.4.8: for a function $G$ of the curvature and $g_{\mathrm{w}}\in W_{\mathrm{af}}$ one has $G^{\mathrm{gw}}(F^{\mathrm{u}})=G(g_{\mathrm{w}}^*F^{\mathrm{u}})$ if $g_{\mathrm{w}}$ is even and $G^{\mathrm{gw}}(F^{\mathrm{u}})=\overline{G(g_{\mathrm{w}}^*F^{\mathrm{u}})}$ if $g_{\mathrm{w}}$ is odd.
\cite[Lemma~6.4(b)]{QEDII}: if $G\in\mathcal{A}(U)$, then $G^{\mathrm{gw}}\in\mathcal{A}(g_{\mathrm{w}}U)$, for every $g_{\mathrm{w}}\in W_{\mathrm{af}}$.
\item[{[Q8]}] Proposition~46\,(f) at $\Phi=1$;
the first row of the table in Section~I.7.2;
the second and the third sentence of Proposition~46\,(e), for $n_\mu\ge4$;
\cite[Definition~2.11]{QEDII}: $\int_{\mathbb{T}}\exp\left(-S_{\mathbb{T}}(A)\right)=\det D^{\mathbb{T}}_A$, where $D^{\mathbb{T}}_A$ is the operator $D_A$ of \cite[Definition~2.10]{QEDII} on the spinors of the torus $\mathbb{T}$, here $\mathbb{T}=C(n)$;
and \cite[Definition~2.13]{QEDII}: the definitions of $Z_{\mathbb{T}}$ on $\mathcal{A}$ and on non-negative Borel functions $\Phi$ of the curvature agree where both apply.
\end{description}
Definition~39.0 belongs to the data of Proposition~40.
Of Proposition~39\,(a) the third sentence enters as well: $W(\mathcal{Q})$ acts simply transitively on the cells; a proof of this is given in 1.5.
Section~8 rests in addition on the definitions of Sections~I.0.5, I.6.1, I.6.2, I.6.4, I.6.5 and~I.7.6, and Remark~8.5 on the inequality $L>8$, which Proposition~14.2\,(f) derives from (L1).

\subsubsection*{1. Tori, cells, the constant 1, positivity, the group of the walls, mirror copies}

\paragraph{1.1 Tori and cells.}
Let $\mathrm{sv}_1,\mathrm{sv}_2,\mathrm{sv}_3$ be even positive integers with $\mathrm{hw}\cdot\mathrm{sv}_\mu\ge4$, and put $n_\mu:=\mathrm{hw}\cdot\mathrm{sv}_\mu$.
As $2\,\mathrm{hw}$ and $\mathrm{sv}_\mu/2$ are integers, so is $n_\mu=(2\,\mathrm{hw})\cdot(\mathrm{sv}_\mu/2)$; moreover $n_\mu\ge4$ and $n_\mu\in2\,\mathrm{hw}\,\mathbb{Z}$.
Thus $C(n)$ is a torus in the sense of Section~I.7.1, and it carries $\mathcal{Q}$ (Definition~39.0).

By the definition of the open cells, a point $y$ lies in $Q^{\circ}_{\mathfrak{c}}(\mathcal{Q})$ if and only if
\[
o_\nu+\mathfrak{c}_\nu\cdot\mathrm{hw}<y_\nu<o_\nu+(\mathfrak{c}_\nu+1)\cdot\mathrm{hw}\qquad\text{for }\nu=1,2,3.
\]
Since $n_\mu=\mathrm{sv}_\mu\cdot\mathrm{hw}$, the translation by $n_\mu e_\mu$ maps $Q^{\circ}_{\mathfrak{c}}(\mathcal{Q})$ onto $Q^{\circ}_{\mathfrak{c}+\mathrm{sv}_\mu e_\mu}(\mathcal{Q})$.
By Definition~39.0 the cells of $C(n)$ are the orbits of the cells under $\sum_{\mu}n_{\mu}\mathbb{Z}\,e_{\mu}$.
They therefore correspond to the classes of $\mathbb{Z}^3$ modulo $\mathrm{sv}_1\mathbb{Z}\times\mathrm{sv}_2\mathbb{Z}\times\mathrm{sv}_3\mathbb{Z}$, and their number is
\[
\mathrm{sv}_1\cdot\mathrm{sv}_2\cdot\mathrm{sv}_3=\prod_{\mu}\frac{n_\mu}{\mathrm{hw}}=\mathrm{nc}(n).
\]
In particular, $n_r=\mathrm{hw}\cdot\mathrm{ncs}$ is an integer $\ge4$;
for every even integer $\mathrm{sv}$ with $\mathrm{hw}\cdot\mathrm{sv}\ge4$ the cubic torus $\mathrm{hw}\cdot\mathrm{sv}$ carries $\mathcal{Q}$ and has $\mathrm{sv}^3=\mathrm{kr}^3\cdot N_r$ cells;
the torus $n_r$ has $N_r$ cells;
and $\check{j}\cdot\mathrm{ncs}$ is an even integer with $\mathrm{hw}\cdot\check{j}\cdot\mathrm{ncs}=\check{j}\cdot n_r\ge8$.

\emph{Representatives.}
In a product over the cells of $C(n)$, or over a set $\mathrm{Scl}$ of cells of $C(n)$, every cell is represented by one index $\mathfrak{c}\in\mathbb{Z}^3$ of its class, and this index may be chosen arbitrarily.
The facts [Q1]--[Q3] hold for every such choice, and so does everything that follows.
The number $\mathcal{Z}_{\mathrm{sv}}(O)$ of Proposition~39\,(e) is the same for all sets of representatives, since $Z_n$ of a product over the cells of $C(n)$, here of $O^{\sharp}_n$, does not depend on that set (Definition~39.0, the sentence on the sets of representatives).
(The function on the configurations of $C(n)$ that such a product defines depends neither on the choice of the indices nor on the order of the factors; see the paragraph \emph{Families} of \cite[Definition~6.13]{QEDII}, together with \cite[Lemma~6.4(a),(c)]{QEDII}.)

\paragraph{1.2 The constant 1.}
The first part of [Q7] shows that $1\in\mathcal{A}(Q^{\circ}_{\mathfrak{c}}(\mathcal{Q}))$ for every $\mathfrak{c}$.
By the second part of [Q7], $1^{\mathrm{gw}}=1$ for every $g_{\mathrm{w}}\in W_{\mathrm{af}}$, since the constant function $1$ of the curvature is real and independent of its argument.
Hence $1^{[\mathfrak{c}']}=1$ for all $\mathfrak{c}'$, and $1^{\sharp}_n=1$ on every torus $C(n)$ that carries $\mathcal{Q}$.
For $\mathrm{sv}$ as in 1.1 it follows that
\begin{equation}\tag{1.1}
\mathcal{Z}_{\mathrm{sv}}(1)=Z_{\mathrm{hw}\cdot\mathrm{sv}}[1],\qquad\text{in particular}\qquad\mathcal{Z}_{[\mathrm{sv}]}(1)=Z_{[\mathrm{sv}]}.
\end{equation}

\paragraph{1.3 One number.}
For integers $n_\mu\ge4$ the symbol $Z_n[1]$ admits three readings.
\begin{description}
\item[(r1)] The expression (I.7.1) with $G=1$.
\item[(r2)] The number $Z_{\mathbb{T}}[\Phi]$ of \cite[Definition~2.13]{QEDII} for the Borel function $\Phi=1$ of the curvature, with $\mathbb{T}=C(n)$ and the mass parameter $M_{\mathrm{ch}}$.
\item[(r3)] The number $Z_{\mathbb{T}}[G]$ of \cite[Definition~2.13]{QEDII} for the element $G=1$ of $\mathcal{A}(\mathbb{R}^3)$, which has no box and is therefore read on every torus.
\end{description}
The readings (r1) and (r2) agree.
Indeed, by Proposition~46\,(f) at $\Phi=1$, the number (r2) is the number $Z_n[1]$ of \cite[Section~5.2.1]{QEDI}, that is, $\omega_n(1)$ of \cite[Section~21.1.4]{QEDI}, formed with $\det\left(D_A+\boldsymbol{\mu}\right)$ in the place of the Wilson determinant
(\cite[Section~5.2.1]{QEDI} treats odd cubic tori; for the other $n$ the number is that of \cite[Definition~10.1]{QEDI}, as in Proposition~46\,(f)).
This expression is (I.7.1), by the first row of the table in Section~I.7.2: it has the same sum over the closed $\mathfrak{z}$, the same box $[0,r_q)^{\mathrm{links}}$, the same weight $\exp\left(-\frac{1}{2}\|F^{\mathrm{u}}\|^2\right)$ and the same curvature $F^{\mathrm{u}}=dA+r_q\,\mathfrak{z}$.
The readings (r2) and (r3) agree because $\int_{\mathbb{T}}\exp\left(-S_{\mathbb{T}}(A)\right)=\det D^{\mathbb{T}}_A$ ([Q8]).
The fact [Q5] concerns the reading (r1), the fact [Q6] the reading (r2)~$=$~(r3), and the facts [Q1]--[Q4], together with (1.1), the reading (r3).
From now on $Z_n[1]$ denotes this single number.
It does not depend on the point $v_0\in\mathrm{Chm}_0$ from which the fermion sites are formed in \cite{QEDII}:
in (r2) the fermion enters only through $\det D^{\mathbb{T}}_A=\det\left(D_A+\boldsymbol{\mu}\right)$, and this determinant is independent of $v_0$ (Proposition~46\,(e)).

\paragraph{1.4 Positivity; the numbers of the statement are defined.}
By [Q5], $\log Z_n[1]$ is a real number for all integers $n_\mu\ge4$.
Hence $Z_{[\mathrm{sv}]}>0$, and $\log Z_{[\mathrm{sv}]}$, $D_{\check{j}}$ and $\delta_{\check{j}}$ are real numbers.

Let $\mathrm{sv}$ be as in 1.1, let $n=\mathrm{hw}\cdot\mathrm{sv}$, and let $O\in\mathcal{A}(Q^{\circ}_0(\mathcal{Q}))$.
The class of $0$ is represented by the index $0$; by 1.1 this choice does not affect $\mathcal{Z}_{\mathrm{sv}}(O)$.
By 1.2, setting $O_0:=O$ and $O_{\mathfrak{c}}:=1$ for the other cells $\mathfrak{c}$ of $C(n)$ defines a family as in [Q1].
The first assertion of [Q2], taken at the cell $0$, gives
\begin{equation}\tag{1.2}
\mathcal{Z}_{\mathrm{sv}}(O)=Z_n\left[O^{\sharp}_n\right]\ge0.
\end{equation}
(That $O^{\sharp}_n$ is read on $C(n)$ and that $Z_n\left[O^{\sharp}_n\right]$ converges absolutely is also the last assertion of 1.5.)
Consequently
\[
\|O\|_*=\left(\frac{\mathcal{Z}_{[\mathrm{ncs}]}(O)}{Z_{[\mathrm{ncs}]}}\right)^{\frac{1}{N_r}}
\]
is a real number $\ge0$, and by the definition of $E^{\mathrm{af}}_n$ on $\mathcal{A}$ in 0.1 it equals $\left(E_{[\mathrm{ncs}]}O^{\sharp}\right)^{\frac{1}{N_r}}$.

\paragraph{1.5 The group of the walls; mirror copies.}
Each wall $(y_\mu,o_\mu+j\cdot\mathrm{hw})$ of $\mathcal{Q}$ is a mirror in the sense of Section~I.4.8, because $o_\mu+j\cdot\mathrm{hw}$ lies in $\frac{1}{2}\mathbb{Z}=\operatorname{Adm}(y_\mu)$.
Its reflection is therefore an element of $W_{\mathrm{af}}$ (Section~I.4.8), and $W(\mathcal{Q})$ is a subgroup of $W_{\mathrm{af}}$.

\emph{The element $g_{\mathrm{w}}$ of Definition~39.0 exists and is unique} (third sentence of Proposition~39\,(a)).
For a mirror $(y_\mu,\mathrm{po})$ of an axis family the vector of the form $y_\mu$ is $e_\mu$, of squared length $1$, and Section~I.4.8 gives
\[
\theta_{y_\mu,\mathrm{po}}(y)=y-2\left(y_\mu-\mathrm{po}\right)e_\mu .
\]
Thus $\theta_{y_\mu,\mathrm{po}}$ replaces the coordinate $y_\mu$ by $2\,\mathrm{po}-y_\mu$ and leaves the two other coordinates unchanged.

Fix $\mu$.
For $i\in\mathbb{Z}$ the open interval $\left(o_\mu+i\cdot\mathrm{hw},\,o_\mu+(i+1)\cdot\mathrm{hw}\right)$ is called the \emph{interval number}~$i$, and the maps
\[
y_\mu\mapsto2\left(o_\mu+i\cdot\mathrm{hw}\right)-y_\mu
\qquad\text{and}\qquad
y_\mu\mapsto y_\mu+2i\cdot\mathrm{hw}
\]
are called the \emph{reflection number}~$i$ and the \emph{translation number}~$i$, respectively.
For $i,i'\in\mathbb{Z}$ substitution yields the following compositions, in which the map named first is applied last:
\begin{align*}
\text{reflection }i,\ \text{reflection }i'\colon\quad&y_\mu\mapsto2\left(o_\mu+i\cdot\mathrm{hw}\right)-2\left(o_\mu+i'\cdot\mathrm{hw}\right)+y_\mu&&(\text{the translation number }i-i'),\\
\text{translation }i,\ \text{reflection }i'\colon\quad&y_\mu\mapsto2\left(o_\mu+i'\cdot\mathrm{hw}\right)-y_\mu+2i\cdot\mathrm{hw}&&(\text{the reflection number }i'+i),\\
\text{reflection }i',\ \text{translation }i\colon\quad&y_\mu\mapsto2\left(o_\mu+i'\cdot\mathrm{hw}\right)-y_\mu-2i\cdot\mathrm{hw}&&(\text{the reflection number }i'-i),\\
\text{translation }i,\ \text{translation }i'\colon\quad&y_\mu\mapsto y_\mu+2i'\cdot\mathrm{hw}+2i\cdot\mathrm{hw}&&(\text{the translation number }i+i').
\end{align*}
Hence the set of these maps is closed under composition.
It is closed under taking inverses as well:
the reflection number $i$ is its own inverse (first line with $i'=i$, the translation number $0$ being the identity), and the inverse of the translation number $i$ is the translation number $-i$ (fourth line).
By the first line, the translation number $i$ is the composition of the reflection number $i$ with the reflection number $0$.
Therefore this set is the group generated by the reflections number $i$, $i\in\mathbb{Z}$.

The translation number $i$ maps the interval number $0$, which is $(o_\mu,o_\mu+\mathrm{hw})$, onto $\left(o_\mu+2i\cdot\mathrm{hw},\,o_\mu+(2i+1)\cdot\mathrm{hw}\right)$, the interval number $2i$.
The reflection number $i$ maps it onto $\left(o_\mu+(2i-1)\cdot\mathrm{hw},\,o_\mu+2i\cdot\mathrm{hw}\right)$, the interval number $2i-1$.
Every integer is of exactly one of the two forms $2i$ and $2i-1$, with exactly one $i$.
So exactly one element of this group maps the interval number $0$ onto the interval number $\mathfrak{c}_\mu$.

The generators $\theta_{y_\mu,o_\mu+j\cdot\mathrm{hw}}$ of $W(\mathcal{Q})$ act on the coordinate $y_\mu$ as the reflection number $j$ and leave the two other coordinates unchanged.
Maps acting on different coordinates commute.
Hence $W(\mathcal{Q})$ consists of the maps
\[
g_{\mathrm{w}}\colon y\mapsto\left(g_{\mathrm{w},1}(y_1),\,g_{\mathrm{w},2}(y_2),\,g_{\mathrm{w},3}(y_3)\right),
\]
where $g_{\mathrm{w},\mu}$ is an element of the group of the coordinate $\mu$; and $g_{\mathrm{w}}$ determines the triple $(g_{\mathrm{w},1},g_{\mathrm{w},2},g_{\mathrm{w},3})$.
The cell $Q^{\circ}_{\mathfrak{c}}(\mathcal{Q})$ is the product of the intervals number $\mathfrak{c}_1$, $\mathfrak{c}_2$, $\mathfrak{c}_3$ of the three coordinates, and $g_{\mathrm{w}}Q^{\circ}_0(\mathcal{Q})$ is the product of the images of the three intervals number $0$ under $g_{\mathrm{w},1}$, $g_{\mathrm{w},2}$, $g_{\mathrm{w},3}$.
Two products of three non-empty intervals coincide if and only if their factors coincide.
Therefore $g_{\mathrm{w}}Q^{\circ}_0(\mathcal{Q})=Q^{\circ}_{\mathfrak{c}}(\mathcal{Q})$ holds if and only if $g_{\mathrm{w},\mu}$ maps the interval number $0$ onto the interval number $\mathfrak{c}_\mu$ for $\mu=1,2,3$.
It follows that for every $\mathfrak{c}\in\mathbb{Z}^3$ exactly one element of $W(\mathcal{Q})$ maps $Q^{\circ}_0(\mathcal{Q})$ onto $Q^{\circ}_{\mathfrak{c}}(\mathcal{Q})$; for $\mathfrak{c}=0$ this element is the identity.

Now let $\mathfrak{c},\mathfrak{c}'$ be two cells, and let $g_{\mathrm{w}}'$, $g_{\mathrm{w}}''$ be the elements of $W(\mathcal{Q})$ that map $Q^{\circ}_0(\mathcal{Q})$ onto $Q^{\circ}_{\mathfrak{c}}(\mathcal{Q})$ and onto $Q^{\circ}_{\mathfrak{c}'}(\mathcal{Q})$, respectively.
Then $g_{\mathrm{w}}''\,{g_{\mathrm{w}}'}^{-1}$ maps $Q^{\circ}_{\mathfrak{c}}(\mathcal{Q})$ onto $Q^{\circ}_{\mathfrak{c}'}(\mathcal{Q})$.
If, conversely, $g_{\mathrm{w}}\in W(\mathcal{Q})$ maps $Q^{\circ}_{\mathfrak{c}}(\mathcal{Q})$ onto $Q^{\circ}_{\mathfrak{c}'}(\mathcal{Q})$, then $g_{\mathrm{w}}\,g_{\mathrm{w}}'$ maps $Q^{\circ}_0(\mathcal{Q})$ onto $Q^{\circ}_{\mathfrak{c}'}(\mathcal{Q})$;
hence $g_{\mathrm{w}}\,g_{\mathrm{w}}'=g_{\mathrm{w}}''$ and $g_{\mathrm{w}}=g_{\mathrm{w}}''\,{g_{\mathrm{w}}'}^{-1}$.
(For $\mathcal{Q}_*$, that is for $o=0$ and $\mathrm{hw}=\mathrm{per}/2$, the translation number $i$ is the map $y_\mu\mapsto i\cdot\mathrm{per}+y_\mu$ and the reflection number $i+1$ is the map $y_\mu\mapsto(i+1)\cdot\mathrm{per}-y_\mu$;
these are the two maps by which Section~I.6.4 describes $\mathsf{g}_{\mathfrak{c}}$ for $\mathfrak{c}_\mu=2i$ and for $\mathfrak{c}_\mu=2i+1$.)

\emph{Mirror copies.}
Let $O\in\mathcal{A}(Q^{\circ}_0(\mathcal{Q}))$ and $\mathfrak{c}\in\mathbb{Z}^3$.
By Definition~39.0, $O^{[\mathfrak{c}]}=O^{\mathrm{gw}}$, where $g_{\mathrm{w}}$ is the element of $W(\mathcal{Q})$ with $g_{\mathrm{w}}Q^{\circ}_0(\mathcal{Q})=Q^{\circ}_{\mathfrak{c}}(\mathcal{Q})$.
The third part of [Q7] places $O^{[\mathfrak{c}]}$ in $\mathcal{A}\left(g_{\mathrm{w}}Q^{\circ}_0(\mathcal{Q})\right)=\mathcal{A}(Q^{\circ}_{\mathfrak{c}}(\mathcal{Q}))$.
Consequently, if $C(n)$ carries $\mathcal{Q}$, the copies $O^{[\mathfrak{c}']}$, one for each cell $\mathfrak{c}'$ of $C(n)$, form a family as in [Q1]:
their product $O^{\sharp}_n$ is read on $C(n)$, and $Z_n\left[O^{\sharp}_n\right]$ converges absolutely.

\subsubsection*{2. Two lemmas on functions with (m), (cv)}

Throughout this section $\mathrm{Kr}$ is an arbitrary subset of $(0,\infty)$, and $\check{j}$, $\mathrm{kr}$ are real numbers; $\check{j}$ need not be an integer.
For a function $\mathrm{wr}\colon\mathrm{Kr}\to\mathbb{R}$ the two properties in question are
\begin{equation}\tag{m}
\frac{\mathrm{wr}(\mathrm{kr}'')}{\mathrm{kr}''}\le\frac{\mathrm{wr}(\mathrm{kr}')}{\mathrm{kr}'}\qquad\text{for }\mathrm{kr}'\le\mathrm{kr}''\text{ in }\mathrm{Kr},
\end{equation}
\begin{equation}\tag{cv}
\begin{split}
(\mathrm{kr}'''-\mathrm{kr}')\cdot\mathrm{wr}(\mathrm{kr}'')&\le(\mathrm{kr}'''-\mathrm{kr}'')\cdot\mathrm{wr}(\mathrm{kr}')+(\mathrm{kr}''-\mathrm{kr}')\cdot\mathrm{wr}(\mathrm{kr}''')\\
&\qquad\text{for }\mathrm{kr}'<\mathrm{kr}''<\mathrm{kr}'''\text{ in }\mathrm{Kr}.
\end{split}
\end{equation}

\begin{named}{Lemma 1}[one variable]
Let $\mathrm{wr}\colon\mathrm{Kr}\to\mathbb{R}$ have the properties (m) and (cv), and let $1$ be a point of $\mathrm{Kr}$.
Then (U) holds.
If $\check{j}$ and $\mathrm{kr}$ are points of $\mathrm{Kr}$ with $1<\check{j}\le\mathrm{kr}$, and $\alpha:=(\mathrm{kr}-1)/(\check{j}-1)$, then $\alpha\ge1$ and (L) holds.
\begin{equation}\tag{U}
\begin{split}
&\mathrm{wr}(\mathrm{kr}'')\le\frac{\mathrm{kr}''}{\mathrm{kr}'}\cdot\mathrm{wr}(\mathrm{kr}')\qquad\text{for }\mathrm{kr}'\le\mathrm{kr}''\text{ in }\mathrm{Kr};\\
&\text{in particular}\quad\mathrm{wr}(\mathrm{kr}'')\le\mathrm{kr}''\cdot\mathrm{wr}(1)\qquad\text{for }\mathrm{kr}''\in\mathrm{Kr},\ \mathrm{kr}''\ge1;
\end{split}
\end{equation}
\begin{equation}\tag{L}
\mathrm{wr}(\mathrm{kr})\ge\alpha\cdot\mathrm{wr}(\check{j})-(\alpha-1)\cdot\mathrm{wr}(1).
\end{equation}
\end{named}

\begin{proof}
From $\mathrm{kr}-1\ge\check{j}-1>0$ one gets $\alpha\ge1$.
Multiplication of (m) by $\mathrm{kr}''>0$ gives (U).
As to (L), if $\mathrm{kr}=\check{j}$, then $\alpha=1$ and (L) is an equality.
If $\mathrm{kr}>\check{j}$, the inequality (cv) at the three points $1<\check{j}<\mathrm{kr}$ reads
\[
(\mathrm{kr}-1)\cdot\mathrm{wr}(\check{j})\le(\mathrm{kr}-\check{j})\cdot\mathrm{wr}(1)+(\check{j}-1)\cdot\mathrm{wr}(\mathrm{kr}).
\]
Dividing by $\check{j}-1>0$ one obtains
\[
\alpha\cdot\mathrm{wr}(\check{j})\le\frac{\mathrm{kr}-\check{j}}{\check{j}-1}\cdot\mathrm{wr}(1)+\mathrm{wr}(\mathrm{kr}),
\]
and $(\mathrm{kr}-\check{j})/(\check{j}-1)=\alpha-1$.
\end{proof}

\begin{named}{Lemma 2}[three variables; the floor from the two values at $(1,1,1)$ and $(\check{j},\check{j},\check{j})$]
Let $\mathrm{fz}\colon\mathrm{Kr}\times\mathrm{Kr}\times\mathrm{Kr}\to\mathbb{R}$ be such that, for $\mu=1,2,3$ and for every choice of the two other arguments in $\mathrm{Kr}$, the function of the argument number $\mu$ has the properties (m) and (cv).
Let $1$ and $\check{j}$ be points of $\mathrm{Kr}$ with $1<\check{j}$, and put $D_{\check{j}}:=\check{j}^3\cdot\mathrm{fz}(1,1,1)-\mathrm{fz}(\check{j},\check{j},\check{j})$.
Then (i) holds.
If $\mathrm{kr}$ is a point of $\mathrm{Kr}$ with $\check{j}\le\mathrm{kr}$, and $\alpha:=(\mathrm{kr}-1)/(\check{j}-1)$, then (ii) holds.
\begin{equation}\tag{i}
\begin{split}
&\mathrm{fz}(\mathrm{kr}'',\mathrm{kr}'',\mathrm{kr}'')\le\left(\frac{\mathrm{kr}''}{\mathrm{kr}'}\right)^3\cdot\mathrm{fz}(\mathrm{kr}',\mathrm{kr}',\mathrm{kr}')\qquad\text{for }\mathrm{kr}'\le\mathrm{kr}''\text{ in }\mathrm{Kr};\\
&\text{in particular}\quad D_{\check{j}}\ge0;
\end{split}
\end{equation}
\begin{equation}\tag{ii}
\mathrm{fz}(\mathrm{kr},\mathrm{kr},\mathrm{kr})\ge\mathrm{kr}^3\cdot\mathrm{fz}(1,1,1)-\alpha^3\cdot D_{\check{j}}.
\end{equation}
\end{named}

\begin{proof}
(i) Apply (U) in the first, in the second and in the third variable, multiplying each time by a positive number:
\begin{align*}
\mathrm{fz}(\mathrm{kr}'',\mathrm{kr}'',\mathrm{kr}'')
&\le\frac{\mathrm{kr}''}{\mathrm{kr}'}\cdot\mathrm{fz}(\mathrm{kr}',\mathrm{kr}'',\mathrm{kr}'')
\le\left(\frac{\mathrm{kr}''}{\mathrm{kr}'}\right)^2\cdot\mathrm{fz}(\mathrm{kr}',\mathrm{kr}',\mathrm{kr}'')\\
&\le\left(\frac{\mathrm{kr}''}{\mathrm{kr}'}\right)^3\cdot\mathrm{fz}(\mathrm{kr}',\mathrm{kr}',\mathrm{kr}').
\end{align*}
For $\mathrm{kr}'=1$ and $\mathrm{kr}''=\check{j}$ this is $\mathrm{fz}(\check{j},\check{j},\check{j})\le\check{j}^3\cdot\mathrm{fz}(1,1,1)$.

(ii) Apply (L) in the first variable, the two other arguments being $\mathrm{kr},\mathrm{kr}$;
in the second variable, the two other arguments being $\check{j},\mathrm{kr}$;
and in the third variable, the two other arguments being $\check{j},\check{j}$:
\begin{align}
\mathrm{fz}(\mathrm{kr},\mathrm{kr},\mathrm{kr})&\ge\alpha\cdot\mathrm{fz}(\check{j},\mathrm{kr},\mathrm{kr})-(\alpha-1)\cdot\mathrm{fz}(1,\mathrm{kr},\mathrm{kr}),\tag{2.1}\\
\mathrm{fz}(\check{j},\mathrm{kr},\mathrm{kr})&\ge\alpha\cdot\mathrm{fz}(\check{j},\check{j},\mathrm{kr})-(\alpha-1)\cdot\mathrm{fz}(\check{j},1,\mathrm{kr}),\tag{2.2}\\
\mathrm{fz}(\check{j},\check{j},\mathrm{kr})&\ge\alpha\cdot\mathrm{fz}(\check{j},\check{j},\check{j})-(\alpha-1)\cdot\mathrm{fz}(\check{j},\check{j},1).\tag{2.3}
\end{align}
The inequality (U), applied from the point $1$ in the variables indicated, gives
\begin{align}
\mathrm{fz}(1,\mathrm{kr},\mathrm{kr})&\le\mathrm{kr}\cdot\mathrm{fz}(1,1,\mathrm{kr})\le\mathrm{kr}^2\cdot\mathrm{fz}(1,1,1)&&\text{(second variable, third variable)},\tag{2.4}\\
\mathrm{fz}(\check{j},1,\mathrm{kr})&\le\mathrm{kr}\cdot\mathrm{fz}(\check{j},1,1)\le\mathrm{kr}\cdot\check{j}\cdot\mathrm{fz}(1,1,1)&&\text{(third variable, first variable)},\tag{2.5}\\
\mathrm{fz}(\check{j},\check{j},1)&\le\check{j}\cdot\mathrm{fz}(1,\check{j},1)\le\check{j}^2\cdot\mathrm{fz}(1,1,1)&&\text{(first variable, second variable)}.\tag{2.6}
\end{align}
Insert (2.3) into (2.2), whose coefficient $\alpha$ is positive, and the result into (2.1), whose coefficient $\alpha$ is positive as well:
\begin{multline*}
\mathrm{fz}(\mathrm{kr},\mathrm{kr},\mathrm{kr})\ge\alpha^3\cdot\mathrm{fz}(\check{j},\check{j},\check{j})-\alpha^2(\alpha-1)\cdot\mathrm{fz}(\check{j},\check{j},1)\\
-\alpha(\alpha-1)\cdot\mathrm{fz}(\check{j},1,\mathrm{kr})-(\alpha-1)\cdot\mathrm{fz}(1,\mathrm{kr},\mathrm{kr}).
\end{multline*}
The three coefficients $-\alpha^2(\alpha-1)$, $-\alpha(\alpha-1)$ and $-(\alpha-1)$ are $\le0$, so that (2.6), (2.5) and (2.4) yield
\[
\mathrm{fz}(\mathrm{kr},\mathrm{kr},\mathrm{kr})\ge\alpha^3\cdot\mathrm{fz}(\check{j},\check{j},\check{j})-\left[\alpha^2(\alpha-1)\cdot\check{j}^2+\alpha(\alpha-1)\cdot\check{j}\cdot\mathrm{kr}+(\alpha-1)\cdot\mathrm{kr}^2\right]\cdot\mathrm{fz}(1,1,1).
\]
Substituting $\mathrm{fz}(\check{j},\check{j},\check{j})=\check{j}^3\cdot\mathrm{fz}(1,1,1)-D_{\check{j}}$ one arrives at
\[
\mathrm{fz}(\mathrm{kr},\mathrm{kr},\mathrm{kr})\ge\left\{\alpha^3\check{j}^3-\alpha^2(\alpha-1)\cdot\check{j}^2-\alpha(\alpha-1)\cdot\check{j}\cdot\mathrm{kr}-(\alpha-1)\cdot\mathrm{kr}^2\right\}\cdot\mathrm{fz}(1,1,1)-\alpha^3\cdot D_{\check{j}}.
\]
The brace equals $\mathrm{kr}^3$.
Indeed, $\alpha\check{j}-(\alpha-1)=\alpha(\check{j}-1)+1=(\mathrm{kr}-1)+1=\mathrm{kr}$, and this identity, used three times, gives
\begin{align*}
\alpha^3\check{j}^3-\alpha^2(\alpha-1)\cdot\check{j}^2&=\alpha^2\check{j}^2\cdot\left(\alpha\check{j}-(\alpha-1)\right)=\alpha^2\check{j}^2\cdot\mathrm{kr},\\
\alpha^2\check{j}^2\cdot\mathrm{kr}-\alpha(\alpha-1)\cdot\check{j}\cdot\mathrm{kr}&=\alpha\check{j}\cdot\mathrm{kr}\cdot\left(\alpha\check{j}-(\alpha-1)\right)=\alpha\check{j}\cdot\mathrm{kr}^2,\\
\alpha\check{j}\cdot\mathrm{kr}^2-(\alpha-1)\cdot\mathrm{kr}^2&=\mathrm{kr}^2\cdot\left(\alpha\check{j}-(\alpha-1)\right)=\mathrm{kr}^3.
\end{align*}
At no point has the sign of a value of $\mathrm{fz}$ been used: every inequality was multiplied by a coefficient of the sign indicated.
\end{proof}

\subsubsection*{3. The function $\mathrm{fz}$}

\paragraph{3.1 Definition and the two properties.}
From now on
\[
\mathrm{Kr}:=\left\{\mathrm{kr}'\in(0,\infty):\ \mathrm{kr}'\cdot n_r\text{ is an integer}\ge4\right\},
\qquad
\mathrm{fz}(\mathrm{kr}_1,\mathrm{kr}_2,\mathrm{kr}_3):=\log Z_{(\mathrm{kr}_1\cdot n_r,\,\mathrm{kr}_2\cdot n_r,\,\mathrm{kr}_3\cdot n_r)}[1]
\]
for $\mathrm{kr}_1,\mathrm{kr}_2,\mathrm{kr}_3\in\mathrm{Kr}$; by 1.4 this is a real number.
The point $1$ lies in $\mathrm{Kr}$, by 1.1.
If $\mathrm{sv}$ is an even integer with $\mathrm{hw}\cdot\mathrm{sv}\ge4$, then $\mathrm{kr}=\mathrm{sv}/\mathrm{ncs}$ lies in $\mathrm{Kr}$, because $\mathrm{kr}\cdot n_r=\mathrm{hw}\cdot\mathrm{sv}$, and $\mathrm{fz}(\mathrm{kr},\mathrm{kr},\mathrm{kr})=\log Z_{[\mathrm{sv}]}$.
In particular $\check{j}$ lies in $\mathrm{Kr}$ (take $\mathrm{sv}=\check{j}\cdot\mathrm{ncs}$), and
\begin{equation}\tag{3.1}
\mathrm{fz}(1,1,1)=\log Z_{[\mathrm{ncs}]},\qquad\mathrm{fz}(\check{j},\check{j},\check{j})=\log Z_{[\check{j}\cdot\mathrm{ncs}]},\qquad D_{\check{j}}=\check{j}^3\cdot\mathrm{fz}(1,1,1)-\mathrm{fz}(\check{j},\check{j},\check{j});
\end{equation}
thus the number $D_{\check{j}}$ of Lemma~2, formed with this function $\mathrm{fz}$, is the number $D_{\check{j}}$ of Proposition~40.

\begin{named}{Claim}
For $\mu=1,2,3$ and for every choice of the two other arguments in $\mathrm{Kr}$, the function of the argument number $\mu$ has the properties (m) and (cv) on $\mathrm{Kr}$.
\end{named}

\begin{proof}
Fix $\mu\in\{1,2,3\}$, and let $\nu,\lambda$ be the two other directions.
Fix $\mathrm{kr}_\nu,\mathrm{kr}_\lambda\in\mathrm{Kr}$ and put $n_\nu:=\mathrm{kr}_\nu\cdot n_r$ and $n_\lambda:=\mathrm{kr}_\lambda\cdot n_r$, which are integers $\ge4$.
For an integer $i\ge4$ let $n(i)$ be the triple with the entry $i$ in the place $\mu$ and the entries $n_\nu$, $n_\lambda$ in the places $\nu$, $\lambda$, and put $\mathrm{wz}(i):=\log Z_{n(i)}[1]$.
The function of the argument number $\mu$ is then $\mathrm{wr}(\mathrm{kr}')=\mathrm{wz}(\mathrm{kr}'\cdot n_r)$, $\mathrm{kr}'\in\mathrm{Kr}$.

\emph{The property (m).}
By [Q6]\,(a), the map $i\mapsto f(n(i))=\mathrm{wz}(i)/(i\cdot n_\nu\cdot n_\lambda)$ is non-increasing on $\{4,5,\dots\}$.
Since the number $n_\nu\cdot n_\lambda>0$ is fixed, the map $i\mapsto\mathrm{wz}(i)/i$ is non-increasing as well.
For $\mathrm{kr}'\le\mathrm{kr}''$ in $\mathrm{Kr}$ it follows that
\[
\frac{\mathrm{wr}(\mathrm{kr}'')}{\mathrm{kr}''}=n_r\cdot\frac{\mathrm{wz}(\mathrm{kr}''\cdot n_r)}{\mathrm{kr}''\cdot n_r}\le n_r\cdot\frac{\mathrm{wz}(\mathrm{kr}'\cdot n_r)}{\mathrm{kr}'\cdot n_r}=\frac{\mathrm{wr}(\mathrm{kr}')}{\mathrm{kr}'}.
\]

\emph{The property (cv).}
By [Q6]\,(b), the function $\mathrm{wz}$ is convex on $\{4,5,\dots\}$:
\cite[Proposition~6.19(b)]{QEDII} states that it is the restriction to these integers of a convex function on $[1,\infty)$.
Consequently
\[
(i''-i)\cdot\mathrm{wz}(i')\le(i''-i')\cdot\mathrm{wz}(i)+(i'-i)\cdot\mathrm{wz}(i'')\qquad\text{for all integers }4\le i<i'<i'';
\]
this is the condition (Cvx) of \cite[Lemma~6.18]{QEDII}, which the proof of \cite[Proposition~6.19(c),(d)]{QEDII} obtains from the clause (b).
The same inequality follows already from the inequalities $\mathrm{wz}(i+1)-\mathrm{wz}(i)\ge\mathrm{wz}(i)-\mathrm{wz}(i-1)$, $i\ge5$, between consecutive differences.
Indeed, by these the differences $\mathrm{wz}(j+1)-\mathrm{wz}(j)$, $j\ge4$, are non-decreasing in $j$, so that for $4\le i<i'<i''$
\begin{align*}
\mathrm{wz}(i')-\mathrm{wz}(i)&=\sum_{j=i}^{i'-1}\left(\mathrm{wz}(j+1)-\mathrm{wz}(j)\right)\le(i'-i)\cdot\left(\mathrm{wz}(i')-\mathrm{wz}(i'-1)\right),\\
\mathrm{wz}(i'')-\mathrm{wz}(i')&=\sum_{j=i'}^{i''-1}\left(\mathrm{wz}(j+1)-\mathrm{wz}(j)\right)\ge(i''-i')\cdot\left(\mathrm{wz}(i'+1)-\mathrm{wz}(i')\right)\\
&\ge(i''-i')\cdot\left(\mathrm{wz}(i')-\mathrm{wz}(i'-1)\right).
\end{align*}
Multiplying the first of these relations by $i''-i'>0$ and the second by $i'-i>0$ one obtains
\[
(i''-i')\cdot\left(\mathrm{wz}(i')-\mathrm{wz}(i)\right)\le(i'-i)\cdot\left(\mathrm{wz}(i'')-\mathrm{wz}(i')\right),
\]
which is the asserted inequality with its terms rearranged.
Now let $\mathrm{kr}'<\mathrm{kr}''<\mathrm{kr}'''$ be points of $\mathrm{Kr}$, and take $i=\mathrm{kr}'\cdot n_r$, $i'=\mathrm{kr}''\cdot n_r$ and $i''=\mathrm{kr}'''\cdot n_r$.
Division by $n_r>0$ gives (cv) for $\mathrm{wr}$.
\end{proof}

\begin{named}{Remark 3.2}[a second source of (m), (cv) at the points that are used]
The set
\[
\mathrm{Krv}:=\left\{\frac{\mathrm{sv}'}{\mathrm{ncs}}:\ \mathrm{sv}'\text{ an even integer},\ \mathrm{hw}\cdot\mathrm{sv}'\ge4\right\}
\]
is a subset of $\mathrm{Kr}$; it contains $1$, $\check{j}$ and the ratio $\mathrm{kr}$ of every torus that occurs in Proposition~40.
For $\mathrm{kr}_\mu=\mathrm{sv}_\mu/\mathrm{ncs}$ in $\mathrm{Krv}$ one has $\mathrm{fz}(\mathrm{kr}_1,\mathrm{kr}_2,\mathrm{kr}_3)=\log\mathcal{Z}_{\mathrm{sv}}(1)$ by (1.1), and $\mathcal{Z}_{\mathrm{sv}}(1)>0$ by 1.4.
Proposition~39\,(e) at $O=1$ asserts that $\mathrm{sv}_\mu\mapsto\mathcal{Z}_{\mathrm{sv}}(1)^{\frac{1}{\mathrm{sv}_\mu}}$ is non-increasing, which is (m) on $\mathrm{Krv}$ (take logarithms, by (pw4), and multiply by $\mathrm{ncs}$);
it asserts moreover that $\mathrm{sv}_\mu\mapsto\log\mathcal{Z}_{\mathrm{sv}}(1)$ is convex on $\{\mathrm{sv}_\mu\text{ even},\ \mathrm{hw}\cdot\mathrm{sv}_\mu\ge4\}$, which is (cv) on $\mathrm{Krv}$ (by the computation of 3.1, with steps of length $2$).
Lemmas~1 and~2 hold for every subset of $(0,\infty)$.
Hence, in Sections~5, 6 and~10, the set $\mathrm{Kr}$ may be replaced by $\mathrm{Krv}$, and [Q6] by Proposition~39\,(e) at $O=1$.
\end{named}

\subsubsection*{4. Proof of (b)}

\begin{proof}[Proof of Proposition~40\,(b)]
Let $\mathrm{sv}\ge\mathrm{ncs}$ be an even integer, and let $O\in\mathcal{A}\left(Q^{\circ}_{0}(\mathcal{Q})\right)$.
Then $\mathrm{hw}\cdot\mathrm{sv}\ge\mathrm{hw}\cdot\mathrm{ncs}=n_r\ge4$.
Each of the four triples
\[
(\mathrm{sv},\mathrm{sv},\mathrm{sv}),\qquad
(\mathrm{ncs},\mathrm{sv},\mathrm{sv}),\qquad
(\mathrm{ncs},\mathrm{ncs},\mathrm{sv}),\qquad
(\mathrm{ncs},\mathrm{ncs},\mathrm{ncs})
\]
consists of even positive integers $\mathrm{sv}_\mu$ with $\mathrm{hw}\cdot\mathrm{sv}_\mu\ge4$.
Consequently, by Section~1.1, each of the four tori with the sides $\mathrm{hw}\cdot\mathrm{sv}_1$, $\mathrm{hw}\cdot\mathrm{sv}_2$, $\mathrm{hw}\cdot\mathrm{sv}_3$ carries $\mathcal{Q}$;
by the last sentence of [Q1] (Proposition~39\,(b)), one and the same $O$ is read on all four of them;
and by (1.2) the four numbers $\mathcal{Z}_{(\mathrm{sv}_1,\mathrm{sv}_2,\mathrm{sv}_3)}(O)$ are defined and non-negative.
By definition,
\[
\mathcal{Z}_{[\mathrm{sv}]}(O)=\mathcal{Z}_{(\mathrm{sv},\mathrm{sv},\mathrm{sv})}(O),
\qquad
\mathcal{Z}_{[\mathrm{ncs}]}(O)=\mathcal{Z}_{(\mathrm{ncs},\mathrm{ncs},\mathrm{ncs})}(O).
\]

Apply the monotonicity of [Q4] (Proposition~39\,(e)) in the direction $\mu=1$, the two other entries being fixed at $(\mathrm{sv}_2,\mathrm{sv}_3)=(\mathrm{sv},\mathrm{sv})$, at the two points $\mathrm{ncs}\le\mathrm{sv}$:
\[
\mathcal{Z}_{(\mathrm{sv},\mathrm{sv},\mathrm{sv})}(O)^{\frac{1}{\mathrm{sv}}}\le\mathcal{Z}_{(\mathrm{ncs},\mathrm{sv},\mathrm{sv})}(O)^{\frac{1}{\mathrm{ncs}}}.
\]
Both sides are raised to the power $\mathrm{sv}$.
The inequality is preserved by (pw1), and (pw2), applied on both sides together with $\mathrm{sv}/\mathrm{ncs}=\mathrm{kr}$, turns it into
\[
\mathcal{Z}_{(\mathrm{sv},\mathrm{sv},\mathrm{sv})}(O)\le\mathcal{Z}_{(\mathrm{ncs},\mathrm{sv},\mathrm{sv})}(O)^{\mathrm{kr}}.
\]
Next, [Q4] in the direction $\mu=2$ with the fixed entries $(\mathrm{sv}_1,\mathrm{sv}_3)=(\mathrm{ncs},\mathrm{sv})$, and in the direction $\mu=3$ with the fixed entries $(\mathrm{sv}_1,\mathrm{sv}_2)=(\mathrm{ncs},\mathrm{ncs})$, each time at the two points $\mathrm{ncs}\le\mathrm{sv}$, gives
\[
\mathcal{Z}_{(\mathrm{ncs},\mathrm{sv},\mathrm{sv})}(O)^{\frac{1}{\mathrm{sv}}}\le\mathcal{Z}_{(\mathrm{ncs},\mathrm{ncs},\mathrm{sv})}(O)^{\frac{1}{\mathrm{ncs}}},
\qquad
\mathcal{Z}_{(\mathrm{ncs},\mathrm{ncs},\mathrm{sv})}(O)^{\frac{1}{\mathrm{sv}}}\le\mathcal{Z}_{(\mathrm{ncs},\mathrm{ncs},\mathrm{ncs})}(O)^{\frac{1}{\mathrm{ncs}}}.
\]
Raising both of them to the power $\mathrm{sv}$, again by (pw1) and (pw2), one obtains
\[
\mathcal{Z}_{(\mathrm{ncs},\mathrm{sv},\mathrm{sv})}(O)\le\mathcal{Z}_{(\mathrm{ncs},\mathrm{ncs},\mathrm{sv})}(O)^{\mathrm{kr}},
\qquad
\mathcal{Z}_{(\mathrm{ncs},\mathrm{ncs},\mathrm{sv})}(O)\le\mathcal{Z}_{(\mathrm{ncs},\mathrm{ncs},\mathrm{ncs})}(O)^{\mathrm{kr}}.
\]
As $\mathrm{kr}>0$, the map $\mathrm{bz}\mapsto\mathrm{bz}^{\mathrm{kr}}$ is non-decreasing on $[0,\infty)$ by (pw1).
The three inequalities can therefore be chained, and with (pw2) in the last step it follows that
\[
\mathcal{Z}_{[\mathrm{sv}]}(O)\le\left(\left(\mathcal{Z}_{[\mathrm{ncs}]}(O)^{\mathrm{kr}}\right)^{\mathrm{kr}}\right)^{\mathrm{kr}}=\mathcal{Z}_{[\mathrm{ncs}]}(O)^{\mathrm{kr}^3}.
\]
This is the first assertion of (b).

\emph{The second assertion of (b).}
Take $O=1$, an element of $\mathcal{A}\left(Q^{\circ}_{0}(\mathcal{Q})\right)$ by Section~1.2.
By (1.1), the inequality just proved becomes $Z_{[\mathrm{sv}]}\le Z_{[\mathrm{ncs}]}^{\mathrm{kr}^3}$, and both of its sides are positive by Section~1.4.
The logarithm is increasing, and $\log\left(Z_{[\mathrm{ncs}]}^{\mathrm{kr}^3}\right)=\mathrm{kr}^3\cdot\log Z_{[\mathrm{ncs}]}$ by (pw4).
Hence
\[
\log Z_{[\mathrm{sv}]}\le\mathrm{kr}^3\cdot\log Z_{[\mathrm{ncs}]}.
\]
The same inequality is also Lemma~2\,(i) for the function $\mathrm{fz}$ of Section~3, taken at $\mathrm{kr}'=1$ and $\mathrm{kr}''=\mathrm{kr}$.
\end{proof}

\subsubsection*{5. Proof of (a)}

\begin{proof}[Proof of Proposition~40\,(a)]
By Section~1.1, the number $\check{j}\cdot\mathrm{ncs}$ is an even integer; it is $\ge\mathrm{ncs}$, and the ratio that belongs to it is $\mathrm{kr}=\check{j}\cdot\mathrm{ncs}/\mathrm{ncs}=\check{j}$.
The second assertion of (b), applied with $\mathrm{sv}=\check{j}\cdot\mathrm{ncs}$, therefore yields
\[
\log Z_{[\check{j}\cdot\mathrm{ncs}]}\le\check{j}^3\cdot\log Z_{[\mathrm{ncs}]},
\]
and this is the inequality $D_{\check{j}}\ge0$.
As $(\check{j}-1)^3\cdot N_r>0$, also $\delta_{\check{j}}=D_{\check{j}}/\left((\check{j}-1)^3\cdot N_r\right)\ge0$.
\end{proof}

\subsubsection*{6. Proof of (c)}

\begin{proof}[Proof of Proposition~40\,(c)]
\emph{The case $\mathrm{sv}=\mathrm{ncs}$.}
Here $\mathrm{kr}=1$.
The first inequality of (c) takes the form $\log Z_{[\mathrm{ncs}]}\ge\log Z_{[\mathrm{ncs}]}-0$.
The second one takes the form
\[
\log Z_{[\mathrm{ncs}]}\ge\log Z_{[\mathrm{ncs}]}-\frac{D_{\check{j}}}{(\check{j}-1)^3},
\]
which is true by (a).

\emph{The case of an even integer $\mathrm{sv}\ge\check{j}\cdot\mathrm{ncs}$.}
In this case $\mathrm{hw}\cdot\mathrm{sv}\ge\check{j}\cdot n_r\ge8$.
By Section~3.1 the three points $1$, $\check{j}$ and $\mathrm{kr}=\mathrm{sv}/\mathrm{ncs}$ belong to $\mathrm{Kr}$, one has $1<2\le\check{j}\le\mathrm{kr}$, and $\mathrm{fz}$ has the properties (m) and (cv) in each variable.
Lemma~2\,(ii), combined with (3.1), with $\mathrm{fz}(\mathrm{kr},\mathrm{kr},\mathrm{kr})=\log Z_{[\mathrm{sv}]}$ and with $\alpha=(\mathrm{kr}-1)/(\check{j}-1)$, gives
\[
\log Z_{[\mathrm{sv}]}\ge\mathrm{kr}^3\cdot\log Z_{[\mathrm{ncs}]}-\left(\frac{\mathrm{kr}-1}{\check{j}-1}\right)^3\cdot D_{\check{j}}.
\]
This is the first inequality of (c).
Since $0\le\mathrm{kr}-1\le\mathrm{kr}$, and since $D_{\check{j}}\ge0$ by (a),
\[
\left(\frac{\mathrm{kr}-1}{\check{j}-1}\right)^3\cdot D_{\check{j}}\le\frac{\mathrm{kr}^3\cdot D_{\check{j}}}{(\check{j}-1)^3},
\]
and the second inequality of (c) follows.
\end{proof}

\begin{named}{Remark}[The tori in the second case of the proof of (c)]
Let $\mathrm{sv}\ge\check{j}\cdot\mathrm{ncs}$ be an even integer.
In this case (c) is derived from Lemma~2\,(ii), whose proof uses the values of $\mathrm{fz}$ at the eleven triples
\begin{gather*}
(\mathrm{kr},\mathrm{kr},\mathrm{kr}),\quad
(\check{j},\mathrm{kr},\mathrm{kr}),\quad
(1,\mathrm{kr},\mathrm{kr}),\quad
(\check{j},\check{j},\mathrm{kr}),\quad
(\check{j},1,\mathrm{kr}),\quad
(\check{j},\check{j},\check{j}),\\
(\check{j},\check{j},1),\quad
(1,1,\mathrm{kr}),\quad
(\check{j},1,1),\quad
(1,\check{j},1),\quad
(1,1,1).
\end{gather*}
The corresponding tori are rectangular;
their sides are among $n_r$, $\check{j}\cdot n_r$ and $\mathrm{hw}\cdot\mathrm{sv}$, which are integers $\ge4$ of either parity, as in [Q6].
\end{named}

\subsubsection*{7. Proof of (d)}

\begin{proof}[Proof of Proposition~40\,(d)]
Let either $\mathrm{sv}=\mathrm{ncs}$, or let $\mathrm{sv}$ be an even integer with $\mathrm{sv}\ge\check{j}\cdot\mathrm{ncs}$.
Put $n:=\mathrm{hw}\cdot\mathrm{sv}$, and let $C(n)$ be the cubic torus with $n$ sites per side.
Let $\mathrm{Scl}$ be a set of distinct cells of $C(n)$, and let $O_{\mathfrak{c}}\in\mathcal{A}\left(Q^{\circ}_{0}(\mathcal{Q})\right)$ for $\mathfrak{c}\in\mathrm{Scl}$.
By Section~1.1, the torus $C(n)$ carries $\mathcal{Q}$ and has $\mathrm{nc}(n)=\mathrm{sv}^3=\mathrm{kr}^3\cdot N_r$ cells.
In particular $\mathrm{Scl}$ is finite, and $\#|\mathrm{Scl}|\le\mathrm{sv}^3$.

\paragraph{7.1. The family.}
Define
\[
\mathrm{Og}_{\mathfrak{c}}:=O_{\mathfrak{c}}^{[\mathfrak{c}]}\quad\text{for }\mathfrak{c}\in\mathrm{Scl},
\qquad
\mathrm{Og}_{\mathfrak{c}}:=1\quad\text{for the }\mathrm{sv}^3-\#|\mathrm{Scl}|\text{ other cells }\mathfrak{c}\text{ of }C(n).
\]
By Sections~1.5 and~1.2, $\mathrm{Og}_{\mathfrak{c}}$ lies in $\mathcal{A}\left(Q^{\circ}_{\mathfrak{c}}(\mathcal{Q})\right)$ for every cell $\mathfrak{c}$ of $C(n)$;
so $(\mathrm{Og}_{\mathfrak{c}})$ is a family as in [Q1].
The constant $1$ is the neutral element of the product; hence
\[
\prod_{\mathfrak{c}}\mathrm{Og}_{\mathfrak{c}}=\prod_{\mathfrak{c}\in\mathrm{Scl}}O_{\mathfrak{c}}^{[\mathfrak{c}]},
\]
where the right side is $1$ if $\mathrm{Scl}$ is empty.
By [Q1], this product is read on $C(n)$, and $Z_n[\,\cdot\,]$ of it converges absolutely.
By [Q3] (Proposition~39\,(d)), $(\mathrm{Og}_{\mathfrak{c}})^{\sharp}_n=(O_{\mathfrak{c}})^{\sharp}_n$ for $\mathfrak{c}\in\mathrm{Scl}$;
by Section~1.2, $(\mathrm{Og}_{\mathfrak{c}})^{\sharp}_n=1$ for the remaining cells.
Therefore
\[
Z_n\left[(\mathrm{Og}_{\mathfrak{c}})^{\sharp}_n\right]=\mathcal{Z}_{[\mathrm{sv}]}(O_{\mathfrak{c}})\quad\text{for }\mathfrak{c}\in\mathrm{Scl},
\qquad
Z_n\left[(\mathrm{Og}_{\mathfrak{c}})^{\sharp}_n\right]=Z_{[\mathrm{sv}]}\quad\text{for the other cells}.
\]

\paragraph{7.2. Chessboard.}
The second assertion of [Q2] (Proposition~39\,(c)) is applied to the family $(\mathrm{Og}_{\mathfrak{c}})$, with $\mathrm{nc}(n)=\mathrm{sv}^3$ and with the values found in Section~7.1.
A cell $\mathfrak{c}\in\mathrm{Scl}$ contributes the factor $\mathcal{Z}_{[\mathrm{sv}]}(O_{\mathfrak{c}})^{\frac{1}{\mathrm{sv}^3}}$.
Each of the $\mathrm{sv}^3-\#|\mathrm{Scl}|$ other cells contributes the factor $Z_{[\mathrm{sv}]}^{\frac{1}{\mathrm{sv}^3}}$, and by (pw4) the product of these equal factors is $Z_{[\mathrm{sv}]}^{\frac{\mathrm{sv}^3-\#|\mathrm{Scl}|}{\mathrm{sv}^3}}$;
if $\mathrm{Scl}$ is the set of all cells, this is the empty product $1=Z_{[\mathrm{sv}]}^{0}$.
One obtains
\[
\left|Z_n\left[\prod_{\mathfrak{c}\in\mathrm{Scl}}O_{\mathfrak{c}}^{[\mathfrak{c}]}\right]\right|
\le\prod_{\mathfrak{c}\in\mathrm{Scl}}\mathcal{Z}_{[\mathrm{sv}]}(O_{\mathfrak{c}})^{\frac{1}{\mathrm{sv}^3}}\cdot Z_{[\mathrm{sv}]}^{\frac{\mathrm{sv}^3-\#|\mathrm{Scl}|}{\mathrm{sv}^3}}.
\]
This is divided by $Z_{[\mathrm{sv}]}=Z_{[\mathrm{sv}]}^{1}$, a positive number by Section~1.4.
By (pw4),
\[
\frac{Z_{[\mathrm{sv}]}^{\frac{\mathrm{sv}^3-\#|\mathrm{Scl}|}{\mathrm{sv}^3}}}{Z_{[\mathrm{sv}]}}
=Z_{[\mathrm{sv}]}^{-\frac{\#|\mathrm{Scl}|}{\mathrm{sv}^3}}
=\prod_{\mathfrak{c}\in\mathrm{Scl}}Z_{[\mathrm{sv}]}^{-\frac{1}{\mathrm{sv}^3}}.
\]
Moreover, (pw4) and (pw3) give, for every $\mathrm{bz}\ge0$ and every $\mathrm{xe}>0$,
\begin{equation}
\mathrm{bz}^{\mathrm{xe}}\cdot Z_{[\mathrm{sv}]}^{-\mathrm{xe}}
=\mathrm{bz}^{\mathrm{xe}}\cdot\left(\frac{1}{Z_{[\mathrm{sv}]}}\right)^{\mathrm{xe}}
=\left(\frac{\mathrm{bz}}{Z_{[\mathrm{sv}]}}\right)^{\mathrm{xe}}.
\tag{7.0}
\end{equation}
Take $\mathrm{bz}=\mathcal{Z}_{[\mathrm{sv}]}(O_{\mathfrak{c}})$ and $\mathrm{xe}=1/\mathrm{sv}^3$.
Since $E_{[\mathrm{sv}]}=Z_n[\,\cdot\,]/Z_{[\mathrm{sv}]}$ (Section~0.1), it follows that
\begin{equation}
\left|E_{[\mathrm{sv}]}\prod_{\mathfrak{c}\in\mathrm{Scl}}O_{\mathfrak{c}}^{[\mathfrak{c}]}\right|
\le\prod_{\mathfrak{c}\in\mathrm{Scl}}\left(\frac{\mathcal{Z}_{[\mathrm{sv}]}(O_{\mathfrak{c}})}{Z_{[\mathrm{sv}]}}\right)^{\frac{1}{\mathrm{sv}^3}}.
\tag{7.1}
\end{equation}

\paragraph{7.3. One factor.}
Let $O\in\mathcal{A}\left(Q^{\circ}_{0}(\mathcal{Q})\right)$.
From (b), from (pw1) with the exponent $1/\mathrm{sv}^3=1/\left(\mathrm{kr}^3\cdot N_r\right)$, and from (pw2), one gets
\begin{equation}
\mathcal{Z}_{[\mathrm{sv}]}(O)^{\frac{1}{\mathrm{sv}^3}}
\le\left(\mathcal{Z}_{[\mathrm{ncs}]}(O)^{\mathrm{kr}^3}\right)^{\frac{1}{\mathrm{kr}^3\cdot N_r}}
=\mathcal{Z}_{[\mathrm{ncs}]}(O)^{\frac{1}{N_r}}.
\tag{7.2}
\end{equation}
By (c), the first member of its chain is at least the last one, that is, $\log Z_{[\mathrm{sv}]}\ge\mathrm{kr}^3\cdot\left[\log Z_{[\mathrm{ncs}]}-D_{\check{j}}/(\check{j}-1)^3\right]$ (the second form of (c)).
Division by $\mathrm{kr}^3\cdot N_r>0$ gives
\begin{equation}
\frac{\log Z_{[\mathrm{sv}]}}{\mathrm{kr}^3\cdot N_r}
\ge\frac{\log Z_{[\mathrm{ncs}]}}{N_r}-\frac{D_{\check{j}}}{(\check{j}-1)^3\cdot N_r}
=\frac{\log Z_{[\mathrm{ncs}]}}{N_r}-\delta_{\check{j}}.
\tag{7.3}
\end{equation}
Powers of a positive base are defined in Section~0.2 by $\mathrm{bz}^{\mathrm{xe}}=\exp\left(\mathrm{xe}\cdot\log\mathrm{bz}\right)$.
By this definition, and because $\exp$ is increasing,
\[
Z_{[\mathrm{sv}]}^{-\frac{1}{\mathrm{sv}^3}}
=\exp\left(-\frac{\log Z_{[\mathrm{sv}]}}{\mathrm{kr}^3\cdot N_r}\right)
\le\exp\left(-\frac{\log Z_{[\mathrm{ncs}]}}{N_r}+\delta_{\check{j}}\right)
=Z_{[\mathrm{ncs}]}^{-\frac{1}{N_r}}\cdot\exp(\delta_{\check{j}}).
\]
This inequality is multiplied by (7.2); all factors are non-negative.
On the left side (7.0) is used, and on the right side (7.0) with $\mathrm{ncs}$ in the place of $\mathrm{sv}$ and with $\mathrm{xe}=1/N_r$.
With the definition of $\|O\|_*$ this gives
\begin{equation}
\left(\frac{\mathcal{Z}_{[\mathrm{sv}]}(O)}{Z_{[\mathrm{sv}]}}\right)^{\frac{1}{\mathrm{sv}^3}}
\le\left(\frac{\mathcal{Z}_{[\mathrm{ncs}]}(O)}{Z_{[\mathrm{ncs}]}}\right)^{\frac{1}{N_r}}\cdot\exp(\delta_{\check{j}})
=\exp(\delta_{\check{j}})\cdot\|O\|_*.
\tag{7.4}
\end{equation}

\paragraph{7.4. Conclusion.}
For each $\mathfrak{c}\in\mathrm{Scl}$, take (7.4) with $O=O_{\mathfrak{c}}$ and insert it in the right side of (7.1).
Since all factors are non-negative, the result is the inequality of (d).
If $\mathrm{Scl}$ is empty, both sides of that inequality are equal to $1$.
\end{proof}

\begin{named}{Remark 7.5}[One cell]
By \cite[Definition~2.12]{QEDII}, $\mathcal{A}\left(Q^{\circ}_{0}(\mathcal{Q})\right)$ is a linear space that contains the constant $1$.
The identity map is an element of $W(\mathcal{Q})$ which maps $Q^{\circ}_{0}(\mathcal{Q})$ onto itself;
by the uniqueness established in Section~1.5, it is the element $g_{\mathrm{w}}$ of Definition~39.0 for $\mathfrak{c}=\mathfrak{c}'=0$.
It is represented by the empty word of reflections, and so $O^{[0]}=O$ (Section~I.4.8).
The functional $Z_n[\,\cdot\,]$ is linear \cite[Definition~2.13(a)]{QEDII}.
Now let $\mathrm{sv}$ be as in (d), let $O\in\mathcal{A}\left(Q^{\circ}_{0}(\mathcal{Q})\right)$, and let $\mathrm{cz}$ be a complex number.
Then (d), applied with $\mathrm{Scl}=\{0\}$ and $O_0:=O-\mathrm{cz}$, reads
\[
\left|E_{[\mathrm{sv}]}O-\mathrm{cz}\right|\le\exp(\delta_{\check{j}})\cdot\|O-\mathrm{cz}\|_*.
\]
\end{named}

\subsubsection*{8. Proof of (e)}

\begin{proof}[Proof of Proposition~40\,(e)]
Assume the setting of (e), and let $\mathcal{Q}=\mathcal{Q}_*$, $\mathrm{hw}=\mathrm{hc}$, $\mathrm{ncs}=2L$ and $\check{j}=L$ as chosen there.

\paragraph{8.1. The data of Proposition~40.}
By Section~I.0.5, $L$ is an odd integer $\ge5$, $M$ is a power of $L$ with $M\ge L^2$, $\bar{m}\in(0,1]$, and $h=L^{\left\lfloor\log_L\frac{1}{\bar{m}}\right\rfloor}$ is a power of $L$ with a non-negative exponent.
By Section~I.6.1, $r_K(e)$ is a power $L^i$ with $i\ge0$.
Hence $M\cdot r_K(e)$ is a power of $L$ which is $\ge L^2$, and $T^{*}(e)=\log_L\left(M\cdot r_K(e)\right)-1$ is an integer $\ge1$ (Section~I.6.2).
The number $N$ is an integer $\ge0$ (Section~I.0.5).
Consequently,
\begin{gather*}
n_*=h\cdot L^{T^{*}(e)+N}\ \text{is an integer}\ \ge L\ge5,\\
\mathrm{per}=\frac{n_*}{L}=h\cdot L^{T^{*}(e)+N-1}\ \text{is an integer}\ \ge1,\\
\mathrm{hc}=\frac{\mathrm{per}}{2}\in\frac{1}{2}\mathbb{Z},\qquad\mathrm{hc}>0.
\end{gather*}
So $\mathcal{Q}_*=(\mathfrak{F}_{\mathrm{c}};0,0,0;\mathrm{hc})$ of Section~I.6.4 is a coordinate cell structure in the sense of Definition~39.0, with $o=0$ and $\mathrm{hw}=\mathrm{hc}$.
Its open cells $Q^{\circ}_{\mathfrak{c}}$, its group $W(\mathcal{Q}_*)$, the elements $\mathsf{g}_{\mathfrak{c}}$, the copies $O^{[\mathfrak{c}]}=O^{\mathsf{g}_{\mathfrak{c}}}$ and the products $O^{\sharp}$ of Section~I.6.4 are the objects of Definition~39.0 for this cell structure;
indeed, by Definition~39.0, $\mathcal{Q}_*$ is the case $o=0$, $\mathrm{hw}=\mathrm{hc}$, and there, for $\mathfrak{c}=0$, the element $g_{\mathrm{w}}$ with $g_{\mathrm{w}}Q^{\circ}_{0}=Q^{\circ}_{\mathfrak{c}'}$ is $\mathsf{g}_{\mathfrak{c}'}$.
Moreover,
\begin{gather*}
\mathrm{ncs}:=2L\ \text{is an even integer}\ \ge2,\qquad
n_r=\mathrm{hc}\cdot2L=\mathrm{per}\cdot L=n_*\ge4,\\
N_r=(2L)^3=8L^3=N_{*},\qquad
\check{j}:=L\ \text{is an integer}\ \ge2,\\
q=e\cdot\eta^{\frac{1}{2}}>0,\qquad
\boldsymbol{\mu}\ \text{is real}.
\end{gather*}
(In the setting of (e) the number $\mathrm{per}$ is in fact at least $L^2$.
Indeed, $N\ge N_0(e)=N_{\min}\left(R^{*}(e),e\right)$ means that $n_*=R^{*}(e)\cdot L^{N}\ge L\cdot M\cdot r_0$ (Section~I.6.2), and $r_0\ge1$ (Section~I.6.1);
hence $\mathrm{per}=n_*/L\ge M\cdot r_0\ge L^2$.)

\paragraph{8.2. The numbers.}
One has $Z_{[\mathrm{ncs}]}=Z_{\mathrm{hc}\cdot2L}[1]=Z_{n_*}[1]$.
Further, $\check{j}\cdot\mathrm{ncs}=2L^2$ and $\mathrm{hc}\cdot2L^2=L\cdot n_*$, so that $Z_{[\check{j}\cdot\mathrm{ncs}]}=Z_{L n_*}[1]$.
Hence
\[
D_{\check{j}}=L^3\cdot\log Z_{n_*}[1]-\log Z_{L n_*}[1]=D_L,
\qquad
\delta_{\check{j}}=\frac{D_L}{(L-1)^3\cdot N_{*}}=\delta,
\]
the numbers of Section~I.7.6, formed with the lattice mass $\boldsymbol{\mu}$.
Also $E_{[\mathrm{ncs}]}=E^{\mathrm{af}}_{n_*}$, and $O^{\sharp}_{n_*}$ is the product of the $O^{[\mathfrak{c}]}$ over the $N_{*}$ cells of the torus $\mathrm{kt}=1$;
thus $\|O\|_*=\left(E^{\mathrm{af}}_{n_*}O^{\sharp}\right)^{\frac{1}{N_{*}}}$.
By (a), $D_L\ge0$.

\paragraph{8.3. The sides.}
Let $n$ be a positive integer.
Since $\mathrm{hc}\cdot\mathrm{sv}=\mathrm{per}\cdot(\mathrm{sv}/2)$, an even integer $\mathrm{sv}$ with $\mathrm{hc}\cdot\mathrm{sv}=n$ exists if and only if $\mathrm{per}$ divides $n$.
In that case $\mathrm{sv}=2n/\mathrm{per}$ and $\mathrm{kr}=\mathrm{sv}/\mathrm{ncs}=n/(L\cdot\mathrm{per})=n/n_*$, and for such $n$
\begin{gather*}
\mathrm{sv}=\mathrm{ncs}\iff n=n_*,\qquad
\mathrm{sv}\ge\mathrm{ncs}\iff n\ge n_*,\\
\mathrm{sv}\ge\check{j}\cdot\mathrm{ncs}=2L^2\iff n\ge L^2\cdot\mathrm{per}=L\cdot n_*;
\end{gather*}
in these cases $\mathrm{hw}\cdot\mathrm{sv}=n\ge n_*\ge4$.
So the hypothesis of (c) and (d), namely ``$\mathrm{sv}=\mathrm{ncs}$, or $\mathrm{sv}$ is an even integer $\ge\check{j}\cdot\mathrm{ncs}$'', says that $n$ belongs to
\[
\mathrm{Sd}:=\{n_*\}\cup\left\{n\in\mathrm{per}\,\mathbb{Z}:\ n\ge L\cdot n_*\right\},
\]
and for $n\in\mathrm{Sd}$ the hypothesis of (b) is satisfied as well.
Hence (b), (c) and (d) hold on the cubic torus with $n$ sites per side for every $n\in\mathrm{Sd}$, with $\mathrm{kr}=n/n_*$ and $\mathrm{sv}=2n/\mathrm{per}$.

\paragraph{8.4. The list.}
\begin{itemize}
\item[(1)]
Let $\mathrm{kt}\ge1$ be an integer.
Then $\mathrm{kt}\cdot n_*=\mathrm{kt}\cdot L\cdot\mathrm{per}$ lies in $\mathrm{per}\,\mathbb{Z}$, and it lies in $\mathrm{Sd}$ if and only if $\mathrm{kt}=1$ or $\mathrm{kt}\ge L$.
Thus the multiples of $n_*$ in $\mathrm{Sd}$ are the sides of the tori $\mathrm{kt}$ of Section~I.6.4 with $\mathrm{kt}=1$ or $\mathrm{kt}\ge L$.
\item[(2)]
The torus $\mathrm{kt}=L$ is the second reference torus (Sections~I.6.2 and~I.6.4).
\item[(3)]
Let $T\ge T^{*}(e)+1$ and $n=8hL^{T+N}$ (Section~I.6.5).
Then
\[
n=8L^{T-T^{*}(e)}\cdot hL^{T^{*}(e)+N}=\mathrm{kt}\cdot n_*
\qquad\text{with the integer}\qquad
\mathrm{kt}=8L^{T-T^{*}(e)}\ge8L\ge L;
\]
hence $n\in\mathrm{Sd}$, and $n$ is even.
\item[(4)]
The number $n:=(L^2+1)\cdot\mathrm{per}$ lies in $\mathrm{per}\,\mathbb{Z}$ and satisfies $n\ge L^2\cdot\mathrm{per}=L\cdot n_*$; hence $n\in\mathrm{Sd}$.
It is not a multiple of $n_*=L\cdot\mathrm{per}$, because $L\ge2$ does not divide $L^2+1$.
\end{itemize}
\end{proof}

\begin{named}{Remark 8.5}[The sides $8hL^{T+N}$ that lie in $\mathrm{Sd}$]
Let $T\ge0$ be an integer and $n=8hL^{T+N}=8L^{T-T^{*}(e)}\cdot n_*$.
If $T\ge T^{*}(e)+1$, then $n\in\mathrm{Sd}$ by 8.4\,(3).
If $T=T^{*}(e)$, then $n=8n_*$ is the side of the torus $\mathrm{kt}=8$, and by 8.4\,(1) it lies in $\mathrm{Sd}$ if and only if $8\ge L$.
If $T<T^{*}(e)$, then $n\le8n_*/L$.
Assume now that $L>8$ (Proposition~14.2\,(f), consequences of (L1)).
Then $1<8<L$, so that $8n_*\notin\mathrm{Sd}$;
moreover $8n_*/L<n_*$, whereas every element of $\mathrm{Sd}$ is $\ge n_*$.
Consequently $n\in\mathrm{Sd}$ if and only if $T\ge T^{*}(e)+1$.
For $T\ge T^{*}(e)$ this is the equivalence ``$\mathrm{kt}\ge L$ if and only if $T\ge T^{*}(e)+1$'' of Section~I.6.5 and of Proposition~47.2\,(a).
\end{named}

\subsubsection*{9. What the bounds depend on}

\begin{itemize}
\item[(1)]
The clauses (a)--(d), (c+) and (d+) hold for every real $\boldsymbol{\mu}$, every $q>0$, every coordinate cell structure $\mathcal{Q}$ of Definition~39.0, every even $\mathrm{ncs}$ with $\mathrm{hw}\cdot\mathrm{ncs}\ge4$ and every integer $\check{j}\ge2$;
in particular, the sign of $\boldsymbol{\mu}$ is arbitrary.
Apart from $\rho$ and the weights, which enter through $Z_n[\,\cdot\,]$, no constant of Section~I.8 occurs in them.
\item[(2)]
The numbers $Z_{[\mathrm{sv}]}$, $D_{\check{j}}$ and $\delta_{\check{j}}$ are formed with the constant $1$.
The numbers $Z_{[\mathrm{sv}]}$ and $D_{\check{j}}$ depend on $\rho$, the weights, $q$ and $\boldsymbol{\mu}$, and on the sides $\mathrm{hw}\cdot\mathrm{sv}$, $n_r$, $\check{j}\cdot n_r$;
the number $\delta_{\check{j}}=D_{\check{j}}/\left((\check{j}-1)^3\cdot N_r\right)$ contains in addition $N_r=\mathrm{ncs}^3$, that is, the width $\mathrm{hw}=n_r/\mathrm{ncs}$.
None of them depends on the offsets $o$.
The quantity $\|O\|_*$ depends on $O$, on $\mathcal{Q}$ and on the single torus $C(n_r)$.
\item[(3)]
The right side of (d) contains the torus $C(\mathrm{hw}\cdot\mathrm{sv})$ only through the index set $\mathrm{Scl}$:
the factor $\exp(\delta_{\check{j}})\cdot\|O_{\mathfrak{c}}\|_*$ of a cell is formed on the two tori $C(n_r)$ and $C(\check{j}\cdot n_r)$.
Therefore (d) is uniform in $\mathrm{sv}$ and in $\mathrm{Scl}$.
In the setting of (e), for fixed $(e,N,\boldsymbol{\mu})$ the factor $\exp(\delta)\cdot\|O\|_*$ does not depend on $n\in\mathrm{Sd}$, that is, neither on $\mathrm{kt}$ nor on $T$.
\end{itemize}

\subsubsection*{10. Proof of (c+), (d+), (e+)}

In this section the data and the notation are those of Proposition~40.

\begin{proof}[Proof of Proposition~40\,(c+), (d+), (e+)]
(c+)
Let $\mathrm{sv}$ be an even integer with $\mathrm{ncs}\le\mathrm{sv}\le\check{j}\cdot\mathrm{ncs}$.
By Section~3.1, the numbers $\mathrm{kr}=\mathrm{sv}/\mathrm{ncs}$ and $\check{j}$ lie in $\mathrm{Kr}$, and $1\le\mathrm{kr}\le\check{j}$.
Lemma~2\,(i) for the function $\mathrm{fz}$ of Section~3, at $\mathrm{kr}'=\mathrm{kr}$ and $\mathrm{kr}''=\check{j}$, gives
\[
\mathrm{fz}(\check{j},\check{j},\check{j})\le\left(\frac{\check{j}}{\mathrm{kr}}\right)^3\cdot\mathrm{fz}(\mathrm{kr},\mathrm{kr},\mathrm{kr}).
\]
Together with (3.1) this yields
\begin{align*}
\log Z_{[\mathrm{sv}]}=\mathrm{fz}(\mathrm{kr},\mathrm{kr},\mathrm{kr})
&\ge\left(\frac{\mathrm{kr}}{\check{j}}\right)^3\cdot\mathrm{fz}(\check{j},\check{j},\check{j})
=\left(\frac{\mathrm{kr}}{\check{j}}\right)^3\cdot\left(\check{j}^3\cdot\log Z_{[\mathrm{ncs}]}-D_{\check{j}}\right)\\
&=\mathrm{kr}^3\cdot\left[\log Z_{[\mathrm{ncs}]}-\frac{D_{\check{j}}}{\check{j}^3}\right].
\end{align*}
The second inequality of (c+) follows from $D_{\check{j}}\ge0$, which is (a), and from $\check{j}^3\ge(\check{j}-1)^3>0$.

(d+)
Let $\mathrm{sv}\ge\mathrm{ncs}$ be an even integer, and let $\mathrm{Scl}$ and the $O_{\mathfrak{c}}$ be as in (d+).
In Section~7 the hypothesis ``$\mathrm{sv}=\mathrm{ncs}$, or $\mathrm{sv}\ge\check{j}\cdot\mathrm{ncs}$'' enters at one place only, namely in (7.3), through the second form of (c).
By (c) and (c+), this second form,
\[
\log Z_{[\mathrm{sv}]}\ge\mathrm{kr}^3\cdot\left[\log Z_{[\mathrm{ncs}]}-\frac{D_{\check{j}}}{(\check{j}-1)^3}\right],
\]
holds for every even $\mathrm{sv}\ge\mathrm{ncs}$.
The constructions and estimates of Sections~7.1 and~7.2, and the inequality (7.2), hold for every even $\mathrm{sv}\ge\mathrm{ncs}$.
Hence (7.3), (7.4) and the conclusion of Section~7.4 hold for such $\mathrm{sv}$, and they give the inequality of (d).

(e+)
By Section~8.3, for $n\in\mathrm{per}\,\mathbb{Z}$ the condition $n\ge n_*$ says that $\mathrm{sv}=2n/\mathrm{per}$ is an even integer $\ge\mathrm{ncs}$;
so (d+) applies on the cubic torus with $n$ sites per side.
For every integer $\mathrm{kt}\ge1$, the number $\mathrm{kt}\cdot n_*$ lies in $\mathrm{per}\,\mathbb{Z}$ and is $\ge n_*$.
For $T\ge T^{*}(e)$ one has $8hL^{T+N}=8L^{T-T^{*}(e)}\cdot n_*$, a multiple of $n_*$ of this kind.
\end{proof}

\begin{named}{Remark 10.1}[The first form of (c) between $\mathrm{ncs}$ and $\check{j}\cdot\mathrm{ncs}$; one variable]
Let $\mathrm{wr}$ be a real function on a subset of $(0,\infty)$ which has the property (cv) there, let $1<\mathrm{kr}<\check{j}$ be points of this subset, and put $\alpha:=(\mathrm{kr}-1)/(\check{j}-1)$.
Then $0<\alpha<1$, and
\[
\mathrm{wr}(\mathrm{kr})\le\mathrm{kr}\cdot\mathrm{wr}(1)-\alpha\cdot\left(\check{j}\cdot\mathrm{wr}(1)-\mathrm{wr}(\check{j})\right).
\]
Indeed, the identity $\alpha\cdot(\check{j}-1)=\mathrm{kr}-1$ gives $\mathrm{kr}-\alpha\cdot\check{j}=1-\alpha$, so that the right side equals $(1-\alpha)\cdot\mathrm{wr}(1)+\alpha\cdot\mathrm{wr}(\check{j})$.
For the three points $1<\mathrm{kr}<\check{j}$ the property (cv) reads
\[
(\check{j}-1)\cdot\mathrm{wr}(\mathrm{kr})\le(\check{j}-\mathrm{kr})\cdot\mathrm{wr}(1)+(\mathrm{kr}-1)\cdot\mathrm{wr}(\check{j});
\]
division by $\check{j}-1>0$, together with $(\check{j}-\mathrm{kr})/(\check{j}-1)=1-\alpha$, gives $\mathrm{wr}(\mathrm{kr})\le(1-\alpha)\cdot\mathrm{wr}(1)+\alpha\cdot\mathrm{wr}(\check{j})$.
Thus, in one variable, the inequality that corresponds to the first form of (c) holds between $1$ and $\check{j}$ with the opposite sign.
\end{named}

\subsubsection*{Conditions on the constants}
\leavevmode\par\nopagebreak

\noindent
\begin{tabular}{p{0.2\textwidth}p{0.27\textwidth}p{0.4\textwidth}}
condition & line of Section~I.8.2; Proposition~14.2 & where it is used\\
\hline
$\rho\in(0,1)$; the standard weights
& (C0)
& standing data of [Q1]--[Q4] (Proposition~39), of [Q5] (Proposition~50\,(e)) and of [Q6] (the standing hypotheses $-1\le\rho\le1$, $w_i>0$ of \cite[Section~6]{QEDII}); they enter only through these statements\\
$L$ an integer $\ge5$
& Section~I.0.5; (C3); Proposition~14.2\,(b), (L5), a consequence of (L1)
& Section~8.1: $\check{j}=L\ge2$, $n_*\ge L\ge4$; 8.4\,(3): $8L\ge L$; 8.4\,(4): $L\ge2$\\
$M$ a power of $L$, $M\ge L^2$
& Section~I.0.5; (C5); Proposition~14.2\,(b), (M1)
& Section~8.1: $T^{*}(e)$ is an integer $\ge1$, so that $\mathrm{per}$ is an integer\\
$\bar{m}\in(0,1]$
& (C10)
& Section~8.1: $h$ is a power of $L$ with a non-negative exponent\\
$L>8$
& (C3); Proposition~14.2\,(f), consequences of (L1); Annex~B, Section~G6.Z4, first row
& Remark~8.5\\
\end{tabular}

\medskip
All conditions of the table are among those of Section~I.8.2 and Proposition~14.2, at the places named in the second column.
The clauses (a)--(d), (c+) and (d+) require no condition on a constant beyond the standing data of the first row, $\rho\in(0,1)$ and the standard weights, and these data enter them only through [Q1]--[Q6].
The clauses (e) and (e+) require in addition the conditions of the second, third and fourth rows: $L$ is an integer $\ge5$, $M$ is a power of $L$ with $M\ge L^2$, and $\bar{m}\in(0,1]$.
By Section~8.1 these three conditions make $\mathcal{Q}=\mathcal{Q}_*$, $\mathrm{hw}=\mathrm{hc}$, $\mathrm{ncs}=2L$ and $\check{j}=L$ admissible data of Proposition~40:
$\mathrm{per}$ is an integer $\ge1$, $\mathrm{hc}\in\frac{1}{2}\mathbb{Z}$, $n_r=n_*\ge4$ and $\check{j}=L\ge2$.
In the proof of (e) they enter through Section~8.1, on which Sections~8.2 and~8.3 rest, and through 8.4\,(3) and~(4);
in the proof of (e+) they enter only through Sections~8.1 and~8.3.
Every clause of Proposition~40 holds without the condition $L>8$ of the last row, which is assumed in Remark~8.5 alone.

\section*{Chapter 42. The defect of the two reference tori and the transfer inequality for Wick powers: proofs of Propositions 42.1 and 42.2}

\subsection*{0. Conventions of this chapter}

\paragraph*{0.1. Setting.}
Everything takes place at level~$0$ and in lattice units (Sections~I.2.5 and~I.7.1): $\eta_0=1$, the charge is $q$, and $r_q=\frac{2\pi}{q}$.
The tori are the cubic tori $C(n,n,n)$ with $n\ge4$ photon sites per side, of either parity, and the spinors on them are antiperiodic (Section~I.6.6).
The fermion operator enters only through the number $\det\left(D_A+\boldsymbol{\mu}\right)$, the determinant of a linear map of the space $\mathfrak{S}^{\mathrm{ch}}(1;n)$, of dimension $192\cdot n^3$, into itself (Proposition~46\,(e)).
In turn, this number enters only the functionals $Z_n[\,\cdot\,]$ and $Z^{(w)}_n[\,\cdot\,]$ (Sections~I.7.2 and~I.7.3) and the numbers $\mathrm{zr}_n$ and $\mathring{\Phi}$ (Sections~I.7.3 and~I.7.6).
The Wilson parameter $\rho\in(0,1)$ and the standard weights are those of Section~I.0.5, and $M_{\mathrm{ch}}=\boldsymbol{\mu}+\frac{\rho}{2}\cdot\mathrm{sw}$.

\paragraph*{0.2. The objects of \cite{QEDII} on the cubic torus.}
Throughout, $\mathbb{T}=C(n,n,n)$.
The mass parameter of \cite[Definition~2.10]{QEDII} is $M_{\mathrm{ch}}$, and the curvature $F$ of \cite{QEDII} is $F^{\mathrm{u}}$.
The functional $Z_{\mathbb{T}}$ and the expectation $\mathbb{E}_{\mathbb{T}}$ of \cite[Definition~2.13]{QEDII} coincide with $Z_n$ and $E^{\mathrm{af}}_n$ (Proposition~46\,(f); Section~I.7.2).
In \cite[Definition~13.7]{QEDII} the coordinate two-tori are written $\Sigma_{\mu\nu}$ and the flux is written $k=(k_{23},k_{13},k_{12})$;
the indicator of the flux sector $k$ is $\Phi_k$, its sector functional is $Z^{(k)}_{\mathbb{T}}$, and its weight is the number $\mathbb{E}_{\mathbb{T}}[\Phi_k]$ of \cite[Definition~13.7(a)]{QEDII};
the two-cochain whose value at every plaquette $(x;\mu\nu)$ is $\frac{2\pi}{q}\cdot\frac{k_{\mu\nu}}{n_\mu n_\nu}$ is written $\mathfrak{h}_k$ \cite[Definition~13.7(d)]{QEDII}.
Both $\Sigma_{\mu\nu}$ and the set $\Pi_{\mu\nu}$ of Section~I.7.1 are the set of the $n_\mu n_\nu$ plaquettes in one coordinate plane parallel to $e_\mu$ and $e_\nu$, and in both places the flux is the vector of the sums of the integer two-cochain over these sets, in the order $(23,13,12)$;
so the flux $k$ of a configuration is its flux $w$ of Section~I.7.1, and $\Phi_w=\mathbf{1}[\mathrm{flux}=w]$ (Section~I.7.1).
By Proposition~46\,(f+), for every $w\in\mathbb{Z}^3$ the sector functional $Z^{(w)}_{\mathbb{T}}$ coincides with $Z^{(w)}_n$, the weight $\mathbb{E}_{\mathbb{T}}[\Phi_w]$ equals $P_n(w)$, and $\mathfrak{h}_w$ is the cochain $h_w$ of Section~I.7.1.
For a box $\mathrm{Bx}$, the sets $\mathrm{Ad}(\mathrm{Bx})$ and $\mathrm{Ad}^{\mathrm{ev}}(\mathrm{Bx})$ are the set $\mathcal{G}(B)$ of \cite[Definition~2.12]{QEDII} and the set $\mathcal{G}^{\mathrm{ev}}(B)$ of its even elements, for $B=\mathrm{Bx}$.
The standing hypotheses of \cite[Section~13.3]{QEDII}, namely $-1\le\rho\le1$, weights $w_i>0$, a real mass parameter and $q>0$, together with a torus $C(n_1,n_2,n_3)$ with $n_\mu\ge4$, are the hypotheses of \cite[Definition~13.7, Lemma~13.12 and Theorem~6.27]{QEDII}.
The data of 0.1 satisfy them.

\paragraph*{0.3. Notation.}
Throughout this chapter $\mathrm{xq}$, $\mathrm{tq}$ and $\mathrm{uq}$, as well as $\mathrm{yq}$ with or without an index, denote real numbers; $\mathrm{Pol}_{m,\mathrm{Scl}}$ is the polynomial defined in (42.5); $\hat{\mathfrak{a}}^{\mathrm{o}}$ is the one-cochain on $\mathbb{Z}^3$ belonging to the data of Proposition~41; and $\mathrm{gc}$ is a real two-cochain.
The letter $i$ denotes an integer, serving as an index of summation or as an exponent, and $z$ denotes a real variable, as in Section~I.7.6.
In Appendix~A, $\mathrm{fq}$ denotes a function, $\mathrm{oq}$, $\mathrm{Mk}$ and $\mathrm{cj}$ denote real numbers, $\mathrm{ks}$ a pair of integers and $\mathrm{zeq}$ a root of unity, while $k$ is an index of summation: a class in $\mathbb{Z}/n_1$ in (Ax1), and elsewhere the triple $(k_1,k_2,k_3)$ of the display that defines $\mathrm{Ld}(n)$ in Proposition~38.2\,(c).
The exponent denoted $r$ in \cite[Theorem~6.27(b)]{QEDII} is written $i$ here; the letter $r$ alone denotes the integer of the line (C7) of Section~I.8.2.
The letter $v$ alone denotes the placement; the condition (e8) of Proposition~14.2\,(b), which is formulated there with a real variable, is evaluated at the number $v_K$.
For an integer $i\ge1$ and a Borel function $G$ on $C^2(n_*)$ with $E^{\mathrm{af}}_{n_*}|G|^i<\infty$ let
\begin{equation}
\|G\|_{[i]}:=\left(E^{\mathrm{af}}_{n_*}|G|^i\right)^{\frac{1}{i}}.
\tag{0.1}
\end{equation}
The following numerical bounds hold:
\begin{gather*}
2<\mathrm{e}<2.7183,\qquad\mathrm{e}^2<7.3892,\qquad3<\pi,\qquad\pi^2<10,\qquad\log3>1,\\
3.1415<\pi<3.1416,\qquad1.4142<2^{\frac12}<1.41422,\qquad1.7320<3^{\frac12}.
\end{gather*}

\subsection*{1. Proposition 42.1}

\begin{proof}[Proof of Proposition~42.1]
Both $n$ and $\check{j}\cdot n$ are integers $\ge4$.
The statements applied below, namely Proposition~38.2\,(c), Proposition~46\,(f+), Proposition~50\,(e) and \cite[Lemma~13.12(a)]{QEDII}, hold on every cubic torus with at least four sites per side, of either parity, for every $q>0$ and every real $\boldsymbol{\mu}$.
They therefore apply to the side $n$ and to the side $\check{j}\cdot n$, with one and the same $q$ and $\boldsymbol{\mu}$.

\emph{(a), the first identity.}
Set $\eta:=1$, $m_{\mathrm{in}}:=\boldsymbol{\mu}-\bar{m}$ and $\varepsilon^{\mathrm{v}}_{\mathrm{in}}:=0$.
Then $\left(m_{\mathrm{in}},\varepsilon^{\mathrm{v}}_{\mathrm{in}}\right)$ is a real pair with $\bar{m}\cdot\eta+m_{\mathrm{in}}=\boldsymbol{\mu}$, and Proposition~38.2\,(c) yields
\begin{equation}
Z^{(0)}_n[1]=Z^{\mathrm{ph}}_n\cdot\mathrm{zr}_n,\qquad\mathrm{zr}_n>0,\qquad Z^{\mathrm{ph}}_n=r_q^{n^3+2}\cdot(2\pi)^{n^3-1}\cdot\exp\left(-\frac{\mathrm{Ld}(n)}{2}\right)>0.
\tag{1.1}
\end{equation}

\emph{(a), the second identity.}
Proposition~50\,(e) gives $0<Z^{(0)}_n[1]\le Z_n[1]<\infty$, and Section~I.7.3 defines $P_n(0)=\frac{Z^{(0)}_n[1]}{Z_n[1]}$.
Combined with (1.1), this yields
\begin{equation}
0<P_n(0)\le1,\qquad Z_n[1]=\frac{Z^{(0)}_n[1]}{P_n(0)}=\frac{Z^{\mathrm{ph}}_n\cdot\mathrm{zr}_n}{P_n(0)}.
\tag{1.2}
\end{equation}

\emph{(a), the third assertion.}
For $\mathbb{T}=C(n,n,n)$, \cite[Lemma~13.12(a)]{QEDII} asserts that $\mathbb{E}_{\mathbb{T}}[\Phi_k]\le\exp\left(-\frac12\|\mathfrak{h}_k\|^2\right)$ for every $k\in\mathbb{Z}^3$.
Let $w\in\mathbb{Z}^3$ and take $k=w$.
By Proposition~46\,(f+) (paragraph~0.2), $\mathbb{E}_{\mathbb{T}}[\Phi_w]=P_n(w)$, the cochain $\mathfrak{h}_w$ is the cochain $h_w$ of Section~I.7.1, and
\[
\exp\left(-\frac12\|h_w\|^2\right)=\exp\left(-\frac{r_q^2}{2|n|}\cdot\sum_{\mu<\nu}n_{\mathfrak{l}}^2\cdot w_{\mu\nu}^2\right),
\]
where $\mathfrak{l}$ is the direction different from $\mu$ and $\nu$.
On the cubic torus $|n|=n^3$ and $n_{\mathfrak{l}}=n$ for $\mathfrak{l}=1,2,3$; consequently
\begin{align}
&\frac{r_q^2}{2n^3}\cdot n^2\cdot\left(w_{23}^2+w_{13}^2+w_{12}^2\right)=\left(\frac{2\pi}{q}\right)^2\cdot\frac{|w|_2^2}{2n}=\frac{2\pi^2\cdot|w|_2^2}{q^2\cdot n},\notag\\
&\text{and hence}\quad P_n(w)\le\mathfrak{q}_w(n)\quad\text{for every }w\in\mathbb{Z}^3.\tag{1.3}
\end{align}

\emph{(a), the fourth assertion.}
The upper bound $P_n(0)\le1$ is part of (1.2).
A closed $\mathfrak{z}\in C^2(n;\mathbb{Z})$ has exactly one flux $w\in\mathbb{Z}^3$ (Section~I.7.1), so that $\sum_{w\in\mathbb{Z}^3}\mathbf{1}[\mathrm{flux}=w]=1$ at every configuration.
According to Proposition~50\,(e), the sum over $\mathfrak{z}$ and the integral over $A$ in (I.7.1) converge absolutely for $G=1$.
The terms of an absolutely convergent series may be grouped, and the group formed by the terms $\mathfrak{z}$ of flux $w$ is $Z^{(w)}_n[1]=Z_n\left[\mathbf{1}[\mathrm{flux}=w]\right]$ (Section~I.7.3).
Thus $Z_n[1]=\sum_wZ^{(w)}_n[1]$, and division by $Z_n[1]$ gives $\sum_wP_n(w)=1$.
Together with (1.3) this implies
\begin{equation}
P_n(0)=1-\sum_{w\ne0}P_n(w)\ge1-\sum_{w\ne0}\mathfrak{q}_w(n)=1-\varsigma_n.
\tag{1.4}
\end{equation}
The series $\varsigma_n$ converges by (1.5).

\emph{(b).}
Let $\mathrm{xq}:=\exp\left(-\frac{2\pi^2}{q^2\cdot n}\right)\in(0,1)$.
Then $\mathfrak{q}_w(n)=\mathrm{xq}^{w_{23}^2}\cdot\mathrm{xq}^{w_{13}^2}\cdot\mathrm{xq}^{w_{12}^2}$; since all factors are non-negative, the sum of this product over $w\in\mathbb{Z}^3$ is the product of three equal sums.
As $i^2\ge i$ and $0<\mathrm{xq}<1$, one obtains
\begin{align}
\varsigma_n&=\left(\sum_{i\in\mathbb{Z}}\mathrm{xq}^{i^2}\right)^3-1,\notag\\
\sum_{i\in\mathbb{Z}}\mathrm{xq}^{i^2}&=1+2\sum_{i\ge1}\mathrm{xq}^{i^2}\le1+2\sum_{i\ge1}\mathrm{xq}^i=1+\mathrm{tq},\qquad\mathrm{tq}:=\frac{2\cdot\mathrm{xq}}{1-\mathrm{xq}}.\tag{1.5}
\end{align}
Assume now that $q^2\cdot n\le1$.
Then $\frac{2\pi^2}{q^2\cdot n}\ge2\pi^2>18$, and since $\mathrm{e}>2$ it follows that $\mathrm{xq}<2^{-18}$.
For $0\le z\le\frac12$ one has $(1+2z)\cdot(1-z)=1+z\cdot(1-2z)\ge1$, that is, $\frac{1}{1-z}\le1+2z$.
Hence $\frac{1}{1-\mathrm{xq}}\le1+2^{-17}<1.0005$ and $\mathrm{tq}\le2.001\cdot\mathrm{xq}<10^{-5}$, so that
\begin{equation}
\varsigma_n\le(1+\mathrm{tq})^3-1=\mathrm{tq}\cdot\left(3+3\cdot\mathrm{tq}+\mathrm{tq}^2\right)\le2.001\cdot\mathrm{xq}\cdot3.0001\le7\cdot\mathrm{xq},\qquad7\cdot\mathrm{xq}<7\cdot2^{-18}<\frac18.
\tag{1.6}
\end{equation}

\emph{(c), the identity.}
Applied to the sides $n$ and $\check{j}\cdot n$, (1.1) and (1.2) show that the numbers $Z^{\mathrm{ph}}$, $\mathrm{zr}$ and $P(0)$ of both sides are positive, and that
\[
\log Z_n[1]=\log Z^{\mathrm{ph}}_n+\log\mathrm{zr}_n-\log P_n(0),\qquad\log Z_{\check{j}\cdot n}[1]=\log Z^{\mathrm{ph}}_{\check{j}\cdot n}+\log\mathrm{zr}_{\check{j}\cdot n}-\log P_{\check{j}\cdot n}(0).
\]
Hence
\begin{align}
\check{j}^3\cdot\log Z_n[1]-\log Z_{\check{j}\cdot n}[1]={}&\left[\check{j}^3\cdot\log Z^{\mathrm{ph}}_n-\log Z^{\mathrm{ph}}_{\check{j}\cdot n}\right]+\left[\check{j}^3\cdot\log\mathrm{zr}_n-\log\mathrm{zr}_{\check{j}\cdot n}\right]\notag\\
&+\left[-\check{j}^3\cdot\log P_n(0)+\log P_{\check{j}\cdot n}(0)\right].\tag{1.7}
\end{align}
The second bracket is $D^{\mathrm{mat}}_{\check{j}}$, and the third is $D^{\mathrm{fl}}$.
The third relation in (1.1) gives $\log Z^{\mathrm{ph}}_n=\left(n^3+2\right)\cdot\log r_q+\left(n^3-1\right)\cdot\log(2\pi)-\frac12\mathrm{Ld}(n)$, and the same formula holds for the side $\check{j}\cdot n$, with $(\check{j}\cdot n)^3=\check{j}^3\cdot n^3$.
It follows that
\begin{align*}
\check{j}^3\cdot\log Z^{\mathrm{ph}}_n-\log Z^{\mathrm{ph}}_{\check{j}\cdot n}={}&\left[\check{j}^3\cdot\left(n^3+2\right)-\left(\check{j}^3\cdot n^3+2\right)\right]\cdot\log r_q+\left[\check{j}^3\cdot\left(n^3-1\right)-\left(\check{j}^3\cdot n^3-1\right)\right]\cdot\log(2\pi)\\
&-\frac12\left[\check{j}^3\cdot\mathrm{Ld}(n)-\mathrm{Ld}(\check{j}\cdot n)\right]\\
={}&\left(\check{j}^3-1\right)\cdot\left[2\log r_q-\log(2\pi)\right]+\frac12\left[\mathrm{Ld}(\check{j}\cdot n)-\check{j}^3\cdot\mathrm{Ld}(n)\right],
\end{align*}
and $2\log r_q-\log(2\pi)=\log\frac{(2\pi)^2}{2\pi\cdot q^2}=\log\frac{2\pi}{q^2}$.
So the first bracket in (1.7) is $D^{\mathrm{ph}}$.

\emph{(c), the bound on $D^{\mathrm{ph}}$.}
The inequality of \cite[Lemma~21.31]{QEDI} (Lemma~A of Appendix~A) holds for all integers $n\ge2$ and $j\ge2$.
In it, the number $\mathrm{Ld}(n)$ of \cite[Section~21.8.1]{QEDI} is the number $\mathrm{Ld}(n)$ of Proposition~38.2\,(c), and at $j=\check{j}$ the constants $G_j$ and $K_j$ of \cite[Section~21.1.7]{QEDI} are the numbers $G_{\check{j}}$ and $K_{\check{j}}$ of Proposition~42.1.
The inequality reads
\[
\frac12\left[\mathrm{Ld}(\check{j}\cdot n)-\check{j}^3\cdot\mathrm{Ld}(n)\right]\le-\left(\check{j}^3-1\right)\cdot\log n+K_{\check{j}},
\]
and it implies
\[
D^{\mathrm{ph}}\le\left(\check{j}^3-1\right)\cdot\left[\log\frac{2\pi}{q^2}-\log n\right]+K_{\check{j}}=\left(\check{j}^3-1\right)\cdot\log\frac{2\pi}{q^2\cdot n}+K_{\check{j}}.
\]

\emph{(c), the bound on $D^{\mathrm{fl}}$.}
Suppose that $\varsigma_n<1$.
For the side $\check{j}\cdot n$, (1.2) gives $\log P_{\check{j}\cdot n}(0)\le0$.
By (1.4), $P_n(0)\ge1-\varsigma_n>0$, whence $-\log P_n(0)\le-\log\left(1-\varsigma_n\right)$.
Therefore $D^{\mathrm{fl}}\le-\check{j}^3\cdot\log\left(1-\varsigma_n\right)$.
\end{proof}

\paragraph*{Dependence.}
The relations (1.1)--(1.5) and (1.7) hold for every quadruple $(n,\check{j},q,\boldsymbol{\mu})$, and (1.6) holds whenever $q^2\cdot n\le1$.
Neither the numerals $7$ and $\frac18$ nor the numbers $G_{\check{j}}$ and $K_{\check{j}}$ contain $n$, $q$, $\boldsymbol{\mu}$, $\rho$ or the weights.

\subsection*{2. Proposition 42.2}

\noindent\textit{Proof of Proposition~42.2.}
Here (2.1)--(2.4$'$) are the displays of the statement of Proposition~42.2, and part~0.1, \dots, part~0.7 are the seven parts of Step~0 below.

\textbf{Step 0. Preliminaries.}

\emph{0.1. The range of the charge.}
By Section~I.0.5, $h=L^{\left\lfloor\log_L\frac{1}{\bar{m}}\right\rfloor}$ with $0<\bar{m}\le1$; hence $h=L^i$ for an integer $i\ge0$, $h\ge1$ and $\log h\ge0$.
The functions $e\mapsto e\cdot h^{\frac12}$ and $e\mapsto e\cdot h^{\frac12}\cdot\log h$ are therefore continuous and non-decreasing on $(0,\infty)$.
Consider the set
\[
\left\{e>0:\ e\cdot h^{\frac12}\le e_1,\ e\cdot h^{\frac12}\cdot\log h\le1\right\},
\]
whose supremum is $e_G$ (line (C10) of Section~I.8.2; Proposition~14.2\,(a), row $e_G$).
It is non-empty, because $e_1>0$ (Section~I.0.5, Proposition~14.3\,(f)) and both functions tend to $0$ as $e\to0$; it is downward closed, closed in $(0,\infty)$, and bounded by $e_1\cdot h^{-\frac12}$.
Consequently it is the interval $(0,e_G]$.
As moreover $e_1\le\frac13$ (Proposition~14.2\,(f)), the charge $e$ of the statement satisfies
\begin{align}
&e\cdot h^{\frac12}\le e_1,\qquad e\cdot h^{\frac12}\cdot\log h\le1\quad\text{(this is the condition (Z8.w) of Proposition~14.2\,(d))},\notag\\
&e\le e\cdot h^{\frac12}\le e_1\le\frac13.\tag{42.1}
\end{align}
Thus $e$ and $N$ are admissible data for Section~I.6.4, for Proposition~40\,(e) and for Proposition~41: $e\cdot h^{\frac12}\le e_1$ and $N\ge N_0(e)$.

\emph{0.2. The condition (W3).}
Let $v_K=-\log\left(e\cdot h^{\frac12}\right)$ as in Section~I.6.1.
From (42.1) one gets $v_K\ge-\log e_1\ge\log3>1$.
By Section~I.6.1, $r_K=\min\left\{L^i\ge v_K^r:\ i\ge0\text{ an integer}\right\}$, where $r\ge2$ is the integer of the line (C7).
Because $v_K^r>1=L^0$, this minimum is attained at some $i\ge1$, and for this $i$ one has $L^{i-1}<v_K^r$; so $r_K<L\cdot v_K^r$.
Section~I.6.2 gives $R^*=\frac{R_{\mathrm{top}}(e)}{L}=\frac{h\cdot M\cdot r_K}{L}<h\cdot M\cdot v_K^r$.
Furthermore $e^2\cdot h=\exp\left(-2v_K\right)$ and $\log\frac1e=v_K+\frac12\log h>0$.
With the second inequality of (42.1) in the last step, this gives
\begin{align*}
e^2\cdot R^*\cdot\log\frac1e&<M\cdot v_K^r\cdot\exp\left(-2v_K\right)\cdot\left(v_K+\frac12\log h\right)\\
&=M\cdot v_K^r\cdot\exp\left(-v_K\right)\cdot\left(v_K\cdot\exp\left(-v_K\right)+\frac12e\cdot h^{\frac12}\cdot\log h\right)\\
&\le M\cdot v_K^r\cdot\exp\left(-v_K\right)\cdot\left(v_K\cdot\exp\left(-v_K\right)+\frac12\right).
\end{align*}
The condition (e8) of Proposition~14.2\,(b) is imposed at every real number $\ge-\log e_1$, and by Proposition~14.3\,(f) it holds at every such number.
Since $v_K\ge-\log e_1$, it holds at $v_K$, where it states that
\[
6N_*\cdot M\cdot v_K^r\cdot\exp\left(-v_K\right)\cdot\left(v_K\cdot\exp\left(-v_K\right)+\frac12\right)\le\pi^2.
\]
Therefore
\begin{equation}
6N_*\cdot e^2\cdot R^*\cdot\log\frac1e\le\pi^2.
\tag{W3}
\end{equation}

\emph{0.3. Arithmetic.}
As $L\ge5$ (Section~I.0.5), $N_*=8L^3\ge1000$; and $\log\frac1e\ge\log3>1$ by (42.1).
Division of (W3) by $6N_*\cdot\log\frac1e$, and division of (W3) by $e^2\cdot R^*$, give respectively
\begin{equation}
e^2\cdot R^*\le\frac{\pi^2}{6N_*}\cdot\left(\log\frac1e\right)^{-1}<\frac{10}{6000}<1,\qquad\frac{\pi^2}{e^2\cdot R^*}\ge6N_*\cdot\log\frac1e.
\tag{42.2}
\end{equation}
Let $\mathrm{xq}:=\exp\left(-\frac{2\pi^2}{e^2\cdot R^*}\right)$.
The second inequality of (42.2) yields
\begin{equation}
\mathrm{xq}\le\exp\left(-12N_*\cdot\log\frac1e\right)=e^{12N_*},\qquad\mathrm{xq}^{\frac{1}{2N_*}}=\exp\left(-\frac{\pi^2}{N_*\cdot e^2\cdot R^*}\right)\le\exp\left(-6\log\frac1e\right)=e^6.
\tag{42.3}
\end{equation}

\emph{0.4. Sides and cells.}
Since $M$ and $r_K$ are powers of $L$ with $M\ge L^2$ and $r_K\ge1$, the number $T^*(e)=\log_L\left(M\cdot r_K\right)-1$ is an integer $\ge1$ (Section~I.6.2).
Moreover $h$ is a power $L^i$ with an integer $i\ge0$ (part~0.1), and $N\ge0$.
Hence $n_*=\frac{R^*}{\eta}=h\cdot L^{T^*(e)+N}$ is an integer $\ge L\ge5$, the number $\mathrm{per}=\frac{n_*}{L}=h\cdot L^{T^*(e)+N-1}$ is an integer $\ge1$, and
\[
\mathrm{hc}=\frac{\mathrm{per}}{2}\ge\frac12,\qquad n_*=L\cdot\mathrm{per}=2L\cdot\mathrm{hc},\qquad q^2\cdot n_*=\frac{e^2\cdot\eta\cdot R^*}{\eta}=e^2\cdot R^*.
\]
Both $n_*$ and $L\cdot n_*$ are $\ge4$.
A side $n$ of the statement, that is, $n=n_*$ or $n\in\mathrm{per}\cdot\mathbb{Z}$ with $n\ge L\cdot n_*$, lies in $\mathrm{per}\cdot\mathbb{Z}$ and satisfies $n\ge n_*$; so it is a side admitted in Proposition~41, and it is among the sides named in Proposition~40\,(e).
The cells of the torus $n$ are the classes of the $\mathfrak{c}\in\mathbb{Z}^3$ modulo $\frac{2n}{\mathrm{per}}\cdot\mathbb{Z}^3$; the torus $n_*$ has $(2L)^3=N_*$ cells, and their set is $\mathrm{Cel}$ (Section~I.6.4).

\emph{0.5. The cell observables.}
The data of Proposition~41 and of Proposition~38.3\,(e) are at hand: (42.1) holds, $N\ge N_0(e)$, $\beta$ fits $\left(N,T^*(e)\right)$, $v$ is a placement, and $n$ is a side as in part~0.4.
By Proposition~41\,(d), $O_m=e^{-m}\cdot H_m(X)-c_m$ is real and is a bosonic element of $\mathrm{Ad}\left(\mathrm{Bx}_0\right)$, hence of $\mathcal{A}\left(Q^{\circ}_0\right)$; moreover
\begin{equation}
O_m^{[\mathfrak{c}]}=e^{-m}\cdot H_m\left(X_{\mathfrak{c}}\right)-c_m\qquad\text{for every }\mathfrak{c}\in\mathbb{Z}^3\quad(m=2,4).
\tag{42.4}
\end{equation}
The map $\mathsf{g}_0$ is the identity (Section~I.6.4 at $\mathfrak{c}=0$), so that $\hat{\mathfrak{a}}_0=\hat{\mathfrak{a}}$, $X_0=X$ and $O_m^{[0]}=O_m$.
On the torus $n$, the function $X_{\mathfrak{c}}$ depends only on the cell of $\mathfrak{c}$.
To see this, let $\mathfrak{c}'=\mathfrak{c}+\frac{2n}{\mathrm{per}}\cdot\mathfrak{c}''$ with $\mathfrak{c}''\in\mathbb{Z}^3$.
As $\frac{2n}{\mathrm{per}}$ is even, every coordinate $\mathfrak{c}_\mu$ keeps its parity, while the integer $i$ of Section~I.6.4 (given by $\mathfrak{c}_\mu=2i$ or $\mathfrak{c}_\mu=2i+1$) changes by $\frac{n}{\mathrm{per}}\cdot\mathfrak{c}''_\mu$.
So $\mathsf{g}_{\mathfrak{c}'}$ is $\mathsf{g}_{\mathfrak{c}}$ followed by the translation by $n\cdot\mathfrak{c}''$, and $\left(\mathsf{g}_{\mathfrak{c}'}\right)_*\hat{\mathfrak{a}}=\left(\mathsf{g}_{\mathfrak{c}}\right)_*\hat{\mathfrak{a}}$, since the cochain $\left(\mathsf{g}_{\mathfrak{c}}\right)_*\hat{\mathfrak{a}}$ is periodic under $n\cdot\mathbb{Z}^3$.
Proposition~41\,(c) gives $\left\|d\hat{\mathfrak{a}}_{\mathfrak{c}}\right\|^2=\boldsymbol{\sigma}^2$ on every torus $n$ as in part~0.4 and for every $\mathfrak{c}$, where $\boldsymbol{\sigma}^2$ depends neither on $n$, nor on $\mathfrak{c}$, nor on the placement.
Hence $H_2$ and $H_4$ are the same two polynomials on all these tori.
The value $\boldsymbol{\sigma}=0$ is not excluded; in that case $d\hat{\mathfrak{a}}_{\mathfrak{c}}=0$ and $X_{\mathfrak{c}}=0$ for all $\mathfrak{c}$, and all displays below hold with $\boldsymbol{\sigma}=0$.
According to the data of Proposition~41, on every side $n$ as in part~0.4 the cochain $\hat{\mathfrak{a}}$ is the periodisation of a single cochain $\hat{\mathfrak{a}}^{\mathrm{o}}$ on $\mathbb{Z}^3$, which is formed without reference to a torus.
All plaquettes of $\operatorname{supp}d\hat{\mathfrak{a}}^{\mathrm{o}}$ lie in $\mathrm{Bx}_0$ (Proposition~41\,(b)), and $\mathrm{Bx}_0+n\cdot\mathfrak{c}''$ is disjoint from $\mathrm{Bx}_0$ for $\mathfrak{c}''\in\mathbb{Z}^3\setminus\{0\}$ (part~0.6\,(i) at $\mathfrak{c}=0$).
Therefore
\[
\left(d\hat{\mathfrak{a}}\right)_p=\sum_{\mathfrak{c}''\in\mathbb{Z}^3}\left(d\hat{\mathfrak{a}}^{\mathrm{o}}\right)_{p+n\cdot\mathfrak{c}''}=\left(d\hat{\mathfrak{a}}^{\mathrm{o}}\right)_p\qquad\text{for every plaquette }p\text{ in }\mathrm{Bx}_0.
\]
As $e$, $c_m$ and $\boldsymbol{\sigma}^2$ do not depend on $n$ either, $O_m$ is one and the same element
\[
\hat{A}\mapsto e^{-m}\cdot H_m\left(\sum_{p\text{ a plaquette in }\mathrm{Bx}_0}\left(d\hat{\mathfrak{a}}^{\mathrm{o}}\right)_p\cdot\left(d\hat{A}\right)_p\right)-c_m
\]
of $\mathrm{Ad}\left(\mathrm{Bx}_0\right)$, and for every side $n$ as in part~0.4 the function $e^{-m}\cdot H_m\left(\left\langle d\hat{\mathfrak{a}},F^{\mathrm{u}}\right\rangle\right)-c_m$ on $C^2(n)$ is its reading on $C(n)$ (Proposition~41\,(d); Proposition~39\,(b), last sentence: one and the same $O$ is read on all tori that carry $\mathcal{Q}_*$).

\emph{0.6. Reading.}
Let $n$ be a side as in part~0.4, and let $\mathrm{Scl}$ be a set of distinct cells of the torus $n$, each given by one representative $\mathfrak{c}\in\mathbb{Z}^3$; by part~0.5 the function $X_{\mathfrak{c}}$ on $C^2(n)$ does not depend on the choice of the representative.
The torus $C(n,n,n)$ carries $\mathcal{Q}_*$, since $n\in\mathrm{per}\cdot\mathbb{Z}=2\mathrm{hc}\cdot\mathbb{Z}$ and $n\ge4$ (Section~I.6.4; Definition~39.0).

\emph{(i) One factor.}
Let $\mathfrak{c}\in\mathbb{Z}^3$ and $\mathrm{Bx}:=\mathsf{g}_{\mathfrak{c}}\mathrm{Bx}_0$.
By Proposition~41\,(c), $\mathrm{Bx}$ is a box in $Q^{\circ}_{\mathfrak{c}}$, and every plaquette of $\operatorname{supp}d\hat{\mathfrak{a}}_{\mathfrak{c}}$ lies in $\mathrm{Bx}$ modulo $n\cdot\mathbb{Z}^3$.
The box $\mathrm{Bx}$ embeds in $n\cdot\mathbb{Z}^3$ in the sense of \cite[Definition~2.12]{QEDII}.
Indeed, in each coordinate $\mathsf{g}_{\mathfrak{c}}$ acts as a translation or as a reflection (Section~I.6.4), so that it maps a coordinate cube onto a coordinate cube of the same side.
Consequently $\mathrm{Bx}+[-1,1]^3$ is a closed cube whose side is
\[
2r_0+2=\mathrm{hc}+5-\frac4\eta\le\mathrm{hc}+1,
\]
because $r_0=\frac{\mathrm{hc}}{2}-\frac1\eta\cdot\left(2-\frac{3\eta}{2}\right)$ (Section~I.7.6) and $\eta\le1$.
Moreover $\mathrm{hc}+1<10\cdot\mathrm{hc}\le2L\cdot\mathrm{hc}=n_*\le n$, as $\mathrm{hc}\ge\frac12$ and $L\ge5$ (parts~0.4 and~0.3); hence the translates of $\mathrm{Bx}+[-1,1]^3$ by $n\cdot\mathbb{Z}^3$ are pairwise disjoint.
A bosonic element of $\mathrm{Ad}(\mathrm{Bx})$ is a function of the numbers $\left(d\hat{A}\right)_p$, $p$ a plaquette in $\mathrm{Bx}$.
On $C(n)$ it is evaluated at a local lift $\hat{A}$, and $d\hat{A}=F^{\mathrm{u}}$ on the plaquettes of $\mathrm{Bx}$ \cite[Definition~2.12]{QEDII}; thus it is evaluated at the curvature of the configuration.
By (42.4), $O_m^{[\mathfrak{c}]}$ is the function
\[
u\mapsto e^{-m}\cdot H_m\left(\left\langle d\hat{\mathfrak{a}}_{\mathfrak{c}},u\right\rangle\right)-c_m
\]
of $u=F^{\mathrm{u}}\in C^2(n)$, where $\hat{\mathfrak{a}}_{\mathfrak{c}}$ is the periodic cochain $\left(\mathsf{g}_{\mathfrak{c}}\right)_*\hat{\mathfrak{a}}$.
This function is a polynomial, hence Borel, and it involves only the coordinates $u_p$ with $p$ a plaquette in $\mathrm{Bx}$; as a function of $\hat{A}$ it depends on $d\hat{A}$ alone, and it is therefore invariant under $\hat{A}\mapsto\hat{A}-d\boldsymbol{\lambda}$.
Let
\[
|u|_{\mathrm{Bx}}:=\left(\sum_{p\text{ a plaquette in }\mathrm{Bx}}u_p^2\right)^{\frac12}.
\]
Since the translates of $\mathrm{Bx}$ are disjoint, distinct plaquettes in $\mathrm{Bx}$ are inequivalent modulo $n\cdot\mathbb{Z}^3$, so that, by part~0.5,
\[
\left\langle d\hat{\mathfrak{a}}_{\mathfrak{c}},u\right\rangle=\sum_{p\text{ a plaquette in }\mathrm{Bx}}\left(d\hat{\mathfrak{a}}_{\mathfrak{c}}\right)_p\cdot u_p,\qquad\sum_{p\text{ a plaquette in }\mathrm{Bx}}\left(d\hat{\mathfrak{a}}_{\mathfrak{c}}\right)_p^2=\left\|d\hat{\mathfrak{a}}_{\mathfrak{c}}\right\|^2=\boldsymbol{\sigma}^2.
\]
The Cauchy--Schwarz inequality gives $\left|\left\langle d\hat{\mathfrak{a}}_{\mathfrak{c}},u\right\rangle\right|\le\boldsymbol{\sigma}\cdot|u|_{\mathrm{Bx}}$, and hence $1+\left|\left\langle d\hat{\mathfrak{a}}_{\mathfrak{c}},u\right\rangle\right|\le(1+\boldsymbol{\sigma})\cdot\left(1+|u|_{\mathrm{Bx}}\right)$.
For real $z$,
\begin{align*}
\left|H_4(z)\right|&\le|z|^4+6\boldsymbol{\sigma}^2\cdot|z|^2+3\boldsymbol{\sigma}^4\le\left(1+6\boldsymbol{\sigma}^2+3\boldsymbol{\sigma}^4\right)\cdot(1+|z|)^4,\\
\left|H_2(z)\right|&\le|z|^2+\boldsymbol{\sigma}^2\le\left(1+6\boldsymbol{\sigma}^2+3\boldsymbol{\sigma}^4\right)\cdot(1+|z|)^2.
\end{align*}
Put $\mathrm{cf}:=\left(e^{-m}\cdot\left(1+6\boldsymbol{\sigma}^2+3\boldsymbol{\sigma}^4\right)+\left|c_m\right|\right)\cdot(1+\boldsymbol{\sigma})^m$.
Using $z^m\le m!\cdot\exp(z)$ for $z\ge0$ and $z\le1+\frac{z^2}{4}$ for real $z$, one obtains
\begin{align*}
\left|e^{-m}\cdot H_m\left(\left\langle d\hat{\mathfrak{a}}_{\mathfrak{c}},u\right\rangle\right)-c_m\right|&\le\mathrm{cf}\cdot\left(1+|u|_{\mathrm{Bx}}\right)^m\le\mathrm{cf}\cdot m!\cdot\exp\left(1+|u|_{\mathrm{Bx}}\right)\\
&\le\mathrm{cf}\cdot m!\cdot\exp\left(2+\frac14|u|_{\mathrm{Bx}}^2\right).
\end{align*}
This is the growth bound in the definition of $\mathrm{Ad}(\mathrm{Bx})$ in \cite[Definition~2.12]{QEDII}, with $\frac14=\frac12-\frac14$ in the exponent.
Scalars are even.
Hence $O_m^{[\mathfrak{c}]}$ is a bosonic element of $\mathrm{Ad}^{\mathrm{ev}}(\mathrm{Bx})$, where $\mathrm{Bx}$ is a box in $Q^{\circ}_{\mathfrak{c}}$, and consequently an element of $\mathcal{A}\left(Q^{\circ}_{\mathfrak{c}}\right)$ (paragraph `Admissible products' of \cite[Definition~2.12]{QEDII}, with one factor in a box and $\Phi=1$).

\emph{(ii) The product.}
Proposition~39\,(b) is applied to the family
\[
O_{\mathfrak{c}}:=O_m^{[\mathfrak{c}]}\quad\text{for }\mathfrak{c}\in\mathrm{Scl},\qquad O_{\mathfrak{c}}:=1\quad\text{for the cells }\mathfrak{c}\text{ of }C(n)\text{ outside }\mathrm{Scl}.
\]
Here the constant $1$ belongs to $\mathcal{A}\left(Q^{\circ}_{\mathfrak{c}}\right)$: it is the product of the paragraph `Admissible products' of \cite[Definition~2.12]{QEDII} with no factor in a box and with $\Phi=1$; it has no box, and it leaves the product unchanged.
By (i), every $O_{\mathfrak{c}}$ of the family is a single product in the sense of \cite[Definition~2.12]{QEDII}, so that the expansion of $\prod_{\mathfrak{c}}O_{\mathfrak{c}}$ considered in Proposition~39\,(b) consists of one term.
Proposition~39\,(b) asserts that the boxes $\mathsf{g}_{\mathfrak{c}}\mathrm{Bx}_0$, $\mathfrak{c}\in\mathrm{Scl}$, together with all their translates by $n\cdot\mathbb{Z}^3$, are pairwise disjoint, and that $\prod_{\mathfrak{c}\in\mathrm{Scl}}O_m^{[\mathfrak{c}]}$ is read on $C(n)$ in the sense of \cite[Definition~2.13(a)]{QEDII}.
This is the meaning of the left side of the inequality in Proposition~40\,(d).
(By (i) the boxes embed disjointly in $n\cdot\mathbb{Z}^3$; hence the product embeds in $C(n)$, and $Z_{\mathbb{T}}$ of the product is also given by the first display of \cite[Definition~2.13]{QEDII}.)
By (i), the product is the function
\begin{equation}
\mathrm{Pol}_{m,\mathrm{Scl}}(u):=\prod_{\mathfrak{c}\in\mathrm{Scl}}\left(e^{-m}\cdot H_m\left(\left\langle d\hat{\mathfrak{a}}_{\mathfrak{c}},u\right\rangle\right)-c_m\right),\qquad u\in C^2(n),
\tag{42.5}
\end{equation}
which is a real polynomial of degree at most $m\cdot\#|\mathrm{Scl}|$ in the coordinates of $u$.
A polynomial of degree at most $i$ in the coordinates of $u$ is bounded by $\mathrm{cf}'\cdot(1+\|u\|)^i\le\mathrm{cf}'\cdot i!\cdot\exp\left(1+\|u\|\right)$, where $\mathrm{cf}'$ is the sum of the moduli of its coefficients, because every coordinate has modulus at most $\|u\|$; so it belongs to $\mathfrak{F}(n)$.
For a scalar integrand, the Berezin integral in \cite[Definition~2.13]{QEDII} is $\int_{\mathbb{T}}\exp\left(-S_{\mathbb{T}}(A)\right)$, the determinant of \cite[Definition~2.11]{QEDII} on $\mathbb{T}$, which equals $\det\left(D_A+\boldsymbol{\mu}\right)$ (Proposition~46\,(e)); and the integral over the configuration space of \cite[Definition~2.6]{QEDII} is the sum over the closed $\mathfrak{z}$ together with the integral over $A\in[0,r_q)^{\mathrm{links}}$.
Hence, for $\mathbb{T}=C(n,n,n)$, the number $Z_{\mathbb{T}}\left[\prod_{\mathfrak{c}\in\mathrm{Scl}}O_m^{[\mathfrak{c}]}\right]$ of \cite[Definition~2.13]{QEDII} is the number $Z_n\left[\mathrm{Pol}_{m,\mathrm{Scl}}\right]$ of (I.7.1) (Proposition~46\,(f)), and the number $E^{\mathrm{af}}_n\prod_{\mathfrak{c}\in\mathrm{Scl}}O_m^{[\mathfrak{c}]}$ of Proposition~40\,(d) satisfies
\begin{align}
&E^{\mathrm{af}}_n\prod_{\mathfrak{c}\in\mathrm{Scl}}O_m^{[\mathfrak{c}]}=\frac{Z_n\left[\mathrm{Pol}_{m,\mathrm{Scl}}\right]}{Z_n[1]},\ \text{the expectation on the left side of (2.4)};\notag\\
&\text{on }n=n_*:\quad O_m^{\sharp}=\mathrm{Pol}_{m,\mathrm{Cel}}.\tag{42.6}
\end{align}

\emph{0.7. The law.}
By Proposition~50\,(e), $E^{\mathrm{af}}_n=\frac{Z_n[\,\cdot\,]}{Z_n[1]}$ is the expectation with respect to a probability law of $F^{\mathrm{u}}$ on $C^2(n)$, under which every element of $\mathfrak{F}(n)$ is integrable.
Polynomials in the coordinates of $u$, the $i$-th powers ($i\ge1$) of their moduli, and the products of such functions with bounded Borel functions belong to $\mathfrak{F}(n)$.
Hence all expectations below exist, and the inequalities of Cauchy--Schwarz, H\"older and Minkowski apply to them.

\textbf{Step 1: the bound (2.3).}
By Proposition~40\,(e), $D_L\ge0$ (this is Proposition~40\,(a) at $\check{j}=L$); as $(L-1)^3\cdot N_*>0$, it follows that $\delta\ge0$.
Proposition~42.1 is now applied with $\check{j}=L\ge2$, $n=n_*\ge4$, and with the charge $q$ and the mass $\boldsymbol{\mu}$ of the statement.
At $\check{j}=L$, the numbers $G_{\check{j}}$ and $K_{\check{j}}$ of Proposition~42.1 are $\left(1+\frac{L}{3.2}\right)\cdot\frac{2L^2}{2.56}$ and $\left(1+L+L^2\right)\cdot\left(\log L+L\cdot G_L\right)$; these are the constants $G_L$ and $K_L$ of the line (C4) of Section~I.8.2 (Proposition~14.2\,(a), rows $G_L$ and $K_L$), and $K_L$ is the number that appears in (2.3).
Moreover $D^{\mathrm{mat}}_L=D^{\mathrm{mat}}$, where $\mathrm{zr}_{n_*}>0$ and $\mathrm{zr}_{L\cdot n_*}>0$ by Proposition~38.2\,(c), that is, by (1.1) at the sides $n_*$ and $L\cdot n_*$.
Proposition~42.1\,(c) and $q^2\cdot n_*=e^2\cdot R^*$ give
\begin{equation}
D_L=L^3\cdot\log Z_{n_*}[1]-\log Z_{L\cdot n_*}[1]=D^{\mathrm{ph}}+D^{\mathrm{mat}}+D^{\mathrm{fl}},\qquad D^{\mathrm{ph}}\le\left(L^3-1\right)\cdot\log\frac{2\pi}{e^2\cdot R^*}+K_L.
\tag{42.7}
\end{equation}
By (42.2), $q^2\cdot n_*\le1$; therefore Proposition~42.1\,(b) applies and gives, with the number $\mathrm{xq}=\exp\left(-\frac{2\pi^2}{q^2\cdot n_*}\right)$ of part~0.3,
\begin{equation}
\varsigma_{n_*}\le7\cdot\mathrm{xq}\le\frac18.
\tag{42.8}
\end{equation}
For $0\le z<1$ one has $-\log(1-z)=\log\left(1+\frac{z}{1-z}\right)\le\frac{z}{1-z}$.
At $z=\varsigma_{n_*}\le\frac18$ this yields $\frac{1}{1-z}\le\frac87$, and hence $-\log\left(1-\varsigma_{n_*}\right)\le\frac87\varsigma_{n_*}\le8\cdot\mathrm{xq}$.
The last assertion of Proposition~42.1\,(c), the identity $8L^3=N_*$, (42.3) and $e\le\frac13$ now give
\begin{equation}
D^{\mathrm{fl}}\le-L^3\cdot\log\left(1-\varsigma_{n_*}\right)\le8L^3\cdot\mathrm{xq}=N_*\cdot\mathrm{xq}\le N_*\cdot e^{12N_*}\le N_*\cdot3^{-12N_*}\le1;
\tag{42.9}
\end{equation}
in the last step, $3^i\ge1+2i>i$ for every integer $i\ge1$, and $i=12N_*\ge N_*$.
From (42.7) and (42.9),
\[
D_L\le\left(L^3-1\right)\cdot\log\frac{2\pi}{e^2\cdot R^*}+K_L+D^{\mathrm{mat}}+1.
\]
Dividing by $(L-1)^3\cdot N_*$ and using $\frac{L^3-1}{(L-1)^3}=a_L$, one obtains the upper bound in (2.3).

\textbf{Step 2: moments on the torus $n_*$.}
Fix $\mathfrak{c}\in\mathrm{Cel}$.
On $C(n_*,n_*,n_*)$, \cite[Theorem~6.27(b)]{QEDII} is applied to the real two-cochain $\mathrm{gc}:=d\hat{\mathfrak{a}}_{\mathfrak{c}}$, for which $\|\mathrm{gc}\|^2=\boldsymbol{\sigma}^2$ (part~0.5) and $\left\langle\mathrm{gc},F^{\mathrm{u}}\right\rangle=X_{\mathfrak{c}}$; it gives
\begin{equation}
E^{\mathrm{af}}_{n_*}X_{\mathfrak{c}}^{2i}\le\left(2i\cdot\mathrm{e}\cdot\boldsymbol{\sigma}^2\right)^i\qquad\text{for every integer }i\ge1.
\tag{42.10}
\end{equation}
In terms of the norm (0.1), at $i=4N_*$ and at $i=2N_*$ this means
\begin{align*}
\left\|X_{\mathfrak{c}}^4\right\|_{[2N_*]}&=\left(E^{\mathrm{af}}_{n_*}X_{\mathfrak{c}}^{8N_*}\right)^{\frac{1}{2N_*}}\le\left(8N_*\cdot\mathrm{e}\cdot\boldsymbol{\sigma}^2\right)^{\frac{4N_*}{2N_*}}=64\mathrm{e}^2\cdot N_*^2\cdot\boldsymbol{\sigma}^4,\\
\left\|X_{\mathfrak{c}}^2\right\|_{[2N_*]}&=\left(E^{\mathrm{af}}_{n_*}X_{\mathfrak{c}}^{4N_*}\right)^{\frac{1}{2N_*}}\le\left(4N_*\cdot\mathrm{e}\cdot\boldsymbol{\sigma}^2\right)^{\frac{2N_*}{2N_*}}=4\mathrm{e}\cdot N_*\cdot\boldsymbol{\sigma}^2.
\end{align*}
Minkowski's inequality, the formulas $H_4(z)=z^4-6\boldsymbol{\sigma}^2\cdot z^2+3\boldsymbol{\sigma}^4$ and $H_2(z)=z^2-\boldsymbol{\sigma}^2$, and $N_*\ge1$ give
\begin{align*}
\left\|H_4\left(X_{\mathfrak{c}}\right)\right\|_{[2N_*]}&\le64\mathrm{e}^2\cdot N_*^2\cdot\boldsymbol{\sigma}^4+6\boldsymbol{\sigma}^2\cdot4\mathrm{e}\cdot N_*\cdot\boldsymbol{\sigma}^2+3\boldsymbol{\sigma}^4\le\left(64\mathrm{e}^2+24\mathrm{e}+3\right)\cdot N_*^2\cdot\boldsymbol{\sigma}^4\le600N_*^2\cdot\boldsymbol{\sigma}^4,\\
\left\|H_2\left(X_{\mathfrak{c}}\right)\right\|_{[2N_*]}&\le4\mathrm{e}\cdot N_*\cdot\boldsymbol{\sigma}^2+\boldsymbol{\sigma}^2\le\left(4\mathrm{e}+1\right)\cdot N_*\cdot\boldsymbol{\sigma}^2\le12N_*\cdot\boldsymbol{\sigma}^2,
\end{align*}
where $64\cdot7.3892+24\cdot2.7183+3=541.148\le600$ and $4\cdot2.7183+1=11.8732\le12$.
A second application of Minkowski's inequality, together with (42.4), yields
\begin{equation}
\left\|O_m^{[\mathfrak{c}]}\right\|_{[2N_*]}\le e^{-m}\cdot\left\|H_m\left(X_{\mathfrak{c}}\right)\right\|_{[2N_*]}+\left|c_m\right|\le\mathrm{Mw}_m\qquad\text{for every }\mathfrak{c}\in\mathrm{Cel},\ m=2,4.
\tag{42.11}
\end{equation}

\textbf{Step 3: the bound (2.4$'$).}
In this step all objects are those of the torus $n_*$.
By \cite[Definition~13.7]{QEDII}, written in the notation of paragraph~0.2 of Section~0 ($\Sigma_{\mu\nu}=\Pi_{\mu\nu}$, $k=w$, $F=F^{\mathrm{u}}$, $\frac{2\pi}{q}=r_q$), one has $\sum_{p\in\Pi_{\mu\nu}}F^{\mathrm{u}}_p=r_q\cdot w_{\mu\nu}$ for every configuration $(A,\mathfrak{z})$ and all $\mu<\nu$, since the sum of $dA$ over $\Pi_{\mu\nu}$ vanishes.
Hence $\mathbf{1}[\mathrm{flux}=0]=\prod_{\mu<\nu}\mathbf{1}\left[\sum_{p\in\Pi_{\mu\nu}}F^{\mathrm{u}}_p=0\right]$ and $\mathbf{1}[\mathrm{flux}\ne0]=1-\mathbf{1}[\mathrm{flux}=0]$ are bounded Borel functions of $F^{\mathrm{u}}$, and
\begin{equation}
E^{\mathrm{af}}_{n_*}O_m^{\sharp}=E^{\mathrm{af}}_{n_*}\left[O_m^{\sharp}\cdot\mathbf{1}[\mathrm{flux}=0]\right]+E^{\mathrm{af}}_{n_*}\left[O_m^{\sharp}\cdot\mathbf{1}[\mathrm{flux}\ne0]\right].
\tag{42.12}
\end{equation}

\emph{The first term.}
Recall the definitions of Section~I.7.3: $Z^{(0)}_n[G]=Z_n\left[\mathbf{1}[\mathrm{flux}=0]\cdot G\right]$, $P_n(0)=\frac{Z^{(0)}_n[1]}{Z_n[1]}$ and $E^{\mathrm{av}}_n=\frac{Z^{(0)}_n[\,\cdot\,]}{Z^{(0)}_n[1]}$.
Here $Z_{n_*}[1]>0$ (Proposition~50\,(e)) and $\frac78\le1-\varsigma_{n_*}\le P_{n_*}(0)\le1$ (Proposition~42.1\,(a) and (42.8)), so that $Z^{(0)}_{n_*}[1]=P_{n_*}(0)\cdot Z_{n_*}[1]>0$ and
\[
E^{\mathrm{af}}_{n_*}\left[O_m^{\sharp}\cdot\mathbf{1}[\mathrm{flux}=0]\right]=\frac{Z^{(0)}_{n_*}\left[O_m^{\sharp}\right]}{Z_{n_*}[1]}=P_{n_*}(0)\cdot E^{\mathrm{av}}_{n_*}O_m^{\sharp}.
\]
By (42.6) and (42.5), $O_m^{\sharp}=\prod_{\mathfrak{c}\in\mathrm{Cel}}\left(e^{-m}\cdot H_m\left(X_{\mathfrak{c}}\right)-c_m\right)$; therefore Proposition~38.3\,(e), applied with the real numbers $c_m$ of the statement, gives $\left|E^{\mathrm{av}}_{n_*}O_m^{\sharp}\right|=\left(\Theta^{(0)}_m\right)^{N_*}$.
Since $P_{n_*}(0)\le1$,
\begin{equation}
\left|E^{\mathrm{af}}_{n_*}\left[O_m^{\sharp}\cdot\mathbf{1}[\mathrm{flux}=0]\right]\right|\le\left(\Theta^{(0)}_m\right)^{N_*}.
\tag{42.13}
\end{equation}

\emph{The second term.}
By Proposition~42.1\,(a), $E^{\mathrm{af}}_{n_*}\mathbf{1}[\mathrm{flux}\ne0]=1-P_{n_*}(0)\le\varsigma_{n_*}$.
The Cauchy--Schwarz inequality gives
\[
\left|E^{\mathrm{af}}_{n_*}\left[O_m^{\sharp}\cdot\mathbf{1}[\mathrm{flux}\ne0]\right]\right|\le\left(E^{\mathrm{af}}_{n_*}\mathbf{1}[\mathrm{flux}\ne0]\right)^{\frac12}\cdot\left(E^{\mathrm{af}}_{n_*}\left|O_m^{\sharp}\right|^2\right)^{\frac12}\le\varsigma_{n_*}^{\frac12}\cdot\left(E^{\mathrm{af}}_{n_*}\left|O_m^{\sharp}\right|^2\right)^{\frac12}.
\]
The function $\left|O_m^{\sharp}\right|^2$ is the product of the $N_*$ functions $\left|O_m^{[\mathfrak{c}]}\right|^2$, $\mathfrak{c}\in\mathrm{Cel}$.
H\"older's inequality with $N_*$ exponents all equal to $N_*$, followed by (42.11), gives
\[
E^{\mathrm{af}}_{n_*}\prod_{\mathfrak{c}\in\mathrm{Cel}}\left|O_m^{[\mathfrak{c}]}\right|^2\le\prod_{\mathfrak{c}\in\mathrm{Cel}}\left(E^{\mathrm{af}}_{n_*}\left|O_m^{[\mathfrak{c}]}\right|^{2N_*}\right)^{\frac{1}{N_*}}=\prod_{\mathfrak{c}\in\mathrm{Cel}}\left\|O_m^{[\mathfrak{c}]}\right\|_{[2N_*]}^2\le\mathrm{Mw}_m^{2N_*}.
\]
Consequently
\begin{equation}
\left|E^{\mathrm{af}}_{n_*}\left[O_m^{\sharp}\cdot\mathbf{1}[\mathrm{flux}\ne0]\right]\right|\le\varsigma_{n_*}^{\frac12}\cdot\mathrm{Mw}_m^{N_*}=\left(\varsigma_{n_*}^{\frac{1}{2N_*}}\cdot\mathrm{Mw}_m\right)^{N_*}.
\tag{42.14}
\end{equation}

\emph{The root.}
The function $O_m^{\sharp}$ is real, and $E^{\mathrm{af}}_{n_*}O_m^{\sharp}\ge0$ by the first inequality of Proposition~39\,(f).
That inequality is used with $\mathcal{Q}=\mathcal{Q}_*$, $\mathrm{hw}=\mathrm{hc}$ and $n=n_*\in2\mathrm{hc}\cdot\mathbb{Z}$, $n_*\ge4$, for the family $O_0:=O_m\in\mathcal{A}\left(Q^{\circ}_0\right)$ (part~0.5) and $O_{\mathfrak{c}}:=1\in\mathcal{A}\left(Q^{\circ}_{\mathfrak{c}}\right)$ for the other cells (part~0.6\,(ii)), and it is taken at $\mathfrak{c}=0$, where $\left(O_0\right)^{\sharp}_{n_*}=\prod_{\mathfrak{c}\in\mathrm{Cel}}O_m^{[\mathfrak{c}]}=O_m^{\sharp}$ (Definition~39.0).
So the $N_*$-th root of $E^{\mathrm{af}}_{n_*}O_m^{\sharp}$ is a real number $\ge0$, namely the number $\left\|O_m\right\|_*$ of Proposition~40\,(e).
By (42.12)--(42.14),
\[
0\le E^{\mathrm{af}}_{n_*}O_m^{\sharp}\le\mathrm{yq}_1^{N_*}+\mathrm{yq}_2^{N_*},\qquad\mathrm{yq}_1:=\Theta^{(0)}_m\ge0,\qquad\mathrm{yq}_2:=\varsigma_{n_*}^{\frac{1}{2N_*}}\cdot\mathrm{Mw}_m\ge0.
\]
The binomial formula, all of whose terms are non-negative here, gives $\mathrm{yq}_1^{N_*}+\mathrm{yq}_2^{N_*}\le\left(\mathrm{yq}_1+\mathrm{yq}_2\right)^{N_*}$.
Therefore
\begin{equation}
\left\|O_m\right\|_*=\left(E^{\mathrm{af}}_{n_*}O_m^{\sharp}\right)^{\frac{1}{N_*}}\le\Theta^{(0)}_m+\varsigma_{n_*}^{\frac{1}{2N_*}}\cdot\mathrm{Mw}_m.
\tag{42.15}
\end{equation}
From (42.8), (42.3) and $N_*\ge1000$,
\[
\varsigma_{n_*}^{\frac{1}{2N_*}}\le7^{\frac{1}{2N_*}}\cdot\mathrm{xq}^{\frac{1}{2N_*}}\le7^{\frac{1}{2000}}\cdot e^6\le2e^6.
\]
Combined with (42.15), this gives $\left(E^{\mathrm{af}}_{n_*}O_m^{\sharp}\right)^{\frac{1}{N_*}}\le\Theta^{(0)}_m+2e^6\cdot\mathrm{Mw}_m=\mathrm{Tw}_m$, which is (2.4$'$).

\textbf{Step 4: the bound (2.4).}
Let $n$ be a side of the statement and $\mathrm{Scl}$ a set of distinct cells of the torus $n$.
By parts~0.1 and~0.4 the setting of Proposition~40\,(e) is present, with $\delta_{\check{j}}=\delta$.
Proposition~40\,(d), applied with $O_{\mathfrak{c}}:=O_m\in\mathcal{A}\left(Q^{\circ}_0\right)$ for every $\mathfrak{c}\in\mathrm{Scl}$ (a single element, read on $C(n)$ and on $C(n_*)$: the end of part~0.5), followed by (2.4$'$), gives
\[
\left|E^{\mathrm{af}}_n\prod_{\mathfrak{c}\in\mathrm{Scl}}O_m^{[\mathfrak{c}]}\right|\le\prod_{\mathfrak{c}\in\mathrm{Scl}}\exp(\delta)\cdot\left\|O_m\right\|_*=\left(\exp(\delta)\cdot\left\|O_m\right\|_*\right)^{\#|\mathrm{Scl}|}\le\left(\exp(\delta)\cdot\mathrm{Tw}_m\right)^{\#|\mathrm{Scl}|}=\mathsf{u}_m^{\#|\mathrm{Scl}|}.
\]
By (42.6) and (42.5), the left side here is the left side of (2.4).
For the empty set both sides of (2.4) are equal to $1$.

\textbf{Step 5: the bounds (2.1) and (2.2).}
Let $n$ be a side of the statement.
For the set $\mathrm{Scl}=\{0\}$, which consists of one cell, and since $X_0=X$ (part~0.5), the bound (2.4) reads
\begin{equation}
\left|e^{-m}\cdot E^{\mathrm{af}}_nH_m(X)-c_m\right|\le\mathsf{u}_m\qquad(m=2,4).
\tag{42.16}
\end{equation}
Let $i\ge1$ be an odd integer.
The function $\Phi(u):=\left\langle d\hat{\mathfrak{a}},u\right\rangle^i$ belongs to $\mathfrak{F}(n)$ and is odd: $\Phi(-u)=-\Phi(u)$.
Proposition~50\,(g$'$) gives $Z_n\left[\Phi(-\,\cdot\,)\right]=Z_n[\Phi]$, so that $Z_n[\Phi]=-Z_n[\Phi]=0$.
Taking $i=1$ and $i=3$, one finds $E^{\mathrm{af}}_nX=E^{\mathrm{af}}_nX^3=0$.
For a real random variable $X$ with finite fourth moment, the second and the fourth cumulant are
\begin{align*}
\kappa_2&=E^{\mathrm{af}}_nX^2-\left(E^{\mathrm{af}}_nX\right)^2,\\
\kappa_4&=E^{\mathrm{af}}_nX^4-4E^{\mathrm{af}}_nX\cdot E^{\mathrm{af}}_nX^3-3\left(E^{\mathrm{af}}_nX^2\right)^2+12\left(E^{\mathrm{af}}_nX\right)^2\cdot E^{\mathrm{af}}_nX^2-6\left(E^{\mathrm{af}}_nX\right)^4.
\end{align*}
Hence $\kappa_2=E^{\mathrm{af}}_nX^2$ and $\kappa_4=E^{\mathrm{af}}_nX^4-3\left(E^{\mathrm{af}}_nX^2\right)^2$.
Let $\mathrm{uq}:=E^{\mathrm{af}}_nH_2(X)=E^{\mathrm{af}}_nX^2-\boldsymbol{\sigma}^2=\kappa_2-\boldsymbol{\sigma}^2$.
From the identity $z^4=H_4(z)+6\boldsymbol{\sigma}^2\cdot z^2-3\boldsymbol{\sigma}^4$ one obtains

\begin{equation}
\tag{42.17}
\begin{aligned}
E^{\mathrm{af}}_n X^4&=E^{\mathrm{af}}_n H_{4}(X)+6\boldsymbol{\sigma}^2\cdot\left(\mathrm{uq}+\boldsymbol{\sigma}^2\right)-3\boldsymbol{\sigma}^4=E^{\mathrm{af}}_n H_{4}(X)+6\boldsymbol{\sigma}^2\cdot\mathrm{uq}+3\boldsymbol{\sigma}^4,\\
\kappa_4&=E^{\mathrm{af}}_n H_{4}(X)+6\boldsymbol{\sigma}^2\cdot\mathrm{uq}+3\boldsymbol{\sigma}^4-3\left(\mathrm{uq}+\boldsymbol{\sigma}^2\right)^2=E^{\mathrm{af}}_n H_{4}(X)-3\,\mathrm{uq}^2.
\end{aligned}
\end{equation}
The first line holds because $E^{\mathrm{af}}_n X^2=\mathrm{uq}+\boldsymbol{\sigma}^2$, the second because $\kappa_4=E^{\mathrm{af}}_n X^4-3\left(E^{\mathrm{af}}_n X^2\right)^2$.

\emph{Proof of (2.1).}
Since $\kappa_2-\boldsymbol{\sigma}^2=\mathrm{uq}=E^{\mathrm{af}}_n H_{2}(X)$,
\[
e^{-2}\cdot\left(\kappa_2-\boldsymbol{\sigma}^2\right)-c_2=e^{-2}\cdot E^{\mathrm{af}}_n H_{2}(X)-c_2,
\]
and by (42.16) at $m=2$ the modulus of this number is at most $\mathsf{u}_2$.

\emph{Proof of (2.2).}
Let $\mathrm{yq}:=e^{-2}\cdot\mathrm{uq}$.
By (42.17),
\[
e^{-4}\cdot\kappa_4-\left(c_4-3c_2^2\right)=\left(e^{-4}\cdot E^{\mathrm{af}}_n H_{4}(X)-c_4\right)-3\left(\mathrm{yq}^2-c_2^2\right).
\]
By (42.16) at $m=4$ and at $m=2$,
\[
\left|e^{-4}\cdot E^{\mathrm{af}}_n H_{4}(X)-c_4\right|\le\mathsf{u}_4,\qquad\left|\mathrm{yq}-c_2\right|\le\mathsf{u}_2.
\]
The function $z\mapsto z\cdot\left(2|c_2|+z\right)$ is non-decreasing on $[0,\infty)$; together with $|\mathrm{yq}+c_2|\le|\mathrm{yq}-c_2|+2|c_2|$ this gives
\[
\left|\mathrm{yq}^2-c_2^2\right|=\left|\mathrm{yq}-c_2\right|\cdot\left|\mathrm{yq}+c_2\right|\le\left|\mathrm{yq}-c_2\right|\cdot\left(2|c_2|+\left|\mathrm{yq}-c_2\right|\right)\le\mathsf{u}_2\cdot\left(2|c_2|+\mathsf{u}_2\right).
\]
Consequently
\[
\left|e^{-4}\cdot\kappa_4-\left(c_4-3c_2^2\right)\right|\le\mathsf{u}_4+3\mathsf{u}_2\cdot\left(2|c_2|+\mathsf{u}_2\right),
\]
which is (2.2).
This completes the proof of Proposition~42.2.\qed

\begin{named}{Remark}[Dependence of the bounds]
The right sides of (2.1), (2.2) and (2.4) do not contain the side $n$, and they depend on the set $\mathrm{Scl}$ only through $\#|\mathrm{Scl}|$.
They are formed from $e$, $c_2$, $c_4$, $\boldsymbol{\sigma}^2$, $N_{*}$, $\delta$ and $\Theta^{(0)}_m$, that is, from numbers attached to the cutoff $N$, to the probe and to the two tori $n_*$ and $L\cdot n_*$.
Each of the three bounds therefore holds with one and the same right side for all sides $n=n_*$ and $n\in\mathrm{per}\,\mathbb{Z}$ with $n\ge L\cdot n_*$, of either parity.
The numerals $600$ and $12$, the factor $2$ in $2e^6$ and the summand $1$ in (2.3) are absolute numbers, while $N_{*}$, $a_{L}$ and $K_{L}$ are functions of $L$; none of these numbers contains $\rho$, the weights, $\boldsymbol{\mu}$, $\bar{m}$, $N$, $e$, $\beta$, $v$, $c_2$, $c_4$ or $n$.
The placement $v$ enters only through Proposition~41\,(b), (c), (d) and Proposition~38.3\,(e), and these hold for every placement.
No sign condition on $\boldsymbol{\mu}$ or on $E^{\mathrm{av}}_{n_*}O_m^{\sharp}$ is needed, and of the invariances of the law only the invariance under $F^{\mathrm{u}}\mapsto-F^{\mathrm{u}}$ enters.
\end{named}

\subsection*{3. Conditions on the constants}

In the following table 0.1--0.7 are the parts of Step~0 of the proof of Proposition~42.2, and Steps~1, 3, 4 are steps of the same proof.

\noindent
\begin{tabular}{p{0.27\textwidth}p{0.4\textwidth}p{0.2\textwidth}}
\hline
condition & used in & line of Section~I.8.2; label in Proposition~14.2\\
\hline
$\rho\in(0,1)$; the standard weights
& proofs of Propositions~42.1 and~42.2, only through the hypotheses of the statements applied there: Propositions~38.2, 38.3, 39, 40 and~50, Proposition~46\,(e), (f), (f+), which are stated for real $\rho$ and the standard weights, and \cite[Theorem~6.27]{QEDII}, \cite[Lemma~13.12]{QEDII} (the last two require $-1\le\rho\le1$)
& (C0)\\
\hline
$L$ an odd integer, $L\ge5$
& proof of Proposition~42.2: $\check{j}=L\ge2$ (Steps~1 and~4); $n_*\ge L\ge4$ (0.4); $10\,\mathrm{hc}\le2L\cdot\mathrm{hc}$ (0.6\,(i)); $N_{*}=8L^3\ge1000$ in (42.2) and in $7^{\frac{1}{2N_{*}}}\le2$ (Step~3)
& (C3); Proposition~14.2\,(b), (L5)\\
\hline
$M$ a power of $L$, $M\ge L^2$
& proof of Proposition~42.2, 0.4: $T^{*}(e)$ is an integer $\ge1$
& (C5); Proposition~14.2\,(b), (M1)\\
\hline
$r$ an integer, $r\ge2$
& proof of Proposition~42.2, 0.2 (only $r>0$ enters: it gives $v_K^r>1$ for $v_K>1$)
& (C7)\\
\hline
$e_1>0$
& proof of Proposition~42.2, 0.1: the set whose supremum is $e_G$ is not empty
& (C9); Section~I.0.5; Proposition~14.3\,(f)\\
\hline
$e_1\le\frac{1}{3}$
& proof of Proposition~42.2, (42.1): $e\le\frac{1}{3}$; 0.2: $v_K\ge\log3>1$
& (C9); Proposition~14.2\,(f)\\
\hline
\end{tabular}

\medskip
\noindent
\begin{tabular}{p{0.27\textwidth}p{0.4\textwidth}p{0.2\textwidth}}
\hline
condition & used in & line of Section~I.8.2; label in Proposition~14.2\\
\hline
(e8) at the number $v_K$: $6N_{*}\cdot M\cdot v_K^r\cdot\exp(-v_K)\cdot\left(v_K\cdot\exp(-v_K)+\frac{1}{2}\right)\le\pi^2$ (the condition is imposed at every real number $\ge-\log e_1$, and by Proposition~14.3\,(f) it holds at every such number; $v_K\ge-\log e_1$ by (42.1))
& proof of Proposition~42.2, 0.2: (W3)
& (C9); Proposition~14.2\,(b), (e8); Proposition~14.3\,(f)\\
\hline
$\bar{m}\in(0,1]$; hence $h=L^{\left\lfloor\log_L\frac{1}{\bar{m}}\right\rfloor}=L^i$ with an integer $i\ge0$, $h\ge1$ and $\log h\ge0$
& proof of Proposition~42.2, 0.1: $e\mapsto e\cdot h^{\frac{1}{2}}\cdot\log h$ is non-decreasing, and $e\le e\cdot h^{\frac{1}{2}}$ in (42.1); 0.2: $\log\frac{1}{e}=v_K+\frac{1}{2}\log h>0$; 0.4: $n_*$ and $\mathrm{per}$ are integers
& (C10); Section~I.0.5; Proposition~14.2\,(a), rows $\bar{m}$, $h$\\
\hline
$e\cdot h^{\frac{1}{2}}\le e_1$ and (Z8.w) $e\cdot h^{\frac{1}{2}}\cdot\log h\le1$
& proof of Proposition~42.2, (42.1); both conditions are contained in the hypothesis $e\le e_G$
& (C10); Proposition~14.2\,(a), row $e_G$, and (d), (Z8.w)\\
\hline
(W3): $6N_{*}\cdot e^2\cdot R^{*}\cdot\log\frac{1}{e}\le\pi^2$
& proof of Proposition~42.2, (42.2), (42.3): $e^2\cdot R^{*}<1$; $D^{\mathrm{fl}}\le1$; $\varsigma_{n_*}^{\frac{1}{2N_{*}}}\le2e^6$. It is derived in 0.2 from five of the rows above ($r$; $e_1\le\frac{1}{3}$; (e8); $\bar{m}$; $e\cdot h^{\frac{1}{2}}\le e_1$ with (Z8.w))
& (C10); Proposition~14.2\,(d), (W3); Annex~B, Section~G6.Z4\\
\hline
\end{tabular}

\medskip
The numbers $N_{*}=8L^3$, $a_{L}=\frac{L^3-1}{(L-1)^3}$, $G_{L}$ and $K_{L}$ are given by formulas in $L$ (line (C4); Proposition~14.2\,(a)), and no condition is imposed on them.
Proposition~42.1 requires no condition on the constants besides the data of the statements applied in its proof.
The constants among these data are $\rho\in(0,1)$ and the standard weights (the first row of the table above), which enter only through the hypotheses of Proposition~38.2\,(c), Proposition~46\,(f+), Proposition~50\,(e) and \cite[Lemma~13.12(a)]{QEDII}, and the number $\bar{m}$, which enters only as a datum of Proposition~38.2\,(c).
Indeed, $\bar{m}$ occurs in the proof of Proposition~42.1 only in $m_{\mathrm{in}}:=\boldsymbol{\mu}-\bar{m}$, and what Proposition~38.2\,(c) asserts on $Z^{(0)}_n[1]$, $Z^{\mathrm{ph}}_n$ and $\mathrm{zr}_n$ involves $n$, $q$ and $\boldsymbol{\mu}$ only.
No condition on $L$, $M$, $r$, $e_1$ or $e_G$ enters the proof of Proposition~42.1.
No condition on the constants other than those imposed in Sections~I.8.2 and~I.8.4 and in Proposition~14.2 enters the proofs of Propositions~42.1 and~42.2; hence both propositions hold for every choice of the constants that satisfies these conditions.

\subsection*{Appendix A. A bound for the lattice sum $\mathrm{Ld}$}

For integers $n_1,n_2,n_3\ge2$ let
\[
\mathrm{Ld}(n_1,n_2,n_3):=\sum_{\substack{k\in\prod_{\mu}\mathbb{Z}/n_\mu\mathbb{Z}\\ k\ne0}}\log\left(\sum_{\mu=1}^{3}\left(2-2\cos\frac{2\pi\cdot k_\mu}{n_\mu}\right)\right),\qquad\mathrm{Ld}(n):=\mathrm{Ld}(n,n,n).
\]
For $k\ne0$ at least one of the three summands is positive, so that every logarithm is defined; $\mathrm{Ld}(n)$ is the number that occurs in Proposition~38.2\,(c).
The function $\mathrm{Ld}$ is symmetric in its three arguments, since a permutation of the coordinates of $k$ permutes the terms of the sum.

\begin{named}{Lemma A}[{\cite[Lemma~21.31]{QEDI}}]
For all integers $n\ge2$ and $\check{j}\ge2$,
\[
\frac{1}{2}\left[\mathrm{Ld}(\check{j}\cdot n)-\check{j}^3\cdot\mathrm{Ld}(n)\right]\le-\left(\check{j}^3-1\right)\cdot\log n+K_{\check{j}},
\]
where $G_{\check{j}}$ and $K_{\check{j}}$ are the numbers of Proposition~42.1.
\end{named}

\begin{proof}
\emph{(Ax1) Summation in one coordinate direction.}
Let $n_1\ge1$ be an integer, let $\mathrm{zeq}:=\exp\left(\frac{2\pi\cdot\sqrt{-1}}{n_1}\right)$, a primitive $n_1$-th root of $1$, and let $\mathrm{oq}\ge0$ be a real number.
For every integer $k$,
\begin{align*}
\left(\exp(\mathrm{oq})-\mathrm{zeq}^k\right)\cdot\left(1-\exp(-\mathrm{oq})\cdot\mathrm{zeq}^{-k}\right)&=\exp(\mathrm{oq})+\exp(-\mathrm{oq})-\mathrm{zeq}^k-\mathrm{zeq}^{-k}\\
&=2\cosh\mathrm{oq}-2\cos\frac{2\pi\cdot k}{n_1}.
\end{align*}
The polynomial identity $\prod_{k\in\mathbb{Z}/n_1\mathbb{Z}}\left(z-\mathrm{zeq}^k\right)=z^{n_1}-1$ gives, with all products extending over $k\in\mathbb{Z}/n_1\mathbb{Z}$ (as $k$ ranges over $\mathbb{Z}/n_1\mathbb{Z}$, so does $-k$),
\begin{align*}
\prod_k\left(\exp(\mathrm{oq})-\mathrm{zeq}^k\right)&=\exp\left(n_1\cdot\mathrm{oq}\right)-1,\\
\prod_k\left(1-\exp(-\mathrm{oq})\cdot\mathrm{zeq}^{-k}\right)&=\exp\left(-n_1\cdot\mathrm{oq}\right)\cdot\prod_k\left(\exp(\mathrm{oq})-\mathrm{zeq}^{-k}\right)=1-\exp\left(-n_1\cdot\mathrm{oq}\right).
\end{align*}
Multiplication of the two identities yields
\begin{align*}
\prod_{k\in\mathbb{Z}/n_1\mathbb{Z}}\left(2\cosh\mathrm{oq}-2\cos\frac{2\pi\cdot k}{n_1}\right)&=\left(\exp\left(n_1\cdot\mathrm{oq}\right)-1\right)\cdot\left(1-\exp\left(-n_1\cdot\mathrm{oq}\right)\right)\\
&=2\cosh\left(n_1\cdot\mathrm{oq}\right)-2=4\sinh^2\left(\frac{n_1\cdot\mathrm{oq}}{2}\right).
\end{align*}
Moreover $\prod_{k\ne0}\left(z-\mathrm{zeq}^k\right)=\frac{z^{n_1}-1}{z-1}=1+z+\dots+z^{n_1-1}$; at $z=1$ this gives $\prod_{k\ne0}\left(1-\mathrm{zeq}^k\right)=n_1$, and since $2-2\cos\frac{2\pi\cdot k}{n_1}=\left|1-\mathrm{zeq}^k\right|^2$,
\[
\prod_{k\ne0}\left(2-2\cos\frac{2\pi\cdot k}{n_1}\right)=\prod_{k\ne0}\left|1-\mathrm{zeq}^k\right|^2=n_1^2.
\]
Write $k=(k_1,\mathrm{ks})$ with $\mathrm{ks}=(k_2,k_3)\in\mathbb{Z}/n_2\mathbb{Z}\times\mathbb{Z}/n_3\mathbb{Z}$.
For $\mathrm{ks}\ne0$ let
\[
\mathrm{Mk}(\mathrm{ks}):=\sum_{\mu=2,3}\left(2-2\cos\frac{2\pi\cdot k_\mu}{n_\mu}\right)>0,
\]
and let $\mathrm{oq}(\mathrm{ks})>0$ be the number with $2\cosh\mathrm{oq}(\mathrm{ks})=2+\mathrm{Mk}(\mathrm{ks})$; it does not depend on $n_1$.
In the sum that defines $\mathrm{Ld}(n_1,n_2,n_3)$ the terms with $\mathrm{ks}=0$ have $k_1\ne0$, and by the last identity they add up to $\log n_1^2$.
For a fixed $\mathrm{ks}\ne0$ the argument of the logarithm is $2\cosh\mathrm{oq}(\mathrm{ks})-2\cos\frac{2\pi\cdot k_1}{n_1}$, and by the product formula at $\mathrm{oq}=\mathrm{oq}(\mathrm{ks})$ the terms with this $\mathrm{ks}$ add up to $\log\left(4\sinh^2\left(\frac{n_1\cdot\mathrm{oq}(\mathrm{ks})}{2}\right)\right)$.
Therefore
\begin{equation}
\mathrm{Ld}(n_1,n_2,n_3)=2\log n_1+2\sum_{\mathrm{ks}\ne0}\log\left(2\sinh\left(\frac{n_1\cdot\mathrm{oq}(\mathrm{ks})}{2}\right)\right).\tag{A.1}
\end{equation}

\emph{(Ax2) One side multiplied by $\check{j}$.}
For $z>0$ let $\mathrm{fq}(z):=-\log\left(1-\exp(-2z)\right)>0$.
Since $2\sinh z=\exp(z)\cdot\left(1-\exp(-2z)\right)$,
\begin{align*}
\log\left(2\sinh\left(\check{j}\cdot z\right)\right)-\check{j}\cdot\log\left(2\sinh z\right)&=\left[\check{j}\cdot z+\log\left(1-\exp\left(-2\check{j}\cdot z\right)\right)\right]-\check{j}\cdot\left[z-\mathrm{fq}(z)\right]\\
&=\log\left(1-\exp\left(-2\check{j}\cdot z\right)\right)+\check{j}\cdot\mathrm{fq}(z)\le\check{j}\cdot\mathrm{fq}(z).
\end{align*}
Apply (A.1) to the triples $(\check{j}\cdot n_1,n_2,n_3)$ and $(n_1,n_2,n_3)$, which have the same set of indices $\mathrm{ks}$ and the same numbers $\mathrm{oq}(\mathrm{ks})$, and use the last inequality at $z=\frac{n_1\cdot\mathrm{oq}(\mathrm{ks})}{2}$ for every $\mathrm{ks}\ne0$.
This gives
\begin{equation}
\mathrm{Ld}(\check{j}\cdot n_1,n_2,n_3)-\check{j}\cdot\mathrm{Ld}(n_1,n_2,n_3)\le2\log\left(\check{j}\cdot n_1\right)-2\check{j}\cdot\log n_1+2\check{j}\cdot\sum_{\mathrm{ks}\ne0}\mathrm{fq}\left(\frac{n_1\cdot\mathrm{oq}(\mathrm{ks})}{2}\right).\tag{A.2}
\end{equation}

\emph{(Ax3) The sum over $\mathrm{ks}$.}
Let $n_1=n$ and $n_2,n_3\in\{n,\check{j}\cdot n\}$.
Represent each class $k_\mu$ ($\mu=2,3$) by an integer with $|k_\mu|\le\frac{n_\mu}{2}$, and put $p_\mu:=\frac{2\pi\cdot k_\mu}{n_\mu}\in[-\pi,\pi]$ and $p:=(p_2,p_3)$.
Different classes $\mathrm{ks}\ne0$ have different representatives in $\mathbb{Z}^2\setminus\{0\}$.
\begin{itemize}
\item[(i)] Since $\sin z\ge\frac{2z}{\pi}$ on $\left[0,\frac{\pi}{2}\right]$, one has $2-2\cos p_\mu=4\sin^2\frac{p_\mu}{2}\ge\frac{4p_\mu^2}{\pi^2}$; moreover $2-2\cos p_\mu\le4$.
Hence $\frac{4}{\pi^2}\cdot|p|_2^2\le\mathrm{Mk}\le8$.
\item[(ii)] One has $2\sinh^2\frac{\mathrm{oq}}{2}=\cosh\mathrm{oq}-1=\frac{\mathrm{Mk}}{2}$; so $\mathrm{oq}=2\operatorname{arcsinh}\left(\frac{\sqrt{\mathrm{Mk}}}{2}\right)$ with $\frac{\sqrt{\mathrm{Mk}}}{2}\in\left(0,\sqrt{2}\right]$.
The function $\operatorname{arcsinh}$ is concave on $[0,\infty)$ and vanishes at $0$; hence $\operatorname{arcsinh}z\ge z\cdot\frac{\operatorname{arcsinh}\sqrt{2}}{\sqrt{2}}$ on $\left[0,\sqrt{2}\right]$.
Further $\operatorname{arcsinh}\sqrt{2}=\log\left(\sqrt{2}+\sqrt{3}\right)\ge0.81\cdot\sqrt{2}$: indeed $0.81\cdot\sqrt{2}<1.1456$, and
\[
\exp(1.1456)=\mathrm{e}\cdot\exp(0.1456)<2.7183\cdot1.15674<3.1445<1.4142+1.7320<\sqrt{2}+\sqrt{3},
\]
where $\exp(0.1456)<1.15674$ holds because $\exp(z)\le1+z+\frac{z^2}{2}+\frac{2z^3}{3\cdot(4-z)}$ for $0<z<1$ and the right side is $<1.15674$ at $z=0.1456$.
Hence
\[
\mathrm{oq}\ge0.81\cdot\sqrt{\mathrm{Mk}}\ge0.81\cdot\frac{2}{\pi}\cdot|p|_2\ge0.51\cdot|p|_2,
\]
because $\frac{1.62}{3.1416}>0.51$.
\item[(iii)] Since $\frac{n}{n_\mu}\ge\frac{1}{\check{j}}$,
\[
\frac{n\cdot\mathrm{oq}}{2}\ge0.255\,n\cdot|p|_2=0.51\pi\cdot\left|\left(\frac{k_2\cdot n}{n_2},\frac{k_3\cdot n}{n_3}\right)\right|_2\ge\frac{0.51\pi}{\check{j}}\cdot|\mathrm{ks}|_2\ge\mathrm{cj}\cdot|\mathrm{ks}|_2,\qquad\mathrm{cj}:=\frac{1.6}{\check{j}},
\]
where $|\mathrm{ks}|_2$ is the Euclidean norm of the representative $(k_2,k_3)\in\mathbb{Z}^2$, and $0.51\cdot3.1415>1.6$.
\item[(iv)] For $0\le z<1$ one has $-\log(1-z)\le\frac{z}{1-z}$.
For $z>0$ the inequality $\exp(z)\ge1+z$ gives $\frac{1}{1-\exp(-z)}=1+\frac{1}{\exp(z)-1}\le1+\frac{1}{z}$.
Hence
\[
\mathrm{fq}(z)\le\frac{\exp(-2z)}{1-\exp(-2z)}\le\exp(-2z)\cdot\left(1+\frac{1}{2z}\right)\qquad(z>0),
\]
and the last function is decreasing in $z>0$.
\item[(v)] For $\mathrm{ks}\in\mathbb{Z}^2\setminus\{0\}$ one has $|\mathrm{ks}|_2\ge1$ and $|\mathrm{ks}|_2\ge\frac{|k_2|+|k_3|}{\sqrt{2}}$.
By (iii) and (iv),
\begin{align*}
\mathrm{fq}\left(\frac{n\cdot\mathrm{oq}(\mathrm{ks})}{2}\right)&\le\exp\left(-2\mathrm{cj}\cdot|\mathrm{ks}|_2\right)\cdot\left(1+\frac{1}{2\mathrm{cj}\cdot|\mathrm{ks}|_2}\right)\\
&\le\left(1+\frac{1}{2\mathrm{cj}}\right)\cdot\exp\left(-\sqrt{2}\cdot\mathrm{cj}\cdot\left(|k_2|+|k_3|\right)\right).
\end{align*}
For $y>0$ one has $\sum_{i\in\mathbb{Z}}\exp\left(-y\cdot|i|\right)=\frac{1+\exp(-y)}{1-\exp(-y)}=\coth\frac{y}{2}$; moreover $\coth^2y-1=\frac{1}{\sinh^2y}$ for $y>0$, and $\sinh y\ge y$ for $y\ge0$.
\end{itemize}
Summing the bound of (v) over the distinct representatives $\mathrm{ks}\ne0$, and enlarging the sum to all of $\mathbb{Z}^2\setminus\{0\}$, one obtains
\begin{equation}
\tag{A.3}
\begin{split}
\sum_{\mathrm{ks}\ne0}\mathrm{fq}\left(\frac{n\cdot\mathrm{oq}(\mathrm{ks})}{2}\right)&\le\left(1+\frac{1}{2\mathrm{cj}}\right)\cdot\left[\left(\sum_{i\in\mathbb{Z}}\exp\left(-\sqrt{2}\cdot\mathrm{cj}\cdot|i|\right)\right)^2-1\right]\\
&=\left(1+\frac{1}{2\mathrm{cj}}\right)\cdot\left(\sinh\frac{\mathrm{cj}}{\sqrt{2}}\right)^{-2}\\
&\le\left(1+\frac{1}{2\mathrm{cj}}\right)\cdot\frac{2}{\mathrm{cj}^2}=\left(1+\frac{\check{j}}{3.2}\right)\cdot\frac{2\check{j}^2}{2.56}=G_{\check{j}}.
\end{split}
\end{equation}

\emph{(Ax4) All three sides.}
By the symmetry of $\mathrm{Ld}$ one has $\mathrm{Ld}(\check{j}\cdot n,n,\check{j}\cdot n)=\mathrm{Ld}(n,\check{j}\cdot n,\check{j}\cdot n)$ and $\mathrm{Ld}(\check{j}\cdot n,n,n)=\mathrm{Ld}(n,n,\check{j}\cdot n)$; hence
\begin{multline*}
\mathrm{Ld}(\check{j}\cdot n,\check{j}\cdot n,\check{j}\cdot n)-\check{j}^3\cdot\mathrm{Ld}(n,n,n)=\left[\mathrm{Ld}(\check{j}\cdot n,\check{j}\cdot n,\check{j}\cdot n)-\check{j}\cdot\mathrm{Ld}(n,\check{j}\cdot n,\check{j}\cdot n)\right]\\
+\check{j}\cdot\left[\mathrm{Ld}(\check{j}\cdot n,n,\check{j}\cdot n)-\check{j}\cdot\mathrm{Ld}(n,n,\check{j}\cdot n)\right]+\check{j}^2\cdot\left[\mathrm{Ld}(\check{j}\cdot n,n,n)-\check{j}\cdot\mathrm{Ld}(n,n,n)\right].
\end{multline*}
The three brackets are the left side of (A.2) with $n_1=n$ and $(n_2,n_3)=(\check{j}\cdot n,\check{j}\cdot n)$, $(n,\check{j}\cdot n)$, $(n,n)$, respectively.
By (A.2) and (A.3) each of them is at most
\[
2\log\left(\check{j}\cdot n\right)-2\check{j}\cdot\log n+2\check{j}\cdot G_{\check{j}}=2\log\check{j}+2\left(1-\check{j}\right)\cdot\log n+2\check{j}\cdot G_{\check{j}}.
\]
The coefficients $1$, $\check{j}$, $\check{j}^2$ of the brackets are positive, and $\left(1+\check{j}+\check{j}^2\right)\cdot\left(1-\check{j}\right)=1-\check{j}^3$.
It follows that
\begin{align*}
\mathrm{Ld}(\check{j}\cdot n)-\check{j}^3\cdot\mathrm{Ld}(n)&\le\left(1+\check{j}+\check{j}^2\right)\cdot\left[2\log\check{j}+2\check{j}\cdot G_{\check{j}}\right]-2\left(\check{j}^3-1\right)\cdot\log n\\
&=2K_{\check{j}}-2\left(\check{j}^3-1\right)\cdot\log n,
\end{align*}
and division by $2$ gives the assertion.
\end{proof}

\section*{Chapter 43. Reference data and the ultraviolet proposition: proofs of Propositions 43.1, 43.2}

\subsection*{0. Conventions of this chapter; the facts on the constants; three facts on $h$, $e_G$, $C_{\mathrm{mat}}$}

The notation is that of Part~I and of the definitions of Part~II whose section numbers begin with G5 or G6.
The scales $h$, $e_K$, $v_K$, $\mathrm{lg}$, $r_K=r_K(e)$ (the number $r_k$ of Section~I.6.1 at $k=K$), $R_{\mathrm{top}}(e)$, $R^{*}(e)$, $T^{*}(e)$, $\mathrm{Tor}(e)$, $N_{\min}(R,e)$, $N_0(e)$,
and the numbers $\eta$, $q$, $n_*$, $\mathrm{per}$, $\ell_c$, $N_{*}$, $a_{L}$, are those of Sections~I.0.5, I.2.5, I.6.1, I.6.2, I.6.4 and~I.7.6.
In this chapter the bare letter $v$ denotes a placement and nothing else (Section~I.7.6).
The real variable of the conditions on $e_1$, which is called $v$ in Proposition~14.2\,(b), in Proposition~14.3\,(f) and in the assertion (Sc6) of Section~G4.0.2, is written $v_k$.

\paragraph{The facts on the constants.}
The arguments below rest on the following facts.

\emph{From Part~I.}
\begin{itemize}
\item[--] Section~I.0.5: $\rho\in(0,1)$; $w_3=6$; $L$ is an odd integer with $L\ge5$; $M$ is a power of $L$ with $M\ge L^2$.
\item[--] (I.O8): $\lambda_{\mathrm{int}}>0$.
\item[--] (I.U4)$^{\mathrm{ch}}$ of Section~I.8.1: a constant is a number.
The finiteness of $K_{\mathrm{Str}}$ and of $C_{6}$ is asserted in Proposition~14.2\,(f), paragraph \emph{Ranges}.
\item[--] Section~I.8.2:
$\bar{\varepsilon}=1/200$ (line (C0));
the formulas of the line (C0r), which enter (0.M) below;
$\alpha<1$ (line (C1));
$N_{*}$, $a_{L}$ and $K_{L}$ are constants of the line (C4), and $K_{L}$ is the number $K_j$ of \cite[Section~21.1.7]{QEDI} at $j=L$, a function of $L$; by the row $K_{L}$ of Proposition~14.2\,(a),
\[
K_{L}=\left(1+L+L^2\right)\cdot\left(\log L+L\cdot G_{L}\right),\qquad G_{L}=\left(1+\frac{L}{3.2}\right)\cdot\frac{2L^2}{2.56};
\]
$r$ is an integer with $r\ge2$, and $1-\alpha<\varepsilon\le\bar{\varepsilon}$, $0<s<1$ (line (C7));
$e_1\le1/3$ (line (C9));
$e_G$ is the supremum of the line (C10),
\[
e_G=\sup\left\{e>0:\ e\cdot h^{\frac12}\le e_1,\ e\cdot h^{\frac12}\cdot\log h\le1\right\}.
\]
\item[--] Section~I.8.5: the exponents $\mathfrak{e}$ and $\sigma$, and $\mathfrak{e}_2=0$.
\end{itemize}

\emph{From Proposition~14.2.}
\begin{itemize}
\item[--] (a), the rows $\mathfrak{e}$ and $\sigma$:
$\mathfrak{e}=s\left(\frac14-9\varepsilon\right)-12\varepsilon$ and $\sigma=\mathfrak{e}-\frac{2a_{L}}{N_{*}}$.
\item[--] (b), the rows (e1), (e3), (e8) and (S2).
The first and the third inequality of the row (e1), the third inequality of the row (e3), the row (e8) and the row (S2) read
\begin{gather*}
v_k\ge8,\qquad\left(1+\frac{3\,\mathrm{lg}}{v_k}\right)^r<L,\qquad\exp\left(-\left(\frac14-9\varepsilon\right)v_k\right)\le\frac18,\\
6N_{*}\cdot M\cdot v_k^r\cdot\exp(-v_k)\cdot\left(v_k\cdot\exp(-v_k)+\frac12\right)\le\pi^2,\\
\sigma=s\left(\frac14-9\varepsilon\right)-12\varepsilon-\frac{2a_{L}}{N_{*}}>0.
\end{gather*}
By the sentence that opens Proposition~14.2\,(b), the conditions on $e_1$, among them (e1), (e3) and (e8), are imposed for \emph{all} $v_k\ge-\log e_1$.
\item[--] (d): the conditions (Z8.w) and (W3), namely
\[
e\cdot h^{\frac12}\cdot\log h\le1,\qquad 6N_{*}\cdot e^2\cdot R^{*}(e)\cdot\log\frac1e\le\pi^2.
\]
\end{itemize}

\emph{From Proposition~14.3\,(f).}
One has $e_1>0$, and the rows (e1)--(e8) hold for all $v_k\ge-\log e_1$;
the latter fact is called below the second assertion of Proposition~14.3\,(f).

The labels (G6.T-14), (G6.T-35), (G6.T-36), (G6.T-37) refer to the forms collected in Section~G6.Z1 of Annex~A.
The labels (Z10.s1), (Z10.s2), (Z10.s4) of Section~I.8.4 denote the three conditions in the definition of $e_2$ in Proposition~43.2, that is,
\[
e_2\le e_G,\qquad e_2\le e_{\sharp}(\beta_{\mathrm{ix}})\quad\left(\mathrm{ix}\in\{0,+,-\}\right),\qquad\frac52\,C'_{\bar{\kappa}}\cdot e_2^{\sigma}\le\frac{b(f)}{4}.
\]

\paragraph{Notation.}
For $\mathrm{ix}\in\{0,+,-\}$ put
\[
s_{\mathrm{ix}}:=\|g_{\mathrm{ix}}\|_2+\|\beta_{\mathrm{ix}}\|_{H^4};
\]
this is the number $s_f$ of Proposition~47.1\,(b5) for the probe $\beta_{\mathrm{ix}}$.
In Step~5 of Section~2.2 the letters $\mathrm{Rv}_1,\dots,\mathrm{Rv}_4$ denote random variables, and $\mathrm{kq}_2$, $\mathrm{kq}_4$ denote two multilinear forms.

The result \cite[Proposition~19.1]{QEDI} is applied with $h'=h$ and $m_0=0$ (Section~I.6.1); it does not involve the number $e_1$.
The result \cite[Proposition~19.2]{QEDI} is formulated with the number $e_1$ of \cite[Section~19.1.2]{QEDI}.
In its stead one uses the assertion (Sc6) of Section~G4.0.2, that is, Proposition~14.2\,(c), (f), together with the row (e1) at $v_k=v_K$, which holds by the second assertion of Proposition~14.3\,(f).
Here $e_1$ is the number of Section~I.8.2, line (C9), whose definition is given in Proposition~14.2\,(a), line (C9), and (b), rows (e1)--(e8).
The first three assertions of (Sc6) have a single datum, a number $e>0$ with $e_K=e\cdot h^{\frac12}\le e_1$ (Section~G4.0.2).
Of the number $e_1$ they require only $e_1\le1/3$ and the row (e1) at $v_k=v_K\ge-\log e_1$;
the row (e1) involves the constants $r$, $L$, $p_0$ only.

\begin{named}{(0.H)}
One has $h\ge1$ and $\frac1L<\bar{m}\cdot h\le1$.
\end{named}

\begin{proof}
As $\bar{m}\le1$, the number $\log_L(1/\bar{m})$ is non-negative; hence so is its integer part, and $h=L^{\left\lfloor\log_L\frac{1}{\bar{m}}\right\rfloor}\ge1$.
Moreover
\[
\log_L\left(\bar{m}\cdot h\right)=\left\lfloor\log_L\frac{1}{\bar{m}}\right\rfloor-\log_L\frac{1}{\bar{m}}\in(-1,0].
\]
The second assertion is the relation $\mu_K=\bar{m}\cdot h\in(1/L,1]$ of Section~I.6.1.
\end{proof}

\begin{named}{(0.G)}
One has $e_G>0$, and
\[
\left\{e'>0:\ e'\cdot h^{\frac12}\le e_1,\ e'\cdot h^{\frac12}\cdot\log h\le1\right\}=\left(0,e_G\right].
\]
Hence every $e\in(0,e_G]$ satisfies $e\cdot h^{\frac12}\le e_1$ and (Z8.w), that is, $e\cdot h^{\frac12}\cdot\log h\le1$; moreover $e_G\le e_1\cdot h^{-\frac12}$.
\end{named}

\begin{proof}
By (0.H), $\log h\ge0$.
The first condition on $e'$ reads $e'\le e_1\cdot h^{-\frac12}$.
If $h=1$, the second condition holds for every $e'$; if $h>1$, it reads $e'\le\left(h^{\frac12}\cdot\log h\right)^{-1}$.
Let $e''$ be the smallest of these upper bounds.
Since $e_1>0$ (Proposition~14.3\,(f)) and $1\le h<\infty$, one has $e''>0$.
Consequently the set is the interval $(0,e'']$, and its supremum, which is $e_G$ by Section~I.8.2, line (C10), equals $e''$.
\end{proof}

\begin{named}{(0.M)}
One has $\kappa_{(0)}>0$, and $\varepsilon_5(1/L)$ is a real number; consequently
\[
C_{\mathrm{mat}}=\left(L^3+1\right)\left(\varepsilon_5(1/L)+1\right)
\]
is a real number.
It is formed with $(L,C_{6},\kappa_{(0)})$, that is, with $L$ and $\rho$.
\end{named}

\begin{proof}
The formulas used are those of Section~I.8.2, line (C0r).
Since $\rho>0$ (Section~I.0.5) and $\lambda_{\mathrm{int}}>0$ (by (I.O8)), one has $\hat{g}=\rho\cdot\lambda_{\mathrm{int}}/2>0$.
The number $\theta_{0}$ is defined by $3(1+\rho)\left(\exp(\theta_{0})-1\right)=\hat{g}/2>0$;
therefore $\exp(\theta_{0})>1$, that is, $\theta_{0}>0$.
The number $c_{\mathrm{hl}}=\sqrt3\,(1+\rho)$ is finite and $\hat{g}>0$;
hence $0<K_{\rho}=\rho+2c_{\mathrm{hl}}^2/\hat{g}<\infty$, and $t_{0}=1/\left(4K_{\rho}^2\right)>0$.
Further $c_{*}=\min(w_3,\lambda_{\mathrm{int}})/18>0$, and it follows that
\[
c_{\mathrm{el}}=\min\left\{0.22,\ \frac{\rho\cdot c_{*}\cdot t_{0}}{\pi\sqrt3}\right\}>0.
\]
By Section~I.8.1, $K_{\mathrm{Str}}$ is a number (its finiteness is asserted in Proposition~14.2\,(f), \emph{Ranges}), and $K_{\mathrm{Str}}\ge0$, since it bounds a norm.
Hence $c_{\mathrm{el}}\cdot\theta_{0}/\left(12K_{\mathrm{Str}}\right)>0$, this quotient being read as $+\infty$ if $K_{\mathrm{Str}}=0$, and
\[
\kappa_{(0)}=\min\left\{\frac{\theta_{0}}{2},\ \frac{c_{\mathrm{el}}\cdot\theta_{0}}{12K_{\mathrm{Str}}}\right\}>0.
\]
It follows that $0<\exp\left(-\kappa_{(0)}/L\right)<1$, so that the logarithm of $\left(1-\exp\left(-\kappa_{(0)}/L\right)\right)^{-1}$ is a real number.
As $C_{6}$ is a number (Section~I.8.1; Proposition~14.2\,(f), \emph{Ranges}),
\[
\varepsilon_5(1/L)=C_{6}\cdot\log\left(\left(1-\exp\left(-\frac{\kappa_{(0)}}{L}\right)\right)^{-1}\right)
\]
is a real number.
By Section~I.8.2, line (C0r), the constants $C_{6}$ and $\kappa_{(0)}$ depend on $\rho$ only.
\end{proof}

\subsection*{1. Proposition 43.1}

\subsubsection*{1.1. The statement}

The statement proved below is Proposition~43.1 of Part~II.
Its content is the following.
The numbers $T^{*}(e)$, $R^{*}(e)$, $N_0(e)$, $m_{0,N}(e)$ and the lattice mass $\boldsymbol{\mu}_N=\bar{m}\cdot\eta+m_{0,N}(e)$ are those of Sections~I.6.2 and~I.2.5, and
\[
C_{\mathrm{mat}}=\left(L^3+1\right)\cdot\left(\varepsilon_5\left(\frac1L\right)+1\right).
\]
For every $e>0$ with $e\cdot h^{\frac12}\le e_1$ the clauses (R0)--(R2) hold.
\begin{itemize}
\item[(R0)] For every integer $N\ge0$ the number $m_{0,N}(e)$ is real, and $\left|m_{0,N}(e)\right|<\frac18$.
These numbers depend on $(N,e,\bar{m},\rho)$ and on the scheme (Section~I.6.2), and they depend neither on a torus nor on a twist;
consequently $N\mapsto m_{0,N}(e)$ is a mass assignment in the sense of Section~I.7.5.
\item[(R1)] One has $e\le\frac13$.
The number $T^{*}(e)$ is an integer with $T^{*}(e)\ge1$, one has $R^{*}(e)=h\cdot L^{T^{*}(e)}$, and $R^{*}(e)$ and $L\cdot R^{*}(e)$ both belong to $\mathrm{Tor}(e)$.
The length $\ell_c=\frac{R^{*}(e)}{2L}$ tends to $\infty$ as $e\to0$, and $\exp\left(-\bar{m}\cdot R^{*}(e)\right)\le e\cdot h^{\frac12}$.
If $e\le e_G$, then the inequality (W3) of Proposition~14.2\,(d) holds, namely
\[
6N_{*}\cdot e^2\cdot R^{*}(e)\cdot\log\frac1e\le\pi^2.
\]
\item[(R2)] For every $N\ge N_0(e)$, with $q=e\cdot\eta^{\frac12}$ and $\boldsymbol{\mu}=\boldsymbol{\mu}_N$, one has $D^{\mathrm{mat}}\le C_{\mathrm{mat}}$.
\end{itemize}
The clause (R6) does not contain $e$; it states that the exponents satisfy
\[
0<\mathfrak{e}\le\frac14,\qquad0\le\mathfrak{e}_2<1,\qquad0<\sigma=\mathfrak{e}-\frac{2a_{L}}{N_{*}}\le\frac14.
\]
Further, for every probe $\beta$ there are numbers $e_{\sharp}(\beta)\in\left(0,e_1\cdot h^{-\frac12}\right]$ and $C^{G}_{\beta}<\infty$, which depend only on the constants of the lines (C0)--(C9), on $\bar{m}$ and on $\beta$, with the property stated in (R3)--(R5) below.
These two numbers are fixed first, at the line (C10) of Section~I.8.2, and the charge $e$ is chosen after them, at the line (C11).
Let $0<e\le\min\left(e_G,e_{\sharp}(\beta)\right)$, let $N\ge N_0(e)$, and let $v$ be any placement;
put $\boldsymbol{\mu}=\boldsymbol{\mu}_N$ and $c_m=N_0^{(m)}(0)$, where $N_0$ is the free ratio on $C(n_*)$ (Section~I.7.6).
Then the following clauses hold.
\begin{itemize}
\item[(R3)] The probe $\beta$ fits $\left(N,T^{*}(e)\right)$, and
\[
\Theta^{(0)}_4\le C^{G}_{\beta}\cdot e^{\mathfrak{e}},\qquad\left(1+|c_{2}|\right)\cdot\Theta^{(0)}_2\le C^{G}_{\beta}\cdot e^{\mathfrak{e}}.
\]
\item[(R4)] The numbers $c_{2}$ and $c_{4}$ are real, they do not depend on the placement, and
\[
|c_{2}|+|c_{4}|^{\frac12}\le C^{G}_{\beta}\cdot e^{-\mathfrak{e}_2}.
\]
\item[(R5)] One has $c_{4}-3c_{2}^2=b_{\mathrm{ch}}\left(\eta,R^{*}(e)\right)$, and
\[
\left|c_{4}-3c_{2}^2-b\left(d\,d^{*}\beta\right)\right|\le C^{G}_{\beta}\cdot\left(e+\eta^{\frac14}\right).
\]
\end{itemize}
The constant is
\[
C^{G}_{\beta}=\max\left\{C_{35}(\beta),\ C_{\mathrm{R4}},\ C_{1}^{\mathrm{N2}},\ C_{2}^{\mathrm{N2}}\cdot\exp\left(\bar{m}\cdot D_{\mathrm{be}}\right)\cdot h^{\frac12}\right\},
\]
where
\begin{align*}
C_{35}(\beta)&=3\cdot10^5\,\mathrm{Cb}^4\cdot h^{\frac18},\\
C_{\mathrm{R4}}&=\mathrm{Cpv}_{\beta}+\left(\left|b\left(d\,d^{*}\beta\right)\right|+C_{1}^{\mathrm{N2}}+C_{2}^{\mathrm{N2}}+3\,\mathrm{Cpv}_{\beta}^2\right)^{\frac12}.
\end{align*}
The first of these is the constant of Proposition~35.9\,(b).
The second is the constant of Proposition~37.3\,(b);
as $R^{*}(e)\ge D_{\mathrm{be}}+1$, one has $\exp\left(-\bar{m}\cdot\left(R^{*}(e)-D_{\mathrm{be}}\right)\right)\le1$.
The numbers $\mathrm{Cb}$, $\mathrm{Cpv}_{\beta}$, $C_{1}^{\mathrm{N2}}$, $C_{2}^{\mathrm{N2}}$ and $D_{\mathrm{be}}$ are those of the probe $\beta$ in Section~G5.1.3, Proposition~36\,(g) and Proposition~37.2\,(d), (e).

\subsubsection*{1.2. Proof}

\begin{proof}[Proof of Proposition~43.1]
Throughout, $e>0$ is a charge with $e_K=e\cdot h^{\frac12}\le e_1$.

\textbf{Step 1 (scales; (R1)).}

\emph{(1.1)} By (0.H), $h\ge1$ and $\frac1L<\bar{m}\cdot h\le1$.

\emph{(1.2) Claim.} One has $e\le\frac13$, $v_K\ge-\log e_1$, $v_K\ge8$, and
\[
\left(1+\frac{\mathrm{lg}}{v_K}\right)^r\le\left(1+\frac{3\,\mathrm{lg}}{v_K}\right)^r<L.
\]
\emph{Proof of (1.2).}
By (0.H), $e=e_K\cdot h^{-\frac12}\le e_K\le e_1$, and $e_1\le\frac13$ by Section~I.8.2, line (C9).
The inequality $e_K\le e_1$ gives $v_K=-\log e_K\ge-\log e_1$.
The third assertion, and the last inequality of the fourth, are contained in the assertion (Sc6) of Section~G4.0.2, that is, in Proposition~14.2\,(c), (f) together with the row (e1) at $v_k=v_K$, which holds by the second assertion of Proposition~14.3\,(f) (Section~0):
they are the first and the third inequality of the row (e1) at $v_k=v_K$.
Since $\mathrm{lg}>0$ (because $L>1$), $v_K>0$ and $r\ge1$, one has $\left(1+\frac{\mathrm{lg}}{v_K}\right)^r\le\left(1+\frac{3\,\mathrm{lg}}{v_K}\right)^r$.

\emph{Consequence.}
The rows (e3) and (e8) of Proposition~14.2\,(b), which are imposed for all $v_k\ge-\log e_1$ (Section~0), hold at $v_k=v_K$.

\emph{(1.3) Claim.} One has $v_K^r\le r_K<L\cdot v_K^r$ and $r_K\ge v_K$.

\emph{Proof of (1.3).}
By its definition, $r_K\ge v_K^r$.
Since $r\ge2$ (Section~I.8.2, line (C7)) and $v_K\ge8$ by (1.2), one has $v_K^r\ge64>1=L^0$.
Hence $r_K=L^i$ with an integer $i\ge1$, and the minimality of $r_K$ gives $L^{i-1}<v_K^r$, that is, $r_K<L\cdot v_K^r$.
Finally, $v_K\ge1$ and $r\ge1$ imply $v_K^r\ge v_K$.

\emph{(1.4) Claim.} $T^{*}(e)$ is an integer with $T^{*}(e)\ge1$; $R^{*}(e)=h\cdot L^{T^{*}(e)}$; both $R^{*}(e)$ and $L\cdot R^{*}(e)$ belong to $\mathrm{Tor}(e)$; and $R^{*}(e)\ge L$.

\emph{Proof of (1.4).}
As $M$ is a power of $L$ with $M\ge L^2$ (Section~I.0.5), $\log_LM$ is an integer with $\log_LM\ge2$; and $\log_Lr_K$ is a non-negative integer.
Therefore $T^{*}(e)=\log_LM+\log_Lr_K-1$ is an integer with $T^{*}(e)\ge1$.
Moreover
\begin{gather*}
R^{*}(e)=\frac{h\cdot M\cdot r_K}{L}=h\cdot L^{\log_L(M\cdot r_K)-1}=h\cdot L^{T^{*}(e)},\\
L\cdot R^{*}(e)=R_{\mathrm{top}}(e)=h\cdot L^{T^{*}(e)+1}.
\end{gather*}
Both numbers are of the form $h\cdot L^T$ with an integer $T\ge0$, and neither of them exceeds $R_{\mathrm{top}}(e)$; hence both lie in $\mathrm{Tor}(e)$.
From $h\ge1$ and $T^{*}(e)\ge1$ one obtains $R^{*}(e)\ge L$.

\emph{(1.5) Claim.} $\ell_c\to\infty$ as $e\to0$.

\emph{Proof of (1.5).}
By (1.3),
\[
\ell_c=\frac{R^{*}(e)}{2L}=\frac{h\cdot M\cdot r_K}{2L^2}\ge\frac{h\cdot M\cdot v_K^r}{2L^2}.
\]
The numbers $h$, $M$, $L$, $r$ do not depend on $e$, whereas $v_K=-\log\left(e\cdot h^{\frac12}\right)\to\infty$ as $e\to0$.

\emph{(1.6) Claim.} $\exp\left(-\bar{m}\cdot R^{*}(e)\right)<e\cdot h^{\frac12}$.

\emph{Proof of (1.6).}
By (0.H), by $M\ge L^2$ (Section~I.0.5) and by (1.3),
\[
\bar{m}\cdot R^{*}(e)=\left(\bar{m}\cdot h\right)\cdot\frac{M\cdot r_K}{L}>\frac{M\cdot r_K}{L^2}\ge r_K\ge v_K.
\]
Therefore $\exp\left(-\bar{m}\cdot R^{*}(e)\right)<\exp(-v_K)=e_K$.

\emph{(1.7) Claim.} The numbers $N_{\min}\left(R_{\mathrm{top}}(e),e\right)$ and $N_0(e)$ are finite non-negative integers, and $N_{\min}\left(R_{\mathrm{top}}(e),e\right)\le N_0(e)$.

\emph{Proof of (1.7).}
By (1.2), $v_K\ge1$ and $\left(1+\frac{\mathrm{lg}}{v_K}\right)^r\le L$.
These are the hypotheses of \cite[Proposition~19.1]{QEDI}, which therefore applies to the tori $R^{*}(e)$ and $R_{\mathrm{top}}(e)$ of $\mathrm{Tor}(e)$ (see (1.4)).
By its clauses (b) and (c),
\[
N_{\min}(R,e)=\max\left\{0,\ \mathrm{js}(R,e)+1-\log_Lh\right\},
\]
where $\mathrm{js}(R,e)\ge0$ is an integer, the integer $j_{*}(R,e)$ of \cite[Proposition~19.1(b)]{QEDI}; and $\log_Lh$ is an integer.
By its clause (d), applied to $R^{*}(e)\le R_{\mathrm{top}}(e)$,
\[
N_{\min}\left(R_{\mathrm{top}}(e),e\right)\le N_{\min}\left(R^{*}(e),e\right)=N_0(e).
\]

\emph{Consequence} (Section~I.6.2).
For every $N\ge N_0(e)$ the pair $P_N(e)$ is defined, and $m_{0,N}(e)$ is its first entry.

\emph{(1.8) Claim.} If $e\le e_G$, then (W3) holds.

\emph{Proof of (1.8).}
By (0.G), $e_K\cdot\log h\le1$.
By (1.3) and (1.2), the latter giving $e\le\frac13<1$,
\begin{gather*}
e^2\cdot R^{*}(e)=\frac{e^2\cdot h\cdot M\cdot r_K}{L}=\frac{e_K^2\cdot M\cdot r_K}{L}<e_K^2\cdot M\cdot v_K^r,\\
0<\log\frac1e=-\log\left(e_K\cdot h^{-\frac12}\right)=v_K+\frac12\log h.
\end{gather*}
It follows that
\begin{align*}
6N_{*}\cdot e^2\cdot R^{*}(e)\cdot\log\frac1e
&<6N_{*}\cdot M\cdot v_K^r\cdot e_K\cdot\left(e_K\cdot v_K+\frac12\,e_K\cdot\log h\right)\\
&\le6N_{*}\cdot M\cdot v_K^r\cdot e_K\cdot\left(e_K\cdot v_K+\frac12\right)\\
&=\left.6N_{*}\cdot M\cdot v_k^r\cdot\exp(-v_k)\cdot\left(v_k\cdot\exp(-v_k)+\frac12\right)\right|_{v_k=v_K}\le\pi^2,
\end{align*}
where the last inequality is the row (e8) at $v_k=v_K$ (Consequence of (1.2)).
This is (W3).

Together, the claims (1.2), (1.4), (1.5), (1.6) and (1.8) are the assertions of (R1).

\textbf{Step 2 ((R6)).}
By Section~I.8.2, lines (C0), (C1) and (C7), one has $\alpha<1$ and $1-\alpha<\varepsilon\le\bar{\varepsilon}=1/200$; hence $0<\varepsilon\le1/200$.
Moreover $0<s<1$.
The rows $\mathfrak{e}$ and $\sigma$ of Proposition~14.2\,(a), which agree with Section~I.8.5, give
\[
\mathfrak{e}=s\left(\frac14-9\varepsilon\right)-12\varepsilon,\qquad\sigma=\mathfrak{e}-\frac{2a_{L}}{N_{*}}.
\]
By the row (S2) of Proposition~14.2\,(b), a condition that can be met by Proposition~14.3\,(d), one has $\sigma>0$.
Since $a_{L}>0$ and $N_{*}>0$, it follows that $\mathfrak{e}=\sigma+2a_{L}/N_{*}>\sigma>0$.
From $\varepsilon>0$ and $s>0$ one obtains $\mathfrak{e}<s/4<1/4$.
Hence $0<\sigma<\mathfrak{e}<1/4$.
Finally $\mathfrak{e}_2=0$ (Section~I.8.5), so that $0\le\mathfrak{e}_2<1$.

\textbf{Step 3 ((R0)).}
By (1.4) the side $R_{\mathrm{top}}(e)$ belongs to $\mathrm{Tor}(e)$, and by (1.7) the number $N_{\min}\left(R_{\mathrm{top}}(e),e\right)$ is finite.
Consequently, for every integer $N\ge N_{\min}\left(R_{\mathrm{top}}(e),e\right)$, the quadruple $\left(\bar{m},e,R_{\mathrm{top}}(e),N\right)$ is a set of data (FL) in the sense of Section~I.1.8, and Proposition~28.3 applies to it.

Let $N\ge N_{\min}\left(R_{\mathrm{top}}(e),e\right)$.
By Proposition~28.3\,(c), the pair $P_N(e)$ is real, and it is a function of $(N,e,\bar{m})$ and of the scheme, that is, of the constants of the lines (C0)--(C9) (among them $\rho$ and the weights), of the basis $\{e_i^0\}$ and of the systems of primitives on $T_k\left(R_{\mathrm{top}}(e)\right)$, $k<k_{\mathrm{top}}$.
No torus is among the arguments of $P_N(e)$;
in particular the datum $R$ of (FL) is not one of them: for given $(N,e,\bar{m})$ the pair $P_N(e)$ is the same for all $R\in\mathrm{Tor}(e)$ with $N\ge N_{\min}(R,e)$.
The only torus that occurs in the definition of $P_N(e)$ is $R_{\mathrm{top}}(e)$, which lies in $\mathrm{Tor}(e)$ by (1.4) and is itself a function of $(e,\bar{m})$ and of the constants $L$, $M$, $r$.
No twist enters the definition of $P_N(e)$.
By Section~I.6.2, $m_{0,N}(e)$ is the first entry of $P_N(e)$.

Let $0\le N<N_{\min}\left(R_{\mathrm{top}}(e),e\right)$.
Then $m_{0,N}(e)=0$ (Section~I.6.2).
The threshold $N_{\min}\left(R_{\mathrm{top}}(e),e\right)$ is formed with $(e,h,L,M,r)$ (Sections~I.6.1 and~I.6.2), that is, with $e$, $\bar{m}$ and constants of the scheme.

Hence, for every $N\ge0$, the number $m_{0,N}(e)$ is real and is a function of $(N,e,\bar{m},\rho)$ and of the scheme; no torus and no twist is among its arguments.
By Proposition~28.3\,(d),
\[
\left|m_{0,N}(e)\right|\le\frac{w^{\mathrm{m}}_0}{8}<\frac18\qquad\text{for all }N\ge0
\]
(for $N<N_{\min}\left(R_{\mathrm{top}}(e),e\right)$ the left side vanishes), and $N\mapsto m_{0,N}(e)$ is a real sequence that does not depend on a torus.
In particular it does not depend on the index $T$ of the volume (Section~I.6.5), and so it is a mass assignment in the sense of Section~I.7.5.

\textbf{Step 4 ((R2)).}
Let $N\ge N_0(e)$, $q=e\cdot\eta^{\frac12}$ and $\boldsymbol{\mu}=\boldsymbol{\mu}_N=\bar{m}\cdot\eta+m_{0,N}(e)$.
By (1.4), $n_*=h\cdot L^{T^{*}(e)+N}$ is an integer with $n_*\ge L$, and $L\ge5$ (Section~I.0.5); so $n_*$ and $L\cdot n_*$ are integers $\ge4$.
By Step~3, $m_{0,N}(e)$ is real.
Proposition~38.2\,(c), applied on the cubic tori $n=n_*$ and $n=L\cdot n_*$ with the real pair $\left(m_{\mathrm{in}},\varepsilon^{\mathrm{v}}_{\mathrm{in}}\right)=\left(m_{0,N}(e),0\right)$, gives
\[
\mathrm{zr}_n\left(q,\boldsymbol{\mu}_N\right)>0.
\]
Hence both logarithms in $D^{\mathrm{mat}}=L^3\log\mathrm{zr}_{n_*}-\log\mathrm{zr}_{L n_*}$ (Section~I.7.6) are defined.

By (1.7), $m_{0,N}(e)$ is the first entry of $P_N(e)$;
therefore the number $\mathrm{zr}_n=\mathrm{zr}_n\left(q,\boldsymbol{\mu}_N\right)$ of Section~I.7.6 is the value $\mathrm{zr}_n\left(e\cdot\eta^{\frac12},\bar{m}\cdot\eta+m_{0,N}(e)\right)$ that occurs in Proposition~32.3.
By (10.3) of Proposition~32.3, for $N\ge N_0(e)$,
\begin{align*}
\left|D^{\mathrm{mat}}\right|
&\le\left(L^3+1\right)\left(\varepsilon_5\left(\bar{m}\cdot R^{*}(e)\right)+3e_K^{\frac14-9\varepsilon}\right)\\
&\le\left(L^3+1\right)\left(\varepsilon_5(1/L)+1\right)=C_{\mathrm{mat}}.
\end{align*}
(The second inequality is part of (10.3).
It rests on two facts.
First, $\varepsilon_5$ is a decreasing function (Section~G4.0.3), and $\bar{m}\cdot R^{*}(e)>r_K\ge1>1/L$ by (1.6) and (1.3), so that $\varepsilon_5\left(\bar{m}\cdot R^{*}(e)\right)\le\varepsilon_5(1/L)$.
Second, the third inequality of the row (e3) at $v_k=v_K$ (Consequence of (1.2)) reads $e_K^{\frac14-9\varepsilon}\le\frac18$, so that $3e_K^{\frac14-9\varepsilon}\le\frac38\le1$.)
In particular $D^{\mathrm{mat}}\le C_{\mathrm{mat}}$.
By (0.M), $C_{\mathrm{mat}}$ is a real number formed with $(L,C_{6},\kappa_{(0)})$; it contains neither $e$ nor $N$.

\textbf{Step 5 (the numbers $e_{\sharp}(\beta)$ and $C^{G}_{\beta}$).}
Let $\beta$ be a probe.

The number $e_{\sharp}(\beta)$ is that of Section~G5.1.9: one half of a supremum of numbers in $\left(0,e_1\cdot h^{-\frac12}\right]$.
Hence $e_{\sharp}(\beta)\le e_1\cdot h^{-\frac12}/2$, and $e_{\sharp}(\beta)>0$ by Proposition~35.10.
By Proposition~35.10, the number $e_{\sharp}(\beta)$ depends on the constants of the lines (C0)--(C9) of Section~I.8.2, the line (C6$'$) among them, on $\bar{m}$, and on $\beta$ through $\left(C_{\beta},\mathbf{D},D_\beta\right)$;
it depends neither on $N$, nor on a torus, nor on a placement, nor on a twist.

By its definition in Proposition~43.1 (paragraph \emph{The constant}), $C^{G}_{\beta}$ is the maximum of the four numbers
\[
C_{35}(\beta),\qquad C_{\mathrm{R4}},\qquad C_{1}^{\mathrm{N2}},\qquad C_{2}^{\mathrm{N2}}\cdot\exp\left(\bar{m}\cdot D_{\mathrm{be}}\right)\cdot h^{\frac12}.
\]
These numbers are finite, and they depend only on the following quantities.
\begin{itemize}
\item[--] $C_{35}(\beta)=3\cdot10^5\,\mathrm{Cb}^4\cdot h^{\frac18}$, with $\mathrm{Cb}=1+h^{\frac12}\cdot C_{\beta}$;
by Proposition~35.3\,(c1), $C_{\beta}\le2h^{\frac12}\cdot\left(1+h+3h^2\right)\cdot|\beta|_{C^3}<\infty$, and $C_{\beta}$ depends on $(\beta,h,L,\alpha)$.
\item[--] $C_{\mathrm{R4}}=\mathrm{Cpv}_{\beta}+\left(\left|b\left(d\,d^{*}\beta\right)\right|+C_{1}^{\mathrm{N2}}+C_{2}^{\mathrm{N2}}+3\,\mathrm{Cpv}_{\beta}^2\right)^{\frac12}$;
here $\mathrm{Cpv}_{\beta}=\mathrm{Cpv}_{\beta}(\rho,\bar{m},\beta)<\infty$ by (G6.T-36),
while $C_{1}^{\mathrm{N2}}=C_{1}^{\mathrm{N2}}(\rho,\bar{m},\beta)<\infty$ and $C_{2}^{\mathrm{N2}}=C_{2}^{\mathrm{N2}}(\bar{m},\beta)<\infty$, and $b\left(d\,d^{*}\beta\right)$ is a real number formed with $(\bar{m},\beta)$, by (G6.T-37) and Section~I.7.6.
\item[--] $C_{1}^{\mathrm{N2}}$.
\item[--] $C_{2}^{\mathrm{N2}}\cdot\exp\left(\bar{m}\cdot D_{\mathrm{be}}\right)\cdot h^{\frac12}$, where $D_{\mathrm{be}}$ is a function of $\beta$ (Section~I.7.6).
\end{itemize}
Since $h$ is a function of $(L,\bar{m})$, the number $C^{G}_{\beta}<\infty$ depends on constants of the lines (C0)--(C9) (namely $\rho$, $\alpha$, $L$), on $\bar{m}$ and on $\beta$;
it contains neither $e$, nor $N$, nor a placement, nor a torus.

From now on let
\[
0<e\le\min\left(e_G,e_{\sharp}(\beta)\right),\qquad N\ge N_0(e),\qquad\boldsymbol{\mu}=\boldsymbol{\mu}_N.
\]
By Proposition~35.10, for these $e$ and $N$:
\begin{itemize}
\item[(5.A)] $e\cdot h^{\frac12}\le e_1$;
\item[(5.B)] every centred placement (Section~G5.1.3) is admitted;
\item[(5.C)] $D_\beta+3\eta\le\ell_c$;
\item[(5.D)] $R^{*}(e)\ge D_{\mathrm{be}}+1$.
\end{itemize}
A placement in the sense of Section~I.7.6 is a vector $v\in\eta\mathbb{Z}^3$ with
\[
\left|v+y_\beta-\frac{\ell_c}{2}(1,1,1)\right|\le\frac{\eta}{2},
\]
where $|\cdot|$ is the sup norm (Section~I.1.1); this is the definition of a centred placement in Section~G5.1.3.
Placements exist, since every point of $\mathbb{R}^3$ has sup distance at most $\eta/2$ from $\eta\mathbb{Z}^3$.
Let $v$ be any placement.
The lattice probe $a_{\eta}$ of Section~I.7.6, formed with $\tau_v\beta$ on $\Lambda_{\eta,R^{*}(e)}$, is the cochain $a_{\eta}[\tau_v\beta]$ of Section~G5.1.3, the two defining displays being the same.
Put $\hat{\mathfrak{a}}=\eta^{\frac12}\cdot a_{\eta}$ and, on $C(n_*)$,
\[
N_0(t)=\frac{\det\left(D_{(t/e)\hat{\mathfrak{a}}}+\bar{m}\cdot\eta\right)}{\det\left(D_0+\bar{m}\cdot\eta\right)}
\]
(Sections~I.7.6 and~G5.1.5);
by Proposition~46\,(g+), this is the function $N_0$ of Section~G5.6.0 for the probe $\tau_v\beta$ on the torus of side $R^{*}(e)$.
Finally, $c_{m}=N_0^{(m)}(0)$.

\textbf{Step 6 ((R3)).}
By (5.C), the probe $\beta$ fits $\left(N,T^{*}(e)\right)$ (Section~I.7.6).
The data (DAT) of Section~G5.4.0 are at hand.
Indeed, the setting (ST*)$^{\mathrm{ch}}$ of Section~G5.1.1 consists of the constants, of $\bar{m}$, of a charge $e$ with $e\cdot h^{\frac12}\le e_1$ (by (5.A)), of the torus $R^{*}(e)$, of an integer $N\ge N_0(e)$, and of the pair $P_N(e)$, whose first entry is $m_{0,N}(e)$ by (1.7), so that the lattice mass of (ST*)$^{\mathrm{ch}}$ is $\boldsymbol{\mu}_N$;
moreover $e\le e_{\sharp}(\beta)$, and $v$ is admitted by (5.B).
The number $c_{2}$ of Proposition~35.9 is $N_0''(0)$ (Section~G5.1.5), that is, the number $c_{2}$ defined above; in this form the two bounds are stated in (G6.T-35).
By Proposition~35.9\,(b) and the last sentence of Proposition~35.9 (for a centred placement, $\mathring{\Phi}$, $N_0$, $c_{m}$ and $\Theta^{(0)}_m$ are the objects of Section~I.7.6),
\begin{gather*}
\Theta^{(0)}_4\le3\cdot10^5\,\mathrm{Cb}^4\cdot h^{\frac18}\cdot e^{\mathfrak{e}}=C_{35}(\beta)\cdot e^{\mathfrak{e}}\le C^{G}_{\beta}\cdot e^{\mathfrak{e}},\\
\left(1+|c_{2}|\right)\cdot\Theta^{(0)}_2\le3\cdot10^4\,\mathrm{Cb}^4\cdot h^{\frac18}\cdot e^{\mathfrak{e}}\le C_{35}(\beta)\cdot e^{\mathfrak{e}}\le C^{G}_{\beta}\cdot e^{\mathfrak{e}}.
\end{gather*}

\textbf{Step 7 ((R4)).}
Apply Proposition~50\,(h) with $n=(n_*,n_*,n_*)$, with $\mathbf{m}=\bar{m}\cdot\eta>0$, and with the real one-cochain $\hat{\mathfrak{a}}/e$ in place of $a$:
the function $N_0$ is entire and even, it is real on the real axis, and $N_0(0)=1$.
Hence all derivatives of $N_0$ at $0$ are real;
in particular $c_{2}$ and $c_{4}$ are real, and $c_{1}=c_{3}=0$.
The data of Proposition~37.3 are at hand:
$e\cdot h^{\frac12}\le e_1$ by (5.A), $R^{*}(e)\ge D_{\mathrm{be}}+1$ by (5.D), $N\ge N_0(e)$, and $v\in\eta\mathbb{Z}^3$.
By Proposition~37.3\,(a), $N_0$ does not depend on $v\in\eta\mathbb{Z}^3$; consequently $c_{2}$ and $c_{4}$ do not depend on the placement.
By Proposition~37.3\,(b) and $\mathfrak{e}_2=0$,
\begin{align*}
|c_{2}|+|c_{4}|^{\frac12}
&\le\mathrm{Cpv}_{\beta}+\left(\left|b\left(d\,d^{*}\beta\right)\right|+C_{1}^{\mathrm{N2}}+C_{2}^{\mathrm{N2}}+3\,\mathrm{Cpv}_{\beta}^2\right)^{\frac12}\\
&=C_{\mathrm{R4}}\le C^{G}_{\beta}=C^{G}_{\beta}\cdot e^{-\mathfrak{e}_2}.
\end{align*}

\textbf{Step 8 ((R5)).}
By Step~7, near $t=0$,
\[
N_0(t)=1+\frac{c_{2}}{2}\,t^2+\frac{c_{4}}{24}\,t^4+O\left(t^6\right),
\]
and hence
\[
\log N_0(t)=\frac{c_{2}}{2}\,t^2+\left(\frac{c_{4}}{24}-\frac{c_{2}^2}{8}\right)t^4+O\left(t^6\right).
\]
Therefore the number $b_{\mathrm{ch}}\left(\eta,R^{*}(e)\right)=\left.\left(\frac{d}{dt}\right)^4\log N_0(t)\right|_{t=0}$ of Section~I.7.6 equals
\[
24\left(\frac{c_{4}}{24}-\frac{c_{2}^2}{8}\right)=c_{4}-3c_{2}^2.
\]
(By Proposition~37.3\,(a), this number is also the number $b_{\mathrm{ch}}\left(\eta,R^{*}(e)\right)[\beta]$ of Section~G5.6.0, formed with $\beta$ itself.)
By Proposition~37.3\,(c), with $f=d\,d^{*}\beta$,
\begin{align*}
\left|c_{4}-3c_{2}^2-b\left(d\,d^{*}\beta\right)\right|
&\le C_{1}^{\mathrm{N2}}\cdot\eta^{\frac14}+C_{2}^{\mathrm{N2}}\cdot\exp\left(\bar{m}\cdot D_{\mathrm{be}}\right)\cdot h^{\frac12}\cdot e\\
&\le C^{G}_{\beta}\cdot\left(\eta^{\frac14}+e\right).
\end{align*}

\textbf{Step 9 (uniformity).}
The numbers $e_{\sharp}(\beta)$ and $C^{G}_{\beta}$ were fixed in Step~5, before $e$, $N$ and $v$ were chosen.
Hence the bounds of Steps~6--8 hold, with these two numbers, for every $e\in\left(0,\min\left(e_G,e_{\sharp}(\beta)\right)\right]$, every $N\ge N_0(e)$ and every placement $v$.
Likewise $C_{\mathrm{mat}}$ was fixed before $e$ and $N$, and (R2) holds for every $e$ with $e\cdot h^{\frac12}\le e_1$ and every $N\ge N_0(e)$.
\end{proof}

\subsection*{2. Proposition 43.2}

\subsubsection*{2.1. The statement}

The statement proved below is Proposition~43.2 of Part~II, together with its data and its constants.
They read as follows.

\emph{Data.}
Let $\beta$ be a probe, and let $r_\beta>0$ be a number such that $\operatorname{supp}\beta$ lies in a ball of radius $r_\beta$.
This number is a datum in addition to $(\beta,x_0)$;
it enters only $d_0$ and the inequality displayed below, and none of the constants below depends on it.
Let
\[
a=d^{*}\beta,\qquad\|a\|_\infty=\sup_{y\in\mathbb{R}^3}|a(y)|_2,\qquad f=d\,d^{*}\beta,
\]
and assume $b(f)>0$; such probes exist by \cite[Theorem~21.1(a)]{QEDI}.
Let $x_0\in\mathbb{R}^3$ be a point with $d_0:=|x_0|_2-2r_\beta>0$ for which the condition of \cite[Section~2.8]{QEDI} holds:
\begin{multline*}
\frac{6}{\bar{m}^2}\cdot\|a\|_\infty^4\cdot\left(\frac{4\pi\cdot r_\beta^3}{3}\right)^2\cdot\left(\bar{m}^2+\left(\bar{m}+\frac{1}{d_0}\right)^2\right)\\
\cdot\frac{\exp\left(-2\bar{m}\cdot d_0\right)}{4\pi^2\cdot d_0^2}\le\frac{b(f)}{12}.
\end{multline*}
The three probes are $\beta_0=\beta$ and $\beta_\pm=\beta\pm\tau_{x_0}\beta$.
The tests are $g_{\mathrm{ix}}=d\,d^{*}\beta_{\mathrm{ix}}$ for $\mathrm{ix}\in\{0,+,-\}$, so that $g_0=f$ and $g_\pm=f\pm\tau_{x_0}f$.
For $0<\eta\le1$, $g^{(\eta)}_{\mathrm{ix}}$ is the two-form $f^{(\eta)}$ defined at the beginning of Proposition~47.1, formed with the probe $\beta_{\mathrm{ix}}$ and not with the test $g_{\mathrm{ix}}$;
below, $\eta=L^{-N}$.

\emph{Constants.}
The following constants belong to the line (C10) and are fixed in the order in which they are listed;
each of them depends only on the preceding ones and on $(\bar{m},\beta,x_0)$:
\begin{gather*}
s_{B}=\max_{\mathrm{ix}}\left(\|g_{\mathrm{ix}}\|_2+\|\beta_{\mathrm{ix}}\|_{H^4}\right),\qquad C_{3\mathrm{p}}=\max_{\mathrm{ix}}C^{G}_{\beta_{\mathrm{ix}}},\\
C_{\Gamma}^{(16)}=(2\pi)^{\frac{a_{L}}{N_{*}}}\cdot\exp\left(\frac{K_{L}+C_{\mathrm{mat}}+1}{(L-1)^3\cdot N_{*}}\right),
\end{gather*}
\begin{multline*}
C_{U}=\max\left\{C_{\Gamma}^{(16)}\cdot\left(C_{3\mathrm{p}}+1200N_{*}^2\cdot s_{B}^4+2C_{3\mathrm{p}}^2\right),\right.\\
\left.C_{\Gamma}^{(16)}\cdot\left(C_{3\mathrm{p}}+2\left(12N_{*}\cdot s_{B}^2+C_{3\mathrm{p}}\right)\cdot\left(1+C_{3\mathrm{p}}\right)\right),\ C_{3\mathrm{p}},\ 60\,s_{B}^4\right\},
\end{multline*}
\[
C_{\bar{\kappa}}=2\left(8C_{U}+3C_{U}^2\right),\qquad C'_{\bar{\kappa}}=8C_{\bar{\kappa}}+3C_{\bar{\kappa}}^2;
\]
and $e_2$ is the largest number such that
\[
e_2\le e_G,\qquad e_2\le e_{\sharp}(\beta_{\mathrm{ix}})\quad\text{for }\mathrm{ix}\in\{0,+,-\},\qquad\frac52\,C'_{\bar{\kappa}}\cdot e_2^{\sigma}\le\frac{b(f)}{4}.
\]

\emph{Assertions.}
Let $e\in(0,e_2]$, and let $N\ge N_0(e)$ and $T\ge T^{*}(e)+1$ be integers.
Put
\[
\eta=L^{-N},\qquad n=8h\cdot L^{T+N},\qquad R=8h\cdot L^{T},\qquad q=e\cdot\eta^{\frac12},\qquad\boldsymbol{\mu}=\boldsymbol{\mu}_N=\bar{m}\cdot\eta+m_{0,N}(e),
\]
and consider the law $E^{\mathrm{af}}_n$, under which $X_{\eta,R}(\,\cdot\,)$ has the law given by $P_{N,T}$.
For $\mathrm{ix}\in\{0,+,-\}$ let $v$ be a placement of $\beta_{\mathrm{ix}}$, and let
\[
c_m=c_m(N,e,\beta_{\mathrm{ix}})=N_0^{(m)}(0),\qquad\boldsymbol{\sigma}^2=\boldsymbol{\sigma}^2(N,\beta_{\mathrm{ix}}),\qquad\mathsf{u}_m=\mathsf{u}_m(N,e,\beta_{\mathrm{ix}};v)\qquad(m=2,4)
\]
be the numbers of Proposition~42.2;
here $\boldsymbol{\sigma}^2=\|d\hat{\mathfrak{a}}\|^2$ is the number of Section~I.7.6, formed with $\tau_v\beta_{\mathrm{ix}}$, and not the square of the exponent $\sigma$ of Section~I.8.5;
through $\Theta^{(0)}_m$ the number $\mathsf{u}_m$ may depend on the placement $v$.
For every placement $v$ of $\beta_{\mathrm{ix}}$ the following clauses hold.
\begin{itemize}
\item[(U-a)] One has
\begin{multline*}
\left|\kappa_4\left(X_{\eta,R}(g_{\mathrm{ix}})\right)-e^4\cdot b(g_{\mathrm{ix}})\right|\le\frac{C_{\bar{\kappa}}}{2}\cdot e^{4+\sigma}+C_{U}\cdot\left(\eta^2+e^4\cdot\eta^{\frac14}\right)\\
\le C_{\bar{\kappa}}\cdot\left(e^{4+\sigma}+\eta^{\frac14}\right).
\end{multline*}
\item[(U-a$'$)] One has
\begin{gather*}
\left|\kappa_4\left(X_{\eta,R}(g_{\mathrm{ix}})\right)-e^4\cdot b_{\mathrm{ch}}\left(\eta,R^{*}(e)\right)[\beta_{\mathrm{ix}}]\right|\le\left(7C_{U}+3C_{U}^2\right)\cdot e^{4+\sigma}+C_{U}\cdot\eta^2,\\
\left|\kappa_2\left(X_{\eta,R}(g_{\mathrm{ix}})\right)-\boldsymbol{\sigma}^2-e^2\cdot c_{2}\right|\le e^2\cdot\mathsf{u}_2+2s_{B}^2\cdot\eta^2.
\end{gather*}
\item[(U-b)] For every finite subset $\mathrm{Ypt}$ of $\frac{R^{*}(e)}{L}\,\mathbb{Z}^3$ whose points are pairwise incongruent modulo $R\mathbb{Z}^3$, and for $m=2,4$,
\[
\left|E^{\mathrm{af}}_n\prod_{y\in\mathrm{Ypt}}\mathrm{Pw}_m\left(X_{\eta,R}\left(\tau_y\,g^{(\eta)}_{\mathrm{ix}}\right);\boldsymbol{\sigma}^2,c_m\right)\right|\le\mathsf{u}_m^{\#|\mathrm{Ypt}|}.
\]
\item[(U-c)] One has
\begin{gather*}
\mathsf{u}_4\le C_{U}\cdot e^{\sigma},\qquad\left(1+|c_{2}|\right)\cdot\mathsf{u}_2\le C_{U}\cdot e^{\sigma},\\
|c_{2}|\le C_{U}\cdot e^{-\mathfrak{e}_2},\qquad\left|c_{4}-3c_{2}^2-b(g_{\mathrm{ix}})\right|\le C_{U}\cdot\left(e^{\sigma}+\eta^{\frac14}\right).
\end{gather*}
\end{itemize}
In consequence, the following holds.
The exponents $\sigma$ and $\mathfrak{e}_2$ of Section~I.8.5 satisfy $\sigma>0$ and $\mathfrak{e}_2\in[0,1)$;
the constants above satisfy $C_{\bar{\kappa}}<\infty$, $e_2\in(0,1]$ and $\frac52\,C'_{\bar{\kappa}}\cdot e_2^{\sigma}\le\frac{b(f)}{4}$.
Let $e\in(0,e_2]$.
Then $N\mapsto m_{0,N}(e)$ is a mass assignment, which does not depend on $T$;
$N_0(e)$ and $T^{*}(e)+1$ are integers;
and the number $\frac{R^{*}(e)}{L}>0$ satisfies $\frac{R^{*}(e)}{L\cdot\eta}\in\mathbb{Z}$ for $N\ge N_0(e)$.
For $N\ge N_0(e)$ and $\mathrm{ix}\in\{0,+,-\}$ let one placement $v(N,e,\mathrm{ix})$ of $\beta_{\mathrm{ix}}$ be fixed, and let $c_m$ and $\mathsf{u}_m$, $m=2,4$, be formed with it.
Then $c_m$ is real, $\mathsf{u}_m\ge0$ and $\boldsymbol{\sigma}^2\ge0$;
the test $g_N:=g^{(\eta)}_{\mathrm{ix}}$ is real, lies in $C_c^\infty(\mathbb{R}^3;\Lambda^2)$ and is supported in the cube $[-r_g,r_g]^3$, where
\[
r_g=|y_{\beta_{\mathrm{ix}}}|_\infty+\frac{D_{\beta_{\mathrm{ix}}}}{2}
\]
does not depend on $N$;
and $\|g_N-g_{\mathrm{ix}}\|_2\to0$ and $\boldsymbol{\sigma}^2\to\|g_{\mathrm{ix}}\|_2^2$ as $N\to\infty$.
Finally, let $N\ge N_0(e)$ and $T\ge T^{*}(e)+1$.
Then
\[
\left|\kappa_4\left(X_{\eta,R}(g_{\mathrm{ix}})\right)-e^4\cdot b(g_{\mathrm{ix}})\right|\le C_{\bar{\kappa}}\cdot\left(e^{4+\sigma}+\eta^{\frac14}\right);
\]
for every integer $J\ge1$, all pairwise distinct points $y_1,\dots,y_J\in\frac{R^{*}(e)}{L}\,\mathbb{Z}^3$ with $R>2\max_i|y_i|_\infty+2r_g$, and $m=2,4$,
\[
\left|E^{\mathrm{af}}_n\prod_{i=1}^{J}\mathrm{Pw}_m\left(X_{\eta,R}\left(\tau_{y_i}g_N\right);\boldsymbol{\sigma}^2,c_m\right)\right|\le\mathsf{u}_m^{J};
\]
and
\begin{gather*}
\mathsf{u}_4\le C_{\bar{\kappa}}\cdot e^{\sigma},\qquad\left(1+|c_{2}|\right)\cdot\mathsf{u}_2\le C_{\bar{\kappa}}\cdot e^{\sigma},\\
|c_{2}|\le C_{\bar{\kappa}}\cdot e^{-\mathfrak{e}_2},\qquad\left|c_{4}-3c_{2}^2-b(g_{\mathrm{ix}})\right|\le C_{\bar{\kappa}}\cdot\left(e^{\sigma}+\eta^{\frac14}\right).
\end{gather*}

\subsubsection*{2.2. Proof}

\textbf{Step 0 (the data, the constants, the objects).}

\emph{(0.0) Data exist.}
By \cite[Theorem~21.1(a)]{QEDI} there are real two-forms $\beta\in C_c^\infty(\mathbb{R}^3;\Lambda^2)$, supported in a ball, with $b\left(d\,d^{*}\beta\right)>0$.
For such a $\beta$, let $r_\beta$ be the radius of such a ball.
Consider the inequality on $x_0$ in the data of Proposition~43.2, which is the inequality of \cite[Section~2.8]{QEDI}:
\begin{multline*}
\frac{6}{\bar{m}^2}\cdot\|a\|_\infty^4\cdot\left(\frac{4\pi\cdot r_\beta^3}{3}\right)^2\cdot\left(\bar{m}^2+\left(\bar{m}+\frac{1}{d_0}\right)^2\right)\\
\cdot\frac{\exp\left(-2\bar{m}\cdot d_0\right)}{4\pi^2\cdot d_0^2}\le\frac{b(f)}{12}.
\end{multline*}
Its left side is a continuous function of $d_0>0$ that tends to $0$ as $d_0\to\infty$, and its right side is positive.
Hence the inequality holds for all $x_0$ with $|x_0|_2$ larger than some number.

\emph{(0.1) The three probes.}
The two-form $\beta$ is real, smooth and non-zero, and it is supported in a ball of radius $r_\beta$;
the two-form $\tau_{x_0}\beta$ is supported in the translate of this ball by $x_0$.
Any point of the first ball and any point of the second have Euclidean distance at least $|x_0|_2-2r_\beta=d_0>0$.
Therefore the two supports are disjoint, and $\beta_\pm=\beta\pm\tau_{x_0}\beta$ coincides with $\beta$ on $\operatorname{supp}\beta$; in particular $\beta_\pm\ne0$.
Moreover $\beta_\pm$ is real, smooth, antisymmetric in its indices and compactly supported.
Hence $\beta_0$, $\beta_+$, $\beta_-$ are probes (Section~I.7.6).
The operators $d$ and $d^{*}$ have constant coefficients, so that $d\,d^{*}\tau_{x_0}\beta=\tau_{x_0}d\,d^{*}\beta$;
thus $g_{\mathrm{ix}}=d\,d^{*}\beta_{\mathrm{ix}}$ equals $f$ for $\mathrm{ix}=0$, and $f\pm\tau_{x_0}f$ for $\mathrm{ix}=\pm$.
For each of the three probes one centre $y_{\beta_{\mathrm{ix}}}$ of $Q_{\beta_{\mathrm{ix}}}$ is fixed (Section~I.7.6).

\emph{(0.2) The constants are finite, and each is formed with earlier ones.}
The forms $\beta_{\mathrm{ix}}$ and $g_{\mathrm{ix}}$ lie in $C_c^\infty$; hence $s_{\mathrm{ix}}<\infty$ and $s_{B}=\max_{\mathrm{ix}}s_{\mathrm{ix}}<\infty$.
By Proposition~43.1, for every probe $\beta$ the numbers $e_{\sharp}(\beta)>0$ and $C^{G}_{\beta}<\infty$ exist and do not depend on $e$, and (R6) does not involve $e$;
in particular $C_{3\mathrm{p}}<\infty$.
The numbers $a_{L}=\frac{L^3-1}{(L-1)^3}$, $N_{*}=8L^3$ and $K_{L}$ are constants of the line (C4) of Section~I.8.2 and functions of $L$ (Section~0);
the number $C_{\mathrm{mat}}$ is real and formed with $L$ and $\rho$, by (0.M);
and $(L-1)^3\cdot N_{*}>0$, because $L\ge5$ (Section~I.0.5).
Hence $C_{\Gamma}^{(16)}$ is finite.
The number $C_{U}$ is the largest of four polynomials in $C_{\Gamma}^{(16)}$, $C_{3\mathrm{p}}$, $s_{B}$, $N_{*}$;
$C_{\bar{\kappa}}$ is a polynomial in $C_{U}$, and $C'_{\bar{\kappa}}$ a polynomial in $C_{\bar{\kappa}}$.
None of these numbers contains $e$, $N$, $T$ or a placement.
Since $\mathrm{Cb}\ge1$ (Section~G5.1.3) and $h\ge1$ (by (0.H)),
\[
C_{U}\ge C_{3\mathrm{p}}\ge C^{G}_{\beta}\ge C_{35}(\beta)=3\cdot10^5\,\mathrm{Cb}^4\cdot h^{\frac18}\ge3\cdot10^5.
\]
Consequently
\[
0<C_{U}\le C_{\bar{\kappa}}=16C_{U}+6C_{U}^2\le C'_{\bar{\kappa}}=8C_{\bar{\kappa}}+3C_{\bar{\kappa}}^2<\infty.
\]

\emph{(0.3) The number $e_2$.}
As $\sigma>0$ by (R6), $b(f)>0$ and $0<C'_{\bar{\kappa}}<\infty$, the condition $\frac52\,C'_{\bar{\kappa}}\cdot e_2^{\sigma}\le\frac{b(f)}{4}$ is equivalent to
\[
e_2\le\left(\frac{b(f)}{10\,C'_{\bar{\kappa}}}\right)^{\frac{1}{\sigma}}.
\]
All conditions on $e_2$ are upper bounds, and the largest number that satisfies all of them is
\[
e_2=\min\left\{e_G,\ e_{\sharp}(\beta_0),\ e_{\sharp}(\beta_+),\ e_{\sharp}(\beta_-),\ \left(\frac{b(f)}{10\,C'_{\bar{\kappa}}}\right)^{\frac{1}{\sigma}}\right\}.
\]
Since $e_G>0$ by (0.G) and $e_{\sharp}(\beta_{\mathrm{ix}})>0$ by Proposition~43.1, one has $e_2>0$.
From $e_2\le e_G$ and (0.G) one obtains $e_2\cdot h^{\frac12}\le e_1$; hence (R1) of Proposition~43.1, at $e=e_2$, gives $e_2\le\frac13$.

The radius $r_\beta$ enters only $d_0$ and the inequality displayed in (0.0).
The number $s_{B}$ is formed with the forms $\beta_{\mathrm{ix}}$ and $g_{\mathrm{ix}}$, that is, with $(\beta,x_0)$, by (0.1);
$C^{G}_{\beta_{\mathrm{ix}}}$ and $e_{\sharp}(\beta_{\mathrm{ix}})$ depend on constants of the lines (C0)--(C9), on $\bar{m}$ and on $\beta_{\mathrm{ix}}$ (Proposition~43.1);
$C_{\Gamma}^{(16)}$ is formed with $L$ and $\rho$, by (0.2);
$e_G$ is formed with $e_1$ and $h$, by (0.G);
$b(f)$ is formed with $(\bar{m},\beta)$ (Section~I.7.6);
and $\sigma$ is a constant of the line (C7).
Therefore none of the numbers $s_{B}$, $C_{3\mathrm{p}}$, $C_{\Gamma}^{(16)}$, $C_{U}$, $C_{\bar{\kappa}}$, $C'_{\bar{\kappa}}$, $e_2$ contains $r_\beta$.

From now on $e\in(0,e_2]$, $N\ge N_0(e)$, $T\ge T^{*}(e)+1$ and $\mathrm{ix}\in\{0,+,-\}$ are fixed, and
\[
\eta=L^{-N},\qquad n=8h\cdot L^{T+N},\qquad R=8h\cdot L^T,\qquad q=e\cdot\eta^{\frac12},\qquad\boldsymbol{\mu}=\boldsymbol{\mu}_N.
\]

\emph{(0.4) What Proposition~43.1 yields.}
From $e\le e_2$ one has $e\le e_G$, hence $e\cdot h^{\frac12}\le e_1$ by (0.G), and $e\le\min\left(e_G,e_{\sharp}(\beta_{\mathrm{ix}})\right)$.
Therefore Proposition~43.1 holds for this $e$ and for the probe $\beta_{\mathrm{ix}}$, with $C^{G}_{\beta_{\mathrm{ix}}}\le C_{3\mathrm{p}}$:
\begin{itemize}
\item[(In-1)] $e\le\frac13$; $T^{*}(e)$ is an integer with $T^{*}(e)\ge1$; and $R^{*}(e)=h\cdot L^{T^{*}(e)}\ge1$ (by (R1) and (0.H));
\item[(In-2)] $\boldsymbol{\mu}_N$ is real (by (R0)), and $D^{\mathrm{mat}}\le C_{\mathrm{mat}}$ (by (R2));
\item[(In-3)] $\beta_{\mathrm{ix}}$ fits $\left(N,T^{*}(e)\right)$, that is, $D_{\beta_{\mathrm{ix}}}+3\eta\le\ell_c=R^{*}(e)/(2L)$ (by (R3));
\item[(In-4)] for every placement $v$ of $\beta_{\mathrm{ix}}$,
$\Theta^{(0)}_4\le C_{3\mathrm{p}}\cdot e^{\mathfrak{e}}$ and $\left(1+|c_{2}|\right)\cdot\Theta^{(0)}_2\le C_{3\mathrm{p}}\cdot e^{\mathfrak{e}}$ (by (R3));
\item[(In-5)] $c_{2}$ and $c_{4}$ are real, they do not depend on the placement, and
$|c_{2}|\le C_{3\mathrm{p}}\cdot e^{-\mathfrak{e}_2}$, $|c_{4}|\le C_{3\mathrm{p}}^2\cdot e^{-2\mathfrak{e}_2}$
(by (R4), since each of the two non-negative terms $|c_{2}|$ and $|c_{4}|^{\frac12}$ is at most their sum);
\item[(In-6)] $c_{4}-3c_{2}^2=b_{\mathrm{ch}}\left(\eta,R^{*}(e)\right)$ and $\left|c_{4}-3c_{2}^2-b(g_{\mathrm{ix}})\right|\le C_{3\mathrm{p}}\cdot\left(e+\eta^{\frac14}\right)$ (by (R5), with $g_{\mathrm{ix}}=d\,d^{*}\beta_{\mathrm{ix}}$);
\item[(In-7)] $0<\mathfrak{e}\le\frac14$, $0\le\mathfrak{e}_2<1$, and $0<\sigma=\mathfrak{e}-\frac{2a_{L}}{N_{*}}\le\frac14$ (by (R6)).
\end{itemize}
Placements of $\beta_{\mathrm{ix}}$ exist, because every point of $\mathbb{R}^3$ has sup distance at most $\eta/2$ from $\eta\mathbb{Z}^3$.
Let $v$ be one of them.

\emph{(0.5) The torus.}
One has
\begin{gather*}
n=8h\cdot L^{T+N}=8L^{T-T^{*}(e)}\cdot h\cdot L^{T^{*}(e)+N}=\mathrm{kt}\cdot n_*,\\
\mathrm{kt}=8L^{T-T^{*}(e)}\ge8L,
\end{gather*}
where $\mathrm{kt}$ is an integer (Proposition~47.2\,(a)).
Hence $n=\left(\mathrm{kt}\cdot L\right)\cdot\mathrm{per}$ lies in $\mathrm{per}\,\mathbb{Z}$, and $n\ge L\cdot n_*$, $n\ge4$;
so $n$ is one of the sides admitted in Proposition~42.2.
Moreover $R=n\cdot\eta\ge8L\cdot R^{*}(e)$, and by (In-3)
\[
D_{\beta_{\mathrm{ix}}}<\ell_c<R^{*}(e)<R.
\]
Therefore $(\eta,R)$ and $\left(\eta,R^{*}(e)\right)$ are data of grade~0 of Proposition~47.1 for the probe $\beta_{\mathrm{ix}}$
(one has $0<\eta\le1$, the side exceeds $D_{\beta_{\mathrm{ix}}}$, and its quotient by $\eta$ is an integer $\ge2$),
and $(e,N,T)$ are data of Proposition~47.2.

\emph{(0.6) The objects.}
The objects $\hat{\mathfrak{a}}$, $\boldsymbol{\sigma}^2$, $X=\langle d\hat{\mathfrak{a}},F^{\mathrm{u}}\rangle$, $\hat{\mathfrak{a}}_{\mathfrak{c}}$, $X_{\mathfrak{c}}$, $H_{2}$, $H_{4}$ of Section~I.7.6 are formed with $\tau_v\beta_{\mathrm{ix}}$,
and $g_{\mathrm{ix}}^{(\eta)}$ is the two-form $f^{(\eta)}$ of Proposition~47.1 for the probe $\beta_{\mathrm{ix}}$.
Here $\boldsymbol{\sigma}^2=\|d\hat{\mathfrak{a}}\|^2$ is the number so denoted in Section~I.7.6, and $\boldsymbol{\sigma}=\|d\hat{\mathfrak{a}}\|\ge0$ is its non-negative square root.
By Proposition~47.1\,(b5$'$), for $v\in\eta\mathbb{Z}^3$ the number $\boldsymbol{\sigma}^2=\|d\hat{\mathfrak{a}}\|^2$ depends neither on $v$ nor on the side of the torus on which $\hat{\mathfrak{a}}$ is formed.
Thus one and the same number $\boldsymbol{\sigma}^2$ belongs to the torus $n_*$ and to the torus $n$;
it is a function of $(\eta,\beta_{\mathrm{ix}})$, that is, of $(N,\beta_{\mathrm{ix}})$.
By Proposition~47.1\,(b5) on $\Lambda_{\eta,R}$,
\begin{equation}
\boldsymbol{\sigma}\le\left\|g_{\mathrm{ix}}^{(\eta)}\right\|_2\le\|g_{\mathrm{ix}}\|_2+\|\beta_{\mathrm{ix}}\|_{H^4}=s_{\mathrm{ix}}\le s_{B}.
\tag{0.a}
\end{equation}
The data of Proposition~42.2 are at hand:
$0<e\le e_G$;
$N\ge N_0(e)$;
the real number $\boldsymbol{\mu}=\boldsymbol{\mu}_N$;
the probe $\beta_{\mathrm{ix}}$, which fits by (In-3);
the placement $v$;
the real numbers $c_{2}$, $c_{4}$ of (In-5);
and the side $n$ of (0.5).
With these data, for $m=2,4$,
\begin{gather*}
\mathrm{Mw}_4=600N_{*}^2\cdot\boldsymbol{\sigma}^4\cdot e^{-4}+|c_{4}|,\qquad\mathrm{Mw}_2=12N_{*}\cdot\boldsymbol{\sigma}^2\cdot e^{-2}+|c_{2}|,\\
\mathrm{Tw}_m=\Theta^{(0)}_m+2e^6\cdot\mathrm{Mw}_m,\qquad\mathsf{u}_m=\exp(\delta)\cdot\mathrm{Tw}_m\ge0,
\end{gather*}
where $D^{\mathrm{mat}}$, $\delta$ and $\mathring{\Phi}$ are formed with $\boldsymbol{\mu}_N$, and $\Theta^{(0)}_m$ is formed with this $\mathring{\Phi}$ and with the numbers $c_{m}=N_0^{(m)}(0)$ of (In-5).
So $D^{\mathrm{mat}}$ and $\Theta^{(0)}_m$ are the numbers of Section~I.7.6, that is, those of (In-2) and (In-4).

\emph{The placement in $\mathsf{u}_m$.}
The number $\mathsf{u}_m$ formed here is the number $\mathsf{u}_m(N,e,\beta_{\mathrm{ix}};v)$ of Proposition~43.2.
The numbers $\boldsymbol{\sigma}^2$ and $c_{m}$ do not depend on the placement (Proposition~47.1\,(b5$'$) and (In-5)), and $\delta$ contains no probe (Section~I.7.6).
Hence $\mathsf{u}_m$ depends on $v$ through $\Theta^{(0)}_m$ only.
A second placement $v'$, for which $v'-v\in\eta\mathbb{Z}^3$ by Section~I.7.6, replaces $\hat{\mathfrak{a}}_{\mathfrak{c}}=\left(\mathsf{g}_{\mathfrak{c}}\right)_*\hat{\mathfrak{a}}$ by its translate by $r_{\mathfrak{c}}(v'-v)/\eta$,
where $r_{\mathfrak{c}}=\operatorname{diag}\left((-1)^{\mathfrak{c}_\mu}\right)$ is the linear part of $\mathsf{g}_{\mathfrak{c}}$ (Sections~I.6.4 and~I.7.6).
The direction of the shift varies with the cell;
so in general the second family $\left(\hat{\mathfrak{a}}_{\mathfrak{c}}\right)_{\mathfrak{c}}$ is not a lattice translate of the first one, and $\Theta^{(0)}_m$ may depend on $v$.
The estimates below are established for the fixed placement $v$; they do not require that $\Theta^{(0)}_m$ or $\mathsf{u}_m$ take the same value for two placements.

\textbf{Step 1 ($\exp(\delta)$).}
By (2.3) of Proposition~42.2,
\[
\delta\le\frac{a_{L}}{N_{*}}\cdot\log\frac{2\pi}{e^2\cdot R^{*}(e)}+\frac{K_{L}+D^{\mathrm{mat}}+1}{(L-1)^3\cdot N_{*}}.
\]
As $R^{*}(e)\ge1$ by (In-1), one has $\log\frac{2\pi}{e^2\cdot R^{*}(e)}\le\log(2\pi)+2\log\frac1e$;
moreover $a_{L}/N_{*}>0$, and $D^{\mathrm{mat}}\le C_{\mathrm{mat}}$ by (In-2).
Therefore
\begin{equation}
\begin{aligned}
\exp(\delta)
&\le(2\pi)^{\frac{a_{L}}{N_{*}}}\cdot e^{-\frac{2a_{L}}{N_{*}}}\cdot\exp\left(\frac{K_{L}+C_{\mathrm{mat}}+1}{(L-1)^3\cdot N_{*}}\right)\\
&=C_{\Gamma}^{(16)}\cdot e^{-\frac{2a_{L}}{N_{*}}}.
\end{aligned}
\tag{1.a}
\end{equation}

\textbf{Step 2 ($\mathrm{Tw}_4$ and $\mathrm{Tw}_2$).}
By (In-1), $0<e\le1$; so a power of $e$ does not decrease when its exponent is replaced by a number that is not larger.

For $m=4$, (In-4), (0.a) and (In-5) give
\begin{align*}
\mathrm{Tw}_4
&=\Theta^{(0)}_4+1200N_{*}^2\cdot\boldsymbol{\sigma}^4\cdot e^2+2e^6\cdot|c_{4}|\\
&\le C_{3\mathrm{p}}\cdot e^{\mathfrak{e}}+1200N_{*}^2\cdot s_{B}^4\cdot e^2+2C_{3\mathrm{p}}^2\cdot e^{6-2\mathfrak{e}_2}.
\end{align*}
By (In-7), $2\ge\mathfrak{e}$ and $6-2\mathfrak{e}_2>4\ge\mathfrak{e}$.
Hence
\begin{equation}
\mathrm{Tw}_4\le\left(C_{3\mathrm{p}}+1200N_{*}^2\cdot s_{B}^4+2C_{3\mathrm{p}}^2\right)\cdot e^{\mathfrak{e}}.
\tag{2.a}
\end{equation}

For $m=2$, (In-4), (0.a) and (In-5) give
\begin{align*}
\left(1+|c_{2}|\right)\cdot\mathrm{Tw}_2
&=\left(1+|c_{2}|\right)\cdot\Theta^{(0)}_2+2e^6\left(12N_{*}\cdot\boldsymbol{\sigma}^2\cdot e^{-2}+|c_{2}|\right)\left(1+|c_{2}|\right)\\
&\le C_{3\mathrm{p}}\cdot e^{\mathfrak{e}}+2e^6\left(12N_{*}\cdot s_{B}^2\cdot e^{-2}+C_{3\mathrm{p}}\cdot e^{-\mathfrak{e}_2}\right)\left(1+C_{3\mathrm{p}}\cdot e^{-\mathfrak{e}_2}\right).
\end{align*}
Since $\mathfrak{e}_2<1$, one has $e^{-\mathfrak{e}_2}\le e^{-1}\le e^{-2}$, and $1\le e^{-1}$.
Hence the first bracket is at most $\left(12N_{*}\cdot s_{B}^2+C_{3\mathrm{p}}\right)e^{-2}$, the second is at most $\left(1+C_{3\mathrm{p}}\right)e^{-1}$, and $2e^6\cdot e^{-2}\cdot e^{-1}=2e^3\le2e^{\mathfrak{e}}$.
It follows that
\begin{equation}
\left(1+|c_{2}|\right)\cdot\mathrm{Tw}_2\le\left(C_{3\mathrm{p}}+2\left(12N_{*}\cdot s_{B}^2+C_{3\mathrm{p}}\right)\left(1+C_{3\mathrm{p}}\right)\right)\cdot e^{\mathfrak{e}}.
\tag{2.b}
\end{equation}

\textbf{Step 3 ((U-c)).}
Multiply (2.a) and (2.b) by (1.a), and use $\mathfrak{e}-\frac{2a_{L}}{N_{*}}=\sigma$.
The constants that arise are the first two entries of the maximum that defines $C_{U}$:
\begin{align*}
\mathsf{u}_4=\exp(\delta)\cdot\mathrm{Tw}_4
&\le C_{\Gamma}^{(16)}\cdot\left(C_{3\mathrm{p}}+1200N_{*}^2\cdot s_{B}^4+2C_{3\mathrm{p}}^2\right)\cdot e^{\sigma}\le C_{U}\cdot e^{\sigma},\\
\left(1+|c_{2}|\right)\cdot\mathsf{u}_2=\exp(\delta)\cdot\left(1+|c_{2}|\right)\cdot\mathrm{Tw}_2
&\le C_{\Gamma}^{(16)}\cdot\left(C_{3\mathrm{p}}+2\left(12N_{*}\cdot s_{B}^2+C_{3\mathrm{p}}\right)\left(1+C_{3\mathrm{p}}\right)\right)\cdot e^{\sigma}\\
&\le C_{U}\cdot e^{\sigma}.
\end{align*}
By (In-5) and $C_{3\mathrm{p}}\le C_{U}$, one has $|c_{2}|\le C_{U}\cdot e^{-\mathfrak{e}_2}$.
By (In-6), by $e\le e^{\sigma}$ (which holds as $e\le1$ and $\sigma\le1$) and by $C_{3\mathrm{p}}\le C_{U}$,
\[
\left|c_{4}-3c_{2}^2-b(g_{\mathrm{ix}})\right|\le C_{U}\cdot\left(e^{\sigma}+\eta^{\frac14}\right).
\]
These are the four inequalities of (U-c).

\textbf{Step 4 (the cumulants of $X$).}
For every real $\boldsymbol{\mu}$, $E^{\mathrm{af}}_n$ is the expectation of a probability law of $F^{\mathrm{u}}$ (Section~I.7.2);
under it every variable $\langle\mathrm{gc},F^{\mathrm{u}}\rangle$ with $\mathrm{gc}\in C^2(n)$ has finite moments of all orders (see (5.2) below).
Let $\kappa_2(X)$ and $\kappa_4(X)$ be the second and the fourth cumulant of $X=\langle d\hat{\mathfrak{a}},F^{\mathrm{u}}\rangle$ under $E^{\mathrm{af}}_n$.
Multiplying (2.1) and (2.2) of Proposition~42.2 by $e^2$ and by $e^4$, respectively, one obtains
\begin{equation}
\begin{aligned}
&\left|\kappa_2(X)-\boldsymbol{\sigma}^2-e^2\cdot c_{2}\right|\le e^2\cdot\mathsf{u}_2,\\
&\left|\kappa_4(X)-e^4\left(c_{4}-3c_{2}^2\right)\right|\le e^4\left[\mathsf{u}_4+6\,\mathsf{u}_2\cdot|c_{2}|+3\,\mathsf{u}_2^2\right].
\end{aligned}
\tag{4.a}
\end{equation}
By Step~3,
$\mathsf{u}_4\le C_{U}\cdot e^{\sigma}$;
$\mathsf{u}_2\cdot|c_{2}|\le\left(1+|c_{2}|\right)\cdot\mathsf{u}_2\le C_{U}\cdot e^{\sigma}$;
and $\mathsf{u}_2\le\left(1+|c_{2}|\right)\cdot\mathsf{u}_2\le C_{U}\cdot e^{\sigma}$, whence $\mathsf{u}_2^2\le C_{U}^2\cdot e^{2\sigma}\le C_{U}^2\cdot e^{\sigma}$.
Consequently
\begin{equation}
\left|\kappa_4(X)-e^4\left(c_{4}-3c_{2}^2\right)\right|\le\left(7C_{U}+3C_{U}^2\right)\cdot e^{4+\sigma}.
\tag{4.b}
\end{equation}


\textbf{Step 5 (from $X$ to $X_{\eta,R}(g_{\mathrm{ix}})$).}

\emph{(5.1) $X$ as a test variable; removal of the placement.}
On $\Lambda_{\eta,R}$, Proposition~47.1\,(b2$'$) gives
\[
X=X_{\eta,R}\left(\tau_v g_{\mathrm{ix}}^{(\eta)}\right)
\]
as an identity between functions of $F^{\mathrm{u}}\in C^{2}(n)$.
The set $\{0\}$ is a subset of $\left(R^{*}(e)/L\right)\mathbb{Z}^3$,
and Proposition~47.2\,(d), applied with $\mathrm{Ypt}=\{0\}$ and with the polynomials $\mathrm{Pol}(z)=z^j$, yields
\[
E^{\mathrm{af}}_n\left[X_{\eta,R}\left(\tau_v g_{\mathrm{ix}}^{(\eta)}\right)^j\right]
=E^{\mathrm{af}}_n\left[X_{\eta,R}\left(g_{\mathrm{ix}}^{(\eta)}\right)^j\right],
\qquad j=1,2,3,4.
\]
A cumulant of order at most four of a random variable is a polynomial in the moments of order at most four of that variable.
Hence
\begin{equation}\tag{5.a}
\kappa_m(X)=\kappa_m\left(X_{\eta,R}\left(g_{\mathrm{ix}}^{(\eta)}\right)\right)\qquad(m=2,4).
\end{equation}

\emph{(5.2) Multilinear bounds.}
The expectation $E^{\mathrm{af}}_n$ is the expectation $\mathbb{E}_{C(n,n,n)}$ of \cite[Definition~2.13]{QEDII},
formed on the torus $C(n,n,n)$ with the mass parameter $M_{\mathrm{ch}}=\boldsymbol{\mu}_N+\frac{\rho}{2}\cdot\mathrm{sw}$,
and the curvature $F$ of that definition is $F^{\mathrm{u}}$ (Sections~I.7.2 and~I.9.3; Proposition~46\,(d), (f)).
The assertions of \cite[Theorem~6.27]{QEDII} hold on every torus $C(n)$ with $n_\mu\ge4$, for every real mass parameter and every $\rho$ with $-1\le\rho\le1$.
These hypotheses are fulfilled here:
$n\ge4$ by (0.5);
$M_{\mathrm{ch}}$ is real, because $\boldsymbol{\mu}_N$ is real by (In-2);
and $0<\rho<1$ by Section~I.0.5.
Let $\mathrm{gc},\mathrm{gc}_1,\dots,\mathrm{gc}_4\in C^{2}(n)$ be real two-cochains.
By \cite[Theorem~6.27(b)]{QEDII}, every moment of $\langle\mathrm{gc},F^{\mathrm{u}}\rangle$ is finite, and
\begin{gather*}
E^{\mathrm{af}}_n\langle\mathrm{gc},F^{\mathrm{u}}\rangle=0,\qquad
E^{\mathrm{af}}_n\langle\mathrm{gc},F^{\mathrm{u}}\rangle^2\le\|\mathrm{gc}\|^2,\\
\left|\kappa_4\left(\langle\mathrm{gc}_1,F^{\mathrm{u}}\rangle,\dots,\langle\mathrm{gc}_4,F^{\mathrm{u}}\rangle\right)\right|
\le15\,\|\mathrm{gc}_1\|\cdot\|\mathrm{gc}_2\|\cdot\|\mathrm{gc}_3\|\cdot\|\mathrm{gc}_4\|.
\end{gather*}
Here, for centred random variables $\mathrm{Rv}_1,\dots,\mathrm{Rv}_4$ with finite fourth moments, the joint cumulants are
\begin{align*}
\kappa_2(\mathrm{Rv}_1,\mathrm{Rv}_2)&=E^{\mathrm{af}}_n\left[\mathrm{Rv}_1\cdot\mathrm{Rv}_2\right],\\
\kappa_4(\mathrm{Rv}_1,\dots,\mathrm{Rv}_4)&=E^{\mathrm{af}}_n\left[\mathrm{Rv}_1\cdot\mathrm{Rv}_2\cdot\mathrm{Rv}_3\cdot\mathrm{Rv}_4\right]
-E^{\mathrm{af}}_n\left[\mathrm{Rv}_1\cdot\mathrm{Rv}_2\right]\cdot E^{\mathrm{af}}_n\left[\mathrm{Rv}_3\cdot\mathrm{Rv}_4\right]\\
&\quad-E^{\mathrm{af}}_n\left[\mathrm{Rv}_1\cdot\mathrm{Rv}_3\right]\cdot E^{\mathrm{af}}_n\left[\mathrm{Rv}_2\cdot\mathrm{Rv}_4\right]
-E^{\mathrm{af}}_n\left[\mathrm{Rv}_1\cdot\mathrm{Rv}_4\right]\cdot E^{\mathrm{af}}_n\left[\mathrm{Rv}_2\cdot\mathrm{Rv}_3\right],
\end{align*}
and $\kappa_m(\mathrm{Rv}_1)=\kappa_m(\mathrm{Rv}_1,\dots,\mathrm{Rv}_1)$.
Define two forms on $C^{2}(n)$ by
\begin{align*}
\mathrm{kq}_2(\mathrm{gc}_1,\mathrm{gc}_2)&:=\kappa_2\left(\langle\mathrm{gc}_1,F^{\mathrm{u}}\rangle,\langle\mathrm{gc}_2,F^{\mathrm{u}}\rangle\right),\\
\mathrm{kq}_4(\mathrm{gc}_1,\dots,\mathrm{gc}_4)&:=\kappa_4\left(\langle\mathrm{gc}_1,F^{\mathrm{u}}\rangle,\dots,\langle\mathrm{gc}_4,F^{\mathrm{u}}\rangle\right).
\end{align*}
Both forms are symmetric and linear in each argument.
The second of the bounds above, together with the Cauchy--Schwarz inequality, gives
$\left|\mathrm{kq}_2(\mathrm{gc}_1,\mathrm{gc}_2)\right|\le\|\mathrm{gc}_1\|\cdot\|\mathrm{gc}_2\|$.

\emph{(5.3) The two cochains.}
Let $f'\in\mathcal{S}$ be a real test.
By Section~I.7.5,
\[
X_{\eta,R}(f')=\langle f'_\eta,F\rangle_\eta=\left\langle(f'_\eta)^{\mathrm{u}},F^{\mathrm{u}}\right\rangle,\qquad
(f'_\eta)^{\mathrm{u}}=\eta^{\frac32}\cdot f'_\eta\in C^{2}(n),\qquad
\left\|(f'_\eta)^{\mathrm{u}}\right\|=|f'_\eta|_\eta,
\]
and the map $f'\mapsto f'_\eta$ is linear.
Put
\[
\mathrm{gc}:=\left((g_{\mathrm{ix}})_\eta\right)^{\mathrm{u}},\qquad
\mathrm{gc}':=\left(\left(g_{\mathrm{ix}}^{(\eta)}\right)_\eta\right)^{\mathrm{u}}.
\]
Then
\begin{gather*}
\mathrm{gc}-\mathrm{gc}'=\left(\left(g_{\mathrm{ix}}-g_{\mathrm{ix}}^{(\eta)}\right)_\eta\right)^{\mathrm{u}},\\
\kappa_m\left(X_{\eta,R}(g_{\mathrm{ix}})\right)=\mathrm{kq}_m(\mathrm{gc},\dots,\mathrm{gc}),\qquad
\kappa_m\left(X_{\eta,R}\left(g_{\mathrm{ix}}^{(\eta)}\right)\right)=\mathrm{kq}_m(\mathrm{gc}',\dots,\mathrm{gc}')\qquad(m=2,4).
\end{gather*}
The support of $g_{\mathrm{ix}}$ lies in $\operatorname{supp}\beta_{\mathrm{ix}}$, which is a subset of $Q_{\beta_{\mathrm{ix}}}$;
by Proposition~47.2\,(e), the two-form $g_{\mathrm{ix}}^{(\eta)}$ is real, belongs to $C^{\infty}_{\mathrm{c}}$ and is supported in $Q_{\beta_{\mathrm{ix}}}$.
As $D_{\beta_{\mathrm{ix}}}<R$ by (0.5), the cube
\[
Q_{\beta_{\mathrm{ix}}}=y_{\beta_{\mathrm{ix}}}+\left[-\frac{D_{\beta_{\mathrm{ix}}}}{2},\frac{D_{\beta_{\mathrm{ix}}}}{2}\right]^3
\]
is contained in $y_{\beta_{\mathrm{ix}}}+(-R/2,R/2)^3$.
Consequently Proposition~47.1\,(b0), with $y_0=y_{\beta_{\mathrm{ix}}}$, applies to each of the three tests
$g_{\mathrm{ix}}$, $g_{\mathrm{ix}}^{(\eta)}$ and $g_{\mathrm{ix}}-g_{\mathrm{ix}}^{(\eta)}$.
Combined with (0.a) and with Proposition~47.1\,(b3) at the Sobolev order $s=0$, it gives
\begin{equation}\tag{5.b}
\begin{aligned}
&\|\mathrm{gc}\|\le\|g_{\mathrm{ix}}\|_2\le s_{\mathrm{ix}},\qquad
\|\mathrm{gc}'\|\le\left\|g_{\mathrm{ix}}^{(\eta)}\right\|_2\le s_{\mathrm{ix}},\\
&\|\mathrm{gc}-\mathrm{gc}'\|\le\left\|g_{\mathrm{ix}}-g_{\mathrm{ix}}^{(\eta)}\right\|_2\le\eta^2\cdot\|\beta_{\mathrm{ix}}\|_{H^4}.
\end{aligned}
\end{equation}

\emph{(5.4) Telescoping.}
Multilinearity gives
\begin{align*}
&\mathrm{kq}_4(\mathrm{gc},\mathrm{gc},\mathrm{gc},\mathrm{gc})-\mathrm{kq}_4(\mathrm{gc}',\mathrm{gc}',\mathrm{gc}',\mathrm{gc}')\\
&\qquad=\mathrm{kq}_4(\mathrm{gc}-\mathrm{gc}',\mathrm{gc},\mathrm{gc},\mathrm{gc})+\mathrm{kq}_4(\mathrm{gc}',\mathrm{gc}-\mathrm{gc}',\mathrm{gc},\mathrm{gc})\\
&\qquad\quad+\mathrm{kq}_4(\mathrm{gc}',\mathrm{gc}',\mathrm{gc}-\mathrm{gc}',\mathrm{gc})+\mathrm{kq}_4(\mathrm{gc}',\mathrm{gc}',\mathrm{gc}',\mathrm{gc}-\mathrm{gc}'),\\
&\mathrm{kq}_2(\mathrm{gc},\mathrm{gc})-\mathrm{kq}_2(\mathrm{gc}',\mathrm{gc}')
=\mathrm{kq}_2(\mathrm{gc}-\mathrm{gc}',\mathrm{gc})+\mathrm{kq}_2(\mathrm{gc}',\mathrm{gc}-\mathrm{gc}').
\end{align*}
By the bounds of (5.2) and by (5.b),
each of the four terms on the right of the first identity has modulus at most $15\,\eta^2\cdot\|\beta_{\mathrm{ix}}\|_{H^4}\cdot s_{\mathrm{ix}}^3$,
and each of the two terms on the right of the second identity has modulus at most $\eta^2\cdot\|\beta_{\mathrm{ix}}\|_{H^4}\cdot s_{\mathrm{ix}}$.
Moreover $\|\beta_{\mathrm{ix}}\|_{H^4}\le s_{\mathrm{ix}}\le s_{B}$,
and $60\,s_{B}^4\le C_{U}$, the number $60\,s_{B}^4$ being the fourth entry of the maximum that defines $C_{U}$.
Therefore
\begin{equation}\tag{5.c}
\begin{aligned}
\left|\kappa_4\left(X_{\eta,R}(g_{\mathrm{ix}})\right)-\kappa_4\left(X_{\eta,R}\left(g_{\mathrm{ix}}^{(\eta)}\right)\right)\right|
&\le60\,\eta^2\cdot\|\beta_{\mathrm{ix}}\|_{H^4}\cdot s_{\mathrm{ix}}^3\le60\,s_{B}^4\cdot\eta^2\le C_{U}\cdot\eta^2,\\
\left|\kappa_2\left(X_{\eta,R}(g_{\mathrm{ix}})\right)-\kappa_2\left(X_{\eta,R}\left(g_{\mathrm{ix}}^{(\eta)}\right)\right)\right|
&\le2\,\eta^2\cdot\|\beta_{\mathrm{ix}}\|_{H^4}\cdot s_{\mathrm{ix}}\le2\,s_{B}^2\cdot\eta^2.
\end{aligned}
\end{equation}

\textbf{Step 6 ((U-a$'$), (U-a)).}
By (In-6), $c_{4}-3c_{2}^2=b_{\mathrm{ch}}(\eta,R^{*}(e))$, the number of Section~I.7.6 for the probe $\beta_{\mathrm{ix}}$, formed with $\tau_v\beta_{\mathrm{ix}}$;
by (In-5) it does not depend on the placement.
The hypotheses of Proposition~37.3 are satisfied for the probe $\beta_{\mathrm{ix}}$:
\begin{itemize}
\item[--] $e\cdot h^{\frac12}\le e_1$, by (0.4);
\item[--] $N\ge N_0(e)$;
\item[--] $v\in\eta\mathbb{Z}^3$, since $v$ is a placement (Section~I.7.6);
\item[--] $R^{*}(e)=2L\cdot\ell_c\ge\ell_c\ge D_{\beta_{\mathrm{ix}}}+3\eta>D_{\mathrm{be}}+1$, by (In-3) and by $D_\beta=D_{\mathrm{be}}+4$ (Section~I.7.6), where $D_{\mathrm{be}}$ and $D_\beta$ are formed with $\beta_{\mathrm{ix}}$;
\item[--] $\hat{\mathfrak{a}}=\eta^{\frac12}\cdot a_{\eta}$ on $C(n_*)$, where the lattice probe $a_{\eta}$ of Section~I.7.6, formed with $\tau_v\beta_{\mathrm{ix}}$ on $\Lambda_{\eta,R^{*}(e)}$, is the cochain $a_{\eta}[\tau_v\beta_{\mathrm{ix}}]$ of Section~G5.1.3, the two defining displays being the same;
\item[--] $c_{m}=N_0^{(m)}(0)$, with the free ratio $N_0$ of Section~I.7.6 on $C(n_*)$.
\end{itemize}
Hence Proposition~37.3\,(a) gives
$c_{4}-3c_{2}^2=b_{\mathrm{ch}}(\eta,R^{*}(e))[\beta_{\mathrm{ix}}]$,
the number of Section~G5.6.0 formed with $\beta_{\mathrm{ix}}$ itself.
Combining (4.b) with (5.a) and (5.c), and the first inequality of (4.a) with (5.a) and (5.c), one obtains
\begin{align*}
\left|\kappa_4\left(X_{\eta,R}(g_{\mathrm{ix}})\right)-e^4\cdot b_{\mathrm{ch}}(\eta,R^{*}(e))[\beta_{\mathrm{ix}}]\right|
&\le\left(7C_{U}+3C_{U}^2\right)\cdot e^{4+\sigma}+C_{U}\cdot\eta^2,\\
\left|\kappa_2\left(X_{\eta,R}(g_{\mathrm{ix}})\right)-\boldsymbol{\sigma}^2-e^2\cdot c_{2}\right|
&\le e^2\cdot\mathsf{u}_2+2s_{B}^2\cdot\eta^2.
\end{align*}
These are the two inequalities of (U-a$'$).
Adding $e^4$ times the fourth inequality of (U-c) to the first of them, one finds
\begin{align*}
\left|\kappa_4\left(X_{\eta,R}(g_{\mathrm{ix}})\right)-e^4\cdot b(g_{\mathrm{ix}})\right|
&\le\left(7C_{U}+3C_{U}^2\right)\cdot e^{4+\sigma}+C_{U}\cdot\eta^2+C_{U}\cdot e^{4+\sigma}+C_{U}\cdot e^4\cdot\eta^{\frac14}\\
&=\left(8C_{U}+3C_{U}^2\right)\cdot e^{4+\sigma}+C_{U}\cdot\left(\eta^2+e^4\cdot\eta^{\frac14}\right),
\end{align*}
and $8C_{U}+3C_{U}^2=\frac{C_{\bar{\kappa}}}{2}$.
This is the first inequality of (U-a).
For the second one, note that $\eta\le1$ implies $\eta^2\le\eta^{\frac14}$ and that $e\le1$ implies $e^4\le1$.
Therefore
\[
C_{U}\cdot\left(\eta^2+e^4\cdot\eta^{\frac14}\right)\le2C_{U}\cdot\eta^{\frac14}\le C_{\bar{\kappa}}\cdot\eta^{\frac14},
\]
because $C_{\bar{\kappa}}=16C_{U}+6C_{U}^2\ge2C_{U}$;
moreover $\frac{C_{\bar{\kappa}}}{2}\cdot e^{4+\sigma}\le C_{\bar{\kappa}}\cdot e^{4+\sigma}$.

\textbf{Step 7 ((U-b)).}
Let $m\in\{2,4\}$, and let $\mathrm{Ypt}$ be a finite subset of $\left(R^{*}(e)/L\right)\mathbb{Z}^3$ whose points are pairwise incongruent modulo $R\mathbb{Z}^3$.
(If $\mathrm{Ypt}$ is empty, both sides of the inequality of (U-b) equal $1$.)
Recall that $R^{*}(e)/L=\mathrm{per}\cdot\eta$, and put
\[
\mathfrak{c}(y):=\frac{2y}{\mathrm{per}\cdot\eta}\in2\mathbb{Z}^3\qquad(y\in\mathrm{Ypt}).
\]
By Proposition~47.2\,(c), the indices $\mathfrak{c}(y)$ and $\mathfrak{c}(y')$ denote the same cell of the torus $n$ if and only if $y-y'\in R\mathbb{Z}^3$.
So $\mathrm{Scl}:=\{\mathfrak{c}(y):y\in\mathrm{Ypt}\}$ is a set of $\#|\mathrm{Ypt}|$ distinct cells of the torus $n$.
By Proposition~47.2\,(b),
\[
X_{\mathfrak{c}(y)}=X_{\eta,R}\left(\tau_{v+y}g_{\mathrm{ix}}^{(\eta)}\right).
\]
The definitions of $H_m$ and $\mathrm{Pw}_m$ in Section~I.7.6 give
\begin{align*}
e^{-2}H_{2}(z)-c_{2}&=e^{-2}\left(z^2-\boldsymbol{\sigma}^2\right)-c_{2}=\mathrm{Pw}_2\left(z;\boldsymbol{\sigma}^2,c_{2}\right),\\
e^{-4}H_{4}(z)-c_{4}&=e^{-4}\left(z^4-6\boldsymbol{\sigma}^2\cdot z^2+3\boldsymbol{\sigma}^4\right)-c_{4}=\mathrm{Pw}_4\left(z;\boldsymbol{\sigma}^2,c_{4}\right).
\end{align*}
Accordingly, Proposition~42.2\,(2.4) for the set $\mathrm{Scl}$ takes the form
\[
\left|E^{\mathrm{af}}_n\prod_{y\in\mathrm{Ypt}}\mathrm{Pw}_m\left(X_{\eta,R}\left(\tau_{v+y}g_{\mathrm{ix}}^{(\eta)}\right);\boldsymbol{\sigma}^2,c_{m}\right)\right|
\le\mathsf{u}_m^{\#|\mathrm{Ypt}|}.
\]
The function
\[
\mathrm{Pol}\left((z_y)_{y\in\mathrm{Ypt}}\right):=\prod_{y\in\mathrm{Ypt}}\mathrm{Pw}_m\left(z_y;\boldsymbol{\sigma}^2,c_{m}\right)
\]
is a polynomial in $\#|\mathrm{Ypt}|$ variables.
By Proposition~47.2\,(d), the expectation keeps its value when $\tau_{v+y}$ is replaced by $\tau_y$ for every $y\in\mathrm{Ypt}$.
This is (U-b).

\textbf{Step 8 (the last paragraph of Proposition~43.2).}
The setting of Proposition~43.2 is an instance of the setting of \cite[Definition~2.36]{QEDII}.
Indeed, $\rho\in(0,1)$, the standard weights, $L$, which is an odd integer $\ge5$ and hence $\ge3$, and $\bar{m}\in(0,1]$ are standing data of Proposition~43.2 (Section~G6.A1; for the ranges, Section~I.0.5).
The number $L$ is the constant of the line (C3) of Section~I.8.2,
and \cite[Definition~2.36]{QEDII} admits every odd $L\ge3$, in particular this one.
The scale $L^{\lfloor\log_L\bar{m}^{-1}\rfloor}$ of \cite[Definition~2.36]{QEDII} is $h$.
One has $n=8h\cdot L^{T+N}$ and $R=n\cdot\eta$ (Section~I.6.5; \cite[Definition~2.36]{QEDII}).
The polynomials $\mathrm{Pw}_2$, $\mathrm{Pw}_4$ that occur in the last paragraph of Proposition~43.2 are those of Section~I.7.6,
and by Proposition~46\,(a) nothing depends on the point $v_0$ of \cite[Definition~2.36]{QEDII}.
The proof of the assertions of that paragraph follows the order in which the paragraph states them;
it is divided into ten parts, labelled (PU-1)--(PU-10).

\emph{(PU-1) The numbers $\sigma$, $\mathfrak{e}_2$, $C_{\bar{\kappa}}$ and $e_2$.}
By (In-7), $\sigma>0$ and $\mathfrak{e}_2\in[0,1)$.
By (0.2), $C_{\bar{\kappa}}<\infty$.
By (0.3) and by the definition of $e_2$, one has $e_2\in(0,1]$ and $\frac52\,C'_{\bar{\kappa}}\cdot e_2^{\sigma}\le\frac{b(f)}{4}$, where $C'_{\bar{\kappa}}=8C_{\bar{\kappa}}+3C_{\bar{\kappa}}^2$.

From now on let $e\in(0,e_2]$.

\emph{(PU-2) The mass assignment.}
By Proposition~43.1\,(R0), $N\mapsto m_{0,N}(e)$ is a real sequence that does not depend on $T$.
It is therefore a mass assignment in the sense of \cite[Definition~2.36]{QEDII}.

\emph{(PU-3) The numbers $N_0(e)$ and $T^{*}(e)+1$.}
The number $N_0(e)$ is an integer $\ge0$:
this is the claim (1.7) of Section~1.2, which holds for every $e>0$ with $e\cdot h^{\frac12}\le e_1$ and whose proof rests on (1.2), (1.4) and \cite[Proposition~19.1(b), (c)]{QEDI}.
By (In-1), the number $T^{*}(e)+1$ is an integer.

\emph{(PU-4) The spacing $R^{*}(e)/L$.}
By (In-1), $R^{*}(e)=h\cdot L^{T^{*}(e)}\ge1$, so that $R^{*}(e)/L>0$.
Let $N\ge N_0(e)$.
Any $T\ge T^{*}(e)+1$, any $\mathrm{ix}\in\{0,+,-\}$ and any placement of $\beta_{\mathrm{ix}}$ form, together with $e$ and $N$, data of Proposition~47.2, by (0.4) and (0.5);
clause (e) of that proposition shows that $\left(R^{*}(e)/L\right)/\eta=\mathrm{per}$ is an integer, a number in which $T$, $\mathrm{ix}$ and the placement do not occur.

\emph{(PU-5) The placement $v(N,e,\mathrm{ix})$ of the last paragraph of Proposition~43.2, and the numbers $c_{m}$, $\mathsf{u}_m$, $\boldsymbol{\sigma}^2$.}
Let $N\ge N_0(e)$ and $\mathrm{ix}\in\{0,+,-\}$.
The set of the placements of $\beta_{\mathrm{ix}}$ is finite and not empty, because in each coordinate $v_\mu$ lies in $\eta\mathbb{Z}$ and in a closed interval of length $\eta$.
Let $v=v(N,e,\mathrm{ix})$ be an arbitrary element of this set; one may take, for instance, its first element in the lexicographic order.
Since (In-5) and (0.6) hold for every placement of $\beta_{\mathrm{ix}}$, for this placement $c_{m}$ is real by (In-5), $\mathsf{u}_m\ge0$ by (0.6), and $\boldsymbol{\sigma}^2\ge0$.
None of these three numbers depends on $T$:
$c_{m}$, $\Theta^{(0)}_m$ and $\delta$ are formed on the tori $n_*$ and $L\cdot n_*$ (Section~I.7.6),
and $\boldsymbol{\sigma}^2$ does not depend on the torus, by (0.6).
By (0.6), $\boldsymbol{\sigma}^2$ is a function of $(N,\beta_{\mathrm{ix}})$;
it contains neither $e$ nor $T$, although the choice of the placement $v(N,e,\mathrm{ix})$ may depend on $e$.
The placement $v$ of (0.4) is an arbitrary placement of $\beta_{\mathrm{ix}}$;
hence (U-a)--(U-c), which Steps~1--7 establish for that placement, hold for $v=v(N,e,\mathrm{ix})$, and so do (PU-8)--(PU-10) below, which are drawn from them.

\emph{(PU-6) The forms $g_N$.}
By Proposition~47.2\,(e), the form $g_N:=g_{\mathrm{ix}}^{(\eta)}$ is real, lies in $C^{\infty}_{\mathrm{c}}(\mathbb{R}^3;\Lambda^2)$ and is supported in $Q_{\beta_{\mathrm{ix}}}$,
which is contained in the cube $[-r_g,r_g]^3$, $r_g=|y_{\beta_{\mathrm{ix}}}|_\infty+D_{\beta_{\mathrm{ix}}}/2$, a cube that does not depend on $N$;
moreover $\|g_N-g_{\mathrm{ix}}\|_2\le\eta^2\cdot\|\beta_{\mathrm{ix}}\|_{H^4}\to0$ and $\boldsymbol{\sigma}^2\to\|g_{\mathrm{ix}}\|_2^2$ as $N\to\infty$.

From now on let, in addition, $N\ge N_0(e)$ and $T\ge T^{*}(e)+1$.

\emph{(PU-7) The law $P_{N,T}$.}
$P_{N,T}$ is the law $P^{\mathrm{ch}}_{N,T}$ of \cite[Definition~2.36]{QEDII} (Section~I.7.5; Proposition~46\,(d), (f)).
It is the law of $F^{\sharp}$ when $F^{\mathrm{u}}$ has the law $E^{\mathrm{af}}_n$, and $F^{\sharp}(f')=X_{\eta,R}(f')$ (Section~I.7.5).
Hence the expectation under $P^{\mathrm{ch}}_{N,T}$ of a polynomial in finitely many of the variables $X_{\eta,R}(f')$ equals its expectation under $E^{\mathrm{af}}_n$.

\emph{(PU-8) The fourth cumulant.}
The bound on the fourth cumulant in the last paragraph of Proposition~43.2,
\[
\left|\kappa_4\left(X_{\eta,R}(g_{\mathrm{ix}})\right)-e^4\cdot b(g_{\mathrm{ix}})\right|\le C_{\bar{\kappa}}\cdot\left(e^{4+\sigma}+\eta^{\frac14}\right),
\]
is the second inequality of (U-a).

\emph{(PU-9) Products at distinct points.}
The bound of the last paragraph of Proposition~43.2 on the expectation of a product of the polynomials $\mathrm{Pw}_m$ at pairwise distinct points follows from (U-b).
Let $m\in\{2,4\}$, and let $y_1,\dots,y_j\in\left(R^{*}(e)/L\right)\mathbb{Z}^3$ be pairwise distinct, with $R>2\max_i|y_i|_\infty+2r_g$.
For $i\neq i'$ one has
\[
0<|y_i-y_{i'}|_\infty\le2\max_i|y_i|_\infty<R,
\]
so that $y_i-y_{i'}\notin R\mathbb{Z}^3$.
Hence $\mathrm{Ypt}:=\{y_1,\dots,y_j\}$ consists of $j$ points that are pairwise incongruent modulo $R\mathbb{Z}^3$,
and (U-b) gives the bound $\mathsf{u}_m^{j}$.

\emph{(PU-10) The bounds with $C_{\bar{\kappa}}$.}
The four bounds with the constant $C_{\bar{\kappa}}$ that close the last paragraph of Proposition~43.2 follow from (U-c), because $C_{U}\le C_{\bar{\kappa}}$ by (0.2).

This completes the proof of the last paragraph of Proposition~43.2.

\textbf{Step 9 (uniformity).}
The numbers $s_{B}$, $C_{3\mathrm{p}}$, $C_{\Gamma}^{(16)}$, $C_{U}$, $C_{\bar{\kappa}}$, $C'_{\bar{\kappa}}$, $e_2$ are fixed in Step~0, before $e$, $N$, $T$, $\mathrm{ix}$ and $v$.
They are functions of constants of the lines (C0)--(C9) and of $(\bar{m},\beta,x_0)$, and by (0.3) none of them contains $r_\beta$.
Steps~1--8 hold for every $e\in(0,e_2]$, every $N\ge N_0(e)$, every $T\ge T^{*}(e)+1$, every $\mathrm{ix}$ and every placement $v$, Step~8 for every choice of the placement $v(N,e,\mathrm{ix})$;
so all clauses of Proposition~43.2, its last paragraph included, hold for every placement, with $\mathsf{u}_m=\mathsf{u}_m(N,e,\beta_{\mathrm{ix}};v)$.
The right sides of (U-a), (U-a$'$), (U-b) and (U-c) do not contain $T$, and they contain $N$ only through $\eta$ and $\mathsf{u}_m$.
\qed

\subsection*{3. Conditions on the constants}

The proofs of Propositions~43.1 and~43.2 rest on the following conditions on the constants.
Each condition is followed by the places where it is imposed and by the places where it enters.

\begin{description}
\item[$0<\rho<1$.]
Imposed in Section~I.0.5; Section~I.8.2, line (C0); Proposition~14.2\,(a), row $\rho$.
The inequality $\rho>0$ enters (0.M).
The range $0<\rho<1$ enters Section~2.2, (5.2), where $\rho$ has to lie in the range of \cite[Theorem~6.27]{QEDII},
and Section~2.2, Step~8, in its opening paragraph and in (PU-7), where $\rho$ is a datum of \cite[Definition~2.36]{QEDII}.

\item[$L\ge5$.]
Imposed in Section~I.0.5; Section~I.8.2, line (C3); Proposition~14.2\,(b), (L5).
Used in Section~1.2, Step~4 ($n_*\ge4$).
The inequality $L\ge3$ enters Section~2.2, Step~8, in its opening paragraph and in (PU-7), as a datum of \cite[Definition~2.36]{QEDII}.
The inequality $L>1$ enters (0.H) and (0.M);
Section~1.2 at (1.2) ($\mathrm{lg}>0$), (1.3), (1.4), Step~2 ($a_{L}>0$) and Step~4 ($1/L<1$);
and Section~2.2 at (0.2), (In-1), (0.5), in Steps~1 and~6, and in Step~8, (PU-6), where $\eta=L^{-N}\to0$ as $N\to\infty$.

\item[$L$ odd.]
Imposed in Section~I.0.5; Section~I.8.2, line (C3); Proposition~14.2\,(b), (L5).
Used in Section~2.2, Step~8, in its opening paragraph and in (PU-7), as a datum of \cite[Definition~2.36]{QEDII}.

\item[$M$ a power of $L$, $M\ge L^2$.]
Imposed in Section~I.0.5; Section~I.8.2, line (C5); Proposition~14.2\,(b), (M1).
Used in Section~1.2 at (1.4) and (1.6).

\item[$r\ge2$.]
Imposed in Section~I.8.2, line (C7); Proposition~14.2\,(a), row $r$.
Used in Section~1.2 at (1.3); the inequality $r\ge1$ enters Section~1.2 at (1.2);
the inequality $r>0$ enters Section~1.2 at (1.5), where it gives $v_K^r\to\infty$ as $e\to0$.

\item[$\alpha<1$, $1-\alpha<\varepsilon$ (hence $\varepsilon>0$), $0<s<1$.]
Imposed in Section~I.8.2, lines (C1) and (C7); Proposition~14.2\,(a), rows $\alpha$, $\varepsilon$, $s$, and (b), (S1).
Used in Section~1.2, Step~2;
through (R6) ($\mathfrak{e}\le1/4$, $\sigma\le1/4$) in Section~2.2, Steps~2 and~3.

\item[$\sigma=s\left(\frac14-9\varepsilon\right)-12\varepsilon-\frac{2a_{L}}{N_{*}}>0$.]
Imposed in Section~I.8.2, line (C7); Proposition~14.2\,(b), (S2).
Used in Section~1.2, Step~2;
through (R6) in Section~2.2 at (0.3), in Step~4 and in Step~8, (PU-1).
The inequality $\sigma\le1$, a consequence of (R6), enters Section~2.2, Step~3.

\item[$\mathfrak{e}_2=0$.]
Imposed in Section~I.8.5; Proposition~14.2\,(a), last line.
Used in Section~1.2, Steps~2 and~7;
through (R6) ($0\le\mathfrak{e}_2<1$) in Section~2.2, Step~2 and Step~8, (PU-1).

\item[$e_1>0$.]
Imposed in Section~I.0.5; Section~I.8.2, line (C9); Proposition~14.3\,(f).
Used in (0.G), and through (0.G) in Section~2.2 at (0.3).

\item[$e_1\le1/3$.]
Imposed in Section~I.8.2, line (C9); Proposition~14.2\,(a), row $e_1$, and (f).
Used in Section~1.2 at (1.2), and through (R1) in Section~2.2 at (0.3) and (0.4).

\item[$v_k\ge8$ and $\left(1+\frac{3\,\mathrm{lg}}{v_k}\right)^r<L$ for all $v_k\ge-\log e_1$.]
Imposed in Section~I.8.2, line (C9); Proposition~14.2\,(b), row (e1), first and third inequality.
Used in Section~1.2 at (1.2), and through (1.2) at (1.3) and (1.7).

\item[$\exp\left(-\left(\frac14-9\varepsilon\right)v_k\right)\le\frac18$ for all $v_k\ge-\log e_1$.]
Imposed in Section~I.8.2, line (C9); Proposition~14.2\,(b), row (e3), third inequality.
Used in Section~1.2, Step~4, for the last member of Proposition~32.3\,(10.3).

\item[$6N_{*}M\cdot v_k^r\cdot\exp(-v_k)\cdot\left(v_k\cdot\exp(-v_k)+\frac12\right)\le\pi^2$ for all $v_k\ge-\log e_1$.]
Imposed in Section~I.8.2, line (C9); Proposition~14.2\,(b), row (e8).
Used in Section~1.2 at (1.8).

\item[(Z8.w) $e\cdot h^{\frac12}\cdot\log h\le1$.]
Imposed in Section~I.8.2, line (C10), where it is contained in $e\le e_G$; Proposition~14.2\,(d).
Used in Section~1.2 at (1.8).

\item[(g1)--(g9).]
Imposed in Section~I.8.2, line (C10), where they are contained in $e\le e_{\sharp}(\beta)$; Section~G5.1.9, that is, Proposition~14.2\,(d).
They enter through Propositions~35.9\,(b) and~35.10, which are applied in Section~1.2, Steps~5 and~6.

\item[(Pb) $b(f)>0$, $d_0>0$ and the inequality on $x_0$.]
Imposed in Section~I.8.2, line (C10); Section~I.8.4; Proposition~14.2\,(d).
The inequality on $x_0$ is the one displayed among the data of Proposition~43.2, and it enters as one of these data;
it is the condition on $x_0$ in \cite[Section~2.8]{QEDI}.
The condition $b(f)>0$ is used in Section~2.2 at (0.0) and at (0.3), and $d_0>0$ in Section~2.2 at (0.1).
In (0.0) the whole of (Pb) is shown to be satisfiable;
there $b(f)>0$ is the positivity of the right side of the inequality on $x_0$, which is considered on the range $d_0>0$.

\item[(Z10.s1) $e_2\le e_G$.]
Imposed in Section~I.8.2, line (C10); Section~I.8.4; Proposition~14.2\,(d).
Used in Section~2.2 at (0.3) and (0.4).

\item[(Z10.s2) $e_2\le e_{\sharp}(\beta_{\mathrm{ix}})$ for $\mathrm{ix}=0,+,-$.]
Imposed in Section~I.8.2, line (C10); Section~I.8.4; Proposition~14.2\,(d).
Used in Section~2.2 at (0.4).

\item[(Z10.s4) $\frac52\,C'_{\bar{\kappa}}\cdot e_2^{\sigma}\le\frac{b(f)}{4}$.]
Imposed in Section~I.8.2, line (C10); Section~I.8.4; Proposition~14.2\,(d).
Used in Section~2.2, Step~8, (PU-1).
\end{description}

All of these conditions are imposed in Section~I.0.5 or in Section~I.8,
and the proofs of Propositions~43.1 and~43.2 impose no further condition on $\rho$, $\alpha$, $\kappa$, $L$, $M$, $\varepsilon$, $s$, $e_0$, $e_1$.
The inequalities (0.H) rest on $L>1$ and on the range $\bar{m}\in(0,1]$, which is one of the standing data (Section~I.0.5).

\section*{Chapter 45-2. The ergodic decomposition, tensor tests, one additional mirror: proofs of Propositions 45.5, 45.6, 45.7}

\subsection*{0. Conventions of this chapter: setting, notation, auxiliary results, transport by a rotation}

\subsubsection*{0.1. Setting}

The three statements proved in this chapter concern only probability measures on $(\mathcal S',\Sigma)$; they and their proofs are independent of the lattice constructions, of the charge and of the constants of Section~I.8.

\subsubsection*{0.2. Notation}

The notation is that of (n1)--(n8) of Part~II (the notation for laws on $\mathcal S'$) and of the head of the Propositions~45.m.
For $y\in\mathbb R^3$ the norms $|y|_2$, $|y|_\infty$, $|y|_1$ are those of Section~I.1.1.
For a point $y=(y_1,\dots,y_n)$ of $\mathbb R^{3n}=(\mathbb R^3)^n$, with $y_a\in\mathbb R^3$ and the coordinates $(y_a)_\ell$, $\ell\in\{1,2,3\}$, as in Section~2, the same symbols denote the norms formed over the $3n$ coordinates:
\[
|y|_\infty:=\max_{a,\ell}\left|(y_a)_\ell\right|,\qquad|y|_1:=\sum_{a,\ell}\left|(y_a)_\ell\right|,\qquad|y|_2:=\left(\sum_{a,\ell}(y_a)_\ell^2\right)^{\frac12};
\]
these are the sup norm, the sum norm and the Euclidean norm of the vector of the $3n$ coordinates, and the same formulas, with $\mathrm{zl}$ or $\mathrm{kf}$ in the place of $y$, define them on $\mathbb Z^{3n}$.
Unless the contrary is said, test forms $f$, $f_p$, $f'$ are real elements of $\mathcal S=\mathcal S(\mathbb R^3;\Lambda^2)$, and
\[
\langle f,f'\rangle:=\sum_{i<\ell}\int\mathrm{d}y\,f_{i\ell}(y)\cdot f'_{i\ell}(y),\qquad\|f\|_2:=\langle f,f\rangle^{\frac12}.
\]

\emph{Exponentials.}
The exponential function is written $\exp(\cdot)$, also where its argument begins with $\pm i$, as in $\exp\left(iX(f)\right)$.
Accordingly, a law $\xi$ has the property (G) of (n3) if
$\xi\left(\exp X(f)\right)\le\exp\left(\frac{\|f\|_2^2}{2}\right)$ for every real $f\in\mathcal S$.

\emph{Letters with two meanings.}
The letter $i$ is the imaginary unit, except where it is the first index of a pair $i<\ell$ of coordinates (as above and in 2.1) and where it is the index of a coordinate in the proof of (in2); no other index is denoted by $i$.
In Sections~0, 1 and~3 the letter $a$ is a vector of $\mathbb R^3$ (as in $\tau_a$, $U(a)$ of (n2)).
In Section~2 it is the number of an argument, $y=(y_1,\dots,y_n)$ with $y_a\in\mathbb R^3$, as in Proposition~45.5\,(c), (d), and it has no other meaning there.
Throughout this chapter the letter $b$ denotes a vector of $\mathbb R^3$, as in (L3) and in Propositions~45.1\,(5) and~45.2\,(4).

\emph{Indices and parameters.}
In this chapter $k$ is an index, and $t$, $t'$, $t_p$, $t'_p$ are real numbers, as in Propositions~45.1\,(3) and~45.3\,(1), (3).
The letter $\ell$ is the number of a coordinate: the second index of a pair $i<\ell$, as in (n3), and $\ell\in\{1,2,3\}$ in $(y_a)_\ell$, $\partial_{a,\ell}$, $(\mathrm{kf}_a)_\ell$.
The letters $m$, $n$, $n'$ denote integers $\ge0$: the number of the arguments of a tensor test and the index of a norm $q_m$, as in (n5) and Proposition~45.5, and also an exponent (in (qc5) and in Step~9 of Section~1).
The letters $p$, $q$ are indices, as in (n3); in (qc4)\,(iii) the letter $p$ is the exponent of $L^p(\xi)$.
The letters $s$, $s'$ are real numbers in $\tau_{s\cdot e_j}$, $U(s\cdot e_j)$, $U(s\cdot\hat n)$, as in Proposition~45.7\,(c); in the proof of (La1), in Step~2 of Section~1 and in 2.6, $s$ is the number of the terms of a finite sum, as in (L4).
The sets in which $\mathrm{zl}$, $\mathrm{kf}$, $\mathrm{ki}$ vary ($\mathbb Z$, $\mathbb Z^3$ or $\mathbb Z^{3n}$) are named at each place.

\emph{Laws.}
$\mathrm{Law}$ denotes the set of the laws (n1), and $B_{\mathrm{Law}}$ the $\sigma$-algebra of (n7) on $\mathrm{Law}$, that is, the $\sigma$-algebra generated by the maps $\xi\mapsto\xi(\mathrm{Ev})$, $\mathrm{Ev}\in\Sigma$.
$\mathrm{Law}_{\mathrm{erg}}$ denotes the set of the $\mathbb R^3$-ergodic laws (n3).
Points of $\mathrm{Law}$ are written $\xi$, $\xi'$; in the statements the same point, regarded as a law, is written $\mathrm{lw}_\xi$ (n7), so that $\mathrm{lw}_\xi=\xi$.
For a measurable map $\mathrm{Gi}$ of $\mathcal S'$ into itself, $\mathrm{Gi}_*\xi$ is the image of $\xi$:
$(\mathrm{Gi}_*\xi)(\mathrm{Ev})=\xi\{\omega:\mathrm{Gi}\,\omega\in\mathrm{Ev}\}$.

\emph{Almost every.}
As in (n7), ``for $\nu$-almost every $\xi$'' means: for all $\xi$ in a set $\mathrm{Ls}\in B_{\mathrm{Law}}$ with $\nu(\mathrm{Ls})=1$.

\emph{Correspondence with \cite[Section~8]{QEDI}.}
The properties of laws introduced in \cite[Section~8.1.3]{QEDI} carry the following names there:
$(T_{123})$ is the invariance under $\mathbb R^3$;
$(T_j)$ is the invariance under $\mathbb R e_j$, that is, under all $\tau_{s\cdot e_j}$, $s\in\mathbb R$;
$(R_j)$ is the invariance under $\theta_j$;
$(RP_j)$ is the reflection positivity at $\{y_j=0\}$.
The reflection $\theta_{j,c}$ of \cite[Section~8.1.2]{QEDI} is $r_{e_j,\mathrm{po}}$ with $\mathrm{po}=c$, and the operator $\Theta_j$ of \cite[Section~8.1.2]{QEDI} is the complex linear operator $\hat\theta_j$.
The other symbols of \cite[Section~8]{QEDI} correspond to those of this chapter as follows.

\begin{tabular}{ll}
in \cite[Section~8]{QEDI} & in this chapter\\
\hline
$Q$ & $\mathcal S'$\\
$\rho$ & $\xi$\\
$m$, $\mu$ & $\mathrm{lw}$\\
$\mathcal M$, $\mathcal B_{\mathcal M}$, $\mathcal M^{123}_{\mathrm{erg}}$ & $\mathrm{Law}$, $B_{\mathrm{Law}}$, $\mathrm{Law}_{\mathrm{erg}}$\\
$\mathcal D$, $\mathcal D(O)$ & $C_c^\infty(\mathbb R^3;\Lambda^2)$, $C_c^\infty(O;\Lambda^2)$\\
$\mathcal D_0$, $\mathcal D_0(O)$ & $\mathrm{Dq}$, $\mathrm{Dq}(O)$\\
$\mathcal A^0(O)$, $\mathcal A^0$ & $\mathrm{Aq}(O)$, $\mathrm{Aq}$\\
$F$, $G$ & $\mathrm{Ga}$, $\mathrm{Gb}$\\
$A$ (an event) & $\mathrm{Ev}$\\
$\mathcal I^{123}$, its elements & $I$, its elements $\mathrm{Ie}$\\
$g$, $R$ & $\mathrm{Gi}$, $\mathrm{Rl}$\\
an element of $O(3)$ & $\mathrm{gr}$\\
$n$, $c$, $r_{n,c}$ & $\hat n$, $\mathrm{po}$, $r_{\hat n,\mathrm{po}}$\\
$A^{123}_TF$, $T$ & $\mathrm{Avg}_{\mathrm{sl}}\mathrm{Ga}$, $\mathrm{sl}$
\end{tabular}

The space $\mathrm{Sw}_n$ of (n5) is the space $\mathscr S_n$ of \cite[Section~2.5]{QEDI}; for $n\ge1$ it is also the space $\mathcal S_n$ of \cite[Section~7.6]{QEDI}.

\emph{Further notation.}
The letter $\varphi$ denotes a function on $\mathbb C$ (in (qc1)), $\mathrm{Ls}$ a set in $B_{\mathrm{Law}}$, $\mathrm{Tg}$ the map of 0.5, and $\mathrm{Gu}$ a function as in Proposition~45.6\,(a$'$) (Step~2 of Section~1).
The letters $\alpha$, $\beta$, $\gamma$ denote multi-indices ($\alpha$ as in (n5)), and $\mathrm{cf}$, $\mathrm{cf}_p$, $\mathrm{cf}_n$, $\mathrm{cf}_{n,m}$ denote numbers.
The letter $\mathbb Q$ denotes the field of the rational numbers; it occurs only in $\mathbb Q^3$ and in $\mathbb Q+i\mathbb Q$.

\subsubsection*{0.3. Auxiliary results}

The clauses (qc1)--(qc9) below are results of \cite[Section~8]{QEDI}, written in the notation of 0.2 and specialised, where \cite{QEDI} admits an index set $J\subset\{1,2,3\}$, to $J=\{1,2,3\}$; the clause (qc10) consists of two results of \cite{Kallenberg2002}.

\begin{itemize}
\item[(qc1)] \cite[Section~8.1.2]{QEDI}.
Let $\mathrm{Gi}$, $\mathrm{Gi}'$ be Euclidean maps, $f\in\mathcal S$, and $O\subset\mathbb R^3$ open.
Then
\[
\widehat{\mathrm{Gi}}\,X(f)=X(\mathrm{Gi}\cdot f),\qquad\left(\mathrm{Gi}\,\mathrm{Gi}'\right)^{\wedge}=\widehat{\mathrm{Gi}}\,\widehat{\mathrm{Gi}'},\qquad\|\mathrm{Gi}\cdot f\|_2=\|f\|_2 ,
\]
where the operator that (n2) attaches to a product of Euclidean maps, such as $\mathrm{Gi}\,\mathrm{Gi}'$, is written $\left(\mathrm{Gi}\,\mathrm{Gi}'\right)^{\wedge}$.
The operator $\widehat{\mathrm{Gi}}$ maps the $\Sigma_O$-measurable functions onto the $\Sigma_{\mathrm{Gi}\,O}$-measurable functions.
It is linear and multiplicative, and it commutes with complex conjugation and with $\mathrm{Ga}\mapsto\varphi\circ\mathrm{Ga}$ for every function $\varphi$ on $\mathbb C$.
If $\mathrm{Gi}_*\xi=\xi$, then $\widehat{\mathrm{Gi}}$ is unitary on $L^2(\xi)$.
Moreover $U(a)X(f)=X(\tau_a\cdot f)$, where $(\tau_a\cdot f)(y)=f(y-a)$.

The second identity rests on the two identities
\begin{align*}
\left(\mathrm{Gi}\cdot(\mathrm{Gi}'\cdot f)\right)(y)&=\left(\wedge^2\mathrm{Rl}\right)\,\left(\wedge^2\mathrm{Rl}'\right)\,f\left(\mathrm{Gi}'^{-1}\,\mathrm{Gi}^{-1}\,y\right)=\left((\mathrm{Gi}\,\mathrm{Gi}')\cdot f\right)(y),\\
\left(\mathrm{Gi}\,(\mathrm{Gi}'\,\omega)\right)(f)&=(\mathrm{Gi}'\,\omega)\left(\mathrm{Gi}^{-1}\cdot f\right)=\omega\left(\mathrm{Gi}'^{-1}\cdot\left(\mathrm{Gi}^{-1}\cdot f\right)\right)\\
&=\omega\left((\mathrm{Gi}\,\mathrm{Gi}')^{-1}\cdot f\right)=\left((\mathrm{Gi}\,\mathrm{Gi}')\,\omega\right)(f),
\end{align*}
which follow from (n2); here $\mathrm{Rl}$ and $\mathrm{Rl}'$ are the linear parts of $\mathrm{Gi}$ and $\mathrm{Gi}'$.
For $f\in\mathcal S$ one has $X(f)(\mathrm{Gi}\,\omega)=\omega\left(\mathrm{Gi}^{-1}\cdot f\right)=X\left(\mathrm{Gi}^{-1}\cdot f\right)(\omega)$.
Since $\Sigma$ is generated by the functions $X(f)$, $f\in\mathcal S$ (n1), the bijection $\omega\mapsto\mathrm{Gi}\,\omega$ of $\mathcal S'$ is $\Sigma$-measurable, and so is its inverse $\omega\mapsto\mathrm{Gi}^{-1}\omega$.
Consequently $\mathrm{Gi}\,\mathrm{Ev}\in\Sigma$ for $\mathrm{Ev}\in\Sigma$, the image $\mathrm{Gi}_*\xi$ of a law is a law, and $(\mathrm{Gi}\,\mathrm{Gi}')_*\xi=\mathrm{Gi}_*(\mathrm{Gi}'_*\xi)$.

\item[(qc2)] \cite[Lemma~8.11]{QEDI}.
The \emph{seven sets} are $\mathbb R^3$, $\{y_j>0\}$ and $\{y_j<0\}$, $j=1,2,3$.
\begin{itemize}
\item[(i)] There is a countable vector space $\mathrm{Dq}$ over the rational numbers, contained in $C_c^\infty(\mathbb R^3;\Lambda^2)$ and stable under $O_{\mathrm{cub}}$ and under $\tau_a$ for $a\in\mathbb Q^3$ (the points with rational coordinates), such that for each of the seven sets $O$ the space $\mathrm{Dq}(O):=\mathrm{Dq}\cap C_c^\infty(O;\Lambda^2)$ is dense, for the topology of $\mathcal S$, in $\{f\in\mathcal S:\operatorname{supp}f\subset\bar O\}$.
A space $\mathrm{Dq}$ with these properties is fixed.
\item[(ii)] $\Sigma_{\mathbb R^3}=\Sigma$.
\item[(iii)] $(\mathcal S',\Sigma)$ is a standard Borel space.
\end{itemize}

\item[(qc3)] \cite[Lemma~8.12]{QEDI}.
For each of the seven sets $O$ let $\mathrm{Aq}(O)$ be the countable set of the finite linear combinations of the functions $\exp\left(iX(f)\right)$, $f\in\mathrm{Dq}(O)$, with coefficients in $\mathbb Q+i\mathbb Q$ (the complex numbers with rational real and imaginary parts), and let $\mathrm{Aq}:=\mathrm{Aq}(\mathbb R^3)$.
Then for every law $\xi$:
\begin{itemize}
\item[(i)] $\mathrm{Aq}(O)$ is dense in $L^2(\Sigma_O,\xi)$;
\item[(ii)] $\xi$ is determined by the numbers $\xi\left(\exp\left(iX(f)\right)\right)$, $f\in\mathrm{Dq}$.
\end{itemize}

\item[(qc4)] \cite[Lemma~8.13]{QEDI}, with the index set $\{1,2,3\}$.
Here $I$ is the $\sigma$-algebra of the events $\mathrm{Ev}\in\Sigma$ with $\tau_a\mathrm{Ev}=\mathrm{Ev}$ for all $a\in\mathbb R^3$, as in Proposition~45.6.
\begin{itemize}
\item[(i)] The map $(a,\omega)\mapsto\tau_a\omega$ from $\mathbb R^3\times\mathcal S'$ to $\mathcal S'$ is measurable.
For a measurable function $\mathrm{Ga}$ on $\mathcal S'$ that is bounded or non-negative, and for $\mathrm{sl}>0$, the function
\[
(\mathrm{Avg}_{\mathrm{sl}}\mathrm{Ga})(\omega):=\mathrm{sl}^{-3}\cdot\int_{[0,\mathrm{sl}]^3}\mathrm{d}a\,\mathrm{Ga}(\tau_{-a}\omega)
\]
is measurable.
\end{itemize}
Let now $\xi$ be invariant under $\mathbb R^3$.
\begin{itemize}
\item[(iii)] For $\mathrm{Ga}\in L^2(\xi)$ the function $\mathrm{Avg}_{\mathrm{sl}}\mathrm{Ga}$ is defined $\xi$-almost everywhere, $\|\mathrm{Avg}_{\mathrm{sl}}\mathrm{Ga}\|_{L^p(\xi)}\le\|\mathrm{Ga}\|_{L^p(\xi)}$ for $1\le p<\infty$, and $\mathrm{Avg}_{\mathrm{sl}}\mathrm{Ga}\to P^{123}_\xi\mathrm{Ga}$ in $L^2(\xi)$ as $\mathrm{sl}\to\infty$.
\item[(iv)] $\left\{\mathrm{Ga}\in L^2(\xi):U(a)\mathrm{Ga}=\mathrm{Ga}\text{ for all }a\in\mathbb R^3\right\}=L^2(I,\xi)$, and $P^{123}_\xi=E_\xi[\,\cdot\mid I\,]$.
Here $L^2(I,\xi)$ is the closed subspace of $L^2(\xi)$ formed by the classes that have an $I$-measurable representative; it is defined as $L^2(\Sigma_O,\xi)$ in (n1), with the $\sigma$-algebra $I$ in the place of $\Sigma_O$, and it is the space $L^2(\mathcal G,\rho)$ of \cite[Section~8.1.1]{QEDI} with $\mathcal G=I$ and $\rho=\xi$.
\end{itemize}

\item[(qc5)] \cite[Lemma~8.15(i)]{QEDI}.
If $\xi$ has (G), then for $f\in\mathcal S$ and every integer $n\ge1$
\[
\xi\left(|X(f)|^n\right)\le n\cdot2^{\frac n2}\cdot\Gamma\left(\frac n2\right)\cdot\|f\|_2^n .
\]

\item[(qc6)] \cite[Lemma~8.16]{QEDI}.
Let $\xi$ be invariant under $\theta_j$.
The following conditions are equivalent:
\begin{itemize}
\item[(1)] $\xi$ is reflection positive at $\{y_j=0\}$;
\item[(2)] $\xi\left(\overline{\hat\theta_j\mathrm{Ga}}\cdot\mathrm{Ga}\right)\ge0$ for all $\mathrm{Ga}\in\mathrm{Aq}(\{y_j>0\})$;
\item[(3)] $\left\langle\hat\theta_j\mathrm{Ga},\mathrm{Ga}\right\rangle_\xi\ge0$ for all $\mathrm{Ga}\in L^2(\Sigma_{\{y_j>0\}},\xi)$.
\end{itemize}
If $\xi$ is in addition invariant under $\mathbb R e_j$, these conditions are also equivalent to
\begin{itemize}
\item[(5)] $\xi$ is reflection positive at every plane $\{y_j=\mathrm{po}\}$, $\mathrm{po}\in\mathbb R$.
\end{itemize}

\item[(qc7)] \cite[Lemma~8.17]{QEDI}.
The following sets of laws belong to $B_{\mathrm{Law}}$: the set of the laws that are invariant under $\mathbb R^3$; the set of those invariant under $\theta_j$; the set of those reflection positive at $\{y_j=0\}$; the set of those with (G); and $\mathrm{Law}_{\mathrm{erg}}$.
More precisely, for every law $\xi$:
\begin{itemize}
\item $\xi$ is invariant under $\mathbb R^3$ if and only if $\xi\left(\exp\left(iX(\tau_a\cdot f)\right)\right)=\xi\left(\exp\left(iX(f)\right)\right)$ for all $a\in\mathbb Q^3$ and $f\in\mathrm{Dq}$;
\item $\xi$ is invariant under $\theta_j$ if and only if $\xi\left(\exp\left(iX(\theta_j\cdot f)\right)\right)=\xi\left(\exp\left(iX(f)\right)\right)$ for all $f\in\mathrm{Dq}$;
\item $\xi$ is reflection positive at $\{y_j=0\}$ if and only if condition (2) of (qc6) holds, whether or not $\xi$ is invariant under $\theta_j$;
\item $\xi$ has (G) if and only if $\xi\left(\exp X(f)\right)\le\exp\left(\frac{\|f\|_2^2}{2}\right)$ for all $f\in\mathrm{Dq}$.
\end{itemize}

\item[(qc8)] \cite[Lemma~8.18(iii)]{QEDI}.
Let $\xi$ be invariant under $\mathbb R e_j$ and under $\theta_j$, and reflection positive at $\{y_j=0\}$.
Let $\mathrm{Gb}\in L^2(\xi)$ be bounded with $U(s\cdot e_j)\mathrm{Gb}=\mathrm{Gb}$ for all $s\in\mathbb R$.
Then
\[
\xi\left(\mathrm{Gb}\cdot\hat\theta_j\mathrm{Ga}\right)=\xi\left(\mathrm{Gb}\cdot\mathrm{Ga}\right)\qquad\text{for }\mathrm{Ga}\in L^1(\xi),
\]
and, if $\mathrm{Gb}\ge0$,
\[
\xi\left(\mathrm{Gb}\cdot\overline{\hat\theta_j\mathrm{Ga}}\cdot\mathrm{Ga}\right)\ge0\qquad\text{for }\mathrm{Ga}\in L^2(\Sigma_{\{y_j>0\}},\xi).
\]

\item[(qc9)] \cite[Proposition~8.20]{QEDI}, the last assertion.
Let $\xi$ be as in (qc8).
Then $U(s\cdot e_j)\to P^j_\xi$ weakly on $L^2(\xi)$ as $|s|\to\infty$.

\item[(qc10)] (i) \cite[Theorem~6.3]{Kallenberg2002}.
A \emph{Borel space} is a measurable space that is isomorphic to a Borel subset of $[0,1]$ endowed with the trace of the Borel $\sigma$-algebra.
Let a random element of a Borel space and a random element of a measurable space be defined on one probability space.
Then there is a probability kernel from the second space to the first whose value at the second random element is a regular version of the conditional distribution of the first random element given the second one.

(ii) \cite[Theorem~A1.2]{Kallenberg2002}.
A Polish space, endowed with its Borel $\sigma$-algebra, is a Borel space.
\end{itemize}

\subsubsection*{0.4. Three facts on $(\mathrm{Law},B_{\mathrm{Law}})$: Proposition~45.6\,(d0)}

The three sentences of Proposition~45.6\,(d0) are denoted (La1), (La2), (La3).
In them the set of the laws is $\mathrm{Law}$, the $\sigma$-algebra of (n7) is $B_{\mathrm{Law}}$, and $\mathrm{lw}_\xi=\xi$ (0.2); no particular law $\mathrm{lw}$ is given.
For a $\Sigma$-measurable $\Phi\colon\mathcal S'\to\mathbb C$ put
\[
\mathrm{Law}(\Phi):=\left\{\xi\in\mathrm{Law}:\xi(|\Phi|)<\infty\right\}.
\]
In this notation the three sentences read as follows.
\begin{itemize}
\item[(La1)] For every $\Sigma$-measurable $\Phi\colon\mathcal S'\to[0,\infty]$ the map $\xi\mapsto\xi(\Phi)$ from $\mathrm{Law}$ to $[0,\infty]$ is $B_{\mathrm{Law}}$-measurable.
\item[(La2)] For every $\Sigma$-measurable $\Phi\colon\mathcal S'\to\mathbb C$ the set $\mathrm{Law}(\Phi)$ belongs to $B_{\mathrm{Law}}$, and $\xi\mapsto\xi(\Phi)$ is $B_{\mathrm{Law}}$-measurable on $\mathrm{Law}(\Phi)$.
\item[(La3)] Every one-point set $\{\xi_0\}\subset\mathrm{Law}$ belongs to $B_{\mathrm{Law}}$.
\end{itemize}

\begin{proof}[Proof of (La1)]
For $\Phi=\mathbf 1_{\mathrm{Ev}}$ with $\mathrm{Ev}\in\Sigma$ the map $\xi\mapsto\xi(\mathrm{Ev})$ is one of the maps that define $B_{\mathrm{Law}}$.
For $\Phi=\sum_{p\le s}t_p\cdot\mathbf 1_{\mathrm{Ev}_p}$ with $t_p\ge0$ and $\mathrm{Ev}_p\in\Sigma$ one has $\xi(\Phi)=\sum_pt_p\cdot\xi(\mathrm{Ev}_p)$, a finite sum of measurable functions of $\xi$.
A general $\Phi$ is the pointwise limit of an increasing sequence $(\Phi_k)$ of functions of this kind, and $\xi(\Phi)=\lim_k\xi(\Phi_k)$ for every $\xi$ by monotone convergence; a pointwise limit of measurable functions is measurable.
\end{proof}

\begin{proof}[Proof of (La2)]
By (La1) applied to $|\Phi|$, the map $\xi\mapsto\xi(|\Phi|)$ is measurable, so that $\mathrm{Law}(\Phi)\in B_{\mathrm{Law}}$.
For $\xi\in\mathrm{Law}(\Phi)$ the four numbers $\xi\left((\operatorname{Re}\Phi)^+\right)$, $\xi\left((\operatorname{Re}\Phi)^-\right)$, $\xi\left((\operatorname{Im}\Phi)^+\right)$, $\xi\left((\operatorname{Im}\Phi)^-\right)$ are finite, they are measurable functions of $\xi$ by (La1), and
\[
\xi(\Phi)=\xi\left((\operatorname{Re}\Phi)^+\right)-\xi\left((\operatorname{Re}\Phi)^-\right)+i\cdot\xi\left((\operatorname{Im}\Phi)^+\right)-i\cdot\xi\left((\operatorname{Im}\Phi)^-\right).
\]
\end{proof}

\begin{proof}[Proof of (La3)]
By (qc3)\,(ii),
\[
\{\xi_0\}=\bigcap_{f\in\mathrm{Dq}}\left\{\xi:\xi\left(\exp\left(iX(f)\right)\right)=\xi_0\left(\exp\left(iX(f)\right)\right)\right\}.
\]
The intersection is countable by (qc2)\,(i), and each of the sets belongs to $B_{\mathrm{Law}}$ by (La2), applied to the bounded function $\exp\left(iX(f)\right)$, for which $\mathrm{Law}\left(\exp\left(iX(f)\right)\right)=\mathrm{Law}$.
\end{proof}

\subsubsection*{0.5. Transport of a law by a rotation}

\begin{named}{Transport Lemma}
Let $\mathrm{gr}\in O(3)$, regarded as the Euclidean map $y\mapsto\mathrm{gr}\,y$, put $\hat n:=\mathrm{gr}\,e_1$, let $\xi'$ be a law, and put $\xi:=\mathrm{gr}_*\xi'$.
\begin{itemize}
\item[(tr1)] $\xi\left(\widehat{\mathrm{gr}}\,\Phi\right)=\xi'(\Phi)$ for every $\Sigma$-measurable $\Phi\ge0$ and for every bounded $\Sigma$-measurable complex $\Phi$.
The operator $\widehat{\mathrm{gr}}$ maps classes modulo $\xi'$ to classes modulo $\xi$; it is a unitary map of $L^2(\xi')$ onto $L^2(\xi)$ with the inverse $\widehat{\mathrm{gr}^{-1}}$; and for every open set $O$ it maps $L^2(\Sigma_O,\xi')$ onto $L^2(\Sigma_{\mathrm{gr}\,O},\xi)$.
\item[(tr2)] For every Euclidean map $\mathrm{Gi}$ one has $\left(\mathrm{gr}\,\mathrm{Gi}\,\mathrm{gr}^{-1}\right)_*\xi=\mathrm{gr}_*(\mathrm{Gi}_*\xi')$.
Hence $\xi$ is invariant under $\mathrm{gr}\,\mathrm{Gi}\,\mathrm{gr}^{-1}$ if and only if $\xi'$ is invariant under $\mathrm{Gi}$.
As operators on functions on $\mathcal S'$, $\widehat{\mathrm{gr}}\,\widehat{\mathrm{Gi}}\,\widehat{\mathrm{gr}^{-1}}=\left(\mathrm{gr}\,\mathrm{Gi}\,\mathrm{gr}^{-1}\right)^{\wedge}$, in the notation of (qc1).
\item[(tr3)] $\mathrm{gr}\,\tau_a\,\mathrm{gr}^{-1}=\tau_{\mathrm{gr}\,a}$ for $a\in\mathbb R^3$.
For real $\mathrm{po}$ one has $\mathrm{gr}\,r_{e_1,\mathrm{po}}\,\mathrm{gr}^{-1}=r_{\hat n,\mathrm{po}}$ and $\mathrm{gr}\{y_1>\mathrm{po}\}=\{\hat n\cdot y>\mathrm{po}\}$, and $r_{\hat n,\mathrm{po}}=\tau_{\mathrm{po}\cdot\hat n}\,r_{\hat n,0}\,\tau_{-\mathrm{po}\cdot\hat n}$.
\item[(tr4)] $\xi$ is invariant under $\mathbb R^3$ if and only if $\xi'$ is; $\xi\in\mathrm{Law}_{\mathrm{erg}}$ if and only if $\xi'\in\mathrm{Law}_{\mathrm{erg}}$.
\item[(tr5)] For real $\mathrm{po}$, $\xi$ is reflection positive at $\{\hat n\cdot y=\mathrm{po}\}$ if and only if $\xi'$ is reflection positive at $\{y_1=\mathrm{po}\}$.
\item[(tr6)] The map $\mathrm{Tg}\colon\mathrm{Law}\to\mathrm{Law}$, $\mathrm{Tg}(\xi'):=\mathrm{gr}_*\xi'$, is a bijection; it and its inverse are $B_{\mathrm{Law}}$-measurable, and $\mathrm{Tg}(\mathrm{Law}_{\mathrm{erg}})=\mathrm{Law}_{\mathrm{erg}}$.
\end{itemize}
For every unit vector $\hat n$ there is $\mathrm{gr}\in O(3)$ with $\mathrm{gr}\,e_1=\hat n$: complete $\hat n$ to an orthonormal basis $(\hat n,\hat n',\hat n'')$ of $\mathbb R^3$, and let $\mathrm{gr}$ be the matrix with these columns.
\end{named}

\begin{proof}
(tr1)
By (n2), $\left(\widehat{\mathrm{gr}}\,\Phi\right)(\mathrm{gr}\,\omega)=\Phi\left(\mathrm{gr}^{-1}\,\mathrm{gr}\,\omega\right)=\Phi(\omega)$.
The change of variables for the image measure $\xi=\mathrm{gr}_*\xi'$ therefore gives
\[
\xi\left(\widehat{\mathrm{gr}}\,\Phi\right)=\int\xi'(\mathrm{d}\omega)\,\left(\widehat{\mathrm{gr}}\,\Phi\right)(\mathrm{gr}\,\omega)=\xi'(\Phi).
\]
Let $\mathrm{Ga}=\mathrm{Gb}$ outside an event $\mathrm{Ev}$ with $\xi'(\mathrm{Ev})=0$.
Then $\widehat{\mathrm{gr}}\,\mathrm{Ga}=\widehat{\mathrm{gr}}\,\mathrm{Gb}$ outside $\mathrm{gr}\,\mathrm{Ev}$, and $\xi(\mathrm{gr}\,\mathrm{Ev})=\xi\left(\widehat{\mathrm{gr}}\,\mathbf 1_{\mathrm{Ev}}\right)=\xi'(\mathrm{Ev})=0$.
With $\Phi=|\mathrm{Ga}|^2$ and $\left|\widehat{\mathrm{gr}}\,\mathrm{Ga}\right|^2=\widehat{\mathrm{gr}}\,|\mathrm{Ga}|^2$ (qc1) one obtains $\left\|\widehat{\mathrm{gr}}\,\mathrm{Ga}\right\|_{L^2(\xi)}=\|\mathrm{Ga}\|_{L^2(\xi')}$.
Thus $\widehat{\mathrm{gr}}$ is an isometry of $L^2(\xi')$ into $L^2(\xi)$.
The argument so far applies to every element of $O(3)$ and every law.
Applied to $\mathrm{gr}^{-1}$, $\xi$ and $\left(\mathrm{gr}^{-1}\right)_*\xi=\xi'$ in the places of $\mathrm{gr}$, $\xi'$ and $\xi$, it shows that $\widehat{\mathrm{gr}^{-1}}$ is an isometry of $L^2(\xi)$ into $L^2(\xi')$.
By (qc1), $\widehat{\mathrm{gr}}\,\widehat{\mathrm{gr}^{-1}}=1$ and $\widehat{\mathrm{gr}^{-1}}\,\widehat{\mathrm{gr}}=1$; hence $\widehat{\mathrm{gr}}$ is unitary.
If $\mathrm{Ga}$ has a $\Sigma_O$-measurable representative $\mathrm{Ga}_0$, then $\widehat{\mathrm{gr}}\,\mathrm{Ga}_0$ is a $\Sigma_{\mathrm{gr}\,O}$-measurable representative of $\widehat{\mathrm{gr}}\,\mathrm{Ga}$ (qc1); and $\widehat{\mathrm{gr}^{-1}}$ maps $\Sigma_{\mathrm{gr}\,O}$-measurable functions to $\Sigma_O$-measurable functions.

(tr2)
By (qc1),
\[
\left(\mathrm{gr}\,\mathrm{Gi}\,\mathrm{gr}^{-1}\right)_*\,\mathrm{gr}_*\xi'=\left(\mathrm{gr}\,\mathrm{Gi}\,\mathrm{gr}^{-1}\,\mathrm{gr}\right)_*\xi'=(\mathrm{gr}\,\mathrm{Gi})_*\xi'=\mathrm{gr}_*(\mathrm{Gi}_*\xi').
\]
If $\mathrm{Gi}_*\xi'=\xi'$, the right side equals $\xi$.
Conversely, if $\left(\mathrm{gr}\,\mathrm{Gi}\,\mathrm{gr}^{-1}\right)_*\xi=\xi$, then $\mathrm{gr}_*(\mathrm{Gi}_*\xi')=\mathrm{gr}_*\xi'$, and the application of $\left(\mathrm{gr}^{-1}\right)_*$ gives $\mathrm{Gi}_*\xi'=\xi'$.
The identity between the operators is the second assertion of (qc1), applied twice.

(tr3)
For $y\in\mathbb R^3$, $\mathrm{gr}\,\tau_a\,\mathrm{gr}^{-1}\,y=\mathrm{gr}\left(\mathrm{gr}^{-1}y+a\right)=y+\mathrm{gr}\,a$.
Next, $r_{\hat n,\mathrm{po}}(y)=y-2\left(\hat n\cdot y-\mathrm{po}\right)\cdot\hat n$.
Since $\hat n\cdot(\mathrm{gr}\,y)=(\mathrm{gr}\,e_1)\cdot(\mathrm{gr}\,y)=y_1$, it follows that
\[
r_{\hat n,\mathrm{po}}(\mathrm{gr}\,y)=\mathrm{gr}\,y-2\left(y_1-\mathrm{po}\right)\cdot\mathrm{gr}\,e_1=\mathrm{gr}\left(y-2\left(y_1-\mathrm{po}\right)\cdot e_1\right)=\mathrm{gr}\,r_{e_1,\mathrm{po}}(y).
\]
Moreover
\[
\mathrm{gr}\{y_1>\mathrm{po}\}=\left\{y:\left(\mathrm{gr}^{-1}y\right)_1>\mathrm{po}\right\}=\left\{y:e_1\cdot\left(\mathrm{gr}^{-1}y\right)>\mathrm{po}\right\}=\{y:\hat n\cdot y>\mathrm{po}\},
\]
and
\[
\tau_{\mathrm{po}\cdot\hat n}\,r_{\hat n,0}\,\tau_{-\mathrm{po}\cdot\hat n}\,y=\left(y-\mathrm{po}\cdot\hat n\right)-2\left(\hat n\cdot y-\mathrm{po}\right)\cdot\hat n+\mathrm{po}\cdot\hat n=r_{\hat n,\mathrm{po}}(y).
\]

(tr4)
The first assertion is (tr2) with $\mathrm{Gi}=\tau_a$, together with (tr3), because $a\mapsto\mathrm{gr}\,a$ is a bijection of $\mathbb R^3$.
Let $\xi'\in\mathrm{Law}_{\mathrm{erg}}$, and let $\mathrm{Ga}\in L^2(\xi)$ satisfy $U(b)\mathrm{Ga}=\mathrm{Ga}$ for all $b\in\mathbb R^3$.
Put $\mathrm{Gb}:=\widehat{\mathrm{gr}^{-1}}\,\mathrm{Ga}\in L^2(\xi')$ (tr1).
For $a\in\mathbb R^3$, (tr2) and (tr3) give
\[
U(a)\mathrm{Gb}=\widehat{\mathrm{gr}^{-1}}\left(\widehat{\mathrm{gr}}\,U(a)\,\widehat{\mathrm{gr}^{-1}}\right)\mathrm{Ga}=\widehat{\mathrm{gr}^{-1}}\,U(\mathrm{gr}\,a)\,\mathrm{Ga}=\widehat{\mathrm{gr}^{-1}}\,\mathrm{Ga}=\mathrm{Gb}.
\]
Hence $\mathrm{Gb}$ is constant, and $\mathrm{Ga}=\widehat{\mathrm{gr}}\,\mathrm{Gb}$ is the same constant; so $\xi\in\mathrm{Law}_{\mathrm{erg}}$.
For the converse, apply what has been proved to $\mathrm{gr}^{-1}$, $\xi$ and $\xi'=\left(\mathrm{gr}^{-1}\right)_*\xi$ in the places of $\mathrm{gr}$, $\xi'$ and $\xi$.

(tr5)
The map $f\mapsto f':=\mathrm{gr}^{-1}\cdot f$ is a bijection of $C_c^\infty(\{\hat n\cdot y>\mathrm{po}\};\Lambda^2)$ onto $C_c^\infty(\{y_1>\mathrm{po}\};\Lambda^2)$, because $\operatorname{supp}\left(\mathrm{gr}^{-1}\cdot f\right)=\mathrm{gr}^{-1}\operatorname{supp}f$ and by (tr3).
Let $f_p$ belong to the first space, and $f'_p:=\mathrm{gr}^{-1}\cdot f_p$.
By (qc1) and (tr3),
\[
\widehat{\mathrm{gr}}\,X(f'_q)=X(\mathrm{gr}\cdot f'_q)=X(f_q),\qquad\widehat{\mathrm{gr}}\,X\left(r_{e_1,\mathrm{po}}\cdot f'_p\right)=X\left(\left(\mathrm{gr}\,r_{e_1,\mathrm{po}}\,\mathrm{gr}^{-1}\right)\cdot f_p\right)=X\left(r_{\hat n,\mathrm{po}}\cdot f_p\right),
\]
and therefore
\[
\widehat{\mathrm{gr}}\,\exp\left(i\left(X(f'_q)-X\left(r_{e_1,\mathrm{po}}\cdot f'_p\right)\right)\right)=\exp\left(i\left(X(f_q)-X\left(r_{\hat n,\mathrm{po}}\cdot f_p\right)\right)\right).
\]
By (tr1),
\[
\sum_{p,q}\overline{c_p}\cdot c_q\cdot\xi\left(\exp\left(i\left(X(f_q)-X\left(r_{\hat n,\mathrm{po}}\cdot f_p\right)\right)\right)\right)=\sum_{p,q}\overline{c_p}\cdot c_q\cdot\xi'\left(\exp\left(i\left(X(f'_q)-X\left(r_{e_1,\mathrm{po}}\cdot f'_p\right)\right)\right)\right).
\]
Hence the two families of inequalities of (n3) that define the two reflection positivities are the same family.

(tr6)
For $\mathrm{Ev}\in\Sigma$ one has $\mathrm{Tg}(\xi')(\mathrm{Ev})=\xi'\left(\mathrm{gr}^{-1}\mathrm{Ev}\right)$, and $\mathrm{gr}^{-1}\mathrm{Ev}\in\Sigma$ (qc1).
So the composition of $\mathrm{Tg}$ with the map $\xi\mapsto\xi(\mathrm{Ev})$ is the map $\xi'\mapsto\xi'\left(\mathrm{gr}^{-1}\mathrm{Ev}\right)$, which is one of the maps that define $B_{\mathrm{Law}}$; hence $\mathrm{Tg}$ is measurable.
By (qc1), the map $\xi\mapsto\left(\mathrm{gr}^{-1}\right)_*\xi$ is the inverse of $\mathrm{Tg}$.
For $\mathrm{Ev}\in\Sigma$ one has $\left(\left(\mathrm{gr}^{-1}\right)_*\xi\right)(\mathrm{Ev})=\xi(\mathrm{gr}\,\mathrm{Ev})$ with $\mathrm{gr}\,\mathrm{Ev}\in\Sigma$; so $\xi\mapsto\left(\left(\mathrm{gr}^{-1}\right)_*\xi\right)(\mathrm{Ev})$ is again one of the maps that define $B_{\mathrm{Law}}$, and the inverse of $\mathrm{Tg}$ is measurable.
Finally, $\mathrm{Tg}(\mathrm{Law}_{\mathrm{erg}})=\mathrm{Law}_{\mathrm{erg}}$ is (tr4).
\end{proof}

\subsection*{1. Proposition 45.6. The ergodic decomposition, in the form that is used}

\begin{proof}[Proof of Proposition~45.6\,(a), (a$'$), (a$''$)]
Let $\mathrm{lw}$ be a law that is invariant under $\mathbb R^3$, and let $I$ be the family of the events $\mathrm{Ev}\in\Sigma$ with $\tau_a\mathrm{Ev}=\mathrm{Ev}$ for all $a\in\mathbb R^3$.

\textbf{Step 1 (the kernel; the first sentence of (a)).}
The family $I$ is a $\sigma$-algebra contained in $\Sigma$.
Indeed, $\mathcal S'\in I$; and if $\mathrm{Ie}$ and $\mathrm{Ie}_1,\mathrm{Ie}_2,\dots$ belong to $I$ and $a\in\mathbb R^3$, then, since $\omega\mapsto\tau_a\omega$ is a bijection of $\mathcal S'$,
\[
\tau_a\left(\mathcal S'\setminus\mathrm{Ie}\right)=\mathcal S'\setminus\tau_a\mathrm{Ie}=\mathcal S'\setminus\mathrm{Ie},\qquad\tau_a\left(\bigcup_k\mathrm{Ie}_k\right)=\bigcup_k\tau_a\mathrm{Ie}_k=\bigcup_k\mathrm{Ie}_k .
\]
By (qc2)\,(iii) the measurable space $(\mathcal S',\Sigma)$ is standard Borel, that is, isomorphic to a Polish space with its Borel $\sigma$-algebra.
By (qc10)\,(ii) a Polish space with its Borel $\sigma$-algebra is isomorphic to a Borel subset of $[0,1]$ with the trace $\sigma$-algebra, and such an isomorphism carries the Borel subsets of the Polish space onto Borel subsets of $[0,1]$; consequently every Borel subset of a Polish space, endowed with the trace $\sigma$-algebra, is a Borel space as well.
In particular $(\mathcal S',\Sigma)$ is a Borel space.
Apply (qc10)\,(i) on the probability space $(\mathcal S',\Sigma,\mathrm{lw})$ to the two random elements $\mathrm{id}\colon(\mathcal S',\Sigma)\to(\mathcal S',\Sigma)$ and $\mathrm{id}\colon(\mathcal S',\Sigma)\to(\mathcal S',I)$; the second one is measurable because $I\subset\Sigma$, and the $\sigma$-algebra generated by it is $I$.
This yields a map $(\omega,\mathrm{Ev})\mapsto\mathrm{lw}_\omega(\mathrm{Ev})$ on $\mathcal S'\times\Sigma$ with the following properties.
\begin{itemize}
\item[(kn1)] For every $\omega\in\mathcal S'$, without exception, $\mathrm{Ev}\mapsto\mathrm{lw}_\omega(\mathrm{Ev})$ is a probability measure on $(\mathcal S',\Sigma)$.
\item[(kn2)] For every $\mathrm{Ev}\in\Sigma$ the function $\omega\mapsto\mathrm{lw}_\omega(\mathrm{Ev})$ is $I$-measurable.
\item[(kn3)] For every $\mathrm{Ev}\in\Sigma$, $\mathrm{lw}_\omega(\mathrm{Ev})=E_{\mathrm{lw}}\left[\mathbf 1_{\mathrm{Ev}}\mid I\right](\omega)$ for $\mathrm{lw}$-almost every $\omega$.
\end{itemize}
By (kn1) and (kn2) the map is a probability kernel from $(\mathcal S',I)$ to $(\mathcal S',\Sigma)$; with (kn3) this proves the first sentence of (a).
In all the following steps, $(\mathrm{lw}_\omega)$ denotes an arbitrary map with the properties (kn1)--(kn3).
For a function $\Phi$ on $\mathcal S'$ put
\[
\Phi^I(\omega):=\mathrm{lw}_\omega(\Phi)=\int\mathrm{lw}_\omega(\mathrm{d}\omega')\,\Phi(\omega')
\]
wherever the integral is defined.

\textbf{Step 2 ((a$'$), including its last sentence).}
\emph{Conditional expectations with values in $[0,\infty]$.}
Let $\Phi\colon\mathcal S'\to[0,\infty]$ be $\Sigma$-measurable.
As in the last sentence of Proposition~45.6\,(a$'$), $E_{\mathrm{lw}}[\Phi\mid I]$ stands for any $I$-measurable function $\mathrm{Gu}\colon\mathcal S'\to[0,\infty]$ with
\[
\int_{\mathrm{Ie}}\mathrm{lw}(\mathrm{d}\omega)\,\mathrm{Gu}(\omega)=\int_{\mathrm{Ie}}\mathrm{lw}(\mathrm{d}\omega)\,\Phi(\omega)\qquad\text{for all }\mathrm{Ie}\in I.
\]
Any two such functions $\mathrm{Gu}$, $\mathrm{Gu}'$ agree $\mathrm{lw}$-almost surely.
Indeed, for $k\ge1$ the event $\mathrm{Ie}_k:=\{\mathrm{Gu}<\mathrm{Gu}',\ \mathrm{Gu}\le k\}$ belongs to $I$, and $\int_{\mathrm{Ie}_k}\mathrm{lw}(\mathrm{d}\omega)\,\mathrm{Gu}(\omega)\le k$ is finite.
The integrals of $\mathrm{Gu}$ and of $\mathrm{Gu}'$ over $\mathrm{Ie}_k$ both equal the integral of $\Phi$ over $\mathrm{Ie}_k$; therefore $\int_{\mathrm{Ie}_k}\mathrm{lw}(\mathrm{d}\omega)\,\left(\mathrm{Gu}'(\omega)-\mathrm{Gu}(\omega)\right)=0$, while $\mathrm{Gu}'-\mathrm{Gu}>0$ on $\mathrm{Ie}_k$.
It follows that $\mathrm{lw}(\mathrm{Ie}_k)=0$, and that $\{\mathrm{Gu}<\mathrm{Gu}'\}=\bigcup_k\mathrm{Ie}_k$ is a null set.
The hypothesis is symmetric in $\mathrm{Gu}$ and $\mathrm{Gu}'$; so $\{\mathrm{Gu}'<\mathrm{Gu}\}$ is a null set as well.
For $\mathrm{lw}$-integrable $\Phi\ge0$ this is the conditional expectation of an integrable function.

\emph{Non-negative $\Phi$.}
Let first $\Phi=\sum_{p\le s}t_p\cdot\mathbf 1_{\mathrm{Ev}_p}$ with $t_p\ge0$ and $\mathrm{Ev}_p\in\Sigma$.
Then $\Phi^I(\omega)=\sum_pt_p\cdot\mathrm{lw}_\omega(\mathrm{Ev}_p)$ is $I$-measurable by (kn2), and for $\mathrm{Ie}\in I$, by (kn3),
\[
\int_{\mathrm{Ie}}\mathrm{lw}(\mathrm{d}\omega)\,\Phi^I(\omega)=\sum_pt_p\cdot\int_{\mathrm{Ie}}\mathrm{lw}(\mathrm{d}\omega)\,E_{\mathrm{lw}}\left[\mathbf 1_{\mathrm{Ev}_p}\mid I\right](\omega)=\sum_pt_p\cdot\mathrm{lw}(\mathrm{Ie}\cap\mathrm{Ev}_p)=\int_{\mathrm{Ie}}\mathrm{lw}(\mathrm{d}\omega)\,\Phi(\omega).
\]
A general $\Sigma$-measurable $\Phi\ge0$ is the pointwise limit of an increasing sequence $(\Phi_k)$ of functions of this kind.
Monotone convergence under each $\mathrm{lw}_\omega$ (kn1) gives $\Phi^I(\omega)=\lim_k\Phi_k^I(\omega)$ in $[0,\infty]$ for every $\omega$; hence $\Phi^I$ is $I$-measurable.
Monotone convergence under $\mathrm{lw}$, applied to $\mathbf 1_{\mathrm{Ie}}\cdot\Phi_k^I$ and to $\mathbf 1_{\mathrm{Ie}}\cdot\Phi_k$, gives
\begin{equation}\tag{1.1}
\int_{\mathrm{Ie}}\mathrm{lw}(\mathrm{d}\omega)\,\Phi^I(\omega)=\int_{\mathrm{Ie}}\mathrm{lw}(\mathrm{d}\omega)\,\Phi(\omega)\qquad\text{for all }\mathrm{Ie}\in I.
\end{equation}
So $\Phi^I$ is an $I$-measurable function $\mathrm{Gu}$ with the property displayed at the beginning of this step, that is, $\Phi^I=E_{\mathrm{lw}}[\Phi\mid I]$ in the sense described there; the choice $\mathrm{Ie}=\mathcal S'$ gives $\mathrm{lw}(\Phi)=\int\mathrm{lw}(\mathrm{d}\omega)\,\mathrm{lw}_\omega(\Phi)$.

\emph{Integrable $\Phi$.}
Let $\Phi\colon\mathcal S'\to\mathbb C$ be $\Sigma$-measurable with $\mathrm{lw}(|\Phi|)<\infty$.
By (1.1), $\int\mathrm{lw}(\mathrm{d}\omega)\,|\Phi|^I(\omega)=\mathrm{lw}(|\Phi|)$ is finite.
Hence the event $\mathrm{Ie}_\Phi:=\{\omega:\mathrm{lw}_\omega(|\Phi|)<\infty\}$, which belongs to $I$, satisfies $\mathrm{lw}(\mathrm{Ie}_\Phi)=1$.
On $\mathrm{Ie}_\Phi$ the function
\[
\Phi^I=\left((\operatorname{Re}\Phi)^+\right)^I-\left((\operatorname{Re}\Phi)^-\right)^I+i\cdot\left((\operatorname{Im}\Phi)^+\right)^I-i\cdot\left((\operatorname{Im}\Phi)^-\right)^I
\]
is defined, finite and $I$-measurable.
The four parts have finite integrals, and (1.1), applied to each of them, gives $\int_{\mathrm{Ie}}\mathrm{lw}(\mathrm{d}\omega)\,\Phi^I(\omega)=\int_{\mathrm{Ie}}\mathrm{lw}(\mathrm{d}\omega)\,\Phi(\omega)$ for all $\mathrm{Ie}\in I$.
Consequently $\Phi^I=E_{\mathrm{lw}}[\Phi\mid I]$ $\mathrm{lw}$-almost surely, and $\mathrm{lw}(\Phi)=\int\mathrm{lw}(\mathrm{d}\omega)\,\mathrm{lw}_\omega(\Phi)$.
If $\Phi'=\Phi$ $\mathrm{lw}$-almost everywhere, then $\int\mathrm{lw}(\mathrm{d}\omega)\,|\Phi-\Phi'|^I(\omega)=\mathrm{lw}(|\Phi-\Phi'|)=0$ by (1.1).
So $\mathrm{lw}_\omega(|\Phi-\Phi'|)=0$, and hence $\Phi'^I(\omega)=\Phi^I(\omega)$, for $\mathrm{lw}$-almost every $\omega$: the class of $\Phi^I$ depends only on the class of $\Phi$ in $L^1(\mathrm{lw})$.

\emph{Bounded $\Phi$.}
If $|\Phi|\le\mathrm{cf}$ everywhere, then $\left|\Phi^I(\omega)\right|\le\mathrm{cf}$ for every $\omega$, without exception, by (kn1), and $\mathrm{Ie}_\Phi=\mathcal S'$.

\emph{$I$-measurable $\Phi$.}
If $\Phi$ is bounded and $I$-measurable, then $E_{\mathrm{lw}}[\Phi\mid I]=\Phi$; therefore
\begin{equation}\tag{1.2}
\mathrm{lw}_\omega(\Phi)=\Phi(\omega)\qquad\text{for }\mathrm{lw}\text{-almost every }\omega.
\end{equation}

\textbf{Step 3 (the second sentence of (a)).}
Let $(\mathrm{lw}_\omega)$ and $(\mathrm{lw}'_\omega)$ both have the properties (kn1)--(kn3).
For $f\in\mathrm{Dq}$, Step~2 applied to the bounded function $\exp\left(iX(f)\right)$ gives
$\mathrm{lw}_\omega\left(\exp\left(iX(f)\right)\right)=E_{\mathrm{lw}}\left[\exp\left(iX(f)\right)\mid I\right](\omega)=\mathrm{lw}'_\omega\left(\exp\left(iX(f)\right)\right)$ for $\mathrm{lw}$-almost every $\omega$.
The event
\[
\left\{\omega:\mathrm{lw}_\omega\left(\exp\left(iX(f)\right)\right)=\mathrm{lw}'_\omega\left(\exp\left(iX(f)\right)\right)\text{ for all }f\in\mathrm{Dq}\right\}
\]
belongs to $I$, since it is defined by countably many $I$-measurable functions (Step~2), and it has $\mathrm{lw}$-measure $1$, since $\mathrm{Dq}$ is countable (qc2)\,(i).
For $\omega$ in this event, (qc3)\,(ii) gives $\mathrm{lw}_\omega=\mathrm{lw}'_\omega$.

\textbf{Step 4 (the measure $\nu$; it represents $\mathrm{lw}$; the second sentence of (a$''$)).}
The map $\omega\mapsto\mathrm{lw}_\omega$ from $(\mathcal S',I)$ to $(\mathrm{Law},B_{\mathrm{Law}})$ is measurable: $B_{\mathrm{Law}}$ is generated by the maps $\xi\mapsto\xi(\mathrm{Ev})$, and $\omega\mapsto\mathrm{lw}_\omega(\mathrm{Ev})$ is $I$-measurable (kn2).
Let $\nu$ be the image of $\mathrm{lw}$ under this map:
\begin{equation}\tag{1.3}
\nu(\mathrm{Ls})=\mathrm{lw}\{\omega:\mathrm{lw}_\omega\in\mathrm{Ls}\}\qquad\text{for }\mathrm{Ls}\in B_{\mathrm{Law}}.
\end{equation}
By Step~3, two kernels with (kn1)--(kn3) define the same measure $\nu$.
For $\mathrm{Ev}\in\Sigma$, the change of variables for the image measure and (1.1) with $\Phi=\mathbf 1_{\mathrm{Ev}}$, $\mathrm{Ie}=\mathcal S'$ give
\[
\int\nu(\mathrm{d}\xi)\,\xi(\mathrm{Ev})=\int\mathrm{lw}(\mathrm{d}\omega)\,\mathrm{lw}_\omega(\mathrm{Ev})=\mathrm{lw}(\mathrm{Ev}).
\]
This is the representation $\mathrm{lw}=\int\nu(\mathrm{d}\xi)\,\mathrm{lw}_\xi$.
Let a property define a set $\mathrm{Ls}\in B_{\mathrm{Law}}$.
By (1.3), $\nu(\mathrm{Ls})=1$ holds if and only if $\mathrm{lw}_\omega\in\mathrm{Ls}$ for $\mathrm{lw}$-almost every $\omega$.
This is the second sentence of (a$''$).

\textbf{Step 5 ($\mathrm{lw}_\omega$ is invariant under $\mathbb R^3$ for $\mathrm{lw}$-almost every $\omega$).}
For $\mathrm{Ie}\in I$ and $a\in\mathbb R^3$ one has
$\left(U(a)\mathbf 1_{\mathrm{Ie}}\right)(\omega)=\mathbf 1_{\mathrm{Ie}}(\tau_{-a}\omega)=\mathbf 1_{\tau_a\mathrm{Ie}}(\omega)=\mathbf 1_{\mathrm{Ie}}(\omega)$.
For a bounded measurable $\mathrm{Gb}$,
\[
\mathrm{lw}\left(U(a)\mathrm{Gb}\right)=\int\mathrm{lw}(\mathrm{d}\omega)\,\mathrm{Gb}(\tau_{-a}\omega)=\left((\tau_{-a})_*\mathrm{lw}\right)(\mathrm{Gb})=\mathrm{lw}(\mathrm{Gb}).
\]
Let $f\in\mathrm{Dq}$.
Since $U(a)$ is multiplicative and $U(a)\exp\left(iX(f)\right)=\exp\left(iX(\tau_a\cdot f)\right)$ (qc1),
\[
\mathrm{lw}\left(\mathbf 1_{\mathrm{Ie}}\cdot\exp\left(iX(\tau_a\cdot f)\right)\right)=\mathrm{lw}\left(U(a)\left(\mathbf 1_{\mathrm{Ie}}\cdot\exp\left(iX(f)\right)\right)\right)=\mathrm{lw}\left(\mathbf 1_{\mathrm{Ie}}\cdot\exp\left(iX(f)\right)\right)\qquad\text{for all }\mathrm{Ie}\in I.
\]
Hence $E_{\mathrm{lw}}\left[\exp\left(iX(\tau_a\cdot f)\right)\mid I\right]=E_{\mathrm{lw}}\left[\exp\left(iX(f)\right)\mid I\right]$, and Step~2 gives $\mathrm{lw}_\omega\left(\exp\left(iX(\tau_a\cdot f)\right)\right)=\mathrm{lw}_\omega\left(\exp\left(iX(f)\right)\right)$ for $\mathrm{lw}$-almost every $\omega$.
The pairs $(a,f)\in\mathbb Q^3\times\mathrm{Dq}$ are countably many.
So for $\mathrm{lw}$-almost every $\omega$ the equality holds for all of them simultaneously, and by (qc7) the law $\mathrm{lw}_\omega$ is invariant under $\mathbb R^3$.

\textbf{Step 6 ($\mathrm{lw}_\omega$ is $\mathbb R^3$-ergodic for $\mathrm{lw}$-almost every $\omega$).}
Enumerate $\mathrm{Aq}=\{\mathrm{Ga}_1,\mathrm{Ga}_2,\dots\}$ (qc3).
Every $\mathrm{Ga}_p$ is a bounded $\Sigma$-measurable function; put $\mathrm{cf}_p:=\sup|\mathrm{Ga}_p|$.
By (qc4)\,(i), $\mathrm{Avg}_{\mathrm{sl}}\mathrm{Ga}_p$ is a measurable function on $\mathcal S'$, defined without reference to any law, with $\left|\mathrm{Avg}_{\mathrm{sl}}\mathrm{Ga}_p\right|\le\mathrm{cf}_p$.
By Step~2, $\mathrm{Ga}_p^I$ is $I$-measurable with $\left|\mathrm{Ga}_p^I\right|\le\mathrm{cf}_p$ everywhere.
By (qc4)\,(iii), (iv) and Step~2,
\[
\left\|\mathrm{Avg}_{\mathrm{sl}}\mathrm{Ga}_p-\mathrm{Ga}_p^I\right\|_{L^2(\mathrm{lw})}\to0\qquad(\mathrm{sl}\to\infty).
\]
Choose $\mathrm{sl}_1<\mathrm{sl}_2<\cdots$ with $\mathrm{sl}_k\ge k$ and
\[
\left\|\mathrm{Avg}_{\mathrm{sl}_k}\mathrm{Ga}_p-\mathrm{Ga}_p^I\right\|^2_{L^2(\mathrm{lw})}\le2^{-k}\qquad\text{for all }p\le k.
\]
Put $\Phi_p:=\sum_{k\ge1}\left|\mathrm{Avg}_{\mathrm{sl}_k}\mathrm{Ga}_p-\mathrm{Ga}_p^I\right|^2\colon\mathcal S'\to[0,\infty]$.
By monotone convergence,
\[
\mathrm{lw}(\Phi_p)\le\sum_{k<p}4\cdot\mathrm{cf}_p^2+\sum_{k\ge p}2^{-k}<\infty .
\]
Two consequences follow.
\begin{itemize}
\item[(er1)] By (1.1), $\int\mathrm{lw}(\mathrm{d}\omega)\,\mathrm{lw}_\omega(\Phi_p)=\mathrm{lw}(\Phi_p)$ is finite.
Hence, for $\mathrm{lw}$-almost every $\omega$ and all $p$, the sum $\sum_k\mathrm{lw}_\omega\left(\left|\mathrm{Avg}_{\mathrm{sl}_k}\mathrm{Ga}_p-\mathrm{Ga}_p^I\right|^2\right)$, which equals $\mathrm{lw}_\omega(\Phi_p)$ by monotone convergence under $\mathrm{lw}_\omega$, is finite, and therefore
\[
\mathrm{lw}_\omega\left(\left|\mathrm{Avg}_{\mathrm{sl}_k}\mathrm{Ga}_p-\mathrm{Ga}_p^I\right|^2\right)\to0\qquad(k\to\infty).
\]
\item[(er2)] $\Phi_p$ is finite $\mathrm{lw}$-almost everywhere.
Hence the event
\[
\mathrm{Ev}_1:=\left\{\omega':\mathrm{Avg}_{\mathrm{sl}_k}\mathrm{Ga}_p(\omega')\to\mathrm{Ga}_p^I(\omega')\text{ as }k\to\infty,\text{ for all }p\right\},
\]
which belongs to $\Sigma$, satisfies $\mathrm{lw}(\mathrm{Ev}_1)=1$.
\end{itemize}
The bounded functions $\mathrm{Ga}_p^I$ and $\left|\mathrm{Ga}_p^I\right|^2$ are $I$-measurable.
By (1.2), for $\mathrm{lw}$-almost every $\omega$ and all $p$ one has $\mathrm{lw}_\omega\left(\mathrm{Ga}_p^I\right)=\mathrm{Ga}_p^I(\omega)$ and $\mathrm{lw}_\omega\left(\left|\mathrm{Ga}_p^I\right|^2\right)=\left|\mathrm{Ga}_p^I(\omega)\right|^2$, whence
\[
\mathrm{lw}_\omega\left(\left|\mathrm{Ga}_p^I-\mathrm{Ga}_p^I(\omega)\right|^2\right)=\mathrm{lw}_\omega\left(\left|\mathrm{Ga}_p^I\right|^2\right)-2\operatorname{Re}\left(\overline{\mathrm{Ga}_p^I(\omega)}\cdot\mathrm{lw}_\omega\left(\mathrm{Ga}_p^I\right)\right)+\left|\mathrm{Ga}_p^I(\omega)\right|^2=0.
\]
Together with (er1) and the triangle inequality in $L^2(\mathrm{lw}_\omega)$ this shows that, for $\mathrm{lw}$-almost every $\omega$ and all $p$,
\[
\mathrm{Avg}_{\mathrm{sl}_k}\mathrm{Ga}_p\to\mathrm{Ga}_p^I(\omega)\cdot1=\mathrm{lw}_\omega(\mathrm{Ga}_p)\cdot1\qquad\text{in }L^2(\mathrm{lw}_\omega).
\]
On the other hand, Step~5 and (qc4)\,(iii) for the law $\xi=\mathrm{lw}_\omega$ give, for $\mathrm{lw}$-almost every $\omega$ and all $p$,
\[
\mathrm{Avg}_{\mathrm{sl}_k}\mathrm{Ga}_p\to P^{123}_{\mathrm{lw}_\omega}\mathrm{Ga}_p\qquad\text{in }L^2(\mathrm{lw}_\omega).
\]
Hence $P^{123}_{\mathrm{lw}_\omega}\mathrm{Ga}_p=\mathrm{lw}_\omega(\mathrm{Ga}_p)\cdot1$ for all $p$.
The maps $\mathrm{Ga}\mapsto P^{123}_{\mathrm{lw}_\omega}\mathrm{Ga}$ and $\mathrm{Ga}\mapsto\mathrm{lw}_\omega(\mathrm{Ga})\cdot1$ are bounded and linear on $L^2(\mathrm{lw}_\omega)$, and $\mathrm{Aq}$ is dense in $L^2(\Sigma_{\mathbb R^3},\mathrm{lw}_\omega)=L^2(\mathrm{lw}_\omega)$ by (qc3)\,(i) and (qc2)\,(ii).
So $P^{123}_{\mathrm{lw}_\omega}\mathrm{Ga}=\mathrm{lw}_\omega(\mathrm{Ga})\cdot1$ for all $\mathrm{Ga}\in L^2(\mathrm{lw}_\omega)$: every vector that is invariant under all $U(a)$ is constant.
Thus $\mathrm{lw}_\omega\in\mathrm{Law}_{\mathrm{erg}}$ for $\mathrm{lw}$-almost every $\omega$.
Since $\mathrm{Law}_{\mathrm{erg}}\in B_{\mathrm{Law}}$ (qc7), (1.3) gives $\nu(\mathrm{Law}_{\mathrm{erg}})=1$.

\textbf{Step 7 (uniqueness of $\nu$; the last sentence of (a$''$)).}
Let $\nu'$ be a probability measure on $(\mathrm{Law},B_{\mathrm{Law}})$ with $\nu'(\mathrm{Law}_{\mathrm{erg}})=1$ and $\mathrm{lw}(\mathrm{Ev})=\int\nu'(\mathrm{d}\xi)\,\xi(\mathrm{Ev})$ for all $\mathrm{Ev}\in\Sigma$.
For the event $\mathrm{Ev}_1$ of (er2) one has $\int\nu'(\mathrm{d}\xi)\,\xi(\mathrm{Ev}_1)=\mathrm{lw}(\mathrm{Ev}_1)=1$ and $\xi(\mathrm{Ev}_1)\le1$.
Hence the set $\mathrm{Ls}_1:=\{\xi\in\mathrm{Law}_{\mathrm{erg}}:\xi(\mathrm{Ev}_1)=1\}$, which belongs to $B_{\mathrm{Law}}$, satisfies $\nu'(\mathrm{Ls}_1)=1$.
Let $\xi\in\mathrm{Ls}_1$.
As $\xi$ is ergodic, $P^{123}_\xi\mathrm{Ga}_p=\left\langle1,\mathrm{Ga}_p\right\rangle_\xi\cdot1=\xi(\mathrm{Ga}_p)\cdot1$; by (qc4)\,(iii), $\mathrm{Avg}_{\mathrm{sl}_k}\mathrm{Ga}_p\to\xi(\mathrm{Ga}_p)\cdot1$ in $L^2(\xi)$, and therefore $\xi$-almost everywhere along a subsequence.
On $\mathrm{Ev}_1$, that is, $\xi$-almost everywhere, $\mathrm{Avg}_{\mathrm{sl}_k}\mathrm{Ga}_p\to\mathrm{Ga}_p^I$.
So $\mathrm{lw}_{\omega'}(\mathrm{Ga}_p)=\mathrm{Ga}_p^I(\omega')=\xi(\mathrm{Ga}_p)$ for $\xi$-almost every $\omega'$.
There are countably many $p$, and $\mathrm{Aq}$ contains the functions $\exp\left(iX(f)\right)$, $f\in\mathrm{Dq}$; so (qc3)\,(ii) gives
\begin{equation}\tag{1.4}
\text{for every }\xi\in\mathrm{Ls}_1:\qquad\mathrm{lw}_{\omega'}=\xi\quad\text{for }\xi\text{-almost every }\omega'.
\end{equation}
Let $\mathrm{Ls}\in B_{\mathrm{Law}}$.
The event $\{\omega':\mathrm{lw}_{\omega'}\in\mathrm{Ls}\}$ belongs to $I$ (Step~4), and (1.3), the representation of $\mathrm{lw}$ by $\nu'$ and (1.4) give
\[
\nu(\mathrm{Ls})=\mathrm{lw}\{\omega':\mathrm{lw}_{\omega'}\in\mathrm{Ls}\}=\int\nu'(\mathrm{d}\xi)\,\xi\{\omega':\mathrm{lw}_{\omega'}\in\mathrm{Ls}\}=\int_{\mathrm{Ls}_1}\nu'(\mathrm{d}\xi)\,\mathbf 1_{\mathrm{Ls}}(\xi)=\nu'(\mathrm{Ls}).
\]
Together with Steps~4 and~6 this proves the first sentence of (a$''$).

Let now $\mathrm{lw}$ be $\mathbb R^3$-ergodic.
The point mass at $\mathrm{lw}$, $\mathrm{Ls}\mapsto\mathbf 1_{\mathrm{Ls}}(\mathrm{lw})$, is a probability measure on $(\mathrm{Law},B_{\mathrm{Law}})$; it gives the measure $1$ to $\mathrm{Law}_{\mathrm{erg}}$, and the integral of $\xi(\mathrm{Ev})$ with respect to it is $\mathrm{lw}(\mathrm{Ev})$.
By the uniqueness just proved, it is $\nu$.
(By (La3), $\{\mathrm{lw}\}\in B_{\mathrm{Law}}$, and $\nu(\{\mathrm{lw}\})=1$.)
\end{proof}

\paragraph{Step 8 ((b)).}
Assume that $\mathrm{lw}$ is invariant under $\theta_j$ and reflection positive at $\{y_j=0\}$.
Being invariant under $\mathbb{R}^3$, the law $\mathrm{lw}$ is in particular invariant under $\mathbb{R}e_j$.
For $\mathrm{Ie}\in I$ the function $\mathbf{1}_{\mathrm{Ie}}$ is bounded and non-negative, and $U(se_j)\mathbf{1}_{\mathrm{Ie}}=\mathbf{1}_{\mathrm{Ie}}$ for every real $s$ (Step 5).
Hence (qc8) applies with $\xi=\mathrm{lw}$ and $\mathrm{Gb}=\mathbf{1}_{\mathrm{Ie}}$.

\textit{Invariance.}
Let $f\in\mathrm{Dq}$.
By (qc1), $\hat\theta_j\exp\left(iX(f)\right)=\exp\left(iX(\theta_j\cdot f)\right)$, and the first assertion of (qc8) gives
\[
\mathrm{lw}\left(\mathbf{1}_{\mathrm{Ie}}\cdot\exp\left(iX(\theta_j\cdot f)\right)\right)=\mathrm{lw}\left(\mathbf{1}_{\mathrm{Ie}}\cdot\exp\left(iX(f)\right)\right)\qquad(\mathrm{Ie}\in I).
\]
Consequently $E_{\mathrm{lw}}\left[\exp\left(iX(\theta_j\cdot f)\right)\mid I\right]=E_{\mathrm{lw}}\left[\exp\left(iX(f)\right)\mid I\right]$, and Step 2 yields
\[
\mathrm{lw}_\omega\left(\exp\left(iX(\theta_j\cdot f)\right)\right)=\mathrm{lw}_\omega\left(\exp\left(iX(f)\right)\right)\qquad\text{for $\mathrm{lw}$-almost every }\omega.
\]
As $\mathrm{Dq}$ is countable, (qc7) shows that $\mathrm{lw}_\omega$ is invariant under $\theta_j$ for $\mathrm{lw}$-almost every $\omega$.

\textit{Positivity.}
Let $\mathrm{Ga}\in\mathrm{Aq}(\{y_j>0\})$.
Being a finite linear combination of functions $\exp\left(iX(f)\right)$ with $f\in C_c^\infty(\{y_j>0\};\Lambda^2)$, the function $\mathrm{Ga}$ is bounded and $\Sigma_{\{y_j>0\}}$-measurable.
Put
\[
\mathrm{Gk}:=\left(\overline{\hat\theta_j\mathrm{Ga}}\cdot\mathrm{Ga}\right)^I;
\]
by Step 2 this is a bounded $I$-measurable complex function.
The identity (1.1), in its form for integrable functions, and the second assertion of (qc8) give
\[
\mathrm{lw}\left(\mathbf{1}_{\mathrm{Ie}}\cdot\mathrm{Gk}\right)=\mathrm{lw}\left(\mathbf{1}_{\mathrm{Ie}}\cdot\overline{\hat\theta_j\mathrm{Ga}}\cdot\mathrm{Ga}\right)\ge0\qquad(\mathrm{Ie}\in I);
\]
in particular the left side is real.
Take first $\mathrm{Ie}=\{\operatorname{Im}\mathrm{Gk}>0\}$, an event of $I$.
Then $0=\operatorname{Im}\mathrm{lw}\left(\mathbf{1}_{\mathrm{Ie}}\cdot\mathrm{Gk}\right)=\mathrm{lw}\left(\mathbf{1}_{\mathrm{Ie}}\cdot\operatorname{Im}\mathrm{Gk}\right)$, while $\operatorname{Im}\mathrm{Gk}>0$ on $\mathrm{Ie}$; hence $\mathrm{lw}(\mathrm{Ie})=0$.
For $\mathrm{Ie}=\{\operatorname{Im}\mathrm{Gk}<0\}$ one has likewise $0=\mathrm{lw}\left(\mathbf{1}_{\mathrm{Ie}}\cdot\operatorname{Im}\mathrm{Gk}\right)$ with $\operatorname{Im}\mathrm{Gk}<0$ on $\mathrm{Ie}$, so that $\mathrm{lw}(\mathrm{Ie})=0$.
For $\mathrm{Ie}=\{\operatorname{Re}\mathrm{Gk}<0\}$ one has $0\le\operatorname{Re}\mathrm{lw}\left(\mathbf{1}_{\mathrm{Ie}}\cdot\mathrm{Gk}\right)=\mathrm{lw}\left(\mathbf{1}_{\mathrm{Ie}}\cdot\operatorname{Re}\mathrm{Gk}\right)$ with $\operatorname{Re}\mathrm{Gk}<0$ on $\mathrm{Ie}$, and again $\mathrm{lw}(\mathrm{Ie})=0$.
It follows that
\[
\mathrm{lw}_\omega\left(\overline{\hat\theta_j\mathrm{Ga}}\cdot\mathrm{Ga}\right)=\mathrm{Gk}(\omega)\ge0\qquad\text{for $\mathrm{lw}$-almost every }\omega.
\]
The set $\mathrm{Aq}(\{y_j>0\})$ is countable.
Therefore, for $\mathrm{lw}$-almost every $\omega$, the law $\mathrm{lw}_\omega$ satisfies (qc6)\,(2), and by (qc7) it is reflection positive at $\{y_j=0\}$.

\textit{Conclusion.}
Let $\mathrm{Ls}_b$ be the set of the laws that are invariant under $\mathbb{R}^3$ and under $\theta_j$ and reflection positive at $\{y_j=0\}$.
By (qc7) it belongs to $B_{\mathrm{Law}}$.
Step 5 and the two preceding paragraphs show that $\mathrm{lw}_\omega\in\mathrm{Ls}_b$ for $\mathrm{lw}$-almost every $\omega$, and (1.3) gives $\nu(\mathrm{Ls}_b)=1$.
Every $\xi\in\mathrm{Ls}_b$ is invariant under $\theta_j$ and under $\mathbb{R}e_j$; by the implication (1)\,$\Rightarrow$\,(5) of (qc6), $\xi$ is reflection positive at every plane $\{y_j=\mathrm{po}\}$, $\mathrm{po}\in\mathbb{R}$.

\paragraph{Step 9 ((c)).}
Assume that $\mathrm{lw}$ has the property (G).
Fix $f\in\mathcal{S}$ and put
\[
\mathrm{Ga}:=\exp X(f)>0,\qquad\mathrm{bd}:=\exp\left(\frac{1}{2}\|f\|_2^2\right),\qquad\mathrm{ov}(y):=\langle f,\tau_y\cdot f\rangle\quad(y\in\mathbb{R}^3).
\]
For every integer $n\ge1$, the property (G) for the test form $nf$ gives
\[
\|\mathrm{Ga}\|_{L^n(\mathrm{lw})}^n=\mathrm{lw}\left(\exp X(nf)\right)\le\exp\left(\frac{n^2}{2}\|f\|_2^2\right).
\]
By Step 2 the function $\mathrm{Ga}^I\colon\mathcal{S}'\to[0,\infty]$ is $I$-measurable, and by (1.1) with $\Phi=\mathrm{Ga}$ its integrals over all $\mathrm{Ie}\in I$ equal those of $\mathrm{Ga}$.
Thus $\mathrm{Ga}^I$ represents $E_{\mathrm{lw}}[\mathrm{Ga}\mid I]$ in the sense of the last sentence of Proposition~45.6\,(a$'$), as explained at the beginning of Step 2.

\textit{(9a) The function $\mathrm{ov}$.}
For $a,b\in\mathbb{R}^3$ one has
\[
\langle\tau_a\cdot f,\tau_b\cdot f\rangle=\int\mathrm{d}y\,f(y-a)\cdot f(y-b)=\mathrm{ov}(b-a).
\]
The Cauchy--Schwarz inequality and (qc1) give $|\mathrm{ov}(y)|\le\|f\|_2^2$.
For $y,y'\in\mathbb{R}^3$ at least one of the inequalities $|y'|_2\ge|y|_2/2$ and $|y'-y|_2\ge|y|_2/2$ holds.
Split the integral $\mathrm{ov}(y)=\int\mathrm{d}y'\,f(y')\cdot f(y'-y)$ into the part over $\{|y'|_2\ge|y|_2/2\}$ and the part over $\{|y'|_2<|y|_2/2\}$, and apply the Cauchy--Schwarz inequality to each part.
This gives
\[
|\mathrm{ov}(y)|\le2\|f\|_2\cdot\left\|f\cdot\mathbf{1}[\,|\cdot|_2\ge|y|_2/2\,]\right\|_2,
\]
and by dominated convergence the right side tends to $0$ as $|y|_2\to\infty$.

\textit{(9b) Moments of the averages.}
Let $n\ge2$ and $\mathrm{sl}>0$.
By (qc1), $\mathrm{Ga}(\tau_{-a}\omega)=\left(U(a)\mathrm{Ga}\right)(\omega)=\exp\left(X(\tau_a\cdot f)(\omega)\right)$, and by (qc4)\,(i) the map $(a,\omega)\mapsto\mathrm{Ga}(\tau_{-a}\omega)$ is measurable.
Since $\omega$ is linear,
\[
\prod_{p=1}^n\mathrm{Ga}(\tau_{-a_p}\omega)=\exp X\left(\sum_{p=1}^n\tau_{a_p}\cdot f\right)(\omega).
\]
Tonelli's theorem, used first for the $n$-th power of the integral over $[0,\mathrm{sl}]^3$ and then for the interchange of $\mathrm{lw}$ with the integral over $([0,\mathrm{sl}]^3)^n$, yields
\[
\mathrm{lw}\left((\mathrm{Avg}_{\mathrm{sl}}\mathrm{Ga})^n\right)=\mathrm{sl}^{-3n}\cdot\int_{([0,\mathrm{sl}]^3)^n}\mathrm{d}a_1\cdots\mathrm{d}a_n\,\mathrm{lw}\left(\exp X\left(\sum_{p=1}^n\tau_{a_p}\cdot f\right)\right).
\]
By (G) for the test form $\sum_p\tau_{a_p}\cdot f$ and by (9a),
\[
\mathrm{lw}\left(\exp X\left(\sum_p\tau_{a_p}\cdot f\right)\right)\le\exp\left(\frac{1}{2}\left\|\sum_p\tau_{a_p}\cdot f\right\|_2^2\right)=\exp\left(\frac{n}{2}\|f\|_2^2+\frac{1}{2}\sum_{p\ne q}\mathrm{ov}(a_q-a_p)\right)\le\exp\left(\frac{n^2}{2}\|f\|_2^2\right).
\]
Let $\varepsilon>0$, and choose, by (9a), a number $\delta>0$ such that $|\mathrm{ov}(y)|\le\varepsilon$ whenever $|y|_\infty\ge\delta$.
The set of the points $(a_1,\dots,a_n)\in([0,\mathrm{sl}]^3)^n$ with $|a_p-a_q|_\infty<\delta$ for some pair $p<q$ has Lebesgue measure at most $\frac{n(n-1)}{2}\cdot(2\delta)^3\cdot\mathrm{sl}^{3(n-1)}$: for a fixed pair, $a_q$ and the remaining $n-2$ points vary in $[0,\mathrm{sl}]^3$, and $a_p$ varies in a cube of side $2\delta$.
Outside this set the exponent is at most $\frac{n}{2}\|f\|_2^2+\frac{n(n-1)\varepsilon}{2}$.
Hence
\begin{equation}\tag{1.5}
\mathrm{lw}\left((\mathrm{Avg}_{\mathrm{sl}}\mathrm{Ga})^n\right)\le\mathrm{bd}^n\cdot\exp\left(\frac{n(n-1)\varepsilon}{2}\right)+\frac{n(n-1)}{2}\cdot(2\delta)^3\cdot\mathrm{sl}^{-3}\cdot\exp\left(\frac{n^2}{2}\|f\|_2^2\right).
\end{equation}

\textit{(9c) The limit $\mathrm{sl}\to\infty$.}
The function $\mathrm{Ga}$ lies in $L^2(\mathrm{lw})$ and in $L^{2n}(\mathrm{lw})$.
By (qc4)\,(iii), (iv) and Step 2, $\mathrm{Avg}_{\mathrm{sl}}\mathrm{Ga}\to\mathrm{Ga}^I$ in $L^2(\mathrm{lw})$ as $\mathrm{sl}\to\infty$, and $\|\mathrm{Avg}_{\mathrm{sl}}\mathrm{Ga}\|_{L^{2n}(\mathrm{lw})}\le\|\mathrm{Ga}\|_{L^{2n}(\mathrm{lw})}$.
For every $\omega$, H\"older's inequality under the probability measure $\mathrm{lw}_\omega$ (kn1) gives
\[
\left(\mathrm{Ga}^I(\omega)\right)^{2n}=\left(\mathrm{lw}_\omega(\mathrm{Ga})\right)^{2n}\le\mathrm{lw}_\omega\left(\mathrm{Ga}^{2n}\right),
\]
and $\int\mathrm{d}\mathrm{lw}(\omega)\,\mathrm{lw}_\omega\left(\mathrm{Ga}^{2n}\right)=\mathrm{lw}\left(\mathrm{Ga}^{2n}\right)$ by (1.1).
Thus $\|\mathrm{Ga}^I\|_{L^{2n}(\mathrm{lw})}\le\|\mathrm{Ga}\|_{L^{2n}(\mathrm{lw})}$, and Minkowski's inequality gives $\|\mathrm{Avg}_{\mathrm{sl}}\mathrm{Ga}-\mathrm{Ga}^I\|_{L^{2n}(\mathrm{lw})}\le2\|\mathrm{Ga}\|_{L^{2n}(\mathrm{lw})}$.
Put $\mathrm{hx}:=1/(n-1)$; then $\frac{1}{n}=\frac{\mathrm{hx}}{2}+\frac{1-\mathrm{hx}}{2n}$.
Let $n>2$ and $\mathrm{Gb}:=\mathrm{Avg}_{\mathrm{sl}}\mathrm{Ga}-\mathrm{Ga}^I$.
Writing $|\mathrm{Gb}|^n=|\mathrm{Gb}|^{\mathrm{hx}\cdot n}\cdot|\mathrm{Gb}|^{(1-\mathrm{hx})n}$ and applying H\"older's inequality with the exponents $2/(\mathrm{hx}\cdot n)$ and $2/(1-\mathrm{hx})$, whose reciprocals add up to $n\left(\frac{\mathrm{hx}}{2}+\frac{1-\mathrm{hx}}{2n}\right)=1$, one obtains $\|\mathrm{Gb}\|_{L^n(\mathrm{lw})}\le\|\mathrm{Gb}\|_{L^2(\mathrm{lw})}^{\mathrm{hx}}\cdot\|\mathrm{Gb}\|_{L^{2n}(\mathrm{lw})}^{1-\mathrm{hx}}$; therefore
\[
\left\|\mathrm{Avg}_{\mathrm{sl}}\mathrm{Ga}-\mathrm{Ga}^I\right\|_{L^n(\mathrm{lw})}\le\left\|\mathrm{Avg}_{\mathrm{sl}}\mathrm{Ga}-\mathrm{Ga}^I\right\|_{L^2(\mathrm{lw})}^{\mathrm{hx}}\cdot\left(2\|\mathrm{Ga}\|_{L^{2n}(\mathrm{lw})}\right)^{1-\mathrm{hx}}\to0\qquad(\mathrm{sl}\to\infty).
\]
For $n=2$ one has $\mathrm{hx}=1$, and the two sides of this inequality coincide.
By Minkowski's inequality, $\|\mathrm{Avg}_{\mathrm{sl}}\mathrm{Ga}\|_{L^n(\mathrm{lw})}\to\|\mathrm{Ga}^I\|_{L^n(\mathrm{lw})}$.
Letting first $\mathrm{sl}\to\infty$ in (1.5) and then $\varepsilon\to0$, one finds
\[
\|\mathrm{Ga}^I\|_{L^n(\mathrm{lw})}\le\mathrm{bd}\qquad\text{for every }n\ge2.
\]
For $t>0$ Chebyshev's inequality gives
\[
\mathrm{lw}\left\{\mathrm{Ga}^I\ge(1+t)\cdot\mathrm{bd}\right\}\le\left((1+t)\cdot\mathrm{bd}\right)^{-n}\cdot\mathrm{lw}\left((\mathrm{Ga}^I)^n\right)\le(1+t)^{-n},
\]
and the right side tends to $0$ as $n\to\infty$; so $\mathrm{lw}\left\{\mathrm{Ga}^I\ge(1+t)\cdot\mathrm{bd}\right\}=0$ for every $t>0$.
The event $\left\{\mathrm{Ga}^I>\mathrm{bd}\right\}$ is the union over $k\ge1$ of the events $\left\{\mathrm{Ga}^I\ge(1+1/k)\cdot\mathrm{bd}\right\}$, a countable union of $\mathrm{lw}$-null sets.
Therefore $\mathrm{Ga}^I\le\mathrm{bd}$ holds $\mathrm{lw}$-almost surely, that is,
\[
\mathrm{lw}_\omega\left(\exp X(f)\right)\le\exp\left(\frac{1}{2}\|f\|_2^2\right)\qquad\text{for $\mathrm{lw}$-almost every }\omega.
\]

\textit{(9d) All $f$ simultaneously.}
Applying (9c) to each of the countably many $f\in\mathrm{Dq}$, one obtains for $\mathrm{lw}$-almost every $\omega$ the inequality $\mathrm{lw}_\omega\left(\exp X(f)\right)\le\exp\left(\frac{1}{2}\|f\|_2^2\right)$ for all $f\in\mathrm{Dq}$ at once; by (qc7) the law $\mathrm{lw}_\omega$ then has (G), for all $f\in\mathcal{S}$.
By (qc7) the set of the laws with (G) belongs to $B_{\mathrm{Law}}$, and by (1.3) its $\nu$-measure is $1$.

\paragraph{Step 10 ((c$'$)).}
Let $\mathrm{Ev}\in\Sigma$ with $\mathrm{lw}(\mathrm{Ev})=1$.
The set $\mathrm{Ls}:=\{\xi:\xi(\mathrm{Ev})=1\}$ belongs to $B_{\mathrm{Law}}$.
By Step 4, $\int\mathrm{d}\nu(\xi)\,\xi(\mathrm{Ev})=\mathrm{lw}(\mathrm{Ev})=1$, whereas $\xi(\mathrm{Ev})\le1$ for every $\xi$.
Hence $\nu(\mathrm{Ls})=1$.

\paragraph{Step 10$'$ ((d0)).}
The three sentences of (d0) are the facts (La1), (La2), (La3), which are proved in 0.4.

\paragraph{Step 11 ((d), including its assertions on the representative of $\mathrm{Ga}$).}
Let $\mathrm{Ga}_0$ be a real $\Sigma$-measurable representative of $\mathrm{Ga}$.
By the Cauchy--Schwarz inequality, $\mathrm{lw}(|\mathrm{Ga}_0|)\le\|\mathrm{Ga}\|_{L^2(\mathrm{lw})}<\infty$.
By (La2) the set
\[
\mathrm{Ls}_d:=\left\{\xi\in\mathrm{Law}(\mathrm{Ga}_0):|\xi(\mathrm{Ga}_0)|\le\mathrm{cf}\right\}
\]
belongs to $B_{\mathrm{Law}}$.
By Step 2, for $\mathrm{lw}$-almost every $\omega$ the number $\mathrm{lw}_\omega(|\mathrm{Ga}_0|)$ is finite, i.e. $\mathrm{lw}_\omega\in\mathrm{Law}(\mathrm{Ga}_0)$, and $\mathrm{lw}_\omega(\mathrm{Ga}_0)=E_{\mathrm{lw}}[\mathrm{Ga}\mid I](\omega)$, a number of modulus at most $\mathrm{cf}$ by hypothesis.
So $\mathrm{lw}_\omega\in\mathrm{Ls}_d$ for $\mathrm{lw}$-almost every $\omega$, and (1.3) gives $\nu(\mathrm{Ls}_d)=1$: for $\nu$-almost every $\xi$ the function $\mathrm{Ga}_0$ is $\mathrm{lw}_\xi$-integrable and $|\mathrm{lw}_\xi(\mathrm{Ga}_0)|\le\mathrm{cf}$.
Let $\mathrm{Ga}_1$ be another such representative.
Then $\mathrm{lw}\{\mathrm{Ga}_0\ne\mathrm{Ga}_1\}=0$, and Step 10 gives $\mathrm{lw}_\xi\{\mathrm{Ga}_0\ne\mathrm{Ga}_1\}=0$, hence $\mathrm{lw}_\xi(\mathrm{Ga}_1)=\mathrm{lw}_\xi(\mathrm{Ga}_0)$, for $\nu$-almost every $\xi$.
\qed

\begin{named}{Remark}[one axis]
Let $\xi$ be a law with (G) that is invariant under $\mathbb{R}e_j$, and let $f\in\mathcal{S}$.
Then $P^{j}_{\xi}X(f)=0$.
In the proof of \cite[Theorem~8.3(iv)]{QEDI} this identity is obtained from the decomposition of \cite[Theorem~8.1]{QEDI} for the index set $\{j\}$; it follows from \cite[Lemma~8.13(iii)]{QEDI} and \cite[Lemma~8.15(i)]{QEDI} alone, without a decomposition.
\end{named}

\begin{proof}
By (qc5), $X(f)\in L^2(\xi)$.
The hypothesis of \cite[Lemma~8.13]{QEDI} for the index set $J=\{j\}$ is the invariance under $\mathbb{R}e_j$; by \cite[Lemma~8.13(iii)]{QEDI} for this index set, $P^{j}_{\xi}X(f)$ is the limit in $L^2(\xi)$, as $\mathrm{sl}\to\infty$, of the functions
\[
\omega\mapsto\frac{1}{\mathrm{sl}}\int_0^{\mathrm{sl}}\mathrm{d}s\,X(f)(\tau_{-se_j}\omega)=\frac{1}{\mathrm{sl}}\int_0^{\mathrm{sl}}\mathrm{d}s\,X(\tau_{se_j}\cdot f)(\omega)=X(f_{\mathrm{sl}})(\omega),\qquad f_{\mathrm{sl}}:=\frac{1}{\mathrm{sl}}\int_0^{\mathrm{sl}}\mathrm{d}s\,\tau_{se_j}\cdot f\in\mathcal{S}.
\]
Here $s\mapsto\tau_{se_j}\cdot f$ is a continuous map of $[0,\mathrm{sl}]$ into $\mathcal{S}$, so that the Riemann sums of the last integral converge in $\mathcal{S}$, and $\omega$ is linear and continuous on $\mathcal{S}$.
By \cite[Lemma~8.15(i)]{QEDI}, in the form $\xi\left(X(g)^2\right)\le\|g\|_2^2$ for $g\in\mathcal{S}$, one has $\|X(f_{\mathrm{sl}})\|_{L^2(\xi)}\le\|f_{\mathrm{sl}}\|_2$.
Let $\mathrm{ov}$ be the function of Step 9; the assertions of (9a) involve no law.
Then
\[
\|f_{\mathrm{sl}}\|_2^2=\mathrm{sl}^{-2}\cdot\int_0^{\mathrm{sl}}\mathrm{d}s\int_0^{\mathrm{sl}}\mathrm{d}s'\,\mathrm{ov}\left((s'-s)e_j\right)\to0\qquad(\mathrm{sl}\to\infty),
\]
because $|\mathrm{ov}|\le\|f\|_2^2$ and $\mathrm{ov}(y)\to0$ as $|y|_2\to\infty$, by (9a).
\end{proof}

\subsection*{2. Proposition 45.5. Tensor tests}

\subsubsection*{2.1. Components, norms, auxiliary facts}

\textit{Components.}
Let $\mathrm{TI}:=\{(1,2),(1,3),(2,3)\}$.
For $\mathrm{ti}=(i,\ell)\in\mathrm{TI}$ put $\mathrm{d}y_{\mathrm{ti}}:=\mathrm{d}y_i\wedge\mathrm{d}y_\ell$; these three forms constitute the orthonormal basis of $\Lambda^2$ of \cite[Section~8.1.1]{QEDI}.
For $\mathrm{tj}=(\mathrm{ti}_1,\dots,\mathrm{ti}_n)\in\mathrm{TI}^n$ put $\mathrm{d}y_{\mathrm{tj}}:=\mathrm{d}y_{\mathrm{ti}_1}\otimes\dots\otimes\mathrm{d}y_{\mathrm{ti}_n}$; these $3^n$ tensors form a basis of $(\Lambda^2)^{\otimes n}$.
The space of the complex Schwartz functions on $\mathbb{R}^{3n}=(\mathbb{R}^3)^n$ is denoted $\mathcal{S}_{\mathbb{C}}(\mathbb{R}^{3n})$.
A point of $\mathbb{R}^{3n}$ is written $y=(y_1,\dots,y_n)$, with the \emph{arguments} $y_a\in\mathbb{R}^3$, $a=1,\dots,n$.
As in Proposition~45.2\,(5), the coordinates of the argument $y_a$ are written $(y_a)_1,(y_a)_2,(y_a)_3$, and $\partial_{a,\ell}$ denotes the derivative with respect to the coordinate $(y_a)_\ell$ ($a=1,\dots,n$; $\ell=1,2,3$).
The same conventions apply to the points $\mathrm{zl}=(\mathrm{zl}_1,\dots,\mathrm{zl}_n)$ and $\mathrm{kf}=(\mathrm{kf}_1,\dots,\mathrm{kf}_n)$ of $\mathbb{Z}^{3n}=(\mathbb{Z}^3)^n$.
Where $y$ stands for a single point of $\mathbb{R}^3$, namely in the sets $\{y_j>0\}$, $\{y_1>\mathrm{po}\}$, $\{\hat{n}\cdot y>\mathrm{po}\}$ and in the definition of $\mathrm{pun3}$, $\mathrm{cof3}$ in 2.2, an index at $y$ is the number of a coordinate.
On $\mathcal{S}_{\mathbb{C}}(\mathbb{R}^{3n})$ the norm $q_m$ is given by the formula of (n5),
\[
q_m(\Psi)=\max_{|\alpha|\le m}\sup_y\left(1+|y|_2^2\right)^{\frac{m}{2}}\cdot|\partial^\alpha\Psi(y)|.
\]
Every $\Psi\in\mathrm{Sw}_n$ can be written $\Psi=\sum_{\mathrm{tj}}\Psi_{\mathrm{tj}}\,\mathrm{d}y_{\mathrm{tj}}$ with components $\Psi_{\mathrm{tj}}\in\mathcal{S}_{\mathbb{C}}(\mathbb{R}^{3n})$, and
\[
q_m(\Psi)=\max_{\mathrm{tj}}q_m(\Psi_{\mathrm{tj}}),\qquad\overline{\Psi}:=\sum_{\mathrm{tj}}\overline{\Psi_{\mathrm{tj}}}\,\mathrm{d}y_{\mathrm{tj}}.
\]
For $\mathrm{sf}_1,\dots,\mathrm{sf}_n\in\mathcal{S}_{\mathbb{C}}(\mathbb{R}^3)$ let $(\mathrm{sf}_1\otimes\dots\otimes\mathrm{sf}_n)(y):=\prod_a\mathrm{sf}_a(y_a)$.
For $f_a=\sum_{\mathrm{ti}}f_{a,\mathrm{ti}}\,\mathrm{d}y_{\mathrm{ti}}\in\mathcal{S}$, $a=1,\dots,n$, one has
\begin{equation}\tag{2.1}
f_1\otimes\dots\otimes f_n=\sum_{\mathrm{tj}\in\mathrm{TI}^n}\left(f_{1,\mathrm{ti}_1}\otimes\dots\otimes f_{n,\mathrm{ti}_n}\right)\mathrm{d}y_{\mathrm{tj}}.
\end{equation}
An \emph{elementary tensor} is an element $(\mathrm{sf}_1\otimes\dots\otimes\mathrm{sf}_n)\,\mathrm{d}y_{\mathrm{tj}}$ of $\mathrm{Sw}_n$ with $\mathrm{sf}_a\in\mathcal{S}_{\mathbb{C}}(\mathbb{R}^3)$ and $\mathrm{tj}\in\mathrm{TI}^n$.
For $\Phi\in\mathrm{Sw}_m$ and $\Psi\in\mathrm{Sw}_n$ the tensor $\Phi\otimes\Psi\in\mathrm{Sw}_{m+n}$ has the components $(\Phi\otimes\Psi)_{(\mathrm{tj}',\mathrm{tj})}(y',y)=\Phi_{\mathrm{tj}'}(y')\cdot\Psi_{\mathrm{tj}}(y)$.
For $\Psi_0\in\mathrm{Sw}_0=\mathbb{C}$ one has $\Psi_0\otimes\Psi=\Psi_0\Psi$ by (n8), and one puts $\Psi\otimes\Psi_0:=\Psi_0\Psi$.

\textit{Topology.}
One has $q_0\le q_1\le q_2\le\cdots$.
The topology of $\mathrm{Sw}_n$ is that of the Schwartz space; it is defined by the seminorms $\max_{\mathrm{tj}}\sup_y|y^\beta\partial^\alpha\Psi_{\mathrm{tj}}(y)|$, where $\alpha,\beta$ are multi-indices.
The norms $q_m$ define the same topology.
On the one hand, $|y^\beta|\le\left(1+|y|_2^2\right)^{\frac{|\beta|}{2}}$, so that $\sup_y|y^\beta\partial^\alpha\Psi_{\mathrm{tj}}(y)|\le q_m(\Psi)$ whenever $m\ge\max\{|\alpha|,|\beta|\}$.
On the other hand,
\[
\left(1+|y|_2^2\right)^{\frac{m}{2}}\le\left(1+|y|_2^2\right)^m=\left(1+\sum_{a=1}^n\sum_{\ell=1}^3(y_a)_\ell^2\right)^m,
\]
and this power, once multiplied out, is a sum of $(1+3n)^m$ monomials $y^\beta$ with $|\beta|\le2m$; hence
\[
q_m(\Psi)\le(1+3n)^m\cdot\max_{\mathrm{tj}}\max_{|\alpha|\le m,\,|\beta|\le2m}\sup_y|y^\beta\partial^\alpha\Psi_{\mathrm{tj}}(y)|.
\]
Consequently, a linear map $\mathrm{Lm}$ of $\mathrm{Sw}_n$ into a normed space is continuous as soon as $\|\mathrm{Lm}(\Psi)\|\le\mathrm{cf}\cdot q_m(\Psi)$ for some $m$, some $\mathrm{cf}$ and all $\Psi$; and a continuous $\mathrm{Lm}$ satisfies $\mathrm{Lm}(\Psi_k)\to\mathrm{Lm}(\Psi)$ whenever $q_m(\Psi_k-\Psi)\to0$ for every $m$.

\textit{Peetre's inequality.}
For $y,y'\in\mathbb{R}^{3n}$,
\begin{equation}\tag{2.2}
1+|y|_2^2\le2\left(1+|y-y'|_2^2\right)\left(1+|y'|_2^2\right),
\end{equation}
since $|y|_2^2\le2|y-y'|_2^2+2|y'|_2^2$.

\begin{named}{(in1)}
$q_m(\Phi\otimes\Psi)\le q_m(\Phi)\cdot q_m(\Psi)$ for $\Phi\in\mathrm{Sw}_{n'}$ and $\Psi\in\mathrm{Sw}_n$.
\end{named}

\begin{proof}
A derivative of order at most $m$ of $\Phi_{\mathrm{tj}'}(y')\cdot\Psi_{\mathrm{tj}}(y)$ is of the form $(\partial^{\alpha'}\Phi_{\mathrm{tj}'})(y')\cdot(\partial^\alpha\Psi_{\mathrm{tj}})(y)$ with $|\alpha'|\le m$ and $|\alpha|\le m$, and
\[
\left(1+|y'|_2^2+|y|_2^2\right)^{\frac{m}{2}}\le\left(1+|y'|_2^2\right)^{\frac{m}{2}}\cdot\left(1+|y|_2^2\right)^{\frac{m}{2}}.
\]
\end{proof}

\begin{named}{(in2)}
Let $\mathrm{Gi}\colon y\mapsto\mathrm{Rl}\,y+b$ be a Euclidean map ($\mathrm{Rl}\in O(3)$, $b\in\mathbb{R}^3$), and let $(\mathrm{Gi}\cdot)^{\otimes n}$ be the map of (n8),
\[
\left((\mathrm{Gi}\cdot)^{\otimes n}\Psi\right)(y_1,\dots,y_n)=(\wedge^2\mathrm{Rl})^{\otimes n}\,\Psi(\mathrm{Gi}^{-1}y_1,\dots,\mathrm{Gi}^{-1}y_n).
\]
It is a linear map of $\mathrm{Sw}_n$ into itself, and
\[
(\mathrm{Gi}\cdot)^{\otimes n}(f_1\otimes\dots\otimes f_n)=(\mathrm{Gi}\cdot f_1)\otimes\dots\otimes(\mathrm{Gi}\cdot f_n),\qquad(\mathrm{Gi}\cdot)^{\otimes n}\,\overline{\Psi}=\overline{(\mathrm{Gi}\cdot)^{\otimes n}\Psi},
\]
\[
q_m\left((\mathrm{Gi}\cdot)^{\otimes n}\Psi\right)\le3^{n+m}\cdot\left(2\left(1+n|b|_2^2\right)\right)^{\frac{m}{2}}\cdot q_m(\Psi).
\]
For $\mathrm{Rl}=1$, that is $\mathrm{Gi}=\tau_b$, it is the map $\tau_b^{\otimes n}$ of (n8): $(\tau_b^{\otimes n}\Psi)(y_1,\dots,y_n)=\Psi(y_1-b,\dots,y_n-b)$.
\end{named}

\begin{proof}
The second assertion is (n2), applied factor by factor.
In the basis $(\mathrm{d}y_{\mathrm{ti}})$ the matrix of $\wedge^2\mathrm{Rl}$ has as its entries the $2\times2$ minors of $\mathrm{Rl}$; they are real and of modulus at most $1$.
Indeed, for real numbers $t_1,t_2,t_1',t_2'$ the inequality
\[
|t_1t_2'-t_2t_1'|\le\left(t_1^2+t_2^2\right)^{\frac{1}{2}}\left((t_1')^2+(t_2')^2\right)^{\frac{1}{2}}
\]
is the Cauchy--Schwarz inequality for the vectors $(t_1,t_2)$ and $(t_2',-t_1')$, and the rows of $\mathrm{Rl}$ have norm $1$.
Hence the matrix of $(\wedge^2\mathrm{Rl})^{\otimes n}$ in the basis $(\mathrm{d}y_{\mathrm{tj}})$ is real, which gives the third assertion, and each of its rows consists of $3^n$ entries of modulus at most $1$.
Let $y':=(\mathrm{Gi}^{-1}y_1,\dots,\mathrm{Gi}^{-1}y_n)$, whose arguments are $y'_a=\mathrm{Gi}^{-1}y_a=\mathrm{Rl}^{-1}(y_a-b)$.
For a smooth function $\mathrm{sf}$ on $\mathbb{R}^{3n}$ the chain rule shows that the derivative of $y\mapsto\mathrm{sf}(y')$ with respect to the coordinate $(y_a)_i$ equals
\[
\sum_{\ell=1}^3(\mathrm{Rl}^{-1})_{\ell i}\cdot(\partial_{a,\ell}\mathrm{sf})(y'),
\]
a sum of three terms whose coefficients have modulus at most $1$.
By induction on $|\alpha|$, every derivative of order $|\alpha|\le m$ of $y\mapsto\mathrm{sf}(y')$ is bounded by $3^{|\alpha|}\max_{|\beta|=|\alpha|}|(\partial^\beta\mathrm{sf})(y')|$.
Finally $y_a=\mathrm{Rl}\,y'_a+b$, so that $|y|_2^2=\sum_{a=1}^n|\mathrm{Rl}\,y'_a+b|_2^2\le2|y'|_2^2+2n|b|_2^2$ and $1+|y|_2^2\le2\left(1+n|b|_2^2\right)\left(1+|y'|_2^2\right)$.
Together these facts give, for $|\alpha|\le m$ and every $\mathrm{tj}$,
\begin{align*}
\left(1+|y|_2^2\right)^{\frac{m}{2}}\left|\partial^\alpha\left((\mathrm{Gi}\cdot)^{\otimes n}\Psi\right)_{\mathrm{tj}}(y)\right|&\le3^n\cdot3^m\cdot\left(2\left(1+n|b|_2^2\right)\right)^{\frac{m}{2}}\cdot\max_{\mathrm{tj}'}\max_{|\beta|\le m}\left(1+|y'|_2^2\right)^{\frac{m}{2}}\left|(\partial^\beta\Psi_{\mathrm{tj}'})(y')\right|\\
&\le3^{n+m}\cdot\left(2\left(1+n|b|_2^2\right)\right)^{\frac{m}{2}}\cdot q_m(\Psi).
\end{align*}
As this holds for every $m$, the function $(\mathrm{Gi}\cdot)^{\otimes n}\Psi$ belongs to $\mathrm{Sw}_n$.
\end{proof}

\begin{named}{(in3)}
For $\Psi\in\mathrm{Sw}_n$ and $b\in\mathbb{R}^3$ with $|b|_2\le1$,
\[
q_m\left(\tau_b^{\otimes n}\Psi-\Psi\right)\le\sqrt{3}\cdot n\cdot|b|_2\cdot\left(2(1+n)\right)^{\frac{m}{2}}\cdot q_{m+1}(\Psi).
\]
\end{named}

\begin{proof}
Let $b^n:=(b,\dots,b)\in\mathbb{R}^{3n}$; then $|b^n|_1=n|b|_1\le\sqrt{3}\cdot n\cdot|b|_2$ and $|b^n|_2^2=n|b|_2^2\le n$.
Let $|\alpha|\le m$, and let $\Psi_{\mathrm{tj}}$ be a component of $\Psi$; the corresponding component of $\tau_b^{\otimes n}\Psi$ is $y\mapsto\Psi_{\mathrm{tj}}(y-b^n)$, by (in2).
The fundamental theorem of calculus, applied to $t\mapsto\partial^\alpha\Psi_{\mathrm{tj}}(y-tb^n)$, gives
\[
\partial^\alpha\Psi_{\mathrm{tj}}(y-b^n)-\partial^\alpha\Psi_{\mathrm{tj}}(y)=-\int_0^1\mathrm{d}t\,b^n\cdot\left(\mathrm{grad}\,\partial^\alpha\Psi_{\mathrm{tj}}\right)(y-tb^n),
\]
whose modulus is at most $|b^n|_1\cdot\sup_{0\le t\le1}\max_{a,\ell}\left|\partial_{a,\ell}\partial^\alpha\Psi_{\mathrm{tj}}(y-tb^n)\right|$.
By (2.2), $\left(1+|y|_2^2\right)^{\frac{m}{2}}\le\left(2(1+n)\right)^{\frac{m}{2}}\left(1+|y-tb^n|_2^2\right)^{\frac{m}{2}}$ for $0\le t\le1$, and $\left(1+|y'|_2^2\right)^{\frac{m}{2}}\left|\partial_{a,\ell}\partial^\alpha\Psi_{\mathrm{tj}}(y')\right|\le q_{m+1}(\Psi)$ for all $y'$.
\end{proof}

\begin{named}{(in4)}
$\|\mathrm{sf}\|_{L^2(\mathbb{R}^3)}\le\pi\cdot q_2(\mathrm{sf})$ for $\mathrm{sf}\in\mathcal{S}_{\mathbb{C}}(\mathbb{R}^3)$.
\end{named}

\begin{proof}
One has $|\mathrm{sf}(y)|^2\le q_2(\mathrm{sf})^2\cdot\left(1+|y|_2^2\right)^{-2}$, and in polar coordinates
\[
\int_{\mathbb{R}^3}\mathrm{d}y\,\left(1+|y|_2^2\right)^{-2}=4\pi\int_0^\infty\mathrm{d}t\,t^2(1+t^2)^{-2}=4\pi\cdot\frac{\pi}{4}=\pi^2;
\]
indeed, the function $\frac{1}{2}\left(\arctan t-\frac{t}{1+t^2}\right)$ has the derivative $t^2(1+t^2)^{-2}$, vanishes at $t=0$ and tends to $\pi/4$ as $t\to\infty$.
\end{proof}

\begin{named}{(FS)}
Let $\mathrm{sl}>0$, and let $\mathrm{sf}$ be a continuous function on $\mathbb{R}^{3n}$, periodic with respect to $\mathrm{sl}\,\mathbb{Z}^{3n}$, whose Fourier coefficients
\[
\mathrm{fd}_{\mathrm{kf}}:=\mathrm{sl}^{-3n}\cdot\int_{(-\mathrm{sl}/2,\mathrm{sl}/2)^{3n}}\mathrm{d}y\,\mathrm{sf}(y)\cdot\exp\left(-\frac{2\pi i\,\mathrm{kf}\cdot y}{\mathrm{sl}}\right),\qquad\mathrm{kf}\in\mathbb{Z}^{3n},
\]
satisfy $\sum_{\mathrm{kf}}|\mathrm{fd}_{\mathrm{kf}}|<\infty$.
Then
\[
\mathrm{sf}(y)=\sum_{\mathrm{kf}}\mathrm{fd}_{\mathrm{kf}}\cdot\exp\left(\frac{2\pi i\,\mathrm{kf}\cdot y}{\mathrm{sl}}\right)\qquad\text{for every }y.
\]
\end{named}

\begin{proof}
The series converges uniformly.
Its sum $\mathrm{sf}'$ is continuous and periodic, and by the orthogonality relations and the uniform convergence its Fourier coefficients are the numbers $\mathrm{fd}_{\mathrm{kf}}$.
Therefore all Fourier coefficients of the continuous periodic function $\mathrm{sf}'':=\mathrm{sf}-\mathrm{sf}'$ vanish.
By the Stone--Weierstrass theorem there are trigonometric polynomials $\mathrm{sf}_k$, that is, finite linear combinations of the functions $\exp\left(\frac{2\pi i\,\mathrm{kf}\cdot y}{\mathrm{sl}}\right)$, such that $\mathrm{sf}_k\to\mathrm{sf}''$ uniformly.
Then $\int\mathrm{d}y\,\overline{\mathrm{sf}_k}\cdot\mathrm{sf}''=0$ over the cube of periods for every $k$, and the limit $k\to\infty$ gives $\int\mathrm{d}y\,|\mathrm{sf}''|^2=0$; hence $\mathrm{sf}''=0$.
\end{proof}

\begin{named}{(abs)}
Let $(\Psi_p)$ be a countable family in $\mathcal{S}_{\mathbb{C}}(\mathbb{R}^{3n})$ with $\sum_pq_m(\Psi_p)<\infty$ for every $m$, and let $\Psi$ be a function on $\mathbb{R}^{3n}$ with $\Psi(y)=\sum_p\Psi_p(y)$ for every $y$.
Then $\Psi\in\mathcal{S}_{\mathbb{C}}(\mathbb{R}^{3n})$, and for every finite set $\mathrm{Fs}$ of indices and every $m$
\[
q_m\left(\Psi-\sum_{p\in\mathrm{Fs}}\Psi_p\right)\le\sum_{p\notin\mathrm{Fs}}q_m(\Psi_p),
\]
which tends to $0$ as $\mathrm{Fs}$ exhausts the set of the indices.
\end{named}

\begin{proof}
For every multi-index $\alpha$ one has $\sup|\partial^\alpha\Psi_p|\le q_{|\alpha|}(\Psi_p)$, so that the series $\sum_p\partial^\alpha\Psi_p$ converges uniformly.
Hence $\Psi$ is smooth, and $\partial^\alpha\Psi=\sum_p\partial^\alpha\Psi_p$ by differentiation term by term.
For $|\alpha|\le m$,
\[
\left(1+|y|_2^2\right)^{\frac{m}{2}}\left|\partial^\alpha\left(\Psi-\sum_{p\in\mathrm{Fs}}\Psi_p\right)(y)\right|\le\sum_{p\notin\mathrm{Fs}}\left(1+|y|_2^2\right)^{\frac{m}{2}}|\partial^\alpha\Psi_p(y)|\le\sum_{p\notin\mathrm{Fs}}q_m(\Psi_p).
\]
For $\mathrm{Fs}=\emptyset$ this shows that $q_m(\Psi)$ is finite for every $m$.
\end{proof}

\subsubsection*{2.2. Two lemmas on test functions}

\textit{Fixed functions.}
Let
\[
\mathrm{wz}(t):=\exp\left(-\left(\frac{9}{16}-t^2\right)^{-1}\right)\quad\text{for }|t|<\frac{3}{4},\qquad\mathrm{wz}(t):=0\quad\text{for }|t|\ge\frac{3}{4}.
\]
It is the composition of the polynomial $t\mapsto\frac{9}{16}-t^2$ with the function that equals $\exp\left(-\frac{1}{t'}\right)$ for $t'>0$ and $0$ for $t'\le0$, which belongs to $C^\infty(\mathbb{R})$.
So $\mathrm{wz}$ is a function in $C^\infty(\mathbb{R})$, $\mathrm{wz}\ge0$, $\operatorname{supp}\mathrm{wz}=[-\frac{3}{4},\frac{3}{4}]$, and $\mathrm{wz}>0$ on $[-\frac{1}{2},\frac{1}{2}]$.
Let
\[
\mathrm{Wp}(t):=\sum_{\mathrm{zl}\in\mathbb{Z}}\mathrm{wz}(t-\mathrm{zl}).
\]
The sum is locally finite; $\mathrm{Wp}$ is smooth, periodic with period $1$, and positive, because every real $t$ admits an integer $\mathrm{zl}$ with $|t-\mathrm{zl}|\le\frac{1}{2}$.
Put
\[
\mathrm{pun}:=\mathrm{wz}/\mathrm{Wp},\qquad\mathrm{cof}(t):=\mathrm{pun}(t+1)+\mathrm{pun}(t)+\mathrm{pun}(t-1).
\]
Then:
\begin{itemize}
\item[(fn1)] $\mathrm{pun}\in C^\infty(\mathbb{R})$, $\mathrm{pun}\ge0$, $\operatorname{supp}\mathrm{pun}=[-\frac{3}{4},\frac{3}{4}]$, and, since $\mathrm{Wp}(t-\mathrm{zl})=\mathrm{Wp}(t)$,
\[
\sum_{\mathrm{zl}\in\mathbb{Z}}\mathrm{pun}(t-\mathrm{zl})=\mathrm{Wp}(t)^{-1}\sum_{\mathrm{zl}\in\mathbb{Z}}\mathrm{wz}(t-\mathrm{zl})=1\qquad\text{for all }t;
\]
\item[(fn2)] $\mathrm{cof}\in C^\infty(\mathbb{R})$, $0\le\mathrm{cof}\le1$, $\operatorname{supp}\mathrm{cof}\subset[-\frac{7}{4},\frac{7}{4}]$, and $\mathrm{cof}=1$ on $[-1,1]$: for $|t|\le1$, a term $\mathrm{pun}(t-\mathrm{zl})\ne0$ has $|t-\mathrm{zl}|<\frac{3}{4}$, hence $|\mathrm{zl}|<\frac{7}{4}$ and $\mathrm{zl}\in\{-1,0,1\}$, so that $\mathrm{cof}(t)=\sum_{\mathrm{zl}\in\mathbb{Z}}\mathrm{pun}(t-\mathrm{zl})=1$.
\end{itemize}
For a single point $y$ of $\mathbb{R}^3$ with the coordinates $y_1,y_2,y_3$ let
\[
\mathrm{pun3}(y):=\mathrm{pun}(y_1)\cdot\mathrm{pun}(y_2)\cdot\mathrm{pun}(y_3),\qquad\mathrm{cof3}(y):=\mathrm{cof}(y_1)\cdot\mathrm{cof}(y_2)\cdot\mathrm{cof}(y_3).
\]
For a point $y=(y_1,\dots,y_n)$ of $\mathbb{R}^{3n}$ with the arguments $y_a\in\mathbb{R}^3$ let
\[
\mathrm{PU}_n(y):=\prod_a\mathrm{pun3}(y_a),\qquad\mathrm{CO}_n(y):=\prod_a\mathrm{cof3}(y_a).
\]
Then
\[
\sum_{\mathrm{zl}\in\mathbb{Z}^{3n}}\mathrm{PU}_n(y-\mathrm{zl})=1,\qquad\operatorname{supp}\mathrm{PU}_n\subset\left[-\tfrac{3}{4},\tfrac{3}{4}\right]^{3n},
\]
\[
0\le\mathrm{CO}_n\le1,\qquad\mathrm{CO}_n=1\ \text{on}\ [-1,1]^{3n},\qquad\operatorname{supp}\mathrm{CO}_n\subset\left[-\tfrac{7}{4},\tfrac{7}{4}\right]^{3n}.
\]
Further, let
\[
\mathrm{pu}_m:=\max_{|\beta|\le m}\sup|\partial^\beta\mathrm{PU}_n|,\qquad\mathrm{cm}_m:=\max_{|\beta|\le m}\sup|\partial^\beta\mathrm{CO}_n|\ge1;
\]
these numbers depend on $n$ and $m$.

\textit{Building functions and coefficients.}
For $\mathrm{zl},\mathrm{kf}\in\mathbb{Z}^3$ let
\[
\mathrm{bfn}_{\mathrm{zl},\mathrm{kf}}(y):=\mathrm{cof3}(y-\mathrm{zl})\cdot\exp\left(\frac{i\pi}{2}\,\mathrm{kf}\cdot(y-\mathrm{zl})\right),
\]
an element of $C_c^\infty(\mathbb{R}^3;\mathbb{C})$.
For $\mathrm{zl}=(\mathrm{zl}_1,\dots,\mathrm{zl}_n)$ and $\mathrm{kf}=(\mathrm{kf}_1,\dots,\mathrm{kf}_n)$ in $(\mathbb{Z}^3)^n=\mathbb{Z}^{3n}$ let
\[
\mathrm{bfn}_{\mathrm{zl},\mathrm{kf}}:=\mathrm{bfn}_{\mathrm{zl}_1,\mathrm{kf}_1}\otimes\dots\otimes\mathrm{bfn}_{\mathrm{zl}_n,\mathrm{kf}_n},\quad\text{that is,}\quad\mathrm{bfn}_{\mathrm{zl},\mathrm{kf}}(y)=\mathrm{CO}_n(y-\mathrm{zl})\cdot\exp\left(\frac{i\pi}{2}\,\mathrm{kf}\cdot(y-\mathrm{zl})\right),
\]
and for $\Psi\in\mathcal{S}_{\mathbb{C}}(\mathbb{R}^{3n})$ let
\[
\mathrm{fc}_{\mathrm{zl},\mathrm{kf}}(\Psi):=4^{-3n}\cdot\int_{\mathbb{R}^{3n}}\mathrm{d}y\,\mathrm{PU}_n(y)\cdot\Psi(y+\mathrm{zl})\cdot\exp\left(-\frac{i\pi}{2}\,\mathrm{kf}\cdot y\right).
\]

\begin{named}{Expansion Lemma}[expansion in products]
Let $n\ge1$ and $\Psi\in\mathcal{S}_{\mathbb{C}}(\mathbb{R}^{3n})$.
\begin{itemize}
\item[(xp1)] $\mathrm{fc}_{\mathrm{zl},\mathrm{kf}}$ is a linear functional on $\mathcal{S}_{\mathbb{C}}(\mathbb{R}^{3n})$, and
\[
\mathrm{fc}_{\mathrm{zl},\mathrm{kf}}(\mathrm{sf}_1\otimes\dots\otimes\mathrm{sf}_n)=\prod_{a=1}^n\mathrm{fc}_{\mathrm{zl}_a,\mathrm{kf}_a}(\mathrm{sf}_a)\qquad\text{for }\mathrm{sf}_a\in\mathcal{S}_{\mathbb{C}}(\mathbb{R}^3).
\]
\item[(xp2)] For every $m\ge0$,
\[
\sum_{\mathrm{zl},\mathrm{kf}\in\mathbb{Z}^{3n}}\left|\mathrm{fc}_{\mathrm{zl},\mathrm{kf}}(\Psi)\right|\cdot q_m(\mathrm{bfn}_{\mathrm{zl},\mathrm{kf}})\le\mathrm{cf}_{n,m}\cdot q_{2m+6n}(\Psi),
\]
with
\[
\mathrm{cf}_{n,m}:=2^{-3n}\cdot\left(8(1+3n)^2\right)^{m+3n}\cdot\mathrm{pu}_{2m+6n}\cdot\pi^m\cdot\mathrm{cm}_m\cdot\left(2(1+12n)\right)^{\frac{m}{2}}\cdot5^{6n}.
\]
\item[(xp3)] $\Psi=\sum_{\mathrm{zl},\mathrm{kf}}\mathrm{fc}_{\mathrm{zl},\mathrm{kf}}(\Psi)\cdot\mathrm{bfn}_{\mathrm{zl},\mathrm{kf}}$, the series converging absolutely in every norm $q_m$.
\item[(xp4)] $q_0(\mathrm{bfn}_{\mathrm{zl},\mathrm{kf}})=1$; for $n=1$, $\|\mathrm{bfn}_{\mathrm{zl},\mathrm{kf}}\|_{L^2(\mathbb{R}^3)}=\|\mathrm{cof}\|_{L^2(\mathbb{R})}^3\le8$.
\end{itemize}
\end{named}

\begin{proof}
\textit{(xp1).}
The linearity is evident from the formula.
For $\Psi=\mathrm{sf}_1\otimes\dots\otimes\mathrm{sf}_n$ the integrand is the product over $a$ of the functions $\mathrm{pun3}(y_a)\cdot\mathrm{sf}_a(y_a+\mathrm{zl}_a)\cdot\exp\left(-\frac{i\pi}{2}\,\mathrm{kf}_a\cdot y_a\right)$, and Fubini's theorem applies.

\textit{(xp4).}
One has $|\mathrm{bfn}_{\mathrm{zl},\mathrm{kf}}(y)|=\mathrm{CO}_n(y-\mathrm{zl})$, whose maximum is $1$.
For $n=1$, $\int\mathrm{d}y\,|\mathrm{bfn}_{\mathrm{zl},\mathrm{kf}}|^2=\left(\int\mathrm{d}t\,\mathrm{cof}(t)^2\right)^3$, and $\int\mathrm{d}t\,\mathrm{cof}(t)^2\le\frac{7}{2}\le4$ by (fn2).

\textit{(xp2): the coefficients.}
Put $\Psi_{\mathrm{zl}}(y):=\mathrm{PU}_n(y)\cdot\Psi(y+\mathrm{zl})$, an element of $C_c^\infty\left((-1,1)^{3n};\mathbb{C}\right)$.
Let $\mathrm{Lap}$ be the Laplace operator of $\mathbb{R}^{3n}$, and let $\mathrm{pw}\ge0$ be an integer.
Since
\[
(1-\mathrm{Lap})\exp\left(-\frac{i\pi}{2}\,\mathrm{kf}\cdot y\right)=\left(1+\left(\frac{\pi}{2}\right)^2|\mathrm{kf}|_2^2\right)\exp\left(-\frac{i\pi}{2}\,\mathrm{kf}\cdot y\right),
\]
integration by parts, which produces no boundary terms because $\Psi_{\mathrm{zl}}$ has compact support, gives
\[
\left(1+\left(\frac{\pi}{2}\right)^2|\mathrm{kf}|_2^2\right)^{\mathrm{pw}}\mathrm{fc}_{\mathrm{zl},\mathrm{kf}}(\Psi)=4^{-3n}\cdot\int\mathrm{d}y\,\left((1-\mathrm{Lap})^{\mathrm{pw}}\Psi_{\mathrm{zl}}\right)(y)\cdot\exp\left(-\frac{i\pi}{2}\,\mathrm{kf}\cdot y\right).
\]
The support of $\Psi_{\mathrm{zl}}$ has volume at most $2^{3n}$.
The operator $(1-\mathrm{Lap})^{\mathrm{pw}}=\left(1-\sum_{a=1}^n\sum_{\ell=1}^3\partial_{a,\ell}^2\right)^{\mathrm{pw}}$ is a sum of $(1+3n)^{\mathrm{pw}}$ terms $\pm\partial^\alpha$ with $|\alpha|\le2\mathrm{pw}$.
By Leibniz' rule,
\[
|\partial^\alpha\Psi_{\mathrm{zl}}(y)|\le2^{|\alpha|}\cdot\mathrm{pu}_{2\mathrm{pw}}\cdot\max_{|\gamma|\le2\mathrm{pw}}|\partial^\gamma\Psi(y+\mathrm{zl})|.
\]
For $y\in[-1,1]^{3n}$ one has $|y|_2^2\le3n$, and (2.2) gives $1+|\mathrm{zl}|_2^2\le2(1+3n)\left(1+|y+\mathrm{zl}|_2^2\right)$; hence, for $|\gamma|\le2\mathrm{pw}$,
\[
|\partial^\gamma\Psi(y+\mathrm{zl})|\le q_{2\mathrm{pw}}(\Psi)\cdot\left(1+|y+\mathrm{zl}|_2^2\right)^{-\mathrm{pw}}\le q_{2\mathrm{pw}}(\Psi)\cdot\left(2(1+3n)\right)^{\mathrm{pw}}\cdot\left(1+|\mathrm{zl}|_2^2\right)^{-\mathrm{pw}}.
\]
Combining these bounds and using $1+\left(\frac{\pi}{2}\right)^2|\mathrm{kf}|_2^2\ge1+|\mathrm{kf}|_2^2$, one obtains
\begin{equation}\tag{2.3}
\left|\mathrm{fc}_{\mathrm{zl},\mathrm{kf}}(\Psi)\right|\le2^{-3n}\cdot\left(8(1+3n)^2\right)^{\mathrm{pw}}\cdot\mathrm{pu}_{2\mathrm{pw}}\cdot\left(1+|\mathrm{kf}|_2^2\right)^{-\mathrm{pw}}\cdot\left(1+|\mathrm{zl}|_2^2\right)^{-\mathrm{pw}}\cdot q_{2\mathrm{pw}}(\Psi).
\end{equation}

\textit{(xp2): the building functions.}
For $|\alpha|\le m$, Leibniz' rule gives
\[
\partial^\alpha\mathrm{bfn}_{\mathrm{zl},\mathrm{kf}}(y)=\sum_{\beta\le\alpha}\binom{\alpha}{\beta}\left(\partial^{\alpha-\beta}\mathrm{CO}_n\right)(y-\mathrm{zl})\cdot\left(\frac{i\pi}{2}\,\mathrm{kf}\right)^\beta\cdot\exp\left(\frac{i\pi}{2}\,\mathrm{kf}\cdot(y-\mathrm{zl})\right),
\]
where $\left|\left(\frac{i\pi}{2}\,\mathrm{kf}\right)^\beta\right|\le\left(\frac{\pi}{2}\right)^m\left(1+|\mathrm{kf}|_2^2\right)^{\frac{m}{2}}$ and $\sum_{\beta\le\alpha}\binom{\alpha}{\beta}=2^{|\alpha|}\le2^m$.
On the support one has $|y-\mathrm{zl}|_2^2\le3n\cdot4$, and (2.2) gives $1+|y|_2^2\le2(1+12n)\left(1+|\mathrm{zl}|_2^2\right)$.
Hence
\begin{equation}\tag{2.4}
q_m(\mathrm{bfn}_{\mathrm{zl},\mathrm{kf}})\le\pi^m\cdot\mathrm{cm}_m\cdot\left(2(1+12n)\right)^{\frac{m}{2}}\cdot\left(1+|\mathrm{kf}|_2^2\right)^{\frac{m}{2}}\cdot\left(1+|\mathrm{zl}|_2^2\right)^{\frac{m}{2}}.
\end{equation}

\textit{(xp2): the sum.}
Take $\mathrm{pw}:=m+3n$ in (2.3).
The product of (2.3) and (2.4) contains the factor $\left(1+|\mathrm{kf}|_2^2\right)^{-\mathrm{pw}+\frac{m}{2}}\left(1+|\mathrm{zl}|_2^2\right)^{-\mathrm{pw}+\frac{m}{2}}$, and $\mathrm{pw}-\frac{m}{2}\ge3n$.
Since $\left(1+|\mathrm{kf}|_2^2\right)^{3n}\ge\prod_{a=1}^n\prod_{\ell=1}^3\left(1+(\mathrm{kf}_a)_\ell^2\right)$,
\[
\sum_{\mathrm{kf}\in\mathbb{Z}^{3n}}\left(1+|\mathrm{kf}|_2^2\right)^{-3n}\le\left(\sum_{\mathrm{ki}\in\mathbb{Z}}\left(1+\mathrm{ki}^2\right)^{-1}\right)^{3n}\le5^{3n},
\]
because
\[
\sum_{\mathrm{ki}\in\mathbb{Z}}\left(1+\mathrm{ki}^2\right)^{-1}\le1+2\sum_{\mathrm{ki}\ge1}\mathrm{ki}^{-2}\le1+2\left(1+\sum_{\mathrm{ki}\ge2}\left(\frac{1}{\mathrm{ki}-1}-\frac{1}{\mathrm{ki}}\right)\right)=5.
\]
The sum over $\mathrm{zl}\in\mathbb{Z}^{3n}$ is bounded by the same number.
This proves (xp2).

\textit{(xp3).}
Fix $\mathrm{zl}$.
The periodic extension of $\Psi_{\mathrm{zl}}$ with respect to $4\,\mathbb{Z}^{3n}$ is smooth, because $\operatorname{supp}\Psi_{\mathrm{zl}}\subset[-\frac{3}{4},\frac{3}{4}]^{3n}$ lies in the interior of the cube $(-2,2)^{3n}$.
Its Fourier coefficients in the sense of (FS), with $\mathrm{sl}=4$, are the numbers $\mathrm{fc}_{\mathrm{zl},\mathrm{kf}}(\Psi)$, $\mathrm{kf}\in\mathbb{Z}^{3n}$, and they are summable by (2.3) with $\mathrm{pw}=3n$.
By (FS),
\[
\Psi_{\mathrm{zl}}(y)=\sum_{\mathrm{kf}}\mathrm{fc}_{\mathrm{zl},\mathrm{kf}}(\Psi)\cdot\exp\left(\frac{i\pi}{2}\,\mathrm{kf}\cdot y\right)\qquad\text{for }y\in(-2,2)^{3n}.
\]
Multiply by $\mathrm{CO}_n(y)$.
As $\mathrm{CO}_n=1$ on $\operatorname{supp}\Psi_{\mathrm{zl}}$ and $\mathrm{CO}_n=0$ outside $(-2,2)^{3n}$,
\[
\Psi_{\mathrm{zl}}(y)=\sum_{\mathrm{kf}}\mathrm{fc}_{\mathrm{zl},\mathrm{kf}}(\Psi)\cdot\mathrm{CO}_n(y)\cdot\exp\left(\frac{i\pi}{2}\,\mathrm{kf}\cdot y\right)\qquad\text{for \emph{all} }y\in\mathbb{R}^{3n}.
\]
Now replace $y$ by $y-\mathrm{zl}$ and sum over $\mathrm{zl}$.
Since $\sum_{\mathrm{zl}}\Psi_{\mathrm{zl}}(y-\mathrm{zl})=\sum_{\mathrm{zl}}\mathrm{PU}_n(y-\mathrm{zl})\cdot\Psi(y)=\Psi(y)$, it follows that
\[
\Psi(y)=\sum_{\mathrm{zl}}\sum_{\mathrm{kf}}\mathrm{fc}_{\mathrm{zl},\mathrm{kf}}(\Psi)\cdot\mathrm{bfn}_{\mathrm{zl},\mathrm{kf}}(y)\qquad\text{for every }y.
\]
By (xp2), the family $\left(\mathrm{fc}_{\mathrm{zl},\mathrm{kf}}(\Psi)\cdot\mathrm{bfn}_{\mathrm{zl},\mathrm{kf}}\right)$, indexed by $(\mathrm{zl},\mathrm{kf})\in\mathbb{Z}^{3n}\times\mathbb{Z}^{3n}$, satisfies the hypothesis of (abs); hence the order of summation is immaterial, and (abs) gives the convergence in every norm $q_m$.
\end{proof}

\textit{Consequence.}
Let $\Psi=\sum_{\mathrm{tj}}\Psi_{\mathrm{tj}}\,\mathrm{d}y_{\mathrm{tj}}\in\mathrm{Sw}_n$.
Applying (xp3) to each component, one obtains
\begin{equation}\tag{2.5}
\Psi=\sum_{\mathrm{tj}\in\mathrm{TI}^n}\sum_{\mathrm{zl},\mathrm{kf}}\mathrm{fc}_{\mathrm{zl},\mathrm{kf}}(\Psi_{\mathrm{tj}})\cdot\mathrm{bfn}_{\mathrm{zl},\mathrm{kf}}\,\mathrm{d}y_{\mathrm{tj}},
\end{equation}
the series converging absolutely in every norm $q_m$.
Each $\mathrm{bfn}_{\mathrm{zl},\mathrm{kf}}\,\mathrm{d}y_{\mathrm{tj}}$ is an elementary tensor.
Hence the linear span of the elementary tensors is dense in $\mathrm{Sw}_n$.

\begin{named}{(in5)}
For $\Phi\in\mathrm{Sw}_n$ and $k\ge1$,
\[
q_m\left(\mathrm{CO}_n(\cdot/k)\,\Phi\right)\le2^m\cdot\mathrm{cm}_m\cdot q_m(\Phi),\qquad q_m\left(\left(1-\mathrm{CO}_n(\cdot/k)\right)\Phi\right)\le2^m\cdot\mathrm{cm}_m\cdot\left(1+k^2\right)^{-\frac{1}{2}}\cdot q_{m+1}(\Phi).
\]
\end{named}

\begin{proof}
For $|\alpha|\le m$, Leibniz' rule gives
\[
\partial^\alpha\left(\mathrm{CO}_n(y/k)\cdot\Phi_{\mathrm{tj}}(y)\right)=\sum_{\beta\le\alpha}\binom{\alpha}{\beta}k^{-|\beta|}\cdot\left(\partial^\beta\mathrm{CO}_n\right)(y/k)\cdot\partial^{\alpha-\beta}\Phi_{\mathrm{tj}}(y),
\]
and the first bound follows.
The function $1-\mathrm{CO}_n(y/k)$ takes its values in $[0,1]$ and vanishes, together with all its derivatives, on $\{|y|_\infty\le k\}$.
So $\partial^\alpha\left(\left(1-\mathrm{CO}_n(y/k)\right)\cdot\Phi_{\mathrm{tj}}(y)\right)$ vanishes for $|y|_\infty\le k$, and for $|y|_\infty>k$ it is bounded by $2^m\cdot\mathrm{cm}_m\cdot\max_{|\gamma|\le m}|\partial^\gamma\Phi_{\mathrm{tj}}(y)|$.
For such $y$ one has $|y|_2>k$ and
\[
\left(1+|y|_2^2\right)^{\frac{m}{2}}|\partial^\gamma\Phi_{\mathrm{tj}}(y)|\le\left(1+|y|_2^2\right)^{-\frac{1}{2}}\cdot q_{m+1}(\Phi)\le\left(1+k^2\right)^{-\frac{1}{2}}\cdot q_{m+1}(\Phi).
\]
\end{proof}

\begin{named}{Support Lemma}[products with supports in an open set]
Let $O$ be an open subset of $\mathbb{R}^3$, let $n\ge1$, and let $\Psi\in C_c^\infty(O^n;\mathbb{C})$.
There are functions $\mathrm{sf}_{a,\mathrm{ki}}\in C_c^\infty(O;\mathbb{C})$ ($a=1,\dots,n$; $\mathrm{ki}\in\mathbb{Z}^3$) and complex numbers $\mathrm{fd}_{\mathrm{kf}}$ ($\mathrm{kf}=(\mathrm{kf}_1,\dots,\mathrm{kf}_n)\in(\mathbb{Z}^3)^n$) such that
\[
\sum_{\mathrm{kf}}|\mathrm{fd}_{\mathrm{kf}}|\cdot q_m\left(\mathrm{sf}_{1,\mathrm{kf}_1}\otimes\dots\otimes\mathrm{sf}_{n,\mathrm{kf}_n}\right)<\infty\qquad\text{for every }m
\]
and
\[
\Psi=\sum_{\mathrm{kf}}\mathrm{fd}_{\mathrm{kf}}\cdot\mathrm{sf}_{1,\mathrm{kf}_1}\otimes\dots\otimes\mathrm{sf}_{n,\mathrm{kf}_n},
\]
the series converging in every norm $q_m$.
\end{named}

\noindent\textit{Proof.}
If $\Psi=0$, take $\mathrm{fd}_{\mathrm{kf}}=0$ and $\mathrm{sf}_{a,\mathrm{ki}}=0$.
Let $\Psi\ne0$, and let $\mathrm{Kp}_a$ be the image of $\operatorname{supp}\Psi$, a compact subset of $O^n$, under the continuous map $(y_1,\dots,y_n)\mapsto y_a$.
It is a non-empty compact subset of $O$, and $\operatorname{supp}\Psi\subset\mathrm{Kp}_1\times\dots\times\mathrm{Kp}_n$.

\textit{Cut-off functions.}
If $O=\mathbb{R}^3$, let $\delta_a:=1$; otherwise let $\delta_a$ be the Euclidean distance between $\mathrm{Kp}_a$ and $\mathbb{R}^3\setminus O$, which is positive, the first set being compact, the second closed, and the two disjoint.
Choose $\mathrm{mk}>0$ with $2\sqrt{3}\cdot\mathrm{mk}<\delta_a$ for all $a$, and put
\[
\mathrm{cq}_a(y):=\sum\mathrm{pun3}\left(y/\mathrm{mk}-\mathrm{zl}\right),
\]
the sum extending over the finitely many $\mathrm{zl}\in\mathbb{Z}^3$ for which the closed cube $\mathrm{mk}\cdot\mathrm{zl}+[-\mathrm{mk},\mathrm{mk}]^3$ meets $\mathrm{Kp}_a$.
The function $\mathrm{cq}_a$ is smooth and takes its values in $[0,1]$, being part of the sum $\sum_{\mathrm{zl}\in\mathbb{Z}^3}\mathrm{pun3}(y/\mathrm{mk}-\mathrm{zl})=1$ of non-negative terms.
Moreover $\mathrm{cq}_a=1$ on $\mathrm{Kp}_a$: if $y\in\mathrm{Kp}_a$ and $\mathrm{pun3}(y/\mathrm{mk}-\mathrm{zl})\ne0$, then $|y/\mathrm{mk}-\mathrm{zl}|_\infty<1$, so that $y$ lies in $\mathrm{mk}\cdot\mathrm{zl}+[-\mathrm{mk},\mathrm{mk}]^3$; this cube meets $\mathrm{Kp}_a$, and the term belongs to the sum defining $\mathrm{cq}_a$.
The support of $\mathrm{cq}_a$ is contained in the union of the finitely many closed cubes just described.
Every point of such a cube has Euclidean distance at most $2\sqrt{3}\cdot\mathrm{mk}<\delta_a$ from a point of $\mathrm{Kp}_a$, and therefore lies in $O$.
So $\mathrm{cq}_a\in C_c^\infty(O)$.
Put $\mathrm{CQ}(y):=\prod_a\mathrm{cq}_a(y_a)$; then $\mathrm{CQ}=1$ on $\operatorname{supp}\Psi$.

\textit{Fourier series.}
Choose $\mathrm{sl}>0$ such that $\operatorname{supp}\mathrm{cq}_a\subset(-\mathrm{sl}/2,\mathrm{sl}/2)^3$ for all $a$.
Then $\operatorname{supp}\Psi$ is contained in the open cube $(-\mathrm{sl}/2,\mathrm{sl}/2)^{3n}$, and the periodic extension of $\Psi$ with respect to $\mathrm{sl}\,\mathbb{Z}^{3n}$ is smooth.
Its Fourier coefficients $\mathrm{fd}_{\mathrm{kf}}$ in the sense of (FS) satisfy, by integration by parts, for every integer $\mathrm{pw}\ge0$,
\[
\left(1+\left(\frac{2\pi}{\mathrm{sl}}\right)^2|\mathrm{kf}|_2^2\right)^{\mathrm{pw}}|\mathrm{fd}_{\mathrm{kf}}|\le\sup\left|(1-\mathrm{Lap})^{\mathrm{pw}}\Psi\right|.
\]
Since $1+\left(\frac{2\pi}{\mathrm{sl}}\right)^2|\mathrm{kf}|_2^2\ge\min\left\{1,\left(\frac{2\pi}{\mathrm{sl}}\right)^2\right\}\left(1+|\mathrm{kf}|_2^2\right)$, this gives
\begin{equation}\tag{2.6}
|\mathrm{fd}_{\mathrm{kf}}|\le\max\left\{1,\frac{\mathrm{sl}}{2\pi}\right\}^{2\mathrm{pw}}\cdot\sup\left|(1-\mathrm{Lap})^{\mathrm{pw}}\Psi\right|\cdot\left(1+|\mathrm{kf}|_2^2\right)^{-\mathrm{pw}}.
\end{equation}
For $\mathrm{pw}=3n$ the coefficients are summable, because $\sum_{\mathrm{kf}}\left(1+|\mathrm{kf}|_2^2\right)^{-3n}\le5^{3n}$, as shown in the proof of (xp2).
By (FS), $\Psi(y)=\sum_{\mathrm{kf}}\mathrm{fd}_{\mathrm{kf}}\exp\left(\frac{2\pi i\,\mathrm{kf}\cdot y}{\mathrm{sl}}\right)$ on the cube.
Multiplication by $\mathrm{CQ}(y)$, which equals $1$ on $\operatorname{supp}\Psi$ and $0$ outside the cube, gives
\[
\Psi(y)=\sum_{\mathrm{kf}}\mathrm{fd}_{\mathrm{kf}}\cdot\mathrm{CQ}(y)\cdot\exp\left(\frac{2\pi i\,\mathrm{kf}\cdot y}{\mathrm{sl}}\right)\qquad\text{for \emph{all} }y\in\mathbb{R}^{3n}.
\]
For $\mathrm{ki}\in\mathbb{Z}^3$ let
\[
\mathrm{sf}_{a,\mathrm{ki}}(y_a):=\mathrm{cq}_a(y_a)\cdot\exp\left(\frac{2\pi i\,\mathrm{ki}\cdot y_a}{\mathrm{sl}}\right),
\]
an element of $C_c^\infty(O;\mathbb{C})$.
With these functions the term of index $\mathrm{kf}$ is $\mathrm{fd}_{\mathrm{kf}}\cdot\mathrm{sf}_{1,\mathrm{kf}_1}\otimes\dots\otimes\mathrm{sf}_{n,\mathrm{kf}_n}$.

\emph{Norms.}
Let $|\alpha|\le m$.
Leibniz' rule gives
\[
\partial^{\alpha}\left(\mathrm{CQ}(y)\cdot\exp\left(\frac{2\pi i\,\mathrm{kf}\cdot y}{\mathrm{sl}}\right)\right)
=\sum_{\beta\le\alpha}\binom{\alpha}{\beta}\cdot\left(\partial^{\alpha-\beta}\mathrm{CQ}\right)(y)\cdot\left(\frac{2\pi i\,\mathrm{kf}}{\mathrm{sl}}\right)^{\beta}\cdot\exp\left(\frac{2\pi i\,\mathrm{kf}\cdot y}{\mathrm{sl}}\right),
\]
and here
\[
\left|\left(\frac{2\pi i\,\mathrm{kf}}{\mathrm{sl}}\right)^{\beta}\right|\le\max\left\{1,\frac{2\pi}{\mathrm{sl}}\right\}^{m}\cdot\left(1+|\mathrm{kf}|_2^2\right)^{\frac{m}{2}},
\qquad
\sum_{\beta\le\alpha}\binom{\alpha}{\beta}\le2^{m}.
\]
On the support of $\mathrm{CQ}$ one has $|y|_2^2\le\frac{3n\,\mathrm{sl}^2}{4}$.
Consequently
\[
q_m\left(\mathrm{sf}_{1,\mathrm{kf}_1}\otimes\cdots\otimes\mathrm{sf}_{n,\mathrm{kf}_n}\right)
\le2^{m}\cdot\max_{|\beta|\le m}\sup\left|\partial^{\beta}\mathrm{CQ}\right|\cdot\max\left\{1,\frac{2\pi}{\mathrm{sl}}\right\}^{m}\cdot\left(1+\frac{3n\,\mathrm{sl}^2}{4}\right)^{\frac{m}{2}}\cdot\left(1+|\mathrm{kf}|_2^2\right)^{\frac{m}{2}}.
\]
Take $\mathrm{pw}=m+3n$ in (2.6); then $\mathrm{pw}-\frac{m}{2}\ge3n$, and the series
\[
\sum_{\mathrm{kf}}\left|\mathrm{fd}_{\mathrm{kf}}\right|\cdot q_m\left(\mathrm{sf}_{1,\mathrm{kf}_1}\otimes\cdots\otimes\mathrm{sf}_{n,\mathrm{kf}_n}\right)
\]
is bounded by a finite multiple of $\sum_{\mathrm{kf}}\left(1+|\mathrm{kf}|_2^2\right)^{-3n}\le5^{3n}$.
By (abs), the series $\sum_{\mathrm{kf}}\mathrm{fd}_{\mathrm{kf}}\cdot\mathrm{sf}_{1,\mathrm{kf}_1}\otimes\cdots\otimes\mathrm{sf}_{n,\mathrm{kf}_n}$ therefore converges to $\Psi$ in every norm $q_m$.
\qed

\subsubsection*{2.3. Proof of (a)}

\begin{proof}[Proof of Proposition~45.5\,(a)]
\emph{The bound.}
Since $\frac{1}{2}=n\cdot\frac{1}{2n}$, H\"older's inequality for $n$ factors, each taken with the exponent $2n$, yields for $\mathrm{Ga}_1,\dots,\mathrm{Ga}_n\in L^{2n}(\mathrm{lw})$
\begin{equation}\tag{2.7}
\left\|\mathrm{Ga}_1\cdots\mathrm{Ga}_n\right\|_{L^2(\mathrm{lw})}\le\prod_{a=1}^{n}\left\|\mathrm{Ga}_a\right\|_{L^{2n}(\mathrm{lw})}.
\end{equation}
As $\mathrm{lw}$ has (G), the bound (qc5) holds; with $2n$ in place of $n$ it reads
\[
\left\|X(f)\right\|_{L^{2n}(\mathrm{lw})}\le\left(2n\cdot2^{n}\cdot\Gamma(n)\right)^{\frac{1}{2n}}\cdot\|f\|_2\qquad(f\in\mathcal{S}).
\]
For $f_1,\dots,f_n\in\mathcal{S}$ take $\mathrm{Ga}_a=X(f_a)$ in (2.7).
The product of the $n$ factors $\left(2n\cdot2^{n}\cdot\Gamma(n)\right)^{\frac{1}{2n}}$ is $\left(2n\cdot2^{n}\cdot\Gamma(n)\right)^{\frac{1}{2}}$, and one obtains the first assertion of (a).

\emph{Complex scalar tests.}
For $\mathrm{sf}\in\mathcal{S}_{\mathbb{C}}(\mathbb{R}^3)$ and $\mathrm{ti}\in\mathrm{TI}$ define
\[
X_{\mathrm{ti}}(\mathrm{sf}):=X\left((\operatorname{Re}\mathrm{sf})\,\mathrm{d}y_{\mathrm{ti}}\right)+i\,X\left((\operatorname{Im}\mathrm{sf})\,\mathrm{d}y_{\mathrm{ti}}\right).
\]
Every $\omega\in\mathcal{S}'$ is real-linear; hence $\mathrm{sf}\mapsto X_{\mathrm{ti}}(\mathrm{sf})$ is additive.
It is even complex-linear, since for real $t,t'$
\[
X_{\mathrm{ti}}\left((t+it')\,\mathrm{sf}\right)=X\left((t\operatorname{Re}\mathrm{sf}-t'\operatorname{Im}\mathrm{sf})\,\mathrm{d}y_{\mathrm{ti}}\right)+i\,X\left((t\operatorname{Im}\mathrm{sf}+t'\operatorname{Re}\mathrm{sf})\,\mathrm{d}y_{\mathrm{ti}}\right)=(t+it')\cdot X_{\mathrm{ti}}(\mathrm{sf}).
\]
The form $\mathrm{d}y_{\mathrm{ti}}$ is a unit vector of $\Lambda^2$, so that $\left\|(\operatorname{Re}\mathrm{sf})\,\mathrm{d}y_{\mathrm{ti}}\right\|_2=\|\operatorname{Re}\mathrm{sf}\|_{L^2(\mathbb{R}^3)}$, and likewise for the imaginary part.
Minkowski's inequality in $L^{2n}(\mathrm{lw})$, the bound (qc5) with $2n$ in place of $n$ for the two real forms $(\operatorname{Re}\mathrm{sf})\,\mathrm{d}y_{\mathrm{ti}}$ and $(\operatorname{Im}\mathrm{sf})\,\mathrm{d}y_{\mathrm{ti}}$, and the inequality $\|\operatorname{Re}\mathrm{sf}\|_{L^2(\mathbb{R}^3)}+\|\operatorname{Im}\mathrm{sf}\|_{L^2(\mathbb{R}^3)}\le\sqrt{2}\cdot\|\mathrm{sf}\|_{L^2(\mathbb{R}^3)}$ together give
\begin{equation}\tag{2.8}
\left\|X_{\mathrm{ti}}(\mathrm{sf})\right\|_{L^{2n}(\mathrm{lw})}\le\left(2n\cdot2^{n}\cdot\Gamma(n)\right)^{\frac{1}{2n}}\cdot\sqrt{2}\cdot\|\mathrm{sf}\|_{L^2(\mathbb{R}^3)}.
\end{equation}
For $f=\sum_{\mathrm{ti}}f_{\mathrm{ti}}\,\mathrm{d}y_{\mathrm{ti}}\in\mathcal{S}$ one has $X(f)=\sum_{\mathrm{ti}}X_{\mathrm{ti}}(f_{\mathrm{ti}})$.

\emph{Definition of $X^{(n)}$.}
Let $n\ge1$ and $\Psi=\sum_{\mathrm{tj}}\Psi_{\mathrm{tj}}\,\mathrm{d}y_{\mathrm{tj}}\in\mathrm{Sw}_n$, and write $\mathrm{tj}=(\mathrm{ti}_1,\dots,\mathrm{ti}_n)$.
Put
\[
X^{(n)}(\Psi):=\sum_{\mathrm{tj}\in\mathrm{TI}^n}\ \sum_{\mathrm{zl},\mathrm{kf}\in(\mathbb{Z}^3)^n}\mathrm{fc}_{\mathrm{zl},\mathrm{kf}}(\Psi_{\mathrm{tj}})\cdot\prod_{a=1}^{n}X_{\mathrm{ti}_a}\left(\mathrm{bfn}_{\mathrm{zl}_a,\mathrm{kf}_a}\right).
\]
By (2.7), (2.8) and (xp4),
\[
\left\|\prod_{a=1}^{n}X_{\mathrm{ti}_a}\left(\mathrm{bfn}_{\mathrm{zl}_a,\mathrm{kf}_a}\right)\right\|_{L^2(\mathrm{lw})}\le\left(2n\cdot2^{n}\cdot\Gamma(n)\right)^{\frac{1}{2}}\cdot\left(8\sqrt{2}\right)^{n},
\]
and by (xp2) with $m=0$ together with (xp4),
\[
\sum_{\mathrm{zl},\mathrm{kf}}\left|\mathrm{fc}_{\mathrm{zl},\mathrm{kf}}(\Psi_{\mathrm{tj}})\right|\le\mathrm{cf}_{n,0}\cdot q_{6n}(\Psi_{\mathrm{tj}})\le\mathrm{cf}_{n,0}\cdot q_{6n}(\Psi).
\]
Hence the series that defines $X^{(n)}(\Psi)$ converges absolutely in the Banach space $L^2(\mathrm{lw})$, and, since $\#|\mathrm{TI}^n|=3^n$,
\begin{equation}\tag{2.9}
\left\|X^{(n)}(\Psi)\right\|_{L^2(\mathrm{lw})}\le\mathrm{cf}_n\cdot q_{6n}(\Psi),\qquad\mathrm{cf}_n:=\left(2n\cdot2^{n}\cdot\Gamma(n)\right)^{\frac{1}{2}}\cdot\left(24\sqrt{2}\right)^{n}\cdot\mathrm{cf}_{n,0}.
\end{equation}
The map $X^{(n)}$ is linear, because $\Psi\mapsto\Psi_{\mathrm{tj}}$ and the functionals $\mathrm{fc}_{\mathrm{zl},\mathrm{kf}}$ are linear (xp1).
By (2.9) and the paragraph \emph{Topology} of 2.1, it is continuous.

\emph{Values on elementary tensors.}
For $\mathrm{sf}_1,\dots,\mathrm{sf}_n\in\mathcal{S}_{\mathbb{C}}(\mathbb{R}^3)$ and $\mathrm{tj}\in\mathrm{TI}^n$ one has
\begin{equation}\tag{2.10}
X^{(n)}\left((\mathrm{sf}_1\otimes\cdots\otimes\mathrm{sf}_n)\,\mathrm{d}y_{\mathrm{tj}}\right)=\prod_{a=1}^{n}X_{\mathrm{ti}_a}(\mathrm{sf}_a).
\end{equation}

\emph{Proof of (2.10).}
Fix $a$.
The Expansion Lemma is applied here on $\mathbb{R}^3$, that is, with $1$ in place of its $n$, whereas the exponent $2n$ in (2.7) and (2.8) is left unchanged.
By (xp2) on $\mathbb{R}^3$ with $m=0$ and by (xp4),
\[
\sum_{\mathrm{zl}_a,\mathrm{kf}_a\in\mathbb{Z}^3}\left|\mathrm{fc}_{\mathrm{zl}_a,\mathrm{kf}_a}(\mathrm{sf}_a)\right|\le\mathrm{cf}_{1,0}\cdot q_6(\mathrm{sf}_a),
\]
and by (2.8) and (xp4),
\[
\left\|X_{\mathrm{ti}_a}\left(\mathrm{bfn}_{\mathrm{zl}_a,\mathrm{kf}_a}\right)\right\|_{L^{2n}(\mathrm{lw})}\le\left(2n\cdot2^{n}\cdot\Gamma(n)\right)^{\frac{1}{2n}}\cdot8\sqrt{2}.
\]
The series
\[
\sum_{\mathrm{zl}_a,\mathrm{kf}_a\in\mathbb{Z}^3}\mathrm{fc}_{\mathrm{zl}_a,\mathrm{kf}_a}(\mathrm{sf}_a)\cdot X_{\mathrm{ti}_a}\left(\mathrm{bfn}_{\mathrm{zl}_a,\mathrm{kf}_a}\right)
\]
therefore converges absolutely in $L^{2n}(\mathrm{lw})$.
Its sum is $X_{\mathrm{ti}_a}(\mathrm{sf}_a)$.
Indeed, let $\mathrm{Fs}$ be a finite subset of $\mathbb{Z}^3\times\mathbb{Z}^3$, and let $\mathrm{sf}'$ be the sum of the functions $\mathrm{fc}_{\mathrm{zl}_a,\mathrm{kf}_a}(\mathrm{sf}_a)\cdot\mathrm{bfn}_{\mathrm{zl}_a,\mathrm{kf}_a}$ over $(\mathrm{zl}_a,\mathrm{kf}_a)\in\mathrm{Fs}$.
By the linearity of $X_{\mathrm{ti}_a}$, the difference between $X_{\mathrm{ti}_a}(\mathrm{sf}_a)$ and the partial sum over $\mathrm{Fs}$ is $X_{\mathrm{ti}_a}(\mathrm{sf}_a-\mathrm{sf}')$.
By (2.8) and (in4), its norm in $L^{2n}(\mathrm{lw})$ does not exceed
\[
\left(2n\cdot2^{n}\cdot\Gamma(n)\right)^{\frac{1}{2n}}\cdot\sqrt{2}\cdot\pi\cdot q_2(\mathrm{sf}_a-\mathrm{sf}'),
\]
and this tends to $0$ as $\mathrm{Fs}$ exhausts $\mathbb{Z}^3\times\mathbb{Z}^3$, by (xp3) on $\mathbb{R}^3$.
By (2.7), the multiplication $(\mathrm{Ga}_1,\dots,\mathrm{Ga}_n)\mapsto\mathrm{Ga}_1\cdots\mathrm{Ga}_n$ is an $n$-linear map $L^{2n}(\mathrm{lw})^n\to L^2(\mathrm{lw})$ of norm at most $1$.
Hence the product of the $n$ sums is the sum of the $n$-fold series, and this series converges absolutely in $L^2(\mathrm{lw})$:
\begin{align*}
\prod_{a=1}^{n}X_{\mathrm{ti}_a}(\mathrm{sf}_a)
&=\sum_{\mathrm{zl},\mathrm{kf}\in(\mathbb{Z}^3)^n}\ \prod_{a=1}^{n}\left(\mathrm{fc}_{\mathrm{zl}_a,\mathrm{kf}_a}(\mathrm{sf}_a)\cdot X_{\mathrm{ti}_a}\left(\mathrm{bfn}_{\mathrm{zl}_a,\mathrm{kf}_a}\right)\right)\\
&=\sum_{\mathrm{zl},\mathrm{kf}\in(\mathbb{Z}^3)^n}\mathrm{fc}_{\mathrm{zl},\mathrm{kf}}(\mathrm{sf}_1\otimes\cdots\otimes\mathrm{sf}_n)\cdot\prod_{a=1}^{n}X_{\mathrm{ti}_a}\left(\mathrm{bfn}_{\mathrm{zl}_a,\mathrm{kf}_a}\right),
\end{align*}
the second equality being (xp1).
Every component of the element $(\mathrm{sf}_1\otimes\cdots\otimes\mathrm{sf}_n)\,\mathrm{d}y_{\mathrm{tj}}$ of $\mathrm{Sw}_n$ other than the one at the index $\mathrm{tj}$ vanishes, and the component at $\mathrm{tj}$ is $\mathrm{sf}_1\otimes\cdots\otimes\mathrm{sf}_n$.
So the right-hand side is $X^{(n)}\left((\mathrm{sf}_1\otimes\cdots\otimes\mathrm{sf}_n)\,\mathrm{d}y_{\mathrm{tj}}\right)$, and (2.10) is proved.

\emph{Values on $f_1\otimes\cdots\otimes f_n$.}
Let $f_a=\sum_{\mathrm{ti}}f_{a,\mathrm{ti}}\,\mathrm{d}y_{\mathrm{ti}}\in\mathcal{S}$ for $a=1,\dots,n$.
By (2.1), the linearity of $X^{(n)}$ and (2.10),
\[
X^{(n)}(f_1\otimes\cdots\otimes f_n)=\sum_{\mathrm{tj}}\prod_{a=1}^{n}X_{\mathrm{ti}_a}\left(f_{a,\mathrm{ti}_a}\right)=\prod_{a=1}^{n}\left(\sum_{\mathrm{ti}\in\mathrm{TI}}X_{\mathrm{ti}}\left(f_{a,\mathrm{ti}}\right)\right)=\prod_{a=1}^{n}X(f_a).
\]

\emph{Uniqueness.}
Let $\mathrm{Lm}\colon\mathrm{Sw}_n\to L^2(\mathrm{lw})$ be a continuous linear map with $\mathrm{Lm}(f_1\otimes\cdots\otimes f_n)=X(f_1)\cdots X(f_n)$ for all $f_1,\dots,f_n\in\mathcal{S}$.
For $\mathrm{sf}_a\in\mathcal{S}_{\mathbb{C}}(\mathbb{R}^3)$ write $\mathrm{sf}_a^{(0)}:=\operatorname{Re}\mathrm{sf}_a$ and $\mathrm{sf}_a^{(1)}:=\operatorname{Im}\mathrm{sf}_a$.
Then
\[
(\mathrm{sf}_1\otimes\cdots\otimes\mathrm{sf}_n)\,\mathrm{d}y_{\mathrm{tj}}=\sum_{\mathrm{bi}\in\{0,1\}^n}i^{\mathrm{bi}_1+\cdots+\mathrm{bi}_n}\cdot\left(\mathrm{sf}_1^{(\mathrm{bi}_1)}\,\mathrm{d}y_{\mathrm{ti}_1}\right)\otimes\cdots\otimes\left(\mathrm{sf}_n^{(\mathrm{bi}_n)}\,\mathrm{d}y_{\mathrm{ti}_n}\right)
\]
is a linear combination, with complex coefficients, of tensor products of elements of $\mathcal{S}$, and the linearity of $\mathrm{Lm}$ gives
\[
\mathrm{Lm}\left((\mathrm{sf}_1\otimes\cdots\otimes\mathrm{sf}_n)\,\mathrm{d}y_{\mathrm{tj}}\right)=\sum_{\mathrm{bi}\in\{0,1\}^n}i^{\mathrm{bi}_1+\cdots+\mathrm{bi}_n}\cdot\prod_{a=1}^{n}X\left(\mathrm{sf}_a^{(\mathrm{bi}_a)}\,\mathrm{d}y_{\mathrm{ti}_a}\right)=\prod_{a=1}^{n}X_{\mathrm{ti}_a}(\mathrm{sf}_a).
\]
By (2.10), $\mathrm{Lm}$ and $X^{(n)}$ therefore agree on the elementary tensors.
Both maps are linear and continuous, and by (2.5) every $\Psi\in\mathrm{Sw}_n$ is the sum of a series of multiples of elementary tensors that converges in every norm $q_m$.
Hence $\mathrm{Lm}(\Psi)=X^{(n)}(\Psi)$ for all $\Psi\in\mathrm{Sw}_n$.

The two remaining clauses of (a), namely $X^{(0)}(\Psi_0):=\Psi_0$ for $\Psi_0\in\mathbb{C}$ and $S_n(\Psi):=\mathrm{lw}\left(X^{(n)}(\Psi)\right)$, are definitions.
The integral exists because $L^2(\mathrm{lw})\subset L^1(\mathrm{lw})$.
\end{proof}

\subsubsection*{2.4. Proof of (b)}

\begin{proof}[Proof of Proposition~45.5\,(b)]
\emph{Products.}
If $m=0$ or $n=0$, the identity $X^{(m)}(\Phi)\cdot X^{(n)}(\Psi)=X^{(m+n)}(\Phi\otimes\Psi)$ is the linearity of $X^{(n)}$, respectively of $X^{(m)}$.
Let $m,n\ge1$, $\Phi\in\mathrm{Sw}_m$ and $\Psi\in\mathrm{Sw}_n$.
The Cauchy--Schwarz inequality and (2.9) give (2.11); the inequality $\|\cdot\|_{L^1(\mathrm{lw})}\le\|\cdot\|_{L^2(\mathrm{lw})}$, (2.9) and (in1) give (2.12):
\begin{align}
\left\|X^{(m)}(\Phi)\cdot X^{(n)}(\Psi)\right\|_{L^1(\mathrm{lw})}&\le\mathrm{cf}_m\cdot\mathrm{cf}_n\cdot q_{6m}(\Phi)\cdot q_{6n}(\Psi),\tag{2.11}\\
\left\|X^{(m+n)}(\Phi\otimes\Psi)\right\|_{L^1(\mathrm{lw})}&\le\mathrm{cf}_{m+n}\cdot q_{6(m+n)}(\Phi)\cdot q_{6(m+n)}(\Psi).\tag{2.12}
\end{align}
Thus both sides of the identity are bilinear maps $\mathrm{Sw}_m\times\mathrm{Sw}_n\to L^1(\mathrm{lw})$, continuous in each variable.
Let $\Phi=(\mathrm{sf}'_1\otimes\cdots\otimes\mathrm{sf}'_m)\,\mathrm{d}y_{\mathrm{tj}'}$, with $\mathrm{tj}'=(\mathrm{ti}'_1,\dots,\mathrm{ti}'_m)$, and $\Psi=(\mathrm{sf}_1\otimes\cdots\otimes\mathrm{sf}_n)\,\mathrm{d}y_{\mathrm{tj}}$ be elementary tensors.
The tensor
\[
\Phi\otimes\Psi=(\mathrm{sf}'_1\otimes\cdots\otimes\mathrm{sf}'_m\otimes\mathrm{sf}_1\otimes\cdots\otimes\mathrm{sf}_n)\,\mathrm{d}y_{(\mathrm{tj}',\mathrm{tj})}
\]
is again elementary, and by (2.10) both sides of the identity are equal to
\[
\prod_{a=1}^{m}X_{\mathrm{ti}'_a}(\mathrm{sf}'_a)\cdot\prod_{a=1}^{n}X_{\mathrm{ti}_a}(\mathrm{sf}_a).
\]
Fix first an elementary tensor $\Psi$.
As functions of $\Phi$, the two sides are continuous linear maps that coincide on the elementary tensors, hence on $\mathrm{Sw}_m$ by (2.5).
Fix next an arbitrary $\Phi\in\mathrm{Sw}_m$.
As functions of $\Psi$, the two sides are continuous linear maps that coincide on the elementary tensors, hence on $\mathrm{Sw}_n$ by (2.5).

\emph{Conjugation.}
Let $n\ge1$ and put $\mathrm{Lm}(\Psi):=\overline{X^{(n)}\left(\overline{\Psi}\right)}$ for $\Psi\in\mathrm{Sw}_n$.
Being composed of two antilinear maps and a linear one, $\mathrm{Lm}$ is complex-linear.
It is continuous by (2.9), because $q_m\left(\overline{\Psi}\right)=q_m(\Psi)$.
For $f_1,\dots,f_n\in\mathcal{S}$ the tensor $f_1\otimes\cdots\otimes f_n$ is real and the $X(f_a)$ are real functions; therefore
\[
\mathrm{Lm}(f_1\otimes\cdots\otimes f_n)=\overline{X(f_1)\cdots X(f_n)}=X(f_1)\cdots X(f_n).
\]
The uniqueness in (a) gives $\mathrm{Lm}=X^{(n)}$, and on replacing $\Psi$ by $\overline{\Psi}$ one obtains $\overline{X^{(n)}(\Psi)}=X^{(n)}\left(\overline{\Psi}\right)$.
For $n=0$ both sides are equal to $\overline{\Psi_0}$.

\emph{Euclidean maps.}
Let $\mathrm{Gi}$ be a Euclidean map with $\mathrm{Gi}_*\mathrm{lw}=\mathrm{lw}$.
By (qc1), $\widehat{\mathrm{Gi}}$ is unitary on $L^2(\mathrm{lw})$, with inverse $\widehat{\mathrm{Gi}^{-1}}$.
The map
\[
\mathrm{Lm}(\Psi):=\widehat{\mathrm{Gi}^{-1}}\,X^{(n)}\left((\mathrm{Gi}\cdot)^{\otimes n}\Psi\right)
\]
is linear, and it is continuous by (2.9) and (in2).
For $f_1,\dots,f_n\in\mathcal{S}$, (in2), (a) and (qc1) give
\[
X^{(n)}\left((\mathrm{Gi}\cdot)^{\otimes n}(f_1\otimes\cdots\otimes f_n)\right)=X(\mathrm{Gi}\cdot f_1)\cdots X(\mathrm{Gi}\cdot f_n)=\left(\widehat{\mathrm{Gi}}\,X(f_1)\right)\cdots\left(\widehat{\mathrm{Gi}}\,X(f_n)\right)=\widehat{\mathrm{Gi}}\left(X(f_1)\cdots X(f_n)\right).
\]
Hence $\mathrm{Lm}(f_1\otimes\cdots\otimes f_n)=X(f_1)\cdots X(f_n)$, and $\mathrm{Lm}=X^{(n)}$ by the uniqueness in (a).
Applying $\widehat{\mathrm{Gi}}$ to this identity one arrives at $X^{(n)}\left((\mathrm{Gi}\cdot)^{\otimes n}\Psi\right)=\widehat{\mathrm{Gi}}\,X^{(n)}(\Psi)$.
For $n=0$, with $(\mathrm{Gi}\cdot)^{\otimes0}:=1$ on $\mathbb{C}$, both sides are the constant $\Psi_0$.
\end{proof}

\subsubsection*{2.5. Proof of (c)}

\begin{proof}[Proof of Proposition~45.5\,(c)]
\emph{First clause.}
Let $O\subset\mathbb{R}^3$ be open and $\Psi=\sum_{\mathrm{tj}}\Psi_{\mathrm{tj}}\,\mathrm{d}y_{\mathrm{tj}}$ with $\Psi_{\mathrm{tj}}\in C_c^\infty(O^n;\mathbb{C})$.
By the Support Lemma, each $\Psi_{\mathrm{tj}}$ is the sum of a series, convergent in every norm $q_m$, whose terms are multiples of functions $\mathrm{sf}_1\otimes\cdots\otimes\mathrm{sf}_n$ with $\mathrm{sf}_a\in C_c^\infty(O;\mathbb{C})$.
For such functions (2.10) gives
\[
X^{(n)}\left((\mathrm{sf}_1\otimes\cdots\otimes\mathrm{sf}_n)\,\mathrm{d}y_{\mathrm{tj}}\right)=\prod_{a=1}^{n}\left(X\left((\operatorname{Re}\mathrm{sf}_a)\,\mathrm{d}y_{\mathrm{ti}_a}\right)+i\,X\left((\operatorname{Im}\mathrm{sf}_a)\,\mathrm{d}y_{\mathrm{ti}_a}\right)\right).
\]
The forms $(\operatorname{Re}\mathrm{sf}_a)\,\mathrm{d}y_{\mathrm{ti}_a}$ and $(\operatorname{Im}\mathrm{sf}_a)\,\mathrm{d}y_{\mathrm{ti}_a}$ are real elements of $C_c^\infty(O;\Lambda^2)$.
Hence this product is $\Sigma_O$-measurable, by the definition of $\Sigma_O$ in (n1), and it lies in $L^2(\mathrm{lw})$ by (2.7) and (2.8).
Since $X^{(n)}$ is linear and continuous, $X^{(n)}(\Psi)$ is a limit in $L^2(\mathrm{lw})$ of finite linear combinations of such products.
As $L^2(\Sigma_O,\mathrm{lw})$ is a closed subspace of $L^2(\mathrm{lw})$ (n1), $X^{(n)}(\Psi)$ belongs to it.

\emph{Second clause.}
Let $\hat{n}$ be a unit vector and $\mathrm{po}\in\mathbb{R}$, put $O:=\{y\in\mathbb{R}^3:\hat{n}\cdot y>\mathrm{po}\}$, and let $\Psi\in\mathrm{Sw}_n$ have its support in $\{\hat{n}\cdot y_a\ge\mathrm{po}\text{ for all }a\}$.
For $k\ge1$ set
\[
\Psi_k(y):=\mathrm{CO}_n\left(\frac{y}{k}\right)\cdot\left(\tau_{\hat{n}/k}^{\otimes n}\Psi\right)(y)=\mathrm{CO}_n\left(\frac{y}{k}\right)\cdot\Psi\left(y_1-\frac{\hat{n}}{k},\dots,y_n-\frac{\hat{n}}{k}\right).
\]
If $\Psi\left(y_1-\frac{\hat{n}}{k},\dots,y_n-\frac{\hat{n}}{k}\right)\ne0$, then $\hat{n}\cdot\left(y_a-\frac{\hat{n}}{k}\right)\ge\mathrm{po}$, that is $\hat{n}\cdot y_a\ge\mathrm{po}+\frac{1}{k}$, for every $a$.
So the support of $\Psi_k$ is contained in the closed set $\left\{\hat{n}\cdot y_a\ge\mathrm{po}+\frac{1}{k}\text{ for all }a\right\}$ and, since $\mathrm{CO}_n$ vanishes outside $\left[-\frac{7}{4},\frac{7}{4}\right]^{3n}$, in the compact set $[-2k,2k]^{3n}$.
Being closed and bounded, it is a compact subset of $O^n$, and $\Psi_k\in C_c^\infty(O^n)$, with values in $(\Lambda^2)^{\otimes n}\otimes\mathbb{C}$.
By the first clause, $X^{(n)}(\Psi_k)\in L^2(\Sigma_O,\mathrm{lw})$.
Moreover
\[
\Psi_k-\Psi=\mathrm{CO}_n\left(\frac{\cdot}{k}\right)\cdot\left(\tau_{\hat{n}/k}^{\otimes n}\Psi-\Psi\right)-\left(1-\mathrm{CO}_n\left(\frac{\cdot}{k}\right)\right)\cdot\Psi,
\]
and (in5), together with (in3) for $b=\frac{\hat{n}}{k}$, $|b|_2=\frac{1}{k}$, gives for every $m$
\[
q_m(\Psi_k-\Psi)\le2^{m}\cdot\mathrm{cm}_m\cdot\left(\sqrt{3}\cdot n\cdot k^{-1}\cdot\left(2(1+n)\right)^{\frac{m}{2}}+\left(1+k^2\right)^{-\frac{1}{2}}\right)\cdot q_{m+1}(\Psi)\longrightarrow0\qquad(k\to\infty).
\]
By the continuity of $X^{(n)}$, $X^{(n)}(\Psi_k)\to X^{(n)}(\Psi)$ in $L^2(\mathrm{lw})$.
The subspace $L^2(\Sigma_O,\mathrm{lw})$ is closed, so it contains $X^{(n)}(\Psi)$.
\end{proof}

\subsubsection*{2.6. Proof of (d)}

\begin{proof}[Proof of Proposition~45.5\,(d)]
Let $\mathrm{lw}$ be invariant under $r_{\hat{n},0}$ and reflection positive at $\{\hat{n}\cdot y=0\}$, and write $\mathrm{rn}:=r_{\hat{n},0}$.
Choose $\mathrm{gr}\in O(3)$ with $\mathrm{gr}\,e_1=\hat{n}$ and put $\mathrm{lw}':=(\mathrm{gr}^{-1})_*\mathrm{lw}$.
Then $\mathrm{lw}=\mathrm{gr}_*\mathrm{lw}'$ by (qc1), and the Transport Lemma applies with $\xi'=\mathrm{lw}'$ and $\xi=\mathrm{lw}$.

\emph{First inequality.}
By (tr3), $\mathrm{gr}\,\theta_1\,\mathrm{gr}^{-1}=\mathrm{rn}$, where $\theta_1=r_{e_1,0}$.
Since $\mathrm{lw}$ is invariant under $\mathrm{rn}$, (tr2) with $\mathrm{Gi}=\theta_1$ shows that $\mathrm{lw}'$ is invariant under $\theta_1$.
Since $\mathrm{lw}$ is reflection positive at $\{\hat{n}\cdot y=0\}$, (tr5) with $\mathrm{po}=0$ shows that $\mathrm{lw}'$ is reflection positive at $\{y_1=0\}$.
The implication (1)$\Rightarrow$(3) of (qc6), for the law $\mathrm{lw}'$ and $j=1$, now yields
\[
\left\langle\hat\theta_1\,\mathrm{Gb},\mathrm{Gb}\right\rangle_{\mathrm{lw}'}\ge0\qquad\text{for all }\mathrm{Gb}\in L^2\left(\Sigma_{\{y_1>0\}},\mathrm{lw}'\right).
\]
Let $\mathrm{Ga}\in L^2\left(\Sigma_{\{\hat{n}\cdot y>0\}},\mathrm{lw}\right)$.
By (tr1) and (tr3), $\mathrm{Gb}:=\widehat{\mathrm{gr}^{-1}}\,\mathrm{Ga}$ lies in $L^2\left(\Sigma_{\{y_1>0\}},\mathrm{lw}'\right)$, and $\mathrm{Ga}=\widehat{\mathrm{gr}}\,\mathrm{Gb}$.
By (tr2), $\widehat{\mathrm{rn}}\,\widehat{\mathrm{gr}}=\widehat{\mathrm{gr}}\,\hat\theta_1$.
Since $\widehat{\mathrm{gr}}$ is a unitary map of $L^2(\mathrm{lw}')$ onto $L^2(\mathrm{lw})$ (tr1),
\[
\left\langle\widehat{\mathrm{rn}}\,\mathrm{Ga},\mathrm{Ga}\right\rangle_{\mathrm{lw}}=\left\langle\widehat{\mathrm{rn}}\,\widehat{\mathrm{gr}}\,\mathrm{Gb},\widehat{\mathrm{gr}}\,\mathrm{Gb}\right\rangle_{\mathrm{lw}}=\left\langle\widehat{\mathrm{gr}}\,\hat\theta_1\,\mathrm{Gb},\widehat{\mathrm{gr}}\,\mathrm{Gb}\right\rangle_{\mathrm{lw}}=\left\langle\hat\theta_1\,\mathrm{Gb},\mathrm{Gb}\right\rangle_{\mathrm{lw}'}\ge0.
\]

\emph{Second inequality.}
Let $\Psi_0\in\mathbb{C}$, and let $\Psi_n\in\mathrm{Sw}_n$, $1\le n\le s$, be supported in $\{\hat{n}\cdot y_a\ge0\text{ for all }a\}$.
Put $\mathrm{Ga}:=\sum_{n=0}^{s}X^{(n)}(\Psi_n)$.
By (c) with $\mathrm{po}=0$, each $X^{(n)}(\Psi_n)$ with $n\ge1$ lies in $L^2\left(\Sigma_{\{\hat{n}\cdot y>0\}},\mathrm{lw}\right)$, and so does the constant $X^{(0)}(\Psi_0)=\Psi_0$.
By (b), applied with $\mathrm{Gi}=\mathrm{rn}$ (note that $\mathrm{rn}_*\mathrm{lw}=\mathrm{lw}$), and by the definition of $\Theta^{\hat{n}}_m$ in (n5) and (n8),
\[
\overline{\widehat{\mathrm{rn}}\,X^{(m)}(\Psi_m)}=\overline{X^{(m)}\left((\mathrm{rn}\cdot)^{\otimes m}\Psi_m\right)}=X^{(m)}\left(\overline{(\mathrm{rn}\cdot)^{\otimes m}\Psi_m}\right)=X^{(m)}\left(\Theta^{\hat{n}}_m\Psi_m\right),
\]
and
\[
X^{(m)}\left(\Theta^{\hat{n}}_m\Psi_m\right)\cdot X^{(n)}(\Psi_n)=X^{(m+n)}\left(\Theta^{\hat{n}}_m\Psi_m\otimes\Psi_n\right)\qquad\text{in }L^1(\mathrm{lw}).
\]
Together with the first inequality this gives
\begin{align*}
0\le\left\langle\widehat{\mathrm{rn}}\,\mathrm{Ga},\mathrm{Ga}\right\rangle_{\mathrm{lw}}
&=\sum_{m,n=0}^{s}\mathrm{lw}\left(\overline{\widehat{\mathrm{rn}}\,X^{(m)}(\Psi_m)}\cdot X^{(n)}(\Psi_n)\right)\\
&=\sum_{m,n=0}^{s}\mathrm{lw}\left(X^{(m+n)}\left(\Theta^{\hat{n}}_m\Psi_m\otimes\Psi_n\right)\right)\\
&=\sum_{m,n=0}^{s}S_{m+n}\left(\Theta^{\hat{n}}_m\Psi_m\otimes\Psi_n\right).
\end{align*}
This is the second inequality.
\end{proof}

\subsection*{3. Proposition 45.7. One additional mirror passes to the components}

Let $\mathrm{lw}$ be a law in the class $\mathcal{C}$ of \cite[Section~8.1.3]{QEDI}, let $\nu$ be the measure of Proposition~45.6 for $\mathrm{lw}$, and let $\hat{n}$ be a unit vector such that $(r_{\hat{n},0})_*\mathrm{lw}=\mathrm{lw}$ and $\mathrm{lw}$ is reflection positive at $\{\hat{n}\cdot y=0\}$.
These are the hypotheses of Proposition~45.7.

\subsubsection*{Proof of (a), (b), (c)}

\begin{proof}[Proof of Proposition~45.7\,(a), (b), (c)]
The hypothesis that $\mathrm{lw}$ lies in the class $\mathcal{C}$ enters the following argument only through the invariance of $\mathrm{lw}$ under $\mathbb{R}^3$.

\textbf{Step 1 (transport).}
By the Transport Lemma there is $\mathrm{gr}\in O(3)$ with $\mathrm{gr}\,e_1=\hat{n}$.
Fix such a $\mathrm{gr}$ and put $\mathrm{lw}':=(\mathrm{gr}^{-1})_*\mathrm{lw}$.
Then $\mathrm{lw}=\mathrm{gr}_*\mathrm{lw}'$ by (qc1), and the Transport Lemma applies with $\xi'=\mathrm{lw}'$ and $\xi=\mathrm{lw}$.
It yields three properties of $\mathrm{lw}'$:
\begin{itemize}
\item $\mathrm{lw}'$ is invariant under $\mathbb{R}^3$, by (tr4), since $\mathrm{lw}$ is;
\item $\mathrm{lw}'$ is invariant under $\theta_1=r_{e_1,0}$, by (tr2) with $\mathrm{Gi}=\theta_1$ and the identity $\mathrm{gr}\,\theta_1\,\mathrm{gr}^{-1}=r_{\hat{n},0}$ of (tr3), since $\mathrm{lw}$ is invariant under $r_{\hat{n},0}$;
\item $\mathrm{lw}'$ is reflection positive at $\{y_1=0\}$, by (tr5) with $\mathrm{po}=0$, since $\mathrm{lw}$ is reflection positive at $\{\hat{n}\cdot y=0\}$.
\end{itemize}

\textbf{Step 2 (the two decomposition measures).}
Since $\mathrm{lw}'$ is invariant under $\mathbb{R}^3$ (Step 1), Proposition~45.6\,(a$''$) provides its decomposition measure $\nu'$.
One has $\nu'(\mathrm{Law}_{\mathrm{erg}})=1$ and
\[
\mathrm{lw}'(\mathrm{Ev})=\int\mathrm{d}\nu'(\xi')\,\xi'(\mathrm{Ev})\qquad(\mathrm{Ev}\in\Sigma).
\]
Let $\mathrm{Tg}$ be the map of (tr6), and let $\nu''$ be the image of $\nu'$ under $\mathrm{Tg}$, a probability measure on $(\mathrm{Law},B_{\mathrm{Law}})$.
By (tr6),
\[
\nu''(\mathrm{Law}_{\mathrm{erg}})=\nu'\left(\mathrm{Tg}^{-1}\mathrm{Law}_{\mathrm{erg}}\right)=\nu'(\mathrm{Law}_{\mathrm{erg}})=1.
\]
For $\mathrm{Ev}\in\Sigma$ the set $\mathrm{gr}^{-1}\mathrm{Ev}$ lies in $\Sigma$ (qc1), and
\[
\mathrm{lw}(\mathrm{Ev})=(\mathrm{gr}_*\mathrm{lw}')(\mathrm{Ev})=\mathrm{lw}'\left(\mathrm{gr}^{-1}\mathrm{Ev}\right)=\int\mathrm{d}\nu'(\xi')\,\xi'\left(\mathrm{gr}^{-1}\mathrm{Ev}\right)=\int\mathrm{d}\nu'(\xi')\,\mathrm{Tg}(\xi')(\mathrm{Ev})=\int\mathrm{d}\nu''(\xi)\,\xi(\mathrm{Ev}),
\]
where the last equality is the change of variables for the image measure $\nu''$.
Thus $\nu''$ is carried by the $\mathbb{R}^3$-ergodic laws and represents $\mathrm{lw}$.
The uniqueness in Proposition~45.6\,(a$''$), applied to the law $\mathrm{lw}$, gives $\nu=\nu''$, that is,
\begin{equation}\tag{3.1}
\nu(\mathrm{Ls})=\nu'\left(\mathrm{Tg}^{-1}\mathrm{Ls}\right)\qquad(\mathrm{Ls}\in B_{\mathrm{Law}}).
\end{equation}

\textbf{Step 3 (a set of full measure).}
By Step 1, the law $\mathrm{lw}'$ and the index $j=1$ satisfy the hypotheses of Proposition~45.6\,(b).
By that statement, read with the meaning of ``almost every'' fixed in (n7), there is a set $\mathrm{Ls}'_0\in B_{\mathrm{Law}}$ with $\nu'(\mathrm{Ls}'_0)=1$ such that every $\xi'\in\mathrm{Ls}'_0$ is invariant under $\theta_1$ and reflection positive at every plane $\{y_1=\mathrm{po}\}$, $\mathrm{po}\in\mathbb{R}$.
Put $\mathrm{Ls}':=\mathrm{Ls}'_0\cap\mathrm{Law}_{\mathrm{erg}}$.
This set belongs to $B_{\mathrm{Law}}$ (qc7), and $\nu'(\mathrm{Ls}')=1$ (Step 2).
By (n3), every $\xi'\in\mathrm{Ls}'$ is, in addition, invariant under $\mathbb{R}^3$.
Put $\mathrm{Ls}:=\mathrm{Tg}(\mathrm{Ls}')$.
Since $\mathrm{Tg}^{-1}$ is measurable (tr6), $\mathrm{Ls}$ belongs to $B_{\mathrm{Law}}$, and (3.1) gives
\[
\nu(\mathrm{Ls})=\nu'\left(\mathrm{Tg}^{-1}\mathrm{Tg}\,\mathrm{Ls}'\right)=\nu'(\mathrm{Ls}')=1.
\]

\textbf{Step 4 ((a) and (b)).}
Let $\xi\in\mathrm{Ls}$, so that $\xi=\mathrm{gr}_*\xi'$ with $\xi'\in\mathrm{Ls}'$.
The Transport Lemma, applied to this pair, shows:
\begin{itemize}
\item $\xi$ is invariant under $\mathbb{R}^3$ (tr4);
\item $\xi$ is invariant under $\mathrm{gr}\,\theta_1\,\mathrm{gr}^{-1}=r_{\hat{n},0}$, by (tr2) with $\mathrm{Gi}=\theta_1$ and by (tr3);
\item for every real $\mathrm{po}$, $\xi$ is reflection positive at $\{\hat{n}\cdot y=\mathrm{po}\}$, by (tr5), because $\xi'$ is reflection positive at $\{y_1=\mathrm{po}\}$ (Step 3).
\end{itemize}
For real $\mathrm{po}$, the last identity of (tr3) and (qc1) give
\[
(r_{\hat{n},\mathrm{po}})_*\xi=(\tau_{\mathrm{po}\cdot\hat{n}})_*(r_{\hat{n},0})_*(\tau_{-\mathrm{po}\cdot\hat{n}})_*\xi=\xi.
\]
Since $\nu(\mathrm{Ls})=1$, these properties are the assertions (a) and (b).

\textbf{Step 5 ((c)).}
By Step 1, $\mathrm{lw}'$ satisfies the hypotheses of (qc9) with $j=1$.
Hence $U(s\cdot e_1)\to P^{1}_{\mathrm{lw}'}$ weakly on $L^2(\mathrm{lw}')$ as $|s|\to\infty$.
By (tr1), $\widehat{\mathrm{gr}}$ is a unitary map of $L^2(\mathrm{lw}')$ onto $L^2(\mathrm{lw})$ with inverse $\widehat{\mathrm{gr}^{-1}}$, and by (tr2), (tr3)
\[
\widehat{\mathrm{gr}}\,U(s\cdot e_1)\,\widehat{\mathrm{gr}^{-1}}=U(s\cdot\mathrm{gr}\,e_1)=U(s\cdot\hat{n}).
\]
For $\mathrm{Ga},\mathrm{Gb}\in L^2(\mathrm{lw})$ it follows that, as $|s|\to\infty$,
\begin{align*}
\left\langle\mathrm{Ga},U(s\cdot\hat{n})\,\mathrm{Gb}\right\rangle_{\mathrm{lw}}
&=\left\langle\widehat{\mathrm{gr}^{-1}}\,\mathrm{Ga},U(s\cdot e_1)\,\widehat{\mathrm{gr}^{-1}}\,\mathrm{Gb}\right\rangle_{\mathrm{lw}'}\\
&\longrightarrow\left\langle\widehat{\mathrm{gr}^{-1}}\,\mathrm{Ga},P^{1}_{\mathrm{lw}'}\,\widehat{\mathrm{gr}^{-1}}\,\mathrm{Gb}\right\rangle_{\mathrm{lw}'}=\left\langle\mathrm{Ga},\widehat{\mathrm{gr}}\,P^{1}_{\mathrm{lw}'}\,\widehat{\mathrm{gr}^{-1}}\,\mathrm{Gb}\right\rangle_{\mathrm{lw}}.
\end{align*}
As the conjugate of an orthogonal projection by a unitary map, $\widehat{\mathrm{gr}}\,P^{1}_{\mathrm{lw}'}\,\widehat{\mathrm{gr}^{-1}}$ is the orthogonal projection of $L^2(\mathrm{lw})$ onto the image under $\widehat{\mathrm{gr}}$ of the range of $P^{1}_{\mathrm{lw}'}$.
This range is $\left\{\mathrm{Gb}'\in L^2(\mathrm{lw}'):U(s\cdot e_1)\,\mathrm{Gb}'=\mathrm{Gb}'\text{ for all }s\right\}$.
For $\mathrm{Ga}=\widehat{\mathrm{gr}}\,\mathrm{Gb}'$ one has $U(s\cdot\hat{n})\,\mathrm{Ga}=\widehat{\mathrm{gr}}\,U(s\cdot e_1)\,\mathrm{Gb}'$, so that $U(s\cdot\hat{n})\,\mathrm{Ga}=\mathrm{Ga}$ if and only if $U(s\cdot e_1)\,\mathrm{Gb}'=\mathrm{Gb}'$.
The image in question is therefore $\left\{\mathrm{Ga}\in L^2(\mathrm{lw}):U(s\cdot\hat{n})\,\mathrm{Ga}=\mathrm{Ga}\text{ for all }s\right\}$, and (c) follows.
\end{proof}

\subsubsection*{Proof of (d)}

\begin{proof}[Proof of Proposition~45.7\,(d)]
Let $\hat{n}$ be a unit vector, and let $\mathrm{Ls}_{\hat{n}}$ be the set of the laws that are invariant under $\mathbb{R}^3$ and under $r_{\hat{n},0}$ and reflection positive at $\{\hat{n}\cdot y=0\}$.
Clause (d) makes three assertions on this set:
\begin{itemize}
\item[(su1)] $\mathrm{Ls}_{\hat{n}}$ belongs to $B_{\mathrm{Law}}$, the $\sigma$-algebra of (n7);
\item[(su2)] every law in $\mathrm{Ls}_{\hat{n}}$ is invariant under $r_{\hat{n},\mathrm{po}}$ and reflection positive at $\{\hat{n}\cdot y=\mathrm{po}\}$ for every real $\mathrm{po}$, so that $\mathrm{Ls}_{\hat{n}}$ is the set of all laws that are invariant under $\mathbb{R}^3$, invariant under $r_{\hat{n},\mathrm{po}}$ for every real $\mathrm{po}$, and reflection positive at every plane $\{\hat{n}\cdot y=\mathrm{po}\}$;
\item[(su3)] under the hypotheses of Proposition~45.7, $\nu(\mathrm{Ls}_{\hat{n}})=1$.
\end{itemize}
The assertions (su1) and (su2) concern the set $\mathrm{Ls}_{\hat{n}}$ alone; no law $\mathrm{lw}$ is involved in them.
For $\hat{n}=e_j$ the set $\mathrm{Ls}_{\hat{n}}$ is the set $\mathrm{Ls}_b$ of Step~8 of Section~1.

Choose $\mathrm{gr}\in O(3)$ with $\mathrm{gr}\,e_1=\hat{n}$, and let $\mathrm{Tg}$ be the map of (tr6).
Let $\mathrm{Ls}'_b$ be the set of the laws that are invariant under $\mathbb{R}^3$ and under $\theta_1$ and reflection positive at $\{y_1=0\}$.
It is the intersection of three sets that belong to $B_{\mathrm{Law}}$ (qc7).

\emph{Proof of (su1).}
Let $\xi'$ be a law and $\xi:=\mathrm{Tg}(\xi')=\mathrm{gr}_*\xi'$.
By (tr4), $\xi$ is invariant under $\mathbb{R}^3$ if and only if $\xi'$ is.
By (tr2) with $\mathrm{Gi}=\theta_1$ and the identity $\mathrm{gr}\,\theta_1\,\mathrm{gr}^{-1}=r_{\hat{n},0}$ of (tr3), $\xi$ is invariant under $r_{\hat{n},0}$ if and only if $\xi'$ is invariant under $\theta_1$.
By (tr5) with $\mathrm{po}=0$, $\xi$ is reflection positive at $\{\hat{n}\cdot y=0\}$ if and only if $\xi'$ is reflection positive at $\{y_1=0\}$.
Hence $\xi\in\mathrm{Ls}_{\hat{n}}$ if and only if $\xi'\in\mathrm{Ls}'_b$.
Since $\mathrm{Tg}$ is a bijection of $\mathrm{Law}$ onto itself (tr6), it follows that
\[
\mathrm{Ls}_{\hat{n}}=\mathrm{Tg}(\mathrm{Ls}'_b).
\]
This is the inverse image of $\mathrm{Ls}'_b$ under the map $\mathrm{Tg}^{-1}$, which is measurable (tr6); so $\mathrm{Ls}_{\hat{n}}\in B_{\mathrm{Law}}$.

\emph{Proof of (su2).}
Let $\xi\in\mathrm{Ls}_{\hat{n}}$; by the proof of (su1), $\xi=\mathrm{gr}_*\xi'$ with $\xi'\in\mathrm{Ls}'_b$.
The law $\xi'$ is invariant under $\theta_1$ and under $\mathbb{R}e_1$, and reflection positive at $\{y_1=0\}$.
By the implication (1)$\Rightarrow$(5) of (qc6), it is reflection positive at every plane $\{y_1=\mathrm{po}\}$, $\mathrm{po}\in\mathbb{R}$.
By (tr5), $\xi$ is reflection positive at every plane $\{\hat{n}\cdot y=\mathrm{po}\}$.
For real $\mathrm{po}$, the last identity of (tr3) and (qc1) give
\[
(r_{\hat{n},\mathrm{po}})_*\xi=(\tau_{\mathrm{po}\cdot\hat{n}})_*(r_{\hat{n},0})_*(\tau_{-\mathrm{po}\cdot\hat{n}})_*\xi=\xi.
\]
Conversely, a law that is invariant under $\mathbb{R}^3$, invariant under $r_{\hat{n},\mathrm{po}}$ for every real $\mathrm{po}$ and reflection positive at every plane $\{\hat{n}\cdot y=\mathrm{po}\}$ has, for $\mathrm{po}=0$, the three properties that define $\mathrm{Ls}_{\hat{n}}$.

\emph{Proof of (su3).}
By Step~4, every $\xi$ in the set $\mathrm{Ls}$ of Step~3 is invariant under $\mathbb{R}^3$ and under $r_{\hat{n},0}$, and reflection positive at $\{\hat{n}\cdot y=0\}$.
Thus $\mathrm{Ls}\subset\mathrm{Ls}_{\hat{n}}$.
Since $\nu(\mathrm{Ls})=1$ (Step~3) and $\mathrm{Ls}_{\hat{n}}\in B_{\mathrm{Law}}$ by (su1), one concludes that $\nu(\mathrm{Ls}_{\hat{n}})=1$.
\end{proof}

\subsection*{4. Conditions on the constants}

Propositions~45.5, 45.6 and 45.7 and their proofs are independent of the constants of Section~I.8 and of Proposition~14.2, and the three Propositions hold for all values of these constants.
The numbers $\mathrm{cf}_{n,m}$, $\mathrm{cf}_n$, $\mathrm{pu}_m$, $\mathrm{cm}_m$ of Section~2 are auxiliary quantities of the proofs.
They depend only on $n$, $m$ and on the two functions $\mathrm{pun}$, $\mathrm{cof}$, which are given by explicit formulas in 2.2.

\section*{Chapter 45-1. Single laws on $\mathcal{S}'$ with (L1)--(L4), and soft laws: proofs of Propositions 45.1--45.4}

\subsection*{0. Conventions of this chapter}

\paragraph{Notation.}
Laws on $\mathcal{S}'$ are described in the notation (n1)--(n8) of Part~II, which is set up there for the Propositions 45.m and for Theorem~44.
The two-form $d^{*}\mathrm{gs}$, the function $\Phi_{\mathrm{lw}}$, the basis $(h_\kappa)_{\kappa}$ with the numbers $\lambda_\kappa$, the numbers $\mathrm{ZH}_s$ and $\sigma_n$, and the properties (L1)--(L4) of a law are defined at the head of the Propositions 45.m in Part~II.
The groups $O_{\mathrm{cub}}$ and $\{\pm1\}^3$ are those of Section~I.1.2.

\paragraph{Complex test functions.}
Let $\mathrm{gs}=\mathrm{gs}'+i\,\mathrm{gs}''$ be a complex element of $\mathcal{S}(\mathbb{R}^3)$, with real $\mathrm{gs}'$ and $\mathrm{gs}''$.
As at the head of the Propositions 45.m, one puts $d^{*}\mathrm{gs}:=d^{*}\mathrm{gs}'+i\,d^{*}\mathrm{gs}''$.

\paragraph{Letters.}
In this chapter $\mathrm{bz}$ and $\mathrm{gv}$ are multi-indices, used beside the multi-index $\alpha$ of (n5); the notation $\partial^\alpha$, $|\alpha|$, $y^{\mathrm{bz}}$, $\binom{\mathrm{gv}}{\mathrm{gv}'}$ for multi-indices is defined in B.4.
The letters $\mathrm{dn}$, $\mathrm{dn}'$ denote integers used as bound variables; $\epsilon$ denotes a positive number; $\mathrm{rv}$, $\mathrm{rv}'$ denote real and $\mathrm{zv}$ denotes a complex bound variable.
The letters $\mathrm{va}$, $\mathrm{vb}$ (vectors of $\mathbb{R}^3$) and $\mathrm{bk}$ (a block of a partition) have the meanings of (n8) and (n5).
In B.3, in B.8 and in (F.1), $(\mathrm{Prs},\mathrm{Pb})$ is a probability space with the expectation $\mathrm{Eb}$, $\mathrm{Ts}$ is a real vector space with a Hilbert seminorm, and $\mathrm{ts}$ is an element of $\mathrm{Ts}$.
A Hilbert seminorm on a real vector space $\mathrm{Ts}$ is a map $\mathrm{ts}\mapsto\langle\mathrm{ts},\mathrm{ts}\rangle_{\mathrm{Ts}}^{\frac12}$, where $\langle\cdot,\cdot\rangle_{\mathrm{Ts}}$ is a positive semidefinite symmetric bilinear form on $\mathrm{Ts}$; if this form is positive definite, the map is a Hilbert norm.
In Part~E, Step~7, $(\mathrm{Mq},\mathrm{mq})$ is a measure space.
The letters $\mathrm{fs}$, $\mathrm{fs}'$ denote functions on $\mathbb{R}^3$ or on $\mathbb{R}^{3n}$, and $\mathrm{cf}$ denotes a finite number.
Without an index, $\Delta$ is the Laplace operator $\sum_\ell\partial_\ell^2$ of $\mathbb{R}^3$ or of $\mathbb{R}^{3n}$.

The letters $a$, $\ell$, $\ell'$, $p$, $q$ serve as indices.
In addition, $a$ denotes a real number in (He3) and in the conclusion (3) of \cite[Proposition~7.10]{QEDI} (see B.3), and a vector of $\mathbb{R}^3$ in F\,(c), as in (n2) and (n3); $p$ denotes a real exponent in the same conclusion (3), and a momentum.
The letters $b$, $b'$, $b_0$ denote vectors of $\mathbb{R}^3$; in the proof of Lemma~B.4, $b$ is an index.
The $c_p$ are complex coefficients; $m$, $m'$, $n$, $n'$ are integers $\ge0$; $r$ and $w$ are elements of $O_{\mathrm{cub}}$; $s$ is a real number.
The letters $t$, $t'$, $t_0$ denote real numbers, with two exceptions: as the exponent in $(\mathrm{Osc}^t)^{\otimes n}$, in B.2, B.4, B.9 and in Lemma~D.1, $t$ is an integer $\ge0$, and in the proof of Lemma~B.1 it is a vector of $\mathbb{R}^m$.
The letter $z$ denotes a complex number or a vector of $\mathbb{C}^m$ or of $\mathbb{C}^n$, and $O$ denotes an open subset of $\mathbb{R}^3$.
With an argument, $T(t)$ denotes the family of operators of \cite[Proposition~8.19\,(ii)]{QEDI}; it appears only in E\,(4).

\paragraph{Labels.}
The labels (Nm0)--(Nm8), (He1)--(He3), (Eu1)--(Eu7) and the labels (B.1a), \dots, (F.2) of displays refer to this chapter.
The display (3.1) is the hypothesis of the proposition recalled in B.3, and its conclusions (1)--(7) are cited as \cite[Proposition~7.10(1)]{QEDI}, \dots, \cite[Proposition~7.10(7)]{QEDI}.
From Part~C on, such a citation means the conclusion in question for the probability space $(\mathcal{S}',\Sigma,\mathrm{lw})$, the space $L^2(\mathbb{R}^3;\Lambda^2)$ with the norm $\|\cdot\|_2$, and the map $X$ of Proposition~45.1\,(1), as fixed at the beginning of the proof in Part~C.
The letters B, C, D, E, F, G, without brackets, denote the parts of this chapter; thus E\,(4) is the proof of the clause (4) in Part~E.

\paragraph{Constants.}
Every number in the four Propositions 45.1--45.4 is absolute: it depends on no lattice, level, side, torus, cutoff, mass or charge, and on none of the constants of Section~I.8.

\subsection*{B. Quoted statements, preliminaries}

\subsubsection*{B.1. Real and complex spaces}

Let $\Lambda^2$ be the space of the real alternating two-tensors on $\mathbb{R}^3$, with the orthonormal basis $e_{\mathrm{il}}:=\mathrm{d}y_i\wedge\mathrm{d}y_\ell$, $\mathrm{il}\in\{12,13,23\}$; its elements are written $f=\sum_{\mathrm{il}}f_{\mathrm{il}}\cdot e_{\mathrm{il}}$.
The Euclidean norm on $\mathbb{R}^3$ and on $\mathbb{R}^{3n}$ is denoted by $|\cdot|_2$.
By (n1) the space $\mathcal{S}$ is real; $\mathcal{S}+i\mathcal{S}$ is its complexification.
The topology of $\mathcal{S}$ is that of the norms $q_m$, $m\ge0$, that is, the topology induced by $\mathrm{Sw}_1=\mathcal{S}+i\mathcal{S}$.
By (n8), $L^2(\mathbb{R}^3;\Lambda^2)$ is a real Hilbert space, with
\[
\langle f,f'\rangle:=\sum_{\mathrm{il}}\int\mathrm{d}y\,f_{\mathrm{il}}(y)\cdot f'_{\mathrm{il}}(y),\qquad\|f\|_2:=\langle f,f\rangle^{\frac12};
\]
on $L^2(\mathbb{R}^3;\Lambda^2)\otimes\mathbb{C}$ one puts $\|f+if'\|_2^2:=\|f\|_2^2+\|f'\|_2^2$.
An element $\omega$ of $\mathcal{S}'$ is a continuous real-linear functional on $\mathcal{S}$; for $f,f'\in\mathcal{S}$ one puts
\[
\omega(f+if'):=\omega(f)+i\,\omega(f'),\qquad X(f+if'):=X(f)+i\,X(f').
\]
The space $L^2(\mathrm{lw})$ is complex, and $\|\mathrm{Ga}\|_{L^2(\mathrm{lw})}:=\langle\mathrm{Ga},\mathrm{Ga}\rangle_{\mathrm{lw}}^{\frac12}$.

For $n\ge1$ the space $\mathrm{Sw}_n$ consists of the functions
\[
\Psi=\sum_{\mathrm{il}_1,\dots,\mathrm{il}_n}\Psi_{\mathrm{il}_1\dots\mathrm{il}_n}\cdot e_{\mathrm{il}_1}\otimes\cdots\otimes e_{\mathrm{il}_n}
\]
on $\mathbb{R}^{3n}$ whose $3^n$ components $\Psi_{\mathrm{il}_1\dots\mathrm{il}_n}$ belong to the complex space $\mathcal{S}(\mathbb{R}^{3n})$.
Its topology is that of the norms $q_m$, $m\ge0$; and $\mathrm{Sw}_1=\mathcal{S}+i\mathcal{S}$.
One defines
\begin{gather*}
\langle\Phi,\Psi\rangle_{L^2}:=\sum_{\mathrm{il}_1,\dots,\mathrm{il}_n}\int\mathrm{d}y\,\overline{\Phi_{\mathrm{il}_1\dots\mathrm{il}_n}(y)}\cdot\Psi_{\mathrm{il}_1\dots\mathrm{il}_n}(y),\qquad\|\Psi\|_{L^2}:=\langle\Psi,\Psi\rangle_{L^2}^{\frac12},\\
(f_1\otimes\cdots\otimes f_n)(y_1,\dots,y_n):=f_1(y_1)\otimes\cdots\otimes f_n(y_n),\\
(\Psi\otimes\Phi)(y_1,\dots,y_{m+n}):=\Psi(y_1,\dots,y_m)\otimes\Phi(y_{m+1},\dots,y_{m+n}),
\end{gather*}
the last one for $\Psi\in\mathrm{Sw}_m$ and $\Phi\in\mathrm{Sw}_n$.
For scalar functions on $\mathbb{R}^3$ the norms $q_m$ are given by the same formula.
For $n=0$ the conventions of (n8) apply; among them are, for $\Psi_0\in\mathrm{Sw}_0=\mathbb{C}$, $\Psi\in\mathrm{Sw}_n$ and every Euclidean map $\mathrm{Gi}$,
\[
\Psi\otimes\Psi_0:=\Psi_0\cdot\Psi,\qquad(\mathrm{Gi}\cdot)^{\otimes0}\Psi_0:=\Psi_0.
\]

\paragraph{The exterior square.}
Let $e_1,e_2,e_3$ be the unit vectors of $\mathbb{R}^3$.
For $\mathrm{va},\mathrm{vb}\in\mathbb{R}^3$ put
\[
\mathrm{va}\wedge\mathrm{vb}:=\sum_{i<\ell}\left(\mathrm{va}_i\cdot\mathrm{vb}_\ell-\mathrm{va}_\ell\cdot\mathrm{vb}_i\right)\cdot e_{\mathrm{il}}.
\]
The map $(\mathrm{va},\mathrm{vb})\mapsto\mathrm{va}\wedge\mathrm{vb}$ from $\mathbb{R}^3\times\mathbb{R}^3$ to $\Lambda^2$ is bilinear and alternating, that is, $\mathrm{vb}\wedge\mathrm{va}=-\mathrm{va}\wedge\mathrm{vb}$ and $\mathrm{va}\wedge\mathrm{va}=0$; moreover $e_i\wedge e_\ell=e_{\mathrm{il}}$ for $i<\ell$.
For a real $3\times3$ matrix $\mathrm{Rl}$ let $\wedge^2\mathrm{Rl}$ be the linear map of $\Lambda^2$ into itself with $(\wedge^2\mathrm{Rl})\,e_{\mathrm{il}}:=\mathrm{Rl}\,e_i\wedge\mathrm{Rl}\,e_\ell$ for $i<\ell$; by (B.1a) below, it is the map $\wedge^2\mathrm{Rl}$ of (n8).
Let $\langle\cdot,\cdot\rangle$ be the scalar product of $\Lambda^2$ for which the $e_{\mathrm{il}}$ are orthonormal.
For all $\mathrm{va},\mathrm{vb},\mathrm{va}',\mathrm{vb}'\in\mathbb{R}^3$ and all real $3\times3$ matrices $\mathrm{Rl}$, $\mathrm{Rl}'$ one has
\begin{align}
(\wedge^2\mathrm{Rl})(\mathrm{va}\wedge\mathrm{vb})&=\mathrm{Rl}\,\mathrm{va}\wedge\mathrm{Rl}\,\mathrm{vb},\tag{B.1a}\\
\wedge^2(\mathrm{Rl}\,\mathrm{Rl}')&=(\wedge^2\mathrm{Rl})(\wedge^2\mathrm{Rl}'),\qquad\wedge^2 1=1,\tag{B.1b}\\
\langle\mathrm{va}\wedge\mathrm{vb},\mathrm{va}'\wedge\mathrm{vb}'\rangle&=(\mathrm{va}\cdot\mathrm{va}')\cdot(\mathrm{vb}\cdot\mathrm{vb}')-(\mathrm{va}\cdot\mathrm{vb}')\cdot(\mathrm{vb}\cdot\mathrm{va}'),\tag{B.1c}
\end{align}
and, if $\mathrm{Rl}$ lies in $O(3)$,
\begin{equation}
\left\langle(\wedge^2\mathrm{Rl})\,e_{\mathrm{il}},(\wedge^2\mathrm{Rl})\,e_{\mathrm{il}'}\right\rangle=\langle e_{\mathrm{il}},e_{\mathrm{il}'}\rangle\qquad\left(\mathrm{il},\mathrm{il}'\in\{12,13,23\}\right).\tag{B.1d}
\end{equation}
By (B.1d), for $\mathrm{Rl}\in O(3)$ the matrix of $\wedge^2\mathrm{Rl}$ in the basis $(e_{\mathrm{il}})$ is orthogonal.
Each of its columns is a unit vector, so that its entries have modulus at most $1$; and for $f=\sum_{\mathrm{il}}f_{\mathrm{il}}\cdot e_{\mathrm{il}}\in\Lambda^2$ the squares of the components of $(\wedge^2\mathrm{Rl})f$ add up to $\sum_{\mathrm{il}}f_{\mathrm{il}}^2$.

\begin{proof}
(B.1a).
Since $\mathrm{Rl}\,\mathrm{va}=\sum_i\mathrm{va}_i\cdot\mathrm{Rl}\,e_i$, bilinearity gives
\[
\mathrm{Rl}\,\mathrm{va}\wedge\mathrm{Rl}\,\mathrm{vb}=\sum_{i,\ell=1}^3\mathrm{va}_i\cdot\mathrm{vb}_\ell\cdot\left(\mathrm{Rl}\,e_i\wedge\mathrm{Rl}\,e_\ell\right).
\]
The terms with $i=\ell$ vanish, and $\mathrm{Rl}\,e_\ell\wedge\mathrm{Rl}\,e_i=-\mathrm{Rl}\,e_i\wedge\mathrm{Rl}\,e_\ell$.
Hence the sum equals
\[
\sum_{i<\ell}\left(\mathrm{va}_i\cdot\mathrm{vb}_\ell-\mathrm{va}_\ell\cdot\mathrm{vb}_i\right)\cdot\left(\mathrm{Rl}\,e_i\wedge\mathrm{Rl}\,e_\ell\right),
\]
which is $(\wedge^2\mathrm{Rl})(\mathrm{va}\wedge\mathrm{vb})$ by the definitions of $\mathrm{va}\wedge\mathrm{vb}$ and of $\wedge^2\mathrm{Rl}$.

(B.1b).
By (B.1a),
\[
(\wedge^2\mathrm{Rl})(\wedge^2\mathrm{Rl}')\,e_{\mathrm{il}}=(\wedge^2\mathrm{Rl})\left(\mathrm{Rl}'e_i\wedge\mathrm{Rl}'e_\ell\right)=\mathrm{Rl}\,\mathrm{Rl}'e_i\wedge\mathrm{Rl}\,\mathrm{Rl}'e_\ell=\left(\wedge^2(\mathrm{Rl}\,\mathrm{Rl}')\right)e_{\mathrm{il}},
\]
and $(\wedge^2 1)\,e_{\mathrm{il}}=e_i\wedge e_\ell=e_{\mathrm{il}}$.

(B.1c).
One has
\[
(\mathrm{va}\cdot\mathrm{va}')\cdot(\mathrm{vb}\cdot\mathrm{vb}')-(\mathrm{va}\cdot\mathrm{vb}')\cdot(\mathrm{vb}\cdot\mathrm{va}')=\sum_{i,\ell=1}^3\mathrm{va}_i\cdot\mathrm{vb}_\ell\cdot\left(\mathrm{va}'_i\cdot\mathrm{vb}'_\ell-\mathrm{va}'_\ell\cdot\mathrm{vb}'_i\right).
\]
The bracket vanishes for $i=\ell$, and it changes its sign when $i$ and $\ell$ are exchanged.
So the sum equals
\[
\sum_{i<\ell}\left(\mathrm{va}_i\cdot\mathrm{vb}_\ell-\mathrm{va}_\ell\cdot\mathrm{vb}_i\right)\cdot\left(\mathrm{va}'_i\cdot\mathrm{vb}'_\ell-\mathrm{va}'_\ell\cdot\mathrm{vb}'_i\right),
\]
and this is the scalar product of $\mathrm{va}\wedge\mathrm{vb}$ and $\mathrm{va}'\wedge\mathrm{vb}'$, because the $e_{\mathrm{il}}$ are orthonormal.

(B.1d).
Write $\mathrm{il}=(i,\ell)$ and $\mathrm{il}'=(i',\ell')$.
Apply (B.1c) to the four vectors $\mathrm{Rl}\,e_i$, $\mathrm{Rl}\,e_\ell$, $\mathrm{Rl}\,e_{i'}$, $\mathrm{Rl}\,e_{\ell'}$, use $\mathrm{Rl}\,y\cdot\mathrm{Rl}\,y'=y\cdot y'$ for $\mathrm{Rl}\in O(3)$ and $y,y'\in\mathbb{R}^3$, and apply (B.1c) once more, to $e_i$, $e_\ell$, $e_{i'}$, $e_{\ell'}$:
\[
\left\langle\mathrm{Rl}\,e_i\wedge\mathrm{Rl}\,e_\ell,\mathrm{Rl}\,e_{i'}\wedge\mathrm{Rl}\,e_{\ell'}\right\rangle=(e_i\cdot e_{i'})\cdot(e_\ell\cdot e_{\ell'})-(e_i\cdot e_{\ell'})\cdot(e_\ell\cdot e_{i'})=\left\langle e_i\wedge e_\ell,e_{i'}\wedge e_{\ell'}\right\rangle.
\]
\end{proof}

\subsubsection*{B.2. The Hermite scale}

This section recalls \cite[Section~7.1.5]{QEDI}, together with the display of that section in which the sums $\sum_\kappa\lambda_\kappa^{-a}$ are bounded.
Let $\mathrm{Nat}:=\{0,1,2,\dots\}$.
The Hermite functions $\mathrm{hf}_{\mathrm{kn}}$, $\mathrm{kn}\in\mathrm{Nat}^3$, on $\mathbb{R}^3$ are real elements of $\mathcal{S}(\mathbb{R}^3)$ and form an orthonormal basis of $L^2(\mathbb{R}^3)$; with $\mathrm{Osc}:=1+|y|_2^2-\Delta$ they satisfy
\[
\mathrm{Osc}\,\mathrm{hf}_{\mathrm{kn}}=\left(2|\mathrm{kn}|_1+4\right)\cdot\mathrm{hf}_{\mathrm{kn}}.
\]
Put $\mathrm{Kh}:=\mathrm{Nat}^3\times\{12,13,23\}$, and for $\kappa=(\mathrm{kn},\mathrm{il})\in\mathrm{Kh}$ put
\[
h_\kappa:=\mathrm{hf}_{\mathrm{kn}}\cdot e_{\mathrm{il}},\qquad\lambda_\kappa:=2|\mathrm{kn}|_1+4\ge4.
\]
The family $(h_\kappa)_{\kappa\in\mathrm{Kh}}$ is an orthonormal basis of $L^2(\mathbb{R}^3;\Lambda^2)$.
For $\mathrm{kk}=(\kappa_1,\dots,\kappa_n)\in\mathrm{Kh}^n$ let $h_{\mathrm{kk}}:=h_{\kappa_1}\otimes\cdots\otimes h_{\kappa_n}$.
The functions $h_{\mathrm{kk}}$ are real, and $(h_{\mathrm{kk}})_{\mathrm{kk}\in\mathrm{Kh}^n}$ is an orthonormal basis of the complex Hilbert space of the square integrable functions $\mathbb{R}^{3n}\to(\Lambda^2)^{\otimes n}\otimes\mathbb{C}$, with the scalar product $\langle\cdot,\cdot\rangle_{L^2}$; this is the theorem on products of orthonormal bases \cite[Section~II.4, Proposition~2, with Theorem~II.10]{ReedSimon1980}.
Let $\mathrm{Fin}_n$ be the complex linear span of the $h_{\mathrm{kk}}$, $\mathrm{kk}\in\mathrm{Kh}^n$.

For $\Psi\in\mathrm{Sw}_n$ and $\mathrm{kk}\in\mathrm{Kh}^n$ put $c_{\mathrm{kk}}(\Psi):=\langle h_{\mathrm{kk}},\Psi\rangle_{L^2}$, as in Proposition~45.2\,(2).
For $a=1,\dots,n$ let
\[
\mathrm{Osc}_a:=1+|y_a|_2^2-\sum_{\ell=1}^3\partial_{(y_a)_\ell}^2
\]
act on the $a$-th argument of a function on $\mathbb{R}^{3n}$, component by component.
The operators $\mathrm{Osc}_a$ commute with each other, since they act on different arguments.
For integers $t\ge0$ put $(\mathrm{Osc}^t)^{\otimes n}:=\prod_{a=1}^n\mathrm{Osc}_a^t$.

For real $s$ and $\omega\in\mathcal{S}'$, the quantities $\|\omega\|_s$ and $H_s$ are those of (n8); the number $\|\omega\|_s^2$ lies in $[0,\infty]$.
The function $\omega\mapsto\|\omega\|_s^2$ is $\Sigma$-measurable, being a countable sum of squares of the functions $X(h_\kappa)$; hence $H_s$ belongs to $\Sigma$.
\begin{itemize}
\item[(He1)] $\mathcal{S}'=\bigcup_{s\ge0}H_{-s}$.
\item[(He2)] For every $n\ge1$, $\mathrm{Fin}_n$ is dense in $\mathrm{Sw}_n$.
The real linear span of the $h_\kappa$ is dense in $\mathcal{S}$.
\item[(He3)] $\mathrm{ZH}_a<\infty$ if and only if $a>3$, and
\[
\mathrm{ZH}_4=\frac{3}{32}\sum_{\mathrm{dn}\ge0}\frac{\mathrm{dn}+1}{(\mathrm{dn}+2)^3}\le\frac{3}{32}\left(\frac{\pi^2}{6}-1\right)<\frac{1}{16}.
\]
\end{itemize}

\begin{proof}[Proof of (He3)]
Assertion (He3) is the display of \cite[Section~7.1.5]{QEDI} that concerns these sums.
For $\mathrm{dn}\ge0$ there are $\frac{(\mathrm{dn}+1)(\mathrm{dn}+2)}{2}$ multi-indices $\mathrm{kn}\in\mathrm{Nat}^3$ with $|\mathrm{kn}|_1=\mathrm{dn}$; therefore
\[
\mathrm{ZH}_a=3\sum_{\mathrm{dn}\ge0}\frac{(\mathrm{dn}+1)(\mathrm{dn}+2)}{2}\cdot(2\,\mathrm{dn}+4)^{-a}.
\]
Since $\mathrm{dn}+1\le\mathrm{dn}+2\le2(\mathrm{dn}+1)$, every term of this series lies between two positive multiples of $(\mathrm{dn}+1)^{2-a}$; so the series converges if and only if $a-2>1$.
For $a=4$ the term equals
\[
\frac{3}{32}\cdot\frac{\mathrm{dn}+1}{(\mathrm{dn}+2)^3}\le\frac{3}{32}\cdot(\mathrm{dn}+2)^{-2}.
\]
Moreover $\sum_{\mathrm{dn}\ge0}(\mathrm{dn}+2)^{-2}=\frac{\pi^2}{6}-1<0.65$, because $\pi<3.1416$ and hence $\pi^2<9.87$; and $\frac{3}{32}\cdot0.65=0.0609375<0.0625=\frac{1}{16}$.
\end{proof}

The assertions (He1) and (He2) are proved in B.9.
In \cite[Section~7.1.5]{QEDI}, (He1) and the first sentence of (He2) are obtained from the $N$-representation theorems for $\mathcal{S}$ and for $\mathcal{S}'$; these two theorems are \cite[Theorems~V.13 and~V.14]{ReedSimon1980}.

\subsubsection*{B.3. Three statements of \cite[Section~7]{QEDI}}

The following three statements are those of \cite[Proposition~7.10, Lemma~7.14 and Lemma~7.16]{QEDI}, written in the notation of this chapter.

\begin{named}{Proposition}[{\cite[Proposition~7.10]{QEDI}}]
Let $(\mathrm{Prs},\mathrm{Pb})$ be a probability space with the expectation $\mathrm{Eb}$, let $\mathrm{Ts}$ be a real vector space with a Hilbert seminorm $|\cdot|$, that is, the seminorm of a positive semidefinite symmetric bilinear form on $\mathrm{Ts}$, and let $X$ be a linear map from $\mathrm{Ts}$ to the real random variables such that
\begin{equation}
\mathrm{Eb}\exp X(\mathrm{ts})\le\exp\left(\frac{|\mathrm{ts}|^2}{2}\right)\qquad\left(\mathrm{ts}\in\mathrm{Ts}\right).\tag{3.1}
\end{equation}
Extend $X$ and $|\cdot|$ to $\mathrm{Ts}+i\mathrm{Ts}$, with $|\mathrm{ts}+i\mathrm{ts}'|^2=|\mathrm{ts}|^2+|\mathrm{ts}'|^2$.
Then the following hold.
\begin{itemize}
\item[(1)] $\mathrm{Eb}\exp\left(t\cdot|X(\mathrm{ts})|\right)\le2\exp\left(\frac{t^2\cdot|\mathrm{ts}|^2}{2}\right)$.
For $f_1,\dots,f_m\in\mathrm{Ts}+i\mathrm{Ts}$ the function $z\mapsto\mathrm{Eb}\exp\sum_iz_i\cdot X(f_i)$ is entire on $\mathbb{C}^m$.
One has $\left|\mathrm{Eb}\exp X(f)\right|\le\exp\left(\frac{|\operatorname{Re}f|^2}{2}\right)$; in particular $\left|\mathrm{Eb}\exp\left(z\cdot X(\mathrm{ts})\right)\right|\le\exp\left(\frac{(\operatorname{Re}z)^2\cdot|\mathrm{ts}|^2}{2}\right)$.
\item[(2)] $\mathrm{Eb}X(\mathrm{ts})=0$, $\mathrm{Eb}X(\mathrm{ts})^2\le|\mathrm{ts}|^2$, and $1\le\mathrm{Eb}\exp\left(t\cdot X(\mathrm{ts})\right)\le\exp\left(\frac{t^2\cdot|\mathrm{ts}|^2}{2}\right)$ for real $t$.
\item[(3)] For $t>0$ one has $\mathrm{Pb}\left(|X(\mathrm{ts})|\ge t\right)\le2\exp\left(-\frac{t^2}{2|\mathrm{ts}|^2}\right)$, where the right side is read as $0$ if $|\mathrm{ts}|=0$;
$\mathrm{Eb}|X(\mathrm{ts})|^p\le2\cdot2^{\frac{p}{2}}\cdot\Gamma\left(\frac{p}{2}+1\right)\cdot|\mathrm{ts}|^p$ for $p>0$;
$\mathrm{Eb}\exp\left(a\cdot X(\mathrm{ts})^2\right)\le\left(1-2a\cdot|\mathrm{ts}|^2\right)^{-\frac12}$ for $0\le2a\cdot|\mathrm{ts}|^2<1$.
\item[(4)] $\mathrm{Eb}\prod_{i=1}^n|X(\mathrm{ts}_i)|\le\sigma_n\cdot\prod_i|\mathrm{ts}_i|$ for $\mathrm{ts}_i\in\mathrm{Ts}$, and $\mathrm{Eb}\prod_{i=1}^n|X(f_i)|\le2^{\frac{n}{2}}\cdot\sigma_n\cdot\prod_i|f_i|$ for $f_i\in\mathrm{Ts}+i\mathrm{Ts}$.
\item[(5)] For $\mathrm{ts}\in\mathrm{Ts}$ and $z\in\mathbb{C}$,
\[
\left|\mathrm{Eb}\exp\left(z\cdot X(\mathrm{ts})\right)-1\right|\le2\left(\exp\left(\frac{|z|^2\cdot|\mathrm{ts}|^2}{2}\right)-1\right).
\]
Hence $\left|\mathrm{Eb}\exp\left(z\cdot X(\mathrm{ts})\right)-1\right|\le\frac12$ if $|z|\cdot|\mathrm{ts}|\le\frac23$, and $\mathrm{Eb}\exp\left(z\cdot X(\mathrm{ts})\right)\ne0$ if $|z|\cdot|\mathrm{ts}|<\left(2\log\frac32\right)^{\frac12}=0.9005\ldots$
\item[(6)] For $f\in\mathrm{Ts}+i\mathrm{Ts}$ one has $\left|\mathrm{Eb}\exp X(f)-1\right|\le\exp\left(2|f|^2\right)-1$, and the right side is at most $\frac12$ if $|f|\le\frac25$.
\item[(7)] For $f_1,\dots,f_n\in\mathrm{Ts}+i\mathrm{Ts}$ the function $z\mapsto\log\mathrm{Eb}\exp\sum_iz_i\cdot X(f_i)$, formed with the principal branch of the logarithm, is analytic and of modulus at most $\log2$ on the set $\left\{z:\sum_i|z_i|\cdot|f_i|\le\frac25\right\}$.
Its derivative $\partial_{z_1}\cdots\partial_{z_n}$ at $z=0$, which this proposition denotes by $\kappa_n(f_1,\dots,f_n)$, satisfies
\[
\left|\kappa_n(f_1,\dots,f_n)\right|\le(\log2)\cdot\left(\frac{5n}{2}\right)^n\cdot\prod_i|f_i|\le(\log2)\cdot\left(\frac{5\mathrm{e}}{2}\right)^n\cdot n!\cdot\prod_i|f_i|.
\]
For $\mathrm{ts}\in\mathrm{Ts}$ one has $\left|\kappa_n(\mathrm{ts},\dots,\mathrm{ts})\right|\le(\log2)\cdot n!\cdot\left(\frac32\right)^n\cdot|\mathrm{ts}|^n$, $\kappa_1=0$, and $0\le\kappa_2(\mathrm{ts},\mathrm{ts})\le|\mathrm{ts}|^2$.
\end{itemize}
\end{named}

\begin{named}{Lemma}[{\cite[Lemma~7.14]{QEDI}}]
For every probability measure $\mathrm{lw}$ on $(\mathcal{S}',\Sigma)$ and every real linear subspace $\mathrm{Dq}$ of $\mathcal{S}$, the complex linear span of the functions $\exp\left(i\,X(f)\right)$, $f\in\mathrm{Dq}$, is dense in $L^2\left(\mathcal{S}',\sigma\left(X(f):f\in\mathrm{Dq}\right),\mathrm{lw}\right)$.
\end{named}

\begin{named}{Lemma}[{\cite[Lemma~7.16]{QEDI}}]
For $\Psi\in\mathrm{Sw}_n$,
\[
\sum_{\mathrm{kk}}\left|c_{\mathrm{kk}}(\Psi)\right|\le\mathrm{ZH}_4^{\frac{n}{2}}\cdot\left\|(\mathrm{Osc}^2)^{\otimes n}\Psi\right\|_{L^2}\le4^{-n}\cdot\left(34\pi\sqrt3\right)^n\cdot q_{6n}(\Psi).
\]
\end{named}

No lattice and no law of a lattice model enters these three statements.
The space $L^2\left(\mathcal{S}',\sigma\left(X(f):f\in\mathrm{Dq}\right),\mathrm{lw}\right)$ is identified with the closed subspace of $L^2(\mathrm{lw})$ formed by the classes that have a representative measurable for this $\sigma$-algebra.
For $\mathrm{Dq}=C_c^\infty(O;\Lambda^2)$, with real test functions, it is the space $L^2(\Sigma_O,\mathrm{lw})$ of (n1).
The coefficients $c_{\mathrm{kk}}$ and the operator $(\mathrm{Osc}^2)^{\otimes n}$ are those of B.2.

\subsubsection*{B.4. Inequalities for the norms $q_m$}

Let $n\ge1$, $m\ge0$ and $\Psi\in\mathrm{Sw}_n$.
Let $\alpha$, $\mathrm{bz}$, $\mathrm{gv}$ be multi-indices in $\mathrm{Nat}^{3n}$, as in (n5).
The entries of a multi-index $\alpha=(\alpha_1,\dots,\alpha_{3n})$ correspond to the $3n$ coordinates of $y=(y_1,\dots,y_n)\in\mathbb{R}^{3n}$, the entry $\alpha_{3(a-1)+\ell}$ to the coordinate $(y_a)_\ell$, for $a=1,\dots,n$ and $\ell=1,2,3$; one puts
\[
\partial^\alpha:=\prod_{a=1}^n\prod_{\ell=1}^3\partial_{(y_a)_\ell}^{\alpha_{3(a-1)+\ell}},\qquad y^{\mathrm{bz}}:=\prod_{a=1}^n\prod_{\ell=1}^3(y_a)_\ell^{\mathrm{bz}_{3(a-1)+\ell}},\qquad|\alpha|:=\alpha_1+\dots+\alpha_{3n}.
\]
The number $|\alpha|$ is the order of $\alpha$ and of $\partial^\alpha$; it is the quantity written $|\cdot|_1$ in B.2 and B.9, and not the sup norm of Section~I.1.1.
Relations such as $\mathrm{gv}'\le\mathrm{gv}$ and expressions such as $\mathrm{gv}-\mathrm{gv}'+\alpha$ are meant entry by entry, and for $\mathrm{gv}'\le\mathrm{gv}$ one puts
\[
\binom{\mathrm{gv}}{\mathrm{gv}'}:=\prod_{\ell=1}^{3n}\binom{\mathrm{gv}_\ell}{\mathrm{gv}'_\ell}.
\]
Then the following hold.
\begin{itemize}
\item[(Nm0)] For $y,y'\in\mathbb{R}^{3n}$,
\[
1+|y|_2^2\le2\left(1+|y-y'|_2^2\right)\cdot\left(1+|y'|_2^2\right);
\]
and $\prod_{a=1}^n\left(1+|y_a|_2^2\right)\le\left(1+|y|_2^2\right)^n$ for $y=(y_1,\dots,y_n)$.
\item[(Nm1)] $q_m\left(y^{\mathrm{bz}}\partial^\alpha\Psi\right)\le2^m\cdot|\mathrm{bz}|!\cdot q_{m+|\alpha|+|\mathrm{bz}|}(\Psi)$.
\item[(Nm2)] $\|\Psi\|_{L^2}\le\left(\pi\sqrt3\right)^n\cdot q_{2n}(\Psi)$.
\item[(Nm3)] For every smooth function $\mathrm{wd}$ on $\mathbb{R}^{3n}$ that is bounded together with all its derivatives,
\[
q_m(\mathrm{wd}\cdot\Psi)\le2^m\cdot\max_{|\mathrm{gv}|\le m}\sup\left|\partial^{\mathrm{gv}}\mathrm{wd}\right|\cdot q_m(\Psi).
\]
\item[(Nm4)] For every Euclidean map $\mathrm{Gi}:y\mapsto\mathrm{Rl}\,y+b$,
\[
q_m\left((\mathrm{Gi}\cdot)^{\otimes n}\Psi\right)\le3^{n+m}\cdot\left(2\left(1+n|b|_2^2\right)\right)^{\frac{m}{2}}\cdot q_m(\Psi).
\]
\item[(Nm5)] For $0<\epsilon\le1$ and $j=1,2,3$,
\[
q_m\left((\tau_{\epsilon e_j}\cdot)^{\otimes n}\Psi-\Psi\right)\le\epsilon\cdot n\cdot\left(2(1+n)\right)^{\frac{m}{2}}\cdot q_{m+1}(\Psi).
\]
\item[(Nm6)] For $\Phi\in\mathrm{Sw}_{n'}$ the function $\Psi\otimes\Phi$ lies in $\mathrm{Sw}_{n+n'}$, and $q_m(\Psi\otimes\Phi)\le q_m(\Psi)\cdot q_m(\Phi)$.
\item[(Nm7)] For $a=1,\dots,n$ the function $\mathrm{Osc}_a\Psi$ lies in $\mathrm{Sw}_n$, and $q_m(\mathrm{Osc}_a\Psi)\le10\cdot2^m\cdot q_{m+2}(\Psi)$.
For integers $t\ge0$ the function $(\mathrm{Osc}^t)^{\otimes n}\Psi$ lies in $\mathrm{Sw}_n$, and
\[
q_m\left((\mathrm{Osc}^t)^{\otimes n}\Psi\right)\le10^{nt}\cdot2^{nt(m+nt-1)}\cdot q_{m+2nt}(\Psi),
\]
where $nt$ is the product of $n$ and $t$.
\item[(Nm8)] For integers $t\ge0$ and $\mathrm{kk}=(\kappa_1,\dots,\kappa_n)\in\mathrm{Kh}^n$,
\[
c_{\mathrm{kk}}\left((\mathrm{Osc}^t)^{\otimes n}\Psi\right)=\left(\prod_i\lambda_{\kappa_i}\right)^t\cdot c_{\mathrm{kk}}(\Psi),\qquad\sum_{\mathrm{kk}}\left(\prod_i\lambda_{\kappa_i}\right)^{2t}\cdot\left|c_{\mathrm{kk}}(\Psi)\right|^2=\left\|(\mathrm{Osc}^t)^{\otimes n}\Psi\right\|_{L^2}^2<\infty.
\]
\end{itemize}

\begin{proof}
(Nm0).
From $|y|_2^2\le2|y-y'|_2^2+2|y'|_2^2$ one obtains
\[
1+|y|_2^2\le1+2|y-y'|_2^2+2|y'|_2^2\le2\left(1+|y-y'|_2^2\right)\cdot\left(1+|y'|_2^2\right).
\]
Each factor $1+|y_a|_2^2$ is at most $1+|y|_2^2$.

(Nm1).
Let $|\mathrm{gv}|\le m$.
By the rule of Leibniz,
\[
\partial^{\mathrm{gv}}\left(y^{\mathrm{bz}}\partial^\alpha\Psi\right)=\sum_{\mathrm{gv}'\le\mathrm{gv}}\binom{\mathrm{gv}}{\mathrm{gv}'}\cdot\left(\partial^{\mathrm{gv}'}y^{\mathrm{bz}}\right)\cdot\partial^{\mathrm{gv}-\mathrm{gv}'+\alpha}\Psi.
\]
Here $\sum_{\mathrm{gv}'}\binom{\mathrm{gv}}{\mathrm{gv}'}=2^{|\mathrm{gv}|}\le2^m$.
The factor $\partial^{\mathrm{gv}'}y^{\mathrm{bz}}$ vanishes unless $\mathrm{gv}'\le\mathrm{bz}$, and in that case
\[
\left|\partial^{\mathrm{gv}'}y^{\mathrm{bz}}\right|\le\left(\prod_\ell\mathrm{bz}_\ell!\right)\cdot|y|_2^{|\mathrm{bz}|-|\mathrm{gv}'|}\le|\mathrm{bz}|!\cdot\left(1+|y|_2^2\right)^{\frac{|\mathrm{bz}|}{2}}.
\]
Put $m':=m+|\alpha|+|\mathrm{bz}|$.
Since $|\mathrm{gv}-\mathrm{gv}'+\alpha|\le m'$, every component satisfies
\[
\left|\partial^{\mathrm{gv}-\mathrm{gv}'+\alpha}\Psi(y)\right|\le q_{m'}(\Psi)\cdot\left(1+|y|_2^2\right)^{-\frac{m'}{2}}.
\]
Therefore
\[
\left(1+|y|_2^2\right)^{\frac{m}{2}}\cdot\left|\partial^{\mathrm{gv}}\left(y^{\mathrm{bz}}\partial^\alpha\Psi\right)(y)\right|\le2^m\cdot|\mathrm{bz}|!\cdot\left(1+|y|_2^2\right)^{\frac{m+|\mathrm{bz}|-m'}{2}}\cdot q_{m'}(\Psi)\le2^m\cdot|\mathrm{bz}|!\cdot q_{m'}(\Psi).
\]

(Nm2).
By the definition of $q_{2n}$ and by (Nm0), every component satisfies
\[
\left|\Psi_{\mathrm{il}_1\dots\mathrm{il}_n}(y)\right|^2\le q_{2n}(\Psi)^2\cdot\left(1+|y|_2^2\right)^{-2n}\le q_{2n}(\Psi)^2\cdot\prod_{a=1}^n\left(1+|y_a|_2^2\right)^{-2}.
\]
Moreover
\[
\int_{\mathbb{R}^3}\mathrm{d}y\,\left(1+|y|_2^2\right)^{-2}=4\pi\int_0^\infty\mathrm{d}t\,t^2\cdot\left(1+t^2\right)^{-2}=\pi^2,
\]
and there are $3^n$ components.

(Nm3).
This follows from the rule of Leibniz, with $\sum_{\mathrm{gv}'\le\mathrm{gv}}\binom{\mathrm{gv}}{\mathrm{gv}'}\le2^m$.

(Nm4).
By (B.1d), the matrix of $\wedge^2\mathrm{Rl}$ in the basis $(e_{\mathrm{il}})$ is orthogonal.
Hence each component of the image of a tensor under $(\wedge^2\mathrm{Rl})^{\otimes n}$ is a sum of the $3^n$ components of this tensor, with coefficients of modulus at most $1$.
One has $\mathrm{Gi}^{-1}y_a=\mathrm{Rl}^{-1}(y_a-b)$.
The derivative $\partial_{(y_a)_\ell}$ of $y\mapsto\Psi(\mathrm{Gi}^{-1}y_1,\dots,\mathrm{Gi}^{-1}y_n)$ is the sum over $\ell'=1,2,3$ of $(\mathrm{Rl}^{-1})_{\ell'\ell}$ times the derivative of $\Psi$ in $(y_a)_{\ell'}$ at the transformed point.
Consequently a derivative of order $|\mathrm{gv}|\le m$ is a sum of at most $3^m$ derivatives of $\Psi$ of the same order, with coefficients of modulus at most $1$, taken at the point $y':=(\mathrm{Gi}^{-1}y_a)_a$.
Here $|y'|_2=|y-(b,\dots,b)|_2$ and $|(b,\dots,b)|_2^2=n|b|_2^2$; so (Nm0) gives
\[
\left(1+|y|_2^2\right)^{\frac{m}{2}}\le\left(2\left(1+n|b|_2^2\right)\right)^{\frac{m}{2}}\cdot\left(1+|y'|_2^2\right)^{\frac{m}{2}}.
\]

(Nm5).
One has $\left((\tau_{\epsilon e_j}\cdot)^{\otimes n}\Psi\right)(y)=\Psi(y-\epsilon\,\mathrm{vq})$, where $\mathrm{vq}:=(e_j,\dots,e_j)\in\mathbb{R}^{3n}$.
For $|\mathrm{gv}|\le m$,
\[
\partial^{\mathrm{gv}}\Psi(y-\epsilon\,\mathrm{vq})-\partial^{\mathrm{gv}}\Psi(y)=-\epsilon\int_0^1\mathrm{d}t\,\sum_{a=1}^n\left(\partial_{(y_a)_j}\partial^{\mathrm{gv}}\Psi\right)(y-t\epsilon\,\mathrm{vq}).
\]
Every component satisfies $\left|\left(\partial_{(y_a)_j}\partial^{\mathrm{gv}}\Psi\right)(y')\right|\le q_{m+1}(\Psi)\cdot\left(1+|y'|_2^2\right)^{-\frac{m}{2}}$.
For $y'=y-t\epsilon\,\mathrm{vq}$ one has $|t\epsilon\,\mathrm{vq}|_2^2\le n$, and (Nm0) gives
\[
\left(1+|y|_2^2\right)^{\frac{m}{2}}\le\left(2(1+n)\right)^{\frac{m}{2}}\cdot\left(1+|y'|_2^2\right)^{\frac{m}{2}}.
\]

(Nm6).
The components of $\Psi\otimes\Phi$ are the products of a component of $\Psi$, taken at $(y_1,\dots,y_n)$, and a component of $\Phi$, taken at $(y_{n+1},\dots,y_{n+n'})$.
A derivative of order at most $m$ of such a product is the product of a derivative of order at most $m$ of the first factor and a derivative of order at most $m$ of the second factor.
Moreover $1+|y|_2^2+|y'|_2^2\le\left(1+|y|_2^2\right)\cdot\left(1+|y'|_2^2\right)$.

(Nm7).
The function $q_m$ is a seminorm on $\mathrm{Sw}_n$, and $q_m\le q_{m'}$ for $m\le m'$.
The function $\mathrm{Osc}_a\Psi$ is the sum of the seven functions $\Psi$, $(y_a)_\ell^2\Psi$ and $-\partial_{(y_a)_\ell}^2\Psi$, $\ell=1,2,3$.
One has $q_m(\Psi)\le q_{m+2}(\Psi)$.
By (Nm1) with $\alpha=0$ and $|\mathrm{bz}|=2$, $q_m\left((y_a)_\ell^2\Psi\right)\le2^m\cdot2!\cdot q_{m+2}(\Psi)$.
Further $q_m\left(\partial_{(y_a)_\ell}^2\Psi\right)\le q_{m+2}(\Psi)$, because for $|\mathrm{gv}|\le m$ the derivative $\partial^{\mathrm{gv}}\partial_{(y_a)_\ell}^2$ has order at most $m+2$, and $\left(1+|y|_2^2\right)^{\frac{m}{2}}\le\left(1+|y|_2^2\right)^{\frac{m+2}{2}}$.
Therefore
\[
q_m(\mathrm{Osc}_a\Psi)\le\left(1+3\cdot2^{m+1}+3\right)\cdot q_{m+2}(\Psi)=\left(4+6\cdot2^m\right)\cdot q_{m+2}(\Psi)\le10\cdot2^m\cdot q_{m+2}(\Psi).
\]
As this holds for every $m$, the smooth function $\mathrm{Osc}_a\Psi$ lies in $\mathrm{Sw}_n$.
The operator $(\mathrm{Osc}^t)^{\otimes n}$ is a product of $nt$ operators $\mathrm{Osc}_a$.
Applying the inequality $nt$ times, with the indices $m,m+2,\dots,m+2(nt-1)$, one obtains the factor
\[
\prod_{i=0}^{nt-1}\left(10\cdot2^{m+2i}\right)=10^{nt}\cdot2^{nt(m+nt-1)},
\]
because $\sum_{i=0}^{nt-1}(m+2i)=nt\cdot m+nt\cdot(nt-1)$.
For $t=0$ both sides equal $q_m(\Psi)$.

(Nm8).
For complex $\mathrm{fs},\mathrm{fs}'\in\mathcal{S}(\mathbb{R}^3)$,
\[
\int\mathrm{d}y\,\mathrm{fs}\cdot(\mathrm{Osc}\,\mathrm{fs}')=\int\mathrm{d}y\,(\mathrm{Osc}\,\mathrm{fs})\cdot\mathrm{fs}',
\]
by two integrations by parts in each coordinate; there are no boundary terms, because $\mathrm{fs}$, $\mathrm{fs}'$ and their derivatives decrease faster than every power.
Let $\mathrm{kk}=\left((\mathrm{kn}_a,\mathrm{il}_a)\right)_a$.
The only non-zero component of $h_{\mathrm{kk}}$ is the real function $\prod_a\mathrm{hf}_{\mathrm{kn}_a}(y_a)$, at the place $(\mathrm{il}_1,\dots,\mathrm{il}_n)$; hence
\[
c_{\mathrm{kk}}(\Psi)=\int\mathrm{d}y\,\prod_a\mathrm{hf}_{\mathrm{kn}_a}(y_a)\cdot\Psi_{\mathrm{il}_1\dots\mathrm{il}_n}(y).
\]
By Fubini's theorem, the identity above may be applied in the $a$-th argument, with $\mathrm{fs}=\mathrm{hf}_{\mathrm{kn}_a}$ and with $\mathrm{fs}'$ the function $y_a\mapsto\Psi_{\mathrm{il}_1\dots\mathrm{il}_n}(y_1,\dots,y_n)$, which lies in $\mathcal{S}(\mathbb{R}^3)$ for fixed values of the other arguments.
Since $\mathrm{Osc}\,\mathrm{hf}_{\mathrm{kn}_a}=\lambda_{\kappa_a}\cdot\mathrm{hf}_{\mathrm{kn}_a}$, it follows that
\[
c_{\mathrm{kk}}(\mathrm{Osc}_a\Psi)=\lambda_{\kappa_a}\cdot c_{\mathrm{kk}}(\Psi).
\]
The operator $\mathrm{Osc}_a$ maps $\mathrm{Sw}_n$ into itself, by (Nm7); so iteration gives the first formula.
The second formula is Parseval's identity for the orthonormal basis $(h_{\mathrm{kk}})$ of B.2 and the element $(\mathrm{Osc}^t)^{\otimes n}\Psi$ of $\mathrm{Sw}_n$, which is square integrable by (Nm2).
\end{proof}

By (Nm1), every differential operator with polynomial coefficients maps $\mathrm{Sw}_n$ continuously into itself.

\subsubsection*{B.5. Euclidean maps; the change of sign}

Let $\mathrm{Gi}\colon y\mapsto\mathrm{Rl}\,y+b$ and $\mathrm{Gi}'\colon y\mapsto\mathrm{Rl}'\,y+b'$ be Euclidean maps, with $\mathrm{Rl},\mathrm{Rl}'\in O(3)$;
let $O$ be an open subset of $\mathbb{R}^3$ and $\mathrm{lw}$ a law.
The actions $f\mapsto\mathrm{Gi}.f$, $\omega\mapsto\mathrm{Gi}\,\omega$, $\mathrm{Ga}\mapsto\widehat{\mathrm{Gi}}\,\mathrm{Ga}$ and the image law $\mathrm{Gi}_*\mathrm{lw}$ are defined as in (n2).
There the translation part is called $a$;
here, as in (Nm4), it is denoted $b$, since the letter $a$ indexes the arguments $y_a$.
\begin{itemize}
\item[(Eu1)] By (B.1d) the map $\wedge^2\mathrm{Rl}$ is an orthogonal map of $\Lambda^2$, and by (B.1b) one has $\wedge^2(\mathrm{Rl}\,\mathrm{Rl}')=(\wedge^2\mathrm{Rl})(\wedge^2\mathrm{Rl}')$.
Consequently
\[
\left(\mathrm{Gi}.\left(\mathrm{Gi}'.f\right)\right)(y)=(\wedge^2\mathrm{Rl})(\wedge^2\mathrm{Rl}')\,f\left(\mathrm{Gi}'^{-1}\mathrm{Gi}^{-1}y\right)=\left(\left(\mathrm{Gi}\,\mathrm{Gi}'\right).f\right)(y).
\]
Since $y\mapsto\mathrm{Gi}^{-1}y$ preserves the Lebesgue measure, $\|\mathrm{Gi}.f\|_2=\|f\|_2$;
moreover $\operatorname{supp}(\mathrm{Gi}.f)=\mathrm{Gi}(\operatorname{supp}f)$.
The map $f\mapsto\mathrm{Gi}.f$ is a linear bijection of $\mathcal{S}$, continuous by (Nm4), whose inverse is $f\mapsto\mathrm{Gi}^{-1}.f$;
it maps $C_c^\infty(O;\Lambda^2)$ onto $C_c^\infty(\mathrm{Gi}\,O;\Lambda^2)$.
\item[(Eu2)] For $\omega\in\mathcal{S}'$ the functional $\mathrm{Gi}\,\omega=\omega\circ(\mathrm{Gi}^{-1}.\,)$ belongs to $\mathcal{S}'$ by (Eu1), and
\[
\left(\mathrm{Gi}\left(\mathrm{Gi}'\omega\right)\right)(f)=\left(\mathrm{Gi}'\omega\right)\left(\mathrm{Gi}^{-1}.f\right)=\omega\left(\mathrm{Gi}'^{-1}.\mathrm{Gi}^{-1}.f\right)=\left(\left(\mathrm{Gi}\,\mathrm{Gi}'\right)\omega\right)(f).
\]
Thus $\omega\mapsto\mathrm{Gi}\,\omega$ is a bijection of $\mathcal{S}'$, with the inverse $\omega\mapsto\mathrm{Gi}^{-1}\omega$.
One has $X(f)(\mathrm{Gi}^{-1}\omega)=\omega(\mathrm{Gi}.f)$, that is,
\[
X(f)\circ\mathrm{Gi}^{-1}=X(\mathrm{Gi}.f).
\]
For $f\in C_c^\infty(O;\Lambda^2)$ the function $X(\mathrm{Gi}.f)$ is $\Sigma_{\mathrm{Gi}\,O}$-measurable.
Therefore $\omega\mapsto\mathrm{Gi}^{-1}\omega$ is measurable from $(\mathcal{S}',\Sigma_{\mathrm{Gi}\,O})$ to $(\mathcal{S}',\Sigma_O)$;
the same argument with $f\in\mathcal{S}$ shows that it is measurable from $(\mathcal{S}',\Sigma)$ to $(\mathcal{S}',\Sigma)$.
\item[(Eu3)] One has $\widehat{\mathrm{Gi}}\,\mathrm{Ga}=\mathrm{Ga}\circ\mathrm{Gi}^{-1}$.
Hence $\widehat{\mathrm{Gi}}\,X(f)=X(\mathrm{Gi}.f)$ and $\widehat{\mathrm{Gi}\,\mathrm{Gi}'}=\widehat{\mathrm{Gi}}\,\widehat{\mathrm{Gi}'}$;
the operator $\widehat{\mathrm{Gi}}$ is linear and multiplicative and commutes with complex conjugation;
and, by (Eu2) applied to $\mathrm{Gi}$ and to $\mathrm{Gi}^{-1}$, it maps the $\Sigma_O$-measurable functions onto the $\Sigma_{\mathrm{Gi}\,O}$-measurable functions.
\item[(Eu4)] For $\mathrm{Ev}\in\Sigma$ let $(\mathrm{Gi}_*\mathrm{lw})(\mathrm{Ev}):=\mathrm{lw}(\mathrm{Gi}^{-1}\mathrm{Ev})$, where $\mathrm{Gi}^{-1}\mathrm{Ev}=\{\omega:\mathrm{Gi}\,\omega\in\mathrm{Ev}\}$.
By (Eu2) for $\mathrm{Gi}^{-1}$ this set belongs to $\Sigma$, so that $\mathrm{Gi}_*\mathrm{lw}$ is a law.
For measurable $\mathrm{Ga}\ge0$ and for $\mathrm{Ga}\in L^1(\mathrm{Gi}_*\mathrm{lw})$ one has
\[
(\mathrm{Gi}_*\mathrm{lw})(\mathrm{Ga})=\mathrm{lw}(\mathrm{Ga}\circ\mathrm{Gi})=\mathrm{lw}\left(\widehat{\mathrm{Gi}^{-1}}\,\mathrm{Ga}\right);
\]
moreover $(\mathrm{Gi}\,\mathrm{Gi}')_*\mathrm{lw}=\mathrm{Gi}_*(\mathrm{Gi}'_*\mathrm{lw})$, and
\[
\Phi_{\mathrm{Gi}_*\mathrm{lw}}(f)=\mathrm{lw}\left(\exp\left(iX(f)\circ\mathrm{Gi}\right)\right)=\Phi_{\mathrm{lw}}\left(\mathrm{Gi}^{-1}.f\right)\qquad\text{for real }f\in\mathcal{S}.
\]
Equip the set of the laws with the $\sigma$-algebra of (n7).
For every $\mathrm{Ev}\in\Sigma$ the function $\mathrm{lw}\mapsto(\mathrm{Gi}_*\mathrm{lw})(\mathrm{Ev})=\mathrm{lw}(\mathrm{Gi}^{-1}\mathrm{Ev})$ is one of the functions that generate this $\sigma$-algebra;
hence the map $\mathrm{lw}\mapsto\mathrm{Gi}_*\mathrm{lw}$ of the set of the laws into itself is measurable.
\item[(Eu5)] By (Eu4), the maps $\mathrm{Gi}$ with $\mathrm{Gi}_*\mathrm{lw}=\mathrm{lw}$ form a group.
Suppose that $\mathrm{Gi}_*\mathrm{lw}=\mathrm{lw}$.
Then
\[
\mathrm{lw}\left(\widehat{\mathrm{Gi}}\,\mathrm{Ga}\right)=\left((\mathrm{Gi}^{-1})_*\mathrm{lw}\right)(\mathrm{Ga})=\mathrm{lw}(\mathrm{Ga}).
\]
So $\widehat{\mathrm{Gi}}$ maps functions that agree $\mathrm{lw}$-almost everywhere to functions that agree $\mathrm{lw}$-almost everywhere, acts on the classes, preserves $\|\cdot\|_{L^2(\mathrm{lw})}$ and has the inverse $\widehat{\mathrm{Gi}^{-1}}$:
it is a unitary operator on $L^2(\mathrm{lw})$.
By (Eu3), $\widehat{\mathrm{Gi}}\,L^2(\Sigma_O,\mathrm{lw})=L^2(\Sigma_{\mathrm{Gi}\,O},\mathrm{lw})$.
\item[(Eu6)] The map $\mathrm{sgn}\colon\omega\mapsto-\omega$ is a bijection of $\mathcal{S}'$ with $\mathrm{sgn}\circ\mathrm{sgn}=1$, and $X(f)\circ\mathrm{sgn}=-X(f)=X(-f)$;
since the $X(f)$ generate $\Sigma$, the map $\mathrm{sgn}$ is measurable.
With $(\mathrm{sgn}_*\mathrm{lw})(\mathrm{Ev}):=\mathrm{lw}(\mathrm{sgn}^{-1}\mathrm{Ev})$ one has $\Phi_{\mathrm{sgn}_*\mathrm{lw}}(f)=\Phi_{\mathrm{lw}}(-f)$.
\item[(Eu7)] For real $t$ one has $\theta_j\tau_{te_j}=\tau_{-te_j}\theta_j$, because $\theta_j(y+te_j)=\theta_jy-te_j$;
for $y\in e_j^\perp$ one has $\theta_j\tau_y=\tau_y\theta_j$, because $\theta_jy=y$.
Moreover
\[
r_{e_j,\mathrm{po}}=\tau_{\mathrm{po}\,e_j}\,\theta_j\,\tau_{-\mathrm{po}\,e_j}:
\]
the right side sends $y$ to $y-\mathrm{po}\,e_j$, then to $\theta_jy+\mathrm{po}\,e_j$, and finally to $\theta_jy+2\,\mathrm{po}\,e_j$, which replaces $y_j$ by $2\,\mathrm{po}-y_j$.
Finally, $\tau_{te_j}\{y_j>\mathrm{po}\}=\{y_j>\mathrm{po}+t\}$.
\end{itemize}

\subsubsection*{B.6. Two lemmas on the $\sigma$-algebra}

For a real linear subspace $\mathrm{Dq}$ of $\mathcal{S}$, denote by $\mathrm{Cyl}(\mathrm{Dq})$ the family of the sets
\[
\left\{\omega\in\mathcal{S}':\left(\omega(f_1),\dots,\omega(f_m)\right)\in\mathrm{Bq}\right\},
\]
where $m\ge1$, $f_1,\dots,f_m\in\mathrm{Dq}$, and $\mathrm{Bq}$ ranges over the Borel sets of $\mathbb{R}^m$.

\begin{named}{Lemma B.1}
Let $\mathrm{lw}$ and $\mathrm{lw}'$ be probability measures on $(\mathcal{S}',\Sigma)$, and let $\mathrm{Dq}$ be a real linear subspace of $\mathcal{S}$ such that $\Phi_{\mathrm{lw}}(f)=\Phi_{\mathrm{lw}'}(f)$ for all $f\in\mathrm{Dq}$.
Then $\mathrm{lw}(\mathrm{Ev})=\mathrm{lw}'(\mathrm{Ev})$ for every $\mathrm{Ev}\in\sigma\left(X(f):f\in\mathrm{Dq}\right)$.
\end{named}

\begin{proof}
Let $f_1,\dots,f_m\in\mathrm{Dq}$.
For $t\in\mathbb{R}^m$, the characteristic function of the random vector $\left(X(f_1),\dots,X(f_m)\right)$ under $\mathrm{lw}$, evaluated at $t$, is
\[
\mathrm{lw}\left(\exp\left(i\sum_at_aX(f_a)\right)\right)=\Phi_{\mathrm{lw}}\left(\sum_at_af_a\right),
\]
and $\sum_at_af_a$ belongs to $\mathrm{Dq}$.
The vector therefore has the same characteristic function under $\mathrm{lw}$ and under $\mathrm{lw}'$, and so, by the uniqueness theorem for characteristic functions, the same law on $\mathbb{R}^m$.
In other words, $\mathrm{lw}$ and $\mathrm{lw}'$ agree on $\mathrm{Cyl}(\mathrm{Dq})$.
The family $\mathrm{Cyl}(\mathrm{Dq})$ is stable under finite intersections:
the intersection of the set formed with $(f_1,\dots,f_m;\mathrm{Bq})$ and the set formed with $(f'_1,\dots,f'_{m'};\mathrm{Bq}')$ is the set formed with $(f_1,\dots,f_m,f'_1,\dots,f'_{m'};\mathrm{Bq}\times\mathrm{Bq}')$.
It generates $\sigma\left(X(f):f\in\mathrm{Dq}\right)$, since it contains the sets $\{X(f)\in\mathrm{Bq}\}$ and consists of elements of this $\sigma$-algebra.
The sets $\mathrm{Ev}$ of this $\sigma$-algebra with $\mathrm{lw}(\mathrm{Ev})=\mathrm{lw}'(\mathrm{Ev})$ include $\mathcal{S}'$ and are stable under complements and under countable disjoint unions.
By Dynkin's theorem they exhaust the $\sigma$-algebra.
\end{proof}

Fix a function $\mathrm{bmp}\in C_c^\infty(\mathbb{R}^3)$ with $0\le\mathrm{bmp}\le1$ and $\mathrm{bmp}=1$ on $\{|y|_2\le1\}$.

\begin{named}{Lemma B.2}
\begin{itemize}
\item[(a)] Let $f\in\mathcal{S}$, or let $f$ be a scalar function in $\mathcal{S}(\mathbb{R}^3)$, and for $\mathrm{dn}=1,2,\dots$ put $f_{\mathrm{dn}}:=\mathrm{bmp}(\cdot/\mathrm{dn})\,f$.
Then $f_{\mathrm{dn}}$ is smooth with compact support, $\operatorname{supp}f_{\mathrm{dn}}\subset\operatorname{supp}f$, and, for every integer $m\ge0$,
\[
q_m(f-f_{\mathrm{dn}})\le2^m\cdot\mathrm{cf}_m\cdot q_{m+1}(f)\cdot\left(1+\mathrm{dn}^2\right)^{-\frac12},\qquad\mathrm{cf}_m:=\max\left(1,\max_{1\le|\mathrm{gv}|\le m}\sup\left|\partial^{\mathrm{gv}}\mathrm{bmp}\right|\right).
\]
\item[(b)] $\Sigma_{\mathbb{R}^3}=\Sigma$.
\end{itemize}
\end{named}

\begin{proof}
(a) The function $1-\mathrm{bmp}(\cdot/\mathrm{dn})$ vanishes on $\{|y|_2\le\mathrm{dn}\}$ and takes its values in $[0,1]$;
since $\mathrm{dn}\ge1$, its derivatives of the orders $1,\dots,m$ are bounded by $\mathrm{cf}_m$.
For $|\mathrm{gv}|\le m$ the rule of Leibniz therefore gives
\begin{align*}
\left(1+|y|_2^2\right)^{\frac m2}\left|\partial^{\mathrm{gv}}(f-f_{\mathrm{dn}})(y)\right|
&\le2^m\cdot\mathrm{cf}_m\cdot\mathbf{1}\left[|y|_2\ge\mathrm{dn}\right]\cdot\left(1+|y|_2^2\right)^{\frac m2}\max_{|\mathrm{gv}'|\le m}\left|\partial^{\mathrm{gv}'}f(y)\right|\\
&\le2^m\cdot\mathrm{cf}_m\cdot q_{m+1}(f)\cdot\left(1+\mathrm{dn}^2\right)^{-\frac12}.
\end{align*}
(b) Let $f\in\mathcal{S}$.
By (a) and the continuity of $\omega$, one has $\omega(f)=\lim_{\mathrm{dn}}\omega(f_{\mathrm{dn}})$ for every $\omega\in\mathcal{S}'$.
Thus $X(f)$ is a pointwise limit of $\Sigma_{\mathbb{R}^3}$-measurable functions.
As the $X(f)$, $f\in\mathcal{S}$, generate $\Sigma$, this gives $\Sigma\subset\Sigma_{\mathbb{R}^3}$;
the reverse inclusion holds because $C_c^\infty(\mathbb{R}^3;\Lambda^2)\subset\mathcal{S}$.
\end{proof}

\subsubsection*{B.7. Closed forms}

Put
\[
\mathrm{Ecl}:=\left\{\omega\in\mathcal{S}':\omega(d^{*}\mathrm{gs})=0\text{ for every real }\mathrm{gs}\in\mathcal{S}(\mathbb{R}^3)\right\};
\]
this is the event $\{dF=0\text{ in }\mathcal{S}'\}$ of (n8).
Up to signs, the components of $d^{*}\mathrm{gs}$ are the three first derivatives of $\mathrm{gs}$.
Hence $q_m(d^{*}\mathrm{gs})\le q_{m+1}(\mathrm{gs})$, and $\mathrm{gs}\mapsto d^{*}\mathrm{gs}$ is continuous from $\mathcal{S}(\mathbb{R}^3)$ to $\mathcal{S}$.

\begin{named}{Lemma B.3}
There are real functions $\mathrm{gs}_1,\mathrm{gs}_2,\dots$ in $C_c^\infty(\mathbb{R}^3)$ such that
\[
\mathrm{Ecl}=\bigcap_a\left\{X(d^{*}\mathrm{gs}_a)=0\right\}.
\]
In particular, $\mathrm{Ecl}$ belongs to $\Sigma$.
\end{named}

\begin{proof}
Let $\{\mathrm{gs}_a\}$ be the countable set of the functions $\mathrm{bmp}(\cdot/\mathrm{dn})\,\mathrm{gs}'$, where $\mathrm{dn}=1,2,\dots$ and $\mathrm{gs}'$ is a finite linear combination of the $\mathrm{hf}_{\mathrm{kn}}$ with rational coefficients.
These are real elements of $C_c^\infty(\mathbb{R}^3)$.
Let $\mathrm{gs}\in\mathcal{S}(\mathbb{R}^3)$ be real, and let $m\ge1$.
By (He2) there is a real linear combination $f'$ of the $h_\kappa$ with $q_m(\mathrm{gs}\,e_{12}-f')<\frac{1}{3m}$;
its component $f'_{12}$ is a real linear combination of the $\mathrm{hf}_{\mathrm{kn}}$, and $q_m(\mathrm{gs}-f'_{12})<\frac{1}{3m}$.
On replacing the finitely many coefficients of $f'_{12}$ by rational numbers close to them one obtains a function $\mathrm{gs}'$ with $q_m(f'_{12}-\mathrm{gs}')<\frac{1}{3m}$.
By Lemma B.2\,(a) there is $\mathrm{dn}$ with $q_m\left(\mathrm{gs}'-\mathrm{bmp}(\cdot/\mathrm{dn})\,\mathrm{gs}'\right)<\frac{1}{3m}$.
Consequently there is an index $a(m)$ with $q_m(\mathrm{gs}-\mathrm{gs}_{a(m)})<\frac1m$.
Since $q_{m'}\le q_m$ for $m'\le m$, it follows that $\mathrm{gs}_{a(m)}\to\mathrm{gs}$ in $\mathcal{S}(\mathbb{R}^3)$ as $m\to\infty$.
Now let $\omega\in\mathcal{S}'$ satisfy $\omega(d^{*}\mathrm{gs}_a)=0$ for all $a$.
Then $d^{*}\mathrm{gs}_{a(m)}\to d^{*}\mathrm{gs}$ in $\mathcal{S}$, and by the continuity of $\omega$
\[
\omega(d^{*}\mathrm{gs})=\lim_m\omega(d^{*}\mathrm{gs}_{a(m)})=0.
\]
So the intersection is contained in $\mathrm{Ecl}$.
The reverse inclusion is immediate from the definition of $\mathrm{Ecl}$.
\end{proof}

\subsubsection*{B.8. The sub-Gaussian estimates for maps into classes}

\begin{named}{Lemma B.4}
Let $\mathrm{lw}$ be a law, $\mathrm{Ts}$ a real vector space with a Hilbert seminorm $|\cdot|$, and $X$ a real-linear map from $\mathrm{Ts}$ into the space of the classes, modulo $\mathrm{lw}$-null sets, of real measurable functions on $(\mathcal{S}',\Sigma)$, such that
\[
\mathrm{lw}\left(\exp X(\mathrm{ts})\right)\le\exp\left(\frac{|\mathrm{ts}|^2}{2}\right)\qquad\text{for }\mathrm{ts}\in\mathrm{Ts}.
\]
Then the conclusions (1)--(7) of \cite[Proposition~7.10]{QEDI} hold for $(\mathcal{S}',\Sigma,\mathrm{lw})$, $\mathrm{Ts}$, $|\cdot|$ and $X$.
\end{named}

\begin{proof}
Each of the conclusions is a statement on finitely many elements $f_1,\dots,f_n$ of $\mathrm{Ts}+i\,\mathrm{Ts}$;
it involves $X$ only through integrals of functions of $X(f_1),\dots,X(f_n)$, and these integrals do not depend on the choice of representatives.
Let $\mathrm{Ts}_0$ be the real linear span of the $\operatorname{Re}f_i$ and $\operatorname{Im}f_i$, let $\mathrm{ts}_1,\dots,\mathrm{ts}_{\mathrm{dn}}$ be a basis of $\mathrm{Ts}_0$, let $\mathrm{Yq}_1,\dots,\mathrm{Yq}_{\mathrm{dn}}$ be real measurable functions in the classes $X(\mathrm{ts}_1),\dots,X(\mathrm{ts}_{\mathrm{dn}})$, and define
\[
X_0\left(\sum_bt_b\,\mathrm{ts}_b\right):=\sum_bt_b\,\mathrm{Yq}_b.
\]
Then $X_0$ is a linear map from $\mathrm{Ts}_0$ to real random variables, and, $X$ being linear, $X_0(\mathrm{ts})$ belongs to the class $X(\mathrm{ts})$ for every $\mathrm{ts}\in\mathrm{Ts}_0$.
Hence the exponential bound required in \cite[Proposition~7.10]{QEDI} holds for $X_0$ on $(\mathrm{Ts}_0,|\cdot|)$, and that proposition, applied to $(\mathrm{Ts}_0,|\cdot|,X_0)$, yields the conclusion.
\end{proof}

\subsubsection*{B.9. (He1) and (He2), proved from the ladder relations}

The Hermite functions $\mathrm{hf}^1_{\mathrm{dn}}$, $\mathrm{dn}\in\mathrm{Nat}$, on $\mathbb{R}$ satisfy the ladder relations \cite[Appendix to Section~V.3]{ReedSimon1980}
\[
(y+\partial)\,\mathrm{hf}^1_{\mathrm{dn}}=(2\,\mathrm{dn})^{\frac12}\,\mathrm{hf}^1_{\mathrm{dn}-1},\qquad(y-\partial)\,\mathrm{hf}^1_{\mathrm{dn}}=(2\,\mathrm{dn}+2)^{\frac12}\,\mathrm{hf}^1_{\mathrm{dn}+1}\qquad\left(\mathrm{hf}^1_{-1}:=0\right),
\]
and $\mathrm{hf}_{\mathrm{kn}}(y)=\prod_{\ell=1}^3\mathrm{hf}^1_{\mathrm{kn}_\ell}(y_\ell)$ for $\mathrm{kn}\in\mathrm{Nat}^3$.

\emph{Notation.}
Let $n\ge1$.
The index $\ell$ ranges over the $3n$ coordinates of $\mathbb{R}^{3n}$, and $\mathrm{ue}_\ell$ denotes the element of $\mathrm{Nat}^{3n}$ whose entry at the place $\ell$ is $1$ and whose other entries are $0$.
For $\mathrm{km}\in\mathrm{Nat}^{3n}$ put $\mathrm{hf}_{\mathrm{km}}(y):=\prod_\ell\mathrm{hf}^1_{\mathrm{km}_\ell}(y_\ell)$;
for $\mathrm{km}=(\mathrm{kn}_1,\dots,\mathrm{kn}_n)$ this is $\mathrm{hf}_{\mathrm{kn}_1}(y_1)\cdots\mathrm{hf}_{\mathrm{kn}_n}(y_n)$.
The functions $\mathrm{hf}_{\mathrm{km}}$ are real and form an orthonormal basis of $L^2(\mathbb{R}^{3n})$ (products of orthonormal bases, as in Section~B.2).
For $\mathrm{fs}$ in the complex space $\mathcal{S}(\mathbb{R}^{3n})$ and integers $\mathrm{ro}\ge0$ put
\[
\mathrm{xq}_{\mathrm{km}}(\mathrm{fs}):=\int\mathrm{d}y\,\mathrm{hf}_{\mathrm{km}}(y)\,\mathrm{fs}(y),\qquad\|\mathrm{fs}\|_{(\mathrm{ro})}^2:=\sum_{\mathrm{km}}\left(1+|\mathrm{km}|_1\right)^{\mathrm{ro}}\left|\mathrm{xq}_{\mathrm{km}}(\mathrm{fs})\right|^2.
\]
By Parseval's identity $\|\mathrm{fs}\|_{(0)}=\|\mathrm{fs}\|_{L^2}$, and $\|\mathrm{fs}\|_{(\mathrm{ro})}\le\|\mathrm{fs}\|_{(\mathrm{ro}')}$ for $\mathrm{ro}\le\mathrm{ro}'$.
Let
\[
\mathrm{Aq}_\ell:=\frac{y_\ell+\partial_\ell}{\sqrt2},\qquad\mathrm{Cq}_\ell:=\frac{y_\ell-\partial_\ell}{\sqrt2}.
\]
Then $y_\ell=\frac{\mathrm{Aq}_\ell+\mathrm{Cq}_\ell}{\sqrt2}$ and $\partial_\ell=\frac{\mathrm{Aq}_\ell-\mathrm{Cq}_\ell}{\sqrt2}$, and
\[
\mathrm{Aq}_\ell\,\mathrm{hf}_{\mathrm{km}}=\mathrm{km}_\ell^{\frac12}\,\mathrm{hf}_{\mathrm{km}-\mathrm{ue}_\ell}\quad\left(=0\text{ if }\mathrm{km}_\ell=0\right),\qquad\mathrm{Cq}_\ell\,\mathrm{hf}_{\mathrm{km}}=\left(\mathrm{km}_\ell+1\right)^{\frac12}\,\mathrm{hf}_{\mathrm{km}+\mathrm{ue}_\ell}.
\]

\begin{named}{Lemma B.5}
Let $n\ge1$ and $m\ge0$, put $\mathrm{wl}:=6n+3m$, and let $t$ be an integer with $2t\ge\mathrm{wl}$.
Then, for all $\Psi\in\mathrm{Sw}_n$,
\[
q_m(\Psi)\le\mathrm{cf}'_{n,m}\cdot\left\|(\mathrm{Osc}^t)^{\otimes n}\Psi\right\|_{L^2},
\]
where
\[
\mathrm{cf}'_{n,m}:=(3n+1)^{m+3n}\cdot2^{\frac{\mathrm{wl}(\mathrm{wl}+1)}{4}}\cdot(2\pi)^{-\frac{3n}{2}}\cdot\left(\int_{\mathbb{R}^{3n}}\mathrm{d}p\,\left(1+|p|_2^2\right)^{-6n}\right)^{\frac12}<\infty.
\]
\end{named}

\begin{proof}
\emph{Step 1.}
Let $\mathrm{fs}=\Psi_{\mathrm{il}_1\cdots\mathrm{il}_n}$ be a component of $\Psi$.
For $\mathrm{km}=(\mathrm{kn}_1,\dots,\mathrm{kn}_n)$ and $\mathrm{kk}:=\left((\mathrm{kn}_a,\mathrm{il}_a)\right)_a$ one has $\mathrm{xq}_{\mathrm{km}}(\mathrm{fs})=c_{\mathrm{kk}}(\Psi)$ and
\[
1+|\mathrm{km}|_1=1+\sum_a|\mathrm{kn}_a|_1\le\prod_a\left(1+|\mathrm{kn}_a|_1\right)\le\prod_a\lambda_{\kappa_a}.
\]
For $\mathrm{ro}\le2t$, (Nm8) and the inequality $\prod_i\lambda_{\kappa_i}\ge1$ therefore give
\begin{equation}\tag{B.9a}
\|\mathrm{fs}\|_{(\mathrm{ro})}^2\le\sum_{\mathrm{kk}}\left(\prod_i\lambda_{\kappa_i}\right)^{2t}\left|c_{\mathrm{kk}}(\Psi)\right|^2=\left\|(\mathrm{Osc}^t)^{\otimes n}\Psi\right\|_{L^2}^2<\infty.
\end{equation}
Every $\mathrm{fs}\in\mathcal{S}(\mathbb{R}^{3n})$ is a component of an element of $\mathrm{Sw}_n$;
hence $\|\mathrm{fs}\|_{(\mathrm{ro})}<\infty$ for all $\mathrm{fs}\in\mathcal{S}(\mathbb{R}^{3n})$ and all $\mathrm{ro}$.

\emph{Step 2.}
Let $\mathrm{fs}\in\mathcal{S}(\mathbb{R}^{3n})$.
One integration by parts, which produces no boundary terms, together with the ladder relations gives
\begin{gather*}
\mathrm{xq}_{\mathrm{km}}(\mathrm{Aq}_\ell\,\mathrm{fs})=\int\mathrm{d}y\,\left(\mathrm{Cq}_\ell\,\mathrm{hf}_{\mathrm{km}}\right)\mathrm{fs}=\left(\mathrm{km}_\ell+1\right)^{\frac12}\,\mathrm{xq}_{\mathrm{km}+\mathrm{ue}_\ell}(\mathrm{fs}),\\
\mathrm{xq}_{\mathrm{km}}(\mathrm{Cq}_\ell\,\mathrm{fs})=\int\mathrm{d}y\,\left(\mathrm{Aq}_\ell\,\mathrm{hf}_{\mathrm{km}}\right)\mathrm{fs}=\mathrm{km}_\ell^{\frac12}\,\mathrm{xq}_{\mathrm{km}-\mathrm{ue}_\ell}(\mathrm{fs})\quad\left(=0\text{ if }\mathrm{km}_\ell=0\right).
\end{gather*}
In the first of the two sums below substitute $\mathrm{km}'=\mathrm{km}+\mathrm{ue}_\ell$, so that $1+|\mathrm{km}|_1=|\mathrm{km}'|_1$ and $\mathrm{km}_\ell+1=\mathrm{km}'_\ell\le1+|\mathrm{km}'|_1$;
in the second substitute $\mathrm{km}'=\mathrm{km}-\mathrm{ue}_\ell$, so that $1+|\mathrm{km}|_1=2+|\mathrm{km}'|_1\le2\left(1+|\mathrm{km}'|_1\right)$ and $\mathrm{km}_\ell=\mathrm{km}'_\ell+1\le1+|\mathrm{km}'|_1$.
This gives
\begin{gather*}
\|\mathrm{Aq}_\ell\,\mathrm{fs}\|_{(\mathrm{ro})}^2=\sum_{\mathrm{km}}\left(1+|\mathrm{km}|_1\right)^{\mathrm{ro}}\left(\mathrm{km}_\ell+1\right)\left|\mathrm{xq}_{\mathrm{km}+\mathrm{ue}_\ell}(\mathrm{fs})\right|^2\le\|\mathrm{fs}\|_{(\mathrm{ro}+1)}^2,\\
\|\mathrm{Cq}_\ell\,\mathrm{fs}\|_{(\mathrm{ro})}^2=\sum_{\mathrm{km}:\,\mathrm{km}_\ell\ge1}\left(1+|\mathrm{km}|_1\right)^{\mathrm{ro}}\,\mathrm{km}_\ell\left|\mathrm{xq}_{\mathrm{km}-\mathrm{ue}_\ell}(\mathrm{fs})\right|^2\le2^{\mathrm{ro}}\,\|\mathrm{fs}\|_{(\mathrm{ro}+1)}^2.
\end{gather*}
Consequently, for $\mathrm{Zl}=y_\ell$ and for $\mathrm{Zl}=\partial_\ell$,
\[
\|\mathrm{Zl}\,\mathrm{fs}\|_{(\mathrm{ro})}\le\frac{\|\mathrm{Aq}_\ell\,\mathrm{fs}\|_{(\mathrm{ro})}+\|\mathrm{Cq}_\ell\,\mathrm{fs}\|_{(\mathrm{ro})}}{\sqrt2}\le\left(1+2^{\frac{\mathrm{ro}}{2}}\right)2^{-\frac12}\,\|\mathrm{fs}\|_{(\mathrm{ro}+1)}\le2^{\frac{\mathrm{ro}+1}{2}}\,\|\mathrm{fs}\|_{(\mathrm{ro}+1)}.
\]

\emph{Step 3.}
Let $\mathrm{Zl}_1,\dots,\mathrm{Zl}_{\mathrm{dn}'}$, with $\mathrm{dn}'\le\mathrm{wl}$, be operators of the form $y_\ell$ or $\partial_\ell$.
They map $\mathcal{S}(\mathbb{R}^{3n})$ into itself.
Applying Step 2 $\mathrm{dn}'$ times, with $\mathrm{ro}=0,1,\dots,\mathrm{dn}'-1$, one obtains
\begin{align*}
\left\|\mathrm{Zl}_1\cdots\mathrm{Zl}_{\mathrm{dn}'}\,\mathrm{fs}\right\|_{L^2}
&\le2^{\frac12}\left\|\mathrm{Zl}_2\cdots\mathrm{Zl}_{\mathrm{dn}'}\,\mathrm{fs}\right\|_{(1)}\le\dots\le2^{\frac{1+2+\dots+\mathrm{dn}'}{2}}\,\|\mathrm{fs}\|_{(\mathrm{dn}')}\\
&\le2^{\frac{\mathrm{wl}(\mathrm{wl}+1)}{4}}\,\|\mathrm{fs}\|_{(\mathrm{wl})}.
\end{align*}

\emph{Step 4.}
For $\mathrm{fs}'\in\mathcal{S}(\mathbb{R}^{3n})$ one has
\[
\mathrm{fs}'(y)=(2\pi)^{-3n}\int\mathrm{d}p\,\exp(ip\cdot y)\,\widehat{\mathrm{fs}'}(p).
\]
The Cauchy--Schwarz inequality, the identity $\left(1+|p|_2^2\right)^{3n}\widehat{\mathrm{fs}'}=\widehat{(1-\Delta)^{3n}\mathrm{fs}'}$ and Plancherel's theorem on $\mathbb{R}^{3n}$ give
\begin{align*}
\sup|\mathrm{fs}'|
&\le(2\pi)^{-3n}\left(\int\mathrm{d}p\,\left(1+|p|_2^2\right)^{-6n}\right)^{\frac12}\left(\int\mathrm{d}p\,\left(1+|p|_2^2\right)^{6n}\left|\widehat{\mathrm{fs}'}(p)\right|^2\right)^{\frac12}\\
&=(2\pi)^{-\frac{3n}{2}}\left(\int\mathrm{d}p\,\left(1+|p|_2^2\right)^{-6n}\right)^{\frac12}\left\|(1-\Delta)^{3n}\mathrm{fs}'\right\|_{L^2};
\end{align*}
the integral is finite because $12n>3n$.
The operator $(1-\Delta)^{3n}=\left(1-\sum_\ell\partial_\ell^2\right)^{3n}$ is a sum of $(3n+1)^{3n}$ operators $\pm\partial^{\mathrm{gv}}$ with $|\mathrm{gv}|\le6n$.

\emph{Step 5.}
Let $|\alpha|\le m$.
One has $\left(1+|y|_2^2\right)^{\frac m2}\le\left(1+|y|_2^2\right)^m=\left(1+\sum_\ell y_\ell^2\right)^m$, and the right side is a sum of $(3n+1)^m$ monomials $y^{\mathrm{bz}}\ge0$ with $|\mathrm{bz}|\le2m$.
Step 4 for $\mathrm{fs}'=y^{\mathrm{bz}}\partial^\alpha\mathrm{fs}$, combined with Step 3 for the words $\partial^{\mathrm{gv}}y^{\mathrm{bz}}\partial^\alpha$, whose length is at most $6n+2m+m=\mathrm{wl}$, yields
\begin{align*}
\sup\left(1+|y|_2^2\right)^{\frac m2}\left|\partial^\alpha\mathrm{fs}\right|
&\le\sum_{\mathrm{bz}}\sup\left|y^{\mathrm{bz}}\partial^\alpha\mathrm{fs}\right|\\
&\le(3n+1)^m\cdot(2\pi)^{-\frac{3n}{2}}\cdot\left(\int\mathrm{d}p\,\left(1+|p|_2^2\right)^{-6n}\right)^{\frac12}\cdot(3n+1)^{3n}\cdot2^{\frac{\mathrm{wl}(\mathrm{wl}+1)}{4}}\cdot\|\mathrm{fs}\|_{(\mathrm{wl})}\\
&=\mathrm{cf}'_{n,m}\cdot\|\mathrm{fs}\|_{(\mathrm{wl})}.
\end{align*}
Together with (B.9a) for $\mathrm{ro}=\mathrm{wl}\le2t$, and after taking the maximum over the components and over $\alpha$, this is the assertion.
\end{proof}

\begin{proof}[Proof of (He2), first sentence]
Let $\Psi\in\mathrm{Sw}_n$.
For $\mathrm{dn}=1,2,\dots$ let $\Psi_{\mathrm{dn}}$ be the sum of the terms $c_{\mathrm{kk}}(\Psi)\,h_{\mathrm{kk}}$ over the finite set $\left\{\mathrm{kk}:\lambda_{\kappa_i}\le\mathrm{dn}\text{ for all }i\right\}$;
it belongs to $\mathrm{Fin}_n$.
By orthonormality, $c_{\mathrm{kk}}(\Psi-\Psi_{\mathrm{dn}})=0$ for $\mathrm{kk}$ in this set, and $c_{\mathrm{kk}}(\Psi-\Psi_{\mathrm{dn}})=c_{\mathrm{kk}}(\Psi)$ for $\mathrm{kk}$ outside it.
Let $m\ge0$, and let $t$ be as in Lemma B.5.
Lemma B.5 and (Nm8) for $\Psi-\Psi_{\mathrm{dn}}$ give
\[
q_m(\Psi-\Psi_{\mathrm{dn}})^2\le\left(\mathrm{cf}'_{n,m}\right)^2\sum_{\mathrm{kk}\text{ outside the set}}\left(\prod_i\lambda_{\kappa_i}\right)^{2t}\left|c_{\mathrm{kk}}(\Psi)\right|^2.
\]
By (Nm8) for $\Psi$ the series over all $\mathrm{kk}$ converges, and the sets increase to $\mathrm{Kh}^n$.
Hence $q_m(\Psi-\Psi_{\mathrm{dn}})\to0$ as $\mathrm{dn}\to\infty$, for every $m$.
\end{proof}

\begin{proof}[Proof of (He2), second sentence]
Let $f\in\mathcal{S}$.
By the first sentence for $n=1$, where $\mathrm{Sw}_1=\mathcal{S}+i\mathcal{S}$, there are $\Psi_{\mathrm{dn}}\in\mathrm{Fin}_1$ with $\Psi_{\mathrm{dn}}\to f$ in $\mathrm{Sw}_1$.
The $h_\kappa$ are real, so the real part of a complex linear combination of the $h_\kappa$ is the combination with the real parts of the coefficients;
thus $\operatorname{Re}\Psi_{\mathrm{dn}}$ lies in the real linear span of the $h_\kappa$.
As $f$ is real,
\[
q_m\left(\operatorname{Re}\Psi_{\mathrm{dn}}-f\right)=q_m\left(\operatorname{Re}(\Psi_{\mathrm{dn}}-f)\right)\le q_m(\Psi_{\mathrm{dn}}-f),
\]
because $\operatorname{Re}$ commutes with the derivatives and $|\operatorname{Re}z|\le|z|$.
\end{proof}

\begin{proof}[Proof of (He1)]
Let $\omega\in\mathcal{S}'$.
Since the norms $q_m$ increase with $m$, the sets $\{f\in\mathcal{S}:q_m(f)<\varepsilon\}$ form a base of the neighbourhoods of $0$ in $\mathcal{S}$.
As $\omega$ is continuous, there are $m$ and $\varepsilon>0$ such that $|\omega(f)|\le1$ whenever $q_m(f)<\varepsilon$;
by homogeneity, $|\omega(f)|\le\frac{2}{\varepsilon}\cdot q_m(f)$ for all $f\in\mathcal{S}$.
Lemma B.5 with $n=1$ and an integer $t$ with $2t\ge6+3m$, together with $\mathrm{Osc}^th_\kappa=\lambda_\kappa^th_\kappa$, gives
\[
|\omega(h_\kappa)|\le\frac{2}{\varepsilon}\cdot\mathrm{cf}'_{1,m}\cdot\left\|\mathrm{Osc}^th_\kappa\right\|_{L^2}=\frac{2}{\varepsilon}\cdot\mathrm{cf}'_{1,m}\cdot\lambda_\kappa^t.
\]
Hence, by (He3),
\[
\|\omega\|_{-s}^2\le\left(\frac{2}{\varepsilon}\cdot\mathrm{cf}'_{1,m}\right)^2\mathrm{ZH}_{2s-2t}<\infty\qquad\text{for }2s-2t>3.
\]
\end{proof}

\subsection*{C. Proposition 45.1}

\subsubsection*{Proof}

Let $\mathrm{lw}$ be a law with (L1)--(L4);
in the proofs of (10) and (11), $\mathrm{lw}$ is an arbitrary law.
In what follows, \cite[Proposition~7.10(1)]{QEDI}, \dots, \cite[Proposition~7.10(7)]{QEDI} stand for the conclusions (1)--(7) of \cite[Proposition~7.10]{QEDI} (see Section~B.3) for the probability space $(\mathcal{S}',\Sigma,\mathrm{lw})$, the space $L^2(\mathbb{R}^3;\Lambda^2)$, the norm $\|\cdot\|_2$ and the map $X$ of clause (1);
by clause (1), these conclusions hold as soon as its proof is complete.

\begin{proof}[Proof of Proposition~45.1\,(1)]
\emph{Step 1.}
The functions $X(f)$, $f\in\mathcal{S}$, are real random variables on $(\mathcal{S}',\Sigma,\mathrm{lw})$, the map $f\mapsto X(f)$ is linear, and $\|\cdot\|_2$ is a Hilbert norm on $\mathcal{S}$;
for this space, norm and map, the exponential bound required in \cite[Proposition~7.10]{QEDI} is (L1).
Conclusion (2) of \cite[Proposition~7.10]{QEDI}, applied to the space $\mathcal{S}$ with the norm $\|\cdot\|_2$ and the map $f\mapsto X(f)$, therefore gives
\begin{equation}\tag{C.1}
\mathrm{lw}\left(X(f)\right)=0,\qquad\mathrm{lw}\left(X(f)^2\right)\le\|f\|_2^2\qquad\text{for }f\in\mathcal{S}.
\end{equation}

\emph{Step 2 (extension).}
The space $\mathcal{S}$ contains $C_c^\infty(\mathbb{R}^3;\Lambda^2)$, which is dense in $L^2(\mathbb{R}^3;\Lambda^2)$.
By (C.1), the real-linear map from $\mathcal{S}$ to $L^2(\mathrm{lw})$ that sends $f$ to the class of $X(f)$ satisfies $\|X(f)\|_{L^2(\mathrm{lw})}\le\|f\|_2$.
It therefore has exactly one continuous extension $X$ to $L^2(\mathbb{R}^3;\Lambda^2)$.
This extension is real-linear and satisfies $\|X(f)\|_{L^2(\mathrm{lw})}\le\|f\|_2$;
moreover $X(f)$ is a class of real functions, because the classes of real functions form a closed subset of $L^2(\mathrm{lw})$.
For $f,f'\in L^2(\mathbb{R}^3;\Lambda^2)$ put $X(f+if'):=X(f)+iX(f')$;
then
\[
\|X(f+if')\|_{L^2(\mathrm{lw})}^2=\|X(f)\|_{L^2(\mathrm{lw})}^2+\|X(f')\|_{L^2(\mathrm{lw})}^2\le\|f+if'\|_2^2.
\]

\emph{Step 3 (the bound).}
Let $f\in L^2(\mathbb{R}^3;\Lambda^2)$, and let $f_a\in\mathcal{S}$ with $\|f_a-f\|_2\to0$.
Then $X(f_a)\to X(f)$ in $L^2(\mathrm{lw})$, and also $\mathrm{lw}$-almost surely along a subsequence.
Along this subsequence, Fatou's lemma and (L1) give
\[
\mathrm{lw}\left(\exp X(f)\right)\le\liminf_a\mathrm{lw}\left(\exp X(f_a)\right)\le\lim_a\exp\left(\frac{\|f_a\|_2^2}{2}\right)=\exp\left(\frac{\|f\|_2^2}{2}\right).
\]

\emph{Step 4.}
By Steps 2 and 3, the hypotheses of Lemma B.4 are satisfied by $\mathrm{Ts}=L^2(\mathbb{R}^3;\Lambda^2)$, $|\cdot|=\|\cdot\|_2$ and $X$.
Hence the conclusions (1)--(7) of \cite[Proposition~7.10]{QEDI} hold.
\end{proof}

\begin{proof}[Proof of Proposition~45.1\,(3)]
Let $s>\frac32$.
By monotone convergence, \cite[Proposition~7.10(2)]{QEDI} and $\|h_\kappa\|_2=1$,
\[
\mathrm{lw}\left(\|\omega\|_{-s}^2\right)=\sum_\kappa\lambda_\kappa^{-2s}\cdot\mathrm{lw}\left(X(h_\kappa)^2\right)\le\sum_\kappa\lambda_\kappa^{-2s}=\mathrm{ZH}_{2s},
\]
and the right side is finite by (He3), as $2s>3$.
Hence $\|\omega\|_{-s}<\infty$ for $\mathrm{lw}$-almost every $\omega$, that is, $\mathrm{lw}(H_{-s})=1$;
here $H_{-s}$ belongs to $\Sigma$ by Section~B.2.
Let $0\le2t\cdot\mathrm{ZH}_{2s}<1$, and put $\mathrm{pw}_\kappa:=\frac{\lambda_\kappa^{-2s}}{\mathrm{ZH}_{2s}}$;
these are positive numbers with the sum $1$.
For every $\omega\in H_{-s}$, Jensen's inequality for the convex function $\exp$ and the probability weights $\mathrm{pw}_\kappa$ on $\mathrm{Kh}$ gives
\[
\exp\left(t\|\omega\|_{-s}^2\right)=\exp\left(\sum_\kappa\mathrm{pw}_\kappa\cdot t\cdot\mathrm{ZH}_{2s}\cdot\omega(h_\kappa)^2\right)\le\sum_\kappa\mathrm{pw}_\kappa\cdot\exp\left(t\cdot\mathrm{ZH}_{2s}\cdot\omega(h_\kappa)^2\right),
\]
the sum in the exponent being finite.
The complement of $H_{-s}$ is an $\mathrm{lw}$-null set.
Integrating over $H_{-s}$ (by monotone convergence) and applying \cite[Proposition~7.10(3)]{QEDI} with $a=t\cdot\mathrm{ZH}_{2s}$, $\|h_\kappa\|_2=1$ and $0\le2a<1$, one obtains
\[
\mathrm{lw}\left(\exp\left(t\|\omega\|_{-s}^2\right)\right)\le\sum_\kappa\mathrm{pw}_\kappa\cdot\left(1-2t\cdot\mathrm{ZH}_{2s}\right)^{-\frac12}=\left(1-2t\cdot\mathrm{ZH}_{2s}\right)^{-\frac12}.
\]
\end{proof}

\begin{proof}[Proof of Proposition~45.1\,(4)]
The Fourier transform is $\widehat{\mathrm{fs}}(p):=\int\mathrm{d}y\,\exp(-ip\cdot y)\,\mathrm{fs}(y)$ (Section~I.9.7), and Plancherel's theorem reads
\[
\int\mathrm{d}p\,\left|\widehat{\mathrm{fs}}(p)\right|^2=(2\pi)^3\int\mathrm{d}y\,\left|\mathrm{fs}(y)\right|^2.
\]
For $\mathrm{fs}\in L^2(\mathbb{R}^3)$, $\widehat{\mathrm{fs}}$ denotes the image of $\mathrm{fs}$ under the continuous extension to $L^2(\mathbb{R}^3)$ of the map $\mathrm{fs}\mapsto\widehat{\mathrm{fs}}$ on $\mathcal{S}(\mathbb{R}^3)$;
by Plancherel's theorem this extension exists and satisfies the same identity.
The relation $\widehat{\mathrm{bv}}(-p)=\overline{\widehat{\mathrm{bv}}(p)}$ for real $\mathrm{bv}$ remains valid for the extension, by continuity.
The same holds for the relation $\widehat{\overline{\mathrm{bv}}}(p)=\overline{\widehat{\mathrm{bv}}(-p)}$ for complex $\mathrm{bv}$;
since the extension is injective, it follows that a square integrable $\mathrm{bv}$ with $\widehat{\mathrm{bv}}(-p)=\overline{\widehat{\mathrm{bv}}(p)}$ for almost every $p$ is real.

\emph{Step 1 (closedness).}
By Lemma B.3, $\mathrm{Ecl}$ is the intersection of countably many events $\{X(d^{*}\mathrm{gs}_a)=0\}$ with real $\mathrm{gs}_a\in C_c^\infty(\mathbb{R}^3)$.
By (L2), each of them has $\mathrm{lw}$-measure $1$.
Hence $\mathrm{lw}(\mathrm{Ecl})=1$, which is the statement $\mathrm{lw}\{dF=0\text{ in }\mathcal{S}'\}=1$ (Section~B.7).
In particular,
\begin{equation}\tag{C.2}
X(d^{*}\mathrm{gs})=0\quad\mathrm{lw}\text{-almost surely, for every real }\mathrm{gs}\in\mathcal{S}(\mathbb{R}^3).
\end{equation}

\emph{Step 2 (the projection).}
The map $\ast\colon\Lambda^2\to\mathbb{R}^3$, $\ast f=(f_{23},-f_{13},f_{12})$, is an isometry.
For $p\ne0$ the matrix
\[
\Pi(p)=1-\frac{p\,p^T}{|p|_2^2}
\]
is real and symmetric, and $\Pi(p)^2=\Pi(p)$, $\Pi(-p)=\Pi(p)$.
By Plancherel's theorem, the operator $\Pi_{\mathrm{c}}$ defined by $\widehat{\ast\Pi_{\mathrm{c}}f}(p)=\Pi(p)\,\widehat{\ast f}(p)$ is an orthogonal projection of $L^2(\mathbb{R}^3;\Lambda^2)\otimes\mathbb{C}$.
A square integrable function $\mathrm{bv}\colon\mathbb{R}^3\to\mathbb{C}^3$ is real if and only if $\widehat{\mathrm{bv}}(-p)=\overline{\widehat{\mathrm{bv}}(p)}$ for almost every $p$, and multiplication by $\Pi(p)$ preserves this relation.
Hence $\Pi_{\mathrm{c}}$ maps $L^2(\mathbb{R}^3;\Lambda^2)$ into itself, and
\begin{equation}\tag{C.3}
\|\Pi_{\mathrm{c}}f\|_2^2=\langle f,\Pi_{\mathrm{c}}\,f\rangle\le\|f\|_2^2\qquad\text{for }f\in L^2(\mathbb{R}^3;\Lambda^2).
\end{equation}

\emph{Step 3 ($(1-\Pi_{\mathrm{c}})f$ is a limit of forms $d^{*}\mathrm{gs}$).}
Let $f\in L^2(\mathbb{R}^3;\Lambda^2)$, put $\mathrm{bv}:=\ast f$, and for $p\ne0$ let
\[
\mathrm{uq}(p):=-i\,\frac{p\cdot\widehat{\mathrm{bv}}(p)}{|p|_2^2}.
\]
Then
\begin{gather*}
\widehat{\ast(1-\Pi_{\mathrm{c}})f}(p)=\frac{p\left(p\cdot\widehat{\mathrm{bv}}(p)\right)}{|p|_2^2}=ip\,\mathrm{uq}(p),\\
\int\mathrm{d}p\,|p|_2^2\left|\mathrm{uq}(p)\right|^2=\int\mathrm{d}p\,\frac{\left|p\cdot\widehat{\mathrm{bv}}(p)\right|^2}{|p|_2^2}\le\int\mathrm{d}p\,\left|\widehat{\mathrm{bv}}(p)\right|^2<\infty,\\
\mathrm{uq}(-p)=i\,\frac{p\cdot\widehat{\mathrm{bv}}(-p)}{|p|_2^2}=i\,\frac{p\cdot\overline{\widehat{\mathrm{bv}}(p)}}{|p|_2^2}=\overline{\mathrm{uq}(p)}.
\end{gather*}
Thus $p\mapsto|p|_2\,\mathrm{uq}(p)$ belongs to $L^2(\mathbb{R}^3,\mathrm{d}p)$.
The point $0$ is a null set, and $C_c^\infty(\mathbb{R}^3\setminus\{0\})$ is dense in $L^2(\mathbb{R}^3\setminus\{0\},\mathrm{d}p)$;
choose $\mathrm{uq}'_a\in C_c^\infty(\mathbb{R}^3\setminus\{0\})$ with
\[
\delta_a:=\left(\int\mathrm{d}p\,\left|\mathrm{uq}'_a(p)-|p|_2\,\mathrm{uq}(p)\right|^2\right)^{\frac12}\to0,
\]
and put
\[
\mathrm{uq}_a(p):=\frac{\mathrm{uq}'_a(p)+\overline{\mathrm{uq}'_a(-p)}}{2\,|p|_2}.
\]
Since $p\mapsto\frac{1}{|p|_2}$ is smooth on $\mathbb{R}^3\setminus\{0\}$, the function $\mathrm{uq}_a$ belongs to $C_c^\infty(\mathbb{R}^3\setminus\{0\})$, and $\mathrm{uq}_a(-p)=\overline{\mathrm{uq}_a(p)}$.
As $|p|_2\,\mathrm{uq}(p)=\frac12\left(|p|_2\,\mathrm{uq}(p)+\overline{|p|_2\,\mathrm{uq}(-p)}\right)$, the triangle inequality in $L^2(\mathrm{d}p)$ and the invariance of $\mathrm{d}p$ under $p\mapsto-p$ give
\[
\left(\int\mathrm{d}p\,|p|_2^2\left|\mathrm{uq}_a(p)-\mathrm{uq}(p)\right|^2\right)^{\frac12}\le\frac{\delta_a}{2}+\frac{\delta_a}{2}=\delta_a.
\]
Let
\[
\mathrm{gs}_a(y):=(2\pi)^{-3}\int\mathrm{d}p\,\exp(ip\cdot y)\,\mathrm{uq}_a(p).
\]
This function belongs to $\mathcal{S}(\mathbb{R}^3)$, it is real because $\mathrm{uq}_a(-p)=\overline{\mathrm{uq}_a(p)}$, and $\widehat{\operatorname{grad}\mathrm{gs}_a}(p)=ip\,\mathrm{uq}_a(p)$.
By Plancherel's theorem,
\[
\left\|\operatorname{grad}\mathrm{gs}_a-\ast(1-\Pi_{\mathrm{c}})f\right\|_{L^2}^2=(2\pi)^{-3}\int\mathrm{d}p\,|p|_2^2\left|\mathrm{uq}_a(p)-\mathrm{uq}(p)\right|^2\le(2\pi)^{-3}\delta_a^2\to0.
\]
Since $\ast(d^{*}\mathrm{gs}_a)=-\operatorname{grad}\mathrm{gs}_a$ and $\ast$ is an isometry, it follows that
\[
\left\|-d^{*}\mathrm{gs}_a-(1-\Pi_{\mathrm{c}})f\right\|_2\to0.
\]

\emph{Step 4.}
By (1), $X\left((1-\Pi_{\mathrm{c}})f\right)=\lim_aX(-d^{*}\mathrm{gs}_a)$ in $L^2(\mathrm{lw})$, and $X(d^{*}\mathrm{gs}_a)=0$ in $L^2(\mathrm{lw})$ by (C.2).
Hence $X\left((1-\Pi_{\mathrm{c}})f\right)=0$, that is, $X(f)=X(\Pi_{\mathrm{c}}f)$ $\mathrm{lw}$-almost surely.

\emph{Step 5.}
By Step 4, (1) and (C.3),
\[
\mathrm{lw}\left(\exp X(f)\right)=\mathrm{lw}\left(\exp X(\Pi_{\mathrm{c}}f)\right)\le\exp\left(\frac{\|\Pi_{\mathrm{c}}f\|_2^2}{2}\right)=\exp\left(\frac{\langle f,\Pi_{\mathrm{c}}\,f\rangle}{2}\right).
\]
The map $f\mapsto\langle f,\Pi_{\mathrm{c}}\,f\rangle^{\frac12}=\|\Pi_{\mathrm{c}}f\|_2$ is a Hilbert seminorm on $L^2(\mathbb{R}^3;\Lambda^2)$.
By Lemma B.4 for this seminorm, conclusion (2) of \cite[Proposition~7.10]{QEDI} holds, and it gives $0\le\mathrm{lw}\left(X(f)^2\right)\le\langle f,\Pi_{\mathrm{c}}\,f\rangle$;
finally, $\langle f,\Pi_{\mathrm{c}}\,f\rangle\le\|f\|_2^2$ is (C.3).
\end{proof}

\begin{proof}[Proof of Proposition~45.1\,(5)]
\emph{Step 1 (invariance).}
Let $b\in\mathbb{R}^3$.
By (Eu4) and the first equality of (L3) for $-b$,
\[
\Phi_{(\tau_b)_*\mathrm{lw}}(f)=\Phi_{\mathrm{lw}}(\tau_{-b}.f)=\Phi_{\mathrm{lw}}(f)\qquad\text{for all real }f\in\mathcal{S}.
\]
Lemma B.1 with $\mathrm{Dq}=\mathcal{S}$ shows that $(\tau_b)_*\mathrm{lw}=\mathrm{lw}$ on $\sigma\left(X(f):f\in\mathcal{S}\right)=\Sigma$.

\emph{Step 2 (representation).}
By (Eu5), $U(b)=\hat\tau_b$ is unitary on $L^2(\mathrm{lw})$;
by (Eu3), $U(b)\,U(b')=\widehat{\tau_b\tau_{b'}}=U(b+b')$, and $U(0)=1$.

\emph{Step 3 (continuity).}
Let $f\in\mathcal{S}$.
By (Eu3), $U(b)\exp\left(iX(f)\right)=\exp\left(iX(\tau_b.f)\right)$.
Since $\left|\exp(i\,\mathrm{rv})-\exp(i\,\mathrm{rv}')\right|\le|\mathrm{rv}-\mathrm{rv}'|$ for real $\mathrm{rv},\mathrm{rv}'$, (C.1) gives
\[
\left\|U(b)\exp\left(iX(f)\right)-\exp\left(iX(f)\right)\right\|_{L^2(\mathrm{lw})}^2\le\mathrm{lw}\left(X(\tau_b.f-f)^2\right)\le\|\tau_b.f-f\|_2^2.
\]
From
\[
(\tau_b.f)(y)-f(y)=f(y-b)-f(y)=-\int_0^1\mathrm{d}t\,\sum_\ell b_\ell\,(\partial_\ell f)(y-tb),
\]
Minkowski's inequality and the invariance of the Lebesgue measure under translations one obtains
\[
\|\tau_b.f-f\|_2\le\sum_\ell|b_\ell|\cdot\|\partial_\ell f\|_2\le|b|_2\left(\sum_\ell\|\partial_\ell f\|_2^2\right)^{\frac12}.
\]
Let $\mathrm{Ga}\in L^2(\mathrm{lw})$ and $\varepsilon>0$.
By \cite[Lemma~7.14]{QEDI} with $\mathrm{Dq}=\mathcal{S}$ there is a finite sum $\mathrm{Gb}=\sum_pc_p\exp\left(iX(f_p)\right)$, with $c_p\in\mathbb{C}$ and $f_p\in\mathcal{S}$, such that $\|\mathrm{Ga}-\mathrm{Gb}\|_{L^2(\mathrm{lw})}<\varepsilon$.
Then, the norms being those of $L^2(\mathrm{lw})$,
\begin{align*}
\|U(b)\mathrm{Ga}-\mathrm{Ga}\|
&\le\|U(b)(\mathrm{Ga}-\mathrm{Gb})\|+\|U(b)\mathrm{Gb}-\mathrm{Gb}\|+\|\mathrm{Gb}-\mathrm{Ga}\|\\
&\le2\varepsilon+|b|_2\sum_p|c_p|\left(\sum_\ell\|\partial_\ell f_p\|_2^2\right)^{\frac12}.
\end{align*}
Hence $\limsup_{b\to0}\|U(b)\mathrm{Ga}-\mathrm{Ga}\|\le2\varepsilon$ for every $\varepsilon>0$.
For $b_0\in\mathbb{R}^3$ one has $\|U(b)\mathrm{Ga}-U(b_0)\mathrm{Ga}\|=\|U(b-b_0)\mathrm{Ga}-\mathrm{Ga}\|\to0$ as $b\to b_0$.
\end{proof}

\begin{proof}[Proof of Proposition~45.1\,(6)]
Let $r\in O_{\mathrm{cub}}$.
Then $r^{-1}\in O_{\mathrm{cub}}$, and by (Eu4) and the second equality of (L3),
\[
\Phi_{r_*\mathrm{lw}}(f)=\Phi_{\mathrm{lw}}(r^{-1}.f)=\Phi_{\mathrm{lw}}(f)\qquad\text{for all real }f\in\mathcal{S}.
\]
By Lemma B.1 with $\mathrm{Dq}=\mathcal{S}$, $r_*\mathrm{lw}=\mathrm{lw}$.
\end{proof}

\begin{proof}[Proof of Proposition~45.1\,(7)]
Let $j\in\{1,2,3\}$ and $\mathrm{po}\in\mathbb{R}$, and put $O_{\mathrm{po}}:=\{y_j>\mathrm{po}\}$.

\emph{Step 1 (invariance).}
The map $\theta_j$ is the diagonal matrix with the entry $-1$ at the place $j$ and the entries $+1$ elsewhere;
it belongs to $O_{\mathrm{cub}}$.
By (Eu7), $r_{e_j,\mathrm{po}}=\tau_{\mathrm{po}\,e_j}\,\theta_j\,\tau_{-\mathrm{po}\,e_j}$.
By (5), (6) and (Eu5), according to which the maps that preserve $\mathrm{lw}$ form a group, $(r_{e_j,\mathrm{po}})_*\mathrm{lw}=\mathrm{lw}$.
Hence, by (Eu5), $\hat r_{e_j,\mathrm{po}}$ is a unitary operator on $L^2(\mathrm{lw})$;
it acts on classes.
By (Eu3),
\begin{equation}\tag{C.4}
\hat r_{e_j,\mathrm{po}}=U(\mathrm{po}\,e_j)\,\hat\theta_j\,U(-\mathrm{po}\,e_j).
\end{equation}

\emph{Step 2 ($\mathrm{po}=0$).}
Let $c_p\in\mathbb{C}$ and let $f_p\in C_c^\infty(O_0;\Lambda^2)$ be real, $p$ ranging over a finite index set, and put $\mathrm{Ga}:=\sum_pc_p\exp\left(iX(f_p)\right)$.
By (Eu3), $\hat\theta_j\mathrm{Ga}=\sum_pc_p\exp\left(iX(\theta_j.f_p)\right)$.
Hence, by (L4),
\begin{align*}
\left\langle\hat\theta_j\mathrm{Ga},\mathrm{Ga}\right\rangle_{\mathrm{lw}}
&=\mathrm{lw}\left(\overline{\hat\theta_j\mathrm{Ga}}\cdot\mathrm{Ga}\right)=\sum_{p,q}\overline{c_p}\cdot c_q\cdot\mathrm{lw}\left(\exp\left(iX(f_q-\theta_j.f_p)\right)\right)\\
&=\sum_{p,q}\overline{c_p}\cdot c_q\cdot\Phi_{\mathrm{lw}}\left(f_q-\theta_j.f_p\right)\ge0.
\end{align*}
By \cite[Lemma~7.14]{QEDI}, with $\mathrm{Dq}$ the real space $C_c^\infty(O_0;\Lambda^2)$, these $\mathrm{Ga}$ are dense in $L^2(\Sigma_{O_0},\mathrm{lw})$.
For $\mathrm{Ga},\mathrm{Gb}\in L^2(\mathrm{lw})$, the unitarity of $\hat\theta_j$ gives, with the norms of $L^2(\mathrm{lw})$,
\begin{align*}
\left|\langle\hat\theta_j\mathrm{Ga},\mathrm{Ga}\rangle_{\mathrm{lw}}-\langle\hat\theta_j\mathrm{Gb},\mathrm{Gb}\rangle_{\mathrm{lw}}\right|
&\le\left|\langle\hat\theta_j(\mathrm{Ga}-\mathrm{Gb}),\mathrm{Ga}\rangle_{\mathrm{lw}}\right|+\left|\langle\hat\theta_j\mathrm{Gb},\mathrm{Ga}-\mathrm{Gb}\rangle_{\mathrm{lw}}\right|\\
&\le\|\mathrm{Ga}-\mathrm{Gb}\|\left(\|\mathrm{Ga}\|+\|\mathrm{Gb}\|\right).
\end{align*}
So the function $\mathrm{Ga}\mapsto\langle\hat\theta_j\mathrm{Ga},\mathrm{Ga}\rangle_{\mathrm{lw}}$ is continuous on $L^2(\mathrm{lw})$, and it is $\ge0$ on $L^2(\Sigma_{O_0},\mathrm{lw})$.

\emph{Step 3 (general $\mathrm{po}$).}
Let $\mathrm{Ga}\in L^2(\Sigma_{O_{\mathrm{po}}},\mathrm{lw})$, and put $\mathrm{Gb}:=U(-\mathrm{po}\,e_j)\mathrm{Ga}$.
Since $\tau_{-\mathrm{po}\,e_j}O_{\mathrm{po}}=O_0$ by (Eu7), $\mathrm{Gb}$ belongs to $L^2(\Sigma_{O_0},\mathrm{lw})$ by (Eu5).
By (C.4) and the unitarity of $U(\mathrm{po}\,e_j)$,
\[
\langle\hat r_{e_j,\mathrm{po}}\mathrm{Ga},\mathrm{Ga}\rangle_{\mathrm{lw}}=\langle U(\mathrm{po}\,e_j)\,\hat\theta_j\mathrm{Gb},U(\mathrm{po}\,e_j)\mathrm{Gb}\rangle_{\mathrm{lw}}=\langle\hat\theta_j\mathrm{Gb},\mathrm{Gb}\rangle_{\mathrm{lw}}\ge0.
\]

\emph{Step 4 (representatives).}
Let $f\in L^2(\mathbb{R}^3;\Lambda^2)$ vanish almost everywhere on $\{y_j<\mathrm{po}\}$.
The plane $\{y_j=\mathrm{po}\}$ is a null set, so $f$ vanishes almost everywhere outside $O_{\mathrm{po}}$.
Since $C_c^\infty(O_{\mathrm{po}};\Lambda^2)$ is dense in $L^2(O_{\mathrm{po}};\Lambda^2)$, one may choose $f_a\in C_c^\infty(O_{\mathrm{po}};\Lambda^2)$ with $\|f_a-f\|_2\le2^{-a}$.
By (1) and the Cauchy--Schwarz inequality,
\[
\mathrm{lw}\left(\sum_a\left|X(f_a)-X(f)\right|\right)\le\sum_a\left\|X(f_a-f)\right\|_{L^2(\mathrm{lw})}\le\sum_a2^{-a}<\infty;
\]
hence $X(f_a)\to X(f)$ $\mathrm{lw}$-almost surely.
Let $\mathrm{Yq}(\omega):=\lim_a\omega(f_a)$ where this limit exists in $\mathbb{R}$, and $\mathrm{Yq}(\omega):=0$ elsewhere.
The set on which a sequence of $\Sigma_{O_{\mathrm{po}}}$-measurable functions converges belongs to $\Sigma_{O_{\mathrm{po}}}$;
hence $\mathrm{Yq}$ is $\Sigma_{O_{\mathrm{po}}}$-measurable, and $\mathrm{Yq}=X(f)$ $\mathrm{lw}$-almost surely.

\emph{Step 5 (polynomials, exponentials).}
Let $f_1,\dots,f_m$ be functions as $f$ in Step 4, and let $\mathrm{Yq}_1,\dots,\mathrm{Yq}_m$ be the corresponding representatives.
A polynomial with complex coefficients in $\mathrm{Yq}_1,\dots,\mathrm{Yq}_m$ is $\Sigma_{O_{\mathrm{po}}}$-measurable, and it belongs to $L^2(\mathrm{lw})$:
the square of its modulus is bounded by a finite sum of multiples of products $|X(f_{a_1})|\cdots|X(f_{a_{\mathrm{dn}}})|$, and by \cite[Proposition~7.10(4)]{QEDI}
\[
\mathrm{lw}\left(\prod_{i=1}^{\mathrm{dn}}\left|X(f_{a_i})\right|\right)\le\sigma_{\mathrm{dn}}\prod_i\|f_{a_i}\|_2.
\]
For $z\in\mathbb{C}^m$, the function $\exp\left(\sum_az_a\mathrm{Yq}_a\right)$ is $\Sigma_{O_{\mathrm{po}}}$-measurable, and by (1)
\[
\mathrm{lw}\left(\left|\exp\sum_az_aX(f_a)\right|^2\right)=\mathrm{lw}\left(\exp X\left(2\sum_a(\operatorname{Re}z_a)f_a\right)\right)\le\exp\left(2\left\|\sum_a(\operatorname{Re}z_a)f_a\right\|_2^2\right).
\]
Hence these functions belong to $L^2(\Sigma_{O_{\mathrm{po}}},\mathrm{lw})$, and Step 3 applies to their finite linear combinations.
\end{proof}

\begin{proof}[Proof of Proposition~45.1\,(8)]
Let $f\in C_c^\infty(\mathbb{R}^3;\Lambda^2)$ be real.
For real $\mathrm{rv}$ all terms of the series $\cosh\mathrm{rv}=\sum_{a\ge0}\frac{\mathrm{rv}^{2a}}{(2a)!}$ are $\ge0$.
Monotone convergence, the hypothesis on $\mathrm{lw}'$ (with $2a$ tests equal to $f$) and \cite[Proposition~7.10(4)]{QEDI} give
\[
\mathrm{lw}'\left(\cosh X(f)\right)=\sum_a\frac{\mathrm{lw}'\left(X(f)^{2a}\right)}{(2a)!}=\sum_a\frac{\mathrm{lw}\left(X(f)^{2a}\right)}{(2a)!}\le\sum_a\frac{\sigma_{2a}\cdot\|f\|_2^{2a}}{(2a)!},\qquad\frac{\sigma_{2a}}{(2a)!}=\frac{2\cdot2^a\cdot a!}{(2a)!}.
\]
The quotient of the terms with the indices $a+1$ and $a$ is
\[
\|f\|_2^2\cdot\frac{2(a+1)}{(2a+2)(2a+1)}=\frac{\|f\|_2^2}{2a+1}\to0,
\]
so the series converges.
Consequently $\exp|X(f)|\le2\cosh X(f)$ is $\mathrm{lw}'$-integrable.
The partial sums of $\exp\left(iX(f)\right)=\sum_{n\ge0}\frac{\left(iX(f)\right)^n}{n!}$ are bounded in modulus by $\exp|X(f)|$.
Dominated convergence and the hypothesis give
\[
\Phi_{\mathrm{lw}'}(f)=\sum_n\frac{i^n\cdot\mathrm{lw}'\left(X(f)^n\right)}{n!}=\sum_n\frac{i^n\cdot\mathrm{lw}\left(X(f)^n\right)}{n!}.
\]
Under $\mathrm{lw}$ the function $\exp|X(f)|$ is integrable by \cite[Proposition~7.10(1)]{QEDI}, and dominated convergence gives $\Phi_{\mathrm{lw}}(f)=\sum_n\frac{i^n\cdot\mathrm{lw}\left(X(f)^n\right)}{n!}$.
Thus $\Phi_{\mathrm{lw}'}=\Phi_{\mathrm{lw}}$ on the real linear space $\mathrm{Dq}=C_c^\infty(\mathbb{R}^3;\Lambda^2)$.
By Lemma B.1, $\mathrm{lw}'=\mathrm{lw}$ on $\sigma\left(X(f):f\in\mathrm{Dq}\right)=\Sigma_{\mathbb{R}^3}$, and $\Sigma_{\mathbb{R}^3}=\Sigma$ by Lemma B.2\,(b).
\end{proof}

\begin{proof}[Proof of Proposition~45.1\,(9)]
Let $\Phi_{\mathrm{lw}}(-f)=\Phi_{\mathrm{lw}}(f)$ for every real $f\in\mathcal S$.
By (Eu6), $\Phi_{\mathrm{sgn}_*\mathrm{lw}}(f)=\Phi_{\mathrm{lw}}(-f)=\Phi_{\mathrm{lw}}(f)$, and Lemma~B.1 with $\mathrm{Dq}=\mathcal S$ yields $\mathrm{sgn}_*\mathrm{lw}=\mathrm{lw}$.
Now let $n$ be odd and $f_1,\dots,f_n\in\mathcal S$.
The product $X(f_1)\cdots X(f_n)$ is $\mathrm{lw}$-integrable by \cite[Proposition~7.10\,(4)]{QEDI}, and
\[
\left(X(f_1)\cdots X(f_n)\right)\circ\mathrm{sgn}=(-1)^n\,X(f_1)\cdots X(f_n).
\]
Consequently
\[
S_n(f_1,\dots,f_n)=(\mathrm{sgn}_*\mathrm{lw})\left(X(f_1)\cdots X(f_n)\right)=\mathrm{lw}\left(\left(X(f_1)\cdots X(f_n)\right)\circ\mathrm{sgn}\right)=-S_n(f_1,\dots,f_n),
\]
that is, $S_n(f_1,\dots,f_n)=0$.
By multilinearity this remains true for $f_1,\dots,f_n\in\mathcal S+i\mathcal S$.
For $f_i,f_i'\in L^2(\mathbb R^3;\Lambda^2)\otimes\mathbb C$, multilinearity and \cite[Proposition~7.10\,(4)]{QEDI} give
\begin{align*}
\left|S_n(f_1,\dots,f_n)-S_n(f_1',\dots,f_n')\right|
&\le\sum_{i=1}^n\left|S_n\left(f_1',\dots,f_{i-1}',f_i-f_i',f_{i+1},\dots,f_n\right)\right|\\
&\le2^{\frac n2}\cdot\sigma_n\cdot\sum_{i=1}^n\|f_i-f_i'\|_2\cdot\prod_{\ell<i}\|f_\ell'\|_2\cdot\prod_{\ell>i}\|f_\ell\|_2 .
\end{align*}
As $\mathcal S+i\mathcal S$ is dense in $L^2(\mathbb R^3;\Lambda^2)\otimes\mathbb C$, it follows that $S_n(f_1,\dots,f_n)=0$ for all $f_1,\dots,f_n\in L^2(\mathbb R^3;\Lambda^2)\otimes\mathbb C$.
For tests in $\mathrm{Sw}_n$ the corresponding assertion is the last sentence of Proposition~45.2\,(6).
\end{proof}

\begin{proof}[Proof of Proposition~45.1\,(10)]
In this proof $\mathrm{lw}$ is an arbitrary law.

\emph{(L3) and invariance.}
If $\mathrm{lw}$ has (L3), then Step~1 of the proof of (5) and the proof of (6), which use only (L3), (Eu4) and Lemma~B.1, give $(\tau_b)_*\mathrm{lw}=\mathrm{lw}$ and $r_*\mathrm{lw}=\mathrm{lw}$ for all $b\in\mathbb R^3$ and $r\in O_{\mathrm{cub}}$.
Conversely, let $\mathrm{Gi}_*\mathrm{lw}=\mathrm{lw}$ for $\mathrm{Gi}=\tau_{-b}$ and for $\mathrm{Gi}=r^{-1}$, for all $b\in\mathbb R^3$ and $r\in O_{\mathrm{cub}}$.
By (Eu4), $\Phi_{\mathrm{lw}}(\mathrm{Gi}^{-1}.f)=\Phi_{\mathrm{Gi}_*\mathrm{lw}}(f)=\Phi_{\mathrm{lw}}(f)$ for real $f\in\mathcal S$; that is, $\Phi_{\mathrm{lw}}(\tau_b.f)=\Phi_{\mathrm{lw}}(f)=\Phi_{\mathrm{lw}}(r.f)$.

\emph{(H).}
By (n3), the hypothesis (H) consists of (G), the invariance under $\mathbb R^3$ and under $O_{\mathrm{cub}}$, and the reflection positivity at $\{y_j=0\}$ for $j=1,2,3$.
The condition (G) is literally the condition (L1).
The invariance is equivalent to (L3), as has just been shown.
Reflection positivity at $\{y_j=0\}$ is the condition of (n3) with $\hat{n}=e_j$, $\mathrm{po}=0$ and $r_{e_j,0}=\theta_j$; since $X(f_q)-X(\theta_j.f_p)=X(f_q-\theta_j.f_p)$, it is (L4) for this $j$.

\emph{(L2) and (dF).}
If $\mathrm{lw}(\mathrm{Ecl})=1$, which by B.7 is the property (dF), then $X(d^{*}\mathrm{gs})=0$ on $\mathrm{Ecl}$, hence $\mathrm{lw}$-almost surely, for every real $\mathrm{gs}\in C_c^\infty(\mathbb R^3)$; this is (L2).
If $\mathrm{lw}$ has (L2), then Step~1 of the proof of (4), which uses only (L2) and Lemma~B.3, gives $\mathrm{lw}(\mathrm{Ecl})=1$; the condition (L1) does not enter there.

\emph{The last sentence.}
Let $\mathrm{lw}$ have (H) and (dF).
By the paragraph on (H), it has (L1), (L3) and (L4); by the paragraph on (L2) and (dF), it has (L2).
\end{proof}

\begin{proof}[Proof of Proposition~45.1\,(11)]
Let $\mathrm{lw}$ be an arbitrary law.
The proofs above use the following properties of $\mathrm{lw}$:
\begin{itemize}
\item[(1), (3), (8):] (L1);
\item[(4):] (L1) and (L2);
\item[(5):] (L1) and the first equality of (L3);
\item[(6):] the second equality of (L3);
\item[(7):] (L1), (L4), the first equality of (L3) for $b\in\mathbb R e_j$, and the second one for $r=\theta_j$;
\item[(9):] (L1) and the evenness of $\Phi_{\mathrm{lw}}$.
\end{itemize}
In detail, this is seen as follows.
\begin{itemize}
\item[(1):] (L1) enters in Steps~1 and~3.
\item[(3), (8):] these proofs use clause (1), in the form of the conclusions (1)--(4) of \cite[Proposition~7.10]{QEDI}.
\item[(4):] (L2) enters in Step~1, and clause (1) in Steps~4 and~5.
\item[(5):] the first equality of (L3) enters in Step~1, and (C.1) in Step~3.
\item[(7):] Step~1 uses Step~1 of the proof of (5) for $b=\mathrm{po}\,e_j$ and for $b=-\mathrm{po}\,e_j$, and the proof of (6) for $r=\theta_j=\theta_j^{-1}$.
Step~2 uses (L4), and Step~3 uses what Steps~1 and~2 give.
Steps~4 and~5 use clause (1) together with \cite[Proposition~7.10\,(4)]{QEDI}, hence only (L1), and they prove the second sentence of (7) in the form stated in (11).
The one exception is the final assertion of Step~5, that Step~3 applies to the finite linear combinations of the functions considered there: it uses what Step~3 gives.
\item[(9):] the proof uses \cite[Proposition~7.10\,(4)]{QEDI} and the evenness of $\Phi_{\mathrm{lw}}$.
\end{itemize}
All other steps rest on Part~B, where $\mathrm{lw}$ is an arbitrary law.
Hence each clause holds for every law that has the hypotheses named with it.
\end{proof}

\subsection*{D. Proposition 45.2}

\subsubsection*{Three lemmas}

Let $n\ge1$.
The coefficients $c_{\mathrm{kk}}$ and the operators $\mathrm{Osc}_a$ and $(\mathrm{Osc}^t)^{\otimes n}$ are those of B.2.
Sums over $\mathrm{kk}$ extend over all $\mathrm{kk}=(\kappa_1,\dots,\kappa_n)\in\mathrm{Kh}^n$, and
\[
Q_\Psi(\omega)=\sum_{\mathrm{kk}}c_{\mathrm{kk}}(\Psi)\prod_{i=1}^n\omega(h_{\kappa_i})
\]
is the series of Proposition~45.2\,(2).
No law enters the Lemmas~D.1 and~D.2.

\begin{named}{Lemma D.1}[The series]
Let $\Psi\in\mathrm{Sw}_n$, and let $t\ge0$ be an integer.
\begin{itemize}
\item[(a)] The function $(\mathrm{Osc}^t)^{\otimes n}\Psi$ lies in $\mathrm{Sw}_n$, and $c_{\mathrm{kk}}\left((\mathrm{Osc}^t)^{\otimes n}\Psi\right)=\left(\prod_i\lambda_{\kappa_i}\right)^t\cdot c_{\mathrm{kk}}(\Psi)$.
Moreover
\begin{equation}\tag{D.1}
\sum_{\mathrm{kk}}\left(\prod_i\lambda_{\kappa_i}\right)^{2t}\cdot\left|c_{\mathrm{kk}}(\Psi)\right|^2=\left\|\left(\mathrm{Osc}^t\right)^{\otimes n}\Psi\right\|_{L^2}^2,
\end{equation}
\begin{equation}\tag{D.2}
\sum_{\mathrm{kk}}\left(\prod_i\lambda_{\kappa_i}\right)^{t-2}\cdot\left|c_{\mathrm{kk}}(\Psi)\right|\le\mathrm{ZH}_4^{\frac n2}\cdot\left\|\left(\mathrm{Osc}^t\right)^{\otimes n}\Psi\right\|_{L^2}\le\mathrm{cf}_{n,t}\cdot q_{2n+2nt}(\Psi),
\end{equation}
where
\[
\mathrm{cf}_{n,t}:=4^{-n}\cdot(\pi\sqrt3)^n\cdot10^{nt}\cdot2^{nt(2n+nt-1)}.
\]
\item[(b)] Let $\omega\in\mathcal S'$ and $s\ge0$ with $\|\omega\|_{-s}<\infty$ (such an $s$ exists by (He1)), and let $t\ge s+2$.
Then
\begin{equation}\tag{D.3}
\sum_{\mathrm{kk}}\left|c_{\mathrm{kk}}(\Psi)\right|\cdot\prod_i\left|\omega(h_{\kappa_i})\right|\le\|\omega\|_{-s}^n\cdot\mathrm{ZH}_4^{\frac n2}\cdot\left\|\left(\mathrm{Osc}^t\right)^{\otimes n}\Psi\right\|_{L^2}\le\mathrm{cf}_{n,t}\cdot\|\omega\|_{-s}^n\cdot q_{2n+2nt}(\Psi).
\end{equation}
\item[(c)] For every $\omega\in\mathcal S'$ the series $Q_\Psi(\omega)$ converges absolutely, and $\Psi\mapsto Q_\Psi(\omega)$ is a continuous linear functional on $\mathrm{Sw}_n$.
For every $\Psi$ the function $\omega\mapsto Q_\Psi(\omega)$ is $\Sigma$-measurable.
\end{itemize}
\end{named}

\begin{proof}
(a) The first two assertions and (D.1) are (Nm7) and (Nm8).
By the Cauchy--Schwarz inequality and (D.1),
\[
\sum_{\mathrm{kk}}\left(\prod_i\lambda_{\kappa_i}\right)^{-2}\cdot\left(\prod_i\lambda_{\kappa_i}\right)^t\cdot\left|c_{\mathrm{kk}}(\Psi)\right|\le\left(\sum_{\mathrm{kk}}\prod_i\lambda_{\kappa_i}^{-4}\right)^{\frac12}\cdot\left\|\left(\mathrm{Osc}^t\right)^{\otimes n}\Psi\right\|_{L^2},
\qquad
\sum_{\mathrm{kk}\in\mathrm{Kh}^n}\prod_i\lambda_{\kappa_i}^{-4}=\mathrm{ZH}_4^n .
\]
This is the first inequality of (D.2).
By (Nm2), and by (Nm7) with $m=2n$,
\[
\left\|\left(\mathrm{Osc}^t\right)^{\otimes n}\Psi\right\|_{L^2}\le(\pi\sqrt3)^n\cdot q_{2n}\left(\left(\mathrm{Osc}^t\right)^{\otimes n}\Psi\right)\le(\pi\sqrt3)^n\cdot10^{nt}\cdot2^{nt(2n+nt-1)}\cdot q_{2n+2nt}(\Psi).
\]
Together with $\mathrm{ZH}_4^{\frac n2}<4^{-n}$, which holds by (He3), this gives the second inequality of (D.2).

(b) Each term of the series that defines $\|\omega\|_{-s}^2$ is at most the sum of the series: $\lambda_\kappa^{-2s}\,\omega(h_\kappa)^2\le\|\omega\|_{-s}^2$.
Hence $|\omega(h_\kappa)|\le\|\omega\|_{-s}\cdot\lambda_\kappa^s\le\|\omega\|_{-s}\cdot\lambda_\kappa^{t-2}$, because $\lambda_\kappa\ge1$ and $s\le t-2$.
Inserting this bound into (D.2) gives (D.3).

(c) Let $s$ and $t$ be as in (b).
The absolute convergence and the bound $|Q_\Psi(\omega)|\le\mathrm{cf}_{n,t}\cdot\|\omega\|_{-s}^n\cdot q_{2n+2nt}(\Psi)$ are (D.3).
Each $c_{\mathrm{kk}}$ is linear in $\Psi$; hence $Q_\Psi(\omega)$ is linear in $\Psi$, and the bound expresses its continuity.
For $\mathrm{dn}=1,2,\dots$ the sums over the finite sets $\{\mathrm{kk}:\lambda_{\kappa_i}\le\mathrm{dn}\text{ for all }i\}$ are polynomials in the functions $X(h_\kappa)$, hence $\Sigma$-measurable, and they converge to $Q_\Psi(\omega)$ at every $\omega$.
\end{proof}

\begin{named}{Lemma D.2}[Rules; (i)--(iii), (v), (vi) are the rules of Proposition~45.2\,(6)]
For every $\omega\in\mathcal S'$ the following hold.
\begin{itemize}
\item[(i)] $Q_f(\omega)=\omega(f)$ for $f\in\mathrm{Sw}_1=\mathcal S+i\mathcal S$.
\item[(ii)] $Q_{\Psi\otimes\Phi}(\omega)=Q_\Psi(\omega)\cdot Q_\Phi(\omega)$ for $\Psi\in\mathrm{Sw}_m$, $\Phi\in\mathrm{Sw}_n$, $m,n\ge0$.
Consequently $Q_{f_1\otimes\dots\otimes f_n}(\omega)=\prod_i\omega(f_i)$ for $f_1,\dots,f_n\in\mathcal S+i\mathcal S$.
\item[(iii)] $\overline{Q_\Psi(\omega)}=Q_{\overline{\Psi}}(\omega)$.
\item[(iv)] A continuous linear map from $\mathrm{Sw}_n$ to a normed space that vanishes at all $h_{\mathrm{kk}}$, $\mathrm{kk}\in\mathrm{Kh}^n$, vanishes identically.
\item[(v)] $Q_{(\mathrm{Gi}.)^{\otimes n}\Psi}(\omega)=Q_\Psi(\mathrm{Gi}^{-1}\omega)$ for every Euclidean map $\mathrm{Gi}$; that is, $Q_{(\mathrm{Gi}.)^{\otimes n}\Psi}=\widehat{\mathrm{Gi}}\,Q_\Psi$ as functions on $\mathcal S'$.
Moreover $Q_\Psi(-\omega)=(-1)^n\,Q_\Psi(\omega)$.
\item[(vi)] $\overline{\hat\theta_j\,Q_\Psi}=Q_{\Theta^{e_j}_n\Psi}$ as functions on $\mathcal S'$, for $j=1,2,3$.
\end{itemize}
\end{named}

\begin{proof}
(iv) By linearity the map vanishes on $\mathrm{Fin}_n$, and $\mathrm{Fin}_n$ is dense in $\mathrm{Sw}_n$ by (He2).

(i) By Lemma~D.1\,(c), $f\mapsto Q_f(\omega)$ is a continuous linear functional on $\mathcal S+i\mathcal S$.
So is $f\mapsto\omega(f)$: $\omega$ is continuous on $\mathcal S$, and $q_m(\mathrm{Re}\,f)\le q_m(f)$, $q_m(\mathrm{Im}\,f)\le q_m(f)$.
As $(h_\kappa)$ is orthonormal, $c_\kappa(h_{\kappa'})=1$ for $\kappa=\kappa'$ and $c_\kappa(h_{\kappa'})=0$ otherwise; hence $Q_{h_{\kappa'}}(\omega)=\omega(h_{\kappa'})$.
By (iv) the two functionals agree.

(ii) Let $m,n\ge1$.
Then $\Psi\otimes\Phi$ lies in $\mathrm{Sw}_{m+n}$ by (Nm6).
One has $\mathrm{Kh}^{m+n}=\mathrm{Kh}^m\times\mathrm{Kh}^n$ and $h_{(\mathrm{kk},\mathrm{kk}')}=h_{\mathrm{kk}}\otimes h_{\mathrm{kk}'}$, and by Fubini's theorem $c_{(\mathrm{kk},\mathrm{kk}')}(\Psi\otimes\Phi)=c_{\mathrm{kk}}(\Psi)\cdot c_{\mathrm{kk}'}(\Phi)$.
The product of the two absolutely convergent series $Q_\Psi(\omega)$ and $Q_\Phi(\omega)$ is the absolutely convergent series, over the pairs $(\mathrm{kk},\mathrm{kk}')$, of the products of their terms; this series is $Q_{\Psi\otimes\Phi}(\omega)$.
For $m=0$ or $n=0$ the assertion follows from the conventions for $n=0$ ((n8) and B.1) together with the linearity of $Q$.
The second sentence of (ii) follows from (i) by induction on $n$.

(iii) As $h_{\mathrm{kk}}$ is real, $\overline{c_{\mathrm{kk}}(\Psi)}=c_{\mathrm{kk}}(\overline{\Psi})$; the numbers $\omega(h_\kappa)$ are real.
For $n=0$ both sides equal $\overline{\Psi_0}$, because $Q_{\Psi_0}=\Psi_0$ and $Q_{\overline{\Psi_0}}=\overline{\Psi_0}$ by (n8).

(v) By (Nm4) and Lemma~D.1\,(c), $\Psi\mapsto Q_{(\mathrm{Gi}.)^{\otimes n}\Psi}(\omega)$ and $\Psi\mapsto Q_\Psi(\mathrm{Gi}^{-1}\omega)$ are continuous linear functionals on $\mathrm{Sw}_n$.
Since $(\mathrm{Gi}.)^{\otimes n}h_{\mathrm{kk}}=(\mathrm{Gi}.h_{\kappa_1})\otimes\dots\otimes(\mathrm{Gi}.h_{\kappa_n})$, (ii) and (Eu2) show that the first functional takes at $h_{\mathrm{kk}}$ the value
\[
\prod_i\omega(\mathrm{Gi}.h_{\kappa_i})=\prod_i(\mathrm{Gi}^{-1}\omega)(h_{\kappa_i}),
\]
which is the value of the second one.
By (iv) the two functionals agree.
Each term of the series $Q_\Psi(\omega)$ contains $n$ factors $\omega(h_{\kappa_i})$; this gives the second formula.
For $n=0$ both sides of each of the two formulas equal $\Psi_0$, because $Q_{\Psi_0}=\Psi_0$ and $(\mathrm{Gi}.)^{\otimes0}\Psi_0=\Psi_0$ ((n8) and B.1).

(vi) One has $\theta_j^{-1}=\theta_j=r_{e_j,0}$.
By (v) and (iii),
\[
\overline{\hat\theta_j\,Q_\Psi}=\overline{Q_{(\theta_j.)^{\otimes n}\Psi}}=Q_{\overline{(\theta_j.)^{\otimes n}\Psi}},
\]
and $\overline{(r_{e_j,0}.)^{\otimes n}\Psi}$ is $\Theta^{e_j}_n\Psi$ by (n5).
For $n=0$ both sides equal $\overline{\Psi_0}$.
\end{proof}

\subsubsection*{Proof}

Throughout the proofs below $\mathrm{lw}$ is a law with (L1)--(L4), and $X$ is the map of Proposition~45.1\,(1).
As in Part~C, the conclusions (1)--(7) of \cite[Proposition~7.10]{QEDI} are applied to the probability space $(\mathcal S',\Sigma,\mathrm{lw})$, the space $L^2(\mathbb R^3;\Lambda^2)$ with the norm $\|\cdot\|_2$, and the map $X$; by Proposition~45.1\,(1) they hold in this setting.

\begin{proof}[Proof of Proposition~45.2\,(1)]
One has $S_0=1$, the empty product.
For $f_1,\dots,f_n\in L^2(\mathbb R^3;\Lambda^2)\otimes\mathbb C$ the product $X(f_1)\cdots X(f_n)$ is $\mathrm{lw}$-integrable, and by \cite[Proposition~7.10\,(4)]{QEDI}
\[
\mathrm{lw}\left(\prod_i|X(f_i)|\right)\le
\begin{cases}
\sigma_n\cdot\prod_i\|f_i\|_2&\text{for real }f_i,\\
2^{\frac n2}\cdot\sigma_n\cdot\prod_i\|f_i\|_2&\text{for complex }f_i.
\end{cases}
\]
Since $\sigma_n\le2\left(2^n\cdot n!\right)^{\frac12}$ and $2^{\frac n2}\cdot2\left(2^n\cdot n!\right)^{\frac12}=2\left(4^n\cdot n!\right)^{\frac12}$, these are the two bounds of (1).
As $X$ is complex linear, $S_n$ is multilinear; as the product of functions is commutative, $S_n$ is symmetric; and as $X(f)$ is a class of real functions for real $f$ (Proposition~45.1\,(1)), $S_n$ is real on real tests.
Finally $S_1(f)=\mathrm{lw}(X(f))=0$ for real $f$ by \cite[Proposition~7.10\,(2)]{QEDI}, and for complex $f$ by linearity.
\end{proof}

\begin{proof}[Proof of Proposition~45.2\,(2)]
\emph{Step 1 (the functional $Q_\Psi(\omega)$).}
The absolute convergence for every $\omega$, and the continuity and linearity in $\Psi$, are Lemma~D.1\,(c).
The value $\prod_i\omega(f_i)$ at $f_1\otimes\dots\otimes f_n$, $f_i\in\mathcal S+i\mathcal S$, is Lemma~D.2\,(ii).
If a continuous linear functional on $\mathrm{Sw}_n$ takes these values at all products of real $f_i\in\mathcal S$, then its difference with $\Psi\mapsto Q_\Psi(\omega)$ vanishes at all $h_{\mathrm{kk}}$, and hence vanishes identically by Lemma~D.2\,(iv).

\emph{Step 2 (norms).}
The sum $\sum_{\mathrm{kk}}|c_{\mathrm{kk}}(\Psi)|$ is finite by \cite[Lemma~7.16]{QEDI}.
Write $\mathrm{kk}'=(\kappa_1',\dots,\kappa_n')$.
By Tonelli's theorem and \cite[Proposition~7.10\,(4)]{QEDI}, applied to the $2n$ real tests $h_{\kappa_1},\dots,h_{\kappa_n},h_{\kappa_1'},\dots,h_{\kappa_n'}$, each of norm $1$,
\begin{align*}
\mathrm{lw}\left(|Q_\Psi|^2\right)
&\le\mathrm{lw}\left(\left(\sum_{\mathrm{kk}}\left|c_{\mathrm{kk}}(\Psi)\right|\cdot\prod_i\left|X(h_{\kappa_i})\right|\right)^2\right)\\
&=\sum_{\mathrm{kk},\mathrm{kk}'}\left|c_{\mathrm{kk}}(\Psi)\right|\cdot\left|c_{\mathrm{kk}'}(\Psi)\right|\cdot\mathrm{lw}\left(\prod_i\left|X(h_{\kappa_i})\right|\cdot\prod_i\left|X(h_{\kappa_i'})\right|\right)
\le\sigma_{2n}\cdot\left(\sum_{\mathrm{kk}}\left|c_{\mathrm{kk}}(\Psi)\right|\right)^2 .
\end{align*}
By Tonelli's theorem and \cite[Proposition~7.10\,(4)]{QEDI} for the $n$ tests $h_{\kappa_1},\dots,h_{\kappa_n}$,
\[
\mathrm{lw}\left(|Q_\Psi|\right)\le\sum_{\mathrm{kk}}\left|c_{\mathrm{kk}}(\Psi)\right|\cdot\mathrm{lw}\left(\prod_i\left|X(h_{\kappa_i})\right|\right)\le\sigma_n\cdot\sum_{\mathrm{kk}}\left|c_{\mathrm{kk}}(\Psi)\right|.
\]
Together with \cite[Lemma~7.16]{QEDI} this gives
\begin{equation}\tag{D.4}
\left\|Q_\Psi\right\|_{L^2(\mathrm{lw})}\le\sigma_{2n}^{\frac12}\cdot\sum_{\mathrm{kk}}\left|c_{\mathrm{kk}}(\Psi)\right|\le\sigma_{2n}^{\frac12}\cdot4^{-n}\cdot\left(34\pi\sqrt3\right)^n\cdot q_{6n}(\Psi).
\end{equation}
Hence $\Psi\mapsto(\text{class of }Q_\Psi)$ is a continuous linear map from $\mathrm{Sw}_n$ to $L^2(\mathrm{lw})$.

\emph{Step 3 (the bound on $S_n$).}
The functional $S_n(\Psi):=\mathrm{lw}(Q_\Psi)$ is linear in $\Psi$, and
\[
|S_n(\Psi)|\le\sigma_n\cdot\sum_{\mathrm{kk}}\left|c_{\mathrm{kk}}(\Psi)\right|\le2\left(2^n\cdot n!\right)^{\frac12}\cdot4^{-n}\cdot\left(34\pi\sqrt3\right)^n\cdot q_{6n}(\Psi)=2\,(n!)^{\frac12}\cdot\left(\frac{\sqrt2\cdot34\pi\sqrt3}{4}\right)^n\cdot q_{6n}(\Psi).
\]
Here
\[
\left(\frac{\sqrt2\cdot34\pi\sqrt3}{4}\right)^2=\frac{6\cdot34^2\cdot\pi^2}{16}=433.5\,\pi^2<433.5\cdot9.87<4279<4356=66^2,
\]
because $\pi<3.1416$ and hence $\pi^2<9.87$.
Therefore $|S_n(\Psi)|\le2\cdot66^n\cdot(n!)^{\frac12}\cdot q_{6n}(\Psi)$; in particular $S_n$ is continuous on $\mathrm{Sw}_n$.

\emph{Step 4 (values, uniqueness).}
By Lemma~D.2\,(ii), $S_n(f_1\otimes\dots\otimes f_n)=\mathrm{lw}\left(\prod_iX(f_i)\right)=S_n(f_1,\dots,f_n)$ for $f_i\in\mathcal S+i\mathcal S$.
A continuous linear functional on $\mathrm{Sw}_n$ with these values at the products of real $f_i\in\mathcal S$ agrees with $S_n$ at all $h_{\mathrm{kk}}$, and hence everywhere by Lemma~D.2\,(iv).

\emph{Step 5 ($X^{(n)}$; the last sentence of (2)).}
By (n5), $X^{(n)}$ is the continuous linear map $\mathrm{Sw}_n\to L^2(\mathrm{lw})$ with $X^{(n)}(f_1\otimes\dots\otimes f_n)=X(f_1)\cdots X(f_n)$.
By Step~2 and Lemma~D.2\,(ii), $\Psi\mapsto(\text{class of }Q_\Psi)$ is a continuous linear map with the same values.
Two continuous linear maps $\mathrm{Sw}_n\to L^2(\mathrm{lw})$ with these values agree at all $h_{\mathrm{kk}}$, and hence everywhere by Lemma~D.2\,(iv).
So there is exactly one such map, and $X^{(n)}(\Psi)=Q_\Psi$ in $L^2(\mathrm{lw})$ for every $\Psi\in\mathrm{Sw}_n$.

\emph{Step 6 (odd $n$; the last sentence of (6)).}
Let $\Phi_{\mathrm{lw}}$ be even and $n$ odd.
Then $\mathrm{sgn}_*\mathrm{lw}=\mathrm{lw}$ by Proposition~45.1\,(9), and $Q_\Psi\circ\mathrm{sgn}=-Q_\Psi$ by Lemma~D.2\,(v).
Hence $S_n(\Psi)=\mathrm{lw}(Q_\Psi\circ\mathrm{sgn})=-S_n(\Psi)$, that is, $S_n(\Psi)=0$ for all $\Psi\in\mathrm{Sw}_n$.
\end{proof}

\begin{proof}[Proof of Proposition~45.2\,(4)]
Let $\mathrm{Gi}=\tau_b$ with $b\in\mathbb R^3$, or $\mathrm{Gi}=r\in O_{\mathrm{cub}}$.
Then $\mathrm{Gi}_*\mathrm{lw}=\mathrm{lw}$ by Proposition~45.1\,(5), (6).
By Lemma~D.2\,(v) and (Eu5),
\[
S_n\left((\mathrm{Gi}.)^{\otimes n}\Psi\right)=\mathrm{lw}\left(\widehat{\mathrm{Gi}}\,Q_\Psi\right)=\mathrm{lw}(Q_\Psi)=S_n(\Psi).
\]
Let $\mathrm{gs}\in\mathcal S(\mathbb R^3)$ be real and $\Psi'\in\mathrm{Sw}_{n-1}$.
By Lemma~D.2\,(i), (ii),
\[
Q_{d^{*}\mathrm{gs}\otimes\Psi'}(\omega)=\omega(d^{*}\mathrm{gs})\cdot Q_{\Psi'}(\omega)\qquad\text{for every }\omega\in\mathcal S'.
\]
The right side vanishes for $\omega\in\mathrm{Ecl}$, and $\mathrm{lw}(\mathrm{Ecl})=1$ by Proposition~45.1\,(4).
Hence $S_n(d^{*}\mathrm{gs}\otimes\Psi')=\mathrm{lw}\left(Q_{d^{*}\mathrm{gs}\otimes\Psi'}\right)=0$.
For $\mathrm{gs}=\mathrm{gs}'+i\,\mathrm{gs}''$ with real $\mathrm{gs}',\mathrm{gs}''$ one has $d^{*}\mathrm{gs}=d^{*}\mathrm{gs}'+i\,d^{*}\mathrm{gs}''$ (Section~0), and $S_n$ is linear.
\end{proof}

Fix $\mathrm{stp}\in C^\infty(\mathbb R)$ with $0\le\mathrm{stp}\le1$, $\mathrm{stp}=0$ on $(-\infty,\frac12]$ and $\mathrm{stp}=1$ on $[1,\infty)$.
For $\epsilon>0$ and the direction $j$ put
\[
\mathrm{stp}_\epsilon(y):=\mathrm{stp}\left(\frac{y_j}{\epsilon}\right)\quad(y\in\mathbb R^3),\qquad
\mathrm{stp}_\epsilon^{\otimes n}(y_1,\dots,y_n):=\prod_a\mathrm{stp}_\epsilon(y_a).
\]

\begin{named}{Lemma D.3}[Support; this is Proposition~45.2\,(7)]
Let $j\in\{1,2,3\}$ and $n\ge1$, and let $\Psi\in\mathrm{Sw}_n$ be supported in $\{(y_a)_j\ge0\text{ for all }a\}$.
Then the class of $Q_\Psi$ lies in $L^2(\Sigma_{\{y_j>0\}},\mathrm{lw})$.
\end{named}

\begin{proof}
Let $0<\epsilon\le1$ and $\Psi_\epsilon:=(\tau_{\epsilon e_j}.)^{\otimes n}\Psi$, that is, $\Psi_\epsilon(y_1,\dots,y_n)=\Psi(y_1-\epsilon e_j,\dots,y_n-\epsilon e_j)$.
It is supported in $\{(y_a)_j\ge\epsilon\text{ for all }a\}$, where $\mathrm{stp}_\epsilon^{\otimes n}=1$; hence $\mathrm{stp}_\epsilon^{\otimes n}\Psi_\epsilon=\Psi_\epsilon$.
By (He2) there are $\Psi_{\epsilon,\mathrm{dn}}\in\mathrm{Fin}_n$, $\mathrm{dn}=1,2,\dots$, with $\Psi_{\epsilon,\mathrm{dn}}\to\Psi_\epsilon$ in $\mathrm{Sw}_n$.
The function $\mathrm{stp}_\epsilon^{\otimes n}$ is smooth and bounded together with all its derivatives; by (Nm3),
\[
\mathrm{stp}_\epsilon^{\otimes n}\Psi_{\epsilon,\mathrm{dn}}\to\mathrm{stp}_\epsilon^{\otimes n}\Psi_\epsilon=\Psi_\epsilon\quad\text{in }\mathrm{Sw}_n\text{ as }\mathrm{dn}\to\infty .
\]
One has $\mathrm{stp}_\epsilon^{\otimes n}h_{\mathrm{kk}}=(\mathrm{stp}_\epsilon h_{\kappa_1})\otimes\dots\otimes(\mathrm{stp}_\epsilon h_{\kappa_n})$, and $\mathrm{stp}_\epsilon h_\kappa$ is a real element of $\mathcal S$ that vanishes on $\{y_j\le\epsilon/2\}$.
By Lemma~D.2\,(ii), $Q_{\mathrm{stp}_\epsilon^{\otimes n}\Psi_{\epsilon,\mathrm{dn}}}$ is a polynomial with complex coefficients in finitely many of the functions $X(\mathrm{stp}_\epsilon h_\kappa)$.
Each of these functions is $\Sigma_{\{y_j>0\}}$-measurable.
Indeed, let $f=\mathrm{stp}_\epsilon h_\kappa$.
By Lemma~B.2\,(a), for $\mathrm{dn}'=1,2,\dots$ the real form $\mathrm{bmp}(\cdot/\mathrm{dn}')f$ is smooth, and its support is compact and contained in $\operatorname{supp}f$, hence in $\{y_j\ge\epsilon/2\}$; so $\mathrm{bmp}(\cdot/\mathrm{dn}')f$ lies in $C_c^\infty(\{y_j>0\};\Lambda^2)$, and $\mathrm{bmp}(\cdot/\mathrm{dn}')f\to f$ in $\mathcal S$ as $\mathrm{dn}'\to\infty$.
Hence $\omega(f)=\lim_{\mathrm{dn}'\to\infty}\omega\left(\mathrm{bmp}(\cdot/\mathrm{dn}')f\right)$ for every $\omega\in\mathcal S'$, and $X(f)$ is a pointwise limit of $\Sigma_{\{y_j>0\}}$-measurable functions.
So the polynomial is $\Sigma_{\{y_j>0\}}$-measurable; it lies in $L^2(\mathrm{lw})$ by (D.4), and its class lies in $L^2(\Sigma_{\{y_j>0\}},\mathrm{lw})$.
This subspace of $L^2(\mathrm{lw})$ is closed, and by (D.4) the map $\Psi\mapsto Q_\Psi$ is continuous from $\mathrm{Sw}_n$ to $L^2(\mathrm{lw})$.
Hence the class of $Q_{\Psi_\epsilon}$ lies in this subspace.
By (Nm5), $\Psi_\epsilon\to\Psi$ in $\mathrm{Sw}_n$ as $\epsilon\to0$; hence the class of $Q_\Psi$ lies in it as well.
\end{proof}

\begin{proof}[Proof of Proposition~45.2\,(5)]
Let $\Psi_0\in\mathbb C$, and for $1\le n\le\mathrm{dn}$ let $\Psi_n\in\mathrm{Sw}_n$ be supported in $\{(y_a)_j\ge0\text{ for all }a\}$.
Put $\mathrm{Ga}:=\sum_{n=0}^{\mathrm{dn}}Q_{\Psi_n}$, where $Q_{\Psi_0}=\Psi_0$.
By Lemma~D.3, $\mathrm{Ga}$ lies in $L^2(\Sigma_{\{y_j>0\}},\mathrm{lw})$.
Proposition~45.1\,(7) with $\mathrm{po}=0$, where $r_{e_j,0}=\theta_j$, Lemma~D.2\,(vi), (ii) and the definition of $S_{m+n}$ give
\begin{align*}
0\le\langle\hat\theta_j\mathrm{Ga},\mathrm{Ga}\rangle_{\mathrm{lw}}
&=\sum_{m,n=0}^{\mathrm{dn}}\mathrm{lw}\left(\overline{\hat\theta_j\,Q_{\Psi_m}}\cdot Q_{\Psi_n}\right)
=\sum_{m,n=0}^{\mathrm{dn}}\mathrm{lw}\left(Q_{\Theta^{e_j}_m\Psi_m}\cdot Q_{\Psi_n}\right)\\
&=\sum_{m,n=0}^{\mathrm{dn}}\mathrm{lw}\left(Q_{\Theta^{e_j}_m\Psi_m\otimes\Psi_n}\right)
=\sum_{m,n=0}^{\mathrm{dn}}S_{m+n}\left(\Theta^{e_j}_m\Psi_m\otimes\Psi_n\right).
\end{align*}
\end{proof}

\begin{proof}[Proof of Proposition~45.2\,(6)]
The six rules are Lemma~D.2\,(i), (ii), (iii), (v) (two formulas) and (vi).
The last sentence is Step~6 of the proof of (2).
\end{proof}

\begin{proof}[Proof of Proposition~45.2\,(7)]
This is Lemma~D.3.
\end{proof}

\begin{proof}[Proof of Proposition~45.2\,(8)]
Let $\mathrm{lw}$ be an arbitrary law.
The proofs above use the following properties of $\mathrm{lw}$.
\begin{itemize}
\item[(1), (2):] (L1).
These proofs use Proposition~45.1\,(1), namely the reality of $X(f)$ for real $f$, and the conclusions (2) and (4) of \cite[Proposition~7.10]{QEDI}; by Proposition~45.1\,(11) these hold under (L1).
\item[(4):] the properties that Proposition~45.1\,(4), (5), (6) use.
The proof of (4) also uses clause (2), that is (L1), which Proposition~45.1\,(4), (5) use as well.
\item[(7):] (L1).
Indeed (7) is Lemma~D.3, and of the law $\mathrm{lw}$ its proof uses only (D.4), that is, clause (2).
\item[(5):] the properties that Proposition~45.1\,(7) uses (Proposition~45.1\,(11)).
The proof of (5) also uses Lemma~D.3 and clause (2), that is (L1), which Proposition~45.1\,(7) uses as well.
\item[(6):] no law enters the Lemmas~D.1 and~D.2; Step~6 of the proof of (2) uses clause (2), Lemma~D.2\,(v) and Proposition~45.1\,(9), which holds under (L1) and $\Phi_{\mathrm{lw}}(-f)=\Phi_{\mathrm{lw}}(f)$ (Proposition~45.1\,(11)).
\end{itemize}
\end{proof}

\subsection*{E. Proposition 45.3}

\subsubsection*{Proof}

Throughout, $\mathrm{lw}$ is a law with (L1)--(L4) and $j\in\{1,2,3\}$; the space $\mathrm{Ep}_j$, the form $\mathrm{bf}$, the space $\mathcal{H}_j$, the map $[\cdot]_j$ and the vector $\Omega_j$ are those of Proposition~45.3.
Norms and scalar products without an index are those of $\mathcal{H}_j$; those of $L^2(\mathrm{lw})$ carry the index.

\paragraph{E.0. The form $\mathrm{bf}$ and the space $\mathcal{H}_j$.}
Since $\theta_j\in O_{\mathrm{cub}}$, Proposition~45.1\,(6) and (Eu5) show that $\hat\theta_j$ is unitary on $L^2(\mathrm{lw})$.
By (Eu3), $\hat\theta_j\hat\theta_j=\widehat{\theta_j\theta_j}=1$; hence $\hat\theta_j$ is self-adjoint.
The form $\mathrm{bf}$ is antilinear in the first and linear in the second argument, and for $\mathrm{Ga},\mathrm{Gb}\in\mathrm{Ep}_j$
\begin{equation}\tag{E.1}
\overline{\mathrm{bf}(\mathrm{Gb},\mathrm{Ga})}=\overline{\langle\hat\theta_j\mathrm{Gb},\mathrm{Ga}\rangle_{\mathrm{lw}}}=\langle\mathrm{Ga},\hat\theta_j\mathrm{Gb}\rangle_{\mathrm{lw}}=\langle\hat\theta_j\mathrm{Ga},\mathrm{Gb}\rangle_{\mathrm{lw}}=\mathrm{bf}(\mathrm{Ga},\mathrm{Gb}),
\end{equation}
\begin{equation}\tag{E.2}
\mathrm{bf}(\mathrm{Ga},\mathrm{Ga})\ge0,
\end{equation}
\begin{equation}\tag{E.3}
\left|\mathrm{bf}(\mathrm{Ga},\mathrm{Gb})\right|\le\left\|\hat\theta_j\mathrm{Ga}\right\|_{L^2(\mathrm{lw})}\cdot\|\mathrm{Gb}\|_{L^2(\mathrm{lw})}=\|\mathrm{Ga}\|_{L^2(\mathrm{lw})}\cdot\|\mathrm{Gb}\|_{L^2(\mathrm{lw})};
\end{equation}
here (E.2) is Proposition~45.1\,(7) with $\mathrm{po}=0$, since $r_{e_j,0}=\theta_j$.
So $\mathrm{bf}$ is a positive semidefinite Hermitian form on $\mathrm{Ep}_j$, and the Cauchy--Schwarz inequality holds for it:
\begin{equation}\tag{E.4}
\left|\mathrm{bf}(\mathrm{Ga},\mathrm{Gb})\right|^2\le\mathrm{bf}(\mathrm{Ga},\mathrm{Ga})\cdot\mathrm{bf}(\mathrm{Gb},\mathrm{Gb}).
\end{equation}
The function $\mathrm{Nrm}(\mathrm{Ga}):=\mathrm{bf}(\mathrm{Ga},\mathrm{Ga})^{\frac12}$ is a seminorm on $\mathrm{Ep}_j$, and $\mathrm{Nq}:=\{\mathrm{Ga}\in\mathrm{Ep}_j:\mathrm{Nrm}(\mathrm{Ga})=0\}$ is a linear subspace; by (E.4), $\mathrm{bf}(\mathrm{Ga},\mathrm{Gb})$ does not change when elements of $\mathrm{Nq}$ are added to $\mathrm{Ga}$ and to $\mathrm{Gb}$.
Hence $\langle[\mathrm{Ga}]_j,[\mathrm{Gb}]_j\rangle:=\mathrm{bf}(\mathrm{Ga},\mathrm{Gb})$, as in (n8), is a scalar product on $\mathrm{Ep}_j/\mathrm{Nq}$.
The space $\mathcal{H}_j$ is the completion, the map $[\cdot]_j\colon\mathrm{Ep}_j\to\mathcal{H}_j$ is linear with dense range, and by (E.3)
\begin{equation}\tag{E.5}
\left\|[\mathrm{Ga}]_j\right\|=\mathrm{Nrm}(\mathrm{Ga})\le\|\mathrm{Ga}\|_{L^2(\mathrm{lw})}.
\end{equation}

\emph{$\mathcal{H}_j$ is separable.}
Let $\mathrm{Srt}$ be the countable set of the linear combinations of the $h_\kappa$ with rational coefficients.
By (He2), and because real coefficients are limits of rational ones, $\mathrm{Srt}$ is dense in $\mathcal S$ for the norms $q_m$, and hence for the norm $\|\cdot\|_2\le\pi\sqrt3\cdot q_2$ (Nm2).
For $f,f'\in\mathcal S$, (C.1) gives
\[
\left\|\exp(iX(f))-\exp(iX(f'))\right\|_{L^2(\mathrm{lw})}\le\left\|X(f-f')\right\|_{L^2(\mathrm{lw})}\le\|f-f'\|_2 .
\]
Hence the closed linear span of $\{\exp(iX(f)):f\in\mathrm{Srt}\}$ contains every $\exp(iX(f))$, $f\in\mathcal S$, and therefore equals $L^2(\mathrm{lw})$ by \cite[Lemma~7.14]{QEDI}.
So $L^2(\mathrm{lw})$ is separable; its subset $\mathrm{Ep}_j$ is separable, as is every subset of a separable metric space; and by (E.5) the image under $[\cdot]_j$ of a countable dense subset of $\mathrm{Ep}_j$ is dense in $\mathcal{H}_j$.

\begin{proof}[Proof of Proposition~45.3\,(1)]
\emph{Step 1.}
For $t\ge0$ one has $\tau_{te_j}\{y_j>0\}=\{y_j>t\}\subset\{y_j>0\}$ by (Eu7).
Hence $C_c^\infty(\{y_j>t\};\Lambda^2)\subset C_c^\infty(\{y_j>0\};\Lambda^2)$ and $\Sigma_{\{y_j>t\}}\subset\Sigma_{\{y_j>0\}}$, and by Proposition~45.1\,(5) and (Eu5)
\[
U(te_j)\,\mathrm{Ep}_j=L^2(\Sigma_{\{y_j>t\}},\mathrm{lw}),
\]
a subspace of $\mathrm{Ep}_j$.

\emph{Step 2 (symmetry).}
By (Eu7) and (Eu3), $\hat\theta_j\,U(te_j)=U(-te_j)\,\hat\theta_j$; and $U(-te_j)$ is the inverse, hence the adjoint, of the unitary operator $U(te_j)$ on $L^2(\mathrm{lw})$.
For $\mathrm{Ga},\mathrm{Gb}\in\mathrm{Ep}_j$ and $t\ge0$ it follows that
\begin{equation}\tag{E.6}
\begin{split}
\mathrm{bf}(\mathrm{Ga},U(te_j)\mathrm{Gb})&=\langle\hat\theta_j\mathrm{Ga},U(te_j)\mathrm{Gb}\rangle_{\mathrm{lw}}=\langle U(-te_j)\hat\theta_j\mathrm{Ga},\mathrm{Gb}\rangle_{\mathrm{lw}}\\
&=\langle\hat\theta_jU(te_j)\mathrm{Ga},\mathrm{Gb}\rangle_{\mathrm{lw}}=\mathrm{bf}(U(te_j)\mathrm{Ga},\mathrm{Gb}).
\end{split}
\end{equation}

\emph{Step 3 (contraction).}
By (E.6), $U(te_j)U(te_j)=U(2te_j)$ and (E.4),
\[
\mathrm{Nrm}(U(te_j)\mathrm{Ga})^2=\mathrm{bf}(\mathrm{Ga},U(2te_j)\mathrm{Ga})\le\mathrm{Nrm}(\mathrm{Ga})\cdot\mathrm{Nrm}(U(2te_j)\mathrm{Ga})\qquad\text{for all }t\ge0 .
\]
It follows by induction on $a$ that, for $a=1,2,\dots$ and all $t\ge0$,
\[
\mathrm{Nrm}(U(te_j)\mathrm{Ga})\le\mathrm{Nrm}(\mathrm{Ga})^{1-2^{-a}}\cdot\mathrm{Nrm}\left(U(2^ate_j)\mathrm{Ga}\right)^{2^{-a}}.
\]
For $a=1$ this is the inequality just proved.
If it holds for $a$, then the case $a=1$ with $2^at$ in place of $t$ gives $\mathrm{Nrm}(U(2^ate_j)\mathrm{Ga})\le\mathrm{Nrm}(\mathrm{Ga})^{\frac12}\cdot\mathrm{Nrm}(U(2^{a+1}te_j)\mathrm{Ga})^{\frac12}$, and so
\[
\mathrm{Nrm}(U(te_j)\mathrm{Ga})\le\mathrm{Nrm}(\mathrm{Ga})^{1-2^{-a}+2^{-a-1}}\cdot\mathrm{Nrm}\left(U(2^{a+1}te_j)\mathrm{Ga}\right)^{2^{-a-1}},
\]
which is the inequality for $a+1$.
By (E.5) and the unitarity of $U(2^ate_j)$ on $L^2(\mathrm{lw})$, $\mathrm{Nrm}(U(2^ate_j)\mathrm{Ga})\le\|\mathrm{Ga}\|_{L^2(\mathrm{lw})}$.
Hence $\mathrm{Nrm}(U(te_j)\mathrm{Ga})\le\mathrm{Nrm}(\mathrm{Ga})^{1-2^{-a}}\cdot\|\mathrm{Ga}\|_{L^2(\mathrm{lw})}^{2^{-a}}$ for all $a$.
If $\mathrm{Nrm}(\mathrm{Ga})=0$, the right side is $0$.
Otherwise $\mathrm{Ga}\ne0$, and the right side tends to $\mathrm{Nrm}(\mathrm{Ga})$ as $a\to\infty$.
In both cases
\begin{equation}\tag{E.7}
\mathrm{Nrm}(U(te_j)\mathrm{Ga})\le\mathrm{Nrm}(\mathrm{Ga})\qquad(\mathrm{Ga}\in\mathrm{Ep}_j,\ t\ge0).
\end{equation}

\emph{Step 4 (the operators).}
By (E.7), $U(te_j)$ maps $\mathrm{Nq}$ into $\mathrm{Nq}$; hence $[\mathrm{Ga}]_j\mapsto[U(te_j)\mathrm{Ga}]_j$ is a well defined linear map of the dense subspace $[\mathrm{Ep}_j]_j$ into itself, of norm at most $1$.
It has exactly one extension to a bounded operator $\mathrm{Tq}(t)$ on $\mathcal{H}_j$, and $\|\mathrm{Tq}(t)\|\le1$.
By (E.6), $\langle[\mathrm{Ga}]_j,\mathrm{Tq}(t)[\mathrm{Gb}]_j\rangle=\langle\mathrm{Tq}(t)[\mathrm{Ga}]_j,[\mathrm{Gb}]_j\rangle$; by density and continuity $\mathrm{Tq}(t)$ is symmetric, hence self-adjoint.
On classes, $\mathrm{Tq}(t)\mathrm{Tq}(t')[\mathrm{Ga}]_j=[U(te_j)U(t'e_j)\mathrm{Ga}]_j=\mathrm{Tq}(t+t')[\mathrm{Ga}]_j$ and $\mathrm{Tq}(0)[\mathrm{Ga}]_j=[\mathrm{Ga}]_j$; by density, $\mathrm{Tq}(t)\mathrm{Tq}(t')=\mathrm{Tq}(t+t')$ and $\mathrm{Tq}(0)=1$.

\emph{Step 5 (strong continuity).}
By (E.5) and Proposition~45.1\,(5),
\[
\left\|\mathrm{Tq}(t)[\mathrm{Ga}]_j-[\mathrm{Ga}]_j\right\|\le\left\|U(te_j)\mathrm{Ga}-\mathrm{Ga}\right\|_{L^2(\mathrm{lw})}\to0\qquad\text{as }t\to0 .
\]
For $\mathrm{zt}\in\mathcal{H}_j$ and $\epsilon>0$ choose $\mathrm{Ga}\in\mathrm{Ep}_j$ with $\|\mathrm{zt}-[\mathrm{Ga}]_j\|<\epsilon$; then $\|\mathrm{Tq}(t)\mathrm{zt}-\mathrm{zt}\|\le2\epsilon+\|\mathrm{Tq}(t)[\mathrm{Ga}]_j-[\mathrm{Ga}]_j\|$.
Hence $\mathrm{Tq}(t)\mathrm{zt}\to\mathrm{zt}$ as $t\to0$.
For $t,t_0\ge0$,
\[
\left\|\mathrm{Tq}(t)\mathrm{zt}-\mathrm{Tq}(t_0)\mathrm{zt}\right\|=\left\|\mathrm{Tq}(\min(t,t_0))\left(\mathrm{Tq}(|t-t_0|)\mathrm{zt}-\mathrm{zt}\right)\right\|\le\left\|\mathrm{Tq}(|t-t_0|)\mathrm{zt}-\mathrm{zt}\right\|.
\]

\emph{Step 6 (positivity, injectivity).}
One has $\mathrm{Tq}(t)=\mathrm{Tq}(t/2)^2$ with $\mathrm{Tq}(t/2)$ self-adjoint; hence $\mathrm{Tq}(t)\ge0$.
Let $\mathrm{Tq}(t)\mathrm{zt}=0$ for some $t>0$.
Then $\|\mathrm{Tq}(t/2)\mathrm{zt}\|^2=\langle\mathrm{zt},\mathrm{Tq}(t)\mathrm{zt}\rangle=0$; by induction $\mathrm{Tq}(2^{-a}t)\mathrm{zt}=0$ for all $a$, and $\mathrm{zt}=\lim_{a\to\infty}\mathrm{Tq}(2^{-a}t)\mathrm{zt}=0$ by Step~5.

\emph{Step 7 (the generator).}
The operator $\mathrm{Tq}(1)$ is bounded and self-adjoint on the separable Hilbert space $\mathcal{H}_j$.
By the spectral theorem \cite[Theorem~VII.3 and its Corollary]{ReedSimon1980} there are a finite measure space $(\mathrm{Mq},\mathrm{mq})$, a unitary operator $\mathrm{Uq}\colon\mathcal{H}_j\to L^2(\mathrm{Mq},\mathrm{mq})$ and a bounded real measurable function $\mathrm{Fq}$ on $\mathrm{Mq}$ such that $\mathrm{Uq}\,\mathrm{Tq}(1)\,\mathrm{Uq}^{-1}$ is the multiplication by $\mathrm{Fq}$.
From $0\le\langle\mathrm{zt},\mathrm{Tq}(1)\mathrm{zt}\rangle\le\|\mathrm{zt}\|^2$, that is,
\[
0\le\int\mathrm{d}\mathrm{mq}\;\mathrm{Fq}\cdot\left|\mathrm{Uq}\,\mathrm{zt}\right|^2\le\int\mathrm{d}\mathrm{mq}\;\left|\mathrm{Uq}\,\mathrm{zt}\right|^2\qquad\text{for all }\mathrm{zt}\in\mathcal{H}_j,
\]
it follows that $0\le\mathrm{Fq}\le1$ $\mathrm{mq}$-almost everywhere: take for $\mathrm{Uq}\,\mathrm{zt}$ the indicator functions of $\{\mathrm{Fq}<0\}$ and of $\{\mathrm{Fq}>1\}$.
The indicator function of $\{\mathrm{Fq}=0\}$ lies in $L^2(\mathrm{Mq},\mathrm{mq})$ and is annihilated by the multiplication by $\mathrm{Fq}$; by Step~6 it is $0$, that is, $\mathrm{Fq}>0$ $\mathrm{mq}$-almost everywhere.
Hence $-\log\mathrm{Fq}$ is a measurable function with values in $[0,\infty)$ almost everywhere.
Let
\[
\mathsf{H}_j:=\mathrm{Uq}^{-1}\left(\text{multiplication by }-\log\mathrm{Fq}\right)\mathrm{Uq},
\qquad
\text{with the domain }\left\{\mathrm{zt}:(\log\mathrm{Fq})\cdot\mathrm{Uq}\,\mathrm{zt}\in L^2(\mathrm{Mq},\mathrm{mq})\right\}.
\]
This operator is self-adjoint \cite[Section~VIII.3, Proposition~1]{ReedSimon1980}, and
\[
\langle\mathrm{zt},\mathsf{H}_j\mathrm{zt}\rangle=\int\mathrm{d}\mathrm{mq}\;(-\log\mathrm{Fq})\cdot\left|\mathrm{Uq}\,\mathrm{zt}\right|^2\ge0
\]
for $\mathrm{zt}$ in its domain.

\emph{The symbol $\exp(-tH)$.}
For a self-adjoint operator $H\ge0$ on a separable Hilbert space and $t\ge0$, $\exp(-tH)$ denotes the operator $\mathrm{phq}_t(H)$ of the functional calculus for bounded Borel functions \cite[Theorem~VIII.5]{ReedSimon1980}, where $\mathrm{phq}_t(\mathrm{rv}):=\exp(-t\max(\mathrm{rv},0))$ for real $\mathrm{rv}$, as in (n8).
If $H=\mathrm{Uq}^{-1}(\text{multiplication by }\mathrm{Fq}')\,\mathrm{Uq}$ with a real measurable function $\mathrm{Fq}'$ on $\mathrm{Mq}$ that is finite almost everywhere, then $\mathrm{phq}_t(H)=\mathrm{Uq}^{-1}(\text{multiplication by }\mathrm{phq}_t\circ\mathrm{Fq}')\,\mathrm{Uq}$: this is the definition of the functional calculus that precedes \cite[Theorem~VIII.5]{ReedSimon1980}, and by the uniqueness assertion of that theorem the operator does not depend on the representation of $H$ as a multiplication operator.
For $\mathrm{Fq}'=-\log\mathrm{Fq}\ge0$ one has $\mathrm{phq}_t\circ\mathrm{Fq}'=\exp(t\log\mathrm{Fq})=\mathrm{Fq}^t$.
Hence
\[
\exp(-t\mathsf{H}_j)=\mathrm{Uq}^{-1}\left(\text{multiplication by }\mathrm{Fq}^t\right)\mathrm{Uq}\qquad(t\ge0).
\]

\emph{Step 8 ($\mathrm{Tq}(t)=\exp(-t\mathsf{H}_j)$).}
For $a=0$ the identity $\mathrm{Tq}(2^{-a})=\exp(-2^{-a}\mathsf{H}_j)$ is the choice of $\mathrm{Fq}$, at $t=1$.
Let $\mathrm{Tq}(2^{-a})=\exp(-2^{-a}\mathsf{H}_j)$.
The operators $\mathrm{Tq}(2^{-a-1})$ and $\exp(-2^{-a-1}\mathsf{H}_j)$ are self-adjoint and $\ge0$ (Step~6, and $\mathrm{Fq}^{2^{-a-1}}\ge0$), and the square of each of them is $\mathrm{Tq}(2^{-a})=\exp(-2^{-a}\mathsf{H}_j)$.
A bounded operator $\ge0$ has only one square root that is $\ge0$ \cite[Theorem~VI.9]{ReedSimon1980}.
Hence $\mathrm{Tq}(2^{-a-1})=\exp(-2^{-a-1}\mathsf{H}_j)$.
By the two semigroup laws, $\mathrm{Tq}(t)=\exp(-t\mathsf{H}_j)$ for $t=\mathrm{dn}\,2^{-a}$ with integers $\mathrm{dn},a\ge0$.
The map $t\mapsto\mathrm{Tq}(t)\mathrm{zt}$ is continuous (Step~5), and so is $t\mapsto\exp(-t\mathsf{H}_j)\mathrm{zt}$, because
\[
\left\|\exp(-t\mathsf{H}_j)\mathrm{zt}-\exp(-t_0\mathsf{H}_j)\mathrm{zt}\right\|^2=\int\mathrm{d}\mathrm{mq}\;\left|\mathrm{Fq}^t-\mathrm{Fq}^{t_0}\right|^2\cdot\left|\mathrm{Uq}\,\mathrm{zt}\right|^2\to0\qquad\text{as }t\to t_0,
\]
by dominated convergence ($0<\mathrm{Fq}\le1$).
Hence $\mathrm{Tq}(t)=\exp(-t\mathsf{H}_j)$ for all $t\ge0$.

\emph{Step 9 ($\Omega_j$).}
The constant function $1$ lies in $\mathrm{Ep}_j$, and $\hat\theta_j\,1=1$, $U(te_j)\,1=1$.
Hence $\|\Omega_j\|^2=\mathrm{bf}(1,1)=\mathrm{lw}(1)=1$ and $\mathrm{Tq}(1)\Omega_j=[U(e_j)1]_j=\Omega_j$.
Consequently $(\mathrm{Fq}-1)\cdot\mathrm{Uq}\,\Omega_j=0$ almost everywhere: $\mathrm{Uq}\,\Omega_j$ vanishes almost everywhere outside $\{\mathrm{Fq}=1\}$, where $\log\mathrm{Fq}=0$.
So $\Omega_j$ lies in the domain of $\mathsf{H}_j$, and $\mathsf{H}_j\Omega_j=0$.
\end{proof}

\begin{proof}[Proof of Proposition~45.3\,(2)]
Let $y,y'\in e_j^{\perp}$.
By (Eu7) and (Eu3), $\hat\theta_j\,U(y)=U(y)\,\hat\theta_j$.
Since $\tau_y\{y_j>0\}=\{y_j>0\}$, Proposition~45.1\,(5) and (Eu5) give $U(y)\,\mathrm{Ep}_j=\mathrm{Ep}_j$.
As $U(y)$ is unitary on $L^2(\mathrm{lw})$, one has for $\mathrm{Ga},\mathrm{Gb}\in\mathrm{Ep}_j$
\[
\mathrm{bf}(U(y)\mathrm{Ga},U(y)\mathrm{Gb})=\langle U(y)\hat\theta_j\mathrm{Ga},U(y)\mathrm{Gb}\rangle_{\mathrm{lw}}=\mathrm{bf}(\mathrm{Ga},\mathrm{Gb}).
\]
Hence $U(y)$ maps $\mathrm{Nq}$ into $\mathrm{Nq}$, and $[\mathrm{Ga}]_j\mapsto[U(y)\mathrm{Ga}]_j$ is a well defined linear isometry of $[\mathrm{Ep}_j]_j$ onto itself; it is onto because $U(y)\,\mathrm{Ep}_j=\mathrm{Ep}_j$.
It has exactly one extension to a unitary operator $V_j(y)$ on $\mathcal{H}_j$.
On classes,
\begin{align*}
V_j(y)V_j(y')[\mathrm{Ga}]_j&=[U(y)U(y')\mathrm{Ga}]_j=V_j(y+y')[\mathrm{Ga}]_j,\\
V_j(y)\,\mathrm{Tq}(t)[\mathrm{Ga}]_j&=[U(y)U(te_j)\mathrm{Ga}]_j=[U(te_j)U(y)\mathrm{Ga}]_j=\mathrm{Tq}(t)\,V_j(y)[\mathrm{Ga}]_j,\\
V_j(y)\Omega_j&=[U(y)1]_j=\Omega_j;
\end{align*}
the second line uses Proposition~45.1\,(5): $U$ is a representation of the commutative group $\mathbb R^3$.
By density the first two identities hold on $\mathcal{H}_j$.
By (E.5) and Proposition~45.1\,(5), $\|V_j(y)[\mathrm{Ga}]_j-[\mathrm{Ga}]_j\|\le\|U(y)\mathrm{Ga}-\mathrm{Ga}\|_{L^2(\mathrm{lw})}\to0$ as $y\to0$.
Together with $\|V_j(y)\|=1$ and the density of $[\mathrm{Ep}_j]_j$, this gives $V_j(y)\mathrm{zt}\to\mathrm{zt}$ as $y\to0$ for every $\mathrm{zt}\in\mathcal{H}_j$, and $\|V_j(y)\mathrm{zt}-V_j(y_0)\mathrm{zt}\|=\|V_j(y-y_0)\mathrm{zt}-\mathrm{zt}\|$.
\end{proof}

\begin{proof}[Proof of Proposition~45.3\,(3)]
By Lemma~D.3 the classes of $Q_\Psi$ and $Q_\Phi$ lie in $\mathrm{Ep}_j$; for $m=0$ or $n=0$ they are constants.
Let $t\ge0$.
By Lemma~D.2\,(v) with $\mathrm{Gi}=\tau_{te_j}$, one has $U(te_j)Q_\Phi=Q_{\tau_{te_j}^{\otimes n}\Phi}$.
The definition of $S_{m+n}$ (Proposition~45.2\,(2)), Lemma~D.2\,(ii), (vi) and (1) give
\begin{align*}
S_{m+n}\left(\Theta^{e_j}_m\Psi\otimes\tau_{te_j}^{\otimes n}\Phi\right)
&=\mathrm{lw}\left(Q_{\Theta^{e_j}_m\Psi}\cdot Q_{\tau_{te_j}^{\otimes n}\Phi}\right)
=\mathrm{lw}\left(\overline{\hat\theta_j\,Q_\Psi}\cdot U(te_j)Q_\Phi\right)\\
&=\mathrm{bf}\left(Q_\Psi,U(te_j)Q_\Phi\right)
=\left\langle[Q_\Psi]_j,\exp(-t\mathsf{H}_j)[Q_\Phi]_j\right\rangle .
\end{align*}
In the notation of Step~7 of the proof of (1),
\[
\ker\mathsf{H}_j=\left\{\mathrm{zt}:(\log\mathrm{Fq})\cdot\mathrm{Uq}\,\mathrm{zt}=0\text{ almost everywhere}\right\}=\left\{\mathrm{zt}:\mathrm{Uq}\,\mathrm{zt}=0\text{ almost everywhere outside }\{\mathrm{Fq}=1\}\right\};
\]
hence $E_0=\mathrm{Uq}^{-1}\left(\text{multiplication by }\mathbf{1}[\mathrm{Fq}=1]\right)\mathrm{Uq}$.
For $\mathrm{zt}\in\mathcal{H}_j$,
\[
\left\|\exp(-t\mathsf{H}_j)\mathrm{zt}-E_0\mathrm{zt}\right\|^2=\int\mathrm{d}\mathrm{mq}\;\left|\mathrm{Fq}^t-\mathbf{1}[\mathrm{Fq}=1]\right|^2\cdot\left|\mathrm{Uq}\,\mathrm{zt}\right|^2\to0\qquad\text{as }t\to\infty,
\]
by dominated convergence: the integrand is at most $|\mathrm{Uq}\,\mathrm{zt}|^2$, and $\mathrm{Fq}^t\to\mathbf{1}[\mathrm{Fq}=1]$ wherever $0<\mathrm{Fq}\le1$.
With $\mathrm{zt}=[Q_\Phi]_j$ this gives the limit.
\end{proof}

\begin{proof}[Proof of Proposition~45.3\,(4)]
We use \cite[Proposition~8.19(i)--(iii)]{QEDI}.
In our notation it reads as follows.
Let $\mathrm{lw}$ be a law that is invariant under $\tau_{t\cdot e_j}$ for every real $t$ and under $\theta_j$, and that is reflection positive at $\{y_j=0\}$.
Put $\mathrm{Ep}_j:=L^2(\Sigma_{\{y_j>0\}},\mathrm{lw})$, and let $\mathrm{bf}(\mathrm{Ga},\mathrm{Gb}):=\langle\hat{\theta}_j\mathrm{Ga},\mathrm{Gb}\rangle_{\mathrm{lw}}$ for $\mathrm{Ga},\mathrm{Gb}\in\mathrm{Ep}_j$.
\begin{itemize}
\item[(i)] The form $\mathrm{bf}$ is Hermitian and positive semidefinite, and $\mathrm{bf}(\mathrm{Ga},\mathrm{Ga})\le\|\mathrm{Ga}\|_{L^2(\mathrm{lw})}^2$.
Let $\mathcal{H}_j$ be the completion of the quotient of $\mathrm{Ep}_j$ by the subspace $\{\mathrm{Ga}\in\mathrm{Ep}_j:\mathrm{bf}(\mathrm{Ga},\mathrm{Ga})=0\}$, and let $[\,\cdot\,]_j\colon\mathrm{Ep}_j\to\mathcal{H}_j$ be the canonical map; it is a linear contraction with dense range.
With $\Omega_j:=[1]_j$ one has $\|\Omega_j\|=1$ and $\langle\Omega_j,[\mathrm{Ga}]_j\rangle=\mathrm{lw}(\mathrm{Ga})$.
\item[(ii)] There is exactly one family $T(t)$, $t\ge0$, of bounded operators on $\mathcal{H}_j$ such that $T(t)[\mathrm{Ga}]_j=[U(t\cdot e_j)\mathrm{Ga}]_j$.
This family is a strongly continuous semigroup of self-adjoint contractions, and $T(t)=\exp(-t\cdot\mathsf{H}_{j})$ for exactly one self-adjoint operator $\mathsf{H}_{j}\ge0$; moreover $\mathsf{H}_{j}\Omega_j=0$.
\item[(iii)] If $y\perp e_j$ and $(\tau_y)_*\mathrm{lw}=\mathrm{lw}$, there is exactly one unitary operator $V_j(y)$ on $\mathcal{H}_j$ such that $V_j(y)[\mathrm{Ga}]_j=[U(y)\mathrm{Ga}]_j$; it commutes with every $T(t)$, and $V_j(y)\Omega_j=\Omega_j$.
If $(\tau_y)_*\mathrm{lw}=\mathrm{lw}$ for every $y\perp e_j$, then $y\mapsto V_j(y)$ is a strongly continuous unitary representation of $e_j^{\perp}$.
\end{itemize}
The law $\mathrm{lw}$ of Proposition~45.3 satisfies these hypotheses.
It is invariant under every translation by Proposition~45.1\,(5), and under $\theta_j$ by Proposition~45.1\,(6), since $\theta_j\in O_{\mathrm{cub}}$; it is reflection positive at $\{y_j=0\}$ by (L4).
The space $\mathrm{Ep}_j$, the form $\mathrm{bf}$, the space $\mathcal{H}_j$, the map $[\,\cdot\,]_j$ and the vector $\Omega_j$ of the cited Proposition are given by the same formulas as those of Proposition~45.3.
By Step~4 of the proof of (1), the operators $\mathrm{Tq}(t)$ are bounded and $\mathrm{Tq}(t)[\mathrm{Ga}]_j=[U(t\cdot e_j)\mathrm{Ga}]_j$.
The uniqueness assertion of (ii) therefore gives $\mathrm{Tq}(t)=T(t)$.
The operator $\mathsf{H}_{j}$ constructed in Step~7 is self-adjoint and non-negative, and $T(t)=\exp(-t\cdot\mathsf{H}_{j})$ by Step~8, the symbol $\exp(-tH)$ having the meaning fixed in Step~7.
By the uniqueness assertion of (ii), it is the operator $\mathsf{H}_{j}$ of the cited Proposition.
Finally, the operators $V_j(y)$ of (2) are unitary and satisfy $V_j(y)[\mathrm{Ga}]_j=[U(y)\mathrm{Ga}]_j$.
By the uniqueness assertion of (iii), they are the operators $V_j(y)$ of the cited Proposition.
\end{proof}

\begin{proof}[Proof of Proposition~45.3\,(5)]
Fix $j$.
Of the law $\mathrm{lw}$, the proofs of (1)--(4) for the direction $j$ use only the following facts:
Proposition~45.1\,(5);
the invariance $(\theta_j)_*\mathrm{lw}=\mathrm{lw}$, which E.0 and the proof of (4) obtain from Proposition~45.1\,(6), and which is also the first assertion of Proposition~45.1\,(7) for $\mathrm{po}=0$, because $r_{e_j,0}=\theta_j$;
Proposition~45.1\,(7) for $\mathrm{po}=0$;
the condition (L4) itself, which enters the proof of (4) as the reflection positivity at $\{y_j=0\}$ among the hypotheses of \cite[Proposition~8.19]{QEDI};
the bounds (C.1);
the Lemmas~D.2 and~D.3;
and Proposition~45.2\,(2).
By Proposition~45.1\,(11) and Proposition~45.2\,(8), all of them hold for every law with (L1), with (L4), with the first equality of (L3), and with the second equality of (L3) for $r=\theta_j$.
Consequently (1)--(4) hold, for this $j$, for every law with these four properties.
\end{proof}

\subsection*{F. Proposition 45.4}

\subsubsection*{A lemma on cumulants}

Let $\mathrm{Yq}_1,\dots,\mathrm{Yq}_n$ be complex random variables on a probability space with expectation $\mathrm{Eb}$, such that $\mathrm{Eb}|\mathrm{Yq}_i|^n<\infty$ for $i=1,\dots,n$.
Their joint cumulant is
\begin{equation}\tag{F.1}
\kappa_n(\mathrm{Yq}_1,\dots,\mathrm{Yq}_n):=\sum_{\mathrm{pt}}(-1)^{\#|\mathrm{pt}|-1}\cdot\left(\#|\mathrm{pt}|-1\right)!\cdot\prod_{\mathrm{bk}\in\mathrm{pt}}\mathrm{Eb}\left(\prod_{i\in\mathrm{bk}}\mathrm{Yq}_i\right),
\end{equation}
where the sum extends over all partitions $\mathrm{pt}$ of $\{1,\dots,n\}$ into non-empty blocks $\mathrm{bk}$, $\#|\mathrm{pt}|$ denotes the number of blocks of $\mathrm{pt}$, and $\#|\mathrm{bk}|$ the number of elements of $\mathrm{bk}$.
The products over the blocks are integrable.
Indeed, H\"older's inequality, applied to the $\#|\mathrm{bk}|$ factors $|\mathrm{Yq}_i|$, $i\in\mathrm{bk}$, and to $n-\#|\mathrm{bk}|$ factors $1$, all with the exponent $n$, gives
\[
\mathrm{Eb}\prod_{i\in\mathrm{bk}}|\mathrm{Yq}_i|\le\prod_{i\in\mathrm{bk}}\left(\mathrm{Eb}|\mathrm{Yq}_i|^n\right)^{\frac{1}{n}}.
\]
Now let $\mathrm{lw}$ be a law with (G), and let $f_1,\dots,f_n\in L^2(\mathbb{R}^3;\Lambda^2)\otimes\mathbb{C}$.
The number $\kappa_n(f_1,\dots,f_n)$ of (n5), (n8) is the joint cumulant $\kappa_n(X(f_1),\dots,X(f_n))$ of (F.1) under $\mathrm{lw}$; in (n5) the partitions are denoted by $\mathrm{pn}$ rather than $\mathrm{pt}$.
The moment condition of (F.1) is satisfied, since
\[
\mathrm{lw}\left(|X(f_i)|^n\right)\le2^{\frac{n}{2}}\cdot\sigma_n\cdot\|f_i\|_2^n
\]
by the conclusion (4) of \cite[Proposition~7.10]{QEDI} for the probability space $(\mathcal{S}',\Sigma,\mathrm{lw})$, the space $L^2(\mathbb{R}^3;\Lambda^2)$ with the norm $\|\cdot\|_2$ and the map $X$ of Proposition~45.1\,(1).
This conclusion is part of Proposition~45.1\,(1), and it holds for every law with (G), that is with (L1), by Proposition~45.1\,(11).

\begin{named}{Lemma F.1}
Let $\mathrm{lw}$ be a law with (L1), let $f_1,\dots,f_n\in L^2(\mathbb{R}^3;\Lambda^2)\otimes\mathbb{C}$, and put
\[
\mathrm{Zq}(z):=\mathrm{lw}\left(\exp\sum_{i}z_i\cdot X(f_i)\right)\qquad(z\in\mathbb{C}^n).
\]
Then the derivative $\partial_{z_1}\cdots\partial_{z_n}$ of $\log\mathrm{Zq}(z)$, formed with the principal branch of the logarithm, at $z=0$ is the number $\kappa_n(f_1,\dots,f_n)$ of (F.1).
\end{named}

\begin{proof}
The map $X$ is that of Proposition~45.1\,(1).
The conclusions (1)--(7) of \cite[Proposition~7.10]{QEDI} are taken for the probability space $(\mathcal{S}',\Sigma,\mathrm{lw})$, the space $L^2(\mathbb{R}^3;\Lambda^2)$ with the norm $\|\cdot\|_2$ and this map $X$; they hold by Proposition~45.1\,(1), which holds for every law with (L1) by Proposition~45.1\,(11).
For a subset $\mathrm{Ix}$ of $\{1,\dots,n\}$ put $\partial_{\mathrm{Ix}}:=\prod_{i\in\mathrm{Ix}}\partial_{z_i}$.

\emph{Step~1: $\partial_{\mathrm{Ix}}\mathrm{Zq}(0)=\mathrm{lw}\left(\prod_{i\in\mathrm{Ix}}X(f_i)\right)$.}
Write $f_i=f'_i+i\,f''_i$ with real $f'_i,f''_i$.
Then $|X(f_i)|\le|X(f'_i)|+|X(f''_i)|$.
Jensen's inequality for the exponential function, with the uniform weights on $2n$ points, gives
\[
\exp\left(\sum_i|z_i|\cdot\left(|X(f'_i)|+|X(f''_i)|\right)\right)\le\frac{1}{2n}\sum_i\left(\exp\left(2n\cdot|z_i|\cdot|X(f'_i)|\right)+\exp\left(2n\cdot|z_i|\cdot|X(f''_i)|\right)\right),
\]
and the right side is $\mathrm{lw}$-integrable by \cite[Proposition~7.10(1)]{QEDI}.
It bounds the moduli of the partial sums of the series
\[
\exp\left(\sum_iz_i\cdot X(f_i)\right)=\sum_{\mathrm{dn}\ge0}\frac{1}{\mathrm{dn}!}\left(\sum_iz_i\cdot X(f_i)\right)^{\mathrm{dn}}.
\]
By dominated convergence,
\[
\mathrm{Zq}(z)=\sum_{\mathrm{dn}\ge0}\frac{1}{\mathrm{dn}!}\cdot\mathrm{lw}\left(\left(\sum_iz_i\cdot X(f_i)\right)^{\mathrm{dn}}\right)\qquad\text{for all }z\in\mathbb{C}^n.
\]
The term with the index $\mathrm{dn}$ is a homogeneous polynomial of degree $\mathrm{dn}$ in $z$: every product of $\mathrm{dn}$ factors $X(f_i)$ is $\mathrm{lw}$-integrable by \cite[Proposition~7.10(4)]{QEDI}, applied to $\mathrm{dn}$ tests.
So this series is the expansion of the entire function $\mathrm{Zq}$ (\cite[Proposition~7.10(1)]{QEDI}) into homogeneous polynomials.
To see this, fix $z$; the series above at the point $\mathrm{zv}\cdot z$, and the Taylor series of $\mathrm{Zq}$ about $0$ at the point $\mathrm{zv}\cdot z$, ordered by the degree, are power series in $\mathrm{zv}\in\mathbb{C}$ with the same sum $\mathrm{Zq}(\mathrm{zv}\cdot z)$, so that their coefficients agree.
The monomial $\prod_{i\in\mathrm{Ix}}z_i$ occurs only in the term with $\mathrm{dn}=\#|\mathrm{Ix}|$, and its coefficient in $\left(\sum_iz_i\cdot X(f_i)\right)^{\mathrm{dn}}$ is $\mathrm{dn}!\cdot\prod_{i\in\mathrm{Ix}}X(f_i)$.
Hence its coefficient in $\mathrm{Zq}$ is $\mathrm{lw}\left(\prod_{i\in\mathrm{Ix}}X(f_i)\right)$.
Since all exponents of this monomial are equal to $1$, this coefficient is $\partial_{\mathrm{Ix}}\mathrm{Zq}(0)$.

\emph{Step~2.}
Let $\mathrm{phq}$ be analytic on the disc $\{\mathrm{zv}\in\mathbb{C}:|\mathrm{zv}-1|<1\}$.
We show that on the open set $\{z:|\mathrm{Zq}(z)-1|<1\}$, which contains $0$, one has for every non-empty $\mathrm{Ix}$
\begin{equation}\tag{F.2}
\partial_{\mathrm{Ix}}\left(\mathrm{phq}\circ\mathrm{Zq}\right)=\sum_{\mathrm{pt}\text{ partition of }\mathrm{Ix}}\mathrm{phq}^{(\#|\mathrm{pt}|)}(\mathrm{Zq})\cdot\prod_{\mathrm{bk}\in\mathrm{pt}}\partial_{\mathrm{bk}}\mathrm{Zq}.
\end{equation}
For $\mathrm{Ix}=\{i\}$ this is the chain rule.
Assume (F.2) for $\mathrm{Ix}$, and let $\ell\notin\mathrm{Ix}$.
The product rule and the chain rule give, for every partition $\mathrm{pt}$ of $\mathrm{Ix}$,
\begin{multline*}
\partial_{z_\ell}\left(\mathrm{phq}^{(\#|\mathrm{pt}|)}(\mathrm{Zq})\cdot\prod_{\mathrm{bk}\in\mathrm{pt}}\partial_{\mathrm{bk}}\mathrm{Zq}\right)
=\mathrm{phq}^{(\#|\mathrm{pt}|+1)}(\mathrm{Zq})\cdot\left(\partial_{z_\ell}\mathrm{Zq}\right)\cdot\prod_{\mathrm{bk}\in\mathrm{pt}}\partial_{\mathrm{bk}}\mathrm{Zq}\\
+\mathrm{phq}^{(\#|\mathrm{pt}|)}(\mathrm{Zq})\cdot\sum_{\mathrm{bk}_0\in\mathrm{pt}}\left(\partial_{\mathrm{bk}_0\cup\{\ell\}}\mathrm{Zq}\right)\cdot\prod_{\mathrm{bk}\in\mathrm{pt},\ \mathrm{bk}\ne\mathrm{bk}_0}\partial_{\mathrm{bk}}\mathrm{Zq}.
\end{multline*}
The first term on the right is the term of (F.2) for the partition of $\mathrm{Ix}\cup\{\ell\}$ whose blocks are those of $\mathrm{pt}$ and the block $\{\ell\}$.
The terms of the sum over $\mathrm{bk}_0$ are the terms of (F.2) for the partitions that arise from $\mathrm{pt}$ when $\ell$ is added to the block $\mathrm{bk}_0$.
Every partition of $\mathrm{Ix}\cup\{\ell\}$ arises in exactly one of these ways from exactly one partition $\mathrm{pt}$ of $\mathrm{Ix}$, namely from the partition obtained by removing $\ell$.
Summing over $\mathrm{pt}$ gives (F.2) for $\mathrm{Ix}\cup\{\ell\}$.

\emph{Step~3.}
For $\mathrm{phq}=\log$ one has $\mathrm{phq}^{(\mathrm{dn})}(1)=(-1)^{\mathrm{dn}-1}\cdot(\mathrm{dn}-1)!$.
Since $\mathrm{Zq}(0)=1$, the identity (F.2) at $z=0$ for $\mathrm{Ix}=\{1,\dots,n\}$, combined with Step~1, is (F.1) for the variables $X(f_1),\dots,X(f_n)$ under $\mathrm{lw}$.
\end{proof}

\subsubsection*{Proof}

\begin{proof}[Proof of Proposition~45.4\,(a)]
Let $\mathrm{lw}$ be soft.
By Proposition~45.1\,(10), the hypothesis (H) gives (L1), (L3) and (L4), and (dF) gives (L2).
By (Eu6) and $\mathrm{sgn}_*\mathrm{lw}=\mathrm{lw}$,
\[
\Phi_{\mathrm{lw}}(-f)=\Phi_{\mathrm{sgn}_*\mathrm{lw}}(f)=\Phi_{\mathrm{lw}}(f)\qquad\text{for real }f\in\mathcal{S}.
\]
Therefore Proposition~45.1, including its clause (9), and the Propositions~45.2 and~45.3 hold for $\mathrm{lw}$.

It remains to verify the properties (P1)--(P7) of Theorem~44\,(B1), with $\mathrm{lw}$ in the place of $\mu$.
For each of them we indicate the clauses from which it follows.
The conclusions (1)--(7) of \cite[Proposition~7.10]{QEDI} are taken under $\mathrm{lw}$, as in the proof of Lemma~F.1; they hold by Proposition~45.1\,(1).
\begin{itemize}
\item[(P1)] The two inequalities, for real $f'\in\mathcal{S}$ and for real $f'\in L^2(\mathbb{R}^3;\Lambda^2)$, are Proposition~45.1\,(4).
The extension is Proposition~45.1\,(1).
The bound $\left|\mathrm{lw}\left(\exp(z\cdot X(f'))\right)\right|\le\exp\left(\frac{(\operatorname{Re}z)^2\cdot\|f'\|_2^2}{2}\right)$ is \cite[Proposition~7.10(1)]{QEDI}.
The assertion $\mathrm{lw}\left(\exp(z\cdot X(f'))\right)\ne0$ for $|z|\cdot\|f'\|_2<0.9005$ follows from \cite[Proposition~7.10(5)]{QEDI} together with
\[
0.9005^2=0.81090025<0.81092<4\cdot\left(\frac{1}{5}+\frac{1}{375}+\frac{1}{15625}\right)<2\log\frac{3}{2},
\]
where the last inequality holds because $\log\frac{3}{2}=2\sum_{a\ge0}\frac{5^{-2a-1}}{2a+1}$.
Now let $n\ge1$ be an integer, and let $f_1,\dots,f_n$ lie in $L^2(\mathbb{R}^3;\Lambda^2)$ or in $L^2(\mathbb{R}^3;\Lambda^2)\otimes\mathbb{C}$.
By (n8), $X$ is complex linear on $L^2(\mathbb{R}^3;\Lambda^2)\otimes\mathbb{C}$; hence $\sum_iz_i\cdot X(f_i)=X\left(\sum_iz_i\cdot f_i\right)$ for $z\in\mathbb{C}^n$.
The norm on $L^2(\mathbb{R}^3;\Lambda^2)\otimes\mathbb{C}$ given by $\|g+ih\|_2^2=\|g\|_2^2+\|h\|_2^2$ for real $g,h$ is the norm of this complex Hilbert space, and it is the norm $|\cdot|$ that \cite[Proposition~7.10]{QEDI} puts on the complexification, for the space $L^2(\mathbb{R}^3;\Lambda^2)$ with the norm $\|\cdot\|_2$.
It is homogeneous under complex scalars and satisfies the triangle inequality.
Let $z\in\mathbb{C}^n$ satisfy $\sum_i|z_i|\cdot\|f_i\|_2\le\frac{2}{5}$.
Then
\[
\left\|\sum_iz_i\cdot f_i\right\|_2\le\sum_i|z_i|\cdot\|f_i\|_2\le\frac{2}{5},
\]
and \cite[Proposition~7.10(6)]{QEDI}, applied to $\sum_iz_i\cdot f_i$, gives
\[
\left|\mathrm{lw}\left(\exp\sum_iz_i\cdot X(f_i)\right)-1\right|\le\exp\left(2\cdot\left\|\sum_iz_i\cdot f_i\right\|_2^2\right)-1\le\exp\left(2\cdot\left(\frac{2}{5}\right)^2\right)-1\le\frac{1}{2};
\]
the last inequality holds because $\frac{8}{25}<\frac{2}{5}\le\log\frac{3}{2}$, the series for $\log\frac{3}{2}$ above having positive terms, so that $\exp\left(\frac{8}{25}\right)<\frac{3}{2}$.
Thus $\mathrm{lw}\left(\exp\sum_iz_i\cdot X(f_i)\right)$ lies in the closed disc of radius $\frac{1}{2}$ about $1$, which does not meet $(-\infty,0]$, so that the principal branch of the logarithm is defined at this number.
By \cite[Proposition~7.10(7)]{QEDI}, the function $z\mapsto\log\mathrm{lw}\left(\exp\sum_iz_i\cdot X(f_i)\right)$, formed with this branch, is analytic and of modulus at most $\log2$ on the set of the $z\in\mathbb{C}^n$ with $\sum_i|z_i|\cdot\|f_i\|_2\le\frac{2}{5}$.
\item[(P2)] The bound on $S_n$ is Proposition~45.2\,(1).
The bound on $\kappa_n$ is \cite[Proposition~7.10(7)]{QEDI}, applied to the derivative of the logarithm; by Lemma~F.1 this derivative is $\kappa_n(f_1,\dots,f_n)$.
The extension to $\mathrm{Sw}_n$ and its bound are Proposition~45.2\,(2).
\item[(P3)] The identity $S_0=1$, the symmetry and the reality are Proposition~45.2\,(1).
The invariance under $\omega\mapsto-\omega$ is part of the definition of a soft law.
The vanishing of $S_n$ for odd $n$ is Proposition~45.1\,(9); on $\mathrm{Sw}_n$ it is the last sentence of Proposition~45.2\,(6).
\item[(P4)] The determination of the law by its moments is Proposition~45.1\,(8); the identity $\mathrm{lw}(H_{-s})=1$ for $s>\frac{3}{2}$ is Proposition~45.1\,(3); the property (dF) is Proposition~45.1\,(4).
\item[(P5)] The invariance under $\mathbb{R}^3$, with the strongly continuous unitary group, is Proposition~45.1\,(5); the invariance under $O_{\mathrm{cub}}$ is Proposition~45.1\,(6).
\item[(P6)] The inequality for $\hat{r}_{e_j,\mathrm{po}}$ is Proposition~45.1\,(7); the inequality for the numbers $S_{m+n}$ is Proposition~45.2\,(5).
\item[(P7)] By Proposition~45.3\,(4) the five objects are the reconstruction (n4); the remaining assertions are Proposition~45.3\,(1) and~(2).
\end{itemize}
\end{proof}

\begin{proof}[Proof of Proposition~45.4\,(b)]
By (H), the law $\mathrm{lw}$ is invariant under $\mathbb{R}^3$.
Hence Proposition~45.6 applies to $\mathrm{lw}$, and the measure $\nu$ of Proposition~45.6\,(a$''$) is the measure $\nu$ of (n7): each of the two is the only probability measure on the laws that is carried by the $\mathbb{R}^3$-ergodic laws and represents $\mathrm{lw}$.

\emph{The first statement.}
By Proposition~45.6\,(a$''$), $\nu$ is carried by the $\mathbb{R}^3$-ergodic laws, and by (n3) an $\mathbb{R}^3$-ergodic law is invariant under $\mathbb{R}^3$.

\emph{The second statement.}
The law $\mathrm{lw}$ has (G) by (H); so the statement is Proposition~45.6\,(c).

\emph{The third statement.}
The event $\mathrm{Ecl}$ belongs to $\Sigma$ by Lemma~B.3, and $\mathrm{lw}(\mathrm{Ecl})=1$ by (dF).
Proposition~45.6\,(c$'$) with $\mathrm{Ev}=\mathrm{Ecl}$ gives $\mathrm{lw}_\xi(\mathrm{Ecl})=1$ for $\nu$-almost every $\xi$.

\emph{The fourth to sixth statement.}
Let $j\in\{1,2,3\}$.
As $\theta_j\in O_{\mathrm{cub}}$, the hypothesis (H) shows that $\mathrm{lw}$ is invariant under $\theta_j$; by (H) it is also reflection positive at $\{y_j=0\}$.
So the statement for this $j$ is Proposition~45.6\,(b).
\end{proof}

\begin{proof}[Proof of Proposition~45.4\,(c)]
Let $\xi$ be such that the six statements of (b) hold for $\mathrm{lw}_\xi$.
We recall the definitions of \cite[Section~8.1.3]{QEDI} in our notation.
A law $\mathrm{lw}'$ satisfies $(T_{123})$ if $(\tau_a)_*\mathrm{lw}'=\mathrm{lw}'$ for all $a\in\mathbb{R}^3$; it satisfies $(R_j)$ if $(\theta_j)_*\mathrm{lw}'=\mathrm{lw}'$; it satisfies $(RP_j)$ if it is reflection positive at $\{y_j=0\}$ in the sense of (n3).
The class $\mathcal{C}$ consists of the laws with $(T_{123})$ and with $(R_j)$, $(RP_j)$ for $j=1,2,3$.
By the first statement, $\mathrm{lw}_\xi$ has $(T_{123})$ and is $\mathbb{R}^3$-ergodic; the definition of ergodicity in \cite[Section~8.1.4]{QEDI}, for the index set $\{1,2,3\}$, is that of (n3).
By the last three statements, $\mathrm{lw}_\xi$ has $(R_j)$, and $(RP_j)$, which concerns the plane with $\mathrm{po}=0$, for $j=1,2,3$.
So $\mathrm{lw}_\xi$ belongs to the class $\mathcal{C}$.
By the second and the third statement, $\mathrm{lw}_\xi$ has (G) and (dF).
The group $\{\pm1\}^3$ of the diagonal matrices with entries $\pm1$ is generated by $\theta_1$, $\theta_2$, $\theta_3$, and the Euclidean maps that preserve $\mathrm{lw}_\xi$ form a group by (Eu5); so $\mathrm{lw}_\xi$ is invariant under $\{\pm1\}^3$.

It remains to list the hypotheses of the cited results.
Those of \cite[Theorem~8.3]{QEDI} are: for (i), a law in the class $\mathcal{C}$; for (ii), invariance under $\tau_{t\cdot e_j}$, $t$ real, and under $\theta_j$, together with $(RP_j)$; for (iii), a law in the class $\mathcal{C}$ that is $\mathbb{R}^3$-ergodic; for (iv), a law in the class $\mathcal{C}$ with (G).
Those of \cite[Proposition~8.19]{QEDI} are invariance under $\tau_{t\cdot e_j}$ and under $\theta_j$, the condition $(RP_j)$, and, in (iii), invariance under $\tau_y$ for $y\in e_j^{\perp}$.
The hypothesis of Proposition~45.5 is (G); that of its clause (d), for $\hat{n}=e_j$, is invariance under $r_{e_j,0}=\theta_j$ together with reflection positivity at $\{y_j=0\}$.
The law $\mathrm{lw}_\xi$ satisfies all of them.
\end{proof}

\begin{proof}[Proof of Proposition~45.4\,(d)]
The first sentence is Proposition~45.6\,(c$'$).
For the second sentence we use \cite[Lemma~8.21]{QEDI} for the index set $\{1,2,3\}$.
In our notation it states: let $\mathrm{lw}$ be invariant under $\mathbb{R}^3$, and let $\mathrm{Gi}$ be a Euclidean map with $\mathrm{Gi}_*\mathrm{lw}=\mathrm{lw}$; then the map $\mathrm{lw}'\mapsto\mathrm{Gi}_*\mathrm{lw}'$ preserves $\nu$, where $\nu$ is the image of $\mathrm{lw}$ under $\omega\mapsto\mathrm{lw}_\omega$.
This $\nu$ is the measure of Proposition~45.6\,(a$''$), and hence that of (n7).
The condition of the Lemma on the linear part $\mathrm{Rl}$ of $\mathrm{Gi}$, namely $\mathrm{Rl}\,\mathbb{R}^3=\mathbb{R}^3$, holds for every $\mathrm{Rl}\in O(3)$.
For $w\in O_{\mathrm{cub}}$ one has $w_*\mathrm{lw}=\mathrm{lw}$ by (H).
The map $\mathrm{lw}'\mapsto w_*\mathrm{lw}'$ of the set of the laws into itself is measurable by the last sentence of (Eu4).
So the image of $\nu$ under this map is defined, and by the Lemma it equals $\nu$; that is, $w_*\nu=\nu$ in the sense of (n8).
\end{proof}

\subsection*{G. Conditions on the constants}

The Propositions~45.1--45.4 hold for all values of the constants of Section~I.8, without any condition on them.
Their statements and their proofs involve none of these constants, and they rest on no line of Section~I.8.2 and on no line of Proposition~14.2.
The numbers $\frac{1}{2}$ (in the exponents), $\mathrm{ZH}_{2s}$, $2\cdot\left(2^n\cdot n!\right)^{\frac{1}{2}}$, $2\cdot\left(4^n\cdot n!\right)^{\frac{1}{2}}$, $\sigma_{2n}$ and $2\cdot66^n\cdot(n!)^{\frac{1}{2}}$ are absolute.
They are the same for all laws with (L1)--(L4).

\section*{Chapter 44. The theorem (cutoffs, limit points, volume-first limit points, ergodic components): proof of Theorem~44}

The statement proved in this chapter is Theorem~44 of Part~II.

\subsection*{Proof}

\subsubsection*{0. Conventions of this chapter}

\paragraph{0.1 Notation.}
Throughout the chapter the notation (n1)--(n8) for laws on $\mathcal{S}'$ of Part~II is in force, together with the notation stated with the Propositions~45.m, namely (L1)--(L4), $\Phi_{\mathrm{lw}}$, $\sigma_n$ and $\mathrm{ZH}_s$, the notation of Section~G6.A3 and of Sections~I.6.2, I.6.5, I.7.5 and~I.7.6, and the Data and the Constants of Proposition~43.2.

\emph{Laws on $H_{-4}$.}
By (n8) the set $H_{-4}$ belongs to $\Sigma$, and the trace of $\Sigma$ on $H_{-4}$ is the Borel $\sigma$-algebra of $H_{-4}$.
A Borel probability measure $\mu$ on $H_{-4}$ is identified with the law $\mathrm{Ev}\mapsto\mu(\mathrm{Ev}\cap H_{-4})$ on $(\mathcal{S}',\Sigma)$.
For $f'\in\mathcal{S}$ the functional $\omega\mapsto\omega(f')$ is continuous on $H_{-4}$, since
\[
|\omega(f')|\le\|\omega\|_{-4}\cdot\|f'\|_4
\]
\cite[Section~7.1.5]{QEDI}.

\emph{Almost every.}
The set of the laws is denoted $\mathrm{Lws}$, and $\mathrm{Blw}$ is the $\sigma$-algebra of (n7) on $\mathrm{Lws}$.
As in (n7), the words ``for $\nu$-almost every $\xi$'' mean: for all $\xi$ outside a set of $\mathrm{Blw}$ of $\nu$-measure $0$.

\emph{Cumulants.}
Let $\mathrm{xv}_1,\dots,\mathrm{xv}_{\mathrm{ns}}$ be real random variables on one probability space, with finite moments of all orders.
For a non-empty subset $\mathrm{bk}$ of $\{1,\dots,\mathrm{ns}\}$ let $\mathrm{jm}(\mathrm{bk})$ denote the expectation of $\prod_{a\in\mathrm{bk}}\mathrm{xv}_a$.
The joint cumulant of $\mathrm{xv}_1,\dots,\mathrm{xv}_{\mathrm{ns}}$ is
\[
\kappa_{\mathrm{ns}}(\mathrm{xv}_1,\dots,\mathrm{xv}_{\mathrm{ns}}):=\sum_{\mathrm{pn}}(-1)^{\#|\mathrm{pn}|-1}\cdot\left(\#|\mathrm{pn}|-1\right)!\cdot\prod_{\mathrm{bk}\in\mathrm{pn}}\mathrm{jm}(\mathrm{bk}),
\]
where $\mathrm{pn}$ ranges over the partitions of $\{1,\dots,\mathrm{ns}\}$ into non-empty blocks $\mathrm{bk}$, and $\#|\mathrm{pn}|$ is the number of the blocks of $\mathrm{pn}$.
The number $\kappa_{\mathrm{ns}}(f_1,\dots,f_{\mathrm{ns}})$ of (n5) is the joint cumulant of $X(f_1),\dots,X(f_{\mathrm{ns}})$ under the law at hand, and $\kappa_m\left(X_{\eta,R}(f')\right)$ is the joint cumulant of $m$ copies of $X_{\eta,R}(f')$ under $E^{\mathrm{af}}_n$.
If every $\mathrm{xv}_a$ has mean $0$, each partition with a block of one element contributes $0$.
For $\mathrm{ns}=2$ the partition into one block remains; for $\mathrm{ns}=4$ there remain the partition into one block, with the coefficient $1$, and the three partitions into two pairs, with the coefficient $-1$:
\begin{gather*}
\kappa_2(\mathrm{xv}_1,\mathrm{xv}_2)=\mathrm{jm}(\{1,2\}),\\
\kappa_4(\mathrm{xv}_1,\dots,\mathrm{xv}_4)=\mathrm{jm}(\{1,2,3,4\})-\mathrm{jm}(\{1,2\})\cdot\mathrm{jm}(\{3,4\})-\mathrm{jm}(\{1,3\})\cdot\mathrm{jm}(\{2,4\})-\mathrm{jm}(\{1,4\})\cdot\mathrm{jm}(\{2,3\}).
\end{gather*}

\emph{Further notation.}
In this chapter $(N_i,T_i)$, $i=1,2,\dots$, is a sequence of cutoffs, and
\[
\eta_i:=L^{-N_i},\qquad n_i:=8hL^{T_i+N_i},\qquad R_i:=8hL^{T_i};
\]
$E_i:=E^{\mathrm{af}}_{n_i}$ is formed with $q=e\cdot\eta_i^{\frac12}$ and $\boldsymbol{\mu}=\boldsymbol{\mu}_{N_i}$, and $X_i:=X_{\eta_i,R_i}$.
For $f'\in\mathcal{S}$ and $R>0$ let
\[
p_R(f'):=\|f'\|_2+26\cdot R^{-2}\cdot\left\|\,|y|_2^2\,f'\right\|_{L^2\left(\mathbb{R}^3\setminus[-R/2,R/2)^3\right)},
\]
the seminorm of \cite[Definition~12.4]{QEDII}.
The letter $\mathrm{ns}$ denotes a number of factors or of variables, and in Lemma~2.2 the letter $\alpha$ denotes a multi-index, as in (n5).
The constants $C_0$, $C_+$, $C_-$, $C_*$ in 3.2 are those of \cite[Theorem~9.21]{QEDI}; in 3.2 they take the values $C_0=C_+=C_-=C'_{\bar{\kappa}}$ and $C_*=2C_0+\frac14(C_++C_-)=\frac52\cdot C'_{\bar{\kappa}}$.

\paragraph{0.2 The charge; three inequalities.}
Let $e\in(0,e_2]$.
By the definition of $e_2$ (Proposition~43.2), $e\le e_2\le e_G$.
Since $\bar{m}\le1$ (Section~I.0.5), one has $h=L^{\lfloor\log_L(1/\bar{m})\rfloor}\ge1$, hence $\log h\ge0$.
The two conditions in the definition of $e_G$ (Section~I.8.2, line (C10)) read
\[
e\le e_1\cdot h^{-\frac12}\qquad\text{and, if }h>1,\qquad e\le\frac{1}{h^{\frac12}\cdot\log h}.
\]
Thus $e_G$ is the minimum of these numbers, and every $e\le e_G$ satisfies $e\cdot h^{\frac12}\le e_1$.
As $e_1\le\frac13$ (line (C9)) and $h\ge1$, it follows that $e\le\frac13$.
Moreover $\sigma>0$ (line (C7)) and $\mathfrak{e}_2=0$ (Section~I.8.5).
Consequently
\begin{equation}\tag{0.1}
e^{2\sigma}\le e^{\sigma}\le1\le e^{-\mathfrak{e}_2},\qquad e^{\sigma}\le e_2^{\sigma}.
\end{equation}
By Proposition~43.2 one has $C_U\ge60\,s_B^4\ge0$, $C_{\bar{\kappa}}=16C_U+6C_U^2$ and $C'_{\bar{\kappa}}=8C_{\bar{\kappa}}+3C_{\bar{\kappa}}^2$.
Hence
\begin{equation}\tag{0.2}
C_U\le C_{\bar{\kappa}},\qquad 8C_U+3C_U^2=\frac{C_{\bar{\kappa}}}{2}\le C'_{\bar{\kappa}},\qquad 2C_U\le C'_{\bar{\kappa}}.
\end{equation}
The third condition in the definition of $e_2$, combined with (0.1), gives
\begin{equation}\tag{0.3}
\frac52\cdot C'_{\bar{\kappa}}\cdot e^{\sigma}\le\frac{b(f)}{4},\qquad\text{that is,}\qquad C'_{\bar{\kappa}}\cdot e^{\sigma}\le\frac{b(f)}{10}.
\end{equation}

\paragraph{0.3 The laws $P_{N,T}$ are the laws $P^{\mathrm{ch}}_{N,T}$ of \cite[Definition~2.36]{QEDII}.}
Let $N\mapsto m_{0,N}$ be a sequence of real numbers, let $e>0$ and $N,T\ge0$, and put
\[
n=8hL^{T+N},\qquad q=e\cdot\eta^{\frac12},\qquad\boldsymbol{\mu}=\bar{m}\eta+m_{0,N};
\]
thus $n\ge8$.
By Proposition~46\,(d), (f), the number $Z_n[\Phi]$ of (I.7.1) equals the number $Z_{C(n,n,n)}[\Phi]$ of \cite[Definition~2.13]{QEDII}, formed on the torus $C(n,n,n)$ of \cite[Definition~2.11]{QEDII} with the Wilson parameter $\rho$, the standard weights, the mass parameter $M_{\mathrm{ch}}=\boldsymbol{\mu}+\frac{\rho}{2}\cdot\mathrm{sw}$ and the lattice charge $q$, for every Borel function $\Phi$ of $F^{\mathrm{u}}$ with $|\Phi|\le\mathrm{cf}\cdot\exp\left(\mathrm{cf}\cdot\|F^{\mathrm{u}}\|\right)$.
Hence $E^{\mathrm{af}}_n$ coincides on these functions with the expectation $\mathbb{E}_{C(n,n,n)}$ of \cite[Definition~2.13]{QEDII}, and the law $P_{N,T}$ of Section~I.7.5 is the law $P^{\mathrm{ch}}_{N,T}$ of \cite[Definition~2.36]{QEDII}; the integer $L^{\lfloor\log_L\bar{m}^{-1}\rfloor}$ in the side $n$ of that definition is $h$.
The point $v_0$ of \cite[Definition~2.36]{QEDII} does not enter (Proposition~46\,(a)).
The data required by \cite[Definition~2.36]{QEDII} are: $\rho\in(0,1)$; the standard weights; an odd integer $L\ge3$; $0<\bar{m}\le1$; $e>0$; a real sequence $N\mapsto m_{0,N}$ that does not depend on $T$.
The first four of them are given by Section~I.0.5, where $L$ is an odd integer $\ge5$.
The limit points, the sets $\mathcal{V}_N$ and the volume-first limit points of Section~I.7.5 are those of \cite[Definition~2.37]{QEDII}.
The condition (G), the hypothesis (H), reflection positivity at a plane, the class $\mathcal{C}_9$ and the reconstruction (n4) are those of \cite[Definition~2.38]{QEDII}, and (G), (H), reflection positivity at a plane and the reconstruction are also those of \cite[Sections~8.1.3, 8.1.5 and Proposition~8.19]{QEDI}.
In \cite[Definition~2.38]{QEDII} and in \cite[Section~8.1.3]{QEDI} the invariance under $\mathbb{R}^3$ is the condition $(T_{123})$, the invariance under $\theta_j$ is the condition $(R_j)$, and reflection positivity at $\{y_j=0\}$ is the condition $(RP_j)$.

The statements of \cite{QEDII} are applied in the notation of this chapter:
the variable of integration $F$ of \cite[Proposition~12.13(a)]{QEDII} is written $\omega$;
a point $x$ of $\mathbb{R}^3$ is written $y$;
the position $c$ and the unit normal $n$ of a plane are written $\mathrm{po}$ and $\hat{n}$;
the functions $\Phi$, $\Phi'$ in $L^2$ are written $\mathrm{Ga}$, $\mathrm{Gb}$;
the pair $(\lambda,k)$ of an energy $\lambda$ and a momentum $k$, the variable of the projection-valued measure $\mathfrak{E}_j$ of \cite[Lemma~12.18]{QEDII}, is written $(E,p)$, and $\mathfrak{E}_j$ is written $\mathrm{Es}_j$;
the half-side $r$ of the cubes $(-r,r)^3$ in \cite[Theorem~12.39]{QEDII} is written $\mathrm{vp}$, and the constants $C_0$ and $C_1$ of that theorem are $c_{\mathrm{cl}}$ and $4$;
the measure $\mathfrak{K}_\mu$ of \cite[Theorem~1.4(d)]{QEDII} is $\mathrm{kl}_\mu$, with the variable $t$ in the place of $m$.

\subsubsection*{1. The first sentence; (A)}

\paragraph{1.1}
The first sentence of Theorem~44 is Proposition~43.1\,(R0): for $e>0$ with $e\cdot h^{\frac12}\le e_1$ the numbers $m_{0,N}(e)$, $N\ge0$, are real, $|m_{0,N}(e)|<\frac18$, and they depend on $(N,e,\bar{m},\rho)$ and on the scheme of Section~I.6.2, not on a torus.
The index $T$ is the index of the torus of the side $R=8hL^T$ (Section~I.6.5); so these numbers do not depend on $T$, and $N\mapsto m_{0,N}(e)$ is a mass assignment (Section~I.7.5).
By 0.2 this applies to every $e\in(0,e_2]$, and 0.3 applies to the laws $P_{N,T}$ of Theorem~44.

\paragraph{1.2 (A).}
\begin{proof}[Proof of (A)]
Let $e\in(0,e_2]$, $N\ge N_0(e)$, $T\ge T^{*}(e)+1$ and $\mathrm{ix}\in\{0,+,-\}$.
The clauses (U-a), (U-a$'$), (U-b), (U-c) are the assertions of Proposition~43.2 for these data and for a placement $v$ of $\beta_{\mathrm{ix}}$.
Such a placement exists.
It is a point $v\in\eta\mathbb{Z}^3$ with
\[
\left|v+y_{\beta_{\mathrm{ix}}}-\frac{\ell_c}{2}\cdot(1,1,1)\right|\le\frac{\eta}{2}
\]
(Section~I.7.6), where $|\cdot|$ is the sup norm (Section~I.1.1); for $\mu=1,2,3$ let $v_\mu$ be a point of $\eta\mathbb{Z}$ nearest to the number $\frac{\ell_c}{2}-(y_{\beta_{\mathrm{ix}}})_\mu$, so that its distance from this number is at most $\frac{\eta}{2}$.
The first inequality of (U-a) and the inequality $C_U\le C_{\bar{\kappa}}$ of (0.2) give
\begin{align*}
\left|\kappa_4\left(X_{\eta,R}(g_{\mathrm{ix}})\right)-e^4\cdot b(g_{\mathrm{ix}})\right|
&\le\frac{C_{\bar{\kappa}}}{2}\cdot e^{4+\sigma}+C_U\cdot\left(\eta^2+e^4\cdot\eta^{\frac14}\right)\\
&\le\frac{C_{\bar{\kappa}}}{2}\cdot e^{4+\sigma}+C_{\bar{\kappa}}\cdot\left(\eta^2+e^4\cdot\eta^{\frac14}\right).
\end{align*}
The right side involves $e$, $\eta=L^{-N}$, $\sigma$ and $C_{\bar{\kappa}}$ only; it does not involve $T$.
\end{proof}

\subsubsection*{2. (B0), (B1)}

\begin{named}{Lemma 2.1}[Bounds at one cutoff]
Let $N,T\ge0$, $\eta=L^{-N}$, $n=8hL^{T+N}$, $R=8hL^T$, $e>0$, $q=e\cdot\eta^{\frac12}$ and $\boldsymbol{\mu}=\bar{m}\eta+m_{0,N}$ for a real sequence $N\mapsto m_{0,N}$.
Let $f',f_1,\dots,f_{\mathrm{ns}}$ be real elements of $\mathcal{S}$.
\begin{itemize}
\item[(a)]
One has
\[
E^{\mathrm{af}}_n\exp X_{\eta,R}(f')\le\exp\left(\frac12|f'_\eta|_\eta^2\right).
\]
\item[(b)]
One has $|f'_\eta|_\eta\le p_R(f')\le p_1(f')$, and $p_R(f')\to\|f'\|_2$ as $R\to\infty$; $p_1$ is a seminorm on $\mathcal{S}$; and $|f'_\eta|_\eta\le\|f'\|_2$ if $\operatorname{supp}f'\subset(-R/2,R/2)^3$.
\item[(c)]
One has $E^{\mathrm{af}}_nX_{\eta,R}(f')=0$,
\begin{gather*}
E^{\mathrm{af}}_n\exp\left(t\cdot|X_{\eta,R}(f')|\right)\le2\exp\left(\frac12t^2\cdot|f'_\eta|_\eta^2\right)\qquad(t\ge0),\\
E^{\mathrm{af}}_n\prod_{a\le\mathrm{ns}}|X_{\eta,R}(f_a)|\le\sigma_{\mathrm{ns}}\cdot\prod_{a\le\mathrm{ns}}|f_{a,\eta}|_\eta,
\end{gather*}
and $z\mapsto E^{\mathrm{af}}_n\exp\sum_{a\le\mathrm{ns}}z_aX_{\eta,R}(f_a)$ is an entire function on $\mathbb{C}^{\mathrm{ns}}$.
\end{itemize}
\end{named}

\begin{proof}
(a) One has $X_{\eta,R}(f')=\langle f'_\eta,F\rangle_\eta$ (Section~I.7.5), and $f'_\eta$ is a real two-cochain on $\Lambda_{\eta,R}$; the inequality is \cite[Proposition~12.10(c)]{QEDII}.

(b) This is \cite[Lemma~12.5(c)]{QEDII}, which applies because $R=8hL^T\ge1$.

(c) The map $f'\mapsto f'_\eta$ is linear.
By (a), the hypothesis of \cite[Proposition~7.10]{QEDI} is fulfilled by the probability space formed by $C^{2}(n)$ with the law of $F^{\mathrm{u}}$ under $E^{\mathrm{af}}_n$, by the real vector space $\mathcal{S}$ with the Hilbert seminorm $f'\mapsto|f'_\eta|_\eta$, and by the linear map $f'\mapsto X_{\eta,R}(f')$.
The four assertions of (c) are, in this order, its conclusions (2), (1), (4) and (1).
\end{proof}

\begin{named}{Lemma 2.2}[(P8)]
Let $N_i,T_i\to\infty$ and $P_{N_i,T_i}\to\mu$ weakly on $H_{-4}$, and let $f_1,\dots,f_k\in\mathcal{S}$.
For $z\in\mathbb{C}^k$ put
\[
\mathrm{Gn}_i(z):=E_i\exp\sum_{a\le k}z_aX_i(f_a),\qquad\mathrm{Gn}(z):=\mu\left(\exp\sum_{a\le k}z_aX(f_a)\right).
\]
Then the integral that defines $\mathrm{Gn}(z)$ converges absolutely for every $z$, and $\mathrm{Gn}_i\to\mathrm{Gn}$ uniformly on every bounded subset of $\mathbb{C}^k$.
Every joint moment $E_i\prod_aX_i(f_a)^{\alpha_a}$ converges to $\mu\left(\prod_aX(f_a)^{\alpha_a}\right)$, and every joint cumulant of the variables $X_i(f_a)$ under $E_i$ converges to the joint cumulant of the variables $X(f_a)$ under $\mu$.
\end{named}

\begin{proof}
\begin{itemize}
\item[(1)]
$P_{N,T}$ is the law of $F^{\sharp}$, and $F^{\sharp}(f')=X_{\eta,R}(f')$ (Section~I.7.5).
Hence
\[
E_i\,\varphi\left(X_i(f_1),\dots,X_i(f_k)\right)=\int\mathrm{d}P_{N_i,T_i}(\omega)\,\varphi\left(\omega(f_1),\dots,\omega(f_k)\right)
\]
for every Borel function $\varphi$ on $\mathbb{R}^k$ for which one of the two sides exists.
\item[(2)]
Let $\mathrm{rd}>0$ and $|z_a|\le\mathrm{rd}$ for all $a$.
Then $\left|\exp\sum_az_aX_i(f_a)\right|^2=\exp X_i(f')$ with the real test $f':=2\sum_a(\operatorname{Re}z_a)\,f_a$.
By 2.1 (a), (b),
\[
E_i\left|\exp\sum_az_aX_i(f_a)\right|^2\le\exp\left(\frac12p_1(f')^2\right)\le\exp\left(2\,\mathrm{rd}^2\cdot\left(\sum_ap_1(f_a)\right)^2\right)=:\mathrm{bd}_{\mathrm{rd}}^2.
\]
\item[(3)]
The function $\omega\mapsto\exp\sum_az_a\,\omega(f_a)$ is continuous on $H_{-4}$ (0.1, Laws on $H_{-4}$).
By the mapping theorem \cite[Theorem~2.7]{Billingsley1999}, the laws of its real part and of its imaginary part under $P_{N_i,T_i}$ converge weakly to their laws under $\mu$.
By (2), these variables have second moments at most $\mathrm{bd}_{\mathrm{rd}}^2$ for all $i$, so that they are uniformly integrable.
By \cite[Theorem~3.5]{Billingsley1999} they are integrable under $\mu$, and $\mathrm{Gn}_i(z)\to\mathrm{Gn}(z)$.
\item[(4)]
$\mathrm{Gn}_i$ is entire (2.1 (c)), and $|\mathrm{Gn}_i|\le\mathrm{bd}_{\mathrm{rd}}$ on the polydisc $\{|z_a|\le\mathrm{rd}\text{ for all }a\}$, by (2) and the Cauchy--Schwarz inequality.
Cauchy's estimate on the discs of radius $\mathrm{rd}$ contained in the polydisc of radius $2\,\mathrm{rd}$ gives
\[
\left|\frac{\partial\mathrm{Gn}_i}{\partial z_a}\right|\le\frac{\mathrm{bd}_{2\mathrm{rd}}}{\mathrm{rd}}\qquad\text{on }\{|z_{\mathrm{bq}}|\le\mathrm{rd}\text{ for all }\mathrm{bq}\le k\}.
\]
So all $\mathrm{Gn}_i$ are Lipschitz functions on this convex compact set, for the norm $\sum_a|z_a|$, with the one constant $\mathrm{bd}_{2\mathrm{rd}}/\mathrm{rd}$.
A sequence of functions on a compact set that has a common Lipschitz constant and converges at every point converges uniformly.
Therefore $\mathrm{Gn}_i\to\mathrm{Gn}$ uniformly on every polydisc, $\mathrm{Gn}$ is entire, and, by Cauchy's formula, $\partial^\alpha\mathrm{Gn}_i(0)\to\partial^\alpha\mathrm{Gn}(0)$ for every multi-index $\alpha$.
\item[(5)]
Let $\mathrm{xv}_1,\dots,\mathrm{xv}_k$ be real random variables such that $\exp\left(t\cdot|\mathrm{xv}_a|\right)$ is integrable for all $t\ge0$ and all $a$, and put $\mathrm{xv}^\alpha:=\prod_a\mathrm{xv}_a^{\alpha_a}$.
Then
\[
\sum_\alpha\frac{|z^\alpha\cdot\mathrm{xv}^\alpha|}{\alpha!}=\exp\sum_a|z_a|\cdot|\mathrm{xv}_a|,
\]
which is integrable by H\"older's inequality.
Hence the expectation of $\exp\sum_az_a\mathrm{xv}_a$ is the sum over $\alpha$ of $\frac{z^\alpha}{\alpha!}$ times the expectation of $\mathrm{xv}^\alpha$, and its derivative $\partial^\alpha$ at $z=0$ is the expectation of $\mathrm{xv}^\alpha$.
This applies to $E_i$ and $\mathrm{xv}_a=X_i(f_a)$ (2.1 (c)), and to $\mu$ and $\mathrm{xv}_a=X(f_a)$; in the second case the integrability follows from (3) at the two points $z$ with $z_a=\pm t$ and all other coordinates $0$, since $\exp\left(t\cdot|\mathrm{xv}_a|\right)\le\exp\left(t\cdot\mathrm{xv}_a\right)+\exp\left(-t\cdot\mathrm{xv}_a\right)$.
Together with (4), all joint moments converge.
\item[(6)]
A joint cumulant is a polynomial in the joint moments whose coefficients do not depend on the law (0.1, Cumulants).
By (5), all joint cumulants converge.
\end{itemize}
\end{proof}

\begin{named}{Lemma 2.2a}[The two forms of a cumulant]
Let $\mathrm{xv}_1,\dots,\mathrm{xv}_k$ be as in 2.2 (5), and let $\mathrm{Gj}(z)$ be the expectation of $\exp\sum_{a\le k}z_a\mathrm{xv}_a$.
Then $\partial_{z_1}\cdots\partial_{z_k}\log\mathrm{Gj}$ at $z=0$ is the joint cumulant $\kappa_k(\mathrm{xv}_1,\dots,\mathrm{xv}_k)$ of 0.1.
\end{named}

\begin{proof}
By 2.2 (5), $\mathrm{Gj}$ is the sum of a power series that converges on all of $\mathbb{C}^k$, $\mathrm{Gj}(0)=1$, and for every non-empty subset $\mathrm{bk}$ of $\{1,\dots,k\}$ the derivative
\[
\partial_{\mathrm{bk}}\mathrm{Gj}:=\left(\prod_{a\in\mathrm{bk}}\partial_{z_a}\right)\mathrm{Gj}
\]
takes the value $\mathrm{jm}(\mathrm{bk})$ at $z=0$.
On a neighbourhood of $0$ on which $\mathrm{Gj}\ne0$ one has, for $\mathrm{ns}=1,\dots,k$,
\[
\partial_{z_1}\cdots\partial_{z_{\mathrm{ns}}}\log\mathrm{Gj}=\sum_{\mathrm{pn}}(-1)^{\#|\mathrm{pn}|-1}\cdot\left(\#|\mathrm{pn}|-1\right)!\cdot\mathrm{Gj}^{-\#|\mathrm{pn}|}\cdot\prod_{\mathrm{bk}\in\mathrm{pn}}\partial_{\mathrm{bk}}\mathrm{Gj},
\]
where $\mathrm{pn}$ ranges over the partitions of $\{1,\dots,\mathrm{ns}\}$.
For $\mathrm{ns}=1$ this is $\partial_{z_1}\log\mathrm{Gj}=\mathrm{Gj}^{-1}\cdot\partial_{z_1}\mathrm{Gj}$.
Suppose the formula holds for some $\mathrm{ns}<k$.
By the product rule, the derivative $\partial_{z_{\mathrm{ns}+1}}$ of the term of $\mathrm{pn}$ is a sum of $\#|\mathrm{pn}|+1$ terms.
In $\#|\mathrm{pn}|$ of them the element $\mathrm{ns}+1$ is adjoined to one block of $\mathrm{pn}$, with the same coefficient and the same power of $\mathrm{Gj}$.
In the remaining term $\{\mathrm{ns}+1\}$ is a new block, the power is $\mathrm{Gj}^{-\#|\mathrm{pn}|-1}$, and the coefficient is
\[
(-1)^{\#|\mathrm{pn}|-1}\cdot\left(\#|\mathrm{pn}|-1\right)!\cdot\left(-\#|\mathrm{pn}|\right)=(-1)^{\#|\mathrm{pn}|}\cdot\left(\#|\mathrm{pn}|\right)!.
\]
Each partition of $\{1,\dots,\mathrm{ns}+1\}$ arises in this way exactly once, namely from the partition of $\{1,\dots,\mathrm{ns}\}$ that remains when $\mathrm{ns}+1$ is removed.
This is the formula for $\mathrm{ns}+1$.
For $\mathrm{ns}=k$ and $z=0$ it is the assertion.
\end{proof}

\paragraph{2.3 (B0).}
\begin{proof}[Proof of (B0)]
By 1.1 and 0.3, \cite[Proposition~12.13(a)]{QEDII} applies to the family $(P_{N,T})$.
It gives
\[
\sup_{N,T}\int\mathrm{d}P_{N,T}(\omega)\,\|\omega\|_{-3}^2\le61,
\]
the tightness of the family on $H_{-4}$, and, for every sequence of cutoffs, a subsequence along which the laws converge weakly on $H_{-4}$.
\end{proof}

\paragraph{2.4 (B1).}
\begin{proof}[Proof of (B1)]
Let $\mu$ be a limit point and $(N_i,T_i)$ a defining sequence: $N_i,T_i\to\infty$ and $P_{N_i,T_i}\to\mu$ weakly on $H_{-4}$.
For all $N,T\ge0$ the law of $F=\eta^{-\frac32}\cdot F^{\mathrm{u}}$ under $E^{\mathrm{af}}_n$ has the properties (a)--(e) of \cite[Proposition~12.10]{QEDII}.
The properties (a), (b), (c), (e) are the hypotheses of \cite[Proposition~12.12]{QEDII}, and (d) is the hypothesis of the step (8) of the proof of \cite[Proposition~12.12]{QEDII}.
The steps (3), (5), (6), (7), (8) of that proof show, in this order, that $\mu$ has (L1), that $\mu$ has (L2), that $\mu$ has (L3), that $\mu$ has (L4), and that $\Phi_\mu(-f')=\Phi_\mu(f')$ for all real $f'\in\mathcal{S}$.
Hence Proposition~45.1, with its clause (9), Proposition~45.2 and Proposition~45.3 hold for $\mathrm{lw}=\mu$.
The condition (G) is (L1); the invariance of $\mu$ under $\mathbb{R}^3$ and under $O_{\mathrm{cub}}$ is Proposition~45.1\,(5), (6); reflection positivity at $\{y_j=0\}$, $j=1,2,3$, is (L4); and (dF) is the first assertion of Proposition~45.1\,(4).
Therefore
\begin{equation}\tag{2.1}
\mu\ \text{satisfies (H) and (dF).}
\end{equation}
The properties (P1)--(P8) are obtained clause by clause.
\begin{itemize}
\item[(P1)]
For all real $f'\in L^2$,
\[
\mu\left(\exp X(f')\right)\le\exp\left(\frac12\langle f',\Pi_{\mathrm{c}}\,f'\rangle\right)\qquad\text{and}\qquad\langle f',\Pi_{\mathrm{c}}\,f'\rangle\le\|f'\|_2^2:
\]
this is Proposition~45.1\,(4).
The extension of $X$ and its boundedness are Proposition~45.1\,(1).
By Proposition~45.1\,(1), the conclusions (1)--(7) of \cite[Proposition~7.10]{QEDI} hold under $\mu$ for the space $L^2(\mathbb{R}^3;\Lambda^2)$ with the norm $\|\cdot\|_2$.
Conclusion (1) gives
\[
\left|\mu\left(\exp\left(zX(f')\right)\right)\right|\le\exp\left(\frac12(\operatorname{Re}z)^2\cdot\|f'\|_2^2\right).
\]
Conclusion (5) gives $\mu\left(\exp\left(zX(f')\right)\right)\ne0$ if $|z|\cdot\|f'\|_2<\left(2\log\frac32\right)^{\frac12}$.
Here $\left(2\log\frac32\right)^{\frac12}>0.9005$.
Indeed, $0.9005^2=0.81090025$, and
\[
\log\frac32=\log\left(\left(1+\frac15\right)\cdot\left(1-\frac15\right)^{-1}\right)=2\left(\frac15+\frac13\left(\frac15\right)^3+\frac15\left(\frac15\right)^5+\cdots\right)
\]
is a series of positive terms, so that
\[
2\log\frac32\ge4\left(\frac15+\frac{1}{375}+\frac{1}{15625}\right)=\frac{38012}{46875}>0.8109226>0.81090025.
\]
Let $n\ge1$ be an integer, and let $f_1,\dots,f_n$ lie in $L^2(\mathbb{R}^3;\Lambda^2)$ or in its complexification $L^2(\mathbb{R}^3;\Lambda^2)\otimes\mathbb{C}$.
On the complexification $X$ is extended complex linearly, as in (n8), and the norm is defined by $\|f_i\|_2^2=\|\operatorname{Re}f_i\|_2^2+\|\operatorname{Im}f_i\|_2^2$;
these are the extensions of $X$ and of the norm in \cite[Proposition~7.10]{QEDI}.
The norm so defined is that of the complex Hilbert space $L^2(\mathbb{R}^3;\Lambda^2)\otimes\mathbb{C}$; it is homogeneous under complex scalars and satisfies the triangle inequality.
Let $z\in\mathbb{C}^n$ satisfy $\sum_i|z_i|\cdot\|f_i\|_2\le\frac25$.
Then $\sum_iz_iX(f_i)=X\left(\sum_iz_i\,f_i\right)$ and
\[
\left\|\sum_iz_i\,f_i\right\|_2\le\sum_i|z_i|\cdot\|f_i\|_2\le\frac25 .
\]
Conclusion (6), applied to $\sum_iz_i\,f_i$, gives
\[
\left|\mu\left(\exp\sum_iz_iX(f_i)\right)-1\right|\le\exp\left(2\left\|\sum_iz_i\,f_i\right\|_2^2\right)-1\le\exp\left(2\cdot\left(\frac25\right)^2\right)-1\le\frac12 ;
\]
the last step holds because $\frac{8}{25}<\frac25<\log\frac32$, by the series above, so that $\exp\left(\frac{8}{25}\right)<\frac32$.
Thus $\mu\left(\exp\sum_iz_iX(f_i)\right)$ lies in the closed disc of radius $\frac12$ about $1$, which does not meet $(-\infty,0]$, and the principal branch of the logarithm is defined at this number for every such $z$.
Conclusion (7) gives that the function $z\mapsto\log\mu\left(\exp\sum_iz_iX(f_i)\right)$, formed with the principal branch, is analytic, and of modulus at most $\log2$, on the set $\left\{z\in\mathbb{C}^n:\sum_i|z_i|\cdot\|f_i\|_2\le\frac25\right\}$.
\item[(P2)]
The bound on $S_n(f_1,\dots,f_n)$ is Proposition~45.2\,(1).
The bound on $\kappa_n$ is conclusion (7):
\[
|\kappa_n(f_1,\dots,f_n)|\le(\log2)\cdot\left(\frac{5n}{2}\right)^n\cdot\prod_i\|f_i\|_2\le(\log2)\cdot\left(\frac{5\mathrm{e}}{2}\right)^n\cdot n!\cdot\prod_i\|f_i\|_2.
\]
In conclusion (7) the number $\kappa_n(f_1,\dots,f_n)$ is the derivative $\partial_{z_1}\cdots\partial_{z_n}$ of $\log\mu\left(\exp\sum_iz_iX(f_i)\right)$ at $z=0$.
The functions $\exp\left(t\cdot|X(f_i)|\right)$, $t\ge0$, are integrable by conclusion (1); so, by Lemma~2.2a, this number is the joint cumulant of 0.1 and (n5).
The extension to $\mathrm{Sw}_n$, its uniqueness and the bound by $q_{6n}$ are Proposition~45.2\,(2).
\item[(P3)]
Proposition~45.2\,(1) and Proposition~45.1\,(9).
\item[(P4)]
Proposition~45.1\,(8), (3) and (4).
\item[(P5)]
Proposition~45.1\,(5), (6).
\item[(P6)]
Proposition~45.1\,(7) and Proposition~45.2\,(5).
\item[(P7)]
Proposition~45.3\,(1), (2), (4).
\item[(P8)]
Lemma~2.2.
\end{itemize}
\end{proof}

\subsubsection*{3. (B2), (B3)}

From here to the end of Section~5, $e\in(0,e_2]$, $\mu$ is a limit point of $(P_{N,T})$, and $(N_i,T_i)$ is a defining sequence.
The numbers $N_0(e)$ and $T^{*}(e)$ are integers (Section~I.6.2), and $N_i,T_i\to\infty$; so there is an index $i_0$ with $N_i\ge N_0(e)$ and $T_i\ge T^{*}(e)+1$ for all $i\ge i_0$.
Let $i\ge i_0$ and $\mathrm{ix}\in\{0,+,-\}$.
The probe $\beta_{\mathrm{ix}}$ fits $(N_i,T^{*}(e))$ by the first clause of Proposition~43.1\,(R3), since $e\le e_2\le\min\left(e_G,e_{\sharp}(\beta_{\mathrm{ix}})\right)$ by the definition of $e_2$.
The clauses (U-a), (U-a$'$), (U-b), (U-c) hold at the cutoff $(N_i,T_i)$, with the numbers $c_m$, $\mathsf{u}_m$ and $\boldsymbol{\sigma}^2(N_i,\beta_{\mathrm{ix}})$ that belong to a placement of $\beta_{\mathrm{ix}}$ at the cutoff $N_i$ (1.2).

\paragraph{3.1 (B2).}
\begin{proof}[Proof of (B2)]
Fix $\mathrm{ix}$.
The test $g_{\mathrm{ix}}=dd^{*}\beta_{\mathrm{ix}}$ is a real element of $C_c^\infty(\mathbb{R}^3;\Lambda^2)$, which is a subset of $\mathcal{S}$.
By (P8),
\[
\kappa_4\left(X_i(g_{\mathrm{ix}})\right)\to\kappa_4^\mu(g_{\mathrm{ix}}),\qquad\kappa_2\left(X_i(g_{\mathrm{ix}})\right)\to\kappa_2^\mu(g_{\mathrm{ix}},g_{\mathrm{ix}})=S_2^\mu(g_{\mathrm{ix}},g_{\mathrm{ix}}),
\]
where the last equality holds because $S_1^\mu=0$ (Proposition~45.2\,(1); 0.1).

\emph{Fourth cumulant.}
For $i\ge i_0$ the first inequality of (U-a) reads
\[
\left|\kappa_4\left(X_i(g_{\mathrm{ix}})\right)-e^4\cdot b(g_{\mathrm{ix}})\right|\le\frac{C_{\bar{\kappa}}}{2}\cdot e^{4+\sigma}+C_U\cdot\left(\eta_i^2+e^4\cdot\eta_i^{\frac14}\right).
\]
Since $\eta_i\to0$, it follows, with (0.2), that
\begin{equation}\tag{3.1}
\left|\kappa_4^\mu(g_{\mathrm{ix}})-e^4\cdot b(g_{\mathrm{ix}})\right|\le\frac{C_{\bar{\kappa}}}{2}\cdot e^{4+\sigma}\le C'_{\bar{\kappa}}\cdot e^{4+\sigma}.
\end{equation}

\emph{Second cumulant.}
First, $\mathsf{u}_2\ge0$: for $m=2$ and a set $\mathrm{Ypt}$ of one point, the left side of (U-b) is a modulus, and the right side is then $\mathsf{u}_2$.
The second inequality of (U-a$'$), the bounds $|c_2|\le C_U\cdot e^{-\mathfrak{e}_2}$ and $\mathsf{u}_2\le(1+|c_2|)\cdot\mathsf{u}_2\le C_U\cdot e^{\sigma}$ of (U-c), and (0.1) give, for $i\ge i_0$,
\begin{align*}
\left|\kappa_2\left(X_i(g_{\mathrm{ix}})\right)-\boldsymbol{\sigma}^2(N_i,\beta_{\mathrm{ix}})\right|
&\le e^2\cdot|c_2|+e^2\cdot\mathsf{u}_2+2s_B^2\cdot\eta_i^2\\
&\le e^2\cdot\left(C_U\cdot e^{-\mathfrak{e}_2}+C_U\cdot e^{\sigma}\right)+2s_B^2\cdot\eta_i^2\\
&\le2C_U\cdot e^{2-\mathfrak{e}_2}+2s_B^2\cdot\eta_i^2.
\end{align*}
Moreover $\boldsymbol{\sigma}^2(N_i,\beta_{\mathrm{ix}})\to\|g_{\mathrm{ix}}\|_2^2$ as $i\to\infty$, by Proposition~47.2\,(e) for the probe $\beta_{\mathrm{ix}}$ (see the last paragraph of Proposition~43.2).
Hence, with (0.2),
\begin{equation}\tag{3.2}
\left|S_2^\mu(g_{\mathrm{ix}},g_{\mathrm{ix}})-\|g_{\mathrm{ix}}\|_2^2\right|\le2C_U\cdot e^{2-\mathfrak{e}_2}\le C'_{\bar{\kappa}}\cdot e^{2-\mathfrak{e}_2}.
\end{equation}
\end{proof}

\begin{named}{Lemma 3.2}[Three tests]
Let $e\in(0,e_2]$, and let $\mathrm{lw}$ be a law with (G) that is invariant under $\mathbb{R}^3$ and satisfies
\[
\left|\kappa_4^{\mathrm{lw}}(g_{\mathrm{ix}})-e^4\cdot b(g_{\mathrm{ix}})\right|\le C'_{\bar{\kappa}}\cdot e^{4+\sigma}\qquad(\mathrm{ix}=0,+,-).
\]
Then
\[
\mathrm{Ext}^{\mathrm{lw}}_{x_0}(f)\ge\frac12e^4\cdot b(f)\qquad\text{and}\qquad\kappa_4^{\mathrm{lw}}(f)\ge\frac{9}{10}e^4\cdot b(f)\ge\frac12e^4\cdot b(f)>0.
\]
\end{named}

\begin{proof}
The condition (G) is the hypothesis (H1) of \cite[Section~9.2]{QEDI}.
By \cite[Lemma~9.1(a), (c)]{QEDI}, $\mathrm{lw}\left(X(f')\right)=0$ and all moments of $X(f')$ are finite, for $f'\in\mathcal{S}$.
So the joint cumulants of order four of (n5) are the forms $\kappa_4'$ of \cite[Section~9.3]{QEDI}, by the last display of 0.1 (Cumulants), and $\mathrm{Ext}^{\mathrm{lw}}_{x_0}$ is the function $\mathrm{Ext}'_{x_0}$ of \cite[Section~9.3]{QEDI}.
Now \cite[Theorem~9.21(i)]{QEDI} is applied to the law $\mathrm{lw}$, with $\bar{m}$ in the place of its $m$, $r_\beta$ in the place of its $r$, and the data $\beta$, $a=d^{*}\beta$, $f=dd^{*}\beta$, $x_0$ of Proposition~43.2.
Its hypotheses are satisfied:
\begin{itemize}
\item[(1)]
$\mathrm{lw}$ satisfies (H1) and is invariant under translations.
The first sentence of \cite[Theorem~9.21]{QEDI} names an ergodic decomposition of $\mathrm{lw}$ in the sense of \cite[Definition~9.5]{QEDI}; the decomposition of Proposition~45.6\,(a$''$) is one, by \cite[Definition~9.5]{QEDI} and \cite[Remark~9.6]{QEDI}; the clause (i) does not involve it.
\item[(2)]
$\beta$ is a real element of $C_c^\infty(\mathbb{R}^3;\Lambda^2)$ whose support lies in a ball of radius $r_\beta$, and $b(f)>0$ and $d_0=|x_0|_2-2r_\beta>0$ (Section~I.7.6; the Data of Proposition~43.2).
The inequality on $x_0$ that \cite[Theorem~9.21]{QEDI} requires is, letter by letter, the inequality of \cite[Section~2.8]{QEDI}, which is the inequality in the Data of Proposition~43.2.
\item[(3)]
\emph{The numbers $b(\cdot)$.}
In \cite[Section~9.7.1]{QEDI} one has
\[
b(dd^{*}\beta'):=-6\operatorname{Tr}\left[\left(S_m\,i\gamma\cdot d^{*}\beta'\right)^4\right]
\]
for every real $\beta'\in C_c^\infty(\mathbb{R}^3;\Lambda^2)$.
For $\beta'=\beta_{\mathrm{ix}}$, where $\beta_0=\beta$, $\beta_\pm=\beta\pm\tau_{x_0}\beta$ and $g_{\mathrm{ix}}=dd^{*}\beta_{\mathrm{ix}}$, this is the number $b(g_{\mathrm{ix}})$ of Section~I.7.6, because the objects of the two definitions agree:
$\gamma_\mu=\operatorname{diag}(\sigma_\mu,-\sigma_\mu)$ in Section~I.2.1 and in \cite[Section~9.7.1]{QEDI};
$S_{\bar{m}}=(\bar{m}+\gamma\cdot\partial)^{-1}$ on $L^2(\mathbb{R}^3;\mathbb{C}^4)$ in Section~I.2.6 and there;
$(d^{*}\beta')_\nu=-\sum_\mu\partial_\mu\beta'_{\mu\nu}$ in Section~I.7.6 and there.
\item[(4)]
The condition $(\kappa_4;\,g)$ of \cite[Section~9.2]{QEDI} holds at the three tests $f=g_0$, $f+\tau_{x_0}f=g_+$, $f-\tau_{x_0}f=g_-$, with the exponent $\sigma>0$, the numbers $b(g_{\mathrm{ix}})$ and the three constants $C_0=C_+=C_-=C'_{\bar{\kappa}}$: this is the hypothesis of the lemma.
\item[(5)]
$C_*=2C_0+\frac14(C_++C_-)=\frac52\cdot C'_{\bar{\kappa}}$, and the requirement $C_*\cdot e^{\sigma}\le\frac{b(f)}{4}$ is (0.3).
\end{itemize}
The conclusion (i) is $\mathrm{Ext}^{\mathrm{lw}}_{x_0}(f)\ge\frac12e^4\cdot b(f)$.
For the second assertion, the hypothesis at $\mathrm{ix}=0$ and (0.3) give
\[
\kappa_4^{\mathrm{lw}}(f)\ge e^4\cdot b(f)-C'_{\bar{\kappa}}\cdot e^{4+\sigma}\ge e^4\cdot b(f)-\frac{1}{10}e^4\cdot b(f),
\]
and $b(f)>0$.
\end{proof}

\paragraph{3.3 (B3).}
\begin{proof}[Proof of (B3)]
The law $\mu$ has (G) by (P1) and is invariant under $\mathbb{R}^3$ by (P5), and (3.1) holds at the three tests.
Lemma~3.2 for $\mathrm{lw}=\mu$ gives
\[
\kappa_4^\mu(f)\ge\frac12e^4\cdot b(f)>0\qquad\text{and}\qquad\mathrm{Ext}^\mu_{x_0}(f)\ge\frac12e^4\cdot b(f).
\]

\emph{$\mu$ is not Gaussian.}
A law is Gaussian if every $X(f')$, $f'\in\mathcal{S}$, is a normal variable, possibly degenerate \cite[Section~9.3]{QEDI}.
If $\mu$ were Gaussian, $X(f)$ would be a normal variable with mean $0$ (by (P3), $S_1=0$), and the fourth cumulant of a normal variable vanishes; but $\kappa_4^\mu(f)>0$.

\emph{$\mu$ is not a mixture of translation invariant Gaussian laws with (G).}
Suppose that $\mu=\int\mathrm{d}\,\mathrm{mx}(\mathfrak{l})\,\gamma_{\mathfrak{l}}$ in the sense of (n6), where every $\gamma_{\mathfrak{l}}$ is Gaussian, invariant under $\mathbb{R}^3$ and has (G).
All tests in this paragraph are real elements of $\mathcal{S}$.
\begin{itemize}
\item[(1)]
By monotone classes, $\mathfrak{l}\mapsto\gamma_{\mathfrak{l}}(\Phi)$ is measurable and $\mu(\Phi)=\int\mathrm{d}\,\mathrm{mx}(\mathfrak{l})\,\gamma_{\mathfrak{l}}(\Phi)$ for every measurable $\Phi\ge0$.
Hence the same holds for every real measurable $\Phi$ with $\sup_{\mathfrak{l}}\gamma_{\mathfrak{l}}(|\Phi|)<\infty$: apply the first case to $\max(\Phi,0)$ and to $\max(-\Phi,0)$.
\item[(2)]
By (G) and \cite[Proposition~7.10(2), (4)]{QEDI}, applied to the law $\gamma_{\mathfrak{l}}$, the space $\mathcal{S}$ and the norm $\|\cdot\|_2$, one has $\gamma_{\mathfrak{l}}\left(X(f_1)\right)=0$ and
\[
\gamma_{\mathfrak{l}}\left(\prod_{a\le\mathrm{ns}}|X(f_a)|\right)\le\sigma_{\mathrm{ns}}\cdot\prod_{a\le\mathrm{ns}}\|f_a\|_2
\]
for every $\mathfrak{l}$, so that (1) applies to the products of two and of four of the variables $X(f_a)$.
Moreover $\mathrm{Cv}_{\mathfrak{l}}(f_1,f_2):=\gamma_{\mathfrak{l}}\left(X(f_1)\cdot X(f_2)\right)$ satisfies $|\mathrm{Cv}_{\mathfrak{l}}(f_1,f_2)|\le\|f_1\|_2\cdot\|f_2\|_2$, by the Cauchy--Schwarz inequality and $\gamma_{\mathfrak{l}}\left(X(f_1)^2\right)\le\|f_1\|_2^2$.
By (1), the function $\mathfrak{l}\mapsto\mathrm{Cv}_{\mathfrak{l}}(f_1,f_2)$ is measurable and bounded, and
\[
S_2^\mu(f_1,f_2)=\int\mathrm{d}\,\mathrm{mx}(\mathfrak{l})\,\mathrm{Cv}_{\mathfrak{l}}(f_1,f_2).
\]
\item[(3)]
Every real linear combination of $X(f_1),\dots,X(f_4)$ is of the form $X(f')$, hence normal under $\gamma_{\mathfrak{l}}$.
So these four variables are jointly normal and centred under $\gamma_{\mathfrak{l}}$, and Wick's rule gives
\begin{multline*}
\gamma_{\mathfrak{l}}\left(X(f_1)\cdot X(f_2)\cdot X(f_3)\cdot X(f_4)\right)\\
=\mathrm{Cv}_{\mathfrak{l}}(f_1,f_2)\cdot\mathrm{Cv}_{\mathfrak{l}}(f_3,f_4)+\mathrm{Cv}_{\mathfrak{l}}(f_1,f_3)\cdot\mathrm{Cv}_{\mathfrak{l}}(f_2,f_4)+\mathrm{Cv}_{\mathfrak{l}}(f_1,f_4)\cdot\mathrm{Cv}_{\mathfrak{l}}(f_2,f_3).
\end{multline*}
\item[(4)]
The law $\mu$ is centred.
So, by 0.1,
\begin{multline*}
\kappa_4^\mu(f_1,\dots,f_4)=S_4^\mu(f_1,\dots,f_4)-S_2^\mu(f_1,f_2)\cdot S_2^\mu(f_3,f_4)\\
-S_2^\mu(f_1,f_3)\cdot S_2^\mu(f_2,f_4)-S_2^\mu(f_1,f_4)\cdot S_2^\mu(f_2,f_3).
\end{multline*}
Insert (1), (2) and (3) on the right.
With $\mathrm{Cov}_{\mathrm{mx}}$ and $\mathrm{Var}_{\mathrm{mx}}$ the covariance and the variance of bounded real functions of $\mathfrak{l}$ under $\mathrm{mx}$, this gives
\begin{multline*}
\kappa_4^\mu(f_1,f_2,f_3,f_4)=\mathrm{Cov}_{\mathrm{mx}}\left(\mathrm{Cv}_{\mathfrak{l}}(f_1,f_2),\mathrm{Cv}_{\mathfrak{l}}(f_3,f_4)\right)\\
+\mathrm{Cov}_{\mathrm{mx}}\left(\mathrm{Cv}_{\mathfrak{l}}(f_1,f_3),\mathrm{Cv}_{\mathfrak{l}}(f_2,f_4)\right)+\mathrm{Cov}_{\mathrm{mx}}\left(\mathrm{Cv}_{\mathfrak{l}}(f_1,f_4),\mathrm{Cv}_{\mathfrak{l}}(f_2,f_3)\right).
\end{multline*}
\item[(5)]
Since $\gamma_{\mathfrak{l}}$ is invariant under $\tau_{x_0}$, one has $\mathrm{Cv}_{\mathfrak{l}}(\tau_{x_0}f,\tau_{x_0}f)=\mathrm{Cv}_{\mathfrak{l}}(f,f)$.
So
\[
\kappa_4^\mu(f)=3\,\mathrm{Var}_{\mathrm{mx}}\left(\mathrm{Cv}_{\mathfrak{l}}(f,f)\right),\qquad\kappa_4^\mu(f,f,\tau_{x_0}f,\tau_{x_0}f)=\mathrm{Var}_{\mathrm{mx}}\left(\mathrm{Cv}_{\mathfrak{l}}(f,f)\right)+2\,\mathrm{Var}_{\mathrm{mx}}\left(\mathrm{Cv}_{\mathfrak{l}}(f,\tau_{x_0}f)\right).
\]
Hence
\[
\mathrm{Ext}^\mu_{x_0}(f)=-6\,\mathrm{Var}_{\mathrm{mx}}\left(\mathrm{Cv}_{\mathfrak{l}}(f,\tau_{x_0}f)\right)\le0,
\]
which contradicts $\mathrm{Ext}^\mu_{x_0}(f)\ge\frac12e^4\cdot b(f)>0$.
\end{itemize}
\end{proof}

\subsubsection*{4. (B4)}

\paragraph{4.1. The numbers $c^{\mathrm{inf}}$, $\mathsf{u}^{\mathrm{inf}}$.}
For every $N\ge N_0(e)$ and every $\mathrm{ix}$ fix one placement $v$ of $\beta_{\mathrm{ix}}$ at the cutoff $N$; such a placement exists by 1.2.
For $i\ge i_0$, $m\in\{2,4\}$ and $\mathrm{ix}\in\{0,+,-\}$ let $c_{m,\mathrm{ix}}(i)$, $\mathsf{u}_{m,\mathrm{ix}}(i)$ and $\boldsymbol{\sigma}_{\mathrm{ix}}(i)^2$ denote the numbers $c_m$, $\mathsf{u}_m$ and $\boldsymbol{\sigma}^2$ of Proposition~43.2 at $N=N_i$, formed with the probe $\beta_{\mathrm{ix}}$ and with this placement.
By the last paragraph of Proposition~43.2 these numbers are real.
Moreover $\mathsf{u}_{m,\mathrm{ix}}(i)\ge0$, by (U-b) for a set $\mathrm{Ypt}$ consisting of one point.
From (U-c), (0.1) and $\eta_i\le1$ one obtains
\begin{gather*}
0\le\mathsf{u}_{4,\mathrm{ix}}(i)\le C_{U}\cdot e^{\sigma},\qquad
0\le\mathsf{u}_{2,\mathrm{ix}}(i)\le C_{U}\cdot e^{\sigma},\qquad
\left|c_{2,\mathrm{ix}}(i)\right|\le C_{U}\cdot e^{-\mathfrak{e}_2},\\
\left|c_{4,\mathrm{ix}}(i)\right|\le\left|b(g_{\mathrm{ix}})\right|+3\cdot C_{U}^2\cdot e^{-2\mathfrak{e}_2}+2\cdot C_{U}.
\end{gather*}
One thus has twelve bounded sequences of real numbers.
By the theorem of Bolzano and Weierstrass the sequence $(N_i,T_i)$ has a subsequence along which all twelve converge, and this subsequence is again a defining sequence of $\mu$.
From now on $(N_i,T_i)$ denotes this subsequence, and
\[
c^{\mathrm{inf}}_{m,\mathrm{ix}}:=\lim_{i\to\infty}c_{m,\mathrm{ix}}(i),\qquad
\mathsf{u}^{\mathrm{inf}}_{m,\mathrm{ix}}:=\lim_{i\to\infty}\mathsf{u}_{m,\mathrm{ix}}(i)\ge0,\qquad
\mathrm{Wk}_{m,\mathrm{ix}}:=\mathrm{Pw}_m\left(X(g_{\mathrm{ix}});\|g_{\mathrm{ix}}\|_2^2,c^{\mathrm{inf}}_{m,\mathrm{ix}}\right).
\]
The four inequalities of (U-c) hold for every $i$, and $\eta_i\to0$.
Letting $i\to\infty$ in them gives
\begin{equation}\tag{4.1}
\begin{aligned}
&\mathsf{u}^{\mathrm{inf}}_{4,\mathrm{ix}}\le C_{U}\cdot e^{\sigma},\qquad
\left(1+\left|c^{\mathrm{inf}}_{2,\mathrm{ix}}\right|\right)\cdot\mathsf{u}^{\mathrm{inf}}_{2,\mathrm{ix}}\le C_{U}\cdot e^{\sigma},\\
&\left|c^{\mathrm{inf}}_{2,\mathrm{ix}}\right|\le C_{U}\cdot e^{-\mathfrak{e}_2},\qquad
\left|c^{\mathrm{inf}}_{4,\mathrm{ix}}-3\cdot\left(c^{\mathrm{inf}}_{2,\mathrm{ix}}\right)^2-b(g_{\mathrm{ix}})\right|\le C_{U}\cdot e^{\sigma}.
\end{aligned}
\end{equation}

\paragraph{4.2. The polynomials.}
For real numbers $z$, $z'$, $\mathrm{cf}$ one has, by Section~I.7.6,
\[
\mathrm{Pw}_m(z;z',\mathrm{cf})=\sum_{\mathrm{dq}\in\{0,2,\dots,m\}}\mathrm{pw}_{m,\mathrm{dq}}(z',\mathrm{cf})\cdot z^{\mathrm{dq}},
\]
where
\begin{gather*}
\mathrm{pw}_{2,2}:=e^{-2},\qquad\mathrm{pw}_{2,0}:=-e^{-2}\cdot z'-\mathrm{cf};\\
\mathrm{pw}_{4,4}:=e^{-4},\qquad\mathrm{pw}_{4,2}:=-6\cdot e^{-4}\cdot z',\qquad\mathrm{pw}_{4,0}:=3\cdot e^{-4}\cdot z'^2-\mathrm{cf}.
\end{gather*}
Each $\mathrm{pw}_{m,\mathrm{dq}}$ is a polynomial in $(z',\mathrm{cf})$ whose coefficients depend on $e$ alone.
Let $\mathrm{Ypt}$ be a finite non-empty set and let $(\mathrm{zv}_y)_{y\in\mathrm{Ypt}}$ be real numbers.
Expanding the product, one finds
\begin{equation}\tag{4.2}
\prod_{y\in\mathrm{Ypt}}\mathrm{Pw}_m(\mathrm{zv}_y;z',\mathrm{cf})
=\sum_{\mathrm{dg}}\left[\prod_{y\in\mathrm{Ypt}}\mathrm{pw}_{m,\mathrm{dg}(y)}(z',\mathrm{cf})\right]\cdot\prod_{y\in\mathrm{Ypt}}\mathrm{zv}_y^{\mathrm{dg}(y)},
\end{equation}
where $\mathrm{dg}$ ranges over the $\left(\frac{m}{2}+1\right)^{\#|\mathrm{Ypt}|}$ functions $\mathrm{Ypt}\to\{0,2,\dots,m\}$.

\paragraph{4.3. The limit of (U-b).}
\begin{proof}[Proof of (B4)]
Fix $m$, $\mathrm{ix}$ and a finite non-empty subset $\mathrm{Ypt}$ of $\left(R^{*}(e)/L\right)\mathbb{Z}^3$.
Let
\[
r_{\mathrm{ix}}:=|y_{\beta_{\mathrm{ix}}}|_\infty+D_{\beta_{\mathrm{ix}}}/2;
\]
this is the number $r_g$ of Proposition~43.2.
The test $g_{\mathrm{ix}}$, which is supported in $\operatorname{supp}\beta_{\mathrm{ix}}$, and the test $g^{(\eta)}_{\mathrm{ix}}$ of Proposition~47.2\,(e) are real elements of $C_c^{\infty}(\mathbb{R}^3;\Lambda^2)$ with supports in $Q_{\beta_{\mathrm{ix}}}$, which is a subset of $[-r_{\mathrm{ix}},r_{\mathrm{ix}}]^3$, and
\begin{equation}\tag{4.3}
\left\|g^{(\eta)}_{\mathrm{ix}}-g_{\mathrm{ix}}\right\|_2\le\eta^2\cdot\|\beta_{\mathrm{ix}}\|_{H^4},\qquad
\|g_{\mathrm{ix}}\|_2\le s_{B},\qquad
\left\|g^{(\eta)}_{\mathrm{ix}}\right\|_2\le\|g_{\mathrm{ix}}\|_2+\|\beta_{\mathrm{ix}}\|_{H^4}\le s_{B},
\end{equation}
by Proposition~47.2\,(e), by $\eta\le1$ and by the definition of $s_{B}$.
Since $R_i\to\infty$, the index $i_0$ may be replaced by an index such that, in addition,
\[
R_i>2\cdot\left(\max_{y\in\mathrm{Ypt}}|y|_\infty+r_{\mathrm{ix}}\right)\qquad\text{for all }i\ge i_0.
\]
Let $i\ge i_0$.
\begin{itemize}
\item[(1)]
Let $y$ and $y'$ be two different points of $\mathrm{Ypt}$.
Then $0<|y-y'|_\infty\le2\cdot\max_{y''\in\mathrm{Ypt}}|y''|_\infty<R_i$, whereas every point of $R_i\mathbb{Z}^3$ other than $0$ has sup norm at least $R_i$.
Hence $y-y'\notin R_i\mathbb{Z}^3$: the points of $\mathrm{Ypt}$ are pairwise incongruent modulo $R_i\mathbb{Z}^3$.
Consequently (U-b) holds for $\mathrm{Ypt}$ at the cutoff $(N_i,T_i)$:
\[
\left|E_i\prod_{y\in\mathrm{Ypt}}\mathrm{Pw}_m\left(X_i\left(\tau_yg^{(\eta_i)}_{\mathrm{ix}}\right);\boldsymbol{\sigma}_{\mathrm{ix}}(i)^2,c_{m,\mathrm{ix}}(i)\right)\right|\le\mathsf{u}_{m,\mathrm{ix}}(i)^{\#|\mathrm{Ypt}|}.
\]
\item[(2)]
For $y\in\mathrm{Ypt}$ the supports of $\tau_yg_{\mathrm{ix}}$, of $\tau_yg^{(\eta_i)}_{\mathrm{ix}}$ and of $\tau_y\left(g^{(\eta_i)}_{\mathrm{ix}}-g_{\mathrm{ix}}\right)$ lie in $y+[-r_{\mathrm{ix}},r_{\mathrm{ix}}]^3$, which is a subset of $(-R_i/2,R_i/2)^3$.
By Lemma~2.1\,(b), the seminorm $f'\mapsto|f'_{\eta_i}|_{\eta_i}$ of each of these three tests is therefore at most its norm in $L^2$, and translations leave this norm unchanged.
\item[(3)]
For a function $\mathrm{dg}\colon\mathrm{Ypt}\to\{0,2,\dots,m\}$ put $\mathrm{ns}:=\sum_{y}\mathrm{dg}(y)\le m\cdot\#|\mathrm{Ypt}|$ and
\begin{gather*}
\mathrm{Mq}_i(\mathrm{dg}):=E_i\prod_{y}X_i\left(\tau_yg^{(\eta_i)}_{\mathrm{ix}}\right)^{\mathrm{dg}(y)},\qquad
\mathrm{Mq}'_i(\mathrm{dg}):=E_i\prod_{y}X_i\left(\tau_yg_{\mathrm{ix}}\right)^{\mathrm{dg}(y)},\\
\mathrm{Mq}(\mathrm{dg}):=\mu\left(\prod_{y}X\left(\tau_yg_{\mathrm{ix}}\right)^{\mathrm{dg}(y)}\right).
\end{gather*}
By (4.2), the left side of the inequality in (1) equals
\[
\left|\sum_{\mathrm{dg}}\left[\prod_{y}\mathrm{pw}_{m,\mathrm{dg}(y)}\left(\boldsymbol{\sigma}_{\mathrm{ix}}(i)^2,c_{m,\mathrm{ix}}(i)\right)\right]\cdot\mathrm{Mq}_i(\mathrm{dg})\right|.
\]
\item[(4)]
One has
\[
\left|\mathrm{Mq}_i(\mathrm{dg})-\mathrm{Mq}'_i(\mathrm{dg})\right|\le\mathrm{ns}\cdot\sigma_{\mathrm{ns}}\cdot s_{B}^{\mathrm{ns}-1}\cdot\|\beta_{\mathrm{ix}}\|_{H^4}\cdot\eta_i^2.
\]
To see this, write the $\mathrm{ns}$ factors of the first product as $X_i(\mathrm{fq}_1),\dots,X_i(\mathrm{fq}_{\mathrm{ns}})$, each test $\tau_yg^{(\eta_i)}_{\mathrm{ix}}$ being listed $\mathrm{dg}(y)$ times, and the factors of the second product, in the same order, as $X_i(\mathrm{fq}'_1),\dots,X_i(\mathrm{fq}'_{\mathrm{ns}})$.
Since $X_i$ is linear in the test,
\[
\prod_{a\le\mathrm{ns}}X_i(\mathrm{fq}_a)-\prod_{a\le\mathrm{ns}}X_i(\mathrm{fq}'_a)
=\sum_{\mathrm{bq}=1}^{\mathrm{ns}}\left[\prod_{a<\mathrm{bq}}X_i(\mathrm{fq}'_a)\right]\cdot X_i\left(\mathrm{fq}_{\mathrm{bq}}-\mathrm{fq}'_{\mathrm{bq}}\right)\cdot\left[\prod_{a>\mathrm{bq}}X_i(\mathrm{fq}_a)\right].
\]
By Lemma~2.1\,(c), by (2) and by (4.3), the expectation under $E_i$ of the modulus of each of the $\mathrm{ns}$ terms on the right is at most
\[
\sigma_{\mathrm{ns}}\cdot s_{B}^{\mathrm{ns}-1}\cdot\left\|g^{(\eta_i)}_{\mathrm{ix}}-g_{\mathrm{ix}}\right\|_2\le\sigma_{\mathrm{ns}}\cdot s_{B}^{\mathrm{ns}-1}\cdot\|\beta_{\mathrm{ix}}\|_{H^4}\cdot\eta_i^2.
\]
\item[(5)]
As $i\to\infty$ one has $\mathrm{Mq}'_i(\mathrm{dg})\to\mathrm{Mq}(\mathrm{dg})$; this is (P8) for the finitely many tests $\tau_yg_{\mathrm{ix}}$, $y\in\mathrm{Ypt}$, which lie in $\mathcal{S}$.
Together with (4) and $\eta_i\to0$ it follows that $\mathrm{Mq}_i(\mathrm{dg})\to\mathrm{Mq}(\mathrm{dg})$.
\item[(6)]
By Proposition~47.2\,(e), $\boldsymbol{\sigma}_{\mathrm{ix}}(i)^2\to\|g_{\mathrm{ix}}\|_2^2$; by the definition of $c^{\mathrm{inf}}_{m,\mathrm{ix}}$ and $\mathsf{u}^{\mathrm{inf}}_{m,\mathrm{ix}}$ in 4.1, $c_{m,\mathrm{ix}}(i)\to c^{\mathrm{inf}}_{m,\mathrm{ix}}$ and $\mathsf{u}_{m,\mathrm{ix}}(i)\to\mathsf{u}^{\mathrm{inf}}_{m,\mathrm{ix}}$.
The functions $\mathrm{pw}_{m,\mathrm{dq}}$ are continuous, and the sum over $\mathrm{dg}$ is finite.
Hence, with (5), the inequality of (1), in the form given in (3), survives the limit $i\to\infty$:
\[
\left|\sum_{\mathrm{dg}}\left[\prod_{y}\mathrm{pw}_{m,\mathrm{dg}(y)}\left(\|g_{\mathrm{ix}}\|_2^2,c^{\mathrm{inf}}_{m,\mathrm{ix}}\right)\right]\cdot\mathrm{Mq}(\mathrm{dg})\right|\le\left(\mathsf{u}^{\mathrm{inf}}_{m,\mathrm{ix}}\right)^{\#|\mathrm{Ypt}|}.
\]
\item[(7)]
The operator $U(y)=\hat\tau_y$ is linear and multiplicative on the functions on $\mathcal{S}'$, it fixes the constants, and $U(y)X(f')=X(\tau_y\cdot f')$ \cite[Section~8.1.2]{QEDI}.
Therefore
\[
U(y)\mathrm{Wk}_{m,\mathrm{ix}}=\mathrm{Pw}_m\left(X(\tau_yg_{\mathrm{ix}});\|g_{\mathrm{ix}}\|_2^2,c^{\mathrm{inf}}_{m,\mathrm{ix}}\right),
\]
and by (4.2) the sum inside the modulus in (6) is $\mu\left(\prod_{y\in\mathrm{Ypt}}U(y)\mathrm{Wk}_{m,\mathrm{ix}}\right)$; all monomials that occur are integrable, by (P2).
It follows that
\begin{equation}\tag{4.4}
\left|\mu\left(\prod_{y\in\mathrm{Ypt}}U(y)\mathrm{Wk}_{m,\mathrm{ix}}\right)\right|\le\left(\mathsf{u}^{\mathrm{inf}}_{m,\mathrm{ix}}\right)^{\#|\mathrm{Ypt}|}
\qquad\text{for every finite subset $\mathrm{Ypt}$ of $\left(R^{*}(e)/L\right)\mathbb{Z}^3$;}
\end{equation}
for the empty set both sides equal $1$, the left side being the integral of the empty product and the right side being read with $0^0:=1$.
\end{itemize}
The inequalities (4.1) and (4.4) together are (B4).
\end{proof}

\subsubsection*{5. (B5)}

The law $\mu$ is invariant under $\mathbb{R}^3$, by (P5).
Let $\nu$ be the measure of (n7).
By Proposition~45.6\,(a$''$) it is the image of $\mu$ under the map $\omega\mapsto\mu_\omega$, where $(\omega,\mathrm{Ev})\mapsto\mu_\omega(\mathrm{Ev})$ is the kernel of Proposition~45.6\,(a), and it is carried by the $\mathbb{R}^3$-ergodic laws.
Let $I$ be the $\sigma$-algebra of Proposition~45.6, that is, the $\sigma$-algebra of the events $\mathrm{Ev}\in\Sigma$ with $\tau_a\mathrm{Ev}=\mathrm{Ev}$ for every $a\in\mathbb{R}^3$.
The numbers $c^{\mathrm{inf}}_{m,\mathrm{ix}}$, $\mathsf{u}^{\mathrm{inf}}_{m,\mathrm{ix}}$ and the variables $\mathrm{Wk}_{m,\mathrm{ix}}$ in (c3) belong to the family that is fixed in Theorem~44 after (B4).
This family may be any family with (4.1) and (4.4), for instance the one constructed in 4.1, which is determined by a subsequence of the defining sequence of $\mu$ and by one placement for each pair $(N,\mathrm{ix})$.
Of the family, only the following properties enter: in 5.2, that each $c^{\mathrm{inf}}_{m,\mathrm{ix}}$ is real, that $\mathsf{u}^{\mathrm{inf}}_{m,\mathrm{ix}}\ge0$, and (4.4); in 5.3, the inequalities (4.1) and the statements (c3).
Therefore (B5) holds for every such family.

\paragraph{5.1. (c1).}
\begin{proof}
The first statement: $\nu$ is carried by the $\mathbb{R}^3$-ergodic laws, by Proposition~45.6\,(a$''$), and an $\mathbb{R}^3$-ergodic law is invariant under $\mathbb{R}^3$ by its definition in (n3).
The second statement: $\mu$ has (G) by (P1), and Proposition~45.6\,(c) applies.
The third statement: the event $\{dF=0\text{ in }\mathcal{S}'\}$ belongs to $\Sigma$ \cite[Section~8.1.6]{QEDI}, and its $\mu$-measure is $1$ by (P4); Proposition~45.6\,(c$'$) applies.
\end{proof}

\paragraph{5.2. (c3).}
\begin{proof}
Fix $m\in\{2,4\}$ and $\mathrm{ix}$.
The variable $\mathrm{Wk}_{m,\mathrm{ix}}$ is a polynomial with real coefficients in $X(g_{\mathrm{ix}})$.
By \cite[Proposition~7.10(3)]{QEDI}, which holds under $\mu$ by 2.4 and (P1), it belongs to $L^p(\mu)$ for every $p<\infty$.
\begin{itemize}
\item[(1)]
For $z\in\mathbb{Z}^3$ the map $\omega\mapsto\tau_{(R^{*}(e)/L)z}\,\omega$ of $\mathcal{S}'$ is measurable and preserves $\mu$, by (P5).
These maps form an action of $\mathbb{Z}^3$, and by (n2)
\[
\mathrm{Ga}\circ\tau_{-(R^{*}(e)/L)z}=U\left(\left(R^{*}(e)/L\right)z\right)\mathrm{Ga}.
\]
By (4.4), for every finite set of pairwise distinct points $z$ of $\mathbb{Z}^3$ the modulus of the integral of the product of the functions $U\left(\left(R^{*}(e)/L\right)z\right)\mathrm{Wk}_{m,\mathrm{ix}}$ is at most $\mathsf{u}^{\mathrm{inf}}_{m,\mathrm{ix}}$ raised to the number of these points.
These are the hypotheses of \cite[Lemma~21.33]{QEDI} for an action of $\mathbb{Z}^3$, that is, in dimension three, with the probability space $(\mathcal{S}',\Sigma,\mu)$, the function $\mathrm{Wk}_{m,\mathrm{ix}}$ and the number $\mathsf{u}^{\mathrm{inf}}_{m,\mathrm{ix}}\ge0$.
Let $P^{\mathrm{lat}}_{\mu}$ be the orthogonal projection of $L^2(\mu)$ onto
\[
\left\{\mathrm{Gb}:U(y)\mathrm{Gb}=\mathrm{Gb}\text{ for all }y\in\left(R^{*}(e)/L\right)\mathbb{Z}^3\right\}.
\]
The lemma gives
\[
\left|P^{\mathrm{lat}}_{\mu}\mathrm{Wk}_{m,\mathrm{ix}}\right|\le\mathsf{u}^{\mathrm{inf}}_{m,\mathrm{ix}}\qquad\text{$\mu$-almost everywhere.}
\]
\item[(2)]
A vector that is invariant under all $U(a)$, $a\in\mathbb{R}^3$, is invariant in particular under the operators $U(y)$, $y\in\left(R^{*}(e)/L\right)\mathbb{Z}^3$.
Hence the range of $P^{123}_{\mu}$ is contained in the range of $P^{\mathrm{lat}}_{\mu}$, and $P^{123}_{\mu}=P^{123}_{\mu}P^{\mathrm{lat}}_{\mu}$.
\item[(3)]
One has $P^{123}_{\mu}=E_\mu\left[\,\cdot\mid I\,\right]$, by \cite[Lemma~8.13(iv)]{QEDI} with the index set $\{1,2,3\}$; here $I$ is the $\sigma$-algebra $\mathcal{I}^{123}$ of \cite[Section~8.1.4]{QEDI}.
With (1) and (2) this gives
\[
\left|E_\mu\left[\mathrm{Wk}_{m,\mathrm{ix}}\mid I\right]\right|=\left|E_\mu\left[P^{\mathrm{lat}}_{\mu}\mathrm{Wk}_{m,\mathrm{ix}}\mid I\right]\right|\le E_\mu\left[\left|P^{\mathrm{lat}}_{\mu}\mathrm{Wk}_{m,\mathrm{ix}}\right|\mid I\right]\le\mathsf{u}^{\mathrm{inf}}_{m,\mathrm{ix}}\qquad\text{$\mu$-almost surely.}
\]
\item[(4)]
Apply Proposition~45.6\,(d) with $\mathrm{Ga}=\mathrm{Wk}_{m,\mathrm{ix}}$, a real $\Sigma$-measurable function on $\mathcal{S}'$ that is taken as the representative of its class, and with $\mathrm{cf}=\mathsf{u}^{\mathrm{inf}}_{m,\mathrm{ix}}$.
It yields $\left|\mu_\xi\left(\mathrm{Wk}_{m,\mathrm{ix}}\right)\right|\le\mathsf{u}^{\mathrm{inf}}_{m,\mathrm{ix}}$ for $\nu$-almost every $\xi$.
\end{itemize}
\end{proof}

\paragraph{5.3. The four inequalities, for every $\xi$ for which the nine statements hold.}
\begin{proof}
Let $\xi$ be such a point, put $\mathrm{lw}:=\mu_\xi$, and fix $\mathrm{ix}$.
Of the nine statements only the following enter: $\mathrm{lw}$ is invariant under $\mathbb{R}^3$; $\mathrm{lw}$ has (G); and $\left|\mathrm{lw}\left(\mathrm{Wk}_{m,\mathrm{ix}}\right)\right|\le\mathsf{u}^{\mathrm{inf}}_{m,\mathrm{ix}}$ for $m=2,4$ and the three values of $\mathrm{ix}$.
\begin{itemize}
\item[(1)]
By (G), \cite[Proposition~7.10]{QEDI} applies to $\mathrm{lw}$, to the real space $\mathcal{S}$ and to the norm $\|\cdot\|_2$.
Its assertions (2) and (3) give $\mathrm{lw}\left(X(f')\right)=0$ and the finiteness of all moments of $X(f')$, for $f'\in\mathcal{S}$.
Hence, for the real random variable $\mathrm{xv}:=X(g_{\mathrm{ix}})$ under $\mathrm{lw}$, one has by 0.1
\[
S_2^{\xi}(g_{\mathrm{ix}},g_{\mathrm{ix}})=\mathrm{lw}(\mathrm{xv}^2),\qquad\kappa_4^{\xi}(g_{\mathrm{ix}})=\mathrm{lw}(\mathrm{xv}^4)-3\cdot\mathrm{lw}(\mathrm{xv}^2)^2.
\]
\item[(2)]
For every real number $z'$,
\[
\mathrm{lw}(\mathrm{xv}^4)-3\cdot\mathrm{lw}(\mathrm{xv}^2)^2=\mathrm{lw}\left(\mathrm{xv}^4-6\cdot z'\cdot\mathrm{xv}^2+3\cdot z'^2\right)-3\cdot\left(\mathrm{lw}\left(\mathrm{xv}^2-z'\right)\right)^2;
\]
indeed, the right side equals
\[
\mathrm{lw}(\mathrm{xv}^4)-6\cdot z'\cdot\mathrm{lw}(\mathrm{xv}^2)+3\cdot z'^2-3\cdot\mathrm{lw}(\mathrm{xv}^2)^2+6\cdot z'\cdot\mathrm{lw}(\mathrm{xv}^2)-3\cdot z'^2.
\]
\item[(3)]
Put $z':=\|g_{\mathrm{ix}}\|_2^2$ and $\mathrm{dv}_m:=\mathrm{lw}\left(\mathrm{Wk}_{m,\mathrm{ix}}\right)$, so that $|\mathrm{dv}_m|\le\mathsf{u}^{\mathrm{inf}}_{m,\mathrm{ix}}$.
The definition of $\mathrm{Pw}_m$ in Section~I.7.6 gives
\[
e^{-2}\cdot\mathrm{lw}\left(\mathrm{xv}^2-z'\right)=c^{\mathrm{inf}}_{2,\mathrm{ix}}+\mathrm{dv}_2,\qquad
e^{-4}\cdot\mathrm{lw}\left(\mathrm{xv}^4-6\cdot z'\cdot\mathrm{xv}^2+3\cdot z'^2\right)=c^{\mathrm{inf}}_{4,\mathrm{ix}}+\mathrm{dv}_4.
\]
\item[(4)]
\emph{The second moment.}
By (3), (4.1), (0.1) and (0.2),
\[
\left|S_2^{\xi}(g_{\mathrm{ix}},g_{\mathrm{ix}})-\|g_{\mathrm{ix}}\|_2^2\right|=e^2\cdot\left|c^{\mathrm{inf}}_{2,\mathrm{ix}}+\mathrm{dv}_2\right|\le e^2\cdot\left(C_{U}\cdot e^{-\mathfrak{e}_2}+C_{U}\cdot e^{\sigma}\right)\le2\cdot C_{U}\cdot e^{2-\mathfrak{e}_2}\le C'_{\bar{\kappa}}\cdot e^{2-\mathfrak{e}_2}.
\]
\item[(5)]
\emph{The fourth cumulant.}
By (1), (2) and (3),
\[
e^{-4}\cdot\kappa_4^{\xi}(g_{\mathrm{ix}})=c^{\mathrm{inf}}_{4,\mathrm{ix}}+\mathrm{dv}_4-3\cdot\left(c^{\mathrm{inf}}_{2,\mathrm{ix}}+\mathrm{dv}_2\right)^2,
\]
and therefore
\[
e^{-4}\cdot\kappa_4^{\xi}(g_{\mathrm{ix}})-\left(c^{\mathrm{inf}}_{4,\mathrm{ix}}-3\cdot\left(c^{\mathrm{inf}}_{2,\mathrm{ix}}\right)^2\right)=\mathrm{dv}_4-3\cdot\mathrm{dv}_2\cdot\left(\mathrm{dv}_2+2\cdot c^{\mathrm{inf}}_{2,\mathrm{ix}}\right).
\]
By (4.1) and (0.1), and since
\[
\mathsf{u}^{\mathrm{inf}}_{2,\mathrm{ix}}\cdot\left|c^{\mathrm{inf}}_{2,\mathrm{ix}}\right|\le\left(1+\left|c^{\mathrm{inf}}_{2,\mathrm{ix}}\right|\right)\cdot\mathsf{u}^{\mathrm{inf}}_{2,\mathrm{ix}}\le C_{U}\cdot e^{\sigma},\qquad\mathsf{u}^{\mathrm{inf}}_{2,\mathrm{ix}}\le C_{U}\cdot e^{\sigma},
\]
one obtains
\begin{align*}
\left|\mathrm{dv}_4-3\cdot\mathrm{dv}_2\cdot\left(\mathrm{dv}_2+2\cdot c^{\mathrm{inf}}_{2,\mathrm{ix}}\right)\right|
&\le\mathsf{u}^{\mathrm{inf}}_{4,\mathrm{ix}}+6\cdot\mathsf{u}^{\mathrm{inf}}_{2,\mathrm{ix}}\cdot\left|c^{\mathrm{inf}}_{2,\mathrm{ix}}\right|+3\cdot\left(\mathsf{u}^{\mathrm{inf}}_{2,\mathrm{ix}}\right)^2\\
&\le C_{U}\cdot e^{\sigma}+6\cdot C_{U}\cdot e^{\sigma}+3\cdot C_{U}^2\cdot e^{2\sigma}\le\left(7\cdot C_{U}+3\cdot C_{U}^2\right)\cdot e^{\sigma}.
\end{align*}
Combined with the last inequality of (4.1) and with (0.2), this yields
\[
\left|\kappa_4^{\xi}(g_{\mathrm{ix}})-e^4\cdot b(g_{\mathrm{ix}})\right|\le\left(8\cdot C_{U}+3\cdot C_{U}^2\right)\cdot e^{4+\sigma}=\frac{C_{\bar{\kappa}}}{2}\cdot e^{4+\sigma}\le C'_{\bar{\kappa}}\cdot e^{4+\sigma}.
\]
\item[(6)]
The inequality of (5) holds for $\mathrm{ix}=0,+,-$, and $\mathrm{lw}$ has (G) and is invariant under $\mathbb{R}^3$.
Lemma~3.2 therefore gives
\[
\mathrm{Ext}^{\xi}_{x_0}(f)\ge e^4\cdot b(f)/2\qquad\text{and}\qquad\kappa_4^{\xi}(f)\ge e^4\cdot b(f)/2>0.
\]
\end{itemize}
\end{proof}

\subsubsection*{6. (C)}

In this section $e\in(0,e_2]$, and the laws $P_{N,T}$ are those of Theorem~44.
The assertions of 6.1--6.7 hold, with the same proofs, for every $e>0$ and every sequence $m_{0,N}(e)$ with the properties stated in 1.1, except for the last clause of 6.2, namely that (B) holds for $\mu$;
the parts (B2)--(B5) of (B) are proved for $e\le e_2$ by means of Proposition~43.2.

\paragraph{6.1. Existence.}
\begin{proof}
By the first sentence of \cite[Theorem~1.3(b)]{QEDII}, $\mathcal{V}_N$ is not empty for every $N\ge0$, and every sequence $\mu^{(N_i)}\in\mathcal{V}_{N_i}$ with $N_i\to\infty$ has a subsequence that converges weakly on $H_{-4}$.
Choose $\mu^{(N)}\in\mathcal{V}_N$ for every $N$.
Some subsequence of $\left(\mu^{(N)}\right)_N$ converges weakly on $H_{-4}$, and its limit is a volume-first limit point in the sense of Section~I.7.5.
\end{proof}

\paragraph{6.2. A volume-first limit point is a limit point.}
\begin{proof}
(This is the first assertion of \cite[Proposition~12.16]{QEDII}; see also \cite[Proposition~12.13(d)]{QEDII}, the counterpart of \cite[Remark~21.17]{QEDI}.)
Let $\mu=\lim_i\mu^{(N_i)}$ weakly on $H_{-4}$, with $N_i\to\infty$, and let $\mu^{(N_i)}=\lim_jP_{N_i,T_j(N_i)}$ weakly on $H_{-4}$, where $T_j(N_i)\to\infty$ as $j\to\infty$.
On the separable metric space $H_{-4}$, weak convergence of Borel probability measures is metrised by the Prokhorov metric $d_P$ \cite[Theorem~6.8]{Billingsley1999}.
For every $i$ choose $j_i$ with
\[
T_{j_i}(N_i)\ge i\qquad\text{and}\qquad d_P\left(P_{N_i,T_{j_i}(N_i)},\mu^{(N_i)}\right)\le\frac{1}{i}.
\]
Then
\[
d_P\left(P_{N_i,T_{j_i}(N_i)},\mu\right)\le\frac{1}{i}+d_P\left(\mu^{(N_i)},\mu\right)\to0,
\]
and both $N_i$ and $T_{j_i}(N_i)$ tend to infinity.
Hence $\mu$ is a limit point with the defining sequence $\left(N_i,T_{j_i}(N_i)\right)$, and (B) holds for $\mu$ by Sections~2--5.
\end{proof}

From here on, $\mu$ is a volume-first limit point.

\paragraph{6.3. (C1).}
\begin{proof}
By \cite[Proposition~12.16]{QEDII}, $\mu$ lies in $\mathcal{C}_9$.
\end{proof}

\paragraph{6.4. (C2), (C3), (C7).}
\begin{proof}
By (2.1) and (C1),
\begin{equation}\tag{6.1}
\text{$\mu$ satisfies (H) and (dF), and $\mu$ lies in $\mathcal{C}_9$.}
\end{equation}
Now (C2) is \cite[Theorem~1.4, Part~I\,(a)]{QEDII}, (C3) is \cite[Theorem~1.4, Part~I\,(b)]{QEDII}, and (C7) is \cite[Theorem~1.4, Part~I\,(d)]{QEDII}.
(In (C2) and (C7), $S_n(f_1,\dots,f_n)=\mu\left(X(f_1)\cdots X(f_n)\right)$ for $f_a\in L^2(\mathbb{R}^3;\Lambda^2)$, as in (n5).
The integrand $\hat b(p)^{\dagger}P^{\perp}(p)\,\hat b'(p)$ of \cite[Theorem~1.4, Part~I\,(d)]{QEDII} is $\overline{\hat b'(p)}\cdot\Pi(p)\,\hat b''(p)$ in the notation of (C7), where $b'=\ast f'$ and $b''=\ast f''$.
The display of \cite[Theorem~1.4, Part~I\,(d)]{QEDII} holds for all real $f',f''\in L^2(\mathbb{R}^3;\Lambda^2)$.)
\end{proof}

\paragraph{6.5. (C4).}
\emph{The joint spectral measure.}
Let $\mathrm{lw}$ be a law that is invariant under $\mathbb{R}^3$ and under $\theta_j$ and reflection positive at $\{y_j=0\}$.
The plane $e_j^{\perp}$ is identified with $\mathbb{R}^2$ by means of the coordinates $(y_i)_{i\ne j}$; thus $p\cdot y=\sum_{i\ne j}p_i\cdot y_i$ is the scalar product in $\mathbb{R}^3$ of vectors $p,y$ of $e_j^{\perp}$.
By \cite[Lemma~12.18(b), (c)]{QEDII} there is exactly one projection-valued measure $\mathrm{Es}_j=\mathrm{Es}_j(\mathrm{lw})$ on the Borel sets of $\mathbb{R}\times e_j^{\perp}$ such that
\[
\exp\left(it\cdot\mathsf{H}_j\right)V_j(y)=\int\mathrm{d}\mathrm{Es}_j(E,p)\,\exp\left(it\cdot E+ip\cdot y\right)\qquad\left(t\in\mathbb{R},\ y\in e_j^{\perp}\right);
\]
it is carried by $[0,\infty)\times e_j^{\perp}$,
\[
\exp\left(-t\cdot\mathsf{H}_j\right)V_j(y)=\int\mathrm{d}\mathrm{Es}_j(E,p)\,\exp\left(-t\cdot E+ip\cdot y\right)\qquad(t\ge0),
\]
and $\mathsf{P}_i=\int\mathrm{d}\mathrm{Es}_j(E,p)\,p_i$ for $i\ne j$.
The measure $\mathrm{Es}_j$ is the joint spectral measure of $\left(\mathsf{H}_j,(\mathsf{P}_i)_{i\ne j}\right)$.

\begin{named}{Lemma 6.5}
Let $\mathrm{lw}$ satisfy (H) and (dF) and lie in $\mathcal{C}_9$.
Then, for $j=1,2,3$,
\[
\mathrm{Es}_j(\mathrm{lw})\left(\left\{(E,p):E<|p|_2\right\}\right)=0.
\]
\end{named}

\begin{proof}
By (H), $\mathrm{lw}$ is invariant under $\mathbb{R}^3$ and under $\theta_j$ and reflection positive at $\{y_j=0\}$, for $j=1,2,3$; so the three reconstructions (n4) and the three measures $\mathrm{Es}_j=\mathrm{Es}_j(\mathrm{lw})$ are defined.
By \cite[Theorem~1.4, Part~I\,(a), (f)]{QEDII}, $\mathrm{Gi}_*\mathrm{lw}=\mathrm{lw}$ for every Euclidean map $\mathrm{Gi}$; in particular $\mathrm{gr}_*\mathrm{lw}=\mathrm{lw}$ for every $\mathrm{gr}\in\mathrm{O}(3)$.
\begin{itemize}
\item[(1)]
\emph{Transport.}
Let $j\in\{1,2,3\}$, and let $\mathrm{gr}\in\mathrm{O}(3)$ satisfy $\mathrm{gr}\,e_1=e_j$.
For $y\in\mathbb{R}^3$ one has $(\mathrm{gr}\,y)_j=(\mathrm{gr}\,y)\cdot(\mathrm{gr}\,e_1)=y_1$.
Hence $\mathrm{gr}\{y_1>0\}=\{y_j>0\}$ and $\mathrm{gr}\,e_1^{\perp}=e_j^{\perp}$.
Moreover
\[
\mathrm{gr}\,\theta_1\,\mathrm{gr}^{-1}=\mathrm{gr}\left(1-2e_1e_1^T\right)\mathrm{gr}^T=1-2e_je_j^T=\theta_j,
\]
and $\mathrm{gr}\,\tau_a\,\mathrm{gr}^{-1}=\tau_{\mathrm{gr}\,a}$ for $a\in\mathbb{R}^3$.
By \cite[Section~8.1.2]{QEDI}, $\widehat{\mathrm{gr}}$ is unitary on $L^2(\mathrm{lw})$ and maps $L^2\left(\Sigma_{\{y_1>0\}},\mathrm{lw}\right)$ onto $L^2\left(\Sigma_{\{y_j>0\}},\mathrm{lw}\right)$; moreover $\widehat{\mathrm{gr}}\,\hat\theta_1=\hat\theta_j\,\widehat{\mathrm{gr}}$ and $\widehat{\mathrm{gr}}\,U(a)=U(\mathrm{gr}\,a)\,\widehat{\mathrm{gr}}$.
Consequently, for $\mathrm{Ga},\mathrm{Gb}\in L^2\left(\Sigma_{\{y_1>0\}},\mathrm{lw}\right)$,
\[
\left\langle\hat\theta_j\widehat{\mathrm{gr}}\,\mathrm{Ga},\widehat{\mathrm{gr}}\,\mathrm{Gb}\right\rangle_{\mathrm{lw}}=\left\langle\widehat{\mathrm{gr}}\,\hat\theta_1\mathrm{Ga},\widehat{\mathrm{gr}}\,\mathrm{Gb}\right\rangle_{\mathrm{lw}}=\left\langle\hat\theta_1\mathrm{Ga},\mathrm{Gb}\right\rangle_{\mathrm{lw}}.
\]
By (n4), that is, by \cite[Proposition~8.19(i)--(iii)]{QEDI} for the index $j$ and, in its place, for the index $1$:
$\mathcal{H}_j$ is the completion of $L^2\left(\Sigma_{\{y_j>0\}},\mathrm{lw}\right)$ modulo the null vectors of the form $\left\langle\hat\theta_j\,\cdot\,,\,\cdot\,\right\rangle_{\mathrm{lw}}$;
the classes $[\mathrm{Ga}]_j$ are dense in $\mathcal{H}_j$;
$\exp\left(-t\cdot\mathsf{H}_j\right)[\mathrm{Ga}]_j=\left[U(t\cdot e_j)\mathrm{Ga}\right]_j$ for $t\ge0$;
and $V_j(y)[\mathrm{Ga}]_j=\left[U(y)\mathrm{Ga}\right]_j$ for $y\in e_j^{\perp}$.
It follows that $[\mathrm{Ga}]_1\mapsto\left[\widehat{\mathrm{gr}}\,\mathrm{Ga}\right]_j$ is well defined and isometric from a dense subspace of $\mathcal{H}_1$ onto a dense subspace of $\mathcal{H}_j$, and that it extends to a unitary map $\mathrm{Uw}\colon\mathcal{H}_1\to\mathcal{H}_j$.
For $t\ge0$ and $y\in e_1^{\perp}$,
\begin{align*}
\mathrm{Uw}\exp\left(-t\cdot\mathsf{H}_1\right)V_1(y)[\mathrm{Ga}]_1
&=\mathrm{Uw}\left[U(t\cdot e_1+y)\mathrm{Ga}\right]_1
=\left[\widehat{\mathrm{gr}}\,U(t\cdot e_1+y)\mathrm{Ga}\right]_j
=\left[U(t\cdot e_j+\mathrm{gr}\,y)\,\widehat{\mathrm{gr}}\,\mathrm{Ga}\right]_j\\
&=\exp\left(-t\cdot\mathsf{H}_j\right)V_j(\mathrm{gr}\,y)\left[\widehat{\mathrm{gr}}\,\mathrm{Ga}\right]_j
=\exp\left(-t\cdot\mathsf{H}_j\right)V_j(\mathrm{gr}\,y)\,\mathrm{Uw}[\mathrm{Ga}]_1.
\end{align*}
The operators are bounded and the classes are dense; hence
\[
\mathrm{Uw}\exp\left(-t\cdot\mathsf{H}_1\right)V_1(y)\,\mathrm{Uw}^{-1}=\exp\left(-t\cdot\mathsf{H}_j\right)V_j(\mathrm{gr}\,y).
\]
For a Borel subset $\mathrm{Sb}$ of $[0,\infty)\times e_j^{\perp}$ put
\[
\mathrm{Es}'(\mathrm{Sb}):=\mathrm{Uw}\,\mathrm{Es}_1\left(\left\{(E,p):(E,\mathrm{gr}\,p)\in\mathrm{Sb}\right\}\right)\mathrm{Uw}^{-1}.
\]
The map $(E,p)\mapsto(E,\mathrm{gr}\,p)$ is a homeomorphism of $[0,\infty)\times e_1^{\perp}$ onto $[0,\infty)\times e_j^{\perp}$.
Therefore $\mathrm{Es}'$ is a projection-valued measure, and
\[
\int\mathrm{d}\mathrm{Es}'(E,p')\,\varphi(E,p')=\mathrm{Uw}\left(\int\mathrm{d}\mathrm{Es}_1(E,p)\,\varphi(E,\mathrm{gr}\,p)\right)\mathrm{Uw}^{-1}
\]
for every bounded Borel function $\varphi$.
Let $t\ge0$ and $y'\in e_j^{\perp}$.
Since $(\mathrm{gr}\,p)\cdot y'=p\cdot(\mathrm{gr}^{-1}y')$ and $\mathrm{gr}^{-1}y'\in e_1^{\perp}$,
\begin{align*}
\int\mathrm{d}\mathrm{Es}'(E,p')\,\exp\left(-t\cdot E+ip'\cdot y'\right)
&=\mathrm{Uw}\left(\int\mathrm{d}\mathrm{Es}_1(E,p)\,\exp\left(-t\cdot E+ip\cdot(\mathrm{gr}^{-1}y')\right)\right)\mathrm{Uw}^{-1}\\
&=\mathrm{Uw}\exp\left(-t\cdot\mathsf{H}_1\right)V_1(\mathrm{gr}^{-1}y')\,\mathrm{Uw}^{-1}
=\exp\left(-t\cdot\mathsf{H}_j\right)V_j(y')\\
&=\int\mathrm{d}\mathrm{Es}_j(E,p')\,\exp\left(-t\cdot E+ip'\cdot y'\right).
\end{align*}
Consequently, for every $\psi\in\mathcal{H}_j$, the finite positive measures $\left\langle\psi,\mathrm{Es}'(\,\cdot\,)\psi\right\rangle$ and $\left\langle\psi,\mathrm{Es}_j(\,\cdot\,)\psi\right\rangle$ on $[0,\infty)\times\mathbb{R}^2$ have the same integrals of $\exp\left(-t\cdot E+ip'\cdot y'\right)$ for $t=0,1,2,\dots$ and $y'\in\mathbb{R}^2$; by \cite[Lemma~12.25]{QEDII} they coincide.
A projection is determined by its diagonal matrix elements.
Hence
\begin{equation}\tag{6.2}
\mathrm{Es}_j(\mathrm{Sb})=\mathrm{Uw}\,\mathrm{Es}_1\left(\left\{(E,p):(E,\mathrm{gr}\,p)\in\mathrm{Sb}\right\}\right)\mathrm{Uw}^{-1}\qquad\text{for every Borel subset $\mathrm{Sb}$ of $[0,\infty)\times e_j^{\perp}$.}
\end{equation}
\item[(2)]
\emph{The case $j=1$.}
Let $\hat{n}'$ be a unit vector of $e_1^{\perp}$, and let $\mathrm{gr}$ be the element of $\mathrm{O}(3)$ with $\mathrm{gr}\,e_1=e_1$, $\mathrm{gr}\,e_2=\hat{n}'$, and $\mathrm{gr}\,e_3$ one of the two unit vectors that are orthogonal to $e_1$ and to $\hat{n}'$.
Put
\[
\mathrm{Sb}:=\left\{(E,p'):E<|\hat{n}'\cdot p'|\right\}.
\]
For $p\in e_1^{\perp}$ one has $\hat{n}'\cdot(\mathrm{gr}\,p)=(\mathrm{gr}\,e_2)\cdot(\mathrm{gr}\,p)=p_2$; so $\left\{(E,p):(E,\mathrm{gr}\,p)\in\mathrm{Sb}\right\}=\left\{E<|p_2|\right\}$.
By \cite[Proposition~12.32(a)]{QEDII}, $\mathrm{Es}_1$ is carried by $\left\{E\ge|p|_\infty\right\}$, and $|p_2|\le|p|_\infty$; hence $\mathrm{Es}_1\left(\left\{E<|p_2|\right\}\right)=0$.
By (6.2) with $j=1$,
\[
\mathrm{Es}_1\left(\left\{E<|\hat{n}'\cdot p|\right\}\right)=0\qquad\text{for every unit vector $\hat{n}'$ of $e_1^{\perp}$.}
\]
The set $\left\{E<|p|_2\right\}$ is the union of the sets $\left\{E<|\hat{n}'\cdot p|\right\}$ over a countable dense set of unit vectors $\hat{n}'$ of $e_1^{\perp}$.
Indeed, on the one hand $|\hat{n}'\cdot p|\le|p|_2$; on the other hand, if $E<|p|_2$, then $p\ne0$, and $\hat{n}'\mapsto|\hat{n}'\cdot p|$ is continuous and takes the value $|p|_2$ at $\hat{n}'=p/|p|_2$.
For every $\psi\in\mathcal{H}_1$, $\left\langle\psi,\mathrm{Es}_1(\,\cdot\,)\psi\right\rangle$ is a positive measure; so a countable union of sets on which $\mathrm{Es}_1$ vanishes is again a set on which $\mathrm{Es}_1$ vanishes.
Hence $\mathrm{Es}_1\left(\left\{E<|p|_2\right\}\right)=0$.
\item[(3)]
\emph{The cases $j=2,3$.}
Let $\mathrm{gr}\in O_{\mathrm{cub}}$ be the matrix that exchanges $e_1$ and $e_j$ and fixes the third unit vector.
Then $\mathrm{gr}\,e_1=e_j$ and $|\mathrm{gr}\,p|_2=|p|_2$.
By (6.2) for $\mathrm{Sb}=\left\{E<|p'|_2\right\}$ and by (2),
\[
\mathrm{Es}_j\left(\left\{E<|p'|_2\right\}\right)=\mathrm{Uw}\,\mathrm{Es}_1\left(\left\{E<|p|_2\right\}\right)\mathrm{Uw}^{-1}=0.
\]
\end{itemize}
\end{proof}

\begin{proof}[Proof of (C4)]
By (6.1), (C4) is Lemma~6.5 for $\mathrm{lw}=\mu$.
(Lemma~6.5 is the first clause of \cite[Theorem~1.4, Part~I\,(c), (f)]{QEDII}, the first assertion of \cite[Proposition~14.47(b)]{QEDII} in each of the three pictures, and the assertion on the round forward cone in \cite[Theorem~1.4, Part~II\,(E8)]{QEDII}.)
\end{proof}

\paragraph{6.6. (C5), (C6).}
\begin{proof}
The law $\mu$ lies in $\mathcal{C}_9$.
Hence (C5) is \cite[Proposition~12.32(c)]{QEDII}.
Further, (C6) is the last display of \cite[Theorem~12.39]{QEDII}, that is, the second of its two inequalities, the one that is stated under a lower bound on $|a|_\infty$.
The constants are those of that theorem: $c_1=4$ in this lower bound on $|a|_\infty$, and $c_0=3\cdot C_2$ in the inequality, where
\[
C_2=32\cdot(1+\sqrt{2})\cdot C_1,\qquad C_1=\frac{1}{4\pi^2}\cdot\left(128+64\cdot\left(\frac{16}{\pi}+\frac{64}{\pi^2}\right)\right).
\]
Thus
\[
c_0=\frac{96\cdot(1+\sqrt{2})}{4\pi^2}\cdot\left(128+64\cdot\left(\frac{16}{\pi}+\frac{64}{\pi^2}\right)\right)=c_{\mathrm{cl}}.
\]
(In \cite[Theorem~12.39]{QEDII} the numbers written here $c_1$, $c_0$, $C_1$, $C_2$ are denoted $C_1$, $C_0$, $C_2$, $C_3$.)
\end{proof}

\paragraph{6.7. (C8).}

\begin{named}{Lemma 6.7}
Let $\mathrm{lw}$ be a law with (H) and (dF) that is invariant under every Euclidean map of $\mathbb{R}^3$ and satisfies $\left\langle\widehat{r_{\hat{n},0}}\,\mathrm{Ga},\mathrm{Ga}\right\rangle_{\mathrm{lw}}\ge0$ for every unit vector $\hat{n}$ and all $\mathrm{Ga}\in L^2\left(\Sigma_{\{\hat{n}\cdot y>0\}},\mathrm{lw}\right)$.
Then $\mathrm{lw}$ has (E0$'$), (E1), (E2) and (E3).
\end{named}

\begin{proof}
The law $\mathrm{lw}$ has (G); hence $X^{(n)}$ and $S_n(\Psi)=\mathrm{lw}\left(X^{(n)}(\Psi)\right)$ are defined, by Proposition~45.5\,(a) and (n5).
\begin{itemize}
\item[(E0$'$)]
By (n5), $S_0=1$.
The law $\mathrm{lw}$ has (L1)--(L4):
\begin{itemize}
\item[(L1), (L4)]
(L1) is (G), and (L4) is reflection positivity at $\{y_j=0\}$, $j=1,2,3$; these two identifications are made in the notation stated with the Propositions~45.m, with $\theta_j=r_{e_j,0}$ and with $X\left(f_q-\theta_j\cdot f_p\right)=X(f_q)-X(\theta_j\cdot f_p)$, which holds because $X$ is linear in the test.
Both are parts of (H).
\item[(L3)]
Let $\mathrm{Gi}=\tau_a$ with $a\in\mathbb{R}^3$, or $\mathrm{Gi}=r\in O_{\mathrm{cub}}$, and let $f'\in\mathcal{S}$ be real.
Then $\exp\left(iX(\mathrm{Gi}\cdot f')\right)=\widehat{\mathrm{Gi}}\exp\left(iX(f')\right)$, because $\widehat{\mathrm{Gi}}X(f')=X(\mathrm{Gi}\cdot f')$ and $\widehat{\mathrm{Gi}}$ commutes with the composition with a function \cite[Section~8.1.2]{QEDI}.
For a bounded measurable $\mathrm{Ga}$,
\[
\mathrm{lw}\left(\widehat{\mathrm{Gi}}\,\mathrm{Ga}\right)=\int\mathrm{d}\mathrm{lw}(\omega)\,\mathrm{Ga}\left(\mathrm{Gi}^{-1}\omega\right)=\left(\left(\mathrm{Gi}^{-1}\right)_*\mathrm{lw}\right)(\mathrm{Ga})=\mathrm{lw}(\mathrm{Ga}),
\]
since $\mathrm{Gi}^{-1}$ is again a translation, respectively an element of $O_{\mathrm{cub}}$, and $\mathrm{lw}$ is invariant under $\mathbb{R}^3$ and under $O_{\mathrm{cub}}$ by (H).
Hence $\Phi_{\mathrm{lw}}(\mathrm{Gi}\cdot f')=\Phi_{\mathrm{lw}}(f')$.
\item[(L2)]
Let $\mathrm{gs}\in C_c^{\infty}(\mathbb{R}^3)$ be real.
Then $d^{*}\mathrm{gs}$ is an element of $\mathcal{S}$ with $\ast\left(d^{*}\mathrm{gs}\right)=-\nabla\mathrm{gs}=\nabla(-\mathrm{gs})$, and $-\mathrm{gs}\in\mathcal{S}(\mathbb{R}^3)$.
So $\omega\left(d^{*}\mathrm{gs}\right)=0$ for every closed $\omega$, by the definition of a closed $\omega$ in \cite[Section~8.1.6]{QEDI}; and the set $\{dF=0\text{ in }\mathcal{S}'\}$ of the closed $\omega$, see (n8), has $\mathrm{lw}$-measure $1$ by (dF).
Hence $X\left(d^{*}\mathrm{gs}\right)=0$ $\mathrm{lw}$-almost surely.
\end{itemize}
Therefore Proposition~45.2 holds for $\mathrm{lw}$.
By Proposition~45.2\,(2),
\[
\left|\mathrm{lw}(Q_{\Psi})\right|\le2\cdot66^n\cdot(n!)^{\frac{1}{2}}\cdot q_{6n}(\Psi),
\]
and $Q_{\Psi}=X^{(n)}(\Psi)$ in $L^2(\mathrm{lw})$; so $\mathrm{lw}(Q_{\Psi})=S_n(\Psi)$.
\item[(E1)]
Let $\mathrm{Gi}$ be a Euclidean map.
Then $\mathrm{Gi}_*\mathrm{lw}=\mathrm{lw}$ and $\left(\mathrm{Gi}^{-1}\right)_*\mathrm{lw}=\mathrm{lw}$.
By the last clause of Proposition~45.5\,(b), $X^{(n)}\left((\mathrm{Gi}\,\cdot)^{\otimes n}\Psi\right)=\widehat{\mathrm{Gi}}\,X^{(n)}(\Psi)$.
For $\mathrm{Ga}\in L^1(\mathrm{lw})$,
\[
\mathrm{lw}\left(\widehat{\mathrm{Gi}}\,\mathrm{Ga}\right)=\int\mathrm{d}\mathrm{lw}(\omega)\,\mathrm{Ga}\left(\mathrm{Gi}^{-1}\omega\right)=\left(\left(\mathrm{Gi}^{-1}\right)_*\mathrm{lw}\right)(\mathrm{Ga})=\mathrm{lw}(\mathrm{Ga}).
\]
Hence $S_n\left((\mathrm{Gi}\,\cdot)^{\otimes n}\Psi\right)=S_n(\Psi)$.
\item[(E2)]
Let $\hat{n}$ be a unit vector.
The reflection $r_{\hat{n},0}$ is a Euclidean map; so $\left(r_{\hat{n},0}\right)_*\mathrm{lw}=\mathrm{lw}$.
Let $(c_p)$ be a finite family in $\mathbb{C}$ and $(f_p)$ a finite family of real elements of $C_c^{\infty}\left(\{\hat{n}\cdot y>0\};\Lambda^2\right)$, and let $\mathrm{Ga}:=\sum_pc_p\exp\left(iX(f_p)\right)$.
Then $\mathrm{Ga}\in L^2\left(\Sigma_{\{\hat{n}\cdot y>0\}},\mathrm{lw}\right)$, $\widehat{r_{\hat{n},0}}\exp\left(iX(f_p)\right)=\exp\left(iX(r_{\hat{n},0}\cdot f_p)\right)$ by \cite[Section~8.1.2]{QEDI}, and
\[
0\le\left\langle\widehat{r_{\hat{n},0}}\,\mathrm{Ga},\mathrm{Ga}\right\rangle_{\mathrm{lw}}=\sum_{p,q}\overline{c_p}\cdot c_q\cdot\mathrm{lw}\left(\exp\left(i\left(X(f_q)-X(r_{\hat{n},0}\cdot f_p)\right)\right)\right).
\]
So $\mathrm{lw}$ is reflection positive at $\{\hat{n}\cdot y=0\}$ in the sense of (n3), and the second display of Proposition~45.5\,(d) is (E2).
\item[(E3)]
Let $\mathrm{pm}$ be a permutation of $\{1,\dots,n\}$, and let $\Psi\mapsto\Psi^{\mathrm{pm}}$ be the linear map of $\mathrm{Sw}_n$ that permutes the $n$ arguments together with their tensor indices, so that $\left(f_1\otimes\cdots\otimes f_n\right)^{\mathrm{pm}}=f_{\mathrm{pm}(1)}\otimes\cdots\otimes f_{\mathrm{pm}(n)}$.
Every seminorm $q_m$ of (n5) satisfies $q_m(\Psi^{\mathrm{pm}})=q_m(\Psi)$; hence this map is continuous.
The map $\Psi\mapsto X^{(n)}(\Psi^{\mathrm{pm}})$ is a continuous linear map $\mathrm{Sw}_n\to L^2(\mathrm{lw})$ whose value at $f_1\otimes\cdots\otimes f_n$ is
\[
X(f_{\mathrm{pm}(1)})\cdots X(f_{\mathrm{pm}(n)})=X(f_1)\cdots X(f_n),
\]
because the $X(f_a)$ are functions on $\mathcal{S}'$ and their product does not depend on the order of the factors.
By the uniqueness in Proposition~45.5\,(a), $X^{(n)}(\Psi^{\mathrm{pm}})=X^{(n)}(\Psi)$; therefore $S_n(\Psi^{\mathrm{pm}})=S_n(\Psi)$.
\end{itemize}
\end{proof}

\begin{proof}[Proof of (C8)]
By (6.1), by (C2) and by (C3) at $\mathrm{po}=0$, the law $\mu$ satisfies the hypotheses of Lemma~6.7, and (C8) is Lemma~6.7 for $\mathrm{lw}=\mu$.
(Lemma~6.7 is the assertion (E0$'$), (E1), (E2), (E3) of \cite[Theorem~1.4, Part~II]{QEDII}.)
\end{proof}

\subsubsection*{7. (D)}

Throughout this section $\mu$ is a volume-first limit point, and $\nu$ and the laws $\mu_\xi$ are those of Section~5.
By 6.2 the law $\mu$ is a limit point, so that 5.1 and 5.2 apply to it: the statements (c1) and (c3) hold for $\mu$, formed with the family of numbers $c^{\mathrm{inf}}_{m,\mathrm{ix}}$, $\mathsf{u}^{\mathrm{inf}}_{m,\mathrm{ix}}$ that is fixed after (B4).
A family of this kind exists: 4.1 provides one, formed along a subsequence of the defining sequence of $\mu$ constructed in 6.2.
The set $\Xi_\mu$ is formed with this family.
In 7.2 and 7.3 the family enters only through the properties that enter Section~5, namely that the numbers $c^{\mathrm{inf}}_{m,\mathrm{ix}}$ are real, that $\mathsf{u}^{\mathrm{inf}}_{m,\mathrm{ix}}\ge0$, and (4.1), (4.4); consequently (D) holds for every such family.

The nine unit vectors are
\[
\hat{n}=e_j\quad(j=1,2,3),\qquad
\hat{n}=\frac{e_i-e_l}{\sqrt{2}},\qquad
\hat{n}=\frac{e_i+e_l}{\sqrt{2}}\quad(1\le i<l\le3).
\]
For these vectors the planes $\{y_j=\mathrm{po}\}$, $\{y_i-y_l=\mathrm{po}\}$ and $\{y_i+y_l=\mathrm{po}\}$ are the planes
\[
\{\hat{n}\cdot y=\mathrm{po}\},\qquad\left\{\hat{n}\cdot y=\frac{\mathrm{po}}{\sqrt{2}}\right\},\qquad\left\{\hat{n}\cdot y=\frac{\mathrm{po}}{\sqrt{2}}\right\},
\]
respectively.
As $\mathrm{po}$ ranges over the real numbers, each of the nine families is therefore the family of all planes with the normal $\hat{n}$.

\paragraph{7.1 (c2).}

\begin{proof}
By (C1), $\mu$ is invariant under $\mathbb{R}^3$, and for each of the nine vectors $\hat{n}$ it is invariant under $r_{\hat{n},0}$ and reflection positive at $\{\hat{n}\cdot y=0\}$.
For $\hat{n}=e_1,e_2,e_3$ these properties say that $\mu$ belongs to the class $\mathcal{C}$ of \cite[Section~8.1.3]{QEDI}.
Hence the hypotheses of Proposition~45.7 are satisfied by $\mathrm{lw}=\mu$ and by each of the nine vectors $\hat{n}$.
By Proposition~45.7\,(a), $\nu$-almost every $\mu_\xi$ is invariant under $r_{\hat{n},\mathrm{po}}$ for every real $\mathrm{po}$.
By Proposition~45.7\,(b), $\nu$-almost every $\mu_\xi$ is reflection positive at every plane $\{\hat{n}\cdot y=\mathrm{po}\}$, $\mathrm{po}\in\mathbb{R}$.
The union of the two exceptional sets is a set of $\mathrm{Blw}$ of $\nu$-measure $0$.
One obtains nine almost sure statements, one for each $\hat{n}$.
For $\hat{n}=e_j$ the statement follows also from Proposition~45.6\,(b) together with the first statement of (c1), because
\[
r_{e_j,\mathrm{po}}=\tau_{2\mathrm{po}\cdot e_j}\,\theta_j.
\]
\end{proof}

\begin{named}{Lemma 7.2}
The set $\Xi_\mu$ belongs to $\mathrm{Blw}$, and $\nu(\Xi_\mu)=1$.
\end{named}

\begin{proof}
Here $\mathrm{Lws}$ is the set of the laws and $\mathrm{Blw}$ the $\sigma$-algebra of (n7) on it, as in 0.1.
For every measurable function $\Phi\ge0$ on $\mathcal{S}'$ the map $\mathrm{lw}\mapsto\mathrm{lw}(\Phi)$ from $\mathrm{Lws}$ to $[0,\infty]$ is $\mathrm{Blw}$-measurable, and for every bounded measurable $\Phi$ the map $\mathrm{lw}\mapsto\mathrm{lw}(\Phi)$ is $\mathrm{Blw}$-measurable.
Both facts follow by monotone classes, as at the beginning of the proof of \cite[Lemma~8.17]{QEDI}.
Put
\begin{align*}
\mathrm{Lws}_1&:=\text{the set of the $\mathbb{R}^3$-ergodic laws},\\
\mathrm{Lws}_2&:=\text{the set of the laws with (G)},\\
\mathrm{Lws}_3&:=\left\{\mathrm{lw}:\ \mathrm{lw}\{dF=0\text{ in }\mathcal{S}'\}=1\right\}.
\end{align*}
For each of the nine vectors $\hat{n}$ let $\mathrm{Lws}_{\hat{n}}$ be the set of the laws $\mathrm{lw}$ with $(r_{\hat{n},0})_*\mathrm{lw}=\mathrm{lw}$ that are reflection positive at $\{\hat{n}\cdot y=0\}$ (nine sets).
For $m\in\{2,4\}$ and $\mathrm{ix}\in\{0,+,-\}$ put
\[
\mathrm{Lws}_{m,\mathrm{ix}}:=\left\{\mathrm{lw}:\ \mathrm{lw}\left(\left|\mathrm{Wk}_{m,\mathrm{ix}}\right|\right)<\infty\text{ and }\left|\mathrm{lw}\left(\mathrm{Wk}_{m,\mathrm{ix}}\right)\right|\le\mathsf{u}^{\mathrm{inf}}_{m,\mathrm{ix}}\right\}
\]
(six sets).
\begin{itemize}
\item[(1)]
Let $(R_j)$ and $(RP_j)$ be the conditions of \cite[Section~8.1.3]{QEDI}: invariance under $\theta_j$, and reflection positivity at $\{y_j=0\}$.
By \cite[Lemma~8.17]{QEDI}, the set of the laws with (G), the set of the laws with $(R_j)$, the set of the laws with $(RP_j)$, and the set of the laws that are invariant and ergodic under $\mathbb{R}^3$ belong to $\mathrm{Blw}$.
Hence $\mathrm{Lws}_1$, $\mathrm{Lws}_2$ and $\mathrm{Lws}_{e_j}$, $j=1,2,3$, belong to $\mathrm{Blw}$.
The event $\{dF=0\text{ in }\mathcal{S}'\}$ lies in $\Sigma$ \cite[Section~8.1.6]{QEDI}; therefore $\mathrm{Lws}_3$ belongs to $\mathrm{Blw}$.
The maps $\mathrm{lw}\mapsto\mathrm{lw}\left(\max(\mathrm{Wk}_{m,\mathrm{ix}},0)\right)$ and $\mathrm{lw}\mapsto\mathrm{lw}\left(\max(-\mathrm{Wk}_{m,\mathrm{ix}},0)\right)$ are $\mathrm{Blw}$-measurable with values in $[0,\infty]$.
The set $\mathrm{Lws}_{m,\mathrm{ix}}$ consists of the laws at which both are finite and the modulus of their difference is at most $\mathsf{u}^{\mathrm{inf}}_{m,\mathrm{ix}}$; so it belongs to $\mathrm{Blw}$.
\item[(2)]
Let $\hat{n}$ be one of the six vectors other than $e_1,e_2,e_3$, and let $\mathrm{gr}\in\mathrm{O}(3)$ satisfy $\mathrm{gr}\,e_1=\hat{n}$.
Then
\[
\mathrm{gr}\,\theta_1\,\mathrm{gr}^{-1}=\mathrm{gr}\left(1-2e_1e_1^{T}\right)\mathrm{gr}^{T}=1-2\hat{n}\,\hat{n}^{T}=r_{\hat{n},0}.
\]
By \cite[Lemma~14.12(b)]{QEDII}, a law $\mathrm{lw}$ lies in $\mathrm{Lws}_{\hat{n}}$ if and only if $\left(\mathrm{gr}^{-1}\right)_*\mathrm{lw}$ lies in $\mathrm{Lws}_{e_1}$.
The map $\mathrm{lw}\mapsto\left(\mathrm{gr}^{-1}\right)_*\mathrm{lw}$ of $\mathrm{Lws}$ into itself is $\mathrm{Blw}$-measurable, because
\[
\left(\left(\mathrm{gr}^{-1}\right)_*\mathrm{lw}\right)(\mathrm{Ev})=\mathrm{lw}(\mathrm{gr}\,\mathrm{Ev})
\]
and $\mathrm{gr}\,\mathrm{Ev}$ lies in $\Sigma$ for every $\mathrm{Ev}\in\Sigma$ \cite[Section~8.1.2]{QEDI}.
So $\mathrm{Lws}_{\hat{n}}$, the inverse image of $\mathrm{Lws}_{e_1}$ under this map, belongs to $\mathrm{Blw}$.
\item[(3)]
Let $\mathrm{lw}$ be invariant under $\mathbb{R}^3$, let $\hat{n}$ be a unit vector, and let $\mathrm{po}$ be real.
One has $r_{\hat{n},\mathrm{po}}=\tau_{\mathrm{po}\cdot\hat{n}}\,r_{\hat{n},0}\,\tau_{-\mathrm{po}\cdot\hat{n}}$.
Hence $\mathrm{lw}$ is invariant under $r_{\hat{n},\mathrm{po}}$ if it is invariant under $r_{\hat{n},0}$.
Now let $\mathrm{lw}$ be reflection positive at $\{\hat{n}\cdot y=0\}$, let $(f_p)$ be a finite family of real forms in $C_c^{\infty}(\{\hat{n}\cdot y>\mathrm{po}\};\Lambda^2)$, and let $(c_p)$ be a family of complex numbers with the same indices.
The forms $f'_p:=\tau_{-\mathrm{po}\cdot\hat{n}}\cdot f_p$ lie in $C_c^{\infty}(\{\hat{n}\cdot y>0\};\Lambda^2)$, and
\[
X(f_q)-X(r_{\hat{n},\mathrm{po}}\cdot f_p)=U(\mathrm{po}\cdot\hat{n})\left[X(f'_q)-X(r_{\hat{n},0}\cdot f'_p)\right].
\]
The law $\mathrm{lw}$ is invariant under $\tau_{\mathrm{po}\cdot\hat{n}}$.
Consequently the sum of (n3) formed at the plane $\{\hat{n}\cdot y=\mathrm{po}\}$ with $(f_p)$, $(c_p)$ equals the sum formed at the plane $\{\hat{n}\cdot y=0\}$ with $(f'_p)$, $(c_p)$:
\begin{align*}
&\sum_{p,q}\overline{c_p}\cdot c_q\cdot\mathrm{lw}\left(\exp\left(i\left(X(f_q)-X(r_{\hat{n},\mathrm{po}}\cdot f_p)\right)\right)\right)\\
&\qquad=\sum_{p,q}\overline{c_p}\cdot c_q\cdot\mathrm{lw}\left(\exp\left(i\left(X(f'_q)-X(r_{\hat{n},0}\cdot f'_p)\right)\right)\right)\ge0.
\end{align*}
It follows that, for a law that is invariant under $\mathbb{R}^3$, the property of being invariant under the reflections in all planes with the normal $\hat{n}$ and reflection positive at all of them holds if and only if the law lies in $\mathrm{Lws}_{\hat{n}}$.
\item[(4)]
Under a law with (G) the function $\mathrm{Wk}_{m,\mathrm{ix}}$, a real polynomial in $X(g_{\mathrm{ix}})$, is integrable \cite[Proposition~7.10(3)]{QEDI}.
The three statements (c1) for $\mu_\xi$ say that $\mu_\xi$ lies in $\mathrm{Lws}_1$, $\mathrm{Lws}_2$ and $\mathrm{Lws}_3$.
If $\mu_\xi\in\mathrm{Lws}_1$, then by (3) the statement (c2) for the family of the planes with the normal $\hat{n}$ holds for $\mu_\xi$ if and only if $\mu_\xi\in\mathrm{Lws}_{\hat{n}}$.
If $\mu_\xi\in\mathrm{Lws}_2$, then the statement (c3) for $(m,\mathrm{ix})$ holds for $\mu_\xi$ if and only if $\mu_\xi\in\mathrm{Lws}_{m,\mathrm{ix}}$.
So $\Xi_\mu$ is the intersection of the eighteen sets $\mathrm{Lws}_1$, $\mathrm{Lws}_2$, $\mathrm{Lws}_3$, $\mathrm{Lws}_{\hat{n}}$, $\mathrm{Lws}_{m,\mathrm{ix}}$, and $\Xi_\mu$ belongs to $\mathrm{Blw}$.
\item[(5)]
Each of the eighteen statements (c1), (c2), (c3) holds outside a set of $\mathrm{Blw}$ of $\nu$-measure $0$ (5.1, 7.1, 5.2).
The complement of $\Xi_\mu$ is contained in the union of these eighteen sets.
Hence $\nu(\Xi_\mu)=1$.
\end{itemize}
\end{proof}

\paragraph{7.3 The properties (D1)--(D6), for every $\xi\in\Xi_\mu$.}

\begin{proof}
Let $\xi\in\Xi_\mu$, and put $\mathrm{lw}:=\mu_\xi$.
Every step below is a statement about the single law $\mathrm{lw}$; no exceptional set is removed besides the complement of $\Xi_\mu$.
\begin{itemize}
\item[(D1)]
By (c1), $\mathrm{lw}\{dF=0\text{ in }\mathcal{S}'\}=1$ and $\mathrm{lw}$ has (G); also by (c1), $\mathrm{lw}$ is invariant and ergodic under $\mathbb{R}^3$.
By (c2), $\mathrm{lw}$ is invariant under the reflection in every plane of the nine families and reflection positive at every such plane.
Together with the invariance under $\mathbb{R}^3$, this says that $\mathrm{lw}$ lies in $\mathcal{C}_9$.
For $\hat{n}=\frac{e_i-e_l}{\sqrt{2}}$ one has
\[
r_{\hat{n},0}\,y=y-(y_i-y_l)\cdot(e_i-e_l),
\]
the map that exchanges the coordinates $y_i$ and $y_l$.
The group $O_{\mathrm{cub}}$ is the group of the signed permutation matrices; it is generated by $\theta_1$, $\theta_2$, $\theta_3$ and the three exchanges of two coordinates.
Hence $\mathrm{lw}$ is invariant under $O_{\mathrm{cub}}$.
Moreover, $\mathrm{lw}$ is reflection positive at $\{y_j=0\}$, $j=1,2,3$.
Altogether,
\begin{equation}\tag{7.1}
\text{$\mathrm{lw}$ satisfies (H) and (dF), lies in $\mathcal{C}_9$ and is $\mathbb{R}^3$-ergodic.}
\end{equation}
\item[(D2)]
By (7.1), \cite[Theorem~1.4(f)]{QEDII} applies to the law $\mathrm{lw}$, and so the assertions (a), (b), (d) of \cite[Theorem~1.4]{QEDII} hold for $\mathrm{lw}$.
By 6.4, these are (C2), (C3), (C7) with $\mathrm{lw}$ in the place of $\mu$.
The property (C4) for $\mathrm{lw}$ is Lemma~6.5.
\item[(D3)]
The law $\mathrm{lw}$ is $\mathbb{R}^3$-ergodic; by the last sentence of Proposition~45.6\,(a$''$), the measure of its $\mathbb{R}^3$-ergodic decomposition is the point mass at $\mathrm{lw}$.
By (7.1), $\mathrm{lw}$ satisfies (H) and lies in $\mathcal{C}_9$, so that \cite[Proposition~12.33]{QEDII} applies to the law $\mathrm{lw}$.
A set of $\mathrm{Blw}$ that has the measure $0$ under the point mass at $\mathrm{lw}$ does not contain $\mathrm{lw}$.
Therefore the clauses of \cite[Proposition~12.33]{QEDII} hold for the component $\mathrm{lw}$ itself.
By \cite[Proposition~12.33(d)]{QEDII},
\[
\ker\mathsf{H}_j(\mathrm{lw})=\mathbb{C}\,\Omega_j,\qquad j=1,2,3.
\]
\item[(D4)]
Applied to the law $\mathrm{lw}\in\mathcal{C}_9$, with the constants $4$ and $c_{\mathrm{cl}}$ of 6.6, \cite[Theorem~12.39]{QEDII} gives, for every $\mathrm{vp}>0$, all $\mathrm{Ga},\mathrm{Gb}\in L^2(\Sigma_{(-\mathrm{vp},\mathrm{vp})^3},\mathrm{lw})$ and all $a\in\mathbb{R}^3$ with $|a|_\infty\ge4\,\mathrm{vp}$,
\[
\left|\langle\mathrm{Ga},U(a)\mathrm{Gb}\rangle_{\mathrm{lw}}-\langle\mathrm{Ga},P^{123}_{\mathrm{lw}}\mathrm{Gb}\rangle_{\mathrm{lw}}\right|\le c_{\mathrm{cl}}\cdot\|\mathrm{Ga}\|_2\cdot\|\mathrm{Gb}\|_2\cdot\frac{\mathrm{vp}}{|a|_\infty}.
\]
Since $\mathrm{lw}$ is $\mathbb{R}^3$-ergodic, the assertion of \cite[Theorem~12.39]{QEDII} on ergodic laws gives $\langle\mathrm{Ga},P^{123}_{\mathrm{lw}}\mathrm{Gb}\rangle_{\mathrm{lw}}=\overline{\mathrm{lw}(\mathrm{Ga})}\cdot\mathrm{lw}(\mathrm{Gb})$; moreover $\langle\mathrm{Ga},U(a)\mathrm{Gb}\rangle_{\mathrm{lw}}=\mathrm{lw}\left(\overline{\mathrm{Ga}}\cdot U(a)\mathrm{Gb}\right)$.
With these two identities the last display becomes the first display of (D4).
The second assertion of (D4) is \cite[Proposition~12.33(e)]{QEDII} for the law $\mathrm{lw}$, read as in (D3).
\item[(D5)]
By (7.1) and (D2), $\mathrm{lw}$ has (H) and (dF), is invariant under every Euclidean map ((C2) for $\mathrm{lw}$), and satisfies (C3) at $\mathrm{po}=0$.
Lemma~6.7 gives (E0$'$), (E1), (E2), (E3).
For (E4), let $\Psi\in\mathrm{Sw}_m$, $\Phi\in\mathrm{Sw}_n$ and $a\in\mathbb{R}^3$.
Proposition~45.5\,(b) is applied three times: first its first clause; then its last clause with $\mathrm{Gi}=\tau_a$, which is admissible because $(\tau_a)_*\mathrm{lw}=\mathrm{lw}$, and where $\hat{\tau}_a=U(a)$; finally its second clause.
This yields
\begin{align*}
S_{m+n}\left(\Psi\otimes\tau_a^{\otimes n}\Phi\right)
&=\mathrm{lw}\left(X^{(m)}(\Psi)\cdot X^{(n)}\left(\tau_a^{\otimes n}\Phi\right)\right)\\
&=\mathrm{lw}\left(X^{(m)}(\Psi)\cdot U(a)X^{(n)}(\Phi)\right)\\
&=\left\langle X^{(m)}\left(\overline{\Psi}\right),U(a)X^{(n)}(\Phi)\right\rangle_{\mathrm{lw}}.
\end{align*}
The functions $X^{(m)}(\overline{\Psi})$ and $X^{(n)}(\Phi)$ lie in $L^2(\mathrm{lw})$.
By the second assertion of (D4), as $|a|\to\infty$ in any manner, the right side converges to
\begin{align*}
\overline{\mathrm{lw}\left(X^{(m)}\left(\overline{\Psi}\right)\right)}\cdot\mathrm{lw}\left(X^{(n)}(\Phi)\right)
&=\mathrm{lw}\left(X^{(m)}(\Psi)\right)\cdot\mathrm{lw}\left(X^{(n)}(\Phi)\right)\\
&=S_m(\Psi)\cdot S_n(\Phi).
\end{align*}
\item[(D6)]
The nine statements (c1), (c3) are among the eighteen; so, by 5.3, the four inequalities at the end of (B5) hold for $\xi$.
Let $\mathrm{Gi}$ be a Euclidean map.
By (C2) for $\mathrm{lw}$,
\[
S_n^{\xi}(\mathrm{Gi}\cdot f,\dots,\mathrm{Gi}\cdot f)=S_n^{\xi}(f,\dots,f)\qquad(n=1,\dots,4).
\]
For a test $f'$, the number $\kappa_4^{\xi}(f')$ is a polynomial in the numbers $S_n^{\xi}(f',\dots,f')$, $n\le4$ (0.1, Cumulants).
It follows that
\[
\kappa_4^{\xi}(\mathrm{Gi}\cdot f)=\kappa_4^{\xi}(f)\ge\frac{e^4\cdot b(f)}{2}>0.
\]
The law $\mathrm{lw}$ is not Gaussian: by (G), $\mathrm{lw}\left(X(f)\right)=0$ \cite[Lemma~9.1(a)]{QEDI}, and the fourth cumulant of a normal variable is $0$.
Nor is $\mathrm{lw}$ the point mass at $\omega=0$: under that law $X(f)=0$ almost surely, and its fourth cumulant is $0$.
\end{itemize}
\end{proof}

\subsubsection*{8. The numeral in (E0$'$); the first Remark after Theorem~44}

\begin{proof}
One has $66^2=4356=2\cdot2178$.
For every integer $n\ge0$ the number $\frac{4356^n}{n!}$ is one term of the series
\[
\sum_{k\ge0}\frac{4356^k}{k!}=\exp(4356),
\]
all of whose terms are positive; hence $\frac{4356^n}{n!}\le\exp(4356)$.
Taking square roots,
\[
\frac{66^n}{(n!)^{\frac{1}{2}}}=\left(\frac{4356^n}{n!}\right)^{\frac{1}{2}}\le\exp(2178),
\]
that is,
\[
2\cdot66^n\cdot(n!)^{\frac{1}{2}}\le2\exp(2178)\cdot n!.
\]
This is the inequality in brackets in (E0$'$).
It is also the assertion of the first Remark after Theorem~44 that the sequence $2\cdot66^n\cdot(n!)^{\frac{1}{2}}$ is at most $\alpha\cdot(n!)^{\beta}$ with $\alpha=2\exp(2178)$ and $\beta=1$, where $\alpha$, $\beta$ and $s$ denote the three numbers of that Remark.
The norm in (E0$'$) has the index $6n=ns$ with $s=6$.
\end{proof}

\subsubsection*{9. Uniformity}

\begin{itemize}
\item[(1)]
\emph{In $T$.}
The right sides of (U-a), (U-a$'$), (U-b), (U-c) and of the display of (A) do not contain $T$; they hold for all $T\ge T^{*}(e)+1$ (Proposition~43.2; 1.2).
\item[(2)]
\emph{In $N$.}
These right sides depend on $N$ through $\eta=L^{-N}$ and through the numbers $c_m$, $\mathsf{u}_m$, $\boldsymbol{\sigma}^2$ of the cutoff $N$.
For all $N\ge N_0(e)$, (U-c) bounds $\mathsf{u}_4$, $\left(1+|c_2|\right)\cdot\mathsf{u}_2$, $|c_2|$ and $\left|c_4-3c_2^2-b(g_{\mathrm{ix}})\right|$ by expressions in $C_U$, $e$, $\sigma$, $\mathfrak{e}_2$, $\eta$.
\item[(3)]
\emph{In the charge.}
The numbers $s_B$, $C_{3\mathrm{p}}$, $C_{\Gamma}^{(16)}$, $C_U$, $C_{\bar{\kappa}}$, $C'_{\bar{\kappa}}$, $e_2$ are constants of the line (C10): they depend on the constants of the lines (C0)--(C9) and on $(\bar{m},\beta,x_0)$, and not on $e$, $N$, $T$ (Proposition~43.2).
The inequalities (0.1)--(0.3) hold for all $e\in(0,e_2]$.
\item[(4)]
\emph{In the limit point, the subsequence, the placement and the component.}
The constants in (B2), (B3), (4.1), in the four inequalities of (B5) and in (D6) are $C_U$, $C_{\bar{\kappa}}$, $C'_{\bar{\kappa}}$, $b(f)$, $b(g_{\mathrm{ix}})$.
They are the same for all limit points $\mu$, for all subsequences and placements chosen in 4.1, and for all $\xi$.
The numbers $c^{\mathrm{inf}}_{m,\mathrm{ix}}$, $\mathsf{u}^{\mathrm{inf}}_{m,\mathrm{ix}}$ depend on these choices.
\item[(5)]
The numbers $c_{\mathrm{cl}}$, $61$, $0.9005$, $\frac{2}{5}$, $\log2$, $\frac{5\mathrm{e}}{2}$, $66$, $2178$ are absolute.
\end{itemize}

\subsection*{Conditions on the constants}

The conditions on the constants that enter the proof of Theorem~44 are the following; the second column names the constant on which each bears and the places in Section~I.8 and in Proposition~14.2 where it stands, the third the places of the proof where it enters.

\medskip
\noindent
\begin{tabular}{p{0.27\textwidth}p{0.33\textwidth}p{0.27\textwidth}}
\hline
condition & on which constant; where it stands & used in\\
\hline
$0<\rho<1$
& $\rho$; Section~I.8.2, line (C0)
& 0.3: a datum of \cite[Definition~2.36 and Theorems~1.2, 1.3, 1.4]{QEDII}\\
$L$ odd, $L\ge3$
& $L$; Section~I.0.5, line (C3) ($L$ odd, $L\ge5$)
& 0.3: a datum of \cite[Definition~2.36]{QEDII}\\
$0<\bar{m}\le1$, hence $h\ge1$
& $\bar{m}$; line (C10), Section~I.0.5
& 0.2, 0.3; $R=8h\cdot L^{T}\ge1$ in 2.1\,(b)\\
$\sigma=s\cdot\left(\frac{1}{4}-9\varepsilon\right)-12\varepsilon-\frac{2a_L}{N_{*}}>0$
& $\varepsilon$, $s$; line (C7); Proposition~14.2\,(b), (S2)
& (0.1), (0.3), 3.2\\
$\mathfrak{e}_2=0$, which enters (0.1) only through $\mathfrak{e}_2\ge0$
& Section~I.8.5
& (0.1)\\
$e_1\le\frac{1}{3}$, which enters (0.1) only through $e\le1$
& $e_1$; line (C9); Proposition~14.2\,(f)
& (0.1)\\
(Z10.s1) $e_2\le e_G$
& $e_2$; line (C10); Proposition~14.2\,(d)
& 0.2: $e\cdot h^{\frac{1}{2}}\le e_1$, so that Proposition~43.1\,(R0) and Section~I.6.2 apply; Proposition~43.2\\
(Z10.s2) $e_2\le e_{\sharp}(\beta_{\mathrm{ix}})$, $\mathrm{ix}=0,+,-$
& $e_2$; line (C10); Proposition~14.2\,(d)
& Section~3, first paragraph (the probes fit); Proposition~43.2\\
(Z10.s4) $\frac{5}{2}\cdot C'_{\bar{\kappa}}\cdot e_2^{\sigma}\le\frac{b(f)}{4}$
& $e_2$, after $C'_{\bar{\kappa}}$; line (C10); Proposition~14.2\,(d); Annex~B, Section~G6.Z4
& (0.3): 3.2, hence (B3), (B5), (D6)\\
(Pb), first part: $b(f)>0$
& the probe $\beta$; line (C10); Section~I.8.4; Proposition~14.2\,(d); the Data of Proposition~43.2
& 3.2\\
(Pb), second part: $d_0=|x_0|_2-2r_\beta>0$ and the inequality of \cite[Section~2.8]{QEDI}, displayed in the Data of Proposition~43.2
& $x_0$ and the radius $r_\beta$, after $\beta$; line (C10); Section~I.8.4; Proposition~14.2\,(d); the Data of Proposition~43.2; Annex~B, Section~G6.Z4
& 3.2\\
$C_{\bar{\kappa}}=2\cdot\left(8C_U+3C_U^2\right)$, $C'_{\bar{\kappa}}=8C_{\bar{\kappa}}+3C_{\bar{\kappa}}^2$ (definitions)
& line (C10); Proposition~43.2
& (0.2)\\
\hline
\end{tabular}
\medskip

\begin{thebibliography}{99}

\bibitem{Ahlfors1979}
L.~V.~Ahlfors, \emph{Complex analysis. An introduction to the theory of analytic functions
of one complex variable}, third edition, McGraw-Hill, New York, 1979.

\bibitem{Balaban}
T.~Ba{\l}aban, \emph{(Higgs)$_{2,3}$ quantum fields in a finite volume~I, II, III}, Commun.\ Math.\ Phys.\ 85 (1982) 603--626, 86 (1982) 555--594, 88 (1983) 411--445; \emph{Regularity and decay of lattice Green's functions}, Commun.\ Math.\ Phys.\ 89 (1983) 571--597; \emph{Propagators and renormalization transformations for lattice gauge theories~I, II}, Commun.\ Math.\ Phys.\ 95 (1984) 17--40, 96 (1984) 223--250; \emph{Averaging operations for lattice gauge theories}, Commun.\ Math.\ Phys.\ 98 (1985) 17--51; \emph{Spaces of regular gauge field configurations on a lattice and gauge fixing conditions}, Commun.\ Math.\ Phys.\ 99 (1985) 75--102; \emph{Propagators for lattice gauge theories in a background field}, Commun.\ Math.\ Phys.\ 99 (1985) 389--434; \emph{Ultraviolet stability of three-dimensional lattice pure gauge field theories}, Commun.\ Math.\ Phys.\ 102 (1985) 255--275; \emph{The variational problem and background fields in renormalization group method for lattice gauge theories}, Commun.\ Math.\ Phys.\ 102 (1985) 277--309; \emph{Renormalization group approach to lattice gauge field theories~I, II}, Commun.\ Math.\ Phys.\ 109 (1987) 249--301, 116 (1988) 1--22.

\bibitem{BalabanOCarrollSchor}
T.~Ba{\l}aban, M.~O'Carroll and R.~Schor, \emph{Block renormalization group for Euclidean fermions}, Commun.\ Math.\ Phys.\ 122 (1989) 233--247; \emph{Properties of block renormalization group operators for Euclidean fermions in an external gauge field}, J.\ Math.\ Phys.\ 32 (1991) 3199--3208.

\bibitem{Bhatia1997}
R.~Bhatia, \emph{Matrix analysis},
Springer-Verlag, New York, 1997.

\bibitem{Billingsley1999}
P.~Billingsley, \emph{Convergence of probability measures}, second edition,
John Wiley \& Sons, New York, 1999.

\bibitem{Bogachev2007}
V.~I.~Bogachev, \emph{Measure theory. Vol.~I},
Springer-Verlag, Berlin, 2007.

\bibitem{DimockMB}
J.~Dimock, \emph{Multiscale block averaging for QED in $d=3$}, arXiv:1712.10029v3.

\bibitem{DimockUR}
J.~Dimock, \emph{Ultraviolet regularity for QED in $d=3$}, J.\ Math.\ Phys.\ 59 (2018) 012301, arXiv:1512.04373v3.

\bibitem{DimockUVS}
J.~Dimock, \emph{Ultraviolet stability for QED in $d=3$}, arXiv:2009.01156v2.

\bibitem{FILS78}
J.~Fr\"ohlich, R.~Israel, E.~H.~Lieb and B.~Simon, \emph{Phase transitions and reflection positivity~I}, Commun.\ Math.\ Phys.\ 62 (1978) 1--34.

\bibitem{Froehlich76}
J.~Fr\"ohlich, \emph{The pure phases, the irreducible quantum fields, and dynamical symmetry breaking in Symanzik--Nelson positive quantum field theories}, Ann.\ Phys.\ 97 (1976) 1--54.

\bibitem{FSS76}
J.~Fr\"ohlich, B.~Simon and T.~Spencer, \emph{Infrared bounds, phase transitions and continuous symmetry breaking}, Commun.\ Math.\ Phys.\ 50 (1976) 79--95.

\bibitem{GlimmJaffe}
J.~Glimm and A.~Jaffe, \emph{Quantum physics. A functional integral point of view},
second edition, Springer-Verlag, New York, 1987.

\bibitem{GunningRossi1965}
R.~C.~Gunning and H.~Rossi, \emph{Analytic functions of several complex variables},
Prentice-Hall, Englewood Cliffs, NJ, 1965.

\bibitem{Hormander1990}
L.~H\"ormander, \emph{An introduction to complex analysis in several variables}, third edition,
North-Holland, Amsterdam, 1990.

\bibitem{HornJohnson2013}
R.~A.~Horn and C.~R.~Johnson, \emph{Matrix analysis}, second edition,
Cambridge University Press, Cambridge, 2013.

\bibitem{Humphreys1990}
J.~E.~Humphreys, \emph{Reflection groups and Coxeter groups},
Cambridge University Press, Cambridge, 1990.

\bibitem{Kallenberg2002}
O.~Kallenberg, \emph{Foundations of modern probability}, second edition,
Springer-Verlag, New York, 2002.

\bibitem{Kato1995}
T.~Kato, \emph{Perturbation theory for linear operators}, reprint of the 1980 edition,
Springer-Verlag, Berlin, 1995.

\bibitem{Lang2002}
S.~Lang, \emph{Algebra}, revised third edition,
Springer-Verlag, New York, 2002.

\bibitem{LiebLoss2001}
E.~H.~Lieb and M.~Loss, \emph{Analysis}, second edition,
American Mathematical Society, Providence, RI, 2001.

\bibitem{OsterwalderSchrader1975}
K.~Osterwalder and R.~Schrader, \emph{Axioms for Euclidean Green's functions~II},
Commun.\ Math.\ Phys.\ 42 (1975) 281--305.

\bibitem{QEDI}
\emph{Noncompact lattice QED in three dimensions: limit points of the curvature field
and their fourth cumulant} (companion paper).

\bibitem{QEDII}
\emph{Noncompact lattice QED in three dimensions: chamber fermions,
reflection positivity at all mirrors and rotation invariance} (companion paper).

\bibitem{ReedSimon1980}
M.~Reed and B.~Simon, \emph{Methods of modern mathematical physics.~I. Functional analysis},
second edition, Academic Press, New York, 1980.

\bibitem{Rudin1976}
W.~Rudin, \emph{Principles of mathematical analysis}, third edition,
McGraw-Hill, New York, 1976.

\bibitem{Rudin1987}
W.~Rudin, \emph{Real and complex analysis}, third edition,
McGraw-Hill, New York, 1987.

\bibitem{Simon2005}
B.~Simon, \emph{Trace ideals and their applications}, second edition,
American Mathematical Society, Providence, RI, 2005.

\bibitem{Stanley2012}
R.~P.~Stanley, \emph{Enumerative combinatorics. Vol.~1}, second edition,
Cambridge University Press, Cambridge, 2012.

\bibitem{SteinWeiss1971}
E.~M.~Stein and G.~Weiss, \emph{Introduction to Fourier analysis on Euclidean spaces},
Princeton University Press, Princeton, NJ, 1971.

\bibitem{vanLintWilson2001}
J.~H.~van Lint and R.~M.~Wilson, \emph{A course in combinatorics}, second edition,
Cambridge University Press, Cambridge, 2001.

\end{thebibliography}

\end{document}
